\documentclass[10pt]{amsart}
\usepackage{amsmath,amssymb,amsthm,mathtools,mathrsfs}
\usepackage[a4paper,hmargin=0.76in,vmargin=0.60in]{geometry}
\usepackage{booktabs}
\usepackage{array}
\usepackage{longtable}
\usepackage{tabularx}
\usepackage{xr-hyper}
\usepackage[hypertexnames=false]{hyperref}
\externaldocument[A-]{column-packages-character-families}
\hypersetup{hidelinks}
\usepackage[final,expansion=false]{microtype}
\usepackage{url}
\emergencystretch=2.5em

\newtheorem{theorem}{Theorem}[section]
\newtheorem{proposition}[theorem]{Proposition}
\newtheorem{lemma}[theorem]{Lemma}
\newtheorem{corollary}[theorem]{Corollary}
\theoremstyle{definition}
\newtheorem{definition}[theorem]{Definition}
\newtheorem{problem}[theorem]{Problem}
\theoremstyle{remark}
\newtheorem{remark}[theorem]{Remark}

\newcommand{\C}{\mathbb C}
\newcommand{\R}{\mathbb R}
\newcommand{\N}{\mathbb N}
\newcommand{\Z}{\mathbb Z}
\newcommand{\Q}{\mathbb Q}
\newcommand{\Rel}{\operatorname{Re}}
\newcommand{\Iml}{\operatorname{Im}}
\newcommand{\supp}{\operatorname{supp}}
\newcommand{\TV}{\operatorname{TV}}
\newcommand{\Kp}{\mathcal K^{\mathrm{prim}}_{2T}}
\newcommand{\Kw}{\mathcal K_{2T}}
\newcommand{\PT}{\mathcal P_T}
\newcommand{\IT}{I_T}
\newcommand{\lamP}{\lambda^{\mathrm{prim}}_T}
\newcommand{\lamW}{\lambda^{\mathrm W}_T}
\newcommand{\lamA}{\lambda_{\mathrm A}(T)}
\newcommand{\sT}{\sigma_T}
\newcommand{\gh}[1]{\widehat g_{#1}}
\newcommand{\nuh}{\widehat\nu}
\newcommand{\wh}{\widehat w}
\newcommand{\Fh}{\widehat F}
\newcommand{\cF}{c_F}
\newcommand{\Lev}{\mathfrak E}
\newcommand{\Wcl}{\mathcal W_T}
\newcommand{\NT}{\mathcal N_T}
\newcommand{\Aconst}{\mathsf A_T}
\newcommand{\Ainf}{\mathcal A_\infty}
\newcommand{\PiT}{\Pi_T}
\newcommand{\mT}{\mathfrak m_T}
\newcommand{\laminf}{\lambda^{\infty}_{\min}}
\newcommand{\laminfP}{\lambda^{\infty,\mathrm P}_{\min}}
\newcolumntype{Y}{>{\raggedright\arraybackslash}X}

\newcommand{\anThmA}{\ref{thm:srcA}}\newcommand{\anThmB}{\ref{thm:srcB}}\newcommand{\anThmC}{\ref{thm:srcC}}

\begin{document}
\addtolength{\abovedisplayskip}{-3.4pt}
\addtolength{\belowdisplayskip}{-3.4pt}
\addtolength{\belowdisplayshortskip}{-2pt}
\title[The sign of the radial intersection form]{The sign of the radial intersection form on the idele class group}
\author{David Alejandro Trejo Pizzo}
\date{}
\begin{abstract}
The radial trace on the idele class group of $\mathbb Q$ defines a Hermitian
form on functions supported in $(-T,T)$. We determine which sign of this form
corresponds to the intersection inequality on the kernel of the two Tate
functionals. The sign on every window is equivalent to the statement that
$\zeta$ has no zero in $\Re s>1/2$. We reduce the sign on all windows to
positive semidefiniteness in every degree of a single sequence of explicit
matrices, and on bounded windows we obtain quantitative estimates: for
$T\le\tfrac12\log5$ the lower bounds on unit vectors are $10^{-18}$ for the
complete form and $2.5\cdot10^{-12}$ for its primitive restriction, and a
certificate using finite-height zero data gives an explicit positive lower
bound up to $T=2.8022$. We then prove that no Dirichlet $L$-function
vanishes in $\Re s>1/2$; in particular the Riemann hypothesis holds, and the
form is nonnegative on every window. The proof has three parts: a one-sided
criterion of Landau type for a prime sum weighted by $1-\chi_K$, for
imaginary quadratic fields $K$; the construction of arithmetic data; and
three theorems on integral algebraic modular forms on a chain of theta lifts
between definite unitary groups over a CM field, which give the bound for
the prime sum required by the criterion. Complete proofs are given.
\end{abstract}
\maketitle
\setcounter{tocdepth}{1}\tableofcontents
% intro-new.tex -- Introduction of the paper, rewritten for the new second part.
% Paragraphs 1--3 are those of paper46-original.tex, unchanged.  Paragraphs 4--8 are new.
% References to Theorems A, B, C go through \anThmA, \anThmB, \anThmC (see sec-analytic.tex).
\providecommand{\anThmA}{\ref{thm:A}}
\providecommand{\anThmB}{\ref{thm:B}}
\providecommand{\anThmC}{\ref{thm:C}}
\section{Introduction}\label{sec:intro}
In Weil's argument for a curve over a finite field the two degree
directions pair hyperbolically. The index inequality concerns the
restriction of self-intersection to their common kernel. For the
radial correspondences on $\operatorname{Spec}\mathbb Z$, the
carrier is the colimit of the semilocal idele class groups over the
finite sets of places of $\mathbb Q$ containing infinity. The
involution, trace and Tate functionals are constructed in \cite{P43}.
The degree-$n$ diagonal contact contributes $\Lambda(n)/\sqrt n$
by \cite{P44}; the sum of contacts and the archimedean term has
the zero expansion with polar correction of \cite{Paper45}.

These three constructions are Rows~A, B and C, respectively.
This fourth paper treats Row~D, the index inequality for the resulting
radial intersection pairing on $\operatorname{Spec}\mathbb Z$.
The question is whether every radial class of bidegree zero has nonpositive
self-intersection, for every support radius $T>0$. From Rows~A--C
we use the pairing, its two degree maps and the trace identity
\eqref{comparison}. The remaining assertion is the sign of that
pairing on the common kernel of the degree maps.

The series follows the rows of Weil's proof for a curve over a finite
field, and this paper is its fourth. In \cite{P43} the correspondences
of a curve are modelled by graph bimodules of idele classes over a crossed
product by the positive $S$-units, the graded parts of whose Hochschild
complexes are computed up to homotopy; the sum over $S$ of Weil's local terms is a trace $T$ on
$C_c^\infty(C_S)$, the form $T(F\star G^\dagger)$ is nondegenerate, and the
two Tate characters $F\mapsto\int F|w|_S^{\mp1/2}$ take the place of the two
rulings. In \cite{P44} the graphs of the powers of Frobenius are replaced by
probability measures $\mathrm{Frob}_n$ on $C_S$ with
$\mathrm{Frob}_m\star\mathrm{Frob}_n=\mathrm{Frob}_{mn}$; on test functions
approximating $\mathrm{Frob}_n$ the finite local terms of $T$ sum to
$\Lambda(n)n^{-1/2}h(1/n)$, in the place of the count of points by
intersection with the diagonal. In \cite{Paper45} the limit of the trace
over $S$ is evaluated on radial functions: it is Weil's explicit formula
$T_{\lim}(\Gamma_h)=\Phi_h(0)+\Phi_h(1)-\sum_\rho\Phi_h(\rho)$, the Tate
functionals being the residues at $s=0$ and $s=1$, and its polarization is
\eqref{comparison}. What remains of Weil's argument is the analogue of the
index inequality, Row~D, and that is the subject of this paper.

The arithmetic input of this paper is a zero condition for the trivial
character of conductor one, whose $L$-function is $\zeta$: no zero to
the right of the critical line \eqref{eq:dirichlet-input}. Given it, the
functional equation puts every nontrivial zero on that line, and the zero
expansion \eqref{comparison}, which is Weil's explicit formula
(Section~\ref{sec:trace}), gives a sum of absolute squares for every compactly
supported smooth test. The quantitative Fourier estimates on
fixed supports and the matrix criterion for arbitrary supports give
two further descriptions of the same sign.

The zero condition is reached as the case of the trivial character of a
statement for all Dirichlet $L$-functions: $L(s,\chi)\ne0$ for
$\Re s>\frac12$ (Theorem~\ref{thm:GRH-half}). The chain has four links.
(i)~A criterion of Landau's type (Theorem~\ref{an:thm-criterion}). For a
modulus $q$ and an odd prime $d\nmid q$, with $\chi_K=(\cdot/d)$ (the
character of $K=\mathbb Q(\sqrt{-d})$ when $d\equiv3\pmod4$), let
$R_q(x)=\sum_{n\le x,\ n\equiv1\,(q)}\Lambda(n)(1-\chi_K(n))(1/n-1/x)-\log x/\varphi(q)-c_q$.
Let $a$ be prime to $2q$ and $\frac12<\beta<1$. If, for every prime
$d\equiv a\pmod q$, $R_q(x)\le x^{\beta-1}$ for all large $x$, then no
$L(s,\psi)$, $\psi$ modulo $q$, vanishes in $\Re s>\beta$. (In the proof
of Theorem~\ref{thm:GRH-half} the criterion is applied to the modulus
$q_0=18480q$ and the class $a=q_0/2-1$, so that $d\equiv7\pmod8$; a
character modulo $q$ induces one modulo $q_0$.) The
proof uses Landau's lemma on transforms of one-signed functions and the
existence of primes $d$ for which no twist $L(s,\psi\chi_K)$ vanishes in a
given box $\Re s\ge\beta$, $|\Im s|\le H$; the latter is derived here from
Heath-Brown's large sieve for real characters
(Proposition~\ref{an:prop-density}).
(ii)~Admissible data (Definition~\ref{an:def-data},
Theorem~\ref{an:thm-data}): for every large integer $x$, scales
$x\le x_1\le b$ with $\log b\approx x^{\varepsilon}$, a number
$m\approx\sqrt x$, a CM field $E=K(\zeta_{q_1})$ with
$x^{11}\le f=[E:K]\le x^{12}$,
a sparse set of its infinite places, weights, and three primes.
(iii)~For such data, three theorems on a space of automorphic forms on a
unitary group with integral structures at all primes: Theorem~\anThmA\
(integral forms with short images span), Theorem~\anThmB\ (few translates of
these images span almost everything) and Theorem~\anThmC\ (a lower bound for
covolume plus index).
(iv)~Their combination, the main inequality
$m\cdot\mathrm{rem}\le f\,a\,(\log b)^5$ with $a=xb$
(Theorem~\ref{an:thm-main}), in which
the total remainder $\mathrm{rem}$ is an explicit combination of $R_{q_0}(x)$,
$R_{q_0}(x_1)$ and $R_1(x_1)$ (Lemma~\ref{an:lem-rem}). Division by $m$ and the
prime number theorem for $R_{q_0}$ and $R_1$ at the point $x_1$
(Proposition~\ref{an:prop-pnt}) give
$R_{q_0}(y)\le y^{-1/2+\varepsilon}$ (Theorem~\ref{an:thm-onesided}), which is
the hypothesis of the criterion for every $\beta>\frac12$.

Links (i), (ii) and (iv) are proved in Sections~\ref{an:sec-criterion},
\ref{an:sec-data} and~\ref{an:sec-main}; the deduction (iv) of the main
inequality and of Theorem~\ref{thm:GRH-half} uses Theorems~\anThmA,
\anThmB\ and~\anThmC. These three theorems carry the weight of the argument.
They are stated in \S\ref{ss:au-ABC}; Theorems~\anThmA\ and~\anThmB\ are
deduced in \S\ref{ss:au-proofAB} and Theorem~\anThmC\ in \S\ref{ss:au-proofC}
from a list of propositions on the theta lifts, their divisibility at every
prime, their archimedean size, the second moment of a Poincar\'e map, the
determinant of the chain of lifts and the existence of regular families of
twists. These propositions are proved in
Sections~\ref{sec:proofsA}--\ref{prm:sec}, which form the larger part of
the paper.

This paper is Row~D in full. Everything that the proof of
Theorem~\ref{thm:GRH-half} and of the Riemann hypothesis
(Theorem~\ref{thm:RH}) requires is contained in it, with complete proofs:
the criterion, the data, the three theorems and all the propositions from
which they are deduced. Nothing is taken from the earlier papers of the
series for this purpose; they enter only in the interpretation of the
result as an index inequality. What we use from the literature are
published theorems, cited at the place where they are needed, and the table
in \S\ref{ss:au-remains} collects them for the automorphic part. In the
analytic part they are classical theorems on Dirichlet $L$-functions and
Heath-Brown's large sieve for real characters. In the automorphic part the
main ones are the Weil representation and Howe duality, the Siegel--Weil
formula, the theory of complex multiplication with the Chowla--Selberg
formula, the discrete series of real groups, Green's parametrisation of the
characters of finite general linear groups, standard monomial theory, the
Galois representations attached to automorphic forms on unitary groups,
Bloch--Kato Selmer groups and Poitou--Tate duality, Fontaine's theory of
$p$-adic representations, and Hecke's non-vanishing theorem on $\Re s=1$.
Once the Riemann hypothesis is known, the identity \eqref{comparison}, which
is Weil's explicit formula \cite{Weil52,Bombieri2000}
(Section~\ref{sec:trace}), gives the sign $W_T\ge0$ on the entire form domain
for every $T>0$ (Corollary~\ref{cor:all-windows}). The results are not
effective, because of the theorem of Siegel and Walfisz.

The paper is organised as follows. Sections~\ref{sec:trace}--\ref{sec:ladder}
study the radial trace and its form: the archimedean multiplier, duality and
certificates, the explicit formula on flat-top weights, the two Tate
directions, lower bounds on finite windows, the reduction to piecewise
exponentials and a common family of test matrices for all windows.
Section~\ref{an:sec-criterion} introduces the prime sum $R_q$ and proves the
one-sided criterion, including Landau's lemma, the prime number theorem for
$R_q$ and the choice of the prime $d$. Section~\ref{an:sec-data} defines the
admissible data and proves that they exist. Section~\ref{sec:automorphic}
defines the target forms, covolume and index, states
Theorems~\anThmA, \anThmB\ and~\anThmC, and deduces them from the
propositions listed in \S\ref{ss:au-remains}, which are proved in
Sections~\ref{sec:proofsA}--\ref{prm:sec}. Section~\ref{an:sec-main} proves the
identity for the total remainder, the main inequality, the one-sided bound
and Theorem~\ref{thm:GRH-half}, and discusses the role of $m$.
Section~\ref{sec:conclusion} deduces the Riemann hypothesis and the sign of
the form on every window.

\section{The radial trace}\label{sec:trace}
Write $I_T=(-T,T)$ and extend functions by zero. Our two transforms are
\[
 \widehat F(\tau)=\int F(t)e^{i\tau t}\,dt,\qquad
 \Phi_F(s)=\int F(t)e^{(s-1/2)t}\,dt .
\]
The \emph{Tate moments} are $A(F):=\Phi_F(1)=\int F(t)e^{t/2}\,dt$ and $B(F):=\Phi_F(0)=\int F(t)e^{-t/2}\,dt$.
A class is primitive when both vanish. Put
$c_F(u)=\int F(t)\overline{F(t-u)}\,dt$.
For real $F,G\in C_c^\infty(\IT)$ put $I_\infty(\Gamma_F,\Gamma_G):=-\mathfrak a_T(F,G)$, where
$\mathfrak a_T(F,G)=\frac1{2\pi}\int\sT\,\Fh\,\overline{\widehat G}\,d\tau$ is the polarization of the
truncated Weil form \eqref{eq:form} below; here $\Gamma_F$ is only a label for the test $F$.
Weil's explicit formula \cite{Weil52} (in the form of \cite[\S2, Explicit Formula, p.~186]{Bombieri2000}), applied to $h=|\Fh|^2$,
whose inverse Fourier transform is $c_F$ (smooth, supported in $[-2T,2T]$), gives the identity below:
$h(\pm i/2)=A(F)B(F)$, so the polar terms give $2A(F)B(F)$; the archimedean terms give
$\frac1{2\pi}\int m\,|\Fh|^2$ and the prime terms give $-\sum_{n\in\NT}v_n c_F(\log n)$, with $m$, $v_n$,
$\NT$ as in \eqref{eq:symbol} (after the change of variables $x=e^t$ and the standard rewriting of the
archimedean term through $\psi$). The archimedean and prime terms together are $\mathfrak a_T(F)$, so
$\sum_\rho\Phi_F(\rho)\Phi_F(1-\rho)=2A(F)B(F)+\mathfrak a_T(F)$. Hence
\begin{equation}\label{comparison}
 I_\infty(\Gamma_F,\Gamma_F)
 =2A(F)B(F)-\sum_\rho\Phi_F(\rho)\Phi_F(1-\rho).
\end{equation}
Repeated integration by parts in the strip and $N(H)=O(H\log H)$ give
absolute convergence with multiplicities. The primitive inequality is $I_\infty\le0$.

On real tests the moment map $M=(A,B)$ has pairing
$H((a,b),(a',b'))=ab'+ba'$. This is the symmetric bilinear form whose
quadratic form $2A(F)B(F)$ is the polar term of \eqref{comparison}; for
complex tests the hermitian form is $2\Rel(a\bar b')$, as in \eqref{eq:W}.
An even nonnegative bump of positive mass translated by $r$ has moment
vector proportional to $(e^{r/2},e^{-r/2})$, with proportionality
constant independent of $r$. Two distinct small translates have
independent moments. Their common kernel therefore has codimension
two in each window. Polarizing \eqref{comparison} gives
\[
 I_\infty(\Gamma_F,\Gamma_G)=H(M(F),M(G))-
 \tfrac12\sum_\rho\{\Phi_F(\rho)\Phi_G(1-\rho)+\Phi_G(\rho)\Phi_F(1-\rho)\}.
\]
For $F\in L^2(\IT)$ its autocorrelation is continuous, supported
in $[-2T,2T]$, and has Fourier transform $|\Fh|^2$. Put
\begin{equation}\label{eq:symbol}
m(\tau)=\Rel\psi\!\left(\tfrac14+\tfrac{i\tau}{2}\right)-\log\pi,\qquad
\sT(\tau)=m(\tau)-\sum_{n\in\NT}v_n\cos(\tau\log n),\qquad v_n=\frac{2\Lambda(n)}{\sqrt n},
\end{equation}
where $\psi$ is the digamma function, $\NT=\{n\ge2:\ \log n<2T\}$ is the set of \emph{contacts} and
$\Aconst:=\sum_{n\in\NT}v_n$ is the total mass of the prime comb. The truncated Weil form is
\begin{equation}\label{eq:form}
\mathfrak a_T(F)=\frac1{2\pi}\int_\R\sT(\tau)\,|\Fh(\tau)|^2\,d\tau
\end{equation}
on the form domain $\mathcal D_T=\{F\in L^2(\IT):\int\log(2+|\tau|)|\Fh|^2<\infty\}$. With the Tate
moments $A,B$ defined above, the \emph{primitive subspace} is
$\PT=\{F\in L^2(\IT):A(F)=B(F)=0\}$, and the \emph{complete Weil form} on the window is
\begin{equation}\label{eq:W}
W_T(F)=\mathfrak a_T(F)+2\,\Rel\big(A(F)\overline{B(F)}\big).
\end{equation}
In particular the arithmetic side gives $I_\infty(\Gamma_F,\Gamma_F)=-\mathfrak a_T(F)$ for real smooth tests. The explicit formula identifies $W_T(F)$ with the zero sum in
\eqref{comparison} for real smooth tests \cite{Weil52,Bombieri2000};
Lemma~\ref{lem:roughEF} treats the extension. We write
\[
\lamA=\inf\{\mathfrak a_T(F):F\in\PT\cap\mathcal D_T,\ \|F\|=1\},\qquad
\lamW=\inf\{W_T(F):F\in\mathcal D_T,\ \|F\|=1\};
\]
$\lamW\le\lamA$ because the two forms agree on $\PT$. Window positivity asks whether
\begin{equation}\label{eq:target}
\lamA\ge0\qquad(T>0),
\end{equation}
and complete positivity asks the same of $\lamW$. By monotonicity
(Lemma~\ref{lem:monoT}), \eqref{eq:target} is equivalent to
$\lim_{T\to\infty}\lamA\ge0$.
\subsection{Conventions}\label{subsec:conv}
$\widehat h(\tau)=\int_\R h(u)e^{i\tau u}du$; for even $h$ this is $\int h\cos(\tau u)$. For a finite even
measure $\nu$, $\nuh(\tau)=\int\cos(\tau u)\,d\nu(u)$, and for a probability measure $\mu$ on $\R$,
$\check\mu(u)=\int\cos(\tau u)\,d\mu(\tau)$. The letter $\tau$ always denotes a frequency, and $u$ a
variable dual to it. The right shift is $(V_aF)(t)=F(t-a)$, the reflection is $(JF)(t)=F(-t)$, and $Q_T$
denotes restriction to $\IT$ followed by extension by zero. We write $\mT(F)=\frac1{2\pi}\int m|\Fh|^2$
for the archimedean form, $\PiT(F)=-\sum_{n\in\NT}v_n\Rel\cF(\log n)$ for the prime form, so that
$\mathfrak a_T=\mT+\PiT$, and $\Ainf=\mT+2\Rel(A\bar B)$ for the complete archimedean form. The functions
\begin{equation}\label{eq:g}
g_0(u)=\cosh\tfrac u2\,\mathbf 1_{|u|<2T},\qquad g_1(u)=u\sinh\tfrac u2\,\mathbf 1_{|u|<2T},
\end{equation}
have transforms $\gh0,\gh1$ (Lemma~\ref{lem:gbounds}). For a parity label
$\mathrm{par}\in\{\mathrm{ev},\mathrm{odd}\}$, $\chi_{\mathrm{ev}}=\cosh\frac t2\,\mathbf 1_{\IT}$ and
$\chi_{\mathrm{odd}}=\sinh\frac t2\,\mathbf 1_{\IT}$. 

The variational approach follows \cite{Yoshida1992,Bombieri2000,Weil52,ConnesConsani2021}.

\section{The archimedean multiplier}\label{sec:levy}

\begin{lemma}[the L\'evy density]\label{lem:K}
For $u>0$ put $K(u)=\dfrac{2e^{3u/2}}{e^{2u}-1}=\dfrac{1}{2\sinh(u/2)}+\dfrac{1}{2\cosh(u/2)}
=2\sum_{k\ge0}e^{-(2k+\frac12)u}$. Then $K$ is positive, completely monotone,
$K(u)=u^{-1}+O(1)$ as $u\downarrow0$, $K(u)=2e^{-u/2}(1+O(e^{-u}))$ as $u\to\infty$, and
\begin{equation}\label{eq:m}
\begin{split}
m(\tau)&=\Rel\psi\!\left(\tfrac14+\tfrac{i\tau}2\right)-\log\pi
= C_0+\int_0^\infty(1-\cos\tau u)\,K(u)\,du,\\
C_0&=\psi(\tfrac14)-\log\pi=-5.3721834\ldots
\end{split}
\end{equation}
\end{lemma}
\begin{proof}
Writing the hyperbolic functions as exponentials gives
$K(u)=2e^{-u/2}/(1-e^{-2u})$. Its geometric series gives complete
monotonicity by termwise differentiation locally on $(0,\infty)$;
expansion at zero and infinity gives the stated bounds. In the integral
$\psi(z)=-\gamma+\int_0^\infty(e^{-t}-e^{-zt})/(1-e^{-t})\,dt$
subtract the value at $1/4$, take real parts and put $t=2u$.
The resulting integrand is $(1-\cos\tau u)K(u)$, integrable by
the bounds just obtained.
\end{proof}

\begin{theorem}[complete Bernstein structure]\label{thm:S}
Put $a_k=k+\frac14$ and $\varphi(x)=\sum_{k\ge0}\dfrac{x}{a_k(x+4a_k^2)}$ for $x\ge0$. Then
\emph{(i)} $\Rel\psi(\frac14+\frac{i\tau}2)-\psi(\frac14)=\varphi(\tau^2)$ and $\varphi$ is a
complete Bernstein function; \emph{(ii)} the L\'evy measure of the negative-definite function
$\tau\mapsto\varphi(\tau^2)$ on $\R$ is $\frac12K(|u|)du$, which is completely monotone on each
half-line; since $\varphi$ is a complete Bernstein function by (i), the associated L\'evy process is a pure-jump subordinate Brownian motion
$B_{S_t}$ with a special subordinator $S$ of Laplace exponent $\varphi$; in particular it has no
Gaussian component; \emph{(iii)} $\int_{2T}^\infty K(u)\,du=\log\coth\frac T2+\pi-2\arctan e^{T}$.
\end{theorem}
\begin{proof}
Subtract the digamma series at $1/4$ from its value at
$1/4+i\tau/2$ and take real parts. Each difference equals
$\tau^2/[a_k(4a_k^2+\tau^2)]$. Thus $\varphi$ has the Stieltjes
representation with masses $1/a_k$ at $4a_k^2$, locally convergent
by comparison with $\sum_{k\ge1}k^{-3}$. Lemma~\ref{lem:K} identifies
the jump measure, and $\varphi(x)/x\to0$ excludes a Gaussian part.
The subordination statement follows from \cite[Chs.~6,13]{SSV2012}.
The tail formula follows by evaluating the primitives
$\log\tanh(u/4)$ and $2\arctan e^{u/2}$ of the two hyperbolic kernels.
\end{proof}

\begin{theorem}[the form as a L\'evy energy]\label{thm:A}
For every $F\in L^2(\IT)$, with $\Lev(\Phi):=\int_0^\infty(\Phi(0)-\Rel\Phi(u))K(u)\,du\in[0,\infty]$,
\begin{equation}\label{eq:Aform}
\mathfrak a_T(F)=C_0\|F\|^2+\Lev(\cF)-\sum_{n\in\NT}v_n\Rel\cF(\log n),
\end{equation}
both sides being finite exactly when $\int\log(2+|\tau|)|\Fh|^2<\infty$. Consequently, with
$\Kw=\{\cF:F\in L^2(\IT),\|F\|=1\}$ and the primitive class
\begin{equation}\label{eq:Kprim}
\Kp=\Big\{\Phi\in\Kw:\ \int\Phi(u)\cosh\tfrac u2\,du=0,\ \int\Phi(u)\,u\sinh\tfrac u2\,du=0\Big\},
\end{equation}
one has
\begin{align*}
\lamP&:=\inf_{\Phi\in\Kp}\Big[C_0+\Lev(\Phi)-\sum_{n\in\NT}v_n\Rel\Phi(\log n)\Big]=C_0-\Lambda_T,\\
\Lambda_T&:=\sup_{\Phi\in\Kp}\Big[\sum_{n\in\NT}v_n\Rel\Phi(\log n)-\Lev(\Phi)\Big],
\end{align*}
and $\lamP\le\lamA$.
\end{theorem}
\begin{proof}
Plancherel gives $\int\cos(\tau u)|\Fh|^2d\tau/(2\pi)=\Rel\cF(u)$.
Apply Tonelli to $(1-\cos(\tau u))K(u)|\Fh|^2$ in \eqref{eq:m}.
The resulting integral is finite on the stated domain since
$\int_0^1(1-\cos(\tau u))\,du/u=\log|\tau|+O(1)$ at infinity and
$K(u)-1/u$ is bounded near zero. Expanding the hyperbolic functions
in \eqref{eq:Kprim} gives $\Rel(A\bar B)$ and
$\Rel(A'\bar B-A\bar B')$, where primes insert $t$ into the moment
integrals. Both vanish on $\PT$, giving $\lamP\le\lamA$. For primitive
tests $\widehat{\cF}(z)=\widehat F(z)\overline{\widehat F(\bar z)}$
has double zeros at $\pm i/2$.
\end{proof}

\begin{remark}\label{rem:Kprim}
$\Kp$ is slightly larger than the set of primitive autocorrelations: for real $F$ its two
conditions read $AB=0$ and $A'B-AB'=0$, i.e.\ $F\in\PT$, or $A=A'=0$, or $B=B'=0$. Every
certificate below bounds $\lamP$, hence $\lamA$; the equality $\lamP=\lamA$ is not needed
anywhere and is not asserted. Since $W_T(F)=\mathfrak a_T(F)$ whenever $A\bar B$ is purely
imaginary, also $\lamW\le\lamP$. Whether $\lamP=\lamW$, i.e.\ whether the optimal multipliers of the
primitive problem are the polar ones, is characterized in Theorem~\ref{thm:mult}.
\end{remark}

\begin{corollary}[monotonicity; upper bound]\label{cor:mono}
With the contact set held fixed, $T\mapsto\Lambda_T$ is nondecreasing and $T\mapsto\lamP$
nonincreasing. Moreover
$\lamA\le C_0+\Aconst+\varphi\big((3\pi/2T)^2\big)$, which equals $1.291$ at $T=\frac12\log3$ and
$2.327$ at $T=\log 2$.
\end{corollary}
\begin{proof}
Monotonicity: $\Kp$ grows with $T$ and $\Lev$ does not depend on $T$. Upper bound: the form
$\Lev(\cF)$ is the Dirichlet form of the L\'evy process of Theorem~\ref{thm:S} killed on
leaving $\IT$; its eigenvalues are bounded by those of the subordinate killed Brownian motion,
$\lambda_n\le\varphi((n\pi/2T)^2)$ \cite[Prop.~2.3]{SongVondracek2003}, the Brownian motion being
normalized with characteristic exponent $\xi^2$; $\PT$ has codimension
$2$, so the bottom of $\Lev$ on $\PT$ is at most $\lambda_3$; and $-\sum v_n\cF(\log n)\le\Aconst\|F\|^2$.
\end{proof}

\begin{lemma}[elementary bounds]\label{lem:mbounds}
\emph{(i)} $m$ is increasing on $[0,\infty)$. \emph{(ii)} For every $\tau>0$,
\[
\log\frac{\tau}{2\pi}-\frac1\tau-\frac1{2\tau^2}\ \le\ m(\tau)\ \le\ \log\frac{\tau}{2\pi}+\frac1\tau+\frac12\log\Big(1+\frac1{4\tau^2}\Big),
\]
and the lower bound, which we denote by $m_{\mathrm{lo}}(\tau)$, is increasing in $\tau$. \emph{(iii)}
\[
-\sum_{n\in\NT}v_n\cos(\tau\log n)\ge-\Aconst,\qquad\text{with}\quad
\sup_\tau\sum_{n\in\NT}v_n\cos(\tau\log n)=\Aconst.
\]
\end{lemma}
\begin{proof}
(i) In $\Rel\psi(\frac14+iy)=-\gamma+\sum_{n\ge0}\big[\frac1{n+1}-\frac{n+\frac14}{(n+\frac14)^2+y^2}\big]$
every summand increases with $|y|$. (ii) Use $\psi(z)=\log z-\frac1{2z}-\int_0^\infty g(t)e^{-zt}dt$
for $\Rel z>0$, $g(t)=\frac1{e^t-1}-\frac1t+\frac12$ \cite[5.9.13]{DLMF}, at $z=\frac14+iy$, $y=\tau/2$.
Here $g(0^+)=0$, $g'(t)=t^{-2}-\frac14\sinh^{-2}(t/2)\in(0,\infty)$ and $g(\infty)=\frac12$, so one
integration by parts gives $|\int_0^\infty g e^{-zt}dt|=|z|^{-1}|\int_0^\infty g'e^{-zt}dt|\le\frac1{2|z|}\le\frac1{2y}$.
Also $\Rel\log z=\frac12\log(\frac1{16}+y^2)\in[\log y,\log y+\frac12\log(1+\frac1{16y^2})]$ and
$\Rel\frac1{2z}=\frac{1/8}{1/16+y^2}\in[0,\frac1{8y^2}]$. Substituting $y=\tau/2$ and $-\log\pi$
gives (ii); the derivative of the lower bound is $\tau^{-1}+\tau^{-2}+\tau^{-3}>0$. (iii) Each cosine is at most one, and equality holds simultaneously at $\tau=0$.
\end{proof}

\begin{lemma}[density and continuity in form norm]\label{lem:form-density}
Let $T>0$ and put
$\|F\|_{\mathcal D}^2=\frac1{2\pi}\int_\R\log(2+|\tau|)\,|\Fh(\tau)|^2\,d\tau$
on $\mathcal D_T$.
\begin{enumerate}
\item[(i)] The complex combinations of real functions in $C_c^\infty(\IT)$
are dense in $\mathcal D_T$ for $\|\cdot\|_{\mathcal D}$.
\item[(ii)] There is $C_T$ with $|W_T(F)|\le C_T\|F\|_{\mathcal D}^2$ on
$\mathcal D_T$. Hence $W_T$ and its polarization are continuous for
$\|\cdot\|_{\mathcal D}$.
\end{enumerate}
\end{lemma}
\begin{proof}
(i) Real and imaginary parts of $F\in\mathcal D_T$ lie in $\mathcal D_T$,
since the weight is even; so let $F$ be real. First dilate $F(t)$ to
$F(t/\lambda)$ with $\lambda\uparrow1$; the support lies in
$[-\lambda T,\lambda T]$. Fourier dilation is strongly continuous in
$L^2(\log(2+|\tau|)d\tau)$ by density of compactly supported continuous
functions and doubling of the weight. Mollification by a real mollifier at
scale below $(1-\lambda)T$ preserves interior support and converges in the
same norm by dominated convergence.

(ii) The function $m$ is even and continuous, so Lemma~\ref{lem:mbounds}(ii)
gives a constant $C$ with $|m(\tau)|\le C+\log(2+|\tau|)$, and
Lemma~\ref{lem:mbounds}(iii) bounds the prime part of $\sT$ by $\Aconst$.
With Plancherel and $\log2\,\|F\|_2^2\le\|F\|_{\mathcal D}^2$ this gives
$|\mathfrak a_T(F)|\le(1+(C+\Aconst)/\log2)\|F\|_{\mathcal D}^2$. By
Cauchy--Schwarz, $|A(F)|,|B(F)|\le(2e^T)^{1/2}\|F\|_2$. This proves the bound
for $W_T$; continuity of the quadratic form and of its polarization follows
from the bound and the polarization identity.
\end{proof}

\begin{lemma}[monotonicity across contact thresholds]\label{lem:monoT}
The functions $T\mapsto\lamA$, $T\mapsto\lamP$ and $T\mapsto\lamW$ are nonincreasing on $(0,\infty)$,
including at the points $T=\frac12\log n$ where a contact enters.
\end{lemma}
\begin{proof}
Let $T<T'$ and extend $F\in L^2(\IT)$ by zero to $I_{T'}$. The norm, the Tate moments $A(F),B(F)$ and
the autocorrelation $\cF$ are unchanged, and $\cF$ vanishes on $|u|\ge2T$. A contact
$n\in\mathcal N_{T'}\setminus\NT$ has $\log n\ge2T$, so $\cF(\log n)=0$, and by \eqref{eq:Aform}
$\mathfrak a_{T'}(F)=\mathfrak a_T(F)$. Hence $\PT\subset\mathcal P_{T'}$ isometrically with the same
form values, which gives the statement for $\lamA$; the same computation with $W_T=\mathfrak a_T+2\Rel(A\bar B)$
gives it for $\lamW$, and for $\lamP$ the class $\Kp$ grows with $T$, $\Lev$ is $T$-free and the new
contacts contribute $\Phi(\log n)=0$.
\end{proof}

\begin{lemma}[compactness]\label{lem:compact}
For every $V\in\R$ the set $\{F\in\mathcal D_T:\|F\|=1,\ \mathfrak a_T(F)\le V\}$ is relatively compact in
$L^2(\IT)$. Consequently the form $\mathfrak a_T$ restricted to $\PT\cap\mathcal D_T$ is the form of a
self-adjoint operator $H_T$ on $\PT$ with compact resolvent, $\lamA=\min\operatorname{spec}H_T$ is attained,
and the bottoms of $\mathfrak a_T$ and $W_T$ on the unit sphere of $L^2(\IT)\cap\mathcal D_T$, and on the unit sphere of
$V\cap\mathcal D_T$ for any closed subspace $V$ with $V\cap\mathcal D_T\ne\{0\}$, are attained.
\end{lemma}
\begin{proof}
By Lemma~\ref{lem:mbounds} the shifted form norm is equivalent to the
weighted Fourier norm defining $\mathcal D_T$. This domain is complete
by Plancherel and closedness of the support condition. Its unit ball
has uniformly vanishing Fourier tails since $m(R)\to\infty$; on
$[-R,R]$ the Fourier restriction has square-integrable kernel
$e^{i\tau t}$ on $[-R,R]\times\IT$. Compactness on each band and
the tail bound give compactness of the form-domain embedding.
Closed restriction to $\PT$ and addition of the bounded polar form
preserve this argument; here $\PT\cap\mathcal D_T$ is dense in $\PT$: given $F\in\PT$ and $F_k\in C_c^\infty(\IT)$
with $F_k\to F$ in $L^2$, subtract $c_{1,k}\varphi_1+c_{2,k}\varphi_2$, where $\varphi_1,\varphi_2\in C_c^\infty(\IT)$ have
invertible moment matrix and $(c_{1,k},c_{2,k})$ is chosen to kill $A(F_k),B(F_k)$; the coefficients tend to $0$
because $A,B$ are continuous on $L^2(\IT)$. A minimizing sequence on the unit sphere of $V\cap\mathcal D_T$ lies in
a sublevel set, so it has an $L^2$-convergent subsequence; its limit lies in $V$, has norm one, and lies in
$\mathcal D_T$ with value at most the infimum by lower semicontinuity of the closed form. Compact resolvent and attainment follow.
\end{proof}

On $\PT$ the operator $H_T$ commutes with $J$, since $\sT$ is even and $A(JF)=B(F)$; so its eigenspaces
split into even and odd parts.

\begin{lemma}[continuity in the window]\label{lem:contT}\label{lem:contA}
Both $\lamW$ and $\lamA$ are continuous in $T$, as is the bottom
of the complete archimedean form. For the complete Weil form this
also follows from the parity minima in \cite[Thm.~5, p.~199]{Bombieri2000}.
\end{lemma}
\begin{proof}
Work on $I_{2T}$ with zero extension, so by Lemma~\ref{lem:monoT}
the form of a vector is independent of its supporting window.
For $T_j\downarrow T$, choose unit minimizers. Their energies are
bounded above by the minimum at $T$. The polar term is bounded on the
unit ball, so compactness gives an $L^2$-convergent subsequence.
Its limit is a unit vector supported in $[-T,T]$; the Tate conditions,
when imposed, pass to the limit. Lower semicontinuity on the common
window and monotonicity give right continuity.

For the other direction take a unit minimizer $F$ at $T$ and put
$F_r(t)=r^{-1/2}F(t/r)$, $r\uparrow1$. Then $F_r\to F$ in form norm:
on Fourier transforms the dilations $g(s)\mapsto r^{1/2}g(rs)$
are uniformly bounded on $L^2(\log(2+|s|)\,ds)$ for $1/2<r<1$,
and converge strongly on continuous compactly supported functions,
a dense subspace. In the primitive case correct the two moments
of $F_r$ by a linear combination of two fixed smooth bumps in
$I_{T/2}$ with independent moment vectors. Its coefficients tend
to zero by $L^2$ convergence. The corrected vectors, normalized,
are supported in $I_{rT}$ and their energies converge to the minimum
at $T$. Monotonicity concludes. With the bounded prime comb removed
the same compactness and approximation arguments apply.
\end{proof}

\section{Duality for the window form}\label{sec:duality}

Fix $T$ and the weights $v=(v_n)_{n\in\NT}$ (any nonnegative weights may be used; $v_n=2\Lambda(n)/\sqrt n$
is the case of $\zeta$).

\begin{definition}\label{def:PD}
\textbf{(P)} $\ \lamP=\inf\big\{\int_\R\sT\,d\mu:\ \mu$ an even probability measure on $\R$,
$\check\mu(u)=0$ for $|u|\ge2T$, $\int\gh0\,d\mu=\int\gh1\,d\mu=0\big\}$.

\textbf{(D)} $\ D_T=\sup\big\{\inf_{\tau\in\R}\big[\sT(\tau)+\nuh(\tau)+\beta_0\gh0(\tau)+\beta_1\gh1(\tau)\big]:\
\nu$ a finite even signed Borel measure with $\supp\nu\subseteq\{|u|\ge2T\}$, $\beta\in\R^2\big\}$.
A triple $(\nu,\beta_0,\beta_1)$ with $f:=\sT+\nuh+\beta\cdot\widehat g\ge c$ on $\R$ is a
\emph{certificate of value $c$}.
\end{definition}

\begin{proposition}[the primal is the eigenvalue problem; weak duality]\label{prop:weak}
The value of \textup{(P)} is the $\lamP$ of Theorem~\ref{thm:A}. Every certificate of value $c$
satisfies $c\le\lamP\le\lamA$.
\end{proposition}
\begin{proof}
For $\cF\in\Kp$ take the even part of $|\Fh|^2d\tau/(2\pi)$.
Its cosine transform is $\Rel\cF$; it preserves the two zero moments
and the integral of the even symbol $\sT$. Conversely a feasible
$\mu$ has continuous positive definite transform $\check\mu$ supported
on $[-2T,2T]$ with $\check\mu(0)=1$. Boas--Kac factorization
\cite{BoasKac1945} gives $\check\mu=\cF$ for a unit vector on $\IT$;
the two constraints put $\cF$ in $\Kp$. Finally
$\int\nuh\,d\mu=\int\check\mu\,d\nu=0$ by the support conditions,
so integrating a certificate gives $c\le\int\sT\,d\mu$.
\end{proof}

\begin{lemma}[bounds on the multipliers]\label{lem:gbounds}
The multipliers $\gh0(\tau)=2\int_0^{2T}\cosh\frac u2\cos\tau u\,du$ and $\gh1(\tau)=2\int_0^{2T}u\sinh\frac u2\cos\tau u\,du$
are even, real, bounded, tend to $0$ at infinity, and for $\tau>0$
\[
|\gh0(\tau)|\le k_0(\tau):=\frac{2\cosh T}{\tau}+\frac{\sinh T}{\tau^2},\qquad
|\gh1(\tau)|\le k_1(\tau):=\frac{8T\sinh T}{\tau},
\]
both bounds decreasing in $\tau$. Moreover $\gh0(\tau)=\frac{2\cosh T\sin2T\tau}{\tau}+O(\tau^{-2})$,
and for any $g=y_0g_0+y_1g_1$, $\widehat g(\tau)=\frac{2g(2T^-)\sin2T\tau}{\tau}+\frac{2g'(2T^-)\cos2T\tau}{\tau^2}+O(\tau^{-3})$.
\end{lemma}
\begin{proof}
Integrate by parts: $\gh0=\frac{2\cosh T\sin2T\tau}\tau-\frac1\tau\int_0^{2T}\sinh\frac u2\sin\tau u\,du$,
and once more,
\[\frac1\tau\int_0^{2T}\sinh\frac u2\sin\tau u\,du=-\frac{\sinh T\cos2T\tau}{\tau^2}+\frac{\gh0(\tau)}{4\tau^2}.\]
Solving gives $|\gh0|(1+\frac1{4\tau^2})\le\frac{2\cosh T}\tau+\frac{\sinh T}{\tau^2}$. For $\gh1$, one
integration by parts gives $\gh1=\frac{4T\sinh T\sin2T\tau}{\tau}-\frac2\tau\int_0^{2T}\big(\sinh\frac u2+\frac u2\cosh\frac u2\big)\sin\tau u\,du$,
and the integrand is the derivative of $u\sinh\frac u2$, whose integral over $[0,2T]$ is $2T\sinh T$.
The asymptotic expansion is two integrations by parts for a function real-analytic on $[0,2T]$
with $g'(0)=0$.
\end{proof}


\begin{theorem}[strong duality; finitely atomic dual sequences]\label{thm:D1}
$D_T=\lamP$. Moreover, for every $\varepsilon>0$ there are finitely many atoms
$u_1,\dots,u_N\ge2T$, reals $a_1,\dots,a_N$ and $\beta\in\R^2$ such that
$\sT(\tau)+\sum_j2a_j\cos(u_j\tau)+\beta_0\gh0(\tau)+\beta_1\gh1(\tau)>\lamP-\varepsilon$ for all $\tau\in\R$.
In particular, for every $T$ with $\lamP>0$, $\lamA>0$ is provable by a finite certificate.
\end{theorem}
\begin{proof}
\emph{(i) Finite relaxations.} For a finite set $U=\{u_1<\dots<u_N\}\subset[2T,\infty)$ let
$P(U)=\inf\{\int\sT d\mu:\ \mu\ \text{even probability},\ \check\mu(u_i)=0\ \forall i,\ \int\gh j d\mu=0\}$
and $D(U)$ the supremum in (D) restricted to $\nu$ supported on $\pm U$. Then $P(U)\le\lamP$
(the feasible set of (P) is contained in that of $P(U)$) and $D(U)\le D_T\le\lamP$.

\emph{(ii) $P(U)=D(U)$.} Write $f_0\equiv1$, $f_i(\tau)=\cos(u_i\tau)$ ($1\le i\le N$),
$f_{N+1}=\gh0$, $f_{N+2}=\gh1$, $b=(1,0,\dots,0)\in\R^{N+3}$, and let
$C=\{\sum_kt_kf(\tau_k):\ t_k\ge0\}$ be the convex cone generated by the moment curve
$\tau\mapsto f(\tau)=(f_0,\dots,f_{N+2})(\tau)$. The value function
$w(b')=\inf\{\int\sT d\mu:\ \mu\ge0,\ \int f\,d\mu=b'\}$ is convex, $>-\infty$ everywhere
($\sT\ge-\Aconst-|C_0|$ and $\int f_0d\mu=b'_0$ bounds the mass) and finite on $C$ (finite atomic
$\mu$ have finite value). Its Legendre biconjugate at $b$ is the Lagrangian dual
$w^{**}(b)=\sup\{y_0:\ \sT-\sum_ky_kf_k\ge0\text{ on }\R\}$, which is $D(U)$ after
$a_i=-y_i/2$, $\beta_j=-y_{N+1+j}$. A proper convex function agrees with its biconjugate at every
interior point of its domain \cite[Thm.~7.4, Cor.~12.1.1]{Rockafellar}; so it suffices to show
$b\in\operatorname{int}C$.

Otherwise separation from the interior supplies $y\ne0$ with
$y\cdot z\ge y\cdot b$ for $z\in C$ (if $C$ has empty interior, use
a normal to its span, with sign chosen at $b$). Taking $z=0$ and
$z=tf(\tau)$ gives $y_0\le0$ and $G=y\cdot f\ge0$. The Ces\`aro mean
of $G$ is $y_0$, since the distinct positive frequencies have mean zero
and $\gh j\to0$. Hence $y_0=0$, and the mean of $|G|$ is zero.
Testing against $\cos(u_i\tau)$ gives $y_i/2=0$. There remains
$\widehat g\ge0$ with $g=y_{N+1}g_0+y_{N+2}g_1$. By the expansion in Lemma~\ref{lem:gbounds},
$\widehat g\ge0$ forces $g(2T^-)=0$ and then $g'(2T^-)=0$; but
$g(2T^-)=y_{N+1}\cosh T+2T\,y_{N+2}\sinh T$ and $g'(2T^-)=\frac12y_{N+1}\sinh T+y_{N+2}(\sinh T+T\cosh T)$,
whose determinant is $\sinh T\cosh T+T>0$. So $y=0$, a contradiction; hence $b\in\operatorname{int}C$.

Choose increasing finite $U_n$ with dense union in $[2T,\infty)$
and $1/n$-minimizers $\mu_n$. Since $\sigma_T$ is bounded below and tends to
infinity, these probabilities are tight. In a weak limit every cosine
constraint on the union survives, hence by continuity every constraint
on $[2T,\infty)$; the bounded continuous $\gh j$ likewise preserve the
two moments. Lower semicontinuity gives $\lamP\le\lim P(U_n)\le\lamP$.
Thus $D_T\ge\sup_n D(U_n)=\lamP$, with the required finite certificates.
\end{proof}

\begin{corollary}\label{cor:D1W}
For every $\varepsilon>0$ there are finitely many $u_j\ge2T$ and $a_j\in\R$ with
$\sT+\sum_j2a_j\cos(u_j\tau)+2\gh0>\lamW-\varepsilon$ on $\R$. Thus
$\lamW=\sup_\nu\inf_\tau[\sT+2\gh0+\nuh]$ over finitely atomic even $\nu$ on $\{|u|\ge2T\}$.
\end{corollary}
\begin{proof}
Repeat the proof of Theorem~\ref{thm:D1} with the objective $\sT+2\gh0$ (bounded below, continuous,
tending to $+\infty$) and with the two $\widehat g$ constraints deleted. In (ii) the separation argument
stops at $y_0=0$ and $y_i=0$, so $b$ is interior; the limit step is unchanged. The primal value is then
$\inf\frac1{2\pi}\int(\sT+2\gh0)|\Fh|^2$ over unit $F$, which is $\lamW$ by Proposition~\ref{prop:P}.
\end{proof}
More generally, for every $\beta\in\R^2$ the same proof with the objective $\sT+\beta\cdot\widehat g$
and both moment constraints deleted gives $V_T(\beta)=\sup_\nu\inf_\tau[\sT+\beta\cdot\widehat g+\nuh]$ over
finitely atomic even $\nu$ on $\{|u|\ge2T\}$, where $V_T(\beta)$ is defined in Section~\ref{sec:complete}
(its two moment terms have the multipliers $\gh0,\gh1$); the added symbol is bounded and continuous.


\begin{theorem}[non-attainment]\label{thm:O}
Let $\nu$ be any finite even signed measure with $\supp\nu\subseteq\{|u|\ge2T\}$ and let
$h_0,\dots,h_J$ be finitely many continuous functions on $[-2T,2T]$ that annihilate $\Kp$
(for instance $g_0,g_1$). Then for every $\beta\in\R^{J+1}$
\[
\inf_{\tau\in\R}\Big[\sT(\tau)+\nuh(\tau)+\sum_j\beta_j\widehat h_j(\tau)\Big]<\lamP\quad\text{strictly.}
\]
The statement holds for every finite weight family, including $v=0$.
\end{theorem}
\begin{proof}
A minimizer $\Phi_\star=c_{F_\star}\in\Kp$ of the primal exists: the sublevel sets of
$\mathfrak a_T$ on the unit sphere of $L^2(\IT)$ are compact (Lemma~\ref{lem:compact}),
and, in the language of measures, the sublevel sets $\{\mu\text{ feasible}:\int\sT\,d\mu\le V\}$ are tight (as
$\sT\to+\infty$) and weak-$*$ closed and $\mu\mapsto\int\sT d\mu$ is
weak-$*$ lower semicontinuous, exactly as in the last paragraph (passage to the limit) of the proof of Theorem~\ref{thm:D1}. Put $\varrho=\sT+\nuh+\sum\beta_j\widehat h_j-\lamP$, a
continuous function ($\nu$ finite). By the annihilation properties,
$\frac1{2\pi}\int\varrho|\Fh_\star|^2=\mathfrak a_T(F_\star)-\lamP=0$. If $\varrho\ge0$ then $\varrho=0$ a.e.\ on
$\{\Fh_\star\ne0\}$, whose complement is the discrete zero set of a nonzero entire function of
exponential type; so $\varrho\equiv0$ on $\R$, and by Fourier inversion, as distributions in $u$,
\[
\check\sT+\nu+\sum_j\beta_jh_j=\lamP\,\delta_0 .
\]
Restrict to the open set $0<|u|<2T$, where $\nu=0$: by \eqref{eq:m}, $\check m=-\frac12K(|u|)$
there, and the prime part of $\check\sT$ is $-\sum_{n\in\NT}\frac{v_n}2(\delta_{\log n}+\delta_{-\log n})$;
hence $\sum_j\beta_jh_j(u)=\frac12K(|u|)+\sum_{n\in\NT}\frac{v_n}2(\delta_{\log n}+\delta_{-\log n})$
on $0<|u|<2T$. The left side is a bounded continuous function; the right side is unbounded as
$u\to0$ ($K(u)\sim1/u$) and carries atoms as soon as $\NT\neq\emptyset$. Contradiction.
\end{proof}

Every value below a positive $\lamP$ admits an atomic lower
certificate by Theorem~\ref{thm:D1}. The dual supremum is not
attained (Theorem~\ref{thm:O}); the primal infimum is attained
(Lemma~\ref{lem:compact}; see the first lines of the proof of Theorem~\ref{thm:O}).


\begin{theorem}[finite-interval Wiener--Hopf equation]\label{thm:W}
Let $\mathsf S_T$ denote the convolution operator on $\R$ with symbol $\sT$,
$\widehat{\mathsf S_TF}=\sT\Fh$, acting on $L^2(\IT)$ extended by zero and taking values in tempered
distributions. A unit vector $F\in\PT\cap\mathcal D_T$ is a critical point of $\mathfrak a_T$ on the unit
sphere of $\PT$ with critical value $\lambda$ if and only if there are $\alpha,\beta\in\C$ (real for real
$F$) with
\[
(\mathsf S_T-\lambda)F=\alpha\cosh\tfrac t2+\beta\sinh\tfrac t2\qquad\text{on }\IT,
\]
in the sense of distributions on $\IT$. The symbol $\sT-\lambda$ is real, even, tends to
$+\infty$, and changes sign for every $\lambda>C_0-\Aconst$; in that range it is not elliptic, and no
Wiener--Hopf index is defined for it.
\end{theorem}
\begin{proof}
For $G\in C_c^\infty(\IT)$, $\langle(\mathsf S_T-\lambda)F,G\rangle=\frac1{2\pi}\int(\sT-\lambda)\Fh\,\overline{\widehat G}$
is the polarization of $\mathfrak a_T-\lambda\|\cdot\|^2$. Criticality on the unit sphere of $\PT$
with value $\lambda$ means that this vanishes for all $G$ in $\PT\cap\mathcal D_T$; since
$C_c^\infty(\IT)\cap\PT$ is dense in $\PT\cap\mathcal D_T$ for the form norm (by Lemma~\ref{lem:form-density}(i):
dilate, then mollify, then correct the two Tate moments by a fixed pair in $C_c^\infty(\IT)$ with independent
moment vectors; the coefficients tend to $0$ since $A,B$ are continuous on $L^2(\IT)$), it is equivalent to require it for
$G\in C_c^\infty(\IT)$ with $\int Ge^{t/2}=\int Ge^{-t/2}=0$. This is a subspace of codimension $2$
of $C_c^\infty(\IT)$, and a distribution on $\IT$ that annihilates the common kernel of two continuous
linear functionals is a linear combination of them; the functionals are integration against
$e^{\pm t/2}$, i.e.\ against $\cosh\frac t2\pm\sinh\frac t2$. Conversely, if $(\mathsf S_T-\lambda)F$
is such a combination on $\IT$, it annihilates every $G$ in that subspace, so $F$ is critical. Finally
$\sT(0)=C_0-\Aconst$ and $\sT\to+\infty$ (Lemma~\ref{lem:mbounds}).
\end{proof}


\begin{lemma}[moment interior]\label{lem:W}
For $F\in L^2(\IT)$ real put $Q_0(F)=\int\cF\cosh\frac u2=A B$ and
$Q_1(F)=\int\cF\,u\sinh\frac u2=A'B-AB'$. The image of the set of unit vectors under
$F\mapsto(Q_0,Q_1)$ contains a neighborhood of $0$ in $\R^2$, realised by $C_c^\infty(\IT)$
functions; equivalently, $0$ is an interior point of the (convex) moment image
$\mathfrak m(\mu)=(\int\gh0d\mu,\int\gh1d\mu)$ of the feasible set of (P) without the two
constraints, and the interior is realised by measures of finite value.
\end{lemma}
\begin{proof}
The formulas for $Q_0,Q_1$ follow from expanding $\cosh\frac{t-s}2$ and $(t-s)\sinh\frac{t-s}2$
in exponentials. The functions $e^{t/2},e^{-t/2},te^{t/2},te^{-t/2}$ are linearly independent
on $\IT$, and a nontrivial combination is real-analytic, so it cannot annihilate
$C_c^\infty(\IT)$: the linear map $F\mapsto(A,B,A',B')$ is onto $\R^4$ from $C_c^\infty(\IT)$.
If $0$ were not interior, some $(t_0,t_1)\ne0$ would satisfy $t_0Q_0(F)+t_1Q_1(F)\ge0$ for all
$F$ (normalization only rescales $Q_j$ by a positive factor). Take $(A,B,A',B')=(1,0,0,\mp1)$:
$Q_0=0$, $Q_1=\pm1$, so $t_1=0$; take $(1,\pm1,0,0)$: $Q_1=0$, $Q_0=\pm1$, so $t_0=0$. A finite
convex combination of the four witnesses covers a ball, and each has finite value because
$\Fh$ is Schwartz. For the image of the unit sphere itself: a function with
$(A,B,A',B')=(1,x,0,-y)$ has $(Q_0,Q_1)=(x,y)$; adding an element of the kernel of the moment map
does not change $(Q_0,Q_1)$ and makes the norm any number above the minimal one, and
$(Q_0,Q_1)(F/\|F\|)=(Q_0,Q_1)(F)/\|F\|^2$. Since the minimal norm is bounded on $x^2+y^2=1$, every
point $t(x,y)$ with $0<t\le t_0$, $t_0>0$ uniform, is attained on the unit sphere, and $(0,0)$ is
attained on $\PT$.
\end{proof}

\begin{proposition}[the polar multiplier]\label{prop:P}
For every $F\in L^2(\IT)$ in the form domain,
$\frac1{2\pi}\int\gh0|\Fh|^2=\Rel\big(A(F)\overline{B(F)}\big)$ and hence
\[
\frac1{2\pi}\int\big(\sT+2\gh0\big)|\Fh|^2\,d\tau=W_T(F);
\]
if moreover $F$ is real and $\sum_\rho|\Phi_F(\rho)\Phi_F(1-\rho)|<\infty$ (in particular if
$F\in C_c^\infty(\IT)$ is real), then $W_T(F)=\sum_\rho\Phi_F(\rho)\Phi_F(1-\rho)$
(Lemma~\ref{lem:roughEF} below).
Consequently a certificate with $(\beta_0,\beta_1)=(2,0)$ and value $c$ proves $W_T\ge c$ on all
of $L^2(\IT)$, and $(2,0)$ is the only multiplier for which the integrated certificate equals $W_T$ identically.
\end{proposition}
\begin{proof}
$\frac1{2\pi}\int\gh0|\Fh|^2=\int g_0\cF=\iint\cosh\frac{t-s}2F(t)\overline{F(s)}$ and
$\cosh\frac{t-s}2=\frac12e^{t/2}e^{-s/2}+\frac12e^{-t/2}e^{s/2}$ gives $\frac12(A\bar B+B\bar A)$.
The identification of $W_T$ with the zero sum is the explicit formula for the autocorrelation
$\cF$, \eqref{comparison} \cite{Weil52,Bombieri2000}, whose polar term is
$2\Phi_F(0)\Phi_F(1)=2\Rel(A\bar B)$ for the real forms considered here. For $\beta\ne(2,0)$,
$\frac1{2\pi}\int(\sT+\beta_0\gh0+\beta_1\gh1)|\Fh|^2=W_T(F)+(\beta_0-2)\Rel(A\bar B)+\beta_1\Rel(A'\bar B-A\bar B')$
(Theorem~\ref{thm:A}), and the correction term takes both signs on the unit sphere of $L^2(\IT)$
(Lemma~\ref{lem:W}, witnesses $(1,\pm1,0,0)$ and $(1,0,0,\mp1)$), so it is not identically zero.
\end{proof}

The explicit formula \eqref{comparison} is stated for real $F\in C_c^\infty$. The first equality of
Proposition~\ref{prop:P} holds on the whole form domain and is the only one the certificates use;
wherever the zero sum is used for a rough test (Theorem~\ref{thm:margin}), it enters through the
following lemma.

\begin{lemma}[the explicit formula for rough tests]\label{lem:roughEF}
Let $F\in L^2(\IT)$ be real, with \emph{(a)} $\int\log(2+|\tau|)|\Fh|^2<\infty$ and \emph{(b)}
$\sum_\rho|\Phi_F(\rho)\Phi_F(1-\rho)|<\infty$. Then
\[
\sum_\rho\Phi_F(\rho)\Phi_F(1-\rho)=W_T(F)=2\Phi_F(0)\Phi_F(1)+\mathfrak a_T(F),
\]
the sum being absolutely convergent. In particular $\mathfrak a_T(F)=\sum_\rho\Phi_F(\rho)\Phi_F(1-\rho)$
for $F\in\PT$, and a zero $\rho=\frac12+i\gamma$ on the critical line contributes $|\Fh(\gamma)|^2$.
\end{lemma}
\begin{proof}
For real $F$, $\Phi_F(1)=A(F)$ and $\Phi_F(0)=B(F)$, so $2\Phi_F(0)\Phi_F(1)=2\Rel(A\bar B)$. Let
$\eta_\varepsilon\ge0$, $\eta_\varepsilon\in C_c^\infty(-\varepsilon,\varepsilon)$, $\int\eta_\varepsilon=1$,
and $F_\varepsilon=F*\eta_\varepsilon\in C_c^\infty$, supported in $[-T-\varepsilon,T+\varepsilon]$.
Apply \eqref{comparison} (Weil's explicit formula \cite{Weil52,Bombieri2000}) to $F_\varepsilon$ at the window $T+\varepsilon$. Zero side:
$\Phi_{F_\varepsilon}=\Phi_F\Phi_{\eta_\varepsilon}$, $|\Phi_{\eta_\varepsilon}(s)|\le e^{\varepsilon|\Rel s-\frac12|}\le e^{\varepsilon/2}$
in the critical strip and $\Phi_{\eta_\varepsilon}\to1$ pointwise, so by (b) and dominated convergence
the zero sum (which for $F_\varepsilon$ converges absolutely, so that the symmetric order is
immaterial) tends to $\sum_\rho\Phi_F(\rho)\Phi_F(1-\rho)$. Polar term:
$\Phi_{F_\varepsilon}(0)\Phi_{F_\varepsilon}(1)\to\Phi_F(0)\Phi_F(1)$. Archimedean term:
$|\Fh_\varepsilon|^2=|\Fh|^2|\widehat{\eta_\varepsilon}|^2\le|\Fh|^2$ and $|m(\tau)|\le|C_0|+\log(1+|\tau|)$
(Lemma~\ref{lem:mbounds}(i),(ii); $m\le0$ for $|\tau|\le\tau_0$), so by (a) and dominated convergence
$\frac1{2\pi}\int m|\Fh_\varepsilon|^2\to\frac1{2\pi}\int m|\Fh|^2$. Prime term: $c_{F_\varepsilon}\to\cF$
uniformly, the sum runs over the finitely many $n<e^{2T+2\varepsilon}$, and for
$e^{2T}\le n<e^{2T+2\varepsilon}$ one has $\cF(\log n)=0$; so it tends to $-\sum_{n\in\NT}v_n\cF(\log n)$.
\end{proof}


\section{Atomic certificates}\label{sec:tails}

A certificate is an inequality $f\ge c$ on all of $\R$. The lemmas of this section reduce it to a
bounded interval; they rest on Lemmas~\ref{lem:mbounds} and \ref{lem:gbounds}.

For a feasible probability $\mu$ of value at most $\lamP+\varepsilon$
put $G=\lamP+\varepsilon+\Aconst-C_0$. The pointwise inequalities
$0\le m-C_0\le\sT+\Aconst-C_0$ give
$\int(m-C_0)d\mu\le G$. By monotonicity the mass outside $[-R,R]$
is at most $G/(m(R)-C_0)$. This is at most $\eta\in(0,1)$ if
$R\ge\max(2,2\pi\exp(1+C_0+G/\eta))$ by Lemma~\ref{lem:mbounds}.
For a certificate $f\ge c$, weak duality gives
$\lamP-c\le\int(f-c)d\mu$. Outside the interval,
$f-c\le m-C_0+(E_R+C_0)^+$ with
$E_R=\Aconst+\|\nu\|_{\TV}+\sum_i|\beta_i|k_i(R)-c$.
The omitted integral therefore costs at most
$\int_{|\tau|>R}(m-C_0)d\mu+\eta(E_R+C_0)^+$.
If $f\ge c$ is known only on the interval, set
$V=\Aconst+\|\nu\|_{\TV}+\sum_i|\beta_i|\widehat g_i(0)-C_0$.
Since $g_i\ge0$, one has $f\ge-V$ everywhere. For $c\ge-V$,
integration yields $c(1-\eta)-V\eta\le\int f\,d\mu\le\lamP+\varepsilon$.
For $c<-V$ the same bound follows from $\int f\,d\mu\ge-V$.

\subsection{Lattice atoms: the graded periodic tail}

\begin{lemma}[graded periodic tail]\label{lem:F1}
Let $d=2T/M$, $u_j=jd$ for $M\le j\le J$ (so $u_M=2T$), $\nu=\sum_ja_j(\delta_{u_j}+\delta_{-u_j})$,
$P=2\pi/d$, so that $\nuh(\tau)=\sum_j2a_j\cos(u_j\tau)$ is $P$-periodic. Let
$f=\sT+\nuh+\beta_0\gh0+\beta_1\gh1$. Fix $\tau_1>0$, $N\ge1$, and for $x\in[\tau_1,\tau_1+P)$ put
\begin{align*}
Y_k(x)&=m(x+kP)-\sum_{n\in\NT}v_n\cos\big((x+kP)\log n\big)-|\beta_0|k_0(x+kP)-|\beta_1|k_1(x+kP)\ \ (k<N),\\
Y_N(x)&=m(x+NP)-\Aconst-|\beta_0|k_0(x+NP)-|\beta_1|k_1(x+NP).
\end{align*}
If $\nuh(x)+Y_k(x)\ge c$ for all $x\in[\tau_1,\tau_1+P)$ and all $0\le k\le N$, then $f(\tau)\ge c$
for all $\tau\ge\tau_1$.
\end{lemma}
\begin{proof}
Every $\tau\ge\tau_1$ is uniquely $x+kP$ with $x\in[\tau_1,\tau_1+P)$, $k\ge0$, and
$\nuh(\tau)=\nuh(x)$ since $u_jP=2\pi j$. For $k<N$ the $k$-th row is $f(x+kP)$ with the exact
prime term and the two multiplier terms replaced by the lower bounds of Lemma~\ref{lem:gbounds}.
For $k\ge N$: $m(x+kP)\ge m(x+NP)$, $k_i(x+kP)\le k_i(x+NP)$, and the prime term is $\ge-\Aconst$.
\end{proof}

The prime term is never required to be periodic: it is evaluated exactly on the first $N$
copies of one period of $\nuh$ and replaced by its infimum beyond. The lemma therefore holds
for any finite set of contacts; the incommensurability of $\log2$ and $\log3$ plays no role.
It is a sufficient criterion only (the rows are conservative in $\beta$ by
$|\beta_0|k_0+|\beta_1|k_1$).

\subsection{Exponential columns: the monotone tail}

\begin{proposition}[exponential columns]\label{prop:T0}
For $s>0$ and an integer $p\ge0$ let $d\nu_{s,p}(u)=\frac12(|u|-2T)^pe^{-s(|u|-2T)}\mathbf 1_{|u|\ge2T}du$,
a positive finite measure of mass $p!/s^{p+1}$ supported outside the window. Then
\[
\nuh_{s,p}(\tau)=\Rel\Big[e^{2iT\tau}\frac{p!}{(s-i\tau)^{p+1}}\Big],\qquad
|\nuh_{s,p}(\tau)|\le\operatorname{cap}_{s,p}(\tau):=\frac{p!}{(s^2+\tau^2)^{(p+1)/2}},
\]
and $\operatorname{cap}_{s,p}$ is decreasing in $\tau\ge0$.
\end{proposition}
\begin{proof}
$\nuh_{s,p}(\tau)=\int_{2T}^\infty(u-2T)^pe^{-s(u-2T)}\cos\tau u\,du=\Rel\big[e^{2iT\tau}\int_0^\infty t^pe^{-(s-i\tau)t}dt\big]$.
\end{proof}

\begin{lemma}[monotone tail]\label{lem:T}
Let $\nu=\sum_{k=1}^Kc_k\nu_{s_k,p_k}+\sum_{j=1}^Ja_j(\delta_{u_j}+\delta_{-u_j})$ with $u_j\ge2T$, and
$f=\sT+\nuh+\beta_0\gh0+\beta_1\gh1$. Define for $\tau>0$
\[
L(\tau)=\log\frac\tau{2\pi}-\frac1\tau-\frac1{2\tau^2}-\Aconst-\sum_k|c_k|\operatorname{cap}_{s_k,p_k}(\tau)
-2\sum_j|a_j|-|\beta_0|k_0(\tau)-|\beta_1|k_1(\tau).
\]
Then $f\ge L$ on $(0,\infty)$, $L$ is nondecreasing, and $L(\tau_1)\ge c$ for one $\tau_1$ implies
$f\ge c$ on $[\tau_1,\infty)$.
\end{lemma}
\begin{proof}
Lemma~\ref{lem:mbounds}(ii),(iii), Proposition~\ref{prop:T0}, the trivial bound on the atoms and
Lemma~\ref{lem:gbounds}; every term of $L$ is nondecreasing.
\end{proof}

\begin{corollary}[cost of the tail]\label{cor:T2}
For an atom-free $\nu$ with $A:=\sum_k|c_k|p_k!$ and $A':=|\beta_0|(2\cosh T+\sinh T)+|\beta_1|8T\sinh T$,
the tail inequality $L(\tau_1)\ge c$ holds for every $\tau_1\ge\max\{A+A'+2,\ 2\pi e^{c+\Aconst+1}\}$,
so the tail costs linearly in the $\ell^1$ size of the certificate. For atoms it costs
exponentially, $\tau_1\approx2\pi e^{\Aconst+2\sum|a_j|}$, which is why Lemma~\ref{lem:F1} and not
Lemma~\ref{lem:T} is used for lattice certificates. Neither bound is uniform over the dual class,
since both grow with the size of the certificate.
\end{corollary}
\begin{proof}
For $\tau\ge1$, $\operatorname{cap}_{s,p}(\tau)\le p!/\tau$ and $|\beta_i|k_i\le A'/\tau$, so
$L(\tau)\ge\log\frac\tau{2\pi}-\Aconst-\frac{A+A'+2}{\tau}$.
\end{proof}

\begin{lemma}[hybrid tail]\label{lem:hybrid}
Let $\nu=\nu_{\mathrm{lat}}+\sum_\ell c_\ell\nu_{s_\ell,p_\ell}$ with $\nu_{\mathrm{lat}}$ as in Lemma~\ref{lem:F1}.
For $x\in[\tau_1,\tau_1+P)$ define $H_k(x)=\nuh_{\mathrm{lat}}(x)+Y_k(x)+\sum_\ell c_\ell\nuh_{s_\ell,p_\ell}(x+kP)$
for $0\le k<N$, and $H_N(x)=\nuh_{\mathrm{lat}}(x)+Y_N(x)-\sum_\ell|c_\ell|\operatorname{cap}_{s_\ell,p_\ell}(x+NP)$.
If $H_k\ge c$ on $[\tau_1,\tau_1+P)$ for all $0\le k\le N$ then $f\ge c$ on $[\tau_1,\infty)$.
\end{lemma}
\begin{proof}
As in Lemma~\ref{lem:F1}; on the rows $k<N$ the column term is evaluated exactly, and for
$k\ge N$ it is bounded by $-\sum_\ell|c_\ell|\operatorname{cap}(x+kP)\ge-\sum_\ell|c_\ell|\operatorname{cap}(x+NP)$
since $\operatorname{cap}$ decreases.
\end{proof}



\begin{theorem}[exponential columns exhaust the dual]\label{thm:E}
Let $\mathcal E$ be the real span of $\{\nu_{s,p}:s>0,p\ge0\}$. Then
$\sup\{\inf_\tau[\sT+\nuh+\beta\cdot\widehat g]:\ \nu\in\mathcal E,\beta\in\R^2\}=\lamP$, and the same
holds with $\mathcal E$ replaced by the absolutely continuous measures with bounded density and
compact support in $\{|u|\ge2T\}$.
\end{theorem}
\begin{proof}
Weak duality gives the upper bound. By Theorem~\ref{thm:D1} choose an
atomic certificate within $\varepsilon$ of $\lamP$. Since $\sT\to\infty$,
the tail bound beyond some $R$ holds for every replacement of total variation at most
$\|\nu\|_{\rm TV}+1$ (the bound subtracts only $\Aconst+\|\cdot\|_{\rm TV}+\sum_i|\beta_i|k_i$ from
$m_{\rm lo}$, and $m_{\rm lo}\to\infty$); every replacement below has total variation at most
$\|\nu\|_{\rm TV}+1$: spreading an atom does not increase total variation, and the
density approximation is taken within total-variation distance $\min(1,\varepsilon)$. Spread each positive-location atom uniformly over an
interval of length $\eta$ to its right and reflect at zero. This preserves
the support condition and changes the transform on $[-R,R]$ by at most
$R\eta\|\nu\|_{\rm TV}$. The densities $t^pe^{-t}$ span a dense subspace
of $L^1(0,\infty)$: an annihilator in $L^\infty$ has Laplace transform
with every derivative zero at one, hence is zero by uniqueness.
Translate this approximation by $2T$ and reflect. Total variation error
bounds Fourier error uniformly, proving both assertions.
\end{proof}


\begin{proposition}[the canonical column; band-limitation]\label{prop:M}
Let $d\nu_K=\frac12K(|u|)\mathbf 1_{|u|\ge2T}du$, a positive measure of mass
$c_K=\int_{2T}^\infty K$ (Theorem~\ref{thm:S}(iii)). Then for all $\tau$
\[
\sT(\tau)+\nuh_K(\tau)=\big(C_0+c_K\big)+\int_0^{2T}(1-\cos\tau u)K(u)\,du-\sum_{n\in\NT}v_n\cos(\tau\log n),
\]
a function whose distributional Fourier transform is supported in $[-2T,2T]$ (it still grows
like $\log|\tau|$: the singularity of $K$ at $u=0$ lies inside the window). Moreover
$d\nu_K=2e^{-T}\big(1+O(e^{-2|u|})\big)\,d\nu_{\frac12,0}$:
the archimedean L\'evy density selects the scale $s=\frac12$ among the exponential columns.
\end{proposition}
\begin{proof}
Splitting \eqref{eq:m} at $2T$ gives the identity. The constant,
the prime cosines and each $1-\cos(\tau u)$ with $u\le2T$ have
Fourier support in $[-2T,2T]$. Integration against $K$ preserves
this support by Fubini on Schwartz tests, using
$|1-\cos(\tau u)|\le\min(2,\tau^2u^2)$. The density ratio follows
from $K(u)/2=e^{-u/2}/(1-e^{-2u})$ for $u>0$.
\end{proof}

At $T=\tfrac12\log3$ only the contact at two contributes. For spacing
$\log(3)/24$, Lemma~\ref{lem:F1} reduces the tail test to periodic rows.

\begin{proposition}[second-order cover]\label{prop:TAY}
Write $f=m+q+r$ with $q=-\sum_{n\in\NT}v_n\cos(\tau\log n)+\nuh$ ($\nu$ atomic) and
$r=\beta_0\gh0+\beta_1\gh1$, and put
$M_2=\sum_{n\in\NT}v_n\log^2n+2\sum_j|a_j|u_j^2+|\beta_0|H_0+|\beta_1|H_1$ with
$H_0=\int_{-2T}^{2T}u^2\cosh\frac u2\,du$, $H_1=\int_{-2T}^{2T}|u|^3\sinh\frac{|u|}2\,du$. Then
$\sup|q''+r''|\le M_2$, and on every cell $[x_0-\frac w2,x_0+\frac w2]\subset[0,\infty)$
\[
\min_{\text{cell}}f\ \ge\ m\!\big(x_0-\tfrac w2\big)+(q+r)(x_0)-|(q+r)'(x_0)|\,\tfrac w2-M_2\tfrac{w^2}8 .
\]
The same bound applies to each row $\nuh+Y_k$ of Lemma~\ref{lem:F1}, and to each row $H_k$ of
Lemma~\ref{lem:hybrid} with $M_2$ replaced by
$M_2^{\rm hyb}=M_2+\sum_\ell|c_\ell|\big[\tfrac{p_\ell!(p_\ell+1)(p_\ell+2)}{s_\ell^{p_\ell+3}}+\tfrac{4T(p_\ell+1)!}{s_\ell^{p_\ell+2}}+\tfrac{4T^2p_\ell!}{s_\ell^{p_\ell+1}}\big]$
$\big(=M_2+\sum_\ell|c_\ell|\int u^2\,d\nu_{s_\ell,p_\ell}\big)$; here $m(x+kP)$, $-|\beta_i|k_i(x+kP)$ and, in $H_N$,
the nondecreasing term $-\sum_\ell|c_\ell|\operatorname{cap}_{s_\ell,p_\ell}(x+NP)$ are evaluated at the left endpoint. A cell is certified when the
right side exceeds the target and is bisected otherwise; near a contact of $f$ the admissible width
is $\approx\sqrt{8c/M_2}$, so the number of evaluations grows like $c^{-1/2}$ rather than the
$c^{-1}$ of a Lipschitz cover.
\end{proposition}
\begin{proof}
Taylor's formula with integral remainder for $q+r\in C^2$, and monotonicity of $m$ and of
$-|\beta_i|k_i$. For the rows of Lemma~\ref{lem:hybrid}, $|(\nuh_{s,p})''(\tau)|=|\int u^2\cos(\tau u)\,d\nu_{s,p}|\le\int u^2\,d\nu_{s,p}$,
which is the bracket (write $u=(u-2T)+2T$ and use Gamma integrals), and the cap term in $H_N$ is
nondecreasing by Proposition~\ref{prop:T0}.
\end{proof}



\begin{remark}
For fixed $e\in(0,1)$ and $R\to\infty$, a feasible $\mu$ with mass $e$ outside $[-R,R]$ has value at least
$C_0-\Aconst+e(m(R)-m(0))$. Conversely, for every $\delta>0$ there is a feasible $\mu$ with mass at least
$e(1-o(1))$ outside $[-R+\sqrt R,R-\sqrt R]$ and value at most $\lamP+\delta+e(\log R+O(1))$. So mass at
radius $R$ costs $e\log R$ to leading order.
\end{remark}
\begin{proof}
If a feasible $\mu$ has mass $e$ outside $[-R,R]$, then
$\int\sT\,d\mu\ge C_0-\Aconst+e(m(R)-m(0))$. For the reverse scale
take any feasible measure within $\delta$ of $\lamP$ and mix its
frequency translates by $0,\pm R$ in proportions $1-e,e/2,e/2$.
This multiplies $\check\mu$ by $1-e+e\cos Ru$. The moment defect is
$v=e\,o(1)$ by bounded convergence since $\gh j$ vanish at infinity.
For $v\ne0$ mix in the bounded-cost measure of Lemma~\ref{lem:W}
with moments $-r_0v/|v|$ in proportion $|v|/(r_0+|v|)$.
The corrected measure has mass at least $e(1-o(1))$ outside
$[-R+\sqrt R,R-\sqrt R]$, since $\mu([-\sqrt R,\sqrt R])\to1$.
Its value is at most $\lamP+\delta+e(\log R+O(1))$: the translate
cost follows from $\log(2+|t\pm R|)\le\log(2+R)+\log(2+|t|)$
and the finite log moment of $\mu$.
\end{proof}

The direct tail bound subtracts $\Aconst+\|\nu\|_{\rm TV}
+\sum_i|\beta_i|k_i$ from $m_{\rm lo}$. Partial summation in the prime
number theorem gives $\Aconst=(4+o(1))e^T$.

\section{The explicit formula on flat-top weights}\label{sec:Z}

\begin{definition}\label{def:W}
$\Wcl$ is the set of even measurable $w:\R\to\R$ with $w\equiv1$ on $[-2T,2T]$ and
$|w(u)|\le Ce^{-(\frac12+\delta)|u|}$ for $|u|\ge2T$, for some $C,\delta>0$. $\Wcl^{\mathrm{reg}}\subset\Wcl$
consists of the $w\in C^\infty(\R)$ with compact support. For $w\in\Wcl$ define the even measure
\begin{equation}\label{eq:nuw}
d\nu_w=\Big[\tfrac12\big(1-w(u)\big)K(|u|)+2w(u)\cosh\tfrac u2\Big]\mathbf 1_{|u|\ge2T}\,du
-\sum_{n\notin\NT}\tfrac{v_n}2\,w(\log n)\,\big(\delta_{\log n}+\delta_{-\log n}\big),
\end{equation}
and $f_w:=\sT+\nuh_w+2\gh0$.
\end{definition}
The measure $\nu_w$ is finite ($K$ is integrable on $[2T,\infty)$, $|w|\cosh\frac u2$ is
integrable by the decay condition, and $\sum_nv_n|w(\log n)|\le2C\sum\Lambda(n)n^{-1-\delta}$) and
supported in $\{|u|\ge2T\}$. Hence $(\nu_w,2,0)$ is a dual feasible point, and by
Propositions~\ref{prop:weak} and \ref{prop:P}, $W_T\ge\inf f_w$ on $L^2(\IT)$.

\begin{theorem}\label{thm:Z}
\emph{(i)} For every $w\in\Wcl$ and $\tau\in\R$,
\begin{equation}\label{eq:fw}
\begin{split}
f_w(\tau)=C_0+\int_0^\infty\big(1-w(u)\cos\tau u\big)K(u)\,du&-\sum_{n\ge2}v_n\,w(\log n)\cos(\tau\log n)\\
&+2\int_\R w(u)\cosh\tfrac u2\cos(\tau u)\,du,
\end{split}
\end{equation}
which is the right-hand side of the Weil explicit formula for the even test function
$h_\tau(u)=w(u)\cos(\tau u)$.
\emph{(ii)} For $w\in\Wcl^{\mathrm{reg}}$, with $\wh(t)=\int w(u)e^{itu}du$ (an even entire function
of exponential type),
\begin{equation}\label{eq:Z}
f_w(\tau)=\sum_\rho\wh(\tau-\gamma_\rho),\qquad\rho=\tfrac12+i\gamma_\rho,
\end{equation}
the sum running over the non-trivial zeros of $\zeta$ with multiplicity, and converging absolutely.
A zero on the critical line contributes $\wh(\tau-\gamma_\rho)$ with $\gamma_\rho$ real, so on
the zeros of the critical line $f_w$ is the smoothing of the zero-counting measure by the kernel $\wh$.
\end{theorem}
\begin{proof}
For (i), insert \eqref{eq:nuw} into $f_w$. On $[0,2T]$ one has $w=1$;
on its complement the density $(1-w)K/2$ supplies the correction to
\eqref{eq:m}. The atoms supply exactly the omitted prime powers and the
remaining density completes the polar integral, giving \eqref{eq:fw}.
For (ii), apply the explicit formula \cite{Weil52}, \cite[Thm.~5.11]{IK}
to $h=w\cos(\tau\,\cdot)$. Its transform is
$(\wh(t-\tau)+\wh(t+\tau))/2$; evenness and $\rho\leftrightarrow1-\rho$
give \eqref{eq:Z}. Indeed Fourier inversion against $1$ and $\cos(tu)$
gives $1$ and $w(u)\cos(\tau u)$, respectively, in the archimedean
term. The interchange is justified by
$\int|\wh(t-\tau)|\int_0^\infty\min(2,t^2u^2)K(u)\,du\,dt<\infty$.
Uniform rapid decay on $|\Iml t|\le1/2$ and $N(t)\ll t\log t$
give absolute convergence over zeros.
\end{proof}


\begin{corollary}[certificates are statements about zeros]\label{cor:zeros}
For $w\in\Wcl^{\mathrm{reg}}$, a lower bound $f_w\ge c$ applies
to the full zero sum by \eqref{eq:Z}. Finite symmetric truncations
to real ordinates take negative values by Theorem~\ref{thm:finZ}.
\end{corollary}

\begin{theorem}[exhaustion by flat-top weights]\label{thm:Zexh}
Let $T\ge\frac12\log2$. \emph{(i)} $\sup_{w\in\Wcl^{\mathrm{reg}}}\inf_{\tau\in\R}f_w(\tau)=\lamW$.
\emph{(ii)} $\sup\{\inf_\tau[f_w+(\beta_0-2)\gh0+\beta_1\gh1]:\ w\in\Wcl^{\mathrm{reg}},\beta\in\R^2\}=\lamP$.
\end{theorem}
\begin{proof}
Weak duality gives the upper bounds. Corollary~\ref{cor:D1W} (for (i)) and
Theorems~\ref{thm:D1}, \ref{thm:E} (for (ii)) give finite even atomic certificates $\nu$ approaching
the two suprema. If $2T=\log n$, let
$B=-v_n(\delta_{2T}+\delta_{-2T})/2$; otherwise put $B=0$.
Every flat-top weight gives this same boundary atom since $w(2T)=1$.
Replace each atom of $\nu-B$ by a smooth bump of the same mass,
supported strictly outside the window, respecting reflection. The
resulting density $h$ has $\|h\|_1\le\|\nu-B\|_{\TV}$, and the
transform of $B+h\,du$ approximates $\widehat\nu$ uniformly on
any prescribed compact frequency interval by taking narrow bumps.
Choose that interval first: $m(\tau)\to\infty$ while the prime comb,
moment multipliers and transforms of all these measures are bounded
independently of the bump widths. Thus the lower bound for the
certificate loses at most $\varepsilon$.

To realize the density, set $w=(h-K/2)/(2\cosh(u/2)-K/2)$ on
$u>2T$ and reflect evenly. The denominator is positive there by
monotonicity and $2T\ge\log2$. Interpolate smoothly to $1$ at the
window edge, and cut off beyond a large radius, where $w=O(e^{-u})$.
These changes cost arbitrarily little in the $L^1$ norm of the density
in \eqref{eq:nuw}. Inside the cutoff force $w$ to vanish in small
disjoint neighborhoods of the finitely many $\log n>2T$ by smooth
cutoffs. Their total density cost tends to zero with their lengths.
The resulting atoms are exactly $B$, so the final density error
perturbs the transform uniformly by arbitrarily little.
\end{proof}

\begin{theorem}\label{thm:finZ}
Let $Z$ be a nonempty finite symmetric multiset of real ordinates.
If $w$ is even, continuous, real and integrable, with $w=1$ on
$[-2T,2T]$, $T>0$, then $\sum_{\gamma\in Z}\widehat w(t-\gamma)$
takes negative values.
\end{theorem}
\begin{proof}
Write $g(u)=\sum_{\gamma\in Z}e^{iu\gamma}$ and $f=\widehat{wg}$.
If $f\ge0$, Gaussian damping followed by monotone convergence gives
$\int f(t)\,dt=2\pi g(0)$ by continuity at zero. The measure
$d\nu=f(t)\,dt/(2\pi)$ therefore has Fourier transform $wg$ by
inversion and evenness. On the window this is also the transform
of $\mu=\sum_{\gamma\in Z}\delta_\gamma$.
Let $z_1,\ldots,z_l$ be the distinct ordinates. For $0<a<2T/l$,
the integral of $|\prod_{j=1}^l(e^{iat}-e^{iaz_j})|^2$ uses only
Fourier values at $ka$, $|k|\le l$. It vanishes for $\mu$, hence
for $\nu$. Fix one such $a$. Then $\nu$ is carried by the countable set
$\{t:a(t-z_j)\in2\pi\mathbb Z\text{ for some }j\}$; as $\nu$ is absolutely
continuous, $\nu(\mathbb R)=0$, contradicting $\nu(\mathbb R)=|Z|$.
\end{proof}

For a smooth compactly supported flat-top weight on $[-U,U]$ one
can also control the truncation error directly. If $D^{k+2}w$ is a
measure of total variation $V$, integration by parts gives
$|\widehat w(z)|\le Ve^{U|\Im z|}|z|^{-k-2}$ for $z\ne0$.
For $|t|\le H/2$ and $H\ge2$, group the omitted ordinates
$\gamma_\rho=(\rho-1/2)/i$ in shells
$2^jH<|\gamma_\rho|\le2^{j+1}H$. Their distances from $t$ are
at least $2^{j-1}H$, and $|\Im\gamma_\rho|\le1/2$.
The zero count $N(y)=O(y\log y)$ gives a total error
$O_{U,k}(VH^{-k-1}\log H)$ after summation. For autocorrelations
the on-line zeros instead contribute the nonnegative terms
$|\widehat F(\gamma)|^2$ individually.

\subsection{The Kre\u{\i}n extension problem at one window}\label{subsec:KEP}
Section~\ref{sec:Z} prescribes the Fourier transform of a positive measure on an interval.
This is the distributional analogue of Kre\u{\i}n's extension problem
\cite{Krein1940,Sasvari1994}.

\begin{definition}[the window datum, and $\mathrm{KEP}(T)$]\label{def:KEP}
Let $g$ be the even distribution defined on even $h\in C_c^\infty(\R)$ by
\begin{equation}\label{eq:gdatum}
\langle g,h\rangle=C_0\,h(0)+\int_0^\infty\big(h(0)-h(u)\big)K(u)\,du-\sum_{n\ge2}v_n\,h(\log n)
+2\int_\R h(u)\cosh\tfrac u2\,du,
\end{equation}
Let $\mathfrak S(T)$ consist of the positive
symmetric tempered measures $\nu$ (on the $\tau$-line) satisfying $\int\widehat h\,d\nu=\langle g,h\rangle$
for every even $h\in C_c^\infty(-2T,2T)$.
The \emph{Kre\u{\i}n extension problem} $\mathrm{KEP}(T)$ is \emph{solvable} if
$\mathfrak S(T)\ne\emptyset$, \emph{determinate} if $\#\mathfrak S(T)=1$, and \emph{indeterminate}
if $\#\mathfrak S(T)\ge2$.
\end{definition}
Even tests determine the symmetric part of a measure. By Schwartz--Bochner,
solvability means that the window datum extends to a positive definite tempered distribution.

\begin{lemma}[one negative square in the polar term]\label{lem:polar}
The polar kernel $2\cosh((s-t)/2)$ has one positive and one negative
square on any set of at least two distinct points.
\end{lemma}
\begin{proof}
Writing $a_i=e^{t_i/2}$, $b_i=e^{-t_i/2}$ gives the matrix $ab^t+ba^t$.
Its two nonzero eigenvalues $a\cdot b\pm|a||b|$ have opposite signs
by strict Cauchy--Schwarz.
\end{proof}

\begin{theorem}[a certificate is an absolutely continuous solution]\label{thm:KEP1}
Let $T\ge\frac12\log2$ and let $w\in\Wcl^{\mathrm{reg}}$ satisfy $f_w\ge c$ on $\R$.
\emph{(i)} If $c\ge0$ then $\mu_w:=f_w(\tau)\,d\tau/2\pi$ belongs to $\mathfrak S(T)$; in particular
$\mathrm{KEP}(T)$ is solvable, by an absolutely continuous solution.
\emph{(ii)} If $c>0$ then $\mathfrak S(T)$ contains a continuum of pairwise distinct absolutely
continuous measures; in particular $\mathrm{KEP}(T)$ is indeterminate.
\end{theorem}
\begin{proof}
(i) Theorem~\ref{thm:Z} gives $\widehat{wg}=f_w$. Its logarithmic growth
gives temperedness, evenness gives symmetry, and $wg=g$ on the window.
(ii) Take a nonzero even nonnegative smooth bump $\chi$ supported beyond
$\max(2T,\log2)$ in absolute value and away from prime atoms.
There $\chi g$ has density $\chi(2\cosh(u/2)-K(|u|)/2)$, strictly
positive wherever $\chi>0$ by the proof of Theorem~\ref{thm:Zexh}.
Its Fourier transform $r$ is bounded and nonzero. For
$|\varepsilon|<c/\|r\|_\infty$, the densities
$(f_w+\varepsilon r)/(2\pi)$ give distinct positive measures with
the same window datum, since $\chi g$ vanishes there.
\end{proof}

\begin{theorem}[solvability implies Weil positivity at that window]\label{thm:KEP2}
If $\mathfrak S(T)\ne\emptyset$ then $\lamW\ge0$.
\end{theorem}
\begin{proof}
For real $F\in C_c^\infty(\IT)$ its autocorrelation $h$ is even and
supported in $(-2T,2T)$. By \eqref{eq:gdatum} and \eqref{eq:W},
$W_T(F)=\langle g,h\rangle=\int|\Fh|^2d\nu\ge0$ for $\nu\in\mathfrak S(T)$.
Real and imaginary parts add for the real symmetric form, so $W_T\ge0$ on
complex combinations of real smooth tests. Lemma~\ref{lem:form-density}
extends the inequality to $\mathcal D_T$.
\end{proof}

\begin{theorem}[the dichotomy at one window]\label{thm:KEPdich}
For every $T\ge\frac12\log2$,
\[
\lamW>0\ \Longrightarrow\ \mathrm{KEP}(T)\ \text{is indeterminate}\ \Longrightarrow\ \lamW\ge0 ,
\]
so certificates and indeterminacy coincide except on the set $\{T:\lamW=0\}$.
\end{theorem}
\begin{proof}
If $\lamW>0$, Theorem~\ref{thm:Zexh}(i) gives $w\in\Wcl^{\mathrm{reg}}$ with $\inf f_w>0$, and
Theorem~\ref{thm:KEP1}(ii) applies; the second implication is Theorem~\ref{thm:KEP2}.
\end{proof}

\section{The archimedean form on primitive tests}\label{sec:archstr}

Put, for $F\in L^2(\IT)$,
\[
\Ainf(F)=2\Rel\big(A(F)\overline{B(F)}\big)+\frac1{2\pi}\int_\R m(\tau)|\Fh(\tau)|^2\,d\tau
=2\Rel\big(A(F)\overline{B(F)}\big)+C_0\|F\|^2+\Lev(\cF),
\]
the complete Weil form with the prime comb deleted (the second equality is Theorem~\ref{thm:A} with
$\NT=\emptyset$), and write $\PiT(F)=-\sum_{n\in\NT}v_n\Rel\cF(\log n)$ for the prime form (the real part is redundant for real $F$), so that
$W_T=\Ainf+\PiT$. Let $\laminf(T)$ be the bottom of $\Ainf$ on the unit sphere of $L^2(\IT)$ and
$\laminfP(T)$ its bottom on the unit sphere of $\PT$, where the pole term vanishes identically and
$\Ainf=C_0\|F\|^2+\Lev(\cF)$.

\begin{proposition}[two archimedean thresholds]\label{prop:arch}
(i) Both minima decrease weakly with $T$, and $\laminfP$ decreases strictly.
(ii) $C_0\le\laminfP(T)\le C_0+\varphi((3\pi/2T)^2)$ and
$\laminf(T)\le\laminfP(T)$. In particular $\laminfP(T)\downarrow C_0$.
(iii) For $2T<\log2$, $W_T=\Ainf$ and $\laminf(T)\ge0$.
(iv) The suprema $T^*_\infty,T^*_{\infty,\mathrm P}$ of the radii of
strict positivity for the two minima satisfy
$\tfrac12\log2\le T^*_\infty\le T^*_{\infty,\mathrm P}\le T_0(0)$,
where $\varphi((3\pi/(2T_0(0)))^2)=-C_0$ and $T_0(0)\simeq0.74921$.
\end{proposition}
\begin{proof}
Zero extension enlarges the feasible sets; strictness on primitive tests
is Lemma~\ref{lem:strict}. Inclusion orders the minima, $\Lev\ge0$,
and the three-dimensional trial space of Corollary~\ref{cor:mono}
gives (ii) with the prime comb omitted. Its upper bound tends to $C_0$.
There are no prime contacts below $\tfrac12\log2$, so $W_T=\Ainf$ there. For $2T\le\log2$ the
Weil functional is nonnegative on tests supported in $\IT$ \cite[Thm.~1, p.~283]{Yoshida1992}; by
Lemma~\ref{lem:form-density} this extends to $\mathcal D_T$, so $\laminf(T)\ge0$ for $T\le\frac12\log2$,
which is (iii). The bottom of $\Ainf$ is attained (Lemma~\ref{lem:compact}; the polar term is bounded),
and $\laminf$ decreases strictly (the remark after Lemma~\ref{lem:strict}); hence
$\laminf(T)>\laminf(\frac12\log2)\ge0$ for $T<\frac12\log2$, so $T^*_\infty\ge\frac12\log2$. The other
inequalities in (iv) follow from (ii).
See also \cite{Bombieri2000,ConnesConsani2021,Zhu2026}.
\end{proof}

\begin{proposition}[nodal structure of $\PT$]\label{prop:nodes}
The ground state of $\mT=C_0+\Lev$ on the full window is simple and can
be chosen strictly positive almost everywhere. Every nonzero real primitive
test, including a real primitive ground state, has at least two sign changes:
its signs cannot be ordered on two consecutive subintervals, even allowing an
empty subinterval. For an even test a change of sign on $(0,T)$ is required.
\end{proposition}
\begin{proof}
Zero extension and symmetry give the jump and killing terms
\[
 \Lev(\cF)=\tfrac14\iint_{\IT^2}K(|x-y|)|F(x)-F(y)|^2\,dx\,dy
 +\tfrac12\int_{\IT}|F(x)|^2\!\int_{|y|>T}K(|x-y|)\,dy\,dx.
\]
The Markov contraction decreases this form and $K>0$ connects every
pair of positive-measure sets, so its semigroup is positivity improving.
Compactness follows from fixed support and the uniform Fourier tail
bound $\int_{|\tau|>R}|\Fh|^2\le2\pi\Lev(\cF)/(m(R)-C_0)$;
Perron--Frobenius gives the assertion. If the signs of $F$ were
ordered about $t_0$, then $G=e^{t/2}-e^{t_0}e^{-t/2}$ or $-G$
would have the matching signs and $\int FG\ne0$, contradicting
both Tate conditions. For even $F$ the positive weight $\cosh(t/2)$
already forces both signs on $(0,T)$.
\end{proof}

For $j=\mathrm{ev},\mathrm{odd}$ let $\PT^j$ be the corresponding parity
subspace, $\lambda_j,n_j$ its spectral bottom and negative index
(counting multiplicities), and $T_j^*=\sup\{T:\lambda_j(T)>0\}$.
Its moment condition is orthogonality to $\cosh(t/2)$ or $\sinh(t/2)$,
respectively.

\begin{lemma}[parity splitting]\label{lem:parity}
On $\PT$ the negative indices add and $\laminfP=\min_j\lambda_j$.
\end{lemma}
\begin{proof}
Reflection exchanges the Tate moments and commutes with the real even
multiplier $m$. Its two eigenspaces thus split both the norm and the form
orthogonally, giving the assertions by compact resolvent. Complexification
preserves the eigenvalues and their multiplicities.
\end{proof}

\begin{lemma}[zero extension]\label{lem:zeroext}
In each sector $n_j$ increases weakly and $\lambda_j$ decreases weakly
with the radius.
\end{lemma}
\begin{proof}
Zero extension preserves parity, both moments, the norm and the
autocorrelation on $\mathbb R$, hence also $\mT$.
\end{proof}

\begin{lemma}[strict monotonicity]\label{lem:strict}
For every $T>0$ the bottoms $\lambda_{\mathrm{ev}}(T)$, $\lambda_{\mathrm{odd}}(T)$ and
$\laminfP(T)$ are attained, and all three are strictly decreasing on $(0,\infty)$. Consequently, for each
sector, $\lambda_{\mathrm{par}}(T)>0$ for $T<T^*_{\mathrm{par}}$ and $\lambda_{\mathrm{par}}(T)<0$ for
$T>T^*_{\mathrm{par}}$; $\lambda_{\mathrm{par}}(T)<0$ if and only if $n_{\mathrm{par}}(T)\ge1$;
$T^*_{\infty,\mathrm P}=\min(T^*_{\mathrm{ev}},T^*_{\mathrm{odd}})$; and $\laminfP$ is positive on
$(0,T^*_{\infty,\mathrm P})$ and negative on $(T^*_{\infty,\mathrm P},\infty)$.
\end{lemma}
\begin{proof}
The form domain embeds compactly by Lemma~\ref{lem:compact} (the prime form is bounded, so the
sublevel sets of $\mT$ and of $\mathfrak a_T$ on the unit sphere are comparable). Both
moment constraints and parity are closed in $L^2$, so the direct method
gives a minimizer; real and imaginary parts split and one can take it real.
If a sector bottom agrees at $T'<T$, its normalized minimizer $F^*$ at
$T'$ also minimizes at $T$ by zero extension. Test its eigenvalue equation
against $G_0+\epsilon JG_0$, where $\epsilon$ is its parity and
$G_0\in C_c^\infty(T',T)$ has both moments zero. Disjoint supports and
reflection give $2\mT(F^*,G_0)=0$. Polarizing the jump integral, this
equals $-\iint K(|y-x|)F^*(x)G_0(y)\,dx\,dy$. Hence
$\Phi(y):=\int_{I_{T'}}K(|x-y|)F^*(x)\,dx$ satisfies $\int\Phi G_0=0$ for every such $G_0$, so
$\Phi=\alpha e^{y/2}+\beta e^{-y/2}$ on $(T',T)$. For $y>T'\ge x$, Lemma~\ref{lem:K} and Tonelli give
$\Phi(y)=2\sum_{k\ge0}c_ke^{-(2k+\frac12)y}$ with $c_k=\int F^*(x)e^{(2k+\frac12)x}dx$,
$|c_k|\le\|F^*\|_1e^{(2k+\frac12)T'}$; so $e^{y/2}\Phi(y)=2\sum_kc_kw^k$, $w=e^{-2y}$, is a convergent
power series for $0<w<e^{-2T'}$, and the identity extends to all $y>T'$. As $y\to\infty$, $\Phi\to0$,
so $\alpha=0$; then $2\sum_kc_kw^k\equiv\beta$, so $c_k=0$ for $k\ge1$ (the value $c_0=\beta/2$ is not
needed). With $u=e^{2x}$, $c_k=\int f(u)u^k\,du$ for
$f(u)=\frac12F^*(\frac12\log u)u^{-3/4}\in L^1[e^{-2T'},e^{2T'}]$, so all moments of $uf(u)$ vanish and
$uf=0$ by Weierstrass, whence $F^*=0$, contradicting $\|F^*\|=1$. For $\laminfP$ use Lemma~\ref{lem:parity}.

Strict decrease gives the asserted signs on either side of the supremum.
The negative index is positive exactly when the minimum is negative,
and the primitive threshold is the smaller sector threshold.
\end{proof}
The same argument gives strict decrease of $\laminf$. Its bottom is attained by Lemma~\ref{lem:compact}
(the polar term $2\Rel(A\bar B)$ is a bounded rank-two perturbation). If $\laminf(T')=\laminf(T)$ with
$T'<T$, a real minimizer $F^*$ at $T'$ also minimizes at $T$ by zero extension. For $G_0\in C_c^\infty(T',T)$
with both moments zero, $A(G_0)=B(G_0)=0$ and $\langle F^*,G_0\rangle=0$, so the eigenvalue equation gives
$\mT(F^*,G_0)=0$; by disjoint supports this is $-\frac12\iint K(|y-x|)F^*(x)G_0(y)\,dx\,dy=0$, and the
argument of Lemma~\ref{lem:strict} from the definition of $\Phi$ on gives $F^*=0$, a contradiction. No parity
is needed.

\section{The two Tate directions}\label{sec:complete}

\begin{lemma}[non-attainment for the complete form]\label{lem:Onat}
Let $T\ge\frac12\log2$ and $w\in\Wcl$. Then $\inf_\R f_w<\lamW$.
\end{lemma}
\begin{proof}
Weak duality gives $c=\inf f_w\le\lamW$. Equality would imply
$\int(f_w-c)|\Fh_\star|^2=0$ for a normalized ground state
by Proposition~\ref{prop:P} and attainment of the bottom of $W_T$ (Lemma~\ref{lem:compact}); the $\nuh_w$
part integrates to zero against $|\Fh_\star|^2$, since $\nu_w$ is carried by $\{|u|\ge2T\}$, where $c_{F_\star}$
vanishes.
The nonzero entire transform and continuity then force $f_w\equiv c$.
On $0<u<\min(2T,\log2)$ Fourier inversion gives
$2\cosh(u/2)=K(u)/2$, impossible near zero.
\end{proof}

\begin{theorem}[flat-top certificates and the complete bottom]\label{thm:Al}
\emph{(a)} If $f_w\ge0$ on $\R$ for some $w\in\Wcl$, regular or not, then $W_T\ge0$ on the form domain;
hence $\lamW\ge0$ and $\lamA\ge0$.
\emph{(b)} Let $T\ge\frac12\log2$. Then there is $w\in\Wcl$ with $f_w\ge0$ on $\R$ if and only if
$\lamW>0$, and $w$ can then be taken in $\Wcl^{\mathrm{reg}}$ with $\inf f_w>0$. Likewise, if $\lamP>0$
there are $w\in\Wcl^{\mathrm{reg}}$ and $\beta\in\R^2$ with $\inf_\tau[f_w+(\beta_0-2)\gh0+\beta_1\gh1]>0$.
\end{theorem}
\begin{proof}
(a) $(\nu_w,2,0)$ is dual feasible and $W_T\ge\inf f_w$ on $L^2(\IT)$ (the paragraph after
Definition~\ref{def:W}), and $W_T=\mathfrak a_T$ on $\PT$. (b) ``$\Leftarrow$'' and the regular choice:
Theorem~\ref{thm:Zexh}(i). ``$\Rightarrow$'': $\lamW\ge\inf f_w\ge0$ by weak duality, and $\lamW>\inf f_w$ by
Lemma~\ref{lem:Onat}. The last sentence is Theorem~\ref{thm:Zexh}(ii).
\end{proof}

\begin{definition}[the alignment inequality]\label{def:Al}
For $w\in\Wcl$ let $\mathrm{Al}_T(w)$ be the statement: for every $\tau\in\R$,
\[
-2\Rel\sum_{n\ge2}\frac{\Lambda(n)w(\log n)}{n^{\frac12+i\tau}}\ \ge\
-\Big[C_0+\int_0^\infty\big(1-w(u)\cos\tau u\big)K(u)\,du\Big]-2\int_\R w(u)\cosh\tfrac u2\cos(\tau u)\,du .
\]
The left side is the smoothed prime sum on the critical line, $-\sum_nv_nw(\log n)\cos(\tau\log n)$; by
\eqref{eq:fw}, $\mathrm{Al}_T(w)$ is exactly the statement $f_w\ge0$ on $\R$ of Theorem~\ref{thm:Al}.
\end{definition}

\begin{lemma}[the alignment inequality lives on a band]\label{lem:Alband}
For $w\in\Wcl$ let $\tau_{\mathrm{env}}(w)$ be the unique positive root of
$m_{\mathrm{lo}}(\tau)-\Aconst-\|\nu_w\|_{\TV}-2k_0(\tau)=0$, with $m_{\mathrm{lo}}$ as in Lemma~\ref{lem:mbounds}(ii)
and $k_0$ as in Lemma~\ref{lem:gbounds}. Then $f_w>0$ on $(\tau_{\mathrm{env}}(w),\infty)$, so $\mathrm{Al}_T(w)$ is
equivalent to its restriction to $|\tau|\le\tau_{\mathrm{env}}(w)$.
\end{lemma}
\begin{proof}
$m_{\mathrm{lo}}$ increases from $-\infty$ to $+\infty$ and $k_0$ decreases (Lemmas~\ref{lem:mbounds}, \ref{lem:gbounds}),
so the left side is strictly increasing and the root is unique. The envelope for $(\nu_w,2,0)$ gives $f_w\ge m_{\mathrm{lo}}-\Aconst-\|\nu_w\|_{\TV}-2k_0$, and $f_w$ is even.
\end{proof}

For $T\ge\frac12\log2$ the optimal uniform margin is $\lamW$
by Theorem~\ref{thm:Zexh}(i). For regular weights,
Theorem~\ref{thm:Z}(ii) identifies $\mathrm{Al}_T(w)$ with
$\sum_\rho\wh(\tau-\gamma_\rho)\ge0$.

\begin{remark}[the prime side of a primitive test is the Chebyshev remainder]\label{rem:Cheb}
Let $\psi_\Lambda(x)=\sum_{n\le x}\Lambda(n)$ and $r_\Lambda(u):=\psi_\Lambda(e^u)-e^u$. For real $F\in\PT\cap\mathcal D_T$,
\begin{equation}\label{eq:Cheb}
\mathfrak a_T(F)=C_0\|F\|^2+\Lev(\cF)+2\int_0^{2T}\cF(u)e^{-u/2}\,du-2\int_{(0,2T)}e^{-u/2}\cF(u)\,dr_\Lambda(u).
\end{equation}
The prime sum integrates $2e^{-u/2}\cF$ against $d\psi_\Lambda(e^u)$,
with zero boundary contribution since $\cF(2T)=0$. Split off $e^u du$.
Evenness and $\int_{-2T}^{2T}\cF(u)e^{u/2}du=A(F)B(F)=0$
give \eqref{eq:Cheb}. If $\lamA\ge0$, unit tests approaching the
infimum therefore make the last integral, with its factor two, agree
with the preceding three terms up to any error $\eta>\lamA$.
Corollary~\ref{cor:margin1} gives $\lamA\le4\cdot10^9Te^{-15.449T}$ for $T\ge1$.
\end{remark}

To pass from primitive to complete positivity at a fixed window it
suffices to consider real tests, since real and imaginary parts add.
For $j=\mathrm{ev},\mathrm{odd}$ let $\mathcal D_T^j$ denote the
corresponding parity sector, $\ell_j(F)=\langle F,\chi_j\rangle$
as in \S\ref{subsec:conv}, and $\varepsilon_j=+1,-1$, respectively.
Then $A=\ell_{\rm ev}+\ell_{\rm odd}$, $B=\ell_{\rm ev}-\ell_{\rm odd}$,
and $\mathcal Q_T^j=\ker(\ell_j|\mathcal D_T^j)$.
Write $\lambda_{\rm A}^j(T)$ for the bottom of $\mathfrak a_T$ on
$\mathcal Q_T^j$ and $\lambda^{\rm W}_{T,j}$ for that of $W_T$
on $\mathcal D_T^j$, using unit vectors in both cases.

\begin{theorem}[the parity Schur complement of the complete form]\label{thm:parityschur}
Fix $T>0$. In one parity sector abbreviate its moment, moment kernel,
primitive bottom and complete bottom by $\ell,\mathcal Q,\lambda,\omega$.
\emph{(o)} The sectors are orthogonal for $L^2$ and both forms;
$W_T(F)=\mathfrak a_T(F)+2\varepsilon\ell(F)^2$ on each. The
bottom of either form is the minimum of its two sector bottoms.

\emph{(a)} If zero belongs to the resolvent of the primitive operator,
there is a unique $h$ of moment one form-orthogonal to $\mathcal Q$.
Write $g=\mathfrak g_{\rm par}(T)=W_T(h)$. Then
\emph{(a1)} $W_T(G+xh)=\mathfrak a_T(G)+gx^2$ for $G\in\mathcal Q$.
\emph{(a2)} The negative index increases by $[g<0]$, and the radical
is $\R h$ for $g=0$ and trivial otherwise.
\emph{(a3)} For $\lambda>0$ both forms attain their moment-one
minima at $h$. \emph{(a4)} Put $b=g/\|h\|^2$. Then $\omega\le b$
and $\|h\|\ge1/\|\chi_{\rm par}\|$. For $g,\lambda>0$ one has
$2\omega\ge\min(\lambda,b)$, hence $b\le2\omega$ if $2\omega<\lambda$.

\emph{(b)} If $\lambda=0$, let $N$ be the primitive kernel. For
$z\in N$ write $\mathfrak a_T(z,\cdot)=\alpha(z)\ell$ on the sector.
If $\alpha\ne0$, then $\omega<0$. If $\alpha=0$ (the range condition),
$N$ lies in the radical of $W_T$ and (a) applies on $N^\perp$.

\emph{(c)} Put $\mathfrak G_T=\operatorname{diag}(\mathfrak g_{\rm ev},
\mathfrak g_{\rm odd})$ whenever both entries are defined.
For invertible primitive operator the negative indices of the primitive
form and $\mathfrak G_T$ add to that of $W_T$, and the nullity of
$W_T$ equals that of $\mathfrak G_T$. At every window $\lamW>0$ is
equivalent to $\lamA>0$ and $\mathfrak G_T\succ0$; $\lamW\ge0$ is
equivalent to $\lamA\ge0$, the range condition in each zero sector,
and $\mathfrak G_T\succeq0$.
\end{theorem}
\begin{proof}
Reflection preserves the Fourier multiplier and interchanges $A,B$,
giving (o) since $AB=\ell_{\rm ev}^2-\ell_{\rm odd}^2$. Choose
a smooth $H_0$ in the sector with $\ell(H_0)=1$. Subtracting
$\ell(F)H_0$ from smooth approximants gives density on the moment
kernel. The closed restricted form has compact resolvent by
Lemma~\ref{lem:compact}. If its operator is invertible, its negative
space is finite dimensional and on its positive space the spectral gap
gives coercivity in the shifted form norm. Thus finite-dimensional
inversion and Lax--Milgram give a unique $Y\in\mathcal Q$ with
$\mathfrak a_T(Y,G)=\mathfrak a_T(H_0,G)$ for every $G\in\mathcal Q$.
Here the right side is continuous by shifted form Cauchy--Schwarz.
Now $h=H_0-Y$, and subtracting $\ell(F)h$ gives (a1). Reading the
two blocks gives (a2) and (a3). Testing $h$ gives the upper bound
in (a4); the lower bound follows from
$\|G+xh\|^2\le2(\|G\|^2+x^2\|h\|^2)$.

For (b), compact resolvent gives a finite-dimensional zero eigenspace
and a gap above it. For $z\in N$ the functional vanishes on $\ker\ell$,
so $\alpha(z)=\mathfrak a_T(z,H_0)$. The energy on $H_0+xz$ is
$W_T(H_0)+2x\alpha(z)$, which takes negative values if $\alpha(z)\ne0$.
Otherwise, if $\alpha$ vanishes identically, orthogonal projection off
$N$ preserves the moment and gives the asserted reduction. Adding the
sector decompositions gives (c), using (a4) for strict positivity.
\end{proof}

\begin{proposition}[$\mathfrak G_T$ is the polar capacity]\label{prop:polarcap}
Assume $\lamA\ge0$ (by Theorem~\ref{thm:parityschur}(o), equivalently both primitive sector bottoms are
$\ge0$) and the range condition of Theorem~\ref{thm:parityschur}(b) in every sector whose primitive bottom is
zero; define $\mathfrak g_{\rm par}$ by Theorem~\ref{thm:parityschur}(a) in a sector with positive bottom, and by (a)
applied on $N^\perp$ in a sector with zero bottom. Under these assumptions,
let $E_T(\beta,\alpha)$ be the infimum of $W_T$ at real tests in
$\mathcal D_T$ with $B=\beta$, $A=\alpha$. This infimum is attained and
\[
 E_T(\beta,\alpha)=\tfrac14\bigl[\mathfrak g_{\rm ev}(\alpha+\beta)^2+
 \mathfrak g_{\rm odd}(\alpha-\beta)^2\bigr].
\]
Thus $\lamW\ge0$ exactly when $E_T\ge0$ on $\R^2$. In parity
coordinates the polar correction to the primitive Schur complement
is $D_{\rm pol}=\operatorname{diag}(2,-2)$.
\end{proposition}
\begin{proof}
The parity moments are $c=(\alpha+\beta)/2$ and $s=(\alpha-\beta)/2$.
The decomposition of Theorem~\ref{thm:parityschur} adds
$\mathfrak a_T(G)$, $G\in\PT$, to the displayed value; its minimum
is zero at $G=0$, because $\mathfrak a_T\ge0$ on $\PT$ by $\lamA\ge0$; in a zero sector the kernel $N$ lies in
the radical of $W_T$ and has zero moment, so adding it changes neither $W_T$ nor $(A,B)$
(Theorem~\ref{thm:parityschur}(b)). Without $\lamA\ge0$ the formula fails: scaling a primitive $G$ with
$\mathfrak a_T(G)<0$ gives $E_T=-\infty$. The polar correction is $2c^2-2s^2=2\alpha\beta$.
\end{proof}

To compare the primitive multipliers with the polar value $(2,0)$, put
$A'(F)=\int tF(t)e^{t/2}dt$, $B'(F)=\int tF(t)e^{-t/2}dt$, and set
\[
 \mathfrak m(F)=\big(\Rel(A\bar B),\Rel(A'\bar B-A\bar B')\big),\qquad
 V_T(\beta)=\inf_{F\in\mathcal D_T,\ \|F\|=1}
       [\mathfrak a_T(F)+\beta\cdot\mathfrak m(F)] .
\]
The two components of $\mathfrak m$ have multipliers $\gh0,\gh1$ by
Theorem~\ref{thm:A}, and $V_T(2,0)=\lamW$ by Proposition~\ref{prop:P}.
All three forms add on the real and imaginary parts since the symbols
are even and $\Iml(\widehat u\,\overline{\widehat v})$ is odd for real $u,v$.

\begin{theorem}[the optimal multipliers]\label{thm:mult}
Fix $T>0$.
\emph{(i)} $V_T$ is finite, concave and continuous on $\R^2$, $V_T(\beta)\to-\infty$ as $|\beta|\to\infty$, and
\[
\lamP=\max_{\beta\in\R^2}V_T(\beta).
\]
Let $\mathcal M_T\subset\R^2$ be the set of maximizers; it is nonempty, compact and convex.
\emph{(ii)} If $(\nu_k,\beta^{(k)})$ are certificates of values $c_k\to\lamP$, then
$\operatorname{dist}(\beta^{(k)},\mathcal M_T)\to0$. Hence every such sequence satisfies $\beta^{(k)}\to(2,0)$ if and
only if $\mathcal M_T=\{(2,0)\}$, and $\lamP=\lamW$ if and only if $(2,0)\in\mathcal M_T$.
\emph{(iii)} Let $\mathcal Z_T$ be the set of real unit vectors $F\in\mathcal D_T$ with $W_T(F)=\lamW$; it is nonempty
and compact in $L^2(\IT)$. Then $(2,0)\in\mathcal M_T$ if and only if
$0\in\operatorname{conv}\{\mathfrak m(F):F\in\mathcal Z_T\}$.
\emph{(iv)} Suppose $\lamW<\lamA$. Then every $F\in\mathcal Z_T$ of definite parity has
$\Rel(A\bar B)=\ell_{\mathrm{ev}}(F)^2>0$ if $F$ is even and $\Rel(A\bar B)=-\ell_{\mathrm{odd}}(F)^2<0$ if $F$ is odd.
Consequently, if moreover $\lambda^{\mathrm W}_{T,\mathrm{ev}}\ne\lambda^{\mathrm W}_{T,\mathrm{odd}}$, then $\lamW<\lamP$, and
the multipliers of every maximizing sequence of certificates stay at a positive distance from $(2,0)$.
\emph{(v)} If $\lamW=\lamA$, then $\lamP=\lamW$ and $(2,0)\in\mathcal M_T$.
\end{theorem}
\begin{proof}
(i) By Cauchy--Schwarz the two moments of a unit vector are bounded
by $2\sinh T$ and $4T\sinh T$. Thus the affine functions defining
$V_T$ have uniformly bounded slope; their infimum is finite, concave
and Lipschitz since $\mathfrak a_T$ is bounded below.
The four finite-energy measures of Lemma~\ref{lem:W} realize a ball
of moments of radius $r>0$ by convex combination and the factorization
of Proposition~\ref{prop:weak}. Their energies have a common upper
bound $V_0$, so the moment $-r\beta/|\beta|$ gives
$V_T(\beta)\le V_0-r|\beta|$ for $\beta\ne0$. Feasible measures
have zero moments and hence $V_T(\beta)\le\lamP$. Every certificate
of value $c$ with multiplier $\beta$ gives $c\le V_T(\beta)$:
its annihilating measure integrates to zero against each autocorrelation.
Strong duality (Theorem~\ref{thm:D1}) now gives $\sup V_T=\lamP$.
Continuity, decay at infinity and concavity give a nonempty compact
convex set of maximizers.

(ii) Since $c_k\le V_T(\beta^{(k)})\le\lamP$, the multipliers
remain bounded and all their limit points maximize $V_T$.
For the converse at a fixed maximizer $\beta$, the remark after
Corollary~\ref{cor:D1W} gives certificates with multiplier $\beta$ and
values approaching $V_T(\beta)$. The last equivalence follows
from $V_T(2,0)=\lamW$.
(iii) Unit sublevels of $W_T$ are compact in $L^2$: the polar
term is bounded there, so Lemma~\ref{lem:compact} gives relative
compactness and lower semicontinuity gives closedness. Hence the
bottom is attained. Additivity on real and imaginary parts gives
a real unit minimizer and expresses the moment vector of each
complex unit minimizer as a convex combination of real ones.
Put $\beta_*=(2,0)$. For each minimizer $F$,
$V_T(\beta)\le\lamW+(\beta-\beta_*)\cdot\mathfrak m(F)$.
If their moment vectors have zero in their convex hull, averaging
gives $V_T(\beta)\le V_T(\beta_*)$. Otherwise compactness and
strict separation give $d\cdot\mathfrak m(F)\ge\delta>0$ on all
minimizers. By compactness of sublevels and continuity of the moments,
$d\cdot\mathfrak m(F)\ge\delta/2$ on some unit sublevel
$W_T\le\lamW+\eta$, $\eta>0$. Outside it,
$W_T(F)+t d\cdot\mathfrak m(F)\ge\lamW+\eta-t C_T|d|$,
where $C_T$ bounds the moments in (i). For sufficiently small $t>0$
both regions give a uniform strict improvement over $\lamW$,
so $\beta_*$ cannot maximize $V_T$.

(iv) If $\lamW<\lamA$, a unit minimizer of definite parity
has nonzero parity moment, since otherwise it would be primitive.
Its first moment coordinate is $A B=A^2$ in the even sector and
$-A^2$ in the odd sector. By parity additivity any component
in a sector of strictly larger bottom must vanish in a minimizer.
Thus when the sector bottoms differ the first coordinates of all
real minimizers have one strict sign. Apply (iii) and (ii).
(v) $\lamW\le\lamP\le\lamA$ (Theorem~\ref{thm:A}, Remark~\ref{rem:Kprim}) and (ii).
\end{proof}

\section{Lower bounds on finite windows}\label{sec:log5}

Let $T_5=\tfrac12\log5$. For $\tfrac12\log4<T\le T_5$ the
contacts are $2,3,4$ and $\Aconst=2.94197\ldots$.
For related radii and margins see \cite[Thm.~1.2]{Zhu2026}
and \cite[Thms.~A, B]{Liu2026}.

\begin{theorem}[Lower bounds at $T_5$]\label{thm:log5W}\label{thm:log5A}
At $T_5$ the even and odd primitive minima are at least
$2.5\cdot10^{-12}$ and $3.0\cdot10^{-10}$ respectively.
The complete even and odd minima are at least $10^{-18}$ and
$10^{-16}$. For $0<T\le T_5$ one has $\lamA\ge2.5\cdot10^{-12}$,
$\lamW\ge10^{-18}$ and $\mathfrak G_T\succ0$.
\end{theorem}
\begin{proof}
In Proposition~\ref{prop:lowercert} use $c=0$ on the primitive spaces
and $c=-\varepsilon_{\rm par}/2$ on the complete parity spaces.
The ball enclosures in \S\ref{subsec:encllog5} give the four bounds; this step is a
computation in ball arithmetic, whose program is not reproduced in this paper.
Take the smaller parity minimum by Theorem~\ref{thm:parityschur}(o)
and restrict the support by Lemma~\ref{lem:monoT}.
Theorem~\ref{thm:parityschur}(c) gives the matrix inequality.
\end{proof}

\begin{corollary}[indeterminacy for $T\le\frac12\log5$]\label{cor:KEPlog5}
$\mathrm{KEP}(T)$ is indeterminate for every $0<T\le\frac12\log5$.
\end{corollary}
\begin{proof}
Theorems~\ref{thm:log5W} and \ref{thm:KEPdich} give the claim for
$\frac12\log2\le T\le T_5$. For $0<T<\frac12\log2$, every even $h\in C_c^\infty(-2T,2T)$ lies in
$C_c^\infty(-\log2,\log2)$, so $\mathfrak S(\frac12\log2)\subset\mathfrak S(T)$ and
$\#\mathfrak S(T)\ge\#\mathfrak S(\frac12\log2)\ge2$.
\end{proof}

\begin{lemma}[a Ky Fan ceiling for the time--band limiting operator]\label{lem:KyFan}
Let $S_c=Q_TB_sQ_T$ on $L^2(\IT)$, where $Q_T$ is restriction to $\IT$ and $B_s$ the band projection
to $|\tau|<s$, and put $c=Ts$. Let $j_n$ be the spherical Bessel functions and
$G_k(x)=\sum_{j<k}(2j+1)j_j(x)^2\in[0,1]$. Then, in the orthonormal Legendre basis
$e_j(t)=\sqrt{(2j+1)/2T}\,P_j(t/T)$,
\[
\langle S_ce_j,e_j\rangle=\frac{2(2j+1)}\pi\int_0^cj_j(x)^2dx,\qquad
\sum_{j\ge0}\langle S_ce_j,e_j\rangle=\frac{2c}\pi=\operatorname{tr}S_c ,
\]
and for every $k\ge0$
\[
\lambda_k(S_c)\ \le\ \Lambda_k(c):=\frac2\pi\int_0^c\big(1-G_k(x)\big)\,dx .
\]
In particular $\Lambda_0(c)=2c/\pi$ is the trace and $\Lambda_k(c)=O(c^{2k+1})$ as $c\to0$.
\end{lemma}
\begin{proof}
The Fourier--Legendre identity gives
$|\widehat e_j(\tau)|^2=2T(2j+1)j_j(\tau T)^2$. Plancherel gives
the diagonal entry and $\sum_j(2j+1)j_j(x)^2=1$ gives the trace.
Positivity and Ky Fan's principle give
$\lambda_k\le\sum_{j\ge k}\lambda_j\le\operatorname{tr}S_c-
\sum_{j<k}\langle S_ce_j,e_j\rangle=\Lambda_k(c)$.
\end{proof}

\begin{corollary}[quadrature-free form]\label{cor:P1}
$|j_n(x)|\le x^n/(2n+1)!!$ for $x\ge0$, hence
$\lambda_k(S_c)\le\Lambda_k(c)\le\Lambda^*_k(c):=\frac2\pi\sum_{j\ge k}c^{2j+1}/\big((2j+1)!!\big)^2$, a convergent
series of rationals in $c$: the bound is evaluated by arithmetic alone, with no quadrature.
\end{corollary}


\begin{lemma}[bordered inertia]\label{lem:border}
Let $q$ be a real quadratic form on a real vector space $V$, $f$ a linear functional on
$V$ and $v_0\in V$ with $f(v_0)=1$ and $q(\cdot,v_0)$ finite. On $V\oplus\R$ put
$b(v,\xi)=q[v]+2\xi f(v)$. Then $n_-(b)=n_-(q|_{\ker f})+1$, $n_-$ being the largest dimension of a
subspace on which the form is negative definite.
\end{lemma}
\begin{proof}
Write $v=w+tv_0$ with $w\in\ker f$. Then $b=q[w]+2tq(w,v_0)+t^2q[v_0]+2\xi t$, and the invertible
substitution $\eta=\xi+q(w,v_0)$ gives $b=q[w]+(t^2q[v_0]+2t\eta)$: a $b$-orthogonal sum of
$q|_{\ker f}$ and a binary form of determinant $-1$, whose inertia is $(1,1)$.
\end{proof}

\subsection{The certificates at \texorpdfstring{$T=\frac12\log5$}{T = (1/2) log 5}}\label{subsec:encllog5}
Use the Legendre basis and the band projection of Lemma~\ref{lem:KyFan} in each real parity space. Here the lower form is a frequency truncation of the full symbol $\sT$, and the operator subtracted from a
multiple of the identity is indefinite.

\emph{The lower form.} Fix $\tau_N>0$ and a number $s_\infty\le m(\tau_N)-\Aconst$. Since $m$ is increasing
with $m(0)=C_0$ and the comb is at least $-\Aconst$ (Lemma~\ref{lem:mbounds}(i),(iii)), the even function
$\underline s$ equal to $\sT$ on $|\tau|\le\tau_N$ and to $s_\infty$ on $|\tau|>\tau_N$ satisfies
$\underline s\le\sT$. Write $g=s_\infty-\underline s=(s_\infty-\sT)\mathbf 1_{|\tau|<\tau_N}=g_+-g_-$ with
$g_\pm\ge0$, and $\mathcal G_{(\pm)}=Q_Tg_{(\pm)}(D)Q_T$. Then
\[
\mathfrak a_T(F)\ \ge\ s_\infty\|F\|^2-\langle\mathcal GF,F\rangle,\qquad
0\le\mathcal G_\pm\le\omega_\pm S_{c_*},\quad c_*=T\tau_N,
\]
with $\omega_+=s_\infty-C_0+\Aconst$ and $\omega_-=m(\tau_N)+\Aconst-s_\infty$. For $\mu\in\R$ and
$c\in\{0,\frac12,-\frac12\}$ put, on $L^2_{\mathrm{par}}(\IT)\oplus\R$,
\[
\mathfrak b_{\mu,c}(F,\xi)=(s_\infty-\mu)\|F\|^2-\langle\mathcal GF,F\rangle+2\xi\,\ell_{\mathrm{par}}(F)+c\,\xi^2 .
\]

\begin{proposition}[the lower certificate]\label{prop:lowercert}
Let $k_1<k_2$, let $E$ be the span of the $e_j$, $j<k_1$, of parity $\mathrm{par}$, $Y$ its orthogonal
complement in $L^2_{\mathrm{par}}(\IT)$, and $P_Y$ the orthogonal projection onto $Y$.
\emph{(a)} On $Y$, $\mathfrak b_{\mu,c}\ge q_Y:=s_\infty-\mu-\omega_+\Lambda_{k_1}(c_*)$.
\emph{(b)} Suppose $q_Y>0$, and let $\mathsf R$ be the symmetric matrix on $E\oplus\R$ with
\[
\begin{gathered}
\mathsf R_{ij}=(s_\infty-\mu)\delta_{ij}-\langle\mathcal Ge_i,e_j\rangle
-\frac{\langle P_Y\mathcal Ge_i,P_Y\mathcal Ge_j\rangle}{q_Y},\\
\mathsf R_{i\xi}=\langle e_i,\chi_{\mathrm{par}}\rangle+\frac{\langle P_Y\mathcal Ge_i,P_Y\chi_{\mathrm{par}}\rangle}{q_Y},\qquad
\mathsf R_{\xi\xi}=c-\frac{\|P_Y\chi_{\mathrm{par}}\|^2}{q_Y}.
\end{gathered}
\]
Then $n_-(\mathfrak b_{\mu,c})\le n_-(\mathsf R)$, and the same holds for any symmetric $\mathsf R'\le\mathsf R$.
\emph{(c)} If $n_-(\mathsf R')\le1$ with $c=0$, then $\mathfrak a_T\ge\mu$ on the unit sphere of
$\PT^{\mathrm{par}}$. If $n_-(\mathsf R')=0$ with $c=\frac12$ and $\mathrm{par}=\mathrm{odd}$, or $n_-(\mathsf R')\le1$ with
$c=-\frac12$ and $\mathrm{par}=\mathrm{ev}$, then $W_T\ge\mu\|\cdot\|^2$ on $\mathcal D^{\mathrm{par}}_T$.
\end{proposition}
\begin{proof}\sloppy
(a) $-P_Y\mathcal GP_Y\ge-P_Y\mathcal G_+P_Y$ and
$\|P_Y\mathcal G_+P_Y\|\le\omega_+\operatorname{tr}(P^\perp S_{c_*}P^\perp)=\omega_+\Lambda_{k_1}(c_*)$.
(b) For $i<k_1$ and $y\in Y$, $\mathfrak b_{\mu,c}(e_i,y)=-\langle P_Y\mathcal Ge_i,y\rangle$ and the border pairs
with $y$ through $+\langle P_Y\chi_{\mathrm{par}},y\rangle$. Completing the square in the $Y$ variable gives $\mathfrak b_{\mu,c}\cong\mathfrak s\oplus\mathfrak b_{\mu,c}|_Y$ with
$\mathfrak b_{\mu,c}|_Y\ge q_Y>0$, so $n_-(\mathfrak b_{\mu,c})=n_-(\mathfrak s)$, and the subtracted term of the Schur
complement $\mathfrak s$ is at most $q_Y^{-1}$ times the Gram matrix of the $Y$-columns
$-P_Y\mathcal Ge_i$ and $P_Y\chi_{\mathrm{par}}$; hence $\mathfrak s\ge\mathsf R\ge\mathsf R'$.
(c) Put $\mathfrak q_\mu(F)=\mathfrak b_{\mu,0}(F,0)$, so that $\mathfrak a_T-\mu\|\cdot\|^2\ge\mathfrak q_\mu$. For
$c=0$, Lemma~\ref{lem:border} gives $n_-(\mathfrak q_\mu|_{\ker\ell_{\mathrm{par}}})=n_-(\mathfrak b_{\mu,0})-1=0$. For
$c=\frac12$, $\mathfrak b_{\mu,1/2}=\mathfrak q_\mu-2\ell_{\mathrm{par}}^2+\frac12(\xi+2\ell_{\mathrm{par}})^2$, so
$\mathfrak q_\mu-2\ell_{\mathrm{odd}}^2\ge0$, and $W_T=\mathfrak a_T-2\ell_{\mathrm{odd}}^2$ on odd $F$
(Theorem~\ref{thm:parityschur}(o)). For $c=-\frac12$,
$\mathfrak b_{\mu,-1/2}=\mathfrak q_\mu+2\ell_{\mathrm{par}}^2-\frac12(\xi-2\ell_{\mathrm{par}})^2$; the substitution
$\eta=\xi-2\ell_{\mathrm{par}}(F)$ is a congruence, so
$n_-(\mathfrak q_\mu+2\ell_{\mathrm{ev}}^2)=n_-(\mathfrak b_{\mu,-1/2})-1=0$, and $W_T=\mathfrak a_T+2\ell_{\mathrm{ev}}^2$ on even
$F$. The statements pass from $L^2_{\mathrm{par}}(\IT)$ to the form domain because $\mathfrak q_\mu$ is bounded.
\end{proof}

\emph{The computation.} For $i-l$ even,
\[
\langle\mathcal Ge_i,e_l\rangle=\frac{2T}\pi\sqrt{(2i+1)(2l+1)}\,(-1)^{(l-i)/2}\int_0^{\tau_N}g(\tau)j_i(\tau T)j_l(\tau T)\,d\tau
\]
(proof of Lemma~\ref{lem:KyFan}), and the entries with $i-l$ odd vanish. The integral is computed by composite
Gauss--Legendre quadrature with $46$ nodes on panels of length $\frac12$, with the remainder bound for
functions analytic in a Bernstein ellipse \cite[Thm.~19.3]{Trefethen2013}, $\rho=3.2$. The ellipse lies
within distance $0.439<\frac12$ of the panel, while the poles of $\psi(\frac14\pm\frac{i\tau}2)$ lie at
$|\Iml\tau|\ge\frac12$. The modulus on the ellipse is bounded with $|j_n(z)|\le e^{|\Iml z|}$,
$|\cos(z\log n)|\le\cosh(\Iml z\log n)$ and, for $\psi$, the recurrence $\psi(w)=\psi(w+8)-\sum_{k<8}(w+k)^{-1}$
together with Binet's bound $|\psi(u)|\le\log|u|+\frac\pi2+\frac1{2\Rel u}+\frac1{12(\Rel u)^2}$ for
$\Rel u\ge8$. The resulting remainder, $1.5\cdot10^{-43}$ per entry (below $2\cdot10^{-42}$ under either
point-count convention for the rule), is added as a ball to every entry. The inner products over $Y$ are
exact sums over $k_1\le l<k_2$ of the given parity plus tail boxes: by Cauchy--Schwarz for the positive
operators $\mathcal G_\pm$, $\|P^\perp_{k_2}\mathcal G_\pm e\|^2\le\|P^\perp_{k_2}\mathcal G_\pm P^\perp_{k_2}\|
\langle\mathcal G_\pm e,e\rangle\le\omega_\pm^2\Lambda_{k_2}(c_*)$, so
$\|P^\perp_{k_2}\mathcal Ge\|\le(\omega_++\omega_-)\Lambda_{k_2}(c_*)^{1/2}$, and
$\|P^\perp_{k_2}\chi_{\mathrm{par}}\|^2=\sinh T\pm T-\sum_{l<k_2}\langle e_l,\chi_{\mathrm{par}}\rangle^2$ with
$\langle e_l,\chi_{\mathrm{par}}\rangle=\sqrt{(2l+1)/2T}\,2T\,i_l(T/2)$. The tails $\Lambda_k$ are bounded by
$\Lambda^*_k$ of Corollary~\ref{cor:P1}. Everything is midpoint--radius ball arithmetic at $640$ bits
(Arb \cite{Johansson2017}, through \texttt{python-flint}), including the digamma and Bessel values and the
Gauss--Legendre nodes and weights; every Bessel ball is asserted to have radius below $10^{-60}$, and
comparisons are ball predicates.

For the inertia test, replacing $\|P_Y\chi_{\rm par}\|^2/q_Y$ by its
upper bound $\delta'$ gives $\mathsf R'\le\mathsf R$. Interval Cholesky
proves $n_-(\mathsf R')=0$ in the positive cases. Otherwise choose a
Householder matrix $H$ from the midpoint of the $\xi$ column and certify
$\|H^tH-I\|_\infty<1$. Congruence by $\operatorname{diag}(H,1)$ gives
a block $B$ on $(e_0,\xi)$ with certified $\det B<0$ and positive Schur
complement by interval Cholesky. Inertia additivity gives one negative
eigenvalue \cite{Haynsworth1968}.

\emph{Data.} $T=\frac12\log5$, $\tau_N=240$, $c_*=193.13$, $k_1=280$, $k_2=330$ ($\dim E=140$ in each parity).
From $m(\tau_N)=3.64276\ldots$: $s_\infty=0.700787608\ldots$, $\omega_+=9.0150$, $\omega_-=5.8839$,
$\Lambda^*_{k_1}(c_*)\le4.5\cdot10^{-20}$, $\Lambda^*_{k_2}(c_*)\le3.9\cdot10^{-70}$, $q_Y\ge0.70078$, tail
box $(\omega_++\omega_-)\Lambda^*_{k_2}(c_*)^{1/2}\le3.0\cdot10^{-34}$; the entry radii of $\mathsf R$ are at
most $4.4\cdot10^{-41}$. The verified cases are:
\begin{center}\small
\begin{tabular}{@{}llccc@{}}
\toprule
form & sector & certified $\mu$ & minimal pivot at that $\mu$ & fails at\\
\midrule
$\mathfrak a_T|\PT$ ($c=0$) & even & $2.5\cdot10^{-12}$ & $8.47\cdot10^{-3}$ & $3\cdot10^{-12}$\\
$\mathfrak a_T|\PT$ ($c=0$) & odd & $3.0\cdot10^{-10}$ & $2.03\cdot10^{-2}$ & $4\cdot10^{-10}$\\
$W_T$ ($c=-\frac12$) & even & $10^{-18}$ & $4.47\cdot10^{-4}$ & $10^{-16}$\\
$W_T$ ($c=\frac12$) & odd & $10^{-16}$ & $8.14\cdot10^{-14}$ & $10^{-14}$\\
\bottomrule
\end{tabular}
\end{center}
\subsection{The margin decays}\label{subsec:margin}
An upper bound related to Zhu's conjectural law $-\log\lamW\asymp e^{2T}$ \cite[Conj.~12.1]{Zhu2026}, namely that
the margin tends to $0$ at an explicit rate, can be proved from three facts about the zeros of $\zeta$.
These are also the zero inputs of the subsequent window computation cited in
\S\ref{subsec:power-window}.
\begin{itemize}
\item[(Z1)] $\rho_1=\frac12+i\gamma_1$, $\gamma_1=14.134725141734\ldots$, is a simple zero on the critical line,
and $\zeta$ has no non-trivial zero $\rho\ne\rho_1,\bar\rho_1$ with $|\Iml\rho|\le21$
\cite{RosserYoheSchoenfeld1968,PlattTrudgian2021}. Only $\gamma_1\in[14.13,14.14]$ and the gap $(\gamma_1,21]$
are used.
\item[(Z2)] Every non-trivial zero $\rho$ with $|\Iml\rho|\le H_{\mathrm{PT}}:=3\cdot10^{12}$ lies on the
critical line \cite{PlattTrudgian2021}.
\item[(Z3)] For $t\ge2$, $\big|N(t)-\frac t{2\pi}\log\frac t{2\pi e}-\frac78\big|\le0.137\log t+0.443\log\log t+1.588$,
where $N(t)$ counts the zeros with $0<\Iml\rho\le t$ with multiplicity \cite[Thm.~19]{Rosser1941}.
\end{itemize}
Besides these, we use that the non-trivial zeros lie in $0<\Rel s<1$, and Lemma~\ref{lem:roughEF}.
\begin{lemma}[a crude zero count]\label{lem:Ncount}
$N(t)\le\frac t{2\pi}\log t$ for $t\ge21$.
\end{lemma}
\begin{proof}
$\frac t{2\pi}\log\frac t{2\pi e}=\frac t{2\pi}\log t-\frac{1+\log2\pi}{2\pi}t$ and
$\frac{1+\log2\pi}{2\pi}=0.45166\ldots$, so by (Z3)
$\frac t{2\pi}\log t-N(t)\ge0.45166\,t-\frac78-1.588-0.137\log t-0.443\log\log t$, which is $6.11\ldots$ at
$t=21$ and has derivative $0.45166-0.137/t-0.443/(t\log t)>0$ for $t\ge21$.
\end{proof}

\begin{lemma}[zero sums by counting]\label{lem:zerosum}
For $p\ge2$ and $H\ge21$, with the sum over the ordinates $\gamma=\Iml\rho>0$ of the non-trivial zeros
counted with multiplicity,
\[
\Sigma_p(H):=\sum_{\gamma>H}\gamma^{-p}\ \le\ \frac p{2\pi}\,H^{1-p}\Big[\frac{\log H}{p-1}+\frac1{(p-1)^2}\Big].
\]
\end{lemma}
\begin{proof}
$\sum_{\gamma>H}\gamma^{-p}=\int_{H^+}^\infty t^{-p}dN(t)=\big[t^{-p}N(t)\big]_H^\infty+p\int_H^\infty t^{-p-1}N(t)\,dt$;
the term at $\infty$ vanishes and the one at $H$ is $-H^{-p}N(H)\le0$ by Lemma~\ref{lem:Ncount}, which
also bounds the integral by $\frac p{2\pi}\int_H^\infty t^{-p}\log t\,dt=\frac p{2\pi}H^{1-p}\big[\frac{\log H}{p-1}+\frac1{(p-1)^2}\big]$.
\end{proof}

\begin{lemma}[the test family]\label{lem:testfam}
Fix $T>0$ and an integer $k\ge5$; let $b_k=\frac k2\mathbf 1_{[-1/k,1/k]}$ and $\omega_k=b_k^{*k}$ (the
$k$-fold convolution).
\emph{(i)} $\omega_k\ge0$, $\int\omega_k=1$, $\supp\omega_k\subseteq[-1,1]$, $\omega_k\in C^{k-2}(\R)$ with
$\omega_k^{(j)}(\pm1)=0$ for $j\le k-2$, and $\|\omega_k\|_2^2\ge\frac12$.
\emph{(ii)} $\widehat\omega_k(x)=\big(\sin(x/k)/(x/k)\big)^k$; $|\widehat\omega_k(x)|\le\min\{1,(k/|x|)^k\}$ for real
$x$, and $|\widehat\omega_k(z)|\le(k/|\Rel z|)^ke^{|\Iml z|}$ for complex $z$.
\emph{(iii)} Put $G(t)=\omega_k(t/T)$, $\varpi(c)=(c^2-\frac14)(c^2+\gamma_1^2)$, $\partial=d/dt$, and
\[
F:=\varpi(\partial)G=G''''+(\gamma_1^2-\tfrac14)G''-\tfrac{\gamma_1^2}4G .
\]
Then $F$ is real, supported in $[-T,T]$, bounded, and $F\in\PT$: four integrations by parts, the
boundary terms carrying $G^{(j)}(\pm T)=0$ for $j\le3$, give $\int Fe^{ct}dt=\varpi(c)\int Ge^{ct}dt$
for every $c\in\C$, and $\varpi(\pm\frac12)=0$.
\emph{(iv)} $\Fh(\tau)=(\tau^2+\frac14)(\tau^2-\gamma_1^2)\,T\,\widehat\omega_k(T\tau)$,
$\Phi_F(s)=\varpi(s-\frac12)\Phi_G(s)$ and $\Phi_G(\frac12+\beta+i\gamma)=T\,\widehat\omega_k\big(T(\gamma-i\beta)\big)$;
in particular $\Phi_F(\rho_1)=\Phi_F(\bar\rho_1)=0$.
\emph{(v)} With $Y:=T\gamma_1/\sqrt2$ and $\varepsilon_k(Y):=\frac1\pi k^{2k}Y^{1-2k}/(2k-1)$,
$\|F\|^2\ge\frac{\gamma_1^4}{64}\,T\,\big(\frac12-\varepsilon_k(Y)\big)$ whenever $\varepsilon_k(Y)<\frac12$.
\end{lemma}
\begin{proof}
(i) $\omega_k$ is a spline of degree $k-1$ supported in $[-1,1]$ and vanishing there to order $k-1$;
$1=(\int_{-1}^1\omega_k)^2\le2\int\omega_k^2$. (ii) $\widehat b_k(x)=\sin(x/k)/(x/k)$; $|\sin y/y|\le\min(1,1/|y|)$
on $\R$, and $|\sin(a+ib)|\le\cosh b\le e^{|b|}$. (iii) $G\in C^3$ with $G'''$ absolutely continuous
and $G''''\in L^\infty$ ($k=5$), and $G^{(j)}(\pm T)=0$ for $j\le3$; $\int G''e^{ct}=c^2\int Ge^{ct}$.
(iv) is (iii) with $c=i\tau$, resp.\ $c=s-\frac12$, and $\varpi(\pm i\gamma_1)=0$. (v) On
$|\tau|\le\gamma_1/\sqrt2$ one has $(\tau^2+\frac14)^2\ge\frac1{16}$ and $(\tau^2-\gamma_1^2)^2\ge\gamma_1^4/4$,
so $\|F\|^2=\frac1{2\pi}\int|\Fh|^2\ge\frac{\gamma_1^4}{64}\,T\,\frac1{2\pi}\int_{|x|\le Y}|\widehat\omega_k|^2$, and
$\frac1{2\pi}\int_{|x|\le Y}|\widehat\omega_k|^2\ge\|\omega_k\|_2^2-\frac1\pi\int_Y^\infty(k/x)^{2k}dx\ge\frac12-\varepsilon_k(Y)$.
\end{proof}

\begin{theorem}[the margin decays]\label{thm:margin}
Let $T>0$, $k\ge5$ an integer, $p:=2k-8$, $Y:=T\gamma_1/\sqrt2$, and suppose $\varepsilon_k(Y)<\frac12$. Then
\begin{equation}\label{eq:Bmargin}
\lamW\ \le\ \lamP\ \le\ \lamA\ \le\ \mathcal B(T,k):=
\frac{2\cdot1.0012\cdot T^2(k/T)^{2k}\big[\Sigma_p(21)+e^T\,\Sigma_p(H_{\mathrm{PT}})\big]}
{(14.13^4\,T/64)\,\big(\frac12-\varepsilon_k(Y)\big)},
\end{equation}
with $\Sigma_p(H)$ bounded by Lemma~\ref{lem:zerosum}. Consequently $\limsup_{T\to\infty}\lamA\le0$, and the
same for $\lamP$ and $\lamW$.
\end{theorem}
\begin{proof}
Take $F$ of Lemma~\ref{lem:testfam}. (1) $F\in\PT$, so $\lamA\le\mathfrak a_T(F)/\|F\|^2$, and
$\lamW\le\lamP\le\lamA$ by Theorem~\ref{thm:A} and Remark~\ref{rem:Kprim}.
(2) For $|\tau|\ge21$, $(\tau^2+\frac14)^2\le1.0012\,\tau^4$ and $(\tau^2-\gamma_1^2)^2\le\tau^4$, so
$|\Fh(\tau)|^2\le1.0012\,\tau^8T^2(k/(T|\tau|))^{2k}=O(|\tau|^{8-2k})$ with $8-2k\le-2$: $F$ is in the
form domain, which is hypothesis (a) of Lemma~\ref{lem:roughEF}; hypothesis (b) is verified in step (6).
Hence $\mathfrak a_T(F)=\sum_\rho\Phi_F(\rho)\Phi_F(1-\rho)$.
(3) $\rho_1,\bar\rho_1$ contribute $0$ (Lemma~\ref{lem:testfam}(iv)); by (Z1) no other zero has $|\Iml\rho|\le21$.
(4) A zero with $21<|\gamma|\le H_{\mathrm{PT}}$ lies on the line by (Z2) and contributes
$|\Fh(\gamma)|^2\le1.0012\,T^2(k/T)^{2k}|\gamma|^{8-2k}$.
(5) A zero $\rho=\frac12+\beta+i\gamma$ with $|\gamma|>H_{\mathrm{PT}}$ has $|\beta|<\frac12$;
$|\varpi(\beta+i\gamma)|\le(\gamma^2+\beta^2+\frac14)(\gamma^2+\beta^2+\gamma_1^2)\le(1+10^{-22})\gamma^4$, $\varpi$ is even,
and $|\Phi_G(\rho)|,|\Phi_G(1-\rho)|\le T(k/(T|\gamma|))^ke^{T|\beta|}$ by Lemma~\ref{lem:testfam}(ii),(iv);
so $|\Phi_F(\rho)\Phi_F(1-\rho)|\le1.0012\,T^2(k/T)^{2k}|\gamma|^{8-2k}e^T$.
(6) Summing (4)--(5) over $\pm\gamma$ with Lemma~\ref{lem:zerosum} ($p=2k-8\ge2$) gives
$\sum_\rho|\Phi_F(\rho)\Phi_F(1-\rho)|\le2\cdot1.0012\,T^2(k/T)^{2k}\big[\Sigma_p(21)+e^T\Sigma_p(H_{\mathrm{PT}})\big]<\infty$,
which is hypothesis (b) of Lemma~\ref{lem:roughEF} and bounds $\mathfrak a_T(F)$.
(7) Divide by Lemma~\ref{lem:testfam}(v) with $\gamma_1^4\ge14.13^4$. The limit statement is
Corollary~\ref{cor:margin1}.
\end{proof}

\begin{corollary}[closed form]\label{cor:margin1}
For every $T\ge1$,
\begin{equation}\label{eq:closedmargin}
\lamW\ \le\ \lamP\ \le\ \lamA\ \le\ 4\cdot10^9\,T\,e^{-15.449\,T},\qquad 15.449<\tfrac{42}e=15.4509\ldots\footnotemark
\end{equation}
\footnotetext{The exponent $\frac{42}{e}$ belongs to the test family of Lemma~\ref{lem:testfam}: one factor removes only the first zero, so the cutoff is $21$, just below the second ordinate $\gamma_2=21.022\ldots$, and $\frac{42}e$ comes from optimizing $k$ in $(k/21T)^{2k}$. It is a property of that family, not a sharp rate for $\lamA$.}
\end{corollary}
\begin{proof}
Take $k:=\lceil21T/e\rceil\ge8$ and write $k=21T/e+\theta$, $\theta\in[0,1)$. Then
$(k/21T)^{2k}=\exp\big(2k[\log(1+\theta e/21T)-1]\big)\le\exp(-42T/e+2\theta^2e/21T)\le e^{-15.4494\,T+0.259}$.
Next $(k/T)^{2k}\Sigma_p(21)\le21^9(k/21T)^{2k}\cdot\frac1{2\pi}\frac{2k-8}{2k-9}\big[\log21+\frac1{2k-9}\big]\le21^9\cdot0.5798\,(k/21T)^{2k}$
for $k\ge8$. The bound for the $H_{\mathrm{PT}}$-term is at most $10\,e^T(21/H_{\mathrm{PT}})^{2k-9}$ times that for the first, with
$21/H_{\mathrm{PT}}=7\cdot10^{-12}$ and $2k-9\ge15.449\,T-9$, hence less than $10^{-70}$ of it for $T\ge1$.
Finally $\varepsilon_k(Y)\le0.01567$ for all $T\ge1$: for fixed $k$, $\varepsilon_k$ decreases in $T$, so it
suffices to evaluate it at the left end of each $k$-range, namely $T=1$ for $k=8$ ($\varepsilon_8=0.0060$)
and $T=(k-1)e/21$ for $k\ge9$, where the largest value is $\varepsilon_9=0.015661\ldots$ at $T=8e/21$, the values
decreasing thereafter. So $\frac12-\varepsilon_k\ge0.484$, and
$\mathcal B(T,k)\le\frac{2\cdot1.0012\cdot0.5798\cdot21^9e^{0.259}}{(14.13^4/64)\cdot0.484}\,T\,e^{-15.4494\,T}(1+10^{-70})
\le3.97\cdot10^9\,T\,e^{-15.4494\,T}$.
\end{proof}

For the complete minimum Mori reports the upper bound
$7.58\cdot10^{24}\mu^5(\log\mu)^3e^{-4\pi\mu}$ for $\mu=e^{2T}\ge50$,
using finite-height zero data \cite{MoriUpper}. Corollary~\ref{cor:margin1}
also bounds the primitive minimum.

By monotonicity (Lemma~\ref{lem:monoT}), the three minima have limits
$L^{\rm W}\le L^{\rm prim}\le L^{\rm A}$ in $[-\infty,0]$, since
\eqref{eq:closedmargin} tends to zero. In particular \eqref{eq:target}
is equivalent to $L^{\rm A}=0$.

\subsection{A reported larger window and its inputs}\label{subsec:power-window}
The extended window certificate gives, using (Z1)--(Z3) of \S\ref{subsec:margin},
\begin{equation}\label{eq:power-window}
 W_T(F)\ge e^{-120110026807}\|F\|_2^2,
 \qquad 0<T\le\frac{14011}{5000}=2.8022,
\end{equation}
on the complete form domain with $\supp F\subset[-T,T]$.
The polar term vanishes on $\PT$, giving the same bound for $\lamA$.
The compact spectral filter isolates the first zero pair using (Z1),
uses (Z2) to height $3\cdot10^{12}$, and bounds the tail by (Z3).
With $U=57/20$, $H=2673507086718$, $R=2014985964$, $m=9115062736$,
$b=13$, and filter order $120122483651$, the $17$ checks include the
support, $2048$ continuous Schur rows, the mask and both correction tails.
They were reproduced in $192$-bit Arb arithmetic with python-flint~0.9.0; the program is not
reproduced in this paper, and no other statement of this paper depends on \eqref{eq:power-window}.
Extension by zero gives the smaller radii.

\section{Reduction to piecewise exponentials}\label{sec:coer}

We impose both Tate conditions on each of $k$ blocks. Coercivity on
their common kernel leaves $2k-2$ piecewise exponential directions.

\begin{definition}\label{def:coer-V}
Let $k\ge1$ and let $B_1,\dots,B_k$ be the $k$ consecutive open intervals of length $2r$,
$r:=T/k$, that partition $\IT$ up to endpoints, $B_j=(-T+2(j-1)r,\,-T+2jr)$. Put
\[
\mathcal V_k:=\Big\{F\in L^2(\IT):\ \int_{B_j}F(t)e^{t/2}dt=\int_{B_j}F(t)e^{-t/2}dt=0\ \ (1\le j\le k)\Big\},
\qquad E_k:=\PT\ominus\mathcal V_k ,
\]
and $e_j^\pm:=e^{\pm t/2}\mathbf 1_{B_j}$. (Lemma~\ref{lem:coer-decomp} holds for any partition into
intervals; equal blocks are used from Lemma~\ref{lem:coer-lowpass} on.)
\end{definition}

\begin{lemma}[the decomposition $\PT=\mathcal V_k\oplus E_k$]\label{lem:coer-decomp}
\emph{(i)} $\mathcal V_k$ is a closed subspace of $\PT$. \emph{(ii)} The orthogonal complement of
$\mathcal V_k$ in $L^2(\IT)$ is $\operatorname{span}\{e_j^\pm\}$, of dimension $2k$.
\emph{(iii)} $E_k=\operatorname{span}\{e_j^\pm\}\cap\PT=\{\sum_jc_j^+e_j^++c_j^-e_j^-:\ A=B=0\}$ has
dimension $2k-2$, and $\PT=\mathcal V_k\oplus E_k$ (orthogonal in $L^2(\IT)$).
\emph{(iv)} $E_k\subset\mathcal D_T$, and $\PT\cap\mathcal D_T=(\mathcal V_k\cap\mathcal D_T)\oplus E_k$.
\end{lemma}
\begin{proof}
(i) The $2k$ conditions are continuous linear functionals; summing them over $j$ gives $A(F)=B(F)=0$.
(ii) $\mathcal V_k$ is the joint kernel of the functionals $F\mapsto\langle F,e_j^\pm\rangle$, and the $2k$
functions $e_j^\pm$ are linearly independent (disjoint supports; $e^{t/2},e^{-t/2}$ independent on an
interval). (iii) $E_k=\mathcal V_k^\perp\cap\PT$ since $\mathcal V_k\subset\PT$. The two functionals
$A,B$ restricted to $\operatorname{span}\{e_j^\pm\}$ are independent: on the two-dimensional space
$\operatorname{span}\{e_1^+,e_1^-\}$ their matrix is
$\begin{psmallmatrix}\int_{B_1}e^t&|B_1|\\|B_1|&\int_{B_1}e^{-t}\end{psmallmatrix}$, of determinant
$\int_{B_1}e^t\int_{B_1}e^{-t}-|B_1|^2>0$ (Cauchy--Schwarz, strict since $e^t$ is not constant). Hence
$\dim E_k=2k-2$. If $F\in\PT$ and $F=v+w$ with $v\in\mathcal V_k$, $w\in\mathcal V_k^\perp$, then
$w=F-v\in\PT$, so $w\in E_k$. (iv) $\widehat{e_j^\pm}(s)=\big[e^{(\pm\frac12+is)t}/(\pm\frac12+is)\big]_{\partial B_j}$,
so $|\widehat{e_j^\pm}(s)|\le2e^{T/2}(\frac14+s^2)^{-1/2}$ and $\int\log(2+|s|)|\widehat{e_j^\pm}|^2<\infty$.
Since $\mathcal D_T$ is a linear space containing $E_k$, the $\mathcal V_k$-component of an $F\in\PT\cap\mathcal D_T$
lies in $\mathcal D_T$.
\end{proof}

For a window $r>0$ we write $\lambda_{\mathrm A}(r)$ for the bottom of $\mathfrak a_r$ on the unit sphere of
$\mathcal P_r\cap\mathcal D_r$; it is attained (Lemma~\ref{lem:compact}) and continuous in $r$ (Lemma~\ref{lem:contA}).
For $2r\le\log2$ the contact set $\mathcal N_r$ is empty, so $\mathfrak a_r=\mT$ on $L^2(I_r)$ and $\lambda_{\mathrm A}(r)$ is a
purely archimedean quantity. We use $|\cF(u)|\le\cF(0)=\|F\|^2$ and $\cF(u)=0$ for $|u|\ge2T$.

\begin{lemma}[the reduced kernel]\label{lem:coer-K1}
Put $K_1(u):=K(u)-2e^{-u/2}=\dfrac{2e^{-5u/2}}{1-e^{-2u}}=2\sum_{j\ge1}e^{-(2j+\frac12)u}$ for $u>0$, with $K$ as in
Lemma~\ref{lem:K}. Then $0<K_1(u)\le u^{-1}$ for all $u>0$, and $K_1(u)=u^{-1}-\frac32+O(u)$ as $u\downarrow0$. Let
$F,G\in L^2(\IT)$ be real, $\supp F\subset\overline B$, $\supp G\subset\overline{B'}$, with $B,B'$ disjoint intervals and $B'$
to the right of $B$. Then
\begin{equation}\label{eq:coer-cross}
\mT(F+G)=\mT(F)+\mT(G)-\iint F(x)G(y)K(y-x)\,dx\,dy ,
\end{equation}
and if $\int F(x)e^{x/2}dx=0$ or $\int G(y)e^{-y/2}dy=0$, the kernel $K$ in \eqref{eq:coer-cross} may be replaced by $K_1$.
\end{lemma}
\begin{proof}
The three expressions for $K_1$ agree by Lemma~\ref{lem:K} ($K=2e^{-u/2}\sum_{j\ge0}e^{-2ju}$). $K_1\le u^{-1}$ is
$e^{5u/2}-e^{u/2}\ge2u$: the function $g(u)=e^{5u/2}-e^{u/2}-2u$ has $g(0)=g'(0)=0$ and
$g''=\frac{25}4e^{5u/2}-\frac14e^{u/2}>0$. The expansion follows from $\frac1{2\sinh(u/2)}=u^{-1}-\frac u{24}+O(u^3)$,
$\frac1{2\cosh(u/2)}=\frac12+O(u^2)$ and $2e^{-u/2}=2-u+O(u^2)$. For \eqref{eq:coer-cross}, apply \eqref{eq:Aform} with
$\NT=\emptyset$ to $F+G$: $c_{F+G}(u)=\cF(u)+c_G(u)+\int F(t)G(t-u)dt+\int G(t)F(t-u)dt$; for $u>0$ the first cross
integral vanishes ($t\in B$, $t-u\in B'$ is impossible), while the second equals $\iint F(x)G(y)\mathbf 1_{y-x=u}$; integrating
against $K(u)du$ over $u>0$ gives \eqref{eq:coer-cross}, the term $\|F+G\|^2=\|F\|^2+\|G\|^2$ having no cross part.
Finally $2e^{-(y-x)/2}=2e^{-y/2}e^{x/2}$, so $\iint F(x)G(y)\,2e^{-(y-x)/2}dxdy=2\int Fe^{x/2}\int Ge^{-y/2}$, which vanishes
under either Tate condition.
\end{proof}
The cross term in \eqref{eq:coer-cross} carries the factor $1$ for the unordered pair $\{B,B'\}$; the polarized form is
$\mT(F,G)=-\frac12\iint F(x)G(y)K(y-x)\,dxdy$.

\begin{lemma}[a block sector is a short window]\label{lem:coer-sector}
Let $B\subset\IT$ be an interval of length $2r$ with $2r\le\log2$. Then
\[
\inf\Big\{\mathfrak a_T(F):\ F\in\mathcal D_T,\ \supp F\subset\overline B,\ \int_BFe^{\pm t/2}=0,\ \|F\|=1\Big\}=\lambda_{\mathrm A}(r),
\]
independently of $T$ and of the position of $B$; the infimum is attained. In particular, for $F\in\mathcal V_k$ with
$r=T/k$, $2r\le\log2$, and $F_j:=F\mathbf 1_{B_j}$,
\begin{equation}\label{eq:coer-diag}
\sum_{j=1}^k\mT(F_j)\ \ge\ \lambda_{\mathrm A}(r)\,\|F\|^2 .
\end{equation}
\end{lemma}
\begin{proof}
Let $B=I_r+a$. For $F$ supported in $\overline B$ and $F_a(t):=F(t+a)$, supported in $\overline{I_r}$, one has $c_{F_a}=\cF$,
hence $\mT(F_a)=\mT(F)$ by \eqref{eq:Aform}, and $\int F_ae^{\pm t/2}=e^{\mp a/2}\int Fe^{\pm t/2}$, so the two Tate
conditions on $B$ are those of $\mathcal P_r$ for $F_a$. Since $\cF(\log n)=0$ for $\log n\ge2r$ and $\log2\ge2r$, also
$\PiT(F)=0$, so $\mathfrak a_T(F)=\mT(F)=\mathfrak a_r(F_a)$ and the infimum is $\lambda_{\mathrm A}(r)$, attained since
$\lambda_{\mathrm A}(r)$ is. \eqref{eq:coer-diag} is the sum over the blocks, each $F_j$ satisfying the Tate conditions of
$B_j$, together with $\sum_j\|F_j\|^2=\|F\|^2$.
\end{proof}


\subsection{A lower bound for the short window}\label{subsec:coer-short}

\begin{lemma}\label{lem:coer-K}
$K(u)\ge u^{-1}$ for $0<u\le2$.
\end{lemma}
\begin{proof}
By Lemma~\ref{lem:K}, $K(u)-u^{-1}=\frac{u-2\sinh(u/2)}{2u\sinh(u/2)}+\frac1{2\cosh(u/2)}$. From $\tanh x\le x$,
$2\sinh(u/2)\le u\cosh(u/2)$, so $u-2\sinh(u/2)\ge-u(\cosh(u/2)-1)=-2u\sinh^2(u/4)$, and the first term is
$\ge-\frac{\sinh^2(u/4)}{\sinh(u/2)}=-\frac12\tanh\frac u4\ge-\frac u8$. Hence $K(u)-u^{-1}\ge\frac1{2\cosh(u/2)}-\frac u8\ge0$
whenever $u\cosh(u/2)\le4$, which holds for $u\le2$ ($2\cosh1=3.09$).
\end{proof}

\begin{lemma}[the tail]\label{lem:coer-tail}
For $r>0$, $\int_{2r}^\infty K(u)\,du\ge\log\frac2r+\frac\pi2-r$.
\end{lemma}
\begin{proof}
By Theorem~\ref{thm:S}(iii) the integral is $\log\coth\frac r2+\pi-2\arctan e^r$; $\coth x\ge x^{-1}$ bounds the first
term, and $\frac{d}{dr}(\pi-2\arctan e^r)=-\frac1{\cosh r}\ge-1$ with value $\frac\pi2$ at $r=0$ bounds the second.
\end{proof}

\begin{lemma}[the truncated logarithm]\label{lem:coer-w}
Put $\mathsf w(s):=\int_0^2\frac{1-\cos sv}v\,dv=\log(2|s|)+\gamma-\operatorname{Ci}(2|s|)$ for $s\ne0$, $\mathsf w(0)=0$. Then $\mathsf w$ is
even, increasing in $|s|$, and for every $0<R\le\pi$ and all $s\in\R$, $\mathsf w(s)\ge\mathsf w(R)\min\{1,s^2/R^2\}$.
\end{lemma}
\begin{proof}
$\mathsf w'(s)=\int_0^2\sin(sv)\,dv=\frac{1-\cos2s}s\ge0$ for $s>0$. Next $\mathsf w(s)/s^2=\int_0^2v\,\frac{1-\cos sv}{(sv)^2}dv$, and
$x\mapsto\frac{1-\cos x}{x^2}=\frac12\big(\frac{\sin(x/2)}{x/2}\big)^2$ is decreasing on $(0,2\pi)$; for $0<s\le R\le\pi$ and $v\le2$
one has $sv\le2\pi$, so $\mathsf w(s)/s^2\ge\mathsf w(R)/R^2$. For $s\ge R$ use monotonicity.
\end{proof}

\begin{lemma}[a double Poincar\'e inequality]\label{lem:coer-poinc}
Let $h\in L^2(\R)$ with $\supp h\subset[-1,1]$ and $h''\in L^2(\R)$. Then $\|h''\|^2\ge\pi^2\|h'\|^2\ge\frac{\pi^4}4\|h\|^2$.
\end{lemma}
\begin{proof}
$h,h''\in L^2(\R)$ give $h\in H^2(\R)\subset C^1(\R)$, and the support condition gives $h(\pm1)=h'(\pm1)=0$. The second
inequality is the Dirichlet Poincar\'e inequality for $h\in H^1_0(-1,1)$ (first eigenvalue $(\pi/2)^2$). For the first,
$g:=h'\in H^1_0(-1,1)$ with $\int_{-1}^1g=h(1)-h(-1)=0$. Expand $g=\sum_{n\ge1}c_n\phi_n$ in the orthonormal Dirichlet
eigenbasis $\phi_n(x)=\sin\frac{n\pi(x+1)}2$ of $(-1,1)$, so that $\|g\|^2=\sum c_n^2$ and $\|g'\|^2=\sum(\frac{n\pi}2)^2c_n^2$. Since
$\int\phi_n=\frac4{n\pi}$ for odd $n$ and $0$ for even $n$, the constraint reads $c_1=-\sum_{n\ge3\ \mathrm{odd}}c_n/n$. The
even-indexed terms satisfy $(\frac{n\pi}2)^2\ge\pi^2$, and for the odd-indexed ones
\[
\sum_{n\ \mathrm{odd}}\Big(\frac{n^2\pi^2}4-\pi^2\Big)c_n^2=\frac{\pi^2}4\Big[\sum_{n\ge3\ \mathrm{odd}}(n^2-4)c_n^2-3c_1^2\Big]\ \ge\ 0,
\]
because by Cauchy--Schwarz $c_1^2\le\mathcal S\sum_{n\ge3\ \mathrm{odd}}(n^2-4)c_n^2$ with
$\mathcal S=\sum_{n\ge3\ \mathrm{odd}}\frac1{n^2(n^2-4)}=\frac14\big(\frac43-\frac{\pi^2}8\big)=0.0249<\frac13$ (partial fractions and
$\sum_{n\ \mathrm{odd}}n^{-2}=\pi^2/8$). Hence $\|g'\|^2\ge\pi^2\|g\|^2$.
\end{proof}

\begin{theorem}[the short window]\label{thm:coer-A}
Let $0<r\le1$ with $2r\le\log2$, and put $a:=r^2/4$. For every $R\in(0,\pi]$,
\begin{equation}\label{eq:coer-A}
\lambda_{\mathrm A}(r)\ \ge\ C_0+\log\frac2r+\frac\pi2-r+\mathsf w(R)\Big[1-\frac{16(R^2+a)^3}{27\pi^4R^2}\Big].
\end{equation}
In particular, with $R=2$ ($\mathsf w(2)=\log4+\gamma-\operatorname{Ci}(4)=2.1045$),
$\lambda_{\mathrm A}(r)\ge\log\frac1r-1.209-r$ for $0<r\le0.1$, so $\lambda_{\mathrm A}(r)\ge\log\frac1r-1.31$ on $(0,0.1]$.
\end{theorem}
\begin{proof}
Let $F\in\mathcal P_r\cap\mathcal D_r$, $\|F\|=1$. Since $\mathcal N_r=\emptyset$ and $\cF(u)=0$ for $u\ge2r$, \eqref{eq:Aform} gives
\[
\mathfrak a_r(F)=C_0+\int_{2r}^\infty K+\int_0^{2r}\big(1-\Rel\cF(u)\big)K(u)\,du .
\]
The tail is bounded by Lemma~\ref{lem:coer-tail}. Put $G(y):=r^{1/2}F(ry)$, so that $\supp G\subset[-1,1]$, $\|G\|=1$,
$\cF(u)=c_G(u/r)$, and $\int_0^{2r}(1-\Rel\cF)K=\int_0^2(1-\Rel c_G(v))\,rK(rv)\,dv\ge\int_0^2\frac{1-\Rel c_G(v)}v\,dv$, using
$1-\Rel c_G\ge0$ and Lemma~\ref{lem:coer-K} ($rv\le2r\le2$). By Plancherel, $1-\Rel c_G(v)=\frac1{2\pi}\int(1-\cos sv)|\widehat G(s)|^2ds$,
so by Tonelli the last integral equals $\frac1{2\pi}\int\mathsf w|\widehat G|^2$, and by Lemma~\ref{lem:coer-w}, for $0<R\le\pi$,
\[
\frac1{2\pi}\int\mathsf w|\widehat G|^2\ \ge\ \mathsf w(R)\Big[1-\frac1{2\pi}\int\Big(1-\frac{s^2}{R^2}\Big)^+|\widehat G(s)|^2ds\Big].
\]
The Tate conditions $\int Fe^{\pm t/2}=0$ read $\int Ge^{\pm ry/2}=0$, i.e.\ $\widehat G(\mp ir/2)=0$. $\widehat G$ is entire of
exponential type $\le1$ and in $L^2(\R)$, so $H(s):=\widehat G(s)/(s^2+a)$ is entire of exponential type $\le1$ and in $L^2(\R)$
($s^2+a\ge a>0$ on $\R$); by the Paley--Wiener theorem $H=\widehat h$ with $h\in L^2(\R)$, $\supp h\subset[-1,1]$, and
$G=-h''+ah$, so $h''\in L^2$. Then
$1=\|G\|^2=\frac1{2\pi}\int(s^2+a)^2|H|^2=\|h''\|^2+2a\|h'\|^2+a^2\|h\|^2\ge\|h''\|^2\ge\frac{\pi^4}4\|h\|^2$
(Lemma~\ref{lem:coer-poinc}). On the other hand $(1-\frac{s^2}{R^2})^+(s^2+a)^2\le M_R:=\frac{4(R^2+a)^3}{27R^2}$ (the
maximum of $(1-x/R^2)(x+a)^2$ on $0\le x\le R^2$ is at $x=(2R^2-a)/3$), so
$\frac1{2\pi}\int(1-s^2/R^2)^+|\widehat G|^2\le M_R\|h\|^2\le\frac{4M_R}{\pi^4}$. Collecting terms gives \eqref{eq:coer-A}. For
$R=2$ and $a\le\frac1{400}$ the bracket (decreasing in $a$) is $\ge0.90248$ and $\mathsf w(2)\cdot0.90248\ge1.89926$; with
$C_0+\log2+\frac\pi2=-3.108239\ldots\ge-3.10824$ this gives $\log\frac1r-1.20898-r\ge\log\frac1r-1.209-r$.
\end{proof}

Together with Corollary~\ref{cor:mono} and Lemma~\ref{lem:mbounds}(ii),
the lower estimate gives $\lambda_{\rm A}(r)=\log(1/r)+O(1)$ under the two Tate conditions.

Yoshida proves short-window coercivity after removing finitely many Fourier
modes \cite[Lemma~3, p.~290]{Yoshida1992}; Connes and Consani obtain a finite-codimension
inequality from the archimedean operator \cite[Thm.~3.6]{ConnesConsani2021}.
Here the two Tate conditions give the explicit rate $\log(1/r)$.
For numerical and finite-section comparisons see \cite{Kim2026,CvS2025}.


\subsection{Coercivity of the archimedean form on $\mathcal V_k$}\label{subsec:coer-B}
For $R>0$ let $P_R$ be the orthogonal projection of $L^2(\R)$ onto the band $\{|s|\le R\}$,
$\widehat{P_RF}=\Fh\mathbf 1_{|s|\le R}$.

\begin{lemma}[low-pass leakage of $\mathcal V_k$]\label{lem:coer-lowpass}
Let $F\in\mathcal V_k$ with equal blocks of half-length $r$, and $R>0$. Then
\[
\|P_RF\|^2\le\varepsilon_{\mathrm{lp}}(Rr)\,\|F\|^2,\qquad \varepsilon_{\mathrm{lp}}(\theta):=\frac{\cosh^2r}{6}\Big(\theta^2+\frac{r^2}4\Big)^2 .
\]
\end{lemma}
\begin{proof}
$G:=P_RF$ is entire. Fix a block $B$ with centre $c$ and put, on $B$,
$u(x):=G(c)\cosh\frac{x-c}2+2G'(c)\sinh\frac{x-c}2\in\operatorname{span}\{e^{x/2},e^{-x/2}\}$, so that
$u(c)=G(c)$, $u'(c)=G'(c)$, $u''=u/4$. Then $\varrho:=G-u$ satisfies $\varrho(c)=\varrho'(c)=0$ and
$\varrho''-\frac\varrho4=G''-\frac G4=:H$, whence $\varrho(x)=\int_c^x2\sinh\frac{x-t}2\,H(t)\,dt$. By Cauchy--Schwarz and
$2\sinh\frac z2\le z\cosh\frac z2\le z\cosh r$ ($0\le z\le r$),
$|\varrho(x)|^2\le\cosh^2r\,\frac{|x-c|^3}3\,\|H\|_{L^2(B)}^2$, so $\int_B\varrho^2\le\cosh^2r\,\frac{r^4}6\|H\|_{L^2(B)}^2$.
Since $F_B:=F\mathbf 1_B$ is orthogonal to $e^{\pm x/2}\mathbf 1_B$, $\langle F_B,G\rangle=\langle F_B,\varrho\rangle\le\|F_B\|\,\|\varrho\|_{L^2(B)}$.
Summing over the blocks and using Cauchy--Schwarz once more,
\begin{align*}
\|G\|^2=\langle F,G\rangle&\le\|F\|\Big(\sum_B\|\varrho\|_{L^2(B)}^2\Big)^{1/2}\le\|F\|\cosh r\,\frac{r^2}{\sqrt6}\,\|G''-\tfrac14G\|_{L^2(\R)}\\
&\le\|F\|\cosh r\,\frac{r^2}{\sqrt6}\Big(R^2+\frac14\Big)\|G\|,
\end{align*}
the first equality because $P_R$ is an orthogonal projection, the last inequality because
$\widehat G$ is supported in $|s|\le R$. Squaring gives the claim, as $r^4(R^2+\frac14)^2=((Rr)^2+\frac{r^2}4)^2$.
\end{proof}

\begin{theorem}[coercivity on $\mathcal V_k$]\label{thm:coer-B}
Let $k\ge1$, $r=T/k\le1$, $\kappa:=\cosh^2r$, $a:=r^2/4$, $\varepsilon_{\mathrm{lp}}$ as in
Lemma~\ref{lem:coer-lowpass}, and $\theta_*:=(\sqrt{6/\kappa}-a)^{1/2}$ (the root of $\varepsilon_{\mathrm{lp}}(\theta_*)=1$); $m_{\mathrm{lo}}$ is
the lower bound of Lemma~\ref{lem:mbounds}(ii).
Then for every $F\in\mathcal V_k\cap\mathcal D_T$,
\begin{equation}\label{eq:coer-B}
\mT(F)\ \ge\ \Big(\log\frac1r-C_B(r)\Big)\|F\|^2,\qquad
C_B(r):=\log\frac1r-C_0\,\varepsilon_{\mathrm{lp}}(r)-\frac{2\kappa}3\big[\Phi_r(\theta_*)-\Phi_r(r)\big],
\end{equation}
where $\Phi_r$ is the elementary primitive
\[
\Phi_r(\theta):=\Big(\frac{\theta^4}4+\frac{a\theta^2}2\Big)\big(\log\theta-\log(2\pi r)\big)-\frac{\theta^4}{16}-\frac{a\theta^2}4
-r\Big(\frac{\theta^3}3+a\theta\Big)-\frac{r^2\theta^2}4-\frac{ar^2}2\log\theta
\]
of $(\theta^3+a\theta)\,m_{\mathrm{lo}}(\theta/r)$. Moreover $C_B(r)\to\frac14-\log\frac{6^{1/4}}{2\pi}=1.6399$ as
$r\to0$, and $C_B(0.02)=1.6573$, $C_B(0.0619)=1.6957$, $C_B(0.1)=1.7331$ (evaluations of the closed form).
\end{theorem}
\begin{proof}
Let $\|F\|=1$ and $\mu:=\frac1{2\pi}|\Fh|^2ds$, a probability measure, so that $\mT(F)=\int m\,d\mu$ by
\eqref{eq:form} with $\NT=\emptyset$. By Lemma~\ref{lem:mbounds}(i) and \eqref{eq:m}, $m$ is even, real-analytic,
increasing on $[0,\infty)$ with $m(0)=C_0$; hence $m(s)=C_0+\int_0^{|s|}m'(\rho)d\rho$ with $m'\ge0$ and, by Tonelli,
\[
\int m\,d\mu=C_0+\int_0^\infty m'(\rho)\,\mu(|s|\ge\rho)\,d\rho\ \ge\ C_0+\int_0^{\rho_*}m'(\rho)\big(1-\varepsilon_{\mathrm{lp}}(\rho r)\big)d\rho,
\qquad\rho_*:=\theta_*/r,
\]
since $\mu(|s|\ge\rho)=1-\|P_{\rho}F\|^2+\mu(|s|=\rho)\ge1-\varepsilon_{\mathrm{lp}}(\rho r)$ by Lemma~\ref{lem:coer-lowpass} and the
integrand is nonnegative on $[0,\rho_*]$, where $\varepsilon_{\mathrm{lp}}(\rho r)\le1$. Integrating by parts on $[0,\rho_*]$, with
vanishing boundary terms ($1-\varepsilon_{\mathrm{lp}}(\theta_*)=0$ and $m(0)-C_0=0$),
\[
\int_0^{\rho_*}m'(\rho)(1-\varepsilon_{\mathrm{lp}}(\rho r))d\rho=\int_0^{\theta_*}\big(m(\theta/r)-C_0\big)\varepsilon_{\mathrm{lp}}'(\theta)\,d\theta,
\]
so that
\[
\int m\,d\mu\ \ge\ C_0\,\varepsilon_{\mathrm{lp}}(0)+\int_0^{\theta_*}m(\theta/r)\,\varepsilon_{\mathrm{lp}}'(\theta)\,d\theta ,
\]
using $\int_0^{\theta_*}\varepsilon_{\mathrm{lp}}'=1-\varepsilon_{\mathrm{lp}}(0)$. Here $\varepsilon_{\mathrm{lp}}'(\theta)=\frac{2\kappa}3\theta(\theta^2+a)\ge0$. By
Lemma~\ref{lem:mbounds}(i),(ii), $m\ge\max\{C_0,m_{\mathrm{lo}}\}$ on $(0,\infty)$; we use $m\ge C_0$ on $[0,r]$ (i.e.\ $\rho\le1$)
and $m\ge m_{\mathrm{lo}}$ on $[r,\theta_*]$ (note $m_{\mathrm{lo}}(1)=-3.34>C_0$, and $m_{\mathrm{lo}}$ is increasing, so no further loss occurs
there). Hence
\[
\int m\,d\mu\ \ge\ C_0\,\varepsilon_{\mathrm{lp}}(0)+C_0\big(\varepsilon_{\mathrm{lp}}(r)-\varepsilon_{\mathrm{lp}}(0)\big)+\int_r^{\theta_*}m_{\mathrm{lo}}(\theta/r)\,\varepsilon_{\mathrm{lp}}'(\theta)\,d\theta
=C_0\,\varepsilon_{\mathrm{lp}}(r)+\frac{2\kappa}3\big[\Phi_r(\theta_*)-\Phi_r(r)\big],
\]
which is \eqref{eq:coer-B}; the primitive is checked by differentiation. As $r\to0$: $\kappa\to1$, $a\to0$,
$\theta_*\to6^{1/4}$, $\varepsilon_{\mathrm{lp}}(r)\to0$, $\frac{2\kappa}3\Phi_r(r)\to0$, and
$\frac{2\kappa}3\Phi_r(\theta_*)=\log\theta_*-\log(2\pi r)-\frac14+o(1)$ (the $a$- and $r$-terms are $o(1)$),
whence $C_B(r)\to\log2\pi-\log6^{1/4}+\frac14$. The three values are evaluations of the closed form (double precision).
\end{proof}

Corollary~\ref{cor:coer-L1} uses $C_B$ only at $r=T/k$. The bound
uses the full multiplier; the cost of separating adjacent blocks
is computed in Remark~\ref{rem:coer-nogo}.




\begin{proposition}[the norm of the prime form]\label{prop:coer-prime}
\emph{(i)} $|\PiT(F)|\le\Aconst\|F\|^2$ for all $F\in L^2(\IT)$ and all $T$.
\emph{(ii)} For $\ell>0$ let $L(T,\ell):=\lceil2T/\ell\rceil$. Then $|\Rel\cF(\ell)|\le\cos\frac{\pi}{L(T,\ell)+1}\,\|F\|^2$,
hence
\[
\|\PiT\|\ \le\ \sum_{n\in\NT}v_n\cos\frac\pi{L(T,\log n)+1}
=\begin{cases}
\Aconst/2, & T\le\log2,\\[2pt]
\dfrac{v_2}{\sqrt2}+\dfrac{v_3+v_4}2=1.6740, & \log2<T\le\tfrac12\log5 .
\end{cases}
\]
Here $\frac{v_2}{\sqrt2}+\frac{v_3+v_4}2=\log2+\frac{\log3}{\sqrt3}+\frac{\log2}2$; in the first case the bound is
$\frac{v_2}2=0.4901$ for $T\le\frac12\log3$ and $\frac{v_2+v_3}2=1.1244$ for $\frac12\log3<T\le\log2$. The bounds hold on all of $L^2(\IT)$, in particular on $\mathcal V_k$.
\end{proposition}
\begin{proof}
(i) is $|\cF(\log n)|\le\|F\|^2$. (ii) Let $V_\ell$ be the shift $F\mapsto F(\cdot-\ell)\mathbf 1_{\IT}$ on
$L^2(\IT)$, so that $\Rel\cF(\ell)=\frac12\langle(V_\ell+V_\ell^*)F,F\rangle$. The map
$F\mapsto\big(F(-T+y+j\ell)\big)_{0\le j<n(y)}$, $y\in[0,\ell)$, $n(y):=\#\{j\ge0:-T+y+j\ell<T\}\le L(T,\ell)$, is unitary
from $L^2(\IT)$ onto $\int^\oplus_{[0,\ell)}\C^{n(y)}dy$, and it carries $V_\ell$ to the fibrewise shift $j\mapsto j+1$.
Hence $V_\ell+V_\ell^*$ acts on the fibre $\C^{n}$ as the adjacency matrix of the path with $n$ vertices, whose
eigenvalues are $2\cos\frac{j\pi}{n+1}$, $1\le j\le n$; its norm is $2\cos\frac\pi{n+1}\le2\cos\frac\pi{L+1}$. For
$T\le\log2$ one has $2T/\log n\le2$ for all $n\ge2$, so $L=2$ and $\cos\frac\pi3=\frac12$. For $\log2<T\le\frac12\log5$:
$\NT=\{2,3,4\}$, $L(T,\log2)=3$ since $2<2T/\log2\le\log5/\log2=2.32$, and $L(T,\log3)=L(T,\log4)=2$ since
$2T\le\log5<2\log3$; $\cos\frac\pi4=\frac1{\sqrt2}$. The values use $v_2=\sqrt2\log2$, $v_3=2\log3/\sqrt3$, $v_4=\log2$.
\end{proof}
For $\frac12\log2<T<\frac12\log3$ the bound $\Aconst/2$ is the exact norm: on a set of fibres of positive
measure the fibre has length $2$ and the prime form acts there as $-v_2\begin{psmallmatrix}0&1/2\\1/2&0\end{psmallmatrix}$.
Once $T>\log2$ the orbit of $\log2$ has length $3$ and the path norm exceeds $1$, so the bound
$\|\PiT\|\le\Aconst/2$ is specific to $T\le\log2$.

\begin{corollary}[coercivity of the truncated form on $\mathcal V_k$]\label{cor:coer-L1}
Let $k\ge1$ with $r=T/k\le1$, and put $c(T,k):=\inf\{\mathfrak a_T(F):F\in\mathcal V_k\cap\mathcal D_T,\ \|F\|=1\}$. Then
\begin{equation}\label{eq:coer-L1}
c(T,k)\ \ge\ \log\frac kT-C_B\Big(\frac Tk\Big)-\|\PiT\| ,
\end{equation}
with $\|\PiT\|$ bounded by Proposition~\ref{prop:coer-prime}. In particular (evaluations of the closed form):
\begin{center}
\begin{tabular}{@{}lllll@{}}
\toprule
$T$ & $k$ & $\dim E_k$ & bound for $\|\PiT\|$ & lower bound for $c(T,k)$\\
\midrule
$\frac12\log3$ & $8$ & $14$ & $0.4901$ & $0.486$\\
$\log2$ & $12$ & $22$ & $1.1244$ & $0.035$\\
$0.69$ & $14$ & $26$ & $1.1244$ & $0.201$\\
$\frac12\log5$ & $25$ & $48$ & $1.6740$ & $0.093$\\
$\frac12\log5$ & $26$ & $50$ & $1.6740$ & $0.134$\\
$\frac12\log5$ & $30$ & $58$ & $1.6740$ & $0.281$\\
$\frac12\log5$ & $85$ & $168$ & $\Aconst=2.9420$ & $0.069$\\
\bottomrule
\end{tabular}
\end{center}
Moreover, for every $T>0$ there is $k_0(T)$ such that $c(T,k)>0$ for all $k\ge k_0(T)$; indeed the right-hand side
of \eqref{eq:coer-L1} tends to $+\infty$ with $k$, because $C_B(T/k)\to1.6399$ (Theorem~\ref{thm:coer-B}).
\end{corollary}
\begin{proof}
$\mathfrak a_T=\mT+\PiT$ on $\mathcal D_T$, Theorem~\ref{thm:coer-B} and Proposition~\ref{prop:coer-prime}.
\end{proof}

By \eqref{eq:coer-L1}, $k>Te^{\Aconst+C_B(T/k)}$ suffices.
As $C_B(T/k)\to1.6399$ and $\Aconst\asymp e^T$, this sufficient dimension
grows doubly exponentially. For $T\le\tfrac12\log5$ one can keep $2k-2<60$
by Corollary~\ref{cor:coer-L1}.




\subsection{The finite reduction}\label{subsec:coer-schur}

\begin{theorem}[reduction to the piecewise Tate exponentials]\label{thm:coer-schur}
Fix $T>0$ and $k\ge1$ with $c:=c(T,k)>0$ (Corollary~\ref{cor:coer-L1}: every $k\ge k_0(T)$). Put
$\mathcal V:=\mathcal V_k\cap\mathcal D_T$ and $M:=\Aconst-C_0$.
\emph{(i)} $\mathfrak a_T$ is an inner product on $\mathcal V$ for which $\mathcal V$ is complete, and for every
$e\in E_k$ the functional $v\mapsto\mathfrak a_T(e,v)$ is bounded on $(\mathcal V,\mathfrak a_T)$; hence there is a unique
$\mathsf Ke\in\mathcal V$ with $\mathfrak a_T(\mathsf Ke,v)=\mathfrak a_T(e,v)$ for all $v\in\mathcal V$. The map $\mathsf K:E_k\to\mathcal V$ is linear,
and $\|\mathsf K\|:=\sup\{\|\mathsf Ke\|:\|e\|=1\}<\infty$.
\emph{(ii)} For $v\in\mathcal V$ and $e\in E_k$,
\begin{equation}\label{eq:coer-schur}
\begin{split}
\mathfrak a_T(v+e)&=\mathfrak a_T(v+\mathsf Ke)+\mathfrak s_{T,k}(e),\\
\mathfrak s_{T,k}(e)&:=\mathfrak a_T(e)-\mathfrak a_T(\mathsf Ke)=\mathfrak a_T(e)-\sup_{u\in\mathcal V}\big[2\Rel\mathfrak a_T(e,u)-\mathfrak a_T(u)\big].
\end{split}
\end{equation}
Let $S_{T,k}$ be the matrix of the real quadratic form $\mathfrak s_{T,k}$ in an orthonormal basis of $E_k$ (a
real symmetric $(2k-2)\times(2k-2)$ matrix). For $k=1$, $E_k=0$ and $\lamA=c$.
For $k\ge2$, writing $\lambda_S=\min\sigma(S_{T,k})$, one has
\begin{equation}\label{eq:coer-sandwich}
\begin{aligned}
 \frac{\min\{c,\lambda_S\}}{(1+\|\mathsf K\|)^2}
 &\le\lamA\le\min\{c,\lambda_S\} &&(\lambda_S\ge0),\\
 \lambda_S &\le\lamA\le\frac{\lambda_S}{1+\|\mathsf K\|^2} &&(\lambda_S<0).
\end{aligned}
\end{equation}
\emph{(iii)} Consequently $\lamA>0$ if and only if $S_{T,k}\succ0$, and $\lamA\ge0$ if and only if $S_{T,k}\succeq0$.
\end{theorem}
\begin{proof}
(i) By Lemma~\ref{lem:mbounds}, $\sigma_T\ge C_0-\Aconst$, so $\mathfrak a_T+M\|\cdot\|^2\ge0$ on $\mathcal D_T$, and
by Lemma~\ref{lem:mbounds}(ii) and $m\ge C_0$ there are constants $c_1,c_2>0$ (depending on $T$) with
$c_1\int\log(2+|s|)|\Fh|^2\le\mathfrak a_T(F)+(M+1)\|F\|^2\le c_2\int\log(2+|s|)|\Fh|^2$ on $\mathcal D_T$. On $\mathcal V$,
$\mathfrak a_T\ge c\|\cdot\|^2$ gives $\mathfrak a_T(v)\ge\frac c{c+M+1}\big(\mathfrak a_T(v)+(M+1)\|v\|^2\big)$, so
$\mathfrak a_T$ is equivalent on $\mathcal V$ to the weighted norm $\int\log(2+|s|)|\widehat v|^2$, for which
$\mathcal D_T$ is complete and $\mathcal V_k\cap\mathcal D_T$ is closed ($\mathcal V_k$ is $L^2$-closed and the weighted norm
dominates the $L^2$ norm). For $e\in E_k\subset\mathcal D_T$ (Lemma~\ref{lem:coer-decomp}(iv)), the Cauchy--Schwarz
inequality for the nonnegative form $\mathfrak a_T+M\langle\cdot,\cdot\rangle$ gives
$|\mathfrak a_T(e,v)|\le(\mathfrak a_T(e)+M\|e\|^2)^{1/2}(\mathfrak a_T(v)+M\|v\|^2)^{1/2}+M\|e\|\|v\|\le C_e\,\mathfrak a_T(v)^{1/2}$,
so $\mathsf Ke$ exists by the Riesz representation theorem; $\mathsf K$ is linear, and $\|\mathsf K\|<\infty$ since $\dim E_k<\infty$.
(ii) $\mathfrak a_T(v+\mathsf Ke)=\mathfrak a_T(v)+2\Rel\mathfrak a_T(\mathsf Ke,v)+\mathfrak a_T(\mathsf Ke)=\mathfrak a_T(v)+2\Rel\mathfrak a_T(e,v)+\mathfrak a_T(\mathsf Ke)$,
which rearranges to \eqref{eq:coer-schur}; the variational expression is $\mathfrak a_T(\mathsf Ke)=\sup_u[2\Rel\mathfrak a_T(\mathsf Ke,u)-\mathfrak a_T(u)]$
(maximum at $u=\mathsf Ke$). For \eqref{eq:coer-sandwich}: every $F\in\PT\cap\mathcal D_T$ is $v+e$ (Lemma~\ref{lem:coer-decomp}(iv)),
$\|F\|^2=\|v\|^2+\|e\|^2$, and $v\mapsto v+\mathsf Ke$ is a bijection of $\mathcal V$. Taking $e=0$ gives $\lamA\le c$. For $k\ge2$, take a unit least
eigenvector $e$ and $v=-\mathsf Ke$. The quotient is exactly
$\lambda_S/(1+\|\mathsf Ke\|^2)$, giving the upper bound in each case.
If $\lambda_S\ge0$, the square completion bounds the form below by
$\min(c,\lambda_S)(\|v+\mathsf Ke\|^2+\|e\|^2)$, and the inverse
change of variables on the orthogonal sum has norm at most $1+\|\mathsf K\|$.
If $\lambda_S<0$, drop the nonnegative square to get
$\mathfrak a_T(v+e)\ge\lambda_S\|e\|^2\ge\lambda_S\|v+e\|^2$.
These bounds also prove (iii).
\end{proof}

\begin{remark}[inertia of the two Schur complements]\label{rem:inertiatower}
If $S_{T,k}$ is nonsingular, the change of variables in
\eqref{eq:coer-schur} splits the primitive form into its coercive block
and $S_{T,k}$. Thus its negative index is $n_-(S_{T,k})$ and its radical
is zero. Hence the primitive operator has trivial kernel; since it has compact resolvent
(Lemma~\ref{lem:compact}) its spectrum is discrete, so zero lies in its resolvent set in each
parity sector (Theorem~\ref{thm:parityschur}(o)). Theorem~\ref{thm:parityschur}(c) then gives
$n_-(W_T)=n_-(S_{T,k})+n_-(\mathfrak G_T)$ and
$n_0(W_T)=n_0(\mathfrak G_T)$. The size $2k-2$ counts blocks;
crossing a contact varies the entries continuously at fixed $k$.
\end{remark}



At $T=\tfrac12\log5$ one may take $k=25$ in \eqref{eq:coer-sandwich}.
Galerkin approximation bounds the Schur complement from above; a lower
certificate needs an error bound using $c>0$. For Legendre bases see
\cite[Thms.~1.1, 1.4]{Zhu2026} and \cite{Groskin2026}.

\begin{remark}[the route by operator norms]\label{rem:coer-nogo}
For adjacent blocks the Tate-orthogonal coupling has bilinear norm
$\pi$. In endpoint distances $x,y$ the kernel satisfies
$0<K_1(x+y)\le1/(x+y)$, so Hilbert's inequality gives the upper
bound \cite[Thm.~316]{HLP}. For sharpness take
$f_M(x)=x^{-1/2}\mathbf1_{[2re^{-M},2r]}(x)$ on each block.
Its squared norm is $M$; logarithmic coordinates give the Hilbert
pairing $\int_{[-M,0]^2}(2\cosh((u-v)/2))^{-1}du\,dv=\pi M-O(1)$.
The bounded kernel difference costs $O(r)$ since $\|f_M\|_1^2=O(r)$.
At fixed $r$ the projection onto the two Tate directions has bounded
norm, changing the pairing by $O(\sqrt M)$. Thus in
\eqref{eq:coer-cross} the coupling costs $\pi/2$ for two blocks
and at least this much for $k\ge2$. Subtracting the coupling and
prime norms separately in \eqref{eq:coer-A} requires
$\log(1/r)>1.209+r+\pi/2+1.674$ at $T=\tfrac12\log5$, hence
$k\ge70$ versus $25$ in Theorem~\ref{thm:coer-B}.
\end{remark}


\begin{lemma}[first-contact reduction]\label{lem:firstcontact}
If $\lamA<0$ for some $T$, then there is a window $T_*>\frac12\log5$ and a nonzero $F\in\mathcal P_{T_*}$ such
that $\lambda_{\mathrm A}(T)\ge0$ for $T\le T_*$, $\lambda_{\mathrm A}(T)<0$ for $T>T_*$, and
\[
 H_{T_*}\ge0,\qquad H_{T_*}F=0,\qquad JF=\varepsilon F,\quad \varepsilon\in\{\pm1\}.
\]
\end{lemma}
\begin{proof}
Set $T_*=\inf\{T:\lamA<0\}$. Monotonicity and continuity
(Lemmas~\ref{lem:monoT} and~\ref{lem:contA}) give the asserted signs
and $\lambda_{\mathrm A}(T_*)=0$; Theorem~\ref{thm:log5A} gives the
strict lower bound on $T_*$. Compact resolvent (Lemma~\ref{lem:compact})
gives a zero eigenvector, and reflection preserves its eigenspace.
\end{proof}
\section{A common family of tests for all windows}\label{sec:ladder}
There is one sequence of matrices which suffices to test the sign for every
window. Its entries can be computed before choosing any zeros. For $a>1$
write
\[
 L(a)=\frac1a+\frac1{a-1}-\frac12\log\pi+
 \frac12\psi(a/2)+\frac{\zeta'(a)}{\zeta(a)},\qquad
 K(a,b)=\frac{L(a)+L(b)}{a+b-1}.
\]
For real tests $h$ on the half-line, extended by zero, write $\mathcal W[h]$
for \eqref{eq:W} with all prime powers included. We use it only where
the integral, moments and prime series converge absolutely. Put
$h_a(y)=e^{-(a-1/2)y}\mathbf1_{y>0}$. The polarized form satisfies
$\mathcal W[h_a,h_b]=K(a,b)$. Indeed the positive shifts
in the symmetrized autocorrelation contribute the two Dirichlet series for
$\zeta'/\zeta$, divided by $a+b-1$. The moments contribute
$1/((a-1)b)+1/(a(b-1))$. For the archimedean integral the partial fractions
of the two Fourier transforms give Poisson kernels of widths $a-1/2$ and
$b-1/2$, divided by $a+b-1$. Poisson integration of $m$ supplies the
digamma and constant terms displayed above.

\begin{proposition}\label{ladder}
If $[K(32+i,32+j)]_{0\le i,j\le m}$ is positive semidefinite for every
$m\ge0$, then $W_T\ge0$ on $\mathcal D_T$ for every $T>0$.
\end{proposition}
\begin{proof}
Translate a real smooth compact test $h$ into $(0,\infty)$, preserving its
autocorrelation and moment product. The function $f(x)=x^{-63/2}h(-\log x)$
extends smoothly by zero at both ends of $[0,1]$. Approximate its derivative
by polynomials and integrate to obtain $P_j$ with $\epsilon_j=\|P_j-f\|_{C^1}\to0$.
The functions $e^{-63y/2}P_j(e^{-y})$ belong to the stated spans on $y>0$.
Their errors and first derivatives there are $O(\epsilon_j e^{-63y/2})$;
integration by parts after zero extension, including the boundary value,
bounds the Fourier error by $O(\epsilon_j/(1+|\tau|))$. Hence the
archimedean integral and both Tate moments converge. At $\log n$ the
autocorrelation error is $O(\epsilon_j n^{-63/2})$ by the common exponential
bound, summable against $v_n$. Thus the matrix inequalities pass to $h$,
to complex tests by polarization of the real symmetric form, and to the
form domain by core density (Lemma~\ref{lem:form-density}).
\end{proof}

The following finite-order comparisons are reported without proof; they are not used below. For up to four distinct real exponents $a>1$ the matrix $K$ is positive
definite. For up to thirteen real poles of rates
$a-1/2\ge63/2$ the complete form exceeds $\|h\|_2^2/4000$ on nonzero
rational Laplace data with proper numerator, including repeated poles. Proposition~\ref{ladder} requires the sign also
when the degree grows without bound.

\begin{proposition}[the matrices in every degree]\label{prop:all-degree-gram}
With the arithmetic input \eqref{eq:dirichlet-input}, let $a_j>1$
and $u_j=a_j-\tfrac12$, $1\le j\le d$. Then, for $c_j\in\mathbb C$,
\[
 \sum_{j,k=1}^d c_j\overline{c_k}K(a_j,a_k)
 =\sum_{\gamma}\left|\sum_{j=1}^d\frac{c_j}{u_j-i\gamma}\right|^2 .
\]
Here $\tfrac12+i\gamma$ ranges over the nontrivial zeros with multiplicity.
In particular the hypothesis of Proposition~\ref{ladder} holds in every degree.
\end{proposition}
\begin{proof}
By evenness the order-one Hadamard product pairs the zeros at $\pm i\gamma$
of $\xi(\tfrac12+z)$; $N(H)=O(H\log H)$ gives locally uniform convergence
of the paired logarithmic derivatives away from the zeros. Thus
$L(\tfrac12+u)=\sum_\gamma(u-i\gamma)^{-1}$ in symmetric order.
Pair the terms at $\gamma,-\gamma$ in the two evaluations defining $K$:
$[(u-i\gamma)^{-1}+(v+i\gamma)^{-1}]/(u+v)
=[(u-i\gamma)(v+i\gamma)]^{-1}$. Summing gives the identity.
\end{proof}

% sec-analytic.tex: analytic part.  Labels an:...; macros \an...
\providecommand{\anThmA}{\ref{thm:A}}
\providecommand{\anThmB}{\ref{thm:B}}
\providecommand{\anThmC}{\ref{thm:C}}
\providecommand{\anchi}{\chi_K}
\providecommand{\anLam}{\Lambda_K}
\providecommand{\anlam}{\lambda}
\providecommand{\aneu}{\mathrm{eu}_d}
\providecommand{\anord}{\operatorname{ord}}
\providecommand{\anF}{\mathcal F}
\providecommand{\anPl}{\Phi}
\providecommand{\anPa}{\Phi_a}
\providecommand{\anPb}{\Phi_b}
\providecommand{\ancl}{\mathrm{cl}}
\providecommand{\anrd}{\mathrm{rd}}
\providecommand{\anrem}{\operatorname{rem}}
\providecommand{\anhm}{h_{\mathrm{main}}}
\providecommand{\anhd}{h_{\mathrm{dir}}}
\providecommand{\ansf}{s_{\mathrm{fin}}}
\providecommand{\ansm}{s_{\mathrm{main}}}
\providecommand{\ansi}{s_{\infty}}
\providecommand{\andT}{d_{\mathcal T}}
\providecommand{\anCP}{C_P}
\providecommand{\ancE}{\mathcal E}

\section{A one-sided criterion for Dirichlet \texorpdfstring{$L$}{L}-functions}\label{an:sec-criterion}
The zero condition we need is
\begin{equation}\label{eq:dirichlet-input}
 \zeta(s)\ne0\qquad (\Re s>\tfrac12,\ s\ne1).
\end{equation}
It is the case of the trivial character in Theorem~\ref{thm:GRH-half}: no
Dirichlet $L$-function vanishes in $\Re s>\frac12$. The argument has four
parts. This section proves a criterion of Landau's type
(Theorem~\ref{an:thm-criterion}): a one-sided bound $R_q(x)\le x^{\beta-1}$
for a prime sum $R_q$, attached to a modulus $q$ and to a quadratic character
$\anchi$, excludes the zeros with $\Re s>\beta$ of all $L(s,\psi)$,
$\psi$ modulo $q$, provided $\anchi$ may be chosen with twists
$L(s,\psi\anchi)$ free of zeros in a box. Section~\ref{an:sec-data} defines
the admissible data and proves that they exist. For such data
Section~\ref{sec:automorphic} attaches a space of automorphic forms with
integral structures, and states three theorems
(Theorems~\anThmA, \anThmB\ and~\anThmC) that bound one quantity from above
and from below. Section~\ref{an:sec-main} combines these bounds into
the inequality $m\cdot\anrem\le f\,a\,(\log b)^5$ (Theorem~\ref{an:thm-main}),
identifies $\anrem$ with a combination of values of $R_q$, and deduces the
one-sided bound $R_{q}(y)\le y^{-1/2+\varepsilon}$ and the theorems on zeros.

Everything in this section and in Sections~\ref{an:sec-data}
and~\ref{an:sec-main} is proved in the text, from the following published
results: the large sieve for real characters of Heath-Brown
(Theorem~\ref{an:thm-HB}), the prime number theorem for arithmetic
progressions with de la Vall\'ee Poussin's error term and the theorem of
Siegel and Walfisz \cite[Chapters~20 and~22]{Davenport}, the inequality of
P\'olya and Vinogradov \cite[Chapter~23]{Davenport}, and standard facts on
Dirichlet $L$-functions and on cyclotomic fields. The deduction of
Theorem~\ref{an:thm-main} uses Theorems~\anThmA, \anThmB\ and~\anThmC\ and
nothing else from Section~\ref{sec:automorphic}; those three theorems are
proved there and in the sections that follow it (see \S\ref{ss:au-remains}), and all results of
Section~\ref{an:sec-main} from Theorem~\ref{an:thm-main} on rest on them.

\subsection{The prime sum \texorpdfstring{$R_q$}{R\_q} and its transform}\label{an:ssec-R}
Throughout this section $d$ is an odd prime, $\anchi=(\cdot/d)$ is the
Legendre symbol, a primitive real character modulo $d$, and $q\ge1$ is an
integer not divisible by $d$. When $d\equiv3\pmod4$, $\anchi$ is the
character of the field $K=\mathbb Q(\sqrt{-d})$. Let $\Lambda$ be von
Mangoldt's function and put
\[
 \anlam_q(n)=\begin{cases}\Lambda(n)\,(1-\anchi(n)),& n\equiv1\pmod q,\\ 0,&\text{otherwise,}\end{cases}
 \qquad \anLam=\anlam_1 .
\]
Thus $0\le\anlam_q(n)\le2\Lambda(n)$, and $\anLam$ is supported on the powers
of the primes that do not split in $K$. For a Dirichlet character $\psi$
modulo $q$ the product $\psi\anchi$ is a character modulo $qd$, and it is not
principal, because $\anchi$ is primitive modulo $d$ and $d\nmid q$. By
orthogonality of the characters modulo $q$, for $\Re s>1$,
\begin{equation}\label{an:eq-D}
 D_q(s):=\sum_{n\ge1}\frac{\anlam_q(n)}{n^{s}}
 =\frac1{\varphi(q)}\sum_{\psi\bmod q}\Bigl(\frac{L'}{L}(s,\psi\anchi)-\frac{L'}{L}(s,\psi)\Bigr).
\end{equation}
The right side is meromorphic in $\mathbb C$. At $s=1$ only the principal
character $\psi_0$ contributes a pole, $-(L'/L)(s,\psi_0)=(s-1)^{-1}+O(1)$,
and no $L(s,\chi)$ vanishes on $\Re s\ge1$. Hence
\[
 G_q(s):=D_q(s)-\frac{s}{\varphi(q)\,(s-1)}
\]
is holomorphic in a neighbourhood of the closed half-plane $\Re s\ge1$. It is
real on $(1,\infty)$, since the coefficients $\anlam_q(n)$ are real; hence
\[
 c_q:=G_q(1)\in\mathbb R .
\]
For $q=1$, the expansion $-(\zeta'/\zeta)(s)=(s-1)^{-1}-\gamma_E+O(s-1)$
gives
\begin{equation}\label{an:eq-c1}
 c_1=-1-\aneu,\qquad \aneu:=\gamma_E-\frac{L'}{L}(1,\anchi).
\end{equation}

\begin{definition}\label{an:def-R}
For real $x\ge1$,
\[
 R_q(x)=R_{d,q}(x)=\sum_{n\le x}\anlam_q(n)\Bigl(\frac1n-\frac1x\Bigr)-\frac{\log x}{\varphi(q)}-c_q .
\]
\end{definition}

For $\rho\in\mathbb C$ with $\Re\rho>0$ put
\begin{equation}\label{an:eq-r}
 n_q(\rho)=\sum_{\psi\bmod q}\anord_\rho L(\cdot,\psi),\qquad
 n_q^{K}(\rho)=\sum_{\psi\bmod q}\anord_\rho L(\cdot,\psi\anchi),\qquad
 r_q(\rho)=\frac{n_q(\rho)-n_q^{K}(\rho)}{\varphi(q)\,\rho\,(1-\rho)}\quad(\rho\ne1),
\end{equation}
where $\anord_\rho$ is the order of vanishing at $\rho$, and $r_q(1)=0$.
For $\rho\ne1$ the numbers $n_q(\rho)$ and $n_q^K(\rho)$ are non-negative
integers.

\begin{lemma}\label{an:lem-R}
\begin{enumerate}
\item[(i)] Let $\ancE_q(t)=\sum_{n\le t}\anlam_q(n)-t/\varphi(q)$. Then
$R_q(y)=\int_1^{y}\ancE_q(t)\,t^{-2}\,dt-c_q$ for $y\ge1$. There is an
absolute constant $C_{\mathrm{Ch}}$ with $|\ancE_q(t)|\le C_{\mathrm{Ch}}t$ for $t\ge1$; hence
$|R_q(z)-R_q(y)|\le C_{\mathrm{Ch}}(z-y)/y$ for $1\le y\le z$, and
$|R_q(y)|\le C_{\mathrm{Ch}}\log y+|c_q|$.
\item[(ii)] For $\Re s>1$,
\[
 \anF_q(s):=\int_0^\infty R_q(e^{v})\,e^{(1-s)v}\,dv=\frac{G_q(s)/s-c_q}{s-1}.
\]
\item[(iii)] $\anF_q$ extends to a meromorphic function on $\mathbb C$,
holomorphic on a neighbourhood of $\Re s\ge1$. Its poles in $\Re s>0$ are
simple, they lie at zeros of the functions $L(s,\psi)$ and $L(s,\psi\anchi)$,
$\psi$ modulo $q$, and for every $\rho$ with $\Re\rho>0$
\[
 \lim_{s\to\rho}(s-\rho)\,\anF_q(s)=r_q(\rho).
\]
\end{enumerate}
\end{lemma}
\begin{proof}
(i) Since $\int_n^{y}t^{-2}dt=1/n-1/y$, the sum in
Definition~\ref{an:def-R} is $\int_1^{y}\bigl(\sum_{n\le t}\anlam_q(n)\bigr)t^{-2}dt$,
and $\int_1^{y}(t/\varphi(q))t^{-2}dt=\log y/\varphi(q)$. Chebyshev's bound
$\sum_{n\le t}\Lambda(n)\ll t$ and $0\le\anlam_q\le2\Lambda$ give
$|\ancE_q(t)|\le C_{\mathrm{Ch}}t$. The other two statements follow by integration.

(ii) By (i) the integral converges absolutely for $\Re s>1$. With $x=e^{v}$,
\[
 \int_1^\infty\sum_{n\le x}\anlam_q(n)\Bigl(\frac1n-\frac1x\Bigr)x^{-s}\,dx
 =\sum_{n\ge1}\anlam_q(n)\,n^{-s}\Bigl(\frac1{s-1}-\frac1s\Bigr)=\frac{D_q(s)}{s(s-1)},
\]
and $\int_1^\infty\log x\,x^{-s}dx=(s-1)^{-2}$,
$\int_1^\infty x^{-s}dx=(s-1)^{-1}$. Inserting
$D_q(s)=s/(\varphi(q)(s-1))+G_q(s)$, the terms with $(s-1)^{-2}$ cancel and
\[
 \anF_q(s)=\frac{D_q(s)}{s(s-1)}-\frac1{\varphi(q)(s-1)^2}-\frac{c_q}{s-1}
 =\frac{G_q(s)}{s(s-1)}-\frac{c_q}{s-1}.
\]

(iii) $G_q$ is meromorphic on $\mathbb C$ by \eqref{an:eq-D}, and holomorphic
near $\Re s\ge1$. At $s=1$ the numerator $G_q(s)/s-c_q$ vanishes, by the
definition of $c_q$; so $\anF_q$ is holomorphic there. Let $\Re\rho>0$,
$\rho\ne1$. The logarithmic derivative of a meromorphic function $F\not\equiv0$ has
at most a simple pole at $\rho$, with residue $\anord_\rho F$. By
\eqref{an:eq-D}, $(s-\rho)G_q(s)\to(n_q^K(\rho)-n_q(\rho))/\varphi(q)$, and
division by $\rho(\rho-1)$ gives $r_q(\rho)$.
\end{proof}


\subsection{The prime number theorem for \texorpdfstring{$R_q$}{R\_q}}\label{an:ssec-pnt}
\begin{proposition}\label{an:prop-pnt}
There are constants $c>0$ and $C$, depending on $d$ and $q$, such that
\[
 |R_q(y)|\le C\exp\bigl(-c\sqrt{\log y}\bigr)\qquad(y\ge2).
\]
In particular, for every $B>0$, $|R_q(y)|\le(\log y)^{-B}$ for all large $y$.
\end{proposition}
\begin{proof}
For a character $\chi$ put $\Sigma(t,\chi)=\sum_{n\le t}\Lambda(n)\chi(n)$. By
orthogonality,
\[
 \sum_{n\le t}\anlam_q(n)=\frac1{\varphi(q)}\sum_{\psi\bmod q}\bigl(\Sigma(t,\psi)-\Sigma(t,\psi\anchi)\bigr).
\]
The moduli $q$ and $qd$ are fixed. By the prime number theorem for
arithmetic progressions \cite[Chapter~20]{Davenport}, for a character $\chi$
of fixed modulus $\Sigma(t,\chi)=\epsilon_\chi t+O(t\exp(-c_1\sqrt{\log t}))$,
with $\epsilon_\chi=1$ if $\chi$ is principal and $0$ otherwise; a possible
real zero close to $1$ is a fixed number $\beta_1<1$ for fixed modulus, and
its term $t^{\beta_1}$ is absorbed. No $\psi\anchi$ is principal, so
$\ancE_q(t)\ll t\exp(-c_1\sqrt{\log t})$. Hence
$I=\int_1^\infty\ancE_q(t)t^{-2}dt$ converges absolutely, and by
Lemma~\ref{an:lem-R}(i), with $\ell=I-c_q$,
\[
 R_q(y)-\ell=-\int_y^\infty\ancE_q(t)\,t^{-2}\,dt\ll\int_{\log y}^\infty e^{-c_1\sqrt u}\,du\ll\exp\bigl(-c\sqrt{\log y}\bigr)
\]
for any $c<c_1$. It remains to see that $\ell=0$. For real $\sigma>1$,
\[
 \anF_q(\sigma)=\frac{\ell}{\sigma-1}+\int_0^\infty\bigl(R_q(e^{v})-\ell\bigr)e^{-(\sigma-1)v}\,dv ,
\]
and the last integral is bounded uniformly in $\sigma>1$, because
$R_q(e^{v})-\ell\ll e^{-c\sqrt v}$ is integrable. By
Lemma~\ref{an:lem-R}(iii), $\anF_q$ is holomorphic at $1$, so $\anF_q(\sigma)$
is bounded as $\sigma\to1^{+}$. Therefore $\ell=0$.
\end{proof}
The constant $c_q=G_q(1)$ is thus the unique constant for which $R_q\to0$.


\subsection{Landau's lemma}\label{an:ssec-landau}
The following is Landau's theorem on integrals of one-signed functions
(see \cite[\S15.1, Lemma~15.1]{MV}), in the form adapted to transforms with simple
poles. We give the proof.

\begin{lemma}[Landau]\label{an:lem-landau}
Let $E\colon[0,\infty)\to\mathbb R$ be measurable with $|E(v)|\le C(1+v)$, and
suppose that $\anF(s)=\int_0^\infty E(v)e^{(1-s)v}dv$, $\Re s>1$, extends to a
meromorphic function on $\Re s>0$ whose poles on the positive real axis are
simple. Let $\beta>0$ and $v_0\ge0$ be such that
\[
 E(v)\le e^{(\beta-1)v}\qquad(v\ge v_0).
\]
Then one of the following holds:
\begin{enumerate}
\item[(a)] $\anF$ is holomorphic in $\Re s>\beta$;
\item[(b)] there is a real $x_0>\beta$ such that $\anF$ is holomorphic in
$\Re s>x_0$ and has a pole at $x_0$ whose residue is real and negative.
\end{enumerate}
\end{lemma}
\begin{proof}
Put $f(v)=e^{(\beta-1)v}-E(v)$; then $f\ge0$ on $[v_0,\infty)$. For
$\Re s>\max(1,\beta)$,
\[
 I(s):=\int_{v_0}^\infty f(v)e^{(1-s)v}\,dv=\frac{e^{(\beta-s)v_0}}{s-\beta}+\int_0^{v_0}E(v)e^{(1-s)v}\,dv-\anF(s)=:\Psi(s).
\]
The second term is entire, so $\Psi$ is meromorphic in $\Re s>0$; its poles
at real points $x>\beta$ are those of $\anF$, they are simple, and
$\operatorname{Res}_{x}\Psi=-\operatorname{Res}_{x}\anF$. Let $\sigma_c$ be
the infimum of the real $\sigma$ for which
$\int_{v_0}^\infty f(v)e^{(1-\sigma)v}dv<\infty$. Since $f\ge0$, the integral
$I(s)$ converges absolutely and is holomorphic in $\Re s>\sigma_c$, and it
equals $\Psi$ there by the identity theorem; in particular $\Psi$, and so
$\anF$, is holomorphic in $\Re s>\max(\sigma_c,\beta)$.

If $\sigma_c\le\beta$ we have (a). Let $\sigma_c>\beta$, and suppose first
that $\Psi$ is holomorphic at $\sigma_c$. Put $a=\sigma_c+1$. The closed disc
$|s-a|\le1$ lies in $\{\Re s>\sigma_c\}\cup\{\sigma_c\}$, so $\Psi$ is
holomorphic in a disc $|s-a|<1+2\delta$, $\delta>0$, and its Taylor series
at $a$ converges at $a-1-\delta$:
\[
 \Psi(a-1-\delta)=\sum_{k\ge0}\frac{(1+\delta)^k}{k!}(-1)^k\Psi^{(k)}(a)
 =\sum_{k\ge0}\frac{(1+\delta)^k}{k!}\int_{v_0}^\infty f(v)\,v^{k}e^{(1-a)v}\,dv .
\]
All terms are non-negative, so sum and integral may be exchanged, and
$\int_{v_0}^\infty f(v)e^{(1-a+1+\delta)v}dv<\infty$. This is convergence at
$\sigma_c-\delta$, contrary to the definition of $\sigma_c$. Hence $\Psi$
has a pole at $\sigma_c$, simple by hypothesis, with residue
$\kappa=\lim_{x\to\sigma_c^+}(x-\sigma_c)I(x)$. As $I(x)\ge0$ for real
$x>\sigma_c$, $\kappa$ is real and $\ge0$, and $\kappa\ne0$. So $\anF$ has a
pole at $x_0=\sigma_c$ with residue $-\kappa<0$, which is (b).
\end{proof}


\subsection{The choice of the prime \texorpdfstring{$d$}{d}}\label{an:ssec-choice}
The published input is the large sieve for real characters. We write
$S(Q)$ for the set of all real primitive Dirichlet characters of conductor at
most $Q$, and $\square$ for an unspecified perfect square.

\begin{theorem}[Heath-Brown {\cite[Corollary~2]{HeathBrown1995}}]\label{an:thm-HB}
Let $N,Q$ be positive integers and $a_1,\dots,a_N$ arbitrary complex numbers. For every
$\epsilon>0$,
\[
 \sum_{\chi\in S(Q)}\Bigl|\sum_{n\le N}a_n\chi(n)\Bigr|^2
 \ \ll_\epsilon\ (QN)^{\epsilon}(Q+N)\sum_{\substack{n_1,n_2\le N\\ n_1n_2=\square}}|a_{n_1}a_{n_2}| .
\]
\end{theorem}
We use it only for the characters $\anchi=\chi_d=(\cdot/d)$ with $d\le Q$ an odd
prime: the Legendre symbol modulo $d$ is a real primitive character of
conductor $d$, so it belongs to $S(Q)$; by positivity the other
terms may be dropped. The coefficients $a_n$ are arbitrary, and this is
what allows the twist by a fixed character $\xi$ below. The zero-density
theorems of \cite{HeathBrown1995} are stated for real characters; the
statement we need concerns the characters $\xi\chi_d$, and we derive it
directly from Theorem~\ref{an:thm-HB}.

\begin{proposition}[zeros of the twists are rare]\label{an:prop-density}
Let $N\ge1$, $\frac12<\beta<1$ and $H\ge0$. Put $\theta=\beta-\frac12$ and
$\eta=\theta^2/8$. There is a constant $C=C(N,\beta,H)$ such that for all
$X\ge2$ the number of odd primes $d\le X$, $d\nmid N$, for which some
$L(s,\xi\chi_d)$, $\xi$ a Dirichlet character modulo $N$, has a zero in the
box
\[
 \mathcal B=\{s:\ \Re s\ge\beta,\ |\Im s|\le H\}
\]
is at most $C\,X^{1-\eta}$.
\end{proposition}
\begin{proof}
Let $\sigma_0=\frac12+\frac\theta2$ and $r_0=\beta-\sigma_0=\theta/2$; thus
$2\sigma_0-1=\theta<\frac12$. Let $X$ be large and put
\[
 Y=X^{\theta/4},\qquad Z=X^{1-\theta/4},\qquad YZ=X,\quad Y\le Z .
\]
Fix $\xi$ modulo $N$ and an odd prime $d\le X$, $d\nmid N$, and write
$\chi=\xi\chi_d$, a non-principal character modulo $Nd$. All constants below
depend on $N,\beta,H$ only.

\emph{Step 1: truncation.} Let $S(u)=\sum_{n\le u}\chi(n)$. By the
inequality of P\'olya and Vinogradov \cite[Chapter~23]{Davenport},
$|S(u)-S(u')|\le P:=c_0\sqrt X\log X$ for all $u,u'$. (For a primitive
character $\chi^{*}$ modulo $q^{*}$ the inequality reads
$|\sum_{u'<n\le u}\chi^{*}(n)|\ll\sqrt{q^{*}}\log q^{*}$. Our $\chi$ is induced by
a primitive character $\chi^{*}$ of conductor $q^{*}\le NX$, and
$\sum_{n\le u}\chi(n)=\sum_{e\mid N}\mu(e)\chi^{*}(e)\sum_{k\le u/e}\chi^{*}(k)$,
which gives the bound with a constant depending on $N$.) For $\Re s=\sigma>0$
partial summation gives
\begin{equation}\label{an:eq-trunc}
 L(s,\chi)=\sum_{n\le Z}\frac{\chi(n)}{n^{s}}+E(s),\qquad
 E(s)=s\int_Z^\infty\bigl(S(u)-S(Z)\bigr)u^{-s-1}\,du,\qquad |E(s)|\le\frac{|s|}{\sigma}\,P\,Z^{-\sigma}.
\end{equation}

\emph{Step 2: the mollified function.} Let
$M(s)=\sum_{k\le Y}\mu(k)\chi(k)k^{-s}$ and $g(s)=g_{d,\xi}(s)=1-L(s,\chi)M(s)$,
an entire function. Since $\chi$ is completely multiplicative,
\[
 \sum_{n\le Z}\frac{\chi(n)}{n^{s}}\cdot M(s)=\sum_{n\le X}\frac{b_n\chi(n)}{n^{s}},\qquad
 b_n=\sum_{\substack{k\mid n,\ k\le Y\\ n/k\le Z}}\mu(k).
\]
If $n\le Y$, both conditions hold for all divisors $k$ of $n$, so $b_1=1$ and
$b_n=0$ for $1<n\le Y$; and $|b_n|\le\tau(n)$ always, $\tau$ being the
divisor function. With \eqref{an:eq-trunc},
\begin{equation}\label{an:eq-g}
 g(s)=-\sum_{Y<n\le X}\frac{b_n\,\xi(n)\,\chi_d(n)}{n^{s}}-E(s)M(s).
\end{equation}

\emph{Step 3: detection of a zero.} Let
$\mathcal R=\{s:\sigma_0\le\Re s\le2,\ |\Im s|\le H+1\}$. Suppose
$L(\rho,\chi)=0$ with $\rho\in\mathcal B$. Then $\Re\rho<1$, the disc
$|s-\rho|\le r_0$ lies in $\mathcal R$, and $g(\rho)=1$. By the mean value
property of the holomorphic function $g$ and the Cauchy--Schwarz inequality,
\begin{equation}\label{an:eq-detect}
 1=|g(\rho)|^2\le\frac1{\pi r_0^2}\iint_{|s-\rho|\le r_0}|g(s)|^2\,dA(s)\le\frac1{\pi r_0^2}\iint_{\mathcal R}|g_{d,\xi}(s)|^2\,dA(s).
\end{equation}
The rectangle $\mathcal R$ does not depend on $d$ or on the zero.

\emph{Step 4: the mean square over $d$.} Fix $s\in\mathcal R$, so
$\sigma=\Re s\ge\sigma_0$ and $|s|\le H+3$. By \eqref{an:eq-g},
$|g|^2\le2|\Sigma_1|^2+2|EM|^2$, where $\Sigma_1$ is the sum in
\eqref{an:eq-g}. For the second term, $|M(s)|\le\sum_{k\le Y}k^{-1/2}\le2\sqrt Y$,
and \eqref{an:eq-trunc} gives $|E(s)M(s)|\ll\sqrt X\log X\cdot Z^{-\sigma_0}\sqrt Y$.
Hence
\[
 \sum_{d\le X}|E(s)M(s)|^2\ll X^{2}\,Z^{-2\sigma_0}\,Y\,\log^2X=X^{1-e_1}\log^2X,\qquad
 e_1=(1+\theta)\Bigl(1-\frac\theta4\Bigr)-1-\frac\theta4=\frac{\theta(2-\theta)}4\ \ge\frac{3\theta}8 .
\]
For the first term we apply Theorem~\ref{an:thm-HB} with $Q=\lfloor X\rfloor$ and sum length $\lfloor X\rfloor$ (the letter $N$ of that theorem is not the modulus $N$ of $\xi$), restricted to the characters $\chi_d$, $d\le X$ an odd prime, and with
$a_n=b_n\xi(n)n^{-s}$ for $Y<n\le X$ and $a_n=0$ for $n\le Y$. Then
$|a_n|\le\tau(n)n^{-\sigma_0}$, and
\[
 \sum_{d\le X}|\Sigma_1|^2\ll_\epsilon X^{1+2\epsilon}\,T_Y,\qquad
 T_Y=\sum_{\substack{n_1,n_2>Y\\ n_1n_2=\square}}\frac{\tau(n_1)\tau(n_2)}{(n_1n_2)^{\sigma_0}} .
\]
Two integers have a square product exactly when they have the same
squarefree part: $n_i=ku_i^2$ with $k$ squarefree. With
$\tau(n)\ll_\epsilon n^{\epsilon}$ and $\alpha=2\sigma_0-2\epsilon=1+\theta-2\epsilon$,
where $0<\epsilon\le\theta/4$,
\[
 T_Y\ll_\epsilon\sum_{k\ge1}k^{-\alpha}\Bigl(\sum_{u\ge\max(1,\sqrt{Y/k})}u^{-\alpha}\Bigr)^2
 \ll\sum_{k\le Y}k^{-\alpha}\Bigl(\frac Yk\Bigr)^{1-\alpha}+\sum_{k>Y}k^{-\alpha}\ll Y^{1-\alpha}\log(2Y),
\]
because $\sum_{u\ge U}u^{-\alpha}\ll U^{1-\alpha}$ for $U\ge1$, with a
constant depending on $\theta$. So, for large $X$,
\[
 \sum_{d\le X}|\Sigma_1|^2\ll_\epsilon X^{1+2\epsilon}\,Y^{-\theta+2\epsilon}\log X\le X^{1-\theta^2/4+3\epsilon}.
\]
Taking $\epsilon=\theta^2/32$ and using $\theta^2/8<3\theta/8$,
\begin{equation}\label{an:eq-ms}
 \sum_{\substack{d\le X\\ d\nmid 2N\ \mathrm{prime}}}|g_{d,\xi}(s)|^2\le C_1X^{1-\theta^2/8}\qquad(s\in\mathcal R),
\end{equation}
uniformly in $s\in\mathcal R$.

\emph{Step 5.} Sum \eqref{an:eq-detect} over the primes $d$ in question,
choosing for each one a character $\xi$ and a zero, and bound the result by
the sum over all $\xi$ modulo $N$ and all $d$; exchanging sum and integral
and using \eqref{an:eq-ms}, the number of such primes is at most
$(\pi r_0^2)^{-1}\varphi(N)\,\mathrm{area}(\mathcal R)\,C_1X^{1-\theta^2/8}$.
For bounded $X$ the statement is trivial.
\end{proof}

\begin{corollary}[choice of the prime]\label{an:cor-choice}
Let $N\ge1$, let $a$ be an integer prime to $N$, let $\beta>\frac12$ and
$H\ge0$. There are infinitely many primes $d\equiv a\pmod N$ such that, with
$\anchi=(\cdot/d)$,
\[
 L(s,\xi\anchi)\ne0\qquad\text{for every character $\xi$ modulo $N$ and every $s$ with }\Re s\ge\beta,\ |\Im s|\le H .
\]
\end{corollary}
\begin{proof}
No Dirichlet $L$-function vanishes on $\Re s\ge1$, so we may assume
$\beta<1$. By Dirichlet's theorem in its quantitative form
\cite[Chapter~20]{Davenport} the number of primes $d\le X$ with
$d\equiv a\pmod N$ is asymptotic to $X/(\varphi(N)\log X)$. By
Proposition~\ref{an:prop-density} all of them, apart from $O(X^{1-\eta})$,
have the required property.
\end{proof}

\begin{remark}\label{an:rem-choice}
(1) The range $\beta>\frac12$ is the full range in which
Theorem~\ref{an:thm-HB} saves a power of $X$ in the proof above. A density
theorem for all characters of all moduli up to $X$, of Montgomery's type,
would give a result of this kind only for $\beta$ bounded away from
$\frac12$; the mean value over real characters is essential here.
(2) Nothing beyond non-vanishing on $\Re s\ge1$ is used about the zeros of
the functions $L(s,\xi\chi_d)$, so there is no circularity with
Theorem~\ref{thm:GRH-half}. (3) The corollary is not effective in $d$ as
written, only because the constant of Theorem~\ref{an:thm-HB} is not made
explicit.
\end{remark}


\subsection{The criterion}\label{an:ssec-criterion}
\begin{theorem}[one-sided criterion]\label{an:thm-criterion}
Let $q\ge1$, let $a$ be an integer prime to $2q$, and let $\frac12<\beta<1$.
Suppose that for every prime $d\equiv a\pmod{q}$ with $d\nmid2q$ there is
$x_0(d)$ such that
\[
 R_{d,q}(x)\le x^{\beta-1}\qquad(x\ge x_0(d)).
\]
Then no $L(s,\psi)$, $\psi$ a Dirichlet character modulo $q$, has a zero with
$\Re s>\beta$.
\end{theorem}
\begin{proof}
Let $\Re\rho>\beta$. Since $L(s,\psi)\ne0$ for $\Re s\ge1$ we may assume
$\Re\rho<1$. By Corollary~\ref{an:cor-choice}, with $N=q$ and
$H=|\Im\rho|$, there is a prime $d\equiv a\pmod q$, $d\nmid 2q$, such that
\begin{equation}\label{an:eq-notwist}
 L(s,\psi\anchi)\ne0\qquad\text{for all $\psi$ modulo $q$ and all $s$ with }\Re s\ge\beta,\ |\Im s|\le|\Im\rho| .
\end{equation}
Fix it, and apply Lemma~\ref{an:lem-landau} to $E(v)=R_{d,q}(e^{v})$ and
$\anF=\anF_q$; the hypotheses hold by Lemma~\ref{an:lem-R} and by the
assumption of the theorem, with $v_0=\log x_0(d)$. By \eqref{an:eq-notwist},
$n_q^{K}(s)=0$ in the box, so that by \eqref{an:eq-r}
\[
 r_q(s)=\frac{n_q(s)}{\varphi(q)\,s\,(1-s)}\qquad(\Re s\ge\beta,\ |\Im s|\le|\Im\rho|,\ s\ne1).
\]
The box contains the real segment $[\beta,1)$, and there $s(1-s)>0$ and
$n_q(s)\ge0$, so the residue of $\anF_q$ at a real point $x\in(\beta,1)$ is
$\ge0$; at real points $x\ge1$, $\anF_q$ is holomorphic. So alternative (b)
of Lemma~\ref{an:lem-landau} does not occur, and $\anF_q$ is holomorphic in
$\Re s>\beta$. In particular $r_q(\rho)=0$, and since $\rho$ lies in the
box, $n_q(\rho)=0$: no $L(s,\psi)$, $\psi$ modulo $q$, vanishes at $\rho$.
\end{proof}
Two points of the proof deserve emphasis. The weight $1-\anchi(n)$ brings in
the zeros of the twists $L(s,\psi\anchi)$ with the sign opposite to that of
the zeros of the $L(s,\psi)$; an upper bound for $R_{d,q}$ forbids only a
real pole of negative residue, and this is why the twists must have no zero
on the real segment $[\beta,1)$, which the box contains. Outside the box
cancellation between the two families of zeros is not excluded, but the
conclusion is drawn only at $\rho$, inside the box. The prime $d$ depends on
$\rho$ through $H=|\Im\rho|$, which is why the bound is required for every
prime of the class.


\section{Admissible data}\label{an:sec-data}
\subsection{Parameters, the field \texorpdfstring{$E$}{E} and its places}\label{an:ssec-param}
\begin{definition}\label{an:def-param}
\emph{Parameters} are a pair $(q_0,d)$, where $q_0$ is a positive multiple of
$18480=2^4\cdot3\cdot5\cdot7\cdot11$ and $d$ is a prime with
$d\equiv q_0/2-1\pmod{q_0}$.
\end{definition}
For parameters $(q_0,d)$ we have $q_0\mid2(d+1)$ and $8\mid d+1$, so
$d\equiv7\pmod8$, $d\nmid q_0$, and $K=\mathbb Q(\sqrt{-d})$ has discriminant
$-d$ and character $\anchi=(\cdot/d)$. Every prime $p\mid q_0$ satisfies
$4p\mid d+1$, hence $\anchi(p)=1$: for $p=2$ because $d\equiv7\pmod 8$, and
for odd $p$ by quadratic reciprocity,
$(p/d)=(-1)^{(p-1)/2}(d/p)=(-1)^{(p-1)/2}(-1/p)=1$. Thus
\begin{equation}\label{an:eq-split}
 \text{every prime dividing $q_0$ splits in $K$; a prime $p$ with $\anchi(p)\ne1$ satisfies $p\ge13$.}
\end{equation}
For an integer $s\ge1$ put
\[
 l_0(s)=3^{\lfloor\log_3s\rfloor+4},\qquad T(s)=3\,l_0(s)\,2^{s},\qquad q_1(s)=q_0\,T(s),\qquad q_2(s)=(q_0d)^{s+6}.
\]
Then $27s<l_0(s)\le81s$; the primes dividing $q_1$ are those dividing $q_0$;
$2^{s+4}\mid q_1$; and $q_1d$ divides $q_2$, because $3l_0(s)2^{s}$ divides
$6^{s+5}$ and $6\mid q_0$.

Let $E=E_s=K(\zeta)$, $\zeta$ a primitive $q_1$-th root of unity. Since $d$
is prime to $q_1$, $E$ is the compositum of the linearly disjoint fields $K$
and $F=\mathbb Q(\zeta)$, whose discriminants are coprime. It is a CM field,
abelian over $\mathbb Q$, containing $i$ (as $4\mid q_1$), with complex
conjugation $z\mapsto z^{*}$, and
\[
 [E:K]=\varphi(q_1)=\varphi(q_0)\,T(s)=:f .
\]
Let $\anPl=\anPl_s$ be the set of infinite places of $E$; $\#\anPl=f$. Fix an
embedding $K\subset\mathbb C$. Each place is induced by exactly one embedding
of $E$ extending it, so $\anPl$ is a principal homogeneous space for
$\operatorname{Gal}(E/K)\cong(\mathbb Z/q_1)^{\times}$; fixing a base place we
identify $\anPl$ with the group $(\mathbb Z/q_1)^{\times}$. For an integer $j$
prime to $q_1$ we write $u(j)\in\anPl$ for the place of class $j$, and for a
divisor $m'$ of $q_1$ we write
$\ancl_{m'}\colon\anPl\to(\mathbb Z/m')^{\times}$ for the reduction. (If $j$ is
not prime to $q_1$ we put $u(j)=u(1)$; by \eqref{an:eq-split} this case does
not occur in the sums below.) Let
\[
 \anPa=\{\tau\in\anPl:\ \ancl_{q_0}(\tau)=1\},\qquad\#\anPa=T(s).
\]

\begin{definition}\label{an:def-sparse}
Let $L\ge1$. A set $\mathsf T\subseteq\anPl_s$ is \emph{sparse for the primes
up to $L$} if for every prime $\ell\le L$ and all integers $0<i<j<l_0(s)$
there are no $\sigma,\tau,\rho\in\mathsf T$ and signs
$\epsilon,\epsilon'\in\{\pm1\}$ with
\[
 \ancl_{q'}(\tau/\sigma)=\epsilon\,\ell^{\,i},\qquad \ancl_{q'}(\rho/\sigma)=\epsilon'\,\ell^{\,j}\qquad\text{in }(\mathbb Z/q')^{\times},
\]
where $q'$ is the largest divisor of $q_1(s)$ prime to $\ell$.
\end{definition}

\subsection{The definition}\label{an:ssec-data}
The parameters $(q_0,d)$ determine the constant
\begin{equation}\label{an:eq-CP}
 \anCP=\anCP(q_0,d)\ =\ q_0+2d+6+|\aneu|+C_0(q_0,d)+q_0\log q_0 ,
\end{equation}
where $C_0(q_0,d)\ge4$ is the number defined in \eqref{pz:eq-C0}: it is built from the prime
factors of $q_0$ and from the finitely many normalised lattices with multiplication by the integers
of $K=\mathbb Q(\sqrt{-d})$, and from nothing else (Lemma~\ref{pz:C0}(a)); in particular it is
defined before, and independently of, the data of Definition~\ref{an:def-data}. In the present
sections only the inequality $\anCP\ge q_0+2d+6+|\aneu|$ is used, and the fact that $\anCP$ depends
on $(q_0,d)$ alone; the term $C_0$ is used in the proof of Proposition~\ref{au:kernel}, in the form
$C_0\le\log\log b$, and the term $q_0\log q_0$ only in (S7) of Lemma~\ref{an:lem-sizes}.

\begin{definition}[admissible data]\label{an:def-data}\label{an:data}
Let $(q_0,d)$ be parameters. \emph{Admissible data} $D$ for $(q_0,d)$ consist
of the following objects, subject to the conditions \textup{(D1)--(D8)}.
\begin{enumerate}
\item[(D1)] Integers $1\le x\le x_1\le b$, the \emph{scales}.
\item[(D2)] They satisfy
\begin{gather*}
 10^{12}\,\anCP\le\log\log b,\qquad 100\log\log b\le\log x,\qquad 10^{6}\log x\le\log x_1,\\
 10^{6}\log x_1\le\log b,\qquad 6\log\log x_1\le\log\log b .
\end{gather*}
\item[(D3)] Integers $m$, $r$, $s$: $m$ is a power of $2$ with
$\sqrt x\le m\le2\sqrt x$; $r$ is a power of $3$ with $x^{5}\le r\le3x^{5}$;
$2^{s}=m^{22}$; and $s\le32\log x$. The field of the data is $E=E_s$, with
$q_1=q_1(s)$, $q_2=q_2(s)$, $f=[E:K]$, $\anPl=\anPl_s$.
\item[(D4)] $q_1x_1$ divides $b$.
\item[(D5)] A set $\anPb\subseteq\anPl$ (the \emph{$b$-places}) with
$\ancl_{q_0}(\sigma)=-1$ for all $\sigma\in\anPb$, with
$T(s)\le\#\anPb\cdot\log b$, and sparse for the primes up to $5r(r+1-m)$.
\item[(D6)] The weight $\kappa\colon\anPl\to\mathbb N$: $\kappa(\tau)=b$ for
$\tau\in\anPb$ and $\kappa(\tau)=xb/x_1$ otherwise.
\item[(D7)] Two distinct primes $p_0,p_1\equiv1\pmod{q_2^{2}}$ in the interval
$[b^{10},b^{11}]$ and an integer $\pi_M$ of $K$ with
$\pi_M\overline{\pi_M}=p_0p_1=:M$.
\item[(D8)] A prime $N\equiv1\pmod{q_2^{2}}$ with
\[
 \log N_-\le\log N\le\log N_-+\sqrt{\log b},\qquad
 \log N_-:=\frac mf\sum_{\tau\in\anPl}\log\kappa(\tau)+(m-1)\log\frac{M\,\anrd_E}{2\pi}-\frac{15}2\log b ,
\]
where $\anrd_E=|\operatorname{disc}E|^{1/[E:\mathbb Q]}$ is the root
discriminant, and an integer $\pi_N$ of $E$ with $\pi_N\pi_N^{*}=N$.
\end{enumerate}
We write $a=xb$ and $k=r+1-m$; the sets $\anPa$ and $\anPb$ are disjoint.
\end{definition}

\subsection{Sizes}\label{an:ssec-sizes}
\begin{lemma}\label{an:lem-sizes}
For admissible data the following hold.
\begin{enumerate}
\item[(S1)] $x^{10^{6}}\le x_1$, $x_1^{10^{6}}\le b$,
$4\log x_1\le(\log x_1)^{6}\le\log b$,
$(\log b)^{100}\le x$, $\log b\ge128$, and
$\log x\ge10^{14}(q_0+2d+|\aneu|)$.
\item[(S2)] $x^{11}\le2^{s}\le2^{22}x^{11}$; $l_0(s)\le2592\log x$;
$2m\le r$, so $r/2\le k\le r$.
\item[(S3)] $x^{11}\le T(s)=\#\anPa\le f\le q_1\le x^{12}$, and
$f=\varphi(q_0)\,\#\anPa$; moreover $1\le\#\anPa/\#\anPb\le\log b$.
\item[(S4)] $\anrd_E=\sqrt d\;q_1\prod_{p\mid q_0}p^{-1/(p-1)}$, and
$x^{10}\le\anrd_E\le x^{12}$.
\item[(S5)] $1\le xb/x_1\le\kappa(\tau)\le b\le a\le b^{2}$ for all $\tau$;
$b^{20}\le M\le b^{22}$;
$20\log b\le\log(M\anrd_E/2\pi)\le23\log b$; and
$m\log b\le\log N\le25\,m\log b$.
\item[(S6)] With $\Omega:=e^{\aneu}\cdot2\pi/\anrd_E$: $0<\Omega\le1$ and
$20\log b\le\log(M/\Omega)\le23\log b$.
\item[(S7)] With $n=m/2$ and $\varrho=n-1$ (so that $r=k+2\varrho+1$): $r$ is odd; every prime
dividing $q_1$ divides $q_0$; $q_0^{\,q_0}\le x$; $\varrho\ge10^{5}$;
$k\ge512000\,n^{9}$ and $k\ge1000\,(2\varrho+2)^{9}$; $k\ge2\varrho^{2}+3$; $18\varrho f\le b^{8}$;
$f\le b/1000$; $\anrd_E\le b$; and $N\le b^{60n}$. Moreover, with $\Delta=2r^{8}$:
\begin{itemize}
\item[(a)] $\varrho\ge10^{8}$ (so $n\ge100$); $n^{7}\le k$; $12n^{3}+500\le k$; $k-3n-6>10^{9}n^{3}$;
\item[(b)] $r\ge100$ (so $2r^{8}\ge500r^{3}+10r^{5}$); $3r^{3}f\le r^{7}$; $8nr\le b$; $kf\ge40$; $f\ge4000\,n$;
\item[(c)] $2f\le2^{2\varrho}$ and $2^{-\varrho}\le b^{-6}$;
\item[(d)] $(200\,f\,k\,\anrd_E)^{1000}\le b$ and $(25\,n^{3}f\,\anrd_E)^{1000}\le b$;
\item[(e)] $\sqrt d\le\anrd_E$, and $n,r,f,\sqrt d\le b$; $M\le b^{23}$;
\item[(f)] $\binom{a+k-1}{k-1}^{f}\le b^{2kf}$ (this bounds the dimension $d_\infty=\dim W_{\lambda^X}$ of
\S\ref{ss:au-spaces}, a product of at most $f$ such binomial coefficients);
\item[(g)] $m\Delta+r\le\kappa(\tau)$ for all $\tau$; $k+2b+2m\Delta\le b^{2}$ (this bounds every
entry of $2\varsigma_\tau$, so that $L_\tau:=\max_i|2\varsigma_{\tau,i}|\le b^{2}$); and
$(r+1)^{4}(b^{2}+r+1+2n^{2})\le b^{5}$.
\end{itemize}
\end{enumerate}
\end{lemma}
\begin{proof}
(S1) Exponentiate (D2); $\log\log b\ge10^{12}\anCP\ge10^{12}$ gives
$\log b\ge128$, and $\log x\ge100\log\log b\ge10^{14}\anCP$ with
\eqref{an:eq-CP}. (S2) $2^{s}=m^{22}$ and $\sqrt x\le m\le2\sqrt x$;
$l_0(s)\le81s\le81\cdot32\log x$; $2m\le4\sqrt x\le x^{5}\le r$.
(S3) $T(s)=3l_0(s)2^{s}\ge2^{s}$, $f=\varphi(q_0)T(s)\le q_0T(s)=q_1$, and
$q_1\le q_0\cdot3\cdot2592\log x\cdot2^{22}x^{11}\le x^{12}$ by (S1), which
gives $x\ge q_0\,2^{35}\log x$. The class map $\ancl_{q_0}$ is onto, so each
of its fibres has $T(s)$ elements; $\anPb$ lies in the fibre of $-1$, and
(D5) gives the last statement.
(S4) The Dirichlet characters attached to the abelian field $E$ are the
characters $\psi$ attached to $F$ and their products $\psi\anchi$, of
conductors $\mathfrak f_\psi$ and $d\,\mathfrak f_\psi$. By the
conductor--discriminant formula \cite[Theorem~3.11]{Washington},
$|\operatorname{disc}E|=d^{\varphi(q_1)}\prod_\psi\mathfrak f_\psi^{2}=d^{\varphi(q_1)}|\operatorname{disc}F|^{2}$,
so $\anrd_E=\sqrt d\,\anrd_F$; and
$\anrd_F=q_1\prod_{p\mid q_1}p^{-1/(p-1)}$ \cite[Proposition~2.7]{Washington}.
Since $\prod_{p\mid q_0}p^{1/(p-1)}\le q_0$ we get
$T(s)\le\anrd_E\le\sqrt d\,q_1$; now use $T(s)\ge x^{11}$ and, as in (S3),
$\sqrt d\,q_1\le x^{12}$.
(S5) $x\le x_1\le b$ give $1\le xb/x_1\le b$; $a=xb\le b^{2}$;
$M=p_0p_1$ with (D7). Put $\mu=f^{-1}\sum_\tau\log\kappa(\tau)\in[0,\log b]$
and $\lambda=\log(M\anrd_E/2\pi)$. By (S4) and (S1),
$2\pi\le\anrd_E\le b$, so $20\log b\le\lambda\le23\log b$, and, as $m\ge2$,
\[
 m\log b\le20(m-1)\log b-\tfrac{15}2\log b\le\log N_-=m\mu+(m-1)\lambda-\tfrac{15}2\log b\le24\,m\log b .
\]
With (D8) and $\sqrt{\log b}\le m\log b$ this gives the bounds for $\log N$.
(S6) By (S4) and (S1), $\anrd_E\ge x^{10}\ge2\pi e^{|\aneu|}$, so
$\Omega\le1$; then $\log(M/\Omega)\ge\log M\ge20\log b$ and
$\log(M/\Omega)\le22\log b+12\log x+|\aneu|\le23\log b$.
(S7) $r$ is a power of $3$. By definition (\S\ref{an:ssec-param}), $q_1/q_0=T(s)=3\,l_0(s)\,2^{s}$
has only the prime factors $2$ and $3$, which divide $q_0$. By (S1) and \eqref{an:eq-CP},
$\log x\ge10^{14}\anCP\ge q_0\log q_0$. Next $n=m/2\ge\sqrt x/2$ and $n\le\sqrt x$, and
$x\ge(\log b)^{100}\ge128^{100}$, so $\varrho\ge10^{5}$; by (S2), $k\ge r/2\ge x^{5}/2$, while
$512000\,n^{9}\le512000\,x^{9/2}$ and $1000\,(2\varrho+2)^{9}=1000\,m^{9}\le512000\,x^{9/2}$, and
$512000\,x^{9/2}\le x^{5}/2$ because $\sqrt x\ge1024000$; also $2\varrho^{2}+3\le2x+3\le x^{5}/2$.
By (S3) and (S1), $f\le x^{12}\le b/1000$ and $18\varrho f\le18x^{13}\le b^{8}$. By (S5),
$\log N\le25\,m\log b=50\,n\log b$, and by (S4), $\anrd_E\le x^{12}\le b$.

For (a)--(g) we use only: $x\ge(\log b)^{100}\ge128^{100}$ and $x^{10^6}\le x_1\le b$ (S1), so that
$\log b\le x^{1/100}$ and $x^{c}\le b$ for $c\le10^{6}$; $\sqrt x/2\le n\le\sqrt x$ and $\varrho\ge\sqrt x/2-1$;
$x^{5}\le r\le3x^{5}$ and $x^{5}/2\le k\le r$ (S2); $x^{11}\le f\le x^{12}$ (S3); $x^{10}\le\anrd_E\le x^{12}$ (S4);
$xb/x_1\le\kappa(\tau)\le b$ and $M\le b^{22}$ (S5); $a=xb$.
(a) $\varrho\ge128^{50}/2-1\ge10^{8}$. Since $x^{7/2}\ge128^{350}$, each of $n^{7}\le x^{7/2}$,
$12n^{3}+500\le13x^{3/2}$ and $3n+6+10^{9}n^{3}\le2\cdot10^{9}x^{3/2}$ is $<x^{5}/2\le k$.
(b) $r\ge x^{5}\ge100$, and then $500r^{3}+10r^{5}\le510r^{5}\le2r^{8}$; $3f\le3x^{12}\le x^{20}\le r^{4}$; $8nr\le24x^{11/2}\le b$; $kf\ge k\ge40$;
$4000n\le4000\sqrt x\le x^{11}\le f$.
(c) $2^{2\varrho}\ge2^{\sqrt x-2}\ge2x^{12}\ge2f$, as $\sqrt x\ge128^{50}$; and
$\varrho\log2\ge(\sqrt x/2-1)\log2\ge6x^{1/100}\ge6\log b$.
(d) $200fk\,\anrd_E\le600\,x^{29}\le x^{30}$ and $25n^{3}f\anrd_E\le25x^{51/2}\le x^{26}$; and
$x^{30000}\le x^{10^{6}}\le b$.
(e) By (S4), $\anrd_E=\sqrt d\,q_1\prod_{p\mid q_0}p^{-1/(p-1)}\ge\sqrt d\,q_1/q_0\ge\sqrt d$, since
$\prod_{p\mid q_0}p^{1/(p-1)}\le q_0\le q_1$. Then $\sqrt d\le x^{12}\le b$, and $n,r,f\le x^{12}\le b$;
$M\le b^{22}\le b^{23}$.
(f) $\binom{a+k-1}{k-1}=\prod_{i=1}^{k-1}\frac{a+i}{i}\le(a+1)^{k-1}\le(2xb)^{k}\le b^{2k}$, as $2x\le b$.
(g) $m\Delta+r\le2\sqrt x\cdot2(3x^{5})^{8}+3x^{5}\le x^{42}\le b^{1-10^{-6}}\le xb/x_1\le\kappa(\tau)$, since
$x_1\le b^{10^{-6}}$. Hence $k+2b+2m\Delta\le3b\le b^{2}$, and the entries of $2\varsigma_\tau$, namely
$2(\frac k2-j)$ for $1\le j\le k-1$ and $-2\kappa(\tau)-2j\Delta$ for $1\le j\le m$
(\S\ref{ss:au-spaces}), have absolute value at most $k+2b+2m\Delta$. Finally
$(r+1)^{4}(b^{2}+r+1+2n^{2})\le(4x^{5})^{4}\cdot2b^{2}\le b^{5}$.
\end{proof}

\begin{remark}[sizes]\label{an:rem-sizes}
The data live on three scales: $\log x_1$ and $\log b$ exceed every fixed
power of $\log x$, and $\log b\ge10^{6}\log x_1$, while $m$, $r$, $k$, $f$,
$q_1$, $\anrd_E$ are powers of $x$ with the exponents $\frac12$, $5$, $5$,
between $11$ and $12$. For the data constructed in
Theorem~\ref{an:thm-data}, $\log x_1\approx x^{\varepsilon/6}$ and
$\log b\approx x^{\varepsilon}$. Lemma~\ref{an:lem-sizes} contains every
comparison of sizes used in this section and in
Section~\ref{an:sec-main}, and it contains the inequalities
\eqref{eq:au-scales} and \eqref{eq:au-sizes} quoted at the beginning of
Section~\ref{sec:automorphic}; all hold with a large margin, and the
constants $10^{6}$, $10^{12}$, $100$ in (D2) are far from critical for them.
\end{remark}

\subsection{A sparse set of places}\label{an:ssec-sparse}
The set $\anPb$ is obtained from Behrend's construction \cite{Behrend1946} of
large sets without three-term progressions, in a form that also excludes the
weighted relations $(\alpha+\alpha')e_2=\alpha e_1+\alpha'e_3$.

\begin{lemma}\label{an:lem-behrend}
Let $n,H,\Lambda_0\ge1$ and $o$ be integers with $(\Lambda_0H)^{n}\le o$.
There is a set $S\subseteq\{0,1,\dots,o-1\}$ with $\#S\ge H^{n-2}/n$ such
that for all $e_1,e_2,e_3\in S$ and all integers $\alpha,\alpha'\ge1$ with
$\alpha+\alpha'\le\Lambda_0$,
\[
 (\alpha+\alpha')\,e_2\equiv\alpha e_1+\alpha'e_3\pmod o\quad\Longrightarrow\quad e_1=e_2=e_3 .
\]
\end{lemma}
\begin{proof}
The function $v\mapsto\sum_lv_l^2$ takes at most $nH^{2}$ values on
$\{0,\dots,H-1\}^{n}$, so one of its fibres, $\Sigma$, has at least
$H^{n}/(nH^{2})$ elements. Let $\beta(v)=\sum_{l<n}v_l(\Lambda_0H)^{l}$ and
$S=\beta(\Sigma)$; $\beta$ is injective on vectors with entries in
$[0,\Lambda_0H)$ and maps them into $[0,(\Lambda_0H)^{n})\subseteq[0,o)$. If
$u,v,w\in\Sigma$ satisfy the congruence for $e_1=\beta(u)$, $e_2=\beta(v)$,
$e_3=\beta(w)$, then the vectors $(\alpha+\alpha')v$ and
$\alpha u+\alpha'w$ have entries in $[0,\Lambda_0H)$ and images under $\beta$
congruent modulo $o$, hence equal; so
$(\alpha+\alpha')v=\alpha u+\alpha'w$. As $u,v,w$ lie on one sphere, strict
convexity gives $u=v=w$: indeed
$\alpha\alpha'|u-w|^2=(\alpha+\alpha')(\alpha|u|^2+\alpha'|w|^2)-|(\alpha+\alpha')v|^2=0$.
\end{proof}

\begin{proposition}\label{an:prop-sparse}
Let $(q_0,d)$ be parameters and let $s,n,H,L\ge1$ be integers with
\[
 (2\,l_0(s)\,H)^{n}\le2^{s}\qquad\text{and}\qquad(L+1)\,l_0(s)\le2^{s}.
\]
There is a set $\mathsf T\subseteq\anPl_s$ with $\ancl_{q_0}(\sigma)=-1$ for
all $\sigma\in\mathsf T$, with $\#\mathsf T\ge H^{n-2}/n$, and sparse for the
primes up to $L$.
\end{proposition}
\begin{proof}
Write $l_0=l_0(s)$ and $q_0=2^{v}w$ with $w$ odd; $v\ge4$. Let
$g=u(1+3l_0q_0)$. Since $q_1=2^{s+v}\cdot3l_0w$ and $1+3l_0q_0=1+2^{v}w'$ with
$w'$ odd, the class $\ancl_{2^{s+v}}(g)$ has order exactly $2^{s}$; and
$\ancl_{3l_0q_0}(g)=1$. Let $S$ be the set of
Lemma~\ref{an:lem-behrend} for $o=2^{s}$ and $\Lambda_0=2l_0$, and
\[
 \mathsf T=\{\,u(-1)\,g^{e}:\ e\in S\,\}.
\]
The map $e\mapsto u(-1)g^{e}$ is injective on $[0,2^{s})$, so
$\#\mathsf T=\#S$, and $\ancl_{q_0}=-1$ on $\mathsf T$.

Let $\ell\le L$ be a prime, $0<i<j<l_0$, and suppose that
$\sigma=u(-1)g^{e_1}$, $\tau=u(-1)g^{e_2}$, $\rho=u(-1)g^{e_3}$ in
$\mathsf T$ satisfy the two relations of Definition~\ref{an:def-sparse}.

\emph{The case $\ell=2$.} Then $3l_0=3^{c}$ divides $q'$, and
$\ancl_{3^{c}}(\tau/\sigma)=\ancl_{3^c}(g)^{e_2-e_1}=1$. So
$\epsilon2^{i}\equiv1$ and $4^{i}\equiv1\pmod{3^{c}}$. The order of $4$
modulo $3^{c}$ is $3^{c-1}=l_0$, so $l_0\mid i$, contrary to $0<i<l_0$.

\emph{The case $\ell$ odd.} Then $2^{s+v}$ divides $q'$. Squaring the two
relations modulo $2^{s+v}$ removes the signs:
$\ancl(g)^{2(e_2-e_1)}=\ell^{2i}$ and $\ancl(g)^{2(e_3-e_1)}=\ell^{2j}$, hence
$\ancl(g)^{2j(e_2-e_1)}=\ell^{2ij}=\ancl(g)^{2i(e_3-e_1)}$. As $\ancl(g)$ has
order $2^{s}$,
\[
 2j\,e_2\equiv(2j-2i)\,e_1+2i\,e_3\pmod{2^{s}},
\]
a relation of Lemma~\ref{an:lem-behrend} with $\alpha=2j-2i\ge1$,
$\alpha'=2i\ge1$ and $\alpha+\alpha'=2j\le2l_0$. Therefore $e_1=e_2=e_3$,
so $\tau=\sigma$ and $\ell^{2i}\equiv1\pmod{2^{s+1}}$. Now
$v_2(\ell^{2i}-1)=v_2(\ell-1)+v_2(\ell+1)+v_2(i)$, and one of
$v_2(\ell\pm1)$ equals $1$ while the other is at most $\log_2(\ell+1)$; so
$s+1\le1+\log_2((\ell+1)i)$, that is
$2^{s}\le(\ell+1)\,i<(L+1)\,l_0\le2^{s}$, a contradiction.
\end{proof}

\subsection{Existence}\label{an:ssec-exist}
\begin{lemma}[primes in windows]\label{an:lem-windows}
Let $A\ge1$. There is $w_0(A)$ such that for all real $W\ge w_0(A)$, all
$U\in[W,W^{A}]$ and all integers $1\le Q\le W^{A}$ there is a prime
$p\equiv1\pmod Q$ with $U\le\log p\le U+W$.
\end{lemma}
\begin{proof}
Put $y_1=e^{U}$, $y_2=e^{U+W}$; then
$Q\le W^{A}\le(\log y_1)^{A}\le(\log y_2)^{A}$. By the theorem of Siegel and
Walfisz \cite[Chapter~22]{Davenport}, applied at $y_1$ and at $y_2$,
\[
 \psi(y_2;Q,1)-\psi(y_1;Q,1)=\frac{y_2-y_1}{\varphi(Q)}+O_A\bigl(y_2\exp(-c_A\sqrt U)\bigr)\ \ge\ \frac{y_2}{2W^{A}}-C_A\,y_2\exp(-c_A\sqrt W)\ \ge\ \frac{y_2}{4W^{A}}
\]
for $W\ge w_0(A)$. The prime powers $p^{k}\le y_2$, $k\ge2$, contribute at
most $O(\sqrt{y_2}\,(\log y_2)^2)=O(\sqrt{y_2}\,W^{2A})$ to the left side,
which is smaller than $y_2/(4W^{A})$ for large $W$ because $\sqrt{y_2}\ge e^{W}$. So there
is a prime $p\equiv1\pmod Q$ with $y_1<p\le y_2$.
\end{proof}

\begin{theorem}[existence of admissible data]\label{an:thm-data}
Let $(q_0,d)$ be parameters and $0<\varepsilon<\frac1{100}$. There is
$x_0=x_0(q_0,d,\varepsilon)$ such that for every integer $x\ge x_0$ there are
admissible data $D$ for $(q_0,d)$ with scale $x$ and with
\[
 x_1=\bigl\lfloor\exp(x^{\varepsilon/6})\bigr\rfloor,\qquad x^{\varepsilon}\le\log b\le x^{\varepsilon}+1,\qquad\sqrt x\le m\le2\sqrt x .
\]
\end{theorem}
\begin{proof}
Let $W=x^{\varepsilon/6}$ and $n=\lceil44/\varepsilon\rceil$. ``Large''
means large in terms of $(q_0,d,\varepsilon)$; recall that $\anCP$ and the
class number $h_K$ depend on $(q_0,d)$ only.

\emph{Ranks and field.} Let $m$ be the least power of $2$ with $m\ge\sqrt x$,
$r$ the least power of $3$ with $r\ge x^{5}$, and $s=22\log_2m$. Then
$m\le2\sqrt x$, $r\le3x^{5}$, $2^{s}=m^{22}\in[x^{11},2^{22}x^{11}]$ and
$s\le11\log_2x+22\le32\log x$. This is (D3). As in
Lemma~\ref{an:lem-sizes}, $l_0(s)\le2592\log x$ and $q_1\le x^{12}$; and
$q_2^{2}=(q_0d)^{2(s+6)}\le x^{c_2}$ with $c_2=66\log(q_0d)$.

\emph{Scales.} Let $x_1=\lfloor e^{W}\rfloor$ and let $b$ be the least
multiple of $q_1x_1$ exceeding $\exp(x^{\varepsilon})$. Since
$q_1x_1\le x^{12}e^{W}\le\exp(x^{\varepsilon})$, we have
$\exp(x^\varepsilon)<b\le2\exp(x^{\varepsilon})$, so
$x^{\varepsilon}\le\log b\le x^{\varepsilon}+1$; also
$W-1\le\log x_1\le W$ and $x\le x_1\le b$. This gives (D1) and (D4). For
(D2): $\log\log b\ge\varepsilon\log x\ge10^{12}\anCP$;
$100\log\log b\le100(\varepsilon\log x+1)\le\log x$ because
$\varepsilon<\frac1{100}$; $10^{6}\log x\le W-1\le\log x_1$;
$10^{6}\log x_1\le10^{6}x^{\varepsilon/6}\le x^{\varepsilon}\le\log b$;
and $6\log\log x_1\le6\log W=\varepsilon\log x\le\log\log b$.

\emph{Places.} Let $H=\lfloor2^{s/n}/(2l_0(s))\rfloor$. As
$2^{s/n}\ge x^{11/n}$ and $l_0(s)\le2592\log x$, $H\ge2^{s/n}/(4l_0(s))\ge1$,
and $(2l_0(s)H)^{n}\le2^{s}$. Let $L=5r(r+1-m)\le45x^{10}$; then
$(L+1)l_0(s)\le46x^{10}\cdot2592\log x\le x^{11}\le2^{s}$.
Proposition~\ref{an:prop-sparse} gives $\anPb:=\mathsf T$ with
$\ancl_{q_0}=-1$, sparse for the primes up to $L$, and
\[
 \#\anPb\ \ge\ \frac{H^{n-2}}n\ \ge\ \frac{2^{s}}{2^{2s/n}\,n\,(4l_0(s))^{n-2}} .
\]
Since $2s/n\le s\varepsilon/22$, $2^{2s/n}\le(2^{22}x^{11})^{\varepsilon/22}\le2x^{\varepsilon/2}$,
and $T(s)=3l_0(s)2^{s}$, we get
\[
 \frac{T(s)}{\#\anPb}\le6\,n\,(4\cdot2592\log x)^{n-1}\,x^{\varepsilon/2}\le x^{\varepsilon}\le\log b .
\]
This is (D5); (D6) defines $\kappa$, and $1\le\kappa\le b$.

\emph{Windows.} Let $A=6(c_2+1)/\varepsilon$, so that $W^{A}=x^{c_2+1}\ge\max(q_2^{2},x)$.
By Lemma~\ref{an:lem-windows} with $Q=q_2^{2}$: for every $U$ with
$W\le U\le x$ there is a prime $p\equiv1\pmod{q_2^{2}}$ with
$U\le\log p\le U+W$. We use it for $U\ge10\log b\ge W$ and
$U\le24m\log b\le48\sqrt x(x^{\varepsilon}+1)\le x$.

\emph{The primes $p_0,p_1$.} For $i=0,\dots,h_K$ take a prime
$p^{(i)}\equiv1\pmod{q_2^{2}}$ with
$\log p^{(i)}\in[10\log b+2iW,\,10\log b+(2i+1)W]$. These $h_K+1$ primes are
distinct and lie in $[b^{10},b^{11}]$, because $(2h_K+1)W\le\log b$. Each is
$\equiv1\pmod d$, so it splits in $K$: $(p^{(i)})=\mathfrak p_i\overline{\mathfrak p_i}$.
Two of the ideal classes of the $\mathfrak p_i$ coincide, say for $i\ne j$;
then $\mathfrak p_i\overline{\mathfrak p_j}=(\pi_M)$ is principal, and
$\pi_M\overline{\pi_M}=p^{(i)}p^{(j)}$. Put $p_0=p^{(i)}$, $p_1=p^{(j)}$.
This is (D7).

\emph{The level $N$.} With $M=p_0p_1$ and $\log N_-$ as in (D8), the proof of
Lemma~\ref{an:lem-sizes}(S4),(S5) uses only (D1)--(D7) and shows
$m\log b\le\log N_-\le24m\log b$. Take a prime $N\equiv1\pmod{q_2^{2}}$ with
$\log N_-\le\log N\le\log N_-+W$; and $W\le\sqrt{\log b}$. Since $4\mid q_2$,
$N\equiv1\pmod4$, so $N=A_0^2+B_0^2$ with integers $A_0,B_0$, and
$\pi_N=A_0+B_0i$ is an integer of $E$ with $\pi_N\pi_N^{*}=N$. This is (D8).
\end{proof}
\begin{remark}\label{an:rem-ineff}
The threshold $x_0(q_0,d,\varepsilon)$ is not effective, because of the
theorem of Siegel and Walfisz. This does not affect what follows.
\end{remark}


% sec-automorphic.tex
% Macros of this section (prefix \au). Only amsmath/amssymb/amsthm are used.
% Reference to the definition of admissible data in the analytic section (one place to change).
\providecommand{\auDataRef}{\ref{an:data}}
\providecommand{\auO}{\mathcal O}
\providecommand{\auI}{\mathcal I}
\providecommand{\auQlb}{\overline{\Q}_\ell}
\providecommand{\auAut}{\operatorname{Aut}}
\providecommand{\auNm}{\operatorname{Nm}}
\providecommand{\auGram}{\operatorname{Gram}}
\providecommand{\aucol}{\operatorname{col}}
\providecommand{\auU}{\mathrm U}
\providecommand{\auGL}{\mathrm{GL}}
\providecommand{\auHom}{\operatorname{Hom}}
\providecommand{\auEnd}{\operatorname{End}}
\providecommand{\auA}{\mathbb A}
\providecommand{\auF}{\mathbb F}
\providecommand{\auPl}{\Phi}
\providecommand{\auPa}{\Phi_a}
\providecommand{\auPb}{\Phi_b}
\providecommand{\auS}{\mathcal S}
\providecommand{\auSram}{\mathcal S^{\mathrm{ram}}}
\providecommand{\auT}{\mathcal T}
\providecommand{\auAz}{\mathcal A_0}
\providecommand{\auAm}{\mathcal A_m}
\providecommand{\auG}{\overline G}
\providecommand{\auP}{\mathcal P}
\providecommand{\aupp}{\mathfrak p}
\providecommand{\auLat}{\mathcal L}
\providecommand{\audT}{d_{\auT}}
\providecommand{\auDX}{D_X}
\providecommand{\augain}{\operatorname{gain}}
\providecommand{\auerr}{\operatorname{err}}
\providecommand{\auhmain}{h_{\mathrm{main}}}
\providecommand{\auhdir}{h_{\mathrm{dir}}}
\providecommand{\aurem}{\operatorname{rem}}
\providecommand{\ausinf}{s_\infty}
\providecommand{\ausfin}{s_{\mathrm{fin}}}
\providecommand{\ausmain}{s_{\mathrm{main}}}
\providecommand{\aueu}{\varepsilon_{\mathrm u}}
\providecommand{\aurd}{\mathrm{rd}_E}
\providecommand{\aucB}{c_B}
\providecommand{\aubud}{\mathrm{bud}}
\providecommand{\auchi}{\chi_K}
\providecommand{\aubar}[1]{\overline{#1}}
\providecommand{\auf}[1]{\texttt{#1}}
\providecommand{\ausource}[1]{}

\section{Integral forms on a chain of unitary groups}\label{sec:automorphic}

This section states and proves the three theorems on automorphic forms
from which the one-sided bound for the prime sum is deduced in the analytic section. The objects
are algebraic modular forms on definite unitary groups over a CM field $E$ of large degree, a
chain of $m$ theta lifts, and the lattice of forms all of whose lifts have integral values. The
three theorems say that this lattice contains many short vectors (Theorems~\ref{thm:srcA}
and~\ref{thm:srcB}) and that it has no short vectors beyond what its covolume allows
(Theorem~\ref{thm:srcC}); the two amounts differ by the prime sum.

\emph{Where the proofs are.} Each statement of this section is either proved where it stands or
followed by a short proof naming the results from which it follows and the sections in which these
are proved. The table of \S\ref{ss:au-remains} lists, for each input of Theorems~\ref{thm:srcA},
\ref{thm:srcB} and~\ref{thm:srcC}, the section of its proof and the published results it uses.

\subsection{Notation and the data}\label{ss:au-notation}

Throughout, $D$ denotes admissible data in the sense of Definition~\auDataRef. We recall the
notation. The parameters are $q_0$ (a multiple of $18480$) and a prime $d\equiv q_0/2-1
\pmod{q_0}$; $K=\Q(\sqrt{-d})$, with quadratic character $\auchi$; all primes dividing $q_0$ split
in $K$. The scales are integers $1<x<x_1<b$ (condition (D1) reads $1\le x\le x_1\le b$, and the
inequalities are strict by (D2): $\log x\ge100\log\log b>0$, $\log x_1\ge10^6\log x>\log x$ and
$\log b\ge10^6\log x_1>\log x_1$) with
\begin{equation}\label{eq:au-scales}
(\log b)^{100}\le x,\qquad x^{10^6}\le x_1,\qquad x_1^{10^6}\le b,\qquad(\log x_1)^6\le\log b,
\end{equation}
and $a=xb$. The number of lifts $m$ is a power of $2$ with $\sqrt x\le m\le2\sqrt x$; the rank $r$
is a power of $3$ with $x^5\le r\le3x^5$; $k=r+1-m$, $t=5r$, $n_j=k+j(t+1)$ and $\Delta=2r^8$. The
integer $q_1=q_1(s)$ is the one defined before Definition~\ref{an:data}; its prime divisors are
those of $q_0$, so they split in $K$, and $d\nmid q_1$. The
field is $E=K(\zeta_{q_1})$, a CM field with maximal totally real subfield $E^+$ and complex
conjugation $c$; it is abelian over $\Q$, contained in $\Q(\zeta_{q_1d})$, contains a primitive
$5040$-th root of unity, and $E/E^+$ is unramified at all finite places. The set $\auPl$ of
infinite places of $E$ has $f$ elements and is identified with $\operatorname{Gal}(E/K)=
(\Z/q_1)^\times$; for an integer $\nu$ prime to $q_1$ we write $u(\nu)\in\auPl$ for its class.
The $a$-places $\auPa$ are the places of class $1$ modulo $q_0$; the $b$-places $\auPb$ form a
sparse set of places of class $-1$ modulo $q_0$ with $\#\auPa\le\#\auPb\cdot\log b$. The weight is
$\kappa(\tau)=b$ for $\tau\in\auPb$ and $\kappa(\tau)=a/x_1$ otherwise, and
$\kappa_j(\tau)=\kappa(\tau)+(j+\tfrac12)\Delta$. The level primes are $p_0\ne p_1$ in
$[b^{10},b^{11}]$ and $N$, all $\equiv1\pmod{q_2^2}$ with $q_1d\mid q_2$, so that they split
completely in $E$; $M=p_0p_1$; $\pi_M\in\auO_K$ and $\pi_N\in\auO_E$ satisfy
$\pi_M\aubar{\pi_M}=M$, $\pi_N\aubar{\pi_N}=N$. The primes of $E$ dividing $\pi_M$ are the
\emph{chosen places}; there are $f$ of them above each of $p_0$, $p_1$, and $f$ primes divide
$\pi_N$. With $\aurd$ the root discriminant of $E$ and
\[
\Omega=\exp\Bigl(\gamma-\Rel\frac{L'}{L}(1,\auchi)\Bigr)\cdot\frac{2\pi}{\aurd},\qquad
c_1=-1-\gamma+\Rel\frac{L'}{L}(1,\auchi),
\]
the level $N$ satisfies
\begin{equation}\label{eq:au-level}
\log N_-\le\log N\le\log N_-+\sqrt{\log b},\qquad
\log N_-=\frac mf\sum_{\tau\in\auPl}\log\kappa(\tau)+(m-1)\log\frac{M\,\aurd}{2\pi}-\frac{15}{2}\log b .
\end{equation}
We use freely the following consequences of the definition of the data:
\begin{equation}\label{eq:au-sizes}
\begin{gathered}
f\le x^{12},\qquad x^{10}\le\aurd\le x^{12},\qquad 20\log b\le\log\frac M\Omega\le23\log b,\\
\Bigl|\log\frac{M\aurd}{2\pi}\Bigr|\le23\log b,\qquad m\log b\le\log N\le25\,m\log b .
\end{gathered}
\end{equation}
The \emph{error unit} is $\aueu=f\,a\,(\log b)^3$. Every quantity that is polynomial in $x$ and in
$\log b$ is negligible against $b$; inequalities of this kind are used without comment and are
called \emph{comparisons of sizes}.

\subsection{Heights and a product formula with a twist for each prime}\label{ss:au-heights}

A \emph{place datum} is a pair $\iota=(\ell,\iota)$ of a prime $\ell=\ell(\iota)$ and a field
isomorphism $\iota\colon\C\to\auQlb$; $|z|_\iota$ is the $\ell$-adic absolute value of $\iota(z)$,
with $|\ell|_\iota=\ell^{-1}$, and $\auO_\iota=\{z\in\C:|z|_\iota\le1\}$, a valuation ring of $\C$.
We write $\auI$ for the set of place data and $\aubar\iota$ for $\iota$ composed with complex
conjugation. A \emph{family} is a function $y\colon\auI\to\C$.

\begin{definition}\label{def:au-height}
A family $y$ \emph{has height at most $c\in\R$} if there are $H\in\C^\times$, $\alpha\in\R$ and
reals $h_\ell\ge0$, one for each prime, such that $|\sigma H|\le e^\alpha$ for all
$\sigma\in\auAut(\C)$, $|H\,y_\iota|_\iota\le\ell(\iota)^{h_{\ell(\iota)}}$ for all $\iota$, and
$\alpha+\sum_{\ell\in T}h_\ell\log\ell\le c$ for every finite set $T$ of primes.
\end{definition}

The multiplier $H$ is the same for all $\iota$, whereas $y_\iota$ may depend on $\iota$. Heights
are subadditive under products of families: multiply the multipliers and add the exponents. If
$N_\iota$ is a positive integer prime to $\ell(\iota)$ and $e_\iota\in\Q$, then
$\iota\mapsto N_\iota^{e_\iota}y_\iota$ has height at most $c$ when $y$ has. If $H\in\C^\times$
has $|\sigma H|\le e^{c}$ for all $\sigma$, then the constant family $H^{-1}$ has height at most
$c$.

\begin{proposition}[Product formula]\label{au:product}
If the constant family $1$ has height at most $c$, then $c\ge0$. Consequently, if $z\in\C^\times$,
the constant family $z$ has height at most $c$, and $|\sigma z|\le e^{c'}$ for all
$\sigma\in\auAut(\C)$, then $c+c'\ge0$.
\end{proposition}
\begin{proof}
Let $(H,\alpha,h)$ be as in Definition~\ref{def:au-height} for $y=1$. Then
$|H|_\iota\le2^{h_2}$ for every $\iota$ with $\ell(\iota)=2$. A transcendental number is carried
by suitable automorphisms of $\C$ to $2^{-n}H'$ for any $n$ and any transcendental $H'$; so $H$ is
algebraic. Let $n=[\Q(H):\Q]$ and $q=\auNm_{\Q(H)/\Q}H\in\Q^\times$. Every embedding of $\Q(H)$
in $\auQlb$, respectively in $\C$, extends to an isomorphism $\C\to\auQlb$, respectively to an
automorphism of $\C$. Hence $|q|_\ell\le\ell^{nh_\ell}$ for every $\ell$ and $|q|\le e^{n\alpha}$.
The product formula for $q$ gives $0=\log|q|+\sum_\ell\log|q|_\ell\le n(\alpha+\sum_{\ell\in T}
h_\ell\log\ell)\le nc$, with $T$ the set of primes dividing the numerator or the denominator of
$q$. For the second statement, the constant family $1=z^{-1}z$ has height at most $c'+c$.
\end{proof}

\begin{definition}[Twisted lattice; index; covolume]\label{def:au-lattice}
Let $V$ be a finite-dimensional complex inner product space. A \emph{twisted lattice} $\auLat$ on
$V$ consists of a set $\Xi$ of \emph{twists} with a distinguished element $1$; for each
$\xi\in\Xi$ an inner product space $H_\xi$ and an injective linear map
$\Theta_\xi\colon V\to H_\xi$; a relation ``$\xi$ is allowed at $\ell$''; for each $\xi$ and
$\iota$ an $\auO_\iota$-submodule $\Lambda_{\xi,\iota}\subseteq H_\xi$; and a linear isomorphism
$\operatorname{coord}$ of $H_1$ onto a subspace of $\C^{\mathsf C}$, $\mathsf C$ finite, with a
positive definite matrix $G\in M_{\mathsf C}(\Q)$ such that
$\langle\varphi,\varphi'\rangle=\aubar{\operatorname{coord}(\varphi)}^{\,t}G\operatorname{coord}(\varphi')$
on $H_1$. The condition is
\begin{equation}\label{eq:au-pairing}
\langle\varphi,\varphi'\rangle\in\auO_\iota\qquad\text{for }\varphi\in\Lambda_{\xi,\aubar\iota},\
\varphi'\in\Lambda_{\xi,\iota},\ \xi\text{ allowed at }\ell(\iota).
\end{equation}
\begin{enumerate}
\item[(a)] The module of \emph{vectors of $V$ integral at $\iota$} is
$\auLat_\iota=\{y\in V:\Theta_\xi(y)\in\Lambda_{\xi,\iota}\text{ for every }\xi\text{ allowed at }
\ell(\iota)\}$.
\item[(b)] A finite family $y=(y_1,\dots,y_n)$ in $V$ \emph{has index at most $\delta$} if there
are reals $e_\ell\ge0$ with $\sum_{\ell\in T}e_\ell\log\ell\le\delta$ for all finite $T$ such
that for every $\iota$ there is $u\in M_n(\auO_\iota)$ with $|\det u|_\iota\ge
\ell(\iota)^{-e_{\ell(\iota)}}$ and $\sum_ku_{ik}y_k\in\auLat_\iota$ for all $i$.
\item[(c)] For $\sigma\in\auAut(\C)$ and $y\in V$ let $\|\Theta_1y\|_\sigma^2$ be the value of the
form $G$ at the vector $\sigma(\operatorname{coord}(\Theta_1y))$, where $\sigma$ is applied to
each coordinate. The family $y$ \emph{has covolume at most $c$} if, for every $\sigma$, the Gram
determinant for $G$ of the vectors $\sigma(\operatorname{coord}(\Theta_1y_i))$ is at most $e^c$.
\item[(d)] $Z(\xi)=\det(\Theta_\xi^*\Theta_\xi)>0$.
\end{enumerate}
\end{definition}

Thus the matrix $u$ in (b) may depend on $\iota$ but serves all twists allowed at $\ell(\iota)$,
and the ``covolume'' in (c) is a bound for a squared volume, read for the trivial twist and at all
conjugates. The index of a union of two families is at most the sum of the indices (block diagonal
matrices). By Hadamard's inequality, if $\|\Theta_1y_i\|_\sigma^2\le e^h$ for all $i$ and $\sigma$,
then $y$ has covolume at most $nh$.

\begin{theorem}[A twist for each prime]\label{au:twistprime}
Let $\auLat$ be a twisted lattice on $V$ and let $y_1,\dots,y_n,z_1,\dots,z_{n'}$ be a basis of
$V$. For each $\iota\in\auI$ let $\theta_\iota\in\Xi$ be allowed at $\ell(\iota)$. Let
$\gamma\colon\Xi\to\C^\times$ and $c,\delta,b_{\mathrm{short}},b_{\mathrm{den}},b_{\mathrm{vol}}
\in\R$ be such that
\begin{enumerate}
\item[(i)] $y$ has covolume at most $c$ and index at most $\delta$;
\item[(ii)] $\gamma(\xi)\Theta_\xi(z_i)\in\Lambda_{\xi,\iota}$ for all $i$, all $\iota$ and all
$\xi$ allowed at $\ell(\iota)$;
\item[(iii)] $\|\Theta_1z_i\|_\sigma^2\le\exp(b_{\mathrm{short}})$ for all $i$ and $\sigma$;
\item[(iv)] the family $\iota\mapsto|\gamma(\theta_\iota)|^{-2}$ has height at most
$b_{\mathrm{den}}$, and $\iota\mapsto Z(1)/Z(\theta_\iota)$ has height at most $b_{\mathrm{vol}}$.
\end{enumerate}
Then $0\le c+2\delta+n'(b_{\mathrm{short}}+b_{\mathrm{den}})+b_{\mathrm{vol}}$.
\end{theorem}
\begin{proof}
Write $w=(y_1,\dots,z_{n'})$ and $\Gamma(\xi)=\det\auGram(\Theta_\xi w)$. Since $w$ is a basis,
$\Gamma(\xi)=Z(\xi)\det\auGram(w)$, so that for every $\iota$
\begin{equation}\label{eq:au-factor}
\Gamma(1)=\Bigl(|\gamma(\theta_\iota)|^{2n'}\,\Gamma(\theta_\iota)\Bigr)\cdot
\Bigl(|\gamma(\theta_\iota)|^{-2}\Bigr)^{n'}\cdot\frac{Z(1)}{Z(\theta_\iota)} .
\end{equation}
\emph{The first factor has height at most $2\delta$.} Fix $\iota$, with prime $\ell$, and put
$\xi=\theta_\iota$, $\gamma=\gamma(\xi)$. Let $e$ be as in Definition~\ref{def:au-lattice}(b), and
let $u$, $u'$ be the matrices of the index at $\iota$ and at $\aubar\iota$; both have the prime
$\ell$, and both serve for $\xi$. Put $v=\operatorname{diag}(u,\gamma1_{n'})$ and
$v'=\operatorname{diag}(u',\gamma1_{n'})$. By (ii) and Definition~\ref{def:au-lattice}(a), the
vectors $\sum_kv_{ik}\Theta_\xi(w_k)$ lie in $\Lambda_{\xi,\iota}$ and the vectors
$\sum_kv'_{ik}\Theta_\xi(w_k)$ lie in $\Lambda_{\xi,\aubar\iota}$. By \eqref{eq:au-pairing} the
matrix of their mutual inner products, which is $\aubar{v'}\,\auGram(\Theta_\xi w)\,v^{t}$, has
entries in $\auO_\iota$. Taking determinants,
$|\aubar{\det v'}|_\iota\,|\Gamma(\xi)|_\iota\,|\det v|_\iota\le1$. Now
$|\aubar{\det v'}|_\iota=|\det u'|_{\aubar\iota}\,|\aubar\gamma|^{n'}_\iota\ge
\ell^{-e_\ell}|\aubar\gamma|^{n'}_\iota$ and $|\det v|_\iota\ge\ell^{-e_\ell}|\gamma|_\iota^{n'}$,
so $\bigl||\gamma|^{2n'}\Gamma(\xi)\bigr|_\iota\le\ell^{2e_\ell}$. This is the bound of
Definition~\ref{def:au-height} with $H=1$, $\alpha=0$ and $h_\ell=2e_\ell$.

\emph{Conclusion.} By (iv) and the rules for heights, the constant family $\Gamma(1)$ has height at
most $2\delta+n'b_{\mathrm{den}}+b_{\mathrm{vol}}$. It is not zero, because $\Theta_1$ is injective.
For $\sigma\in\auAut(\C)$ put $\sigma'=c\sigma c$. As $G$ is rational,
$\sigma\langle\Theta_1v,\Theta_1w\rangle=\aubar{\sigma'(\operatorname{coord}\Theta_1v)}^{\,t}G\,
\sigma(\operatorname{coord}\Theta_1w)$; by the Cauchy--Schwarz inequality in the top exterior
power, $|\sigma\Gamma(1)|^2$ is at most the product of the Gram determinants for $G$ of the
conjugates of the family $w$ by $\sigma'$ and by $\sigma$. Each is at most $e^{c}\,
e^{n'b_{\mathrm{short}}}$, by Fischer's inequality for the blocks $y$, $z$, by (i), and by
Hadamard's inequality and (iii). Hence $|\sigma\Gamma(1)|\le\exp(c+n'b_{\mathrm{short}})$ for all
$\sigma$, and Proposition~\ref{au:product} gives the claim.
\end{proof}

The theorem is a slope inequality in the style of Bost~\cite{BostFoliations} and of the
interpolation determinants of Laurent~\cite{LaurentInterp}; its one new feature is that the
integral structure at $\ell$ is tested after a twist $\theta_\iota$ that may depend on the place,
the cost being the height of $Z(1)/Z(\theta_\iota)$.


\subsection{The spaces, the groups and the forms}\label{ss:au-spaces}

\subsubsection*{Hermitian spaces}
All hermitian spaces are over $E$, for the conjugation $c$, and carry a lattice that is unimodular
at every finite place. Put $\varrho=m/2-1$, so that $m=2\varrho+2$ and $r=k+2\varrho+1$.
\begin{itemize}
\item $X_j$ ($0\le j\le m$) is $E^{n_j}$ with the form $\sum x_i\aubar{x_i}$ and the lattice
$\auO_E^{n_j}$; it is positive definite at every infinite place. The last basis vector of $X_j$ is
written $e_{\mathrm{last}}$.
\item $Y_0=X_0\perp\mathbb H^{\varrho}\perp\langle-N\rangle$, where $\mathbb H$ is the hyperbolic
plane; it has rank $r$ and signature $(k+\varrho,\varrho+1)$ at every infinite place. For
$0<j<m$, $Y_j=X_j\perp S_j$, where the \emph{string} $S_j$ has rank $m-1-j$ and consists of
hyperbolic planes and, if $m-1-j$ is odd, a line $\langle-N\rangle$; the rank of $Y_j$ is
$r+jt$ and it has $\lceil(m-1-j)/2\rceil$ negative signs at every infinite place. Finally
$Y_m=e_{\mathrm{last}}^{\perp}\subset X_m$, positive definite of rank $n_m-1=r+mt$. The lattice of
$Y_j$ is the lattice $\mathsf M_j$ of \eqref{ps:eq-lattice} (\S\ref{ps:ss-Y}); all objects on $Y_0$
below (the level $J_M$, the forms, the action of $\auG$, the Poincar\'e map) are formed with
$\mathsf M_0$. The proofs for Theorem~\ref{thm:srcB} are written with another unimodular lattice $L$ of
$Y_0$, and Lemma~\ref{pb:lattice} carries their results to $\mathsf M_0$.
\item $W_j$ ($1\le j\le m$) is $E^r$ with the diagonal form with $k+j-1$ entries $+1$ followed by
$m-j$ entries $-1$; $V_j$ is the positive definite standard space of rank $r+t$.
\end{itemize}
The unitary groups are regarded as groups over $E^+$. Thus $\auU(X_j)$, $\auU(Y_m)$, $\auU(V_j)$
and $\auU(W_m)$ are definite, while $\auU(Y_j)$ ($j<m$) and $\auU(W_j)$ ($j<m$) are not. Two
chains of dual pairs are used: the \emph{upper chain} $\auU(X_0),\auU(X_1),\dots,\auU(X_m)$ and
the \emph{lower chain} $\auU(Y_0),\auU(Y_1),\dots,\auU(Y_m)$; consecutive groups form a dual
reductive pair of unitary groups in the sense of Howe, and the rank grows by $t+1$, respectively
by $t$, at each step. The inclusions $X_j\subset Y_j$ ($j<m$) and $Y_m\subset X_m$ relate them: the
string loses one vector at each step, and at the last step the roles are exchanged. The pairs
$(\auU(W_j),\auU(V_j))$ are used only in the lower bound.

\subsubsection*{Algebraic modular forms and their integral structures}
Let $X$ be one of the positive definite standard spaces, of rank $n$, with $G=\auU(X)$,
$G_f=G(\auA_{E^+,f})$, $\Gamma=G(E^+)$, and $K^\circ\subset G_f$ the stabiliser of $\auO_E^n$. A
\emph{weight} $\lambda$ assigns to each $\tau\in\auPl$ a Young diagram with fewer than $n$ rows,
given by the numbers $\lambda_\tau(h)$ of its columns of height $h<n$, and an integer $e_\tau$.
Let $W_\lambda=\bigotimes_\tau W_{\lambda_\tau}\otimes\det^{e_\tau}$ be the corresponding
irreducible representation of $\prod_\tau\auGL_n(\C)$, in its model as the dual of the Schur module
spanned by the products of minors: $W_\lambda(\Z)$ is the lattice with the basis $(e_T)$ dual to the
standard products of minors, and the matrix of $g$ in this basis is the Schur matrix $S_\lambda(g)$
(\S\ref{pl:ss-norm}, (N1)); $W_\lambda$ carries an invariant hermitian form given on this basis by a
positive definite matrix with integer entries. The group $\Gamma$ acts on $W_\lambda$ through the
embeddings $\tau\circ c$: $\gamma$ acts by $\bigotimes_\tau\rho_\tau(\tau(\aubar\gamma))$, where
$\tau(\aubar\gamma)={}^t\tau(\gamma)^{-1}$. For an open
subgroup $J\subseteq K^\circ$ and a character $\chi$ of $J$ of finite order, the \emph{algebraic
modular forms} of weight $\lambda$ and level $(J,\chi)$, in the sense of Gross~\cite{GrossAMF},
are
\[
\mathcal A_\lambda(J,\chi)=\{\varphi\colon G_f\to W_\lambda:\ \varphi(\gamma gk)=\chi(k)\,\gamma\,
\varphi(g)\ \text{for }\gamma\in\Gamma,\ k\in J\},
\]
a finite-dimensional space with the inner product
$\langle\varphi,\varphi'\rangle=\sum_{g\in\Gamma\backslash G_f/J}\#(\Gamma\cap gJg^{-1})^{-1}
\langle\varphi(g),\varphi'(g)\rangle_{W_\lambda}$. For $\sigma\in\auAut(\C)$,
$\|\varphi\|_\sigma^2$ is the same sum for the vectors $\sigma(\varphi(g))$, with $\sigma$ applied
to the coordinates.

Let $\ell$ be a prime and $j\colon\C\to\C_\ell$ a ring homomorphism. Through the embeddings
$j\circ\tau\circ c$, that is, through its components at the primes
$w_\tau=\aubar{\mathfrak l_\tau}$ conjugate to the primes $\mathfrak l_\tau$ induced by $j\circ\tau$,
the group $G(E^+\otimes\Q_\ell)$ acts on $W_\lambda\otimes_{\C,j}\C_\ell$ (\S\ref{pl:ss-norm}, (N2)). The
\emph{gauge} at $(j,g)$ of $\xi\in W_\lambda$ is the largest absolute value of a coordinate of
$g_\ell^{-1}\cdot j(\xi)$; the form $\varphi$ is \emph{integral at $j$} if every value
$\varphi(g)$ has gauge at most $1$ at $(j,g)$, that is, $j(\varphi(g))\in g_\ell\,W_\lambda
(\auO_{\C_\ell})$. This is the integral structure of \cite{GrossAMF}; it is well defined on
$\Gamma\backslash G_f/J$. More generally $\varphi$ \emph{has gauge at most $C$ at $j$} if all its
values have. For $\iota\in\auI$ we say ``at $\iota$'' for the homomorphism $\C\to\auQlb\subset
\C_\ell$. If all groups $\Gamma\cap gJg^{-1}$ are trivial (the level is \emph{free}) then
$\langle\varphi,\varphi'\rangle\in\auO_\iota$ whenever $\varphi$ is integral at $\aubar\iota$ and
$\varphi'$ at $\iota$. Indeed, let $G_\lambda$ be the integral Gram matrix. Since $S_\lambda$ has
integer coefficients, invariance of the hermitian form says ${}^tS_\lambda(\aubar B)\,G_\lambda\,
S_\lambda(B)=G_\lambda$ for $B$ unitary, that is ${}^tS_\lambda({}^tB^{-1})\,G_\lambda\,S_\lambda(B)
=G_\lambda$; both sides are regular functions on $\auGL_n$ and the unitary group is Zariski dense
in $\auGL_n(\C)$ (it is a real form), so this holds for all $B\in\auGL_n$. Now $\aubar\iota\circ\tau$
induces the prime $\aubar{w_\tau}$, and for $g$ unitary $g_{\aubar{w_\tau}}={}^tg_{w_\tau}^{-1}$
under the identifications; hence $\iota\langle\varphi(g),\varphi'(g)\rangle={}^ta\,G_\lambda\,b$
with $a,b$ the integral coordinate vectors of $g_\ell^{-1}\cdot\varphi(g)$ (at $\aubar\iota$) and of
$g_\ell^{-1}\cdot\varphi'(g)$ (at $\iota$), and the weights $\#(\Gamma\cap gJg^{-1})^{-1}$ equal
$1$. Moreover evaluation at representatives of the classes identifies
$\mathcal A_\lambda(J,\chi)$ with a subspace of $W_\lambda^{\,\Gamma\backslash G_f/J}$, with the
block diagonal integral Gram matrix.

\begin{lemma}[Free levels]\label{au:free}
Let $n\le2n_m$ and let $J\subseteq K^\circ$ consist of elements congruent to $1$ modulo every
chosen place. Then $\Gamma\cap gJg^{-1}=1$ for all $g\in G_f$.
\end{lemma}
\begin{proof}
An element $\gamma$ of the intersection has finite order, since $\Gamma$ is discrete and $gJg^{-1}$
compact. Let $w$ be a chosen place, above $p_0$. As $p_0$ splits completely in $E$, $E_w=\Q_{p_0}$
and $g_w^{-1}\gamma g_w\in1+p_0M_n(\Z_{p_0})$. So every eigenvalue $\zeta$ of $\gamma$ is a root of
unity of degree at most $n$ over $\Q_{p_0}$ with $|\zeta-1|_{p_0}<1$. If $\zeta\ne1$ it has order a
power of $p_0$ and degree at least $p_0-1\ge b^{10}-1>2n_m$ over $\Q_{p_0}$. Hence all eigenvalues
are $1$, and a unipotent element of finite order is trivial.
\end{proof}

\subsubsection*{Source forms}
The space $\auS$ of \emph{source forms} is $\mathcal A_{\lambda^X}(J_X,1)$ on $\auU(X_0)$, of rank
$k$. The diagram of $\lambda^X$ at $\tau$ has $a$ columns of height $1$ if $\tau\in\auPa$ and none
otherwise, so that $W_{\lambda^X}=\bigotimes_{\tau\in\auPa}V_\tau$, where $V_\tau$ is the dual of the space of
forms of degree $a$ in $k$ variables with the basis $(e_\beta)$ dual to the monomials; there is no
power of the determinant on $X_0$ ($e_{0,\tau}=0$), and for $j\ge1$ the determinant exponent of
$\lambda(X_j)$ is $e_{j,\tau}=n_{j-1}/2-\kappa_j(\tau)$ (\S\ref{ps:ss-weil}). The level $J_X$ consists of the $g\in K^\circ$ that are $\equiv1$ modulo every chosen
place and whose last row is $\equiv(0,\dots,0,1)$ modulo $\pi_N$; the level $J_{X_j}$ of $X_j$
($0\le j\le m$) is defined in the same way, and $J_X=J_{X_0}$ (\S\ref{pz:ss-levels}). It is free by
Lemma~\ref{au:free}. Put $\auDX=\dim\auS$, and for $j\colon\C\to\C_\ell$ let $\auS_j$ be the
$\auO_j$-module of source forms integral at $j$.

For each prime $w\mid\pi_N$ of $E$ (its residue field is $\auF_N$) let $z_w\in\auA_{E,f}^\times$
be a unit idele of norm one to $E^+$, equal to $1$ outside $\{w,\aubar w\}$, whose component at
$w$ generates $\auF_N^\times$. The scalar matrix $z_w$ lies in $K^\circ$, is central in $G_f$,
normalises $J_X$, and $z_w^{\,i}\in J_X$ if and only if $(N-1)\mid i$. So $z_w$ acts on $\auS$
with $z_w^{N-1}=1$. The \emph{ramified} source forms are
\[
\auSram=\bigcap_{w\mid\pi_N}\ker\Bigl(\sum_{i=0}^{N-2}z_w^{\,i}\Bigr)=
\Bigl(\prod_{w\mid\pi_N}(1-z_w)\Bigr)\auS ,
\]
the forms with no component on which some $z_w$ acts trivially; and
$\auSram_{\mathrm{int}}=\{\varphi\in\auSram:\varphi\in\auS_\iota\text{ for all }\iota\in\auI\}$.

\subsubsection*{Target forms}
Let $G_0=\auU(Y_0)$. At each $\tau\in\auPl$ let $\varsigma_\tau$ be the Harish-Chandra parameter
\[
\varsigma_\tau=\Bigl(\tfrac k2-1,\tfrac k2-2,\dots,1-\tfrac k2\,;\ -\kappa(\tau)-\Delta,\
-\kappa(\tau)-2\Delta,\ \dots,\ -\kappa(\tau)-m\Delta\Bigr)
\]
($k-1$ entries, then $m$ entries), which is regular with gaps at least $\Delta$ between the two
groups of entries and inside the second. Let $J_M\subset G_0(\auA_{E^+,f})$ be the principal
congruence subgroup of level $\pi_M$ in the stabiliser of the lattice of $Y_0$. At each $\tau$ let
$\pi_\tau$ be the discrete series representation of $G_0(E^+\otimes_\tau\R)\cong\auU(k+\varrho,\varrho+1)$
with Harish-Chandra parameter $\varsigma_\tau$ in the chamber in which the first $k-1$ entries and the
entries $-\kappa(\tau)-j\Delta$ with $j$ odd are attached to the positive part and the entries with $j$
even to the negative part (the signs fixed below); let $c_\tau=\langle v_0,\pi_\tau(\cdot)v_0\rangle$ for
a unit highest weight vector $v_0$ of its lowest $K$-type, and $e_\tau=\deg_\tau\overline{c_\tau}$, with
$\deg_\tau$ the formal degree. Let $\auAz$ be the space of the bounded continuous functions $y$ on
$G_0(E^+)\backslash G_0(\auA_{E^+})/J_M$ with $\int e_\infty(a)y(ga)\,da=y(g)$,
$e_\infty=\bigotimes_\tau e_\tau$. This is the space of the forms of type $e_0$ and level $J_M$ of
Definition~\ref{pz:def-type}; these forms are cuspidal and $\auAz$ is finite-dimensional
(Proposition~\ref{pz:types}(a)). The proofs for Theorem~\ref{thm:srcB} (\S\S\ref{pb:sec}--\ref{pm:sec},
\ref{pds:sec}--\ref{pfg:sec}) use only the reproducing equation and the boundedness of the forms. The
finite group
\[
\auG=\prod_{w\text{ chosen}}\auGL_r(\auF_w),\qquad\auF_w=\auF_{p_0}\text{ or }\auF_{p_1},
\]
the quotient of the stabiliser at the chosen places by $J_M$, acts unitarily on $\auAz$; write
$\varrho(s)$ for the action of $s\in\auG$.

A \emph{label} over a finite field $\auF$ is a monic squarefree polynomial of degree $r$ over
$\auF$ with non-zero constant term. In Green's parametrisation~\cite{Green1955} a label $c$ names
an irreducible character $\Xi_c$ of $\auGL_r(\auF)$ and hence a central primitive idempotent $e_c$ of
$\C[\auGL_r(\auF)]$; if $c$ is irreducible, $\Xi_c$ is a cuspidal character (property (G) of
\S\ref{pm:ss-mult}, the only property of the parametrisation used in the proof of
Proposition~\ref{au:mult}). A family $(c_w)_w$ of labels, one for each chosen place, is \emph{good} if no
$c_w$ has a root in $\auF_w$ that is a $q_2$-th root of unity and if, above each of $p_0$ and
$p_1$, at least one $c_w$ is irreducible with a root generating $\auF_{p_i^r}^\times$. Then
\[
e_0=\sum_{(c_w)\text{ good}}\ \bigotimes_we_{c_w}\ \in\C[\auG]
\]
is a central idempotent.

\begin{definition}[Target forms]\label{def:au-target}
The space of \emph{target forms} is $\auT=\varrho(e_0)\auAz$, the \emph{good forms} on $Y_0$; it
is a finite-dimensional inner product space, and $\audT=\dim\auT$.
\end{definition}

The same definitions on $W_j$, with the parameter $\varsigma$ and the level $J_M$, give spaces
$\mathcal A(W_j)$ and their good parts $e_0\mathcal A(W_j)$; for $j=m$ these are algebraic modular
forms on a definite group. In all cases the discrete series at $\tau$ is the one with
Harish-Chandra parameter $\varsigma_\tau$ for the following signs of the coordinates
(\eqref{pm:eq-signs}): the $k-1$ entries $\tfrac k2-1,\dots,1-\tfrac k2$ are on positive
coordinates; the entry $-\kappa(\tau)-a\Delta$, $1\le a\le m$, is on a positive coordinate of $W_j$
if and only if $a\le j$, and on a positive coordinate of $Y_0$ if and only if $a$ is odd. The
lattices of $Y_0$ and of the $W_j$ are unimodular at every finite place, so that these spaces are
isometric at all finite places by isometries of the lattices; this identifies their finite adelic
groups, the levels $J_M$ and the actions of $\auG$ (\S\ref{pm:ss-mult}).

\subsubsection*{Twists}
A \emph{twist} $\vartheta$ is a unitary automorphic character of $\auU(1)(\auA_{E^+})$, trivial at
the infinite places, unramified outside a finite set $Q=Q_\vartheta$ of primes of $E$ of norm
prime to $MN$ with $Q\cap\aubar Q=\emptyset$, and tamely ramified with non-trivial tame part
$\psi_w$ at each $w\in Q$. It is \emph{$m$-fold} if $\psi_w^n\ne1$ for $1\le n\le m$ and all
$w\in Q$. Put $\auNm Q_\vartheta=\prod_{w\in Q}\auNm w$. The trivial twist is $1$, with
$Q=\emptyset$. Such characters exist with any prescribed tame parts that are trivial on the roots
of unity of $E$ of norm one; this is class field theory.

\subsubsection*{The lower chain and the lattice of target forms}
Fix, at each chosen place $w$, a character $\omega_w$ of $\auF_w^\times$ of order $q_2$ (it exists
since $q_2^2\mid p_i-1$), and characters $\chi_M$ of $\auU(1)(\auA_{E^+})$ of weight $0$ with tame part
$\omega_w$ at every chosen place and unramified elsewhere, and $\mu^\flat$ of weight
$-\kappa-\frac\Delta2$ unramified at all finite places; put $\mu_j^0=\mu^\flat(\upsilon^{-\Delta}\chi_M)^j$,
of weight $-\kappa_j$, with tame part $\omega_w^j$ at the chosen places and unramified elsewhere
(Lemma~\ref{pz:mu}, Definition~\ref{pfq:def-chars}, Proposition~\ref{pfq:tame}). For a twist
$\vartheta$ put $\mu_j=\mu_j^0\vartheta^j$. For $0\le j<m$ let
\[
\Theta^Y_{j,\vartheta}\colon\{\text{forms on }\auU(Y_j)\}\to\{\text{forms on }\auU(Y_{j+1})\}
\]
be the theta lift for the dual pair $(\auU(Y_j),\auU(Y_{j+1}))$, with splitting characters built
from $\mu_{j+1}$, and with the following Schwartz function on $\auHom_E(Y_j,Y_{j+1})(\auA_{E^+})$:
at an infinite place the Gaussian times the polynomial joining the two archimedean types; at a
chosen place and at a place of $Q$ the indicator of the homomorphisms that reduce to the inclusion
of the lattices; at every other finite place the indicator of the integral homomorphisms. The lift
is normalised by the index of its level. The \emph{lower chain} is
$\iota_\vartheta=\Theta^Y_{m-1,\vartheta}\circ\dots\circ\Theta^Y_{0,\vartheta}$. On $\auAz$ it
takes values in a space $\auAm(\vartheta)=\mathcal A_{\lambda^\circ}(J_\vartheta,\chi_\vartheta)$
of algebraic modular forms on the definite group $\auU(Y_m)$, where $\lambda^\circ$ is a fixed
weight, $J_\vartheta$ consists of the elements of $K^\circ$ that are $\equiv1$ at the chosen places
and block upper triangular modulo $Q$, and $\chi_\vartheta$ is given by the tame characters of the
lifts on the diagonal blocks. By Lemma~\ref{au:free} the level $J_\vartheta$ is free.

\begin{definition}[The lattice of target forms]\label{def:au-targetlattice}
The twisted lattice $\auLat_D$ on $\auT$ has as twists the $m$-fold twists; the maps
$\Theta_\vartheta=\iota_\vartheta|_{\auT}\colon\auT\to\auAm(\vartheta)$; $\vartheta$ is allowed at
$\ell$ if $\ell\nmid\auNm Q_\vartheta$; $\Lambda_{\vartheta,\iota}$ is the module of forms of
$\auAm(\vartheta)$ integral at $\iota$; and the coordinates on $\auAm(1)$ are the coordinates of
the values at representatives of the classes. Thus a target form $y$ is \emph{integral at $\iota$}
if $\Theta_\vartheta y$ has integral values at $\iota$ for every $m$-fold twist $\vartheta$ of
modulus prime to $\ell(\iota)$; the \emph{index} and the \emph{covolume} of a family of target
forms are those of Definition~\ref{def:au-lattice}; and $Z(\vartheta)=\det(\Theta_\vartheta^*
\Theta_\vartheta)$ is the \emph{chain determinant}.
\end{definition}

Condition~\eqref{eq:au-pairing} holds because the levels $J_\vartheta$ are free. The injectivity
of $\Theta_\vartheta$ is Proposition~\ref{au:inj} below. Note that integrality of a target form is
a condition on infinitely many lifts; that the integral forms span $\auT$ is a consequence of
Theorems~\ref{thm:srcA} and~\ref{thm:srcB}, not of the definition.

\subsubsection*{The Poincar\'e map}
Let $P\subset G_0$ be the parabolic subgroup stabilising a full flag of isotropic subspaces in
$\mathbb H^\varrho$, with unipotent radical $U$ and Levi factor
$\auGL_1(E)^\varrho\times\auU(X_0\perp\langle-N\rangle)$, let $\psi_U$ be the generic character
\eqref{pb:eq-psi} of $U(E^+)\backslash U(\auA)$, which is invariant under conjugation by the subgroup
$\auU(X_0)$ of the Levi factor, and let
$R_0=U\rtimes\auU(X_0)$ be the corresponding Bessel subgroup. Let $g_0\in
G_0(E^+\otimes\R)$ be a point and $\mathsf t$ a linear form on $W_{\lambda^X}$; the definitions that
follow make sense for every pair $(g_0,\mathsf t)$, and Proposition~\ref{au:seesaw} is stated and proved
for every pair. \emph{From now on $(g_0,\mathsf t)$ is the pair of Theorem~\ref{pb:est}}: it is for this
pair, and for the character \eqref{pb:eq-psi}, that Proposition~\ref{au:moment} holds
(Theorem~\ref{pb:thmB}), and the map $\auP$ of Theorems~\ref{thm:srcA}, \ref{thm:srcB} and~\ref{thm:srcC}, the
constant $c_{\mathrm{ss}}$ of \S\ref{ss:au-proofAB} and the constant $c_P$ of
Proposition~\ref{au:moment} are those of this pair. For $\varphi\in\auS$ let
$\mathcal W_\varphi(uh)=\psi_U(u)\,\mathsf t\bigl(h_\infty^{-1}\varphi(h_f)\bigr)$ on
$R_0(\auA)$. The \emph{Poincar\'e map} $\auP\colon\auS\to\auAz$ is defined by
\[
\langle y,\auP\varphi\rangle=\int_{R_0(E^+)\backslash R_0(\auA)}\mathcal W_\varphi(r)\,
\aubar{y(rg_0)}\,dr\qquad(y\in\auAz);
\]
so $\auP\varphi$ is the projection to $\auAz$ of a Poincar\'e series of Bessel type, and it
represents the Bessel period. For $s\in\auG$ put $\auP_s=\#\auG\cdot\varrho(e_0)\varrho(s)\auP
\colon\auS\to\auT$.

\subsubsection*{The upper chain and the composite lift}
Let $0\le j<m$ and write $\mathcal H_j=\auHom_E(X_j,X_{j+1})(\auA_{E^+,f})$. For a
Schwartz--Bruhat function $\phi_j$ on $\mathcal H_j$ let $\Theta^X_j(\phi_j)$ be the theta lift from algebraic
modular forms on $\auU(X_j)$ to those on $\auU(X_{j+1})$, with the splitting characters of the
lower chain and, at the infinite places, the pluriharmonic polynomial joining the weights
(both groups being compact there, the kernel is a finite-dimensional space of Taylor coefficients
of theta functions; see \S\ref{ss:au-size}). The weights along the upper chain have, at $\tau$,
$a\cdot\mathbf 1_{\tau\in\auPa}$ columns of height $1$, $\kappa_1(\tau)-\lfloor n_1/2\rfloor$
columns of height $k$ and $\Delta-(t+1)$ columns of each height $n_i$, $0<i<j$. The \emph{vacuum}
$\phi^\circ_j$ is the indicator of the integral homomorphisms, suitably scaled; its translates by
the finite Heisenberg group are indexed by $x_j\in\mathcal H_j$. For a
finitely supported function $\Psi$ on tuples $x=(x_0,\dots,x_{m-1})$ let
\[
C_{\vartheta,\Psi}=\sum_x\Psi(x)\ \Theta^X_{m-1}(\phi^\circ_{m-1,x_{m-1}})\circ\dots\circ
\Theta^X_0(\phi^\circ_{0,x_0}),
\]
and let $\operatorname{res}$ be the restriction of forms from $\auU(X_m)$ to $\auU(Y_m)$, followed
by the branching projection $W_{\lambda(X_m)}\to W_{\lambda^\circ}$ and by multiplication by
$\prod_\tau\tau(\aubar{\pi_N})^{-n_\tau}$, where $n_\tau=|\lambda(X_m)_\tau|-|\lambda^\circ_\tau|$
is the difference of the numbers of boxes; one has
\begin{equation}\label{eq:au-branch}
\sum_{\tau\in\auPl}n_\tau\ \ge\ \#\auPa\cdot a-f\,m\,\Delta/2 .
\end{equation}
The \emph{composite lift} is $B_{\vartheta,\Psi}=\operatorname{res}\circ C_{\vartheta,\Psi^{\mathrm
g}}\colon\auS\to\{\text{forms on }\auU(Y_m)\}$, where $\Psi^{\mathrm g}$ is obtained from $\Psi$ by
a fixed \emph{gluing} operation on the data at the places above $N$ in each pair of consecutive
lifts. A \emph{chain datum} for the twist $\vartheta$ is a $\Psi$ such that: its support consists
of tuples that vanish at all places prime to $M\auNm Q_\vartheta$ and satisfy an admissibility
condition there; its values are algebraic numbers integral at all primes not dividing
$\auNm Q_\vartheta$; if $Q_\vartheta=\emptyset$ then $\sum_x|\sigma\Psi(x)|\le M^{3fmn_m^2}$ for
every $\sigma\in\auAut(\C)$; and the glued tensor of data is invariant under the levels of
$X_0,\dots,X_{m-1}$.


\subsection{Prime sums, the gain at a prime, and the main amount}\label{ss:au-primes}

Let $\Lambda_K(\nu)=\Lambda(\nu)(1-\auchi(\nu))$. For a prime $p$ put $v(\pi_p)=1$ if $p\ne d$ and
$v(\pi_d)=\tfrac12$, and let $I_p$ be the set of odd positive integers if $p$ is inert in $K$, the
set of all positive integers if $p=d$, and empty if $p$ splits.

\begin{definition}[Prime sums, gain, main amount, remainder]\label{def:au-hmain}
Let $y\ge0$ be real and $\sigma\in\auPl$.
\begin{enumerate}
\item[(a)] The \emph{local prime sum} at $p$ is
$s_p(y,\sigma)=v(\pi_p)\sum_{e\in I_p}\bigl(y/p^e-\kappa(u(-p^e)\sigma)\bigr)^+$, and the
\emph{prime sum} is
$\ausfin(y,\sigma)=\sum_{1\le\nu\le a}\Lambda_K(\nu)\bigl(y/\nu-\kappa(u(-\nu)\sigma)\bigr)^+$.
\item[(b)] The \emph{gain} at $p$ is
\[
\augain_p=m\Bigl(\sum_{\tau\in\auPa}s_p(a,\tau)+k\sum_{\sigma\in\auPl}s_p(\kappa(\sigma),\sigma)
\Bigr).
\]
\item[(c)] The size of a derivative of order $j$ is $\ausinf(j)=\log j!+j\log(M/\Omega)$; the
\emph{direction amount} is $\auhdir=\sum_{\sigma\in\auPl}\bigl(\ausinf(\kappa(\sigma))-
\ausfin(\kappa(\sigma),\sigma)\bigr)$; the \emph{main amount} is $\auhmain=m\,k\,\auhdir$.
\item[(d)] The \emph{expected value} is
$\ausmain=a\bigl(\log a-\tfrac1f\sum_{\tau\in\auPl}\log\kappa(\tau)+c_1\bigr)$, and the
\emph{total remainder} is $\aurem=\sum_{\tau\in\auPa}\bigl(\ausfin(a,\tau)-\ausmain\bigr)$.
\end{enumerate}
\end{definition}

The place $u(-p^e)\sigma$ is the \emph{partner} of $\sigma$ at $p^e$: for $p$ non-split in $K$
and $e\in I_p$ the element $c\circ\mathrm{Frob}_p^{\,e}$ of $\operatorname{Gal}(E/\Q)$ lies in
$\operatorname{Gal}(E/K)$ and is the class of $-p^e$. In the analytic section it is shown that
$\aurem=a\bigl(\#\auPb\,(R_{q_0}(x)-R_{q_0}(x_1))+\#\auPa\,R_1(x_1)\bigr)$, where $R_q$ is the
remainder of the weighted prime sum over $\nu\equiv1\pmod q$: for $\tau\in\auPa$ the partner at
$\nu$ is a $b$-place exactly for $\nu\equiv1\pmod{q_0}$, and then the term is
$b\,(x/\nu-1)^+$.

\begin{lemma}[The gains and the prime sums]\label{au:gains}
Let $0\le y\le a$ and $\sigma\in\auPl$.
\begin{enumerate}
\item[(a)] $\ausfin(y,\sigma)=\sum_p2\log p\cdot s_p(y,\sigma)$, and $s_p(y,\sigma)=0$ if $p$
splits in $K$ or $p\ge x_1$. Hence $\augain_p=0$ for such $p$.
\item[(b)] With $\mathcal N$ the set of primes $p<x_1$ that do not split in $K$,
\[
2\sum_{p\in\mathcal N}\augain_p\log p=m\sum_{\tau\in\auPa}\ausfin(a,\tau)+m\,k\sum_{\sigma\in\auPl}
\ausfin(\kappa(\sigma),\sigma),
\]
and
\begin{equation}\label{eq:au-cancel}
\auhmain-m\cdot\aurem+2\sum_{p\in\mathcal N}\augain_p\log p=m\,k\sum_{\sigma\in\auPl}\ausinf
(\kappa(\sigma))+m\cdot\#\auPa\cdot\ausmain .
\end{equation}
\item[(c)] $\ausinf(a)-\ausmain\le a\bigl(\tfrac1f\sum_\tau\log\kappa(\tau)+\log\frac{M\aurd}
{2\pi}\bigr)+1+\tfrac12\log a$, and
\begin{equation}\label{eq:au-excess}
m\bigl(\ausinf(a)-\ausmain\bigr)-a\log N\le32\,a\log b .
\end{equation}
\end{enumerate}
\end{lemma}
\begin{proof}
(a) $\Lambda_K$ is supported on powers $\nu=p^e$. If $p$ splits, $\Lambda_K(p^e)=0$. If $p$ is
inert, $\Lambda_K(p^e)=(1-(-1)^e)\log p$, which is $2\log p$ for odd $e$ and $0$ otherwise. If
$p=d$, $\Lambda_K(d^e)=\log d=2v(\pi_d)\log d$. A non-split prime does not divide $q_1$, so
$u(-p^e)$ is defined. Terms with $p^e>a$ vanish on both sides because $y/p^e<1\le\kappa$, and if
$p\ge x_1$ then $y/p^e\le a/x_1\le\kappa$. (b) The first identity is (a) summed with the weights
of Definition~\ref{def:au-hmain}(b); the second follows by inserting the definitions of
$\auhmain$ and $\aurem$: the two sums of $\ausfin$ cancel. (c) By Stirling's formula in the form
$\log a!\le a\log a-a+1+\tfrac12\log a$, and since $\log\Omega=-1-c_1+\log(2\pi/\aurd)$,
\[
\ausinf(a)\le a(\log a+c_1)+a\log\frac{M\aurd}{2\pi}+1+\tfrac12\log a ;
\]
subtracting $\ausmain$ gives the first claim, in which $c_1$ has cancelled. Write $w$ for the mean
of $\log\kappa$ and $G=\log(M\aurd/2\pi)$. By \eqref{eq:au-level},
$\log N\ge mw+(m-1)G-\tfrac{15}2\log b$, so the left side of \eqref{eq:au-excess} is at most
$aG+\tfrac{15}2a\log b+m(1+\tfrac12\log a)\le(23+\tfrac{15}2+1)\,a\log b$, by
\eqref{eq:au-sizes} and a comparison of sizes.
\end{proof}

Identity \eqref{eq:au-cancel} is the bookkeeping behind the whole argument: its right-hand side
contains no primes. Inequality \eqref{eq:au-excess} is the reason for the lower end of the window \eqref{eq:au-level} for the level $N$. The upper end
of the window is needed in Theorem~\ref{thm:srcB}: there the terms of the expansion other than
the identity term have to be small against the identity term, and this is the upper bound for the
ratio of the two dimensions in Lemma~\ref{pb:sizes}, used in Proposition~\ref{pbe:two}. The two
requirements are compatible: both place $\log N$ at $m\log(\bar\kappa M\aurd/2\pi)$ up to
$O(\log b)$, the first from the size of the lifts and the second from the dimensions of the spaces.

\subsubsection*{The local computation at a non-split prime}
The prime sums enter through the $p$-adic valuation of Taylor coefficients of theta functions
with complex multiplication at a prime $p$ that does not split in $K$. The statement used is the
following (\emph{the count}); we state it and then prove the model case that
shows where the shape $(y/p^e-\kappa)^+$ comes from.

Let $p$ be non-split, fix $\C\to\C_p$, let $v$ be the valuation with $v(p)=1$, and put
$\rho_p=v(\pi_p)/(q_p-1)$ with $q_p$ the norm of the prime of $K$ above $p$ (so
$\rho_p=1/(p^2-1)$ for inert $p$). Let $z=(z_{\tau,i})$ be variables indexed by $\tau\in\auPl$ and
$i$ in a finite set. For $e\ge0$ the \emph{coupling} is $\Gamma_e=\sum_{\tau,i}z_{\tau,i}^{\,p^e}
\,z_{u(-p^e)\tau,i}$. A power series $F\in\C_p[[z]]$ \emph{has the normal form} if, for some
$\varpi$ with $v(\varpi)=\rho_p$,
\begin{equation}\label{eq:au-normal}
F(\varpi z)=W(z)\cdot\exp\Bigl(\frac{\varpi^{p+1}}p\sum_{e\in I_p^+}\ \sum_{i\ge0}
\lambda_{e,i}\,\Gamma_e^{\,p^i}\Bigr),\qquad W\in\auO_{\C_p}[[z]],\ \lambda_{e,i}\in\auO_{\C_p},
\end{equation}
where $I_p^+=I_p$ if $p\ne d$ and $I_p^+=I_p\cup\{0\}$ if $p=d$.

\begin{proposition}[The count]\label{au:count}
Let $F$ have the normal form. Let $\lambda$ be a weight whose columns have height less than $n'$, with $a_\tau$ columns of
height $1$, $a'_\tau$ other columns and $l_\tau$ boxes in these, and $B\ge1$ boxes in all. Let $Q$ be either
the product $Q_{T,y}$, for a filling $T$ of $\lambda$ and integral directions $y$, of the minors of
$z_\tau y_\tau$ on the rows of each column and the first $h$ columns ($h$ the height of the column),
as in Proposition~\ref{pa:count}, or an entry $S_\lambda(z)_{T',T}$ of the Schur matrix with $T'$
standard (Lemma~\ref{pa:schur}). Then the jet
$Q(\partial)F(0)$ has valuation at least
\[
v(\pi_p)\sum_{e\in I_p}\sum_{\tau}\Bigl(\bigl(\tfrac{a_\tau}{p^e}-a_{\tau'}-a'_{\tau'}\bigr)^+
+\bigl(\tfrac{l_\tau}{p^e}-l_{\tau'}\bigr)^+-\tfrac{a'_\tau}{p^e}\Bigr)-f\,n'^3\bigl(1+\log_p
B\bigr),\qquad\tau'=u(-p^e)\tau.
\]
For $p=d$ the same bound holds after subtracting in
addition the term $\tfrac12\sum_\tau\mathrm{ded}_0(\tau)=v(\pi_d)\sum_\tau\mathrm{ded}_0(\tau)$ for $e=0$, with
$\mathrm{ded}_0$ as in Proposition~\ref{pa:count}, which is the precise statement in all cases.
\end{proposition}
\begin{proof}
This is Proposition~\ref{pa:count}, proved in \S\ref{pa:ss-count} for exactly these jets
(with the explicit term for $p=d$), together with Lemma~\ref{pa:schur} (entries of the Schur matrix). The only published
result used is standard monomial theory~\cite{DEP}. The model case below contains the whole idea.
\end{proof}

\begin{proposition}[Model case: one pair of places, one variable each]\label{au:local}
Let $p$ be an odd prime, $v(\varpi)=1/(p^2-1)$, and let $F\in\C_p[[z,w]]$ satisfy
\[
F(\varpi z,\varpi w)=W(z,w)\,\exp\Bigl(\lambda\,\frac{\varpi^{p+1}}p\,z^pw\Bigr),\qquad
W\in\auO_{\C_p}[[z,w]],\ \lambda\in\auO_{\C_p}.
\]
Then for all integers $y,y'\ge0$,
\[
v\bigl(\partial_z^{\,y}\partial_w^{\,y'}F(0)\bigr)\ \ge\ \Bigl(\frac yp-y'\Bigr)^+
+\frac{p}{p^2-1}\,y'+\frac{y}{p(p^2-1)}-\log_p(y+1)-\log_p(y'+1).
\]
Moreover the loss $\min(\lfloor y/p\rfloor,y')$ that is responsible for the positive part is
attained up to $\log_p(y+1)$ when $\lambda$ and $W_{y-pj_0,y'-j_0}$ are units, where
$j_0=\min(\lfloor y/p\rfloor,y')$, and $W_{y-pj,y'-j}=0$ for the other $j\le j_0$.
\end{proposition}
\begin{proof}
Write $F=\sum f_{y,y'}z^yw^{y'}$ and $c=\lambda\varpi^{p+1}/p$. Comparing coefficients,
\[
f_{y,y'}\,\varpi^{y+y'}=\sum_{0\le j\le\min(\lfloor y/p\rfloor,\,y')}W_{y-pj,\,y'-j}\,
\frac{c^j}{j!} .
\]
Now $v(c)\ge(p+1)/(p^2-1)-1=1/(p-1)-1$ and $v(j!)=(j-s_p(j))/(p-1)$, with $s_p(j)$ the sum of the
digits of $j$ in base $p$; hence $v(c^j/j!)\ge-j+s_p(j)/(p-1)\ge-j$, with equality up to
$s_p(j)/(p-1)$ if $\lambda$ is a unit. Therefore $v(f_{y,y'})\ge-(y+y')/(p^2-1)-\min(\lfloor
y/p\rfloor,y')$. The jet is $y!\,y'!\,f_{y,y'}$, and $v(y!)\ge y/(p-1)-\log_p(y+1)$. With
$\gamma_p=1/(p-1)-1/(p^2-1)=p/(p^2-1)$ this gives
\[
v\bigl(\partial_z^{\,y}\partial_w^{\,y'}F(0)\bigr)\ge\gamma_p(y+y')-\min\bigl(\tfrac yp,y'\bigr)
-\log_p(y+1)-\log_p(y'+1),
\]
and $-\min(y/p,y')=(y/p-y')^+-y/p$, $\gamma_p-1/p=1/(p(p^2-1))$. For the last assertion take
$j_0=\min(\lfloor y/p\rfloor,y')$ and $W$ with $W_{y-pj_0,y'-j_0}$ a unit and the other
coefficients in the sum zero.
\end{proof}

\begin{remark}[Reading of the model]\label{au:local-remark}
(1) In the application $z$ and $w$ are the Taylor variables of a theta function at two infinite
places $\tau$ and $\tau'=u(-p)\tau$, $y$ and $y'$ are the orders of differentiation, and $y'$ is
the weight $\kappa(\tau')$. The normal form expresses two facts about an abelian variety with
multiplication by $\auO_E$ at a prime $p$ inert in $K$: its formal group has height two on each
factor, so the theta function converges only on the disc of radius $p^{-\rho_p}$ and each
derivative gains only $\gamma_p$ instead of $1/(p-1)$; and the Frobenius at $p$ composed with $c$
carries $\tau$ to $\tau'$, so that the hermitian exponential factor of the theta function becomes
$\exp(\text{const}\cdot z^pw)$. The terms proportional to $y$ and $y'$ are \emph{smooth}: summed
over $p$ with weights $\log p$ they are absorbed in the normalisation by the period (they are the
$p$-adic counterpart of $y\log(M/\Omega)$, and after rescaling the series by the period they
do not appear; compare Proposition~\ref{au:count}). What remains is $(y/p-y')^+$: \emph{the
$y$-th derivative at $\tau$ is divisible by $p^{(y/p-\kappa(\tau'))^+}$ beyond the smooth
amount}. With the higher couplings $\Gamma_e$ this becomes $\sum_{e\in I_p}(y/p^e-\kappa(u(-p^e)
\tau))^+=s_p(y,\tau)$, and $\sum_p2\log p\cdot s_p=\ausfin$ by Lemma~\ref{au:gains}(a).

(2) For $E=K$ and one variable these are the classical divisibilities of Bernoulli--Hurwitz
numbers and of Taylor coefficients of elliptic theta functions at supersingular primes
(Katz~\cite{KatzDivisibilities}, Chellali~\cite{Chellali},
Bannai--Kobayashi~\cite{BannaiKobayashi}). What is used here goes beyond them in three ways: many
places coupled by the partner rule, exact constants in the weight aspect, and uniformity in the
rank.

(3) The gain of Definition~\ref{def:au-hmain}(b) adds these exponents over what the $m$ lifts of
the upper chain differentiate: at each lift, a derivative of order $a$ at each $a$-place (the
first row of the diagram) and derivatives of order about $\kappa(\sigma)$ in $k-1$ further
directions at every place $\sigma$ (counted as $k$).
\end{remark}


\subsection{The three theorems}\label{ss:au-ABC}

Let $\mathcal N$ be the set of primes $p<x_1$ not split in $K$. For $s\in\auG$ and a linear map
$\aupp\colon\auSram\to\auT$ write $s\aupp=\varrho(s)\circ\aupp$.

\begin{theorem}[Theorem A: short integral forms]\label{thm:srcA}
There are $c_{\mathrm{ss}}\in\C^\times$ and a positive rational number $\aucB$, a product of
powers of the primes in $\mathcal N\cup\{N\}$, such that the map
\[
\aupp=\aucB\,c_{\mathrm{ss}}^{-1}\,\auP_1|_{\auSram}\colon\auSram\to\auT
\]
has the following properties.
\begin{enumerate}
\item[(a)] \emph{(Size.)} $\auSram$ is spanned over $\C$ by the forms
$\varphi\in\auSram_{\mathrm{int}}$ such that, for all $s\in\auG$ and $\sigma\in\auAut(\C)$,
\[
\|\Theta_1(s\aupp\,\varphi)\|_\sigma^2\le\exp\bigl(\auhmain-m\cdot\aurem+f\,a\,(\log b)^4\bigr).
\]
\item[(b)] \emph{(Index.)} For every finite set $\mathsf t\subseteq\auG$, every finite family of
target forms each of which lies in $s\aupp(\auSram_{\mathrm{int}})$ for some $s\in\mathsf t$ has
index at most $\#\mathsf t\cdot\dim\auSram\cdot2\aueu$ in $\auLat_D$.
\end{enumerate}
\end{theorem}

\begin{theorem}[Theorem B: enough forms]\label{thm:srcB}
$\audT>0$, and there is a finite set $\mathsf t\subseteq\auG$ with
\[
\#\mathsf t\cdot\dim\auSram\le5\,\audT\log(1+2b^2)\qquad\text{and}\qquad
\dim\sum_{s\in\mathsf t}s\aupp(\auSram)\ \ge\ \Bigl(1-\frac4{b^2}\Bigr)\audT .
\]
\end{theorem}

\begin{theorem}[Theorem C: integral forms are long]\label{thm:srcC}
Let $y_1,\dots,y_n\in\auT$ be linearly independent with $n\ge(1-4/b^2)\audT$, of index at most
$\delta$ and covolume at most $c$ in $\auLat_D$. Then
\[
c+2\delta\ \ge\ \audT\cdot\bigl(\auhmain-3\,f\,a\,(\log b)^4\bigr).
\]
\end{theorem}

\begin{proof}
Theorems~\ref{thm:srcA} and~\ref{thm:srcB} are deduced in \S\ref{ss:au-proofAB} from
Propositions~\ref{au:seesaw}, \ref{au:divis}, \ref{au:kernel}, \ref{au:restrict},
\ref{au:expansion}, \ref{au:equiv} and~\ref{au:moment}. Of these, Propositions~\ref{au:restrict},
\ref{au:expansion} and~\ref{au:equiv} are proved in \S\ref{ps:sec}; Propositions~\ref{au:divis}
and~\ref{au:kernel} in \S\S\ref{pa:sec}--\ref{pf:sec}; Proposition~\ref{au:seesaw} in
\S\S\ref{ps:sec}, \ref{pk:sec}, \ref{pf:sec}, using that the lower chain is well defined
and consists of cusp forms, which is Proposition~\ref{pz:types}, deduced from the statements
(A1)--(A3) of \S\ref{pz:ss-types}; these are proved in \S\ref{par:sec} and \S\ref{pco:sec} from
published results on discrete series and on the archimedean Weil representation.
Proposition~\ref{au:moment} is proved in \S\ref{pb:sec} and the sections after it, the covolume of
Lemma~\ref{pb:covol} being Theorem~\ref{pbv:covol}. The published results used are listed for
Theorem~\ref{thm:srcA} in \S\ref{pz:ss-audit}, and for Theorem~\ref{thm:srcB} in \S\ref{pb:ss-remains}.
Theorem~\ref{thm:srcC} is deduced in \S\ref{ss:au-proofC} (in more detail:
Proposition~\ref{pc:C}) from Propositions~\ref{au:inj}, \ref{au:aux}, \ref{au:factor},
\ref{au:regular}, \ref{au:onelift}, \ref{au:hecke} and~\ref{au:mult}, each of which is followed by
a proof naming the sections where it is proved; the first factors of the proof of
Proposition~\ref{au:onelift} are treated in \S\ref{pfq:sec} (Theorem~\ref{pfq:ff}), and the tame
characters at the chosen places in Definition~\ref{pfq:def-chars} and Proposition~\ref{pfq:tame}.
The table of \S\ref{ss:au-remains} lists the published results on which these inputs rest.
\end{proof}

In the analytic section the three theorems are combined: a basis of
$\sum_{s\in\mathsf t}s\aupp(\auSram)$ chosen among the short vectors of (a) has, by Hadamard, covolume at
most $n(\auhmain-m\aurem+fa\log^4b)$ and, by (b) and Theorem~\ref{thm:srcB}, index at most
$10\,\audT\log(1+2b^2)\aueu$; Theorem~\ref{thm:srcC} then gives $m\cdot\aurem\le f\,a\,(\log b)^5$.

\begin{remark}[One statement]\label{au:merged}
The three theorems are equivalent in use to the following single statement about the lattice of
target forms. \emph{Every family of independent target forms with $n\ge(1-4/b^2)\audT$ members
satisfies $c+2\delta\ge\audT(\auhmain-3fa\log^4b)$; and there is such a family with
$c\le n(\auhmain-m\cdot\aurem+fa\log^4b)$ and $\delta\le10\,\audT\log(1+2b^2)\,\aueu$.} In the
language of slopes: the adelic hermitian lattice of integral target forms is semistable up to
$O(fa\log^5b)$ per dimension, with slope $\auhmain$, and it contains a sublattice of almost full
rank and slope at most $\auhmain-m\cdot\aurem$. We keep three theorems because their proofs are
independent and of different nature: an upper bound for a period (A), a second moment (B), and a
congruence statement (C).
\end{remark}

\subsection{The mechanism}\label{ss:au-mechanism}

\subsubsection*{The inequality}
The product formula says that an integral vector cannot be short: for a non-zero algebraic number
the archimedean size is at least the sum of the $\ell$-adic divisibilities.
Theorem~\ref{au:twistprime} is this statement for the Gram determinant of the lifts of a basis of
$\auT$. Its two sides are computed by two different means.

\emph{Upper side (Theorem A).} A target form of the shape $\auP_s\varphi$ can be lifted in a second
way. By the seesaw identity (Proposition~\ref{au:seesaw}) its image under the lower chain is the
restriction to $\auU(Y_m)$ of the image of $\varphi$ under the upper chain
$\auU(X_0)\to\dots\to\auU(X_m)$, all of whose groups are compact at infinity. The upper chain is
explicit: the value of a lift is a finite sum of values of $\varphi$ against Taylor coefficients
(jets) of theta functions of an abelian variety with complex multiplication by $E$, of orders
given by the rows of the weights: $a$ at each $a$-place and about $\kappa(\sigma)$ in $k$
directions at every place. Two things are known about such a jet of order $j$ at a place: all
its conjugates have size at most about $\exp(\ausinf(j))$, where $\ausinf(j)=\log j!+j\log(M/
\Omega)$ involves the Chowla--Selberg period $\Omega$ of $K$ (Proposition~\ref{au:kernel}; the curve
is $\C/\Omega_0\auO_K$, normalised with the admissible point of order two, \S\ref{pf:ss-adm}); and at
a non-split prime $p$ it is divisible by $p^{s_p(j,\sigma)}$ beyond the smooth amount
(Propositions~\ref{au:count} and~\ref{au:local}). Multiplying the lifted form by
$\aucB\approx\prod_pp^{-\augain_p}$ one obtains integral vectors whose squared norm has
logarithm at most
\[
m\,k\sum_\sigma\ausinf(\kappa(\sigma))+m\,\#\auPa\,\ausmain-2\sum_p\augain_p\log p=
\auhmain-m\cdot\aurem ,
\]
by \eqref{eq:au-cancel}, up to the errors. Here $m\,\#\auPa\,\ausmain$ replaces the true size
$m\,\#\auPa\,\ausinf(a)$; the difference is a multiple of $a$ by \eqref{eq:au-excess} and is
removed by the factor $\prod_\tau\tau(\aubar{\pi_N})^{-n_\tau}$ in the restriction to $Y_m$, which
is why the level is $N\approx(\bar\kappa M\aurd/2\pi)^m$.

\emph{Lower side (Theorem C).} The Gram determinant of the lifts of a basis is
$Z(\vartheta)\cdot\det\auGram(\text{basis})$, and $Z(\vartheta)$ is computed by a Rallis inner
product formula for each of the $m$ lifts: $Z$ factors as a product of determinants $Z_j$ of
single lifts $W_j\to V_j$ (Proposition~\ref{au:factor}), and each $Z_j(\vartheta)$ is, up to a
congruence number, the $\audT$-th power of a product of critical values of Hecke $L$-functions of
$E$ twisted by $\vartheta^j$ (Proposition~\ref{au:onelift}). The untwisted values have
archimedean size $k\sum_\sigma\ausinf(\kappa(\sigma))$ per lift. Their $\ell$-adic valuations are
not controlled; but for each $\ell$ one may choose a twist $\theta_\ell$ of conductor prime to
$\ell$ for which the twisted value has exactly the valuation forced by the divisibility of jets,
namely $k\sum_\sigma s_\ell(\kappa(\sigma),\sigma)$ (Propositions~\ref{au:regular}
and~\ref{au:hecke}). Thus $Z(1)/Z(\theta_\ell)$ has height at most $-\audT\,m\,k\,\auhdir=
-\audT\auhmain$, and Theorem~\ref{au:twistprime} gives covolume $+\,2\cdot$index
$\ge\audT\auhmain$ for every almost full family.

\emph{Comparison.} The family of Theorems A and B has $c+2\delta\le n(\auhmain-m\aurem)+$errors.
Hence $n\,m\,\aurem\le\audT\cdot$errors: \emph{a positive remainder would produce integral
vectors shorter than the product formula allows.} A negative remainder produces no contradiction,
which is why the result is one-sided.

\subsubsection*{Where the primes enter}
Only through $s_p$. The upper side uses divisibility of jets of two kinds: of order $a$ at the
$a$-places (the term $\sum_{\tau\in\auPa}s_p(a,\tau)$ of the gain) and of order $\kappa(\sigma)$
in $k$ further directions (the term $k\sum_\sigma s_p(\kappa(\sigma),\sigma)$). The lower side
uses only the second kind, through the valuations of twisted $L$-values. The second kind is
therefore charged on both sides and cancels: it is the content of $\auhmain$. The first kind
appears on the upper side only, and its sum over $p$ is $\sum_{\tau\in\auPa}\ausfin(a,\tau)=
\#\auPa\,\ausmain+\aurem$. The weights are chosen so that this isolates a progression: at an
$a$-place the partner at $\nu$ has weight $b=a/x$ if $\nu\equiv1\pmod{q_0}$ and $a/x_1$ otherwise,
so $(a/\nu-\kappa)^+$ equals $b(x/\nu-1)^+$ or $b(x/\nu-x/x_1)^+$, a sum over $\nu\le x$ in the
progression plus a sum over $\nu\le x_1$ that is evaluated by the prime number theorem. The factor
$1-\auchi(\nu)$ is forced, since split primes give no divisibility.

\subsubsection*{Why the main terms cancel}
The cancellation of $\auhmain$ is not a coincidence between two computations, since $\auhmain$ is
defined as size minus guaranteed divisibility, per direction; what must be true is that each side
attains it. On the upper side this is \eqref{eq:au-cancel}, an identity, together with the
archimedean bound and the divisibility with the same normalisations: the constant $c_1$ occurs in
$\Omega$ (size) and in $\ausmain$ (primes) and cancels in Lemma~\ref{au:gains}(c). On the lower
side it is a statement of sharpness: for a suitable twist the $L$-values are not more divisible
than the jets are known to be. The relative precision required is extreme: $\auhmain\ge x^4fa$,
while the admitted error is $fa\log^5b$ per dimension, so both sides must be exact to a factor
$1+O(x^{-4+\epsilon})$ in the exponent, to leading order in the weight $b$. Everything
polynomial in $x$ (class numbers, numbers of terms, ranks, discriminants) is negligible against
$b$; only terms linear in the weights matter.

\subsubsection*{What grows with $m$, and why $m\approx\sqrt x$}
The conclusion is $m\cdot\aurem\le$ error, and $\aurem\approx a\,\#\auPb\,R_{q_0}(x)$, so the bound
for $R_{q_0}(x)$ is the error divided by $m\,a\,\#\auPb$. If the error stayed $O(fa\log^5b)$ for
$m=x^{1/2+\delta}$ one would get $R_{q_0}(x)\le x^{-1/2-\delta+\epsilon}$, a one-sided bound
beyond what the generalised Riemann hypothesis gives; as explained in \S\ref{an:ssec-m}, it would
force a systematic vanishing of the twisted $L$-functions at the centre, which is not expected. The
argument does contain errors that grow faster than linearly in $m$. There are two, both quadratic, one on each side.

\emph{(1) The level prime, in Theorem A.} By \eqref{eq:au-excess} the level must satisfy
$\log N\approx m\log(\bar\kappa M\aurd/2\pi)\in[m\log b,25m\log b]$: each lift exceeds the expected
size by $a\log(\bar\kappa M\aurd/2\pi)$ at each $a$-place, and $a\log N$ pays for all $m$ lifts.
At the prime $N$ itself nothing is gained, and each of the $m$ lifts may lose up to $fb$ powers of
$N$ (Proposition~\ref{au:divis}(iii)). The cost is
\[
\auerr_N\cdot\log N=2\,m\,f\,b\cdot\log N\le60\,m^2\,f\,b\log b ,
\]
to be compared with the error unit $\aueu=f\,a\log^3b=x\,f\,b\log^3b$. It fits if and only if
$m^2\ll x\log^2b$ (Lemma~\ref{au:errsum}).

\emph{(2) The comparison operators, in Theorem C.} For each lift the determinant $Z_j(\vartheta)$
is related to $L$-values through an operator whose determinant at the trivial twist has all
conjugates at most $\exp(3\,\audT\,f\,m\,b\log^2b)$ (Proposition~\ref{au:onelift}). Over $m$ lifts this
is $3\,m^2f\,b\log^2b$ per dimension, against the allowance $f\,a\log^4b=x\,f\,b\log^4b$; again
$m^2\ll x\log^2b$ (proof of Proposition~\ref{au:Zj}).

With a general $m$ the two sides give $m\cdot\aurem\ll fa\log^5b+m^2fb\log^2b$, that is,
\[
R_{q_0}(x)\ \ll_{q_0}\ \frac{(\log b)^6}{m}+\frac{m\,(\log b)^3}{x},
\]
which is optimal at $m\approx\sqrt x$ and gives $x^{-1/2+\epsilon}$; for larger $m$ the bound
deteriorates and is trivial at $m\approx x$. All other errors are linear in $m$ per lift, hence
bounded by the error unit: the slack of the rows ($m\cdot O(fb\log b)$, which needs only
$m\le x\log^2b$), the non-split primes $p<x_1$ ($\sum_p\auerr_p\log p\ll fa\log^2b\log^2x_1$, the
factor $m$ at the primes $p\equiv\pm1\pmod m$ being offset by their rarity), the denominators of
the projections to lifts (budget $\aueu/(mk)$ per lift and per unit of rank), and the ratios of
$L$-values ($m^2t\,x_1\le b$).

\begin{remark}\label{au:barrier-remark}
The first quadratic term has a clear origin: the lifts are not independent, because one level
$N$ serves all of them, its size is proportional to $m$, and the price of this prime of bad
reduction is paid by each of the $m$ lifts. The two estimates above are the inequality
$2mfb\log N\le60\,m^2fb\log b\le240\,fa\log b$ in the bound for the losses (Lemma~\ref{au:errsum}) and the
inequality $3fmb(\log b)^2+10mx_1\le k\,\aubud(\log b-8)$ in the cost of one lift (proof of
Proposition~\ref{au:Zj}).
\end{remark}

\subsection{Deduction of Theorems A and B}\label{ss:au-proofAB}

We first fix the two constants of Theorem~\ref{thm:srcA}. For a commutative ring $\auO$, modules
$S\subseteq M$ and $z\in\auO$ we write $\aucol(M,S)\mid z$ if for every finite family
$\varphi_1,\dots,\varphi_n$ in $M$ there is $\mu\in M_n(\auO)$ with $\det\mu\mid z$ and
$\sum_k\mu_{ik}\varphi_k\in S$ for all $i$ (over a discrete valuation ring and for finitely
generated $M$ this says that the length of $M/S$ is at most that of $\auO/z$). The \emph{losses}
are
\begin{equation}\label{eq:au-err}
\auerr_p=\begin{cases}
2f(\log b)^2\cdot\dfrac ap\cdot\epsilon_m(p)&p\in\mathcal N,\quad\epsilon_m(p)=m\text{ if }p\equiv
\pm1\!\!\pmod m,\ \epsilon_m(p)=1\text{ otherwise},\\[4pt]
2\,m\,f\,b&p=N,\\
0&\text{otherwise.}
\end{cases}
\end{equation}
Let $c_{\mathrm{ss}}$ be the constant of Proposition~\ref{au:seesaw} for the point $g_0$ and the
linear form $\mathsf t$ in the definition of $\auP$, and put
\[
g_\ell=\lfloor\augain_\ell-\auerr_\ell\rfloor\in\Z,\qquad\aucB=\prod_{\ell\in\mathcal N\cup\{N\}}
\ell^{-g_\ell} .
\]

\begin{lemma}[The losses are affordable]\label{au:errsum}
$\sum_{p\in T}\auerr_p\log p\le\aueu$ for every finite set $T$ of primes, and
$\log\aucB\le-\sum_{p\in\mathcal N}\augain_p\log p+\aueu+b$.
\end{lemma}
\begin{proof}
For $p\in\mathcal N$ with $p\not\equiv\pm1\pmod m$ the term is at most
$2fa(\log b)^2\log x_1\cdot p^{-1}$. If $p\equiv\pm1\pmod m$, write $p=mc\pm1$ with $c\ge1$; then
$m\le2p/c$ and the term is at most $4fa(\log b)^2\log x_1\cdot c^{-1}$, and each $c\le x_1$
occurs for at most two primes. As $\sum_{n\le x_1}1/n\le1+\log x_1$, these terms add up to at most
$10fa(\log b)^2\log x_1(1+\log x_1)$. The term at $N$ is $2mfb\log N\le60\,m^2fb\log b\le
240\,fa\log b$, by \eqref{eq:au-sizes} and $m^2\le4x$. By \eqref{eq:au-scales},
$(\log x_1)^2\le(\log b)^{1/3}$, and the total is at most $fa(\log b)^3$. For the second claim,
$-g_\ell<-\augain_\ell+\auerr_\ell+1$, $\augain_N=0$ (Lemma~\ref{au:gains}(a)), and
$\sum_{\ell\in\mathcal N\cup\{N\}}\log\ell\le x_1^2+\log N\le b$.
\end{proof}

The inputs are the following; each is stated precisely and discussed in
\S\S\ref{ss:au-inputsA}--\ref{ss:au-inputsB}.
\begin{itemize}
\item[(I1)] \emph{Seesaw} (Proposition~\ref{au:seesaw}): every twist $\vartheta$ has a chain datum
$\Psi$ with $\iota_\vartheta(\auP\varphi)=c_{\mathrm{ss}}B_{\vartheta,\Psi}\varphi$ for all
$\varphi\in\auSram$.
\item[(I2)] \emph{Divisibility} (Proposition~\ref{au:divis}): for every prime $\ell$ and
$j\colon\C\to\C_\ell$ there are $i\le\auerr_\ell\auDX$ and $S\subseteq\auS_j\cap\auSram$ with
$\aucol(\auS_j\cap\auSram,S)\mid\ell^i$, such that $B_{\vartheta,\Psi}\varphi$ has gauge at most
$\ell^{\auerr_\ell-\augain_\ell}$ at $j$ for all $\varphi\in S$, all twists $\vartheta$ with
$\ell\nmid\auNm Q_\vartheta$ and all chain data $\Psi$ of $\vartheta$.
\item[(I3)] \emph{Short forms span} (Proposition~\ref{au:short}): $\auSram$ is spanned by forms
that are integral at every $j$ and all of whose values have all conjugates of squared norm at most
$e^{\aueu}$.
\item[(I4)] \emph{Size} (Proposition~\ref{au:size}): if $\Psi$ is a chain datum of the trivial
twist and the values of $\sigma\varphi$ have squared norm at most $e^B$, then the values of
$\sigma(B_{1,\Psi}\varphi)$ have squared norm at most $\exp(B+\auhmain-m\aurem+2\sum_{p\in
\mathcal N}\augain_p\log p+\aueu)$.
\item[(I5)] \emph{The good part} (Proposition~\ref{au:good}): passing from $z\in\auAz$ to
$\#\auG\varrho(e_0)\varrho(s)z$ keeps the integrality of $c\cdot\iota_\vartheta z$ at every $j$,
and enlarges a bound $e^B$ for the values of $\sigma(\iota_1z)$ to the bound $\exp(x^{99}+B)$ for
$\|\iota_1(\cdot)\|_\sigma^2$.
\item[(I6)] \emph{Count} (Lemma~\ref{au:ramcount}): $\dim\auSram\ge(1-b^{-8})\auDX$.
\item[(I7)] \emph{Spreading} (Proposition~\ref{au:spread} and Lemma~\ref{au:few}): $\audT>0$, and
for every subspace $F\subseteq\auT$ there is $s\in\auG$ with
$\dim\auP_s^{-1}(F)\le\bigl(\tfrac12(1+\dim F/\audT)+\tfrac3{16b^2}\bigr)\auDX$.
\end{itemize}

\begin{proof}[Proof of Theorem~\ref{thm:srcA} from (I1)--(I6)]
Put $\auP'_s=\auP_s\circ c_{\mathrm{ss}}^{-1}$ on $\auSram$, so that $s\aupp=\aucB\auP'_s$, since
$e_0$ is central. By (I1), every twist $\vartheta$ has a chain datum $\Psi$ with
\begin{equation}\label{eq:au-ssused}
\iota_\vartheta\bigl(\auP(c_{\mathrm{ss}}^{-1}\varphi)\bigr)=B_{\vartheta,\Psi}\varphi\qquad
(\varphi\in\auSram).
\end{equation}

(a) By (I3) it suffices to treat a form $\varphi\in\auSram$ that is integral everywhere, so
$\varphi\in\auSram_{\mathrm{int}}$, and whose values have all conjugates of squared norm at most
$e^{\aueu}$. Fix $s$ and $\sigma$, and let $\beta=\auhmain-m\aurem+2\sum_{\mathcal N}\augain_p
\log p+\aueu$. By \eqref{eq:au-ssused} for $\vartheta=1$ and (I4), the values of the
$\sigma$-conjugate of $\iota_1(\auP(c_{\mathrm{ss}}^{-1}\varphi))$ have squared norm at most
$\exp(\aueu+\beta)$. By (I5), $\|\Theta_1(\auP'_s\varphi)\|_\sigma^2\le\exp(x^{99}+\aueu+\beta)$.
Multiplying by the rational number $\aucB$ and using Lemma~\ref{au:errsum},
\[
\log\|\Theta_1(s\aupp\,\varphi)\|_\sigma^2\le2\Bigl(-\sum_{\mathcal N}\augain_p\log p+\aueu+b
\Bigr)+x^{99}+\aueu+\beta=\auhmain-m\cdot\aurem+x^{99}+4\aueu+2b ;
\]
the gains have cancelled, and $x^{99}+4\aueu+2b\le fa(\log b)^4$.

(b) Fix $\iota$ with prime $\ell$, and let $i$, $S$ be as in (I2) for $\iota$. We claim
$s\aupp(S)\subseteq\auLat_{D,\iota}$ for every $s\in\auG$. Let $\varphi\in S$ and let $\vartheta$
be an $m$-fold twist with $\ell\nmid\auNm Q_\vartheta$. By \eqref{eq:au-ssused} and (I2),
$\iota_\vartheta(\auP(c_{\mathrm{ss}}^{-1}\varphi))$ has gauge at most
$\ell^{\auerr_\ell-\augain_\ell}$. The $\ell$-adic absolute value of $\aucB$ is $\ell^{g_\ell}$
(if $\ell\notin\mathcal N\cup\{N\}$ then $\augain_\ell=\auerr_\ell=g_\ell=0$), and
$g_\ell\le\augain_\ell-\auerr_\ell$; so $\aucB\,\iota_\vartheta(\auP(c_{\mathrm{ss}}^{-1}\varphi))$
is integral at $\iota$, and by (I5) so is $\Theta_\vartheta(s\aupp\,\varphi)$. This proves the
claim. Now let $y_1,\dots,y_n\in s\aupp(\auSram_{\mathrm{int}})$, $y_k=s\aupp\,\varphi_k$. The
$\varphi_k$ lie in $\auS_\iota\cap\auSram$, so there is $\mu\in M_n(\auO_\iota)$ with
$\det\mu\mid\ell^i$ and $\sum_k\mu_{ik}\varphi_k\in S$; then $\sum_k\mu_{ik}y_k\in\auLat_{D,\iota}$
and $|\det\mu|_\iota\ge\ell^{-i}\ge\ell^{-\auerr_\ell\auDX}$. By Lemma~\ref{au:errsum} and (I6),
$\sum_\ell\auerr_\ell\auDX\log\ell\le\auDX\aueu\le2\dim\auSram\cdot\aueu$. So a family contained
in one $s\aupp(\auSram_{\mathrm{int}})$ has index at most $2\dim\auSram\cdot\aueu$; a family
contained in the union over $s\in\mathsf t$ is a union of $\#\mathsf t$ such families, and indices
add.
\end{proof}

\begin{proof}[Proof of Theorem~\ref{thm:srcB} from (I6), (I7)]
Let $\varepsilon_0=(8b^2)^{-1}$. By (I7) the maps $\auP_s$, $s\in\auG$, spread with error
$\tfrac32\varepsilon_0$ in the sense of Lemma~\ref{au:few}, which gives a finite
$\mathsf t\subseteq\auG$ with
\[
\#\mathsf t\cdot\auDX\le\tfrac52\bigl(1+\log(b^2/2)\bigr)\audT\le5\,\audT\log(1+2b^2),\qquad
\dim\sum_{s\in\mathsf t}\auP_s(\auS)\ge\Bigl(1-\frac2{b^2}\Bigr)\audT .
\]
Since $\dim\auSram\le\auDX$ the first claim follows. Restricting each $\auP_s$ to $\auSram$ lowers
the dimension of the sum of the images by at most $\#\mathsf t(\auDX-\dim\auSram)\le
b^{-8}\cdot5\,\audT\log(1+2b^2)\le2\audT/b^2$, by (I6). Finally
$s\aupp(\auSram)=\auP_s(\auSram)$, the scalars being non-zero.
\end{proof}

\subsection{The inputs of Theorem A}\label{ss:au-inputsA}

\subsubsection*{The seesaw identity}

\begin{proposition}[Seesaw identity for the two chains]\label{au:seesaw}
For every point $g_0\in G_0(E^+\otimes\R)$ and every linear form $\mathsf t$ on $W_{\lambda^X}$
there is $c_{\mathrm{ss}}\in\C^\times$ such that every twist $\vartheta$ has a chain datum $\Psi$
with
\[
\iota_\vartheta(\auP\varphi)=c_{\mathrm{ss}}\cdot B_{\vartheta,\Psi}\varphi\qquad\text{for all }
\varphi\in\auSram .
\]
The constant does not depend on $\vartheta$. The statement does not assert that either side is
non-zero.
\end{proposition}
\begin{proof}
The proof is given in \S\ref{ps:ss-seesaw}, with the complements of \S\ref{pk:sec}:
Proposition~\ref{ps:onestep} is proved in \S\ref{pk:ss-onestep}, Lemma~\ref{ps:definite} is
Proposition~\ref{pk:definite}, and the identity $\chi_0=1$ is Lemma~\ref{pk:chi0}. It uses the standard
facts (W0)--(W5) on the Weil representation (Lemma~\ref{ps:weil}; \cite{WeilActa},
\cite{KudlaSplitting}, \cite{HKS}, and \S\ref{pwf:sec}), the fact (W6) of \S\ref{pk:ss-conv}, which is
Proposition~\ref{pf:flip}, and Witt's theorem. It also uses that the lifts of the lower chain
and their transposes preserve the archimedean type of the chain and that these forms are cusp
forms; this is Proposition~\ref{pz:types}, deduced from three statements (A1)--(A3) on
discrete series of real unitary groups (\S\ref{pz:ss-types}): (A1) is Theorem~\ref{par:A1}, (A2) is
Corollary~\ref{par:A2}, and (A3) is Theorem~\ref{par:A3}, proved from the published results
(D1)--(D3) of \S\ref{par:ss-rigid} and the normalisation (V), which is Theorem~\ref{pco:V}.
\end{proof}

\noindent\emph{Outline.} Pair $\iota_\vartheta(\auP\varphi)$ with a form $\eta$ on the definite
group $\auU(Y_m)$ and move the theta kernels to $\eta$: the pairing becomes the Bessel period of
the $m$-fold descent of $\eta$ to $\auU(Y_0)$ against $\mathcal W_\varphi$. The Bessel subgroup of
$\auU(Y_j)=\auU(X_j\perp S_j)$ attached to the string $S_j$ has reductive part $\auU(X_j)$, and one
theta lift carries the Bessel (for $S_j$ of odd rank) or Fourier--Jacobi (even rank) period for
$(Y_{j+1},S_{j+1})$ to the one for $(Y_j,S_j)$ of a theta lift on the smaller groups
$\auU(X_j)\subset\auU(X_{j+1})$; this is the transfer of generalised Whittaker models under the
theta correspondence (Gomez--Zhu~\cite{GomezZhu} locally; the global statement is a seesaw in the
sense of Kudla~\cite{KudlaSeesaw} for $\auU(Y_j)\supset\auU(X_j)\times\auU(S_j)$ combined with an
unfolding along the unipotent radical). After $m$ steps the string is exhausted and one is left
with the upper chain applied to $\varphi$, evaluated on $Y_m\subset X_m$. Each unfolding produces
an integral over a real group, which gives a constant independent of $\vartheta$ and of
$\varphi$ (and this is the origin of $c_{\mathrm{ss}}$), times a finite sum over cosets of a
level, weighted by reciprocals of stabiliser indices that divide powers of $\auNm Q_\vartheta$;
this sum defines the Schwartz data of the upper chain, which are combinations of translates of the
vacuum, supported at the places dividing $M\auNm Q_\vartheta$, with coefficients integral away
from $\auNm Q_\vartheta$. This is the chain datum. The hypothesis that $\varphi$ is ramified
enters through the data at the primes above $N$ (the gluing, \S\ref{ps:ss-glue}). The
convergence of the unfolded integrals is part of Proposition~\ref{ps:onestep}; the constant is
independent of $\vartheta$ by Lemmas~\ref{ps:double} and~\ref{ps:fin}; and the invariance of the
resulting data under the levels $J_{X_j}$ is proved in Step~5 of the proof in \S\ref{ps:ss-seesaw}.

\subsubsection*{Divisibility of the composite lift}

For a prime $p$ put $\mathfrak r_p=\lceil\log_pm\rceil$. A prime $p$ is \emph{exceptional} if it
is inert in $K$, $p\not\equiv\pm1\pmod m$, and the order of $-p$ in $(\Z/q_0)^\times$ is odd and
at most $\mathfrak r_p$; then $u(-p^e)$ is an $a$-place for some odd $e$ with $p^{e-1}<m$, so
that the partner of an $a$-place at $p^e$ is again an $a$-place (a \emph{collision}).

\begin{proposition}[Divisibility of the composite lift]\label{au:divis}
Let $\ell$ be a prime and $j\colon\C\to\C_\ell$. There are $i\in\N$ and an $\auO_j$-submodule
$S\subseteq\auS_j\cap\auSram$ with $i\le\auerr_\ell\cdot\auDX$ and
$\aucol(\auS_j\cap\auSram,S)\mid\ell^{\,i}$ such that for every twist $\vartheta$ with
$\ell\nmid\auNm Q_\vartheta$, every chain datum $\Psi$ of $\vartheta$ and every $\varphi\in S$ the
form $B_{\vartheta,\Psi}\varphi$ has gauge at most $\ell^{\,\auerr_\ell-\augain_\ell}$ at $j$,
with $\auerr$ as in \eqref{eq:au-err}. More precisely:
\begin{enumerate}
\item[(i)] if $\ell\ne N$ and $\ell$ splits in $K$ or $\ell\ge x_1$, then $S=\auS_j\cap\auSram$
and the composite lift preserves integrality (on all of $\auS_j$: the hypothesis that $\varphi$ is
ramified is not used);
\item[(ii)] if $\ell\in\mathcal N$ is not exceptional, then $S=\auS_j\cap\auSram$ (and the bound
holds on all of $\auS_j$);
\item[(iii)] if $\ell=N$, the composite lift loses at most $2mfb$ powers of $N$ on the ramified
forms $\varphi$ such that $U_N^{m/2}\varphi$ is integral for the lattice $L_0$; the operator $U_N$
\eqref{pl:eq-U} and the lattice $L_0$ are defined in \S\ref{pl:ss-level}, at the primes of $E$
above $N$ conjugate to those induced by $j$ (Lemma~\ref{pl:level});
\item[(iv)] if $\ell\in\mathcal N$ is exceptional, $S=(Y_\Lambda^{\,m}\auS_j+\ell^{\,j_\ell}\auS_j)
\cap\auSram$, where $j_\ell=\lceil2mfa/\ell\rceil$ and $Y_\Lambda$ is the product of the normalised
Hecke operators $Y_{\mathfrak l}$ at the primes $\mathfrak l$ of a full mark $\Lambda$ of primes of
$E$ above $\ell$ (\S\ref{pa:ss-divis}, Theorem~\ref{pa:divis}(iv)); the operators $Y_{\mathfrak l}$
are not claimed to preserve $\auSram$.
\end{enumerate}
\end{proposition}
\begin{proof}
This is Theorem~\ref{pa:divis}, proved in \S\ref{pa:ss-divis}; its last assertion is the
first assertion here. The ingredients are: Proposition~\ref{au:count} (\S\ref{pa:ss-count}); the
normal form of the theta products (Theorem~\ref{pa:onecurve}, Proposition~\ref{pa:product}), from
published results on elliptic curves with complex multiplication; Lemma~\ref{pa:kernel}, proved in
\S\ref{pk:ss-kernel} and, for the integrality of the jets at the primes $\ell\mid2q_1M$ and at
$\ell=N$ (Lemma~\ref{pl:kernelN}), in Corollary~\ref{pf:K1}; for case (iv), Lemma~\ref{pa:hecke}
(\S\ref{pl:ss-hecke}) and Lemma~\ref{pa:eldiv} (\S\ref{ph:ss-ded}, from Theorem~\ref{ph:thm}, whose
three inputs are proved in \S\ref{pe:sec} and whose normalisation and use of weak approximation are
justified in Remark~\ref{pl:N} and Corollary~\ref{pl:wacor}); for case (iii), Lemma~\ref{pl:level}.
\end{proof}

\noindent\emph{Outline.} A value of the upper chain is a finite sum, with coefficients integral
away from $\auNm Q_\vartheta$ (Proposition~\ref{au:expansion}), of products of $m$ jets of theta
series of the lattices $\auHom(\auO_E^{n_i},\auO_E^{n_{i+1}})$ translated by torsion points.
(a) \emph{Normal form.} The Taylor series at the origin of such a theta series, in coordinates
scaled by the period, has the normal form \eqref{eq:au-normal} at a non-split prime: this is a
statement about the formal group of an abelian variety with multiplication by $\auO_E$ and its
canonical theta function (compare \cite{BannaiKobayashi}
for one elliptic curve). (b) \emph{Count.} Proposition~\ref{au:count} then gives the valuation of
each jet, and the sum over the $m$ lifts of the predicted divisibilities is at least
$\augain_p$ minus $m$ times the sum of $a/p^e$ over the collisions $(\tau,e)$, minus an error.
(c) If $p$ is not exceptional, the collisions are few enough (by the definition of exceptional primes,
they occur only for $p^e$ at least about $m$, or for $p\equiv\pm1\pmod m$) that the loss
\eqref{eq:au-err} pays for them. (d) If $p$ is exceptional the
collisions cost up to $m\cdot fa/p$, which is not affordable; the lost divisibility is recovered on
the image of suitable Hecke operators at primes above $p$, applied once per lift and moved through
the lifts, and the theorem on elementary divisors shows that this image, enlarged by
$p^{j_p}\auS_j$, has colength at most $p^{\auerr_p\auDX}$. (e) At $\ell=N$ the data are those of
the level; denominators are traded for the operator $U_N$ and each lift loses at most $fb$ powers
of $N$. (f) At split primes and at primes $\ge x_1$ the jets are integral and nothing is claimed
beyond integrality.

\subsubsection*{Short forms span}

\begin{lemma}[Short bases]\label{au:shortbasis}
Let $n\ge1$ and let $c_0$ be a positive integer in the different of $E$. For every
$g\in\auU(E^n)(\auA_{E^+,f})$ there is $Y\in\auGL_n(E)$ with $g^{-1}Y\in M_n(\widehat{\auO}_E)$
all of whose entries have all conjugates of absolute value at most $C_{\mathrm{sb}}=6c_0(fn)^2$.
\end{lemma}
\begin{proof}
This is proved in \S\ref{ps:ss-short}, with the constant $6c_0(fn)^2$, by an elementary
transference argument (Lemma~\ref{ps:transfer}); Banaszczyk's theorem~\cite[Theorem~2.1]{Banaszczyk} would give
a better constant and is not needed.
\end{proof}

\begin{proposition}[Short ramified forms span]\label{au:short}
$\auSram$ is spanned by the forms $\varphi\in\auSram$ that are integral at every
$j\colon\C\to\C_\ell$, for every $\ell$, and satisfy $\|\sigma(\varphi(g))\|^2\le e^{\aueu}$ for
all $\sigma\in\auAut(\C)$ and all $g$.
\end{proposition}
\begin{proof}
\emph{Step 1: $\auS$ is spanned by forms as in the statement with the bound $e^{\aueu}/4^f$.}
The level $J_X$ is free, so $\auS$ is the direct sum, over the classes $x\in\Gamma\backslash G_f/
J_X$, of the spaces of forms supported on $\Gamma xJ_X$, and $\xi\mapsto\delta_{x,\xi}$, with
$\delta_{x,\xi}(\gamma xk)=\gamma\xi$, identifies $W_{\lambda^X}$ with the forms supported on the
class of $x$. The form $\delta_{x,\xi}$ is integral at $j$ if and only if $\xi$ has gauge at most
$1$ at $(j,x)$, and its values are the vectors $\gamma\xi$, whose conjugates have the norm of
those of $\xi$. So it suffices to span $W_{\lambda^X}$ by vectors $\xi$ with gauge at most $1$ at
$(j,x)$ for every $j$ and with all conjugates of squared norm at most $e^{\aueu}/4^f$. Let $Y$ be
as in Lemma~\ref{au:shortbasis} for $g=x$ and $c_0=q_1d$, and let $Y$ act on $W_{\lambda^X}$
by $\bigotimes_\tau\rho_\tau(\tau(\aubar Y))$, the formula defining the action of $\Gamma$, extended
to $\auGL_k(E)$ (that is, polynomially through the components $Y_{w_\tau}$ at the primes $w_\tau$
of (N2), the same components through which $x_\ell^{-1}$ acts); let $\mathrm{ev}(\zeta)\in W_{\lambda^X}(\Z[\zeta])$ be the evaluation of
the products of minors at a family $\zeta$ of matrices whose entries are roots of unity. The
vectors $\mathrm{ev}(\zeta)$ span $W_{\lambda^X}$ (a polynomial vanishing at all points with
coordinates roots of unity is zero), and $Y$ is invertible, so the vectors $\xi=Y\mathrm{ev}
(\zeta)$ span. They have gauge at most $1$ at $(j,x)$: under $j$ the action of $Y$ is through its images
$Y_{w_\tau}\in\auGL_k(E_{w_\tau})$, so $x_\ell^{-1}\cdot j(Y\mathrm{ev}(\zeta))=\bigotimes_\tau
S_{\lambda^X_\tau}\bigl((x^{-1}Y)_{w_\tau}\bigr)\,j(\mathrm{ev}(\zeta))$, and $(x^{-1}Y)_{w_\tau}\in
M_k(\auO_{w_\tau})$ acts by matrices with integral entries in the basis of products of minors (the weight $\lambda^X$ has no
determinant part). The coordinates of
$\sigma(\xi)$ are polynomials of degree at most the number of boxes, at most $fa$, in the entries
of $\sigma Y$ and of $\sigma\zeta$, with integer coefficients, and the entries have absolute value
at most $C_{\mathrm{sb}}$; by Lemma~\ref{pz:short}, $\|\sigma\xi\|^2\le C_{\mathrm{sb}}^{\,5fa}$, and
$4^fC_{\mathrm{sb}}^{\,5fa}\le e^{\aueu}$, because
$\log C_{\mathrm{sb}}=\log(6q_1d(fk)^2)\le100\log x$ (Lemma~\ref{pz:sizes}), $\log x\le\log b$, and
$\aueu=fa(\log b)^3$.

\emph{Step 2.} For $k\in K^\circ$, right translation by $k$ preserves integrality at every $j$
(the gauge at $(j,g)$ and at $(j,gk)$ agree) and the norms of the conjugates of the values. Hence
$1-z_w$ maps a form as in Step 1 with bound $B$ to one with bound $4B$, and
$\prod_{w\mid\pi_N}(1-z_w)$, a product of $f$ factors, maps the spanning set of Step 1 into the
set of the proposition. Its image is $\auSram$.
\end{proof}

\subsubsection*{The archimedean size of the composite lift}\label{ss:au-size}

Since both groups of a pair $(\auU(X_j),\auU(X_{j+1}))$ are compact at infinity, the Weil
representation at an infinite place, restricted to them, is the space of polynomials on
$\auHom(X_j,X_{j+1})\otimes_{E,\tau}\C$ (Fock model), in which the two weights occur jointly with
multiplicity one (Howe~\cite{HoweTheta}, Kashiwara--Vergne~\cite[Ch.~III, Proposition~5.6]{KashiwaraVergne}). The theta
kernel for a Schwartz--Bruhat function $\phi$ is therefore a function
$\theta_j(\phi;g,g')\in\auHom(W_{\lambda(X_j)},W_{\lambda(X_{j+1})})$ of the finite adelic
variables whose matrix entries are jets $Q(\partial)\vartheta_{\phi,g,g'}(0)$ of theta series of
positive definite hermitian $\auO_E$-lattices along the entries $Q$ of a Schur matrix, and
$(\Theta^X_j(\phi)\varphi)(g')=\int_{\Gamma\backslash G_f}\theta_j(\phi;g,g')\varphi(g)\,dg$ for the
Haar measure giving the level $J_{X_j}$ the measure $1$ (\S\ref{pl:ss-norm}, (N4)); this is
$\sum_{g\in\Gamma\backslash G_f/J_{X_j}}\theta_j(\phi;g,g')\varphi(g)$ when $\phi$ and $\varphi$ are
fixed by $J_{X_j}$, the level being free. The lengths of the rows of the diagram of the $j$-th lift are: $a$ in
the first row at the $a$-places; $\kappa(\tau)$, up to a \emph{slack}, in the next $k-1$ rows at
every $\tau$; and slack only in the remaining rows; the total slack at one place is at most
$b+a/x_1+n_m\Delta^2/4$ boxes. (We count this description of the rows as part of
the statement of Proposition~\ref{au:kernel}.)

\begin{proposition}[One kernel]\label{au:kernel}
Let $0\le j<m$, let $\phi$ be a translate of the vacuum and $\sigma\in\auAut(\C)$. Then for all
$g,g'$ and all $\xi$,
\[
\|\sigma(\theta_j(\phi;g,g')\xi)\|\le\exp\bigl(x^{50}+\tfrac12\mathrm{rs}_j\bigr)\,\|\sigma(\xi)\|,
\qquad\mathrm{rs}_j=\sum_{\tau\in\auPl}\ \sum_{\text{rows }i}\ausinf(\text{length of row }i
\text{ at }\tau).
\]
\end{proposition}
\begin{proof}
The proof is given in \S\ref{ps:ss-kernel}. It rests on Lemma~\ref{ps:conj}, which is
Proposition~\ref{pk:conj} together with the reciprocity law for the translates by torsion points of
the theta function of one curve with complex multiplication (Proposition~\ref{pk:curve}(C3)(ii),
proved as Theorem~\ref{pf:recip}). The other ingredients are the expression of a conjugate jet as a
Gaussian integral, the formula $|t_\sigma|^2=1/\Omega$ (Proposition~\ref{pk:curve}(C4), from the
Chowla--Selberg formula~\cite{ChowlaSelberg}), and the uniform bound with the explicit constant
$C_0(q_0,d)$ of \eqref{pz:eq-C0}, which is at most $\log\log b$ by Lemma~\ref{pz:C0}.
\end{proof}

\begin{proposition}[Restriction to $Y_m$]\label{au:restrict}
For every form $\psi$ on $\auU(X_m)$ of weight $\lambda(X_m)$, every $\sigma\in\auAut(\C)$ and
every $y\in\auU(Y_m)(\auA_{E^+,f})$,
$\|\sigma((\operatorname{res}\psi)(y))\|^2\le\exp\bigl(x^{70}-(\sum_\tau n_\tau)\log N\bigr)\cdot
\|\sigma(\psi(\operatorname{diag}(y,1)))\|^2$.
\end{proposition}
\begin{proof}
The proof is given in \S\ref{ps:ss-restrict}, for the branching projection
$\mathrm{br}=c_{\mathrm{br}}\,i^*$ of Lemma~\ref{ps:branching}; part (c) of that lemma (the
integrality of the matrix, with $c_{\mathrm{br}}\le L^{3L}$) is Proposition~\ref{pk:branching}.
It uses the irreducibility of the Schur modules~\cite{FultonYT} and standard
monomial theory~\cite{DEP}. The inequality \eqref{eq:au-branch} holds with equality, see
\eqref{ps:eq-ntau}. That the restricted form remains integral at the primes above $N$ after
division by the powers of $\aubar{\pi_N}$ is part of Proposition~\ref{au:divis}(iii).
\end{proof}

\begin{proposition}[Expansion of a chain datum]\label{au:expansion}
Let $\Psi$ be a chain datum of the twist $\vartheta$. There is a finitely supported function
$\Psi'$ on tuples of translations with $B_{\vartheta,\Psi}=\operatorname{res}\circ C_{\vartheta,
\Psi'}$ (no gluing), whose values are integral at the primes not dividing $N\auNm Q_\vartheta$,
and with $\sum_x|\sigma\Psi'(x)|\le\exp(x^{60})$ for all $\sigma\in\auAut(\C)$ if
$Q_\vartheta=\emptyset$.
\end{proposition}
\begin{proof}
The proof is given in \S\ref{ps:ss-glue}, for the gluing and the chain data defined
there by (C1)--(C4).
\end{proof}

\begin{proposition}[Size of the composite lift]\label{au:size}
Let $\Psi$ be a chain datum of the trivial twist, $\sigma\in\auAut(\C)$, $\varphi\in\auS$ and
$B\in\R$ with $\|\sigma(\varphi(g))\|^2\le e^B$ for all $g$. Then for all $y$,
\[
\|\sigma((B_{1,\Psi}\varphi)(y))\|^2\le\exp\Bigl(B+\auhmain-m\cdot\aurem+2\sum_{p\in\mathcal N}
\augain_p\log p+\aueu\Bigr).
\]
\end{proposition}
\begin{proof}
Let $\Psi'$ be as in Proposition~\ref{au:expansion}. The levels of the $X_j$ are free
(Lemma~\ref{au:free}) and have at most $\exp(x^{30})$ classes each (Proposition~\ref{pf:classes}, Lemma~\ref{pz:sizes}). A value of $C_{1,\Psi'}\varphi$ is the sum, over the tuples $x$ in the
support of $\Psi'$ and over the tuples of classes of $X_0,\dots,X_{m-1}$, of $\Psi'(x)$ times a
product of $m$ kernels applied to a value of $\varphi$. By Proposition~\ref{au:kernel},
\[
\|\sigma((C_{1,\Psi'}\varphi)(g))\|\le\exp(x^{30})^m\cdot\exp(x^{60})\cdot\prod_{j<m}\exp\bigl(
x^{50}+\tfrac12\mathrm{rs}_j\bigr)\cdot e^{B/2}.
\]
Each box of slack costs at most $\log(2a+1)+\log(M/\Omega)\le25\log b$ in $\ausinf$, so by the
description of the rows
\[
\sum_{j<m}\mathrm{rs}_j\le m\Bigl(\#\auPa\,\ausinf(a)+k\sum_\tau\ausinf(\kappa(\tau))\Bigr)+
25\,m\,f\log b\,\bigl(b+a/x_1+n_m\Delta^2/4\bigr).
\]
By Proposition~\ref{au:restrict} and \eqref{eq:au-branch}, the logarithm of the left side of the
claim is at most
\[
B+m\,k\sum_\tau\ausinf(\kappa(\tau))+\#\auPa\bigl(m\,\ausinf(a)-a\log N\bigr)+\mathcal E,
\]
\[
\mathcal E=x^{70}+\tfrac12fm\Delta\log N+2mx^{30}+2x^{60}+2mx^{50}+25mf\log b\,(b+a/x_1+
n_m\Delta^2/4).
\]
By \eqref{eq:au-excess}, $\#\auPa(m\ausinf(a)-a\log N)\le m\#\auPa\ausmain+32fa\log b$, and by
\eqref{eq:au-cancel} the main terms are $\auhmain-m\aurem+2\sum\augain_p\log p$. Finally
$\mathcal E+32fa\log b\le\aueu$ by comparison of sizes; the largest terms are $25mfb\log b\le
50fa\log b/\sqrt x$ and $32fa\log b$.
\end{proof}

\subsubsection*{The good part}

\begin{proposition}[Equivariance of the lower chain]\label{au:equiv}
For every tame character $\vartheta$ and every $u\in\auG$ there is $\beta_\vartheta(u)$ in the
stabiliser $K^\circ$ of the lattice of $Y_m$ such
that $(\iota_\vartheta(\varrho(u)z))(g)=(\iota_\vartheta z)(g\beta_\vartheta(u))$ for all
$z\in\auAz$ and all $g$.
\end{proposition}
\begin{proof}
The proof is given in \S\ref{ps:ss-Y} (after Lemma~\ref{ps:level}); it uses the standard facts (W1),
(W4), (W5) on the Weil representation (Lemma~\ref{ps:weil}; \cite{WeilActa},
\cite{KudlaSplitting}, \cite{HKS}, and \S\ref{pwf:sec}). The element $\beta_\vartheta(u)$ depends on $\vartheta$ through
the places of $Q_\vartheta$; Proposition~\ref{au:good} uses it for one $\vartheta$ at a time.
\end{proof}

\begin{proposition}[Passage to the good part]\label{au:good}
Let $z\in\auAz$, $s\in\auG$, and $z'=\#\auG\,\varrho(e_0)\varrho(s)z$.
\begin{enumerate}
\item[(a)] If $c\in\C$ and $c\cdot\iota_\vartheta z$ is integral at $j\colon\C\to\C_\ell$, so is
$c\cdot\iota_\vartheta z'$.
\item[(b)] If $\sigma\in\auAut(\C)$ and all values of $\sigma(\iota_1z)$ have squared norm at most
$e^B$, then $\|\iota_1z'\|_\sigma^2\le\exp(x^{99}+B)$.
\end{enumerate}
\end{proposition}
\begin{proof}
Write $\#\auG\,e_0=\sum_{u\in\auG}\epsilon(u)u$. As $e_0$ is a sum of central primitive
idempotents, $\epsilon(u)=\sum_\chi\chi(1)\aubar{\chi(u)}$ over a set of irreducible characters;
so $\epsilon(u)$ is an algebraic integer, and every conjugate of it is the value of a character
of degree at most $\#\auG$. By Proposition~\ref{au:equiv},
$(\iota_\vartheta z')(g)=\sum_u\epsilon(u)(\iota_\vartheta z)(g\beta_\vartheta(us))$. (a) The gauge is
ultrametric, homogeneous, and unchanged by right translation by elements of $K^\circ$. (b) Each
value of $\sigma(\iota_1z')$ has norm at most $(\#\auG)^2e^{B/2}$, and $\|\cdot\|_\sigma^2$ is a
sum over at most $\exp(x^{30})$ classes; $\log\#\auG\le2f\,r^2\log p_1\le x^{24}$.
\end{proof}

\subsection{The inputs of Theorem B}\label{ss:au-inputsB}

\begin{lemma}[Count of the ramified forms]\label{au:ramcount}
$\dim\auSram\ge\bigl(1-f/(N-1)\bigr)\auDX\ge(1-b^{-8})\auDX$.
\end{lemma}
\begin{proof}
Let $w\mid\pi_N$ and let $J'$ be the group generated by $J_X$ and $z_w$. The forms of $\auS$ fixed
by $z_w$ are the forms of level $J'$. The group $J'$ consists of elements of $K^\circ$ that are
$\equiv1$ at the chosen places, so it is free by Lemma~\ref{au:free}, and $[J':J_X]=N-1$. For a
free level the dimension of the space of forms is $\dim W_\lambda$ times the number of classes,
and the number of classes of $J_X$ is $[J':J_X]$ times that of $J'$. So
$\dim\ker(1-z_w)=\auDX/(N-1)$. The operators $1-z_w$ commute, and the kernel of a product has
dimension at most the sum of the dimensions of the kernels; there are $f$ primes $w$. Since
$\auSram$ is the image of $\prod_w(1-z_w)$, the first inequality follows; the second holds because
$f\le b$ and $N\ge b^9+1$.
\end{proof}

In the next two lemmas all spaces are finite-dimensional complex inner product spaces, $\|A\|$ is
the Hilbert--Schmidt norm of a linear map, and $\mathrm{pr}_U$ is an orthogonal projection.

\begin{lemma}[One translate]\label{au:onetranslate}
Let a finite group $\auG$ act unitarily on $T$, let $T_0\subseteq T$ be a stable subspace with
projection $\mathrm{pr}_0$, let $B\colon S\to T$ be linear and $\varepsilon\ge0$. Suppose
\begin{enumerate}
\item[(H1)] $\|\mathrm{pr}_0B\,\mathrm{pr}_U\|^2\ge\dim U-\varepsilon\dim S$ for every subspace
$U\subseteq S$;
\item[(H2)] $\dim T_0\cdot\sum_{s\in\auG}\|B^*sB\|^2\le\#\auG\,(1+\varepsilon)(\dim S)^2$.
\end{enumerate}
Then for every subspace $F\subseteq T_0$, with $\alpha=\dim F/\dim T_0$, there is $s\in\auG$ with
$\dim(\mathrm{pr}_0sB)^{-1}(F)\le\bigl(\tfrac12(1+\alpha)+\tfrac32\varepsilon\bigr)\dim S$.
\end{lemma}
\begin{proof}
We have $\sum_s\|\mathrm{pr}_FsB\|^2=\langle\mathrm{pr}_F,\Sigma\rangle$ with
$\Sigma=\sum_ssBB^*s^{-1}$ and the Hilbert--Schmidt inner product. Now $\|\mathrm{pr}_F\|^2=\dim
F$, and $\|\Sigma\|^2=\sum_{s,s'}\operatorname{tr}(sBB^*s^{-1}s'BB^*s'^{-1})=\#\auG\sum_u
\operatorname{tr}\bigl((B^*uB)(B^*uB)^*\bigr)=\#\auG\sum_u\|B^*uB\|^2$. By Cauchy--Schwarz and
(H2),
\[
\Bigl(\sum_s\|\mathrm{pr}_FsB\|^2\Bigr)^2\le\dim F\cdot\#\auG\sum_u\|B^*uB\|^2\le\alpha(1+
\varepsilon)(\#\auG\dim S)^2\le\Bigl(\frac{\alpha+1+\varepsilon}2\Bigr)^2(\#\auG\dim S)^2 .
\]
So some $s$ has $\|\mathrm{pr}_FsB\|^2\le\tfrac12(\alpha+1+\varepsilon)\dim S$. Let
$U=(\mathrm{pr}_0sB)^{-1}(F)$. For $x\in U$, $\mathrm{pr}_FsBx=\mathrm{pr}_F\mathrm{pr}_0sBx=
\mathrm{pr}_0sBx=s\,\mathrm{pr}_0Bx$, because $F\subseteq T_0$ and $\mathrm{pr}_0$ commutes with
$s$; hence $\|\mathrm{pr}_0Bx\|=\|\mathrm{pr}_FsBx\|$ and
$\|\mathrm{pr}_0B\,\mathrm{pr}_U\|^2\le\|\mathrm{pr}_FsB\|^2$. With (H1),
$\dim U-\varepsilon\dim S\le\tfrac12(\alpha+1+\varepsilon)\dim S$.
\end{proof}

\begin{lemma}[Few translates]\label{au:few}
Let $B_s\colon V\to W$, $s\in\Sigma$, be linear maps with $V\ne0$, $\dim W>0$, and
$0<\varepsilon<\tfrac1{16}$. Suppose the $B_s$ \emph{spread with error $\tfrac32\varepsilon$}: for
every subspace $F\subseteq W$, with $\alpha=\dim F/\dim W$, some $s$ has
$\dim B_s^{-1}(F)\le(\tfrac12(1+\alpha)+\tfrac32\varepsilon)\dim V$. Then there is a finite
$\mathsf t\subseteq\Sigma$ with
\[
\#\mathsf t\cdot\dim V\le\tfrac52\bigl(1+\log\tfrac1{16\varepsilon}\bigr)\dim W,\qquad
\dim\sum_{s\in\mathsf t}B_s(V)\ge(1-16\varepsilon)\dim W .
\]
Moreover, for a subspace $U\subseteq V$,
$\dim\sum_{s\in\mathsf t}B_s(V)\le\dim\sum_{s\in\mathsf t}B_s(U)+\#\mathsf t(\dim V-\dim U)$.
\end{lemma}
\begin{proof}
Let $F\subseteq W$ have $1-\alpha>16\varepsilon$ and let $s$ be as in the hypothesis. Then
$\tfrac12(1+\alpha)+\tfrac32\varepsilon\le1-(\tfrac12-\tfrac3{32})(1-\alpha)\le1-\tfrac25(1-
\alpha)$, and $\dim(F+B_sV)=\dim F+\dim V-\dim B_s^{-1}(F)\ge\dim F+\tfrac25(1-\alpha)\dim V$. So
adding $B_s(V)$ multiplies the codimension by at most $1-\rho\le e^{-\rho}$, where
$\rho=2\dim V/(5\dim W)$; applied to $F=0$ this shows $\rho\le1$. Starting from $F=0$ and
repeating while the codimension exceeds $16\varepsilon\dim W$, after
$l=\lceil\rho^{-1}\log(1/16\varepsilon)\rceil$ steps at most the codimension is at most
$16\varepsilon\dim W$. And $l\rho\le1+\log(1/16\varepsilon)$, so
$\#\mathsf t\dim V\le l\dim V=\tfrac52l\rho\dim W$. For the last claim choose a complement $C$ of
$U$: $\sum_sB_s(V)=\sum_sB_s(U)+\sum_sB_s(C)$.
\end{proof}

\begin{proposition}[Second moment of the Poincar\'e map]\label{au:moment}
Let $\deg$ be the product over $\tau\in\auPl$ of the formal degrees of the discrete series
$\varsigma_\tau$, and $\mathrm{mass}_r$ the mass of the genus of the standard hermitian lattice of
rank $r$. There are an inner product on $\auS$ and a scalar $c_P>0$ such that $P=c_P\auP$
satisfies:
\begin{enumerate}
\item[(i)] $0<\audT\le\#\auG\cdot\deg\cdot\mathrm{mass}_r$;
\item[(ii)] $\|\varrho(e_0)P\,\mathrm{pr}_U\|^2\ge\dim U-5b^{-4}\auDX$ for every subspace
$U\subseteq\auS$;
\item[(iii)] $\sum_{s\in\auG}\|P^*\varrho(s)P\|^2\le(1+9b^{-4})\,\auDX^2/(\deg\cdot
\mathrm{mass}_r)$.
\end{enumerate}
\end{proposition}
\begin{proof}
The proof is given in \S\ref{pb:sec} (Theorem~\ref{pb:thmB}, Propositions~\ref{pb:i}, \ref{pb:ii},
\ref{pb:iii}), from the four inequalities of Theorem~\ref{pb:est}, which are proved in
\S\S\ref{pb:sec}--\ref{pbm:sec} and \S\S\ref{pds:sec}--\ref{pfg:sec}, and from
Proposition~\ref{au:mult}; the value of the covolume $\nu$ (Lemma~\ref{pb:covol})
is Theorem~\ref{pbv:covol}. Those sections work with the lattice $L$ of $Y_0$ of
\S\ref{pb:ss-setting}; the passage to the lattice $\mathsf M_0$ of \eqref{ps:eq-lattice}, with
which $\auAz$, $\auG$ and $\auP$ are defined here, is Lemma~\ref{pb:lattice} (the two Poincar\'e
maps correspond under a unitary isomorphism of the targets commuting with $\auG$ and a unitary
automorphism of $\auS$, and the same $g_0$, $\mathsf t$, $\psi_U$ serve for both). The dependency tree, with the published results used, is in
\S\ref{pb:ss-remains}.
\end{proof}

\begin{remark}[The two lattices of $Y_0$]\label{au:two-lattices}
Theorems~\ref{thm:srcA} and~\ref{thm:srcC} are proved with the lattice $\mathsf M_0$ of $Y_0$, and the
sections on Theorem~\ref{thm:srcB} with the lattice $L$. These are not the same lattice, and the
difference is not a matter of notation. In the basis of \S\ref{pb:ss-setting}, $\mathsf M_0$ is
obtained from $L$ by multiplying the vectors of the string (the vectors $e_a$, $e_a^*$ and $z$) by a
finite idele $u_0$ of $E$ with $u_0\aubar{u_0}=1$ which is not a unit: it is non-trivial at the chosen
places and at the places above $q_1$ at which $\alpha/\aubar\alpha$ is not a unit. On the side of
Theorem~\ref{thm:srcA} the idele $u_0$ cannot be removed: the relations
$T_j(\mathsf M_j^+)=\mathfrak a\,\mathsf M_{j+1}^-$ of the chain and the unimodularity of the lattices
determine one such idele for each string (\S\ref{ps:ss-Y}). The stabilisers of $L$ and of
$\mathsf M_0$ are different open compact subgroups of $\auU(Y_0)(\auA_{E^+,f})$, and nothing in the
data makes the ideal of $u_0$ principal with a generator of norm one, so the two are not exchanged by
an element of $\auU(Y_0)(E^+)$ in general. A reader comparing \S\ref{ps:ss-Y} with
\S\ref{pb:ss-setting} will therefore find two different spaces of forms on $\auU(Y_0)$, and the three
theorems are combined in \S\ref{ss:au-ABC} for a single space. We explain why this does not affect
the argument.

The multiplication by $u_0$ on the string is the product of two commuting elements: the scalar
$u_0$ on all of $Y_0$, and the element $\mathsf h_{u_0}$ which acts by $u_0^{-1}$ on $X_0$ and trivially
on the string. A scalar of norm one lies in the centre of $\auU(Y_0)(\auA_{E^+,f})$ and does not change
the stabiliser of any lattice. The element $\mathsf h_{u_0}$ lies in the centre of
$\auU(X_0)(\auA_{E^+,f})$, hence in $R_0(\auA)$, the group over which the Poincar\'e map integrates, and
it commutes with the real point $g_0$. So $\mathsf M_0=u_0\,\mathsf h_{u_0}L$, the two stabilisers are
conjugate by $\mathsf h_{u_0}$, right translation by $\mathsf h_{u_0}$ is a unitary isomorphism between
the two spaces of forms which commutes with the action of $\auG$, and the substitution
$r\mapsto r\,\mathsf h_{u_0}^{-1}$ in the integral over $[R_0]$ shows that the two Poincar\'e maps differ
by this isomorphism on the target and by a unitary automorphism of the source $\auS$ (translation by
the central element $\mathsf h_{u_0}$); this is Lemma~\ref{pb:lattice}(a),(b). The data $g_0$,
$\mathsf t$, $\psi_U$, the Haar measures and the constant $c_P$ are the same for both.

Theorem~\ref{thm:srcB} enters the proof of the main inequality only through
Proposition~\ref{au:moment}, and the three assertions of that proposition are a dimension, the
Hilbert--Schmidt norms of the compressions $\varrho(e_0)P\,\mathrm{pr}_{\mathsf U}$ over subspaces
$\mathsf U\subseteq\auS$ of given dimension, and the Hilbert--Schmidt norms of $P^*\varrho(s)P$ for
$s\in\auG$. Each of them is unchanged when $P$ is composed on the left with a unitary isomorphism
commuting with $\auG$ and on the right with a unitary automorphism of $\auS$
(Lemma~\ref{pb:lattice}(c)). Hence every statement proved for $L$ in
\S\S\ref{pb:sec}--\ref{pfg:sec} holds verbatim, with the same constants, for $\mathsf M_0$; no step
of those sections has to be repeated, and no estimate changes. In particular the space $\auT$, the
group $\auG$ and the map $\auP$ of Theorems~\ref{thm:srcA}, \ref{thm:srcB} and~\ref{thm:srcC} are the same
objects, formed with $\mathsf M_0$, and the deduction of the main inequality in \S\ref{ss:au-ABC} is
not affected.
\end{remark}

\noindent\emph{Outline.} (iii) is a Petersson formula: $\langle\auP\varphi,\varrho(s)\auP\varphi'
\rangle$ is the Bessel period of a Poincar\'e series, which unfolds into a sum over the double
cosets $R_0(E^+)\backslash G_0(E^+)/R_0(E^+)$ of products of local orbital integrals. The
identity coset gives the main term $\auDX^2/(\deg\cdot\mathrm{mass}_r)$ after summation over $s$.
The remaining cosets are bounded by the decay of the Bessel functions of discrete series whose
parameter has a large gap (here the gap $\Delta=2r^8$ and the size of the weights are used;
\S\S\ref{pbe:sec}, \ref{pds:sec}, \ref{pbp:sec}), by a count of the cosets
that meet the support at the finite places (\S\S\ref{ppi:sec}--\ref{pbo:sec}), and by bounds for sums of
Kloosterman type at the level primes (\S\S\ref{pbm:sec}, \ref{pfr:sec}--\ref{pfg:sec}); the total is at most
$9b^{-4}$ times the main term, uniformly in the rank $r\approx x^5$. (ii) follows from the same
formula for the diagonal terms together with the fact that $1-e_0$ has a small coefficient at the
identity, at most $b^{-8}/2$ for each level prime: the labels that are not good are a small
proportion for the Plancherel measure of $\auGL_r(\auF_{p_i})$~\cite{Green1955}. (i) is a count on
the definite inner form $\auU(W_m)$, where the level is free and the dimension is given by the
mass formula, transferred to $Y_0$ by Proposition~\ref{au:mult}. No Petersson or relative trace
formula of this uniformity is in the literature, and we regard (iii) as a new theorem.

\begin{proposition}[Spreading]\label{au:spread}
$\audT>0$, and for every subspace $F\subseteq\auT$, with $\alpha=\dim F/\audT$, there is
$s\in\auG$ with $\dim\auP_s^{-1}(F)\le\bigl(\tfrac12(1+\alpha)+\tfrac32(8b^2)^{-1}\bigr)\auDX$.
\end{proposition}
\begin{proof}
Apply Lemma~\ref{au:onetranslate} to $T=\auAz$, $T_0=\auT$, $\mathrm{pr}_0=\varrho(e_0)$ (an
orthogonal projection commuting with $\auG$, as $e_0$ is a central idempotent and $\varrho$ is
unitary), $B=P$ and $\varepsilon=(8b^2)^{-1}$. (H1) is (ii), as $5b^{-4}\le\varepsilon$. (H2)
follows from (i) and (iii), as $9b^{-4}\le\varepsilon$. The maps $\mathrm{pr}_0\varrho(s)P$ and
$\auP_s$ differ by a non-zero scalar and have the same preimages.
\end{proof}

\subsection{Theorem C}\label{ss:au-proofC}

\subsubsection*{The inputs}

\begin{proposition}[Injectivity of the lower chain]\label{au:inj}
For every $m$-fold twist $\vartheta$ the map $\iota_\vartheta\colon\auAz\to\auAm(\vartheta)$ is
injective.
\end{proposition}
\begin{proof}
The proof is given in \S\ref{pc:ss-triangle}, from Lemma~\ref{pc:arch} (archimedean types; proved as
Proposition~\ref{ptv:arch}), Theorem~\ref{pc:tri}(a),(b) (the squares of the triangle commute up to
constants; proved in \S\S\ref{ptr:sec}--\ref{ptw:sec} from the exchange identity
Theorem~\ref{pc:X}, which is Theorem~\ref{pcx:X}, proved in \S\S\ref{pcx:sec}--\ref{pcn:sec}) and
Weil's Siegel formula (SW) of Lemma~\ref{pc:SW}. The published results used are listed in
\S\ref{pct:ss-status} and \S\ref{ptq:ss-status}.

\emph{Outline.} Only the first lift of each row of a triangle of lifts is in Weil's range; it is
injective because the series of Rallis' formula is close to the identity (Lemma~\ref{pc:series},
Proposition~\ref{pc:rallis}); the injectivity of the chain follows from the commutation relations
alone (Proposition~\ref{pc:tri-inj}), whatever the constants of the squares. No $L$-value is used.
\end{proof}

\begin{proposition}[Equality of multiplicities]\label{au:mult}
For $1\le j\le m$ and every element $h$ of the spherical Hecke algebra away from $M$ (the
compactly supported functions on the finite adelic group outside the chosen places that are
bi-invariant under the stabiliser of the lattice, at all places, split or not),
$\operatorname{tr}(h\,|\,e_0\mathcal A(W_j))=\operatorname{tr}(h\,|\,\auT)$. In particular
$\dim e_0\mathcal A(W_j)=\audT$, and each system of eigenvalues of the Hecke operators at the split
places prime to $M$ has the same multiplicity in $e_0\mathcal A(W_j)$ and in $\auT$. (The
hypothesis of Proposition~\ref{pm:mult} holds for the data: at every $\tau$, two entries of
$\varsigma_\tau$ on coordinates of different signs \eqref{pm:eq-signs} differ by at least
$\Delta=2r^8\ge500r^3+10r^5$.)
\end{proposition}
\begin{proof}
The proof is given in \S\ref{pm:sec} (Proposition~\ref{pm:mult}), from Theorem~\ref{pm:arch}, two
statements on the coefficients of discrete series of $\auU(p,q)$ (decay, and Harish-Chandra's
formula for elliptic orbital integrals), proved in Remark~\ref{pm:rem-other}(1) and
\S\ref{pbr:sec} from Schur orthogonality, admissibility, Schmid's theorem on $K$-types and the
theory of compact Lie groups, and from the fact
that irreducible labels name cuspidal characters (\S\ref{pm:ss-mult}, (G)). The proof of
Proposition~\ref{pm:mult} is written for operators $T=\varphi_S\otimes1$; it applies without change
to every compactly supported bi-invariant $\varphi_S$ at any finite set $S$ of finite places away
from the chosen ones, split or not, because Theorem~\ref{pm:main} allows an arbitrary function
away from the place $v_0$; this gives the statement for the full spherical algebra (Remark~\ref{pm:rem-full}). No
stabilised comparison is used: the test function at infinity is the integrable kernel of the
discrete series, the terms of the reducible characteristic polynomials
vanish, and for irreducible ones the comparison is class field theory for a norm one torus.
\end{proof}

A \emph{critical pair} is a pair $(\nu,z)$ of an integer $1\le z\le t$ and an automorphic
character $\nu$ of $\auU(1)$ of weight $-c$ with $z\le c(\tau)\le4b$ at every infinite place and
tamely ramified at chosen places only; it is \emph{ramified} if $\nu$ is ramified at every chosen
place. Let $L(s,\nu^{-1})$ be the Hecke $L$-function of the character $x\mapsto\nu^{-1}(x/\aubar
x)$ of $E$, without the Euler factors at the level. A \emph{ramified ratio} is a function of the
twist of the form
\[
\vartheta\longmapsto\prod_{i=1}^{t'}\frac{|L(z_i,(\nu_i\vartheta^{\pm n_i})^{-1})|}{|L(z_i,
\nu_i^{-1})|},\qquad t'\le t,\ 1\le n_i\le m,\ (\nu_i,z_i)\text{ ramified critical pairs.}
\]

\begin{proposition}[Factorisation of the chain determinant]\label{au:factor}
For $1\le j\le m$ and a twist $\vartheta$ let $Z_j(\vartheta)=\det(u^*u)>0$, where
$u=\Theta_{j,\vartheta}$ is the theta lift from $e_0\mathcal A(W_j)$ to algebraic modular forms on
$\auU(V_j)$ with the character $\mu_j^0\vartheta$, normalised as in \S\ref{pc:ss-rallis} (joint
harmonic at infinity; at the chosen places the indicator of the inclusion; at $w\in Q_\vartheta$
the function $x\mapsto\psi_w(\det x^{\mathrm{top}})$; values in the algebraic modular forms of level
$J^V_\vartheta$ with their inner product). There are $e\in\Z$ and the ramified ratios $\rho_{ij}$,
$0<i<j\le m$, of \eqref{pc:eq-rho}, such that
for every $m$-fold twist $\vartheta$
\[
\frac{Z(1)}{Z(\vartheta)}=\Bigl((\auNm Q_\vartheta)^{e/2}\prod_{0<i<j\le m}\rho_{ij}(\vartheta)
\Bigr)^{-\audT}\cdot\prod_{j=1}^m\frac{Z_j(1)}{Z_j(\vartheta^j)} .
\]
\end{proposition}
\begin{proof}
The proof is given in \S\ref{pc:ss-triangle}, from Lemma~\ref{pc:arch}, Theorem~\ref{pc:tri}, Weil's Siegel formula and Proposition~\ref{au:mult} (in the form $(\mathrm M^+)$ of Lemma~\ref{pc:detcompare}); Lemma~\ref{pc:arch} is Proposition~\ref{ptv:arch}, Theorem~\ref{pc:tri} is proved in \S\S\ref{ptr:sec}--\ref{ptw:sec} (part (c) through (RS), Theorem~\ref{ptw:RS}), and Theorem~\ref{pc:X} is Theorem~\ref{pcx:X}. The published results used are listed in the table of \S\ref{ss:au-remains}.

\emph{Outline.} On $\auT$
the operator $\iota_\vartheta^*\iota_\vartheta$ is a scalar times a product of $m$ commuting
operators, the $j$-th being the operator of the Rallis inner product formula for the single lift
$W_j\to V_j$. This is proved by embedding the chain in a triangle of spaces joined by lifts and
transposed lifts and commuting them square by square ($m^2$ squares); each square is an exchange
identity of Siegel--Weil type between theta integrals over two unitary groups of complementary
rank, valid up to a constant whose modulus is a ratio of Hecke $L$-values and a power of
$\auNm Q$. Proposition~\ref{au:mult} is used to identify the good forms on the $W_j$ with $\auT$
as Hecke modules.
\end{proof}

\begin{proposition}[Auxiliary forms]\label{au:aux}
There are $e'\in\Z$, ramified ratios $\rho'_{ij}$ ($0\le i,j<m$) and
$\gamma\colon\{m\text{-fold twists}\}\to\C^\times$ with $\gamma(1)=1$ and
$|\gamma(\vartheta)|=(\auNm Q_\vartheta)^{e'/2}\prod_{i,j}\rho'_{ij}(\vartheta)$, such that $\auT$
is spanned by forms $y$ with: $\gamma(\vartheta)\Theta_\vartheta y$ is integral at every $\iota$
with $\ell(\iota)\nmid\auNm Q_\vartheta$, for every $m$-fold twist $\vartheta$; and
$\|\Theta_1y\|_\sigma^2\le\exp(2x_1^2b^2)$ for all $\sigma\in\auAut(\C)$.
\end{proposition}
\begin{proof}
The proof is given in \S\ref{pc:ss-aux}, from Theorem~\ref{pc:auxcore} and Proposition~\ref{au:good}; Theorem~\ref{pc:auxcore} is proved in \S\S\ref{pgr:sec}--\ref{pgm:sec} (dependencies in \S\ref{pgr:ss-status}); the ratios are the $\rho'_{ij}$ of \eqref{pc:eq-rhoprime}. The published results used are listed in the table of \S\ref{ss:au-remains}.

\emph{Outline.} Forms on $Y_0$ are obtained by transposed lifts from short integral forms on a second,
untwisted chain of definite spaces; moving a path through the two grids by $m^2$ turns expresses
their images under $\iota_\vartheta$ through integral forms, up to the scalar $\gamma(\vartheta)$;
the deduction for the good part is as in Proposition~\ref{au:good}.
\end{proof}

\begin{definition}[Regular families of twists]\label{def:au-regular}
Put $\aubud=\aueu/(mk)$. A family $\theta\colon\auI\to\{m\text{-fold twists}\}$ is \emph{regular}
if:
\begin{enumerate}
\item[(R1)] $\ell(\iota)\nmid\auNm Q_{\theta_\iota}$ for all $\iota$;
\item[(R2)] for every ramified critical pair $(\nu,z)$, $1\le n\le m$ and sign $\pm$, the family
$\iota\mapsto|L(z,\nu^{-1})|/|L(z,(\nu\theta_\iota^{\pm n})^{-1})|$ has height at most $x_1$;
\item[(R3)] for $1\le j\le m$ there are $\beta_\ell\ge0$ with $\sum_{\ell\in T}\beta_\ell\le
\aubud$ for all finite $T$ such that, for every $\iota$ with prime $\ell$: (i) the twist
$\theta_\iota^{\,j}$ \emph{has denominators at most $\beta_\ell$ at $\iota$ for the $j$-th lift}
in the sense of Definition~\ref{po:denom} (the false lifts being those of
Definition~\ref{po:false}: the forms in the image of $e_0$ that are orthogonal to the lifts and
are killed by the Hecke polynomial $P_w(X_w)$ of Lemma~\ref{po:facts}(L3) at one split place $w$
off the level; a form is orthogonal to the false lifts if it is so for every such $w$); and (ii)
for every main pair $p$ for $j$ (Definition~\ref{po:main}) and all test forms $\Xi,\Xi'$ at
$\ell$ for $(p,\theta_\iota^{\,j})$ (Definition~\ref{po:R3}), the inner product of the orthogonal
projections of $\Xi$ and $\Xi'$ to the line of the normalised lift
$\varphi_{p,\theta_\iota^{\,j}}$ has absolute value at most $e^{\beta_\ell}$ at $\iota$.
\end{enumerate}
\end{definition}

Condition (R3) says that the congruence module between lifts and non-lifts is small at $\iota$
after the twist. By (R2) and the rules for heights, for a regular family the inverse of
$(\auNm Q_{\theta_\iota})^{e/2}$ times a product of $n$ ramified ratios has height at most
$n\,t\,x_1$.

\begin{proposition}[Regular families exist]\label{au:regular}
For all admissible data there is a regular family of $m$-fold twists.
\end{proposition}
\begin{proof}
The proof is given in \S\ref{pg:sec} (Theorem~\ref{pg:main} and Lemma~\ref{pg:sums}); summary in
\S\ref{pg:ss-summary}. The sums over $\ell$ are Lemma~\ref{pg:sums}; the
reduction is to Theorems~\ref{pg:I}, \ref{pg:II}, \ref{pg:III}. Theorem~\ref{pg:I} is proved in
\S\ref{pga:sec} (with Proposition~\ref{pax:sign}) from published theorems on Galois representations
attached to automorphic forms on unitary groups; Theorem~\ref{pg:II} in \S\ref{pse:sec} from
statement {\rm(IM)}, which is Propositions~\ref{pex:IM} and~\ref{pex:IM1}; Theorem~\ref{pg:III} in
\S\S\ref{pri:sec}--\ref{prm:sec} (Propositions~\ref{pri:odd}, \ref{pri:kummer}, \ref{pri:dict} included),
with the local loss Proposition~\ref{pg:loss} proved in \S\ref{pll:sec} and its explicit exponents,
Theorem~\ref{pll:exponent}, in \S\ref{pki:ss-deduce}. The proof also uses Proposition~\ref{au:mult}.
The published results used are listed in the table of \S\ref{ss:au-remains}.
\end{proof}

\noindent\emph{Outline.} Fix $\iota$ with prime $\ell$. A non-trivial denominator in the
projection to the lifts gives a congruence modulo a power of $\ell$ between a lift and a form that
is not a lift, hence, by the method of Ribet~\cite{Ribet1976} applied to the Galois
representations attached to the eigenforms on $\auU(V_j)$, a non-trivial class in a Selmer group
of a character of $\operatorname{Gal}(\overline E/E)$ built from $\mu_j^0\vartheta^j$ and powers
of the cyclotomic character, with crystalline conditions above $\ell$. One then shows that there
is an anticyclotomic twist $\vartheta$ of $\ell$-power order, tamely ramified at auxiliary primes
of norm $\equiv1$ modulo a large power of $\ell$, for which these Selmer groups, for the first $m$ powers of
$\vartheta$ simultaneously, have length at most $(4f+3)\,r\,v_\ell(N_{\mathrm{ch}})+
2^{-\ell}$, where $N_{\mathrm{ch}}$ is the product of the orders of $\auGL_{rt}$ over the residue
fields of the chosen places (Theorem~\ref{pg:II}); this bound does not involve $p$-adic Hodge
theory beyond the dimension formula of Bloch and Kato. The classes produced by Ribet's method
modulo a power of $\ell$ need not satisfy the crystalline condition; their distance from it, the
\emph{local loss}, is bounded by integral $p$-adic Hodge theory, and here the
sparseness of the set $\auPb$ is used (Proposition~\ref{pg:loss}); summed over $\ell$ the losses
are within $\aubud$ (Lemma~\ref{pg:sums}). The extensions between the powers of the cyclotomic
character that Ribet's method also produces are controlled, without any bound for class groups of
real cyclotomic fields, by auxiliary places of Taylor--Wiles type (Proposition~\ref{pg:triangle}). The same twist must satisfy (R2); this is an upper bound for the
$\ell$-adic valuation of the ratio at one twist of $\ell$-power order near the trivial one
(Proposition~\ref{pse:stick}), not a non-vanishing theorem modulo $\ell$ for a family of twists. The background required is: Galois
representations attached to automorphic forms on unitary groups with local--global compatibility
at the places used; integral comparison theorems above $\ell$; bounds for class groups of
cyclotomic type. The published results of this kind (for instance Hida~\cite{HidaNonvanishing},
Finis~\cite{Finis}, Prasanna~\cite{Prasanna}) assume $\ell$ odd, unramified and split in the CM
extension; here every $\ell$ is needed, including $\ell=2$, the primes dividing the levels, and
the primes inert in $K$, which carry the prime sum; the proof in \S\ref{pg:sec} therefore
does not rest on those results.

\begin{proposition}[One lift and Hecke values]\label{au:onelift}
Let $1\le j\le m$ and $d_j=\dim e_0\mathcal A(W_j)$. There are a constant $C_j\ne0$, for every twist
$\vartheta$ a number $\Delta_j(\vartheta)\ne0$ (the determinant of a \emph{comparison operator}
compressed to the space of lifts) and the positive number $H_j(\vartheta)$ of \eqref{po:eq-H}: for the three critical pairs $(\nu_i,z_i)$ of \S\ref{po:ss-hecke} it is the product of $(\auNm\mathsf S_\vartheta)^{z_i^2}\prod_{1\le z\le z_i}|L^{\mathsf S_\vartheta}(z,(\nu_i\vartheta^{\pm1})^{-1})|^2$, such that
\[
Z_j(\vartheta)\cdot\Delta_j(\vartheta)=C_j\cdot\bigl(H_j(\vartheta)/H_j(1)\bigr)^{d_j},
\]
and
\begin{enumerate}
\item[(a)] $|\sigma\Delta_j(1)|\le\exp(3d_j\cdot f\,m\,b\,(\log b)^2)$ for all
$\sigma\in\auAut(\C)$;
\item[(b)] if $\ell(\iota)\nmid\auNm Q_\vartheta$ and $\vartheta$ has denominators at most $\beta$
at $\iota$ for the $j$-th lift, then $|\Delta_j(\vartheta)|_\iota\le\exp(d_j\,r\,\beta)$.
\end{enumerate}
\end{proposition}
\begin{proof}
The proof is given in \S\ref{po:sec}. Its inputs (L4), Lemmas~\ref{po:ff}, \ref{po:classes} and Propositions~\ref{po:exchange}, \ref{po:span}, \ref{po:periods} are proved in \S\S\ref{plf:sec}, \ref{pte:sec}--\ref{ptf:sec}, \ref{pfq:sec}, \ref{psp:sec}--\ref{psv:sec}, \ref{ppe:sec}: Lemma~\ref{po:ff} is Theorem~\ref{pfq:ff}, for the first factors of Definition~\ref{pfq:def-ff}, for which Proposition~\ref{po:exchange} holds by Proposition~\ref{pfq:exchange}; the characters $\chi_M$, $\mu_j^0$ and their tame parts $\omega^j$ at the chosen places (used in \S\ref{po:ss-hecke} and in Lemma~\ref{ppe:three}) are Definition~\ref{pfq:def-chars} and Proposition~\ref{pfq:tame}, with existence by Lemma~\ref{pz:mu}; the positivity of $H_j(\vartheta)$ is Lemma~\ref{po:Hpos}. The published results used are listed in the table of \S\ref{ss:au-remains}.

\emph{Outline.} The comparison operator is $U_\vartheta=c\cdot\mathrm{glue}_\vartheta
\mathrm{glue}_\vartheta^*$, where $\mathrm{glue}_\vartheta$ maps tuples of forms on a smaller
definite space to forms on $V_j$ by multiplication with theta series; it takes lifts to forms
orthogonal to the false lifts; a tight frame of integral tuples writes it as a sum of operators of
rank one, so it preserves integrality away from $Q_\vartheta$ and $U_1$ is short; and
$Z_j(\vartheta)\det(\mathrm{pr}\,U_\vartheta)$ depends on $\vartheta$ only through $L$-values, by
three exchange identities.
\end{proof}

\begin{proposition}[Heights of twisted Hecke values]\label{au:hecke}
For a regular family $\theta$ and $1\le j\le m$ the family $\iota\mapsto H_j(1)/H_j(\theta_\iota^
{\,j})$ has height at most $10\,m\,x_1+k\,(6\,\aubud-\auhdir)$.
\end{proposition}
\begin{proof}
The proof is given in \S\ref{po:ss-hecke}, from Propositions~\ref{po:rallis} (= Proposition~\ref{pex:rallis}), \ref{po:arch} (= Theorem~\ref{psu:arch}, with (CL) = Theorem~\ref{pcj:CL}) and \ref{po:small} (= Theorem~\ref{psu:small}, with the constant $3x^{35}$). It uses the characters $\mu_j^0$ of \S\ref{po:ss-hecke}, whose tame parts are given by Proposition~\ref{pfq:tame}. The published results used are listed in the table of \S\ref{ss:au-remains}.

\emph{Outline.}
$H_j$ is a product of the values at the main pairs $(\nu,z)$, $m<z\le r$, and of the $5m$ remaining
ramified pairs (Lemma~\ref{po:Hratio}). A ramified pair costs $x_1$ by (R2). The value at a main pair is the squared norm
of the theta lift of $\nu\vartheta$ from $\auU(1)$ to $\auU(2z)$, whose values are jets of order
$\kappa(\sigma)$, up to slack, of theta functions with complex multiplication. Its archimedean
size is $\sum_\sigma\ausinf(\kappa(\sigma))$ up to the error; its $\ell$-adic valuation is at
least $\sum_\sigma s_\ell(\kappa(\sigma),\sigma)$ by the count, and for the twist $\theta_\iota$
it is at most this plus a multiple of $\beta_\ell$, by (R3) and by the sharpness of the count (the valuation
predicted by Proposition~\ref{au:count} is attained by some jet). Summing over $\ell$ with
Lemma~\ref{au:gains}(a) gives $-\auhdir$ for each main pair; there are $k/2$ main pairs, each entering squared (Lemma~\ref{po:Hratio}), and a main pair costs $6\aubud-\auhdir$ as a whole (Proposition~\ref{po:mainheight}).
\end{proof}

\subsubsection*{The deductions}

\begin{proposition}[The cost of one lift]\label{au:Zj}
For a regular family $\theta$ and $1\le j\le m$ the family $\iota\mapsto Z_j(1)/Z_j(\theta_\iota^
{\,j})$ has height at most $\audT\bigl(f\,a\,(\log b)^4/m-k\,\auhdir\bigr)$.
\end{proposition}
\begin{proof}
By Proposition~\ref{au:mult}, $d_j=\audT$. By Proposition~\ref{au:onelift},
\[
\frac{Z_j(1)}{Z_j(\theta_\iota^{\,j})}=\frac{\Delta_j(\theta_\iota^{\,j})}{\Delta_j(1)}\cdot
\Bigl(\frac{H_j(1)}{H_j(\theta_\iota^{\,j})}\Bigr)^{\audT} .
\]
With the multiplier $\Delta_j(1)$, (a) and (b) of that proposition and (R3) show that the first
factor has height at most $\audT(3A+r\,\aubud)$, $A=f\,m\,b(\log b)^2$. With
Proposition~\ref{au:hecke} the product has height at most
$\audT\bigl(3A+r\aubud+10mx_1+6k\aubud-k\auhdir\bigr)$. Now $r\le2k$ and
$k\,\aubud\log b=fa(\log b)^4/m$, so it suffices that $3A+10mx_1\le k\,\aubud(\log b-8)$, that
is $3fm^2b(\log b)^2+10m^2x_1\le fa(\log b)^3(\log b-8)$; this holds because $m^2\le4x$ and
$a=xb$.
\end{proof}

\begin{proof}[Proof of Theorem~\ref{thm:srcC} from Propositions~\ref{au:inj}, \ref{au:aux},
\ref{au:factor}, \ref{au:regular} and~\ref{au:Zj}]
Let $\theta$ be a regular family (Proposition~\ref{au:regular}), and let $\gamma$ be as in
Proposition~\ref{au:aux}. Complete $y_1,\dots,y_n$ to a basis of $\auT$ by auxiliary forms
$z_1,\dots,z_{n'}$ from the spanning set of that proposition; $n'=\audT-n\le4\audT/b^2$. We apply
Theorem~\ref{au:twistprime} to the lattice $\auLat_D$, whose maps are injective by
Proposition~\ref{au:inj}, with the twists $\theta_\iota$, allowed by (R1). Hypotheses (ii) and
(iii) hold with $b_{\mathrm{short}}=2x_1^2b^2$. The family $|\gamma(\theta_\iota)|^{-1}$ is the
inverse of a power of $\auNm Q_{\theta_\iota}$ times $m^2$ ramified ratios, so
$|\gamma(\theta_\iota)|^{-2}$ has height at most $2m^2tx_1\le b=b_{\mathrm{den}}$. By
Proposition~\ref{au:factor}, $Z(1)/Z(\theta_\iota)$ is the $\audT$-th power of the inverse of a
power of $\auNm Q$ times fewer than $m^2$ ramified ratios, of height at most $m^2tx_1\le b$,
times the product over $j$ of the families of Proposition~\ref{au:Zj}; so, since
$m\cdot k\,\auhdir=\auhmain$,
\[
b_{\mathrm{vol}}=\audT\bigl(b+f\,a\,(\log b)^4-\auhmain\bigr).
\]
The theorem gives $0\le c+2\delta+n'(2x_1^2b^2+b)+\audT(b+fa(\log b)^4-\auhmain)$, and
$n'(2x_1^2b^2+b)\le4\audT(2x_1^2+b^{-1})\le\audT b$. Hence
$\audT(\auhmain-fa(\log b)^4-2b)\le c+2\delta$, and $2b\le2fa(\log b)^4$.
\end{proof}


\subsection{Where the inputs are proved}\label{ss:au-remains}

The table gives, for each input of the three theorems, the place of its proof and the
published results on which it rests (the precise forms and citations are given in the sections
named in the last column). The statements of this section proved where they stand
(Theorem~\ref{au:twistprime}, Lemmas~\ref{au:gains}, \ref{au:errsum}, the deductions of
the three theorems, Propositions~\ref{au:short}, \ref{au:size}, \ref{au:good}, \ref{au:spread},
\ref{au:Zj} and the lemmas of linear algebra) are not listed.

\begin{center}
\small
\begin{tabular}{llp{4.6cm}}
\hline
Statement&Used in&Proof; published results\\
\hline
Prop.~\ref{au:seesaw}&A&\S\ref{ps:ss-seesaw}, \S\ref{pk:sec}, Prop.~\ref{pf:flip}, Prop.~\ref{pz:types}; Weil representation (W0)--(W5), Witt; (A1)--(A3) in \S\S\ref{par:sec}, \ref{pco:sec}\\
Prop.~\ref{au:divis}&A&\S\ref{pa:ss-divis}, Cor.~\ref{pf:K1}; CM elliptic curves, standard monomial theory\\
\quad Prop.~\ref{au:count}&A, C&\S\ref{pa:ss-count}; standard monomial theory\\
Prop.~\ref{au:kernel}&A&\S\ref{ps:ss-kernel}, Thm.~\ref{pf:recip}; Chowla--Selberg, CM theory\\
Prop.~\ref{au:restrict}&A&\S\ref{ps:ss-restrict}; Schur modules\\
Prop.~\ref{au:expansion}&A&\S\ref{ps:ss-glue}\\
Prop.~\ref{au:equiv}&A, C&\S\ref{ps:ss-Y}; (W1), (W4), (W5)\\
Lemma~\ref{au:shortbasis}&A&\S\ref{ps:ss-short}\\
Prop.~\ref{au:moment}&B&\S\ref{pb:sec}, Thm.~\ref{pbv:covol}; list in \S\ref{pb:ss-remains}\\
Prop.~\ref{au:mult}&B, C&\S\S\ref{pm:sec}, \ref{pbr:sec}; discrete series (Schur orthogonality, Schmid), Green\\
Prop.~\ref{au:inj}&C&\S\ref{pc:ss-triangle}, Prop.~\ref{ptv:arch}, \S\S\ref{ptr:sec}--\ref{ptw:sec}, Thm.~\ref{pcx:X}; Weil's Siegel formula (SW); list in \S\ref{pct:ss-status}\\
Prop.~\ref{au:factor}&C&as Prop.~\ref{au:inj}, with Thm.~\ref{ptw:RS} and $(\mathrm M^+)$\\
Prop.~\ref{au:aux}&C&\S\ref{pc:ss-aux}, Thm.~\ref{pc:auxcore} (\S\ref{pgr:ss-status}); Howe duality, Weyl branching, class numbers\\
Prop.~\ref{au:regular}&C&\S\ref{pg:sec}, \S\S\ref{pga:sec}--\ref{prm:sec}, \S\ref{pll:sec}, Thm.~\ref{pll:exponent}; Galois representations (Labesse, BLGGT, Caraiani), Bloch--Kato, Poitou--Tate, Fontaine\\
Prop.~\ref{au:onelift}&C&\S\ref{po:sec}, \S\ref{pfq:sec}\\
Prop.~\ref{au:hecke}&C&\S\ref{po:ss-hecke}, Prop.~\ref{pex:rallis}, Thms.~\ref{psu:arch}, \ref{psu:small}, \ref{pcj:CL}; Hecke's non-vanishing on $\Rel s=1$, Siegel--Weil\\
\hline
\end{tabular}
\end{center}

\noindent Two points concerning Proposition~\ref{au:onelift} are proved in separate sections: the integrality and size of the first factors for which
the exchange identity is proved (Theorem~\ref{pfq:ff}, Proposition~\ref{pfq:exchange}), and the tame
parts of the characters $\mu_j^0$ at the chosen places (Definition~\ref{pfq:def-chars},
Proposition~\ref{pfq:tame}, Lemma~\ref{pz:mu}).

\section{Proofs of the statements used for Theorems A, B and C}\label{sec:proofsA}

%% ===== proofs/A-divisibility.tex =====
% proofs/A-divisibility.tex
% Proofs of Propositions au:count and au:divis (to be \input after the statement of Theorem A).
% Uses only amsmath/amssymb/amsthm and the theorem environments of the paper.
\providecommand{\auO}{\mathcal O}
\providecommand{\auPl}{\Phi}
\providecommand{\auPa}{\Phi_a}
\providecommand{\auPb}{\Phi_b}
\providecommand{\auS}{\mathcal S}
\providecommand{\auSram}{\mathcal S^{\mathrm{ram}}}
\providecommand{\auDX}{D_X}
\providecommand{\augain}{\operatorname{gain}}
\providecommand{\auerr}{\operatorname{err}}
\providecommand{\auNm}{\operatorname{Nm}}
\providecommand{\auU}{\mathrm U}
\providecommand{\auGL}{\mathrm{GL}}
\providecommand{\auA}{\mathbb A}
\providecommand{\auF}{\mathbb F}
\providecommand{\aucol}{\operatorname{col}}
\providecommand{\aubar}[1]{\overline{#1}}
\providecommand{\auf}[1]{\texttt{#1}}
% macros of this file (prefix \pa)
\providecommand{\pao}{\mathfrak o}
\providecommand{\pam}{\mathfrak m}
\providecommand{\paCp}{\C_p}
\providecommand{\paI}{I_p^{\circ}}
\providecommand{\pavp}{\varpi}
\providecommand{\paded}{\operatorname{ded}}
\providecommand{\paPred}{\operatorname{Pred}}
\providecommand{\paloss}{\operatorname{loss}}
\providecommand{\panot}[1]{\par\smallskip{\raggedright\noindent\textbf{[Not proved here:} #1\textbf{]}\par}\smallskip}

\subsection{Divisibility: proofs of Propositions~\ref{au:count} and~\ref{au:divis}}\label{pa:sec}

This subsection gives a complete proof of Proposition~\ref{au:count} (\S\ref{pa:ss-count}); a
complete proof, from standard facts on elliptic curves with complex multiplication, that products of
theta functions of a CM curve along a resolution of the identity have the normal form
\eqref{eq:au-normal} (\S\S\ref{pa:ss-curve}--\ref{pa:ss-prod}); and the deduction of
Proposition~\ref{au:divis} from these (\S\ref{pa:ss-divis}). The deduction is complete for the
primes that are not exceptional and differ from $N$, given one structural statement on the
theta kernels (Lemma~\ref{pa:kernel}). This lemma is proved in \S\ref{pk:ss-kernel}
(Propositions~\ref{pk:K0}, \ref{pk:std}, \ref{pk:K12}) and in Corollary~\ref{pf:K1} (the integrality of the
jets in (K1) at the primes $\ell\mid2q_1M$). For the exceptional primes
and for $N$ the deduction rests on three further statements
(Lemmas~\ref{pa:hecke}, \ref{pa:eldiv}, \ref{pa:level}), proved in \S\S\ref{ph:sec}, \ref{pe:sec}
and~\ref{pl:sec}; the case of $N$ uses Lemma~\ref{pa:kernel} in the extended form (K1$_N$), which
is Lemma~\ref{pl:kernelN} and is proved in Corollary~\ref{pf:K1}(b).

\subsubsection{Conventions}\label{pa:ss-conv}
Throughout, $p$ is a prime that does not split in $K$; thus $p$ is odd ($2$ splits, as
$d\equiv7\pmod 8$), $p\nmid q_1$, and either $p$ is inert or $p=d$. We fix a ring homomorphism
$\C\to\paCp$, write $\pao$ for the valuation ring of $\paCp$, $\pam$ for its maximal ideal and $v$
for the valuation with $v(p)=1$; numbers and power series with coefficients in $\pao$ are called
\emph{integral}, and congruences are between integral objects. Let $\pi=\pi_p\in\auO_K$ be $p$ if
$p$ is inert and $\sqrt{-d}$ if $p=d$, and $q=q_p=\auNm(\pi)$, so $q=p^2$ or $q=p$, $v(\pi)=1$ or
$\tfrac12$, and $q/\pi^2$ is a unit. Put
\[
\rho=\rho_p=\frac{v(\pi)}{q-1},\qquad\gamma=\gamma_p=\frac1{p-1}-\rho,\qquad
\paI=I_p\cup(\{0\}\cap\{p=d\}),
\]
so that $\paI$ is the set of odd positive integers if $p$ is inert and $\N$ if $p=d$. In both
cases
\begin{equation}\label{pa:eq-const}
(p+1)\rho-1+v(\pi)=\frac1{p-1},\qquad\gamma=v(\pi)\sum_{e\in I_p}p^{-e}.
\end{equation}
We use the following elementary facts. For $n\ge1$, with $s_p(n)$ the sum of the digits of $n$ in
base $p$,
\begin{equation}\label{pa:eq-legendre}
\frac n{p-1}-1-\log_pn\ \le\ v(n!)=\frac{n-s_p(n)}{p-1}\ \le\ \frac{n-1}{p-1}.
\end{equation}
\begin{itemize}
\item[(E1)] If $h$ is an integral power series in any set of variables with $h(0)=0$ and
$v(c)\ge\frac1{p-1}$, then $\exp(ch)$ is integral: $v(c^k/k!)\ge0$ by \eqref{pa:eq-legendre}.
\item[(E2)] If moreover $\beta\in\pao$ has $v(\beta)\ge\tfrac12$, then
$\exp(\beta h)\in1+\beta h\,\pao[[\cdot]]$ and $\log(1+\beta h)\in\beta h\,\pao[[\cdot]]$: for
$k\ge1$, $v(\beta^{k-1}/k!)\ge(k-1)(\tfrac12-\tfrac1{p-1})\ge0$ and
$v(\beta^{k-1}/k)\ge\tfrac{k-1}2-v(k)\ge0$, because $p\ge3$.
\item[(E3)] In a commutative $\pao$-algebra, $(A+B)^{p^e}\equiv A^{p^e}+B^{p^e}\pmod p$.
\end{itemize}

\subsubsection{The count}\label{pa:ss-count}
Let $\auPl$ be a finite set with $f$ elements, $\mathsf C$ a finite set, and
$z=(z_{\tau,i})_{\tau\in\auPl,\,i\in\mathsf C}$ variables. Let $\mathrm{pr}_e\colon\auPl\to\auPl$
($e\in\paI$) be arbitrary maps (\emph{partner maps}); in the application
$\mathrm{pr}_e(\tau)=u(-p^e)\tau$. The \emph{couplings} are
$\Gamma_e=\sum_{\tau,i}z_{\tau,i}^{\,p^e}z_{\mathrm{pr}_e\tau,i}$, homogeneous of degree $p^e+1$.
A series $F\in\paCp[[z]]$ \emph{has the normal form} if, for some $\pavp$ with $v(\pavp)=\rho$,
\begin{equation}\label{pa:eq-nf}
F(\pavp z)=W(z)\exp\bigl(c_\pavp\,\Xi(z)\bigr),\qquad c_\pavp=\frac{\pavp^{p+1}}p,\qquad
\Xi=\sum_{e\in\paI}\sum_{j\ge0}\lambda_{e,j}\Gamma_e^{\,p^j},
\end{equation}
with $W$ integral and $\lambda_{e,j}\in\pao$; this is \eqref{eq:au-normal}. By
\eqref{pa:eq-const}, $v(c_\pavp\pi)=\frac1{p-1}$. Hence, by (E1), \emph{it suffices to know the
exponent modulo $\pi$}: if $F(\pavp z)=W\exp(c_\pavp\Xi')$ with $W,\Xi'$ integral,
$\Xi'(0)=0$ and $\Xi'\equiv\Xi\pmod\pi$, then $F$ has the normal form (replace $W$ by
$W\exp(c_\pavp(\Xi'-\Xi))$). Also, if $F$ has the normal form and $\alpha\in\pao$, then so has
$F(\alpha z)$, with $\lambda_{e,j}$ replaced by $\alpha^{(p^e+1)p^j}\lambda_{e,j}$, the couplings
being homogeneous.

Now let $\mathsf C=[n']\times[n]$, so that for each $\tau$ the variables form an $n'\times n$
matrix $z_\tau=(z_{\tau,b,a})$. A \emph{shape} $\kappa$ gives numbers $\kappa_\tau(h)\ge0$ of
columns of height $h$ at $\tau$, for $1\le h<n'$. Put
\[
a_\tau=\kappa_\tau(1),\qquad a'_\tau=\sum_{h\ge2}\kappa_\tau(h),\qquad
l_\tau=\sum_{h\ge2}h\,\kappa_\tau(h),\qquad B=\sum_\tau(a_\tau+l_\tau),
\]
the numbers of columns of height one, of the other (\emph{low}) columns, of the boxes in these, and
of all boxes. A \emph{filling} $T$ of $\kappa$ assigns to each column $\mathfrak c=(\tau,h,c)$,
$1\le c\le\kappa_\tau(h)$, a sequence $T(\mathfrak c)=(b_1,\dots,b_h)$ of rows in $[n']$. For
matrices $y_\tau\in M_{n\times(n'-1)}(\pao)$ (\emph{directions}) let
\[
Q_{T,y}(z)=\prod_{\mathfrak c=(\tau,h,c)}\det\bigl((z_\tau y_\tau)_{b_i,k}\bigr)_{1\le i,k\le h},
\qquad(b_1,\dots,b_h)=T(\mathfrak c),
\]
the product over the columns of the minors of $z_\tau y_\tau$ on the rows of the column and on the
first $h$ columns; it is homogeneous of degree $B$ with integral coefficients. For $e\in\paI$ and
$\tau\in\auPl$, with $\tau'=\mathrm{pr}_e\tau$, the \emph{deduction} is
\[
\paded_e(\tau)=\min\Bigl(\frac{a_\tau}{p^e},\,a_{\tau'}+a'_{\tau'}\Bigr)+\frac{a'_\tau}{p^e}
+\min\Bigl(\frac{l_\tau}{p^e},\,l_{\tau'}\Bigr),
\]
and the \emph{predicted divisibility} of $\kappa$ is
\[
\paPred(\kappa)=v(\pi)\sum_{e\in I_p}\sum_{\tau}\Bigl(\bigl(\tfrac{a_\tau}{p^e}-a_{\tau'}-
a'_{\tau'}\bigr)^+ +\bigl(\tfrac{l_\tau}{p^e}-l_{\tau'}\bigr)^+-\tfrac{a'_\tau}{p^e}\Bigr)
-\mathbf 1_{p=d}\cdot\tfrac12\sum_\tau\paded_0(\tau).
\]
Since $u-\min(u,u')=(u-u')^+$, the second identity of \eqref{pa:eq-const} gives
\begin{equation}\label{pa:eq-pred}
\paPred(\kappa)=\gamma B-v(\pi)\sum_{e\in\paI}\sum_\tau\paded_e(\tau).
\end{equation}

\begin{proposition}[The count; Proposition~\ref{au:count}]\label{pa:count}
Let $F\in\paCp[[z]]$ have the normal form \eqref{pa:eq-nf}, let $\kappa$ be a shape with $B\ge1$
boxes, $T$ a filling of $\kappa$ and $y$ integral directions. Then
\[
v\bigl(Q_{T,y}(\partial)F(0)\bigr)\ \ge\ \gamma B-v(\pi)\!\!\sum_{e\in\paI,\ p^e\le B}\ \sum_\tau
\paded_e(\tau)-f\,n'^3\,(1+\log_pB)\ \ge\ \paPred(\kappa)-f\,n'^3\,(1+\log_pB).
\]
\end{proposition}

For $p$ inert the right-hand side is the expression of Proposition~\ref{au:count}; for $p=d$ the
term subtracted in addition is the $e=0$ term $v(\pi_d)\sum_\tau\paded_0(\tau)=\tfrac12\sum_\tau\paded_0(\tau)$. Together with
Lemma~\ref{pa:schur}, the bound holds for the jets along $Q_{T,y}$ and along the entries
$S_\kappa(z)_{T',T}$ of the Schur matrix; these are the only jets to which it is applied, and no
claim is made for products of minors on non-nested column sets. The proof needs one lemma of
linear algebra.

\begin{lemma}[Elimination]\label{pa:elim}
Let $h\ge1$, let $\ell_1,\dots,\ell_s\colon\pao^h\to\pao$ be linear forms and
$\alpha_1,\dots,\alpha_s\in\pao$. There is $g\in M_h(\pao)$ with $\det g\mid\alpha_1\cdots\alpha_s$
such that $\alpha_i\mid\ell_i(r)$ for every row $r$ of $g$ and every $i$.
\end{lemma}
\begin{proof}
Induction on $s$, starting from $g=1$. Let $g$ serve $\ell_1,\dots,\ell_{s-1}$, with rows
$r_1,\dots,r_h$. If all $\ell_s(r_k)$ vanish, $g$ serves. Otherwise let $k_0$ be an index for which
$v(\ell_s(r_{k_0}))$ is least, replace $r_k$ by $r_k-\bigl(\ell_s(r_k)/\ell_s(r_{k_0})\bigr)r_{k_0}$
for $k\ne k_0$, and $r_{k_0}$ by $\alpha_sr_{k_0}$. The new rows are integral combinations of the
old ones, so the earlier divisibilities persist; $\ell_s$ vanishes on the first kind and is
divisible by $\alpha_s$ on the last; and the determinant has been multiplied by $\alpha_s$.
\end{proof}

\begin{proof}[Proof of Proposition~\ref{pa:count}]
The second inequality is \eqref{pa:eq-pred}, the deductions being non-negative. For the first, let
$S_B=\{e\in\paI:p^e\le B\}$.

\emph{Step 1: a linear substitution.} Introduce new variables $s_{\tau,h,k,b}$ for $\tau\in\auPl$,
$1\le k\le h<n'$ and $b\in[n']$; there are at most $f\,n'^3$ of them. For matrices
$g_{\tau,h}\in M_h(\pao)$, to be chosen, define the \emph{new directions}
$\beta_{\tau,h,k}=\sum_{k'\le h}g_{\tau,h}(k,k')\,y_\tau(\cdot,k')\in\pao^n$ and the substitution
\[
z_{\tau,b,a}(s)=\sum_{h,k}\beta_{\tau,h,k}(a)\,s_{\tau,h,k,b}.
\]
Let $G(s)=F(\pavp z(s))$ and $P_*(s)=\prod_{\mathfrak c=(\tau,h,c)}\det(s_{\tau,h,k,b_i})_{1\le
i,k\le h}$, a homogeneous polynomial of degree $B$ with integer coefficients. For a linear
substitution $z=Ms$ and a polynomial $P$ one has $P(\partial_s)(F\circ M)=(\widetilde
P(\partial_z)F)\circ M$, where $\widetilde P$ is obtained from $P$ by replacing each $s_o$ by the
linear form $\sum_xM_{x,o}z_x$. Here $s_{\tau,h,k,b}$ is replaced by
$\sum_a\beta_{\tau,h,k}(a)z_{\tau,b,a}=\sum_{k'}g_{\tau,h}(k,k')(z_\tau y_\tau)_{b,k'}$, so the
matrix of a column $\mathfrak c$ becomes the corresponding matrix of $z_\tau y_\tau$ times the
transpose of $g_{\tau,h}$. Taking determinants, and using that $Q_{T,y}$ is homogeneous of degree
$B$,
\begin{equation}\label{pa:eq-chain}
P_*(\partial)G(0)=\pavp^{B}\prod_{\tau,h}\det(g_{\tau,h})^{\kappa_\tau(h)}\cdot Q_{T,y}(\partial)F(0).
\end{equation}

\emph{Step 2: the couplings after substitution.} For $\beta,\beta'\in\pao^n$ put
$[\beta,\beta']_e=\sum_a\beta(a)^{p^e}\beta'(a)$. By (E3), for $\tau'=\mathrm{pr}_e\tau$,
\[
\Gamma_e(z(s))\equiv\sum_{\tau,b}\ \sum_{h,k}\sum_{h',k'}[\beta_{\tau,h,k},\beta_{\tau',h',k'}]_e\;
s_{\tau,h,k,b}^{\,p^e}\,s_{\tau',h',k',b}\pmod p .
\]
We claim that the $g_{\tau,h}$ can be chosen such that
\begin{equation}\label{pa:eq-goal}
\pi\mid[\beta_{\tau,h,k},\beta_{\mathrm{pr}_e\tau,h',k'}]_e\ \ \text{for all }e\in S_B,\ \tau,\
h,k,h',k',\qquad
\sum_{\tau,h}\kappa_\tau(h)\,v(\det g_{\tau,h})\le v(\pi)\sum_{e\in S_B}\sum_\tau\paded_e(\tau).
\end{equation}
For $e\in S_B$ fix $\pi_e\in\pao$ with $\pi_e^{\,p^e}=\pi$, and for $w\in\pao^n$ a vector
$w^{1/p^e}$ of $p^e$-th roots of its coordinates. By (E3),
\begin{equation}\label{pa:eq-root}
\pi_e\mid\textstyle\sum_aw^{1/p^e}(a)\beta(a)\ \Longrightarrow\ \pi\mid[\beta,w]_e .
\end{equation}
For each pair $(\tau,e)$, with $\tau'=\mathrm{pr}_e\tau$, we impose the following congruences on
the rows of the matrices at $\tau$ (the \emph{source}) and at $\tau'$; recall that a new direction
of height $1$ at a place $\sigma$ is a multiple of $y_\sigma(\cdot,1)$.
\begin{itemize}
\item[(a)] On every row $r$ of $g_{\tau,h}$, $h\ge2$: one congruence modulo $\pi_e$, namely
$\pi_e\mid\sum_a y_{\tau'}(a,1)^{1/p^e}\beta_r(a)$, where $\beta_r=\sum_{k'}r_{k'}y_\tau(\cdot,k')$.
By \eqref{pa:eq-root} it gives $\pi\mid[\beta,\beta']_e$ for $\beta$ of height $h\ge2$ at $\tau$
and $\beta'$ of height $1$ at $\tau'$.
\item[(b)] If $a_\tau/p^e\le a_{\tau'}+a'_{\tau'}$: $\pi_e$ divides the $1\times1$ matrix
$g_{\tau,1}$. Otherwise: on every row $r$ of every $g_{\tau',h'}$, $h'\ge1$, one congruence modulo
$\pi$, namely $\pi\mid[y_\tau(\cdot,1),\beta'_r]_e$, which is linear in $r$. In both cases
$\pi\mid[\beta,\beta']_e$ for $\beta$ of height $1$ at $\tau$ and all $\beta'$ at $\tau'$.
\item[(c)] If $l_\tau/p^e\le l_{\tau'}$: $\pi_e$ divides all entries of $g_{\tau,h}$, $h\ge2$
($h$ congruences on each row). Otherwise: $\pi$ divides all entries of $g_{\tau',h'}$, $h'\ge2$. In
the first case $\pi\mid[\beta,\beta']_e$ for all $\beta$ of height $\ge2$ at $\tau$; in the second,
for all $\beta'$ of height $\ge2$ at $\tau'$.
\end{itemize}
Every pair of heights $(h,h')$ is covered: $h=1$ by (b); $h\ge2$, $h'=1$ by (a); $h,h'\ge2$ by
(c) or, in its second case, by the congruences on $\beta'$. The conditions imposed on the rows of
one matrix $g_{\sigma,h}$ by all pairs $(\tau,e)$ with $\tau=\sigma$ or $\mathrm{pr}_e\tau=\sigma$
are finitely many congruences of linear forms, so Lemma~\ref{pa:elim} gives $g_{\sigma,h}$ whose
determinant has valuation at most the sum of the valuations of the moduli. A modulus $\pi_e$ costs
$v(\pi)/p^e$ and a modulus $\pi$ costs $v(\pi)$. Counting each matrix $\kappa_\sigma(h)$ times, the
pair $(\tau,e)$ contributes: from (a), $v(\pi)a'_\tau/p^e$; from (b),
$v(\pi)\min(a_\tau/p^e,a_{\tau'}+a'_{\tau'})$; from (c), $v(\pi)\min(l_\tau/p^e,l_{\tau'})$. The
sum is $v(\pi)\paded_e(\tau)$, which proves \eqref{pa:eq-goal}.

\emph{Step 3: $G$ is integral in degrees at most $B$.} By \eqref{pa:eq-nf},
$G=W(z(s))\exp(c_\pavp\Xi(z(s)))$. For $e\in\paI\setminus S_B$ the series $\Gamma_e(z(s))^{p^j}$ has
no terms of degree $\le B$. For $e\in S_B$, Step~2 and $\pi\mid p$ give
$\pi\mid\Gamma_e(z(s))$. Hence $\Xi(z(s))=\pi h+h'$ with $h,h'$ integral, $h(0)=0$, and $h'$ of
order $>B$. As $v(c_\pavp\pi)=\frac1{p-1}$, (E1) shows that $\exp(c_\pavp\pi h)$ is integral; and
$G$ agrees with $W(z(s))\exp(c_\pavp\pi h)$ in degrees $\le B$.

\emph{Step 4: Legendre.} Write $P_*=\sum_\delta c_\delta s^\delta$ with $c_\delta\in\Z$,
$|\delta|=B$. Then $P_*(\partial)G(0)=\sum_\delta c_\delta\,\delta!\,G_\delta$, where $G_\delta$ is
the coefficient of $s^\delta$ in $G$, integral by Step~3. By \eqref{pa:eq-legendre}, applied to each
of the at most $fn'^3$ exponents of $\delta$,
$v(\delta!)\ge\frac B{p-1}-fn'^3(1+\log_pB)$. With \eqref{pa:eq-chain}, \eqref{pa:eq-goal} and
$v(\pavp^B)=\rho B$,
\[
v\bigl(Q_{T,y}(\partial)F(0)\bigr)\ge\frac B{p-1}-fn'^3(1+\log_pB)-\rho B-v(\pi)\sum_{e\in S_B}
\sum_\tau\paded_e(\tau),
\]
and $\frac1{p-1}-\rho=\gamma$.
\end{proof}

\begin{remark}[Marked places]\label{pa:charges}
A variant is used at the exceptional primes (Lemma~\ref{pa:hecke}); it is proved as
Proposition~\ref{pl:charges}. Let
$\mathsf E\subseteq\auPl$ and let integers $w_{\tau,h,b}\ge0$ (\emph{charges}) be such that, for
$e\in\paI$, $\tau'=\mathrm{pr}_e\tau$, every row $b$ and every height $h'$: if $\tau\in\mathsf E$,
or $h'=1$ and $\tau'\in\mathsf E$, then $w_{\tau,1,b}>0$ or $w_{\tau',h',b}>0$. Substitute
$z_{\tau,b,a}(s)=\sum_{h,k}p^{w_{\tau,h,b}}\beta_{\tau,h,k}(a)s_{\tau,h,k,b}$. The pairs of heights
$(1,h')$ of the hypothesis then need no congruence, since the coefficient of
$s_{\tau,1,k,b}^{p^e}s_{\tau',h',k',b}$ acquires the factor $p^{p^ew_{\tau,1,b}+w_{\tau',h',b}}$,
which is divisible by $p$; so the columns of height one at the places of $\mathsf E$ may be left out of the
deductions, and the result is: $v(Q_{T,y}(\partial)F(0))+\sum_{\text{boxes}}w\ge
\paPred(\kappa^\circ)+\gamma(B-B^\circ)-fn'^3(1+\log_pB)$, where $\kappa^\circ$ is $\kappa$
without its columns of height one at $\mathsf E$, $B^\circ$ its number of boxes, and the sum is
over the boxes of $T$, a box in row $b$ of a column of height $h$ at $\tau$ counting
$w_{\tau,h,b}$.
\end{remark}

\subsubsection*{From entries of the Schur matrix to products of minors}
Let $n\le n'$ and let $\kappa$ be a shape all of whose columns have height at most $n$. For a
generic $n\times\infty$ matrix $Y$ and a filling $T$ of a diagram by elements of $[n]$ let
$\nabla_T(Y)$ be the product over the columns of the minors of $Y$ on the rows of the column and
the first $h$ columns. By standard monomial theory~\cite[\S\S1--2, Thm.~2.1]{DEP}, \cite[\S8.1, Thm.~1 and its corollary]{FultonYT} the $\nabla_T$ with $T$
standard form a $\Z$-basis of the span of all $\nabla_T$ of a given shape, and they remain linearly
independent over every field; the lattice $W_\kappa(\Z)$ of \S\ref{ss:au-spaces} is the dual of the tensor product over $\tau$ of
these spans, with the basis $(e_T)$ dual to the standard $\nabla_T$ (\S\ref{pl:ss-norm}, (N1)). For $z_\tau\in M_{n'\times n}$ and a standard filling $T'$ by
elements of $[n']$, the Cauchy--Binet formula and the straightening law give
\[
\prod_\tau\nabla_{T'_\tau}(z_\tau Y_\tau)=\sum_{T\text{ standard}}S_\kappa(z)_{T',T}\,\prod_\tau
\nabla_{T_\tau}(Y_\tau),\qquad S_\kappa(z)_{T',T}\in\Z[z];
\]
$S_\kappa(z)$ is the \emph{Schur matrix}, the matrix of the map $W_\kappa(\C^n)\to W_\kappa(\C^{n'})$,
$\mathrm{ev}_Y\mapsto\mathrm{ev}_{zY}$, in the bases $(e_T)$.

\begin{lemma}\label{pa:schur}
For $F\in\paCp[[z]]$ and $T'$ standard,
$\min_Tv\bigl(S_\kappa(\partial)_{T',T}F(0)\bigr)\ge\inf_yv\bigl(Q_{T',y}(\partial)F(0)\bigr)$,
the infimum over all integral directions $y$.
\end{lemma}
\begin{proof}
Put $J_T=S_\kappa(\partial)_{T',T}F(0)$. Specialising $Y_\tau$ to $y_\tau$ in the identity above,
$Q_{T',y}(\partial)F(0)=\sum_TJ_T\nabla_T(y)$. Let $T_0$ have $|J_{T_0}|$ largest; we may assume
$J_{T_0}\ne0$. The polynomial $\sum_T(J_T/J_{T_0})\nabla_T(Y)$ is integral and its reduction modulo
$\pam$ is not zero, by the linear independence of the $\nabla_T$ over the residue field. This
field is infinite, so the reduction does not vanish at some point, which lifts to an integral $y$
with $v(Q_{T',y}(\partial)F(0))=v(J_{T_0})$.
\end{proof}

\subsubsection{The theta function of a CM curve at a non-split prime}\label{pa:ss-curve}
Let $\Lambda\subset\C$ be a lattice with $\auO_K\Lambda=\Lambda$, of area $A$, with Weierstrass
functions $\sigma,\wp$, and let $t\in\tfrac12\Lambda\setminus\Lambda$ be the \emph{admissible} point: the one with
$\omega t\notin\Lambda$ and $\aubar\omega t\notin\Lambda$, $\omega=(1+\sqrt{-d})/2$ (\S\ref{pf:ss-adm}). After scaling
$\Lambda$ we may assume $\wp''(t)=2$. Put
\[
x(z)=\wp(z+t)-\wp(t),\qquad a_1^2=12\,\wp(t),\qquad y(z)=\tfrac12\bigl(\wp'(z+t)-a_1x(z)\bigr).
\]
With $e_1=\wp(t)$ and $e_2,e_3$ the other roots of the cubic of $\wp$ one has $e_2+e_3=-e_1$ and
$(e_1-e_2)(e_1-e_3)=\wp''(t)/2=1$, so $\wp'^2=4x(x^2+3e_1x+1)$ at $z+t$ and
\begin{equation}\label{pa:eq-curve}
y^2+a_1xy=x^3+x,\qquad x'=2y+a_1x,\qquad y'=1+3x^2-a_1y,
\end{equation}
the last from $\wp''=6\wp^2-g_2/2=6x^2+a_1^2x+2$. The point $z=0$ is the point $P_0=(0,0)$ of
order two of the curve $\mathcal E\colon y^2+a_1xy=x^3+x$, the point at infinity is $z\equiv t$,
and $y=z+O(z^2)$. All Taylor series at $z=0$ below have algebraic coefficients and are read in
$\paCp[[z]]$.

\begin{lemma}\label{pa:model}
$a_1$ is an algebraic integer and the discriminant $a_1^4-64$ of $\mathcal E$ is a unit at every
odd prime. The reduction $\widetilde{\mathcal E}$ modulo $\pam$ is a supersingular elliptic curve,
that is, it has no point of order $p$ over the algebraic closure of the residue field.
\end{lemma}
\begin{proof}
The invariants of $\mathcal E$ are $c_4=a_1^4-48$, $\Delta=a_1^4-64$, so
$(a_1^4-48)^3=j\cdot(a_1^4-64)$. As $j$ is an algebraic integer~\cite[Chapter~II, Theorem~6.1]{SilvermanAT},
so is $a_1$. If a prime above an odd prime divided $a_1^4-64$ it would divide $a_1^4-48$, hence
$16$. The curve has complex multiplication by $K$ and good reduction, and $p$ does not split in
$K$; by Deuring's criterion~\cite[Chapter~13, Theorem~12]{LangEF} the reduction is supersingular (for a curve over a field
of characteristic $p$, having no point of order $p$ over the algebraic closure is equivalent to
the other usual definitions of supersingularity~\cite[Chapter~V, Theorem~3.1(a)]{SilvermanAEC}).
\end{proof}
For the admissible point $t$, $a_1$ is a unit at the places above $2$ as well, so that $a_1^4-64$
is a unit at every prime (Lemma~\ref{pf:two}).

\emph{The ring $R$.} Let $R=\pao[[y]]\subset\paCp[[z]]$ be the ring of the series $\sum c_ny(z)^n$
with $c_n\in\pao$; it is closed under coefficientwise limits. Solving \eqref{pa:eq-curve} for $x$
by iteration, $x=y^2+a_1xy-x^3\in y^2(1+yR)$; hence $R$ contains $x$ and, by
\eqref{pa:eq-curve}, is stable under $d/dz$. Geometrically $R$ is the completion of
$\mathcal E_{/\pao'}$ along the section $P_0$, for $\pao'\subset\pao$ a discrete valuation ring
containing $a_1$, tensored with $\pao$.

\emph{Endomorphisms.} Let $\alpha\in\auO_K$, $\alpha\equiv1\pmod{2\auO_K}$, $\alpha\ne0$. Since
$(\alpha-1)t\in\Lambda$, multiplication by $\alpha$ on $\C/\Lambda$ fixes $P_0$ and is $z\mapsto
\alpha z$ in the coordinate $z$. The smooth Weierstrass scheme $\mathcal E_{/\pao'}$ is the N\'eron
model of its generic fibre~\cite[\S1.2, Definition~1, p.~12, and Proposition~8, p.~15]{BLR} (\cite[IV, Corollary~6.3, p.~329, and the N\'eron mapping property, p.~319]{SilvermanAT}), so $[\alpha]$ extends to it and preserves the
section $P_0$. Therefore
\begin{equation}\label{pa:eq-endo}
G\in R\ \Longrightarrow\ G(\alpha z)\in R\qquad(\alpha\in1+2\auO_K).
\end{equation}
Let $K_p$ be the closure of $K$ in $\paCp$. The coefficients of $G(\alpha z)$ in powers of $y$ are
polynomials in $\alpha$, and $1+2\auO_K$ is dense in $\auO_{K_p}$ because $p$ is odd; so
\eqref{pa:eq-endo} holds for all $\alpha\in\auO_{K_p}$. In particular it holds for the group
$\mu_{q-1}\subset\auO_{K_p}$ of roots of unity of order prime to $p$. Note that $\pi\equiv1
\pmod{2\auO_K}$ in both cases, since $(\sqrt{-d}-1)/2\in\auO_K$.

\begin{lemma}[The parameter]\label{pa:param}
Let $T(z)=\frac1{q-1}\sum_{\zeta\in\mu_{q-1}}\zeta^{-1}y(\zeta z)$, the part of $y$ in the degrees
$\equiv1\pmod{q-1}$.
\begin{enumerate}
\item[(a)] $T\in y+y^2R$, so $R=\pao[[T]]$; and $T(\pi z)=\pi T+\sum_{k\ge1}b_kT^{1+k(q-1)}$ with
$b_k\in\pao$ and $b_1\in\pao^\times$.
\item[(b)] Let $\pavp^{q-1}=-\pi/b_1$; then $v(\pavp)=\rho$. The series $Y(z)=T(\pavp z)/\pavp$ is
integral, and $Y\equiv\sum_{j\ge0}z^{q^j}\pmod\pi$.
\item[(c)] For every $G\in R$ the series $G(\pavp z)$ is integral, and the coefficients of $G$
itself in the degrees $<q$ in $z$ are integral. In particular this holds for $x$ and $y$.
\end{enumerate}
\end{lemma}
\begin{proof}
(a) By \eqref{pa:eq-endo}, $T\in R$, and $T=z+\dots=y+\dots$. Also $T(\zeta z)=\zeta T(z)$, and
$T(\pi z)\in R=\pao[[T]]$ by \eqref{pa:eq-endo}; writing $T(\pi z)=P(T)$ we get $P(\zeta T)=\zeta
P(T)$, which gives the shape of the expansion, with first coefficient $\pi$. Modulo $\pam$, $P$
is the action of the reduction of $[\pi]$ on a parameter at the point $\widetilde P_0$ of
$\widetilde{\mathcal E}$. This reduction is an isogeny of degree $\auNm(\pi)=q$ (reduction of homomorphisms between curves with good reduction preserves degrees~\cite[Chapter~II, Proposition~4.4]{SilvermanAT}), a power of $p$, and
$\widetilde{\mathcal E}$ has no point of order $p$ (Lemma~\ref{pa:model}). Let $\widetilde O$ be the origin of $\widetilde{\mathcal E}$. The kernel of
$\widetilde{[\pi]}$ consists of points killed by $\widetilde{[\aubar\pi]}\circ\widetilde{[\pi]}=[q]$;
if it contained a point $Q\ne\widetilde O$, then for the least $k$ with $p^kQ=\widetilde O$ the
point $p^{k-1}Q$ would have order $p$. So $\ker\widetilde{[\pi]}=\{\widetilde O\}$. Since
$\pi\equiv1\pmod2$ and $P_0$ has order $2$, $\widetilde{[\pi]}(\widetilde P_0)=\widetilde P_0$, so the
fibre of $\widetilde{[\pi]}$ over $\widetilde P_0$ is $\widetilde P_0+\ker\widetilde{[\pi]}=\{\widetilde P_0\}$. For a nonconstant morphism $\phi$ of smooth projective curves the
ramification indices over any point add up to $\deg\phi$~\cite[Chapter~II, Proposition~2.6(a)]{SilvermanAEC}; hence the
ramification index of $\widetilde{[\pi]}$ at $\widetilde P_0$ is $q$, which says exactly that
$P\bmod\pam$, the pull-back of the uniformiser $T$, has order exactly $q$ in $T$. Hence $b_1$ is a unit.

(b) Replacing $z$ by $\pavp z$ in (a) and dividing by $\pavp$,
\[
Y(\pi z)=\pi\bigl(Y-Y^q+\pi E(Y)\bigr),\qquad E(Y)=\sum_{k\ge2}b_k(-b_1)^{-k}\pi^{k-2}\,
Y^{1+k(q-1)}\in Y^{2q-1}\pao[[Y]].
\]
Write $Y=\sum_{n\ge1}c_nz^n$, $c_1=1$. Comparing the coefficients of $z^n$, $n\ge2$:
$(1-\pi^{n-1})c_n=[Y^q]_n-\pi[E(Y)]_n$, and the right side involves only $c_k$ with $k<n$. By
induction all $c_n$ are integral, and $c_n-[Y^q]_n\in\pi\pao$ for $n\ge2$, that is,
$Y\equiv z+Y^q\pmod\pi$. Iterating with (E3) gives $Y\equiv\sum_jz^{q^j}$.

(c) If $G=\sum g_nT^n$ with $g_n\in\pao$, then $G(\pavp z)=\sum g_n\pavp^nY^n$ is integral. By
(b), $T(z)=\sum_nc_n\pavp^{1-n}z^n$ with $\pi\mid c_n$ for $1<n<q$, and
$v(\pi)-(n-1)\rho=(q-n)\rho>0$; so $T\equiv z$ modulo $\pam$ and $z^q$, and the second claim
follows.
\end{proof}

\emph{The theta function.} Let $H(u,w)=\boldsymbol\pi\,\aubar uw/A$, where
$\boldsymbol\pi=3.14\ldots$ (elsewhere $\pi$ is $\pi_p$), and $s_2=\lim_{s\to0^+}\sum_{\omega\in\Lambda\setminus0}\omega^{-2}
|\omega|^{-2s}$. The function $\vartheta(z)=e^{-s_2z^2/2}\sigma(z)$ satisfies
$\vartheta(z+\omega)=\varepsilon(\omega)e^{H(\omega,z+\omega/2)}\vartheta(z)$ for $\omega\in\Lambda$,
with $\varepsilon(\omega)=1$ if $\omega\in2\Lambda$ and $-1$ otherwise~\cite[\S1.2]{BannaiKobayashi}.
Put
\[
\theta(z)=\frac{\vartheta(z+t)}{\vartheta(t)}\,e^{-H(t,z)},\qquad g=-\log\theta\in z\,\C[[z]].
\]
Then $\theta(0)=1$, the zeros of $\theta$ are simple and are the points of $t+\Lambda$, and
\begin{equation}\label{pa:eq-theta}
\theta(z+\omega)=\chi(\omega)\,e^{H(\omega,z+\omega/2)}\,\theta(z),\qquad
\chi(\omega)=\varepsilon(\omega)e^{H(\omega,t)-H(t,\omega)}\in\{\pm1\},
\end{equation}
where $\chi(\omega)$ depends only on $\omega$ modulo $2\Lambda$ (as $2t\in\Lambda$ and the
imaginary part of $H$ is $\boldsymbol\pi$ times an integral form on $\Lambda$). The function
$\theta(-z)$ satisfies \eqref{pa:eq-theta} and has the same zeros, so $\theta$ is even. Finally
\begin{equation}\label{pa:eq-g2}
g''=x+2g_2,\qquad g_2=\tfrac12(s_2+\wp(t))\ \text{the coefficient of }z^2\text{ in }g,
\end{equation}
because $-(\log\sigma)''=\wp$.

\begin{lemma}[Units]\label{pa:units}
\begin{enumerate}
\item[(a)] For $\alpha\in1+2\auO_K$, $\alpha\ne0$, the function $f_\alpha(z)=\theta(\alpha z)/
\theta(z)^{\auNm\alpha}$ is a polynomial $R_\alpha(x)$ in $x$ of degree $(\auNm\alpha-1)/2$ with
$R_\alpha\in\pao[x]$, $R_\alpha(0)=1$. In particular $f_\alpha\in1+yR$.
\item[(b)] For $\zeta\in\mu_{q-1}$, $\exp\bigl(\zeta^{p+1}g(z)-g(\zeta z)\bigr)\in1+yR$.
\item[(c)] $\pi g_2\in\pao$, and $\pi g_n\in\pao$ for all $n<q$, where $g=\sum g_nz^n$.
\end{enumerate}
\end{lemma}
\begin{proof}
(a) By \eqref{pa:eq-theta}, $\chi(\alpha\omega)=\chi(\omega)=\chi(\omega)^{\auNm\alpha}$
($\auNm\alpha$ is odd) and $H(\alpha\omega,\alpha w)=\auNm\alpha\cdot H(\omega,w)$; so $f_\alpha$ is
$\Lambda$-periodic. In the variable $u=z+t$ its zeros are the points of
$\alpha^{-1}\Lambda\setminus\Lambda$ (as $\alpha t\equiv t$), all simple, its only pole is $u=0$,
of order $\auNm\alpha-1$, and it is even. Hence $f_\alpha=\prod_s(1-x/x_s)$, the product over the
non-zero points $s$ of $\alpha^{-1}\Lambda/\Lambda$ up to sign, with $x_s=\wp(s)-\wp(t)$. We claim
$|x_s|\ge1$. If $|x_s|<1$, then the point $s$ of $\mathcal E$ has $|y(y+a_1x)|=|x(1+x^2)|<1$, hence
$|y|<1$, and reduces to $\widetilde P_0$; as reduction is a homomorphism, $s-t$ lies in its kernel
$E_1$, and so does $[\auNm\alpha](s-t)=-t$ ($\auNm\alpha$ is odd and kills $s$). But $t$ does not
reduce to the origin. This proves the claim, and (a) follows since $x\in y^2R$.

(b) Choose $\alpha_n\in1+2\auO_K$ with $\alpha_n\to\zeta$ in $\auO_{K_p}$. Then
$\auNm\alpha_n=\alpha_n\aubar{\alpha_n}\to\zeta\cdot\zeta^{p}$ if $p$ is inert (the conjugation of
$K_p$ is the Frobenius) and $\to\zeta^2=\zeta^{p+1}$ if $p=d$ ($\zeta\in\Z_p$). The coefficients of
$f_{\alpha_n}=\exp(\auNm\alpha_n\,g(z)-g(\alpha_nz))$ are polynomials in
$(\alpha_n,\auNm\alpha_n)$, and $1+yR$ is closed.

(c) By (a) and Lemma~\ref{pa:param}(c), the coefficient of $z^2$ in $\log f_\alpha=\auNm\alpha\,
g(z)-g(\alpha z)$ is integral; it is $\alpha(\aubar\alpha-\alpha)g_2$. For
$\alpha=2+\sqrt{-d}$ this is $-2\sqrt{-d}\,(2+\sqrt{-d})g_2$, and $2(2+\sqrt{-d})$ is a unit at
$p$: if $p=d$ its norm $4(4+d)$ is prime to $p$, and if $p$ is inert then $p$ is a prime element
of $\auO_K$ that does not divide $2+\sqrt{-d}$. So $\sqrt{-d}\,g_2$ is integral, and hence
$\pi g_2$ is. For $3\le n<q$, \eqref{pa:eq-g2} gives
$n(n-1)g_n=x_{n-2}$, the coefficient of $z^{n-2}$ in $x$, integral by Lemma~\ref{pa:param}(c); and
$v(n(n-1))\le v(\pi)$ for $n<q$. Finally $g_1=0$.
\end{proof}

\begin{proposition}[Frobenius congruence]\label{pa:frob}
$q\,g(z)-g(\pi z)=\pi H$ for some $H\in yR$.
\end{proposition}
\begin{proof}
By Lemma~\ref{pa:units}(a), $\exp(q\,g(z)-g(\pi z))=f_\pi=R_\pi(x)$ with $R_\pi=1+\sum_{i=1}^{(q-1)/2}
r_ix^i$, $r_i\in\pao$. We show $\pi\mid r_i$ for all $i$. The coefficient of $z^n$ in
$q\,g(z)-g(\pi z)$ is $(q-\pi^n)g_n=(q/\pi-\pi^{n-1})\cdot\pi g_n$; for $2\le n<q$ it is divisible
by $\pi$ by Lemma~\ref{pa:units}(c), and it vanishes for $n=1$. By (E2) with $\beta=\pi$,
$R_\pi(x(z))\equiv1$ modulo $\pi$ in all degrees $<q$ in $z$. By Lemma~\ref{pa:param}(c) and
$x=z^2+\dots$, the series $x(z)^i$ is integral in degrees $<q$ with lowest term $z^{2i}$, and
$2i\le q-1$. Comparing the coefficients of $z^2,z^4,\dots,z^{q-1}$ in turn gives $\pi\mid r_1$, then
$\pi\mid r_2$, and so on. Thus $R_\pi(x)=1+\pi H_0$ with $H_0\in x\pao[x]\subset yR$, and
$q\,g(z)-g(\pi z)=\log(1+\pi H_0)\in\pi\,yR$ by (E2).
\end{proof}

\begin{theorem}[Normal form on one curve]\label{pa:onecurve}
Let $\pavp$ be as in Lemma~\ref{pa:param}(b) and $c_\pavp=\pavp^{p+1}/p$. There are integral series
$W,\Xi_1\in\pao[[z]]$, $\Xi_1(0)=0$, and $\lambda\in\pao$ such that
\[
\theta(\pavp z)=W(z)\exp\bigl(c_\pavp\,\Xi_1(z)\bigr),\qquad
\Xi_1\equiv\sum_{e\in\paI}\sum_{j\ge0}\lambda_{e,j}\bigl(z^{p^e}\cdot z\bigr)^{p^j}\pmod\pi ,
\]
where $\lambda_{e,j}=\lambda$ if $p$ is inert, and $\lambda_{0,0}=\lambda/2$,
$\lambda_{0,j}=\lambda$ $(j\ge1)$, $\lambda_{e,j}=2\lambda$ $(e\ge1)$ if $p=d$.
\end{theorem}
\begin{proof}
\emph{Step 1: averaging.} For a series $G$ let $G^\flat=\frac1{q-1}\sum_{\zeta\in\mu_{q-1}}
\zeta^{-(p+1)}G(\zeta z)$ be its part in the degrees $\equiv p+1\pmod{q-1}$. Then
\[
\exp(g^\flat-g)=\prod_{\zeta}\Bigl(\exp\bigl(\zeta^{p+1}g(z)-g(\zeta z)\bigr)\Bigr)^{\epsilon_\zeta},
\qquad\epsilon_\zeta=-\frac{\zeta^{-(p+1)}}{q-1}\in\Z_p ,
\]
because $\zeta^{p+1}\in\mu_{p-1}\subset\Z_p$. A $p$-adic integral power of an element of $1+yR$ is
in $1+yR$, so $U:=\theta\exp(g^\flat)\in1+yR$ by Lemma~\ref{pa:units}(b). Put $W(z)=U(\pavp z)$,
integral by Lemma~\ref{pa:param}(c), and $\Xi_1(z)=-g^\flat(\pavp z)/c_\pavp$. Then
$\theta(\pavp z)=W\exp(c_\pavp\Xi_1)$.

\emph{Step 2: a functional equation.} The operation $\flat$ commutes with $z\mapsto\pi z$ and
preserves $R$, by \eqref{pa:eq-endo}; so Proposition~\ref{pa:frob} gives
$q\,g^\flat(z)-g^\flat(\pi z)=\pi H^\flat$ with $H^\flat\in R$. Expanding $H^\flat$ in powers of
$T$, which has only degrees $\equiv1\pmod{q-1}$, only the powers $T^n$ with $n\equiv p+1$ occur.
Let $i_0$ be the least positive integer $\equiv p+1\pmod{q-1}$: $i_0=p+1$ if $p$ is inert, $i_0=2$
if $p=d$ (then $q-1=d-1\ge6$). So $H^\flat=h\,T^{i_0}+T^{i_0+q-1}H_1$ with $h\in\pao$, $H_1\in R$.
Replace $z$ by $\pavp z$, use $T(\pavp z)=\pavp Y$ and divide by $-q\,c_\pavp$:
\[
\Xi_1(z)-q^{-1}\,\Xi_1(\pi z)=-u\bigl(h\,Y^{i_0}+\pavp^{q-1}Y^{i_0+q-1}H_1(\pavp z)\bigr),\qquad
u=\frac{\pi\,\pavp^{i_0}}{q\,c_\pavp}.
\]
One checks $v(u)=1-v(\pi)+(i_0-p-1)\rho=0$ in both cases, and $v(\pavp^{q-1})=v(\pi)$. With
$\lambda=-uh$, for every $n$,
\begin{equation}\label{pa:eq-fe}
\bigl(1-\pi^n/q\bigr)\,[\Xi_1]_n\equiv\lambda\,[Y^{i_0}]_n\pmod{\pi\pao}.
\end{equation}

\emph{Step 3: solution.} $\Xi_1$ has only terms of degree $n\ge i_0$, $n\equiv i_0\pmod{q-1}$. For
$n\ge3$, $v(\pi^n/q)=(n-2)v(\pi)\ge v(\pi)$, so \eqref{pa:eq-fe} shows that $[\Xi_1]_n$ is integral
and $\equiv\lambda[Y^{i_0}]_n$. If $p$ is inert, all $n$ are $\ge p+1\ge4$; thus
$\Xi_1\equiv\lambda Y^{p+1}$, and by Lemma~\ref{pa:param}(b) and (E3)
\[
Y^{p+1}\equiv\Bigl(\sum_{i\text{ odd}}z^{p^i}\Bigr)\Bigl(\sum_{i\text{ even}}z^{p^i}\Bigr)
=\sum_{e\text{ odd}}\ \sum_{j\ge0}\bigl(z^{p^e}z\bigr)^{p^j},
\]
each pair (odd $i$, even $i'$) being written once as $\{j,j+e\}$. If $p=d$, then for $n=2$ one has
$\pi^2/q=-1$ and $[Y^2]_2=1$, so $[\Xi_1]_2\equiv\lambda/2$; hence
$\Xi_1\equiv\lambda(Y^2-\tfrac12z^2)$ and
$Y^2\equiv\sum_iz^{2p^i}+2\sum_{i<i'}z^{p^i+p^{i'}}$, which gives the stated coefficients with
$j=i$, $e=i'-i$.
\end{proof}

\begin{remark}
For results of this kind for one curve at supersingular primes (radius of convergence of the
$p$-adic theta function, divisibility of its coefficients) see \cite{BannaiKobayashi} and
\cite{KatzDivisibilities}; we do not use them. Lemma~\ref{pa:param} says that $[\pi]$ is, in the
parameter $T$, an endomorphism of the kind studied by Lubin and Tate~\cite{LubinTate}.
\end{remark}

\subsubsection{Products of theta functions along a resolution of the identity}\label{pa:ss-prod}
We identify $\auPl$ with $\operatorname{Gal}(E/K)=(\Z/q_1)^\times$: $\tau_\nu=\tau_1\circ\sigma_\nu$,
where $\tau_1$ is the given embedding and $\sigma_\nu(\zeta_{q_1})=\zeta_{q_1}^{\nu}$; then
$u(\nu)\tau=\tau\circ\sigma_\nu$. Each $\tau$ gives $E\to\C\to\paCp$.

\begin{lemma}[Powers and partners]\label{pa:partner}
Let $e\in\paI$ and let $a\in E$ be integral at all primes above $p$. For $\tau\in\auPl$ and
$\tau'=u(-p^e)\tau$,
\[
\bigl(\aubar{\tau(a)}\bigr)^{p^e}\equiv\tau'(a)\pmod{\pi\pao}.
\]
\end{lemma}
\begin{proof}
Multiplying by an integer prime to $p$ we may assume $a\in\auO_E=\auO_K[\zeta_{q_1}]$ (the
discriminants of $K$ and of $\Q(\zeta_{q_1})$ are coprime). Write $a=\sum_ic_i\zeta_{q_1}^i$ with
$c_i\in\auO_K$. If $p$ is inert and $e$ is odd, $c_i^{p^e}\equiv\aubar{c_i}\pmod p$; if $p=d$,
$c_i^{p^e}\equiv c_i\equiv\aubar{c_i}\pmod{\sqrt{-d}}$. By (E3), in both cases
$a^{p^e}\equiv\sum_i\aubar{c_i}\,\zeta_{q_1}^{ip^e}=\aubar{\sigma_{-p^e}(a)}\pmod{\pi\auO_E}$.
Apply this to $\aubar a$ and then $\tau$: $\tau(\aubar a)^{p^e}\equiv\tau(\sigma_{-p^e}(a))=
\tau'(a)$, and $\tau(\aubar a)=\aubar{\tau(a)}$ because $E$ is a CM field.
\end{proof}

Let $\mathsf C$ be a finite set. For $u\in E^{\mathsf C}$ let $u^\dagger$ be the linear form
$u^\dagger(z)=\sum_{\tau\in\auPl,\,i\in\mathsf C}\aubar{\tau(u_i)}\,z_{\tau,i}$. A \emph{resolution
of the identity} is a finite family $u_1,\dots,u_J\in E^{\mathsf C}$ with multiplicities
$m_1,\dots,m_J\in\N$ such that
\begin{equation}\label{pa:eq-res}
\sum_{j}m_j\,\tau(u_{j,i})\cdot u_j^\dagger(z)=z_{\tau,i}\qquad\text{for all }\tau\in\auPl,\
i\in\mathsf C .
\end{equation}

\begin{proposition}[Normal form of theta products]\label{pa:product}
Let $(u_j,m_j)$ be a resolution of the identity by vectors whose coordinates are integral at all
primes above $p$, let $\mu_j=u_j^\dagger$, and let $G$ be an integral power series in the series
$y(\mu_j(z))$. Then, for every $c\in\pao$, the series
\[
F(z)=F_0(cz),\qquad F_0(z)=\prod_j\theta(\mu_j(z))^{m_j}\cdot G(z),
\]
has the normal form \eqref{pa:eq-nf}, with $\pavp$ as in Lemma~\ref{pa:param}(b), the index set
$\paI$ and the partner maps $\mathrm{pr}_e(\tau)=u(-p^e)\tau$.
\end{proposition}
\begin{proof}
By the remark after \eqref{pa:eq-nf} we may take $c=1$. The forms $\mu_j$ have integral
coefficients. By Theorem~\ref{pa:onecurve} and Lemma~\ref{pa:param}(c),
\[
F_0(\pavp z)=W_0(z)\exp\Bigl(c_\pavp\sum_jm_j\,\Xi_1(\mu_j(z))\Bigr),\qquad
W_0=\prod_jW(\mu_j(z))^{m_j}\cdot G(\pavp z)\ \text{integral}.
\]
It suffices to compute the exponent modulo $\pi$. By Theorem~\ref{pa:onecurve}, (E3) and
$m^{p}\equiv m\pmod p$,
\[
\sum_jm_j\,\Xi_1(\mu_j)\equiv\sum_{e,k}\lambda_{e,k}\Bigl(\sum_jm_j\,\mu_j^{\,p^e}\mu_j\Bigr)^{p^k}
\pmod\pi .
\]
By (E3) and Lemma~\ref{pa:partner}, $\mu_j^{\,p^e}\equiv\sum_{\tau,i}\tau'(u_{j,i})\,
z_{\tau,i}^{\,p^e}$ with $\tau'=u(-p^e)\tau$; hence, by \eqref{pa:eq-res} at the place $\tau'$,
\[
\sum_jm_j\,\mu_j^{\,p^e}\mu_j\equiv\sum_{\tau,i}z_{\tau,i}^{\,p^e}\sum_jm_j\,\tau'(u_{j,i})\,
u_j^\dagger(z)=\sum_{\tau,i}z_{\tau,i}^{\,p^e}\,z_{\tau',i}=\Gamma_e\pmod\pi .\qedhere
\]
\end{proof}

\subsubsection{The composite lift}\label{pa:ss-divis}
We now prove Proposition~\ref{au:divis} as far as stated in the introduction of this subsection.
Let $\ell$ be a prime and $\C\to\C_\ell$ a homomorphism (when $\ell=p$ is non-split, the
conventions of \S\ref{pa:ss-conv} apply). Let $0\le j<m$. The weight $\lambda_{j+1}=\lambda(X_{j+1})$
has, at $\tau$ (see \S\ref{ss:au-spaces}), $a\cdot\mathbf 1_{\tau\in\auPa}$ columns of height $1$,
$\kappa_1(\tau)-\lfloor n_1/2\rfloor$ columns of height $k$ and $\Delta-(t+1)$ columns of each of
the heights $n_1,\dots,n_j$; all heights are $<n_{j+1}$ and $\le n_j$. In the notation of
\S\ref{pa:ss-count}, with $n'=n_{j+1}$, $n=n_j$:
\begin{equation}\label{pa:eq-diag}
a_\tau=a\,\mathbf 1_{\tau\in\auPa},\qquad\kappa(\tau)\le a'_\tau\le\kappa(\tau)+x^{60},\qquad
k\,\kappa(\tau)\le l_\tau\le k\,\kappa(\tau)+x^{60},
\end{equation}
because $a'_\tau-\kappa(\tau)=\tfrac32\Delta-\lfloor n_1/2\rfloor+j(\Delta-t-1)$ and
$l_\tau-k\kappa(\tau)=k(\tfrac32\Delta-\lfloor n_1/2\rfloor)+(\Delta-t-1)\sum_{i\le j}n_i$ are
non-negative ($n_1=k+t+1\le\Delta$) and at most $x^{60}$ ($\Delta=2r^8\le2\cdot3^8x^{40}$,
$n_m\le40x^6$). We write $\paPred_j(p)$ for the
predicted divisibility of this shape, $B_j$ for its number of boxes, and
\[
\paloss_p=f\,(\log b)^2\,\frac ap\,\epsilon_m(p)\qquad(p\in\mathcal N),
\]
so that $\auerr_p=2\paloss_p$ in \eqref{eq:au-err}.

\begin{lemma}[The kernels]\label{pa:kernel}
Let $\ell\ne N$, let $\vartheta$ be a twist with $\ell\nmid\auNm Q_\vartheta$, $\Psi$ a chain datum
of $\vartheta$, and $\Psi'$ as in Proposition~\ref{au:expansion}, so that
$B_{\vartheta,\Psi}=\operatorname{res}\circ C_{\vartheta,\Psi'}$.
\begin{enumerate}
\item[(K0)] For $\phi\in\auS$ and $g_m\in\auU(X_m)(\auA_{E^+,f})$,
\[
(C_{\vartheta,\Psi'}\phi)(g_m)=\sum_{g_0,\dots,g_{m-1}}\ \sum_x\Psi'(x)\;\theta_{m-1}
(\phi^\circ_{m-1,x_{m-1}};g_{m-1},g_m)\cdots\theta_0(\phi^\circ_{0,x_0};g_0,g_1)\,\phi(g_0),
\]
where $g_i$ runs over representatives of the classes of the level of $X_i$ and $\theta_i$ is the
kernel of \S\ref{ss:au-size}.
\item[(K1)] For every $x$ in the support of $\Psi'$, every $j<m$ and all $g,g'$ there is a power
series $F=F_{x,j,g,g'}$ in the variables $z_{\tau,b,a}$ ($\tau\in\auPl$, $b\le n_{j+1}$, $a\le n_j$)
such that the matrix of $g'^{-1}_\ell\circ\theta_j(\phi^\circ_{j,x_j};g,g')\circ g_\ell$ in the
bases $(e_T)$ of \S\ref{pl:ss-norm} is, up to multiplication on both sides by matrices that are
invertible over $\auO_{\C_\ell}$, the matrix
$\bigl(S_{\lambda_{j+1}}(\partial)_{T',T}F(0)\bigr)_{T',T}$ of jets along the entries of the Schur
matrix. If $\ell\ne d$, all jets $Q(\partial)F(0)$ with $Q\in\auO_{\C_\ell}[z]$ are integral
(Proposition~\ref{pk:K12} for $\ell\nmid2q_1M$; Corollary~\ref{pf:K1}(a) for all $\ell\ne d$).
\item[(K2)] If $\ell=p\in\mathcal N$, then $F$ is one of the series of
Proposition~\ref{pa:product}, for $\mathsf C=[n_{j+1}]\times[n_j]$, the resolution of the identity
$u_{k,i}=(\zeta_{q_1}^k/q_1)e_i$ $(k<q_1,\ i\in\mathsf C)$ with multiplicities $q_1$, the lattice
$\Lambda=\Omega_0\auO_K$ and $c=\pi_M\varpi_{\mathrm{vac}}/\mathsf u$.
\item[(K1$_N$)] Let $\ell=N$ and $\iota\colon\C\to\C_N$. Let $z\in\operatorname{Hom}(X_j,X_{j+1})(\auA_{E^+,f})$
be integral at the primes above $N$ that are conjugate to the primes induced by the $\iota\circ\tau$,
and such that $Nz$ is integral at the primes induced by the $\iota\circ\tau$. Then for all $g,g'$ the
matrix of $g'^{-1}_N\circ\theta_j(\phi^\circ_{j,z};g,g')\circ g_N$ in the bases $(e_T)$ has entries in
$\auO_\iota$. (This is Lemma~\ref{pl:kernelN}.)
\end{enumerate}
\end{lemma}
\begin{proof} (K0) is Proposition~\ref{pk:K0}; the form of the matrix in (K1), and (K2), are
Proposition~\ref{pk:K12}; the integrality in (K1) at every prime $\ell\ne d$ and (K1$_N$) are
Corollary~\ref{pf:K1}.\end{proof}


\begin{lemma}[Gauge of a chain]\label{pa:chain}
Assume Lemma~\ref{pa:kernel}. Let $\ell\ne N$, $\ell\nmid\auNm Q_\vartheta$, and let $r_j\in\R$
$(0\le j<m)$ be such that all the jets $S_{\lambda_{j+1}}(\partial)_{T',T}F_{x,j,g,g'}(0)$ of (K1)
have valuation at least $r_j$. Then for every $\phi\in\auS$ integral at the given place,
$B_{\vartheta,\Psi}\phi$ has gauge at most $\ell^{-\sum_jr_j}$.
\end{lemma}
\begin{proof}
That $\phi(g_0)$ has gauge at most $1$ means that $g_{0,\ell}^{-1}\phi(g_0)$ has integral
coordinates. Insert $g_{i,\ell}\,g_{i,\ell}^{-1}$ between consecutive kernels in (K0): by (K1) and
the hypothesis the $j$-th factor becomes a matrix with entries of valuation $\ge r_j$. The
coefficients $\Psi'(x)$ are integral at $\ell$ (Proposition~\ref{au:expansion}, as $\ell\nmid
N\auNm Q_\vartheta$), and the valuation is ultrametric; so $g_{m,\ell}^{-1}(C_{\vartheta,\Psi'}
\phi)(g_m)$ has coordinates of valuation $\ge\sum_jr_j$. The restriction is the composition of the
evaluation at $\operatorname{diag}(y,1)$, of the branching projection, which is
$\auGL_{n_m-1}$-equivariant and has an integral matrix in the standard bases
(Proposition~\ref{au:restrict}), and of the multiplication by $\prod_\tau\tau(\aubar{\pi_N})^{-
n_\tau}$, an $\ell$-adic unit because $\ell\ne N$.
\end{proof}

\begin{lemma}[One lift gains its share]\label{pa:gain}
Let $p\in\mathcal N$ and $0\le j<m$. Call $(\tau,e)$ a \emph{collision} if $e\in I_p$ and both
$\tau$ and $u(-p^e)\tau$ are $a$-places, and let $C_p=\sum_{(\tau,e)\text{ collision}}a/p^e$. Then
\[
\paPred_j(p)-f\,n_{j+1}^3\,(1+\log_pB_j)\ \ge\ \frac{\augain_p}m-C_p-\frac{f\,a\,(\log b)^2}{p\,m}.
\]
\end{lemma}
\begin{proof}
By Definition~\ref{def:au-hmain}, $\augain_p/m=v(\pi)\sum_{e\in I_p}\bigl(\sum_{\tau\in\auPa}
(a/p^e-\kappa(\tau'))^++k\sum_\tau(\kappa(\tau)/p^e-\kappa(\tau'))^+\bigr)$, $\tau'=u(-p^e)\tau$.
All these terms vanish if $p^e\ge x_1$, since then $a/p^e\le a/x_1\le\kappa$ and
$\kappa(\tau)/p^e\le b/x_1<a/x_1$. We compare with the terms of $\paPred_j(p)$, using
\eqref{pa:eq-diag} and writing $s=x^{60}$. Let $e\in I_p$.

(1) Let $\tau\in\auPa$. If $\tau'\notin\auPa$, then $(a_\tau/p^e-a_{\tau'}-a'_{\tau'})^+\ge
(a/p^e-\kappa(\tau')-s)^+\ge(a/p^e-\kappa(\tau'))^+-s$. If $\tau'\in\auPa$ (a collision), the left
side is $\ge0\ge(a/p^e-\kappa(\tau'))^+-a/p^e$.

(2) For every $\tau$, $(l_\tau/p^e-l_{\tau'})^+\ge(k\kappa(\tau)/p^e-k\kappa(\tau')-s)^+\ge
k(\kappa(\tau)/p^e-\kappa(\tau'))^+-s$.

(3) By \eqref{pa:eq-const}, $v(\pi)\sum_{e\in I_p}a'_\tau/p^e=\gamma a'_\tau\le\frac2p(b+s)$.

\noindent In (1) and (2) the loss $s$ is needed only for the at most $2\log x_1$ exponents with
$p^e<x_1$. If $p=d$ there is the further term $\tfrac12\sum_\tau\paded_0(\tau)$, with
$\tau'=u(-1)\tau$. The places $\tau$ and $u(-1)\tau$ are neither both $a$-places nor both
$b$-places (their classes modulo $q_0$ are $\pm1$ and $\mp1$); so
$\min(a_\tau,a_{\tau'}+a'_{\tau'})\le b+s$, $a'_\tau\le b+s$ and $\min(l_\tau,l_{\tau'})\le
ka/x_1+s$. Altogether the left side of the lemma is at least $\augain_p/m-C_p$ minus
\[
f\Bigl(4s\log x_1+\frac{2(b+s)}p+\mathbf 1_{p=d}\bigl(b+2s+\tfrac{ka}{x_1}\bigr)+n_m^3(1+\log_pB_j)
\Bigr).
\]
Here $p<x_1$, $m\le2\sqrt x$, $a=xb$, $B_j\le b^3$, $d\le\log\log b$ by (D2), and $s$, $n_m$,
$k$ are polynomial in $x$; so each term is at most a quarter of $a(\log b)^2/(pm)\ge
\sqrt x\,b\,(\log b)^2/(2p)$, by a comparison of sizes.
\end{proof}

\begin{lemma}[Collisions]\label{pa:coll}
Let $p\in\mathcal N$ and let $o_0$ be the order of $-p$ in $(\Z/q_0)^\times$. If $p=d$ or $o_0$ is
even, there are no collisions. Otherwise the collisions are the pairs $(\tau,e)$ with
$\tau\in\auPa$, $e$ odd and $o_0\mid e$, and $C_p\le2fa/p^{o_0}$. If $p$ is not exceptional, then
$m\,C_p\le\paloss_p$. In all cases $m\,C_p\le j_p:=\lceil2mfa/p\rceil$.
\end{lemma}
\begin{proof}
For $\tau\in\auPa$ the place $u(-p^e)\tau$ is an $a$-place if and only if $-p^e\equiv1\pmod{q_0}$.
If $p=d$: $d\equiv q_0/2-1$ and $16\mid q_0$ give $d^2\equiv1\pmod{q_0}$, so $-d^e$ is $-1$ or $-d$
modulo $q_0$, and neither is $1$. If $p$ is inert, $e$ is odd and $-p^e=(-p)^e$, which is $1$
modulo $q_0$ exactly when $o_0\mid e$; this needs $o_0$ odd. Then
$C_p\le f\,a\sum_{i\ge0}p^{-o_0(1+2i)}\le2fa/p^{o_0}\le2fa/p$, which gives the last claim. Let $p$
be inert and not exceptional, with collisions. Either $p\equiv\pm1\pmod m$, and then
$\paloss_p=mf(\log b)^2a/p\ge2mfa/p$; or $o_0>\mathfrak r_p$, and then $p^{o_0}\ge p\,m$ and
$m\,C_p\le2fa/p\le\paloss_p$.
\end{proof}

\subsubsection*{Exceptional primes}
Let $p\in\mathcal N$ be exceptional. Each $\tau\in\auPl$ induces a prime $\mathfrak l_\tau$ of $E$
above $p$, the centre of the valuation of $\paCp$ on $\tau(E)$. A \emph{full mark} is a set
$\Lambda$ of primes of $E$ above $p$ with $\Lambda\cap\aubar\Lambda=\emptyset$, each of the form
$\mathfrak l_\tau$ with $\tau\in\auPa$, and containing $\mathfrak l_\tau$ or
$\aubar{\mathfrak l_\tau}$ for every $\tau\in\auPa$.

\begin{lemma}\label{pa:mark}
A full mark exists.
\end{lemma}
\begin{proof}
It suffices that no prime $\mathfrak l$ of $E$ above $p$ equals its conjugate; then one chooses
one prime from each pair $\{\mathfrak l,\aubar{\mathfrak l}\}$ that contains some
$\mathfrak l_\tau$, $\tau\in\auPa$, taking a member of that form. The decomposition group of
$\mathfrak l$ in $\operatorname{Gal}(E/\Q)$ is generated by the Frobenius $\varphi$, which is not
trivial on $K$. If $\aubar{\mathfrak l}=\mathfrak l$, then $c=\varphi^i$ with $i$ odd, so
$c\varphi^i\in\operatorname{Gal}(E/K)$ is trivial; it is the class of $-p^i$, so
$p^i\equiv-1\pmod{q_1}$. Now $m\mid q_1$ ($m$ is a power of $2$ with $m\le2\sqrt x$, while $2^{s+4}\mid q_1$ and $2^s\ge x^{11}/(3l_0(s))$ by (S3) of Lemma~\ref{an:lem-sizes}), and $m$ is a power of $2$, so $(\Z/m)^\times$ is a $2$-group and raising to the
odd power $i$ is a bijection of it; as $(-1)^i=-1$, we get $p\equiv-1\pmod m$, which is excluded
for an exceptional prime.
\end{proof}

For $\mathfrak l\in\Lambda$ the group $\auU(X_0)$ at the place of $E^+$ below $\mathfrak l$ is
$\auGL_k(E_{\mathfrak l})$. Let $T_{\mathfrak l}$ be the Hecke operator on $\auS$ of the double coset of
$\operatorname{diag}(p1_{\mathfrak r_p},1_{k-\mathfrak r_p})$, $(T_{\mathfrak l}\varphi)(g)=\sum_h
\varphi(gh)$, in the identification by the component at $\mathfrak l$ (\S\ref{pl:ss-norm}, (N3)), and
\[
Y_{\mathfrak l}=p^{A_{\mathfrak l}}\,T_{\mathfrak l},\qquad
A_{\mathfrak l}=a\cdot\#\{\tau\in\auPa:\mathfrak l_\tau=\mathfrak l\},\qquad
Y_\Lambda=\prod_{\mathfrak l\in\Lambda}Y_{\mathfrak l}.
\]
Write $\auS_\iota$ for the module of source forms integral at the given place.

\begin{lemma}[Valuation on the image of the Hecke operators]\label{pa:hecke}
Let $p$ be exceptional, $\Lambda$ a full mark, $\vartheta$ a twist with $p\nmid\auNm Q_\vartheta$
and $\Psi$ a chain datum. For every $\phi\in\auS_\iota$ the form
$B_{\vartheta,\Psi}(Y_\Lambda^{\,m}\phi)$ has gauge at most $p^{\,2\paloss_p-\augain_p}$.
\end{lemma}
\noindent Proved in \S\ref{pl:ss-hecke}, from Lemma~\ref{pa:kernel} (K1), (K2).

\begin{lemma}[Elementary divisors]\label{pa:eldiv}
Let $p$ be exceptional, $\Lambda$ a full mark and $\mathfrak l\in\Lambda$. Then
$Y_{\mathfrak l}\auS_\iota\subseteq\auS_\iota$, and for every integer $j_0\ge0$ there is an integer
$i\le\bigl(12A_{\mathfrak l}/(mp)+j_0/(m^2f)\bigr)\auDX$ with
$\aucol\bigl(\auS_\iota,\,Y_{\mathfrak l}\auS_\iota+p^{j_0}\auS_\iota\bigr)\mid p^{\,i}$.
\end{lemma}
\noindent Proved in \S\ref{ph:ss-ded}, from Theorem~\ref{ph:thm}; the three inputs of that theorem
(Propositions~\ref{ph:root}, \ref{ph:count}, \ref{ph:av}) are proved in \S\ref{pe:sec}, the
normalisation used is justified in Remark~\ref{pl:N}, and weak approximation is
Corollary~\ref{pl:wacor}.

\subsubsection*{The level prime}
\begin{lemma}[The prime $N$]\label{pa:level}
Let $\ell=N$, and let $U_N$, $L_0$ and $\mathcal A'=\{\varphi\in\auS:U_N^{\,m/2}\varphi\text{ is
integral for }L_0\}$ be as in \S\ref{pl:ss-level}. Then {\rm(a)} and {\rm(b)} of Lemma~\ref{pl:level}
hold.
\begin{enumerate}
\item[(a)] for every twist $\vartheta$, every chain datum $\Psi$ and every
$\phi\in\mathcal A'\cap\auSram$, the form $B_{\vartheta,\Psi}\phi$ has gauge at most $N^{2mfb}$;
\item[(b)] $\aucol(\auS_\iota,\mathcal A'\cap\auS_\iota)\mid N^{\,i}$ for some
$i\le m\,f\,b\,\auDX$.
\end{enumerate}
\end{lemma}
\noindent Proved in \S\ref{pl:ss-level}, from Lemma~\ref{pl:kernelN} (which is
Corollary~\ref{pf:K1}(b)) and Lemma~\ref{ps:branching}(c) (which is Proposition~\ref{pk:branching}).

\subsubsection*{The theorem}

\begin{theorem}[Proposition~\ref{au:divis}]\label{pa:divis}
(The statements used in the proof, Lemma~\ref{pa:kernel} including (K1$_N$),
Lemma~\ref{ps:branching}(c) and Proposition~\ref{au:expansion}, are proved in \S\ref{pk:ss-kernel}
and Corollary~\ref{pf:K1}, in Proposition~\ref{pk:branching} and in \S\ref{ps:ss-glue}.) Let $\ell$ be a prime with a
homomorphism $\C\to\C_\ell$, and let $\auS_\iota$ be the module of source forms integral there.
\begin{enumerate}
\item[(i)] If $\ell\ne N$ and $\ell\notin\mathcal N$, then for every twist $\vartheta$ with
$\ell\nmid\auNm Q_\vartheta$, every chain datum $\Psi$ and every $\phi\in\auS_\iota$ the form
$B_{\vartheta,\Psi}\phi$ is integral.
\item[(ii)] If $\ell=p\in\mathcal N$ is not exceptional, then under the same conditions
$B_{\vartheta,\Psi}\phi$ has gauge at most $p^{\,2\paloss_p-\augain_p}$.
\item[(iii)] If $\ell=N$, let $\mathcal A'$ be as in Lemma~\ref{pl:level}; then $S_0=\mathcal A'\cap\auS_\iota$ has
$\aucol(\auS_\iota,S_0)\mid N^{\,i}$ with $i\le mfb\,\auDX$, and $B_{\vartheta,\Psi}\phi$ has gauge
at most $N^{2mfb}$ for $\phi\in S_0\cap\auSram$.
\item[(iv)] If $\ell=p\in\mathcal N$ is exceptional, let $\Lambda$ be a full mark and $j_p=\lceil2mfa/p\rceil$. Then
$S_0=Y_\Lambda^{\,m}\auS_\iota+p^{j_p}\auS_\iota$ is contained in $\auS_\iota$,
$\aucol(\auS_\iota,S_0)\mid p^{\,i}$ with $i\le\paloss_p\auDX$, and $B_{\vartheta,\Psi}\phi$ has
gauge at most $p^{\,2\paloss_p-\augain_p}$ for $\phi\in S_0$.
\end{enumerate}
Consequently, in each case, with $S_0=\auS_\iota$ in cases (i), (ii), the module
$S=S_0\cap\auSram$ and the number $i$ satisfy the first assertion of
Proposition~\ref{au:divis}: $S\subseteq\auS_\iota\cap\auSram$, $i\le\auerr_\ell\auDX$,
$\aucol(\auS_\iota\cap\auSram,S)\mid\ell^{\,i}$, and $B_{\vartheta,\Psi}\phi$ has gauge at most
$\ell^{\,\auerr_\ell-\augain_\ell}$ for $\phi\in S$.
\end{theorem}
\begin{proof}
(i) Here $\ell\ne d$, because $d\in\mathcal N$. By (K1) and Lemma~\ref{pa:chain} with $r_j=0$,
$B_{\vartheta,\Psi}\phi$ is integral. (And $\augain_\ell=\auerr_\ell=0$ by
Lemma~\ref{au:gains}(a).)

(ii) Let $x$, $j$, $g$, $g'$ be as in (K1). By (K2) and Proposition~\ref{pa:product} the series
$F_{x,j,g,g'}$ has the normal form. By Lemma~\ref{pa:schur} and Proposition~\ref{pa:count}, every
jet along an entry of the Schur matrix $S_{\lambda_{j+1}}$ has valuation at least
$\paPred_j(p)-fn_{j+1}^3(1+\log_pB_j)$, hence, by Lemma~\ref{pa:gain}, at least
$r_j:=\augain_p/m-C_p-\paloss_p/m$. By Lemma~\ref{pa:chain} the gauge of $B_{\vartheta,\Psi}\phi$
is at most $p^{-\sum r_j}=p^{\,mC_p+\paloss_p-\augain_p}$, for every $p\in\mathcal N$ and
$\phi\in\auS_\iota$; and $mC_p\le\paloss_p$ by Lemma~\ref{pa:coll}.

(iii) is Lemma~\ref{pl:level}; $\augain_N=0$ because $N\ge x_1$, and $i\le mfb\,\auDX\le\auerr_N\auDX$.

(iv) We first note a rule for colengths: if $\gamma$ is an endomorphism of a module $M$,
$\alpha$ a scalar, and $\aucol(M,\gamma M+\alpha M)\mid z$, $\aucol(M,M'+\alpha M)\mid z'$ for a
submodule $M'\subseteq M$, then $\aucol(M,\gamma M'+\alpha M)\mid zz'$. Indeed, for a family
$\varphi$ in $M$ there is $\mu$ with $\det\mu\mid z$ and $\mu\varphi=\gamma\psi+\alpha\varrho$, and
$\mu'$ with $\det\mu'\mid z'$ and $\mu'\psi=\nu+\alpha\varrho'$, where $\psi,\varrho,\varrho'$ are
families in $M$ and $\nu$ in $M'$; then $\mu'\mu\varphi=\gamma\nu+\alpha(\gamma\varrho'+\mu'
\varrho)$. Applying this $m\cdot\#\Lambda$ times with Lemma~\ref{pa:eldiv} for $j_0=j_p$ gives
$S_0\subseteq\auS_\iota$ and $\aucol(\auS_\iota,S_0)\mid p^{\,i}$ with
\[
i\le m\sum_{\mathfrak l\in\Lambda}\Bigl(\frac{12A_{\mathfrak l}}{mp}+\frac{j_p}{m^2f}\Bigr)\auDX
\le\Bigl(\frac{12af}p+\frac{j_p}m\Bigr)\auDX\le\Bigl(\frac{14af}p+\frac1m\Bigr)\auDX\le\paloss_p
\auDX,
\]
since $\sum_{\mathfrak l}A_{\mathfrak l}\le af$ and $\#\Lambda\le f$. For the gauge, let
$\phi=Y_\Lambda^{\,m}\phi_1+p^{j_p}\phi_2$ with $\phi_1,\phi_2\in\auS_\iota$. The first term is
treated by Lemma~\ref{pa:hecke}. For the second, the bound proved in (ii) for all $p\in\mathcal N$
gives gauge at most $p^{-j_p+mC_p+\paloss_p-\augain_p}\le p^{\,\paloss_p-\augain_p}$, by
Lemma~\ref{pa:coll}.

For the last statement: $i\le\auerr_\ell\auDX$ in each case, as $\auerr_p=2\paloss_p$ and
$\auerr_N=2mfb$. If a family lies in $\auS_\iota\cap\auSram$, the matrix $\mu$ given by
$\aucol(\auS_\iota,S_0)\mid\ell^{\,i}$ carries it into $S_0$ and, $\auSram$ being a subspace, into
$S_0\cap\auSram$.
\end{proof}

\begin{remark}[Use in Theorem~\ref{thm:srcA}]
The deduction of Theorem~\ref{thm:srcA} in \S\ref{ss:au-proofAB} uses Proposition~\ref{au:divis} only
through input (I2), which is the last statement of Theorem~\ref{pa:divis}; so it goes through
unchanged. Compared with the items of Proposition~\ref{au:divis}: in (iv) the module is
$(Y_\Lambda^{\,m}\auS_\iota+p^{j_p}\auS_\iota)\cap\auSram$ rather than a sum of submodules of
$\auS_\iota\cap\auSram$ (the operators $Y_{\mathfrak l}$ are not claimed to preserve $\auSram$); in
(i) and (ii) nothing is used of the hypothesis that $\phi$ is ramified. What Theorem~\ref{thm:srcA} then depends on, for this input, is exactly: Lemma~\ref{pa:kernel}
including (K1$_N$), Lemma~\ref{ps:branching}(c), and the properties (W4), (W5) of the Weil
representation in the form (N5) of \S\ref{pl:ss-norm}. Of these, Lemma~\ref{ps:branching}(c) is
Proposition~\ref{pk:branching}, and Lemma~\ref{pa:kernel} is proved in \S\ref{pk:ss-kernel} and in
Corollary~\ref{pf:K1} (the integrality of the jets in (K1) at the primes $\ell\mid2q_1M$, and
(K1$_N$)); nothing remains open for this input.
\end{remark}


%% ===== proofs/A-seesaw-size.tex =====
% proofs/A-seesaw-size.tex
% Proofs of au:seesaw, au:kernel, au:restrict, au:expansion, au:equiv and of the constant in
% au:shortbasis. Labels of this file start with ps:. Only amsmath/amssymb/amsthm are used.
\providecommand{\auO}{\mathcal O}
\providecommand{\auU}{\mathrm U}
\providecommand{\auGL}{\mathrm{GL}}
\providecommand{\auHom}{\operatorname{Hom}}
\providecommand{\auA}{\mathbb A}
\providecommand{\auF}{\mathbb F}
\providecommand{\auPl}{\Phi}
\providecommand{\auPa}{\Phi_a}
\providecommand{\auS}{\mathcal S}
\providecommand{\auSram}{\mathcal S^{\mathrm{ram}}}
\providecommand{\auAz}{\mathcal A_0}
\providecommand{\auG}{\overline G}
\providecommand{\auP}{\mathcal P}
\providecommand{\auNm}{\operatorname{Nm}}
\providecommand{\auAut}{\operatorname{Aut}}
\providecommand{\aubar}[1]{\overline{#1}}
\providecommand{\ausinf}{s_\infty}
\providecommand{\aurd}{\mathrm{rd}_E}
\providecommand{\psM}{\mathsf M}
\providecommand{\psHS}{\mathscr H}
\providecommand{\psW}{\mathcal W}
\providecommand{\psZ}{Z}
\providecommand{\psvac}{\mathbf v}
\providecommand{\psrho}{\varrho}
\providecommand{\psE}{\mathcal E}
\providecommand{\psGlue}{\operatorname{Glue}}
\providecommand{\psfa}{\mathfrak a}
\providecommand{\psnot}[1]{\par\smallskip{\raggedright\noindent\textsc{Not proved here.} [#1]\par}\smallskip}
\renewcommand{\psHS}{\mathcal H}

\subsection{Proofs for the upper bound: the seesaw identity, sizes, restriction}\label{ps:sec}

This subsection proves Lemma~\ref{au:shortbasis} with its constant and
Propositions~\ref{au:equiv}, \ref{au:restrict}, \ref{au:kernel}, \ref{au:expansion}
and~\ref{au:seesaw}. To do so it must first make precise several objects that
\S\ref{ss:au-spaces} only names: the model of the Weil representation, the vacuum and the theta
kernels with their determinant twists (\S\ref{ps:ss-weil}), the lattices of the spaces $Y_j$
(\S\ref{ps:ss-Y}), the gluing and the chain data (\S\ref{ps:ss-glue}). Three statements are stated
here and proved in \S\ref{pk:sec}: the integrality of the branching map
(Lemma~\ref{ps:branching}(c); Proposition~\ref{pk:branching}), the identification of the definite
theta lift with the lift by jets (Lemma~\ref{ps:definite}; Proposition~\ref{pk:definite}), and the
law of Galois conjugation of jets (Lemma~\ref{ps:conj}), which is deduced in
Proposition~\ref{pk:conj} from a reciprocity law on one elliptic curve
(Proposition~\ref{pk:curve}(C3)(ii)), which is Theorem~\ref{pf:recip}. The
standard properties of the Weil representation are collected in Fact~\ref{ps:weil} with references.

\subsubsection{The Weil representation, the Fock model, the vacuum and the kernels}\label{ps:ss-weil}

\emph{Heisenberg representation.} Put $\jmath=-1/\sqrt{-d}$, so that
$\tau(\jmath)=i/\sqrt d$ for $\tau\in\auPl$, and let $\psi=\psi_\Q\circ
\operatorname{Tr}_{E/\Q}$ be the standard character of $\auA_E/E$, with
$\psi_{\Q,\infty}(t)=e^{-2\boldsymbol\pi it}$; fix a finite idele $\alpha$ such that $\widehat\auO_E$ is its
own dual for $(z,u)\mapsto\psi(\jmath\alpha\aubar\alpha\,\aubar zu)$, and a unitary Hecke
character $\upsilon$ of $E$ whose restriction to $\auA_{E^+}^\times$ is the quadratic character of
$E/E^+$; $\upsilon$ is unramified at every finite place ($\upsilon(x)=\prod_\tau x_\tau/|x_\tau|$
when the finite part of $x$ is a unit; such a character exists by Lemma~\ref{ps:upsilon} below, and
we fix one), and $\alpha$ is $1/P'(\zeta_{q_1})$ at the places dividing
$(1+\sqrt{-d})/2$ and $1$ elsewhere ($P$ the minimal polynomial of $\zeta_{q_1}$ over $K$, as in
\S\ref{pk:ss-conv}). We also fix the finite idele $a$ with $a_w=\alpha_w$ if $\alpha_w$ is not a unit
and $a_w=1$ otherwise; so $a\widehat\auO_E=\alpha\widehat\auO_E$, $a/\alpha$ is a unit idele, and
$a_w\ne1$ only at the primes $w\mid\omega$ above $q_1$ (Proposition~\ref{pgv:alpha}). The
coordinates $\tilde x=a^{-1}x$ in which the data at split places are written are described in
Lemma~\ref{pgx:dict}. All hermitian forms $\langle\,,\rangle$ are conjugate-linear in the first variable. For
a hermitian space $W$ over $E$ put $\beta(x,y)=\psi(\tfrac12\jmath\langle y,x\rangle)$ on
$W(\auA)=W\otimes_{E^+}\auA_{E^+}$ and $[x,y]=\beta(x,y)\beta(y,x)^{-1}$. Let $(\psrho,\psHS(W))$
be the smooth irreducible representation of the Heisenberg group of $W(\auA)$ with
$\psrho(x)\psrho(y)=\beta(x,y)\psrho(x+y)$; it is the restricted tensor product
$\psHS_\infty(W)\otimes\psHS_f(W)$, and $\omega=\omega_W$ denotes the Weil representation of
$\auU(W)(\auA)$ on it defined with $\upsilon$ (Weil~\cite{WeilActa}, Kudla~\cite{KudlaSplitting}),
normalised at every place as in Definition~\ref{pwf:def} (at the places of $E^+$ split in $E$, which
the cited works do not treat, by (D3) there).
For hermitian spaces $V,V'$ the space $\auHom_E(V,V')$ carries the form
$\langle w_1,w_2\rangle=\operatorname{tr}(w_1^\dagger w_2)$, where $w^\dagger\colon V'\to V$ is the
adjoint, and $\auU(V)\times\auU(V')$ acts on it by $(g,h)w=hwg^{-1}$. For a character $\mu$ of
$\auU(1)(E^+)\backslash\auU(1)(\auA_{E^+})$ we write $\omega_\mu(g,h)=\mu(\det h)\mu(\det g)^{-1}
\upsilon_{n}(\det h)\upsilon_{n'}(\det g)^{-1}\,\omega_{\auHom(V,V')}(g,h)$, where $n,n'$ are the
ranks of $V,V'$ and $\upsilon_k=(\upsilon|_{\auU(1)})^{\lfloor k/2\rfloor}$. (These are the
\emph{determinant twists}.)

\begin{lemma}[The splitting character]\label{ps:upsilon}
Let $E^+$ be the maximal real subfield of $E$ and $\epsilon$ the quadratic character of
$\auA_{E^+}^\times/E^{+\times}$ attached to $E/E^+$.
\begin{enumerate}
\item[(a)] $E/E^+$ is unramified at every finite place, the group of roots of unity of $E$ is
$\mu_{q_1}$, and it is generated by $-\zeta_{q_1}=(1-\zeta_{q_1})/(1-\aubar{\zeta_{q_1}})$, where
$1-\zeta_{q_1}$ is a unit of $E$ with $\prod_{\tau\in\auPl}\tau(1-\zeta_{q_1})=1$.
\item[(b)] There is a continuous character $\upsilon\colon\auA_E^\times/E^\times\to\{z\in\C:|z|=1\}$
such that $\upsilon|_{\auA_{E^+}^\times}=\epsilon$ and $\upsilon(x)=\prod_{\tau\in\auPl}x_\tau/|x_\tau|$
for every idele $x$ whose finite part lies in $\widehat\auO_E^\times$; here $x_\tau\in\C^\times$ is the
component of $x$ at the place $\tau$, read through the embedding of $E$ which induces $\tau$ and
extends the fixed embedding of $K$.
\item[(c)] For such a character: $\upsilon_w$ is unramified at every finite place $w$ of $E$;
$\upsilon_\tau(z)=z/|z|$ for $\tau\in\auPl$, so that $\upsilon_\tau(z)=z$ for $|z|=1$;
$\upsilon(x\aubar x)=1$ for $x\in\auA_E^\times$; if a place $v$ of $E^+$ splits into $w,\aubar w$,
then $\upsilon_{\aubar w}=\upsilon_w^{-1}$ on $E^{+\times}_v$, and $\upsilon(t,t^{-1})=\upsilon_w(t)^2$
for the element $(t,t^{-1})$ of $\auU(1)(E^+_v)\subset E_w^\times\times E_{\aubar w}^\times$; and
$\upsilon|\cdot|^{-1/2}$ is an algebraic Hecke character, with component $z\mapsto\aubar z^{\,-1}$ at
every $\tau\in\auPl$.
\end{enumerate}
\end{lemma}
\begin{proof}
(a) The field $E$ contains $i$ (as $4\mid q_1$) and $\sqrt{-d}$, and neither lies in $E^+$; so
$E=E^+(i)=E^+(\sqrt{-d})$. The polynomial $X^2+1$ has discriminant $-4$, so $E/E^+$ is unramified
outside $2$; and $\frac12(1+\sqrt{-d})$ is a root of $X^2-X+\frac14(1+d)$, which has integral
coefficients and odd discriminant $-d$ because $d\equiv7\pmod 8$, so $E/E^+$ is unramified above $2$.
Let $\mu_n$ be the group of roots of unity of $E$; then $q_1\mid n$ and $\Q(\mu_n)\subseteq E$. Let $p$
be a prime, $p^c$ and $p^e$ the exact powers of $p$ dividing $n$ and $q_1$, and suppose $c>e$. The
ramification index of $p$ in $E$ is at least $\varphi(p^c)$. If $p\ne d$, then $p$ is unramified in $K$
and its ramification index in $E=K\Q(\zeta_{q_1})$ is $\varphi(p^e)$; but $\varphi(p^c)>\varphi(p^e)$,
because $e\ge1$ if $p=2$. If $p=d$, then $e=0$, $d$ is unramified in $\Q(\zeta_{q_1})$ and its
ramification index in $E$ is $2$; but $\varphi(d^c)\ge d-1\ge6$. So $n=q_1$. Write $\zeta=\zeta_{q_1}$.
We have $-\zeta=\zeta^{1+q_1/2}$, and $1+q_1/2$ is prime to $q_1$ because $q_1/2$ is even and is
divisible by every prime factor of $q_1$; so $-\zeta$ generates $\mu_{q_1}$. Clearly
$(1-\zeta)/(1-\zeta^{-1})=-\zeta$. The places $\tau\in\auPl$ correspond to the embeddings of $E$
extending that of $K$, that is, to $\operatorname{Gal}(E/K)$, which restricts isomorphically to
$\operatorname{Gal}(\Q(\zeta)/\Q)$; so $\prod_\tau\tau(1-\zeta)=\prod_{j\in(\Z/q_1)^\times}(1-\zeta^j)=
\Phi^{\mathrm{cyc}}_{q_1}(1)$, the value at $1$ of the $q_1$-th cyclotomic polynomial, which is $1$ because $q_1$ is divisible by $2$
and by $3$ and so is not a power of a prime. In particular $1-\zeta$ is a unit.

(b) Let $\mathsf U=(E\otimes\R)^\times\times\widehat\auO_E^\times$, an open subgroup of $\auA_E^\times$,
and $\upsilon_{\mathsf U}(x)=\prod_\tau x_\tau/|x_\tau|$ for $x\in\mathsf U$. By (a), $\epsilon$ is
unramified at every finite place, so it is trivial on $\widehat\auO_{E^+}^\times$; at the real place of
$E^+$ below $\tau$ it is the sign. As $\mathsf U\cap\auA_{E^+}^\times=(E^+\otimes\R)^\times\times
\widehat\auO_{E^+}^\times$ and $x_\tau/|x_\tau|$ is the sign of $x_\tau$ for real $x_\tau$, the
characters $\upsilon_{\mathsf U}$ and $\epsilon$ agree on $\mathsf U\cap\auA_{E^+}^\times$. Hence there
is a unique character $\chi$ of the open subgroup $\mathsf U\cdot\auA_{E^+}^\times$ with
$\chi(ua)=\upsilon_{\mathsf U}(u)\epsilon(a)$; it is continuous, being continuous on the open subgroup
$\mathsf U$.

We claim that $\chi(e)=1$ for every $e\in E^\times\cap\mathsf U\auA_{E^+}^\times$. Write $e=ua$. Then
$e/\aubar e=u/\aubar u$ is an element of $E^\times$ which is a unit at every finite place and has
absolute value $1$ at every infinite place; so it is a root of unity (Kronecker), and by (a)
$e/\aubar e=(-\zeta)^j$ for an integer $j$. Then $e^+=e\,(1-\zeta)^{-j}$ satisfies $\aubar{e^+}=e^+$,
that is, $e^+\in E^{+\times}$. As $(1-\zeta)^j\in\auO_E^\times\subset\mathsf U$ and
$e^+\in\auA_{E^+}^\times$,
\[
\chi(e)=\upsilon_{\mathsf U}\bigl((1-\zeta)^j\bigr)\,\epsilon(e^+)=\Bigl(\prod_\tau\frac{\tau(1-\zeta)}
{|\tau(1-\zeta)|}\Bigr)^{j}=1,
\]
because $\epsilon$ is trivial on $E^{+\times}$ and $\prod_\tau\tau(1-\zeta)=1$ by (a). This proves
the claim.

Consequently $\chi_1(eua)=\chi(ua)$ ($e\in E^\times$, $u\in\mathsf U$, $a\in\auA_{E^+}^\times$) is a
well-defined character of the open subgroup $G_1=E^\times\mathsf U\auA_{E^+}^\times$ of
$\auA_E^\times$, trivial on $E^\times$ and continuous. The quotient $\auA_E^\times/G_1$ is a quotient of
the ideal class group $\auA_E^\times/E^\times\mathsf U$ of $E$, hence finite. A character of a
subgroup of an abelian group with values in the divisible group $\{|z|=1\}$ extends to the group; an
extension $\upsilon$ of $\chi_1$ to $\auA_E^\times$ is continuous because $G_1$ is open, and it has the
properties stated.

(c) The first two assertions are the formula for $\upsilon$ on $\mathsf U$. Next,
$x\aubar x\in\auA_{E^+}^\times$ is a norm from $\auA_E^\times$, and $\epsilon$ is trivial on norms. If $v$
splits into $w,\aubar w$, then $t\in E^{+\times}_v$ is the idele of $E$ with components $t$ at $w$ and at
$\aubar w$, and $\epsilon_v=1$; so $\upsilon_w(t)\upsilon_{\aubar w}(t)=1$, and $\upsilon(t,t^{-1})=
\upsilon_w(t)\upsilon_{\aubar w}(t)^{-1}=\upsilon_w(t)^2$. Finally the idele norm of $E$ is $|z|^2$ at a
complex place, so the component of $\upsilon|\cdot|^{-1/2}$ at $\tau$ is $z|z|^{-2}=\aubar z^{\,-1}$.
\end{proof}

The character $\upsilon$ is determined by the conditions of Lemma~\ref{ps:upsilon}(b) up to a
character of the finite group $\auA_E^\times/E^\times\mathsf U\auA_{E^+}^\times$, where
$\mathsf U=(E\otimes\R)^\times\times\widehat\auO_E^\times$; any choice will do,
and it is fixed once and for all.

\begin{lemma}[Standard facts]\label{ps:weil}
Let $W$ be a hermitian space over $E$.
\begin{enumerate}
\item[(W1)] $\omega(g)\psrho(w)\omega(g)^{-1}=\psrho(gw)$. There is a functional $\theta=\theta_W$
on $\psHS(W)$, unique up to a scalar, with $\theta\circ\psrho(w)=\theta$ for $w\in W(E)$; it
satisfies $\theta\circ\omega(\gamma)=\theta$ for $\gamma\in\auU(W)(E^+)$.
\item[(W2)] Let $I\subset W$ be an isotropic $E$-subspace, $W_I=I^\perp/I$, and $P_I$ the stabiliser
of $I$ in $\auU(W)$. There is a continuous linear map $\psZ_I\colon\psHS(W)\to\psHS(W_I)$ (the
\emph{zero term}) with $\theta_{W_I}(\psZ_IF)=\int_{I(E)\backslash I(\auA)}\theta(\psrho(i)F)\,di$,
$\psZ_I\psrho(w)=\psrho(\bar w)\psZ_I$ for $w\in I^\perp(\auA)$ with image $\bar w$, and
$\psZ_I\omega(p)=\chi_I(p)\,\omega_{W_I}(\bar p)\psZ_I$ for $p\in P_I(\auA)$, where $\chi_I(p)=
(\upsilon|\cdot|^{-1/2})(\det p|_I)^{-1}=\upsilon(\det p|_I)^{-1}|\det p|_I|^{1/2}$, with
$\det p|_I\in\auA_E^\times$ the determinant on the $\auA_E$-module $I(\auA)$ and $|\cdot|$ the
idele norm of $E$; so $\chi_I=1$ on unipotent elements and on $P_I(E^+)$.
For isotropic $I\subset I'$ one has $\psZ_{I'}=\psZ_{I'/I}\psZ_I$, and $\psZ_I=\psZ_{I,\infty}
\otimes\psZ_{I,f}$.
\item[(W3)] For $F\in\psHS(W)$, $\theta(F)=\sum_{w\in W(E)/I^\perp(E)}\theta_{W_I}(\psZ_I\psrho(w)F)$,
the series converging absolutely. If $W=\auHom(V,V')$, $F=\Psi\otimes f$ with $\Psi$ smooth, and
$\mathcal C$ is a compact set of finite adelic points of $\auU(V)\times\auU(V')$, there are $C,A$ with
$\sum_w|\theta_{W_I}(\psZ_I\psrho(w)\omega(g,h)F)|\le C(\|g_\infty\|\|h_\infty\|)^A$ whenever
$(g_f,h_f)\in\mathcal C$.
\item[(W4)] Let $L\subset W(\auA_f)$ be a lattice that is its own dual for $[\,,]$. The vectors of
$\psHS_f(W)$ on which $\psrho(l)$, $l\in L$, acts by a fixed character form a line. If $\psvac_L$
spans the line of the character $\epsilon_L$ with $\epsilon_L(l)\epsilon_L(l')=\beta(l,l')
\epsilon_L(l+l')$ fixed below, and $\psrho^\circ(x)=\epsilon(x)\psrho(x)$ for a fixed function
$\epsilon$ with values in roots of unity extending $\epsilon_L^{-1}$, then
$\psrho^\circ(l)\psvac_L=\psvac_L$, $\psrho^\circ(x+l)\psvac_L=\zeta\psrho^\circ(x)\psvac_L$ with a
root of unity $\zeta$, the $\psrho^\circ(x)\psvac_L$, $x\in W(\auA_f)/L$, form a basis of
$\psHS_f(W)$, and $\omega(g)\psvac_L\in\C^\times\psvac_{gL}$ for $g\in\auU(W)(\auA_f)$. If
$g\in\auU(W)(\auA_f)$ stabilises $L$ then $\omega(g)\psvac_L=\psvac_L$.
\item[(W5)] At a place $v=w\aubar w$ of $E^+$ split in $E$, $\psHS_v(\auHom(V,V'))$ is the space
of Schwartz--Bruhat functions on $\auHom(V_w,V'_w)$, with $(\omega_\mu(g,h)\Psi)(x)=\mu_w(\det h_w)
\mu_w(\det g_w)^{-1}\upsilon'_w(g_w,h_w)\,\Psi(h_w^{-1}xg_w)$, where, for even ranks $n,n'$, $\mu_w\upsilon'_w$ is given by \eqref{pl:eq-w5}: the factor is
$\chi_w(\det h_w)|\det h_w|_w^{-n/2}\chi_w(\det g_w)^{-1}|\det g_w|_w^{n'/2}$ with $\chi$ the base
change of $\mu$; for arbitrary ranks the untwisted operator is
\[\Psi\mapsto\upsilon_w(\det h_w)^{-n}\upsilon_w(\det g_w)^{n'}|\det h_w|_w^{-n/2}|\det g_w|_w^{n'/2}\Psi(h_w^{-1}xg_w)\]
(Proposition~\ref{pwf:W5}), so that all factors multiplying $\Psi(h_w^{-1}xg_w)$ are built from $\mu_w$, $\upsilon_w$ and $|\cdot|_w$ evaluated at $\det g_w,\det h_w$; $\psrho(a_w,0)\Psi(x)=\Psi(x+a_w)$
and $\psrho(0,a_{\bar w})$ is multiplication by the character $x\mapsto[x,a_{\bar w}]$.
\item[(W0), (W6)] Two further facts are used; they are stated in full in \S\ref{pk:ss-conv}.
(W0): $\theta$ is defined on the smooth vectors of the Fock space tensored with $\psHS_f(W)$ and is
continuous in the first factor. (W6): for hermitian spaces $V,V'$ of ranks $n,n'$ the map
$w\mapsto w^\dagger$ identifies $\psHS(\auHom(V,V'))$ with $\psHS(\auHom(V',V^-))$, carrying
$\omega_{\auHom(V,V')}(g,g')$ to $\upsilon(\det g)^{n'}\upsilon(\det g')^{-n}\,
\omega_{\auHom(V',V^-)}(g',g)$. It carries the self-dual lattice $\alpha\auHom(\psM,\psM')$ to
$\aubar\alpha\auHom(\psM',\psM)$, which differs from $\alpha\auHom(\psM',\psM)$ at the places above $q_1$
(Remark~\ref{pf:rem-flip-lattice}).
\end{enumerate}
\end{lemma}
\begin{proof}[References]
(W1): the commutation relation is part of the definition (Definition~\ref{pwf:def}, with
\cite[Th\'eor\`eme~1 and no.~34]{WeilActa}); existence and uniqueness of $\theta$ are
Lemma~\ref{pwf:unique}, and its invariance under $\auU(W)(E^+)$ for the splitting by $\upsilon$ is
Theorem~\ref{pwf:inv} (by Corollary~\ref{pwf:rk}, $\omega$ agrees on the rational points with Weil's
lift $r_k$, for which the invariance is \cite[Th\'eor\`eme~6]{WeilActa}). (W2), (W3) are the
properties of the mixed model attached to $W=I\oplus I^*\oplus W_I$: $\psHS(W)=\auS(I^*(\auA))
\otimes\psHS(W_I)$, $\psZ_I$ is evaluation at $0$, $\theta(F)=\sum_{\xi\in I^*(E)}\theta_{W_I}(F(\xi))$;
they are proved in Propositions~\ref{pwf:W2} and~\ref{pwf:W3} (Kudla's explicit formulas \cite[\S5]{KudlaSplitting} are
for a Lagrangian $I$); the absolute
convergence and the bound in (W3) are Lemma~\ref{pz:absgrowth} (which uses (W0), (W1), (W2), (W4) and
Lemma~\ref{pz:growth}, not (W3)). (Weil~\cite[Th\'eor\`eme~6]{WeilActa} gives the continuity of the theta
distribution and its invariance under the rational points, used in (W1), but no bound of this kind.)
(W4) is the lattice model; for the last assertion, the vector
$\omega(g)\psvac_L$ lies on the line of $\psvac_L$ and the scalar is $1$ for the splitting used
(it is $\upsilon(\det u)|\det u|^{-1/2}$ with $uL=L\cap g^{-1}L$). (W5) is Proposition~\ref{pwf:W5}
(the cited works treat only places where $E^+_v\otimes E$ is a field); see \S\ref{pl:ss-norm}, (N5). For (W0) see \S\ref{pk:ss-conv}; (W6) is proved in
Proposition~\ref{pf:flip} from (W1), (W4) and the local form of (W2).
\end{proof}

\emph{The Fock model and the theta function.} Let $X,X'$ be positive definite of ranks $n<n'$ and
$W=\auHom(X,X')$, so that $W\otimes_{E^+,\tau}\R=M_{n'\times n}(\C)$ for $\tau\in\auPl$. We use
the coordinates $x_\tau=c_\psi\cdot\tau(w)\in M_{n'\times n}(\C)$ with
$c_\psi=\sqrt{2\boldsymbol\pi/\sqrt d}$, in which $\beta_\infty(x,y)=\exp(+i\,\Iml\sum_\tau
\operatorname{tr}(y_\tau^*x_\tau))$; the variable $z$ of the Fock space is the complex conjugate of
$x$. Then $\psHS_\infty(W)$ is the Fock space of entire functions $F$ on
$\prod_\tau M_{n'\times n}(\C)$ with $\int|F|^2e^{-|z|^2}<\infty$, with
\[
(\psrho_\infty(y)F)(z)=e^{z\cdot y-|y|^2/2}F(z-\aubar y),\qquad z\cdot y=\sum_\tau
\operatorname{tr}(y_\tau^{\,t}z_\tau),
\]
the \emph{real vacuum} is $F=1$, and $\auU(X)(\R)\times\auU(X')(\R)$ acts by
$F\mapsto F(h^{t}z\,\aubar g)$ times the determinant twists
$\prod_\tau\det(h_\tau)^{e_{j+1,\tau}}\det(g_\tau)^{-(e_{j,\tau}-c_\tau(n_j))}$. For $f\in\psHS_f(W)$ the \emph{theta function} of $f$ is the entire
function
\begin{equation}\label{ps:eq-theta}
\Theta_f(z)=\theta\bigl(e^{\langle\,\cdot\,,\bar z\rangle}\otimes f\bigr)=e^{|z|^2/2}\,\theta\bigl(
\psrho_\infty(z)(1\otimes f)\bigr)\qquad(z\in\textstyle\prod_\tau M_{n'\times n}(\C)),
\end{equation}
and for a polynomial $Q$ in the entries of $z$ and $c\in\C^\times$ the \emph{jet of $f$ along $Q$
at scale $c$} is $Q(\partial)[\Theta_f(c\,\cdot)](0)$. By \eqref{ps:eq-theta} and the reproducing
property of the Fock space, with $d\mu(y)=\pi^{-fnn'}dy$,
\begin{equation}\label{ps:eq-bargmann}
Q(\partial)[\Theta_f(c\,\cdot)](0)=\theta\bigl(Q(c\,\cdot)\otimes f\bigr)=\int Q(c\,\bar y)\,
e^{-|y|^2/2}\,\theta\bigl(\psrho_\infty(y)(1\otimes f)\bigr)\,d\mu(y);
\end{equation}
indeed $\partial^\alpha H(0)=\int\bar y^\alpha H(y)e^{-|y|^2}d\mu(y)$ for entire $H$ with
$H(y)e^{-|y|^2/2}$ bounded, and the theta values are bounded.

\emph{The vacuum and its normalisation.} The lattice of a standard space $X$ is
$\widehat\auO_E^{\,n}$, and $L=\alpha\auHom(\widehat\auO_E^{\,n},\widehat\auO_E^{\,n'})$ is
self-dual in $W(\auA_f)$. The character $\epsilon_L$ is $l\mapsto\psi_f(-\tfrac12\omega\,\jmath\langle l,l\rangle)$
with $\omega=(1+\sqrt{-d})/2$, an element of trace one (then $\epsilon_L(l)\epsilon_L(l')=\beta(l,l')
\epsilon_L(l+l')$, and $\epsilon_L=\pm1$: \S\ref{pf:ss-adm}, Lemma~\ref{pf:eps}), the \emph{vacuum} is the vector $\psvac=\psvac_L$ of (W4) normalised by
$\theta(1\otimes\psvac)$ having mean square $1$ over $W(E)\backslash W(\auA)$, and the
\emph{integral vacuum} is $\psvac^\circ=c_{\mathrm{vac}}^{\,nn'}\psvac$, where $c_{\mathrm{vac}}\psvac$, for the space of rank one, is the vector \eqref{pk:eq-vac}: a sum of
$q_1^{2f}$ translates, with coefficients roots of unity, of the vector $F_0$ whose theta function
is the product \eqref{pk:eq-product} of theta functions of the curve $\C/\Omega_0\auO_K$. (No
closed formula for $c_{\mathrm{vac}}$ is needed; \eqref{pk:eq-vac} defines it.) The translates $\phi^\circ_x=\psrho^\circ(x)\psvac^\circ$, $x\in W(\auA_f)$, are
the ``translates of the vacuum'' of \S\ref{ss:au-spaces}.

\emph{Weights and kernels.} At $\tau\in\auPl$ let $c_\tau(h)$ be the number of columns of height
$h$ fixed in \S\ref{ss:au-spaces}: $c_\tau(1)=a\mathbf 1_{\tau\in\auPa}$, $c_\tau(k)=\kappa_1(\tau)
-n_1/2$, $c_\tau(n_i)=\Delta-(t+1)$ for $i\ge1$, and $c_\tau(h)=0$ otherwise. The weight
$\lambda(X_j)$ has the columns of height $<n_j$ and the determinant exponent $e_{j,\tau}=0$ for
$j=0$ and $e_{j,\tau}=n_{j-1}/2-\kappa_j(\tau)$ for $j\ge1$; thus its rows at $\tau$ are
$\lambda(X_j)_{\tau,p}=e_{j,\tau}+\sum_{p\le h<n_j}c_\tau(h)$ ($1\le p\le n_j$), and
$e_{j,\tau}-c_\tau(n_j)=n_{j+1}/2-\kappa_{j+1}(\tau)=e_{j+1,\tau}+\tfrac12(n_{j+1}-n_j)$. For a matrix $z_\tau\colon\C^{n_j}\to\C^{n_{j+1}}$ let
$S_\tau(z_\tau)\colon W_{\lambda(X_j)_\tau}\to W_{\lambda(X_{j+1})_\tau}$ be the Schur functor of
the diagram with $c_\tau(h)$ columns of each height $h<n_{j+1}$, applied to $z_\tau$ (on $\C^{n_j}$
the columns of height $n_j$ are powers of the determinant), and $S(z)=\bigotimes_\tau S_\tau
(z_\tau)$; its matrix entries in the standard bases are polynomials in $z$ with integer
coefficients, homogeneous of degree $\sum_pl_{\tau,p}$ in $z_\tau$, where $l_{\tau,p}=\sum_{p\le h<
n_{j+1}}c_\tau(h)$, $1\le p\le n_j$, are the rows of the diagram of the $j$-th lift. With these
coordinates the rational points act on $W_\lambda$ through the conjugate embeddings,
$\gamma\cdot\xi=\rho_\lambda(\aubar{\tau(\gamma)})\xi$, and $g_\infty\cdot\xi=\rho_\lambda(
\aubar{g_\infty})\xi$; the expressions $g_\infty^{-1}\psi(g_f)$ below are to be read with this
action. Let
$s_0=(\varpi_{\mathrm{vac}}/\mathsf u)\,c_\psi/\Omega_0$ be the scale \eqref{pk:eq-s0};
$|s_0|^2=1/\Omega$. For $\Phi\in\psHS_f(\auHom(X_j,X_{j+1}))$
and finite adelic points $g,g'$ the \emph{kernel} is
\begin{equation}\label{ps:eq-kernel}
\theta_j(\Phi;g,g')=S(\partial)\bigl[\Theta_{\omega_{\mu_{j+1}}(g,g')\Phi}(\pi_Ms_0\,\cdot)\bigr](0)
\ \in\auHom(W_{\lambda(X_j)},W_{\lambda(X_{j+1})}),
\end{equation}
that is, the matrix of the jets at scale $\pi_Ms_0$ along the entries of $S$. The kernels have the
two properties (T1), (T2) of \S\ref{pl:ss-norm}, (N4): equivariance under the rational points,
$\theta_j(\Phi;\gamma g,\gamma'g')=\gamma'\circ\theta_j(\Phi;g,g')\circ\gamma^{-1}$, and under the
finite adelic points, $\theta_j(\omega_{\mu_{j+1}}(h,h')\Phi;g,g')=\theta_j(\Phi;gh,g'h')$. The lift is
$(\Theta^X_j(\Phi)\varphi)(g')=\int_{\Gamma\backslash G_f}\theta_j(\Phi;g,g')\varphi(g)\,dg$ for the
Haar measure giving volume $1$ to $J_{X_j}$; this is $[J_{X_j}:J]^{-1}\sum_{g\in\Gamma\backslash
G_f/J}\theta_j(\Phi;g,g')\varphi(g)$ for every free level $J\subseteq J_{X_j}$ fixing $\varphi$ and
$\omega_{\mu_{j+1}}(\cdot,1)\Phi$. This is the definition used in the rest of the paper; that it has
values in algebraic modular forms, and its relation to theta integrals, is the following.

\begin{lemma}[The definite lift]\label{ps:definite}
Let $X,X',X''$ be consecutive spaces $X_{j},X_{j+1},X_{j+2}$. For every smooth vector $u$ of the Fock
space of $\auHom(X',X^-\oplus X'')\otimes\R$ and every linear form $l$ on $W_{\lambda(X)}$ there is
a linear form $l''$ on $W_{\lambda(X'')}$ such that for all $\vartheta$, all
$\Phi\in\psHS_f(\auHom(X,X'))\otimes\psHS_f(\auHom(X',X''))$, all algebraic modular forms $\psi$ on
$\auU(X)$ of weight $\lambda(X)$ and all $k\in\auU(X'')(\auA)$,
\[
\int_{[\auU(X')]}\int_{[\auU(X)]}l\bigl(g_\infty^{-1}\psi(g_f)\bigr)\,\theta\bigl(\omega(g,r,k)(u\otimes
\Phi)\bigr)\,dg\,dr=l''\Bigl(k_\infty^{-1}\bigl(\Theta^X_{j+1}\Theta^X_j(\Phi)\psi\bigr)(k_f)\Bigr),
\]
where $[G]=G(E^+)\backslash G(\auA)$ has volume $1$, $\omega(g,r,k)=\omega_{\mu_{j+1}}(g,r)\otimes
\omega_{\mu_{j+2}}(r,k)$ on $\psHS(\auHom(X,X'))\otimes\psHS(\auHom(X',X''))=\psHS(\auHom(X',
X^-\oplus X''))$, and $\Theta^X_{j+1}\Theta^X_j(\Phi)$ is defined by linearity from pure tensors.
\end{lemma}
\begin{proof} See Proposition~\ref{pk:definite}. \end{proof}

\subsubsection{Short bases: proof of Lemma~\ref{au:shortbasis}}\label{ps:ss-short}

\begin{lemma}\label{ps:transfer}
Let $\Lambda$ be a lattice in a euclidean space of dimension $m\ge1$ whose dual lattice has no
non-zero vector shorter than $\rho>0$. Then every point of the space is within $m^2/\rho$ of
$\Lambda$, and the vectors of $\Lambda$ of length at most $3m^2/\rho$ span the space.
\end{lemma}
\begin{proof}
By Minkowski's convex body theorem applied to the cube $[-t,t]^m$, $t=(\det\Lambda)^{1/m}$, $\Lambda$ has a non-zero vector $v$ with $|v|\le\sqrt m\,
(\det\Lambda)^{1/m}$, and $\det\Lambda=(\det\Lambda^*)^{-1}\le(\sqrt m/\rho)^m$ by the same bound
for $\Lambda^*$; so $|v|\le m/\rho$. We prove the first claim by induction on $m$. For $m=1$,
$\Lambda=\Z v$, $1/|v|\ge\rho$, and the covering radius is $|v|/2$. For $m>1$ let $v$ be a
shortest vector and $\Lambda'$ the projection of $\Lambda$ to $v^\perp$; its dual is
$\Lambda^*\cap v^\perp$, so the hypothesis holds for $\Lambda'$, and the covering radius of
$\Lambda$ is at most that of $\Lambda'$ plus $|v|/2$, hence at most $(m-1)^2/\rho+m/(2\rho)\le
m^2/\rho$. For the second claim put $r=m^2/\rho$: if the vectors of length $\le3r$ lay in a
hyperplane $H$, a point at distance $2r$ from $H$ on the normal through $0$ would be within $r$ of
some $\lambda\in\Lambda$, with $|\lambda|\le3r$, so $\lambda\in H$, which is absurd.
\end{proof}

\begin{proof}[Proof of Lemma~\ref{au:shortbasis}, with $C_{\mathrm{sb}}=6c_0(fn)^2$]
Let $L=E^n\cap g\widehat\auO_E^{\,n}$, a lattice in $E^n$ that is unimodular for the standard
hermitian form $h$, and embed $E^n$ in $\C^{\auPl\times n}=\R^{2fn}$ by $y\mapsto(\tau(y_i))$, with
$\|y\|^2=\sum_\tau\tau(h(y,y))$; the real inner product is $(x,y)=\tfrac12\operatorname{Tr}_{E/\Q}
h(x,y)$. If $y\in L$ is not zero then $h(y,y)$ is a totally positive integer of $E^+$, so
$\|y\|^2\ge f\ge1$. Put $\Lambda=(2/c_0)L$. If $x$ lies in the dual of $\Lambda$, then
$\operatorname{Tr}_{E/\Q}h(\xi,x/c_0)\in\Z$ for all $\xi\in L$; since $L$ is unimodular this means
$x/c_0\in\mathfrak d_E^{-1}L$, so $x\in c_0\mathfrak d_E^{-1}L\subseteq L$ because $c_0$ lies in
the different $\mathfrak d_E$. Hence the dual of $\Lambda$ has no non-zero vector shorter than $1$,
and by Lemma~\ref{ps:transfer} with $m=2fn$ the vectors of $\Lambda$ of length at most $3(2fn)^2$
span $\R^{2fn}$. Multiplying by $c_0/2$: the vectors of $L$ of length at most $6c_0(fn)^2$ span
$\R^{2fn}$ over $\R$, hence $E^n$ over $E$, and $n$ of them form a basis; let $Y$ be the matrix of
its columns. Then $g^{-1}Y$ is integral, and every entry $y$ of $Y$ lies in $E$ and satisfies
$|\tau(y)|\le6c_0(fn)^2$ for $\tau\in\auPl$, hence at every embedding of $E$ in $\C$, hence
$|\sigma y|\le6c_0(fn)^2$ for every $\sigma\in\auAut(\C)$.
\end{proof}
The sharper transference theorem of Banaszczyk~\cite[Theorem~2.1]{Banaszczyk} would give a constant linear in
$fn$; the lemma is stated with the constant that the elementary argument gives.

\subsubsection{The spaces $Y_j$, their lattices and data; proof of Proposition~\ref{au:equiv}}
\label{ps:ss-Y}

Let $0\le j<m$ and $s_j=m-1-j$. The space $Y_j$ has a basis consisting of the standard basis of
$X_j$, with last vector $y=e_{\mathrm{last}}$, and of a \emph{string} $f_0,\dots,f_{s_j-1}$ orthogonal
to $X_j$ with $\langle f_a,f_b\rangle=0$ unless $a+b=s_j-1$, $\langle f_a,f_b\rangle=1$ if
$a+b=s_j-1$, $a\ne b$, and $\langle f_a,f_a\rangle=-N$ for the \emph{middle vector} $z=f_a$,
$2a=s_j-1$ (which exists when $j$ is even). The \emph{grade} is $0$ on $X_j$ and $s_j-1-2a$ on
$f_a$. Put $Y_m=e_{\mathrm{last}}^\perp\subset X_m$, with no string. Let $c_M$ be the idele that
generates $(\pi_M)$ and is $1$ away from the chosen places, $\mathsf a=\alpha/\aubar{c_M}$,
$\psfa$ the ideal of $\mathsf a$, $u_j=(\mathsf a/\aubar{\mathsf a})^{\lceil(m-j)/2\rceil}$, and
$c=-(\jmath M\delta_E)^{-1}\in E^+$ with $\delta_E=\varpi$ the purely imaginary generator of the different of $E$ fixed in Lemma~\ref{pds:kappa} (the one used in \eqref{pb:eq-psi});
$c$ has the valuations of $\mathsf a\aubar{\mathsf a}$ at all finite places. The lattice
$\psM_j$ of $Y_j$ is spanned over $\widehat\auO_E$ by
\begin{equation}\label{ps:eq-lattice}
\begin{cases}
e_i\ (i<n_j),\quad y/\aubar{\pi_N}+u_jz/N,\quad u_jf_a\ (\text{all }a)&(j\text{ even}),\\
e_i\ (i\le n_j),\quad u_j\mathsf a^{-1}f_a\ (2a<s_j),\quad u_j\aubar{\mathsf a}f_a\ (2a>s_j)&(j
\text{ odd}),
\end{cases}
\end{equation}
and $\psM_m=\widehat\auO_E^{\,n_m-1}$. It is unimodular: for $j$ even the vector
$y/\aubar{\pi_N}+u_jz/N$ is isotropic and pairs to $-1$ with $u_jz$. The idele $u_j$ satisfies
$u_j\aubar{u_j}=1$ and is not a unit: with $e=\lceil(m-j)/2\rceil$, at a chosen place $w$ it is
$c_{M,w}^{\,e}$, a generator of $\pi_M^{\,e}\auO_w$, and at $\aubar w$ it is $\aubar{c_{M,w}}^{\,-e}$, at a place $w\mid\omega$ above $q_1$
it is $\alpha_w^{\,e}$ and at $\aubar w$ it is $\aubar{\alpha_w}^{\,-e}$, and it is $1$ elsewhere
(in particular above $N$). For $j=0$ this is the lattice $L_{u_0}=u_0\cdot\mathsf h_{u_0}L$ of
Lemma~\ref{pb:lattice}, where $L$ is the lattice of $Y_0$ used in \S\ref{pb:sec} (the same basis
with $u_0$ replaced by $1$) and $\mathsf h_{u_0}$ is the central element of $\auU(X_0)(\auA_f)$ acting by $u_0^{-1}$;
all the objects on $Y_0$ of \S\ref{ss:au-spaces} (the level $J_M$, the space $\auAz$, the action
of $\auG$, the map $\auP$) are formed with $\psM_0$, and Lemma~\ref{pb:lattice} carries the
statements of Proposition~\ref{au:moment} from $L$ to $\psM_0$. For $j$ even, $\psM_j\cap X_j=
\widehat\auO^{\,n_j-1}\oplus\pi_N\widehat\auO y$, and $\psM_j$ is not the sum of its intersections
with $X_j$ and with the string at the places above $N$ (it is \emph{glued} there). We write
$\psM_j^-$ for the intersection of $\psM_j$ with the span of the string and $\psM_j^+$ for its
projection to that span; they are dual to each other, and equal except that $\psM_j^+\ni u_jz/N$
for $j$ even. Let $G_j=\auU(Y_j)$, $K_j\subset G_j(\auA_f)$ the stabiliser of $\psM_j$,
$J_j\subseteq K_j$ the subgroup of elements $\equiv1$ at the chosen places and at
the places of $Q=Q_\vartheta$ (the principal congruence subgroup modulo $(\pi_M)\prod_{w\in Q}w$; for
$j=m$ it is contained in, and for $Q\ne\emptyset$ strictly smaller than, the level $J_\vartheta$ of
$\auAm(\vartheta)$, which is block upper triangular modulo $Q$, Lemma~\ref{pz:lastlevel}), and $J_M\supseteq J_0$ the level of $\auAz$. The frame of $\psM_j$
used to write reductions is \eqref{ps:eq-lattice} up to an element of $\auGL(\widehat\auO)$ that
is $\equiv1$ at the chosen places, numbered so that the frame of $Y_j$ is the beginning of the
frame of $Y_{j+1}$; $\iota_j$ denotes the corresponding inclusion matrix. (Any numbering with
this property will do; nothing below depends on it.)

\emph{The string map.} Let $T_j\colon Y_j\to Y_{j+1}$ be zero on $X_j$ and $T_jf_a=t_af'_a$
($a\le s_j-2$), $T_jf_{s_j-1}=0$, where $f'_a$ is the string of $Y_{j+1}$ and, for $j$ even,
$t_a=1,\,-Nc,\,c$ according as the grade of $f_a$ is positive, zero, negative, and for $j$ odd
$t_a=c,\,1$ according as it is positive, negative. Then $T_j$ lowers grades by $1$, and a direct
computation with \eqref{ps:eq-lattice} gives
\begin{equation}\label{ps:eq-TM}
T_j(\psM_j^+)=\psfa\,\psM_{j+1}^-,\qquad T_j^\dagger T_j=\mathfrak e_j,\qquad T_jT_j^\dagger=
\mathfrak e_{j+1},
\end{equation}
where $\mathfrak e_j$ is a self-adjoint nilpotent endomorphism of $Y_j$ that kills $X_j$ and maps
$f_a$ to a non-zero multiple of $f_{a+1}$. (For $j$ even both sides of the first identity are
$u_j(\bigoplus_{a<p}\widehat\auO f'_a\oplus c\bigoplus_{a\ge p}\widehat\auO f'_a)$, $2p=s_j-1$,
because $u_{j+1}=u_j$; for $j$ odd both are $u_j\aubar{\mathsf a}\bigoplus_a\widehat\auO f'_a$,
because $u_j\aubar{\mathsf a}=u_{j+1}\mathsf a$.)

\emph{Bessel groups.} Let $N_j\subset G_j$ be the group of $u$ such that $u-1$ raises grades,
$N_j^{\ge2}$ the subgroup where it raises them by at least $2$ (so $N_j^{\ge2}=N_j$ for $j$ even),
and $R_j=N_j\rtimes\auU(X_j)$, where $\auU(X_j)$ acts on $X_j$ and fixes the string; put
$R_m=G_m$. Equivalently $R_j$ is the group of $g$ with $gf_a\in f_a+\langle f_b:b<a\rangle$ for
all $f_a$ of grade $\ge0$. Let $\psi_j(u)=\psi(\tfrac12\jmath\operatorname{tr}(\mathfrak e_j(u-1)))$,
a character of $N_j^{\ge2}(\auA)$ trivial on $N_j^{\ge2}(E^+)$ and invariant under conjugation by
$R_j(\auA)$. A continuous function $\psW$ on $R_j(E^+)\backslash R_j(\auA)$ is
\emph{equivariant} if $\psW(ur)=\psi_j(u)\psW(r)$ for $u\in N_j^{\ge2}(\auA)$. The group $R_0$ is
the Bessel subgroup of \S\ref{ss:au-spaces}, $\psi_0=\psi_U$, and $\psW_\varphi$ is equivariant.
(For the character \eqref{pb:eq-psi}: with $p=(s_0-1)/2$ and $u_{a,a+1}$ the coefficient of $f_a$ in
$uf_{a+1}$, \eqref{ps:eq-TM} gives $\mathfrak e_0f_a=cf_{a+1}$ for $a\ne p$ and $\mathfrak e_0f_p=-Ncf_{p+1}$,
and the unitarity of $u$ gives $u_{s_0-2-a,s_0-1-a}=-\aubar{u_{a,a+1}}$ for $a<p-1$ and
$u_{p,p+1}=\aubar{u_{p-1,p}}/N$; hence $\operatorname{tr}(\mathfrak e_0(u-1))=c\,(\sigma(u)-\aubar{\sigma(u)})$
with $\sigma(u)=\sum_{a<p}u_{a,a+1}$ the sum in \eqref{pb:eq-psi}, and $\psi_0(u)=\psi(\jmath c\,\sigma(u))=
\psi(-\sigma(u)/(M\delta_E))$.)

\emph{The data.} The lift $\Theta^Y_{j,\vartheta}$ is $[K_j:J_j]$ times the theta lift for
$\omega_{\mu_{j+1}}$ on $\psHS(\auHom(Y_j,Y_{j+1}))$ with the datum $\Upsilon_j=P_j\otimes
\Upsilon_{j,f}$: $P_j$ is the joint harmonic of \S\ref{ss:au-spaces}, and $\Upsilon_{j,f}$ is the
vacuum $\psvac_{L_j}$, $L_j=\alpha\auHom(\psM_j,\psM_{j+1})$, at the places outside the level
places, and at a level place $w$ (chosen, or in $Q$) the function
$\mathbf 1_{a_w(\iota_j+\varpi_w\auHom(\psM_{j,w},\psM_{j+1,w}))}$ in the model (W5), written in the
frames, with $a_w$ the component of the idele $a$ of \S\ref{ps:ss-weil} (so $a_w=1$ except at the
primes of $Q$ above $q_1$ dividing $\omega$; in the coordinates $\tilde x=a_w^{-1}x$ of
Lemma~\ref{pgx:dict} this is $\mathbf 1_{\iota_j+\varpi_w\auHom}$, the datum $\delta_\iota^\natural$ of
Definition~\ref{pgx:conv}). By (W5), Lemma~\ref{pgx:vac} ($L_{j,w}=a_w\auHom(\psM_{j,w},\psM_{j+1,w})$ and
$L_{j,\bar w}$ are dual for $[\,,]$) and finite Fourier inversion on $L_{j,w}/\varpi_wL_{j,w}$, with
$\iota^\natural_j=a_w\iota_j$ at each level place,
\begin{equation}\label{ps:eq-upsilon}
\Upsilon_{j,f}=\prod_w\frac1{\#_w}\sum_{a}[\iota^\natural_j,a]^{-1}\,\psrho^\circ(a)\psvac_{L_j},\qquad
\psrho(l)\Upsilon_{j,f}\in\C^\times\Upsilon_{j,f}\ \ (l\in L'_j),
\end{equation}
where $a$ runs over $\varpi_{\bar w}^{-1}L_{j,\bar w}/L_{j,\bar w}$ for each level place $w$,
$\#_w$ is the number of such $a$, and $L'_j$ is the self-dual lattice equal to $L_j$ outside the
level places and to $\varpi_wL_{j,w}\oplus\varpi_{\bar w}^{-1}L_{j,\bar w}$ at them. The forms on
$G_j$ \emph{of the archimedean type of the chain} are defined in Definition~\ref{pz:def-type}. We
use the following facts, which are Proposition~\ref{pz:types} (it is proved in \S\ref{pz:ss-types}
from three statements (A1)--(A3) on discrete series of real unitary groups): $\Theta^Y_{j,\vartheta}$
maps the forms on $G_j$ of the archimedean type of the chain to those on $G_{j+1}$; the
transposed lift $\Theta^{Y,*}_{j,\vartheta}z(g)=[K_j:J_j]\int_{[G_{j+1}]}\theta(\omega_{\mu_{j+1}}
(g,h)\Upsilon_j)z(h)\,dh$ maps the complex conjugates of these forms on $G_{j+1}$ to those on
$G_j$; and all these forms are cusp forms, hence rapidly decreasing.

\begin{lemma}\label{ps:level}
Let $k\in K_j$, $k'\in K_{j+1}$ have reductions $\bar k,\bar k'$ at the level places with
$\bar k'\iota_j=\iota_j\bar k$. Then $\omega_{\mu_{j+1}}(k,k')\Upsilon_{j,f}=\prod_w\psi_{j+1,w}
(\det\bar k'_w)\psi_{j+1,w}(\det\bar k_w)^{-1}\cdot\Upsilon_{j,f}$, where $\psi_{j+1,w}$ is the
tame part of $\mu_{j+1}$ at $w$. In particular $J_j\times J_{j+1}$ fixes $\Upsilon_{j,f}$.
\end{lemma}
\begin{proof}
Outside the level places $k,k'$ stabilise $L_j$ and all characters are unramified, so (W4)
applies. At a level place $w$, by (W5) the function $\mathbf 1_{a_w(\iota+\varpi\auHom)}$ is carried
to $x\mapsto\mathbf 1(k_w'^{-1}xk_w\in a_w(\iota+\varpi\auHom))=\mathbf 1(x\in a_w(k'_w\iota k_w^{-1}+\varpi
\auHom))$ ($a_w$ is a scalar), which is the same function by hypothesis, times $\mu_{j+1,w}(\det k'_w)\mu_{j+1,w}(\det
k_w)^{-1}$; the factor $\upsilon'_w$ is trivial on units: by Proposition~\ref{pwf:W5} (valid for ranks of either parity) it is built from $\upsilon_w$ and $|\cdot|_w$ evaluated at $\det k_w,\det k'_w\in\mathcal O_w^\times$, and $\upsilon$ is unramified at every
finite place.
\end{proof}

\begin{proof}[Proof of Proposition~\ref{au:equiv}]
The statement is proved in the following form, which is what the proof of
Proposition~\ref{au:good} uses: \emph{for every tame character $\vartheta$ and every $u\in\auG$
there is $\beta_\vartheta(u)\in K_m$ with $(\iota_\vartheta(\varrho(u)z))(g)=(\iota_\vartheta z)
(g\beta_\vartheta(u))$ for all $z\in\auAz$ and $g$.} (The element depends on $\vartheta$ through
the places of $Q_\vartheta$.) The level places are split, so the reduction $K_i\to\prod_w\auGL
(\psM_i/w)$ is onto. Choose $k_0\in K_0$ with reduction $u$ at the chosen places and $1$ at $Q$, and
$k_i\in K_i$ ($1\le i\le m$) with reduction $\operatorname{diag}(u,1)$ at the chosen places and $1$
at $Q$, in the frames; then $\bar k_{i+1}\iota_i=\iota_i\bar k_i$. By (W1) and
Lemma~\ref{ps:level}, for a form $F$ on $G_i$,
\[
(\Theta^Y_{i,\vartheta}F)(hk_{i+1})=[K_i:J_i]\!\int_{[G_i]}\!\theta\bigl(\omega(g,h)\omega(k_i,k_{i+1})
\omega(k_i^{-1},1)\Upsilon_i\bigr)F(g)\,dg=\varepsilon_i\,\bigl(\Theta^Y_{i,\vartheta}(F(\,\cdot\,k_i))
\bigr)(h),
\]
after the change of variable $g\mapsto gk_i$, with $\varepsilon_i=\prod_w\psi_{i+1,w}(\det\bar
k_{i+1,w})\psi_{i+1,w}(\det\bar k_{i,w})^{-1}=1$, because $\det\operatorname{diag}(u_w,1)=\det u_w$.
By induction $(\iota_\vartheta z)(gk_m)=\iota_\vartheta(z(\,\cdot\,k_0))(g)$, and
$z(\,\cdot\,k_0)=\varrho(u)z$ by the definition of $\varrho$. Put $\beta_\vartheta(u)=k_m$.
\end{proof}

\subsubsection{Branching and restriction: proof of Proposition~\ref{au:restrict}}
\label{ps:ss-restrict}

\emph{Conventions for $W_\lambda$.} For $\auGL_n(\C)$ and a diagram with $c(h)$ columns of height
$h<n$, the module $W_\lambda$ is the span, in the polynomials on $M_n(\C)$, of the products of
minors $\prod_{\text{columns}}\Delta_{I,[h]}$ (rows $I$, $|I|=h$, columns $1,\dots,h$), with
$(gp)(x)=p(g^tx)$, twisted by $\det^e$. It is the image of $\bigotimes_{\text{columns}}\Lambda^h
\C^n$ under the multiplication map $\mathsf m$, which is the orthogonal projection onto the Cartan
component when $W_\lambda$ carries the invariant form for which the highest weight vector
$\prod\Delta_{[h],[h]}^{c(h)}$ has norm $1$; this is the hermitian form of \S\ref{ss:au-spaces}.
Hence a product of minors has norm at most $1$. The products of minors span a lattice
$\Sigma_\lambda$, with the standard ones as basis; $W_\lambda(\Z)$ is the dual lattice
$\Sigma_\lambda^\#$, with the dual basis, whose Gram matrix is integral (it is the matrix of
coefficients of the reproducing kernel $\prod\Delta_{[h],[h]}(\bar y^tx)^{c(h)}$ in the products
$b_i(x)\overline{b_j(y)}$). For a linear map $Z\colon\C^n\to\C^{n'}$ the Schur functor
$S(Z)\colon W_\lambda(\C^n)\to W_\lambda(\C^{n'})$ is the restriction of $\bigotimes\Lambda^hZ$; its
weight spaces for a diagonal $Z$ are orthogonal, so, with $s_1\ge s_2\ge\dots$ the singular values
of $Z$ and $l_p$ the rows of the diagram,
\begin{equation}\label{ps:eq-schurnorm}
\|S(Z)\|_{\mathrm{op}}=\prod_ps_p^{\,l_p},\qquad\sum_ps_p^2=\|Z\|_{\mathrm{HS}}^2 .
\end{equation}

\emph{The weight $\lambda^\circ$.} Put $\mathrm{kept}(h)=\tfrac12(\Delta-(t+1))$ if $h\in\{k,n_1,
\dots,n_{m-1}\}$ and $0$ otherwise; then $\mathrm{kept}(h)\le c_\tau(h)$. The weight
$\lambda^\circ$ on $\auU(Y_m)$ has, at $\tau$, $\mathrm{kept}(h)+c_\tau(h+1)-\mathrm{kept}(h+1)$
columns of height $h<n_m-1$ and the determinant exponent $e_{m,\tau}$; that is,
$\lambda^\circ_{\tau,p}=\lambda(X_m)_{\tau,p+1}+\mathrm{kept}(p)$: of the columns of
$\lambda(X_m)_\tau$ of height $h$, $\mathrm{kept}(h)$ keep their height and the others lose a box.
The two weights interlace, and
\begin{equation}\label{ps:eq-ntau}
n_\tau=|\lambda(X_m)_\tau|-|\lambda^\circ_\tau|=\sum_h\bigl(c_\tau(h)-\mathrm{kept}(h)\bigr)+e_{m,\tau}
=a\mathbf 1_{\tau\in\auPa}-\frac{m\Delta}2 .
\end{equation}
Indeed the sum is $a\mathbf 1_{\auPa}(\tau)+\kappa_1(\tau)-\tfrac{n_1}2+(m-1)(\Delta-(t+1))-\tfrac m2
(\Delta-(t+1))$ and $e_{m,\tau}=\tfrac{n_{m-1}}2-\kappa_1(\tau)-(m-1)\Delta$, with
$n_{m-1}-n_1=(m-2)(t+1)$. Summing over $\tau$ gives \eqref{eq:au-branch}, with equality.

\begin{lemma}[The branching map]\label{ps:branching}
Let $L$ be the number of boxes of $\lambda^\circ$, summed over $\tau$, in the columns of height at
least $k$. There is a non-zero linear map $\mathrm{br}\colon W_{\lambda(X_m)}\to W_{\lambda^\circ}$
with:
\begin{enumerate}
\item[(a)] $\mathrm{br}\circ\rho(\operatorname{diag}(g,\zeta))=\prod_\tau\zeta_\tau^{\,n_\tau}\cdot
\rho^\circ(g)\circ\mathrm{br}$ for $g\in\prod_\tau\auGL_{n_m-1}(\C)$, $\zeta\in(\C^\times)^{\auPl}$;
\item[(b)] $\|\mathrm{br}\|_{\mathrm{op}}\le c_{\mathrm{br}}$ for a positive integer
$c_{\mathrm{br}}$;
\item[(c)] the matrix of $\mathrm{br}$ in the bases of $W_{\lambda(X_m)}(\Z)$ and
$W_{\lambda^\circ}(\Z)$ has integer entries, and $c_{\mathrm{br}}\le L^{3L}$.
\end{enumerate}
\end{lemma}
\begin{proof}[Proof of (a), (b)]
Let $v\in W_{\lambda(X_m)}$ be the product, over $\tau$ and over the columns, of the minors
$\Delta_{[h],[h]}$ for the columns that keep their height and $\Delta_{\{n_m\}\cup[h-1],[h]}$ for
those that lose a box. It is a highest weight vector of weight $\lambda^\circ$ for
$\prod_\tau\auGL_{n_m-1}$, of weight $\prod\zeta_\tau^{n_\tau}$ for the last diagonal entries by
\eqref{ps:eq-ntau} (the exponent $e_{m,\tau}$ comes from the twist), and $\|v\|\le1$. Let
$i\colon W_{\lambda^\circ}\to W_{\lambda(X_m)}$ be the equivariant map sending the highest weight
vector to $v$; it has norm $\|v\|\le1$, since both modules are irreducible and unitary. Put
$\mathrm{br}=c_{\mathrm{br}}\,i^*$ with $c_{\mathrm{br}}=\prod_\tau H_{\lambda^{\circ\flat}_\tau}$,
the product of the hook numbers (Lemma~\ref{pk:hook}) of the diagrams $\lambda^\circ_\tau$ without
their first $k-1$ rows.
\end{proof}
\noindent Part (c) is Proposition~\ref{pk:branching}.

\begin{proof}[Proof of Proposition~\ref{au:restrict}, given Lemma~\ref{ps:branching}]
The restriction is $(\operatorname{res}\psi)(y)=\prod_\tau\tau(\aubar{\pi_N})^{-n_\tau}\cdot
\mathrm{br}\bigl(\psi(\operatorname{diag}(y,1))\bigr)$; by (a) it is an algebraic modular form of
weight $\lambda^\circ$. Let $\sigma\in\auAut(\C)$. As $\mathrm{br}$ has an integral matrix,
$\sigma(\mathrm{br}\,\xi)=\mathrm{br}(\sigma\xi)$, so $\|\sigma(\mathrm{br}\,\xi)\|\le
c_{\mathrm{br}}\|\sigma\xi\|$. For every embedding $\tau'$ of $E$ in $\C$,
$|\tau'(\aubar{\pi_N})|^2=\tau'(\pi_N\aubar{\pi_N})=N$, because complex conjugation induces $c$ on
the CM field $E$; hence $|\sigma\prod_\tau\tau(\aubar{\pi_N})^{-n_\tau}|^2=N^{-\sum n_\tau}$.
Finally $L$ is small (at most $\Delta$ columns of each of at most $2m$ heights
$\ge k$, by the definition of $\lambda^\circ$; here $L\le2f\,n_m\,m\Delta$), so $L\le x^{60}$ and $c_{\mathrm{br}}^2\le
L^{6L}\le\exp(x^{70})$.
\end{proof}

\subsubsection{The size of one kernel: proof of Proposition~\ref{au:kernel}}\label{ps:ss-kernel}

\begin{lemma}[Conjugates of jets]\label{ps:conj}
Let $X,X'$ be standard of ranks $n<n'$, $W=\auHom(X,X')$, and $\sigma\in\auAut(\C)$ with
$\sigma|_K=1$. There are $t_\sigma\in\C$ with $|t_\sigma|^2=\Omega^{-1}$, a permutation
$\varpi_\sigma$ of $\auPl$, and a $\sigma$-semilinear map $\Lambda_\sigma$ of $\psHS_f(W)$ such
that for all $f\in\psHS_f(W)$:
\begin{enumerate}
\item[(i)] the Taylor series of $\Theta_{\Lambda_\sigma f}(t_\sigma\,\cdot)$ is obtained from that
of $\Theta_f(s_0\,\cdot)$ by applying $\sigma$ to the coefficients and $\varpi_\sigma$ to the
variables $z_\tau$;
\item[(ii)] if $f=\psrho^\circ(x)\omega(g)\psvac^\circ$ with $x\in W(\auA_f)$, $g\in\auU(W)
(\auA_f)$, then $|\theta(\psrho_\infty(y)(1\otimes\Lambda_\sigma f))|\le\exp\bigl(C_0(q_1nn'+1)^2
\bigr)$ for all $y$, where $C_0$ depends only on $q_0$ and $d$.
\end{enumerate}
\end{lemma}
\begin{proof} Proposition~\ref{pk:conj}, with $t_\sigma=\sigma(\varpi_{\mathrm{vac}}/\mathsf u)c_\psi/
\Omega_\sigma$ and an explicit $C_0(q_0,d)$, and Theorem~\ref{pf:recip}.\end{proof}

\begin{lemma}\label{ps:monomial}
Let $L\ge1$, $u\in\R$, $R\ge0$, let $\rho_i\ge0$ and $l_i\in\N$ ($i$ in a finite set) with
$\sum_i\rho_i\le e^u(2L+2)R$. Then $\prod_i\rho_i^{\,l_i}\le\exp\bigl(\sum_i(l_i/L+\log l_i!+l_iu)
+2LR\bigr)$.
\end{lemma}
\begin{proof}
With $v=u+1/L$ and $x^l\le l!\,e^x$ for $x=\rho_ie^{-v}$: $\rho_i^{l_i}\le l_i!\,e^{l_iv}\exp
(\rho_ie^{-v})$, and $\sum_i\rho_ie^{-v}\le(2L+2)R\,e^{-1/L}\le(2L+2)R/(1+1/L)=2LR$.
\end{proof}

\begin{proof}[Proof of Proposition~\ref{au:kernel}, given Lemma~\ref{ps:conj}]
Here $X=X_j$, $X'=X_{j+1}$, $n=n_j$, $n'=n_{j+1}$, $\phi=\phi^\circ_x$, the twist is trivial, and
$\theta_j(\phi;g,g')$ is \eqref{ps:eq-kernel}. Since $|\sigma w|=|\sigma'w|$ for all $w$, where
$\sigma'$ is $\sigma$ or its composite with complex conjugation, and one of them fixes $K$, we may
assume $\sigma|_K=1$.

\emph{Step 1: the datum.} The element $(g,g')$ acts on $W$ through $k\in\auU(W)(\auA_f)$, and
$\omega_{\mu_{j+1}}(g,g')\phi^\circ_x=\epsilon\,\psrho^\circ(kx)\omega(k)\psvac^\circ$ by (W1), where
$\epsilon$ is the product of a root of unity and of the values of the characters of
\S\ref{ps:ss-weil} at $\det g$, $\det g'$. These characters have algebraic values on finite adeles
and $|\sigma\epsilon|=1$: the values are products of roots of unity and of numbers
$\prod_\tau\tau(\nu)^{e_\tau}$ with $\nu\in E$, $\nu\aubar\nu=1$, all of whose conjugates have
absolute value $1$. So it suffices to bound $\sigma$ of the jets of $f=\psrho^\circ(x')\omega(k)
\psvac^\circ$ at scale $\pi_Ms_0$.

\emph{Step 2: a conjugate jet is a Gaussian integral.} Let $Q$ be a polynomial, $Q^\sigma$ the
polynomial obtained by applying $\sigma$ to its coefficients and $\varpi_\sigma$ to its variables.
By Lemma~\ref{ps:conj}(i) and \eqref{ps:eq-bargmann} for $\Lambda_\sigma f$ at scale $t_\sigma$,
\[
\sigma\bigl(Q(\partial)[\Theta_f(\pi_Ms_0\,\cdot)](0)\bigr)=\int Q^\sigma\bigl(\sigma(\pi_M)\,t_\sigma
\bar y\bigr)\,e^{-|y|^2/2}\,\theta\bigl(\psrho_\infty(y)(1\otimes\Lambda_\sigma f)\bigr)\,d\mu(y).
\]
\emph{Step 3: growth of the Schur matrix.} Let $\xi\in W_{\lambda(X_j)}$ and let $\eta$ be a
covector on $W_{\lambda(X_{j+1})}$ with $\|\sigma\eta\|\le1$ for the dual norm, and $Q=\eta\cdot
S(z)\xi$. The entries of $S$ have integer coefficients, so $Q^\sigma(z)=\sigma\eta\cdot S(z')
\sigma\xi$ with $z'$ a permutation of $z$, and by \eqref{ps:eq-schurnorm} $|Q^\sigma(z)|\le
\|\sigma\xi\|\prod_{\tau,p}s_{\tau,p}^{\,l_{\tau,p}}$ with $\sum_{\tau,p}s_{\tau,p}^2=|z|^2$. For
$z=\sigma(\pi_M)t_\sigma\bar y$ one has $|z|^2=(M/\Omega)|y|^2$, as $|\sigma\pi_M|^2=M$. Apply
Lemma~\ref{ps:monomial} to $\rho=s^2$ with $u=\log(M/\Omega)$, $R=|y|^2/(2L+2)$:
\[
|Q^\sigma(\sigma(\pi_M)t_\sigma\bar y)|\le\|\sigma\xi\|\exp\Bigl(\tfrac12\sum_{\tau,p}\bigl(
\tfrac{l_{\tau,p}}L+\ausinf(l_{\tau,p})\bigr)+\frac{L|y|^2}{2L+2}\Bigr),
\]
since $\log l!+l\log(M/\Omega)=\ausinf(l)$.

\emph{Step 4: conclusion.} By Steps 2, 3 and Lemma~\ref{ps:conj}(ii) the integrand is at most
$\|\sigma\xi\|\exp(C_0(q_1nn'+1)^2+\tfrac12\sum(l/L+\ausinf(l)))\cdot\exp(-|y|^2/(2L+2))$, and
$\int\exp(-|y|^2/(2L+2))\,d\mu=(2L+2)^{fnn'}$. Taking the supremum over $\eta$,
\[
\|\sigma(\theta_j(\phi;g,g')\xi)\|\le\exp\Bigl(C_0(q_1nn'+1)^2+fnn'\log(2L+2)+\frac1{2L}\sum_{\tau,p}
l_{\tau,p}+\tfrac12\mathrm{rs}_j\Bigr)\|\sigma\xi\| .
\]
Take $L=\sum_{\tau,p}l_{\tau,p}\le fn_m\cdot2a$. Then the third term is $\tfrac12$, and by
Lemma~\ref{an:lem-sizes} ($f\le q_1\le x^{12}$, $nn'\le n_m^2\le1600x^{12}$, $\log(2L+2)\le3\log
b$) and $C_0\le\log\log b$ (Lemma~\ref{pz:C0}) the first two are at most $x^{50}-1$; see \S\ref{pz:ss-const}.
\end{proof}
The rows $l_{\tau,p}$ are: $a\mathbf 1_{\auPa}(\tau)+\kappa_1(\tau)-\tfrac{n_1}2+j(\Delta-(t+1))$
for $p=1$, the same without $a$ for $1<p\le k$, and at most $m\Delta$ for $p>k$; this is the
description of the rows in \S\ref{ss:au-size}, with slack at most $\kappa_1(\tau)-\kappa(\tau)+j\Delta\le(m+1)\Delta$ per row.

\subsubsection{Gluing and chain data; proof of Proposition~\ref{au:expansion}}\label{ps:ss-glue}

Write $W_j=\auHom(X_j,X_{j+1})$, $L^X_j=\alpha\auHom(\widehat\auO^{\,n_j},\allowbreak \widehat\auO^{\,n_{j+1}})$,
and let $\mathsf P$ be the set of the $f$ primes of $E$ dividing $\pi_N$. The operators
$\psrho^\circ$ of (W4) are products of local operators, so $\psrho^\circ(x+x')=\psrho^\circ(x)
\psrho^\circ(x')$ when $x,x'$ have disjoint supports. A \emph{translation tuple} is
$x=(x_0,\dots,x_{m-1})$, $x_j\in W_j(\auA_f)$; it is \emph{admissible for $q$} if
$\aubar{\pi_M}x_j\in L^X_j$ at all places not dividing $q$. For $0\le i<m/2$ and a family
$\mathsf y=(\mathsf y_{b,w})$ of integers in $[0,N)$ indexed by the basis vectors $e'_b$ of
$X_{2i+1}$ and $w\in\mathsf P$, let $p_{\mathsf y}\in W_{2i}(\auA_f)$ be the map that kills
$e_1,\dots,e_{n_{2i}-1}$ and sends $e_{\mathrm{last}}$ to $-\alpha N^{-1}\sum_b\mathsf y_{b,w}e'_b$
at each $w\in\mathsf P$ and to $0$ at the other places, and $p'_{\mathsf y}\in W_{2i+1}(\auA_f)$
the map $e'_b\mapsto-\alpha N^{-1}\mathsf y_{b,w}\,e_{\mathrm{last}}$ at $\aubar w$, $w\in\mathsf P$,
and $0$ elsewhere. The \emph{gluing} of the $i$-th pair of lifts is
\[
\psGlue_i=\sum_{\mathsf y}\psrho(p_{\mathsf y})\otimes\psrho(p'_{\mathsf y})\quad
\text{on }\psHS_f(W_{2i})\otimes\psHS_f(W_{2i+1}),
\]
and $\psGlue=\bigotimes_i\psGlue_i$; the points $p_{\mathsf y}$, $p'_{\mathsf y}$ are supported on
$\mathsf P$, respectively $\overline{\mathsf P}$, and $\psrho$ of such points is given by
\eqref{pl:eq-heis}. For a tuple $x$ put $\phi^\circ_x=\bigotimes_j\phi^\circ_{j,
x_j}$, and for a finitely supported $\Psi$ put $C_{\vartheta,\Phi}=\Theta^X_{m-1}\cdots\Theta^X_0
(\Phi)$ for a tensor $\Phi$ (by linearity). A \emph{chain datum} for $\vartheta$, with
$q=\auNm Q_\vartheta$, is a finitely supported $\Psi$ such that: (C1) every $x$ in its support is
admissible for $q$ and vanishes at the places not dividing $Mq$; (C2) every $\Psi(x)$ is an
algebraic number integral at all primes not dividing $q$; (C3) if $Q_\vartheta=\emptyset$ then
$\sum_x|\sigma\Psi(x)|\le M^{3fmn_m^2}$ for all $\sigma$; (C4) the tensor $\psGlue(\sum_x\Psi(x)
\phi^\circ_x)$ is fixed by $\bigotimes_j\omega_{\mu_{j+1}}(b_j,b_{j+1})$ for all $b_j\in J_{X_j}$
($j<m$), $b_m=1$. The composite lift is $B_{\vartheta,\Psi}=\operatorname{res}\circ C_{\vartheta,
\psGlue(\sum\Psi(x)\phi^\circ_x)}$.

\begin{proof}[Proof of Proposition~\ref{au:expansion}]
Let $x$ be in the support of $\Psi$. By (C1) $x$ vanishes above $N$, while $p_{\mathsf y}$,
$p'_{\mathsf y}$ are supported above $N$; so $\psGlue_i(\phi^\circ_{2i,x_{2i}}\otimes\phi^\circ_{2i+1,
x_{2i+1}})=\sum_{\mathsf y}\zeta_{\mathsf y}\,\phi^\circ_{2i,x_{2i}+p_{\mathsf y}}\otimes
\phi^\circ_{2i+1,x_{2i+1}+p'_{\mathsf y}}$ with roots of unity $\zeta_{\mathsf y}$ ($\psrho^\circ$
and $\psrho$ differ by roots of unity). Hence $\psGlue(\sum_x\Psi(x)\phi^\circ_x)=\sum_{x'}\Psi'(x')
\phi^\circ_{x'}$ with
$\Psi'(x')=\sum\zeta\,\Psi(x)$ over the pairs $(x,(\mathsf y^{(i)})_i)$ with $x'=x+(p_{\mathsf y^{(i)}},
p'_{\mathsf y^{(i)}})_i$, $\zeta$ the product of the $\zeta_{\mathsf y^{(i)}}$, and
$B_{\vartheta,\Psi}=\operatorname{res}\circ C_{\vartheta,\Psi'}$. The
values of $\Psi'$ are sums of roots of unity times values of $\Psi$, so integral outside $q$; the tensor is the same,
so it is invariant under the levels; the tuples $x'$ are admissible for $qN$. The number of
families $(\mathsf y^{(i)})$ is $N^{f\sum_in_{2i+1}}\le N^{fmn_m/2}$, so for $Q_\vartheta=\emptyset$,
$\sum_{x'}|\sigma\Psi'(x')|\le M^{3fmn_m^2}N^{fmn_m/2}\le\exp(x^{60})$ by \eqref{eq:au-sizes}.
\end{proof}

\subsubsection{The seesaw identity: proof of Proposition~\ref{au:seesaw}}\label{ps:ss-seesaw}

We keep the notation of \S\ref{ps:ss-Y}. For gradings on $V$ and $V'$ a map $w\in\auHom(V,V')$
\emph{has degree $\ge d$} if it maps the part of grade $a$ into the parts of grade $\ge a+d$;
$\auHom_{>0}$, $\auHom_{\ge0}$ denote the maps of degree $\ge1$, $\ge0$. The space
$I=\auHom_{>0}(V,V')$ is isotropic and $I^\perp=\auHom_{\ge0}(V,V')$. For $F\in\psHS(\auHom(V,V'))$
and $w\in\auHom(V,V')(\auA)$ put
\[
c_w(F)=\theta_{I^\perp/I}\bigl(\psZ_I\psrho(w)F\bigr)=\int_{[I]}\theta(\psrho(i)\psrho(w)F)\,di .
\]
By (W1), (W2): $c_w$ depends only on $w$ modulo $I(\auA)+I^\perp(E)$ up to the factors of
$\beta$; $c_{w+\delta}(F)=\beta(\delta,w)^{-1}c_w(F)$ for $\delta\in I(\auA)$; and if $p$ is a
unipotent element of the stabiliser of $I$ acting trivially on $I^\perp/I$, then
$c_w(\omega(p)F)=c_{p^{-1}w}(F)$.

\begin{proposition}[One step]\label{ps:onestep}
Let $0\le j<m$, $V=Y_j$, $V'=Y_{j+1}$, $T=T_j$, $F\in\psHS(\auHom(V,V'))$, $\psW$ an equivariant
function on $[R_j]$, and
\[
\psW'(h)=\int_{[R_j]}\psW(r)\,c_T\bigl(\omega_{\mu_{j+1}}(r,h)F\bigr)\,dr\qquad(h\in G_{j+1}(\auA)).
\]
Then $\psW'$ is continuous, left invariant under $R_{j+1}(E^+)$, and $\psW'(n'h)=\psi_{j+1}(n')
\psW'(h)$ for $n'\in N_{j+1}^{\ge2}(\auA)$; and for every rapidly decreasing automorphic function
$z$ on $G_{j+1}$, the integral $\int_{R_{j+1}(E^+)\backslash G_{j+1}(\auA)}\psW'z$ converges
absolutely and
\[
\int_{[R_j]}\psW(r)\int_{[G_{j+1}]}\theta\bigl(\omega_{\mu_{j+1}}(r,h)F\bigr)z(h)\,dh\,dr=
\int_{R_{j+1}(E^+)\backslash G_{j+1}(\auA)}\psW'(h)z(h)\,dh .
\]
\end{proposition}
\begin{proof}
For $j=m-1$ there are no strings on $V'$ and at most the vector $z$ on $V$: $I=0$, $T=0$, $c_T=
\theta$, and there is nothing to prove. Let $j<m-1$, $s=s_j$. Write $V=V^+\oplus V^0\oplus V^-$
by the sign of the grade, $p=\dim V^+$, and $T_+=T|_{V^+}$; $T_+$ is injective, and by
\eqref{ps:eq-TM} $R_{j+1}$ is the group of $h\in G_{j+1}$ that preserve $T(V^+)$ and act on it by a
unitriangular matrix in the basis $Tf_0,\dots,Tf_{p-1}$. Let $N_L\subset N_j$ be the unitriangular
group of $\auGL(V^+)$, $U^+$ the group of $u\in G_j$ acting trivially on $V^+$, on
$V^{+\perp}/V^+$ and on $V/V^{+\perp}$, and $\mathcal Z\subset U^+$ the subgroup acting trivially on
$V^0$; then $N_j=N_L\ltimes U^+$, and $N_L,\mathcal Z\subseteq N_j^{\ge2}$.

\emph{(1) Invariance.} For $n\in N_j^{\ge2}(\auA)$, $Tn-T=T(n-1)\in I(\auA)$, so
$c_T(\omega(n,1)F')=c_{Tn}(F')=\beta(T(n-1),T)^{-1}c_T(F')=\psi_j(n)^{-1}c_T(F')$, since
$\langle T,T(n-1)\rangle=\operatorname{tr}(\mathfrak e_j(n-1))$. Likewise $c_T(\omega(1,n')F')=
\psi_{j+1}(n')c_T(F')$ for $n'\in N^{\ge2}_{j+1}(\auA)$, and for $\delta\in R_{j+1}(E^+)$,
$\gamma\in R_j(E^+)$ the map $\delta^{-1}T\gamma-T$ has degree $\ge0$, so
$c_T(\omega(\gamma,\delta)F')=c_T(F')$ by (W1). This gives the properties of $\psW'$ (continuity
is clear, $[R_j]$ being compact).

\emph{(2) Coarse expansion.} Let $I_0\subseteq I$ be the maps vanishing on $V^+\oplus V^0$; then
$I_0^\perp$ is the space of maps vanishing on $V^+$, it is stable under $G_j'=P^+\times G_{j+1}$
($P^+$ the stabiliser of $V^+$), and by (W3)
$\theta(F')=\sum_{w\in\auHom(V^+,V')(E)}c^0_w(F')$, $c^0_w=\theta\psZ_{I_0}\psrho(\tilde w)$,
$\tilde w=w$ on $V^+$ and $0$ on $V^0\oplus V^-$, with $c^0_w(\omega(\gamma,\delta)F')=
c^0_{\delta^{-1}w\gamma}(F')$ for rational $\gamma\in R_j$, $\delta\in G_{j+1}$. For $u\in U^+
(\auA)$, $\tilde wu=\tilde w+\tilde w(u-1)$ and $\tilde w(u-1)\in I_0^\perp(\auA)$ has image
$w\circ\mathrm b(u)\in\auHom(V^0,V')$ in $I_0^\perp/I_0$, where $\mathrm b(u)\colon V^0\to V^+$ is
the component of $u-1$; so by (W2)
\begin{equation}\label{ps:eq-coarse}
c^0_w(\omega(u,1)F')=\psi\bigl(\tfrac12\jmath\operatorname{tr}(\tilde w^\dagger\tilde w(u-1))
\bigr)^{-1}\,\theta\bigl(\psrho(w\mathrm b(u))\,\psZ_{I_0}\psrho(\tilde w)F'\bigr),
\end{equation}
and $c^0_w(\omega(a,1)F')=c^0_w(F')$ for $a\in N_L(\auA)$ with $wa=w$.

\emph{(3) Only the orbit of $T_+$ contributes} to the integral over $[R_j]$ of the coarse
expansion; see below.

\emph{(4) The orbit of $T_+$.} The stabiliser $\mathsf M$ of $T_+$ in $G_{j+1}$ is contained in
$R_{j+1}$, and $\delta\mapsto a_\delta$, $\delta T_+=T_+a_\delta$, identifies
$\mathsf M\backslash R_{j+1}$ with $N_L$ (Witt). Hence $\{\delta^{-1}T_+:\delta\in\mathsf M(E^+)
\backslash R_{j+1}(E^+)\}=T_++J(E)$, $J=\auHom_{>0}(V^+,V')$ (this is correct only because the non-zero grades of $V$ and
$V'$ have different parities, remark (G1) of \S\ref{pk:ss-onestep}), and by (3)
\[
\int_{[R_j]}\psW(r)\theta(\omega(r,h)F)\,dr=\sum_{\gamma\in R_{j+1}(E^+)\backslash G_{j+1}(E^+)}
\int_{[R_j]}\psW(r)\sum_{w\in T_++J(E)}c^0_w(\omega(r,\gamma h)F)\,dr .
\]
\emph{(3), (5), (6).} Steps (3), (5) and (6) are carried out in \S\ref{pk:ss-onestep}
(Lemmas~\ref{pk:vanish}, \ref{pk:coarse}, \ref{pk:fine} and the proof given there).
\end{proof}

The statement is the global transfer of Bessel and Fourier--Jacobi periods under a theta lift for
a hook-type nilpotent orbit; compare Gomez--Zhu~\cite{GomezZhu} for the local statement. (See also B.~Wang, \emph{Global theta lifting and automorphic periods associated to
nilpotent orbits}, J.\ Number Theory 271 (2025), for related global statements.)

\emph{The functions $\psW_j$.} Fix $\vartheta$, $\varphi\in\auS$, and points $g_i\in G_i(\auA)$,
$0\le i<m$. Put $\psW_0=\psW_\varphi$ on $R_0(\auA)$ and, for $0\le j<m$,
\[
\psW_{j+1}[g_0,\dots,g_j](h)=\int_{[R_j]}\psW_j[g_0,\dots,g_{j-1}](rg_j)\,c_{T_j}\bigl(
\omega_{\mu_{j+1}}(rg_j,h)\Upsilon_j\bigr)\,dr\qquad(h\in G_{j+1}(\auA)),
\]
where for $j=0$ the first factor is $\psW_\varphi(r)$. By Proposition~\ref{ps:onestep}(1),
$r\mapsto\psW_j[\dots](rg_j)$ is equivariant on $[R_j]$, and by Lemma~\ref{ps:level} $\psW_{j+1}$ is
right invariant under $J_{j+1}$ and depends on $g_{i,f}$ only through $g_{i,f}J_i$.

\begin{lemma}[Two steps]\label{ps:double}
Let $j=2i<m$, $X=X_j$, $X'=X_{j+1}$, $X''=X_{j+2}$. There is a linear map $\psE_i\colon
\psHS_f(\auHom(Y_j,Y_{j+1}))\otimes\psHS_f(\auHom(Y_{j+1},Y_{j+2}))\to\psHS_f(W_j)\otimes
\psHS_f(W_{j+1})$, independent of $\vartheta$, with the following properties.
\begin{enumerate}
\item[(a)] For every linear form $l$ on $W_{\lambda(X)}$ and all $\xi\in G_j(\R)\times G_{j+1}(\R)
\times G_{j+2}(\R)$ there is a linear form $l''$ on $W_{\lambda(X'')}$ such that: if
$\psW_j[\dots](\iota(g)g_j)=l(g_\infty^{-1}\psi(g_f))$ for $g\in\auU(X)(\auA)$ and an algebraic
modular form $\psi$, then for all $k\in\auU(X'')(\auA)$, with $\iota(k)=k$ if $j+2=m$,
\[
\psW_{j+2}[\dots,g_j,g_{j+1}](\iota(k)g_{j+2})=l''\Bigl(k_\infty^{-1}\bigl(\Theta^X_{j+1}\Theta^X_j
(\Phi)\psi\bigr)(k_f)\Bigr),\quad\Phi=\psE_i\bigl(\omega(g_j,g_{j+1})\Upsilon_{j}\otimes\omega(g_{j+1},
g_{j+2})\Upsilon_{j+1}\bigr)_f,
\]
where $\xi$ is the archimedean part of $(g_j,g_{j+1},g_{j+2})$ and $g_m=1$.
\item[(b)] $\psE_i\circ(\omega(\iota b,\iota b')\otimes\omega(\iota b',\iota b''))=(\omega(b,b')
\otimes\omega(b',b''))\circ\psE_i$ for finite adelic points $b,b',b''$ of $\auU(X)$, $\auU(X')$,
$\auU(X_{j+2})$ (of $\auU(Y_m)$ if $j+2=m$).
\end{enumerate}
\end{lemma}
\begin{proof}
Let $V=Y_j$, $V'=Y_{j+1}$, $V''=Y_{j+2}$, $\mathbb V=V^-\oplus V''$ ($V^-$ is $V$ with the opposite
form), with the grading of $V$ and $V''$; $s=s_j$ is odd. By (W6) the map $w\mapsto w^\dagger$ identifies
$\psHS(\auHom(V,V'))$ with $\psHS(\auHom(V',V^-))$, carrying $\omega_{\auHom(V,V')}(g,g')$ to
$\upsilon(\det g)^{N'}\upsilon(\det g')^{-N}\omega_{\auHom(V',V^-)}(g',g)$ ($N,N'$ the ranks of
$V,V'$); so $\psHS(\auHom(V,V'))\otimes\psHS
(\auHom(V',V''))=\psHS(\auHom(V',\mathbb V))$ with this character, $\theta\otimes\theta$ is a non-zero multiple of
$\theta$, and $c_T(F_1)c_{T'}(F_2)$ is a multiple of $c_S(F_1\otimes F_2)$ for the isotropic space
$\mathbb I=\auHom_{>0}(V',\mathbb V)$ and $S=(T^\dagger,T')$, $T=T_j$, $T'=T_{j+1}$. By
\eqref{ps:eq-TM}, $S^\dagger S=-TT^\dagger+T'^\dagger T'=0$. The maps of degree $0$ are
$\auHom(X',\mathbb V_0)$; if $s\ge3$, $\mathbb V_0=X^-\oplus X_{j+2}\oplus H$ with $H$ the
hyperbolic plane spanned by the middle vectors $z,z''$, and $\ell=E\cdot Sf'_{p-1}$ ($f'_{p-1}$
the vector of grade $1$ of $V'$) is an isotropic line of $H$; if $s=1$, $\mathbb V_0=X^-\oplus
X_m$, where $X_m=Y_m\oplus E\,z/\pi_N$ and $z/\pi_N$, of norm $1$ in $V^-$, is
$e_{\mathrm{last}}$. Let $\mathbb J=\mathbb I\oplus\auHom(X',\ell)$ (resp.\ $\mathbb I$), so that
$\mathbb J^\perp/\mathbb J=\auHom(X',X^-\oplus X'')=W_j\oplus W_{j+1}$, and define
$\psE_i=\psZ_{\mathbb J,f}\circ\psrho_f(S)$.

By definition, $\psW_{j+2}(h)=\int_{[R_{j+1}]}\int_{[R_j]}\psW_j(rg_j)c_T(\omega(rg_j,r'g_{j+1})
\Upsilon_j)c_{T'}(\omega(r'g_{j+1},h)\Upsilon_{j+1})\,dr\,dr'$. Put $F=\omega(g_j,g_{j+1})\Upsilon_j
\otimes\omega(g_{j+1},g_{j+2})\Upsilon_{j+1}$ and $h=\iota(k)g_{j+2}$. By
Proposition~\ref{ps:onestep}(1) the integrand is invariant under $r\mapsto nr$, $n\in N_j(\auA)$
(here $N_j=N_j^{\ge2}$), so the integral over $[R_j]$ is one over $[\auU(X)]$. For $n'\in
N_{j+1}(\auA)$, $Sn'-S=S(n'-1)$ has degree $\ge0$, its component of degree $0$ is the map
$\lambda(n')\in\auHom(X',\ell)$, $x\mapsto(\text{coefficient of }f'_{p-1}\text{ in }n'x)\,Sf'_{p-1}$,
and $\beta(S(n'-1),S)=1$ because $S^\dagger S=0$. Hence, by (W2), the product of the two constant
terms at $(g,n'r',k)$, $g\in\auU(X)$, $r'\in\auU(X')$, is a constant times $\theta(\psrho(\lambda
(n'))\psZ_{\mathbb I}\psrho(S)\omega(g,r',k)F)$; $n'\mapsto\lambda(n')$ maps $[N_{j+1}]$ onto
$[\auHom(X',\ell)]$ with the Haar measure as image, so integration over $[N_{j+1}]$ gives
$\theta_{W_j\oplus W_{j+1}}(\psZ_{\mathbb J}\psrho(S)\omega(g,r',k)F)$. The elements $g,r',k$ fix $S$
and lie in the Levi subgroup of the stabiliser of $\mathbb J$; by (W2), $\psZ_{\mathbb J}\omega(\iota
g,\iota r',\iota k)=\chi_0(g,r',k)\,\omega(g,r',k)\psZ_{\mathbb J}$, where $\chi_0$ is a character of
the three determinants, built from $\upsilon$ and the ranks only. One has $\chi_0=1$ (Lemma~\ref{pk:chi0}). By (W2), $\psZ_{\mathbb J}
\psrho(S)F=u_\xi\otimes\psE_i(F_f)$ with a Fock vector $u_\xi$ depending only on $\xi$ and the
$P_j$. So
\[
\psW_{j+2}(\iota(k)g_{j+2})=\mathrm{const}\int_{[\auU(X')]}\int_{[\auU(X)]}l(g_\infty^{-1}\psi(g_f))\,
\theta\bigl(\omega(g,r',k)(u_\xi\otimes\Phi)\bigr)\,dg\,dr',
\]
and (a) is Lemma~\ref{ps:definite}. (b) is the same intertwining property of $\psZ_{\mathbb J,f}$,
with $\iota(b)S\iota(b')^{-1}=S$.
\end{proof}

\begin{remark}\label{ps:chi0}
The identity $\chi_0=1$ is Lemma~\ref{pk:chi0}; it uses that all $n_j$ are even and that $s_j$ is
odd for $j$ even.
\end{remark}

\begin{lemma}[Support]\label{ps:support}
Let $\varphi\in\auSram$, $g_{0,f}\in J_M$, $1\le j<m$, and let $\mathfrak B_j$ be the set of
$x\in G_j(\auA_f)$ for which there is $r\in R_j(\auA_f)$ such that $\Lambda'=rx\psM_j$ is the
direct sum of its intersections with the grades and $\mathsf q^{\,j}\Lambda'^-\subseteq\psM_j^-$,
$\aubar{\mathsf q}^{\,j}\psM_j^-\subseteq\Lambda'^-$, where $\mathsf q$ is an idele generating
$\prod_{w\in Q}w$ and $\Lambda'^-$ is the intersection with the span of the string. If
$\psW_j[g_0,\dots,g_{j-1}](h)\ne0$ then $h_f\in\mathfrak B_j$. Moreover $\mathfrak B_j=R_j(\auA_f)
\mathsf C_jK_j$ for a finite set $\mathsf C_j$ of elements that are $1$ at all places not dividing
$q$, and $\mathfrak B_j=R_j(\auA_f)K_j$ if $Q=\emptyset$.
\end{lemma}
\begin{proof}
Induction on $j$; put $\mathfrak B_0=J_M$. If $\psW_j(h)\ne0$ there is $r$ with
$\psW_{j-1}(rg_{j-1})\ne0$ and $c_T(\omega(rg_{j-1},h)\Upsilon_{j-1})\ne0$, $T=T_{j-1}$. By induction
$(rg_{j-1})_f=r_1g^\circ$ with $r_1\in R_{j-1}(\auA_f)$ and $L=g^\circ\psM_{j-1}$ graded with the
bounds. By the Iwasawa decomposition for the parabolic subgroup of $G_j$ with unipotent radical
$N_j$ (see (R1) of \S\ref{pk:ss-rest}; $K_{j,v}$ is hyperspecial because $E/E^+$ is unramified and
$\psM_j$ is unimodular), $h_f=n\,l\,k$ with $n\in N_j(\auA_f)$, $k\in K_j$ and $l$ preserving the grades; so
$\Lambda'=l\psM_j$ is graded. As in Proposition~\ref{ps:onestep}(1), $r_1$ and $n$ change $T$ by
maps of degree $\ge0$ and act by invertible operators after $\psZ_I$, so
$\psZ_{I,f}\psrho(T)\omega(g^\circ,lk)\Upsilon_{j-1,f}\ne0$. By Lemma~\ref{ps:level} and
\eqref{ps:eq-upsilon}, $\omega(g^\circ,lk)\Upsilon_{j-1,f}$ is a multiple of $\psrho(x_0)
\psvac_{\mathcal L}$ with $\mathcal L=lL'_{j-1}g^{\circ-1}$ and $x_0\in\alpha\auHom(L,\Lambda')$; and
$\psZ_{I,f}\psrho(y)\psvac_{\mathcal L}=0$ unless $y\in\mathcal L+I^\perp(\auA_f)$ (for $i\in
I\cap\mathcal L$, $\psrho(i)$ fixes $\psvac_{\mathcal L}$ up to a scalar and $\psZ_I\psrho(i)=\psZ_I$,
while $\psZ_I\psrho(i)\psrho(y)=[i,y]\psZ_I\psrho(y)\psrho(i)$; so $[i,y]=1$ on $I\cap\mathcal L$,
and $I\cap\mathcal L$ is the annihilator of $\mathcal L+I^\perp$). Hence $T\in\alpha\,\aubar{c_M
\mathsf q}^{-1}\auHom(L,\Lambda')+\auHom_{\ge0}$; since $L$, $\Lambda'$ are graded and $T$ is
homogeneous, $T(L)\subseteq\psfa\,\aubar{\mathsf q}^{-1}\Lambda'$, that is $T(L^+)\subseteq\psfa\,
\aubar{\mathsf q}^{-1}\Lambda'^-$. Now $L^+=(L^-)^\#\supseteq\aubar{\mathsf q}^{\,j-1}\psM_{j-1}^+$
by the induction hypothesis, and \eqref{ps:eq-TM} gives $\aubar{\mathsf q}^{\,j}\psM_j^-\subseteq
\Lambda'^-$. Dually $\Lambda'^-\subseteq\Lambda'^+\subseteq\mathsf q^{-j}\psM_j^+$, and
$\psM_j^+=\psM_j^-$ unless $j$ is even, where they differ by $u_jz/N$ at the places above $N$. Let
$j$ be even and $v$ a place of $E^+$ below $w\in\mathsf P$. The line $\Lambda'\cap E_vz$ contains
$\widehat\auO u_jz$ and is integral; if it is larger, it is unimodular, so the part of grade $0$ of
$\Lambda'_v$ is the sum of its intersections with $X_j$ and with $E_vz$, and the scalar
$\iota(\mathsf z_w)$ ($\mathsf z_w$ the idele $z_w$ of \S\ref{ss:au-spaces}) stabilises
$\Lambda'$; then $u=h_f^{-1}n\iota(\mathsf z_w)n^{-1}h_f\in J_j$, and by right invariance
$\psW_j(n\iota(\mathsf z_w^{\,e})n^{-1}h)=\psW_j(h)$ for all $e$. On the other hand, by
Lemma~\ref{ps:double}(a) applied $j/2$ times, $\psW_j(\gamma\iota(g)h)=\psi_j(\gamma)l(g_\infty^{-1}
(C\varphi)(g_f))$ for $\gamma\in N_j(\auA)$, $g\in\auU(X_j)(\auA)$, with $C$ a composite of lifts;
the central element $\mathsf z_w$ commutes with the lifts (it acts on $\auHom(X_i,X_{i+1})$
trivially, and the characters are unramified above $N$), so $\sum_{e=0}^{N-2}\psW_j(n\iota(\mathsf
z_w^{\,e})n^{-1}h)=\psi\cdot l((C\sum_e\mathsf z_w^{\,e}\varphi)(1))=0$ for ramified $\varphi$
(see (R2) of \S\ref{pk:ss-rest}: $n\iota(\mathsf z_w^{\,e})n^{-1}=n'\iota(\mathsf z_w^{\,e})$ with
$n'\in N_j(\auA)=N_j^{\ge2}(\auA)$ and $\psi_j(n')=1$). So $(N-1)\psW_j(h)=0$, a contradiction; hence $\Lambda'\cap E_vz=\widehat\auO u_jz$
and $\mathsf q^{\,j}\Lambda'^-\subseteq\psM_j^-$. The last assertion is Lemma~\ref{pf:orbits}.
\end{proof}

\begin{lemma}[Finite data]\label{ps:fin}
There is $C_1\in\C$, independent of $\vartheta$, with the following property. Let
$q=\auNm Q_\vartheta$ and let $\mathsf y=(\mathsf y_0,\dots,\mathsf y_m)$, $\mathsf y_j\in
G_j(\auA_f)$, be integral at all places not dividing $q$, with $\mathsf y_m=1$. Then
\[
\Phi(\mathsf y):=\bigotimes_{i<m/2}\psE_i\bigl(\omega(\mathsf y_{2i},\mathsf y_{2i+1})\Upsilon_{2i,f}
\otimes\omega(\mathsf y_{2i+1},\mathsf y_{2i+2})\Upsilon_{2i+1,f}\bigr)=C_1\sum_xr_{\mathsf y}(x)\,
\psGlue(\phi^\circ_x),
\]
a finite sum over tuples $x$ as in (C1), where $B^{2m}r_{\mathsf y}(x)$, $B=M^{fn_m^2}$, is an
algebraic number integral at the primes not dividing $q$, and $\sum_x|\sigma r_{\mathsf y}(x)|\le1$
for all $\sigma\in\auAut(\C)$ if $Q_\vartheta=\emptyset$.
\end{lemma}
\begin{proof}
It suffices to treat one pair $i$, with $j=2i$ and the notation of the proof of
Lemma~\ref{ps:double}. \emph{(1)} By \eqref{ps:eq-upsilon}, (W1), (W4), each
$\omega(\mathsf y_j,\mathsf y_{j+1})\Upsilon_{j,f}$ is a combination $\sum_ar(a)\psrho^\circ(a)
\psvac_{L_j}$ with $\aubar{\pi_M}a\in L_j$ outside $q$: at the places outside $q$ the $\mathsf y$
stabilise $L_j$ and permute the translates up to roots of unity and unit values of characters; at
the places of $q$, $\omega(g)\psvac_L$ is a multiple of $\psvac_{gL}=\sum_{r\in gL/(L\cap gL)}
\zeta_r\psrho^\circ(r)\psvac_L$ with a factor whose denominator divides a power of $q$. The
coefficients $r(a)$ are roots of unity times $\prod_w\#_w^{-1}$ outside $q$; $\prod_{w\text{
chosen}}\#_w$ divides $B$; and if $Q=\emptyset$ then $\sum|\sigma r(a)|\le1$. Under the
identification (W6) of the first factor, $w\mapsto w^\dagger$ carries $L_j$ to
$L_j^\dagger=\aubar\alpha\auHom(\psM_{j+1},\psM_j^-)=\alpha\auHom(\psM_{j+1},u\psM_j^-)$, $u=\aubar a/a$
(Remark~\ref{pf:rem-flip-lattice}; $u_v=1$ at the places not above $q_1$), and $\mathrm{Fl}$, normalised
locally to carry $\psvac_{L_{j,v}}$ to $\psvac_{L^\dagger_{j,v}}$ (possible by that remark), carries
translates to translates up to roots of unity. So the tensor of two translates is a root of unity
times a translate $\psrho^\circ(a)\psvac_{\mathbb L}$,
$\mathbb L=L_j^\dagger\oplus L_{j+1}=\alpha\auHom(\psM_{j+1},u\psM_j^-\oplus\psM_{j+2})$.

\emph{(2)} $\psE_i(\psrho^\circ(a)\psvac_{\mathbb L})=\psZ_{\mathbb J,f}\psrho(S+a)\psvac_{\mathbb L}$
up to a root of unity. As in the proof of Lemma~\ref{ps:support}, this vanishes unless
$S+a\in\mathbb L+\mathbb J^\perp$, and then it is a root of unity times
$\psrho^\circ(a_0)\psZ_{\mathbb J,f}\psvac_{\mathbb L}$, $a_0$ the image in $W_j\oplus W_{j+1}$ of
an element of $(S+a+\mathbb L)\cap\mathbb J^\perp$. The vector $\psZ_{\mathbb J,f}\psvac_{\mathbb L}$
is not zero and is fixed, up to scalars, by $\psrho(\bar l)$, $l\in\mathbb L\cap\mathbb J^\perp$; so
it is $C\,\psvac_{\mathbb L_0}$ with $\mathbb L_0$ the image of $\mathbb L\cap\mathbb J^\perp$ (a
self-dual lattice), and $C$ depends on the pair $i$ only.

\emph{(3) The lattice $\mathbb L_0$.} Since $\psM_{j+1}$ is the sum of its parts in $X'$ and in
the string, $\mathbb L_0=\alpha\auHom(\widehat\auO^{\,n'},\Lambda_0)$ with $\Lambda_0\subset
(X^-\oplus X'')(\auA_f)$ the image of the vectors of $u\psM_j^-\oplus\psM_{j+2}$ without components
of negative grade and with component of grade $0$ in $X^-\oplus X_{j+2}\oplus\ell$. At the places
above $q_1$ (where $u_v\ne1$ is possible) there is no prime above $N$ and no chosen place
(Lemma~\ref{pgx:dict}(a)), the parts of $\psM_j$, $\psM_{j+2}$ in $X_j$, $X_{j+2}$ are the standard
lattices (and for $j+2=m$ the factor $u\,u_{m-2}$ of the vector $z/\pi_N$ is a unit) and $\Lambda_{0,v}$ is
$u_v$ times the standard lattice of $X^-$ plus the standard lattice of $X''$, so that $\mathbb
L_{0,v}=(L^X_{j,v})^\dagger\oplus L^X_{j+1,v}$ is exactly the image of $L^X_{j,v}\oplus L^X_{j+1,v}$
under $w\mapsto w^\dagger$ on the first factor. Elsewhere $u_v=1$ and, by
\eqref{ps:eq-lattice}, with $y,y''$ the last vectors of $X,X''$: outside $N$ and the chosen
places, $\Lambda_0$ is the standard lattice; at the chosen places it is the standard lattice
except that for $j+2=m$ the vector $y''=z/\pi_N$ is multiplied by the norm-one idele $u_{m-2}$
(whose denominator divides $\aubar{\pi_M}$) (see (R3) of \S\ref{pk:ss-rest}: at the chosen places
$\psM_j$ is the orthogonal sum of its parts in $X_j$ and in the string, so the image is standard for
$j+2<m$ whatever $u_j,u_{j+2}$ are); and at a place $v=w\aubar w$, $w\in\mathsf P$, where
$u_j=u_{j+2}=1$ and $Sf'_{p-1}=c(z+z'')$ (resp.\ $\psM_{m-2}\ni(y+z/\pi_N)/\aubar{\pi_N}$, $z$),
\[
\Lambda_{0,v}=\bigl\{\textstyle\sum a_ie_i+\beta y+\sum a''_ie''_i+\beta''y'':\ \beta''\equiv\beta\
(\varpi_w)\text{ at }w,\ \ (\beta,\beta'')\in\auO_{\bar w}^2+\varpi_{\bar w}^{-1}\auO_{\bar w}(1,1)
\text{ at }\aubar w\bigr\},
\]
the neighbour of the standard lattice along the isotropic vector $y+y''$. Hence
$\mathbb L_0/(\mathbb L_0\cap L^X)$, $L^X=L^X_j\oplus L^X_{j+1}$ (read in $\auHom(X',X^-\oplus X'')$ through $w\mapsto w^\dagger$ on the first factor), is represented by the sums
$t+\mathsf p_{\mathsf y}$, where $t$ runs over a finite set $\mathsf T_i$ of elements vanishing
above $N$ (trivial unless $j+2=m$, of order dividing $B$), and $\mathsf p_{\mathsf y}$ is the map
$x'\mapsto-\alpha N^{-1}\langle\mathsf y,x'\rangle(y+y'')$ at the $\aubar w$; under
$W_j\oplus W_{j+1}=\auHom(X',X^-\oplus X'')$, $w\mapsto w^\dagger$, $\mathsf p_{\mathsf y}$ is the
pair $(p_{\mathsf y},p'_{\mathsf y})$ of \S\ref{ps:ss-glue}. The $\mathsf p_{\mathsf y}$ are
supported at the $\aubar w$, so they are mutually isotropic, $\psrho^\circ=\psrho$ on them, and
they commute with the $\psrho^\circ(t)$. Therefore $\sum_{t,\mathsf y}\zeta_t\psrho^\circ(t)
\psrho^\circ(\mathsf p_{\mathsf y})\psvac_{L^X}$, for suitable roots of unity $\zeta_t$, is fixed up
to characters by $\mathbb L_0$ and is not zero (the translates are independent): it is a multiple
of $\psvac_{\mathbb L_0}$, and
\[
\psZ_{\mathbb J,f}\psvac_{\mathbb L}=C'\,\frac1{\#\mathsf T_i}\sum_{t\in\mathsf T_i}\zeta_t\,
\psGlue_i\bigl(\phi^\circ_t\bigr),\qquad\phi^\circ_t=\psrho^\circ(t)(\psvac^\circ\otimes\psvac^\circ).
\]
\emph{(4)} In (2), $a_0$ is integral for $\mathbb L_0$ above $N$, so $\psrho^\circ(a_0)
\psvac_{\mathbb L_0}$ is a root of unity times $\psrho^\circ(a_0')\psvac_{\mathbb L_0}$ with $a_0'$
vanishing at all places not dividing $Mq$; $\psrho^\circ(a_0')$ commutes with the
$\psrho^\circ(\mathsf p_{\mathsf y})$, and $\psrho^\circ(a'_0)\psrho^\circ(t)=\zeta\psrho^\circ
(a_0'+t)$. Altogether $\psE_i(\dots)$ is $C'$ times a combination of the $\psGlue_i(\phi^\circ_x)$
with coefficients that are roots of unity times products of the $r(a)$ and of $1/\#\mathsf T_i$;
the product of the denominators over the $m/2$ pairs divides $B^{m+1}$, hence $B^{2m}$ (see (R4)
of \S\ref{pk:ss-rest}), and the sizes multiply.
\end{proof}

\begin{proof}[Proof of Proposition~\ref{au:seesaw}]
Fix $g_0\in G_0(\R)$ and $\mathsf t$. Let $\eta$ be a form on $G_m$ of level $J_m$ and weight
$\lambda^\circ$, $z_m$ the complex conjugate of the corresponding function on $[G_m]$, and
$z_j=\Theta^{Y,*}_{j,\vartheta}z_{j+1}$. By Fubini (the forms are rapidly decreasing and the theta
kernels of moderate growth) $\int_{[G_m]}z_m\,\iota_\vartheta(\auP\varphi)=\int_{[G_0]}z_0\,\auP
\varphi$. \emph{Step 1.} The function $z_0$ has level $J_0$ (Lemma~\ref{ps:level}); its trace
$\sum_{\gamma\in J_M/J_0}z_0(\,\cdot\,\gamma)$ is the conjugate of an element of $\auAz$, so by the
definition of $\auP$
\[
[J_M:J_0]\int_{[G_0]}z_0\,\auP\varphi=\sum_{\gamma\in J_M/J_0}\int_{[R_0]}\psW_\varphi(r)z_0(rg_0
\gamma)\,dr .
\]
\emph{Step 2.} Put $\Pi_j(x_0,\dots,x_j)=\int_{[R_j]}\psW_j[x_0,\dots,x_{j-1}](rx_j)z_j(rx_j)\,dr$
for $j<m$. By Proposition~\ref{ps:onestep} with $F=\omega(x_j,1)\Upsilon_j$ and
$\psW=\psW_j(\,\cdot\,x_j)$,
\begin{align*}
\Pi_j(x_0,\dots,x_j)&=[K_j:J_j]\int_{R_{j+1}(E^+)\backslash G_{j+1}(\auA)}\psW_{j+1}[x_0,\dots,x_j]
\,z_{j+1}\\&=[K_j:J_j]\int_{R_{j+1}(\auA)\backslash G_{j+1}(\auA)}\Pi_{j+1}(x_0,\dots,x_{j+1})\,dx_{j+1}
\end{align*}
for $j+1<m$, absolutely convergent, and for $j=m-1$ the first integral is over $[G_m]$.
\emph{Step 3 (the top).} By Lemma~\ref{ps:double}(a), applied $m/2$ times starting from
$\psW_\varphi(\iota(g))=\mathsf t(g_\infty^{-1}\varphi(g_f))$, for $h\in G_m(\auA)$
\[
\psW_m[x_0,\dots,x_{m-1}](h)=l_\xi\bigl(h_\infty^{-1}(C_{\vartheta,\Phi(\mathsf y)}\varphi)
(\operatorname{diag}(h_f,1))\bigr),
\]
where $\xi$, $\mathsf y$ are the archimedean and finite parts of $(x_0,\dots,x_{m-1},1)$ and
$l_\xi$ is a linear form on $W_{\lambda(X_m)}$ depending on $\xi$, $g_0$, $\mathsf t$ only. The
group $G_m(\R)$ is compact and $W_{\lambda^\circ}$ occurs once in $W_{\lambda(X_m)}$ restricted to
it; by Schur orthogonality there is a linear form $\mathsf c$ on the dual of $W_{\lambda(X_m)}$,
independent of everything else, with $\int_{[G_m]}\psW_mz_m=\mathsf c(l_\xi)\,\langle
\operatorname{res}C_{\vartheta,\Phi(\mathsf y)}\varphi,\eta\rangle$ (the factor
$\prod\tau(\aubar{\pi_N})^{n_\tau}$ and the dimension are put in $\mathsf c$). \emph{Step 4 (the
outer integrals).} Thus $\Pi_{m-1}=[K_{m-1}:J_{m-1}]\,\mathfrak a_{m-1}(\xi)\,S_{m-1}(\mathsf y)$ with
$\mathfrak a_{m-1}=\mathsf c(l_\xi)$ and $S_{m-1}(\mathsf y)=\langle\operatorname{res}C_{\vartheta,
\Phi(\mathsf y)}\varphi,\eta\rangle$. Suppose $\Pi_{j+1}=[K_{j+1}:J_{j+1}]\mathfrak a_{j+1}(\xi)
S_{j+1}(\mathsf y)$, where $\mathfrak a_{j+1}$ does not depend on $\vartheta,\varphi,\eta$. The
function $\Pi_{j+1}$ is left invariant under $R_{j+1}(\auA)$ in $x_{j+1}$, right invariant under
$J_{j+1}$, and vanishes unless $x_{j+1,f}\in\mathfrak B_{j+1}$ (Lemma~\ref{ps:support}; for
$x_{0,f}\in J_M$). Put $H_{j+1}=R_{j+1}(\auA_f)\cap K_{j+1}$. The group
$R_{j+1}(\auA_f)$ has finitely many orbits in $\mathfrak B_{j+1}/J_{j+1}=R_{j+1}(\auA_f)\mathsf
C_{j+1}K_{j+1}/J_{j+1}$, because $\mathsf C_{j+1}$ and $K_{j+1}/J_{j+1}$ are finite. In each orbit
choose a point of $\mathsf C_{j+1}K_{j+1}/J_{j+1}$, and let $\Omega_{j+1}$ be the union of the
$H_{j+1}$-orbits of the chosen points: a finite subset of $\mathfrak B_{j+1}/J_{j+1}$, stable under
left multiplication by $H_{j+1}$, contained in $H_{j+1}\mathsf C_{j+1}K_{j+1}/J_{j+1}$, which every
$R_{j+1}(\auA_f)$-orbit meets in exactly one $H_{j+1}$-orbit. (For $Q=\emptyset$, $\mathsf
C_{j+1}=\{1\}$ and $\Omega_{j+1}$ is a subset of $K_{j+1}/J_{j+1}$, in general not contained in
$H_{j+1}J_{j+1}/J_{j+1}$.) For $t\in\Omega_{j+1}$ let $m(t)=[\mathsf S_t:\mathsf S_t\cap K_{j+1}]$,
$\mathsf S_t$ the stabiliser of $t$ in $R_{j+1}(\auA_f)$. With the measures giving volume $1$ to
$K_{j+1}$ and to $H_{j+1}$, every function $\mathsf b$ on $R_{j+1}(\auA_f)\backslash
\mathfrak B_{j+1}/J_{j+1}$ has $\int\mathsf b=[K_{j+1}:J_{j+1}]^{-1}\sum_{t\in\Omega_{j+1}}\mathsf
b(t)/m(t)$: the orbit of $t$ has volume $[K:J]^{-1}\operatorname{vol}(\mathsf S_t)^{-1}$ (the
groups $G_{j+1}(\auA_f)$ and $R_{j+1}(\auA_f)$ are unimodular and $tJ_{j+1}$ has volume
$[K:J]^{-1}$), and $\operatorname{vol}(\mathsf S_t)^{-1}=[H_{j+1}:\mathsf S_t\cap K_{j+1}]/m(t)$ is
the number of elements of $\Omega_{j+1}$ in the orbit, which is the number of elements of
$H_{j+1}t$, divided by $m(t)$. Therefore
\[
\Pi_j=[K_j:J_j]\Bigl(\int_{R_{j+1}(\R)\backslash G_{j+1}(\R)}\mathfrak a_{j+1}\,d\xi_{j+1}\Bigr)
\sum_{t\in\Omega_{j+1}}\frac{S_{j+1}(\mathsf y_0,\dots,\mathsf y_j,t)}{m(t)}
\]
whenever the sum is not zero, the integral being then absolutely convergent by Step~2; define
$\mathfrak a_j(\xi_0,\dots,\xi_j)$ as this integral when it converges absolutely and as $0$
otherwise, and $S_j$ as the sum. Then $\Pi_j=[K_j:J_j]\mathfrak a_jS_j$ in all cases. By descending
induction and Step~1,
\[
\int_{[G_m]}z_m\,\iota_\vartheta(\auP\varphi)=c_0\,\bigl\langle\operatorname{res}C_{\vartheta,
\Phi_\vartheta}\varphi,\eta\bigr\rangle,\quad\Phi_\vartheta=\sum_{\gamma\in J_M/J_0}\,\sum_{t_1\in
\Omega_1}\!\cdots\!\sum_{t_{m-1}\in\Omega_{m-1}}\frac{\Phi(\gamma,t_1,\dots,t_{m-1},1)}{m(t_1)\cdots
m(t_{m-1})},
\]
with $c_0=[K_0:J_M]\,\mathfrak a_0(g_0)$ independent of $\vartheta$, $\varphi$, $\eta$. As $\eta$
is arbitrary, $\iota_\vartheta(\auP\varphi)=c_0\operatorname{res}C_{\vartheta,\Phi_\vartheta}\varphi$
for all $\varphi\in\auSram$.

\emph{Step 5 (the chain datum).} The elements of $\Omega_j$ are integral outside $q$, and $m(t)$
divides a power of $q$: $\mathsf S_t$ is a product of local groups, contained in $K_j$ outside $q$,
and at $v\mid q$ it is conjugate to a subgroup of the pro-$p$ group $J_{j,v}$. If
$Q=\emptyset$ then $\Omega_j\subseteq K_j/J_j$, $m=1$ and $\#\Omega_j\le B$. Indeed, for
$Q=\emptyset$ the group $J_j$ is the kernel of the reduction $K_j\to\auGL(\psM_j/\pi_M\psM_j)$, so
$[K_j:J_j]\le\#\operatorname{End}(\psM_j/\pi_M\psM_j)=(\#\auO_E/\pi_M\auO_E)^{(\dim Y_j)^2}=
M^{f(\dim Y_j)^2}$, since $|\sigma\pi_M|^2=M$ for every embedding $\sigma$; and
$\dim Y_j\le\dim Y_m=n_m-1$, the frame of $Y_j$ being the beginning of that of $Y_{j+1}$. Hence
$\#\Omega_j\le[K_j:J_j]\le B$, and likewise $[J_M:J_0]\le[K_0:J_0]\le B$, so $\Phi_\vartheta$ is a
sum of at most $B^m$ terms $\Phi(\mathsf y)/(m(t_1)\cdots m(t_{m-1}))$ with all $m(t_j)=1$. By
Lemma~\ref{ps:fin}, applied to each term, $\Phi_\vartheta=C_1\sum_xr_\vartheta(x)\psGlue(\phi^\circ_x)$ with
$B^{2m}r_\vartheta$ integral outside $q$ and $\sum_x|\sigma r_\vartheta(x)|\le B^{m}$ if
$Q=\emptyset$. Put $\Psi=B^{2m}r_\vartheta$. Then (C1)--(C3) hold, since $B^{3m}=M^{3fmn_m^2}$. For
(C4), let $b_j\in J_{X_j}$, $b_m=1$. The element $\iota(b_j)$ lies in $R_j(\auA_f)\cap K_j$ and is
$\equiv1$ at the chosen places: for $j$ odd this is clear from \eqref{ps:eq-lattice}; for $j$ even,
the last row of $b_j$ is $\equiv e_{\mathrm{last}}$ modulo $\pi_N$ and, $b_j^{-1}=b_j^*$ being in
$J_{X_j}$ too, its last column is $\equiv e_{\mathrm{last}}$ modulo $\aubar{\pi_N}$; so
$b_je_i\in\widehat\auO^{\,n_j-1}\oplus\pi_N\widehat\auO y\subset\psM_j$ and $(b_j-1)y/\aubar{\pi_N}
\in\widehat\auO^{\,n_j-1}\oplus\pi_N\widehat\auO y$, whence $\iota(b_j)\psM_j=\psM_j$. By
Lemma~\ref{ps:double}(b), $\bigotimes_j\omega(b_j,b_{j+1})\Phi(\mathsf y)=\Phi((\iota(b_j)\mathsf
y_j)_j)$; and $t\mapsto\iota(b_j)t$ permutes $\Omega_j$ (and $J_M/J_0$) and preserves $m$. So
$\Phi_\vartheta$ is invariant, and $\psGlue(\sum\Psi(x)\phi^\circ_x)=B^{2m}C_1^{-1}\Phi_\vartheta$
if $C_1\ne0$. If $c_0C_1\ne0$ put $c_{\mathrm{ss}}=c_0C_1B^{-2m}$: then $\iota_\vartheta(\auP
\varphi)=c_{\mathrm{ss}}B_{\vartheta,\Psi}\varphi$. If $c_0C_1=0$ then $\iota_\vartheta(\auP\varphi)
=0$ for all $\vartheta$ and $\varphi\in\auSram$, and the statement holds with $c_{\mathrm{ss}}=1$
and $\Psi=0$.
\end{proof}


%% ===== proofs/A-seesaw-complete.tex =====
% proofs/A-seesaw-complete.tex
% Complements to proofs/A-seesaw-size.tex and proofs/A-divisibility.tex: proofs of ps:definite,
% of chi_0 = 1 (ps:chi0), of the steps of ps:onestep, of ps:branching(c), of ps:conj and of
% pa:kernel, as far as stated in \S pk:sec. Labels of this file start with pk:.
% Only amsmath/amssymb/amsthm are used.
\providecommand{\auO}{\mathcal O}
\providecommand{\auU}{\mathrm U}
\providecommand{\auGL}{\mathrm{GL}}
\providecommand{\auHom}{\operatorname{Hom}}
\providecommand{\auA}{\mathbb A}
\providecommand{\auF}{\mathbb F}
\providecommand{\auPl}{\Phi}
\providecommand{\auPa}{\Phi_a}
\providecommand{\auS}{\mathcal S}
\providecommand{\auSram}{\mathcal S^{\mathrm{ram}}}
\providecommand{\auAz}{\mathcal A_0}
\providecommand{\auNm}{\operatorname{Nm}}
\providecommand{\auAut}{\operatorname{Aut}}
\providecommand{\aubar}[1]{\overline{#1}}
\providecommand{\aurd}{\mathrm{rd}_E}
\providecommand{\psM}{\mathsf M}
\providecommand{\psHS}{\mathcal H}
\providecommand{\psW}{\mathcal W}
\providecommand{\psZ}{Z}
\providecommand{\psvac}{\mathbf v}
\providecommand{\psrho}{\varrho}
\providecommand{\psE}{\mathcal E}
\providecommand{\pao}{\mathfrak o}
\providecommand{\pam}{\mathfrak m}
\providecommand{\pkI}{\mathsf I}
\providecommand{\pkm}{\mathsf m}
\providecommand{\pkev}{\operatorname{ev}}
\providecommand{\pktr}{\operatorname{tr}}
\providecommand{\pkLW}{\mathrm{LW}}

\subsection{The upper bound: definite lifts, unfolding and kernels}\label{pk:sec}

This subsection follows \S\ref{ps:sec} and uses its notation. It proves the statements of \S\ref{ps:sec}
whose proofs are not given there, and Lemma~\ref{pa:kernel}, as follows.
\begin{itemize}
\item \emph{Proved here}, from Lemma~\ref{ps:weil} (with the conventions of
\S\ref{pk:ss-conv}) and the classical results named at the place of use: the definite lift
(Lemma~\ref{ps:definite}; Proposition~\ref{pk:definite}); the identity $\chi_0=1$ of
Remark~\ref{ps:chi0} (Lemma~\ref{pk:chi0}); Proposition~\ref{ps:onestep}, with complete proofs of
its steps (3), (5) and (6) (\S\ref{pk:ss-onestep}); Lemma~\ref{ps:branching}(c)
(Proposition~\ref{pk:branching}); Lemma~\ref{pa:kernel}(K0) (Proposition~\ref{pk:K0}).
\item \emph{Proved here from published results}: Lemma~\ref{pa:kernel}(K2), and (K1) at the
primes $\ell\nmid2q_1Md$ (all primes $\ell\ne d$ in \S\ref{pf:sec}) (Propositions~\ref{pk:std}, \ref{pk:K12}), from the properties (C1), (C2),
(C3)(i), (C4) of one curve with complex multiplication (Proposition~\ref{pk:curve}).
\item \emph{Proved in \S\ref{pf:sec}}: Proposition~\ref{pk:curve}(C3)(ii) (Theorem~\ref{pf:recip}), hence
Lemma~\ref{ps:conj}; the integrality in Lemma~\ref{pa:kernel}(K1) at the primes $\ell\mid2q_1M$ and at
$\ell=N$ (Corollary~\ref{pf:K1}); the fact (W6) (Proposition~\ref{pf:flip}).
\end{itemize}


\subsubsection{Conventions}\label{pk:ss-conv}

We fix the following conventions for \S\ref{ps:ss-weil} and this subsection; they are the ones
for which the statements of \S\ref{pa:ss-prod} on the linear forms $u^\dagger$ hold.

\emph{The additive character and $\jmath$.} $\psi=\psi_\Q\circ\pktr_{E/\Q}$ with
$\psi_{\Q,\infty}(t)=e^{-2\boldsymbol\pi it}$ and $\psi_{\Q,p}(t)=e^{2\boldsymbol\pi i\{t\}_p}$,
where $\boldsymbol\pi=3.14\ldots$; and $\jmath=-1/\sqrt{-d}$, so that $\tau(\jmath)=i/\sqrt d$
for every $\tau\in\auPl$ (the elements of $\auPl$ extend the given embedding of $K$). The idele
$\alpha$ is $1/P'(\zeta_{q_1})$ at the places of $E$ dividing
$\omega=(1+\sqrt{-d})/2$, and $1$ elsewhere, where $P$ is the minimal polynomial of $\zeta_{q_1}$
over $K$; every prime dividing $q_1$ divides $\auNm\omega=(1+d)/4$ and splits in $K$, so that
$\mathfrak a_0=E\cap\alpha\widehat\auO_E$ satisfies $\mathfrak a_0\aubar{\mathfrak a_0}=
\mathfrak d_{E/K}^{-1}$ and $\auO_E\subseteq\mathfrak a_0\subseteq\mathfrak d_{E/K}^{-1}$. The
character $\upsilon$ is unramified at \emph{every} finite place: $\upsilon(x)=\prod_\tau
x_\tau/|x_\tau|$ if the finite part of $x$ is a unit. Its existence, and the properties of it that
are used ($\upsilon_\tau(z)=z/|z|$; $\upsilon_{\aubar w}=\upsilon_w^{-1}$ on $E^{+\times}_v$ at a split
place; $\upsilon|\cdot|^{-1/2}$ algebraic), are Lemma~\ref{ps:upsilon}.

\emph{Fock coordinates.} For $\tau\in\auPl$ and $x,y\in W\otimes_{E^+,\tau}\R$,
\[
\psi_\tau\bigl(\tfrac12\jmath\langle y,x\rangle\bigr)=\exp\bigl(-2\boldsymbol\pi i\,\Rel(i\,d^{-1/2}
\langle y,x\rangle)\bigr)=\exp\bigl(i\,c_\psi^2\,\Iml\langle y,x\rangle\bigr),\qquad
c_\psi=\sqrt{2\boldsymbol\pi/\sqrt d}\,.
\]
Thus, in the coordinates $x_\tau=c_\psi\,\tau(w)$, $\beta_\infty(x,y)=\exp(+i\,\Iml\sum_\tau
\pktr(y_\tau^*x_\tau))$, the sign used in \S\ref{ps:ss-weil}. The Fock
space $\mathfrak F$ of $W=\auHom(X,X')$ consists of the entire functions $F$ of
$z=(z_\tau)_\tau\in\prod_\tau M_{n'\times n}(\C)$ with $\int|F|^2e^{-|z|^2}<\infty$; the variable
$z$ is the \emph{complex conjugate} of the coordinate $x$, and
\begin{equation}\label{pk:eq-fock}
(\psrho_\infty(y)F)(z)=e^{z\cdot y-|y|^2/2}F(z-\aubar y),\qquad z\cdot y=\sum_{\tau}\pktr(y_\tau^{\,t}
z_\tau),
\end{equation}
\[
(\omega_\infty(g,h)F)(z)=F(h^{t}z\,\aubar g)\qquad(g\in\auU(X)(\R),\ h\in\auU(X')(\R)),
\]
times the determinant twists (there is no further multiplier, $W$ being positive definite).
Indeed the exponents of $\psrho_\infty(x)\psrho_\infty(y)$ and of $\psrho_\infty(x+y)$ differ by
$-\aubar x\cdot y+\Rel(\aubar x\cdot y)=i\,\Iml(y^*x)$, and for a unitary matrix $u$ acting on $W
\otimes\R$ one has $(\aubar u^{-1}z)\cdot y=z\cdot(uy)$, so that $F\mapsto F(\aubar u^{-1}z)$
conjugates $\psrho_\infty(y)$ into $\psrho_\infty(uy)$; for $u\colon w\mapsto hwg^{-1}$,
$\aubar u^{-1}z=h^tz\aubar g$. With $E_z(w)=e^{w\cdot z}=e^{|z|^2/2}(\psrho_\infty(z)1)(w)$ the
theta function of $f\in\psHS_f(W)$ is
\begin{equation}\label{pk:eq-theta}
\Theta_f(z)=\theta(E_z\otimes f)=e^{|z|^2/2}\,\theta\bigl(\psrho_\infty(z)(1\otimes f)\bigr),
\end{equation}
which agrees with \eqref{ps:eq-theta}, and \eqref{ps:eq-bargmann} reads
\begin{equation}\label{pk:eq-bargmann}
Q(\partial)[\Theta_f(c\,\cdot)](0)=\theta\bigl(Q(c\,\cdot)\otimes f\bigr)=\int Q(c\,\bar y)\,
e^{-|y|^2/2}\,\theta\bigl(\psrho_\infty(y)(1\otimes f)\bigr)\,d\mu(y).
\end{equation}
The first equality follows from
\eqref{pk:eq-theta} and (W0) below, because $E_z=\sum_\alpha z^\alpha w^\alpha/\alpha!$ in
$\mathfrak F$.

\emph{The action of $\Gamma$ on the weights.} With these coordinates the rational points act on
$W_\lambda$ through the conjugate embeddings: $\gamma\cdot\xi=\rho_\lambda(\aubar{\tau(\gamma)})\xi$,
and for a real point $g_\infty\cdot\xi=\rho_\lambda(\aubar{g_\infty})\xi$ (see
\eqref{pk:eq-schuract}). The sentence ``$\Gamma$ acts on $W_\lambda$ through the embeddings $\tau$''
of \S\ref{ss:au-spaces}, and the gauge there, refer to the embeddings $\aubar\tau=\tau\circ c$; this
is harmless, the set of data being stable under $\auPl\leftrightarrow\aubar\auPl$.

\emph{Three additions to Lemma~\ref{ps:weil}.}
\begin{itemize}
\item[(W0)] $\theta$ is defined on $\mathfrak F^{\mathrm{sm}}\otimes\psHS_f(W)$, where
$\mathfrak F^{\mathrm{sm}}$ is the space of smooth vectors, and is continuous in the first factor
for its Fr\'echet topology. In particular, for a compact subgroup $C$ of $\auU(W)(\R)$ and a
continuous function $a$ on $C$, $\theta\bigl(\int_Ca(c)\omega(c)F\,dc\bigr)=\int_Ca(c)\theta(\omega(c)F)\,dc$.
\item[(W2$'$)] In (W2) the character is $\chi_I(p)=\bigl(\upsilon|\cdot|^{-1/2}\bigr)(\det p|_I)^{-1}=
\upsilon(\det p|_I)^{-1}|\det p|_I|^{1/2}$, where $\det p|_I\in\auA_E^\times$ is the determinant of
$p$ on the $\auA_E$-module $I(\auA)$ and $|\cdot|$ is the idele norm of $E$. (This is Proposition~\ref{pwf:W2}(c).)
\item[(W6)] Let $V,V'$ be hermitian spaces of ranks $n,n'$ and $V^-$ the space $V$ with the opposite
form. The map $w\mapsto w^\dagger$ is an $E$-antilinear isomorphism $\auHom(V,V')\to\auHom(V',V^-)$
with $\langle w_1^\dagger,w_2^\dagger\rangle=-\langle w_2,w_1\rangle$; hence it preserves $\beta$,
and there is an isomorphism $\mathrm{Fl}\colon\psHS(\auHom(V,V'))\to\psHS(\auHom(V',V^-))$,
unique up to a scalar, with $\mathrm{Fl}\,\psrho(w)=\psrho(w^\dagger)\mathrm{Fl}$. It satisfies
\[
\mathrm{Fl}\circ\omega_{\auHom(V,V')}(g,g')=\upsilon(\det g)^{n'}\upsilon(\det g')^{-n}\cdot
\omega_{\auHom(V',V^-)}(g',g)\circ\mathrm{Fl}\qquad(g\in\auU(V)(\auA),\ g'\in\auU(V')(\auA)),
\]
and $\theta\circ\mathrm{Fl}$ is a multiple of $\theta$. The map $\mathrm{Fl}$ carries the vacuum of
$\alpha\auHom(\psM,\psM')$ to a multiple of the vacuum of its image $\aubar\alpha\auHom(\psM',\psM)$, which
is $\omega(1,u)$ applied to the vacuum of $\alpha\auHom(\psM',\psM)$, up to a scalar, with the norm-one
central idele $u=\aubar a/a$ of $\auU(V)$, trivial at the places not above $q_1$
(Remark~\ref{pf:rem-flip-lattice}).
\end{itemize}
(W0) is part of the construction of $\theta$ in \cite[no.~39, p.~189, and Th\'eor\`eme~6, p.~193]{WeilActa}. For (W6): the first assertions
are immediate ($(w^\dagger)^\dagger=-w$ for the pair $(V',V^-)$, and $\psi(\aubar a)=\psi(a)$); the
two sides of the displayed identity are operators with the same commutation relations (W1), so
they differ by a character of $\auU(V)(\auA)\times\auU(V')(\auA)$, and the assertion is that this
character is $\upsilon(\det g)^{n'}\upsilon(\det g')^{-n}$. 
(W6) is proved in Proposition~\ref{pf:flip}.

\emph{The lifts of the upper chain.} The measure on $\auU(X_j)(\auA_f)$ gives volume $1$ to the
level $J_{X_j}$, and for $\Phi\in\psHS_f(W_j)$
\begin{equation}\label{pk:eq-lift}
(\Theta^X_j(\Phi)\varphi)(g')=\int_{\Gamma\backslash\auU(X_j)(\auA_f)}\theta_j(\Phi;g,g')\varphi(g)\,dg;
\end{equation}
this is the class sum over $\Gamma\backslash\auU(X_j)(\auA_f)/J$ divided by $[J_{X_j}:J]$, for
every free level $J\subseteq J_{X_j}$ fixing $\varphi$ and $\omega_{\mu_{j+1}}(\cdot,1)\Phi$; it
makes precise the words ``for a free level $J$ fixing $\Phi$'' after \eqref{ps:eq-kernel}.

\emph{The constants.} $c_\psi=\sqrt{2\boldsymbol\pi/\sqrt d}$. Let $\Omega_0$ be a complex number
such that the lattice $\Lambda=\Omega_0\auO_K$ admits a point $t$ as in \S\ref{pa:ss-curve} with
$\wp''(t)=2$ (the \emph{period}) and $t$ admissible; let $\mathsf u$ be the number of
Proposition~\ref{pk:curve}(C4) for this lattice (the \emph{Chowla--Selberg unit}); let
$\varpi_{\mathrm{vac}}=(\upsilon(\alpha)|\alpha|^{-1/2})^{-1/f}$ (principal branch). Then
\begin{equation}\label{pk:eq-s0}
s_0=\frac{\varpi_{\mathrm{vac}}}{\mathsf u}\cdot\frac{c_\psi}{\Omega_0},\qquad|\sigma
\varpi_{\mathrm{vac}}|^2=\frac{\aurd}{\sqrt d}\ \ (\sigma\in\auAut(\C/K)),\qquad|s_0|^2=\frac1\Omega .
\end{equation}
(For the second formula: $|\alpha|=\auNm(\mathfrak d_{E/K})^{1/2}$, the discriminant of $E$ is
$d^f\auNm(\mathfrak d_{E/K})$, and the conjugates over $K$ of the value $\upsilon(\alpha)$ of the
algebraic Hecke character $\upsilon$ have absolute value $1$. The third follows from the second
and (C4), since the area of $\Lambda$ is $|\Omega_0|^2\sqrt d/2$.) The integral vacuum of rank one
is defined in \eqref{pk:eq-vac}; it is $c_{\mathrm{vac}}\psvac$ for a constant
$c_{\mathrm{vac}}\ne0$ which is thereby determined but for which we have no closed formula.

\subsubsection{The definite lift: proof of Lemma~\ref{ps:definite}}\label{pk:ss-definite}

Let $X,X'$ be consecutive spaces $X_j,X_{j+1}$, of ranks $n<n'$, $\lambda=\lambda(X)$,
$\lambda'=\lambda(X')$, $e=e_{j,\tau}$, $e'=e_{j+1,\tau}$, $c(n)=c_\tau(n_j)$. Let $\rho_X$ be the
representation of $\prod_\tau\auGL_n(\C)$ on $W_\lambda$ and $\rho_{X'}$ that on $W_{\lambda'}$. The
determinant twists of $\omega_{\mu_{j+1}}$ at the infinite places are $\prod_\tau\det(h_\tau)^{e'}
\det(g_\tau)^{-(e-c(n))}$: indeed $\mu_{j+1}$ has weight $-\kappa_{j+1}$, $\upsilon|_{\auU(1)}$ has
weight $1$, $n,n'$ are even, and $e'=n/2-\kappa_{j+1}$, $e-c(n)=n'/2-\kappa_{j+1}$ (for $j=0$:
$e=0$, $c_\tau(k)=\kappa_1-n_1/2$). All $n_j$ are even, because $k=r+1-m$ and $t+1=5r+1$ are. By
the multiplicativity of Schur functors, $S(AzB)=\rho^{\circ}_{X'}(A)S(z)\rho^\circ_X(B)$, where
$\rho^\circ_{X'}=\rho_{X'}\otimes\det^{-e'}$ and $\rho_X^\circ=\rho_X\otimes\det^{-(e-c(n))}$ (the
columns of height $n$ give $\det^{c(n)}$). Hence, for a covector $A$ on $W_{\lambda'}$ and
$\xi\in W_\lambda$, as $\det\aubar g=\det g^{-1}$ and $h^t=\aubar h^{-1}$,
\begin{equation}\label{pk:eq-schuract}
\omega_{\mu_{j+1},\infty}(g,h)\bigl[A\cdot S(z)\xi\bigr]=A\cdot\rho_{X'}(\aubar h)^{-1}S(z)\,\rho_X(\aubar
g)\xi .
\end{equation}

\begin{lemma}[Cauchy decomposition]\label{pk:cauchy}
Let $\mathcal P$ be the space of polynomials in $z$. As a module over $\prod_\tau\auU(n)$ acting by
$\omega_{\mu_{j+1},\infty}(g,1)$, $\mathcal P$ is the direct sum of the spaces $\mathcal P_\nu$
spanned by the polynomials $A\cdot S_\nu(z)\xi$, where $\nu=(\nu_\tau)$ runs over the families of
diagrams with at most $n$ rows and $S_\nu$ is the Schur functor of $\nu$; the modules
$\mathcal P_\nu$ are isotypic of pairwise different irreducible types, and $(A,\xi)\mapsto A\cdot
S_\nu(z)\xi$ is injective on $W^*_{\nu}(\C^{n'})\otimes W_\nu(\C^n)$.
\end{lemma}
\begin{proof}
This is the decomposition of the polynomials on $M_{n'\times n}(\C)$ under $\auGL_{n'}\times\auGL_n$
into the irreducible modules $W_\nu(\C^{n'})^*\otimes W_\nu(\C^n)^*$, each once, the summand of
$\nu$ being the span of the matrix coefficients of $S_\nu(z)$
(\cite{HoweTheta}, \cite[Theorem~5.6.7]{GoodmanWallach}; for the unitary groups see
\cite[Ch.~III, Proposition~5.6]{KashiwaraVergne}), applied at each $\tau$.
\end{proof}

\begin{lemma}[One lift]\label{pk:single}
For every $P\in\mathcal P$ and every linear form $l$ on $W_\lambda$ there is a linear form $l'$ on
$W_{\lambda'}$ such that for all twists, all $\Phi\in\psHS_f(\auHom(X,X'))$, all algebraic modular
forms $\psi$ on $\auU(X)$ of weight $\lambda$ and all $h\in\auU(X')(\auA)$,
\[
\int_{[\auU(X)]}l\bigl(g_\infty^{-1}\cdot\psi(g_f)\bigr)\,\theta\bigl(\omega_{\mu_{j+1}}(g,h)(P\otimes
\Phi)\bigr)\,dg=l'\Bigl(h_\infty^{-1}\cdot\bigl(\Theta^X_j(\Phi)\psi\bigr)(h_f)\Bigr).
\]
The form $l'$ depends linearly on $P$, and $l'=0$ if $P$ is orthogonal to $\mathcal P_{\nu^0}$,
where $\nu^0$ is the diagram of $\lambda$ with its $c(n)$ columns of height $n$.
\end{lemma}
\begin{proof}
Write $G_\infty=\auU(X)(\R)$, compact, and $[\auU(X)]=\Gamma\backslash(G_\infty\times G_f)$. Since
$[\auU(X)]$ has volume $1$ and $J_X$ is free, the measure is the product of the measure of $G_f$
with volume $1$ on $J_X$ and of the Haar measure of $G_\infty$ of total mass
$m_\infty=\#(\Gamma\backslash G_f/J_X)^{-1}$. For fixed $g_f$ put $v=\psi(g_f)$ and
$\Phi'=\omega_{\mu_{j+1},f}(g_f,h_f)\Phi$. By (W0) the integral over $G_\infty$ is
$\theta(P_v\otimes\Phi')$ with
\[
P_v=\int_{G_\infty}l\bigl(\rho_X(\aubar g)^{-1}v\bigr)\,\omega_\infty(g,1)P'\,dg\in\mathcal P,\qquad
P'=\omega_{\mu_{j+1},\infty}(1,h_\infty)P .
\]
The map $v\mapsto P_v$ is linear and $P_{\rho_X(\aubar g_0)v}=\omega_\infty(g_0,1)P_v$. The module
$(W_\lambda,\ g\mapsto\rho_X(\aubar g))$ is irreducible, of the type of $\mathcal P_{\nu^0}$ by
\eqref{pk:eq-schuract}. By Lemma~\ref{pk:cauchy} and Schur's lemma there is a unique covector
$A=A(l,P')$ on $W_{\lambda'}$ with $P_v=A\cdot S(z)v$ for all $v$; it is linear in $P'$ and
vanishes if $P'\perp\mathcal P_{\nu^0}$. As $\omega_\infty(1,h_\infty)$ commutes with
$\omega_\infty(g,1)$, \eqref{pk:eq-schuract} gives $A(l,P')=A(l,P)\circ\rho_{X'}(\aubar h_\infty)^{-1}$.
The entries of $S$ are homogeneous of degree $L=\sum_{\tau,p}l_{\tau,p}$, so by
\eqref{pk:eq-bargmann} and \eqref{ps:eq-kernel}
\[
\theta(P_v\otimes\Phi')=A(l,P')\bigl(\theta(S(z)\otimes\Phi')v\bigr)=(\pi_Ms_0)^{-L}A(l,P)\Bigl(
h_\infty^{-1}\cdot\theta_j(\Phi;g_f,h_f)\psi(g_f)\Bigr).
\]
Integrating over $\Gamma\backslash G_f$ and using \eqref{pk:eq-lift} gives the lemma with
$l'=m_\infty(\pi_Ms_0)^{-L}A(l,P)$.
\end{proof}

\begin{proposition}[Lemma~\ref{ps:definite}]\label{pk:definite}
Lemma~\ref{ps:definite} holds, for every $u\in\mathfrak F^{\mathrm{sm}}(\auHom(X',X^-\oplus X''))$,
with $g_\infty^{-1}\psi(g_f)$ and $k_\infty^{-1}(\cdots)$ read as $g_\infty^{-1}\cdot\psi(g_f)$,
$k_\infty^{-1}\cdot(\cdots)$, and the tensor product of the two Weil representations transported
by (W6).
\end{proposition}
\begin{proof}
Denote the left side by $\mathcal L(u,l)$; it is linear in $u$ and $l$ and, by (W0) and the
moderate growth of theta functions, continuous in $u$. \emph{Reduction to polynomials.} Replacing
$u$ by $\omega_\infty(g_0,r_0,1)u$ with $g_0\in\auU(X)(\R)$, $r_0\in\auU(X')(\R)$ and substituting
$g\mapsto gg_0^{-1}$, $r\mapsto rr_0^{-1}$, we get $\mathcal L(\omega_\infty(g_0,r_0,1)u,l)=
\mathcal L(u,l\circ g_0)$. Apply this to the scalar matrices $g_0=(e^{i\vartheta_\tau})$,
$r_0=(e^{i\vartheta'_\tau})$, which act on $W_\lambda$ by characters: by (W6), $\mathfrak F
(\auHom(X',X^-\oplus X''))$ is the completed tensor product of the Fock spaces of $\auHom(X,X')$
and $\auHom(X',X'')$, with variables $z_1,z_2$, and by \eqref{pk:eq-fock} these scalars multiply
the part of $u$ of bidegree $(d_1(\tau),d_2(\tau))_\tau$ in $(z_1,z_2)$ by a character of
$(\vartheta,\vartheta')$ from which $d_1$ and $d_1-d_2$ can be read. Hence $\mathcal L(u,l)=
\mathcal L(u^\flat,l)$, where $u^\flat$ is the component of $u$ of one bidegree, a polynomial.
\emph{Polynomials.} By linearity let $u=P_1\otimes P_2$ and $\Phi=\Phi_1\otimes\Phi_2$. Then
$\theta(\omega(g,r,k)(u\otimes\Phi))$ is a constant multiple of $\theta(\omega_{\mu_{j+1}}(g,r)
(P_1\otimes\Phi_1))\cdot\theta(\omega_{\mu_{j+2}}(r,k)(P_2\otimes\Phi_2))$ by (W6). Integrate over
$g$ by Lemma~\ref{pk:single}, which gives $l'(r_\infty^{-1}\cdot(\Theta^X_j(\Phi_1)\psi)(r_f))$,
and then over $r$ by the same lemma for the pair $(X',X'')$, the form $l'$ and the form
$\Theta^X_j(\Phi_1)\psi$. The resulting form $l''$ depends only on $(l,P_1,P_2)$.
\end{proof}

\subsubsection{The character $\chi_0$ is trivial}\label{pk:ss-chi0}

\begin{lemma}\label{pk:chi0}
In the proof of Lemma~\ref{ps:double}, $\chi_0=1$: with the determinant twists of
\S\ref{ps:ss-weil} on the spaces $Y$ and on the spaces $X$, and the identifications {\rm(W6)},
\[
\psZ_{\mathbb J}\circ\bigl(\omega_{\mu_{j+1}}(\iota g,\iota r')\otimes\omega_{\mu_{j+2}}(\iota r',
\iota k)\bigr)=\bigl(\omega_{\mu_{j+1}}(g,r')\otimes\omega_{\mu_{j+2}}(r',k)\bigr)\circ\psZ_{\mathbb J}
\]
for $g\in\auU(X)(\auA)$, $r'\in\auU(X')(\auA)$, $k\in\auU(X'')(\auA)$ (resp.\ $\auU(Y_m)(\auA)$).
Hence Proposition~\ref{au:seesaw} holds with $\chi_0=1$, the upper chain being defined by $\omega_\mu$.
\end{lemma}
\begin{proof}
Let $n,n',n''$ be the ranks of $X,X',X''$ and $N=n+s$, $N'=n'+s-1$, $N''$ the ranks of
$V,V',V''$; so $N''=n''+s-2$ if $s\ge3$ and $N''=n''-1$ if $s=1$. Here $s$ is odd and
$n,n',n''$ are even; thus $N,N''$ are odd and $N'$ is even. Write $\upsilon(g)$ for
$\upsilon(\det g)$. By the definition of $\omega_\mu$ and (W6), the left operator, before
$\psZ_{\mathbb J}$, is
\[
\frac{\mu_{j+1}(r')\mu_{j+2}(k)}{\mu_{j+1}(g)\mu_{j+2}(r')}\cdot\upsilon(r')^{\frac{N-1}2}\upsilon(g)^{-
\frac{N'}2}\cdot\upsilon(k)^{\frac{N'}2}\upsilon(r')^{-\frac{N''-1}2}\cdot\upsilon(g)^{N'}\upsilon(r')^{-N}
\cdot\omega_{\auHom(V',\mathbb V)}(r',g\oplus k),
\]
and the right operator is the same expression with $(n,n',n'')$ in place of $(N,N',N'')$, the
floors being $n/2,n'/2,n''/2$, and $\auHom(X',X^-\oplus X'')$ in place of $\auHom(V',\mathbb V)$.
By (W2$'$), $\psZ_{\mathbb J}\,\omega_{\auHom(V',\mathbb V)}(p)=\upsilon(\det p|_{\mathbb J})^{-1}
\omega_{\auHom(X',X^-\oplus X'')}(p)\,\psZ_{\mathbb J}$ for $p=(r',g\oplus k)$; the absolute value
is $1$ because $\det g$, $\det r'$, $\det k$ are ideles of norm one to $E^+$.

\emph{The determinant on $\mathbb J$.} For $s=1$ all grades are $0$, $\mathbb J=0$ and the
determinant is $1$. Let $s\ge3$. The element $p$ acts on $w\in\auHom(V',\mathbb V)$ by
$w\mapsto(g\oplus k)wr'^{-1}$; it acts by $r'$ on $X'$, by $g\oplus k$ on $X\oplus X''$ and
trivially on all vectors of the strings. On a block $\auHom(A,B)$ the determinant is
$\det(p|_B)^{\dim A}\det(p|_A)^{-\dim B}$. The string of $V'$ has $s-1$ vectors, of odd grades,
$(s-1)/2$ of them negative; the strings of $V$ and $V''$ have $(s-1)/2$ and $(s-3)/2$ vectors of
positive grade. The blocks of $\mathbb J=\auHom_{>0}(V',\mathbb V)\oplus\auHom(X',\ell)$ on which
$p$ is not the identity are therefore: $\auHom(V'_a,X\oplus X'')$ for the $(s-1)/2$ negative
grades $a$; $\auHom(X',\mathbb V_b)$ for the $(s-1)/2+(s-3)/2=s-2$ string vectors of positive
grade $b$; and $\auHom(X',\ell)$. So
\[
\det p|_{\mathbb J}=(\det g\det k)^{\frac{s-1}2}(\det r')^{-(s-1)}\ \ (s\ge3),\qquad\det p|_{\mathbb
J}=1\ \ (s=1).
\]
\emph{Comparison of exponents} (the characters $\mu$ agree on both sides). For $s\ge3$: the
exponent of $\upsilon(g)$ on the left is $-\frac{N'}2+N'-\frac{s-1}2=\frac{n'}2$, and on the right
$-\frac{n'}2+n'=\frac{n'}2$. The exponent of $\upsilon(k)$ is $\frac{N'}2-\frac{s-1}2=\frac{n'}2$
on the left and $\frac{n'}2$ on the right. The exponent of $\upsilon(r')$ on the left is
$\frac{N-1}2-\frac{N''-1}2-N+(s-1)=\frac{n-n''}2+1-n-1=-\frac{n+n''}2$, and on the right
$\frac n2-\frac{n''}2-n=-\frac{n+n''}2$. For $s=1$: $N=n+1$, $N'=n'$, $N''=n''-1$, and the exponents
on the left are $\frac{n'}2$ for $\upsilon(g)$ and $\upsilon(k)$, and $\frac n2-(\frac{n''}2-1)-
(n+1)=-\frac{n+n''}2$ for $\upsilon(r')$; they agree with those on the right ($\det\operatorname{diag}
(k,1)=\det k$).
\end{proof}

\subsubsection{One step of the unfolding: proof of Proposition~\ref{ps:onestep}}\label{pk:ss-onestep}

We give a complete proof of Proposition~\ref{ps:onestep}. Steps (1), (2) and (4) of the proof in
\S\ref{ps:ss-seesaw} are kept; steps (3), (5), (6) are replaced by
Lemmas~\ref{pk:vanish}--\ref{pk:fine} and the argument at the end. The notation is that of
Proposition~\ref{ps:onestep}: $0\le j<m-1$, $V=Y_j$, $V'=Y_{j+1}$, $s=s_j\ge1$,
$W=\auHom(V,V')$, $G=G_j$, $G'=G_{j+1}$, $R=R_j$, $R'=R_{j+1}$, $N=N_j$, $T=T_j$,
$\mathfrak e=\mathfrak e_j=T^\dagger T$, $\omega=\omega_{\mu_{j+1}}$. We write $V_g$ for the part of
grade $g$ of $V$, $g_a=s-1-2a$ for the grade of $f_a$, $g'_b=s-2-2b$ for that of $f'_b$, and
$V^+,V^0,V^-$ for the parts of positive, zero, negative grade; $p=\dim V^+=\lfloor s/2\rfloor$.
Recall $Tf_a=t_af'_a$ with $t_a\ne0$ ($a\le s-2$), $\mathfrak ef_a=\mathsf e_af_{a+1}$ with
$\mathsf e_a\ne0$ ($a\le s-2$), and $T(X_j)=0$.

\emph{Three remarks on grades.} (G1) The non-zero grades of $V$ are $\equiv s-1$ and those of $V'$
are $\equiv s\pmod2$. Hence a map $V\to V'$ of degree $\ge0$ has degree $\ge1$ on $V^+\oplus V^-$,
and $\auHom_{\ge0}(V^+,V')=\auHom_{>0}(V^+,V')=:J$; the maps in $J$ have image in
$V'^{+}\subseteq T(V^+)$. (G2) $T(V^+)$ is the span of $f'_0,\dots,f'_{p-1}$; the form
$Q_T(v,v')=\langle Tv,Tv'\rangle$ on $V^+$ is $0$ if $s$ is odd, and if $s$ is even its radical is
$\langle f_0,\dots,f_{p-2}\rangle$ and $Q_T(f_{p-1},f_{p-1})\ne0$ (because $\langle f'_a,f'_b\rangle
\ne0$ only for $a+b=s-2$, and $a,b\le p-1$). (G3) For $y\in V$, $y'\in V'$ let $y'\otimes y^\flat$
be the map $x\mapsto\langle y,x\rangle y'$; then $\pktr(w^\dagger(y'\otimes y^\flat))=\langle wy,y'
\rangle$ for $w\in W$, and $f_a^\flat$ vanishes outside $V_{-g_a}$.

\emph{Constant terms along an isotropic space.} Let $\pkI\subseteq W$ be one of the two isotropic
spaces
\[
I_0=\{w:\ w|_{V^+\oplus V^0}=0\},\qquad I=\auHom_{>0}(V,V'),
\]
so that $I_0^\perp=\{w:w|_{V^+}=0\}$ and $I^\perp=\auHom_{\ge0}(V,V')$. Both are stable under
$w\mapsto w\gamma$, $\gamma\in R$, and $I_0$ is stable under $w\mapsto\delta w$, $\delta\in G'$.
Moreover
\begin{equation}\label{pk:eq-H}
\pkI^\perp\cdot(n-1)\subseteq\pkI\qquad(n\in N),
\end{equation}
because $n-1$ has degree $\ge1$ and maps $V^+\oplus V^0$ into $V^+$. For $w\in W(\auA)$ and
$F\in\psHS(W)$ put $c^{\pkI}_w(F)=\theta_{W_{\pkI}}(\psZ_{\pkI}\psrho(w)F)$; thus $c^I=c$ and
$c^{I_0}=c^0$ in the notation of \S\ref{ps:ss-seesaw}. By (W1), (W2):
\begin{itemize}
\item[(c1)] $c^{\pkI}_{w+i}=\beta(i,w)^{-1}c^{\pkI}_w$ for $i\in\pkI(\auA)$, and $c^{\pkI}_{w+\delta}=
c^{\pkI}_w$ for $w\in W(E)$, $\delta\in\pkI^\perp(E)$;
\item[(c2)] $c^{\pkI}_w(\omega(\gamma,1)F)=c^{\pkI}_{w\gamma}(F)$ for $\gamma\in R(E^+)$, and
$c^{\pkI}_w(\omega(n,1)F)=c^{\pkI}_{wn}(F)$ for $n\in N(\auA)$; for $\pkI=I_0$ also
$c^{0}_w(\omega(1,\delta)F)=c^{0}_{\delta^{-1}w}(F)$ for $\delta\in G'(E^+)$.
\end{itemize}
(For (c2): $\psrho(w)\omega(g,h)=\omega(g,h)\psrho(h^{-1}wg)$; $\psZ_{\pkI}\omega(p)=\chi_{\pkI}(p)
\omega_{W_{\pkI}}(\bar p)\psZ_{\pkI}$; for rational $p$, $\chi_{\pkI}(p)=1$ and
$\theta_{W_{\pkI}}\circ\omega_{W_{\pkI}}(\bar p)=\theta_{W_{\pkI}}$; for $p=(n,1)$, $\chi_{\pkI}(p)=1$
and $\bar p=1$ by \eqref{pk:eq-H}; and the determinant twists are trivial on rational and on
unipotent elements.)

For $w\in W(E)$ let
\[
A^{\pkI}_w=\{n\in N^{\ge2}:\ w(n-1)\in\pkI\},\qquad\chi^{\pkI}_w(n)=\beta\bigl(w(n-1),w\bigr)^{-1}=
\psi\bigl(\tfrac12\jmath\,\pktr(w^\dagger w(n-1))\bigr)^{-1}.
\]

\begin{lemma}\label{pk:char}
{\rm(a)} $A^{\pkI}_w$ is a unipotent algebraic subgroup of $N^{\ge2}$ over $E^+$, and
$\chi^{\pkI}_w$ is a character of $A^{\pkI}_w(\auA)$, trivial on $A^{\pkI}_w(E^+)$; both depend
only on the class of $w$ modulo $\pkI^\perp(E)$. {\rm(b)} $c^{\pkI}_w(\omega(n,1)F)=
\chi^{\pkI}_w(n)\,c^{\pkI}_w(F)$ for $n\in A^{\pkI}_w(\auA)$. {\rm(c)} For $\gamma\in R(E^+)$,
$A^{\pkI}_{w\gamma}=\gamma^{-1}A^{\pkI}_w\gamma$ and $\chi^{\pkI}_{w\gamma}(\gamma^{-1}n\gamma)=
\chi^{\pkI}_w(n)$; for $n\in A^{\pkI}_w(E^+)$, $wn\equiv w$ modulo $\pkI^\perp(E)$.
\end{lemma}
\begin{proof}
Write $\nu=n-1$, $\nu'=n'-1$; $\nu^\dagger=n^{-1}-1$. (a) $w(nn'-1)=w\nu+w\nu'+(w\nu)\nu'$ and
$(w\nu)\nu'\in\pkI\nu'\subseteq\pkI$ by \eqref{pk:eq-H}; $w(n^{-1}-1)=-w\nu-(w\nu)\nu^\dagger\in
\pkI$. So $A^{\pkI}_w$ is a group; it is defined by linear conditions on $\nu$. Since
$\langle w,u\nu'\rangle=\pktr(\nu'w^\dagger u)=\langle w\nu'^\dagger,u\rangle$, we get $\langle w,
(w\nu)\nu'\rangle=\langle w(n'^{-1}-1),w\nu\rangle=0$, both entries being in $\pkI$; so
$\chi^{\pkI}_w$ is multiplicative. It is trivial on rational points because $\psi$ is trivial on
$E$. If $\delta\in\pkI^\perp(E)$ then $\delta\nu\in\pkI$ by \eqref{pk:eq-H}, so
$A^{\pkI}_{w+\delta}=A^{\pkI}_w$, and $\langle w+\delta,(w+\delta)\nu\rangle-\langle w,w\nu\rangle=
\langle\delta,w\nu\rangle+\langle w\nu^\dagger,\delta\rangle+\langle\delta,\delta\nu\rangle=0$, each
term pairing $\pkI$ with $\pkI^\perp$. (b) By (c2) and (c1), $c^{\pkI}_w(\omega(n,1)F)=
c^{\pkI}_{w+w\nu}(F)=\beta(w\nu,w)^{-1}c^{\pkI}_w(F)$. (c) $w\gamma(\gamma^{-1}n\gamma-1)=(w\nu)
\gamma$ and $\pkI\gamma=\pkI$; $\langle w\gamma,(w\nu)\gamma\rangle=\langle w,w\nu\rangle$ as
$\gamma$ is unitary; and $wn-w=w\nu\in\pkI(E)$.
\end{proof}

A class $w\in W(E)/\pkI^\perp(E)$ is \emph{dead} (for $\pkI$) if the character
$\psi_j\chi^{\pkI}_w$ of $A^{\pkI}_w(\auA)$ is not trivial. By Lemma~\ref{pk:char}(c) and the
invariance of $\psi_j$ under conjugation by $R(\auA)$, the set of dead classes is stable under
$R(E^+)$.

\begin{lemma}[Vanishing]\label{pk:vanish}
Let $\mathcal O$ be a set of dead classes for $\pkI$, stable under $R(E^+)$, $\psW$ an equivariant
function on $[R]$ and $F\in\psHS(W)$. Then $\sum_{w\in\mathcal O}|c^{\pkI}_w(\omega(r,1)F)|$
converges uniformly for $r$ in compact sets and
\[
\int_{[R]}\psW(r)\sum_{w\in\mathcal O}c^{\pkI}_w(\omega(r,1)F)\,dr=0 .
\]
\end{lemma}
\begin{proof}
The convergence is (W3). Since $[R]$ is compact it suffices to treat one orbit $\mathcal O=w_0
R(E^+)$. Let $S\subseteq R(E^+)$ be the stabiliser of the class $w_0$, $A=A^{\pkI}_{w_0}$,
$\chi=\chi^{\pkI}_{w_0}$. By (c2), $\sum_{w\in\mathcal O}c^{\pkI}_w(\omega(r,1)F)=\sum_{\gamma\in
S\backslash R(E^+)}\phi(\gamma r)$ with $\phi(r)=c^{\pkI}_{w_0}(\omega(r,1)F)$, and $\psW$ is left
invariant under $R(E^+)$; so the integral is $\int_{S\backslash R(\auA)}\psW\phi$, absolutely
convergent. By Lemma~\ref{pk:char}(c), $A(E^+)\subseteq S$, $S\cap A(\auA)=A(E^+)$, and $S$
normalises $A$. The group $S\,A(\auA)$ is closed in $R(\auA)$ (it is $S\cdot\mathsf K$ with $S$
discrete and $\mathsf K\subset A(\auA)$ compact, $[A]$ being compact), and it is unimodular, like
$R(\auA)$, since both contain discrete subgroups with compact quotient. Hence
\[
\int_{S\backslash R(\auA)}\psW\phi=\int_{SA(\auA)\backslash R(\auA)}\int_{[A]}\psW(ar)\phi(ar)\,da\,
d\dot r,
\]
and $\psW(ar)\phi(ar)=\psi_j(a)\chi(a)\psW(r)\phi(r)$ by the equivariance of $\psW$ and
Lemma~\ref{pk:char}(b). The inner integral is $0$ because $\psi_j\chi$ is a non-trivial character
of the compact group $[A]$.
\end{proof}

\begin{lemma}[Coarse classes]\label{pk:coarse}
Let $w\in\auHom(V^+,V')(E)=W(E)/I_0^\perp(E)$. If $w\notin G'(E^+)\,T|_{V^+}$, then $w$ is dead
for $I_0$.
\end{lemma}
\begin{proof}
Let $\tilde w$ be the extension of $w$ by $0$ on $V^0\oplus V^-$, and $Q_w(v,v')=\langle wv,wv'
\rangle$ on $V^+$. For $n\in N^{\ge2}$, $n\in A^{I_0}_w$ means $w(n-1)|_{V^+\oplus V^0}=0$.

(a) \emph{$Q_w=Q_T$ unless $w$ is dead.} Let $\nu\colon V\to V$ be zero on $V^+\oplus V^0$, map
$V^-$ into $V^+$, and satisfy $\nu^\dagger=-\nu$; then $n=1+\nu\in N^{\ge2}$ and $\tilde w\nu$
vanishes on $V^+\oplus V^0$, so $n\in A^{I_0}_w$. In the basis $f_{s-1-a}$ ($a<p$) of $V^-$, dual
to the basis $f_a$ of $V^+$, write $\nu f_{s-1-a}=\sum_{b<p}\nu_{ba}f_b$; the condition
$\nu^\dagger=-\nu$ is $\nu_{ab}=-\aubar{\nu_{ba}}$, that is, $(\nu_{ba})=\jmath\mathsf H$ with
$\mathsf H$ hermitian. Now $\pktr(\tilde w^\dagger\tilde w\nu)=\sum_{a<p}\langle f_a,\tilde
w^\dagger\tilde w\nu f_{s-1-a}\rangle=\sum_{a,b}\nu_{ba}Q_w(f_a,f_b)$, and the same with $T$. Hence
$\psi_j\chi^{I_0}_w(n)=\psi\bigl(\tfrac12\jmath^2\pktr(\mathsf H\mathsf D)\bigr)$ with
$\mathsf D=(Q_T-Q_w)(f_a,f_b)$, a hermitian matrix over $E$. If $\mathsf D\ne0$, $\mathsf
H\mapsto\jmath^2\pktr(\mathsf H\mathsf D)$ is a surjective $\auA_{E^+}$-linear map from the
hermitian adelic matrices to $\auA_{E^+}$, and $x\mapsto\psi(\tfrac12x)$ is a non-trivial character
of $\auA_{E^+}$; so $w$ is dead.

(b) \emph{$w$ is injective unless it is dead.} Let $Q_w=Q_T$ and $\mathsf K=\ker w\ne0$; then
$\mathsf K$ lies in the radical of $Q_T$. \emph{First case:} $\mathsf K$ contains a non-zero
vector of $\langle f_0,\dots,f_{p-2}\rangle$. Choose such a vector $\kappa=\sum_{b<i}\kappa_bf_b$
with $\kappa_{i-1}\ne0$, $1\le i\le p-1$. For $\mathsf t\in\auA_E$ let $\nu_+$ be the map
$f_i\mapsto\mathsf t\kappa$, zero on the other basis vectors of $V$, and $n=1+\nu_+-\nu_+^\dagger$.
As $\nu_+^2=0=\nu_+\nu_+^\dagger=\nu_+^\dagger\nu_+$, $n$ is unitary; $\nu_+$ raises grades by
$2(i-b)\ge2$, so $n\in N^{\ge2}$; $\nu_+^\dagger$ vanishes on $V^+\oplus V^0$ and $\tilde w\nu_+=0$,
so $\tilde w(n-1)=0$: thus $n\in A^{I_0}_w$ and $\chi^{I_0}_w(n)=1$. Further
$\pktr(\mathfrak e\nu_+)=\mathsf t\,\mathsf e_{i-1}\kappa_{i-1}=:\varsigma$ and
$\pktr(\mathfrak e\nu_+^\dagger)=\aubar\varsigma$ ($\mathfrak e$ is self-adjoint), so
$\psi_j(n)=\psi(\tfrac12\jmath(\varsigma-\aubar\varsigma))=\psi(\jmath\varsigma)$, which is not $1$
for a suitable $\mathsf t$. \emph{Second case:} $\mathsf K\cap\langle f_0,\dots,f_{p-2}\rangle=0$.
By (G2) $s$ is odd, $V^0=X_j\oplus Ez$ with $z=f_p$, and $\mathsf K=E\kappa$ with
$\kappa_{p-1}\ne0$. Let $b\colon V^0\to V^+$ be $z\mapsto\mathsf t\kappa$, $X_j\mapsto0$, extended
by $0$, and $n=\exp(b-b^\dagger)=1+b-b^\dagger-\tfrac12bb^\dagger$. All grades being even,
$n\in N^{\ge2}$; $\tilde wb=0$ and $\tilde wb^\dagger=0$ ($b^\dagger$ maps $V^-$ to $V^0$), so
$\tilde w(n-1)=0$, $n\in A^{I_0}_w$, $\chi^{I_0}_w(n)=1$. And $\pktr(\mathfrak eb)=\mathsf t\,
\mathsf e_{p-1}\kappa_{p-1}=:\varsigma$, $\pktr(\mathfrak eb^\dagger)=\aubar\varsigma$,
$\pktr(\mathfrak ebb^\dagger)=0$ ($\mathfrak ebb^\dagger$ maps $V^-$ into $V^+\oplus V^0$ and kills
the rest). So $\psi_j(n)=\psi(\jmath\varsigma)\ne1$ for a suitable $\mathsf t$.

(c) If $w$ is not dead, it is injective with $Q_w=Q_T$; then $Tv\mapsto wv$ is an isometry between
subspaces of $V'$, which by Witt's theorem is induced by an element of $G'(E^+)$.
\end{proof}

\begin{lemma}[Fine classes]\label{pk:fine}
Let $w=T+v$ with $v\in W(E)$, $v|_{V^+}=0$, $v\notin\auHom_{\ge0}(V,V')$. Then the class of $w$ is
dead for $I$.
\end{lemma}
\begin{proof}
Write $v$ in the bases: for a basis vector $\kappa$ of $V^0\oplus V^-$ (a vector of the basis of
$X_j$, or $f_a$ with $a\ge p$), of grade $g(\kappa)\le0$, $v\kappa=\sum_bv_{b,\kappa}f'_b+
(\text{a vector of }X_{j+1})$. An entry $v_{b,\kappa}$ has negative degree if $g'_b<g(\kappa)$;
by hypothesis some entry of negative degree is not zero (an entry with values in $X_{j+1}$ has
degree $-g(\kappa)\ge0$). Among the non-zero entries of negative degree choose one,
$v_{c,\kappa_0}$, with $g'_c$ least, and put $g'=g'_c<g(\kappa_0)\le0$ and $a=s-2-c$. Then
$g_a=1-g'\ge2$, so $f_a\in V^+$; and $Tf_a=t_af'_a$ pairs to zero with all basis vectors of $V'$
except $f'_c$.

For $\mathsf t\in\auA_E$ let $\kappa=\mathsf t\kappa_0$ and
\[
n=1+f_a\otimes\kappa^\flat-\kappa\otimes f_a^\flat-\tfrac12\langle\kappa,\kappa\rangle f_a\otimes
f_a^\flat .
\]
Since $\langle f_a,f_a\rangle=0=\langle f_a,\kappa_0\rangle$ (the only basis vector not orthogonal
to $f_a$ is $f_{s-1-a}$, of grade $g'-1<g(\kappa_0)$), a direct computation gives $\langle nx,ny
\rangle=\langle x,y\rangle$. By (G3) the three terms of $n-1$ map $V_{-g(\kappa_0)}\to V_{g_a}$,
$V_{-g_a}\to V_{g(\kappa_0)}$, $V_{-g_a}\to V_{g_a}$, of degrees $g_a+g(\kappa_0)=1-g'+g(\kappa_0)
\ge2$ (twice) and $2g_a\ge2$; so $n\in N^{\ge2}(\auA)$. Next, $wf_a=Tf_a$ and
\[
w(n-1)=Tf_a\otimes\kappa^\flat-w\kappa\otimes f_a^\flat-\tfrac12\langle\kappa,\kappa\rangle
Tf_a\otimes f_a^\flat .
\]
The first and third terms have degrees $g_a-1+g(\kappa_0)\ge1$ and $2g_a-1\ge1$. The second maps
$V_{-g_a}$ to the line of $w\kappa_0=T\kappa_0+v\kappa_0$; $T\kappa_0$ has grade
$g(\kappa_0)-1\ge g'=1-g_a$, the components of $v\kappa_0$ of negative degree have grades $\ge g'$
by the choice of $c$, and the others have grades $\ge g(\kappa_0)>g'$. So the second term has
degree $\ge g'+g_a=1$, and $n\in A^I_w(\auA)$. By (G3),
\[
\pktr(w^\dagger w(n-1))=\langle w\kappa,Tf_a\rangle-\langle Tf_a,w\kappa\rangle-\tfrac12\langle
\kappa,\kappa\rangle\langle Tf_a,Tf_a\rangle ,
\]
and the same holds with $w$ replaced by $T$. Subtracting,
$\pktr((\mathfrak e-w^\dagger w)(n-1))=\langle Tf_a,v\kappa\rangle-\langle v\kappa,Tf_a\rangle=
\varsigma-\aubar\varsigma$ with $\varsigma=\mathsf t\,\aubar{t_a}\,v_{c,\kappa_0}\langle f'_a,f'_c
\rangle$, and $\langle f'_a,f'_c\rangle\ne0$ as $a+c=s-2$. So $\psi_j\chi^I_w(n)=\psi(\jmath
\varsigma)$, which is not $1$ for a suitable $\mathsf t$.
\end{proof}

\begin{proof}[Proof of Proposition~\ref{ps:onestep}]
The case $j=m-1$ and step (1) (the properties of $\psW'$) are as in \S\ref{ps:ss-seesaw}. If
$s=1$ then $V^+=0$, $I=0$, $T=0$, $R'=G'$ and there is nothing to prove; let $s\ge2$. Fix
$h\in G'(\auA)$, write $T_+=T|_{V^+}$, and let $\mu(h)$ be a function of moderate growth on
$G'(\auA)$ with
\begin{equation}\label{pk:eq-major}
\sum_{w\in\auHom(V^+,V')(E)}\bigl|c^0_w(\omega(r,h)F)\bigr|\le\mu(h)\qquad(r\in\mathsf K_R),
\end{equation}
where $\mathsf K_R\subset R(\auA)$ is a compact set with $R(E^+)\mathsf K_R=R(\auA)$; it exists by
(W3) applied to $I_0$. In detail: by the law $c^0_w(\omega(\gamma,\delta)F')=c^0_{\delta^{-1}w\gamma}(F')$ of
\S\ref{ps:ss-seesaw} (with $\gamma=1$), the left side of \eqref{pk:eq-major} is unchanged when $h$ is
replaced by $\delta h$, $\delta\in G'(E^+)$, the twisting characters being trivial on rational points.
By Lemma~\ref{pcx:reduce} write $h=\delta^{-1}x$ with $\delta$ rational, $x_f$ in a fixed compact set and
$\|x_\infty\|\le C\,\mathrm{ht}(h)$ (Definition~\ref{pcx:def-ht}); by Lemma~\ref{pz:absgrowth}(b) (for
$F$ a sum of tensors $\Psi\otimes f$ with $\Psi$ smooth, $(r,x)$ with $(r_f,x_f)$ in a compact set and
$r_\infty$ bounded) the left side is at most $C'\,\mathrm{ht}(h)^{A_0}$, where $A_0=4(D_W+1)$ and $2D_W=\dim_\R\auHom(V,V')_\infty$ (the exponent $4s$ of that lemma; we rename it, $s$ being the rank here). So $\mu=C'\,\mathrm{ht}^{A_0}$ will do.

\emph{(i) Coarse expansion.} By (W3), $\theta(\omega(r,h)F)=\sum_wc^0_w(\omega(r,h)F)$. The set of
$w\notin G'(E^+)T_+$ is stable under $R(E^+)$, which acts through its restriction to $V^+$ (the
orbit $G'(E^+)T_+$ is stable by (ii) below), and consists of dead classes
(Lemma~\ref{pk:coarse}). By Lemma~\ref{pk:vanish} and (c2),
\[
\int_{[R]}\psW(r)\theta(\omega(r,h)F)\,dr=\sum_{\delta\in\mathsf M(E^+)\backslash G'(E^+)}\Psi(\delta h),
\qquad\Psi(h)=\int_{[R]}\psW(r)\,c^0_{T_+}(\omega(r,h)F)\,dr,
\]
where $\mathsf M$ is the stabiliser of $T_+$ in $G'$, and $\sum_\delta|\Psi|(\delta h)\le\|\psW
\|_\infty\mu(h)$, where $|\Psi|$ is defined with absolute values inside the integral.

\emph{(ii) The orbit of $T_+$ under $R'$.} By the description of $R'$ in \S\ref{ps:ss-Y} and (G2),
$R'$ is the group of $\epsilon\in G'$ with $\epsilon^{-1}Tf_i\in Tf_i+\langle Tf_b:b<i\rangle$ for
$i<p$; so $\mathsf M\subseteq R'$ and $\epsilon^{-1}T_+-T_+\in J(E)$ for $\epsilon\in R'(E^+)$.
Conversely, for $\eta\in J(E)$ the map $T_++\eta$ is injective and, by (G1), (G2),
$Q_{T_++\eta}=Q_T$; by Witt's theorem $T_++\eta=\epsilon^{-1}T_+$ with $\epsilon\in G'(E^+)$, and
then $\epsilon\in R'$. Hence $\epsilon\mapsto\epsilon^{-1}T_+$ is a bijection $\mathsf M(E^+)
\backslash R'(E^+)\to T_++J(E)$, and by (c2)
\begin{equation}\label{pk:eq-ii}
\sum_{\epsilon\in\mathsf M(E^+)\backslash R'(E^+)}\Psi(\epsilon h)=\int_{[R]}\psW(r)\sum_{\eta\in
J(E)}c^0_{T_++\eta}(\omega(r,h)F)\,dr .
\end{equation}
\emph{(iii) From coarse to fine classes.} Let $F'\in\psHS(W)$, let $\tilde T_+$ and $\tilde\eta$
be the extensions by $0$, and $\tilde J=\{\tilde\eta\}\subseteq I$. We claim
\begin{equation}\label{pk:eq-poisson}
\sum_{\eta\in J(E)}c^0_{T_++\eta}(F')=\sum_{w\in\mathcal C}c_w(F'),\qquad\mathcal C=\bigl(T+I_0^\perp
(E)\bigr)\big/\bigl(I_0^\perp(E)\cap I^\perp(E)\bigr),
\end{equation}
the set of classes modulo $I^\perp(E)$ of the maps $T+v$, $v|_{V^+}=0$. First, $I_0\cap\tilde
J^\perp=I_0\cap I$: for $i\in I_0$, $\langle\tilde\eta,i\rangle$ pairs the components of $\eta$
from $V_g$ to $V'_{g'}$, $g'>g>0$, with the components of $i$ from $V_{-g}$ to $V'_{-g'}$, and
these are the components of $i$ of negative degree; by (G1) a map of degree $\ge0$ on $V^-$ has
degree $>0$. Put $I_0'=I_0\cap I$. Since $\tilde J\cap I_0^\perp=0$, the pairing identifies
$\tilde J(E)$ with the character group of the compact group $\mathsf B=I_0(\auA)/(I_0(E)+
I_0'(\auA))$. The function $\phi(i)=\theta(\psrho(i)\psrho(\tilde T_+)F')$ on $[I_0]$ is
continuous, and by (W1), as $\psrho(i)\psrho(\tilde\eta)=[i,\tilde\eta]\psrho(\tilde\eta)\psrho(i)$
and $\theta\circ\psrho(\tilde\eta)=\theta$, one has $c^0_{T_++\eta}(F')=\int_{[I_0]}[i,\tilde\eta]
\phi(i)\,di$. These are the Fourier coefficients of the continuous function
$\phi_1(\mathsf b)=\int_{[I_0']}\phi(\mathsf b+i')\,di'$ on $\mathsf B$; they are absolutely
summable by (W3), so their sum is $\phi_1(0)=\theta_{W_{I_0'}}(\psZ_{I_0'}\psrho(\tilde T_+)F')$.
Now apply (W3) to the space $W_{I_0'}=I_0'^\perp/I_0'$ and its isotropic subspace $I/I_0'$, and
(W2) ($\psZ_I=\psZ_{I/I_0'}\psZ_{I_0'}$): the sum is $\sum_uc_{\tilde T_++u}(F')$ over
$u\in I_0'^\perp(E)/I^\perp(E)$. Finally $I_0'^\perp=I_0^\perp+I^\perp$ and $T-\tilde T_+\in
I_0^\perp$, which gives \eqref{pk:eq-poisson}.

\emph{(iv) Only the class of $T$ survives.} The set $\mathcal C\smallsetminus\{T\}$ is stable
under $R(E^+)$: for $\gamma\in R(E^+)$, $T\gamma-T$ has degree $\ge0$, and $v\gamma$ vanishes on
$V^+$ and has degree $\ge0$ if and only if $v$ has. Its elements are dead for $I$
(Lemma~\ref{pk:fine}). By Lemma~\ref{pk:vanish}, \eqref{pk:eq-ii} and \eqref{pk:eq-poisson} with
$F'=\omega(r,h)F$,
\[
\sum_{\epsilon\in\mathsf M(E^+)\backslash R'(E^+)}\Psi(\epsilon h)=\int_{[R]}\psW(r)\,c_T(\omega(r,h)F)
\,dr=\psW'(h).
\]
\emph{(v) Conclusion and convergence.} By (i) and (iv), $\int_{[R]}\psW(r)\theta(\omega(r,h)F)\,dr=
\sum_{\gamma\in R'(E^+)\backslash G'(E^+)}\psW'(\gamma h)$ and $\sum_\gamma|\psW'(\gamma h)|\le
\sum_\delta|\Psi|(\delta h)\le\|\psW\|_\infty\mu(h)$. Let $z$ be rapidly decreasing on $[G']$. Then
$\int_{R'(E^+)\backslash G'(\auA)}|\psW'z|=\int_{[G']}\sum_\gamma|\psW'(\gamma h)|\,|z(h)|\,dh\le
\|\psW\|_\infty\int_{[G']}\mu|z|<\infty$. (Here ``rapidly decreasing'' means: on a Siegel set
$\mathfrak S=\mathfrak S_\infty\mathsf cK'$ of $G'$ with $G'(\auA)=G'(E^+)\mathfrak S$, as in the proof of
Proposition~\ref{pz:types}(a), $|z(g)|\le C_N\|g_\infty\|^{-N}$ for every $N$. On $\mathfrak S$,
$\mathrm{ht}(g)\le c_0^{-1}\|g_\infty\|$: for $g=g_\infty ck'$ and $0\ne\gamma\in\Lambda_g$, the vector
$c^{-1}\gamma$ lies in one of finitely many lattices, so $\|\gamma\|\ge c_0>0$ and $|\gamma|_g=
\|g_\infty^{-1}\gamma\|\ge c_0\|g_\infty\|^{-1}$; this is the paragraph ``Heights on Siegel sets'' in the
proof of Proposition~\ref{ptr:conv}. Hence $\mu|z|=C'\mathrm{ht}^{A_0}|z|$ is bounded on $\mathfrak S$, hence
on $G'(\auA)$ by left $G'(E^+)$-invariance, and $[G']$ has finite volume.) And since $|\theta(\omega(r,h)F)|\le\mu(h)$ for
$r\in\mathsf K_R$, Fubini's theorem gives
\[
\int_{[R]}\psW(r)\int_{[G']}\theta(\omega(r,h)F)z(h)\,dh\,dr=\int_{[G']}z(h)\sum_\gamma\psW'(\gamma h)
\,dh=\int_{R'(E^+)\backslash G'(\auA)}\psW'z .\qedhere
\]
\end{proof}

\subsubsection{Support and finite data in the seesaw identity}\label{pk:ss-rest}

We give the verifications that complete Lemmas~\ref{ps:support} and~\ref{ps:fin} and the proof
of Proposition~\ref{au:kernel}.

\emph{(R1) Iwasawa decomposition (used in Lemma~\ref{ps:support}).} Let $v$ be a finite place of
$E^+$, $\Lambda$ a unimodular lattice in $Y_j\otimes E^+_v$, $K_v$ its stabiliser and $P$ the
stabiliser of the flag of the spaces $\langle f_0,\dots,f_i\rangle$, $i<p$. Then
$G_j(E^+_v)=P(E^+_v)K_v$. Indeed $E/E^+$ is unramified at $v$, so $K_v$ is hyperspecial and this
is the Iwasawa decomposition \cite[\S3.3.2, p.~51]{TitsCorvallis}. (Directly: $K_v$ is transitive on the
isotropic flags of a given type, because for a totally isotropic direct summand $\Lambda_1$ of
$\Lambda$ there is a totally isotropic $\Lambda_1^*\subset\Lambda$ with $\Lambda=(\Lambda_1\oplus
\Lambda_1^*)\perp\Lambda_2$, the trace $\auO_{E,v}\to\auO_{E^+,v}$ being surjective, and
$\auGL(\Lambda_1)\subset K_v$ is transitive on the full flags of $\Lambda_1$ with saturated
members.)

\emph{(R2) The step with $\mathsf z_w$ in Lemma~\ref{ps:support}.} Let $j$ be even, $h=n\,l\,k$,
$\Lambda'=l\psM_j$, and suppose that $\Lambda'\cap E_vz$ is larger than $\widehat\auO u_jz$ at a
place $v=w\aubar w$, $w\mid\pi_N$. As shown there, the part of grade $0$ of $\Lambda'_v$ is then the
orthogonal sum of its intersections with $X_j$ and $E_vz$, so $\iota(\mathsf z_w)$, which is a
unit scalar on $X_j$ and the identity elsewhere, stabilises $\Lambda'$. Hence
$\mathsf u_e=h_f^{-1}n\iota(\mathsf z_w^{\,e})n^{-1}h_f=k^{-1}(l^{-1}\iota(\mathsf z_w^{\,e})l)k$ lies
in $K_j$; it is $1$ outside $v$, so $\mathsf u_e\in J_j$, and $\psW_j(n\iota(\mathsf z_w^{\,e})n^{-1}
h)=\psW_j(h\mathsf u_e)=\psW_j(h)$. On the other hand $n\iota(\mathsf z_w^{\,e})n^{-1}=n'\iota(
\mathsf z_w^{\,e})$ with $n'=n\cdot\iota(\mathsf z_w^{\,e})n^{-1}\iota(\mathsf z_w^{\,e})^{-1}\in
N_j(\auA)=N_j^{\ge2}(\auA)$ ($j$ is even) and $\psi_j(n')=\psi_j(n)\psi_j(n)^{-1}=1$, because
$\psi_j$ is invariant under conjugation by $R_j(\auA)$; so $\psW_j(n\iota(\mathsf z_w^{\,e})n^{-1}h)=
\psW_j(\iota(\mathsf z_w^{\,e})h)$. By Lemma~\ref{ps:double}(a), applied $j/2$ times,
$\psW_j(\iota(g)h)=l_h(g_\infty^{-1}\cdot(C_h\varphi)(g_f))$ for $g\in\auU(X_j)(\auA)$, where $C_h$
is a composite of lifts \eqref{pk:eq-lift}. Each lift commutes with the translation by the central
element $\mathsf z_w$: the pair $(\mathsf z_w,\mathsf z_w)$ acts trivially on $\auHom(X_i,X_{i+1})$,
and the determinant twists are trivial on it since $\mu_{i+1}$ and $\upsilon$ are unramified above
$N$. Therefore $(N-1)\psW_j(h)=\sum_{e=0}^{N-2}\psW_j(\iota(\mathsf z_w^{\,e})h)=l_h\bigl((C_h
\sum_e\mathsf z_w^{\,e}\varphi)(1)\bigr)=0$ for $\varphi\in\auSram$.

\emph{(R3) The lattice $\Lambda_0$ at the chosen places (Lemma~\ref{ps:fin}(3)).} Let $w$ be a
chosen place. There $N$ is a unit, $\alpha=1$, $\mathsf a_w=1$,
$\mathsf a_{\bar w}=\varpi^{-1}$, so $u_j=\varpi^{e_j}$ at $w$ and $\varpi^{-e_j}$ at $\aubar w$,
$e_j=\lceil(m-j)/2\rceil$. By \eqref{ps:eq-lattice}, for $j$ even $y/\aubar{\pi_N}+u_jz/N$ and
$u_jz$ span the same module as $y$ and $u_jz$; so $\psM_j=\widehat\auO^{\,n_j}\oplus u_j\bigoplus_a
\widehat\auO f_a$ at $w$ and $\aubar w$, the orthogonal sum of its parts in $X_j$ and in the
string, and likewise for $\psM_{j+2}$ if $j+2<m$. The lattice $\Lambda_0$ is the image in
$X^-\oplus X''$ of the vectors of $\psM_j\oplus\psM_{j+2}$ without components of negative grade and
with component of grade $0$ in $X^-\oplus X''\oplus\ell$; for a split lattice this image is
$\widehat\auO^{\,n_j}\oplus\widehat\auO^{\,n_{j+2}}$, whatever $u_j,u_{j+2}$ are. Thus $\Lambda_0$
is standard at the chosen places for all pairs with $j+2<m$. For $j+2=m$: $\psM_{m-2}=\widehat
\auO^{\,n_{m-2}}\oplus\widehat\auO u_{m-2}z$ and $z=\pi_Ne_{\mathrm{last}}$ with $\pi_N$ a unit, so
$\Lambda_0=\widehat\auO^{\,n_{m-2}}\oplus\widehat\auO^{\,n_m-1}\oplus u_{m-2}\widehat\auO
e_{\mathrm{last}}$, as stated in Lemma~\ref{ps:fin}.

\emph{(R4) The denominators in Lemma~\ref{ps:fin}.} For a chosen place $w$ above $p\in\{p_0,p_1\}$,
$\#_w=p^{\,\operatorname{rk}Y_j\operatorname{rk}Y_{j+1}}$ divides $p^{n_m^2}$; there are $f$ chosen
places above each of $p_0,p_1$, so $\prod_{w\text{ chosen}}\#_w$ divides $M^{fn_m^2}=B$ for each of
the $m$ data $\Upsilon_j$. By (R3), $\mathsf T_i$ is trivial unless $2i+2=m$, and then
$\#\mathsf T_i=\prod_{w}p^{\,n_{m-1}}$ divides $B$. So the product of all denominators at the primes
not dividing $q$ divides $B^{m+1}$, hence $B^{2m}$.

\emph{(R5) Constants in the proof of Proposition~\ref{au:kernel}.} $C_0$ is the constant of
Proposition~\ref{pk:conj}(ii); it depends on $q_0,d$ only, and it is a term of the constant $C_P$ in
the definition of admissible data, so that $C_0\le\log\log b$ (\S\ref{pz:ss-const},
Lemma~\ref{pz:C0}; there is no circularity, $C_0$ being defined from $q_0$ and $K$ alone). One
may assume $\sigma|_K=1$: if $\sigma|_K\ne1$, then $c\circ\sigma$ fixes $K$ and $|c\sigma(w)|=|\sigma(w)|$. The
bound $q_1\le x^{12}$ is Lemma~\ref{an:lem-sizes}(S3).

\emph{(R6) Lemma~\ref{ps:level}.} The factor $\upsilon'_w$ is trivial on units because $\upsilon$
is unramified at every finite place (\S\ref{pk:ss-conv}).

\subsubsection{Integrality of the branching map: proof of Lemma~\ref{ps:branching}(c)}
\label{pk:ss-branch}

All modules below are tensor products over $\tau\in\auPl$ and all statements are multiplicative
in $\tau$; we write them for one $\tau$ and one group $\auGL_N$. Let $\mu$ be a diagram with
columns $\mathfrak c$ of heights $h_{\mathfrak c}\le N$ and $|\mu|=\sum h_{\mathfrak c}$ boxes.

\emph{The tensor model.} Let $\mathbb T_\mu(R)=\bigotimes_{\mathfrak c}\Lambda^{h_{\mathfrak c}}R^N$
for a ring $R$, with the basis $e_U=\bigotimes_{\mathfrak c}e_{U_{\mathfrak c}}$ indexed by the
fillings $U$ (a subset $U_{\mathfrak c}\subseteq[N]$ of $h_{\mathfrak c}$ rows for each column), and
the symmetric bilinear form $(\ ,\,)$ for which this basis is orthonormal. For a matrix $y$ with
columns $y_1,\dots,y_N\in R^N$ put $\pkev_y=\bigotimes_{\mathfrak c}y_1\wedge\dots\wedge
y_{h_{\mathfrak c}}$; then $(e_U,\pkev_y)=b_U(y):=\prod_{\mathfrak c}\Delta_{U_{\mathfrak c},
[h_{\mathfrak c}]}(y)$, the product of minors of \S\ref{ps:ss-restrict}. So the multiplication map
is $\pkm(u)(y)=(u,\pkev_y)$, and it is equivariant: $\pkm(gu)(y)=\pkm(u)(g^ty)$. Let
$\mathsf C_\mu\subseteq\mathbb T_\mu(\C)$ be the span of the $\pkev_y$, $y\in M_N(\C)$. It is
defined over $\Q$, so $\mathbb T_\mu(\C)=\mathsf C_\mu\oplus\mathsf C_\mu^\perp$, and
$\ker\pkm=\mathsf C_\mu^\perp$ by the formula for $\pkm$; let $\pi_\mu$ be the orthogonal projection
onto $\mathsf C_\mu$. Thus $\pkm=\pkm|_{\mathsf C_\mu}\circ\pi_\mu$ and $\pkm\colon\mathsf C_\mu\to
W_\mu$ is an isomorphism. The pull-back of the invariant form of $W_\mu$ to $\mathsf C_\mu$ is
invariant under $\auU(N)$, and $W_\mu$ is irreducible under $\auGL_N(\C)$~\cite[\S8.1, Thm.~1 and its corollary; \S8.2, Thm.~2]{FultonYT}, hence under its Zariski-dense subgroup $\auU(N)$; so it is a multiple of
the standard hermitian form, and the multiple is $1$ because $e_{U_0}=\pkev_1$ ($U_0$ the filling
by the leading rows) has norm $1$ and $\pkm(e_{U_0})=f_\mu:=\prod\Delta_{[h_{\mathfrak c}],
[h_{\mathfrak c}]}$ has norm $1$. Consequently, for $w\in\mathsf C_\mu$,
$\langle\pkm w,b_U\rangle=\langle w,\pi_\mu e_U\rangle=(w,e_U)$, and
\begin{equation}\label{pk:eq-dual}
\Sigma_\mu=\pkm(\mathbb T_\mu(\Z)),\qquad W_\mu(\Z)=\Sigma_\mu^\#=\pkm\bigl(\mathsf C_\mu\cap
\mathbb T_\mu(\Z)\bigr).
\end{equation}
We also use that the standard products of minors are a $\Z$-basis of $\Sigma_\mu$ and remain
linearly independent modulo every prime~\cite[Thm.~2.1]{DEP}, \cite[\S8.1, Thm.~1 and its corollary]{FultonYT}; hence
\begin{equation}\label{pk:eq-sat}
\Sigma_\mu=(\Sigma_\mu\otimes\Q)\cap\Z[y_{ab}] .
\end{equation}

\begin{lemma}[The hook number]\label{pk:hook}
Let $\langle P,Q\rangle_F=(P(\partial)\aubar Q)(0)=\sum_\alpha\alpha!\,P_\alpha\aubar{Q_\alpha}$ be the
Fischer form on polynomials in the entries of $y$, and $H_\mu=\langle f_\mu,f_\mu\rangle_F$. Then
$H_\mu$ is a positive integer, $H_\mu\le|\mu|^{3|\mu|}$ $(0^0=1)$, the invariant form of $W_\mu$ is
$H_\mu^{-1}\langle\ ,\,\rangle_F$, and $H_\mu\Sigma_\mu\subseteq\Sigma_\mu^\#$.
\end{lemma}
\begin{proof}
The Fischer form is positive definite and invariant under unitary substitutions of the variables,
in particular under $P\mapsto P(g^ty)$, $g\in\auU(N)$; on the irreducible module $W_\mu$ it is
therefore $H_\mu$ times the invariant form. Products of minors have integer coefficients, so
$\langle b,b'\rangle_F\in\Z$ and $\langle H_\mu b,b'\rangle=\langle b,b'\rangle_F\in\Z$ for
$b,b'\in\Sigma_\mu$. Finally $f_\mu$ is homogeneous of degree $L=|\mu|$ and the sum of the absolute
values of its coefficients is at most $\prod_{\mathfrak c}h_{\mathfrak c}!\le\prod h_{\mathfrak
c}^{h_{\mathfrak c}}\le L^L$; so $H_\mu\le L!\,(L^L)^2\le L^{3L}$.
\end{proof}

\begin{lemma}[Leading rows]\label{pk:leading}
Let $0\le p<N$ and let $\mu^\flat$ be the diagram, for $\auGL_{N-p}$, with the columns of heights
$\max(h_{\mathfrak c}-p,0)$. Let $(r_1,\dots,r_p)$ be distinct rows and $U$ a filling of $\mu$ such
that $U_{\mathfrak c}\supseteq\{r_1,\dots,r_p\}$ if $h_{\mathfrak c}\ge p$ and $U_{\mathfrak c}=
\{r_1,\dots,r_{h_{\mathfrak c}}\}$ if $h_{\mathfrak c}<p$. Then $H_{\mu^\flat}\,b_U\in\Sigma_\mu^\#$.
\end{lemma}
\begin{proof}
Permutation matrices preserve $\Sigma_\mu$, $\Sigma_\mu^\#$ and the form, and map $b_U$ to $\pm$
the product of minors of the permuted filling; so let $r_i=i$. Then $U=\pkLW(T)$ for a unique
filling $T$ of $\mu^\flat$: $U_{\mathfrak c}=[p]\cup(p+T_{\mathfrak c})$ if $h_{\mathfrak c}\ge p$.
Define $\pkLW\colon\mathbb T_{\mu^\flat}(\Z^{N-p})\to\mathbb T_\mu(\Z^N)$ by $e_T\mapsto
e_{\pkLW(T)}$, for all fillings $T$ of $\mu^\flat$. (1) For $y\in M_{N-p}$,
$\pkLW(\pkev_y)=\pkev_{\operatorname{diag}(1_p,y)}$, since the first $h$ columns of
$\operatorname{diag}(1_p,y)$ have exterior product $e_{[h]}$ for $h\le p$ and $e_{[p]}\wedge(y_1
\wedge\dots\wedge y_{h-p})$ for $h>p$. Hence $\pkLW(\mathsf C_{\mu^\flat})\subseteq\mathsf C_\mu$.
(2) For $x=\bigl(\begin{smallmatrix}A&B\\C&D\end{smallmatrix}\bigr)$ with $A\in\auGL_p$ and
$S(x)=D-CA^{-1}B$, the formula for the determinant of a block matrix gives
$\Delta_{[p]\cup(p+I),[p+k]}(x)=\det A\cdot\Delta_{I,[k]}(S(x))$ for $I\subseteq[N-p]$, $|I|=k$ (the
first $k$ columns of $S(x)$ depend only on the first $p+k$ columns of $x$). Hence, for every
filling $T$ of $\mu^\flat$, $b_{\pkLW(T)}(x)=\mathsf F(x)\,b^\flat_T(S(x))$ with $\mathsf F(x)=
\prod_{h_{\mathfrak c}<p}\Delta_{[h_{\mathfrak c}],[h_{\mathfrak c}]}(x)\cdot\det A^{\,\#\{\mathfrak
c:h_{\mathfrak c}\ge p\}}$ independent of $T$; by linearity $\pkm(\pkLW u)(x)=\mathsf F(x)\,
\pkm^\flat(u)(S(x))$ for all $u$, on a dense open set. So $\pkLW(\ker\pkm^\flat)\subseteq\ker\pkm$,
that is $\pkLW(\mathsf C_{\mu^\flat}^\perp)\subseteq\mathsf C_\mu^\perp$; with (1),
$\pi_\mu\circ\pkLW=\pkLW\circ\pi_{\mu^\flat}$. (3) By Lemma~\ref{pk:hook} and \eqref{pk:eq-dual} for
$\mu^\flat$, $H_{\mu^\flat}\pi_{\mu^\flat}(e_T)\in\mathsf C_{\mu^\flat}\cap\mathbb T_{\mu^\flat}(\Z)$.
Applying $\pkLW$, which preserves integrality: $H_{\mu^\flat}\pi_\mu(e_U)\in\mathsf C_\mu\cap\mathbb
T_\mu(\Z)$, and $\pkm$ of this vector is $H_{\mu^\flat}b_U$; conclude by \eqref{pk:eq-dual}.
\end{proof}

\begin{proposition}[Lemma~\ref{ps:branching}(c)]\label{pk:branching}
Let $\lambda=\lambda(X_m)$, $n=n_m$, and let $\lambda^{\circ\flat}_\tau$ be $\lambda^\circ_\tau$
without its first $k-1$ rows. Put $c_{\mathrm{br}}=\prod_\tau H_{\lambda^{\circ\flat}_\tau}$ and
$\mathrm{br}=c_{\mathrm{br}}\,i^*$, with $i$ as below. Then {\rm(a), (b), (c)} of
Lemma~\ref{ps:branching} hold, and $c_{\mathrm{br}}\le L^{3L}$.
\end{proposition}
\begin{proof}
We write the proof for one $\tau$. Call a column of $\lambda$ \emph{kept} or \emph{lost}
according as it keeps its height in $\lambda^\circ$ or loses a box; the kept columns have heights
$\ge k$. Let $\varepsilon\colon\mathbb T_{\lambda^\circ}(\Z^{n-1})\to\mathbb T_\lambda(\Z^n)$ send
$e_{U^\circ}$ to the tensor with factor $e_{U^\circ_{\mathfrak c}}$ in a kept column and
$e_n\wedge e_{U^\circ_{\mathfrak c}}$ in a lost one. It is isometric, integral and equivariant for
$\auGL_{n-1}$. Put $i=\pkm\circ\varepsilon\circ(\pkm^\circ|_{\mathsf C_{\lambda^\circ}})^{-1}\colon
W_{\lambda^\circ}\to W_\lambda$. Then $i$ is $\auGL_{n-1}$-equivariant, $i(f_{\lambda^\circ})=\pm v$
with $v$ the vector of the proof of Lemma~\ref{ps:branching}(a), and under the identifications
\eqref{pk:eq-dual} $i=\pi_\lambda\varepsilon|_{\mathsf C_{\lambda^\circ}}$, $i^*=\pi_{\lambda^\circ}
\varepsilon^t|_{\mathsf C_\lambda}$; so $\|i^*\|\le1$. This gives (a) and (b) as in
\S\ref{ps:ss-restrict}. By duality, (c) is equivalent to
\begin{equation}\label{pk:eq-c}
c_{\mathrm{br}}\cdot i(b^\circ_{T'})\in\Sigma_\lambda\qquad\text{for every filling }T'\text{ of }
\lambda^\circ .
\end{equation}
\emph{The kernel.} For $x\in M_{n-1}(R)$, whose columns are regarded as vectors of $R^n$ with last
coordinate $0$, and $y\in M_n(R)$ let $\mathcal K(x,y)=(\varepsilon\,\pkev^\circ_x,\pkev_y)$. Then
\begin{equation}\label{pk:eq-kernelbr}
\mathcal K(x,y)=(\pkev^\circ_x,\pi_{\lambda^\circ}\varepsilon^t\pkev_y),\qquad
(\pi_{\lambda^\circ}\varepsilon^t\pkev_y,e_{T'})=(\pkev_y,\varepsilon\pi_{\lambda^\circ}e_{T'})=
i(b^\circ_{T'})(y).
\end{equation}
\emph{Special $y$.} All columns of $\lambda$ have height $\le n_{m-1}\le n-2$. Let the last row
of $y$ be $(\mathsf t,0,\dots,0)$ and let $g\in M_{n-1}(R)$ be the matrix of the first $n-1$ rows
and columns of $y$, with columns $g_1,\dots,g_{n-1}$ (the last column of $y$ does not occur in
$\pkev_y$). In a kept column of height $h$ the factor of
$\mathcal K$ is $\det(x_a\cdot g_b)_{a,b\le h}=\Delta_{[h],[h]}(g^tx)$; in a lost column of height
$h$ it is the determinant of the matrix with first row $(\mathsf t,0,\dots,0)$ and rows
$(x_a\cdot g_b)_{b\le h}$, $a<h$, that is $\mathsf t\,\Delta_{\{2,\dots,h\},[h-1]}(g^tx)$. So
\[
\mathcal K(x,y)=\mathsf t^{\,\ell}\,\mathsf P(g^tx),\qquad\mathsf P=\prod_{\text{kept}}\Delta_{[h],[h]}
\cdot\prod_{\text{lost}}\Delta_{\{2,\dots,h\},[h-1]}\ \in\Sigma_{\lambda^\circ},
\]
with $\ell$ the number of lost columns. The product $\mathsf P$ is $b^\circ_U$ for a filling $U$
of $\lambda^\circ$ to which Lemma~\ref{pk:leading} applies with $p=k-1$ and the rows
$(2,3,\dots,k)$: a kept column has height $h\ge k$ and rows $[h]\supseteq\{2,\dots,k\}$; a lost
column, of height $h-1$ in $\lambda^\circ$, has rows $\{2,\dots,h\}$, which contain $\{2,\dots,k\}$
if $h-1\ge k-1$ and are the first $h-1$ of the rows $(2,\dots,k)$ otherwise. Hence $c\,\mathsf P=
\pkm^\circ(w_0)$ with $c=H_{\lambda^{\circ\flat}}$ and $w_0\in\mathsf C_{\lambda^\circ}\cap\mathbb
T_{\lambda^\circ}(\Z)$, and $c\,\mathcal K(x,y)=\mathsf t^{\,\ell}(w_0,\pkev^\circ_{g^tx})=(\mathsf
t^{\,\ell}\,g\,w_0,\pkev^\circ_x)$, where $g$ acts on the tensors.

\emph{General $y$.} Work over $R=\Z[y_{ab}]$ and let $y'$ have the columns $y'_1=y_1$,
$y'_b=y_{n1}y_b-y_{nb}y_1$ $(b\ge2)$. Its last row is $(y_{n1},0,\dots,0)$, and
$\pkev_{y'}=y_{n1}^{\,\mathsf M}\pkev_y$ with $\mathsf M=\sum_{\mathfrak c}(h_{\mathfrak c}-1)$. By
the special case and \eqref{pk:eq-kernelbr}, with $g'$ the matrix of the first $n-1$ rows and
columns of $y'$, $y_{n1}^{\,\mathsf M}\,c\,\pi_{\lambda^\circ}
\varepsilon^t\pkev_y$ and $y_{n1}^{\,\ell}g'w_0$ are two vectors of $\mathsf C_{\lambda^\circ}
\otimes\Q(y)$ with the same pairing against all $\pkev^\circ_x$; so they are equal, and the second
has coordinates in $R$. Therefore $y_{n1}^{\,\mathsf M}\cdot c\,i(b^\circ_{T'})(y)\in\Z[y_{ab}]$,
and as $c\,i(b^\circ_{T'})\in\Q[y_{ab}]$ and $y_{n1}$ is a prime element of $\Z[y_{ab}]$,
$c\,i(b^\circ_{T'})\in\Z[y_{ab}]$. Since $i(b^\circ_{T'})\in W_\lambda=\Sigma_\lambda\otimes\C$ has
rational coordinates, \eqref{pk:eq-sat} gives \eqref{pk:eq-c}.

\emph{The size.} $|\lambda^{\circ\flat}_\tau|$ is at most the number $L_\tau$ of boxes of
$\lambda^\circ_\tau$ in the columns of height $\ge k$, so $c_{\mathrm{br}}\le\prod_\tau
L_\tau^{3L_\tau}\le L^{3L}$ by Lemma~\ref{pk:hook}.
\end{proof}

\subsubsection{One curve with complex multiplication}\label{pk:ss-curve}

We use the notation of \S\ref{pa:ss-curve}: $\Lambda$ with $\auO_K\Lambda=\Lambda$, of area $A$,
the point $t$ with $\wp''(t)=2$, $a_1$, the curve $\mathcal E$, $x,y$, $\theta$, $g=-\log\theta$,
$H(u,w)=\boldsymbol\pi\aubar uw/A$, and $s_2$. The transformation law of $\vartheta$ quoted there
is equivalent to $\zeta(z+\omega)=\zeta(z)+\eta(\omega)$ with
$\eta(\omega)=s_2\omega+(\boldsymbol\pi/A)\aubar\omega$. For $u\in\C$ the \emph{twisted translate} is
\[
\theta_u(z)=\exp\bigl(-\tfrac{\boldsymbol\pi}A(z\aubar u+\tfrac12|u|^2)\bigr)\,\theta(z+u);
\]
$|\theta_u(0)|$ is $\Lambda$-periodic in $u$, and $\theta_{u+\omega}=\chi(\omega)e^{i(\boldsymbol\pi/A)
\Iml(\aubar\omega u)}\theta_u$ for $\omega\in\Lambda$ by \eqref{pa:eq-theta}. We call
$(\Lambda,t,a_1)$ a \emph{normalised lattice} if $t$ is the admissible point of order two
(\S\ref{pf:ss-adm}). A power series $F(z)$ in one or several variables
is \emph{$\partial$-integral} at a homomorphism $\C\to\C_\ell$ if all its derivatives at $0$ are in
$\pao$; sums and products of $\partial$-integral series are $\partial$-integral, and so is
$\Phi\circ F$ if $F(0)=0$ and $\Phi\in\pao[[X]]$ or $\Phi=\exp$.

\begin{proposition}\label{pk:curve}
Let $(\Lambda,t,a_1)$ be a normalised lattice.
\begin{enumerate}
\item[(C1)] {\rm(Conjugation.)} For $\sigma\in\auAut(\C/K)$ there is a normalised lattice
$(\Lambda^\sigma,t^\sigma,\sigma a_1)$, unique, such that $\sigma$ maps the Taylor series at $0$ of
$x$, $y$, $\theta$ to those of $\Lambda^\sigma$. One has $g_i(\Lambda^\sigma)=\sigma g_i(\Lambda)$
$(i=2,3)$ and $s_2(\Lambda^\sigma)=\sigma s_2(\Lambda)$.
\item[(C2)] {\rm(Integrality.)} For every prime $\ell\ne d$ and every $\C\to\C_\ell$, $x$, $y$ and
$\theta$ are $\partial$-integral.
\item[(C3)] {\rm(Translates.)} Let $u\in\C$ with $\mathsf Nu\in\Lambda$, $\mathsf N\ge1$,
$u\notin t+\Lambda$.
\begin{itemize}
\item[(i)] Let $\ell$ be an odd prime not dividing $\mathsf N$ and $\C\to\C_\ell$. Then
$\theta_u(0)\in\pao$ (this holds at every prime, Proposition~\ref{pf:values}) and $\theta_u(z)=\theta_u(0)\,\theta(z)\,\mathsf G_u(z)$ with
$\mathsf G_u=\Gamma_u(y(z))$, $\Gamma_u\in1+Y\pao[[Y]]$; and $\mathsf G_u$ is $\partial$-integral.
If $\ell\ne d$, $\theta_u$ is $\partial$-integral.
\item[(ii)] Let $\Lambda=\Omega_0\auO_K$ and $\sigma\in\auAut(\C/K)$. There are $\Omega_\sigma\in\C^\times$ and an ideal $\mathfrak c$ of $\auO_K$ with
$\Lambda^\sigma=\Omega_\sigma\mathfrak a$, $\mathfrak a=\mathfrak c\aubar{\mathfrak c}^{-1}$, and a finite idele
$\beta_\sigma$ of $K$ generating $\mathfrak a\widehat\auO_K$, such that for every $\mathsf x\in K$, every $n\ge1$
with $n\mathsf x\in\auO_K$ and every $\mathsf x'\in K$ with $\mathsf x'-\beta_\sigma\mathsf x\in2n\mathfrak a\widehat\auO_K$,
$\sigma$ maps the Taylor series of $\theta_{\mathsf x\Omega_0}$ to that of the twisted translate
$\theta^\sigma_{\mathsf x'\Omega_\sigma}$ of the theta function of $\Lambda^\sigma$.
\end{itemize}
\item[(C4)] {\rm(Chowla--Selberg.)} There is an algebraic number $\mathsf u\ne0$, a unit at every prime, with $\operatorname{area}(\Lambda^\sigma)\,|\sigma\mathsf u|^2=2\boldsymbol\pi^2e^{c_{CS}}/
\sqrt d$ for all $\sigma\in\auAut(\C/K)$, where $c_{CS}=\gamma-\Rel\frac{L'}L(1,\chi_K)$.
\end{enumerate}
\end{proposition}
\begin{proof}
(C1). Let $\Lambda^\sigma$ be the lattice with invariants $\sigma g_2,\sigma g_3$ (it exists, as
$\sigma(g_2^3-27g_3^2)\ne0$). The number $e_1=\wp(t)$ is a root of $4X^3-g_2X-g_3$ with
$6e_1^2-g_2/2=\wp''(t)=2$; so $\sigma e_1=\wp_{\Lambda^\sigma}(t^\sigma)$ for a point $t^\sigma$ of
order two with $\wp''(t^\sigma)=2$, and $(\sigma a_1)^2=12\sigma e_1$. For $\alpha\in\auO_K$ there
is a rational function $R_\alpha$ with $\wp(\alpha z)=R_\alpha(\wp(z))$; it is determined by the
Laurent series of $\wp$, whose coefficients are polynomials in $g_2,g_3$ over $\Q$, so
$\sigma R_\alpha$ satisfies the same identity for $\Lambda^\sigma$ ($\sigma\alpha=\alpha$), and
$\alpha\Lambda^\sigma\subseteq\Lambda^\sigma$. The point $t^\sigma$ is admissible
(\S\ref{pf:ss-adm}): for $\beta\in\{\omega,\aubar\omega\}\subset\auO_K$, $\beta t\notin\Lambda$
means $\wp(\beta t)\ne\infty$, i.e.\ $R_\beta(e_1)\ne\infty$ ($\wp(t)=e_1$ is finite as $t\notin\Lambda$);
applying $\sigma$ to the numerator and denominator of $R_\beta$, this is equivalent to
$(\sigma R_\beta)(\sigma e_1)\ne\infty$, i.e.\ $\wp_{\Lambda^\sigma}(\beta t^\sigma)\ne\infty$, i.e.\
$\beta t^\sigma\notin\Lambda^\sigma$. By \eqref{pa:eq-curve}, $x$ and $y$ are the solutions
vanishing at $0$ of a differential system with coefficients in $\Z[a_1]$; so their Taylor
coefficients lie in $\Q[a_1]$ and are mapped by $\sigma$ to those of $\Lambda^\sigma$. By
\eqref{pa:eq-g2}, $\theta=\exp(-g)$ with $g''=x+s_2+a_1^2/12$, $g(0)=g'(0)=0$, and it remains to
show $\sigma s_2(\Lambda)=s_2(\Lambda^\sigma)$. This follows from the identity
\begin{equation}\label{pk:eq-s2}
\alpha(\alpha-\aubar\alpha)\,s_2=\sum_{0\ne u\in\alpha^{-1}\Lambda/\Lambda}\wp(u)\qquad(\alpha\in
\auO_K,\ \alpha\ne0),
\end{equation}
because the right side is minus the second coefficient of the monic denominator
$\prod_{u\ne0}(X-\wp(u))$ of $R_\alpha$, and $\alpha(\alpha-\aubar\alpha)\ne0$ for
$\alpha=\sqrt{-d}$. To prove \eqref{pk:eq-s2}: the function $\alpha\zeta(\alpha z)-\sum_{u\in
\alpha^{-1}\Lambda/\Lambda}\zeta(z+u)$ is entire, and changes under $z\mapsto z+\omega$ by
$\alpha\eta(\alpha\omega)-\auNm\alpha\cdot\eta(\omega)=s_2(\alpha^2-\alpha\aubar\alpha)\omega$; so
it is $s_2\alpha(\alpha-\aubar\alpha)z$ plus a constant. Differentiate and let $z\to0$; the terms
$-\alpha^2\wp(\alpha z)+\wp(z)$ tend to $0$. Uniqueness of $(\Lambda^\sigma,t^\sigma)$ is clear
from the series of $x$.

(C2). $a_1$ is an algebraic integer (Lemma~\ref{pa:model}), and $d/dz$ preserves $\Z[a_1][x,y]$ by
\eqref{pa:eq-curve}; so $x$, $y$ are $\partial$-integral at every prime. Since $g^{(n)}(0)=
x^{(n-2)}(0)$ for $n\ge3$ and $\theta=\exp(-g)$, it suffices to show $g''(0)=s_2+a_1^2/12\in\pao$.
By \eqref{pk:eq-s2} with $\alpha=\sqrt{-d}$, $\wp(u)=x(u-t)+a_1^2/12$ and $\auNm\alpha=d$,
\[
-2d\,g''(0)=\sum_{0\ne u\in\alpha^{-1}\Lambda/\Lambda}x(u-t)\ -\ \frac{d+1}{12}\,a_1^2 .
\]
The numbers $x(u-t)$ are the $x$-coordinates of the non-zero points of order $d$ of $\mathcal E$;
they are $\ell$-integral for $\ell\ne d$, because the points with non-integral $x$ form a group
isomorphic to the group of points of a formal group over the integers of a finite extension of
$\Q_\ell$, which has no torsion of order prime to $\ell$~\cite[Chapter~VII, Proposition~2.2, and Chapter~IV, Proposition~3.2(b)]{SilvermanAEC}.
They are equal in pairs ($x(-u-t)=x(u-t)$, $u\not\equiv-u$), and $24\mid d+1$ because $d\equiv
q_0/2-1\pmod{q_0}$ and $48\mid q_0$. So $-d\,g''(0)\in\pao$, and $g''(0)\in\pao$ for $\ell\ne d$.

(C3)(i). Put $R_u=\theta_u/\theta$. By \eqref{pa:eq-theta}, $R_u(z+\omega)=e^{2i(\boldsymbol\pi/A)
\Iml(\aubar\omega u)}R_u(z)$, and $\Iml(\aubar\omega\mathsf Nu)\in A\Z$; so $R_u^{\mathsf N}$ is
$\Lambda$-periodic. Its only poles are at $z\in t+\Lambda$, the origin $O$ of $\mathcal E$, of order $\mathsf N$, and its only zeros are at $z\in t-u+\Lambda$, a point $Q_1$ of $\mathcal E$
whose order divides $\mathsf N$, and they have order $\mathsf N$ (if $u\in\Lambda$, $R_u$ is
constant). Thus $R_u^{\mathsf N}=\theta_u(0)^{\mathsf N}\mathsf f(x,y)/
\mathsf f(0,0)$ for a polynomial $\mathsf f$ with divisor $\mathsf N[Q_1]-\mathsf N[O]$; here
$\mathsf f(0,0)\ne0$ because $z=0$ is the point $P_0=(0,0)\ne Q_1$ ($u\notin t+\Lambda$). Scale
$\mathsf f$ so that its coefficients are integral and not all in $\pam$. The curve has good
reduction at $\ell$ (Lemma~\ref{pa:model}), so the divisor of the reduction $\tilde{\mathsf f}\ne0$
is the reduction $\mathsf N[\tilde Q_1]-\mathsf N[\tilde O]$ of the divisor of $\mathsf f$. The
point $Q_1-P_0$ is not zero and has order prime to $\ell$ (as $\ell\nmid2\mathsf N$), so it does
not reduce to the origin; hence $\tilde Q_1\ne\tilde P_0$ and $\mathsf f(0,0)$ is a unit. As $x\in
y^2\pao[[y]]$, $\mathsf f(x,y)/\mathsf f(0,0)\in1+y\pao[[y]]$, and $\Gamma_u$ is its $\mathsf N$-th
root with constant term $1$, which has coefficients in $\pao$ because $\mathsf N$ is a unit at
$\ell$. It is $\partial$-integral in $z$ because $y$ is (C2). Finally $\theta_u(0)\in\pao$ by Proposition~\ref{pf:values}.
The last assertion of (i) follows from (C2).

(C4). Let $h$ be the class number of $K$; it is odd, so every ideal class contains an ideal
$\mathfrak c\aubar{\mathfrak c}^{-1}$. Choose fractional ideals $\mathfrak c_1,\dots,\mathfrak c_h$
representing the classes with $\prod_i\mathfrak c_i=(1)$, numbers $\mu_i\in\C^\times$ such that
$\Lambda_i=\mu_i\mathfrak c_i\Lambda$ is the lattice of a normalised lattice $(\Lambda_i,t_i,a_{1,i})$,
and let $\mathsf u$ be a root of $\mathsf u^{12h}=\prod_i\mu_i^{12}\Delta(\Lambda_i)$, where
$\Delta=g_2^3-27g_3^2=a_1^4-64$. The numbers $\mu_i^2$ are algebraic. For $\sigma\in\auAut(\C/K)$,
$\Lambda_i^\sigma=\sigma(\mu_i)\mathfrak c_i\Lambda^\sigma$: if $\lambda\in\mu_i\mathfrak c_i$,
multiplication by $\lambda$ maps $\Lambda$ into $\Lambda_i$, so $\wp_{\Lambda_i}(\lambda z)$ is a
rational function of $\wp_\Lambda(z)$ determined by the Laurent series, and conjugating as in
(C1) gives $\sigma(\lambda)\Lambda^\sigma\subseteq\Lambda_i^\sigma$; the same for the inverse
ideal gives equality. The lattices $\Lambda_i^\sigma$ represent the $h$ classes of lattices with
multiplication by $\auO_K$. For such a lattice $\Z\omega_1+\Z\omega_2$, $\tau=\omega_2/\omega_1$,
one has $\operatorname{area}\cdot|\Delta|^{1/6}=4\boldsymbol\pi^2\,\Iml\tau\,|\eta(\tau)|^4$, and
Kronecker's limit formula together with $\zeta_K=\zeta\cdot L(\cdot,\chi_K)$ and the class number
formula gives the following: by \cite[\S12, the display preceding~(2)]{ChowlaSelberg}, with
$\tau_j=(b_j+i\sqrt d)/(2a_j)$ running over the reduced forms of discriminant $-d$ and $w=2$,
$\sum_j\log|\Delta(\tau_j)|=6\bigl(h\gamma+\log(\prod_ja_j/d^h)\bigr)-\frac{3w}{\boldsymbol\pi}\sqrt d\,
L'(1,\chi_K)$, where $\Delta=\eta^{24}$; since $\Iml\tau_j=\sqrt d/(2a_j)$ and
$L(1,\chi_K)=2\boldsymbol\pi h/(w\sqrt d)$, dividing by $6h$ and adding
$\frac1h\sum_j\log\Iml\tau_j=\log(\sqrt d/2)-\frac1h\log\prod_ja_j$ gives $\gamma-\frac{w\sqrt d}{2\boldsymbol\pi h}L'(1,\chi_K)-\log d+\log(\sqrt d/2)=\gamma-\frac{L'}{L}(1,\chi_K)-\log(2\sqrt d)$, the terms in $a_j$ cancelling; here $\log|\eta(\tau)|^4=\frac16\log|\Delta(\tau)|$, and $\Iml\tau\,|\eta(\tau)|^4$ is invariant under $\mathrm{SL}_2(\Z)$ and under $\tau\mapsto-\aubar\tau$, so it depends only on the homothety class of $\Z+\Z\tau$, and the reduced forms $[a_j,\pm b_j,c_j]$ of discriminant $-d$ correspond bijectively to the $h$ homothety classes of lattices with multiplication by $\auO_K=\Z[\frac{1+\sqrt{-d}}2]$ via $\tau_j\mapsto\Z+\Z\tau_j$. This proves (for Kronecker's limit formula see also
\cite[Chapter~20, \S4]{LangEF})
\[
\frac1h\sum_{\text{classes}}\log\bigl(\Iml\tau\,|\eta(\tau)|^4\bigr)=c_{CS}-\log(2\sqrt d).
\]
Hence $\prod_i\operatorname{area}(\Lambda_i^\sigma)|\Delta(\Lambda_i^\sigma)|^{1/6}=(2\boldsymbol\pi^2
e^{c_{CS}}/\sqrt d)^h$; the left side is $\prod_i|\sigma\mu_i|^2\auNm(\mathfrak c_i)\operatorname{area}
(\Lambda^\sigma)|\sigma\Delta(\Lambda_i)|^{1/6}=(\operatorname{area}(\Lambda^\sigma)|\sigma\mathsf
u|^2)^h$. For the unit property let $\ell$ be any prime. The $\Delta(\Lambda_i)$ are units
(Lemmas~\ref{pa:model}, \ref{pf:two}), the $t_i$ being admissible. If $\lambda\Lambda\subseteq\Lambda'$ for normalised lattices, then
$z\mapsto\lambda z$ is a homomorphism between elliptic curves with good reduction which multiplies
the invariant differentials $dx/(2y+a_1x)=dz$ by $\lambda$; these differentials generate the
modules of invariant differentials of the smooth models, and the homomorphism extends to these
models, which are N\'eron models~\cite[\S1.2, Definition~1, p.~12, and Proposition~8, p.~15]{BLR} (\cite[IV, Corollary~6.3, p.~329, and the N\'eron mapping property, p.~319]{SilvermanAT}), so $\lambda\in\pao$. Applying this to
$\lambda\in\mu_i\mathfrak c_i$ and to $\lambda\in\mu_i^{-1}\mathfrak c_i^{-1}$ gives
$|\mu_i|_\ell=|\mathfrak c_i|_\ell^{-1}$ in the obvious sense, and $\prod_i\mathfrak c_i=(1)$ gives
$\prod_i|\mu_i|_\ell=1$.
\end{proof}

(C3)(ii) is Theorem~\ref{pf:recip}; that $\mathsf u$ is a unit at $2$ is Corollary~\ref{pf:CS}.

\subsubsection{Theta products, the integral vacuum, and conjugate jets: Lemma~\ref{ps:conj}}
\label{pk:ss-conj}

Let $W$ be a standard space $E^{\mathsf C}$ with the form $\sum\aubar x_ix_i$ ($\mathsf C$ a finite
set; for $W=\auHom(X,X')$, $\mathsf C=[n']\times[n]$), $L=\alpha\widehat\auO_E^{\,\mathsf C}$,
$\Lambda_L=W(E)\cap L=\mathfrak a_0^{\mathsf C}$. A point $\gamma\in W(E)$ has the coordinate
$\gamma^\natural=(c_\psi\,\tau(\gamma_i))_{\tau,i}$. Note that the argument $z$ of $\Theta_f$ in
\eqref{pk:eq-theta} is a point of $W\otimes\R$ in these coordinates (it enters through
$\psrho_\infty(z)$), while the variable of the Fock space is the conjugate coordinate.

\begin{lemma}[Dictionary]\label{pk:dict}
The map $f\mapsto\Theta_f$ is a bijection from $\psHS_f(W)$ onto the space of entire functions
$\Theta$ on $\C^{\auPl\times\mathsf C}$ for which there are an integer $\mathsf N\ge1$ and a function
$\epsilon$ on $\mathsf N\Lambda_L$, with values in the roots of unity, with
\[
\Theta(z+\lambda^\natural)=\epsilon(\lambda)\,e^{\langle z,\lambda^\natural\rangle+|\lambda^\natural|^2
/2}\,\Theta(z)\qquad(\lambda\in\mathsf N\Lambda_L),\qquad\langle z,w\rangle=\sum\aubar w z .
\]
For $\gamma\in W(E)$, $\Theta_{\psrho_f(\gamma)f}(z)$ is a root of unity times
$e^{\langle z,\gamma^\natural\rangle-|\gamma^\natural|^2/2}\,\Theta_f(z-\gamma^\natural)$.
\end{lemma}
\begin{proof}
For $\gamma\in W(E)$, $\theta\circ\psrho(-\gamma)=\theta$ and $\psrho(-\gamma)=\psrho_\infty(
-\gamma^\natural)\psrho_f(-\gamma)$. So $\theta(\psrho_\infty(z)(1\otimes\psrho_f(\gamma)f))$ is a
root of unity (a value of $\beta_f$) times $\beta_\infty(-\gamma^\natural,z)\,\theta(\psrho_\infty(z-
\gamma^\natural)(1\otimes f))$, and $\beta_\infty(-x,z)e^{|z|^2/2-|z-x|^2/2}=e^{\langle z,x\rangle-
|x|^2/2}$; with \eqref{pk:eq-theta} this is the last formula. If $f$ is fixed up to characters by
$\psrho_f(l)$, $l\in\mathsf NL$ (every $f$ is, for some $\mathsf N$, by (W4)), the formula with
$\gamma=-\lambda$ gives the functional equation, with $\epsilon(\lambda)$ the inverse of the
eigenvalue of $\psrho_f(\lambda)$ on $f$ (the root of unity $\beta_f(\lambda,-\lambda)$ is $1$,
because $\pktr_{E/E^+}(\jmath\langle\lambda,\lambda\rangle)=0$); this eigenvalue is a root of unity,
the action of $\psrho_f$ on $\mathsf NL$ being smooth and $\beta_f$ having values in the roots of
unity. The map is injective: if $\Theta_f=0$ then, by
the last formula and the density of $W(E)$ in $W(\auA_f)$, $\theta(\psrho(w)(1\otimes f))=0$ for
all $w\in W(\auA)$, so $1\otimes f$ lies in the kernel of $F\mapsto\theta(\psrho(\cdot)F)$, a
proper closed invariant subspace of an irreducible representation. It is surjective: for fixed
$\mathsf N$ the functions with a given $\epsilon$ form a space of dimension
$[\mathsf N^{-1}\Lambda_L^{*}:\mathsf N\Lambda_L]^{1/2}=\mathsf N^{2f\#\mathsf C}$, the hermitian form
$|z|^2$ being unimodular on $\Lambda_L^\natural$ by the self-duality of
$L$~\cite[\S3, Proposition, p.~26]{MumfordAV}; the admissible $\epsilon$ with values in the roots
of unity differ from one another by characters of the finitely generated group
$\mathsf N\Lambda_L$ with values in the roots of unity, that is, by the characters of
$\mathsf N\Lambda_L$ of finite order (without the condition on the values the statement would
fail: the functions $z\mapsto e^{-\langle z,c\rangle}\Theta_f(z+c)$, $c\in\C^{\auPl\times\mathsf C}$,
satisfy the functional equation with $\epsilon$ multiplied by the unitary character
$\lambda\mapsto e^{2i\,\Iml\langle c,\lambda^\natural\rangle}$, which is arbitrary);
these are produced by the translations
$\psrho_f(x)$, $x\in W(\auA_f)$; and the space of $f$ on which $\psrho_f(\mathsf NL)$ acts by a
given character has the same dimension, by (W4).
\end{proof}

\emph{The theta product.} Let $\Lambda=\Omega_0\auO_K$, $\theta$ its theta function, $u_k=
\zeta_{q_1}^k/q_1$ $(0\le k<q_1)$ and $u_k^\dagger(z_i)=\sum_\tau\aubar{\tau(u_k)}z_{\tau,i}$. Since
$\sum_k\tau(\zeta^k)\aubar{\tau'(\zeta^k)}=q_1\delta_{\tau\tau'}$, the family $(u_k)$ with
multiplicities $q_1$ is a resolution of the identity \eqref{pa:eq-res}. Put
\begin{equation}\label{pk:eq-product}
\Pi(z)=\prod_{i\in\mathsf C}\prod_{k<q_1}\theta\bigl(\tfrac{\Omega_0}{c_\psi}\,u_k^\dagger(z_i)\bigr)^{q_1}.
\end{equation}
For $\lambda\in q_1\Lambda_L$ the arguments change by $\Omega_0\,\pktr_{E/K}(\aubar{u_k}\lambda_i)\in
\Omega_0\auO_K$ (as $\mathfrak a_0\subseteq\mathfrak d_{E/K}^{-1}$), and by \eqref{pa:eq-theta},
\eqref{pa:eq-res} and $A=|\Omega_0|^2\sqrt d/2$, $\Pi$ satisfies the functional equation of
Lemma~\ref{pk:dict} with $\mathsf N=q_1$, $\epsilon=\pm1$. So $\Pi=\Theta_{F_0}$ for a unique
$F_0\in\psHS_f(W)$, and $\theta(1\otimes F_0)=\Pi(0)=1$. By Lemma~\ref{pk:dict}, for $\gamma\in
W(E)$,
\begin{equation}\label{pk:eq-transl}
\Theta_{\psrho_f(\gamma)F_0}(z)=\zeta_\gamma\prod_{i,k}\theta_{\mathsf x_{k,i}\Omega_0}\bigl(\tfrac
{\Omega_0}{c_\psi}u_k^\dagger(z_i)\bigr)^{q_1},\qquad\mathsf x_{k,i}=-\pktr_{E/K}(\aubar{u_k}\gamma_i)
\in K,
\end{equation}
with a root of unity $\zeta_\gamma$: by \eqref{pa:eq-res}, $\langle z,x\rangle=\sum_{i,k}q_1\,H(a_{k,i}
(x),a_{k,i}(z))$ for $a_{k,i}(z)=\frac{\Omega_0}{c_\psi}u_k^\dagger(z_i)$, since
$\boldsymbol\pi|\Omega_0|^2/(Ac_\psi^2)=1$, and $a_{k,i}(\gamma^\natural)=-\mathsf x_{k,i}\Omega_0$. The vectors $\psrho_f(x)F_0$, $x\in W(\auA_f)$, span
$\psHS_f(W)$ (irreducibility), and each is a root of unity times $\psrho_f(\gamma)F_0$ for a
rational $\gamma$ close to $x$.

\emph{The integral vacuum.} The character $\epsilon_L$ is trivial on $q_1L$. On the space of the vectors fixed by
$\psrho(q_1L)$ the operator $\sum_{l\in L/q_1L}\epsilon_L(l)^{-1}\psrho(l)$ is a multiple of the
projector onto the line of $\psvac_L$, and this space is spanned by the translates of $F_0$ that
it contains, which are of the form $\psrho(b)F_0$ with $q_1b\in L$ (the function $\epsilon$ of
$F_0$ is $\epsilon_L$ on $q_1\Lambda_L$, Lemma~\ref{pf:semi}). So there are $b\in q_1^{-1}L$ and roots of unity $\zeta_l$ such that
\begin{equation}\label{pk:eq-vac}
\psvac^\circ:=\sum_{l\in L/q_1L}\zeta_l\,\psrho(l)\psrho(b)F_0\ \in\C^\times\psvac_L
\end{equation}
(some such translate has a non-zero projection). For
$\mathsf C$ a point this defines the integral vacuum of rank one and $c_{\mathrm{vac}}=\psvac^\circ
/\psvac$; in general $\Pi$, $F_0$ and $\psvac_L$ are tensor products over $\mathsf C$, and
$\psvac^\circ=\bigotimes_{\mathsf C}\psvac^\circ_1=c_{\mathrm{vac}}^{\#\mathsf C}\psvac$. It is a sum
of $q_1^{2f\#\mathsf C}$ translates of $F_0$ with coefficients roots of unity.

\begin{proposition}[Lemma~\ref{ps:conj}]\label{pk:conj}
Lemma~\ref{ps:conj} holds, with
$t_\sigma=\sigma(\varpi_{\mathrm{vac}}/\mathsf u)\,c_\psi/\Omega_\sigma$, $\varpi_\sigma$ the
permutation $\tau\mapsto\sigma\circ\tau$ of $\auPl$, and $C_0=C_0(q_0,d)$ given at the end of the proof.
\end{proposition}
\begin{proof}
\emph{(i).} Let $F_0^\sigma$ be the vector attached by \eqref{pk:eq-product} to the lattice
$\Lambda^\sigma=\Omega_\sigma\mathfrak c\aubar{\mathfrak c}^{-1}$ and the number $\Omega_\sigma$; the
functional equation holds with $\mathsf N=q_1\auNm\mathfrak c$, because $\auNm\mathfrak c\in
\mathfrak c\aubar{\mathfrak c}^{-1}$ and $\operatorname{area}(\Lambda^\sigma)=|\Omega_\sigma|^2\sqrt
d/2$. At the scale $s_0$ the arguments of $\theta$ in \eqref{pk:eq-transl} are
$(\varpi_{\mathrm{vac}}/\mathsf u)\,u_k^\dagger(z_i)$, and at the scale $t_\sigma$ those of the
product for $\Lambda^\sigma$ are $\sigma(\varpi_{\mathrm{vac}}/\mathsf u)\,u_k^\dagger(z_i)$. Since
$\sigma$ fixes $K$ and commutes with the complex conjugation of the CM field $E$,
$\sigma(\aubar{\tau(u_k)})=\aubar{(\sigma\tau)(u_k)}$; so $\sigma$ maps the linear form
$(\varpi_{\mathrm{vac}}/\mathsf u)u_k^\dagger$ to $\sigma(\varpi_{\mathrm{vac}}/\mathsf u)u_k^\dagger$
with the variables permuted by $\varpi_\sigma$. Let $\mathsf B$ be $\beta_\sigma$ regarded in
$\auA_{E,f}$; it is an idele and $\mathsf B\aubar{\mathsf B}$ is a unit (Theorem~\ref{pf:recip}(a)), so
$\mathsf BL$ is self-dual; for $\gamma'\in W(E)$ close to $\mathsf B\gamma$ the numbers
$-\pktr_{E/K}(\aubar{u_k}\gamma'_i)$ are close to $\beta_\sigma\mathsf x_{k,i}$. By (C1) and
(C3)(ii), for $\gamma\in W(E)$ and $\gamma'$ close enough to $\mathsf B\gamma$, the Taylor series
of $\Theta_{\psrho_f(\gamma')F_0^\sigma}(t_\sigma\,\cdot)$ is $\zeta'$ times the image under
$\sigma$ and $\varpi_\sigma$ of that of $\Theta_{\psrho_f(\gamma)F_0}(s_0\,\cdot)$, with $\zeta'$ a
root of unity. Define $\Lambda_\sigma$ on the spanning family by $\Lambda_\sigma(\sum_ic_i
\psrho_f(\gamma_i)F_0)=\sum_i\sigma(c_i)\zeta_i'^{-1}\psrho_f(\gamma'_i)F_0^\sigma$. It is well
defined and satisfies (i), because $f\mapsto$ (Taylor series of $\Theta_f$) is injective
(Lemma~\ref{pk:dict}) and both sides of (i) are $\sigma$-semilinear in $f$. Moreover
$|t_\sigma|^2=1/\Omega$ by \eqref{pk:eq-s0}, (C4) and the area of $\Lambda^\sigma$.

\emph{Commutation rules.} Comparing Taylor series with Lemma~\ref{pk:dict},
$\Lambda_\sigma\psrho_f(x)=\zeta_x\,\psrho_f(\mathsf Bx)\Lambda_\sigma$ with roots of unity
$\zeta_x$ $(x\in W(\auA_f))$; in particular $\Lambda_\sigma$ is injective and maps the line of
$\psvac_L$ to the line of a character of the self-dual lattice $L'=\mathsf BL$. For
$g\in\auU(W)(\auA_f)$ the linear operator $\mathsf D=\omega(g)^{-1}\Lambda_\sigma\omega(g)
\Lambda_\sigma^{-1}$ on $\Lambda_\sigma\psHS_f$ satisfies $\mathsf D\psrho_f(y)\mathsf D^{-1}\in
\C^\times\psrho_f(y)$ for all $y$ by (W1) ($\mathsf B$ is a scalar and commutes with $g$); by the
irreducibility of the Heisenberg representation, $\mathsf D=\mathsf c_g\psrho_f(y_g)$ for some
$y_g$. Thus $\Lambda_\sigma\omega(g)\psvac^\circ=\mathsf c_g\,\omega(g)\psrho_f(y_g)\Lambda_\sigma
\psvac^\circ$, and $|\mathsf c_g|=1$: by (W4), $\omega(g)\psvac_L=\delta_g\sum_r\zeta_r\psrho^\circ(r)
\psvac_L$, a sum over the $\mathsf I_g=[gL:L\cap gL]$ classes $r$, where $\delta_g$ is a root of
unity times $\mathsf I_g^{-1/2}$ times a value of $\upsilon^{\pm1}$ on finite ideles, all of whose
conjugates have absolute value $1$; so $|\sigma\delta_g|=|\delta_g|=\mathsf I_g^{-1/2}$, the left
side is a sum of $\mathsf I_g$ orthogonal vectors of norm $\mathsf I_g^{-1/2}\|\Lambda_\sigma
\psvac^\circ\|$, and the right side has norm $|\mathsf c_g|\,\|\Lambda_\sigma\psvac^\circ\|$.

\emph{(ii).} We normalise $\theta$ by $\int_{W(E)\backslash W(\auA)}|\theta(\psrho(w)F)|^2dw=\|F
\|^2$ (self-dual measure; this is the normalisation of the vacuum in \S\ref{ps:ss-weil}). Then
$\sum_{\gamma\in W(E)}\langle F',\psrho(\gamma)F\rangle=\theta(F)\aubar{\theta(F')}$ for
$F,F'\in1\otimes\psHS_f$: the left side converges, is invariant under $W(E)$ in $F$ and in $F'$,
hence is a multiple of the right side by (W1), and the multiple is $1$ by the normalisation. With
$F'=F=\psrho(w)(1\otimes\psvac_{L''})$, $L''$ self-dual, this gives
\begin{equation}\label{pk:eq-sup}
\sup_{w\in W(\auA)}\bigl|\theta\bigl(\psrho(w)(1\otimes\psvac_{L''})\bigr)\bigr|^2\le\sum_{\gamma\in
W(E)\cap L''}e^{-|\gamma^\natural|^2/2},
\end{equation}
because $|\langle\psrho(\gamma)\psrho(w)\mathbf v'',\psrho(w)\mathbf v''\rangle|=|\langle\psrho(\gamma)
\mathbf v'',\mathbf v''\rangle|=e^{-|\gamma^\natural|^2/2}\mathbf 1_{L''}(\gamma)$ for $\mathbf v''=1\otimes
\psvac_{L''}$. For $\gamma\ne0$ in $W(E)\cap L''$, $\nu=\langle\gamma,\gamma\rangle$ is a totally
positive element of $\mathfrak d_{E/K}^{-1}$ (self-duality), so $\sum_\tau\tau(\nu)\ge
f\auNm(\mathfrak d_{E/K})^{-1/(2f)}$ by the inequality of the means ($\nu\in E^+$ is totally
positive and $N_{E^+/\Q}(\nu)\ge\auNm(\mathfrak d_{E/K})^{-1/2}$). Here $E=K\otimes\Q(\zeta_{q_1})$
with $f=\varphi(q_1)$, and $d\nmid q_1$ (every prime dividing $q_1$ divides $(1+d)/4$), so the
discriminants of $K$ and $\Q(\zeta_{q_1})$ are coprime, the discriminant of $E$ is
$d^f\operatorname{disc}(\Q(\zeta_{q_1}))^2$, and $\auNm(\mathfrak d_{E/K})^{1/(2f)}$ is the root
discriminant $q_1\prod_{p\mid q_1}p^{-1/(p-1)}$ of $\Q(\zeta_{q_1})$. Since $q_1$ and $q_0$ have the same
prime divisors ($q_0\mid q_1$, and Lemma~\ref{an:lem-sizes}(S7)), this gives
$f\auNm(\mathfrak d_{E/K})^{-1/(2f)}=\prod_{p\mid q_1}(1-\tfrac1p)p^{1/(p-1)}=\mathsf m:=\prod_{p\mid q_0}
(1-\tfrac1p)p^{1/(p-1)}$; hence the
points $\gamma^\natural$ are separated by $\delta=c_\psi\sqrt{\mathsf m}$, the number of them with
$k\delta\le|\gamma^\natural|<(k+1)\delta$ is at most $(2k+3)^{\mathsf D}$, $\mathsf D=2f\#\mathsf
C\le2q_1\#\mathsf C$, and the right side of \eqref{pk:eq-sup} is at most $\sum_{k\ge0}(2k+3)^{
\mathsf D}e^{-\delta^2k^2/2}\le(1+4/\delta)\exp(2\mathsf D+\mathsf D^2/\delta^2)$.

Let $\mathsf M_d=\max_{\Lambda'}\sup_u|\theta^{\Lambda'}_u(0)|$, the maximum over the finitely many
normalised lattices $\Lambda'$ with multiplication by $\auO_K$; it is finite because
$|\theta_u(0)|$ is continuous and periodic in $u$. By \eqref{pk:eq-vac}, (i) and
\eqref{pk:eq-transl} for $\Lambda^\sigma$, $\Lambda_\sigma\psvac^\circ$ is a sum of $q_1^{2f\#
\mathsf C}$ translates of $F_0^\sigma$ with coefficients roots of unity, and every value at $0$ of
the theta function of a translate of $F_0^\sigma$ is a product of $q_1^2\#\mathsf C$ numbers of
absolute value $\le\mathsf M_d$. By Lemma~\ref{pk:dict} and continuity,
$\sup_w|\theta(\psrho(w)(1\otimes\Lambda_\sigma\psvac^\circ))|\le(q_1^{2f}\mathsf M_d^{q_1^2})^{\#
\mathsf C}$, and by the normalisation of $\theta$ the norm $\|\Lambda_\sigma\psvac^\circ\|$ is at
most this number. Now let $f=\psrho^\circ(x)\omega(g)\psvac^\circ$. By the commutation rules,
$\Lambda_\sigma f$ is a number of absolute value $\|\Lambda_\sigma\psvac^\circ\|$ times
$\psrho_f(x')\omega(g)$ applied to a unit vector on the line of a character of $L'$, that is,
times a translate of a unit vector on the line of $\psvac_{gL'}$; so by \eqref{pk:eq-sup}, for
all $y$,
\[
\bigl|\theta\bigl(\psrho_\infty(y)(1\otimes\Lambda_\sigma f)\bigr)\bigr|\le(q_1^{2f}\mathsf
M_d^{q_1^2})^{\#\mathsf C}\,(1+4/\delta)^{1/2}\exp(\mathsf D+\mathsf D^2/2\delta^2)\le\exp\bigl(C_0
(q_1\#\mathsf C+1)^2\bigr),
\]
with $C_0=4+\log^+\mathsf M_d+\sqrt d/(\boldsymbol\pi\mathsf m)+\log(1+4/\delta)$, which depends
on $q_0$ and $d$ only.
\end{proof}

\subsubsection{The kernels: Lemma~\ref{pa:kernel}}\label{pk:ss-kernel}

\begin{proposition}[Lemma~\ref{pa:kernel}(K0)]\label{pk:K0}
With the lifts defined by \eqref{pk:eq-lift}, {\rm(K0)} holds.
\end{proposition}
\begin{proof}
For a pure tensor $\Phi=\bigotimes_j\Phi_j$ and $\phi\in\auS$, \eqref{pk:eq-lift} gives
$(C_{\vartheta,\Phi}\phi)(g_m)=\int\cdots\int\mathsf F_\Phi(g_0,\dots,g_{m-1})\,dg_0\cdots dg_{m-1}$
over $\prod_j\Gamma_j\backslash\auU(X_j)(\auA_f)$, with $\mathsf F_\Phi=\theta_{m-1}(\Phi_{m-1};
g_{m-1},g_m)\cdots\theta_0(\Phi_0;g_0,g_1)\phi(g_0)$; both sides are linear in $\Phi$. By
\eqref{ps:eq-kernel}, $\theta_j(\Phi_j;g_jb_j,g_{j+1}b_{j+1})=\theta_j(\omega_{\mu_{j+1}}(b_j,
b_{j+1})\Phi_j;g_j,g_{j+1})$. Hence, for the tensor $\Phi=\sum_x\Psi'(x)\phi^\circ_x$, which is
fixed by $\bigotimes_j\omega_{\mu_{j+1}}(b_j,b_{j+1})$ for $b_j\in J_{X_j}$, $b_m=1$ ((C4) of
\S\ref{ps:ss-glue}; the tensor is the same for $\Psi$ and $\Psi'$), and for $\phi$ of level
$J_{X_0}$, $\mathsf F_\Phi$ is invariant under $g_j\mapsto g_jb_j$ in all variables independently.
The levels are free and have volume $1$, so the iterated integral is the sum over the
representatives of the classes, which is (K0).
\end{proof}

A \emph{standard datum at $\ell$} on $W=\auHom(X_j,X_{j+1})$ is a combination, with coefficients
in $\pao\cap\C$, of the vectors $\psrho^\circ(x)\psvac^\circ$ with $\aubar{\pi_M}x$ integral at
the places above $\ell$.

\begin{proposition}\label{pk:std}
Let $\ell\ne N$ be a prime and $\C\to\C_\ell$.
\begin{enumerate}
\item[(S1)] If $g\in\auU(W)(\auA_f)$ is integral at the places above $\ell$, then $\omega(g)$ maps
standard data at $\ell$ to standard data at $\ell$.
\item[(S2)] If $\ell\ne d$ and $f$ is a standard datum at $\ell$, then
$\Theta_f(\pi_Ms_0\,\cdot)$ is $\partial$-integral.
\item[(S3)] If $\ell=p$ does not split in $K$ (so $p\nmid2q_1M$) and $f$ is a standard datum at
$p$, then $\Theta_f(\pi_Ms_0\,z)=\prod_{i,k}\theta(\mu_{k,i}(z))^{q_1}\cdot G(z)$, where
$\mu_{k,i}=\mathsf c\,u_k^\dagger(z_i)$ with $\mathsf c=\pi_M\varpi_{\mathrm{vac}}/\mathsf u\in
\pao$, and $G$ is an integral power series in the series $y(\mu_{k,i}(z))$; so it is one of the
series of Proposition~\ref{pa:product}, for the resolution of the identity $(u_ke_i)_{k,i}$ of
$E^{\mathsf C}$ with multiplicities $q_1$.
\end{enumerate}
\end{proposition}
\begin{proof}
(S1). By (W1), $\omega(g)\psrho^\circ(x)\psvac^\circ$ is a root of unity times $\psrho^\circ(gx)
\omega(g)\psvac^\circ$, and $\aubar{\pi_M}gx$ is integral above $\ell$. By (W4),
$\omega(g)\psvac_L=\delta_g\sum_{r\in gL/(L\cap gL)}\zeta_r\psrho^\circ(r)\psvac_L$ with roots of
unity $\zeta_r$, where $\delta_g$ is a root of unity times $[gL:L\cap gL]^{-1/2}$ times the value
of $\upsilon^{\pm1}$ at an idele $\det\mathsf u$ with $\mathsf uL=L\cap g^{-1}L$. This idele is a
unit above $\ell$; so the index is prime to $\ell$, the value of the algebraic Hecke character
$\upsilon$ is a unit at $\ell$, and the $r$ can be taken integral above $\ell$.

(S2), (S3). By \eqref{pk:eq-vac} and \eqref{pk:eq-transl}, and since a translation by $x$ agrees
up to a root of unity with the translation by a rational $\gamma$ close to $x$, $f$ is a
combination with coefficients in $\pao$ of vectors $\psrho_f(\gamma)F_0$ with $\gamma\in W(E)$,
$2q_1\aubar{\pi_M}\gamma$ integral above $\ell$, and
\[
\Theta_{\psrho_f(\gamma)F_0}(\pi_Ms_0z)=\zeta_\gamma\prod_{i,k}\theta_{\mathsf x_{k,i}\Omega_0}\bigl(
\mathsf c\,u_k^\dagger(z_i)\bigr)^{q_1},\qquad\mathsf c=\tfrac{\Omega_0}{c_\psi}\pi_Ms_0=
\frac{\pi_M\varpi_{\mathrm{vac}}}{\mathsf u}.
\]
The number $\mathsf c$ is integral at $\ell\nmid2q_1$: $\mathsf u$ is a unit (C4), and
$\varpi_{\mathrm{vac}}^{2f}=\upsilon(\alpha)^{-2}\auNm(\mathfrak d_{E/K})^{1/2}$ is a unit at the
primes not dividing $q_1$. The coefficients $\aubar{\tau(u_k)}$ are integral at $\ell\nmid q_1$,
and the points $\mathsf x_{k,i}\Omega_0$ are torsion points of order prime to $\ell$, as
$\ell\nmid2q_1M$. Let $u$ be one of them. If $u\notin t+\Lambda$, (C3)(i) gives
$\theta_u(\mu)=\theta_u(0)\,\theta(\mu)\,\Gamma_u(y(\mu))$ with $\theta_u(0)\in\pao$,
$\Gamma_u\in\pao[[Y]]$. If $u\in t+\Lambda$, then $\theta_u$ is a constant multiple of $\theta_t$
and $(\theta_t/\theta)^2$ is $\Lambda$-periodic with divisor $2[P_0]-2[O]$, hence equal to $\mathsf
c_tx$; as $x=y^2\mathsf e(y)$ with $\mathsf e\in1+y\pao[[y]]$, $\theta_t=\mathsf c_t^{1/2}\,\theta
\cdot y\,\mathsf e(y)^{1/2}$, and $\mathsf c_t^{1/2}=\theta_t'(0)$ is a unit (proof of
Corollary~\ref{pf:K1}(c)). In both cases $\theta_u(\mu)=\theta(\mu)\cdot\Gamma(y(\mu))$ with
$\Gamma\in\pao[[Y]]$. Hence $\Theta_{\psrho_f(\gamma)F_0}(\pi_Ms_0z)=\prod_{i,k}\theta(\mu_{k,i}(z))^{
q_1}\cdot G_\gamma(z)$ with $G_\gamma$ an integral power series in the $y(\mu_{k,i}(z))$; the first
factor does not depend on $\gamma$, so the same holds for $f$, with $G=\sum_\gamma c_\gamma
G_\gamma$. This is (S3). For (S2): $\theta$ and $y$ are $\partial$-integral at $\ell\ne d$ by (C2),
the $\mu_{k,i}$ have coefficients in $\pao$, and compositions and products preserve
$\partial$-integrality. This argument gives (S2) for $\ell\nmid2q_1Md$. For $\ell\mid2q_1M$, (S2) is Theorem~\ref{pf:std} with
$\delta=\aubar{\pi_M}$.
\end{proof}


\begin{proposition}[Lemma~\ref{pa:kernel}(K1), (K2)]\label{pk:K12}
{\rm(K1)} and {\rm(K2)} of Lemma~\ref{pa:kernel} hold.
The integrality of the jets in (K1) holds for every $\ell\ne d$, by (S2) and Theorem~\ref{pf:std}.
\end{proposition}
\begin{proof}
Let $x$ be in the support of $\Psi'$; then $\aubar{\pi_M}x_j$ is integral at the places above
$\ell$ ($\ell\nmid N\auNm Q_\vartheta$), so $\phi^\circ_{j,x_j}$ is a standard datum at $\ell$. Let
$g\in\auU(X_j)(\auA_f)$, $g'\in\auU(X_{j+1})(\auA_f)$. Unitary groups are rational varieties (Cayley transform) and therefore satisfy weak
approximation~\cite[\S7.1, Propositions~7.3 and~7.4]{PlatonovRapinchuk}, so there are $\gamma\in\auU(X_j)(E^+)$, $\gamma'\in\auU(X_{j+1})(E^+)$ with $\mathsf k=\gamma^{-1}g$,
$\mathsf k'=\gamma'^{-1}g'$ so close to $1$ at the places above $\ell$ that they lie in the
stabilisers of the standard lattices there. By the equivariance of the kernels for the rational
points (which follows from \eqref{pk:eq-schuract} and $\theta\circ\omega(\gamma,\gamma')=\theta$),
$\theta_j(\Phi;g,g')=\gamma'\cdot\theta_j(\Phi;\mathsf k,\mathsf k')\cdot\gamma^{-1}$, and
therefore
\[
g'^{-1}_\ell\circ\theta_j(\Phi;g,g')\circ g_\ell=\mathsf k'^{-1}_\ell\circ\theta_j(\Phi;\mathsf k,
\mathsf k')\circ\mathsf k_\ell ,
\]
where $\mathsf k_\ell,\mathsf k'_\ell$ act on $W_\lambda(\pao)$ by matrices invertible over
$\pao$. Thus the two invertible matrices of (K1) are $\rho(\mathsf k'_\ell)^{-1}$ on the left and
$\rho(\mathsf k_\ell)$ on the right; if the components of $g$ and $g'$ above $\ell$ are trivial one
may take $\gamma=\gamma'=1$, so that $\mathsf k=g$, $\mathsf k'=g'$, both matrices are the identity,
and the matrix of $\theta_j(\Phi;g,g')$ in the bases $(e_T)$ is exactly the matrix of jets below. By \eqref{ps:eq-kernel}, $\theta_j(\Phi;\mathsf k,\mathsf k')$ is the matrix of jets along
$S$ of $F=\Theta_f(\pi_Ms_0\,\cdot)$ with $f=\omega_{\mu_{j+1}}(\mathsf k,\mathsf k')
\phi^\circ_{j,x_j}$, which is a unit multiple of a standard datum at $\ell$ by (S1) (the values of
$\mu_{j+1}$ and $\upsilon$ at ideles that are units above $\ell$ are units at $\ell$). This is the
series $F$ of (K1); its jets are integral by (S2) if $\ell\ne d$, and (K2) is (S3): a prime
that does not split in $K$ is odd and does not divide $q_1M$.
\end{proof}

\subsubsection{The prime $N$: a second proof of Lemma~\ref{pl:kernelN} under a condition}\label{pk:ss-N}

We use the normalisations (N1)--(N5) of \S\ref{pl:ss-norm}; they agree with those of
\S\ref{pk:ss-conv} ($W_\lambda(\Z)=\Sigma_\lambda^\#$ with the dual basis, the action through
$\tau\circ c$, the lifts \eqref{pk:eq-lift}; Lemma~\ref{pl:sum} is Proposition~\ref{pk:K0}). Lemma~\ref{pl:kernelN} is proved in Corollary~\ref{pf:K1}(b); this subsection gives a second proof under the condition (Gen) of Proposition~\ref{pk:K1N}, and is not used elsewhere. Let
$\ell$ be an odd prime, $\iota\colon\C\to\C_\ell$, and $E_1\subset\C/\Lambda$ the group of the
torsion points that reduce to the origin of $\mathcal E$ (the kernel of reduction). Put
$\mathrm E_1^*(v)=\zeta(v)-s_2v-(\boldsymbol\pi/A)\aubar v$; it is $\Lambda$-periodic and odd, and
$\theta_u'(0)/\theta_u(0)=\mathrm E_1^*(u+t)$ by the definitions of $\theta$ and $\theta_u$.

\begin{lemma}\label{pk:E1}
Let $v\in\C$ be a torsion point, $v\notin E_1$, and suppose that there is $\alpha\in\auO_K$ with
$\alpha v\equiv v$ modulo $\Lambda$, $\iota(\alpha(1-\aubar\alpha))$ a unit and $\auNm\alpha$ prime
to $\ell$. Then $\mathrm E_1^*(v)\in\pao$. If $\ell$ splits in $K$ and $\mathfrak n$ is the prime
of $K$ induced by $\iota$, such $\alpha$ exists whenever the annihilator of $v$ in $\auO_K$ is
prime to $\aubar{\mathfrak n}$.
\end{lemma}
\begin{proof}
As in the proof of \eqref{pk:eq-s2}, $\alpha\mathrm E_1^*(\alpha z)-\sum_{w\in\alpha^{-1}\Lambda/
\Lambda}\mathrm E_1^*(z+w)$ is constant, and it is odd, so it is $0$. For $w\notin\Lambda$,
$\mathrm E_1^*(z+w)-\mathrm E_1^*(z)-\mathrm E_1^*(w)=\tfrac12(\wp'(z)-\wp'(w))/(\wp(z)-\wp(w))=:
\lambda(z,w)$, and $\sum_{w\ne0}\mathrm E_1^*(w)=0$. Hence $\alpha\mathrm E_1^*(\alpha z)-\auNm
\alpha\,\mathrm E_1^*(z)=\sum_{w\ne0}\lambda(z,w)$, and for $z=v$:
$\alpha(1-\aubar\alpha)\mathrm E_1^*(v)=\sum_{0\ne w\in\alpha^{-1}\Lambda/\Lambda}\lambda(v,w)$. In
the coordinates of $\mathcal E$, $\lambda(v,w)=(Y_v-Y_w)/(X_v-X_w)+a_1/2$. The points $v$ and $w$
do not reduce to the origin ($w$ has order prime to $\ell$), so their coordinates are integral;
and $X_v-X_w$ is a unit unless $v\mp w\in E_1$. This is impossible: $E_1$ is stable under
$\auO_K$ and contains no non-zero point of order prime to $\ell$, and $(\alpha-1)(v\mp w)\equiv
\pm w$ has order prime to $\ell$; so $v\mp w\in E_1$ would give $w=0$. For the last assertion:
choose $\alpha\equiv1$ modulo $\mathfrak n$ and modulo the annihilator of $v$, and $\alpha\equiv2$
modulo $\aubar{\mathfrak n}$; then $\iota(\alpha)\equiv1$ and $\iota(\aubar\alpha)\equiv2$ modulo
$\pam$, so $\iota(\alpha(1-\aubar\alpha))$ is a unit, and $\auNm\alpha$ is prime to $\ell$.
\end{proof}

\begin{lemma}\label{pk:formal}
Let $\ell\ne d$ be odd, $u$ a torsion point with $u\notin t+\Lambda+E_1$, and suppose that
$v=u+t$ satisfies the hypothesis of Lemma~\ref{pk:E1}. Then $\theta_u$ is $\partial$-integral.
\end{lemma}
\begin{proof}
As in the proof of Proposition~\ref{pk:curve}(C3)(i), $R_u=\theta_u/\theta$ satisfies
$R_u'/R_u=\mathsf C+\mathsf h$ with $\mathsf h=(y-y_Q)/(x-x_Q)-(\text{its value at }z=0)$,
$(x_Q,y_Q)$ the point of parameter $u-t$, and $\mathsf C=\mathrm E_1^*(u+t)\in\pao$
(Lemma~\ref{pk:E1}). If $Q$ does not reduce to the origin, $x_Q$ is a unit (it reduces to $\tilde
P_0$ only if $u-t\in E_1$) and $\mathsf h$ is $\partial$-integral, $d/dz$ preserving the local
ring of $\pao[x,y]$ at $x-x_Q$. If $Q$ reduces to the origin, with parameter $\mathsf t_Q=-x_Q/y_Q
\in\pam$, then $x_Q^{-1}\in\mathsf t_Q^2\pao$, $y_Q/x_Q\in\mathsf t_Q^{-1}\pao^\times$ and
$(y-y_Q)/(x-x_Q)=(y_Q/x_Q-y/x_Q)\sum_{k\ge0}(x/x_Q)^k$, whose terms other than the constant
$y_Q/x_Q$ have coefficients in $\pao$; so $\mathsf h$ is $\partial$-integral again. Hence
$R_u=\theta_u(0)\exp(\mathsf Cz+\int_0^z\mathsf h)$ is $\partial$-integral, $\theta_u(0)$ being
integral as in (C3)(i), and so is $\theta_u=\theta R_u$ by (C2).
\end{proof}

\begin{proposition}[Lemma~\ref{pl:kernelN}, with a restriction]\label{pk:K1N}
Let $\iota\colon\C\to\C_N$, and let $z$ be as in Lemma~\ref{pl:kernelN}: integral at the primes
induced by $\iota\circ\tau\circ c$ and with $Nz$ integral at the primes $\mathfrak l_\tau$ induced
by $\iota\circ\tau$ $(\tau\in\auPl)$, which are the primes of $E$ above the prime $\mathfrak n$ of
$K$ induced by $\iota$. Assume the following condition {\rm(Gen)}: in the expansion of
$\omega_{\mu_{j+1}}(\mathsf k,\mathsf k')\phi^\circ_{j,z}$ in translates $\psrho_f(\gamma)F_0$
used in the proof of Proposition~\ref{pk:std}, no point $\mathsf x_{k,i}\Omega_0$ lies in
$t+\Lambda+E_1$ outside $t+\Lambda$. Then the conclusion of Lemma~\ref{pl:kernelN} holds.
\end{proposition}
\begin{proof}
The reduction of Proposition~\ref{pk:K12} applies with $\mathsf k,\mathsf k'$ integral above $N$
((G2) of \S\ref{pl:ss-norm}), and (S1) holds for the data that are combinations of translates by
$x$ with $x$ integral above $\aubar{\mathfrak n}$ and $Nx$ integral above $\mathfrak n$ (the
elements integral above $N$ preserve these conditions). As in the proof of (S2), such a datum is
a combination with coefficients in $\pao$ of translates $\psrho_f(\gamma)F_0$ whose theta
functions at the scale $\pi_Ms_0$ are $\zeta_\gamma\prod\theta_{\mathsf x\Omega_0}(\mathsf c\,
u_k^\dagger(z_i))^{q_1}$ with $\mathsf c$ a unit at $N$ ($N\nmid2q_1M$) and $\mathsf x=-\pktr_{E/K}(
\aubar{u_k}\gamma_i)\in K$ with denominator prime to $\aubar{\mathfrak n}$: the trace of an element
of $E$ that is integral above $\aubar{\mathfrak n}$ is integral at $\aubar{\mathfrak n}$. So the
annihilator of $\mathsf x\Omega_0+t$ is prime to $\aubar{\mathfrak n}$ ($N$ is odd). By (Gen) each
factor is either a multiple of $\theta_t$, treated in the proof of (S3), or satisfies the
hypotheses of Lemma~\ref{pk:formal}; so each factor, and the product, is $\partial$-integral.
\end{proof}
Lemma~\ref{pl:kernelN} is proved in full, without the condition (Gen), in Corollary~\ref{pf:K1}(b):
the translates $\theta_u$ with $u\in t+E_1$ are $\partial$-integral as well (Remark~\ref{pf:gen}). So
Lemmas~\ref{pk:E1}, \ref{pk:formal} and Proposition~\ref{pk:K1N} are not used in the sequel.


%% ===== proofs/A-hecke-eldiv.tex =====
% proofs/A-hecke-eldiv.tex
% Lemmas pa:hecke and pa:eldiv of proofs/A-divisibility.tex (to be \input after that file).
\providecommand{\auO}{\mathcal O}
\providecommand{\auPl}{\Phi}
\providecommand{\auPa}{\Phi_a}
\providecommand{\auS}{\mathcal S}
\providecommand{\auDX}{D_X}
\providecommand{\auGL}{\mathrm{GL}}
\providecommand{\auU}{\mathrm U}
\providecommand{\auF}{\mathbb F}
\providecommand{\aucol}{\operatorname{col}}
\providecommand{\aubar}[1]{\overline{#1}}
\providecommand{\auf}[1]{\texttt{#1}}
\providecommand{\pao}{\mathfrak o}
\providecommand{\pam}{\mathfrak m}
% macros of this file (prefix \ph)
\providecommand{\phB}{\mathfrak B}
\providecommand{\phA}{\mathfrak A}
\providecommand{\phk}{\boldsymbol k}
\providecommand{\phdep}{\operatorname{dp}}
\providecommand{\phdef}{\operatorname{def}}
\providecommand{\phim}{\operatorname{im}}
\providecommand{\phcd}{c_{\mathrm{div}}}

\subsection{Elementary divisors of a scaled Hecke operator: Lemmas~\ref{pa:eldiv}
and~\ref{pa:hecke}}\label{ph:sec}

This subsection concerns the two lemmas on which case (iv) of Theorem~\ref{pa:divis} rests. The
Lemma~\ref{pa:eldiv} is a special case of a theorem on a
local ring with a cycle of Frobenius-linked embeddings (Theorem~\ref{ph:thm}), which has nothing to
do with theta lifts. We give the reduction of Lemma~\ref{pa:eldiv} to Theorem~\ref{ph:thm}
(\S\ref{ph:ss-ded}, using the normalisation of \S\ref{pl:ss-norm}) and the proof of Theorem~\ref{ph:thm}
(\S\S\ref{ph:ss-la}--\ref{ph:ss-proof}). That proof uses three statements on forms in
$N$ groups of $k$ variables, Propositions~\ref{ph:root}, \ref{ph:count} and~\ref{ph:av}, which are
stated here and proved in \S\ref{pe:sec}. Lemma~\ref{pa:hecke} is proved in \S\ref{pl:ss-hecke}.

\subsubsection{The setting and the theorem}\label{ph:ss-set}
Let $\pao$ be a valuation domain of characteristic $0$ with maximal ideal $\pam$ and residue field
$\phk$, $p$ a prime that is not a unit of $\pao$, and $w\in\pao$ with $w^2=p$. Let $W$ be a local
ring with residue field $\auF_q$. A \emph{Frobenius cycle} consists of an even integer $N\ge2$, an
integer $e\ge1$ with $P=p^e\ge5$ and $q=P^N$, ring homomorphisms $\iota_s\colon W\to\pao$
$(s\in\Z/N)$ with
\[
\iota_{s+1}(y)\equiv\iota_s(y)^P\pmod{p\pao}\qquad(y\in W),
\]
and a parity on $\Z/N$ (the residue of $s$ modulo $2$). Then $\iota_s(y)\equiv\iota_s(y)^q$, so for
$y\in\pam_W$ the element $p$ divides $\iota_s(y)\,\iota_s(1-y^{q-1})$, whose second factor is a
unit:
\begin{equation}\label{ph:eq-max}
\iota_s(\pam_W)\subseteq p\pao .
\end{equation}
Hence $\bar\iota_s=\iota_s\bmod\pam$ is an embedding $\auF_q\to\phk$ and
$\bar\iota_{s+1}=\bar\iota_s^{\,P}$.

Let $k\ge2$, $n\ge1$, $G=\auGL_k(W)$. For $g\in G$ put $g_s=\iota_s(g)$ if $s$ is even and
$g_s=\iota_s({}^tg^{-1})$ if $s$ is odd. The \emph{variables} are $x_{b,s}$ $(1\le b\le k,\
s\in\Z/N)$; $s$ is the \emph{letter}. $G$ acts on polynomials over $\pao$ by the substitution
$x_{c,s}\mapsto\sum_b(g_s)_{bc}x_{b,s}$. Let $\phB$ be the set of exponents $\beta$ of degree $n$
in each letter, and $\rho(g)\in\auGL_\phB(\pao)$ the matrix of $g$ on coefficient vectors of forms
$\sum_\beta h_\beta x^\beta$. Fix $r$ with $0<2r\le k$. The variable $x_{b,s}$ is \emph{deep} if
$s$ is even and $b>r$, or $s$ is odd and $b\le r$; $\phdep(\beta)$ is the degree of $x^\beta$ in
the deep variables, and $\Delta=\operatorname{diag}(w^{\phdep\beta})$. Let $K_r\subset G$ be the
subgroup of the $g$ with $g_{ij}\in\pam_W$ for $i\le r<j$, and $L\subset K_r$ the block diagonal
subgroup $\auGL_r(W)\times\auGL_{k-r}(W)$. The index of $K_r$ is the number of $r$-dimensional
subspaces of $\auF_q^k$.

\begin{lemma}\label{ph:para}
For $u\in K_r$ the matrix $\Delta\rho(u)\Delta^{-1}$ has entries in $\pao$; for $u\in L$ it equals
$\rho(u)$.
\end{lemma}
\begin{proof}
$\Delta\rho(u)\Delta^{-1}$ is the matrix of the substitution $x_{c,s}\mapsto\sum_bw^{\epsilon(b,s)
-\epsilon(c,s)}(u_s)_{bc}x_{b,s}$, with $\epsilon$ the indicator of the deep variables. If
$x_{c,s}$ is deep and $x_{b,s}$ is not, then $(u_s)_{bc}$ is $\iota_s(u_{bc})$ with $b\le r<c$ ($s$
even) or $\iota_s((u^{-1})_{cb})$ with $c\le r<b$ ($s$ odd), and is divisible by $p=w^2$ by
\eqref{ph:eq-max}; for $u\in L$ it is zero, and so are the entries with $x_{b,s}$ deep and
$x_{c,s}$ not.
\end{proof}

\emph{Edges and the operator.} Let $X$ and $H$ be finite sets; put $\mathcal E=X\times G/K_r$ (the
\emph{edges}; $e=(x,\varkappa)$ is an edge \emph{into} $x$). Let there be given maps
\[
o\colon\mathcal E\to X,\qquad\gamma\colon\mathcal E\to G,\qquad t\colon G/K_r\to G\ \text{with }
t(\varkappa)^{-1}K_r=\varkappa,\qquad u\colon\mathcal E\to\auGL_H(\pao),
\]
such that
\begin{equation}\label{ph:eq-free}
o(e)=o(e')\ \text{and}\ {}^t\gamma(e)\,{}^t\gamma(e')^{-1}\in K_r\ \Longrightarrow\ e=e'.
\end{equation}
The \emph{edge Hecke operator} on $\mathcal F_X=\{F\colon X\to\pao^H\otimes\pao^\phB\}$ is
\begin{equation}\label{ph:eq-A}
(AF)(x)=\sum_{\varkappa\in G/K_r}\Bigl(u(e)\otimes{}^t\bigl(\rho(\gamma(e))\,\Delta^2\,\rho(t(
\varkappa))\bigr)\Bigr)F(o(e)),\qquad e=(x,\varkappa).
\end{equation}
For submodules $M,M'$ of an $\pao$-module and $z\in\pao$ we write $\aucol(M,M')\mid z$ if for every
finite family $\varphi=(\varphi_1,\dots,\varphi_\nu)$ in $M$ there is $\mu\in M_\nu(\pao)$ with
$\det\mu\mid z$ and all $\sum_j\mu_{ij}\varphi_j\in M'$ (this is the notion used in
Theorem~\ref{pa:divis}).

\begin{theorem}[Elementary divisors of an edge Hecke operator]\label{ph:thm}
Assume Propositions~\ref{ph:root}, \ref{ph:count} and~\ref{ph:av}. Let $0<\theta\le1$ and
\begin{equation}\label{ph:eq-large}
n^4\le q,\qquad\bigl(Nk\log q/\theta\bigr)^5P^{3r}\le n .
\end{equation}
Let $A$ be as in \eqref{ph:eq-A}, let $f\colon\mathcal F_X\to M$ be an $\pao$-linear map and $Y$
an endomorphism of $M$ with $Yf=fA$, and $\Lambda=f(\mathcal F_X)$. Then $Y\Lambda\subseteq
\Lambda$, and for every integer $j\ge0$ there is an integer
\[
l\le\bigl(3\phcd+\theta j\bigr)\,|X|\,|H|\,|\phB|,\qquad\phcd=Nn\Bigl(\frac1{P^{2r}}+\frac{r^2}{kP}
\Bigr),\qquad\text{with}\quad\aucol(\Lambda,\,Y\Lambda+p^j\Lambda)\mid p^{\,l}.
\]
\end{theorem}

Thus, up to a share $\theta$, the elementary divisors of $Y$ on $\Lambda$ have mean exponent at
most $3\phcd$, to be compared with the exponent $nN/2$ of the scaling. The proof occupies
\S\S\ref{ph:ss-la}--\ref{ph:ss-proof}. Its plan: with a pairing for which $\rho(g)$ and
$\rho({}^tg)$ are adjoint, $A$ is the composite of a restriction of functions on $X$ to functions
on edges with values in a lattice ``halfway'' along the edge, and of the transpose of the same
restriction for the reversed edges (\S\ref{ph:ss-edges}). Each restriction has small colength in
its saturation, because few forms modulo $\pam$ have all their translates divisible by $w$ halfway
along every edge (\S\ref{ph:ss-depth}; a count of zeros with multiplicity). Between the two
saturations the composite is nearly onto, up to few generators, because the values of the saturated
functions satisfy congruences that cut the dimension by the number of edges at a vertex, and an
approximate inverse with values satisfying the same congruences exists (\S\ref{ph:ss-inputs}).

\subsubsection{Linear algebra over a valuation domain}\label{ph:ss-la}
Every matrix over $\pao$ is $U\,D\,V$ with $U,V$ invertible and $D$ diagonal (choose an entry
dividing all others, clear its row and column, and repeat). The following rules are immediate
from the definition: (C1) if $\aucol(M,M')\mid z$ and $\aucol(M',M'')\mid z'$ then
$\aucol(M,M'')\mid zz'$; (C2) if $\aucol(M,M')\mid z$ and $\phi$ is linear with
$M_1\subseteq\phi(M)+M_1'$ and $\phi(M')\subseteq M_1'$, then $\aucol(M_1,M_1')\mid z$.

\begin{lemma}\label{ph:col}
Let $\nu\ge1$. {\rm(a)} For $B\in M_\nu(\pao)$, $\aucol(\pao^\nu,B\pao^\nu)\mid\det B$.
{\rm(b)} If $M'\subseteq\pao^\nu$ is a submodule, $s$ a finite subset and $z\in\pao$, then
$\aucol(M'+\langle s\rangle,\,M'+z\pao^\nu)\mid z^{|s|}$.
{\rm(c)} Let $Z,d,d'\in M_\nu(\pao)$, $a,z\in\pao$ and $s\subset\pao^\nu$ finite with
$a\pao^\nu\subseteq Z\pao^\nu+\langle s\rangle$. Then
$\aucol\bigl(\pao^\nu,\ d'Zd\,\pao^\nu+z\pao^\nu\bigr)\mid\det d'\cdot a^\nu\cdot\det d\cdot
z^{|s|}$.
\end{lemma}
\begin{proof}
(a) Write $B=UDV$, $D=\operatorname{diag}(t_i)$. A family $\varphi$ has coordinates
$c_j=U^{-1}\varphi_j$, and a combination $\sum_j\mu_{ij}\varphi_j$ lies in $B\pao^\nu=UD\pao^\nu$ if
and only if $t_h$ divides $\sum_j\mu_{ij}c_{j,h}$ for every $h$. By Lemma~\ref{pa:elim}, applied to
these $\nu$ linear forms in the row of $\mu$, there is $\mu$ with $\det\mu\mid\prod t_h$. (b) Write
$\varphi_j=m_j+\sum_{\sigma\in s}c_{j\sigma}\sigma$ and apply Lemma~\ref{pa:elim} to the $|s|$
forms $\sum_j\mu_{ij}c_{j\sigma}$ with moduli $z$. (c) By (a) for $a\cdot1$ and the hypothesis,
$\aucol(\pao^\nu,Z\pao^\nu+\langle s\rangle)\mid a^\nu$. By (a) for $d$ and (C2) with $\phi=Z$,
$\aucol(Z\pao^\nu+\langle s\rangle,Zd\pao^\nu+\langle s\rangle)\mid\det d$. By (b),
$\aucol(Zd\pao^\nu+\langle s\rangle,Zd\pao^\nu+z\pao^\nu)\mid z^{|s|}$. Chain these by (C1), apply
(C2) with $\phi=d'$, and chain with (a) for $d'$.
\end{proof}

\begin{lemma}[Near inverse]\label{ph:near}
Let $\mathcal X,I,\mathcal J$ be finite sets, $Q\colon\pao^{\mathcal J}\to\pao^{\mathcal X\times
I}$ with a left inverse $R$, $\mathcal A\colon\pao^{\mathcal X\times I}\to\pao^{\mathcal J}$ and
$\mathcal B\colon\pao^{\mathcal J}\to\pao^{\mathcal X\times I}$ linear, and $a,u\in\pao$,
$\varpi\in\pam$ with $a\varpi\mid u$ and $\mathcal A\mathcal B\equiv a\pmod u$. Let
$C\subseteq\pao^I$ be a submodule such that all $\psi$ in the images of $Q$ and of $\mathcal B$
have $\psi(x,\cdot)\in C$ for all $x\in\mathcal X$, and $S\subseteq\phk^I$ a subspace containing
$h\bmod\pam$ whenever $uh\in C$. Then there is a set $s\subset\pao^{\mathcal J}$ with
\[
|\mathcal J|+|s|\le|\mathcal X|\dim S,\qquad a\,\pao^{\mathcal J}\subseteq\mathcal AQ\,
\pao^{\mathcal J}+\langle s\rangle .
\]
\end{lemma}
\begin{proof}
Put $\pi=QR$, a projection onto $\phim Q$. The square matrix $(1-\pi)\mathcal BR$ has the image
$(1-\pi)\phim\mathcal B$; write it as $U\operatorname{diag}(t_j)V$. With $U_j$ the columns of $U$,
\[
\phim Q+\phim\mathcal B=\phim Q+\langle t_jU_j\rangle_j,\qquad RU_j=0\ \text{if }t_j\ne0 .
\]
Let $J=\{j:t_j\ne0,\ t_j\mid u\}$; for $j\notin J$, $u\mid t_j$. For $j\in J$, $uU_j=(u/t_j)t_jU_j$
lies in $\phim Q+\phim\mathcal B$, so its blocks lie in $C$ and those of $U_j$ reduce into $S$; the
same holds for the columns $Q_i$ of $Q$. These $|\mathcal J|+|J|$ vectors are independent modulo
$\pam$: applying $R$ to a relation kills the coefficients of the $Q_i$, and $U$ is invertible. They
lie in $S^{\mathcal X}$ modulo $\pam$, which gives the count for $s=\{\mathcal AU_j:j\in J\}$. Now
$\mathcal A(\phim Q+\phim\mathcal B)\subseteq\mathcal G+u\pao^{\mathcal J}$ with $\mathcal
G=\phim\mathcal AQ+\langle s\rangle$, and $a\eta=\mathcal A\mathcal B\eta-u\eta'$; so
$a\pao^{\mathcal J}\subseteq\mathcal G+\varpi\,a\pao^{\mathcal J}$, and Nakayama's lemma for the
finitely generated module $(a\pao^{\mathcal J}+\mathcal G)/\mathcal G$ gives the claim.
\end{proof}

\subsubsection{Edges, halfway lattices, and the factorisation of $A$}\label{ph:ss-edges}
\emph{Restrictions.} Let $\mathsf E$ be a finite set with maps $\sigma\colon\mathsf E\to X$,
$g\colon\mathsf E\to G$ such that $e\mapsto(\sigma(e),g_e^{-1}K_r)$ is a bijection onto
$X\times G/K_r$. The \emph{restriction} is
\[
\operatorname{Res}=\operatorname{Res}_{\sigma,g}\colon\mathcal F_X\to\mathcal F_{\mathsf E}
=\{\psi\colon\mathsf E\to\pao^H\otimes\pao^\phB\},\qquad(\operatorname{Res}\varphi)(e)=
\bigl(1\otimes\Delta\rho(g_e)\bigr)\varphi(\sigma(e)),
\]
and its \emph{saturation} is $\operatorname{Sat}=\{\psi:w^\mu\psi\in\phim\operatorname{Res}\text{
for some }\mu\}$. (The form $\rho(g_e)\varphi$ is $\varphi$ in a basis adapted to the edge; the
vectors $\Delta h$, $h\in\pao^\phB$, are the coefficients, in the normalised monomials
$w^{-\phdep\beta}x^\beta$, of the integral forms.) Let
\[
\phA=\bigl\{v\in\pao^\phB:\ \exists\,\mu\ge0,\ h\in\pao^\phB\ \text{with}\ w^\mu v=\Delta h\
\text{and}\ w^\mu\mid\Delta\rho(g)h\ \text{for all }g\in G\bigr\},
\]
the forms lying halfway along \emph{all} edges from a vertex.

\begin{lemma}\label{ph:sat}
{\rm(a)} If $z\in\pao$ divides $\operatorname{Res}\varphi$, then $z$ divides $(1\otimes\Delta
\rho(g))\varphi(y)$ for all $y\in X$, $g\in G$. {\rm(b)} For $\psi\in\operatorname{Sat}$ and
$e\in\mathsf E$, $\psi(e)\in\pao^H\otimes\phA$.
\end{lemma}
\begin{proof}
(a) There is $e$ with $\sigma(e)=y$ and $g_e^{-1}K_r=g^{-1}K_r$, so $g=ug_e$ with $u\in K_r$ and
$\Delta\rho(g)\varphi(y)=(\Delta\rho(u)\Delta^{-1})\,(\operatorname{Res}\varphi)(e)$; use
Lemma~\ref{ph:para}. (b) If $w^\mu\psi=\operatorname{Res}\varphi$, take
$h=\rho(g_e)\varphi(\sigma(e))$ (componentwise in $\pao^H$) and apply (a) with $z=w^\mu$ to
$\Delta\rho(g)h=\Delta\rho(gg_e)\varphi(\sigma(e))$.
\end{proof}

\emph{The pairing.} Since $n<q=p^{eN}$, the multinomial coefficient $\binom n{\beta_s}$ of each
$\beta_s$ divides $p^{c_0}$, $c_0=k^2eN$ (it is a product of $k-1$ binomial coefficients with upper
index $<p^{eN}$, each of $p$-adic valuation $\le eN$ by Kummer's theorem). Put
\[
b_\beta=p^{Nc_0}\prod_s\binom n{\beta_s}^{-1},\qquad b'_\beta=\prod_s\binom n{\beta_s},\qquad
b_\beta b'_\beta=p^{Nc_0},\qquad\mathsf b=\operatorname{diag}(b_\beta).
\]
On forms of degree $n$ in $k$ variables the bilinear form with $\langle x^\alpha,x^\beta\rangle=
\delta_{\alpha\beta}\binom n\alpha^{-1}$ satisfies $\langle f(x\,M),f'\rangle=\langle f,f'(x\,{}^tM)
\rangle$ for every matrix $M$ (it suffices to check it on powers of linear forms, where
$\langle(x\cdot\xi)^n,(x\cdot\eta)^n\rangle=(\xi\cdot\eta)^n$). As $({}^tg)_s={}^t(g_s)$ for both
parities,
\begin{equation}\label{ph:eq-adj}
\mathsf b\,\rho({}^tg)={}^t\rho(g)\,\mathsf b\qquad(g\in G).
\end{equation}

\emph{The two sets of edges.} On $\mathsf E=\mathcal E$ consider $(\sigma,g)=(o,{}^t\gamma)$ and
$(\sigma',g')=(\mathrm{pr}_X,\ t\circ\mathrm{pr}_{G/K_r})$. The map $e\mapsto(\sigma'(e),
g_e'^{-1}K_r)$ is the identity. The map $e\mapsto(o(e),{}^t\gamma(e)^{-1}K_r)$ is injective by
\eqref{ph:eq-free}, hence bijective, the set being finite. Let $\operatorname{Res}$,
$\operatorname{Res}'$ be the two restrictions and $\mathsf u=\bigoplus_eu(e)$ on
$\mathcal F_{\mathcal E}$.

\begin{lemma}[Factorisation]\label{ph:fact}
$A\,(1\otimes\mathsf b)={}^t\operatorname{Res}'\circ(\mathsf u\otimes\mathsf b)\circ
\operatorname{Res}$. In particular $A$ has entries in $\pao$.
\end{lemma}
\begin{proof}
At $x$ the right side is $\sum_\varkappa u(e)\otimes{}^t(\Delta\rho(t(\varkappa)))\,\mathsf b\,
\Delta\rho({}^t\gamma(e))$ applied to $F(o(e))$. As $\mathsf b$ and $\Delta$ are diagonal,
\eqref{ph:eq-adj} gives ${}^t\rho(t)\Delta\,\mathsf b\,\Delta\rho({}^t\gamma)={}^t\rho(t)\Delta^2\,
{}^t\rho(\gamma)\,\mathsf b={}^t(\rho(\gamma)\Delta^2\rho(t))\,\mathsf b$. The entries of $A$ are
those of the matrices in \eqref{ph:eq-A}.
\end{proof}

\subsubsection{Few forms are deep for all edges}\label{ph:ss-depth}
Over $\phk$ the group $G$ acts on forms through the $\bar\iota_s$; write $\bar\rho$. For
$\delta\ge1$ let $D_\delta\subseteq\phk^\phB$ be the space of the forms $F$ such that, for every
$g\in G$, all monomials of $\bar\rho(g)F$ have depth at least $\delta$.

\begin{lemma}[Zeros with multiplicity]\label{ph:null}
Let $\mathsf K$ be a field, $\mathsf I$ a finite set of variables with integers $\kappa_v\ge1$; the
\emph{cost} of a monomial $X^\alpha$ is $\sum_v\alpha_v/\kappa_v$. Let $Z\subset\mathsf K^{\mathsf
I}$ be a set of at least $\prod_v\kappa_v$ points such that, for each $v$, there is a polynomial
$g_v$ all of whose monomials have cost $<1$ with $g_v(z)=z_v^{\kappa_v}$ for all $z\in Z$. If
$f\ne0$ vanishes to order at least $\delta$ at every point of $Z$, then $f$ has a monomial of cost
at least $\delta$.
\end{lemma}
\begin{proof}
Put $\psi_v=X_v^{\kappa_v}-g_v$ and $b(\alpha,\gamma)=X^\alpha\psi^\gamma$; call it \emph{reduced}
if $\alpha_v<\kappa_v$ for all $v$. \emph{Claim:} every $b(\alpha,\gamma)$ is a combination of
reduced $b(\alpha',\gamma')$ with $|\gamma'|+\operatorname{cost}(\alpha')\le|\gamma|+
\operatorname{cost}(\alpha)$. Indeed, if $\alpha_v\ge\kappa_v$ and $\alpha=\alpha_0+\kappa_ve_v$,
then $b(\alpha,\gamma)=b(\alpha_0,\gamma+e_v)+g_v\,b(\alpha_0,\gamma)$, and the terms on the right
have $\alpha$ of cost smaller by $1$, respectively smaller, with $|\gamma|$ larger by $1$,
respectively equal; costs lie in a discrete set, so induction on $\operatorname{cost}(\alpha)$
applies. Let $\mathcal J\colon\mathsf K[X]\to\prod_{z\in Z}\mathsf K[X]/\mathfrak m_z^\delta$ be
the jet map; it is onto by the Chinese remainder theorem, its target has dimension
$|Z|\cdot\#\{\gamma:|\gamma|<\delta\}$, and it kills $b(\alpha,\gamma)$ when $|\gamma|\ge\delta$,
because $\psi_v\in\mathfrak m_z$. By the claim the images of the reduced $b(\alpha,\gamma)$ with
$|\gamma|<\delta$ span the target; there are $\prod\kappa_v\cdot\#\{\gamma\}\le\dim$ of them, so
they are independent modulo $\ker\mathcal J$. If all monomials $X^\alpha=b(\alpha,0)$ of $f$ had
cost $<\delta$, then by the claim $f$ would be a combination of these, and $\mathcal J(f)=0$ would
give $f=0$.
\end{proof}

\begin{proposition}[Depth count]\label{ph:depth}
$\displaystyle\sum_{\delta\ge1}\dim_{\phk} D_\delta\le\phcd\,|\phB|$.
\end{proposition}
\begin{proof}
Give a deep variable $v=x_{b,s}$ the number $\kappa_v=P^{2r}$ if $b>2r$ and $\kappa_v=P$ if
$b\le2r$. We show that a non-zero $F\in D_\delta$ has a monomial $\beta$ with
$\operatorname{wt}(\beta):=\sum_{v\text{ deep}}\beta_v/\kappa_v\ge\delta$. Then the projection of
$D_\delta$ to the coordinates $\beta$ with $\operatorname{wt}(\beta)\ge\delta$ is injective, and
\[
\sum_\delta\dim D_\delta\le\sum_\beta\operatorname{wt}(\beta)=|\phB|\,\frac nk\sum_{v\text{ deep}}
\frac1{\kappa_v}\le|\phB|\,\frac nk\cdot\frac N2\Bigl(\frac{k}{P^{2r}}+\frac{2r}P\Bigr)\le\phcd
|\phB|,
\]
because the mean of $\beta_v$ over $\phB$ is $n/k$ by symmetry and each index $b$ is deep in $N/2$
letters.

Let $\mathcal P$ be the set of the positions $(a,a+r)$, $a\le r$, and $(a,c)$, $a\le r<2r<c\le k$;
$|\mathcal P|=r(k-2r+1)$. For $y\in\auF_q^{\mathcal P}$ let $\mathsf A\in M_k(W)$ have entries
lifting $y$ at $\mathcal P$ and $0$ elsewhere; $\mathsf A^2=0$, $g=1+\mathsf A\in G$,
$g^{-1}=1-\mathsf A$. The substitution of $g^{-1}$ fixes the variables that are not deep and sends
a deep variable $X_v$ to $X_v-z_v(y)$, where
\[
z_{(c,s)}(y)=\sum_{a:(a,c)\in\mathcal P}\bar\iota_s(y_{ac})\,x_{a,s}\ \ (s\text{ even},\ c>r),\qquad
z_{(a,s)}(y)=-\sum_{c:(a,c)\in\mathcal P}\bar\iota_s(y_{ac})\,x_{c,s}\ \ (s\text{ odd},\ a\le r)
\]
are linear forms in variables that are not deep. Let $\mathsf K$ be the field of rational
functions over $\phk$ in the variables that are not deep, and $f\in\mathsf K[X_v:v\text{ deep}]$
the polynomial $F$; $f\ne0$. As $F=\bar\rho(g)^{-1}(\bar\rho(g)F)$ and the monomials of
$\bar\rho(g)F$ have depth $\ge\delta$, $f$ vanishes to order $\ge\delta$ at the point $z(y)$. We
apply Lemma~\ref{ph:null} to $Z=\{z(y)\}$; note $\prod_v\kappa_v=P^{\frac N2(2r(k-2r)+2r)}
=q^{|\mathcal P|}$.

For a letter $i$ let $\Lambda_i$ be the $\mathsf K$-span of the functions $y\mapsto z_v(y)^{P^j}$
with $v=(b,s)$ deep, $s+j=i$ and $P^j<\kappa_v$. Since $\bar\iota_s^{\,P^j}=\bar\iota_{s+j}$,
$z_v(y)^{P^j}$ is the linear form in the $\bar\iota_i(y_t)$ whose coefficients are the $P^j$-th
powers of those of $z_v$. \emph{Every function $y\mapsto\bar\iota_i(y_t)$, $t\in\mathcal P$, lies
in $\Lambda_i$.} For $t=(a_0,c)$ with $c>2r$: the $r$ integers $j\in[0,2r)$ with $i-j$ even give
the forms $z_{(c,i-j)}^{P^j}=\sum_{a\le r}\bar\iota_i(y_{ac})\,x_{a,i-j}^{\,P^j}$, and the
determinant of $(x_{a,i-j}^{\,P^j})_{j,a}$ is a signed sum of $r!$ distinct monomials, hence not
zero. For $t=(a,a+r)$: if $i$ is even, $z_{(a+r,i)}=\bar\iota_i(y_{a,a+r})x_{a,i}$; if $i$ is odd,
$z_{(a,i)}=-\bar\iota_i(y_{a,a+r})x_{a+r,i}-\sum_{c>2r}\bar\iota_i(y_{ac})x_{c,i}$ and the last
terms are in $\Lambda_i$. This proves the statement. Consequently $y\mapsto z(y)$ is injective, so
$|Z|=q^{|\mathcal P|}$; and $z_v(y)^{\kappa_v}$, a linear form in the $\bar\iota_{s+j_0}(y_t)$
($\kappa_v=P^{j_0}$), equals $g_v(z(y))$ for a combination $g_v$ of monomials $X_{v'}^{P^j}$ of cost
$P^j/\kappa_{v'}<1$. Lemma~\ref{ph:null} gives a monomial of $f$ of cost $\ge\delta$; it is the
deep part of a monomial $\beta$ of $F$, and its cost is $\operatorname{wt}(\beta)$.
\end{proof}

\begin{proposition}[The ends]\label{ph:ends}
Let $\operatorname{Res}=\operatorname{Res}_{\sigma,g}$ be a restriction and $\nu=|X||H||\phB|$.
There are $Q\colon\pao^\nu\to\mathcal F_{\mathsf E}$ with a left inverse and $d\in M_\nu(\pao)$
with $\operatorname{Res}=Qd$, $\phim Q=\operatorname{Sat}$ and $\det d\mid w^{\,\lambda}$,
$\lambda\le\phcd\,\nu$.
\end{proposition}
\begin{proof}
Write $\operatorname{Res}=U\binom D0V$ with $U,V$ invertible, $D=\operatorname{diag}(t_i)$; let $Q$
consist of the first $\nu$ columns of $U$ and $d=DV$. $\operatorname{Res}$ is injective, so
$t_i\ne0$; if $w^\mu\psi=Qd\varphi$ then $\psi=QR\psi$ for the left inverse $R$, so
$\operatorname{Sat}\subseteq\phim Q$. Let $\varphi_i=V^{-1}e_i$; then $\operatorname{Res}\varphi_i=
t_iQe_i$, and by Lemma~\ref{ph:sat}(a) $t_i$ divides $w^{\phdep\beta}(\rho(g)\varphi_i(y))_\beta$
for all $g,y,\beta$ (and all components in $\pao^H$). With $g=1$: $t_i$ divides $w^{nN}$ times
every entry of $\varphi_i$, one of which is a unit; so $t_i\mid w^{nN}$ and $\phim
Q\subseteq\operatorname{Sat}$. If $t_i\nmid w^\delta$, then $t_i\in w^\delta\pam$ and the entries
of $\rho(g)\varphi_i(y)$ at the $\beta$ of depth $\le\delta$ lie in $\pam$: the reductions of the
$\varphi_i(y)$ lie in $\phk^H\otimes D_{\delta+1}$. The $\varphi_i$ are independent modulo $\pam$,
so $\#\{i:t_i\nmid w^\delta\}\le|X||H|\dim D_{\delta+1}$. Finally $\prod_it_i$ divides
$w^\lambda$ with $\lambda=\sum_{\delta\ge0}\#\{i:t_i\nmid w^\delta\}$ (each $t_i$ divides $w$ to
the number of $\delta$ with $t_i\nmid w^\delta$), and Proposition~\ref{ph:depth} bounds $\lambda$.
\end{proof}

\subsubsection{The three inputs on forms}\label{ph:ss-inputs}
Let $1\le a\ne a'\le k$. In the letter $s$ put $(X_s,Y_s)=(x_{a,s},x_{a',s})$ if $s$ is even and
$(x_{a',s},x_{a,s})$ if $s$ is odd. The \emph{defect} of an exponent $\alpha$ at the letter $i$
along $(a,a')$ is
\[
\phdef_i^{a,a'}(\alpha)=\alpha(Y_i)+\frac1{q-1}\sum_{s\in\Z/N}P^{[s-i]}\bigl(\alpha(Y_s)-\alpha(X_s)
\bigr),\qquad[s-i]\in\{0,\dots,N-1\}.
\]
(The sum is the weight of $x^\alpha$, modulo $q-1$ and up to the factor $P^{-i}$, for the torus of
the root.) Let $\mathrm{SL}_2^{a,a'}(W)\subseteq G$ be the copy of $\mathrm{SL}_2(W)$ in the rows
and columns $a,a'$.

\begin{proposition}[Root lemma]\label{ph:root}
Let $R$ be $\pao$ or $\phk$ (with $w=0$ in $\phk$), let $c$ be real weights of the variables with
$c(Y_s)-c(X_s)\in[\tfrac12,1]$ for all $s$, $f$ a polynomial over $R$, $\alpha$ an exponent and
$\lambda\ge0$ an integer. Suppose that for every $g\in\mathrm{SL}_2^{a,a'}(W)$ and every monomial
$\beta$, the coefficient of $x^\beta$ in $g\cdot f$ is divisible by $w^{\lceil\lambda+c\cdot\alpha
-c\cdot\beta\rceil^+}$, and that $w^{\lambda+1}$ does not divide the coefficient of $x^\alpha$ in
$f$. Then $\phdef_i^{a,a'}(\alpha)\le0$ for every letter $i$.
\end{proposition}
\noindent Proved in \S\ref{pe:ss-root} (Proposition~\ref{pe:root}).

Now let $(a,a')=(r,r+1)$ and write $\phdef_i=\phdef_i^{r,r+1}$; so $Y_s$ is the deep one of the two
variables. For an integer $c\ge1$ put $c^*=3kc$ and
\begin{align*}
C_c&=\bigl\{v\in\pao^\phB:\ w^c\mid(\rho(g)v)_\beta\ \text{for all }g\in L\text{ and all }\beta
\text{ with }\max_i\phdef_i(\beta)>3c^*\bigr\},\\
S_c&=\bigl\{v\in\phk^\phB:\ (\bar\rho(g)v)_\beta=0\ \text{for all }g\in L\text{ and all }\beta
\text{ with }\max_i\phdef_i(\beta)>3c^*\bigr\}.
\end{align*}
Clearly $w^c\pao^\phB\subseteq C_c$, and if $w^{c-1}h\in C_c$ then $h\bmod\pam\in S_c$.

\begin{lemma}\label{ph:AinC}
$\phA\subseteq C_c$ for every $c$; more precisely, for
$v\in\phA$, $g\in L$ and $\phdef_i(\beta)\ge3j$, $w^j$ divides $(\rho(g)v)_\beta$.
\end{lemma}
\noindent Proved in \S\ref{pe:sec} (Lemma~\ref{pe:AinC}).

\begin{proposition}[Count of the forms of bounded defect]\label{ph:count}
Let $c\ge1$, $0<\theta\le1$ and $24\,Nk^2c\,P^{2r}\le\theta n$. Then
$[G:K_r]\cdot\dim_{\phk} S_c\le(1+\theta)|\phB|$.
\end{proposition}
\noindent Proved in \S\ref{pe:ss-count} (Proposition~\ref{pe:count}), from the exponential formula
for permutations~\cite[Cor.~5.1.9]{StanleyEC2}, Euler's product identity~\cite{AndrewsPartitions}, \cite[Ch.~I, \S2, Example~4]{Macdonald}
and the convergence of lattice point counts to the volume for compact convex sets; the hypothesis
is derived there, see \eqref{pe:eq-hyp}.

\begin{proposition}[Averaging]\label{ph:av}
Let $N=2\ell$ and $Nk(\ell+2)n+n+k<P^\ell$, and $c_0=k^2eN$. For every $c\ge1$ there is
$\mathrm{Av}\in M_\phB(\pao)$ such that $\mathrm{Av}\,h\in C_c$ for all $h\in\pao^\phB$, and
\[
{}^th'\,\mathsf b\,\mathrm{Av}\,h\equiv p^{Nc_0}\;{}^th'\,\mathsf b\,h\pmod{w^c}\qquad(h\in
\pao^\phB,\ h'\in\phA).
\]
\end{proposition}
\noindent Proved in \S\ref{pe:ss-av} (Proposition~\ref{pe:av}).

\subsubsection{Proof of Theorem~\ref{ph:thm}}\label{ph:ss-proof}
First $Y\Lambda=fA\mathcal F_X\subseteq\Lambda$. By (C2) with $\phi=f$ it suffices to prove
$\aucol(\mathcal F_X,A\mathcal F_X+p^j\mathcal F_X)\mid p^l$, and since
$A(1\otimes\mathsf b)\mathcal F_X\subseteq A\mathcal F_X$ it suffices to prove it for
$A'=A(1\otimes\mathsf b)$. Let $\nu=|X||H||\phB|$.

\emph{Sizes.} From \eqref{ph:eq-large}: $k\ge2$, $N\ge2$, $\log q\ge1$, so $n\ge(Nk)^5\ge4^5$,
$n<q$, $q\ge n^4$ gives $\log q\ge27$ and $P^{\ell}=\sqrt q\ge n^2>Nk(\ell+2)n+n+k$ for $\ell=N/2$
(as $n\ge(Nk)^5$ and $\ell+2\le N+1$). With $c_0=k^2eN\le\tfrac32k^2\log q$ and $c=5Nc_0$: $Nc_0
\ge2$; $Nc_0\le Nn/P^{2r}\le\phcd$; and $24Nk^2cP^{2r}\le180\,N^2k^4\log q\cdot P^{2r}\le\theta^5(Nk\log q/\theta)^5P^{3r}
\le\theta n$. So
Propositions~\ref{ph:count} and~\ref{ph:av} apply with this $c$, and $n<q$ gives the pairing of
\S\ref{ph:ss-edges}.

\emph{Factorisation.} By Proposition~\ref{ph:ends} for the two restrictions,
$\operatorname{Res}=Qd$, $\operatorname{Res}'=Q'd'$, and by Lemma~\ref{ph:fact}
\[
A'={}^td'\,Z\,d,\qquad Z=\mathcal AQ,\qquad\mathcal A={}^tQ'(\mathsf u\otimes\mathsf b)\colon
\mathcal F_{\mathcal E}\to\pao^\nu .
\]

\emph{The middle.} Let $R'$ be a left inverse of $Q'$ and $\mathcal B=(\mathsf u^{-1}\otimes
\mathrm{Av}\,\mathsf b')\,{}^tR'$, with $\mathsf b'=\operatorname{diag}(b'_\beta)$. Then
$\mathcal A\mathcal B={}^tQ'(1\otimes\mathsf b\,\mathrm{Av}\,\mathsf b')\,{}^tR'$. The columns of
$Q'$ have their values in $\pao^H\otimes\phA$ (Lemma~\ref{ph:sat}(b)), so by
Proposition~\ref{ph:av} and $\mathsf b\mathsf b'=p^{Nc_0}$, ${}^tQ'R'^{\,t}=1$,
\[
\mathcal A\mathcal B\equiv p^{Nc_0}\;{}^tQ'(1\otimes\mathsf b\mathsf b')\,{}^tR'=p^{2Nc_0}
\pmod{w^c}.
\]
We apply Lemma~\ref{ph:near} with $\mathcal X=\mathcal E\times H$, $I=\phB$, $a=p^{2Nc_0}=w^{4Nc_0}$,
$u=w^{c-1}$ (so $aw\mid u$, as $c-1\ge4Nc_0+1$), $C=C_c$ and $S=S_c$: the values of the columns of
$Q$ lie in $\phA\subseteq C_c$ (Lemmas~\ref{ph:sat}, \ref{ph:AinC}), those of $\mathcal B$ lie in
$C_c$ by Proposition~\ref{ph:av} ($C_c$ is a module and $\mathsf u^{-1}$ acts on the factor
$\pao^H$), and the condition on $S$ was noted after the definition. By
Proposition~\ref{ph:count}, $|\mathcal X|\dim S=|X||H|[G:K_r]\dim S_c\le(1+\theta)\nu$. Hence there
is $s$ with $|s|\le\theta\nu$ and $p^{2Nc_0}\pao^\nu\subseteq Z\pao^\nu+\langle s\rangle$.

\emph{Conclusion.} By Lemma~\ref{ph:col}(c) with $z=p^j$,
$\aucol(\pao^\nu,A'\pao^\nu+p^j\pao^\nu)$ divides $\det d'\cdot p^{2Nc_0\nu}\cdot\det d\cdot
p^{j|s|}$, where $\det d\mid w^\lambda$, $\det d'\mid w^{\lambda'}$ with integers $\lambda,\lambda'\le\phcd
\nu$. This divides $p^l$ with $l=\max(\lambda,\lambda')+2Nc_0\nu+j|s|\le(3\phcd+\theta j)\nu$.
\qed

\subsubsection{Deduction of Lemma~\ref{pa:eldiv}}\label{ph:ss-ded}
Let $p$ be exceptional, $\Lambda$ a full mark, $\mathfrak l\in\Lambda$, and write
$\mathfrak r=\mathfrak r_p$, $\pao=\auO_\iota$. We use the normalisation (N1)--(N3) of \S\ref{pl:ss-norm}, which for $\lambda^X$ reads
(Remark~\ref{pl:N}):
\begin{itemize}
\item[(N)] $W_{\lambda^X}=\bigotimes_{\tau\in\auPa}V_\tau$, where $V_\tau$ has the basis
$(e_\beta)$ indexed by the monomials of degree $a$ in $k$ variables, there is no power of the
determinant, and $\gamma\in\auGL_k(E)$ acts on $V_\tau$ by the transpose of the matrix, in the
monomial basis, of the substitution $x\mapsto x\cdot{}^t\tau(\aubar\gamma)$ on forms of degree
$a$. The Hecke operator is $(T_{\mathfrak l}\varphi)(g)=\sum_{h}\varphi(gh)$, $h$ running over
$K^\circ_{\mathfrak l}\,\mathsf d\,K^\circ_{\mathfrak l}/K^\circ_{\mathfrak l}$,
$\mathsf d=\operatorname{diag}(p1_{\mathfrak r},1_{k-\mathfrak r})$, in the identification of
$\auU(X_0)$ at the place below $\mathfrak l$ with $\auGL_k(E_{\mathfrak l})$ by the coordinates at
$\mathfrak l$.
\end{itemize}
For $\gamma$ unitary, ${}^t\tau(\aubar\gamma)=\tau(\gamma)^{-1}$.

\emph{The cycle.} $p$ is inert in $K$, so $p\ge13$ (the primes below $13$ divide $q_0$ and split).
Let $\varphi$ be the Frobenius of $p$ in $\operatorname{Gal}(E/\Q)$, which generates the
decomposition group of every prime above $p$. For $\tau\in\auPl$ and $\nu\in(\Z/q_1)^\times$ the
place $u(\nu)\tau=\tau\circ\sigma_\nu$ induces $\mathfrak l_\tau$ if and only if $\sigma_\nu\in
\langle\varphi\rangle$, and $\aubar{\mathfrak l_\tau}$ if and only if $c\sigma_\nu\in\langle
\varphi\rangle$; as in the proof of Lemma~\ref{pa:mark}, these are the classes $p^{i}$, $i$ even,
and $-p^{i}$, $i$ odd, that is, together the group $\langle-p\rangle$. Let $o$ be the order of
$-p$ modulo $q_1$ and $o_0$ its order modulo $q_0$, odd and $\le\mathfrak r$ by hypothesis. Fix
$\tau_0\in\auPa$ with $\mathfrak l_{\tau_0}=\mathfrak l$. The $a$-places inducing $\mathfrak l$ or
$\aubar{\mathfrak l}$ are the $\tau_s=u((-p)^{o_0s})\tau_0=u(-p^{o_0})^s\tau_0$, $s\in\Z/N$,
$N=o/o_0$; $\tau_s$ induces $\mathfrak l$ exactly for even $s$. Since $-p\not\equiv1\pmod m$ and
$m\mid q_1$ is a power of $2$, $o$ is even, hence so is $N$; and the residue field of $\mathfrak
l$ has $p^{\operatorname{ord}\varphi}=p^{o}$ elements, because $\operatorname{ord}\varphi=2
\operatorname{ord}(p^2)=o$. Let $W$ be the localisation of $\auO_E$ at $\mathfrak l$, and
$\iota_s=\tau_s|_W$ for even $s$, $\iota_s=(\tau_s\circ c)|_W$ for odd $s$; both map $W$ to
$\pao$. By Lemma~\ref{pa:partner} with $e=o_0$, $\tau_{s+1}(y)\equiv\tau_s(\aubar y)^{P}$ modulo
$p$ for $y\in\auO_E$, $P=p^{o_0}$; this is $\iota_{s+1}\equiv\iota_s^P$ for both parities, and it
extends to $W$. So we have a Frobenius cycle with $q=P^N=p^o$, $P\ge13$, and $A_{\mathfrak
l}=aN/2$.

\emph{The edges.} Let $G_f,\Gamma,J=J_X$ be as in \S\ref{ss:au-spaces}; $J$ is free and is the
full group $\auGL_k(\auO_{E_{\mathfrak l}})$ at $\mathfrak l$. Choose representatives $x$ of
$X=\Gamma\backslash G_f/J$ whose components at all primes above $p$ are integral; this is possible by Corollary~\ref{pl:wacor}(a). Let $v_0$ be the place of $E^+$ below
$\mathfrak l$. We may moreover take $x_{v_0}=1$ for every representative: given an integral
representative $x$, the set of $y\in\prod_{v\mid p}\auU(X_0)(E^+_v)$ with $y_{v_0}x_{v_0}\in J_{v_0}$ and
$y_vx_v$ integral for the other $v\mid p$ is open and contains $(x_v^{-1})_{v\mid p}$, so by
Proposition~\ref{pl:wa} it contains some $\gamma\in\Gamma$; then $\gamma x$ represents the same
class, is integral above $p$, and $\gamma x\,j$ with $j\in J$ equal to $(\gamma x_{v_0})^{-1}$ at
$v_0$ and to $1$ elsewhere represents it as well, is integral above $p$ and has component $1$ at
$v_0$. Then $\varphi\in\auS$ is integral at $\iota$ if and only
if all $\varphi(x)$ have coordinates in $\pao$ (Corollary~\ref{pl:wacor}(b)), and $\varphi\mapsto(\varphi(x))_x$ identifies
$\auS_\iota$ with the image of $f\colon\mathcal F_X\to\auS$, where
$\phB\times H$ indexes the basis of $W_{\lambda^X}$: $\phB$ for the factors at
$\tau_s$, $s\in\Z/N$, and $H$ for the other $a$-places. One has
$K^\circ_{\mathfrak l}\mathsf dK^\circ_{\mathfrak l}=\bigsqcup_{\varkappa}\beta_\varkappa\mathsf
dK^\circ_{\mathfrak l}$, $\varkappa=\beta_\varkappa K_{\mathfrak r}\in G/K_{\mathfrak r}$,
$\beta_\varkappa\in G=\auGL_k(W)$, because $\mathsf d^{-1}\beta\mathsf d$ is integral exactly for
$\beta\in K_{\mathfrak r}$. For an edge $e=(x,\varkappa)$ write
\[
x\,\beta_\varkappa\mathsf d=\gamma_e\cdot o(e)\cdot j_e,\qquad\gamma_e\in\Gamma,\ o(e)\in X,\
j_e\in J;\qquad\gamma(e):=\gamma_e^{-1}\beta_\varkappa\mathsf d\in\auGL_k(E),\qquad
t(\varkappa)=\beta_\varkappa^{-1}.
\]
At $v_0$ both $x$ and $o(e)$ have component $1$, so $\beta_\varkappa\mathsf d=\gamma_e\,j_{e,v_0}$
and $\gamma(e)=j_{e,v_0}$ is integral with integral inverse at $\mathfrak l$, so $\gamma(e)\in G$, and $\gamma_e=\beta_\varkappa\mathsf d\gamma(e)^{-1}$. Then
$(T_{\mathfrak l}\varphi)(x)=\sum_\varkappa\gamma_e\cdot\varphi(o(e))$. By (N), $\gamma_e$ acts on
$V_{\tau_s}$ through the substitution by $\tau_s(\gamma_e)^{-1}=\iota_s(\gamma(e)\mathsf d^{-1}
\beta_\varkappa^{-1})$ for even $s$, and by ${}^t(\tau_sc)(\gamma_e)=\iota_s({}^t\gamma(e)^{-1}
\mathsf d\,{}^t\beta_\varkappa)$ for odd $s$. Now $p\,\mathsf d^{-1}$ and $\mathsf d$ are the
diagonal matrices with $p$ at the deep variables of an even, respectively odd, letter (for
$r=\mathfrak r$). Hence $p^{aN/2}$ times the matrix of the substitution of $\gamma_e$ on
$\bigotimes_sV_{\tau_s}$ is $\rho(\gamma(e))\Delta^2\rho(t(\varkappa))$, and $\gamma_e$ acts by
its transpose. At an $a$-place $\tau$ not among the $\tau_s$ the primes $\mathfrak
l_\tau,\aubar{\mathfrak l_\tau}$ differ from $\mathfrak l,\aubar{\mathfrak l}$, and there
$\gamma_e=x\,j_e^{-1}o(e)^{-1}$ is integral with integral inverse; so $\gamma_e$ acts on
$\bigotimes_{\tau\ne\tau_s}V_\tau$ by a matrix $u(e)\in\auGL_H(\pao)$. Therefore
$p^{A_{\mathfrak l}}T_{\mathfrak l}\circ f=f\circ A$ with $A$ as in \eqref{ph:eq-A}.

\emph{Condition \eqref{ph:eq-free}.} Let $o(e)=o(e')=y$ and ${}^t\gamma(e)\,{}^t\gamma(e')^{-1}\in
K_{\mathfrak r}$, that is, $v=\gamma(e)^{-1}\gamma(e')$ has its lower left block in $\pam_W$. Then
$\eta=\gamma_e\gamma_{e'}^{-1}=\beta_\varkappa\mathsf dv\mathsf d^{-1}\beta_{\varkappa'}^{-1}\in
G$. The element $x^{-1}\eta x'=\beta_\varkappa\mathsf dj_e^{-1}j_{e'}\mathsf d^{-1}
\beta_{\varkappa'}^{-1}$ of $G_f$ (its component at $v_0$, where $x_{v_0}=x'_{v_0}=1$) agrees
with $j_e^{-1}j_{e'}$ away from $v_0$ and equals $\eta\in G$ at $v_0$, hence is integral with
integral inverse there; so it lies in $J$, $x=x'$, and
$\eta\in\Gamma\cap xJx^{-1}=1$. Then $\beta_\varkappa^{-1}\beta_{\varkappa'}=\mathsf dv\mathsf
d^{-1}$ is integral with upper right block in $\pam_W$, so $\varkappa=\varkappa'$.

\emph{Sizes.} Take $r=\mathfrak r$, $n=a$, $\theta=1/(m^2f)$. Then $0<2r\le k$ since
$\mathfrak r\le1+\log_2m$; $N\le f$; $o\ge m$, because an element of $(\Z/2^{j'})^\times$ that is
not $\pm1$ modulo $m=2^{j}$ has order at least $2^{j'-j}$, and $2^{j'}\mid q_1$ for $2^{j'}=16\,m^{22}$ by the
definition of the data; so
$q=p^o\ge13^m\ge a^4$ by \eqref{eq:au-scales}. Further $P^{3r}\le p^{3\mathfrak r^2}\le(pm)^{3
\mathfrak r}\le b^{1/2}$, as $p<x_1$ and $(\log x_1)^6\le\log b$; and, using $f\le x^{12}$,
$k\le3x^5$, $m\le2\sqrt x$ and $\log q\le f\log x_1$, $(Nk\log q/\theta)^5\le b^{1/2}$. With
$q\ge a^4$ these give \eqref{ph:eq-large}; the computation is written out in
\S\ref{pe:ss-sizes}, items (a)--(e). Finally $P^{2r}\ge p^{\mathfrak r}\cdot p\ge mp$ and $r^2/(kP)\le1/(mp)$, so
$3\phcd\le6Na/(mp)=12A_{\mathfrak l}/(mp)$. Theorem~\ref{ph:thm} with $j=j_0$,
$|X||H||\phB|=\auDX$, gives Lemma~\ref{pa:eldiv}: $Y_{\mathfrak l}\auS_\iota\subseteq\auS_\iota$
and $\aucol(\auS_\iota,Y_{\mathfrak l}\auS_\iota+p^{j_0}\auS_\iota)\mid p^i$ with
$i\le(12A_{\mathfrak l}/(mp)+j_0/(m^2f))\auDX$. \qed

\begin{remark}\label{ph:status-eldiv}
Thus Lemma~\ref{pa:eldiv}, as stated in \S\ref{pa:ss-divis} and used in the proof of
Theorem~\ref{pa:divis}(iv), follows from Propositions~\ref{ph:root}, \ref{ph:count}, \ref{ph:av}
(proved in \S\ref{pe:sec}); the normalisation (N) and weak approximation are \S\S\ref{pl:ss-norm},
\ref{pl:ss-wa}. Related results on
slopes of Hecke operators on forms of large weight are in \cite{LTXZ}; we know no published
statement that implies Theorem~\ref{ph:thm}.
\end{remark}


%% ===== proofs/A-eldiv-inputs.tex =====
% proofs/A-eldiv-inputs.tex
% Proofs of Propositions ph:root, ph:count, ph:av of proofs/A-hecke-eldiv.tex (to be \input after it).
\providecommand{\auGL}{\mathrm{GL}}
\providecommand{\auF}{\mathbb F}
\providecommand{\aucol}{\operatorname{col}}
\providecommand{\pao}{\mathfrak o}
\providecommand{\pam}{\mathfrak m}
\providecommand{\phB}{\mathfrak B}
\providecommand{\phA}{\mathfrak A}
\providecommand{\phk}{\boldsymbol k}
\providecommand{\phdep}{\operatorname{dp}}
\providecommand{\phdef}{\operatorname{def}}
\providecommand{\phcd}{c_{\mathrm{div}}}
% macros of this file (prefix \pe)
\providecommand{\peX}{\mathbf X}
\providecommand{\peG}{\mathbf G}
\providecommand{\peS}{\mathbf S}
\providecommand{\peht}{\operatorname{ht}}
\providecommand{\perw}{\operatorname{rw}}
\providecommand{\pecw}{\operatorname{cw}}
\providecommand{\pevol}{\operatorname{vol}}
\providecommand{\petr}{\operatorname{tr}}
\providecommand{\pesgn}{\operatorname{sgn}}
\providecommand{\peC}{\mathcal C}
\providecommand{\pecl}[1]{\lceil #1\rceil^{+}}
\providecommand{\pesh}{\sharp}

\subsection{The three inputs on forms: proofs of Propositions~\ref{ph:root}, \ref{ph:count}
and~\ref{ph:av}}\label{pe:sec}

We keep the setting and the notation of \S\ref{ph:ss-set} and \S\ref{ph:ss-inputs}. This subsection
proves Propositions~\ref{ph:root}, \ref{ph:count}, \ref{ph:av} and
Lemma~\ref{ph:AinC} (Lemma~\ref{pe:AinC}), derives the hypothesis of Proposition~\ref{ph:count} (it comes
out as stated, see \eqref{pe:eq-hyp}), and writes out the comparisons of sizes
(\S\ref{pe:ss-sizes}). The results used without proof are: Kummer's theorem on binomial
coefficients, Nakayama's lemma, the exponential formula for permutations, Euler's product identity
\eqref{pe:eq-euler}, and the convergence of lattice point counts to the volume for compact convex
sets.

Throughout, $R$ is $\pao$ or $\phk$, with $w=0$ and $p=0$ in $\phk$; over $\phk$ the group
$G$ acts through the $\bar\iota_s$, and $\bar\iota_{s+1}=\bar\iota_s^{\,P}$. For $x\in\R$,
$\pecl x=\max(\lceil x\rceil,0)$; one has $\pecl{x+y}\le\pecl x+\pecl y$. For $j\in\Z/N$, $[j]\in
\{0,\dots,N-1\}$ is its representative. Recall $P\ge5$, $q=P^N$, $N$ even.

\subsubsection{Levels}\label{pe:ss-lev}
Let $Z$ be a finite set of variables with real weights $c\colon Z\to\R$ (negative weights are
allowed), and $c\cdot\mu=\sum_zc(z)\mu_z$ for an exponent $\mu$. A polynomial $f=\sum_\mu f_\mu
z^\mu\in R[Z]$ has \emph{level at least $b$} (for $c$), $b\in\R$, if $w^{\pecl{b-c\cdot\mu}}$
divides $f_\mu$ for every $\mu$. Over $\phk$ this says that every monomial of $f$ has weight at
least $b$.

\begin{lemma}\label{pe:lev}
{\rm(a)} If $f$, $f'$ have levels at least $b$, $b'$, then $ff'$ has level at least $b+b'$; sums of
polynomials of level at least $b$ have level at least $b$. {\rm(b)} Let $(Z',c')$ be a second set
of weighted variables, $\chi,\chi'\colon R[Z]\to R[Z']$ homomorphisms of $R$-algebras and
$\delta\ge0$, such that for every $z\in Z$ the polynomial $\chi(z)$ has level at least $c(z)$ and
$\chi'(z)-\chi(z)$ has level at least $c(z)+\delta$. If $f$ has level at least $b$, then $\chi(f)$
and $\chi'(f)$ have level at least $b$ and $\chi'(f)-\chi(f)$ has level at least $b+\delta$.
\end{lemma}
\begin{proof}
(a) A coefficient of $ff'$ is a sum of products $f_\mu f'_{\mu'}$, and
$\pecl{b+b'-c\cdot(\mu+\mu')}\le\pecl{b-c\cdot\mu}+\pecl{b'-c\cdot\mu'}$. (b) It suffices to treat
$f=f_\mu z^\mu$; the constant $f_\mu$ has level at least $b-c\cdot\mu$. By (a), $\chi(z^\mu)$ has
level at least $c\cdot\mu$, and so has $\chi'(z^\mu)$ because $\chi'(z)=\chi(z)+(\chi'(z)-\chi(z))$
has level at least $c(z)$. Finally $u_1\cdots u_n-v_1\cdots v_n=\sum_ju_1\cdots u_{j-1}(u_j-v_j)
v_{j+1}\cdots v_n$ shows that $\chi'(z^\mu)-\chi(z^\mu)$ has level at least $c\cdot\mu+\delta$.
\end{proof}

\subsubsection{The root lemma}\label{pe:ss-root}

\begin{proposition}[Root lemma; Proposition~\ref{ph:root}]\label{pe:root}
Let $R$ be $\pao$ or $\phk$ (with $w=0$ in $\phk$), let $c$ be real weights of the variables with
$c(Y_s)-c(X_s)\in[\tfrac12,1]$ for all $s$, $f$ a polynomial over $R$, $\alpha$ an exponent and
$\lambda\ge0$ an integer. Suppose that for every $g\in\mathrm{SL}_2^{a,a'}(W)$ and every monomial
$\beta$, the coefficient of $x^\beta$ in $g\cdot f$ is divisible by $w^{\lceil\lambda+c\cdot\alpha
-c\cdot\beta\rceil^+}$, and that $w^{\lambda+1}$ does not divide the coefficient of $x^\alpha$ in
$f$. Then $\phdef_i^{a,a'}(\alpha)\le0$ for every letter $i$.
\end{proposition}
\begin{proof}
The hypothesis says that every $g\cdot f$ has level at least $b:=\lambda+c\cdot\alpha$ for $c$. Fix
$i$ and put $P_s=P^{[s-i]}$, so $P_i=1$ and $5\le P_s\le q/5$ for $s\ne i$; let $\epsilon_s=-1$ if
$s\equiv i\pmod2$ and $\epsilon_s=1$ otherwise. With
\[
\nu=q\,\alpha(Y_i)+\sum_{s\ne i}P_s\,\alpha(Y_s),\qquad\nu'=\sum_{s}P_s\,\alpha(X_s),\qquad
\text{one has}\quad\phdef_i^{a,a'}(\alpha)=\frac{\nu-\nu'}{q-1},
\]
and we have to show $\nu\le\nu'$.

\emph{The elements $g_r$.} For $r\in W$ let $g_r\in\mathrm{SL}_2^{a,a'}(W)$ have the block
$\bigl(\begin{smallmatrix}0&1\\-1&x\end{smallmatrix}\bigr)$ in the rows and columns $a,a'$, with
$x=-r$ if $i$ is even and $x=r$ if $i$ is odd. The inverse transpose of this block is
$\bigl(\begin{smallmatrix}x&1\\-1&0\end{smallmatrix}\bigr)$. Reading the columns of $(g_r)_s$ for
the two parities of $s$ one finds, for every letter $s$,
\begin{equation}\label{pe:eq-gr}
g_r\cdot X_s=\pm Y_s,\qquad g_r\cdot Y_s=\pm\bigl(X_s+\epsilon_s\,\iota_s(r)\,Y_s\bigr),
\end{equation}
the sign being $-$ for the variables of index $a$ and $+$ for those of index $a'$; the other
variables are fixed.

\emph{The substitution.} Let $t$ be a new variable and $\chi,\chi_0\colon R[x]\to R[t][x]$ the
homomorphisms of $R$-algebras that fix the variables $x_{b,s}$ with $b\ne a,a'$ and are given by
\begin{align*}
\chi_0&\colon& Y_i&\mapsto X_i,& X_s&\mapsto t^{qP_s}X_s\ \ (\text{all }s),& Y_s&\mapsto Y_s\ \
(s\ne i);\\
\chi&\colon& Y_i&\mapsto X_i,\quad X_i\mapsto t^{q}X_i,& X_s&\mapsto t^{qP_s}X_s-\epsilon_st^{P_s}
Y_s,& Y_s&\mapsto Y_s-\epsilon_sX_s\ \ (s\ne i).
\end{align*}
Let $\tau_r$ be the shift $t\mapsto t+\iota_i(r)$ and $e_j(t)=(t+\iota_i(r))^{P^j}-t^{P^j}-
\iota_{i+j}(r)$. Since $(u+v)^{P}\equiv u^P+v^P$ and $\iota_{s+1}\equiv\iota_s^P$ modulo $p$, and
since $u\equiv v$ implies $u^P\equiv v^P$ modulo $p$, we have $e_j\in pR[t]=w^2R[t]$. From
\eqref{pe:eq-gr}, with $d=[s-i]$, so that $\iota_{i+d}=\iota_{i+N+d}=\iota_s$:
\begin{equation}\label{pe:eq-key}
\begin{aligned}
\tau_r\chi(g_r\cdot X_i)&=\pm X_i,\qquad\tau_r\chi(g_r\cdot X_s)=\pm(Y_s-\epsilon_sX_s)\quad(s\ne
i),\\
\tau_r\chi(g_r\cdot Y_i)&=\pm\bigl((t+\iota_i(r))^q-\iota_i(r)\bigr)X_i=\pm\bigl(t^q+e_N\bigr)X_i,\\
\tau_r\chi(g_r\cdot Y_s)&=\pm\bigl((t^{qP_s}+e_{N+d})X_s-\epsilon_s(t^{P_s}+e_d)Y_s\bigr)\quad(s\ne
i).
\end{aligned}
\end{equation}
(For the last line, $\chi(X_s+\epsilon_s\iota_s(r)Y_s)=(t^{qP_s}-\iota_s(r))X_s-\epsilon_s(t^{P_s}-
\iota_s(r))Y_s$, as $\epsilon_s^2=1$.)

\emph{Two systems of weights.} Let $\kappa=(c(Y_i)-c(X_i))/q$; so $\tfrac12\le q\kappa\le1$,
$P_s\kappa\le\tfrac15$, and $qP_s\kappa\ge\tfrac52$ for $s\ne i$. On the variables $t$ and
$x_{b,s}$ define $c^\infty$ and $c^0$ as $c$, except that
\begin{align*}
c^\infty(t)&=-\kappa,& c^\infty(X_s)&=c(X_s)+qP_s\kappa\ \ (\text{all }s);\\
c^0(t)&=\kappa,& c^0(Y_i)&=c(X_i),\qquad c^0(Y_s)=c(Y_s)-P_s\kappa\ \ (s\ne i).
\end{align*}
(i) \emph{For $c^\infty$, each $\chi_0(z)$ has level at least $c(z)$ and $\chi(z)-\chi_0(z)$ has
level at least $c(z)+\tfrac3{10}$.} Indeed $c^\infty(X_i)=c(Y_i)$ and $c^\infty(t^{qP_s}X_s)=
c(X_s)$; the differences are $-\epsilon_st^{P_s}Y_s$ for $z=X_s$, of weight $c(Y_s)-P_s\kappa\ge
c(X_s)+\tfrac12-\tfrac15$, and $-\epsilon_sX_s$ for $z=Y_s$, of weight $c(X_s)+qP_s\kappa\ge
c(X_s)+\tfrac52\ge c(Y_s)+\tfrac32$ $(s\ne i)$.

(ii) \emph{For $c^0$, each $\tau_r\chi(g_r\cdot z)$ has level at least $c(z)$.} We go through
\eqref{pe:eq-key}. $c^0(X_i)=c(X_i)$; $c^0(Y_s)\ge c(X_s)+\tfrac3{10}$ and $c^0(X_s)=c(X_s)$;
$c^0(t^qX_i)=q\kappa+c(X_i)=c(Y_i)$; $c^0(t^{qP_s}X_s)\ge c(X_s)+\tfrac52\ge c(Y_s)$;
$c^0(t^{P_s}Y_s)=c(Y_s)$. The polynomials $e_j$ are divisible by $w^2$ and $c^0(t)\ge0$, so they
have level at least $2$; and $2+c^0(X_s)\ge c(Y_s)$, $2+c^0(Y_s)\ge c(Y_s)$.

\emph{The auxiliary polynomial.} Put $\mathsf p=\chi(f)\in R[t][x]$. As $f$ has level at least
$b$, Lemma~\ref{pe:lev}(b) and (i) give: $\mathsf p$ has level at least $b$ and $\mathsf p-\chi_0(f)$
has level at least $b+\tfrac3{10}$, for $c^\infty$. As $f'=g_r^{-1}\cdot f$ has level at least $b$
and $h\mapsto\tau_r\chi(g_r\cdot h)$ is a homomorphism of $R$-algebras, Lemma~\ref{pe:lev}(b) and
(ii) give: $\tau_r\mathsf p=\tau_r\chi(g_r\cdot f')$ has level at least $b$ for $c^0$. Let
$\bar\alpha$ be the exponent of the monomial obtained from $x^\alpha$ by writing $X_i$ for $Y_i$,
and $\Phi\in R[t]$ the coefficient of $x^{\bar\alpha}$ in $\mathsf p$. Then
\[
c^\infty\cdot\bar\alpha=c\cdot\alpha+q\nu'\kappa,\qquad c^0\cdot\bar\alpha=c\cdot\alpha-\nu\kappa,
\]
because $c^\infty(X_i)=c(Y_i)$, $c^0(X_i)=c^0(Y_i)$. Hence the coefficient of $t^j$ in $\Phi$ is
divisible by $w^{\pecl{\lambda+(j-q\nu')\kappa}}$, so by $w^{\lambda+1}$ if $j>q\nu'$; and the
coefficient of $t^j$ in $\Phi(t+\iota_i(r))$ is divisible by $w^{\pecl{\lambda+(\nu-j)\kappa}}$,
so by $w^{\lambda+1}$ if $j<\nu$. Moreover the coefficients of $t^{q\nu'}x^{\bar\alpha}$ in
$\mathsf p$ and in $\chi_0(f)$ differ by a multiple of $w^{\pecl{\lambda+3/10}}=w^{\lambda+1}$.
Now $\chi_0(x^\beta)=t^{q\sum_sP_s\beta(X_s)}x^{\bar\beta}$, and $\bar\beta=\bar\alpha$ with
$\sum_sP_s\beta(X_s)=\nu'$ forces $\beta=\alpha$ (the two exponents agree outside $X_i,Y_i$, then
$\beta(X_i)=\alpha(X_i)$ as $P_i=1$). So that coefficient of $\chi_0(f)$ is $f_\alpha$.

Let $\bar R=R/w^{\lambda+1}R$ and $\bar\Phi\in\bar R[t]$ the image of $\Phi$. We have shown:
$\deg\bar\Phi=q\nu'$, with leading coefficient the class of $f_\alpha$, which is not zero; and
$(t-a_r)^\nu$ divides $\bar\Phi$ for every $r\in W$, $a_r$ the class of $\iota_i(r)$. Let $r$ run
through representatives of the $q$ residue classes of $W$. The differences $a_r-a_{r'}$ are images
of units, so the ideals $((t-a_r)^\nu)$ are pairwise comaximal and the monic polynomial
$\prod_r(t-a_r)^\nu$, of degree $q\nu$, divides $\bar\Phi\ne0$. Therefore $q\nu\le q\nu'$.
\end{proof}

\begin{lemma}[Lemma~\ref{ph:AinC}]\label{pe:AinC}
$\phA\subseteq C_c$ for every $c$; more precisely, for $v\in\phA$, $g\in L$ and
$\phdef_i(\beta)\ge3j$, $w^j$ divides $(\rho(g)v)_\beta$.
\end{lemma}
\begin{proof}
Let $w^\mu v=\Delta h$ as in the definition of
$\phA$, $F=\sum h_\beta x^\beta$, and $\omega=\phdep-\tfrac13\phdef_i$, a weight of the variables
with $\omega(Y_s)-\omega(X_s)=1-\tfrac13\bigl(\mathbf 1_{s=i}+2P^{[s-i]}/(q-1)\bigr)\in[\tfrac12,1]$
(the bracket lies in $[0,\tfrac32]$ because $5P^{[s-i]}\le q$ for $s\ne i$ and $q\ge5$). By the
definition of $\phA$, every $g\cdot F$, $g\in G$, has level at least $\mu$ for $\phdep$. We claim
that it has level at least $\mu$ for $\omega$. Let $\mathcal T$ be the set of pairs
$(\beta,\lambda)$, $\lambda\ge0$ an integer, with $\lambda+\omega\cdot\beta<\mu$, for which some
$g\in G$ has $w^{\lambda+1}\nmid(g\cdot F)_\beta$. It is finite, since $\beta$ has the degrees of
$F$ and $\lambda<\mu-\omega\cdot\beta$. If the claim fails, $\mathcal T\ne\emptyset$ (take
$\lambda=\pecl{\mu-\omega\cdot\beta}-1$); choose $(\alpha,\lambda)\in\mathcal T$ with
$\lambda+\omega\cdot\alpha$ least, and $g_0$ with $w^{\lambda+1}\nmid(g_0\cdot F)_\alpha$. Every
$g\cdot F$ has level at least $\lambda+\omega\cdot\alpha$ for $\omega$: if $w^e\nmid(g\cdot
F)_\beta$ with $e=\pecl{\lambda+\omega\cdot\alpha-\omega\cdot\beta}$, then $e\ge1$ and
$(\beta,e-1)\in\mathcal T$ with $e-1+\omega\cdot\beta<\lambda+\omega\cdot\alpha$, against the
choice. So Proposition~\ref{pe:root} applies to $f=g_0\cdot F$, $c=\omega$, and gives
$\phdef_i(\alpha)\le0$. Then $\lambda+\phdep(\alpha)\le\lambda+\omega\cdot\alpha<\mu$, both sides
are integers, and $w^{\mu-\phdep\alpha}$, which divides $(g_0\cdot F)_\alpha$, is a multiple of
$w^{\lambda+1}$: a contradiction. Now let $g\in L$ and $\phdef_i(\beta)\ge3j$. By
Lemma~\ref{ph:para}, $w^\mu\rho(g)v=\Delta\rho(g)h$, so $w^\mu(\rho(g)v)_\beta=w^{\phdep\beta}
(g\cdot F)_\beta$, and by the claim $(g\cdot F)_\beta$ is divisible by $w^{\pecl{\mu-\phdep\beta
+j}}$. In both cases, $\mu-\phdep\beta+j\ge0$ and $\mu-\phdep\beta+j<0$, cancelling $w^\mu$ in the
domain $\pao$ shows $w^j\mid(\rho(g)v)_\beta$ (in the second case $(\rho(g)v)_\beta$ is even
divisible by $w^{\phdep\beta-\mu}$, and $\phdep\beta-\mu>j$).
\end{proof}

\subsubsection{The count of the forms of bounded defect}\label{pe:ss-count}
In this paragraph everything is over $\phk$. For $\beta\in\R^{[k]\times\Z/N}$ the defects
$\phdef^{b,b'}_i(\beta)$ are defined by the same formula, and $|\beta_s|=\sum_b\beta_{b,s}$ is the
degree in the letter $s$. As $N$ is even, $(-1)^s$ makes sense for $s\in\Z/N$. For
$\sigma=(\sigma_s)_{s\in\Z/N}\in\mathfrak S_k^N$ let $\mathsf S_\sigma$ be the linear map of
$\R^{[k]\times\Z/N}$
\[
(\mathsf S_\sigma\nu)_{b,s}=(-1)^s\bigl(P\,\nu_{b,s+1}-\nu_{\sigma_s(b),s}\bigr),\qquad\mathsf
S=\mathsf S_1,\qquad\bar\sigma=\sigma_0\sigma_1\cdots\sigma_{N-1}.
\]
Let
\[
\peC=\bigl\{\beta:\phdef_i^{b,b+1}(\beta)\le0\ \text{for all }i\text{ and }1\le b<k\bigr\},\qquad
\Sigma_D=\bigl\{\beta\ge0:\ D_s-1\le|\beta_s|\le D_s\ \text{for all }s\bigr\}
\]
for $D\in\R^{\Z/N}$.

\begin{lemma}\label{pe:coord}
{\rm(a)} For $a\ne a'$, a letter $i$ and $\nu\in\R^{[k]\times\Z/N}$:
$\phdef_i^{a,a'}(\mathsf S\nu)=(\mathsf S\nu)(Y_i)-(\nu_{a,i}-\nu_{a',i})$. Hence $\mathsf S\nu\in
\peC$ if and only if, for all $b<k$: $P\nu_{b+1,i+1}\le\nu_{b,i}$ for even $i$, and
$\nu_{b+1,i}\le P\nu_{b,i+1}$ for odd $i$.
{\rm(b)} $|\det\mathsf S_\sigma|=P^{kN}\Delta(\bar\sigma)$, where $\Delta(\tau)=\prod_{\gamma}(1-
q^{-|\gamma|})$, the product over the cycles $\gamma$ of $\tau\in\mathfrak S_k$.
{\rm(c)} For $\beta\ge0$ with $|\beta_s|\le d$ for all $s$, $|\phdef_i^{a,a'}(\beta)|\le2d$.
{\rm(d)} There is $\varsigma\in\Z_{\ge0}^{[k]\times\Z/N}$ with $|\varsigma_s|\le P^{2r}-1$,
$\phdef_i^{r,r+1}(\varsigma)\le-1$ and $\phdef_i^{b,b+1}(\varsigma)\le0$ for all $i$ and $b\ne r$.
\end{lemma}
\begin{proof}
(a) Let $\beta=\mathsf S\nu$ and $d_s=\nu_{a,s}-\nu_{a',s}$. For both parities of $s$ one finds
$\beta(Y_s)-\beta(X_s)=d_s-Pd_{s+1}$. Hence $\sum_sP^{[s-i]}(\beta(Y_s)-\beta(X_s))=\sum_{j=0}^{N-1}
(P^jd_{i+j}-P^{j+1}d_{i+j+1})=(1-q)d_i$. The second statement follows: $Y_i=x_{b+1,i}$ for even
and $x_{b,i}$ for odd $i$.

(b) Write $\mathsf S_\sigma=\mathsf E\cdot P\,\mathsf{Sh}\cdot(1-\mathsf T)$ with $\mathsf E$ the
diagonal matrix of the signs, $(\mathsf{Sh}\,\nu)_{b,s}=\nu_{b,s+1}$ and $(\mathsf T\nu)_{b,s}=
P^{-1}\nu_{\sigma_{s-1}(b),s-1}$. The map $\mathsf T$ sends the coordinates of the letter $s-1$ to
those of the letter $s$, and $(\mathsf T^N\nu)_{b,0}=q^{-1}\nu_{\bar\sigma(b),0}$. If $\zeta$ is a
primitive $N$-th root of unity and $\mathsf D$ multiplies the letter $s$ by $\zeta^s$, then
$\mathsf D\mathsf T\mathsf D^{-1}=\zeta\mathsf T$; so for a variable $z$,
$\det(1-z\mathsf T)^N=\prod_j\det(1-\zeta^jz\mathsf T)=\det(1-z^N\mathsf T^N)=\det(1-z^N\mathsf
M)^N$, where $\mathsf M=\mathsf T^N$ on the letter $0$ (the restrictions of $\mathsf T^N$ to the
letters are products of the same $N$ maps in cyclically permuted order and have the same
characteristic polynomial). Two polynomials with constant term $1$ and equal $N$-th powers are
equal; at $z=1$, $\det(1-\mathsf T)=\det(1-\mathsf M)=\prod_\gamma(1-q^{-|\gamma|})$, the product
over the cycles of $\bar\sigma$.

(c) $|\phdef_i^{a,a'}(\beta)|\le d+\frac d{q-1}\sum_{j<N}P^j=d+\frac d{P-1}$.

(d) Let $\mu_b=P^{2(r-b)}$ for $b\le r$ and $\mu_b=0$ for $b>r$, so $\mu_b=P^2\mu_{b+1}+\mathbf
1_{b=r}$; let $\nu_{b,s}=\mu_b$ for even and $P\mu_b$ for odd $s$, and $\varsigma=\mathsf S\nu$.
Then $\varsigma_{b,s}=(P^2-1)\mu_b$ for even $s$ and $0$ for odd $s$, with $|\varsigma_s|\le
P^{2r}-1$. By (a), $\phdef_i^{b,b+1}(\varsigma)=P^2\mu_{b+1}-\mu_b=-\mathbf 1_{b=r}$ for even $i$,
and $=-P(\mu_b-\mu_{b+1})\le-\mathbf 1_{b=r}$ for odd $i$.
\end{proof}

\begin{lemma}[Interlacing]\label{pe:inter}
Let $x_1\ge\dots\ge x_k$ and $y_1\ge\dots\ge y_k$ be real. The determinant of the matrix
$(\mathbf 1_{y_c<x_b})_{b,c}$ is $1$ if $x_1>y_1\ge x_2>y_2\ge\dots\ge x_k>y_k$, and $0$ otherwise.
\end{lemma}
\begin{proof}
Induction on $k$. If $y_1<x_b$ for some $b\ge2$, the rows $1$ and $b$ consist of ones, the
determinant is $0$, and $y_1\ge x_2$ fails. Otherwise the first column is $(\mathbf
1_{y_1<x_1},0,\dots,0)$; expand along it.
\end{proof}

\begin{lemma}[The share of the cone]\label{pe:vol}
For every $D$, $\displaystyle\pevol(\peC\cap\Sigma_D)=\prod_{l=1}^{k}\frac{q-1}{q^l-1}\cdot
\pevol(\Sigma_D)$.
\end{lemma}
\begin{proof}
Let $\nu$ lie outside the finitely many hyperplanes $\nu_{b,s}=\nu_{b',s}$ $(b\ne b')$ and
$P\nu_{b,s+1}=\nu_{b',s}$, let $\tau_s\in\mathfrak S_k$ be the permutation with $\nu'_{b,s}:=
\nu_{\tau_s(b),s}$ decreasing in $b$. \emph{Claim:}
\begin{equation}\label{pe:eq-point}
\sum_{\sigma\in\mathfrak S_k^N}\pesgn(\bar\sigma)\,\mathbf 1\bigl[\mathsf S_\sigma\nu>0\bigr]=
\mathbf 1\bigl[\mathsf S\nu'>0\ \text{and}\ \mathsf S\nu'\in\peC\bigr],
\end{equation}
where $>0$ is meant in every coordinate. The left side is $\prod_s\det\mathsf M_s$ for the
matrices $\mathsf M_s(b',b)=\mathbf 1[(-1)^s(P\nu_{b,s+1}-\nu_{b',s})>0]$, because
$\pesgn\bar\sigma=\prod_s\pesgn\sigma_s$ and $\mathbf 1[\mathsf S_\sigma\nu>0]=\prod_s\prod_b
\mathsf M_s(\sigma_s(b),b)$. Permuting the rows by $\tau_s$ and the columns by $\tau_{s+1}$
replaces $\nu$ by $\nu'$ in $\mathsf M_s$ and multiplies the determinant by
$\pesgn\tau_s\pesgn\tau_{s+1}$; these signs cancel in the product over $s\in\Z/N$. For even $s$ the
new matrix is the transpose of the matrix of Lemma~\ref{pe:inter} for $x=P\nu'_{\cdot,s+1}$,
$y=\nu'_{\cdot,s}$, and its determinant is $1$ exactly if $P\nu'_{b,s+1}>\nu'_{b,s}\ge
P\nu'_{b+1,s+1}$ for all $b$; for odd $s$ it is that matrix for $x=\nu'_{\cdot,s}$,
$y=P\nu'_{\cdot,s+1}$, with determinant $1$ exactly if $\nu'_{b,s}>P\nu'_{b,s+1}\ge\nu'_{b+1,s}$.
In both cases the strict inequalities say $(\mathsf S\nu')_{\cdot,s}>0$ and the others are the
conditions of Lemma~\ref{pe:coord}(a) at the letter $s$. This proves \eqref{pe:eq-point}.

Integrate \eqref{pe:eq-point} over the set of the $\nu$ for which the degrees of $\mathsf S\nu$ lie
in $[D_s-1,D_s]$; $\mathsf S_\sigma\nu$ and $\mathsf S\nu'$ have the same degrees as $\mathsf
S\nu$. On the left, the term of $\sigma$ is $\pesgn(\bar\sigma)$ times the volume of $\mathsf
S_\sigma^{-1}(\Sigma_D\cap\{\beta>0\})$, which is $\pevol(\Sigma_D)/(P^{kN}\Delta(\bar\sigma))$ by
Lemma~\ref{pe:coord}(b); each $\tau\in\mathfrak S_k$ is $\bar\sigma$ for $(k!)^{N-1}$ families
$\sigma$. On the right, let $\mathcal U=\mathsf S^{-1}(\peC\cap\Sigma_D\cap\{\beta>0\})$. For
$\nu\in\mathcal U$ each $\nu_{\cdot,s}$ is strictly decreasing, because $\nu_{b,s}-\nu_{b+1,s}\ge
(\mathsf S\nu)(Y_s)>0$ by Lemma~\ref{pe:coord}(a). So the $(k!)^N$ images of $\mathcal U$ under
permutations of the indices in each letter are disjoint, and the right side of
\eqref{pe:eq-point} is the indicator of their union. Hence
\[
(k!)^N\,\frac{\pevol(\peC\cap\Sigma_D)}{P^{kN}\Delta(1)}=(k!)^{N-1}\,\frac{\pevol(\Sigma_D)}{P^{kN}}
\sum_{\tau\in\mathfrak S_k}\frac{\pesgn\tau}{\Delta(\tau)} .
\]
By the exponential formula for permutations~\cite[Cor.~5.1.9]{StanleyEC2}, as an $l$-cycle has sign
$(-1)^{l-1}$,
\begin{equation}\label{pe:eq-euler}
\sum_{k\ge0}\frac{z^k}{k!}\sum_{\tau\in\mathfrak S_k}\frac{\pesgn\tau}{\Delta(\tau)}
=\exp\Bigl(\sum_{l\ge1}\frac{(-1)^{l-1}z^l}{l\,(1-q^{-l})}\Bigr)=\prod_{j\ge0}(1+q^{-j}z)
=\sum_{k\ge0}\frac{q^{-k(k-1)/2}\,z^k}{\prod_{l=1}^k(1-q^{-l})},
\end{equation}
the second equality by $(1-q^{-l})^{-1}=\sum_jq^{-jl}$ and the third being Euler's
identity~\cite{AndrewsPartitions}, \cite[Ch.~I, \S2, Example~4]{Macdonald} with $q$ replaced by $q^{-1}$; see also \cite[(17.5.1)]{pe-DLMF17} with $z\mapsto -z$. With $\Delta(1)=(1-q^{-1})^k$ this gives the
share $(1-q^{-1})^kq^{-k(k-1)/2}\prod_l(1-q^{-l})^{-1}=\prod_l(q-1)/(q^l-1)$.
\end{proof}

Let $I\subseteq G$ be the subgroup of the matrices with $g_{bc}\in\pam_W$ for $b<c$. The
\emph{height} of $x_{b,s}$ is $b$ for even and $-b$ for odd $s$, and $\peht(\beta)$ is the height
of $x^\beta$.

\begin{lemma}[Flags and volume]\label{pe:flag}
Let $c\ge1$, $c^*=3kc$, $\varsigma$ as in Lemma~\ref{pe:coord}(d) and $D_s=n+k+1+4c^*|\varsigma_s|$.
Then $\dim_{\phk}S_c\le[L:L\cap I]\cdot\pevol(\peC\cap\Sigma_D)$.
\end{lemma}
\begin{proof}
Let $m=q^j$, $j\ge1$, and $[m]=\{0,\dots,m-1\}$. \emph{Packing.} For $u\in[m]^{[k]\times\Z/N}$ let
$\gamma_u$ be $u$ with $m(k-\lfloor|u_s|/m\rfloor)$ added at the index $(1,s)$, for each $s$; so
$\gamma_u\equiv u$ modulo $m$ and $mk\le|\gamma_{u,s}|\le mk+m$. To a family $(F_u)_u$ of elements
of $S_c$ attach
\[
\mathbb F=\sum_uF_u(x^m)\,x^{\gamma_u}\in\phk[x] .
\]
The coefficient of $\mathbb F$ at $m\delta+\gamma_u$ is that of $F_u$ at $\delta$, as the class
modulo $m$ of an exponent determines $u$; so $(F_u)\mapsto\mathbb F$ is injective. The elements
of $\bar\iota_s(\auF_q)$ are fixed by the $m$-th power, which is additive; hence
$g\cdot(F(x^m))=(g\cdot F)(x^m)$ and $g\cdot\mathbb F=\sum_u(g\cdot F_u)(x^m)\,(g\cdot
x^{\gamma_u})$ for $g\in G$. The action preserves the degree in each letter. Therefore every
monomial $\alpha$ of $g\cdot\mathbb F$ is $m\delta+\varepsilon$, with $\delta$ a monomial of some
$g\cdot F_u$ and $mk\le|\varepsilon_s|\le mk+m$; in particular $m(n+k)\le|\alpha_s|\le m(n+k+1)$.
If $g\in L$, then $\phdef_i(\delta)\le3c^*$ by the definition of $S_c$, and with
Lemma~\ref{pe:coord}(c)
\begin{equation}\label{pe:eq-a}
\phdef_i^{r,r+1}(\alpha)\le m\,(3c^*+2k+2)\le4mc^*\qquad\text{for all }i,
\end{equation}
since $c^*\ge3k\ge2k+2$.

\emph{Least height.} Let $y\in I$. Over $\phk$ the matrix $y_s$ is lower triangular for even $s$
and upper triangular for odd $s$, so $y\cdot x_{c,s}$ is a non-zero multiple of $x_{c,s}$ plus a
combination of variables $x_{b,s}$ of greater height. Thus $y\cdot x^\beta$ is a non-zero multiple
of $x^\beta$ plus terms of greater height. Let $\mathcal R\subseteq L$ represent the cosets
$(L\cap I)\backslash L$, and let $\mathbb F\ne0$. Among the monomials of all $g\cdot\mathbb F$,
$g\in L$ (finitely many exponents), let $\alpha$ have least height, in $g\cdot\mathbb F$ say, and
write $g=yx$, $y\in L\cap I$, $x\in\mathcal R$. No monomial of $x\cdot\mathbb F$ is lower than
$\alpha$, so the coefficient of $x^\alpha$ in $y\cdot(x\cdot\mathbb F)$ is a non-zero multiple of
that in $x\cdot\mathbb F$: \emph{$\alpha$ is a monomial of $x\cdot\mathbb F$}. Let $b\ne r$, $b<k$.
Then $\mathrm{SL}_2^{b,b+1}(W)\subseteq L$, the height rises by $1$ from $X_s$ to $Y_s$ for the
root $(b,b+1)$, and all monomials of the translates of $f=x\cdot\mathbb F$ by this group have
height at least $\peht(\alpha)$. Proposition~\ref{pe:root} with $R=\phk$, $c=\peht$, $\lambda=0$
gives
\begin{equation}\label{pe:eq-b}
\phdef_i^{b,b+1}(\alpha)\le0\qquad\text{for all }i\text{ and all }b\ne r .
\end{equation}

\emph{The count.} By \eqref{pe:eq-a}, \eqref{pe:eq-b} and Lemma~\ref{pe:coord}(d), the point
$y=\alpha/m+4c^*\varsigma$ lies in $\Lambda_m:=\peC\cap\Sigma_D\cap m^{-1}\Z^{[k]\times\Z/N}$.
Consequently the linear map
\[
S_c^{\,[m]^{[k]\times\Z/N}}\to\phk^{\mathcal R\times\Lambda_m},\qquad(F_u)\mapsto\bigl(\text{the
coefficient of }x\cdot\mathbb F\text{ at the exponent }m(y-4c^*\varsigma)\bigr)_{x,y}
\]
(zero if this is not an exponent) is injective, and $m^{kN}\dim S_c\le[L:L\cap I]\cdot\#\Lambda_m$.
The set $\peC\cap\Sigma_D$ is compact and convex, hence Jordan measurable, and
$\#\Lambda_m/m^{kN}$ tends to its volume as $j\to\infty$.
\end{proof}

\begin{proposition}[Proposition~\ref{ph:count}]\label{pe:count}
Let $c\ge1$, $0<\theta\le1$ and $24\,Nk^2c\,P^{2r}\le\theta n$. Then
$[G:K_r]\cdot\dim_{\phk} S_c\le(1+\theta)|\phB|$.
\end{proposition}
\begin{proof}
Reduction modulo $\pam_W$ maps $G$ onto $\auGL_k(\auF_q)$ and $I$ is the preimage of the lower
triangular group, so $[G:I]=\prod_{l=1}^k(q^l-1)/(q-1)$. As $I\subseteq K_r$ and $L\subseteq K_r$,
$[G:K_r]\,[L:L\cap I]\le[G:K_r]\,[K_r:I]=[G:I]$. With Lemmas~\ref{pe:flag} and~\ref{pe:vol},
$[G:K_r]\dim S_c\le\pevol(\Sigma_D)$. The set $\{y\in\R_{\ge0}^k:d-1\le\sum y_b\le d\}$ has volume
$(d^k-\max(d-1,0)^k)/k!\le d^{k-1}/(k-1)!$, and $|\phB|=\binom{n+k-1}{k-1}^N\ge(n^{k-1}/(k-1)!)^N$.
Put $E=4c^*P^{2r}=12kc\,P^{2r}$; then $D_s\le n+k+1+4c^*(P^{2r}-1)\le n+E$, and
\[
\pevol(\Sigma_D)\le\prod_s\frac{D_s^{k-1}}{(k-1)!}\le|\phB|\Bigl(1+\frac En\Bigr)^{N(k-1)}\le|\phB|
\,e^{x},\qquad x=\frac{N(k-1)E}{n} .
\]
The hypothesis is
\begin{equation}\label{pe:eq-hyp}
2NkE=24\,Nk^2c\,P^{2r}\le\theta n ,
\end{equation}
so $x\le\theta/2\le\tfrac12$ and $e^x\le1+x+x^2\le1+2x\le1+\theta$.
\end{proof}

\subsubsection{The averaging map}\label{pe:ss-av}
\emph{An invariant form on polynomials in a matrix.} Let $\iota$ be one of the $\iota_j$, $\peX=
(X_{bc})$ a generic $k\times k$ matrix, and $\pao[\peX]_d$ the forms of degree $d$. For $a\in L$
put $A=\iota(a)$ and $c_A\psi=\psi(A^{-1}\peX A)$. Fix $\varrho\in\Z^k$; the \emph{row weight} and
the \emph{column weight} of $X^\alpha$ are $\perw(\alpha)=\sum\alpha_{bc}\varrho_b$ and
$\pecw(\alpha)=\sum\alpha_{bc}\varrho_c$. Recall $c_0=k^2eN$ and $q=p^{eN}$.

\begin{lemma}[Average]\label{pe:avg}
Let $t\in\Z$ and $D\ge0$ with $D+k\le q$. There is an $\pao$-linear form $\mathsf I$ on $\pao[\peX]$
such that: {\rm(I1)} $\mathsf I(c_A\psi)=\mathsf I(\psi)$ for all $a\in L$; {\rm(I2)} $\mathsf
I(X^\alpha)=0$ if $\perw(\alpha)>t$; {\rm(I3)} if $\deg\psi\le D$ and $\mathfrak a\subseteq\pao$ is an ideal
containing the coefficient of $X^\alpha$ in $c_A\psi$ for all $a\in L$ and all $\alpha$ with
$\perw(\alpha)>t$, then $\mathsf I(\psi)\equiv p^{c_0}\psi(1)\pmod{\mathfrak a}$. Moreover {\rm(I4)} every
$\pao$-linear form with {\rm(I1)} vanishes on the $X^\alpha$ with $\perw(\alpha)\ne\pecw(\alpha)$.
\end{lemma}
\begin{proof}
(I4) Let $u=1+p\in W^\times$ and $a=\operatorname{diag}(u^{\varrho_b})\in L$. Then
$c_AX^\alpha=(1+p)^{\pecw(\alpha)-\perw(\alpha)}X^\alpha$, and $(1+p)^e\ne1$ in $\pao$ for
$e\ne0$.

\emph{The algebra $\mathcal S_d$.} Fix $d\le D$. Let $U=(\pao^k)^{\otimes d}$ with the basis
$e_j=e_{j_1}\otimes\dots\otimes e_{j_d}$, $j\in[k]^d$, of weight $\sum_\pi\varrho_{j_\pi}$. Let
$\peG\in M_k(\pao[\peX])$ have the entry $X_{cb}$ at $(b,c)$ if $b$ and $c$ lie in the same block
of $L$, and $0$ otherwise; so $\peG(B)={}^t(B^L)$ for a matrix $B$ with block diagonal part $B^L$.
For a linear form $l$ on $\pao[\peX]_d$ let $\Theta(l)=l(\peG^{\otimes d})\in\operatorname{End}
(U)$, $l$ being applied to the entries $\prod_\pi\peG_{i_\pi j_\pi}$, and $\mathcal S_d=\{\Theta(l)\}$. For the evaluation $\mathrm{ev}_B$ at $B\in M_k(\pao)$, $\Theta(\mathrm{ev}_B)=({}^t
B^L)^{\otimes d}$.

\emph{{\rm(S1)} $\mathcal S_d$ is the $\pao$-span of the $\iota(a)^{\otimes d}$, $a\in L$; in
particular it is a ring.} It suffices that the vectors $(B^\alpha)_{|\alpha|=d}$, for $B\in M_k(W)$
with $\det B^L\in W^\times$, span the $W$-module of all vectors indexed by the monomials of degree
$d$: then every $l$ is a combination of evaluations at such $\iota(B)$, and ${}^t(B^L)\in L$. If
they did not, by Nakayama's lemma their reductions would lie in a hyperplane, that is, a non-zero
form $F$ of degree $d$ with coefficients in $\phk$ would vanish at all points $\bar\iota(\bar B)$,
$\bar B\in M_k(\auF_q)$ with $\det\bar B^L\ne0$ (each lifts to $W$); here $\bar\iota$ is the
reduction of $\iota$ and $\bar\iota(\auF_q)\subset\phk$ is a subfield with $q$ elements. Then $F\cdot\det\peX^L$ is a non-zero form of degree $d+k\le q$
vanishing on $\bar\iota(\auF_q)^{k^2}$. This is impossible: reducing a form of degree at most $q$ modulo the
$X_v^q-X_v$ only replaces its monomials $X_v^q$ by $X_v$, which creates no cancellation (the form is homogeneous of degree $d+k\ge2$, so it contains no
monomial of degree $1$, distinct $X_v^q$ go to distinct $X_v$, and the monomials $X_v^q$ are the
only ones with an exponent $\ge q$),
and a non-zero polynomial of degree $<q$ in each variable does not vanish identically on $\bar\iota(\auF_q)^{k^2}$.

\emph{The operator $\psi^\pesh$.} For $|\alpha|=d$ the multinomial coefficient $\binom d\alpha$ is
a product of $k^2-1$ binomial coefficients with upper index at most $d<p^{eN}$, so it divides
$p^{c_0}$ in $\pao$ (Kummer; integers prime to $p$ are units of $\pao$); put $\eta_\alpha=p^{c_0}
\binom d\alpha^{-1}$ and $\langle X^\alpha,X^\beta\rangle=\delta_{\alpha\beta}\eta_\alpha$ on
$\pao[\peX]_d$. As in \S\ref{ph:ss-edges}, $\langle\psi\circ T,\varphi\rangle=\langle\psi,\varphi
\circ{}^tT\rangle$ for every linear map $T$ of the space $M_k$ of the variables, ${}^tT$ being the
transpose in the basis of the elementary matrices. For $\psi\in\pao[\peX]$ with component $\psi_d$
of degree $d$ let $\psi^\pesh=\Theta(\langle\psi_d,\cdot\rangle)\in\mathcal S_d$. Then:

{\rm(S2)} $(c_A\psi)^\pesh=A^{\otimes d}\,\psi^\pesh\,(A^{-1})^{\otimes d}$ for $a\in L$. Indeed the
transpose of $\peX\mapsto A^{-1}\peX A$ is $\peX\mapsto{}^tA^{-1}\peX\,{}^tA$, and $\peG({}^tA^{-1}
\peX\,{}^tA)=A\,\peG\,A^{-1}$ because $A$ is block diagonal.

{\rm(S3)} $\petr(\psi^\pesh)=p^{c_0}\psi_d(1)$: $\petr\peG^{\otimes d}=(\sum_bX_{bb})^d=\sum\binom
d\alpha X^\alpha$ over the monomials in the diagonal variables, and $\psi_d(1)$ is the sum of the
coefficients of $\psi$ at these.

{\rm(S4)} $(X^\alpha)^\pesh e_j=0$ unless $e_j$ has weight $\perw(\alpha)$: the entry $(i,j)$ of
$(X^\alpha)^\pesh$ is $\eta_\alpha$ times the coefficient of $X^\alpha$ in $\prod_\pi\peG_{i_\pi j_\pi}$,
which vanishes unless $X^\alpha=\prod_\pi X_{j_\pi i_\pi}$.

\emph{The form.} Let $U^+=\sum\mathcal S_d\,e_j$ over the $j$ of weight $>t$. By (S1) it is stable
under $\mathcal S_d$ and spanned by the $A^{\otimes d}e_j$ ($a\in L$, $e_j$ of weight $>t$). It is
a finitely generated submodule of $U$, hence free (\S\ref{ph:ss-la}). Put
\[
\mathsf I_d(\psi)=\petr(\psi^\pesh\mid U)-\petr(\psi^\pesh\mid U^+),\qquad\mathsf I=\sum_{d\le
D}\mathsf I_d .
\]
(I1) follows from (S2), as $A^{\otimes d}$ and its inverse lie in $\mathcal S_d$ and preserve
$U^+$. (I2): if $\perw(\alpha)>t$, then $(X^\alpha)^\pesh U\subseteq U^+$ by (S4), and an
endomorphism of $U$ with image in $U^+$ has the same trace on $U$ and on $U^+$. (I3): by (S3),
$\mathsf I(\psi)-p^{c_0}\psi(1)=-\sum_{d\le D}\petr(\psi^\pesh\mid U^+)$, and it suffices to show
$\psi^\pesh U^+\subseteq\mathfrak aU^+$ (then the trace on $U^+$ lies in $\mathfrak a$, and $A^{\otimes d}\mathfrak aU^+\subseteq\mathfrak aU^+$). By (S2) for $a^{-1}$, $\psi^\pesh A^{\otimes d}e_j=A^{\otimes d}
(c_{A^{-1}}\psi)^\pesh e_j$, and for $e_j$ of weight $>t$ the vector $(c_{A^{-1}}\psi)^\pesh e_j$ is
the sum over $\alpha$ of the coefficient of $X^\alpha$ in $c_{A^{-1}}\psi$ times $(X^\alpha)^\pesh
e_j\in U^+$; by (S4) only $\alpha$ with $\perw(\alpha)>t$ contribute, and their coefficients lie in
$\mathfrak a$.
\end{proof}

\emph{The substitution of Frobenius factors.} Let $N=2\ell$. Take generic matrices $\peX_j$,
$j\in\Z/N$ (the \emph{copies}), and $\pao[\peX]=\bigotimes_j\pao[\peX_j]$. For $M\in M_k$ let
$\operatorname{Sym}^n(M)$ be the matrix, on coefficient vectors of forms of degree $n$, of the
substitution $x_c\mapsto\sum_bM_{bc}x_b$; it is multiplicative in $M$, $\rho(g)=\bigotimes_s
\operatorname{Sym}^n(g_s)$, and each term of an entry in the row $\beta$ of
$\operatorname{Sym}^n(M)$ is a product of entries of $M$, $\beta_b$ of them from the row $b$. For a
letter $s$ put $M^\natural=M$ if $s$ is even and $M^\natural={}^t\!\operatorname{adj}(M)$, the
matrix of cofactors, if $s$ is odd. For $h\in L$, $d\ge0$ and $m\ge0$ let
\[
\Psi^h_{s,d}=\iota_s(h)\,\bigl(\iota_{s-d}(h)^{-1}\peX_{s-d}\,\iota_{s-d}(h)\bigr)^{(P^d)}\,\iota_s
(h)^{-1},\qquad\Pi^h_{s,m}=\Psi^h_{s,m-1}\cdots\Psi^h_{s,1}\Psi^h_{s,0},
\]
where $M^{(P^d)}$ is the matrix of the $P^d$-th powers of the entries. Let $\peS^h_s$ be the matrix
whose row $b$ is the row $b$ of $(\Pi^h_{s,\ell+1})^\natural$ if $x_{b,s}$ is deep and of
$(\Pi^h_{s,\ell})^\natural$ if not, and
\[
\peS^h=\bigotimes_{s\in\Z/N}\operatorname{Sym}^n(\peS^h_s)\in M_\phB(\pao[\peX]),\qquad\peS=\peS^1.
\]
Let $c_h$ be the automorphism of $\pao[\peX]$ with $\peX_j\mapsto\iota_j(h)^{-1}\peX_j\iota_j(h)$
for all $j$.

\begin{lemma}\label{pe:subst}
Let $h\in L$. {\rm(a)} $\peS=1$ at $\peX_j=1$ for all $j$. {\rm(b)} $c_h(\peS)=\rho(h)^{-1}\,
\peS^h\,\rho(h)$. {\rm(c)} $\Psi^h_{s,d}\equiv\peX_{s-d}^{(P^d)}$ modulo $p\,\pao[\peX]$, and the
entries of $\Psi^h_{s,d}$ are forms of degree $P^d$ in $\peX_{s-d}$. {\rm(d)} Every monomial of an
entry of $\peS$ has degree at most $Nnk\,P^\ell$ in each copy.
\end{lemma}
\begin{proof}
(a) is clear. (b) $c_h(\Psi^1_{s,d})=\iota_s(h)^{-1}\Psi^h_{s,d}\iota_s(h)$, hence
$c_h(\Pi^1_{s,m})=\iota_s(h)^{-1}\Pi^h_{s,m}\iota_s(h)$ and, for both parities,
$c_h(\Pi^1_{s,m})^\natural=h_s^{-1}(\Pi^h_{s,m})^\natural h_s$ (for odd $s$,
${}^t\!\operatorname{adj}(B^{-1}MB)={}^tB\;{}^t\!\operatorname{adj}(M)\,{}^tB^{-1}$ and
$h_s={}^t\iota_s(h)^{-1}$). The matrix $h_s^{-1}$ is block diagonal, and the number of factors,
$\ell+1$ or $\ell$, is the same for all rows of a block; so the row $b$ of $h_s^{-1}(\Pi^h_{s,m_b}
)^\natural$ is the row $b$ of $h_s^{-1}\peS^h_s$, and $c_h(\peS^1_s)=h_s^{-1}\peS^h_sh_s$. Apply
$\bigotimes_s\operatorname{Sym}^n$. (c) Modulo $p$ the $P^d$-th power is a ring homomorphism
$\phi$ with $\phi\circ\iota_{s-d}\equiv\iota_s$, so $(\iota_{s-d}(h)^{-1}\peX_{s-d}\iota_{s-d}(h)
)^{(P^d)}\equiv\iota_s(h)^{-1}\peX_{s-d}^{(P^d)}\iota_s(h)$. (d) The copy $j$ occurs in the letter
$s$ only in the factor $d=[s-j]$, if $d\le\ell$ (as $\ell<N$), with degree $P^d\le P^\ell$; a
cofactor has degree at most $(k-1)P^\ell$ in it, an entry of $\operatorname{Sym}^n$ at most
$n(k-1)P^\ell$, and there are $N$ letters.
\end{proof}

Now let $\varrho=e_{r+1}-e_r\in\Z^k$, and fix a copy $i$. Give the variable $X_{j;bc}$ (the entry
$(b,c)$ of $\peX_j$) the weight
\[
\omega_i(X_{i;bc})=\frac{\varrho_b}{P^\ell},\qquad\omega_i(X_{j;bc})=\frac{P^{[j-i]}}{q}\cdot
\frac{\varrho_b-\varrho_c}{P^\ell}\quad(j\ne i);\qquad|\omega_i|\le\frac2{P^\ell}.
\]
So for a monomial $\nu=(\nu_j)_j$, $\omega_i\cdot\nu=P^{-\ell}\bigl(\perw(\nu_i)+\sum_{j\ne
i}P^{[j-i]}q^{-1}(\perw(\nu_j)-\pecw(\nu_j))\bigr)$.

\begin{lemma}[Weight and defect]\label{pe:wd}
Let $Nnk(\ell+2)\le P^\ell$, $h\in L$, $c\ge1$, and let $\nu$ be a monomial whose coefficient in
some entry of the row $\beta$ of $\peS^h$ is not divisible by $w^c$. Then $|\phdef_{i+\ell}(\beta)
-\omega_i\cdot\nu|\le2c$.
\end{lemma}
\begin{proof}
Let $\varepsilon=\pm\tfrac12$; levels are taken for the weights $\varepsilon\omega_i$ on
$\pao[\peX]$, of size at most $P^{-\ell}$. Put $u=\tfrac12P^{-\ell}$ and define the
\emph{potentials}
\[
\alpha_{s,d}(b)=\varepsilon\,\frac{\varrho_b}{P^\ell}\cdot\frac{P^{\,d+[s-d-i]}}{q}\qquad(s\in\Z/N,\
d\ge0,\ 1\le b\le k);\qquad\sum_b\alpha_{s,d}(b)=0,\quad|\alpha_{s,0}(b)|\le u .
\]
\emph{Step 1: for $d\le\ell$ the entry $(b,c)$ of $\Psi^h_{s,d}$ has level at least
$\alpha_{s,d+1}(b)-\alpha_{s,d}(c)-u$.} Let $\vartheta=\varepsilon P^d\omega_i(X_{s-d;bc})$, the
weight of $X_{s-d;bc}^{P^d}$; $|\vartheta|\le1$. By Lemma~\ref{pe:subst}(c) the entry is
$X_{s-d;bc}^{P^d}$ plus a multiple of $w^2$ all of whose monomials have degree $P^d$, hence weight
at least $-1$; as $\pecl{\vartheta+1}\le2$, the entry has level at least $\vartheta$. If $s-d\ne
i$, then $[s-d-1-i]=[s-d-i]-1$ and $\vartheta=\alpha_{s,d+1}(b)-\alpha_{s,d}(c)$. If $s-d=i$, then
$\vartheta=\varepsilon\varrho_bP^{d-\ell}=\alpha_{s,d+1}(b)$ and $|\alpha_{s,d}(c)|\le\tfrac12
P^{d-\ell}q^{-1}\le u$.

\emph{Step 2.} By Lemma~\ref{pe:lev}(a) the potentials telescope: for $m\le\ell+1$ the entry
$(b,c)$ of $\Pi^h_{s,m}$ has level at least $\alpha_{s,m}(b)-\alpha_{s,0}(c)-mu\ge\alpha_{s,m}(b)-
(\ell+2)u$. A cofactor in the row $b$ is a signed sum of products of entries, one from each other
row; as $\sum_{b'}\alpha_{s,m}(b')=0$, it has level at least $-\alpha_{s,m}(b)-(k-1)(\ell+2)u$. So
the entries of the row $b$ of $\peS^h_s$ have level at least $(-1)^s\alpha_{s,m_b}(b)-k(\ell+2)u$,
$m_b=\ell+1$ or $\ell$ as $x_{b,s}$ is deep or not, and the entries of the row $\beta$ of $\peS^h$
have level at least $\sum_{b,s}\beta_{b,s}(-1)^s\alpha_{s,m_b}(b)-Nnk(\ell+2)u$.

\emph{Step 3: $(-1)^s\alpha_{s,m_b}(b)$ is $\varepsilon(1-q^{-1})$ times the defect weight of
$x_{b,s}$ at the letter $i'=i+\ell$ along $(r,r+1)$.} That weight, times $1-q^{-1}$, is
$P^{1+[s-1-i']}/q$ at $Y_s$, $-P^{[s-i']}/q$ at $X_s$, and $0$ elsewhere. For even $s$:
$Y_s=x_{r+1,s}$ is deep and $\varrho_{r+1}=1$, so $\alpha_{s,\ell+1}(r+1)=\varepsilon
P^{1+[s-1-i']}/q$; $X_s=x_{r,s}$ is not deep and $\alpha_{s,\ell}(r)=-\varepsilon P^{[s-i']}/q$.
For odd $s$ the variables $Y_s=x_{r,s}$ (deep) and $X_s=x_{r+1,s}$ exchange their roles and the
sign $(-1)^s$ compensates $\varrho_r=-1$.

Let $\Delta=\phdef_{i+\ell}(\beta)$. By Steps 2 and 3 the coefficient of $\nu$ is divisible by
$w^{\pecl y}$ with $y=\varepsilon((1-q^{-1})\Delta-\omega_i\cdot\nu)-\tfrac12Nnk(\ell+2)P^{-\ell}$;
so $y\le c-1$, and for the two signs of $\varepsilon$, $|(1-q^{-1})\Delta-\omega_i\cdot\nu|\le
2c-1$. Finally $|\Delta|\le2n\le q$ (Lemma~\ref{pe:coord}(c), and $6n\le Nnk(\ell+2)\le
P^\ell=\sqrt q$ since $N=2\ell\ge2$), so $|\Delta-\omega_i\cdot\nu|\le2c-1+|\Delta|/q\le2c$.
\end{proof}

\begin{proposition}[Averaging; Proposition~\ref{ph:av}]\label{pe:av}
Let $N=2\ell$ and $Nk(\ell+2)n+n+k<P^\ell$, and $c_0=k^2eN$. For every $c\ge1$ there is
$\mathrm{Av}\in M_\phB(\pao)$ such that $\mathrm{Av}\,h\in C_c$ for all $h\in\pao^\phB$, and
\[
{}^th'\,\mathsf b\,\mathrm{Av}\,h\equiv p^{Nc_0}\;{}^th'\,\mathsf b\,h\pmod{w^c}\qquad(h\in
\pao^\phB,\ h'\in\phA).
\]
\end{proposition}
\begin{proof}
Put $c^*=3kc$, $t=2c^*P^\ell$ and $D=Nnk(\ell+2)P^\ell$. Then $D+k<(Nk(\ell+2)n+1)P^\ell\le
P^{2\ell}=q$. For each $j$ let $\mathsf I_j$ be a form as in Lemma~\ref{pe:avg} for $\iota=
\iota_j$, $\varrho=e_{r+1}-e_r$, $t$ and $D$, and $\mathsf e=p^{c_0}\mathrm{ev}_1$; $\mathsf e$
also satisfies (I1). For a family $l=(l_j)_{j\in\Z/N}$ of linear forms on the $\pao[\peX_j]$ write
$l^\otimes=\bigotimes_jl_j$ on $\pao[\peX]$, and for $J\subseteq\Z/N$ let $l^J_j=\mathsf I_j$ for
$j\in J$ and $l^J_j=\mathsf e$ otherwise. Define
\[
\mathrm{Av}_l=l^\otimes(\peS)\in M_\phB(\pao),\qquad\mathrm{Av}_J=\mathrm{Av}_{l^J},\qquad
\mathrm{Av}=\mathrm{Av}_{\Z/N};\qquad\mathrm{Av}_\emptyset=p^{Nc_0}
\]
by Lemma~\ref{pe:subst}(a).

\emph{Frames.} Let $g\in L$, let $l$ be a family and $l'_j=l_j\circ c_{\iota_j(g)}$. Then
$l'^{\otimes}=l^\otimes\circ c_g$, and Lemma~\ref{pe:subst}(b) gives
\begin{equation}\label{pe:eq-frame}
\rho(g)\,\mathrm{Av}_{l'}=l^\otimes(\peS^g)\,\rho(g).
\end{equation}
Fix a copy $i$, a family $l$ all of whose members satisfy (I1), and $v\in\pao^\phB$. For $g\in L$
and $\beta\in\phB$ let
\[
F^g_\beta=\sum_{\gamma}(\rho(g)v)_\gamma\cdot\Bigl(\bigotimes_{j\ne i}l_j\Bigr)\bigl(\peS^g_{\beta
\gamma}\bigr)\in\pao[\peX_i],
\]
the forms $l_j$, $j\ne i$, being applied in the copies $j\ne i$. By \eqref{pe:eq-frame}, applied
to the family $l$ with $l_i$ replaced by a linear form $\mathsf m$, and (I1) for the $l_j$:
\begin{equation}\label{pe:eq-m}
\mathsf m(F^g_\beta)=\bigl(\rho(g)\,\mathrm{Av}_{l[i:\,\mathsf m\circ c_{\iota_i(g)}]}\,v\bigr)_\beta,
\qquad\text{hence}\qquad c_{\iota_i(g)}F^1_\gamma=\sum_\beta\rho(g^{-1})_{\gamma\beta}\,F^g_\beta ,
\end{equation}
the second identity because both sides have the same image under every $\mathsf m$ (take $g=1$ on
the left). By Lemma~\ref{pe:subst}(d), $\deg F^1_\beta\le D$.

\emph{Weight and defect in one copy.} Suppose that the coefficient of $X^\alpha$ in $F^g_\beta$ is
not divisible by $w^c$. Then there are $\gamma$ and a monomial $\nu$ of $\pao[\peX]$ with
$\nu_i=\alpha$, with coefficient in $\peS^g_{\beta\gamma}$ not divisible by $w^c$, and with
$l_j(\peX_j^{\nu_j})\ne0$ for all $j\ne i$. By (I4), $\perw(\nu_j)=\pecw(\nu_j)$ for $j\ne i$, so
$\omega_i\cdot\nu=\perw(\alpha)/P^\ell$, and Lemma~\ref{pe:wd} gives (as $2c\le c^*$)
\begin{equation}\label{pe:eq-AB}
\phdef_{i+\ell}(\beta)>3c^*\ \Rightarrow\ \perw(\alpha)>t;\qquad\perw(\alpha)>t\ \Rightarrow\
\phdef_{i+\ell}(\beta)>c^* .
\end{equation}

\emph{$\mathrm{Av}\,v\in C_c$.} Let $l=l^{\Z/N}$, $g\in L$, and $\beta$ with $\phdef_{i'}(\beta)>
3c^*$; write $i'=i+\ell$. By \eqref{pe:eq-m} with $\mathsf m=\mathsf I_i$, which is invariant,
$(\rho(g)\mathrm{Av}\,v)_\beta=\mathsf I_i(F^g_\beta)$. By \eqref{pe:eq-AB} every monomial
$X^\alpha$ of $F^g_\beta$ with coefficient not divisible by $w^c$ has $\perw(\alpha)>t$, and
$\mathsf I_i$ vanishes on it by (I2). So $w^c\mid(\rho(g)\mathrm{Av}\,v)_\beta$.

\emph{The pairing.} It suffices to show, for $i\notin J$ and $h'\in\phA$, that ${}^th'\mathsf b\,
\mathrm{Av}_{J\cup\{i\}}v\equiv{}^th'\mathsf b\,\mathrm{Av}_Jv$ modulo $w^c$; adding these along a
chain from $\emptyset$ to $\Z/N$ gives the statement. Let $l=l^J$ and $\psi=\sum_\beta h'_\beta
b_\beta F^1_\beta$. By \eqref{pe:eq-m} with $g=1$, the two sides are $\mathsf I_i(\psi)$ and
$\mathsf e(\psi)=p^{c_0}\psi(1)$, and $\deg\psi\le D$. By (I3) with $\mathfrak a=w^c\pao$ it suffices that
for $g\in L$ and $\perw(\alpha)>t$ the coefficient of $X^\alpha$ in $c_{\iota_i(g)}\psi$ is
divisible by $w^c$. By \eqref{pe:eq-m} this coefficient is
\[
\sum_\beta\bigl(\text{coefficient of }X^\alpha\text{ in }F^g_\beta\bigr)\cdot\bigl({}^th'\,\mathsf
b\,\rho(g^{-1})\bigr)_\beta,\qquad\bigl({}^th'\,\mathsf b\,\rho(g^{-1})\bigr)_\beta=b_\beta\,\bigl(
\rho({}^tg^{-1})h'\bigr)_\beta
\]
by \eqref{ph:eq-adj}. If the first factor is not divisible by $w^c$, then $\phdef_{i+\ell}(\beta)>
c^*\ge3c$ by \eqref{pe:eq-AB}, and $w^c$ divides the second factor by Lemma~\ref{pe:AinC}, as
${}^tg^{-1}\in L$.
\end{proof}

\subsubsection{Sizes}\label{pe:ss-sizes}
\emph{In the proof of Theorem~\ref{ph:thm}.} Assume \eqref{ph:eq-large}; put $\Lambda_q=\log q$.
As $P\ge5$ and $N,k\ge2$, we have $q\ge25$ and $\Lambda_q>3$; so $n\ge(Nk\Lambda_q)^5\ge12^5$,
and then $\Lambda_q\ge4\log n\ge49$ by $q\ge n^4$.
(1) $P^\ell=\sqrt q\ge n^2$, and $Nk(\ell+2)n+n+k\le(Nk(N+1)+2)\,n\le(Nk)^5n<n^2$; so
Proposition~\ref{ph:av} applies, and $n<q$. (2) $eN=\Lambda_q/\log p\le\tfrac32\Lambda_q$, so
$c_0=k^2eN\le\tfrac32k^2\Lambda_q$, and with $c=5Nc_0$:
\[
24\,Nk^2c\,P^{2r}\le180\,N^2k^4\Lambda_q\,P^{2r}\le(Nk\Lambda_q)^5P^{3r}\le\theta^5n\le\theta n,
\]
because $180\le N^3k\Lambda_q^4$; this is \eqref{pe:eq-hyp}. (3) $Nc_0\le\tfrac32Nk^2\Lambda_q\le
N(Nk\Lambda_q)^5\le Nn/P^{3r}\le\phcd$ and $Nc_0\ge2$, so $c-1\ge4Nc_0+1$. (4) In the conclusion,
$\det d'\det d\mid w^{\lambda+\lambda'}\mid p^{\max(\lambda,\lambda')}$, whence $l\le\phcd\nu+2Nc_0
\nu+j\theta\nu\le(3\phcd+\theta j)\nu$.

\emph{In \S\ref{ph:ss-ded}.} There $p$ is exceptional, $13\le p<x_1$, $\mathfrak r=\lceil\log_p
m\rceil$, $P=p^{o_0}$ with $1\le o_0\le\mathfrak r$, $q=p^o$ with $m\le o\le f$, $N=o/o_0\le f$,
$n=a=xb$, $r=\mathfrak r$ and $\theta=1/(m^2f)$; we use \eqref{eq:au-scales}, $\sqrt x\le m\le
2\sqrt x$, $x^5-2\sqrt x\le k\le3x^5$ and $f\le x^{12}$ \eqref{eq:au-sizes}. (a) $p^{\mathfrak r-1}<m$ gives
$\mathfrak r\le1+\log_2m\le\log x$ ($x$ is large), so $0<2\mathfrak r\le k$. (b) $q\ge13^m$ and $m\ge(\log
b)^{50}\ge4\log b$, so $q\ge b^{8}\ge a^4$. (c) $P^{3\mathfrak r}\le p^{3\mathfrak r^2}\le(pm)^{3
\mathfrak r}\le x_1^{6\log x}$, and $6\log x\log x_1\le(\log x_1)^2\le\tfrac12\log b$; so
$P^{3\mathfrak r}\le b^{1/2}$. (d) $\log q=o\log p\le f\log x_1$, so $Nk\log q/\theta\le x^{12}
\cdot3x^5\cdot x^{12}\log x_1\cdot4x^{13}\le x_1$ and $(Nk\log q/\theta)^5\le x_1^5\le b^{1/2}$.
By (c) and (d), $(Nk\log q/\theta)^5P^{3\mathfrak r}\le b\le a$, which with (b) is
\eqref{ph:eq-large}. (e) $P^{2\mathfrak r}\ge p^{2\mathfrak r}\ge mp$ and $\mathfrak r^2/(kP)\le
(\log x)^2/(kp)\le1/(mp)$, so $\phcd\le2Na/(mp)$.

\begin{remark}\label{pe:status}
Propositions~\ref{pe:root}, \ref{pe:count}, \ref{pe:av} and Lemma~\ref{pe:AinC} remove the
hypotheses of Theorem~\ref{ph:thm}. Besides these, Lemma~\ref{pa:eldiv} uses the normalisation
{\rm(N)} and weak approximation for the unitary group.
\end{remark}


%% ===== proofs/A-hecke-level.tex =====
% proofs/A-hecke-level.tex
% Lemmas pa:hecke and pa:level of proofs/A-divisibility.tex, the normalisation (N) of
% proofs/A-hecke-eldiv.tex, and weak approximation. To be \input after A-divisibility.tex,
% A-seesaw-size.tex and A-hecke-eldiv.tex. Labels of this file start with pl:.
\providecommand{\auO}{\mathcal O}
\providecommand{\auPl}{\Phi}
\providecommand{\auPa}{\Phi_a}
\providecommand{\auPb}{\Phi_b}
\providecommand{\auS}{\mathcal S}
\providecommand{\auSram}{\mathcal S^{\mathrm{ram}}}
\providecommand{\auDX}{D_X}
\providecommand{\augain}{\operatorname{gain}}
\providecommand{\auerr}{\operatorname{err}}
\providecommand{\auNm}{\operatorname{Nm}}
\providecommand{\auU}{\mathrm U}
\providecommand{\auGL}{\mathrm{GL}}
\providecommand{\auHom}{\operatorname{Hom}}
\providecommand{\psrho}{\varrho}
\providecommand{\psGlue}{\operatorname{Glue}}
\providecommand{\auA}{\mathbb A}
\providecommand{\auF}{\mathbb F}
\providecommand{\aucol}{\operatorname{col}}
\providecommand{\aubar}[1]{\overline{#1}}
\providecommand{\auf}[1]{\texttt{#1}}
\providecommand{\pao}{\mathfrak o}
\providecommand{\paCp}{\C_p}
\providecommand{\paI}{I_p^{\circ}}
\providecommand{\paded}{\operatorname{ded}}
\providecommand{\paPred}{\operatorname{Pred}}
\providecommand{\paloss}{\operatorname{loss}}
% macros of this file (prefix \pl)
\providecommand{\plone}{\mathbf 1}
\providecommand{\plD}{\mathsf D}
\providecommand{\pld}{\mathsf d}
\providecommand{\ple}{\mathsf e}
\providecommand{\plP}{\mathsf P}
\providecommand{\plSi}{\Sigma_N}
\providecommand{\plMir}{\mathrm{Mir}}
\providecommand{\pldir}{\plone^{\mathrm{dir}}}
\providecommand{\plrow}{\plone^{\mathrm{row}}}
\providecommand{\plcol}{\plone^{\mathrm{col}}}
\providecommand{\pltr}{\operatorname{tr}}
\providecommand{\plvol}{\operatorname{vol}}
\providecommand{\plr}{\mathfrak r}
\providecommand{\plE}{\mathsf E}
\providecommand{\plH}{\mathcal H}
\providecommand{\plgb}[2]{\|#2\|_{#1}}

\subsection{Hecke operators through the lifts, and the level prime: Lemmas~\ref{pa:hecke}
and~\ref{pa:level}}\label{pl:sec}

This subsection proves Lemma~\ref{pa:hecke} (\S\S\ref{pl:ss-count}--\ref{pl:ss-hecke}) and
Lemma~\ref{pa:level} (\S\ref{pl:ss-level}); the latter is first given a closed statement
(Lemma~\ref{pl:level}), for which the lattices at $N$, the operator $U_N$ and the gluing are
defined. Both proofs use Lemma~\ref{pa:kernel} only as stated, except that at $\ell=N$, which
Lemma~\ref{pa:kernel} excludes, the same integrality statement is needed and is isolated as
Lemma~\ref{pl:kernelN}; it is proved in Corollary~\ref{pf:K1}(b). \S\ref{pl:ss-norm} fixes the normalisations that
\S\ref{ss:au-spaces} leaves open and checks that they are the ones under which
Theorem~\ref{ph:thm} was applied in \S\ref{ph:ss-ded}; \S\ref{pl:ss-wa} proves weak approximation
in the form used.

\subsubsection{Normalisations}\label{pl:ss-norm}
\emph{{\rm(N1)} Weight modules.} For $\auGL_n$ and a diagram with $c(h)$ columns of height $h<n$
let $\Sigma_\lambda$ be the lattice spanned by the products of minors $\nabla_T(Y)$ of a generic
$n\times\infty$ matrix $Y$, with the standard ones as basis (\S\ref{pa:ss-count},
\S\ref{ps:ss-restrict}). We put
\[
W_\lambda=\operatorname{Hom}_\Z(\Sigma_\lambda,\C),\qquad W_\lambda(\Z)=\operatorname{Hom}_\Z
(\Sigma_\lambda,\Z),\qquad(e_T)\ \text{the basis dual to }(\nabla_T)_{T\text{ standard}},
\]
and $g\in\auGL_n(\C)$ acts by $(\rho(g)\xi)(P)=(\det g)^{e}\,\xi\bigl(P(g\,Y)\bigr)$. Thus the
evaluations satisfy $\rho(g)\mathrm{ev}_Y=(\det g)^e\mathrm{ev}_{gY}$, and the matrix of $\rho(g)$
in the basis $(e_T)$ is $(\det g)^eS_\lambda(g)$, where $\nabla_{T'}(gY)=\sum_TS_\lambda(g)_{T',T}
\nabla_T(Y)$ is the Schur matrix of \S\ref{pa:ss-count}; the same formula defines $S_\lambda(z)$
for a rectangular matrix $z$ over any ring, with entries that are integral polynomials in $z$. For
$g=\operatorname{diag}(u_1,\dots,u_n)$,
\begin{equation}\label{pl:eq-diag}
\rho(g)\,e_T=(u_1\cdots u_n)^{e}\prod_{\text{boxes }\square\text{ of }T}u_{T(\square)}\cdot e_T .
\end{equation}
For a single row of length $a$ ($a$ columns of height $1$), $\Sigma_\lambda$ is the lattice of
the forms of degree $a$ in the first column $y$ of $Y$ with its basis of monomials, and the matrix
of $\rho(g)$ in $(e_\beta)$ is the transpose of the matrix, in the monomials, of the substitution
$P(y)\mapsto P(gy)$. This is the lattice $\Sigma_\lambda^{\#}$ of \S\ref{ps:ss-restrict} and the
basis in which Lemma~\ref{pa:kernel}(K1) is stated.

\emph{{\rm(N2)} The action of the groups, and the gauge.} For $G=\auU(X)$, $X$ standard of rank
$n$, the group $\Gamma$ acts on $W_\lambda=\bigotimes_\tau W_{\lambda_\tau}$ through the embeddings
$\tau\circ c$: $\gamma$ acts by $\bigotimes_\tau\rho_\tau(\tau(\aubar\gamma))$, and
$\tau(\aubar\gamma)={}^t\tau(\gamma)^{-1}$. Let $\ell$ be a prime and $\iota\colon\C\to\C_\ell$ a
homomorphism, with valuation ring $\pao$. For $\tau\in\auPl$ let $\mathfrak l_\tau$ be the prime of
$E$ induced by $\iota\circ\tau$ and $w_\tau=\aubar{\mathfrak l_\tau}$ the prime induced by
$\iota\circ\tau\circ c$, and let $\hat\iota_\tau\colon E_{w_\tau}\to\C_\ell$ be the continuous
extension of $\iota\circ\tau\circ c$. An element $g\in G_f$ acts on $W_\lambda\otimes_{\C,\iota}
\C_\ell$ by $\rho(g_\ell)=\bigotimes_\tau\rho_\tau(\hat\iota_\tau(g_{w_\tau}))$, where $g_w\in\auGL_n
(E_w)$ is the component of $g$ at $w$; this extends the action of $\Gamma$. With $\|\cdot\|$ the
largest absolute value of a coordinate in $(e_T)$, the gauge at $(\iota,g)$ of $\xi$ is
$\plgb g\xi:=\|\rho(g_\ell)^{-1}\iota(\xi)\|$. A form $\varphi$ \emph{has gauge at most $C$ at
the base point $\pld\in G_f$} if $\plgb{g\pld}{\varphi(g)}\le C$ for all $g$; for $\pld=1$ this is
the notion of \S\ref{ss:au-spaces}. We use:
\begin{itemize}
\item[(G1)] $\plgb{\gamma g}{\gamma\xi}=\plgb g\xi$ for $\gamma\in\Gamma$; the gauge is
ultrametric in $\xi$;
\item[(G2)] if all components of $k\in G_f$ above $\ell$ are integral with integral inverse
(``$k$ is integral above $\ell$''), then $\rho(k_\ell)$ has integral entries and integral inverse,
so $\plgb{gk}\xi=\plgb g\xi$ and $\|\rho(k_\ell)\eta\|=\|\eta\|$.
\end{itemize}

\emph{{\rm(N3)} Hecke operators.} Let $\mathfrak l\ne\aubar{\mathfrak l}$ be a prime of $E$ with
uniformiser $\varpi$. The component at $\mathfrak l$ identifies $G$ at the place below $\mathfrak
l$ with $\auGL_n(E_{\mathfrak l})$ (the component at $\aubar{\mathfrak l}$ being ${}^t\aubar
g^{-1}$); for $h\in\auGL_n(E_{\mathfrak l})$ we denote by the same letter the element of $G_f$ that
is $h$ there and $1$ at the other places. Let $K_{\mathfrak l}=\auGL_n(\auO_{\mathfrak l})$ and
$d_i=\operatorname{diag}(\varpi1_i,1_{n-i})$. On a representation $(\pi,V)$ of $G_f$ and a vector
fixed by $K_{\mathfrak l}$,
\[
T_{\mathfrak l,i}\,v=\sum_{h\in K_{\mathfrak l}d_iK_{\mathfrak l}/K_{\mathfrak l}}\pi(h)v ;\qquad
\text{on forms,}\quad(T_{\mathfrak l,i}\varphi)(g)=\sum_h\varphi(gh).
\]
The operator $T_{\mathfrak l}$ of \S\ref{pa:ss-divis} is $T_{\mathfrak l,\plr_p}$ on $\auS$, with
$\varpi=p$.

\emph{{\rm(N4)} Lifts with arbitrary data.} Let $G=\auU(X_j)$, $G'=\auU(X_{j+1})$ and let $dg$ be
the Haar measure of $G_f$ for which $J_{X_j}$ has measure $1$. For $\Phi\in\plH_f(\auHom(X_j,
X_{j+1}))$ and a smooth form $\varphi$ on $X_j$ of weight $\lambda(X_j)$ put
\[
(\Theta^X_j(\Phi)\varphi)(g')=\int_{\Gamma\backslash G_f}\theta_j(\Phi;g,g')\,\varphi(g)\,dg ,
\]
a finite sum; here $\theta_j$ is \eqref{ps:eq-kernel}. We use the following properties of the
kernels, which express that $\theta$ is invariant under rational points (Lemma~\ref{ps:weil}(W1))
and that $\theta_j(\Phi;g,g')$ is defined through $\omega_{\mu_{j+1}}(g,g')\Phi$:
\begin{itemize}
\item[(T1)] $\theta_j(\Phi;\gamma g,\gamma'g')=\gamma'\circ\theta_j(\Phi;g,g')\circ\gamma^{-1}$ for
$\gamma\in\Gamma$, $\gamma'\in\Gamma'$;
\item[(T2)] $\theta_j(\omega_{\mu_{j+1}}(h,h')\Phi;g,g')=\theta_j(\Phi;gh,g'h')$ for $h\in G_f$,
$h'\in G'_f$.
\end{itemize}
By (T1) the integrand is a function on $\Gamma\backslash G_f$ and $\Theta^X_j(\Phi)\varphi$ is a
form on $X_{j+1}$; by (T2), with $R$ the right translation,
\begin{equation}\label{pl:eq-equiv}
R_{h'}\circ\Theta^X_j(\Phi)=\Theta^X_j\bigl(\omega_{\mu_{j+1}}(h,h')\Phi\bigr)\circ R_h .
\end{equation}
If $\Phi$ and $\varphi$ are fixed by the free level $J_{X_j}$, this is the sum of
\S\ref{ss:au-size}. For a tensor of data, $C_\Phi=\Theta^X_{m-1}\cdots\Theta^X_0(\Phi)$ is defined
by linearity from pure tensors, as in \S\ref{ps:ss-glue}.

\begin{lemma}[Chains are plain sums]\label{pl:sum}
Let $\Phi_{x,j}\in\plH_f(\auHom(X_j,X_{j+1}))$ $(x$ in a finite set, $0\le j<m)$ and $c_x\in\C$ be
such that the tensor $\sum_xc_x\bigotimes_j\Phi_{x,j}$ is fixed by $\bigotimes_j\omega_{\mu_{j+1}}
(b_j,b_{j+1})$ for all $b_j\in J_{X_j}$ $(j<m)$, $b_m=1$. Then for $\varphi\in\auS$ and $g_m$,
\[
\Bigl(\sum_xc_x\,\Theta^X_{m-1}(\Phi_{x,m-1})\cdots\Theta^X_0(\Phi_{x,0})\varphi\Bigr)(g_m)=
\sum_{g_0,\dots,g_{m-1}}\sum_xc_x\,\theta_{m-1}(\Phi_{x,m-1};g_{m-1},g_m)\cdots\theta_0(\Phi_{x,0};
g_0,g_1)\,\varphi(g_0),
\]
where $g_i$ runs over representatives of $\Gamma_i\backslash\auU(X_i)(\auA_{E^+,f})/J_{X_i}$.
\end{lemma}
\begin{proof}
The left side is the integral over $\prod_{j<m}\Gamma_j\backslash\auU(X_j)(\auA_f)$ of the function
$F(g_0,\dots,g_{m-1})$ given by the inner sum on the right. By (T1) it is invariant on the left
under $\prod\Gamma_j$. Replacing every $g_j$ by $g_jb_j$ with $b_j\in J_{X_j}$ replaces, by (T2),
each $\Phi_{x,j}$ by $\omega(b_j,b_{j+1})\Phi_{x,j}$ and $\varphi(g_0)$ by $\varphi(g_0b_0)=
\varphi(g_0)$; the sum over $x$ is the value of a multilinear function at the tensor, which is
unchanged. So $F$ is right invariant under $\prod J_{X_j}$, and each double coset has measure
$1$ because the levels are free (Lemma~\ref{au:free}).
\end{proof}
Formula (K0) of Lemma~\ref{pa:kernel} is the case $\Phi_{x,j}=\phi^\circ_{j,x_j}$, $c_x=\Psi'(x)$.

\emph{{\rm(N5)} The Weil representation at split places.} Let $\mathsf S$ be a finite set of primes
of $E$ with $\mathsf S\cap\aubar{\mathsf S}=\emptyset$, $E_{\mathsf S}=\prod_{w\in\mathsf S}E_w$,
and let $n=n_j$, $n'=n_{j+1}$ (both are even). By Lemma~\ref{ps:weil}(W5), $\plH_f(\auHom(X_j,
X_{j+1}))=\auS(M_{n'\times n}(E_{\mathsf S}))\otimes\plH^{\mathsf S}$, where the first factor is the
space of Schwartz--Bruhat functions of $T=(a_w^{-1}x_w)_{w\in\mathsf S}$ ($x_w$ the matrix in the frames,
$a$ the idele of \S\ref{ps:ss-weil}: $a_w=\alpha_w$ at the primes $w\mid\omega$ above $q_1$ and $a_w=1$
elsewhere; these are the coordinates $\tilde x$ of Lemma~\ref{pgx:dict}) and the second is the
representation at the other places. In these coordinates the vacuum at $\mathsf S\cup\aubar{\mathsf
S}$ is the indicator $\plone$ of $M_{n'\times n}(\auO_{\mathsf S})$, and we use (W5) in the
following explicit form: for $g\in\auGL_n(E_{\mathsf S})$, $h\in\auGL_{n'}(E_{\mathsf S})$, read as
elements of the unitary groups supported on $\mathsf S\cup\aubar{\mathsf S}$,
\begin{equation}\label{pl:eq-w5}
\omega_{\mu_{j+1}}(g,h)(\Psi\otimes v)=\frac{\chi_{n}(h)}{\chi_{n'}(g)}\,\Psi(h^{-1}Tg)\otimes v,
\qquad\chi_r(g)=\prod_{w\in\mathsf S}\chi_{j+1,w}(\det g_w)\,|\det g_w|_w^{-r/2},
\end{equation}
where $\chi_{j+1}=\mu_{j+1}\circ(x\mapsto x/\aubar x)$ is the base change of $\mu_{j+1}$ to a Hecke
character of $E$; and for $a$ supported on $\mathsf S$, $a'$ supported on $\aubar{\mathsf S}$ (in
the coordinates $T$, and with $E_{\aubar w}$ identified with $E_w$ by $c$),
\begin{equation}\label{pl:eq-heis}
\psrho(a)(\Psi\otimes v)=\Psi(T+a)\otimes v,\qquad\psrho(a')(\Psi\otimes v)=\lambda\bigl(\pltr
({}^ta'T)\bigr)\Psi\otimes v,\qquad\lambda=\prod_{w\in\mathsf S}\lambda_w,
\end{equation}
with $\lambda_w(t)=\psi_w(\jmath(a\aubar a)_wt)$, a character of $E_w$ that is trivial on
$\auO_w$ and not on $\varpi_w^{-1}\auO_w$ (self-duality of $\widehat\auO_E$ for
$\psi(\jmath\alpha\aubar\alpha\,\cdot)$, \S\ref{ps:ss-weil}, and $a\aubar a/\alpha\aubar\alpha$ a unit);
here $a'$ is written in the coordinate $a_{\aubar w}^{-1}x_{\aubar w}$ of the $\aubar w$-component. Formula \eqref{pl:eq-heis} is $[x,y]=\psi(\jmath\pltr(y^\dagger x))$ for the
polarisation $\auHom(X_{j,w},X_{j+1,w})\oplus\auHom(X_{j,\aubar w},X_{j+1,\aubar w})$; these
operators do not depend on the model, both summands being isotropic. Formula \eqref{pl:eq-w5} is
Proposition~\ref{pwf:W5}, for the normalisation (D3) of Definition~\ref{pwf:def} at split places;
for even ranks the powers of $\upsilon$ in the splitting and in the determinant twists $\upsilon_n$
of \S\ref{ps:ss-weil} cancel, leaving the base change of $\mu_{j+1}$. (Only two properties of $\chi_{j+1}$
are used below: it is unramified where $\mu_{j+1}$ is, and Lemma~\ref{pl:char}.)

\begin{remark}[The normalisation (N) of \S\ref{ph:ss-ded}]\label{pl:N}
For $\lambda^X$, (N1)--(N3) are the normalisation (N): $V_\tau=W_{\lambda^X_\tau}$ has the basis
$(e_\beta)$ dual to the monomials, there is no power of the determinant, $\gamma$ acts by
$\rho_\tau(\tau(\aubar\gamma))$, whose matrix is the transpose of that of the substitution
$x\mapsto x\cdot{}^t\tau(\aubar\gamma)$ on row vectors, and $T_{\mathfrak l}$ is the operator of
(N3). With this choice the module $\auS_\iota$ of \S\ref{pa:ss-divis} is the image of
$f\colon\mathcal F_X\to\auS$, $F\mapsto(\varphi\colon\varphi(x)=\sum F(x)_{\eta,\beta}\,e_\eta
\otimes e_\beta)$, by (G1), (G2) and the choice of representatives $x$ integral above $p$
(Corollary~\ref{pl:wacor}(a)); and for $\gamma_e=\beta_\varkappa\mathsf d\gamma(e)^{-1}$ the
substitution is by $\tau_s(\gamma_e)^{-1}$, that is, by $\iota_s(\gamma(e))\cdot\iota_s(\mathsf
d^{-1})\cdot\iota_s(\beta_\varkappa^{-1})$ for even $s$ and by $\iota_s({}^t\gamma(e)^{-1})\cdot
\mathsf d\cdot\iota_s({}^t\beta_\varkappa)$ for odd $s$; substitutions compose in the order
written, and $p\,\mathsf d^{-1}$, $\mathsf d$ multiply the deep variables of an even, respectively
odd, letter by $p$; so $p^{A_{\mathfrak l}}$ times its matrix on $\bigotimes_sV_{\tau_s}$ is
$\rho(\gamma(e))\Delta^2\rho(t(\varkappa))$, and $\gamma_e$ acts on coordinates by the transpose. This is \eqref{ph:eq-A}; so Theorem~\ref{ph:thm}
is applied in \S\ref{ph:ss-ded} to the lattice $\auS_\iota$ of Theorem~\ref{pa:divis}. Had
$W_\lambda(\Z)$ been the lattice $\Sigma_\lambda$ itself, or had $\Gamma$ acted through $\tau$,
the operator would not have the form \eqref{ph:eq-A}.
\end{remark}

\subsubsection{Weak approximation}\label{pl:ss-wa}
\begin{proposition}\label{pl:wa}
Let $F$ be a number field, $E/F$ a quadratic extension with conjugation $c$, $V$ a non-degenerate
hermitian space over $E$, $G=\auU(V)$ and $\mathsf T$ a finite set of places of $F$. Then $G(F)$ is
dense in $\prod_{v\in\mathsf T}G(F_v)$.
\end{proposition}
\begin{proof}
For an $F$-algebra $R$ let $\mathfrak s(R)$ be the $R$-module of the $X\in\operatorname{End}_{E
\otimes R}(V\otimes R)$ with $X^\dagger=-X$, and $\mathfrak s^\circ$, $G^\circ$ the subsets of
$\mathfrak s$, $G$ where $1+X$, respectively $1+u$, is invertible. The Cayley transform
$X\mapsto(1-X)(1+X)^{-1}$ is a bijection $\mathfrak s^\circ(R)\to G^\circ(R)$ with inverse
$u\mapsto(1-u)(1+u)^{-1}$: if $X^\dagger=-X$ then $u^\dagger u=(1-X)^{-1}(1+X)(1-X)(1+X)^{-1}=1$
and $1+u=2(1+X)^{-1}$; if $u^\dagger=u^{-1}$ then $X^\dagger=(1+u^{-1})^{-1}(1-u^{-1})=-X$. For
$R=F_v$ it is a homeomorphism. Now $\mathfrak s(F)$ is a finite-dimensional $F$-vector space, so it
is dense in $\prod_{v\in\mathsf T}\mathfrak s(F_v)$ (weak approximation in $F$), and
$\prod_v\mathfrak s^\circ(F_v)$ is open; hence $G^\circ(F)$ is dense in $\prod_vG^\circ(F_v)$. It
remains to see that $G^\circ(F_v)$ is dense in $G(F_v)$. Let $u\in G(F_v)$. The group $\{\zeta\in
(E\otimes F_v)^\times:\zeta\aubar\zeta=1\}$ contains the elements $x/\aubar x$, which form an
infinite set accumulating at $1$; all but finitely many of them satisfy $\det(1+\zeta u)\ne0$, and
$\zeta u\in G(F_v)$.
\end{proof}

\begin{corollary}\label{pl:wacor}
Let $X$ be a standard space, $G=\auU(X)$, $\ell$ a prime.
{\rm(a)} For $g,d\in G_f$ there is $\gamma\in\Gamma$ such that $\gamma g$ and $d\,\gamma g\,d^{-1}$
are integral above $\ell$. In particular every class of $\Gamma\backslash G_f/J$ has a
representative that is integral above $\ell$.
{\rm(b)} If $\varphi$ is a form of level $J\subseteq K^\circ$ and $(x)$ a system of representatives
that are integral above $\ell$, and if $J$ is the full group $\auGL_n(\auO_w)$ at the places above
$\ell$, then $\varphi$ is integral at $\iota$ if and only if all $\iota(\varphi(x))$ have
coordinates in $\pao$.
\end{corollary}
\begin{proof}
(a) The set of the $y\in\prod_{v\mid\ell}G(E^+_v)$ with $yg_\ell$ and $d_\ell yg_\ell d_\ell^{-1}$
integral is open and contains $g_\ell^{-1}$; apply Proposition~\ref{pl:wa}. (b) By (G1), (G2),
$\plgb{\gamma xk}{\varphi(\gamma xk)}=\plgb{xk}{\varphi(x)}=\|\iota(\varphi(x))\|$.
\end{proof}
This is the form of weak approximation used in \S\ref{ph:ss-ded}.

\subsubsection{Values of the characters}\label{pl:ss-char}
\begin{lemma}\label{pl:char}
Let $1\le i\le m$, $\chi=\chi_i$ the base change of $\mu_i=\mu_i^0\vartheta^i$, and $w\ne\aubar w$
a prime of $E$ above $\ell$, unramified over $\Q$, not chosen and not in $Q_\vartheta\cup\aubar{
Q_\vartheta}$. Then $\chi_w$ is unramified, $\chi_w(\ell)$ is an algebraic number, and for
$\iota\colon\C\to\C_\ell$, with $v$ the valuation normalised by $v(\ell)=1$,
\[
\bigl|v\bigl(\iota\,\chi_w(\ell)\bigr)\bigr|=\Bigl|\sum_{\tau:\,\mathfrak l_\tau=w}\kappa_i(\tau)-
\sum_{\tau:\,\mathfrak l_\tau=\aubar w}\kappa_i(\tau)\Bigr| .
\]
\end{lemma}
\begin{proof}
$\mu_i$ is unramified at the place below $w$, so $\chi_w$ is trivial on $\auO_w^\times$. Let $h$ be
the class number of $E$ and $w^h=(\beta)$, $\beta\in\auO_E$. As $\chi$ is trivial on $E^\times$,
$\chi_w(\ell)^h=\chi_w(\beta)=\chi_\infty(\beta)^{-1}\prod_{w'\ne w}\chi_{w'}(\beta)^{-1}$. The
element $\beta$ is a unit at every $w'\ne w$, and a continuous character of the compact totally
disconnected group $\prod_{w'\ne w}\auO_{w'}^\times$ has finite image; so the last product is a root
of unity (its factors are $1$ at almost all $w'$). The character $\mu_i$ has weight $-\kappa_i$ at the infinite
places, so $\chi_\infty(\beta)=\prod_\tau(\tau(\beta)/\tau(\aubar\beta))^{\mp\kappa_i(\tau)}$ with
one sign for all $\tau$. Finally $v(\iota\tau(\beta))$ is $h$ if $\mathfrak l_\tau=w$ and $0$
otherwise, and $v(\iota\tau(\aubar\beta))$ is $h$ if $\mathfrak l_\tau=\aubar w$ and $0$
otherwise.
\end{proof}

\subsubsection{The count with charges, and the gain with marked places}\label{pl:ss-count}
We prove the variant of Proposition~\ref{pa:count} announced in Remark~\ref{pa:charges}. Notation
and hypotheses are those of \S\ref{pa:ss-count}.

\begin{proposition}[Count with charges]\label{pl:charges}
Let $F$ have the normal form \eqref{pa:eq-nf}, let $\kappa$ be a shape with $B\ge1$ boxes, $T$ a
filling and $y$ integral directions. Let $\plE\subseteq\auPl$ and let integers $w_{\tau,h,b}\ge0$
$(\tau\in\auPl$, $1\le h<n'$, $b\in[n'])$ be such that for all $e\in\paI$, $\tau$, $b$, $h'$, with
$\tau'=\mathrm{pr}_e\tau$:
\begin{equation}\label{pl:eq-chargehyp}
\tau\in\plE,\ \text{or}\ \ (h'=1\ \text{and}\ \tau'\in\plE)\quad\Longrightarrow\quad
w_{\tau,1,b}>0\ \ \text{or}\ \ w_{\tau',h',b}>0 .
\end{equation}
Let $\kappa^\circ$ be $\kappa$ without its columns of height one at the places of $\plE$, with
$B^\circ$ boxes, and $w(T)=\sum w_{\tau,h,b}$, the sum over the boxes of $T$, a box in row $b$ of a
column of height $h$ at $\tau$ counting $w_{\tau,h,b}$. Then
\[
v\bigl(Q_{T,y}(\partial)F(0)\bigr)+w(T)\ \ge\ \paPred(\kappa^\circ)+\gamma\,(B-B^\circ)-f\,n'^3\,
(1+\log_pB).
\]
\end{proposition}
\begin{proof}
We follow the four steps of the proof of Proposition~\ref{pa:count}, with the substitution
\[
z_{\tau,b,a}(s)=\sum_{h,k}p^{\,w_{\tau,h,b}}\,\beta_{\tau,h,k}(a)\,s_{\tau,h,k,b}.
\]
\emph{Step 1.} Now $s_{\tau,h,k,b}$ is replaced by $p^{w_{\tau,h,b}}\sum_{k'}g_{\tau,h}(k,k')(z_\tau
y_\tau)_{b,k'}$, so the matrix of a column $(\tau,h,c)$ with rows $b_1,\dots,b_h$ is multiplied in
addition on the left by $\operatorname{diag}(p^{w_{\tau,h,b_i}})$, and \eqref{pa:eq-chain} becomes
$P_*(\partial)G(0)=\varpi^{B}p^{\,w(T)}\prod_{\tau,h}\det(g_{\tau,h})^{\kappa_\tau(h)}\cdot Q_{T,y}
(\partial)F(0)$.

\emph{Step 2.} Modulo $p$, the coefficient of $s_{\tau,h,k,b}^{\,p^e}s_{\tau',h',k',b}$ in
$\Gamma_e(z(s))$ is $p^{\,p^ew_{\tau,h,b}+w_{\tau',h',b}}[\beta_{\tau,h,k},\beta_{\tau',h',k'}]_e$.
By \eqref{pl:eq-chargehyp} it is divisible by $p$, hence by $\pi$, for the pairs of heights
$(1,h')$ if $\tau\in\plE$, and for the pair $(1,1)$ if $\tau'\in\plE$; call these pairs
\emph{free}. For the others we impose the congruences (a), (b), (c) of the proof of
Proposition~\ref{pa:count}, with the following changes in (b), for a pair $(\tau,e)$: if
$\tau\in\plE$ nothing is imposed; if $\tau\notin\plE$ and $\tau'\in\plE$, either $\pi_e$ divides
$g_{\tau,1}$, or the congruence modulo $\pi$ is imposed on the rows of the $g_{\tau',h'}$ with
$h'\ge2$ only, whichever is cheaper. Every pair of heights that is not free is covered as before.
The cost of (b) for $(\tau,e)$ is now $v(\pi)\min(a^\circ_\tau/p^e,a^\circ_{\tau'}+a'_{\tau'})$,
where $a^\circ_\sigma=a_\sigma$ for $\sigma\notin\plE$ and $a^\circ_\sigma=0$ for
$\sigma\in\plE$; the costs of (a) and (c) are unchanged. So \eqref{pa:eq-goal} holds with the
deductions $\paded^\circ_e(\tau)$ of the shape $\kappa^\circ$ (whose numbers are $a^\circ_\tau$,
$a'_\tau$, $l_\tau$).

\emph{Steps 3, 4} are unchanged and give $v(P_*(\partial)G(0))\ge\frac B{p-1}-fn'^3(1+\log_pB)$.
Hence
\[
v\bigl(Q_{T,y}(\partial)F(0)\bigr)+w(T)\ge\gamma B-v(\pi)\sum_{e\in\paI}\sum_\tau\paded^\circ_e
(\tau)-fn'^3(1+\log_pB),
\]
and $\gamma B^\circ-v(\pi)\sum_{e,\tau}\paded_e^\circ(\tau)=\paPred(\kappa^\circ)$ by
\eqref{pa:eq-pred}.
\end{proof}

For $p\in\mathcal N$, $0\le j<m$ and $\plE\subseteq\auPl$ let $\kappa^\circ$ be the shape of the
$j$-th lift (\S\ref{pa:ss-divis}) without its columns of height one at $\plE$, and
\[
D_j(\plE)=\paPred(\kappa^\circ)+\gamma_p\,a\,\#(\auPa\cap\plE)-f\,n_{j+1}^3(1+\log_pB_j),\qquad
C_p(\plE)=\sum_{(\tau,e)}\frac a{p^e},
\]
the last sum over the pairs with $e\in I_p$ and $\tau,u(-p^e)\tau\in\auPa\setminus\plE$.

\begin{lemma}\label{pl:gain}
Let $p\in\mathcal N$ be inert in $K$. Then
$D_j(\plE)\ge\dfrac{\augain_p}m-C_p(\plE)+\#(\auPa\cap\plE)\,\dfrac a{x_1}-\dfrac{f\,a\,(\log b)^2}
{p\,m}$.
\end{lemma}
\begin{proof}
This is the proof of Lemma~\ref{pa:gain} with two changes; recall $v(\pi)=1$, $\gamma_p=\sum_{e\in
I_p}p^{-e}$, and that the numbers of $\kappa^\circ$ are $a^\circ_\tau$, $a'_\tau$, $l_\tau$. In
(1), for $\tau\in\auPa\setminus\plE$: if $\tau'\notin\auPa\setminus\plE$ then $a^\circ_{\tau'}=0$
and $(a/p^e-a^\circ_{\tau'}-a'_{\tau'})^+\ge(a/p^e-\kappa(\tau'))^+-s$; otherwise $(\tau,e)$ is one
of the pairs of $C_p(\plE)$ and the loss is at most $a/p^e$. For $\tau\in\auPa\cap\plE$ the shape
$\kappa^\circ$ contributes nothing, but the term $\gamma_pa$ of $D_j(\plE)$ gives
\[
\sum_{e\in I_p}\frac a{p^e}\ \ge\ \sum_{e\in I_p}\Bigl(\frac a{p^e}-\kappa(\tau')\Bigr)^++\min\Bigl(
\frac ap,\kappa(u(-p)\tau)\Bigr)\ \ge\ \sum_{e\in I_p}\Bigl(\frac a{p^e}-\kappa(\tau')\Bigr)^++\frac
a{x_1},
\]
because $1\in I_p$, $p<x_1$ and $\kappa\ge a/x_1$. Steps (2), (3) and the comparison of sizes are
those of Lemma~\ref{pa:gain}.
\end{proof}

\subsubsection{Hecke operators through one lift and through the chain}\label{pl:ss-transfer}
Let $\mathfrak l\ne\aubar{\mathfrak l}$ be a prime of $E$, with residue field of $q$ elements and
uniformiser $\varpi$, at which $\mu_{j+1}$ is unramified. We use (N5) for $\mathsf S=\{\mathfrak
l\}$; write $K=\auGL_n(\auO_{\mathfrak l})$, $K'=\auGL_{n'}(\auO_{\mathfrak l})$, $n'=n+k_0$ with
$k_0=t+1$, and on data fixed by $\omega(K,K')$ put
\[
T'_i=\sum_{h\in K'd_iK'/K'}\omega_{\mu_{j+1}}(1,h),\qquad T^-_i=\sum_{g\in Kd_i^{-1}K/K}
\omega_{\mu_{j+1}}(g,1),\qquad P'(X)=\sum_{i\ge0}(-X)^iq^{\binom i2}T'_i,
\]
and $P^-(X)$ likewise with $T^-_i$. For a representation as in (N3) let $P_{\mathfrak l}(X)=\sum_i
(-X)^iq^{\binom i2}T_{\mathfrak l,i}$.

\begin{lemma}[The vacuum]\label{pl:vac}
Let $\alpha=\chi_{j+1,\mathfrak l}(\varpi)q^{n/2}$ and $\beta=\chi_{j+1,\mathfrak l}(\varpi)
q^{n'/2}$. For every $v$, the vector $\plone\otimes v$ is fixed by $\omega(K,K')$, and
\[
P'(X)(\plone\otimes v)=\prod_{l<k_0}(1-q^l\alpha X)\cdot P^-\bigl(q^{k_0}\alpha\beta^{-1}X\bigr)
(\plone\otimes v).
\]
\end{lemma}
\begin{proof}
The first claim is \eqref{pl:eq-w5}, $\chi_{j+1,\mathfrak l}$ being unramified. For $h\in
K'd_iK'$, \eqref{pl:eq-w5} gives $\omega(1,h)\plone=\alpha^i\,\plone(h^{-1}T)$, and $\plone(h^{-1}
T)=1$ if and only if $T$ is integral and the image of $T$ modulo $\varpi$ lies in the image of
$h\auO^{n'}$ in $\auF_q^{n'}$, a subspace of codimension $i$; every such subspace occurs for
exactly one coset. So $T'_i\plone$ is $\alpha^i\plone$ times the number of subspaces of dimension
$i$ of a space of dimension $n'-\operatorname{rk}(T\bmod\varpi)$. Similarly, for $g\in
Kd_i^{-1}K$, $\omega(g,1)\plone=\beta^i\plone(Tg)$, and $\plone(Tg)=1$ if and only if $T$ is
integral and $T$ modulo $\varpi$ kills the image of $\varpi g\auO^n$ in $\auF_q^n$, a subspace of
dimension $i$; so $T^-_i\plone$ is $\beta^i\plone$ times the number of subspaces of dimension $i$
of $\ker(T\bmod\varpi)$, of dimension $n-\operatorname{rk}(T\bmod\varpi)$. For a space $V$ of
dimension $d$ over $\auF_q$,
\begin{equation}\label{pl:eq-qbin}
\sum_{U\subseteq V}(-X)^{\dim U}q^{\binom{\dim U}2}=\prod_{l<d}(1-q^lX),
\end{equation}
by induction on $d$: of the subspaces of dimension $i$, those in a fixed hyperplane $H$ are
counted by the case $d-1$, each subspace of dimension $i-1$ of $H$ is $U\cap H$ for exactly
$q^{d-i}$ subspaces $U\not\subseteq H$ of dimension $i$, and $\binom i2-\binom{i-1}2=i-1$; so the
sum for $V$ is $(1-q^{d-1}X)$ times the sum for $H$. Hence, with
$r=\operatorname{rk}(T\bmod\varpi)$, $P'(X)\plone=\prod_{l<n'-r}(1-q^l\alpha X)\plone$ and
$P^-(Y)\plone=\prod_{l<n-r}(1-q^l\beta Y)\plone$ as functions of $T$, and $n'-r=k_0+(n-r)$.
\end{proof}

\begin{proposition}[One lift]\label{pl:onelift}
Let $\Phi=\plone\otimes v$ be a datum that is the vacuum at $\mathfrak l,\aubar{\mathfrak l}$, and
let $\varphi$ be a smooth form on $X_j$ fixed by $K$. Then $\Theta^X_j(\Phi)\varphi$ is fixed by
$K'$, and with $z=q^{1/2}$ and $\alpha$ as in Lemma~\ref{pl:vac},
\[
P_{\mathfrak l}(X)\,\Theta^X_j(\Phi)\varphi=\prod_{l<t+1}(1-q^l\alpha X)\cdot\Theta^X_j(\Phi)\bigl(
P_{\mathfrak l}(z^{\,t+1}X)\varphi\bigr).
\]
\end{proposition}
\begin{proof}
By \eqref{pl:eq-equiv}, $R_{h'}\Theta^X_j(\Phi)\varphi=\Theta^X_j(\omega(1,h')\Phi)\varphi$; this
gives the first claim and $T_{\mathfrak l,i}\Theta^X_j(\Phi)\varphi=\Theta^X_j(T'_i\Phi)\varphi$.
Also by \eqref{pl:eq-equiv}, $\Theta^X_j(\omega(g,1)\Phi)\varphi=\Theta^X_j(\Phi)(R_{g^{-1}}
\varphi)$, and $\Theta^X_j(\Phi)\circ R_k=\Theta^X_j(\Phi)$ for $k\in K$. Therefore, with $dy$ the
Haar measure of $\auGL_n(E_{\mathfrak l})$ giving $K$ the measure $1$,
\[
\Theta^X_j(T^-_i\Phi)\varphi=\Theta^X_j(\Phi)\Bigl(\sum_{x\in K\backslash Kd_iK}\int_KR_{kx}\varphi
\,dk\Bigr)=\Theta^X_j(\Phi)\Bigl(\int_{Kd_iK}R_y\varphi\,dy\Bigr)=\Theta^X_j(\Phi)(T_{\mathfrak l,i}
\varphi).
\]
Apply $\Theta^X_j(\,\cdot\,)\varphi$ to Lemma~\ref{pl:vac}; $q^{k_0}\alpha/\beta=q^{k_0/2}$.
\end{proof}

From now on $p$ is an exceptional prime, $\Lambda$ a full mark, $\plr=\plr_p$, $\vartheta$ a twist
with $p\nmid\auNm Q_\vartheta$, $\Psi$ a chain datum and $\Psi'$ its expansion
(Proposition~\ref{au:expansion}). Every $x$ in the support of $\Psi'$ is integral above $p$, since
$p\nmid MN\auNm Q_\vartheta$; so $\phi^\circ_{j,x_j}$ is the vacuum at all places above $p$, and
\begin{equation}\label{pl:eq-fix}
\omega_{\mu_{j+1}}(k,k')\,\phi^\circ_{j,x_j}=\phi^\circ_{j,x_j}\qquad\text{for }k,k'\text{ in }
K^\circ\text{ and supported above }p,
\end{equation}
by Lemma~\ref{ps:weil}(W4) ($\omega(k,k')$ fixes the vacuum of a lattice that it stabilises, and
commutes with the translation by the part of $x_j$ away from $p$) and because $\mu_{j+1}$ is
unramified above $p$. We recall from \S\ref{ph:ss-ded}, for $\mathfrak l\in\Lambda$, $\mathfrak
l=\mathfrak l_{\tau_0}$ with $\tau_0\in\auPa$: $p$ is a uniformiser at $\mathfrak l$; the places
inducing $\mathfrak l$ or $\aubar{\mathfrak l}$ are the $u((-p)^i)\tau_0$, those with even $i$
inducing $\mathfrak l$; the order $o$ of $-p$ modulo $q_1$ is even; and the residue field of
$\mathfrak l$ has $q=p^{o}$ elements, $o\le f$.

\begin{lemma}[Balance]\label{pl:balance}
Let $\mathfrak l\in\Lambda$. No place inducing $\mathfrak l$ or $\aubar{\mathfrak l}$ is a
$b$-place; so $\kappa(\tau)=a/x_1$, $\kappa_i(\tau)$ and $e_{i,\tau}$ are the same for all these
places, and as many of them induce $\mathfrak l$ as $\aubar{\mathfrak l}$. Consequently
$\iota(\chi_{i,\mathfrak l}(p))$ is a unit for $1\le i\le m$. Moreover, for odd $e$ the place
$u(-p^e)\tau$ induces $\aubar{\mathfrak l_\tau}$.
\end{lemma}
\begin{proof}
The classes modulo $q_0$ of these places are powers of $-p$, which has odd order modulo $q_0$; so
none is $-1$, the class of the $b$-places. The second assertion holds because $o$ is even, and the
third follows from Lemma~\ref{pl:char}. The last is the description just recalled, as $-p^e=
(-p)^e$.
\end{proof}

\begin{corollary}[The chain]\label{pl:chainT}
Let $\mathfrak l\in\Lambda$ and $\nu_p=f\,\plr\,(n_m+\plr)$. For $x$ in the support of $\Psi'$ and
$0\le e\le m$ let $C_{x,e}=\Theta^X_{e-1}(\phi^\circ_{e-1,x_{e-1}})\cdots\Theta^X_0(\phi^\circ_{0,
x_0})$. For each $1\le e\le m$ there are $c_i\in\pao$ $(0\le i\le\plr)$, independent of $x$, such
that for every smooth form $\varphi$ on $X_0$ fixed by $\auGL_k(\auO_{\mathfrak l})$,
\[
\sum_{i\le\plr}c_i\,T_{\mathfrak l,i}\,C_{x,e}\varphi=p^{\,\nu_p}\,C_{x,e}\bigl(T_{\mathfrak l,
\plr}\varphi\bigr).
\]
\end{corollary}
\begin{proof}
\emph{One lift.} In Proposition~\ref{pl:onelift}, $R(X)=\prod_{l<t+1}(1-q^l\alpha X)$ has
coefficients in $\pao$ by Lemma~\ref{pl:balance} and $R(0)=1$. Dividing by $R$ in $\pao[[X]]$ and
comparing the coefficients of $X^i$: for $i\le n_j$ there are $c_{i,i'}\in\pao$, depending only on
$j$ and $\mathfrak l$, with
\[
z^{(t+1)i}q^{\binom i2}\,\Theta^X_j(\Phi)(T_{\mathfrak l,i}\varphi)=\sum_{i'\le i}c_{i,i'}\,
T_{\mathfrak l,i'}\,\Theta^X_j(\Phi)\varphi .
\]
Here $z^{(t+1)i}q^{\binom i2}=p^{\,oi(t+i)/2}$ divides $p^{\,f\plr t}$ for $i\le\plr$, because
$o\le f$ and $\plr\le t$. \emph{Induction on $e$:} we claim $\sum_{i\le\plr}c^{(e)}_iT_{\mathfrak
l,i}C_{x,e}\varphi=p^{\,ef\plr t}C_{x,e}(T_{\mathfrak l,\plr}\varphi)$ with $c_i^{(e)}\in\pao$.
For $e=0$ take $c^{(0)}_i=\delta_{i\plr}$. The form $C_{x,e}\varphi$ is fixed by the integral
points at $\mathfrak l$ (Proposition~\ref{pl:onelift}), so applying $p^{f\plr t}\Theta^X_e
(\phi^\circ_{e,x_e})$ to the identity for $e$ and using the case of one lift for each $i\le\plr\le
k$ gives the identity for $e+1$. Finally multiply by $p^{\,\nu_p-ef\plr t}$, an integer because
$e\,t\le m\,t\le n_m$.
\end{proof}

\subsubsection{A lift at a neighbour; proof of Lemma~\ref{pa:hecke}}\label{pl:ss-hecke}
Let $\plE=\{\tau:\mathfrak l_\tau\in\Lambda\}$ and $\aubar\plE=\{\tau:\aubar{\mathfrak l_\tau}\in
\Lambda\}$; they are disjoint and $\Phi_a\subseteq\plE\cup\aubar\plE$, $\Lambda$ being a full
mark. With the numbers $a'_\tau$ of \eqref{pa:eq-diag} for the $j$-th lift put
\[
\mathrm L_j=\plr\sum_{\tau\in\plE}a'_\tau-D_j(\plE),\qquad\sum_{\mathfrak l\in\Lambda}A_{\mathfrak l}=a\cdot
\#(\auPa\cap\plE).
\]

\begin{lemma}[One lift at a neighbour]\label{pl:neigh}
Assume Lemma~\ref{pa:kernel}. Let $0\le j<m$, let $x$ be in the support of $\Psi'$, and let $h\in
\auU(X_{j+1})(\auA_{E^+,f})$ be supported on $\Lambda\cup\aubar\Lambda$ with component in
$K'd_{i_{\mathfrak l}}K'$, $i_{\mathfrak l}\le\plr$, at each $\mathfrak l\in\Lambda$. Then for all
$g,g'$ and $\xi$,
\[
\plgb{g'}{\theta_j\bigl(\omega_{\mu_{j+1}}(1,h)\phi^\circ_{j,x_j};g,g'\bigr)\xi}\ \le\ p^{\,\sum_{
\mathfrak l\in\Lambda}A_{\mathfrak l}+\mathrm L_j}\ \plgb g\xi .
\]
\end{lemma}
\begin{proof}
\emph{Reduction to a diagonal element.} By the theory of elementary divisors $h=u\,d\,u'$ with
$u,u'$ integral and supported on $\Lambda\cup\aubar\Lambda$, and $d$ diagonal: at $\mathfrak
l\in\Lambda$ its entries are $p$ on a set $J(\mathfrak l)\subseteq[n_{j+1}]$ of $i_{\mathfrak l}$
indices and $1$ elsewhere, and at $\aubar{\mathfrak l}$ they are the inverses. By (T2), the
left side concerns $\theta_j(\phi^\circ;g,g'h)$, where $\phi^\circ=\phi^\circ_{j,x_j}$. By
Corollary~\ref{pl:wacor}(a) there are $\gamma\in\Gamma$, $\varepsilon\in\Gamma'$ such that
$g_0=\gamma g$, $g_1=\varepsilon g'ud$ and $dg_1d^{-1}$ are integral above $p$. By (T1),
$\theta_j(\phi^\circ;g,g'h)=\varepsilon^{-1}\theta_j(\phi^\circ;g_0,g_1u')\gamma$. Let $\dot g_0$,
$\dot g_1$ be $g_0$, $g_1u'$ with their components above $p$ replaced by $1$. By (T2) and
\eqref{pl:eq-fix}, $\theta_j(\phi^\circ;g_0,g_1u')=\theta_j(\phi^\circ;\dot g_0,\dot g_1)=:\theta$.
Since $(\varepsilon g')_p^{-1}=u_p\,d_p\,g_{1,p}^{-1}$, and $\rho(d)\rho(g_1)^{-1}=\rho(dg_1^{-1}
d^{-1})\rho(d)$, (G1) and (G2) give
\[
\plgb{g'}{\theta_j(\phi^\circ;g,g'h)\xi}=\bigl\|\rho(d_p)\,\theta\,\rho(g_{0,p})\,\eta\bigr\|,
\qquad\eta=\rho(g_{0,p})^{-1}\iota(\gamma\xi),\quad\|\eta\|=\plgb g\xi .
\]
The components of $\dot g_0,\dot g_1$ above $p$ are trivial. Hence, by the definition
\eqref{ps:eq-kernel} of the kernel, the matrix of $\theta$ in the bases $(e_T)$ is the matrix
$(J_{T',T})=(S_{\lambda_{j+1}}(\partial)_{T',T}F(0))$ of the jets of the series $F=\Theta_f(\pi_M
s_0\,\cdot)$, $f=\omega_{\mu_{j+1}}(\dot g_0,\dot g_1)\phi^\circ$, itself, with no further factors:
for such points one takes $\gamma=\gamma'=1$ in the proof of Proposition~\ref{pk:K12}, so that the
two invertible matrices of (K1), which are the matrices of the components above $p$, are the
identity. (This is needed here, because $\rho(d_p)$ is not integral and could not be moved across an
arbitrary invertible matrix.) The datum $f$ is a unit multiple of a standard datum at $p$ by
Proposition~\ref{pk:std}(S1), so $F$ has the normal form by Proposition~\ref{pk:std}(S3) and
Proposition~\ref{pa:product}, as in (K2). As $\rho(g_{0,p})$ is integral and $\rho(d_p)$
is diagonal, it suffices to show, for all standard $T'$ and all $T$,
\begin{equation}\label{pl:eq-entry}
v\bigl(\rho(d_p)_{T'}\bigr)+v(J_{T',T})\ \ge\ -\sum_{\mathfrak l\in\Lambda}A_{\mathfrak l}-\mathrm L_j ,
\end{equation}
where $\rho(d_p)_{T'}$ is the eigenvalue \eqref{pl:eq-diag} on $e_{T'}$.

\emph{The eigenvalue.} At $\tau$ the action is through the component at $w_\tau=\aubar{\mathfrak
l_\tau}$. So $\hat\iota_\tau(d_{w_\tau})$ is diagonal with entries $p^{-1}$ on $J(\mathfrak l_\tau)$
if $\tau\in\plE$, $p$ on $J(\aubar{\mathfrak l_\tau})$ if $\tau\in\aubar\plE$, and it is $1$ for
the other $\tau$. Call an index $b$ \emph{marked at $\tau$} if it lies in this set $J$. The
determinants contribute $\sum_{\mathfrak l\in\Lambda}\#J(\mathfrak l)\bigl(\sum_{\tau:\aubar{
\mathfrak l_\tau}=\mathfrak l}e_{j+1,\tau}-\sum_{\tau:\mathfrak l_\tau=\mathfrak l}e_{j+1,\tau}
\bigr)=0$ to the valuation, by Lemma~\ref{pl:balance}. Hence
\[
v\bigl(\rho(d_p)_{T'}\bigr)=\#\{\text{boxes of }T'\text{ at }\aubar\plE\text{ with marked entry}\}-
\#\{\text{boxes of }T'\text{ at }\plE\text{ with marked entry}\}.
\]
\emph{The jets.} Define charges: $w_{\tau,h,b}=1$ if $\tau\in\plE$, $h=1$ and $b$ is not marked at
$\tau$, or if $\tau\in\aubar\plE$ and $b$ is marked at $\tau$; and $w_{\tau,h,b}=0$ otherwise.
They satisfy \eqref{pl:eq-chargehyp} for $\plE$ (an exceptional prime $p$ is inert in $K$, so
$p\ne d$ and $\paI=I_p$ is the set of odd positive integers): let $e\in I_p$, $\tau'=u(-p^e)\tau$; by
Lemma~\ref{pl:balance} one of $\tau,\tau'$ is in $\plE$ exactly when the other is in $\aubar\plE$,
and then the same indices are marked at both. If $\tau\in\plE$: a marked $b$ has $w_{\tau',h',b}=
1$ and an unmarked one has $w_{\tau,1,b}=1$. If $\tau'\in\plE$ and $h'=1$: a marked $b$ has
$w_{\tau,1,b}=1$ and an unmarked one has $w_{\tau',1,b}=1$. By Lemma~\ref{pa:schur} and
Proposition~\ref{pl:charges}, $v(J_{T',T})\ge D_j(\plE)-w(T')$.

\emph{Count of boxes.} A box at $\aubar\plE$ with marked entry counts $+1$ in $v(\rho(d_p)_{T'})$
and $1$ in $w(T')$; a box at $\plE$ in a column of height one counts $-1$ or $0$ in the first and
$0$ or $1$ in the second, according as its entry is marked or not; in a column of height $h\ge2$
at $\plE$ the entries are distinct, so at most $\plr$ of them are marked. Therefore
\[
v\bigl(\rho(d_p)_{T'}\bigr)-w(T')\ \ge\ -\sum_{\tau\in\plE}\bigl(a_\tau+\plr\,a'_\tau\bigr)=-\sum_{
\mathfrak l\in\Lambda}A_{\mathfrak l}-\plr\sum_{\tau\in\plE}a'_\tau ,
\]
which with the bound for the jets is \eqref{pl:eq-entry}.
\end{proof}

\begin{lemma}\label{pl:arith}
For $0\le j<m$: $\ \#\Lambda\cdot\nu_p+\mathrm L_j\le(2\paloss_p-\augain_p)/m$.
\end{lemma}
\begin{proof}
Recall that $p$ is inert in $K$, being exceptional, so $\paI=I_p$ is the set of odd positive
integers. There is no pair $(\tau,e)$ in $C_p(\plE)$: if $\tau\in\auPa\setminus\plE$ then $\tau\in\aubar
\plE$, and $u(-p^e)\tau\in\plE$ by Lemma~\ref{pl:balance}. So Lemma~\ref{pl:gain} gives
$D_j(\plE)\ge\augain_p/m+\#(\auPa\cap\plE)a/x_1-\paloss_p/m$ (here $\epsilon_m(p)=1$). It remains
to show
\[
\#\Lambda\,\nu_p+\plr\sum_{\tau\in\plE}a'_\tau\ \le\ \frac{\paloss_p}m+\#(\auPa\cap\plE)\frac
a{x_1}.
\]
By \eqref{pa:eq-diag} and Lemma~\ref{pl:balance}, $a'_\tau\le a/x_1+s$ for $\tau\in\plE$,
$s=x^{60}$. If $p\le m$, then $\plr(a/x_1+s)\le a/(pm)$, because $\plr\le2\log m$, $pm\le4x$ and
$x_1\ge x^{10^6}$. If $p>m$, then $\plr=1$, so the order of $-p$ modulo $q_0$ is $1$, all places of
$\plE$ are $a$-places, and $s\le a/(pm)$. In both cases $\plr\sum_{\tau\in\plE}a'_\tau\le
fa/(pm)+\#(\auPa\cap\plE)a/x_1$. Finally $\#\Lambda\,\nu_p\le f^2\plr(n_m+\plr)\le fa/(pm)$ and
$2fa/(pm)\le fa(\log b)^2/(pm)=\paloss_p/m$, by a comparison of sizes ($p<x_1$, $m\le2\sqrt x$,
$f\le x^{12}$, $n_m\le40x^6$, $a\ge x_1^{10^6}$).
\end{proof}

\begin{proof}[Proof of Lemma~\ref{pa:hecke}, given Lemma~\ref{pa:kernel}]
Let $\mathcal V_j$ be the space of the data of the $j$-th lift fixed by $\omega(1,K'_{\mathfrak
l})$ for all $\mathfrak l\in\Lambda$. For $\mathfrak l\in\Lambda$ let $c^{(j+1,\mathfrak l)}_i$
be the numbers of Corollary~\ref{pl:chainT} for $e=j+1$, and on $\mathcal V_j$ put
\[
t_{j,\mathfrak l}=p^{\,A_{\mathfrak l}-\nu_p}\sum_{i\le\plr}c^{(j+1,\mathfrak l)}_i\,T'_{\mathfrak
l,i},\qquad t_j=\prod_{\mathfrak l\in\Lambda}t_{j,\mathfrak l},
\]
with $T'_{\mathfrak l,i}$ as in \S\ref{pl:ss-transfer}; these operators preserve $\mathcal V_j$
and commute.

\emph{{\rm(i)} $\Theta^X_j(t_j\phi^\circ_{j,x_j})\,C_{x,j}\,y=C_{x,j+1}(Y_\Lambda y)$ for
$y\in\auS$ and $x$ in the support of $\Psi'$.} For $\Phi\in\mathcal V_j$ and any $\psi$,
$\Theta^X_j(T'_{\mathfrak l,i}\Phi)\psi=T_{\mathfrak l,i}\Theta^X_j(\Phi)\psi$ by
\eqref{pl:eq-equiv}. So if $\Theta^X_j(\Phi)C_{x,j}y=C_{x,j+1}(y')$ with $y'\in\auS$, then by
Corollary~\ref{pl:chainT} ($y'$ is fixed by $\auGL_k(\auO_{\mathfrak l})\subseteq J_X$)
\[
\Theta^X_j(t_{j,\mathfrak l}\Phi)\,C_{x,j}\,y=p^{\,A_{\mathfrak l}-\nu_p}\sum_ic_i\,T_{\mathfrak
l,i}\,C_{x,j+1}(y')=C_{x,j+1}\bigl(p^{A_{\mathfrak l}}T_{\mathfrak l,\plr}\,y'\bigr)=C_{x,j+1}
(Y_{\mathfrak l}\,y').
\]
Start from $\Phi=\phi^\circ_{j,x_j}$, $y'=y$, and apply this for the primes of $\Lambda$ in turn;
the operators $Y_{\mathfrak l}$ preserve $\auS$.

\emph{{\rm(ii)}} By (i) and induction on $j$, $\Theta^X_{j-1}(t_{j-1}\phi^\circ_{j-1,x_{j-1}})
\cdots\Theta^X_0(t_0\phi^\circ_{0,x_0})\,\varphi=C_{x,j}(Y_\Lambda^{\,j}\varphi)$. Hence
\[
C_{\vartheta,\Psi'}(Y_\Lambda^{\,m}\varphi)=\sum_x\Psi'(x)\,\Theta^X_{m-1}(t_{m-1}\phi^\circ_{m-1,
x_{m-1}})\cdots\Theta^X_0(t_0\phi^\circ_{0,x_0})\,\varphi .
\]
\emph{{\rm(iii)} The tensor $\sum_x\Psi'(x)\bigotimes_jt_j\phi^\circ_{j,x_j}$ is invariant under
the levels.} The tensor $\sum_x\Psi'(x)\bigotimes_j\phi^\circ_{j,x_j}$ is (proof of
Proposition~\ref{au:expansion}) and lies in $\bigotimes_j\mathcal V_j$ by \eqref{pl:eq-fix}. For
$b=(b_j,b_{j+1})$ in the levels, $\omega(b)$ preserves $\mathcal V_j$ and $\omega(b)T'_{\mathfrak
l,i}=T'_{\mathfrak l,i}\omega(b)$ on $\mathcal V_j$: the part of $b_{j+1}$ away from $\mathfrak l$
commutes with the elements $h$, and conjugation by its component at $\mathfrak l$, which lies in
$K'_{\mathfrak l}$, permutes the cosets $hK'_{\mathfrak l}$ of the double coset.

\emph{{\rm(iv)} Conclusion.} By (iii), Lemma~\ref{pl:sum} applies. The datum
$t_j\phi^\circ_{j,x_j}$ is $p^{\sum_{\mathfrak l}A_{\mathfrak l}-\#\Lambda\nu_p}$ times a
combination with coefficients in $\pao$ of data $\omega(1,h)\phi^\circ_{j,x_j}$ with $h$ as in
Lemma~\ref{pl:neigh}. By that lemma and the ultrametric inequality,
\[
\plgb{g'}{\theta_j(t_j\phi^\circ_{j,x_j};g,g')\xi}\le p^{\,\#\Lambda\nu_p+\mathrm L_j}\,\plgb g\xi .
\]
The numbers $\Psi'(x)$ are integral at $p$. As in the proof of Lemma~\ref{pa:chain}, the formula of
Lemma~\ref{pl:sum} and the restriction to $Y_m$ give, for $\varphi\in\auS_\iota$, that
$B_{\vartheta,\Psi}(Y_\Lambda^{\,m}\varphi)$ has gauge at most $p^{\sum_{j<m}(\#\Lambda\nu_p+\mathrm L_j)}
\le p^{\,2\paloss_p-\augain_p}$, by Lemma~\ref{pl:arith}.
\end{proof}

\subsubsection{The level prime}\label{pl:ss-level}
Let $\ell=N$ and $\iota\colon\C\to\C_N$. Let $\plP$ be the set of the $f$ primes of $E$
dividing $\pi_N$; the primes above $N$ are those of $\plP$ and of $\aubar\plP$.

\emph{The places $\plSi$.} The primes $\mathfrak l_\tau$ $(\tau\in\auPl)$ lie above one prime
$\mathfrak n$ of $K$, and $\mathfrak n\ne\aubar{\mathfrak n}$; as $N$ splits completely, $\tau
\mapsto\mathfrak l_\tau$ is a bijection of $\auPl$ onto the primes above $\mathfrak n$. Let
\[
\plSi=\{w_\tau=\aubar{\mathfrak l_\tau}:\tau\in\auPl\},
\]
the primes above $\aubar{\mathfrak n}$: it contains one prime of each pair $w,\aubar w$ above $N$;
for $w\in\plSi$ we write $\tau_w$ for the place with $w_{\tau_w}=w$.
Put $E_\Sigma=\prod_{w\in\plSi}E_w=\Q_N^{\plSi}$, $\auO_\Sigma=\Z_N^{\plSi}$, $\auF=\auO_\Sigma/N=
\auF_N^{\plSi}$. By (N2) the component at $\plSi$ identifies $\auU(X_j)(E^+\otimes\Q_N)$ with
$G_j=\auGL_{n_j}(E_\Sigma)$, and $g\in G_j$ acts on $W_{\lambda(X_j)}\otimes\C_N$ by $\bigotimes_\tau
\rho_\tau(g_{w_\tau})$. We write $e$ for the last basis vector, of any $X_j$.

\emph{Lattices.} For $S\subseteq\plSi$ let $\plD_S\in G_j$ be $\operatorname{diag}(1,\dots,1,N)$ at
the places of $S$ and $1$ at the others, and
\[
\plD=\plD_{\plSi},\qquad\pld_j=\plD_{\plSi\cap\plP},\qquad\ple_j=\plD_{\plSi\setminus\plP},\qquad
\plD=\pld_j\ple_j .
\]
The \emph{lattice of $X_j$ at $N$} is $\mathcal L_j=\pld_j\auO_\Sigma^{\,n_j}$, the integral vectors whose
last coordinate is divisible by $\pi_N$ (at $w\in\plSi\cap\plP$ the number $N$ is a unit times
$\pi_N$, at the other $w$ the number $\pi_N$ is a unit); it is \emph{marked} by the class
$\varepsilon_j$ of $\pld_je$ in $\mathcal L_j/N\mathcal L_j$; and $\mathcal L_j^+=\pld_j\plD^{-1}\auO_\Sigma^{\,n_j}=\mathcal L_j+
N^{-1}\auO_\Sigma\,\pld_je$. Let
\[
\plMir=\{k\in\auGL_{n_j}(\auO_\Sigma):k\,e\equiv e\pmod N\},
\]
the stabiliser of the standard lattice marked by $e$.

\begin{lemma}\label{pl:mir}
The group of the components at $\plSi$ of the elements of $J_{X_j}$ is $\pld_j\,\plMir\,
\pld_j^{-1}$, the stabiliser of $(\mathcal L_j,\varepsilon_j)$. For $k\in\plMir$, $\plD k\plD^{-1}$ is
integral with integral inverse and has last row $\equiv{}^te$ modulo $N$.
\end{lemma}
\begin{proof}
Above $N$, $J_{X_j}$ is defined by: $g$ integral with integral inverse, and the last row of $g_u$
is $\equiv{}^te$ modulo $u$ for $u\in\plP$. For $w\in\plSi\cap\plP$ this is a condition on $g_w$.
For $w\in\plSi\setminus\plP$ we have $\aubar w\in\plP$ and $g_{\aubar w}={}^t\aubar{g_w}^{-1}$, so
the condition is that the last column of $g_w^{-1}$, equivalently of $g_w$, is $\equiv e$ modulo
$w$; there $\pld_j=1$. For $w\in\plSi\cap\plP$, conjugation
$k\mapsto\operatorname{diag}(1,\dots,1,N)\,k\,\operatorname{diag}(1,\dots,1,N)^{-1}$
multiplies the last row by $N$ and divides the last column by $N$, outside the corner: it
exchanges ``integral with integral inverse, last row $\equiv{}^te$'' and ``integral with integral
inverse, last column $\equiv e$''. The last claim is the same computation.
\end{proof}

\emph{The representations $\pi_j$.} Let $\mathcal A(X_j)$ be the space of smooth forms on $X_j$ of
weight $\lambda(X_j)$. For $a\in G_j$, read as the element of $\auU(X_j)(\auA_{E^+,f})$ supported
above $N$ with these components at $\plSi$, put $\pi_j(a)=R_{\pld_ja\pld_j^{-1}}$ on $\mathcal
A(X_j)$ (``$a$ is read in a basis of $\mathcal L_j$''). By Lemma~\ref{pl:mir} the forms of level $J_{X_j}$
are those fixed by $\pi_j(\plMir)$ and by the part of $J_{X_j}$ away from $N$. A vector $\varphi$
of $\pi_j$ is \emph{ramified} if for each $w\in\plSi$ there is $x\in\auO_\Sigma^\times$ with
$x_{w'}\equiv1\pmod N$ for $w'\ne w$ and $\varphi\in(1-\pi_j(x\cdot1))\mathcal A(X_j)$. The forms
of $\auSram$ are ramified: for $w\in\plSi$ one of $w,\aubar w$ is a prime $u\in\plP$, and
$\varphi\in(1-z_u)\auS$, where $z_u$ is a scalar whose component at $w$ generates $\auF_N^\times$
and which is $1$ at the other places of $\plSi$.

\emph{The character $\psi_N$ and the operator $U_N$.} Let $\psi_N\colon\auF\to\C^\times$,
$\psi_N(c)=\lambda(-\tilde c/N)$ for a lift $\tilde c\in\auO_\Sigma$, with $\lambda$ as in
\eqref{pl:eq-heis} for $\mathsf S=\plSi$; it is a character, non-trivial on each factor $\auF_N$.
For $y=(y',u)\in\auF^{n_j}$ with last entry $u\in\auF^\times$ let $k_y\in\auGL_{n_j}(\auO_\Sigma)$
be a lift of
\[
\begin{pmatrix}1+u^{-1}y'\,{}^ty'&y'\\{}^ty'&u\end{pmatrix}=\begin{pmatrix}1&y'\\0&u\end{pmatrix}
\begin{pmatrix}1&0\\u^{-1}\,{}^ty'&1\end{pmatrix},
\]
whose last column and last row are both $y$. The operator on $\mathcal A(X_j)$
\begin{equation}\label{pl:eq-U}
U=U^{(j)}=\pi_j(\plD^{-1})\sum_{y\in\auF^{n_j},\ y_{\mathrm{last}}\in\auF^\times}\psi_N(y_{
\mathrm{last}})\,\pi_j(k_y),\qquad U_N:=U^{(0)},
\end{equation}
replaces $\mathcal L_j$ by $\mathcal L_j^+$ and sums against $\psi_N$ over the vectors of $\mathcal L_j/N\mathcal L_j$ with unit
coordinate along $\varepsilon_j$. On vectors fixed by $\pi_j(\plMir)$ the term of $y$ depends only
on $k_y\plMir$, that is, on $y$.

\begin{lemma}\label{pl:Uprop}
{\rm(a)} $U$ commutes with $\pi_j(x\cdot1)$, $x\in\auO_\Sigma^\times$, and with the right
translations by elements that are trivial above $N$; it preserves the ramified vectors.
{\rm(b)} $U$ preserves the vectors fixed by $\pi_j(\plMir)$. In particular $U_N\auS\subseteq\auS$
and $U_N\auSram\subseteq\auSram$.
\end{lemma}
\begin{proof}
(a) is clear. (b) Let $\varphi$ be fixed by $\plMir$ and $g\in\plMir$. By Lemma~\ref{pl:mir},
$k=\plD g\plD^{-1}$ is integral with last row $\equiv{}^te$. Then $kk_y$ has last column $ky$, and
$(ky)_{\mathrm{last}}=y_{\mathrm{last}}$; so $kk_y\in k_{ky}\plMir$ and
\[
\pi_j(g)U\varphi=\pi_j(\plD^{-1})\sum_y\psi_N(y_{\mathrm{last}})\pi_j(kk_y)\varphi=\pi_j(\plD^{-1})
\sum_y\psi_N((ky)_{\mathrm{last}})\pi_j(k_{ky})\varphi=U\varphi .
\]
For $\auSram$ use (a): $U_N$ commutes with the $z_u$.
\end{proof}

\emph{Data at $\plSi$ and the lift law.} We use (N5) for $\mathsf S=\plSi$. Let $x$ be a
translation tuple vanishing above $N$; then $\phi^\circ_{j,x_j}=\plone\otimes v_{x,j}$. For a
function $\Psi$ on $M_{n_{j+1}\times n_j}(E_\Sigma)$ put
\[
L_j(\Psi)=L^x_j(\Psi)=\Theta^X_j\bigl(\omega_{\mu_{j+1}}(\pld_j,\pld_{j+1})(\Psi\otimes v_{x,j})
\bigr)\colon\mathcal A(X_j)\to\mathcal A(X_{j+1}).
\]
We consider only functions $\Psi$ supported on the integral matrices and depending on $T$ modulo
$N$; with $e,e'$ the last basis vectors of $X_j$, $X_{j+1}$, and $T$ integral:
\[
\plcol_y(T)=[\,Te\equiv y\,],\quad\plrow(T)=[\,{}^te'T\equiv{}^te\,],\quad\psi_y(T)=\psi_N({}^te'Ty),
\quad\pldir=\plcol_{e'},
\]
for $y\in\auF^{n_{j+1}}$, respectively $y\in\auF^{n_j}$ in $\psi_y$. Thus $\pldir$, the
\emph{direct datum}, is the indicator of the integral $T$ with $Te\equiv e'$; read in bases of
$\mathcal L_j$ and $\mathcal L_{j+1}$, as in the definition of $L_j(\Psi)$, it is the indicator of the maps
$(\mathcal L_j,\varepsilon_j)\to(\mathcal L_{j+1},\varepsilon_{j+1})$ of marked lattices.

\begin{lemma}[Lift law]\label{pl:law}
For $g\in G_j$, $h\in G_{j+1}$ and $\Psi$, with $\chi_r$ as in \eqref{pl:eq-w5} for $\plSi$,
\[
L_j\bigl(\Psi(h^{-1}\,\cdot\,g)\bigr)=\frac{\chi_{n_{j+1}}(g)}{\chi_{n_j}(h)}\ \pi_{j+1}(h)\circ
L_j(\Psi)\circ\pi_j(g)^{-1}.
\]
The characters $\chi_r$ are trivial on integral elements. In particular, with $\gamma_j=\chi_{n_{
j+1}}(\plD)/\chi_{n_j}(\plD)=N^{f(t+1)/2}\in\Z$ (the first $\plD$ in $G_j$, the second in
$G_{j+1}$),
\begin{equation}\label{pl:eq-rowdir}
\pi_{j+1}(\plD^{-1})\,L_j(\plrow)=\gamma_j\,L_j(\pldir)\,\pi_j(\plD^{-1}).
\end{equation}
\end{lemma}
\begin{proof}
By \eqref{pl:eq-w5}, $\Psi(h^{-1}Tg)\otimes v=\chi_{n_{j+1}}(g)\chi_{n_j}(h)^{-1}\omega(g,h)(\Psi
\otimes v)$, and
\[
\omega(\pld_j,\pld_{j+1})\,\omega(g,h)=\omega\bigl(\pld_jg\pld_j^{-1},\pld_{j+1}h\pld_{j+1}^{-1}
\bigr)\,\omega(\pld_j,\pld_{j+1});
\]
conclude by \eqref{pl:eq-equiv}. The characters $\chi_{j+1,
w}$ are unramified at $w\mid N$. For \eqref{pl:eq-rowdir} take $g=\plD^{-1}$, $h=\plD^{-1}$:
$\plrow(\plD T\plD^{-1})=\pldir(T)$, as in the proof of Lemma~\ref{pl:mir}; and $\chi_r(\plD)=
\prod_w\chi_{j+1,w}(N)\cdot N^{fr/2}$, with $t+1$ even.
\end{proof}

\begin{lemma}[Calculus of direct data]\label{pl:calc}
Let $\varphi\in\mathcal A(X_j)$ be ramified. Write $L=L_j$, $L'=L_{j+1}$, $\pi=\pi_j$, etc.
\begin{enumerate}
\item[(a)] If a function $\mathsf f$ of $T$ modulo $N$ satisfies $\mathsf f(cT)=\mathsf f(T)$
for all $c\in\auF^\times$ that are $1$ outside one place $w$, then $L(\mathsf f)\varphi=0$. For
every $\mathsf f$, $L(\mathsf f)\varphi$ is ramified.
\item[(b)] $\pi'(\plD^{-1})L(\psi_e)\varphi=U'\bigl(L(\pldir)\varphi\bigr)=\gamma_j\,L(\pldir)
(U\varphi)$.
\item[(c)] $\sum_{y\in\auF^{n_{j+1}}}L'(\psi_y)L(\plcol_y)\varphi=\#\mathrm{Prim}\cdot L'(\psi_{e'})
L(\pldir)\varphi$, where $\mathrm{Prim}$ is the set of the $y$ that are non-zero at every place.
\item[(d)] $\pi''(\plD^{-1})\sum_yL'(\psi_y)L(\plcol_y)\varphi=\#\mathrm{Prim}\cdot\gamma_{j+1}
\gamma_j\cdot L'(\pldir)L(\pldir)(U\varphi)$.
\end{enumerate}
\end{lemma}
\begin{proof}
(a) Let $\varphi=\psi-\pi(x)\psi$ with $x$ as in the definition, for the place $w$. By
Lemma~\ref{pl:law} with $g=x^{-1}\cdot1$, $L(\mathsf f(\,\cdot\,x^{-1}))=L(\mathsf f)\pi(x)$, and
$\mathsf f(Tx^{-1})=\mathsf f(T)$; so $L(\mathsf f)\pi(x)\psi=L(\mathsf f)\psi$. For the second
claim, Lemma~\ref{pl:law} with $g=x$, $h=x$ gives $L(\mathsf f)\pi(x)=\pi'(x)L(\mathsf f)$.

(b) Sort $\psi_e$ by the corner entry $c$ of $T$: $\psi_e=\sum_{c\in\auF}\psi_N(c)[T_{\mathrm{
last,last}}\equiv c]$. If $c$ vanishes at a place $w$, its term is invariant as in (a) and
contributes nothing. Sorting the remaining terms by the last column,
\[
L(\psi_e)\varphi=\sum_{y_{\mathrm{last}}\in\auF^\times}\psi_N(y_{\mathrm{last}})\,L(\plcol_y)
\varphi,\qquad\plcol_y(T)=\pldir(k_y^{-1}T),\qquad L(\plcol_y)=\pi'(k_y)L(\pldir),
\]
by Lemma~\ref{pl:law}; applying $\pi'(\plD^{-1})$ gives $U'(L(\pldir)\varphi)$. Sorting instead by
the last row, $L(\psi_e)\varphi=\sum_y\psi_N(y_{\mathrm{last}})L([{}^te'T\equiv{}^ty])\varphi$, and
$[{}^te'T\equiv{}^ty]=\plrow(Tk_y^{-1})$ because the last row of $k_y$ is ${}^ty$; so
$L([{}^te'T\equiv{}^ty])=L(\plrow)\pi(k_y)$ and, by \eqref{pl:eq-rowdir},
\[
\pi'(\plD^{-1})L(\psi_e)\varphi=\gamma_j\,L(\pldir)\,\pi(\plD^{-1})\sum_y\psi_N(y_{\mathrm{last}})
\pi(k_y)\varphi=\gamma_jL(\pldir)(U\varphi).
\]
(c) If $y$ vanishes at a place, $\plcol_y$ is invariant as in (a) and its term is zero. If
$y\in\mathrm{Prim}$, there is $k\in\auGL_{n_{j+1}}(\auO_\Sigma)$ with $ke'\equiv y$; then
$\plcol_y(T)=\pldir(k^{-1}T)$ and $\psi_y(R)=\psi_{e'}(Rk)$, so $L(\plcol_y)=\pi'(k)L(\pldir)$ and
$L'(\psi_y)=L'(\psi_{e'})\pi'(k)^{-1}$ by Lemma~\ref{pl:law}. (d) By (a), $L(\pldir)\varphi$ is
ramified; apply (b) to $L'$ and then (b) to $L$.
\end{proof}

\emph{The gluing.} The gluing of \S\ref{ps:ss-glue} is $\psGlue_i=\sum_{\mathsf y}\psrho(p_{\mathsf
y})\otimes\psrho(p'_{\mathsf y})$, with the Heisenberg operators $\psrho$ themselves (the points
$p_{\mathsf y}$, $p'_{\mathsf y}$ are supported on $\plP$, respectively $\aubar\plP$, where
$\psrho$ is given by \eqref{pl:eq-heis} without any choice); on data that are the vacuum above
$N$ the term of $\mathsf y$ depends only on $\mathsf y$ modulo $N$. In the coordinates $T$, the
point $p_{\mathsf y}$ has last column $-\mathsf y/N$ at the places of $\plP$, and $p'_{\mathsf y}$
has last row $-\mathsf y/N$ at the places of $\aubar\plP$; all other entries vanish.

\begin{lemma}[The glued double lift]\label{pl:glue}
Let $x$ vanish above $N$ and $0\le i<m/2$. Then
\[
\Theta^X_{2i+1}\Theta^X_{2i}\bigl(\psGlue_i(\phi^\circ_{2i,x_{2i}}\otimes\phi^\circ_{2i+1,x_{2i+1}})
\bigr)=\mathsf g_i\cdot\pi_{2i+2}(\plD^{-1})\circ\sum_{z\in\auF^{n_{2i+1}}}L^x_{2i+1}(\psi_z)\circ
L^x_{2i}(\plcol_z),
\]
where $\mathsf g_i=\chi^{(2i+2)}_{n_{2i+1}}(\ple_{2i+2})\cdot\chi^{(2i+1)}_{n_{2i+1}}(\pld_{2i})$,
the upper index being that of the character $\mu$.
\end{lemma}
\begin{proof}
We compute $\sum_{\mathsf y}\psrho(p_{\mathsf y})\plone\otimes\psrho(p'_{\mathsf y})\plone$ as a
function of $(T,R)$, at $\plSi$, by \eqref{pl:eq-heis}. At $w\in\plSi\cap\plP$: $\psrho(p_{\mathsf
y})\plone(T)=\plone(T+p_{\mathsf y})$ is $1$ if and only if $T\pld_{2i}$ is integral with last
column $\equiv\mathsf y$; and $\psrho(p'_{\mathsf y})$ multiplies $\plone(R)$ by $\lambda_w(-{}^te
R\mathsf y/N)=\psi_N({}^teR\mathsf y)$. At $w\in\plSi\setminus\plP$: $\aubar w\in\plP$, so
$\psrho(p_{\mathsf y})$ multiplies $\plone(T)$ by $\psi_N({}^t\mathsf yTe)$, and $\psrho(p'_{
\mathsf y})\plone(R)=\plone(R+p'_{\mathsf y})$ is $1$ if and only if $\ple_{2i+2}R$ is integral
with last row $\equiv{}^t\mathsf y$; the sum over $\mathsf y$ at this place is $\psi_N$ of the
corner entry of $(\ple_{2i+2}R)T$. In both cases the sum is
\[
\sum_{z\in\auF^{n_{2i+1}}}\plcol_z(T\pld_{2i})\,\psi_z(\ple_{2i+2}R).
\]
By \eqref{pl:eq-w5} and \eqref{pl:eq-equiv}, as in the proof of Lemma~\ref{pl:law},
\begin{align*}
\Theta^X_{2i}\bigl(\plcol_z(\,\cdot\,\pld_{2i})\otimes v\bigr)&=\chi^{(2i+1)}_{n_{2i+1}}(\pld_{2i})\,
R_{\pld_{2i+1}}^{-1}\,L_{2i}(\plcol_z),\\
\Theta^X_{2i+1}\bigl(\psi_z(\ple_{2i+2}\,\cdot\,)\otimes v'\bigr)&=\chi^{(2i+2)}_{n_{2i+1}}
(\ple_{2i+2})\,R_{\plD^{-1}}\,L_{2i+1}(\psi_z)\,R_{\pld_{2i+1}},
\end{align*}
because $\ple_{2i+2}\pld_{2i+2}=\plD$; and $R_{\plD^{-1}}=\pi_{2i+2}(\plD^{-1})$.
\end{proof}

\begin{proposition}[All lifts become direct]\label{pl:direct}
There is $c\in\C^\times$, algebraic, with $|\iota(c)|\le N^{mfb}$, such that for every $x$
vanishing above $N$ and every $\varphi\in\auSram$,
\[
C_{\psGlue(\phi^\circ_x)}\varphi=c\cdot L^x_{m-1}(\pldir)\circ\cdots\circ L^x_0(\pldir)\bigl(U_N^{\,
m/2}\varphi\bigr).
\]
\end{proposition}
\begin{proof}
Write $\mathsf L_j=L^x_j(\pldir)$. By Lemmas~\ref{pl:glue} and~\ref{pl:calc}(d), on ramified vectors of
$\pi_{2i}$ the $i$-th glued double lift is
\[
\mathsf g_i\cdot\#\mathrm{Prim}_i\cdot\gamma_{2i+1}\gamma_{2i}\cdot\mathsf L_{2i+1}\mathsf L_{2i}\,
U^{(2i)} .
\]
By Lemma~\ref{pl:calc}(a),(b), for ramified $\psi$ the vector $U^{(2i)}\mathsf L_{2i-1}\cdots\mathsf L_0
\psi$ is $\prod_{l<2i}\gamma_l$ times $\mathsf L_{2i-1}\cdots\mathsf L_0U^{(0)}\psi$; and $U^{(0)}$
preserves the ramified vectors. Induction on the number of double lifts gives the formula with
\[
c=\prod_{i<m/2}\Bigl(\mathsf g_i\cdot\#\mathrm{Prim}_i\cdot\prod_{l<2i+2}\gamma_l\Bigr).
\]
The $\gamma_l$ and $\#\mathrm{Prim}_i$ are integers. Let $w\in\plSi$; then $\mathfrak l_\tau=\aubar
w$ for $\tau=\tau_w$ only, and $\mathfrak l_\tau=w$ for no $\tau$, so by Lemma~\ref{pl:char}
$|\iota\chi^{(l)}_w(N)|\le N^{\kappa_l(\tau_w)}\le N^{2b}$ for $l\le m$. Each $w\in\plSi$ occurs
in exactly one of the two factors of $\mathsf g_i$, and the factors $|N|_w^{-n_{2i+1}/2}=N^{n_{2i+
1}/2}$ are integers. So $|\iota(\mathsf g_i)|\le N^{2fb}$ and $|\iota(c)|\le N^{2fb\cdot m/2}$.
The constant $c$ does not depend on $x$; it depends on the twist $\vartheta$, through the numbers
$\chi^{(l)}_w(N)$ in the $\mathsf g_i$ only (a twist is unramified above $N$ and of finite order, so
it changes these numbers by roots of unity), and the bound for $|\iota(c)|$, which is all that is
used below, is the same for all twists.
\end{proof}

\begin{lemma}[The kernels at $N$]\label{pl:kernelN}
Let $0\le j<m$ and let $z\in\auHom(X_j,X_{j+1})(\auA_{E^+,f})$ be integral at the places of
$\plSi$ and such that $Nz$ is integral at the places of $\aubar{\plSi}$ $($the primes $\mathfrak
l_\tau)$. Then for all $g,g'$ the matrix of $\rho(g'_N)^{-1}\circ\theta_j(\phi^\circ_{j,z};g,g')
\circ\rho(g_N)$ in the bases $(e_T)$ has entries in $\auO_\iota$.
\end{lemma}
\begin{proof} This is Corollary~\ref{pf:K1}(b): with $\delta\in\mathfrak n\setminus(\mathfrak n^2\cup\aubar{\mathfrak n})$,
$\mathfrak n$ the prime of $K$ induced by $\iota$, the datum $\phi^\circ_{j,z}$ is a standard datum of level $\delta$ at $N$;
$\iota(\aubar\delta)$ is a unit, and Theorem~\ref{pf:std} applies with $c=\iota(\pi_M/\aubar\delta)$.\end{proof}

\begin{lemma}[The prime $N$; closed form of Lemma~\ref{pa:level}]\label{pl:level}
Let $\ell=N$, $\iota\colon\C\to\C_N$, and let $\plSi$, $\mathcal L_0$, $\pld_0$ and $U_N$ be as above. Say
that $\psi\in\auS$ is \emph{integral for $\mathcal L_0$} if it has gauge at most $1$ at the base point
$\pld_0$, that is, $\iota(\psi(g))\in\rho(g_N\pld_0)W_{\lambda^X}(\auO_\iota)$ for all $g$. Let
$\mathcal A'=\{\varphi\in\auS:U_N^{\,m/2}\varphi\text{ is integral for }\mathcal L_0\}$. Then:
\begin{enumerate}
\item[(a)] for every twist $\vartheta$, every chain datum $\Psi$
and every $\varphi\in\mathcal A'\cap\auSram$, the form $B_{\vartheta,\Psi}\varphi$ has gauge at
most $N^{2mfb}$;
\item[(b)] $\aucol(\auS_\iota,\mathcal A'\cap\auS_\iota)\mid N^{\,i}$ for an integer $i\le
m\,f\,b\,\auDX$.
\end{enumerate}
\end{lemma}
\begin{proof}[Proof of {\rm(a)}]
\emph{Step 1.} By (C1) every $x$ in the support of $\Psi$ vanishes above $N$. Let $\psi=U_N^{\,
m/2}\varphi$; it lies in $\auS$ (Lemma~\ref{pl:Uprop}) and is integral for $\mathcal L_0$. By
Proposition~\ref{pl:direct},
\[
C_{\psGlue(\sum_x\Psi(x)\phi^\circ_x)}\varphi=c\sum_x\Psi(x)\,\Theta^X_{m-1}(A_{m-1}\phi^\circ_{m-1,
x_{m-1}})\cdots\Theta^X_0(A_0\phi^\circ_{0,x_0})\,\psi,\qquad A_j(\plone\otimes v)=\omega(\pld_j,
\pld_{j+1})(\pldir\otimes v).
\]
\emph{Step 2: the tensor $\sum_x\Psi(x)\bigotimes_jA_j\phi^\circ_{j,x_j}$ is invariant under the
levels.} First, $\sum_x\Psi(x)\bigotimes_j\phi^\circ_{j,x_j}$ is: let $\mathrm{pr}_j$ be the
projection of $\plH_f$ onto the vectors that are the vacuum above $N$ (the mean of $\psrho^\circ
(l)$ over the integral $l$ supported above $N$). It commutes with $\omega(b_j,b_{j+1})$ for $b$ in
the levels, which are integral; and $\bigotimes_j\mathrm{pr}_j$ maps $\psGlue(\phi^\circ_x)$ to
$\phi^\circ_x$, because $\mathrm{pr}(\psrho(p_{\mathsf y})\phi^\circ)=0$ unless $p_{\mathsf y}$ is
integral, that is, $\mathsf y\equiv0$. So the claim follows from (C4). Secondly, for $b_j\in
J_{X_j}$, $b_{j+1}\in J_{X_{j+1}}$ the operator $\omega(b_j,b_{j+1})$ commutes with $A_j$ on the
vectors that are the vacuum above $N$: by Lemma~\ref{pl:mir} the components at $\plSi$ of $\pld_j^{
-1}b_j\pld_j$ and $\pld_{j+1}^{-1}b_{j+1}\pld_{j+1}$ lie in $\plMir$, and $\pldir(k'^{-1}Tk)=
\pldir(T)$ for $k,k'\in\plMir$; use \eqref{pl:eq-w5}.

\emph{Step 3: the direct datum loses at most $N^{fn_{j+1}}$.} We claim that for $x$ vanishing
above $N$ and all $g,g'$ the matrix of $\rho(g'_N)^{-1}\theta_j(\pldir\otimes v_{x,j};g,g')\rho
(g_N)$ has entries of absolute value at most $N^{fn_{j+1}}$. Indeed, by Fourier inversion on
$\auF^{n_{j+1}}$ and \eqref{pl:eq-heis},
\[
\pldir(T)=N^{-fn_{j+1}}\sum_{\mathsf y\in\auF^{n_{j+1}}}\psi_N(-{}^t\mathsf ye')\,\psi_N({}^t\mathsf
yTe)\,\plone(T),\qquad\psi_N({}^t\mathsf yTe)\,\plone=\psrho(a_{\mathsf y})\plone,
\]
where $a_{\mathsf y}$ is supported on $\aubar{\plSi}$ and has last column $-\mathsf y/N$. So
$\pldir\otimes v_{x,j}$ is $N^{-fn_{j+1}}$ times a sum of roots of unity times data $\phi^\circ_{j,
x_j+a_{\mathsf y}}$, to which Lemma~\ref{pl:kernelN} applies.

\emph{Step 4.} By (T2), $\theta_j(A_j\phi^\circ_{j,x_j};g,g')=\theta_j(\pldir\otimes v_{x,j};g
\pld_j,g'\pld_{j+1})$; so by Step~3 this kernel maps vectors of gauge $C$ at $g\pld_j$ to vectors
of gauge at most $N^{fn_{j+1}}C$ at $g'\pld_{j+1}$. By Step~2 and Lemma~\ref{pl:sum}, and since the
$\Psi(x)$ are integral at $N$ and $\psi$ has gauge at most $1$ at $\pld_0$, the form
$F=\sum_x\Psi(x)\Theta^X_{m-1}(A_{m-1}\phi^\circ)\cdots\Theta^X_0(A_0\phi^\circ)\psi$ has gauge at
most $N^{f\sum_jn_{j+1}}\le N^{mfb}$ at the base point $\pld_m$ (comparison of sizes: $n_m\le b$).

\emph{Step 5: restriction.} $(\operatorname{res}F)(y)=\prod_\tau\tau(\aubar{\pi_N})^{-n_\tau}
\mathrm{br}(F(\operatorname{diag}(y,1)))$. The element $(\operatorname{diag}(y,1)\pld_m)^{-1}$ is
$\operatorname{diag}(y^{-1},s)$ with $s_\tau=N^{-1}$ if $w_\tau\in\plP$ and $s_\tau=1$ otherwise,
so by Lemma~\ref{ps:branching}(a),(c)
\[
\Bigl\|\prod_\tau s_\tau^{\,n_\tau}\cdot\rho^\circ(y_N)^{-1}\mathrm{br}\,F(\operatorname{diag}(y,1))
\Bigr\|=\bigl\|\mathrm{br}\,\rho(\operatorname{diag}(y,1)\pld_m)^{-1}F(\operatorname{diag}(y,1))
\bigr\|\le N^{mfb}.
\]
On the other hand $|\iota\tau(\aubar{\pi_N})|=N^{-1}$ if $\pi_N\in\aubar{\mathfrak l_\tau}=w_\tau$,
that is, if $w_\tau\in\plP$, and $=1$ otherwise; so $|\iota\tau(\aubar{\pi_N})^{-n_\tau}|=
|s_\tau^{\,n_\tau}|$ for every $\tau$, and $\operatorname{res}F$ has gauge at most $N^{mfb}$ at
the base point $1$. With $|\iota(c)|\le N^{mfb}$ this proves (a).
\end{proof}

\begin{proof}[Proof of {\rm(b)}]
Let $\pld=\pld_0$, $\pld^+=\pld_0\plD^{-1}=\ple_0^{-1}$, and let $\mathcal A^L$, $\mathcal A^+$ be
the modules of the $\psi\in\auS$ with gauge at most $1$ at the base points $\pld$, $\pld^+$
(integral for $\mathcal L_0$, for $\mathcal L_0^+$).

\emph{{\rm(i)} $\mathcal A^L\subseteq\auS_\iota\subseteq\mathcal A^+$}, because $\rho(\pld)$ and
$\rho(\pld^+)^{-1}=\rho(\ple_0)$ have integral entries: the weight $\lambda^X$ has no power of the
determinant.

\emph{{\rm(ii)} $\aucol(\mathcal A^+,\mathcal A^L)\mid N^{\,l}$ with $l\le\auDX fa/k$.} By
Corollary~\ref{pl:wacor}(a) (for $g=g_0\pld^+$, $g_0$ any representative of the class, $d=\plD^{-1}$;
put $r=\gamma g_0$) every class has a representative $r$ such that $r\pld^+$ and $\plD^{-1}r\pld$ are integral above $N$. For $b\in J_X$,
$\pld^{-1}b\pld\in\plMir$ and $(\pld^+)^{-1}b\pld^+=\plD(\pld^{-1}b\pld)\plD^{-1}$ are integral
above $N$ (Lemma~\ref{pl:mir}). So by (G1), (G2), for $\psi\in\auS$ and $g=\gamma rb$,
\[
\plgb{g\pld^+}{\psi(g)}=\|\iota\psi(r)\|,\qquad\plgb{g\pld}{\psi(g)}=\bigl\|\rho(\plD)^{-1}\iota
\psi(r)\bigr\| .
\]
Thus evaluation at the representatives identifies $\mathcal A^+$ with $\prod_rW_{\lambda^X}
(\auO_\iota)$ and $\mathcal A^L$ with $\prod_r\rho(\plD)W_{\lambda^X}(\auO_\iota)$. By
\eqref{pl:eq-diag}, $\rho(\plD)$ is diagonal with entries $N^{\deg_e\beta}$, where $\deg_e\beta$
is the total degree of the basis vector $e_\beta$, $\beta=(\beta_\tau)_{\tau\in\auPa}$, in the
last variable. By Lemma~\ref{ph:col}(a) the colength divides $N^{\,l}$, $l=\#\{r\}\cdot\sum_\beta
\deg_e\beta$; by the symmetry of the $k$ variables the mean of $\deg_e\beta$ is $a\#\auPa/k$.

\emph{{\rm(iii)} $U_N\mathcal A^L\subseteq\mathcal A^+$.} For $\psi\in\mathcal A^L$, $U_N\psi\in
\auS$ and $(U_N\psi)(g)=\sum_y\psi_N(y_{\mathrm{last}})\psi(gg_y)$ with $g_y=\pld^+k_y\pld^{-1}$.
As $gg_y\pld=g\pld^+k_y$ and $k_y$ is integral, $\plgb{g\pld^+}{\psi(gg_y)}=\plgb{gg_y\pld}{\psi(g
g_y)}\le1$.

\emph{{\rm(iv)}} Let $\mathcal A'_i=\{\varphi\in\auS:U_N^{\,i}\varphi\in\mathcal A^L\}$. We show
$\aucol(\auS_\iota,\mathcal A'_i)\mid N^{\,l(i+1)}$ by induction. For $i=0$ this is (i), (ii). If
a family in $\auS_\iota$ has been carried into $\mathcal A'_i$ by a matrix of determinant dividing
$N^{l(i+1)}$, the image $(\varphi_k)$ has $U_N^{\,i+1}\varphi_k\in\mathcal A^+$ by (iii); by (ii)
a matrix $\mu$ with $\det\mu\mid N^{\,l}$ carries the family $(U_N^{\,i+1}\varphi_k)$ into
$\mathcal A^L$, that is, $(\varphi_k)$ into $\mathcal A'_{i+1}$. For $i=m/2$: the matrices are
integral, so the families are carried into $\mathcal A'\cap\auS_\iota$, and the exponent is
$l(m/2+1)\le\auDX fa(m/2+1)/k\le mfb\,\auDX$, because $a=xb$ and $x(m/2+1)\le mk$.
\end{proof}

\emph{Case (iii) of Theorem~\ref{pa:divis}.} With Lemma~\ref{pl:level} in place of
Lemma~\ref{pa:level}, the deduction in the proof of Theorem~\ref{pa:divis} stands: $S_0=\mathcal A'
\cap\auS_\iota$ has $\aucol(\auS_\iota,S_0)\mid N^{\,i}$, $i\le mfb\,\auDX\le\auerr_N\auDX$, the
form $B_{\vartheta,\Psi}\varphi$ has gauge at most $N^{2mfb}=N^{\auerr_N-\augain_N}$ for
$\varphi\in S_0\cap\auSram$, and the last paragraph of that proof gives $\aucol(\auS_\iota\cap
\auSram,S_0\cap\auSram)\mid N^{\,i}$. The hypotheses of case (iii) are: Lemma~\ref{pl:kernelN} and
the integrality of the branching map (Lemma~\ref{ps:branching}(c)); neither
Proposition~\ref{au:expansion} nor (K0)--(K2) is used.

\begin{remark}[Dependencies]\label{pl:status}
Lemma~\ref{pa:hecke} is proved from Lemma~\ref{pa:kernel} (K1), (K2),
Proposition~\ref{au:expansion}, the facts (W4), (W5) on the Weil representation in the form (N5),
the equivariance (T1) of the kernels under rational points (which is part of their definition as
maps to algebraic modular forms, \S\ref{ps:ss-weil}), and Lemma~\ref{ps:branching}(c) (through the
restriction, as in Lemma~\ref{pa:chain}). The transfer
of Hecke operators (Proposition~\ref{pl:onelift}) is the unramified theta correspondence for
$(\auGL_n,\auGL_{n'})$ at a split place, in the form of Rallis~\cite{RallisFunctoriality}; the
proof given is self-contained. Lemma~\ref{pa:level} is proved, in the form of
Lemma~\ref{pl:level}, from Lemma~\ref{pl:kernelN} and the same facts.
\end{remark}


%% ===== proofs/A-final.tex =====
% proofs/A-final.tex
% The statements of Theorem A left open by proofs/A-seesaw-complete.tex: pk:curve(C3)(ii), the
% integrality (K1) at all primes (including 2, the primes of q_1, p_0, p_1 and N), (W6), the orbit
% classification of ps:support, the Chowla--Selberg unit at 2, the number of classes.
% Labels of this file start with pf:. Only amsmath/amssymb/amsthm are used.
\providecommand{\auO}{\mathcal O}
\providecommand{\auU}{\mathrm U}
\providecommand{\auGL}{\mathrm{GL}}
\providecommand{\auHom}{\operatorname{Hom}}
\providecommand{\auA}{\mathbb A}
\providecommand{\auF}{\mathbb F}
\providecommand{\auPl}{\Phi}
\providecommand{\auS}{\mathcal S}
\providecommand{\auNm}{\operatorname{Nm}}
\providecommand{\auAut}{\operatorname{Aut}}
\providecommand{\aubar}[1]{\overline{#1}}
\providecommand{\psM}{\mathsf M}
\providecommand{\psHS}{\mathcal H}
\providecommand{\psZ}{Z}
\providecommand{\psvac}{\mathbf v}
\providecommand{\psrho}{\varrho}
\providecommand{\pao}{\mathfrak o}
\providecommand{\pam}{\mathfrak m}
\providecommand{\pktr}{\operatorname{tr}}
\providecommand{\pfT}{\operatorname{Tr}}
\providecommand{\pfv}{\operatorname{ord}}
\providecommand{\pfbr}[1]{\theta[#1]}

\subsection{The last statements of Theorem~A}\label{pf:sec}

This subsection follows \S\ref{pk:sec} and uses its notation and that of \S\ref{pa:ss-curve}. It
proves the following statements used there: Proposition~\ref{pk:curve}(C3)(ii)
(Theorem~\ref{pf:recip}); the integrality in Lemma~\ref{pa:kernel}(K1) at every prime $\ell\ne d$,
and Lemma~\ref{pl:kernelN} without any restriction (Theorem~\ref{pf:std},
Corollary~\ref{pf:K1}); the fact (W6) (Proposition~\ref{pf:flip}); the classification of orbits at
the end of Lemma~\ref{ps:support} (Lemma~\ref{pf:orbits}); the unit property at $2$ of the
Chowla--Selberg unit (Corollary~\ref{pf:CS}); and the bound for the number of classes used in
Propositions~\ref{au:size} and~\ref{au:good} (Proposition~\ref{pf:classes}). Two conventions, which matter
only at the prime $2$, are recalled in \S\ref{pf:ss-adm}: the point $t$ of order two is the
\emph{admissible} one, and the character $\epsilon_L$ of the vacuum is defined with the element
$\omega=(1+\sqrt{-d})/2$. Remark~\ref{pf:gen} records that the translates $\theta_u$ with
$u\in t+E_1$ are also $\partial$-integral.

Throughout, $\ell$ is a prime, $\iota\colon\C\to\C_\ell$ a ring homomorphism, $\pao$ the valuation
ring of $\C_\ell$, $\pam$ its maximal ideal; a complex number $c$ is \emph{integral} (a
\emph{unit}) if $\iota(c)\in\pao$ ($\in\pao^\times$), and $|c|$ means $|\iota(c)|$. We write
$\omega=(1+\sqrt{-d})/2$, so that $\auO_K=\Z[\omega]$, $\omega+\aubar\omega=1$ and
$\omega\aubar\omega=(1+d)/4$, an even number divisible by every prime dividing $q_1$; and
$E(u,v)=\Iml(\aubar uv)/A$, so that $H(u,v)-H(v,u)=2\boldsymbol\pi i\,E(u,v)$ and $E$ is integral
on $\Lambda$ and unimodular.

\subsubsection{The admissible point of order two}\label{pf:ss-adm}

Let $\Lambda$ be a lattice with $\auO_K\Lambda=\Lambda$. The prime $2$ splits in $K$,
$2\auO_K=\mathfrak p_2\aubar{\mathfrak p_2}$, and $\tfrac12\Lambda/\Lambda$ is free of rank one over
$\auO_K/2=\auF_2\times\auF_2$. Its three non-zero points are the generator of the part killed by
$\mathfrak p_2$, the generator of the part killed by $\aubar{\mathfrak p_2}$, and their sum. The
last one is the \emph{admissible point}: it is the only $t\in\tfrac12\Lambda\setminus\Lambda$ with
$\omega t\notin\Lambda$ and $\aubar\omega t\notin\Lambda$ (exactly one of $\omega,\aubar\omega$ lies
in each of the two primes). \emph{From now on a normalised lattice $(\Lambda,t,a_1)$ is one in which
$t$ is the admissible point}; this is needed at the prime $2$ only. For $\Lambda=\Omega_0\auO_K$ this fixes $\Omega_0$ up to a fourth root of unity.

\begin{lemma}\label{pf:chi}
Let $(\Lambda,t,a_1)$ be a normalised lattice and $\chi$ the function of \eqref{pa:eq-theta}.
\begin{enumerate}
\item[(a)] For $\lambda\in\Lambda$: $\chi(\lambda)=-1$ if $\lambda/2\in t+\Lambda$, and
$\chi(\lambda)=1$ otherwise. For any choice of $t\in\frac12\Lambda\setminus\Lambda$ this holds.
\item[(b)] $\chi(\aubar\delta\lambda)=\chi(\lambda)^{\auNm\delta}$ for all $\delta\in\auO_K$. If
$\Lambda=\Omega_0\auO_K$, then $\chi(\Omega_0a)=(-1)^{\auNm a}$ for $a\in\auO_K$.
\item[(c)] For $u\in\C$ and $\lambda\in\Lambda$, $\theta_u(z+\lambda)=\chi(\lambda)e^{2\boldsymbol
\pi iE(\lambda,u)}e^{H(\lambda,z+\lambda/2)}\theta_u(z)$ and $\theta_{u+\lambda}=\chi(\lambda)
e^{\boldsymbol\pi iE(\lambda,u)}\theta_u$. In particular $\theta_{u+2n\lambda}=\theta_u$ if
$n\ge1$ and $nu\in\Lambda$, and $\theta_\lambda=\chi(\lambda)\theta$.
\item[(d)] Let $\delta\in\auO_K$, $\delta\ne0$, and $\delta u\in\Lambda$. Then $z\mapsto\theta_u
(\aubar\delta z)$ satisfies $G(z+\lambda)=\chi(\lambda)^{m}e^{mH(\lambda,z+\lambda/2)}G(z)$
$(\lambda\in\Lambda)$ with $m=\auNm\delta$.
\end{enumerate}
\end{lemma}
\begin{proof}
(a) $\chi(\lambda)=\varepsilon(\lambda)e^{2\boldsymbol\pi iE(\lambda,t)}$. If $\lambda\in2\Lambda$
both factors are $1$. Otherwise $\varepsilon(\lambda)=-1$ and $e^{2\boldsymbol\pi iE(\lambda,t)}=
(-1)^{E(\lambda,2t)}$; the form $E$ modulo $2$ is the non-degenerate alternating form on
$\Lambda/2\Lambda\cong\auF_2^2$, so $E(\lambda,2t)$ is even exactly when $\lambda\equiv2t$ modulo
$2\Lambda$. (b) Identify $\Lambda/2\Lambda$ with $\auF_2\times\auF_2$ as an $\auO_K/2$-module so
that $2t=(1,1)$, and let $\aubar\delta$ act by $(a,b)$. By (a), $\chi(\aubar\delta\lambda)=-1$ if
and only if $a=b=1$ and $\lambda\equiv2t$; and $a=b=1$ means that $\auNm\delta$ is odd. For
$\Lambda=\Omega_0\auO_K$ the class of $\Omega_0a$ is $(1,1)$ exactly when $a\notin\mathfrak p_2
\cup\aubar{\mathfrak p_2}$, that is, when $\auNm a$ is odd. (c) is a direct computation from
\eqref{pa:eq-theta} and the definition of $\theta_u$; for the last assertions, $\chi(2n\lambda)=1$
by (a) and $E(2n\lambda,u)=2E(\lambda,nu)\in2\Z$. (d) By (c), $G(z+\lambda)=\chi(\aubar\delta
\lambda)e^{2\boldsymbol\pi iE(\aubar\delta\lambda,u)}e^{H(\aubar\delta\lambda,\aubar\delta z+
\aubar\delta\lambda/2)}G(z)$; here $E(\aubar\delta\lambda,u)=E(\lambda,\delta u)\in\Z$,
$H(\aubar\delta\lambda,\aubar\delta w)=\auNm\delta\,H(\lambda,w)$, and (b) applies.
\end{proof}

\begin{lemma}[Good reduction everywhere]\label{pf:two}
Let $(\Lambda,t,a_1)$ be a normalised lattice. Then $a_1$ is a unit at every place above $2$.
Consequently $a_1^4-64$ is a unit at every prime, the scheme $y^2+a_1xy=x^3+x$ over the ring of
integers of $\Q(a_1)$ is an elliptic curve, and at every $\iota$ the point $P_0=(0,0)$ does not
reduce to the origin.
\end{lemma}
\begin{proof}
Let $\ell=2$ and let $L\subset\C_2$ be a finite extension of $\Q_2$ over which $\mathcal E$ has
all its points of order $2$ and has good reduction; it exists because $j$ is integral
\cite[Chapter~VII, Proposition~5.5]{SilvermanAEC}. Let $y'^2+a_1'x'y'+a_3'y'=x'^3+a_2'x'^2+a_4'x'+a_6'$ be a model with
integral coefficients and unit discriminant $\Delta'$, with formal parameter $\mathsf t=-x'/y'$. As
$2$ splits in $K$, the reduction is ordinary (Deuring's criterion~\cite[Chapter~13, Theorem~12]{LangEF}), i.e.\ the reduced
curve has a point of order $2$ over the algebraic closure of the residue field; this is the
definition of ``ordinary'' used here, and we prove the direction needed. Enlarge $L$ so that it
contains $K$ and the fields of definition of the endomorphisms $[\alpha]$, $\alpha\in\auO_K$ (as
$\auO_K\Lambda=\Lambda$, these are the maps $z\mapsto\alpha z$ of $\C/\Lambda$). The smooth model
is the N\'eron model of its generic fibre~\cite[\S1.2, Definition~1, p.~12, and Proposition~8, p.~15]{BLR} (\cite[IV, Corollary~6.3, p.~329, and the N\'eron mapping property, p.~319]{SilvermanAT}), so each $[\alpha]$ extends to it and
reduces to an endomorphism $\widetilde{[\alpha]}$ of $\widetilde{\mathcal E}$ of the same degree
$\auNm\alpha$ (the curve and the $[\alpha]$ are defined over a number field and have good
reduction at the prime above $2$ in question, and reduction of homomorphisms then preserves
degrees~\cite[Chapter~II, Proposition~4.4]{SilvermanAT}). The invariant differential $\omega'=dx'/(2y'+a_1'x'+a_3')$ is a
nonzero multiple of $dz$, so $[\alpha]^*\omega'=\alpha\omega'$; it reduces to the invariant
differential $\tilde\omega'=(1+\dots)d\tilde{\mathsf t}$ of $\widetilde{\mathcal E}$, so
$\widetilde{[\alpha]}^*\tilde{\mathsf t}=\tilde\alpha\tilde{\mathsf t}+O(\tilde{\mathsf t}^2)$, with
$\tilde\alpha$ the residue of $\alpha$. Let $\mathfrak p=\{\alpha\in\auO_K:\tilde\alpha=0\}$, one of
the two primes above $2$, and $h$ the class number of $K$; then $\aubar{\mathfrak p}^h=(\beta)$ with
$\beta\notin\mathfrak p$ (else $\aubar{\mathfrak p}\subseteq\mathfrak p$), so $\tilde\beta\neq0$ and
$\widetilde{[\beta]}$ has ramification index $1$ at the origin $O$. For $R$ in its kernel,
$\widetilde{[\beta]}\circ\tau_R=\widetilde{[\beta]}$ with $\tau_R$ the translation, so the index is
$1$ at every point of the kernel; as the indices over $O$ add up to the degree~\cite[Chapter~II, Proposition~2.6(a)]{SilvermanAEC},
the kernel has $\auNm\beta=2^h\ge2$ points. They are killed by $[2^h]$, so some multiple of a nonzero
one has order $2$. In characteristic $2$ a point
of order $2$ satisfies $\tilde y=-\tilde y-\tilde a_1'\tilde x-\tilde a_3'$, i.e.\ $\tilde a_1'\tilde
x+\tilde a_3'=0$; if $\tilde a_1'=0$ then $\tilde b_2'=\tilde b_4'=0$, $\tilde b_6'=\tilde a_3'^2$,
$\tilde\Delta'=-27\tilde a_3'^4=\tilde a_3'^4\ne0$, and no point satisfies this. So $a_1'$ is a unit. A
point $T\ne O$ reduces to the origin iff $v(x'(T))<0$ (if $v(x')\ge0$ then $y'$, a root of a monic
polynomial over the integers, is integral; if $v(x')<0$ then $v(y')=\tfrac32v(x')$ and
$[x':y':1]=[x'/y':1:1/y']$ reduces to $[0:1:0]$). The $x'$-coordinates of the points of order $2$ are the
roots of $g(X)=4X^3+b_2'X^2+2b_4'X+b_6'$, with $v(4)=2$, $v(b_2')=v(a_1'^2+4a_2')=0$, $v(2b_4')\ge1$,
$v(b_6')\ge0$ (normalising $v(2)=1$). The Newton polygon of $g$ has the segment from $(2,0)$ to $(3,2)$ and
all other slopes $\le0$; so exactly one root has negative valuation, namely $-2$. Hence the kernel of reduction
contains exactly one point $T_c$ of order two, and $v(x'(T_c))=-2$. This kernel is stable under $\auO_K$, so
$\{O,T_c\}$ is the part of $\mathcal E[2]$ killed by one of the two primes above $2$; let $T_e$
generate the other part. The admissible point is $P_0=T_c+T_e$ (the point $z=0$ is the point $t$
of $\C/\Lambda$). So $P_0$ and $T_e$ do not reduce to the origin, their coordinates $x'$ are
integral, and $v(x'(T_c)-x'(P_0))=-2$. The numbers $x'(T)$, $T\in\mathcal E[2]\setminus O$, are the
roots of $4X^3+b_2'X^2+2b_4'X+b_6'$, and the square of the product of their differences is
$\Delta'/16$ (for a cubic $aX^3+bX^2+cX+d$ with roots $e_i$, $a^4\prod_{i<j}(e_i-e_j)^2=b^2c^2-4ac^3-4b^3d
-27a^2d^2+18abcd$; with $(a,b,c,d)=(4,b_2',2b_4',b_6')$ this is $4(b_2'^2b_4'^2-32b_4'^3-b_2'^3b_6'-108b_6'^2
+36b_2'b_4'b_6')$, and by the definition $\Delta'=-b_2'^2b_8'-8b_4'^3-27b_6'^2+9b_2'b_4'b_6'$ with
$4b_8'=b_2'b_6'-b_4'^2$ the bracket is $4\Delta'$, so the right side is $16\Delta'$; dividing by $a^4=256$ gives $\Delta'/16$; see also
\cite[Chapter~III, \S1]{SilvermanAEC}); hence $2(-2)+2(-2)+2v(x'(T_e)-x'(P_0))=-4$ and
$v(x'(T_e)-x'(P_0))=2$. The number
\[
\varrho=\frac{x'(T_c)-x'(P_0)}{x'(T_e)-x'(P_0)},\qquad v(\varrho)=-4,
\]
does not change under the changes of coordinates $x'=u^2x+r$ between Weierstrass models. In the
model $\mathcal E$ we have $x(P_0)=0$, and $x(T_c),x(T_e)$ are the roots of $X^2+\tfrac14a_1^2X+1$;
so $\varrho+2+\varrho^{-1}=(x(T_c)+x(T_e))^2/(x(T_c)x(T_e))=a_1^4/16$, and $v(a_1^4)=4+v(\varrho)=0$.
The rest follows from Lemma~\ref{pa:model} (the discriminant is $a_1^4-64$; at odd primes $P_0$,
of order two, does not reduce to the origin).
\end{proof}

\begin{corollary}[Proposition~\ref{pk:curve}(C1), (C2), (C4) at every prime]\label{pf:CS}
With normalised lattices in the sense above, {\rm(C1)} holds, {\rm(C2)} holds at every prime
$\ell\ne d$ including $\ell=2$, and the number $\mathsf u$ of {\rm(C4)} is a unit at every prime,
including $2$.
\end{corollary}
\begin{proof}
(C1): the point $t^\sigma$ constructed there is admissible, by Lemma~\ref{pf:transport}(a) below
(the bijection $S_\sigma$ is $\auO_K$-linear and maps $t$ to $t^\sigma$). The proof of (C2) uses
only that $a_1$ is an algebraic integer and that torsion points of order prime to $\ell$ have
integral coordinates; it is valid for $\ell=2$. In the proof of (C4), with the lattices
$(\Lambda_i,t_i,a_{1,i})$ normalised in the present sense (this fixes $\mu_i^{12}$), the only
properties of an odd prime that are used are: $\Delta(\Lambda_i)=a_{1,i}^4-64$ is a unit, and the
curves have good reduction, so that homomorphisms extend to the smooth models and multiply the
invariant differentials $dx/(2y+a_1x)$ by integers. By Lemma~\ref{pf:two} both hold at $2$.
\end{proof}

\emph{The character of the vacuum.} Recall (\S\ref{ps:ss-weil}) that the character $\epsilon_L$ is
\begin{equation}\label{pf:eq-eps}
\epsilon_L(l)=\psi_f\bigl(-\tfrac12\,\omega\,\jmath\,\langle l,l\rangle\bigr)\qquad(l\in L),
\end{equation}
that is, the element ``$t$ with $t+\aubar t=1$'' is $\omega$, and there is a factor $-\frac12$. We check the cocycle relation: $\epsilon_L(l)\epsilon_L(l')
\epsilon_L(l+l')^{-1}=\psi(\frac12\omega\jmath(a+\aubar a))$ with $a=\langle l',l\rangle$, and
$\psi(\frac12\omega\jmath\aubar a)=\psi(-\frac12\aubar\omega\jmath a)$ because $\psi\circ c=\psi$;
so the quotient is $\psi(\frac12\jmath a(\omega-\aubar\omega))=\beta(l,l')\psi(\jmath a(\omega-1))=
\beta(l,l')$, as $\psi(\jmath\langle l',l\rangle\widehat\auO_E)=1$ by self-duality. Another
choice of the element of trace one replaces $\psvac_L$ by its translate by a point of
$\frac12L$; this changes nothing at odd primes, but at $\ell=2$ the statement (K1) would then be
false (the theta function of such a translate contains factors $\theta_s$ with $s$ of order two,
which are not $\partial$-integral at $2$).

\begin{lemma}\label{pf:eps}
Let $W$ be a hermitian space and $L\subset W(\auA_f)$ a self-dual lattice. Then $\epsilon_L$,
defined by \eqref{pf:eq-eps}, has values $\pm1$. If $W=E^{\mathsf C}$ is standard and
$L=\alpha\widehat\auO_E^{\,\mathsf C}$, then $\epsilon_L(\lambda)=(-1)^{\pfT_{E^+/\Q}\langle\lambda,
\lambda\rangle}$ for $\lambda\in\Lambda_L=W(E)\cap L=\mathfrak a_0^{\mathsf C}$.
\end{lemma}
\begin{proof}
Let $\nu=\langle l,l\rangle\in\auA_{E^+,f}$. Since $\omega\jmath=(\jmath-1)/2$ and
$\pfT_{E/E^+}(\jmath\nu)=0$, $\epsilon_L(l)=\psi_f(\nu/4)=\psi_{E^+,f}(\nu/2)$, where
$\psi_{E^+}=\psi_\Q\circ\pfT_{E^+/\Q}$. Since $[x,y]=\beta(x,y)\beta(y,x)^{-1}=\psi(\jmath\langle y,x\rangle)$ and $L$ is a
self-dual $\widehat\auO_E$-module, $\psi_f(\jmath\nu\widehat\auO_E)=1$, that is, $\jmath\nu\in
\mathfrak d_E^{-1}\widehat\auO_E$; as $\jmath^{-1}=-\sqrt{-d}$ is integral, $\nu\in\mathfrak
d_E^{-1}\widehat\auO_E\cap\auA_{E^+,f}=\mathfrak d_{E^+}^{-1}\widehat\auO_{E^+}$, because $E/E^+$ is
unramified at the finite places ($\mathfrak d_E=\mathfrak d_{E^+}\auO_E$).
Hence $\psi_{E^+,f}(\nu)=1$ and $\epsilon_L(l)^2=1$. In the standard case, for
$\lambda\in\Lambda_L$ the number $n=\pfT_{E^+/\Q}\langle\lambda,\lambda\rangle=\pfT_{E/K}
\langle\lambda,\lambda\rangle$ lies in $\pfT_{E/K}(\mathfrak a_0\aubar{\mathfrak a_0})\cap\R=
\pfT_{E/K}(\mathfrak d_{E/K}^{-1})\cap\R\subseteq\Z$, and $\psi_{\Q,f}(n/2)=e^{2\boldsymbol\pi in/2}$.
\end{proof}

\subsubsection{Values at torsion points and the integrality principle}\label{pf:ss-tors}

Let $(\Lambda,t,a_1)$ be a normalised lattice, $\mathcal E$ its curve. We write $Q(z)$ for the
point of $\mathcal E$ with parameter $z$: $Q(z)=(x(z),y(z))$ for $z\notin t+\Lambda$ and $Q(z)=O$
for $z\in t+\Lambda$; thus $z\mapsto Q(z)-P_0$ is the group isomorphism $\C/\Lambda\to\mathcal
E(\C)$ of the Weierstrass uniformisation, in the coordinates $X=x+\wp(t)$, $2Y=2y+a_1x$. Torsion
points have algebraic coordinates; $E_1$ denotes the group of the torsion points $v\in\C/\Lambda$
whose point $Q(v)-P_0$ reduces to the origin at $\iota$ (the kernel of reduction). It is stable
under $\auO_K$ and has no element $\ne0$ of order prime to $\ell$~\cite[Chapter~VII, Proposition~3.1(a)]{SilvermanAEC}; by
Lemma~\ref{pf:two}, $t\notin E_1$. For $u\in\C$ put
\[
\pfbr u=\theta_u(0)=e^{-H(u,u)/2}\theta(u);\qquad\text{then}\qquad e^{-H(w,w)/2}\theta_u(w)=
e^{\boldsymbol\pi iE(w,u)}\,\pfbr{w+u}\quad(u,w\in\C).
\]
For $m\ge0$ let $\mathcal S_m=\mathcal S_m(\Lambda)$ be the space of entire functions $G$ with
$G(z+\lambda)=\chi(\lambda)^me^{mH(\lambda,z+\lambda/2)}G(z)$ for $\lambda\in\Lambda$. If
$G\in\mathcal S_m$, then $G/\theta^m$ is $\Lambda$-periodic with poles only on $t+\Lambda$, of
order at most $m$; so
\begin{equation}\label{pf:eq-sections}
\mathcal S_m=\theta^m\cdot\Bigl\{P(x,y)=\sum_{2a+3e\le m,\ e\le1}c_{a,e}\,x^ay^e\Bigr\},
\end{equation}
the monomials $x^ay^e$, $e\le1$, being a basis of the functions on $\mathcal E$ regular outside
$O$, with a pole of order $2a+3e$ at $O$. In particular $\mathcal S_0=\C$.

\begin{proposition}[Values at torsion points]\label{pf:values}
Let $u\in\Q\Lambda$. If $u\in t+\Lambda$ then $\pfbr u=0$. Otherwise $\pfbr u$ is an algebraic
number and, at every prime $\ell$ and every $\iota$,
\[
|\pfbr u|=1\ \text{ if }u-t\notin E_1,\qquad|\pfbr u|=|x(u)|^{-1/2}<1\ \text{ if }u-t\in E_1 .
\]
In particular $\theta_u(0)$ is an algebraic integer for every torsion point $u$.
\end{proposition}
\begin{proof}
Let $u\notin t+\Lambda$ have order $\mathsf N$ modulo $\Lambda$ and put $n=1+2\mathsf N\ell$. By
Lemma~\ref{pf:chi}(d) the function $\theta(nz)$ lies in $\mathcal S_{n^2}$, and as in
Lemma~\ref{pa:units}(a)
\[
\theta(nz)=\theta(z)^{n^2}R_n(x(z)),\qquad R_n(X)=\prod_s\bigl(1-X/x_s\bigr),
\]
the product over the $(n^2-1)/2$ pairs $\pm s$ of non-zero points of $n^{-1}\Lambda/\Lambda$, with
$x_s=x(s-t)$ the $x$-coordinate of the point $Q(s-t)$; the point $Q(s-t)-P_0$ has parameter $s-t$,
and $Q(s-t)$ reduces to the origin $O=Q(t)$ only if $s=(s-t)-t$ lies in $E_1$, which is impossible since
$s\ne0$ has order dividing $n$, prime to $\ell$. So these points do not reduce to
the origin, so the $x_s$ are integral; and $|x_s|=1$: if $|x_s|<1$ then
$|y(y+a_1x)|=|x^3+x|<1$ at this point, so it reduces to $\widetilde P_0$, that is, $s-t\in E_1$; then $n(s-t)\equiv t$ ($n$ is odd)
would lie in $E_1$, contrary to Lemma~\ref{pf:two}. Since $H(nz,nz)=n^2H(z,z)$, we get
$\pfbr{nz}=\pfbr z^{n^2}R_n(x(z))$ for all $z\notin t+\Lambda$. Take $z=u$. As $nu=u+\lambda$ with
$\lambda=2\ell\mathsf Nu\in2\Lambda$, Lemma~\ref{pf:chi}(c) gives $\pfbr{nu}=\zeta\pfbr u$ with
$\zeta=e^{\boldsymbol\pi iE(\lambda,u)}$ a root of unity, and $\pfbr u\ne0$; so
\[
\pfbr u^{\,n^2-1}=\zeta\cdot\prod_s\frac{x_s}{x_s-x(u)} .
\]
This shows that $\pfbr u$ is algebraic. The point $Q(u)-P_0$ has parameter $u$, so $Q(u)$ reduces to the origin
if and only if $u+t\equiv u-t$ lies in $E_1$, and $x(u)$ is integral if and only if $u-t\notin E_1$. If $u-t\in E_1$, then $|x(u)|>1$, every
factor has absolute value $|x(u)|^{-1}$, and $|\pfbr u|^{n^2-1}=|x(u)|^{-(n^2-1)/2}$. If
$u-t\notin E_1$, every factor has absolute value $\ge1$, with equality unless $x_s\equiv x(u)$
modulo $\pam$, that is, unless $Q(u)$ and $\pm Q(s-t)$ have the same reduction, that is,
$u+t\mp s\in E_1$. Then $(n-1)(u+t\mp s)\equiv\pm s$ (as $(n-1)u,(n-1)t\in\Lambda$, $ns\in\Lambda$)
would be a non-zero point of
$E_1$ of order prime to $\ell$, which is impossible. So $|\pfbr u|=1$.
\end{proof}

\begin{proposition}[Integrality principle]\label{pf:pull}
Let $G\in\mathcal S_m$ and suppose that $e^{-mH(z,z)/2}G(z)$ is integral for all torsion points $z$
of odd order prime to $\ell$. Then $G=\theta^m\cdot\Gamma(y)$ with $\Gamma\in\pao[[Y]]$, as power
series in $z$. In particular, if $\delta\in\auO_K$, $\delta\ne0$, and $\delta u\in\Lambda$, then
\[
\theta_u(\aubar\delta z)=\theta(z)^{\auNm\delta}\,\Gamma_{u,\delta}(y(z)),\qquad\Gamma_{u,\delta}\in
\pao[[Y]],
\]
at every prime $\ell$, and for $\ell\ne d$ the series $\theta_u(\aubar\delta z)$ is
$\partial$-integral. If moreover $\iota(\aubar\delta)$ is a unit and $\ell\ne d$, then $\theta_u$
itself is $\partial$-integral.
\end{proposition}
\begin{proof}
Write $G=\theta^mP(x,y)$ as in \eqref{pf:eq-sections}. Suppose that $P$ is not integral, let
$c_0$ be a coefficient of $P$ of largest absolute value, $|c_0|>1$, and $P_1=P/c_0$. The reduction
$\widetilde P_1$ of $P_1$ modulo $\pam$ is a non-zero function on the reduced curve (the monomials
are linearly independent there), regular outside $\widetilde O$ with a pole of order $\le m$, so
it has at most $m$ zeros. Let $z$ be a torsion point of odd order $n'$
prime to $\ell$. Then $z\notin E_1+t$, since otherwise $n'(z-t)\equiv t$ would lie in $E_1$; so
$Q(z)$ has integral coordinates, and by Proposition~\ref{pf:values} $\pfbr z$ is a unit. Hence
$P(Q(z))=G(z)/\theta(z)^m=e^{-mH(z,z)/2}G(z)\cdot\pfbr z^{-m}$ is integral, and $P_1(Q(z))=P(Q(z))/c_0\in\pam$.
Thus $\widetilde P_1$ vanishes at the reductions of all these points; they are pairwise distinct
(differences of such $z$ have odd order prime to $\ell$) and infinite in number, a contradiction.
So $P$ has coefficients in $\pao$, and $P(x,y)\in\pao[[y]]$ because $x\in y^2\pao[[y]]$
(\S\ref{pa:ss-curve}; this uses only that $a_1$ is integral).

For $G(z)=\theta_u(\aubar\delta z)$ and $m=\auNm\delta$: $G\in\mathcal S_m$ by
Lemma~\ref{pf:chi}(d), and $e^{-mH(z,z)/2}G(z)=e^{-H(w,w)/2}\theta_u(w)$ with $w=\aubar\delta z$ is
a root of unity times $\pfbr{w+u}$, which is integral by Proposition~\ref{pf:values}. The series
$\theta$ and $y$ are $\partial$-integral for $\ell\ne d$ (Corollary~\ref{pf:CS}), and
$\partial$-integral series are stable under products, under substitution into series of
$\pao[[Y]]$, and under $z\mapsto cz$ with $c\in\pao$; for the last assertion take
$c=\iota(\aubar\delta)^{-1}$.
\end{proof}

\begin{remark}[The case excluded in \S\ref{pk:ss-N}]\label{pf:gen}
Proposition~\ref{pf:pull} has no exception: for $u=t+e$ with $0\ne e\in E_1$ of order $\ell$, and
$\delta$ with $\delta e=0$, $\iota(\aubar\delta)\in\pao^\times$, the series $\theta_u$ \emph{is}
$\partial$-integral. Note that $|\theta_u(0)|<1$ (Proposition~\ref{pf:values}) and that $\theta_u'(0)/\theta_u(0)$
is not integral, but this does not prevent integrality of the derivatives: $\theta_u=\theta\cdot
\Gamma(y(\cdot/\aubar\delta))$ with $\Gamma\in\pao[[Y]]$ a series that has a zero in the open unit
disc (at the parameter of the point $-e$), like the series $Y-c$ with $0<|c|<1$. This is why the condition
(Gen) of Proposition~\ref{pk:K1N} is not needed, and Lemmas~\ref{pk:E1}, \ref{pk:formal} are not used. At $\ell=N$ the excluded case is the typical one: the kernel of reduction contains
$\mathfrak n^{-1}\Lambda/\Lambda$ for the prime $\mathfrak n$ induced by $\iota$, which is where the
translation points have their denominators.
\end{remark}

\subsubsection{Conjugation: proof of Proposition~\ref{pk:curve}{\rm(C3)(ii)}}\label{pf:ss-recip}

Let $(\Lambda,t,a_1)$ be a normalised lattice, $\sigma\in\auAut(\C/K)$, and
$(\Lambda^\sigma,t^\sigma,\sigma a_1)$ the lattice of (C1), with functions $x^\sigma,y^\sigma,
\theta^\sigma$, curve $\mathcal E^\sigma$, and $Q^\sigma(z)$. For $G=\theta^mP(x,y)\in\mathcal
S_m(\Lambda)$ put $G^\sigma=(\theta^\sigma)^mP^\sigma(x^\sigma,y^\sigma)\in\mathcal S_m
(\Lambda^\sigma)$, where $P^\sigma$ is $P$ with $\sigma$ applied to the coefficients.

\begin{lemma}[Transport]\label{pf:transport}
\begin{enumerate}
\item[(a)] There is a unique bijection $S_\sigma\colon\C/\Lambda\to\C/\Lambda^\sigma$ with
$Q^\sigma(S_\sigma z)=\sigma(Q(z))$. It is an isomorphism of $\auO_K$-modules, and
$S_\sigma(t)=t^\sigma$.
\item[(b)] $G\mapsto G^\sigma$ is $\sigma$-semilinear, multiplicative, and the Taylor series of
$G^\sigma$ at $0$ is the image under $\sigma$ of that of $G$. If $G\in\mathcal S_m(\Lambda)$ and
$G'$ is an entire function whose Taylor series is the image under $\sigma$ of that of $G$, then
$G'=G^\sigma$.
\item[(c)] For $z\in\C$ and $z'\in\C$ with $z'\equiv S_\sigma(z)$ there is $\mu=\mu(z,z')\in
\C^\times$ such that $G^\sigma(z')=\mu^m\,\sigma(G(z))$ for all $m\ge0$ and all $G\in\mathcal S_m
(\Lambda)$.
\end{enumerate}
\end{lemma}
\begin{proof}
(a) The map $\widetilde S=(Q^\sigma-P_0^\sigma)^{-1}\circ\sigma\circ(Q-P_0)$ is a group isomorphism
$\C/\Lambda\to\C/\Lambda^\sigma$, since $\sigma\colon\mathcal E(\C)\to\mathcal E^\sigma(\C)$ is one.
As $\sigma(P_0)=P_0^\sigma$, $Q^\sigma(\widetilde Sz)=\sigma(Q(z)-P_0)+P_0^\sigma=\sigma Q(z)$; so
$S_\sigma=\widetilde S$, and $S_\sigma(t)=t^\sigma$ because $Q(t)=O$. For $a\in\auO_K$ let $[a]$ be
the endomorphism of $\mathcal E$ induced by $z\mapsto az$; it is characterised by $[a]^*\varpi=
a\varpi$ for the invariant differential $\varpi=dx/(2y+a_1x)=dz$. The conjugate endomorphism
$\sigma[a]$ of $\mathcal E^\sigma$ satisfies $(\sigma[a])^*\varpi^\sigma=\sigma(a)\varpi^\sigma=
a\varpi^\sigma$, so it is $[a]$ on $\mathcal E^\sigma$, and $S_\sigma(az)=aS_\sigma(z)$.
(b) By (C1) the series of $x^\sigma,y^\sigma,\theta^\sigma$ are the conjugates of those of
$x,y,\theta$. The last assertion holds because an entire function is determined by its Taylor
series. (c) If $z\notin t+\Lambda$, then $P^\sigma(Q^\sigma(z'))=\sigma(P(Q(z)))$ and
$\mu=\theta^\sigma(z')/\sigma(\theta(z))$. If $z\in t+\Lambda$, let $\mathsf t=x/y$, a function
with a simple zero at $O$; for $2a+3e=m$ the function $x^ay^e\mathsf t^m=(x^3/y^2)^{(m-e)/2}$
takes the value $1$ at $O$, so $G(z)=c\,\lambda_z^m$, where $c$ is the coefficient in $P$ of the
monomial with $2a+3e=m$ (or $0$) and $\lambda_z=\lim_{w\to z}\theta(w)/\mathsf t(w)\ne0$; likewise
$G^\sigma(z')=\sigma(c)\lambda_{z'}'^{\,m}$, and $\mu=\lambda'_{z'}/\sigma(\lambda_z)$.
\end{proof}

\begin{theorem}[Proposition~\ref{pk:curve}{\rm(C3)(ii)}]\label{pf:recip}
Let $\Lambda=\Omega_0\auO_K$ be normalised and $\sigma\in\auAut(\C/K)$. There are $\Omega_\sigma
\in\C^\times$, an ideal $\mathfrak c\ne0$ of $\auO_K$ and a finite idele $\beta_\sigma$ of $K$ such
that, with $\mathfrak a=\mathfrak c\aubar{\mathfrak c}^{-1}$ and $\widehat{\mathfrak a}=\mathfrak
a\widehat\auO_K$:
\begin{enumerate}
\item[(a)] $\Lambda^\sigma=\Omega_\sigma\mathfrak a$, $\beta_\sigma\widehat\auO_K=\widehat{
\mathfrak a}$ (so $\beta_\sigma\aubar{\beta_\sigma}$ is a unit), and $S_\sigma(\mathsf x\Omega_0)
\equiv\mathsf x'\Omega_\sigma$ modulo $\Lambda^\sigma$ whenever $\mathsf x,\mathsf x'\in K$ and
$\mathsf x'-\beta_\sigma\mathsf x\in\widehat{\mathfrak a}$;
\item[(b)] if $\mathsf x\in K$, $n\ge1$, $n\mathsf x\in\auO_K$, and $\mathsf x'\in K$ satisfies
$\mathsf x'-\beta_\sigma\mathsf x\in2n\widehat{\mathfrak a}$, then $\sigma$ maps the Taylor series
at $0$ of $\theta_{\mathsf x\Omega_0}$ to that of $\theta^\sigma_{\mathsf x'\Omega_\sigma}$.
\end{enumerate}
\end{theorem}
\begin{proof}
(a) $\Lambda^\sigma$ is stable under $\auO_K$, so $\Lambda^\sigma=\Omega\mathfrak b$ for a
fractional ideal $\mathfrak b$. The class number of $K$ is odd. (Proof: only the prime $d$ ramifies in $K$, and
$(\sqrt{-d})$ is the prime above it. Let a class $c$ satisfy $c^2=1$ and let $\mathfrak a\in c$; since
$\mathfrak a\aubar{\mathfrak a}=(\auNm\mathfrak a)$, $\aubar{\mathfrak a}\in c^{-1}=c$, so $\aubar{\mathfrak
a}=\gamma\mathfrak a$ with $\gamma\in K^\times$, and $\gamma\aubar\gamma$ is a positive unit, i.e.\ $1$.
Then $\gamma=\beta/\aubar\beta$ with $\beta=1+\gamma$ if $\gamma\ne-1$ and $\beta=\sqrt{-d}$ if
$\gamma=-1$, and $\mathfrak b=\beta\mathfrak a\in c$ satisfies $\aubar{\mathfrak b}=\aubar\beta\gamma
\mathfrak a=\mathfrak b$. In the prime factorisation of a conjugation-stable fractional ideal, split
primes $\mathfrak p\ne\aubar{\mathfrak p}$ occur with equal exponents and their product is $(p)$, inert
primes are $(p)$, and the ramified prime is $(\sqrt{-d})$; so $\mathfrak b$ is principal and $c=1$. A
finite abelian group without elements of order $2$ has odd order. See also \cite{CoxPrimes}.) So $\mathfrak b=\eta\mathfrak c^2$ with $\eta\in K^\times$, and
$\mathfrak c^2=\auNm(\mathfrak c)\cdot\mathfrak c\aubar{\mathfrak c}^{-1}$. By
Lemma~\ref{pf:transport}(a), $S_\sigma$ induces an isomorphism of $\auO_K$-modules
$K/\auO_K\to K/\mathfrak a$ on the torsion points. Now
$\auHom_{\auO_K}(K/\auO_K,K/\mathfrak a)=\varprojlim_n\auHom(n^{-1}\auO_K/\auO_K,n^{-1}\mathfrak
a/\mathfrak a)=\varprojlim_n\mathfrak a/n\mathfrak a=\widehat{\mathfrak a}$, acting by
multiplication on $K/\auO_K=\auA_{K,f}/\widehat\auO_K$ with values in $\auA_{K,f}/\widehat{
\mathfrak a}=K/\mathfrak a$; the element $\beta_\sigma\in\widehat{\mathfrak a}$ of an isomorphism
generates $\widehat{\mathfrak a}$.

(b) Put $w=\Omega_0/(2n^2)$, choose $\mathsf b\in K$ with $\mathsf b-\beta_\sigma\in2n^2\widehat{
\mathfrak a}$, and put $w'=\mathsf b\Omega_\sigma/(2n^2)$, $u=\mathsf x\Omega_0=2n^2\mathsf xw$,
$u'=2n^2\mathsf xw'=\mathsf x\mathsf b\Omega_\sigma$. By (a), $w'\equiv S_\sigma(w)$, hence
$aw'\equiv S_\sigma(aw)$ for $a\in\auO_K$ and $u'\equiv S_\sigma(u)$. Also $\mathsf b\in\mathfrak
a$, so $nu'\in\Lambda^\sigma$; and $\mathsf x'-\mathsf x\mathsf b\in K\cap2n\widehat{\mathfrak a}=
2n\mathfrak a$, so $\theta^\sigma_{\mathsf x'\Omega_\sigma}=\theta^\sigma_{u'}$ by
Lemma~\ref{pf:chi}(c). Let
\[
g(z)=\theta_u(nz)\in\mathcal S_{n^2}(\Lambda),\qquad g'(z)=\theta^\sigma_{u'}(nz)\in\mathcal S_{n^2}
(\Lambda^\sigma)
\]
(Lemma~\ref{pf:chi}(d) with $\delta=n$). We show $g^\sigma=g'$; by Lemma~\ref{pf:transport}(b)
this is the assertion, since $\sigma(n)=n$.

\emph{$g^\sigma=c\,g'$ with $c\in\C$.} The zeros of $g'$ are simple and are the $z'$ with
$nz'+u'-t^\sigma\in\Lambda^\sigma$. Let $z'$ be one of them and $z\in\C$ with $S_\sigma(z)\equiv
z'$. Then $S_\sigma(nz+u-t)\equiv nz'+u'-t^\sigma\equiv0$, so $nz+u\in t+\Lambda$, $g(z)=0$, and
$g^\sigma(z')=0$ by Lemma~\ref{pf:transport}(c). Hence $g^\sigma/g'$ is entire and lies in
$\mathcal S_0(\Lambda^\sigma)=\C$.

\emph{$c=1$.} Put $p=u/2=n^2\mathsf xw$, $p'=u'/2$, $z_0=-n\mathsf xw$, $z_0'=-n\mathsf xw'$; then
$z_0'\equiv S_\sigma(z_0)$, $nz_0=-p$, $nz_0'=-p'$. Let $\mu=\mu(z_0,z_0')$. From the definition,
$\theta_u(-p)=e^{H(u,p)-H(u,u)/2}\theta(p)=\theta(p)$; so $g(z_0)=\theta(p)$ and $g'(z_0')=
\theta^\sigma(p')$. Let $h(z)=\theta(nz)\in\mathcal S_{n^2}(\Lambda)$; by (C1) and
Lemma~\ref{pf:transport}(b), $h^\sigma(z)=\theta^\sigma(nz)$, and $\theta$ is even, so
Lemma~\ref{pf:transport}(c) gives
\[
\theta^\sigma(p')=h^\sigma(z_0')=\mu^{n^2}\sigma(\theta(p)),\qquad g^\sigma(z_0')=\mu^{n^2}\sigma
(g(z_0))=\mu^{n^2}\sigma(\theta(p))=\theta^\sigma(p')=g'(z_0').
\]
If $\theta(p)\ne0$ this gives $c=1$. If $\theta(p)=0$, then $\theta^\sigma(p')=0$, so
$u=2p\in2t+2\Lambda$ and $u'\in2t^\sigma+2\Lambda^\sigma$; by Lemma~\ref{pf:chi}(a),(c),
$\theta_u=\chi(u)\theta=-\theta$ and $\theta^\sigma_{u'}=-\theta^\sigma$, so $g=-h$,
$g'=-h^\sigma=g^\sigma$.
\end{proof}

\begin{corollary}\label{pf:conj}
Lemma~\ref{ps:conj} and Proposition~\ref{au:kernel} hold.
\end{corollary}
\begin{proof}
Proposition~\ref{pk:conj} deduces Lemma~\ref{ps:conj} from (C3)(ii), which is
Theorem~\ref{pf:recip} (``close enough'' meaning the congruence of (b)); in its proof $\mathsf
B=\beta_\sigma$ is an idele with $\mathsf B\aubar{\mathsf B}$ a unit by (a), so that $\mathsf BL$
is self-dual, as is used there. Proposition~\ref{au:kernel} was proved in \S\ref{ps:ss-kernel} from
Lemma~\ref{ps:conj}.
\end{proof}

\subsubsection{Tight frames, the scale, and standard data at every prime}\label{pf:ss-frames}

Let $h(\mathsf x,\mathsf y)=\pfT_{E/K}(\aubar{\mathsf x}\mathsf y)$ on $E$; the lattice
$\mathfrak a_0$ is unimodular for $h$ over $\auO_K$, because $\mathfrak a_0\aubar{\mathfrak a_0}=
\mathfrak d_{E/K}^{-1}$. It is stable under $\operatorname{Gal}(E/K)$ and under $\auO_E$. A
\emph{tight frame of length $\mathsf b\ge1$} is a finite family $(v_i)_{i\in\mathsf I}$ in
$\mathfrak a_0$ with
\[
\sum_{i\in\mathsf I}\tau(v_i)\,\aubar{\tau'(v_i)}=\mathsf b\,\delta_{\tau\tau'}\qquad(\tau,\tau'\in
\auPl);
\]
then the vectors $u_i=v_i/\mathsf b$ with multiplicities $\mathsf b$, placed in each coordinate
of $E^{\mathsf C}$, form a resolution of the identity \eqref{pa:eq-res}. The family
$(\zeta_{q_1}^k)_{0\le k<q_1}$ is a tight frame of length $q_1$; it is the one of
\eqref{pk:eq-product}.

\begin{lemma}[A frame of length prime to $q_1$]\label{pf:frame}
There is a tight frame of a length $\mathsf h$ prime to $q_1$.
\end{lemma}
\begin{proof}
Let $\zeta=\zeta_{q_1}$, $P$ its minimal polynomial over $K$ (the cyclotomic polynomial, of degree
$f$), and $P(X)/(X-\zeta)=\sum_{i<f}b_iX^i$ with $b_i\in\auO_E$. By Euler's
formula~\cite[III, \S6, Lemma~2, p.~56]{SerreLF}, $\pfT_{E/K}(\zeta^k\,b_i/P'(\zeta))=\delta_{ik}$ for $0\le i,k<f$
(compare the coefficients of $X^i$ in the Lagrange interpolation identity
$\sum_{\sigma}\sigma(\zeta)^kP(X)/\bigl((X-\sigma\zeta)P'(\sigma\zeta)\bigr)=X^k$, $\sigma$ running over the
$K$-embeddings of $E$, which holds since both sides have degree $<f$ and agree at the $f$ roots of $P$),
and $P'(\zeta)$ generates $\mathfrak d_{E/K}$, which divides $q_1$. Choose $J\ge1$ with
$q_1\mid(\omega\aubar\omega)^J$ and put $e_i=\aubar\omega^{\,J}b_i/P'(\zeta)$. Then
$e_i\in\mathfrak a_0$: at the places dividing $\omega$, where $\alpha=1/P'(\zeta)$, the element
$P'(\zeta)e_i$ is integral; at the other places dividing $q_1$, $\omega$ is a unit
($\omega+\aubar\omega=1$) and $\aubar\omega^{\,J}/q_1=(\omega\aubar\omega)^J/(q_1\omega^J)$ is
integral; elsewhere $P'(\zeta)$ is a unit. Transposing the matrix identity that expresses the
duality of the bases $(\zeta^k)$ and $(b_i/P'(\zeta))$, and using that the $\tau\in\auPl$ fix $K$
and commute with $c$,
\begin{equation}\label{pf:eq-cross}
\sum_{i<f}\tau(e_i)\,\aubar{\tau'(\zeta^{-i})}=\sum_{i<f}\tau(e_i)\tau'(\zeta^{i})=\aubar\omega^{\,J}
\delta_{\tau\tau'} .
\end{equation}
For $\mathsf w\in\mathfrak a_0$ the family $(\zeta^k\rho(\mathsf w))_{k<q_1,\ \rho\in\operatorname{
Gal}(E/K)}$ has the matrix $\sum_{k,\rho}\tau(\zeta)^k\tau'(\zeta)^{-k}\tau\rho(\mathsf w)\aubar{
\tau'\rho(\mathsf w)}=q_1\nu(\mathsf w)\delta_{\tau\tau'}$, where $\nu(\mathsf w)=h(\mathsf w,
\mathsf w)\in\auO_K\cap\R_{\ge0}=\Z_{\ge0}$; so it is a tight frame of length $q_1\nu(\mathsf w)$
if $\mathsf w\ne0$. Consider the union of the frame $(\zeta^{-k})_{k<q_1}$ and of these frames for
$\mathsf w=e_0,\dots,e_{f-1}$; in it, for each $i<f$, replace the two members $\zeta^{-i}$ and
$e_i$ $(k=0$, $\rho=1)$ by the single member $\zeta^{-i}+e_i$. By \eqref{pf:eq-cross} and its
conjugate transpose the matrix of the new family is
\[
\Bigl(q_1\bigl(1+\sum_{i<f}\nu(e_i)\bigr)+\omega^J+\aubar\omega^{\,J}\Bigr)\delta_{\tau\tau'}=:
\mathsf h\,\delta_{\tau\tau'} .
\]
The diagonal entries are sums of squares of absolute values, so $\mathsf h$ is a positive integer.
For a prime $p\mid q_1$, $p$ divides $\omega\aubar\omega$; the integers $T_J=\omega^J+\aubar
\omega^{\,J}$ satisfy $T_J=T_{J-1}-\omega\aubar\omega T_{J-2}\equiv T_{J-1}$ and $T_1=1$, so
$\mathsf h\equiv1\pmod p$.
\end{proof}

\begin{lemma}[The scale]\label{pf:scale}
For every prime $\ell$, every $\iota$, every $\tau\in\auPl$ and every $\mathsf v\in\mathfrak a_0$,
the number $\varpi_{\mathrm{vac}}\cdot\aubar{\tau(\mathsf v)}$ is integral. More precisely
$|\varpi_{\mathrm{vac}}|=|\alpha_{w_\tau}|_{w_\tau}^{-1}$, where $w_\tau$ is the place of $E$
induced by $\iota\circ\tau\circ c$ and $|\cdot|_{w_\tau}=|\iota\tau c(\cdot)|$.
\end{lemma}
\begin{proof}
Let $\chi'=\upsilon|\cdot|^{-1/2}$, so that $\varpi_{\mathrm{vac}}^{-f}=\chi'(\alpha)$. It is a
Hecke character of $E$, trivial on $\widehat\auO_E^\times$ ($\upsilon$ is unramified at every
finite place), with $\chi'_\infty(\mathsf a)=\prod_\tau\tau(\mathsf a)/|\tau(\mathsf a)|^2=\prod_\tau
\aubar{\tau(\mathsf a)}^{-1}$. Let $k\ge1$ and $\mathsf a\in E^\times$ with $\alpha^k\in\mathsf a
\widehat\auO_E^\times$. Then $\chi'(\alpha)^k=\chi'_f(\mathsf a)=\chi'_\infty(\mathsf a)^{-1}=
\prod_\tau\aubar{\tau(\mathsf a)}$, and $|\aubar{\tau(\mathsf a)}|=|\mathsf a|_{w_\tau}=|\alpha|_{
w_\tau}^k$. The places $w_\tau$, $\tau\in\auPl$, are the places of $E$ above one place of $K$, they
are permuted transitively by $\operatorname{Gal}(E/K)$, and $\alpha$ is $1$ or a generator of the
Galois-stable ideal $\mathfrak d_{E/K}^{-1}$ at all of them; so $|\alpha|_{w_\tau}$ does not depend
on $\tau$, and $|\varpi_{\mathrm{vac}}|^{-fk}=|\alpha|_{w_\tau}^{fk}$. Finally $|\aubar{\tau(\mathsf
v)}|=|\mathsf v|_{w_\tau}\le|\alpha|_{w_\tau}$ for $\mathsf v\in\mathfrak a_0$.
\end{proof}

Let $W=E^{\mathsf C}$, $L$, $\Lambda_L$ be as in \S\ref{pk:ss-conv}, $\Lambda=\Omega_0\auO_K$
normalised, and let $(v_i)_{i\in\mathsf I}$ be a tight frame of length $\mathsf b$,
$u_i=v_i/\mathsf b$. Put
\[
\Pi^v(z)=\prod_{\mathsf c\in\mathsf C}\prod_{i\in\mathsf I}\theta\bigl(a_{i,\mathsf c}(z)\bigr)^{
\mathsf b},\qquad a_{i,\mathsf c}(z)=\tfrac{\Omega_0}{c_\psi}\,u_i^\dagger(z_{\mathsf c}) ,
\]
so that $a_{i,\mathsf c}(\gamma^\natural)=\Omega_0\pfT_{E/K}(\aubar{u_i}\gamma_{\mathsf c})$ for
$\gamma\in W(E)$, and, by \eqref{pa:eq-res}, as in \S\ref{pk:ss-conj},
\[
\langle z,x\rangle=\sum_{i,\mathsf c}\mathsf b\,H\bigl(a_{i,\mathsf c}(x),a_{i,\mathsf c}(z)\bigr),
\qquad|x|^2=\sum_{i,\mathsf c}\mathsf b\,H\bigl(a_{i,\mathsf c}(x),a_{i,\mathsf c}(x)\bigr).
\]

\begin{lemma}[The product and the vacuum have the same character]\label{pf:semi}
$\Pi^v=\Theta_{F^v}$ for a unique $F^v\in\psHS_f(W)$; $\theta(1\otimes F^v)=1$;
$\psrho_f(l)F^v=\epsilon_L(l)F^v$ for all $l\in\mathsf bL$; and for $\gamma\in W(E)$
\begin{equation}\label{pf:eq-transl}
\Theta_{\psrho_f(\gamma)F^v}(z)=\zeta_\gamma\prod_{i,\mathsf c}\theta_{\mathsf x_{i,\mathsf c}
\Omega_0}\bigl(a_{i,\mathsf c}(z)\bigr)^{\mathsf b},\qquad\mathsf x_{i,\mathsf c}=-\pfT_{E/K}(
\aubar{u_i}\gamma_{\mathsf c}),
\end{equation}
with a root of unity $\zeta_\gamma$.
\end{lemma}
\begin{proof}
For $\lambda\in\mathsf b\Lambda_L$ the numbers $\omega_{i,\mathsf c}=a_{i,\mathsf c}(\lambda^\natural)
=\Omega_0\pfT_{E/K}(\aubar{v_i}\lambda_{\mathsf c}/\mathsf b)$ lie in $\Omega_0\pfT_{E/K}(\mathfrak
d_{E/K}^{-1})\subseteq\Lambda$. By \eqref{pa:eq-theta} and the two identities above,
$\Pi^v(z+\lambda^\natural)=\epsilon(\lambda)e^{\langle z,\lambda^\natural\rangle+|\lambda^\natural
|^2/2}\Pi^v(z)$ with $\epsilon(\lambda)=\prod_{i,\mathsf c}\chi(\omega_{i,\mathsf c})^{\mathsf b}$.
By Lemma~\ref{pf:chi}(b) and \eqref{pa:eq-res},
\[
\epsilon(\lambda)=(-1)^{\sum_{i,\mathsf c}\mathsf b\,|\pfT_{E/K}(\aubar{u_i}\lambda_{\mathsf c})|^2}
=(-1)^{\sum_{\mathsf c}\sum_\tau|\tau(\lambda_{\mathsf c})|^2}=(-1)^{\pfT_{E^+/\Q}\langle\lambda,
\lambda\rangle}=\epsilon_L(\lambda)
\]
(Lemma~\ref{pf:eps}). So Lemma~\ref{pk:dict} gives $F^v$, and \eqref{pf:eq-transl} is proved like
\eqref{pk:eq-transl}. The computation in the proof of Lemma~\ref{pk:dict} shows that
$\Theta_f(z+\lambda^\natural)=\kappa(\lambda)^{-1}e^{\langle z,\lambda^\natural\rangle+|\lambda^
\natural|^2/2}\Theta_f(z)$ whenever $\psrho_f(\lambda)f=\kappa(\lambda)f$, and conversely by the
injectivity of $f\mapsto\Theta_f$; so $\psrho_f(\lambda)F^v=\epsilon_L(\lambda)^{-1}F^v=\epsilon_L
(\lambda)F^v$ for $\lambda\in\mathsf b\Lambda_L$, and for $l\in\mathsf bL$ by continuity.
\end{proof}

\begin{lemma}[Values and expansion of the integral vacuum]\label{pf:fourier}
{\rm(a)} For every $x\in W(\auA_f)$ the number $\theta(1\otimes\psrho^\circ(x)\psvac^\circ)$ is an
algebraic integer.
{\rm(b)} Let $(v_i)$ be a tight frame of length $\mathsf b$. Then $\psvac^\circ=\sum_ja_j\,
\psrho_f(y_j)F^v$, a finite sum with $y_j\in\mathsf b^{-1}L$ and $\mathsf b^{2f\#\mathsf C}a_j$
algebraic integers.
\end{lemma}
\begin{proof}
(a) By \eqref{pk:eq-vac}, $\psrho^\circ(x)\psvac^\circ$ is a sum, with coefficients roots of unity,
of vectors $\psrho_f(\gamma)F_0$, $\gamma\in W(E)$ ($F_0$ is an eigenvector of $\psrho_f(q_1L)$
with eigenvalues $\pm1$, and $W(E)+q_1L=W(\auA_f)$). By \eqref{pk:eq-transl} the value of $\theta$
on $1\otimes\psrho_f(\gamma)F_0$ is a root of unity times a product of numbers $\theta_u(0)$,
$u\in\Q\Lambda$, which are algebraic integers by Proposition~\ref{pf:values}.

(b) Let $\mathsf V$ be the space of the $f\in\psHS_f(W)$ with $\psrho_f(l)f=\epsilon_L(l)f$ for
$l\in\mathsf bL$, that is, $\psrho^\circ(l)f=f$. By (W4) it has the basis $\psrho^\circ(x)\psvac_L$,
$x\in\mathsf b^{-1}L/L$ (these vectors lie in $\mathsf V$ because $[l,x]=1$, and they are
eigenvectors of $\psrho(L)$ for $\mathsf b^{2f\#\mathsf C}=\dim\mathsf V$ distinct characters). The
map $\psrho^\circ$ is a homomorphism on $L$, and $\mathsf P=\mathsf b^{-2f\#\mathsf C}\sum_{l\in
L/\mathsf bL}\psrho^\circ(l)$ is the projector of $\mathsf V$ onto $\C\psvac_L$ along the other
basis vectors. Write $F^v=\sum_xc_x\psrho^\circ(x)\psvac_L$ (Lemma~\ref{pf:semi}). Then $\mathsf
P\psrho^\circ(-x)F^v=\eta_xc_x\psvac_L$ with the root of unity $\eta_x=\epsilon(x)\epsilon(-x)
\beta(-x,x)$. Let $\psvac^\circ=\kappa\psvac_L$ and $V^\circ(x)=\theta(1\otimes\psrho^\circ(x)
\psvac^\circ)$. Applying $\theta$ to the expansion of $F^v$ gives $1=\sum_xc_x\kappa^{-1}V^\circ
(x)$, and therefore
\[
\psvac^\circ=\sum_xV^\circ(x)\,c_x\psvac_L=\sum_{x\in\mathsf b^{-1}L/L}\ \sum_{l\in L/\mathsf bL}
\frac{V^\circ(x)}{\eta_x\,\mathsf b^{2f\#\mathsf C}}\;\psrho^\circ(l)\psrho^\circ(-x)F^v .
\]
Each $\psrho^\circ(l)\psrho^\circ(-x)$ is a root of unity times $\psrho_f(l-x)$, and $V^\circ(x)$
is an algebraic integer by (a).
\end{proof}

For $\delta\in\auO_K$, $\delta\ne0$, a \emph{standard datum of level $\delta$ at $\ell$} is a
combination, with coefficients in $\iota^{-1}(\pao)$, of vectors $\psrho^\circ(x)\psvac^\circ$ with
$\delta x$ integral at the places above $\ell$ (that is, $\delta x_w\in L_w$ for $w\mid\ell$). The
standard data of \S\ref{pk:ss-kernel} are those of level $\aubar{\pi_M}$. \emph{If $g\in\auU(W)(\auA_f)$ is integral with integral inverse at the places above $\ell$, then
$\omega(g)$ preserves the standard data of level $\delta$ at $\ell$, for every $\delta$ and every $\ell$.}
Indeed, by (W1), $\omega(g)\psrho^\circ(x)\psvac^\circ$ is a root of unity times $\psrho^\circ(gx)\,
\omega(g)\psvac^\circ$, and $\delta gx$ is integral above $\ell$ when $\delta x$ is. By (W4),
$\omega(g)\psvac_L=\delta_g\sum_{r\in gL/(L\cap gL)}\zeta_r\psrho^\circ(r)\psvac_L$ with roots of unity
$\zeta_r$, where $\delta_g$ is a root of unity times $[gL:L\cap gL]^{-1/2}$ times the value of
$\upsilon^{\pm1}$ at an idele $\det\mathsf u$ with $\mathsf uL=L\cap g^{-1}L$. As $g_w L_w=L_w$ for $w\mid\ell$,
the index is prime to $\ell$, $\mathsf u$ may be taken to be $1$ above $\ell$ (so the value of the
algebraic Hecke character $\upsilon$ is a unit at $\ell$), and the $r$ may be taken in $L_w$ above
$\ell$, so that $\delta r$ is integral there ($\delta\in\auO_K$). This is the proof of
Proposition~\ref{pk:std}(S1), which does not use the level.

\begin{theorem}[Standard data at every prime]\label{pf:std}
Let $\ell$ be a prime, $\delta\in\auO_K$, $\delta\ne0$, and $f$ a standard datum of level $\delta$
at $\ell$. There are a resolution of the identity $(u'_j,m_j)_j$ of $E^{\mathsf C}$ with
$\alpha^{-1}u'_j$ integral above $\ell$, and an integral power series $\mathsf H$ in the series
$y(\mu_j(z))$, $\mu_j=(\varpi_{\mathrm{vac}}/\mathsf u)\,u_j'^{\dagger}$, such that
\[
\Theta_f\bigl(\aubar\delta\,s_0\,z\bigr)=\Bigl(\prod_j\theta(\mu_j(z))^{m_j}\Bigr)^{\auNm\delta}
\cdot\mathsf H(z),
\]
and the linear forms $\mu_j$ have integral coefficients. If $\ell\ne d$, the series
$\Theta_f(\aubar\delta s_0\,\cdot)$ is $\partial$-integral, and so is $\Theta_f(c\,\aubar\delta s_0
\,\cdot)$ for every $c\in\pao$.
\end{theorem}
\begin{proof}
Let $(v_i)$ be the frame $(\zeta_{q_1}^k)$, $\mathsf b=q_1$, if $\ell\nmid q_1$, and the frame of
Lemma~\ref{pf:frame}, $\mathsf b=\mathsf h$, if $\ell\mid q_1$; so $\mathsf b$ is a unit at
$\ell$. By Lemma~\ref{pf:fourier}(b), $f$ is a combination with integral coefficients of vectors
$\psrho_f(x')F^v$ with $\delta x'$ integral above $\ell$; each of them is a root of unity times
$\psrho_f(\gamma)F^v$ with $\gamma\in W(E)$, $\gamma-x'\in\mathsf bL$ (Lemma~\ref{pf:semi}), and
then $\delta\gamma$ is integral above $\ell$. Let $\mathsf m\ge1$ be an integer prime to $\ell$
with $\mathsf m\delta\gamma\in\Lambda_L$ for all the $\gamma$ that occur, and $\delta'=\mathsf m
\mathsf b\delta$. In \eqref{pf:eq-transl}, $\delta'\mathsf x_{i,\mathsf c}=-\pfT_{E/K}(\aubar{v_i}
\cdot\mathsf m\delta\gamma_{\mathsf c})\in\auO_K$, so $\delta'\cdot\mathsf x_{i,\mathsf c}\Omega_0
\in\Lambda$. Since $s_0\Omega_0/c_\psi=\varpi_{\mathrm{vac}}/\mathsf u=:\mathsf p$, we have
$a_{i,\mathsf c}(\aubar\delta s_0z)=\aubar{\delta'}\cdot\mathsf p\,u_i'^{\dagger}(z_{\mathsf c})$
with $u'_i=v_i/(\mathsf m\mathsf b^2)$, and Proposition~\ref{pf:pull} gives, for each factor,
\[
\theta_{\mathsf x_{i,\mathsf c}\Omega_0}\bigl(a_{i,\mathsf c}(\aubar\delta s_0z)\bigr)=\theta
\bigl(\mu_{i,\mathsf c}(z)\bigr)^{\mathsf m^2\mathsf b^2\auNm\delta}\cdot\Gamma\bigl(y(\mu_{i,
\mathsf c}(z))\bigr),\qquad\mu_{i,\mathsf c}(z)=\mathsf p\,u_i'^{\dagger}(z_{\mathsf c}),\quad
\Gamma\in\pao[[Y]].
\]
Taking the product over $(i,\mathsf c)$ with the exponents $\mathsf b$ we obtain the formula for
$\psrho_f(\gamma)F^v$, with the multiplicities $m_{i,\mathsf c}=\mathsf m^2\mathsf b^3$; the first
factor does not depend on $\gamma$, so the formula holds for $f$. The vectors $u'_i$ with these
multiplicities are a resolution of the identity: $\sum_i\mathsf m^2\mathsf b^3\tau(v_i)\aubar{\tau'
(v_i)}/(\mathsf m\mathsf b^2)^2=\delta_{\tau\tau'}$. The coefficients of $\mu_{i,\mathsf c}$ are
$\varpi_{\mathrm{vac}}\aubar{\tau(v_i)}/(\mathsf u\,\mathsf m\mathsf b^2)$; they are integral by
Lemma~\ref{pf:scale} and Corollary~\ref{pf:CS}. For $\ell\ne d$ the last assertion follows as in
Proposition~\ref{pf:pull}.
\end{proof}

\begin{corollary}[{\rm(K1)} at every prime; the prime $N$]\label{pf:K1}
\begin{enumerate}
\item[(a)] The integrality assertion of Lemma~\ref{pa:kernel}{\rm(K1)} holds for every prime
$\ell\ne d,N$ with $\ell\nmid\auNm Q_\vartheta$, in particular for $\ell=2$, for the primes
dividing $q_1$, and for $p_0,p_1$.
\item[(b)] Lemma~\ref{pl:kernelN} $(=$ {\rm(K1$_N$)}$)$ holds as stated, without the condition
{\rm(Gen)}.
\item[(c)] Proposition~\ref{pk:std}{\rm(S3)}, hence Lemma~\ref{pa:kernel}{\rm(K2)}, holds without
appeal to \cite{KubertLang}.
\end{enumerate}
\end{corollary}
\begin{proof}
(a) By the proof of Proposition~\ref{pk:K12}, $F=\Theta_f(\pi_Ms_0\,\cdot)$ for a unit multiple $f$
of a standard datum of level $\aubar{\pi_M}$ at $\ell$; apply Theorem~\ref{pf:std} with
$\delta=\aubar{\pi_M}$. (b) Let $\ell=N$ and let $\mathfrak n$ be the prime of $K$ induced by
$\iota$; it is the prime below the $\mathfrak l_\tau$, and $\mathfrak n\ne\aubar{\mathfrak n}$.
Choose $\delta\in\mathfrak n\setminus(\mathfrak n^2\cup\aubar{\mathfrak n})$ (possible by the Chinese remainder theorem), so that $\operatorname{ord}_{\mathfrak n}\delta=1$. Then $\iota(\delta)\in\pam$ and
$\iota(\aubar\delta)$ is a unit. The point $z$ of the lemma is integral at the primes above
$\aubar{\mathfrak n}$ and $Nz$ is integral at the primes above $\mathfrak n$, where $N$ and $\delta$
are both uniformisers ($N$ splits completely in $E$, and $\operatorname{ord}_{\mathfrak n}\delta=1$); so $\delta z$ is integral above $N$ and
$\phi^\circ_{j,z}$ is a standard datum of level $\delta$ at $N$. The reduction of
Proposition~\ref{pk:K12}, with $\mathsf k,\mathsf k'$ integral above $N$ ((G2) of
\S\ref{pl:ss-norm}) and (S1) at the level $\delta$, shows that the matrix in question is, up to
matrices invertible over $\pao$, the matrix of the jets along the entries of $S$ of
$\Theta_f(\pi_Ms_0\,\cdot)$ for a unit multiple $f$ of a standard datum of level $\delta$. By
Theorem~\ref{pf:std} with $c=\iota(\pi_M/\aubar\delta)\in\pao$ ($N\ne d$) this series is
$\partial$-integral, and the entries of $S$ are integral polynomials. (c) In the proof of
Proposition~\ref{pk:std}, \cite{KubertLang} is used only for the integrality of $\theta_u(0)$
(through (C3)(i)) and of $\theta_t'(0)$. The first is Proposition~\ref{pf:values}. For the second:
by Lemma~\ref{pf:chi}(c), $\theta_{-t}=\pm\theta_t$, so $\theta_t^2/\theta^2$ is $\Lambda$-periodic
with divisor $2[P_0]-2[O]$ and $\theta_t^2=\mathsf c_t\,x\,\theta^2$, $\mathsf c_t=\theta_t'(0)^2$.
Let $z\ne0$ be a torsion point of odd order prime to $\ell$. Multiplying by $e^{-H(z,z)}$ gives
$\zeta\,\pfbr{z+t}^2=\mathsf c_t\,x(z)\,\pfbr z^2$ with a root of unity $\zeta$. Here $z-t\notin
E_1$ and $z\notin E_1$, so $\pfbr z$ and $\pfbr{z+t}$ are units (Proposition~\ref{pf:values}), and
$x(z)$ is a unit because $Q(z)$ reduces neither to the origin nor to $\widetilde P_0$. Hence
$\theta_t'(0)$ is a unit, at every prime.
\end{proof}

\subsubsection{The fact {\rm(W6)}}\label{pf:ss-flip}

We use (W1), (W4) and (W2$'$) of \S\ref{pk:ss-conv}; the last one also in its local form: for a
finite place $v$ of $E^+$ and an isotropic $E_v$-subspace $I$ of $W_v=W\otimes E^+_v$, the local
representation $\psHS_v(W)$ has the zero term $\psZ_I\colon\psHS_v(W)\to\psHS_v(W_I)$, with
$\psZ_I\psrho(w)=\psrho(\bar w)\psZ_I$ for $w\in I^\perp$ and $\psZ_I\omega(p)=\chi_I(p)\omega_{W_I}
(\bar p)\psZ_I$, $\chi_I(p)=(\upsilon_v|\cdot|^{-1/2})(\det p|_I)^{-1}$, for $p$ in the stabiliser
$P_I(E^+_v)$. (This is how (W2) is obtained: the mixed model $\auS(I^*)\otimes\psHS_v(W_I)$ and
the formulas for $P_I$ are local, see the references given for Lemma~\ref{ps:weil}.) The map
$\psZ_I$ is, up to a scalar, the only non-zero continuous linear map with the first property:
such a map kills $(\psrho(i)-1)\psHS_v(W)$ for $i\in I$, hence factors through the evaluation at
$0$ of the mixed model, and $\psHS_v(W_I)$ is irreducible.

\begin{proposition}[{\rm(W6)}]\label{pf:flip}
Let $W$ be a hermitian space over $E$ that has a self-dual lattice $L\subset W(\auA_f)$, and let
$W'$ be the $E$-vector space $\aubar W$ $(W$ with the conjugate scalar multiplication$)$ with the
form $\langle w_1,w_2\rangle'=-\langle w_2,w_1\rangle$. Then $\beta'=\beta$ on $W(\auA)=W'(\auA)$,
$\auU(W)=\auU(W')$ as groups acting on this set, and if $\mathrm{Fl}\colon\psHS(W)\to\psHS(W')$ is
an isomorphism with $\mathrm{Fl}\,\psrho(w)=\psrho'(w)\,\mathrm{Fl}$ (it exists and is unique up
to a scalar), then
\[
\mathrm{Fl}\circ\omega_W(g)=\upsilon(\det\nolimits_Wg)^{-1}\cdot\omega_{W'}(g)\circ\mathrm{Fl}
\qquad(g\in\auU(W)(\auA)),
\]
where $\det_W$ is the determinant on the $E$-module $W$, and $\theta_{W'}\circ\mathrm{Fl}$ is a
multiple of $\theta_W$. For $W=\auHom(V,V')$ the map $w\mapsto w^\dagger$ identifies $W'$ with
$\auHom(V',V^-)$, and $\det_W(g,g')=\det(g')^{n}\det(g)^{-n'}$; this is {\rm(W6)} of
\S\ref{pk:ss-conv}.
\end{proposition}
\begin{proof}
The form $\langle\,,\rangle'$ is hermitian on $W'$, and $\beta'(x,y)=\psi(\tfrac12\jmath\langle
y,x\rangle')=\psi(-\tfrac12\jmath\langle x,y\rangle)=\psi(\aubar{\tfrac12\jmath\langle y,x\rangle})
=\beta(x,y)$, because $\psi\circ c=\psi$. So $\psrho'$ is another
realisation of the same Heisenberg representation and $\mathrm{Fl}$ exists; $\theta_{W'}\circ
\mathrm{Fl}$ is invariant under $\psrho(W(E))$, hence equals $c_\theta\theta_W$ with $c_\theta\ne0$.
By (W1) for $W$ and $W'$ the operator $\mathrm{Fl}^{-1}\omega_{W'}(g)^{-1}\mathrm{Fl}\,\omega_W(g)$
commutes with all $\psrho(w)$, so it is a scalar $\kappa(g)$, and $\kappa$ is a continuous
character of $\auU(W)(\auA)$. Put $\kappa_0(g)=\kappa(g)\upsilon(\det_Wg)$; we show $\kappa_0=1$.

\emph{Rational points.} For $\gamma\in\auU(W)(E^+)$, $c_\theta\theta_W(F)=c_\theta\theta_W(\omega_W
(\gamma)F)=\kappa(\gamma)\theta_{W'}(\omega_{W'}(\gamma)\mathrm{Fl}F)=\kappa(\gamma)c_\theta\theta_W
(F)$, so $\kappa(\gamma)=1$; and $\upsilon(\det\gamma)=1$ since $\upsilon$ is a Hecke character.

\emph{The stabiliser of $L$.} $L$ is self-dual in $W'$ as well, and $\epsilon'_L(l)=\psi_f(-\frac12\omega\jmath\langle l,l\rangle')=\epsilon_L(l)^{-1}=\epsilon_L(l)$
by Lemma~\ref{pf:eps}. So $\mathrm{Fl}\,\psvac_L\in\C^\times\psvac'_L$, and by (W4), $\kappa(k)=1$ for $k$
in the stabiliser $K_L$ of $L$; and $\upsilon(\det k)=1$ because $\det k$ is a unit idele and
$\upsilon$ is unramified.

\emph{Parabolic elements.} Let $v$ be a finite place, $I\subset W_v$ isotropic and $p\in P_I(E^+_v)$
acting trivially on $I^\perp/I$. The same subset $I$ is an isotropic subspace of $W'_v$, with
$\det(p|_I)$ replaced by its conjugate $\aubar{\mathsf d}$, $\mathsf d=\det\nolimits_W(p|_I)$. By
the uniqueness of the zero term, $\psZ'_I\mathrm{Fl}=c\,\mathrm{Fl}_I\psZ_I$ with $c\ne0$, where
$\mathrm{Fl}_I$ is the map $\mathrm{Fl}$ of $W_I$. Hence, by (W2$'$) for $W$ and $W'$,
\[
\chi_I(p)\,\psZ'_I\mathrm{Fl}=\psZ'_I\mathrm{Fl}\,\omega_W(p)=\kappa(p)\,\psZ'_I\omega_{W'}(p)
\mathrm{Fl}=\kappa(p)\chi'_I(p)\,\psZ'_I\mathrm{Fl},
\]
so $\kappa(p)=\chi_I(p)/\chi'_I(p)=\upsilon(\aubar{\mathsf d})/\upsilon(\mathsf d)=\upsilon(\mathsf
d)^{-2}$, because $\upsilon(\mathsf d\aubar{\mathsf d})=1$ ($\upsilon$ restricted to $\auA_{E^+}^
\times$ is the character of $E/E^+$, trivial on norms). On the other hand $p$ acts on $W/I^\perp$,
the conjugate dual of $I$, with determinant $\aubar{\mathsf d}^{-1}$; so $\det_Wp=\mathsf d/\aubar{
\mathsf d}$ and $\upsilon(\det_Wp)=\upsilon(\mathsf d)^2$. Thus $\kappa_0(p)=1$.

\emph{All finite places.} Let $v$ be a finite place; $E/E^+$ is unramified at $v$. If $v$ is
inert, $L_v$ has an orthonormal basis: if $\langle x,y\rangle$ is a unit and $\langle x,x\rangle$,
$\langle y,y\rangle$ are not, then $x+\lambda y$ has unit norm when $\pfT(\lambda\langle x,y
\rangle)$ is a unit, and the trace of $\auO_{E,v}$ is surjective; a vector of unit norm can be
scaled to norm $1$, the norm of units being surjective; split it off and induct. Let $I_L$ be a
maximal totally isotropic direct summand of $L_v$. By (R1) of \S\ref{pk:ss-rest},
$L_v=(I_L\oplus I_L^*)\perp L_0$ with $I_L^*$ totally isotropic, and $L_0$ has rank $\le1$: a
unimodular lattice of rank $\ge2$ contains the primitive isotropic vector $e_1+\lambda e_2$,
$\lambda\aubar\lambda=-1$, for orthonormal $e_1,e_2$ (if $v$ splits, $L_0=0$). Let $I=E_vI_L$. By
the Iwasawa decomposition ((R1); \cite[\S3.3.2, p.~51]{TitsCorvallis}), $\auU(W)(E^+_v)=P_I(E^+_v)K_{L,v}$,
and $P_I(E^+_v)$ is generated by the elements acting trivially on $I^\perp/I$ and by
$\auU(E_vL_0)$, which consists of unit scalars on $L_0$ and lies in $K_{L,v}$. So $\kappa_0=1$ on
$\auU(W)(E^+_v)$, and by continuity on $\auU(W)(\auA_f)$.

\emph{The infinite places.} $\kappa_0$ is now a continuous character of $\auU(W)(E^+\otimes\R)$
that is trivial on the image of $\auU(W)(E^+)$, which is dense by Proposition~\ref{pl:wa}. So
$\kappa_0=1$. The last assertions are those of \S\ref{pk:ss-conv}: $\langle w_1^\dagger,w_2^\dagger
\rangle=-\langle w_2,w_1\rangle$ and $w\mapsto w^\dagger$ is antilinear.
\end{proof}
The hypothesis on $L$ holds for all spaces $\auHom(V,V')$ to which (W6) is applied (all spaces of
\S\ref{ss:au-spaces} carry unimodular lattices $\psM$, and $\alpha\auHom(\psM,\psM')$ is
self-dual).

\begin{remark}[The lattice preserved by the flip]\label{pf:rem-flip-lattice}
The proof shows that $\mathrm{Fl}$ carries the vacuum of $L$ to a multiple of the vacuum of the
\emph{same} set $L$, viewed in $W'$. For $W=\auHom(V,V')$ and $L=\alpha\auHom(\psM,\psM')$ the
identification $w\mapsto w^\dagger$ of $W'$ with $\auHom(V',V^-)$ carries $L$ to
$L^\dagger=\aubar\alpha\auHom(\psM',\psM)=u\,\alpha\auHom(\psM',\psM)$, $u=\aubar a/a$ (with $a$ as in
\S\ref{ps:ss-weil}; Lemma~\ref{pgx:dict}(f)), and not, in general, to the standard lattice
$\alpha\auHom(\psM',\psM)$ of $\auHom(V',V^-)$. Here $u$ is a finite idele with $u\aubar u=1$, acting on
$\auHom(V',V^-)$ through the centre of $\auU(V^-)=\auU(V)$, and $u_v=1$ at every place $v$ not above
$q_1$. By Lemma~\ref{pgx:dict}(e),
\[
\mathrm{Fl}\,\psvac_{\alpha\auHom(\psM,\psM')}\in\C^\times\,\omega_{\auHom(V',V^-)}(1,u)\,
\psvac_{\alpha\auHom(\psM',\psM)}
\]
($\omega(1,u)$ meaning $u$ acting through $\auU(V^-)$). So $\mathrm{Fl}$ may be normalised to send
vacuum to vacuum at the places not above $q_1$, but at the places above $q_1$ it sends the standard
vacuum to the translate by $\omega(1,u)$ of the standard vacuum. In an integral over $[\auU(V)]$
this translate changes the integral only by the factor of modulus one of Lemma~\ref{pgx:dict}(e); in a
kernel at a fixed point it shifts the argument in $\auU(V)$ by $u$ (as in Proposition~\ref{ptv:rev} and
Lemma~\ref{pgu:line}).
\end{remark}

\subsubsection{The orbits at the end of Lemma~\ref{ps:support}}\label{pf:ss-orbits}

\begin{lemma}\label{pf:orbits}
Let $1\le j<m$ and let $v$ be a finite place of $E^+$. Let $\mathcal X_v$ be the set of the
lattices $\Lambda'_v=y\psM_{j,v}$, $y\in G_j(E^+_v)$, that are the direct sum of their intersections
with the grades and satisfy the bounds of Lemma~\ref{ps:support} at $v$ $(\mathsf q^{\,j}\Lambda'^-\subseteq
\psM_j^-$, $\aubar{\mathsf q}^{\,j}\psM_j^-\subseteq\Lambda'^-$, with $\Lambda'^-$ the intersection
with the span of the whole string$)$. Then $\auU(X_j)(E^+_v)$ has finitely many orbits in
$\mathcal X_v$, and exactly one, that of $\psM_{j,v}$, if $v$ does not divide $q=\auNm Q$. Hence
$\mathfrak B_j=R_j(\auA_f)\mathsf C_jK_j$ with $\mathsf C_j$ finite, consisting of elements that are
$1$ outside $q$, and $\mathsf C_j=\{1\}$ if $Q=\emptyset$.
\end{lemma}
\begin{proof}
Let $\Lambda'=\Lambda'_v\in\mathcal X_v$; it is unimodular. For a grade $g\ne0$ its part
$\Lambda'_g$ is a lattice in the line $E_vf_a$, between $\aubar{\mathsf q}^{\,j}\psM_{j,g}$ and
$\mathsf q^{-j}\psM_{j,g}$: finitely many possibilities, and $\Lambda'_g=\psM_{j,g}$ if $v\nmid q$.
The sum of the parts of non-zero grade is unimodular (it is so for $\psM_j$ and, in general, it is
a sum of pairs of dual lines; unimodularity of $\Lambda'$ forces $\Lambda'_{-g}$ to be the dual of
$\Lambda'_g$), so $\Lambda'=\Lambda'_0\perp\bigoplus_{g\ne0}\Lambda'_g$ with $\Lambda'_0$
unimodular in the space $Y_{j,0}$ of grade $0$. The group $\auU(X_j)$ acts on $Y_{j,0}$ and fixes
the string; so we must classify the $\Lambda'_0$.

\emph{$j$ odd.} $Y_{j,0}=X_j$, and the unimodular lattices of $X_{j,v}$ form one orbit: if $v$
splits, $v=w\aubar w$, they are the pairs $(B,B^\#)$ with $B$ a lattice in $X_{j,w}$, and
$\auU(X_j)(E^+_v)=\auGL(X_{j,w})$ is transitive on lattices; if $v$ is inert, every unimodular
lattice has an orthonormal basis (proof of Proposition~\ref{pf:flip}), and two such lattices in
the same space differ by an isometry.

\emph{$j$ even.} $Y_{j,0}=X_j\perp E z$, $\langle z,z\rangle=-N$. Put $\mathfrak z=\Lambda'_0\cap
E_vz$; the projection of the unimodular lattice $\Lambda'_0$ to $E_vz$ is the dual $\mathfrak
z^\#$ of $\mathfrak z$ in $E_vz$. The bounds leave finitely many possibilities for $\mathfrak z$,
and $\mathfrak z=\psM_j\cap E_vz$ if $v\nmid q$ (the bounds then say $\Lambda'^-=\psM_j^-$). If $v\nmid Nq$, then $\mathfrak z$ is unimodular, so
$\Lambda'_0=(\Lambda'_0\cap X_j)\perp\mathfrak z$ with $\Lambda'_0\cap X_j$ unimodular, and we
conclude as for $j$ odd. Let $v=w\aubar w$ be split (the places above $N$ and above $q$ are). Then
$\Lambda'_0=(B,B^\#)$ with $B$ a lattice in $X_{j,w}\oplus E_wz$, and the conditions are $B\cap
E_wz=\mathfrak z_w$, $\operatorname{pr}_z(B)=\mathfrak z'_w:=(\mathfrak z_{\bar w})^\#$. Such a
$B$ is the same as a triple $(B_1,B_2,\phi)$: lattices $B_1=B\cap X_{j,w}\subseteq B_2=
\operatorname{pr}_X(B)$ and an isomorphism $\phi\colon B_2/B_1\to\mathfrak z'_w/\mathfrak z_w$, with
$B=\{(b,\xi)\in B_2\oplus\mathfrak z'_w:\phi(b)\equiv\xi\}$. The module $\mathfrak z'_w/\mathfrak
z_w$ is cyclic, $\cong\auO_w/\varpi^e$, and $\auGL(X_{j,w})$ is transitive on the pairs
$B_1\subseteq B_2$ with $B_2/B_1\cong\auO_w/\varpi^e$ (elementary divisors); the stabiliser of a
pair maps onto $\operatorname{Aut}(B_2/B_1)=(\auO_w/\varpi^e)^\times$ (choose a basis of $B_2$ whose
last vector generates $B_2/B_1$ and whose other vectors lie in $B_1$, and multiply the last vector
by a unit). So the triples with given $\mathfrak z_w\subseteq\mathfrak z'_w$ form one orbit. For
$v\mid N$, $v\nmid q$, the pair $(\mathfrak z_w,\mathfrak z'_w)$ is that of $\psM_j$ (and $e=1$);
for $v\mid q$ there are finitely many pairs.

\emph{The set $\mathfrak B_j$.} For each $v\mid q$ choose representatives $c_v\psM_{j,v}$,
$c_v\in G_j(E^+_v)$, of the orbits in $\mathcal X_v$, and let $\mathsf C_j$ be the set of the
elements $(c_v)_{v\mid q}$, completed by $1$. If $x\in\mathfrak B_j$ and $r\in R_j(\auA_f)$ is as
in the definition, then $rx\psM_j\in\prod_v\mathcal X_v$, so $u\,rx\psM_j=c\psM_j$ for some $u\in\auU(X_j)(\auA_f)\subset R_j(\auA_f)$ and
$c\in\mathsf C_j$, and $x\in R_j(\auA_f)cK_j$. Conversely $rck\psM_j=r(c\psM_j)$.
\end{proof}

\subsubsection{The number of classes}\label{pf:ss-classes}

\begin{proposition}\label{pf:classes}
Let $X=E^n$ be a standard space with $n\le n_m$, $G=\auU(X)$, $\Gamma=G(E^+)$, $K^\circ$ the
stabiliser of $\widehat\auO_E^{\,n}$ in $G_f$, and for an integer $\mathsf G\ge1$ let $K(\mathsf G)$
be the group of the $k\in K^\circ$ with $k\equiv1$ modulo $\mathsf G$. Then
\[
\#\bigl(\Gamma\backslash G_f/K^\circ\bigr)\le(2\mathsf B+1)^{4fn^2},\qquad\mathsf B=n^2C_{\mathrm{sb}}^2
\cdot2f\sqrt d,\quad C_{\mathrm{sb}}=6q_1d(fn)^2,
\]
and $\#(\Gamma\backslash G_f/J)\le(2\mathsf B+1)^{4fn^2}\mathsf G^{2fn^2}$ for every subgroup
$J\supseteq K(\mathsf G)$ of $G_f$. For $\mathsf G=MN$ this is at most $\exp(x^{30})$. In
particular the levels $J_{X_j}$ $(j\le m)$, and the level $J_1$ of the target forms for the trivial
twist, have at most $\exp(x^{30})$ classes.
\end{proposition}
\begin{proof}
\emph{Two elementary facts.} (i) Every $\mathsf y\in\auA_{E,f}$ is congruent modulo
$\widehat\auO_E$ to an $\mathsf a\in E$ with $|\tau(\mathsf a)|\le\mathsf B_0:=2f\sqrt d$ for all
embeddings $\tau$: indeed $\auA_{E,f}=E+\widehat\auO_E$, and $\auO_E=\auO_K[\zeta_{q_1}]$ has the
$\Z$-basis $\omega^e\zeta_{q_1}^k$ $(e\le1$, $k<f)$, whose conjugates have absolute value at most
$|\omega|\le\sqrt d$; reduce the coordinates of $\mathsf a$ to $[0,1)$. (ii) For $B>0$ the number
of $\mathsf a\in\auO_E$ with $|\tau(\mathsf a)|\le B$ for all $\tau\in\auPl$ is at most
$(2B+1)^{2f}$: two distinct such elements have $\prod_\tau|\tau(\mathsf a-\mathsf a')|^2=\auNm(
\mathsf a-\mathsf a')\ge1$, so the polydiscs of radius $\frac12$ around the points $(\tau(\mathsf
a))_\tau\in\C^{\auPl}$ are disjoint, and they lie in the polydisc of radius $B+\frac12$.

\emph{The invariant of a class.} Let $g\in G_f$. By Lemma~\ref{au:shortbasis} with $c_0=q_1d$ there
is $Y\in\auGL_n(E)$ with $A:=g^{-1}Y\in M_n(\widehat\auO_E)$ and all conjugates of all entries
bounded by $C_{\mathrm{sb}}$. By (i) there is $X\in M_n(E)$ with $A^{-1}-X\in M_n(\widehat\auO_E)$
and entries bounded by $\mathsf B_0$. Put $\mathsf G_g=Y^*Y$ and $\mathsf I(g)=(\mathsf G_g,\mathsf
G_gX)$. As $g$ is unitary, $\mathsf G_g=A^*A$ and $\mathsf G_gX=A^*\bigl(1-A(A^{-1}-X)\bigr)$ have
entries in $E\cap\widehat\auO_E=\auO_E$, and their conjugates are bounded by $nC_{\mathrm{sb}}^2$
and by $n\cdot nC_{\mathrm{sb}}^2\cdot\mathsf B_0=\mathsf B$. By (ii) the number of possible
$\mathsf I(g)$ is at most $(2\mathsf B+1)^{2f\cdot2n^2}$.

\emph{The invariant determines the class.} Let $g,g'$ with data $(Y,A,X)$, $(Y',A',X')$ and
$\mathsf I(g)=\mathsf I(g')$. Then $Y^*Y=Y'^*Y'$, so $\gamma=YY'^{-1}\in\Gamma$; and $X=X'$, so
$A'A^{-1}=A'(A^{-1}-X)+A'X'$ is integral, since $A'X'=1-A'(A'^{-1}-X')$. But $A'A^{-1}=g'^{-1}\gamma^{-1}g=:k$ is unitary, so
$k^{-1}=k^*$ is integral too, $k\in K^\circ$, and $g=\gamma g'k$. This proves the first bound.

\emph{Levels.} The map $\Gamma\backslash G_f/(J\cap K^\circ)\to\Gamma\backslash G_f/J$ is onto, and
the fibres of $\Gamma\backslash G_f/(J\cap K^\circ)\to\Gamma\backslash G_f/K^\circ$ have at most
$[K^\circ:J\cap K^\circ]\le[K^\circ:K(\mathsf G)]\le\#M_n(\auO_E/\mathsf G)=\mathsf G^{2fn^2}$
elements.

\emph{Sizes.} By \eqref{eq:au-sizes} and the definitions, $f\le x^{12}$, $n\le n_m\le40x^6$, and
$6q_1d(fk)^2\le x^{100}$ (the inequality used in the proof of Proposition~\ref{au:short}), and
$C_{\mathrm{sb}}=6q_1d(fn)^2\le x^{102}$ (Lemma~\ref{pz:sizes}(b)), $\mathsf B_0\le C_{\mathrm{sb}}$, $n\le
C_{\mathrm{sb}}$ and $2\mathsf B+1\le C_{\mathrm{sb}}^6$; hence $4fn^2\log(2\mathsf B+1)\le4\cdot
x^{12}\cdot1600x^{12}\cdot612\log x\le x^{26}$. Also $\log(MN)\le22\log b+30m\log b\le x$ by
\eqref{eq:au-scales}, so $2fn^2\log(MN)\le3200x^{25}\le x^{26}$. The levels named contain $K(MN)$:
they are defined inside $K^\circ$ by congruences modulo the chosen places, which divide $M$, and
modulo primes dividing $N$ (\S\ref{ss:au-spaces}).
\end{proof}

\subsubsection{The published results used in this subsection}\label{pf:ss-status}
Proposition~\ref{pk:curve}(C3)(ii) is Theorem~\ref{pf:recip}; Lemma~\ref{ps:conj} and
Proposition~\ref{au:kernel} follow (Corollary~\ref{pf:conj}); Lemma~\ref{pa:kernel} holds in full,
(K1) at every prime $\ell\ne d,N$ and (K1$_N$), that is, Lemma~\ref{pl:kernelN}, without restriction
(Corollary~\ref{pf:K1}); (W6) is Proposition~\ref{pf:flip}; the number of classes is bounded by
Proposition~\ref{pf:classes}. The published results used are: integrality of
$j$ and potential good reduction, Deuring's criterion, the formal group of a Weierstrass equation
and the kernel of reduction~\cite{SilvermanAEC}, \cite[Chapter~13, Theorem~12]{LangEF}; the oddness of the class number of
\cite{CoxPrimes}; Euler's formula for the dual basis~\cite[III, \S6, Lemma~2]{SerreLF}; the Iwasawa
decomposition~\cite[\S3.3.2, p.~51]{TitsCorvallis}; and the facts (W1), (W2$'$) (also in its local form), (W4) on
the Weil representation.



%% ===== proofs/A-closure.tex =====
% proofs/A-closure.tex
% Closing statements for Theorem A: the archimedean types of the lower chain (the assumption made
% in ps:ss-Y), the constant C_0 and the definition of admissible data, the levels of the upper
% chain and the sizes used for the number of classes, the estimate left open in the proof of
% au:short, and the list of what Theorem A depends on.
% Labels of this file start with pz:. Only amsmath/amssymb/amsthm are used.
\providecommand{\auO}{\mathcal O}
\providecommand{\auU}{\mathrm U}
\providecommand{\auGL}{\mathrm{GL}}
\providecommand{\auHom}{\operatorname{Hom}}
\providecommand{\auA}{\mathbb A}
\providecommand{\auF}{\mathbb F}
\providecommand{\auPl}{\Phi}
\providecommand{\auPa}{\Phi_a}
\providecommand{\auS}{\mathcal S}
\providecommand{\auSram}{\mathcal S^{\mathrm{ram}}}
\providecommand{\auAz}{\mathcal A_0}
\providecommand{\auG}{\overline G}
\providecommand{\auP}{\mathcal P}
\providecommand{\auNm}{\operatorname{Nm}}
\providecommand{\auAut}{\operatorname{Aut}}
\providecommand{\aubar}[1]{\overline{#1}}
\providecommand{\aueu}{\varepsilon_{\mathrm u}}
\providecommand{\aurd}{\mathrm{rd}_E}
\providecommand{\psM}{\mathsf M}
\providecommand{\psHS}{\mathcal H}
\providecommand{\psrho}{\varrho}
\providecommand{\pzF}{\mathfrak F}

\subsection{Types, constants and levels for Theorem~A}\label{pz:sec}

This subsection proves the remaining inputs of Theorem~\ref{thm:srcA} not treated in
\S\S\ref{pa:sec}--\ref{pf:sec}. (1) In \S\ref{ps:ss-Y} it was taken ``as part of the definition of the lower
chain'' that the lifts $\Theta^Y_{j,\vartheta}$ and their transposes preserve the archimedean type
and that the forms concerned are cusp forms. In \S\ref{pz:ss-types} the types are defined, and the
assumption is \emph{reduced} to three statements about discrete series of real unitary groups,
(A1)--(A3) below; everything global is proved (Proposition~\ref{pz:types}). Statements (A1) and
(A2) are theorems of Harish-Chandra, Hecht--Schmid and Mili\v ci\'c; (A3) is a case of the
correspondence of discrete series under the theta correspondence. They are treated in \S\ref{par:sec} and \S\ref{pco:sec}: formula \eqref{pds:eq-FJ} is proved (Theorem~\ref{pco:FJ}) and the normalisation (V) is proved from the local properties (W1$_\infty$), (W2$_\infty$), (M), (S) of the archimedean Weil representation (Theorem~\ref{pco:V}). (2) \S\ref{pz:ss-const} makes the constant $C_0$ of
Proposition~\ref{au:kernel} explicit and settles its relation to the definition of admissible data.
(3) \S\ref{pz:ss-levels} defines the levels of the spaces $X_j$, checks the two facts quoted in the
proof of Proposition~\ref{pf:classes}, and proves the estimate used in the proof of
Proposition~\ref{au:short}. (4) \S\ref{pz:ss-audit} lists every statement on
which Theorem~\ref{thm:srcA} depends, with the place of its proof.

\subsubsection{Forms of a discrete series type and the lower chain}\label{pz:ss-types}

\emph{The groups and the parameters.} Let $0\le j\le m$, $G_j=\auU(Y_j)$, $N_j=r+jt$ its rank,
$G_{j,\infty}=G_j(E^+\otimes\R)=\prod_{\tau\in\auPl}G_{j,\tau}$, and $K_{j,\infty}$ the maximal compact
subgroup that stabilises, at every $\tau$, the positive and the negative part of an orthogonal
basis of $Y_j\otimes_{E,\tau}\C$. Through $\tau$ the group $G_{j,\tau}$ is the unitary group of a
hermitian form with $N_j-Q_j$ signs $+$ and $Q_j=\lceil(m-1-j)/2\rceil$ signs $-$
(\S\ref{ss:au-spaces}). For $x\in G_{j,\infty}$ let $\|x\|=\prod_\tau\|x_\tau\|$, the product of the
operator norms on $Y_j\otimes_{E,\tau}\C$ for the positive form attached to that basis; then
$\|x\|\ge1$ and $\|x^{-1}\|=\|x\|$. Put $\kappa_l=\kappa+(l+\tfrac12)\Delta$ as in
\S\ref{ss:au-notation}. The \emph{parameter} of $Y_j$ at $\tau$ is the following set
$\Lambda^{(j)}_\tau$ of $N_j$ numbers, each with a sign:
\begin{equation}\label{pz:eq-param}
\begin{aligned}
&\tfrac k2-1,\ \tfrac k2-2,\ \dots,\ 1-\tfrac k2&&(k-1\text{ numbers, sign }+);\\
&-\kappa(\tau)-i\Delta\quad(1\le i\le m)&&(\text{sign }-\text{ if }i\ge j+2\text{ and }i\equiv j\!\!\pmod 2,\text{ else }+);\\
&-\kappa_l(\tau)+\tfrac{t-1}2-b\quad(1\le l\le j,\ 0\le b<t)&&(\text{sign }+);
\end{aligned}
\end{equation}
all numbers being lowered by $\tfrac12$ if $j$ is odd. For $j=0$ this is the parameter
$\varsigma_\tau$ of \S\ref{ss:au-spaces}; the number of signs $-$ is $Q_j$; and the passage from
$j$ to $j+1$ adds a string of $t$ numbers centred at $-\kappa_{j+1}(\tau)$, shifts by a half, and
changes the sign of the numbers below the centre. Two numbers of $\Lambda^{(j)}_\tau$ with
different signs differ by at least $\tfrac12(\Delta-t)-1\ge\tfrac12r^8$, and any two numbers
differ by a non-zero integer. So $\Lambda^{(j)}_\tau$ is the Harish-Chandra parameter of a discrete
series representation $\pi_{j,\tau}$ of $G_{j,\tau}$ (the signs say which coordinates belong to the
positive part), and for every non-compact root $\beta$
\begin{equation}\label{pz:eq-gap}
|\langle\Lambda^{(j)}_\tau,\beta\rangle|\ \ge\ \tfrac12r^8\ >\ x^{25}+N_j .
\end{equation}
For $j=m$ the group is compact, and $\Lambda^{(m)}_\tau=\lambda^\circ_\tau+\rho$ with $\rho=
(\tfrac{N_m-1}2,\dots,\tfrac{1-N_m}2)$: by \S\ref{ps:ss-restrict} the first $k-1$ rows of
$\lambda^\circ_\tau$ are equal to $-m(t+1)/2$, the $k$-th is $-m(t+1)/2+k/2-\kappa(\tau)-\Delta$,
the next $t$ rows are lower by $\tfrac12(\Delta-t-1)$, and so on; adding $\rho$ gives
\eqref{pz:eq-param}. Thus $\pi_{m,\tau}$ is the representation of highest weight $\lambda^\circ_\tau$.

Let $\pi_j=\bigotimes_\tau\pi_{j,\tau}$, $v_j$ a unit vector which, at every $\tau$, is a highest
weight vector of the minimal $K$-type of $\pi_{j,\tau}$, $d_j$ the formal degree of $\pi_j$, and
\[
c_j(x)=\langle v_j,\pi_j(x)v_j\rangle,\qquad e_j(x)=d_j\,\overline{c_j(x)}\qquad(x\in G_{j,\infty})
\]
(inner products conjugate-linear in the first variable, as in \S\ref{ps:ss-weil}; with this reading $\sigma(e_j)$ is the projector onto the type $\pi_j$; see \eqref{par:eq-coef}).
By the orthogonality relations for square-integrable representations, $e_j(x^{-1})=
\overline{e_j(x)}$, $e_j*e_j=e_j$, and for a unitary representation $\sigma$ of $G_{j,\infty}$ the
operator $\int e_j(x)\sigma(x)\,dx$, when $e_j$ is integrable, is the orthogonal projector onto
the vectors that span with their translates a multiple of $\pi_j$ and correspond to $v_j$.

\begin{definition}\label{pz:def-type}
Let $e=e_j$ or $\overline{e_j}$. For a function $F$ on $G_j(\auA)$ put $(R(e)F)(g)=
\int_{G_{j,\infty}}e(x)F(gx)\,dx$ when the integral converges absolutely. A \emph{form of type $e$}
on $G_j$ is a bounded continuous function $y$ on $G_j(\auA)$, left invariant under $G_j(E^+)$, right
invariant under an open subgroup of $G_j(\auA_f)$ (a \emph{level}), with $R(e)y=y$. The
\emph{forms of the archimedean type of the chain} on $G_j$ are the forms of type $e_j$; their
complex conjugates are the forms of type $\overline{e_j}$. The space $\auAz$ of
\S\ref{ss:au-spaces} is the space of the forms of type $e_0$ and level $J_M$; this is the precise
meaning of ``the fixed vector of the discrete series with parameter $\varsigma_\tau$''. For $j=m$
the forms of type $e_m$ are the functions $g\mapsto l\bigl(g_\infty^{-1}\varphi(g_f)\bigr)$ with
$\varphi$ an algebraic modular form of weight $\lambda^\circ$ and $l$ the linear form given by
$v_m$ (Peter--Weyl); in this way the last lift has values in algebraic modular forms.
\end{definition}

\emph{The Fock space.} Let $0\le j<m$, $W=\auHom(Y_j,Y_{j+1})$, $W_\infty=W\otimes_{E^+}(E^+\otimes\R)$,
of real dimension $2D$, $D=fN_jN_{j+1}$. We realise $\psHS_\infty(W)$ as the Fock space $\pzF$ of
entire functions on $\C^D$ with $\|F\|^2=\boldsymbol\pi^{-D}\int|F(z)|^2e^{-|z|^2}dz$, for a
positive complex structure on $W_\infty$ commuting with $K_{j,\infty}\times K_{j+1,\infty}$, with
$\psrho_\infty$ given by \eqref{pk:eq-fock}; $|\cdot|$ is then a euclidean norm on $W_\infty$, and
$(g,h)\in G_{j,\infty}\times G_{j+1,\infty}$ acts on $W_\infty$ with operator norm at most
$\|g\|\,\|h\|$. Let $\omega_\infty$ be the archimedean part of $\omega_{\mu_{j+1}}$ on $\pzF$. For
$w\in W_\infty$ let $\partial_w$ be the derivative at $s=0$ of $\psrho_\infty(sw)$, an operator on the
space $\pzF^{\mathrm{sm}}$ of smooth vectors. Fix a basis $(\mathsf e_k)_{k\le D}$ of $\C^D$,
orthonormal for the hermitian form, so that $(\mathsf e_k,i\mathsf e_k)_k$ is an orthonormal real
basis $\mathsf B$ of $W_\infty$, and for an integer $s\ge0$ put
\[
\|\Psi\|_s=\max\bigl\{\|\partial_{w_1}\cdots\partial_{w_n}\Psi\|:\ 0\le n\le2s,\ w_1,\dots,w_n\in
\mathsf B\bigr\},
\]
and let $\pzF_s$ be the completion of $\pzF^{\mathrm{sm}}$ for $\|\cdot\|_s$, a Banach space
contained in $\pzF$.

\begin{lemma}[Growth]\label{pz:growth}
Let $s=D+1$.
\begin{enumerate}
\item[(a)] For $g\in G_{j,\infty}$, $h\in G_{j+1,\infty}$ and $\Psi\in\pzF^{\mathrm{sm}}$,
$\|\omega_\infty(g,h)\Psi\|_s\le(2D\|g\|\|h\|)^{2s}\|\Psi\|_s$. Hence $\omega_\infty$ extends to a
strongly continuous representation on $\pzF_s$ with the same bound.
\item[(b)] $(1+|z|)^{2s}\,|\Psi(z)|\,e^{-|z|^2/2}\le(D+1)^{2s}\|\Psi\|_s$ for $\Psi\in\pzF_s$ and
$z\in\C^D$.
\item[(c)] For every $f\in\psHS_f(W)$ there is $C_f$ with $|\theta(\Psi\otimes f)|\le C_f\|\Psi\|_s$
for $\Psi\in\pzF^{\mathrm{sm}}$; so $\theta(\,\cdot\otimes f)$ extends continuously to $\pzF_s$.
\end{enumerate}
\end{lemma}
\begin{proof}
(a) By (W1), $\omega_\infty(g,h)^{-1}\partial_w\,\omega_\infty(g,h)=\partial_{u^{-1}w}$, where $u$ is the
action of $(g,h)$ on $W_\infty$. The coordinates of $u^{-1}w$, $w\in\mathsf B$, in the basis
$\mathsf B$ have absolute value at most $\|u^{-1}\|\le\|g\|\|h\|$. So $\partial_{w_1}\cdots\partial_{w_n}
\omega_\infty(g,h)\Psi$ is $\omega_\infty(g,h)$ applied to a sum of $(2D)^n$ terms $c\,\partial_{w_1'}
\cdots\partial_{w'_n}\Psi$ with $|c|\le(\|g\|\|h\|)^n$, and $\omega_\infty(g,h)$ is unitary on $\pzF$.
The map $(g,h)\mapsto\omega_\infty(g,h)\Psi$ is continuous into $\pzF^{\mathrm{sm}}$ for $\Psi$
smooth (the smooth vectors of the Heisenberg group form the Schwartz space, on which the Weil
representation acts continuously~\cite[no.~35]{WeilActa}), and the bound extends this to $\pzF_s$.
(b) By \eqref{pk:eq-fock}, $(\partial_w\Psi)(z)=(z\cdot w)\Psi(z)-(\mathsf D_{\bar w}\Psi)(z)$, where
$\mathsf D_u$ is the derivative in the direction $u$; as $\Psi$ is holomorphic, $\mathsf D_{iu}=
i\mathsf D_u$, so $(z\cdot w)\Psi=\tfrac12(\partial_w-i\partial_{iw})\Psi$. Hence
$\|z^\beta\Psi\|\le\|\Psi\|_s$ for every monomial $z^\beta$ of degree at most $2s$ in the
coordinates $z_k=z\cdot\mathsf e_k$. The reproducing property gives $|F(z)|\le\|F\|e^{|z|^2/2}$ for
$F\in\pzF$; apply it to $F=z^\beta\Psi$ and expand $(1+\sum_k|z_k|)^{2s}$ into $(D+1)^{2s}$ monomials.
(c) Fix $w_0\in W_\infty$ and $f'\in\psHS_f(W)$ with $\theta(\psrho_\infty(w_0)1\otimes f')\ne0$. They
exist: the vectors $\psrho_\infty(w)1$, $w\in W_\infty$, span a dense subspace of $\pzF^{\mathrm{sm}}$,
and $\theta\ne0$ is continuous in the first factor (W0). There is
an integer $\mathsf N\ge1$ such that $\psrho_f(\mathsf NL)$ acts on $f$ and on $f'$ by characters
((W4), for a self-dual lattice $L$). If $\langle\psrho_f(\gamma)f,f'\rangle\ne0$, comparing the
actions of $\psrho_f(l)$, $l\in\mathsf NL$, on both arguments shows that $l\mapsto[l,\gamma]$ is a
fixed character of $\mathsf NL$; so $\gamma$ lies in one coset $x_0+\mathsf N^{-1}L$ ($x_0\in W(\auA_f)$ depending only on $f,f'$). For $\gamma\in
W(E)$ with coordinate $\gamma^\natural\in W_\infty$, $\langle\psrho_\infty(\gamma^\natural)\Psi,\psrho_\infty(w_0)1\rangle$ is a number of absolute value $1$
times $(\psrho_\infty(v)\Psi)(0)=e^{-|v|^2/2}\Psi(-\aubar v)$, $v=\gamma^\natural-w_0$,
which by (b) is at most $(D+1)^{2s}(1+|v|)^{-2s}\|\Psi\|_s$ in absolute value. The points
$\gamma^\natural$, $\gamma\in W(E)\cap(x_0+\mathsf N^{-1}L)$, form a coset of a lattice in
$W_\infty\cong\R^{2D}$, and $2s>2D$; so
\[
T(\Psi\otimes f)=\sum_{\gamma\in W(E)}\bigl\langle\psrho_\infty(w_0)1\otimes f',\psrho(\gamma)(\Psi\otimes f)\bigr\rangle
\]
converges absolutely (each term is the complex conjugate of the term estimated above), with
$|T(\Psi\otimes f)|\le C\|f\|\|f'\|\,\|\Psi\|_s$. Since the hermitian forms are conjugate-linear in
the first variable, $T$ is linear in $\Psi\otimes f$. The functional $T$ is
invariant under $\psrho(W(E))$, hence a multiple of $\theta$ by (W1); the multiple is
$\overline{\theta(\psrho_\infty(w_0)1\otimes f')}\ne0$, as in the proof of Proposition~\ref{pk:conj}(ii)
(for the normalisation of $\theta$ made there).
\end{proof}

\emph{Other spaces.} The proof of Lemma~\ref{pz:growth} uses only the Fock model of
$\psHS_\infty(W)$ for a positive complex structure, (W0), (W1), (W4), the unitarity of $\omega_\infty$
and the continuity of the Weil representation on smooth vectors; nothing specific to the chain
enters, and (a) holds with every integer $s'\ge0$ in place of $s$. So, for an arbitrary hermitian
space $W$ over $E$ with $W_\infty$ of real dimension $2D$ (and a positive complex structure on
$W_\infty$, $\mathsf B$ and $\|\cdot\|_{s'}$ defined as above), (b) and (c) hold, and (a) holds with
$s'$ for $W=\auHom(X,X')$ and $(g,h)\in\auU(X)(E^+\otimes\R)\times\auU(X')(E^+\otimes\R)$ (operator norms
in coordinates in which the forms are $\mathrm{diag}(\pm1)$).

\begin{lemma}[Absolute convergence along an isotropic subspace]\label{pz:absgrowth}
Let $W$ be a hermitian space over $E$, $2D=\dim_\R W_\infty$, $s=D+1$, with $\|\cdot\|_{s'}$ as in the
preceding paragraph. Let $I\subseteq W$ be an isotropic $E$-subspace and, for a datum $F$ and $w\in
W(E)$, $c_w(F)=\theta_{W_I}(\psZ_I\psrho(w)F)$ \textup{(}notation of \textup{(W2))}; $c_w$ depends only on
the class of $w$ in $W(E)/I^\perp(E)$.
\begin{enumerate}
\item[\textup{(a)}] For every $f\in\psHS_f(W)$ there is $C_f$ with $\sum_{w\in W(E)/I^\perp(E)}|c_w(\Psi\otimes
f)|\le C_f\|\Psi\|_{2s}$ for every smooth $\Psi$.
\item[\textup{(b)}] Let $W=\auHom(X,X')$, $f\in\psHS_f(W)$ and $\mathcal C$ a compact subset of
$\auU(X)(\auA_f)\times\auU(X')(\auA_f)$. There is $C$ such that
\[
\sum_{w\in W(E)/I^\perp(E)}\bigl|c_w\bigl(\omega(g,h)(\Psi\otimes f)\bigr)\bigr|\le C\,(2D\|g_\infty\|\|h_\infty\|)^{4s}
\|\Psi\|_{2s}
\]
for all smooth $\Psi$ and all $(g,h)\in\auU(X)(\auA)\times\auU(X')(\auA)$ with $(g_f,h_f)\in\mathcal C$;
the same holds with $\omega$ replaced by $\omega_{\alpha,\beta}$ or by the twisted $\omega_\mu$, which differ
from it by characters of absolute value $1$.
\item[\textup{(c)}] \textup{(Uniformity under archimedean translates.)} For every $f\in\psHS_f(W)$ there is
$C_f$ with
\[
\sum_{w\in W(E)/I^\perp(E)}|c_w(\psrho_\infty(z)\Psi\otimes f)|\le C_f\|\Psi\|_{2s}
\]
for every
$z\in W_\infty$ and every smooth $\Psi$.
\end{enumerate}
\end{lemma}
\begin{proof}
\emph{Step 1: translates.} For every $f$ there is $C_f$ with $|\theta(\psrho_\infty(z)G\otimes f)|\le
C_f\|G\|_s$ for all $z\in W_\infty$ and all smooth $G$. Indeed, in the proof of Lemma~\ref{pz:growth}(c),
replacing $\Psi$ by $\psrho_\infty(z)G$ replaces $v=\gamma^\natural-w_0$ by $\gamma^\natural+z-w_0$ (up to a
phase, as $\psrho_\infty(\gamma^\natural)\psrho_\infty(z)$ is a phase times $\psrho_\infty(\gamma^\natural+z)$) and
leaves the finite factors $\langle\psrho_f(\gamma)f,f'\rangle$ unchanged; so $|T(\psrho_\infty(z)G\otimes
f)|\le(D+1)^{2s}\|f\|\|f'\|\,\|G\|_s\sup_{u\in\R^{2D}}\sum_\lambda(1+|\lambda+u|)^{-2s}$, the sum over a coset
of a fixed lattice of $W_\infty\cong\R^{2D}$; the supremum is finite because $2s>2D$.

\emph{Step 2: the terms are Fourier coefficients.} Let $F=\Psi\otimes f$ and
$\phi(i)=\theta(\psrho(i)F)$ for $i\in I(\auA)$. On $I(\auA)\times I(\auA)$ the cocycle $\beta$ is $1$, so
$\psrho$ is additive on $I(\auA)$; hence $\phi(\gamma+i)=\phi(i)$ for $\gamma\in I(E)$ (W1), and, $f$ being
smooth, there is an open compact subgroup $K_f\subseteq I(\auA_f)$ with $\psrho_f(\kappa)f=f$, so that
$\phi(i+\kappa)=\phi(i)$ for $\kappa\in K_f$. For $w\in W(E)$, $\psrho(i)\psrho(w)=[i,w]\psrho(w)\psrho(i)$
and $\theta\circ\psrho(w)=\theta$, so by (W2)
\[
c_w(F)=\int_{[I]}\theta(\psrho(i)\psrho(w)F)\,di=\int_{[I]}[i,w]\,\phi(i)\,di .
\]
The character $i\mapsto[i,w]$ of $I(\auA)$ is trivial on $I(E)$, depends only on $w$ modulo
$I^\perp(\auA)$, and determines the class of $w$ modulo $I^\perp(E)$: if $[i,w]=1$ for all $i\in I(\auA)$
then, replacing $i$ by $ai$ ($a\in\auA_E$) and using the non-degeneracy of $\psi$, $\langle w,i\rangle=0$
for all $i$, i.e.\ $w\in I^\perp(E)$. If this character is non-trivial on $K_f$ then $c_w(F)=0$. Since $E$
is dense in $\auA_{E,f}$, $I(\auA_f)=I(E)+K_f$, so $i_\infty\mapsto(i_\infty,0)$ induces an isomorphism of
$I_\infty/\Lambda$ onto $[I]/K_f$, $\Lambda=\{\gamma\in I(E):\gamma\in K_f\}$ (a lattice in $I_\infty$),
compatible with the normalised Haar measures. Write $[i_\infty,w]=e^{2\boldsymbol\pi i\lambda_w(i_\infty)}$
with $\lambda_w$ a real linear form on $I_\infty$. Then, for the $w$ with $c_w(F)\ne0$, $\lambda_w$ lies in the
dual lattice $\Lambda^*$, distinct classes $w$ give distinct $\lambda_w$, and $c_w(F)=\hat\phi(\lambda_w)$, where
$\hat\phi(\lambda)=\int_{I_\infty/\Lambda}e^{2\boldsymbol\pi i\lambda(i)}\phi(i)\,di$. So
$\sum_w|c_w(F)|\le\sum_{\lambda\in\Lambda^*}|\hat\phi(\lambda)|$.

\emph{Step 3: decay of the coefficients.} For $i\in I_\infty$, $\phi(i)=\theta(\psrho_\infty(i)\Psi\otimes
f)$. The orbit map of the smooth vector $\Psi$ is smooth for the topology of smooth vectors, in which
$\|\cdot\|_s$ is continuous, $\psrho_\infty$ is additive on $I_\infty$, and $\theta(\,\cdot\otimes f)$ is
continuous for $\|\cdot\|_s$ (Lemma~\ref{pz:growth}(c)); so $\phi$ is smooth on $I_\infty$ and
$\mathsf D_{u_1}\cdots\mathsf D_{u_n}\phi(i)=\theta(\psrho_\infty(i)\partial_{u_1}\cdots\partial_{u_n}\Psi\otimes f)$
for $u_k\in I_\infty$. Let $\varkappa=\dim_\R I_\infty$; as $I$ is isotropic, $\varkappa\le D<2s$. Let
$u_1,\dots,u_\varkappa$ be an orthonormal basis of $I_\infty$ and $\Delta=\sum_k\mathsf D_{u_k}^2$. Each $u_k$ is a
combination of the elements of $\mathsf B$ with coefficients of absolute value $\le1$, and $w\mapsto
\partial_w$ is linear, so a product of $2s$ operators $\partial_{u_k}$ is a sum of at most $(2D)^{2s}$
products of $2s$ operators $\partial_b$, $b\in\mathsf B$, with coefficients of absolute value $\le1$; by
Step~1,
\[
\sup_{i}|\Delta^s\phi(i)|\le\varkappa^s(2D)^{2s}C_f\max_{b_1,\dots,b_{2s}\in\mathsf B}\|\partial_{b_1}\cdots
\partial_{b_{2s}}\Psi\|_s\le\varkappa^s(2D)^{2s}C_f\|\Psi\|_{2s}.
\]
For $0\ne\lambda\in\Lambda^*$, integration by parts on the torus gives $\hat\phi(\lambda)=(-4\boldsymbol\pi^2
|\lambda|^2)^{-s}\int e^{2\boldsymbol\pi i\lambda(i)}\Delta^s\phi(i)\,di$, and $|\hat\phi(0)|\le\sup|\phi|\le C_f
\|\Psi\|_s$ by Step~1. Since $2s>\varkappa$, $\sum_{0\ne\lambda\in\Lambda^*}|\lambda|^{-2s}<\infty$, and (a)
follows with $C_f(1+\varkappa^s(2D)^{2s}(2\boldsymbol\pi)^{-2s}\sum_{\lambda\ne0}|\lambda|^{-2s})$ (note
$\|\Psi\|_s\le\|\Psi\|_{2s}$).

(b) The vector $f$ is fixed, up to scalars of absolute value $1$, by an open subgroup $K'$ of
$\auU(X)(\auA_f)\times\auU(X')(\auA_f)$; $\mathcal C$ is covered by finitely many cosets $c_jK'$, so
$\omega_f(g_f,h_f)f$ lies, up to such scalars, in the finite set $\{\omega_f(c_j)f\}$. Now
$\omega(g,h)(\Psi\otimes f)=\omega_\infty(g_\infty,h_\infty)\Psi\otimes\omega_f(g_f,h_f)f$, the vector
$\omega_\infty(g_\infty,h_\infty)\Psi$ is smooth, and $\|\omega_\infty(g_\infty,h_\infty)\Psi\|_{2s}\le
(2D\|g_\infty\|\|h_\infty\|)^{4s}\|\Psi\|_{2s}$ by Lemma~\ref{pz:growth}(a) with $s'=2s$. Apply (a) to the
finitely many $\omega_f(c_j)f$.

(c) Fix $z\in W_\infty$ and put $F_z=\psrho_\infty(z)\Psi\otimes f$, $\phi_z(i)=\theta(\psrho(i)F_z)$. Steps~1
and~2 apply to $F_z$ (they use only that $F_z$ is a datum on which $\theta$ is continuous, and
$\psrho_\infty(z)$ preserves $\pzF_s$ by Lemma~\ref{pcx:superpos}(a)): with the same $K_f$, $\Lambda$ and
$\lambda_w$ (which depend on $f$ only), $\sum_w|c_w(F_z)|\le\sum_{\lambda\in\Lambda^*}|\hat\phi_z(\lambda)|$.
Step~3 applied to $F_z$ would bound the coefficients by $\|\psrho_\infty(z)\Psi\|_{2s}$, which grows with
$z$ (the derivatives of $\phi_z$ pick up factors $\kappa(u,z)$); we remove this dependence by a unitary
character. For $i\in I_\infty$,
$\psrho_\infty(i)\psrho_\infty(z)=e^{\kappa(i,z)}\psrho_\infty(z)\psrho_\infty(i)$ with $\kappa(i,z)\in i\R$
real bilinear (Lemma~\ref{pcx:superpos}); write $e^{\kappa(i,z)}=e^{2\boldsymbol\pi i\mu_z(i)}$, $\mu_z$ a real
linear form on $I_\infty$. Then
\[
\phi_z(i)=e^{2\boldsymbol\pi i\mu_z(i)}\varphi_z(i),\qquad\varphi_z(i)=\theta\bigl(\psrho_\infty(z)\psrho_\infty(i)
\Psi\otimes f\bigr).
\]
The orbit map $i\mapsto\psrho_\infty(i)\Psi$ is smooth into the smooth vectors, hence into $\pzF_s$, with
derivatives $\psrho_\infty(i)\partial_{u_1}\cdots\partial_{u_n}\Psi$; $\psrho_\infty(z)$ is bounded on
$\pzF_s$ and $\theta(\,\cdot\otimes f)$ is continuous on $\pzF_s$. So $\varphi_z$ is smooth and
$\mathsf D_{u_1}\cdots\mathsf D_{u_n}\varphi_z(i)=\theta(\psrho_\infty(z)\psrho_\infty(i)\partial_{u_1}\cdots
\partial_{u_n}\Psi\otimes f)$; as $\psrho_\infty(z)\psrho_\infty(i)$ is a phase times
$\psrho_\infty(z+i)$, Step~1 gives, exactly as in Step~3 and now uniformly in $z$,
\[
\sup_i|\varphi_z(i)|\le C_f\|\Psi\|_s,\qquad\sup_i|\Delta^s\varphi_z(i)|\le\varkappa^s(2D)^{2s}C_f\|\Psi\|_{2s}.
\]
Since $\phi_z$ is $\Lambda$-periodic, $\varphi_z(i+\gamma)=e^{-2\boldsymbol\pi i\mu_z(\gamma)}\varphi_z(i)$ for
$\gamma\in\Lambda$, and the same holds for every derivative of $\varphi_z$. For $\lambda\in\Lambda^*$ put
$\nu=\lambda+\mu_z$, so that $\hat\phi_z(\lambda)=\int_{I_\infty/\Lambda}e^{2\boldsymbol\pi i\nu(i)}\varphi_z(i)
\,di$, the integrand being $\Lambda$-periodic. If $\nu\ne0$, $e^{2\boldsymbol\pi i\nu(i)}=(-4\boldsymbol\pi^2
|\nu|^2)^{-1}\Delta e^{2\boldsymbol\pi i\nu(i)}$, and Green's formula on a fundamental parallelotope has no
boundary terms, because every product of a derivative of $e^{2\boldsymbol\pi i\nu(\cdot)}$ with a derivative
of $\varphi_z$ is $\Lambda$-periodic (the multipliers $e^{\pm2\boldsymbol\pi i\mu_z(\gamma)}$ cancel, and
$\lambda(\gamma)\in\Z$); iterating, $\hat\phi_z(\lambda)=(-4\boldsymbol\pi^2|\nu|^2)^{-s}\int e^{2\boldsymbol\pi
i\nu(i)}\Delta^s\varphi_z(i)\,di$. Hence for all $\lambda$ (covolume $1$ for the normalised measure)
\[
|\hat\phi_z(\lambda)|\le\min\bigl(C_f\|\Psi\|_s,\ (2\boldsymbol\pi|\lambda+\mu_z|)^{-2s}\varkappa^s(2D)^{2s}
C_f\|\Psi\|_{2s}\bigr)\le C'_f\,(1+|\lambda+\mu_z|)^{-2s}\|\Psi\|_{2s},
\]
with $C'_f=2^{2s}C_f\max(1,\varkappa^s(2D)^{2s}(2\boldsymbol\pi)^{-2s})$ (use $\|\Psi\|_s\le\|\Psi\|_{2s}$ and
$(1+r)^{2s}\le2^{2s}\max(1,r^{2s})$). Finally $\sup_{\mu}\sum_{\lambda\in\Lambda^*}(1+|\lambda+\mu|)^{-2s}<\infty$,
the supremum over real linear forms $\mu$, since $2s>\varkappa$ (compare the sum with an integral over
$\R^\varkappa$, as in Step~1). This gives (c), with a constant depending on $f$ only.

All three bounds extend to $\Psi\in\pzF_{2s}$: by the formula of Step~2 and Step~1, each $c_w(\,\cdot\otimes
f)$ is continuous for $\|\cdot\|_s$, and so is $\Psi\mapsto c_w(\psrho_\infty(z)\Psi\otimes f)$ for fixed
$z$ (Lemma~\ref{pcx:superpos}(a)), so the bounds for the finite partial sums extend by density, and
$\omega_\infty$ extends to $\pzF_{2s}$ with the same bound.
\end{proof}

\emph{The three archimedean statements.} Put $B=x^{25}$; then $2D+2\le B$, since $f\le x^{12}$ and
$N_j\le n_m\le40x^6$. Let $0\le j<m$. The \emph{joint harmonic} $P_j$ of \S\ref{ps:ss-Y} is a
non-zero polynomial in $\pzF$, harmonic for $G_{j,\infty}$ and for $G_{j+1,\infty}$, which is a
highest weight vector for $K_{j,\infty}$ and for $K_{j+1,\infty}$ with the weights of the minimal
$K$-types of the two types (for the action $\omega_\infty$, which contains the determinant twists
of \S\ref{ps:ss-weil}); we write $\Phi_j$ for this vector of $\pzF$.
\begin{enumerate}
\item[(A1)] \emph{(Integrability.)} For $0\le j\le m$ the function $c_j$ is smooth, left and right
$K_{j,\infty}$-finite and an eigenfunction of the centre of the enveloping algebra, and
$\int_{G_{j,\infty}}|\mathsf Xc_j(x)|\,\|x\|^{B}dx<\infty$ for every differential
operator $\mathsf X$ that is invariant under left or under right translations (including
$\mathsf X=1$).
\item[(A2)] \emph{(Cusp property.)} For $0\le j<m$, every proper parabolic subgroup of
$G_{j,\infty}$ of the form $\prod_\tau\mathsf P_\tau$ with $\mathsf P_\tau$ proper for all $\tau$,
with unipotent radical $\mathsf N$, and all $a,b\in G_{j,\infty}$: $\int_{\mathsf N}c_j(a\,n\,b)\,dn=0$,
the integral being absolutely convergent.
\item[(A3)] \emph{(Correspondence.)} For $0\le j<m$ such a polynomial $P_j$ exists, it is unique up
to a scalar, and in $\pzF$
\[
\int_{G_{j,\infty}}e_j(x)\,\omega_\infty(x^{-1},1)\Phi_j\,dx=\int_{G_{j+1,\infty}}e_{j+1}(x)\,
\omega_\infty(1,x)\Phi_j\,dx .
\]
\end{enumerate}
In words, (A3) says that the component of $\Phi_j$ of the type of $Y_j$ (in the first variable)
equals its component of the type of $Y_{j+1}$ (in the second). It holds if the closed subspace of
$\pzF$ on which $G_{j,\infty}$ acts by multiples of the contragredient of $\pi_j$ is, under
$G_{j+1,\infty}$, a multiple of $\pi_{j+1}$ in which the vector attached to $v_j$ and to $\Phi_j$
is the vector $v_{j+1}$: this is the correspondence of the discrete series with the parameters
\eqref{pz:eq-param} together with the correspondence of their minimal $K$-types in the space of
joint harmonics.

\smallskip\noindent (A1) is Theorem~\ref{par:A1} and (A2) is Corollary~\ref{par:A2}; both are proved from the
decay bound \eqref{par:eq-decay}, which holds by Theorem~\ref{pds:decay} and formula
\eqref{pds:eq-FJ} (Theorem~\ref{pco:FJ}). (A3) is Theorem~\ref{par:A3}, proved modulo the published
results (D1)--(D3) of \S\ref{par:ss-rigid} and the normalisation (V) (Lemma~\ref{par:vac}), which is
Theorem~\ref{pco:V}, proved modulo (W1$_\infty$), (W2$_\infty$), (M), and, for $j=m-1$, (S). The
identification of the type on $Y_m$ with $W_{\lambda^\circ}$ is Remark~\ref{par:Ym}. Everything below
is proved from (A1)--(A3).\par\smallskip

\begin{proposition}[The lower chain is well defined]\label{pz:types}
Assume {\rm(A1)--(A3)}. Let $0\le j\le m$.
\begin{enumerate}
\item[(a)] Every form $y$ of type $e_j$ or $\overline{e_j}$ on $G_j$ is smooth, $K_{j,\infty}$-finite
and finite under the centre of the enveloping algebra; it is a cusp form; it is rapidly
decreasing on Siegel sets, in the form: for every Siegel set $\mathfrak S_\infty\subset G_{j,\infty}$, finite set $\mathsf c\subset G_j(\auA_f)$, open subgroup $K'$ of $G_j(\auA_f)$ and $N\ge0$ there is $C_N$ with $|y(g)|\le C_N\|g_\infty\|^{-N}$ for $g\in\mathfrak S_\infty\mathsf cK'$; and the forms of a given level form a finite-dimensional space.
\item[(b)] Let $j<m$, $f\in\psHS_f(W)$ and let $y$ be a form of type $e_j$ on $G_j$. Then the
integral $F(h)=\int_{[G_j]}\theta\bigl(\omega_{\mu_{j+1}}(g,h)(\Phi_j\otimes f)\bigr)y(g)\,dg$ converges
absolutely and $F$ is a form of type $e_{j+1}$ on $G_{j+1}$.
\item[(c)] Let $j<m$, $f\in\psHS_f(W)$ and let $z$ be a form of type $\overline{e_{j+1}}$ on
$G_{j+1}$. Then $F^*(g)=\int_{[G_{j+1}]}\theta\bigl(\omega_{\mu_{j+1}}(g,h)(\Phi_j\otimes f)\bigr)z(h)
\,dh$ converges absolutely and $F^*$ is a form of type $\overline{e_j}$ on $G_j$.
\item[(d)] The space $\auAz$ is finite-dimensional and stable under the right translations by
the stabiliser of the lattice at the chosen places, so $\auG$ acts on it; the Poincar\'e map
$\auP$ of \S\ref{ss:au-spaces} is well defined; and $\iota_\vartheta$ maps $\auAz$ to algebraic
modular forms of weight $\lambda^\circ$ on $\auU(Y_m)$.
\end{enumerate}
In particular the assumption made in \S\ref{ps:ss-Y} holds: the lifts $\Theta^Y_{j,\vartheta}$ and
their transposes $\Theta^{Y,*}_{j,\vartheta}$ preserve the types, and all the forms concerned are
cusp forms and rapidly decreasing.
\end{proposition}
\begin{proof}
We write $G=G_j$, $e=e_j$; the case of $\overline{e_j}$ is the same.

(a) Let $y=R(e)y$ be bounded. For $g=(g_\infty,g_f)$, $y(g)=\int e(g_\infty^{-1}x)\,y(x,g_f)\,dx$.
By (A1) the derivatives of $e$ under left translations are integrable, so $y$ is smooth in
$g_\infty$ and its derivatives are bounded; the left translates of $e$ under $K_{j,\infty}$ span a finite-dimensional
space, so $y$ is right $K_{j,\infty}$-finite; and an element of the centre of the enveloping
algebra acts on $y$ through its action on $e$, by a scalar. So $y$ is an automorphic form, and it
is bounded. For $j=m$ the quotient $[G]$ is compact, there are no proper parabolic subgroups, and
the finiteness follows from the Peter--Weyl theorem. Let $j<m$,
let $\mathsf P$ be a proper parabolic subgroup of $G$ over $E^+$ with unipotent radical
$\mathsf N$, and $y_{\mathsf P}(g)=\int_{\mathsf N(E^+)\backslash\mathsf N(\auA)}y(ng)\,dn$. It is
bounded, left invariant under $\mathsf N(\auA)$, and $R(e)y_{\mathsf P}=y_{\mathsf P}$ (Fubini).
Fix $g_f$ and put $\phi(x)=y_{\mathsf P}(x,g_f)$ on $G_{j,\infty}$; it is left invariant under
$\mathsf N_\infty=\mathsf N(E^+\otimes\R)$, the unipotent radical of $\prod_\tau\mathsf P(E^+_\tau)$,
and
\[
y_{\mathsf P}(g_\infty,g_f)=\int_{G_{j,\infty}}e(g_\infty^{-1}x)\phi(x)\,dx=\int_{\mathsf N_\infty
\backslash G_{j,\infty}}\phi(x)\Bigl(\int_{\mathsf N_\infty}e(g_\infty^{-1}nx)\,dn\Bigr)dx=0
\]
by (A2), the function $e$ being integrable and $\phi$ bounded. So $y$ is a cusp form. A cuspidal
automorphic form is rapidly decreasing on Siegel sets: on $\mathfrak S_\infty=\omega A_tK$ ($\omega$, $K$ compact) one has $|y(g)|\le C_NH(a)^{-N}$ for $g=uak\,ck'$, where $H(a)=\max_\alpha a^\alpha$ over the simple roots; and $\|g_\infty\|\le C\|a\|\le C''H(a)^{c_0}$ for a fixed $c_0>0$, because $\omega,K$ are compact, $\|a\|$ is the largest value of a weight of the standard representation, a non-negative combination of the simple roots on $A_t$, and $a^\alpha\ge t$; since $\|g_\infty\|\ge1$ ($|\det g_\infty|=1$), this gives the bound $|y(g)|\le C_N\|g_\infty\|^{-N}$ stated in (a). Moreover, the cusp forms of a given level,
$K_\infty$-type and infinitesimal character form a finite-dimensional space (Gelfand and
Piatetski-Shapiro, Harish-Chandra; see \cite[\S1.8, p.~192]{BorelJacquet}, \cite[Chapter~I, Theorem~1, p.~8]{HarishChandraAF}).

(b) Let $s=D+1$, so $2s\le B$. By reduction theory (finiteness of the class number \cite[Theorem~5.1]{BorelFin}; Siegel sets as fundamental sets for the arithmetic subgroups, \cite[Theorem~6.5, and Theorem~12.3 over a number field]{BorelHC1962}; finite Haar measure of Siegel sets, \cite[Proposition~4.3]{BorelHC1962}, stated there for an open subgroup of a real algebraic semisimple group; it applies to the derived group $G^1$ of $G_{j,\infty}$, a product of groups $\mathrm{SU}(p,q)$, and hence to $G_{j,\infty}=Z\cdot G^1$, $Z$ the compact centre: an Iwasawa decomposition $KAN$ of $G_{j,\infty}$ has $A,N\subset G^1$ and $K=Z\cdot(K\cap G^1)$, so a Siegel domain of $G_{j,\infty}$ is $Z$ times a Siegel domain of $G^1$, and its Haar measure is at most $\operatorname{vol}(Z)$ times that of the latter for compatible normalisations, the multiplication map $Z\times G^1\to G_{j,\infty}$ being a surjective homomorphism with finite kernel) there are a Siegel set
$\mathfrak S_\infty\subset G_{j,\infty}$ of finite measure, a finite set $\mathsf c\subset G(\auA_f)$
and an open subgroup $K'$ of $G(\auA_f)$ fixing $f$ under $\omega_{\mu_{j+1}}(\cdot,1)$ and fixing
$y$, such that $G(\auA)=G(E^+)\,\mathfrak S_\infty\,\mathsf c\,K'$. For $g=g_\infty ck'$ in this set
and $\Psi\in\pzF_s$, Lemma~\ref{pz:growth} gives
\[
\bigl|\theta\bigl(\omega(g,h)(\Psi\otimes f)\bigr)\bigr|\le C(h_f)\,(2D\|g_\infty\|\|h_\infty\|)^{2s}
\|\Psi\|_s,\qquad C(h_f)=\max_{c\in\mathsf c}C_{\omega_f(c,h_f)f},
\]
and by (a), with $N=2s$, $|y(g)|\le C'\|g_\infty\|^{-2s}$ on this set. The integrand is left invariant under
$G(E^+)$, by (W1). Hence
\[
L_h(\Psi)=\int_{[G]}\theta\bigl(\omega(g,h)(\Psi\otimes f)\bigr)y(g)\,dg
\]
converges absolutely, and $|L_h(\Psi)|\le C''(h_f)\|h_\infty\|^{2s}\|\Psi\|_s$. The function
$F=L_\bullet(\Phi_j)$ is continuous, left invariant under $G_{j+1}(E^+)$ by (W1), and right
invariant under an open subgroup fixing $f$. For $x\in G_{j+1,\infty}$, $F(hx)=L_h(\omega_\infty(1,x)
\Phi_j)$, so $|F(hx)|\le C(h)\|x\|^{2s}$, and by (A1) the following integrals converge absolutely,
the vector-valued ones in $\pzF_s$:
\[
\mathsf P'\Psi=\int_{G_{j+1,\infty}}e_{j+1}(x)\,\omega_\infty(1,x)\Psi\,dx,\qquad\mathsf P\Psi=
\int_{G_{j,\infty}}e_j(x)\,\omega_\infty(x^{-1},1)\Psi\,dx,\qquad R(e_{j+1})F=L_\bullet(\mathsf P'\Phi_j).
\]
Since $y=R(e_j)y$, Fubini and the change of variable
$g\mapsto gx^{-1}$ give, for every $\Psi\in\pzF_s$,
\[
L_h(\Psi)=\int_{G_{j,\infty}}e_j(x)\int_{[G]}\theta\bigl(\omega(gx^{-1},h)(\Psi\otimes f)\bigr)y(g)\,dg
\,dx=L_h(\mathsf P\Psi).
\]
By (A3), $\mathsf P\Phi_j=\mathsf P'\Phi_j$. Therefore
\[
R(e_{j+1})F=L_\bullet(\mathsf P'\Phi_j)=L_\bullet(\mathsf P\Phi_j)=L_\bullet(\Phi_j)=F .
\]
It remains to see that $F$ is bounded. Let $\mathcal V$ be the space of the forms of type
$e_{j+1}$ on $G_{j+1}$ of a level $K''$ fixing $F$; it is finite-dimensional by (a). For $T>0$ let
$F_T=F$ where $|F|\le T$ and $F_T=0$ elsewhere. It is bounded, measurable, left invariant under
$G_{j+1}(E^+)$ and right invariant under $K''$; so $R(e_{j+1})F_T$ is bounded, continuous
(translation is continuous in $L^1$), has the same invariance, and is fixed by $R(e_{j+1})$,
because $e_{j+1}*e_{j+1}=e_{j+1}$. Thus $R(e_{j+1})F_T\in\mathcal V$. For $T\to\infty$,
$R(e_{j+1})F_T\to R(e_{j+1})F=F$ pointwise, by dominated convergence. A finite-dimensional space of
functions is closed for pointwise convergence; so $F\in\mathcal V$.

(c) The same proof, with the two groups exchanged. With $L^*_g(\Psi)=\int_{[G_{j+1}]}\theta(\omega
(g,h)(\Psi\otimes f))z(h)\,dh$ one finds, using $\overline{e}(x^{-1})=e(x)$: $L^*_g(\Psi)=L^*_g(
\mathsf P'\Psi)$ from $z=R(\overline{e_{j+1}})z$, and $R(\overline{e_j})F^*=L^*_\bullet(\mathsf P
\Phi_j)$. So $R(\overline{e_j})F^*=L^*_\bullet(\mathsf P'\Phi_j)=F^*$ by (A3), and $F^*$ is bounded by
the truncation argument.

(d) The first assertion is (a); right translation by an element of $G_0(\auA_f)$ commutes with
$R(e_0)$. The group $R_0=U\rtimes\auU(X_0)$ has compact quotient $[R_0]$ ($\auU(X_0)$ is
anisotropic), the function $\mathcal W_\varphi$ is bounded on it and the elements of $\auAz$ are
bounded and continuous; so $y\mapsto\int_{[R_0]}\mathcal W_\varphi(r)\overline{y(rg_0)}\,dr$ is an
antilinear form on the finite-dimensional inner product space $\auAz$, which defines $\auP\varphi
\in\auAz$. The last assertion is (b) applied $m$ times with $f=\Upsilon_{j,f}$
(Lemma~\ref{ps:level} gives the level) and Definition~\ref{pz:def-type} for $j=m$.
\end{proof}

\begin{lemma}[Level and character on $Y_m$]\label{pz:lastlevel}
Assume {\rm(A1)--(A3)}. Let $\vartheta$ be a twist, $z\in\auAz$, and let $k\in K_m$ have, at every
level place $w$, a reduction $\bar k_w$ that is block upper triangular for the flag
$\psM_0\subset\psM_1\subset\dots\subset\psM_m$ of the frames, with diagonal blocks
$\bar k_w^{(0)},\dots,\bar k_w^{(m)}$, and $\bar k_w^{(0)}=1$ at the chosen places. Then
$(\iota_\vartheta z)(gk)=\prod_w\prod_{i=1}^m\psi_{i,w}(\det\bar k_w^{(i)})\cdot(\iota_\vartheta z)(g)$,
where $\psi_{i,w}$ is the tame part of $\mu_i$ at $w$. So $\iota_\vartheta z$ has the level
$J_\vartheta$ and the character $\chi_\vartheta$ of \S\ref{ss:au-spaces}.
\end{lemma}
\begin{proof}
Choose $k_i\in K_i$ $(0\le i\le m)$, $k_m=k$, whose reduction at each level place is the upper
left block of $\bar k_w$ of the size of $Y_i$ (the reduction maps are onto, the level places being
split). Then $\bar k_{i+1}\iota_i=\iota_i\bar k_i$, and as in the proof of
Proposition~\ref{au:equiv}, by (W1) and Lemma~\ref{ps:level},
$(\Theta^Y_{i,\vartheta}F)(hk_{i+1})=\varepsilon_i\,\Theta^Y_{i,\vartheta}(F(\,\cdot\,k_i))(h)$ with
$\varepsilon_i=\prod_w\psi_{i+1,w}(\det\bar k_{i+1,w})\psi_{i+1,w}(\det\bar k_{i,w})^{-1}=\prod_w
\psi_{i+1,w}(\det\bar k_w^{(i+1)})$. By induction $(\iota_\vartheta z)(gk)=\prod_i\varepsilon_i\cdot
\iota_\vartheta(z(\,\cdot\,k_0))(g)$, and $z(\,\cdot\,k_0)=z$ because $k_0\in J_M$.
\end{proof}

\begin{lemma}[The characters of the lifts exist]\label{pz:mu}
Let $S$ be a finite set of primes of $E$, no two above the same place of $E^+$, each above a place
split in $E$; let $(c_\tau)_{\tau\in\auPl}$ be integers and $\psi_w$ ($w\in S$) characters of
$\auF_w^\times$ with $\prod_\tau\tau(\xi)^{c_\tau}\prod_{w\in S}\psi_w(\xi\bmod w)=1$ for every root of
unity $\xi\in E$. There is a unitary character $\mu$ of $\auU(1)(E^+)\backslash\auU(1)(\auA_{E^+})$ with
$\mu_\tau(u)=\tau(u)^{c_\tau}$ at the infinite places, with tame part $\psi_w$ at $w\in S$, and
unramified at all other finite places. If the reduction of the roots of unity of $E$ modulo some
$w_0\in S$ is injective (for instance a chosen place $w_0$), then for all $(c_\tau)$ and all $\psi_w$,
$w\ne w_0$, there is $\psi_{w_0}$ satisfying the condition. In particular the characters $\chi_M$,
$\mu^\flat$ and $\mu_j^0$ of \S\ref{ss:au-spaces} exist (Definition~\ref{pfq:def-chars},
Proposition~\ref{pfq:tame}).
\end{lemma}
\begin{proof}
Let $\mathsf U=\auU(1)(E^+\otimes\R)\times\prod_v\auU(1)(\auO_v)$, an open subgroup of
$\auU(1)(\auA_{E^+})$. Its intersection with $\auU(1)(E^+)$ consists of the $\xi\in E$ with
$\xi\aubar\xi=1$ integral at every finite place; all conjugates of such $\xi$ have absolute value
$1$, so (Kronecker) it is the group $\mu_E$ of roots of unity of $E$. At $w\in S$, above the split
place $v$, $\auU(1)(\auO_v)=\auO_w^\times\to\auF_w^\times$. The formula
$\mu_{\mathsf U}(u)=\prod_\tau\tau(u_\tau)^{c_\tau}\prod_{w\in S}\psi_w(u_w\bmod w)$ defines a continuous
unitary character of $\mathsf U$, and the hypothesis says that it is trivial on $\mu_E$ (a rational
point $\xi$ has components $\tau(\xi)$ at the infinite places and $\xi$ at the finite ones). Hence
$\gamma u\mapsto\mu_{\mathsf U}(u)$ is a well-defined character of the open subgroup
$\auU(1)(E^+)\mathsf U$, trivial on $\auU(1)(E^+)$. A unitary character of an open (hence closed)
subgroup of a locally compact abelian group extends to the whole group (Pontryagin duality: the
restriction map of dual groups is onto); the extension is trivial on $\auU(1)(E^+)$ and is the
required $\mu$. For the last assertion, $\xi\mapsto\prod_\tau\tau(\xi)^{-c_\tau}\prod_{w\ne
w_0}\psi_w(\xi\bmod w)^{-1}$ is a character of the subgroup $\mu_E$ of the cyclic group
$\auF_{w_0}^\times$, and extends to it; at a chosen place $w_0$ the reduction of $\mu_E$ is injective
because $p_0,p_1>b^{10}$ do not divide $\#\mu_E$. The characters of \S\ref{ss:au-spaces}: the
condition is checked in Definition~\ref{pfq:def-chars} and Proposition~\ref{pfq:tame}(T1).
\end{proof}

\subsubsection{The constant $C_0$ and the definition of admissible data}\label{pz:ss-const}

Proposition~\ref{au:kernel} was proved in \S\ref{ps:ss-kernel} with the inequality $C_0\le\log\log
b$, justified in (R5) of \S\ref{pk:ss-rest} by ``taking the constant in the definition of
admissible data at least $C_0$''. We make this precise. Let $(q_0,d)$ be parameters in the sense
of Definition~\ref{an:def-param}, $K=\Q(\sqrt{-d})$, and
\begin{equation}\label{pz:eq-C0}
C_0(q_0,d)=4+\log^+\mathsf M_d+\frac{\sqrt d}{\boldsymbol\pi\,\mathsf m}+\log\Bigl(1+\frac4\delta\Bigr),
\qquad\mathsf m=\prod_{p\mid q_0}\Bigl(1-\frac1p\Bigr)p^{1/(p-1)},\quad\delta=\Bigl(\frac{2\boldsymbol
\pi\,\mathsf m}{\sqrt d}\Bigr)^{1/2},
\end{equation}
where $\mathsf M_d=\max_{\Lambda'}\sup_{u\in\C}|\theta^{\Lambda'}_u(0)|$, the maximum over the
normalised lattices $(\Lambda',t',a_1')$ with $\auO_K\Lambda'=\Lambda'$ (\S\ref{pk:ss-curve}, with
$t'$ admissible).

\begin{lemma}\label{pz:C0}
\begin{enumerate}
\item[(a)] $C_0(q_0,d)$ is a finite real number that depends only on $q_0$ and $d$: not on the
scales $x,x_1,b$, nor on $s$, $q_1$, $E$ or any other part of admissible data.
\item[(b)] It is the constant $C_0$ of Proposition~\ref{pk:conj} and of Lemma~\ref{ps:conj}{\rm(ii)}.
\item[(c)] With the constant $C_P$ of \eqref{an:eq-CP} in condition {\rm(D2)} of
Definition~\ref{an:data}, which contains $d$ and $C_0(q_0,d)$ as summands (all summands are
non-negative), all admissible data satisfy
$C_0\le10^{-12}\log\log b$ and $d\le10^{-12}\log\log b$.
\end{enumerate}
\end{lemma}
\begin{proof}
(a) The lattices $\Lambda'$ with $\auO_K\Lambda'=\Lambda'$ form $h_K$ classes up to homothety. In
each class the admissible point is determined modulo the lattice (\S\ref{pf:ss-adm}), and the
condition $\wp''(t')=2$ fixes the lattice up to a fourth root of unity, since
$\wp''_{c\Lambda}(ct)=c^{-4}\wp''_\Lambda(t)$; and $a_1'$ is determined up to sign. So there are
finitely many normalised lattices, and they depend on $d$ only. For each of them $u\mapsto|\theta_u
(0)|$ is continuous and $\Lambda'$-periodic (\S\ref{pk:ss-curve}), hence bounded. So $\mathsf M_d$
is finite and depends on $d$ only. The primes dividing $q_1$ are those dividing $q_0$
(\S\ref{an:ssec-param}), so the number $\mathsf m$ of the proof of Proposition~\ref{pk:conj},
defined there as a product over the primes dividing $q_0$, depends on $q_0$ only, and $\delta=
c_\psi\sqrt{\mathsf m}$ with $c_\psi=\sqrt{2\boldsymbol\pi/\sqrt d}$. (b) This is the last line of
the proof of Proposition~\ref{pk:conj}. (c) The number $c_1=-1-\gamma+\Rel\frac{L'}L(1,\chi_K)$
depends on $d$ only. Definition~\ref{an:def-param} and \eqref{pz:eq-C0} precede and do not involve
the data, so there is no circularity: first $(q_0,d)$, then $C_P$, then the conditions (D1)--(D8).
By (D2), $10^{12}C_P\le\log\log b$, and $C_0\le C_P$, $d\le C_P$.
\end{proof}

Consequently, in the proof of Proposition~\ref{au:kernel} the terms $C_0(q_1nn'+1)^2+fnn'\log(2L+2)$
are at most $x^{50}-1$: by Lemma~\ref{an:lem-sizes}, $q_1\le x^{12}$, $f\le x^{12}$, $nn'\le
n_m^2\le1600x^{12}$, $\log\log b\le\tfrac1{100}\log x$ and $\log(2L+2)\le3\log b\le3x^{1/100}$. The
existence theorem for admissible data (Theorem~\ref{an:thm-data}) uses of $C_P$ only that it
depends on $(q_0,d)$ alone, and is not affected.

\subsubsection{The levels of the upper chain; sizes; short vectors}\label{pz:ss-levels}

\emph{The levels.} For $0\le j\le m$ the level $J_{X_j}$ of $X_j$ is defined as $J_X$ is for
$j=0$: it is the group of the $g$ in the stabiliser $K^\circ$ of $\widehat\auO_E^{\,n_j}$ in
$\auU(X_j)(\auA_{E^+,f})$ that are $\equiv1$ modulo every chosen place and whose last row is
$\equiv(0,\dots,0,1)$ modulo $\pi_N$. This is the definition used in Step~5 of the proof of
Proposition~\ref{au:seesaw} and in \S\ref{pl:ss-level}. For the trivial twist the level of
$\mathcal A_m(1)$ is the group $J_1$ of the elements of the stabiliser of $\psM_m=\widehat\auO_E^{\,
n_m-1}$ in $\auU(Y_m)(\auA_{E^+,f})$ that are $\equiv1$ modulo every chosen place.

\begin{lemma}\label{pz:sizes}
\begin{enumerate}
\item[(a)] The levels $J_{X_j}$ $(0\le j\le m)$ and $J_1$ contain the principal congruence subgroup
$K(MN)$ of the corresponding standard space, of rank $n_j\le n_m$, respectively $n_m-1$; they are
free.
\item[(b)] $n_m\le40x^6$, $6q_1d(fk)^2\le x^{100}$, and $6q_1d(fn)^2\le x^{102}$ for $n\le n_m$.
\item[(c)] $\log(MN)\le x$.
\end{enumerate}
So the hypotheses of Proposition~\ref{pf:classes} hold for these levels, and each of them has at
most $\exp(x^{30})$ classes.
\end{lemma}
\begin{proof}
(a) The chosen places divide $\pi_M$, which divides $M$, and $\pi_N$ divides $N$; so an element
$\equiv1$ modulo $MN$ satisfies all the congruences. The levels are free by Lemma~\ref{au:free}.
(b) By Lemma~\ref{an:lem-sizes}, $q_1\le x^{12}$, $f\le x^{12}$, $m\le2\sqrt x$, $x^5/2\le k\le r\le
3x^5$, and $d\le\log\log b\le\log x$. So $n_m=k+m(5r+1)\le3x^5+2\sqrt x\,(15x^5+1)\le40x^6$,
$6q_1d(fk)^2\le54\,x^{46}\log x\le x^{100}$, and for $n\le n_m$, $6q_1d(fn)^2\le6\cdot1600\,x^{48}
\log x\le x^{102}$. (c) $\log M\le22\log b$ and $\log N\le25m\log b$ (Lemma~\ref{an:lem-sizes}(S5)),
and $(22+50\sqrt x)\log b\le x$ since $(\log b)^{100}\le x$.
\end{proof}

\emph{The estimate in the proof of Proposition~\ref{au:short}.} The proof of
Proposition~\ref{au:short} uses the following bound.

\begin{lemma}\label{pz:short}
Let $Y\in\auGL_k(E)$ have entries all of whose conjugates are at most $C\ge k$ in absolute value,
$\zeta=(\zeta_\tau)_{\tau\in\auPa}$ vectors of $\C^k$ whose coordinates are roots of unity, and
$\xi=Y\mathrm{ev}(\zeta)\in W_{\lambda^X}$ as in the proof of Proposition~\ref{au:short}. Then
$\|\sigma\xi\|^2\le C^{\,5fa}$ for all $\sigma\in\auAut(\C)$.
\end{lemma}
\begin{proof}
$W_{\lambda^X}=\bigotimes_{\tau\in\auPa}V_\tau$, where $V_\tau$ is the dual of the space of forms of
degree $a$ in $k$ variables, with the basis $(e_\beta)$ dual to the monomials. For $w\in\C^k$ the
evaluation $\mathrm{ev}_w$ at $w$ is $\sum_\beta w^\beta e_\beta$. By \S\ref{ps:ss-restrict} the
Gram matrix of $(e_\beta)$ is the matrix of the coefficients of the reproducing kernel
$(\sum_ix_i\bar y_i)^a$ in the products of monomials $x^\beta\bar y^{\beta'}$; hence the value of the
hermitian form at a vector with coordinates $(w^\beta)_\beta$ is $(\sum_i|w_i|^2)^a$. An element
$g$ of $\auGL_k$ acts on $V_\tau$ by $\mathrm{ev}_w\mapsto\mathrm{ev}_{g'w}$, where $g'$ is $g$, its
transpose, or the complex conjugate of one of these, according to the conventions (N1), (N2) of
\S\ref{pl:ss-norm}. So the coordinates of $\xi$ are the numbers $\prod_\tau w_\tau^{\beta_\tau}$ with
$w_\tau=\tau(Y)'\zeta_\tau$, the coordinates of $\sigma\xi$ are $\prod_\tau(\sigma w_\tau)^{\beta_\tau}$,
and the entries of $\sigma w_\tau$ are sums of $k$ products of a conjugate of an entry of $Y$ and a
root of unity, of absolute value at most $kC$. Therefore
\[
\|\sigma\xi\|^2=\prod_{\tau\in\auPa}\Bigl(\sum_i|(\sigma w_\tau)_i|^2\Bigr)^a\le(k\cdot k^2C^2)^{a\#\auPa}
\le C^{\,5fa}.\qedhere
\]
\end{proof}
With $C=C_{\mathrm{sb}}=6q_1d(fk)^2\ge k$ this is the estimate needed: $4^fC_{\mathrm{sb}
}^{\,5fa}\le\exp(f\log4+500fa\log x)\le e^{\aueu}$ by Lemma~\ref{pz:sizes}(b), since
$\aueu=fa(\log b)^3$ and $\log x\le\log b$.

\subsubsection{The inputs of Theorem~A}\label{pz:ss-audit}

We list the statements used in the proof of Theorem~\ref{thm:srcA}, from the top, with the place of
the written proof. ``Published'' means that the proof quotes results from the literature, named at
the place of use.
\begin{itemize}
\item \emph{The deduction} (\S\ref{ss:au-proofAB}) uses Lemmas~\ref{au:gains} and~\ref{au:errsum},
which are proved in \S\ref{sec:automorphic}, and the inputs (I1)--(I6). The objects in the
statement are defined in \S\ref{ss:au-spaces}, with the complements of \S\ref{ps:ss-weil},
\S\ref{ps:ss-Y}, \S\ref{ps:ss-restrict}, \S\ref{ps:ss-glue}, \S\ref{pl:ss-norm} and of
Definition~\ref{pz:def-type}, Lemma~\ref{pz:mu} and \S\ref{pz:ss-levels}; that the lower chain and
the Poincar\'e map are well defined is Proposition~\ref{pz:types}, which rests on (A1)--(A3).
\item \emph{{\rm(I1)}, Proposition~\ref{au:seesaw}}: \S\ref{ps:ss-seesaw}, from
Lemma~\ref{ps:weil} (W0)--(W5) (published: Weil, Kudla, Harris--Kudla--Sweet), (W6)
(Proposition~\ref{pf:flip}), Propositions~\ref{ps:onestep} (\S\ref{pk:ss-onestep}; Witt's theorem),
\ref{pk:definite} (Cauchy decomposition, published), Lemmas~\ref{pk:chi0}, \ref{ps:level},
\ref{ps:double}, \ref{ps:support} (with (R1), (R2) and Lemma~\ref{pf:orbits}; Iwasawa decomposition,
published), \ref{ps:fin} (with (R3), (R4)), and Proposition~\ref{pz:types} for the types, the
cuspidality and the rapid decrease of the forms $z_j$.
\item \emph{{\rm(I2)}, Proposition~\ref{au:divis}}: Theorem~\ref{pa:divis}, from
Propositions~\ref{pa:count}, \ref{pa:product}, Theorem~\ref{pa:onecurve} (published results on
elliptic curves with complex multiplication), Lemma~\ref{pa:kernel} (Propositions~\ref{pk:K0},
\ref{pk:std}, \ref{pk:K12}, Theorem~\ref{pf:std}, Corollary~\ref{pf:K1}), Lemma~\ref{pa:hecke}
(\S\ref{pl:ss-hecke}), Lemma~\ref{pa:eldiv} (Theorem~\ref{ph:thm}, \S\ref{pe:sec}),
Lemma~\ref{pl:level} with Lemma~\ref{pl:kernelN} (Corollary~\ref{pf:K1}(b)),
Proposition~\ref{au:expansion} (\S\ref{ps:ss-glue}) and Proposition~\ref{pk:branching}.
\item \emph{{\rm(I3)}, Proposition~\ref{au:short}}: \S\ref{ss:au-inputsA}, from
Lemma~\ref{au:shortbasis} (\S\ref{ps:ss-short}) and Lemma~\ref{pz:short}.
\item \emph{{\rm(I4)}, Proposition~\ref{au:size}}: \S\ref{ss:au-inputsA}, from
Proposition~\ref{au:kernel} (\S\ref{ps:ss-kernel}; Lemma~\ref{ps:conj} is Proposition~\ref{pk:conj}
with Theorem~\ref{pf:recip}; the constant is Lemma~\ref{pz:C0}; Chowla--Selberg formula,
published), Proposition~\ref{au:restrict} (\S\ref{ps:ss-restrict}, Proposition~\ref{pk:branching}),
Proposition~\ref{au:expansion}, and Proposition~\ref{pf:classes} with Lemma~\ref{pz:sizes}.
\item \emph{{\rm(I5)}, Proposition~\ref{au:good}}: from Proposition~\ref{au:equiv} (\S\ref{ps:ss-Y})
and Proposition~\ref{pf:classes} with Lemma~\ref{pz:sizes}.
\item \emph{{\rm(I6)}, Lemma~\ref{au:ramcount}}: proved in \S\ref{ss:au-inputsB}.
\end{itemize}
\emph{What is not proved in this paper, for Theorem~\ref{thm:srcA}.} After this subsection the proof
of Theorem~\ref{thm:srcA} is written out: the archimedean statements (A1), (A2), (A3) of
\S\ref{pz:ss-types} are proved in \S\ref{par:sec} and \S\ref{pco:sec}
(Theorems~\ref{par:A1}, \ref{par:A3}, \ref{pco:FJ}, \ref{pco:V}, Corollary~\ref{par:A2}) from the
published results named there: (D1)--(D3) on discrete series, and the local properties
(W1$_\infty$), (W2$_\infty$), (M), (S) of the archimedean Weil representation. They enter through
Proposition~\ref{pz:types} only: in the definition of
$\iota_\vartheta$ and $\auP$, and in Proposition~\ref{au:seesaw}. All statements have a
written proof, complete modulo the published results named in them. The explicit form
\eqref{pl:eq-w5} of (W5), with its factor, is Proposition~\ref{pwf:W5}. The absolute convergence and the bound in (W3) are proved here
(Lemma~\ref{pz:absgrowth}).


%% ===== proofs/A-archimedean.tex =====
% proofs/A-archimedean.tex
% The archimedean statements (A1), (A2), (A3) of \S pz:ss-types.
% Labels of this file start with par:. Only amsmath/amssymb/amsthm are used.
\providecommand{\auU}{\mathrm U}
\providecommand{\auPl}{\Phi}
\providecommand{\auHom}{\operatorname{Hom}}
\providecommand{\psrho}{\varrho}
\providecommand{\pzF}{\mathfrak F}
\providecommand{\parsgn}{\operatorname{sgn}}

\subsection{The archimedean statements (A1)--(A3)}\label{par:sec}

This subsection treats the three statements (A1)--(A3) of \S\ref{pz:ss-types}. \emph{Dependencies.}
(A2) is proved from the decay bound \eqref{par:eq-decay} below (Corollary~\ref{par:A2}). (A1) is
proved from the same bound (Theorem~\ref{par:A1}); the bound rests on the explicit formula
\eqref{pds:eq-FJ}, which is proved in Theorem~\ref{pco:FJ} (Remark~\ref{par:FJ}). (A3) is proved
(Theorem~\ref{par:A3}) modulo the published theorems (D1)--(D3) of \S\ref{par:ss-rigid} and the
normalisation (V) of the vacuum character (Lemma~\ref{par:vac}), which is proved in
Theorem~\ref{pco:V} from the local properties (W1$_\infty$), (W2$_\infty$), (M), (S) of the
archimedean Weil representation.

\emph{Conventions.} We fix $\tau\in\auPl$ and $0\le j\le m$. $G=G_{j,\tau}$ is identified with
$\auU(p,q)=\{g:g^*Jg=J\}$, $J=\operatorname{diag}(\varepsilon_1,\dots,\varepsilon_n)$, $n=N_j$, and
$K=G\cap\auU(n)$; $T\subset K$ is the diagonal torus and $\mathfrak g_\C=\mathfrak{gl}_n(\C)$ with
matrix units $E_{ab}$. For a parameter $\lambda\in(\tfrac{n-1}2+\Z)^n$ with distinct entries,
the positive roots are the $e_a-e_b$ with $\lambda_a>\lambda_b$; $\rho$, $\rho_{\mathrm c}$,
$\rho_{\mathrm n}$ are as in \S\ref{pds:ss-setup}, $\Lambda^K(\lambda)=\lambda+\rho_{\mathrm n}-
\rho_{\mathrm c}$, and $\pi_\lambda$ is the discrete series with Harish-Chandra parameter $\lambda$.
We write $\Omega=\sum_{a,b}E_{ab}E_{ba}$; it acts on a representation with infinitesimal character
$\lambda$ by $\kappa_\lambda=|\lambda|^2-|\rho|^2$. The \emph{gap} of $\lambda$ is
$\min\{|\lambda_a-\lambda_b|:\varepsilon_a\ne\varepsilon_b\}$. By \eqref{pz:eq-gap} and
\S\ref{ss:au-notation} ($r\ge x^5$, $N_j\le40x^6$) the parameters $\Lambda^{(j)}_\tau$ have gap at
least $\tfrac12r^8\ge\tfrac12x^{40}$, which exceeds $500(N_j+N_{j+1})^3+10N_j^5+x^{25}$.

\emph{The coefficient.} Inner products are conjugate-linear in the first variable, as in
\S\ref{ps:ss-weil}. We read the coefficient of \S\ref{pz:ss-types} as
\begin{equation}\label{par:eq-coef}
c_j(x)=\langle v_j,\pi_j(x)v_j\rangle,\qquad e_j=d_j\,\overline{c_j}.
\end{equation}
With this reading, Schur's orthogonality relations give, for every unitary representation
$\sigma$, that $\sigma(e_j)=\int e_j(x)\sigma(x)\,dx$ is the orthogonal projector onto the vectors
of type $\pi_j$ that correspond to $v_j$. This is the property used in \S\ref{pz:ss-types}. (With
the other reading, $\sigma(e_j)$ projects onto the type $\overline{\pi_j}$, and the rule
\eqref{pz:eq-param} would no longer agree with Lemma~\ref{par:vac}; see Remark~\ref{par:rem-match}.)

\subsubsection{The decay bound}\label{par:ss-decay}
We say that $\lambda$ \emph{satisfies} \textup{(Dec$_c$)} if the coefficient
$c_\lambda(g)=\langle v,\pi_\lambda(g)v\rangle$, with $v$ a unit highest weight vector of the
lowest $K$-type, satisfies
\begin{equation}\label{par:eq-decay}
|c_\lambda(g)|\le(2/\|g\|)^{c}\qquad(g\in G).
\end{equation}

\begin{remark}\label{par:FJ}
Theorem~\ref{pds:decay} proves \eqref{par:eq-decay} for every integer $c$ up to the gap, for any
pair $(\lambda,\varepsilon)$, \emph{assuming} the integral formula \eqref{pds:eq-FJ} for $c_\lambda$.
The proof of Theorem~\ref{pds:decay} (Lemmas~\ref{pds:blockdet} and \ref{pds:fischer}) uses only
the formula and the gap, not the signature of $Y_0$: Lemma~\ref{pds:blockdet} uses only $g^*Jg=J$ and that the maximal compact subgroup commutes with $J$; Lemma~\ref{pds:fischer} uses only Fischer's inequality, the formula and $c\le h$; and the displayed estimate in the proof of Theorem~\ref{pds:decay} uses only \eqref{pds:eq-FJ} and $\det a(u^{-1}zu)=\det a(z)$ for $u$ in the maximal compact subgroup. Hence Theorem~\ref{pds:decay} and its proof hold verbatim with $N$ replaced by $n=N_j$ and $J$ by the diagonal matrix of signs of $Y_{j,\tau}$ (of any signature), for every regular $\lambda\in(\frac{n-1}2+\Z)^n$ with gap $h$ for which \eqref{pds:eq-FJ} holds.
Formula \eqref{pds:eq-FJ} is Theorem~\ref{pco:FJ}, proved for every regular $\lambda\in
(\frac{n-1}2+\Z)^n$ with gap $h\ge n^2+1$, in particular for every $\Lambda^{(j)}_\tau$: steps (i)
(Lemma~\ref{pco:K}, Proposition~\ref{pco:casimir}), (ii) (Lemma~\ref{pco:decay}) and (iii)
(Proposition~\ref{pco:rigid} and the proof of Theorem~\ref{pco:FJ}) are carried out there. Hence every
$\Lambda^{(j)}_\tau$ satisfies (Dec$_c$) for every integer $c$ up to the gap.
\end{remark}

\subsubsection{(A1) and (A2)}\label{par:ss-A12}

\begin{theorem}[(A1)]\label{par:A1}
Assume that every $\Lambda^{(j)}_\tau$ satisfies \textup{(Dec$_c$)} with $c=B+N_j^2+1$, $B=x^{25}$.
Then for $0\le j\le m$ the function $c_j=\bigotimes_\tau c_{j,\tau}$ is smooth, left and right
$K_{j,\infty}$-finite, an eigenfunction of the centre of the enveloping algebra, and
$\int_{G_{j,\infty}}|\mathsf Xc_j(x)|\,\|x\|^B\,dx<\infty$ for every left- or right-invariant
differential operator $\mathsf X$.
\end{theorem}
\begin{proof}
Smoothness, $K$-finiteness and the eigenfunction property hold for every $K$-finite coefficient
of an irreducible unitary representation: $K$-finite vectors are smooth (indeed analytic)
\cite[Proposition~8.5, p.~208]{Knapp}, and the centre acts on the smooth vectors of an irreducible representation
by scalars. Fix $\tau$ and write $c=c_{j,\tau}$, $e=d\bar c$. By Lemma~\ref{pm:kernel}(b), whose proof
uses only that $c$ is a $K$-finite coefficient of an irreducible unitary representation, there
is $\alpha\in C_c^\infty(G)$ with $e=\alpha*e=e*\alpha^*$; so $\mathsf Xe=(\mathsf X\alpha)*e$ for
right-invariant and $\mathsf Xe=e*(\mathsf X\alpha^*)$ for left-invariant $\mathsf X$. With
$\|gh\|\le\|g\|\|h\|$ and \eqref{par:eq-decay},
$|\mathsf Xe(g)|\le C_{\mathsf X}\|g\|^{-c}\int\|h\|^c|\mathsf X\alpha(h)|dh$. By
Lemma~\ref{pm:kernel}(c), $\int_G\|g\|^{B-c}dg<\infty$ since $c-B>N_j^2$. The norm $\|x\|$ of
$G_{j,\infty}$ is the product of the local norms and $c_j$ is a product, so Fubini gives the
statement on $G_{j,\infty}$ (an invariant operator on the product is a sum of products of local
ones).
\end{proof}

\begin{corollary}[(A2)]\label{par:A2}
Assume \textup{(Dec$_c$)} with $c=5N_j^3$ for every $\Lambda^{(j)}_\tau$. Let $\mathsf P=\prod_\tau
\mathsf P_\tau$ be a parabolic subgroup of $G_{j,\infty}$ with $\mathsf P_\tau$ proper for at least
one $\tau$, and $\mathsf N$ its unipotent radical. Then $\int_{\mathsf N}|c_j(anb)|dn<\infty$ and
$\int_{\mathsf N}c_j(anb)\,dn=0$ for all $a,b$.
\end{corollary}
\begin{proof}
The proof of Lemma~\ref{pm:cusp} uses, for the function $e_\tau$, only: (a) the bound
$|e_\tau(g)|\le C\|g\|^{-5n^3}$; (b) the derivative bounds of Lemma~\ref{pm:kernel}(b); (c) that
$e_\tau$ is an eigenfunction of the centre. Here (a) is \eqref{par:eq-decay}, (b) follows from it as
in the proof of Theorem~\ref{par:A1}, and (c) holds as there. The signature and the gap condition
\eqref{pm:eq-gap} enter Lemma~\ref{pm:cusp} only through (a). So for each $\tau$ with $\mathsf
P_\tau$ proper, $\int_{\mathsf N_\tau}|c_{j,\tau}(a_\tau nb_\tau)|dn<\infty$ and the integral vanishes;
for $\mathsf P_\tau=G_{j,\tau}$, $\mathsf N_\tau=1$. Since $c_j$ is a product and $\mathsf
N=\prod_\tau\mathsf N_\tau$, Fubini gives absolute convergence, and the integral is the product of
the local integrals, one of which is $0$.
\end{proof}
This is stronger than (A2), which assumes every $\mathsf P_\tau$ proper. The hypotheses of
Theorem~\ref{par:A1} and Corollary~\ref{par:A2} hold as soon as \eqref{pds:eq-FJ} holds, since the gap
exceeds both exponents.

\subsubsection{Rigidity of the lowest $K$-type}\label{par:ss-rigid}
We use three published results, for $G=\auU(p,q)$ and a regular parameter $\lambda$ as above.
\begin{enumerate}
\item[(D1)] \emph{(Dirac inequality; Parthasarathy's formula for the square of the Dirac operator \cite[Proposition~3.1]{Parthasarathy72}, there in the geometric setting of spinors on $G/K$, in the algebraic form of \cite[Proposition~1.1 and \S2]{HuangPandzic}.)} Let $X$ be an
irreducible unitary $(\mathfrak g,K)$-module on which $\Omega$ acts by $\omega_X$, $S$ the spin module
of $\mathfrak p$, $\mu$ a $K$-type of $X$ and $\tau_K$ a $K$-type of $\mu\otimes S$. Then
$\omega_X+|\rho|^2\le|\tau_K+\rho_{\mathrm c}|^2$; if equality holds, the Dirac operator kills the
$\tau_K$-isotypic part of $X\otimes S$, so the Dirac cohomology of $X$ contains $\tau_K$. Moreover, for
$\mu=\Lambda^K(\lambda)$, the $K$-type $\tau_K=\mu-\rho_{\mathrm n}=\lambda-\rho_{\mathrm c}$ occurs in
$\mu\otimes S$ (same reference).
\item[(D2)] \emph{(Huang--Pand\v zi\'c \cite[Theorem~2.3]{HuangPandzic}.)} If
the Dirac cohomology of an irreducible $(\mathfrak g,K)$-module with infinitesimal character
$\chi$ contains $\tau_K$, then $\chi$ is $W$-conjugate to $\tau_K+\rho_{\mathrm c}$.
\item[(D3)] \emph{(Lowest $K$-type and cohomology.)} Let $\Psi=\{\alpha:\langle\alpha,\lambda\rangle>0\}$; $K$-types are labelled by highest weights for $\Psi_{\mathrm c}=\Psi\cap\Delta(\mathfrak k)$.
(a) \emph{(Schmid \cite[Theorem~(1.3)]{Schmid75}.)} Up to infinitesimal equivalence, $\pi_\lambda$ is the only irreducible representation that contains the $K$-type $\Lambda^K(\lambda)$ and contains no $K$-type of highest weight $\Lambda^K(\lambda)-A$, $A$ a non-empty sum of elements of $\Psi$.
(b) \emph{(Schmid \cite[Corollary~(3.9), Corollary~(4.19) and the proof of Theorem~(1.3) on p.~560]{Schmid75}.)} $\Lambda^K(\lambda)$ occurs in $\pi_\lambda$ with multiplicity one.
(c) \emph{(Vogan--Zuckerman \cite[Proposition~5.4, Theorems~5.3 and~5.6]{VoganZuckerman}, for a connected semisimple group $G_1$ with finite centre, maximal compact $K_1$, $\mathfrak g_1=\mathfrak k_1+\mathfrak p_1$.)} Let $X$ be the Harish-Chandra module of an irreducible unitary representation of $G_1$ and $F$ an irreducible finite-dimensional representation of $G_1$ such that the Casimir operator acts by the same scalar on $X$ and on $F$. Then $H^*(\mathfrak g_1,K_1;X\otimes F)\cong\operatorname{Hom}_{K_1}(\wedge^*\mathfrak p_1,X\otimes F)$; if this is non-zero, there are a $\theta$-stable parabolic subalgebra $\mathfrak q=\mathfrak l+\mathfrak u$, defined by an element $x_{\mathfrak q}\in i\mathfrak t_{1}$ ($\mathfrak t_1$ a Cartan subalgebra of $\mathfrak k_1$; $\mathfrak u$, $\mathfrak l$ the sums of the positive, zero eigenspaces of $\operatorname{ad}x_{\mathfrak q}$), and a weight $\lambda_{\mathfrak q}$, the differential of a unitary character of the centraliser $L$ of $x_{\mathfrak q}$, with $\langle\alpha,\lambda_{\mathfrak q}\rangle\ge0$ for $\alpha\in\Delta(\mathfrak u)$, such that, for any positive system $\Delta^+(\mathfrak l\cap\mathfrak k_1)$, with $\Delta^+(\mathfrak k_1)=\Delta^+(\mathfrak l\cap\mathfrak k_1)\cup\Delta(\mathfrak u\cap\mathfrak k_1)$ used to label $K_1$-types: every $K_1$-type of $X$ has highest weight $\lambda_{\mathfrak q}+2\rho(\mathfrak u\cap\mathfrak p_1)+\sum_{\beta\in\Delta(\mathfrak u\cap\mathfrak p_1)}n_\beta\beta$, $n_\beta\in\Z_{\ge0}$; and the centre of $U(\mathfrak g_1)$ acts on $X$ by the character with Harish-Chandra parameter $\lambda_{\mathfrak q}+\rho'$, $\rho'$ the half-sum of $\Delta^+(\mathfrak l)\cup\Delta(\mathfrak u)$ for any positive system $\Delta^+(\mathfrak l)\supseteq\Delta^+(\mathfrak l\cap\mathfrak k_1)$ (here $\operatorname{rk}\mathfrak k_1=\operatorname{rk}\mathfrak g_1$, so $\mathfrak t_1$ is a Cartan subalgebra of $\mathfrak g_1$).
\end{enumerate}
We also use that $G$ is of type I, so every unitary representation on a separable Hilbert space
is a direct integral of irreducible ones \cite{Dixmier}, and that irreducible unitary
representations are admissible with smooth $K$-finite vectors \cite[Theorem~8.1 and Proposition~8.5]{Knapp}. (D1), (D2)
are stated for reductive groups with these normalisations in \cite[Proposition~1.1 and \S2]{HuangPandzic}; the centre of
$\auU(p,q)$ acts by a character and does not affect them.

\emph{On (D3).} Schmid's paper is written for a connected semisimple linear group with $\operatorname{rk}K=\operatorname{rk}G$, and covers \emph{all} discrete series, not only the holomorphic ones; it applies to $\auU(p,q)=Z\cdot\mathrm{SU}(p,q)$ ($Z$ the compact centre, acting by a character) through $\mathrm{SU}(p,q)$. The discrete series used in this paper are in general neither holomorphic nor antiholomorphic: in \eqref{pz:eq-param} the signs of the entries $-\kappa(\tau)-i\Delta$ alternate, so there are positive noncompact roots of both kinds $e_+-e_-$ and $e_--e_+$; hence no elementary (generalised Verma module) argument replaces Schmid's theorem here. What \cite{Schmid75} proves is: (i) $\pi_\lambda$ contains $\Lambda^K(\lambda)$ exactly once, and every $K$-type of $\pi_\lambda$ has highest weight $\Lambda^K(\lambda)+\sum n_\beta\beta$, $n_\beta\in\Z_{\ge0}$, over the noncompact roots $\beta$ with $\langle\beta,\lambda\rangle>0$ (Lemmas~(2.22), (3.8), Corollary~(3.9) and, for the identification of $\pi_\lambda$ with the irreducible quotient of the universal module, Corollary~(4.19)); (ii) is (D3)(a). Schmid's characterisation is by $K$-types, not by the infinitesimal character; the form used below (an irreducible unitary module with infinitesimal character $\lambda$ containing $\Lambda^K(\lambda)$ is $\pi_\lambda$) is Lemma~\ref{par:D3}, deduced from (D3)(a) and (D3)(c). In \cite{VoganZuckerman}, $G$ is connected semisimple with finite centre; Theorems~5.3 and~5.6 and Proposition~5.4 are in \S5, where the authors indicate the proofs as extensions of those given in full for trivial coefficients (Theorems~2.5 and~4.1). The positive system $\Delta^+(\mathfrak g)=\Delta^+(\mathfrak l)\cup\Delta(\mathfrak u)$ entering the infinitesimal character is the one fixed before (2.6) there, and the Harish-Chandra parametrisation of infinitesimal characters is (2.9) there: the characters with parameters $\nu,\nu'$ agree if and only if $\nu'\in W\nu$.

\begin{lemma}\label{par:D3}
Let $X$ be the Harish-Chandra module of an irreducible unitary representation of $G=\auU(p,q)$ with infinitesimal character $\lambda$ \textup{(}regular, $\lambda\in(\frac{n-1}2+\Z)^n$\textup{)} that contains the $K$-type $\Lambda^K(\lambda)$. Then every $K$-type of $X$ has highest weight $\Lambda^K(\lambda)+\sum n_\beta\beta$ with $n_\beta\in\Z_{\ge0}$ and $\beta$ running over noncompact roots in $\Psi$; in particular $X$ contains no $K$-type $\Lambda^K(\lambda)-A$ as in \textup{(D3)(a)}, and $X$ is the Harish-Chandra module of $\pi_\lambda$.
\end{lemma}
\begin{proof}
Write $\Lambda=\Lambda^K(\lambda)$. If $p=0$ or $q=0$, $G$ is compact, there are no noncompact roots, $\Lambda=\lambda-\rho$, and $X$ is the irreducible representation with highest weight $\Lambda$, which is $\pi_\lambda$. Let $p,q\ge1$.

\emph{Reduction to $G_1=\mathrm{SU}(p,q)$.} $G=Z\cdot G_1$ with $Z$ the centre, which acts on $X$ by a character (Schur's lemma); so $X$ is the Harish-Chandra module of an irreducible unitary representation of $G_1$, for $K_1=K\cap G_1$, and $K$-types of $G$ restrict to irreducible $K_1$-types. We work with $\mathfrak t_1=\mathfrak t\cap\mathfrak g_1$ and restrict all weights to it; roots, $\Psi$, $\rho$, $\rho_{\mathrm c}$, $\rho_{\mathrm n}$ live on $\mathfrak t_1$ already, and the conclusion for $G$ follows from the one for $G_1$, since a $K$-type is determined by its restriction to $K_1$ and the central character, which is common to all $K$-types of $X$.

\emph{Non-vanishing of cohomology.} Since $\lambda\in(\frac{n-1}2+\Z)^n$ and $\rho\in(\frac{n-1}2+\Z)^n$, $\lambda-\rho\in\Z^n$; for $\alpha\in\Psi$ simple, $\langle\lambda-\rho,\alpha^\vee\rangle=\langle\lambda,\alpha^\vee\rangle-1\ge0$, because $\langle\lambda,\alpha^\vee\rangle$ is a positive integer. Let $F^*$ be the irreducible finite-dimensional representation of $G$ (restricted from $\mathrm{GL}_n(\C)$) with $\Psi$-highest weight $\lambda-\rho$, and $F$ its dual. As $\mathfrak g_1$-modules ($\mathfrak g_\C=\mathfrak g_{1,\C}\oplus\C\cdot1$, and $Z(U(\mathfrak g_\C))=Z(U(\mathfrak g_{1,\C}))\otimes\C[1]$), $X$ and $F^*$ both have infinitesimal character with parameter $\lambda|_{\mathfrak t_1}$, so the Casimir operator of $\mathfrak g_1$ acts on them by the same scalar; it acts on $F$ as on $F^*$ (its scalar on the module with highest weight $\nu$ is $|\nu+\rho|^2-|\rho|^2$, and $F$ has highest weight $-w_0(\lambda-\rho)$, $w_0$ the longest element, with $|-w_0(\lambda-\rho)+\rho|=|\lambda-\rho+\rho|$ as $-w_0\rho=\rho$). So the Casimir operator acts on $F$ as on $X$. Let $\mathfrak u_{\mathrm n}=\bigoplus_{\beta\in\Psi_{\mathrm n}}\mathfrak g_\beta\subset\mathfrak p$, $R=\dim\mathfrak u_{\mathrm n}$. Since $\Psi$ is closed under addition of roots, $[\mathfrak g_\alpha,\mathfrak u_{\mathrm n}]\subset\mathfrak u_{\mathrm n}$ for $\alpha\in\Psi_{\mathrm c}$, so the line $\wedge^R\mathfrak u_{\mathrm n}\subset\wedge^R\mathfrak p$ has weight $2\rho_{\mathrm n}$ and is killed by every $\mathfrak g_\alpha$, $\alpha\in\Psi_{\mathrm c}$. With $v$ a highest weight vector of $F^*$, $e=(\wedge^R\mathfrak u_{\mathrm n})\otimes v$ is therefore a $\Psi_{\mathrm c}$-highest weight vector of weight $2\rho_{\mathrm n}+\lambda-\rho=\Lambda$ in $\wedge^R\mathfrak p\otimes F^*$, and generates a $K$-type $\Lambda$. As $X$ contains $\Lambda$, $\operatorname{Hom}_{K_1}(\wedge^R\mathfrak p\otimes F^*,X)\cong\operatorname{Hom}_{K_1}(\wedge^R\mathfrak p,X\otimes F)\ne0$. By (D3)(c), $H^R(\mathfrak g_1,K_1;X\otimes F)\ne0$, and we obtain $\mathfrak q=\mathfrak l+\mathfrak u$, $x_{\mathfrak q}$ and $\lambda_{\mathfrak q}$ as there.

\emph{Normalising $\mathfrak q$.} All Cartan subalgebras of $\mathfrak k_1$ are $K_1$-conjugate, and conjugating $(\mathfrak q,x_{\mathfrak q},\lambda_{\mathfrak q},\Delta^+(\mathfrak l\cap\mathfrak k_1))$ by an element $k\in K_1$ transforms the conclusions of (D3)(c) into the same conclusions for the conjugated data (a $K_1$-type with highest weight $\delta$ for a positive system $P$ has highest weight $w\delta$ for $wP$, $w$ the image of $k$ in the Weyl group when $k$ normalises $\mathfrak t_1$). So we may assume $x_{\mathfrak q}\in i\mathfrak t_1$ and, choosing $\Delta^+(\mathfrak l\cap\mathfrak k_1)$ and then conjugating by $N_{K_1}(\mathfrak t_1)$, that $\Delta^+(\mathfrak l\cap\mathfrak k_1)\cup\Delta(\mathfrak u\cap\mathfrak k_1)=\Psi_{\mathrm c}$. Then $\alpha(x_{\mathfrak q})>0$ for $\alpha\in\Delta(\mathfrak u)$, $\alpha(x_{\mathfrak q})=0$ for $\alpha\in\Delta(\mathfrak l)$, and $\Psi_{\mathrm c}\cap\Delta(\mathfrak u)=\Delta(\mathfrak u\cap\mathfrak k_1)$, $\Psi_{\mathrm c}\cap\Delta(\mathfrak l)=\Delta^+(\mathfrak l\cap\mathfrak k_1)$.

\emph{The comparison.} Fix $\Delta^+(\mathfrak l)\supseteq\Delta^+(\mathfrak l\cap\mathfrak k_1)$, put $P'=\Delta^+(\mathfrak l)\cup\Delta(\mathfrak u)$, $\rho'$ its half-sum, and $\lambda'=\lambda_{\mathfrak q}+\rho'$. For $\alpha$ simple in $P'$, $\langle\lambda_{\mathfrak q},\alpha\rangle\ge0$ (it is $0$ for $\alpha\in\Delta(\mathfrak l)$, since $\lambda_{\mathfrak q}$ is the differential of a character of $L$, and $\ge0$ for $\alpha\in\Delta(\mathfrak u)$), and $\langle\rho',\alpha^\vee\rangle=1$; so $\lambda'$ is $P'$-dominant, and $x_{\mathfrak q}$ is $P'$-dominant ($\alpha(x_{\mathfrak q})\ge0$ on $P'$). By (D3)(c) and (2.9) of \cite{VoganZuckerman}, $\lambda=w\lambda'$ for some $w$ in the Weyl group $W$. For dominant $\lambda'$, $\lambda'-w\lambda'$ is a sum of elements of $P'$ with non-negative integer coefficients (induction on the length: if $w=s_\alpha w_1$, $\alpha$ simple, $\ell(w)=\ell(w_1)+1$, then $w_1^{-1}\alpha\in P'$ and $\lambda'-w\lambda'=(\lambda'-w_1\lambda')+\langle\lambda',w_1^{-1}\alpha^\vee\rangle\alpha$). Hence
\begin{equation}\label{par:eq-D3a}
\lambda(x_{\mathfrak q})\le\lambda'(x_{\mathfrak q})=\lambda_{\mathfrak q}(x_{\mathfrak q})+\rho(\mathfrak u)(x_{\mathfrak q}),
\end{equation}
as $\rho'-\rho(\mathfrak u)$ is a combination of roots of $\mathfrak l$, which vanish at $x_{\mathfrak q}$. On the other hand $\Lambda$ is a $K_1$-type of $X$, so by (D3)(c), $\Lambda=\lambda_{\mathfrak q}+2\rho(\mathfrak u\cap\mathfrak p_1)+B_{\mathfrak q}$ with $B=\sum_{\beta\in\Delta(\mathfrak u\cap\mathfrak p_1)}n_\beta\beta$, $n_\beta\ge0$. Now $\lambda=\Lambda+\rho_{\mathrm c}-\rho_{\mathrm n}$, and: $2\rho_{\mathrm c}(x_{\mathfrak q})=\sum_{\alpha\in\Psi_{\mathrm c}}\alpha(x_{\mathfrak q})=2\rho(\mathfrak u\cap\mathfrak k_1)(x_{\mathfrak q})$ by the normalisation; for each noncompact $\beta$ with $\beta(x_{\mathfrak q})>0$ exactly one of $\pm\beta$ lies in $\Psi_{\mathrm n}$, and noncompact roots with $\beta(x_{\mathfrak q})=0$ contribute nothing, so $2\rho_{\mathrm n}(x_{\mathfrak q})\le\sum_{\beta\in\Delta(\mathfrak u\cap\mathfrak p_1)}\beta(x_{\mathfrak q})=2\rho(\mathfrak u\cap\mathfrak p_1)(x_{\mathfrak q})$, with equality if and only if $\Delta(\mathfrak u\cap\mathfrak p_1)\subseteq\Psi$. Therefore
\[
\lambda(x_{\mathfrak q})\ge\lambda_{\mathfrak q}(x_{\mathfrak q})+2\rho(\mathfrak u\cap\mathfrak p_1)(x_{\mathfrak q})+B_{\mathfrak q}(x_{\mathfrak q})+\rho(\mathfrak u\cap\mathfrak k_1)(x_{\mathfrak q})-\rho(\mathfrak u\cap\mathfrak p_1)(x_{\mathfrak q})=\lambda_{\mathfrak q}(x_{\mathfrak q})+\rho(\mathfrak u)(x_{\mathfrak q})+B_{\mathfrak q}(x_{\mathfrak q}).
\]
With \eqref{par:eq-D3a}, $B_{\mathfrak q}(x_{\mathfrak q})\le0$; since $\beta(x_{\mathfrak q})>0$ for $\beta\in\Delta(\mathfrak u\cap\mathfrak p_1)$, all $n_\beta=0$, so $B_{\mathfrak q}=0$, and equality holds throughout, so $\Delta(\mathfrak u\cap\mathfrak p_1)\subseteq\Psi$. Thus $\Lambda=\lambda_{\mathfrak q}+2\rho(\mathfrak u\cap\mathfrak p_1)$, and by (D3)(c) every $K_1$-type of $X$ has highest weight $\Lambda+\sum n_\beta\beta$ with $\beta\in\Delta(\mathfrak u\cap\mathfrak p_1)\subseteq\Psi_{\mathrm n}$, $n_\beta\ge0$. This is the first assertion.

If $X$ contained a $K$-type $\Lambda-A$ with $A$ a non-empty sum of elements of $\Psi$, then $-A=\sum n_\beta\beta$ with $\beta\in\Psi$, $n_\beta\ge0$; pairing with $\lambda$, which is positive on $\Psi$, gives $\langle\lambda,-A\rangle<0\le\langle\lambda,\sum n_\beta\beta\rangle$, a contradiction. By (D3)(a) (an irreducible unitary representation is quasi-simple, and infinitesimal equivalence means isomorphic Harish-Chandra modules), $X$ is the Harish-Chandra module of $\pi_\lambda$.
\end{proof}

\begin{lemma}[Rigidity]\label{par:rigid}
Let $\sigma$ be a unitary representation of $G$ on a separable Hilbert space and $y$ a smooth
vector with: $\sigma(t)y=t^{\Lambda^K(\lambda)}y$ for $t\in T$; $d\sigma(E_{ab})y=0$ for
$\varepsilon_a=\varepsilon_b$, $\lambda_a>\lambda_b$; $d\sigma(\Omega)y=\kappa_\lambda y$. Put
$e_\lambda=d_\lambda\overline{c_\lambda}$. Then $\langle y,\sigma(g)y\rangle=\|y\|^2c_\lambda(g)$ and
$\sigma(e_\lambda)y=y$, provided $e_\lambda\in L^1(G)$.
\end{lemma}
\begin{proof}
Let $H_y$ be the closed span of $\sigma(G)y$ and $H_y=\int^\oplus H_\zeta\,d\mu(\zeta)$ a
decomposition into irreducibles, $y=\int y_\zeta$. The three conditions on $y$ are equivalent to
countably many identities $\sigma(f)y=\sigma(f')y$ with $f,f'\in C_c^\infty(G)$ (for $\Omega$: $\sigma(f)
d\sigma(\Omega)y=\sigma(\Omega f)y$, with $\Omega$ acting on $f$ as a bi-invariant operator; for the
others, the projector onto $T$-weight $\Lambda^K$ and integration by parts against $f$). These
identities decompose, so for almost all $\zeta$ they hold for $y_\zeta$. Such a $y_\zeta$ lies in a
$K$-isotypic subspace, which is finite-dimensional, so $y_\zeta$ is smooth and satisfies the three
conditions in the derived action. If $y_\zeta\ne0$, then $H_\zeta$ contains the $K$-type
$\Lambda^K(\lambda)$ and $\Omega$ acts by $\kappa_\lambda$; with $\tau_K=\lambda-\rho_{\mathrm c}$,
$|\tau_K+\rho_{\mathrm c}|^2=|\lambda|^2=\kappa_\lambda+|\rho|^2$, so equality holds in (D1), and (D2)
gives the infinitesimal character $\lambda$. By Lemma~\ref{par:D3}, $H_\zeta\cong\pi_\lambda$, and $y_\zeta$, a
highest weight vector of the $K$-type of multiplicity one ((D3)(b)), is $a(\zeta)U_\zeta v$ for a unitary
$U_\zeta:\pi_\lambda\to H_\zeta$ and a scalar $a(\zeta)$. Hence $\langle y,\sigma(g)y\rangle=
\int|a|^2d\mu\cdot c_\lambda(g)=\|y\|^2c_\lambda(g)$. By Schur's orthogonality,
$\int e_\lambda(h)c_\lambda(g^{-1}h)\,dh=c_\lambda(g^{-1})$, so $\langle\sigma(g)y,\sigma(e_\lambda)y-
y\rangle=0$ for all $g$; since $\sigma(e_\lambda)y\in H_y$, $\sigma(e_\lambda)y=y$.
\end{proof}


\subsubsection{The Fock model at one place}\label{par:ss-fock}
Let $0\le j<m$. At $\tau$ let $R$ (resp.\ $C$) index the orthogonal basis of
$Y_j\otimes_{E,\tau}\C$ (resp.\ $Y_{j+1}$), with signs $\varepsilon_r$ (resp.\ $\varepsilon'_c$),
$n=|R|$, $n'=|C|=n+t$. The polynomial vectors of $\pzF_\tau$ are the polynomials in variables
$z_{rc}$; the variable $z_{rc}$ is the conjugate of the coordinate $x_{rc}$ if
$\varepsilon_r\varepsilon'_c=1$ and the coordinate itself otherwise (the positive complex structure
of \S\ref{pz:ss-types}). Put $\eta_{rc}=-\varepsilon_r\varepsilon'_c$ and define operators
$a^+_{rc}=z_{rc}$, $a^-_{rc}=\partial_{rc}$ if $\eta_{rc}=1$, and $a^+_{rc}=-\partial_{rc}$,
$a^-_{rc}=z_{rc}$ if $\eta_{rc}=-1$; then $[a^-_i,a^+_k]=\delta_{ik}$ and the $a^+$ (resp.\ $a^-$)
commute. For integers or half-integers $\gamma,\gamma'$ put, with $s_{ab}=1$ if
$\varepsilon_a=\varepsilon_b$ and $-1$ otherwise,
\begin{equation}\label{par:eq-osc}
E_{ab}=s_{ab}\sum_{c}a^+_{ac}a^-_{bc}+\gamma\delta_{ab},\qquad
E'_{cd}=-s'_{cd}\sum_ra^-_{rc}a^+_{rd}+\gamma'\delta_{cd}.
\end{equation}
These are commuting actions of $\mathfrak{gl}_n$ and $\mathfrak{gl}_{n'}$ (direct computation from the
commutation relations). For $\varepsilon_a\ne\varepsilon_b$ let $\Delta_{ab}=\sum_{c:\eta_{ac}=-1}
\partial_{ac}\partial_{bc}$ (the degree-lowering part of $E_{ab}$), and $\Delta'_{cd}$ likewise for
columns. On polynomials we use the Fischer inner product, for which $z_i$ is adjoint to $\partial_i$;
it is the inner product of $\pzF$.

\begin{lemma}[The derived action]\label{par:vac}
After conjugation by an element of $T\times T'$, the derived actions of $\omega_\infty(\cdot,1)$ and
$\omega_\infty(1,\cdot)$ at $\tau$ on polynomial vectors are \eqref{par:eq-osc}, with
\[
\gamma=\tau_0+\tfrac{n'}2,\quad\gamma'=\tau'_0+\tfrac n2,\qquad
\tau_0=\kappa_{j+1}(\tau)+\tfrac12[j\text{ odd}],\quad\tau'_0=-\kappa_{j+1}(\tau)-\tfrac12[j\text{ even}],
\]
provided the following holds.
\begin{enumerate}
\item[(V)] At $\tau$, the operator $\omega_{\auHom(Y_j,Y_{j+1})}(g,h)$ of \S\ref{ps:ss-weil} is the
genuine oscillator action of the metaplectic covers (Weyl-ordered quadratic operators) times
$\det(g)^{n'/2}\det(h)^{-n/2}$, where $\upsilon_\tau$ restricted to $\auU(1)$ is $z\mapsto z$.
\end{enumerate}
\end{lemma}
\begin{proof}
$K\times K'$ acts by linear substitutions, so the variable $z_{rc}$ has weight $\eta_{rc}e_r$ for
$T$ and $-\eta_{rc}e'_c$ for $T'$ (compare $F(h^tz\bar g)$ in the definite case of
\S\ref{pk:ss-definite}); hence the compact part is \eqref{par:eq-osc} up to the constants
$\gamma,\gamma'$. For all of $\mathfrak{gl}_n$ and $\mathfrak{gl}_{n'}$, compact and noncompact parts alike, we compute the operator $d\omega^{\mathrm{Mp}}$ of \S\ref{pco:ss-V} directly from its characterisation there: it is the unique element of the span $\mathcal V$ of the symmetrised quadratic operators in the $a^\pm_{rc}$ whose commutator with $d\varrho_\tau(w)$ is $d\varrho_\tau$ of the image of $w$. Let $x_{rc}$ be the coordinates of $w\in W_\tau$ (so that $we_r=\sum_cx_{rc}e'_c$ in the orthogonal bases, scaled as in \S\ref{pk:ss-conv}); the hermitian form is $\sum\varepsilon_r\varepsilon'_c|x_{rc}|^2$, the coordinates for the positive complex structure are $y_{rc}=x_{rc}$ if $\eta_{rc}=-1$ and $y_{rc}=\overline{x_{rc}}$ if $\eta_{rc}=1$, and by \eqref{pk:eq-fock} $d\varrho_\tau(w)=\sum_{r,c}(y_{rc}z_{rc}-\overline{y_{rc}}\partial_{rc})=\sum_{r,c}(\overline{x_{rc}}\,a^+_{rc}-\eta_{rc}x_{rc}\,a^-_{rc})$. Hence
\[
[a^-_{rc},d\varrho_\tau(w)]=\overline{x_{rc}},\qquad[a^+_{rc},d\varrho_\tau(w)]=\eta_{rc}\,x_{rc}.
\]
Let $X\in\mathfrak u(V)$; it acts on $W_\tau$ by $w\mapsto-wX$, that is, $x_{rc}\mapsto-\sum_{r'}X_{r'r}x_{r'c}$, and $\overline{X_{r'r}}=-s_{rr'}X_{rr'}$. So
\[
d\varrho_\tau(-wX)=\sum_{r,r',c}\bigl(s_{rr'}X_{rr'}\,\overline{x_{r'c}}\,a^+_{rc}+\eta_{rc}X_{r'r}\,x_{r'c}\,a^-_{rc}\bigr)
=\sum_{r,r'}s_{rr'}X_{rr'}\Bigl[\sum_ca^+_{rc}a^-_{r'c},\,d\varrho_\tau(w)\Bigr],
\]
where we inserted the two commutators, used $\eta_{rc}\eta_{r'c}=s_{rr'}$, and exchanged $r,r'$ in the second sum. The operator $\sum_{r,r'}s_{rr'}X_{rr'}\sum_c\frac12(a^+_{rc}a^-_{r'c}+a^-_{r'c}a^+_{rc})$ lies in $\mathcal V$ and differs from the one in the bracket by a scalar; therefore it is $d\omega^{\mathrm{Mp}}(X)$, and by complex linearity $d\omega^{\mathrm{Mp}}(E_{ab})=s_{ab}\sum_ca^+_{ac}a^-_{bc}+\frac{n'}2\delta_{ab}$ for all $a,b$. In the same way, $Y\in\mathfrak u(V')$ acts by $x_{rc}\mapsto\sum_{c'}Y_{cc'}x_{rc'}$ with $\overline{Y_{cc'}}=-s'_{cc'}Y_{c'c}$, and one finds $d\varrho_\tau(Yw)=-\sum_{c,c'}s'_{cc'}Y_{c'c}[\sum_ra^+_{rc}a^-_{rc'},d\varrho_\tau(w)]$, so $d\omega^{\mathrm{Mp}}(E'_{cd})=-s'_{cd}\sum_r\frac12(a^-_{rc}a^+_{rd}+a^+_{rd}a^-_{rc})$. Thus no conjugation is needed: the element of $T\times T'$ in the statement may be taken to be $1$. (For the compact parts this agrees with the weights found above.) These are the quadratic parts of \eqref{par:eq-osc}, with the constants $\tfrac{n'}2\delta_{ab}$ and $\tfrac n2\delta_{cd}$: for $\varepsilon_a=\varepsilon_b$ and $\varepsilon'_c=\varepsilon'_d$ the Weyl-ordered operators are $\sum_ca^+_{ac}a^-_{bc}+\tfrac{n'}2\delta_{ab}$ and
$-\sum_ra^-_{rc}a^+_{rd}+\tfrac n2\delta_{cd}$. The twists of $\omega_{\mu_{j+1}}$ contribute
$\det(g)^{\kappa_{j+1}}\det(g)^{-\lfloor n'/2\rfloor}\det(h)^{-\kappa_{j+1}}\det(h)^{\lfloor n/2\rfloor}$
($\mu_{j+1}$ has weight $-\kappa_{j+1}$, $\upsilon_k=\upsilon^{\lfloor k/2\rfloor}$), and (V) contributes
$\det(g)^{n'/2}\det(h)^{-n/2}$. So $\tau_0=\kappa_{j+1}+n'/2-\lfloor n'/2\rfloor$ and
$\tau'_0=-\kappa_{j+1}+\lfloor n/2\rfloor-n/2$; and $N_j=r+jt$ is odd exactly when $j$ is even.
\end{proof}
(V) is proved in Theorem~\ref{pco:V} (in the derived form (V$'$), which is what the proof above
uses), for all pairs, from (W1$_\infty$), (W2$_\infty$) and (M) of \S\ref{pco:ss-V}, and for the
definite pairs (only $j=m-1$, Remark~\ref{pco:rem-V}) also from the compatibility (S) with
orthogonal sums.

\emph{Parameters.} The row representation is $\overline{\pi_j}\cong\pi_\lambda$ with
$\lambda=-\Lambda^{(j)}_\tau$, and the column representation is $\pi_{\lambda'}$ with
$\lambda'=\Lambda^{(j+1)}_\tau$. Put $\xi_r=\lambda_r-\gamma+n'/2$ and $\xi'_c=\lambda'_c-\gamma'+n/2$.
We say that $(\lambda,\lambda')$ \emph{match} if there are a subset $B\subset C$ and a bijection
$r\mapsto c(r)$ of $R$ onto $C\smallsetminus B$ with
\begin{equation}\label{par:eq-match}
\xi'_{c(r)}=-\xi_r,\quad|\xi_r|>\tfrac t2,\quad\eta_{r\,c(r)}=\parsgn\xi_r;\qquad
\{\xi'_c:c\in B\}=\{\tfrac{t-1}2-b:0\le b<t\},\ \varepsilon'\text{ constant on }B.
\end{equation}

\begin{lemma}\label{par:match}
Under Lemma~\ref{par:vac}, $(-\Lambda^{(j)}_\tau,\Lambda^{(j+1)}_\tau)$ match.
\end{lemma}
\begin{proof}
Here $\xi_r=-\Lambda_r-\tau_0$, $\xi'_c=\Lambda'_c-\tau'_0$, and $\xi'_{c(r)}=-\xi_r$ means
$\Lambda'_{c(r)}=\Lambda_r+\tau_0+\tau'_0$, a shift by $-\tfrac12$ for $j$ even and $+\tfrac12$ for $j$
odd. By \eqref{pz:eq-param} the passage $j\to j+1$ shifts the old numbers by exactly this amount and
adds the string $-\kappa_{j+1}+\tfrac{t-1}2-b$ lowered by $\tfrac12[j\text{ even}]$, that is, centred
at $\tau'_0$: this is $B$, with signs $+$. The condition $\eta_{rc(r)}=\parsgn\xi_r$ reads
$\varepsilon'_{c(r)}=\varepsilon_r\parsgn(\Lambda'_{c(r)}-\tau'_0)$: the sign is kept above the centre
and changed below it, which is the rule of \eqref{pz:eq-param}; one checks directly from
\eqref{pz:eq-param} that the numbers $-\kappa-i\Delta$ with $i\ge j+2$ are exactly those below the
centre, and that the resulting signs are those of $\Lambda^{(j+1)}$. Finally the old numbers are at
distance at least $\tfrac12(\Delta-t)>t/2$ from the centre.
\end{proof}

\begin{remark}\label{par:rem-match}
With the other reading of $c_j$, the row parameter would be $\Lambda^{(j)}$ and the column parameter
$-\Lambda^{(j+1)}$; the first condition of \eqref{par:eq-match} would put the new string at $-\tau'_0$,
near $+\kappa_{j+1}$, contradicting \eqref{pz:eq-param}. This is why \eqref{par:eq-coef} is the
reading that makes \S\ref{pz:ss-types} consistent.
\end{remark}

\subsubsection{The joint harmonic}\label{par:ss-harm}
In this subsection $(\lambda,\lambda')$ is any matching pair of parameters with distinct entries.
Write $\Lambda=\Lambda^K(\lambda)$, $\Lambda'=\Lambda^K(\lambda')$. The vacuum $1$ has $E_{aa}$-weight
$v_a=\gamma-\#\{c:\eta_{ac}=-1\}$ and $E'_{cc}$-weight $v'_c=\gamma'-\#\{r:\eta_{rc}=1\}$. Put
$w_r=\Lambda_r-v_r$ and $w'_c=\Lambda'_c-v'_c$.

\begin{definition}\label{par:def-harm}
A \emph{joint harmonic} is a polynomial $\Phi\ne0$ with $E_{aa}\Phi=\Lambda_a\Phi$,
$E'_{cc}\Phi=\Lambda'_c\Phi$, $E_{ab}\Phi=0$ ($\varepsilon_a=\varepsilon_b$, $\lambda_a>\lambda_b$),
$E'_{cd}\Phi=0$ ($\varepsilon'_c=\varepsilon'_d$, $\lambda'_c>\lambda'_d$), $\Delta_{ab}\Phi=0$
($\varepsilon_a\ne\varepsilon_b$) and $\Delta'_{cd}\Phi=0$ ($\varepsilon'_c\ne\varepsilon'_d$). A
\emph{pre-joint harmonic} satisfies the same conditions except $E'_{cc}\Phi=\Lambda'_c\Phi$.
\end{definition}
In the notation of \S\ref{pz:ss-types}, $\Phi_j=\bigotimes_\tau\Phi_\tau$ with $\Phi_\tau$ a joint
harmonic at $\tau$ for $(-\Lambda^{(j)}_\tau,\Lambda^{(j+1)}_\tau)$. Being highest for the order of
$-\Lambda^{(j)}$ is being lowest for the order of $\Lambda^{(j)}$: $\Phi_j$ has the $K_j$-weight
$-\Lambda^K(\Lambda^{(j)})$ of the vector $\overline{v_j}$ of $\overline{\pi_j}$.

\begin{lemma}\label{par:relw}
\textup{(i)} $w_r\parsgn\xi_r\ge0$ for all $r$. \textup{(ii)} $w_r=-w'_{c(r)}$, and $w'_c=0$ for $c\in B$.
\textup{(iii)} If $\varepsilon_a=\varepsilon_b$ and $\lambda_a>\lambda_b$ then $\Lambda_a\ge\Lambda_b$;
likewise for $\Lambda'$.
\end{lemma}
\begin{proof}
From $2(\rho_{\mathrm n}-\rho_{\mathrm c})_r=\sum_k\varepsilon_r\varepsilon_k\parsgn(\lambda_k-\lambda_r)=:A_r$
and $2\#\{c:\eta_{rc}=-1\}=n'-\sum_c\eta_{rc}$ we get $2w_r=2\xi_r+A_r-\beta_r$, $\beta_r=\sum_c\eta_{rc}$.
By \eqref{par:eq-match}, $\varepsilon'_{c(k)}=-\varepsilon_k\parsgn\xi_k$, so
$\beta_r=\varepsilon_r\sum_k\varepsilon_k\parsgn\xi_k-\varepsilon_r\sum_{c\in B}\varepsilon'_c$. Let $\xi_r>0$
(the case $\xi_r<0$ is symmetric). Comparing $\parsgn(\xi_k-\xi_r)$ with $\parsgn\xi_k$ term by term,
\[
2w_r=2\xi_r-1-2\!\!\sum_{0<\xi_k<\xi_r}\!\!\varepsilon_r\varepsilon_k+\varepsilon_r\sum_{c\in B}\varepsilon'_c .
\]
All $\xi$ and $\xi'$ lie in $\tfrac{t-1}2+\Z$, so $|\xi_k|>t/2$ gives $|\xi_k|\ge\tfrac{t+1}2$, and at most
$\xi_r-\tfrac{t+1}2$ indices $k$ have $0<\xi_k<\xi_r$. Hence $2w_r\ge2\xi_r-1-(2\xi_r-t-1)-t=0$: this is
(i). For (ii), the same computation for the column $c=c(r)$ (with $\xi'_c=-\xi_r$, signs
$-\eta_{rc}$ for the column action, and $|\xi'_d|<t/2<\xi_r$ for $d\in B$) gives
$2w'_c=-2\xi_r+1+2\sum_{0<\xi_k<\xi_r}\varepsilon_r\varepsilon_k-\varepsilon_r\sum_{B}\varepsilon'_d=-2w_r$. For
$c\in B$ the matched columns contribute $\varepsilon'_c\sum_k\varepsilon_k$, which cancels the column
analogue of $\beta$, and what remains is $2\xi'_c+\sum_{d\in B}\parsgn(\xi'_d-\xi'_c)=0$ because
$\varepsilon'$ is constant on $B$ and $B$ is a string centred at $0$. (iii): $\Lambda_a-\Lambda_b=
\lambda_a-\lambda_b-1-\sum\varepsilon_a\varepsilon_k$ over the $k$ with $\lambda_b<\lambda_k<\lambda_a$, of
which there are at most $\lambda_a-\lambda_b-1$.
\end{proof}

\begin{proposition}[Existence]\label{par:exist}
For $r\in R$ and an integer $0\le\ell<|w_r|$, call $(r,c(r))$ a pair of \emph{level} $\ell$. For each
$\ell$ and each pair of signs $(\epsilon,\epsilon')$ let $L_{\ell,\epsilon,\epsilon'}$ be the set of
rows $r$ having a pair of level $\ell$ with $\varepsilon_r=\epsilon$, $\varepsilon'_{c(r)}=\epsilon'$, and
$M_{\ell,\epsilon,\epsilon'}=\det(z_{r\,c(r')})_{r,r'\in L_{\ell,\epsilon,\epsilon'}}$. Then
$\Phi=\prod_{\ell,\epsilon,\epsilon'}M_{\ell,\epsilon,\epsilon'}$ is a joint harmonic. Moreover no monomial
of $\Phi$ contains two variables of one column from rows of opposite signs, or of one row from columns
of opposite signs.
\end{proposition}
\begin{proof}
$\Phi\ne0$: each minor is a generic determinant in distinct variables. A variable $z_{rc}$ occurs only
if $c=c(r')$ with $r,r'$ in a common layer, so $\varepsilon_r=\varepsilon_{r'}$ and the column sign is
$\varepsilon'_{c(r)}$; this gives the last assertion, hence $\Delta\Phi=0=\Delta'\Phi$. All variables of
row $r$ in $\Phi$ have sign $\eta_{rc(r)}=\parsgn\xi_r$ (\eqref{par:eq-match}).

For $\varepsilon_a=\varepsilon_b$, \eqref{par:eq-osc} gives
$E_{ab}=D^+_{a\leftarrow b}-D^-_{b\leftarrow a}+\delta_{ab}v_a$ with
$D^s_{a\leftarrow b}=\sum_{c:\eta_{bc}=s}z_{ac}\partial_{bc}$, a derivation. \emph{Weight.} $D^s_{a\leftarrow a}$ is
the degree in the sign-$s$ variables of row $a$, which is the number $|w_a|$ of layers containing
$a$ if $s=\parsgn\xi_a$ and $0$ otherwise; by Lemma~\ref{par:relw}(i) $E_{aa}\Phi=(w_a+v_a)\Phi=
\Lambda_a\Phi$. \emph{Raising operators.} Let $\lambda_a>\lambda_b$. If $D^+_{a\leftarrow b}$ acts
non-trivially on a minor, $b$ lies in that layer with $\parsgn\xi_b=+1$; then $\xi_a>\xi_b>0$, $a$ is
matched with a column of the same sign as $c(b)$ (by $\eta_{ac(a)}=\parsgn\xi_a$), and
$w_a\ge w_b$ by Lemma~\ref{par:relw}(iii) ($v_a=v_b$); so $a$ is in the same layer, and the derivation
produces a determinant with two equal rows. The same argument applies to $D^-_{b\leftarrow a}$ with
the roles of $a$, $b$ and of the signs exchanged. So $E_{ab}\Phi=0$. \emph{Columns.} The data are
symmetric under exchanging rows and columns (which replaces $\eta$ by $-\eta$); by
Lemma~\ref{par:relw}(ii) the levels of $(r,c(r))$ seen from the column are the same, columns in $B$
carry no variable and have $w'_c=0$, so the same argument gives the column conditions.
\end{proof}

\begin{lemma}[Polarisations]\label{par:polar}
Let $P$ be killed by the row raising operators and by the $\Delta'_{cd}$. Let $\varepsilon_a=\varepsilon_b$.
\textup{(a)} $\langle D^{-s}_{b\leftarrow a}P,D^{s}_{a\leftarrow b}P\rangle=0$ for $s=\pm1$.
\textup{(b)} If $s\lambda_a>s\lambda_b$ then $D^s_{a\leftarrow b}P=0$.
\textup{(c)} If $E_{aa}P=\Lambda_aP$ and $k=s\,w_a\ge0$, then $D^s_{a\leftarrow a}P=kP$ and
$\partial_{ac}P=0$ whenever $\eta_{ac}=-s$.
\end{lemma}
\begin{proof}
Put $u_{acd}=\partial_{ac}\partial_{ad}P$ and $N_s(a,b)=\sum_{\eta_{ad}=-s,\ \eta_{ac}=s}\langle u_{acd},u_{bcd}
\rangle$. Adjointness gives $\langle D^{-s}_{b\leftarrow a}P,D^s_{a\leftarrow b}P\rangle=N_s(a,b)$ (no
commutator term, as $c\ne d$), and $N_{-s}=N_s$. If $\lambda_a>\lambda_b$, $E_{ab}P=0$ gives
$D^+_{a\leftarrow b}P=D^-_{b\leftarrow a}P$, so $N_1(a,b)\ge0$; the cases $a=b$ and $\lambda_a<\lambda_b$
are similar, so $N_1(a,b)\ge0$ for all $b$ with $\varepsilon_b=\varepsilon_a$. For fixed $c,d$ with
$\eta_{ac}=1=-\eta_{ad}$, $\sum_{b:\varepsilon_b=\varepsilon_a}\partial_{bc}\partial_{bd}P=\pm\Delta'_{cd}P=0$, so
$\sum_bN_1(a,b)=0$, hence all $N_s(a,b)=0$: (a). (b) follows from (a) and the equality of the two
polarisations. (c): $E_{aa}P=\Lambda_aP$ gives $D^s_{a\leftarrow a}P=D^{-s}_{a\leftarrow a}P+kP$; by (a)
with $b=a$, $\|D^{-s}_{a\leftarrow a}P\|^2+k\sum_{\eta_{ac}=-s}\|\partial_{ac}P\|^2=0$.
\end{proof}

\begin{proposition}[Uniqueness]\label{par:unique}
Two pre-joint harmonics are proportional. In particular the joint harmonic is unique up to a scalar.
\end{proposition}
\begin{proof}
Let $x_0$ be the matrix with $z_{rc}=1$ if $c=c(r)$ and $0$ otherwise. We show that a pre-joint
harmonic $F$ has $F(x_0)\ne0$; then for pre-joint harmonics $P,Q$, $P-\tfrac{P(x_0)}{Q(x_0)}Q$
vanishes at $x_0$ and is, if non-zero, a pre-joint harmonic, which is absurd. Suppose $F(x_0)=0$.
Expanding $F$ at $x_0$ and taking a monomial of least positive degree, we find a variable $i$ with
$\partial_iF\ne0$ such that $F$ is an eigenvector of no vector field $\sum_kz_{a_k}\partial_{b_k}$ whose
only term not vanishing at $x_0$ is a multiple of $\partial_i$ (compare the coefficients of that
monomial). Let $i=(r,c)$, $s=\eta_{rc}$. Since $|\xi_r|>t/2$, either $-s\xi_r>t/2$, and then
Lemma~\ref{par:polar}(c) with $-s$ (and Lemma~\ref{par:relw}(i)) gives $\partial_{rc}F=0$; or
$s\xi_r>t/2$, so $\eta_{rc(r)}=s$ and $\varepsilon'_{c(r)}=\varepsilon'_c$. Three cases. If $c=c(r)$, $F$ is an
eigenvector of $D^s_{r\leftarrow r}$ by Lemma~\ref{par:polar}(c), whose only term non-zero at $x_0$
is $z_{rc}\partial_{rc}$. If $s\lambda'_{c(r)}<s\lambda'_c$, the column polarisation
$\sum_{a:\eta_{ac}=s}z_{a\,c(r)}\partial_{ac}$ kills $F$ (Lemma~\ref{par:polar}(b) for columns), and its only
term non-zero at $x_0$ is $a=r$. If $s\lambda'_{c(r)}>s\lambda'_c$, then $|\xi'_c|>t/2$, so $c=c(a)$ with
$s\lambda_a>s\lambda_r$, $\varepsilon_a=\varepsilon_r$, and $D^s_{a\leftarrow r}F=0$ by Lemma~\ref{par:polar}(b),
with only term $d=c$ non-zero at $x_0$. Each case contradicts the choice of $i$.
\end{proof}

\subsubsection{The Casimir comparison and the proof of (A3)}\label{par:ss-A3}

\begin{lemma}\label{par:casimir}
Put $\Omega=\sum E_{ab}E_{ba}$, $\Omega'=\sum E'_{cd}E'_{dc}$, $\mathcal N=\sum_{a,c}a^+_{ac}a^-_{ac}=
\sum_aE_{aa}-n\gamma$ and $s_0=\gamma+\gamma'-\tfrac{n+n'}2$. On polynomials,
$\Omega'=\Omega-2s_0\mathcal N+n^2n'-2\gamma'nn'+n'\gamma'^2-n\gamma^2$. If $(\lambda,\lambda')$ match,
then on every polynomial $\Phi$ with $\sum_aE_{aa}\Phi=(\sum_a\lambda_a)\Phi$ the scalars satisfy
$\kappa_{\lambda'}=\kappa_\lambda-2s_0(\sum_a\lambda_a-n\gamma)+n^2n'-2\gamma'nn'+n'\gamma'^2-n\gamma^2$.
\end{lemma}
\begin{proof}
The signs $s_{ab}$ drop out since $s_{ab}s_{ba}=1$. Write $A=a^+$, $B=a^-$ and
$Q=\sum A_{ac}A_{bd}B_{bc}B_{ad}$. Normal ordering gives $\sum_{a,b}(\sum_cA_{ac}B_{bc})(\sum_dA_{bd}B_{ad})
=Q+n\mathcal N$, and with $F_{cd}=\sum_rB_{rc}A_{rd}$, $\sum_{c,d}F_{cd}F_{dc}=Q+(2n+n')\mathcal N+n^2n'$,
while $\sum_cF_{cc}=\mathcal N+nn'$. For the second count, moving the $A$'s to the left with $[B_i,A_k]=\delta_{ik}$ gives
\[
B_{rc}A_{rd}B_{r'd}A_{r'c}=A_{r'c}A_{rd}B_{rc}B_{r'd}+\delta_{rr'}A_{rd}B_{rd}+\delta_{cd}A_{r'c}B_{r'd}+\delta_{cd}A_{rd}B_{rc}+\delta_{cd},
\]
and summing over $r,r',c,d$ the four correction terms give $n'\mathcal N$, $n\mathcal N$, $n\mathcal N$ and $n^2n'$. Adding the shifts gives the first identity. For the second,
$\lambda'=\{-\lambda_r+s_0\}\cup\{\xi_0+\tfrac{t-1}2-b\}$ with $\xi_0=\gamma'-\tfrac n2$, by
\eqref{par:eq-match}; so $|\lambda'|^2=|\lambda|^2-2s_0\sum\lambda+ns_0^2+t\xi_0^2+\tfrac{t(t^2-1)}{12}$, and
$|\rho'|^2-|\rho|^2=\tfrac{n'(n'^2-1)-n(n^2-1)}{12}$. Substituting $\gamma=s_0-\xi_0+n'/2$,
$\gamma'=\xi_0+n/2$, both sides reduce to $ns_0^2+t\xi_0^2-nn't/4$.
\end{proof}
(Here $\sum\Lambda_a=\sum\lambda_a$, since $\rho_{\mathrm n}-\rho_{\mathrm c}$ is a combination of roots.)

\begin{theorem}[(A3)]\label{par:A3}
Assume Lemma~\ref{par:vac} (that is, (V)), the conclusion of Theorem~\ref{par:A1}, and
\textup{(D1)--(D3)}. Let $0\le j<m$. A joint harmonic $\Phi_j=\bigotimes_\tau\Phi_\tau$ exists and is
unique up to a scalar, and
\[
\int_{G_{j,\infty}}e_j(x)\,\omega_\infty(x^{-1},1)\Phi_j\,dx=\int_{G_{j+1,\infty}}e_{j+1}(x)\,\omega_\infty(1,x)
\Phi_j\,dx .
\]
\end{theorem}
\begin{proof}
Existence and uniqueness at each $\tau$ are Lemma~\ref{par:match} and
Propositions~\ref{par:exist}, \ref{par:unique}. For the product: the conditions of
Definition~\ref{par:def-harm} are conditions on each tensor factor, and on each finite-dimensional
space of polynomials of bounded degree the common kernel of operators acting on separate factors is
the tensor product of the kernels. Both sides are tensor products of the local integrals (Fubini;
$e_j\in L^1$, $\omega_\infty$ unitary), so it suffices to prove the local identity
$P\Phi=P'\Phi$, $P=\int e_{j,\tau}(x)\omega(x^{-1},1)dx$, $P'=\int e_{j+1,\tau}(x)\omega(1,x)dx$, which
are bounded commuting operators on $\pzF_\tau$.

Since $e_{j,\tau}(x^{-1})=\overline{e_{j,\tau}(x)}$ and $\overline{c_{j,\tau}}=c_{-\Lambda^{(j)}}$ (the
conjugate of $\pi_j$ is $\pi_{-\Lambda^{(j)}}$, with $\overline{v_j}$ a highest weight vector for the
order of $-\Lambda^{(j)}$ of its lowest $K$-type), $P=\sigma(e_\lambda)$ for $\sigma=\omega(\cdot,1)$ and
$\lambda=-\Lambda^{(j)}_\tau$; likewise $P'=\sigma'(e_{\lambda'})$, $\lambda'=\Lambda^{(j+1)}_\tau$.

Let $u=P\Phi$. It is smooth for $\sigma'$, and $d\sigma'(Y)u=Pd\sigma'(Y)\Phi$, because $P$ is
bounded and commutes with $\sigma'$ and $\Phi$ is a smooth vector. So $u$ has $T'$-weight $\Lambda'$
and is killed by the compact raising operators of the columns. By Lemma~\ref{par:casimir},
$d\sigma'(\Omega')u=P\Omega'\Phi=P\Omega\Phi+c\,P\Phi$ with $c=-2s_0(\sum\lambda_a-n\gamma)+\dots$, and
$P\,d\sigma(\Omega)\Phi=\kappa_\lambda P\Phi$: integrate by parts twice, using that $e_\lambda$ is an
eigenfunction of the bi-invariant operator $\Omega$ with eigenvalue $\kappa_\lambda$ and that its
derivatives are integrable (Theorem~\ref{par:A1}). So $d\sigma'(\Omega')u=\kappa_{\lambda'}u$, and
Lemma~\ref{par:rigid} for $\sigma'$ gives $P'u=u$. The symmetric argument for $y=P'\Phi$ (with
Lemma~\ref{par:casimir} read backwards) gives $Py=y$. Hence $P\Phi=P'P\Phi=PP'\Phi=P'\Phi$.
\end{proof}

\begin{remark}[The type on $Y_m$]\label{par:Ym}
In Definition~\ref{pz:def-type}, $G_{m,\infty}$ acts on $W_{\lambda^\circ}$ through $g\mapsto\rho(\bar g)$
(\S\ref{ps:ss-weil}). For $y_l(g)=l(\rho(\bar g_\infty)^{-1}\varphi(g_f))$ one has $y_l(gx)=y_{\sigma(x)l}(g)$
with $\sigma(x)l=l\circ\rho(\bar x)^{-1}$. For $t$ in the torus, $\bar t^{-1}=t$, so the dual basis vector
$e_T^*$ has $\sigma$-weight equal to the weight of $e_T$. Thus $\sigma\cong W_{\lambda^\circ}$, the
highest weight vector being $e^*_{T}$ for the highest $T$, and with \eqref{par:eq-coef}
$R(e_m)y_l=y_l$ exactly when $l$ is proportional to that vector. So the type on $Y_m$ is
$W_{\lambda^\circ}$ and not its dual, as Definition~\ref{pz:def-type} states. This agrees with the
definite case \eqref{pk:eq-schuract}, where $h$ acts on covectors $A$ by $A\mapsto A\rho(\bar h)^{-1}$.
\end{remark}


%% ===== proofs/A-coefficient.tex =====
% proofs/A-coefficient.tex
% Proof of the integral formula pds:eq-FJ for the lowest-K-type coefficient, and of the
% normalisation (V) of Lemma par:vac. Labels of this file start with pco:.
% Only amsmath/amssymb/amsthm are used.
\providecommand{\auU}{\mathrm U}
\providecommand{\auHom}{\operatorname{Hom}}
\providecommand{\pcom}{\mathsf m}
\providecommand{\pcoD}{\mathsf D}
\providecommand{\pcoH}{\mathsf H}
\providecommand{\pcoB}{\mathsf B}
\providecommand{\pcotr}{\operatorname{tr}}

\subsection{The coefficient of the lowest $K$-type and the vacuum character}\label{pco:sec}

This subsection proves the two statements needed from \S\ref{par:sec}: the integral formula
\eqref{pds:eq-FJ} for the coefficient of the lowest $K$-type (Theorem~\ref{pco:FJ}), and the
normalisation (V) of Lemma~\ref{par:vac} (Theorem~\ref{pco:V}). \emph{Dependencies.}
Theorem~\ref{pco:FJ} is proved in full from the published results (D1)--(D3) of
\S\ref{par:ss-rigid}, the lemmas \ref{pds:blockdet}, \ref{pds:fischer} and the argument of \ref{pm:kernel}(c),
and standard facts recalled where they are used (highest weight theory for compact groups,
uniqueness of invariant measures on a quotient by a compact subgroup, admissibility). Theorem~\ref{pco:V}
is proved modulo four exactly stated properties of the archimedean Weil representation:
the equivariance (W1$_\infty$), the archimedean form (W2$_\infty$) of (W2), the formula (M) for the metaplectic representation
in a mixed model, and, for definite pairs only, the compatibility (S) with orthogonal sums;
see \S\ref{pco:ss-V}.

\subsubsection{Notation and statement}\label{pco:ss-stat}
We use the conventions of \S\ref{par:sec}: $G=\auU(p,q)=\{g:g^*Jg=J\}$ with
$J=\operatorname{diag}(\varepsilon_1,\dots,\varepsilon_n)$, $K=G\cap\auU(n)$, $T\subset K$ the
diagonal torus, $\lambda\in(\frac{n-1}2+\Z)^n$ with distinct entries, the positive roots
$e_a-e_b$ with $\lambda_a>\lambda_b$, $\rho,\rho_{\mathrm c},\rho_{\mathrm n}$,
$\Lambda=\Lambda^K(\lambda)=\lambda+\rho_{\mathrm n}-\rho_{\mathrm c}$,
\[
\mu=\lambda+\rho=\Lambda+2\rho_{\mathrm c},\qquad\kappa_\lambda=|\lambda|^2-|\rho|^2,
\]
and $h$ is the gap of $\lambda$. Since $\rho_a=\frac12(\#\{b:\lambda_b<\lambda_a\}-\#\{b:\lambda_b>
\lambda_a\})\in\frac{n-1}2+\Z$, $\mu\in\Z^n$, and $\Lambda\in\Z^n$. For $t\in T$ (or any
diagonal matrix $t$) and $w\in\Z^n$ we write $t^w=\prod_at_{aa}^{w_a}$. The coefficient is
$c_\lambda(g)=\langle v,\pi_\lambda(g)v\rangle$, $v$ a unit highest weight vector of the lowest
$K$-type $\Lambda$ of $\pi_\lambda$ (\S\ref{par:ss-decay}); it does not depend on the choice of
$v$, which is unique up to a scalar of modulus one by (D3)(b).

For $x\in M_n(\C)$ and $S\subseteq\{1,\dots,n\}$, $\Delta_S(x)$ is the principal minor on $S$
($\Delta_\emptyset=1$), and for $w\in\Z^n$ and invertible $x$, as in \eqref{pds:eq-m},
\[
\pcom_w(x)=\prod_{t\in\Z}\Delta_{S_t(w)}(x)\,(\det x)^{-[t\le0]},\qquad S_t(w)=\{a:w_a\ge t\}.
\]
Every $g\in G$ is uniquely $g=kz$ with $k\in K$ and $z\in G$ positive definite
(\S\ref{pds:ss-FJ}). Let $du$ be the Haar measure of mass $1$ on $K$, and
$d_\Lambda=(\int_K|\pcom_\Lambda|^2du)^{-1}$ (note $\pcom_\Lambda(1)=1$). Define
\begin{equation}\label{pco:eq-F}
F(kz)=d_\Lambda\int_K\frac{\pcom_\Lambda(ku)\,\overline{\pcom_\Lambda(u)}}{\pcom_{-\mu}(u^{-1}zu)}\,du
\qquad(k\in K,\ z\in G\text{ positive definite}),
\end{equation}
the right side of \eqref{pds:eq-FJ}; the denominator is positive by Lemma~\ref{pds:fischer}.

\begin{theorem}\label{pco:FJ}
Assume $h\ge n^2+1$. Then $F=c_\lambda$ on $G$; that is, \eqref{pds:eq-FJ} holds.
\end{theorem}

The hypothesis holds for every $\Lambda^{(j)}_\tau$ (the gap is at least $\frac12x^{40}$ and
$n=N_j\le40x^6$, \S\ref{par:sec}), and for the parameters of \S\ref{pds:ss-setup}
($h\ge2r^8$, $N=r$). The proof occupies \S\S\ref{pco:ss-minors}--\ref{pco:ss-ident}. It follows
the outline of Remark~\ref{par:FJ}: $F$ is real analytic, right $K$-highest of weight $\Lambda$,
an eigenfunction of $\Omega$ with eigenvalue $\kappa_\lambda$ (\S\ref{pco:ss-casimir}),
integrable and square integrable (by the proof of Theorem~\ref{pds:decay}); a weak form of
Lemma~\ref{par:rigid} and the left $K$-behaviour of $F$ then give $F=c_\lambda$
(\S\ref{pco:ss-ident}). The analytic difficulty is that \eqref{pco:eq-F} uses the polar
decomposition, which is not holomorphic; it is removed by showing that the polar factor $k$ may be
replaced by any block diagonal $a$ for which $a^{-1}g$ is accretive (Proposition~\ref{pco:indep}).

\subsubsection{Products of minors}\label{pco:ss-minors}
Let $\pcoD\subset M_n(\C)$ be the algebra of block diagonal matrices (those commuting with $J$),
$\pcoH=\pcoD\cap\mathrm{GL}_n(\C)\cong\mathrm{GL}_p(\C)\times\mathrm{GL}_q(\C)$, so that $K=\pcoH\cap
\auU(n)$. Let $\pcoB\subset\pcoH$ be the subgroup of $b\in\pcoH$ with $b_{ab}=0$ whenever
$\lambda_a<\lambda_b$ (upper triangular in the order of decreasing $\lambda$), and $\pcoB^-=
\{b^*:b\in\pcoB\}$. A matrix $Z$ is \emph{accretive} if $Z+Z^*$ is positive definite.

\begin{lemma}\label{pco:minors}
Let $w\in\Z^n$ and let $x$ be invertible.
\begin{enumerate}
\item[(a)] $\pcom_{-w}(x)=\pcom_w(x^{-1})$, and $\overline{\pcom_w(x)}=\pcom_w(x^*)$.
\item[(b)] Let $b$ be invertible, with $\Delta_S(b)=\prod_{a\in S}b_{aa}$ for every $S$. If
$b_{ab}=0$ whenever $w_a<w_b$, then $\pcom_w(xb)=b^w\pcom_w(x)$. If $b_{ab}=0$ whenever
$w_b<w_a$, then $\pcom_w(bx)=b^w\pcom_w(x)$. The condition on the minors holds for diagonal $b$,
for $b=1+sE_{ab}$ \textup($a\ne b$, $s\in\C$\textup), and for $b\in\pcoB\cup\pcoB^-$.
\item[(c)] If $\varepsilon_a=\varepsilon_b$ and $\lambda_a>\lambda_b$, then $\Lambda_a\ge\Lambda_b$.
For all $a,b$: $\mu_a>\mu_b$ if and only if $\lambda_a>\lambda_b$.
\item[(d)] If $Z$ is accretive and $y$ invertible, then $y^*Zy$ and $Z^{-1}$ are accretive, and
every principal minor of $Z$ is nonzero; so $\pcom_w(Z)$ is defined and nonzero.
\end{enumerate}
\end{lemma}
\begin{proof}
(a) Jacobi's identity $\Delta_{S^c}(x)=\det x\cdot\Delta_S(x^{-1})$ and
$\{a:-w_a\ge t\}=S_{1-t}(w)^c$ give, with $s=1-t$, the factor $\Delta_{S_s(w)}(x^{-1})
(\det x)^{1-[s\ge1]}=\Delta_{S_s(w)}(x^{-1})(\det x^{-1})^{-[s\le0]}$ of $\pcom_w(x^{-1})$. The
second identity holds factor by factor, since $\overline{\Delta_S(x)}=\Delta_S(\bar x^{\,t})$.

(b) If $b_{cd}=0$ for $c\notin S$, $d\in S$, then $(xb)_{SS}=x_{SS}b_{SS}$, so $\Delta_S(xb)=
\Delta_S(x)\Delta_S(b)$. For $S=S_t(w)$, $c\notin S$, $d\in S$ means $w_c<t\le w_d$, so the
hypothesis applies. With $\Delta_S(b)=\prod_{a\in S}b_{aa}$ (also for $S$ the full set, so
$\det b=\prod b_{aa}$), the factors of $\pcom_w(xb)/\pcom_w(x)$ multiply to $\prod_ab_{aa}^{e_a}$
with $e_a=\#\{t\in(-B,B]:t\le w_a\}-\#\{t\in(-B,B]:t\le0\}=w_a$ for $B$ large. The left version
is the same with rows and columns exchanged. Diagonal and elementary matrices have the stated
minors; a matrix triangular for a total order has every principal submatrix triangular.

(c) Let $m_{\mathrm c}$ (resp.\ $m_{\mathrm n}$) be the number of $c$ with $\varepsilon_c=
\varepsilon_a$ (resp.\ $\ne$) and $\lambda_b<\lambda_c<\lambda_a$. Counting from the definitions,
$\rho_{\mathrm c,a}-\rho_{\mathrm c,b}=1+m_{\mathrm c}$ and $\rho_{\mathrm n,a}-
\rho_{\mathrm n,b}=m_{\mathrm n}$. The entries of $\lambda$ are distinct and congruent modulo $\Z$,
so $m_{\mathrm c}+m_{\mathrm n}\le\lambda_a-\lambda_b-1$, and $\Lambda_a-\Lambda_b=
(\lambda_a-\lambda_b-1-m_{\mathrm c}-m_{\mathrm n})+2m_{\mathrm n}\ge0$. For the second claim,
$\rho_a-\rho_b$ has the sign of $\lambda_a-\lambda_b$.

(d) $v^*(y^*Zy+y^*Z^*y)v>0$ for $v\ne0$; $Z^{-1}+Z^{-*}=Z^{-1}(Z^*+Z)Z^{-*}$. A principal submatrix
$Z_{SS}$ is accretive, and $Z_{SS}v=0$ forces $v^*(Z_{SS}+Z_{SS}^*)v=0$, so $v=0$.
\end{proof}

\begin{corollary}\label{pco:equiv}
Let $Y$ be invertible.
\begin{enumerate}
\item[(a)] For $b\in\pcoB$: $\pcom_\Lambda(Yb)=b^\Lambda\pcom_\Lambda(Y)$, $\pcom_\mu(Yb)=
b^\mu\pcom_\mu(Y)$, $\pcom_\Lambda(b^*Y)=\overline{b^\Lambda}\pcom_\Lambda(Y)$ and $\pcom_\mu(b^*Y)
=\overline{b^\mu}\pcom_\mu(Y)$.
\item[(b)] If $\varepsilon_a=\varepsilon_b$ and $\lambda_a>\lambda_b$, then $\pcom_\Lambda(Y(1+sE_{ab}))
=\pcom_\Lambda(Y)=\pcom_\Lambda((1+sE_{ba})Y)$ for $s\in\C$. For $t$ diagonal,
$\pcom_\Lambda(Yt)=t^\Lambda\pcom_\Lambda(Y)=\pcom_\Lambda(tY)$.
\item[(c)] Put $\psi(Y)=\pcom_{-\mu}(Y)^{-1}$ where $\pcom_{-\mu}(Y)\ne0$. If $\lambda_a<\lambda_b$
\textup(any signs\textup), $\psi(Y(1+sE_{ab}))=\psi(Y)$; for $t$ diagonal, $\psi(Yt)=t^\mu\psi(Y)$.
\end{enumerate}
\end{corollary}
\begin{proof}
Apply Lemma~\ref{pco:minors}(b). (a) An off-diagonal entry $b_{cd}\ne0$ of $b\in\pcoB$ has
$\varepsilon_c=\varepsilon_d$ and $\lambda_c>\lambda_d$, hence $\Lambda_c\ge\Lambda_d$ and $\mu_c>
\mu_d$ by (c); so $b_{cd}=0$ whenever $\Lambda_c<\Lambda_d$ or $\mu_c<\mu_d$. The matrix $b^*$
satisfies the transposed condition, and $(b^*)^w=\overline{b^w}$. (b) The entry $(a,b)$ of
$1+sE_{ab}$ is allowed since $\Lambda_a\ge\Lambda_b$, and likewise the entry $(b,a)$ of $1+sE_{ba}$
on the left. (c) For $w=-\mu$ the entry $(a,b)$ is allowed on the right unless $-\mu_a<-\mu_b$,
i.e.\ unless $\lambda_a>\lambda_b$; and $t^{-(-\mu)}=t^\mu$.
\end{proof}

\subsubsection{Densities on the flag variety of $\pcoH$}\label{pco:ss-density}
Let $T_{\pcoH}=\pcoB\cap K$; a unitary triangular matrix is diagonal, so $T_{\pcoH}=T$. Put
\[
\delta(b)=\bigl|b^{2\rho_{\mathrm c}}\bigr|^2=\prod_{\varepsilon_a=\varepsilon_b,\ \lambda_a>\lambda_b}
|b_{aa}/b_{bb}|^2\qquad(b\in\pcoB).
\]

\begin{lemma}\label{pco:density}
Let $f\colon\pcoH\to\C$ be continuous with $f(yb)=\delta(b)^{-1}f(y)$ for $y\in\pcoH$, $b\in\pcoB$.
Then $\int_Kf(xu)\,du=\int_Kf(u)\,du$ for every $x\in\pcoH$.
\end{lemma}
\begin{proof}
\emph{Step 1 (Iwasawa).} Let $\pcoB_+\subset\pcoB$ be the elements with positive diagonal. The map
$K\times\pcoB_+\to\pcoH$, $(u,b)\mapsto ub$, is a homeomorphism: in each of the two blocks, list
the indices in the order of decreasing $\lambda$; then $\pcoB$ is the upper triangular group, and
the inverse map is the Gram--Schmidt process applied to the columns. In particular $\pcoH=K\pcoB$,
and $ub=u'b'$ ($u,u'\in K$, $b,b'\in\pcoB$) forces $b'b^{-1}=u'^{-1}u\in T$.

\emph{Step 2 (Haar measures of $\pcoB$).} Write $b=tn$ with $t$ diagonal invertible and $n$
unipotent in $\pcoB$; let $d^\times t$ and $dn$ be Haar measures of the two groups (both
unimodular). Left multiplication by $t_0n_0$ sends $(t,n)$ to $(t_0t,(t^{-1}n_0t)n)$, so
$d_lb=d^\times t\,dn$ is a left Haar measure. Right multiplication by $t_0$ sends $(t,n)$ to
$(tt_0,t_0^{-1}nt_0)$, and $n\mapsto t_0^{-1}nt_0$ multiplies $dn$ by $|\det\operatorname{Ad}(t_0)|
_{\mathfrak n}|^{-1}=\delta(t_0)^{-1}$ (real determinant on the Lie algebra $\mathfrak n$ of the
unipotent group, on which $\operatorname{Ad}(t_0)E_{ab}=(t_{0,a}/t_{0,b})E_{ab}$); right
multiplication by $n_0$ is $(t,n)\mapsto(t,nn_0)$. Hence $d_rb=\delta(b)\,d_lb$ is a right Haar
measure, where $\delta(tn)=\delta(t)$.

\emph{Step 3 (an integral formula).} Let $dh$ be a Haar measure of $\pcoH$, which is unimodular. For
$\varphi\in C_c(\pcoH)$ put $\Lambda'(\varphi)=\int_K\int_{\pcoB}\varphi(ub)\,d_rb\,du$; the inner
integrand has compact support $\pcoB\cap u^{-1}\operatorname{supp}\varphi$ and the result is
continuous in $u$, so $\Lambda'$ is a nonzero positive Radon functional. It is invariant under
$\varphi\mapsto\varphi(k^{-1}\cdot\,b_0)$ for $k\in K$, $b_0\in\pcoB$. The group $K\times\pcoB$ acts
on $\pcoH$ by $(k,b)\cdot y=kyb^{-1}$, transitively by Step 1, with stabiliser $\{(t,t):t\in T\}$ of
$1$, which is compact; by Step 1 the orbit map identifies $(K\times\pcoB)/T$ with $\pcoH$
homeomorphically. Both $dh$ and $\Lambda'$ are invariant measures on this homogeneous space, and
an invariant Radon measure on the quotient of a locally compact group $L$ by a compact subgroup $M$ is
unique up to a positive scalar. (If $\mu$ is one, $f\mapsto\int_{L/M}\int_Mf(xm)\,dm\,d\mu(xM)$ is a
left Haar measure on $L$; the map $f\mapsto\int_Mf(\cdot\,m)\,dm$ from $C_c(L)$ onto $C_c(L/M)$ is
surjective, since $\varphi=\int_M(\chi\cdot\varphi\circ p)(\cdot\,m)\,dm$ for $\chi\in C_c(L)$, $\chi\ge0$, with
$\int_M\chi(xm)dm=1$ on $p^{-1}(\operatorname{supp}\varphi)$, $p$ the projection; so uniqueness of Haar measure
gives that of $\mu$. See also \cite[\S2.6, Proposition~(2.48) and Theorem~(2.49), pp.~56--57]{FollandAHA}.) Normalising,
$\int_{\pcoH}\varphi\,dh=\int_K\int_{\pcoB}\varphi(ub)\,d_rb\,du$.

\emph{Step 4.} Choose $\psi\in C_c(\pcoB)$, $\psi\ge0$, $\psi\ne0$, with $\psi(tb)=\psi(b)$ for $t\in T$
(for instance a function of $|b_{aa}|$ and $|b_{ab}|$ only), and put $\Psi(ub)=\psi(b)$; by Step 1
this is well defined and continuous, with compact support $K\cdot T\operatorname{supp}\psi$. Let
$c_\psi=\int\psi\,d_lb>0$. For $x\in\pcoH$ let $A(x)=\int_{\pcoH}f(xh)\Psi(h)\,dh$. By Step 3 and
$f(xub)\psi(b)\,d_rb=f(xu)\psi(b)\,d_lb$,
\[
A(x)=\int_K\int_{\pcoB}f(xub)\psi(b)\,d_rb\,du=c_\psi\int_Kf(xu)\,du .
\]
By left invariance, $A(x)=\int_{\pcoH}f(h)\Psi(x^{-1}h)\,dh=\int_Kf(u)\bigl[\int_{\pcoB}\Psi(x^{-1}ub)
\delta(b)^{-1}d_rb\bigr]du$. Write $x^{-1}u=u'b'$ (Step 1); then $\Psi(x^{-1}ub)=\psi(b'b)$ and the
bracket is $\int\psi(b'b)\,d_lb=c_\psi$. Comparing the two expressions gives the claim.
\end{proof}

\subsubsection{Independence of the block diagonal factor}\label{pco:ss-indep}
For $x\in\mathrm{GL}_n(\C)$ let $A_x=\{a\in\pcoD:x^{-1}a\text{ is accretive}\}$; its elements are
invertible, $A_x$ is open and convex in $\pcoD$ (the accretive matrices form an open convex cone),
and $a\in A_x$ if and only if $a^{-1}x$ is accretive (Lemma~\ref{pco:minors}(d)). For $a\in A_x$ put
\begin{equation}\label{pco:eq-I}
I(x;a)=\int_K\frac{\pcom_\Lambda(au)\,\overline{\pcom_\Lambda(u)}}{\pcom_{-\mu}(u^{-1}a^{-1}xu)}\,du ;
\end{equation}
the denominator is nonzero by Lemma~\ref{pco:minors}(d), and continuous in $u$. If $x=kz\in G$ is
the polar decomposition, then $k\in A_x$ ($k^{-1}x=z$ is positive definite) and
$F(x)=d_\Lambda I(x;k)$.

\begin{proposition}\label{pco:indep}
For every $x\in\mathrm{GL}_n(\C)$, $a\mapsto I(x;a)$ is constant on $A_x$. Moreover the function
$\mathcal J(p,r)$ defined in the proof is locally constant on the open set $U_x$ defined there; in
particular $\mathcal J(p,r)=\mathcal J(1,1)$ whenever $(p,r)$ is joined to $(1,1)$ by a path in $U_x$.
\end{proposition}
\begin{proof}
By Lemma~\ref{pco:minors}(a), $\overline{\pcom_\Lambda(u)}=\pcom_\Lambda(u^*)$ and
$\pcom_{-\mu}(u^{-1}a^{-1}xu)=\pcom_\mu(u^*x^{-1}au)$. Let $U_x$ be the open set of pairs
$(p,r)\in\pcoH\times\pcoH$ with $rx^{-1}p$ accretive, and for $(p,r)\in U_x$ and $y\in\pcoH$ put
\[
f_{p,r}(y)=\frac{\pcom_\Lambda(py)\,\pcom_\Lambda(y^*r)}{\pcom_\mu(y^*rx^{-1}py)},\qquad
\mathcal J(p,r)=\int_Kf_{p,r}(u)\,du .
\]
By Lemma~\ref{pco:minors}(d), $f_{p,r}$ is defined and continuous on all of $\pcoH$, and
$I(x;a)=\mathcal J(a,1)$ for $a\in A_x$.

(i) \emph{$f_{p,r}$ is a density.} For $b\in\pcoB$, Corollary~\ref{pco:equiv}(a) gives
$f_{p,r}(yb)=b^\Lambda\,\overline{b^\Lambda}\,(\overline{b^\mu}b^\mu)^{-1}f_{p,r}(y)=|b^{\Lambda-\mu}|^2
f_{p,r}(y)=\delta(b)^{-1}f_{p,r}(y)$, since $\mu-\Lambda=2\rho_{\mathrm c}$.

(ii) \emph{Invariance.} For $y\in\pcoH$, $(py,y^*r)\in U_x$ and $f_{py,y^*r}(u)=f_{p,r}(yu)$, so
Lemma~\ref{pco:density} gives $\mathcal J(py,y^*r)=\mathcal J(p,r)$.

(iii) \emph{Holomorphy.} The integrand is jointly continuous on $U_x\times K$ and holomorphic in
$(p,r)$ (the $\pcom_w$ are rational functions whose denominators do not vanish there); since $K$ is
compact, $\mathcal J$ is holomorphic on $U_x$ and may be differentiated under the integral sign
(Cauchy's formula on small polydiscs and Fubini).

(iv) \emph{Polarisation.} Fix $(p,r)\in U_x$ and $X\in\pcoD$. For $s\in\C$ near $0$, (ii) with
$y=e^{sX}$ gives $\mathcal J(pe^{sX},e^{\bar sX^*}r)=\mathcal J(p,r)$. Write $\partial_1,\partial_2$
for the holomorphic differentials in $p$ and $r$. Differentiating along $s=\sigma\in\R$ and along
$s=i\sigma$ at $\sigma=0$,
\[
\partial_1\mathcal J[pX]+\partial_2\mathcal J[X^*r]=0,\qquad i\,\partial_1\mathcal J[pX]-
i\,\partial_2\mathcal J[X^*r]=0,
\]
so $\partial_1\mathcal J(p,r)[pX]=0$. As $X$ runs over $\pcoD$, so does $pX$; hence $\partial_1
\mathcal J=0$ on $U_x$. The function $a\mapsto\mathcal J(a,1)$ on the convex open set $A_x$ is
therefore holomorphic with vanishing differential, hence constant. Dividing the second identity by
$i$ and comparing with the first gives also $\partial_2\mathcal J(p,r)[X^*r]=0$; as $X$ runs over
$\pcoD$, so does $X^*r$ ($\pcoD$ is stable under $X\mapsto X^*$ and $r$ is invertible in $\pcoD$), so
$\partial_2\mathcal J=0$. Both holomorphic differentials of $\mathcal J$ vanish on $U_x$, so
$\mathcal J$ is locally constant there, hence constant along every path in $U_x$.
\end{proof}

\subsubsection{Local holomorphic form, $K$-behaviour and the Casimir eigenvalue}\label{pco:ss-casimir}
For $Z\in\mathfrak{gl}_n(\C)=\mathfrak g\otimes_\R\C$ and a $C^1$ function $\phi$ on $G$ let $R(Z)\phi$ be
the right derivative, $(R(X)\phi)(x)=\frac d{ds}\phi(xe^{sX})|_{s=0}$ for $X\in\mathfrak g$, extended
complex-linearly; then $R(Z_1)R(Z_2)-R(Z_2)R(Z_1)=R([Z_1,Z_2])$, and $\Omega=\sum_{a,b}E_{ab}E_{ba}$
acts by $R(\Omega)=\sum_{a,b}R(E_{ab})R(E_{ba})$. For a holomorphic function $\Phi$ on an open
subset of $\mathrm{GL}_n(\C)$ let $(\widetilde Z\Phi)(x)=d\Phi(x)[xZ]=\frac d{ds}\Phi(x(1+sZ))|_{s=0}$;
again $[\widetilde Z_1,\widetilde Z_2]=\widetilde{[Z_1,Z_2]}$. If $\phi=\Phi|_G$ near a point, then
$R(Z_1)\cdots R(Z_k)\phi=(\widetilde Z_1\cdots\widetilde Z_k\Phi)|_G$ there: for $Z\in\mathfrak g$ this is
the chain rule, and both sides are complex-linear in each $Z_i$.

\begin{lemma}\label{pco:local}
Let $g_0=k_0z_0\in G$ \textup(polar decomposition\textup) and let $W$ be the open set of
$x\in\mathrm{GL}_n(\C)$ with $k_0^{-1}x$ accretive. Then $\widetilde F(x)=d_\Lambda I(x;k_0)$ is
holomorphic on $W$, and $F=\widetilde F$ on $W\cap G$, a neighbourhood of $g_0$ in $G$. In particular
$F$ is real analytic on $G$, and $F(1)=1$.
\end{lemma}
\begin{proof}
For $x\in W\cap G$ with polar decomposition $x=kz$, both $k$ and $k_0$ lie in $A_x$, so
$F(x)=d_\Lambda I(x;k)=d_\Lambda I(x;k_0)$ by Proposition~\ref{pco:indep}. Holomorphy: the integrand
of $I(x;k_0)$ is jointly continuous on $W\times K$ and holomorphic in $x$, as in
Proposition~\ref{pco:indep}(iii). Finally $F(1)=d_\Lambda\int_K|\pcom_\Lambda|^2du/\pcom_{-\mu}(1)=1$.
\end{proof}

\begin{lemma}[$K$-behaviour]\label{pco:K}
Let $x\in G$, $a\in A_x$ and $k_1\in K$. Then
\begin{align}
F(xk_1)&=d_\Lambda\int_K\frac{\pcom_\Lambda(au)\,\pcom_\Lambda(u^*k_1)}{\pcom_{-\mu}(u^{-1}a^{-1}xu)}\,du,
\label{pco:eq-right}\\
F(k_1x)&=d_\Lambda\int_K\frac{\pcom_\Lambda(k_1au)\,\overline{\pcom_\Lambda(u)}}{\pcom_{-\mu}(u^{-1}a^{-1}xu)}
\,du.\label{pco:eq-left}
\end{align}
Consequently $F(xt)=t^\Lambda F(x)=F(tx)$ for $t\in T$, and if $\varepsilon_a=\varepsilon_b$,
$\lambda_a>\lambda_b$, then $R(E_{ab})F=0$.
\end{lemma}
\begin{proof}
$ak_1\in A_{xk_1}$ because $(ak_1)^{-1}xk_1=k_1^{-1}(a^{-1}x)k_1$; by Lemma~\ref{pco:local} and
Proposition~\ref{pco:indep} (applied at the point $xk_1$), $F(xk_1)=d_\Lambda I(xk_1;ak_1)$. Substitute
$u\mapsto k_1^{-1}u$ and use $\overline{\pcom_\Lambda(k_1^{-1}u)}=\pcom_\Lambda(u^*k_1)$
(Lemma~\ref{pco:minors}(a)). Similarly $k_1a\in A_{k_1x}$ and $(k_1a)^{-1}k_1x=a^{-1}x$, which gives
\eqref{pco:eq-left}. The weights follow from Corollary~\ref{pco:equiv}(b). For the derivative, fix
$x$ and $a\in A_x$; the right side of \eqref{pco:eq-right} is the restriction to $K$ of a
holomorphic function of $k_1\in\pcoH$, so $R(E_{ab})F(x)$ is its derivative along
$k_1=1+sE_{ab}$ at $s=0$, which vanishes by Corollary~\ref{pco:equiv}(b).
\end{proof}

\begin{proposition}[Casimir]\label{pco:casimir}
$R(\Omega)F=\kappa_\lambda F$ on $G$.
\end{proposition}
\begin{proof}
Fix $g_0$ and use Lemma~\ref{pco:local}: near $g_0$, $F=\widetilde F|_G$ with
$\widetilde F(x)=d_\Lambda\int_K\alpha(u)\,\psi(L_u(x))\,du$, where $\alpha(u)=\pcom_\Lambda(k_0u)
\overline{\pcom_\Lambda(u)}$, $L_u(x)=u^{-1}k_0^{-1}xu$ and $\psi=\pcom_{-\mu}^{-1}$ as in
Corollary~\ref{pco:equiv}(c). By Proposition~\ref{pco:indep}(iii) (applied to $\widetilde F$),
holomorphic derivatives of $\widetilde F$ may be taken under the integral sign. Since
$L_u(x(1+sZ))=L_u(x)(1+s\,u^{-1}Zu)$, $\widetilde Z(\psi\circ L_u)=(\widetilde{Z^u}\psi)\circ L_u$ with
$Z^u=u^{-1}Zu$, and so
\[
\widetilde\Omega(\psi\circ L_u)=\Bigl(\sum_{a,b}\widetilde{E_{ab}^u}\,\widetilde{E_{ba}^u}\,\psi\Bigr)\circ L_u
=(\widetilde\Omega\psi)\circ L_u,
\]
because $\sum_{a,b}c^{-1}E_{ab}c\otimes c^{-1}E_{ba}c=\sum_{a,b}E_{ab}\otimes E_{ba}$ for $c\in
\mathrm{GL}_n(\C)$ (as an operator on $\C^n\otimes\C^n$ this tensor is the flip, which commutes with
$c\otimes c$). It remains to show $\widetilde\Omega\psi=\kappa_\lambda\psi$ on the open set where
$\pcom_{-\mu}\ne0$. By Corollary~\ref{pco:equiv}(c), $\widetilde E_{aa}\psi=\mu_a\psi$ and
$\widetilde E_{ab}\psi=0$ if $\lambda_a<\lambda_b$. For a pair $\lambda_a<\lambda_b$,
$\widetilde E_{ba}\widetilde E_{ab}\psi=0$ and $\widetilde E_{ab}\widetilde E_{ba}\psi=\widetilde E_{ba}
\widetilde E_{ab}\psi+(\widetilde E_{aa}-\widetilde E_{bb})\psi=(\mu_a-\mu_b)\psi$. Summing,
\[
\widetilde\Omega\psi=\Bigl(\sum_a\mu_a^2-\sum_{\lambda_a<\lambda_b}(\mu_b-\mu_a)\Bigr)\psi
=(|\mu|^2-2\langle\mu,\rho\rangle)\psi=(|\lambda|^2-|\rho|^2)\psi,
\]
since $2\rho=\sum_{\lambda_a<\lambda_b}(e_b-e_a)$ and $\mu=\lambda+\rho$.
\end{proof}

\begin{lemma}[Decay]\label{pco:decay}
For every integer $0\le c\le h$ and $g\in G$, $|F(g)|\le(2/\|g\|)^c$. If $h\ge n^2+1$, then $F\in
L^1(G)\cap L^2(G)\cap L^\infty(G)$.
\end{lemma}
\begin{proof}
The proof of Theorem~\ref{pds:decay} uses \eqref{pds:eq-FJ} only through its right side, which is
$F$, together with Lemmas~\ref{pds:blockdet} and \ref{pds:fischer}; these are stated for an
arbitrary $\auU(J)$, an arbitrary regular $\lambda\in(\frac{n-1}2+\Z)^n$ and its gap. So it proves
the bound for $F$. With $c=h$, $\int_G\|g\|^{-c}dg<\infty$ by the proof of
Lemma~\ref{pm:kernel}(c), which is stated there for $G_\tau$ but applies verbatim to any $\auU(p,q)$,
$n=p+q$: in $G=KA^+K$, with coordinates $t_1\ge\dots\ge t_{\min(p,q)}\ge0$ on $A^+$, one has
$\|kak'\|=e^{t_1}$ and Haar density $\delta(a)\le e^{n(n-1)t_1}$ (\cite[Ch.~I, \S5, no.~2]{Helgason}, \cite[Lemma~2.4.2]{WallachRRG1}: $\delta(a)=\prod_{\alpha>0}\sinh\alpha(\log a)$, with multiplicities, up to a constant), so the integral
is bounded by a constant times $\int_{A^+}e^{(n(n-1)-c)t_1}\,da<\infty$ for $c>n(n-1)$, in particular
for $c=h\ge n^2+1$. Also $|F|\le2^c$.
\end{proof}

\subsubsection{Identification with $c_\lambda$}\label{pco:ss-ident}
Let $(\rho_\Lambda,V_\Lambda)$ be the irreducible representation of $K$ of highest weight $\Lambda$
for the positive compact roots $e_a-e_b$, $\varepsilon_a=\varepsilon_b$, $\lambda_a>\lambda_b$, with
unit highest weight vector $v_\Lambda$, and $\beta(k)=\dim V_\Lambda\cdot\overline{\langle v_\Lambda,
\rho_\Lambda(k)v_\Lambda\rangle}$. We use two standard facts on a unitary representation $\sigma$ of
$K$, with $\sigma(\beta)=\int_K\beta(k)\sigma(k)\,dk$.
\begin{enumerate}
\item[(K1)] $\sigma(\beta)$ is the orthogonal projection onto the space of vectors of the
$\Lambda$-isotypic subspace that are highest weight vectors of weight $\Lambda$. (Peter--Weyl and
Schur orthogonality: on an irreducible summand $\cong V_\Lambda$, $\sigma(\beta)u=\langle v_\Lambda,u
\rangle v_\Lambda$ in the identification, and on other types $\sigma(\beta)=0$.)
\item[(K2)] If $\sigma$ is finite-dimensional and $y$ satisfies $\sigma(t)y=t^\Lambda y$ ($t\in T$) and
$d\sigma(E_{ab})y=0$ for the positive compact roots, then $\sigma(\beta)y=y$. (Decompose $y$ along
irreducible summands; in an irreducible module a vector killed by the positive root vectors is a
highest weight vector \cite[Theorem~5.5(c), p.~279]{KnappBeyond}; then apply (K1).)
\end{enumerate}

\begin{lemma}\label{pco:Kfin}
For every $x\in G$, $\int_K\beta(k)F(xk)\,dk=F(x)$ and $\int_K\beta(k)\,\overline{F(k^{-1}x)}\,dk=
\overline{F(x)}$.
\end{lemma}
\begin{proof}
Let $\mathcal P$ be the space of functions on $K$ of the form $(\det k)^{-e}P(k)$, with $P$ a
polynomial of degree $\le D$ in the entries of $k$ ($e,D$ fixed and large enough). It is
finite-dimensional and stable under right translation $\sigma(k_1)\phi=\phi(\cdot\,k_1)$, since the
entries of $kk_1$ are linear in those of $k$ and $\det(kk_1)=\det k\det k_1$. Let $\mathcal
V\subset\mathcal P$ be the span of the functions $k\mapsto\pcom_\Lambda(Yk)$, $Y\in
\mathrm{GL}_n(\C)$; it is finite-dimensional, hence closed in $\mathcal P$. By
Corollary~\ref{pco:equiv}(b) every $\phi\in\mathcal V$ satisfies $\sigma(t)\phi=t^\Lambda\phi$
($t\in T$) and $d\sigma(E_{ab})\phi=0$ (the derivative of the holomorphic extension along
$1+sE_{ab}$); these are pointwise properties of $\phi$. By (K2) applied to $\sigma$ on $\mathcal P$,
$\int_K\beta(k)\phi(k_1k)\,dk=\phi(k_1)$ for $\phi\in\mathcal V$. By \eqref{pco:eq-right} the function
$k_1\mapsto F(xk_1)$ is an integral of elements of $\mathcal V$ against a continuous function of
$u$, hence lies in $\mathcal V$; evaluating at $k_1=1$ gives the first identity. By
\eqref{pco:eq-left} and Lemma~\ref{pco:minors}(a), $\phi(k_1)=\overline{F(k_1^{-1}x)}=d_\Lambda\int_K
\pcom_\Lambda(u^*a^*k_1)\,(\cdots)\,du$ also lies in $\mathcal V$, which gives the second.
\end{proof}

\begin{proposition}[Rigidity for $F$]\label{pco:rigid}
Let $R$ be the right regular representation of $G$ on $L^2(G)$. Then $\langle F,R(g)F\rangle=
\|F\|^2c_\lambda(g)$ for $g\in G$; consequently $c_\lambda\in L^1(G)$, and with $e_\lambda=d_\lambda
\overline{c_\lambda}$ \textup($d_\lambda$ the formal degree\textup), $R(e_\lambda)F=F$.
\end{proposition}
\begin{proof}
We follow the proof of Lemma~\ref{par:rigid}, with its hypotheses on $y=F$ replaced by the
following three properties, which hold by Lemmas~\ref{pco:decay}, \ref{pco:Kfin} and
Proposition~\ref{pco:casimir}: (1) $F\in L^2(G)$ and $F=R(\beta)F$, where $R(\beta)=\int_K\beta(k)R(k)dk$
(by Lemma~\ref{pco:Kfin}, pointwise, hence in $L^2$); (2) for every $\varphi\in C_c^\infty(G)$,
$\kappa_\lambda R(\varphi)F=R(\Omega\varphi)F$, where $\Omega\varphi=R(\Omega)\varphi$. For (2):
$(R(\varphi)F)(x)=\int\varphi(h)F(xh)\,dh$, and for $X\in\mathfrak g$, substituting $h\mapsto he^{-sX}$
and differentiating under the integral sign (compact support, $F\in C^\infty$),
$R(\varphi)(R(X)F)=-R(R(X)\varphi)F$; iterating, $R(\varphi)(R(\Omega)F)=R(\sum_{a,b}R(E_{ba})R(E_{ab})
\varphi)F=R(\Omega\varphi)F$, and $R(\Omega)F=\kappa_\lambda F$.

Let $H_F$ be the closed span of $R(G)F$, $H_F=\int^\oplus H_\zeta\,d\nu(\zeta)$ a decomposition into
irreducible unitary representations \cite{Dixmier}, $F=\int^\oplus y_\zeta$. Fix a sequence
$\varphi_i\in C_c^\infty(G)$, $\varphi_i\ge0$, $\int\varphi_i=1$, with supports shrinking to $\{1\}$.
The identities $F=R(\beta)F$ and $\kappa_\lambda R(\varphi_i)F=R(\Omega\varphi_i)F$ ($i\ge1$) are
countably many and decompose; so for almost every $\zeta$, $y_\zeta=\sigma_\zeta(\beta)y_\zeta$ and
$\kappa_\lambda\sigma_\zeta(\varphi_i)y_\zeta=\sigma_\zeta(\Omega\varphi_i)y_\zeta$ for all $i$. By (K1),
$y_\zeta$ is a highest weight vector of weight $\Lambda$ in the $\Lambda$-isotypic subspace of
$H_\zeta$, which is finite-dimensional (admissibility, \cite[Theorem~8.1, p.~205]{Knapp}); so $y_\zeta$ is
$K$-finite, hence a smooth vector, and $\sigma_\zeta(\Omega\varphi_i)y_\zeta=\sigma_\zeta(\varphi_i)
d\sigma_\zeta(\Omega)y_\zeta$ by the same integration by parts. Letting $i\to\infty$,
$d\sigma_\zeta(\Omega)y_\zeta=\kappa_\lambda y_\zeta$. From here the proof of Lemma~\ref{par:rigid}
applies word for word: if $y_\zeta\ne0$, (D1) and (D2) give the infinitesimal character $\lambda$,
Lemma~\ref{par:D3} gives $H_\zeta\cong\pi_\lambda$ and (D3)(b) gives $y_\zeta=a(\zeta)U_\zeta v$, and $\langle F,R(g)F\rangle=
\int|a|^2d\nu\cdot c_\lambda(g)=\|F\|^2c_\lambda(g)$.

Since $F\in L^1$, $\int_G|\langle F,R(g)F\rangle|\,dg\le\int\!\!\int|F(x)||F(xg)|\,dx\,dg=\|F\|_1^2$,
and $\|F\|\ne0$ ($F$ is continuous with $F(1)=1$); so $c_\lambda\in L^1$. The last assertion is the
last step of the proof of Lemma~\ref{par:rigid}: by Schur orthogonality $\int e_\lambda(h)c_\lambda
(g^{-1}h)\,dh=c_\lambda(g^{-1})$, so $\langle R(g)F,R(e_\lambda)F-F\rangle=0$ for all $g$, and
$R(e_\lambda)F\in H_F$.
\end{proof}

\begin{proof}[Proof of Theorem~\ref{pco:FJ}]
Write $\pi=\pi_\lambda$, $c=c_\lambda$, $e=e_\lambda=d\bar c$ with $d=d_\lambda$.

\emph{Step 1: $F(g)=d\int_GF(x)\,c(x^{-1}g)\,dx$.} By Proposition~\ref{pco:rigid}, $F=R(e)F$ in
$L^2$. Since $F$ is bounded and continuous and $e\in L^1$, $g\mapsto\int e(h)F(gh)\,dh$ is
continuous (dominated convergence), so $F(g)=\int e(h)F(gh)\,dh$ for every $g$. Insert
$\overline{c(h)}=\|F\|^{-2}\overline{\langle F,R(h)F\rangle}=\|F\|^{-2}\int F(x)\overline{F(xh)}\,
dx$; the double integral converges absolutely (it is bounded by $\|F\|_\infty\|F\|_1^2$), and by
Fubini
\[
F(g)=\frac d{\|F\|^2}\int F(x)\int\overline{F(xh)}F(gh)\,dh\,dx=\frac d{\|F\|^2}\int F(x)\langle R(x)F,
R(g)F\rangle\,dx=d\int F(x)\,c(x^{-1}g)\,dx .
\]

\emph{Step 2: $F(g)=\langle w,\pi(g)v\rangle$ with $w=d\int_G\overline{F(x)}\,\pi(x)v\,dx$.} The
integral converges in norm since $F\in L^1$, and $c(x^{-1}g)=\langle\pi(x)v,\pi(g)v\rangle$, with inner
products conjugate-linear in the first variable.

\emph{Step 3: $w$ is a multiple of $v$.} By Step 2 and the second identity of Lemma~\ref{pco:Kfin},
for all $g$,
\[
\langle\pi(\beta)w,\pi(g)v\rangle=\int_K\overline{\beta(k)}\langle w,\pi(k^{-1}g)v\rangle dk
=\overline{\int_K\beta(k)\overline{F(k^{-1}g)}\,dk}=F(g)=\langle w,\pi(g)v\rangle .
\]
Since $\pi$ is irreducible, the $\pi(g)v$ span a dense subspace, so $\pi(\beta)w=w$. By (K1), $w$ is
a highest weight vector of weight $\Lambda$ in the $\Lambda$-isotypic subspace of $\pi$, which is
irreducible by (D3)(b); so $w=\alpha v$ for some $\alpha\in\C$.

\emph{Step 4.} Hence $F=\bar\alpha\,c$. Evaluating at $g=1$, $1=F(1)=\bar\alpha\,c(1)=\bar\alpha$
(Lemma~\ref{pco:local}). So $F=c$.
\end{proof}

\begin{remark}\label{pco:rem-FJ}
The proof uses: Lemmas~\ref{pds:blockdet}, \ref{pds:fischer} (Fischer's inequality, Jacobi's
identity), Lemma~\ref{pm:kernel}(c), the published results (D1)--(D3) and \cite{Dixmier} as in
Lemma~\ref{par:rigid}, Schur orthogonality for square-integrable representations, admissibility
and smoothness of $K$-finite vectors \cite[Theorem~8.1 and Proposition~8.5]{Knapp}, highest weight theory for $K$, and
uniqueness of invariant measures on a quotient by a compact subgroup. No property of the signature
or of the special parameters of \S\ref{pds:ss-setup} is used beyond $h\ge n^2+1$. The
non-vanishing of $F$ and the normalisation at the identity are both $F(1)=1$, which is immediate
from \eqref{pco:eq-F}; the scalar $\alpha$ in Step~3 is not assumed nonzero.
\end{remark}

\subsubsection{The vacuum character (V)}\label{pco:ss-V}
Fix $0\le j<m$ and $\tau$, and let $V=Y_j\otimes_{E,\tau}\C$, $V'=Y_{j+1}\otimes_{E,\tau}\C$, of
signatures $(p,q)$, $(p',q')$ and ranks $n,n'$, and $W=\auHom(V,V')$ with the form
$\langle w_1,w_2\rangle=\pcotr(w_1^\dagger w_2)$ of \S\ref{ps:ss-weil}; it has signature
$(P,Q)=(pp'+qq',pq'+qp')$. Let $\iota\colon\auU(V)\times\auU(V')\to\auU(W)$, $\iota(g,h)w=hwg^{-1}$;
then $\det_W\iota(g,h)=\det(h)^n\det(g)^{-n'}$, and on Lie algebras
$\pcotr_W\iota(X,Y)=n\,\pcotr Y-n'\pcotr X$. Let $\omega_\infty$ be the component at $\tau$ of the
Weil representation $\omega_W$ of $\auU(W)$ defined with $\upsilon$ (\S\ref{ps:ss-weil}), acting on
the Fock space $\mathfrak F_\tau$ of the Heisenberg representation $\varrho_\tau$, so that
$\omega_{\auHom(V,V')}(g,h)=\omega_\infty(\iota(g,h))$ at $\tau$. Let $\mathrm{Mp}$ be the metaplectic
double cover of the symplectic group of $W$ viewed as a real symplectic space (symplectic form
$\Iml$ of the hermitian form, as in $\beta_\infty$), $\omega^{\mathrm{Mp}}$ the metaplectic
(oscillator) representation on $\mathfrak F_\tau$, characterised up to a scalar on each element by
$\omega^{\mathrm{Mp}}(\tilde u)\varrho_\tau(w)\omega^{\mathrm{Mp}}(\tilde u)^{-1}=\varrho_\tau(uw)$
\cite[Th\'eor\`eme~1 and no.~34]{WeilActa}, and $\widetilde{\auU}(W)$ the inverse image of $\auU(W)(\R)$ in $\mathrm{Mp}$.

We state (V) in the form in which Lemma~\ref{par:vac} uses it, the derived action on polynomial
vectors:
\begin{enumerate}
\item[(V$'$)] For $X\in\mathfrak u(V)$, $Y\in\mathfrak u(V')$, on polynomial vectors,
$d\omega_\infty(\iota(X,Y))=d\omega^{\mathrm{Mp}}(\iota(X,Y))+\tfrac{n'}2\pcotr X-\tfrac n2\pcotr Y$.
\end{enumerate}
Here $d\omega^{\mathrm{Mp}}$ on polynomial vectors is given by the Weyl-ordered quadratic operators
in the $a^\pm_{rc}$, in the following sense. Let $\mathcal V$ be the span of the operators $\frac12(a_ia_k+a_ka_i)$ on polynomials, $a_i,a_k$ running over the $a^\pm_{rc}$. By Howe \cite[p.~537, (2.2), and (2.3)]{HoweJAMS}, the polynomials are the $\widetilde K$-finite vectors of $\omega^{\mathrm{Mp}}$ for the maximal compact subgroup $\widetilde K$ of $\mathrm{Mp}$ (so they are smooth vectors, stable under $d\omega^{\mathrm{Mp}}$), and $d\omega^{\mathrm{Mp}}$ maps the complexified Lie algebra onto $\mathcal V$ (there: the spans of $\frac12(z_i\partial_j+\partial_jz_i)$, $z_iz_j$, $\partial_i\partial_j$, which together span $\mathcal V$). Differentiating Weil's relation on polynomial vectors (which are smooth for $\varrho_\tau$, with $d\varrho_\tau$ given by linear combinations of the $a^\pm_{rc}$ and scalars) gives $[d\omega^{\mathrm{Mp}}(X),d\varrho_\tau(w)]=d\varrho_\tau(Xw)$. This determines $d\omega^{\mathrm{Mp}}(X)$ within $\mathcal V$: if $Q\in\mathcal V$ commutes with all $a^\pm_{rc}$, then $Q$ commutes with all $z_{rc}$ and $\partial_{rc}$, so $Q1$ is a constant $c$ and $Qf=fQ1=cf$; and a scalar lies in $\mathcal V$ only if it is $0$ (writing $\mathcal V$ as the span of $z_iz_j$, $\partial_i\partial_j$, $z_i\partial_j+\frac12\delta_{ij}$, up to signs: if a combination of these equals $c$, applying it to the polynomials of degree $\le2$ shows that all coefficients vanish, hence $c=0$). So $d\omega^{\mathrm{Mp}}(X)$ is the unique element of $\mathcal V$ whose commutator with $d\varrho_\tau(w)$ is $d\varrho_\tau(Xw)$; this is what ``Weyl-ordered quadratic operator'' means here. Exponentiating, (V$'$) says that
$\omega_\infty(\iota(g,h))=\det(g)^{n'/2}\det(h)^{-n/2}\,\omega^{\mathrm{Mp}}(\widetilde{\iota(g,h)})$,
the product of the two genuine factors being a representation of the linear group; this is (V).

We use the following properties of the archimedean Weil representation, for every nondegenerate
hermitian space $W$ over $E_\tau$ (not only those of the form $\auHom(V,V')$); Proposition~\ref{pwf:W2}(c) and Lemma~\ref{pwf:sum} of \S\ref{pwf:sec}, on which they rest, are
stated for an arbitrary nondegenerate hermitian space.
\begin{enumerate}
\item[(W1$_\infty$)] $\omega_\infty(u)\varrho_\tau(w)\omega_\infty(u)^{-1}=\varrho_\tau(uw)$ for
$u\in\auU(W)(\R)$; and $\omega_\infty$ is strongly continuous. (The component at $\tau$ of (W1).)
\item[(W2$_\infty$)] Let $I\subset W$ be an isotropic complex line, $I^*$ an isotropic line with
$\langle I^*,I\rangle\ne0$, $W_I=(I+I^*)^\perp$, and $P_I$ the stabiliser of $I$ in $\auU(W)(\R)$.
There is a nonzero continuous linear map $\mathsf Z\colon\mathfrak F_\tau^{\mathrm{sm}}\to
\mathfrak F_\tau(W_I)^{\mathrm{sm}}$ (evaluation at $0$ in the mixed model $\mathcal S(I^*)\otimes
\mathfrak F_\tau(W_I)$) with $\mathsf Z\,\omega_\infty(p)=\chi_{I,\tau}(p)\,\omega_{W_I,\infty}(\bar p)
\mathsf Z$ for $p\in P_I$, where $\chi_{I,\tau}(p)=\upsilon_\tau(a)^{-1}|a|$ if $p$ acts on $I$ by
$a\in\C^\times$. (The component at $\tau$ of (W2) and (W2$'$); it is proved in
Proposition~\ref{pwf:W2}(c), local statement, from the Levi action of the Siegel parabolic in
Kudla's Schr\"odinger model \cite[Theorem~3.1 and \S5]{KudlaSplitting} (Lemma~\ref{pwf:siegel}) and
the compatibility with orthogonal sums (Lemma~\ref{pwf:sum}).)
\item[(M)] In the same mixed model, for $X$ in the Lie algebra of the subgroup of $P_I$ that acts
trivially on $W_I$ and preserves $I$ and $I^*$, $\mathsf Z\,d\omega^{\mathrm{Mp}}(X)=\pm\frac12
\pcotr_\R(X|_{I})\,\mathsf Z$ on smooth vectors, where $\pcotr_\R$ is the trace on $I$ viewed as a
real plane. (The Levi subgroup $\mathrm{GL}(I)$ of the Siegel parabolic of $\mathrm{Sp}(I+I^*)$ acts in
the Schr\"odinger model by $\phi(x)\mapsto|\det_\R a|^{\pm1/2}\phi(a^{\mp1}x)$ up to a scalar, these operators being intertwiners for Weil's relation, as in the proof of Lemma~\ref{pwf:siegel}. Infinitesimally: by the description of $d\omega^{\mathrm{Mp}}$ just given, transported to the Schr\"odinger model on the first factor (the commutation relations determine the element of the span of symmetrised quadratics in $x_i$, $\partial_i$ in the same way), if $X$ acts on the coordinates of $I^*$ by a real matrix $A$, then $d\omega^{\mathrm{Mp}}(X)=\mp\frac12\sum_{i,k}A_{ik}(x_i\partial_k+\partial_kx_i)=\mp(\sum_{i,k}A_{ik}x_i\partial_k+\frac12\pcotr A)$; evaluation at $0$ kills the vector field.)
\item[(S)] (Only for definite $W$.) For a hermitian line $W'$ over $E_\tau$, under
$\mathfrak F_\tau(W\oplus W')=\mathfrak F_\tau(W)\hat\otimes\mathfrak F_\tau(W')$,
$\omega_{W\oplus W',\infty}(u\oplus1)=\omega_{W,\infty}(u)\otimes1$ for $u\in\auU(W)(\R)$
(compatibility of the $\upsilon$-splittings with orthogonal sums, Lemma~\ref{pwf:sum}; cf.\ \cite[p.~389]{KudlaSplitting}).
\end{enumerate}

\begin{theorem}\label{pco:V}
Assume \textup{(W1$_\infty$)}, \textup{(W2$_\infty$)}, \textup{(M)} for every nondegenerate hermitian
space over $E_\tau$ (in particular for $W$ and, if $W$ is definite, for $W\oplus W'$ as in Step~3),
and, if $W$ is definite, \textup{(S)}. Then \textup{(V$'$)} holds; that is, \textup{(V)} of Lemma~\ref{par:vac} holds at
$\tau$, with $\upsilon_\tau|_{\auU(1)}(z)=z$.
\end{theorem}
\begin{proof}
\emph{Step 1: the ratio is a character.} For $\tilde u\in\widetilde{\auU}(W)$ over $u$, by
(W1$_\infty$) the operator $\omega_\infty(u)\omega^{\mathrm{Mp}}(\tilde u)^{-1}$ commutes with every
$\varrho_\tau(w)$; $\varrho_\tau$ is irreducible, so it is a scalar $\chi(\tilde u)$ of modulus one.
Then $\chi$ is a continuous character of $\widetilde{\auU}(W)$, and on smooth (in particular
polynomial) vectors $d\omega_\infty(X)=d\omega^{\mathrm{Mp}}(X)+d\chi(X)$ for $X\in\mathfrak u(W)$.
Since $d\chi$ vanishes on $[\mathfrak u(W),\mathfrak u(W)]=\mathfrak{su}(W)$, $d\chi=c\cdot\pcotr_W$
for some $c\in\C$. By the trace formula for $\iota$, (V$'$) is equivalent to $c=-\frac12$.

\emph{Step 2: $W$ indefinite.} Then $P,Q\ge1$ and there are $I,I^*$ as in (W2$_\infty$). The
elements $m_\theta\in\auU(W)(\R)$ acting by $e^{i\theta}$ on $I+I^*$ and trivially on $W_I$ lie in
$P_I$, with $a=e^{i\theta}$ and $\bar m_\theta=1$; their generator $X_0$ has $\pcotr_WX_0=2i$ and
$\pcotr_\R(X_0|_I)=0$ (multiplication by $i$ on a real plane). Differentiating
(W2$_\infty$) along $m_\theta$ at $\theta=0$ on a smooth vector $\Phi$ (continuity of $\mathsf Z$),
with $\upsilon_\tau(e^{i\theta})=e^{i\theta}$, gives $\mathsf Z\,d\omega_\infty(X_0)\Phi=-i\,\mathsf Z\Phi$;
and (M) gives $\mathsf Z\,d\omega^{\mathrm{Mp}}(X_0)\Phi=0$. By Step 1, $d\chi(X_0)\mathsf Z\Phi=-i
\,\mathsf Z\Phi$; choosing $\Phi$ with $\mathsf Z\Phi\ne0$, $2ic=-i$, so $c=-\frac12$.

\emph{Step 3: $W$ definite.} Let $W'$ be a line with form of the sign opposite to that of $W$; then
$W\oplus W'$ is indefinite, and Step~2 gives $d\chi_{W\oplus W'}=-\frac12\pcotr_{W\oplus W'}$. By
(S), $d\omega_{W,\infty}(X)\otimes1=d\omega_{W\oplus W',\infty}(X\oplus0)$, and the Weyl-ordered
operator of $X\oplus0$ is that of $X$ tensored with $1$ (it involves only the variables of $W$); so
$d\chi_W(X)=d\chi_{W\oplus W'}(X\oplus0)=-\frac12\pcotr_WX$.
\end{proof}

\begin{remark}\label{pco:rem-V}
(a) $W_\tau$ is indefinite as soon as $Y_j$ or $Y_{j+1}$ is indefinite at $\tau$: if
$p,q\ge1$ and $n'\ge1$, then $P,Q\ge1$; if $V$ is definite and $p',q'\ge1$, likewise. By
\S\ref{ps:ss-Y}, $Y_j$ contains the string $f_0,\dots,f_{s_j-1}$, which is indefinite for
$s_j\ge1$, i.e.\ $j\le m-2$; for $j=m-1$ both $Y_{m-1}=X_{m-1}$ and $Y_m$ are positive definite, so
among the pairs $(Y_j,Y_{j+1})$, (S) is used exactly for the last pair $j=m-1$. (In general (S) is
needed exactly when $W_\tau$ is definite.)
For definite pairs (V$'$) is the formula of \S\ref{pk:ss-conv} (``no further multiplier''); the
argument there identifies $\omega_\infty$ with $F\mapsto F(h^tz\bar g)$ only up to a character,
and Step~3 supplies the character.

(b) The alternative of \emph{defining} $\omega_\infty$ by the Fock-model formulas with the
determinant twists, as an explicit representation, does not remove the need for a comparison:
(W1) requires $\theta\circ\omega(\gamma)=\theta$ for $\gamma\in\auU(W)(E^+)$, and multiplying
$\omega_\infty$ by a character $\prod_\tau\det(g_\tau)^{a_\tau}$ multiplies $\theta\circ\omega(\gamma)$
by $\prod_\tau\tau(\det\gamma)^{a_\tau}$, which is not identically $1$ on $\auU(W)(E^+)$ unless all
$a_\tau=0$. So with that definition (W1) would have to be proved, and that is equivalent to
Theorem~\ref{pco:V}. We therefore keep the definition of \S\ref{ps:ss-weil} and derive (V) from
the local properties (W1$_\infty$), (W2$_\infty$), (M), (S).

(c) (W2$_\infty$) is the only input through which the value of $\upsilon_\tau$ enters the
computation of $c$; (M) and the Weyl-ordered form of $d\omega^{\mathrm{Mp}}$ are statements about the
metaplectic representation alone, and (S) only transports $c$ from $W\oplus W'$ to $W$. The signs
in (M) do not matter, because $\pcotr_\R(X_0|_I)=0$.
\end{remark}


%% ===== proofs/A-weil-foundations.tex =====
% proofs/A-weil-foundations.tex
% Foundations of the Weil representation used in Lemma ps:weil: the global invariance of theta
% under rational points (M4), the mixed model and the zero term (W2), (W3) (M3), the model at
% split places (W5) (M2), the formal degree of Proposition pds:haar (M6), and the archimedean
% compact scalar of C-threeex2 (two proofs of e_+=0). Labels of this file start with pwf:.
% Only amsmath/amssymb/amsthm are used.
\providecommand{\auA}{\mathbb A}
\providecommand{\auU}{\mathrm U}
\providecommand{\auGL}{\mathrm{GL}}
\providecommand{\auHom}{\operatorname{Hom}}
\providecommand{\auS}{\mathcal S}
\providecommand{\aubar}[1]{\overline{#1}}
\providecommand{\psHS}{\mathcal H}
\providecommand{\psrho}{\varrho}
\providecommand{\psZ}{Z}
\providecommand{\pwfA}{\mathsf A}
\providecommand{\pwfL}{\mathsf L}
\providecommand{\pwfM}{\mathsf M}
\providecommand{\pwfev}{\operatorname{ev}_0}

\subsection{Foundations of the Weil representation}\label{pwf:sec}

Lemma~\ref{ps:weil} collects the properties of the Weil representation on which Theorems A, B, C
rest. For several of them the works cited there do not contain the statement: the invariance of
$\theta$ under the rational points for the splitting by characters (the cited works on the
splitting are local), the mixed model attached to a non-Lagrangian isotropic subspace, and the model
at places of $E^+$ split in $E$. This subsection proves these statements from the definition of
$\omega$ and from statements that are contained, with the locators given, in Weil~\cite{WeilActa}
and Kudla~\cite{KudlaSplitting}. It also replaces the appeal to Harish-Chandra's formula for the
formal degree in Proposition~\ref{pds:haar} by a result proved in the paper, and proves the
archimedean normalisation \eqref{ptf:eq-compact} by a computation on spaces of rank one.

\emph{Summary.} Theorem~\ref{pwf:inv} (the theta functional is invariant under $\auU(W)(E^+)$) is
proved in full from Definition~\ref{pwf:def} and the published statements (P1)--(P5) below.
Propositions~\ref{pwf:W2} and~\ref{pwf:W3} prove (W2), (W2$'$) and the identity of (W3) in full
(the bound of moderate growth in (W3) is Lemma~\ref{pz:absgrowth}(b), proved from (W0), (W1), (W2),
(W4), not from (W3)). Proposition~\ref{pwf:W5} proves (W5) and
\eqref{pl:eq-w5}, for the normalisation at split places fixed in Definition~\ref{pwf:def}(D3).
Proposition~\ref{pwf:haar} proves Proposition~\ref{pds:haar} for the parameters to which it is
applied. Lemma~\ref{pwf:delta-red} and Proposition~\ref{pwf:eplus} prove the archimedean
normalisation \eqref{ptf:eq-compact}.

\subsubsection{Conventions, published inputs, and the definition of $\omega$}\label{pwf:ss-def}

\emph{Notation.} $k=E^+$; for a place $v$ of $k$ put $F=k_v$ and $E_v=E\otimes_kF$, which is a
quadratic field extension of $F$ or $F\times F$. All hermitian forms are conjugate-linear in the
first variable. For a non-degenerate hermitian space $W$ over $E$ (or $E_v$) put
\[
\pwfA(x,y)=\operatorname{Tr}_{E/k}\bigl(\jmath\langle y,x\rangle\bigr),\qquad
\beta(x,y)=\psi\bigl(\tfrac12\jmath\langle y,x\rangle\bigr)=\psi_k\bigl(\tfrac12\pwfA(x,y)\bigr),\qquad
[x,y]=\beta(x,y)\beta(y,x)^{-1}=\psi_k\bigl(\pwfA(x,y)\bigr),
\]
where $\psi=\psi_k\circ\operatorname{Tr}_{E/k}$ and $\psi_k=\psi_\Q\circ\operatorname{Tr}_{k/\Q}$
(\S\ref{ps:ss-weil}); the second equality holds because $\operatorname{Tr}(\jmath\langle x,y\rangle)
=\operatorname{Tr}(\aubar{\jmath\langle x,y\rangle})=-\operatorname{Tr}(\jmath\langle y,x\rangle)$.
Thus $\pwfA$ is a non-degenerate alternating $k$-bilinear form, $\beta$ is bimultiplicative, and
$\beta(x,y)=1$ whenever $\pwfA(x,y)=0$. A subspace $X$ is \emph{$\pwfA$-isotropic} if $\pwfA$
vanishes on $X\times X$; then $\beta=1$ on $X\times X$. Every $E$-subspace isotropic for
$\langle\,,\rangle$ is $\pwfA$-isotropic. $\pwfA$ is unchanged if $\langle\,,\rangle$ is replaced by
$\jmath\langle\,,\rangle$ and $\jmath$ by $1$; so $W$ with the skew-hermitian form
$\jmath\langle\,,\rangle$ is the space $\mathbb W=V\otimes_EW$ of~\cite{KudlaSplitting} for the
hermitian line $V=E$ (case $3_+$, $m=\dim V=1$), and $\auU(W)$ is the group $G=\auU(W)$ there.

A \emph{Heisenberg representation} of $W_v$ (resp.\ $W(\auA)$) is a pair $(\psHS,\psrho)$ of a
Hilbert space and a strongly continuous map $\psrho$ to unitary operators with
$\psrho(x)\psrho(y)=\beta(x,y)\psrho(x+y)$, irreducible; $\psHS^\infty$ denotes its smooth vectors.
An \emph{intertwiner} between two of them (possibly for different spaces, along an isometry
$\sigma$) is a unitary $T$ with $T\psrho(x)=\psrho'(\sigma x)T$.

\emph{Published inputs.} We use exactly the following.
\begin{itemize}
\item[(P1)] (Stone--von Neumann, and \cite[Th\'eor\`eme~1 and no.~34]{WeilActa}.) Up to
isomorphism there is one Heisenberg representation of $W_v$, and one of $W(\auA)$; an
intertwiner, if it exists, is unique up to a scalar of modulus $1$. For $g\in\mathrm{Sp}(W_v,\pwfA)$
there is a unitary $T$, unique up to a scalar, with $T\psrho(x)T^{-1}=\psrho(gx)$; it preserves
$\psHS^\infty$ and is continuous on it \cite[no.~35]{WeilActa}. A continuous linear operator on
$\psHS^\infty$ commuting with all $\psrho(x)$ is a scalar.
\item[(P2)] (\emph{Schr\"odinger model}.) Let $W=X\oplus Y$ with $X,Y$ $\pwfA$-isotropic. On
$\auS(X)$ (Schwartz--Bruhat functions) put $\psrho(x_0)\varphi(x)=\varphi(x+x_0)$ ($x_0\in X$),
$\psrho(y_0)\varphi(x)=[x,y_0]\varphi(x)$ ($y_0\in Y$) and $\psrho(x_0+y_0)=\beta(x_0,y_0)^{-1}
\psrho(x_0)\psrho(y_0)$. This is a Heisenberg representation, with $\psHS^\infty=\auS(X)$. (That
$\psrho(z)\psrho(z')=\beta(z,z')\psrho(z+z')$ is the computation $\psrho(y)\psrho(x')=[y,x']
\psrho(x')\psrho(y)$, $\beta(x+x',y+y')\beta(x,y)^{-1}\beta(x',y')^{-1}[y,x']=\beta(x,y')\beta(y,x')
=\beta(z,z')$, using $\beta=1$ on $X\times X$ and $Y\times Y$; irreducibility and the smooth
vectors are the classical facts of the Schr\"odinger model.) The same holds over $\auA$ with
$\auS(X(\auA))$.
\item[(P3)] (\emph{Kudla's splitting for split spaces}, \cite[Theorem~3.1, p.~378; (1.11)--(1.13),
p.~370; (1.14), p.~371; \S5, pp.~399--400]{KudlaSplitting}.) Let $F\ne\C$ be a local field of characteristic $0$, $E_v$ a field,
$W$ split, $Y\subset W$ a Lagrangian $E_v$-subspace and $X$ a complementary Lagrangian. Fix a
character $\xi$ of $E_v^\times$ with $\xi|_{F^\times}=\epsilon_{E_v/F}$. Kudla defines unitary
operators $\omega_{Y,\xi}(g)=\beta_{Y}(g)\,r_Y(g)$ on $\auS(X)$, where $r_Y(\sigma)\varphi(x)=
\int_{Y_\sigma}f_\sigma(x+y)\varphi(x\alpha+y\gamma)\,d\mu_\sigma(y)$ is Rao's operator for
$\sigma=\bigl(\begin{smallmatrix}\alpha&\beta\\\gamma&\delta\end{smallmatrix}\bigr)$ (with a
positive Haar measure $\mu_\sigma$ on $Y_\sigma=Y/\ker\gamma$, normalised so that $r_Y(\sigma)$ is
unitary), and $\beta_Y(g)=\xi(x(g))\gamma_F(\eta\circ RV)^{-j(g)}$ (Kudla's $\beta_V(g)$, case $3_+$,
with $\eta(t)=\psi_K(\frac12t)$ for Kudla's additive character $\psi_K$ of $F$); here $g\in P_Y\tau_jP_Y$ (Bruhat
decomposition (1.11)) and, for $g=p_1\tau_Sp_2$, $x(g)$ is the reduced norm, here the determinant, of
$p_1p_2|_Y$ (1.13)--(1.14). The map
$g\mapsto\omega_{Y,\xi}(g)$ is a representation of $\auU(W)$ on $\auS(X)$ \cite[p.~400]{KudlaSplitting}.
Kudla's groups act on the right (\cite[p.~362]{KudlaSplitting}: rows times matrices), and Rao's
operators normalise the Heisenberg operators $\psrho(x_0)\varphi(x)=\varphi(x+x_0)$,
$\psrho(y_0)\varphi(x)=\psi_K(\langle\!\langle x,y_0\rangle\!\rangle)\varphi(x)$ by
$\omega_{Y,\xi}(g)\psrho(w)\omega_{Y,\xi}(g)^{-1}=\psrho(w\cdot g^{-1})$ (for the generators of the
Siegel parabolic this is the direct computation from the formula for $r_Y$; with the opposite sign in
$\psrho(y_0)$ the unipotent and Levi elements would conjugate in opposite directions). For $W$ with
the form $\jmath\langle\,,\rangle$ and the hermitian line $V=E$, $(v_1,v_2)=\aubar v_1v_2$, Kudla's
symplectic form is $\langle\!\langle\,,\rangle\!\rangle=\frac12\pwfA$ \cite[(0.8), (0.9), p.~364]{KudlaSplitting},
so these Heisenberg operators are those of (P2) exactly when $\psi_K(t)=\psi_k(2t)$; this is the
character with which (P3) is read.
\item[(P4)] (\emph{Kudla's splitting for non-split spaces}, \cite[(4.1)--(4.5), p.~389, and
Proposition~4.1, p.~390]{KudlaSplitting}.) For $W$ not split, $W\oplus W^-$ is split ($W^-$ is $W$ with
the form negated), $\auU(W)$ embeds by $g\mapsto(g,1)$, and the splitting of $\auU(W)$ is the
restriction of the splitting (P3) of $\auU(W\oplus W^-)$, read in the decomposition
$\mathrm{Mp}(W)\times\mathrm{Mp}(W^-)\to\mathrm{Mp}(W\oplus W^-)$ attached to a polarisation
$Y'=Y'_1+Y'_2$ with $Y'_1\subset W$, $Y'_2\subset W^-$, along which the pull-back of Rao's cocycle
of $W\oplus W^-$ is the product of those of $W$ and $W^-$ \cite[p.~389]{KudlaSplitting}.
\item[(P5)] (\emph{Characters and Fourier analysis}.) $\psi$ is trivial on $E$ and non-trivial on
$\auA_E$; $\upsilon$ and $|\cdot|_{\auA_E}$ are trivial on $E^\times$; for a $k$-vector space $X$
the pairing $\psi_k(\pwfA(\cdot,\cdot))$ identifies $Y(k)$ with the character group of
$X(\auA)/X(k)$ when $X\oplus Y$ is a Lagrangian decomposition (self-duality of $\auA_k/k$); Fourier
series of smooth functions on a real torus converge in the $C^\infty$ topology; a distribution $u$
on $\R^N$ with $x_iu=0$ for all $i$ is a multiple of $\delta_0$.
\end{itemize}

\emph{Left actions.} Kudla's groups act on the right (P3); the paper's act on the left. We pass from
one to the other by letting $g$ act on the left by $g\cdot w:=w\cdot g^{-1}$, which is the unique
way to keep the group law and the identity $\omega(g)\psrho(w)\omega(g)^{-1}=\psrho(g\cdot w)$ of
(W1). Under this dictionary, for $p$ in the stabiliser $P_Y$ of $Y$, the endomorphism $p_1p_2|_Y$
of (1.13) with $p_1=p$, $p_2=1$ is the inverse of the left action of $p$ on $Y$. Hence, in the
paper's notation,
\begin{equation}\label{pwf:eq-xP}
\beta_Y(p)=\xi\bigl(\det(p|_Y)\bigr)^{-1}\qquad(p\in P_Y),
\end{equation}
since $j(p)=0$. (The value of $\xi$ on $\det(p|_Y)$ only matters modulo norms, and $\xi(\aubar
z^{-1})=\xi(z)\xi(z\aubar z)^{-1}=\xi(z)$; so the same value is obtained if $\det(p|_Y)$ is replaced
by the determinant of $p$ on a complement of $Y$ dual to $Y$.)

\begin{definition}[The Weil representation]\label{pwf:def}
Let $W$ be a non-degenerate hermitian space over $E$. For each place $v$ of $k$ a homomorphism
$\omega_v=\omega_{W,v}$ from $\auU(W)(k_v)$ to the unitary operators of a Heisenberg representation
$\psHS_v(W)$, with $\omega_v(g)\psrho(w)\omega_v(g)^{-1}=\psrho(gw)$, is fixed as follows.
\begin{itemize}
\item[(D1)] $E_v$ a field and $W_v$ split: $\omega_v=\omega_{Y,\upsilon_v}$ of (P3) for a Lagrangian
$Y\subset W_v$ (by Lemma~\ref{pwf:indep} it does not depend on $Y$).
\item[(D2)] $E_v$ a field and $W_v$ not split: the splitting (P4) with $\xi=\upsilon_v$, characterised
intrinsically as follows. The space $W_v\oplus W_v^-$ is split, and $\omega_{W\oplus W^-,v}(g,1)$ is
defined by (D1), independently of the Lagrangian (Lemma~\ref{pwf:indep}). Since $(g,1)$ fixes $W^-$
pointwise, this operator commutes with $1\otimes\psrho(w')$ for $w'\in W^-_v$ on
$\psHS_v(W\oplus W^-)=\psHS_v(W)\widehat\otimes\psHS_v(W^-)$; as $\psrho|_{W^-}$ is irreducible (P1),
the commutant of $1\otimes\psrho(W^-_v)$ is the set of operators $A\otimes1$, so
$\omega_{W\oplus W^-,v}(g,1)=A_g\otimes1$ for a unique bounded $A_g$, and we set $\omega_v(g):=A_g$.
As $A\mapsto A\otimes1$ is an injective, multiplicative, $*$-preserving map, $g\mapsto A_g$ is a
unitary representation, and $A_g\psrho(w)A_g^{-1}=\psrho(gw)$ follows from the same identity for
$(g,1)$ on $w\in W_v$. This is the operator of (P4) \cite[(4.1)--(4.5), Proposition~4.1]{KudlaSplitting},
and it involves no auxiliary choice.
\item[(D3)] $v=w\aubar w$ split in $E$: $W_v=W_w\oplus W_{\aubar w}$ with both summands
$\pwfA$-isotropic, $\auU(W)(k_v)=\auGL(W_w)$ (restriction to $W_w$), and on the model (P2)
$\auS(W_w)$,
\[
(\omega_v(g)\Psi)(x)=\upsilon_w(\det g_w)^{-1}|\det g_w|_w^{-1/2}\,\Psi(g_w^{-1}x).
\]
\end{itemize}
The global representation is $\omega=\bigotimes_v\omega_v$ on $\psHS(W)=\bigotimes'_v\psHS_v(W)$, as
in \S\ref{ps:ss-weil} and (W4).
\end{definition}

In (D3) the formula defines a representation with the required property: $g_w$ preserves $W_w$
and $W_{\aubar w}$ (as $g$ is $E_v$-linear), $\Psi\mapsto|\det g_w|^{-1/2}\Psi(g_w^{-1}\cdot)$ is
unitary, conjugates translation by $x_0$ into translation by $g_wx_0$, and conjugates multiplication
by $[\cdot,y_0]$ into multiplication by $[g_w^{-1}\cdot,y_0]=[\cdot,gy_0]$ ($g$ is an isometry);
the factor $\upsilon_w(\det g_w)^{-1}$ is a character. The cited works treat only the case of a
field $E_v$; (D3) is the normalisation used in the paper at split places (it is the one for which
\eqref{pl:eq-w5} holds, Proposition~\ref{pwf:W5}, and the one compatible with (D1), (D2) on the
rational points, Lemma~\ref{pwf:siegel}). In (D1), (D2) the identity $\omega_v(g)\psrho(w)
\omega_v(g)^{-1}=\psrho(gw)$ is (P3) read with left actions.

\begin{lemma}[Naturality; independence of the Lagrangian]\label{pwf:indep}
In case (D1), let $\sigma\colon W\to W'$ be an isometry of split spaces and $T\colon\psHS_v(W)\to
\psHS_v(W')$ an intertwiner along $\sigma$. Then $T\omega_{Y,\xi}(g)T^{-1}=\omega_{\sigma Y,\xi}
(\sigma g\sigma^{-1})$. Moreover $\omega_{Y,\xi}=\omega_{Y',\xi}$ for any two Lagrangians $Y,Y'$.
\end{lemma}
\begin{proof}
All ingredients of (P3) ($\psi$, $\xi$, the Bruhat cells and the invariants $x$, $j$, Rao's
operators and their measures, the Weil index of the form $RV$) are defined
(here $x(g)$ is the reduced norm of the $E_v$-linear endomorphism $p_1p_2|_Y$, which $\sigma$
preserves; the Weil index $\gamma_F(\eta\circ RV)$ depends only on $V$, the same on both sides; and
Rao's measures are fixed by unitarity of his operators, so they are carried over as well) from $(W,Y)$ by
formulas that are carried to those for $(W',\sigma Y)$ by $\sigma$ and by the transport of functions
on $X$ to functions on $\sigma X$; the transport of functions is an intertwiner along $\sigma$, and
by (P1) any other intertwiner $T$ differs from it by a scalar, which cancels in $T(\cdot)T^{-1}$.
This is the first claim. For the second, $Y'=\alpha Y$ for some $\alpha\in\auU(W)$ (Witt's theorem).
Apply the first claim with $W'=W$, $\sigma=\alpha$ and $T=\omega_{Y,\xi}(\alpha)$, which is an
intertwiner along $\alpha$: $\omega_{Y',\xi}(\alpha g\alpha^{-1})=\omega_{Y,\xi}(\alpha)\omega_{Y,\xi}(g)
\omega_{Y,\xi}(\alpha)^{-1}=\omega_{Y,\xi}(\alpha g\alpha^{-1})$ for all $g$.
\end{proof}

\subsubsection{Local lemmas: invariant functionals, the Siegel parabolic, orthogonal sums}\label{pwf:ss-local}

\begin{lemma}[Distributions supported at $0$]\label{pwf:dist}
Let $v$ be a place, $W_v=X\oplus Y\oplus W'$ with $X,Y$ $\pwfA$-isotropic, $X\oplus Y$ orthogonal
to $W'$ for $\pwfA$, and $(\psHS',\psrho')$ a Heisenberg representation of $W'$ (or $\psHS'=\C$ if
$W'=0$). Realise $\psHS_v(W)$ as $\auS(X)\widehat\otimes\psHS'$, with smooth vectors
$\auS(X;\psHS'^\infty)$, where $X\oplus Y$ acts by \textup{(P2)} on the first factor and $W'$ by
$\psrho'$ on the second. If $T\colon\auS(X;\psHS'^\infty)\to\psHS'^\infty$ is continuous linear,
$T\circ\psrho(y)=T$ for $y\in Y$ and $T\circ\psrho(w')=\psrho'(w')\circ T$ for $w'\in W'$, then
$T=c\,\pwfev$ for a constant $c$, where $\pwfev F=F(0)$.
\end{lemma}
\begin{proof}
(This is a realisation of $\psHS_v(W)$: $\beta(z,w')=1$ for $z\in X\oplus Y$, $w'\in W'$, so the
operators commute and $\psrho(z+w')=\psrho(z)\psrho(w')$; irreducibility is (P1).) Let $F\in
\auS(X;\psHS'^\infty)$ with $F(0)=0$; we show $TF=0$. If $v$ is finite, $F$ is locally constant
with compact support, so $F=F\cdot(1-\mathbf 1_U)$ for an open compact subgroup $U\subset X$; let
$M\subset Y$ be the open compact subgroup $\{y:[U',y]=1\}$ for $U'\subset U$ small enough, with
$\int_M[x,y]\,dy=\mathbf 1_{M^\perp}(x)$ ($dy$ of mass $1$) and $M^\perp\subset U$. Then
$F=F-\mathbf 1_{M^\perp}F=\int_M(F-\psrho(y)F)\,dy$, a finite average (the integrand is locally
constant in $y$), so $TF=\int_M(TF-T\psrho(y)F)dy=0$. If $v$ is real or complex, differentiate
$T(\psrho(ty)F)=T(F)$ at $t=0$: $T(\pwfA(x,y)F)=0$ for all $y$, i.e.\ $T(x_iF)=0$ for linear
coordinates $x_i$ on $X$. For a continuous functional $\lambda$ on $\psHS'^\infty$ and $h\in
\psHS'^\infty$ the distribution $u(\varphi)=\lambda(T(\varphi h))$ on $X$ satisfies $x_iu=0$, so
$u=c(\lambda,h)\delta_0$ (P5); hence $T(\varphi h)=\varphi(0)T(\phi_0h)$ for any $\phi_0$ with
$\phi_0(0)=1$, and by density of $\auS(X)\otimes\psHS'^\infty$ and continuity, $TF=T_0(F(0))$ with
$T_0(h)=T(\phi_0h)$. In both cases $T=T_0\circ\pwfev$ with $T_0$ continuous on $\psHS'^\infty$, and
$T_0\psrho'(w')=\psrho'(w')T_0$; by (P1), $T_0$ is a scalar.
\end{proof}

In the rest of this paragraph $W$ is split over $E_v$, $W=I\oplus I^*$ with $I,I^*$ Lagrangian
$E_v$-subspaces, and $\psHS_v(W)=\auS(I^*)$ by (P2) ($X=I^*$, $Y=I$). For $a\in\auGL(I)$ let $m(a)$
act by $a$ on $I$ and by the inverse adjoint $(a^*)^{-1}$ on $I^*$ (adjoint for the perfect pairing
$\langle\,,\rangle\colon I\times I^*\to E_v$), and for $b\in\auHom_{E_v}(I^*,I)$ with
$\langle bx,y\rangle+\langle x,by\rangle=0$ let $n(b)$ be the identity on $I$ and $x\mapsto x+bx$ on
$I^*$. These are elements of $\auU(W)$, $n(b)n(b')=n(b+b')$, and $m(a)n(b)m(a)^{-1}=n(m(a)\,b\,
m(a)^{-1})$. Put
\[
(\pwfL_a\varphi)(x)=|\det a|_{E_v}^{1/2}\varphi(m(a)^{-1}x),\qquad(\pwfM_b\varphi)(x)=\beta(x,bx)^{-1}
\varphi(x),
\]
where $\det a=\det_{E_v}(a|_I)$ and $|\cdot|_{E_v}$ is the module of $E_v$.

\begin{lemma}[The Siegel parabolic]\label{pwf:siegel}
At every place $v$, $\omega_v(m(a))=\upsilon_v(\det a)^{-1}\pwfL_a$ and $\omega_v(n(b))=\pwfM_b$.
\end{lemma}
\begin{proof}
\emph{The operators.} $m(a)$ acts on $I^*$ with module $|\det a|_{E_v}^{-1}$ (its
$E_v$-determinant in dual bases is $\aubar{\det a}^{-1}$), so $\pwfL_a$ is unitary; it conjugates
translation by $x_0$ into translation by $m(a)x_0$ and multiplication by $[\cdot,y_0]$ into
multiplication by $[m(a)^{-1}\cdot,y_0]=[\cdot,m(a)y_0]$, so it is an intertwiner along $m(a)$. For
$\pwfM_b$: the condition on $b$ gives $\pwfA(x,by)=\pwfA(y,bx)$, hence $\beta(x,by)=\beta(y,bx)$; with
$q(x)=\beta(x,bx)$ one gets $q(x+x_0)=q(x)q(x_0)[x,bx_0]$ and $[x_0,bx_0]=q(x_0)^2$. Therefore
$\pwfM_b\psrho(x_0)\pwfM_b^{-1}\varphi(x)=q(x_0)[x,bx_0]\varphi(x+x_0)=\beta(x_0,bx_0)^{-1}[x+x_0,bx_0]
\varphi(x+x_0)=\psrho(x_0+bx_0)\varphi(x)$, while $\pwfM_b$ commutes with $\psrho(y_0)$, $y_0\in I$;
as $n(b)x_0=x_0+bx_0$ and $n(b)y_0=y_0$, $\pwfM_b$ is an intertwiner along $n(b)$. Moreover
$\pwfM_b\pwfM_{b'}=\pwfM_{b+b'}$ and $\pwfL_a\pwfM_b\pwfL_a^{-1}=\pwfM_{m(a)bm(a)^{-1}}$ (as
$\beta(m^{-1}x,bm^{-1}x)=\beta(x,mbm^{-1}x)$).

\emph{The unipotent elements.} By (P1), $\omega_v(n(b))=c(b)\pwfM_b$ and $\omega_v(m(a))=\kappa(a)
\pwfL_a$ with scalars of modulus $1$. Since $\omega_v$ is a homomorphism, $c$ is a character of the
vector group of the $b$'s, and conjugating by $\omega_v(m(a))$ gives $c(m(a)bm(a)^{-1})=c(b)$. For
$a=2\cdot\mathrm{id}_I$, $m(a)bm(a)^{-1}=4b$; so $c(3b)=1$ for all $b$, and $c=1$.

\emph{The Levi elements.} $\kappa$ is a character and it remains to find it. Case (D1): $W_v$ is
split, $\omega_v=\omega_{I,\upsilon_v}$ (Lemma~\ref{pwf:indep}). For $\sigma=m(a)$ the lower left
block $\gamma$ vanishes, $Y_\sigma=0$, $f_\sigma=1$ (the upper right block is $0$), so $r_I(m(a))$ is a
positive multiple of $\varphi\mapsto\varphi(m(a)^{-1}\cdot)$ (read with left actions), i.e.\ of
$\pwfL_a$; and $\beta_I(m(a))=\upsilon_v(\det a)^{-1}$ by \eqref{pwf:eq-xP}. So $\kappa(a)=
\upsilon_v(\det a)^{-1}$. Case (D2) does not occur ($W_v$ is split). Case (D3): $v=w\aubar w$. The
idempotents $e_w,e_{\aubar w}$ of $E_v$ satisfy $\aubar{e_w}=e_{\aubar w}$, so
$\langle e_wx,e_wy\rangle=e_{\aubar w}e_w\langle x,y\rangle=0$: $\langle\,,\rangle$ pairs $W_w$ with
$W_{\aubar w}$, and $I^*_w$ with $I_{\aubar w}$. Let $\Phi\colon\auS(W_w)\to\auS(I^*)$ be an
intertwiner (P1) between the models (P2) for $W_w\oplus W_{\aubar w}$ and $I^*\oplus I$, and
$\pwfL'_a\Psi=|\det(m(a)|_{W_w})|^{-1/2}\Psi(m(a)^{-1}\cdot)$. Both $\Phi\pwfL'_a\Phi^{-1}$ and
$\pwfL_a$ are intertwiners along $m(a)$, so $\Phi\pwfL'_a\Phi^{-1}=\kappa'(a)\pwfL_a$ with
$|\kappa'(a)|=1$. To see $\kappa'(a)=1$, consider the functionals $\lambda_1(\Psi)=\int_{I_w}\Psi$ on
$\auS(W_w)$ and $\lambda_2=\pwfev$ on $\auS(I^*)$. Both are invariant under $\psrho(z)$, $z\in I$:
for $z\in I_w\subset W_w$, $\psrho(z)$ is a translation along $I_w$; for $z\in I_{\aubar w}$ it is
multiplication by $[x,z]$, which is $1$ on $I_w$ since $I$ is isotropic; for $\lambda_2$ by (P2). By
Lemma~\ref{pwf:dist} (with $W'=0$, in the model $\auS(I^*)$), $\lambda_1\circ\Phi^{-1}=c\lambda_2$,
and $c\ne0$. Now $\lambda_1\circ\pwfL'_a$ and $\lambda_2\circ\pwfL_a$ are positive multiples of
$\lambda_1$, $\lambda_2$ (change of variables), and $\lambda_1\circ\pwfL'_a=c\lambda_2\circ\Phi\pwfL'_a=
c\,\kappa'(a)\lambda_2\circ\pwfL_a\circ\Phi$; evaluating on $\Psi$ with $\lambda_1(\Psi)\ne0$ gives
$\kappa'(a)>0$, so $\kappa'(a)=1$. By (D3), $\omega_v(m(a))$ is $\upsilon_w(\det(m(a)|_{W_w}))^{-1}
\pwfL'_a$ in the first model. On $W_w=I_w\oplus I^*_w$, $m(a)$ acts by $a_w$ and by the inverse
transpose of $a_{\aubar w}$, so $\det(m(a)|_{W_w})=\det a_w/\det a_{\aubar w}$. As $\upsilon$ is
trivial on $k_v^\times$ ($\epsilon_{E/k}$ is trivial at split places), $\upsilon_{\aubar w}=
\upsilon_w^{-1}$ on $k_v^\times$, and $\upsilon_w(\det a_w/\det a_{\aubar w})^{-1}=\upsilon_w(\det
a_w)^{-1}\upsilon_{\aubar w}(\det a_{\aubar w})^{-1}=\upsilon_v(\det a)^{-1}$. Transporting by
$\Phi$: $\kappa(a)=\upsilon_v(\det a)^{-1}$.
\end{proof}

\begin{lemma}[Orthogonal sums and doubling]\label{pwf:sum}
Let $W=W_1\oplus W_2$ be an orthogonal sum of non-degenerate hermitian spaces over $E_v$, and
identify $\psHS_v(W)=\psHS_v(W_1)\widehat\otimes\psHS_v(W_2)$ by an intertwiner (P1). Then
$\omega_{W,v}(g_1,g_2)=\omega_{W_1,v}(g_1)\otimes\omega_{W_2,v}(g_2)$ for $g_i\in\auU(W_i)(k_v)$. In
particular $\omega_{W\oplus W^-,v}(g,1)=\omega_{W,v}(g)\otimes1$.
\end{lemma}
\begin{proof}
The identity does not depend on the intertwiner chosen (P1). \emph{Case (D3).} $W_w=W_{1,w}\oplus
W_{2,w}$, $\auS(W_w)=\auS(W_{1,w})\widehat\otimes\auS(W_{2,w})$ is an intertwiner,
$\det((g_1,g_2)_w)=\det g_{1,w}\det g_{2,w}$, and the formula of (D3) factors.

\emph{Case $E_v$ a field, $W_1,W_2$ split.} Choose Lagrangians $Y_i\subset W_i$ with complements
$X_i$; then $Y=Y_1\oplus Y_2$ is a Lagrangian of $W$, and by Lemma~\ref{pwf:indep} we may compute
$\omega_{W,v}$ with $Y$ on $\auS(X_1\oplus X_2)=\auS(X_1)\widehat\otimes\auS(X_2)$, which is an
intertwiner. If $g_i=p_i\tau_{S_i}p'_i$ are Bruhat decompositions in $\auU(W_i)$ (for bases adapted
to $W_i=X_i\oplus Y_i$), then $(g_1,g_2)=(p_1,p_2)\tau_{S_1\sqcup S_2}(p'_1,p'_2)$ is one in
$\auU(W)$; so $j(g_1,g_2)=j(g_1)+j(g_2)$ and $x(g_1,g_2)=x(g_1)x(g_2)$ by (1.14), and, the space $V$
and the Weil index being the same, $\beta_Y(g_1,g_2)=\beta_{Y_1}(g_1)\beta_{Y_2}(g_2)$. For
$\sigma=(\sigma_1,\sigma_2)$ block diagonal, $Y_\sigma=Y_{1,\sigma_1}\times Y_{2,\sigma_2}$ and
$f_\sigma=f_{\sigma_1}\otimes f_{\sigma_2}$, so $r_Y(\sigma)$ is a positive multiple of
$r_{Y_1}(\sigma_1)\otimes r_{Y_2}(\sigma_2)$ (the measures being positive); both are unitary, so the
multiple is $1$. Hence the identity.

\emph{Doubling, $E_v$ a field.} If $W$ is split, so is $W\oplus W^-$, and the case just
treated gives $\omega_{W\oplus W^-,v}(g,1)=\omega_{W,v}(g)\otimes1$ when $W$ and $W^-$ are split.
If $W$ is not split, this identity holds by the definition (D2): $\omega_{W\oplus W^-,v}(g,1)$,
which does not depend on the Lagrangian of the split space $W\oplus W^-$ (Lemma~\ref{pwf:indep}),
commutes with $1\otimes\psrho(W^-_v)$ and hence equals $A_g\otimes1$ with $A_g=\omega_{W,v}(g)$.
So in all cases
\begin{equation}\label{pwf:eq-double}
\omega_{W\oplus W^-,v}(g,1)=\omega_{W,v}(g)\otimes1\qquad(g\in\auU(W)(k_v)).
\end{equation}

\emph{General case, $E_v$ a field.} Let $D_i=W_i\oplus W_i^-$ (split), and $\sigma\colon
W\oplus W^-\to D_1\oplus D_2$ the isometry $(w_1,w_2,w'_1,w'_2)\mapsto(w_1,w'_1,w_2,w'_2)$; it
carries $((g_1,g_2),1)$ to $((g_1,1),(g_2,1))$. By \eqref{pwf:eq-double} for $W$, by
Lemma~\ref{pwf:indep} for $\sigma$, by the split case for $D_1\oplus D_2$, and by
\eqref{pwf:eq-double} for $W_1$, $W_2$:
\[
\omega_{W,v}(g_1,g_2)\otimes1=\omega_{W\oplus W^-,v}((g_1,g_2),1)\ \hat=\ \omega_{D_1,v}(g_1,1)\otimes
\omega_{D_2,v}(g_2,1)=\bigl(\omega_{W_1,v}(g_1)\otimes1\bigr)\otimes\bigl(\omega_{W_2,v}(g_2)\otimes1
\bigr),
\]
where $\hat=$ means equality after the rearrangement of tensor factors induced by $\sigma$ (an
intertwiner along $\sigma$). Comparing the factors on $\psHS_v(W)$ gives the claim; the last
assertion is \eqref{pwf:eq-double}.
\end{proof}

\subsubsection{Uniqueness of the theta functional}\label{pwf:ss-unique}

Globally, $\psHS(W)^\infty$ denotes the smooth vectors (at the infinite places for the action of
the Heisenberg group, at the finite places the algebraic restricted tensor product), with its
natural topology; in a Schr\"odinger model (P2) over $\auA$ it is $\auS(X(\auA))$. A \emph{theta
functional} is a non-zero continuous linear functional $\theta$ on $\psHS(W)^\infty$ with
$\theta\circ\psrho(w)=\theta$ for $w\in W(k)$; this is the functional of (W0), (W1).

\begin{lemma}[Uniqueness]\label{pwf:unique}
Let $W(k)=X(k)\oplus Y(k)$ with $X,Y$ $\pwfA$-isotropic $k$-subspaces (such a decomposition always
exists, $\pwfA$ being alternating). In the model $\auS(X(\auA))$ every continuous linear functional
$T$ with $T\circ\psrho(w)=T$ for $w\in W(k)$ is $T(\varphi)=c\sum_{\xi\in X(k)}\varphi(\xi)$. In
particular a theta functional exists and is unique up to a scalar, in every model.
\end{lemma}
\begin{proof}
The functional $\vartheta(\varphi)=\sum_{\xi\in X(k)}\varphi(\xi)$ is continuous and invariant:
translations by $X(k)$ permute the terms; multiplication by $[\cdot,y]$, $y\in Y(k)$, does not
change them, since $[\xi,y]=\psi_k(\pwfA(\xi,y))=1$ (P5); and $\beta(x_0,y_0)=1$ for rational $x_0,
y_0$. Conversely let $T$ be invariant. Let $C=X(\auA)/X(k)$, compact. Fix an open compact subgroup
$K_f\subset X(\auA_f)$ and $\Lambda=X(k)\cap K_f$, a lattice in $X_\infty$, with $X(\auA)=X(k)+(X_\infty
\times K_f)$; choose $u_\infty\in C_c^\infty(X_\infty)$ with $\sum_{\lambda\in\Lambda}u_\infty(x+
\lambda)=1$ and put $u=u_\infty\otimes\mathbf 1_{K_f}$. Then $\sum_{\xi\in X(k)}u(x+\xi)=1$ for all
$x\in X(\auA)$.

(i) \emph{Multiplication.} For a $K_f$-invariant smooth function $f$ on $C$ (a smooth function on the
real torus $X_\infty/\Lambda$, the $K_f$-invariant part of $C$) let $\tilde f$ be its lift to
$X(\auA)$. The characters of $C$ are the $[\cdot,y]$, $y\in Y(k)$ (P5), and $T([\cdot,y]\varphi)=
T(\varphi)$. Hence $T(\tilde f\varphi)=f(0)T(\varphi)$ for trigonometric polynomials $f$, and by
$C^\infty$ convergence of Fourier series (P5) and continuity of $f\mapsto\tilde f\varphi$ into
$\auS(X(\auA))$, for all such $f$.

(ii) \emph{Periodisation.} For $\varphi\in\auS(X(\auA))$, $\varphi=\sum_{\xi}\varphi\cdot u(\cdot-\xi)$
with convergence in $\auS(X(\auA))$ (only finitely many classes of $\xi$ modulo $\Lambda$ meet the
compact support of the finite part, and the sum over $\Lambda$ converges in $\auS(X_\infty)$ because
$u_\infty$ has compact support). By invariance under translation by $\xi$,
$T(\varphi\cdot u(\cdot-\xi))=T(\varphi(\cdot+\xi)u)$, so $T(\varphi)=T\bigl((P\varphi)u\bigr)$ with
$P\varphi(x)=\sum_\xi\varphi(x+\xi)$, a smooth function on $C$, $K_f'$-invariant for some open
$K'_f$ (apply (i) with $K'_f\subset K_f$). By (i), $T(\varphi)=P\varphi(0)\,T(u)=T(u)\,\vartheta(\varphi)$.
The last assertion follows because any two Heisenberg representations are intertwined (P1), and an
intertwiner is a topological isomorphism of the smooth vectors.
\end{proof}

\subsubsection{The theta functional is invariant under the rational points}\label{pwf:ss-M4}

\begin{lemma}[Generation]\label{pwf:gen}
Let $W=I\oplus I^*$ be split over $E$ with $I,I^*$ Lagrangian. Then $\auU(W)(k)$ is generated by the
$m(a)$, $n(b)$ of \S\ref{pwf:ss-local} and the elements $n^-(c)$ that are the identity on $I^*$ and
map $y\in I$ to $y+cy$, $c\in\auHom(I,I^*)$ with $\langle cy,y'\rangle+\langle y,cy'\rangle=0$. With
$\varepsilon\in E^\times$, $\aubar\varepsilon=-\varepsilon$, the element $w_\varepsilon$ that maps
$y\in I$ to $-\varepsilon^{-1}\iota(y)$ and $x\in I^*$ to $\varepsilon\iota^{-1}(x)$, where
$\iota\colon I\to I^*$ is the isomorphism attached to dual bases, lies in $\auU(W)(k)$, satisfies
$w_\varepsilon=n(b)n^-(c)n(b)$ for suitable $b,c$, and $w_\varepsilon n(b')w_\varepsilon^{-1}$ runs
over all $n^-(c)$ as $b'$ runs over the admissible $b'$.
\end{lemma}
\begin{proof}
Choose dual bases of $I$, $I^*$; the Gram matrix is $J=\bigl(\begin{smallmatrix}0&1\\1&0
\end{smallmatrix}\bigr)$ and $g=\bigl(\begin{smallmatrix}a&b\\c&d\end{smallmatrix}\bigr)\in\auU(W)$
iff $g^*Jg=J$, equivalently $gJg^*=J$; in particular $cd^*+dc^*=0$. Then $m(a)=\mathrm{diag}(a,
a^{*-1})$, $n(b)=\bigl(\begin{smallmatrix}1&b\\0&1\end{smallmatrix}\bigr)$ and $n^-(c)=
\bigl(\begin{smallmatrix}1&0\\c&1\end{smallmatrix}\bigr)$ with $b,c$ skew-hermitian. The matrix
identities
\[
\begin{pmatrix}1&\varepsilon\\0&1\end{pmatrix}\begin{pmatrix}1&0\\-\varepsilon^{-1}&1\end{pmatrix}
\begin{pmatrix}1&\varepsilon\\0&1\end{pmatrix}=\begin{pmatrix}0&\varepsilon\\-\varepsilon^{-1}&0
\end{pmatrix}=w_\varepsilon,\qquad w_\varepsilon n(b')w_\varepsilon^{-1}=n^-(-\varepsilon^{-2}b')
\]
(scalars times the identity of the blocks; $\pm\varepsilon^{\pm1}$ are skew-hermitian and
$\varepsilon^2\in k$) prove the assertions on $w_\varepsilon$.

Let $g\in\auU(W)(k)$. \emph{Step 1: there is a skew-hermitian $y$ with $d+cy$ invertible.} Replacing
$g$ by $m(\alpha)gm(\beta)$ replaces $(c,d)$ by $(\alpha^{*-1}c\beta,\alpha^{*-1}d\beta^{*-1})$ and
$d+cy$ by $\alpha^{*-1}(d+c\,\beta y\beta^*)\beta^{*-1}$, with $\beta y\beta^*$ skew-hermitian; so we may
assume $c=\mathrm{diag}(1_s,0)$. Writing $d$ in blocks, $cd^*+dc^*=0$ gives $d_{21}=0$ and
$d_{11}+d_{11}^*=0$; the rows of $(c\ d)$ are independent, so $d_{22}$ is invertible. Take
$y=\mathrm{diag}(\varepsilon1_s-d_{11},0)$, skew-hermitian; then $d+cy=\bigl(\begin{smallmatrix}
\varepsilon1_s&d_{12}\\0&d_{22}\end{smallmatrix}\bigr)$ is invertible.
\emph{Step 2.} Put $g'=gn(y)$, whose lower right block $d'=d+cy$ is invertible. From $c'd'^*+d'c'^*
=0$, $d'^{-1}c'$ is skew-hermitian, so $n^-(-d'^{-1}c')\in\auU(W)(k)$, and $g'n^-(-d'^{-1}c')=
\bigl(\begin{smallmatrix}*&*\\0&d'\end{smallmatrix}\bigr)$ lies in the stabiliser $P$ of $I$. An
element $p=\bigl(\begin{smallmatrix}a&b\\0&d\end{smallmatrix}\bigr)\in P$ has $d=a^{*-1}$ and
$p\,m(a)^{-1}=\bigl(\begin{smallmatrix}1&ba^*\\0&1\end{smallmatrix}\bigr)$, which is unitary, hence an
$n(b')$. So $g\in\langle m,n,n^-\rangle$.
\end{proof}

\begin{theorem}[Invariance of $\theta$; (W1)]\label{pwf:inv}
Let $W$ be a non-degenerate hermitian space over $E$, $\omega$ as in Definition~\ref{pwf:def} and
$\theta$ a theta functional. Then $\theta\circ\omega(\gamma)=\theta$ on $\psHS(W)^\infty$ for every
$\gamma\in\auU(W)(k)$.
\end{theorem}
\begin{proof}
For $\gamma\in\auU(W)(k)$, $\theta\circ\omega(\gamma)$ is again a theta functional:
$\theta(\omega(\gamma)\psrho(w)F)=\theta(\psrho(\gamma w)\omega(\gamma)F)=\theta(\omega(\gamma)F)$ for
$w\in W(k)$, and $\omega(\gamma)$ is continuous on $\psHS(W)^\infty$ (P1). By Lemma~\ref{pwf:unique},
\begin{equation}\label{pwf:eq-char}
\theta\circ\omega(\gamma)=\kappa(\gamma)\,\theta,\qquad\kappa(\gamma)\in\C^\times,
\end{equation}
and $\kappa$ is a character of $\auU(W)(k)$ (as $\omega$ is a homomorphism and $\theta\ne0$).

\emph{Step 1: $W$ split.} Write $W=I\oplus I^*$ with rational Lagrangians and use the model
$\auS(I^*(\auA))=\bigotimes'_v\auS(I^*_v)$ (P2), in which $\theta=c\sum_{\xi\in I^*(k)}\varphi(\xi)$
by Lemma~\ref{pwf:unique} ($X=I^*$, $Y=I$). For $a\in\auGL(I)(E)$, by Lemma~\ref{pwf:siegel} at all
places, $\omega(m(a))\varphi=\prod_v\upsilon_v(\det a)^{-1}|\det a|_{E_v}^{1/2}\,\varphi(m(a)^{-1}\cdot)
=\varphi(m(a)^{-1}\cdot)$ by (P5) (almost all factors are $1$), and $m(a)^{-1}$ permutes $I^*(k)$:
$\kappa(m(a))=1$. For rational $b$, $\omega(n(b))$ is multiplication by $\beta(x,bx)^{-1}$, which is
$1$ at $x\in I^*(k)$ (P5): $\kappa(n(b))=1$. For $n^-(c)$ write $n^-(c)=w_\varepsilon n(b')
w_\varepsilon^{-1}$ (Lemma~\ref{pwf:gen}); $\kappa$ being a character, $\kappa(n^-(c))=\kappa(n(b'))=1$.
By Lemma~\ref{pwf:gen}, $\kappa=1$.

\emph{Step 2: general $W$.} The space $W\oplus W^-$ is split: $I_\Delta=\{(x,x)\}$ and
$I_\nabla=\{(x,-x)\}$ are complementary Lagrangians. Choose rational $\pwfA$-isotropic
decompositions $W=X\oplus Y$, $W^-=X^-\oplus Y^-$ (Lemma~\ref{pwf:unique}); then $(X\oplus X^-)\oplus
(Y\oplus Y^-)$ is one of $W\oplus W^-$, and in the Schr\"odinger models
$\auS((X\oplus X^-)(\auA))\supset\auS(X(\auA))\otimes\auS(X^-(\auA))$, which realise
$\psHS(W\oplus W^-)=\psHS(W)\widehat\otimes\psHS(W^-)$. Let $\theta^-$, $\Theta$ be theta functionals
for $W^-$ and $W\oplus W^-$. By Lemma~\ref{pwf:unique} in these models, $\Theta(F\otimes F')=
c\,\theta(F)\theta^-(F')$ for $F,F'$ smooth, with $c\ne0$. By \eqref{pwf:eq-double} at every place,
$\omega_{W\oplus W^-}(\gamma,1)=\omega_W(\gamma)\otimes1$. By Step 1 applied to $W\oplus W^-$ and
$(\gamma,1)\in\auU(W\oplus W^-)(k)$:
\[
c\,\theta(\omega_W(\gamma)F)\theta^-(F')=\Theta\bigl(\omega_{W\oplus W^-}(\gamma,1)(F\otimes F')\bigr)
=\Theta(F\otimes F')=c\,\theta(F)\theta^-(F').
\]
Choosing $F'$ with $\theta^-(F')\ne0$ gives $\theta\circ\omega_W(\gamma)=\theta$.
\end{proof}

\begin{remark}\label{pwf:rem-inv}
The proof uses neither Weil's theorem on $\Theta$ nor the product formula for Weil indices: only
the values of the local splittings on the Levi subgroup of a Siegel parabolic
(Lemma~\ref{pwf:siegel}), whose product over the places is the Hecke character
$\upsilon(\det)^{-1}|\det|^{1/2}$, trivial on $\auGL(I)(E)$. Together with Lemma~\ref{pwf:unique}
it proves all of (W1): the commutation relation is part of Definition~\ref{pwf:def}, and existence,
uniqueness and invariance of $\theta$ are Lemma~\ref{pwf:unique} and Theorem~\ref{pwf:inv}.
\end{remark}

\begin{corollary}[Agreement with Weil's lift of rational points]\label{pwf:rk}
Realise $\psHS(W)$ as $\auS(X(\auA))$ for a rational $\pwfA$-isotropic decomposition $W=X\oplus Y$,
and let $r_k\colon\mathrm{Sp}(W)(k)\to\mathrm{Mp}(W)_\auA$ be Weil's lift of the rational points
\cite[no.~40]{WeilActa} (defined there on $Ps(X)_k$, which is identified with $\mathrm{Sp}(W)(k)$ when the
characteristic is not $2$, \cite[p.~181]{WeilActa}). Then $\omega(\gamma)=r_k(\gamma)$ as operators on $\auS(X(\auA))$ for
$\gamma\in\auU(W)(k)$.
\end{corollary}
\begin{proof}
Both operators induce $\gamma$ on the Heisenberg group, so by (P1) $\omega(\gamma)=\lambda(\gamma)
r_k(\gamma)$ with $|\lambda(\gamma)|=1$. The functional $\vartheta(\varphi)=\sum_{\xi\in X(k)}\varphi(\xi)$
is invariant under $r_k(\gamma)$ by \cite[Th\'eor\`eme~6]{WeilActa} (evaluate the function $\Theta$
there at the identity, for $S=r_k(\gamma)$ and every $\varphi$), and under $\omega(\gamma)$ by
Theorem~\ref{pwf:inv} and Lemma~\ref{pwf:unique}. Since $\vartheta\ne0$, $\lambda(\gamma)=1$.
\end{proof}

In particular, Weil's Siegel formula, stated in~\cite[Th\'eor\`eme~5]{WeilSiegel} for the lift $r_k$, applies
verbatim to $\omega$ restricted to the rational points, as the paper uses it.

\subsubsection{The mixed model and the zero term: (W2), (W2$'$), (W2$_\infty$), (W3)}\label{pwf:ss-mixed}

Let $I\subset W$ be a non-zero isotropic $E$-subspace (over $k$, or over $k_v$ in the local
statements). Choose an isotropic $I^*$ with $\langle\,,\rangle\colon I\times I^*\to E$ perfect and
put $W_I=(I\oplus I^*)^\perp$; then $I^\perp=I\oplus W_I$, and $W_I\to I^\perp/I$ is an isometry, by
which we identify $W_I$ with $I^\perp/I$. Let $P_I$ be the stabiliser of $I$; $p\in P_I$ preserves
$I^\perp$ and induces $\bar p\in\auU(W_I)$. Let $M_I\subset P_I$ be the stabiliser of $I$ and $I^*$
(it preserves $W_I$, and $M_I=\auGL(I)\times\auU(W_I)$, $(a,h)$ acting by $m(a)$ on $I\oplus I^*$ and
by $h$ on $W_I$), and $N_I$ the subgroup of $P_I$ acting trivially on $I$, on $I^\perp/I$ and on
$W/I^\perp$.

\begin{lemma}[Levi decomposition]\label{pwf:levi}
$P_I=N_IM_I$.
\end{lemma}
\begin{proof}
Let $p\in P_I$ and $J=pI^*$: isotropic, of the dimension of $I^*$, with $J\cap I^\perp=p(I^*\cap
I^\perp)=0$. As $W=I^\perp\oplus I^*$, $J$ is the graph $\{x+\alpha x+\beta x\}$ of a map $(\alpha,
\beta)\colon I^*\to I\oplus W_I$, and isotropy of $J$ says $\langle x,\alpha y\rangle+\langle\alpha x,
y\rangle+\langle\beta x,\beta y\rangle=0$. Define $u$ by $u(i)=i$ on $I$, $u(x)=x+\alpha x+\beta x$ on
$I^*$, $u(w)=w+\gamma w$ on $W_I$, where $\gamma\colon W_I\to I$ is defined by $\langle\gamma w,x\rangle
=-\langle w,\beta x\rangle$. The isotropy relation and the definition of $\gamma$ give $\langle uz,uz'
\rangle=\langle z,z'\rangle$ on all pairs of basis vectors, so $u\in\auU(W)$, and $u\in N_I$. Then
$u^{-1}pI^*=I^*$, so $u^{-1}p\in M_I$.
\end{proof}

\emph{The mixed model.} By Lemma~\ref{pwf:sum} for $W=(I\oplus I^*)\oplus W_I$ and the model (P2)
$\auS(I^*)$ of $\psHS(I\oplus I^*)$, $\psHS(W)^\infty=\auS(I^*;\psHS(W_I)^\infty)$, locally and over
$\auA$, where $I\oplus I^*$ acts by (P2) on the variable and $W_I$ by $\psrho_{W_I}$ on the values,
and $\omega_W(m)=\omega_{I\oplus I^*}(m(a))\otimes\omega_{W_I}(h)$ for $m=(a,h)\in M_I$. Over $\auA$
the realisation is unique up to a scalar (P1); we fix it by requiring
\begin{equation}\label{pwf:eq-thetamixed}
\theta_W(F)=\sum_{\xi\in I^*(k)}\theta_{W_I}\bigl(F(\xi)\bigr)\qquad(F\in\auS(I^*(\auA);\psHS(W_I)^\infty)),
\end{equation}
for the theta functionals $\theta_W$, $\theta_{W_I}$ fixed in \S\ref{ps:ss-weil}. This is possible:
the right side converges absolutely ($\theta_{W_I}\circ F\in\auS(I^*(\auA))$ by continuity of
$\theta_{W_I}$), is continuous in $F$, non-zero, and invariant under $\psrho(w)$, $w\in W(k)$ (for
$w\in I(k)$ it multiplies $F(\xi)$ by $[\xi,w]=1$; for $w\in I^*(k)$ it translates the variable by
$w$; for $w\in W_I(k)$ it applies $\psrho_{W_I}(w)$ to the values, and $\theta_{W_I}$ is invariant;
the factors $\beta$ between rational points are $1$); so it is a multiple of $\theta_W$ by
Lemma~\ref{pwf:unique}, and the scalar is absorbed in the realisation.

\begin{definition}[Zero term]\label{pwf:defZ}
$\psZ_I\colon\psHS(W)^\infty\to\psHS(W_I)^\infty$ is $\psZ_IF=F(0)$ in the realisation just fixed;
locally, $\psZ_{I,v}F=F(0)$ in the local mixed model (defined up to a scalar).
\end{definition}

\begin{lemma}[Characterisation of the zero term]\label{pwf:Zunique}
Let $T\colon\psHS(W)^\infty\to\psHS(W_I)^\infty$ be continuous linear with $T\psrho(w)=
\psrho(\bar w)T$ for all $w\in I^\perp$ (over $\auA$, resp.\ over $k_v$), $\bar w$ the image in
$W_I=I^\perp/I$. Then $T=c\,\psZ_I$ (resp.\ $c\,\psZ_{I,v}$).
\end{lemma}
\begin{proof}
Locally this is Lemma~\ref{pwf:dist} with $X=I^*$, $Y=I$, $W'=W_I$: the hypothesis for $w\in I$ is
$T\psrho(w)=T$, and for $w\in W_I$ it is $T\psrho(w)=\psrho_{W_I}(w)T$. Over $\auA$: the local
argument in the proof of Lemma~\ref{pwf:dist} uses only the invariance under $\psrho(I_v)$, so for
each $v$ and each fixed vector $F^v$ away from $v$, $T(\varphi_v\otimes F^v)=\varphi_v(0)\,
T(\varphi_v^0\otimes F^v)$, where $\varphi^0_v$ is the standard vector (with $\varphi^0_v(0)=1$) at
the finite places outside a finite set and any fixed vector with value $1$ at $0$ elsewhere. As every
element of $\auS(I^*(\auA))\otimes\psHS(W_I)^\infty$ is a finite sum of tensors that are standard
outside a finite set of places, $TF=T_0(F(0))$ on these, with $T_0(h)=T(\otimes_v\varphi^0_v\otimes h)$,
and by continuity everywhere. $T_0$ commutes with $\psrho_{W_I}(W_I(\auA))$, so it is a scalar (P1).
\end{proof}

\begin{proposition}[(W2), (W2$'$), (W2$_\infty$)]\label{pwf:W2}
With $\theta_W,\theta_{W_I}$ as in \S\ref{ps:ss-weil} and $di$ the Haar measure of mass $1$ on
$[I]=I(k)\backslash I(\auA)$:
\begin{itemize}
\item[(a)] $\theta_{W_I}(\psZ_IF)=\int_{[I]}\theta_W(\psrho(i)F)\,di$; this property together with
\textup{(b)} characterises $\psZ_I$.
\item[(b)] $\psZ_I\psrho(w)=\psrho(\bar w)\psZ_I$ for $w\in I^\perp(\auA)$.
\item[(c)] $\psZ_I\,\omega_W(p)=\chi_I(p)\,\omega_{W_I}(\bar p)\,\psZ_I$ for $p\in P_I(\auA)$, with
$\chi_I(p)=\upsilon(\det p|_I)^{-1}|\det p|_I|^{1/2}$ ($\det p|_I\in\auA_E^\times$, $|\cdot|$ the
idele norm of $E$). Locally, $\psZ_{I,v}\omega_{W,v}(p)=\chi_{I,v}(p)\omega_{W_I,v}(\bar p)
\psZ_{I,v}$ with $\chi_{I,v}(p)=\upsilon_v(\det p|_I)^{-1}|\det p|_I|_{E_v}^{1/2}$; at a real place
and for a line $I$ on which $p$ acts by $a\in\C^\times$, $\chi_{I,\tau}(p)=\upsilon_\tau(a)^{-1}|a|$
\textup{(W2$_\infty$)}. In particular $\chi_I=1$ on $N_I(\auA)$ and on $P_I(k)$.
\item[(d)] For isotropic $I\subset I'$, $\psZ_{I'}=\psZ_{I'/I}\circ\psZ_I$, and $\psZ_I=\psZ_{I,\infty}
\otimes\psZ_{I,f}$ (for suitable normalisations of the local maps).
\end{itemize}
\end{proposition}
\begin{proof}
(a) By (P2), $(\psrho(i)F)(x)=[x,i]F(x)$ for $i\in I(\auA)$, so by \eqref{pwf:eq-thetamixed}
$\theta_W(\psrho(i)F)=\sum_\xi[\xi,i]\,\theta_{W_I}(F(\xi))$, with $\sum_\xi|\theta_{W_I}(F(\xi))|<\infty$
independent of $i$. For $\xi\in I^*(k)$, $i\mapsto[\xi,i]$ is a character of $[I]$, non-trivial if
$\xi\ne0$ (P5). Integrating term by term leaves $\theta_{W_I}(F(0))$. If $T$ satisfies (a) and (b),
then $T=c\psZ_I$ by Lemma~\ref{pwf:Zunique}, and (a) for both gives $c\,\theta_{W_I}\circ\psZ_I=
\theta_{W_I}\circ\psZ_I\ne0$, so $c=1$.

(b) For $w=i+w'$, $i\in I$, $w'\in W_I$: $\beta(i,w')=1$, $\psrho(w)=\psrho(i)\psrho_{W_I}(w')$, and
$(\psrho(w)F)(0)=[0,i]\psrho_{W_I}(w')F(0)=\psrho(\bar w)F(0)$.

(c) \emph{Local statement.} For $p\in P_I(k_v)$ the operator $T=\omega_{W_I}(\bar p)^{-1}\psZ_{I,v}
\omega_W(p)$ satisfies the hypothesis of Lemma~\ref{pwf:Zunique}: for $w\in I^\perp$, $T\psrho(w)=
\omega(\bar p)^{-1}\psZ_{I,v}\psrho(pw)\omega(p)=\omega(\bar p)^{-1}\psrho(\bar p\bar w)\psZ_{I,v}
\omega(p)=\psrho(\bar w)T$, using (b) locally and $\overline{pw}=\bar p\bar w$. So $\psZ_{I,v}\omega(p)=
\chi(p)\omega(\bar p)\psZ_{I,v}$, and $\chi$ is a continuous character of $P_I(k_v)$ ($\psZ_{I,v}$ is
onto). It suffices to compute $\chi$ on $M_I$ and $N_I$ (Lemma~\ref{pwf:levi}). For $m=(a,h)\in
M_I$, Lemma~\ref{pwf:sum} and Lemma~\ref{pwf:siegel} give $\omega_W(m)F=\upsilon_v(\det a)^{-1}
\pwfL_a\otimes\omega_{W_I}(h)F$, and $(\pwfL_a\varphi)(0)=|\det a|_{E_v}^{1/2}\varphi(0)$; as
$\bar m=h$ and $\det m|_I=\det a$, $\chi(m)=\chi_{I,v}(m)$. On $N_I$, $\chi$ is a character with
$\chi(qnq^{-1})=\chi(n)$ for $q\in P_I$. Write $n\in N_I$ as in the proof of
Lemma~\ref{pwf:levi}: $n(x)=x+\alpha x+\beta x$ on $I^*$; then $n\mapsto\beta\in\auHom(I^*,W_I)$ is a
homomorphism (as $(nn')(x)=x+(\alpha+\alpha'+\gamma\beta')x+(\beta+\beta')x$) whose kernel $Z$ is the
abelian group of the $n$ with $\beta=0$, parametrised additively by the skew maps $\alpha$. Let
$t=m(2\cdot\mathrm{id}_I)\in M_I$ ($2$ on $I$, $\frac12$ on $I^*$, $1$ on $W_I$). Then $tnt^{-1}$
has parameters $(4\alpha,2\beta)$ ($tnt^{-1}(x)=t\,n(2x)=x+4\alpha x+2\beta x$). On $Z$:
$\chi(\alpha)=\chi(4\alpha)=\chi(\alpha)^4$, so $\chi^3=1$ on the vector group $Z$, and $\chi(\alpha)=
\chi(\alpha/3)^3=1$. Hence $\chi$ factors through the abelian group $N_I/Z$, on which
$\chi(\beta)=\chi(2\beta)=\chi(\beta)^2$, so $\chi=1$. Since $\bar n=1$ and $\det n|_I=1$,
$\chi=\chi_{I,v}$ on $N_I$, hence on $P_I(k_v)$. At a real place $\tau$ of $k$ one has $E_\tau\cong\C$
with module $|a|_{E_\tau}=|a|^2$, so $|a|_{E_\tau}^{1/2}=|a|$, which is (W2$_\infty$).

\emph{Global statement.} By (P1) the global realisation is the tensor product of local ones up to a
scalar, so $\psZ_I=c\bigotimes_v\psZ_{I,v}$, and $\omega_W=\bigotimes_v\omega_{W,v}$,
$\omega_{W_I}=\bigotimes_v\omega_{W_I,v}$; so (c) holds with $\chi_I=\prod_v\chi_{I,v}$, which is the
stated formula (almost all factors are $1$ on a given $p$). It is $1$ on $N_I(\auA)$ by the local
case and on $P_I(k)$ by (P5).

(d) The second statement is the construction ($\psZ_I=c\bigotimes_v\psZ_{I,v}$). For the first,
$I'/I$ is an isotropic subspace of $W_I$ with $W_{I'}=(I'/I)^\perp/(I'/I)$. The map
$T=\psZ_{I'/I}\psZ_I$ satisfies the hypothesis of Lemma~\ref{pwf:Zunique} for $I'$: for $w\in
I'^\perp\subset I^\perp$, $\psZ_I\psrho(w)=\psrho(\bar w)\psZ_I$ with $\bar w\in(I'/I)^\perp$, and then
(b) for $I'/I$. So $T=c\psZ_{I'}$. By (a) three times and Fubini (choose a complement $J$ of $I$ in
$I'$; $\psrho(i)\psrho(j)=\psrho(i+j)$ as $\beta=1$ on the isotropic $I'$, and $[I']$ fibres over
$[J]=[I'/I]$ with fibre $[I]$, all measures of mass $1$),
\[
\theta_{W_{I'}}(TF)=\int_{[J]}\theta_{W_I}\bigl(\psZ_I\psrho(j)F\bigr)dj=\int_{[J]}\int_{[I]}
\theta_W(\psrho(i)\psrho(j)F)\,di\,dj=\int_{[I']}\theta_W(\psrho(i')F)\,di'=\theta_{W_{I'}}(\psZ_{I'}F),
\]
using (b) for $\psZ_I\psrho(j)=\psrho(\bar j)\psZ_I$. So $c=1$.
\end{proof}

\begin{proposition}[The identity of (W3)]\label{pwf:W3}
For $F\in\psHS(W)^\infty$, $\theta_W(F)=\sum_{w\in W(k)/I^\perp(k)}\theta_{W_I}(\psZ_I\psrho(w)F)$; the
terms do not depend on the representatives, and the series converges absolutely.
\end{proposition}
\begin{proof}
$I^*(k)$ is a set of representatives of $W(k)/I^\perp(k)$, and $\psZ_I\psrho(x)F=F(x)$ for $x\in
I^*(k)$ by (P2); so the series is \eqref{pwf:eq-thetamixed}, absolutely convergent. For $u\in
I^\perp(k)$, $\psrho(u)\psrho(w)=\beta(u,w)\psrho(u+w)$ with $\beta(u,w)=1$ for rational $u,w$ (P5); so $\theta_{W_I}(\psZ_I\psrho(w+u)F)=
\theta_{W_I}(\psZ_I\psrho(u)\psrho(w)F)=\theta_{W_I}(\psrho(\bar u)\psZ_I\psrho(w)F)=\theta_{W_I}(\psZ_I
\psrho(w)F)$, by (b) and the invariance of $\theta_{W_I}$ under $\psrho(\bar u)$, $\bar u\in W_I(k)$.
\end{proof}

\begin{remark}\label{pwf:rem-W3}
The moderate growth of $(g,h)\mapsto\sum_w|\theta_{W_I}(\psZ_I\psrho(w)\omega(g,h)F)|$ asserted in
(W3) is not proved in this subsection: it is Lemma~\ref{pz:absgrowth}(b) (for $F=\Psi\otimes f$, with
$A=4s$ and the constant $C(2D)^{4s}\|\Psi\|_{2s}$ there), whose proof uses (W0), (W1), (W2), (W4) but
not (W3); its part (a) gives again the absolute convergence of Proposition~\ref{pwf:W3}. Proposition~\ref{pwf:W2}(c) and Lemma~\ref{pwf:Zunique} also give the uniqueness
of the zero term used in \S\ref{pf:ss-flip} and the identity $\psZ_{L_0}\circ\omega(\gamma)=
\psZ_{\gamma^{-1}L_0}$ for rational $\gamma$ and a Lagrangian $L_0$: both sides
satisfy (a) and (b) for $\gamma^{-1}L_0$, by Theorem~\ref{pwf:inv} (this is the step in the proof of
Lemma~\ref{pc:clean}).
\end{remark}

\subsubsection{The model at split places: (W5) and \eqref{pl:eq-w5}}\label{pwf:ss-split}

For hermitian spaces $V,V'$ of ranks $n,n'$ and $W=\auHom(V,V')$, $\omega_{\auHom(V,V')}(g,h)$ is
$\omega_W$ of the image of $(g,h)$ under $\auU(V)\times\auU(V')\to\auU(W)$, $(g,h)w=hwg^{-1}$, and
$\omega_\mu$ is its determinant twist (\S\ref{ps:ss-weil}).

\begin{proposition}[(W5)]\label{pwf:W5}
Let $v=w\aubar w$ be split in $E$. Then $\psHS_v(\auHom(V,V'))$ is realised on
$\auS(\auHom(V_w,V'_w))$, with $\psrho(a_w,0)\Psi(x)=\Psi(x+a_w)$ and $\psrho(0,a_{\aubar w})\Psi(x)=
[x,a_{\aubar w}]\Psi(x)$, and
\[
(\omega_{\auHom(V,V')}(g,h)\Psi)(x)=\upsilon_w(\det h_w)^{-n}\upsilon_w(\det g_w)^{n'}\,
|\det h_w|_w^{-n/2}|\det g_w|_w^{n'/2}\,\Psi(h_w^{-1}xg_w).
\]
If $n,n'$ are even, then $(\omega_\mu(g,h)\Psi)(x)=\chi_w(\det h_w)\chi_w(\det g_w)^{-1}|\det
h_w|_w^{-n/2}|\det g_w|_w^{n'/2}\Psi(h_w^{-1}xg_w)$, where $\chi=\mu\circ(z\mapsto z/\aubar z)$; this
is \eqref{pl:eq-w5}.
\end{proposition}
\begin{proof}
$W_v=W_w\oplus W_{\aubar w}$ with $W_w=\auHom(V_w,V'_w)$, both summands isotropic
(\S\ref{pwf:ss-local}, proof of Lemma~\ref{pwf:siegel}); the realisation and the Heisenberg operators
are (P2) with $X=W_w$, $Y=W_{\aubar w}$. The element $(g,h)$ acts on $W_w$ by $x\mapsto h_wxg_w^{-1}$,
whose determinant is $\det(h_w)^n\det(g_w)^{-n'}$. So (D3) gives the first formula. For the second,
$\auU(1)(k_v)=\{(t,t^{-1})\}\cong k_v^\times$ (component at $w$), and $\upsilon(t,t^{-1})=\upsilon_w(t)
\upsilon_{\aubar w}(t)^{-1}=\upsilon_w(t)^2$ because $\upsilon_{\aubar w}=\upsilon_w^{-1}$ on $k_v^\times$
($\upsilon|_{\auA_k^\times}=\epsilon_{E/k}$ is trivial at $v$). Hence, for even $n$,
$\upsilon_n(\det h)=(\upsilon|_{\auU(1)})^{n/2}(\det h)=\upsilon_w(\det h_w)^n$, and similarly for
$g$; the determinant twists $\upsilon_n(\det h)\upsilon_{n'}(\det g)^{-1}$ cancel the powers of
$\upsilon_w$ in the first formula. Finally $\mu(\det h)=\mu(\det h_w,(\det h_w)^{-1})$, and for
$z\in E_w^\times\subset E_v^\times$ one has $z/\aubar z=(z,z^{-1})$, so $\chi_w(z)=\mu(z,z^{-1})$:
$\mu(\det h)\mu(\det g)^{-1}=\chi_w(\det h_w)\chi_w(\det g_w)^{-1}$.
\end{proof}

\begin{remark}\label{pwf:rem-W5}
The normalisation (D3) is the only one of the form $\nu(\det g_w)|\det g_w|^{-1/2}$ for which the
Siegel Levi acts by the same Hecke character $\upsilon(\det)^{-1}|\det|^{1/2}$ at all places
(Lemma~\ref{pwf:siegel}), which is what Theorem~\ref{pwf:inv} needs. With it, the powers of
$\upsilon$ cancel in \eqref{pl:eq-w5}, as stated in (N5).
\end{remark}

\subsubsection{The formal degree in Proposition~\ref{pds:haar}}\label{pwf:ss-haar}

\begin{proposition}\label{pwf:haar}
Proposition~\ref{pds:haar} holds, with the same statement, for every parameter to which it is
applied in the paper. More precisely: for $(\lambda,s)$ satisfying the gap condition
\eqref{pbr:eq-gap} and $dg$ the Haar measure compatible with the Haar measure of mass $1$ on
$\auU(N)$, $\int_G|c|^2dg=\bigl(\prod_{i<j}(\lambda_i-\lambda_j)/(j-i)\bigr)^{-1}$ for the coefficient
$c$ of a unit vector of the lowest $K$-type; and the parameters of \S\ref{pds:sec} satisfy
\eqref{pbr:eq-gap}.
\end{proposition}
\begin{proof}
The first assertion is Corollary~\ref{pbr:degree}: the formal degree $d$ of $\pi(\lambda,s)$ for
$dg$ is defined there by $\int_G|c|^2dg=d^{-1}$ with $c$ the coefficient of a unit highest weight
vector of the lowest $K$-type, and it is proved that $d=\prod_{i<j}(\lambda_i-\lambda_j)/(j-i)$; the
proof does not use Harish-Chandra's formula. The coefficient $c_\tau$ of \S\ref{pds:sec} is of this
kind ($|c_\tau|^2$ does not depend on the convention for inner products), with $N=r$ and $\lambda=
\varsigma_\tau$. Entries of $\varsigma_\tau$ attached to coordinates of opposite signs differ by at
least $\Delta=2r^8$ (\S\ref{pbe:ss-ds}: the first $k-1$ entries are $>-k/2$, the others are
$-\kappa(\tau)-j\Delta$, $1\le j\le m$, with $\kappa(\tau)>0$, and consecutive $j$ have opposite
signs; an entry $>-k/2$ and an entry $-\kappa(\tau)-j\Delta$ with $j\ge2$ differ by more than
$2\Delta-k/2>\Delta$), and $2r^8\ge500r^3+10r^5$ because $r\ge100$ (Lemma~\ref{an:lem-sizes},
(S7)(b)): indeed $500r^3+10r^5\le510r^5\le2r^8$ as soon as $r^3\ge255$. So \eqref{pbr:eq-gap} holds
with $n=r\ge2$. (The other standing hypothesis of \S\ref{pbr:sec}, $\lambda\in(\frac{n-1}2+\Z)^n$ with
distinct entries, holds for $\lambda=\varsigma_\tau$ because it is the Harish-Chandra parameter of a
discrete series of $\auU(k+\varrho,\varrho+1)$.)
\end{proof}

\subsubsection{The archimedean scalar on the maximal compact subgroup}\label{pwf:ss-delta}

Let $\tau$ be a real place, $W$ of signature $(P,Q)$ at $\tau$, $W_\tau=W^+\oplus W^-$ the
decomposition into a positive and a negative definite part, and $K=\auU(W^+)\times\auU(W^-)$. Put
$s=1/\sqrt d$, so that $\tau(\jmath)=is$, and let $\|u\|^2=\langle u^+,u^+\rangle-\langle u^-,u^-\rangle$
be the majorant ($u=u^++u^-$); it is $K$-invariant.

\emph{The vacuum.} \S\ref{ps:ss-weil} defines the Fock model, and its vacuum $1$, for positive
definite spaces. For any non-degenerate $W_\tau$ we call \emph{vacuum} a unit vector $v\in\psHS_\tau(W)$
with
\begin{equation}\label{pwf:eq-vaccoef}
\langle\psrho(u)v,v\rangle=e^{-\pi s\|u\|^2}\qquad(u\in W_\tau).
\end{equation}
(a) \emph{Uniqueness.} A unit vector $v$ of an irreducible unitary representation is determined up
to a scalar of modulus $1$ by $u\mapsto\langle\psrho(u)v,v\rangle$: if $v,v'$ have the same coefficient,
then $\psrho(u)v\mapsto\psrho(u)v'$ extends (by $\|\sum c_i\psrho(u_i)v\|^2=\sum c_i\aubar c_j\beta(\cdot)
\langle\psrho(u_i-u_j)v,v\rangle$, the same for $v'$) to an isometry of $\psHS_\tau(W)$ onto itself
commuting with all $\psrho(u)$, a scalar by (P1).
(b) \emph{Existence.} For positive definite $W$ the Fock vacuum $1$ satisfies \eqref{pwf:eq-vaccoef}:
by \eqref{pk:eq-fock}, $\langle\psrho(y)1,1\rangle=e^{-|y|^2/2}\int e^{z\cdot y}e^{-|z|^2}d\mu(z)=e^{-|y|^2/2}$
(mean value property), and $|y|^2/2=\pi s\|u\|^2$ for $y=c_\psi u$, $c_\psi^2=2\pi s$. For negative
definite $W$, $W^-$ is positive definite and $\beta_W=\aubar{\beta_{W^-}}$ (as $\psi(-t)=\aubar{\psi(t)}$);
so the complex conjugate Hilbert space $\aubar{\psHS_\tau(W^-)}$ with the same operators $\psrho(u)$ is a
Heisenberg representation of $W$, and the vector $1$ in it has coefficient $\aubar{e^{-\pi s\|u\|^2}}=
e^{-\pi s\|u\|^2}$ (the majorants of $W$ and $W^-$ coincide). In this realisation the Fock vectors of
$W$ are the \emph{antiholomorphic} functions of the Fock variable of $W^-$, i.e.\ the holomorphic
functions of the coordinate $w$; this is the convention of \S\ref{ptf:ss-fock} and of (A$_\infty$) in
\S\ref{ptf:ss-conj} (complex structure $-J_\tau$ for $W^-$). For general $W_\tau=W^+\oplus W^-$, $\psHS_\tau(W)=\psHS_\tau(W^+)\widehat
\otimes\psHS_\tau(W^-)$ (an intertwiner, $\beta(W^+,W^-)=1$), and $v^+\otimes v^-$ satisfies
\eqref{pwf:eq-vaccoef}, the majorant being additive. Thus the vacuum exists, is unique up to a
scalar, and is compatible with orthogonal sums of definite spaces of either sign.
(c) \emph{$K$ preserves the vacuum line.} For $\gamma\in K$, $\omega_\tau(\gamma)v$ has coefficient
$\langle\psrho(u)\omega(\gamma)v,\omega(\gamma)v\rangle=\langle\psrho(\gamma^{-1}u)v,v\rangle=
e^{-\pi s\|\gamma^{-1}u\|^2}=e^{-\pi s\|u\|^2}$, so by (a) $\omega_\tau(\gamma)v=\Delta(\gamma)v$ with
$|\Delta(\gamma)|=1$, and $\Delta=\Delta_W$ is a continuous character of $K$. The operator $F\mapsto
F\circ\gamma^{-1}$ on Fock vectors (\S\ref{ptf:ss-fock}) also conjugates $\psrho(u)$ into $\psrho(\gamma u)$
and fixes the vacuum, so it equals $\Delta(\gamma)^{-1}\omega_\tau(\gamma)$ by (P1). Hence
\eqref{ptf:eq-compact} is the assertion $\Delta(\gamma)=\aubar{\det\gamma^-}$.
(d) \emph{Naturality.} An isometry $\sigma\colon W_\tau\to W'_\tau$ carries $\omega_{W,\tau}$ to
$\omega_{W',\tau}$ (Lemma~\ref{pwf:indep} for split spaces; for non-split spaces, Lemma~\ref{pwf:indep} applied to the
isometry $\sigma\oplus\sigma^-$ of the split doubles, together with the characterisation (D2)), and
an intertwiner along $\sigma$ carries vacuum to vacuum (by (a), the majorant being transported); so
$\Delta_{W'}(\sigma\gamma\sigma^{-1})=\Delta_W(\gamma)$.

\begin{lemma}[Reduction to lines]\label{pwf:delta-red}
There are integers $e_+,e_-$, independent of $W$, with $\Delta(\gamma^+,\gamma^-)=\det(\gamma^+)^{e_+}
\det(\gamma^-)^{e_-}$ for all $W$; and $e_++e_-=-1$. Hence \eqref{ptf:eq-compact} holds for all $W$
if and only if $e_+=0$, i.e.\ if and only if, for a positive definite line $\ell$, $\omega_{\ell,\tau}(z)$
fixes the vacuum for $z\in\auU(1)(\R)$.
\end{lemma}
\begin{proof}
By Lemma~\ref{pwf:sum} for $W_\tau=W^+\oplus W^-$ and for orthogonal decompositions of $W^\pm$ into
lines, and by (b) above (the vacuum of an orthogonal sum is the tensor product of the vacua),
$\Delta$ restricted to the diagonal torus of $K$ is the product of the scalars $\Delta_\pm(z)=
z^{e_\pm}$ of the positive and negative lines (continuous characters of $\auU(1)$; by (d) they do not
depend on the line, all positive, resp.\ negative, lines being isometric); a character of $\auU(P)\times\auU(Q)$ is determined by its
restriction to the torus, which gives the first assertion. For the relation, take $W=\ell\oplus\ell^-$
with $\ell$ a positive line; it is split, with Lagrangians $I_\Delta=\{(x,x)\}$, $I_\nabla=
\{(x,-x)\}$. The element $(z,z)\in K$ equals $m(z)$ in the notation of Lemma~\ref{pwf:siegel}
($I=I_\Delta$; on $I_\nabla$ it acts by $z=\aubar z^{-1}$). By Lemma~\ref{pwf:siegel},
$\omega_\tau(m(z))=\upsilon_\tau(z)^{-1}\pwfL_z=z^{-1}\pwfL_z$ in the model $\auS(I_\nabla)$, where
$(\pwfL_z\varphi)(x)=\varphi(z^{-1}x)$ ($|z|=1$). In this model the vacuum is a multiple of the radial
Gaussian $\varphi_0$ of \eqref{pwf:eq-phi0} below (that identification uses only (a) and a Gaussian
integral, not this lemma). As $\pwfL_z\varphi_0=\varphi_0$, $\omega_\tau(m(z))\varphi_0=z^{-1}\varphi_0$.
Hence $\Delta(z,z)=z^{-1}$, i.e.\ $e_++e_-=-1$.
\end{proof}

\begin{proposition}[$e_+=0$; \eqref{ptf:eq-compact}]\label{pwf:eplus}
Let $\ell$ be a positive definite hermitian line at the real place $\tau$ and $\ell^-$ the negative
line. For $(z_1,z_2)\in\auU(\ell)(\R)\times\auU(\ell^-)(\R)$ the operator $\omega_{\ell\oplus\ell^-,\tau}
(z_1,z_2)$ multiplies the Fock vacuum by $z_2^{-1}$. Hence $e_+=0$, $e_-=-1$ in
Lemma~\ref{pwf:delta-red}, and \eqref{ptf:eq-compact} holds: $\Delta(\gamma)=\aubar{\det\gamma^-}$.
\end{proposition}
\begin{proof}
By Lemma~\ref{pwf:sum}, $\omega_{\ell\oplus\ell^-,\tau}(z_1,z_2)=\omega_{\ell,\tau}(z_1)\otimes
\omega_{\ell^-,\tau}(z_2)$, and the vacuum of $\ell\oplus\ell^-$ is the tensor product of the vacua; so the
scalar in question is $\Delta_{\ell}(z_1)\Delta_{\ell^-}(z_2)=z_1^{e_+}z_2^{e_-}$; so it suffices to
compute it for $z_1\ne z_2$, and then $e_+=0$, $e_-=-1$, and the last assertion is
Lemma~\ref{pwf:delta-red}. Write $s=1/\sqrt d$, so that $\tau(\jmath)=is$
(\S\ref{ps:ss-weil}), and $\psi_{k,\tau}(t)=e^{-2\pi it}$.

\emph{Kudla's coordinates.} Let $W=\ell\oplus\ell^-=\C^2$ with $\langle u,w\rangle=\aubar u_1w_1-\aubar
u_2w_2$, and Kudla's skew-hermitian form $\langle u,w\rangle_K=\jmath\langle w,u\rangle$ (linear in $u$).
The vectors $e=(2\jmath)^{-1}(1,-1)$, $f=(1,1)$ span $X=I_\nabla$, $Y=I_\Delta$ and satisfy
$\langle e,e\rangle_K=\langle f,f\rangle_K=0$, $\langle e,f\rangle_K=1$, so in the coordinates
$xe+yf\leftrightarrow(x,y)$ the form is Kudla's (0.6) with $\epsilon=1$ \cite[p.~364]{KudlaSplitting}.
With $V=E$, $(v_1,v_2)=\aubar v_1v_2$ (so $m=1$, signature $(p,q)=(1,0)$) and $\kappa=\frac12$,
$\langle\!\langle(x,0),(0,y)\rangle\!\rangle=\operatorname{Re}(\aubar xy)$, and (P3) is read with
$\psi_K(t)=\psi_{k,\tau}(2t)=e^{-4\pi it}$, so $\eta(t)=e^{-2\pi it}$, and with $\xi=\upsilon_\tau$,
$\xi(z)=z/|z|$. The paper's element $g=(z_1,z_2)$ is Kudla's element acting on rows by the map
$g^{-1}$ (\S\ref{pwf:ss-def}, left actions). With $a_i=\aubar z_i=z_i^{-1}$, $g^{-1}f=(a_1,a_2)$ and
$g^{-1}e=(2\jmath)^{-1}(a_1,-a_2)$, whence Kudla's matrix
\[
\begin{pmatrix}\alpha&\beta\\\gamma&\delta\end{pmatrix}=\begin{pmatrix}\frac{a_1+a_2}2&\frac{a_1-a_2}{4\jmath}\\
\jmath(a_1-a_2)&\frac{a_1+a_2}2\end{pmatrix}.
\]
As $\gamma\ne0$, $\ker\gamma=0$ and $g$ lies in the cell $j=1$ (1.11). Writing
$g=p_1\tau p_2$ with $p_i=\bigl(\begin{smallmatrix}\alpha_i&\beta_i\\0&\delta_i\end{smallmatrix}\bigr)\in
P_Y$ and $\tau=\bigl(\begin{smallmatrix}0&-1\\1&0\end{smallmatrix}\bigr)$ (1.12), one finds
$\gamma=\delta_1\alpha_2$, and $\alpha_2\aubar\delta_2=1$ (unitarity of $p_2$); so $x(g)=\delta_1\delta_2=
\gamma/|\alpha_2|^2$ and $\xi(x(g))=\gamma/|\gamma|$.

\emph{The Weil index.} For Kudla's character $e(t)=e^{2\pi it}$, comparison of Theorem~3.1 with
Corollary~3.2 \cite[pp.~378--379]{KudlaSplitting} gives $\gamma_\R(\eta_0\circ RV)=i^{p-q}=i$ for
$\eta_0(t)=e(t/2)$. Our $\eta(t)=\aubar{\eta_0(2t)}$. Kudla's $\gamma_F(\eta\circ L)$ is Weil's
$\gamma(f)$ for the character of second degree $f(x)=\eta((x,x)_L)$ \cite[p.~364]{KudlaSplitting}, and
Weil's $\gamma(f)$ is the number of modulus $1$ such that $\int(\int\Phi(x-y)f(y)\,dy)\,dx=\gamma(f)
|\varrho_f|^{-1/2}\int\Phi(x)\,dx$ for all $\Phi\in\auS(G)$, $|\varrho_f|>0$ the module of the
symmetric isomorphism attached to $f$ \cite[Th\'eor\`eme~2, p.~161, and Corollaire~2, p.~162]{WeilActa}.
Two consequences, for $G=\R^N$: (i) if $f_\lambda(x)=f(\lambda^{1/2}x)$ with $\lambda>0$ (this is the
character $\eta(\lambda\,\cdot)\circ L$), the substitutions $y\mapsto\lambda^{-1/2}y$, $x\mapsto
\lambda^{-1/2}x$ turn the left side for $(f_\lambda,\Phi)$ into $\lambda^{-N}$ times the left side for
$(f,\Phi(\lambda^{-1/2}\cdot))$, so $\gamma(f_\lambda)|\varrho_{f_\lambda}|^{-1/2}=\lambda^{-N/2}\gamma(f)
|\varrho_f|^{-1/2}$; comparing arguments (take $\int\Phi\ne0$), $\gamma(f_\lambda)=\gamma(f)$. (ii)
Conjugating the identity for $(f,\aubar\Phi)$ gives it for $(\aubar f,\Phi)$ with the factor
$\aubar{\gamma(f)}|\varrho_f|^{-1/2}$, and $\aubar f$ is attached to $-\varrho_f$, of the same module;
so $\gamma(\aubar f)=\aubar{\gamma(f)}$. Hence $\gamma_\R(\eta\circ RV)=\aubar{\gamma_\R(\eta_0\circ RV)}=
-i$ and $\beta_Y(g)=\xi(x(g))(-i)^{-1}=i\,\gamma/|\gamma|$.

\emph{The vacuum in this model.} Let
\begin{equation}\label{pwf:eq-phi0}
\varphi_0(x)=e^{-\pi|x|^2/s}\qquad(x\in X=\C,\ \text{the point }xe).
\end{equation}
With
$\psrho(xe)\varphi(u)=\varphi(u+x)$, $\psrho(yf)\varphi(u)=e^{-4\pi i\operatorname{Re}(\aubar uy)}\varphi(u)$
and $\psrho(xe+yf)=\beta(xe,yf)^{-1}\psrho(xe)\psrho(yf)$ (P2), a Gaussian integral gives
$\langle\psrho(xe+yf)\varphi_0,\varphi_0\rangle/\|\varphi_0\|^2=\exp(-\pi|x|^2/(2s)-2\pi s|y|^2)$
(the phase $[\frac x2,y]$ produced by centring the integral is $\beta(xe,yf)$ and cancels). In the
paper's coordinates $xe+yf=(\frac x{2\jmath}+y,-\frac x{2\jmath}+y)$, and $|x|^2/(2s)+2s|y|^2=s(|u_1|^2+
|u_2|^2)=s\|u\|^2$ for $u=xe+yf$. By (a), (b) of \S\ref{pwf:ss-delta} (the vacuum of $\ell\oplus\ell^-$
is $v_0\otimes v_0^-$, with coefficient $e^{-\pi s\|u\|^2}$, $v_0^-$ the vacuum of the negative line
defined there), $\varphi_0/\|\varphi_0\|$ is the vacuum in the model $\auS(X)$, and by (c) there
$\omega(g)\varphi_0=c\,\varphi_0$ with $|c|=1$.

\emph{The scalar.} By Rao's formula at $x=0$ \cite[p.~400]{KudlaSplitting},
$c=(\omega(g)\varphi_0)(0)=\beta_Y(g)\int_\C\psi_K\bigl(\tfrac12\langle\!\langle y\gamma,y\delta\rangle\!\rangle
\bigr)\varphi_0(y\gamma)\,d\mu(y)$ with $\mu$ a positive Haar measure. Put $w=z_1\aubar z_2$. Then
$\langle\!\langle y\gamma,y\delta\rangle\!\rangle=|y|^2\operatorname{Re}(\aubar\gamma\delta)$ and
$\aubar\gamma\delta=-\frac\jmath2(z_1-z_2)(\aubar z_1+\aubar z_2)=s\operatorname{Im}w$, while
$|\gamma|^2=s^2|z_1-z_2|^2=s^2|1-w|^2$. So the integrand is $\exp(-\pi s|y|^2(|1-w|^2+2i\operatorname{Im}w))
=\exp(-2\pi s|y|^2(1-\aubar w))$, and the integral is a positive multiple of $(1-\aubar w)^{-1}$,
whose phase is $(1-w)/|1-w|$. On the other hand $\gamma/|\gamma|=i(\aubar z_1-\aubar z_2)/|z_1-z_2|=
-i\aubar z_2(1-\aubar w)/|1-w|$, so $\beta_Y(g)=\aubar z_2(1-\aubar w)/|1-w|$. Therefore $c$ is a
positive multiple of $\aubar z_2(1-\aubar w)(1-w)/|1-w|^2=\aubar z_2$, and $|c|=1$ gives $c=z_2^{-1}$.
\end{proof}

\begin{proof}[Second proof, without the Weil index and Rao's formula]
This proof uses only Lemma~\ref{pwf:siegel}, the homomorphism property of $\omega$ and (P1). Keep
$W=\ell\oplus\ell^-$, $I=I_\Delta=\C f$, $I^*=I_\nabla$, but write the points of $I^*$ as $\xi e_0$ with
$e_0=(1,-1)=2\jmath e$, so that $\varphi_0(\xi e_0)=e^{-4\pi s|\xi|^2}$ and, from $\langle yf,\xi e_0\rangle=
2\aubar y\xi$, $[\xi e_0,yf]=e^{8\pi is\operatorname{Im}(\aubar y\xi)}$. For real $t$ let $b_t\colon\xi e_0
\mapsto it\xi f$ (admissible: $it$ is imaginary). Then $\beta(\xi e_0,b_t\xi e_0)=e^{-4\pi ist|\xi|^2}$,
so $\omega(n(b_t))$ is multiplication by $e^{4\pi ist|\xi|^2}$ (Lemma~\ref{pwf:siegel}). The element
$w_0=(1,-1)\in\auU(W)$ exchanges $e_0$ and $f$, and $n^-_t:=w_0n(b_t)w_0^{-1}$ is the identity on $I^*$
and maps $yf\mapsto yf+ity\,e_0$. The operator $(T\varphi)(\xi)=\int_\C\varphi(\xi')[\xi'e_0,w_0^{-1}\xi
e_0]\,d\xi'=\int_\C\varphi(\xi')e^{8\pi is\operatorname{Im}(\aubar\xi\xi')}\,d\xi'$ is, up to a positive
factor, unitary and an intertwiner along $w_0$ (translation by $\xi_0$ becomes multiplication by
$[\xi_0e_0,w_0^{-1}\xi e_0]^{-1}=[\xi e_0,w_0\xi_0e_0]$, and multiplication by $[\cdot,y_0f]$ becomes
translation by $w_0y_0f$, both by bimultiplicativity of $[\,,]$ and $\pwfA(w_0x,x')=\pwfA(x,w_0^{-1}x')$).
By (P1), $\omega(w_0)=\lambda T$, so $\omega(n^-_t)=T\,\omega(n(b_t))\,T^{-1}$ exactly. With
$\int_\C e^{-4\pi a|\xi'|^2+8\pi is\operatorname{Im}(\aubar\xi\xi')}d\xi'=(4a)^{-1}e^{-4\pi s^2|\xi|^2/a}$ for
$\operatorname{Re}a>0$ (a two-dimensional Gaussian integral, holomorphic in $a$), the same with
$-8\pi is$ in place of $8\pi is$, and Fourier inversion ($T^{-1}$ is $16s^2$ times the operator with
the conjugate kernel):
\[
T^{-1}\varphi_0=4s\,\varphi_0,\qquad \omega(n^-_t)\varphi_0=T\bigl(4s\,e^{-4\pi s(1-it)|\xi|^2}\bigr)
=(1-it)^{-1}\exp\bigl(-4\pi s|\xi|^2/(1-it)\bigr).
\]
Now let $g=(z_1,z_2)$ with $z_1\ne-z_2$, $D=\frac{z_1+z_2}2$, $B=\frac{z_1-z_2}2$, $it=B/D$ (imaginary,
as $|z_1|=|z_2|=1$) and $a=\aubar D^{-1}$. In the basis $(f,e_0)$, $g=\bigl(\begin{smallmatrix}D&B\\B&D
\end{smallmatrix}\bigr)$, and, using $D^2-B^2=z_1z_2=D/\aubar D$,
\[
n(b_t)\,m(a)\,n^-_t=\begin{pmatrix}1&it\\0&1\end{pmatrix}\begin{pmatrix}\aubar D^{-1}&0\\0&D\end{pmatrix}
\begin{pmatrix}1&0\\it&1\end{pmatrix}=\begin{pmatrix}\aubar D^{-1}+B^2/D&B\\B&D\end{pmatrix}=g .
\]
By Lemma~\ref{pwf:siegel}, $\omega(m(a))\varphi(\xi)=\upsilon_\tau(a)^{-1}|a|\varphi(\aubar a\xi)$, and
$\upsilon_\tau(a)^{-1}|a|=\aubar a=D^{-1}$. Hence
\[
\omega(g)\varphi_0(\xi)=D^{-1}(1-it)^{-1}\exp\Bigl(-4\pi s|\xi|^2\Bigl(\frac{|a|^2}{1-it}-it\Bigr)\Bigr).
\]
Here $1-it=(D-B)/D=z_2/D$, so $D^{-1}(1-it)^{-1}=z_2^{-1}$; and with $w=z_1\aubar z_2$,
$|a|^2/(1-it)=1/(\aubar Dz_2)=2/(1+\aubar w)=2w/(1+w)$ and $it=(w-1)/(w+1)$, so the bracket is $1$.
Thus $\omega(g)\varphi_0=z_2^{-1}\varphi_0$, as in the first proof.
\end{proof}

\begin{remark}\label{pwf:delta-open}
The result depends on the signs fixed in \S\ref{ps:ss-weil}: with $\psi_{\Q,\infty}(t)=e^{+2\pi it}$
(or with $\tau(\jmath)=-i/\sqrt d$) the same computation gives $c=z_1^{-1}$, i.e.\ $e_+=-1$, $e_-=0$.
For the conventions of the paper, $\Delta(\gamma)=\aubar{\det\gamma^-}$ as used in \S\ref{ptf:ss-fock}
and in \S\ref{ps:ss-weil} (where a positive definite space carries no scalar on the compact group).
\end{remark}


%% ===== proofs/B-second-moment.tex =====
% proofs/B-second-moment.tex
% The second moment of the Poincare map: Proposition au:moment and Theorem B of
% sec-automorphic.tex. Labels of this file start with pb:.
\providecommand{\auO}{\mathcal O}
\providecommand{\auPl}{\Phi}
\providecommand{\auPa}{\Phi_a}
\providecommand{\auPb}{\Phi_b}
\providecommand{\auS}{\mathcal S}
\providecommand{\auSram}{\mathcal S^{\mathrm{ram}}}
\providecommand{\auT}{\mathcal T}
\providecommand{\auAz}{\mathcal A_0}
\providecommand{\auG}{\overline G}
\providecommand{\auP}{\mathcal P}
\providecommand{\auDX}{D_X}
\providecommand{\audT}{d_{\auT}}
\providecommand{\auNm}{\operatorname{Nm}}
\providecommand{\auU}{\mathrm U}
\providecommand{\auGL}{\mathrm{GL}}
\providecommand{\auA}{\mathbb A}
\providecommand{\auF}{\mathbb F}
\providecommand{\aurd}{\mathrm{rd}_E}
\providecommand{\aubar}[1]{\overline{#1}}
\providecommand{\auf}[1]{\texttt{#1}}
% macros of this file (prefix \pb)
\providecommand{\pbR}{\overline R_0}
\providecommand{\pbUb}{\overline U}
\providecommand{\pbHb}{\overline H}
\providecommand{\pbtr}{\operatorname{tr}}
\providecommand{\pbvol}{\operatorname{vol}}
\providecommand{\pbDid}{\mathcal D_{\mathrm{id}}}
\providecommand{\pbDmain}{\mathcal D_{\mathrm{main}}}
\providecommand{\pbdeg}{\operatorname{deg}}
\providecommand{\pbmass}{\operatorname{mass}}
\providecommand{\pbFix}{\operatorname{Fix}}
\providecommand{\pbpr}{\mathrm{pr}}
\providecommand{\pbrpsi}{r_\psi}
\providecommand{\pbW}{\mathbb W}
\providecommand{\pbg}{\mathsf g}
\providecommand{\pbLM}{\mathrm{LM}}
\providecommand{\pbnat}{\natural}
\providecommand{\pbone}{\mathbf 1}
\providecommand{\pbHS}[1]{\|#1\|}
\providecommand{\pbip}[2]{\langle#1,#2\rangle}

\section{The second moment of the Poincar\'e map}\label{pb:sec}

This section concerns Proposition~\ref{au:moment} (the second moment of the Poincar\'e map) and,
through it, Theorem~\ref{thm:srcB}. We use the notation of \S\S\ref{ss:au-notation}--\ref{ss:au-spaces}.

\emph{Structure of the proof.} The proof has three layers.
\begin{itemize}
\item[(I)] \emph{From four inequalities to Theorem~\ref{thm:srcB}} (\S\S\ref{pb:ss-setting}--\ref{pb:ss-thmB}).
The Poincar\'e map $\auP$ intertwines an explicit unitary representation $\pbrpsi$ of a subgroup
$\pbR$ of $\auG$ on the source with the action of $\auG$ on the forms on $Y_0$
(Lemma~\ref{pb:equiv}). We state four inequalities for the operators
$P_s=P^*\varrho(s)P$, $P=c_P\auP$: \textup{(E1)} $P_s$ is close to $\pbrpsi(s)$ near $\pbR$;
\textup{(E2)}, \textup{(E3)} two bounds for sums of $\pbHS{P_s}^2$ over subgroups of $\auG$ that fix
vectors of the frame; \textup{(E4)} a bound for the character induced from the unipotent
character. From these we deduce, with complete proofs,
Proposition~\ref{au:moment}(ii) and (iii) (Propositions~\ref{pb:ii} and~\ref{pb:iii}). Part (i) is
deduced from Proposition~\ref{au:mult} (Proposition~\ref{pb:i}). Together with the deduction of
\S\ref{ss:au-proofAB} this gives Theorem~\ref{thm:srcB} (Theorem~\ref{pb:thmB}).
\item[(II)] \emph{From amplitudes to} (E1)--(E3) (\S\ref{pb:ss-layer2}). Assuming that $\auP$ is
a Poincar\'e series of a Bessel function with the usual properties
(Hypothesis~\ref{pb:bessel}), we prove the Petersson formula for $P_s$
(Proposition~\ref{pb:unfold}), identify the contribution of the Bessel subgroup with $\pbrpsi(s)$
(Proposition~\ref{pb:besselterms}), compare the two normalising constants $\pbDmain$ and $\pbDid$
(Lemma~\ref{pb:sizes}; this is where the window~\eqref{eq:au-level} for the level $N$ is used),
and deduce (E1)--(E3) from three statements about sums of amplitudes
(Proposition~\ref{pb:reduction}).
\item[(III)] \emph{The analytic core}: the existence of the Bessel functions, the three statements
about amplitudes, the inequality (E4), and Proposition~\ref{au:mult}. They are stated precisely here
(Hypothesis~\ref{pb:bessel}, Theorem~\ref{pb:amp}, Theorem~\ref{pb:est}(E4), Proposition~\ref{au:mult})
and proved in \S\S\ref{pbe:sec}--\ref{pm:sec} and \S\S\ref{pds:sec}--\ref{pfg:sec}: the Bessel functions in Theorem~\ref{pbe:bessel},
\textup{(X)} and \textup{(L1)} in Theorem~\ref{pbe:XL1}, \textup{(L2)} in Theorem~\ref{pbm:L2},
\textup{(E4)} in Theorem~\ref{pbm:E4}, and Proposition~\ref{au:mult} in Proposition~\ref{pm:mult}.
\end{itemize}
So Theorem~\ref{thm:srcB} is proved modulo the published results named in the proofs; the value of the
covolume $\nu$ (Lemma~\ref{pb:covol}) is proved in \S\ref{pbv:sec}. (Part~(b) of Theorem~\ref{pm:arch}, Harish-Chandra's formula for stable
elliptic orbital integrals of discrete series coefficients, is proved in \S\ref{pbr:sec}.) The complete dependency tree is in \S\ref{pb:ss-remains}.

\subsection{The setting}\label{pb:ss-setting}

\subsubsection*{The group and the frame}
Put $n=\varrho+1$, so that $r=k+2\varrho+1$. The space $Y_0=\mathbb H^{\varrho}\perp X_0\perp
\langle-N\rangle$ has the basis
\[
e_1,\dots,e_\varrho;\quad x_1,\dots,x_k,\ z;\quad e_1^*,\dots,e_\varrho^*,
\]
with $h(e_a,e_b^*)=\delta_{ab}$, $h(x_i,x_j)=\delta_{ij}$, $h(z,z)=-N$ and all other products
zero. The \emph{frame} is $(f_0,\dots,f_\varrho)=(e_1,\dots,e_\varrho,z)$; $A=\langle e_1,\dots,
e_\varrho\rangle$ is isotropic and $A^\perp/A=X_0\perp Ez$. Write $G=\auU(Y_0)$, a group over
$E^+$; $\Gamma=G(E^+)$; $G_\infty=G(E^+\otimes\R)=\prod_{\tau\in\auPl}G_\tau$;
$G_f=G(\auA_{E^+,f})$. Let $K\subset G_f$ be the stabiliser of the lattice $L$ of $Y_0$ described in the proof of
Lemma~\ref{pb:KR} (for its relation to the lattice $\mathsf M_0$ of \eqref{ps:eq-lattice} see
Lemma~\ref{pb:lattice}) and $J=J_M
\subseteq K$ the principal congruence subgroup of level $\pi_M$; it is normal in $K$. Each chosen
place $w$ is split over $E^+$, so the component of $K$ at the place below $w$ is
$\auGL_r(\auO_{E,w})$, and reduction modulo the chosen places identifies $K/J$ with
$\auG=\prod_w\auGL_r(\auF_w)$. We write $\auG=G_0\times G_1$, where $G_i$ is the product over the
$f$ chosen places above $p_i$. For $s\in\auG$, $\varrho(s)$ is right translation by a lift of $s$
to $K$; it is a unitary operator of $\auAz$, for the inner product $\pbip y{y'}=\int_{\Gamma
\backslash G(\auA)}\aubar y\,y'$. \emph{In this section all inner products are conjugate-linear in
the first variable and linear in the second}; this is the convention for which the map $\auP$ of
\S\ref{ss:au-spaces}, defined by $\pbip y{\auP\varphi}=\int_{[R_0]}\mathcal W_\varphi(r)\aubar{y(rg_0)}
\,dr$, is linear.

\subsubsection*{Three conventions}\label{pb:ss-conv}
The following conventions hold in all the sections on Theorem~\ref{thm:srcB} (\S\S\ref{pb:sec}--\ref{pm:sec}
and \S\S\ref{pds:sec}--\ref{pfg:sec}); sections that use another one say so and point here.
\begin{enumerate}
\item[(Cv1)] \emph{Hermitian forms.} Every hermitian form $h$ on a space over $E$ or over a completion
of $E$ ($Y_0$, $X_0$, the $W_j$, $Y_\tau$, \dots) is \emph{linear in the first variable} and
conjugate-linear in the second: $h(\lambda u,v)=\lambda h(u,v)=h(u,\aubar\lambda v)$, $h(v,u)=
\aubar{h(u,v)}$. Thus $h(v,e_a^*)$ is the coefficient of $e_a$ in $v$. (The inner products of
functions and of the spaces $\pbW_\tau$ above follow the opposite convention; they are not
hermitian forms on the spaces $Y_j$. For representation coefficients: $c_\tau(g)=\langle v_0,\pi(g)v_0\rangle$
of \S\S\ref{pbe:ss-ds}, \ref{pco:sec}, with the inner product conjugate-linear in the first variable, is
the same function as $\langle\pi(g)v_0,v_0\rangle'$ for the inner product $\langle u,w\rangle'=\langle w,u\rangle$,
linear in the first variable, used in \S\S\ref{pm:sec}, \ref{pbr:sec}; no conjugation separates them.)
\item[(Cv2)] \emph{The additive character.} $\psi_E=\psi_\Q\circ\operatorname{Tr}_{E/\Q}$, where
$\psi_\Q$ is the standard character of $\auA_\Q/\Q$: $\psi_{\Q,\infty}(t)=e^{-2\pi it}$ and
$\psi_{\Q,p}(t)=e^{2\pi i\{t\}_p}$. So at every complex place $\tau$ of $E$ the component of $\psi_E$ is
$w\mapsto e^{-2\pi i(w+\aubar w)}$, and $\psi_E$ is trivial exactly on the inverse different
$\varpi^{-1}\widehat\auO_E$ at the finite places. (By Lemma~\ref{pbp:sign} the opposite sign at
infinity would change nothing in the results.)
\item[(Cv3)] \emph{The Haar measure of $U_\infty$.} On $U_\tau$, $\tau\in\auPl$, $du_\tau$ is
Lebesgue measure in the free entries of the rows $\tilde e_a$ of $u$ in the normalised basis of
\S\ref{pbe:ss-real} (the complex entries above the diagonal block that are not fixed by unitarity,
and the imaginary parts of the entries $(\tilde e_a,\tilde e_a^*)$); equivalently, $du_\tau$ is the Haar
measure with $\int_{U_\tau}G_\varrho(u^{-1})\,du_\tau=1$ for the Gaussian $G_\varrho$ of
\S\ref{prp:ss-haar} (Lemma~\ref{prp:heis}). On $U_\infty$, $du=\prod_\tau du_\tau$, and on
$U_\tau\backslash G_\tau$ the quotient measure. This is the measure of \S\S\ref{pbe:ss-real},
\ref{pbs:ss-measures} and of Lemma~\ref{pfr:haar}, and the constant $\nu$ of Lemma~\ref{pb:ortho}
is the covolume $\pbvol(U_c\backslash U_\infty)$ for it.
\end{enumerate}

Let $U\subset G$ be the unipotent radical of the parabolic subgroup $P$ of \S\ref{ss:au-spaces}:
the elements that preserve the flag $\langle e_1\rangle\subset\langle e_1,e_2\rangle\subset\dots
\subset A$, act trivially on its graded pieces and act trivially on $A^\perp/A$. Let
$H=\auU(X_0)\subset G$, acting trivially on $A$, $z$ and the $e_a^*$, and $R_0=U\rtimes H$. For
$u\in U$ put
\[
\sigma(u)=\sum_{a=1}^{\varrho-1}h(ue_{a+1},e_a^*)+h(uz,e_\varrho^*),
\]
the sum of the entries of $u$ just above the diagonal of its block on the frame; $\sigma$ is
additive on $U$ and invariant under conjugation by $H$. The character $\psi_U$ of
\S\ref{ss:au-spaces} is taken to be
\begin{equation}\label{pb:eq-psi}
\psi_U(u)=\psi_E\Bigl(-\frac{\sigma(u)}{M\varpi}\Bigr)\qquad(u\in U(\auA_{E^+})),
\end{equation}
where $\psi_E$ is the standard additive character of $\auA_E/E$ and $\varpi=\gamma\,\Phi'_{q_1}(\zeta)\sqrt{-d}\in\auO_E$ is the purely imaginary generator of the different of $E$ of Lemma~\ref{pds:kappa} ($\zeta$ a primitive $q_1$-th root of unity, $\gamma$ a unit with $\gamma\aubar\gamma=1$ making $\gamma\Phi'_{q_1}(\zeta)$ real). Thus
$\psi_U$ is trivial on $U(E^+)$, invariant under conjugation by $H(\auA_{E^+})$, and trivial on
$U(\auA_{E^+,f})\cap J$: for $u\in J$ the entries of $u-1$ in the basis of the lattice vanish
modulo every chosen place and its conjugate, so they lie in $M\widehat{\auO}_E$, and $\psi_E$ is
trivial on $\varpi^{-1}\widehat{\auO}_E$. We write $\psi_\infty$, $\psi_f$ for the restrictions of
$\psi_U$ to $U_\infty=U(E^+\otimes\R)$ and to $U(\auA_{E^+,f})$; $\psi_\infty(u)\psi_f(u)=1$ for
$u\in U(E^+)$. (This is the character with which $\auP$ is defined in \S\ref{ss:au-spaces}, and it
is the character $\psi_0$ of \S\ref{ps:ss-Y}; likewise the pair $(g_0,\mathsf t)$ with which $\auP$
is defined there is the pair of Theorem~\ref{pb:est} below.)

\subsubsection*{The finite groups}
Put $K_U=U(\auA_{E^+,f})\cap K$, $K_H=H(\auA_{E^+,f})\cap K$ and $K_R=R_0(\auA_{E^+,f})\cap K$.

\begin{lemma}\label{pb:KR}
$K_R=K_U\rtimes K_H$. The group $K_H$ consists of the elements of the stabiliser $K^\circ$ of
$\auO_E^k$ in $H(\auA_{E^+,f})$ whose last row is $\equiv(0,\dots,0,1)$ modulo $\pi_N$, and
$K_H\cap J=J_X$.
\end{lemma}
\begin{proof}
Let $g=uh\in K$ with $u\in U(\auA_f)$, $h\in H(\auA_f)$, and let $L$ be the lattice of $Y_0$. Then
$g$ preserves $L\cap A$ and $L\cap A^\perp$, hence the lattice $(L\cap A^\perp)/(L\cap A)$ of
$A^\perp/A$, on which $u$ acts trivially; so $h$ preserves this lattice. As $h$ fixes the $e_a$
and the $e_a^*$, and $L$ is the sum of its intersections with $A\oplus A^*$ and with $X_0\perp Ez$,
$h$ preserves $L$. So $h\in K_H$ and $u=gh^{-1}\in K_U$. The lattice of $X_0\perp Ez$ is spanned
by $x_1,\dots,x_{k-1}$, $x_k/\aubar{\pi_N}+z/N$ and $z$ (it is unimodular: the last two vectors
have the Gram matrix $\bigl(\begin{smallmatrix}0&-1\\-1&-N\end{smallmatrix}\bigr)$), and its
intersection with $X_0$ is $L_1=\auO x_1+\dots+\auO x_{k-1}+\pi_N\auO x_k$. An element $h$ of
$H(\auA_f)$, which fixes $z$, preserves this lattice if and only if $hL_1=L_1$ and $(hx_k-x_k)/
\aubar{\pi_N}\in L_1$. Away from $N$ this says $h\in K^\circ$. At a prime $w\mid\pi_N$ the group
$H$ is $\auGL_k(E_w)$, the component at $\aubar w$ being the inverse adjoint, and $\aubar{\pi_N}$
is a unit at $w$. Read at $w$, the second condition says $h_we_k\in e_k+L_{1,w}=\auO_w^{k-1}\oplus
(1+\pi_N\auO_w)e_k$, and then the first, applied to $e_1,\dots,e_{k-1}\in L_{1,w}$, says that the
entries $h_{ij}$ ($j<k$) are integral with $h_{kj}\in\pi_N\auO_w$. Thus $h_w\in\auGL_k(\auO_w)$ with
last row $\equiv(0,\dots,0,1)$ modulo $w$; conversely such an $h_w$ satisfies both conditions at $w$.
Read at $\aubar w$, the two conditions give the same congruence for the last row of $h_w^{-1}$,
which is equivalent.
Finally $K_H\cap J$ is defined by the congruence $\equiv1$ at the chosen places, which
gives the level $J_X$ of \S\ref{ss:au-spaces}.
\end{proof}

Let $\pbUb$, $\pbHb$, $\pbR$ be the images of $K_U$, $K_H$, $K_R$ in $\auG$. By the lemma
$\pbR=\pbUb\rtimes\pbHb$, and each of the three groups is the product of its components at the
chosen places: $\pbHb\cong K_H/J_X\cong\prod_w\auGL_k(\auF_w)$, and $\pbUb_w$ is the unipotent
radical of the parabolic subgroup of $\auGL_r(\auF_w)$ with blocks $1,\dots,1,k+1,1,\dots,1$, of
order $(\auNm w)^{\varrho(2k+2\varrho+1)}$. (For $\pbUb$ one uses that $U$ is a split unipotent
group, so that $U(\auO_w)\to U(\auF_w)$ is onto.) Hence
\begin{equation}\label{pb:eq-orders}
\#\pbUb=M^{f\varrho(2k+2\varrho+1)},\qquad\#\pbR=\#\pbUb\cdot\#\pbHb .
\end{equation}
By the remarks after \eqref{pb:eq-psi}, $\psi_f$ induces a character of $\pbUb$, again written
$\psi_U$, which is invariant under conjugation by $\pbHb$. Let $\tilde\psi\colon\auG\to\C$ be
$\psi_U$ on $\pbUb$ and $0$ elsewhere, and
\[
\psi^{\pbnat}(s)=\frac1{\#\auG}\sum_{g\in\auG}\tilde\psi(gsg^{-1})\qquad(s\in\auG).
\]

\subsubsection*{The source and the representation \texorpdfstring{$\pbrpsi$}{r}}
The source $\auS=\mathcal A_{\lambda^X}(J_X,1)$ carries the inner product of
\S\ref{ss:au-spaces}, $\pbip\varphi{\varphi'}=\sum_{c}\pbip{\varphi(h_c)}{\varphi'(h_c)}_{W_{\lambda^X}}$,
where $c$ runs over the set $\mathcal C=H(E^+)\backslash H(\auA_f)/J_X$ of \emph{classes} and $h_c$
is a representative of $c$ (the level is free by Lemma~\ref{au:free}, so there are no weights), and
where the hermitian form on $W_{\lambda^X}$ is invariant under $H_\infty=H(E^+\otimes\R)$. Since
the level is free, $\varphi\mapsto(\varphi(h_c))_c$ is an isomorphism $\auS\to W_{\lambda^X}^{\,
\mathcal C}$, and
\begin{equation}\label{pb:eq-DX}
\auDX=\#\mathcal C\cdot d,\qquad d=\dim W_{\lambda^X}=\prod_{\tau\in\auPa}\binom{a+k-1}{k-1}.
\end{equation}
The group $K_H$ normalises $J_X$, so $\pbHb=K_H/J_X$ acts on $\auS$ by right translation,
$(R(t)\varphi)(g)=\varphi(gk_t)$; this action is unitary. For $s\in\auG$ put
\[
\pbrpsi(s)=\begin{cases}\psi_U(\bar u)\,R(\bar t)&\text{if }s=\bar u\,\bar t\in\pbR,\
\bar u\in\pbUb,\ \bar t\in\pbHb,\\[2pt] 0&\text{if }s\notin\pbR.\end{cases}
\]
Because $\psi_U$ is a character of $\pbUb$ invariant under $\pbHb$, the restriction of $\pbrpsi$ to
$\pbR$ is a unitary representation of $\pbR$ on $\auS$, and
\begin{equation}\label{pb:eq-rext}
\pbrpsi(st)=\pbrpsi(s)\pbrpsi(t),\quad\pbrpsi(ts)=\pbrpsi(t)\pbrpsi(s),\quad\pbrpsi(s)^*=\pbrpsi
(s^{-1})\qquad(t\in\pbR,\ s\in\auG).
\end{equation}

\begin{lemma}[Equivariance]\label{pb:equiv}
For every choice of $g_0$ and $\mathsf t$ in the definition of $\auP$, and every $t\in\pbR$,
$\varrho(t)\circ\auP=\auP\circ\pbrpsi(t)$.
\end{lemma}
\begin{proof}
The group $R_0(\auA)$ is unimodular: $U$ is unipotent, and the determinant of the action of $H$
on $\operatorname{Lie}U$ is an algebraic character $H\to\mathbb G_m$ over $E^+$, hence trivial,
because the characters of the unitary group $H$ factor through $\det\colon H\to\auU(1)$ and
$\auU(1)$ is anisotropic. So conjugation by $H(\auA)$ preserves the Haar measure of $U(\auA)$, and
the measure on $[R_0]=R_0(E^+)\backslash R_0(\auA)$ is invariant under right
translations. Let $k=u_kh_k\in K_R$ be a lift of $t$, $y\in\auAz$ and $\varphi\in\auS$. As $k$ is
a finite adele and $g_0$ is real, they commute, and the substitution $r\mapsto rk$ gives
\[
\pbip y{\varrho(k)\auP\varphi}=\pbip{\varrho(k)^{-1}y}{\auP\varphi}=\int_{[R_0]}\mathcal W_\varphi(r)
\,\aubar{y(rg_0k^{-1})}\,dr=\int_{[R_0]}\mathcal W_\varphi(rk)\,\aubar{y(rg_0)}\,dr .
\]
For $r=uh$ we have $rk=u\cdot(hu_kh^{-1})\cdot hh_k$, so, by the invariance of $\psi_U$ under $H$,
$\mathcal W_\varphi(rk)=\psi_U(u)\psi_U(u_k)\,\mathsf t\bigl(h_\infty^{-1}\varphi(h_fh_k)\bigr)=
\psi_U(u_k)\,\mathcal W_{R(h_k)\varphi}(r)$. Hence $\pbip y{\varrho(k)\auP\varphi}=\psi_U(u_k)\pbip y
{\auP R(h_k)\varphi}=\pbip y{\auP\,\pbrpsi(t)\varphi}$, as $\auP$ is linear and the inner product is
linear in the second variable.
\end{proof}

\begin{lemma}[Trace of $\pbrpsi$]\label{pb:trace}
Let $s\in\auG$ have trivial components at all chosen places above one of the primes $p_0$, $p_1$.
Then $\pbtr\pbrpsi(s)=\auDX\cdot\tilde\psi(s)$.
\end{lemma}
\begin{proof}
If $s\notin\pbR$ both sides vanish. Let $s=\bar u\bar t$. As $\pbR=\prod_w\pbUb_w\rtimes\pbHb_w$,
the components of $\bar t$ above that prime, $p_j$ say, are trivial. Let $k\in K_H$ be a lift of
$\bar t$. In the decomposition $\auS\cong W_{\lambda^X}^{\,\mathcal C}$, $(R(\bar t)\varphi)(h_c)=
\varphi(h_ck)=\eta\,\varphi(h_{c'})$ if $h_ck=\eta h_{c'}j$ with $\eta\in H(E^+)$, $j\in J_X$. So
$R(\bar t)$ permutes the blocks, and a block contributes to the trace only if $c'=c$. In that case
$\eta\in H(E^+)\cap h_cK_Hh_c^{-1}$ has finite order and $h_c^{-1}\eta h_c=kj^{-1}$ is $\equiv1$
modulo a chosen place above $p_j$. The argument of Lemma~\ref{au:free}, which uses only one chosen
place, gives $\eta=1$, so $k\in J_X$ and $\bar t=1$. Thus the trace is $0$ unless $\bar t=1$, that
is $s\in\pbUb$, and then $\pbrpsi(s)=\psi_U(s)\cdot1$.
\end{proof}

\subsubsection*{The constants}
For a positive definite standard space of rank $l$ let $\pbmass_l=\sum_{[g]}\#(\Gamma_l\cap
gK^\circ_lg^{-1})^{-1}$ be the mass of the genus of $\auO_E^l$, the sum being over $\auU_l(E^+)
\backslash\auU_l(\auA_f)/K_l^\circ$. If $J'\subseteq K_l^\circ$ is a subgroup of finite index which is free, then
\begin{equation}\label{pb:eq-classes}
\#\bigl(\Gamma_l\backslash\auU_l(\auA_f)/J'\bigr)=[K_l^\circ:J']\cdot\pbmass_l ,
\end{equation}
because the finite group $\Gamma_l\cap gK_l^\circ g^{-1}$ acts freely on $gK_l^\circ/J'$, by left
multiplication. For $\tau\in\auPl$ let
\[
\pbdeg_\tau=\prod_{1\le i<j\le r}\frac{\varsigma_{\tau,i}-\varsigma_{\tau,j}}{j-i},\qquad
\pbdeg=\prod_{\tau\in\auPl}\pbdeg_\tau ,
\]
where $\varsigma_\tau=(\varsigma_{\tau,1}>\dots>\varsigma_{\tau,r})$ is the parameter of
\S\ref{ss:au-spaces}. By Weyl's formula $\pbdeg_\tau$ is the dimension of the irreducible
representation of the compact group $\auU(r)$ with infinitesimal character $\varsigma_\tau$; it is
the formal degree of the discrete series $\varsigma_\tau$ for the Haar measure that corresponds
to the measure of total mass $1$ on $\auU(r)$. This is the number $\deg$ of
Proposition~\ref{au:moment}. Put
\begin{equation}\label{pb:eq-D}
\pbDid=\#\pbR\cdot\auDX,\qquad\pbDmain=\frac{\auDX^2}{\pbdeg\cdot\pbmass_r}.
\end{equation}

\subsubsection*{Vectors of the frame modulo the chosen places}
Let $w$ be a chosen place and $V_w=\auF_w^{\,r}$ the reduction modulo $w$ of the lattice of $Y_0$,
in a basis of the lattice containing the $e_a$, $z$ and the $e_a^*$. The hermitian form pairs the
completions at $w$ and at $\aubar w$; in this basis it gives a non-degenerate symmetric bilinear
form on $V_w$, with matrix $\mathsf J$. An element $s\in\auG$ acts on $V_w$ by $s_w$, and we let it
act on a second copy $V_{\aubar w}$ of $V_w$ by the contragredient $s_w^{\vee}=\mathsf J^{-1}\,
{}^ts_w^{-1}\mathsf J$; this is the action modulo $\aubar w$ of a lift of $s$ to $K$. For
$\mathfrak p\in\{w,\aubar w\}$ let $F_{\mathfrak p}\subset V_{\mathfrak p}$ be the span of the
frame, of dimension $n$. A \emph{frame vector at $w$} is a non-zero vector $\mathsf d$ of $F_w$ or
of $F_{\aubar w}$; $\pbFix(\mathsf d)\subseteq\auG$ is its stabiliser. For a set $\mathsf A$ of
chosen places above $p_0$ and a family $\mathsf d=(\mathsf d_w)_{w\mid p_0}$ of frame vectors put
$\pbFix(\mathsf A,\mathsf d)=\bigcap_{w\in\mathsf A}\pbFix(\mathsf d_w)$; it is a subgroup of
$\auG$, equal to $\auG$ if $\mathsf A=\emptyset$. For $g\in\auGL_r(\auF_w)$ let
\[
Z_w(g)=\bigl({}^tf_j\,\mathsf J\,g\,f_l\bigr)_{0\le j,l\le\varrho}\in M_n(\auF_w),\qquad
\delta_l(g)=\det\bigl(Z_w(g)_{j,l'}\bigr)_{0\le j,l'<l}\quad(0\le l\le n),
\]
the matrix of the products of the frame with its image, and its leading minors. An element
$s\in\auG$ is \emph{open} if there is a chosen place $w$ above $p_0$ with $\delta_l(s_w)\ne0$ for
$1\le l\le\varrho$.

\subsection{The four inequalities}\label{pb:ss-est}

All norms of linear maps are Hilbert--Schmidt norms. For a linear map $P\colon\auS\to\auAz$ and
$s\in\auG$ we write $P_s=P^*\varrho(s)P\in\operatorname{End}(\auS)$.

\begin{theorem}[The four inequalities]\label{pb:est}
There are a point $g_0\in G_\infty$, a linear form $\mathsf t$ on $W_{\lambda^X}$ and a scalar
$c_P>0$ such that the map $P=c_P\auP$, with $\auP$ defined by $g_0$, $\mathsf t$ and the character
\eqref{pb:eq-psi}, has the following properties.
\begin{enumerate}
\item[(E1)] For $i\in\{0,1\}$ let $\Omega_i$ be the set of the $s\in\auG$ whose components at the
chosen places above $p_i$ are those of an element of $\pbR$. Then
$\sum_{s\in\Omega_i}\pbHS{P_s-\pbrpsi(s)}^2\le b^{-12}\,\pbDid$.
\item[(E2)] For every family $\mathsf d$ of frame vectors at the $f$ chosen places above $p_0$,
$\sum_{s\in\pbFix(\mathrm{all},\mathsf d)}\pbHS{P_s}^2\le b^{-6kf}\,\pbDmain$.
\item[(E3)] For every set $\mathsf A$ of chosen places above $p_0$ and every family $\mathsf d$,
\[
\sum_{s\in\pbFix(\mathsf A,\mathsf d),\ s\text{ open}}\pbHS{P_s}^2\le\bigl(1+4b^{-4}\bigr)\,b^{-6k
\#\mathsf A}\,\pbDmain .
\]
\end{enumerate}
Moreover
\begin{enumerate}
\item[(E4)] $\sum_{s\in\auG}|\psi^{\pbnat}(s)|^2\le2$.
\end{enumerate}
\end{theorem}
\smallskip\noindent(E1)--(E3) are deduced in \S\ref{pb:ss-layer2} (Proposition~\ref{pb:reduction}) from
Hypothesis~\ref{pb:bessel} and Theorem~\ref{pb:amp}, proved in \S\S\ref{pbe:sec}, \ref{pbm:sec} and
\S\S\ref{pds:sec}--\ref{pfg:sec}. (E4) is a statement
about the finite group $\auG$ and the character $\psi_U$ of $\pbUb$ only; it is proved in \S\ref{pbm:ss-E4} (Theorem~\ref{pbm:E4}), using $k\ge2\varrho^2+3$, $p_i\ge b^{10}$ and $18\varrho f\le b^8$ (Lemma~\ref{an:lem-sizes}(S7)).\par\smallskip

\begin{remark}\label{pb:rem-E4}
(E4) says that the representation of $\auG$ induced from the character $\psi_U$ of $\pbUb$ is
almost as free of multiplicities as its dimension allows: a computation with double cosets gives
$\sum_s|\psi^{\pbnat}(s)|^2=(\#\pbUb)^2(\#\auG)^{-1}\cdot\dim\operatorname{End}_{\auG}
\bigl(\operatorname{Ind}_{\pbUb}^{\auG}\psi_U\bigr)$, and $\#\auG/(\#\pbUb)^2$ is the number of
double cosets of $\pbUb$ of full size, if all had full size.
\end{remark}

\subsection{Lemmas on finite groups}\label{pb:ss-finite}

\begin{lemma}[Central idempotents]\label{pb:central}
Let $\mathsf G$ be a finite group and $x=\sum_sx(s)s$ a central idempotent of $\C[\mathsf G]$.
\begin{enumerate}
\item[(a)] $x(s^{-1})=\aubar{x(s)}$, and $\sum_s|x(s)|^2=x(1)\in[0,1]$.
\item[(b)] For a unitary representation $\varrho$ of $\mathsf G$, $\varrho(x)=\sum_sx(s)\varrho(s)$
is an orthogonal projection that commutes with $\varrho(\mathsf G)$.
\item[(c)] If $\mathsf G$ is a subgroup of a finite group $\mathsf G'$ and $F$ is a function on
$\mathsf G'$ with values in a Hilbert space, then
$\sum_{s\in\mathsf G'}\bigl\|\sum_{t\in\mathsf G}x(t)F(st)\bigr\|^2\le\sum_{s\in\mathsf G'}\|F(s)
\|^2$.
\end{enumerate}
\end{lemma}
\begin{proof}
A central idempotent is a sum of the idempotents $e_\chi=\chi(1)(\#\mathsf G)^{-1}\sum_s
\aubar{\chi(s)}s$ over a set of irreducible characters $\chi$. This gives the first claim of (a)
and $x(1)=\sum_\chi\chi(1)^2/\#\mathsf G\in[0,1]$; and $x(1)=(x\cdot x)(1)=\sum_sx(s)x(s^{-1})=
\sum_s|x(s)|^2$. (b) $\varrho(x)^*=\sum_s\aubar{x(s)}\varrho(s^{-1})=\varrho(x)$ by (a), and
$\varrho(x)^2=\varrho(x)$. (c) Apply (b) to the unitary representation $(R_tF)(s)=F(st)$ of
$\mathsf G$ on the functions on $\mathsf G'$: $R(x)$ is an orthogonal projection.
\end{proof}

A real function $\beta$ on a finite group $\mathsf G$ is \emph{of positive type} if there are
vectors $u_a$ ($a\in\mathsf G$) of a complex inner product space with $\pbip{u_a}{u_b}=\beta(a^{-1}
b)$ for all $a,b$. Then $\beta(s)\le\beta(1)$ for all $s$, by the Cauchy--Schwarz inequality, as
$\|u_a\|^2=\beta(1)$. The restriction of $\beta$ to a subgroup is of positive type.

\begin{lemma}[Singular minors]\label{pb:singular}
Let $\mathsf G$ be a finite group, $\auF$ a field with $q$ elements, $\pi\colon\mathsf G\to\auGL_r
(\auF)$ a homomorphism, $A_0,A_1$ two $r\times l$ matrices over $\auF$, $\beta\ge0$ a function of
positive type on $\mathsf G$ and $Y\ge0$. Suppose that for every non-zero $v\in\auF^l$
\[
\sum_{s:\ \pi(s)A_0v=A_0v}\beta(s)\le Y\qquad\text{and}\qquad\sum_{s:\ {}^t\pi(s)^{-1}A_1v=A_1v}
\beta(s)\le Y .
\]
Then
\[
q\sum_{s:\ \det({}^tA_1\pi(s)A_0)=0}\beta(s)\ \le\ 2\sum_{s\in\mathsf G}\beta(s)+q^{\,l+r/2}\,Y .
\]
\end{lemma}
\begin{proof}
Let $\psi$ be a non-trivial additive character of $\auF$, $Z(s)={}^tA_1\pi(s)A_0$, and for $y,v\in
\auF^l$ let $\Sigma(y,v)=\sum_s\psi({}^tyZ(s)v)\beta(s)$. If $\det Z(s)=0$ the left kernel of
$Z(s)$ has at least $q$ elements, and by orthogonality of characters $q^{-l}\sum_v\psi({}^tyZ(s)
v)$ is $1$ if ${}^tyZ(s)=0$ and $0$ otherwise. So
\[
q\sum_{\det Z(s)=0}\beta(s)\le\sum_s\#\{y:{}^tyZ(s)=0\}\,\beta(s)=\sum_y\,q^{-l}\sum_v\Rel\Sigma(y,
v).
\]
If $y=0$ or $v=0$ then $\Sigma(y,v)=\sum_s\beta(s)$; this happens for $q^l$ pairs with $y=0$ and
for one $v$ for each of the other $y$, so these pairs contribute at most $2\sum_s\beta(s)$. It
remains to prove $\Rel\Sigma(y,v)\le q^{r/2}Y$ for $y,v\ne0$; then the remaining pairs contribute at
most $q^l\cdot q^{r/2}Y$.

Put $a=A_1y$, $z=A_0v$, so that ${}^tyZ(s)v={}^ta\,\pi(s)z$. Let $\mathcal H$ be the space of the
$u_a$ and $\ell^2(\auF^r)$ the space of functions on $\auF^r$, with the orthonormal bases
$(\epsilon_\xi)_\xi$ of delta functions and $(c_\xi)_\xi$, $c_\xi(\zeta)=q^{-r/2}\psi(-{}^t\xi
\zeta)$. Let
\[
X_1=\sum_{g\in\mathsf G}c_{{}^t\pi(g)^{-1}a}\otimes u_g,\qquad X_0=\sum_{g\in\mathsf G}\epsilon_{\pi(g)
z}\otimes u_g\qquad\text{in }\ell^2(\auF^r)\otimes\mathcal H .
\]
Since $\pbip{c_\xi}{\epsilon_\zeta}=\aubar{c_\xi(\zeta)}=q^{-r/2}\psi({}^t\xi\zeta)$ and ${}^t\bigl(
{}^t\pi(g)^{-1}a\bigr)\pi(g')z={}^ta\,\pi(g^{-1}g')z$, we get
\[
\pbip{X_1}{X_0}=q^{-r/2}\sum_{g,g'}\psi\bigl({}^ta\,\pi(g^{-1}g')z\bigr)\beta(g^{-1}g')=q^{-r/2}\,\#
\mathsf G\cdot\Sigma(y,v) .
\]
In the same way, by orthonormality, $\|X_0\|^2=\sum_{g,g'}[\pi(g)z=\pi(g')z]\,
\beta(g^{-1}g')=\#\mathsf G\sum_{s:\pi(s)z=z}\beta(s)\le\#\mathsf G\cdot Y$, and $\|X_1\|^2=\#
\mathsf G\sum_{s:{}^t\pi(s)^{-1}a=a}\beta(s)\le\#\mathsf G\cdot Y$. Hence $\Rel\Sigma(y,v)=q^{r/2}
(\#\mathsf G)^{-1}\Rel\pbip{X_1}{X_0}\le q^{r/2}(\#\mathsf G)^{-1}\|X_1\|\|X_0\|\le q^{r/2}Y$.
\end{proof}

\subsubsection*{Labels and bad types}
Let $\auF$ be a field with $q$ elements. For a label $c$ of degree $r$ over $\auF$
(\S\ref{ss:au-spaces}), with irreducible factors $c_1,\dots,c_h$ of degrees $d_1,\dots,d_h$, the
irreducible character $\Xi_c$ of $\auGL_r(\auF)$ named by $c$ is, in Green's
parametrisation~\cite[p.~403 and Thm.~13]{Green1955} (see also~\cite[Ch.~IV, (6.7)--(6.8)]{Macdonald}), the character induced from the
parabolic subgroup with blocks $d_1,\dots,d_h$ of the product of the cuspidal characters of
$\auGL_{d_j}(\auF)$ attached to the roots of $c_j$ (through a fixed embedding of $\overline{\auF}{}^
\times$ in $\C^\times$). We use the following facts of~\cite[Lemma~7.4 and Thms.~12--13]{Green1955}: the cuspidal characters of
$\auGL_d(\auF)$ have degree $\prod_{i=1}^{d-1}(q^i-1)$; $\Xi_c$ is irreducible, because the
$c_j$ are distinct; and $\Xi_c\ne\Xi_{c'}$ for $c\ne c'$. Counting the cosets of the parabolic
subgroup,
\begin{equation}\label{pb:eq-green}
\Xi_c(1)=\frac{\prod_{i=1}^r(q^i-1)}{\prod_{j=1}^h(q^{d_j}-1)} ,\qquad\mu(c):=\frac{\Xi_c(1)^2}{\#
\auGL_r(\auF)}=\frac{\prod_{i=1}^r(q^i-1)}{q^{r(r-1)/2}\prod_j(q^{d_j}-1)^2}\ \ge\ \frac{1-r/q}{
q^r} .
\end{equation}
The idempotent of $c$ is $e_c=\Xi_c(1)(\#\auGL_r(\auF))^{-1}\sum_g\aubar{\Xi_c(g)}g$, with
coefficient $\mu(c)$ at $1$.

\begin{lemma}[Bad types are rare]\label{pb:bad}
For $i\in\{0,1\}$ let $\mathsf g_i=\sum\bigotimes_{w\mid p_i}e_{c_w}\in\C[G_i]$, the sum over the
families $(c_w)_{w\mid p_i}$ of labels such that no $c_w$ has a root in $\auF_w$ that is a
$q_2$-th root of unity and at least one $c_w$ is irreducible with a root generating $\auF_{p_i^r}^
\times$. Then $\mathsf g_i$ is a central idempotent of $\C[G_i]$, $e_0=\mathsf g_0\mathsf g_1$ in
$\C[\auG]$, and the coefficient of $1-\mathsf g_i$ at $1$ is at most $b^{-8}/2$.
\end{lemma}
\begin{proof}
The $\bigotimes_we_{c_w}$ are the idempotents of distinct irreducible characters of $G_i$, so
$\mathsf g_i$ is a central idempotent; $e_0=\mathsf g_0\mathsf g_1$ is the definition of $e_0$.
Write $q=p_i$ and call a monic polynomial of degree $r$ over $\auF_q$ \emph{allowed} if it is a
label without a root that is a $q_2$-th root of unity, and \emph{primitive} if it is irreducible
with a root of order $q^r-1$. As $r>1$, primitive polynomials are allowed. Let $\mathsf a$ be the
number of allowed and $\mathsf p$ the number of primitive polynomials, and $m_0=(1-r/q)q^{-r}$. The
coefficient of $\mathsf g_i$ at $1$ is $\sum\prod_w\mu(c_w)$ over the good families, which by
\eqref{pb:eq-green} is at least $m_0^f\bigl(\mathsf a^f-(\mathsf a-\mathsf p)^f\bigr)$. So the
coefficient of $1-\mathsf g_i$ at $1$ is at most
\[
1-(m_0\mathsf a)^f+\bigl(m_0(q^r-\mathsf p)\bigr)^f .
\]
A monic polynomial of degree $r\ge2$ that is not allowed is not squarefree ($q^{r-1}$
polynomials), or vanishes at $0$ ($q^{r-1}$), or vanishes at one of at most $q_2$ roots of unity
($q_2q^{r-1}$). So $m_0\mathsf a\ge(1-r/q)(1-(2+q_2)/q)\ge1-\varepsilon$ with $\varepsilon=(1+q_2+
2r)/q$, and $1-(m_0\mathsf a)^f\le f\varepsilon$. The roots of the primitive polynomials are the
generators of $\auF_{q^r}^\times$, $r$ for each polynomial, so $\mathsf p=\varphi(q^r-1)/r$. For an
integer $\mathsf n\ge2$ with $\omega$ prime factors, $\varphi(\mathsf n)/\mathsf n\ge\prod_{j=2}^{
\omega+1}(1-1/j)=1/(\omega+1)\ge1/(2\log\mathsf n+1)$. With $\mathsf n=q^r-1\ge q^r/2$ this gives
$\mathsf p/q^r\ge\theta:=\bigl(2r(2r\log q+1)\bigr)^{-1}$, hence $m_0(q^r-\mathsf p)\le1-\theta\le
e^{-\theta}$. Altogether the coefficient is at most $f\varepsilon+e^{-f\theta}$. By the definition
of the data, $x^{11}\le f\le x^{12}$, $r\le3x^5$, $q_2\le b^{1/100}$, and
$b^{10}\le q\le b^{11}$; with \eqref{eq:au-scales} this gives $f\varepsilon\le b^{-9}$ and
$f\theta\ge9\log b$, and $2b^{-9}\le b^{-8}/2$. Here $q_2\le b^{1/100}$ holds because
$q_2=(q_0d)^{s+6}$ with $s\le32\log x$ by (D3), $\log(q_0d)\le2\log C_P\le2\log\log\log b$ (as
$C_P\ge q_0+2d$ and $10^{12}C_P\le\log\log b$), and $\log x\le10^{-6}\log x_1\le10^{-6}(\log b)^{1/6}$
by (D2); so $\log q_2\le70\log x\cdot\log\log\log b\le10^{-4}(\log b)^{1/6}\log\log\log b\le
(\log b)/100$.
\end{proof}

\subsection{The sum over the group: Proposition~\ref{au:moment}(iii)}\label{pb:ss-iii}

\begin{lemma}\label{pb:basic}
Let $P=c_P\auP$ be as in Theorem~\ref{pb:est}. Then:
\begin{enumerate}
\item[(a)] $P_{st}=P_s\,\pbrpsi(t)$ and $P_{ts}=\pbrpsi(t)\,P_s$ for $t\in\pbR$, $s\in\auG$;
\item[(b)] $\pbHS{P^*P-1}^2\le b^{-12}\auDX$ and $\pbHS{P^*P}^2\le2\auDX$;
\item[(c)] $\beta(s)=\pbHS{P_s}^2$ is of positive type on $\auG$, and $\beta(s)\le2\auDX$ for all
$s$.
\end{enumerate}
\end{lemma}
\begin{proof}
(a) By Lemma~\ref{pb:equiv}, $P_{st}=P^*\varrho(s)\varrho(t)P=P^*\varrho(s)P\pbrpsi(t)$. As
$\pbrpsi(t)$ is unitary,
\[
P_{ts}=(\varrho(t^{-1})P)^*\varrho(s)P=(P\pbrpsi(t^{-1}))^*\varrho(s)P=\pbrpsi(t)P_s .
\] (b) For $s\in\pbR$, by (a) and \eqref{pb:eq-rext}, $P_s-\pbrpsi(s)=(P^*P
-1)\pbrpsi(s)$, which has the norm of $P^*P-1$. As $\pbR\subseteq\Omega_0$, (E1) gives $\#\pbR
\cdot\pbHS{P^*P-1}^2\le b^{-12}\#\pbR\cdot\auDX$. Then
$\pbHS{P^*P}\le\pbHS{P^*P-1}+\pbHS1\le(b^{-6}+1)\sqrt{\auDX}$. (c) Let $u_a=\varrho(a)PP^*\varrho(a)^{-1}$, in the space of endomorphisms of
$\auAz$ with the Hilbert--Schmidt inner product. Then
\[
\pbip{u_a}{u_b}=\pbtr\bigl(\varrho(a)PP^*\varrho(a^{-1}b)PP^*\varrho(b^{-1})\bigr)=\pbtr\bigl(P_{a^{-1}
b}\,P_{a^{-1}b}^{\,*}\bigr)=\beta(a^{-1}b) .
\]
So $\beta(s)\le\beta(1)=\pbHS{P^*P}^2\le2\auDX$.
\end{proof}

\begin{proposition}[Proposition~\ref{au:moment}(iii)]\label{pb:iii}
Let $P$ be as in Theorem~\ref{pb:est}. Then $\sum_{s\in\auG}\pbHS{P^*\varrho(s)P}^2\le(1+9b^{-4})\,
\pbDmain$.
\end{proposition}
\begin{proof}
Let $\beta(s)=\pbHS{P_s}^2$ and, for a set $\mathsf A$ of chosen places above $p_0$ and a family
$\mathsf d$ of frame vectors, $m(\mathsf A,\mathsf d)=\sum_{s\in\pbFix(\mathsf A,\mathsf d)}
\beta(s)$. Put $\Theta=b^{6k}$ and choose $(\mathsf A,\mathsf d)$ such that $\Theta^{\#\mathsf A}
m(\mathsf A,\mathsf d)$ is largest. Since $m(\emptyset,\mathsf d)=\sum_{s\in\auG}\beta(s)$, it
suffices to prove
\begin{equation}\label{pb:eq-max}
m(\mathsf A,\mathsf d)\le(1+9b^{-4})\,\Theta^{-\#\mathsf A}\,\pbDmain .
\end{equation}
If $\mathsf A$ is the set of all chosen places above $p_0$, this is (E2). Otherwise let $w\notin
\mathsf A$ be a chosen place above $p_0$, $\Gamma_0=\pbFix(\mathsf A,\mathsf d)$ and $m=m(\mathsf A,
\mathsf d)$. For a frame vector $\mathsf d'$ at $w$, the group $\Gamma_0\cap\pbFix(\mathsf d')$ is
$\pbFix(\mathsf A\cup\{w\},\cdot)$ for the family with $\mathsf d'$ at $w$, so by the choice of
$(\mathsf A,\mathsf d)$
\begin{equation}\label{pb:eq-stab}
\sum_{s\in\Gamma_0\cap\pbFix(\mathsf d')}\beta(s)\le m/\Theta .
\end{equation}
Let $1\le l\le\varrho$. Apply Lemma~\ref{pb:singular} to $\mathsf G=\Gamma_0$, $\auF=\auF_w$ (so
$q=p_0$), $\pi(s)=s_w$, the matrix $A_0$ with columns $f_0,\dots,f_{l-1}$ and $A_1=\mathsf JA_0$,
and $Y=m/\Theta$. For $v\ne0$ the vector $A_0v$ is a non-zero vector of the span of the frame;
$\pi(s)A_0v=A_0v$ means $s\in\pbFix(\mathsf d')$ for $\mathsf d'=A_0v\in F_w$, and ${}^t\pi(s)^{-1}
A_1v=A_1v$ means $s_w^\vee A_0v=A_0v$, that is $s\in\pbFix(\mathsf d')$ for $\mathsf d'=A_0v\in
F_{\aubar w}$. So the hypotheses hold by \eqref{pb:eq-stab}, and since ${}^tA_1\pi(s)A_0$ is the
leading block of size $l$ of $Z_w(s_w)$,
\[
p_0\sum_{s\in\Gamma_0,\ \delta_l(s_w)=0}\beta(s)\le2m+p_0^{\,l+r/2}\,m/\Theta\le3m ,
\]
because $p_0^{\,l+r/2}\le b^{11(\varrho+r/2)}\le b^{6k}$ (as $11(4\varrho+1)\le k$). An element of
$\Gamma_0$ that is not open has $\delta_l(s_w)=0$ for some $1\le l\le\varrho$. So, with
$p_0\ge b^{10}$ and $3\varrho\le b^6$,
\[
\sum_{s\in\Gamma_0,\ s\text{ not open}}\beta(s)\le\varrho\cdot3m\,b^{-10}\le b^{-4}m .
\]
With (E3) for the open elements, $m\le(1+4b^{-4})\Theta^{-\#\mathsf A}\pbDmain+b^{-4}m$. As $(1+9
\epsilon)(1-\epsilon)\ge1+4\epsilon$ for $0\le\epsilon\le\frac1{10}$, \eqref{pb:eq-max} follows.
\end{proof}

\subsection{The good part: Proposition~\ref{au:moment}(ii)}\label{pb:ss-ii}

\begin{lemma}[One prime]\label{pb:oneprime}
Let $P$ be as in Theorem~\ref{pb:est}, $i\in\{0,1\}$, and let $x$ be a central idempotent of
$\C[G_i]$ with $x(1)\le b^{-8}/2$. Then $\pbHS{\varrho(x)P}^2\le2b^{-4}\auDX$.
\end{lemma}
\begin{proof}
Regard $G_i$ as a subgroup of $\auG$. By Lemma~\ref{pb:central}(b), $\varrho(x)$ is an orthogonal
projection, so
\[
\pbHS{\varrho(x)P}^2=\pbtr\bigl(P^*\varrho(x)P\bigr)=\sum_{t\in G_i}x(t)\pbtr\pbrpsi(t)+\pbtr P'_x,
\qquad P'_x=\sum_{t\in G_i}x(t)\bigl(P_t-\pbrpsi(t)\bigr).
\]
\emph{The first term.} The elements of $G_i$ have trivial components above the other prime, so by
Lemma~\ref{pb:trace} the first term is $\auDX\sum_{t\in G_i}x(t)\tilde\psi(t)$. For
$g\in\auG$ and $t\in G_i$, $gtg^{-1}=g_itg_i^{-1}\in G_i$, and $x$ is a class function on $G_i$; so
the sum does not change if $\tilde\psi(t)$ is replaced by $\tilde\psi(gtg^{-1})$, hence by
$\psi^{\pbnat}(t)$. By the Cauchy--Schwarz inequality, Lemma~\ref{pb:central}(a) and (E4),
\[
\Bigl|\sum_{t\in G_i}x(t)\psi^{\pbnat}(t)\Bigr|^2\le x(1)\sum_{t\in\auG}|\psi^{\pbnat}(t)|^2
\le b^{-8} .
\]
\emph{The second term.} Put $P'_s=P_s-\pbrpsi(s)$. Let $i'$ be the other index. For $s\in\pbR$ and
$t\in G_i$ we have $st\in\Omega_{i'}$, and $P'_{st}=\pbrpsi(s)P'_t$ by Lemma~\ref{pb:basic}(a) and
\eqref{pb:eq-rext}. Let $F(s)=P'_s$ for $s\in\Omega_{i'}$ and $F(s)=0$ otherwise. Then, for
$s\in\pbR$, $\sum_tx(t)F(st)=\pbrpsi(s)P'_x$ has the norm of $P'_x$, and by
Lemma~\ref{pb:central}(c) and (E1)
\[
\#\pbR\cdot\pbHS{P'_x}^2\le\sum_{s\in\auG}\Bigl\|\sum_{t\in G_i}x(t)F(st)\Bigr\|^2\le\sum_{s\in
\Omega_{i'}}\pbHS{P'_s}^2\le b^{-12}\,\#\pbR\cdot\auDX .
\]
So $|\pbtr P'_x|\le\pbHS1\,\pbHS{P'_x}\le b^{-6}\auDX$. Together, $\pbHS{\varrho(x)P}^2\le(b^{-4}+
b^{-6})\auDX$.
\end{proof}

\begin{proposition}[Proposition~\ref{au:moment}(ii)]\label{pb:ii}
Let $P$ be as in Theorem~\ref{pb:est}. Then $\pbHS{\varrho(e_0)P\,\pbpr_{\mathsf U}}^2\ge\dim
\mathsf U-5b^{-4}\auDX$ for every subspace $\mathsf U\subseteq\auS$. In particular $\audT>0$.
\end{proposition}
\begin{proof}
Let $\epsilon_i=\varrho(\mathsf g_i)$, with $\mathsf g_i$ as in Lemma~\ref{pb:bad}; these are
commuting orthogonal projections and $\varrho(e_0)=\epsilon_0\epsilon_1$. By Lemmas~\ref{pb:bad}
and~\ref{pb:oneprime} applied to $x=1-\mathsf g_i$, $\pbHS{(1-\epsilon_i)P}^2\le2b^{-4}\auDX$. Since
$1-\epsilon_0\epsilon_1=(1-\epsilon_0)+\epsilon_0(1-\epsilon_1)$ with orthogonal ranges,
$\pbHS{(1-\varrho(e_0))P}^2\le\pbHS{(1-\epsilon_0)P}^2+\pbHS{(1-\epsilon_1)P}^2\le4b^{-4}\auDX$. Now
\[
\pbHS{\varrho(e_0)P\pbpr_{\mathsf U}}^2=\pbHS{P\pbpr_{\mathsf U}}^2-\pbHS{(1-\varrho(e_0))P\pbpr_{
\mathsf U}}^2\ge\pbHS{P\pbpr_{\mathsf U}}^2-4b^{-4}\auDX ,
\]
and $\pbHS{P\pbpr_{\mathsf U}}^2=\pbtr(\pbpr_{\mathsf U}P^*P)=\dim\mathsf U+\pbtr\bigl(\pbpr_{\mathsf
U}(P^*P-1)\bigr)\ge\dim\mathsf U-\sqrt{\dim\mathsf U}\,\pbHS{P^*P-1}\ge\dim\mathsf U-b^{-6}\auDX$ by
Lemma~\ref{pb:basic}(b). For $\mathsf U=\auS$ this gives $\varrho(e_0)P\ne0$, so $\auT\ne0$.
\end{proof}

\subsection{The dimension of the target, and Theorem B}\label{pb:ss-thmB}

\begin{proposition}[Proposition~\ref{au:moment}(i)]\label{pb:i}
If Proposition~\ref{au:mult} holds for $j=m$, then $\audT\le\#\auG\cdot\pbdeg\cdot\pbmass_r$.
\end{proposition}
\begin{proof}
By Proposition~\ref{au:mult}, $\audT=\dim e_0\mathcal A(W_m)\le\dim\mathcal A(W_m)$. The space
$W_m$ is the positive definite standard space of rank $r$ and $\auU(W_m)$ is compact at the
infinite places, where the discrete series with parameter $\varsigma_\tau$ is the irreducible
representation $\mathsf F_\tau$ of dimension $\pbdeg_\tau$. So $\mathcal A(W_m)$ is the space of
algebraic modular forms on $\auU(W_m)$ with values in $\bigotimes_\tau\mathsf F_\tau$ and level
$J_M$. This level is free (Lemma~\ref{au:free}, as $r\le2n_m$), so by \eqref{pb:eq-classes} its
dimension is $[K^\circ_r:J_M]\cdot\pbmass_r\cdot\pbdeg=\#\auG\cdot\pbmass_r\cdot\pbdeg$.
\end{proof}

\subsubsection*{The lattice of $Y_0$}
In this section and in all the sections on Theorem~\ref{thm:srcB}, $L$ is the lattice of $Y_0$ with the
basis $e_a,e_a^*$ $(1\le a\le\varrho)$, $x_1,\dots,x_{k-1}$, $x_k/\aubar{\pi_N}+z/N$, $z$ (proof of
Lemma~\ref{pb:KR}), and $K$, $J$, $\auAz$, $\varrho$, $\auP$ are formed with it. The lattice of $Y_0$
with which the lower chain is defined is the lattice $\mathsf M_0$ of \eqref{ps:eq-lattice}; in the
notation there the string is $f_{a-1}=e_a$, $f_\varrho=z$, $f_{2\varrho+1-a}=e_a^*$ $(1\le a\le\varrho)$
and $y=x_k$, the group $R_0$ and the character $\psi_0=\psi_U$ are those of the present section
(\S\ref{ps:ss-Y}), and $\mathsf M_0$ is spanned by $x_1,\dots,x_{k-1}$, $x_k/\aubar{\pi_N}+u_0z/N$,
$u_0z$, $u_0e_a$, $u_0e_a^*$, where $u_0$ is a finite idele with $u_0\aubar{u_0}=1$ that is not a unit
(at the chosen places, and at the places above $q_1$ where $\alpha/\aubar\alpha$ is not a unit). The
two lattices have different stabilisers, so the spaces of forms are different spaces of functions.
The next lemma shows that all the statements of Proposition~\ref{au:moment} pass from one to the
other; it is applied with $\mathsf u=u_0$.

\begin{lemma}[Change of the lattice]\label{pb:lattice}
Let $\mathsf u$ be a finite idele of $E$ with $\mathsf u\aubar{\mathsf u}=1$, let $\mathsf h_{\mathsf u}$
be the element of the centre of $H(\auA_{E^+,f})$ that acts on $X_0$ by the scalar $\mathsf u^{-1}$
(and trivially on $A$, on $z$ and on the $e_a^*$), and let $L_{\mathsf u}$ be the lattice of $Y_0$
spanned over $\widehat\auO_E$ by $x_1,\dots,x_{k-1}$, $x_k/\aubar{\pi_N}+\mathsf uz/N$, $\mathsf uz$,
$\mathsf ue_a$, $\mathsf ue_a^*$ $(1\le a\le\varrho)$.
\begin{enumerate}
\item[(a)] $L_{\mathsf u}=\mathsf u\cdot\mathsf h_{\mathsf u}L$. So $L_{\mathsf u}$ is unimodular, its
stabiliser is $K_{\mathsf u}=\mathsf h_{\mathsf u}K\mathsf h_{\mathsf u}^{-1}$, the principal congruence
subgroup of level $\pi_M$ in it is $J_{\mathsf u}=\mathsf h_{\mathsf u}J\mathsf h_{\mathsf u}^{-1}$, and
the reduction of $k\in K_{\mathsf u}$ at the chosen places, in the basis of $L_{\mathsf u}$ just given, is
the image of $\mathsf h_{\mathsf u}^{-1}k\,\mathsf h_{\mathsf u}\in K$ in $\auG$.
\item[(b)] Let $\auAz^{\mathsf u}$ be the space $\auAz$ of \S\ref{ss:au-spaces} formed with
$J_{\mathsf u}$ in place of $J$, with the action $\varrho_{\mathsf u}$ of $\auG=K_{\mathsf u}/J_{\mathsf u}$
by right translation, and let $\auP_{\mathsf u}\colon\auS\to\auAz^{\mathsf u}$ be the Poincar\'e map
defined by the same $g_0$, $\mathsf t$, $\psi_U$ and Haar measure of $[R_0]$ as $\auP$. Then
$(\mathsf j_{\mathsf u}y)(g)=y(g\,\mathsf h_{\mathsf u}^{-1})$ defines a unitary isomorphism
$\mathsf j_{\mathsf u}\colon\auAz^{\mathsf u}\to\auAz$ with $\mathsf j_{\mathsf u}\varrho_{\mathsf u}(s)=
\varrho(s)\mathsf j_{\mathsf u}$ for $s\in\auG$, the map $(R_{\mathsf u}\varphi)(h)=\varphi(h\,\mathsf
h_{\mathsf u}^{-1})$ is a unitary automorphism of $\auS$, and
\[
\auP_{\mathsf u}=\mathsf j_{\mathsf u}^{-1}\circ\auP\circ R_{\mathsf u} .
\]
\item[(c)] If {\rm(i)}, {\rm(ii)}, {\rm(iii)} of Proposition~\ref{au:moment} hold for $P=c_P\auP$,
the space $\auAz$ and $\auT=\varrho(e_0)\auAz$, then they hold, with the same inner product on $\auS$
and the same $c_P$, for $P_{\mathsf u}=c_P\auP_{\mathsf u}$, the space $\auAz^{\mathsf u}$ and
$\auT^{\mathsf u}=\varrho_{\mathsf u}(e_0)\auAz^{\mathsf u}$; and $\dim\auT^{\mathsf u}=\dim\auT$.
\end{enumerate}
Another basis of $L_{\mathsf u}$ changes the identification of $K_{\mathsf u}/J_{\mathsf u}$ with $\auG$ by
an inner automorphism of $\auG$; the central idempotent $e_0$, the space $\auT^{\mathsf u}$ and the
statements {\rm(i)--(iii)} do not depend on it.
\end{lemma}
\begin{proof}
(a) $\mathsf h_{\mathsf u}L$ is spanned by $\mathsf u^{-1}x_i$ $(i<k)$, $\mathsf u^{-1}x_k/\aubar{\pi_N}+
z/N$, $z$, $e_a$, $e_a^*$; multiplying by $\mathsf u$ gives the basis of $L_{\mathsf u}$. As
$\mathsf u\aubar{\mathsf u}=1$, the scalar $\mathsf u$ and $\mathsf h_{\mathsf u}$ are isometries of
$Y_0\otimes\auA_{E,f}$, so $L_{\mathsf u}$ is unimodular; a scalar does not change the stabiliser of a
lattice, so $K_{\mathsf u}$ is the stabiliser of $\mathsf h_{\mathsf u}L$, which is $\mathsf h_{\mathsf u}
K\mathsf h_{\mathsf u}^{-1}$, and the congruence subgroups correspond. The matrix of $k$ in the basis
$\mathsf u\,\mathsf h_{\mathsf u}(\mathsf b)$, $\mathsf b$ the basis of $L$, is the matrix of $\mathsf
h_{\mathsf u}^{-1}k\,\mathsf h_{\mathsf u}$ in $\mathsf b$, the scalar $\mathsf u$ being central.

(b) If $y$ is right invariant under $J_{\mathsf u}$, then for $j\in J$, $(\mathsf j_{\mathsf u}y)(gj)=
y(g\,\mathsf h_{\mathsf u}^{-1}\cdot\mathsf h_{\mathsf u}j\mathsf h_{\mathsf u}^{-1})=(\mathsf j_{\mathsf u}y)
(g)$. Right translation by a finite adelic element preserves the invariance under $\Gamma$ on the
left, the boundedness, the equation $\int e_\infty(a)y(ga)\,da=y(g)$ (the two translations commute)
and the inner product $\int_{\Gamma\backslash G(\auA)}\aubar yy'$; so $\mathsf j_{\mathsf u}$ is a
unitary isomorphism, with inverse $y\mapsto y(\,\cdot\,\mathsf h_{\mathsf u})$. For $k\in K_{\mathsf u}$,
$(\mathsf j_{\mathsf u}\varrho_{\mathsf u}(k)y)(g)=y(g\,\mathsf h_{\mathsf u}^{-1}k)=(\varrho(\mathsf h_{\mathsf
u}^{-1}k\,\mathsf h_{\mathsf u})\mathsf j_{\mathsf u}y)(g)$, which by (a) is the stated intertwining.
The element $\mathsf h_{\mathsf u}$ is central in $H(\auA_{E^+,f})$, so $R_{\mathsf u}\varphi$ is again
right invariant under $J_X$ and $R_{\mathsf u}$ maps $\auS$ to itself; $c\mapsto c\,\mathsf h_{\mathsf
u}^{-1}$ permutes the classes, and if $h_c\mathsf h_{\mathsf u}^{-1}=\eta h_{c'}j$ with $\eta\in H(E^+)$,
$j\in J_X$, then $(R_{\mathsf u}\varphi)(h_c)=\eta\,\varphi(h_{c'})$ with $\eta$ acting unitarily on
$W_{\lambda^X}$; so $R_{\mathsf u}$ is unitary. Now let $y\in\auAz^{\mathsf u}$ and $\varphi\in\auS$. The
element $\mathsf h_{\mathsf u}$ is a finite adele and $g_0$ is real, so they commute, and
$y(rg_0)=(\mathsf j_{\mathsf u}y)(rg_0\mathsf h_{\mathsf u})=(\mathsf j_{\mathsf u}y)(r\,\mathsf h_{\mathsf u}g_0)$.
As $\mathsf h_{\mathsf u}\in R_0(\auA)$ and the measure of $[R_0]$ is invariant under right translations
(proof of Lemma~\ref{pb:equiv}), the substitution $r\mapsto r\,\mathsf h_{\mathsf u}^{-1}$ gives
\[
\pbip y{\auP_{\mathsf u}\varphi}=\int_{[R_0]}\mathcal W_\varphi(r)\,\aubar{y(rg_0)}\,dr=\int_{[R_0]}
\mathcal W_\varphi(r\,\mathsf h_{\mathsf u}^{-1})\,\aubar{(\mathsf j_{\mathsf u}y)(rg_0)}\,dr ,
\]
and for $r=uh$, $\mathcal W_\varphi(uh\,\mathsf h_{\mathsf u}^{-1})=\psi_U(u)\,\mathsf t\bigl(h_\infty^{-1}
\varphi(h_f\mathsf h_{\mathsf u}^{-1})\bigr)=\mathcal W_{R_{\mathsf u}\varphi}(uh)$. Hence $\pbip y{\auP_{\mathsf
u}\varphi}=\pbip{\mathsf j_{\mathsf u}y}{\auP R_{\mathsf u}\varphi}=\pbip y{\mathsf j_{\mathsf u}^{-1}\auP
R_{\mathsf u}\varphi}$.

(c) By (b), $\mathsf j_{\mathsf u}$ carries $\auT^{\mathsf u}$ onto $\auT$, which gives the equality of
dimensions and (i). For a subspace $\mathsf U\subseteq\auS$, $\varrho_{\mathsf u}(e_0)P_{\mathsf u}
\pbpr_{\mathsf U}=\mathsf j_{\mathsf u}^{-1}\,\varrho(e_0)P\,\pbpr_{R_{\mathsf u}\mathsf U}\,R_{\mathsf u}$,
which has the same Hilbert--Schmidt norm as $\varrho(e_0)P\,\pbpr_{R_{\mathsf u}\mathsf U}$, and
$\dim R_{\mathsf u}\mathsf U=\dim\mathsf U$; this gives (ii). Finally $P_{\mathsf u}^*\varrho_{\mathsf u}(s)
P_{\mathsf u}=R_{\mathsf u}^*\,P^*\varrho(s)P\,R_{\mathsf u}$ has the norm of $P^*\varrho(s)P$, which gives
(iii). The last assertion of the lemma is clear, the sum in (iii) being over all of $\auG$ and
$e_0$ being central.
\end{proof}

\noindent For $\mathsf u=u_0$ the lattice $L_{u_0}$ is $\mathsf M_0$. (Compare the remark in
\textup{(B4)} of Definition~\ref{pb:bessel}: a translation of the source by a central element of
$H(\auA_{E^+,f})$ is all that distinguishes the two Poincar\'e maps.) The scalar $\mathsf u$ on the
string alone, $d=\mathsf u\cdot\mathsf h_{\mathsf u}$, also carries $L$ to $L_{\mathsf u}$; it is the product
of an element of the centre of $G(\auA_{E^+,f})$ and of $\mathsf h_{\mathsf u}$, and no rational element
with the ideal of $\mathsf u$ is needed.

\begin{theorem}\label{pb:thmB}
Assume Theorem~\ref{pb:est} and Proposition~\ref{au:mult} (for $j=m$). Then
Proposition~\ref{au:moment} holds for the pair $(g_0,\mathsf t)$ and the scalar $c_P$ of
Theorem~\ref{pb:est}, the character \eqref{pb:eq-psi}, the inner product of
\S\ref{pb:ss-setting} on $\auS$, and $\deg=\pbdeg$; and Theorem~\ref{thm:srcB} holds for the map
$\auP$ defined by these choices. Both statements hold for the lattice $L$ of this section and
for the lattice $\mathsf M_0$ of \eqref{ps:eq-lattice}, with the same $(g_0,\mathsf t)$, $c_P$,
$\psi_U$ and inner product.
\end{theorem}
\begin{proof}
For the lattice $L$: parts (i), (ii), (iii) of Proposition~\ref{au:moment} are Propositions~\ref{pb:i}, \ref{pb:ii}
and~\ref{pb:iii}, and $\audT>0$ by Proposition~\ref{pb:ii}; recall $\pbDmain=\auDX^2/(\pbdeg\cdot
\pbmass_r)$. For the lattice $\mathsf M_0=L_{u_0}$ they follow by Lemma~\ref{pb:lattice}(c), and
$\dim\auT^{u_0}=\audT>0$. Theorem~\ref{thm:srcB} follows in either case by the argument of \S\ref{ss:au-proofAB}, which uses
Proposition~\ref{au:moment} through Proposition~\ref{au:spread} and
Lemmas~\ref{au:onetranslate}, \ref{au:few} and~\ref{au:ramcount}; these are proved there and do not
depend on the choice of $(g_0,\mathsf t)$ or of $\psi_U$.
\end{proof}

The error $\frac32(8b^2)^{-1}$ in Proposition~\ref{au:spread}, hence the numbers $4/b^2$ and
$5\log(1+2b^2)$ of Theorem~\ref{thm:srcB}, are thus far from what the method gives: (ii) and (iii)
hold with $5b^{-4}$ and $9b^{-4}$, and these come from two places only, the factor $1+4b^{-4}$ of
(E3) and the weight $b^{-4}$ of the elements that are not open (the union over $\varrho$ minors of
sets of weight $3b^{-10}$), on one side, and the share $b^{-8}/2$ of the bad types on the other.
In the part of the argument that is proved here, the rank enters only through the inequalities
$11(4\varrho+1)\le k$ and $3\varrho\le b^6$ (Proposition~\ref{pb:iii}) and $1<r<p_i$,
$18\,r(22\,r\log b+1)\log b\le f$ (Lemma~\ref{pb:bad}).

\subsection{Second layer: the Poincar\'e series and the inequalities (E1)--(E3)}\label{pb:ss-layer2}

\subsubsection*{Bessel functions}
Let $H_\infty=\prod_\tau H_\tau$, $U_\infty=\prod_\tau U_\tau$, $R_\infty=U_\infty\rtimes H_\infty$.
Fix Haar measures on $G_\tau$ and $U_\tau$, hence invariant measures on $U_\tau\backslash G_\tau$;
they are invariant under left multiplication by $H_\tau$, which normalises $U_\tau$ and is
compact. Let $\pbW=\bigotimes_\tau\pbW_\tau$ be the hermitian space dual to $W_{\lambda^X}$, with
the action $\tau_\infty$ of $H_\infty$ and an orthonormal basis indexed by a set $I$, $\#I=d$, a
product of orthonormal bases of the $\pbW_\tau$. We write $\pbip w\xi=\sum_{i\in I}w_i\xi_i$ for
the pairing of $w\in\pbW$ and $\xi\in W_{\lambda^X}$, in dual bases; $\pbip{\tau_\infty(h)w}{h\xi}=
\pbip w\xi$.

\begin{definition}[Hypothesis on the Bessel functions]\label{pb:bessel}
We say that \emph{Bessel functions exist} for the data if there are a point $g_0\in G_\infty$, a
linear form $\mathsf t$ on $W_{\lambda^X}$, a constant $c_W>0$ and continuous functions
$W_\tau,W_\tau^\flat\colon G_\tau\to\pbW_\tau$ $(\tau\in\auPl)$ with the following properties, where
$W(y)=\bigotimes_\tau W_\tau(y_\tau)$ and $W^\flat(y)=\bigotimes_\tau W^\flat_\tau(y_\tau)$ for
$y\in G_\infty$.
\begin{enumerate}
\item[(B1)] $W_\tau(uhy)=\psi_\tau(u)\,\tau_\tau(h)\,W_\tau(y)$ for $u\in U_\tau$, $h\in H_\tau$,
where $\psi_\tau$ is the component of $\psi_\infty$; and the same for $W_\tau^\flat$.
\item[(B2)] $\|W_\tau\|$ is integrable and square-integrable on $U_\tau\backslash G_\tau$, and
$\int_{U_\tau\backslash G_\tau}\|W_\tau\|^2=\dim\pbW_\tau$.
\item[(B3)] For every class $c$ and $i\in I$ the series $\sum_{\gamma\in U_c\backslash\Gamma}
\bigl|W_i(\gamma_\infty g_\infty)\bigr|\,\pbone_{h_cJ}(\gamma g_f)$ is bounded on $G(\auA)$, where
$U_c=U(E^+)\cap h_cJh_c^{-1}$; and the function $P_{c,i}(g)=\sum_{\gamma\in U_c\backslash\Gamma}W_i(
\gamma_\infty g_\infty)\pbone_{h_cJ}(\gamma g_f)$ belongs to $\auAz$.
\item[(B4)] For every $\varphi\in\auS$, as functions on $G(\auA)$,
\[
\auP\varphi(g)=c_W\sum_{c\in\mathcal C}\ \sum_{\gamma\in U_c\backslash\Gamma}\pbip{W(\gamma_\infty
g_\infty)}{\varphi(h_c)}\,\pbone_{h_cJ}(\gamma g_f),
\]
where $\auP$ is defined by $g_0$, $\mathsf t$ and \eqref{pb:eq-psi}. (Replacing $\varphi(h_c)$ by
the value at $h_c$ of the translate of $\varphi$ by a fixed central element of $H(\auA_f)$ composes
$\auP$ with a unitary automorphism of $\auS$ that commutes with $\pbrpsi$, and changes nothing below.)
\item[(B5)] $W^\flat_\tau=\chi_\tau W_\tau$ with
$\chi_\tau(y)=\max\bigl(0,\min(1,2-2\ell_\tau(y)/b^{6n+2})\bigr)$, where
$\ell_\tau(y)=\inf_{s\in R_\tau}(\|sy\|_\tau+1)$ and $\|\cdot\|_\tau$ is the operator norm of
Definition~\ref{pbe:length}. So $\chi_\tau=1$ where $\ell_\tau\le\frac12b^{6n+2}$, $\chi_\tau=0$
where $\ell_\tau\ge b^{6n+2}$, and $\chi_\tau$ is continuous, $0\le\chi_\tau\le1$, and left
invariant under $R_\tau$; in particular the support of $W^\flat_\tau$ is compact modulo $R_\tau$.
\item[(B5$'$)] For each $\tau$ there is a compact set $C_\tau\subset G_\tau$ with
$\operatorname{supp}W^\flat_\tau\subseteq U_\tau C_\tau$. (This follows from (B5): $\chi_\tau(y)\ne0$ gives
$s=uh\in R_\tau$ with $\|sy\|_\tau<b^{6n+2}$, so $y\in h^{-1}u^{-1}B=u'h^{-1}B$ with $u'=h^{-1}u^{-1}h\in
U_\tau$ and $B$ the closed ball of radius $b^{6n+2}$ in $G_\tau$; take $C_\tau=H_\tau B$, compact as
$H_\tau$ is compact.)
\end{enumerate}
\end{definition}

The terms of the series in (B3), (B4) are well defined: by (B1), $W(uy)=\psi_\infty(u)W(y)$, and
$\psi_\infty(u)=\aubar{\psi_f(u)}=1$ for $u\in U_c$, because $\psi_f$ is invariant under
conjugation by $h_c\in H(\auA_f)$ and trivial on $U(\auA_f)\cap J$. By Theorem~\ref{pbe:bessel}
(with Theorem~\ref{pbe:ds}, proved in \S\S\ref{pds:sec}, \ref{pco:sec}, \ref{pbp:sec}) Bessel functions
exist: $g_0=(x_\tau)_\tau$, $\mathsf t=\bigotimes_\tau v_\tau$, and $W_\tau$ is, up to a positive constant,
$y\mapsto\int_{R_\tau}c_\tau(x_\tau^{-1}sy)\,(\psi_\tau\otimes\tau_\tau)(s)^{-1}v_\tau\,ds$, where
$c_\tau=\langle v_0,\pi_\tau(\cdot)v_0\rangle$ is the coefficient of a unit highest weight vector of the
lowest $K$-type of the discrete series (Proposition~\ref{pbe:local}); (B4) is the unfolding of the Bessel
period against the reproducing kernel, (B2) and (B3) rest on the integrability of the discrete
series with a large gap and on the lower bound \textup{(DS4)} for the period, and $\chi_\tau$ is a cutoff
in the size of $y$ modulo $R_\tau$ at the radius $b^{6n+2}$.

\begin{lemma}\label{pb:ortho}
Assume \textup{(B1)}, \textup{(B2)}. Then $\int_{U_\infty\backslash G_\infty}W_i\,\aubar{W_{i'}}=
\delta_{ii'}$ for $i,i'\in I$. Moreover $\nu:=\pbvol(U_c\backslash U_\infty)$ is finite and does
not depend on $c$.
\end{lemma}
\begin{proof}
The hermitian matrix $\mathsf M_\tau=\int_{U_\tau\backslash G_\tau}W_\tau(y)W_\tau(y)^*\,dy$
satisfies $\tau_\tau(h)\mathsf M_\tau\tau_\tau(h)^*=\mathsf M_\tau$ for $h\in H_\tau$, by (B1) and
the invariance of the measure under $y\mapsto hy$. As $\pbW_\tau$ is irreducible, $\mathsf M_\tau$
is a scalar, and its trace is $\dim\pbW_\tau$ by (B2); so $\mathsf M_\tau=1$. Take the product
over $\tau$. For the second claim: $U$ is a split unipotent group, so $U(E^+)$ is discrete and
cocompact in $U(\auA)$ and $U(E^+)U_\infty$ is dense in $U(\auA)$ (strong approximation for the
additive group, and induction on a composition series). Hence, for the open compact subgroup
$\mathsf K_c=U(\auA_f)\cap h_cJh_c^{-1}$, $U_c\backslash U_\infty\cong U(E^+)\backslash U(\auA)/
\mathsf K_c$ and $\nu\cdot\pbvol(\mathsf K_c)=\pbvol(U(E^+)\backslash U(\auA))$; and $\pbvol(
\mathsf K_c)=\pbvol(U(\auA_f)\cap J)$ because conjugation by $h_c$ preserves the Haar measure of
$U(\auA_f)$ (proof of Lemma~\ref{pb:equiv}).
\end{proof}

The explicit value of $\nu$ enters the main term (\S\ref{pbm:ss-inputs}) and the open cell
(\S\ref{pbs:ss-open}). Put $\pbvol(\widehat\auO_E)=(2/\aurd)^f$ and, with $s_{\tau,a}=\kappa_\tau^{\varrho-a}/
\sqrt N$ as in \eqref{pbe:eq-scale},
\[
\nu_*=\frac{M^{f\varrho(2k+2\varrho+1)}}{N^{kf}}\cdot\frac{\aurd^{\,f\varrho/2}}{\pbvol(\widehat\auO_E)^{\varrho(\varrho+k)}\,
\Lambda},\qquad\Lambda=\prod_\tau\prod_{a=0}^{\varrho}s_{\tau,a}^{\,2(k+2(\varrho-a))}.
\]

\begin{lemma}[The covolume]\label{pb:covol}
For the Haar measure $du$ of \textup{(Cv3)}, $\nu=\nu_*$.
\end{lemma}
\begin{proof}
This is Theorem~\ref{pbv:covol} in \S\ref{pbv:sec}: in the coordinates of (Cv3) a product of fundamental
parallelotopes of $2\varrho$ explicit lattices is a fundamental domain for $U_c$, and the product of their
covolumes is $\nu_*$. Theorem~\ref{pbm:L2} and Lemma~\ref{pbs:open-arith} use $\nu=\nu_*$ (the three forms of
the constant agree, \S\ref{pbv:ss-uses}).
\end{proof}

\subsubsection*{The Petersson formula}
For classes $c,c'$ and $s\in\auG$ the \emph{terms over} $(c,s,c')$ are the elements of
\[
\Gamma_{c,s,c'}=\{\gamma\in\Gamma:\ h_c^{-1}\gamma h_{c'}\in K\text{ and its image in }\auG\text{ is
}s\}.
\]
This set is stable under $\gamma\mapsto u\gamma u'$ for $u\in U_c$, $u'\in U_{c'}$, because $J$ is
normal in $K$. A subset $T\subseteq\Gamma$ is \emph{stable} if $U(E^+)\,T\,U(E^+)=T$. For a stable
$T$ and $\flat\in\{\ ,\flat\}$ the \emph{amplitude sum} is the matrix, indexed by $I\times I$,
\begin{equation}\label{pb:eq-amp}
A^{\flat}_T(c,s,c')=\sum_{\gamma\in U_c\backslash(T\cap\Gamma_{c,s,c'})/U_{c'}}\ \int_{(U_c\cap\gamma
U_{c'}\gamma^{-1})\backslash G_\infty}W^{\flat}(y)\,W^{\flat}(\gamma^{-1}y)^*\,dy ,
\end{equation}
whenever the sum of the integrals of the norms converges. The integrand is invariant under
$U_c\cap\gamma U_{c'}\gamma^{-1}$ and the integral depends only on the double coset of $\gamma$,
since $\psi_\infty$ is trivial on $U_c$ and on $U_{c'}$. We write $A_T$ for the sum formed with
$W$ and $A_T^\flat$ for the one formed with $W^\flat$.

\begin{proposition}[Petersson formula]\label{pb:unfold}
Assume \textup{(B1)--(B4)}, and let $P=(\pbvol(J)\nu)^{-1/2}c_W^{-1}\,\auP$. Then the sums
\eqref{pb:eq-amp} for $W$ converge absolutely for every stable $T$, and for $\varphi,\varphi'\in\auS$
and $s\in\auG$
\[
\pbip{P\varphi}{\varrho(s)P\varphi'}=\frac1\nu\sum_{c,c'\in\mathcal C}\ \sum_{i,i'\in I}\aubar{\varphi(
h_c)_i}\,\varphi'(h_{c'})_{i'}\,\aubar{A_\Gamma(c,s,c')_{i,i'}} .
\]
Consequently $\pbHS{P_s}^2=\nu^{-2}\sum_{c,c'}\|A_\Gamma(c,s,c')\|^2$, with the Euclidean norm of
the matrices.
\end{proposition}
\begin{proof}
Let $F_{c,i}(g)=W_i(g_\infty)\pbone_{h_cJ}(g_f)$, so that $P_{c,i}(g)=\sum_{\gamma\in U_c\backslash
\Gamma}F_{c,i}(\gamma g)$ and $\auP\varphi=c_W\sum_{c,i}\varphi(h_c)_iP_{c,i}$ by (B4). By (B3) the
series converges absolutely, $|P_{c,i}|\le\sum_\gamma|F_{c,i}(\gamma g)|\le\mathsf m$ for a
constant $\mathsf m$, and $P_{c,i}\in\auAz$. Let $k\in K$ be a lift of $s$. Unfolding the first
series,
\[
\pbip{P_{c,i}}{\varrho(k)P_{c',i'}}=\int_{U_c\backslash G(\auA)}\aubar{F_{c,i}(g)}\,P_{c',i'}(gk)\,dg
=\pbvol(J)\int_{U_c\backslash G_\infty}\aubar{W_i(y)}\,P_{c',i'}(y,h_ck)\,dy ;
\]
this is justified because $\int_{U_c\backslash G(\auA)}|F_{c,i}|\,\mathsf m\,dg=\mathsf m\,\pbvol(J)
\,\nu\int_{U_\infty\backslash G_\infty}|W_i|<\infty$ by (B2), and we used that $g\mapsto P_{c',i'}
(gk)$ is right invariant under $J$. Next, $\gamma h_ck\in h_{c'}J$ if and only if $\gamma^{-1}\in
\Gamma_{c,s,c'}$, so $P_{c',i'}(y,h_ck)=\sum_{\gamma\in\Gamma_{c,s,c'}/U_{c'}}W_{i'}(\gamma^{-1}y)$,
and
\[
\int_{U_c\backslash G_\infty}\aubar{W_i(y)}\sum_{\gamma\in\Gamma_{c,s,c'}/U_{c'}}W_{i'}(\gamma^{-1}y)
\,dy=\sum_{\gamma\in U_c\backslash\Gamma_{c,s,c'}/U_{c'}}\int_{(U_c\cap\gamma U_{c'}\gamma^{-1})
\backslash G_\infty}\aubar{W_i(y)}\,W_{i'}(\gamma^{-1}y)\,dy ,
\]
by grouping the $\gamma$ in orbits of $U_c$; the sum of the integrals of the absolute values is at
most $\mathsf m\,\nu\int_{U_\infty\backslash G_\infty}|W_i|$. This proves the convergence (for
$T=\Gamma$, hence for every stable $T$) and $\pbip{P_{c,i}}{\varrho(s)P_{c',i'}}=\pbvol(J)\,\aubar{A_\Gamma
(c,s,c')_{i,i'}}$. The formula follows by sesquilinearity. For the last claim, the vectors with one
non-zero coordinate $\varphi(h_c)_i=1$ form an orthonormal basis of $\auS$.
\end{proof}

\begin{proposition}[The terms in the Bessel subgroup]\label{pb:besselterms}
Assume \textup{(B1)--(B4)} and let $P$ be as in Proposition~\ref{pb:unfold}. For all $s\in\auG$ and
$\varphi,\varphi'\in\auS$,
\[
\frac1\nu\sum_{c,c',i,i'}\aubar{\varphi(h_c)_i}\,\varphi'(h_{c'})_{i'}\,\aubar{A_{R_0(E^+)}(c,s,c')_{
i,i'}}=\pbip\varphi{\pbrpsi(s)\varphi'} .
\]
Hence $\pbHS{P_s-\pbrpsi(s)}^2=\nu^{-2}\sum_{c,c'}\|A_{\Gamma\smallsetminus R_0(E^+)}(c,s,c')\|^2$.
\end{proposition}
\begin{proof}
Let $\gamma=u_0\eta\in R_0(E^+)\cap\Gamma_{c,s,c'}$ with $u_0\in U(E^+)$, $\eta\in H(E^+)$. Then
$h_c^{-1}\gamma h_{c'}=(h_c^{-1}u_0h_c)(h_c^{-1}\eta h_{c'})\in K_R$, so by Lemma~\ref{pb:KR}
$s=\bar u\bar t\in\pbR$, where $\bar u$ is the image of $h_c^{-1}u_0h_c\in K_U$ and $\bar t$ that
of $k_t:=h_c^{-1}\eta h_{c'}\in K_H$. In particular there are no such terms if $s\notin\pbR$, and
both sides vanish. Let $s=\bar u\bar t\in\pbR$ and fix $c$. For a lift $k\in K_H$ of $\bar t$, $h_ck\in H(E^+)h_{c'}
J_X$ for one class $c'$, which gives a pair $(\eta,c')$ with $h_c^{-1}\eta h_{c'}\in K_H$ of image
$\bar t$. The pair is unique: if $(\eta_2,c'')$ is another one, then $h_{c'}^{-1}\eta^{-1}\eta_2h_{
c''}\in K_H\cap J=J_X$, so $c''=c'$ and $\eta^{-1}\eta_2\in H(E^+)\cap h_{c'}J_Xh_{c'}^{-1}=1$, the
level being free. The elements $u_0$ that occur form one coset
$U_cu_0$, non-empty by strong approximation for $U$. Finally $U_cu_0\eta=u_0\eta\,U_{c'}$, because
$K$ normalises $J$. So over $(c,s,\cdot)$ there is exactly one double coset of terms in
$R_0(E^+)$, it lies over $(c,s,c')$, and $U_c\cap\gamma U_{c'}\gamma^{-1}=U_c$. By (B1),
$W(\gamma^{-1}y)=\aubar{\psi_\infty(u_0)}\,\tau_\infty(\eta)^{-1}W(y)$, so by Lemma~\ref{pb:ortho}
\[
A_{R_0(E^+)}(c,s,c')=\psi_\infty(u_0)\int_{U_c\backslash G_\infty}W(y)W(y)^*\,dy\cdot\tau_\infty(\eta)
=\nu\,\psi_\infty(u_0)\,\tau_\infty(\eta) .
\]
Here $\aubar{\psi_\infty(u_0)}=\psi_f(u_0)=\psi_U(\bar u)$. The representation $\tau_\infty$ is
the contragredient of the unitary action of $H_\infty$ on $W_{\lambda^X}$, so in dual orthonormal
bases $\aubar{\tau_\infty(\eta)_{i,i'}}$ is the matrix of $\eta$ on $W_{\lambda^X}$; and $\eta\,
\varphi'(h_{c'})=\varphi'(\eta h_{c'})=\varphi'(h_ck_t)$. So the left side is
\[
\psi_U(\bar u)\sum_c\pbip{\varphi(h_c)}{\eta\,\varphi'(h_{c'})}_{W_{\lambda^X}}=\psi_U
(\bar u)\sum_c\pbip{\varphi(h_c)}{(R(\bar t)\varphi')(h_c)}_{W_{\lambda^X}}=\pbip\varphi{\pbrpsi(s)
\varphi'} .
\]
The last claim follows from Proposition~\ref{pb:unfold} and $A_\Gamma=A_{R_0(E^+)}+A_{\Gamma
\smallsetminus R_0(E^+)}$.
\end{proof}

\subsubsection*{The two normalising constants}

\begin{lemma}[Ratio of masses]\label{pb:massratio}
For $l\ge1$,
\[
\frac{\pbmass_{l+1}}{\pbmass_l}=\Bigl(\frac{\aurd^{\,l+1/2}\;l!}{(2\pi)^{l+1}}\Bigr)^{f}\cdot L\bigl(
l+1,\chi_{E/E^+}^{\,l+1}\bigr),
\]
where $\chi_{E/E^+}$ is the quadratic character of $E/E^+$ and $L(s,\chi^j_{E/E^+})$ means
$\zeta_{E^+}(s)$ for even $j$. In particular $\bigl|\log L(l+1,\cdot)\bigr|\le f\,2^{1-l}$ for
$l\ge2$.
\end{lemma}
\begin{proof}
This is the mass formula for the genus of the unimodular lattice $\auO_E^l$ in the form of
Gan, Hanke and Yu~\cite[Propositions~4.4 and 4.5]{GHY}, $\pbmass_l=\mathsf c^{\,l}\prod_{j=1}^l|L(1-j,\chi_{E/E^+}^{\,j})|$
up to a factor independent of $l$, with all local factors equal to $1$ because $E/E^+$ is
unramified at the finite places and the lattice is unimodular; together with the functional
equations of $\zeta_{E^+}$ and $L(s,\chi_{E/E^+})$, which give $|L(1-j,\chi^j)|=\bigl(2\,\aurd^{\,j-
1/2}(j-1)!/(2\pi)^j\bigr)^f L(j,\chi^j)$ (the discriminant of $E^+$ is $\aurd^{\,f}$ and the
conductor of $\chi_{E/E^+}$ is trivial); the formula of the lemma corresponds to $\mathsf c=2^{-f}$.
The bound for $\log L$ follows from $|\log L(s,\cdot)|\le f\log\zeta(s)$. In the notation of~\cite[(4.2), (4.3), Propositions~4.4--4.5]{GHY},
$\mathrm{Mass}=2^{-lf}\,\tau(G)\,L(M)\prod_{v\in S}\lambda_v$ with $\tau(G)=2$, $L(M)=\prod_{j\le l}L(1-j,\chi^j)$ and $S=\emptyset$; hence $\mathsf c=2^{-f}$. A constant $\mathsf c=C^f$ with another
absolute $C$ would multiply the ratio by a factor $(2C)^f$, which changes nothing in the proof of Lemma~\ref{pb:sizes}. The exact constant $\mathsf c=2^{-f}$ is used in Lemma~\ref{pbm:const}\,\textup{(C0)}, where the ratio $\pbmass_r/\pbmass_k$ has to be known up to a factor $1+b^{-10}$.
\end{proof}

\begin{lemma}[Sizes]\label{pb:sizes}
$b^{11kf}\,\pbDid\le\pbDmain\le(M/b^6)^{kf}\,\pbDid$.
\end{lemma}
\begin{proof}
By Lemma~\ref{pb:KR} and \eqref{pb:eq-classes}, $\#\mathcal C=[K^\circ_k:J_X]\,\pbmass_k=\#\pbHb
\cdot(N^k-1)^f\cdot\pbmass_k$: the index of $K_H$ in $K^\circ_k$ is the number of possible last
rows modulo the $f$ primes dividing $\pi_N$. With \eqref{pb:eq-orders}, \eqref{pb:eq-DX} and
\eqref{pb:eq-D},
\[
\mathcal R:=\frac{\pbDmain}{\pbDid}=\frac{\auDX}{\#\pbR\cdot\pbdeg\cdot\pbmass_r}=\frac{d\,(N^k-1)^f
\,\pbmass_k}{M^{f\varrho(2k+2\varrho+1)}\,\pbdeg\,\pbmass_r} .
\]
\emph{The degree.} Put $\mathsf L=\sum_\tau\log\kappa(\tau)$; $0\le\mathsf L\le f\log b$. The first
$k-1$ entries of $\varsigma_\tau$ are consecutive, so the pairs $i<j\le k-1$ contribute $1$ to
$\pbdeg_\tau$, and $\log\pbdeg=\sum_\tau\sum_{(i,j)}\log(\varsigma_{\tau,i}-\varsigma_{\tau,j})-f
\sum_{l=k-1}^{r-1}\log l!$, the sum over the pairs $i<j$ with $j\ge k$. There are $m(k-1)$ such
pairs with $i\le k-1$, each with a difference between $\kappa(\tau)$ and $\kappa(\tau)+m\Delta+k\le
\kappa(\tau)(1+1/x)$, and the other pairs have a difference between $\Delta$ and $m\Delta\le\kappa(
\tau)$; there are at most $mr$ pairs. So
\[
m(k-1)\mathsf L\ \le\ \log\pbdeg+f\sum_{l=k-1}^{r-1}\log l!\ \le\ mr(\mathsf L+f/x).
\]
\emph{The masses.} By Lemma~\ref{pb:massratio}, with $\mathsf S=\sum_{l=k}^{r-1}(l+\frac12)=(m-1)(k+
\frac{m-1}2)$,
\[
\log\frac{\pbmass_r}{\pbmass_k}=f\Bigl(\mathsf S\log\aurd+\sum_{l=k}^{r-1}\log l!-\bigl(\mathsf S+
\tfrac{m-1}2\bigr)\log2\pi\Bigr)+\theta,\qquad|\theta|\le1 .
\]
\emph{The level.} By \eqref{eq:au-level}, $kf\log N=km\mathsf L+kf(m-1)\log\frac{M\aurd}{2\pi}-
\frac{15}2kf\log b+\theta'$ with $0\le\theta'\le kf\sqrt{\log b}$, and $kf\log N-1\le f\log(N^k-1)
\le kf\log N$. Also $\varrho(2k+2\varrho+1)=(m-2)k+\frac{(m-2)(m-1)}2$.

Inserting these, the factorials cancel except $(k-1)!^f$, the coefficient of $\log M$ is $kf-f
\frac{(m-2)(m-1)}2$, that of $\log\aurd$ is $-f\frac{(m-1)^2}2$, that of $\log2\pi$ is $f\frac{m(m-
1)}2$, and
\[
\log\mathcal R=\log\bigl(d\,(k-1)!^f\bigr)+kf\bigl(\log M-\tfrac{15}2\log b\bigr)+\mathrm{Err},
\]
\[
-m^2\mathsf L-\frac{mrf}x-\frac{fm^2}2\log(M\aurd)-2\ \le\ \mathrm{Err}\ \le\ m\mathsf L+kf\sqrt{
\log b}+\frac{fm^2}2\log2\pi+1 .
\]
As $m^2\le4x$, $r\le3x^5$, $k\ge x^5/2$, $\log(M\aurd)\le24\log b$ and $\log b$ is large,
$|\mathrm{Err}|\le\frac14kf\log b$ (a comparison of sizes). By \eqref{pb:eq-DX}, $1\le d\le(a+k)^{
kf}$, so $0\le\log(d\,(k-1)!^f)\le kf(\log(2xb)+\log k)\le\frac54kf\log b$. Finally $20\log b\le
\log M$. Hence $\log\mathcal R\ge kf(20-\frac{15}2-\frac14)\log b\ge11kf\log b$ and $\log\mathcal
R\le kf(\log M-\frac{15}2\log b+\frac54\log b+\frac14\log b)=kf(\log M-6\log b)$.
\end{proof}

Lemma~\ref{pb:sizes} is the place where the level $N$ is tuned: $\pbDmain$ is the square of the
dimension of the source divided by (a bound for) the dimension of the target per element of
$\auG$, $\pbDid$ is what the terms in the Bessel subgroup contribute to the second moment, and the
window~\eqref{eq:au-level} puts their ratio between two powers of $b^{kf}$.

\subsubsection*{Links, meeting places, and the three statements on amplitudes}
For $\gamma\in\Gamma$ let $Z(\gamma)=\bigl(h(\gamma f_l,f_j)\bigr)_{0\le j,l\le\varrho}\in M_n(E)$,
with leading minors $\delta_l(\gamma)$ ($0\le l\le n$), and
\[
\pbg(\gamma)=\delta_n(\gamma)\aubar{\delta_n(\gamma)}-N^2\,\delta_{n-1}(\gamma)\aubar{\delta_{n-1}
(\gamma)}\in E^+ .
\]
The \emph{link} of $\gamma$ is the subspace of $(A^\perp/A)^2$ of the pairs $([t],[\gamma^{-1}t])$
for $t\in Y_0$ orthogonal to $A$ and to $\gamma A$. Both $\pbg$ and the link are invariant under
$\gamma\mapsto u\gamma u'$, $u,u'\in U(E^+)$, so that for a subspace $\Lambda$ the set
$T_\Lambda=\{\gamma:\pbg(\gamma)\ne0,\ \text{the link of }\gamma\text{ is }\Lambda\}$ is stable. For
$\Omega\subseteq\auG$ the \emph{link moment} is
\[
\pbLM(\Omega)=\sum_{c,c'\in\mathcal C}\ \sum_{s\in\Omega}\ \sum_\Lambda\bigl\|A^\flat_{T_\Lambda}(c,
s,c')\bigr\|^2\in[0,\infty].
\]
An element $s\in\auG$ \emph{meets} at a place $\mathfrak p\in\{w,\aubar w\}$, $w$ chosen, if
$F_{\mathfrak p}\cap sF_{\mathfrak p}\ne0$; its \emph{meeting norm} $\mathsf{mn}(s)$ is the product
of $\auNm\mathfrak p$ over the places at which it meets. A \emph{scalar coset} is a set $\zeta\pbR$,
where $\zeta\in\auO_E$, $\zeta\aubar\zeta=1$, acts on each $V_w$ by its residue; $\Omega$
\emph{avoids the scalar cosets} if $\Omega\cap\zeta\pbR=\emptyset$ for all such $\zeta\ne1$.

\begin{theorem}[Amplitudes]\label{pb:amp}
Bessel functions exist for the data (Definition~\ref{pb:bessel}), and they can be chosen such
that, with $\nu$ as in Lemma~\ref{pb:ortho} and $P$ as in Proposition~\ref{pb:unfold}:
\begin{enumerate}
\item[(X)] \emph{(Expansion of the squares.)} If $\Omega$ avoids the scalar cosets and $\pbLM(
\Omega)\le C\nu^2$ with $C\ge0$, then
\[
\Bigl(\sum_{s\in\Omega}\pbHS{P_s-\pbrpsi(s)}^2\Bigr)^{1/2}\le\sqrt C+b^{-15}\sqrt{\pbDid} .
\]
\item[(L1)] \emph{(Links and meeting places.)} If $\mathsf z\ge1$ and $\mathsf{mn}(s)\ge\mathsf z$
for all $s\in\Omega$, then $\pbLM(\Omega)\le\bigl(b^{2f}/\mathsf z\bigr)^k\,\pbDmain\,\nu^2$.
\item[(L2)] \emph{(The main term.)} For every set $\mathsf A$ of chosen places above $p_0$ and every
family $\mathsf d$ of frame vectors,
$\pbLM\bigl(\{s\in\pbFix(\mathsf A,\mathsf d):s\text{ open}\}\bigr)\le(1+b^{-4})\,b^{-6k\#\mathsf A}\,
\pbDmain\,\nu^2$.
\end{enumerate}
\end{theorem}
\smallskip\noindent The existence of the Bessel functions is Theorem~\ref{pbe:bessel}; {\rm(X)} and
{\rm(L1)} are Theorem~\ref{pbe:XL1} (from Theorems~\ref{pbe:ds} and~\ref{pbe:pillars}, proved in
\S\S\ref{pds:sec}--\ref{pbo:sec}); {\rm(L2)} is Theorem~\ref{pbm:L2}, deduced in \S\ref{pbm:sec} from
Theorem~\ref{pbm:inputs}, which is proved in \S\S\ref{pfr:sec}--\ref{pfg:sec}. All three statements hold
for the same Bessel functions, those of Theorem~\ref{pbe:bessel}.\par\smallskip

\begin{lemma}\label{pb:elem}
\begin{enumerate}
\item[(a)] At most $M^f\,\#\pbR$ elements of $\auG$ lie in a scalar coset.
\item[(b)] An element of $\pbFix(\mathsf d')$, $\mathsf d'$ a frame vector in $F_{\mathfrak p}$,
meets at $\mathfrak p$. Every $s\in\Omega_i$ meets at the $2f$ places $w,\aubar w$ with $w\mid
p_i$ chosen.
\item[(c)] No element of a scalar coset is open.
\item[(d)] $\Omega_i$ avoids the scalar cosets.
\end{enumerate}
\end{lemma}
\begin{proof}
(a) $\zeta\pbR$ depends only on the residues of $\zeta$ at the $2f$ chosen places, at most $\prod_i
p_i^f=M^f$ possibilities. (b) The first claim is clear. An element of $R_0$ fixes $e_1=f_0$; so
does its adjoint inverse. Hence an element of $\pbR$ fixes $f_0$ in $V_w$ and in $V_{\aubar w}$,
and $s\in\Omega_i$ does so for $w\mid p_i$. (c) For $s\in\zeta\pbR$, $s_wf_0=\zeta f_0$, and $f_0=
e_1$ is orthogonal to the frame as $\varrho\ge1$; so the first column of $Z_w(s_w)$ vanishes and
$\delta_1(s_w)=0$ for every $w$. (d) Let $s\in\Omega_i\cap\zeta\pbR$. Comparing the action on $f_0$
modulo the chosen places $w\mid p_i$, by (b), gives $\zeta\equiv1$ modulo $w$. As $\zeta$ is an
algebraic integer of absolute value $1$ at every infinite place of the CM field $E$, it is a root
of unity (Kronecker), and a root of unity $\ne1$ that is $\equiv1$ modulo a prime above $p_i$ has
order divisible by $p_i$ (the product of the $1-\zeta^j$, $0<j<\operatorname{ord}\zeta$, is
$\operatorname{ord}\zeta$), which is impossible in $E$ because $p_i$ is unramified in $E$. So
$\zeta=1$.
\end{proof}

\begin{proposition}[From amplitudes to (E1)--(E3)]\label{pb:reduction}
Theorem~\ref{pb:amp} implies \textup{(E1)}, \textup{(E2)}, \textup{(E3)} of Theorem~\ref{pb:est},
with $c_P=(\pbvol(J)\nu)^{-1/2}c_W^{-1}$.
\end{proposition}
\begin{proof}
Write $\mathsf E(\Omega)=\sum_{s\in\Omega}\pbHS{P_s-\pbrpsi(s)}^2$.

(E1) By Lemma~\ref{pb:elem}(b),(d), $\Omega_i$ avoids the scalar cosets and $\mathsf{mn}\ge p_i^{2f}
\ge(M/b)^f$ on $\Omega_i$ (as $p_{1-i}\le b^{11}\le b\,p_i$). By (L1) and Lemma~\ref{pb:sizes},
\[
\pbLM(\Omega_i)\,\nu^{-2}\le\Bigl(\frac{b^{3f}}{M^f}\Bigr)^k\pbDmain\le\Bigl(\frac{b^{3}}{M}\Bigr)^{
kf}\Bigl(\frac M{b^6}\Bigr)^{kf}\pbDid=b^{-3kf}\pbDid\le b^{-30}\pbDid .
\]
By (X), $\sqrt{\mathsf E(\Omega_i)}\le2b^{-15}\sqrt{\pbDid}$, so $\mathsf E(\Omega_i)\le4b^{-30}\pbDid
\le b^{-12}\pbDid$.

In particular Lemma~\ref{pb:basic} applies: $\beta(s)=\pbHS{P_s}^2\le2\auDX$ for all $s$.

(E3) Let $\Omega=\{s\in\pbFix(\mathsf A,\mathsf d):s\text{ open}\}$ and $\mathsf T=b^{-6k\#\mathsf
A}\pbDmain$. By Lemma~\ref{pb:elem}(c), $\Omega$ meets no scalar coset; in particular $\Omega\cap\pbR
=\emptyset$, so $\pbrpsi$ vanishes on $\Omega$ and $\sum_{s\in\Omega}\beta(s)=\mathsf E(\Omega)$. By
(L2) and (X), $\sqrt{\mathsf E(\Omega)}\le\sqrt{(1+b^{-4})\mathsf T}+b^{-15}\sqrt{\pbDid}$. By
Lemma~\ref{pb:sizes}, $\pbDid\le b^{-11kf}\pbDmain\le\mathsf T$. So $\sqrt{\mathsf E(\Omega)}\le(1+
\frac12b^{-4}+b^{-15})\sqrt{\mathsf T}\le(1+b^{-4})\sqrt{\mathsf T}$ and $\mathsf E(\Omega)\le(1+
b^{-4})^2\mathsf T\le(1+4b^{-4})\mathsf T$.

(E2) Let $\mathsf T=b^{-6kf}\pbDmain$ and split $\pbFix(\mathrm{all},\mathsf d)$ into the elements
that lie in a scalar coset and the set $\Omega$ of the others. By Lemma~\ref{pb:elem}(a) and
$\beta\le2\auDX$, the first part contributes at most $2M^f\,\#\pbR\,\auDX=2M^f\pbDid\le2b^{23f}
b^{-11kf}\pbDmain\le\frac14\mathsf T$. The set $\Omega$ meets no scalar coset, so $\sum_{s\in\Omega}
\beta(s)=\mathsf E(\Omega)$ as before, and $\mathsf{mn}\ge p_0^f\ge b^{10f}$ on $\Omega$ by
Lemma~\ref{pb:elem}(b). By (L1), $\pbLM(\Omega)\nu^{-2}\le(b^{2f}/b^{10f})^k\pbDmain=b^{-2kf}\mathsf
T$, and by (X) and $\pbDid\le\mathsf T$,
\[
\mathsf E(\Omega)\le2\bigl(b^{-2kf}\mathsf T+b^{-30}\pbDid\bigr)\le\tfrac12\mathsf T .
\]
So the sum over $\pbFix(\mathrm{all},\mathsf d)$ is at most $\mathsf T$.
\end{proof}

\subsection{The dependency tree}\label{pb:ss-remains}

Theorem~\ref{thm:srcB} is proved above from the following statements; we list where each is proved.

\begin{enumerate}
\item[(R1)] \emph{The discrete series coefficient}, Theorem~\ref{pbe:ds}: \textup{(DS1)}--\textup{(DS3)} in
\S\ref{pds:sec}, the explicit formula \eqref{pds:eq-FJ} behind \textup{(DS2)} in Theorem~\ref{pco:FJ}, and
\textup{(DS4)} in \S\ref{pbp:sec} (Theorem~\ref{pbp:main}). The Bessel functions and the Poincar\'e series
are treated in \S\S\ref{pbe:ss-real}--\ref{pbe:ss-global} (Theorem~\ref{pbe:bessel}).
\item[(R2)] \emph{Size, gain and count}, Theorem~\ref{pbe:pillars} in the form of
Theorem~\ref{ppi:pillars}: \textup{(R)} in Proposition~\ref{ppi:cell}, \textup{(G$'$)} in
\S\ref{pbg:sec}, \textup{(C1)}, \textup{(C2)} in \S\ref{pbc:sec}, \textup{(S)} in \S\ref{pbs:sec} from the
open-cell bound \textup{(O1)} of \S\ref{pbo:sec} and Lemma~\ref{pb:covol}. With them \textup{(X)} and
\textup{(L1)} follow (Theorem~\ref{pbe:XL1}).
\item[(R4)] \emph{The main term}, Theorem~\ref{pb:amp}\textup{(L2)}: Theorem~\ref{pbm:L2}, from
Theorem~\ref{pbm:inputs} and Lemma~\ref{pb:covol}; \textup{(R)} is Theorem~\ref{pfr:R}
(\S\S\ref{pfr:sec}, \ref{pfs:sec}, \ref{prp:sec}), \textup{(F)} and \textup{(S)} are Theorems~\ref{pff:F}
and~\ref{pff:S} (\S\S\ref{pff:sec}, \ref{pfg:sec}); the constant is evaluated in Lemma~\ref{pbm:const}.
\item[(R5)] \emph{The induced character}, Theorem~\ref{pb:est}\textup{(E4)}: Theorem~\ref{pbm:E4}.
\item[(R6)] \emph{Equality of multiplicities}, Proposition~\ref{au:mult}: Proposition~\ref{pm:mult}, from
Theorem~\ref{pm:arch}, whose part (a) is Remark~\ref{pm:rem-other}(1) and part (b) is
Theorem~\ref{pbr:main}.
\end{enumerate}
Every statement in this tree is proved in this paper; the exact value of the covolume in
Lemma~\ref{pb:covol} is Theorem~\ref{pbv:covol}. (Part~(b) of Theorem~\ref{pm:arch} is proved in \S\ref{pbr:sec}.) In addition, Lemma~\ref{pb:massratio} rests on the mass formula of Gan, Hanke
and Yu~\cite[Propositions~4.4--4.5]{GHY}, and Lemma~\ref{pb:bad} on Green's description of
the characters of $\auGL_r(\auF_q)$.


%% ===== proofs/B-covolume.tex =====
% proofs/B-covolume.tex
% The covolume of the unipotent congruence groups (Lemma pb:covol), and the completion of
% some short steps of Theorem B. Labels of this file start with pbv:.
\providecommand{\auO}{\mathcal O}
\providecommand{\auPl}{\Phi}
\providecommand{\auU}{\mathrm U}
\providecommand{\auGL}{\mathrm{GL}}
\providecommand{\auA}{\mathbb A}
\providecommand{\aurd}{\mathrm{rd}_E}
\providecommand{\aubar}[1]{\overline{#1}}
\providecommand{\pbvol}{\operatorname{vol}}
\providecommand{\pbUb}{\overline U}
\providecommand{\auNm}{\operatorname{Nm}}
\providecommand{\pbg}{\mathsf g}
\providecommand{\pffOT}{\auO_T}
\providecommand{\pfrG}{\mathsf G}
\providecommand{\pfrN}{\mathsf N}
\providecommand{\pgnD}{\mathfrak d}
\providecommand{\pgnd}{\operatorname{d}}
% macros of this file (prefix \pbv)
\providecommand{\pbvO}{\widehat{\auO}}
\providecommand{\pbvL}{\mathsf L}
\providecommand{\pbvcov}{\operatorname{covol}}
\providecommand{\pbvE}{\mathsf E}

\section{The covolume of the unipotent congruence groups}\label{pbv:sec}

This section proves Lemma~\ref{pb:covol}: for the Haar measure $du$ of convention (Cv3) of
\S\ref{pb:ss-setting}, the covolume $\nu=\pbvol(U_c\backslash U_\infty)$ of Lemma~\ref{pb:ortho} equals
the number $\nu_*$ defined before Lemma~\ref{pb:covol}, for every class $c$. The computation is
elementary: the coordinates of (Cv3) (the free entries of the rows $\tilde e_a$ and the imaginary
parts of the entries $(\tilde e_a,\tilde e_a^*)$) are additive along a triangular order
(Lemma~\ref{pbv:add}); the group $U_c$ translates each coordinate, once the earlier ones vanish, by an
explicit lattice, and every point of that lattice is attained by an Eichler transformation
(Lemma~\ref{pbv:latt}); so a product of fundamental parallelotopes is a fundamental domain
(Lemma~\ref{pbv:box}), and $\nu$ is the product of the covolumes of $2\varrho$ lattices: $\varrho$
lattices of linear forms on the spaces $Y_a$ (one for each floor of the tower of Heisenberg groups
of $U$), and $\varrho$ copies, rescaled, of the lattice of skew-hermitian integers of $E$
(Lemmas~\ref{pbv:vol-a}, \ref{pbv:vol-skew}). In \S\ref{pbv:ss-uses} we check that the three forms
of the constant used in the text are the same number. In \S\ref{pbv:ss-weak} we write out in detail four short
steps of Theorem~\ref{thm:srcB}.

\emph{What is proved.} Lemma~\ref{pb:covol} is proved in full (Theorem~\ref{pbv:covol}), using two
classical facts quoted with their hypotheses: Minkowski's formula for the covolume of the ring of
integers of a number field, and the tower formula for discriminants (used only to see that
$|D_E|=D_{E^+}^2$, since $E/E^+$ is unramified at all finite places).

\subsection{Coordinates on the unipotent group}\label{pbv:ss-coord}

We use the notation of \S\ref{pb:ss-setting} and convention (Cv1): $h$ is linear in the first
variable, and $h(v,e_a^*)$ is the coefficient of $e_a$ in $v$. Let $F$ be either $E$, or $\C$ with
$Y_F=Y_\tau=Y_0\otimes_{E,\tau}\C$ for some $\tau\in\auPl$ (with the form $h_\tau(v\otimes1,w\otimes
1)=\tau(h(v,w))$). Write $\aubar x$ for the conjugation of $F$ and $F^+=\{x:\aubar x=x\}$. Let
$U(F)$ be the group of $F$-linear isometries of $Y_F$ that preserve the flag $\langle e_1\rangle
\subset\langle e_1,e_2\rangle\subset\dots\subset A$, act trivially on its graded pieces and act
trivially on $A^\perp/A$; so $U(E)=U(E^+)$ and $U(\C)=U_\tau$ in the notation of
\S\ref{pb:ss-setting}. For $0\le a\le\varrho$ put
\[
A_a=\langle e_1,\dots,e_a\rangle,\qquad Y_a=\langle e_b,e_b^*\ (b>a);\ x_1,\dots,x_k;\ z\rangle,
\qquad m_a=\dim Y_a=2(\varrho-a)+k+1 .
\]
So $A_a$ is isotropic, $A_a^\perp=A_a\oplus Y_a$, $A_{a-1}^\perp=A_a^\perp\oplus Fe_a^*$, $A_0^\perp=Y_F$,
$A_\varrho=A$, and $Y_a$ is non-degenerate. Every $u\in U(F)$ satisfies $(u-1)e_b\in A_{b-1}$ for all
$b$, preserves each $A_a$, hence each $A_a^\perp$, and $(u-1)A_\varrho^\perp\subseteq A_\varrho$.

\subsubsection*{The coordinates}
For $u\in U(F)$ and $1\le a\le\varrho$ let the \emph{row form} of $u$ at $a$ be the linear form
\[
\mathsf r_a(u)\colon Y_a\to F,\qquad \mathsf r_a(u)(v)=h(uv,e_a^*)=h((u-1)v,e_a^*),
\]
the coefficient of $e_a$ in $uv$ (the two expressions agree because $Y_a\perp e_a^*$), and let
$\mathsf y_a(u)=h(ue_a^*,e_a^*)$, the coefficient of $e_a$ in $ue_a^*$. In a basis of $Y_a$, the values
of $\mathsf r_a(u)$ on the basis vectors are the entries of the row of $e_a$ of the matrix of $u$ in the
columns of $Y_a$; these are the columns $e_b,e_b^*$ $(b>a)$, $x_i$, $z$, that is, exactly the entries
of the row of $e_a$ above the diagonal block that are not fixed by unitarity. (The other entries of
that row are: $1$ at $e_a$ and $0$ at $e_b$, $b\ne a$, $b\le a$, by unipotence; the entry
$\mathsf y_a(u)$ at $e_a^*$, whose real part is determined by the other entries through
$h(ue_a^*,ue_a^*)=0$ (see the proof of Lemma~\ref{pbv:rows}); and the entries at $e_b^*$, $b<a$, which are determined by the rows $e_{b'}$,
$b'>b$.)

At a place $\tau$ we use the normalised basis of \S\ref{pbe:ss-real}: $\tilde e_a=\lambda_{\tau,a}e_a$,
$\tilde e_a^*=\lambda_{\tau,a}^{-1}e_a^*$, $x_i$, $\tilde z=s_{\tau,\varrho}z$, where
\begin{equation}\label{pbv:eq-lambda}
\lambda_{\tau,a}=\frac{\kappa_\tau^{\,\varrho-a+1}}{\sqrt N}=s_{\tau,a-1}\quad(1\le a\le\varrho),\qquad
s_{\tau,\varrho}=\frac1{\sqrt N},
\end{equation}
with $s_{\tau,a}=\kappa_\tau^{\varrho-a}/\sqrt N$ as in \eqref{pbe:eq-scale} (so $\tilde f_a=s_{\tau,a}f_a$).
Let $v_1,\dots,v_{m_a}$ be the basis $e_b,e_b^*$ $(b>a)$, $x_1,\dots,x_k$, $z$ of $Y_a$ and $\tilde
v_j=\lambda_{\tau,j}v_j$ the corresponding normalised vectors ($\lambda_{\tau,j}=\lambda_{\tau,b}$,
$\lambda_{\tau,b}^{-1}$, $1$, $s_{\tau,\varrho}$ respectively). The coordinates of (Cv3) of $u\in U_\tau$
are
\begin{equation}\label{pbv:eq-coord}
\rho_a(u)=\bigl(h(u\tilde v_j,\tilde e_a^*)\bigr)_{j\le m_a}=\Bigl(\frac{\lambda_{\tau,j}}{\lambda_{\tau,a}}
\,\mathsf r_a(u)(v_j)\Bigr)_j\in\C^{m_a},\qquad t_a(u)=\Iml h(u\tilde e_a^*,\tilde e_a^*)=
\lambda_{\tau,a}^{-2}\Iml\mathsf y_a(u)\in\R
\end{equation}
$(1\le a\le\varrho)$, and $du_\tau$ is the measure that corresponds to Lebesgue measure under
$c_\tau=(\rho_a,t_a)_{a\le\varrho}\colon U_\tau\to V:=\prod_{a=1}^{\varrho}(\C^{m_a}\times\R)$, Lebesgue
measure on $\C$ being $d\Rel\,d\Iml$. That $c_\tau$ is a bijection is proved in
Corollary~\ref{pbv:bij} below; it is continuous (polynomial), so it is a Borel isomorphism (a
continuous injection of a Polish space is a Borel isomorphism onto its image, by the Lusin--Souslin
theorem), and $du_\tau(S)=\pbvol(c_\tau(S))$ for Borel $S$. We order the $2\varrho$ coordinates as
\begin{equation}\label{pbv:eq-order}
\rho_\varrho<t_\varrho<\rho_{\varrho-1}<t_{\varrho-1}<\dots<\rho_1<t_1 ,
\end{equation}
and we call $u$ \emph{low for} a coordinate if all coordinates of $u$ before it vanish. Since the
$\lambda$'s are non-zero reals, the vanishing of $\rho_a(u)$ means $\mathsf r_a(u)=0$, and the vanishing of
$t_a(u)$ means $\mathsf y_a(u)\in\R$; the same definitions are used over $E$, with
\emph{$\mathsf r_a=0$ and $\mathsf y_a\in E^+$} in place of $\rho_a=0$, $t_a=0$.

\begin{lemma}[Vanishing rows]\label{pbv:rows}
Let $u\in U(F)$ and $1\le a\le\varrho$.
\begin{enumerate}
\item[(i)] If $\mathsf r_d(u)=0$ and $\mathsf y_d(u)\in F^+$ for all $d$ with $a<d\le\varrho$, then
$(u-1)A_a^\perp\subseteq A_a$.
\item[(ii)] If moreover $\mathsf r_a(u)=0$, then $(u-1)A_a^\perp\subseteq A_{a-1}$, $(u-1)e_a^*\in A_a$, and
$\mathsf y_a(u)+\aubar{\mathsf y_a(u)}=0$.
\item[(iii)] If $\mathsf r_d(u)=0$ and $\mathsf y_d(u)\in F^+$ for all $d\le\varrho$, then $u=1$.
\end{enumerate}
\end{lemma}
\begin{proof}
We show by descending induction on $d$ that $(u-1)A_d^\perp\subseteq A_d$ for $a\le d\le\varrho$, and
along the way the claims of (ii). For $d=\varrho$ this is the definition of $U$. Let $a<d\le\varrho$, or
$d=a$ in the situation of (ii), and assume $(u-1)A_d^\perp\subseteq A_d$ and $\mathsf r_d(u)=0$. For $v\in
A_d^\perp=A_d\oplus Y_d$ the vector $(u-1)v$ lies in $A_d$, and its coefficient at $e_d$ is $0$: for
$v=e_b$, $b\le d$, because $(u-1)e_b\in A_{b-1}\subseteq A_{d-1}$; for $v\in Y_d$, because it is
$\mathsf r_d(u)(v)=0$. Hence $(u-1)A_d^\perp\subseteq A_{d-1}$. Put $y=(u-1)e_d^*$. For $v\in A_d^\perp$
write $uv=v+\alpha$ with $\alpha\in A_{d-1}$; as $u$ is an isometry and $\alpha\perp e_d^*$,
\[
h(e_d^*,v)=h(ue_d^*,uv)=h(e_d^*,v)+h(e_d^*,\alpha)+h(y,uv)=h(e_d^*,v)+h(y,uv),
\]
so $h(y,uv)=0$. Since $uA_d^\perp=A_d^\perp$, $y\in(A_d^\perp)^\perp=A_d$. The coefficient of $e_d$ in $y$ is
$h(y,e_d^*)=\mathsf y_d(u)$ (as $h(e_d^*,e_d^*)=0$), and, $A_d$ being isotropic,
\[
0=h(ue_d^*,ue_d^*)=h(e_d^*+y,e_d^*+y)=h(e_d^*,y)+h(y,e_d^*)=\aubar{\mathsf y_d(u)}+\mathsf y_d(u).
\]
This proves (ii) when $d=a$. If $d>a$, then $\mathsf y_d(u)\in F^+$ by hypothesis, so $\mathsf y_d(u)=0$, $y\in
A_{d-1}$, and $(u-1)A_{d-1}^\perp\subseteq A_{d-1}$ because $A_{d-1}^\perp=A_d^\perp+Fe_d^*$. This completes the
induction and proves (i). For (iii), the case $d=1$ of the induction gives $(u-1)A_0^\perp\subseteq
A_0=0$, and $A_0^\perp=Y_F$.
\end{proof}

\begin{lemma}[Additivity]\label{pbv:add}
Let $x,u\in U(F)$ and $1\le a\le\varrho$, and assume $\mathsf r_d(u)=0$ and $\mathsf y_d(u)\in F^+$ for
$a<d\le\varrho$. Then $\mathsf r_a(xu)=\mathsf r_a(x)+\mathsf r_a(u)$. If moreover $\mathsf r_a(u)=0$, then
$\mathsf y_a(xu)=\mathsf y_a(x)+\mathsf y_a(u)$. In particular, at a place $\tau$, the coordinate of
$xu$ for which $u$ is low is the sum of the coordinates of $x$ and of $u$.
\end{lemma}
\begin{proof}
Every element of $U(F)$ maps $e_b$ into $e_b+A_{b-1}$; so for $\alpha\in A_a$ the coefficient of $e_a$ in
$x\alpha$ is that in $\alpha$. Let $v\in Y_a\subseteq A_a^\perp$. By Lemma~\ref{pbv:rows}(i), $uv=v+\alpha$ with
$\alpha\in A_a$, whose coefficient at $e_a$ is $\mathsf r_a(u)(v)$; and $xuv=xv+x\alpha$. Taking
coefficients at $e_a$ gives the first claim. If $\mathsf r_a(u)=0$, then by Lemma~\ref{pbv:rows}(ii)
$ue_a^*=e_a^*+y$ with $y\in A_a$ of coefficient $\mathsf y_a(u)$ at $e_a$, and $xue_a^*=xe_a^*+xy$; take
coefficients at $e_a$. The last sentence follows from \eqref{pbv:eq-coord} and the order
\eqref{pbv:eq-order}: $u$ low for $\rho_a$ means $\mathsf r_d(u)=0$, $\mathsf y_d(u)\in\R$ for $d>a$;
low for $t_a$ means this and $\mathsf r_a(u)=0$.
\end{proof}

\begin{lemma}[Eichler transformations]\label{pbv:eichler}
Let $1\le a\le\varrho$, $w\in Y_a$ \textup{(}a vector of $Y_F$\textup{)} and $q\in F$ with $q+\aubar q=
-h(w,w)$. Then
\[
\pbvE_a(w,q)\colon v\longmapsto v+h(v,w)\,e_a-h(v,e_a)\,w+q\,h(v,e_a)\,e_a
\]
lies in $U(F)$, and: $\mathsf r_a(\pbvE_a(w,q))=h(\cdot,w)|_{Y_a}$, $\mathsf y_a(\pbvE_a(w,q))=q$, and
$\mathsf r_d(\pbvE_a(w,q))=0$, $\mathsf y_d(\pbvE_a(w,q))=0$ for $d\ne a$. Every linear form on $Y_a$ is
$h(\cdot,w)$ for exactly one $w\in Y_a$.
\end{lemma}
\begin{proof}
Write $\gamma=\pbvE_a(w,q)$, and for $v\in Y_F$ put $\alpha=h(v,w)$, $\beta=h(v,e_a)$, so $\gamma v=v+
\alpha e_a-\beta w+q\beta e_a$. Using $h(e_a,e_a)=h(e_a,w)=0$,
\[
h(\gamma v,\gamma v)-h(v,v)=(\aubar\alpha\beta+\alpha\aubar\beta)-(\aubar\beta\alpha+\beta\aubar\alpha)+
(\aubar q+q)\beta\aubar\beta+\beta\aubar\beta\,h(w,w)=0 .
\]
By polarisation ($2\ne0$ in $F$) $\gamma$ preserves $h$, so it is injective, hence an isometry of $Y_F$. For
$v\in A^\perp$ we have $\beta=0$, so $(\gamma-1)v=\alpha e_a\in A$; for $v=e_b$, $\alpha=h(e_b,w)$ vanishes
if $b\le a$ (as $w\in Y_a$), so $(\gamma-1)e_b\in Fe_a\subseteq A_{b-1}$ for $b>a$ and $=0$ for $b\le a$.
Thus $\gamma$ preserves the flag, acts trivially on its graded pieces and on $A^\perp/A$: $\gamma\in U(F)$.
For $v\in Y_a$, $\beta=0$ and the coefficient of $e_a$ in $\gamma v$ is $\alpha=h(v,w)$; and $\gamma e_a^*=e_a^*-w+
qe_a$ ($\alpha=h(e_a^*,w)=0$, $\beta=1$), with coefficient $q$ at $e_a$. Let $d\ne a$. The coefficient of $e_d$
in $(\gamma-1)v$ is $-\beta\,h(w,e_d^*)$. If $d<a$ it is $0$, as $w\in Y_a$ has no component at $e_d$. If $d>a$
and $v\in Y_d$ or $v=e_d^*$, then $\beta=h(v,e_a)=0$. So $\mathsf r_d(\gamma)=0$ and $\mathsf y_d(\gamma)=0$.
The last claim holds because $h$ is non-degenerate on $Y_a$ and conjugate-linear in the second variable.
\end{proof}

\subsection{A box lemma}\label{pbv:ss-box}

\begin{lemma}\label{pbv:box}
Let $G$ be a group, $\Gamma\subseteq G$ a subgroup, and $c_i\colon G\to V_i$ $(1\le i\le m)$ maps into abelian
groups with $c_i(1)=0$; call $g$ \emph{low for $i$} if $c_j(g)=0$ for all $j<i$. Let $L_i\subseteq V_i$ be
subgroups and $F_i\subseteq V_i$ subsets such that for every $v\in V_i$ exactly one $l\in L_i$ has $l+v\in
F_i$. Assume:
\begin{enumerate}
\item[(T)] if $\gamma\in\Gamma$ is low for $i$, then $c_i(x\gamma)=c_i(\gamma)+c_i(x)$ for all $x\in G$;
\item[(M)] if $\gamma\in\Gamma$ is low for $i$, then $c_i(\gamma)\in L_i$;
\item[(A)] for every $i$ and $l\in L_i$ there is $\gamma\in\Gamma$, low for $i$, with $c_i(\gamma)=l$;
\item[(K)] if $\delta\in\Gamma$ and $c_i(\delta)=0$ for all $i$, then $\delta=1$.
\end{enumerate}
Then for every $x\in G$ there is exactly one $\gamma\in\Gamma$ with $c_i(x\gamma)\in F_i$ for all $i$.
\end{lemma}
\begin{proof}
\emph{Existence.} Put $\gamma_0=1$. Given $\gamma_{i-1}\in\Gamma$ with $c_j(x\gamma_{i-1})\in F_j$ for $j<i$,
choose $l\in L_i$ with $l+c_i(x\gamma_{i-1})\in F_i$ and, by (A), $\delta\in\Gamma$ low for $i$ with
$c_i(\delta)=l$, and put $\gamma_i=\gamma_{i-1}\delta$. For $j\le i$, $\delta$ is low for $j$, so by (T)
$c_j(x\gamma_i)=c_j(\delta)+c_j(x\gamma_{i-1})$; this is $c_j(x\gamma_{i-1})\in F_j$ for $j<i$ (as
$c_j(\delta)=0$) and $l+c_i(x\gamma_{i-1})\in F_i$ for $j=i$. Then $\gamma=\gamma_m$ works.

\emph{Uniqueness.} Let $\gamma,\gamma'$ both work and $\delta=\gamma^{-1}\gamma'$. Suppose $c_j(\delta)=0$ for
$j<i$. By (T), $c_i(x\gamma')=c_i(\delta)+c_i(x\gamma)$, and by (M) $c_i(\delta)\in L_i$; as both
$0+c_i(x\gamma)$ and $c_i(\delta)+c_i(x\gamma)$ lie in $F_i$, uniqueness of $l$ gives $c_i(\delta)=0$. By
induction all $c_i(\delta)$ vanish, and $\delta=1$ by (K).
\end{proof}

\begin{corollary}\label{pbv:bij}
For every $\tau$, $c_\tau=(\rho_a,t_a)_a\colon U_\tau\to V$ is bijective.
\end{corollary}
\begin{proof}
Apply Lemma~\ref{pbv:box} to $G=\Gamma=U_\tau$, the coordinates in the order \eqref{pbv:eq-order}, $L_i=V_i$ and
$F_i=\{z_i\}$ for a given $z\in V$, and $x=1$; it gives exactly one $\gamma$ with $c_\tau(\gamma)=z$. (T) is
Lemma~\ref{pbv:add}; (M) is empty; (K) is Lemma~\ref{pbv:rows}(iii). (A): for $\rho_a$ and a target
$(\rho_j)_j\in\C^{m_a}$ take $\gamma=\pbvE_a(w,-\frac12h(w,w))$ with $w\in Y_a\otimes_{E,\tau}\C$ the unique vector such that $h(v_j,w)=
\lambda_{\tau,a}\lambda_{\tau,j}^{-1}\rho_j$ for all $j$ (it exists and is unique by the last claim of Lemma~\ref{pbv:eichler}, as $(v_j)$ is a basis of $Y_a$); for $t_a$ and a target
$t\in\R$ take $\gamma=\pbvE_a(0,i\lambda_{\tau,a}^2t)$. By Lemma~\ref{pbv:eichler} and \eqref{pbv:eq-coord}
these have the required coordinate, and all their rows other than the row $a$ vanish; the second has
also $\rho_a=0$. So they are low for the coordinate in question.
\end{proof}

\begin{lemma}[The measure of \textup{(Cv3)} is a Haar measure]\label{pbv:haar}
For every $\tau$ and every $x\in U_\tau$ the left translation $u\mapsto xu$ preserves $du_\tau$. So $du_\tau$
is a Haar measure of $U_\tau$ (a nilpotent Lie group, hence unimodular), and it is the Haar measure of
Lemma~\ref{prp:heis}\,\textup{(c)}: the two descriptions of $du_\tau$ in \textup{(Cv3)} agree.
\end{lemma}
\begin{proof}
Let $x,u\in U_\tau$ and $1\le a\le\varrho$. As $u$ preserves $A_a^\perp=A_a\oplus Y_a$ and $A_a$, for $v\in Y_a$
we have $uv=\bar u_av+\alpha$ with $\bar u_av\in Y_a$, $\alpha\in A_a$, where $\bar u_a$ is the map induced by
$u$ on $A_a^\perp/A_a\cong Y_a$, and the coefficient of $e_a$ in $\alpha$ is $\mathsf r_a(u)(v)$. Since $x$ maps
$e_b$ into $e_b+A_{b-1}$, the coefficient of $e_a$ in $x\alpha$ is that of $\alpha$; so
\begin{equation}\label{pbv:eq-left1}
\mathsf r_a(xu)(v)=\mathsf r_a(u)(v)+\mathsf r_a(x)(\bar u_av)\qquad(v\in Y_a).
\end{equation}
Next, $ue_a^*\in A_{a-1}^\perp=A_a\oplus Y_a\oplus\C e_a^*$, and its coefficient at $e_a^*$ is $1$, because
$h(ue_a^*,e_a)=h(e_a^*,u^{-1}e_a)=1$ ($u^{-1}e_a\in e_a+A_{a-1}$ and $A_{a-1}\perp e_a^*$). So $ue_a^*=e_a^*+y'+
\alpha'$ with $y'\in Y_a$ and $\alpha'\in A_a$ of coefficient $\mathsf y_a(u)$ at $e_a$, and taking the
coefficient of $e_a$ in $xue_a^*=xe_a^*+xy'+x\alpha'$,
\begin{equation}\label{pbv:eq-left2}
\mathsf y_a(xu)=\mathsf y_a(u)+\mathsf y_a(x)+\mathsf r_a(x)(y').
\end{equation}
We claim that $\bar u_a$ is a function of the coordinates of $u$ that precede $\rho_a$ in the order
\eqref{pbv:eq-order}, and $y'$ a function of those that precede $t_a$. By the proof of
Lemma~\ref{pbv:box} (existence, with $x=1$) and of Corollary~\ref{pbv:bij}, $u=\delta_1\delta_2\cdots\delta_{2\varrho}$,
where $\delta_i$ is the Eichler transformation $\pbvE_d(w,q)$ of the floor $d$ of the $i$-th coordinate, and
is determined by the coordinates of $u$ up to the $i$-th. For $v\in A_a^\perp$ and $d\le a$, $h(v,e_d)=0$ and
$(\pbvE_d(w,q)-1)v=h(v,w)e_d\in A_a$; so the factors of floor $d\le a$ act trivially on $A_a^\perp/A_a$, and
$\bar u_a$ is induced by the product of the factors of floor $d>a$, which come first: this is the first
claim. For $v\in Y_a$, the isometry $u$ gives $0=h(e_a^*,v)=h(ue_a^*,uv)=h(e_a^*+y'+\alpha',\bar u_av+\alpha)=
\aubar{\mathsf r_a(u)(v)}+h(y',\bar u_av)$ ($A_a$ is isotropic and orthogonal to $Y_a$, and $h(e_a^*,\alpha)$ is
the conjugate of the coefficient of $e_a$ in $\alpha$). As $\bar u_a$ is bijective and $h$ is non-degenerate
on $Y_a$, $y'$ is determined by $\bar u_a$ and $\mathsf r_a(u)$: the second claim.

By \eqref{pbv:eq-coord}, \eqref{pbv:eq-left1} and \eqref{pbv:eq-left2}, in the coordinates $c_\tau$ the map
$u\mapsto xu$ is therefore $c_i\mapsto c_i+\phi_i$, where $\phi_i$ is a Borel function of $x$ and of the coordinates
$c_j$, $j<i$ (for $\rho_a$ the scaled values of $\mathsf r_a(x)\circ\bar u_a$; for $t_a$, $t_a(x)+\lambda_{\tau,a}^{-2}
\Iml\mathsf r_a(x)(y')$). Such a map preserves Lebesgue measure of $V$ (Fubini, integrating the last
coordinate first). So $du_\tau$ is left invariant. Finally, in the normalised basis the Gram matrix
$\mathsf J_1$ is a signed permutation matrix exchanging $\tilde e_a$ and $\tilde e_a^*$, and $u^{-1}=\mathsf J_1
u^*\mathsf J_1$; hence the entries of $u^{-1}$ in the column $\tilde e_a^*$ and the rows $\tilde v_j$ are, up to
signs, conjugates and a permutation of $j$, the entries $\rho_a(u)_j$, and $\Iml(u^{-1})_{\tilde e_a,\tilde e_a^*}=
-t_a(u)$. So the Gaussian $G_\varrho$ of \S\ref{prp:ss-haar} satisfies $G_\varrho(u^{-1})=\exp(-\pi|c_\tau(u)|^2)$,
and $\int_{U_\tau}G_\varrho(u^{-1})\,du_\tau=\int_Ve^{-\pi|c|^2}dc=1$. The Haar measure of Lemma~\ref{prp:heis}\,(c)
has the same integral (Lemma~\ref{prp:heis}\,(b), and invariance under $u\mapsto u^{-1}$); two Haar measures
with the same integral of a positive function coincide.
\end{proof}

\subsection{The lattices of a class}\label{pbv:ss-lattices}

Let $L\subseteq Y_0$ be the lattice of \S\ref{pb:ss-setting}: $L=\bigoplus_{a\le\varrho}(\auO e_a\oplus\auO
e_a^*)\perp L^{\mathrm{mid}}$, where $L^{\mathrm{mid}}\subseteq X_0\perp Ez$ is spanned by $x_1,\dots,x_{k-1}$,
$x_k/\aubar{\pi_N}+z/N$ and $z$ (proof of Lemma~\ref{pb:KR}); $L$ is integral and unimodular. For a class
$c$ with representative $h_c\in H(\auA_f)$ let $L_c=Y_0\cap h_c(L\otimes\pbvO)$. As $h_c$ fixes the $e_a$,
the $e_a^*$ and $z$ and preserves $X_0$,
\[
L_c=\bigoplus_{a\le\varrho}(\auO e_a\oplus\auO e_a^*)\perp L_c^{\mathrm{mid}},\qquad L_c^{\mathrm{mid}}=
(X_0\perp Ez)\cap h_c(L^{\mathrm{mid}}\otimes\pbvO),
\]
an integral unimodular lattice. For $1\le a\le\varrho$ put $L_{c,a}=L_c\cap Y_a=\bigoplus_{b>a}(\auO e_b
\oplus\auO e_b^*)\perp L_c^{\mathrm{mid}}$; it is an orthogonal direct summand of $L_c$, hence unimodular,
and we let
\[
L_{c,a}^*=\{\ell\in\operatorname{Hom}_E(Y_a,E):\ \ell(L_{c,a})\subseteq\auO\}=\{h(\cdot,w)|_{Y_a}:\ w\in
L_{c,a}\}
\]
(the second description because $L_{c,a}$ is unimodular). Finally $\auO^-=\{\sigma\in\auO:\aubar\sigma=
-\sigma\}$.

\begin{lemma}[The congruence group]\label{pbv:Uc}
$U_c=\{\gamma\in U(E^+):\ (\gamma-1)L_c\subseteq M L_c\}$.
\end{lemma}
\begin{proof}
By definition $U_c=U(E^+)\cap h_cJh_c^{-1}$, and $h_cJh_c^{-1}$ is the principal congruence subgroup of
level $\pi_M$ of the stabiliser of $h_c(L\otimes\pbvO)$. At a place $v$ of $E^+$ not below a chosen place,
the condition is $\gamma(L_c\otimes\auO_v)=L_c\otimes\auO_v$, and $M$ is a unit at $v$. At a place $v$ below a
chosen place $w$ (above $p_0$ or $p_1$; every place of $E$ above $p_0p_1$ is a chosen place or the
conjugate of one, and $M$ has valuation $1$ at all of them), $G(E^+_v)=\auGL_r(E_w)$, and the condition
$\equiv1$ modulo $w$ at $w$ is equivalent, for an isometry, to $\equiv1$ modulo $\aubar w$ at $\aubar w$ (the
component at $\aubar w$ is the inverse adjoint); together, $(\gamma-1)(L_c\otimes\auO_v)\subseteq M(L_c
\otimes\auO_v)$. Conversely, if $(\gamma-1)L_c\subseteq ML_c$ for $\gamma\in U(E^+)$, then $\gamma L_c\subseteq L_c$
and, as $\gamma-1$ is nilpotent, $\gamma^{-1}-1=\sum_{i\ge1}(1-\gamma)^i$ also maps $L_c$ into $ML_c$; so
$\gamma$ stabilises $L_c$ and is $\equiv1$ modulo $M$, that is, $\gamma\in h_cJh_c^{-1}$.
\end{proof}

\begin{lemma}[Rows of the congruence group]\label{pbv:latt}
Let $1\le a\le\varrho$.
\begin{enumerate}
\item[(M1)] For $\gamma\in U_c$, $M^{-1}\mathsf r_a(\gamma)\in L_{c,a}^*$.
\item[(M2)] If $\gamma\in U_c$, $\mathsf r_d(\gamma)=0$ and $\mathsf y_d(\gamma)\in E^+$ for $d>a$, and $\mathsf r_a(\gamma)=0$,
then $\mathsf y_a(\gamma)\in M\auO^-$.
\item[(A1)] For every $\ell\in L_{c,a}^*$ there is $\gamma\in U_c$ with $\mathsf r_a(\gamma)=M\ell$ and $\mathsf
r_d(\gamma)=0$, $\mathsf y_d(\gamma)=0$ for all $d\ne a$.
\item[(A2)] For every $\sigma\in\auO^-$ there is $\gamma\in U_c$ with $\mathsf r_a(\gamma)=0$, $\mathsf y_a(\gamma)=M\sigma$
and $\mathsf r_d(\gamma)=0$, $\mathsf y_d(\gamma)=0$ for all $d\ne a$.
\end{enumerate}
\end{lemma}
\begin{proof}
(M1) For $v\in L_{c,a}$, $(\gamma-1)v\in ML_c$ by Lemma~\ref{pbv:Uc}, and $e_a^*\in L_c$; as $L_c$ is
integral, $\mathsf r_a(\gamma)(v)=h((\gamma-1)v,e_a^*)\in M\auO$. (M2) Similarly $\mathsf y_a(\gamma)=
h((\gamma-1)e_a^*,e_a^*)\in M\auO$ (here $h(e_a^*,e_a^*)=0$), and $\mathsf y_a(\gamma)+\aubar{\mathsf y_a
(\gamma)}=0$ by Lemma~\ref{pbv:rows}(ii) over $F=E$.

(A1), (A2) The field $E$ contains a primitive sixth root of unity $\zeta_6$ ($6\mid5040$), and $\tau_0=\zeta_6
\in\auO$ satisfies $\tau_0+\aubar{\tau_0}=\zeta_6+\zeta_6^{-1}=1$ (as $\zeta_6^2-\zeta_6+1=0$). Let
$w\in L_{c,a}$ and $\sigma\in\auO^-$, and put
\[
\gamma=\pbvE_a(Mw,q),\qquad q=-M^2\tau_0\,h(w,w)+M\sigma .
\]
As $h(w,w)\in E^+$ and $M\in\Z$, $q+\aubar q=-M^2h(w,w)=-h(Mw,Mw)$, so $\gamma\in U(E^+)$ by
Lemma~\ref{pbv:eichler}, with $\mathsf r_a(\gamma)=Mh(\cdot,w)$, $\mathsf y_a(\gamma)=q$ and $\mathsf
r_d(\gamma)=\mathsf y_d(\gamma)=0$ for $d\ne a$. Moreover
\[
M^{-1}(\gamma-1)v=h(v,w)\,e_a-h(v,e_a)\,w+\bigl(-M\tau_0h(w,w)+\sigma\bigr)h(v,e_a)\,e_a ,
\]
which lies in $L_c$ for $v\in L_c$, since $w,e_a\in L_c$, $L_c$ is integral and $-M\tau_0h(w,w)+\sigma\in
\auO$. By Lemma~\ref{pbv:Uc}, $\gamma\in U_c$. For (A1) take $w$ with $\ell=h(\cdot,w)$ and $\sigma=0$
(then $\mathsf r_a(\gamma)=M\ell$); for (A2) take $w=0$ (then $q=M\sigma$ and $\mathsf r_a(\gamma)=0$).
\end{proof}

Let $\iota_c\colon U_c\to U_\infty=\prod_\tau U_\tau$ be the diagonal embedding. By \eqref{pbv:eq-coord}, for
$\gamma\in U_c$ the coordinates of $\iota_c(\gamma)$ are $\rho_a(\gamma)=\bigl((\lambda_{\tau,j}/\lambda_{\tau,a})
\,\tau(\mathsf r_a(\gamma)(v_j))\bigr)_{\tau,j}$ and $t_a(\gamma)=\bigl(\lambda_{\tau,a}^{-2}\Iml\tau(\mathsf y_a
(\gamma))\bigr)_\tau$. Accordingly define the subgroups
\begin{align*}
\Lambda_a&=\Bigl\{\Bigl(\frac{M\lambda_{\tau,j}}{\lambda_{\tau,a}}\,\tau(\ell(v_j))\Bigr)_{\tau\in\auPl,\,j\le m_a}:\
\ell\in L_{c,a}^*\Bigr\}\subseteq\prod_\tau\C^{m_a},\\
\Lambda'_a&=\Bigl\{\Bigl(\frac M{\lambda_{\tau,a}^2}\,\Iml\tau(\sigma)\Bigr)_{\tau\in\auPl}:\ \sigma\in\auO^-
\Bigr\}\subseteq\R^{\auPl}.
\end{align*}

\begin{lemma}[An adelic covolume]\label{pbv:adelic}
Give $\auA_{E,f}$ the Haar measure with $\pbvol(\pbvO)=(2/\aurd)^f$ and $E_\infty=\prod_{\tau\in\auPl}\C$ the
Lebesgue measure $\prod_\tau d\Rel\,d\Iml$. Let $m\ge1$ and let $\widehat W\subseteq\auA_{E,f}^m$ be an open
compact subgroup. Then $E^m\cap\widehat W$, embedded in $E_\infty^m$ by $(\tau)_\tau$, is a lattice of
covolume $\pbvol(\widehat W)^{-1}$.
\end{lemma}
\begin{proof}
For $\widehat W=\pbvO^m$ the group is $\auO^m$. The image of $\auO$ in $E_\infty$ has
covolume $2^{-f}|D_E|^{1/2}$ for the Lebesgue measure (Lemma~\ref{pbc:tb-covol}; see also \cite[Chapter~I, (5.1) and Proposition~(5.2)]{Neukirch}); as
$|D_E|=\aurd^{2f}$ this is $(\aurd/2)^f=\pbvol(\pbvO)^{-1}$, and the claim follows for $\pbvO^m$. In general
put $\widehat W_0=\widehat W\cap\pbvO^m$, open of finite index in both. As $E$ is dense in $\auA_{E,f}$,
every coset of $\widehat W_0$ in $\widehat W$ (an open set) meets $E^m$, so $[E^m\cap\widehat W:E^m\cap
\widehat W_0]=[\widehat W:\widehat W_0]$, and likewise with $\pbvO^m$ in place of $\widehat W$. Hence
$E^m\cap\widehat W$ is a lattice and
\[
\pbvcov(E^m\cap\widehat W)=\pbvcov(\auO^m)\,\frac{[\pbvO^m:\widehat W_0]}{[\widehat W:\widehat W_0]}=
\pbvol(\pbvO)^{-m}\,\frac{\pbvol(\pbvO^m)}{\pbvol(\widehat W)}=\pbvol(\widehat W)^{-1}.\qedhere
\]
\end{proof}

\begin{lemma}[The lattices of the rows]\label{pbv:vol-a}
$\Lambda_a$ is a lattice in $\prod_\tau\C^{m_a}$ of covolume
$\displaystyle\Bigl(\prod_\tau\frac{M^2}{\lambda_{\tau,a}^2}\Bigr)^{m_a}\pbvol(\pbvO)^{-m_a}$.
\end{lemma}
\begin{proof}
Identify $\ell$ with $y=(\ell(v_j))_j\in E^{m_a}$. Let $B\in\auGL_{m_a}(\auA_{E,f})$ be the matrix whose
columns are the coordinates, in the basis $(v_j)$, of an $\pbvO$-basis of $L_{c,a}\otimes\pbvO$. Since
$\auO=E\cap\pbvO$, $\ell\in L_{c,a}^*$ if and only if $yB\in\pbvO^{m_a}$, so $L_{c,a}^*$ corresponds to
$E^{m_a}\cap\widehat W$ with $\widehat W=\pbvO^{m_a}B^{-1}$, and $\pbvol(\widehat W)=|\det B|_f^{-1}\pbvol(
\pbvO)^{m_a}$. One may take $B=\operatorname{diag}(1,h_cB^{\mathrm{mid}})$, where $1$ is the identity on the
coordinates $e_b,e_b^*$ and $B^{\mathrm{mid}}$ has the columns $x_1,\dots,x_{k-1}$, $x_k/\aubar{\pi_N}+z/N$,
$z$ in the basis $x_1,\dots,x_k,z$; it is triangular with diagonal $1,\dots,1,\aubar{\pi_N}^{\,-1},1$. Now
$|\det h_c|_f=1$, because $\det h_c\cdot\aubar{\det h_c}=1$ and $|\aubar x|_f=|x|_f$; and $|\aubar{\pi_N}^{\,-1}|_f
=|\operatorname N_{E/\Q}(\pi_N)|=N^f$ by the product formula ($\pi_N\aubar{\pi_N}=N$ gives $\operatorname N_{E/\Q}
(\pi_N)^2=N^{2f}$). So $\pbvol(\widehat W)=N^{-f}\pbvol(\pbvO)^{m_a}$, and by Lemma~\ref{pbv:adelic} the
image of $L_{c,a}^*$ in $E_\infty^{m_a}$ is a lattice of covolume $N^f\pbvol(\pbvO)^{-m_a}$. The map to
$\Lambda_a$ multiplies the coordinate $(\tau,j)$ by the non-zero real number $M\lambda_{\tau,j}/\lambda_{\tau,a}$,
so it multiplies volumes by $\prod_{\tau,j}(M\lambda_{\tau,j}/\lambda_{\tau,a})^2$. For each $\tau$,
$\prod_j\lambda_{\tau,j}^2=s_{\tau,\varrho}^2=N^{-1}$: the factors of $e_b$ and $e_b^*$ cancel, those of the
$x_i$ are $1$, and that of $z$ is $s_{\tau,\varrho}^2$. Hence the covolume is $\prod_\tau\bigl((M^2/
\lambda_{\tau,a}^2)^{m_a}N^{-1}\bigr)\cdot N^f\pbvol(\pbvO)^{-m_a}$, as claimed.
\end{proof}

\begin{lemma}[The skew lattices]\label{pbv:vol-skew}
$\Lambda'_a$ is a lattice in $\R^{\auPl}$ of covolume $\displaystyle\prod_\tau\frac M{\lambda_{\tau,a}^2}\cdot
|D_E|^{1/4}=\prod_\tau\frac M{\lambda_{\tau,a}^2}\cdot\aurd^{\,f/2}$.
\end{lemma}
\begin{proof}
$E$ contains $i=\zeta_4$ ($4\mid5040$), a unit with $\aubar i=-i$. So $\sigma\in\auO^-$ if and only if
$\sigma=i\alpha$ with $\alpha\in\auO\cap E^+=\auO_{E^+}$, and $\Iml\tau(i\alpha)=\epsilon_\tau\tau(\alpha)$ with
$\epsilon_\tau=\pm1$ ($\tau(i)=\pm i$, $\tau(\alpha)\in\R$). The restrictions of the $\tau\in\auPl$ to $E^+$ are
the $f$ real embeddings of $E^+$, each once. By Minkowski's theorem for the totally real field $E^+$,
$\{(\tau(\alpha))_\tau:\alpha\in\auO_{E^+}\}$ is a lattice of covolume $|D_{E^+}|^{1/2}$. The tower formula
for discriminants, $|D_E|=D_{E^+}^{2}\,\operatorname N_{E^+/\Q}(\mathfrak d_{E/E^+})$, and the triviality of the
relative discriminant $\mathfrak d_{E/E^+}$ ($E/E^+$ is unramified at all finite places,
\S\ref{ss:au-notation}) give $|D_{E^+}|^{1/2}=|D_E|^{1/4}=\aurd^{f/2}$. Scaling the coordinate $\tau$ by
$\epsilon_\tau M\lambda_{\tau,a}^{-2}$ multiplies the covolume by $\prod_\tau M\lambda_{\tau,a}^{-2}$.
\end{proof}

\subsection{The value of the covolume}\label{pbv:ss-value}

\begin{theorem}[Lemma~\ref{pb:covol}]\label{pbv:covol}
For every class $c$, and for the Haar measure $du=\prod_\tau du_\tau$ of \textup{(Cv3)},
\[
\nu=\pbvol(U_c\backslash U_\infty)=\prod_{a=1}^{\varrho}\pbvcov(\Lambda_a)\,\pbvcov(\Lambda'_a)=
\frac{M^{f\varrho(2k+2\varrho+1)}}{N^{kf}}\cdot\frac{\aurd^{\,f\varrho/2}}{\pbvol(\pbvO)^{\varrho(\varrho+k)}\,
\Lambda}=\nu_*,
\]
with $\Lambda=\prod_\tau\prod_{a=0}^{\varrho}s_{\tau,a}^{\,2(k+2(\varrho-a))}$ and $\pbvol(\pbvO)=(2/\aurd)^f$. (The constant $\Lambda$ is a positive real number and has no relation to the lattices $\Lambda_a,\Lambda'_a$.)
\end{theorem}
\begin{proof}
\emph{A fundamental domain.} Apply Lemma~\ref{pbv:box} to $G=U_\infty$, $\Gamma=\iota_c(U_c)$, the
$2\varrho$ coordinates $\rho_a=(\rho_a^{(\tau)})_\tau$, $t_a=(t_a^{(\tau)})_\tau$ in the order
\eqref{pbv:eq-order}, with values in $\prod_\tau\C^{m_a}$ and $\R^{\auPl}$, the subgroups $L_{\rho_a}=\Lambda_a$,
$L_{t_a}=\Lambda'_a$, and for $F_i$ a half-open fundamental parallelotope of the lattice $L_i$ (spanned by
a $\Z$-basis; Lemmas~\ref{pbv:vol-a} and \ref{pbv:vol-skew}), which meets every coset of $L_i$ exactly once.
Here: (T) is Lemma~\ref{pbv:add}, applied at each $\tau$ (for $\gamma\in U_c$, being low for a coordinate
at all $\tau$ is the same as the corresponding condition over $E$: $\mathsf r_d(\gamma)=0$ and $\mathsf
y_d(\gamma)\in E^+$, as $\Iml\tau(y)=0$ for all $\tau$ means $y\in E^+$). (M) is Lemma~\ref{pbv:latt}(M1)
for $\rho_a$ (it needs no lowness) and (M2) for $t_a$. (A) is Lemma~\ref{pbv:latt}(A1), (A2): the elements
given there have vanishing rows $d\ne a$, and the element of (A2) has also $\mathsf r_a=0$, so they are low
for $\rho_a$, respectively $t_a$. (K) is Lemma~\ref{pbv:rows}(iii) over $E$ (and $\iota_c$ is injective).
Hence $B=c^{-1}(\prod_iF_i)$, $c=(c_\tau)_\tau$, meets every coset $x\Gamma$ in exactly one point.

\emph{The volume.} By Corollary~\ref{pbv:bij} and the definition of $du$, $du(B)=\prod_i\pbvol(F_i)=
\prod_a\pbvcov(\Lambda_a)\pbvcov(\Lambda'_a)$. The group $U_\infty$ is nilpotent, hence unimodular, and $du$ is a
Haar measure (Lemma~\ref{pbv:haar}), so it is invariant under $u\mapsto u^{-1}$; $B^{-1}$ meets every coset $\Gamma x$ exactly
once, and $\nu=\pbvol(U_c\backslash U_\infty)=du(B^{-1})=du(B)$. (This also shows again that $\nu$ does not
depend on $c$.)

\emph{The product.} By Lemmas~\ref{pbv:vol-a} and \ref{pbv:vol-skew},
\[
\nu=\prod_{a=1}^{\varrho}\Bigl[\Bigl(\prod_\tau\frac{M^2}{\lambda_{\tau,a}^2}\Bigr)^{m_a}\pbvol(\pbvO)^{-m_a}\cdot
\prod_\tau\frac M{\lambda_{\tau,a}^2}\cdot\aurd^{\,f/2}\Bigr].
\]
The power of $M$ is $f\sum_{a=1}^{\varrho}(2m_a+1)=f\sum_{a=1}^\varrho\bigl(2k+3+4(\varrho-a)\bigr)=f\varrho(2k+2\varrho
+1)$; that of $\pbvol(\pbvO)$ is $-\sum_am_a=-\varrho(k+1)-\varrho(\varrho-1)=-\varrho(\varrho+k)$; that of $\aurd$ is
$f\varrho/2$. For the scales, $\lambda_{\tau,a}=s_{\tau,a-1}$ and $m_a+1=k+2(\varrho-(a-1))$, so
\[
\prod_{a=1}^{\varrho}\lambda_{\tau,a}^{-2(m_a+1)}=\prod_{a'=0}^{\varrho-1}s_{\tau,a'}^{-2(k+2(\varrho-a'))}=
s_{\tau,\varrho}^{2k}\prod_{a'=0}^{\varrho}s_{\tau,a'}^{-2(k+2(\varrho-a'))}=N^{-k}\prod_{a'=0}^{\varrho}
s_{\tau,a'}^{-2(k+2(\varrho-a'))},
\]
and the product over $\tau$ is $N^{-kf}\Lambda^{-1}$. Collecting the factors gives $\nu_*$.
\end{proof}

\begin{remark}[Structure behind the computation]\label{pbv:rem-tower}
The coordinates $(\rho_a,t_a)$ are those of the tower $U=N_1\ltimes(N_2\ltimes(\dots\ltimes N_\varrho))$, where
$N_a$ is the Heisenberg group of the floor $a$ (the unipotent radical of the stabiliser of $e_a$ in
$\auU(A_{a-1}^\perp/A_{a-1})$), with centre the skew line $\{e_a^*\mapsto e_a^*+te_a\}$ and abelian quotient
$Y_a$; by Lemma~\ref{pbv:rows}(i) the elements low for $\rho_a$ act trivially on $A_a^\perp/A_a$, that is, they lie in $N_1\cdots N_a$, and
Lemma~\ref{pbv:eichler} gives the elements of $N_a$. Thus $\nu$ is the product over the floors of the
covolume of a lattice in the vector space $Y_a\otimes_{E^+}\R$ (the lattice $M L_{c,a}^*$, which carries the
level $M$, the scales $\lambda_{\tau,j}/\lambda_{\tau,a}$ of the normalised basis and the factor $N^f$ of the
non-trivial basis of $L^{\mathrm{mid}}$ at $N$) and of the covolume of the skew-hermitian lattice $M\auO^-$
of the centre (which carries the different through $|D_E|^{1/4}$, i.e.\ through $\varpi$: $|D_E|=\prod_\tau
|\tau(\varpi)|^2$). We did not need the tower explicitly.
\end{remark}

\subsection{The three forms of the constant}\label{pbv:ss-uses}

The constant appears in three places, written in three ways; they are the same number, so that
Theorem~\ref{pbv:covol} gives $\nu$ in each of them.
\begin{enumerate}
\item[(1)] Before Lemma~\ref{pb:covol}: $\nu_*=M^{f\varrho(2k+2\varrho+1)}N^{-kf}\aurd^{\,f\varrho/2}\pbvol(\pbvO)^{-\varrho
(\varrho+k)}\Lambda^{-1}$ with $\Lambda=\prod_\tau\prod_{a=0}^\varrho s_{\tau,a}^{2(k+2(\varrho-a))}$, $s_{\tau,a}=
\kappa_\tau^{\varrho-a}/\sqrt N$, $\pbvol(\pbvO)=(2/\aurd)^f$. This is the value of Theorem~\ref{pbv:covol}.
\item[(2)] In \eqref{pbm:eq-consts} (used in Theorem~\ref{pbm:L2} and Lemma~\ref{pbm:const}):
$\nu_*=\#\pbUb\,N^{-kf}\aurd^{\,f\varrho/2}\mathsf v^{-\varrho(\varrho+k)}\Lambda_*^{-1}$ with $\mathsf v=(2/\aurd)^f$ and
$\Lambda_*=\prod_\tau\prod_as_{\tau,a}^{k+2(\varrho-a)}$, where in \S\ref{pbm:sec} the letter $s_{\tau,a}$ denotes
$\kappa_\tau^{2(\varrho-a)}/N$, the \emph{square} of the scale \eqref{pbe:eq-scale}. So $\Lambda_*=\Lambda$, and
$\#\pbUb=M^{f\varrho(2k+2\varrho+1)}$ by \eqref{pb:eq-orders}: the two expressions coincide.
\item[(3)] In Remark~\ref{pbs:rem-nu} (used in Lemmas~\ref{pbs:O2} and \ref{pbs:open-arith}):
$\nu_{\mathrm{src}}=(M^{2d+\varrho}/N^k)^f\pbvol(\pbvO)^{-d}|D_E|^{\varrho/4}/\Lambda$ with $d=\varrho(\varrho+k)$ and
$\Lambda=\prod_\tau\prod_a(\kappa_\tau^{2(\varrho-a)}/N)^{k+2(\varrho-a)}$, which is the $\Lambda$ of (1). Here
$2d+\varrho=\varrho(2k+2\varrho+1)$ and $|D_E|^{\varrho/4}=\aurd^{\,2f\varrho/4}=\aurd^{\,f\varrho/2}$, so $\nu_{\mathrm{src}}=\nu_*$; and
the identity used in Step~2 of Lemma~\ref{pbs:O2},
\[
\nu_{\mathrm{src}}\Lambda=\#\pbUb\,N^{-kf}\pbvol(\pbvO)^{-\varrho(\varrho+k)}|D_E|^{\varrho/4},
\]
is the definition multiplied by $\Lambda$.
\end{enumerate}
Thus $\nu=\nu_*=\nu_{\mathrm{src}}$ for the measure (Cv3), which is the measure for which the local inputs
(Theorems~\ref{pfr:R} and \ref{pbo:O1}) are proved. In the factorisation of Theorem~\ref{pbv:covol},
$\#\pbUb=\prod_aM^{f(2m_a+1)}$ collects the level $M$ in the $2m_a+1$ real dimensions of each floor, $N^{-kf}$ comes from the scale
$s_{\tau,\varrho}^{2k}$ of $\tilde z$ left over in $\Lambda$ (the factor $N^f$ of $L^{\mathrm{mid}}$ at the
primes above $N$ cancels the scale $s_{\tau,\varrho}^2$ of $\tilde z$ in each row), $\pbvol(\pbvO)^{-\varrho(\varrho+k)}$
from the $\sum_am_a$ coordinates in $E$, and $|D_E|^{\varrho/4}$ from the $\varrho$ skew coordinates.

\subsection{Four short steps of Theorem~\ref{thm:srcB}}\label{pbv:ss-weak}

\subsubsection*{(i) The rearrangement of the Heisenberg measures (Lemma~\ref{pfr:rearr})}
The last step of the proof of Lemma~\ref{pfr:rearr} is the identity
\begin{equation}\label{pbv:eq-rearr}
\eta_j*(\eta_0*\dots*\eta_{j-1})=\eta_0*\dots*\eta_{j-1}*\eta_j ,
\end{equation}
where $(\mu*\mu')(\Phi)=\iint\Phi(gg')\,d\mu(g)\,d\mu'(g')$. It follows from (N3), without any statement
about Haar measures of products. Put $\nu'=\eta_0*\dots*\eta_{j-1}$. For $n\in\pfrN_j\subseteq\pfrG^{(j)}$
and $a<j$, conjugation by $n$ maps $\pfrN_a$ to itself by $n\,n_{\omega,\beta}\,n^{-1}=n_{n\omega,\beta}$ and
preserves $\eta_a$: in the coordinates $(\omega,\beta)$ it is $\omega\mapsto n\omega$, a linear automorphism of the
$\omega$-space (the part of $V_{a+1}$ orthogonal to $f_a,f_a^*$, which $n$ preserves) whose modulus is
$|\det n|$, equal to $1$ because $\det n\cdot\aubar{\det n}=1$, with $\beta$ unchanged. Since $x\mapsto nxn^{-1}$ is
a group automorphism, it maps $\nu'$, the image of $\eta_0\otimes\dots\otimes\eta_{j-1}$ under the product
map, to itself. Hence, for $\Phi\ge0$ measurable, by Tonelli (all measures are $\sigma$-finite),
\[
\begin{aligned}\int\Phi\,d(\eta_j*\nu')&=\iint\Phi(nx)\,d\nu'(x)\,d\eta_j(n)=\iint\Phi\bigl((nxn^{-1})\,n\bigr)\,d\nu'(x)\,
d\eta_j(n)\\&=\iint\Phi(xn)\,d\nu'(x)\,d\eta_j(n)=\int\Phi\,d(\nu'*\eta_j),\end{aligned}
\]
which is \eqref{pbv:eq-rearr} (associativity of convolution gives the right side).

\subsubsection*{(ii) The basis \textup{(L1)} at $T$ (Lemma~\ref{pff:L1})}
The explicit basis used in the proof of Lemma~\ref{pff:L1} is as follows (forms linear in the first variable, (Cv1); $c,\tau\in\pffOT$ with
$c\aubar c=-1$, $\tau+\aubar\tau=1$ as in the proof; note $\aubar c^{\,-1}=-c$).
\begin{itemize}
\item From $v=\aubar{\pi_N}^{\,-1}x_{k_0}+N^{-1}z$ and $z$ ($h(v,v)=0$, $h(v,z)=-1$, $h(z,z)=-N$): the pair
$p_0=-v$, $q_0=z-N\tau v$, with $h(p_0,q_0)=1$, $h(q_0,q_0)=-N+N\aubar\tau+N\tau=0$; and
$\bigl(\begin{smallmatrix}p_0&q_0\end{smallmatrix}\bigr)=\bigl(\begin{smallmatrix}v&z\end{smallmatrix}\bigr)
\bigl(\begin{smallmatrix}-1&-N\tau\\0&1\end{smallmatrix}\bigr)$, determinant $-1$; $z=q_0-N\tau p_0$.
\item From two orthonormal vectors $x,x'$: $p=x+cx'$, $q=\aubar\tau x-c\tau x'$. Then $h(p,p)=1+c\aubar c=0$,
$h(q,q)=\tau\aubar\tau+c\aubar c\,\tau\aubar\tau=0$, $h(p,q)=\tau-c\aubar c\,\aubar\tau=\tau+\aubar\tau=1$; and
$\bigl(\begin{smallmatrix}p&q\end{smallmatrix}\bigr)=\bigl(\begin{smallmatrix}x&x'\end{smallmatrix}\bigr)
\bigl(\begin{smallmatrix}1&\aubar\tau\\c&-c\tau\end{smallmatrix}\bigr)$, determinant $-c(\tau+\aubar\tau)=-c$, a unit.
\end{itemize}
Take $m'+1$ disjoint pairs $(x_{i},x_{i'})$ among the $x_i$, $i\ne k_0$, and keep the other $\iota$ of them.
The basis $e_1,\dots,e_\varrho;\ p_0,q_0;\ p_1,q_1,\dots,p_{m'+1},q_{m'+1};\ x_i\ (\iota\text{ of them});\
e_1^*,\dots,e_\varrho^*$ has the Gram matrix with blocks $\bigl(\begin{smallmatrix}0&1\\1&0\end{smallmatrix}\bigr)$ on
each pair $(e_a,e_a^*)$, $(p_l,q_l)$ and $1$ on the kept $x_i$; the change of basis from the lattice basis is
block diagonal with the blocks above and identity blocks, integral with determinant $\pm\prod_l c\in
\pffOT^\times$, so its inverse is integral. This is the basis of (L1) (the frame vector $z$ lies in the span
of the pair $p_0,q_0$ with integral coefficients, as displayed; $k+1=\iota+2(m'+1)+2$).

\subsubsection*{(iii) The open cell}
The three steps of \S\ref{pbo:sec} used here are proved there: Lemma~\ref{pbo:tube} rests on the gauge $N(u)$ with $\sqrt{N(uu')}\le\sqrt{N(u)}+\sqrt{N(u')}$ and $N(u^{-1})=N(u)$, the scaling of the
balls $B(\rho)$ under $a_t$, and the bound $(1+1/Q)^Q\le e\le3$; Lemma~\ref{pbo:sobolev} on
the one-variable inequality, the recursion for $I_j(w)$ over the letters, and the binomial sum; and
Lemma~\ref{pbo:sizes}(c) contains the explicit polynomial comparison $\|M(Z,W)\|\le15P^6\le P^{12}$ for $P\ge8$.
We add one remark on (c): the Cauchy estimate is applied on the product of the two
operator-norm balls of radius $T^{-6}$ about $Z$ and about $Z^*$ (on which $M$ is holomorphic in $(Z,W)$ and
bounded by $(2T)^{12}$ by (b)), in one complex variable along each direction $H$ with $\|H\|=1$; the real
derivative of $Z\mapsto M(Z,Z^*)$ in the direction $H$ is $\partial_ZM[H]+\partial_WM[H^*]$, each term of norm
at most $(2T)^{12}T^6$, which gives the bound $2(2T)^{12}T^6\le T^{25}$ stated there.

\subsubsection*{(iv) The convention for $h$ in \S\ref{pbg:sec}}
The sentence at the beginning of \S\ref{pbg:sec} is correct, with the following argument. The proofs of
\S\ref{pbg:sec} do use forms linear in the \emph{second} variable: the expansion $x=\sum_kR^\#_kh(R_k,x)$ in the
proof of Lemma~\ref{pbg:jacobi} and the formula $\mathfrak p=\aubar{\mathfrak a(L)}\mathfrak a(L')\det q(\ell,\ell')$
in Lemma~\ref{pbg:basic}(c) require it. Given the form $h$ of (Cv1), put $h'(x,y)=h(y,x)=\aubar{h(x,y)}$, which is
linear in the second variable. Then: orthogonality, integrality ($h(x,L)\subseteq\auO$), dual lattices and
unimodularity are the same for $h$ and $h'$, since $\aubar\auO=\auO$; for a tuple $\ell$, $h'(\ell,\ell)$ is the
transpose of $h(\ell,\ell)$, so $\pgnD_{h'}(L)=\pgnD_h(L)$ and $\pgnd_{h'}=\pgnd_h$; mixed determinants satisfy
$\det h'(\ell,\ell')=\aubar{\det h(\ell,\ell')}$, so the ideal $\mathfrak p$ of Lemma~\ref{pbg:basic}(c) is replaced by
$\aubar{\mathfrak p}$, of the same norm, and $\mathfrak p\subseteq\auO$ iff $\aubar{\mathfrak p}\subseteq\auO$; positivity at the
complex places is unchanged; the matrix $Z(\gamma)$ is replaced by its entrywise conjugate, so the minors
$\delta_l$ by $\aubar{\delta_l}$, and $\pbg(\gamma)$, the norms $\auNm\delta_l$ and all indices are unchanged; in
Lemma~\ref{pbg:jacobi} the form $\tilde h$ built from $(h',c)$ is the conjugate of the one built from $(h,\aubar c)$,
and $\auNm c=\auNm\aubar c$ (the lemma is applied with $c=N\in\Z$). Hence every statement of \S\ref{pbg:sec},
being about lattices, indices and absolute norms, holds verbatim for $h$ in convention (Cv1).


%% ===== proofs/B-bessel-expansion.tex =====
% proofs/B-bessel-expansion.tex
% Bessel functions, the Poincare series, and the expansion of the squares:
% Hypothesis pb:bessel and Theorem pb:amp (X), (L1) of proofs/B-second-moment.tex.
% Labels of this file start with pbe:.
\providecommand{\auO}{\mathcal O}
\providecommand{\auPl}{\Phi}
\providecommand{\auPa}{\Phi_a}
\providecommand{\auS}{\mathcal S}
\providecommand{\auAz}{\mathcal A_0}
\providecommand{\auG}{\overline G}
\providecommand{\auP}{\mathcal P}
\providecommand{\auDX}{D_X}
\providecommand{\auNm}{\operatorname{Nm}}
\providecommand{\auU}{\mathrm U}
\providecommand{\auGL}{\mathrm{GL}}
\providecommand{\auA}{\mathbb A}
\providecommand{\auF}{\mathbb F}
\providecommand{\aurd}{\mathrm{rd}_E}
\providecommand{\aubar}[1]{\overline{#1}}
\providecommand{\auf}[1]{\texttt{#1}}
\providecommand{\panot}[1]{\par\smallskip{\raggedright\noindent\textbf{[Not proved here:} #1\textbf{]}\par}\smallskip}
\providecommand{\pbR}{\overline R_0}
\providecommand{\pbUb}{\overline U}
\providecommand{\pbHb}{\overline H}
\providecommand{\pbvol}{\operatorname{vol}}
\providecommand{\pbDid}{\mathcal D_{\mathrm{id}}}
\providecommand{\pbDmain}{\mathcal D_{\mathrm{main}}}
\providecommand{\pbdeg}{\operatorname{deg}}
\providecommand{\pbrpsi}{r_\psi}
\providecommand{\pbW}{\mathbb W}
\providecommand{\pbg}{\mathsf g}
\providecommand{\pbLM}{\mathrm{LM}}
\providecommand{\pbone}{\mathbf 1}
\providecommand{\pbHS}[1]{\|#1\|}
\providecommand{\pbip}[2]{\langle#1,#2\rangle}
% macros of this file (prefix \pbe)
\providecommand{\pbel}{\ell}
\providecommand{\pbels}{\tilde\ell}
\providecommand{\pbeL}{\mathsf L}
\providecommand{\pbeq}{\mathsf q}
\providecommand{\pbeB}{\mathsf B}
\providecommand{\pbeN}{\mathcal N}
\providecommand{\pbet}{\beta_*}
\providecommand{\pbenp}{\mathsf n_{\mathrm p}}
\providecommand{\pbesz}{\operatorname{size}}
\providecommand{\pbeGs}{\Gamma^*}
\providecommand{\pbeAp}{A^{+}}
\providecommand{\pbeCL}{\mathrm{CL}}

\section{Bessel functions, Poincar\'e series and the expansion of the squares}\label{pbe:sec}

This section concerns three statements of \S\ref{pb:sec} that were left without proof there:
Hypothesis~\ref{pb:bessel} (Bessel functions exist and $\auP$ is their Poincar\'e series) and parts
\textup{(X)} and \textup{(L1)} of Theorem~\ref{pb:amp}. We use the notation of \S\ref{pb:sec}:
$n=\varrho+1$, $r=k+2\varrho+1$, the frame $(f_0,\dots,f_\varrho)=(e_1,\dots,e_\varrho,z)$,
$A=\langle e_1,\dots,e_\varrho\rangle$, $U$, $H$, $R_0=U\rtimes H$, the classes $c\in\mathcal C$ with
representatives $h_c\in H(\auA_f)$, the groups $U_c=U(E^+)\cap h_cJh_c^{-1}$, the terms
$\Gamma_{c,s,c'}$ and the amplitude sums \eqref{pb:eq-amp}. Inner products are conjugate-linear in the
first variable.

\emph{What is proved.}
\begin{itemize}
\item[(1)] \emph{Bessel functions} (\S\S\ref{pbe:ss-real}--\ref{pbe:ss-global}).
Hypothesis~\ref{pb:bessel} is proved, with \textup{(B5)} in its exact form and with the
quantitative properties that the later estimates use (Theorem~\ref{pbe:bessel}), \emph{from four
properties \textup{(DS1)--(DS4)} of the discrete series coefficient at a real place}
(Theorem~\ref{pbe:ds}). \textup{(DS1)} is Harish-Chandra's theory; \textup{(DS2)}--\textup{(DS4)}
are quantitative statements about one explicit function (Flensted-Jensen's coefficient of the lowest
$K$-type), namely its decay, the size of the formal degree against the volume of balls, and a lower
bound for one Bessel period. \textup{(DS1)}--\textup{(DS3)} are proved in \S\ref{pds:sec}, the explicit
formula behind \textup{(DS2)} in Theorem~\ref{pco:FJ}, and \textup{(DS4)} in \S\ref{pbp:sec}. Given
them, everything is proved here: convergence of the local period, equivariance, the reproducing
property, the identity ``the norm is a period'', the decay below an invariant majorant, the cost
of the cut in $L^2$, the count of the terms of a Poincar\'e series (Proposition~\ref{pbe:count}),
its boundedness, and the identification with $\auP$.
\item[(2)] \emph{Expansion of the squares, first two steps} (\S\S\ref{pbe:ss-char}--\ref{pbe:ss-cut}).
Vanishing by character (Lemma~\ref{pbe:vanish}), the structure of the agreeing terms that preserve
$A$ (Lemma~\ref{pbe:scalar}), and the cost of the cut (Proposition~\ref{pbe:cut}) are proved in
full, the last one given (1).
\item[(3)] \emph{Pairs of terms} (\S\S\ref{pbe:ss-pillars}--\ref{pbe:ss-XL1}).
The third step of \textup{(X)} and the whole of \textup{(L1)} are \emph{reduced}, by complete
arguments, to three estimates for the terms with $\pbg\ne0$ or with $\gamma A\ne A$: a bound for one
block by sizes of orbits \textup{(S)}, a bound for the discriminant of the lattice common to two
terms \textup{(G)}, and a count of pairs by their common lattice \textup{(C)}, together with a structural property \textup{(R)}
(Theorem~\ref{pbe:pillars}). These estimates are proved in \S\S\ref{ppi:sec}, \ref{pbg:sec},
\ref{pbc:sec}, \ref{pbs:sec}, \ref{pbo:sec}: (R) in Proposition~\ref{ppi:cell}, (G) from (G$'$) of
\S\ref{pbg:sec}, (C) from (C1), (C2) of \S\ref{pbc:sec}, and (S) in \S\ref{pbs:sec} from the open-cell
bound of \S\ref{pbo:sec}. The linear algebra
of links that the reduction needs is proved (Lemma~\ref{pbe:link}).
\end{itemize}
\S\ref{pbe:ss-remains} lists the statements on which this section rests.

\subsubsection*{Sizes}
Besides \eqref{eq:au-scales} and \eqref{eq:au-sizes} we use the following consequences of the
definition of the data, all of them in Lemma~\ref{an:lem-sizes}(S7): $n=m/2\ge23$ and $12n^3+500\le k$
((S7)(a)); for Theorem~\ref{pbe:pillars}, $\varrho\ge10^5$ and $k\ge1000(2n)^9$ \eqref{ppi:eq-sizes}, and
for \S\ref{pbs:ss-open}, $(200fk\,\aurd)^{1000}\le b$ ((S7)(d)); $r/2\le k<r$ (S2), so that $(r+k)^2\le
4r^2$, and $r\ge100$, $3r^3f\le r^7$, $kf\ge40$ ((S7)(b)); $n,r,f,\sqrt d\le b$, and $\sqrt d\le\aurd$
($\sqrt d$ is the root discriminant of $K\subseteq E$) ((S7)(e)); $b^{20}\le M\le b^{22}$ (S5);
$N\le b^{60n}$; $d_\infty:=\dim W_{\lambda^X}\le b^{2kf}$ ((S7)(f)); and, for \S\ref{pbo:sec},
$8r\le b$ and $(r+1)^4(L_\tau+r+1+2n^2)\le b^5$ with $L_\tau=\max_i|2\varsigma_{\tau,i}|\le b^2$
((S7)(b), (g)). One more inequality concerns the generator $\varpi$ of the
different in \eqref{pb:eq-psi}: for every $\tau\in\auPl$ the real number
\begin{equation}\label{pbe:eq-kappa}
\kappa_\tau:=\frac{M\,\tau(\varpi)}{2\pi i}\qquad\text{satisfies}\qquad1\le|\kappa_\tau|\le b^{23}.
\end{equation}
(It is real because $\varpi$ is purely imaginary. The inequality is false for some purely imaginary generators; it holds for the generator $\varpi=\gamma\Phi'_{q_1}(\zeta)\sqrt{-d}$ of Lemma~\ref{pds:kappa}, which is the one fixed in \eqref{pb:eq-psi}; the two conditions on the data used there are in Lemma~\ref{an:lem-sizes}(S7).)

\subsection{One real place: norms, lengths and the cut}\label{pbe:ss-real}

Fix $\tau\in\auPl$. Let $Y_\tau=Y_0\otimes_{E,\tau}\C$ and $G_\tau=\auU(Y_\tau)\cong\auU(k+\varrho,
\varrho+1)$. The \emph{normalised basis} of $Y_\tau$ is
\[
\tilde e_a=\frac{\kappa_\tau^{\,\varrho-a+1}}{\sqrt N}\,e_a,\qquad x_1,\dots,x_k,\qquad\tilde z=
\frac z{\sqrt N},\qquad\tilde e_a^*=\frac{\sqrt N}{\kappa_\tau^{\,\varrho-a+1}}\,e_a^*\qquad(1\le
a\le\varrho).
\]
In it the hermitian form has the matrix $\mathsf J_1$ with $h(\tilde e_a,\tilde e_b^*)=\delta_{ab}$,
$h(x_i,x_j)=\delta_{ij}$, $h(\tilde z,\tilde z)=-1$ and all other products zero; $\mathsf J_1$ is
hermitian and unitary. The \emph{normalised frame} is $(\tilde f_0,\dots,\tilde f_\varrho)=(\tilde
e_1,\dots,\tilde e_\varrho,\tilde z)$, so that $\tilde f_a=s_a f_a$ with
\begin{equation}\label{pbe:eq-scale}
s_a=\kappa_\tau^{\,\varrho-a}/\sqrt N\quad(0\le a\le\varrho),\qquad|s_a/s_{a'}|=|\kappa_\tau|^{a'-a}
\le1\ \text{ for }a'\le a .
\end{equation}
We give $Y_\tau$ the Euclidean norm $\|\cdot\|=\|\cdot\|_\tau$ for which the normalised basis is
orthonormal, and $\operatorname{End}(Y_\tau)$ the operator norm. For $g\in G_\tau$ we have $g^{-1}=
\mathsf J_1g^*\mathsf J_1$, hence
\begin{equation}\label{pbe:eq-norm}
\|g^{-1}\|=\|g\|\ge1,\qquad|h(v,w)|=|\pbip{\mathsf J_1w}{v}|\le\|v\|\,\|w\|\qquad(g\in G_\tau,\ v,w\in
Y_\tau).
\end{equation}
The groups $U_\tau$, $H_\tau$, $R_\tau=U_\tau\rtimes H_\tau$ are the real points of $U$, $H$, $R_0$
at $\tau$. An element of $R_\tau$ maps $\tilde f_a$ to $\tilde f_a$ plus a combination of the
$\tilde f_{a'}$ with $a'<a$, and $\tilde e_a$ to $\tilde e_a$ plus a combination of the $\tilde
e_{a'}$ with $a'<a$ ($U$ is unitriangular on the flag and trivial on $A^\perp/A$, and $H$ fixes the
$e_a$ and $z$).

\begin{definition}[Length and cut]\label{pbe:length}
For $y\in G_\tau$ the \emph{length} is $\pbel_\tau(y)=\inf_{s\in R_\tau}\bigl(\|sy\|+1\bigr)\ge2$,
and the \emph{scaled length} is $\pbels_\tau(y)=\max\bigl(1,\pbel_\tau(y)/b^{6n}\bigr)$. The
\emph{cutoff} is
\begin{equation}\label{pbe:eq-chi}
\chi_\tau(y)=\max\Bigl(0,\ \min\Bigl(1,\ 2-\frac{2\,\pbel_\tau(y)}{b^{6n+2}}\Bigr)\Bigr),
\end{equation}
so that $\chi_\tau(y)=1$ if $\pbel_\tau(y)\le\frac12b^{6n+2}$ and $\chi_\tau(y)=0$ if $\pbel_\tau(y)
\ge b^{6n+2}$. For $y=(y_\tau)\in G_\infty$ put $\pbeL(y)=\prod_\tau\pbels_\tau(y_\tau)\ge1$.
\end{definition}

\begin{lemma}\label{pbe:len}
Let $y,g\in G_\tau$ and $s\in R_\tau$.
\begin{enumerate}
\item[(a)] $\pbel_\tau(sy)=\pbel_\tau(y)$, and $\pbel_\tau(yg)\le\|g\|\,\pbel_\tau(y)$. Hence
$\pbel_\tau$, $\pbels_\tau$ and $\chi_\tau$ are continuous and left invariant under $R_\tau$.
\item[(b)] If $\chi_\tau(y)<1$ then $\pbels_\tau(y)\ge b$.
\item[(c)] Let $w_a=y^{-1}\tilde f_a$ $(0\le a\le\varrho)$ and let $o_0,\dots,o_\varrho$ be the
Gram--Schmidt orthogonalisation of $w_0,\dots,w_\varrho$, so that $w_a=o_a+\sum_{a'<a}t_{a'a}o_{a'}$.
The vectors $o_a$ do not change if $y$ is replaced by $sy$, and $\|o_a\|<\pbel_\tau(y)$.
\end{enumerate}
\end{lemma}
\begin{proof}
(a) The first claim is clear. For the second, $\|s yg\|+1\le\|sy\|\,\|g\|+1\le\|g\|(\|sy\|+1)$ by
\eqref{pbe:eq-norm}. Applied to $g$ and $g^{-1}$ this gives $\|g\|^{-1}\pbel_\tau(y)\le\pbel_\tau(yg)
\le\|g\|\pbel_\tau(y)$, hence the continuity. (b) $\chi_\tau(y)<1$ means $\pbel_\tau(y)>\frac12b^{6n+
2}$, so $\pbels_\tau(y)>b^2/2\ge b$. (c) $(sy)^{-1}\tilde f_a=y^{-1}s^{-1}\tilde f_a$ is $w_a$ plus
a combination of the $w_{a'}$, $a'<a$; so the flags of the two families agree and the $o_a$ are the
same. Hence $\|o_a\|\le\|(sy)^{-1}\tilde f_a\|\le\|(sy)^{-1}\|=\|sy\|$ for every $s$, and the
infimum over $s$ is $\pbel_\tau(y)-1$.
\end{proof}

\subsubsection*{Measures}
We fix on $G_\tau$ the Haar measure $dg$ for which the formula \textup{(DS1)(iii)} below holds, on
the compact group $H_\tau$ the Haar measure of mass $1$, and on $U_\tau$ the Haar measure $du$ of
convention (Cv3) of \S\ref{pb:ss-setting} (Lebesgue measure in the free entries of the rows $\tilde e_a$,
as in \S\ref{prp:ss-haar}). The
group $R_\tau$ is unimodular with Haar measure $ds=du\,dh$ ($s=uh$), because conjugation by the
compact group $H_\tau$ preserves $du$. With the quotient measure on $U_\tau\backslash G_\tau$,
\begin{equation}\label{pbe:eq-unfoldR}
\int_{U_\tau\backslash G_\tau}\int_{R_\tau}F(sy)\,ds\,dy=\int_{G_\tau}F(g)\,dg\qquad(F\ge0\text{
measurable}),
\end{equation}
since $\int_{U_\tau\backslash G_\tau}\int_{U_\tau}F(uhy)\,du\,dy=\int_{G_\tau}F(hg)\,dg=\int_{G_\tau}F$
for each $h\in H_\tau$.

\subsection{The discrete series coefficient: what is used}\label{pbe:ss-ds}

Let $\pi_\tau$ be the discrete series representation of $G_\tau$ with Harish-Chandra parameter
$\varsigma_\tau$ (\S\ref{ss:au-spaces}) \emph{in the following chamber}: in a basis of $Y_\tau$ in
which the form is diagonal with $k+\varrho$ signs $+$ and $\varrho+1$ signs $-$, the $k-1$ entries
$\frac k2-1,\dots,1-\frac k2$ and the entries $-\kappa(\tau)-j\Delta$ with $j$ odd are attached to
coordinates of sign $+$, and the entries $-\kappa(\tau)-j\Delta$ with $j$ even to coordinates of sign
$-$. (The infinitesimal character $\varsigma_\tau$ alone does not determine the representation; this
is the choice made in the definition of $\auAz$ in \S\ref{ss:au-spaces}.) Entries attached to coordinates of opposite signs differ by at least
$\Delta=2r^8\ge500r^3+r^7$. Let $K_\tau=G_\tau\cap\auU(r)$ and let $v_0$ be a unit vector of highest
weight in the lowest $K_\tau$-type of $\pi_\tau$. Put
\[
c_\tau(g)=\pbip{v_0}{\pi_\tau(g)v_0},\qquad\pbdeg_\tau=\prod_{i<j}\frac{\varsigma_{\tau,i}-
\varsigma_{\tau,j}}{j-i},\qquad e_\tau=\pbdeg_\tau\cdot\aubar{c_\tau},\qquad\pbeq=9r^2 .
\]
Let $\pbW_\tau$ be the representation of $H_\tau=\auU(k)$ of \S\ref{pb:ss-layer2} (irreducible, of
dimension $d_\tau$; it is trivial if $\tau\notin\auPa$), with its action $\tau_\tau$, and let
$\psi_\tau$ be the component of $\psi_\infty$. Then
\[
\pi_R(uh)=\psi_\tau(u)\,\tau_\tau(h)\qquad(u\in U_\tau,\ h\in H_\tau)
\]
is a unitary representation of $R_\tau$ on $\pbW_\tau$, because $\psi_\tau$ is invariant under
conjugation by $H_\tau$. For $x\in G_\tau$ let $\pbeB_x=\{s\in R_\tau:\|x^{-1}sx-1\|\le1\}$, a
compact neighbourhood of $1$ in $R_\tau$ (the map $s\mapsto x^{-1}sx$ is a homeomorphism of $R_\tau$
onto a closed subgroup of $G_\tau$); so $0<\pbvol(\pbeB_x)<\infty$.

\begin{theorem}[Properties of the discrete series coefficient]\label{pbe:ds}
For every $\tau\in\auPl$:
\begin{enumerate}
\item[(DS1)] \textup{(i)} $c_\tau$ is continuous, $c_\tau(g^{-1})=\aubar{c_\tau(g)}$ and $|c_\tau|
\le1$. \textup{(ii)} If $c_\tau$ is integrable, then $\int_{G_\tau}e_\tau(a)\,c_\tau(ga)\,da=c_\tau(g)$
for all $g\in G_\tau$, for the Haar measure $da$ for which $\int_{G_\tau}|c_\tau|^2=\pbdeg_\tau^{-1}$.
\textup{(iii)} This Haar measure is the one that corresponds to the measure of mass $1$ on the
compact form $\auU(r)$.
\item[(DS2)] $|c_\tau(g)|\le\bigl(2/\|g\|\bigr)^{\,j+\pbeq}$ for all $g\in G_\tau$ and all integers
$0\le j\le r^7$.
\item[(DS3)] $\pbdeg_\tau\cdot\int_{G_\tau}\bigl(32/\|g\|\bigr)^{\pbeq}\,dg\le b^{99r^2}$.
\item[(DS4)] There are a point $x_\tau\in G_\tau$ with $\pbel_\tau(x_\tau)\le b^{5n}$ and a unit
vector $v_\tau\in\pbW_\tau$ such that
\[
\Bigl|\int_{R_\tau}c_\tau(x_\tau^{-1}s\,x_\tau)\,\pbip{v_\tau}{\pi_R(s)^{-1}v_\tau}\,ds\Bigr|\ \ge\
b^{-60(r+k)^2}\,\pbvol(\pbeB_{x_\tau}).
\]
\end{enumerate}
\end{theorem}
\smallskip\noindent\textup{(DS1)}--\textup{(DS3)} are proved in \S\ref{pds:sec}: \textup{(DS1)(i),(ii)} by Schur orthogonality, \textup{(DS1)(iii)} from Corollary~\ref{pbr:degree} (Propositions~\ref{pds:haar} and~\ref{pwf:haar}), \textup{(DS2)} from the integral formula \eqref{pds:eq-FJ} for $c_\tau$ (Theorems~\ref{pds:decay} and~\ref{pco:FJ}), \textup{(DS3)} from \textup{(DS1)(ii)} and the Casimir values (Theorem~\ref{pds:tail}). \textup{(DS4)} is proved in \S\ref{pbp:sec} (Theorem~\ref{pbp:main}).\par\smallskip

\begin{lemma}[Covering]\label{pbe:cover}
For all $x,y\in G_\tau$, $\displaystyle\int_{R_\tau}\Bigl(\frac2{\|x^{-1}sy\|}\Bigr)^{\pbeq}ds\le16^{
\pbeq}\,\pbvol(\pbeB_x)$.
\end{lemma}
\begin{proof}
Let $m(s)=x^{-1}sy$ and, for $T\ge1$, $S_T=\{s\in R_\tau:\|m(s)\|\le T\}$. Choose a maximal family
$s_1,\dots,s_\nu$ in $S_T$ with $\|x^{-1}s_is_j^{-1}x-1\|>\frac12$ for $i\ne j$. For $i\ne j$,
$m(s_i)-m(s_j)=(x^{-1}s_is_j^{-1}x-1)\,m(s_j)$, so by \eqref{pbe:eq-norm} $\|m(s_i)-m(s_j)\|\ge
\frac12\|m(s_j)^{-1}\|^{-1}\ge(2T)^{-1}$. The $m(s_i)$ lie in the ball of radius $T$ of the space of
$r\times r$ matrices, of real dimension $2r^2$; comparing volumes of balls, $\nu\le(1+4T^2)^{2r^2}
\le(5T^2)^{2r^2}$. By maximality $S_T\subseteq\bigcup_i\pbeB_x s_i$, and the Haar measure is right
invariant, so $\pbvol(S_T)\le5^{2r^2}T^{4r^2}\pbvol(\pbeB_x)$. As $\|m(s)\|\ge1$,
\[
\int_{R_\tau}\Bigl(\frac2{\|m(s)\|}\Bigr)^{\pbeq}ds=2^{\pbeq}\int_1^\infty\pbeq\,T^{-\pbeq-1}\pbvol
(S_T)\,dT\le2^{\pbeq}5^{2r^2}\frac{\pbeq}{\pbeq-4r^2}\pbvol(\pbeB_x)\le16^{\pbeq}\pbvol(\pbeB_x).
\qedhere
\]
\end{proof}

\subsection{The Bessel function at a real place}\label{pbe:ss-local}

In this subsection $\tau$ is fixed and dropped from the notation where possible; $x=x_\tau$ and
$v=v_\tau$ are as in \textup{(DS4)}.

\begin{proposition}[The local Bessel function]\label{pbe:local}
Assume Theorem~\ref{pbe:ds}. For $y\in G_\tau$ the integral
\[
W^\circ(y)=\int_{R_\tau}c_\tau(x^{-1}sy)\,\pi_R(s)^{-1}v\,ds\ \in\pbW_\tau
\]
converges absolutely, and:
\begin{enumerate}
\item[(a)] $W^\circ$ is continuous, $\|W^\circ(y)\|\le16^{\pbeq}\pbvol(\pbeB_x)$, and $W^\circ(sy)=
\pi_R(s)W^\circ(y)$ for $s\in R_\tau$.
\item[(b)] $c_\tau$ is integrable on $G_\tau$, $\int_{G_\tau}e_\tau(a)W^\circ(ya)\,da=W^\circ(y)$ for
all $y$, and $\int_{U_\tau\backslash G_\tau}\|W^\circ\|\le\int_{G_\tau}|c_\tau|$.
\item[(c)] \emph{(The norm is a period.)} $\displaystyle\pbdeg_\tau\int_{U_\tau\backslash G_\tau}\|W^
\circ(y)\|^2dy=\pbip v{W^\circ(x)}=\int_{R_\tau}c_\tau(x^{-1}sx)\pbip v{\pi_R(s)^{-1}v}\,ds$. In
particular this number is real and, by \textup{(DS4)}, at least $b^{-60(r+k)^2}\pbvol(\pbeB_x)>0$.
\end{enumerate}
Let $c_0>0$ be such that $W_\tau:=c_0W^\circ$ has $\int_{U_\tau\backslash G_\tau}\|W_\tau\|^2=d_\tau$,
and $W_\tau^\flat:=\chi_\tau W_\tau$ with $\chi_\tau$ as in \eqref{pbe:eq-chi}. Put
$\mathsf A_\tau(y)=c_0\int_{R_\tau}(2/\|x^{-1}sy\|)^{\pbeq}ds$ and $\mathsf a_\tau=c_0\,16^{\pbeq}\pbvol(
\pbeB_x)$. Then:
\begin{enumerate}
\item[(d)] $\mathsf A_\tau$ is continuous and left invariant under $R_\tau$, $\mathsf A_\tau\le\mathsf
a_\tau$, and
\[
\pbels_\tau(y)^{3r^3f}\,\|W_\tau(y)\|\le\mathsf A_\tau(y)\quad(y\in G_\tau),\qquad\mathsf a_\tau\int_{
U_\tau\backslash G_\tau}\mathsf A_\tau\le b^{400r^2}\,d_\tau .
\]
\item[(e)] $\|W^\flat_\tau\|\le\|W_\tau\|$, $W^\flat_\tau(y)=0$ unless $\pbel_\tau(y)\le b^{6n+2}$, and
$\displaystyle\int_{U_\tau\backslash G_\tau}\|W_\tau^\flat-W_\tau\|^2\le b^{-2r^3f}\,d_\tau$.
\end{enumerate}
\end{proposition}
\begin{proof}
(a) By \textup{(DS2)} with $j=0$ and Lemma~\ref{pbe:cover}, $\int_{R_\tau}|c_\tau(x^{-1}sy)|ds\le16^{
\pbeq}\pbvol(\pbeB_x)$; as $\pi_R$ is unitary and $\|v\|=1$ this gives the convergence and the bound.
For $y'$ in a compact neighbourhood of $y$, $\|x^{-1}sy\|\le\|x^{-1}sy'\|\,\|y'^{-1}y\|$, so
$|c_\tau(x^{-1}sy')|\le C(2/\|x^{-1}sy\|)^{\pbeq}$ with a constant $C$, an integrable function of $s$;
the continuity follows by dominated convergence. The substitution $s\mapsto ss'^{-1}$ gives
$W^\circ(s'y)=\int c_\tau(x^{-1}sy)\pi_R(ss'^{-1})^{-1}v\,ds=\pi_R(s')W^\circ(y)$.

(b) By \textup{(DS2)} with $j=0$ and \textup{(DS3)}, $\int_{G_\tau}|c_\tau|<\infty$. By (a) the
double integral of $|e_\tau(a)|\,|c_\tau(x^{-1}sya)|$ over $R_\tau\times G_\tau$ is at most
$\|e_\tau\|_1\,16^{\pbeq}\pbvol(\pbeB_x)$, so Fubini applies, and by \textup{(DS1)(ii)}
\[
\int_{G_\tau}e_\tau(a)W^\circ(ya)\,da=\int_{R_\tau}\Bigl(\int_{G_\tau}e_\tau(a)c_\tau(x^{-1}sya)\,da
\Bigr)\pi_R(s)^{-1}v\,ds=W^\circ(y).
\]
By \eqref{pbe:eq-unfoldR}, $\int_{U_\tau\backslash G_\tau}\|W^\circ\|\le\int_{U_\tau\backslash G_\tau}
\int_{R_\tau}|c_\tau(x^{-1}sy)|\,ds\,dy=\int_{G_\tau}|c_\tau(x^{-1}g)|\,dg$.

(c) Insert the definition of $W^\circ$ in the first variable of $\pbip{W^\circ(y)}{W^\circ(y)}$ and
use (a):
\[
\pbip{W^\circ(y)}{W^\circ(y)}=\int_{R_\tau}\aubar{c_\tau(x^{-1}sy)}\,\pbip{\pi_R(s)^{-1}v}{W^\circ(y)}
\,ds=\int_{R_\tau}\aubar{c_\tau(x^{-1}sy)}\,\pbip v{W^\circ(sy)}\,ds .
\]
Let $\Phi(g)=e_\tau(x^{-1}g)\pbip v{W^\circ(g)}$; it is integrable on $G_\tau$, since $e_\tau$ is and
$W^\circ$ is bounded. By the identity above and \eqref{pbe:eq-unfoldR} (applied to the positive and
negative parts of the real and imaginary parts of $\Phi$),
\[
\pbdeg_\tau\int_{U_\tau\backslash G_\tau}\|W^\circ\|^2=\int_{U_\tau\backslash G_\tau}\int_{R_\tau}\Phi(
sy)\,ds\,dy=\int_{G_\tau}\Phi(g)\,dg=\int_{G_\tau}e_\tau(a)\pbip v{W^\circ(xa)}\,da=\pbip v{W^\circ
(x)},
\]
by (b). The last expression of (c) is the definition of $W^\circ(x)$.

(d) The invariance of $\mathsf A_\tau$ under $y\mapsto s'y$ is the right invariance of $ds$; the
continuity is proved as in (a); $\mathsf A_\tau\le\mathsf a_\tau$ is Lemma~\ref{pbe:cover}. Let $s\in
R_\tau$ and $g=x^{-1}sy$. By Lemma~\ref{pbe:len}(a), $\pbel_\tau(y)=\pbel_\tau(sy)=\pbel_\tau(xg)\le
\|g\|\pbel_\tau(x)\le\|g\|b^{5n}$, so $\pbel_\tau(y)/b^{6n}\le\|g\|b^{-n}\le\|g\|/2$ and
$\pbels_\tau(y)\le\max(1,\|g\|/2)$. If $\pbels_\tau(y)=1$,
\textup{(DS2)} with $j=0$ gives $\pbels_\tau(y)^{3r^3f}|c_\tau(g)|\le(2/\|g\|)^{\pbeq}$. Otherwise
$\pbels_\tau(y)\le\|g\|/2$, and \textup{(DS2)} with $j=3r^3f\le r^7$ gives $\pbels_\tau(y)^{j}|c_\tau(
g)|\le(\|g\|/2)^j(2/\|g\|)^{j+\pbeq}=(2/\|g\|)^{\pbeq}$. Integrating over $s$ and multiplying by
$c_0$ we get the first inequality of (d). For the second, by \eqref{pbe:eq-unfoldR}, (c), the
definition of $c_0$ and \textup{(DS3)},
\begin{align*}
\mathsf a_\tau\int_{U_\tau\backslash G_\tau}\mathsf A_\tau&=c_0^2\,16^{\pbeq}\pbvol(\pbeB_x)\int_{G_
\tau}\Bigl(\frac2{\|g\|}\Bigr)^{\pbeq}dg\le c_0^2\,b^{60(r+k)^2}\pbdeg_\tau\!\int_{U_\tau\backslash
G_\tau}\!\!\|W^\circ\|^2\cdot\int_{G_\tau}\Bigl(\frac{32}{\|g\|}\Bigr)^{\pbeq}dg\\
&=d_\tau\,b^{60(r+k)^2}\,\pbdeg_\tau\int_{G_\tau}\Bigl(\frac{32}{\|g\|}\Bigr)^{\pbeq}dg\le d_\tau\,
b^{240r^2+99r^2}\le d_\tau\,b^{400r^2}.
\end{align*}

(e) The first two claims follow from $0\le\chi_\tau\le1$ and \eqref{pbe:eq-chi}. Where $W_\tau^\flat
\ne W_\tau$ we have $\pbels_\tau\ge b$ by Lemma~\ref{pbe:len}(b), so by (d)
$b^{3r^3f}\|W_\tau^\flat-W_\tau\|\le b^{3r^3f}\|W_\tau\|\le\mathsf A_\tau\le\mathsf a_\tau$ everywhere,
hence $b^{6r^3f}\|W^\flat_\tau-W_\tau\|^2\le\mathsf a_\tau\mathsf A_\tau$. Integrating,
$\int\|W_\tau^\flat-W_\tau\|^2\le b^{-6r^3f}b^{400r^2}d_\tau\le b^{-2r^3f}d_\tau$, because $400r^2\le4r^3
\le4r^3f$.
\end{proof}

\begin{remark}\label{pbe:rem-local}
The whole local theory is thus four lines: the average of the coefficient over $R_\tau$ is
reproduced by the kernel; therefore its $L^2$-norm on $U_\tau\backslash G_\tau$ is its own value at
$x$, a period; the period is not small by \textup{(DS4)}; and the decay of the coefficient passes
to the average through the covering Lemma~\ref{pbe:cover}.
\end{remark}

\subsection{The Poincar\'e series}\label{pbe:ss-global}

Put $W=\bigotimes_\tau W_\tau$, $W^\flat=\bigotimes_\tau W_\tau^\flat$ on $G_\infty$ with values in
$\pbW=\bigotimes_\tau\pbW_\tau$, $g_0=(x_\tau)_\tau$, $e_\infty=\bigotimes_\tau e_\tau$, $\mathsf
A(y)=\prod_\tau\mathsf A_\tau(y_\tau)$ and $\mathsf a_\infty=\prod_\tau\mathsf a_\tau$. By
Proposition~\ref{pbe:local}(d), for $y\in G_\infty$,
\begin{equation}\label{pbe:eq-decay}
\|W(y)\|\le\mathsf A(y)\,\pbeL(y)^{-3r^3f}\le\mathsf a_\infty\,\pbeL(y)^{-3r^3f},\qquad\mathsf a_\infty
\int_{U_\infty\backslash G_\infty}\mathsf A\le b^{400r^2f}\,d_\infty .
\end{equation}

\subsubsection*{The space $\auAz$}
By the definition of \S\ref{ss:au-spaces}, $\auAz$ is the space of the bounded continuous functions $y$ on $G(\auA)$, left invariant under
$\Gamma$ and right invariant under $J$, with
\begin{equation}\label{pbe:eq-repro}
\int_{G_\infty}e_\infty(a)\,y(ga)\,da=y(g)\qquad(g\in G(\auA)).
\end{equation}
This is the meaning of ``the archimedean component is the vector singled out by the reproducing
kernel''. The quotient $\Gamma\backslash G(\auA)$ has finite volume~\cite[Theorem~5.6(i)]{BorelFin} (the group has no rational characters, its centre being anisotropic; see also \cite[Theorem~9.4]{BorelHC1962} for arithmetic subgroups of the real group), so $\auAz$
carries the inner product of \S\ref{pb:ss-setting}; as in \S\ref{ss:au-spaces} we use that $\auAz$
is finite-dimensional. The linear form $\varphi\mapsto\int_{[R_0]}\mathcal W_\varphi(r)\aubar{y(rg_0)}
\,dr$ is defined for $y\in\auAz$ because $[R_0]$ is compact ($U$ is unipotent and $H$ is
anisotropic) and $y$ is continuous.

\subsubsection*{Counting the terms of a series}

\begin{lemma}\label{pbe:eichler}
Let $c\in\mathcal C$. For every strictly upper triangular matrix $(C_{a'a})_{0\le a',a\le\varrho}$
with entries in $M\auO_E$ there is $u\in U_c$ with $uf_a=f_a+\sum_{a'<a}C_{a'a}f_{a'}$ for $0\le a\le
\varrho$.
\end{lemma}
\begin{proof}
Let $L'=h_cL$, where $L$ is the lattice of $Y_0$; as $h_c\in H(\auA_f)$, $L'=L_{\mathrm{hyp}}\oplus
L'_{\mathrm{mid}}$ with $L_{\mathrm{hyp}}=\bigoplus_a(\widehat{\auO}e_a\oplus\widehat{\auO}e_a^*)$ and
$L'_{\mathrm{mid}}$ a unimodular lattice of $X_0\perp Ez$ containing $z$, and $U_c$ consists of the
$u\in U(E^+)$ with $(u-1)L'\subseteq ML'$ (the condition at the conjugates of the chosen places is
the same condition, read on the adjoint).
\emph{First}, let $a_1$ be the unitriangular endomorphism of $A$ with $a_1e_{a}=e_a+\sum_{a'<a}C_{a'-1,
a-1}e_{a'}$ and $\mathsf m(a_1)\in G(E^+)$ the element that acts by $a_1$ on $A$, by the inverse
adjoint of $a_1$ on $\langle e_a^*\rangle$ and trivially on $X_0\perp Ez$. It lies in $U(E^+)$,
commutes with $h_c$, and $\mathsf m(a_1)-1$ maps $L_{\mathrm{hyp}}$ into $ML_{\mathrm{hyp}}$, since
$a_1^{-1}-1$ is a polynomial without constant term in $a_1-1$. So $\mathsf m(a_1)\in U_c$.
\emph{Second}, let $v\in M\bigoplus_a\auO_Ee_a$. The map $q\mapsto h(q,z)$ from $L'_{\mathrm{mid}}\cap
(X_0\perp Ez)$ to $\auO_E$ is onto: its image is an ideal whose completion is $h(L'_{\mathrm{mid}},z)
=\widehat{\auO}$, because $L'_{\mathrm{mid}}$ is unimodular and $z$ is primitive in it ($h(x_k/\aubar{
\pi_N}+z/N,z)=-1$, and $h_c$ fixes $z$). Choose $q$ with $h(z,q)=1$, and $\omega\in\auO_E$ with
$\omega+\aubar\omega=1$ ($E/E^+$ is unramified at all finite places, so the trace is onto). Then
\[
\mathsf E(t)=t+h(t,q)\,v-h(t,v)\,q-\omega\,h(q,q)\,h(t,v)\,v
\]
is an isometry of $Y_0$ (a direct computation, using $h(v,v)=h(v,q)=0$); it fixes $A$ pointwise and
maps $t\in A^\perp$ to $t+h(t,q)v\in t+A$, so $\mathsf E\in U(E^+)$; $\mathsf E(z)=z+v$; and $(\mathsf
E-1)L'\subseteq ML'$ because $q\in L'$, $v\in ML'$. So $\mathsf E\in U_c$. The element $u=\mathsf
m(a_1)\mathsf E$ with $v=a_1^{-1}\sum_{a'<\varrho}C_{a'\varrho}e_{a'+1}$ has the required property.
\end{proof}

\begin{proposition}[Count of terms]\label{pbe:count}
Let $c\in\mathcal C$, $s\in\auG$ and $y\in G_\infty$. Then
\[
\sum_{c'\in\mathcal C}\ \sum_{\gamma\in\Gamma_{c,s,c'}/U_{c'}}\pbeL(\gamma^{-1}y)^{-r^3f}\ \le\ 2\,b^{
12n^3rf}.
\]
\end{proposition}
The summand depends only on the coset $\gamma U_{c'}$, by Lemma~\ref{pbe:len}(a).
\begin{proof}
For $R\ge1$ let $\pbeN_R$ be the set of the pairs $(c',\gamma U_{c'})$ with $\gamma\in\Gamma_{c,s,c'}$
and $\pbeL(\gamma^{-1}y)\le R$. We show
\begin{equation}\label{pbe:eq-pack}
\#\pbeN_R\le\bigl(b^{6n^2}R^2\bigr)^{D},\qquad D=2nrf .
\end{equation}
Then, with $i=r^3f\ge4D$ (as $r^2\ge8n$), the sum of the proposition is
$\int_1^\infty iT^{-i-1}\#\pbeN_T\,dT\le b^{6n^2D}\,i/(i-2D)\le2b^{12n^3rf}$.

Put $\rho=b^{37n}R$. For $(c',\gamma U_{c'})\in\pbeN_R$ and every $\tau$ we have $\pbels_\tau\le\pbeL$,
hence
\begin{equation}\label{pbe:eq-lenR}
\pbel_\tau\bigl((\gamma^{-1}y)_\tau\bigr)\le b^{6n}R .
\end{equation}
\emph{Step 1: reduced representatives.} Each element of $\pbeN_R$ has a representative $\gamma$ with
$\|y_\tau^{-1}\gamma\tilde f_a\|\le\rho$ for all $\tau$ and $a$. Indeed let $z=\gamma^{-1}y$, and at
each $\tau$ let $w_a=z_\tau^{-1}\tilde f_a$, $o_a$ and $t_{a'a}$ be as in Lemma~\ref{pbe:len}(c). For
$u\in U_{c'}$ with $uf_a=f_a+\sum_{a'<a}C_{a'a}f_{a'}$, replacing $\gamma$ by $\gamma u$ replaces
$w_a$ by $z_\tau^{-1}u\tilde f_a=w_a+\sum_{a'<a}\tau(C_{a'a})(s_a/s_{a'})w_{a'}$, does not change the
$o_a$, and replaces $t_{a'a}(s_{a'}/s_a)$ by
\[
t_{a'a}\frac{s_{a'}}{s_a}+\tau(C_{a'a})+\sum_{a'<a''<a}t_{a'a''}\frac{s_{a'}}{s_{a''}}\,\tau(C_{a''a}).
\]
The lattice $M\auO_E$ of $\prod_\tau\C$ contains $M\auO_K[\zeta_{q_1}]$, which has the $\Z$-basis
$M\zeta_{q_1}^i$, $M\omega_K\zeta_{q_1}^i$ ($0\le i<f$, $\omega_K=\frac12(1+\sqrt{-d})$ or $\sqrt{-d}$)
of vectors with all coordinates of absolute value at most $M\sqrt d$; so every point of
$\prod_\tau\C$ is at distance at most $fM\sqrt d\le b^{26}$, in every coordinate, from a point of $M
\auO_E$. By Lemma~\ref{pbe:eichler} we can therefore choose the $C_{a'a}\in M\auO_E$, by induction
on $a-a'$, such that after the replacement $|t_{a'a}(s_{a'}/s_a)|\le b^{26}$ at every $\tau$; then
$|t_{a'a}|\le b^{26}$ by \eqref{pbe:eq-scale}, and by Lemma~\ref{pbe:len}(c) and \eqref{pbe:eq-lenR}
\[
\|y_\tau^{-1}\gamma\tilde f_a\|=\|w_a\|\le\sum_{a'\le a}|t_{a'a}|\,\|o_{a'}\|\le n\,b^{26}\,b^{6n}R\le\rho .
\]
\emph{Step 2: distinct elements have distinct frames.} Let $\gamma\in\Gamma_{c,s,c'}$, $\gamma'\in
\Gamma_{c,s,c''}$ with $\gamma f_a=\gamma'f_a$ for all $a$. Then $g=\gamma^{-1}\gamma'$ fixes the
frame, so $g=u\eta$ with $\eta\in H(E^+)$ and $u\in U(E^+)$ fixing the frame, and
$h_{c'}^{-1}gh_{c''}=(h_{c'}^{-1}\gamma^{-1}h_c)(h_c^{-1}\gamma'h_{c''})$ lies in $K$ and has reduction
$s^{-1}s=1$; so it lies in $J\cap R_0(\auA_f)$. By Lemma~\ref{pb:KR} and $\pbUb\cap\pbHb=1$,
$J\cap R_0(\auA_f)=(J\cap U(\auA_f))\rtimes J_X$. Hence $h_{c'}^{-1}\eta h_{c''}\in J_X$, so $c''=c'$
and $\eta\in H(E^+)\cap h_{c'}J_Xh_{c'}^{-1}=1$, the level being free; and then $u\in U_{c'}$. So
$\gamma'U_{c'}=\gamma U_{c'}$.

\emph{Step 3: rigidity.} Let $\gamma,\gamma'$ be reduced representatives, as in Step~1, of two
elements of $\pbeN_R$, and suppose that
\begin{equation}\label{pbe:eq-close}
\|y_\tau^{-1}\gamma'\tilde f_a-y_\tau^{-1}\gamma\tilde f_a\|<(N\rho)^{-1}\qquad\text{for all }\tau,a .
\end{equation}
We claim $\gamma'f_a=\gamma f_a$ for all $a$. Let $g=\gamma^{-1}\gamma'$; then $h_{c'}^{-1}gh_{c''}\in
K$, and the $h_c$ fix the vectors $e_a$, $e_a^*$, $z$ of the lattice, so the numbers $h(gf_a,f_b)$
and $h(gf_a,e_x^*)$ are algebraic integers of $E$. First, by \eqref{pbe:eq-norm}, \eqref{pbe:eq-close}
and $|s_a|\ge N^{-1/2}$,
\[
\bigl|\tau\bigl(h(gf_a,f_b)-h(f_a,f_b)\bigr)\bigr|=\frac{|h(\gamma'\tilde f_a-\gamma\tilde f_a,\gamma
\tilde f_b)|}{|s_a||s_b|}\le N\,\|y_\tau^{-1}(\gamma'\tilde f_a-\gamma\tilde f_a)\|\,\|y_\tau^{-1}
\gamma\tilde f_b\|<1
\]
for every $\tau$; an algebraic integer all of whose conjugates have absolute value less than $1$ is
$0$, so $h(gf_a-f_a,f_b)=0$ for all $a,b$. Thus $gf_a-f_a$ is orthogonal to $\pbeAp=A+Ez$, that is,
$gf_a=f_a+\alpha_a+\xi_a$ with $\alpha_a\in A$, $\xi_a\in X_0$. As $g$ is an isometry and $A\oplus
X_0$ is orthogonal to the frame, $h(f_a,f_b)=h(gf_a,gf_b)=h(f_a,f_b)+h(\xi_a,\xi_b)$; $X_0$ is
definite, so $\xi_a=0$. Second, write $\alpha_a=\sum_xc_{xa}e_x$ with $c_{xa}=h(gf_a-f_a,e_x^*)\in\auO_E$
and suppose $\alpha_a\ne0$; let $x$ be the largest index with $c_{xa}\ne0$. At a place $\tau$,
$g\tilde f_a-\tilde f_a=\sum_{x'}\tau(c_{x'a})(s_a/s'_{x'})\tilde e_{x'}=:\tilde\alpha$, where
$s'_{x'}$ is the scale of $e_{x'}$ and $|s_a/s'_{x'}|\ge|\kappa_\tau|^{-\varrho}\ge b^{-23\varrho}$
by \eqref{pbe:eq-kappa}. Let $z=\gamma^{-1}y$. For $s\in R_\tau$ the vector $s\tilde\alpha$ has
the same coefficient at $\tilde e_x$ as $\tilde\alpha$, and $s\tilde\alpha=(sz_\tau)\,z_\tau^{-1}\tilde
\alpha$, so
\[
b^{-23\varrho}|\tau(c_{xa})|\le\|s\tilde\alpha\|\le\|sz_\tau\|\,\|z_\tau^{-1}(g\tilde f_a-\tilde f_a)\|
<\|sz_\tau\|\,(N\rho)^{-1}
\]
by \eqref{pbe:eq-close}. Taking the infimum over $s$ and using \eqref{pbe:eq-lenR}, $|\tau(c_{xa})|<
b^{23\varrho+6n}R\,(Nb^{37n}R)^{-1}<1$ for every $\tau$, so $c_{xa}=0$, a contradiction. This proves
the claim.

\emph{Step 4: packing.} To each element of $\pbeN_R$ attach the point $(y_\tau^{-1}\gamma\tilde f_a)_{
\tau,a}$ of $V=(\C^r)^{n\times f}$, for a reduced representative $\gamma$; $V$ has real dimension
$D$ and we give it the norm $\max_{\tau,a}\|\cdot\|$. By Step~1 the points lie in the ball of radius
$\rho$; by Steps~2 and~3 the points of two different elements are at distance at least $\delta=(N
\rho)^{-1}$. Comparing volumes, $\#\pbeN_R\le(1+2\rho/\delta)^D=(1+2N\rho^2)^D\le(3b^{134n}R^2)^D
\le(b^{6n^2}R^2)^D$, as $n\ge23$. This is \eqref{pbe:eq-pack}.
\end{proof}

\subsubsection*{Existence of the Bessel functions}

\begin{theorem}[Hypothesis~\ref{pb:bessel} holds]\label{pbe:bessel}
Assume Theorem~\ref{pbe:ds}. Let $g_0=(x_\tau)_\tau$, let $\mathsf t=\bigotimes_\tau v_\tau\in\pbW$,
regarded as a linear form on $W_{\lambda^X}$, and let $W_\tau$, $W_\tau^\flat$, $\chi_\tau$ be as in
Proposition~\ref{pbe:local}. Then \textup{(B1)--(B4)} of Definition~\ref{pb:bessel} hold, with
\[
c_W=\frac{\pbvol(R_0(\auA_f)\cap J)}{\pbvol(J)}\prod_\tau\frac{\pbdeg_\tau}{c_{0,\tau}}>0,
\]
and \textup{(B5)} holds in the following exact form:
\begin{enumerate}
\item[(B5)] $W_\tau^\flat=\chi_\tau W_\tau$, where $\chi_\tau(y)=\max\bigl(0,\min(1,2-2\pbel_\tau(y)/
b^{6n+2})\bigr)$ and $\pbel_\tau(y)=\inf_{s\in R_\tau}(\|sy\|_\tau+1)$ is the length of
Definition~\ref{pbe:length}.
\end{enumerate}
Moreover:
\begin{enumerate}
\item[(B6)] \emph{(Decay.)} \eqref{pbe:eq-decay} holds, with $\mathsf A$ continuous and left
invariant under $R_\infty$.
\item[(B7)] \emph{(Reproduction.)} $\int_{G_\tau}e_\tau(a)W_\tau(ya)\,da=W_\tau(y)$ for all $\tau$
and $y\in G_\tau$.
\end{enumerate}
\end{theorem}
\begin{proof}
(B1), (B2), (B5), (B6), (B7) are Proposition~\ref{pbe:local} and Lemma~\ref{pbe:len}(a) (for
$W_\tau^\flat$ in (B1)).

(B3) Fix $g=(g_\infty,g_f)$. If no $\gamma_1\in\Gamma$ has $\gamma_1g_f\in h_cJ$, the series is
empty. Otherwise the $\gamma$ that occur are those of $\Gamma_c\gamma_1$, $\Gamma_c=\Gamma\cap h_cJh_c
^{-1}=\Gamma_{c,1,c}$, and with $y=\gamma_{1,\infty}g_\infty$, \eqref{pbe:eq-decay} and
Proposition~\ref{pbe:count} (for $s=1$, keeping only $c'=c$, and $\gamma\mapsto\gamma^{-1}$),
\[
\sum_{\gamma\in U_c\backslash\Gamma}\|W(\gamma_\infty g_\infty)\|\,\pbone_{h_cJ}(\gamma g_f)=\sum_{
\gamma\in\Gamma_c/U_c}\|W(\gamma^{-1}y)\|\le\mathsf a_\infty\sum_{\gamma\in\Gamma_c/U_c}\pbeL(\gamma^{
-1}y)^{-r^3f}\le2\,\mathsf a_\infty\,b^{12n^3rf}.
\]
So the series is bounded, and $P_{c,i}$ is a bounded function, left invariant under $\Gamma$ and
right invariant under $J$. By (B7), Fubini and the bound just proved (with $\|e_\infty\|_1<\infty$),
$P_{c,i}$ satisfies \eqref{pbe:eq-repro}; and a bounded function that satisfies \eqref{pbe:eq-repro}
is continuous, because $g_\infty\mapsto e_\infty(g_\infty^{-1}\,\cdot\,)$ is continuous into
$L^1(G_\infty)$. So $P_{c,i}\in\auAz$.

(B4) Let $y\in\auAz$ and $\varphi\in\auS$. The function $F(r)=\mathcal W_\varphi(r)\aubar{y(rg_0)}$ on
$R_0(\auA)$ is left invariant under $R_0(E^+)$ and right invariant under $K_R'=R_0(\auA_f)\cap J$
(for $\mathcal W_\varphi$ this is the computation in the proof of Lemma~\ref{pb:equiv}, with
$\psi_f$ trivial on $U(\auA_f)\cap J$ and $\varphi$ right invariant under $J_X$). By strong
approximation for $U$ and the definition of the classes, $R_0(\auA)=\bigsqcup_cR_0(E^+)\,(R_\infty
\times h_cK_R')$, and the stabiliser $R_0(E^+)\cap h_cK_R'h_c^{-1}$ is $U_c$ (as in Step~2 of the
proof of Proposition~\ref{pbe:count}). Hence
\[
\pbip y{\auP\varphi}=\int_{[R_0]}F=\pbvol(K_R')\sum_c\int_{U_c\backslash R_\infty}\mathcal W_\varphi(s,
h_c)\,\aubar{y(sg_0,h_c)}\,ds .
\]
By \eqref{pbe:eq-repro} and the substitution $a=g_0^{-1}s^{-1}g$, $\aubar{y(sg_0,h_c)}=\int_{G_\infty}
\aubar{e_\infty(g_0^{-1}s^{-1}g)}\,\aubar{y(g,h_c)}\,dg$. For $s=uh$, $\mathcal W_\varphi(s,h_c)=
\psi_\infty(u)\,\mathsf t(h^{-1}\varphi(h_c))=\pbip{\pi_R(s)\mathsf t}{\varphi(h_c)}$ in the pairing of
$\pbW$ with $W_{\lambda^X}$ (here $\pi_R=\bigotimes_\tau\pi_{R,\tau}$ and $\tau_\infty$ is the
contragredient action). As $y$ is bounded, $e_\infty$ integrable and $y(ug,h_c)=y(g,h_c)$ for $u\in
U_c$, we may exchange the integrals and fold the integral over $G_\infty$:
\[
\pbip y{\auP\varphi}=\pbvol(K_R')\sum_c\int_{U_c\backslash G_\infty}\aubar{y(g,h_c)}\,\Bigl\langle\int_{
R_\infty}\aubar{e_\infty(g_0^{-1}s^{-1}g)}\,\pi_R(s)\mathsf t\,ds,\ \varphi(h_c)\Bigr\rangle\,dg .
\]
Since $\aubar{e_\tau}=\pbdeg_\tau c_\tau$, the substitution $s\mapsto s^{-1}$ shows that the inner
integral is $\bigotimes_\tau\pbdeg_\tau W_\tau^\circ(g_\tau)=\prod_\tau(\pbdeg_\tau/c_{0,\tau})\cdot
W(g)$. On the other hand, by (B3) the series $P_\varphi=\sum_{c,i}\varphi(h_c)_iP_{c,i}$ lies in
$\auAz$, and unfolding it as in the proof of Proposition~\ref{pb:unfold} gives $\pbip y{P_\varphi}=
\pbvol(J)\sum_c\int_{U_c\backslash G_\infty}\aubar{y(g,h_c)}\pbip{W(g)}{\varphi(h_c)}\,dg$. So
$\pbip y{\auP\varphi}=c_W\pbip y{P_\varphi}$ for all $y\in\auAz$, and $\auP\varphi=c_WP_\varphi$.
\end{proof}

\subsection{Vanishing by character, and the terms that preserve \texorpdfstring{$A$}{A}}
\label{pbe:ss-char}

For $\gamma\in\Gamma$ let $V_\gamma=U\cap\gamma U\gamma^{-1}$, a unipotent group over $E^+$, and
$\chi_\gamma(u)=\psi_\infty(u)\,\aubar{\psi_\infty(\gamma^{-1}u\gamma)}$, a character of $V_\gamma(
E^+\otimes\R)$. We say that $\gamma$ \emph{agrees} if $\chi_\gamma$ is trivial on $V_\gamma(E^+)$;
since $\psi_\infty\psi_f=1$ on $U(E^+)$, this means $\psi_f(u)=\psi_f(\gamma^{-1}u\gamma)$ for all
$u\in U(E^+)\cap\gamma U(E^+)\gamma^{-1}$. As $\psi_\infty$ is a character of $U_\infty$, $\chi_{u_1
\gamma u_2}(u)=\chi_\gamma(u_1^{-1}uu_1)$ for $u_1,u_2\in U(E^+)$, so the set of agreeing elements is
stable.

\begin{lemma}[Vanishing by character]\label{pbe:vanish}
Assume \textup{(B1)--(B4)}. Let $\flat\in\{\ ,\flat\}$ and let $\gamma\in\Gamma_{c,s,c'}$ not agree.
Then the double coset of $\gamma$ contributes $0$ to \eqref{pb:eq-amp}:
\[
\int_{(U_c\cap\gamma U_{c'}\gamma^{-1})\backslash G_\infty}W^{\flat}(y)\,W^{\flat}(\gamma^{-1}y)^*\,dy=0.
\]
Consequently $A_T^\flat(c,s,c')=0$ for every stable $T$ without agreeing elements.
\end{lemma}
\begin{proof}
Let $\Lambda=U_c\cap\gamma U_{c'}\gamma^{-1}$ and $V_\infty=V_\gamma(E^+\otimes\R)$. The integral
converges absolutely (Proposition~\ref{pb:unfold} and $\|W^\flat\|\le\|W\|$). The group $\Lambda$
is an arithmetic subgroup of the unipotent group $V_\gamma$, hence discrete and cocompact in
$V_\infty$, and $V_\infty$ is unimodular. By (B1) and Lemma~\ref{pbe:len}(a) the integrand $F(y)$
satisfies $F(uy)=\chi_\gamma(u)F(y)$ for $u\in V_\infty$, and it is invariant under $\Lambda$
(\S\ref{pb:ss-layer2}). If $F=0$ there is nothing to prove. Otherwise $\chi_\gamma$ is trivial on
$\Lambda$, and the integral equals $\int_{V_\infty\backslash G_\infty}F(y)\bigl(\int_{\Lambda\backslash
V_\infty}\chi_\gamma(u)\,du\bigr)dy$. The measure of $\Lambda\backslash V_\infty$ is invariant under
right translations, so $\int_{\Lambda\backslash V_\infty}\chi_\gamma(u)\,du=\chi_\gamma(u_0)\int_{
\Lambda\backslash V_\infty}\chi_\gamma(u)\,du$ for every $u_0\in V_\infty$; and $\chi_\gamma(u_0)\ne1$
for some $u_0\in V_\gamma(E^+)\subseteq V_\infty$, as $\gamma$ does not agree. So the inner integral
is $0$.
\end{proof}

\begin{lemma}[Agreeing terms that preserve $A$]\label{pbe:scalar}
Let $\gamma\in\Gamma$ agree and $\gamma A=A$. Then $\gamma=\zeta\gamma_0$ with $\gamma_0\in R_0(E^+)$
and $\zeta\in E$, $\zeta\aubar\zeta=1$. If moreover $\gamma\in\Gamma_{c,s,c'}$,
then $\zeta\in\auO_E$ and $s\in\zeta\pbR$.
\end{lemma}
\begin{proof}
Let $P_A$ be the stabiliser of $A$, with unipotent radical $N_A$ (the elements acting trivially on
$A$, $A^\perp/A$ and $Y_0/A^\perp$). Then $N_A\subseteq U$, $\gamma$ normalises $N_A$, and $\gamma$
induces $a\in\auGL(A)$ and an isometry $\mathsf m$ of $Q=A^\perp/A=X_0\perp Ez$. Write $A_j=\langle
e_1,\dots,e_j\rangle$ and $e_j^*$ also for the linear form $h(\cdot,e_j^*)$ on $A$. We use the
following remark. Let $\psi_{E,f}$ be the finite part of $\psi_E$, so that $\psi_f(u)=\psi_{E,f}(-
\sigma(u)/(M\varpi))$. If $\mathsf l_1,\mathsf l_2$ are $E$-linear forms on an $E$-vector space with
$\psi_{E,f}(\mathsf l_1(\phi)/(M\varpi))=\psi_{E,f}(\mathsf l_2(\phi)/(M\varpi))$ for all $\phi$, then
$\mathsf l_1=\mathsf l_2$; for otherwise $\mathsf l_1-\mathsf l_2$ would be onto $E$ and $\psi_{E,f}$
would be trivial on $E$, which is dense in $\auA_{E,f}$.

\emph{The radical.} An element $\nu\in N_A(E^+)$ induces $\phi_\nu\in\operatorname{Hom}_E(Q,A)$,
$\phi_\nu(t)=\nu t-t$, every $\phi$ occurs, $\sigma(\nu)=e_\varrho^*(\phi_\nu(z))$, and $\phi_{\gamma^{
-1}\nu\gamma}=a^{-1}\phi_\nu\mathsf m$. As $\gamma$ agrees, the remark gives $e_\varrho^*(\phi(z))=e_
\varrho^*(a^{-1}\phi(\mathsf mz))$ for all $\phi$. With $\phi=\lambda(\cdot)v$, $\lambda$ a linear
form on $Q$ and $v\in A$: $\lambda(z)e_\varrho^*(v)=\lambda(\mathsf mz)e_\varrho^*(a^{-1}v)$. If
$\mathsf mz\notin Ez$ we could take $\lambda(z)=0\ne\lambda(\mathsf mz)$ and get $e_\varrho^*\circ a^{
-1}=0$, which is absurd. So $\mathsf mz=\zeta z$, $\zeta\aubar\zeta=1$ because $h(z,z)\ne0$, and
$e_\varrho^*\circ a^{-1}=\zeta^{-1}e_\varrho^*$ on $A$: $a$ preserves $A_{\varrho-1}$ and acts by
$\zeta$ on $A/A_{\varrho-1}$.

\emph{Descent on the flag.} Suppose, for some $1\le j<\varrho$, that $a$ preserves $A_i$ and acts
by $\zeta$ on $A_{i+1}/A_i$ for $j\le i<\varrho$. For $v\in A_j$ let $u_v\in U(E^+)$ act on $A$ by
$1+\phi$, $\phi(e_{j+1})=v$, $\phi(e_i)=0$ for $i\ne j+1$, by the inverse adjoint on $\langle
e_a^*\rangle$ and trivially on $X_0\perp Ez$. Then $\gamma^{-1}u_v\gamma$ is conjugate under $N_A
\subseteq U$ to the element of the same kind that acts on $A$ by $1+a^{-1}\phi a$, and $a^{-1}\phi a$
maps $A$ into $A_j$ and kills $A_j$; so this element lies in $U$, $u_v\in V_\gamma(E^+)$, and, $\psi$
being invariant under conjugation by $U$, $\sigma(\gamma^{-1}u_v\gamma)$ may be computed on it:
\[
\sigma(u_v)=e_j^*(v),\qquad\sigma(\gamma^{-1}u_v\gamma)=\sum_{i<\varrho}e_i^*\bigl(a^{-1}\phi(ae_{i+1})
\bigr)=e_j^*\bigl(a^{-1}\phi(ae_{j+1})\bigr)=\zeta\,e_j^*(a^{-1}v).
\]
(For $i>j$ the term vanishes because $a^{-1}\phi(\cdot)\in A_j$; for $i<j$ because $ae_{i+1}\in A_j$
is killed by $\phi$; and $ae_{j+1}\in\zeta e_{j+1}+A_j$.) By the remark, $e_j^*\circ a^{-1}=\zeta^{-1}
e_j^*$ on $A_j$: $a$ preserves $A_{j-1}$ and acts by $\zeta$ on $A_j/A_{j-1}$. By descending
induction, $\zeta^{-1}a$ is unitriangular on the flag.

So $\gamma_0=\zeta^{-1}\gamma$ is unitriangular on $A$ and induces on $Q$ an isometry that fixes
$z$, that is, an element $\eta$ of $H(E^+)$; then $\gamma_0\eta^{-1}\in U(E^+)$ and $\gamma_0\in
R_0(E^+)$. If $\gamma\in\Gamma_{c,s,c'}$, then $k=h_c^{-1}\gamma h_{c'}\in K$ and $ke_1=\zeta e_1$, so
$\zeta$ and $\zeta^{-1}$ are integral; $\zeta^{-1}k=h_c^{-1}\gamma_0h_{c'}\in K\cap R_0(\auA_f)$ has
its reduction in $\pbR$, and $s$ is this reduction times the residues of $\zeta$.
\end{proof}

Let
\[
\pbeGs=\{\gamma\in\Gamma:\pbg(\gamma)\ne0\}\ \cup\ \{\gamma\in\Gamma:\ \gamma\text{ agrees},\ \pbg(
\gamma)=0,\ \gamma A\ne A\},
\]
a stable set (for the last condition note $u\gamma u'A=u\gamma A$ and $uA=A$).

\begin{corollary}\label{pbe:star}
Assume \textup{(B1)--(B4)}. Let $\Omega$ avoid the scalar cosets, $s\in\Omega$ and $\flat\in\{\ ,
\flat\}$. Then $A^\flat_{\Gamma\smallsetminus R_0(E^+)}(c,s,c')=A^\flat_{\pbeGs}(c,s,c')$ for all
$c,c'$.
\end{corollary}
\begin{proof}
If $\gamma A=A$ then the first $\varrho$ columns of $Z(\gamma)$ vanish, so $\pbg(\gamma)=0$; hence
$\pbeGs\cap R_0(E^+)=\emptyset$. The set $(\Gamma\smallsetminus R_0(E^+))\smallsetminus\pbeGs$ is the
union of the stable set $T_1$ of its elements that do not agree, which contributes $0$ by
Lemma~\ref{pbe:vanish}, and the set $T_2$ of the $\gamma\notin R_0(E^+)$ that agree and have
$\gamma A=A$. By Lemma~\ref{pbe:scalar} an element of $T_2\cap\Gamma_{c,s,c'}$ lies in $\zeta R_0(E^+)$
with $s\in\zeta\pbR$; as $s\in\Omega$, $\zeta=1$, a contradiction. So $T_2\cap\Gamma_{c,s,c'}=\emptyset$.
\end{proof}

\subsection{The cost of the cut}\label{pbe:ss-cut}

\begin{proposition}[Cost of the cut]\label{pbe:cut}
Assume Theorem~\ref{pbe:ds} and let $W,W^\flat$ be as in Theorem~\ref{pbe:bessel}. For every stable
$T\subseteq\Gamma$ and every $\Omega\subseteq\auG$,
\[
\sum_{c,c'\in\mathcal C}\ \sum_{s\in\Omega}\bigl\|A_T(c,s,c')-A_T^\flat(c,s,c')\bigr\|^2\le b^{-40}\,
\pbDid\,\nu^2 .
\]
\end{proposition}
\begin{proof}
Fix $c$ and $s$. Grouping the terms in orbits of $U_c$ as in the proof of
Proposition~\ref{pb:unfold},
\[
\sum_{c'}\|A_T(c,s,c')-A_T^\flat(c,s,c')\|\le\int_{U_c\backslash G_\infty}\ \sum_{c'}\ \sum_{\gamma\in
\Gamma_{c,s,c'}/U_{c'}}\bigl\|W(y)W(\gamma^{-1}y)^*-W^\flat(y)W^\flat(\gamma^{-1}y)^*\bigr\|\,dy .
\]
The difference in the integrand is $(1-\chi(y)\chi(\gamma^{-1}y))\,W(y)W(\gamma^{-1}y)^*$ with $\chi=
\prod_\tau\chi_\tau$. It vanishes unless $\chi_\tau<1$ at some component of $y$ or of $\gamma^{-1}y$,
in which case $\pbeL(y)\pbeL(\gamma^{-1}y)\ge b$ by Lemma~\ref{pbe:len}(b); and by \eqref{pbe:eq-decay}
its norm is at most
\[
\mathsf A(y)\,\mathsf a_\infty\,\bigl(\pbeL(y)\pbeL(\gamma^{-1}y)\bigr)^{-2r^3f}\,\pbeL(\gamma^{-1}y)^{
-r^3f}\le b^{-2r^3f}\,\mathsf a_\infty\,\mathsf A(y)\,\pbeL(\gamma^{-1}y)^{-r^3f}.
\]
By Proposition~\ref{pbe:count}, the invariance of $\mathsf A$ under $U_\infty$, $\pbvol(U_c\backslash
U_\infty)=\nu$ and \eqref{pbe:eq-decay},
\[
\sum_{c'}\|A_T(c,s,c')-A_T^\flat(c,s,c')\|\le b^{-2r^3f}\cdot2b^{12n^3rf}\cdot\nu\,\mathsf a_\infty\int_{
U_\infty\backslash G_\infty}\mathsf A\le\nu\,b^{-2r^3f+12n^3rf+400r^2f+2kf+1}\le\nu\,b^{-k^3f};
\]
the last step is $12n^3r+400r^2+2k+1+k^3\le2r^3$, which holds because $k<r$ and $12n^3+500\le r$.
The sum of the squares over $c'$ is at most the square of the sum. Summing over $c$ and over all
$s\in\auG$, with $\#\auG\le M^{r^2f}\le b^{22r^2f}$, $\#\mathcal C\le\pbDid$ and $2k^3f\ge22r^2f+40$ (as
$k\ge r/2$, $r\ge100$), we get the proposition.
\end{proof}

\subsection{Links, and the three estimates for pairs of terms}\label{pbe:ss-pillars}

Let $Q=A^\perp/A=X_0\perp Ez$ and let $Q^+=X_0\perp Ez^+$ be the same vector space with the positive
definite form in which $z^+$ has norm $N$. Write $\pbeAp=A+Ez$, the span of the frame. Recall that
the link of $\gamma$ is $\Lambda_\gamma=\{([t],[\gamma^{-1}t]):t\in(A+\gamma A)^\perp\}\subseteq Q^2$.

\begin{lemma}[Links]\label{pbe:link}
Let $\gamma,\gamma'\in\Gamma$.
\begin{enumerate}
\item[(a)] $\pbg(\gamma)$ is, up to sign, the Gram determinant of the $2n$ vectors $f_0,\dots,
f_\varrho,\gamma f_0,\dots,\gamma f_\varrho$, and $\pbg(\gamma)=0$ if and only if $\pbeAp\cap\gamma
\pbeAp\ne0$. Both $\pbg$ and the link are invariant under $\gamma\mapsto u\gamma u'$, $u,u'\in
U(E^+)$.
\item[(b)] $\dim\Lambda_\gamma=k+1$, and $\Lambda_\gamma$ is the set of the pairs $(y_0+\lambda'z,\
y_0'+\lambda z)$ with $\Upsilon_\gamma(y_0+\lambda z^+)=y_0'+\lambda'z^+$, for an isometry
$\Upsilon_\gamma$ of $Q^+$. So $\Lambda_\gamma=\Lambda_{\gamma'}$ if and only if $\Upsilon_\gamma=
\Upsilon_{\gamma'}$.
\item[(c)] Let $X_0^\gamma=X_0\cap\Upsilon_\gamma^{-1}X_0$. Then $\dim X_0^\gamma=k-1+\dim(\pbeAp\cap
\gamma\pbeAp)-\dim(A\cap\gamma A)$. In particular $\dim X_0^\gamma=k-1$ if $\pbg(\gamma)\ne0$, and
$\dim X_0^\gamma=k$ if $\pbg(\gamma)=0$ and $\pbeAp\cap\gamma\pbeAp\not\subseteq A$.
\item[(d)] Let $\pbeCL(\gamma,\gamma')$ be the space of the $x\in X_0$ such that $(x,x')\in\Lambda_
\gamma\cap\Lambda_{\gamma'}$ for some $x'\in X_0$, and $i(\gamma,\gamma')=k+1-\dim(\Lambda_\gamma\cap
\Lambda_{\gamma'})$. Then $k-\dim\pbeCL(\gamma,\gamma')\le i(\gamma,\gamma')+k-\dim X_0^\gamma$.
\end{enumerate}
\end{lemma}
\begin{proof}
(a) The Gram matrix of the $2n$ vectors is $\bigl(\begin{smallmatrix}\mathsf D&Z^*\\Z&\mathsf D
\end{smallmatrix}\bigr)$ with $Z=Z(\gamma)$ up to transposition and $\mathsf D$ the Gram matrix of
the frame, whose only non-zero entry is $-N$ in the last place. For square matrices $\mathsf B,
\mathsf C$ one has $\det\bigl(\begin{smallmatrix}\mathsf B&\mathsf D\\\mathsf D&\mathsf C
\end{smallmatrix}\bigr)=\det\mathsf B\det\mathsf C-N^2\delta_{n-1}(\mathsf B)\delta_{n-1}(\mathsf C)$
(for invertible $\mathsf B$ this is $\det\mathsf B\cdot\det(\mathsf C-\mathsf D\mathsf B^{-1}\mathsf
D)$, where $\mathsf D\mathsf B^{-1}\mathsf D$ has the only entry $N^2\delta_{n-1}(\mathsf B)/\det
\mathsf B$ in the last place; the general case follows, both sides being polynomials). So the Gram
determinant is $\pm(\det Z\,\aubar{\det Z}-N^2\delta_{n-1}\aubar{\delta_{n-1}})=\pm\pbg(\gamma)$. If $\pbeAp\cap\gamma\pbeAp\ne0$ the $2n$ vectors are dependent and
the determinant is $0$. Conversely, let it be $0$ with $\pbeAp\cap\gamma\pbeAp=0$; then the span
$\mathsf V$ of the $2n$ vectors is degenerate: there is $v\ne0$ in $\mathsf V\cap\mathsf V^\perp$.
At a real place $\pbeAp$ is a negative semidefinite subspace of dimension $n$, the largest
dimension possible in signature $(k+\varrho,n)$; $v$ is orthogonal to $\pbeAp$ and to itself, so
$\pbeAp+\C v$ is negative semidefinite and $v\in\pbeAp$. In the same way $v\in\gamma\pbeAp$, a
contradiction. For the invariance: $uf$ and $u'f$ are obtained from $f$ by unitriangular matrices,
so the Gram determinant of $(f,u\gamma u'f)$, which is that of $(u^{-1}f,\gamma u'f)$, is that of
$(f,\gamma f)$; and for $t\perp A+u\gamma A$ the vector $t_1=u^{-1}t$ is orthogonal to $A+\gamma A$,
$[t_1]=[t]$ and $[(u\gamma u')^{-1}t]=[\gamma^{-1}t_1]$ in $Q$.

(b) Let $I=A\cap\gamma A$. The space $(A+\gamma A)^\perp$ has dimension $r-2\varrho+\dim I=k+1+\dim
I$ and contains $I$, which is the kernel of $t\mapsto([t],[\gamma^{-1}t])$. So $\dim\Lambda_\gamma=k+1$.
Since $h(t,t)=h(\gamma^{-1}t,\gamma^{-1}t)$, a pair $(y_0+\lambda'z,y_0'+\lambda z)$ of $\Lambda_\gamma$
satisfies $|y_0|^2-N|\lambda'|^2=|y_0'|^2-N|\lambda|^2$, that is, $y_0+\lambda z^+$ and $y_0'+\lambda'
z^+$ have the same norm in $Q^+$. So the two projections of the $(k+1)$-dimensional space of the
pairs $(y_0+\lambda z^+,y_0'+\lambda'z^+)$ to $Q^+$ are injective, hence bijective, and this space is
the graph of an isometry.

(c) $X_0^\gamma$ is the set of the first members of the pairs of $\Lambda_\gamma$ with $\lambda=
\lambda'=0$, that is, with $t\perp z$ and $\gamma^{-1}t\perp z$; so it is the image of $(\pbeAp+
\gamma\pbeAp)^\perp$ under $t\mapsto[t]$, whose kernel is $I$ (note $I\perp\pbeAp+\gamma\pbeAp$; and if $[t]=0$ then
$[\gamma^{-1}t]=0$ by (b): the pair $([t],[\gamma^{-1}t])$ has $y_0=0$ and $\lambda=-h(t,\gamma z)/N=0$ because $t\perp\gamma\pbeAp$, so it is $0$ by the injectivity of the first projection; hence $t\in I$).
Hence $\dim X_0^\gamma=r-2n+\dim(\pbeAp\cap\gamma\pbeAp)-\dim I$, and $r-2n=k-1$. The last claims
follow from (a) and from $\dim X_0^\gamma\le k$.

(d) In the graph of $\Upsilon_\gamma$, of dimension $k+1$, the graph of $\Upsilon_\gamma|X_0^\gamma$
and the image of $\Lambda_\gamma\cap\Lambda_{\gamma'}$ meet in a space of dimension at least $\dim
X_0^\gamma+(k+1-i)-(k+1)$, and the first projection maps it injectively into $\pbeCL(\gamma,
\gamma')$.
\end{proof}

Put
\[
\pbet=b^{kf/4},\qquad\pbenp=\frac{\#\pbR\cdot\#\mathcal C}{N^{kf}},\qquad\mathcal D(\gamma)=\begin{cases}
\pbDmain&\text{if }\pbg(\gamma)\ne0,\\ \pbet^{-16}\,\pbDid&\text{if }\pbg(\gamma)=0.\end{cases}
\]
Let $U^0\subseteq U$ be the subgroup that fixes the frame and $U_c^0=U_c\cap U^0$.

\begin{theorem}[Size, gain and count]\label{pbe:pillars}
Let $W^\flat$ be as in Theorem~\ref{pbe:bessel}. For all classes $c,c'$ there are a set $\pbeN(c,c')
\subseteq\pbeGs$ that contains at most one element of each double coset $U_c^0\gamma U_{c'}^0$, a
function $\mathsf G=\mathsf G_{c,c'}\colon\pbeN(c,c')\to(0,b^{n^3f}]$ (the \emph{gauge}), and for
$\gamma,\gamma'\in\pbeN(c,c')$ over a common block a real number $\mathsf d(\gamma,\gamma')=\mathsf d(
\gamma',\gamma)\ge1$ (the discriminant of the common lattice), with the following properties. Put
$\pbesz(\gamma)=\pbet^2\bigl(\mathcal D(\gamma)/(\pbenp\,\mathsf G(\gamma)^k)\bigr)^{1/2}$.
\begin{enumerate}
\item[(R)] If $\gamma\in\pbeN(c,c')$ and $\pbg(\gamma)=0$, then $\pbeAp\cap\gamma\pbeAp\not\subseteq
A$.
\item[(S)] \emph{(Size.)} For every stable $T\subseteq\pbeGs$ and every block $(c,s,c')$,
\[
\bigl\|A_T^\flat(c,s,c')\bigr\|\le\nu\sum_{\gamma\in\pbeN(c,c')\cap T\cap\Gamma_{c,s,c'}}\pbesz(\gamma).
\]
\item[(G)] \emph{(Gain.)} For $\gamma,\gamma'\in\pbeN(c,c')\cap\Gamma_{c,s,c'}$, with $i=i(\gamma,
\gamma')$,
\[
\bigl(\mathsf d(\gamma,\gamma')\cdot\mathsf x\bigr)^2\le\Bigl(\frac2M\Bigr)^{2if}\,\mathsf G(\gamma)\,
\mathsf G(\gamma'),\qquad\mathsf x=\begin{cases}\mathsf{mn}(s)&\text{if }i=0\text{ and }\pbg(\gamma)
\ne0,\\1&\text{otherwise.}\end{cases}
\]
\item[(C)] \emph{(Count.)} For every integer $0\le l\le n^3$,
\[
\sum_{c,c'\in\mathcal C}\ \sum_{(\gamma,\gamma')}\mathsf d(\gamma,\gamma')^{-k}\le\pbet^{\,l+2}\,\pbenp ,
\]
the inner sum over the pairs of elements of $\pbeN(c,c')$ over a common block with $k-\dim\pbeCL(
\gamma,\gamma')=l$ and $\mathsf d(\gamma,\gamma')\le b^{n^3f}$.
\end{enumerate}
\end{theorem}
\smallskip\noindent The exact form of the theorem, with the set $\pbeN(c,c')$ restricted to agreeing blocks and with the gauge, the common lattice and nearness defined, is Theorem~\ref{ppi:pillars}; every use made below of the present theorem is covered by that form: Proposition~\ref{ppi:check} shows that Lemmas~\ref{pbe:rank} and~\ref{pbe:pairs}, Proposition~\ref{pbe:two} and Theorem~\ref{pbe:XL1} hold with Theorem~\ref{ppi:pillars} in its place, with the same proofs. In particular \textup{(S)} is used only for the stable sets $T_0$, $T_\Lambda$ and $\pbeGs$, and in that form the sum in \textup{(S)} runs over elements of agreeing blocks only; the elements of a non-agreeing block contribute $0$ to $A^\flat_T$ by Lemma~\ref{pbe:vanish}. \textup{(R)} is proved in Proposition~\ref{ppi:cell}, \textup{(G)} in \S\ref{pbg:sec}, \textup{(C)} in \S\ref{pbc:sec}, and \textup{(S)} in \S\ref{pbs:sec} from the open-cell bound (Theorem~\ref{pbo:O1}).\par\smallskip

\subsection{Proof of \textup{(X)} and \textup{(L1)} from Theorem~\ref{pbe:pillars}}\label{pbe:ss-XL1}

We say that $\gamma,\gamma'$ are \emph{in the same link class} if $\pbg(\gamma)\ne0$, $\pbg(\gamma')
\ne0$ and $\Lambda_\gamma=\Lambda_{\gamma'}$.

\begin{lemma}\label{pbe:rank}
Assume Theorem~\ref{pbe:pillars}. Let $\gamma,\gamma'\in\pbeN(c,c')\cap\Gamma_{c,s,c'}$, $i=i(\gamma,
\gamma')$ and $l=k-\dim\pbeCL(\gamma,\gamma')$. Then:
\begin{enumerate}
\item[(a)] $l\le i+1$, and $l\le i$ if $\pbg(\gamma)=0$;
\item[(b)] if $i=0$, then $\pbg(\gamma)=0$ if and only if $\pbg(\gamma')=0$;
\item[(c)] $\mathsf d(\gamma,\gamma')\le b^{n^3f}$ and $l\le n^3$.
\end{enumerate}
\end{lemma}
\begin{proof}
(a) By Lemma~\ref{pbe:link}(c),(d) and (R), $k-\dim X_0^\gamma$ is $1$ if $\pbg(\gamma)\ne0$ and $0$
otherwise. (b) If $i=0$ then $\Upsilon_\gamma=\Upsilon_{\gamma'}$ and $X_0^\gamma=X_0^{\gamma'}$, and
the dimension of this space decides whether $\pbg$ vanishes. (c) By (G), with $\mathsf x\ge1$ and
$\mathsf G\le b^{n^3f}$, $\mathsf d^2\le(2/M)^{2if}b^{2n^3f}$. So $\mathsf d\le b^{n^3f}$, and, as
$\mathsf d\ge1$ and $M/2\ge b^{18}$, $18i\le n^3$; hence $l\le i+1\le n^3$.
\end{proof}

\begin{lemma}[Sums over pairs]\label{pbe:pairs}
Assume Theorem~\ref{pbe:pillars}. For all $c,c'$ let $\mathcal Q(c,c')$ be a set of pairs $(\gamma,
\gamma')$ of elements of $\pbeN(c,c')$ over a common block, and let $B\ge0$. Suppose that for every
such pair there are $\eta>0$ and $\mathsf D\ge\max(\mathcal D(\gamma),\mathcal D(\gamma'))$ with
\[
(\mathsf d(\gamma,\gamma')\,\eta)^2\le\mathsf G(\gamma)\mathsf G(\gamma')\qquad\text{and}\qquad\pbet^{\,
l+6}\,\mathsf D\,\eta^{-k}\le B,\quad l=k-\dim\pbeCL(\gamma,\gamma').
\]
Then $\sum_{c,c'}\sum_{(\gamma,\gamma')\in\mathcal Q(c,c')}\pbesz(\gamma)\pbesz(\gamma')\le\pbet\,B$.
\end{lemma}
\begin{proof}
For such a pair, $\pbesz(\gamma)\pbesz(\gamma')\le\pbet^4\mathsf D\,\pbenp^{-1}(\mathsf G(\gamma)
\mathsf G(\gamma'))^{-k/2}\le\pbet^4\mathsf D\eta^{-k}\pbenp^{-1}\mathsf d^{-k}\le B\,\pbet^{-l-2}
\pbenp^{-1}\,\mathsf d(\gamma,\gamma')^{-k}$. By Lemma~\ref{pbe:rank}(c) every pair is counted in (C)
for its $l\le n^3$, so the pairs with a given $l$ contribute at most $B$, and the total is at most
$(n^3+1)B\le b^4B\le\pbet B$.
\end{proof}

\begin{proposition}\label{pbe:two}
Assume Theorem~\ref{pbe:pillars}.
\begin{enumerate}
\item[(a)] The sum of $\pbesz(\gamma)\pbesz(\gamma')$ over all $c,c'$ and all pairs of elements of
$\pbeN(c,c')$ over a common block that are \emph{not} in the same link class is at most $b^{-kf}
\pbDid$.
\item[(b)] Let $\mathsf z\ge1$. The same sum over the pairs \emph{in} the same link class and over a
block $(c,s,c')$ with $\mathsf{mn}(s)\ge\mathsf z$ is at most $(b^{2f}/\mathsf z)^k\,\pbDmain$.
\end{enumerate}
\end{proposition}
\begin{proof}
We apply Lemma~\ref{pbe:pairs}. Recall $\pbet^4=b^{kf}$, $b^{11kf}\pbDid\le\pbDmain\le(M/b^6)^{kf}\pbDid$
(Lemma~\ref{pb:sizes}); in particular $\mathcal D(\gamma)\le\pbDmain$ always.

(a) Let $B=\pbet^{-10}\pbDid$; the bound will be $\pbet^{-9}\pbDid\le\pbet^{-4}\pbDid$. Let $(\gamma,
\gamma')$ be a pair, $i=i(\gamma,\gamma')$. \emph{If $i=0$}, the links are equal, so by
Lemma~\ref{pbe:rank}(b) and the hypothesis $\pbg(\gamma)=\pbg(\gamma')=0$, and $l=0$ by
Lemma~\ref{pbe:rank}(a). Take $\eta=1$ (by (G)) and $\mathsf D=\pbet^{-16}\pbDid$: $\pbet^6\mathsf D=B$.
\emph{If $i\ge1$}, take $\eta=(M/2)^{if}$ (by (G)) and $\mathsf D=\pbDmain$; $l\le i+1$, and with
$\mathsf y=\pbet(2/M)^{kf}\le1$,
\[
\pbet^{\,l+6}\,\pbDmain\,\eta^{-k}\le\pbet^7\mathsf y^i\pbDmain\le\pbet^7\mathsf y\Bigl(\frac M{b^6}
\Bigr)^{kf}\pbDid=\pbet^8\Bigl(\frac2{b^6}\Bigr)^{kf}\pbDid\le\pbet^{8-20}\pbDid\le B .
\]
(b) Here $i=0$ and $\pbg(\gamma)\ne0$, so (G) gives $(\mathsf d\,\mathsf z)^2\le(\mathsf d\,\mathsf{
mn}(s))^2\le\mathsf G(\gamma)\mathsf G(\gamma')$, and $l\le1$. Take $\eta=\mathsf z$, $\mathsf D=
\pbDmain$, $B=\pbet^7\pbDmain\mathsf z^{-k}$; the bound is $\pbet^8\mathsf z^{-k}\pbDmain=(b^{2f}/
\mathsf z)^k\pbDmain$.
\end{proof}

\begin{theorem}[Theorem~\ref{pb:amp}\textup{(X)} and \textup{(L1)}]\label{pbe:XL1}
Assume Theorems~\ref{pbe:ds} and~\ref{pbe:pillars}. Then, for the Bessel functions of
Theorem~\ref{pbe:bessel}, parts \textup{(X)} and \textup{(L1)} of Theorem~\ref{pb:amp} hold.
\end{theorem}
\begin{proof}
(L1) Let $\mathsf{mn}\ge\mathsf z$ on $\Omega$. The sets $T_\Lambda$ are stable subsets of $\pbeGs$.
By (S),
\[
\|A^\flat_{T_\Lambda}(c,s,c')\|^2\le\nu^2\sum_{\gamma,\gamma'\in\pbeN(c,c')\cap T_\Lambda\cap\Gamma_{c,
s,c'}}\pbesz(\gamma)\pbesz(\gamma').
\]
The pairs that occur are in the same link class and over the block $(c,s,c')$, and a pair determines
$s$ and $\Lambda$; so, summing over $c,c'$, $s\in\Omega$ and $\Lambda$, $\pbLM(\Omega)\le\nu^2(b^{2f}/
\mathsf z)^k\pbDmain$ by Proposition~\ref{pbe:two}(b).

(X) Let $\Omega$ avoid the scalar cosets and $\pbLM(\Omega)\le C\nu^2$. Fix a block $(c,s,c')$ with
$s\in\Omega$. The set $\pbeGs$ is the disjoint union of the stable sets $T_\Lambda$ and $T_0=\{\gamma
\in\pbeGs:\pbg(\gamma)=0\}$. The amplitude sums converge absolutely, so $A^\flat_{\pbeGs}=\sum_\kappa
A^\flat_{T_\kappa}$ over the pieces $\kappa\in\{0\}\cup\{\Lambda\}$, and
\[
\|A^\flat_{\pbeGs}(c,s,c')\|^2\le\Bigl(\sum_\kappa\|A^\flat_{T_\kappa}\|\Bigr)^2=\sum_\Lambda\|A^\flat_{
T_\Lambda}\|^2+\sum_{(\kappa,\kappa')\ne(\Lambda,\Lambda)}\|A^\flat_{T_\kappa}\|\,\|A^\flat_{T_{\kappa'}}\|.
\]
By (S) each product in the last sum is at most $\nu^2\sum\pbesz(\gamma)\pbesz(\gamma')$ over $\gamma\in
\pbeN(c,c')\cap T_\kappa\cap\Gamma_{c,s,c'}$, $\gamma'\in\pbeN(c,c')\cap T_{\kappa'}\cap\Gamma_{c,s,c'}$;
these pairs are not in the same link class, and a pair determines $s$, $\kappa$ and $\kappa'$. (Here (S) is applied to $T_\kappa$, which may contain elements of non-agreeing blocks; by Lemma~\ref{pbe:vanish} these contribute $0$ to $A^\flat_{T_\kappa}$, so only elements of $\pbeN(c,c')$, which lie over agreeing blocks, enter the bound.)
Summing over the blocks and using Proposition~\ref{pbe:two}(a),
\[
\sum_{c,c'}\sum_{s\in\Omega}\|A^\flat_{\pbeGs}(c,s,c')\|^2\le\pbLM(\Omega)+\nu^2b^{-kf}\pbDid\le\nu^2\bigl(C
+b^{-kf}\pbDid\bigr).
\]
By Proposition~\ref{pb:besselterms}, Corollary~\ref{pbe:star}, the triangle inequality in $\ell^2$
and Proposition~\ref{pbe:cut} (for $T=\pbeGs$),
\[
\Bigl(\sum_{s\in\Omega}\pbHS{P_s-\pbrpsi(s)}^2\Bigr)^{1/2}=\frac1\nu\Bigl(\sum_{c,c',s}\|A_{\pbeGs}\|^2
\Bigr)^{1/2}\le b^{-20}\sqrt{\pbDid}+\sqrt{C+b^{-kf}\pbDid}\le\sqrt C+b^{-15}\sqrt{\pbDid},
\]
as $kf\ge40$.
\end{proof}

\subsection{The inputs of this section}\label{pbe:ss-remains}

Hypothesis~\ref{pb:bessel} and Theorem~\ref{pb:amp}\textup{(X)}, \textup{(L1)} rest on the following
statements, proved elsewhere in the paper:
\begin{enumerate}
\item[(i)] Theorem~\ref{pbe:ds}: proved in \S\S\ref{pds:sec}, \ref{pco:sec}, \ref{pbp:sec}, modulo the published results cited there.
\item[(ii)] Theorem~\ref{pbe:pillars}: \textup{(O1)} (Theorem~\ref{pbo:O1}, for \textup{(S)}), and the statements \textup{(G$'$)}, \textup{(C1)}, \textup{(C2)} of \S\ref{ppi:sec}, proved in \S\S\ref{pbo:sec}, \ref{pbg:sec}, \ref{pbc:sec}; \textup{(S)} also uses Lemma~\ref{pb:covol}.
\item[(iii)] The inequality \eqref{pbe:eq-kappa} for the generator $\varpi$ (Lemma~\ref{pds:kappa}), and the
finite dimension of $\auAz$ (Proposition~\ref{pz:types}(a)); $\auAz$ is by definition the space of bounded
solutions of \eqref{pbe:eq-repro} (\S\ref{ss:au-spaces}).
\end{enumerate}


%% ===== proofs/B-main-term.tex =====
% proofs/B-main-term.tex
% The induced character (Theorem pb:est (E4)) and the main term (Theorem pb:amp (L2)).
% Labels of this file start with pbm:.
\providecommand{\auO}{\mathcal O}
\providecommand{\auPl}{\Phi}
\providecommand{\auS}{\mathcal S}
\providecommand{\auG}{\overline G}
\providecommand{\auDX}{D_X}
\providecommand{\auNm}{\operatorname{Nm}}
\providecommand{\auU}{\mathrm U}
\providecommand{\auGL}{\mathrm{GL}}
\providecommand{\auA}{\mathbb A}
\providecommand{\auF}{\mathbb F}
\providecommand{\aurd}{\mathrm{rd}_E}
\providecommand{\aubar}[1]{\overline{#1}}
\providecommand{\auf}[1]{\texttt{#1}}
\providecommand{\panot}[1]{\par\smallskip{\raggedright\noindent\textbf{[Not proved here:} #1\textbf{]}\par}\smallskip}
\providecommand{\pbR}{\overline R_0}
\providecommand{\pbUb}{\overline U}
\providecommand{\pbHb}{\overline H}
\providecommand{\pbvol}{\operatorname{vol}}
\providecommand{\pbDid}{\mathcal D_{\mathrm{id}}}
\providecommand{\pbDmain}{\mathcal D_{\mathrm{main}}}
\providecommand{\pbdeg}{\operatorname{deg}}
\providecommand{\pbmass}{\operatorname{mass}}
\providecommand{\pbFix}{\operatorname{Fix}}
\providecommand{\pbg}{\mathsf g}
\providecommand{\pbLM}{\mathrm{LM}}
\providecommand{\pbnat}{\natural}
\providecommand{\pbip}[2]{\langle#1,#2\rangle}
% macros of this file (prefix \pbm)
\providecommand{\pbmD}{\mathsf D}
\providecommand{\pbmPhi}{\varphi}
\providecommand{\pbmGf}{G^{\flat}}
\providecommand{\pbmXp}{X^{+}}
\providecommand{\pbmOrb}{\mathcal O}
\providecommand{\pbmKl}{\mathsf{Kl}}
\providecommand{\pbmT}{\mathcal T}
\providecommand{\pbmAT}{\mathbb A_T}
\providecommand{\pbmRT}{R_T}
\providecommand{\pbmE}{\mathsf E}
\providecommand{\pbmn}{\mathsf n}
\providecommand{\pbmden}{\mathsf{den}}
\providecommand{\pbmvol}{\mathsf v}
\providecommand{\pbmdec}{\mathsf d}
\providecommand{\pbmaS}{\mathsf a_*}
\providecommand{\pbmrS}{\mathsf r_*}
\providecommand{\pbmnuS}{\nu_*}
\providecommand{\pbmcG}{c_G}
\providecommand{\pbmTh}{\Theta}
\providecommand{\pbmSt}{\operatorname{Stab}}

\section{The induced character and the main term}\label{pbm:sec}

This section concerns two of the statements to which \S\ref{pb:sec} reduced Theorem~\ref{thm:srcB}:
the inequality \textup{(E4)} of Theorem~\ref{pb:est}, and the statement \textup{(L2)} of
Theorem~\ref{pb:amp}. We use the notation of \S\ref{pb:ss-setting}.

\emph{What is proved.}
\begin{itemize}
\item[(I)] \emph{The inequality \textup{(E4)} is proved in full} (\S\ref{pbm:ss-E4},
Theorem~\ref{pbm:E4}). The proof is short and does not use the analysis of the big cell: the
function that counts the elements of $\pbUb\cap h^{-1}\pbUb h$ on which $\psi_U$ and its conjugate
by $h$ agree is of positive type, so Lemma~\ref{pb:singular} applies to it; two exact recursions,
which remove one hyperbolic pair, or one vector of the frame, at a time, and a crude count do the
rest.
\item[(II)] \emph{The statement \textup{(L2)}} is deduced, by complete arguments, from one precisely
stated theorem on local objects (Theorem~\ref{pbm:inputs}, proved in \S\S\ref{pfr:sec}--\ref{pfg:sec})
and the value of the covolume $\nu$ (Lemma~\ref{pb:covol}): the
existence, at each real place and at each finite set of finite places, of a system of forms on
pairs of heads with explicit bounds, and an inequality (\emph{Siegel's step}) that bounds sums of
squares of class sums by the top forms. We prove the description of the terms over
open elements (big cell, heads, links; Propositions~\ref{pbm:cell} and~\ref{pbm:classsum}); the
identification of an orbit sum with a Kloosterman sum times an orbital integral
(Proposition~\ref{pbm:orbit}); the sieve by one-sided Poisson summation, with all its constants
(Lemma~\ref{pbm:poisson}, Theorem~\ref{pbm:sieve}); the evaluation of the constant of the main
term, which is where $\pbDmain=\auDX^2/(\pbdeg\cdot\pbmass_r)$ comes from (Lemma~\ref{pbm:const});
and the deduction of \textup{(L2)} (Theorem~\ref{pbm:L2}).
\end{itemize}
\subsection{The induced character: proof of \textup{(E4)}}\label{pbm:ss-E4}

\subsubsection*{Chain data}
Let $\auF$ be a finite field with $q$ elements and $\psi_1,\psi_2$ two non-trivial additive
characters of $\auF$. Let $i,i'\ge0$ and $m\ge1$ be integers. A \emph{chain datum of type $(i,i')$}
is a vector space $V$ over $\auF$ with a basis
\[
e_1,\dots,e_i;\qquad\text{a basis of a subspace }X'\text{ of dimension }m;\qquad e^*_{i'},\dots,e^*_1,
\]
together with a vector $z\in X'$, $z\ne0$, and a linear form $\zeta\ne0$ on $X'$, extended by zero
on the $e_a$ and the $e_a^*$. Write $r=\dim V=i+m+i'$, $\epsilon_a$ and $\epsilon_a^*$ for the
coordinate forms of $e_a$ and $e_a^*$, $A_a=\langle e_1,\dots,e_a\rangle$ $(0\le a\le i)$, $A=A_i$
and $B_a=A\oplus X'\oplus\langle e^*_{i'},\dots,e^*_{a+1}\rangle$ $(0\le a\le i')$, so that
$B_{i'}=A\oplus X'$ and $B_0=V$. Let $U=U_\pbmD$ be the group of the $u\in\auGL(V)$ with
\[
(u-1)A_a\subseteq A_{a-1}\ (1\le a\le i),\qquad(u-1)B_{i'}\subseteq A,\qquad(u-1)B_a\subseteq B_{a+1}\
(0\le a<i');
\]
it is the unipotent radical of the parabolic subgroup with blocks $1,\dots,1,m,1,\dots,1$. The
function
\[
\psi_\pbmD(u)=\psi_1\Bigl(\sum_{a=1}^{i-1}\epsilon_a(ue_{a+1})+\epsilon_i(uz)\Bigr)\,\psi_2\Bigl(\sum_{a=1}^{
i'-1}\epsilon^*_{a+1}(ue^*_a)+\zeta(ue^*_{i'})\Bigr)
\]
(the first factor is $1$ if $i=0$, the second is $1$ if $i'=0$) is a character of $U$: its
arguments are sums of entries of $u$ just above the diagonal blocks, and these are additive on
$U$. We call $e_1$ (or $z$, if $i=0$) the \emph{first vector of the left chain}. For $h\in\auGL(V)$
put
\[
f_\pbmD(h)=\sum_{u\in U\cap h^{-1}Uh}\psi_\pbmD(huh^{-1})\,\aubar{\psi_\pbmD(u)},\qquad\pbmPhi(\pbmD)=\frac1{\#
\auGL(V)}\sum_{h\in\auGL(V)}f_\pbmD(h).
\]

\begin{lemma}\label{pbm:pos}
Let $\pbmD$ be a chain datum, $U=U_\pbmD$, $\psi=\psi_\pbmD$, $f=f_\pbmD$ and $\mathsf G=\auGL(V)$.
\begin{enumerate}
\item[(a)] $f(h)$ is $\#(U\cap h^{-1}Uh)$ or $0$; in particular $f\ge0$. Moreover $f(uhu')=f(h)$
for $u,u'\in U$ and $f(h^{-1})=f(h)$.
\item[(b)] $f$ is of positive type on $\mathsf G$ (in the sense of \S\ref{pb:ss-finite}).
\item[(c)] Let $\tilde\psi$ be $\psi$ on $U$ and $0$ on $\mathsf G\smallsetminus U$, and
$\psi^{\pbnat}(s)=(\#\mathsf G)^{-1}\sum_{g\in\mathsf G}\tilde\psi(gsg^{-1})$. Then
$\sum_{s\in\mathsf G}|\psi^{\pbnat}(s)|^2=\pbmPhi(\pbmD)$.
\end{enumerate}
\end{lemma}
\begin{proof}
(a) $u\mapsto\psi(huh^{-1})\aubar{\psi(u)}$ is a character of the group $U\cap h^{-1}Uh$, and the
sum of a character over a finite group is the order of the group or $0$. The substitutions
$u\mapsto u'^{-1}uu'$ and the invariance of $\psi$ under conjugation by $U$ give $f(uhu')=f(h)$.
With $u=h^{-1}u_1h$ we get $f(h^{-1})=\sum_{u_1\in U\cap h^{-1}Uh}\psi(u_1)\aubar{\psi(hu_1h^{-1})}=
\aubar{f(h)}=f(h)$.
(b) In the space of functions on $\mathsf G$, with the orthonormal basis of the delta functions
$\delta_g$, let $v_a=\sum_{u\in U}\aubar{\psi(u)}\,\delta_{aua^{-1}}$ for $a\in\mathsf G$. Then, with
$s=b^{-1}a$,
\[
\pbip{v_a}{v_b}=\sum_{u,u'\in U}\psi(u)\aubar{\psi(u')}\,[aua^{-1}=bu'b^{-1}]=\sum_{u\in U\cap s^{-1}Us}
\psi(u)\aubar{\psi(sus^{-1})}=\aubar{f(b^{-1}a)}=f(a^{-1}b)
\]
by (a). (c) Expanding the square and substituting $s=g'^{-1}ug'$ with $u\in U$,
\[
\sum_s|\psi^{\pbnat}(s)|^2=\frac1{(\#\mathsf G)^2}\sum_{g,g'}\sum_{u\in U}\tilde\psi\bigl(gg'^{-1}u\,g'g^{
-1}\bigr)\aubar{\psi(u)}=\frac1{\#\mathsf G}\sum_{h\in\mathsf G}\ \sum_{u\in U,\ huh^{-1}\in U}\psi(huh^{-1})
\aubar{\psi(u)} .\qedhere
\]
\end{proof}

\begin{lemma}[Removing a hyperbolic pair]\label{pbm:pair}
Let $\pbmD$ be a chain datum of type $(i,i')$ with $i,i'\ge1$ on $V$, and let $\pbmD'$ be the chain
datum of type $(i-1,i'-1)$ on $V'=\langle e_2,\dots,e_i\rangle\oplus X'\oplus\langle e^*_{i'},\dots,
e^*_2\rangle$ with the same $z$, $\zeta$, $\psi_1$, $\psi_2$. Let $\Omega$ be the set of the
$h\in\auGL(V)$ with $\epsilon_1^*(he_1)\ne0$ and $\epsilon_1^*(h^{-1}e_1)\ne0$. Then
\[
\sum_{h\in\Omega}f_\pbmD(h)=\#\Omega\cdot\pbmPhi(\pbmD').
\]
\end{lemma}
\begin{proof}
Let $H_1=\ker\epsilon_1^*=B_1$, let $P_1$ be the stabiliser of the flag $\auF e_1\subset H_1$,
$N_1$ its unipotent radical (the $n$ with $(n-1)V\subseteq H_1$, $(n-1)H_1\subseteq\auF e_1$,
$ne_1=e_1$) and $L_1=\auF^\times\times\auGL(V')\times\auF^\times$ the subgroup of the $l=(t,g',t')$
with $le_1=te_1$, $l|_{V'}=g'$, $le_1^*=t'e_1^*$; so $P_1=N_1\rtimes L_1$. Let $w_1$ exchange $e_1$
and $e_1^*$ and fix $V'$; it normalises $L_1$, and $P_1\cap w_1P_1w_1^{-1}=L_1$.

\emph{$\Omega=N_1w_1L_1N_1$, and $(n,l,n')\mapsto nw_1ln'$ is a bijection $N_1\times L_1\times N_1\to
\Omega$.} For $h\in\Omega$ we have $he_1\notin H_1$ and $e_1\notin hH_1$, so $V=\auF e_1\oplus(H_1
\cap hH_1)\oplus\auF he_1$ and $H_1=\auF e_1\oplus(H_1\cap hH_1)$. Let $p\in\auGL(V)$ fix $e_1$, map
$H_1\cap hH_1$ onto $V'$ and $he_1$ to $e_1^*$. Then $p\in P_1$, and $phe_1=w_1e_1$, $phH_1=V'\oplus
\auF e_1^*=w_1H_1$, so that $w_1^{-1}ph\in P_1$. Hence $\Omega\subseteq P_1w_1P_1=N_1w_1L_1N_1$; the
other inclusion is clear, and $\#(P_1w_1P_1)=(\#P_1)^2/\#L_1=(\#N_1)^2\,\#L_1$ gives the bijection.

\emph{$U_\pbmD=N_1\rtimes U_{\pbmD'}$}, where $\auGL(V')$ is embedded in $\auGL(V)$ by letting it
fix $e_1$ and $e_1^*$. Indeed $N_1\subseteq U_\pbmD$ by the definitions ($(n-1)$ maps $V$ into
$H_1=B_1$ and $H_1$ into $\auF e_1=A_1$), every $u\in U_\pbmD$ lies in $P_1$, and its component in
$L_1$ is $(1,u',1)$ with $u'=n^{-1}u\in U_\pbmD\cap\auGL(V')=U_{\pbmD'}$. Moreover
$\psi_\pbmD|_{U_{\pbmD'}}=\psi_{\pbmD'}$: for $u'\in\auGL(V')$ the entries $\epsilon_1(u'e_2)$ (or
$\epsilon_1(u'z)$) and $\epsilon_2^*(u'e_1^*)$ (or $\zeta(u'e_1^*)$) vanish.

\emph{$f_\pbmD(nw_1ln')=f_{\pbmD'}(g')$ for $l=(t,g',t')$.} By Lemma~\ref{pbm:pos}(a) we may take
$n=n'=1$; put $h=w_1l$. Let $u=\nu u'$ with $\nu\in N_1$, $u'\in U_{\pbmD'}$. Then $huh^{-1}=(h\nu
h^{-1})(g'u'g'^{-1})$ with $h\nu h^{-1}\in w_1N_1w_1^{-1}$, the unipotent radical opposite to $N_1$,
and $g'u'g'^{-1}\in\auGL(V')\subseteq P_1$. If $huh^{-1}\in U_\pbmD\subseteq P_1$ then $h\nu h^{-1}\in
w_1N_1w_1^{-1}\cap P_1=1$, so $\nu=1$ and $g'u'g'^{-1}\in U_\pbmD\cap\auGL(V')=U_{\pbmD'}$. Hence
$f_\pbmD(h)=\sum_{u'\in U_{\pbmD'}\cap g'^{-1}U_{\pbmD'}g'}\psi_{\pbmD'}(g'u'g'^{-1})\aubar{\psi_{\pbmD'}(u')}=
f_{\pbmD'}(g')$.

Summing over $\Omega$ gives $(\#N_1)^2(q-1)^2\sum_{g'}f_{\pbmD'}(g')=(\#N_1)^2\,\#L_1\cdot\pbmPhi(\pbmD')$.
\end{proof}

\begin{lemma}[Fixing the first vector of the left chain]\label{pbm:first}
Let $\pbmD$ be a chain datum of type $(i,i')$ with $i\ge1$ on $V$, of dimension $r$, and let
$\pbmD_1$ be the chain datum of type $(i-1,i')$ on $V_1=\langle e_2,\dots,e_i\rangle\oplus X'\oplus
\langle e^*_{i'},\dots,e^*_1\rangle$ with the same $z$, $\zeta$, $\psi_1$, $\psi_2$. Let $v_2$ be
the first vector of the left chain of $\pbmD_1$ ($v_2=e_2$ if $i\ge2$, $v_2=z$ if $i=1$). Then
\[
\sum_{h\in\auGL(V),\ he_1=e_1}f_\pbmD(h)=q^{2(r-1)}\sum_{g\in\auGL(V_1),\ gv_2=v_2}f_{\pbmD_1}(g).
\]
\end{lemma}
\begin{proof}
Let $S$ be the stabiliser of $e_1$ and $N=\{1+e_1\otimes\xi:\xi\in V^*,\ \xi(e_1)=0\}\cong V_1^*$, of
order $q^{r-1}$; then $S=N\rtimes\auGL(V_1)$, where $\auGL(V_1)$ fixes $e_1$. As in the proof of
Lemma~\ref{pbm:pair}, $N\subseteq U_\pbmD\subseteq S$ and $U_\pbmD=N\rtimes U_{\pbmD_1}$, so that an
element of $S$ lies in $U_\pbmD$ if and only if its component in $\auGL(V_1)$ lies in
$U_{\pbmD_1}$; and $\psi_\pbmD|_{U_{\pbmD_1}}=\psi_{\pbmD_1}$, $\psi_\pbmD(1+e_1\otimes\xi)=\psi_1(\xi(
v_2))$. Every element of $U_{\pbmD_1}$ fixes $v_2$.

Let $h=ng\in S$ and $u=n_1u_1\in U_\pbmD$, with $n,n_1\in N$, $g\in\auGL(V_1)$, $u_1\in U_{\pbmD_1}$,
$n_1=1+e_1\otimes\xi_1$. Then $huh^{-1}\in U_\pbmD$ if and only if $u_1':=gu_1g^{-1}\in U_{\pbmD_1}$,
and in that case
\[
huh^{-1}=n\cdot gn_1g^{-1}\cdot u_1'n^{-1}u_1'^{-1}\cdot u_1',\qquad\psi_\pbmD(huh^{-1})=\psi_1\bigl(
\xi_1(g^{-1}v_2)\bigr)\,\psi_{\pbmD_1}(u_1'),
\]
because $gn_1g^{-1}=1+e_1\otimes\xi_1g^{-1}$ and because $\psi_\pbmD(u_1'n^{-1}u_1'^{-1})=\psi_\pbmD(n)^{-1}$,
as $u_1'$ fixes $v_2$. Summing over $\xi_1\in V_1^*$, by the orthogonality of the characters of
$V_1^*$ ($\psi_1$ is not trivial),
\[
f_\pbmD(ng)=\sum_{\xi_1}\psi_1\bigl(\xi_1(g^{-1}v_2-v_2)\bigr)\cdot f_{\pbmD_1}(g)=q^{r-1}\,[gv_2=v_2]\,
f_{\pbmD_1}(g).
\]
The sum over $n\in N$ gives another factor $q^{r-1}$.
\end{proof}

\begin{lemma}[A crude bound]\label{pbm:crude}
Let $\pbmD$ be a chain datum of type $(0,i')$ on $W$, with $\dim W=n=m+i'$ and $m>i'$. Then
\[
\sum_{g\in\auGL(W),\ gz=z}f_\pbmD(g)\ \le\ 16\,(i'+1)\,q^{\,i'(i'+1)}\cdot\frac{\#\auGL(W)}{q^n-1}.
\]
\end{lemma}
\begin{proof}
Here $A=0$, so every $u\in U=U_\pbmD$ is the identity on $X'$. If $u$ and $gug^{-1}$ lie in $U$,
then $u-1$ vanishes on $X'+g^{-1}X'$; hence, by Lemma~\ref{pbm:pos}(a),
\[
f_\pbmD(g)\le\#(U\cap g^{-1}Ug)\le q^{\,c(g)\,n},\qquad c(g)=\dim W-\dim(X'+g^{-1}X')=\dim(X'\cap g^{-1}
X')-(m-i').
\]
Let $g$ run over the stabiliser $S_z$ of $z$, of order $\#\auGL(W)/(q^n-1)$. Both $X'$ and
$g^{-1}X'$ contain $z$; let $\bar X=X'/\auF z\subset\bar W=W/\auF z$ and let $\bar g$ be the map
induced by $g$. The homomorphism $S_z\to\auGL(\bar W)$ is onto, and $c(g)\ge c$ means that
$\bar X\cap\bar g^{-1}\bar X$ contains a subspace $D$ of dimension $d=m-1-i'+c$. For a fixed $D$ of
dimension $d$ in $\bar X$, the subspace $\bar gD$ is equidistributed among the subspaces of
dimension $d$ of $\bar W$. With the Gaussian binomial coefficients, which satisfy $q^{d(a-d)}\le
\binom ad_q\le4q^{d(a-d)}$, the proportion of the $g\in S_z$ with $c(g)\ge c$ is therefore at most
\[
\binom{m-1}d_q^{\,2}\Big/\binom{n-1}d_q\ \le\ 16\,q^{\,2d(m-1-d)-d(n-1-d)}=16\,q^{-c\,(m-1-i'+c)} .
\]
As $0\le c(g)\le i'$,
\[
\sum_{g\in S_z}f_\pbmD(g)\le\#S_z\sum_{c=0}^{i'}16\,q^{-c(m-1-i'+c)}\,q^{\,c(m+i')}=16\,\#S_z\sum_{c=0}^{i'}
q^{\,c(2i'+1-c)}\le16\,(i'+1)\,q^{\,i'(i'+1)}\,\#S_z .\qedhere
\]
\end{proof}

\begin{proposition}\label{pbm:Phi}
Let $\pbmD$ be a chain datum of type $(\varrho,\varrho)$ with $m\ge2\varrho^2+4$ and $q\ge8$. Then
\[
1\le\pbmPhi(\pbmD)\le\Bigl(1-\frac4q\Bigr)^{-\varrho}\Bigl(1+\frac{128\,\varrho(\varrho+1)}{q^2}\Bigr).
\]
\end{proposition}
\begin{proof}
The lower bound is the term $u=1$. For $0\le j\le\varrho$ let $\pbmD^{(j)}$ be the chain datum of
type $(j,j)$ obtained from $\pbmD$ by removing the vectors $e_a,e_a^*$ with $a\le\varrho-j$, on the
space $V^{(j)}$ of dimension $r_j=m+2j$, with group $\mathsf G_j=\auGL(V^{(j)})$. For $j=0$ the
group $U$ is trivial and $\pbmPhi(\pbmD^{(0)})=1$. Let $1\le j\le\varrho$, write $f=f_{\pbmD^{(j)}}$ and
let $e_1,e_1^*$ denote the first vectors of the two chains of $\pbmD^{(j)}$. By
Lemma~\ref{pbm:pair}, with its set $\Omega$,
\begin{equation}\label{pbm:eq-rec}
\#\mathsf G_j\cdot\pbmPhi(\pbmD^{(j)})\le\#\Omega\cdot\pbmPhi(\pbmD^{(j-1)})+\sum_{\epsilon_1^*(he_1)=0}f(h)
+\sum_{\epsilon_1^*(h^{-1}e_1)=0}f(h),
\end{equation}
and the last two sums are equal, because $f(h^{-1})=f(h)$. By Lemma~\ref{pbm:pos}, $f\ge0$ is of
positive type. We apply Lemma~\ref{pb:singular} with $\pi$ the identity of $\mathsf G_j$, $l=1$,
$A_0$ the column $e_1$ and $A_1$ the column of the linear form $\epsilon_1^*$, so that
${}^tA_1hA_0=\epsilon_1^*(he_1)$:
\begin{equation}\label{pbm:eq-fourier}
\sum_{\epsilon_1^*(he_1)=0}f(h)\le\frac2q\sum_{h\in\mathsf G_j}f(h)+q^{\,r_j/2}\,Y,\qquad
Y=\max\Bigl(\sum_{he_1=e_1}f(h),\ \sum_{\epsilon_1^*\circ h=\epsilon_1^*}f(h)\Bigr)
\end{equation}
(the stabiliser of a non-zero multiple of $e_1$, or of $\epsilon_1^*$, is that of $e_1$, or of
$\epsilon_1^*$).

\emph{The bound for $Y$.} By Lemma~\ref{pbm:first}, applied $j$ times (to the data of types $(j,j),
(j-1,j),\dots,(1,j)$ on spaces of dimensions $r_j,r_j-1,\dots,r_j-j+1$), and by
Lemma~\ref{pbm:crude} for the datum of type $(0,j)$ on a space of dimension $n=m+j$,
\[
\sum_{he_1=e_1}f(h)=q^{\,2\sum_{t=1}^j(r_j-t)}\sum_{g\in\auGL_n(\auF),\ gz=z}f_{(0,j)}(g)\le q^{\,2\sum_{t=1}^j(
r_j-t)}\cdot16(j+1)q^{\,j(j+1)}\,\frac{\#\auGL_n(\auF)}{q^n-1}.
\]
Since $\#\mathsf G_j=\#\auGL_n(\auF)\prod_{t=0}^{j-1}(q^{r_j-t}-1)q^{r_j-t-1}$, and $\prod_{s\ge2}(1-
q^{-s})^{-1}\le2$, $(q^n-1)^{-1}\le2q^{-n}$,
\[
\frac1{\#\mathsf G_j}\sum_{he_1=e_1}f(h)\le16(j+1)q^{\,j(j+1)}\cdot2q^{-n}\cdot\prod_{t=0}^{j-1}\frac{q^{-1}}{1-
q^{-(r_j-t)}}\le64\,(j+1)\,q^{\,j^2-m-j}.
\]
The second sum in $Y$ is of the same kind for the \emph{dual datum}: let $\auGL(V)$ act on $V^*$ by
$a\mapsto a\circ h^{-1}$, and give $V^*$ the basis $\epsilon_1^*,\dots,\epsilon_j^*$; the dual of
the basis of $X'$; $\epsilon_j,\dots,\epsilon_1$, with the vector $\zeta\in X'^*$ and the linear form
``evaluation at $z$'' on $X'^*$. In this basis $U$ is again the group of a chain datum of type
$(j,j)$, and, since $\epsilon(u^{-1}y)=-\epsilon(uy)$ for the entries just above the diagonal
blocks, $\psi_{\pbmD^{(j)}}$ is the character of this datum for the pair of characters $(\psi_2^{-1},
\psi_1^{-1})$. Its function $f$ is the same function on the same group, and the stabiliser of the
form $\epsilon_1^*$ is the stabiliser of the first vector of its left chain. So the bound holds
for both sums, and, as $r_j/2=m/2+j$ and $m\ge2j^2+4$,
\[
q^{\,r_j/2}\,Y\le64\,(j+1)\,q^{\,j^2-m/2}\cdot\#\mathsf G_j\le64\,(\varrho+1)\,q^{-2}\cdot\#\mathsf G_j .
\]
\emph{Conclusion.} By \eqref{pbm:eq-rec}, \eqref{pbm:eq-fourier} and $\#\Omega\le\#\mathsf G_j$,
\[
\Bigl(1-\frac4q\Bigr)\pbmPhi(\pbmD^{(j)})\le\pbmPhi(\pbmD^{(j-1)})+\frac{128\,(\varrho+1)}{q^2},
\]
and the proposition follows by induction on $j$.
\end{proof}

\begin{lemma}[The character $\psi_U$ at a chosen place]\label{pbm:chain}
Let $w$ be a chosen place, $\auF=\auF_w$, and $V_w$, $\mathsf J$ as in \S\ref{pb:ss-setting}, with
the basis $e_1,\dots,e_\varrho$; the reduction $X'$ of the lattice of $X_0\perp Ez$;
$e^*_\varrho,\dots,e^*_1$. Then $\pbUb_w$ is the group, and the component $\psi_{U,w}$ of $\psi_U$ on
$\pbUb_w$ the character, of the chain datum of type $(\varrho,\varrho)$ on $V_w$ with this basis,
$m=k+1$, the vector $z$, the linear form $\zeta(y)={}^tz\,\mathsf Jy$ on $X'$, and two non-trivial
characters $\psi_1,\psi_2$ of $\auF_w$.
\end{lemma}
\begin{proof}
The group $\pbUb_w$ is the unipotent radical of the parabolic subgroup with blocks $1,\dots,1,k+1,1,
\dots,1$ that stabilises the reduction of the flag $A_1\subset\dots\subset A\subset A^\perp\subset
A_{\varrho-1}^\perp\subset\dots$; in the basis above this is $U_\pbmD$ (\S\ref{pb:ss-setting}). The
vector $z$ lies in the lattice, $z\ne0$ in $X'$, and $\zeta(z)=h(z,z)=-N$ is not zero modulo $w$,
as $w\nmid N$. Let $u\in K_U$ have the components $u_w$ and $u_{\aubar w}=\mathsf J^{-1}\,{}^tu_w^{-1}
\mathsf J$ modulo $w$ and $\aubar w$. By \eqref{pb:eq-psi}, $\psi_U(u)=\prod_v\psi_{E,v}(-\sigma(u)_v/
(M\varpi))$; the factors at the places $v\nmid M$ are $1$, because $\sigma(u)_v$ is integral and
$M\varpi$ generates the different at $v$. At $v\in\{w,\aubar w\}$ the place is unramified over $\Q$,
$M$ is a uniformiser and $\varpi$ a unit, so $x\mapsto\psi_{E,v}(-x/(M\varpi))$ is a non-trivial
character of $\auO_v/v=\auF_w$; call these $\psi_1$ (at $w$) and $\psi_2'$ (at $\aubar w$). The
component of $\sigma(u)$ at $w$ is $\sum_{a<\varrho}\epsilon_a(u_we_{a+1})+\epsilon_\varrho(u_wz)$.
At $\aubar w$, with $\langle y,y'\rangle={}^ty\,\mathsf Jy'$ and $\langle u_{\aubar w}y,y'\rangle=
\langle y,u_w^{-1}y'\rangle$, it is
\[
\sum_{a<\varrho}\langle e_{a+1},u_w^{-1}e_a^*\rangle+\langle z,u_w^{-1}e_\varrho^*\rangle=\sum_{a<\varrho}
\epsilon^*_{a+1}(u_w^{-1}e_a^*)+\zeta(u_w^{-1}e^*_\varrho)=-\Bigl(\sum_{a<\varrho}\epsilon^*_{a+1}(u_we_a^*)+
\zeta(u_we^*_\varrho)\Bigr),
\]
because $u_w^{-1}=1-(u_w-1)+(u_w-1)^2-\dots$ and $(u_w-1)^2$ contributes nothing to these entries.
So $\psi_{U,w}=\psi_\pbmD$ with $\psi_2=\psi_2'^{-1}$.
\end{proof}

\begin{theorem}[Theorem~\ref{pb:est}\,\textup{(E4)}]\label{pbm:E4}
Assume $b\ge3$, $k\ge2\varrho^2+3$, $\auNm w\ge b^{10}$ for the chosen places $w$, and $18\,\varrho f\le b^8$
(these hold for the data of \S\ref{pb:ss-setting}: $\varrho<2\sqrt x$, $k\ge x^5/2$, $f\le x^{12}$, $p_i\ge b^{10}$, and $b\ge1000$ by Lemma~\ref{an:lem-sizes}(S1)). Then
$\sum_{s\in\auG}|\psi^{\pbnat}(s)|^2\le2$.
\end{theorem}
\begin{proof}
$\auG$, $\pbUb$ and $\psi_U$ are products over the $2f$ chosen places $w$ (above $p_0$ and $p_1$)
of $\auGL_r(\auF_w)$, $\pbUb_w$ and $\psi_{U,w}$; hence $\tilde\psi$ and $\psi^{\pbnat}$ are products
of the corresponding functions of the factors, and by Lemmas~\ref{pbm:pos}(c)
and~\ref{pbm:chain}, $\sum_{s\in\auG}|\psi^{\pbnat}(s)|^2=\prod_w\pbmPhi(\pbmD_w)$ for chain data
$\pbmD_w$ of type $(\varrho,\varrho)$ with $m=k+1$ over fields with $q_w\ge b^{10}$ elements. By
Proposition~\ref{pbm:Phi}, and $-\log(1-4/q)\le8/q$, and $128(\varrho+1)\le128(b^8/18+1)\le b^{10}\le q_w$ (as $b\ge3$),
\[
\log\pbmPhi(\pbmD_w)\le\frac{8\varrho}{q_w}+\frac{128\,\varrho(\varrho+1)}{q_w^2}\le\frac{9\varrho}{b^{10}},
\qquad\sum_w\log\pbmPhi(\pbmD_w)\le\frac{18\,\varrho f}{b^{10}}\le b^{-2}\le\log2 .\qedhere
\]
\end{proof}

\begin{remark}\label{pbm:rem-E4}
(a) The proof does not use $\zeta(z)\ne0$. (b) The same method bounds the sum of $f_\pbmD$ over
the stabiliser of any non-zero vector of the frame. The gain $b^{-6k}$ for each fixed vector in
Theorem~\ref{pbm:inputs}(F) is proved differently (Theorem~\ref{pff:fixed}): the stabiliser of a vector
has index $p_0^r-1$, which makes its volume small, while its mass integral is controlled by its main
part. (c) \textup{(E4)} could also be deduced from a bound for $\int|\int_K\tilde\psi(yxy^{-1})\,dy|^2dx$
over the whole group at the level primes, which is part of the local theory behind \textup{(L2)}; the
argument above is independent of it.
\end{remark}

\subsection{The terms over open elements}\label{pbm:ss-cell}

We turn to \textup{(L2)}. Put $n=\varrho+1$. We index the frame by $0\le a\le\varrho$:
$f_a=e_{a+1}$ for $a<\varrho$ and $f_\varrho=z$, and write $f_a^*=e_{a+1}^*$ for $a<\varrho$. Let
$\pbmXp=X_0\perp Ez$, the orthogonal complement of $A\oplus A^*$, where $A^*=\langle e_1^*,\dots,
e_\varrho^*\rangle$, and let $\pbmGf=\auU(\pbmXp)\subset G$ be the group of the elements that fix
the $e_a$ and the $e_a^*$; so $H\subset\pbmGf$ is the stabiliser of $z$. A \emph{head} is a tuple
$p=(p_0,\dots,p_{\varrho-1})$ of non-zero elements of $E$; its \emph{Weyl element} $w_p\in\Gamma$
is defined by
\[
w_pf_a=p_a\,f_a^*,\qquad w_pf_a^*=\aubar{p_a}^{\,-1}f_a\quad(a<\varrho),\qquad w_p=1\text{ on }\pbmXp .
\]
For $\sigma\in\pbmGf(E^+)$ the \emph{last entry} is $t(\sigma)=h(\sigma z,z)$. As $\sigma z=x-
(t/N)z$ with $x\in X_0$ and $h(\sigma z,\sigma z)=-N$,
\begin{equation}\label{pbm:eq-last}
h(x,x)=\frac{t\aubar t-N^2}N,\qquad\text{so}\qquad|\tau(t)|\ge N\ \text{ for all }\tau\in\auPl ,
\end{equation}
with equality for one $\tau$ if and only if $x=0$. For $\gamma=w_p\sigma$ we have $h(\gamma f_l,f_j)
=p_l\delta_{jl}$ for $j,l<\varrho$, $h(\gamma z,z)=t(\sigma)$ and $h(\gamma f_l,z)=h(\gamma z,f_j)=0$,
so
\begin{equation}\label{pbm:eq-Z}
Z(w_p\sigma)=\operatorname{diag}(p_0,\dots,p_{\varrho-1},t(\sigma)),\qquad\delta_l(w_p\sigma)=p_0\cdots
p_{l-1}\ (l\le\varrho),\qquad\pbg(w_p\sigma)=|\delta_\varrho|^2\bigl(t\aubar t-N^2\bigr),
\end{equation}
where $|\delta|^2=\delta\aubar\delta$.

\begin{proposition}[The big cell]\label{pbm:cell}
Let $\gamma$ be a point of $G$ with values in $E^+$ or in a completion $E^+_v$; its matrix has
entries in $F=E$ or $F=E\otimes_{E^+}E^+_v$. Suppose that $\delta_l(\gamma)$ is invertible in $F$ for
$1\le l\le\varrho$. Then there are unique $u,u'\in U$, a
head $p$ and $\sigma\in\pbmGf$ with $\gamma=u\,w_p\,\sigma\,u'$; moreover $p_a=\delta_{a+1}(\gamma)/
\delta_a(\gamma)$ and $t(\sigma)=\delta_n(\gamma)/\delta_\varrho(\gamma)$. In particular $U\cap\gamma
U\gamma^{-1}=1$. The maps $\gamma\mapsto u,u',p,\sigma$ are continuous on the open set where the
$\delta_l$, $l\le\varrho$, are invertible. If $v$ is a finite place, $\gamma$ stabilises the lattice
of $Y_0$ at $v$ and the $\delta_l(\gamma)$, $l\le\varrho$, are units at $v$, then $u$, $u'$,
$\sigma$ stabilise the lattice at $v$.
\end{proposition}
\begin{proof}
The leading minors $\delta_l$ are invariant under $\gamma\mapsto u\gamma u'$: $u'f_l$ is $f_l$ plus
a combination of the $f_{l'}$ with $l'<l$, and the same holds for $u^{-1}f_j$, so $Z(u\gamma u')=
\mathsf L\,Z(\gamma)\,\mathsf R$ with $\mathsf L$ lower and $\mathsf R$ upper unitriangular. Let
$U_Q\subseteq U$ be the subgroup of the elements that act trivially on $A$, on $A^\perp/A$ and on
$Y_0/A^\perp$, and $N_A\subseteq U$ the subgroup of the elements that preserve $A^*$ and $\pbmXp$;
then $U=U_Q\rtimes N_A$, and $N_A$ acts on $A$ by all the unitriangular matrices of the basis
$(e_a)$. The group $U_Q$ acts simply transitively on the totally isotropic subspaces $A'$ of
dimension $\varrho$ with $A'\cap A^\perp=0$: such an $A'$ has a unique basis $(e_a')$ with $h(e_b,
e_a')=\delta_{ab}$, and $e_a'-e_a^*\in A^\perp$; the linear map $u$ with $ue_a=e_a$, $ue_a^*=e_a'$
and $ux=x-\sum_a\aubar{h(e_a',x)}\,e_a$ for $x\in\pbmXp$ is unitary (the correction makes $ux$
orthogonal to the $e_a'$, and $A$ is orthogonal to $A^\perp$), lies in $U_Q$ and maps $A^*$ onto
$A'$; an element of $U_Q$ with this property must map $e_a^*$ to $e_a'$ and $x$ to the element of
$x+A$ orthogonal to $A'$, so it is $u$.

\emph{Existence.} As $\delta_\varrho(\gamma)\ne0$, the pairing of $A$ with $\gamma A$ is perfect, so
$\gamma A\cap A^\perp=0$ and there is $u_1\in U_Q$ with $u_1A^*=\gamma A$. Then $\gamma_1=u_1^{-1}
\gamma$ maps $A$ onto $A^*$; the pairing of $A$ with $\gamma_1^{-1}A$ is the adjoint of that of
$\gamma_1A=A^*$ with $A$, hence perfect, and there is $u_2\in U_Q$ with $u_2A^*=\gamma_1^{-1}A$. Now
$\gamma_2=\gamma_1u_2$ maps $A$ onto $A^*$ and $A^*$ onto $A$, hence preserves $\pbmXp$. The block
$(h(\gamma_2f_l,f_j))_{j,l<\varrho}$ has the leading minors $\delta_l(\gamma)\ne0$, so it is
$\mathsf L\,\mathsf D\,\mathsf R$ with unique unitriangular $\mathsf L$, $\mathsf R$ and diagonal
$\mathsf D=\operatorname{diag}(\delta_{a+1}/\delta_a)$. Multiplying $\gamma_2$ on both sides by the
elements of $N_A$ that induce $\mathsf L^{-1}$ and $\mathsf R^{-1}$ we get $\gamma_3$ with $\gamma_3
f_a=p_af_a^*$; by unitarity $\gamma_3f_a^*=\aubar{p_a}^{-1}f_a$, so $\gamma_3=w_p\sigma$ with
$\sigma=\gamma_3|_{\pbmXp}\in\pbmGf$. The value of $t(\sigma)$ follows from \eqref{pbm:eq-Z}.

\emph{Uniqueness.} Let $u\,w_p\sigma\,u'=w_q\sigma'$. The minors give $p=q$. Applying both sides to
$A$ gives $uA^*=A^*$, so $u\in N_A$, and in the same way $u'\in N_A$; then $\mathsf L\mathsf D\mathsf
R=\mathsf D$ forces $u=u'=1$ and $\sigma=\sigma'$. For $u\in U\cap\gamma U\gamma^{-1}$ apply this to
$u\gamma=\gamma u'$.

\emph{Continuity and integrality.} Every step is the solution of a linear system whose determinant
is a product of the $\delta_l$; this gives the continuity, and, over the ring of integers at $v$
when the $\delta_l$ are units and $\gamma$ preserves the lattice (which is the sum of its
intersections with $A\oplus A^*$ and $\pbmXp$, with the basis $(e_a,e_a^*)$ of the first), the same
construction gives integral $u_1,u_2,\mathsf L,\mathsf R$, hence integral $u$, $u'$, $\sigma$.
\end{proof}

\subsubsection*{Orbit sums}
We use the cut Bessel functions $W^\flat$ of Definition~\ref{pb:bessel}, and assume in addition:
\begin{enumerate}
\item[(B5$'$)] for each $\tau$ there is a compact set $C_\tau\subset G_\tau$ with $\operatorname{supp}
W_\tau^\flat\subseteq U_\tau C_\tau$.
\end{enumerate}
(This holds for the cutoff of Theorem~\ref{pbe:bessel}: there $W^\flat_\tau$ vanishes unless
$\|sy\|_\tau<b^{6n+2}$ for some $s\in R_\tau=U_\tau\rtimes H_\tau$, and $H_\tau$ is compact.) For
$\gamma\in\Gamma$ in the big cell, that is, with $\delta_l(\gamma)\ne0$ for $1\le l\le\varrho$, the
\emph{orbital integral} is
\[
O(\gamma)=\int_{G_\infty}W^\flat(y)\,W^\flat(\gamma^{-1}y)^*\,dy .
\]
The integrand is continuous with compact support: if it does not vanish at $y$, then for each
$\tau$, $y_\tau=u_1c_1$ and $\gamma^{-1}y_\tau=u_2c_2$ with $u_i\in U_\tau$, $c_i\in C_\tau$, so that
$u_1^{-1}\gamma u_2=c_1c_2^{-1}$ lies in the compact set $C_\tau C_\tau^{-1}$; by
Proposition~\ref{pbm:cell} the set $U_\tau\gamma U_\tau$ is closed in $G_\tau$ (it is defined by
the equations $\delta_l=\delta_l(\gamma)$, $l\le\varrho$, and $\sigma(\cdot)=\sigma(\gamma)$, where
$\sigma(\cdot)$ is continuous on the big cell) and $(u_1,u_2)\mapsto u_1^{-1}\gamma u_2$ is a
homeomorphism of $U_\tau^2$ onto it; hence $u_1$ stays in a compact set, and so does $y_\tau$. By
Proposition~\ref{pb:unfold} and $\|W^\flat\|\le\|W\|$ the amplitude sums \eqref{pb:eq-amp} for
$W^\flat$ converge absolutely. For classes
$c,c'$, $s\in\auG$ and $\gamma_0$ in the big cell the set $U(E^+)\gamma_0U(E^+)$ is stable; its amplitude sum is the \emph{orbit sum}
$\pbmOrb_{c,s,c'}(\gamma_0)=A^\flat_{U(E^+)\gamma_0U(E^+)}(c,s,c')$.

\begin{proposition}[Orbit sums are Kloosterman sums]\label{pbm:orbit}
Let $\gamma_0$ be in the big cell and
\[
\pbmKl_{c,s,c'}(\gamma_0)=\sum_{\substack{u\in U(E^+)/U_c,\ u'\in U(E^+)/U_{c'}\\ u^{-1}\gamma_0u'\in\Gamma_{c,s,
c'}}}\psi_f(u)\,\aubar{\psi_f(u')} .
\]
If $O(\gamma_0)\ne0$ the sum is finite and $\pbmOrb_{c,s,c'}(\gamma_0)=\pbmKl_{c,s,c'}(\gamma_0)\cdot
O(\gamma_0)$; if $O(\gamma_0)=0$ then $\pbmOrb_{c,s,c'}(\gamma_0)=0$.
\end{proposition}
\begin{proof}
By Proposition~\ref{pbm:cell}, $(u,u')\mapsto u^{-1}\gamma_0u'$ is a bijection of $U(E^+)^2$ onto
the orbit, the stabilisers $U_c\cap\gamma U_{c'}\gamma^{-1}$ are trivial, and the double cosets
$U_c\gamma U_{c'}$ in the orbit correspond to the pairs $(uU_c,u'U_{c'})$. For $\gamma=u^{-1}\gamma_0
u'$ the substitution $y\mapsto u^{-1}y$ and (B1) give
\[
O(\gamma)=\int_{G_\infty}W^\flat(u^{-1}y)\,W^\flat(u'^{-1}\gamma_0^{-1}y)^*\,dy=\aubar{\psi_\infty(u)}\,\psi_\infty(
u')\,O(\gamma_0)=\psi_f(u)\,\aubar{\psi_f(u')}\,O(\gamma_0),
\]
since $\psi_\infty\psi_f=1$ on $U(E^+)$. All terms of the amplitude sum have the norm $\|O(\gamma_0)
\|$, and the sum of the norms converges; so the number of terms is finite if $O(\gamma_0)\ne0$.
\end{proof}

By strong approximation for $U$, $U(E^+)/U_c=U(\auA_f)/(U(\auA_f)\cap h_cJh_c^{-1})$; so
$\pbmKl_{c,s,c'}(\gamma_0)$ is a sum over the pairs $(u,u')$ of finite adelic points with $h_c^{-1}
u^{-1}\gamma_0u'h_{c'}\in K$ of reduction $s$: a Kloosterman sum for the group $U$, the character
$\psi_f$ and the level $J$.

\subsubsection*{Heads, links and class sums}
By (B5$'$) and the invariance of the $\delta_l$ there is a constant $C_0$ such that
\begin{equation}\label{pbm:eq-bdd}
O(\gamma)\ne0\ \Longrightarrow\ |\tau(\delta_l(\gamma))|\le C_0\quad(\tau\in\auPl,\ 1\le l\le n):
\end{equation}
if $O(\gamma)\ne0$ there are $y_\tau=u_1c_1$ and $\gamma^{-1}y_\tau=u_2c_2$ with $u_i\in U_\tau$,
$c_i\in C_\tau$, so $\tau(\delta_l(\gamma))=\delta_l(c_1c_2^{-1})$, which is bounded on the compact
set $C_\tau C_\tau^{-1}$. Let $\mathsf B$ be the finite set of the non-zero integers of $E$ all of
whose conjugates have absolute value at most $C_0$, and
\[
\pbmT=\{\delta/\delta':\delta,\delta'\in\mathsf B\}\subset E^\times ,
\]
the \emph{set of entries}; it is finite and depends only on the data and on the $C_\tau$.

\begin{proposition}[Class sums over open elements]\label{pbm:classsum}
Let $s\in\auG$ be open, $c,c'$ classes and $\Lambda$ a link.
\begin{enumerate}
\item[(a)] Every $\gamma\in\Gamma_{c,s,c'}$ lies in the big cell, and $\delta_l(\gamma)\in\auO_E$ for
$l\le n$.
\item[(b)] For $\gamma=u\,w_p\,\sigma\,u'$ as in Proposition~\ref{pbm:cell} the link of $\gamma$ is
the graph $\{(x,\sigma^{-1}x):x\in\pbmXp\}$ of $\sigma^{-1}$, and $\pbg(\gamma)\ne0$ if and only if
$|\tau(t(\sigma))|>N$ for all $\tau$.
\item[(c)] If $A^\flat_{T_\Lambda}(c,s,c')\ne0$, there is a unique $\sigma=\sigma_\Lambda\in\pbmGf(E^+)$
whose inverse has the graph $\Lambda$; its last entry $t=t(\sigma)$ lies in $\pbmT$ and
satisfies $|\tau(t)|>N$ for all $\tau$; and
\[
A^\flat_{T_\Lambda}(c,s,c')=\sum_{p\in\pbmT^\varrho}\pbmOrb_{c,s,c'}(w_p\sigma).
\]
Moreover, if $v$ is a finite place at which all elements of $\mathsf B$ are units, then
$h_c^{-1}\sigma h_{c'}$ stabilises the lattice at $v$.
\end{enumerate}
\end{proposition}
\begin{proof}
(a) Let $\kappa=h_c^{-1}\gamma h_{c'}\in K$. As $h_c,h_{c'}\in H(\auA_f)$ are unitary and fix the
frame, $h(\gamma f_l,f_j)=h(\kappa f_l,f_j)$ is integral at all finite places (the lattice is
unimodular and contains the frame); so $Z(\gamma)$ has entries in $\auO_E$. Its reduction modulo a
chosen place $w$ is the matrix $Z_w(s_w)$ of \S\ref{pb:ss-setting} (up to transposition, which does
not change leading minors). As $s$ is open, $\delta_l(s_w)\ne0$ for $1\le l\le\varrho$ at some $w$,
hence $\delta_l(\gamma)\ne0$. (b) The link is invariant under $U(E^+)$ on both sides, so we may take
$\gamma=w_p\sigma$. Then $\gamma A=A^*$, the vectors orthogonal to $A$ and $\gamma A$ are those of
$\pbmXp$, and $\gamma^{-1}x=\sigma^{-1}x$ for $x\in\pbmXp$. The second claim is \eqref{pbm:eq-Z},
\eqref{pbm:eq-last} and the invariance of $\pbg$: $t\aubar t-N^2\in E^+$ is totally non-negative.
(c) By (a), (b) and Proposition~\ref{pbm:cell}, $T_\Lambda\cap\Gamma_{c,s,c'}$ is the union, over
the heads $p$, of the sets $U(E^+)w_p\sigma U(E^+)\cap\Gamma_{c,s,c'}$ for the unique $\sigma$ with
graph of $\sigma^{-1}$ equal to $\Lambda$ (if there is none, or if $|\tau(t(\sigma))|=N$, the set is
empty). So the amplitude sum is the sum of the orbit sums, each orbit being stable. If
$\pbmOrb_{c,s,c'}(w_p\sigma)\ne0$, the orbit contains a $\gamma\in\Gamma_{c,s,c'}$ with $O(\gamma)\ne
0$; by (a) and \eqref{pbm:eq-bdd} its minors $\delta_l(\gamma)$, $1\le l\le n$, lie in $\mathsf B$
(note $\delta_n=\delta_\varrho\,t\ne0$), so $p_a=\delta_{a+1}/\delta_a\in\pbmT$ and $t=\delta_n/
\delta_\varrho\in\pbmT$. Finally let $v$ be as in the last claim and $\gamma$ as before. Then
$\kappa_v=(h_c^{-1}\gamma h_{c'})_v$ stabilises the lattice and has unit minors, so $\kappa_v=u_v
w_p\sigma_vu_v'$ with $\sigma_v$ integral, by Proposition~\ref{pbm:cell}. As $h_c$, $h_{c'}$ commute
with $w_p$ and normalise $U$, $\gamma=(h_cu_vh_c^{-1})\,w_p\,(h_c\sigma_vh_{c'}^{-1})\,(h_{c'}u_v'h_{c'}^{
-1})$, and the uniqueness in Proposition~\ref{pbm:cell} gives $h_c^{-1}\sigma h_{c'}=\sigma_v$.
\end{proof}

\subsection{Sieve systems and one-sided Poisson summation}\label{pbm:ss-sieve}

\subsubsection*{Sieve systems}
A \emph{place} is a quadruple $\mathfrak P=(C,dc,\chi,\iota)$: a locally compact ring $C$ with a
Haar measure $dc$ such that $C^\times$ is open with complement of measure $0$, a continuous unitary
character $\chi$ of $(C,+)$, and a map $\iota\colon C^\times\to C$. A \emph{head of length $i$} at
$\mathfrak P$ is an element of $(C^\times)^i$; for heads $p,q$ of length $i$ and $c\in C^\times$,
$(p,c)$ is the head of length $i+1$.

\begin{definition}[Sieve system]\label{pbm:system}
Let $\mathfrak P$ be a place, $\pbmE\ge0$, $\beta\colon C^\times\to[0,\infty)$ and $\Lambda_i\colon
(C^\times)^i\to[0,\infty)$ for $0\le i\le\varrho$. A \emph{sieve system} at $\mathfrak P$ with mass
$\pbmE$, weights $\Lambda$ and decay $\beta$ is a family of functions $J_i\colon(C^\times)^i\times
(C^\times)^i\to\C$, $0\le i\le n=\varrho+1$, such that:
\begin{enumerate}
\item[(S1)] for $0\le i\le\varrho$, heads $p,q$ of length $i$ and $\xi\in C$, the function $c\mapsto
\chi(\xi c)J_{i+1}((p,c),(q,c))$ is integrable on $C$; $\int_CJ_{i+1}((p,c),(q,c))\,dc=J_i(p,q)$; and
$\bigl|\int_C\chi(\xi c)J_{i+1}((p,c),(q,c))\,dc\bigr|\le\beta(\xi)\Lambda_i(p)\Lambda_i(q)$ for
$\xi\in C^\times$;
\item[(S2)] for heads $p^{(1)},\dots,p^{(\kappa)}$ of length $\varrho$ and $t\in C^\times$ the matrix
$\bigl(J_n((p^{(a)},t),(p^{(b)},t))\bigr)_{a,b}$ is hermitian and positive semidefinite;
\item[(S3)] $J_n$ is continuous;
\item[(S4)] $|J_0|\le\pbmE$;
\item[(S5)] for $1\le i\le\varrho$ and heads $p,q$ of length $i$ such that $\xi=\iota(p_{i-1})-\iota(
q_{i-1})$ is a unit, $|J_i(p,q)|\le\beta(\xi)\Lambda_i(p)\Lambda_i(q)$.
\end{enumerate}
\end{definition}

\begin{lemma}\label{pbm:sysprop}
\begin{enumerate}
\item[(a)] \emph{(Products.)} Let $(J^v_i)$ be sieve systems at places $\mathfrak P_v=(C_v,dc_v,\chi_v,
\iota_v)$, $v$ in a finite set, with masses $\pbmE_v$, weights $\Lambda^v$ and decays $\beta_v$. Then
$J_i(p,q)=\prod_vJ^v_i(p_v,q_v)$ is a sieve system at the product place $(\prod C_v,\prod dc_v,\prod
\chi_v,\prod\iota_v)$ with mass $\prod\pbmE_v$, weights $\prod_v\Lambda^v_i(p_v)$ and decay $\prod_v
\beta_v(\xi_v)$.
\item[(b)] \emph{(Positivity and semicontinuity.)} Let $(J_i)$ be a sieve system, $0\le i\le\varrho$
and $\mathcal P$ a finite set of heads of length $i$. Then $D(x)=\sum_{p,q\in\mathcal P}J_{i+1}((p,x),
(q,x))$ is real, non-negative and lower semicontinuous on $C^\times$.
\end{enumerate}
\end{lemma}
\begin{proof}
(a) The units of the product are the families of units. (S1) follows from Fubini's theorem: the
function is a product of integrable functions of the separate variables. (S2) is Schur's theorem
that the entrywise product of positive semidefinite hermitian matrices is positive semidefinite.
(S3)--(S5) are clear. (b) For $0\le j\le\varrho-i$ and $c\in(C^\times)^j$ let $D_j(x,c)=\sum_{p,q\in
\mathcal P}J_{i+1+j}((p,x,c),(q,x,c))$. For $j=\varrho-i$ this is the sum of all entries of a matrix
as in (S2), so it is $\ge0$, and it is continuous by (S3). By (S1), $D_j(x,c)=\int_CD_{j+1}(x,c,c')
\,dc'$; so, by descending induction on $j$ and Fatou's lemma, $D_j\ge0$ is lower semicontinuous in
$(x,c)$. The case $j=0$ is the claim.
\end{proof}

\subsubsection*{One-sided Poisson summation}

\begin{lemma}\label{pbm:poisson}
Let $\mathbb A$ be a locally compact abelian group with a Haar measure, $L\subset\mathbb A$ a
discrete subgroup with compact quotient of volume $\mathsf c_L$, and $L^\perp$ the group of the
characters of $\mathbb A$ trivial on $L$. Let $F\ge0$ be integrable on $\mathbb A$ and lower
semicontinuous at every point of a finite set $V\subseteq L$. Then
\[
\sum_{t\in V}F(t)\le\frac1{\mathsf c_L}\sum_{\eta\in L^\perp}\Bigl|\int_{\mathbb A}F(x)\,\eta(x)\,dx\Bigr| .
\]
\end{lemma}
\begin{proof}
Let $\mathsf S$ be the right side; we may assume $\mathsf S<\infty$. The function $PF(x)=\sum_{t\in L}
F(x+t)$ is integrable on the compact group $\mathbb A/L$, whose dual is $L^\perp$, and its Fourier
coefficients, for the measure of mass $1$, are $\mathsf c_L^{-1}\int_{\mathbb A}F\aubar\eta$. They are
absolutely summable, so $PF$ agrees almost everywhere with a continuous function of absolute value
at most $\mathsf S$ (two integrable functions with the same Fourier coefficients agree almost
everywhere). Let $\epsilon>0$. By lower semicontinuity there is a neighbourhood $\mathsf N$ of $0$
with $F(t+x)>F(t)-\epsilon$ for $t\in V$, $x\in\mathsf N$. As $\mathsf N$ has positive measure,
there is $x\in\mathsf N$ with $\mathsf S\ge PF(x)\ge\sum_{t\in V}F(t+x)\ge\sum_{t\in V}F(t)-\epsilon\,
\#V$.
\end{proof}

\subsubsection*{The global place}
Let $T$ be a finite set of finite places of $E$, stable under conjugation, containing the places
above $M$. Let $E_T=\prod_{w\in T}E_w$, with integers $\auO_T$ and the Haar measure for which
$\auO_T$ has volume $1$; $|\cdot|_T=\prod_{w\in T}|\cdot|_w$ is the module, and for $\xi\in E_T^\times$
\[
\pbmn(\xi)=|\xi|_T^{-1},\qquad\pbmden(\xi)=[\auO_T:\auO_T\cap\xi^{-1}\auO_T]=\prod_{w\in T}\max(1,|\xi|_w).
\]
Let $\pbmRT\subset E$ be the ring of the elements integral outside $T$, and $\psi_T(x)=\psi_{E,f}(-x/
\varpi)$ for $x\in E_T$, where $\psi_{E,f}$ is the finite part of $\psi_E$. Put
\[
\pbmAT=\C^{\auPl}\times E_T,\qquad\chi_T(x)=\prod_{\tau\in\auPl}e^{x_\tau-\aubar{x_\tau}}\cdot\psi_T(x_T),
\]
with Lebesgue measure on each $\C$. Recall $\kappa_\tau=M\tau(\varpi)/(2\pi i)\in\R$ and put
\begin{equation}\label{pbm:eq-scales}
s_{\tau,a}=\frac{\kappa_\tau^{\,2(\varrho-a)}}N\quad(0\le a\le\varrho),\qquad[t]_a=\bigl((s_{\tau,a}\tau(
t))_\tau,\ t\bigr)\in\pbmAT\quad(t\in E),
\end{equation}
the element $t$ \emph{read as an entry at position $a$} ($s_{\tau,a}$ is the square of the scale of
$f_a$ in the normalised frame of \S\ref{pbe:ss-real}; thus this $s_{\tau,a}$ is the square of the $s_{\tau,a}$ of \eqref{pbe:eq-scale} and of Lemma~\ref{pb:covol}, which is why $\Lambda_*=\Lambda$ below). For a tuple $p=(p_0,\dots,p_{i-1})$ of
non-zero elements of $E$ the head $[p]=([p_0]_0,\dots,[p_{i-1}]_{i-1})$ has components $[p]_\tau\in
(\C^\times)^i$ and $[p]_T=p\in(E_T^\times)^i$. Let
\[
\pbmvol=(2/\aurd)^f,\qquad c_a=\frac{\pbmvol}{\prod_\tau s_{\tau,a}^{\,2}}\quad(0\le a\le\varrho).
\]

\begin{lemma}[Poisson summation in one entry]\label{pbm:poissonT}
Let $F\ge0$ be integrable on $\pbmAT$ and lower semicontinuous at the points $[t]_a$, $t\in V$, for
a finite set $V\subset\pbmRT$. Then
\[
\sum_{t\in V}F([t]_a)\le c_a\sum_{y\in\pbmRT}\Bigl|\int_{\pbmAT}F(x)\,\chi_T(\xi_a(y)x)\,dx\Bigr|,\qquad
\xi_a(y)=\Bigl(\Bigl(\frac{2\pi i}{\tau(\varpi)\,s_{\tau,a}}\,\tau(y)\Bigr)_\tau,\ y\Bigr).
\]
\end{lemma}
\begin{proof}
Embed $\pbmRT$ in $\pbmAT$ by $t\mapsto((\tau(t))_\tau,t)$. It is discrete with compact quotient:
$\pbmAT\times\prod_{v\notin T}\auO_v$ is an open subgroup of $\auA_E$ that meets $E$ in $\pbmRT$ and,
by strong approximation for the additive group, generates $\auA_E$ together with $E$. So $\pbmAT/
\pbmRT$ is the quotient of $\auA_E/E$ by the image of $\prod_{v\notin T}\auO_v$, and its characters
are the characters $x\mapsto\psi_E(y'x)$, $y'\in E$, of $\auA_E/E$ that are trivial on $\prod_{v
\notin T}\auO_v$, that is, those with $y'\in\varpi^{-1}\pbmRT$ (the different is generated by
$\varpi$). As the components of $\psi_E$ at the infinite places are $x\mapsto e^{-2\pi i(x+\aubar x)}$
(convention (Cv2) of \S\ref{pb:ss-setting}), these are, with $y=-\varpi y'\in\pbmRT$, the characters $x\mapsto\prod_\tau e^{\omega_\tau-\aubar{
\omega_\tau}}\cdot\psi_T(yx_T)$ with $\omega_\tau=2\pi i\,\tau(y/\varpi)\,x_\tau$, up to the sign of
$y$. The covolume of $\auO_E$ in $\C^{\auPl}$ is
$2^{-f}|\operatorname{disc}E|^{1/2}=\pbmvol^{-1}$, and $\pbmRT+\auO_T=E_T$, $\pbmRT\cap\auO_T=\auO_E$; so
$\mathsf c_L=\pbmvol^{-1}$. Apply Lemma~\ref{pbm:poisson} to $F_a(x)=F((s_{\tau,a}x_\tau)_\tau,x_T)$;
the substitution $x_\tau\mapsto x_\tau/s_{\tau,a}$ has the Jacobian $\prod_\tau s_{\tau,a}^{-2}$ and
turns the character of $y$ into $\chi_T(\xi_a(y)\,\cdot\,)$.
\end{proof}

\subsubsection*{The sieve}
We consider sieve systems at the place $(\pbmAT,dx,\chi_T,\iota)$ with
\begin{equation}\label{pbm:eq-iota}
\iota\bigl((x_\tau)_\tau,x_T\bigr)=\bigl((x_\tau^{-1})_\tau,\ M^{-1}x_T^{-1}\bigr).
\end{equation}

\begin{theorem}[The sieve, abstract form]\label{pbm:abstract}
Let $(J_i)$ be a sieve system at $(\pbmAT,dx,\chi_T,\iota)$ with mass $\pbmE$, weights $\Lambda$ and
decay $\beta$. Let $\pbmT\subset E^\times$ be a finite set with $t,t^{-1}\in\pbmRT$ for $t\in\pbmT$,
and $\epsilon,\Gamma\ge0$ such that:
\begin{enumerate}
\item[(D)] for every $a\le\varrho$, every family $(c_\tau)$ of complex numbers with $|c_\tau|^{-1}\le
|\kappa_\tau|/M$ and every finite set $X\subset\pbmRT\smallsetminus\{0\}$,
$\sum_{y\in X}\beta\bigl((c_\tau\tau(y)/s_{\tau,a})_\tau,\,y\bigr)\le\epsilon$;
\item[(W)] $\bigl(\sum_{p\in\pbmT^i}\Lambda_i([p])\bigr)^2\le\Gamma\prod_{a<i}c_a$ for $0\le i\le
\varrho$.
\end{enumerate}
Assume $|\kappa_\tau|\ge1$ for all $\tau$. Then for every subset $\mathsf U\subseteq\pbmT$
\[
\sum_{t\in\mathsf U}\ \sum_{p,q\in\pbmT^\varrho}\Rel J_n\bigl([p,t],[q,t]\bigr)\le\prod_{a=0}^{\varrho}c_a
\cdot\bigl(\pbmE+2n\,\Gamma\,\epsilon\bigr).
\]
\end{theorem}
\begin{proof}
For $0\le i\le\varrho$ and $x\in\pbmAT^\times$ let $D_i(x)=\sum_{p,q\in\pbmT^i}J_{i+1}(([p],x),([q],x))$,
and put $\Sigma_0=\Rel J_0$, $\Sigma_{i+1}=\sum_{t\in\pbmT}D_i([t]_i)$ and $W_i=\sum_{p\in\pbmT^i}\Lambda_i
([p])$. By Lemma~\ref{pbm:sysprop}(b), $D_i\ge0$ is lower semicontinuous at the points $[t]_i$, so
the left side of the theorem is at most $\Sigma_n$. We show
\begin{equation}\label{pbm:eq-step}
\Sigma_{i+1}\le c_i\bigl(\Sigma_i+2\epsilon W_i^2\bigr)\qquad(0\le i\le\varrho);
\end{equation}
with (W) and (S4) this gives $\Sigma_i\le\prod_{a<i}c_a\cdot(\pbmE+2i\Gamma\epsilon)$ by induction.

\emph{Poisson summation.} By (S1), $D_i$ is integrable, $\int D_i=\sum_{p,q\in\pbmT^i}J_i([p],[q])$,
and $|\int D_i(x)\chi_T(\xi x)dx|\le\beta(\xi)W_i^2$ for units $\xi$. For $y\in\pbmRT$, $y\ne0$, the
frequency $\xi_i(y)$ of Lemma~\ref{pbm:poissonT} is a unit and has the form in (D) with $c_\tau=2\pi
i/\tau(\varpi)$, for which $|c_\tau|^{-1}=|\kappa_\tau|/M$. By Lemma~\ref{pbm:poissonT} (the sums
over finite sets of frequencies being bounded by (D)),
\[
\Sigma_{i+1}\le c_i\Bigl(\sum_{p,q\in\pbmT^i}\Rel J_i([p],[q])+\epsilon W_i^2\Bigr).
\]
\emph{Unequal last entries.} For $i=0$ the sum is $\Sigma_0$. Let $i\ge1$. The pairs $(p,q)$ with
$p_{i-1}=q_{i-1}$ contribute $\Sigma_i$, by definition. For a pair with $p_{i-1}=t\ne\tilde t=q_{i-1}$ put
$y=M^{-1}(t^{-1}-\tilde t^{-1})\in\pbmRT\smallsetminus\{0\}$; then $\iota([t]_{i-1})-\iota([\tilde t]_{i-1})=((M\tau(
y)/s_{\tau,i-1})_\tau,y)$ is a unit of the form in (D) with $c_\tau=M$, and $M^{-1}\le|\kappa_\tau|/
M$. By (S5) and (D) with $X=\{y\}$, $\Rel J_i([p],[q])\le\epsilon\Lambda_i([p])\Lambda_i([q])$. The
sum over these pairs is at most $\epsilon W_i^2$, which proves \eqref{pbm:eq-step}.
\end{proof}

\subsubsection*{The explicit weights and decay}
Let $\ell_K(r)=r^K\min(1,r)^6$, let $\pbmTh>0$ be a constant and
\begin{equation}\label{pbm:eq-beta}
\begin{gathered}
\beta^\infty(\xi)=\pbmTh\,b^{2000n^9}\,\ell_{k-n^7}(|\xi|^{-1})\quad(\xi\in\C^\times),\qquad
\beta^T(\xi)=\frac{\pbmn(\xi)^{(k-n^7)/2}}{\pbmden(\xi)^3}\,b^{\,n^9f}\quad(\xi\in E_T^\times),\\
\Lambda^\infty_i(p)=\prod_{l=1}^i\bigl[\,|\delta_l|\le b^{60n^2}\bigr]\max(1,|\delta_l|^{-1})^{n^6},\qquad
\Lambda^T_i(p)=\prod_{l=1}^i\bigl[\delta_l\in\auO_T\bigr]\,\pbmn(\delta_l)^{n^6/2},
\end{gathered}
\end{equation}
where $\delta_l=p_0\cdots p_{l-1}$ and $[\,\cdot\,]$ is $1$ or $0$ according as the condition holds
or not. Let
\[
\pbmdec=\prod_{\tau\in\auPl}\pbmTh\Bigl(\frac{|\kappa_\tau|^{2\varrho+1}}{MN}\Bigr)^{k-n^7}.
\]

\begin{lemma}[Sums over frequencies]\label{pbm:freq}
For $y\in E^\times$ let $\mathfrak b_y=\{a\in\auO_E:ay\in\auO_E\}$. For every finite $X\subset E^\times$
\[
\sum_{y\in X}\ \prod_{\tau\in\auPl}\max(1,|\tau(y)|)^{-6}\cdot\auNm(\mathfrak b_y)^{-3}\le162^f .
\]
\end{lemma}
\begin{proof}
The $y$ with $\mathfrak b_y=\mathfrak b$ lie in $\mathfrak b^{-1}$, which is the union of $\auNm
\mathfrak b$ classes modulo $\auO_E$. Cut each $\C$ into the half-open squares $Q_{\mathsf m}$ of
side $\frac12$ with corners in $\frac12\Z[i]$. Two elements of one class that lie in the same
product of squares differ by an integer all of whose conjugates have absolute value $<1$, so they
are equal. Hence the sum over one class is at most $C_1^f$ with $C_1=\sum_{\mathsf m}\sup_{Q_{\mathsf
m}}\max(1,|z|)^{-6}\le1+\sum_{j\ge1}8j\max(1,\frac{j-1}2)^{-6}<54$. So the left side is at most
$C_1^f\sum_{\mathfrak b}\auNm(\mathfrak b)^{-2}\le(54\,\zeta(2)^2)^f$.
\end{proof}

\begin{theorem}[The sieve]\label{pbm:sieve}
Assume $n\ge3$, $b\ge1000$, $2n\le b$, $1\le|\kappa_\tau|\le b^{23}$ for all $\tau$, $N\le b^{60n}$
and $\aurd\le b$. Let $T$ and $\pbmT$ be as in Theorem~\ref{pbm:abstract}, with $M^{-1}\in\pbmRT$. For
each $\tau\in\auPl$ let $(J^\tau_i)$ be a sieve system at $(\C,\text{Lebesgue},e^{c-\aubar c},w\mapsto
w^{-1})$ with mass $\pbmE_\tau$, weights $\Lambda^\infty$ and decay $\beta^\infty$, and let $(J^T_i)$
be a sieve system at $(E_T,dx,\psi_T,w\mapsto M^{-1}w^{-1})$ with mass $\pbmE_T$, weights $\Lambda^T$
and decay $\beta^T$. Then for every $\mathsf U\subseteq\pbmT$
\[
\sum_{t\in\mathsf U}\ \sum_{p,q\in\pbmT^\varrho}\Rel\Bigl(\prod_\tau J^\tau_n\bigl([p,t]_\tau,[q,t]_\tau
\bigr)\cdot J^T_n\bigl((p,t),(q,t)\bigr)\Bigr)\le\prod_{a=0}^{\varrho}c_a\cdot\Bigl(\prod_\tau\pbmE_\tau
\cdot\pbmE_T+b^{2403\,n^9f}\,\pbmdec\Bigr).
\]
\end{theorem}
\begin{proof}
By Lemma~\ref{pbm:sysprop}(a) the products $J_i=\prod_\tau J^\tau_i\cdot J^T_i$ form a sieve system
at $(\pbmAT,dx,\chi_T,\iota)$ with mass $\prod\pbmE_\tau\cdot\pbmE_T$, weights $\Lambda_i(p)=\prod_\tau
\Lambda^\infty_i(p_\tau)\Lambda^T_i(p_T)$ and decay $\beta(\xi)=\prod_\tau\beta^\infty(\xi_\tau)\beta^T(
\xi_T)$. We verify (D) with $\epsilon=\pbmdec\,b^{2002n^9f}$ and (W) with $\Gamma=b^{400n^9f}$;
Theorem~\ref{pbm:abstract} and $2n\le b^{n^9f}$ then give the claim. We use the product formula in
the form
\begin{equation}\label{pbm:eq-pf}
\pbmn(y)^{1/2}\le\prod_\tau|\tau(y)|\qquad(y\in\pbmRT,\ y\ne0),
\end{equation}
which holds because $|y|_w\le1$ for $w\notin T$.

\emph{(D).} Let $K=k-n^7$, $x_\tau=|\kappa_\tau|^{2\varrho+1}/(MN)$ and $B=b^{23(2\varrho+1)}\ge\max(1,
x_\tau)$. Let $a$, $(c_\tau)$ be as in (D) and $y\in\pbmRT$, $y\ne0$. Put $r_\tau=s_{\tau,a}/(|c_\tau|
|\tau(y)|)$; then $r_\tau|\tau(y)|\le x_\tau$, as $s_{\tau,a}\le|\kappa_\tau|^{2\varrho}/N$, and
$\min(1,r_\tau)\le B\max(1,|\tau(y)|)^{-1}$. With \eqref{pbm:eq-pf} and $\pbmden(y)=\auNm(\mathfrak
b_y)$,
\[
\prod_\tau\ell_K(r_\tau)\cdot\frac{\pbmn(y)^{K/2}}{\pbmden(y)^3}\le\prod_\tau\bigl(r_\tau|\tau(y)|\bigr)^K
\min(1,r_\tau)^6\cdot\auNm(\mathfrak b_y)^{-3}\le\prod_\tau x_\tau^KB^6\max(1,|\tau(y)|)^{-6}\cdot\auNm(
\mathfrak b_y)^{-3}.
\]
By Lemma~\ref{pbm:freq} the sum over $y\in X$ of $\beta$ is at most $\pbmdec\cdot(b^{2000n^9}\cdot
b^{n^9}\cdot162B^6)^f\le\pbmdec\,b^{2002n^9f}$, as $162\,b^{138(2\varrho+1)}\le b^{n^9}$ for $n\ge3$.

\emph{(W).} Let $i\le\varrho$, $p\in\pbmT^i$ with $\Lambda_i([p])\ne0$, and $\delta_l=p_0\cdots p_{l-1}
\in\pbmRT$ for $1\le l\le i$. Then $\delta_l\in\auO_T\cap\pbmRT=\auO_E$, and $|\sigma_{\tau,l}\tau(\delta_l
)|\le b^{60n^2}$ with $\sigma_{\tau,l}=\prod_{a<l}s_{\tau,a}\ge N^{-l}\ge b^{-60n^2}$; so $|\tau(\delta_l)|
\le R:=b^{120n^2}$. By \eqref{pbm:eq-pf} the factor of $\Lambda_i([p])$ belonging to $l$ is at most
\[
\Bigl(\prod_\tau\max\bigl(1,|\sigma_{\tau,l}\tau(\delta_l)|^{-1}\bigr)|\tau(\delta_l)|\Bigr)^{n^6}=\Bigl(
\prod_\tau\max\bigl(|\tau(\delta_l)|,\sigma_{\tau,l}^{-1}\bigr)\Bigr)^{n^6}\le R^{\,fn^6},
\]
so $\Lambda_i([p])\le R^{fn^7}=b^{120n^9f}$. The map $p\mapsto(\delta_1,\dots,\delta_i)$ is injective,
and by the argument with squares of Lemma~\ref{pbm:freq} there are at most $(4R+3)^{2f}$ integers of
$E$ with all conjugates of absolute value at most $R$; so at most $(4R+3)^{2fn}\le b^{242n^3f}$
heads have a non-zero weight. Hence $W_i\le b^{130n^9f}$. On the other hand $s_{\tau,a}\le b^{46
\varrho}$ and $\aurd\le b$ give $c_a\ge b^{-93nf}$ and $\prod_{a<i}c_a\ge b^{-93n^2f}$. So $W_i^2\le
b^{400n^9f}\prod_{a<i}c_a$.
\end{proof}

\subsection{The local systems and Siegel's step: statement}\label{pbm:ss-inputs}

\subsubsection*{The constants}
Let $d=\dim W_{\lambda^X}$, $\pbmcG=\pbvol\auU(k)^2/(\pbvol\auU(r)\pbvol\auU(k-1))$ with $\pbvol
\auU(j)=\prod_{l=1}^j2\pi^l/(l-1)!$ (the product of the areas of the unit spheres of $\C^l$), and
$\pbmTh=\pi^{(2\varrho+1)k}\pbvol\auU(k)/\pbvol\auU(r)$ in \eqref{pbm:eq-beta}. Put
$\Lambda_*=\prod_\tau\prod_{a=0}^{\varrho}s_{\tau,a}^{\,k+2(\varrho-a)}$ and
\begin{equation}\label{pbm:eq-consts}
\begin{gathered}
\pbmaS=\#\pbR\cdot\#\pbUb\cdot\#\mathcal C\cdot\Bigl(\frac{\pi^k}{(k-1)!\,N^k}\Bigr)^f\aurd^{\,f/2}\,\pbmvol^{
\,k},\qquad\pbmrS=\frac1d\prod_\tau\Bigl(\sqrt{\pbmcG}\prod_{a=0}^{\varrho}s_{\tau,a}^{\,k-1+2(\varrho-a)}
\Bigr),\\
\pbmnuS=\frac{\#\pbUb}{N^{kf}}\cdot\frac{\aurd^{\,f\varrho/2}}{\pbmvol^{\,\varrho(\varrho+k)}\,\Lambda_*} .
\end{gathered}
\end{equation}
The statement \textup{(L2)} is homogeneous in the Haar measure of $U_\infty$: if this measure is
multiplied by $\lambda$, then, to keep (B2), $W$ and $W^\flat$ are multiplied by $\lambda^{1/2}$,
the orbital integrals over $G_\infty$ by $\lambda$, and $\nu$ by $\lambda$; so both sides of
\textup{(L2)} are multiplied by $\lambda^2$. (They do not depend on the Haar measure of $G_\infty$.)
The local inputs (R) and (S) below are, however, proved for the measure $du$ of convention (Cv3) of
\S\ref{pb:ss-setting}, which we use; for it $\nu=\pbmnuS$ (Lemma~\ref{pb:covol}; $\pbmnuS$ is the
constant $\nu_*$ there, as $\Lambda_*=\Lambda$).

\begin{theorem}[Local systems and Siegel's step]\label{pbm:inputs}
Bessel functions exist for the data (Definition~\ref{pb:bessel}), with \textup{(B5$'$)}, and they
can be chosen, together with functions $J^\tau_i$ $(\tau\in\auPl)$ and $J^T_i$ $(0\le i\le n$; the
$J^T_i$ depend on $T$, $\mathsf A$, $\mathsf d$ as in \textup{(F)}$)$, such that the following
holds, for the entry set $\pbmT$ of \S\ref{pbm:ss-cell}.
\begin{enumerate}
\item[(R)] \emph{(Real places.)} For every $\tau\in\auPl$, $(J^\tau_i)_{0\le i
\le n}$ is a sieve system at the place $(\C,\text{Lebesgue},e^{c-\aubar c},w\mapsto w^{-1})$ with mass $(1+b^{-10})/
\pbdeg_\tau$, weights $\Lambda^\infty$ and decay $\beta^\infty$; and its top form is the Gram form
$J^\tau_n(x,x')=\langle\mathsf N(x),\mathsf N(x')\rangle$ of the normalised orbital integrals of
\S\ref{pbm:ss-remains} (Theorem~\ref{pfr:R}\,(R-b)).
\item[(F)] \emph{(Finite places.)} Let $T$ be a finite set of finite places of $E$, stable under
conjugation, that contains the places above $MN$ and those at which some element of the set
$\mathsf B$ of \S\ref{pbm:ss-cell} is not a unit. For every set $\mathsf A$ of chosen places above
$p_0$ and every family $\mathsf d$ of frame vectors, $(J^{T}_i)_{0\le i\le
n}$ is a sieve system at the place $(E_T,dx,\psi_T,w\mapsto M^{-1}w^{-1})$ with mass $(1+b^{-6})\,b^{-6k\#\mathsf A}$,
weights $\Lambda^T$ and decay $\beta^T$.
\item[(S)] \emph{(Siegel's step.)} Let $T$, $\mathsf A$, $\mathsf d$ be as in \textup{(F)}, and let $J^T_n$
be the top form of \textup{(F)} for the same $T$, $\mathsf A$, $\mathsf d$ (it is built from a group
$\mathsf Z(\mathsf A,\mathsf d)$ that contains the components at $T$ of the integral elements whose
reductions lie in $\pbFix(\mathsf A,\mathsf d)$; see Theorem~\ref{pff:F}). Let
$(c_j,c_j',s_j,\sigma_j)_{j\in\mathsf I}$ be a finite family of pairwise distinct quadruples of
classes $c_j,c_j'$, elements $s_j\in\pbFix(\mathsf A,\mathsf d)$ and $\sigma_j\in\pbmGf(E^+)$ such
that the last entry $t_j=t(\sigma_j)$ lies in $\pbmT$, $|\tau(t_j)|>N$ for all $\tau$, and
$h_{c_j}^{-1}\sigma_jh_{c_j'}$ stabilises the lattice at every finite place outside $T$. Then
\[
\pbmrS^{\,2}\sum_{j\in\mathsf I}\Bigl\|\sum_{p\in\pbmT^\varrho}\pbmOrb_{c_j,s_j,c_j'}(w_p\sigma_j)\Bigr\|^2\le
(1+b^{-6})\,\pbmaS\sum_{t\in\{t_j\}}\ \sum_{p,q\in\pbmT^\varrho}\Rel\Bigl(\prod_\tau J^\tau_n\bigl([p,t]_\tau,
[q,t]_\tau\bigr)\,J^T_n\bigl((p,t),(q,t)\bigr)\Bigr).
\]
\end{enumerate}
\end{theorem}
\smallskip\noindent The existence of the Bessel functions with \textup{(B5$'$)} is
Theorem~\ref{pbe:bessel}. \textup{(R)} is Theorem~\ref{pfr:R}, proved in \S\S\ref{pfr:sec},
\ref{pfs:sec}, \ref{prp:sec} (its hypotheses (H0)--(H5) are checked for the Bessel functions of
Theorem~\ref{pbe:bessel} in \S\ref{prp:sec}); \textup{(F)} and \textup{(S)} are Theorems~\ref{pff:F}
and~\ref{pff:S}, proved in \S\S\ref{pff:sec}, \ref{pfg:sec}. \S\ref{pbm:ss-remains} describes the forms.\par\smallskip

\subsection{The constant of the main term}\label{pbm:ss-const}

\begin{lemma}[Constants]\label{pbm:const}
Assume $k\ge512000\,n^9$, $n\ge100$, $b\ge1000$, $f\le b/1000$, $1\le|\kappa_\tau|$ for all
$\tau$, $M\le b^{23}$, $N\le b^{60n}$, $d\le b^{2kf}$ and $\prod_\tau|\kappa_\tau|=(M\aurd/2\pi)^f$.
Then:
\begin{enumerate}
\item[(C0)] $\displaystyle\pbmaS\prod_{a=0}^{\varrho}c_a\le(1+b^{-9})\,\pbmrS^{\,2}\,\pbdeg\cdot\pbDmain\,
\pbmnuS^{\,2}$;
\item[(C1)] $\displaystyle\pbmaS\prod_{a=0}^{\varrho}c_a\cdot b^{2403\,n^9f}\,\pbmdec\le b^{-10}\,b^{-6kf}\,
\pbmrS^{\,2}\,\pbDmain\,\pbmnuS^{\,2}$.
\end{enumerate}
\end{lemma}
\begin{proof}
From the definitions, $\prod_ac_a/\pbmrS^{\,2}=\pbmvol^{\,n}d^2/(\pbmcG^{\,f}\Lambda_*^2)$. Recall
$\pbDmain=(\#\mathcal C\,d)^2/(\pbdeg\,\pbmass_r)$, $\#\pbR=\#\pbUb\,\#\pbHb$ and $\#\mathcal C=\#\pbHb
(N^k-1)^f\pbmass_k$ (proof of Lemma~\ref{pb:sizes}).

(C0) After cancelling $d^2\,\#\pbUb^2\,\#\pbHb\,\#\mathcal C\,\Lambda_*^{-2}$ the inequality reads
\[
\Bigl(\frac{\pi^k}{(k-1)!\,N^k}\Bigr)^f\aurd^{\,f/2}\,\pbmvol^{\,k+n+2\varrho(\varrho+k)}\le(1+b^{-9})\,\pbmcG^{
\,f}\,\frac{(N^k-1)^f}{N^{2kf}}\cdot\frac{\pbmass_k}{\pbmass_r}\,\aurd^{\,f\varrho}.
\]
By Lemma~\ref{pb:massratio}, $\pbmass_r/\pbmass_k=\bigl(\prod_{l=k}^{r-1}\aurd^{\,l+1/2}\,l!/(2\pi)^{l+1}
\bigr)^f\cdot\mathsf L$ with $\mathsf L=\prod_{l=k}^{r-1}L(l+1,\chi^{\,l+1}_{E/E^+})\le e^{f2^{2-k}}$. We
have $\sum_{l=k}^{r-1}(l+\frac12)=(2\varrho+1)(k+\varrho+\frac12)$ and $\sum_{l=k}^{r-1}l=(2\varrho+1)k+
\varrho(2\varrho+1)$, and
\[
\pbmcG\,\frac{(k-1)!}{\pi^k}=\frac{2\,\pbvol\auU(k)}{\pbvol\auU(r)}=2\prod_{l=k}^{r-1}\frac{l!}{2\pi^{l+1}}=
2^{\,1+\sum_{l=k}^{r-1}l}\,\prod_{l=k}^{r-1}\frac{l!}{(2\pi)^{l+1}} .
\]
Inserting $\pbmvol=(2/\aurd)^f$, the powers of $\aurd$ on the two sides are $\aurd^{f(1/2-k-n-2\varrho(
\varrho+k))}$ and $\aurd^{f(\varrho-(2\varrho+1)(k+\varrho+1/2))}$, which are equal; the powers of $2$
are $2^{f(k+n+2\varrho(\varrho+k))}$ and $2^{f(1+(2\varrho+1)k+\varrho(2\varrho+1))}$, which are equal;
the factorials, the powers of $\pi$ and one factor $N^{-kf}$ cancel. What remains is $\mathsf L\le
(1+b^{-9})(1-N^{-k})^f$, which holds because $\log\mathsf L+2fN^{-k}\le f2^{3-k}\le b^{-10}$. (Unlike
Lemma~\ref{pb:sizes}, this uses the exact constant of the mass formula in
Lemma~\ref{pb:massratio}, that is, $\pbmass_{l+1}/\pbmass_l=2^{-f}|L(-l,\chi_{E/E^+}^{\,l+1})|$.)

(C1) By the same computation, with $\mathsf s=\aurd^{1/2}$, $\mathsf Q=M\mathsf s^2/2\pi$, $x=
\mathsf Q^{2\varrho+1}/(MN)$, so that $\pbmdec=(\pbmTh\,x^{k-n^7})^f$, and $\mathsf D=\varrho(2k+2
\varrho+1)$,
\[
\pbmaS\,\frac{\prod_ac_a}{\pbmrS^{\,2}}\,\pbmdec=\mathsf c\cdot\mathsf A^f,\qquad d\cdot\pbDid\cdot\pbmnuS^{
\,2}=\mathsf c\cdot\mathsf B^f,\qquad\mathsf c=\#\pbR\,\#\mathcal C\,d^2\Lambda_*^{-2},
\]
\[
\mathsf A=M^{\mathsf D}\,\frac{\pi^k\,\mathsf s}{(k-1)!\,N^k}\Bigl(\frac2{\mathsf s^2}\Bigr)^{k+n}\frac{
\pbmTh}{\pbmcG}\,x^{k-n^7},\qquad\mathsf B=\Bigl(\frac{M^{\mathsf D}}{N^k}\Bigl(\frac{\mathsf s^2}2
\Bigr)^{\varrho(k+\varrho)}\mathsf s^{\varrho}\Bigr)^2 .
\]
As $\pbmTh/\pbmcG=\pi^{(2\varrho+1)k}(k-1)!/(2\pi^k)$, a direct computation gives
\[
\mathsf A=\mathsf B\cdot\Bigl(\frac2{\mathsf s^2}\Bigr)^{n+2\varrho^2}\frac{\mathsf s^{1-2\varrho}}{2\,M^{
\varrho(2\varrho+1)}}\cdot x^{-n^7}\le\mathsf B\cdot b^{\,n+2\varrho^2}\,(MN)^{n^7}\le\mathsf B\cdot b^{\,n^9},
\]
because $\mathsf s\ge1$, $\mathsf Q\ge1$ and $MN\le b^{23+60n}$. Hence the left side of (C1) is at
most $b^{2404n^9f}\,d\,\pbDid\,\pbmrS^{\,2}\pbmnuS^{\,2}$. By Lemma~\ref{pb:sizes}, $\pbDid\le b^{-11kf}
\pbDmain$, and $d\le b^{2kf}$; as $2404n^9f+10\le3kf$, (C1) follows.
\end{proof}

The identity behind \textup{(C0)} is the reason for the shape of the main term: the constant of
Siegel's step, times the constants of the $n$ Poisson summations, divided by the normalisations of
the orbital integrals, is $\auDX^2\nu^2/\pbmass_r$ up to the factor $\mathsf L(1-N^{-k})^{-f}$; the
factor $1/\pbdeg$ is the mass of the local systems at the real places.

\subsection{Proof of \textup{(L2)} from Theorem~\ref{pbm:inputs}}\label{pbm:ss-L2}

\begin{theorem}[Theorem~\ref{pb:amp}\,\textup{(L2)}]\label{pbm:L2}
Assume Theorem~\ref{pbm:inputs} and the inequalities on the data listed in Theorem~\ref{pbm:sieve}
and Lemma~\ref{pbm:const}. Then, for the Bessel functions of Theorem~\ref{pbm:inputs}, every set
$\mathsf A$ of chosen places above $p_0$ and every family $\mathsf d$ of frame vectors,
\[
\pbLM\bigl(\{s\in\pbFix(\mathsf A,\mathsf d):s\text{ open}\}\bigr)\le(1+b^{-4})\,b^{-6k\#\mathsf A}\,\pbDmain
\,\nu^2 .
\]
\end{theorem}
\begin{proof}
By Lemma~\ref{pb:covol}, $\nu=\pbmnuS$. The left side is a sum of non-negative terms $\|A^\flat_{
T_\Lambda}(c,s,c')\|^2$ over the \emph{blocks} $(c,s,c',\Lambda)$ with $s\in\pbFix(\mathsf A,\mathsf d)$
open; it suffices to bound the sum over a finite set $\mathsf I$ of blocks with non-zero terms.
Let $T$ be as in Theorem~\ref{pbm:inputs}(F); then $M^{-1},N^{-1}\in\pbmRT$ and $t,t^{-1}\in\pbmRT$
for $t\in\pbmT$. By Proposition~\ref{pbm:classsum}(c) each block $j=(c_j,s_j,c_j',\Lambda_j)$ of
$\mathsf I$ has an element $\sigma_j\in\pbmGf(E^+)$ that determines $\Lambda_j$, so the quadruples
$(c_j,c_j',s_j,\sigma_j)$ are pairwise distinct; $t_j=t(\sigma_j)\in\pbmT$, $|\tau(t_j)|>N$,
$h_{c_j}^{-1}\sigma_jh_{c_j'}$ is integral outside $T$, and the term of the block is the square of
the norm of $\sum_{p\in\pbmT^\varrho}\pbmOrb_{c_j,s_j,c_j'}(w_p\sigma_j)$. By Siegel's step (S) and then
Theorem~\ref{pbm:sieve}, applied to the systems of (R) and (F) and to $\mathsf U=\{t_j\}$,
\[
\pbmrS^{\,2}\sum_{j\in\mathsf I}\|A^\flat_{T_{\Lambda_j}}(c_j,s_j,c_j')\|^2\le(1+b^{-6})\,\pbmaS\prod_{a=0}^{
\varrho}c_a\cdot\Bigl(\frac{(1+b^{-10})^f}{\pbdeg}\,(1+b^{-6})\,b^{-6k\#\mathsf A}+b^{2403n^9f}\,\pbmdec
\Bigr).
\]
By Lemma~\ref{pbm:const}, $(1+b^{-10})^f\le e^{fb^{-10}}\le1+b^{-9}$ and $\#\mathsf A\le f$, the right
side is at most
\[
(1+b^{-6})\bigl((1+b^{-9})^2(1+b^{-6})+b^{-10}\bigr)\,b^{-6k\#\mathsf A}\,\pbmrS^{\,2}\,\pbDmain\,\pbmnuS^{\,2}
\le(1+b^{-4})\,b^{-6k\#\mathsf A}\,\pbmrS^{\,2}\,\pbDmain\,\pbmnuS^{\,2},
\]
since $(1+u^6)((1+u^9)^2(1+u^6)+u^{10})\le1+u^4$ for $0\le u\le10^{-3}$. Divide by $\pbmrS^{\,2}>0$.
\end{proof}

The inequalities on the data used in Theorems~\ref{pbm:sieve}, \ref{pbm:L2} and
Lemma~\ref{pbm:const} are: $n=\varrho+1\ge100$ (so $m=2n$ in the notation of Lemma~\ref{an:lem-sizes}); $k\ge512000\,n^9$; $b\ge1000$, $f\le b/1000$,
$\aurd\le b$, $N\le b^{60n}$; $M\le b^{23}$; $d\le b^{2kf}$ (all in Lemma~\ref{an:lem-sizes}(S1), (S5),
(S7), with (S7)(a), (e), (f) for $n\ge100$, $M\le b^{23}$ and $d\le b^{2kf}$); for
Theorem~\ref{pbm:inputs}\,(R), the hypotheses (H0)--(H5) of Theorem~\ref{pfr:R}, among them
$\varrho\ge10^8$, $k\ge n^7$, $k>10^9n^3+3n+6$ (Lemma~\ref{an:lem-sizes}(S7)(a)) and $b^{6n+2}\ge8r$
((S7)(b)); and the condition \eqref{pbe:eq-kappa} on the generator $\varpi$, $1\le|\kappa_\tau|\le b^{23}$
(Lemma~\ref{pds:kappa}), together with $\prod_\tau|\kappa_\tau|=(M\aurd/2\pi)^f$, which holds because
$\prod_\tau|\tau(\varpi)|^2$ is the norm of the different, $\aurd^{\,2f}$.

\subsection{The forms of Theorem~\ref{pbm:inputs}}\label{pbm:ss-remains}

After this section, \textup{(L2)} rests on Theorem~\ref{pbm:inputs}, on Lemma~\ref{pb:covol}, and,
through Lemma~\ref{pbm:const}, on the mass formula in the form of Lemma~\ref{pb:massratio}. We
describe the forms of the three statements and how they are proved.

\emph{The forms.} At a real place $\tau$ write $\tilde G=G_\tau$ in the normalised basis of
\S\ref{pbe:ss-real}, in which a head $x=(x_0,\dots,x_\varrho)$ of complex numbers defines a Weyl
element $w_x$ and $x_\varrho$ is a last entry, $|x_\varrho|>1$. For $\sigma\in\pbmGf_\tau$ with last
entry $x_\varrho$ the top form is the Gram matrix of normalised orbital integrals,
\begin{gather*}
J^\tau_n(x,x')=\pbip{\mathsf N(x)}{\mathsf N(x')},\\
\mathsf N(x)=\frac{\sqrt{\pbmcG}\,(|x_\varrho|^2-1)^{(k-1)/2}}{\dim\mathbb W_\tau}\prod_{a<\varrho}|x_a|^{
k+1+2(\varrho-a-1)}\int_{G_\tau}W^\flat_\tau(y)\otimes\aubar{W^\flat_\tau((w_x\sigma)^{-1}y)}\,dy,
\end{gather*}
which depends only on $x$, $x'$ when $x_\varrho=x'_\varrho$ (by (B1) for $H_\tau$). At $T$ the
Bessel function is replaced by the vector of the translates, by the cosets of the principal
congruence subgroup in a group $\mathsf Z$ between it and the stabiliser of the lattice, of the
function $\chi_M$ that is $\psi_f(u)$ on $u(K\cap J)$, $u\in U(E_T)$, and $0$ elsewhere; its
orbital integrals are the Kloosterman sums of Proposition~\ref{pbm:orbit}, and $J^T_n(p,q)$ is the
integral, over the points with last entry $p_\varrho$ of the quadric of the vectors of norm $-N$
of $\pbmXp\otimes E_T$, of the pairing of two such orbital integrals, divided by $\pbvol(\mathsf Z)
\pbvol(K\cap J)\,\#\pbUb^{\,2}$ in a suitable normalisation. The group $\mathsf Z$ is the group of
the elements whose reductions lie in a subgroup of $\auG$ containing $\pbFix(\mathsf A,\mathsf d)$.
The lower forms are defined by integration over the entries, as (S1) requires: $J_i$ is the form
$J_n$ integrated over the last $n-i$ entries; by Fourier inversion along the Heisenberg groups of
the frame it is an integral, over a position $g$ and two frames, of the product of two
correlations of the Bessel function, and $J_0$ is the \emph{mass} $\int|\int_{\mathsf Z}\chi_M(ygy^{
-1})dy|^2dg$ (compare Lemma~\ref{pbm:pos}(c)).

\begin{enumerate}
\item[(R)] For the cut Bessel function of a discrete series whose parameter has a large gap: the
absolute convergence of the forms; the mass is at most $(1+b^{-10})/\pbdeg_\tau$ (the positions near the
identity give $1/\pbdeg_\tau$, by (B2) and the reproducing property; the remote positions are small);
the decay $\beta^\infty$ in the frequency, with the gain $(\pi/|\xi|)^{k-n^7}$, and the same gain for two
heads with different entries. This is Theorem~\ref{pfr:R}, proved in \S\S\ref{pfr:sec}, \ref{pfs:sec},
\ref{prp:sec} (sizes and lengths, Plancherel's formula on quadrics, depth of frame vectors, separation
of the two frames by a weight made of depths, the gains at one position, the sum over dyadic shells
of positions), with the bounds for the Bessel function on the open cell of \S\ref{pbo:sec}.
\item[(F)] The same at the finite places, for the function $\chi_M$, with the decay $\beta^T$, and
in addition: the mass of the stabiliser of the lattice is at most $1+b^{-6}$ times its main part
(Corollary~\ref{pff:massK}; the positions outside $K(M)$ are handled by the sum over dyadic shells, and
Theorem~\ref{pbm:E4} is not used), and the mass of $\mathsf Z$ is smaller by $b^{-6k}$ for each fixed
vector (Theorem~\ref{pff:fixed}, through the index of the stabiliser of a vector). This is
Theorem~\ref{pff:F}.
\item[(S)] Siegel's step. The quadruples $(c,c',s,\sigma)$ with a given last entry $t$ are
parametrised, through $\sigma\mapsto\sigma z$, by classes and rational points $x$ of the sphere
$h(x,x)=(t\aubar t-N^2)/N$ of $X_0$, see \eqref{pbm:eq-last}; by Proposition~\ref{pbm:orbit} the
squares of the class sums are values at these points of a non-negative function on the adelic
sphere, which is a pairing of orbital integrals at $T$ times the Gram matrix of orbital integrals
at the real places. The sum of such a function over the rational points of the sphere, and over the
classes, is bounded by $(1+b^{-6})$ times the constant $\pbmaS$ times its integral: Siegel's formula
for the representations of a number by the genus of the hermitian lattice $\auO_E^k$, derived with
all constants from Weil's Siegel--Weil theorem. This is Theorem~\ref{pff:S}, with
Proposition~\ref{pfg:siegel}.
\end{enumerate}


%% ===== proofs/B-multiplicities.tex =====
% proofs/B-multiplicities.tex
% Equality of multiplicities: Proposition au:mult of sec-automorphic.tex.
% Labels of this file start with pm:.
\providecommand{\auO}{\mathcal O}
\providecommand{\auPl}{\Phi}
\providecommand{\auT}{\mathcal T}
\providecommand{\auAz}{\mathcal A_0}
\providecommand{\auG}{\overline G}
\providecommand{\audT}{d_{\auT}}
\providecommand{\auU}{\mathrm U}
\providecommand{\auGL}{\mathrm{GL}}
\providecommand{\auA}{\mathbb A}
\providecommand{\auF}{\mathbb F}
\providecommand{\aubar}[1]{\overline{#1}}
\providecommand{\auf}[1]{\texttt{#1}}
% macros of this file (prefix \pm)
\providecommand{\pmA}{\mathcal A}
\providecommand{\pmtr}{\operatorname{tr}}
\providecommand{\pmvol}{\operatorname{vol}}
\providecommand{\pmsupp}{\operatorname{supp}}
\providecommand{\pmN}{\operatorname{N}}
\providecommand{\pmTr}{\operatorname{Tr}}
\providecommand{\pmAd}{\operatorname{Ad}}
\providecommand{\pmch}{\operatorname{ch}}
\providecommand{\pmS}{\mathfrak S}
\providecommand{\pmn}{\mathfrak n}
\providecommand{\pmg}{\mathfrak g}
\providecommand{\pmGa}{\Gamma}
\providecommand{\pmell}{\mathrm{ell}}
\providecommand{\pmQ}{[G]}
\providecommand{\pmP}{\mathsf P}
\providecommand{\pmE}{\mathfrak E}

\section{Equality of multiplicities}\label{pm:sec}

This section proves Proposition~\ref{au:mult}, in the precise form of
Proposition~\ref{pm:mult}, from two statements on discrete series of the real groups
$\auU(p,q)$ (Theorem~\ref{pm:arch}), known theorems of Harish-Chandra and of
Trombi--Varadarajan and Mili\v{c}i\'c which are proved here for the coefficients we use
(Remark~\ref{pm:rem-other}(1) and \S\ref{pbr:sec}), and from the fact that the irreducible labels
name cuspidal characters (\S\ref{pm:ss-mult}, (G)). Everything else is proved here or quoted from
named results.

\emph{The argument.} Let $V$, $V'$ be two hermitian spaces over $E$ of the same rank that are
isometric at all finite places; the groups $\auU(V)$ and $\auU(V')$ then have the same finite
adelic points $G_f$. At the infinite places we use, instead of a pseudo-coefficient, the kernel
$e=d_\pi\,\aubar{\langle\pi(\cdot)v,v\rangle}$ of a discrete series: it is integrable, because the
parameter is far from the walls, and $R(e)$ is the projection onto the forms of type $\pi$, so that
no other representation with the same infinitesimal character interferes. For a test function
$\varphi$ on $G_f$ the operator $R(e\otimes\varphi)$ has finite rank and its trace is the
integral over the diagonal of the kernel $\sum_\gamma(e\otimes\varphi)(x^{-1}\gamma y)$
(Proposition~\ref{pm:trace}); no truncation is needed, because $e$ is a cusp form in the sense of
Harish-Chandra and this makes all partial sums of the kernel bounded
(Proposition~\ref{pm:bound}). The trace is thus a sum over the characteristic polynomials $R$
of the rational points. If $\varphi$ has a supercuspidal component at a split place, the terms
of the $R$ that are products of two coprime factors vanish (Proposition~\ref{pm:red}; an
argument of Deligne--Kazhdan type), and the terms of the proper powers are
removed by a condition at an auxiliary place (cf.~\cite[Theorem~1]{FlickerKazhdan}, where a
condition at an auxiliary split place removes the regular non-elliptic terms), which is afterwards lifted by linear algebra in a
parahoric Hecke algebra (\S\ref{pm:ss-aux}). For irreducible $R$ the centraliser is the norm one
torus of a field; the conjugacy classes with polynomial $R$ in $\auU(V)$ and in $\auU(V')$ are
described by class field theory, and the sum of the orbital integrals of $e$ over the classes
in a stable class is, by Harish-Chandra, a character of the compact form, which does not see the
signs (Proposition~\ref{pm:ell}). There is no endoscopic term: a torus of this kind lies in no
proper endoscopic group. This is the comparison of Labesse~\cite{Labesse} in a case in which it can
be done by hand.

\subsection{Setting and statements}\label{pm:ss-setting}

\subsubsection*{Groups}
$E$ is a CM field, $F=E^+$, $f=[F:\Q]$, and $\auPl$ is the set of infinite places of $F$, each
with a fixed extension $\tau\colon E\to\C$. Let $(V,h)$ be a non-degenerate hermitian space over
$E$ of dimension $n\ge2$, $G=G_V=\auU(V)$, a reductive group over $F$ with anisotropic centre,
$\pmGa=G(F)$, $G_\infty=\prod_{\tau}G_\tau$ with $G_\tau=\auU(V\otimes_{E,\tau}\C)$, and
$G_f=G(\auA_{F,f})$, the restricted product for the stabilisers of an $\auO_E$-lattice of $V$.
We write $\pmQ=\pmGa\backslash G(\auA_F)$. A finite place $v$ of $F$ that splits in $E$ comes
with a place of $E$ above it; then $G(F_v)=\auGL(V\otimes_EF_v)$, which we write $H_v$, and a basis
identifies $H_v$ with $\auGL_n(F_v)$. For $y\in G_f$, $y_v$ is its component at $v$.
The characteristic polynomial of $\gamma\in\pmGa$, an element of $E[X]$, is $\pmch(\gamma)$; for
split $v$, $\pmch(y_v)\in F_v[X]$ and $E\subset F_v$.

We fix a sign vector $s\in\{\pm1\}^n$ with $p$ entries $+1$ and $q$ entries $-1$, the same for all
$\tau$, and for each $\tau$ an orthogonal basis of $V_\tau=V\otimes_{E,\tau}\C$ on which $h$ is
$\operatorname{diag}(s)$. This identifies $G_\tau$ with $G_s=\{g\in\auGL_n(\C):g^*I_sg=I_s\}$,
$I_s=\operatorname{diag}(s)$; a change of the basis changes the identification by an inner
automorphism. Let $K_s=G_s\cap\auU(n)$, $T\subset K_s$ the diagonal torus, $\|g\|$ the operator
norm of $g\in G_s$ on $\C^n$ (so $\|g\|=\|g^{-1}\|\ge1$, as $g^{-1}=I_sg^*I_s$, and $\|kgk'\|=\|g\|$
for $k,k'\in K_s$), and for $g\in G_\infty$ put $\|g\|=\max_\tau\|g_\tau\|$ and
$\pmP(g)=\prod_\tau\|g_\tau\|$. The dimension of $G_s$ is $n^2$.

\subsubsection*{Measures}
On each $G_\tau$, on $G_f$ and on every other unimodular group we fix a Haar measure; compact
groups at the infinite places have mass $1$. $G(\auA_F)$ has the product measure, $\pmGa$ the
counting measure, and $\pmQ$ the quotient measure, which is finite~\cite[Theorem~5.6(i)]{BorelFin}. The right
regular representation of $G(\auA_F)$ on $L^2(\pmQ)$ is $R$, and
$R(f)\psi(x)=\int_{G(\auA_F)}f(g)\psi(xg)\,dg$ for $f\in L^1$; thus $R(f_1*f_2)=R(f_1)R(f_2)$ with
$(f_1*f_2)(g)=\int f_1(h)f_2(h^{-1}g)\,dh$, and $R(f)^*=R(f^*)$ with $f^*(g)=\aubar{f(g^{-1})}$.
$C_c^\infty(G_f)$ is the space of locally constant functions of compact support (\emph{test
functions}).

\subsubsection*{The input at the infinite places}
A \emph{parameter} is a vector $\lambda\in(\tfrac{n-1}2+\Z)^n$ with distinct entries. The pair
$(\lambda,s)$ satisfies the \emph{gap condition} if
\begin{equation}\label{pm:eq-gap}
|\lambda_i-\lambda_j|\ \ge\ 500\,n^3+10\,n^5\qquad\text{whenever }s_i\ne s_j .
\end{equation}
For a parameter $\lambda$ let $\chi_\lambda$ be the character of the irreducible representation of
the compact group $\auU(n)$ with Harish-Chandra parameter $\lambda$; on
$t=\operatorname{diag}(z_1,\dots,z_n)$ with distinct $z_i$ it is $\det(z_i^{\lambda_j})/\det(z_i^{
\varrho_j})$, up to the sign that makes it the character of a representation, where
$\varrho=(\tfrac{n-1}2,\dots,\tfrac{1-n}2)$. It is a symmetric function of $z_1,\dots,z_n$. For
$\sigma\in S_n$ put $t^\sigma=\operatorname{diag}(z_{\sigma(1)},\dots,z_{\sigma(n)})$.

\begin{theorem}[Harish-Chandra; Trombi--Varadarajan, Mili\v{c}i\'c]\label{pm:arch}
Let $(\lambda,s)$ satisfy \eqref{pm:eq-gap}. There are an irreducible unitary representation
$(\pi,\mathsf H)=\pi(\lambda,s)$ of $G_s$, square-integrable with formal degree $d_\pi$, and a
unit vector $v\in\mathsf H$ contained in a finite-dimensional $K_s$-stable subspace, such that
$c(g)=\langle\pi(g)v,v\rangle$ has the following properties.
\begin{enumerate}
\item[(a)] $\sup_{g\in G_s}\|g\|^{5n^3}\,|c(g)|<\infty$.
\item[(b)] For every $t\in T$ with distinct eigenvalues the function $x\mapsto c(x^{-1}tx)$ is
integrable on $G_s$, and
\[
\sum_{\sigma\in S_n}d_\pi\int_{G_s}\aubar{c(x^{-1}t^\sigma x)}\,dx=p!\,q!\,(-1)^{pq}\,
\aubar{\chi_\lambda(t)} .
\]
\end{enumerate}
\end{theorem}
\noindent Here $\pi(\lambda,s)$ is the discrete series of $G_s\cong\auU(p,q)$ with Harish-Chandra
parameter $\lambda$, for the positive system in which $e_i-e_j$ is positive when
$\lambda_i>\lambda_j$, $v$ is a unit highest weight vector of its lowest $K_s$-type
$\Lambda=\lambda+\rho_{\mathrm n}-\rho_{\mathrm c}$, and inner products are linear in the first
variable, so that $c$ is the function $c_\lambda$ of \S\ref{pco:sec} (written there with the
opposite convention); if $q=0$ it is the representation with character $\chi_\lambda$.
\emph{Proof.} (a) is Remark~\ref{pm:rem-other}(1) (Theorem~\ref{pco:FJ} and Lemma~\ref{pco:decay}).
(b) is Theorem~\ref{pbr:main}, which proves it, together with the unstable form
$d_\pi\int_{G_s}c(x^{-1}tx)\,dx=(-1)^{pq}\sum_{w\in S_p\times S_q}\operatorname{sgn}(w)t^{\kappa\circ w}
/\prod_{\lambda_a>\lambda_b}(t_a-t_b)$, $\kappa=\lambda+\frac{n-1}2\mathbf 1$, from
Theorem~\ref{pco:FJ}, Schur orthogonality, admissibility, Schmid's theorem on the $K$-types of a
discrete series, and the theory of compact Lie groups; Harish-Chandra's character theory is not used.
The product $d_\pi\,dx$ does not depend on the choice of the Haar measure $dx$.

\begin{remark}\label{pm:rem-arch}
(1) By Harish-Chandra, $\Theta_\pi(t)=(-1)^{pq}\sum_{w\in S_p\times S_q}\operatorname{sgn}(w)\,
t^{w\lambda}/D_\lambda(t)$ on the regular set of $T$, with $D_\lambda$ the Weyl denominator of the
positive system of $\lambda$; for the orbital integrals of $c$ this is \eqref{pbr:eq-unstable},
proved in Theorem~\ref{pbr:main}. As $D_\lambda(t^\sigma)=\operatorname{sgn}(\sigma)D_\lambda(t)$, the sum
of $\Theta_\pi(t^\sigma)$ over $\sigma\in S_n$ is $p!q!(-1)^{pq}\sum_{w\in S_n}\operatorname{sgn}(w)
t^{w\lambda}/D_\lambda(t)=p!q!(-1)^{pq}\chi_\lambda(t)$. This is the form (b). (2) The elements
$t^\sigma$, $\sigma\in S_n/(S_p\times S_q)$, represent the conjugacy classes of $G_s$ in the set of
elements with the eigenvalues of $t$: two of them are conjugate in $G_s$ if and only if the same
eigenvalues occur on the positive coordinates.
\end{remark}

\begin{remark}[Relation to other sections]\label{pm:rem-other}
(1) Part (a) of Theorem~\ref{pm:arch} holds for every pair $(\lambda,s)$ satisfying
\eqref{pm:eq-gap}: by Theorem~\ref{pco:FJ} (whose hypothesis $h\ge n^2+1$ follows from
\eqref{pm:eq-gap}) and Lemma~\ref{pco:decay}, $|c(g)|\le(2/\|g\|)^{c}$ for every integer $c$ up to the
gap, which exceeds $5n^3$. Here $c$ is the coefficient of a unit highest weight vector of the lowest
$K$-type, with $\pi(\lambda,s)$ the discrete series of parameter $\lambda$; with inner products
linear in the first variable, $c(g)=\langle\pi(g)v,v\rangle$ is exactly this function $c_\lambda(g)=
\langle v,\pi_\lambda(g)v\rangle$ of \S\ref{pco:sec}, where the opposite convention is used.
(2) The space $\pmA(V;K)$ of Definition~\ref{pm:def-forms} below is defined in $L^2$. By
Proposition~\ref{pm:trace}(a) its elements are bounded and continuous, so for $V=Y_0$ and $K=J_M$
it is the space of forms of type $e_0$ and level $J_M$ of Definition~\ref{pz:def-type}, with the
same sign vector \eqref{pm:eq-signs} as in \eqref{pz:eq-param}; this gives the boundedness that
Proposition~\ref{pds:A0} asks for, under Theorem~\ref{pm:arch}. (3) Lemma~\ref{pm:cusp} proves
the cusp property (A2) of \S\ref{pz:ss-types} for the coefficient $c$ of Theorem~\ref{pm:arch} from
the bound (a).
\end{remark}

In the rest of this section we \emph{assume Theorem~\ref{pm:arch}}. Let $\lambda=(\lambda_\tau)_{
\tau\in\auPl}$ be parameters such that every $(\lambda_\tau,s)$ satisfies \eqref{pm:eq-gap}; let
$\pi_\tau$, $v_\tau$, $c_\tau$ be as in the theorem, transported to $G_\tau$, and
\begin{equation}\label{pm:eq-e}
e_\tau=d_{\pi_\tau}\,\aubar{c_\tau},\qquad e=e_V=\bigotimes_{\tau\in\auPl}e_\tau\colon G_\infty\to\C,
\qquad B=5n^3 .
\end{equation}

\begin{definition}\label{pm:def-forms}
For a compact open subgroup $K\subset G_f$, the space of \emph{forms of type $(\lambda,s)$ and
level $K$} on $V$ is
$\pmA(V;K)=\{\psi\in L^2(\pmQ):R(e)\psi=\psi,\ R(k)\psi=\psi\text{ for }k\in K\}$.
\end{definition}

By Lemma~\ref{pm:kernel}(a), $R(e)$ is an orthogonal projection. For $\psi=R(e)\psi$ the map
$\sum_ia_i\pi_\infty(g_i)v\mapsto\sum_ia_iR(g_i)\psi$, with $\pi_\infty=\bigotimes_\tau\pi_\tau$ and
$v=\bigotimes_\tau v_\tau$, is well defined and is $\|\psi\|$ times an isometry, because
$\langle R(g)\psi,\psi\rangle=\langle R(e)R(g)R(e)\psi,\psi\rangle$ and $e*\delta_g*e=
\langle\pi_\infty(g)v,v\rangle\,e$ by the orthogonality relations. Hence
$\pmA(V;K)\cong\operatorname{Hom}_{G_\infty}(\pi_\infty,L^2(\pmQ)^K)$, $\psi=A(v)\leftrightarrow A$:
the space, with the action of the Hecke algebra of $(G_f,K)$, does not depend on the choice of
$v$, and it is the space of the automorphic forms on $V$ of level $K$ that generate $\pi_\infty$
at infinity, evaluated at the vector $v$. This is the space meant in \S\ref{ss:au-spaces}.

\subsubsection*{The comparison}
A \emph{comparison datum} is $(V,V',\iota)$ with $V$, $V'$ hermitian spaces of dimension $n$, with
sign vectors $s$, $s'$ at every infinite place, and $\iota\colon V\otimes_F\auA_{F,f}\to V'\otimes_F
\auA_{F,f}$ an isometry of $\auA_{E,f}$-modules that carries a lattice to a lattice at almost all
places. It identifies $G_{V'}(\auA_{F,f})$ with $G_f=G_V(\auA_{F,f})$, and $H_v$ for $V'$ with
$H_v$ for $V$. The parameters $\lambda_\tau$ are the same for $V$ and $V'$, and we assume
\eqref{pm:eq-gap} for $(\lambda_\tau,s)$ and for $(\lambda_\tau,s')$; objects of $V'$ carry a
prime. A function $\beta\in C_c^\infty(H_v)$ is a \emph{supercusp function} if
$\int_N\beta(xny)\,dn=0$ for all $x,y\in H_v$ and the unipotent radical $N$ of every proper
parabolic subgroup of $H_v$.

\begin{theorem}\label{pm:main}
Assume Theorem~\ref{pm:arch} and let $f=[F:\Q]$ be even. Let $(V,V',\iota)$ be a comparison datum,
$v_0$ a finite place of $F$ that splits in $E$, and $\varphi=\varphi^{v_0}\otimes\beta\in
C_c^\infty(G_f)$ with $\beta$ a supercusp function on $H_{v_0}$. Then $R(e\otimes\varphi)$ and
$R'(e'\otimes\varphi)$ have finite rank and
\[
\pmtr\bigl(R(e_V\otimes\varphi)\mid L^2(\pmQ)\bigr)=\pmtr\bigl(R'(e_{V'}\otimes\varphi)\mid
L^2([G'])\bigr).
\]
If $\varphi$ is bi-invariant under a compact open $K\subset G_f$, these are the traces of
$R(\varphi)$ on $\pmA(V;K)$ and of $R'(\varphi)$ on $\pmA(V';K)$.
\end{theorem}

The proof occupies \S\S\ref{pm:ss-kernel}--\ref{pm:ss-aux}. We now state the consequence for the
data. In the notation of \S\ref{ss:au-spaces}, $n=r$, and the parameter is $\varsigma_\tau$, whose
entries we name: $\varsigma_\tau(i)=\tfrac k2-i$ for $1\le i\le k-1$ (the \emph{small} entries) and
$\varsigma_\tau(k-1+a)=-\kappa(\tau)-a\Delta$ for $1\le a\le m$. The sign vectors are
\begin{equation}\label{pm:eq-signs}
s_{W_j}(k-1+a)=\begin{cases}+1,&a\le j,\\-1,&a>j,\end{cases}\qquad
s_{Y_0}(k-1+a)=(-1)^{a-1},\qquad s(i)=+1\ \text{for }i\le k-1 ,
\end{equation}
so that $W_j$ has $m-j$ and $Y_0$ has $m/2=\varrho+1$ negative signs, as in
\S\ref{ss:au-spaces}; the discrete series at $\tau$ is $\pi(\varsigma_\tau,s)$. The spaces
$\auAz$ and $\pmA(W_j)$ of \S\ref{ss:au-spaces} are $\pmA(Y_0;J_M)$ and $\pmA(W_j;J_M)$. Let
$\mathcal H^M$ be the algebra of the functions in $C_c^\infty(G_f)$ of the form
$\varphi_S\otimes1_{J_M^{S}}$, with $S$ a finite set of places of $F$ that split in $E$ and do not
divide $M$, $K^\circ\supset J_M$ the stabiliser of the lattice, $\varphi_S$ bi-invariant under
$K^\circ_S=J_{M,S}$, and $J_M^S$ the part of $J_M$ outside $S$; it is commutative (the spherical
Hecke algebras of the groups $\auGL_r(F_v)$ are), stable under $\varphi\mapsto\varphi^*$, and it acts
on $\pmA(V;J_M)$ by $\varphi\mapsto\pmvol(J_M)^{-1}R(\varphi)$, commuting with $\auG$. These are the Hecke operators
at the split places prime to $M$.

\begin{proposition}[Proposition~\ref{au:mult}, precise form]\label{pm:mult}
Assume Theorem~\ref{pm:arch} and the property \textup{(G)} of \S\ref{pm:ss-mult}. Assume that $f$
is even and that for each of the sign vectors \eqref{pm:eq-signs} and each $\tau$ the entries of
$\varsigma_\tau$ on coordinates of different signs differ by at least $500r^3+10r^5$. (Both hold
for the data: $f=\#(\Z/q_1)^\times$, and these entries differ by at least $\Delta=2r^8\ge500r^3+10r^5$.) Then for
$1\le j\le m$ and every $T\in\mathcal H^M$
\[
\pmtr\bigl(T\mid e_0\pmA(W_j)\bigr)=\pmtr\bigl(T\mid\auT\bigr).
\]
In particular $\dim e_0\pmA(W_j)=\audT$, and every character of $\mathcal H^M$ has the same
multiplicity in $e_0\pmA(W_j)$ and in $\auT$.
\end{proposition}

\subsection{The kernel at infinity}\label{pm:ss-kernel}

\begin{lemma}\label{pm:kernel}
Let $e_\tau$ be as in \eqref{pm:eq-e}, on $G_\tau\cong G_s$.
\begin{enumerate}
\item[(a)] $e_\tau$ is continuous, $e_\tau^*=e_\tau$, $|e_\tau(g)|\le C\|g\|^{-B}$,
$e_\tau\in L^1$, and $e_\tau*e_\tau=e_\tau$.
\item[(b)] There is $\alpha_\tau\in C_c^\infty(G_\tau)$ with $\alpha_\tau*e_\tau=e_\tau$. Hence
$e_\tau=e_\tau*\alpha_\tau^*$, $e_\tau$ is smooth, and $|(De_\tau)(g)|\le C_D\|g\|^{-B}$ for every
left-invariant differential operator $D$.
\item[(c)] $\int_{G_\tau}\|g\|^{-B'}dg<\infty$ for $B'>n^2$.
\end{enumerate}
The same holds for $e$ on $G_\infty$ with $\alpha=\bigotimes\alpha_\tau$, with $\pmP(g)^{-B}$ in place
of $\|g\|^{-B}$.
\end{lemma}
\begin{proof}
(c) In the decomposition $G_s=K_sA^+K_s$, with $A\cong\R^{\min(p,q)}$ and coordinates
$t_1\ge\dots\ge t_{\min(p,q)}\ge0$ on $A^+$, one has $\|kak'\|=e^{t_1}$ and the Haar measure is
$\delta(a)\,dk\,da\,dk'$ with $\delta(a)\le e^{2\rho(\log a)}\le e^{n(n-1)t_1}$
(\cite[Ch.~I, \S5, no.~2]{Helgason}, \cite[Lemma~2.4.2]{WallachRRG1}: up to a constant, $\delta(a)$ is the product of $\sinh\alpha(\log a)$ over the positive restricted roots with multiplicities). (a) Continuity and $e^*=e$ are clear, the bound is
Theorem~\ref{pm:arch}(a), integrability follows from (c), and $e*e=e$ is the orthogonality relation
$\int c(h)c(h^{-1}g)\,dh=d_\pi^{-1}c(g)$ for square-integrable
representations~\cite[Proposition~9.6, p.~284]{Knapp}. (b) Let $\mathsf H_0$ be the sum of the $K_\tau$-isotypic
subspaces of the types occurring in the span of $K_\tau v$; it is finite-dimensional, since
irreducible unitary representations are admissible \cite[Theorem~8.1, p.~205]{Knapp}, and the projection
onto it is $\pi(\chi)$ for a central function $\chi$ on $K_\tau$. The subspace
$\{\pi(\chi)\pi(\alpha_0)v:\alpha_0\in C_c^\infty(G_\tau)\}$ of $\mathsf H_0$ contains $v$ in its
closure, so it contains $v$: $\pi(\alpha_1)v=v$ with $\alpha_1=\chi*\alpha_0\in C_c^\infty$. Put
$\alpha=\alpha_1$. Then, as $e(h^{-1}y)=d_\pi\aubar{\langle\pi(y)v,\pi(h)v\rangle}$ and the inner product is conjugate-linear in the second variable, $(\alpha*e)(y)=d_\pi\aubar{\langle\pi(y)v,\pi(\alpha_1)v\rangle}=e(y)$.
Taking adjoints, $e=e*\alpha^*$, so $De=e*D\alpha^*$, and $|(e*\beta)(g)|\le
\int|e(gh^{-1})||\beta(h)|dh\le C\|g\|^{-B}\int\|h\|^B|\beta(h)|dh$.
\end{proof}

\begin{lemma}[$e$ is a cusp form]\label{pm:cusp}
Let $N'$ be the unipotent radical of a proper parabolic subgroup of $G_\tau$. Then
$\int_{N'}|e_\tau(x\nu y)|\,d\nu<\infty$ and $\int_{N'}e_\tau(x\nu y)\,d\nu=0$ for all $x,y\in G_\tau$.
\end{lemma}
\begin{proof}
This is Harish-Chandra's theorem that the coefficients of the discrete series are cusp forms; we
give his proof under the bound of Lemma~\ref{pm:kernel}. Write $e=e_\tau$ and $P'=M'A'N'$ for the
Langlands decomposition of the parabolic subgroup, with $A'=\exp\mathfrak a'\ne1$ consisting of
positive hermitian matrices, and $\mathfrak n'$ for the Lie algebra of $N'$. For
$\nu=\exp X\in N'$ we have $(\nu-1)^{n}=0$, so $|X|\le C\|\nu\|^{n-1}$ and
$\int_{N'}\|\nu\|^{-B/4}d\nu\le C\int_{\mathfrak n'}(1+|X|)^{-B/(4(n-1))}dX<\infty$, as
$B/(4(n-1))>n^2\ge\dim\mathfrak n'$. For $a\in A'$ the eigenvalues of $a\nu$ are those of $a$, so
$\|a\nu\|\ge\|a\|$, and $\|\nu\|\le\|a^{-1}\|\|a\nu\|\le\|a\nu\|^2$; hence
$\|a\nu\|\ge\|a\|^{1/2}\|\nu\|^{1/4}$. Fix $x$ and put $\Phi(g)=\int_{N'}e(x\nu g)\,d\nu$. As
$\|a\nu\|\le\|x\|\|xa\nu g\|\|g\|$, the substitution $\nu\mapsto a\nu a^{-1}$ gives
\begin{equation}\label{pm:eq-cuspdecay}
|\Phi(ag)|=\delta(a)\Bigl|\int_{N'}e(xa\nu g)\,d\nu\Bigr|\le C\|x\|^B\|g\|^B\,\|a\|^{2n^2-B/2}
\int_{N'}\|\nu\|^{-B/4}d\nu ,
\end{equation}
where $\delta(a)=|\det\pmAd(a)|_{\mathfrak n'}|\le\|a\|^{2n^2}$; in particular the integral defining
$\Phi$ converges absolutely, and by Lemma~\ref{pm:kernel}(b) it may be differentiated under the
integral sign. Thus $\Phi$ is smooth, left invariant under $N'$, and an eigenfunction of the
centre $Z(\pmg)$ of the enveloping algebra acting on the right, because $e$ is ($\pi$ is
irreducible, so $Z(\pmg)$ acts by scalars on the smooth vector $v$). For central $z$ the right
action of $z$ is the left action of its transpose, which is central. By Harish-Chandra, every $z\in Z(\pmg)$ lies in $\mu(z)+U(\pmg)\mathfrak n'$ with
$\mu(z)\in Z(\mathfrak m'\oplus\mathfrak a')$, and $Z(\mathfrak m'\oplus\mathfrak a')$ is a finitely
generated module over $\mu(Z(\pmg))$. (In \cite[proof of Proposition~8.22, p.~225]{Knapp}, $Z(\pmg)\subseteq\mathfrak n'U(\pmg)\oplus U(\mathfrak m'\oplus\mathfrak a')$, and the projection $\mu$ to the second summand satisfies $\gamma'_{\mathfrak a'\oplus\mathfrak m'}\circ\mu=\gamma'$, (8.34) there; applying the principal antiautomorphism of $U(\pmg)$, which preserves $Z(\pmg)$ and $U(\mathfrak m'\oplus\mathfrak a')$, gives the form with $U(\pmg)\mathfrak n'$. For the finiteness: via these maps, with a Cartan subalgebra $\mathfrak h$ of $\mathfrak m'\oplus\mathfrak a'$, $\mu(Z(\pmg))$ and $Z(\mathfrak m'\oplus\mathfrak a')$ correspond to the images of $U(\mathfrak h)^W$ and $U(\mathfrak h)^{W_{M'}}$ under translation automorphisms $\sigma_1,\sigma_2$ of $U(\mathfrak h)$ \cite[Theorem~8.18]{Knapp}; $U(\mathfrak h)$ is a finitely generated module over $U(\mathfrak h)^W$ \cite[Theorem~8.19]{Knapp}, hence over the Noetherian ring $\sigma_1(U(\mathfrak h)^W)$, so its submodule $\sigma_2(U(\mathfrak h)^{W_{M'}})$ is finitely generated.) Elements of $U(\pmg)\mathfrak n'$, acting on the left, kill
left $N'$-invariant functions. So $\Phi$ is an eigenfunction of the left action of
$\mu(Z(\pmg))$, and so are its left derivatives along $H\in\mathfrak a'$, which are again left
$N'$-invariant; as $H\in Z(\mathfrak m'\oplus\mathfrak a')$ is integral over $\mu(Z(\pmg))$, the
function $u(r)=\Phi(\exp(rH)g)$ satisfies a linear differential equation with constant
coefficients: it is an exponential polynomial. By \eqref{pm:eq-cuspdecay} and $2n^2<B/2$ it tends
to $0$ for $r\to\pm\infty$ if $H\ne0$, so $u=0$ and $\Phi(g)=u(0)=0$.
\end{proof}

\subsection{Reduction theory and bounds for the kernel}\label{pm:ss-bound}

Let $\varrho_V$ be the Witt index of $V$. If $\varrho_V=0$ then $\pmQ$ is compact~\cite[Theorem~5.6(ii)]{BorelFin} and
everything in this subsection is trivial or void; we assume $\varrho_V\ge1$ and fix a Witt
decomposition $V=\langle e_1,\dots,e_{\varrho_V}\rangle\oplus W\oplus\langle e_1',\dots,
e'_{\varrho_V}\rangle$, with $h(e_i,e'_j)=\delta_{ij}$, $e_i$, $e_i'$ isotropic and orthogonal to the
anisotropic space $W$. Let $V_i=\langle e_1,\dots,e_i\rangle$, $P_i$ the stabiliser of $V_i$ (the
maximal standard parabolic subgroups), $N_i$ its unipotent radical, $P_0=\bigcap_iP_i$. Every proper
parabolic $F$-subgroup of $G$ is conjugate under $\pmGa$ to a subgroup of some $P_i$, and an element
of $\pmGa$ lies in a proper parabolic $F$-subgroup if and only if it stabilises a non-zero isotropic
subspace of $V$. For $a=(a_1,\dots,a_{\varrho_V})\in\R_{>0}^{\varrho_V}$ we also write $a$ for the
element of $G_\infty$ that, at every $\tau$, multiplies $e_i$ by $a_i$, $e'_i$ by $a_i^{-1}$ and is
the identity on $W$. Put $\beta_i(a)=a_i/a_{i+1}$ for $i<\varrho_V$, $\beta_{\varrho_V}(a)=
a_{\varrho_V}$, $\beta(a)=\max_i\beta_i(a)$. We choose the bases at infinity (that is, the maximal
compact subgroup $K_\infty=\prod K_\tau$ and the norms) so that $e_i,e'_i$ and an orthogonal basis
of $W_\tau$ are orthonormal for the norm $|\cdot|_\tau$ of $\C^n$; then $a$ is positive and
diagonal, $|h_\tau(u,u')|\le|u|_\tau|u'|_\tau$, and $\|a\|=\max_i\max(a_i,a_i^{-1})$. Let $K_f^\circ$
be the stabiliser of a lattice $\Lambda\subset V$ containing the Witt basis.

\emph{Siegel sets} (\cite[Th\'eor\`eme~7]{Godement63}; \cite[\S I.2.1, p.~20]{MW}). There are a compact set
$\Omega\subset P_0(\auA_F)$ and $t\in(0,1]$ such that $G(\auA_F)=\pmGa\,\pmS$ with
\[
\pmS=\{\omega ak:\ \omega\in\Omega,\ a\in\R_{>0}^{\varrho_V},\ \beta_i(a)\ge t\ (1\le i\le\varrho_V),\
k\in K_\infty K_f^\circ\}.
\]
For such $a$: $a_i\ge t^{n}$ for all $i$, and $\|a\|\le t^{-n}\max(1,\beta(a))^{n/2}$. For
$x=\omega ak\in\pmS$ we let $i_0=i_0(x)$ be an index with $\beta_{i_0}(a)=\beta(a)$ and put
$P_x=P_{i_0}$, $N_x=N_{i_0}$. In the estimates below, constants depend on $V$, $\Omega$, $t$, the
function $e$ and the test function, but not on $x$.

\begin{lemma}\label{pm:siegel}
Let $x=\omega ak\in\pmS$, $P=P_x$, $N=N_x$.
\begin{enumerate}
\item[(a)] Let $\Omega_f\subset G_f$ be compact. There is $c>0$ such that
$\pmP\bigl((x^{-1}\gamma x)_\infty\bigr)\ge c\,\beta(a)^{f}$ for every $\gamma\in\pmGa\smallsetminus
P(F)$ with $(x^{-1}\gamma x)_f\in\Omega_f$.
\item[(b)] Let $K'\subset G_f$ be compact open. There is $C$ such that $N(\auA_F)\,x\subset
N(F)\,x\,(Z_\delta\times K')$ with $\delta=C\beta(a)^{-1}$, where $Z_\delta=\{z\in G_\infty:
\|z_\tau-1\|\le\delta\text{ for all }\tau\}$.
\item[(c)] Let $B'>n^2$ and $M_{B'}(y)=\pmP(y_\infty)^{-B'}1_{\Omega_f}(y_f)$. There is $C$ such that
$\sum_{\gamma\in\pmGa}M_{B'}(g^{-1}\gamma x)\le C\|a\|^{2fn^2}$ for all $g\in G(\auA_F)$. For $x$ in a compact subset of $G(\auA_F)$, not necessarily in $\pmS$, the series converges
uniformly in $(g,x)$.
\end{enumerate}
\end{lemma}
\begin{proof}
Let $\mathcal B$ be the basis of $V$ formed by the $e_i$, a basis of $W$ in $\Lambda$ and the
$e'_i$, and for $b\in\mathcal B$ let $\chi_b(a)\in\{a_i,1,a_i^{-1}\}$ be the eigenvalue of $a$. Split
$\mathcal B=\mathcal B_1\sqcup\mathcal B_2\sqcup\mathcal B_3$ with $\mathcal B_1=\{e_l:l\le i_0\}$,
a basis of $V_{i_0}$, $\mathcal B_3=\{e'_l:l\le i_0\}$, and $\mathcal B_1\sqcup\mathcal B_2$ a basis
of $V_{i_0}^\perp$. Put $\theta=t^{-n}a_{i_0}^{-1}$ and $\Theta=t^{-n}\max(1,a_{i_0+1})$, with
$a_{\varrho_V+1}:=1$. As $a_l\ge t^na_{i_0}$ for $l\le i_0$ and $t^n\le a_l\le t^{-n}a_{i_0+1}$ for
$l>i_0$,
\begin{equation}\label{pm:eq-contract}
\chi_b^{-1}\le\theta\ (b\in\mathcal B_1),\qquad\chi_b\le\theta\ (b\in\mathcal B_3),\qquad
\chi_b^{\pm1}\le\Theta\ (b\in\mathcal B_2),\qquad\theta\max(\theta,\Theta)\le t^{-5n}\beta(a)^{-1},
\end{equation}
the last because $a_{i_0}=\beta_{i_0}(a)a_{i_0+1}\ge t^n\beta(a)$ and $\beta(a)\ge t$.

(a) As $\gamma V_{i_0}\not\subset V_{i_0}=(V_{i_0}^\perp)^\perp$, there are $b\in\mathcal B_1$ and
$b'\in\mathcal B_1\sqcup\mathcal B_2$ with $\xi=h(b',\gamma b)\ne0$. Let $\Omega^{\mathrm{fin}}$ be
the set of finite components of $\Omega$ and $\Omega'_f=\Omega^{\mathrm{fin}}K_f^\circ$, a compact
set containing $x_f$. The finite component of $\gamma$ lies in the compact set
$\Omega'_f\Omega_f\Omega_f'^{-1}$, so $\gamma\Lambda\subset d^{-1}\Lambda$ and $\xi\in d^{-1}\auO_E$
for an integer $d\ge1$ independent of $x$ and $\gamma$; hence $\prod_\tau|\tau(\xi)|\ge d^{-f}$. On
the other hand
$|\tau(\xi)|=|h_\tau(x_\tau^{-1}b',(x^{-1}\gamma x)_\tau x_\tau^{-1}b)|\le|x_\tau^{-1}b'||x_\tau^{-1}b|
\,\|(x^{-1}\gamma x)_\tau\|$, and $x_\tau^{-1}b=k_\tau^{-1}a^{-1}\omega_\tau^{-1}b$, where
$\omega_\tau^{-1}b$ is a bounded combination of the vectors of $\mathcal B_1$, since $P_0$
preserves $V_{i_0}$; likewise $\omega_\tau^{-1}b'$ is a bounded combination of the vectors of
$\mathcal B_1\sqcup\mathcal B_2$. By \eqref{pm:eq-contract},
$|x_\tau^{-1}b'||x_\tau^{-1}b|\le C\theta\max(\theta,\Theta)\le C'\beta(a)^{-1}$, and the product over
$\tau$ gives (a).

(b) $N(\auA_F)=N(F)\,(C_\infty\times U_f)$ for every open compact subgroup $U_f$ of
$N(\auA_{F,f})$ and a suitable compact $C_\infty\subset N(F\otimes\R)$ (strong approximation for
unipotent groups). Take $U_f\subset\bigcap_{y\in\Omega'_f}yK'y^{-1}$. For $\nu\in N(F)$,
$c\in C_\infty$, $u\in U_f$ we have $\nu cu\,x=\nu x\cdot(x^{-1}cux)$, the finite component of
$x^{-1}cux$ lies in $K'$, and its infinite component is $k^{-1}a^{-1}n'ak$ with
$n'=\omega_\infty^{-1}c\,\omega_\infty$ in a compact subset of $N(F\otimes\R)$. As $n'-1$ maps
$V$ into $V_{i_0}^\perp$, $V_{i_0}^\perp$ into $V_{i_0}$ and $V_{i_0}$ to $0$, its matrix in the basis
$\mathcal B$ has a non-zero entry from $b$ to $b'$ only if $b'\in\mathcal B_1$,
$b\in\mathcal B_2\sqcup\mathcal B_3$, or $b'\in\mathcal B_2$, $b\in\mathcal B_3$; conjugation by
$a^{-1}$ multiplies this entry by $\chi_b(a)\chi_{b'}(a)^{-1}\le\theta\max(\theta,\Theta)$, by
\eqref{pm:eq-contract}. So $\|a^{-1}n'a-1\|\le C\beta(a)^{-1}$, and $k$ is unitary.

(c) Let $U_x=\{u\in G_\infty:\|u_\tau-1\|\le r_0\|a\|^{-2}\}\times K''$ with $r_0\le1$ and a
compact open subgroup $K''$. Then $xU_xU_x^{-1}x^{-1}$ is contained in $\bigcup_{\omega}\omega U_0U_0^{-1}\omega^{-1}\times
\bigcup_{y}yK''y^{-1}$, the unions over the infinite components $\omega$ of the elements of $\Omega$
and over $y\in\Omega'_f$, where $U_0=\{\|u_\tau-1\|\le r_0\}$, because
$a\,u\,a^{-1}-1=a(u-1)a^{-1}$ has norm at most $\|a\|^2\|u-1\|$. These sets shrink to $\{1\}$ with
$r_0$ and $K''$, and $\pmGa$ is discrete; so for small $r_0$, $K''$ the sets $g^{-1}\gamma xU_x$,
$\gamma\in\pmGa$, are disjoint for all $g$ and all $x\in\pmS$. As $\pmP(yu)\le2^f\pmP(y)$ for
$u\in U_x$, we get $M_{B'}(y)\le2^{fB'}M'(yu)$ with $M'=\pmP^{-B'}\otimes1_{\Omega_fK''}$, and
\[
\sum_\gamma M_{B'}(g^{-1}\gamma x)\le\frac{2^{fB'}}{\pmvol(U_x)}\sum_\gamma\int_{U_x}M'(g^{-1}\gamma
xu)\,du\le\frac{2^{fB'}}{\pmvol(U_x)}\int_{G(\auA_F)}M'.
\]
The integral is finite by Lemma~\ref{pm:kernel}(c), and $\pmvol(U_x)\ge c(r_0\|a\|^{-2})^{fn^2}$,
since $\dim G_\infty=fn^2$. The same computation, with $U_x$ replaced by a fixed small neighbourhood of $1$ and for the sum
over the complement of a finite set of $\gamma$, shows that the series converges uniformly for
$g\in G(\auA_F)$ and $x$ in any compact subset of $G(\auA_F)$.
\end{proof}

\begin{lemma}[Decay in a cusp]\label{pm:decay}
For every integer $j\ge0$ there is $\varepsilon_0\in(0,1]$ with the following property. Let $N$
be the unipotent radical of a proper parabolic $F$-subgroup of $G$, $x\in G(\auA_F)$,
$K'\subset G_f$ compact open and $\delta>0$ with $N(\auA_F)x\subset N(F)\,x\,(Z_\delta\times K')$.
Let $\mathsf F$ be a continuous function on $G(\auA_F)$ with $|\mathsf F|\le T$, left invariant
under $N(F)$, right invariant under $K'$, and with $\mathsf F(g)=\int_{G_\infty}\mathsf F(gh)e(h)\,dh$
for all $g$. Then $|\mathsf F(x)|\le(\delta/\varepsilon_0)^j\,T$.
\end{lemma}
\begin{proof}
\emph{The constant term vanishes.} Let $\mathsf F_N(g)=\int_{N(F)\backslash N(\auA_F)}\mathsf F(ng)\,
dn$ (mass $1$). It is bounded, left invariant under $N(\auA_F)$, and again
$\mathsf F_N(g)=\int\mathsf F_N(gh)e(h)dh$. Fix $\tau$ and write $h=h^\tau h_\tau$. For fixed $g$,
$h^\tau$ the function $\Psi(h_\tau)=\mathsf F_N(gh^\tau h_\tau)$ on $G_\tau$ is left invariant under
$N'=g_\tau^{-1}N(F_\tau)g_\tau$, the unipotent radical of a proper parabolic subgroup of $G_\tau$, so
$\int_{G_\tau}\Psi(h_\tau)e_\tau(h_\tau)dh_\tau=\int_{N'\backslash G_\tau}\Psi(h)\int_{N'}e_\tau(nh)\,dn\,dh
=0$ by Lemma~\ref{pm:cusp}. Hence $\mathsf F_N=0$.

\emph{Oscillation.} By Lemma~\ref{pm:kernel}(b), $\mathsf F(g)=\int\mathsf F(gh)\alpha(h)\,dh$.
Let $\Omega_\alpha=Z_1\pmsupp\alpha$, $c_\alpha=\sup_{\Omega_\alpha}\|h\|$, and
$\omega(\varepsilon)=\sup_{z\in Z_\varepsilon}\int|\alpha(z^{-1}h)-\alpha(h)|dh$, which tends to $0$
with $\varepsilon$. For $z\in Z_\varepsilon$, $\varepsilon\le1$,
\begin{equation}\label{pm:eq-osc}
|\mathsf F(gz)-\mathsf F(g)|=\Bigl|\int\mathsf F(gh)\bigl(\alpha(z^{-1}h)-\alpha(h)\bigr)dh\Bigr|
\le\omega(\varepsilon)\sup_{h\in\Omega_\alpha}|\mathsf F(gh)| .
\end{equation}
Choose $\varepsilon_0$ with $\omega(\varepsilon_0)c_\alpha^{2j}\le\tfrac12$. Let
$m(h)=|\mathsf F(xh)|$ for $h\in G_\infty$ and $C^*=\sup_hm(h)\|h\|^{-2j}\le T$. For $h\in G_\infty$,
$N(\auA_F)xh\subset N(F)\,xh\,(Z_{\delta\|h\|^2}\times K')$, as $h^{-1}Z_\delta h\subset
Z_{\delta\|h\|^2}$. If $\delta\|h\|^2\le\varepsilon_0$, then, since $\mathsf F_N=0$ and by
\eqref{pm:eq-osc},
\[
m(h)=|\mathsf F(xh)-\mathsf F_N(xh)|\le\sup_{z\in Z_{\delta\|h\|^2}}|\mathsf F(xhz)-\mathsf F(xh)|
\le\omega(\varepsilon_0)\sup_{h'\in\Omega_\alpha}m(hh')\le\tfrac12C^*\|h\|^{2j}.
\]
If $\delta\|h\|^2>\varepsilon_0$, then $m(h)\le T\le T(\delta/\varepsilon_0)^j\|h\|^{2j}$. Therefore
$C^*\le\max(\tfrac12C^*,T(\delta/\varepsilon_0)^j)$, that is, $C^*\le T(\delta/\varepsilon_0)^j$, and
$|\mathsf F(x)|=m(1)\le C^*$.
\end{proof}

A subset $S\subset\pmGa$ is \emph{admissible} if it is stable under conjugation by $\pmGa$ and
$N(F)\cdot(S\cap P(F))=S\cap P(F)$ for every proper parabolic $F$-subgroup $P$ with unipotent
radical $N$. For a set $\mathcal R$ of polynomials, $\pmGa_{\mathcal R}=\{\gamma:\pmch(\gamma)\in
\mathcal R\}$ is admissible, because $\gamma\in P(F)$ and $\nu\gamma$, $\nu\in N(F)$, induce the
same maps on the graded pieces of the flag stabilised by $P$. Let $\pmGa_\pmell$ be the set of
the $\gamma$ that stabilise no non-zero isotropic subspace; it contains every $\gamma$ with
irreducible $\pmch(\gamma)$.

\begin{proposition}\label{pm:bound}
Let $\varphi'\in C_c^\infty(G_f)$ and $\mathsf f'=e\otimes\varphi'$.
\begin{enumerate}
\item[(a)] The series $k_{\mathsf f'}(x,y)=\sum_{\gamma\in\pmGa}\mathsf f'(x^{-1}\gamma y)$ converges absolutely and
uniformly on compact sets, and is continuous on $\pmQ\times\pmQ$.
\item[(b)] $\sup_{x\in G(\auA_F)}\sum_{\gamma\in\pmGa_\pmell}|\mathsf f'(x^{-1}\gamma x)|<\infty$.
\item[(c)] There is $C$ such that $\bigl|\sum_{\gamma\in S}\mathsf f'(x^{-1}\gamma x)\bigr|\le C$ for all
$x\in G(\auA_F)$ and all admissible $S$.
\end{enumerate}
\end{proposition}
\begin{proof}
Let $\Omega_f=\pmsupp\varphi'$ and let $K'$ be a compact open subgroup with
$\varphi'(k^{-1}y)=\varphi'(y)$ for $k\in K'$. By Lemma~\ref{pm:kernel}, $|\mathsf f'|\le CM_B$, and
$M_B\le c^{-j}\beta^{-fj}M_{B-j}$ on the set where $\pmP(y_\infty)\ge c\beta^f$. (a) follows from
Lemma~\ref{pm:siegel}(c). The sums in (b), (c) are left invariant under $\pmGa$ as functions of
$x$, so we may take $x=\omega ak\in\pmS$. If $\beta(a)\le1$, then $x$
is in a compact set and the bound is (a). Let $\beta(a)>1$ and $P=P_x$, $N=N_x$. Take
$j=n^3$; then $B-j>n^2$, and by Lemma~\ref{pm:siegel}(a),(c)
\begin{equation}\label{pm:eq-far}
\sum_{\gamma\in\pmGa\smallsetminus P(F)}|\mathsf f'(x^{-1}\gamma x)|\le C\beta(a)^{-fn^3}\sum_\gamma M_{B-j}
(x^{-1}\gamma x)\le C'\beta(a)^{-fn^3}\|a\|^{2fn^2}\le C'' ,
\end{equation}
as $\|a\|\le t^{-n}\beta(a)^{n/2}$. This proves (b), since $\pmGa_\pmell\cap P(F)=\emptyset$. For
(c) it remains to bound $\mathsf F(x)$, where $\mathsf F(g)=\sum_{\gamma\in S\cap P(F)}\mathsf f'(g^{-1}\gamma
x)$. By Lemma~\ref{pm:siegel}(c), $\mathsf F$ is continuous and $|\mathsf F|\le T=C\|a\|^{2fn^2}\le
C_1\beta(a)^{fn^3}$. It is right invariant under $K'$, left invariant under $N(F)$ because $S$ is
admissible, and $\int\mathsf F(gh)e(h)dh=\mathsf F(g)$ because $e*e=e$ (the exchange of sum and
integral is allowed by the bound for $\sum|\mathsf f'|$). By Lemma~\ref{pm:siegel}(b) and
Lemma~\ref{pm:decay} with $j=fn^3$, $|\mathsf F(x)|\le(C\beta(a)^{-1}/\varepsilon_0)^{fn^3}C_1
\beta(a)^{fn^3}$, a constant.
\end{proof}

\subsection{Forms, the trace, and the terms of the characteristic polynomials}\label{pm:ss-trace}

\begin{proposition}\label{pm:trace}
\begin{enumerate}
\item[(a)] For every compact open $K\subset G_f$ the space $\pmA(V;K)$ is finite-dimensional, and
its elements are bounded continuous functions.
\item[(b)] For $\varphi\in C_c^\infty(G_f)$ and $\mathsf f=e\otimes\varphi$, $R(\mathsf f)$ has finite rank and
$\pmtr R(\mathsf f)=\int_{\pmQ}k_{\mathsf f}(x,x)\,dx$. If $\varphi$ is bi-invariant under $K$, then
$\pmtr R(\mathsf f)=\pmtr(R(\varphi)\mid\pmA(V;K))$.
\item[(c)] For a polynomial $R$ put $J_R(\varphi)=J_R^V(\varphi)=\int_{\pmQ}\sum_{\pmch(\gamma)=R}
\mathsf f(x^{-1}\gamma x)\,dx$. This is defined for every $\varphi\in C_c^\infty(G_f)$; $J_R$ is linear,
$J_R(\varphi^g)=J_R(\varphi)$ for $g\in G_f$ and $\varphi^g(y)=\varphi(g^{-1}yg)$; for split $v$,
$J_R(\varphi)=0$ if $\varphi$ vanishes on $\{y\in G_f:\pmch(y_v)=R\}$; and
$\pmtr R(\mathsf f)=\sum_RJ_R(\varphi)$, the sum over all polynomials being unconditionally convergent.
\end{enumerate}
\end{proposition}
\begin{proof}
(a) Let $\mathsf f_0=\pmvol(K)^{-1}e\otimes1_K$. Then $\mathsf f_0^*=\mathsf f_0=\mathsf f_0*\mathsf f_0$, and $R(\mathsf f_0)$ is the orthogonal
projection onto $\pmA(V;K)$. For $\psi\in C_c(\pmQ)$, unfolding gives $R(\mathsf f_0)\psi(x)=\int_{\pmQ}
k_{\mathsf f_0}(x,y)\psi(y)\,dy$, and $\int_{\pmQ}|k_{\mathsf f_0}(x,y)|^2dy=k_{\mathsf f_0*\mathsf f_0^*}(x,x)=k_{\mathsf f_0}(x,x)\le C$ by
Proposition~\ref{pm:bound}(c) with $S=\pmGa$. So $|R(\mathsf f_0)\psi(x)|\le C^{1/2}\|\psi\|_2$; by
continuity every $\psi\in\pmA(V;K)$ is represented by $x\mapsto\langle\psi,\aubar{k_{\mathsf f_0}(x,\cdot)}
\rangle$, a continuous function bounded by $C^{1/2}\|\psi\|_2$. For orthonormal $\psi_1,\dots,\psi_d
\in\pmA(V;K)$, Bessel's inequality gives $\sum_i|\psi_i(x)|^2\le k_{\mathsf f_0}(x,x)\le C$, so
$d\le C\pmvol(\pmQ)$.

(b) Let $\varphi$ be invariant under $K$ on both sides (shrink $K$). Then $R(\mathsf f)=R(\mathsf f)R(\mathsf f_0)=R(\mathsf f_0)
R(\mathsf f)$, so $R(\mathsf f)$ has finite rank, with image and coimage in $\pmA(V;K)$; with an orthonormal
basis $(\psi_i)$ of $\pmA(V;K)$, the function $\sum_i(R(\mathsf f)\psi_i)(x)\aubar{\psi_i(y)}$ is a
continuous kernel of $R(\mathsf f)$. So is $k_{\mathsf f}$ (Proposition~\ref{pm:bound}(a) and unfolding). Two
continuous kernels of the same operator are equal; on the diagonal, $\int k_{\mathsf f}(x,x)dx=\sum_i\langle
R(\mathsf f)\psi_i,\psi_i\rangle=\pmtr R(\mathsf f)$. The last assertion follows from $R(\mathsf f)=R(\varphi)R(e)$, as
$R(e)$ and $R(\varphi)$ commute and $R(\varphi)$ preserves $L^2(\pmQ)^K$, on whose orthogonal complement
it vanishes.

(c) The integrand of $J_R$ is bounded by Proposition~\ref{pm:bound}(c), and $\pmQ$ has finite
volume. Invariance is the substitution $x\mapsto xg$, and the support condition holds because
$\pmch((x^{-1}\gamma x)_v)=\pmch(\gamma)$. Finally $k_{\mathsf f}(x,x)$ is the sum over $R$ of the integrands,
the partial sums over finite sets $\mathcal R$ of polynomials are bounded by the constant of
Proposition~\ref{pm:bound}(c), and dominated convergence applies.
\end{proof}

\subsection{Polynomials that are not irreducible}\label{pm:ss-red}

For a finite split place $u$ of $F$ let $Z_u\subset H_u\cong\auGL_n(F_u)$ be the set of matrices
with entries in $\auO_u$ whose reduction has characteristic polynomial with a \emph{simple} root
$0$. A test function $\varphi$ \emph{satisfies $(Z)$ at $u$} if $y_u\in Z_u$ for all
$y\in\pmsupp\varphi$.

\begin{proposition}\label{pm:red}
Let $\varphi=\varphi^{v_0}\otimes\beta$ with $\beta$ a supercusp function on $H_{v_0}$.
\begin{enumerate}
\item[(a)] If $R=AB$ with $A,B\in E[X]$ monic, coprime and not constant, then $J_R(\varphi)=0$.
\item[(b)] If $\varphi$ satisfies $(Z)$ at some $u$ and $R=Q^j$ with $j\ge2$, then $J_R(\varphi)=0$.
\end{enumerate}
Hence, if $\varphi$ satisfies $(Z)$ at some $u$, then $J_R(\varphi)=0$ for every monic $R$ that
is not irreducible.
\end{proposition}
\begin{proof}
(b) If $g\in Z_u$ had $\pmch(g)=Q^j$ over $F_u$, then $Q$ would be monic with coefficients in
$\auO_u$ (Gauss) and every root of the reduction of $\pmch(g)$ would have multiplicity divisible
by $j$. So $\varphi$ vanishes on $\{y:\pmch(y_u)=R\}$, and Proposition~\ref{pm:trace}(c) applies.

(a) Write $H=H_{v_0}$, $\mathsf C=\{h\in H:\pmch(h)=AB\}$, a closed subset, and
$D(\beta')=J_R(\varphi^{v_0}\otimes\beta')$ for $\beta'\in C_c^\infty(H)$. By
Proposition~\ref{pm:trace}(c), $D$ is invariant under conjugation by $H$ and vanishes on the
functions that vanish on $\mathsf C$; so it is an invariant distribution on $\mathsf C$
(restriction $C_c^\infty(H)\to C_c^\infty(\mathsf C)$ is surjective~\cite[\S1.8, Proposition]{BZ76}). We show
that every such distribution kills supercusp functions. Let $d=\deg A$, $\epsilon$ the projector
of $F_{v_0}^n$ onto the span of the first $d$ basis vectors, $M\cong\auGL_d\times\auGL_{n-d}$ its
centraliser in $H$, $N$ the group of the $1+\epsilon z(1-\epsilon)$, $P=MN$, $K_0=\auGL_n(
\auO_{v_0})$, and $\mathsf X=\{(m_1,m_2)\in M:\pmch(m_1)=A,\ \pmch(m_2)=B\}$. Choose
$a,b\in F_{v_0}[X]$ with $aA+bB=1$. For $h\in\mathsf C$, $(bB)(h)$ is the projector onto
$\ker A(h)$ along $\ker B(h)$, of rank $d$, and it depends continuously on $h$. Hence the map
$\varpi\colon H\times\mathsf X\to\mathsf C$, $(g,m)\mapsto gmg^{-1}$, is surjective, its fibres are
the orbits of $M$ acting by $(g,m)\mu=(g\mu,\mu^{-1}m\mu)$, and it has continuous local sections
(because $H\to H/M$, $g\mapsto g\epsilon g^{-1}$, has: $H/M$ is homeomorphic to the orbit of
$\epsilon$ under conjugation \cite[\S1.6]{BZ76}, and for a projector $e$ in that orbit near
$\epsilon$ the element $s(e)=e\epsilon+(1-e)(1-\epsilon)$ is invertible, depends continuously on
$e$, and satisfies $s(e)\epsilon=e\epsilon=e\,s(e)$, so $s(e)\epsilon s(e)^{-1}=e$; translating by $H$
gives local sections everywhere). For
$\Psi\in C_c^\infty(H\times\mathsf X)$ put $\varpi_*\Psi(gmg^{-1})=\int_M\Psi(g\mu,\mu^{-1}m\mu)d\mu$,
an element of $C_c^\infty(\mathsf C)$. The functional $\Psi\mapsto D(\varpi_*\Psi)$ is invariant
under left translations of the first variable; as $C_c^\infty(H\times\mathsf X)=C_c^\infty(H)\otimes
C_c^\infty(\mathsf X)$ and a left-invariant functional on $C_c^\infty(H)$ is a multiple of the
Haar integral, it is of the form $\Psi\mapsto\nu\bigl(m\mapsto\int_H\Psi(g,m)dg\bigr)$. Thus
\begin{equation}\label{pm:eq-fibre}
\int_H\Psi(g,m)\,dg=0\ \text{ for all }m\in\mathsf X\quad\Longrightarrow\quad D(\varpi_*\Psi)=0 .
\end{equation}
Now let $\beta$ be a supercusp function. Replacing it by its mean over the conjugates by $K_0$
changes neither $D(\beta)$ nor the supercusp property, so let $\beta$ be invariant under
conjugation by $K_0$, and put $\Psi(g,m)=1_{K_0N}(g)\,\beta(gmg^{-1})$. For $g=k\nu$ with $k\in K_0$,
$\nu\in N$, $\Psi(g,m)=\beta(\nu m\nu^{-1})=\beta(m\cdot\nu_m)$, $\nu_m=m^{-1}\nu m\nu^{-1}$. For
$m\in\mathsf X$ the map $\nu\mapsto\nu_m$ is a homeomorphism of $N$ which multiplies the Haar
measure by a constant, and it depends continuously on $m$: in the coordinate $z$ it is
$z\mapsto m_1^{-1}zm_2-z$, an invertible linear map because $m_1$, $m_2$ have no common eigenvalue.
Hence $\Psi$ is locally constant with compact support ($K_0N$ is open and closed, and on the
support of $\Psi$ the element $m\nu_m$ stays in the compact set $\pmsupp\beta$, so $m$, $\nu_m$,
and $\nu$ modulo $N\cap K_0$, stay in compact sets), and $\int_H\Psi(g,m)dg$ is a constant times
$\int_N\beta(mn)\,dn=0$. Moreover $\varpi_*\Psi=\pmvol(M\cap K_0)\cdot\beta|_{\mathsf C}$: for
$g=k\nu\mu_0$ (Iwasawa decomposition $H=K_0NM$) the set of $\mu\in M$ with $g\mu\in K_0N$ is
$\mu_0^{-1}(M\cap K_0)$, because $P\cap K_0N=(M\cap K_0)N$. By \eqref{pm:eq-fibre}, $D(\beta)=0$.

For the last statement, a monic $R$ that is neither irreducible nor a proper power has two
distinct monic irreducible factors, and is of the form $AB$ as in (a).
\end{proof}

\subsection{Irreducible polynomials}\label{pm:ss-ell}

Let $R\in E[X]$ be monic and irreducible of degree $n$ with $\pmGa_R=\{\gamma\in\pmGa:
\pmch(\gamma)=R\}\ne\emptyset$ (for $V$ or for $V'$). Let $L=E[X]/(R)$, $y$ the class of $X$. For
$\gamma\in\pmGa_R$ the adjoint involution of $V$ induces on $E[\gamma]\cong L$ the involution
$*$ with $y^*=y^{-1}$ extending the conjugation of $E$; it depends only on $R$. Let $L^+$ be its
fixed field, $\pmN(z)=zz^*$ the norm of $L/L^+$, and $T_R$ the torus over $F$ with
$T_R(F)=\{z\in L^\times:\pmN(z)=1\}$; it is anisotropic, so $T_R(F)\backslash T_R(\auA_F)$ is
compact. The centraliser of $\gamma$ in $G$ is $T_R$, through $L\cong E[\gamma]$.

\emph{The invariant.} $V$ is a line over $E[\gamma]$. Choosing $v_0\ne0$ we get
$h(a(\gamma)v_0,b(\gamma)v_0)=\pmTr_{L/E}(ab^*c)$ for a unique $c\in L^{+\times}$, and changing
$v_0$ multiplies $c$ by a norm: $\operatorname{inv}(\gamma)\in L^{+\times}/\pmN(L^\times)$ is well
defined. Write $V_c$ for the space $L$ with the form $\pmTr_{L/E}(ab^*c)$, on which $y$ is an
isometry. Two elements $\gamma_1,\gamma_2\in\pmGa_R$ are conjugate in $\pmGa$ if and only if
$\operatorname{inv}(\gamma_1)=\operatorname{inv}(\gamma_2)$: an isometry $g$ of $V$ with
$g\gamma_1g^{-1}=\gamma_2$ is the same as $z\in L^\times$ with $c_1=c_2\pmN(z)$, by
$g(a(\gamma_1)v_1)=a(\gamma_2)zv_2$. The same holds over every completion $F_w$ of $F$, with
$L\otimes_FF_w$ in place of $L$, and over $\auA_{F,f}$ if $z$ is taken in $(L\otimes\auA_{F,f})^\times$.

\emph{Signs.} Assume that for every $\tau$ all roots of $\tau(R)$ have absolute value $1$. Then
$*$ is the complex conjugation in every embedding, $L$ is a CM field and $L^+$ is totally real;
its real places $u$ correspond to the pairs $(\tau,z)$ with $z$ a root of $\tau(R)$, and
$h_\tau$ has the sign of $u(c)$ on the $z$-eigenline of $\gamma$ in $V_\tau$. We write
$\varepsilon(\gamma)=(\operatorname{sgn}u(c))_u$ and let $\pmE_q$ be the set of the families
$\varepsilon=(\varepsilon_u)\in\{\pm1\}^{\{u\}}$ with exactly $q$ entries $-1$ above each $\tau$.
Thus $\varepsilon(\gamma)\in\pmE_q$ for $\gamma\in\pmGa_R$.

\begin{lemma}\label{pm:classes}
Let $f$ be even, $R$ as above with all roots of all $\tau(R)$ of absolute value $1$, and let
$X_R(V)$ be the set of $\pmGa$-conjugacy classes in $\pmGa_R$. The map
\[
X_R(V)\longrightarrow\{\text{conjugacy classes of }G_f\}\times\pmE_q,\qquad
[\gamma]\mapsto\bigl([\gamma]_{G_f},\varepsilon(\gamma)\bigr)
\]
is injective, and its image is $\mathsf C_R\times\pmE_q$ for a set $\mathsf C_R$ of conjugacy
classes of $G_f$. The set $\mathsf C_R$ is the same for $V$ and for $V'$ (with $G_{V'}(\auA_{F,f})$
identified with $G_f$ by $\iota$).
\end{lemma}
\begin{proof}
\emph{Class field theory.} (i) An element of $L^{+\times}$ that is a norm from $L\otimes_{L^+}L^+_u$
at every place $u$ of $L^+$ is a norm from $L$ (Hasse's norm theorem for the quadratic extension
$L/L^+$, \cite[Chapter~VI, Corollary~(4.5)]{Neukirch}). (ii) If $\eta_u\in\{\pm1\}$ is given at every real place of
$L^+$ with $\prod_u\eta_u=1$, there are $x\in L^{+\times}$ and $z\in(L\otimes\auA_{\Q,f})^\times$
with $\operatorname{sgn}u(x)=\eta_u$ for all real $u$ and $x\,\pmN(z)=1$ in $L^+\otimes\auA_{\Q,f}$:
the idele of $L^+$ with the components $\eta_u$ at the real places and $1$ elsewhere has trivial
Artin symbol for $L/L^+$, so it is the product of an element $x$ of $L^{+\times}$ and a norm
$\pmN(z')$ of an idele of $L$ (the kernel of the global norm residue symbol is $L^{+\times}\pmN(\auA_L^\times)$, \cite[Chapter~VI, Theorem~(5.5)]{Neukirch}); take for $z$ the finite part of $z'$; note that $\pmN(z'_u)>0$ at the real places. (iii) Two hermitian
spaces over $E$ that are isometric at all places of $F$ are isometric
(Landherr~\cite{Landherr}; \cite[Chapter~10, Theorem~6.1]{Scharlau}).

\emph{Injectivity.} Let $\gamma_1,\gamma_2\in\pmGa_R$ have the same image, with invariants
$c_1,c_2$. Then $c_1/c_2$ is positive at all real places and, being $\pmN(z)$ for some
$z\in(L\otimes\auA_{F,f})^\times$, a local norm at all finite places. By (i) it is a norm, and
$\gamma_1,\gamma_2$ are conjugate.

\emph{Change of signs and of space.} Let $\gamma\in\pmGa_R$ with invariant $c$, and
$\varepsilon'\in\pmE_{q'}$. We claim: there is $\gamma'\in G_{V'}(F)$ with $\pmch(\gamma')=R$,
$\varepsilon(\gamma')=\varepsilon'$, and $\gamma'$ conjugate in $G_f$ to $\iota\gamma\iota^{-1}$.
Indeed $\prod_u\varepsilon(\gamma)_u\varepsilon'_u=(-1)^{fq}(-1)^{fq'}=1$, as $f$ is even; let
$x$, $z$ be as in (ii) for $\eta=\varepsilon(\gamma)\varepsilon'$. The space $V_{cx}$ has the
signature of $V'$ at every $\tau$, and multiplication by $z$ is an isometry $V_c\otimes\auA_{F,f}
\to V_{cx}\otimes\auA_{F,f}$ commuting with $y$, because $\pmN(z)x=1$; so $V_{cx}$ is isometric to $V_c\cong V$, hence to
$V'$, at every finite place. By (iii) there is an isometry $j\colon V_{cx}\to V'$ over $E$, and
$\gamma'=jyj^{-1}$ has the three properties: the element $j\circ z\circ(\text{the isometry }V\cong
V_c)\circ\iota^{-1}$ of $G_{V'}(\auA_{F,f})$ conjugates $\iota\gamma\iota^{-1}$ into $\gamma'$.

Applied to $V'=V$, $\iota=1$, the claim shows that the image of the map is a product
$\mathsf C_R(V)\times\pmE_q$. Applied to $(V,V')$ and to $(V',V,\iota^{-1})$ it shows
$\mathsf C_R(V)=\mathsf C_R(V')$.
\end{proof}

\begin{lemma}[Selberg's principle]\label{pm:selberg}
Let $R$ be irreducible and suppose that some $\tau(R)$ has a root of absolute value $\ne1$. Then
$J_R(\varphi)=0$ for all $\varphi\in C_c^\infty(G_f)$.
\end{lemma}
\begin{proof}
By Proposition~\ref{pm:bound}(b), sum and integral in $J_R(|\varphi|)$ (with $|e|$) converge, and
unfolding gives, with $\mathsf f=e\otimes\varphi$, $J_R(\varphi)=\sum_{[\gamma]\in X_R(V)}\int_{T_R(F)\backslash G(\auA_F)}
\mathsf f(x^{-1}\gamma x)dx$, each integral being absolutely convergent. Fix $\gamma$. The integral is
$\pmvol(T_R(F)\backslash T_R(\auA_F))$ times the product of $\int_{T_R(\auA_{F,f})\backslash G_f}
\varphi(x^{-1}\gamma x)dx$ and of the integrals $I_{\tau'}=\int_{T_R(F_{\tau'})\backslash G_{\tau'}}
e_{\tau'}(x^{-1}\gamma x)\,dx$, and the same holds with absolute values, in $[0,\infty]$; so either one of the factors with
absolute values is $0$, and then the term of $\gamma$ is $0$, or every $I_{\tau'}$ converges
absolutely. In the second case we show $I_\tau=0$. The eigenvalues of
$\gamma$ on $V_\tau$ are distinct, and eigenvectors for $z,z'$ are orthogonal unless
$z\aubar{z'}=1$. So the span $V_>$ of the eigenlines with $|z|>1$ is isotropic and not zero, its
stabiliser $P_1=M_1N_1$ is a proper parabolic subgroup of $G_\tau$, where we take for $M_1$ the
stabiliser of the decomposition $V_\tau=V_>\oplus V_=\oplus V_<$ into the spans of the eigenlines
with $|z|>1$, $|z|=1$, $|z|<1$; this is a Levi factor of $P_1$, since $V_>^\perp=V_>\oplus V_=$ and
the hermitian form puts $V_<$ in duality with $V_>$. The element $\gamma$ and its centraliser
$T_R(F_\tau)$ preserve every eigenline of $\gamma$, hence lie in $M_1$; and $\pmAd(\gamma)$ has no eigenvalue $1$ on $\pmn_1$ (its eigenvalues
are quotients $z/z'$ of eigenvalues of $\gamma$ with $|z|>|z'|$). By the Iwasawa decomposition $G_\tau=M_1N_1K_\tau$,
\[
I_\tau=\int_{T_R(F_\tau)\backslash M_1}\int_{K_\tau}\int_{N_1}e_\tau\bigl(k^{-1}n^{-1}\gamma_mn\,k\bigr)
\,dn\,dk\,\delta(m)dm,\qquad\gamma_m=m^{-1}\gamma m,
\]
and $n\mapsto\gamma_m^{-1}n^{-1}\gamma_mn$ is a bijection of $N_1$ onto itself that multiplies the
Haar measure by $|\det(\pmAd(\gamma_m^{-1})-1)|_{\pmn_1}|\ne0$ (induction on the descending
central series of $N_1$). So the inner integral is a multiple of $\int_{N_1}e_\tau(k^{-1}\gamma_m
nk)\,dn$, which is $0$ by Lemma~\ref{pm:cusp}.
\end{proof}

\begin{proposition}\label{pm:ell}
Let $f$ be even and $R$ monic irreducible of degree $n$. Then
$J_R^V(\varphi)=J_R^{V'}(\varphi)$ for every $\varphi\in C_c^\infty(G_f)$.
\end{proposition}
\begin{proof}
If $\pmGa_R$ and $\pmGa'_R$ are empty, both sides vanish. If some $\tau(R)$ has a root off the
unit circle, both vanish by Lemma~\ref{pm:selberg}. Otherwise, as in the proof of that lemma,
\[
J_R^V(\varphi)=v_R\sum_{[\gamma]\in X_R(V)}O_{[\gamma]_{G_f}}(\varphi)\prod_{\tau}\int_{G_\tau}
e_\tau(x^{-1}\gamma x)\,dx,
\]
where $v_R=\pmvol(T_R(F)\backslash T_R(\auA_F))$ and $O_\kappa(\varphi)=\int_{T_R(\auA_{F,f})
\backslash G_f}\varphi(x^{-1}\gamma_\kappa x)dx$ for $\gamma_\kappa$ in the class $\kappa$; the
measure of $T_R(\auA_{F,f})$ is fixed once and transported to the centraliser of $\gamma_\kappa$
through $L$ (this is compatible with conjugation), and $T_R(F\otimes\R)$, a product of circles,
has mass $1$. So $v_R$ and $O_\kappa$ are the same for $V$ and $V'$. In the identification
$G_\tau\cong G_s$ the element $\gamma$ is conjugate to $t^\sigma$, where
$t=\operatorname{diag}(z_1,\dots,z_n)$ lists the roots of $\tau(R)$ and $\sigma$ is any
permutation such that $s_i$ is the sign of $h_\tau$ on the eigenline of $z_{\sigma(i)}$; for each
$\varepsilon_\tau$ there are $p!q!$ such $\sigma$. Write $I_\tau(\varepsilon_\tau)$ for the integral over $G_\tau$ when $\gamma$ has the signs
$\varepsilon_\tau$ above $\tau$. By Lemma~\ref{pm:classes} the sum over $X_R(V)$ is a sum over
$\mathsf C_R\times\pmE_q$, absolutely convergent also with $|\varphi|$ and $|e_\tau|$, and by
Theorem~\ref{pm:arch}(b)
\[
\mathsf X:=\sum_{\varepsilon\in\pmE_q}\prod_\tau I_\tau(\varepsilon_\tau)=\prod_\tau\Bigl(\frac1{p!q!}
\sum_{\sigma\in S_n}\int_{G_s}e_\tau(x^{-1}t_\tau^\sigma x)\,dx\Bigr)=(-1)^{pqf}\prod_\tau
\aubar{\chi_{\lambda_\tau}(t_\tau)} .
\]
If the products $\prod_\tau I_\tau(\varepsilon_\tau)$, $\varepsilon\in\pmE_q$, do not all vanish,
then $\sum_{\kappa\in\mathsf C_R}O_\kappa(|\varphi|)<\infty$ and $J_R^V(\varphi)=v_R\bigl(\sum_\kappa
O_\kappa(\varphi)\bigr)\mathsf X$; if they all vanish, $J_R^V(\varphi)=0$ and $\mathsf X=0$. As $f$
is even, $\mathsf X=\prod_\tau\aubar{\chi_{\lambda_\tau}(t_\tau)}$ is the same number for $V$ and
for $V'$. If it is $0$, both $J_R^V(\varphi)$ and $J_R^{V'}(\varphi)$ vanish. If not, the first case
holds for both spaces, with the same $v_R$, $\mathsf C_R$ and $O_\kappa$.
\end{proof}

\begin{remark}\label{pm:rem-even}
The hypothesis that $f$ is even is used twice, in the proof of Lemma~\ref{pm:classes} and at the
end of the last proof, to make $(-1)^{f(q+q')}$ and $(-1)^{pqf}(-1)^{p'q'f}$ equal to $1$. Both hold
for every comparison datum: the quotient of the discriminants of $V$ and $V'$ is an element of
$F^\times$ that is a norm from $E$ at all finite places and has the sign $(-1)^{q-q'}$ at all
$f$ real places, so $f(q-q')$ is even by the product formula for the norm residue symbol of
$E/F$; and $pq\equiv(n-1)q\pmod 2$. We keep the hypothesis because it holds for the data.
\end{remark}

\begin{corollary}\label{pm:simple}
Theorem~\ref{pm:main} holds for the test functions $\varphi=\varphi^{v_0}\otimes\beta$ that
satisfy $(Z)$ at some finite split place $u$.
\end{corollary}
\begin{proof}
By Proposition~\ref{pm:trace}, $\pmtr R(e\otimes\varphi)=\sum_RJ_R^V(\varphi)$. The $R$ that are
not monic of degree $n$ do not occur; those that are not irreducible contribute $0$ by
Proposition~\ref{pm:red}; the others have $J^V_R(\varphi)=J^{V'}_R(\varphi)$ by
Proposition~\ref{pm:ell}.
\end{proof}

\subsection{The auxiliary place: end of the proof of Theorem~\ref{pm:main}}\label{pm:ss-aux}

\begin{lemma}\label{pm:line}
Let $\mathsf B$, $\mathsf B'$ be finite-dimensional Hilbert spaces with operators $(A,T,\epsilon)$
and $(A',T',\epsilon')$ such that $\epsilon^2=\epsilon$, $TT^*=d(1+c\,\epsilon)$ with real $d,c>0$,
$A$ commutes with $T$ and $T^*$, and likewise for the primed operators with the same $d$, $c$. If
$\pmtr(A^lT^i)=\pmtr(A'^lT'^i)$ for all $l,i\ge1$, then $\pmtr(A\epsilon)=\pmtr(A'\epsilon')$.
\end{lemma}
\begin{proof}
$T$ is invertible, and $\epsilon$ is a polynomial in $TT^*$, so $A$ commutes with $\epsilon$. For
$\alpha\in\C$ let $\mathsf B_\alpha$ be the generalised eigenspace of $A$; it is stable under $T$,
$T^*$, $\epsilon$, and $T^*|\mathsf B_\alpha$ is the adjoint of $T|\mathsf B_\alpha$. On
$\mathsf B_\alpha$, $A-\alpha$ is nilpotent and commutes with $T$, so $\pmtr(A^lT^i)=\sum_\alpha
\alpha^l\pmtr(T^i|\mathsf B_\alpha)$. As the sequences $(\alpha^l)_{l\ge1}$, $\alpha\ne0$, are
linearly independent, $\pmtr(T^i|\mathsf B_\alpha)=\pmtr(T'^i|\mathsf B'_\alpha)$ for all
$\alpha\ne0$, $i\ge1$. The power sums of positive index of the eigenvalues of an invertible
operator determine these eigenvalues with their multiplicities; hence
$\dim\mathsf B_\alpha=\dim\mathsf B'_\alpha$ and $|\det(T|\mathsf B_\alpha)|=|\det(T'|\mathsf B'_
\alpha)|$. Now $|\det(T|\mathsf B_\alpha)|^2=\det(TT^*|\mathsf B_\alpha)=d^{\dim\mathsf B_\alpha}
(1+c)^{\operatorname{rk}(\epsilon|\mathsf B_\alpha)}$, so $\epsilon$ and $\epsilon'$ have the same
rank on $\mathsf B_\alpha$ and $\mathsf B'_\alpha$, for $\alpha\ne0$. Finally
$\pmtr(A\epsilon)=\sum_\alpha\alpha\operatorname{rk}(\epsilon|\mathsf B_\alpha)$.
\end{proof}

\begin{lemma}\label{pm:hecke}
Let $u$ be a finite split place of $F$, $\mathbb F$ its residue field, $\mathsf q=\#\mathbb F$,
$K_u=\auGL_n(\auO_u)\subset H_u$, $J\subset K_u$ the subgroup of the matrices whose reduction maps
the first basis vector to a multiple of itself, and $t=\operatorname{diag}(\varpi,1,\dots,1)$ with a
uniformiser $\varpi$. Then
\begin{enumerate}
\item[(a)] $1_{JtJ}*1_{Jt^{-1}J}=\pmvol(J)\,\mathsf q^{n-2}\bigl(1_J+(\mathsf q-1)1_{K_u}\bigr)$;
\item[(b)] for $i\ge1$ the $i$-fold convolution power of $1_{JtJ}$ is supported in $Z_u$.
\end{enumerate}
\end{lemma}
\begin{proof}
(a) $tJt^{-1}$ is the subgroup $J'$ of the matrices in $K_u$ whose first row is
$(\ast,0,\dots,0)$ modulo $\varpi$. So $JtJt^{-1}J=JJ'J\subset K_u$; as $J$ is the preimage of a
maximal parabolic subgroup $\mathsf P$ of $\auGL_n(\mathbb F)$, $\mathsf P\backslash\auGL_n(
\mathbb F)/\mathsf P$ has two elements, and $J'\not\subset J$, we get $JJ'J=K_u$. The convolution
is bi-invariant under $J$, so it is $x_11_J+x_21_{K_u\smallsetminus J}$. Its value at $1$ is
$\pmvol(JtJ)=\pmvol(J)[J:J\cap J']=\pmvol(J)\mathsf q^{n-1}$, and its integral is
$\pmvol(JtJ)^2$; with $[K_u:J]=(\mathsf q^n-1)/(\mathsf q-1)$ this gives $x_1=\pmvol(J)\mathsf
q^{n-1}$ and $x_2=\pmvol(J)\mathsf q^{n-2}(\mathsf q-1)$. (b) The reduction of an element of $JtJ$
kills the first basis vector and induces an invertible map of $\mathbb F^n$ modulo that vector.
Matrices with entries in $\auO_u$ with this property are stable under products, and the
characteristic polynomial of their reduction is $X$ times a polynomial with
non-zero constant term.
\end{proof}

\begin{proof}[Proof of Theorem~\ref{pm:main}]
The finiteness of the rank and the last assertion are Proposition~\ref{pm:trace}(b). Choose a
compact open subgroup $K=\prod_wK_w$ of $G_f$ such that $\varphi$ is bi-invariant under $K$, and a
finite place $u\ne v_0$ of $F$, split in $E$, with $K_u=\auGL_n(\auO_u)$ in a suitable basis and
$\varphi=\varphi^u\otimes1_{K_u}$; almost all split places qualify. Let $J$, $t$ be as in
Lemma~\ref{pm:hecke} and $K^\flat=K^u\times J$. On $\mathsf B=\pmA(V;K^\flat)$, which is
finite-dimensional, consider
\[
A=\frac{R(\varphi^u\otimes1_J)}{\pmvol(K^u)\pmvol(J)},\qquad T=\frac{R(1_{JtJ})}{\pmvol(J)},\qquad
\epsilon=\frac{R(1_{K_u})}{\pmvol(K_u)},
\]
where functions on $H_u$ act through $R$ restricted to $H_u$. Then $\epsilon$ is the orthogonal
projection of $\mathsf B$ onto $\pmA(V;K)$, $T^*=R(1_{Jt^{-1}J})/\pmvol(J)$, $A$ commutes with $T$
and $T^*$, and by Lemma~\ref{pm:hecke}(a), $TT^*=\mathsf q^{n-2}(1+(\mathsf q^n-1)\epsilon)$ on
$\mathsf B$, since $R(1_J)=\pmvol(J)$ there and $\pmvol(K_u)=\pmvol(J)(\mathsf q^n-1)/(\mathsf q-1)$.
For $l,i\ge1$ the operator $A^lT^i$ is a positive multiple, depending only on $l,i$ and the
measures, of $R(\varphi_{l,i})$ on $\mathsf B$, with $\varphi_{l,i}=(\varphi^u)^{*l}\otimes(1_{JtJ})^{*i}$.
This test function is bi-invariant under $K^\flat$, satisfies $(Z)$ at $u$ by
Lemma~\ref{pm:hecke}(b), and is of the form $\varphi_{l,i}^{v_0}\otimes\beta^{*l}$; a convolution
power of a supercusp function is a supercusp function, because
$\int_N(\beta*\beta')(xny)dn=\int\beta'(h)\int_N\beta(xnyh^{-1})\,dn\,dh$. By
Corollary~\ref{pm:simple} and Proposition~\ref{pm:trace}(b), $\pmtr(A^lT^i|\mathsf B)=\pmtr(A'^l
T'^i|\mathsf B')$, where the primed objects are those of $V'$. Lemma~\ref{pm:line} gives
$\pmtr(A\epsilon|\mathsf B)=\pmtr(A'\epsilon'|\mathsf B')$. On the image $\pmA(V;K)$ of $\epsilon$,
$A=\pmvol(K)^{-1}R(\varphi)$. Hence $\pmtr(R(\varphi)|\pmA(V;K))=\pmtr(R'(\varphi)|\pmA(V';K))$.
\end{proof}

\subsection{Proof of Proposition~\ref{pm:mult}}\label{pm:ss-mult}

The following property of Green's parametrisation~\cite{Green1955} is the only one used; we
derive it below from the statements of~\cite{Green1955}.
\begin{enumerate}
\item[(G)] If the label $c$ over $\auF$ is an irreducible polynomial, the character $\Xi_c$ of
$\auGL_r(\auF)$ named by $c$ is cuspidal: the sum of $\Xi_c(xuy)$ over $u\in\mathsf U$ is zero for
all $x,y$ and the unipotent radical $\mathsf U$ of every proper parabolic subgroup.
\end{enumerate}
Green does not use the word cuspidal, so we derive (G) from statements of~\cite{Green1955}. For
$c$ irreducible of degree $r$, $\Xi_c$ is Green's primary character $(g^{(1)})$ of the $r$-simplex
$g$ attached to $c$ \cite[Lemma~7.7 and Thm.~12, pp.~438--439]{Green1955}. Let $\pi$ realise
$\Xi_c$. Since $\Xi_c$ is a class function, $\sum_u\Xi_c(xuy)=\sum_u\Xi_c(yxu)=\#\mathsf U\,
\operatorname{tr}(\pi(yx)E)$, with $E$ the projection onto $\pi^{\mathsf U}$; as $\pi$ is
irreducible, the $\pi(h)$ span $\operatorname{End}(\pi)$, so (G) says $\pi^{\mathsf U}=0$. A larger
parabolic has a smaller unipotent radical and every parabolic is conjugate to a standard one, so it
suffices to take $\mathsf U$ the radical of the standard parabolic with Levi
$\auGL_{r_1}\times\auGL_{r_2}$, $r_1+r_2=r$, $r_i\ge1$. If $\pi^{\mathsf U}\ne0$, it contains an
irreducible representation $\sigma_1\otimes\sigma_2$ of the Levi, and by Frobenius reciprocity
$\Xi_c$ occurs in the character induced from the parabolic, which is Green's $\circ$-product
$\chi_1\circ\chi_2$ of the characters $\chi_i$ of $\sigma_i$ \cite[p.~403]{Green1955}. By
\cite[Thm.~13, p.~439]{Green1955}, $\chi_i=\prod_{g'}(g'^{\nu_i(g')})$ ($\circ$-product) with $\nu_i$
supported on simplexes of degree $\le r_i<r$. The $\circ$-product is commutative and associative,
and for each simplex $g'$ of degree $s$ the map $P_e\mapsto Z_e$ is a ring homomorphism from the
symmetric functions into Green's character algebra sending the Schur function $\{\lambda\}$ to
$\pm(g'^{\lambda})$ \cite[Lemma~7.2, p.~436, Def.~7.3, p.~437, and p.~439]{Green1955}. Since a
product of Schur functions of degrees $a,b$ is a linear combination of Schur functions of degree
$a+b$, $(g'^{\nu_1(g')})\circ(g'^{\nu_2(g')})$ is a combination of the $(g'^{\lambda})$ with
$|\lambda|=|\nu_1(g')|+|\nu_2(g')|$. Hence $\chi_1\circ\chi_2$ is a linear combination of the
characters $(\cdots g'^{\nu(g')}\cdots)$ of Theorem~13 with $\nu$ supported on simplexes of degree
$<r$. These are irreducible and distinct from $\Xi_c=(g^{(1)})$, whose function is supported on
$g$ of degree $r$ \cite[Thm.~13 and p.~440]{Green1955}; by orthogonality of irreducible characters
$\langle\Xi_c,\chi_1\circ\chi_2\rangle=0$, a contradiction. So $\pi^{\mathsf U}=0$, which is (G).

\begin{lemma}\label{pm:supercusp}
Let $v$ be a finite split place of $F$ with residue field $\auF$, and $\Xi$ a cuspidal class
function on $\auGL_n(\auF)$ in the sense of \textup{(G)}. The function $\beta$ on $H_v$ that is
$\Xi(\bar k)$ at $k\in K_v=\auGL_n(\auO_v)$, $\bar k$ the reduction of $k$, and $0$ outside $K_v$, is
a supercusp function.
\end{lemma}
\begin{proof}
Every parabolic subgroup is conjugate under $K_v$ to a standard one, $P=MN$, and $H_v=K_vP=PK_v$.
For $x=kp$, $y=p'k'$ we have $pnp'=\mu n'$ with $\mu\in M$ depending only on $p,p'$, and $n\mapsto
n'$ multiplies the Haar measure of $N$ by a constant. Now $k\mu n'k'\in K_v$ if and only if
$\mu\in M\cap K_v$ and $n'\in N\cap K_v$. So $\int_N\beta(xny)\,dn$ is $0$ if $\mu\notin K_v$, and
otherwise a multiple of $\int_{N\cap K_v}\Xi(\bar k\bar\mu\bar n'\bar k')\,dn'$, a multiple of the
sum of $\Xi$ over a coset of the unipotent radical of a proper parabolic subgroup of
$\auGL_n(\auF)$.
\end{proof}

\begin{proof}[Proof of Proposition~\ref{pm:mult}]
Let $V=Y_0$ and $V'=W_j$, with the sign vectors \eqref{pm:eq-signs} and $\lambda_\tau=\varsigma_\tau$;
\eqref{pm:eq-gap} holds by hypothesis. (For the data: two entries on coordinates of different
signs are either two entries $-\kappa(\tau)-a\Delta$ with different $a$, or one of these and a
small entry; in both cases they differ by at least $\Delta$.) Both spaces carry a lattice
that is unimodular at every finite place, and $E/F$ is unramified at all finite places; hence
at every finite place $w$ of $F$ the two lattices are isometric (for $w$ inert this is the
uniqueness of unimodular hermitian lattices over an unramified extension, i.e.\ the existence of an orthonormal basis (proof: the reduction of the lattice is a non-degenerate hermitian space over the quadratic extension of residue fields, so it has a vector $\bar x$ with $\bar h(\bar x,\bar x)\ne0$ (otherwise polarisation and the surjectivity of the residue trace give $\bar h=0$); a lift $x$ has $h(x,x)$ a unit of the base ring, hence the norm of a unit since the extension is unramified, so after scaling $h(x,x)=1$; then the lattice is $\mathcal O x\perp(x^\perp\cap\text{lattice})$ with unimodular second summand, and one inducts on the rank; see also~\cite{Jacobowitz}); for
$w$ split it is trivial). Choosing such isometries gives a comparison datum $(Y_0,W_j,\iota)$
for which $\iota$ identifies the stabilisers $K^\circ$ of the lattices, the levels $J_M$, the
groups $\auG=K^\circ/J_M=\prod_w\auGL_r(\auF_w)$ and their actions $\varrho$, and the algebra
$\mathcal H^M$. Each chosen place $w$ lies above a place of $F$ above $p_0$ or $p_1$ that splits
in $E$, and distinct chosen places lie above distinct places of $F$.

Let $(c_w)$ be a good family of labels and $T=\varphi_S\otimes1_{J_M^{S}}\in\mathcal H^M$. For
$s\in\auG$ the idempotent $\bigotimes_we_{c_w}$ has the coefficient $\prod_w\Xi_{c_w}(1)\Xi_{c_w}
(s_w^{-1})/\#\auGL_r(\auF_w)$. Hence $\pmvol(J_M)^{-1}R(T)\circ\varrho(\bigotimes_we_{c_w})$ is
a constant multiple, the same for $Y_0$ and $W_j$, of $R(\varphi)$ on $\pmA(V;J_M)$, where
\[
\varphi=\varphi_S\otimes\bigotimes_{w\text{ chosen}}\beta_w\otimes1_{K^{\circ}}^{\,\text{other places}},
\qquad\beta_w(k)=\Xi_{c_w}(\bar k^{-1})\ \text{on }\auGL_r(\auO_w),
\]
a test function bi-invariant under $J_M$. As the family is good, some $c_{w_0}$ is irreducible;
by (G) and Lemma~\ref{pm:supercusp}, $\beta_{w_0}$ is a supercusp function (if $\Xi$ is cuspidal,
so is $g\mapsto\Xi(g^{-1})$). Theorem~\ref{pm:main} gives
$\pmtr(T\varrho(\bigotimes e_{c_w})\mid\pmA(Y_0;J_M))=\pmtr(T\varrho(\bigotimes e_{c_w})\mid
\pmA(W_j;J_M))$. Summing over the good families, and using that $\varrho(e_0)$ is a projection
commuting with $T$, we get $\pmtr(T\mid\auT)=\pmtr(T\mid e_0\pmA(W_j))$.

For $T=1$ this is the equality of dimensions. $\mathcal H^M$ is commutative and stable under
adjoints, so both spaces are direct sums of eigenspaces $\mathsf X_\chi$ for characters $\chi$ of
$\mathcal H^M$, and $\pmtr(T|\mathsf X)=\sum_\chi\chi(T)\dim\mathsf X_\chi$. Distinct characters
are linearly independent, so the multiplicities are equal.
\end{proof}

\begin{remark}[The full spherical algebra]\label{pm:rem-full}
The trace identity holds for every $T=\varphi_S\otimes1_{J_M^S}$ with $S$ any finite set of finite places
of $F$ not below a chosen place, split or inert, and $\varphi_S$ bi-invariant under $K^\circ_S$. Indeed the
comparison datum of the proof identifies, through $\iota$, the stabilisers $K^\circ_w$ at \emph{every}
finite place (for inert $w$ by~\cite{Jacobowitz}), hence the bi-$K^\circ_S$-invariant functions on the two
groups; and Theorem~\ref{pm:main} allows an arbitrary function $\varphi^{v_0}$ away from $v_0$. The first two
paragraphs of the proof therefore apply verbatim. (The last paragraph, on multiplicities, uses only the
split places, where $\mathcal H^M$ is commutative.) This is the form of Proposition~\ref{au:mult}.
\end{remark}


%% ===== proofs/B-orbital.tex =====
% proofs/B-orbital.tex
% Proof of Theorem pm:arch(b) of proofs/B-multiplicities.tex: the stable sum of the elliptic
% orbital integrals of the lowest-K-type coefficient of a discrete series of U(p,q).
% Labels of this file start with pbr:. Only amsmath/amssymb/amsthm are used.
\providecommand{\auU}{\mathrm U}
\providecommand{\aubar}[1]{\overline{#1}}
\providecommand{\pcom}{\mathsf m}
\providecommand{\pcoD}{\mathsf D}
\providecommand{\pcoH}{\mathsf H}
% macros of this file (prefix \pbr)
\providecommand{\pbrB}{\mathsf B}
\providecommand{\pbrN}{\mathsf N}
\providecommand{\pbrn}{\mathfrak n}
\providecommand{\pbrK}{K}
\providecommand{\pbrW}{W_K}
\providecommand{\pbrlev}{\ell}
\providecommand{\pbrsgn}{\operatorname{sgn}}
\providecommand{\pbrAlt}{\mathsf A}
\providecommand{\pbrTp}{T^{+}}
\providecommand{\pbrTh}{\Theta}
\providecommand{\pbrE}{\mathcal E}
\providecommand{\pbrdc}{\delta_{\mathrm c}}
\providecommand{\pbrdD}{\delta_{D}}
\providecommand{\pbrtop}{\tau}
\providecommand{\pbrLo}{\Lambda^{\circ}}
\providecommand{\pbrNo}{N^{\circ}}
\providecommand{\pbrO}{\mathcal O}
\providecommand{\pbrre}{\operatorname{Re}}

\section{Elliptic orbital integrals of the discrete series coefficient}\label{pbr:sec}

This section proves part (b) of Theorem~\ref{pm:arch}. Part (a) is Remark~\ref{pm:rem-other}(1).
The result is Harish-Chandra's formula for the elliptic orbital integrals of a discrete series
coefficient, summed over a stable class; we prove it for $\auU(p,q)$, for parameters satisfying the
gap condition \eqref{pm:eq-gap}, and for the coefficient of the lowest $K$-type, from the explicit
formula of Theorem~\ref{pco:FJ}. \emph{Dependencies:} the proofs use the published results listed
in \S\ref{pbr:ss-inputs} (Schur orthogonality and admissibility, already used in
\S\ref{pm:sec}; the two statements of Schmid's theorem on the $K$-types of a discrete series, of
which the first is (D3)(b) of \S\ref{par:ss-rigid}; and the classical theory of compact connected Lie
groups). Harish-Chandra's character theory is not used, and neither is his formula for the formal
degree.

\subsection{Statement, conventions and inputs}\label{pbr:ss-statement}

\subsubsection*{The group and the parameter}
Let $n\ge2$, $s=(\varepsilon_1,\dots,\varepsilon_n)\in\{\pm1\}^n$ with $p$ entries $+1$ and $q$
entries $-1$, $J=I_s=\operatorname{diag}(s)$, $G=G_s=\{g:g^*Jg=J\}$, $K=G\cap\auU(n)$ (the block
diagonal elements of $G$), and $T\subset K$ the diagonal torus. We write $x\in T$ as
$x=\operatorname{diag}(x_1,\dots,x_n)$, and for $\sigma\in S_n$, $x^\sigma=\operatorname{diag}(
x_{\sigma(1)},\dots,x_{\sigma(n)})$, as in \S\ref{pm:ss-setting}. Let $\lambda\in(\frac{n-1}2+\Z)^n$
have distinct entries, and let
\[
h_\lambda=\min\{|\lambda_a-\lambda_b|:\varepsilon_a\ne\varepsilon_b\}
\]
be its gap ($h_\lambda=\infty$ if $pq=0$). Throughout, $(\lambda,s)$ satisfies \eqref{pm:eq-gap}:
\begin{equation}\label{pbr:eq-gap}
h_\lambda\ \ge\ h_0:=500\,n^3+10\,n^5 .
\end{equation}
As in \S\ref{pco:ss-stat}: the positive roots are the $e_a-e_b$ with $\lambda_a>\lambda_b$, compact
if $\varepsilon_a=\varepsilon_b$ and noncompact otherwise; $\rho,\rho_{\mathrm c},\rho_{\mathrm n}$
are the half sums; $\Lambda=\lambda+\rho_{\mathrm n}-\rho_{\mathrm c}$ and $\mu=\lambda+\rho$, both in
$\Z^n$. A weight $\nu\in\Z^n$ is \emph{$K$-dominant} if $\nu_a\ge\nu_b$ whenever $\varepsilon_a=
\varepsilon_b$ and $\lambda_a>\lambda_b$; the irreducible representations of $K\cong\auU(p)\times
\auU(q)$ are the $(\sigma_\nu,V_\nu)$ with highest weight $\nu$ for the positive compact roots,
$\nu$ $K$-dominant, with characters $\chi^K_\nu$. By Lemma~\ref{pco:minors}(c), $\Lambda$ is
$K$-dominant. We also put
\[
\pbrLo=\lambda-\rho ,
\]
the highest weight, for the order of $\lambda$, of the representation of $\auU(n)$ with character
$\chi_\lambda$ (\S\ref{pm:ss-setting}); $\pbrLo$ is dominant: $\pbrLo_a-\pbrLo_b=(\lambda_a-
\lambda_b)-(\rho_a-\rho_b)\ge0$ for $\lambda_a>\lambda_b$, because $\rho_a-\rho_b-1$ is the number of
entries strictly between $\lambda_b$ and $\lambda_a$, which is at most $\lambda_a-\lambda_b-1$.

\subsubsection*{The representation, the coefficient and the measure}
Inner products are linear in the \emph{first} variable. Let $(\pi,\mathsf H)=\pi(\lambda,s)$ be the
discrete series of $G$ with Harish-Chandra parameter $\lambda$ and $v$ a unit highest weight vector
of its lowest $K$-type $\Lambda$ (unique up to a scalar of modulus one by (D3)); put
\begin{equation}\label{pbr:eq-c}
c(g)=\langle\pi(g)v,v\rangle .
\end{equation}
In the convention of \S\ref{pco:sec} (inner products conjugate-linear in the first variable) this is
the function $c_\lambda(g)=\langle v,\pi_\lambda(g)v\rangle$, so by Theorem~\ref{pco:FJ} (its
hypothesis $h_\lambda\ge n^2+1$ holds) $c=F$ is given by \eqref{pco:eq-F}, and by
Lemma~\ref{pco:decay}
\begin{equation}\label{pbr:eq-decay}
|c(g)|\le(2/\|g\|)^{j}\qquad(g\in G,\ j\in\Z,\ 0\le j\le h_\lambda).
\end{equation}
Let $dg$ be the Haar measure on $G$ compatible with the Haar measure of mass $1$ on $\auU(n)$
(\S\ref{pds:ss-ds1}(iii), Remark~\ref{pds:cayley}), and $d=d_\pi$ the formal degree of $\pi$ for
$dg$: $\int_G|c|^2dg=d^{-1}$. For another Haar measure $c'dg$ the formal degree is $d/c'$, so the
measure $d\,dg$ and every statement below are independent of this choice. Compact groups carry
their Haar measure of mass $1$.

\begin{theorem}\label{pbr:main}
Assume \eqref{pbr:eq-gap}. For every $t\in T$ with distinct eigenvalues the function
$g\mapsto c(g^{-1}tg)$ is integrable on $G$, and
\begin{equation}\label{pbr:eq-unstable}
d\int_Gc(g^{-1}tg)\,dg=(-1)^{pq}\,\frac{\sum_{w\in S_n,\ s\circ w=s}\pbrsgn(w)\,t^{\kappa\circ w}}
{\prod_{\lambda_a>\lambda_b}(t_a-t_b)},\qquad\kappa=\lambda+\tfrac{n-1}2\mathbf 1,
\end{equation}
where $(\kappa\circ w)_a=\kappa_{w(a)}$ and $t^\gamma=\prod_at_a^{\gamma_a}$. Consequently
\begin{equation}\label{pbr:eq-stable}
\sum_{\sigma\in S_n}d\int_G\aubar{c(g^{-1}t^\sigma g)}\,dg=p!\,q!\,(-1)^{pq}\,\aubar{\chi_\lambda(t)} ,
\end{equation}
which is Theorem~\ref{pm:arch}(b).
\end{theorem}

Formula \eqref{pbr:eq-unstable} is Harish-Chandra's formula for the character of $\pi$ on $T$
(\cite[Theorem~16, p.~96]{HC66}; compare Remark~\ref{pm:rem-arch}(1)); we do not use it, we prove it.

\begin{remark}[The convention]\label{pbr:rem-conv}
For $q=0$ ($G=\auU(n)$, $\pi$ the representation with character $\chi_\lambda$, $d=\dim\pi$),
$\int_{\auU(n)}\pi(g^{-1}tg)\,dg=d^{-1}\chi_\lambda(t)$ by Schur's lemma, so $d\int c(g^{-1}tg)dg=
\chi_\lambda(t)$ and \eqref{pbr:eq-stable} reads $n!\,\aubar{\chi_\lambda(t)}$: the complex conjugate
in Theorem~\ref{pm:arch}(b) is correct for the convention \eqref{pbr:eq-c}, linear in the first
variable, in which $e=d\,\aubar c$ is the kernel of the projection onto the $\pi$-isotypic part.
With the opposite convention one has to drop both conjugations in \eqref{pbr:eq-stable}. Only the
fact that the right side of \eqref{pbr:eq-stable} does not depend on $s$ beyond the sign
$(-1)^{pq}$ is used in Proposition~\ref{pm:ell}, so either reading suffices there; the convention of
\S\ref{pm:sec} is stated in Theorem~\ref{pm:arch} and Remark~\ref{pm:rem-other}(1).
\end{remark}

\subsubsection*{Published inputs}\label{pbr:ss-inputs}
\begin{enumerate}
\item[(R1)] \emph{Square-integrable representations.} Schur orthogonality $\int_G\langle\pi(g)a,b
\rangle\aubar{\langle\pi(g)a',b'\rangle}dg=d^{-1}\langle a,a'\rangle\aubar{\langle b,b'\rangle}$
(\cite[Proposition~9.6, p.~284]{Knapp}); admissibility: each $K$-isotypic subspace $\mathsf H_\nu$ of $\pi$ is
finite-dimensional (\cite[Theorem~8.1, p.~205]{Knapp}). These are used in \S\ref{pm:sec} already.
\item[(R2)] \emph{Compact groups} (\cite{KnappBeyond}: Schur orthogonality, Corollary~4.10; orthonormality of characters, Corollary~4.16; Peter--Weyl, Theorem~4.20 and Corollary~4.21; the remaining classical statements below without locator). For a
unitary representation $\sigma$ of $K$, $(\dim V_\nu)\int_K\aubar{\chi^K_\nu(k)}\sigma(k)\,dk$ is the
orthogonal projection onto the $\nu$-isotypic subspace. Peter--Weyl: the $\chi^K_\nu$ form an
orthonormal basis of the space of class functions in $L^2(K)$. Every element of $K$ is
$K$-conjugate to an element of $T$, and a class function is determined by its restriction to $T$,
which is invariant under $\pbrW=\{w\in S_n:s\circ w=s\}\cong S_p\times S_q$. Weyl's integration
formula: $\int_K\phi\,dk=|\pbrW|^{-1}\int_T\phi(x)|V_{\mathrm c}(x)|^2dx$ for continuous class
functions $\phi$, and Weyl's character formula $V_{\mathrm c}\chi^K_\nu=\pbrAlt_{\pbrW}(\nu+\pbrdc)$
on $T$, with the notation of \S\ref{pbr:ss-laurent}; the same two statements for $\auU(n)$, where the
character with highest weight $\pbrLo$ is $\chi_\lambda$ and $V\chi_\lambda=\pbrAlt_{S_n}(\kappa)$.
Schur orthogonality for $\auU(n)$, and the holomorphic extension of its representations to
$\mathrm{GL}_n(\C)$ (Weyl's unitary trick).
\item[(R3)] \emph{Schmid's theorem} (\cite[Lemma~(2.22), Lemma~(3.8), Corollary~(3.9), Corollary~(4.19) and the proof of Theorem~(1.3) on p.~560]{Schmid75}: $\pi$ is the irreducible quotient of a module whose $K$-types satisfy both assertions; valid for every discrete series, see the paragraph ``On (D3)'' in \S\ref{par:ss-rigid}; see also
\cite[Theorem~9.20(b),(c), p.~310]{Knapp}, and \cite[Proposition~9.6, Theorem~8.1]{Knapp} in (R1): these are stated for linear connected semisimple groups and apply to
$\auU(p,q)=Z\cdot\mathrm{SU}(p,q)$, $Z\cong\auU(1)$ the compact centre, $Z\cap\mathrm{SU}(p,q)$ finite: an
irreducible unitary (square-integrable) representation of $\auU(p,q)$ restricts irreducibly to one of
$\mathrm{SU}(p,q)$, $Z$ acting by a character, $K=Z\cdot(K\cap\mathrm{SU}(p,q))$, so $K$-isotypic
subspaces, $K$-types with their multiplicities and the lowest $K$-type correspond, and formal degrees
agree up to the normalisation of Haar measures; see the paragraph ``On (D3)'' in
\S\ref{par:ss-rigid}). The $K$-type $\Lambda$ occurs in $\pi$ with
multiplicity one, and every $K$-type $\nu$ of $\pi$ has the form $\nu=\Lambda+\sum_\beta n_\beta
\beta$ with integers $n_\beta\ge0$, the sum over the positive noncompact roots $\beta$. The first
half is (D3)(b) of \S\ref{par:ss-rigid}.
\item[(R4)] \emph{Linear algebra.} Fischer's inequality $\Delta_{S\cup S'}(X)\le\Delta_S(X)\Delta_{S'}
(X)$ and Jacobi's identity $\Delta_S(X^{-1})=\Delta_{S^c}(X)/\det X$ (as in Lemma~\ref{pds:fischer});
for positive semidefinite $X\le Y$ (that is, $Y-X\ge0$), $\Delta_S(X)\le\Delta_S(Y)$
\cite[Corollary~7.7.4(e)]{HornJohnson}, applied to the principal submatrices on $S$; $|\Delta_S(X)|\le\|X\|^{|S|}$.
\end{enumerate}
Real-analytic functions on a connected real-analytic manifold that are not identically zero have
null zero sets; holomorphic functions on a connected open set agreeing on a set with an
accumulation point agree. We use the minors $\Delta_S$ and the functions $\pcom_w$ of
\S\ref{pco:ss-stat}, and Lemma~\ref{pco:minors}.

\emph{Plan.} \S\ref{pbr:ss-int} proves the integrability, uniformly on $K$ after multiplication by a
power of the noncompact Weyl denominator $D$. \S\ref{pbr:ss-mult} shows, from (R1)--(R3), that
$V_{\mathrm c}D^{M}\cdot d\int_Gc(g^{-1}xg)dg$ is a Laurent polynomial $P$ whose leading part is
controlled. \S\S\ref{pbr:ss-flag}--\ref{pbr:ss-open} compute the orbital integral on an open part of
$T$ as a Weyl character of $\auU(n)$ minus an error $\pbrE$, using the formula \eqref{pco:eq-F} and
an integral formula for densities on the flag variety. \S\ref{pbr:ss-error} shows that $\pbrE$ does
not affect the monomials of low level, and \S\ref{pbr:ss-proof} identifies $P$ by comparing those
monomials.

\subsection{Laurent polynomials, levels and alternants}\label{pbr:ss-laurent}

A \emph{Laurent polynomial} is $P=\sum_{\gamma\in\Z^n}P_\gamma x^\gamma$, finitely many $P_\gamma\ne0$,
$x^\gamma=\prod_ax_a^{\gamma_a}$; its \emph{exponents} are the $\gamma$ with $P_\gamma\ne0$. It is
determined by its values on $T$. A \emph{formal series} is any function $\gamma\mapsto t_\gamma$ on
$\Z^n$, and for a Laurent polynomial $Q$ we put $(Qt)_\gamma=\sum_\beta Q_\beta t_{\gamma-\beta}$.
For $w\in S_n$ and $\gamma\in\Z^n$ let $(\gamma\circ w)_a=\gamma_{w(a)}$; then
$P(x^w)=\sum_\beta P_{\beta\circ w}x^\beta$. We call $P$ (or a formal series) \emph{$\pbrW$-symmetric}
if $P_{\gamma\circ w}=P_\gamma$, and \emph{$\pbrW$-antisymmetric} if $P_{\gamma\circ w}=\pbrsgn(w)
P_\gamma$, for all $w\in\pbrW$ and $\gamma$; on $T$ this means $P(x^w)=P(x)$, resp.\ $\pbrsgn(w)P(x)$.
For a subgroup $W'\subseteq S_n$ the \emph{alternant} is $\pbrAlt_{W'}(\gamma)=\sum_{w\in W'}\pbrsgn(w)
x^{\gamma\circ w}$; if $\gamma$ has distinct entries, its coefficient at $\gamma\circ w$ is
$\pbrsgn(w)$ for $w\in W'$, and it has no other exponents.

\subsubsection*{Notation}
For $a\in\{1,\dots,n\}$ put
\begin{gather*}
o_a=\#\{b:\lambda_b>\lambda_a\},\qquad \pbrdc(a)=\#\{b:\varepsilon_b=\varepsilon_a,\ \lambda_b<
\lambda_a\},\qquad \pbrdD(a)=\#\{b:\varepsilon_b\ne\varepsilon_a,\ \lambda_b>\lambda_a\},\\
\pbrlev(\gamma)=\sum_ao_a\gamma_a\quad(\text{the \emph{level} of }\gamma\in\Z^n),\qquad
N=4n^2,\qquad M=N+1,
\end{gather*}
and recall $\kappa=\lambda+\frac{n-1}2\mathbf 1\in\Z^n$. The vector $o$ strictly reverses the order
of $\lambda$: $\lambda_a>\lambda_b$ implies $o_a<o_b$; and $0\le o_a\le n-1$. Let
\[
V=\prod_{\lambda_a>\lambda_b}(x_a-x_b),\qquad V_{\mathrm c}=\prod_{\substack{\lambda_a>\lambda_b\\
\varepsilon_a=\varepsilon_b}}(x_a-x_b),\qquad V_{\mathrm n}=\prod_{\substack{\lambda_a>\lambda_b\\
\varepsilon_a\ne\varepsilon_b}}(x_a-x_b),\qquad D=\prod_{\substack{\lambda_a<\lambda_b\\
\varepsilon_a\ne\varepsilon_b}}(x_a-x_b).
\]

\begin{lemma}\label{pbr:basic}
\begin{enumerate}
\item[(a)] $V=\pbrAlt_{S_n}(\delta)$ with $\delta_a=\#\{b:\lambda_b<\lambda_a\}$, and $V_{\mathrm c}=
\pbrAlt_{\pbrW}(\pbrdc)$; $V$ is $S_n$-antisymmetric and $V_{\mathrm c}$ is $\pbrW$-antisymmetric.
\item[(b)] $V=V_{\mathrm c}V_{\mathrm n}$ and $D=(-1)^{pq}V_{\mathrm n}$. There is $\zeta\in\{\pm1\}$
with $D=\zeta\prod_{\varepsilon_a=1,\ \varepsilon_b=-1}(x_a-x_b)$; so $D$ is $\pbrW$-symmetric, and
$D(x)\ne0$ if the $x_a$ are distinct.
\item[(c)] For $k\ge0$, $D^k$ has coefficient $1$ at $k\pbrdD$, and every other exponent $\beta$ of
$D^k$ has $\pbrlev(\beta)<\pbrlev(k\pbrdD)$. The exponents of $V$ and of $D^k$ have nonnegative
entries, hence nonnegative level.
\item[(d)] $\Lambda+\pbrdc+\pbrdD=\kappa$ and $\kappa-\mu=o$.
\item[(e)] \textup{(Rearrangement with a gain.)} Let $b\in\Z^n$ be such that $b_j<b_i$ implies
$o_i<o_j$. Then $\pbrlev(b)+|b_i-b_{w(i)}|\le\pbrlev(b\circ w)$ for every $w\in S_n$ and every $i$.
\end{enumerate}
\end{lemma}
\begin{proof}
(a) List the indices as $a_1,\dots,a_n$ with $\lambda$ decreasing; Vandermonde's determinant gives
$\prod_{i<j}(x_{a_i}-x_{a_j})=\det(x_{a_i}^{n-j})_{i,j}$, which is $\pbrAlt_{S_n}(\delta)$ since
$\delta(a_i)=n-i$. The same in each of the two blocks of equal sign gives $V_{\mathrm c}$. The
antisymmetry is that of determinants under permutation of rows.
(b) The first identity sorts the pairs. $D$ is $V_{\mathrm n}$ with each of its $pq$ factors
$x_b-x_a$ ($\lambda_b>\lambda_a$) replaced by $x_a-x_b$. In $D$ the factor attached to a pair
$\varepsilon_a=1$, $\varepsilon_b=-1$ is $\pm(x_a-x_b)$, with a sign that depends only on the pair; the
product of the $x_a-x_b$ over these pairs is invariant under $\pbrW$, which permutes the indices of
each sign among themselves.
(c) Each factor $x_a-x_b$ of $D$ has $\lambda_a<\lambda_b$, so $o_a>o_b$: the monomial $x_a$ has the
larger level, and $\pbrdD$ is the sum of these $e_a$. Expanding $D^k$ as a product of $kpq$ factors,
a monomial is obtained by choosing $x_a$ or $-x_b$ in each factor; its level is at most
$\pbrlev(k\pbrdD)$, with equality only if $x_a$ is chosen every time, which happens for exactly one
choice, with coefficient $1$. The last sentence is clear since $o_a\ge0$.
(d) Write $\mathrm{sb},\mathrm{sa},\mathrm{ob},\mathrm{oa}$ for the numbers of $b\ne a$ with
$\varepsilon_b=\varepsilon_a$ and $\lambda_b<\lambda_a$, resp.\ $\varepsilon_b=\varepsilon_a$ and
$\lambda_b>\lambda_a$, resp.\ $\varepsilon_b\ne\varepsilon_a$, $\lambda_b<\lambda_a$, resp.\
$\varepsilon_b\ne\varepsilon_a$, $\lambda_b>\lambda_a$. Then $2\rho_{\mathrm c,a}=\mathrm{sb}-
\mathrm{sa}$, $2\rho_{\mathrm n,a}=\mathrm{ob}-\mathrm{oa}$, $\pbrdc(a)=\mathrm{sb}$, $\pbrdD(a)=
\mathrm{oa}$, so $\Lambda_a+\pbrdc(a)+\pbrdD(a)-\lambda_a=\frac12(\mathrm{sb}+\mathrm{sa}+\mathrm{ob}
+\mathrm{oa})=\frac{n-1}2$. Similarly $\kappa_a-\mu_a=\frac{n-1}2-\rho_a=\frac12(n-1-\mathrm{sb}-
\mathrm{ob}+\mathrm{sa}+\mathrm{oa})=\mathrm{sa}+\mathrm{oa}=o_a$.
(e) If $o'\in\Z^n$ satisfies $(o'_i-o'_j)(b_i-b_j)\le0$ for all $i,j$, then $\sum_ao'_ab_a\le\sum_a
o'_ab_{w(a)}$ for every $w$ (rearrangement inequality). Since $o$ is integral and strictly reverses
the order of $b$, so does weakly $o'=o+\eta e_i$ for $\eta=\pm1$. This gives $\pbrlev(b)+\eta b_i\le
\pbrlev(b\circ w)+\eta b_{w(i)}$ for both signs, which is the claim.
\end{proof}

\begin{lemma}\label{pbr:top}
Let $k\ge0$, $a\in\Z^n$, and let $t$ be a formal series with $t_a=1$ such that every $\gamma$ with
$t_\gamma\ne0$ is $\gamma=\gamma'\circ w$ with $w\in\pbrW$ and $\pbrlev(\gamma')\le\pbrlev(a)$. Put
$b=a+k\pbrdD$ and assume that $\lambda_j<\lambda_i$ implies $b_j<b_i$. Then $(D^kt)_b=1$, and for
every $\gamma$ with $(D^kt)_\gamma\ne0$ there is $w\in\pbrW$ with $\pbrlev(\gamma\circ w)\le
\pbrlev(b)$.
\end{lemma}
\begin{proof}
Let $(D^k)_\beta t_{\gamma-\beta}\ne0$, and write $\gamma-\beta=\gamma'\circ w$ as in the hypothesis.
Then $\pbrlev(\gamma\circ w^{-1})=\pbrlev(\beta\circ w^{-1})+\pbrlev(\gamma')\le\pbrlev(\beta\circ
w^{-1})+\pbrlev(a)$. By Lemma~\ref{pbr:basic}(b), $(D^k)_{\beta\circ w^{-1}}=(D^k)_\beta\ne0$, so by
(c) $\pbrlev(\beta\circ w^{-1})\le\pbrlev(k\pbrdD)$, with equality only if $\beta\circ w^{-1}=
k\pbrdD$. Hence
\begin{equation}\label{pbr:eq-toplevel}
\pbrlev(\gamma\circ w^{-1})+[\beta\circ w^{-1}\ne k\pbrdD]\le\pbrlev(b),
\end{equation}
which gives the second claim. For $\gamma=b$: $b$ has the order of $\lambda$, so $o$ strictly
reverses it and Lemma~\ref{pbr:basic}(e) gives $\pbrlev(b\circ w^{-1})\ge\pbrlev(b)+|b_i-b_{w^{-1}
(i)}|$ for all $i$. With \eqref{pbr:eq-toplevel}, $b_{w^{-1}(i)}=b_i$ for all $i$, so $w=1$ ($b$ has
distinct entries), and $\beta=k\pbrdD$. Thus the only term of $(D^kt)_b$ is $(D^k)_{k\pbrdD}
t_a=1$.
\end{proof}

\begin{lemma}\label{pbr:tau}
Put $\pbrtop=\kappa+N\pbrdD$. If $\lambda_j<\lambda_i$ then $\pbrtop_j<\pbrtop_i$. Moreover
\[
\pbrlev(o+N\pbrdD)<h_0,\qquad \pbrlev(N\pbrdD)<h_0 .
\]
\end{lemma}
\begin{proof}
$\kappa_i-\kappa_j=\lambda_i-\lambda_j>0$ and $0\le\pbrdD\le n-1$. If $\lambda_i-\lambda_j>Nn$, then
$\pbrtop_i-\pbrtop_j\ge(\lambda_i-\lambda_j)-N(n-1)>0$. Otherwise $\lambda_i-\lambda_j\le Nn=4n^3<h_0$,
so every $a$ with $\lambda_j\le\lambda_a\le\lambda_i$ has $\varepsilon_a=\varepsilon_i$ (else
$|\lambda_a-\lambda_i|\ge h_\lambda\ge h_0$); in particular $\varepsilon_j=\varepsilon_i$, and an index $b$
of the other sign with $\lambda_b>\lambda_j$ has $\lambda_b>\lambda_i$. So $\pbrdD(j)=\pbrdD(i)$ and
$\pbrtop_i-\pbrtop_j=\lambda_i-\lambda_j>0$. For the bounds, $o_a,\pbrdD(a)\le n-1$, so
$\pbrlev(o+N\pbrdD)\le n(n-1)^2(1+N)<n^3+4n^5<h_0$, and $\pbrlev(N\pbrdD)\le Nn(n-1)^2<h_0$.
\end{proof}

\subsection{Integrability}\label{pbr:ss-int}

Write $+$ and $-$ for the sets of indices of sign $+1$ and $-1$, $X_{++}$, $X_{--}$ for the
corresponding blocks of a matrix $X$, and $\Delta_+=\Delta_{+}$ for the principal minor on $+$.

\begin{lemma}\label{pbr:minor}
Let $u\in G$, $x\in T$ and $g=u^{-1}xu$. Then
$(\|u\|/2)^2\,\bigl|\prod_{a\in+,\,b\in-}(x_a-x_b)\bigr|\le(2\|g\|)^{pq}$.
\end{lemma}
\begin{proof}
If $pq=0$, $G$ is compact, $\|u\|=1$ and the claim is $\frac14\le1$. Let $\Pi=\prod_{b\in-}(x-x_b)$,
a diagonal matrix with $\Pi_{aa}=\prod_{b\in-}(x_a-x_b)$, which is $0$ for $a\in-$. Then
$u^{-1}\Pi u=\prod_{b\in-}(g-x_b)$, so $\|u^{-1}\Pi u\|\le(\|g\|+1)^q\le(2\|g\|)^q$ (as $\|g\|\ge1$),
and $|\Delta_+(u^{-1}\Pi u)|\le(2\|g\|)^{pq}$ by (R4). Since $\Pi$ vanishes off the block $++$,
$(u^{-1}\Pi u)_{++}=(u^{-1})_{++}\Pi_{++}u_{++}$, and $(u^{-1})_{++}=(Ju^*J)_{++}=(u_{++})^*$; hence
$|\Delta_+(u^{-1}\Pi u)|=|\det u_{++}|^2\prod_{a\in+}|\Pi_{aa}|$. By Jacobi's identity
$\Delta_+(u^{-1})=\det u_{--}/\det u$ and $|\det u|=1$, so $|\det u_{++}|=|\det u_{--}|$ and
$|\det u_{++}|^2=|\det u_{++}\det u_{--}|=|\det a(u)|\ge(\|u\|/2)^2$ by Lemma~\ref{pds:blockdet}.
\end{proof}

For $k\in K$ write $k=\operatorname{diag}(k_+,k_-)$ in the blocks, and put $\widetilde D(k)=
\zeta(-1)^{pq}\det\chi_{k_+}(k_-)$, where $\chi_{k_+}(z)=\det(z-k_+)$ and $\zeta$ is as in
Lemma~\ref{pbr:basic}(b). It is a continuous class function on $K$, and $\widetilde D(x)=
\zeta\prod_{b\in-}\prod_{a\in+}(x_a-x_b)=D(x)$ for $x\in T$.

\begin{proposition}\label{pbr:int}
The map $\Psi\colon G\to C(K)$, $\Psi(g)(k)=\widetilde D(k)^Mc(g^{-1}kg)$, is continuous, and
$\int_G\sup_K|\Psi(g)|\,dg<\infty$; so $\Psi$ is Bochner integrable in the Banach space $C(K)$ with the
supremum norm. For $x\in T$ with distinct entries, $g\mapsto c(g^{-1}xg)$ is integrable on $G$.
\end{proposition}
\begin{proof}
Continuity: $(g,k)\mapsto c(g^{-1}kg)$ is continuous and $K$ is compact. Let $k=yxy^{-1}$ with
$y\in K$, $x\in T$ (R2), and $u=y^{-1}g$, so $c(g^{-1}kg)=c(u^{-1}xu)$ and $\|u\|=\|g\|$. By
Lemma~\ref{pbr:minor} and \eqref{pbr:eq-decay} with $j=pqM$ (note $pqM\le\frac{n^2}4(4n^2+1)\le
h_0\le h_\lambda$),
\[
|\widetilde D(k)|^M|c(u^{-1}xu)|\le\bigl(2\|u^{-1}xu\|\bigr)^{pqM}(2/\|u\|)^{2M}\bigl(2/\|u^{-1}xu\|
\bigr)^{pqM}=4^{pqM}(2/\|g\|)^{2M}.
\]
Since $2M=8n^2+2>n^2$, $\int_G\|g\|^{-2M}dg<\infty$ by Lemma~\ref{pm:kernel}(c). For the last claim
evaluate at $k=x$, where $\widetilde D(x)=D(x)\ne0$ (Lemma~\ref{pbr:basic}(b)).
\end{proof}

\subsection{Multiplicities of $K$-types and the Laurent polynomial}\label{pbr:ss-mult}

For a $K$-dominant $\nu$ let $m(\nu)$ be the multiplicity of $V_\nu$ in $\pi|_K$, finite by (R1).

\begin{lemma}\label{pbr:mult}
Let $\nu$ be $K$-dominant and $R_\nu(g)=\int_K\aubar{\chi^K_\nu(k)}\,c(g^{-1}kg)\,dk$. Then
$R_\nu\ge0$, $R_\nu\in L^1(G)$ and $d\int_GR_\nu\,dg=m(\nu)$. Moreover $m(\Lambda)=1$, and
$m(\nu)\ne0$ implies $\pbrlev(\nu)\le\pbrlev(\Lambda)$.
\end{lemma}
\begin{proof}
$c(g^{-1}kg)=\langle\pi(k)\pi(g)v,\pi(g)v\rangle$, so by (R2) $R_\nu(g)=(\dim V_\nu)^{-1}\langle
E_\nu\pi(g)v,\pi(g)v\rangle$, with $E_\nu$ the projection onto $\mathsf H_\nu$. Let $e_1,\dots,e_r$ be
an orthonormal basis of $\mathsf H_\nu$, $r=m(\nu)\dim V_\nu$. Then $R_\nu(g)=(\dim V_\nu)^{-1}
\sum_i|\langle\pi(g)v,e_i\rangle|^2$, and by Schur orthogonality (R1) $d\int_G|\langle\pi(g)v,e_i
\rangle|^2dg=1$; so $d\int_GR_\nu=r/\dim V_\nu=m(\nu)$. By (R3), $m(\Lambda)=1$, and if $m(\nu)\ne0$ then
$\nu-\Lambda$ is a sum of positive noncompact roots $e_a-e_b$, $\lambda_a>\lambda_b$, each of level
$o_a-o_b<0$.
\end{proof}

For a continuous function $f$ on $T$ and $\gamma\in\Z^n$ let $\hat f(\gamma)=\int_Tf(x)x^{-\gamma}dx$.
A vector $\gamma\in\Z^n$ is \emph{$\pbrW$-regular} if $\gamma_a\ne\gamma_b$ whenever $a\ne b$ have
the same sign; each $\pbrW$-regular $\gamma$ is $(\nu+\pbrdc)\circ w$ for exactly one $K$-dominant
$\nu$ and one $w\in\pbrW$ (sort the entries in each block in the order of $\lambda$), and the
$\nu+\pbrdc$, $\nu$ $K$-dominant, are exactly the $\pbrW$-regular vectors that are decreasing in the
order of $\lambda$ on each block. Hence the formal series
\[
t=\sum_{\nu\ K\text{-dominant}}m(\nu)\,\pbrAlt_{\pbrW}(\nu+\pbrdc),\qquad t_\gamma=\pbrsgn(w)m(\nu)
\text{ if }\gamma=(\nu+\pbrdc)\circ w,\quad t_\gamma=0\text{ if $\gamma$ is not $\pbrW$-regular},
\]
is well defined and $\pbrW$-antisymmetric.

\begin{proposition}\label{pbr:laurent}
Let $h=d\int_G\Psi(g)\,dg\in C(K)$ (Proposition~\ref{pbr:int}), and for $x\in T$ such that
$g\mapsto c(g^{-1}xg)$ is integrable put $\pbrTh(x)=d\int_Gc(g^{-1}xg)\,dg$.
\begin{enumerate}
\item[(a)] $h$ is a class function on $K$, and $h(x)=D(x)^M\pbrTh(x)$ whenever $\pbrTh(x)$ is
defined.
\item[(b)] There is a $\pbrW$-antisymmetric Laurent polynomial $P$ with $P(x)=V_{\mathrm c}(x)h(x)$
for $x\in T$, and $P=D^Mt$.
\item[(c)] $t_{\Lambda+\pbrdc}=1$, and every $\gamma$ with $t_\gamma\ne0$ is $\gamma'\circ w$ with
$w\in\pbrW$ and $\pbrlev(\gamma')\le\pbrlev(\Lambda+\pbrdc)$.
\end{enumerate}
\end{proposition}
\begin{proof}
(a) Evaluation at $k$ is a continuous functional on $C(K)$, so $h(k)=d\int_G\widetilde D(k)^Mc(g^{-1}
kg)\,dg$; for $y\in K$, the substitution $g\mapsto yg$ gives $h(yky^{-1})=h(k)$. For $k=x$ with
$\pbrTh(x)$ defined, split off the constant $\widetilde D(x)^M=D(x)^M$.

\emph{Step 1: coefficients against characters.} Let $f$ be a continuous class function on $K$ and
$\langle\chi^K_\nu,f\rangle=\int_K\aubar{\chi^K_\nu}f\,dk$. The function $V_{\mathrm c}f$ on $T$ is
$\pbrW$-antisymmetric, so $\widehat{V_{\mathrm c}f}(\gamma\circ w)=\pbrsgn(w)\widehat{V_{\mathrm c}f}
(\gamma)$ for $w\in\pbrW$ (substitute $x\mapsto x^w$ in the integral, using
$x^{\gamma\circ w}=(x^{w^{-1}})^\gamma$); in particular $\widehat{V_{\mathrm c}f}(\gamma)=0$ if
$\gamma$ is not $\pbrW$-regular (use the transposition of two equal entries). By (R2), since
$\aubar{x^\gamma}=x^{-\gamma}$ on $T$,
\[
\langle\chi^K_\nu,f\rangle=\frac1{|\pbrW|}\int_T\aubar{V_{\mathrm c}\chi^K_\nu}\,V_{\mathrm c}f\,dx
=\frac1{|\pbrW|}\sum_{w\in\pbrW}\pbrsgn(w)\,\widehat{V_{\mathrm c}f}\bigl((\nu+\pbrdc)\circ w\bigr)
=\widehat{V_{\mathrm c}f}(\nu+\pbrdc).
\]

\emph{Step 2: the coefficients of $h$ are integers.} Let $\bar D^M$ be the Laurent polynomial with
$(\bar D^M)_\beta=(D^M)_{-\beta}$; as $D$ has integer coefficients, $\aubar{D(x)^M}=\bar D^M(x)$ on
$T$. Fix $\nu$ and let $f_\nu=\chi^K_\nu\,\aubar{\widetilde D}{}^M$, a continuous class function. By
Step 1 and (R2), for $K$-dominant $\nu'$,
\begin{equation}\label{pbr:eq-n}
n_{\nu\nu'}:=\langle\chi^K_{\nu'},f_\nu\rangle=\bigl(\pbrAlt_{\pbrW}(\nu+\pbrdc)\bar D^M\bigr)_{\nu'+
\pbrdc}=\sum_\beta(D^M)_\beta\,\pbrAlt_{\pbrW}(\nu+\pbrdc)_{\nu'+\pbrdc+\beta}\in\Z,
\end{equation}
and $n_{\nu\nu'}\ne0$ for finitely many $\nu'$ only. By Peter--Weyl, $f_\nu=\sum_{\nu'}n_{\nu\nu'}
\chi^K_{\nu'}$ (the difference is a continuous class function orthogonal to all characters), so
$\aubar{\chi^K_\nu}\widetilde D^M=\sum_{\nu'}n_{\nu\nu'}\aubar{\chi^K_{\nu'}}$. Pairing with
$\chi^K_\nu$ is a continuous functional on $C(K)$, hence commutes with the Bochner integral, and by
Lemma~\ref{pbr:mult}
\[
\langle\chi^K_\nu,h\rangle=d\int_G\int_K\aubar{\chi^K_\nu(k)}\widetilde D(k)^Mc(g^{-1}kg)\,dk\,dg=
\sum_{\nu'}n_{\nu\nu'}\,d\int_GR_{\nu'}\,dg=\sum_{\nu'}n_{\nu\nu'}m(\nu')\in\Z .
\]

\emph{Step 3: $h$ is a finite sum of characters.} By Bessel's inequality $\sum_\nu|\langle\chi^K_\nu,h
\rangle|^2\le\|h\|_{L^2}^2<\infty$; integers with a finite sum of squares vanish with finitely many
exceptions. By (a) and Peter--Weyl, $h=\sum_\nu\langle\chi^K_\nu,h\rangle\chi^K_\nu$, a finite sum, and
by Weyl's character formula $P:=\sum_\nu\langle\chi^K_\nu,h\rangle\pbrAlt_{\pbrW}(\nu+\pbrdc)$ is a
Laurent polynomial with $P=V_{\mathrm c}h$ on $T$. It is $\pbrW$-antisymmetric, and $P_\gamma=
\widehat{V_{\mathrm c}h}(\gamma)$.

\emph{Step 4: $P=D^Mt$.} Both sides are $\pbrW$-antisymmetric ($D^M$ is $\pbrW$-symmetric, so
$(D^Mt)_{\gamma\circ w}=\sum_\beta(D^M)_{\beta\circ w}t_{(\gamma-\beta)\circ w}=\pbrsgn(w)(D^Mt)_\gamma$),
so both vanish at the $\gamma$ that are not $\pbrW$-regular, and it suffices to compare them at
$\gamma=\nu+\pbrdc$, $\nu$ $K$-dominant. There $P_\gamma=\langle\chi^K_\nu,h\rangle$ by Steps 1 and 3,
which by Step 2 and \eqref{pbr:eq-n} is $\sum_{\nu'}m(\nu')\sum_{w\in\pbrW}\pbrsgn(w)(D^M)_{(\nu+
\pbrdc)\circ w-(\nu'+\pbrdc)}$. By the symmetry of $D^M$ under $w^{-1}$,
$(D^M)_{(\nu+\pbrdc)\circ w-(\nu'+\pbrdc)}=(D^M)_{(\nu+\pbrdc)-(\nu'+\pbrdc)\circ w^{-1}}$, and
$\sum_{w}\pbrsgn(w)(D^M)_{\gamma-(\nu'+\pbrdc)\circ w^{-1}}=\sum_\beta(D^M)_\beta\,\pbrAlt_{\pbrW}
(\nu'+\pbrdc)_{\gamma-\beta}$. Hence $P_\gamma=\sum_\beta(D^M)_\beta\sum_{\nu'}m(\nu')\pbrAlt_{\pbrW}
(\nu'+\pbrdc)_{\gamma-\beta}=(D^Mt)_\gamma$; the sums are finite (finitely many $\beta$, and for each
$\beta$ at most one $\nu'$ contributes).

(c) $t_{\Lambda+\pbrdc}=m(\Lambda)=1$, and if $t_\gamma\ne0$ then $\gamma=(\nu+\pbrdc)\circ w$ with
$m(\nu)\ne0$, so $\pbrlev(\nu+\pbrdc)\le\pbrlev(\Lambda+\pbrdc)$ by Lemma~\ref{pbr:mult}.
\end{proof}

\subsection{Densities on the flag variety}\label{pbr:ss-flag}

Let $\pbrB=\{b\in\mathrm{GL}_n(\C):b_{ij}=0\text{ whenever }\lambda_i>\lambda_j\}$ (upper triangular
when the indices are listed with $\lambda$ increasing), $\pbrN=\{\nu:\nu_{aa}=1,\ \nu_{ij}=0\text{ for
}i\ne j\text{ unless }\lambda_i>\lambda_j\}$, $\pbrn$ its Lie algebra, and for $Y\in M_n(\C)$ let
$Y_-\in\pbrn$ be the matrix of the entries $Y_{ij}$ with $\lambda_i>\lambda_j$. For $b\in\pbrB$ and
$w\in\Z^n$ put $b^w=\prod_ab_{aa}^{w_a}$, and recall $2\rho_a=\#\{b:\lambda_b<\lambda_a\}-\#\{b:
\lambda_b>\lambda_a\}$. The Haar measure $d\nu$ of the unipotent group $\pbrN$ is Lebesgue measure in
the entries $\nu_{ij}$, $\lambda_i>\lambda_j$. A \emph{density} is a Borel function $F$ on
$\mathrm{GL}_n(\C)$, with values in $\C$ or in $[0,\infty]$, such that $F(yb)=|b^{2\rho}|^2F(y)$ for
$y\in\mathrm{GL}_n(\C)$ and $b\in\pbrB$.

\begin{lemma}\label{pbr:jac}
For $b\in\pbrB$ and $X\in M_n(\C)$, $(bXb^{-1})_-=(bX_-b^{-1})_-$, and the $\C$-linear map
$Y\mapsto(bYb^{-1})_-$ of $\pbrn$ has determinant $b^{2\rho}$.
\end{lemma}
\begin{proof}
$X-X_-$ has its entries at $\lambda_i\le\lambda_j$, and so has $b(X-X_-)b^{-1}$, as $\pbrB$ is a group;
this gives the first claim. Let $r(a)$ be the rank of $\lambda_a$ among the entries of $\lambda$. For
$\lambda_i>\lambda_j$, $(bYb^{-1})_{ij}=\sum_{k,l}b_{ik}Y_{kl}(b^{-1})_{lj}$, and a nonzero term has
$\lambda_i\le\lambda_k$, $\lambda_l\le\lambda_j$, $\lambda_k>\lambda_l$, so $r(k)-r(l)\ge r(i)-r(j)$
with equality only for $(k,l)=(i,j)$. Ordering the pairs by decreasing $r(i)-r(j)$, the matrix of
the map is triangular with diagonal entries $b_{ii}(b^{-1})_{jj}=b_{ii}/b_{jj}$, and
$\prod_{\lambda_i>\lambda_j}b_{ii}/b_{jj}=\prod_ab_{aa}^{\#\{j:\lambda_j<\lambda_a\}-\#\{i:\lambda_i>
\lambda_a\}}=b^{2\rho}$.
\end{proof}

By Gaussian elimination, $g\in\mathrm{GL}_n(\C)$ lies in $\pbrN\pbrB$ if and only if the principal
minors $\Delta_{L_r}(g)$, $L_r$ the set of the $r$ indices with the smallest $\lambda_a$, are nonzero
for $1\le r\le n$; then $g=\nu(g)b(g)$ uniquely, with $\nu(g)$, $b(g)$ rational in $g$. A group
$H=\auU(J')=\{h:h^*J'h=J'\}$, $J'$ a diagonal sign matrix, carries the Haar measure compatible with
the Haar measure of mass $1$ on $\auU(n)$ (\S\ref{pds:ss-ds1}(iii)); for $J'=\pm1$ it is that measure.

\begin{proposition}\label{pbr:flag}
There is a constant $j>0$ such that for every $H=\auU(J')$ as above and every density $F$ which is
nonnegative, or such that $|F|$ is integrable on $H$,
\[
\int_HF(h)\,dh=j^{-1}\int_{\pbrN\cap H\pbrB}F(\nu)\,d\nu .
\]
\end{proposition}
\begin{proof}
(i) $H\cap\pbrB=T$: if $b\in H\cap\pbrB$ then $b^{-1}=J'b^*J'$ has its nonzero entries at
$\lambda_i\ge\lambda_j$ and also, being in $\pbrB$, at $\lambda_i\le\lambda_j$; so $b$ is diagonal,
and $b^*J'b=J'$ gives $|b_{aa}|=1$.

(ii) $H\smallsetminus\pbrN\pbrB$ is null: it is the zero set of the real-analytic function
$\prod_r\Delta_{L_r}$ on the connected group $H$, which is $1$ at $1$.

(iii) Let $\varphi(h)=\nu(h)$ on $H\cap\pbrN\pbrB$. Then $\varphi(ht)=\varphi(h)$ for $t\in T$, and if
$\varphi(h)=\varphi(h')$ then $h^{-1}h'=b(h)^{-1}b(h')\in H\cap\pbrB=T$. The image of $\varphi$ is
$\pbrN\cap H\pbrB$. For $X\in\mathfrak h=\operatorname{Lie}H$ and $h=\nu b$, $he^{sX}=\nu(1+sY+O(s^2))
b$ with $Y=bXb^{-1}$, and $1+sY=(1+sY_-)(1+s(Y-Y_-))+O(s^2)$ with $1+s(Y-Y_-)$ triangular, so
\[
d\varphi_h(hX)=\nu\cdot(bXb^{-1})_-=\nu\cdot(bX_-b^{-1})_-
\]
by Lemma~\ref{pbr:jac}. If this vanishes then $X_-=0$; then $X\in\mathfrak h$ has its nonzero entries
at $\lambda_i\le\lambda_j$, and $X=-J'X^*J'$ has them at $\lambda_i\ge\lambda_j$, so $X\in
\mathfrak t=\operatorname{Lie}T$. As $\dim_\R\mathfrak h-\dim\mathfrak t=n^2-n=\dim_\R\pbrn$, the map
$\bar\varphi\colon(H\cap\pbrN\pbrB)/T\to\pbrN$ is an injective local diffeomorphism onto the open set
$\pbrN\cap H\pbrB$.

(iv) Let $\mathfrak m\subset\mathfrak h$ be the subspace of matrices with zero diagonal, a complement
of $\mathfrak t$. The Haar measure of $H$ corresponds to a Lebesgue measure on $\mathfrak h$ through
the exponential map near $0$, and that of $T$ (mass $1$) to one on $\mathfrak t$; let $d_{\mathfrak m}$
be the quotient Lebesgue measure on $\mathfrak m\cong\mathfrak h/\mathfrak t$. The invariant measure
on $H/T$ with $\int_Hf=\int_{H/T}\int_Tf(ht)\,dt$ has, in the chart $X\mapsto h\exp(X)T$ of a
neighbourhood of $hT$, a continuous density with respect to $d_{\mathfrak m}$ equal to $1$ at $X=0$.
Let $j_H=|\det_\R(X\mapsto X_-)|$ for the map $\mathfrak m\to\pbrn$ and the measures
$d_{\mathfrak m}$, $d\nu$. By (iii), Lemma~\ref{pbr:jac} and the left invariance of $d\nu$, the
Jacobian of $\bar\varphi$ at $hT$ is $j_H|b(h)^{2\rho}|^2$ (it depends only on $hT$ since
$|t^{2\rho}|=1$). By the change of variables formula and (ii), for a density $F\ge0$, using
$F(ht)=F(h)$ and $F(h)=F(\nu(h)b(h))=|b(h)^{2\rho}|^2F(\nu(h))$,
\[
\int_HF\,dh=\int_{(H\cap\pbrN\pbrB)/T}F(h)\,d(hT)=\int_{\pbrN\cap H\pbrB}\frac{|b^{2\rho}|^2F(\nu)}
{j_H|b^{2\rho}|^2}\,d\nu=j_H^{-1}\int_{\pbrN\cap H\pbrB}F(\nu)\,d\nu .
\]
A complex density with $|F|$ integrable is a combination of four nonnegative densities
($(\pbrre F)^\pm$, $(\operatorname{Im}F)^\pm$).

(v) $j_H$ does not depend on $J'$. By definition of compatibility, the map $c\colon X\mapsto X_{\mathrm
d}+iX_{\mathrm o}$ from $\mathfrak h$ to $\mathfrak u(n)$, where $X_{\mathrm d}$ (resp.\ $X_{\mathrm o}$)
keeps the entries with $\varepsilon'_i=\varepsilon'_j$ (resp.\ $\ne$), carries the Lebesgue measure of
$\mathfrak h$ to that of $\mathfrak u(n)$. It is the identity on $\mathfrak t$ and maps $\mathfrak m$
onto the corresponding subspace of $\mathfrak u(n)$, so it carries $d_{\mathfrak m}$ to the measure
used for $\auU(n)$. Moreover $(cX)_-=\iota(X_-)$, where $\iota$ multiplies the entries $(i,j)$ with
$\varepsilon'_i\ne\varepsilon'_j$ by $i$, a $\C$-linear map with $|\det_\R\iota|=1$. Hence
$j_H=j_{\auU(n)}=:j$.
\end{proof}

\begin{corollary}\label{pbr:flagG}
The set $G\pbrB$ is open. Let $F$ be a density that is continuous on $\mathrm{GL}_n(\C)$. Then
$F|_G\in L^1(G)$ and
\[
\int_GF(g)\,dg=\int_{\auU(n)\cap G\pbrB}F(k)\,dk=\int_{\auU(n)}F(k)\,dk-\int_{\auU(n)\smallsetminus
G\pbrB}F(k)\,dk .
\]
\end{corollary}
\begin{proof}
The map $G\times\pbrB\to\mathrm{GL}_n(\C)$, $(g,b)\mapsto gb$, is a submersion: its differential at
$(g,b)$ has image $gb\,(b^{-1}\mathfrak gb+\mathfrak b)$, and $b^{-1}\mathfrak gb\cap\mathfrak b=
b^{-1}(\mathfrak g\cap\mathfrak b)b=b^{-1}\mathfrak tb$ has dimension $n$ (as in (i) above), so the
image has real dimension $n^2+(n^2+n)-n=2n^2$. So $G\pbrB$ is open. The function $F1_{G\pbrB}$ is a
density, since $G\pbrB\pbrB=G\pbrB$. By Gram--Schmidt, $\auU(n)\pbrB=\mathrm{GL}_n(\C)$, so
$\pbrN\cap\auU(n)\pbrB=\pbrN$. Proposition~\ref{pbr:flag} for $H=\auU(n)$ and for $H=G$ gives
$\int_{\auU(n)}|F|1_{G\pbrB}=j^{-1}\int_{\pbrN\cap G\pbrB}|F|=\int_G|F|$, finite because $\auU(n)$ is
compact; and then the same with $F$ in place of $|F|$.
\end{proof}

\subsection{The orbital integral on an open part of the torus}\label{pbr:ss-open}

Let $\pbrTp=\{x\in T:\pbrre(\varepsilon_ax_a)>0\text{ for all }a\}$, the set of $x\in T$ with $Jx$
accretive. For an accretive matrix $A$ and $y\in\mathrm{GL}_n(\C)$ put
\[
F_A(y)=\frac{\pcom_{-\pbrLo}(y^*Jy)}{\pcom_{-\mu}(y^*Ay)},\qquad \pbrE(A)=\int_{\auU(n)\smallsetminus
G\pbrB}F_A(k)\,dk .
\]

\begin{lemma}\label{pbr:FA}
$F_A$ is a density, continuous on $\mathrm{GL}_n(\C)$, and $A\mapsto\pbrE(A)$ is holomorphic on the
open set of accretive matrices.
\end{lemma}
\begin{proof}
$y^*Jy$ is invertible and $y^*Ay$ is accretive (Lemma~\ref{pco:minors}(d)), so the numerator is
defined and continuous (a polynomial in the entries divided by a power of $\det(y^*Jy)$) and the
denominator is continuous and nonzero. Let $b\in\pbrB$ and $w\in\{-\pbrLo,-\mu\}$. If $w_c<w_d$ then
$\lambda_c>\lambda_d$ (for $w=-\pbrLo$ because $\pbrLo$ is dominant, for $w=-\mu$ by
Lemma~\ref{pco:minors}(c)), so $b_{cd}=0$; likewise $(b^*)_{cd}=\aubar{b_{dc}}=0$ whenever $w_d<w_c$.
Lemma~\ref{pco:minors}(b) gives $\pcom_w(b^*Xb)=\aubar{b^w}\,b^w\pcom_w(X)$ for invertible $X$, so
$F_A(yb)=|b^{-\pbrLo}|^2|b^{-\mu}|^{-2}F_A(y)=|b^{2\rho}|^2F_A(y)$, as $\mu-\pbrLo=2\rho$. The integrand
of $\pbrE$ is jointly continuous in $(A,k)$ and holomorphic in $A$, and the domain is compact.
\end{proof}

Let $(\sigma,V_\sigma)$ be the irreducible representation of $\auU(n)$ with highest weight $\pbrLo$
(character $\chi_\lambda$), extended holomorphically to $\mathrm{GL}_n(\C)$, with an invariant inner
product and a unit highest weight vector $v^\circ$, and $\pbrNo=(\dim V_\sigma)^{-1}$.

\begin{lemma}\label{pbr:hw}
$\pcom_{\pbrLo}(g)=\langle\sigma(g)v^\circ,v^\circ\rangle$ for $g\in\mathrm{GL}_n(\C)$; hence
$\int_{\auU(n)}|\pcom_{\pbrLo}|^2=\pbrNo$ and $\int_{\auU(n)}\pcom_{\pbrLo}(y^{-1}xy)\,dy=\pbrNo
\chi_\lambda(x)$. For $z,a\in\mathrm{GL}_n(\C)$ with $a^{-1}z$ accretive,
\begin{equation}\label{pbr:eq-FJn}
\pcom_{\pbrLo}(z)=\frac1{\pbrNo}\int_{\auU(n)}\frac{\pcom_{\pbrLo}(au)\,\aubar{\pcom_{\pbrLo}(u)}}
{\pcom_{-\mu}(u^{-1}a^{-1}zu)}\,du .
\end{equation}
\end{lemma}
\begin{proof}
$\sigma(g)^*=\sigma(g^*)$: both sides are antiholomorphic in $g$ and agree on $\auU(n)$, and a
holomorphic function on $\mathrm{GL}_n(\C)$ vanishing on $\auU(n)$ vanishes (its derivatives at $1$
along $\mathfrak u(n)$, hence along $\mathfrak{gl}_n=\mathfrak u(n)\oplus i\mathfrak u(n)$, vanish).
Let $g=\nu't\eta$ with $\nu'$ unipotent with nonzero off-diagonal entries only at $\lambda_a<
\lambda_b$, $t$ diagonal, $\eta$ unipotent with nonzero off-diagonal entries only at $\lambda_a>
\lambda_b$; such products form a dense open subset. Since $v^\circ$ is killed by the $d\sigma(E_{ab})$,
$\lambda_a>\lambda_b$, $\sigma(\eta)v^\circ=v^\circ=\sigma(\nu'^*)v^\circ$, so $\langle\sigma(g)v^\circ,
v^\circ\rangle=\langle\sigma(t)v^\circ,\sigma(\nu'^*)v^\circ\rangle=t^{\pbrLo}$. By
Lemma~\ref{pco:minors}(b) (if $\pbrLo_a<\pbrLo_b$ then $\lambda_a<\lambda_b$, so $\eta_{ab}=0$;
similarly for $\nu'$), $\pcom_{\pbrLo}(g)=t^{\pbrLo}$. Both sides are holomorphic on
$\mathrm{GL}_n(\C)$, so they agree. The two integrals follow from Schur orthogonality on $\auU(n)$ and
from $\int\sigma(y^{-1}xy)\,dy=\pbrNo\chi_\lambda(x)\cdot1$.

Proposition~\ref{pco:indep} and its proof apply verbatim to the signature $(n,0)$, that is, to $J=1$,
$\pcoD=M_n(\C)$, $\pcoH=\mathrm{GL}_n(\C)$, $K=\auU(n)$; then $\rho_{\mathrm n}=0$, $\Lambda=\pbrLo$,
and $\mu$ is unchanged (no gap condition enters there). So the right side of \eqref{pbr:eq-FJn} does
not depend on $a$ with $a^{-1}z$ accretive. For $a=z$ it is $(\pbrNo)^{-1}\int\pcom_{\pbrLo}(zu)
\aubar{\pcom_{\pbrLo}(u)}\,du$ ($\pcom_{-\mu}(1)=1$), and with $\pcom_{\pbrLo}(zu)=\langle\sigma(u)
v^\circ,\sigma(z^*)v^\circ\rangle$ Schur orthogonality gives $(\pbrNo)^{-1}\cdot\pbrNo\,\aubar{\langle
\sigma(z^*)v^\circ,v^\circ\rangle}=\langle\sigma(z)v^\circ,v^\circ\rangle=\pcom_{\pbrLo}(z)$.
\end{proof}

\begin{lemma}\label{pbr:Un}
For $x\in\pbrTp$, $\int_{\auU(n)}F_{Jx}(k)\,dk=\pbrNo\chi_\lambda(x)$.
\end{lemma}
\begin{proof}
Write $m=\pcom_{\pbrLo}$ and $\psi(Y)=\pcom_{-\mu}(Y)^{-1}$. For $y\in\auU(n)$ apply
\eqref{pbr:eq-FJn} to $z=y^{-1}xy$ and $a=y^{-1}Jy$; then $a^{-1}z=y^*(Jx)y$ is accretive. With
Lemma~\ref{pbr:hw},
\[
\pbrNo\chi_\lambda(x)=\int m(y^{-1}xy)\,dy=\frac1{\pbrNo}\iint m(y^{-1}Jyk)\,\aubar{m(k)}\,
\psi\bigl((yk)^*Jx(yk)\bigr)\,dk\,dy .
\]
The substitution $(y,k)=(vu,u^{-1})$ preserves the product measure on $\auU(n)^2$; it gives
$y^{-1}Jyk=u^{-1}v^*Jv$, $\aubar{m(u^{-1})}=m(u)$ and $yk=v$. By Schur orthogonality, with
$b=v^*Jv$, $\int m(u^{-1}b)m(u)\,du=\int\langle\sigma(u)v^\circ,v^\circ\rangle\allowbreak \aubar{\langle\sigma(u)
v^\circ,\sigma(b)v^\circ\rangle}du=\pbrNo\,m(b)$. Hence $\pbrNo\chi_\lambda(x)=\int m(v^*Jv)\psi(v^*Jxv)
\,dv$, and $m(v^*Jv)=\pcom_{-\pbrLo}((v^*Jv)^{-1})=\pcom_{-\pbrLo}(v^*Jv)$ by
Lemma~\ref{pco:minors}(a), since $(v^*Jv)^{-1}=v^*Jv$.
\end{proof}

\begin{proposition}\label{pbr:orbit}
Let $x\in\pbrTp$ and $\sigma_\lambda=\prod_a\varepsilon_a^{\Lambda_a+\pbrLo_a}\in\{\pm1\}$. Then
$g\mapsto c(g^{-1}xg)$ is integrable on $G$ and
\[
\sigma_\lambda\int_Gc(g^{-1}xg)\,dg=\pbrNo\chi_\lambda(x)-\pbrE(Jx).
\]
\end{proposition}
\begin{proof}
Let $g\in G$ and $y=g^{-1}xg\in G$. Since $g^{-1}=Jg^*J$, $J^{-1}y=g^*(Jx)g$ is accretive, so $J\in
A_y$ in the notation of \S\ref{pco:ss-indep}. By Theorem~\ref{pco:FJ}, $c(y)=F(y)=d_\Lambda I(y;k)$
with $y=kz$ the polar decomposition, and by Proposition~\ref{pco:indep}, $c(y)=d_\Lambda I(y;J)$. In
\eqref{pco:eq-I}, $\pcom_\Lambda(Ju)=J^\Lambda\pcom_\Lambda(u)$ (Corollary~\ref{pco:equiv}(b)) and
$u^{-1}J^{-1}yu=(gu)^*Jx(gu)$. On the other hand $(gu)^*J(gu)=J$ for $u\in K$, so the numerator of
$F_{Jx}(gu)$ is $\pcom_{-\pbrLo}(J)=J^{-\pbrLo}=J^{\pbrLo}$. Hence
\[
\sigma_\lambda\,c(g^{-1}xg)=\int_Kd_\Lambda|\pcom_\Lambda(u)|^2F_{Jx}(gu)\,du .
\]
By Lemma~\ref{pbr:FA} and Corollary~\ref{pbr:flagG}, $F_{Jx}\in L^1(G)$, and $\int_G|F_{Jx}(gu)|dg=
\int_G|F_{Jx}|$ for each $u$. By Tonelli the right side is integrable over $G$, and by Fubini and
$d_\Lambda\int_K|\pcom_\Lambda|^2=1$, $\sigma_\lambda\int_Gc(g^{-1}xg)\,dg=\int_GF_{Jx}$. Conclude with
Corollary~\ref{pbr:flagG} and Lemma~\ref{pbr:Un}.
\end{proof}

\subsection{The error term}\label{pbr:ss-error}

\begin{lemma}\label{pbr:gain}
Let $u\in\auU(n)$, $f\in\Z^n$, $D_f=\operatorname{diag}(2^{f_1},\dots,2^{f_n})$ and $X_f=u^*D_fu$. Let
$w\in\Z^n$, $m\in\Z$ and $h'\in\N$ be such that $f_a+|m|\le f_c$ whenever $w_c<w_a$, and let $i,j$ satisfy
$|w_i-w_j|\ge h'$. Then $|\pcom_w(D_f)|\,(|u_{ij}|^22^m)^{h'}\le|\pcom_w(X_f)|$.
\end{lemma}
\begin{proof}
If $w_i<w_j$, replace $(w,f,X_f)$ by $(-w,-f,X_f^{-1}=u^*D_{-f}u)$: the hypotheses are preserved (with
the same $i,j$), $\pcom_w(X_f)=\pcom_{-w}(X_f^{-1})$ by Lemma~\ref{pco:minors}(a), and $|\pcom_{-w}
(D_{-f})|=|\pcom_w(D_f)|$. So assume $w_j\le w_i-h'$. Let $w_0=\min w$, $w_1=\max w$, $S_t=\{a:w_a\ge
t\}$. From the definition, $\pcom_w(X)=(\det X)^{w_0}\prod_{w_0<t\le w_1}\Delta_{S_t}(X)$ for invertible
$X$, and $\det X_f=\det D_f$; all minors involved are positive. Since $w_0\le w_j\le w_i-h'$, the $h'$
integers $t\in(w_i-h',w_i]$ lie in $(w_0,w_1]$, and it suffices to show, for $w_0<t\le w_1$, $S=S_t$,
\[
g_t\,\Delta_S(D_f)\le\Delta_S(X_f),\qquad g_t=|u_{ij}|^22^m\text{ if }w_i-h'<t\le w_i,\quad g_t=1
\text{ otherwise}.
\]
$S^c=\{a:w_a<t\}$ is nonempty; let $c\in S^c$ minimise $f$ on $S^c$. For $a\in S$, $w_a\ge t>w_c$, so
$f_a+|m|\le f_c$. Put $g=\min(f,f_c\mathbf 1)$ entrywise: $g_a=f_a$ on $S$ and $g_a=f_c$ on $S^c$. Let
$Q=\{j\}$ if $w_i-h'<t\le w_i$ (then $i\in S$, $j\in S^c$) and $Q=\emptyset$ otherwise, and $p=S^c
\smallsetminus Q$. (1) $g_t\le2^{f_c|Q|}\Delta_Q(X_g^{-1})$: for $Q=\{j\}$, $(X_g^{-1})_{jj}=(u^*D_{-g}
u)_{jj}=\sum_a|u_{aj}|^22^{-g_a}\ge|u_{ij}|^22^{-f_i}$ and $f_c-f_i\ge|m|\ge m$. (2) By Jacobi, Fischer
($Q^c=p\sqcup S$), and the monotonicity of minors ($X_g\le X_{f_c\mathbf 1}=2^{f_c}\cdot1$ and $X_g\le X_f$
since $g\le f_c\mathbf1$, $g\le f$),
\[
\Delta_Q(X_g^{-1})=\frac{\Delta_{Q^c}(X_g)}{\det X_g}\le\frac{\Delta_p(X_g)\Delta_S(X_g)}{\det X_g}
\le\frac{2^{f_c|p|}\Delta_S(X_f)}{2^{\sum_{a\in S}f_a+|S^c|f_c}} .
\]
(3) As $|p|+|Q|=|S^c|$ and $\Delta_S(D_f)=2^{\sum_{a\in S}f_a}$, (1) and (2) give the claim.
\end{proof}

\begin{lemma}\label{pbr:compact}
For every $h'\in\N$ there is $\eta>0$ such that every $k\in\auU(n)\smallsetminus G\pbrB$ has an entry
$k_{ij}$ with $\varepsilon_i\ne\varepsilon_j$ and $\eta\,|\pcom_{-\pbrLo}(k^*Jk)|<|k_{ij}|^{2h'}$.
\end{lemma}
\begin{proof}
$C=\auU(n)\smallsetminus G\pbrB$ is compact (Corollary~\ref{pbr:flagG}), and $\phi(k)=\pcom_{-\pbrLo}
(k^*Jk)$ is continuous on $\auU(n)$. A unitary $k$ with $k_{ij}=0$ whenever $\varepsilon_i\ne
\varepsilon_j$ lies in $K\subset G\pbrB$. So each $k_0\in C$ has such an entry $k_{0,ij}\ne0$, and
$\{k:|k_{ij}|^{2h'}>\eta_0|\phi(k)|\}$, with $\eta_0=|k_{0,ij}|^{2h'}/(1+2|\phi(k_0)|)$, is an open
neighbourhood of $k_0$. Take a finite subcover and the least $\eta_0$.
\end{proof}

\begin{proposition}\label{pbr:Ebound}
Let $e\in\Z^n$. There is $C_e$ such that $|\pbrE(\operatorname{diag}(2^{-(ro+e)}))|\le C_e\,2^{-r(
\pbrlev(\mu)+h_0)}$ for all large integers $r$.
\end{proposition}
\begin{proof}
Let $C\ge2\max_a|e_a|$ be an integer, $r>C$, $f=-(ro+e)$, $A=D_f$ (positive definite, so accretive).
We apply Lemma~\ref{pbr:gain} with $w=-\mu$, $m=r-C$, $h'=h_0$. If $w_c<w_a$ then $\mu_a<\mu_c$, so
$\lambda_a<\lambda_c$ and $o_a\ge o_c+1$, whence $f_c-f_a=r(o_a-o_c)+e_a-e_c\ge r-C=|m|$. If
$\varepsilon_i\ne\varepsilon_j$ then $|\mu_i-\mu_j|=|\lambda_i-\lambda_j|+|\rho_i-\rho_j|\ge h_0$ ($\rho$
is ordered like $\lambda$). Also $|\pcom_{-\mu}(D_f)|=2^{\sum_a\mu_a(ro_a+e_a)}=2^{r\pbrlev(\mu)+e\cdot
\mu}$. Let $k\in\auU(n)\smallsetminus G\pbrB$ and $i,j$ as in Lemma~\ref{pbr:compact} with $h'=h_0$
(so $k_{ij}\ne0$). Then
\[
|F_A(k)|\le\frac{\eta^{-1}|k_{ij}|^{2h_0}}{2^{r\pbrlev(\mu)+e\cdot\mu}(|k_{ij}|^22^{r-C})^{h_0}}=
\eta^{-1}2^{Ch_0-e\cdot\mu}\,2^{-r(\pbrlev(\mu)+h_0)},
\]
and the domain of integration has mass at most $1$.
\end{proof}

For $\gamma\in\Z^n$ write $s^\gamma=\prod_a\varepsilon_a^{\gamma_a}\in\{\pm1\}$.

\begin{lemma}\label{pbr:low}
Let $P'$, $R$ be Laurent polynomials with $P'(x)=R(x)\,\pbrE(Jx)$ for all $x\in\pbrTp$, and let
$\gamma\in\Z^n$ satisfy $\pbrlev(\gamma)<\pbrlev(\mu)+\pbrlev(\beta)+h_0$ for every exponent $\beta$ of
$R$. Then $P'_\gamma=0$.
\end{lemma}
\begin{proof}
\emph{Continuation.} Let $f\in\Z^n$, $f\ne0$, and $\Sigma=\{z=\varrho e^{i\theta}:\varrho>0,\ |\theta|<
\pi/(2\max|f_a|)\}$. For $z\in\Sigma$, $\operatorname{diag}(z^{f_a})$ is accretive, and both sides of
\[
\sum_{\gamma'}P'_{\gamma'}s^{\gamma'}z^{f\cdot\gamma'}=\Bigl(\sum_\beta R_\beta s^\beta z^{f\cdot\beta}
\Bigr)\pbrE\bigl(\operatorname{diag}(z^{f_1},\dots,z^{f_n})\bigr)
\]
are holomorphic on the connected set $\Sigma$ (Lemma~\ref{pbr:FA}). For $z=e^{i\theta}\in\Sigma$ the
point $x$ with $x_a=\varepsilon_az^{f_a}$ lies in $\pbrTp$ and $Jx=\operatorname{diag}(z^{f_a})$, so
the identity holds on an arc through $1$, hence on $\Sigma\supset(0,\infty)$; for $f=0$ it holds
trivially.

\emph{Vanishing.} Suppose $P'_\gamma\ne0$, and let $\beta_0$ be the least level of an exponent of $P'$;
then $\beta_0\le\pbrlev(\gamma)<\pbrlev(\mu)+\pbrlev(\beta)+h_0$ for all exponents $\beta$ of $R$. Fix
$e\in\Z^n$ and take $z=\frac12$, $f=ro+e$ ($r\in\N$), multiplied by $2^{r\beta_0}$:
\[
\sum_{\gamma'}P'_{\gamma'}s^{\gamma'}2^{-e\cdot\gamma'}2^{-r(\pbrlev(\gamma')-\beta_0)}=\Bigl(\sum_\beta
R_\beta s^\beta2^{-e\cdot\beta}2^{-r(\pbrlev(\beta)-\beta_0)}\Bigr)\pbrE\bigl(\operatorname{diag}(
2^{-(ro+e)})\bigr).
\]
As $r\to\infty$ the left side tends to $\sum_{\pbrlev(\gamma')=\beta_0}P'_{\gamma'}s^{\gamma'}2^{-e\cdot
\gamma'}$, and by Proposition~\ref{pbr:Ebound} the right side is $O\bigl(\sum_\beta2^{-r(\pbrlev(\beta)
+\pbrlev(\mu)+h_0-\beta_0)}\bigr)\to0$. So $\sum_{\pbrlev(\gamma')=\beta_0}P'_{\gamma'}s^{\gamma'}
2^{-e\cdot\gamma'}=0$ for all $e\in\Z^n$. The maps $e\mapsto2^{-e\cdot\gamma'}$ are distinct characters
of $\Z^n$, hence linearly independent, so $P'_{\gamma'}=0$ whenever $\pbrlev(\gamma')=\beta_0$,
contradicting the choice of $\beta_0$.
\end{proof}

\subsection{Proof of Theorem~\ref{pbr:main}}\label{pbr:ss-proof}

Let $P$ be the Laurent polynomial of Proposition~\ref{pbr:laurent}, $\pbrtop=\kappa+N\pbrdD$, and
\[
Q=D^N\pbrAlt_{\pbrW}(\kappa),\qquad Q^{(1)}=D^N\pbrAlt_{S_n}(\kappa).
\]

\emph{Step 1: leading terms.} By Lemma~\ref{pbr:basic}(d), $(\Lambda+\pbrdc)+M\pbrdD=\kappa+N\pbrdD=
\pbrtop$, and $\pbrtop$ is ordered like $\lambda$ (Lemma~\ref{pbr:tau}). Lemma~\ref{pbr:top} with
$k=M$, $a=\Lambda+\pbrdc$ and the series $t$ (Proposition~\ref{pbr:laurent}(b),(c)) gives $P_\pbrtop=1$,
and every exponent $\gamma$ of $P$ has $w\in\pbrW$ with $\pbrlev(\gamma\circ w)\le\pbrlev(\pbrtop)$.
Lemma~\ref{pbr:top} with $k=N$, $a=\kappa$ and $t=\pbrAlt_{\pbrW}(\kappa)$ (whose exponents are the
$\kappa\circ w$, $w\in\pbrW$, as $\kappa$ has distinct entries) gives the same two statements for
$Q$. Both $P$ and $Q$ are $\pbrW$-antisymmetric.

\emph{Step 2: the relation on $\pbrTp$.} Let $x\in\pbrTp$. By Proposition~\ref{pbr:orbit},
$\pbrTh(x)$ is defined and $\pbrTh(x)=d\sigma_\lambda(\pbrNo\chi_\lambda(x)-\pbrE(Jx))$. By
Proposition~\ref{pbr:laurent}(a),(b) and Lemma~\ref{pbr:basic}(b), $P(x)=V_{\mathrm c}D^M\pbrTh(x)=
(-1)^{pq}VD^N\pbrTh(x)$, and $V\chi_\lambda=\pbrAlt_{S_n}(\kappa)$ (R2). Hence
\[
P(x)-s_0\,Q^{(1)}(x)=R(x)\,\pbrE(Jx),\qquad s_0=(-1)^{pq}d\,\sigma_\lambda\pbrNo,\quad R=-(-1)^{pq}
d\,\sigma_\lambda\,VD^N .
\]

\emph{Step 3: low levels.} Let $\pbrlev(\gamma)\le\pbrlev(\pbrtop)$. The exponents $\beta$ of $R$ have
$\pbrlev(\beta)\ge0$ (Lemma~\ref{pbr:basic}(c)), and $\pbrtop-\mu=o+N\pbrdD$ (Lemma~\ref{pbr:basic}(d)),
so by Lemma~\ref{pbr:tau}
\[
\pbrlev(\gamma)\le\pbrlev(\mu)+\pbrlev(o+N\pbrdD)<\pbrlev(\mu)+h_0\le\pbrlev(\mu)+\pbrlev(\beta)+h_0 .
\]
By Lemma~\ref{pbr:low}, $P_\gamma=s_0Q^{(1)}_\gamma$. Next, $Q^{(1)}_\gamma=\sum_\beta(D^N)_\beta
\sum_w\pbrsgn(w)$, the inner sum over the $w\in S_n$ with $\kappa\circ w=\gamma-\beta$. If such a $w$
is not in $\pbrW$, some $i$ has $\varepsilon_{w(i)}\ne\varepsilon_i$, so $|\kappa_i-\kappa_{w(i)}|=
|\lambda_i-\lambda_{w(i)}|\ge h_0$, and Lemma~\ref{pbr:basic}(e) ($o$ strictly reverses the order of
$\kappa$) gives $\pbrlev(\kappa\circ w)\ge\pbrlev(\kappa)+h_0$; but $\pbrlev(\kappa\circ w)=\pbrlev
(\gamma)-\pbrlev(\beta)\le\pbrlev(\pbrtop)=\pbrlev(\kappa)+\pbrlev(N\pbrdD)<\pbrlev(\kappa)+h_0$ by
Lemmas~\ref{pbr:basic}(c) and \ref{pbr:tau}. So only $w\in\pbrW$ occur, and $Q^{(1)}_\gamma=Q_\gamma$.
Thus $P_\gamma=s_0Q_\gamma$ whenever $\pbrlev(\gamma)\le\pbrlev(\pbrtop)$.

\emph{Step 4: $P=Q$.} Let $\gamma\in\Z^n$ with $P_\gamma\ne0$ or $Q_\gamma\ne0$. By Step 1 there is
$w\in\pbrW$ with $\pbrlev(\gamma\circ w)\le\pbrlev(\pbrtop)$, and by antisymmetry and Step 3,
$\pbrsgn(w)P_\gamma=P_{\gamma\circ w}=s_0Q_{\gamma\circ w}=s_0\pbrsgn(w)Q_\gamma$. So $P=s_0Q$, and
comparing the coefficients at $\pbrtop$ (Step 1), $s_0=1$ and $P=Q$.

\emph{Step 5: conclusion.} Let $t\in T$ have distinct eigenvalues. By Proposition~\ref{pbr:int},
$g\mapsto c(g^{-1}tg)$ is integrable, and by Proposition~\ref{pbr:laurent} and Step 4,
$(-1)^{pq}V(t)D(t)^N\pbrTh(t)=P(t)=D(t)^N\pbrAlt_{\pbrW}(\kappa)(t)$. As $D(t)\ne0\ne V(t)$, this is
\eqref{pbr:eq-unstable}. For $\sigma\in S_n$, $t^\sigma$ again has distinct eigenvalues,
$(t^\sigma)^\gamma=t^{\gamma\circ\sigma^{-1}}$ and $V(t^\sigma)=\pbrsgn(\sigma)V(t)$, so
\[
\sum_{\sigma\in S_n}\pbrTh(t^\sigma)=\frac{(-1)^{pq}}{V(t)}\sum_{\sigma\in S_n}\sum_{w\in\pbrW}
\pbrsgn(w\sigma^{-1})\,t^{\kappa\circ(w\sigma^{-1})}=(-1)^{pq}|\pbrW|\,\frac{\pbrAlt_{S_n}(\kappa)(t)}
{V(t)}=(-1)^{pq}p!\,q!\,\chi_\lambda(t),
\]
since for fixed $w$, $\sigma\mapsto w\sigma^{-1}$ is a bijection of $S_n$, and by (R2). Taking complex
conjugates gives \eqref{pbr:eq-stable}. \qed

\begin{corollary}[Formal degree]\label{pbr:degree}
Under \eqref{pbr:eq-gap}, $d=\dim V_\sigma=\prod_{i<j}(\lambda_i-\lambda_j)/(j-i)$ (entries listed in
decreasing order), and $\sigma_\lambda=(-1)^{pq}$.
\end{corollary}
\begin{proof}
By Step 4, $1=s_0=(-1)^{pq}\sigma_\lambda\,d/\dim V_\sigma$ with $d>0$; Weyl's dimension formula (R2)
for the highest weight $\pbrLo$ gives $\dim V_\sigma=\prod_{i<j}(\kappa_i-\kappa_j)/(j-i)$, and
$\kappa_i-\kappa_j=\lambda_i-\lambda_j$.
\end{proof}
This is Proposition~\ref{pds:haar} for parameters satisfying \eqref{pbr:eq-gap}, now without
appeal to Harish-Chandra's formula.

\begin{remark}[What is used]\label{pbr:rem-used}
The proof uses: Theorem~\ref{pco:FJ}, Lemma~\ref{pco:decay}, Lemma~\ref{pco:minors},
Corollary~\ref{pco:equiv}, Proposition~\ref{pco:indep} (also in the signature $(n,0)$),
Lemmas~\ref{pds:blockdet} and \ref{pm:kernel}(c), and the inputs (R1)--(R4). The gap condition
enters through Lemma~\ref{pco:decay} (with exponent $pqM\le h_0$), through Lemma~\ref{pbr:tau}
($4n^3<h_0$, $n^3+4n^5<h_0$), and through Proposition~\ref{pbr:Ebound}. Schmid's theorem (R3) enters
only in Lemma~\ref{pbr:mult}, to place the $K$-types of $\pi$ in $\Lambda$ plus the cone of positive
noncompact roots; this is what makes $\Lambda+\pbrdc$ the leading exponent of the series $t$.
\end{remark}


%% ===== proofs/C-determinant.tex =====
% proofs/C-determinant.tex
% The determinant part of Theorem C: deduction of Theorem C from its inputs, the triangle of lifts,
% the factorisation of the chain determinant (au:factor), injectivity (au:inj), auxiliary forms
% (au:aux). Labels of this file start with pc:. Only amsmath/amssymb/amsthm are used.
\providecommand{\auO}{\mathcal O}
\providecommand{\auI}{\mathcal I}
\providecommand{\auQlb}{\overline{\Q}_\ell}
\providecommand{\auAut}{\operatorname{Aut}}
\providecommand{\auNm}{\operatorname{Nm}}
\providecommand{\auU}{\mathrm U}
\providecommand{\auGL}{\mathrm{GL}}
\providecommand{\auHom}{\operatorname{Hom}}
\providecommand{\auEnd}{\operatorname{End}}
\providecommand{\auA}{\mathbb A}
\providecommand{\auF}{\mathbb F}
\providecommand{\auPl}{\Phi}
\providecommand{\auT}{\mathcal T}
\providecommand{\auAz}{\mathcal A_0}
\providecommand{\auAm}{\mathcal A_m}
\providecommand{\auG}{\overline G}
\providecommand{\auLat}{\mathcal L}
\providecommand{\audT}{d_{\auT}}
\providecommand{\auhmain}{h_{\mathrm{main}}}
\providecommand{\auhdir}{h_{\mathrm{dir}}}
\providecommand{\aueu}{\varepsilon_{\mathrm u}}
\providecommand{\aubud}{\mathrm{bud}}
\providecommand{\aubar}[1]{\overline{#1}}
\providecommand{\auf}[1]{\texttt{#1}}
\providecommand{\psHS}{\mathcal H}
\providecommand{\psrho}{\varrho}
\providecommand{\psvac}{\mathbf v}
\providecommand{\psZ}{Z}
% macros of this file
\providecommand{\pcF}{\mathcal F}
\providecommand{\pcG}{\mathsf G}
\providecommand{\pcm}{\mathfrak m}
\providecommand{\pcb}{\mathfrak b}
\providecommand{\pcR}{R}
\providecommand{\pcH}{\mathcal H^{M}}
\providecommand{\pcRal}{\mathcal R}
\providecommand{\pcmc}{\operatorname{mc}}
\providecommand{\pcvol}{\operatorname{vol}}
\providecommand{\pctr}{\operatorname{tr}}
\providecommand{\pcMat}{\mathrm M}
\providecommand{\pcW}{\mathbb W}
\providecommand{\pcFF}{\mathfrak F}
\providecommand{\pop}{\mathsf p}

\section{The determinant of the chain of lifts: proofs for Theorem C}\label{pc:sec}

This section treats the part of Theorem~\ref{thm:srcC} that does not concern a single lift
$W_j\to V_j$: the deduction of the theorem from its inputs, the injectivity of the lower chain
(Proposition~\ref{au:inj}), the factorisation of the chain determinant
(Proposition~\ref{au:factor}) and the auxiliary forms (Proposition~\ref{au:aux}). The notation is
that of \S\S\ref{ss:au-notation}--\ref{ss:au-spaces} and of \S\ref{ps:ss-weil}; inner products are
conjugate-linear in the first variable.

\emph{Organisation.} The section has two layers.
\begin{itemize}
\item[(I)] \emph{Proofs given here.} Theorem~\ref{thm:srcC} from Propositions~\ref{au:inj},
\ref{au:aux}, \ref{au:factor}, \ref{au:regular} and~\ref{au:Zj} (Proposition~\ref{pc:C}); the
linear algebra of a \emph{triangle} of maps, which gives injectivity and the Gram form of a chain
from commutation relations in squares (Propositions~\ref{pc:tri-inj}, \ref{pc:tri-gram}); the
series of Rallis' formula, its convergence and its injectivity
(Lemma~\ref{pc:series}); the comparison of its determinants on inner forms
(Lemma~\ref{pc:detcompare}); the inner product formula for a theta lift with \emph{sharp} data in
Weil's range (Proposition~\ref{pc:rallis}), from Weil's form of the Siegel formula and the facts
(W1)--(W5), (W7) (Lemmas~\ref{ps:weil}, \ref{pc:SW}); the finite algebra of matched functions
(Lemmas~\ref{pc:pointlaw}, \ref{pc:next}); the bookkeeping of levels and indices
(Lemmas~\ref{pc:algpet}, \ref{pc:index}); and the deductions of Propositions~\ref{au:inj},
\ref{au:factor} and~\ref{au:aux} from the four statements of layer (II).
\item[(II)] \emph{Statements proved in later sections.} Lemma~\ref{pc:arch} (archimedean types;
Proposition~\ref{ptv:arch}), Theorem~\ref{pc:tri} (the lifts form a triangle whose squares commute
up to constants of known modulus; \S\S\ref{ptr:sec}--\ref{ptw:sec}), Theorem~\ref{pc:X} (the
exchange identity, on which Theorem~\ref{pc:tri} rests; Theorem~\ref{pcx:X}, \S\S\ref{pcx:sec},
\ref{pcu:sec}, \ref{pct:sec}, \ref{pcn:sec}) and Theorem~\ref{pc:auxcore} (integral and short lifts
through the second grid; \S\S\ref{pgr:sec}--\ref{pgm:sec}). In addition Proposition~\ref{au:mult} is
needed in the form $(\mathrm M^+)$ of \S\ref{pc:ss-rallis}, which is its statement (for the full
spherical Hecke algebra away from $M$).
\end{itemize}
\S\ref{pc:ss-remains} lists where the inputs are proved and the published results on which they
rest.

\subsection{Theorem C from its inputs}\label{pc:ss-thmC}

We recall the rules for heights of \S\ref{ss:au-heights} in the form used.

\begin{lemma}\label{pc:heights}
\begin{enumerate}
\item[(a)] If the families $y,y'$ have heights at most $c,c'$, then $yy'$ has height at most
$c+c'$, and $y^n$ has height at most $nc$ for every integer $n\ge0$. A family of height at most $c$
has height at most $c''$ for every $c''\ge c$.
\item[(b)] Let $N_\iota$ be positive integers with $\ell(\iota)\nmid N_\iota$ and let
$e\in\tfrac12\Z$. If $y$ has height at most $c$, so has $\iota\mapsto N_\iota^{\,e}y_\iota$.
\item[(c)] Let $\theta$ be a regular family of $m$-fold twists
\textup{(Definition~\ref{def:au-regular})}, let $\rho_1,\dots,\rho_n$ be ramified ratios in the
sense of \S\ref{ss:au-proofC}, and let $e\in\Z$. Then the family
\[
\iota\longmapsto\Bigl((\auNm Q_{\theta_\iota})^{e/2}\prod_{a=1}^n\rho_a(\theta_\iota)\Bigr)^{-1}
\]
is well defined and has height at most $n\,t\,x_1$.
\end{enumerate}
\end{lemma}
\begin{proof}
(a) If $(H,\alpha,h)$ and $(H',\alpha',h')$ are as in Definition~\ref{def:au-height} for $y$ and
$y'$, then $(HH',\alpha+\alpha',h+h')$ serves for $yy'$; the last assertion is clear from the
definition. (b) The positive real number $N_\iota^{\,e}$ is an algebraic number whose square is an
integral power of $N_\iota$; as $\ell(\iota)\nmid N_\iota$, $|N_\iota^{\,e}|_\iota=1$, and the
triple $(H,\alpha,h)$ of $y$ serves. (c) By definition $\rho_a(\vartheta)=\prod_{i=1}^{t'}
|L(z_i,(\nu_i\vartheta^{\pm n_i})^{-1})|/|L(z_i,\nu_i^{-1})|$ with $t'\le t$, $1\le n_i\le m$ and
ramified critical pairs $(\nu_i,z_i)$. By (R2) each family $\iota\mapsto|L(z_i,\nu_i^{-1})|/
|L(z_i,(\nu_i\theta_\iota^{\pm n_i})^{-1})|$ is defined and has height at most $x_1$. By (a),
$\iota\mapsto\rho_a(\theta_\iota)^{-1}$ has height at most $t'x_1\le tx_1$, and the product over
$a$ has height at most $ntx_1$. By (R1), $\ell(\iota)\nmid\auNm Q_{\theta_\iota}$, and (b) applies
with the exponent $-e/2$.
\end{proof}

\begin{proposition}[Theorem C from its inputs]\label{pc:C}
Assume Propositions~\ref{au:inj}, \ref{au:aux}, \ref{au:factor}, \ref{au:regular}
and~\ref{au:Zj}. Then Theorem~\ref{thm:srcC} holds: if $y_1,\dots,y_n\in\auT$ are linearly independent
with $n\ge(1-4/b^2)\audT$, of index at most $\delta$ and covolume at most $c$ in $\auLat_D$, then
$c+2\delta\ge\audT(\auhmain-3fa(\log b)^4)$.
\end{proposition}
\begin{proof}
\emph{Step 0: the lattice.} By Proposition~\ref{au:inj} the maps $\Theta_\vartheta=\iota_\vartheta|
_{\auT}$ are injective for all $m$-fold twists, and \eqref{eq:au-pairing} holds because the levels
$J_\vartheta$ are free (Lemma~\ref{au:free}); so $\auLat_D$ is a twisted lattice in the sense of
Definition~\ref{def:au-lattice}, and Theorem~\ref{au:twistprime} applies to it.

\emph{Step 1: completion to a basis.} Let $e'$, $\rho'_{ij}$, $\gamma$ and the spanning set
$\mathfrak A\subseteq\auT$ be as in Proposition~\ref{au:aux}. As $y_1,\dots,y_n$ are independent
and $\mathfrak A$ spans $\auT$, the exchange lemma of Steinitz gives $z_1,\dots,z_{n'}\in
\mathfrak A$, $n'=\audT-n$, such that $y_1,\dots,y_n,z_1,\dots,z_{n'}$ is a basis of $\auT$. By
hypothesis $n'\le4\audT/b^2$.

\emph{Step 2: the twists.} Let $\theta$ be a regular family (Proposition~\ref{au:regular}). By (R1)
the twist $\theta_\iota$ is allowed at $\ell(\iota)$ for every $\iota\in\auI$.

\emph{Step 3: the hypotheses of Theorem~\ref{au:twistprime}.} (i) is the hypothesis on $y$. (ii)
For every $i$, every $\iota$ and every $m$-fold twist $\xi$ with $\ell(\iota)\nmid\auNm Q_\xi$, the
form $\gamma(\xi)\Theta_\xi z_i$ is integral at $\iota$, by Proposition~\ref{au:aux}; and
$\gamma(\xi)\ne0$. (iii) holds with $b_{\mathrm{short}}=2x_1^2b^2$, by the same proposition. (iv)
First, $|\gamma(\theta_\iota)|^{-2}$ is the square of the inverse of $(\auNm Q_{\theta_\iota})^{e'/2}$
times $m^2$ ramified ratios; by Lemma~\ref{pc:heights}(a),(c) it has height at most
$2m^2tx_1$, and $2m^2tx_1\le b$ because $m^2\le4x$, $t=5r\le15x^5$ and $x_1^{10^6}\le b$; so
$b_{\mathrm{den}}=b$ serves. Secondly, by Proposition~\ref{au:factor},
\[
\frac{Z(1)}{Z(\theta_\iota)}=\biggl(\Bigl((\auNm Q_{\theta_\iota})^{e/2}\prod_{0<i<j\le m}\rho_{ij}
(\theta_\iota)\Bigr)^{-1}\biggr)^{\audT}\cdot\prod_{j=1}^m\frac{Z_j(1)}{Z_j(\theta_\iota^{\,j})} .
\]
The first factor is the $\audT$-th power of a family which, by Lemma~\ref{pc:heights}(c) with
fewer than $m^2$ ratios, has height at most $m^2tx_1\le b$. By Proposition~\ref{au:Zj} each of the
$m$ remaining factors has height at most $\audT(fa(\log b)^4/m-k\auhdir)$. By
Lemma~\ref{pc:heights}(a) and $m\,k\,\auhdir=\auhmain$
(Definition~\ref{def:au-hmain}(c)), the family $Z(1)/Z(\theta_\iota)$ has height at most
\[
b_{\mathrm{vol}}=\audT\bigl(b+fa(\log b)^4-\auhmain\bigr).
\]

\emph{Step 4: conclusion.} Theorem~\ref{au:twistprime} gives
\[
0\le c+2\delta+n'(2x_1^2b^2+b)+\audT\bigl(b+fa(\log b)^4-\auhmain\bigr),
\]
and $n'(2x_1^2b^2+b)\le4\audT b^{-2}(2x_1^2b^2+b)=4\audT(2x_1^2+b^{-1})\le\audT b$. Hence
$\audT(\auhmain-fa(\log b)^4-2b)\le c+2\delta$, and $2b\le2fa(\log b)^4$.
\end{proof}

Thus, after Proposition~\ref{au:Zj} (proved in \S\ref{ss:au-proofC} from
Propositions~\ref{au:onelift}, \ref{au:hecke} and~\ref{au:mult}), Theorem~\ref{thm:srcC} rests on
Propositions~\ref{au:inj}, \ref{au:factor}, \ref{au:aux} and~\ref{au:regular}. The rest of this
section concerns the first three.

\subsection{Linear algebra: volumes, moduli, triangles}\label{pc:ss-la}

\begin{lemma}[Volumes]\label{pc:vol}
Let $u\colon U\to V$ be a linear map of finite-dimensional inner product spaces, $S\in\auEnd(U)$
and $a\in\C$ with $\langle ux,uy\rangle=a\langle x,Sy\rangle$ for all $x,y\in U$. Then
$\det(u^*u)=a^{\dim U}\det S$.
\end{lemma}
\begin{proof}
The hypothesis says $u^*u=aS$.
\end{proof}

\begin{definition}\label{pc:modulus}
Let $c$ be a complex function and $f$ a real function on the set of $m$-fold twists. We say that
$c$ \emph{has modulus $f$ up to norm powers} if there are a real $C\ge0$ and $k\in\Z$ with
$|c(\vartheta)|^2=C\,(\auNm Q_\vartheta)^k f(\vartheta)^2$ for all $\vartheta$.
\end{definition}

\begin{lemma}\label{pc:mod}
\begin{enumerate}
\item[(a)] Constants have modulus $1$; if $c,c'$ have moduli $f,f'$ up to norm powers, then $cc'$
has modulus $ff'$; if moreover $c$ has no zero, $c^{-1}$ has modulus $f^{-1}$ and $\aubar c$ has
modulus $f$.
\item[(b)] If $c$ has modulus $f$ up to norm powers, $f\ge0$, $f(1)=1$ and $c(1)\ne0$, then there
is $e\in\Z$ with $|c(\vartheta)/c(1)|=(\auNm Q_\vartheta)^{e/2}f(\vartheta)$ for all $\vartheta$.
\end{enumerate}
\end{lemma}
\begin{proof}
(a) Multiply the identities. (b) At $\vartheta=1$, $Q=\emptyset$ and $\auNm Q=1$, so $C=|c(1)|^2>0$;
divide and take the non-negative square root; $e=k$.
\end{proof}

\begin{definition}[Triangle]\label{pc:triangle}
Let $m\ge1$, let $I$ be a set, $V$ a complex inner product space, and for $i,n\ge0$, $i+n\le m$,
let $\pcF_{i,n}$ be complex inner product spaces (not necessarily complete). A \emph{triangle} of
size $m$ over $I$ consists of linear isometries $j_i\colon V\to\pcF_{i,0}$ $(0\le i\le m)$ and, for
every $\vartheta\in I$, of numbers $w_n=w_n^\vartheta>0$ $(0\le n\le m)$ and linear maps
\[
a_{i,n}=a_{i,n}^\vartheta\colon\pcF_{i,n}\to\pcF_{i,n+1},\qquad
d_{i,n}=d_{i,n}^\vartheta\colon\pcF_{i,n+1}\to\pcF_{i+1,n}\qquad(i+n<m)
\]
such that:
\begin{enumerate}
\item[(T1)] \emph{(squares)} for $i+n+2\le m$ there is $\sigma_{i,n}\colon I\to\C$ with
$w_n\,a_{i+1,n}\circ d_{i,n}=\sigma_{i,n}(\vartheta)\,w_{n+2}\,d_{i,n+1}\circ a_{i,n+1}$ on
$\pcF_{i,n+1}$;
\item[(T2)] \emph{(first adjoint)} for $i<m$, $y\in V$ and $z\in\pcF_{i,1}$,
$\langle z,a_{i,0}j_iy\rangle=\langle d_{i,0}z,j_{i+1}y\rangle$;
\item[(T3)] for $i<m$, $d_{i,0}a_{i,0}j_i(V)\subseteq j_{i+1}(V)$, and $d_{i,0}a_{i,0}j_i$ is
injective.
\end{enumerate}
The triangle \emph{has the adjoint law} if moreover
\begin{enumerate}
\item[(T4)] for $i+n+2\le m$ there is $\kappa_{i,n}\colon I\to\C$ such that every
$z\in\pcF_{i,n+2}$ has an $x\in\pcF_{i,n+1}$ with $\langle a_{i,n+1}v,z\rangle=\langle v,x\rangle$
for all $v\in\pcF_{i,n+1}$ and
$\langle d_{i,n}x,v'\rangle=\kappa_{i,n}(\vartheta)\langle d_{i,n+1}z,a_{i+1,n}v'\rangle$ for all
$v'\in\pcF_{i+1,n}$.
\end{enumerate}
The \emph{chains} are $X_{i,n}=X_{i,n}^\vartheta=(w_{n-1}a_{i,n-1})\circ\dots\circ(w_0a_{i,0})\colon
\pcF_{i,0}\to\pcF_{i,n}$, with $X_{i,0}=1$.
\end{definition}

In the application $\pcF_{i,n}$ is a space of automorphic forms on the unitary group of a space
with a stem and the blocks $i+1,\dots,i+n$; $a$ adds the last block, $d$ removes the first one, and
(T1) says that these two operations commute up to a constant.

\begin{proposition}[Chains are injective]\label{pc:tri-inj}
In a triangle of size $m$, for $i+n<m$ and all $\vartheta$,
\begin{equation}\label{pc:eq-tau}
X_{i+1,n}\circ d_{i,0}\circ a_{i,0}=\tau_{i,n}\;d_{i,n}\circ X_{i,n+1},\qquad
\tau_{i,n}=\frac{w_{n+1}}{w_0w_1}\prod_{l<n}\sigma_{i,l}(\vartheta).
\end{equation}
For $i+n\le m$ the map $X_{i,n}\circ j_i\colon V\to\pcF_{i,n}$ is injective. If $V\ne0$, then
$\sigma_{i,n}(\vartheta)\ne0$ for all $i+n+2\le m$ and all $\vartheta$.
\end{proposition}
\begin{proof}
\eqref{pc:eq-tau} for $n=0$ reads $d_{i,0}a_{i,0}=w_0^{-1}d_{i,0}(w_0a_{i,0})$. If it holds for
$n$ and $i+n+2\le m$, then by (T1)
\[
X_{i+1,n+1}d_{i,0}a_{i,0}=w_na_{i+1,n}\,\tau_{i,n}d_{i,n}X_{i,n+1}=\tau_{i,n}\sigma_{i,n}w_{n+2}\,
d_{i,n+1}a_{i,n+1}X_{i,n+1}=\tau_{i,n}\sigma_{i,n}\frac{w_{n+2}}{w_{n+1}}\,d_{i,n+1}X_{i,n+2},
\]
which is \eqref{pc:eq-tau} for $n+1$. Injectivity is proved by induction on $n$, for all $i$ with
$i+n\le m$; the case $n=0$ is clear. Let $i+n+1\le m$ and $X_{i,n+1}j_iy=0$. By (T3) there is
$y'\in V$ with $j_{i+1}y'=d_{i,0}a_{i,0}j_iy$, and by \eqref{pc:eq-tau}
$X_{i+1,n}j_{i+1}y'=\tau_{i,n}d_{i,n}X_{i,n+1}j_iy=0$. By induction $y'=0$, and $y=0$ by (T3).
Finally let $V\ne0$ and suppose $\sigma_{i,n}(\vartheta)=0$ with $i+n+2\le m$. Then
$\tau_{i,n+1}=0$, so $X_{i+1,n+1}d_{i,0}a_{i,0}j_iy=0$ for all $y\in V$; as $(i+1)+(n+1)\le m$, the
injectivity just proved and (T3) give $y=0$ for all $y\in V$, a contradiction.
\end{proof}

\begin{proposition}[The Gram form of a chain]\label{pc:tri-gram}
Let a triangle of size $m$ over $I$ have the adjoint law, let $V\ne0$, and put
$\beta_{i,n}=1/\aubar{\sigma_{i,n}}$. Suppose that for $i<m$ and $\vartheta\in I$ there are
$S_i=S_i^\vartheta\in\auEnd(V)$ and $C_i(\vartheta)\in\C$ with
\begin{equation}\label{pc:eq-first}
w_0\,\langle a_{i,0}j_ix,a_{i,0}j_iy\rangle=C_i(\vartheta)\,\langle x,S_iy\rangle\qquad(x,y\in V).
\end{equation}
Then for all $\vartheta$ and $x,y\in V$
\[
\langle X_{0,m}j_0x,X_{0,m}j_0y\rangle=\frac{w_0w_1^{\,m}}{w_m}\ \prod_{i<m}C_i(\vartheta)\
\prod_{\substack{i,n\ge0\\ i+n+2\le m}}\kappa_{i,n}(\vartheta)\beta_{i,n}(\vartheta)\ \cdot\
\bigl\langle x,S_0S_1\cdots S_{m-1}\,y\bigr\rangle .
\]
\end{proposition}
\begin{proof}
\emph{Step 1.} For $i+n<m$ put $\kappa_i^{(n)}=\prod_{l<n}\kappa_{i,l}$. Then for all
$z\in\pcF_{i,n+1}$ and $y\in V$
\begin{equation}\label{pc:eq-chainadj}
\langle z,a_{i,n}X_{i,n}j_iy\rangle=\kappa_i^{(n)}\,\langle d_{i,n}z,X_{i+1,n}j_{i+1}y\rangle .
\end{equation}
For $n=0$ this is (T2). Let it hold for $n$, let $i+n+2\le m$ and $z\in\pcF_{i,n+2}$, and let $x$
be as in (T4). Since $w_n>0$,
\begin{align*}
\langle z,a_{i,n+1}X_{i,n+1}j_iy\rangle&=\langle x,X_{i,n+1}j_iy\rangle=w_n\langle x,a_{i,n}X_{i,n}
j_iy\rangle=w_n\kappa_i^{(n)}\langle d_{i,n}x,X_{i+1,n}j_{i+1}y\rangle\\
&=w_n\kappa_i^{(n)}\kappa_{i,n}\langle d_{i,n+1}z,a_{i+1,n}X_{i+1,n}j_{i+1}y\rangle=\kappa_i^{(n+1)}
\langle d_{i,n+1}z,X_{i+1,n+1}j_{i+1}y\rangle .
\end{align*}

\emph{Step 2.} Let $i+n<m$, $x,y\in V$, and let $x'\in V$ be the vector with
$j_{i+1}x'=d_{i,0}a_{i,0}j_ix$ (T3). By Proposition~\ref{pc:tri-inj}, $\tau_{i,n}\ne0$ and
$d_{i,n}X_{i,n+1}j_ix=\tau_{i,n}^{-1}X_{i+1,n}j_{i+1}x'$. With \eqref{pc:eq-chainadj} for
$z=X_{i,n+1}j_ix$,
\begin{align*}
\langle X_{i,n+1}j_ix,X_{i,n+1}j_iy\rangle&=w_n\langle X_{i,n+1}j_ix,a_{i,n}X_{i,n}j_iy\rangle
=w_n\kappa_i^{(n)}\langle d_{i,n}X_{i,n+1}j_ix,X_{i+1,n}j_{i+1}y\rangle\\
&=\frac{w_n\,\kappa_i^{(n)}}{\aubar{\tau_{i,n}}}\,\langle X_{i+1,n}j_{i+1}x',X_{i+1,n}j_{i+1}y\rangle,
\qquad\frac1{\aubar{\tau_{i,n}}}=\frac{w_0w_1}{w_{n+1}}\prod_{l<n}\beta_{i,l}.
\end{align*}
Moreover, by (T2) with $z=a_{i,0}j_ix$ and by \eqref{pc:eq-first}, for every $v\in V$
\[
\langle x',v\rangle=\langle d_{i,0}a_{i,0}j_ix,j_{i+1}v\rangle=\langle a_{i,0}j_ix,a_{i,0}j_iv
\rangle=\frac{C_i}{w_0}\langle x,S_iv\rangle .
\]

\emph{Step 3.} We prove by induction on $n$, for $i=m-n$, that
$\langle X_{i,n}j_ix,X_{i,n}j_iy\rangle=K_{i,n}\langle x,S_iS_{i+1}\cdots S_{m-1}y\rangle$ with
$K_{m,0}=1$ and
\[
K_{i,n+1}=\frac{w_1w_n}{w_{n+1}}\,C_i\,\kappa_i^{(n)}\prod_{l<n}\beta_{i,l}\cdot K_{i+1,n}.
\]
Indeed, by Step 2 and the induction hypothesis for $(i+1,n)$ applied to $x'$ and $y$,
$\langle X_{i,n+1}j_ix,X_{i,n+1}j_iy\rangle=w_n\kappa_i^{(n)}\frac{w_0w_1}{w_{n+1}}\prod_{l<n}
\beta_{i,l}\cdot K_{i+1,n}\langle x',S_{i+1}\cdots S_{m-1}y\rangle$, and $\langle x',v\rangle=
(C_i/w_0)\langle x,S_iv\rangle$ with $v=S_{i+1}\cdots S_{m-1}y$. Unfolding the recursion from
$(0,m)$ to $(m,0)$, the weights give $\prod_{p=0}^{m-1}w_1w_{m-p-1}/w_{m-p}=w_1^{\,m}w_0/w_m$, and
each $\kappa_{i,l}$ and each $\beta_{i,l}$ with $i+l+2\le m$ occurs exactly once.
\end{proof}


\subsection{Lifts with data modulo an ideal; the series of Rallis' formula}\label{pc:ss-data}

\subsubsection*{Conventions}
A \emph{modulus} is a non-zero ideal $\pcm$ of $\auO_E$ prime to $\aubar\pcm$; so every prime $w$
dividing $\pcm$ lies above a place $v=w\aubar w$ of $E^+$ that splits in $E$. Put
$\pcR_\pcm=\auO_E/\pcm$, and let $S(\pcm)$ be the set of places of $E^+$ below $\pcm$. For a twist
$\vartheta$ with set $Q=Q_\vartheta$ we write
\[
\pcm_\vartheta=(\pi_M)\prod_{w\in Q}w,\qquad\pcR_\vartheta=\auO_E/\pcm_\vartheta=\prod_{w\mid
\pcm_\vartheta}\auF_w ,
\]
a product of finite fields. If $\vartheta$ is $m$-fold then $Q_{\vartheta^j}=Q_\vartheta$ and
$\pcm_{\vartheta^j}=\pcm_\vartheta$ for $1\le j\le m$.

A \emph{space} is a hermitian space $S$ over $E$ with a basis (its \emph{frame}) whose
$\widehat\auO_E$-span $L_S$ is a unimodular lattice, and with the same signature at all infinite
places. Let $K_S\subset\auU(S)(\auA_{E^+,f})$ be the stabiliser of $L_S$. At $v\in S(\pcm)$ the
frame identifies $\auU(S)(E^+_v)$ with $\auGL_N(E_w)$, $N$ the rank; reduction gives a surjection
$K_S\to\auGL_N(\pcR_\pcm)$, with kernel $J_S(\pcm)$. The Haar measure on $\auU(S)(\auA_{E^+})$ is the
product of a fixed Haar measure on $\auU(S)_\infty=\prod_\tau\auU(S\otimes_{E^+,\tau}\R)$ (of total
mass $1$ if $S$ is definite) and of the measure on the finite adelic points that gives $K_S$ the
volume $1$; $[\auU(S)]=\auU(S)(E^+)\backslash\auU(S)(\auA_{E^+})$. For spaces of the same rank $N$ the
finite adelic groups, with $K_S$ and the reductions, are identified by the frames: the localisations
of the lattices are isometric at every finite place, because $E/E^+$ is unramified there and the
lattices are unimodular, so they have orthonormal bases (proof: the reduction of the lattice is a non-degenerate hermitian space over the quadratic extension of residue fields, so it has a vector $\bar x$ with $\bar h(\bar x,\bar x)\ne0$ (otherwise polarisation and the surjectivity of the residue trace give $\bar h=0$); a lift $x$ has $h(x,x)$ a unit of the base ring, hence the norm of a unit since the extension is unramified, so after scaling $h(x,x)=1$; then the lattice is $\mathcal O x\perp(x^\perp\cap\text{lattice})$ with unimodular second summand, and one inducts on the rank; see also Jacobowitz~\cite{Jacobowitz}; at a split place every lattice is
standard).

Let $S,T$ be spaces of ranks $N,N'$. For automorphic characters $\alpha,\beta$ of $\auU(1)$ put
\[
\omega_{\alpha,\beta}(g,h)=\alpha(\det g)^{-1}\beta(\det h)\,\omega_{\auHom(S,T)}(g,h)\qquad
(g\in\auU(S)(\auA),\ h\in\auU(T)(\auA));
\]
the representation $\omega_\mu$ of \S\ref{ps:ss-weil} is $\omega_{\mu\upsilon_{N'},\mu\upsilon_N}$.
At $v\in S(\pcm)$ we use the model (W5): $\psHS_v(\auHom(S,T))$ is the space of Schwartz--Bruhat
functions on $\pcMat_{N'\times N}(E_w)$, matrices being written in the frames. For a function
$\Psi$ on $\pcMat_{N'\times N}(\pcR_\pcm)$ let $\Psi^\natural$ be the function on
$\prod_{w\mid\pcm}\pcMat_{N'\times N}(E_w)$ that is $\Psi(a^{-1}x\bmod\pcm)$ at $x\in
a\,\pcMat_{N'\times N}(\widehat\auO)$ and $0$ elsewhere, where $a=(a_w)$ is the idele of Definition~\ref{pgx:conv}: $a_w$ is the generator of the ideal
$\alpha_w\auO_w$ fixed in \S\ref{ps:ss-weil} (there $\alpha_w$ denotes the component of a scalar idele,
which has nothing to do with the characters $\alpha,\beta$ of this section), equal to $1$ when
$\alpha_w$ is a unit (Definition~\ref{pgx:conv};
the coordinates are described in Lemma~\ref{pgx:dict}); so
$\mathbf 1^\natural$ is the vacuum of $a\auHom(L_S,L_T)$ at $w$ (Lemma~\ref{pgx:vac}). For a
vector $P$ of the Fock space of $\auHom(S,T)\otimes\R$ (for instance a polynomial)
the \emph{datum} of $(P,\Psi)$ is $[P,\Psi]=P\otimes\Psi^\natural\otimes\psvac^{\pcm}$, where
$\psvac^{\pcm}$ is the vacuum of the lattice $a\auHom(L_S,L_T)$ at the finite places outside
$S(\pcm)$, normalised to have norm $1$, and the norm at $v\in S(\pcm)$ is the $L^2$-norm for the
measure giving $a_w\pcMat(\auO_w)$ volume $1$. The \emph{kernel}, the \emph{lift} and the
\emph{transposed lift} are
\begin{gather*}
\theta_{\alpha,\beta}(g,h;P,\Psi)=\theta\bigl(\omega_{\alpha,\beta}(g,h)[P,\Psi]\bigr),\\
\Theta_{\alpha,\beta}(P,\Psi)\varphi\,(h)=\int_{[\auU(S)]}\theta_{\alpha,\beta}(g,h;P,\Psi)\,
\varphi(g)\,dg,\qquad
\Theta^{\mathsf t}_{\alpha,\beta}(P,\Psi)\varphi'\,(g)=\int_{[\auU(T)]}\theta_{\alpha,\beta}(g,h;
P,\Psi)\,\varphi'(h)\,dh .
\end{gather*}
The \emph{point datum} at a matrix $x$ is the indicator function $\delta_x$ of $\{x\}$; $\iota$
denotes the matrix $\binom10$ of the inclusion of frames. Thus, with the notation of
\S\ref{ps:ss-Y}, $\Theta^Y_{j,\vartheta}=[K_j:J_j]\cdot\Theta_{\mu_{j+1}\upsilon_{r+(j+1)t},\,
\mu_{j+1}\upsilon_{r+jt}}(P_j,\delta_\iota)$ with the modulus $\pcm_\vartheta$: at a level place $w$ the
datum of \S\ref{ps:ss-Y} is $\mathbf 1_{a_w(\iota+\varpi_w\auHom)}$, which is $\delta_\iota^\natural$
(Lemma~\ref{pgx:dict}(c)), including at the primes of $Q_\vartheta$ where $\alpha_w$ is not a unit.

\begin{lemma}\label{pc:level}
Let $\alpha,\beta$ be unramified outside $\pcm$ and tamely ramified at the primes of $\pcm$, with
tame parts giving characters $\alpha_\pcm,\beta_\pcm$ of $\pcR_\pcm^\times$. For $k\in K_S$,
$k'\in K_T$ with reductions $\bar k,\bar k'$ modulo $\pcm$,
\[
\omega_{\alpha,\beta}(k,k')[P,\Psi]=[P,\Psi'],\qquad\Psi'(x)=\alpha_\pcm(\det\bar k)^{-1}\beta_\pcm
(\det\bar k')\,\Psi(\bar k'^{-1}x\bar k).
\]
\end{lemma}
\begin{proof}
As in Lemma~\ref{ps:level}: outside $S(\pcm)$ the pair $(k,k')$ stabilises the lattice and the
characters are unramified, so (W4) applies; at $v\in S(\pcm)$, (W5) gives
$\Psi^\natural(k'^{-1}xk)$ times the characters (as $a$ is a scalar, $k'^{-1}xk\in a\,\pcMat(\widehat\auO)$
if and only if $x$ is, and $a^{-1}k'^{-1}xk=k'^{-1}(a^{-1}x)k$), and the factor $\upsilon'_w$ is trivial
on units.
\end{proof}

\subsubsection*{Forms and types}
Let $S$ be a space, $\pi=\bigotimes_\tau\pi_\tau$ an irreducible square-integrable representation of
$\auU(S)_\infty$ and $v_0\in\pi$ a unit vector. The \emph{type} is the function
$e(g)=d_\pi\aubar{\langle\pi(g)v_0,v_0\rangle}$ on $\auU(S)_\infty$, $d_\pi$ the formal degree. Let
$\pcF(S,e)$ be the space of bounded continuous functions $\varphi$ on $[\auU(S)]$ that are
right-invariant under an open subgroup of $K_S$ and are reproduced by right convolution with $e$
over $\auU(S)_\infty$, that is, $\varphi(x)=\int_{\auU(S)_\infty}\varphi(xg^{-1})e(g)\,dg$
(inner products are conjugate-linear in the first variable, so $e(g^{-1})=d_\pi\aubar{\langle v_0,
\pi(g)v_0\rangle}$, and the condition is $\varphi(x)=\int e(g^{-1})\varphi(xg)\,dg$, in agreement
with \eqref{par:eq-coef} and Definition~\ref{pz:def-type}); it carries the Petersson product $\langle\varphi',\varphi\rangle=
\int_{[\auU(S)]}\aubar{\varphi'}\varphi$ and the unitary action $\rho$ of the finite adelic group
by right translation. The types used in this paper are integrable functions (their parameters are
far from the walls; this is part of Lemma~\ref{pc:arch} below). For an integrable type, by Schur
orthogonality, right convolution with $e$ is the orthogonal projection of $L^2([\auU(S)])$ onto the
image of $\auHom_{\auU(S)_\infty}(\pi,L^2)\otimes v_0$: every $\varphi\in\pcF(S,e)$ is
$\Phi(v_0)$ for an intertwining map $\Phi\colon\pi\to L^2([\auU(S)])$, and
$\langle\Phi'(u),\Phi(v)\rangle=\langle\varphi',\varphi\rangle\langle u,v\rangle$ by Schur's lemma.
In particular for $\varphi,\varphi'\in\pcF(S,e)$ and an integrable function $c$ on
$\auU(S)_\infty$,
\begin{equation}\label{pc:eq-schur}
\int_{\auU(S)_\infty}c(g)\langle\varphi',\rho(g)\varphi\rangle\,dg=\Bigl(\int_{\auU(S)_\infty}c(g)
\langle v_0,\pi(g)v_0\rangle\,dg\Bigr)\langle\varphi',\varphi\rangle .
\end{equation}
The space $\auAz$ of \S\ref{ss:au-spaces} is $\pcF(Y_0,e_0)^{J_M}$ for the type $e_0$ of the
parameter $\varsigma$, where $J_M$ also denotes $J_{Y_0}((\pi_M))$; the spaces $\mathcal A(W_j)$
are $\pcF(W_j,e_{W_j})^{J_M}$.

\begin{lemma}[Archimedean types]\label{pc:arch}
Let $(S,T)$ be one of the pairs $(W_j,V_j)$, $1\le j\le m$, or one of the pairs
$(\pcG_{i,0},\pcG_{i,1})$, $0\le i<m$, of \S\ref{pc:ss-triangle}, with the types $e_S,e_T$ of the
chain and the joint harmonic polynomial $P$ of the pair. Then:
\begin{enumerate}
\item[(A1)] the types $e_S$, $e_T$ are integrable, and every element of $\pcF(S,e_S)$ and of
$\pcF(T,e_T)$ decreases faster than every power of the height on the automorphic quotient;
\item[(A2)] for every modulus $\pcm$, every $\Psi$ and all characters $\alpha,\beta$ of the weights
of the pair, tamely ramified at $\pcm$ and unramified elsewhere, $\Theta_{\alpha,\beta}(P,\Psi)$
maps $\pcF(S,e_S)$ to $\pcF(T,e_T)$;
\item[(A3)] the integral $Z_\infty=\int_{\auU(S)_\infty}\langle P,\omega_\infty(g,1)P\rangle\,\langle
v_0,\pi_S(g)v_0\rangle\,dg$ converges absolutely and is not zero; it does not depend on the twist,
which is trivial at the infinite places.
\end{enumerate}
\end{lemma}
\noindent This is Proposition~\ref{ptv:arch}: (A1), (A2) for the pairs $(\pcG_{i,0},\pcG_{i,1})$ are
Theorem~\ref{ptr:edge}, (A3) is Proposition~\ref{ptr:Z} with Theorem~\ref{ptq:F}; for $(W_j,V_j)$,
(A1) is Lemma~\ref{plf:A12}, and (A2), (A3) follow from Step~1 of the proof of
Proposition~\ref{plf:L4} (modulo (V), (D1)--(D3)).

\subsubsection*{The coefficient of the vacuum and the series}
Let $S$ be a space of rank $n$, $M'\ge1$, and $\alpha$ an automorphic character of $\auU(1)$
unramified outside a modulus $\pcm$. Let $G^\pcm$ be the group of finite adelic points of
$\auU(S)$ with trivial components at $S(\pcm)$, and $K^\pcm=K_S\cap G^\pcm$. For $\gamma\in G^\pcm$
put
\[
\pcmc_{M',\alpha}(\gamma)=\alpha(\det\gamma)^{-1}\,\bigl\langle\psvac^\pcm,\omega_{\auHom(S,T)}
(\gamma,1)\psvac^\pcm\bigr\rangle ,
\]
where $T$ is any space of rank $M'$.

\begin{lemma}[The series]\label{pc:series}
\begin{enumerate}
\item[(a)] $\pcmc_{M',\alpha}$ does not depend on the choice of $T$, is the same for $S$ and for
the space $S^-$ with the negative form, and for all spaces $S$ of rank $n$ under the identification
of their finite adelic groups; it is bi-invariant under $K^\pcm$; and
$|\pcmc_{M',\alpha}(\gamma)|=[L_S:L_S\cap\gamma L_S]^{-M'/2}$.
\item[(b)] If $n\ge2$ and $M'\ge6n$, then
$\sum_{\gamma\in G^\pcm/K^\pcm,\ \gamma\notin K^\pcm}|\pcmc_{M',\alpha}(\gamma)|\le\exp\bigl(2f(4/5)^n
\bigr)-1$.
\item[(c)] Let $n=r$ and $M'=r+t$. The \emph{series}
$\pcRal=\pcRal_{M',\alpha}=\sum_{\gamma\in G^\pcm/K^\pcm}\pcmc_{M',\alpha}(\gamma)\rho(\gamma)$
converges absolutely in the supremum norm on the bounded $K^\pcm$-invariant functions on
$\auU(S)(\auA)$, maps them to themselves, commutes with left translations and with right
translations by $\auU(S)_\infty$ and by the components at $S(\pcm)$, and satisfies
$\|\pcRal y\|_\infty\ge\tfrac12\|y\|_\infty$. In particular $\pcRal$ is injective, and it preserves
$\pcF(S,e)^J$ for every type $e$ and every open $J\supseteq K^\pcm$.
\end{enumerate}
\end{lemma}
\begin{proof}
(a) The Weil representation of $\auU(S)(E^+_v)$ on $\psHS_v(\auHom(S,T))$, with the splitting
defined by $\upsilon$, depends only on the hermitian space $\auHom(S,T)_v$ with its lattice. At a
finite place all spaces of a given rank are isometric with their lattices; and multiplication by a
unit $\epsilon_v$ of $E_v$ with $\epsilon_v\aubar{\epsilon_v}=-1$ (which exists because $E_v/E^+_v$
is unramified or split) is an isometry $S_v^-\to S_v$ that commutes with the group. This gives the
independence. Bi-invariance follows from (W4). For the absolute value put
$\mathbb L=a\auHom(L_S,L_T)$ and $\mathbb L'=\gamma\mathbb L=a\auHom(\gamma L_S,L_T)$, two
self-dual lattices, and $\Lambda=\mathbb L\cap\mathbb L'$. By (W4),
$\omega(\gamma,1)\psvac_{\mathbb L}$ is a unit vector on the line of $\psvac_{\mathbb L'}$. Since
$\epsilon_{\mathbb L}$ and $\epsilon_{\mathbb L'}$ are given by the same formula, they agree on
$\Lambda$, and the vector
$v'=\sum_{l'\in\mathbb L'/\Lambda}\epsilon_{\mathbb L'}(l')^{-1}\psrho(l')\psvac_{\mathbb L}$ is well
defined; the relation $\epsilon(l)\epsilon(l')=\beta(l,l')\epsilon(l+l')$ gives
$\psrho(l_0)v'=\epsilon_{\mathbb L'}(l_0)v'$ for $l_0\in\mathbb L'$, so $v'$ lies on the line of
$\psvac_{\mathbb L'}$. The vectors $\psrho(l')\psvac_{\mathbb L}$, for $l'$ in distinct classes
modulo $\Lambda$, hence modulo $\mathbb L$, are eigenvectors of the operators $\psrho(l)$,
$l\in\mathbb L$, for distinct characters, hence orthogonal, of norm $1$. So
$\|v'\|^2=[\mathbb L':\Lambda]$ and $\langle\psvac_{\mathbb L},v'\rangle=1$, whence
$|\langle\psvac_{\mathbb L},\omega(\gamma,1)\psvac_{\mathbb L}\rangle|=[\mathbb L':\Lambda]^{-1/2}$.
Finally $\Lambda=a\auHom(L_S+\gamma L_S,L_T)$ and $[\mathbb L':\Lambda]=[L_S+\gamma L_S:\gamma
L_S]^{M'}=[L_S:L_S\cap\gamma L_S]^{M'}$.

(b) \emph{Counting.} Let $\mathfrak o$ be the ring of integers of a local field with residue field
of $q$ elements. By the Hermite normal form, the number $s_k$ of sublattices of index $q^k$ in
$\mathfrak o^n$ is $\sum_{k_1+\dots+k_n=k}q^{\sum_i(i-1)k_i}$; hence
\begin{equation}\label{pc:eq-sk}
\sum_{k\ge0}s_kX^k=\prod_{i=0}^{n-1}(1-q^iX)^{-1},\qquad s_k\le\tbinom{k+n-1}{n-1}q^{k(n-1)}\le
2^{k+n-1}q^{k(n-1)} .
\end{equation}
The same numbers count the lattices containing a given lattice with index $q^k$. Every residue
field of $E$ has at least $5$ elements, because $E$ contains the $5040$-th roots of unity: if the
residue characteristic is not $5$ then $5\mid q_w-1$, and otherwise $7\mid q_w-1$.

\emph{One place.} The function $|\pcmc|$ is a product over the places $v\notin S(\pcm)$ of $E^+$
of the functions $\gamma_v\mapsto[L_{S,v}:L_{S,v}\cap\gamma_vL_{S,v}]^{-M'/2}$, and
$G_v/K_v$ is identified with the set of lattices $L'=\gamma_vL_{S,v}$. Let $\varepsilon_v$ be the
sum of $[L:L\cap L']^{-M'/2}$ over $L'\ne L=L_{S,v}$. If $v=w\aubar w$ splits, $L'$ runs through the
lattices of $E_w^n$ (the component at $\aubar w$ is the dual lattice), and the index is $q^{k+k'}$
with $q=q_w$, $q^k=[L_w:L_w\cap L'_w]$, $q^{k'}=[L'_w:L_w\cap L'_w]$. A lattice $L'_w$ is determined
by $L_w\cap L'_w$ and by $L'_w$ above it, so by \eqref{pc:eq-sk}, with $M'/2\ge3n$,
\[
1+\varepsilon_v\le\Bigl(\sum_ks_kq^{-kM'/2}\Bigr)^2=\prod_{i=0}^{n-1}(1-q^{i-M'/2})^{-2}\le
\exp\Bigl(4\sum_{i<n}q^{i-3n}\Bigr)\le\exp(5q^{-2n-1}),
\]
so $\varepsilon_v\le10\,q^{-2n-1}\le2(4/5)^nq_w^{-2}$, because $50\cdot25^{-n}\le2(4/5)^n$ for
$n\ge2$. If $v$ does not split, $L'$ runs through the self-dual lattices of $S_v$; such a lattice
is determined by $\Lambda=L\cap L'$ and by $L'\supset\Lambda$, both of the same index $q_w^k$ since
$L$ and $L'$ have the same volume. So by \eqref{pc:eq-sk}
\[
\varepsilon_v\le\sum_{k\ge1}s_k^2q_w^{-3nk}\le4^{n-1}\sum_{k\ge1}\bigl(4q_w^{-n-2}\bigr)^k\le
2\cdot4^nq_w^{-n-2}\le2(4/5)^nq_w^{-2}.
\]
\emph{All places.} Choose a place $w$ of $E$ above each $v$. There are at most $2f$ places of $E$
above a rational prime $p$, each with $q_w\ge p$, and $\sum_pp^{-2}<\tfrac12$. Hence
$\sum_v\varepsilon_v\le2f(4/5)^n$ and
$\sum_{\gamma\notin K^\pcm}|\pcmc(\gamma)|=\prod_v(1+\varepsilon_v)-1\le\exp(2f(4/5)^n)-1$.

(c) By (b) with $n=r\ge x^5$ and $f\le x^{12}$ \eqref{eq:au-sizes}, the sum of
$|\pcmc(\gamma)|$ over $\gamma\notin K^\pcm$ is at most $e^{1/3}-1<\tfrac12$ (comparison of sizes).
The series is well defined on $K^\pcm$-invariant functions because $\pcmc$ is right
$K^\pcm$-invariant, its values are $K^\pcm$-invariant because $\pcmc$ is left invariant, and
$\pcRal y=y+Ey$ with $\|Ey\|_\infty\le\tfrac12\|y\|_\infty$. The commutation properties are clear,
and they give the last assertion: boundedness, continuity, the reproducing property and
$J$-invariance are preserved.
\end{proof}

For $1\le j\le m$ and a twist $\vartheta$ we write
\begin{equation}\label{pc:eq-Rj}
\pcRal_j(\vartheta)=\pcRal_{r+t,\ \mu_j^0\vartheta\,\upsilon_{r+t}}\qquad\text{with the modulus }
\pcm_\vartheta ,
\end{equation}
acting on $\pcF(S,e)^{J_M}$ for every space $S$ of rank $r$ and every type $e$; here the forms are
also invariant under $K_S$ at the places of $Q_\vartheta$. Let $\pcH$ be the algebra of compactly
supported functions on $G^{(\pi_M)}$ that are bi-invariant under $K^{(\pi_M)}$, the spherical
Hecke algebra away from $M$ at \emph{all} places, split or not; it acts on the spaces
$\pcF(S,e)^{J_M}$ and commutes with the action of $\auG=K_S/J_M$.

\begin{lemma}\label{pc:heckeclosure}
Let $S,S'$ be spaces of rank $r$ with types $e,e'$ such that $\pcF(S,e)^{J_M}$ and
$\pcF(S',e')^{J_M}$ are finite-dimensional. For every $j$ and $\vartheta$ there is $h\in\pcH$ that
acts as $\pcRal_j(\vartheta)$ on both spaces. In particular $\pcRal_j(\vartheta)$ commutes with
$\auG$ and preserves the images of $\varrho(e_0)$.
\end{lemma}
\begin{proof}
Grouping the cosets into double cosets, $\pcRal_j(\vartheta)=\sum_{[\gamma]\in K^\pcm\backslash
G^\pcm/K^\pcm}\pcmc(\gamma)\,T_\gamma$ on $K^\pcm$-invariant functions, where $\pcm=\pcm_\vartheta$
and $T_\gamma$ is the action of the characteristic function of $K^\pcm\gamma K^\pcm$, an element of
$\pcH$ (extended by the unit at the places of $Q_\vartheta$). The series converges absolutely in
$\auEnd(\pcF(S,e)^{J_M})\times\auEnd(\pcF(S',e')^{J_M})$ by Lemma~\ref{pc:series}(c). The image of
$\pcH$ in this finite-dimensional space is a linear subspace, hence closed, and it contains the
partial sums.
\end{proof}

\noindent For every $M'>2n$ the series converges absolutely and is injective by a finer count
(Proposition~\ref{psr:series}); in this section only $M'=6n$ is used, where the count above suffices.

\subsection{The inner product formula for sharp data}\label{pc:ss-rallis}

\begin{lemma}[Two standard facts]\label{pc:SW}
\begin{enumerate}
\item[(W7)] \emph{(complex conjugation).} Let $S,T$ be spaces of ranks $n,M'$ and $S^-$ the space
$S$ with the negative form. There is a conjugate-linear isometry
$\mathsf C\colon\psHS(\auHom(S,T))\to\psHS(\auHom(S^-,T))$ with
\[
\mathsf C\psrho(x)=\psrho^-(x)\mathsf C,\qquad\theta^-(\mathsf CF)=\aubar{\theta(F)},\qquad
\mathsf C\,\omega_{\auHom(S,T)}(g,h)\,\mathsf C^{-1}=\lambda(g,h)\,\omega_{\auHom(S^-,T)}(g,h),
\]
for a character $\lambda$ of $\auU(S)(\auA)\times\auU(T)(\auA)$ that is trivial on the rational
points and is a product of powers of $\upsilon\circ\det$.
\item[(SW)] \emph{(Weil's Siegel formula).} Let $\pcW$ be a hermitian space of rank $2n$ over $E$
that contains an isotropic subspace of dimension $n$, and let $T$ be a hermitian space of rank
$M'$ such that (i) the theta integral converges absolutely ($T$ anisotropic over $E^+$, or
$M'-r_T>n$ with $r_T$ the Witt index of $T$) and (ii) $M'>2n$; both hold if $M'\ge6n$. Let $\auU(T)(\auA)$ act on $\psHS(\auHom(\pcW,T))$ by the Weil representation
$\omega(1,h)$ with the trivial splitting character; for $\pcW=S\oplus S^-$ and
$\psHS(\auHom(\pcW,T))=\psHS(\auHom(S,T))\otimes\psHS(\auHom(S^-,T))$ this is
$\omega_{\auHom(S,T)}(1,h)\otimes\mathsf C\,\omega_{\auHom(S,T)}(1,h)\,\mathsf C^{-1}$. For a
Lagrangian (maximal isotropic subspace) $L\subset\pcW$ let
$I_L=\{w\in\auHom(\pcW,T):w|_L=0\}$, a maximal isotropic subspace of $\auHom(\pcW,T)$, and let
$\psZ_L=\psZ_{I_L}\colon\psHS(\auHom(\pcW,T))\to\C$ be its zero term \textup{(W2)}. Then for every
smooth $\Phi$
\[
\int_{[\auU(T)]}\theta\bigl(\omega(1,h)\Phi\bigr)\,dh=\pcvol([\auU(T)])\sum_L\psZ_L(\Phi),
\]
the sum being over all Lagrangians of $\pcW$; the integral and the sum converge absolutely, and
$\sum_L|\psZ_L(\omega(g,1)\Phi)|$ is bounded by a power of the height of $g\in\auU(\pcW)(\auA)$.
Moreover, for $\pcW=S\oplus S^-$, data $F,F'$ of $\auHom(S,T)$ and characters $\alpha,\beta$ there are
$C,A$ with
\[
\int_{[\auU(T)]}\bigl|\theta(\omega_{\alpha,\beta}(g,h)F)\,\theta(\omega_{\alpha,\beta}(g',h)F')\bigr|\,dh\le
C\,(\mathrm{ht}(g)\mathrm{ht}(g'))^A\qquad(g,g'\in\auU(S)(\auA));
\]
and if $T$ is anisotropic over $E^+$, $\int|\theta(\omega(g,h)\Phi)|\,dh$ is bounded by a power of the
height of $g\in\auU(\pcW)(\auA)$.
\end{enumerate}
\end{lemma}
\begin{proof}[References and comments]
(W7): the cocycle $\beta$ of $S^-$ is the complex conjugate of that of $S$, so the complex
conjugate of the Heisenberg representation of $\auHom(S,T)$ is that of $\auHom(S^-,T)$; this gives
$\mathsf C$, unique up to a scalar of absolute value $1$, which is fixed by the condition on
$\theta$ (the functional $F\mapsto\aubar{\theta(\mathsf C^{-1}F)}$ is invariant under the rational
vectors, hence a multiple of $\theta^-$ by (W1)). By (W1) the operator $\mathsf C\omega(g,h)
\mathsf C^{-1}$ differs from $\omega^-(g,h)$ by a character; that it is a product of powers of
$\upsilon\circ\det$ follows from the form of the splitting~\cite[Theorem~3.1]{KudlaSplitting}, where
$\upsilon$ is replaced by $\aubar\upsilon=\upsilon^{-1}$. The exponents are computed in
Lemma~\ref{ptw:W7}: $\lambda(g,h)=\upsilon(\det g)^{-M'}\upsilon(\det h)^{n}$. (SW) is the Siegel formula of
Weil~\cite[Th\'eor\`eme~5]{WeilSiegel} in the range of absolute convergence, for the dual pair
$(\auU(\pcW),\auU(T))$: the Eisenstein series $\sum_{\gamma\in P(E^+)\backslash\auU(\pcW)(E^+)}
\psZ_{L_0}(\omega(\gamma g,1)\Phi)$ attached to a Lagrangian $L_0$ with stabiliser $P$ converges
for $M'>2n$ (hypothesis (ii)), the theta integral converges under (i) (when $T$ is anisotropic
$[\auU(T)]$ is compact; otherwise this is Weil's criterion that $M'$ minus the Witt index of $T$
exceeds $n$), and the two are equal up to the volume; see also
Ichino~\cite[Introduction, p.~242]{IchinoSW}. The terms are indexed by the Lagrangians $L=\gamma^{-1}L_0$, and
$\psZ_{L_0}\circ\omega(\gamma,1)=\psZ_L$ by the characterisation of the zero term in (W2) and the
invariance of $\theta$ under rational points. In Weil's notation (the hermitian case (I$_2$) of~\cite[no.~27]{WeilSiegel}, $\varepsilon=\frac12$, $m=M'$ the rank of $T$ over $E$,
$n$ the rank of $S$, $r$ the Witt index of $T$): Th\'eor\`eme~5 assumes condition (B) of~\cite[no.~38]{WeilSiegel},
$m>2n+4\varepsilon-2$, i.e.\ $M'>2n$, and the absolute convergence of the theta integral; it also assumes a place $v$ with $U(0)_v\ne\emptyset$ ($U(0)$ is the set of elements of maximal rank whose associated form vanishes) and a subgroup $G'_v$ transitive on each $U(i)_v$; we take $G'=G$, which is transitive on each non-empty $U(i)$ by~\cite[no.~22, Proposition~3]{WeilSiegel}, and $v$ a place of $E^+$ split in $E$, where $T_v$ has Witt index $\lfloor M'/2\rfloor\ge n$, so that a map of rank $n$ onto a totally isotropic subspace exists; the Eisenstein series then converges
absolutely by~\cite[Th\'eor\`eme~1]{WeilSiegel}, and the theta integral converges absolutely when $r=0$ or
$m-r>n+2\varepsilon-1$, i.e.\ $M'-r_T>n$~\cite[Proposition~8]{WeilSiegel}, which follows from $M'>2n$ since $r_T\le M'/2$. Weil's theorems give convergence uniformly on compact sets only; the bounds by a power of the height are proved as follows. For the Eisenstein series it is Lemma~\ref{psr:Egrowth}, which needs only $M'>2n$. For anisotropic $T$, Lemma~\ref{pcx:thetagrowth} for the pair $(\pcW,T)$, with $\mathrm{ht}(h)$ bounded on a compact set of representatives of $[\auU(T)]$, gives the bound for $\int|\theta(\omega(g,h)\Phi)|\,dh$. For the product integral and any $T$ as in (SW): by the Cauchy--Schwarz inequality in $h$ it suffices to bound $\int|\theta(\omega_{\alpha,\beta}(g,h)F)|^2\,dh$ by $C\,\mathrm{ht}(g)^A$. Put $G=\omega_{\alpha,\beta}(g,1)F$; then $\omega_{\alpha,\beta}(g,h)F=\omega_{\alpha,\beta}(1,h)G$ differs from $\omega_{\auHom(S,T)}(1,h)G$ by a character of $h$ of absolute value $1$, and by (W7) and the factorisation $\theta\otimes\theta^-$ of the theta functional of $\auHom(\pcW,T)$ on pure tensors ((W1)), $\theta(\omega(1,h)\Phi_g)=|\theta(\omega_{\auHom(S,T)}(1,h)G)|^2$ for $\Phi_g=G\otimes\mathsf CG$. So $\int|\theta(\omega_{\alpha,\beta}(g,h)F)|^2\,dh=\int\theta(\omega(1,h)\Phi_g)\,dh$, an integral of a non-negative function, which converges under (i) and equals $\pcvol([\auU(T)])\sum_L\psZ_L(\Phi_g)$ by the identity. Up to a factor of absolute value $1$ (the splitting characters and the character $\lambda$ of (W7) are unitary), $\Phi_g=\omega((g,g),1)\Phi_1$ with $\Phi_1=F\otimes\mathsf CF$ and $(g,g)\in\auU(S)\times\auU(S^-)\subseteq\auU(\pcW)$. Hence the integral is at most $\pcvol([\auU(T)])\sum_L|\psZ_L(\omega((g,g),1)\Phi_1)|\le C\,\mathrm{ht}(g,g)^A=C\,\mathrm{ht}(g)^A$, by Lemma~\ref{psr:Egrowth} and the comparison of heights in the comments to Lemma~\ref{psr:SW}. (This is all that Proposition~\ref{pc:rallis} uses for isotropic $T$: absolute convergence of the theta integral, Weil's identity, the bound for $\sum_L|\psZ_L|$ in Step~1 and for the product integral in the Fubini argument of Step~2; no bound for $\int|\theta(\omega(g,h)\Phi)|\,dh$ at a general $g\in\auU(\pcW)(\auA)$ is used.) The statement is used in
this section with $M'=6n$ (Proposition~\ref{pc:rallis}) and elsewhere with a definite $T$ of rank
$M'>2n$ (Lemma~\ref{psr:SW}, Proposition~\ref{pex:rallis}, \S\ref{pgy:sec}).
\end{proof}

\begin{lemma}[Clean data]\label{pc:clean}
Let $S,T$ be as in \textup{(SW)}, $\pcW=S\oplus S^-$, and let $F\in\psHS(\auHom(S,T))$ be a datum
whose component at one split place $v=w\aubar w$ is, in the model \textup{(W5)}, a Schwartz--Bruhat
function supported on injective matrices. Let $F'\in\psHS(\auHom(S,T))$ and
$\Phi=F\otimes\mathsf CF'\in\psHS(\auHom(S,T))\otimes\psHS(\auHom(S^-,T))=\psHS(\auHom(\pcW,T))$.
\begin{enumerate}
\item[(a)] If the Lagrangian $L$ of $\pcW$ is not of the form $L_\delta=\{(\delta s,-s):s\in S\}$
for an element $\delta\in\auU(S)(E^+)$, then $\psZ_L(\Phi)=0$.
\item[(b)] There is a constant $c'\ne0$, depending only on the normalisations, such that for all
$\delta\in\auU(S)(E^+)$ and all $F,F'$
\[
\psZ_{L_\delta}(F\otimes\mathsf CF')=c'\,\bigl\langle F',\omega_{\auHom(S,T)}(\delta^{-1},1)F\bigr
\rangle .
\]
\end{enumerate}
\end{lemma}
\begin{proof}
(a) If the second projection $L\to S$ is injective it is bijective, $L=\{(\delta s,-s)\}$ for a
linear map $\delta$, and the isotropy of $L$ says that $\delta$ is an isometry. Otherwise there is
$\ell\ne0$ in $S$ with $(\ell,0)\in L$. For $z\in T(\auA)$ let $i_z\in\auHom(S,T)(\auA)\subset
\auHom(\pcW,T)(\auA)$ be the map $s\mapsto\langle\ell,s\rangle z$. It vanishes on $L$, because
$\langle\ell,s\rangle=\langle(\ell,0),(s,s')\rangle_{\pcW}=0$ for $(s,s')\in L$; so $i_z\in
I_L(\auA)$ and, by (W2), $\psZ_L(\psrho(i_z)\Phi)=\psZ_L(\Phi)$. The adjoint of $i_z$ is
$t\mapsto\langle z,t\rangle\ell$, so $\langle i_z,x\rangle=\operatorname{tr}(i_z^\dagger x)=
\langle z,x\ell\rangle_T$ for $x\in\auHom(S,T)$. Take $z$ in the component $T_{\aubar w}$ of $T$ at
$v$. By (W5), $\psrho(i_z)$ multiplies the component $F_w$ of $F$ by the function
$x\mapsto[x,i_z]$, which for fixed $x$ is the character $z\mapsto\psi_v(\tfrac12\jmath(\langle z,
x\ell\rangle-\langle x\ell,z\rangle))$ of $T_{\aubar w}$; it is non-trivial when $x\ell\ne0$, since
the form pairs $T_{\aubar w}$ and $T_w$ perfectly. On the compact support of $F_w$ the matrices
are injective, so $x\ell$ stays in a compact subset of $T_w\smallsetminus\{0\}$, and there is a
compact open subgroup $C\subset T_{\aubar w}$ with $\int_C[x,i_z]\,dz=0$ for all $x$ in the support.
Then $\int_C\psrho(i_z)\Phi\,dz=0$ (a finite sum), and $\psZ_L(\Phi)=\pcvol(C)^{-1}\int_C\psZ_L(
\psrho(i_z)\Phi)\,dz=0$.

(b) For $\Delta=L_1=\{(s,-s)\}$, $I_\Delta(\auA)$ consists of the maps $(s,s')\mapsto x(s+s')$,
that is, of the pairs $(x,x)$, and the Heisenberg operator of $(x,x)$ on the tensor product is
$\psrho(x)\otimes\psrho^-(x)$. So by (W2) and (W7),
$\psZ_\Delta(\psrho(x)F\otimes\mathsf C\psrho(x)F')=\psZ_\Delta(F\otimes\mathsf CF')$ for all $x$:
the sesquilinear form $(F',F)\mapsto\psZ_\Delta(F\otimes\mathsf CF')$ is invariant under the
Heisenberg group. As the Heisenberg representation is irreducible, it is a multiple $c'$ of the
inner product, and $c'\ne0$ because $\psZ_\Delta\ne0$. For general $\delta$,
$L_\delta=(\delta,1)\Delta$ and $\psZ_{L_\delta}=\psZ_\Delta\circ\omega((\delta,1)^{-1},1)$, as in
the proof of (SW); the element $(\delta,1)$ acts on $F\otimes\mathsf CF'$ through
$\omega_{\auHom(S,T)}(\delta,1)F$, up to a character that is trivial on rational points.
\end{proof}

\begin{definition}[Sharp functions]\label{pc:sharp}
Let $S,T$ be spaces of ranks $n,M'$, $\pcm$ a modulus, $\alpha$ a character tamely ramified at
$\pcm$, and $J$ a subgroup of $K_S$ containing $J_S(\pcm)$ and the components of $K_S$ outside
$S(\pcm)$, with image $\bar J\subseteq\auGL_n(\pcR_\pcm)$. A function $\Psi\ne0$ on
$\pcMat_{M'\times n}(\pcR_\pcm)$ is \emph{sharp for $J$} if
$\alpha_\pcm(\det\bar k)^{-1}\Psi(x\bar k)=\Psi(x)$ for $\bar k\in\bar J$ and there is a matrix
$A\in\pcMat_{n\times M'}(\pcR_\pcm)$ with $Ax\in\bar J$ for all $x$ in the support of $\Psi$. We put
$\|\Psi\|^2=\sum_x|\Psi(x)|^2$.
\end{definition}

\begin{proposition}[Inner product formula]\label{pc:rallis}
Let $(S,T)$, of ranks $r$ and $r+t$, with types $e_S,e_T$ and polynomial $P$, be a pair for which
\textup{(A1)--(A3)} of Lemma~\ref{pc:arch} hold. Let $\pcm$ be a modulus with $S(\pcm)\ne\emptyset$ \textup(i.e.\ $\pcm\ne(1)$\textup), $\alpha,\beta$ characters
as in \textup{(A2)}, and $\Psi$ sharp for $J$. Put $\Theta=\Theta_{\alpha,\beta}(P,\Psi)$ and
$\pcRal=\pcRal_{r+t,\alpha}$. Then for all $y,y'\in\pcF(S,e_S)^J$
\[
\langle\Theta y',\Theta y\rangle=A_\infty\cdot\frac{\|\Psi\|^2}{\#\pcMat_{(r+t)\times r}(\pcR_\pcm)
\cdot[K_S:J]}\cdot\langle y',\pcRal\,y\rangle,\qquad A_\infty=c'\,Z_\infty\,\pcvol([\auU(T)])\ne0,
\]
where $A_\infty$ depends only on the pair and its archimedean data. Moreover $\Theta$ is injective
on $\pcF(S,e_S)^J$.
\end{proposition}
\begin{proof}
Write $F=[P,\Psi]$ and $\omega_\alpha(g)=\omega_{\alpha,\beta}(g,1)$.

\emph{Step 1: doubling.} The component of $F$ at any $v\in S(\pcm)$ is supported on injective
matrices: if $\Psi(\bar x)\ne0$ then $A\bar x$ is invertible. So is the component of
$\omega_\alpha(g)F$ for every $g\in\auU(S)(\auA)$, since by (W5) $\omega_\alpha(g)$ acts at $v$
through $x\mapsto xg_v$ up to a scalar, and $xg_v$ is injective when $x$ is. The theta functional of
$\auHom(\pcW,T)$ is $\theta\otimes\theta^-$ on pure tensors (it is invariant under the rational
vectors; (W1)). For $g,g'\in\auU(S)(\auA)$ and $h\in\auU(T)(\auA)$, by (W7) and since
$|\beta(\det h)|=1$,
\[
\theta(\omega_{\alpha,\beta}(g,h)F)\,\aubar{\theta(\omega_{\alpha,\beta}(g',h)F)}=\theta_{\auHom(\pcW,
T)}\bigl(\omega(1,h)\,[\omega_\alpha(g)F\otimes\mathsf C\omega_\alpha(g')F]\bigr),
\]
with $\omega(1,h)$ the action of (SW). By (SW) and Lemma~\ref{pc:clean},
\begin{equation}\label{pc:eq-doubling}
\int_{[\auU(T)]}\aubar{\theta(\omega_{\alpha,\beta}(g',h)F)}\,\theta(\omega_{\alpha,\beta}(g,h)F)\,dh
=c'\pcvol([\auU(T)])\sum_{\delta\in\auU(S)(E^+)}\bigl\langle F,\omega_\alpha(g'^{-1}\delta g)F
\bigr\rangle ,
\end{equation}
with absolute convergence and a bound by a power of the heights of $g$ and $g'$.

\emph{Step 2: unfolding.} By (A1) the functions $y,y'$ decrease rapidly, so the following integrals
converge absolutely and Fubini's theorem applies:
\begin{align*}
\langle\Theta y',\Theta y\rangle&=\int_{[\auU(S)]^2}\aubar{y'(g')}y(g)\int_{[\auU(T)]}\aubar{\theta(
\omega(g',h)F)}\theta(\omega(g,h)F)\,dh\,dg\,dg'\\
&=c'\pcvol([\auU(T)])\int_{\auU(S)(\auA)}\langle F,\omega_\alpha(g)F\rangle\,\langle y',\rho(g)y
\rangle\,dg ,
\end{align*}
where the sum over $\delta$ in \eqref{pc:eq-doubling} has been combined with the integral over
$g\in[\auU(S)]$ and $g$ replaced by $g'g$.

\emph{Step 3: the three factors.} The function $g\mapsto\langle F,\omega_\alpha(g)F\rangle$ is the
product of its values on the components $g_\infty$, $g_\pcm$ (at $S(\pcm)$) and $g^\pcm$.
\emph{Outside $\pcm$}: $y$ is $K^\pcm$-invariant and $K^\pcm$ has volume $1$, so
$\int_{G^\pcm}\langle\psvac^\pcm,\omega_\alpha(g)\psvac^\pcm\rangle\rho(g)y\,dg=\pcRal y$.
\emph{At $\pcm$}: by (W5), $\omega_\alpha(g)\Psi^\natural(x)$ is a non-zero multiple of
$\Psi^\natural(xg)$. If $\langle\Psi^\natural,\omega_\alpha(g)\Psi^\natural\rangle\ne0$ there is
$x\in a\pcMat(\widehat\auO)$ with $xg\in a\pcMat(\widehat\auO)$ and the reductions of $a^{-1}x$, $a^{-1}xg$ in
the support of $\Psi$; with an integral lift $\tilde A$ of $A$, the matrices $\tilde Aa^{-1}x$ and
$\tilde Aa^{-1}xg$ are integral with reductions in $\bar J$, hence lie in $J$, and
$g=(\tilde Aa^{-1}x)^{-1}(\tilde Aa^{-1}xg)\in J$ ($a$ is a scalar). For $g\in J$,
$\omega_\alpha(g)\Psi^\natural=\Psi^\natural$ by Lemma~\ref{pc:level} and the first condition of
sharpness. So the coefficient is $\|\Psi^\natural\|^2$ times the indicator function of $J$, with
$\|\Psi^\natural\|^2=\|\Psi\|^2/\#\pcMat_{(r+t)\times r}(\pcR_\pcm)$, and the integral of
$\rho(g_\pcm)$ against it is multiplication by $\|\Psi^\natural\|^2/[K_S:J]$ on $J$-invariant
functions. \emph{At infinity}: $\pcRal y\in\pcF(S,e_S)^J$ by Lemma~\ref{pc:series}(c), and
\eqref{pc:eq-schur} with (A3) gives the factor $Z_\infty$. This proves the formula.

\emph{Step 4: injectivity.} If $\Theta y=0$, the formula with $y'=\pcRal y\in\pcF(S,e_S)^J$ gives
$\langle\pcRal y,\pcRal y\rangle=0$, so $\pcRal y=0$ and $y=0$ by Lemma~\ref{pc:series}(c).
\end{proof}

\subsubsection*{The single lifts}
Let $1\le j\le m$ and let $\vartheta$ be a twist. The lift $u=\Theta_{j,\vartheta}$ of
Proposition~\ref{au:factor} is defined as follows; the normalisation matters, because it depends on
$\vartheta$. The pair is $(W_j,V_j)$, with $\alpha=\mu_j^0\vartheta\,\upsilon_{r+t}$,
$\beta=\mu_j^0\vartheta\,\upsilon_r$, the joint harmonic $P_j^W$ and the modulus $\pcm_\vartheta$.
The finite datum is the \emph{block function} $\Psi^{\mathrm{bl}}_\vartheta$ on
$\pcMat_{(r+t)\times r}(\pcR_\vartheta)$: at a chosen place the indicator of $\iota$; at $w\in
Q_\vartheta$ the function $x\mapsto\psi_w(\det x^{\mathrm{top}})$, where $x^{\mathrm{top}}$ is the
block of the first $r$ rows and $\psi_w$ is the tame part of $\vartheta$ (extended by $0$). Let
$J^V_\vartheta\subseteq K_{V_j}$ consist of the elements whose reduction fixes the first $r$ frame
vectors at the chosen places and preserves the span of the last $t$ frame vectors at the places of
$Q_\vartheta$. By Lemma~\ref{pc:level}, $J^V_\vartheta$ acts on $[P,\Psi^{\mathrm{bl}}_\vartheta]$
by a character, so $\Theta_{\alpha,\beta}(P^W_j,\Psi^{\mathrm{bl}}_\vartheta)$ takes values in the
functions on $\auU(V_j)$ that correspond, by $\varphi\mapsto(h\mapsto\langle\xi_0,\pi^\natural(h_\infty)^{-1}
\varphi(h_f)\rangle)$ with $\xi_0$ the vector of Lemma~\ref{plf:harm} ($\pop_{\xi_0}=P^W_j$) and
$\pi^\natural$ the action of \S\ref{po:ss-set} for which $\xi\mapsto\pop_\xi$ is equivariant, to the space
$\mathcal A(V_j;\vartheta)$ of algebraic modular forms of the weight $\lambda_V$ of \S\ref{po:ss-set},
level $J^V_\vartheta$ and that character, with the inner product of \S\ref{ss:au-spaces}; here $e_{V_j}$
is the type of the representation of highest weight $\Lambda'$ generated by $P^W_j$
(Proposition~\ref{ptw:A2}). Then
\[
\Theta_{j,\vartheta}\colon e_0\mathcal A(W_j)\to\mathcal A(V_j;\vartheta),\qquad
Z_j(\vartheta)=\det(\Theta_{j,\vartheta}^*\Theta_{j,\vartheta}).
\]

\begin{remark}[The two normalisations of the single lift]\label{pc:rem-norm}
In \S\ref{po:ss-set} and \S\ref{pga:ss-L} the single lift is formed with the representation
$(g,h)\mapsto a(\det h)\,b(\det g)^{-1}\omega_{\auHom(W_j,V_j)}(g,h)$, $a=\upsilon^{r_-}$,
$b=\mu\,\upsilon^{\lfloor(r+t)/2\rfloor}$, $\mu=\mu_j^0\vartheta$, where $r_-$ is the number of negative
signs of $W_j$. On $\auU(W_j)$ this is the action used here, since $\alpha=\mu\,\upsilon_{r+t}=b$. On
$\auU(V_j)$ the two differ by the character $\chi\circ\det$, $\chi=\beta a^{-1}=\mu\,
\upsilon^{\lfloor r/2\rfloor-r_-}$, a \emph{unitary} automorphic character of $\auU(1)$. Hence the
lift $\Theta'$ of \S\ref{po:ss-set} and the lift $\Theta$ defined above are related by
$\Theta=\mathsf m_\chi\circ\Theta'$, where $\mathsf m_\chi$ is multiplication by $\chi(\det h)$; this
is an isometry from the forms on $\auU(V_j)$ of level $J^V_\vartheta$ with one character of the level
(and one weight) to those with the character multiplied by $\chi\circ\det$ (and the weight shifted
by that of $\chi$), for the inner products of \S\ref{ss:au-spaces}. Therefore
$\Theta^*\Theta=\Theta'^*\Theta'$, and $Z_j(\vartheta)$, the series of Rallis' formula and
everything in Propositions~\ref{au:inj} and~\ref{au:factor} are the same for both. The
normalisation matters only where the Hecke operators on $\auU(V_j)$ or the integrality of forms on
$\auU(V_j)$ are used: these are Proposition~\ref{pga:L}, Lemma~\ref{po:facts} and
\S\ref{po:sec}, which all use $\Theta'$; under $\mathsf m_\chi$ every root of the Hecke polynomial at
a split place $w$ is multiplied by the value of the base change of $\chi$ at $w$.
\end{remark}

\begin{lemma}[Algebraic forms and the Petersson product]\label{pc:algpet}
Let $S$ be a definite space, $\lambda$ a weight with highest weight vector $v^+$, $J\subseteq K_S$
open and $\chi$ a character of $J$ of finite order. For $\varphi\in\mathcal A_\lambda(J,\chi)$ let
$f_\varphi(h)=\langle v^+,h_\infty^{-1}\varphi(h_f)\rangle$. Then
$\langle\varphi',\varphi\rangle=[K_S:J]\cdot\dim W_\lambda\cdot\|v^+\|^{-2}\cdot\langle f_{\varphi'},
f_\varphi\rangle_{[\auU(S)]}$. The statement and proof hold for every unit vector in place of $v^+$.
\end{lemma}
\begin{proof}
$[\auU(S)]$ is the disjoint union, over $x\in\Gamma\backslash G_f/J$, of the images of
$\auU(S)_\infty\times xJ$, of volume $(\#(\Gamma\cap xJx^{-1})\,[K_S:J])^{-1}$, and
$|\chi|=1$. On each of them Schur orthogonality for the irreducible representation $W_\lambda$ of
the compact group $\auU(S)_\infty$, of total mass $1$, gives
$\int\aubar{\langle v^+,u^{-1}\xi'\rangle}\langle v^+,u^{-1}\xi\rangle\,du=\|v^+\|^2\langle\xi',\xi
\rangle/\dim W_\lambda$.
\end{proof}

\begin{proposition}[Norm of a single lift]\label{pc:single}
For $1\le j\le m$ there is $A_j\in\C^\times$ such that for every twist
$\vartheta$
\[
\Theta_{j,\vartheta}^*\Theta_{j,\vartheta}=A_j\,c_Q(\vartheta)\,\pcRal_j(\vartheta)\ \text{ on }
e_0\mathcal A(W_j),\qquad c_Q(\vartheta)=(\auNm Q_\vartheta)^{rt}\prod_{w\in Q_\vartheta}\
\prod_{i=t+1}^{t+r}\bigl(1-q_w^{-i}\bigr).
\]
Consequently $\Theta_{j,\vartheta}$ is injective, and
$Z_j(\vartheta)=\bigl(A_jc_Q(\vartheta)\bigr)^{d_j}\det\bigl(\pcRal_j(\vartheta)\,|\,e_0\mathcal A
(W_j)\bigr)>0$, $d_j=\dim e_0\mathcal A(W_j)$.
\end{proposition}
\begin{proof}
The function $\Psi^{\mathrm{bl}}_\vartheta$ is sharp for $J=J_M$ (trivial reduction at the chosen
places, all of $\auGL_r(\auF_w)$ at $w\in Q_\vartheta$): the first condition holds at $w\in Q$
because the tame part of $\alpha$ at $w$ is $\psi_w$ and
$\psi_w(\det\bar k)^{-1}\psi_w(\det(x\bar k)^{\mathrm{top}})=\psi_w(\det x^{\mathrm{top}})$; for the
second take $Ax=x^{\mathrm{top}}$. Write $g_N=\#\auGL_N(\auF_w)$ and $q=q_w$. Then
$\|\Psi^{\mathrm{bl}}_\vartheta\|^2=\prod_{w\in Q}g_rq^{rt}$, $[K_{W_j}:J_M]=\prod_{w\text{ chosen}}
g_r$, and Proposition~\ref{pc:rallis} gives, for $y,y'\in\mathcal A(W_j)$,
\[
\langle\Theta y',\Theta y\rangle_{[\auU(V_j)]}=A_\infty\prod_{w\text{ chosen}}\frac1{g_rq^{(r+t)r}}
\prod_{w\in Q}\frac{g_rq^{rt}}{q^{(r+t)r}}\cdot\langle y',\pcRal_j(\vartheta)y\rangle .
\]
By Lemma~\ref{pc:algpet} the inner product in $\mathcal A(V_j;\vartheta)$ is this times a constant
and $[K_{V_j}:J^V_\vartheta]$, which is a factor independent of $\vartheta$ at the chosen places
times $\prod_{w\in Q}g_{r+t}/(g_rg_tq^{rt})$, the index of a parabolic subgroup with blocks $r,t$.
At $w\in Q$ the product of the factors is
$g_{r+t}/(g_tq^{(r+t)r})=q^{rt}\prod_{i=t+1}^{t+r}(1-q^{-i})$. The operator $\pcRal_j(\vartheta)$
preserves $e_0\mathcal A(W_j)$ (Lemma~\ref{pc:heckeclosure}); this gives the identity. Injectivity
is that of Proposition~\ref{pc:rallis}, and the formula for $Z_j$ is Lemma~\ref{pc:vol}; $Z_j>0$
since $\Theta_{j,\vartheta}$ is injective.
\end{proof}

\subsubsection*{Inner forms}
Proposition~\ref{au:mult} is used in the following form.
\begin{quote}
$(\mathrm M^+)$\quad For $1\le j\le m$ and every $h\in\pcH$,\quad
$\pctr\bigl(h\,|\,e_0\mathcal A(W_j)\bigr)=\pctr\bigl(h\,|\,\auT\bigr)$.
\end{quote}
With $h=1$ it gives $\dim e_0\mathcal A(W_j)=\audT$. It involves all places prime to $\pcm_\vartheta$ (as the series $\pcRal_j(\vartheta)$ does), not
only the split ones; this is the statement of Proposition~\ref{au:mult}, proved in
Proposition~\ref{pm:mult} and Remark~\ref{pm:rem-full}.

\begin{lemma}\label{pc:detcompare}
Assume $(\mathrm M^+)$. Then for $1\le j\le m$ and every twist $\vartheta$,
$\det\bigl(\pcRal_j(\vartheta)\,|\,e_0\mathcal A(W_j)\bigr)=\det\bigl(\pcRal_j(\vartheta)\,|\,\auT
\bigr)$.
\end{lemma}
\begin{proof}
By Lemma~\ref{pc:heckeclosure} there is $h\in\pcH$ acting as $\pcRal_j(\vartheta)$ on $\mathcal A
(W_j)$ and on $\auAz$, hence on the images of $\varrho(e_0)$. By $(\mathrm M^+)$ for the powers
$h^k$, $k\ge0$, the two restrictions have the same dimension and the same power sums of
eigenvalues, hence (Newton's identities) the same characteristic polynomial.
\end{proof}


\subsection{The triangle of lifts; proofs of Propositions~\ref{au:inj} and~\ref{au:factor}}
\label{pc:ss-triangle}

\subsubsection*{The spaces and the lifts}
For $i,n\ge0$ with $i+n\le m$ let $\pcG_{i,n}$ be the following space of rank $r+nt$:
$\pcG_{0,n}=Y_n$ with the lattice of \S\ref{ps:ss-Y}; $\pcG_{i,0}$ is $Y_0$ with its form multiplied
by $(-1)^i$; and for $i,n\ge1$, $\pcG_{i,n}=E^{r+nt}$ with the standard lattice and the diagonal
form with entries $\epsilon_1,\dots,\epsilon_r,1,\dots,1$, where
$\epsilon_p=(-1)^{i+a+\min((a-i)^+,\,n)}$, $a=(p-k)^+$. (For $i=0$ these are the signs of $Y_n$.)
Thus $\pcG_{i,n}$ has a \emph{stem} of rank $r$ and $n$ positive definite \emph{blocks} of rank
$t$, numbered $i+1,\dots,i+n$; at most $r$ of its signs are negative. Let $e_{i,n}$ be the type of
the chain at the node $(i,n)$: for $i=0$ the type of the lower chain on $Y_n$, in general the type
whose Harish-Chandra parameter at $\tau$ consists of $\varsigma_\tau$ on the stem and, for the
block $i+l$, of the $t$ numbers $-\kappa_{i+l}(\tau)+\tfrac{t-1}2-b$, $0\le b<t$, all lowered by
$\tfrac12$ if $n$ is odd; and put $\pcF_{i,n}=\pcF(\pcG_{i,n},e_{i,n})$. The groups of
$\pcG_{i,0}$ and of $Y_0$ coincide, and so do the types, so that $\pcF_{i,0}=\pcF_{0,0}\supseteq
\auAz$; let $j_i\colon\auAz\to\pcF_{i,0}$ be the inclusion.

For an $m$-fold twist $\vartheta$ put $\mu_j=\mu_j^0\vartheta^j$,
$w_n^\vartheta=\#\auGL_{r+nt}(\pcR_\vartheta)$, and for $i+n<m$
\[
a_{i,n}^\vartheta=\Theta_{\alpha,\beta}(P_{i,n},\delta_\iota)\colon\pcF_{i,n}\to\pcF_{i,n+1},\qquad
\alpha=\mu_{i+n+1}\upsilon_{r+(n+1)t},\ \ \beta=\mu_{i+n+1}\upsilon_{r+nt},
\]
the lift for the pair $(\pcG_{i,n},\pcG_{i,n+1})$ with the modulus $\pcm_\vartheta$, the point datum
at the inclusion matrix and the joint harmonic $P_{i,n}$ of the two types. For $i=0$ this is
the lift of the lower chain without its normalisation: $[K_n:J_n]=w_n^\vartheta$ and
\begin{equation}\label{pc:eq-iota}
\Theta^Y_{n,\vartheta}=w_n^\vartheta\,a_{0,n}^\vartheta,\qquad\iota_\vartheta=X^\vartheta_{0,m}
\circ j_0 .
\end{equation}
We use the following property of the data (\S\ref{ss:au-spaces}):
$\mu_j^0=\mu^\flat\cdot\nu_1^{\,j}$ with $\nu_1=\upsilon^{-\Delta}\chi_M$ on $\auU(1)$, where
$\chi_M$ has weight $0$, is unramified outside the chosen places, and its tame parts at the chosen
places have order greater than $2m$. Put $\nu_n=\nu_1^{\,n}$, of weight $-n\Delta$. For
$1\le z\le t$ and $1\le n\le 2m$ the pair $(\nu_n,z)$ is a ramified critical pair in the sense of
\S\ref{ss:au-proofC} when $n\Delta\le4b$, which holds for $n\le m$. For $0<i<j\le m$ let
\begin{equation}\label{pc:eq-rho}
\rho_{ij}(\vartheta)=\prod_{z=1}^t\frac{|L(z,(\nu_{j-i}\vartheta^{\,j-i})^{-1})|}{|L(z,\nu_{j-i}^{-1})
|},
\end{equation}
a ramified ratio; the $L$-functions are those of \S\ref{ss:au-proofC}, without the Euler factors at
the primes dividing $\pcm_\vartheta\aubar{\pcm_\vartheta}$, respectively $M$. We write
$L_*(\pcm,\nu,e)=(\prod_{w\mid\pcm}q_w)^{e^2}\prod_{z=1}^e|L^{\pcm}(z,\nu)|^2$ (equal to
$(\auNm\pcm)^{e^2}\prod_{z=1}^e|L^{\pcm}(z,\nu)|^2$ for squarefree $\pcm$), with $L^\pcm$ the Hecke
$L$-function of $x\mapsto\nu(x/\aubar x)$ without the factors at $\pcm\aubar\pcm$.

\begin{theorem}[The lifts form a triangle]\label{pc:tri}
The maps $a^\vartheta_{i,n}$ are well defined and have adjoints, and
there are linear maps $d^\vartheta_{i,n}\colon\pcF_{i,n+1}\to\pcF_{i+1,n}$ $(i+n<m)$, for all
$m$-fold twists $\vartheta$, such that:
\begin{enumerate}
\item[(a)] $d^\vartheta_{i,0}$ is the adjoint of $a^\vartheta_{i,0}$, under $\pcF_{i+1,0}=\pcF_{i,0}$;
\item[(b)] for $i+n+2\le m$ there are $\kappa^\sharp_{i,n}\in\C$ and $C_{i,n}>0$, independent of
$\vartheta$, and numbers $c_{i,n}(\vartheta)\ne0$ with
\[
w_n\,a_{i+1,n}d_{i,n}=\frac{\kappa^\sharp_{i,n}}{c_{i,n}(\vartheta)}\,w_{n+2}\,d_{i,n+1}a_{i,n+1},
\qquad|c_{i,n}(\vartheta)|^2=C_{i,n}\,L_*\bigl(\pcm_\vartheta,(\nu_{n+1}\vartheta^{\,n+1})^{-1},t
\bigr);
\]
\item[(c)] for $i+n+2\le m$ the adjoint law \textup{(T4)} holds with functions $\kappa_{i,n}$ that
have modulus $1$ up to norm powers.
\end{enumerate}
\end{theorem}
\noindent Lemma~\ref{pc:arch} holds: it is Proposition~\ref{ptv:arch}. The maps are well defined with adjoints (Theorem~\ref{ptr:edge}(d), Lemma~\ref{ptv:adj});
(a) is Proposition~\ref{ptv:tri}(ii); (b) is Corollary~\ref{ptr:tri-b} with Theorem~\ref{ptq:H} and
Proposition~\ref{ptv:chars}; (c) is Theorem~\ref{ptv:adjoint}, which uses the square law (RS)
of the reversed squares, Theorem~\ref{ptw:RS} (Corollary~\ref{ptw:adjoint}).

\begin{proposition}[First lifts]\label{pc:first}
For $0\le i<m$ there is $C_i\in\C^\times$ such that for every $m$-fold
twist $\vartheta$ and all $x,y\in\auAz$
\[
w_0^\vartheta\,\langle a^\vartheta_{i,0}x,a^\vartheta_{i,0}y\rangle=\frac{C_i}{\#\pcMat_{(r+t)
\times r}(\pcR_\vartheta)}\,\bigl\langle x,\pcRal_{i+1}(\vartheta^{\,i+1})\,y\bigr\rangle ,
\]
and $a^\vartheta_{i,0}$ is injective on $\auAz$. Moreover, for every $u\in J_M$ there are
$g\in K_{\pcG_{i,1}}$ and $c\in\C$, $|c|=1$, with $a_{i,0}\rho(u)=c\,\rho(g)a_{i,0}$.
\end{proposition}
\begin{proof}
We use Lemma~\ref{pc:arch}, which is Proposition~\ref{ptv:arch}.
The point function $\delta_\iota$ on $\pcMat_{(r+t)\times r}(\pcR_\vartheta)$ is sharp for
$J=J_{\pcG_{i,0}}(\pcm_\vartheta)$, with $A$ the projection to the first $r$ rows: $A\iota=1$. The
elements of $\auAz$ are invariant under $J_M\supseteq J$, where $J_M$ is trivial at the chosen
places and maximal elsewhere. Here $\|\delta_\iota\|^2=1$ and $[K:J]=w_0^\vartheta$, the character
is $\alpha=\mu_{i+1}\upsilon_{r+t}=\mu^0_{i+1}\vartheta^{\,i+1}\upsilon_{r+t}$, and the modulus of
$\vartheta^{\,i+1}$ is $\pcm_\vartheta$; so Proposition~\ref{pc:rallis} gives the formula, with
$C_i=A_\infty$, and the injectivity. (By Lemma~\ref{pc:series}(a) the series is the one of
\eqref{pc:eq-Rj}, although $\pcG_{i,0}$ and $\pcG_{i,1}$ are not $W_{i+1}$ and $V_{i+1}$.) For the
last assertion choose $g$ with reduction $\operatorname{diag}(\bar u,1)$ modulo $\pcm_\vartheta$;
then $\bar g\iota=\iota\bar u$, and by Lemma~\ref{pc:level} the pair $(u^{-1},g^{-1})$ multiplies
the datum $[P,\delta_\iota]$ by a root of unity $c$. Hence
$\theta(g'u^{-1},h;P,\delta_\iota)=c\,\theta(g',hg;P,\delta_\iota)$, and the change of variable
$g'\mapsto g'u$ in the integral defining $a_{i,0}\rho(u)f$ gives the claim.
\end{proof}

\begin{corollary}\label{pc:istriangle}
Let $\auAz\ne0$. Then the spaces
$\pcF_{i,n}$, the maps $j_i$, $a^\vartheta_{i,n}$, $d^\vartheta_{i,n}$ and the weights
$w^\vartheta_n$ form a triangle of size $m$ over the set of $m$-fold twists, with $V=\auAz$ and
with the adjoint law. In it $\sigma_{i,n}=\kappa^\sharp_{i,n}/c_{i,n}$ has no zero,
$\kappa^\sharp_{i,n}\ne0$, and $\beta_{i,n}=1/\aubar{\sigma_{i,n}}$ has modulus $\rho_{i+1,\,i+n+2}$
up to norm powers.
\end{corollary}
\begin{proof}
(T1) and (T4) are (b) and (c). (T2) is (a) restricted to $j_i(\auAz)$. For (T3) let $x\in\auAz$ and
$z=d_{i,0}a_{i,0}x\in\pcF_{i+1,0}=\pcF_{i,0}$. By (a), $\langle z,f\rangle=\langle a_{i,0}x,a_{i,0}f
\rangle$ for all $f\in\pcF_{i,0}$. For $u\in J_M$ let $g,c$ be as in Proposition~\ref{pc:first}
for $u^{-1}$, so that $a_{i,0}\rho(u^{-1})=c\rho(g)a_{i,0}$ and hence
$\rho(g^{-1})a_{i,0}=c\,a_{i,0}\rho(u)$. As $\rho$ is unitary and $\rho(u)x=x$,
\[
\langle\rho(u)z,f\rangle=\langle a_{i,0}x,a_{i,0}\rho(u^{-1})f\rangle=c\langle\rho(g^{-1})a_{i,0}x,
a_{i,0}f\rangle=c\,\aubar c\,\langle a_{i,0}x,a_{i,0}f\rangle=\langle z,f\rangle .
\]
So $z\in\pcF_{i+1,0}^{J_M}=j_{i+1}(\auAz)$. If $z=0$ then $\langle a_{i,0}x,a_{i,0}x\rangle=0$ and
$x=0$ by Proposition~\ref{pc:first}. Thus we have a triangle, and Proposition~\ref{pc:tri-inj}
shows that $\sigma_{i,n}$ has no zero. Finally $|\beta_{i,n}|^2=|c_{i,n}|^2/|\kappa^\sharp_{i,n}|^2$,
and
\[
L_*\bigl(\pcm_\vartheta,(\nu_{n+1}\vartheta^{\,n+1})^{-1},t\bigr)=M^{ft^2}(\auNm Q_\vartheta)^{t^2}
\prod_{z=1}^t|L(z,\nu_{n+1}^{-1})|^2\cdot\rho_{i+1,\,i+n+2}(\vartheta)^2
\]
by \eqref{pc:eq-rho}, because $\auNm\pcm_\vartheta=M^f\auNm Q_\vartheta$; the values
$L(z,\nu_{n+1}^{-1})$ do not depend on $\vartheta$ and are not zero since $c_{i,n}(1)\ne0$.
\end{proof}

\begin{proof}[Proof of Proposition~\ref{au:inj} from Lemma~\ref{pc:arch} and
Theorem~\ref{pc:tri}\textup{(a)}, \textup{(b)}]
We may assume $\auAz\ne0$. Only (T1)--(T3) are needed, that is, parts (a) and (b) of
Theorem~\ref{pc:tri} (in (b) only the existence of the constants) and Proposition~\ref{pc:first}.
By Proposition~\ref{pc:tri-inj} with $(i,n)=(0,m)$ the map $X^\vartheta_{0,m}\circ j_0$ is
injective on $\auAz$, and it is $\iota_\vartheta$ by \eqref{pc:eq-iota}.
\end{proof}

\begin{lemma}[Indices]\label{pc:index}
Let $J_\vartheta\subseteq K_m$ be the level of $\auAm(\vartheta)$ and let
$e_m=mrt+\tfrac12m(m-1)t^2$. There is a constant $C_M>0$, depending only on the chosen places, such
that for every $m$-fold twist $\vartheta$
\[
[K_m:J_\vartheta]\cdot\frac{w_0^\vartheta\,(w_1^\vartheta)^m}{w_m^\vartheta}\cdot\frac1{\bigl(\#
\pcMat_{(r+t)\times r}(\pcR_\vartheta)\bigr)^m}=C_M\,(\auNm Q_\vartheta)^{-e_m}\,c_Q(\vartheta)^m .
\]
\end{lemma}
\begin{proof}
All four quantities are products over the primes $w\mid\pcm_\vartheta$, and the factors at the
chosen places do not depend on $\vartheta$. Let $w\in Q_\vartheta$, $q=q_w$, $g_N=\#\auGL_N(
\auF_w)$. The level $J_\vartheta$ is block upper triangular at $w$, with blocks of sizes
$r,t,\dots,t$ (\S\ref{ss:au-spaces}, Lemma~\ref{pz:lastlevel}; it is not the principal congruence
subgroup $J_m$ of \S\ref{ps:ss-Y}, see Lemma~\ref{pgr:jcirc}), so its index is $g_{r+mt}/(g_rg_t^{\,m}q^{e_m})$, where $e_m$ is the dimension of the
unipotent radical. The factor at $w$ is therefore
\[
\frac{g_{r+mt}}{g_rg_t^{\,m}q^{e_m}}\cdot\frac{g_r\,g_{r+t}^{\,m}}{g_{r+mt}}\cdot q^{-(r+t)rm}=
q^{-e_m}\Bigl(\frac{g_{r+t}}{g_t\,q^{(r+t)r}}\Bigr)^m=q^{-e_m}\Bigl(q^{rt}\prod_{i=t+1}^{t+r}(1-q^{-i})
\Bigr)^m .\qedhere
\]
\end{proof}

\begin{proof}[Proof of Proposition~\ref{au:factor} from Lemma~\ref{pc:arch},
Theorem~\ref{pc:tri} and $(\mathrm M^+)$]
The numbers $Z_j(\vartheta)$ are those defined before Lemma~\ref{pc:algpet}; they are positive by
Proposition~\ref{pc:single}. If $\audT=0$ there is nothing to prove; let $\audT>0$.

\emph{Step 1: the Gram form of the chain.} By Corollary~\ref{pc:istriangle} and
Proposition~\ref{pc:first}, Proposition~\ref{pc:tri-gram} applies with
$S_i=\pcRal_{i+1}(\vartheta^{\,i+1})$ and
$C_i(\vartheta)=C_i/\#\pcMat_{(r+t)\times r}(\pcR_\vartheta)$. With \eqref{pc:eq-iota}, for
$x,y\in\auAz$,
\[
\langle\iota_\vartheta x,\iota_\vartheta y\rangle_{[\auU(Y_m)]}=\frac{w_0w_1^{\,m}}{w_m}\cdot
\frac{\prod_{i<m}C_i}{(\#\pcMat_{(r+t)\times r}(\pcR_\vartheta))^m}\cdot\varkappa(\vartheta)\cdot\bigl
\langle x,\pcRal_1(\vartheta)\pcRal_2(\vartheta^2)\cdots\pcRal_m(\vartheta^m)\,y\bigr\rangle,
\]
where $\varkappa=\prod_{i+n+2\le m}\kappa_{i,n}\beta_{i,n}$. By Lemma~\ref{pc:mod}(a),
Theorem~\ref{pc:tri}(c) and Corollary~\ref{pc:istriangle}, $\varkappa$ has modulus
$\prod_{0<i<j\le m}\rho_{ij}$ up to norm powers: the square $(i,n)$ contributes
$\rho_{i+1,i+n+2}$, and $(i,n)\mapsto(i+1,i+n+2)$ is a bijection onto the pairs $0<i'<j'\le m$.

\emph{Step 2: the inner product of $\auAm(\vartheta)$.} By Lemma~\ref{pc:algpet} for the definite
space $Y_m$ and the level $J_\vartheta$, the inner product of $\auAm(\vartheta)$ is
$c_\infty[K_m:J_\vartheta]$ times the Petersson product, with $c_\infty>0$ independent of
$\vartheta$. With Lemma~\ref{pc:index}, for $x,y\in\auT$,
\[
\langle\Theta_\vartheta x,\Theta_\vartheta y\rangle=\Lambda(\vartheta)\,c_Q(\vartheta)^m\,\bigl
\langle x,\pcRal_1(\vartheta)\cdots\pcRal_m(\vartheta^m)\,y\bigr\rangle,\qquad
\Lambda(\vartheta)=c_\infty C_M\prod_{i<m}C_i\cdot(\auNm Q_\vartheta)^{-e_m}\varkappa(\vartheta).
\]
The operators $\pcRal_j(\vartheta^j)$ preserve $\auT$ (Lemma~\ref{pc:heckeclosure}). By
Lemma~\ref{pc:vol},
\[
Z(\vartheta)=\bigl(\Lambda(\vartheta)c_Q(\vartheta)^m\bigr)^{\audT}\prod_{j=1}^m\det\bigl(\pcRal_j
(\vartheta^j)\,|\,\auT\bigr).
\]

\emph{Step 3: the single lifts.} By $(\mathrm M^+)$, $d_j=\audT$. By Proposition~\ref{pc:single}
for the twist $\vartheta^j$, whose set of places is $Q_\vartheta$, and by
Lemma~\ref{pc:detcompare}, $Z_j(\vartheta^j)=(A_jc_Q(\vartheta))^{\audT}\det(\pcRal_j(\vartheta^j)\,|\,
\auT)$. Dividing,
\[
\frac{Z(\vartheta)}{\prod_jZ_j(\vartheta^j)}=\Bigl(\frac{\Lambda(\vartheta)}{\prod_jA_j}\Bigr)^{\audT}
\]
for every $m$-fold $\vartheta$; the left side is positive, so $\Lambda$ has no zero and the left
side is $|\Lambda(\vartheta)/\prod_jA_j|^{\audT}$.

\emph{Step 4: the modulus.} The function $\Lambda$ has modulus $\prod_{i<j}\rho_{ij}$ up to norm
powers (Step 1 and Lemma~\ref{pc:mod}(a)), this product is non-negative and equal to $1$ at
$\vartheta=1$, and $\Lambda(1)\ne0$. By Lemma~\ref{pc:mod}(b) there is $e\in\Z$ with
$|\Lambda(\vartheta)/\Lambda(1)|=(\auNm Q_\vartheta)^{e/2}\prod_{i<j}\rho_{ij}(\vartheta)$. Dividing
the identity of Step 3 by its case $\vartheta=1$ and inverting gives
\[
\frac{Z(1)}{Z(\vartheta)}=\Bigl((\auNm Q_\vartheta)^{e/2}\prod_{0<i<j\le m}\rho_{ij}(\vartheta)
\Bigr)^{-\audT}\prod_{j=1}^m\frac{Z_j(1)}{Z_j(\vartheta^j)} .\qedhere
\]
\end{proof}


\subsection{The squares: matched functions and the exchange identity}\label{pc:ss-exchange}

This subsection contains the algebra of Theorem~\ref{pc:tri} at the places of the modulus, the
statement of the exchange identity on which the rest depends (proved in \S\ref{pcx:sec}), and an
outline of the deduction of Theorem~\ref{pc:tri}, which is carried out in
\S\S\ref{ptr:sec}--\ref{ptw:sec}.

\subsubsection*{Matched functions over a finite ring}
Let $\pcR$ be a finite commutative ring and $\chi$ a character of $\pcR^\times$, extended by $0$.
Let the finite set $\mathsf N$ be the disjoint union of $\mathsf B$ and $\mathsf A$, of
cardinalities $b$ and $a$. For $k\in\auGL_{\mathsf N}(\pcR)$ let $y_k\in\pcMat_{\mathsf N\times
\mathsf B}(\pcR)$ be the block of the columns $\mathsf B$ of $k$, and $x_k\in\pcMat_{\mathsf A
\times\mathsf N}(\pcR)$ the block of the rows $\mathsf A$ of $k^{-1}$; so $x_ky_k=0$. Functions
$\Phi$ on $\pcMat_{\mathsf A\times\mathsf N}(\pcR)$ and $\Psi$ on $\pcMat_{\mathsf N\times\mathsf B}
(\pcR)$ are \emph{matched} (for $\chi$) if $\Phi$ vanishes off the $x_k$, $\Psi$ vanishes off the
$y_k$, and $\Phi(x_k)=\chi(\det k)\Psi(y_k)$ for all $k$. Functions $L$ on $\pcMat_{\mathsf A\times
\mathsf N}(\pcR)$ and $L'$ on $\pcMat_{\mathsf N\times\mathsf B}(\pcR)$ \emph{satisfy the square law
with constant $c$} if $L(hx)=\chi(\det h)L(x)$ for $h\in\auGL_{\mathsf A}(\pcR)$,
$L'(yh)=\chi(\det h)L'(y)$ for $h\in\auGL_{\mathsf B}(\pcR)$, and
\[
\sum_y\Psi(y)L'(y)=c\sum_x\Phi(x)L(x)\qquad\text{for all matched }\Phi,\Psi .
\]

\begin{lemma}[The law at a point]\label{pc:pointlaw}
Let $L,L'$ satisfy the square law with constant $c$. For $w\in\pcMat_{\mathsf A\times\mathsf B}
(\pcR)$ let $\binom{1}{-w}\in\pcMat_{\mathsf N\times\mathsf B}$ have the blocks $1$ (rows
$\mathsf B$) and $-w$ (rows $\mathsf A$), and $(w\ 1)\in\pcMat_{\mathsf A\times\mathsf N}$ the blocks
$w$ (columns $\mathsf B$) and $1$ (columns $\mathsf A$). Then
\[
\#\auGL_{\mathsf B}(\pcR)\cdot L'\tbinom{1}{-w}=c\cdot\#\auGL_{\mathsf A}(\pcR)\cdot L(w\ 1).
\]
\end{lemma}
\begin{proof}
Let $k_w=\bigl(\begin{smallmatrix}1&0\\-w&1\end{smallmatrix}\bigr)$, so that
$y_{k_w}=\binom1{-w}$, $x_{k_w}=(w\ 1)$ and $\det k_w=1$. Define $\Psi_w(y_{k_w}h)=\chi(\det h)^{-1}$
for $h\in\auGL_{\mathsf B}(\pcR)$ and $\Psi_w=0$ elsewhere, $\Phi_w(gx_{k_w})=\chi(\det g)^{-1}$ for
$g\in\auGL_{\mathsf A}(\pcR)$ and $\Phi_w=0$ elsewhere. They are matched: if $y_k=y_{k_w}h$, then
$k=k_w\bigl(\begin{smallmatrix}h&B\\0&g\end{smallmatrix}\bigr)$ for some $B$ and some
$g\in\auGL_{\mathsf A}(\pcR)$, so $\det k=\det h\det g$ and $x_k=g^{-1}x_{k_w}$; then
$\Phi_w(x_k)=\chi(\det g)=\chi(\det k)\Psi_w(y_k)$. Conversely, $gx_{k_w}=x_k$ for
$k=k_w\bigl(\begin{smallmatrix}1&0\\0&g^{-1}\end{smallmatrix}\bigr)$, and if $x_k=gx_{k_w}$ for some
$k$, then the rows $\mathsf A$ of $k^{-1}k_w$ are $g(0\ 1)=(0\ g)$, so
$k^{-1}k_w=\bigl(\begin{smallmatrix}P_1&P_2\\0&g\end{smallmatrix}\bigr)$ with $\det P_1\det g=\det(k^{-1}k_w)$
a unit, hence $P_1\in\auGL_{\mathsf B}(\pcR)$; then $k=k_w\bigl(\begin{smallmatrix}P_1&P_2\\0&g
\end{smallmatrix}\bigr)^{-1}$ and $y_k=y_{k_w}P_1^{-1}$. So $\Phi_w$ vanishes off the $x_k$, $\Psi_w$
vanishes off the $y_k$, and $x_k\in\auGL_{\mathsf A}(\pcR)x_{k_w}$ if and only if $y_k\in
y_{k_w}\auGL_{\mathsf B}(\pcR)$; in particular $\Phi_w(x_k)=0=\Psi_w(y_k)$ for all other $k$. By the two transformation laws, $\sum_y\Psi_w(y)L'(y)=\#\auGL_{\mathsf B}(\pcR)L'(y_{k_w})$
and $\sum_x\Phi_w(x)L(x)=\#\auGL_{\mathsf A}(\pcR)L(x_{k_w})$.
\end{proof}

The next lemma is the form in which a square of the triangle is used. Let $\mathsf N_0$,
$\mathsf T$, $\mathsf T'$ be finite sets, $\mathsf N_1=\mathsf N_0\sqcup\mathsf T$ (the index set of
$U_1$), $\mathsf U_2=\mathsf N_0\sqcup\mathsf T'$ (that of $U_2$) and $\mathsf A=\mathsf N_1\sqcup
\mathsf T'$ (that of $A$); $\mathsf N_0$ is the index set of $A'$. Take $\mathsf N=\mathsf N_1
\sqcup\mathsf U_2$ and $\mathsf B=\mathsf N_0\subset\mathsf U_2$, so that $\mathsf N=\mathsf B\sqcup
\mathsf A$. For a function $p$ on $\pcMat_{\mathsf N_1\times\mathsf N_0}(\pcR)$ its \emph{extension}
$p^+$ is the function on $\pcMat_{\mathsf A\times\mathsf U_2}(\pcR)$ with
\[
p^+\begin{pmatrix}\eta&s\\0&1\end{pmatrix}=\frac{p(-\eta)}{\#\pcMat_{\mathsf N_1\times\mathsf T'}
(\pcR)}\qquad(\eta\in\pcMat_{\mathsf N_1\times\mathsf N_0},\ s\in\pcMat_{\mathsf N_1\times\mathsf
T'}),
\]
the rows being $\mathsf N_1,\mathsf T'$ and the columns $\mathsf N_0,\mathsf T'$, and $p^+=0$ at all
other matrices.

\begin{lemma}[Extension of a datum]\label{pc:next}
Let $L$ on $\pcMat_{\mathsf A\times\mathsf N}(\pcR)$ and $L'$ on $\pcMat_{\mathsf N\times\mathsf
N_0}(\pcR)$ satisfy the square law with constant $c$. Let $\iota'\in\pcMat_{\mathsf A\times\mathsf
N_1}$ and $\iota\in\pcMat_{\mathsf U_2\times\mathsf N_0}$ be the inclusion matrices. Then for every
$p$
\[
\#\auGL_{\mathsf N_0}(\pcR)\sum_{y\in\pcMat_{\mathsf N_1\times\mathsf N_0}}p(y)\,L'\binom{y}{\iota}=
c\cdot\#\auGL_{\mathsf A}(\pcR)\sum_{x\in\pcMat_{\mathsf A\times\mathsf U_2}}p^+(x)\,L(\iota'\ x),
\]
where $\binom y\iota$ has the rows $\mathsf N_1$ and $\mathsf U_2$, and $(\iota'\ x)$ the columns
$\mathsf N_1$ and $\mathsf U_2$.
\end{lemma}
\begin{proof}
Write matrices with rows and columns in the order $\mathsf N_1,\mathsf N_0,\mathsf T'$. For
$\eta\in\pcMat_{\mathsf N_1\times\mathsf N_0}$ let $w_\eta\in\pcMat_{\mathsf A\times\mathsf N_0}$
have the blocks $\eta$ (rows $\mathsf N_1$) and $0$ (rows $\mathsf T'$). Then
$\binom1{-w_\eta}=\binom{-\eta}{\iota}$, and $(w_\eta\ 1)=\bigl(\begin{smallmatrix}1&\eta&0\\0&0&1
\end{smallmatrix}\bigr)$. For $s\in\pcMat_{\mathsf N_1\times\mathsf T'}$ the matrix
$u_s=\bigl(\begin{smallmatrix}1&s\\0&1\end{smallmatrix}\bigr)\in\auGL_{\mathsf A}(\pcR)$ has
determinant $1$, and $u_s(w_\eta\ 1)=\bigl(\begin{smallmatrix}1&\eta&s\\0&0&1\end{smallmatrix}\bigr)
=(\iota'\ x)$ with $x=\bigl(\begin{smallmatrix}\eta&s\\0&1\end{smallmatrix}\bigr)$. So
$L(\iota'\ x)=L(w_\eta\ 1)$ for all $s$, and by the definition of $p^+$ and
Lemma~\ref{pc:pointlaw}
\[
c\,\#\auGL_{\mathsf A}\sum_xp^+(x)L(\iota'\ x)=c\,\#\auGL_{\mathsf A}\sum_\eta p(-\eta)L(w_\eta\ 1)=
\#\auGL_{\mathsf N_0}\sum_\eta p(-\eta)L'\tbinom{-\eta}{\iota}.\qedhere
\]
\end{proof}

\subsubsection*{The exchange identity}
\begin{theorem}[Exchange identity]\label{pc:X}
Let $W,V,V'$ be spaces of ranks $N=a+b$, $a$, $b$, with $N$ even and $a=b+2e$, $e\ge0$. Let
$q,q_V,q_{V'}$ be their numbers of negative signs, and assume $q=q_V+q_{V'}$ and $2q+2\le N$. Let
$\mu=(\mu_\tau)_{\tau\in\auPl}$ be integers with $N\le2\mu_\tau+2q$ for all $\tau$, and let $v$ be a
vector of the Fock space of $\auHom(V',W)\otimes\R$ all of whose moments $\int(1+|z|)^k|v(z)|\,dz$
are finite. Then there are $C>0$ and a vector $x$ of the Fock space of $\auHom(W,V)\otimes\R$ with
finite moments such that the following holds. Let $\pcm\ne(1)$ be a modulus, $\nu$ an automorphic
character of $\auU(1)$ of weight $\mu$, unramified outside $\pcm$ and tamely ramified with
non-trivial tame part at every prime of $\pcm$, and $\chi$ the character of $\pcR_\pcm^\times$
given by its tame parts. Then there is $c\in\C^\times$ with
\[
|c|^2=C\cdot L_*(\pcm,\nu,e)\ne0
\]
such that for all automorphic characters $\alpha,\alpha'$ of $\auU(1)$ with
$\alpha\alpha'=\gamma:=\nu^{-1}\upsilon^{N/2}$, all functions $\Phi$ on $\pcMat_{a\times N}(\pcR_\pcm)$
and $\Psi$ on $\pcMat_{N\times b}(\pcR_\pcm)$ that are matched for $\chi$, and all
$g\in\auU(W)(\auA)$,
\[
\int_{[\auU(V)]}\theta^{W,V}_{\alpha,\gamma}(g,h;x,\Phi)\,dh=c\int_{[\auU(V')]}\theta^{V',W}_{\gamma,
\alpha'}(h',g;v,\Psi)\,dh',
\]
both integrals being absolutely convergent.
\end{theorem}
\noindent Here the kernels are those of \S\ref{pc:ss-data}, in the convention
$\omega_{\alpha,\beta}(g,h)=\alpha(\det g)^{-1}\beta(\det h)\omega(g,h)$, with the data of
Definition~\ref{pgx:conv}. Theorem~\ref{pc:X} is Theorem~\ref{pcx:X}, proved in \S\S\ref{pcx:sec},
\ref{pcu:sec}, \ref{pct:sec}, \ref{pcn:sec}, with absolutely convergent integrals
(Corollary~\ref{pcx:conv}); the dictionary with the other splitting in use is the transposition of
indices \eqref{pcx:eq-dict}, without any factor, and does not change $c$. The proof unfolds a large
triple along an isotropic frame, builds real vectors with prescribed zero terms, removes hyperbolic
planes one at a time (each plane contributes a ratio of values of $L^\pcm(z,\nu)$ and a power of
$\prod_{w\mid\pcm}q_w$, which build up $L_*(\pcm,\nu,e)$), passes from standard data to all matched
functions by equivariance, and evaluates the constant in equal rank; see \S\ref{pcu:ss-final} for the
list of dependencies.

\subsubsection*{From the exchange identity to Theorem~\ref{pc:tri}}
Fix a square: $U_1=\pcG_{i,n+1}$, $U_2=\pcG_{i+1,n+1}$, $A=\pcG_{i,n+2}$, $A'=\pcG_{i+1,n}$, with
$i+n+2\le m$, and $W=U_1\oplus U_2$, of rank $N=2(r+(n+1)t)$. The hypotheses of
Theorem~\ref{pc:X} hold for $(W,V,V')=(W,A,A')$, with $e=t$ and the weight $\mu_\tau=(n+1)\Delta$ of
$(\nu_{n+1}\vartheta^{\,n+1})^{-1}$: the ranks are $a=r+(n+2)t$, $b=r+nt$; the numbers of negative
signs add up, because the signs $\epsilon_p$ of the stems satisfy
$\#\{\epsilon^{(i,n+1)}=-1\}+\#\{\epsilon^{(i+1,n+1)}=-1\}=\#\{\epsilon^{(i,n+2)}=-1\}+\#\{\epsilon^{
(i+1,n)}=-1\}$ (compare the exponents $i+a+\min((a-i)^+,n)$ entry by entry: for $a\le i$ both sides have the
exponents $\{i+a,\,i+a+1\}$; for $i<a\le i+n+1$ all four exponents are $2a$; for $a\ge i+n+2$ both
sides have $\{i+a+n+1,\,i+a+n+2\}$);
$2q+2\le N$ because $q\le2r<r+t$; and $N\le2(n+1)\Delta$ because $\Delta=2r^8$. The character
$(\nu_{n+1}\vartheta^{\,n+1})^{-1}$ is tamely ramified exactly at the primes of $\pcm_\vartheta$,
because $\vartheta$ is $m$-fold and the tame parts of $\chi_M$ have order greater than $2m$. The
deduction then has the following steps.
\begin{enumerate}
\item[(i)] \emph{Kernels multiply}. $\psHS(\auHom(W,A))=\psHS(\auHom(U_1,A))\otimes
\psHS(\auHom(U_2,A))$, the theta functional is the tensor product, and
$[P_1P_2,\delta_\xi]=[P_1,\delta_{\xi_1}]\otimes[P_2,\delta_{\xi_2}]$ for $\xi=(\xi_1\ \xi_2)$. Hence,
for forms $y$ on $U_1$ and $y_2$ on $U_2$, the pairing of $y_2$ with ``lift $U_1\to A$ with datum
$\xi_1$, then transposed lift $A\to U_2$ with datum $\xi_2$'' applied to $y$ is the integral over
$[\auU(U_1)]\times[\auU(U_2)]$ of $y\otimes y_2$ against the theta integral over $[\auU(A)]$ of the
datum $[P_1P_2,\delta_\xi]$ of the pair $(W,A)$; call it $L(\xi)$. In the same way the path through
$A'$ gives a function $L'(\eta)$ on $\pcMat_{\mathsf N\times\mathsf N_0}(\pcR_\vartheta)$. By
Lemma~\ref{pc:level}, $L$ and $L'$ transform by $\chi\circ\det$, provided the splitting characters
of the four kernels are chosen so that their products are $\gamma$ on $W$; this fixes the
characters of the descents (Definition~\ref{ptr:desc-char}).
\item[(ii)] \emph{Convergence} (Proposition~\ref{ptr:conv}). The integrals over $[\auU(A)]$ and $[\auU(A')]$,
whose groups are not compact, converge because the data vanish on the matrices that are not
injective, respectively surjective, at the primes of $\pcm$, and the forms decrease rapidly.
\item[(iii)] \emph{The infinite places} (Theorem~\ref{ptr:square}, with Theorem~\ref{ptq:H}). Theorem~\ref{pc:X}, applied with $v$ the
vector of the polynomial of the path through $A'$, gives a vector $x$ which is not known. One
needs: for every vector $x$ with finite moments there is $\kappa(x)\in\C$, independent of $\pcm$,
of the characters, of the forms and of $\xi$, such that the function $L_x$ formed with $x$ in
place of $P_1P_2$ equals $\kappa(x)L$. This is a statement of multiplicity one for the types
$e_{i,n+1}$, $e_{i+1,n+1}$ in the Fock space, extended from polynomials to vectors with finite
moments by continuity. With it,
$L,L'$ satisfy the square law with the constant $\kappa^\sharp/c(\vartheta)$,
$\kappa^\sharp=\kappa(x)$.
\item[(iv)] \emph{The descents}. Define $d_{i,n}$ as the transposed lift
along the pair $(\pcG_{i+1,n},\pcG_{i,n+1})$ with a finite datum $\varepsilon_{i,n}$, where
$\varepsilon_{i,0}$ is the datum for which $d_{i,0}=a_{i,0}^*$ (it exists by (W7): the complex
conjugate of the kernel of $a_{i,0}$ is a kernel for the space with the negative form) and
$\varepsilon_{i,n+1}=\varepsilon_{i,n}^+$. Then Lemma~\ref{pc:next}, with
$\#\auGL_{\mathsf N_0}=w_n$ and $\#\auGL_{\mathsf A}=w_{n+2}$, is exactly (b) of
Theorem~\ref{pc:tri}.
\item[(v)] \emph{The adjoint law} (Theorem~\ref{ptv:adjoint} with Theorem~\ref{ptw:RS}; the reversal of
a kernel is Proposition~\ref{ptv:rev}, the Gauss sum Proposition~\ref{ptv:fourier}). The adjoint of a descent is, up to a constant of
modulus $1$ up to norm powers and a translation by a central element, the transposed lift along
the reversed pair with the Fourier transform of the conjugate datum (this needs the reversal
$w\mapsto w^\dagger$ of a kernel, that is, the fact (W6) (Proposition~\ref{pf:flip}), and the Gauss sum of $\chi\circ\det$). It commutes with the lifts $a$ by a second
application of Theorem~\ref{pc:X}, now with $e=0$, where $L_*$ is $1$; so its constant has modulus
$1$ up to norm powers.
\end{enumerate}


\subsection{Auxiliary forms: proof of Proposition~\ref{au:aux}}\label{pc:ss-aux}

For $0\le i,j<m$ let
\begin{equation}\label{pc:eq-rhoprime}
\rho'_{ij}(\vartheta)=\prod_{z=1}^t\frac{|L(z,(\mu^0_{i+1}\mu^0_{j+1}\vartheta^{\,j+1})^{-1})|}{|L(z,(
\mu^0_{i+1}\mu^0_{j+1})^{-1})|}.
\end{equation}
The character $\mu^0_{i+1}\mu^0_{j+1}$ has weight $-(\kappa_{i+1}+\kappa_{j+1})$, with
$t\le\kappa_{i+1}(\tau)+\kappa_{j+1}(\tau)\le4b$, and tame part $\chi_M^{\,i+j+2}$ at the chosen
places, which is non-trivial since $2\le i+j+2\le2m$; so $\rho'_{ij}$ is a ramified ratio.

Let $\widetilde Y$ be the positive definite standard space of rank $\tilde n=r+1+m(t-1)=n_m-m$,
let $\pcb=(\pi_M)^{b}$, and let $J(\pcb)\subset K_{\widetilde Y}$ be the principal congruence
subgroup of level $\pcb$; it is free by Lemma~\ref{au:free}, since $\tilde n\le2n_m$.

\begin{theorem}[Integral and short lifts through the second grid]\label{pc:auxcore}
There are a weight $\tilde\lambda$ of $\auU(\widetilde Y)$, a function $c$ on the $m$-fold twists
with values in $\C^\times$ that has modulus $\prod_{0\le i,j<m}\rho'_{ij}$ up to norm powers, a
surjective linear map $\mathsf r\colon\pcFF\to\auAz$ from
$\pcFF=\mathcal A_{\tilde\lambda}(J(\pcb),1)$, and $s\in\C^\times$, such that for every
$\varphi\in\pcFF$:
\begin{enumerate}
\item[(a)] if $\vartheta$ is an $m$-fold twist, $\iota\in\auI$ with $\ell(\iota)\nmid\auNm
Q_\vartheta$, and $\varphi$ is integral at $\iota$, then
$(c(\vartheta)/c(1))\cdot\iota_\vartheta(s\,\mathsf r\varphi)$ is integral at $\iota$;
\item[(b)] if $\sigma\in\auAut(\C)$ and $\|\sigma(\varphi(g))\|^2\le B$ for all $g$, then
$\|\sigma(\iota_1(s\,\mathsf r\varphi)(g))\|^2\le\exp(2x^{201}b^2)\,B$ for all $g$;
\item[(c)] $\pcFF$ is spanned by forms that are integral at every $\iota\in\auI$ and whose values
have all conjugates of squared norm at most $\exp(x^{99}b^2)$.
\end{enumerate}
\end{theorem}
\noindent Parts (a) and (b) are deduced in \S\ref{pgr:ss-deduce} from Propositions~\ref{pgr:first},
\ref{pgr:turn}, \ref{pgr:last}, \ref{pgr:den}, \ref{pgr:delta}, which are proved in
\S\S\ref{pgr:sec}--\ref{pgm:sec}; (c) is Proposition~\ref{pgr:short}, proved in
Proposition~\ref{pgs:short}. See \S\ref{pgr:ss-status}.

\begin{proof}[Proof of Proposition~\ref{au:aux} from Theorem~\ref{pc:auxcore} and
Proposition~\ref{au:good}]
Put $\gamma(\vartheta)=c(\vartheta)/c(1)$. Then $\gamma(1)=1$, $\gamma$ has no zero, and by
Lemma~\ref{pc:mod}(b), applied to $c$ and $f=\prod_{i,j}\rho'_{ij}$ (which is non-negative and
equal to $1$ at $\vartheta=1$), there is $e'\in\Z$ with
$|\gamma(\vartheta)|=(\auNm Q_\vartheta)^{e'/2}\prod_{0\le i,j<m}\rho'_{ij}(\vartheta)$.

Let $\mathfrak B$ be the set of the forms $z=s\,\mathsf r\varphi\in\auAz$ with $\varphi$ as in
(c). It spans $\auAz$, because $\mathsf r$ is onto and $s\ne0$. For $z\in\mathfrak B$, by (a) the
form $\gamma(\vartheta)\iota_\vartheta z$ is integral at every $\iota$ with
$\ell(\iota)\nmid\auNm Q_\vartheta$, for every $m$-fold $\vartheta$; and by (b) every value of
$\sigma(\iota_1z)$ has squared norm at most $\exp(2x^{201}b^2+x^{99}b^2)$, for every $\sigma$.

Let $y=\#\auG\cdot\varrho(e_0)z\in\auT$. By Proposition~\ref{au:good}(a), with $s=1$ and the
constant $\gamma(\vartheta)$, the form $\gamma(\vartheta)\Theta_\vartheta y=\gamma(\vartheta)
\iota_\vartheta y$ is integral at every $\iota$ with $\ell(\iota)\nmid\auNm Q_\vartheta$. By
Proposition~\ref{au:good}(b),
\[
\|\Theta_1y\|_\sigma^2\le\exp\bigl(x^{99}+2x^{201}b^2+x^{99}b^2\bigr)\le\exp(2x_1^2b^2)
\]
for all $\sigma$, because $x^{10^6}\le x_1$. Finally the forms $y$ span $\varrho(e_0)\auAz=\auT$.
\end{proof}

Proposition~\ref{au:good} is proved in \S\ref{ss:au-inputsA} from Proposition~\ref{au:equiv} and
from the bound $\exp(x^{30})$ for the number of classes of the level of $Y_m$
(Proposition~\ref{pf:classes}).

\medskip\noindent\emph{Outline of the proof of Theorem~\ref{pc:auxcore}\textup{(a)}, \textup{(b)}}
(\S\S\ref{pgr:sec}--\ref{pgm:sec}). There are two grids of spaces, $Y^a_b$ and $\widetilde Y^b_a$
$(0\le a,b\le m)$, defined in Definition~\ref{pgv:def-spaces}: $Y^0_b=Y_b$, and $Y^a_b$ has $a$
further blocks of rank $t$; $\widetilde Y^0_0=Y_0\oplus E\tilde e$ with $\langle\tilde e,\tilde
e\rangle=1$, and $\widetilde Y^b_a$ has $a$ blocks of rank $t+1$ and $b$ blocks of rank $t-1$, so that
$\widetilde Y=\widetilde Y^m_0$; $Y^m_m$ and $\widetilde Y^m_m$ are positive definite, and
$\widetilde Y^l_l=Y^l_l\oplus E\tilde e_l$. The edges carry theta lifts with finite data modulo
$\pcb\prod_{w\in Q}w$. \emph{First path} (\S\ref{pgw:ss-first}): the transposed lifts down the first
column of the second grid, followed by restriction to $\auU(Y_0)$ at a real point and projection to
the type $e_0$, give the surjection $\mathsf r$; its surjectivity uses the injectivity of the
untwisted chain of the second grid (Theorem~\ref{pgy:I}) and a vanishing theorem at the deep level
$\pcb$ by Poincar\'e series (Theorem~\ref{pgq:V}). \emph{Turns} (\S\S\ref{pgr:ss-off},
\ref{pgu:sec}, \ref{pgv:ss-square}): the pairing of $\iota_\vartheta(\mathsf r\varphi)$ with a form on
$Y_m$ is a period along a path through the two grids; each square of a grid moves the path by
one square, at the cost of the constant of Theorem~\ref{pc:X} for that square, of modulus
$\rho'_{ij}$ up to norm powers; on the diagonal the turn uses the reversal of the theta kernel of a
line (Lemma~\ref{pgx:line}). \emph{Last path} (\S\S\ref{pgw:ss-last}, \ref{pgz:sec}, \ref{pgt:sec}):
all spaces are definite, the lifts are algebraic lifts by jets, whose kernels are integral up to a
bounded denominator, which gives (a); the size bounds give (b). Pairing with delta forms reads the
values (Proposition~\ref{pgr:delta}).

\subsection{Where the inputs of Theorem~C are proved}\label{pc:ss-remains}

Proposition~\ref{pc:C} deduces Theorem~\ref{thm:srcC} from Propositions~\ref{au:inj}, \ref{au:aux},
\ref{au:factor}, \ref{au:regular} and~\ref{au:Zj}. In this section, Proposition~\ref{au:inj} is proved from
Lemma~\ref{pc:arch}, Theorem~\ref{pc:tri}(a),(b) and (SW); Proposition~\ref{au:factor} from Lemma~\ref{pc:arch},
Theorem~\ref{pc:tri}, (SW) and $(\mathrm M^+)$; Proposition~\ref{au:aux} from Theorem~\ref{pc:auxcore} and
Proposition~\ref{au:good}. Lemma~\ref{pc:arch} is Proposition~\ref{ptv:arch} (from (V), (D1)--(D3));
Theorem~\ref{pc:tri} is proved in \S\S\ref{ptr:sec}--\ref{ptw:sec} (from Theorem~\ref{pc:X}, (V), (D1)--(D5),
(W7)); Theorem~\ref{pc:X} is Theorem~\ref{pcx:X}, whose inputs are listed in \S\ref{pct:ss-status};
Theorem~\ref{pc:auxcore} is proved in \S\S\ref{pgr:sec}--\ref{pgm:sec} (\S\ref{pgr:ss-status}).
Proposition~\ref{pc:rallis} and Propositions~\ref{pc:single}, \ref{pc:first} use Weil's formula (SW), (W7) and
Lemma~\ref{pc:arch}; Lemma~\ref{pc:detcompare} uses $(\mathrm M^+)$; Lemma~\ref{pc:clean} uses (W2), (W5),
(W7); Lemmas~\ref{pc:level}, \ref{pc:series}, \ref{pc:heckeclosure} use (W4), (W5). Propositions~\ref{au:regular},
\ref{au:onelift}, \ref{au:hecke} and $(\mathrm M^+)$ are proved in the sections named in their proofs in
\S\ref{ss:au-proofC}. The character of (W7) is computed in Lemma~\ref{ptw:W7}.

Two remarks on the statements of \S\ref{ss:au-proofC}. First, $Z_j(\vartheta)$ in
Proposition~\ref{au:factor} must be read with the lift $\Theta_{j,\vartheta}$ normalised as before
Lemma~\ref{pc:algpet} (block function, level $J^V_\vartheta$, inner product of algebraic modular
forms): another normalisation changes $Z_j(\vartheta)$ by the $\audT$-th power of a factor such as
$\prod_{w\in Q}\prod_i(1-q_w^{-i})$, which is not a power of $\auNm Q$ times ramified ratios. With
this normalisation the constant $c_Q(\vartheta)$ of Proposition~\ref{pc:single} cancels against the
index of the level of $Y_m$ (Lemma~\ref{pc:index}). Secondly, Proposition~\ref{au:inj} is not proved through the non-vanishing of
$L$-values in the Rallis inner product formula for each of $m$ lifts: only the first lift of each row of the triangle is in Weil's range,
its injectivity comes from the fact that the series of Lemma~\ref{pc:series} is close to the
identity, and the injectivity of the chain follows from the commutation relations by
Proposition~\ref{pc:tri-inj}, without any information on the constants of the squares.


%% ===== proofs/C-onelift.tex =====
% proofs/C-onelift.tex
% Propositions au:onelift and au:hecke of sec-automorphic.tex: the comparison operator, false lifts,
% Hecke values. To be \input after the A-files (it uses pl:ss-norm, pl:ss-wa, pl:ss-transfer,
% pk:ss-branch). Labels of this file start with po:. Only amsmath/amssymb/amsthm are used.
\providecommand{\auO}{\mathcal O}
\providecommand{\auI}{\mathcal I}
\providecommand{\auU}{\mathrm U}
\providecommand{\auGL}{\mathrm{GL}}
\providecommand{\auHom}{\operatorname{Hom}}
\providecommand{\auEnd}{\operatorname{End}}
\providecommand{\auAut}{\operatorname{Aut}}
\providecommand{\auNm}{\operatorname{Nm}}
\providecommand{\auA}{\mathbb A}
\providecommand{\auF}{\mathbb F}
\providecommand{\auPl}{\Phi}
\providecommand{\auG}{\overline G}
\providecommand{\audT}{d_{\mathcal T}}
\providecommand{\aubud}{\mathrm{bud}}
\providecommand{\auhdir}{h_{\mathrm{dir}}}
\providecommand{\ausinf}{s_\infty}
\providecommand{\ausfin}{s_{\mathrm{fin}}}
\providecommand{\aueu}{\varepsilon_{\mathrm u}}
\providecommand{\aurd}{\mathrm{rd}_E}
\providecommand{\aubar}[1]{\overline{#1}}
\providecommand{\auf}[1]{\texttt{#1}}
\providecommand{\psrho}{\varrho}
\providecommand{\pao}{\mathfrak o}
\providecommand{\plone}{\mathbf 1}
\providecommand{\plgb}[2]{\|#2\|_{#1}}
% macros of this file (prefix \po)
\providecommand{\poA}{\mathcal A}
\providecommand{\poZ}{\mathcal Z}
\providecommand{\poT}{\mathbb T}
\providecommand{\popr}{\operatorname{pr}}
\providecommand{\poglue}{\operatorname{glue}}
\providecommand{\poran}{\operatorname{im}}
\providecommand{\podel}{\delta}
\providecommand{\pop}{\mathsf p}
\providecommand{\pog}{\mathsf g}
\providecommand{\poAA}{A_\ast}

\subsection{One lift: the comparison operator, false lifts and Hecke values}\label{po:sec}

This subsection treats Propositions~\ref{au:onelift} and~\ref{au:hecke}. We fix $1\le j\le m$
throughout and write
\[
W=W_j,\qquad V=V_j,\qquad n=r+t=6r,\qquad r_+=k+j-1,\quad r_-=m-j,\qquad
\poAA=f\,m\,b\,(\log b)^2 .
\]

\emph{What is proved and where.} The two propositions are deduced here, with complete
arguments, from a list of statements about theta kernels and Hecke values which are stated in this
subsection and proved in the subsections that follow. The deduction contains the linear algebra
(the determinant identity, the two eigenvalue bounds, the tight frame, the projection that does not
increase any conjugate norm), the Hecke equation of a lift, the change of lattice by which a lift
is cut into products (``gluing''), and the bookkeeping of heights. The statements and the places
where they are proved:
\begin{center}
\begin{tabular}{lll}
\hline
Statement&Content&Proved in\\
\hline
(L4)&the descent of a form on $V$ is of the type of $\poA(W)$&Prop.~\ref{plf:L4}\\
Prop.~\ref{po:exchange}&three exchanges&\S\ref{pte:sec}, \S\ref{ptf:sec}\\
Prop.~\ref{po:span}&products span&Prop.~\ref{psp:main}, Thm.~\ref{psv:main}\\
Prop.~\ref{po:periods}&periods of false lifts vanish&Thm.~\ref{ppe:main}\\
Lemma~\ref{po:ff}&first factors: integrality, size&Thm.~\ref{pfq:ff}\\
Lemma~\ref{po:classes}&numbers of classes at level $\pi_M^2$&Lemma~\ref{plf:classes}\\
Prop.~\ref{po:rallis}&the norm of a lift from $\auU(1)$&Prop.~\ref{pex:rallis}\\
Prop.~\ref{po:arch}&all conjugates of the untwisted norm&Thm.~\ref{psu:arch}, Thm.~\ref{pcj:CL}\\
Prop.~\ref{po:small}&a test form of least valuation&Thm.~\ref{psu:small}, \ref{psm:attained}\\
\hline
\end{tabular}
\end{center}
Two definitions of \S\ref{ss:au-proofC} have to be made precise for the statements to have a
meaning (false lifts; the clause of (R3) on main pairs); this is done in
Definitions~\ref{po:false}, \ref{po:denom} and~\ref{po:R3}. With these, Propositions~\ref{au:onelift}
and~\ref{au:hecke} hold as stated, and the
proofs of Proposition~\ref{au:Zj} and of Theorem~\ref{thm:srcC} are unchanged.

\subsubsection{The lift $W\to V$: definitions and standing facts}\label{po:ss-set}

\emph{Places.} $\mathsf M$ is the set of chosen places and, for a twist $\vartheta$,
$\mathsf Q=Q_\vartheta$; $\mathsf S_\vartheta=\mathsf M\cup\mathsf Q$. A prime $w$ of $E$ is
\emph{off the level} if $w\ne\aubar w$ and $w,\aubar w\notin\mathsf S_\vartheta$. For
$w\in\mathsf S_\vartheta$, $E_w$ is identified with the completion and $\auF_w$ is the residue
field. We use (N1)--(N5) of \S\ref{pl:ss-norm}, in particular the coordinates $T$ on
$\auHom(W_w,V_w)=M_{n\times r}(E_w)$ at the split places and the formulas \eqref{pl:eq-w5},
\eqref{pl:eq-heis}; the proofs of these formulas are local and apply to the pair $(W,V)$ at every
place $w\ne\aubar w$ (for the odd ranks $r$, $t$ the factor in \eqref{pl:eq-w5} is changed by an
unramified character of $E_w^\times\times E_w^\times$; we use the formula in the form
\eqref{po:eq-w5} below, where the factor on $\auU(V)$ is an unramified character $\eta_w$ that does not
involve $\mu$: on elements of $\auGL(\auO_w)$ it is trivial, and for the non-integral elements $d_1$
(Lemma~\ref{po:facts}) and $\pog_s$ (Lemma~\ref{po:slice}) only $|\eta_w(1,h)|=|\det h|_w^{-r/2}$ and its
independence of $\vartheta$ are used).

\emph{The character, the Weil representation and the datum.} $\mu=\mu_j^0\vartheta$; its tame
part at $w\in\mathsf S_\vartheta$ is a character $\psi_w$ of $\auF_w^\times$, non-trivial for
$w\in\mathsf Q$ by the definition of a twist, and for $w\in\mathsf M$ equal to the tame part of
$\mu_j^0$; we write $\psi_{\mathsf M}(\det s)=\prod_{w\in\mathsf M}\psi_w(\det s_w)$ for
$s\in\auG$. In this subsection $\omega_\mu$ denotes the Weil representation of
$\auU(W)(\auA)\times\auU(V)(\auA)$ on $\mathcal H(\auHom(W,V))$ in the following normalisation: the group $\auU(V)$ acts through the splitting given by
powers of $\upsilon$ alone, and the action of $\auU(W)$ is the one of that splitting multiplied by
$\mu(\det g)^{-1}$. At a place $w\ne\aubar w$ this means, in the coordinates of (N5),
\begin{equation}\label{po:eq-w5}
\omega_\mu(g,h)\Psi(T)=\chi_w(\det g)^{-1}\,\eta_w(g,h)\,\Psi(h^{-1}Tg),
\end{equation}
where $\chi$ is the base change of $\mu$ and $\eta_w$ is an unramified character of
$\auGL_r(E_w)\times\auGL_n(E_w)$ that does not depend on $\mu$, with $|\eta_w(1,h)|=|\det h|_w^{-r/2}$.
In the notation of \S\ref{pga:ss-L} this is the action $a(\det h)\,b(\det g)^{-1}\omega(x\mapsto
hxg^{-1})$ with $a=\upsilon^{r_-}$ and $b=\mu\,\upsilon^{\lfloor n/2\rfloor}$. (It differs from the
convention of \S\ref{ps:ss-weil} for the upper chain, where the character twists both groups, by the
unitary character $\mu(\det h)$ of $\auU(V)(\auA)$; with the present convention the scalars of
the Hecke operators and of the change of lattice on $V$ do not depend on $\vartheta$. The same
representation is used in \S\ref{pga:ss-L}; the lift of \S\ref{pc:ss-rallis} differs from it by multiplication with a unitary character of the determinant on $\auU(V)$, which changes no norm and no determinant $Z_j$: Remark~\ref{pc:rem-norm}.) Let
$\iota_0\in M_{n\times r}$ be the matrix of the inclusion of the first $r$ coordinates, and for
$T\in M_{n\times r}$ let $T'$ be its upper $r\times r$ block. The \emph{block datum} is
$\Phi_\vartheta=\bigotimes_{w\in\mathsf M}\plone[T\equiv\iota_0\bmod w]\otimes
\bigotimes_{w\in\mathsf Q}\plone[T\text{ integral}]\,\psi_w(\det T'\bmod w)\otimes(\text{vacuum
elsewhere})$, with $\psi_w(0)=0$.

\emph{The weight on $V$.} At each $\tau\in\auPl$ there are two Young diagrams $\mu^+_\tau$ and
$\mu^-_\tau$, with columns of heights at most $r_+$, respectively at most $r_-$; they are
determined by the parameter $\varsigma_\tau$ and by $\kappa_j(\tau)$ (explicitly in Lemma~\ref{plf:match}), and they do not depend on
$\vartheta$. We use only: $\mu^+_\tau$ has at most $4b$ columns, and $\mu^-_\tau$ has at most
$\Delta$ columns of each height. In the notation of \S\ref{pk:ss-branch} put
\[
\poT^+(R^n)=\bigotimes_{\mathfrak c\text{ column of }\mu^+}\Lambda^{h_{\mathfrak c}}R^n,\qquad
\poT^-(R^n)=\bigotimes_{\mathfrak c\text{ column of }\mu^-}\Lambda^{h_{\mathfrak c}}(R^n)^\vee
\]
(tensor products over all $\tau$ as well), with the bases of pure tensors of basis vectors and the
hermitian forms for which these are orthonormal. The group $\prod_\tau\auGL_n(\C)$ acts on
$\poT^+\otimes\poT^-$, and $W_{\lambda_V}\subseteq\poT^+(\C^n)\otimes\poT^-(\C^n)$ is the span of
the vectors
\begin{equation}\label{po:eq-ev}
\mathrm{ev}(y,y')=\bigotimes_{\mathfrak c}(y_1\wedge\dots\wedge y_{h_{\mathfrak c}})\otimes
\bigotimes_{\mathfrak c}(y'_1\wedge\dots\wedge y'_{h_{\mathfrak c}}),\qquad y_i\in\C^n,\ y'_i\in
(\C^n)^\vee,\ y'_{i'}(y_i)=0\ \text{for all }i,i'
\end{equation}
(at each $\tau$). This is the irreducible representation with highest weight
$(\mu^+_\tau;0,\dots,0;-\mu^-_\tau)$; in the conventions of (N1) it has columns of heights
$h_{\mathfrak c}$ and $n-h_{\mathfrak c}$ and a power of the determinant, and
$W_{\lambda_V}(\Z)=W_{\lambda_V}\cap(\poT^+\otimes\poT^-)(\Z^n)$, with the restricted hermitian
form, by \eqref{pk:eq-dual} and because $\Lambda^h(\Z^n)^\vee\otimes\det=\Lambda^{n-h}\Z^n$
isometrically. As in \S\ref{pk:ss-branch}, for $\xi\in W_{\lambda_V}\otimes\C_\ell$ the largest
absolute value of a coordinate of $\xi$ in the basis of $W_{\lambda_V}(\Z)$ is the largest
absolute value of a tensor coordinate (the two unit balls coincide, $W_{\lambda_V}(\Z)$ being
saturated in the tensors).

\emph{Forms on $V$.} $K_\vartheta\subseteq K^\circ$ is the group of the $k$ with $k_w\iota_0\equiv
\iota_0\bmod w$ for $w\in\mathsf M$ and with $k_w\bmod w$ block lower triangular, for the blocks
$r+t$, for $w\in\mathsf Q$; by \eqref{po:eq-w5}, $\omega_\mu(1,k)\Phi_\vartheta=\chi_\vartheta(k)^{-1}
\Phi_\vartheta$ for a character $\chi_\vartheta$ of $K_\vartheta$ of finite order, trivial on the
elements that are $\equiv1$ at $\mathsf S_\vartheta$. Put
$\poA(V;\vartheta)=\poA_{\lambda_V}(K_\vartheta,\chi_\vartheta^{-1})$ (so $\Xi(\gamma hk)=\chi_\vartheta(k)^{-1}\gamma\,\Xi(h)$), with the inner product and the
norms $\|\cdot\|_\sigma$ of \S\ref{ss:au-spaces}. The level $K_\vartheta$ is not free. Let
\[
n_{\mathrm{Aut}}=\#\auGL_{2fn}(\auF_3).
\]

\begin{lemma}\label{po:naut}
For every $g\in\auU(V)(\auA_{E^+,f})$ the order of $\Gamma\cap gK^\circ g^{-1}$ divides
$n_{\mathrm{Aut}}$, and $n_{\mathrm{Aut}}\le\exp(x^{40})$. If $\Xi,\Xi'\in\poA(V;\vartheta)$ are
integral at $\aubar\iota$ and at $\iota$ respectively, then $n_{\mathrm{Aut}}\langle\Xi,\Xi'\rangle
\in\auO_\iota$.
\end{lemma}
\begin{proof}
The group is finite and is a subgroup of $\auGL_n(E)\subseteq\auGL_{2fn}(\Q)$. A finite subgroup of
$\auGL_N(\Q)$ stabilises a lattice, so it is conjugate to a subgroup of $\auGL_N(\Z)$, and the
kernel of $\auGL_N(\Z)\to\auGL_N(\auF_3)$ has no torsion (Minkowski): if $g=1+3^aX$ with $a\ge1$,
$X\not\equiv0\bmod3$ and $g^p=1$ for a prime $p$, then $pX\equiv0\bmod3^{a}$ for $p\ne3$, and
$X+3^aX^2+3^{2a-1}X^3=0$ for $p=3$; both are absurd. So the order divides $\#\auGL_N(\auF_3)\le
3^{N^2}$, $N=2fn\le12fr$, and $f\le x^{12}$, $r\le3x^5$. For the last assertion,
$\langle\Xi,\Xi'\rangle=\sum_g\#(\Gamma\cap gK_\vartheta g^{-1})^{-1}\langle\Xi(g),\Xi'(g)\rangle$,
and each $\langle\Xi(g),\Xi'(g)\rangle\in\auO_\iota$ as in \S\ref{ss:au-spaces}.
\end{proof}

\emph{The lift and the descent.} For $x\in\auHom(W,V)\otimes_{E^+,\tau}\R=M_{n\times r}(\C)$ write
$x=(x^+\mid x^-)$ for the columns belonging to the positive and to the negative basis vectors of
$W$. For $\xi\in W_{\lambda_V}$ let $\pop_\xi$ be the vector of the Fock space of
$\auHom(W,V)\otimes\R$ given by the polynomial
\begin{equation}\label{po:eq-poly}
\pop_\xi(x)=\bigl\langle\xi,\ \mathrm{ev}\bigl(\aubar{x^+},(x^-)^\flat\bigr)\bigr\rangle ,
\end{equation}
where $\flat\colon\C^n\to(\C^n)^\vee$ is the conjugate-linear map of the standard form, the
columns of $x^+$ are taken in the fixed numbering of the positive basis vectors and those of $x^-$
in the reverse numbering of the negative ones, and $\langle\ ,\,\rangle$ is the form of
$\poT^+\otimes\poT^-$ (in the Fock model the positive variables are holomorphic and the negative
ones antiholomorphic, so that $\pop_\xi$ is a polynomial vector; it is conjugate-linear in $\xi$
and $\auU(V)(\R)$-equivariant; these are the joint harmonics of~\cite[Ch.~III, \S5]{KashiwaraVergne}). For
$\Phi\in\mathcal H_f(\auHom(W,V))$, $g\in\auU(W)(\auA)$ and $h\in\auU(V)(\auA_f)$ put
\[
k_\Phi(g,h;\xi)=\theta\bigl(\omega_\mu(g,h)(\pop_\xi\otimes\Phi)\bigr),
\]
with $\omega_\mu$ as in \eqref{po:eq-w5}. For a form
$y\in\poA(W_j)$ (\S\ref{ss:au-spaces}; these forms are cuspidal, hence rapidly decreasing, and the kernel has moderate growth by
Lemma~\ref{pcx:thetagrowth}, so the integral converges) the function $\Theta(\Phi)y\colon\auU(V)(\auA_f)\to W_{\lambda_V}$ is
defined by
\begin{equation}\label{po:eq-lift}
\bigl\langle\xi,(\Theta(\Phi)y)(h)\bigr\rangle=\int_{[\auU(W)]}y(g)\,k_\Phi(g,h;\xi)\,dg ,
\end{equation}
and $\Theta_\vartheta=\Theta(\Phi_\vartheta)$; up to a positive constant independent of
$\vartheta$ its restriction to $e_0\poA(W_j)$ is the map $u$ of Proposition~\ref{au:factor}. As in
(N4): $(\Theta(\Phi)y)(\gamma h)=\gamma(\Theta(\Phi)y)(h)$ for $\gamma\in\Gamma$ (invariance of
$\theta$ under rational points and equivariance of $\xi\mapsto\pop_\xi$), and
\begin{equation}\label{po:eq-equiv}
R_{h'}\,\Theta(\Phi)y=\Theta\bigl(\omega_\mu(1,h')\Phi\bigr)y,\qquad
\Theta\bigl(\omega_\mu(g',1)\Phi\bigr)y=\Theta(\Phi)(R_{g'^{-1}}y)
\end{equation}
for finite adelic $g',h'$, with $R$ the right translation. Hence $\Theta_\vartheta y\in\poA(V;
\vartheta)$. The \emph{descent} of $\Xi\in\poA(V;\vartheta)$ is the function on $\auU(W)(\auA)$
\begin{equation}\label{po:eq-descent}
(D_\vartheta\Xi)(g)=\aubar{\sum_{h\in\Gamma\backslash\auU(V)(\auA_f)/K_\vartheta}\frac{k_{\Phi_\vartheta}
(g,h;\Xi(h))}{\#(\Gamma\cap hK_\vartheta h^{-1})}},\qquad\text{so that}\qquad
\langle\Xi,\Theta_\vartheta y\rangle=\int_{[\auU(W)]}y(g)\,\aubar{(D_\vartheta\Xi)(g)}\,dg
\end{equation}
by \eqref{po:eq-lift} and the definition of the inner product.

\emph{The finite group.} For $s\in\auG=\prod_{w\in\mathsf M}\auGL_r(\auF_w)$ let $\tilde s\in
\auU(W)(\auA_f)$ and $\tilde s'\in\auU(V)(\auA_f)$ be integral elements, equal to $1$ outside
$\mathsf M\cup\aubar{\mathsf M}$, that reduce to $s$ and to $\operatorname{diag}(s,1_t)$ at
$\mathsf M$. The element $\tilde s'$ normalises $K_\vartheta$, and
$\varrho_V(s)\Xi=\psi_{\mathsf M}(\det s)\,R_{\tilde s'}\Xi$ defines a representation of $\auG$ on
$\poA(V;\vartheta)$, independent of the choice of $\tilde s'$; it preserves $\|\cdot\|_\sigma$ for
every $\sigma$, the scalar being a root of unity.

\begin{lemma}[Standing facts]\label{po:facts}
Let $\vartheta$ be a twist.
\begin{enumerate}
\item[(L2)] $\Theta_\vartheta\circ\varrho(s)=\varrho_V(s)\circ\Theta_\vartheta$ for $s\in\auG$.
Hence $\Theta_\vartheta(e_0\poA(W_j))\subseteq\varrho_V(e_0)\poA(V;\vartheta)$.
\item[(L3)] Let $w$ be off the level, with residue field of $q$ elements, $P_w(X)=\sum_{i\ge0}
(-X)^iq^{\binom i2}T_{w,i}$ the operator of (N3) on $\poA(V;\vartheta)$, and $\alpha_w$ the number
with $\omega_\mu(1,d_1)\plone=\alpha_w\,\plone(d_1^{-1}T)$ at $w$ (by \eqref{po:eq-w5}, $\alpha_w=\eta_w(1,d_1)$ does not depend on $\vartheta$, and
$|\alpha_w|=q^{r/2}$). Put $X_w=(\alpha_wq^{t-1})^{-1}$. Then $P_w(X_w)\,\Theta_\vartheta y=0$ for
every $y\in\poA(W_j)$.
\end{enumerate}
\end{lemma}
\begin{proof}
(L2) By \eqref{po:eq-w5} for integral elements, $\omega_\mu(\tilde s,\tilde s')\Phi_\vartheta=
\psi_{\mathsf M}(\det s)^{-1}\Phi_\vartheta$: the unramified character is trivial, and
$\plone[\tilde s'^{-1}T\tilde s\equiv\iota_0]=\plone[T\equiv\iota_0]$ because
$\operatorname{diag}(s,1)\iota_0s^{-1}=\iota_0$. So by \eqref{po:eq-equiv}
$R_{\tilde s'}\Theta_\vartheta y=\Theta(\omega_\mu(1,\tilde s')\Phi_\vartheta)y=\psi_{\mathsf M}
(\det s)^{-1}\Theta(\omega_\mu(\tilde s^{-1},1)\Phi_\vartheta)y=\psi_{\mathsf M}(\det s)^{-1}
\Theta_\vartheta(R_{\tilde s}y)$.

(L3) The datum $\Phi_\vartheta$ is the vacuum at $w,\aubar w$, and $\mu$ is unramified there. By
\eqref{po:eq-equiv}, $T_{w,i}\Theta_\vartheta y=\Theta(T'_i\Phi_\vartheta)y$ with $T'_i=\sum_{h\in
K'd_iK'/K'}\omega_\mu(1,h)$, $K'=\auGL_n(\auO_w)$, as in the proof of Proposition~\ref{pl:onelift}.
The computation in the proof of Lemma~\ref{pl:vac} (with $\alpha=\alpha_w$, and $n$, $n'$ there
equal to $r$, $n$ here) gives, as functions of $T\in M_{n\times r}(E_w)$,
\[
\sum_i(-X)^iq^{\binom i2}T'_i\,\plone=\prod_{l<n-\operatorname{rk}(T\bmod w)}(1-q^l\alpha_wX)\cdot
\plone .
\]
Since $\operatorname{rk}(T\bmod w)\le r$, the factor with $l=t-1$ occurs for every $T$, and it
vanishes at $X=X_w$.
\end{proof}

\begin{enumerate}
\item[(L4)] \emph{If $\Xi\in\varrho_V(e_0)\poA(V;\vartheta)$ is orthogonal to
$\Theta_\vartheta(e_0\poA(W_j))$, then $D_\vartheta\Xi=0$.}
\end{enumerate}
\noindent\emph{Proof of {\rm(L4)}.} This is Proposition~\ref{plf:L4}.


\subsubsection{Two facts of linear algebra}\label{po:ss-linalg}

\begin{lemma}\label{po:det}
Let $u\colon U\to H$ be an injective linear map of finite-dimensional complex inner product
spaces, $\popr$ the orthogonal projection of $H$ onto $u(U)$, and $T\in\auEnd(H)$. Then
$\det(u^*u)\cdot\det\bigl(\popr T|_{u(U)}\bigr)=\det(u^*Tu)$.
\end{lemma}
\begin{proof}
Let $e\colon U\to u(U)$ be the isomorphism induced by $u$. Since $u^*$ vanishes on $u(U)^\perp$,
$u^*Tu=u^*\popr Tu=(u^*u)\,e^{-1}(\popr T|_{u(U)})\,e$; take determinants.
\end{proof}

\begin{lemma}\label{po:eig}
Let $L$ be a finite-dimensional complex vector space, $S\in\auEnd(L)$, $\phi\colon\C\to F$ a ring
homomorphism to a field with an absolute value, and $B\ge0$. If $|\phi(\alpha)|\le B$ for every
eigenvalue $\alpha$ of $S$, then $|\phi(\det S)|\le B^{\dim L}$.
\end{lemma}
\begin{proof}
$\det S$ is the product of the roots of the characteristic polynomial, with multiplicities.
\end{proof}

\subsubsection{Conjugate pairings, Hecke operators, false lifts}\label{po:ss-false}

For $\sigma\in\auAut(\C)$ put $c_\sigma=\sigma^{-1}\circ c\circ\sigma$, an involution of the field
$\C$ that inverts roots of unity, and for $\Xi,\Xi'\in\poA(V;\vartheta)$
\[
\langle\Xi,\Xi'\rangle_\sigma=\sigma^{-1}\Bigl(\sum_{h}\#(\Gamma\cap hK_\vartheta h^{-1})^{-1}\,
\aubar{\sigma(\Xi(h))}^{\,t}\,G\,\sigma(\Xi'(h))\Bigr),
\]
where $\sigma$ is applied to the coordinates in the basis of $W_{\lambda_V}(\Z)$ and $G$ is the
integral Gram matrix. This is well defined: $\sigma(\gamma\xi)=\gamma^\sigma\sigma(\xi)$, where
$\gamma^\sigma$ acts through the embeddings $\sigma\circ\tau\circ c$, and these matrices are again
unitary, complex conjugation commuting with all embeddings of the CM field $E$; and $G$ is
invariant under unitary matrices. The form is linear in $\Xi'$, $c_\sigma$-linear in $\Xi$,
$\langle\Xi',\Xi\rangle_\sigma=c_\sigma\langle\Xi,\Xi'\rangle_\sigma$, and
$\sigma\langle\Xi,\Xi\rangle_\sigma=\|\Xi\|_\sigma^2$, which is positive for $\Xi\ne0$. For
$\sigma=1$ it is the inner product. Consequently every operator $A$ has a unique adjoint
$A^{\dagger_\sigma}$ for $\langle\ ,\,\rangle_\sigma$, and
\begin{equation}\label{po:eq-cs}
\bigl|\sigma\langle\Xi,\Xi'\rangle\bigr|\le\|\Xi\|_{c\sigma c}\,\|\Xi'\|_\sigma
\end{equation}
by the Cauchy--Schwarz inequality, since $\sigma(\aubar z)=\aubar{(c\sigma c)(z)}$.

\begin{lemma}\label{po:adj}
Let $w$ be off the level and $\sigma\in\auAut(\C)$.
\begin{enumerate}
\item[(a)] $T_{w,i}^{\dagger_\sigma}=T^*_{w,i}$, the operator of the double coset
$K_wd_i^{-1}K_w$; it does not depend on $\sigma$, and the operators $T_{w,i}$, $T^*_{w,i'}$ all
commute.
\item[(b)] For $X\in\C$ the operator $A=P_w(X)$ commutes with $A^{\dagger_\sigma}$. Hence
$\langle Av,z\rangle_\sigma=0$ for all $v$ whenever $Az=0$; and (for $\sigma=1$) every form
orthogonal to the image of $A$ is killed by $A$.
\item[(c)] For $s\in\auG$, $\varrho_V(s)^{\dagger_\sigma}=\varrho_V(s^{-1})$; for every central
idempotent $e$ of $\C[\auG]$, $\varrho_V(e)^{\dagger_\sigma}=\varrho_V(e)$. The operators
$\varrho_V(s)$ commute with the $T_{w,i}$.
\end{enumerate}
\end{lemma}
\begin{proof}
(a) $(T_{w,i}\Xi)(h)=\sum_{x\in K_wd_iK_w/K_w}\Xi(hx)$ has integer coefficients, so $\sigma\circ
T_{w,i}\Xi=T_{w,i}(\sigma\circ\Xi)$. The sum defining the pairing is a constant times the integral
over $\Gamma\backslash\auU(V)(\auA_f)$ of the function $B(h)=\aubar{\sigma\Xi(h)}^{\,t}G\,
\sigma\Xi'(h)$, and for right $K_w$-invariant functions $a,b$ on $\Gamma\backslash\auU(V)(\auA_f)$,
with the Haar measure of $\auGL_n(E_w)$ giving $K_w$ the volume $1$,
\[
\int\sum_{x\in K_wd_iK_w/K_w}a(hx)\,b(h)\,dh=\int\!\!\int_{K_wd_iK_w}a(hx)b(h)\,dx\,dh=
\int a(h)\sum_{x'\in K_wd_i^{-1}K_w/K_w}b(hx')\,dh .
\]
The Hecke algebra of $(\auGL_n(E_w),K_w)$ is commutative: transposition is an anti-automorphism
that preserves $K_w$ and every double coset (elementary divisors).
(b) By (a), $\langle P_w(X)\Xi,\Xi'\rangle_\sigma=\langle\Xi,A'\Xi'\rangle_\sigma$ with
$A'=\sum_i(-c_\sigma X)^iq^{\binom i2}T^*_{w,i}$, and $A'$ commutes with $A$. If $Az=0$ then
$\langle A'z,A'z\rangle_\sigma=\langle AA'z,z\rangle_\sigma=\langle A'Az,z\rangle_\sigma=0$, so
$A'z=0$ and $\langle Av,z\rangle_\sigma=\langle v,A'z\rangle_\sigma=0$. If $u$ is orthogonal to
the image of $A$, then $\langle Au,Au\rangle=\langle u,A'Au\rangle=\langle u,AA'u\rangle=0$.
(c) Right translation by $\tilde s'$ commutes with $\sigma$ and preserves the sum over the
classes, and $c_\sigma$ inverts the root of unity $\psi_{\mathsf M}(\det s)$; so $\langle\varrho_V(s)\Xi,\varrho_V(s)\Xi'\rangle_\sigma=\langle\Xi,\Xi'\rangle_\sigma$.
A central idempotent is a sum of idempotents $e_\chi=\chi(1)\#\auG^{-1}\sum_s\chi(s^{-1})s$, and
$c_\sigma(\chi(s^{-1}))=\chi(s)$ because $\chi(s^{-1})$ is a sum of roots of unity; hence
$\varrho_V(e_\chi)^{\dagger_\sigma}=\chi(1)\#\auG^{-1}\sum_s\chi(s)\varrho_V(s^{-1})=
\varrho_V(e_\chi)$. The last assertion holds because $\tilde s'$ is $1$ at $w$ and $\aubar w$.
\end{proof}

\begin{definition}[False lifts]\label{po:false}
Let $\poZ\subseteq\poA(V;\vartheta)$ be a subspace and $e$ a central idempotent of $\C[\auG]$. For
$w$ off the level, the \emph{false lifts at $w$ for $(e,\poZ)$} are the elements of
\[
\mathcal F_w(e,\poZ)=\varrho_V(e)\poA(V;\vartheta)\ \cap\ \ker P_w(X_w)\ \cap\ \poZ^\perp .
\]
A form is \emph{orthogonal to the false lifts for $(e,\poZ)$} if it is orthogonal to
$\mathcal F_w(e,\poZ)$ for every $w$ off the level.
\end{definition}

Thus a false lift is a form of the same type as the lifts at the chosen places, with the one
Hecke equation (L3) of a lift at one split place, and orthogonal to $\poZ$. For
$\poZ=\poZ_\vartheta:=\Theta_\vartheta(e_0\poA(W_j))$ and $e=e_0$ this is the notion used in
(R3) of Definition~\ref{def:au-regular}.

\begin{proposition}[The projection does not increase any conjugate norm]\label{po:proj}
Let $e$ be a central idempotent of $\C[\auG]$, $w$ off the level, and $\poZ\subseteq\varrho_V(e)
\poA(V;\vartheta)\cap\ker P_w(X_w)$. Let $y\in\poA(V;\vartheta)$ be orthogonal to
$\mathcal F_w(e,\poZ)$, and let $\popr_\poZ$ be the orthogonal projection onto $\poZ$ (for the inner
product, $\sigma=1$). Then $\|\popr_\poZ y\|_\sigma\le\|y\|_\sigma$ for every $\sigma\in\auAut(\C)$.
\end{proposition}
\begin{proof}
Put $A=P_w(X_w)$, $\varepsilon=\varrho_V(e)$ and $R=\ker\varepsilon+\poran A$. By
Lemma~\ref{po:adj}(c) $\varepsilon$ is self-adjoint, so $(\ker\varepsilon)^\perp=\poran\varepsilon$,
and by Lemma~\ref{po:adj}(b) $(\poran A)^\perp\subseteq\ker A$. Hence
$(\poZ+R)^\perp\subseteq\poZ^\perp\cap\poran\varepsilon\cap\ker A=\mathcal F_w(e,\poZ)$, and
therefore $y\in\mathcal F_w(e,\poZ)^\perp\subseteq(\poZ+R)^{\perp\perp}=\poZ+R$. Write $y=z+x$ with
$z\in\poZ$ and $x=x_1+Av$, $\varepsilon x_1=0$. Let $z'\in\poZ$ and $\sigma\in\auAut(\C)$; then
$z'=\varepsilon z'$ and $Az'=0$, so by Lemma~\ref{po:adj}
\[
\langle x,z'\rangle_\sigma=\langle x_1,\varepsilon z'\rangle_\sigma+\langle Av,z'\rangle_\sigma=
\langle\varepsilon x_1,z'\rangle_\sigma+0=0 .
\]
For $\sigma=1$ this shows $x\perp\poZ$, that is, $z=\popr_\poZ y$. For general $\sigma$ it gives
$\langle z,x\rangle_\sigma=c_\sigma\langle x,z\rangle_\sigma=0$ and
$\langle z+x,z+x\rangle_\sigma=\langle z,z\rangle_\sigma+\langle x,x\rangle_\sigma$; applying
$\sigma$, $\|y\|_\sigma^2=\|z\|_\sigma^2+\|x\|_\sigma^2\ge\|z\|_\sigma^2$.
\end{proof}

\begin{definition}[Denominators]\label{po:denom}
Let $\iota\in\auI$ with prime $\ell$, and $\beta\ge0$. The twist $\vartheta$ \emph{has
denominators at most $\beta$ at $\iota$ for the $j$-th lift} if there is an endomorphism $\Lambda$
of $\poZ_\vartheta=\Theta_\vartheta(e_0\poA(W_j))$ with $|\det\Lambda|_\iota\ge\exp(-\dim
\poZ_\vartheta\cdot r\beta)$ such that $\Lambda(\popr_{\poZ_\vartheta}y)$ is integral at $\iota$
for every $y\in\poA(V;\vartheta)$ that is integral at $\iota$ and orthogonal to the false lifts
for $(e_0,\poZ_\vartheta)$.
\end{definition}

This is the first clause of (R3), with Definition~\ref{po:false} for the words in brackets there.

\subsubsection{Tight frames and products of weights}\label{po:ss-frame}

\begin{lemma}[An integral tight frame]\label{po:frame}
Let $\mu$ be a family of diagrams $(\mu_\tau)_{\tau\in\auPl}$ for $\auGL_N$, with columns of
heights at most $N$, $W_\mu=\bigotimes_\tau W_{\mu_\tau}$ with its lattice $W_\mu(\Z)$ and its
invariant form, and $H=\prod_\tau H_{\mu_\tau}$ the product of the hook numbers of
Lemma~\ref{pk:hook}. There are vectors $\eta_U\in W_\mu(\Z)$, indexed by the fillings $U$ of $\mu$
(for each $\tau$ and each column a set of rows), such that
\[
\sum_U\langle\eta_U,\xi\rangle\,\eta_U=H^2\,\xi\quad(\xi\in W_\mu),\qquad\|\eta_U\|\le H .
\]
The number of fillings is at most $N^{|\mu|}$, and $H\le|\mu|^{3|\mu|}$, where $|\mu|$ is the
total number of boxes.
\end{lemma}
\begin{proof}
In the notation of \S\ref{pk:ss-branch} (for all $\tau$ at once) $W_\mu$ is identified
isometrically, through $\mathsf m$, with $\mathsf C_\mu\subseteq\poT_\mu(\C)$, the vectors $e_U$
form an orthonormal basis of $\poT_\mu(\C)$, $\pi_\mu$ is the orthogonal projection onto
$\mathsf C_\mu$, and $W_\mu(\Z)=\mathsf m(\mathsf C_\mu\cap\poT_\mu(\Z))$ by \eqref{pk:eq-dual}.
Put $\eta_U=H\,\mathsf m(\pi_\mu e_U)=H\,b_U$. By Lemma~\ref{pk:hook}, $H\Sigma_\mu\subseteq
\Sigma_\mu^\#$, so $\eta_U\in W_\mu(\Z)$. For $\xi=\mathsf m(v)$, $v\in\mathsf C_\mu$:
$\sum_U\langle\eta_U,\xi\rangle\eta_U=H^2\,\mathsf m\,\pi_\mu\bigl(\sum_U(e_U,v)e_U\bigr)=H^2\,
\mathsf m(v)$, and $\|\eta_U\|=H\|\pi_\mu e_U\|\le H$. A filling chooses a row for each box.
\end{proof}

\emph{Products.} Let $n=N''+n^\circ$ be a decomposition of the coordinates, $\C^n=\C^{N''}\oplus
\C^{n^\circ}$, and suppose $N''\ge r_+$, $n^\circ\ge r_-$. Let $W''=W_{\mu^+}(\C^{N''})\subseteq
\poT^+(\C^{N''})$ and $W^\circ=W_{\mu^-}(\C^{n^\circ})\subseteq\bigotimes_{\mathfrak c}\Lambda^{
h_{\mathfrak c}}\C^{n^\circ}$ (columns of $\mu^-$), in the tensor models of \S\ref{pk:ss-branch}.
Let $\natural\colon\C^{n^\circ}\to(\C^{n^\circ})^\vee$ be the \emph{linear} map $e_i\mapsto
e_i^\vee$, and $\flat$ the conjugate-linear map of the standard hermitian form; both are applied
to each vector of each factor of a tensor. For $u''\in W''$ and $u^\circ\in W^\circ$ put
\[
u''\odot u^\circ=u''\otimes(u^\circ)^\natural\ \in\poT^+(\C^n)\otimes\poT^-(\C^n),
\]
a bilinear product.

\begin{lemma}\label{po:prod}
\begin{enumerate}
\item[(M1)] $u''\odot u^\circ\in W_{\lambda_V}$; $\|u''\odot u^\circ\|=\|u''\|\,\|u^\circ\|$; and
$\operatorname{diag}(k'',k^\circ)(u''\odot u^\circ)=k''u''\odot\aubar{k^\circ}u^\circ$ for unitary
$k'',k^\circ$.
\item[(M2)] $W''(\Z)\odot W^\circ(\Z)\subseteq W_{\lambda_V}(\Z)$, and for $\sigma\in\auAut(\C)$,
$\sigma(u''\odot u^\circ)=\sigma u''\odot\sigma u^\circ$, so that
$\|\sigma(u''\odot u^\circ)\|=\|\sigma u''\|\,\|\sigma u^\circ\|$.
\item[(M3)] For $x=(x''\,;\,x^\circ)\in M_{n\times r}(\C)$, with $x''$ the rows in $\C^{N''}$ and
$x^\circ$ those in $\C^{n^\circ}$, $\pop_{u''\odot u^\circ}(x)=\pop''_{u''}(x'')\cdot
\pop^\circ_{u^\circ}(x^\circ)$, where $\pop''_{u''}(x'')=\langle u'',\mathrm{ev}^+(\aubar{x''^+})
\rangle$ depends only on the positive columns of $x''$ and $\pop^\circ_{u^\circ}(x^\circ)=\langle
(u^\circ)^\natural,\mathrm{ev}^-(x^{\circ-})^\flat\rangle$ only on the negative columns of
$x^\circ$; both are conjugate-linear in the vector.
\end{enumerate}
\end{lemma}
\begin{proof}
(M1) $W''$ is spanned by the vectors $\bigotimes_{\mathfrak c}y_1\wedge\dots\wedge y_{h_{\mathfrak
c}}$ with $y_i\in\C^{N''}$ and $W^\circ$ by the analogous vectors with $y'_i\in\C^{n^\circ}$; for
such vectors $u''\odot u^\circ=\mathrm{ev}(y,y'^\natural)$, and $y'^\natural_{i'}(y_i)=0$ because
the two sets of coordinates are disjoint. So $u''\odot u^\circ\in W_{\lambda_V}$ by
\eqref{po:eq-ev}. The norm of a tensor product is the product of the norms, $\natural$ maps an
orthonormal basis to an orthonormal basis, and the contragredient of a unitary matrix $k^\circ$
is $\aubar{k^\circ}$, so that $(\aubar{k^\circ}u^\circ)^\natural$ is the image of
$(u^\circ)^\natural$ under $k^\circ$. (M2) The product is given by the same formula over $\Z$;
use the description of the integral structures as intersections with the integral tensors.
(M3) By \eqref{po:eq-poly}, $\pop_{u''\odot u^\circ}(x)=\langle u'',\mathrm{ev}^+(\aubar{x^+})\rangle
\cdot\langle(u^\circ)^\natural,\mathrm{ev}^-(x^-)^\flat\rangle$; the first pairing involves only the
coordinates of $\C^{N''}$ because $u''$ does, the second only those of $\C^{n^\circ}$.
\end{proof}

\subsubsection{Gluing}\label{po:ss-glue}

\emph{The splitting.} Put $N''=r+9m+1$ and $n^\circ=t-9m-1$, and split the coordinates of $V$ as
$V=V''\perp V^\circ$, where $V''$ has the first $N''$ coordinates (so $\iota_0$ maps into $V''$;
write $\iota_0''$ for the corresponding $N''\times r$ matrix) and $V^\circ$ the last $n^\circ$; in
$V^\circ$ we mark the first $r$ coordinates and write $\iota_z$ for the $n^\circ\times r$ matrix
of their inclusion, and $v_z$ for the block of marked rows of an $n^\circ\times r$ matrix $v$. Note
$r\le n^\circ$ and $r_+\le N''$. Let $\auO_{\mathsf M}=\prod_{w\in\mathsf M}\auO_w$; recall that
$\pi_M$ is a uniformiser at every $w\in\mathsf M$ and a unit at $\aubar w$. Let $\Gamma^1$ be the
kernel of $\auGL_r(\auO_E/\pi_M^2)\to\auGL_r(\auO_E/\pi_M)$; for $Y\in\Gamma^1$ let
$\tilde Y\in\auU(W)(\auA_f)$ be an integral lift, equal to $1$ outside $\mathsf M\cup
\aubar{\mathsf M}$; it lies in $J_M$.

For a matrix $s\in M_{3m\times N''}(\auO_E/\pi_M^2)$ let $L''_s=\{u\in\auO_{\mathsf M}^{N''}:
su\equiv0\bmod\pi_M\}$, a lattice with $\pi_M\auO_{\mathsf M}^{N''}\subseteq L''_s$ and with
$\auO_w^{N''}/L''_{s,w}$ of dimension $d_w\le3m$ over $\auF_w$. A matrix $x=(u\,;v)\in
M_{n\times r}(\auO_{\mathsf M})$, with $u$ the rows in $V''$ and $v$ those in $V^\circ$, is
\emph{normalised for $s$} if
\[
v\equiv\iota_z\bmod\pi_M^2\qquad\text{and}\qquad u\in\iota_0''+\pi_ML''_s\ \ (\text{columnwise}).
\]

\begin{lemma}[The glued lattice]\label{po:lattice}
For every $s$ there is $\pog=\pog_s\in\auU(V)(\auA_{E^+,f})$, equal to $1$ outside
$\mathsf M\cup\aubar{\mathsf M}$, such that at the places of $\mathsf M$:
\begin{enumerate}
\item[(a)] $\pog$ and $\pi_M\pog^{-1}$ are integral;
\item[(b)] for $x\in M_{n\times r}(E_{\mathsf M})$: $\pog^{-1}x$ is integral and $\equiv\iota_0
\bmod\pi_M$ if and only if $x$ is integral and $xY$ is normalised for $s$ for some
$Y\in\Gamma^1$; this $Y$ is then unique;
\item[(c)] there is $K_1\in\auGL_{N''}(\auO_{\mathsf M})$ such that, at each $w\in\mathsf M$, all
columns of $\pog^{-1}|_{V''}K_1$ except at most $3m$ are integral, and those have entries in
$\pi_M^{-1}\auO_w$;
\item[(d)] if $k\in\auGL_{N''}(\auO_{\mathsf M})$, $k\equiv1\bmod\pi_M$, $sk\equiv s\bmod\pi_M^2$,
and $k^\circ\in\auGL_{n^\circ}(\auO_{\mathsf M})$, $k^\circ\equiv1\bmod\pi_M^2$, then
$\pog^{-1}\operatorname{diag}(k,k^\circ)\pog$ is integral and $\equiv1\bmod\pi_M$.
\end{enumerate}
\end{lemma}
\begin{proof}
Let $B\in M_{N''}(\auO_{\mathsf M})$ be a matrix whose columns are a basis of $L''_s$, and
$D=\pi_MB^{-1}$; it is integral because $\pi_M\auO^{N''}\subseteq L''_s$, and $Du\equiv0\bmod
\pi_M^2$ if and only if $u\in\pi_ML''_s$. Let $\Delta_z$ be the diagonal $n^\circ\times n^\circ$
matrix with $\pi_M$ at the marked coordinates and $1$ elsewhere, $P$ the permutation matrix that
exchanges the first $r$ coordinates of $V''$ with the marked coordinates of $V^\circ$, and
\[
G=\begin{pmatrix}D&-D\,\iota_0''\,{}^t\iota_z\\0&\Delta_z\end{pmatrix},\qquad
G^{-1}=\begin{pmatrix}D^{-1}&\iota_0''\,{}^t\iota_z\Delta_z^{-1}\\0&\Delta_z^{-1}\end{pmatrix},
\qquad\pog_w=\pi_MG^{-1}P^{-1}\quad(w\in\mathsf M),
\]
with the components at $\aubar{\mathsf M}$ determined by unitarity. (a) $\pi_M\pog^{-1}=PG$ is
integral, and so is $\pi_MG^{-1}$, whose blocks are $B$, $\iota_0''\,{}^t\iota_z$ and
$\pi_M\Delta_z^{-1}$. (b) If $\pog^{-1}x$ is integral then $x=\pog(\pog^{-1}x)$ is integral. For
integral $x=(u;v)$ the condition is $Gx\equiv\pi_MP^{-1}\iota_0\bmod\pi_M^2$, where
$P^{-1}\iota_0=(0;\iota_z)$; that is,
\[
D(u-\iota_0''v_z)\equiv0,\qquad\pi_Mv_z\equiv\pi_M1_r,\qquad v-\iota_zv_z\equiv0\pmod{\pi_M^2}.
\]
If this holds, then $v_z\equiv1\bmod\pi_M$, $Y=v_z^{-1}\bmod\pi_M^2$ lies in $\Gamma^1$, and
$xY=(uY;vY)$ has $vY\equiv\iota_z$ and $D(uY-\iota_0'')\equiv D(u-\iota_0''v_z)Y\equiv0$, so it is
normalised. Conversely, if $xY$ is normalised with $Y\in\Gamma^1$, then $v_z\equiv Y^{-1}\bmod
\pi_M^2$, which determines $Y$, and the three congruences follow in the same way. (c) By the
theorem of elementary divisors $B=K_1\operatorname{diag}(\pi_M1_{d_w},1)K_2$ at $w$ with
$K_1,K_2\in\auGL_{N''}(\auO_w)$; and $\pog^{-1}|_{V''}=P\binom{B^{-1}}0$. (d) Write $k=1+\pi_MX$;
then $sX\equiv0\bmod\pi_M$, so $X=BX_1$ with $X_1$ integral, and $k^\circ=1+\pi_M^2X^\circ$. A
computation with the blocks gives
$\pog^{-1}\operatorname{diag}(k,k^\circ)\pog=1+\pi_M\,P\,G\operatorname{diag}(X,\pi_MX^\circ)
\,\pi_MG^{-1}P^{-1}\cdot\pi_M^{-1}$, and $G\operatorname{diag}(X,\pi_MX^\circ)\pi_MG^{-1}$ has the
blocks $\pi_MX_1B$, $\pi_MX_1\iota_0''\,{}^t\iota_z-\pi_M^2D\iota_0''\,{}^t\iota_zX^\circ
\Delta_z^{-1}$, $0$ and $\pi_M^2\Delta_zX^\circ\Delta_z^{-1}$, all divisible by $\pi_M$.
\end{proof}

\emph{Factorisation of the Weil representation.} $\auHom(W,V)=\auHom(W,V'')\oplus\auHom(W,
V^\circ)$, and
\begin{enumerate}
\item[(W7)] $\mathcal H(\auHom(W,V))=\mathcal H(\auHom(W,V''))\otimes\mathcal H(\auHom(W,
V^\circ))$, $\theta=\theta''\otimes\theta^\circ$, and $\omega_\mu(g,\operatorname{diag}(h,h^\circ))
=\omega''(g,h)\otimes\omega^\circ(g,h^\circ)$, where $\omega''$ and $\omega^\circ$ are Weil
representations of the two pairs: $\omega^\circ=\omega_{\upsilon^{m/2-j},\beta^\circ}$ and
$\omega''=\omega_{\nu''_\vartheta,\beta''}$ in the convention of \S\ref{pc:ss-data}, with
$\nu''_\vartheta=\mu\,\upsilon^{j+(t+k-1)/2}$ and $\beta^\circ,\beta''$ powers of $\upsilon$ times the
characters on $\auU(V)$, so that the characters on $\auU(W)$ multiply to $\mu\upsilon^{\lfloor
n/2\rfloor}$ (Lemma~\ref{psv:char}). In particular $\omega^\circ$ does not depend on $\vartheta$, and
its character on $\auU(W)$ is $\upsilon^{a^\circ}$ with $a^\circ=m/2-j$.
\end{enumerate}
This is the compatibility of the Heisenberg and Weil representations with orthogonal sums, and of
the splittings of~\cite{KudlaSplitting} with it (the splitting over $\auU(W)$ for $\auHom(W,V)$
is the product of those for the two summands, because the characters multiply); the theta
functionals are normalised compatibly. In the Fock model the vector $\pop''_{u''}\pop^\circ_{
u^\circ}$ of Lemma~\ref{po:prod}(M3) is the tensor product of the two polynomial vectors.

\emph{Forms on $V^\circ$, and the kernels.} Let $J^\circ\subset\auU(V^\circ)(\auA_f)$ be the
principal congruence subgroup of level $\pi_M^2$ in the stabiliser of the standard lattice; it is
free by Lemma~\ref{au:free}. Let $\poA^\circ$ be the space of the $\zeta\colon\auU(V^\circ)
(\auA_f)\to W^\circ$ that are right invariant under $J^\circ$ and satisfy $\zeta(\gamma h^\circ)=
\rho_{\mu^-}(\tau(\gamma))\zeta(h^\circ)$ (the rational points acting through the embeddings
$\tau$ themselves); equivalently, the coordinatewise conjugate $\aubar\zeta$ is an algebraic
modular form of weight $\mu^-$ and level $J^\circ$. The inner product is $\langle\zeta,\zeta'
\rangle=\sum_{h^\circ}\langle\zeta(h^\circ),\zeta'(h^\circ)\rangle$ over the classes,
$\|\zeta\|_\sigma^2=\sum_{h^\circ}\|\sigma\zeta(h^\circ)\|^2$, and $\zeta$ is \emph{integral at
$\iota$} if $\aubar\zeta$ is integral at $\aubar\iota$. (By Lemma~\ref{po:prod}(M1) and the
convention (N2), under which $\Gamma$ acts on $W_{\lambda_V}$ through $\tau\circ c$, the block of
$V^\circ$ acts on the second factor of $u''\odot u^\circ$ through $\tau$; this is why the forms on
$V^\circ$ are taken of this type.) Let
\[
\Psi^\circ=\plone[v\equiv\iota_z\bmod\pi_M^2]\otimes(\text{vacuum off }\mathsf M),\qquad
\Psi''_{\vartheta,s}=\plone[u\in\iota_0''+\pi_ML''_s]\otimes\bigotimes_{w\in\mathsf Q}\plone[u
\text{ integral}]\psi_w(\det u')\otimes(\text{vacuum}),
\]
\[
\theta^\circ(g,h^\circ;u^\circ)=\theta^\circ\bigl(\omega^\circ(g,h^\circ)(\pop^\circ_{u^\circ}
\otimes\Psi^\circ)\bigr),\qquad
\theta''_{\vartheta,s}(g,h;u'')=\theta''\bigl(\omega''(g,h)(\pop''_{u''}\otimes\Psi''_{\vartheta,
s})\bigr),
\]
both conjugate-linear in the vector. The \emph{kernel form} $K^\circ(g,\cdot)\in\poA^\circ$ is
defined by $\langle u^\circ,K^\circ(g,h^\circ)\rangle=\theta^\circ(g,h^\circ;u^\circ)$; it does not
depend on $\vartheta$.

\emph{Forms supported on a class.} For $x\in\auU(V)(\auA_f)$ and $v\in W_{\lambda_V}$ let
$\podel[x;v](h)=\sum\chi_\vartheta(k)^{-1}\gamma v$, the sum over the pairs $(\gamma,k)\in\Gamma\times
K_\vartheta$ with $h=\gamma xk$; there are $\#(\Gamma\cap xK_\vartheta x^{-1})$ such pairs if $h$
is in the class of $x$ and none otherwise. Then $\podel[x;v]\in\poA(V;\vartheta)$,
\begin{equation}\label{po:eq-delta}
\langle\podel[x;v],\Xi\rangle=\langle v,\Xi(x)\rangle,\qquad
D_\vartheta\podel[x;v](g)=\aubar{k_{\Phi_\vartheta}(g,x;v)},
\end{equation}
\begin{equation}\label{po:eq-deltasize}
\|\sigma\,\podel[x;v](h)\|\le n_{\mathrm{Aut}}\|\sigma v\|,\qquad\plgb h{\podel[x;v](h)}\le\plgb xv
\quad(\ell\nmid\auNm\mathsf Q),
\end{equation}
and $\podel[\gamma xk;\gamma v]=\chi_\vartheta(k)\,\podel[x;v]$ (the pairs $(\gamma',k')$ for $\gamma xk$ correspond to the pairs $(\gamma'\gamma,kk')$ for $x$). Indeed
$\Xi(x)=\chi_\vartheta(k)^{-1}\gamma\,\Xi(x)$ for each pair with $x=\gamma xk$, so each of the terms of
$\langle\podel[x;v](x),\Xi(x)\rangle$ is $\langle v,\Xi(x)\rangle$; in the same way each term of
$k_{\Phi_\vartheta}(g,x;\podel[x;v](x))$ equals $k_{\Phi_\vartheta}(g,x;v)$, by the invariance of
$\theta$ under $\Gamma$ and the definition of $\chi_\vartheta$; the bounds are (G1), (G2) of
\S\ref{pl:ss-norm}, the unitarity of $\gamma^\sigma$, and $|\sigma\chi_\vartheta(k)|=1=|\chi_\vartheta
(k)|_\iota$.

\begin{lemma}[Slices]\label{po:slice}
Let $\pog=\pog_s$ and $c_\pog=\eta_{\mathsf M}(1,\pog)\ne0$, the scalar of \eqref{po:eq-w5} at
$\mathsf M$; it does not depend on $\vartheta$. For $h\in\auU(V'')(\auA_f)$, $h^\circ\in
\auU(V^\circ)(\auA_f)$, $u''\in W''$, $u^\circ\in W^\circ$ and $g\in\auU(W)(\auA)$,
\[
k_{\Phi_\vartheta}\bigl(g,\operatorname{diag}(h,h^\circ)\pog;u''\odot u^\circ\bigr)=c_\pog\sum_{Y\in
\Gamma^1}\theta''_{\vartheta,s}(g\tilde Y,h;u'')\ \theta^\circ(g\tilde Y,h^\circ;u^\circ).
\]
\end{lemma}
\begin{proof}
By \eqref{po:eq-w5} and Lemma~\ref{po:lattice}(b), at the places of $\mathsf M$,
\begin{align*}
\omega_\mu(1,\pog)\plone[T\equiv\iota_0]&=c_\pog\,\plone[\pog^{-1}T\equiv\iota_0]=c_\pog\sum_{Y\in
\Gamma^1}\plone[TY\text{ normalised for }s]\\
&=c_\pog\sum_{Y\in\Gamma^1}\omega_\mu(\tilde Y,1)\bigl(\plone[u\in
\iota_0''+\pi_ML''_s]\otimes\plone[v\equiv\iota_z\bmod\pi_M^2]\bigr);
\end{align*}
in the last step the scalar of \eqref{po:eq-w5} is $1$ because $\det\tilde Y\equiv1$ at
$\mathsf M$. At $\mathsf Q$ the block datum involves only rows of $V''$, and $\pog$, $\tilde Y$
are $1$ there. So $\omega_\mu(1,\pog)\Phi_\vartheta=c_\pog\sum_Y\omega_\mu(\tilde Y,1)(\Psi''_{
\vartheta,s}\otimes\Psi^\circ)$. Insert this in the definition of $k_{\Phi_\vartheta}$ and use
(W7) and Lemma~\ref{po:prod}(M3).
\end{proof}

\emph{The choices and the first factors.} There is a finite set
$I$ of \emph{choices}, independent of $\vartheta$ (Lemma~\ref{plf:I}); each $i\in I$ has a matrix $s_i\in M_{3m\times
N''}(\auO_E/\pi_M^2)$ (the matrix $\Lambda_i$ of \S\ref{pte:ss-third}), hence $\pog_i=\pog_{s_i}$, $c_i=c_{\pog_i}$, and a level
\[
\begin{aligned}J_{\vartheta,i}=\{k\in K^\circ(V''):\ &k\equiv1\bmod\pi_M,\ s_ik\equiv s_i\bmod\pi_M^2,\ \text{the last }N''-r
\text{ columns of }k\\&\text{are those of }1\bmod\pi_M^2\text{ at }\mathsf M\text{ and}\bmod w\text{ for }w\in\mathsf Q\},\end{aligned}
\]
which is free (Lemma~\ref{au:free}). (This is the group $J'_{\vartheta,i}$ of \S\ref{pte:ss-left}, Remark~\ref{pte:rem-final}(ii).) By Lemma~\ref{po:lattice}(d), $\pog_i^{-1}\operatorname{diag}
(k,k^\circ)\pog_i\in K_\vartheta$ for $k\in J_{\vartheta,i}$, $k^\circ\in J^\circ$; let
$\chi''_{\vartheta}(k)$ be the value of $\chi_\vartheta$ on it (it depends only on $k$ at
$\mathsf Q$). For each twist $\vartheta$ and each $i$ there is a \emph{first factor}
\[
\varphi''_{\vartheta,i}\in\poA_{\mu^+}\bigl(J_{\vartheta,i},\chi''_\vartheta{}^{-1}\bigr)\quad
\text{on }\auU(V''),\qquad\Xi_{\vartheta,i}(g)=\sum_{h\in\Gamma''\backslash\auU(V'')(\auA_f)/
J_{\vartheta,i}}\theta''_{\vartheta,s_i}\bigl(g,h;\varphi''_{\vartheta,i}(h)\bigr),
\]
the second being its theta series on $\auU(W)(\auA)$ (the summand does not depend on the
representative). The first factors $\varphi''_{\vartheta,i}$ are the functions of Definition~\ref{pfq:def-ff} (a
non-zero multiple $\kappa$, depending only on $j$, of the functions of Definition~\ref{ptf:def-ff}, through
Definition~\ref{pte:def-ff}). The character $\chi''_\vartheta$ is
$k\mapsto\chi'_\vartheta(\det\aubar k)$ (Lemma~\ref{ptf:chi}). The \emph{exponent of the denominators} is
\[
l_1=3m\sum_{\tau\in\auPl}\bigl(c^+_\tau+c^-_\tau\bigr),\qquad c^\pm_\tau=\text{the number of
columns of }\mu^\pm_\tau;\qquad l_1\le3mf(4b+m\Delta)\le13fmb .
\]

\begin{lemma}[First factors: integrality and size]\label{po:ff}
The first factors can be chosen so that $\#I\le e^b$ and, for all $i$, $\vartheta$ and $h$:
\begin{enumerate}
\item[(I1)] $\varphi''_{\vartheta,i}$ is integral at every $\iota$ with $\ell(\iota)\nmid\auNm
Q_\vartheta$;
\item[(I2)] $\|\sigma\varphi''_{1,i}(h)\|^2\le\exp(\poAA)$ for every $\sigma\in\auAut(\C)$.
\end{enumerate}
\end{lemma}
\noindent\emph{Proof.} Theorem~\ref{pfq:ff}, for the first factors of Definition~\ref{pfq:def-ff}
(a fixed non-zero multiple, depending only on $j$, of those of Definition~\ref{ptf:def-ff}); $\#I\le e^b$
is Lemma~\ref{plf:I}. Proposition~\ref{po:exchange} holds for these first factors by
Proposition~\ref{pfq:exchange}.


\begin{lemma}[Numbers of classes]\label{po:classes}
Let $X$ be a standard definite space of rank at most $2n_m$ and $J\subseteq K^\circ(X)$ a subgroup
containing the principal congruence subgroup of level $\pi_M^2$. Then $\Gamma\backslash\auU(X)
(\auA_f)/J$ has at most $e^b$ elements.
\end{lemma}
\noindent\emph{Proof.} Lemma~\ref{plf:classes}.


\begin{proposition}[Three exchanges]\label{po:exchange}
Let $H_j(\vartheta)$ be the number \eqref{po:eq-H} below. There are functions $F_i$ $(i\in I)$ on
$\auU(W)(\auA)$, independent of $\vartheta$, a constant $C_{\mathrm x}\ge0$, and for every twist
$\vartheta$ a number $c'_\vartheta\in\C$, such that
\[
c'_\vartheta\,\Xi_{\vartheta,i}=F_i\quad(i\in I),\qquad|c'_\vartheta|^2\,H_j(\vartheta)=
C_{\mathrm x}.
\]
Moreover $yF_iK^\circ(\cdot,h^\circ)$ is integrable on $[\auU(W)]$ for $y\in\poA(W_j)$, and so is
$y\,\Xi_{\vartheta,i}K^\circ(\cdot,h^\circ)$.
\end{proposition}
\begin{proof}
This is proved in \S\ref{pte:sec} and \S\ref{ptf:sec}: the functions $F_i$ are those of \eqref{pte:eq-F}, the
number is $c'_\vartheta=c_3(\vartheta)\gamma(\vartheta)/\#\auGL_{3m}(\auO_E/\pi_M^2)$ and $C_{\mathrm x}=C_3C_{12}/
(\#\auGL_{3m}(\auO_E/\pi_M^2))^2$ (\S\ref{pte:ss-proof}), with the first factors of Definition~\ref{ptf:def-ff}
and the level $J'_{\vartheta,i}$ (Remark~\ref{pte:rem-final}). For the first factors of
Definition~\ref{pfq:def-ff}, which are the ones used from now on, the number is
$c'_\vartheta=c_3(\vartheta)\gamma(\vartheta)/(\aubar\kappa\,\#\auGL_{3m}(\auO_E/\pi_M^2))$ and the constant
$C_{\mathrm x}=C_3C_{12}/(|\kappa|^2(\#\auGL_{3m}(\auO_E/\pi_M^2))^2)$ (Proposition~\ref{pfq:exchange}).
Below, $c'_\vartheta$ and $C_{\mathrm x}$ denote these numbers; the two identities of the proposition
hold for them, and they are the only properties of $c'_\vartheta$, $C_{\mathrm x}$ used.
\end{proof}


\begin{proposition}[Products span]\label{po:span}
If $y\in e_0\poA(W_j)$ and $\int_{[\auU(W)]}y(g)F_i(g)K^\circ(g,h^\circ)\,dg=0$ for all $i\in I$
and all $h^\circ$, then $y=0$.
\end{proposition}
\noindent Proposition~\ref{po:span} is Proposition~\ref{psp:main}, proved in \S\ref{psp:sec} with the inputs
of \S\ref{psq:sec}--\S\ref{psv:sec} (Theorem~\ref{psv:main}).


\begin{proposition}[Periods of false lifts]\label{po:periods}
Let $w$ be off the level and $\Xi\in\mathcal F_w(e_0,\poZ_\vartheta)$ a false lift. Then for every
$s$, every $h\in\auU(V'')(\auA_f)$, $u''\in W''$ and $g\in\auU(W)(\auA)$,
\[
\sum_{h^\circ\in\Gamma^\circ\backslash\auU(V^\circ)(\auA_f)/J^\circ}\bigl\langle u''\odot K^\circ(g,
h^\circ),\ \Xi\bigl(\operatorname{diag}(h,h^\circ)\pog_s\bigr)\bigr\rangle=0 .
\]
\end{proposition}
\begin{proof}[Proof of Proposition~\ref{po:periods}]
This is Theorem~\ref{ppe:main} of \S\ref{ppe:sec}: a false lift has vanishing descent (by
Proposition~\ref{plf:L4}), is killed by the hyperplane sums at the chosen places (dictionary between
labels and characters, Lemma~\ref{ppe:green}), and satisfies one Hecke equation at a split place;
its descents then vanish to $W$ (Proposition~\ref{ppe:bottom}), to $W\oplus\mathbb H^p$ for
$2r+2p+3<n$ (Proposition~\ref{ppe:tower}), and the periods over $\auU(V^\circ)$ vanish
(Proposition~\ref{ppe:rallis}).
\end{proof}


\emph{The gluing map.} $\poA^\oplus=\bigoplus_{i\in I}\poA^\circ$, with $\|\zeta\|_\sigma^2=\sum_i
\|\zeta_i\|_\sigma^2$. For a twist $\vartheta$,
\[
\poglue_\vartheta\colon\poA^\oplus\to\poA(V;\vartheta),\qquad
\poglue_\vartheta(\zeta)=M^{l_1}\sum_{i\in I}\ \sum_{h}\ \sum_{h^\circ}\podel\bigl[\operatorname{diag}
(h,h^\circ)\pog_i;\ \varphi''_{\vartheta,i}(h)\odot\zeta_i(h^\circ)\bigr],
\]
where $h$ runs over $\Gamma''\backslash\auU(V'')(\auA_f)/J_{\vartheta,i}$ and $h^\circ$ over
$\Gamma^\circ\backslash\auU(V^\circ)(\auA_f)/J^\circ$; the summands do not depend on the
representatives, by Lemma~\ref{po:prod}(M1), Lemma~\ref{po:lattice}(d), the rule for
$\podel[\gamma xk;\gamma v]$ and the definition of $\chi''_\vartheta$. Let
\[
\Theta^\circ\colon e_0\poA(W_j)\to\poA^\oplus,\qquad(\Theta^\circ y)_i=M^{l_1}c_i\,\#\Gamma^1
\int_{[\auU(W)]}y(g)\,F_i(g)\,K^\circ(g,\cdot)\,dg .
\]

\begin{proposition}[Properties of the gluing map]\label{po:glue}
Let $e_0\poA(W_j)\ne0$.
\begin{enumerate}
\item[(G1)] $c'_\vartheta\cdot\poglue_\vartheta^*\circ\Theta_\vartheta=\Theta^\circ$ on
$e_0\poA(W_j)$, for every twist $\vartheta$.
\item[(G2)] $\Theta^\circ$ is injective; hence $c'_\vartheta\ne0$, $C_{\mathrm x}>0$, and
$\Theta_\vartheta$ is injective on $e_0\poA(W_j)$.
\item[(G3)] $\poglue_\vartheta(\Theta^\circ y)$ is orthogonal to the false lifts for
$(e_0,\poZ_\vartheta)$, for every $y$.
\item[(G4)] If $\ell(\iota)\nmid\auNm Q_\vartheta$ and $\zeta$ is integral at $\iota$, then
$\poglue_\vartheta(\zeta)$ is integral at $\iota$.
\item[(G5)] $\|\poglue_1(\zeta)\|_\sigma^2\le\exp(2\poAA)\,\|\zeta\|_\sigma^2$ for all $\sigma\in
\auAut(\C)$.
\end{enumerate}
\end{proposition}
\begin{proof}
(G1) Let $\zeta\in\poA^\oplus$. By \eqref{po:eq-descent}, the linearity of $D_\vartheta$,
\eqref{po:eq-delta} and Lemma~\ref{po:slice},
\[
\aubar{D_\vartheta(\poglue_\vartheta\zeta)(g)}=M^{l_1}\sum_ic_i\sum_{Y\in\Gamma^1}\Xi_{\vartheta,i}
(g\tilde Y)\sum_{h^\circ}\theta^\circ\bigl(g\tilde Y,h^\circ;\zeta_i(h^\circ)\bigr)=M^{l_1}\sum_ic_i
\sum_Y\Xi_{\vartheta,i}(g\tilde Y)\,\bigl\langle\zeta_i,K^\circ(g\tilde Y,\cdot)\bigr\rangle .
\]
Since $y$ is right invariant under $\tilde Y\in J_M$ and the measure is invariant,
\[
\langle\zeta,\poglue_\vartheta^*\Theta_\vartheta y\rangle=\langle\poglue_\vartheta\zeta,
\Theta_\vartheta y\rangle=M^{l_1}\sum_ic_i\,\#\Gamma^1\int_{[\auU(W)]}y(g)\,\Xi_{\vartheta,i}(g)\,
\langle\zeta_i,K^\circ(g,\cdot)\rangle\,dg ,
\]
and multiplying by $c'_\vartheta$ gives $\langle\zeta,\Theta^\circ y\rangle$ by
Proposition~\ref{po:exchange}.
(G2) If $\Theta^\circ y=0$ then, $c_i\ne0$, the hypothesis of Proposition~\ref{po:span} holds, so
$y=0$ (Remark~\ref{psv:G2}). If $c'_\vartheta=0$ then $\Theta^\circ=0$ by (G1), against
$e_0\poA(W_j)\ne0$; so $c'_\vartheta\ne0$, and $\Theta_\vartheta$ is injective on $e_0\poA(W_j)$ by
(G1). Finally $C_{\mathrm x}=|c'_\vartheta|^2H_j(\vartheta)$ (Proposition~\ref{po:exchange}) and
$H_j(\vartheta)>0$ (Lemma~\ref{po:Hpos}), so $C_{\mathrm x}>0$. Neither the injectivity of
$\Theta^\circ$ nor Lemma~\ref{po:Hpos} uses $C_{\mathrm x}$, so the argument is not circular.
(G3) Let $\Xi\in\mathcal F_w(e_0,\poZ_\vartheta)$. By \eqref{po:eq-delta},
\[
\langle\poglue_\vartheta\Theta^\circ y,\Xi\rangle=M^{l_1}\sum_{i,h}\ \sum_{h^\circ}\bigl\langle
\varphi''_{\vartheta,i}(h)\odot(\Theta^\circ y)_i(h^\circ),\ \Xi(\operatorname{diag}(h,h^\circ)
\pog_i)\bigr\rangle ,
\]
and $(\Theta^\circ y)_i(h^\circ)$ is an integral over $g$ of multiples of $K^\circ(g,h^\circ)$; for
each $g$ the sum over $h^\circ$ vanishes by Proposition~\ref{po:periods}.
(G4) By \eqref{po:eq-deltasize} and the ultrametric inequality it suffices that
$|M|_\iota^{l_1}\,\plgb{x}{u''\odot u^\circ}\le1$ for $x=\operatorname{diag}(h,h^\circ)\pog_i$,
$u''=\varphi''_{\vartheta,i}(h)$, $u^\circ=\zeta_i(h^\circ)$. Now $\rho(x_\ell)^{-1}=\rho(\pog_{i,
\ell})^{-1}\rho(\operatorname{diag}(h,h^\circ)_\ell)^{-1}$, and by Lemma~\ref{po:prod}(M1) the second
factor maps $\iota(u''\odot u^\circ)$ to a product $a''\odot a^\circ$ in which $\|a''\|=\plgb
h{u''}\le1$ by (I1) and $\|a^\circ\|\le1$ because $\zeta$ is integral at $\iota$ (the block
$h^\circ$ acts on the factor $\poT^-$ through the contragredient, that is, through its components
at the conjugate places; this is why integrality on $V^\circ$ was defined through $\aubar\zeta$
and $\aubar\iota$); the norm of a tensor product of two vectors is the product of the norms.
If $\ell\nmid M$ then $\pog_{i,\ell}=1$ and we are done. Let $\ell\mid M$ and $\tau\in\auPl$. On the
factor $\Lambda^h$ of $\poT^+$ at $\tau$ the element $\pog_i^{-1}$ acts by
$\Lambda^h(\pog_{i,w_\tau}^{-1}|_{V''})$: if $w_\tau\in\mathsf M$ this has entries in
$\pi_M^{-3m}\auO$ by Lemma~\ref{po:lattice}(c) (after the unimodular change $K_1$ at most $3m$
columns are not integral), and otherwise $\pog_{i,w_\tau}^{-1}$ is integral, being $1$ or the
conjugate transpose of $\pog_{i,\aubar w_\tau}$. On a factor $\Lambda^h(\cdot)^\vee$ of $\poT^-$,
$h\le r_-\le m$, it acts by $\Lambda^h({}^t\pog_{i,w_\tau})$, which is integral if
$w_\tau\in\mathsf M$ and has entries in $\aubar{\pi_M}^{-h}\auO$ if $w_\tau\in\aubar{\mathsf M}$.
In all cases a column costs at most $|M|_\iota^{-3m}$, so $\|\rho(\pog_{i,\ell})^{-1}(a''\odot
a^\circ)\|\le|M|_\iota^{-l_1}$.
(G5) For $\vartheta=1$: by Lemma~\ref{po:classes}, $\poA(V;1)$ has at most $e^b$ classes, and
there are at most $\#I\cdot e^b\cdot e^b\le e^{3b}$ terms in $\poglue_1(\zeta)$; by
\eqref{po:eq-deltasize}, Lemma~\ref{po:prod}(M2) and (I2), each has values with
$\|\sigma(\cdot)\|\le n_{\mathrm{Aut}}\,e^{\poAA/2}\|\zeta\|_\sigma$. Hence
$\|\poglue_1\zeta\|_\sigma^2\le e^b\bigl(M^{l_1}e^{3b}n_{\mathrm{Aut}}\bigr)^2e^{\poAA}\|\zeta\|_\sigma^2$,
and $M^{2l_1}e^{7b}n_{\mathrm{Aut}}^2\le\exp(26\cdot22\,fmb\log b+7b+2x^{40})\le\exp(\poAA)$, a
comparison of sizes ($M\le b^{22}$).
\end{proof}

\subsubsection{The comparison operator; proof of Proposition~\ref{au:onelift}}\label{po:ss-proof}

Let $H_\circ=\prod_\tau H_{\mu^-_\tau}$ be the product of the hook numbers of the diagrams
$\mu^-_\tau$ (Lemma~\ref{pk:hook}), and
\[
U_\vartheta=n_{\mathrm{Aut}}\,H_\circ^2\ \poglue_\vartheta\circ\poglue_\vartheta^*\ \in\auEnd\bigl(
\poA(V;\vartheta)\bigr),\qquad\Delta_j(\vartheta)=\det\bigl(\popr_{\poZ_\vartheta}U_\vartheta
|_{\poZ_\vartheta}\bigr),\qquad\poZ_\vartheta=\Theta_\vartheta(e_0\poA(W_j)).
\]
We prove Proposition~\ref{au:onelift} with these $\Delta_j$, with $H_j$ given by \eqref{po:eq-H}
and with $Z_j(\vartheta)=\det(\Theta_\vartheta^*\Theta_\vartheta)$ on $e_0\poA(W_j)$ (the map $u$
of Proposition~\ref{au:factor} is a positive multiple of $\Theta_\vartheta$, independent of
$\vartheta$; this changes only the constant $C_j$). If $d_j=0$ all determinants are $1$ and there
is nothing to prove; let $d_j>0$.

\begin{proof}[Proof of Proposition~\ref{au:onelift}, given the statements of Proposition~\ref{po:glue}]
\emph{Step 1: the identity.} By (G2), $\Theta_\vartheta$ is injective, so $\dim\poZ_\vartheta=d_j$
and Lemma~\ref{po:det} applies; by (G1),
\[
Z_j(\vartheta)\,\Delta_j(\vartheta)=\det\bigl(\Theta_\vartheta^*U_\vartheta\Theta_\vartheta\bigr)=
\det\bigl(n_{\mathrm{Aut}}H_\circ^2\,|c'_\vartheta|^{-2}\,\Theta^{\circ*}\Theta^\circ\bigr)=
\Bigl(\frac{n_{\mathrm{Aut}}H_\circ^2}{C_{\mathrm x}}\Bigr)^{d_j}\det(\Theta^{\circ*}\Theta^\circ)
\cdot H_j(\vartheta)^{d_j},
\]
by Proposition~\ref{po:exchange}. So the identity holds with $C_j=(n_{\mathrm{Aut}}H_\circ^2H_j(1)
/C_{\mathrm x})^{d_j}\det(\Theta^{\circ*}\Theta^\circ)$, which is positive because $\Theta^\circ$
is injective; and $\Delta_j(\vartheta)\ne0$.

\emph{Step 2: a sum of operators of rank one.} Let $\ell$ be a prime. By
Corollary~\ref{pl:wacor}(a) choose representatives $x_c$ of the classes $c$ of
$\Gamma^\circ\backslash\auU(V^\circ)(\auA_f)/J^\circ$ that are integral above $\ell$, and let
$\eta_U\in W^\circ(\Z)$ be the frame of Lemma~\ref{po:frame} for $\mu^-$ and $N=n^\circ$. For
$a=(i,c,U)$ let $f_a\in\poA^\oplus$ be the element whose component $i$ is the form supported on
the class $c$ with value $\eta_U$ at $x_c$, the other components being $0$. Since $J^\circ$ is
free, $\poA^\circ$ is the orthogonal sum of copies of $W^\circ$ indexed by the classes; so
\[
\sum_a\langle f_a,\zeta\rangle f_a=H_\circ^2\,\zeta\quad(\zeta\in\poA^\oplus),\qquad
\|f_a\|_\sigma=\|\eta_U\|\le H_\circ\quad(\sigma\in\auAut(\C)),
\]
and $f_a$ is integral at every $\iota$ with prime $\ell$: $\aubar{f_a}$ has the value $\eta_U$, an
integral vector, at a representative that is integral above $\ell$, and (G1), (G2) of
\S\ref{pl:ss-norm} apply as in Corollary~\ref{pl:wacor}(b). Consequently, for $\Xi\in\poA(V;
\vartheta)$,
\begin{equation}\label{po:eq-rankone}
U_\vartheta\Xi=n_{\mathrm{Aut}}\,\poglue_\vartheta\Bigl(\sum_a\langle f_a,\poglue_\vartheta^*\Xi
\rangle f_a\Bigr)=\sum_an_{\mathrm{Aut}}\,\langle\poglue_\vartheta f_a,\Xi\rangle\ \poglue_\vartheta
f_a .
\end{equation}

\emph{Step 3: $U_\vartheta$ keeps integrality.} Let $\ell=\ell(\iota)\nmid\auNm Q_\vartheta$ and
let $\Xi$ be integral at $\iota$. By (G4) the forms $\poglue_\vartheta f_a$ are integral at $\iota$
and at $\aubar\iota$; by Lemma~\ref{po:naut}, $n_{\mathrm{Aut}}\langle\poglue_\vartheta f_a,\Xi
\rangle\in\auO_\iota$. By \eqref{po:eq-rankone}, $U_\vartheta\Xi$ is integral at $\iota$.

\emph{Step 4: $U_\vartheta$ maps lifts to forms orthogonal to the false lifts.} For
$z=\Theta_\vartheta y\in\poZ_\vartheta$, (G1) gives $U_\vartheta z=n_{\mathrm{Aut}}H_\circ^2
c'^{-1}_\vartheta\,\poglue_\vartheta(\Theta^\circ y)$; apply (G3).

\emph{Step 5: proof of {\rm(b)}.} Let $\ell(\iota)\nmid\auNm Q_\vartheta$ and let $\Lambda$ be as
in Definition~\ref{po:denom}, for $\beta$. Let $\alpha$ be an eigenvalue of the endomorphism
$S=\Lambda\circ\popr_{\poZ_\vartheta}U_\vartheta|_{\poZ_\vartheta}$ of $\poZ_\vartheta$, with
eigenvector $z\ne0$. The gauges of the values of $z$ at representatives of the classes are
absolute values of elements of $\overline\Q_\ell$, not all $0$; multiplying $z$ by a suitable
constant we may assume that $z$ is integral at $\iota$ and that $az$ is integral at $\iota$ only
if $|a|_\iota\le1$. By Steps 3 and 4, $U_\vartheta z$ is integral at $\iota$ and orthogonal to the
false lifts; so $Sz=\alpha z$ is integral at $\iota$, and $|\alpha|_\iota\le1$. By
Lemma~\ref{po:eig}, $|\det\Lambda|_\iota\,|\Delta_j(\vartheta)|_\iota=|\det S|_\iota\le1$, and
$|\det\Lambda|_\iota\ge\exp(-d_jr\beta)$.

\emph{Step 6: proof of {\rm(a)}.} Let $\vartheta=1$ and $\sigma\in\auAut(\C)$. By
\eqref{po:eq-rankone}, \eqref{po:eq-cs}, (G5) and Step 2,
\[
\|U_1\Xi\|_\sigma\le n_{\mathrm{Aut}}\sum_a\|\poglue_1f_a\|_{c\sigma c}\,\|\Xi\|_\sigma\,
\|\poglue_1f_a\|_\sigma\le n_{\mathrm{Aut}}\cdot\#\{a\}\cdot e^{2\poAA}H_\circ^2\cdot\|\Xi\|_\sigma
\le e^{3\poAA}\|\Xi\|_\sigma ;
\]
indeed $\#\{a\}\le\#I\cdot e^b\cdot(n^\circ)^{|\mu^-|}$ by Lemmas~\ref{po:ff}, \ref{po:classes}
and~\ref{po:frame}, $H_\circ\le|\mu^-|^{3|\mu^-|}$, $|\mu^-|\le fm^2\Delta\le x^{60}$, and
$n_{\mathrm{Aut}}\le\exp(x^{40})$; so $n_{\mathrm{Aut}}\#\{a\}H_\circ^2\le e^{3b}\le e^{\poAA}$.
Let $\alpha$ be an eigenvalue of $\popr_{\poZ_1}U_1|_{\poZ_1}$ with eigenvector $z$. There is a
prime $w$ off the level, and $\poZ_1\subseteq\varrho_V(e_0)\poA(V;1)\cap\ker P_w(X_w)$ by
Lemma~\ref{po:facts}; by Step 4, $U_1z$ is orthogonal to $\mathcal F_w(e_0,\poZ_1)$. By
Proposition~\ref{po:proj},
\[
|\sigma\alpha|\,\|z\|_\sigma=\|\popr_{\poZ_1}U_1z\|_\sigma\le\|U_1z\|_\sigma\le e^{3\poAA}
\|z\|_\sigma ,
\]
and $\|z\|_\sigma>0$. By Lemma~\ref{po:eig}, $|\sigma\Delta_j(1)|\le\exp(3d_j\poAA)$.
\end{proof}

\subsubsection{Hecke values; proof of Proposition~\ref{au:hecke}}\label{po:ss-hecke}

\emph{The pairs.} We use the characters fixed with the admissible data in
Definition~\ref{pfq:def-chars}: $\chi_0=\chi_M$, of weight $0$, whose tame part $\omega=(\omega_w)$ is
trivial off $\mathsf M$ and has order $q_2>8m$ at every $w\in\mathsf M$; the tame part of $\mu_j^0$ is
$\omega^j$ (Proposition~\ref{pfq:tame}). For a finite set $S$ of primes of $E$ and a character $\nu$ of $\auU(1)$,
$L^S(s,\nu)$ is the $L$-function of $x\mapsto\nu(x/\aubar x)$ without the Euler factors at
$S\cup\aubar S$ (for $s=1$, the value of the continuation; the characters below are not trivial),
and $\auNm S=\prod_{w\in S}\auNm w$. The three pairs of the $j$-th lift are, with
$n'=j+(t-11m-1)/2$,
\[
(\nu_\ast,z_\ast)=\bigl(\mu_j^0\upsilon^{n'},\ \tfrac{7m}2+\tfrac k2\bigr),\qquad
(\nu_+,z_+)=\bigl(\upsilon^{-m}\chi_0^{\,j},\ m\bigr),\qquad
(\nu_-,z_-)=\bigl(\upsilon^{-m/2}\chi_0^{-j},\ \tfrac m2\bigr),
\]
with the signs $s_\ast=s_+=1$, $s_-=-1$ (here $\upsilon$ is restricted to $\auU(1)$; it has weight
$1$ and trivial tame part). They are critical pairs, of weights $-(\kappa_j(\tau)-n')$, $-m$,
$-m/2$ and tame parts $\omega^j$, $\omega^j$, $\omega^{-j}$. Put
\begin{equation}\label{po:eq-H}
H_j(\vartheta)=\prod_{i\in\{\ast,+,-\}}\Bigl((\auNm\mathsf S_\vartheta)^{z_i^2}\prod_{z=1}^{z_i}
\bigl|L^{\mathsf S_\vartheta}\bigl(z,(\nu_i\vartheta^{s_i})^{-1}\bigr)\bigr|^2\Bigr).
\end{equation}
For a critical pair $p=(\nu,z)$, a twist $\vartheta$ and a sign $s$ the \emph{ratio} is
$\rho_p(\vartheta,s)=|L^{\mathsf S_\vartheta}(z,(\nu\vartheta^{s})^{-1})|/|L^{\mathsf M}(z,
\nu^{-1})|$, and $\rho_p(\vartheta)=\rho_p(\vartheta,1)$; this is the quantity of (R2) (at the
places of $\mathsf Q$ the twisted character is ramified, so no Euler factor is removed there).
More precisely, let $q=(\nu,z)$ be one of the pairs $(\nu_i,z)$ and $1\le j\le m$. The characters
$\nu$ and $\nu\theta_\iota^{\pm j}$ are ramified at every place of $\mathsf S_{\theta_\iota}\cup
\aubar{\mathsf S_{\theta_\iota}}$: at $\mathsf M\cup\aubar{\mathsf M}$ because their tame part is
$\omega^{\pm j}\ne1$, and at $Q\cup\aubar Q$ because $\nu$ is unramified there while $\theta_\iota^{\pm j}$
has tame part $\psi_w^{\pm j}\ne1$, $\theta_\iota$ being an $m$-fold twist (\S\ref{ss:au-spaces}); the
character $x\mapsto\nu(x/\aubar x)$ is ramified at $\aubar w$ exactly when it is ramified at $w$. Hence
$L^{\mathsf M}(z,\nu^{-1})=L(z,\nu^{-1})$ and $L^{\mathsf S_{\theta_\iota}}(z,(\nu\theta_\iota^{\pm j})^{-1})=
L(z,(\nu\theta_\iota^{\pm j})^{-1})$, and $\rho_q(\theta_\iota^{\,j},\pm1)$ is exactly the ratio bounded in (R2)
with $n=j$.

\begin{definition}[Main pairs]\label{po:main}
A critical pair $p=(\nu,z)$, with $\nu$ of weight $-c$, is a \emph{main pair for $j$} if
$m<z\le r$, the tame part of $\nu$ is $\omega^j$, and $\kappa_j(\tau)+r-t\le c(\tau)\le
\kappa_j(\tau)$ for all $\tau$. It is \emph{flat} if $c(\tau)=z$ for all $\tau$.
\end{definition}

\begin{lemma}\label{po:Hratio}
With $e=z_\ast^2+z_+^2+z_-^2$,
\[
\frac{H_j(\vartheta)}{H_j(1)}=(\auNm Q_\vartheta)^{e}\prod_{i\in\{\ast,+,-\}}\ \prod_{z=1}^{z_i}
\rho_{(\nu_i,z)}(\vartheta,s_i)^2 .
\]
All $5m+k/2$ pairs $(\nu_i,z)$ are ramified critical pairs, and the $k/2$ pairs $(\nu_\ast,z)$,
$7m/2<z\le z_\ast$, are main pairs for $j$.
\end{lemma}
\begin{proof}
The identity is the definition, with $\auNm\mathsf S_\vartheta=\auNm\mathsf M\cdot\auNm
Q_\vartheta$. The numbers of pairs are $z_\ast+z_++z_-=5m+k/2$ ($k=r+1-m$ is even). Each
$(\nu_i,z)$ is a critical pair: $1\le z\le z_i\le t$; for $i=\pm$, $c=z_i\ge z$; for $i=\ast$,
$c(\tau)=\kappa_j(\tau)-n'\ge z_\ast$ because $\kappa_j(\tau)=\kappa(\tau)+(j+\frac12)\Delta\ge\Delta\ge3t$ (\S\ref{ss:au-notation}), and $c(\tau)\le
\kappa_j(\tau)\le4b$. They are ramified because $\omega_w^{\pm j}\ne1$ ($1\le j\le m$, the order
of $\omega_w$ exceeds $8m$). For $7m/2<z\le z_\ast$: $m<z\le z_\ast\le r$ (indeed $z_\ast=(r+1+6m)/2\le r$, since $r=k+m-1$ and
$k\ge6m+2$ by (S7) of Lemma~\ref{an:lem-sizes}), the tame part is
$\omega^j$, and $0\le n'\le t-r$.
\end{proof}

\begin{lemma}[Positivity of $H_j$]\label{po:Hpos}
Let $\nu\ne1$ be an automorphic character of $\auU(1)$, $S$ a finite set
of primes of $E$, and $z\ge1$ an integer. Then $L^S(z,\nu^{-1})\ne0$. In particular, for every
twist $\vartheta$ and every $i\in\{\ast,+,-\}$ the values $L^{\mathsf S_\vartheta}(z,(\nu_i
\vartheta^{s_i})^{-1})$, $1\le z\le z_i$, are non-zero, and $H_j(\vartheta)>0$.
\end{lemma}
\begin{proof}
Since $\auU(1)$ is anisotropic over $E^+$, $\auU(1)(E^+)\backslash\auU(1)(\auA_{E^+})$ is compact,
so $\nu$ is unitary, and so is the Hecke character $\chi(x)=\nu^{-1}(x/\aubar x)$ of
$\auA_E^\times/E^\times$. The map $x\mapsto x/\aubar x$ from $E_v^\times=(E\otimes E^+_v)^\times$ to
$\auU(1)(E^+_v)$ is surjective at every place $v$ of $E^+$ (Hilbert's Theorem~90), and maps
$(\auO_E\otimes\auO_{E^+,v})^\times$ onto $\auU(1)(\auO_{E^+,v})$ at every finite $v$ unramified in
$E$ (Hilbert~90 for the unramified extension of local rings); hence $x\mapsto x/\aubar x$ maps
$\auA_E^\times$ onto $\auU(1)(\auA_{E^+})$, and $\chi\ne1$ because $\nu\ne1$. Moreover $\chi$ is not of the
form $|\cdot|_E^{it}$ with $t\in\R$: $\chi$ is trivial on $\auA_{E^+}^\times$, where $x/\aubar x=1$, while
$|\cdot|_E^{it}$ restricts to $|\cdot|_{E^+}^{2it}$ there, which is trivial only for $t=0$; and $\chi\ne1$.

For real $s>1$ the Euler product $L^S(s,\chi)=\prod_{w\notin S,\ \chi\text{ unr.\ at }w}(1-\chi
(\varpi_w)q_w^{-s})^{-1}$ converges absolutely, since $|\chi(\varpi_w)q_w^{-s}|=q_w^{-s}$ and
$\sum_wq_w^{-s}\le2f\zeta(s)<\infty$; an absolutely convergent product of non-zero factors is
non-zero, so $L^S(z,\chi)\ne0$ for $z\ge2$ (Lemma~\ref{psu:euler} gives the quantitative form).
The point $z=1$ is on the edge of the region of absolute convergence. For a unitary Hecke character
$\lambda$ of $E$ write $L(s,\lambda)=\prod_{w\ \text{finite},\ \lambda\text{ unr.\ at }w}(1-\lambda
(\varpi_w)q_w^{-s})^{-1}$ (no Euler factor at the ramified primes, no archimedean factors), and, for a
finite set $R$ of primes, $L_R(s,\lambda)$ for the same product over $w\notin R$. The only analytic input
is Hecke's continuation theorem \cite{Hecke1920}: $L(s,1)=\zeta_E(s)$ extends meromorphically
to $\C$ with a single, simple pole at $s=1$ \cite[Chapter~VII, Corollary~(5.11)]{Neukirch}, and $L(s,\lambda)$
extends to an entire function for every unitary Hecke character $\lambda\ne1$ trivial on
$\auA_{E^+}^\times$ (the only ones used below). Indeed, regard $\lambda$ as a Gr\"o{\ss}encharakter as in
\cite[Chapter~VII, \S6]{Neukirch}; its component at infinity is trivial on the positive reals of
$E^+\otimes\R\subset E\otimes\R$. If $\lambda$ is ramified or its component at infinity is non-trivial, the
completed $L$-function of the associated primitive character is entire by \cite[Chapter~VII,
Theorem~(8.5)]{Neukirch}. Otherwise $\lambda$ is a non-trivial character of the ideal class group, and its
completed $L$-function is entire by \cite[Chapter~VII, the remark after Corollary~(5.10)]{Neukirch}. In both
cases dividing by the archimedean factor, which has no zeros, gives the claim (the primitive and the
original character differ at finitely many Euler factors, which are handled as below). Since the finitely many removed factors
$(1-\lambda(\varpi_w)q_w^{-s})$ are entire and non-zero on $\mathrm{Re}\,s>0$, the same holds for
$L_R$. From this we deduce $L(1,\chi)\ne0$ (this is the case $s=1$ of Hecke's non-vanishing theorem on
$\mathrm{Re}\,s=1$; for characters of finite order it is \cite[Chapter~VII, Lemma~(13.3)]{Neukirch}; see also \cite[Theorem~2.4]{FiteSatoTate}). Let $R$ be the set of
primes at which $\chi$ ramifies. First, $\chi^2$ is not of the form $|\cdot|_E^{it}$ unless $\chi^2=1$, by
the restriction argument above ($\chi^2$ is trivial on $\auA_{E^+}^\times$).

\emph{Case $\chi^2\ne1$.} For real $s>1$ put $Z(s)=L_R(s,1)^3L_R(s,\chi)^4L_R(s,\chi^2)$. Taking
logarithms of the absolutely convergent Euler products, with $\chi(\varpi_w)=e^{i\theta_w}$,
\[
\log|Z(s)|=\sum_{w\notin R}\sum_{k\ge1}\frac{q_w^{-ks}}{k}\bigl(3+4\cos k\theta_w+\cos2k\theta_w\bigr)\ge0,
\]
because $3+4\cos\phi+\cos2\phi=2(1+\cos\phi)^2$. So $|Z(s)|\ge1$ for $s>1$. Near $s=1$, $L_R(s,1)^3$
has a pole of order $3$, $L_R(s,\chi^2)$ is holomorphic, and if $L(1,\chi)=L_R(1,\chi)=0$ then
$L_R(s,\chi)^4$ has a zero of order $\ge4$; then $Z(s)\to0$ as $s\to1^+$, a contradiction.

\emph{Case $\chi^2=1$.} Then $\chi(\varpi_w)\in\{\pm1\}$ for $w\notin R$, and the Euler factor of
$D(s)=L_R(s,1)L_R(s,\chi)$ at $w$ is $(1-x)^{-2}=\sum_k(k+1)x^k$ or $(1-x^2)^{-1}=\sum_kx^{2k}$, with
$x=q_w^{-s}$; both have non-negative coefficients and coefficient $\ge1$ at every even power. Hence
$D(s)=\sum_{n\ge1}a_nn^{-s}$ ($s>1$) with $a_n\ge0$, and, with $d=[E:\Q]$, $a_{m^{2d}}\ge1$ for every
integer $m\ge1$ prime to the primes below $R$ (the ideal $(m)^2$ has norm $m^{2d}$ and the coefficient of
$\mathrm N\mathfrak a^{-s}$ is multiplicative in $\mathfrak a$). Suppose $L(1,\chi)=0$. Then the zero
cancels the simple pole of $L_R(s,1)$ and $D$ is holomorphic on $\mathrm{Re}\,s>0$. Put
$A(x)=\sum_{n\le x}a_n\ge0$; for $s>1$, $D(s)/s=\int_1^\infty A(x)x^{-s-1}dx$. By Landau's theorem
\cite[\S15.1, Lemma~15.1]{MV} (applied to $\int_1^\infty A(x)x^{-u}dx$, $u=s+1$) this integral is not
analytic at its abscissa of convergence, while $D(s)/s$ continues analytically to $\mathrm{Re}\,s>0$;
so the integral converges for every real $s>0$. Since $\int_1^\infty A(x)x^{-s-1}dx\ge\sum_na_n\int_n^\infty
x^{-s-1}dx=s^{-1}\sum_na_nn^{-s}$, the series $\sum_na_nn^{-s}$ converges for $s>0$; but at $s=1/(2d)$
it is $\ge\sum_{m}m^{-1}$ over $m$ prime to a finite set of primes, which diverges. So $L(1,\chi)\ne0$. Finally
$L^S(s,\chi)=L(s,\chi)\prod_{w\in S\cup\aubar S,\ \chi\text{ unr.\ at }w}(1-\chi(\varpi_w)q_w^{-s})$,
and at $s=1$ each removed factor satisfies $|1-\chi(\varpi_w)q_w^{-1}|\ge1-q_w^{-1}>0$; so removing
the Euler factors at $S$ creates no zero, and $L^S(1,\chi)\ne0$.

For the pairs of \eqref{po:eq-H}: $\nu_i\vartheta^{s_i}\ne1$, for two independent reasons. Its weight
is that of $\nu_i$, namely $-c$ with $c(\tau)\ge z_i\ge1$ (the twist has weight $0$; proof of
Lemma~\ref{po:Hratio}), so it is non-trivial on $\auU(1)(E^+\otimes\R)$; and its tame part at
$w\in\mathsf M$ is $\omega_w^{\pm j}\ne1$ (the twist has trivial tame part at $\mathsf M$,
\S\ref{po:ss-set}). The first part applies with $S=\mathsf S_\vartheta$. Since also $\auNm\mathsf S_\vartheta\ge1$, every
factor of \eqref{po:eq-H} is positive.
\end{proof}

\emph{The lift of a character.} Let $p=(\nu,z)$ be a critical pair with $z\ge2$ and $\vartheta$ a
twist. Apply \S\ref{po:ss-set} to the positive definite line $E$ in place of $W$ (so $r=1$, there
is no $\mu^-$, and the unitary group is $\auU(1)$), to the standard space $V^{(z)}$ of rank $2z$,
to the character $\nu\vartheta$, and to the weight with one row of length $l_\tau=c(\tau)-z$ at
$\tau$: this gives the block datum $\Phi^{(1)}_\vartheta$, the level, the space
$\poA(V^{(z)};p,\vartheta)$ with the action of $\auG_1=\prod_{w\in\mathsf M}\auF_w^\times$, the
numbers $X_w$, the lift $\Theta^{(1)}_\vartheta$ of functions on the finite set
$\auU(1)(E^+)\backslash\auU(1)(\auA_{E^+,f})/K^1_M$ ($K^1_M$ the units $\equiv1$ at $\mathsf M$)
and its adjoint, the descent $D^{(1)}_\vartheta$. The \emph{normalised lift} is
\[
\varphi_{p,\vartheta}=c_{\mathrm{vac}}^{\,2z}\,(\pi_Ms_0)^{\sum_\tau l_\tau}\cdot[K^1:K^1_M]\cdot
\Theta^{(1)}_\vartheta(\plone)\ \in\poA(V^{(z)};p,\vartheta),
\]
with the constants $c_{\mathrm{vac}}$, $s_0$ of \S\ref{ps:ss-weil}: its values are sums over the
classes of $\auU(1)$ of vectors of jets at scale $\pi_Ms_0$ of theta functions of translates of
the integral vacuum, as in \eqref{ps:eq-kernel}. Let $e_1\in\C[\auG_1]$ be the mean over $\auG_1$.

\begin{definition}[Test forms; the clause of {\rm(R3)} on main pairs]\label{po:R3}
A form $\Xi\in\poA(V^{(z)};p,\vartheta)$ is a \emph{test form at $\ell$} if it is integral at
every homomorphism $\C\to\C_\ell$ and, for every $w$ off the level, orthogonal to
\[
\varrho(e_1)\poA(V^{(z)};p,\vartheta)\ \cap\ \ker P_w(X_w)\ \cap\ \ker D^{(1)}_\vartheta
\]
(the false lifts for $(e_1,(\ker D^{(1)}_\vartheta)^\perp)$ of Definition~\ref{po:false}). The
second clause of {\rm(R3)} in Definition~\ref{def:au-regular} is: for every $\iota$ with prime
$\ell$, every main pair $p$ for $j$, and all test forms $\Xi,\Xi'$ at $\ell$ for
$(p,\theta_\iota^{\,j})$,
\[
\bigl|\langle\popr_\varphi\Xi,\popr_\varphi\Xi'\rangle\bigr|_\iota\le e^{\beta_\ell},\qquad
\popr_\varphi\text{ the orthogonal projection onto }\C\,\varphi_{p,\theta_\iota^{\,j}} .
\]
\end{definition}

\begin{proposition}[The norm of the lift is an $L$-value]\label{po:rallis}
Let $p=(\nu,z)$ be a critical pair with $z\ge2$. There is $A_p>0$ such that for every twist
$\vartheta$
\[
\|\varphi_{p,\vartheta}\|^2=A_p\,(\auNm Q_\vartheta)^{2z-1}\,L^{\mathsf S_\vartheta}\bigl(z,(\nu
\vartheta)^{-1}\bigr);\qquad\text{hence}\quad\rho_p(\vartheta)=\frac{\|\varphi_{p,\vartheta}
\|^2}{(\auNm Q_\vartheta)^{2z-1}\,\|\varphi_{p,1}\|^2} .
\]
\end{proposition}
\noindent(The $L$-value is positive: the character $x\mapsto(\nu\vartheta)^{-1}(x/\aubar x)$ of
$E$ is carried to its complex conjugate by $c$, the Euler factors at $w\ne\aubar w$ come in
conjugate pairs, and at $w=\aubar w$, which is unramified over $E^+$, the factor is
$(1-\auNm w^{-z})^{-1}$.)
\noindent Proved in Proposition~\ref{pex:rallis}, for all tame characters of modulus $Q$.


\begin{proposition}[The untwisted norm at all conjugates]\label{po:arch}
For a main pair $p$ for $j$ the constant family $\|\varphi_{p,1}\|^2$ has height at most
$\aubud-\sum_{\tau\in\auPl}\ausinf(\kappa(\tau))$.
\end{proposition}
\noindent Proved in Theorem~\ref{psu:arch}, with the law of Galois conjugation (CL) of
\S\ref{psu:ss-arch} proved in Theorem~\ref{pcj:CL}. The multiplier is $I_q(1)/\|\varphi_{p,1}\|^2$ with
$q$ flat of the same rank; its integrality, the archimedean factor $\prod_\tau l_\tau!$ from
Proposition~\ref{pex:rallis} and the Euler product bounds (Lemma~\ref{psu:euler}) are proved there.


\begin{proposition}[A test form with a pairing of least valuation]\label{po:small}
Let $p$ be a main pair for $j$ and $q_2=(\upsilon^{-2}\chi_0,2)$, a flat ramified critical pair.
There is $c_0\in\C^\times$ with $|\sigma c_0|\le\exp(fmb/k+3x^{35})$ for all $\sigma\in\auAut(\C)$
such that: for every prime $\ell$, every homomorphism $\jmath\colon\C\to\C_\ell$ and every twist
$\vartheta$ with $\ell\nmid\auNm Q_\vartheta$ there is a test form $\Xi$ at $\ell$ for
$(p,\vartheta)$ with
\[
\bigl|c_0\,\rho_{q_2}(\vartheta)\bigr|_\jmath\cdot\ell^{-F(\ell)}\le\bigl|\langle\Xi,\varphi_{p,
\vartheta}\rangle\bigr|_\jmath,\qquad F(\ell)=\sum_{\sigma\in\auPl}s_\ell(\kappa(\sigma),\sigma).
\]
\end{proposition}
\noindent\emph{Proof.} Theorem~\ref{psu:small} (with Theorem~\ref{psm:attained} and
Corollary~\ref{pst:attained} for the sharpness of the count). The test form is the triple class sum
\eqref{psu:eq-Xi} of point forms against the normalised lifts of three flat pairs; its pairing with
$\varphi_{p,\vartheta}$ is computed in Proposition~\ref{psu:pairing}, its integrality is
Lemma~\ref{psu:integral}, and its orthogonality to the false lifts (Lemma~\ref{psu:orthfalse}) is the
rank-one case of Theorem~\ref{ppe:main} (Proposition~\ref{psu:periods}, Remark~\ref{ppe:rem-rankone}).
The constant is $3x^{35}$ because one integral value $I_q(1)$ is bounded by $\exp(x^{35})$
(Lemma~\ref{psu:I}(c)).


\begin{proposition}[The height of a main pair]\label{po:mainheight}
Let $\theta$ be a regular
family (Definition~\ref{def:au-regular}, with Definitions~\ref{po:denom} and~\ref{po:R3}) and $p$
a main pair for $j$. Then $\iota\mapsto\rho_p(\theta_\iota^{\,j})^{-1}$ has height at most
$6\,\aubud-\auhdir$.
\end{proposition}
\begin{proof}
Fix $\iota$, with prime $\ell$, and put $\vartheta=\theta_\iota^{\,j}$, $\varphi=\varphi_{p,
\vartheta}$, $\rho_2=\rho_{q_2}(\vartheta)>0$; by (R1), $\ell\nmid\auNm Q_\vartheta$. Let $\Xi$
and $\Xi'$ be the test forms of Proposition~\ref{po:small} for the homomorphisms $\iota$ and
$\aubar\iota$. Since $|\langle\Xi',\varphi\rangle|_{\aubar\iota}=|\langle\varphi,\Xi'\rangle|_\iota$
and $|c_0\rho_2|_{\aubar\iota}=|\aubar{c_0}\rho_2|_\iota$, multiplying the two inequalities gives
\begin{equation}\label{po:eq-two}
\bigl|c_0\aubar{c_0}\,\rho_2^2\bigr|_\iota\,\ell^{-2F(\ell)}\le\bigl|\langle\Xi,\varphi\rangle
\langle\varphi,\Xi'\rangle\bigr|_\iota ;
\end{equation}
in particular the right side is not $0$. Now $\popr_\varphi\Xi=\|\varphi\|^{-2}\langle\varphi,\Xi
\rangle\varphi$, so $\langle\popr_\varphi\Xi,\popr_\varphi\Xi'\rangle=\langle\Xi,\varphi\rangle
\langle\varphi,\Xi'\rangle/\|\varphi\|^2$, and by Proposition~\ref{po:rallis}
\[
\rho_p(\vartheta)^{-1}=(\auNm Q_\vartheta)^{2z-1}\cdot\|\varphi_{p,1}\|^2\cdot\langle\popr_\varphi
\Xi,\popr_\varphi\Xi'\rangle\cdot\rho_2^{-2}\cdot\frac{\rho_2^2}{\langle\Xi,\varphi\rangle\langle
\varphi,\Xi'\rangle} .
\]
We bound the heights of the five families of $\iota$. The first is a positive integer prime to
$\ell(\iota)$ and costs nothing. The second has height at most $\aubud-\sum_\tau\ausinf(\kappa
(\tau))$ (Proposition~\ref{po:arch}). The third has height at most $\aubud$: by the second clause
of (R3) its absolute value at $\iota$ is at most $e^{\beta_\ell}=\ell^{\beta_\ell/\log\ell}$, and
$\sum_{\ell\in T}\beta_\ell\le\aubud$. The fourth has height at most $2x_1$ by (R2), $q_2$ being
ramified. For the fifth take the multiplier $c_0\aubar{c_0}$: its conjugates are at most
$\exp(2fmb/k+6x^{35})$, and by \eqref{po:eq-two} the product has absolute value at most
$\ell^{2F(\ell)}$ at $\iota$; by Lemma~\ref{au:gains}(a), $\sum_{\ell\in T}2F(\ell)\log\ell\le
\sum_\sigma\ausfin(\kappa(\sigma),\sigma)$. Adding, and recalling the definition of $\auhdir$, the
height is at most
\[
2\,\aubud+2x_1+2\bigl(fmb/k+3x^{35}\bigr)-\auhdir\le6\,\aubud-\auhdir ,
\]
because $\aubud=fxb(\log b)^3/(mk)$ exceeds $x_1$ and $fmb/k+3x^{35}$ (Lemma~\ref{psu:sizes}(N6)).
\end{proof}

\begin{proof}[Proof of Proposition~\ref{au:hecke}, given Propositions~\ref{po:rallis}, \ref{po:arch},
\ref{po:small}]
By Lemma~\ref{po:Hratio}, $H_j(1)/H_j(\theta_\iota^{\,j})$ is the product of $(\auNm
Q_{\theta_\iota})^{-e}$, a rational power of a positive integer prime to $\ell(\iota)$ by (R1),
of the squares of $k/2$ inverse ratios of main pairs for $j$, and of the squares of $5m$ inverse
ratios $\rho_q(\theta_\iota^{\,j},\pm1)^{-1}$ of ramified critical pairs. By
Proposition~\ref{po:mainheight} each of the former has height at most $6\aubud-\auhdir$, and by
(R2) each of the latter has height at most $x_1$. Heights are subadditive under products, so the
height is at most $2\cdot\tfrac k2(6\aubud-\auhdir)+2\cdot5m\cdot x_1$.
\end{proof}

\subsubsection{What the two propositions rest on}\label{po:ss-rest}

Proposition~\ref{au:onelift} is proved here from: (L4) (Proposition~\ref{plf:L4}); Lemma~\ref{po:ff}
(Theorem~\ref{pfq:ff}, for the first factors of Definition~\ref{pfq:def-ff}); Lemma~\ref{po:classes} (Lemma~\ref{plf:classes}); Propositions~\ref{po:exchange}
(\S\ref{pte:sec}, \S\ref{ptf:sec}), \ref{po:span} (Proposition~\ref{psp:main}, Theorem~\ref{psv:main}) and
\ref{po:periods} (Theorem~\ref{ppe:main}); Lemma~\ref{po:Hpos}; the fact (W7) with the normalisation
\eqref{po:eq-w5}; and the properties of the weights $\mu^\pm$ (Lemma~\ref{plf:match}) and of the characters of
\S\ref{po:ss-set} and \S\ref{po:ss-hecke} (Definition~\ref{pfq:def-chars}, Proposition~\ref{pfq:tame}). Proposition~\ref{au:hecke} is proved from Propositions~\ref{po:arch}
(Theorem~\ref{psu:arch}) and~\ref{po:small} (Theorem~\ref{psu:small}), Proposition~\ref{po:rallis} being
Proposition~\ref{pex:rallis}; the last contains the sharpness of the count and the rank-one case of
Proposition~\ref{po:periods}. Everything else in this subsection is proved in full:
Lemmas~\ref{po:naut}, \ref{po:facts}, \ref{po:det}, \ref{po:eig}, \ref{po:adj}, \ref{po:frame},
\ref{po:prod}, \ref{po:lattice}, \ref{po:slice}, \ref{po:Hratio}, \ref{po:Hpos},
Propositions~\ref{po:proj}, \ref{po:glue} and~\ref{po:mainheight}, and the two deductions.
The operator is bounded directly from \eqref{po:eq-rankone}, and $n_{\mathrm{Aut}}$ is the integer of
Lemma~\ref{po:naut}.


%% ===== proofs/C-congruences.tex =====
% proofs/C-congruences.tex
% Regular families of twists: Proposition au:regular.
% Labels of this file start with pg:.
\providecommand{\auO}{\mathcal O}
\providecommand{\auI}{\mathcal I}
\providecommand{\auQlb}{\overline{\Q}_\ell}
\providecommand{\auPl}{\Phi}
\providecommand{\auPb}{\Phi_b}
\providecommand{\auT}{\mathcal T}
\providecommand{\auG}{\overline G}
\providecommand{\audT}{d_{\auT}}
\providecommand{\auNm}{\operatorname{Nm}}
\providecommand{\auU}{\mathrm U}
\providecommand{\auGL}{\mathrm{GL}}
\providecommand{\auA}{\mathbb A}
\providecommand{\auF}{\mathbb F}
\providecommand{\aurd}{\mathrm{rd}_E}
\providecommand{\aueu}{\varepsilon_{\mathrm u}}
\providecommand{\aubud}{\mathrm{bud}}
\providecommand{\aubar}[1]{\overline{#1}}
\providecommand{\auf}[1]{\texttt{#1}}
% macros of this file (prefix \pg)
\providecommand{\pgZb}{\overline{\Z}_\ell}
\providecommand{\pgv}{v}
\providecommand{\pgNch}{N_{\mathrm{ch}}}
\providecommand{\pgcodd}{c_{\mathrm{odd}}}
\providecommand{\pgloc}{\mathrm{loc}}
\providecommand{\pgG}{G_E}
\providecommand{\pgSel}{\operatorname{Sel}}
\providecommand{\pglen}{\operatorname{len}}
\providecommand{\pgtr}{\operatorname{tr}}
\providecommand{\pgFrob}{\operatorname{Frob}}
\providecommand{\pgHom}{\operatorname{Hom}}
\providecommand{\pgS}{S_{\mathrm{bad}}}
\providecommand{\pgcyc}{\varepsilon_\ell}
\providecommand{\pgpr}{\operatorname{pr}}

\section{Regular families of twists}\label{pg:sec}

This section concerns Proposition~\ref{au:regular}: for all admissible data there is a regular
family of $m$-fold twists. We state exactly what is asserted for one prime $\ell$ and one
isomorphism $\iota\colon\C\to\auQlb$ (Theorem~\ref{pg:main}) and reduce it to three theorems
(Theorems~\ref{pg:I}, \ref{pg:II}, \ref{pg:III}).

\emph{What is proved and what is not.} Proved in this section are: the deduction of
Proposition~\ref{au:regular} from Theorem~\ref{pg:main}, including all sums over $\ell$
(Lemma~\ref{pg:sums}); the deduction of Theorem~\ref{pg:main} from Theorems~\ref{pg:I}--\ref{pg:III}
(\S\ref{pg:ss-reduction}); the part of Theorem~\ref{pg:I} that makes the argument uniform in $\ell$
(Lemmas~\ref{pg:cyclic} and~\ref{pg:burnside}: the labels at $p_0$ and $p_1$ force every lattice in
a main summand to be residually absolutely irreducible, for every $\ell$); and, inside
Theorem~\ref{pg:II}, the choice of the roots of unity (Proposition~\ref{pg:roots}), the reduction of
all twisted Selmer groups to one system of linear equations (Lemma~\ref{pg:minor}), the relation
between tame inertia and Frobenius (Lemma~\ref{pg:tame}), the count of Greenberg--Wiles type
(Lemma~\ref{pg:count}) and the bound for the constant of a Selmer datum (Lemma~\ref{pg:beta}). The
remaining statements are proved in later sections: Theorem~\ref{pg:I} in \S\ref{pga:sec} from
published theorems on Galois representations; Propositions~\ref{pg:mirror}, \ref{pg:step},
\ref{pg:stick} and statement (A) in \S\ref{pse:sec}, with the input (IM) proved in \S\ref{pex:sec};
Theorem~\ref{pg:III} (Ribet's method with auxiliary primes, Propositions~\ref{pg:ribet},
\ref{pg:triangle}, \ref{pg:odd}, \ref{pg:norm}) in \S\S\ref{pri:sec}, \ref{pod:sec}, \ref{pku:sec},
\ref{pax:sec}, \ref{pay:sec}, \ref{prm:sec}. The only use of integral $p$-adic Hodge theory,
Proposition~\ref{pg:loss}, is proved in \S\ref{pll:sec} from Theorem~\ref{pll:exponent}, which is
proved in \S\ref{pki:ss-deduce} (for $\ell\ge4b+2$, $\ell\nmid q_1dM$ see also
Remark~\ref{pll:FL} for another route). \S\ref{pg:ss-summary}
lists the inputs.

\emph{Conventions.} $\ell$ is a prime and $\iota\colon\C\to\auQlb$ an isomorphism; $\pgZb$ is the
ring of integers of $\auQlb$, and $\pgv$ the valuation of $\auQlb$ with $\pgv(\ell)=1$; for
$z\in\C$ we write $\pgv(z)$ for $\pgv(\iota z)$. A \emph{coefficient ring} is the ring of integers
$\auO$ of a finite Galois extension of $\Q_\ell$ inside $\auQlb$; ``for all large $\auO$'' means
for all $\auO$ containing a given one. Lengths of $\auO$-modules are \emph{normalised} so that
$\auO/\ell$ has length $1$; thus $\pglen\auO/a=\pgv(a)$. $\pgG=\operatorname{Gal}(\overline\Q/E)$,
$c$ is a fixed complex conjugation, $g^c=cgc$ and $g^{-c}=(g^c)^{-1}$ for $g\in\pgG$;
$\pgcyc\colon\pgG\to\Z_\ell^\times$ is the cyclotomic character and $\chi=\pgcyc^{-1}$. For a
finite place $w$ of $E$, $D_w\supseteq I_w$ are a decomposition and an inertia group at a fixed
place of $\overline\Q$ above $w$, and, if $w\nmid\ell$, $\gamma_w\in I_w$ is a \emph{tame
generator}: an element whose image generates the pro-$\ell$ quotient of tame inertia; at the
places $w$ above $p_0p_1$ (the chosen places and their conjugates) we require more, namely that the
image of $\gamma_w$ is a topological generator of the whole tame inertia group $I_w/P_w$ ($P_w$ the
wild inertia group), so that its image under the tame character of every level $d$ generates
$\auF_{q^d}^\times$, $q=\auNm w$ (this is what {\rm(G4)} and {\rm(M3)--(M5)} below use). These
elements are fixed once and for all; $\pgFrob_w$ denotes a geometric Frobenius element in $D_w$. The \emph{bad places} are
the places of $E$ dividing $\ell p_0p_1N$; their set is $\pgS=\pgS(\ell)$. The \emph{chosen places}
are the $2f$ places dividing $\pi_M$.

\subsection{The statement for one prime}\label{pg:ss-statement}

\subsubsection*{Pairs, forms, lifts}
A \emph{pair} $O=(W,V,\mu,Q)$ consists of hermitian spaces $W$, $V$ of ranks $n_W<n_V$ with $V$
positive definite and standard, an automorphic character $\mu$ of $\auU(1)$ tamely ramified at the
chosen places and at a finite set $Q$ of places as in the definition of a twist, and the weight
and level of the forms on $\auU(V)$ that the theta lift for $(\auU(W),\auU(V))$ with the character
$\mu$ produces. We write $\mathcal F(O)$ for this space of algebraic modular forms on $\auU(V)$,
with its inner product, and $t(O)=n_V-n_W$. Two families of pairs occur:
\begin{itemize}
\item for $1\le j\le m$ and a twist $\vartheta$, the pair $O_j(\vartheta)=(W_j,V_j,\mu_j^0\vartheta,
Q_\vartheta)$ of the $j$-th lift, with $n_W=r$, $t(O)=t$; the lift is the map $u$ of
Proposition~\ref{au:factor}, restricted to $e_0\mathcal A(W_j)$, and $\mathcal Z_j(\vartheta)\subseteq
\mathcal F(O_j(\vartheta))$ is its image, the \emph{space of lifts};
\item for a main pair $(\nu,z)$ for $j$ (a critical pair with $m<z\le r$, tame part $\omega_w^{\,j}$
at each chosen place $w$, where $\omega_w$ is the character of $\auF_w^\times$ of order $q_2$ fixed in
\S\ref{ss:au-spaces} (so that $\mu_j^0$ has tame part $\omega_w^{\,j}$), and weight $-c$
with $\kappa_j(\tau)+r-t\le c(\tau)\le\kappa_j(\tau)$) and a twist $\vartheta$, the pair
$O_{\nu,z}(\vartheta)=(W,V,\nu\vartheta,Q_\vartheta)$ with $n_W=1$, $n_V=2z$; the space of lifts
is the line $\C\,L_{\nu,z}(\vartheta)$ spanned by the lift of the constant function.
\end{itemize}
The product $\auG_O=\prod_w\auGL_{n_W}(\auF_w)$ over the chosen places acts on $\mathcal F(O)$
(twisted by the tame part of $\mu$, as in Proposition~\ref{au:equiv}); for a family $c=(c_w)$ of
labels of degree $n_W$, $e_c$ is the corresponding central idempotent. For the pairs $O_j$ we use
the idempotent $e_0$; for the pairs $O_{\nu,z}$ the idempotent $e_1$ of the labels $X-1$ (the mean
over $\prod_w\auF_w^\times$).

Let $w$ be a place of $E$ of degree one with $w\ne\aubar w$ such that neither $w$ nor $\aubar w$
divides $\ell\,M\,\auNm Q$ (a \emph{split place off the level}). The spherical Hecke algebra of
$\auGL_{n_V}(E_w)$ acts on $\mathcal F(O)$; let $T_{w,i}$, $0\le i\le n_V$, be the operator of the
double coset of $\operatorname{diag}(\varpi_w,\dots,\varpi_w,1,\dots,1)$ with $i$ entries
$\varpi_w$, and
\[
H_w(X)=\sum_{i=0}^{n_V}(-1)^i\,(\auNm w)^{i(i-1)/2}\,T_{w,i}\,X^{n_V-i}
\]
the \emph{Hecke polynomial}, a polynomial with coefficients in commuting self-adjoint-up-to-scalar
operators. For an eigenform $y$, $H_w(y;X)\in\C[X]$ is the polynomial of its eigenvalues. An
eigenform $y\in\mathcal F(O)$ \emph{has the Hecke polynomials of a lift of the system}
$q=(q_{w,i})$ (a system of eigenvalues of the operators $T_{w,i}$, $i\le n_W$, on forms on
$\auU(W)$) if for every split place $w$ off the level
\begin{equation}\label{pg:eq-liftpoly}
H_w(y;X)=\Bigl(\text{the polynomial of }q\text{ at }w,\text{ rescaled by }\mu\Bigr)\cdot
\prod_{i=0}^{t(O)-1}\bigl(X-\lambda_w(\auNm w)^{i}\bigr),
\end{equation}
where $\lambda_w=\lambda_w(O)$, the \emph{lift root}, is an explicit number depending on $\mu$ and
$w$ only, and the first factor is the Hecke polynomial of $q$ with its roots multiplied by a
number depending on $\mu$ and $w$ only. (The exact normalisations are those of the unramified
theta correspondence; we do not need them. That theta lifts of eigenforms have these Hecke
polynomials is statement (L) of \S\ref{pg:ss-galois}.)

\begin{definition}[False lifts; denominators]\label{pg:def-den}
Let $O$ be a pair, $e$ a central idempotent of $\C[\auG_O]$, and $\mathcal Z\subseteq\mathcal F(O)$
a subspace.
\begin{enumerate}
\item[(a)] For a split place $w$ off the level, the \emph{false lifts at $w$} are the elements of
\[
\mathcal F^{\mathrm{false}}_w=e\,\mathcal F(O)\ \cap\ \ker H_w(\lambda_w)\ \cap\ \mathcal Z^{\perp},
\]
and $y\in\mathcal F(O)$ is \emph{orthogonal to the false lifts} if $y\perp\mathcal F^{\mathrm{
false}}_w$ for every such $w$.
\item[(b)] Let $\beta\ge0$. The projection to $\mathcal Z$ \emph{has denominator at most $\beta$ at
$\iota$} if there is $\Lambda\in\operatorname{End}(\mathcal Z)$ with
\[
\pgv(\det\Lambda)\le\dim\mathcal Z\cdot n_W\cdot\beta
\]
such that for every $y\in\mathcal F(O)$ that is integral at $\iota$ and orthogonal to the false
lifts, the form $\Lambda\,\pgpr_{\mathcal Z}(y)$ is integral at $\iota$, and
$\langle y',\Lambda\,\pgpr_{\mathcal Z}(y)\rangle\in\pgZb$ for every $y'$ that is integral at every
isomorphism $\C\to\auQlb$.
\end{enumerate}
\end{definition}

\subsubsection*{The constants}
Put
\[
\pgNch=\prod_{w\text{ chosen}}\#\auGL_{rt}(\auF_w),\qquad h_1=\#\bigl(\auU(1)(E^+)\backslash
\auU(1)(\auA_{E^+,f})/J_1\bigr),
\]
where $J_1$ is the group of integral elements that are $\equiv1$ at the chosen places. Let
$\pgcodd$ be a positive integer such that, for every $\ell$, all large $\auO$, all $0\le e<t$ and
all $\nu\ge1$, the group of \emph{odd classes} with coefficients in $(\auO/\ell^\nu)(\chi^e)$
(Proposition~\ref{pg:odd}) has length at most $\pgv(\pgcodd)$. Let $h=r^8=\Delta/2$, and for a prime
$\ell$
\begin{align*}
\lambda_\ell&=\pgv\Bigl((2b)!\cdot\Bigl(\prod_{b-30h^3<n\le b+30h^3}n\Bigr)^{h^3}\Bigr),&
\pgloc_\ell&=4f\cdot10^6\,\frac mk\bigl(b\,\pgv(M)+\lambda_\ell\bigr),\\
e'_\ell&=(4f+3)\,r\,\pgv(\pgNch)+2^{-\ell},&
e_\ell&=t\,\pgv(\pgcodd)+\bigl(8ft+2+(4t+1)h^{22}+h\bigr)\pgv(\pgNch),
\end{align*}
\begin{equation}\label{pg:eq-beta}
\beta_\ell=\pgloc_\ell+(1+2t)\,e'_\ell+e_\ell+\pgv(\pgNch).
\end{equation}
The four terms are: the \emph{local loss} above $\ell$; the allowance for the Selmer groups and
for the ramification of the twist; an error term for the class group, the level and the
congruences between powers of the cyclotomic character; and the order of the group $\auG_O$.

\begin{theorem}[A twist with few congruences at one prime]\label{pg:main}
Let $D$ be admissible data, $\ell$ a prime and $\iota\colon\C\to\auQlb$ an isomorphism. There is an
$m$-fold twist $\tau$ such that:
\begin{enumerate}
\item[(a)] $\ell\nmid\auNm Q_\tau$;
\item[(b)] for every ramified critical pair $(\nu,z)$, every $1\le n\le m$ and both signs,
\[
\pgv\Bigl(\frac{|L(z,(\nu\tau^{\pm n})^{-1})|}{|L(z,\nu^{-1})|}\Bigr)\le e'_\ell;
\]
\item[(c)] for $1\le j\le m$, the projection to $\mathcal Z_j(\tau^j)$ has denominator at most
$\beta_\ell$ at $\iota$, for the idempotent $e_0$;
\item[(d)] for $1\le j\le m$, every main pair $(\nu,z)$ for $j$ and all $\Xi,\Xi'\in\mathcal
F(O_{\nu,z}(\tau^j))$ that are integral at every isomorphism $\C\to\auQlb$ and orthogonal to those
false lifts (for $e_1$) that lie in the kernel of the adjoint of the lift from $\auU(1)$
(\emph{test forms}),
\[
\pgv\bigl(\langle\pgpr\,\Xi,\pgpr\,\Xi'\rangle\bigr)\ge-\beta_\ell-\pgv(h_1),
\]
where $\pgpr$ is the orthogonal projection to the line $\C L_{\nu,z}(\tau^j)$.
\end{enumerate}
\end{theorem}

\begin{remark}[The exact form of (R3)]\label{pg:rem-R3}
Clauses (c), (d) are the exact content of ``$\theta_\iota^{\,j}$ has denominators at most
$\beta$ at $\iota$ for the $j$-th lift'' in Definition~\ref{def:au-regular}(R3), with
$\beta=(\beta_\ell+\pgv(h_1))\log\ell$. The \emph{false lifts} of (R3)(i) are those of
Definition~\ref{pg:def-den}(a) (forms in $e_0\mathcal F$, orthogonal to $\mathcal Z$, killed by
$H_w(\lambda_w)$ at \emph{one} split place $w$; one is orthogonal to them for every $w$); and
(R3)(ii) is (d): a bound $e^{\beta}$ for the absolute value at $\iota$ of
$\langle\pgpr\Xi,\pgpr\Xi'\rangle$ for test forms $\Xi,\Xi'$. The condition in (R3) on the sum is
$\sum_{\ell\in T}\beta(\ell)\le\aubud$ with $\beta(\ell)$ already containing the factor
$\log\ell$. This is the form in which Propositions~\ref{au:onelift}(b), \ref{au:hecke}
and~\ref{au:Zj} quote (R3).
\end{remark}

\begin{remark}[How the primes are treated]\label{pg:rem-primes}
The statement and the argument are the same for every $\ell$: for $\ell=2$, for $\ell\mid q_1d$
(ramified in $E$), for $\ell$ inert or split in $K$, and for $\ell\in\{p_0,p_1,N\}$. No hypothesis
``$\ell$ odd, unramified, ordinary'' is made, and none of the published non-vanishing theorems
modulo $\ell$ for anticyclotomic twists is used. The dependence on $\ell$ is entirely in the
constants: $\pgv(\pgNch)$, $\pgv(\pgcodd)$ and $\pgv(h_1)$ vanish unless $\ell$ divides one of three
fixed integers; $2^{-\ell}$ is summable; and $\pgloc_\ell$ vanishes for $\ell>2b$,
$\ell\nmid M$, while for small $\ell$ it is of size $f(m/k)\,b/(\ell-1)$. The three devices that
make this possible are: (i) the twist has $\ell$-power order and conductor prime to $\ell$, so that
twisting is a deformation of the trivial twist and only upper bounds for valuations at one nearby
point are needed (\S\ref{pg:ss-twisting}); (ii) the labels at the two level primes make the main
summands residually irreducible for all $\ell\ne p_i$, and for $\ell=p_0$ one uses $p_1$
(Lemmas~\ref{pg:cyclic}, \ref{pg:burnside}); (iii) above $\ell$ no ordinarity is used; the failure
of torsion crystalline classes to lift is bounded by $\pgloc_\ell$ (Proposition~\ref{pg:loss}), and
this is what the budget pays for.
\end{remark}

\subsection{The sums over the primes}\label{pg:ss-sums}

\begin{lemma}\label{pg:sums}
Assume that $\pgcodd$ can be taken with $\log\pgcodd\le x^{800}$ (Proposition~\ref{pg:odd}). Then
for every finite set $T$ of primes
\[
\sum_{\ell\in T}e'_\ell\log\ell\le x_1,\qquad
\sum_{\ell\in T}\bigl(\beta_\ell+\pgv(h_1)\bigr)\log\ell\le\aubud=\frac{f\,a\,(\log b)^3}{m\,k}.
\]
\end{lemma}
\begin{proof}
For a positive integer $n$, $\sum_{\ell\in T}\pgv(n)\log\ell\le\log n$; and
$\sum_\ell2^{-\ell}\log\ell\le\sum_{n\ge1}n2^{-n}=2$. We bound the logarithms of the integers that
occur, using \eqref{eq:au-scales}, \eqref{eq:au-sizes} and $\log b\le x^{1/100}$.

\emph{$\pgNch$.} $\#\auGL_n(\auF_p)<p^{n^2}$, there are $2f$ chosen places, $p_i\le b^{11}$ and
$rt=5r^2\le45x^{10}$; so $\log\pgNch\le2f(rt)^2\cdot11\log b\le x^{34}$.

\emph{$h_1$.} Since $E/E^+$ is unramified at all finite places, every element of
$\auU(1)(\auA_{E^+,f})$ is of the form $y/\aubar y$ with $y\in\auA_{E,f}^\times$, and the integral
elements are those with $y$ a unit idele (Hilbert's Theorem~90 at each place, and the vanishing of
$H^1$ of the units in an unramified extension). Hence the class group of $\auU(1)$ at the maximal
level is a quotient of the class group of $E$, and $h_1\le h_E\cdot\prod_{w\text{ chosen}}
(\auNm w-1)$. By Minkowski's bound every ideal class contains an integral ideal of norm at most
$|d_E|^{1/2}$, and the number of integral ideals of norm at most $X$ is at most
$X(1+\log X)^{2f-1}$; with $\log|d_E|=2f\log\aurd\le24x^{12}\log x$ this gives $\log h_E\le
x^{14}$, and $\log h_1\le x^{14}+22f\log b\le x^{15}$.

\emph{The local loss.} $\sum_{\ell\in T}\lambda_\ell\log\ell\le\log(2b)!+h^3\cdot60h^3\log(b+30h^3)
\le2b\log(2b)+60h^6\log(2b)\le3b\log b$, and $\sum_\ell b\,\pgv(M)\log\ell\le b\log M\le22\,b\log
b$. So $\sum_{\ell\in T}\pgloc_\ell\log\ell\le10^8f(m/k)\,b\log b$.

\emph{Conclusion.} $\sum_T e'_\ell\log\ell\le(4f+3)r\,x^{34}+2\le x^{52}\le x_1$. Next,
$\sum_Te_\ell\log\ell\le t\,x^{800}+(8ft+2+(4t+1)h^{22}+h)\,x^{34}\le x^{1000}$, since
$h^{22}=r^{176}\le(3x^5)^{176}$. Therefore
\[
\sum_{\ell\in T}\bigl(\beta_\ell+\pgv(h_1)\bigr)\log\ell\le10^8f\frac mk\,b\log b+(1+2t)x^{52}
+x^{1000}+x^{34}+x^{15}\le2\cdot10^8f\frac mk\,b\log b,
\]
because $b\ge x_1^{10^6}\ge x^{10^{12}}$. Finally $2\cdot10^8\,m^2\log b\le8\cdot10^8x\log b\le
x(\log b)^3$, as $m^2\le4x$ and $(\log b)^2\ge8\cdot10^8$; that is,
$2\cdot10^8f(m/k)b\log b\le f\,xb(\log b)^3/(mk)=\aubud$.
\end{proof}

\begin{proof}[Proof of Proposition~\ref{au:regular} from Theorem~\ref{pg:main} and
Proposition~\ref{pg:odd}]
For each $\iota\in\auI$, with prime $\ell$, let $\theta_\iota$ be the twist $\tau$ of
Theorem~\ref{pg:main}. (R1) is (a). For (R2), let $(\nu,z)$ be a ramified critical pair,
$1\le n\le m$ and a sign be given, and let $y_\iota=|L(z,\nu^{-1})|/|L(z,(\nu\theta_\iota^{\pm
n})^{-1})|$. By (b), $|y_\iota|_\iota\le\ell^{e'_\ell}$; so $y$ has height at most $x_1$ with the
multiplier $H=1$, $\alpha=0$ and the exponents $h_\ell=e'_\ell$, by the first inequality of
Lemma~\ref{pg:sums}. For (R3), put $\beta(\ell)=(\beta_\ell+\pgv(h_1))\log\ell$; the second
inequality of Lemma~\ref{pg:sums} is the condition on the sum, (c) gives the operator $\Lambda$
with $|\det\Lambda|_\iota\ge\exp(-\dim\mathcal Z\cdot r\,\beta_\ell\log\ell)\ge\exp(-\dim\mathcal
Z\cdot r\,\beta(\ell))$, and (d) is the bound for the main pairs (Remark~\ref{pg:rem-R3}).
\end{proof}

\subsection{Reduction to three theorems}\label{pg:ss-reduction}

\subsubsection*{Characters}
For a twist $\vartheta$ let $\psi_\vartheta\colon\pgG\to\pgZb^\times$ be the character of finite
order that corresponds by class field theory, through $\iota$, to the Hecke character
$y\mapsto\vartheta^{-1}(y/\aubar y)$ of $E$. It is \emph{anticyclotomic}:
$\psi_\vartheta(g^c)=\psi_\vartheta(g)^{-1}$; it is unramified outside $Q_\vartheta\cup\aubar{
Q_\vartheta}$, and its restriction to $I_w$, $w\in Q_\vartheta$, is given by the tame part of
$\vartheta$. Conversely, every continuous character $\psi\colon\pgG\to\pgZb^\times$ of finite order
with $\psi(\mathrm{Art}(y))=\eta(y/\aubar y)$ for a character $\eta$ of $\auU(1)(E^+)\backslash
\auU(1)(\auA_{E^+,f})$ that is tamely ramified on a set $Q$ with $Q\cap\aubar Q=\emptyset$ of
places prime to $MN$, and unramified elsewhere, is $\psi_\vartheta$ for the twist
$\vartheta=\iota^{-1}\circ\eta^{-1}$ (trivial at the infinite places); if moreover $\psi(\gamma_w)$
has order greater than $m$ for all $w\in Q$, then $\vartheta$ is $m$-fold.

\begin{definition}[Main summands]\label{pg:def-summand}
Let $n\ge1$ and $t'\ge1$. A continuous representation $R\colon\pgG\to\auGL_n(\auQlb)$ is a
\emph{main summand of rank $n$ and string length $t'$} if:
\begin{enumerate}
\item[(M1)] $R$ is semisimple and $\pgtr R(g^c)=\chi(g)^{t'-1}\,\pgtr R(g^{-1})$ for all $g$;
\item[(M2)] $R$ is unramified outside $\ell p_0p_1$; at $v\mid\ell$ it is crystalline if
$v\nmid p_0p_1$ and potentially crystalline otherwise; its Hodge--Tate weights $J_\sigma$, for
the embeddings $\sigma\colon E\to\C$ (transported by $\iota$), satisfy:
$1\le n\le r$, $t'\le t$; for one $j\le m$ every $u\in J_\sigma$ is \emph{low},
$\Delta/4\le|u|\le2m\Delta$, or \emph{high}, $|u\pm\kappa_j(\sigma)|\le t$ with the sign $+$ for
$\sigma$ in the CM type of the data and $-$ otherwise; at most $nm/k$ elements of $J_\sigma$ have
$|u|\le2h^2$; and $J_\sigma$ and $J_{\aubar\sigma}$ together have at most $n$ negative elements;
\item[(M3)] at every $v\mid p_0p_1$, $v\nmid\ell$: there is $N_v$ prime to $v$ with $R(s^{N_v})=1$
for $s\in I_v$, and an element $\gamma\in I_v$, equal to $\gamma_v$ if $v$ is chosen, with
$\pgv(\det(R(\gamma)-1))\le n\,\pgv(\pgNch)$ and $\pgv(\det(R(\gamma^{-c})-1))\le n\,\pgv(\pgNch)$;
\item[(M4)] for all large $\auO$ and every $\pgG$-stable $\auO$-lattice, with matrices $u(g)$: the
$u(g)$ span $M_n(\auO)$ over $\auO$, and there are $g\in\pgG$ and $i\in\{0,1\}$ with
$\pgcyc(g)=p_i^{\,n}$ and $u(g)$ a scalar;
\item[(M5)] there are a place $v\mid p_0p_1$, $v\nmid\ell$, and $s\in I_v$ such that $R(s)$ is a
scalar $y$ with $y^{\ell^aj}-1$ a unit for all $a\ge0$ and $1\le j\le4n+2$.
\end{enumerate}
\end{definition}

The inequality on negative weights in (M2) says that the Bloch--Kato local conditions are
\emph{balanced}: $\sum_{v\mid\ell}\dim H^1_f(E_v,R)-\dim H^0(E_v,R)\le f\,n$, half of
$[E:\Q]\,n$. Condition (M4) says that every lattice in $R$ is residually absolutely irreducible
(Burnside), so that all lattices are homothetic.

\begin{definition}[Selmer groups]\label{pg:def-sel}
Let $u\colon\pgG\to\auGL_n(\auO)$ be continuous, $Q$ a finite set of places and $\gamma\in\pgG$.
For $v\mid\ell$ let $Z_v(u)$ be the $\auO$-module of \emph{crystalline cocycles}: the continuous
cocycles $D_v\to\auO^n$ whose class in $H^1(E_v,\auO^n\otimes\Q)$ lies in the subspace $H^1_f$ of
Bloch and Kato~\cite[(3.7.2)]{BlochKato}. A function $z\colon\pgG\to B^n$, $B$ an $\auO$-module,
\emph{obeys the crystalline relations} if for every $v\mid\ell$ its restriction to $D_v$ satisfies
every $\auO$-linear relation $\sum_ia_i\,z(g_i)=0$ ($a_i\in M_n(\auO)$, $g_i\in D_v$) that holds
for all $z\in Z_v(u)$. The \emph{Selmer group} $\pgSel_Q(u,\gamma;B)$ is the module of locally
constant cocycles $z\colon\pgG\to B^n$ for $u$ that are unramified outside $\ell p_0p_1$,
$Q\cup\aubar Q$ (that is, $z(I_w)=0$ elsewhere), obey the crystalline relations and vanish at
$\gamma$. We use $B=\auO/a$ and $B=K/\auO$, $K=\operatorname{Frac}\auO$.
\end{definition}

If $u(\gamma)$ is a scalar $\zeta$ with $\zeta-1$ a unit, every class has exactly one cocycle
vanishing at $\gamma$ and $H^0=0$, so $\pgSel_Q(u,\gamma;B)$ is a group of cohomology classes. The
relations define the local condition for every coefficient module at once; for $B=\auO/\ell^M$
the cocycles obeying them contain the reductions of the elements of $Z_v(u)$, and the two coincide
when $Z_v(u)\otimes\auO/\ell^M$ is all of the torsion condition, for instance in the range of
Fontaine and Laffaille.

\begin{theorem}[Main summands exist]\label{pg:I}
Let $\ell$, $\iota$ be given.
\begin{enumerate}
\item[(a)] For $1\le j\le m$ and every system $q$ of Hecke eigenvalues that occurs in
$e_0\mathcal A(W_m)$ there is a main summand $R_{j,q}$ of rank $r$ and string length $t$ such
that, for every twist $\vartheta$, every eigenform $y\in\mathcal F(O_j(\vartheta))$ with the
Hecke polynomials of a lift of $q$ and every Galois representation $S_y$ attached to $y$
(\S\ref{pg:ss-galois}, (G1)),
\begin{equation}\label{pg:eq-trace}
\pgtr\bigl(S_y\otimes\eta_O\bigr)(g)=\psi_\vartheta(g)\,\pgtr R_{j,q}(g)+\sum_{e=0}^{t-1}\chi(g)^e
\qquad(g\in\pgG),
\end{equation}
where $\eta_O$ is the $\ell$-adic character of an explicit algebraic Hecke character depending on
the pair only (the \emph{normalising character}).
\item[(b)] For every $j$ and every main pair $(\nu,z)$ for $j$ there is a main summand
$R_{\nu,z}$ of rank $1$ and string length $2z-1$ with the same property for the pairs
$O_{\nu,z}(\vartheta)$ and the eigenforms with the Hecke polynomials of a lift of the trivial
system.
\end{enumerate}
\end{theorem}

\begin{theorem}[The twisting theorem]\label{pg:II}
Let $\ell$, $\iota$ be given and let $R_1,\dots,R_s$ be main summands, of ranks $n_i\le r$. There
is an $m$-fold twist $\tau$ such that for $1\le j\le m$, with $\vartheta=\tau^j$,
$\psi=\psi_\vartheta$ and $Q=Q_\tau$:
\begin{enumerate}
\item[(a)] the places of $Q$ have prime norm, and $Q$, $\aubar Q$ and $\pgS(\ell)$ are pairwise
disjoint; $\psi$ is trivial on $D_v$ for $v\mid\ell$; and $\sum_{w\in Q}\pgv(\psi(\gamma_w)-1)\le
e'_\ell$;
\item[(b)] for each $i$ there are $\gamma\in I_v$ for some $v\nmid\ell$ and $\zeta\in\pgZb$ with
$\zeta-1$ a unit such that, for all large $\auO$, there is a lattice $u$ in $R_i$ with
$\psi(\gamma)u(\gamma)=\zeta$ and $\pglen\pgSel_Q(\psi u,\gamma;K/\auO)\le e'_\ell$;
\item[(c)] $\pgv\bigl(|L(z,\nu^{-1})|/|L(z,(\nu\vartheta^{\pm1})^{-1})|\bigr)\ge-e'_\ell$ for every
ramified critical pair $(\nu,z)$ and both signs.
\end{enumerate}
\end{theorem}

\begin{theorem}[From Selmer groups to denominators]\label{pg:III}
Let $\ell$, $\iota$ be given, let $\vartheta$ be a twist satisfying
Theorem~\ref{pg:II}(a), $\psi=\psi_\vartheta$, and let $O$ be one of the pairs
$O_j(\vartheta)$, $O_{\nu,z}(\vartheta)$, with its idempotent $e$. Let $\mathcal Z\subseteq
\mathcal F(O)$ be spanned by eigenforms $y$ with the following properties: $y$ lies in
$e_c\mathcal F(O)\subseteq e\,\mathcal F(O)$ for a family $c$ of labels in which, above each of
$p_0$, $p_1$, one label is primitive (if $n_W=r$) or all are $X-1$ (if $n_W=1$), whose rescaled labels do not have the root $1$, and with $n_W$ and $t(O)$ odd (for the data, $r$ is a power of $3$, Lemma~\ref{an:lem-sizes}(S7)); $y$ has the Hecke polynomials of a lift of a system
$q$; and there is a main summand $R$ for which \eqref{pg:eq-trace} holds for all twists, and for
which the conclusion of Theorem~\ref{pg:II}(b) holds for $\psi$. Then the projection to
$\mathcal Z$ has denominator at most $\beta_\ell$ at $\iota$ (Definition~\ref{pg:def-den}).
\end{theorem}

\begin{lemma}[The line of a main pair]\label{pg:line}
Let $O=O_{\nu,z}(\vartheta)$, $L=L_{\nu,z}(\vartheta)$, and suppose that the projection to $\C L$
has denominator at most $\beta$ at $\iota$, for $e_1$. Assume statement \emph{(C)} below. Then
$\pgv(\langle\pgpr\Xi,\pgpr\Xi'\rangle)\ge-\beta-\pgv(h_1)$ for all $\Xi,\Xi'$ as in
Theorem~\ref{pg:main}(d).
\end{lemma}
\smallskip\noindent\textbf{(C)} For $s\in\auU(1)(\auA_{E^+,f})$ let $\varrho(s)$ be the translation of forms on $\auU(V)$ by the
central element $s$, multiplied by the value at $s$ of the character that makes $L$ invariant.
Then $\varrho(s)L=L$, $\langle\varrho(s)y,y'\rangle=\langle y,\varrho(s)^{-1}y'\rangle$; if $s$ is
integral above $\ell$ then $\varrho(s)^{\pm1}$ preserves the forms integral at each isomorphism
$\C\to\auQlb$; every class of $\auU(1)(E^+)\backslash\auU(1)(\auA_{E^+,f})/J_1$ has a
representative integral above $\ell$; and if $\Xi'$ is orthogonal to the false lifts in the kernel
of the adjoint of the lift, and $s_1,\dots,s_{h_1}$ represent the classes, then
$\sum_c\varrho(s_c)^{-1}\Xi'$ is orthogonal to the false lifts in $(\C L)^\perp$. The level $J_1$ is
free by Lemma~\ref{au:free}; this is used in the last assertion. This is proved in Proposition~\ref{pga:C}.\par\smallskip
\begin{proof}
We may assume $L\ne0$. Let $s_1,\dots,s_{h_1}$ be representatives integral above $\ell$ and
$\Xi''=\sum_c\varrho(s_c)^{-1}\Xi'$. By (C), $\Xi''$ is integral at every isomorphism and
orthogonal to the false lifts, and $\langle L,\Xi''\rangle=\sum_c\langle\varrho(s_c)L,\Xi'\rangle
=h_1\langle L,\Xi'\rangle$, so $\pgpr\Xi''=h_1\pgpr\Xi'$. On the line, $\Lambda$ is a scalar
$\mu$ with $\pgv(\mu)\le\beta$ (here $\dim\mathcal Z=n_W=1$). By Definition~\ref{pg:def-den}(b)
applied to $y=\Xi''$ and $y'=\Xi$, $\mu\,h_1\langle\Xi,\pgpr\Xi'\rangle\in\pgZb$. As $\pgpr$ is
self-adjoint and idempotent, $\langle\pgpr\Xi,\pgpr\Xi'\rangle=\langle\Xi,\pgpr\Xi'\rangle$, and
the claim follows.
\end{proof}

\begin{proof}[Proof of Theorem~\ref{pg:main} from Theorems~\ref{pg:I}, \ref{pg:II}, \ref{pg:III},
Proposition~\ref{au:mult} and statements \emph{(L)}, \emph{(C)}]
The set $\mathcal Q$ of systems of Hecke eigenvalues occurring in $e_0\mathcal A(W_m)$ is finite,
and so is the set of main pairs (the weights and the tame parts are bounded and the class group
of $\auU(1)$ at a fixed level is finite). Apply Theorem~\ref{pg:II} to the finitely many main
summands $R_{j,q}$ ($1\le j\le m$, $q\in\mathcal Q$) and $R_{\nu,z}$ of Theorem~\ref{pg:I}. Let
$\tau$ be the resulting twist. (a) holds because $Q_\tau$ is disjoint from the places above
$\ell$. (b) is Theorem~\ref{pg:II}(c) for the powers $\tau^n$.

(c) Fix $j$ and put $\vartheta=\tau^j$. We check that $\mathcal Z=\mathcal Z_j(\vartheta)$
satisfies the hypothesis of Theorem~\ref{pg:III}. The space $e_0\mathcal A(W_j)$ is spanned by
Hecke eigenforms $y_0$ lying in $e_c\mathcal A(W_j)$ for good families $c$ of labels
(definition of $e_0$ in \S\ref{ss:au-spaces}; the Hecke operators at the split places off the
level are normal and commute with $\auG$); put $y=u(y_0)$. A good family has a primitive label above each of $p_0$ and $p_1$ (its root generates
$\auF_{p_i^r}^\times$), and its rescaled labels do not have the root $1$: the rescaling multiplies
the roots of $c_w$ by the element $z_w\in\auF_w^\times$ that names the inverse of the tame part of
$\mu_j^0\vartheta$ at $w$, which is $\omega_w^{-j}$ and has order dividing $q_2$, the order of
$\omega_w$; a rescaled root equal to $1$ would be a root $z_w^{-1}\in\auF_w$ of $c_w$ that is a
$q_2$-th root of unity, which a good family excludes. By (L), $y$ is in $e_c\mathcal F(O)$ and has
the Hecke polynomials of a lift of the system $q$ of $y_0$; by Proposition~\ref{au:mult} (equality
of multiplicities) the system $q$ occurs in $e_0\mathcal A(W_m)$, so
$q\in\mathcal Q$. The main summand $R_{j,q}$ satisfies \eqref{pg:eq-trace} by
Theorem~\ref{pg:I}(a) and the conclusion of Theorem~\ref{pg:II}(b) by the choice of $\tau$. Hence
Theorem~\ref{pg:III} applies and gives (c).

(d) In the same way Theorem~\ref{pg:III} applies to the line $\C L_{\nu,z}(\vartheta)$ of a main
pair for $j$, with the labels $X-1$ and the summand $R_{\nu,z}$: the lift of the constant
function is fixed by $\prod_w\auF_w^\times$ up to the tame part of the character, has the Hecke
polynomials of a lift of the trivial system by (L), and the rescaled label is $X-\xi$ with $\xi$
the element naming $\omega_w^{-j}$ at $w$, which is not $1$ because $\omega_w^{\,j}\ne1$ for $1\le
j\le m$ ($\omega_w$ has order $q_2>m$). Lemma~\ref{pg:line} gives (d).
\end{proof}

\subsection{Galois representations: what is used, and Theorem~\ref{pg:I}}\label{pg:ss-galois}

Let $V$ be a standard positive definite hermitian space of rank $n$ over $E$ and $y$ an algebraic
modular form on $\auU(V)$ of weight $\lambda$ that is an eigenform of the operators $T_{w,i}$ at
almost all split places. The following properties of Galois representations are used in
Theorems~\ref{pg:I} and~\ref{pg:III}, and no others.
\begin{enumerate}
\item[(G1)] \emph{Existence.} There is a continuous semisimple $S_y\colon\pgG\to\auGL_n(\auQlb)$,
unramified at almost all places, such that for almost all split places $w$ at which $y$ is
spherical, $\det(X-S_y(\pgFrob_w))=\iota H_w(y;X)$.
\item[(G2)] \emph{Good places.} If $y$ is fixed by a hyperspecial subgroup at $v$ (for instance by
$\auGL_n(\auO_{E,w})$ at a split place $w$) then $S_y$ is unramified at the places of $E$ above $v$
not dividing $\ell$, with the Frobenius relation of (G1) at split $w$, and crystalline at those
dividing $\ell$.
\item[(G3)] \emph{Weights.} At every place above $\ell$, $S_y$ is de Rham with the Hodge--Tate
weights determined by $\lambda$: at $\sigma\colon E\to\C$ in the CM type the numbers
$i-1-\lambda_{\sigma,i}$ and at $\aubar\sigma$ the numbers $\lambda_{\sigma,i}+n-i$, $1\le i\le n$.
\item[(G4)] \emph{Tame level.} Let $w$ be a split place of residue characteristic $p$ at which $y$
is fixed by the principal congruence subgroup and lies in the isotypic part, for
$\auGL_n(\auF_w)$, of the character named by a label $c$. Then $S_y|_{I_w}$ is tame, of finite
order if $w\nmid\ell$ and potentially crystalline with the corresponding inertial type if
$w\mid\ell$, and its eigenvalues on a generator of tame inertia are the images of the roots of
$c$ under a fixed injective character $\overline{\auF}_p^\times\to\auQlb^\times$ (with the sign
convention $\epsilon=1$ of Proposition~\ref{pga:G45}).
\item[(G5)] \emph{Parahoric level with a character.} If at a split place $w\in Q$, $w\nmid\ell$,
the form $y$ is an eigenvector of a parahoric subgroup of type $(n_W,t)$ with a tame character
$\psi_w$ on the first block and the trivial one on the second, then $(S_y(s)-\psi_w(s))(S_y(s)-1)
=0$ for $s\in I_w$ (after the normalisation; in particular there is no monodromy; with the sign
convention $\epsilon=1$ of Proposition~\ref{pga:G45}).
\end{enumerate}

\emph{Published sources.} By Labesse's base change~\cite{Labesse} (which applies because $E/E^+$
is unramified at all finite places, $[E^+:\Q]\ge2$ and the group is definite), $y$ generates
automorphic representations $\pi$ of $\auU(V)$ with a base change $\Pi$ to $\auGL_n(\auA_E)$, an
isobaric sum of conjugate self-dual discrete representations, compatible with the local base
change at all split places and at all places where $\pi$ is unramified. The Galois
representations attached to the cuspidal constituents of $\Pi$ are those of Harris--Taylor, Shin,
Chenevier--Harris, with local--global compatibility at all finite places, including those above
$\ell$, by Caraiani; all this is collected, for every prime $\ell$ and every regular algebraic
conjugate self-dual cuspidal representation, in~\cite[Theorem~2.1.1]{BLGGT} (above $\ell$ that
theorem covers only components with an Iwahori-fixed vector; above $\ell$ in general we use
\cite[Theorem~A]{BLGGTlp2}, which gives the semisimplified identity and a bound on the monodromy,
and which already yields potential crystallinity with the expected inertial type when the
local parameter has no monodromy, see Proposition~\ref{pga:G45}; see also
\cite{BLGGTlp1,CaraianiANT}). Thus (G1)--(G3)
follow for every $\ell$, including $\ell=2$ and $\ell$ ramified in $E$; (G4) and (G5) follow from
the compatibility at $w$ together with the description of the local Langlands correspondence for
representations of depth zero (a representation containing the inflation of the cuspidal
representation of $\auGL_n(\auF_w)$ of a primitive label corresponds to the representation induced
from the character of the inertia group of the unramified extension of degree $n$ given by a root
of the label) and for representations with Iwahori or parahoric fixed vectors with a tame
character (principal series and their subquotients). \emph{We have not found one published
statement that contains {\rm(G1)--(G5)} for algebraic modular forms on definite unitary groups in
this generality}; each item is a consequence of the cited theorems, but the deduction (passage
from $\Pi$ to its cuspidal constituents, including possible Speh constituents; places of $E^+$
inert in $E$; the normalisations in (G3)) is written in \S\ref{pga:sec}, Propositions~\ref{pga:G123} and~\ref{pga:G45}, with
the twist of Proposition~\ref{pax:sign} for Speh constituents.



\smallskip
The following statement about theta lifts is used; it is proved in Proposition~\ref{pga:L}.
\smallskip\noindent\textbf{(L)} If $y_0$ is a Hecke
eigenform on $\auU(W)$ with system $q$ and its theta lift $y$ to $\auU(V)$ for the character $\mu$
is not zero, then $y$ is an eigenform and has the Hecke polynomials \eqref{pg:eq-liftpoly} of a
lift of $q$ at every split place off the level; and the lift is equivariant for the groups
$\auG_O$ as in Proposition~\ref{au:equiv}, so that $y\in e_c\mathcal F(O)$ if $y_0$ is of type
$e_c$. This is proved in Proposition~\ref{pga:L}, for the Weil representation of \S\ref{po:ss-set}.\par\smallskip

\subsubsection*{The image of inertia at the level primes}

\begin{lemma}\label{pg:cyclic}
Let $p\equiv1\pmod3$ be a prime, $r$ a power of $3$, $\xi\in\overline{\auF}_p^\times$ of order
$p^r-1$, $z\in\auF_p^\times$ and $\xi'=\xi z$. For every prime $\ell\ne p$ and every $u\ge0$ the
elements $(\xi'^{\,p^a})^{\ell^u}$, $0\le a<r$, are pairwise distinct.
\end{lemma}
\begin{proof}
If two of them coincide, for $a<b$, then $\xi^{\ell^up^a(p^{b-a}-1)}=1$, because
$z^{p^a}=z^{p^b}$. So the order $o$ of $\xi^{\ell^u}$, which is $(p^r-1)/\gcd(\ell^u,p^r-1)$,
divides $p^d-1$ with $0<d=b-a<r$, hence divides $p^{\gcd(d,r)}-1$, hence $x-1$ with $x=p^{r/3}$.
(Here we used that $o$ is prime to $p$ and divides $p^r-1$, and that a proper divisor of $r$ divides
$r/3$.) From $o\mid x-1$ and $p^r-1=x^3-1$ we get $x^3-1\mid(x-1)\ell^u$, that is,
$x^2+x+1\mid\ell^u$. But $x\equiv1\pmod3$, so $3\mid x^2+x+1$ and $x^2+x+1\equiv3\pmod9$; a divisor
of $\ell^u$ that is divisible by $3$ and not by $9$ is $3$, and $x^2+x+1>3$. Contradiction.
\end{proof}

\begin{lemma}\label{pg:burnside}
Let $R\colon\pgG\to\auGL_n(\auQlb)$ be continuous, $w\nmid\ell$ a place of $E$ of degree one and
norm $p$, and $s\in I_w$ such that $R|_{I_w}$ is tame of finite order and the eigenvalues of $R(s)$
are $\zeta^{p^a}$, $0\le a<n$, for a root of unity $\zeta$ with the $\zeta^{p^a}$ pairwise
incongruent modulo the maximal ideal of $\pgZb$. Then for every coefficient ring $\auO$ over which
$R$ and $\zeta$ are defined and every stable lattice, with matrices $u(g)$: the $u(g)$ span
$M_n(\auO)$, and $u(\phi^n)$ is a scalar for every arithmetic Frobenius element $\phi\in D_w$;
moreover $\pgcyc(\phi^n)=p^n$.
\end{lemma}
\begin{proof}
Let $k$ be the residue field of $\auO$ and $\bar u$ the reduction. Since the eigenvalues of
$u(s)$ are distinct modulo the maximal ideal, the lattice is the direct sum of the eigenlines
$\Lambda_a$ of $u(s)$, and $k^n$ is the direct sum of the eigenlines $\bar\Lambda_a$ of
$\bar u(s)$ for the distinct eigenvalues $\bar\zeta^{p^a}$. In the tame quotient $\phi
s\phi^{-1}=s^{p}$, and $R$ is trivial on wild inertia; so $u(\phi)$ maps the eigenline of $u(s)$
for an eigenvalue $\alpha$ to the eigenline of $u(s^{p})$ for $\alpha$, which is the eigenline of
$u(s)$ for the eigenvalue $\alpha'$ with $\alpha'^{\,p}=\alpha$. (the $p$-th power map is injective on the eigenvalues, which have order prime to $p$). So
$u(\phi)$ permutes the eigenlines, and the set of eigenvalues is stable under $\alpha\mapsto
\alpha'$; for $\alpha=\zeta$ this gives $\zeta=\zeta^{p^{b+1}}$ with $b<n$, hence $b=n-1$ since
the $\zeta^{p^a}$ are distinct. Thus $u(\phi)$ permutes the $\Lambda_a$ cyclically. A $\bar u$-stable
subspace of $k'^n$, $k'\supseteq k$, is a sum of lines $\bar\Lambda_a\otimes k'$ and is stable
under $\bar u(\phi)$, so it is $0$ or everything: $\bar u$ is absolutely irreducible. By Burnside's
theorem the $\bar u(g)$ span $M_n(k)$, and by Nakayama's lemma the $u(g)$ span $M_n(\auO)$.
Finally $u(\phi^n)$ preserves each $\Lambda_a$ and commutes with $u(\phi)$, which permutes them
transitively; so its eigenvalues on the $\Lambda_a$ are equal, and it is a scalar.
\end{proof}

\begin{proof}[Proof of Theorem~\ref{pg:I}, given {\rm(G1)--(G4)} (Propositions~\ref{pga:G123}, \ref{pga:G45}) and {\rm(L)} (Proposition~\ref{pga:L})]
(a) Let $q\in\mathcal Q$, let $y_0\in e_c\mathcal A(W_m)$ be an eigenform with system $q$, for a
good family $c$ (recall that $\auU(W_m)$ is definite), and let $S_q$ be attached to $y_0$ by (G1).
Put $R_{j,q}=S_q\otimes\eta_j$, where $\eta_j$ is the $\ell$-adic character of the algebraic Hecke
character $y\mapsto(\mu_j^0)^{-1}(y/\aubar y)\,|y|^{(r-t)/2}$ (note that $r\equiv t\pmod2$).

\emph{The trace identity.} Let $\vartheta$, $y$, $S_y$ be as in the statement. By (G1) and
\eqref{pg:eq-liftpoly}, for almost all split places $w$ the characteristic polynomial of
$S_y(\pgFrob_w)$ is the product of that of $S_q(\pgFrob_w)$, with roots rescaled by a number
depending on $\mu_j^0\vartheta$ and $w$, and of $\prod_{i<t}(X-\lambda_w(\auNm w)^i)$. The
rescaling numbers and $\lambda_w$ are values at $\pgFrob_w$ of $\ell$-adic characters of algebraic
Hecke characters built from $\mu_j^0\vartheta$; after multiplication by $\eta_O(\pgFrob_w)$ the
second factor becomes $\prod_{i<t}(X-(\auNm w)^i)=\prod_{i<t}(X-\chi(\pgFrob_w)^i)$ and the first
the characteristic polynomial of $(\psi_\vartheta\otimes R_{j,q})(\pgFrob_w)$; this computation is
carried out in Proposition~\ref{pga:trace} (with $\eta_O$ the character of $(\chi_W(y/\aubar y)\upsilon(y)^{-r}|y|^{-r/2})^{-1}$). The split
places of degree one have density one, so by the Chebotarev density theorem and continuity
\eqref{pg:eq-trace} holds for all $g$.

\emph{{\rm(M1)}.} $S_q$ is semisimple, and $S_q^c\cong S_q^\vee\otimes\chi^{r-1}$
by~\cite[Theorem~2.1.1]{BLGGT} applied to the constituents of the base change (or directly: the
Hecke operators at $w$ and $\aubar w$ are dual to each other, which gives the identity of traces
at Frobenius elements, and then Chebotarev). Twisting by $\eta_j$, which satisfies
$\eta_j^c=\eta_j^{-1}\chi^{t-r}$ (its $\mu$-part is anticyclotomic and $|y|^{(r-t)/2}$
corresponds to $\chi^{(t-r)/2}$), gives $\pgtr R(g^c)=\chi(g)^{t-1}\pgtr R(g^{-1})$.

\emph{{\rm(M2)}.} The forms in $\mathcal A(W_m)$ have level $J_M$: hyperspecial outside $M$ and
principal at the chosen places. So $R_{j,q}$ is unramified outside $\ell p_0p_1$ and crystalline
above $\ell$ if $\ell\nmid M$ by (G2), and potentially crystalline above $\ell=p_i$ by (G4). The
weights: by (G3) and the definition of the parameter $\varsigma_\tau$ of the target forms, the
Hodge--Tate weights of $R_{j,q}$ at $\sigma$ are, up to a shift bounded by $t$ in absolute value
and with $\epsilon=\pm1$ according as $\sigma$ is in the CM type or not,
\[
-\epsilon\bigl(\kappa_j(\sigma)+O(k)\bigr)\quad(k-1\text{ weights}),\qquad
-\epsilon\,(j-i+\tfrac12)\Delta\quad(1\le i\le m).
\]
The first are high. The second are low: $\Delta/4\le|(j-i+\frac12)\Delta|-t$ and $|(j-i+\frac12)\Delta|+t\le2m\Delta$ (as $t=5r\le\Delta/4$). Exactly
$m$ weights have absolute value at most $2h^2$, and $m\le rm/k$. For $\sigma$ in the CM type the
negative weights are the $k-1$ high ones and the $j$ low ones with $i\le j$; for $\aubar\sigma$
they are the $m-j$ low ones with $i>j$; together $k-1+m=r$. The exact weights are computed in Proposition~\ref{pga:weights}: they are
$\frac{t-1}2-\epsilon_\sigma(\varsigma_{\sigma,i}+\kappa_j(\sigma))$.

\emph{{\rm(M3)}, {\rm(M4)}.} Let $i\in\{0,1\}$ with $p_i\ne\ell$ and let $w\mid p_i$ be a chosen
place at which the label $c_w$ is primitive. By (G4), $R_{j,q}|_{I_w}$ is tame of finite order
dividing $p_i^r-1$, and the eigenvalues of $R_{j,q}(\gamma_w)$ are the images of the roots
$\xi'^{\,p^a}$ of the label rescaled by the element $z\in\auF_{p_i}^\times$ naming the restriction
of $\eta_j$ to $I_w$. As $p_i\equiv1\pmod{q_2^2}$ and $3\mid q_0\mid q_2$, Lemma~\ref{pg:cyclic}
applies: for every $u$ the $\ell^u$-th powers of these $r$ roots of unity are distinct, that is, no
quotient of two of them has $\ell$-power order, so they are pairwise incongruent modulo the
maximal ideal. Lemma~\ref{pg:burnside} gives (M4) with $g=\phi^r$,
$\pgcyc(g)=p_i^{\,r}$. For (M3): the places above $p_0p_1$ are the chosen places and their conjugates. At a chosen
$v\nmid\ell$ the eigenvalues of $R(\gamma_v)$ are, by (G4), the images of the roots of the
rescaled label: roots of unity $\omega\ne1$ of order dividing $p_i^r-1$ (not $1$ because the
rescaled labels do not have the root $1$; see the proof of Theorem~\ref{pg:main}). By (M1),
$R(g^c)$ is conjugate to $\chi(g)^{t-1}\,{}^{\mathrm t}R(g)^{-1}$, so $R(\gamma_v^{-c})$ has the
same eigenvalues as $R(\gamma_v)$; this treats the conjugate places. If
$\pgv(\omega-1)>0$ then $\omega$ has order a power $\ell^a>1$ of $\ell$, so
$\pgv(\omega-1)\le1/(\ell-1)\le1$ and $\ell\mid p_i^r-1\mid\pgNch$. Hence
$\pgv(\det(R(\gamma)-1))\le r\le r\,\pgv(\pgNch)$ if some eigenvalue is not a unit, and it is $0$
otherwise.

\emph{{\rm(M5)}.} This is Proposition~\ref{pga:M5}: with $e=(p_i^r-1)/(p_i-1)$, $s=\gamma_w^{\,e}$
acts by a scalar $y$ (the image of the norm $y_0$ of $\xi'$ to $\auF_{p_i}$) whose order is
divisible by $q_2^2$, because $p_i\equiv1\pmod{q_2^2}$; at a chosen place with label $X-1$ (the
case of the twist $\eta_j$, $1\le j\le m$) the inertia generator acts by a scalar of order
$q_2/\gcd(q_2,j)$. In both cases $y_0^x=1$ forces $q_2\mid jx$, and Lemma~\ref{pga:q2} (which
uses $q_2=(q_0d)^{s+6}$ and $2^s=m^{22}$) excludes $q_2\mid j\,u\,\ell^{a'}$ for $1\le u\le4r+2$,
$1\le j\le m$, $a'\ge0$. Hence $y^{\ell^a u}-1$ is a unit for all $a\ge0$ and $1\le u\le4r+2$,
which is (M5); the details are in Proposition~\ref{pga:M5}.

(b) For a main pair, $W$ has rank one, the eigensystem is trivial and $R_{\nu,z}$ is the
$\ell$-adic character of the algebraic Hecke character $y\mapsto\nu^{-1}(y/\aubar y)|y|^{1-z}$. It
is crystalline above $\ell\nmid M$ and potentially crystalline above $\ell\mid M$, with the single
weight $z-1-\epsilon\,c(\sigma)$ at $\sigma$, which is high because $|c-\kappa_j|\le t$; exactly
one of the weights at $\sigma$, $\aubar\sigma$ is negative. (M4) is empty in rank one except for
the scalar, and (M3), (M5) hold as above, the eigenvalue of inertia at a chosen place being the
image of the name of $\omega_w^{-j}$. The trace identity is proved as in (a).
\end{proof}

\noindent Theorem~\ref{pg:I} is proved, modulo the published theorems quoted with their hypotheses in
Theorem~\ref{pga:lit} and Lemmas~\ref{pga:depth0}, \ref{pga:parahoric}; the twist for Speh
constituents is Proposition~\ref{pax:sign} and the convention check is Lemma~\ref{pax:conv}.\par\smallskip

\subsection{The twisting theorem}\label{pg:ss-twisting}

\subsubsection*{Roots of unity}
For a polynomial $g=\sum_iy_iY^i\in\pgZb[Y]$ and $\rho\ge0$ let $\pgv_\rho(g)=\min_i(\pgv(y_i)+
\rho i)$ be the Gauss valuation. It is multiplicative, nondecreasing in $\rho$, and
$\pgv(g(y))\ge\pgv_{\rho}(g)$ if $\pgv(y)\ge\rho$. Put $\rho_h=1/\varphi(\ell^h)$, the valuation
of $\zeta-1$ for $\zeta$ of order $\ell^h$.

\begin{lemma}\label{pg:roots1}
Let $0<\delta<1$, $\rho>0$, $B\ge0$, and let $h_0$ be such that $\delta\sum_h(h+1)^{-1}>B$, the
sum over the $h$ with $1\le h\le h_0$ and $\rho_h\le\rho$. Let $g\in\pgZb[Y]$ with
$\pgv_\rho(g)\le B$. Then there are such an $h$ and a root of unity $\zeta$ of order $\ell^h$ with
$\pgv(g(\zeta-1))\le\pgv_{\rho_h}(g)+\delta$.
\end{lemma}
\begin{proof}
We may assume $g\ne0$; write $g=c\prod_i(Y-\alpha_i)$ over $\auQlb$. By multiplicativity,
$\pgv_{\rho'}(g)=\pgv(c)+\sum_i\min(\rho',\pgv(\alpha_i))$ for every real $\rho'$. For $\pgv(y)=\rho'$ one has $\pgv(y-\alpha_i)=\min(\rho',\pgv(
\alpha_i))$ unless $\pgv(\alpha_i)=\rho'$; hence
\begin{equation}\label{pg:eq-newton}
\pgv(g(y))-\pgv_{\rho'}(g)=\sum_{\alpha\in R(\rho')}\bigl(\pgv(y-\alpha)-\rho'\bigr),
\end{equation}
where $R(\rho')$ is the multiset of roots of valuation exactly $\rho'$. Since $g$ is integral, $\pgv_0(g)\ge0$, so $\pgv_\rho(g)\ge\pgv_\rho(g)-\pgv_0(g)\ge
\sum_{i:\pgv(\alpha_i)>0}\min(\rho,\pgv(\alpha_i))\ge\sum_h\#R(\rho_h)\rho_h$, the last sum over
the $h$ of the statement (the $\rho_h$ are distinct). Hence this sum is at most $B$; so for one
of them $\#R(\rho_h)\,\rho_h\le\delta/(h+1)$. Fix it. For $\alpha\in R(\rho_h)$ let $\zeta_0$ be a
root of unity of order $\ell^h$ nearest to $1+\alpha$. For the other $\zeta$ of order $\ell^h$,
$\pgv(\zeta-\zeta_0)\ge\min(\pgv(\zeta-1-\alpha),\pgv(\zeta_0-1-\alpha))=\pgv(\zeta-1-\alpha)$, and $\pgv(\zeta-\zeta_0)=\rho_{h'}$ if
$\zeta/\zeta_0$ has order $\ell^{h'}$; there are at most $\varphi(\ell^{h'})$ such $\zeta$. So
\[
\sum_{\zeta}\min\bigl(\pgv(\zeta-1-\alpha),1\bigr)\le1+\sum_{h'=1}^{h}\varphi(\ell^{h'})\rho_{h'}
=1+h .
\]
Summing over $\alpha$ and dividing by the number $\varphi(\ell^h)=1/\rho_h$ of roots of unity of
order $\ell^h$, some $\zeta$ has $\sum_{\alpha\in R(\rho_h)}\min(\pgv(\zeta-1-\alpha),1)\le
\#R(\rho_h)\rho_h(h+1)\le\delta<1$. Then no term is $1$, and by \eqref{pg:eq-newton} with
$y=\zeta-1$, $\pgv(g(\zeta-1))-\pgv_{\rho_h}(g)\le\delta$.
\end{proof}

\begin{proposition}[Roots of unity]\label{pg:roots}
Let $s_0,m\ge1$ be integers, $0<\rho\le1$, $\beta\ge0$ and $\varepsilon>0$. There is $h_0$ with the
following property. Let $P_1,\dots,P_K\in\pgZb[X_1,\dots,X_s]$, $s\le s_0$, and
$x_1,\dots,x_K\in\pgZb^{\,s}$ with coordinates of valuation at least $\rho$ and
$\sum_i\pgv(P_i(x_i))\le\beta$. Then there are roots of unity $\zeta_1,\dots,\zeta_s$ of orders
dividing $\ell^{h_0}$ and greater than $m$ such that $\pgv(\zeta_u-1)\le\rho/m$ and, for
$1\le j\le m$ and all $i$,
\[
\pgv\bigl(P_i(\zeta_1^{\,j}-1,\dots,\zeta_s^{\,j}-1)\bigr)\le\pgv(P_i(x_i))+\varepsilon .
\]
\end{proposition}
\begin{proof}
\emph{Reduction to $m=1$.} For $1\le j\le m$ let $P_{i,j}(X)=P_i((1+X_1)^j-1,\dots,(1+X_s)^j-1)$,
and let $x_{i,j}$ have coordinates $y$ with $(1+y)^j=1+x$, $x$ the coordinate of $x_i$, where
$y$ is chosen of maximal valuation among the $j$ roots of $(1+Y)^j-(1+x)$. This polynomial is
monic of degree $j$ with constant term $-x$, so the product of its roots is $\pm x$ and
$\pgv(y)\ge\pgv(x)/j\ge\rho/m$. Then $P_{i,j}(x_{i,j})=P_i(x_i)$. Apply the case $m=1$ to the $Km$
polynomials $P_{i,j}$, the radius $\rho/m$ and the bound $m\beta$. Replacing $\rho$ by $\min(\rho,1/m)$ at the start, we get $1/\varphi(\ell^h)=\pgv(\zeta_u-1)\le
1/m^2$ for the order $\ell^h$ of $\zeta_u$, so this order exceeds $m$.

\emph{The case $m=1$.} Let $\delta=\min(\varepsilon,1)/(2s_0)$ and let $h_0$ be as in
Lemma~\ref{pg:roots1} for $\delta$, $\rho$ and $B=\beta+1$. We replace the coordinates one at a
time. Suppose the first $u-1$ coordinates of all $x_i$ have been replaced by $\zeta_1-1,\dots,
\zeta_{u-1}-1$, giving points $x_i'$ with $\sum_i\pgv(P_i(x'_i))\le\beta+(u-1)\delta$ and
$\pgv(P_i(x_i'))\le\pgv(P_i(x_i))+(u-1)\delta$. Let $g_i(Y)$ be $P_i$ with the $u$-th variable
equal to $Y$ and the others equal to the coordinates of $x_i'$, and $g=\prod_ig_i$. Then
$\pgv_\rho(g)\le\sum_i\pgv(g_i(x'_{i,u}))\le\beta+1$, so Lemma~\ref{pg:roots1} gives $h$ and
$\zeta_u$ of order $\ell^h$, $\rho_h\le\rho$, with $\sum_i\pgv(g_i(\zeta_u-1))\le\sum_i\pgv_{
\rho_h}(g_i)+\delta$. Since $\pgv(g_i(\zeta_u-1))\ge\pgv_{\rho_h}(g_i)$ for each $i$, each
difference is at most $\delta$; and $\pgv_{\rho_h}(g_i)\le\pgv_{\rho}(g_i)\le\pgv(g_i(x'_{i,u}))$.
So the replacement of the $u$-th coordinate raises each valuation by at most $\delta$. After $s$
steps the total increase is at most $s\delta\le\varepsilon$.
\end{proof}

\subsubsection*{Selmer data}
Fix $S_0=\pgS(\ell)$, a chosen place $w_0\nmid\ell$ and $g_0=\gamma_{w_0}$.

\begin{definition}\label{pg:def-datum}
A \emph{Selmer datum} over $\auO$ is a continuous $u\colon\pgG\to\auGL_n(\auO)$, unramified outside
a set $S\subseteq S_0$ containing the places above $\ell$, with: for $v\mid\ell$ a finitely
generated module $Z_v$ of cocycles $D_v\to\auO^n$ containing the coboundaries, with
$\sum_{v\mid\ell}(\operatorname{rk}Z_v-n)\le f\,n$; the $u(g)$ span $M_n(\auO)$; and an element
$\gamma_*\in I_v$, $v\in S$, $v\nmid\ell$, such that $u(\gamma_*)=u(\gamma_*^{-c})=\zeta_*$ is a
scalar with $\zeta_*^2-1$ a unit. Its \emph{constant} is $\beta_T=\mathrm{gl}+\mathrm{pl}+
\mathrm{fp}$, where
\begin{align*}
\mathrm{gl}&=\sup_{a\ne0}\pglen\pgHom_{\pgG}\bigl(T/a,(T/a)^{*c}\bigr),\qquad
\mathrm{pl}=\sum_{v\in S,\,v\nmid\ell}\ \sup_{a\ne0}\pglen H^0\bigl(E_v,(T/a)^*\bigr),\\
\mathrm{fp}&=\pgv\bigl(\det(u(g_0)-1)\bigr)+\pgv\bigl(\det(u(g_0^{-c})-1)\bigr);
\end{align*}
here $T=\auO^n$, $(T/a)^*=\pgHom(T/a,\Q/\Z)(1)$, and $^{c}$ means that $g$ acts as $g^c$. The
local conditions are the relations that hold on the $Z_v$, as in Definition~\ref{pg:def-sel}.
\end{definition}

\begin{lemma}\label{pg:beta}
Let $R$ be a main summand of rank $n$. For all large $\auO$ there is a lattice $u$ in $R$ that,
with $S$ the set of places above $\ell p_0p_1$ and $Z_v=Z_v(u)$ the crystalline cocycles, is a
Selmer datum with $\beta_T\le(4f+3)\,n\,\pgv(\pgNch)$.
\end{lemma}
\begin{proof}
\emph{The rank condition.} For a de Rham representation $V$ of $\operatorname{Gal}(\overline{E_v}
/E_v)$, $\dim H^1_f=\dim D_{\mathrm{dR}}(V)/\mathrm{Fil}^0+\dim H^0$ (Bloch--Kato~\cite[Corollary~3.8.4]{BlochKato}), so $\operatorname{rk}Z_v-n=\dim H^1_f-\dim H^0=\dim_{\Q_\ell}D_{\mathrm{dR}}/
\mathrm{Fil}^0$ counted over $K$, which is the number of negative Hodge--Tate weights at the
embeddings of $E_v$ (with the convention of (G3)). Summing over $v\mid\ell$ gives the number of
negative weights at all $2f$ embeddings of $E$, at most $f\,n$ by (M2).

\emph{The scalar element.} Take $\gamma_*=s$ from (M5), $\zeta_*=y$; $y^2-1$ is a unit. By (M1)
and semisimplicity, $g\mapsto R(g^c)$ is isomorphic to $g\mapsto\chi(g)^{t'-1}\,{}^{\mathrm
t}R(g)^{-1}$; so $R(g^{-c})$ is conjugate, by a matrix independent of $g$, to
$\chi(g)^{1-t'}\,{}^{\mathrm t}R(g)$. As $\chi(\gamma_*)=1$, $R(\gamma_*^{-c})$ is the scalar
$y$.

\emph{The global term.} By (M4) there is $g$ with $\pgcyc(g)=p_i^{\,n}$ and $u(g)=\alpha$ a
scalar; by the same remark, $u(g^{-c})$ is the scalar $\alpha'=\alpha\,\chi(g)^{1-t'}$. On $(T/a)^{*c}$ the
element $g$ acts by $\pgcyc(g^c)\alpha'=\pgcyc(g)\alpha'$. So $\alpha-\pgcyc(g)\alpha'=\alpha(1-
\pgcyc(g)^{t'})=\alpha(1-p_i^{\,nt'})$ kills $\pgHom_{\pgG}(T/a,(T/a)^{*c})$. This module is
cyclic: both $T/a$ and $(T/a)^{*c}$ are free of rank $n$ over $\auO/a$ with residually absolutely
irreducible actions; if $\eta=\varpi^k\eta_0$ is an equivariant map with $k$ least, then $\eta_0$,
an equivariant map modulo $a/\varpi^k$, is not zero modulo $\varpi$, hence an isomorphism, and
every equivariant $\eta'$, being divisible by $\varpi^k$, differs from $\eta$ by an endomorphism
of $T/(a/\varpi^k)$ commuting with $M_n(\auO)$, a scalar. Hence $\mathrm{gl}\le\pgv(p_i^{\,nt'}-1)\le\pgv(\pgNch)$, because
$p^{nt'}-1$ divides $\#\auGL_{rt}(\auF_p)$ for $nt'\le rt$.

\emph{The local terms.} For $v\mid p_0p_1$, $v\nmid\ell$, let $\gamma\in I_v$ be as in (M3). Then
$H^0(E_v,(T/a)^*)$ is contained in the kernel of $u(\gamma)^{-1,\mathrm t}-1$ on
$\pgHom(T/a,\Q/\Z)$ (inertia acts trivially on the Tate twist), whose length is at most
$\pgv(\det(u(\gamma)-1))\le n\,\pgv(\pgNch)$. There are $4f$ such places at most, so
$\mathrm{pl}\le4f\,n\,\pgv(\pgNch)$; and $\mathrm{fp}\le2n\,\pgv(\pgNch)$ by (M3) for the chosen
place $w_0$. If $\pgv(\pgNch)=0$ all terms vanish; otherwise the sum is at most $(1+4f+2)\,n\,
\pgv(\pgNch)$.
\end{proof}

\subsubsection*{One system of equations for all twists}
Let $N\ge M\ge1$. Let $J$ be a finite set of places of $E$ of prime norm $\equiv1\pmod{\ell^N}$
such that $J$, $\aubar J$, $S_0$ are pairwise disjoint, and for $v\in J$ let $\lambda_v\colon\pgG
\to\Z/\ell^N$ be a homomorphism that is anticyclotomic, unramified outside $v,\aubar v$, trivial
on $D_w$ for $w\mid\ell$, with $\lambda_v(\gamma_v)=1$. For a point $x=(x_v)_{v\in J}$ with
coordinates in the ring of integers $\auO'$ of a finite extension of $K$ in $\auQlb$ put
$\psi_x(g)=\prod_v(1+x_v)^{\lambda_v(g)}$ (exponents in $[0,\ell^N)$). If $\pgv(x_v)\ge\rho$ with
$M\le N\rho$, $\rho\le1$, then $(1+x_v)^{\ell^N}\equiv1\pmod{\ell^M}$ (each term
$\binom{\ell^N}{i}x^i$, $i=\ell^ai'$, has valuation at least $N-a+\ell^a\rho\ge N\rho$), so
$\psi_x$ is a character modulo $\ell^M$. For a Selmer datum $T$ let $\pgSel(x)$ be the module of
locally constant cocycles $\pgG\to(\auO'/\ell^M)^n$ for $\psi_xu$ that are unramified outside
$S\cup J\cup\aubar J$, obey the local relations of $T$ and vanish at $\gamma_*$. If $x_v=\zeta_v-1$
for $\ell^N$-th roots of unity $\zeta_v$, then $\psi_x=\prod\zeta_v^{\lambda_v}$ is a character
$\psi$ of $\pgG$ and $\pgSel(x)$ is the Selmer group of the twist $\psi u$ modulo $\ell^M$ (the
local relations are unchanged because $\psi$ is trivial on $D_w$, $w\mid\ell$).

\begin{lemma}\label{pg:minor}
There are a finite set $\mathcal P$ of polynomials in $\auO[X_v:v\in J]$ with the following
properties, for every $\auO'$ and every point $x$ as above. {\rm(i)} For every $P\in\mathcal P$:
$\pglen\pgSel(x)\le\pgv(P(x))$. {\rm(ii)} If $\pglen\pgSel(x)<M$ then some $P\in\mathcal P$ has
$\pgv(P(x))=\pglen\pgSel(x)$. {\rm(iii)} If $x=(\zeta_v-1)$ and $\pgv(P(x))<M$ then, for all large
coefficient rings, $\pglen\pgSel_J(\psi u,\gamma_*;K/\auO)\le\pgv(P(x))$.
\end{lemma}
\begin{proof}
Let $U\subseteq\pgG$ be the kernel of the reduction of $u$ modulo $\ell^M$ and of all $\lambda_v$.
The restriction to $U$ of a cocycle of $\pgSel(x)$ is a homomorphism $U\to(\auO'/\ell^M)^n$
unramified outside $S\cup J\cup\aubar J$, so it factors through the Galois group of the maximal
abelian extension of exponent $\ell^M$ of the fixed field of $U$ unramified outside these places,
which is finite. So there is a finite quotient $\Gamma$ of $\pgG$, independent of $x$ and
$\auO'$, such that $\pgSel(x)$ is the module of functions $z\colon\Gamma\to(\auO'/\ell^M)^n$
satisfying: $z(gh)=z(g)+\psi_x(g)u(g)z(h)$; $z(g)=0$ for $g$ in the images of the inertia groups
outside $S\cup J\cup\aubar J$ and for $g=\gamma_*$; and the local relations, of which finitely
many suffice because the $Z_v$ are finitely generated and only the values on $\Gamma$ matter. These
are linear equations whose coefficients are values at $x$ of polynomials over $\auO$: $\pgSel(x)=
\ker A(x)$ on $(\auO'/\ell^M)^{d}$, $d=n\#\Gamma$, for a matrix $A$ over $\auO[X]$ with $d$
columns. Let $\mathcal P$ be the set of $d\times d$ minors of $A$. (i) The kernel of $A(x)$ is
contained in the kernel of a square submatrix $B$, and for a square matrix over a discrete
valuation ring the kernel of $B$ on $(\auO'/a)^d$ has length $\sum_i\min(\pgv(a),\pgv(d_i))\le
\pgv(\det B)$, the $d_i$ being the elementary divisors. (ii) Let $d_1\mid d_2\mid\dots$ be the
elementary divisors of $A(x)$ over $\auO'$; the length of the kernel is
$\sum_{i\le d}\min(M,\pgv(d_i))$ (with $d_i=0$ beyond the rank). If it is less than $M$ then all
$\pgv(d_i)<M$, the rank is $d$ and the length is $\pgv(d_1\cdots d_d)$, the least valuation of a
$d\times d$ minor. (iii) Multiplication by $\ell^M$ on $\pgSel_J(\psi u,\gamma_*;K/\auO)$ has
kernel $\pgSel(x)$, of length less than $M=\pglen\auO/\ell^M$; if some element of the group were not killed by $\ell^M$, the cyclic submodule it
generates would have $\ell^M$-torsion of length $M$. So the group equals its $\ell^M$-torsion.
\end{proof}

\begin{lemma}[Tame inertia and Frobenius]\label{pg:tame}
Let $v\in J\cup\aubar J$, $\gamma$ a tame generator and $\phi$ a Frobenius element at $v$, and let
$z$ be a cocycle of $\pgSel(x)$, where $0<\rho\le1$, $\pgv(x_w)\ge\rho$ for all $w$ and $M\le N\rho$. Then
\[
\bigl(\psi_x(\phi)u(\phi)-1\bigr)z(\gamma)=\bigl(\psi_x(\gamma)-1\bigr)z(\phi).
\]
\end{lemma}
\begin{proof}
$u$ is unramified at $v$, so $\gamma$ acts by the scalar $\alpha=\psi_x(\gamma)$, an $\ell^N$-th
root of unity modulo $\ell^M$, with $\sum_{j<\ell^N}\alpha^j\equiv0\pmod{\ell^M}$. Indeed, write
$\alpha=1+y$; then $\pgv(y)\ge\rho$, because $\alpha$ is a product of integer powers of the
$1+x_w$, each congruent to $1$ modulo elements of valuation at least $\rho$. Then
$\sum_{j<\ell^N}\alpha^j=\sum_{i=1}^{\ell^N}\binom{\ell^N}{i}y^{i-1}$, and for $i=\ell^ai'$ with
$\ell\nmid i'$ (so $a\le N$) the $i$-th term has valuation at least
$N-a+(\ell^a-1)\rho\ge N\rho\ge M$, since $N-a+(\ell^a-1)\rho-N\rho=(N-a)(1-\rho)+(\ell^a-1-a)\rho\ge0$
for $0<\rho\le1$. Here $\phi$ may be an arithmetic or a geometric Frobenius element (both
occur below, e.g.\ $\phi_v^{-c}$ at $\aubar v$). The cocycle $z$ vanishes on wild inertia (a homomorphism from a pro-$p$ group to an
$\ell$-group), and $\phi\gamma\phi^{-1}=\gamma^{q}$ modulo wild inertia, where $q$ is a positive integer
congruent to $(\auNm v)^{\pm1}$ modulo $\ell^N$ (the sign depending on whether $\phi$ is
arithmetic or geometric); as $\auNm v\equiv1\pmod{\ell^N}$, $q\equiv1\pmod{\ell^N}$. Since
$z(\gamma^{k+\ell^N})=z(\gamma^k)+\alpha^k\bigl(\sum_{j<\ell^N}\alpha^j\bigr)z(\gamma)=z(\gamma^k)$
and $\alpha^{\ell^N}=1$, the value $z(\gamma^k)$ and $\alpha^k$ depend only on $k$ modulo
$\ell^N$. Hence $z(\gamma^q)=z(\gamma)$ and $\alpha^q=\alpha$. Comparing $z(\phi\gamma)=z(\phi)+\psi_x(\phi)u(\phi)z(\gamma)$ with
$z(\gamma^q\phi)=z(\gamma)+\alpha z(\phi)$ gives the relation.
\end{proof}

\subsubsection*{Tied Selmer groups and the descent}
A \emph{set of tied primes} is a set $J$ as above with a \emph{first prime} $P_0\in J$, Frobenius
elements $\phi_v$ and, for each datum $T_i$ of a finite family, \emph{slopes} $s_{v,i}\in\auO$
($v\in J$), such that: $s_{P_0,i}=0$; $u_i(\phi_{P_0})\equiv u_i(g_0)$ and
$u_i(\phi_{P_0}^{-c})\equiv u_i(g_0^{-c})$ modulo $\ell^M$; for $v\ne P_0$,
$\lambda_{P_0}(\phi_v)=1$ and $\lambda_w(\phi_v)=0$ for $w\in J\setminus\{P_0\}$; and
$u_i(\phi_v)\equiv u_i(\phi_v^{-c})\equiv1\pmod{\ell^M}$ whenever $s_{v,i}\ne0$. The \emph{tied
Selmer group} $\pgSel^{\mathrm{tie}}_i(J;a)$, $a\mid\ell^M$, is the module of cocycles
$z\colon\pgG\to(T_i/a)$ for $u_i$ (no twist), unramified outside $S\cup J\cup\aubar J$, obeying
the local relations, vanishing at $\gamma_*$, with the \emph{ties}
\[
z(\gamma_v)=s_{v,i}\,z(\phi_v),\qquad z(\gamma_v^{-c})=s_{v,i}\,z(\phi_v^{-c})\qquad(v\in J).
\]
A tie with slope $0$ says that $z$ is unramified at $v$ and $\aubar v$; a tie with any slope is a
local condition at $v$ of the same size as the unramified one, since a cocycle of the tame
quotient is determined by its values at $\gamma_v$ and $\phi_v$.

\begin{lemma}[The count]\label{pg:count}
Let $J$ be a set of tied primes and $a\mid\ell^M$. Then
\[
\pglen\pgSel^{\mathrm{tie}}_i(J;a)\le\pglen\pgSel^{\mathrm{tie},*}_i(J;a)+\mathrm{pl}(T_i),
\]
where $\pgSel^{\mathrm{tie},*}$ is the Selmer group of $(T_i/a)^*$ with the local conditions
orthogonal, under local Tate duality, to those of $\pgSel^{\mathrm{tie}}$ (no condition at the
places of $S$ prime to $\ell$ becomes the zero condition there).
\end{lemma}
\begin{proof}
By the formula of Greenberg and Wiles for Selmer groups of a finite module $X=T_i/a$ with
arbitrary local conditions $L_v\subseteq H^1(E_v,X)$ (\cite[Proposition~1.6]{WilesFLT} over
$\Q$ for $\ell$ odd; the proof, by the Poitou--Tate sequence and the global Euler characteristic formula, is the
same over a number field, see~\cite[Theorem~8.7.9]{NSW}),
\[
\pglen H^1_L-\pglen H^1_{L^\perp}=\pglen H^0(E,X)-\pglen H^0(E,X^*)+\sum_v\bigl(\pglen L_v-\pglen
H^0(E_v,X)\bigr),
\]
the sum over the places of $S\cup J\cup\aubar J$ and the $f$ complex places, where $L_v=0$ and
the term is $-\pglen X$. Here $H^0(E,X)=0$ because $\zeta_*-1$ is a unit, and the classes are
the cocycles vanishing at $\gamma_*$. We bound the local terms. For $v\mid\ell$: let $C_v$ be the
module of functions $D_v\to X$ obeying the relations of $Z_v$; they are cocycles, because the
cocycle rule is among the relations, and $C_v$ contains the coboundaries, because $Z_v$ does; $L_v$
is the image of $C_v$. Let $\Phi$ be the free $\auO$-module on the symbols $(g,i)$, $g\in D_v$,
$1\le i\le n$, so that functions $D_v\to B^n$ are the linear maps $\Phi\to B$, and let
$N\subseteq\Phi$ be the common kernel of the elements of $Z_v$. Then $\Phi/N$ embeds in
$\pgHom_{\auO}(Z_v,\auO)$, so it is free of rank at most $\operatorname{rk}Z_v$, and
$C_v=\pgHom_{\auO}(\Phi/N,\auO/a)$ has length at most $\operatorname{rk}Z_v\cdot\pgv(a)$. The
coboundaries have length $\pglen X-\pglen H^0(E_v,X)$; so the term is at most
$(\operatorname{rk}Z_v-n)\pgv(a)$, and the sum over $v\mid\ell$ is at most $f\,n\,\pgv(a)=f\pglen
X$, which cancels the complex places. For $v\in S$, $v\nmid\ell$, with no condition:
$\pglen H^1(E_v,X)-\pglen H^0(E_v,X)=\pglen H^2(E_v,X)=\pglen H^0(E_v,X^*)$ by the local Euler
characteristic formula and local duality; the sum is at most $\mathrm{pl}(T_i)$. For $v\in
J\cup\aubar J$: a cocycle of $D_v$ with values in $X$ is a homomorphism on $I_v$ and factors
through the quotient topologically generated by $\gamma_v$ and $\phi_v$; a tied one is determined
by $z(\phi_v)$, so the tied cocycles have length at most $\pglen X$. They contain the
coboundaries (which vanish at $\gamma_v$, and at $\phi_v$ when the slope is not $0$). So
$\pglen L_v\le\pglen X-(\pglen X-\pglen H^0(E_v,X))$ and the term is at most $0$. Finally
$\pglen H^0(E,X^*)\ge0$.
\end{proof}

\begin{proposition}[The dual of a tied Selmer group]\label{pg:mirror}
Let $J$ be a set of tied primes, $a=\varpi^k\mid\ell^M$, $M<N$, and suppose that
$\pgSel_i^{\mathrm{tie}}(J;a)$ contains a cocycle $z$ one of whose values has a unit coordinate.
Suppose that $z$ is balanced at every place $P$ admissible for $T_i$ (\S\ref{pse:ss-aux}): the pair
$(z(\phi_P),z(\phi_P^{-c}))$ is the pair of values at $(\gamma_P,\gamma_P^{-c})$ of a cocycle $z'$
with the local conditions and ties of $J$ that may ramify at $P$ and $\aubar P$. Then
$\pglen\pgSel_i^{\mathrm{tie},*}(J;a)\le\mathrm{gl}(T_i)$.
\end{proposition}
\smallskip\noindent Proved in Proposition~\ref{pse:mirror}.\par\smallskip

\begin{proposition}[One step of the descent]\label{pg:step}
Let $J$ be a set of tied primes for data $T_1,\dots,T_s$ over $\auO$ with uniformiser $\varpi$,
$\varpi^e\mid\ell^M$, and suppose that the residue field of $\auO$ has more than $2n_ie+1$ elements
for all $i$. If for some $i$, some $1\le k\le e$ and some place $P$ that can be added to $J$, an
element of $\pgSel^{\mathrm{tie}}_i(J;\varpi^k)$ is not balanced at $P$, then $P$ can be added to
$J$ with a unit slope for $T_i$ and slope $0$ for the other data so that
$\sum_{i'}\sum_{k'=1}^e\pglen\pgSel_{i'}^{\mathrm{tie}}(J;\varpi^{k'})$ decreases.
\end{proposition}
\smallskip\noindent Proved in Proposition~\ref{pse:step}.\par\smallskip

\begin{proposition}[Polynomials and points]\label{pg:poly}
Assume Propositions~\ref{pg:mirror}, \ref{pg:step} and statement \emph{(A)} below. Let $T_1,\dots,T_s$ be Selmer data over
$\auO$, let $M>\max_i\beta_{T_i}+1$, and suppose the residue field of $\auO$ has more than
$2n_i\pgv(\ell^M)/\pgv(\varpi)+1$ elements. There is an integer $s_0\ge1$ such that for every
rational $0<\rho<1/s_0$ and every large $N$ there are a set $J$ with characters $\lambda_v$ as
above, $2\#J\le s_0$, and for each $i$ a polynomial $P_i\in\pgZb[X_v:v\in J]$ and a point $x_i$
with coordinates of valuation at least $\rho$, such that $P_i$ has the properties {\rm(i)},
{\rm(iii)} of Lemma~\ref{pg:minor} for $T_i$ and $\pgv(P_i(x_i))\le\beta_{T_i}+s_0\rho$.
\end{proposition}
\smallskip\noindent\textbf{(A)} is Proposition~\ref{pse:A}; ``a place that can be added to $J$'' means a place admissible in the sense of \S\ref{pse:ss-aux}.\par\smallskip
\begin{proof}
\emph{The descent.} Start with $J=\{P_0\}$, where $P_0$ is a place of prime norm $\equiv1
\pmod{\ell^N}$ outside $S_0\cup\aubar{S_0}$ whose Frobenius element is so near $g_0$ that it acts
as $g_0$, and $\phi_{P_0}^{-c}$ as $g_0^{-c}$, on all $T_i/\ell^M$, with an anticyclotomic
character $\lambda_{P_0}$ of level $P_0$, $\lambda_{P_0}(\gamma_{P_0})=1$; such places and characters exist by Proposition~\ref{pse:A} (see Remark~\ref{pse:rem-poly}). All slopes are $0$, and the tied groups are the Selmer
groups without auxiliary primes, of total length $b_0$ over the levels $\varpi^k$, $k\le e$,
independent of $N$; here, and only here, $b_0$ is the \emph{unnormalised} length (the number of
composition factors $\auO/\varpi$), an integer, so that a strict decrease is a decrease by at least
$1$. Apply Proposition~\ref{pg:step} as long as it applies; each step lowers the
total length, so after at most $b_0$ steps every element of every tied group modulo every
$\varpi^k$ is balanced at every admissible $P$. Then for each $i$: by passage to the exponent,
there is $k'\le e$ with $\pglen\pgSel^{\mathrm{tie}}_i(J;\varpi^e)\le\pglen\pgSel^{\mathrm{tie}}_i
(J;\varpi^{k'})$ and, if $k'\ne0$, an element of the latter with a unit coordinate (take $k'$
least such that $\varpi^{k'}$ kills the group modulo $\varpi^e$; its elements have values in
$\varpi^{e-k'}T/\varpi^e\cong T/\varpi^{k'}$, and one not killed by $\varpi^{k'-1}$ has, there, a
value outside $\varpi T/\varpi^{k'}$). By Lemma~\ref{pg:count} and Proposition~\ref{pse:mirror},
$\pglen\pgSel^{\mathrm{tie}}_i(J;\ell^M)\le\mathrm{gl}(T_i)+\mathrm{pl}(T_i)$. Put $s_0=2\max(n,2)(1+b_0)$,
$n=\max n_i$.

\emph{Release of the first prime.} Let $\pgSel'_i$ be the group with the ties at $J\setminus
\{P_0\}$ only. By Lemma~\ref{pg:tame} at the point $x=0$ and $v=P_0$, every $z\in\pgSel_i'$ has
$(u_i(g_0)-1)z(\gamma_{P_0})=0$, and likewise at $\aubar{P_0}$; so the quotient
$\pgSel'_i/\pgSel_i^{\mathrm{tie}}$ embeds in the product of the kernels of $u_i(g_0)-1$ and
$u_i(g_0^{-c})-1$ on $T_i/\ell^M$ and has length at most $\mathrm{fp}(T_i)$. Thus
$\pglen\pgSel_i'\le\beta_{T_i}<M$.

\emph{The line.} Fix $i$ and consider the points $x(y)$ with $x_{P_0}=y$, $x_v=s_{v,i}y$
($v\ne P_0$), $\pgv(y)\ge\rho$. For $v\ne P_0$ in $J$ we have $\psi_{x(y)}(\gamma_v)=1+s_{v,i}y$
and $\psi_{x(y)}(\phi_v)=1+y$. Let $z\in\pgSel(x(y))$. By Lemma~\ref{pg:tame},
$((1+y)u_i(\phi_v)-1)z(\gamma_v)=s_{v,i}y\,z(\phi_v)$, and therefore
\[
\bigl((1+y)u_i(\phi_v)-1\bigr)\bigl(z(\gamma_v)-s_{v,i}z(\phi_v)\bigr)=-s_{v,i}(1+y)\bigl(u_i(
\phi_v)-1\bigr)z(\phi_v)=0,
\]
because $u_i(\phi_v)\equiv1$ if $s_{v,i}\ne0$; the same holds at $\aubar v$. Let
$\pgSel^{\mathrm{tie}}(y)\subseteq\pgSel(x(y))$ be the submodule on which the ties at
$J\setminus\{P_0\}$ hold. The quotient embeds in the product over $v\in J\setminus\{P_0\}$ and
$g\in\{\phi_v,\phi_v^{-c}\}$ of the kernels of $(1+y)u_i(g)-1$, so its length is at most
$\pgv(P^1(y))$, $P^1(Y)=\prod_{v,g}\det((1+Y)u_i(g)-1)$. Each factor has the unit coefficient
$\det u_i(g)$ in degree $n_i$; hence $\pgv_\rho(P^1)\le2n_i\#J\rho\le s_0\rho$. The modules
$\pgSel^{\mathrm{tie}}(y)$ are the kernels of one matrix over $\auO[Y]$, as in
Lemma~\ref{pg:minor}; at $y=0$ the module is $\pgSel'_i$, of length less than $M$, so by
Lemma~\ref{pg:minor}(ii) there is a minor $P^0\in\auO[Y]$ with $\pgv(P^0(0))=\pglen\pgSel'_i\le
\beta_{T_i}$ and $\pglen\pgSel^{\mathrm{tie}}(y)\le\pgv(P^0(y))$ for all $y$. In particular
$\pgv_\rho(P^0)\le\pgv(P^0(0))\le\beta_{T_i}$. Choose $y$ of valuation exactly $\rho$ (in a
ramified extension) with $\pgv(P^0P^1(y))=\pgv_\rho(P^0P^1)$; this excludes finitely many residue
classes of $y/\varpi'^{\,\rho}$. Then
\[
\pglen\pgSel(x(y))\le\pgv(P^0(y))+\pgv(P^1(y))\le\beta_{T_i}+s_0\rho<M .
\]
By Lemma~\ref{pg:minor}(ii) there is $P_i\in\mathcal P$ with $\pgv(P_i(x(y)))=\pglen\pgSel(x(y))$;
take $x_i=x(y)$. All its coordinates have valuation at least $\rho$.
\end{proof}

\subsubsection*{The $L$-values}

\begin{proposition}[Polynomials for the ratios of $L$-values]\label{pg:stick}
In the situation of Proposition~\ref{pg:poly}, let finitely many pairs of a ramified critical
pair $(\nu,z)$ and a sign $\delta=\pm1$ be given, with $M>\beta_Y+1$ for each ($\beta_Y$ as in
\S\ref{pse:ss-stick}). Then, after enlarging $s_0$, for every rational $0<\rho<1/s_0$ and every
large $N$ the set $J$ with its characters can be chosen so that moreover, for each of the given
pairs, there are a polynomial $P\in\pgZb[X_v:v\in J]$ and a point $x$ with coordinates of valuation
at least $\rho$ with
\[
\pgv(P(x))\le v_0+4\pgv(\auNm w_0-1)+s_0\rho ,
\]
where $v_0$ is the valuation at $\iota$ of the value at the trivial twist of the measure of
statement (IM) of \S\ref{pse:ss-stick} and $w_0$ is a chosen place, and such that, for all
$\ell^N$-th roots of unity $\zeta_v\ne1$, with $\psi=\prod_v\zeta_v^{\lambda_v}=\psi_\vartheta$,
\[
\pgv\Bigl(\frac{|L(z,(\nu\vartheta^{\delta})^{-1})|}{|L(z,\nu^{-1})|}\Bigr)\le
\max\bigl(0,\ \pgv\bigl(P((\zeta_v-1)_v)\bigr)-v_0\bigr).
\]
\end{proposition}
\smallskip\noindent This is Proposition~\ref{pse:stick}, proved there from statement (IM), which is
proved in Propositions~\ref{pex:IM} and~\ref{pex:IM1}.\par\smallskip

\begin{proof}[Proof of Theorem~\ref{pg:II} from Lemma~\ref{pg:beta} and
Propositions~\ref{pg:roots}, \ref{pg:poly}, \ref{pg:stick}]
Let $c=(4f+3)\,r\,\pgv(\pgNch)$ and $\varepsilon=2^{-\ell}$. By Lemma~\ref{pg:beta} there is
$\auO$ over which each $R_i$ has a lattice $u_i$ that is a Selmer datum $T_i$ with $\beta_{T_i}\le
c$; enlarging $\auO$ (the constants do not grow, lengths being normalised) we may assume its
residue field is as large as Proposition~\ref{pg:poly} requires for an integer $M>\max(c,\max_Y\beta_Y)+1$ (\S\ref{pse:ss-end}). Let $s_0$
be given by Propositions~\ref{pg:poly}, \ref{pg:stick}, choose a rational $\rho$ with
$s_0\rho\le\varepsilon/2$, let $h_0$ be given by Proposition~\ref{pg:roots} for $s_0$, $m$,
$\rho$, the bound $\beta=s(c+1)+\sum_Y(v_0(Y)+4\pgv(\auNm w_0-1)+1)$ and
$\varepsilon/2$, and let $N\ge\max(h_0,M/\rho)$ be large. Take $J$, the $\lambda_v$, the
polynomials and the points of Propositions~\ref{pg:poly} and~\ref{pg:stick}, and let
$\zeta_v$ ($v\in J$) be the roots of unity of Proposition~\ref{pg:roots}; their orders divide
$\ell^{h_0}\mid\ell^N$ and exceed $m$. Let $\psi_1=\prod_v\zeta_v^{\lambda_v}$. It is a character of
$\pgG$ of $\ell$-power order, anticyclotomic, unramified outside $J\cup\aubar J$, trivial on the
decomposition groups above $\ell$, tamely ramified at $v\in J$ with $\psi_1(\gamma_v)=\zeta_v$ of
order greater than $m$. By the remarks at the beginning of \S\ref{pg:ss-reduction},
$\psi_1=\psi_\tau$ for an $m$-fold twist $\tau$ with $Q_\tau=J$, and $\psi_{\tau^j}=\psi_1^{\,j}$
corresponds to the point $(\zeta_v^{\,j}-1)_v$.

(a) $J$, $\aubar J$, $S_0$ are pairwise disjoint and the norms are prime; and
$\sum_{v\in J}\pgv(\zeta_v^{\,j}-1)\le\#J\cdot\rho\le s_0\rho\le\varepsilon$ (if $\zeta$ has order $\ell^h>m\ge j$ then
$\zeta^j$ has order $\ell^{h-a}$, $\ell^a\mid j$, and $\pgv(\zeta^j-1)=\ell^a\pgv(\zeta-1)\le
j\,\rho/m\le\rho$). The same bound holds at the conjugate places.

(b) For each $i$ and $j$: $\pgv(P_i((\zeta_v^{\,j}-1)_v))\le\pgv(P_i(x_i))+\varepsilon/2\le
\beta_{T_i}+\varepsilon<M$; by Lemma~\ref{pg:minor}(iii) the Selmer group of $\psi_1^{\,j}u_i$ with
coefficients in $K/\auO$ has length at most $c+\varepsilon=e'_\ell$ for all large coefficient
rings. Its local conditions are those of Definition~\ref{pg:def-sel}, and $\gamma=\gamma_*$,
$\zeta=\zeta_*$ serve, since $\psi_1(\gamma_*)=1$.

(c) By Proposition~\ref{pse:stick} and \S\ref{pse:ss-end}(4), the valuation of each ratio at $\vartheta=\tau^j$ is at most $4\pgv(\pgNch)+\varepsilon\le c+\varepsilon$; this is the lower bound $-e'_\ell$ for the
valuation of the inverse ratio.
\end{proof}

\subsection{From Selmer groups to denominators}\label{pg:ss-ribet}

Theorem~\ref{pg:III} is proved in \S\ref{pri:sec} (Theorem~\ref{pri:III}) from Propositions~\ref{pri:loss}, \ref{pri:odd}, \ref{pri:kummer}, \ref{pri:dict}, \ref{pri:BL} and Lemmas~\ref{pri:chain}, \ref{pri:block}, \ref{pri:forms}; the propositions below are restated there in exact form. Fix $\ell$, $\iota$, $\vartheta$, $O$, $\mathcal Z$ as in the theorem,
$n=n_V$, $r'=n_W$, $t'=t(O)$, and one generator $y_0$ of $\mathcal Z$, with system $x_0$ and main
summand $R$.

\emph{The mechanism.} Let $y$ be integral at $\iota$ and orthogonal to the false lifts. Its
component orthogonal to $\mathcal Z$ lies, in each type $e_c$, in the span of the eigenforms that
are not \emph{weak lifts} (an eigenform is a weak lift if $\det(1-(S\otimes\eta_O)(g))=0$ for all
$g$); this is where orthogonality to the false lifts is used (Lemma~\ref{pri:weak}). A denominator $\ell^\nu$ in
$\pgpr_{\mathcal Z}(y)$ at the component of $y_0$ is therefore a congruence modulo $\ell^\nu$
between $y_0$ and eigenforms $y_x$ whose normalised Galois representations $S_x$ do not have the
eigenvalue $1$ identically, while $S_{x_0}=\psi R\oplus\bigoplus_{e<t'}\chi^e$ on traces. Ribet's
method produces from such a congruence a non-split extension between two of the $t'+1$ summands
modulo a power of $\ell$. Extensions between $\psi R$ and a character are bounded by the Selmer
group of $\psi R$ (Theorem~\ref{pg:II}(b)), up to the local loss above $\ell$. Extensions between
two characters $\chi^e$, $\chi^{e'}$ are classes in Tate twists of the trivial module: the
\emph{odd} ones are bounded by $\pgv(\pgcodd)$, and for the \emph{even} ones, whose order is
governed by class groups of real cyclotomic fields and cannot be bounded, one passes to
forms of auxiliary level at Taylor--Wiles places $w_j$ chosen in a triangular way, where the
even classes are killed one after another, and works to first order in parameters $\alpha_j$ that
tie tame inertia at $w_j$ to Frobenius.

\begin{proposition}[Ribet's method to first order]\label{pg:ribet}
Let $G$ be a group, $\auO$ a discrete valuation ring with uniformiser $\varpi$, $p=\ell^\nu$, and
$u_x\colon G\to\auGL_n(\auO)$, $x\in Y$ finite, representations such that for $x$ in a subset
$\mathcal X$ the representation $u_x$ is block diagonal with a block $\rho_x$ of rank $r'$ and
the characters $\xi_e$, $e<t'$, $\xi_0=1$. Let $\tau_j,\phi_j\in G$ ($j\in I$), polynomials
$H_j$ with $H_j(1)=1$, and let $\mathfrak a$ be the left ideal of $\auO[X_j:j\in I][G]$ generated
by the elements killed by all $u_x$, by $p$ times the augmentation ideal, and by the \emph{tied
relations} $\tau_j-1-X_j(\phi_j-1)H_j(\phi_j)$. Assume (in the exact form of
Proposition~\ref{pri:ribet}): a triangle
cost $b_e$ for the cocycles of each pair of characters, a bound $b_0$ for the length of the
cocycles of the block with coefficients in $K/\auO$ obeying the relations among the $u_x(g)-1$ and
vanishing at the $\tau_j$, and an element separating the characters up to a cost $C$. Then there
is an element $\mathcal D=\sum_{S\subseteq I}\mathcal D_SX_S$ of the algebra generated over
$\auO[X]$ by the coefficients of the characteristic polynomials of the $u_x(g)$ (as functions of
$x$) such that $\mathcal D\cdot\det(1-u_x(g))$ lies, for every $g$, in the ideal of traces of
$\mathfrak a$; $\mathcal D_S$ vanishes on $\mathcal X$ for $S\ne I$; and $\mathcal D_I$ takes on
$\mathcal X$ a value in $\auO$ of valuation at most $b_0+\sum_eb_e+C$.
\end{proposition}
\smallskip\noindent Exact statement and proof: Proposition~\ref{pri:ribet}.\par\smallskip

\begin{proposition}[Taylor--Wiles places in a triangle]\label{pg:triangle}
For every $\nu\ge1$ and every prescribed open condition on Frobenius elements that is satisfied by
the elements of a non-empty open subset of the kernel of the cyclotomic character modulo
$\ell^\nu$, there are places $w_j\nmid\ell$ of degree one, $j$ in a finite ordered set $I$ whose
size does not depend on $\nu$, with $\auNm w_j\equiv1\pmod{\ell^\nu}$, Frobenius elements $\phi_j$
satisfying the condition, and exponents $0<\mathrm{ty}_j<t'$, such that every even class in
$(\auO/\ell^\nu)(\chi^{\mathrm{ty}_j})$ that vanishes at the earlier $\tau_i$ and at the earlier
$\phi_i$ of the same exponent has its value at $\tau_j$ killed by $\ell^{h'}$, where $h'\le
h\,\pgv(\pgNch)$ and $h'=0$ if $\ell>h^2$. Here a class is \emph{even} if it is invariant under
$g\mapsto g^c$, unramified outside $\ell$ and the $w_j$, and vanishes on the elements of trivial cyclotomic character in the decomposition groups above $\ell$.
\end{proposition}
\smallskip\noindent Proved in Proposition~\ref{pri:triangle}, with Proposition~\ref{pri:kummer} proved in
\S\ref{pku:sec}.\par\smallskip

\begin{proposition}[Odd classes]\label{pg:odd}
There is a positive integer $\pgcodd\le2^{4(q_1d+t+3)^{27}}h_E$ such that for every prime $\ell$,
all large $\auO$, all $0\le e<t$ and $\nu\ge1$ the module of cocycles $\pgG\to(\auO/\ell^\nu)(
\chi^e)$ that, for $e\ne0$, are unramified outside $\ell$, satisfy the local condition above $\ell$ of
Proposition~\ref{pg:triangle} and change sign under $g\mapsto g^c$, and, for $e=0$, the module of
everywhere unramified homomorphisms $\pgG\to\auO/\ell^\nu$, has length at
most $\pgv(\pgcodd)$. In particular $\log\pgcodd\le x^{800}$ (Remark~\ref{pod:rem-size}).
\end{proposition}
\smallskip\noindent Exact statement: Proposition~\ref{pri:odd}; proved in Proposition~\ref{pod:odd}.\par\smallskip

\begin{proposition}[The local loss]\label{pg:loss}
Let $v\mid\ell$, let $\auO$ be large, and let $\mathcal C_v$ be a class of lattices in potentially
crystalline representations of $D_v$ of the following kind: crystalline with Hodge--Tate weights in
prescribed sets, or, if $\ell=p_i$, potentially crystalline with tame inertial type of the shape
allowed by (G4). Let $L$ be a lattice of $\mathcal C_v$ that is the restriction of a lattice in a
main summand twisted as in Theorem~\ref{pg:III}, or of a character $\chi^e$, $e<t'$, where all
weights of all members of $\mathcal C_v$ are weights of these summands. Then for every finite family $u_x$ of global lattices whose restrictions to $D_v$ lie in $\mathcal C_v$, the module of cocycles $\pgG\to L/\ell^M$ obeying all $\auO$-linear relations among the families $(u_x(g)-1)_x$ has, modulo those obeying the crystalline relations of $L$ at $v$, length at most the \emph{defect} $\delta_v(L)$ (the length of the crystalline cocycles of the Tate dual of $L$ modulo those cup-orthogonal to $\mathcal C_v$; see Proposition~\ref{pri:loss}),
and
\[
\sum_{v\mid\ell}\ \sum_{L}\delta_v(L)\le r'\cdot\pgloc_\ell ,
\]
the inner sum over the $t'+1$ summands. In particular there is no loss if $\ell>2b$ and $\ell\nmid M$ (Theorem~\ref{pll:exponent}(d); for $\ell\ge4b+2$, $\ell\nmid q_1dM$ see also Remark~\ref{pll:FL} for another route).
\end{proposition}
\smallskip\noindent\emph{Proof.} The exact form of the statement is Proposition~\ref{pll:main}. The
restrictions to $D_v$ of the cocycles in question come from $\mathcal C_v$ by Lemma~\ref{pll:ribet}
(with $\Gamma=D_v$ inside $\pgG$; its hypothesis is exactly that they obey the $\auO$-linear relations
among the $(u_x(g)-1)_x$), so the first assertion is Proposition~\ref{pll:dual} with $T=L\in\mathcal C_v$. Its
first inequality is Proposition~\ref{pll:dual}; the second is Theorem~\ref{pll:exponent} with
Lemma~\ref{pll:exp} (Remark~\ref{pll:status}), Theorem~\ref{pll:exponent} being proved in
\S\ref{pki:ss-deduce} from Propositions~\ref{pki:generic}--\ref{pki:data}, which are proved in
\S\S\ref{pkc:sec}, \ref{pkd:sec}, \ref{pkj:sec}; the sum over places and summands is
Lemma~\ref{pll:sum}. The clause ``no loss if $\ell>2b$ and $\ell\nmid M$'' is
Theorem~\ref{pll:exponent}(d) (Remark~\ref{pll:status}); for $\ell\ge4b+2$, $\ell\nmid q_1dM$ see
also Remark~\ref{pll:FL} for another route.\par\smallskip

\begin{proposition}[The operators at the auxiliary places]\label{pg:norm}
Let $w_j$, $j\in I$, be places as in Proposition~\ref{pg:triangle} off the level, and consider
the forms of principal level at the $w_j$. Let $\mathbb T$ be the algebra of functions on Hecke
eigensystems that preserve the forms integral at $\iota$. Suppose that at each $w_j$ the Hecke
polynomial of every eigensystem of base level congruent to $x_0$ on which a given $\delta\in
\mathbb T$ does not vanish is, after normalisation, $\prod_{i<t'}(X-(\auNm w_j)^i)$ times a divisor
of $H_j$. Then there are integral operators $\mathcal N_S$, $S\subseteq I$, commuting with
$\mathbb T$, with $\mathcal N_\emptyset$ the sum over $\prod_j\auGL_n(\auF_{w_j})$, such that
$\mathcal N=\sum_S\mathcal N_SX_S$ kills the traces of the ideal $\mathfrak a$ of
Proposition~\ref{pg:ribet} modulo $\ell^\nu/(2\,n!)$ and $\delta^{\#I}\mathcal N$ factors through
$\mathcal N_\emptyset$; consequently, if $\mathcal D$ is as in Proposition~\ref{pg:ribet}, the top
coefficient $\mathcal D_I$, times the least common multiple of the orders of the stabilisers and
times the idempotent of the congruence class of $x_0$, maps $\pgpr_{\mathcal Z}(y)$ to a form
integral at $\iota$, for every $y$ as in Definition~\ref{pg:def-den}.
\end{proposition}
\smallskip\noindent Proved in Lemma~\ref{pri:oneplace} and Proposition~\ref{pri:norm} from
Proposition~\ref{pri:dict} and Lemma~\ref{pri:forms}, which are proved in \S\S\ref{pax:ss-dict},
\ref{pay:ss-cons}, \ref{pax:ss-forms}, \ref{pax:ss-pairing}.\par\smallskip

\emph{How the bound $\beta_\ell$ arises.} In the proof of Theorem~\ref{pri:III}, the eigenvalue at
$y_0$ of the operator
that clears the denominator has valuation at most
\[
\underbrace{e'_\ell}_{\text{Selmer}}+\underbrace{2t'\,e'_\ell}_{\text{ramification of }\psi
\text{ at }Q}+\underbrace{r'\,\pgloc_\ell}_{\text{Proposition \ref{pg:loss}}}+\underbrace{e_\ell}_{
\text{odd classes, level, triangle, characters}},
\]
and one more factor $\pgNch$ accounts for the idempotents; with $r'\ge1$ this is at most
$r'\beta_\ell$, and the determinant of the operator on $\mathcal Z$ has valuation at most
$\dim\mathcal Z\cdot r'\beta_\ell$ (the operator is built generator by generator and is diagonal
in a basis of eigenforms; Steps~1--3 of the proof of Theorem~\ref{pri:III}).

\subsubsection*{The $p$-adic Hodge theoretic input}
The local conditions above $\ell$ enter in exactly two places. (1) In the twisting theorem only
through the \emph{rank} of the module of crystalline cocycles (Lemma~\ref{pg:beta}); this is the
formula of Bloch and Kato and needs nothing integral. (2) In Theorem~\ref{pg:III}: the cocycles
modulo $\ell^M$ that Ribet's method extracts from a congruence are the Ribet cocycles of
Proposition~\ref{pg:loss}, and one needs that they obey the crystalline relations of the
summand, up to a bounded defect. If $\ell$ is unramified in $E$, $\ell\nmid M$ and all Hodge--Tate
weights of all representations involved lie in an interval of length at most $\ell-2$, one
expects defect $0$ from the theory of Fontaine and Laffaille~\cite{FontaineLaffaille} together with
\cite[Lemmas~4.4--4.5]{BlochKato}; this route needs an integral statement on lattices that we do
not claim (Remark~\ref{pll:FL}), and the paper does not use it. Since the weights here are bounded by $2b$ in absolute value, this would concern
the primes $\ell\ge4b+2$ that are prime to $q_1dM$. For the remaining primes ($\ell<4b+2$ or $\ell\mid q_1dM$), in particular all $\ell\le2b$, which include $\ell=2$ and the primes inert in $K$ below $x_1$ that carry the prime
sum, no published result gives the bound: the weights (of size $b$ and $a/x_1$, with gaps
$\Delta$) are far outside the Fontaine--Laffaille range, and results on torsion crystalline
representations in large weight (Liu) give a defect bounded in terms of the weights and the
ramification but not the explicit bound $\pgloc_\ell$, whose shape $b/(\ell-1)$ per unit of rank,
times $m/k$, is what the budget $\aubud$ can afford (Lemma~\ref{pg:sums}). The argument is an
analysis of Kisin modules of the lattices involved (duality: the defect is the length of the
crystalline classes of the Tate dual modulo those orthogonal under cup product to the class
$\mathcal C_v$, Proposition~\ref{pll:dual}; then an exponent of $\ell$ that carries crystalline
cocycles into the orthogonal, Lemma~\ref{pll:exp}, computed on Breuil--Kisin modules, using that the
weights are either small or within a window near $\pm\kappa_j$ and that few are small; the factor
$m/k$ comes from the sparseness of $\auPb$). This is Theorem~\ref{pll:exponent}, proved in
\S\ref{pki:ss-deduce}. \emph{The $p$-adic Hodge theoretic input is not replaced by a citation for any $\ell$}; for
$\ell\ge4b+2$, $\ell\nmid q_1dM$, Remark~\ref{pll:FL} describes a possible alternative route.

\subsection{The inputs of Proposition~\ref{au:regular}}\label{pg:ss-summary}

Proposition~\ref{au:regular} follows from Theorem~\ref{pg:main} by the sums over $\ell$ (Lemma~\ref{pg:sums}),
given $\log\pgcodd\le x^{800}$; Theorem~\ref{pg:main} follows from Theorems~\ref{pg:I}--\ref{pg:III} and
Proposition~\ref{au:mult}. Theorem~\ref{pg:I} is proved in \S\ref{pga:sec} (with Proposition~\ref{pax:sign} and
Lemma~\ref{pax:conv}) from published results; statements (L), (C) are Propositions~\ref{pga:L}, \ref{pga:C}.
Theorem~\ref{pg:II} uses Propositions~\ref{pse:mirror}, \ref{pse:step}, \ref{pse:A}, \ref{pse:stick} of
\S\ref{pse:sec} and (IM), which is Propositions~\ref{pex:IM}, \ref{pex:IM1}, besides published results
(\cite{BlochKato}, \cite{NSW}). Theorem~\ref{pg:III} is proved in \S\S\ref{pri:sec}, \ref{prm:sec} from
Proposition~\ref{pg:loss} (\S\ref{pll:sec}, with Theorem~\ref{pll:exponent} proved in \S\ref{pki:ss-deduce}) and
Proposition~\ref{pg:odd} (Proposition~\ref{pod:odd}); see \S\ref{pri:ss-summary}.



%% ===== proofs/C-selmer.tex =====
% proofs/C-selmer.tex
% Completion of the twisting theorem (Theorem pg:II): the dual of a tied
% Selmer group, the step of the descent, the auxiliary primes, and the systems of the L-values.
% Labels of this file start with pse:. Macros as in Section pg (prefix \pg, \au).
\providecommand{\auO}{\mathcal O}
\providecommand{\auNm}{\operatorname{Nm}}
\providecommand{\auU}{\mathrm U}
\providecommand{\auA}{\mathbb A}
\providecommand{\auF}{\mathbb F}
\providecommand{\aubar}[1]{\overline{#1}}
\providecommand{\auf}[1]{\texttt{#1}}
\providecommand{\panot}[1]{\par\smallskip{\raggedright\noindent\textbf{[Not proved here:} #1\textbf{]}\par}\smallskip}
\providecommand{\pgZb}{\overline{\Z}_\ell}
\providecommand{\pgv}{v}
\providecommand{\pgNch}{N_{\mathrm{ch}}}
\providecommand{\pgG}{G_E}
\providecommand{\pgSel}{\operatorname{Sel}}
\providecommand{\pglen}{\operatorname{len}}
\providecommand{\pgHom}{\operatorname{Hom}}
\providecommand{\pgS}{S_{\mathrm{bad}}}
\providecommand{\pgcyc}{\varepsilon_\ell}
% macros of this file (prefix \pse)
\providecommand{\pseinv}{\operatorname{inv}}
\providecommand{\pseev}{\operatorname{ev}}
\providecommand{\pseCl}{\mathrm{Cl}}
\providecommand{\pseArt}{\operatorname{Art}}

\section{The twisting theorem: Selmer groups and \texorpdfstring{$L$}{L}-values}\label{pse:sec}

This section completes the proof of Theorem~\ref{pg:II}. In \S\ref{pg:ss-twisting} that theorem
was deduced from Propositions~\ref{pg:mirror}, \ref{pg:step}, \ref{pg:stick} and statement~(A).
We prove (A) (\S\ref{pse:ss-aux}), Proposition~\ref{pg:mirror} in the form of Proposition~\ref{pse:mirror}
(\S\ref{pse:ss-mirror}), Proposition~\ref{pg:step} (\S\ref{pse:ss-step}), and
Proposition~\ref{pg:stick} in the form of Proposition~\ref{pse:stick} (\S\ref{pse:ss-stick}) \emph{from one input on
theta lifts}, statement~(IM) of \S\ref{pse:ss-stick}, which is proved in \S\ref{pex:sec}
(Propositions~\ref{pex:IM}, \ref{pex:IM1}). \S\ref{pse:ss-end} collects the choices of constants used in the proof of
Theorem~\ref{pg:II} in \S\ref{pg:ss-twisting}; Theorem~\ref{pg:II} depends on published results and
on~(IM), and on nothing else.

Notation is that of \S\ref{pg:ss-twisting}: $E$ is the CM field, $c$ the fixed complex conjugation
of $\overline{\Q}$, $g^c=cgc$ and $g^{-c}=(g^c)^{-1}$ for $g\in\pgG$; $S_0=\pgS(\ell)$; $w_0\nmid\ell$
is a chosen place and $g_0=\gamma_{w_0}$ its tame generator. Lengths of $\auO$-modules are
normalised as in \S\ref{pg:ss-twisting}. For a finite $\auO$-module $X$ with $\pgG$-action,
$X^*=\pgHom(X,\Q/\Z)(1)$; $g$ acts on $f\in X^*$ by $(gf)(x)=\pgcyc(g)f(g^{-1}x)$.
For a place $B$ of $\overline{\Q}$ over a finite place $P\nmid\ell$ of $E$ we write $I_B\subseteq
D_B$ for the inertia and decomposition groups.

\subsection{Tame generators, and the local pairing at a split prime}\label{pse:ss-tame}

Fix a compatible system $(\xi_m)_{m\ge1}$ of primitive $m$-th roots of unity in $\overline{\Q}$.
Let $P\nmid\ell$ and $B\mid P$. An element $\delta\in I_B$ is a \emph{normalised tame generator for
$m$} (prime to $P$) if for some uniformiser $\varpi$ of $E_P$ and some $m$-th root $\varpi^{1/m}$ in
$\overline{\Q}_B$ one has $\delta(\varpi^{1/m})=\xi_m\varpi^{1/m}$. The condition does not depend on
the uniformiser: for $u\in\mathcal O_{E_P}^\times$ the roots $u^{1/m}$ generate an unramified
extension of $E_P$ ($m$ is prime to $P$), so $\delta$ fixes them and $\delta((u\varpi)^{1/m})/
(u\varpi)^{1/m}=\delta(\varpi^{1/m})/\varpi^{1/m}$. In particular we may take $\varpi\in E$. Normalised tame generators
exist (the tame inertia group $I_B/I_B^{\mathrm{wild}}$ is $\prod_{p\ne q}\Z_p(1)$, $q$ the residue
characteristic, and maps onto $\mu_m$ by $\delta\mapsto\delta(\varpi^{1/m})/\varpi^{1/m}$, independently
of the choice of $\varpi^{1/m}$).

\begin{lemma}\label{pse:conjgen}
If $\delta$ is a normalised tame generator for $m$ at $B$, then $\delta^{-c}$ is one at $cB$, a
place over $\aubar P$.
\end{lemma}
\begin{proof}
$\delta^{-c}=c\delta^{-1}c$ lies in $I_{cB}$. Take the uniformiser $\varpi\in E$ (allowed by the remark after the
definition) and put $\alpha=c(\varpi^{1/m})$, an $m$-th root of the
uniformiser $\aubar\varpi$ of $E_{\aubar P}$. Then $\delta^{-c}(\alpha)=c\,\delta^{-1}(\varpi^{1/m})
=c(\xi_m^{-1}\varpi^{1/m})=\xi_m\alpha$, because $c$ inverts roots of unity.
\end{proof}

\begin{lemma}\label{pse:kpower}
Let $P\nmid\ell$, $B\mid P$, $\gamma\in I_B$ any element, $\delta$ a normalised tame generator for
$\ell^N$ at $B$, and $X$ a finite $\auO$-module killed by $\ell^N$. There is an integer $k$ such that
for every continuous map $z\colon\pgG\to X$ whose restrictions to $I_B$ and to $I_{cB}$ are
homomorphisms one has $z(\gamma)=kz(\delta)$ and $z(\gamma^{-c})=kz(\delta^{-c})$.
\end{lemma}
\begin{proof}
A continuous homomorphism $I_B\to X$ kills wild inertia (a pro-$q$ group, $q\ne\ell$) and factors
through the maximal quotient of exponent $\ell^N$ of tame inertia, which is cyclic of order
$\ell^N$ and generated by the image of $\delta$. Hence $\gamma\equiv\delta^k$ modulo the common
kernel $K_B$ of these homomorphisms, for some $k$. Conjugation by $c$ carries $I_B$, $K_B$ onto
$I_{cB}$, $K_{cB}$, so $\gamma^{c}\equiv(\delta^{c})^k$ modulo $K_{cB}$ and, inverting in the abelian
quotient, $\gamma^{-c}\equiv(\delta^{-c})^k$.
\end{proof}

\begin{lemma}[The pairing at a split prime]\label{pse:pairing}
Let $P\nmid\ell$, $B\mid P$, and let $X$ be a finite $\auO$-module killed by $a\mid\ell^M$ on which
$D_B$ acts trivially, with $\auNm P\equiv1\pmod{\ell^M}$, so that $D_B$ acts trivially on $X^*$ as
well. Let $\phi\in D_B$ be a Frobenius element and $\delta$ a normalised tame generator for
$\ell^M$ at $B$. For a cocycle (here: a homomorphism) $z\colon D_B\to X$ and an unramified
homomorphism $y\colon D_B\to X^*$,
\[
\pseinv_B(z\cup y)=\epsilon\,\bigl\langle z(\delta),y(\phi)\bigr\rangle ,
\]
where $\langle\ ,\ \rangle\colon X\times X^*\to\mu_{\ell^M}\cong\frac1{\ell^M}\Z/\Z$ is the
evaluation (the isomorphism given by $\xi_{\ell^M}\mapsto1/\ell^M$), and $\epsilon\in\{\pm1\}$
depends only on the conventions for cup products and for the local invariant, not on $P$, $B$,
$X$, $z$, $y$.
\end{lemma}
\begin{proof}
Both sides are biadditive in $(z,y)$ and compatible with decompositions of $X$ into cyclic groups;
so we may assume $X=\Z/\ell^s$, $s\le M$, $X^*=\mu_{\ell^s}$. Then $z$ is a character
$\chi\colon D_B\to\Z/\ell^s$, and by Kummer theory $y$ is the class of some $b\in E_P^\times$,
a unit up to $\ell^s$-th powers because $y$ is unramified, with $y(\phi)=\phi(b^{1/\ell^s})/
b^{1/\ell^s}=\bar b^{(\auNm P-1)/\ell^s}$ (residue in $\mu_{\ell^s}$). By local class field theory
$\pseinv(\chi\cup\partial b)=\pm\chi(\pseArt(b))$, and for a unit $b$ the element $\pseArt(b)$ is the
element of tame inertia sending $\varpi^{1/\ell^s}$ to $\bar b^{\mp(\auNm P-1)/\ell^s}
\varpi^{1/\ell^s}$ (the tame symbol); that is, $\pseArt(b)\equiv\delta^{k}$ with $\xi_{\ell^s}^k=
y(\phi)^{\mp1}$. So $\pseinv(z\cup y)=\pm k\chi(\delta)=\pm\langle z(\delta),y(\phi)\rangle$, with a
sign fixed by the conventions.
\end{proof}
\noindent References: local reciprocity and the cup product with a Kummer class, and the tame
symbol: \cite[XIV, \S1, Proposition~3; \S2, Proposition~6; \S3, Proposition~8 and its Corollary]{SerreLF}; \cite[Ch.~I, Remark~2.5(b)]{MilneADT} (the tame symbol).

\subsection{Auxiliary primes: statement (A)}\label{pse:ss-aux}

\begin{proposition}[Statement (A)]\label{pse:A}
Let $N\ge1$, let $S$ be a finite set of finite places of $E$ containing those above $\ell$ and
stable under $c$, let $R$ be a finite set of finite places, let $Z$ be a finite set of finite places
outside $S$ and prime to $\ell$, and let $z$ lie in the normal subgroup of $\pgG$ generated by the
inertia groups $I_B$, $B$ over a place of $Z$. Let $X\subseteq\pgG$ be open with $z\in X$. Then
there are a place $P\notin R\cup\aubar R$ of $E$ of prime norm $q\equiv1\pmod{\ell^N}$ with
$P\ne\aubar P$, a place $B\mid P$, a normalised tame generator $\delta$ for $\ell^N$ at $B$, a
Frobenius element $\phi\in X$ at $B$, and a continuous homomorphism $\lambda\colon\pgG\to\Z/\ell^N$
such that:
\begin{enumerate}
\item[(i)] $\lambda(\pseArt(y))=\theta(y/\aubar y)$ for a character $\theta$ of the finite ideles
trivial on $E^\times$ and on $E_w^\times$ for $w\in S$; hence $\lambda$ is anticyclotomic
($\lambda(g^c)=-\lambda(g)$), unramified outside $P,\aubar P$, and zero on $D_{B'}$ for every
$B'$ over a place of $S$;
\item[(ii)] $\lambda(\delta)=1$ and $\lambda(\phi)=0$;
\item[(iii)] every Frobenius element at $B$ lies in $X$.
\end{enumerate}
\end{proposition}

\begin{proof}
\emph{Step 1: a set $T$.} Choose a finite $c$-stable set $T\supseteq S$ of finite places, disjoint
from $Z\cup\aubar Z$, such that the ideal classes of the places of $T$ generate the class group of
$E$ (possible, since every ideal class contains prime ideals outside any given finite set, by
Chebotarev in the Hilbert class field). Then every finite idele is the product of a principal idele
and of a \emph{$T$-idele}, an idele that is a unit at every place outside $T$. The group
$\mathcal O_{E,T}^\times$ of $T$-units is finitely generated (Dirichlet); let $u_1,\dots,u_k$
generate it.

\emph{Step 2: Chebotarev.} Let $F=E(\mu_{\ell^N},u_1^{1/\ell^N},\dots,u_k^{1/\ell^N})$, a finite
Galois extension, and $U_0=\operatorname{Gal}(\overline{\Q}/F)$. A place $w\in Z$ is prime to
$\ell$ and to every $u_i$ (since $w\notin T$), so $F/E$ is unramified at $w$ and every inertia group
$I_B$ over $w$ lies in the normal subgroup $U_0$; hence $z\in U_0$. Choose an open normal subgroup
$V\subseteq U_0$ of $\pgG$ with $zV\subseteq X$, and let $F'$ be its fixed field. By the Chebotarev
density theorem in $F'/E$, and since places of degree $>1$ over $\Q$ have density $0$, there are
infinitely many places $P$ of $E$ of prime norm, unramified in $F'$, with a place $B\mid P$ whose
Frobenius in $\operatorname{Gal}(F'/E)$ is the image of $z$. Exclude the finitely many $P$ in
$T\cup R\cup\aubar R$ and those ramified in $E/E^+$; a place of prime norm unramified over $E^+$ has
$P\ne\aubar P$ (otherwise $P\cap E^+$ is inert and $\auNm P$ is a square). Every Frobenius element
at $B$ lies in $zV$ (the inertia group $I_B$ lies in $V$, as $P$ is unramified in $F'$); this is
(iii), and every such element lies in $U_0$. Fix one, $\sigma$. As $\sigma$ fixes $\mu_{\ell^N}$,
$q=\auNm P\equiv1\pmod{\ell^N}$; as $\sigma$ fixes $u_i^{1/\ell^N}$ and acts as $x\mapsto x^q$ on the
residue field at $B$, $u_i^{(q-1)/\ell^N}\equiv1\pmod P$: the $T$-units are $\ell^N$-th powers in
$\auF_q^\times$.

\emph{Step 3: the character.} Let $\delta$ be a normalised tame generator for $\ell^N$ at $B$. By
local class field theory $\delta$ restricted to the maximal abelian extension of $E_P$ is the local
Artin symbol of a unit $u_0\in\auO_{E_P}^\times$; by the tame symbol (proof of
Lemma~\ref{pse:pairing}) $\bar u_0^{(q-1)/\ell^N}$ is a primitive $\ell^N$-th root of unity, so $\bar
u_0$ generates $\auF_q^\times/(\auF_q^\times)^{\ell^N}\cong\Z/\ell^N$. Let $\mathrm{lg}\colon
E_P^\times\to\Z/\ell^N$ send $x$ to the class of the residue of $x\varpi^{-v(x)}$ in this quotient,
identified with $\Z/\ell^N$ so that $\bar u_0\mapsto1$; it is continuous, $\mathrm{lg}
(u_0)=1$, and it kills the units $\equiv1\pmod P$ and the $T$-units (which are $\ell^N$-th powers
modulo $P$). Define $\theta$ on finite ideles by $\theta(xa)=\mathrm{lg}(a_P)$ for $x\in E^\times$ and
$a$ a $T$-idele; this is well defined because $E^\times\cap\{T\text{-ideles}\}$ is the group of
$T$-units, and it is continuous, trivial on $E^\times$, on $E_w^\times$ for $w\in T$, and on the
units at places $\ne P$. The character $y\mapsto\theta(y/\aubar y)$ is continuous, of finite order
and trivial on $E^\times$; by the reciprocity law of global class field theory there is a continuous
$\lambda_1\colon\pgG\to\Z/\ell^N$ with $\lambda_1(\pseArt(y))=\theta(y/\aubar y)$. For $y\in
E_w^\times$, $w\in S$, the idele $y/\aubar y$ is supported at $w,\aubar w\in T$, so $\lambda_1$ kills
$D_{B'}$ for $B'$ over $S$. It is unramified outside $P,\aubar P$ because $\theta$ kills the units at
all other places. Since $\pseArt(\aubar y)=c\,\pseArt(y)c$, $\lambda_1$ is anticyclotomic. For the idele
$y=u_0$ at $P$ (and $1$ elsewhere), $y/\aubar y$ has component $u_0$ at $P$ and a unit at $\aubar P\ne
P$, so $\lambda_1(\delta)=\mathrm{lg}(u_0)=1$.

\emph{Step 4: the Frobenius element.} Put $\phi=\sigma\delta^{-\lambda_1(\sigma)}$ (exponent read in
$[0,\ell^N)$). It is a Frobenius element at $B$, so $\phi\in X$ by (iii), and $\lambda_1(\phi)=0$.
Take $\lambda=\lambda_1$.
\end{proof}

\subsubsection*{Admissible places}
We now fix the meaning of ``a place that can be added to $J$'' in Propositions~\ref{pg:mirror},
\ref{pg:step}. Let $T_1,\dots,T_s$ be Selmer data (Definition~\ref{pg:def-datum}) and $M<N$. Let
$U_i\subseteq\pgG$ be the kernel of the action on $T_i/\ell^M$ and on $\mu_{\ell^N}$; it is open and
normal. Let $J$ be a set of tied primes (as in \S\ref{pg:ss-twisting}), with first prime $P_0$; for
$v\in J$ the data are a place $B_v\mid v$, the tame generator $\gamma_v$ (a normalised tame
generator for $\ell^N$ at $B_v$), the Frobenius element $\phi_v$ and $\lambda_v$, the $\lambda_v$
being as in Proposition~\ref{pse:A}(i) for a set $S$ containing the places above $\ell$ and the
places of $J\setminus\{P_0\}$ and their conjugates. A place $P$ with data $(B,\gamma_P,\phi_P,
\lambda_P)$ is \emph{admissible} for $J$ if: $P$, $\aubar P$ are outside $S_0\cup J\cup\aubar J$;
$P$ has prime norm and $P\ne\aubar P$; the data satisfy the conclusions of Proposition~\ref{pse:A}
for $N$ and the set $S=\{w\mid\ell\}\cup(J\setminus\{P_0\})\cup(\aubar J\setminus\{\aubar P_0\})$
(with $\delta=\gamma_P$); and $\lambda_{P_0}(\phi_P)=1$, $\lambda_v(\phi_P)=0$ for $v\in J\setminus
\{P_0\}$. It is \emph{admissible for $T_i$} if moreover $\phi_P,\phi_P^{c}\in U_i$. Then $J\cup
\{P\}$, with slope $0$ for the data for which $P$ is not admissible, and any slopes for the others,
is again a set of tied primes (the condition $\lambda_P(\phi_v)=0$ for $v\in J\setminus\{P_0\}$ holds
because $\lambda_P$ kills $D_{B_v}$).

\begin{lemma}[Admissible places near $\gamma_{P_0}H$]\label{pse:near}
Let $T$ be one of the data, $H\subseteq\pgG$ the subgroup generated by the commutators
$[g,\gamma_*]$ and $[g,\gamma_*^c]$, $g\in\pgG$, where $\gamma_*=\gamma_*(T)$. For every $h\in H$,
every finite set $R$ and every open $X\ni\gamma_{P_0}h$ there is a place $P\notin R\cup\aubar R$
admissible for $T$ with $\phi_P\in X$.
\end{lemma}
\begin{proof}
The element $\gamma_*$ lies in an inertia group at a place $v_*\in S_0$, $v_*\nmid\ell$; so
$[g,\gamma_*]=(g\gamma_*g^{-1})\gamma_*^{-1}$ lies in the normal subgroup generated by inertia at
$v_*$, and $[g,\gamma_*^c]$ in that generated by inertia at $\aubar v_*$. Hence $\gamma_{P_0}h$ lies
in the normal subgroup generated by inertia at the places of $Z=\{P_0,v_*,\aubar v_*\}$, which are
prime to $\ell$ and outside $S=\{w\mid\ell\}\cup(J\setminus\{P_0\})\cup(\aubar J\setminus\{\aubar
P_0\})$ (because $J$, $\aubar J$, $S_0$ are pairwise disjoint). The element $\gamma_{P_0}h$ lies in
the open group $U\cap U^c$, $U=U(T)$: $\gamma_{P_0}$ and $\gamma_{P_0}^c$ are inertial at places
outside $S_0$, where $T$ is unramified and $\mu_{\ell^N}$ too; and $h$, $h^c$ act trivially on $T$
and on $\mu_{\ell^N}$ because both $\gamma_*$ and $\gamma_*^{c}$ act on $T$ by scalars and trivially on
$\mu_{\ell^N}$ (so all commutators with them act trivially). The characters $\lambda_v$, $v\in J$, are
locally constant and kill $H$ (an abelian quotient kills commutators), so on a neighbourhood of
$\gamma_{P_0}h$ they take the values $\lambda_v(\gamma_{P_0})$, which is $1$ for $v=P_0$ and $0$ for
$v\ne P_0$ ($\lambda_v$ is unramified at $P_0$). Apply Proposition~\ref{pse:A} with this $S$ and
$Z$, with $R\cup S_0\cup J$ in place of $R$, and with $X$ replaced by the intersection of $X$, of
$U\cap U^c$ and of that neighbourhood. The place it gives is admissible for $T$.
\end{proof}

\subsection{The dual of a tied Selmer group}\label{pse:ss-mirror}

\subsubsection*{Which dual}
In Proposition~\ref{pg:mirror} and Lemma~\ref{pg:count}, $\pgSel^{\mathrm{tie},*}_i(J;a)$ is the
Selmer group in $H^1(E,(T_i/a)^*)$ with the local conditions orthogonal, under local Tate duality,
to those of the tied group. One may also work with a \emph{module of relations}: equivariant
homomorphisms $y$ on $U_i$ with values in $(T_i/a)^*$, unramified outside a finite set, satisfying
``reciprocity at split primes'' $\sum_jy(\phi_j)(z(\delta_j))=0$ for cocycles $z$ ramified at a
family of split primes. We keep the Tate dual of Lemma~\ref{pg:count}, for which the
Greenberg--Wiles formula is a published theorem, and show (Lemma~\ref{pse:reciprocity}) that
restriction to $U_i$ maps it injectively into the module of relations; the argument below is then
carried out in the module of relations.

Fix $i$, write $T=T_i$, $u=u_i$, $n=n_i$, $\gamma_*$, $\zeta_*$, $U=U_i$, $a=\varpi^k\mid\ell^M$,
$M<N$, $X=T/a$, and $H$ as in Lemma~\ref{pse:near}.

\begin{lemma}[Restriction to $U$]\label{pse:restrict}
Let $y$ be a cocycle representing a class of $\pgSel^{\mathrm{tie},*}(J;a)$. Then $y|_U$ is a
continuous homomorphism $U\to X^*$, independent of the representative, satisfying
$y(gwg^{-1})=g\cdot y(w)$ for $g\in\pgG$, $w\in U$. The class is zero if $y|_U=0$, and $y|_U=0$ if
$y|_H=0$. Moreover $y(\gamma_{P_0})=y(\gamma_{P_0}^c)=0$.
\end{lemma}
\begin{proof}
$U$ acts trivially on $X^*$ (it fixes $X$ and $\mu_{\ell^N}\supseteq\mu_a$), so $y|_U$ is a
homomorphism, coboundaries vanish on $U$, and the cocycle relation gives the equivariance. The image
of $\gamma_*$ in $\pgG/U\subseteq\operatorname{Aut}(T/\ell^M)\times(\Z/\ell^N)^\times$ is central (a
scalar, and $\pgcyc(\gamma_*)=1$ because $\gamma_*$ is inertial at a place prime to $\ell$); it acts on
$X^*$ by $\zeta_*^{-1}$, and $\zeta_*^{-1}-1$ is a unit. By Sah's lemma $H^1(\pgG/U,X^*)=0$, so by
inflation--restriction the class is determined by $y|_U$. If $y|_H=0$, then for $w\in U$,
$y(w)=y([w,\gamma_*])+y(\gamma_*w\gamma_*^{-1})=\gamma_*\cdot y(w)=\zeta_*^{-1}y(w)$ (as $[w,
\gamma_*]\gamma_*w\gamma_*^{-1}=w$), so $(\zeta_*^{-1}-1)y(w)=0$ and $y(w)=0$. Finally $P_0$ has
slope $0$, so the local condition of the tied group at $P_0$ and $\aubar P_0$ is the unramified one;
its orthogonal complement is the unramified condition (local duality at a place prime to $\ell$
where $X$ is unramified); a cocycle unramified at $B_{P_0}$ takes values in $(\gamma-1)X^*=0$ on
$\gamma\in I_{B_{P_0}}$, and likewise at $cB_{P_0}$, which contains $\gamma_{P_0}^c$.
\end{proof}

\begin{lemma}[Reciprocity at an admissible place]\label{pse:reciprocity}
Let $y$ be as in Lemma~\ref{pse:restrict}, and let $P$ be admissible for $T$, with $P,\aubar P$
outside the finite set where $y$ is ramified. Let $z'$ be a cocycle $\pgG\to X$ with the local
conditions and ties of $\pgSel^{\mathrm{tie}}(J;a)$ that may ramify at $P,\aubar P$ (no condition
there). Then
\[
\bigl\langle z'(\gamma_P),y(\phi_P)\bigr\rangle+\bigl\langle z'(\gamma_P^{-c}),y(\phi_P^{c})\bigr
\rangle=0 .
\]
\end{lemma}
\begin{proof}
The cup product $z'\cup y\in H^2(E,\Q/\Z(1))$ has $\sum_w\pseinv_w(z'\cup y)=0$ (reciprocity law of
the Brauer group; \cite[Theorem~8.1.17]{NSW}). Here $S\subseteq S_0$ denotes the ramification set of the datum (Definition~\ref{pg:def-datum}; it contains the places above $\ell$), not the admissibility set above. The places are: those of $S\cup J\cup\aubar J$, where
$z'$ satisfies the condition $L_w$ and $y$ lies in $L_w^\perp$, so the term vanishes; the
archimedean places, which are complex; the places outside $S\cup J\cup\aubar J\cup\{P,\aubar P\}$
where both are unramified and $w\nmid\ell$, where the term vanishes; and $P$, $\aubar P$. At $B$
over $P$ the group $D_B$ acts trivially on $X$ (as $\phi_P\in U$ and $I_B\subseteq U$, $P\notin
S_0$) and $\auNm P\equiv1\pmod{\ell^N}$; by Lemma~\ref{pse:pairing}, $\pseinv_P(z'\cup y)=\epsilon
\langle z'(\delta),y(\phi_P)\rangle$ for a normalised tame generator $\delta$ for $\ell^N$ at $B$ (it is also
normalised for $\ell^M$, by the compatibility of the $\xi_m$). At
$cB$, over $\aubar P$, the element $\phi_P^c$ is a Frobenius element and $\delta^{-c}$ a normalised
tame generator (Lemma~\ref{pse:conjgen}), so $\pseinv_{\aubar P}(z'\cup y)=\epsilon\langle z'(
\delta^{-c}),y(\phi_P^c)\rangle$. Hence $\langle z'(\delta),y(\phi_P)\rangle+\langle z'(\delta^{-c}),
y(\phi_P^c)\rangle=0$. The restrictions of $z'$ to $I_B$ and $I_{cB}$ are homomorphisms ($T$ is
unramified at $P,\aubar P$); by Lemma~\ref{pse:kpower}, applied to $\gamma=\gamma_P$ and $\delta$, $z'(\gamma_P)=k z'(\delta)$ and $z'(\gamma_P^{-c})=kz'(\delta^{-c})$
for one integer $k$. Multiplying by $k$ gives the claim.
\end{proof}

\begin{lemma}[Sah's lemma for the tied group]\label{pse:span}
If $z\in\pgSel^{\mathrm{tie}}(J;a)$ has a value with a unit coordinate, then $z(H)$ spans $X$
over $\auO$.
\end{lemma}
\begin{proof}
As $z(\gamma_*)=0$ and $u(\gamma_*)=\zeta_*$, the cocycle rule gives $z([g,\gamma_*])=(1-\zeta_*)
z(g)$ for all $g$. Since $1-\zeta_*$ is a unit, $z(\pgG)\subseteq\auO z(H)$. The span $Y$ of
$z(\pgG)$ is $\pgG$-stable ($u(g)z(h)=z(gh)-z(g)$), hence stable under the $\auO$-span $M_n(\auO)$ of
the $u(g)$. It contains a vector $x$ not in $\varpi X$; then $M_n(\auO)x=X$. So $\auO z(H)=X$.
\end{proof}

\begin{proposition}[The dual of a tied Selmer group]
\label{pse:mirror}
Let $J$ be a set of tied primes, $a=\varpi^k\mid\ell^M$ with $k\ge1$, $M<N$, and suppose that
$\pgSel^{\mathrm{tie}}_i(J;a)$ contains a cocycle $z$ one of whose values has a unit coordinate.
Suppose that $z$ is \emph{balanced} at every place $P$ admissible for $T_i$: there is a cocycle
$z'$ with the local conditions and ties of $\pgSel^{\mathrm{tie}}_i(J;a)$, allowed to ramify at
$P,\aubar P$, with $z'(\gamma_P)=z(\phi_P)$ and $z'(\gamma_P^{-c})=z(\phi_P^{-c})$. Then
$\pglen\pgSel^{\mathrm{tie},*}_i(J;a)\le\mathrm{gl}(T_i)$.
\end{proposition}

This is the precise form of Proposition~\ref{pg:mirror}: the dual is the Tate dual,
identified through Lemmas~\ref{pse:restrict}, \ref{pse:reciprocity}; ``a place that can be added''
means a place admissible for $T_i$; balance is needed for the one element $z$ only; and $k\ge1$.

\begin{proof}
Let $y$ represent a class of the dual group, and let $R$ be the finite set of places where $y$ is
ramified, enlarged by $S_0\cup J\cup\aubar J$. For $t\in\auO$ and $h,h'\in H$ put
\[
\Psi_t(h,h')=\bigl\langle tz(h),y(h')\bigr\rangle-\bigl\langle tz(h'^c),y(h^c)\bigr\rangle\in\Q/\Z .
\]
Recall $H\subseteq U\cap U^c$ (proof of Lemma~\ref{pse:near}) and $H^c=H$; on $U$, $z$ and $y$ are
homomorphisms. So $\Psi_t$ is biadditive on $H\times H$ and additive in $t$.

\emph{Step 1: $\Psi_t(h,h)=0$.} The tie with slope $0$ at $P_0$ gives $z(\gamma_{P_0})=
z(\gamma_{P_0}^{-c})=0$, hence $z(\gamma_{P_0}^c)=0$; and $y(\gamma_{P_0})=y(\gamma_{P_0}^c)=0$
(Lemma~\ref{pse:restrict}). As $\gamma_{P_0},\gamma_{P_0}^c\in U$, for $g=\gamma_{P_0}h$ we get
$z(g)=z(h)$, $z(g^c)=z(h^c)$, $y(g)=y(h)$, $y(g^c)=y(h^c)$, so $\Psi_t(h,h)=\Psi_t(g,g)$, where we
extend the formula for $\Psi_t(g,g)$ to all $g\in U\cap U^c$. The map $g\mapsto\Psi_t(g,g)$ is
locally constant. By Lemma~\ref{pse:near} there is an admissible $P\notin R\cup\aubar R$ with
$\phi=\phi_P$ so close to $\gamma_{P_0}h$ that $\Psi_t(\phi,\phi)=\Psi_t(\gamma_{P_0}h,
\gamma_{P_0}h)$ and $\phi,\phi^c\in U$. Let $z'$ be the cocycle given by the balance of $z$ at $P$.
Then $z'(\gamma_P)=z(\phi)$ and $z'(\gamma_P^{-c})=z(\phi^{-c})=-u(\phi^{-c})z(\phi^c)=-z(\phi^c)$. So
\[
\Psi_t(\phi,\phi)=\bigl\langle tz'(\gamma_P),y(\phi)\bigr\rangle+\bigl\langle tz'(\gamma_P^{-c}),
y(\phi^c)\bigr\rangle=0
\]
by Lemma~\ref{pse:reciprocity} applied to $tz'$.

\emph{Step 2: $\Psi=0$.} By biadditivity and Step 1 (applied to $hh'$, $h$, $h'$),
$\Psi_t(h,h')+\Psi_t(h',h)=0$. For $h\in H$, $\gamma_*h\gamma_*^{-1}\in H$ ($H$ is normal) and, since
$z(\gamma_*)=0$ and $h$ acts trivially, $z(\gamma_*h\gamma_*^{-1})=\zeta_*z(h)$; by equivariance and
$u(\gamma_*^{-c})=\zeta_*$, $\pgcyc(\gamma_*^c)=1$, $y((\gamma_*h\gamma_*^{-1})^c)=y(h^c)\circ\zeta_*$.
Hence $\Psi_t(\gamma_*h\gamma_*^{-1},h')=\Psi_{\zeta_*t}(h,h')$ and, in the same way,
$\Psi_{\zeta_*t}(h',\gamma_*h\gamma_*^{-1})=\Psi_t(h',h)$. Antisymmetry for the pair
$(\gamma_*h\gamma_*^{-1},h')$ with parameter $\zeta_*s$ gives $\Psi_{\zeta_*^2s}(h,h')+\Psi_s(h',h)=0$;
subtracting antisymmetry for $(h,h')$ with parameter $s$ gives $\Psi_{(\zeta_*^2-1)s}(h,h')=0$. As
$\zeta_*^2-1$ is a unit, $\Psi_t=0$ for all $t$.

\emph{Step 3: the map $\eta$.} Suppose $\sum_ja_jz(h_j)=0$ with $a_j\in\auO$, $h_j\in H$. For every
$h'\in H$ and $t\in\auO$, by $\Psi=0$,
$\sum_ja_j\,y(h_j^c)(tz(h'^c))=\sum_j\langle ta_jz(h_j),y(h')\rangle=0$; as the $z(h'^c)$, $h'\in H$,
span $X$ (Lemma~\ref{pse:span} and $H^c=H$), $\sum_ja_jy(h_j^c)=0$. Hence there is a unique
$\auO$-linear $\eta_y\colon X\to X^*$ with $\eta_y(z(h))=y(h^c)$ for $h\in H$. It is
$\pgG$-equivariant for the action on $X^*$ through $g\mapsto g^c$: for $h\in H$ and $g\in\pgG$,
$u(g)z(h)=z(ghg^{-1})$ (cocycle rule, $h$ and $ghg^{-1}$ in $U$), and $\eta_y(z(ghg^{-1}))=
y(g^ch^cg^{-c})=g^c\cdot y(h^c)$. So $\eta_y\in\pgHom_{\pgG}(X,X^{*c})$.

\emph{Step 4.} The map $y\mapsto\eta_y$ is $\auO$-linear, and injective: if $\eta_y=0$ then $y$
vanishes on $H^c=H$, hence the class of $y$ is $0$ by Lemma~\ref{pse:restrict}. So
$\pglen\pgSel^{\mathrm{tie},*}(J;a)\le\pglen\pgHom_{\pgG}(X,X^{*c})\le\mathrm{gl}(T)$.
\end{proof}

\subsection{The step of the descent}\label{pse:ss-step}

\begin{lemma}[Tilting]\label{pse:tilting}
Let $\auO$ be a discrete valuation ring with residue field $k$, $C$ an $\auO$-module of finite
length, $F=(\auO/I)^m$ with $I\ne0$, and $e_\gamma,e_\phi\colon C\to F$ linear. There is a set
$\mathrm{Bad}\subseteq k$ with at most $m$ elements such that for every $\alpha\in\auO$ with residue
outside $\mathrm{Bad}$,
\[
\pglen\bigl(\ker(\alpha e_\gamma-e_\phi)\bigr)\le\pglen\bigl(\ker e_\gamma\cap e_\phi^{-1}(e_\gamma
C)\bigr).
\]
\end{lemma}
\begin{proof}
Let $C'=\sum_{\alpha\in\auO}\ker(\alpha e_\gamma-e_\phi)\subseteq C$. If $\alpha e_\gamma(x)=
e_\phi(x)$ then $e_\phi(x)=e_\gamma(\alpha x)$ with $\alpha x\in C'$; so $e_\phi(C')\subseteq
e_\gamma(C')$. Let $\Lambda\subseteq\auO^m$ be the preimage of $e_\gamma(C')$; it contains $I\auO^m$,
so it is free of rank $m$; let $\pi\colon\Lambda\to e_\gamma(C')$ be the reduction. As $\Lambda$ is
free there are $\sigma\colon\Lambda\to C'$ with $e_\gamma\sigma=\pi$, and then (since
$e_\phi\sigma(\Lambda)\subseteq e_\gamma(C')=\pi(\Lambda)$) $\mathcal U\colon\Lambda\to\Lambda$ with
$\pi\,\mathcal U=e_\phi\sigma$. Let $\mathrm{Bad}$ be the set of roots in $k$ of the reduction of the
characteristic polynomial of $\mathcal U$. For $\alpha$ with residue outside $\mathrm{Bad}$,
$\det(\alpha-\mathcal U)$ is a unit, $\alpha-\mathcal U$ is bijective, and $(\alpha e_\gamma-e_\phi)
\sigma=\pi(\alpha-\mathcal U)$ maps onto $e_\gamma(C')$. Hence the image of $\alpha e_\gamma-e_\phi$
on $C'$ is at least as long as that of $e_\gamma$, and since
$\pglen C'=\pglen f(C')+\pglen(C'\cap\ker f)$ for every $f$, $\pglen(C'\cap\ker(\alpha e_\gamma-
e_\phi))\le\pglen(C'\cap\ker e_\gamma)$. Finally $\ker(\alpha e_\gamma-e_\phi)\subseteq C'$, and
$C'\cap\ker e_\gamma\subseteq\ker e_\gamma\cap e_\phi^{-1}(e_\gamma C)$ because $e_\phi(C')\subseteq
e_\gamma(C')$.
\end{proof}

\begin{proposition}[One step of the descent; Proposition~\ref{pg:step}]\label{pse:step}
Let $J$ be a set of tied primes for data $T_1,\dots,T_s$ over $\auO$ with uniformiser $\varpi$,
$\varpi^e\mid\ell^M$, $M<N$, and suppose that the residue field of $\auO$ has more than $2n_ie+1$
elements for all $i$. If for some $i$, some $1\le k\le e$ and some place $P$ admissible for $T_i$, an
element of $\pgSel^{\mathrm{tie}}_i(J;\varpi^k)$ is not balanced at $P$, then $P$ can be added to $J$
with a unit slope for $T_i$ and slope $0$ for the other data so that
$\Sigma(J)=\sum_{i'}\sum_{k'=1}^e\pglen\pgSel_{i'}^{\mathrm{tie}}(J;\varpi^{k'})$ decreases.
\end{proposition}
\begin{proof}
\emph{The other data.} With slope $0$ at $P$, a cocycle of the new group for $T_{i'}$, $i'\ne i$,
vanishes at $\gamma_P$ and $\gamma_P^{-c}$; its restrictions to $I_{B}$ and $I_{cB}$ are
homomorphisms killing the generators of their maximal $\ell$-power quotients (Lemma~\ref{pse:kpower}
with $\delta=\gamma_P$, and Lemma~\ref{pse:conjgen}), so it is unramified at $P,\aubar P$ and lies in
the old group. Hence these terms do not increase.

\emph{The datum $T_i$.} For $1\le k'\le e$ let $C_{k'}$ be the module of cocycles modulo
$\varpi^{k'}$ with the local conditions and ties of $\pgSel^{\mathrm{tie}}_i(J;\varpi^{k'})$, allowed
to ramify at $P,\aubar P$; it is finite. Let $e_\gamma,e_\phi\colon C_{k'}\to(T_i/\varpi^{k'})^2$ be
evaluation at $(\gamma_P,\gamma_P^{-c})$ and at $(\phi_P,\phi_P^{-c})$. As just shown, $C_{k'}\cap\ker
e_\gamma=\pgSel_i^{\mathrm{tie}}(J;\varpi^{k'})$, and $\ker e_\gamma\cap e_\phi^{-1}(e_\gamma C_{k'})$
is its submodule $B_{k'}$ of elements balanced at $P$. With slope $s=\alpha^{-1}$, $\alpha$ a unit,
the new group is $C_{k'}\cap\ker(\alpha e_\gamma-e_\phi)$ (the tie $z(\gamma_P)=sz(\phi_P)$, and the
same at $\aubar P$). By Lemma~\ref{pse:tilting} ($m=2n_i$), for $\alpha$ with residue outside a set
of at most $2n_i$ elements its length is at most $\pglen B_{k'}$. The union of these sets over
$k'\le e$, together with $0$, has at most $2n_ie+1$ elements, so some unit $\alpha$ avoids all. For
it, the new group at level $k'$ has length at most $\pglen B_{k'}\le\pglen\pgSel_i^{\mathrm{tie}}
(J;\varpi^{k'})$, and at level $k$ strictly less, since $B_k$ is a proper submodule of a module of
finite length. The enlarged set is a set of tied primes (\S\ref{pse:ss-aux}, admissible places),
the condition $u_i(\phi_P)\equiv u_i(\phi_P^{-c})\equiv1\pmod{\ell^M}$ for the nonzero slope holding
because $\phi_P,\phi_P^c\in U_i$.
\end{proof}

\begin{remark}\label{pse:rem-poly}
In the proof of Proposition~\ref{pg:poly}, ``balanced at every admissible $P$'' is to be read with
the present meaning, and the application of Proposition~\ref{pg:mirror} is replaced by
Proposition~\ref{pse:mirror}: at the end of the descent every element of every tied group at every
level $\varpi^{k'}$, $k'\le e$, is balanced at every place admissible for its datum, in particular the
element with a unit coordinate. The first prime $P_0$ and its character are given by
Proposition~\ref{pse:A} with $S=\{w\mid\ell\}$, $Z=\{w_0\}$, $z=g_0$, $R=S_0$, and $X$ the open set
of $g$ acting as $g_0$ on all $T_i/\ell^M$, with $g^{-c}$ acting as $g_0^{-c}$, and (for
\S\ref{pse:ss-stick}) with $\vartheta_Y(g)\equiv\vartheta_Y(g_0)\pmod{\ell^M}$ for the finitely many
systems $Y$.
\end{remark}

\subsection{The systems of the \texorpdfstring{$L$}{L}-values}\label{pse:ss-stick}

\subsubsection*{The input from theta lifts}
A \emph{tame character of modulus $Q$} is a character $\vartheta$ of $\auU(1)(E^+)\backslash
\auU(1)(\auA_{E^+})$ of weight $0$, trivial at the infinite places, tamely ramified exactly on a
subset of $Q$ and unramified elsewhere, where $Q$ is a finite set of places of $E$ prime to $\ell p_0p_1$ and
to the level, i.e.\ outside $S_0$, with $Q\cap\aubar Q=\emptyset$; it is a twist of modulus $Q$ when it is ramified at every
place of $Q$. For a ramified critical pair $(\nu,z)$ let $r(\vartheta)=|L(z,(\nu\vartheta)^{-1})|/
|L(z,\nu^{-1})|$ for twists $\vartheta$, with the conventions on omitted Euler factors of
Theorem~\ref{pg:main}(b).

\noindent\emph{Statement (IM) (integral measures).} For every ramified critical pair $(\nu,z)$ there
are a finite set $I\subseteq\Z_{\ge0}$ with at most two elements ($I=\{z\}$ if $z\ge2$, $I=\{0,1\}$ if
$z=1$), an integer $m\in\{0,1\}$ ($m=0$ if $z\ge2$, $m=1$ if $z=1$) and a function $F$ on tame
characters such that: (a) $F(1)\ne0$; (b) for every modulus $Q$ there are finitely many
$x_a\in\auU(1)(\auA_{E^+,f})$ and $c_a\in\C$, with $c_a$ integral at every prime not dividing
$\auNm Q$ (under every field embedding into $\C_p$), such that $F(\vartheta)=\sum_ac_a\vartheta(x_a)$
for every tame $\vartheta$ of modulus $Q$; (c) if $v,\aubar v\notin\mathsf M\cup Q\cup\aubar Q$ ($\mathsf M$ the set of chosen places) and
$\vartheta'$ is $\vartheta$ regarded as of modulus $Q\cup\{v\}$, then $F(\vartheta')=F(\vartheta)
\prod_{e\in I}\prod_\pm\bigl(1-\alpha_v^{\pm1}(\auNm v)^{-e}\bigr)$, $\alpha_v$ the value at a
uniformiser of the base change to $E$ of $\nu\vartheta$ at $v$; (d) for every twist $\vartheta$ of
modulus $Q$, $F(\vartheta)=F(1)\,(\auNm Q)^m\,r(\vartheta)^{m+1}$. For $z\ge2$, $F$ is an integral
multiple of the squared norm of the theta lift of $\nu\vartheta$ from $\auU(1)$ to $\auU(2z)$; for
$z=1$, a theta period on $\auU(2)$ times its conjugate.

\smallskip\noindent (IM) is proved in Proposition~\ref{pex:IM} for $z\ge2$ (at every prime $\ell\nmid
\auNm Q$; $\ell=d$ by Lemma~\ref{pex:dint}) and in Proposition~\ref{pex:IM1} for $z=1$ (from the
exchange identity, Theorem~\ref{pcx:X}). The ratio $\rho_p(\vartheta)$ there is $r(\vartheta)$, with
the conventions on omitted Euler factors of Theorem~\ref{pg:main}(b).\par\smallskip

\subsubsection*{From (IM) to a system}
Let $\pseCl_Q$ be the group of finite ideles $x$ with $x\aubar x=1$, modulo $\auU(1)(E^+)$ and the
units $\equiv1$ modulo the places of $Q$. By Hilbert's Theorem~90 for ideles every class is
$y/\aubar y$ for a finite idele $y$. A continuous character $\psi\colon\pgG\to\overline{\Q}_\ell^\times$
is a \emph{character for $Q$} if $\psi(\pseArt(y))=\eta(y/\aubar y)$ for a character $\eta$ of
$\pseCl_Q$ trivial at the places above $\ell$; then $\psi$ is read on $\pseCl_Q$ by
$\psi(y/\aubar y):=\psi(\pseArt(y))$, and extended linearly to the group ring $\pgZb[\pseCl_Q]$. Every
character for $Q$ of finite order is $\psi_\vartheta$ for a tame character $\vartheta$ of modulus $Q$
(\S\ref{pg:ss-reduction}), with $\psi_\vartheta(y/\aubar y)=
\iota\vartheta^{-1}(y/\aubar y)$.

Fix a ramified critical pair $(\nu,z)$, a sign $\delta=\pm1$, and $F$, $I$, $m$ from (IM). Let
$\vartheta_Y\colon\pgG\to\pgZb^\times$ be the continuous character attached through $\iota$ and class
field theory to the algebraic Hecke character $y\mapsto\nu^{-\delta}(y/\aubar y)$ of $E$. It is
unramified outside $S_0$, and its value at a Frobenius element $\phi$ at a place $v\notin S_0$ times
$\psi_\vartheta(\phi)$ is $\iota$ of $\alpha_v^{\pm\delta}$ for the pair $\nu\vartheta^{\delta}$ (one sign for
all $v$), by the definitions, because $\nu^{-\delta}\vartheta^{-1}=(\nu\vartheta^{\delta})^{-\delta}$. At $g_0=\gamma_{w_0}$, $\vartheta_Y(g_0)$ is the value of the tame part
of $\nu^{\mp\delta}$ at a generator of $\auF_{w_0}^\times$, a root of unity $\ne1$ of order dividing
$\auNm w_0-1$, because the pair is ramified at the chosen place $w_0$. Hence
\begin{equation}\label{pse:eq-g0}
\pgv(\vartheta_Y(g_0)-1)\le\pgv(\auNm w_0-1)
\end{equation}
($\prod_{\xi^n=1,\xi\ne1}(1-\xi)=n$ and each factor is integral).

For $Q$ auxiliary (places of prime norm, $Q$, $\aubar Q$, $S_0$ pairwise disjoint) take $x_a$,
$c_a$ from (IM)(b); $\iota(c_a)\in\pgZb$ since $\ell\nmid\auNm Q$. Put $\Theta_Q=\sum_a\iota(c_a)
[x_a^{-\delta}]\in\pgZb[\pseCl_Q]$, and $\Theta_\emptyset$ likewise. Then, for every character
$\psi=\psi_\vartheta$ for $Q$ of finite order,
\[
\psi(\Theta_Q)=\sum_a\iota(c_a)\,\iota\vartheta^{\delta}(x_a)=\iota F(\vartheta^\delta),
\]
and, by (IM)(c), for $v\notin Q$ with $Q\cup\{v\}$ auxiliary (so that $v,\aubar v\notin\mathsf M\cup
Q\cup\aubar Q$, the set $\mathsf M$ of chosen places being contained in $S_0$) and $\phi$ a Frobenius element at a
place over $v$,
\begin{equation}\label{pse:eq-norm}
\psi(\Theta_{Q\cup\{v\}})=\psi(\Theta_Q)\cdot\prod_{e\in I}\prod_\pm\bigl(1-(\vartheta_Y(\phi)
\psi(\phi))^{\pm1}(\auNm v)^{-e}\bigr),
\end{equation}
the product being symmetric under inversion of $\vartheta_Y(\phi)\psi(\phi)$, so that the sign
conventions do not matter. Let $v_0=\pgv(\iota F(1))=\pgv(1(\Theta_\emptyset))$, finite by (IM)(a),
and $\beta_Y=v_0+2\#I\,\pgv(\vartheta_Y(g_0)-1)\le v_0+4\pgv(\auNm w_0-1)$. By (IM)(d), for a twist
$\vartheta$ of auxiliary modulus $Q$ (so that $\auNm Q$ is an $\ell$-adic unit),
\begin{equation}\label{pse:eq-ratio}
\pgv\bigl(\psi_\vartheta(\Theta_Q)\bigr)=v_0+(m+1)\,\pgv\bigl(\iota r(\vartheta^\delta)\bigr),
\quad\text{hence}\quad \pgv\bigl(\iota r(\vartheta^{\delta})\bigr)\le\pgv\bigl(\psi_\vartheta(\Theta_Q)
\bigr)-v_0\ \text{ if }\pgv(\iota r(\vartheta^\delta))\ge0 .
\end{equation}

\subsubsection*{The polynomial and its line}
Let $J$ be a set of tied primes with first prime $P_0$ chosen as in Remark~\ref{pse:rem-poly}, so
that $\vartheta_Y(\phi_{P_0})\equiv\vartheta_Y(g_0)\pmod{\ell^M}$, and let $\lambda_v$, $v\in J$, be
its characters. Put
\[
P_Y=\sum_a\iota(c_a)\prod_{v\in J}(1+X_v)^{\lambda_v(x_a^{-\delta})}\in\pgZb[X_v:v\in J],
\]
the exponents read in $[0,\ell^N)$ and $\lambda_v$ read on $\pseCl_J$ as above (each $\lambda_v$ is
of the form of Proposition~\ref{pse:A}(i), hence a character for $J$). For $\ell^N$-th roots of unity
$\zeta_v$ and $\psi=\prod_v\zeta_v^{\lambda_v}$ one has $P_Y((\zeta_v-1)_v)=\psi(\Theta_J)$.

\begin{lemma}[Euler factors as polynomials]\label{pse:eulerpoly}
Let $\alpha$ and $q$ be $\ell$-adic units, $I$ a finite set, $\rho\ge0$. There is $R\in\pgZb[X]$
with Gauss valuation $\pgv_\rho(R)\le2\#I\rho$ and $R(\zeta-1)=\zeta^{\#I}\prod_{e\in I}\prod_\pm
(1-(\alpha\zeta)^{\pm1}q^{-e})$ for every $\zeta\in\pgZb^\times$.
\end{lemma}
\begin{proof}
For one $e$, with $b=q^{-e}$: $\zeta(1-\alpha b\zeta)(1-\alpha^{-1}b\zeta^{-1})=(1-\alpha b(1+X))
(1+X-\alpha^{-1}b)$ at $X=\zeta-1$. Each factor is linear with unit coefficient of $X$, so has
$\pgv_\rho\le\rho$. Take the product over $e$; $\pgv_\rho$ is additive.
\end{proof}

\begin{lemma}\label{pse:line}
Let $\rho\le1$ be rational and suppose $\beta_Y+1<M\le N\rho$, $s_0\rho<1$ and $4\#J\le s_0$. Then there is a
point $x=(x_v)$ with $x_{P_0}=y$, $\pgv(y)=\rho$, and $x_v=0$ for $v\ne P_0$, such that
$\pgv(P_Y(x))\le\beta_Y+s_0\rho$.
\end{lemma}
\begin{proof}
Write $h=\lambda_{P_0}$, $L(X)=P_Y(X,0,\dots,0)=\sum_a\iota(c_a)(1+X)^{h(x_a^{-\delta})}$ and
$R_0(X)=\sum_a\iota(c_a)(1+X)^{h(x_a^{-\delta})}$ computed with the elements of $\Theta_{\{P_0\}}$
instead of $\Theta_J$. For an $\ell^N$-th root of unity $\zeta$ the character $\psi=\zeta^{h}$ is a
character for $\{P_0\}$ (hence for $J$), with $\psi(\phi_v)=\zeta$ for $v\in J\setminus\{P_0\}$
(tied primes: $\lambda_{P_0}(\phi_v)=1$). By \eqref{pse:eq-norm} applied to the places of
$J\setminus\{P_0\}$ one after another, $L(\zeta-1)=R_0(\zeta-1)\prod_{v\ne P_0}\prod_{e,\pm}(1-
(\vartheta_Y(\phi_v)\zeta)^{\pm1}(\auNm v)^{-e})$. Let $R_v$ be the polynomial of
Lemma~\ref{pse:eulerpoly} for $\alpha=\vartheta_Y(\phi_v)$, $q=\auNm v$, and $R_1=\prod_{v\ne P_0}
R_v$, so $\pgv_\rho(R_1)\le2\#I\#J\rho\le s_0\rho$. Then $(1+X)^{j}L-R_0R_1$, $j=\#I(\#J-1)$,
vanishes at the $\ell^N$ distinct points $\zeta-1$; as $(1+X)^{\ell^N}-1$ is monic with these simple
roots, it divides $(1+X)^jL-R_0R_1$ in $\pgZb[X]$.

$\pgv_\rho(R_0)\le\pgv(R_0(0))=\pgv(1(\Theta_{\{P_0\}}))$, and by \eqref{pse:eq-norm} with $Q=
\emptyset$, $\psi=1$, $1(\Theta_{\{P_0\}})=1(\Theta_\emptyset)\prod_{e,\pm}(1-\vartheta_Y(\phi_{P_0})^{
\pm1}(\auNm P_0)^{-e})$. Since $\auNm P_0\equiv1\pmod{\ell^N}$, $N\ge M$, and $\vartheta_Y(
\phi_{P_0})\equiv\vartheta_Y(g_0)\pmod{\ell^M}$, each factor is $\equiv1-\vartheta_Y(g_0)^{\pm1}
\pmod{\ell^M}$, whose valuation $\pgv(\vartheta_Y(g_0)-1)\le\beta_Y<M$; so each factor has that
valuation and $\pgv_\rho(R_0)\le\beta_Y$. Thus $\pgv_\rho(R_0R_1)\le\beta_Y+s_0\rho<\beta_Y+1<
M\le N\rho$.

Choose $y\in\pgZb$ with $\pgv(y)=\rho$ and $\pgv((R_0R_1)(y))=\pgv_\rho(R_0R_1)$: write $y=\pi w$
with $\pgv(\pi)=\rho$; then $\pgv((R_0R_1)(\pi w))=\pgv_\rho(R_0R_1)$ unless the residue of the unit
$w$ is a root of a nonzero polynomial over the (infinite) residue field. Now $\pgv((1+y)^{\ell^N}-1)
\ge N\rho$ (each term $\binom{\ell^N}{i}y^i$, $i=\ell^bi'$, has valuation $\ge N-b+\ell^b\rho\ge
N\rho$ as $\rho\le1$), which exceeds $\pgv((R_0R_1)(y))$; by the divisibility,
$\pgv((1+y)^jL(y))=\pgv((R_0R_1)(y))$, and $1+y$ is a unit. So $\pgv(P_Y(x))=\pgv(L(y))\le
\beta_Y+s_0\rho$.
\end{proof}

\begin{proposition}[Valuations of the theta polynomials]\label{pse:stick}
Assume (IM). In the situation of Proposition~\ref{pg:poly}, let finitely many pairs of a ramified
critical pair $(\nu,z)$ and a sign $\delta$ be given, with their systems $Y$, and suppose moreover
$M>\beta_Y+1$ for each. Then, with $s_0$ replaced by $\max(s_0,4(1+b_0))$, for every rational
$0<\rho<1/s_0$ and every large $N$ the set $J$ of Proposition~\ref{pg:poly} can be chosen so that
moreover, for each $Y$, the polynomial $P_Y\in\pgZb[X_v:v\in J]$ and a point $x_Y$ with
coordinates of valuation at least $\rho$ satisfy
\[
\pgv(P_Y(x_Y))\le v_0(Y)+4\pgv(\auNm w_0-1)+s_0\rho ,
\]
and, for all $\ell^N$-th roots of unity $\zeta_v\ne1$ with $\psi=\prod_v\zeta_v^{\lambda_v}=
\psi_\vartheta$,
\[
\pgv\bigl(\iota\,r(\vartheta^{\delta})\bigr)\le\max\bigl(0,\ \pgv(P_Y((\zeta_v-1)_v))-v_0(Y)\bigr).
\]
\end{proposition}
\begin{proof}
Choose $P_0$ as in Remark~\ref{pse:rem-poly}; the descent of Proposition~\ref{pg:poly} does not
change $P_0$, and adds at most $b_0$ primes, so $\#J\le1+b_0$ and $4\#J\le s_0$. The first bound is
Lemma~\ref{pse:line} with $\beta_Y\le v_0+4\pgv(\auNm w_0-1)$ (by \eqref{pse:eq-g0}; $M\le N\rho$
for large $N$). For the second: $\vartheta$ is a twist of modulus $J$ (each $\zeta_v\ne1$ and
$\psi(\gamma_v)=\zeta_v$), $J$ is auxiliary, and $P_Y((\zeta_v-1))=\psi(\Theta_J)$; apply
\eqref{pse:eq-ratio}.
\end{proof}

\subsection{Constants in the proof of the twisting theorem}\label{pse:ss-end}

\begin{remark}\label{pse:rem-constants}
The proof of Theorem~\ref{pg:II} in \S\ref{pg:ss-twisting} uses the results of this section with the following choices. (1) $M$ is chosen with
$M>\max(c,\max_Y\beta_Y)+1$, where $Y$ runs over the systems of the finitely many pairs and signs;
the residue field of $\auO$ is enlarged accordingly. (2) $s_0$ is as in Proposition~\ref{pse:stick}.
(3) In Proposition~\ref{pg:roots} the bound is $\beta=s(c+1)+\sum_Y(v_0(Y)+4\pgv(\auNm w_0-1)+1)$,
which is finite; $h_0$ depends on it, which is allowed. (4) In (c): by Proposition~\ref{pg:roots},
$\pgv(P_Y((\zeta_v^j-1)_v))\le\pgv(P_Y(x_Y))+\varepsilon/2\le v_0(Y)+4\pgv(\auNm w_0-1)+\varepsilon$,
so by Proposition~\ref{pse:stick} the ratio at $\vartheta=\tau^j$ has valuation at most
$4\pgv(\auNm w_0-1)+\varepsilon$. As $\auNm w_0-1$ divides $\#\mathrm{GL}_{rt}(\auF_{w_0})$, which
divides $\pgNch$, this is at most $4\pgv(\pgNch)+\varepsilon\le(4f+3)r\pgv(\pgNch)+\varepsilon=
e'_\ell$. (5) Proposition~\ref{pg:mirror} is replaced by Proposition~\ref{pse:mirror} and
Proposition~\ref{pg:step} by Proposition~\ref{pse:step} (Remark~\ref{pse:rem-poly}); statement (A)
is Proposition~\ref{pse:A}.
\end{remark}




%% ===== proofs/C-ribet.tex =====
% proofs/C-ribet.tex
% From Selmer bounds to denominators: Theorem pg:III, except the local loss (Proposition pg:loss,
% used as stated in Proposition pri:loss below).
% Labels of this file start with pri:.
\providecommand{\auO}{\mathcal O}
\providecommand{\auQlb}{\overline{\Q}_\ell}
\providecommand{\auGL}{\mathrm{GL}}
\providecommand{\auF}{\mathbb F}
\providecommand{\auU}{\mathrm U}
\providecommand{\auNm}{\operatorname{Nm}}
\providecommand{\auf}[1]{\texttt{#1}}
\providecommand{\panot}[1]{\par\smallskip{\raggedright\noindent\textbf{[Not proved here:} #1\textbf{]}\par}\smallskip}
\providecommand{\pgZb}{\overline{\Z}_\ell}
\providecommand{\pgv}{v}
\providecommand{\pgNch}{N_{\mathrm{ch}}}
\providecommand{\pgcodd}{c_{\mathrm{odd}}}
\providecommand{\pgloc}{\mathrm{loc}}
\providecommand{\pgG}{G_E}
\providecommand{\pgSel}{\operatorname{Sel}}
\providecommand{\pglen}{\operatorname{len}}
\providecommand{\pgtr}{\operatorname{tr}}
\providecommand{\pgFrob}{\operatorname{Frob}}
\providecommand{\pgpr}{\operatorname{pr}}
\providecommand{\pgcyc}{\varepsilon_\ell}
\providecommand{\pgS}{S_{\mathrm{bad}}}
% macros of this file (prefix \pri)
\providecommand{\priFitt}{\operatorname{Fitt}}
\providecommand{\prilen}{\operatorname{len}}
\providecommand{\prirl}{\operatorname{rl}}
\providecommand{\prival}{\operatorname{val}}
\providecommand{\priAug}{\mathfrak I}
\providecommand{\priT}{\mathbb T}
\providecommand{\priX}{\mathcal X}
\providecommand{\priJ}{\mathcal J}
\providecommand{\prid}{\mathcal D}
\providecommand{\priN}{\mathcal N}
\providecommand{\priZ}{\mathcal Z}
\providecommand{\prics}{\mathrm{cs}}
\providecommand{\prity}{\mathrm{ty}}

\section{From Selmer groups to denominators}\label{pri:sec}

This section proves Theorem~\ref{pg:III} of \S\ref{pg:ss-ribet}, the passage from bounds for Selmer
groups to bounds for the denominators of the projection to a space of lifts, \emph{except} for the
local loss above $\ell$ (Proposition~\ref{pg:loss}), which is used in the exact form stated in
Proposition~\ref{pri:loss} below. The inputs listed in \S\ref{pri:ss-inputs} are proved in
\S\S\ref{pod:sec}, \ref{pku:sec}, \ref{pax:sec}, \ref{pay:sec}, \ref{prm:sec}. The
four propositions of \S\ref{pg:ss-ribet} that are assigned here are restated in their exact form:
Ribet's method to first order (Proposition~\ref{pri:ribet}, replacing~\ref{pg:ribet}), the
Taylor--Wiles places in a triangle (Proposition~\ref{pri:triangle}, replacing~\ref{pg:triangle}),
the odd classes (Proposition~\ref{pri:odd}, replacing~\ref{pg:odd}) and the operators at the
auxiliary places (Proposition~\ref{pri:norm}, replacing~\ref{pg:norm}). The inputs of
Theorem~\ref{pg:III} are listed in \S\ref{pri:ss-summary}.

\emph{Conventions.} Those of \S\ref{pg:sec}. In addition: $\auO$ is a discrete valuation ring with
uniformiser $\varpi$, residue field $k$, fraction field $K$ and additive valuation $\prival=
\prival_\auO$ with $\prival(\varpi)=1$; for an $\auO$-module $M$ of finite length, $\prilen M$ is
its (unnormalised) length, so that $\prilen\auO/a=\prival(a)$. When $\auO$ is a coefficient ring
in the sense of \S\ref{pg:sec}, with ramification index $e_\auO$ over $\Z_\ell$, the normalised
length of \S\ref{pg:sec} is $\prilen/e_\auO$ and $\pgv=\prival/e_\auO$. For submodules $S,T$ of a
module, $\prirl(S,T)=\prilen\bigl(S/(S\cap T)\bigr)$ (possibly $\infty$). For a group $G$ and a
ring $A$, $A[G]$ is the group algebra and $\priAug_A\subseteq A[G]$ its augmentation ideal; for a
finite set $J$, $A[X_J]=A[X_j:j\in J]$, $X_S=\prod_{j\in S}X_j$ for $S\subseteq J$, and the
\emph{top coefficient} of $f\in A[X_J]$ is its coefficient of $X_J$. For $g\in\pgG$,
$g^c=cgc$ and $g^{-c}=(g^c)^{-1}$, as in \S\ref{pg:sec}.

\subsection{The exact statement}\label{pri:ss-statement}

Fix $\ell$, $\iota$ and a twist $\vartheta$ with $\psi=\psi_\vartheta$ and $Q=Q_\vartheta$. Put
$N=\pgNch$ (\S\ref{pg:ss-statement}), $h=r^8$, and let $e'_\ell$, $e_\ell$, $\pgloc_\ell$,
$\beta_\ell$ be as in~\eqref{pg:eq-beta}.

\begin{definition}[Auxiliary twists]\label{pri:def-aux}
The twist $\vartheta$ is \emph{auxiliary with ramification at most $\mathrm{ram}$} if the places
of $Q$ have prime norm, $Q$, $\overline Q$ and $\pgS(\ell)$ are pairwise disjoint, $\psi$ is
anticyclotomic, unramified outside $Q\cup\overline Q$ and trivial on $D_v$ for $v\mid\ell$, and
$\sum_{w\in Q}\pgv(\psi(\gamma_w)-1)\le\mathrm{ram}$ for some choice of tame generators
$\gamma_w$. This is Theorem~\ref{pg:II}(a) with $\mathrm{ram}=e'_\ell$.
\end{definition}

Let $O$ be a pair (\S\ref{pg:ss-statement}), $r'=n_W$, $n=n_V$, $t'=t(O)=n-r'$. A family $c$ of
labels at the chosen places is \emph{in lift setting for $\vartheta$} if: the tame part of the
character of $O$ is $\chi_M^{\,j}$ times that of $\vartheta$ for some $1\le j\le m$; $r'$ and $t'$
are odd, $r'\le r$, $t'\le t$; each $c_w$ is a label of degree $r'$; above each of $p_0,p_1$ there
is a chosen $w$ with $r'=r$ and $c_w$ primitive, or $r'=1$ and $c_w=X-1$; and no rescaled label
has the root $1$. Let $e_c$ be the idempotent of $c$.

\begin{definition}[Lift generators]\label{pri:def-gen}
Let $\mathcal Z\subseteq\mathcal F(O)$ and $e$ a central idempotent. A form $y$ is a \emph{lift
generator of $\mathcal Z$ with Selmer bound $\mathrm{sel}\in[0,\infty]$} if $y\in\mathcal Z$; $y\in
e_c\mathcal F(O)\subseteq e\,\mathcal F(O)$ for a family $c$ in lift setting for $\vartheta$; $y$
has the Hecke polynomials of a lift of a system $q$; and there are a main summand $R$ of rank
$r'$ and string length $t'$ (Definition~\ref{pg:def-summand}) satisfying~\eqref{pg:eq-trace} for
every eigenform with the Hecke polynomials of a lift of $q$, an element $\gamma\in\pgG$ and
$\zeta\in\pgZb$, $\zeta-1$ a unit, such that for all large $\auO$ some lattice $u$ in $R$ has
$\psi(\gamma)u(\gamma)=\zeta$ and $\pglen\pgSel_Q(\psi u,\gamma;K/\auO)\le\mathrm{sel}$.
\end{definition}

\begin{theorem}[From Selmer groups to denominators, exact form]\label{pri:III}
Let $\vartheta$ be auxiliary with ramification at most $e'_\ell$, let $O$ be a pair with central
idempotent $e$, and let $\mathcal Z\subseteq\mathcal F(O)$ be spanned by lift generators of
$\mathcal Z$ with Selmer bound $e'_\ell$. Then there is $\Lambda\in\operatorname{End}(\mathcal Z)$
with
\[
\pgv(\det\Lambda)\le\dim\mathcal Z\cdot r'\cdot\beta_\ell
\]
such that for every $y\in\mathcal F(O)$ integral at $\iota$ and orthogonal to the false lifts
(Definition~\ref{pg:def-den}(a), for $e$ and $\mathcal Z$), $\Lambda\pgpr_{\mathcal Z}(y)$ is
integral at $\iota$ and $\langle y',\Lambda\pgpr_{\mathcal Z}(y)\rangle\in\pgZb$ for every $y'$
integral at every isomorphism $\C\to\auQlb$.
\end{theorem}

\begin{remark}[Comparison with Theorem~\ref{pg:III}]\label{pri:rem-compare}
(1) The conclusion is Definition~\ref{pg:def-den}(b) with $\beta=\beta_\ell$ and $n_W=r'$, i.e.\
the conclusion of Theorem~\ref{pg:III}. (2) The hypotheses of Theorem~\ref{pg:III} give those of
Theorem~\ref{pri:III}: Theorem~\ref{pg:II}(a) is Definition~\ref{pri:def-aux} with
$\mathrm{ram}=e'_\ell$, and every spanning eigenform $y$ of Theorem~\ref{pg:III} is a lift
generator with Selmer bound $e'_\ell$ (Theorem~\ref{pg:II}(b)), \emph{provided} its family $c$ is in
lift setting. Theorem~\ref{pg:III} does not mention two clauses of the lift setting: the tame part
of the character, which holds for $O_j(\vartheta)$ and $O_{\nu,z}(\vartheta)$ by construction, and
the parity of $r'$ and $t'$. For $O_{\nu,z}$, $r'=1$ and $t'=2z-1$ are odd; for $O_j$, $r'=r$ and
$t'=t=5r$ are odd if and only if $r$ is odd, and $r$ is odd by (S7) of Lemma~\ref{an:lem-sizes}.
Theorem~\ref{pg:III} states the parity clause explicitly. With these additions the deduction of
Theorem~\ref{pg:main} in \S\ref{pg:ss-reduction} is unchanged.
\end{remark}

\subsection{What is used from outside this section}\label{pri:ss-inputs}

Besides published results named where they are used, the proof uses the following statements. They
are not proved in this section; each is followed by a pointer to the place where it is proved.

\begin{proposition}[The local loss, form used]\label{pri:loss}
Let $\auO$ be a large coefficient ring, $u_x\colon\pgG\to\auGL_n(\auO)$ ($x\in Y$, finite)
continuous, $\rho\colon\pgG\to\auGL_{d}(\auO)$, $M\ge1$. Let $Z^{\mathrm R}_M(\rho)$ be the module of
cocycles $z\colon\pgG\to(\auO/\ell^M)^d$ for $\rho$ satisfying $\sum_ia_iz(g_i)=0$ whenever
$a_i\in\auO$ and $\sum_ia_i(u_x(g_i)-1)=0$ for all $x$ (\emph{Ribet cocycles}), and
$Z^{\mathrm{cr}}_M(\rho)$ the module of functions obeying the crystalline relations of $\rho$
(Definition~\ref{pg:def-sel}). Put $\prics_M(\rho)=\prirl\bigl(Z^{\mathrm R}_M(\rho),
Z^{\mathrm{cr}}_M(\rho)\bigr)$.
\begin{enumerate}
\item[(a)] If, at each $v\mid\ell$, the restrictions to $D_v$ of all $u_x$ and of $\rho$ lie in a
class $\mathcal C_v$ of local lattices (as in Proposition~\ref{pg:loss}), then
$\prics_M(\rho)\le\sum_{v\mid\ell}\delta_v(\rho)$, where the \emph{defect} $\delta_v(\rho)$ is the
length of the crystalline cocycles of the Tate dual of $\rho|_{D_v}$ modulo those whose cup product
with every crystalline cocycle of every member of $\mathcal C_v$, under every pairing to
$\Z/\ell^m(1)$, is a coboundary.
\item[(b)] Let $c$ be in lift setting, $R$ a main summand of rank $r'$ and string length $t'$,
and $\rho_x$ ($x\in Y$) representations of rank $n$ satisfying the conclusion (c) of
Proposition~\ref{pri:BL} for $c$. For all large $\auO$, all lattices $u_x$ in $\rho_x$, all
lattices $T_0$ in $R$ and $T$ equal to $T_0$ on the $D_v$, $v\mid\ell$, and one $\pi\in Y$ with
$u_\pi=T\oplus\bigoplus_{e<t'}\auO(-e)$, there are classes $\mathcal C_v$ as in (a) containing all
$u_x|_{D_v}$ and the restrictions of the $t'+1$ summands $L$ of $u_\pi$, with
\[
\sum_{L}\sum_{v\mid\ell}\delta_v(L)/e_\auO\le r'\,\pgloc_\ell .
\]
\end{enumerate}
\end{proposition}

This is Proposition~\ref{pg:loss} in the form in which it is used (the passage from the local
statements of \S\ref{pll:sec} to (b), a sum over the local degrees $\le2f$ and over the summands with
weights $1$ and $1/t$, is elementary). Proposition~\ref{pg:loss} is stated in the form (a) with Ribet
cocycles. Its proof is given in \S\ref{pll:sec}: for (a), the restrictions to $D_v$ of the Ribet
cocycles come from $\mathcal C_v$ by Lemma~\ref{pll:ribet} (with $\Gamma=D_v$ inside $\pgG$, the
restrictions of the $u_x$ lying in $\mathcal C_v$ by hypothesis), and Proposition~\ref{pll:dual}
with $T=\rho|_{D_v}$ bounds each local contribution; a cocycle obeys the crystalline relations of
$\rho$ as soon as its restriction to every $D_v$, $v\mid\ell$, does. The explicit bounds of (b) come with the explicit exponents of
Theorem~\ref{pll:exponent} proved in \S\ref{pki:ss-deduce}.

\begin{proposition}[Base level of a label]\label{pri:BL}
Let $c$ be in lift setting for an auxiliary $\vartheta$. There are Hodge--Tate data $J'$ and a
space of forms of ``base level'' (the space $\mathcal F(O)$ of Definition~\ref{prm:def-level}: level
$J_O$ and character $\chi_O$; these forms are of principal level at the chosen places) with: (a) $e_c\mathcal F(O)$ consists of forms of base level; (b) if
$y$ is integral at $\iota$, so is $N\,e_cy$; (c) for every finite set $U$ of places off
$\pgS(\ell)\cup Q\cup\overline Q$ and every system $x$ of Hecke eigenvalues occurring in base level
with principal level at $U$, the normalised Galois representation $\rho_x=S_x\otimes\eta_O$ is
conjugate self-dual with string length $t'$, has Hodge--Tate data $J'$, is unramified outside $\ell
p_0p_1$, $Q\cup\overline Q$, $U\cup\overline U$, inertia at a chosen $w\nmid\ell$ acts through the roots of
the rescaled label times $(X-1)^{t'}$ and semisimply (the monodromy of $\operatorname{rec}(\Pi_w)$
is $0$), so that $o_w(\rho_x(s))(\rho_x(s)-1)=0$ for $s\in I_w$, $o_w$ the monic polynomial whose
roots are the (images of the) roots of the rescaled label, and $(\rho_x(s)-\psi(s))(\rho_x(s)-1)=0$
for $s$ in inertia above $Q\cup\overline Q$.
\end{proposition}
\noindent Proved in \S\ref{pax:ss-BL}, from Propositions~\ref{pga:G123} and \ref{pga:G45} and the
cuspidal-support part of Lemma~\ref{pga:depth0}, with Lemma~\ref{pay:eta}; the semisimplicity at a
chosen $w$ is the paragraph ``No monodromy at a chosen $w$'' of that proof (it uses the clause ``no
rescaled label has the root $1$'' of the lift setting); the base level and the
product labels at the chosen places are Definition~\ref{prm:def-level}, Lemmas~\ref{prm:level},
\ref{prm:label} and Corollary~\ref{prm:inertia}.\par\smallskip

\begin{lemma}[Block-diagonal lattices]\label{pri:block}
Let $\rho_x$ ($x\in Y$ finite) be continuous semisimple representations of rank $n$ and,
for $x\in\priX\subseteq Y$, $\pgtr\rho_x=\psi\pgtr R_x+\sum_{e<t'}\chi^e$ with $R_x$ main summands,
and let $u'_x$ be lattices in $R_x$ over large $\auO$ whose values span $M_{r'}(\auO)$ (M4). Then for
all large $\auO$ there are lattices $u_x$ in the $\rho_x$ with $u_x=\psi u'_x\oplus\bigoplus_{e<t'}
\auO(-e)$ for $x\in\priX$, and the $\psi u'_x(g)$ span $M_{r'}(\auO)$.
\end{lemma}
\noindent Proved in Lemma~\ref{pax:block} (Brauer--Nesbitt in characteristic $0$), where the two
hypotheses used (semisimplicity of the $\rho_x$, spanning (M4) of the $u'_x$) are stated.\par\smallskip

\begin{lemma}[Inputs on forms]\label{pri:forms}
(a) (Higman operators; freeness of the level.) For a finite set of places $w_j$ off the base level,
with $\Gamma=\prod_j\auGL_n(\auF_{w_j})$ acting on the forms of principal level at the $w_j$, and for
every subgroup $\Gamma'$ of a factor, there is an integral operator $\Psi$ with
$\sum_{k\in\Gamma'}k^{-1}\Psi k=c_0$, where $c_0$ is a fixed common multiple of the orders of the
stabilisers of the base level, a divisor of $N$; for $\Gamma'=\Gamma$ this says that $c_0y$ is a norm
$\sum_{k\in\Gamma}k\Psi y$ for $y$ of base level. (b) (Pairings.) If $y'$ is integral at every
isomorphism and $y''$ is integral at $\iota$ then $N\langle y',y''\rangle\in\pgZb$.
\end{lemma}
\noindent Proved in \S\ref{pax:ss-forms} (with $c_0=\#(J_O/J_0)$, $J_0$ the free sublevel of the
base level $J_O$, a divisor of $N$) and \S\ref{pax:ss-pairing}.\par\smallskip

\begin{lemma}[Chain lemma]\label{pri:chain}
Let $y$ be an eigenform in base level whose normalised Galois representation has the eigenvalue
$1$ at every element of $\pgG$ (a \emph{weak lift}). At every split place $w\nmid\ell$ off the level
the rescaled Hecke polynomial $H_w(y;X_0X)$, $X_0$ the value at a geometric Frobenius of
$\iota^{-1}\eta_O^{-1}$, is divisible by $\prod_{i<t'}(X-(\auNm w)^i)$.
\end{lemma}
\noindent Proved in \S\ref{pax:ss-chain}, from (G1) and the spherical chain lemma
(Lemma~\ref{pax:tadic}, proved directly in Corollary~\ref{pay:roots}). The statement holds at every
split $w\nmid\ell$ off the level at which $y$ is spherical.\par\smallskip

\begin{proposition}[Odd classes]\label{pri:odd}
Recall that $E$ is a CM field contained in $\Q(\zeta_{q_1d})$. There is an integer
$0<\pgcodd\le2^{4(q_1d+t+3)^{27}}h_E$ such that for every prime $\ell$, every
coefficient ring $\auO$ containing the values of the Dirichlet characters modulo $q_1d$, all $\nu\ge0$
and $0\le e<t$, the module $\mathrm{Sha}^-_{e,\nu}$ has length at most $\prival_\auO(\pgcodd)$. Here,
for $e\ne0$, $\mathrm{Sha}^-_{e,\nu}$ consists of the continuous cocycles $z\colon\pgG\to
(\auO/\ell^\nu)(-e)$ unramified outside $\ell$, vanishing on the elements of trivial cyclotomic
character in the decomposition groups above $\ell$, with $z(g^c)=-z(g)$; for $e=0$ it consists of
the everywhere unramified homomorphisms $\pgG\to\auO/\ell^\nu$.
\end{proposition}
\noindent Proved in Proposition~\ref{pod:odd}, for every prime $\ell$ including $\ell=2$
(Remark~\ref{pod:rem-two}). The argument is an Euler-system argument with Jacobi-sum units, in the
manner of Kolyvagin and of Rubin's proof of the Main Conjecture for $\Q(\mu_{p^\infty})$ (\cite[Appendix by K.~Rubin, p.~397 ff.]{LangCyclo},
whose \S1 works over $\Q(\mu_m)^+$ for general $m$), with explicit constants.\par\smallskip

\begin{proposition}[Kummer maps and detecting points]\label{pri:kummer}
Let $F=E^+$, $\nu>0$, $G_\nu=\ker(\pgcyc\bmod\ell^\nu)\subseteq\pgG$.
\begin{enumerate}
\item[(a)] Let $x=(w_i)_{i\in I}$ be a finite family of degree-one places $w_i\nmid\ell$ of $E$ above
distinct primes of $F$ with Frobenius elements in $G_\nu$, let $S_0(x)$ be the set of primes of $F$
above $\ell$ or below some $w_i$, let $e\in\Z$, and let $(f_i)_{i\in I}$ be maps $G_\nu\to\auF_\ell$
indexed by the same set. There is a Kummer map $\kappa\colon\pgG\to\Z/\ell^\nu$ of exponent $e$ for
$x$ off $T=S_0(x)$ (Definition~\ref{pku:def-kummer}) whose reduction on $G_\nu$ is not in the span
of the $f_i$. Here a Kummer map is given by a locally constant function $y$ on $G_F$, zero outside
$G_\nu$, which on $G_\nu$ is a homomorphism to $\operatorname{Hom}((\Z/\ell^\nu)((\pgcyc^+)^{-e}),\Q/\Z)$,
equivariant for the action of $G_F$ through $(\pgcyc^+)^{e}\pgcyc$, unramified outside $T$,
written $y(g)(a)=a\kappa(g)/\ell^\nu$, which satisfies the reciprocity law:
$\sum_p y(\phi_p)(z(\gamma_p))=0$ for every finite family of distinct primes $p\notin T$ of $F$ with
Frobenius elements $\phi_p\in G_\nu$ and normalised tame generators $\gamma_p$, and every cocycle
$z$ of $G_F$ in $(\Z/\ell^\nu)((\pgcyc^+)^{-e})$ unramified outside $S_0(x)$ and the $p$, vanishing at
complex conjugations, on the decomposition groups of the places $w_i$, and on the part in
$\ker\pgcyc^+$ of the decomposition groups above $\ell$. ($\pgcyc^+$ extends $\pgcyc$ by $1$ at $c$.)
\item[(b)] Let $g_0$ be a topological generator of $\pgG/(\pgG\cap\ker\pgcyc)$. There is a finite set
$s$ of elements of trivial cyclotomic character such that a locally constant cocycle
$\pgG\to(\auO/\ell^\nu)(-e)$ unramified outside $\ell$ that vanishes at $g_0$ and on $s$ is zero,
for all $\auO$, $\nu$, $e$. If $\ell\le h^2$, $s$ can be chosen with $(t+1)(\#s+1)\le h^{21}$.
\end{enumerate}
\end{proposition}
\noindent Proved in \S\ref{pku:sec}: (a) is Proposition~\ref{pku:main} (from the Greenberg--Wiles
formula, local and global duality and Chebotarev's theorem; with $T=S_0(x)$ the relations are
restrictions of dual Selmer classes), and (b) is Proposition~\ref{pku:b} (from
Lemma~\ref{pod:detect}). In (a) the number of maps $f_i$ is at most the number of places of $x$, and
in (b) the size bound is claimed only for $\ell\le h^2$; these are the forms used in
Proposition~\ref{pri:triangle} (Proposition~\ref{pku:triangle}).\par\smallskip

\begin{proposition}[Dictionary at an auxiliary place]\label{pri:dict}
Let $w$ be a degree-one place off the base level with $\auNm w\equiv1\pmod{\ell^\nu}$, $\gamma$ the
tame generator and $\phi$ a Frobenius element at $w$, and $\Gamma_w=\auGL_n(\auF_w)$ acting on the
forms of principal level at $w$. Fix a uniformiser $\varpi$ and a logarithm $\mathrm{ind}\colon
\auF_w^\times\to\Z/\ell^\nu$ with $\mathrm{ind}(\bar z)=1$ for the unit $z$ whose Artin symbol is
$\gamma$. Write $\langle\gamma\rangle=\mathrm{diag}(\bar z,1,\dots,1)\in\Gamma_w$, $U$ for the
Hecke operator of $\mathrm{diag}(\varpi,1,\dots,1)$ at pro-$p$ Iwahori level, $\tilde U=
\eta_O(\phi)U$, and $\operatorname{tr}[\,\cdot\,]$ for the traces of the normalised Galois
representations as elements of the Hecke algebra $\priT$ at principal level at $w$.
\begin{enumerate}
\item[(a)] On a form on which $\Gamma_w$ acts by $\chi\circ\mathrm{ind}\circ\det$,
$\operatorname{tr}[(\gamma-1)g\,p(\phi)]=(\langle\gamma\rangle-1)\operatorname{tr}[g\,p(\phi)]$ for
all $g$ and $p\in\pgZb[X]$.
\item[(b)] On a form on which the stabiliser of a line in $\Gamma_w$ acts by $\chi(\mathrm{ind}(
k_{11}))$ with $\chi(n!)\ne1$, $\operatorname{tr}[(\gamma-1)g\,p(\phi)]=p(\tilde U)
\operatorname{tr}[(\gamma-1)g]$.
\item[(c)] $\sum_{k\in\Gamma_w}k$ kills $\operatorname{tr}[(\gamma-1)g]$.
\item[(d)] If $\delta\in\priT$ and every base system on which $\delta\ne0$ has, at $w$, a rescaled
Hecke polynomial equal to $\prod_{i<t'}(X-(\auNm w)^i)$ times a divisor of $H$, then $\delta
\sum_{k\in\Gamma}k$ kills the value at the weight $1$ of the operator $B$ of~\eqref{pri:eq-B}.
\end{enumerate}
\end{proposition}
\noindent Proved in \S\ref{pax:ss-dict} (parts (a)--(d)); part (d) uses Lemma~\ref{pax:toproot},
proved in \S\ref{pay:ss-cons}.\par\smallskip

\subsection{Fitting ideals}\label{pri:ss-fitting}

Let $R=\auO[X_J]$, $J$ finite, and $s$ a finite set. For an $R$-submodule $N\subseteq R^s$ let
$\priFitt(N)$ be the ideal generated by the determinants of the $s\times s$ matrices all of whose
rows lie in $N$ (the zeroth Fitting ideal of $R^s/N$). Let $\Phi\colon\auO^n\to R^s$ be
$\auO$-linear, $p\in\auO\smallsetminus\{0\}$ with $\Phi(p\auO^n)\subseteq N$, and $z\colon G\to\auO^n$
a function on a set $G$. The \emph{forms} of $N$ are
\[
\mathrm{F}(N)=\{\Psi\circ z:\ \Psi\colon\auO^n\to\auO/p\ \text{linear},\ \Psi(\Phi^{-1}N)=0\}
\subseteq\operatorname{Map}(G,\auO/p).
\]

\begin{lemma}\label{pri:fitt-basic}
(a) If $N\subseteq N'$ then $\priFitt(N)\subseteq\priFitt(N')$. (b) If $a\in R$, $v\in R^s$ and
$av\in N$, then $a\priFitt(N+Rv)\subseteq\priFitt(N)$. (c) If $w\in\auO^n$ and $\Psi(w)=0$ for every
form-defining $\Psi$ (that is, every $\Psi$ with $\Psi(\Phi^{-1}N)=0$), then $\Phi(w)\in N$.
(d) Let $\mathfrak m=\ker(R\to\auO,\ X_j\mapsto0)$. Then $\priFitt(N+\mathfrak mR^s)\subseteq
\priFitt(N)+\mathfrak m$.
\end{lemma}
\begin{proof}
(a) is clear. (b) A row of $N+Rv$ is $n_i+r_iv$; expand the determinant multilinearly: the terms
with two rows proportional to $v$ vanish, those with no such row lie in $\priFitt(N)$, and in a term
with one such row multiplication by $a$ turns that row into $r_i(av)\in N$. (c) $M=\auO^n/\Phi^{-1}N$
is killed by $p$, hence a finitely generated $\auO/p$-module; $\auO/p$ is self-injective (a
quotient of a discrete valuation ring by a nonzero ideal), so linear maps $M\to\auO/p$ separate
points; thus $w\in\Phi^{-1}N$. (d) Reduce modulo $\mathfrak m$ entrywise: the determinant of a matrix
with rows $n_i+m_i$ is congruent modulo $\mathfrak m$ to that with rows $n_i$.
\end{proof}

\begin{lemma}[Cost of a submodule of forms]\label{pri:fitt-cost}
Let $W\subseteq\operatorname{Map}(G,\auO/p)$ be a submodule. There are $e\in\auO$ with $\prival(e)\le
\prirl(\mathrm F(N),W)$ and $K\supseteq N$ such that $\Phi(w)\in K$ for every $w$ with $\Psi(w)=0$
for all form-defining $\Psi$ with $\Psi\circ z\in W$, and $e\priFitt(K)\subseteq\priFitt(N)$.
\end{lemma}
\begin{proof}
Let $M=\auO^n/\Phi^{-1}N$ (finite length, killed by $p$) and $M^*=\operatorname{Hom}(M,\auO/p)$; the
map $\Psi\mapsto\Psi\circ z$ maps $M^*$ onto $\mathrm F(N)$. Let $M^*_W$ be the preimage of $W$
and $M'=\{m\in M:\Psi(m)=0\ \forall\Psi\in M^*_W\}$. By Matlis duality over the artinian
self-injective ring $\auO/p$, $\prilen(M')=\prilen(M^*/M^*_W)$; and $M^*/M^*_W$ maps onto
$\mathrm F(N)/(\mathrm F(N)\cap W)$ with kernel the image of those $\Psi$ with $\Psi\circ z\in W$
that already lie in $M^*_W$, i.e.\ isomorphically; so $\prilen M'=\prirl(\mathrm F(N),W)$. Choose a
composition series $0=M'_0\subset\dots\subset M'_k=M'$, $k=\prilen M'$, and lifts $w_i\in\auO^n$ of
generators of $M'_i/M'_{i-1}$; then $\varpi w_i\in\auO w_1+\dots+\auO w_{i-1}+\Phi^{-1}N$. Put
$K_i=N+R\Phi(w_1)+\dots+R\Phi(w_i)$, $K=K_k$, $e=\varpi^k$. Then $\varpi\Phi(w_i)\in K_{i-1}$ and
Lemma~\ref{pri:fitt-basic}(b) gives $\varpi\priFitt(K_i)\subseteq\priFitt(K_{i-1})$. Every $w$ as in
the statement has image in $M'$, so $\Phi(w)\in K$.
\end{proof}

\begin{definition}[Triangle cost]\label{pri:def-cost}
Let $\priZ\subseteq\operatorname{Map}(G,B)$ be a submodule, $B$ an $\auO$-module, $\tau,\tau',
\phi,\phi'\colon J\to G$ with $J$ totally ordered, $I\subseteq J$, $C\in\auO$, $b\in[0,\infty]$.
$\priZ$ has \emph{triangle cost $b$ along $I$ with constant $C$} if there is $A\subseteq G$ such
that, with $\priZ_0=\{c\in\priZ:c|_A=0\}$: (a) for $j\in I$ and $c\in\priZ_0$ with $c(\tau_i)=
c(\tau'_i)$ for all $i<j$ and $c(\phi_i)=c(\phi'_i)$ for $i<j$ in $I$, one has $C(c(\tau_j)-
c(\tau'_j))=0$; (b) for every submodule $\priZ_1\subseteq\priZ_0$ whose members vanish at all
$\tau_j,\tau'_j$ and satisfy $c(\phi_j)=c(\phi'_j)$ for $j\in I$, $\prirl(\priZ,\priZ_0)+\prirl(
\priZ_1,\priZ_1[C])\le b$, where $\priZ_1[C]$ is the $C$-torsion.
\end{definition}

\begin{proposition}[Fitting ideal of a triangle]\label{pri:fitt-triangle}
Keep $\Phi$, $p$, $z$, and put $y_g=\Phi(z_g)$. Suppose: modulo $N$, the basis vector $\epsilon_g$
of $R^s$ is $y_g$ for $g\in s\subseteq G$; $y_{\tau_j}-X_jy_{\phi_j}\in N$ and $y_{\tau'_j}-X_j
y_{\phi'_j}\in N$ for all $j$; every $i$ below an element of $I$ lies in $I$ or satisfies $X_iR^s
\subseteq N$; and $\mathrm F(N)$ has triangle cost $b$ along $I$ with constant $C$. Then there is
$F\in R$ with $F\cdot X_I\in\priFitt(N)$ and $\prival(F(0))\le b+(\#I+\#s)\prival(C)$.
\end{proposition}
\begin{proof}
Let $A$ be as in Definition~\ref{pri:def-cost}. \emph{Stage 1.} Lemma~\ref{pri:fitt-cost} with
$W=\{c:c|_A=0\}$ gives $e$, $K\supseteq N$, with $\prival(e)\le\prirl(\mathrm F(N),\mathrm F(N)_0)$
and $e\priFitt(K)\subseteq\priFitt(N)$; every $K'\supseteq K$ has $\mathrm F(K')\subseteq\mathrm
F(N)_0$. \emph{Stage 2.} Put $v_j=\Phi(z_{\phi_j}-z_{\phi'_j})$, so $\Phi(z_{\tau_j}-z_{\tau'_j})
\equiv X_jv_j\pmod K$. Going up $I$ in increasing order, let $T_j=T_{j^-}+Rv_j$ ($T_{\min^-}=K$). We
claim $CX_jv_j\in T_{j^-}$. For $i<j$, $X_iv_i\in T_{j^-}$ (if $i\in I$, $v_i\in T_{j^-}$; otherwise
$X_iR^s\subseteq N$). So every form-defining $\Psi$ of $T_{j^-}$ has $\Psi(z_{\tau_i})=\Psi(
z_{\tau'_i})$ for $i<j$ and $\Psi(z_{\phi_i})=\Psi(z_{\phi'_i})$ for $i<j$ in $I$; as $\Psi\circ z\in
\mathrm F(N)_0$, Definition~\ref{pri:def-cost}(a) gives $C\Psi(z_{\tau_j}-z_{\tau'_j})=0$. By
Lemma~\ref{pri:fitt-basic}(c), $C\Phi(z_{\tau_j}-z_{\tau'_j})\in T_{j^-}$, hence $CX_jv_j\in
T_{j^-}$. Lemma~\ref{pri:fitt-basic}(b) gives $CX_j\priFitt(T_j)\subseteq\priFitt(T_{j^-})$; with
$T=T_{\max}$, $\priFitt(T)\prod_{j\in I}CX_j\subseteq\priFitt(K)$. \emph{Stage 3.} Let $T_0=T+
\mathfrak mR^s$. Lemma~\ref{pri:fitt-cost} for $T_0$ and $W=\{C\text{-torsion}\}$ gives $f$ and
$K'\supseteq T_0$ with $\prival(f)\le\prirl(\mathrm F(T_0),\mathrm F(T_0)[C])$ and $f\priFitt(K')
\subseteq\priFitt(T_0)$. Since $\Psi(Cz_g)=0$ whenever $\Psi\circ z$ is $C$-torsion, $\Phi(Cz_g)\in
K'$, and $C\epsilon_g=\Phi(Cz_g)-C(y_g-\epsilon_g)\in K'$; so $C^{\#s}\in\priFitt(K')$ and $fC^{\#s}\in
\priFitt(T_0)\subseteq\priFitt(T)+\mathfrak m$ (Lemma~\ref{pri:fitt-basic}(d)). Thus some $d\in
\priFitt(T)$ has $d(0)=fC^{\#s}$. \emph{Costs.} The members of $\mathrm F(T_0)$ vanish at all
$\tau_j,\tau'_j$ (as $y_{\tau_j}\in X_jR^s+N\subseteq T_0$) and satisfy $c(\phi_j)=c(\phi'_j)$ for
$j\in I$ (as $v_j\in T$), and lie in $\mathrm F(N)_0$; so by (b) of the definition $\prival(e)+
\prival(f)\le b$. Then $F=eC^{\#I}d$ satisfies $FX_I\in\priFitt(N)$ and $\prival(F(0))\le b+
(\#I+\#s)\prival(C)$.
\end{proof}

\subsection{Ribet's method over an algebra}\label{pri:ss-ribet-alg}

Let $\mathsf R$ be a commutative ring, $\Lambda$ an $\mathsf R$-algebra and $\rho\colon\Lambda\to
M_\sigma(\mathsf L)$ an $\mathsf R$-algebra map to matrices over a commutative $\mathsf
R$-algebra $\mathsf L$, such that the characteristic polynomial of every $\rho(\lambda)$ has
coefficients in the image of $\mathsf R$. For a matrix $P\in M_m(\Lambda)$ let $\rho(P)\in
M_{m\sigma}(\mathsf L)$ be the block matrix $(\rho(P_{ij}))$.

\begin{lemma}[Amitsur]\label{pri:amitsur}
For every $P\in M_m(\Lambda)$ the characteristic polynomial of $\rho(P)$ has coefficients in the
image of $\mathsf R$.
\end{lemma}
\begin{proof}
Write $\rho(P)=\sum_qE_q$ with $E_q$ the block matrix whose only nonzero block row is the $q$-th
row of $\rho(P)$. By Amitsur's formula (Proposition~\ref{pax:amitsur}; \cite{Amitsur}) the coefficients of the characteristic
polynomial of a sum of matrices are integral polynomials in the coefficients of the characteristic
polynomials of the products $E_{q_1}\cdots E_{q_k}$. Such a product is
$\epsilon_{q_1}\,\rho(P_{q_1q_2}P_{q_2q_3}\cdots P_{q_{k-1}q_k})\,\pi_{q_k}$, where $\epsilon_q$ is
the $q$-th block column of $1$ and $\pi_q$ the $q$-th block row of $\rho(P)$; its characteristic
polynomial agrees, up to a power of $X$, with that of $\pi_{q_k}\epsilon_{q_1}\rho(\cdots)=
\rho(P_{q_kq_1}P_{q_1q_2}\cdots P_{q_{k-1}q_k})$, which has coefficients in $\mathsf R$.
(Details in Remark~\ref{pax:rem-amitsur}.)
\end{proof}

\begin{lemma}[Ribet's polynomial]\label{pri:ribet-poly}
Let $\mathfrak a\subseteq\Lambda$ be a two-sided ideal containing $\ker\rho$, $y\in\Lambda^m$,
$I=\sum_k\Lambda y_k$, $\rho_0\in\Lambda$, and $B\in M_m(\Lambda)$ with the entries of $By$ in
$\mathfrak a+I\rho_0$. Then there is $P\in\mathsf R[X]$ with $P(0)=\det\rho(B)$ (up to the image of
$\mathsf R$) and $I\,P(\rho_0)\subseteq\mathfrak a$.
\end{lemma}
\begin{proof}
Write $By\equiv B_1y\rho_0\pmod{\mathfrak a^m}$ with $B_1\in M_m(\Lambda)$. Let $\Lambda[X]$ act on
$\Lambda/\mathfrak a$ with $\Lambda$ acting on the left and $X$ by right multiplication by
$\rho_0$ (well defined as $\mathfrak a$ is two-sided). Then $B_X=B-XB_1$ kills the image of $y$.
Extend $\rho$ to $\rho_X\colon\Lambda[X]\to M_\sigma(\mathsf L[X])$ coefficientwise; as
$\Lambda[X]=\Lambda\otimes_{\mathsf R}\mathsf R[X]$, Lemma~\ref{pri:amitsur} (for $\mathsf R[X]$)
puts the characteristic polynomial $\chi(T)$ of $\rho_X(B_X)$ over $\mathsf R[X]$. By
Cayley--Hamilton, $\chi(\rho_X(B_X))=0$; writing $\chi(T)=TQ(T)+\chi(0)$ and $B'=-Q(B_X)\in
M_m(\Lambda[X])$ (a polynomial in $B_X$ with coefficients in $\mathsf R[X]$), we get $\rho_X(B'B_X)=
P\cdot1$ with $P=\chi(0)=\pm\det\rho_X(B_X)$, so $P(0)=\pm\det\rho(B)$. The entries of $B'B_X-P$
have all coefficients in $\ker\rho\subseteq\mathfrak a$, so $P$ kills the image of $y$ too:
$y_kP(\rho_0)\in\mathfrak a$ for all $k$.
\end{proof}

\begin{proposition}[Ribet element of a matrix of tied rows]\label{pri:ribet-elt}
Let $u_x\colon\auO[G]\to M_\sigma(\auO)$ ($x\in Y$ finite) be representations, $\mathsf L$ an
$\auO$-algebra, $S$ an $\auO$-subalgebra of $\operatorname{Map}(Y,\mathsf L)$ containing the
constants and, as functions of $x$, the coefficients of the characteristic polynomials of the
$u_x(g)$, and $\priJ$ an $S$-submodule of $\operatorname{Map}(Y,\mathsf L)$. Let $\mathfrak a_0$ be
the left ideal of $\auO[G]$ spanned by the $r$ with $(x\mapsto\pgtr u_x(rg))\in\priJ$ for all $g$, let
$y\in\auO[G]^m$ span a left ideal $I\supseteq\priAug$ modulo $\bigcap_x\ker u_x$, and $B\in
M_m(\auO[G])$ with the entries of $By$ in $\mathfrak a_0+I\rho_0$. Suppose that for $x$ in a subset
$\priX$, $u_x$ is (up to a fixed change of basis) $\psi_x\oplus\bigoplus_{e<t'}\eta_e$ with characters
$\eta_e$, that $\rho_1$ acts as $\mu_0$ on $\eta_0$ and as $0$ on the other summands, and that
$\rho_0$ acts as $0$ on $\eta_0$. Then there is $f\in S$ with $f\cdot\det u_x(w)\in\priJ$ (as functions
of $x$) for every $w\in\priAug$, and $f(x)=\mu_0\det\psi_x(B)\prod_e\det\eta_e(B)$ for $x\in\priX$.
\end{proposition}
\begin{proof}
By Lemma~\ref{pri:amitsur} applied to the family, the map $\rho=(u_x)_x$ extends to an
$S$-algebra map $\rho\colon S[G]\to M_\sigma(\operatorname{Map}(Y,\mathsf L))$ all of whose
characteristic polynomials have coefficients in $S$. Let $\mathfrak a$ be the set of
$\lambda\in S[G]$ with $\pgtr\rho(\lambda\mu)\in\priJ$ for all $\mu\in S[G]$: a two-sided ideal
(cyclicity of the trace) containing $\ker\rho$ and the image of $\mathfrak a_0$ (expand $\mu$ in
group elements). Lemma~\ref{pri:ribet-poly} gives $P\in S[X]$ with $P(0)=\det\rho(B)$ and $I\,
P(\rho_0)\subseteq\mathfrak a$. Put $f=\pgtr\rho(\rho_1P(\rho_0))\in S$. For $x\in\priX$,
$\rho_1\rho_0$ acts as $0$, so $\rho_1P(\rho_0)$ acts as $P(0)\rho_1$ and $f(x)=P(0)(x)\mu_0$; and
$\det\rho(B)$ at $x$ is the product of the determinants on the summands. For $w\in\priAug$, Cayley--
Hamilton gives $w''\in S[w]$ with $\rho(w''w)=\det\rho(w)$. Since $\priAug$ is two-sided, $w\rho_1\in
\priAug\subseteq I+\ker\rho$, so $w''w\rho_1P(\rho_0)\in\mathfrak a$, and $\det\rho(w)\,f=\pgtr\rho(
w''w\rho_1P(\rho_0))\in\priJ$.
\end{proof}

\subsection{Ribet's method to first order}\label{pri:ss-ribet}

Fix the following data. $G$ is a group, $\auO$ a discrete valuation ring with fraction field $K$,
$u_x\colon G\to\auGL_\sigma(\auO)$ ($x\in Y$, finite), $S$, $\priJ$ as in
Proposition~\ref{pri:ribet-elt} but with values in $\mathsf L[X_J]$, $J$ a finite totally ordered
set; $p\in\auO\smallsetminus\{0\}$; elements $\tau_k,\phi_k\in G$ for $k\in J\times\{\pm\}$;
polynomials $H_j\in\auO[X]$. \emph{The tied relations} are $\tau_k-1-X_j(\phi_k-1)H_j(\phi_k)$,
$k=(j,\pm)$. We say $\priJ$ is a \emph{trace module} if $S$ contains the $X_j$ and the coefficients
of the characteristic polynomials of the $u_x(g)$, and $\priJ$ contains the families $x\mapsto
p\,\pgtr u_x(g)$ and $x\mapsto\pgtr u_x\bigl((\tau_k-1-X_j(\phi_k-1)H_j(\phi_k))g\bigr)$ for all $g$
and $k$. A \emph{common Ribet element} at $\pi\in\priX\subseteq Y$ with value $d$ is $f\in S$ such
that $f\cdot\det(1-u_x(g))\in\priJ$ for all $g$, no monomial of $f(x)$, $x\in\priX$, properly divides
$X_J$, and the top coefficient of $f(\pi)$ is $d$.

For a character $\xi\colon G\to\auO^\times$, a \emph{character cocycle for $\xi$} (relative to
$\delta_v\in G$, $\gamma$, $a\in G$) is a cocycle $c\colon G\to\auO/p$ for $\xi$ that obeys all
linear relations among the families $(u_x(g)-1)_x$, vanishes at the $\delta_v$, and, if $c(\gamma)
=0$, satisfies $(\xi(g)-1)c(a)=0$ for all $g$.

\begin{proposition}[Ribet's method to first order; exact form of Proposition~\ref{pg:ribet}]
\label{pri:ribet}
Assume:
\begin{enumerate}
\item[(1)] for $x\in\priX$, $u_x=\rho_x\oplus\bigoplus_{e<t'}\xi_e$ (fixed change of basis) with
$\xi_0=1$, and the $\rho_\pi(g)$ span $M_{r'}(\auO)$;
\item[(2)] $\gamma_x\in G$ acts on $\rho_x$ by a scalar $\zeta_x$ with $\zeta_x-1\in\auO^\times$, and
trivially on the $\xi_e$;
\item[(3)] $\priJ$ is a trace module, and $H_j(\xi_e(\phi_k))\equiv1\pmod p$ for all $e,k$;
\item[(4)] $a\in G$ and $C\in\auO$ with $\prod_{e'\ne e}(\xi_e(a)-\xi_{e'}(a))\mid C$ for every $e$;
and $\prity\colon J\to\{0,\dots,t'-1\}$ is monotone;
\item[(5)] the cocycles $G\to(K/\auO)^{r'}$ for $\rho_\pi$ obeying all linear relations among the
$(u_x(g)-1)_x$ and vanishing at the $\tau_k$ and at $\gamma_\pi$ form a module of length $\le b_\rho$;
\item[(6)] $\delta_v\in G$ ($v$ in a finite set) are trivial on the $\xi_e$, and $o_v(u_x(\delta_v))
(u_x(\delta_v)-1)=0$ for all $x$, with $o_v\in\auO[X]$;
\item[(7)] for each $e$, every module of character cocycles for $\xi_e$ has triangle cost $b_e$
along $\{j:\prity(j)=e\}$ with constant $C$ (for $\tau_j=\tau_{(j,+)}$, $\tau'_j=\tau_{(j,-)}$ and
likewise $\phi$), and a finite set $s\subseteq G$ detects it (a member vanishing on $s$ is $0$).
\end{enumerate}
Then there is a common Ribet element at $\pi$ whose value $d$ lies in $\auO$ with
\[
\prival(d)\le b_\rho+\sum_{e<t'}b_e+t'\Bigl(\sum_v\prival(o_v(1))+2\#s\,\prival(C)\Bigr)+(\#J+1)
\prival(C).
\]
\end{proposition}

\begin{proof}
\emph{Generators.} As $Y$ is finite and $\auO$ noetherian, the $\auO$-module spanned by the families
$(u_x(g)-1)_x$ is finitely generated; choose $s_0\subseteq G$ finite so that, with the generators
$\mathrm{gen}=(\delta_v)_v\sqcup s_0\sqcup s$, there are vectors $z_g\in\auO^{\mathrm{gen}}$ with
$u_x(g)-1=\sum_iz_{g,i}(u_x(\mathrm{gen}_i)-1)$ for all $x,g$. A \emph{row} is a vector $r\in
\auO[X_J][G]^{\mathrm{gen}}$, with \emph{value} $\sum_ir_i(\mathrm{gen}_i-1)$. Let $c_\zeta=(\zeta_\pi-1)
^{-1}$, $\nu_0=c_\zeta(\gamma_\pi-1)$ and $\rho_0=\nu_0+(1-\nu_0)(a-1)$. The \emph{tie ideal} is the left
ideal spanned by $\bigcap_x\ker u_x$, $p\priAug$, the tied relations, and $\priAug\rho_0$; a row is
\emph{tied} if its value lies in the tie ideal. Since $\priJ$ is a trace module, $p$ and the tied
relations lie in $\mathfrak a_0$ (Proposition~\ref{pri:ribet-elt}); and $\priAug=\sum_i\auO[X_J][G]
(\mathrm{gen}_i-1)$ modulo $\bigcap\ker u_x$. Thus the values of tied rows lie in $\mathfrak a_0+
I\rho_0$ with $I$ the left ideal of the $\mathrm{gen}_i-1$, and Proposition~\ref{pri:ribet-elt}
applies to every square matrix $B$ of tied rows (with $y_i=\mathrm{gen}_i-1$). The following rows are
tied: the cocycle rows $z_{gh}-z_g-g z_h$; $z_{\mathrm{gen}_i}-\epsilon_i$; $\sum\theta_gz_g$ whenever
$\sum\theta_g(u_x(g)-1)=0$ for all $x$; $p\cdot r$ for any row $r$; the ties $z_{\tau_k}-X_jH_j(
\phi_k)z_{\phi_k}$; $(g-1)(c_\zeta z_{\gamma_\pi}+(1-\nu_0)z_a)$; and the rows $o_v(\delta_v)\epsilon_v$.
(Each value acts as $0$ under every $u_x$, or is one of the generators of the tie ideal.)

\emph{The block.} Read rows through $\rho_\pi$ at $X_J=0$: $g\mapsto\rho_\pi(g)$, so a row gives a
matrix in $M_{r'}(\auO)^{\mathrm{gen}}$, i.e.\ a row of $r'\times r'$ blocks. Let $N_\rho$ be the
$\auO$-span of the images of tied rows. A vector $w\in((K/\auO)^{r'})^{\mathrm{gen}}$ killed by
$N_\rho$ defines $\Psi\colon\auO^{\mathrm{gen}}\to(K/\auO)^{r'}$, and $c=\Psi\circ z$ is a cocycle for
$\rho_\pi$ (cocycle rows), obeys the linear relations, vanishes at the $\tau_k$ (at $X_J=0$ the tie
is the row $z_{\tau_k}$) and at $\gamma_\pi$ (the normaliser row at $X_J=0$, through $\rho_\pi$, is
$c_\zeta z_{\gamma_\pi}$ times $(\rho_\pi(g)-1)$, and the $\rho_\pi(g)-1$ span; with $c_\zeta(\zeta_\pi-1)=
1$). Conversely $c$ determines $w$ (generator rows). By (5) the module of such $w$ has length $\le
b_\rho$, so (structure of finitely generated modules over $\auO$, and $(K/\auO)$-duality) the cokernel of
$N_\rho\to M_{r'}(\auO)^{\mathrm{gen}}$, read as $\auO$-modules of rank $r'\#\mathrm{gen}$, has length
$\le b_\rho$, and there is a square matrix $B_0$ of tied rows whose block determinant through $\rho_\pi$
at $X_J=0$ has valuation $\le b_\rho$.

\emph{The characters.} Fix $e$. Read rows through $\xi_e$ (each $g\mapsto\xi_e(g)\in\auO$),
forgetting the $\delta$-coordinates, and let $N_e\subseteq\auO[X_J]^{s_0\sqcup s}$ be the image of
the tied rows. The forms of $N_e$ (with $\Phi$ the inclusion of $\auO^{s_0\sqcup s}$ and $z$ as
above) are character cocycles for $\xi_e$: cocycle and linearity from the corresponding rows, the
vanishing at $\delta_v$ as explained below, and $(\xi_e(g)-1)(c_\zeta c(\gamma)+c(a))$ from the normaliser rows, as
$\xi_e(\gamma_\pi)=1$. On $\xi_e$ the tie reads $z_{\tau_k}-X_jH_j(\xi_e(\phi_k))z_{\phi_k}$, and
$H_j(\xi_e(\phi_k))=1+pc$ by (3); subtracting the tied row $pX_jcz_{\phi_k}$ gives the relation
$y_{\tau_k}\equiv X_jy_{\phi_k}$ of Proposition~\ref{pri:fitt-triangle}. Since $s$ detects the forms,
the elements of $s_0$ can be eliminated: by Lemma~\ref{pri:fitt-basic}(c) there are tied rows whose
$s_0$-part, on $\xi_e$, is the unit vector $\epsilon_i$; one row can be made to serve all $e$ at once
by correcting it successively on each character with the elements $\pi_e\in\auO[G]$ that are $C$ on
$\xi_e$ and $0$ on the other $\xi_{e'}$ (they exist by (4): $C\prod_{e'\ne e}(a-\xi_{e'}(a))/
\prod_{e'\ne e}(\xi_e(a)-\xi_{e'}(a))$). After elimination, Proposition~\ref{pri:fitt-triangle}
applies with $I=\{\prity=e\}$ to $N'_e=N_e+\mathfrak n'R^{s}$ in place of $N_e$, where $\mathfrak n$ is
the ideal of monomials not dividing $X_J$ and $\mathfrak n'=(\mathfrak n:X_{\{\prity<e\}})$, which
contains every $X_i$ with $\prity(i)<e$: thus $X_iR^s\subseteq N'_e$ for these $i$, the $i$ below
$I$ ($\prity$ is monotone), so the hypothesis of Proposition~\ref{pri:fitt-triangle} holds. Its
other hypotheses pass from $N_e$ to $N'_e$: $\Phi^{-1}N'_e\supseteq\Phi^{-1}N_e$, so the forms of
$N'_e$ are forms of $N_e$, hence character cocycles, with triangle cost $b_e$ by (7). Finally
$\priFitt(N'_e)\subseteq\priFitt(N_e)+\mathfrak n'$, by multilinear expansion of a determinant with
rows $n_i+m_i$, $n_i\in N_e$, $m_i\in\mathfrak n'R^s$. (The vanishing of the forms at $\delta_v$ is
automatic: the tied generator row $z_{\delta_v}-\epsilon_v$ shows that, modulo the image of the tied
rows, $z_{\delta_v}$ has zero component in $\auO^{s_0\sqcup s}$; the price is that the rows $o_v(\delta_v)\epsilon_v$, which read as
$o_v(1)\epsilon_v$ on $\xi_e$, must be used to complete the matrices below, which is the factor
$\prod_vo_v(1)$ in $c_\chi$.) The result: matrices $B_{e,i}$ over $\auO[X_J]$ and
weights $w_{e,i}$ with $\sum_iw_{e,i}\det B_{e,i}=\Gamma_eX_{\{\prity=e\}}+\zeta_e$, $\zeta_e\in(
\mathfrak n:X_{\{\prity<e\}})$, $\prival(\Gamma_e(0))\le b_e+(\#\{\prity=e\}+\#s)\prival(C)$; and
for each choice of one $B_{e,i(e)}$ for every $e$ a single matrix $M_i$ of tied rows with
$\det\xi_e(M_i)=c_\chi\det B_{e,i(e)}$ for all $e$, $c_\chi=C^{\#s}\prod_vo_v(1)$, obtained by the
common-unit-row construction just described (and completing with the rows $o_v(\delta_v)\epsilon_v$,
which makes the matrix block triangular on each $\xi_e$).

\emph{Gluing.} Let $k\in\auO[G]$ be $1$ on the $\xi_e$ and $0$ on every $\rho_x$, $x\in\priX$: take
$k=\prod_x(\zeta_x-1)^{-1}(\zeta_x-\gamma_x)$, each factor being $0$ on $\rho_x$ and $1$ on the
characters. Apply Proposition~\ref{pri:ribet-elt} to the matrices
$(1-k)B_0+kM_i$ (rows tied, as tied rows form a left $\auO[X_J][G]$-module), with $\rho_1=k\pi_0$ ($C$
on $\xi_0$, $0$ on the other $\xi_e$ and on every $\rho_x$; $\rho_0$ is $0$ on $\xi_0=1$), and sum with weights $\prod_ew_{e,i(e)}$.
For $x\in\priX$ the value is $C\det\rho_x(B_0)\cdot c_\chi^{t'}\prod_e(\Gamma_eX_{\{\prity=e\}}+
\zeta_e)$. Since each factor $\Gamma_eX_{\{\prity=e\}}+\zeta_e$ is preceded by the earlier factors,
which are multiples of $X_{\{\prity<e\}}$, the product has no monomial properly dividing $X_J$ and
top coefficient $\prod_e\Gamma_e(0)$ times the constant term of the rest. The top coefficient at
$\pi$ is $d=C\,c_\chi^{t'}\det\rho_\pi(B_0)|_{X=0}\prod_e\Gamma_e(0)$ (the remaining factors are taken
at $X_J=0$). Hence $\prival(d)\le\prival(C)+t'(\#s\prival(C)+\sum_v\prival(o_v(1)))+b_\rho+\sum_eb_e+
(\#J+t'\#s)\prival(C)$, which is the bound stated. Finally $\det(1-u_x(g))=\det u_x(w)$ for
$w=-(g-1)\in\priAug$ (computed through the coordinates $z$), so $f$ is a common Ribet element.
\end{proof}

\begin{remark}
The paragraphs ``The characters'' and ``Gluing'' are written in full in \S\ref{prm:ss-chars}
(Lemmas~\ref{prm:tied}--\ref{prm:coeff}, Propositions~\ref{prm:block}, \ref{prm:chars}). The passage
from coefficients modulo $\ell^M$ to $K/\auO$ in (5) is Lemma~\ref{pri:divisible}.
\end{remark}

\begin{lemma}[Divisible coefficients]\label{pri:divisible}
Let $\mathcal V_B$ be the module of cocycles $G\to B^{r'}$ for $\rho$ obeying the linear relations
among the $(u_x(g)-1)_x$ and vanishing on a set $T$ and at $\gamma$. If $\prilen\mathcal V_{\auO/
p^M}\le b$ for all $M$, with $p$ a nonzero non-unit, then $\prilen\mathcal V_{K/\auO}\le b$.
\end{lemma}
\begin{proof}
Every element of $\mathcal V_{K/\auO}$ is torsion (its values are combinations, with coefficients
$z_g$, of its finitely many values at the generators). With $a=p^{b+1}$, multiplication by $a^{-1}$
identifies $\mathcal V_{\auO/a}$ with the $a$-torsion of $\mathcal V_{K/\auO}$. A torsion module all of
whose $a$-torsion has length $\le b<\prival(a)$ is killed by $\varpi^b$, hence equals its
$a$-torsion.
\end{proof}

\subsection{Even classes, the triangle and the odd classes}\label{pri:ss-triangle}

For a set $S$ of places, an \emph{even class} of exponent $e$ for $S$ is a locally constant cocycle
$z\colon\pgG\to(\auO/\ell^\nu)(-e)$ unramified outside the places above $\ell$, $S$ and $\overline
S$, with $z(g^c)=z(g)$, vanishing on the elements of trivial cyclotomic character in the
decomposition groups above $\ell$.

\begin{definition}[Triangle of primes]\label{pri:def-tri}
Let $J$ be finite totally ordered, $\prity\colon J\to\{1,\dots,t'-1\}$ monotone, $a\in\pgG$, $h_0,
c_{\mathrm{tri}}\in\N$. Degree-one places $w_j\nmid\ell$ with tame generators $\tau_j$ and Frobenius elements
$\phi_j\in G_\nu$ form a \emph{triangle} for $(a,h_0,c_{\mathrm{tri}})$ if: (i) there is a finite set $s$ of
elements of trivial cyclotomic character, with $(\#J+\#s+1)h_0\le c_{\mathrm{tri}}$, detecting (with $a$) the
cocycles unramified outside $\ell$, such that for every $e$ and $g\in s$ some $j$ with $\prity(j)=e$
has $z(\phi_j)=2z(g)$ for all even classes $z$ of exponent $e$ for $\emptyset$; (ii) $\prod_{k\ne e}
(\pgcyc(a)^{-e}-\pgcyc(a)^{-k})\mid\ell^{h_0}$ ($k,e<t'$) and $4(\pgcyc(a)^{-e}-1)^2\mid\ell^{h_0}$
for $0<e<t'$; (iii) if $z$ is an even class of exponent $\prity(j)$ for $\{w_i\}$ with $z(\tau_i)=0$
for $i<j$ and $z(\phi_i)=0$ for $i<j$ with $\prity(i)=\prity(j)$, then $\ell^{h_0}z(\tau_j)=0$.
\end{definition}

\begin{proposition}[Taylor--Wiles places in a triangle; exact form of
Proposition~\ref{pg:triangle}]\label{pri:triangle}
Let $t'\le t$, $h_0\le\pgv(N)$ with $\ell\le h^2$ or $h_0=0$, and $a$ a topological generator of
$\pgG/\ker\pgcyc$ with $\pgcyc(a)^e-1\mid\ell^{h_0}$ for $0<e<t$ (such $a$, $h_0$ exist: below).
Let $\mathcal W$ be a set of Frobenius data such that for every $\nu$ and $g$ there is $\delta$ of
trivial cyclotomic character with $\mathcal W$ containing almost all Frobenius data with Frobenius in
some neighbourhood of every point of a nonempty open subset of a neighbourhood of $\delta^{\ell^\nu}
g$. Then there are $J$, $\prity$, independent of $\nu$, and for every $\nu>0$ places $w_j$ in
$\mathcal W$, above distinct primes of $E^+$, with $\auNm w_j\equiv1\pmod{\ell^\nu}$, forming a
triangle for $(a^{-1},h\,h_0,h^{22}h_0)$.
\end{proposition}

\begin{proof}
\emph{Choice of $a$ and $h_0$}. The image of $\pgcyc$ on $\pgG$ is a closed subgroup of
$\Z_\ell^\times$, which is pro-cyclic for odd $\ell$; for $\ell=2$ it lies in $1+4\Z_2$ (as $E$
contains $\mu_{5040}\supseteq\mu_4$), which is pro-cyclic. So $\pgG/\ker\pgcyc$ is pro-cyclic and has a
topological generator $a$. Every $\pgcyc(g)^e-1$ is a multiple of $\pgcyc(a)^e-1$: write
$\pgcyc(g)=\pgcyc(a)^x$ with $x\in\Z_\ell$ and put $y=\pgcyc(a)^e$. If $\ell$ is odd and
$y\not\equiv1\bmod\ell$, then $y-1$ is a unit. Otherwise ($y\equiv1\bmod\ell$ for odd $\ell$, and
$y\equiv1\bmod4$ for $\ell=2$) one has $\prival(y^x-1)=\prival(y-1)+\prival(x)\ge\prival(y-1)$. If $\ell>h^2$,
then $\ell\nmid q_1d$ (as $q_1d\le h^2$), so an inertia element at $\ell$ whose cyclotomic character
reduces to a generator of $\auF_\ell^\times$ fixes $E\subseteq\Q(\mu_{q_1d})$; since $t<\ell-1$,
$\pgcyc^e-1$ is a unit at it for $0<e<t$, and $h_0=0$. Otherwise a Frobenius element at a chosen place
above $p_i\ne\ell$ has $\pgcyc=p_i$, and $p_i^e-1$ divides $N$ for $0<e\le t\le rt$ (the order of
$\auGL_{rt}(\auF_{p_i})$); take $h_0=\pgv(N)$.

\emph{Separation.} With $u=\pgcyc(a^{-1})$: $u^e-u^k=u^k(u^{e-k}-1)$, a unit times a divisor of
$\ell^{h_0}$, so the product over $k\ne e$, $k<t'$, divides $\ell^{(t'-1)h_0}$; and $4\mid\ell^{h_0}$
unless $h_0=0$, in which case $u^e-1$ is a unit and $\ell$ is odd. As $t+2\le h$ these divide
$\ell^{hh_0}$. This is (ii).

\emph{Test points.} Proposition~\ref{pri:kummer}(b) for $g_0=a^{-1}$ gives $s$, independent of
$\nu$. $J$ is the set of pairs $(e,g)$, $0<e<t'$, $g\in s$, ordered lexicographically with $e$
first, and $\prity(e,g)=e$; so $\#J+\#s+1\le(t+1)(\#s+1)$. If $h_0=0$, $(\#J+\#s+1)hh_0\le h^{22}h_0$
is trivial; if $h_0\ne0$, then $\ell\le h^2$ by hypothesis and the size bound of
Proposition~\ref{pri:kummer}(b) gives $(\#J+\#s+1)hh_0\le h^{22}h_0$.

\emph{Induction}.
Suppose places $w_i$ ($i<j$) with Kummer maps $\kappa_i$ of exponent $\prity(i)$ have been chosen,
the reductions of the $\kappa_i$ with $\prity(i)=e:=\prity(j)$ being linearly independent on $G_\nu$.
By Proposition~\ref{pri:kummer}(a), with $x$ the earlier places of exponent $e$ and $(f_i)$ the
reductions of their Kummer maps, there is a Kummer map $\kappa_0$ of exponent $e$ for $x$, off
$T_j=S_0(x)$, whose reduction is outside the span of the earlier ones. Since the earlier maps of exponent $e$ together with $\kappa_0$ are
additive on $G_\nu$ with independent reductions, they are jointly surjective onto $(\Z/\ell^\nu)^m$
(Nakayama); and conjugation by $a$ followed by a power multiplies a map of exponent $d$ on $G_\nu$
by $\pgcyc(a)^{d}$ (the equivariance), so composing with $\prod_{k\ne e}(\mathrm{conj}_a-\pgcyc(a)^k)$
kills the maps of the other exponents. This gives $u\in G_\nu$ with $\kappa_i(u)=0$ for $i<j$ and
$\kappa_0(u)=\prod_{k\ne e}(\pgcyc(a)^e-\pgcyc(a)^k)$, which divides $\ell^{hh_0}$ by (ii). Let
$g=g_j$ be the test point of $j$ and $\sigma=\delta^{\ell^\nu}\,g\,g^c\,u\,(u^c)^{-1}$. A Kummer map
$\kappa$ of exponent $e$ takes near $\sigma$ the value $2\kappa(u)$ (it is additive on $G_\nu$,
negated by $c$, and $\ell^\nu\kappa(\delta)=0$), and every even class $z$ of exponent $e$ unramified
outside $\ell$ takes near $\sigma$ the value $2z(g)$ ($z$ even: the values at $g,g^c$ add up and
those at $u,u^c$ cancel, as $u\in G_\nu$ acts trivially). By Chebotarev (the hypothesis on $\mathcal
W$) choose a degree-one place $w_j\in\mathcal W$ with Frobenius $\phi_j$ near $\sigma$, above a
prime of $F$ not below the earlier places and not in the $T_i$, $i\le j$. Then $\kappa_i(\phi_j)=0$
for $i<j$, $2\kappa_j(\phi_j)=4\kappa_0(u)=4\prod_{k\ne e}(\pgcyc(a)^e-\pgcyc(a)^k)$, which divides
$\ell^{h_0}\ell^{(t'-1)h_0}\mid\ell^{hh_0}$ in $\Z/\ell^\nu$ (as $4\mid\pgcyc(a)^e-1\mid\ell^{h_0}$),
and $z(\phi_j)=2z(g)$. This gives (i) for $(e,g)$, and $\phi_j\in G_\nu$ gives $\auNm w_j\equiv1\pmod{\ell^\nu}$.

\emph{Property (iii)} . Let $z$ be an even class of exponent $e=\prity(j)$ for the $w_i$,
with $z(\tau_i)=0$ for $i<j$ and $z(\phi_i)=0$ for $i<j$ of exponent $e$. Vanishing at $\tau_i$ and
(by evenness) at $\tau_i^{-c}$ makes $z$ unramified at $w_i,\overline w_i$ for $i<j$ (a tame
generator generates the image of inertia on a module on which inertia acts trivially, and the
cyclotomic character is unramified there). Extend $z$ to $z^+$ on $G_F$ by $z^+(gc)=z(g)$; then
$2z^+$ is a cocycle for $(\pgcyc^+)^{-e}$ with the local conditions of
Proposition~\ref{pri:kummer}(a) at $S_0(x)=T_j$ (at $\ell$ because $g^2\in\pgG$ for $g\in\ker\pgcyc^+$
in a decomposition group, at the earlier places of exponent $e$ because $z$ is unramified there and
vanishes at their Frobenius elements), vanishing at complex conjugations, and unramified outside
$T_j$ and the primes below $w_i$, $i\ge j$ (for $\auO\ne\Z_\ell$ apply this to $f\circ z$ for the
additive maps $f\colon\auO/\ell^\nu\to\Z/\ell^\nu$, which separate points). The reciprocity law of $\kappa_j$ gives $\sum_{i\ge j}\kappa_j(\phi_i)\,
2z(\tau_i)=0$; since $\kappa_j(\phi_i)=0$ for $i>j$ (the later choices), $2\kappa_j(\phi_j)z(
\tau_j)=0$. Values of $z$ at inertia elements at $w_j$ are multiples of $z(\tau_j)$; hence every
multiple of $2\kappa_j(\phi_j)$ in $\Z/\ell^\nu$ kills $z(\tau_j)$, in particular $\ell^{hh_0}$.
\end{proof}

\begin{remark}
The details of the four steps above (choice of $a$ and $h_0$, test points, induction with
$T_j=S_0(x)$, property (iii)) are checked in Proposition~\ref{pku:triangle}. Note that ``even'' here
means invariant under $g\mapsto g^c$ with the local condition at $\ell$ of $\mathrm{Sha}$; this
differs in form from the description in Proposition~\ref{pg:triangle}.
\end{remark}

\begin{lemma}[Cost of a character]\label{pri:cost}
Let the $w_j$ form a triangle for $(a^{-1},h_1,c_1)$ as in Definition~\ref{pri:def-tri}, $0\le e<t'$,
and $\priZ$ a module of cocycles $\pgG\to(\auO/\ell^\nu)(-e)$ unramified outside $\ell$, the $w_j$
and the $\overline w_j$, such that $\pgcyc(a)^{-e}-1$ kills $z(a^{-1})$ for every member unramified
outside $\ell$. Let $b$ be $\prirl(\priZ,Z^{\mathrm{cr}})$ plus the length of $\mathrm{Sha}^-_{e,\nu}$.
Then $\priZ$ has triangle cost $b$ along $\{\prity=e\}$ with constant $\ell^{h_1}$, for $\tau'_j=
\tau_j^{-c}$, $\phi'_j=\phi_j^{-c}$.
\end{lemma}
\begin{proof}
Take $A$ the set on which the local condition at $\ell$ is imposed (elements of trivial cyclotomic
character of the decomposition groups above $\ell$, or inertia above $\ell$ if $e=0$); members of
$\priZ$ obeying the crystalline relations vanish on $A$ (they are combinations of crystalline cocycles
there, which vanish on $A$), so
$\prirl(\priZ,\priZ_0)\le\prirl(\priZ,Z^{\mathrm{cr}})$. Since $\tau_j,\phi_j$ act trivially on
$(\auO/\ell^\nu)(-e)$ (trivial resp.\ $\equiv1\bmod\ell^\nu$ cyclotomic character), $z(g^{-c})=
-z(g^c)$ there, so $z(\tau_j)=z(\tau'_j)$ iff the even part $z^+=z+z\circ c$ vanishes at $\tau_j$,
and likewise for $\phi$. For $e\ne0$, $z^+$ is an even class;
for $e=0$ the condition (a) is vacuous as $\prity>0$. (a) is then (iii) of the triangle applied to
$z^+$. (b): let $\priZ_1$ be as in Definition~\ref{pri:def-cost}. Its members are unramified outside
$\ell$ (vanishing at $\tau_j,\tau'_j$). The map $z\mapsto z\circ c-z$ sends $\priZ_1$ into
$\mathrm{Sha}^-_{e,\nu}$. Let $e\ne0$. If $z\circ c=z$, then $z$ is even, hence an even class of
exponent $e$ for $\emptyset$; as $\phi_j$ acts trivially, $z(\phi'_j)=z(\phi_j^{-c})=-z(\phi_j^c)=
-z(\phi_j)$, so $z(\phi_j)=z(\phi'_j)$ gives $2z(\phi_j)=0$ for $\prity(j)=e$, and by (i) of the
triangle $4z(g)=2z(\phi_j)=0$ for $g\in s$. As $d=\pgcyc(a)^{-e}-1$ kills $z(a^{-1})$, $4dz$ vanishes
at $a^{-1}$ and on $s$, hence $4dz=0$; and $4d\mid4d^2\mid\ell^{h_1}$ by (ii). So the kernel is
$\ell^{h_1}$-torsion and $\prirl(\priZ_1,\priZ_1[\ell^{h_1}])\le\prilen\mathrm{Sha}^-_{e,\nu}$. Let
$e=0$. The coefficients are trivial, so the members of $\priZ_1$ are homomorphisms; they are
unramified outside $\ell$ and vanish on $A$, the inertia groups above $\ell$; so
$\priZ_1\subseteq\mathrm{Sha}^-_{0,\nu}$ (the everywhere unramified homomorphisms,
Proposition~\ref{pri:odd}) and $\prirl(\priZ_1,\priZ_1[\ell^{h_1}])\le\prilen\priZ_1\le\prilen
\mathrm{Sha}^-_{0,\nu}$.
\end{proof}

\subsection{Ribet's element for the Galois representations}\label{pri:ss-columns}

\begin{proposition}[Ribet's element]\label{pri:columns}
Let $c$ be in lift setting for $\vartheta$, auxiliary with ramification at most $\mathrm{ram}$;
let the $w_j$ ($j\in J$) form a triangle for $(a,hh_0,h^{22}h_0)$ with $h_0\le\pgv(N)$ (so $a$ here
is the element written $a^{-1}$ in Proposition~\ref{pri:triangle} and Lemma~\ref{pri:cost}); let
$\rho_x$, $x\in Y$ finite, satisfy the conclusion (c) of Proposition~\ref{pri:BL} with $U=\{w_j\}$;
for $x$ in $\priX\subseteq Y$ let $\pgtr\rho_x=\psi\pgtr R_x+\sum_{e<t'}\chi^e$ with $R_x$ main summands
and $\psi R_x$ having a lattice on which some $\gamma_x$ acts by a scalar $\zeta_x$, $\zeta_x-1$ a
unit; let $\priJ$ be a trace module for the $\rho_x$ with $p=\ell^\nu$, $\tau_{(j,+)}=\tau_j$,
$\tau_{(j,-)}=\tau_j^{-c}$, likewise $\phi$, and $H_j(1)=1$; and at $\pi\in\priX$ let $\pglen
\pgSel_Q(\psi R_\pi,\gamma_\pi;K/\auO)\le\mathrm{sel}<\infty$ (normalised). Then there is a common Ribet
element at $\pi$ with value $d$ satisfying
\[
\pgv(d)+2\pgv(N)\le\mathrm{sel}+2t'\,\mathrm{ram}+r'\pgloc_\ell+e_\ell .
\]
\end{proposition}

\begin{proof}
Choose a large coefficient ring $\auO$ over which: Proposition~\ref{pri:loss}(b) holds; the lengths
of $\mathrm{Sha}^-_{e,\nu}$ are $\le\prival(\pgcodd)$ (Proposition~\ref{pri:odd}); the $H_j$ are
defined; and there are lattices $u_x$ in the $\rho_x$ which, for $x\in\priX$, are block diagonal $T_x
\oplus\bigoplus_{e<t'}\auO(-e)$ with $T_x=\psi u'_x$, $u'_x$ a lattice in $R_x$, $T_x$ spanning
$M_{r'}(\auO)$ (Lemma~\ref{pri:block}, whose hypotheses hold: the $\rho_x$ are semisimple and the $R_x$ satisfy (M4), Remark~\ref{pax:rem-block}), and at $\pi$ the lattice of the Selmer hypothesis. We check the hypotheses of Proposition~\ref{pri:ribet},
with $G=\pgG$, $\xi_e=\chi^e$, $\rho_x=T_x$:

(1), (2): by construction, $\chi(\gamma_x)=1$ since $\gamma_x$ is inertial at a place prime to
$\ell$. (3): by the trace module hypothesis, after doubling the relations by $g\mapsto g^{-c}$ (the
conjugate self-duality turns a relation at $(j,+)$ into one at $(j,-)$ up to multiples of $\ell^\nu
\pgtr$); $\chi^e(\phi_k)\equiv1\bmod\ell^\nu$, so
$H_j(\chi^e(\phi_k))\equiv H_j(1)=1$. (4): $C=\ell^{hh_0}$ by (ii) of the triangle. (6): from (c) of
Proposition~\ref{pri:BL}: there are places covering $Q$ and the places above $p_0p_1$ prime to
$\ell$, with tame generators $\delta_v$ and polynomials $o_v$ with $o_v(\delta_v)(\delta_v-1)=0$ on
every $u_x$ and $\sum_v\pgv(o_v(1))\le2\,\mathrm{ram}+8f\pgv(N)$ (
at $w\in Q$, $o=X-\psi(\delta)$ with $\pgv(o(1))=\pgv(1-\psi(\delta))$, counted for $w$ and $\overline
w$; at a chosen place, $o=o_w=\prod(X-\text{roots of the rescaled label})$ with $o(1)$ dividing $N$ by
the lift setting; at most $8f$ such places). At a chosen place the identity $o_w(\delta_v)(\delta_v-1)=0$
on $u_x$, for every $x\in Y$, is the semisimplicity clause of Proposition~\ref{pri:BL}(c): $\rho_x|_{I_w}$
is semisimple with eigenvalues among the roots of $o_w(X)(X-1)$, the monodromy of
$\operatorname{rec}(\Pi_w)$ being $0$ (paragraph ``No monodromy at a chosen $w$'' of the proof in
\S\ref{pax:ss-BL}); the characteristic polynomial alone would not suffice. (5): Ribet cocycles of $T_\pi$ with values in $K/\auO$
vanishing at the $\tau_k$ and $\gamma_\pi$ are unramified outside $\ell p_0p_1Q\overline Q$ (vanishing
at both $\tau_j$ and $\tau_j^{-c}$ kills inertia at $w_j$ and $\overline w_j$), so modulo $\ell^M$
their length is at most $\prics_M(T_\pi)$ plus the length of the Selmer group; by
Proposition~\ref{pri:loss}(a) and Lemma~\ref{pri:divisible}, $b_\rho=\mathrm{loss}_T+\pglen\pgSel$ with
$\mathrm{loss}_T=\sum_v\delta_v(T_\pi)$. (7): with $s'=\{a\}\cup s\cup\{\tau_k\}$, $\#s'\le2(\#J+\#s+1)$,
Lemma~\ref{pri:cost} gives triangle cost $b_e=\prics_\nu(\auO(-e))+\prilen\mathrm{Sha}^-_{e,\nu}$; the
hypothesis on $z(a^{-1})$ holds because a character cocycle with $c(\gamma_\pi)=0$ (which holds for
those unramified outside $\ell$, as $\gamma_\pi$ is inertial off $\ell$) has $(\chi^e(g)-1)c(a)=0$;
and $s'$ detects (Proposition~\ref{pri:kummer}(b), plus the $\tau_k$ to remove ramification at the
$w_j$).

Proposition~\ref{pri:ribet} gives $d$ with (normalising lengths by $e_\auO$)
\[
\pgv(d)\le\mathrm{loss}_T+\mathrm{sel}+\sum_{e<t'}\bigl(\prics_\nu(\auO(-e))+\pgv(\pgcodd)\bigr)+
t'\bigl(2\mathrm{ram}+8f\pgv(N)+4(\#J+\#s+1)hh_0\bigr)+(\#J+1)hh_0 .
\]
By Proposition~\ref{pri:loss}, $\mathrm{loss}_T+\sum_e\prics_\nu(\auO(-e))\le r'\pgloc_\ell$. With
$(\#J+\#s+1)hh_0\le h^{22}h_0\le h^{22}\pgv(N)$, $t'\le t$ and $h_0\le\pgv(N)$, the remaining terms plus
$2\pgv(N)$ are at most
\[
2t'\mathrm{ram}+t\pgv(\pgcodd)+8ft\pgv(N)+2\pgv(N)+\bigl((4t+1)h^{22}+h\bigr)\pgv(N)=2t'\mathrm{ram}+e_\ell,
\]
which is the definition of $e_\ell$ in~\eqref{pg:eq-beta} .
\end{proof}

\subsection{The operators at the auxiliary places}\label{pri:ss-norm}

Fix places $w_j$ ($j\in J$) off the base level with $\auNm w_j\equiv1\pmod{\ell^\nu}$, the forms of
principal level at the $w_j$, $\Gamma=\prod_j\Gamma_j$, $\Gamma_j=\auGL_n(\auF_{w_j})$,
$\mathrm{Nm}=\sum_{k\in\Gamma}k$, $\mathrm{Nm}_j=\sum_{k\in\Gamma_j}k$, and let $\priT$ be the Hecke
algebra away from the level, acting by integral operators commuting with $\Gamma$. Let $c_0$ be as
in Lemma~\ref{pri:forms}(a). A ring map $\sigma$ from the integral operators is a \emph{reduction
modulo $\ell^\nu/k$} if it kills every $T$ with $kT\in\ell^\nu\cdot(\text{integral operators})$. The
\emph{tied ideal} $\mathfrak a\subseteq\priT[X_J]$ is generated by the $\ell^\nu\pgtr[g]$ and the
$\pgtr[(\gamma_j-1-X_j(\phi_j-1)H_j(\phi_j))g]$, $\gamma_j$ the tame generator and $\phi_j$ a
Frobenius at $w_j$.

\begin{lemma}[One place]
\label{pri:oneplace}
Let $\delta'\in\priT$, $t'\ne0$, and suppose every base system on which $\delta'\ne0$ has at $w=w_j$
a rescaled Hecke polynomial $\prod_{i<t'}(X-(\auNm w)^i)$ times a divisor of $H=H_j$. Let $\sigma$ be a
reduction modulo $\ell^\nu/(2n!)$. There is an integral operator $\partial$ made of translates at $w$
such that for all $g$: $\mathrm{Nm}_j\pgtr[(\gamma-1)g]=0$; $\sigma(\partial\pgtr[(\gamma-1)g])=
\sigma(\mathrm{Nm}_j\pgtr[(\phi-1)H(\phi)g])$; and for each $k\in\Gamma_j$ there is an integral $Z$ with
$\sigma(c_0\partial(k-1))=\sigma(\mathrm{Nm}_jZ)$ and $\delta'\,\mathrm{Nm}\,Z=0$.
\end{lemma}
\begin{proof}
Let $K_1\subseteq\Gamma_j$ be the stabiliser of a line, $\mathrm{ind}_1(k)=\mathrm{ind}(k_{11})$ on
$K_1$, and $\mathrm{Tr}=\sum_hh$, $\mathrm{Tr}'=\sum_hh^{-1}$ over fixed representatives $h$ of
$\Gamma_j/K_1$. For a weight $f\colon\Z/\ell^\nu\to\Z$ put $L(f)=\sum_{k\in K_1}f(\mathrm{ind}_1(k))k$,
and with $p(X)=(X-(\auNm w)^{t'-1})H(X)$ let
\begin{equation}\label{pri:eq-B}
B(f)=p(\tilde U)\,L(f),\qquad\partial=-\mathrm{Tr}\,B(\lambda)\,\mathrm{Tr}',
\end{equation}
$\lambda$ the lift $\Z/\ell^\nu\to\Z$ of least absolute value; $B(f)$ is made of translates at $w$. The first claim is Proposition~\ref{pri:dict}(c).
For the second, modulo $\sigma$: (1) by Proposition~\ref{pri:dict}(b) (decompose under the
characters of $K_1$ and note $\chi(n!)\ne1$ is needed only where the weight is nonzero modulo the
reduction), $B(\lambda)\pgtr[(\gamma-1)g]\equiv L(\lambda)\pgtr[(\gamma-1)g\,p(\phi)]$; (2) the
transfer from $K_1$ to $\Gamma_j$ of $k\mapsto k_{11}$ is $\det$ (an explicit computation on diagonal
matrices and transvections), and for any homomorphism $\alpha$ to
$\Z/\ell^\nu$, $\mathrm{Tr}\,L_\alpha(\lambda)\,\mathrm{Tr}'\equiv L_{\alpha^{\mathrm{tr}}}(\lambda)$
modulo $\ell^\nu$ (the coefficient of $g$ on the left is a sum over representatives of the logarithms
of the components of $g$, congruent to the logarithm of the transfer); so $\mathrm{Tr}\,L(\lambda)\,
\mathrm{Tr}'\equiv L_{\det}(\lambda):=\sum_{k\in\Gamma_j}\lambda(\mathrm{ind}\det k)k$; (3) by
Proposition~\ref{pri:dict}(a), $L_{\det}(\lambda)\pgtr[(\gamma-1)g\,p(\phi)]=(\langle\gamma\rangle-1)
L_{\det}(\lambda)\pgtr[g\,p(\phi)]$; (4) $(\langle\gamma\rangle-1)L_{\det}(\lambda)\equiv-\mathrm{Nm}_j$
modulo $\ell^\nu$, as $\mathrm{ind}\det\langle\gamma\rangle=1$ shifts the weight by $1$ and
$\lambda(x-1)-\lambda(x)\equiv-1$; (5) $p(\phi)=(\phi-(\auNm w)^{t'-1})H(\phi)\equiv(\phi-1)H(\phi)$ as
$\auNm w\equiv1$. Together, $\partial\pgtr[(\gamma-1)g]\equiv\mathrm{Nm}_j\pgtr[(\phi-1)H(\phi)g]$.
For the third claim we use four facts.

(i) \emph{$B(1)$ is fixed by $K_1$ on the left}: $\kappa B(1)=B(1)$ for $\kappa\in K_1$. Indeed the image
of $L(1)$ consists of $K_1$-invariant forms; these are fixed by the pro-$p$ Iwahori subgroup $I_1$
(whose image in $\Gamma_j$ is the upper unipotent group, contained in $K_1$); and $\tilde U$ preserves
the $K_1$-invariant vectors in the $I_1$-invariants, by the argument in the proof of
Lemma~\ref{pax:Ueig} showing that $U$ preserves the vectors of type $(\mathcal P,\tilde\chi(k_{11}))$
(that argument uses only the decomposition of $\mathcal P$ and of $I_1dI_1$, and applies verbatim to
the trivial character). Hence $p(\tilde U)L(1)$ has $K_1$-invariant image. In particular
$L(1)B(1)=\#K_1\cdot B(1)$.

(ii) \emph{Right translation shifts the weight.} As $K_1$ stabilises a line, $k\mapsto k_{11}$ and
hence $\mathrm{ind}_1$ are homomorphisms on $K_1$. For $\kappa\in K_1$ with $a=\mathrm{ind}_1(\kappa)$,
substituting $k'=k\kappa$ gives $L(f)\kappa=\sum_{k'}f(\mathrm{ind}_1(k')-a)k'=L(f(\cdot-a))$.
Since $\lambda(x-a)\equiv\lambda(x)-\lambda(a)\pmod{\ell^\nu}$, the weight $g_a=\ell^{-\nu}(\lambda(\cdot
-a)-\lambda+\lambda(a))$ is $\Z$-valued, and $B(\lambda)(\kappa-1)=-\lambda(a)B(1)+\ell^\nu p(\tilde
U)L(g_a)$, the last term being $\ell^\nu$ times an integral operator, hence killed by $\sigma$.

(iii) \emph{Rewriting $\mathrm{Tr}'(k-1)$.} For each representative $h$, $k^{-1}h$ lies in a unique
coset $h'K_1$; write $h^{-1}k=\kappa_hh'^{-1}$ with $\kappa_h\in K_1$. As $h\mapsto h'$ permutes the
representatives, $\mathrm{Tr}'(k-1)=\sum_h(\kappa_h-1)h'^{-1}$. With $a_h=\mathrm{ind}_1(\kappa_h)$,
(ii) gives
\[
\sigma\bigl(\partial(k-1)\bigr)=\sigma\Bigl(\sum_h\lambda(a_h)\,\mathrm{Tr}\,B(1)\,h'^{-1}\Bigr).
\]

(iv) \emph{Inserting $c_0$.} Lemma~\ref{pri:forms}(a) with $\Gamma'=K_1$ gives an integral $\Psi$ with
$\sum_{\kappa\in K_1}\kappa^{-1}\Psi\kappa=c_0$. By (i), $c_0B(1)=\sum_\kappa\kappa^{-1}\Psi B(1)=
L(1)\Psi B(1)$, and $\mathrm{Tr}\,L(1)=\mathrm{Nm}_j$. Put
\[
Z=\sum_h\lambda(a_h)\,\Psi\,B(1)\,h'^{-1},
\]
an integral operator. Then $\sigma(c_0\partial(k-1))=\sigma(\mathrm{Nm}_jZ)$ by (iii).

It remains to show $\delta'\mathrm{Nm}Z=0$. Write $\mathrm{Nm}=\mathrm{Nm}^{(j)}\mathrm{Nm}_j$, with
$\mathrm{Nm}^{(j)}$ the product of the norms at the other places; it commutes with the operators at
$w$ and with $\delta'$. By the computation in (iv), $\mathrm{Nm}_j\Psi B(1)=\mathrm{Tr}\,L(1)\Psi B(1)=
c_0\,\mathrm{Tr}\,B(1)$. By (i), $\mathrm{Nm}\,B(1)=\mathrm{Nm}^{(j)}\mathrm{Tr}\,L(1)B(1)=\#K_1\cdot
\mathrm{Nm}^{(j)}\mathrm{Tr}\,B(1)$, and $\delta'\mathrm{Nm}B(1)=0$ by Proposition~\ref{pri:dict}(d).
As the operators act on a vector space over a field of characteristic $0$,
$\delta'\mathrm{Nm}^{(j)}\mathrm{Tr}\,B(1)=0$. Hence
$\delta'\mathrm{Nm}Z=c_0\sum_h\lambda(a_h)\,\delta'\mathrm{Nm}^{(j)}\mathrm{Tr}\,B(1)\,h'^{-1}=0$.
(The factor $2n!$ in $\sigma$ is needed only for the second claim.)
\end{proof}

\begin{proposition}[The deformed norm; exact form of Proposition~\ref{pg:norm}]\label{pri:norm}
Let $\delta'\in\priT$, $\delta=c_0^2\delta'$, $t'\ne0$, and suppose the hypothesis of
Lemma~\ref{pri:oneplace} at every $w_j$ (with $H_j$). Let $\sigma$ be a reduction modulo
$\ell^\nu/(2n!)$. There are integral operators $\priN_S$ ($S\subseteq J$), commuting with $\priT$,
and $Y_S$, with $\priN_\emptyset=\mathrm{Nm}$, $Y_\emptyset=\delta^{\#J}$, such that
\begin{enumerate}
\item[(a)] for every $G\in\mathfrak a$, $\sigma\bigl(\sum_S\priN_{J\smallsetminus S}\,G_S\bigr)=0$,
where $G_S$ is the coefficient of $X_S$ in $G$;
\item[(b)] $\sigma(\delta^{\#J}\priN_S)=\sigma(Y_S\,\mathrm{Nm})$ for every $S$.
\end{enumerate}
Consequently: let $\prid\in\priT[X_J]$ with $(\delta_0-n_0)\prid\in\mathfrak a$ for some $\delta_0\in
\priT$ and $n_0\in\pgZb$, let $A\in\priT$ with $A^2=A$, and suppose the hypothesis above holds for
$\delta'=A\delta_0$. Let $l$ be a form of base level with $d_1l$ integral, and $y$ a form with
$A\delta_0l=n_0Ay$, $\delta_0y=n_0y$ and $\prid_Sy=0$ for $S\ne J$. If $2n!(n_0c_0^2)^{\#J+1}$ divides
$\ell^\nu$, then $d_1c_0A\prid_Jy$ is integral.
\end{proposition}

\begin{proof}
For each $j$ let $\partial_j$ be the operator of Lemma~\ref{pri:oneplace}. Operators made of
translates at different places commute, and commute with $\priT$. Let $\priN_S$ be the coefficient of
$X_S$ in $\prod_j(\mathrm{Nm}_j+X_j\partial_j)$, i.e.\ $\priN_S=\prod_{j\in S}\partial_j\prod_{j\notin
S}\mathrm{Nm}_j$; then $\priN_\emptyset=\mathrm{Nm}$, and for $j\notin S$, $\priN_S=\priN'\mathrm{Nm}_j$,
$\priN_{S\cup\{j\}}=\priN'\partial_j$ with the same $\priN'$.

(a) It suffices to treat $G=G'G_0$ with $G_0$ a generator of $\mathfrak a$ and $G'$ arbitrary
(the expression is additive in $G$, and $\mathfrak a$ is spanned by such products). For $G_0=
\ell^\nu\pgtr[g]$, $\sigma$ kills it. For $G_0=t_0-X_jr_0$ with $t_0=\pgtr[(\gamma_j-1)g]$, $r_0=
\pgtr[(\phi_j-1)H_j(\phi_j)g]$: the coefficient of $X_S$ in $G'(t_0-X_jr_0)$ is $G'_St_0$ if $j\notin S$
and $G'_{S}t_0-G'_{S\smallsetminus j}r_0$ if $j\in S$. Pair $S\not\ni j$ with $S\cup\{j\}$: the
complements are $S^c\ni j$ and $S^c\smallsetminus j$, so
\[
\priN_{S^c}G'_St_0+\priN_{S^c\smallsetminus j}(G'_{S\cup j}t_0-G'_Sr_0)=\priN''\bigl(\partial_j
G'_St_0+\mathrm{Nm}_jG'_{S\cup j}t_0-\mathrm{Nm}_jG'_Sr_0\bigr),
\]
and $\mathrm{Nm}_jt_0=0$, $\sigma(\partial_jt_0)=\sigma(\mathrm{Nm}_jr_0)$ (Lemma~\ref{pri:oneplace};
the $G'_S\in\priT$ commute with $\partial_j$). So the sum is $\equiv0$.

(b) is Proposition~\ref{prm:normb}. (The induction must be run for all $S$ simultaneously and at
every place of $S$; the operators $Y_S$ are built from the operator $\Psi$ of
Lemma~\ref{pri:forms}(a), which is not a Hecke operator, so they need not commute with $\priT$, and
this is not used below; see Remark~\ref{prm:rem-norm}.)

\emph{Consequence}. Put $n_1=n_0c_0^2$,
$i_0=\#J+1$ and let $G$ be the top coefficient of $\priN\prid$, i.e.\ $G=\sum_S\priN_{J\smallsetminus
S}\prid_S$. Since $(\delta_0-n_0)\prid\in\mathfrak a$, (a) gives $\sigma(G\delta_0)=\sigma(n_0G)$
after multiplying by the commuting $A$; more precisely $\sigma(\delta G)=\sigma(n_1AG)$ with
$\delta=c_0^2A\delta_0$. Then
\[
\sigma(n_1^{i_0}A G)=\sigma(\delta^{i_0}G)=\sigma\Bigl(\sum_S\delta^{\#J}\priN_{J\smallsetminus
S}\,\delta\prid_S\Bigr)=\sigma\Bigl(\sum_SY_{J\smallsetminus S}\,\mathrm{Nm}\,\delta\prid_S\Bigr).
\]
Apply both sides to the integral form $u=\Psi(d_1l)$, whose norm is $c_0d_1l$
(Lemma~\ref{pri:forms}(a)). Since $\mathrm{Nm}\,\delta\prid_S u=\delta\prid_S(c_0d_1l)=c_0^3d_1
\prid_SA\delta_0l=c_0^3d_1n_0\prid_SAy$, which vanishes for $S\ne J$, only $S=J$ survives, and
$Y_\emptyset=\delta^{\#J}$ acts on $Ay$ as $n_1^{\#J}A$ ($\delta_0y=n_0y$). Hence
$n_1^{i_0}AGu\equiv n_1^{i_0}d_1c_0A\prid_Jy$ modulo $\ell^\nu/(2n!)$ times integral forms; the left
side is $n_1^{i_0}$ times an integral form. As $2n!\,n_1^{i_0}\mid\ell^\nu$, cancelling $n_1^{i_0}$
shows that $d_1c_0A\prid_Jy$ is integral.
\end{proof}


\subsection{Proof of Theorem~\ref{pri:III}}\label{pri:ss-proof}

\begin{lemma}[Orthogonality to the false lifts]\label{pri:weak}
Let $y$ be orthogonal to the false lifts, $c$ in lift setting with $e_c\le e$, and suppose
$\mathcal Z$ is stable under $e_c$ and the Hecke operators at split places off the level. Then
$e_c(y-\pgpr_{\mathcal Z}y)$ lies in the span of the Hecke eigenforms in $e_c\mathcal F(O)$ of base
level that are not weak lifts.
\end{lemma}
\begin{proof}
By Lemma~\ref{pri:chain}, applied to the finitely many weak-lift systems of base level, at all but
finitely many split $w\nmid\ell$ the polynomial $H_w(\cdot;X)$ of each of them vanishes at $\lambda_w$
(the lift root, the largest root of the rescaled chain); fix one such $w$. The Hecke operators are
normal and commute with $e_c$ (which is self-adjoint, being a central idempotent of a unitary group
action); $\mathcal Z^\perp$ is stable under them (the adjoint of a Hecke operator at $w$ is one at
$\overline w$). Let $y'=e_c(y-
\pgpr_{\mathcal Z}y)\in e_c\mathcal F\cap\mathcal Z^\perp$, and $y'_x$ its orthogonal component in a
weak-lift system $x$. Then $y'_x\in e\mathcal F\cap\mathcal Z^\perp\cap\ker H_w(\lambda_w)$ is a false
lift, so $\langle y,y'_x\rangle=0$; also $\langle\pgpr_{\mathcal Z}y,y'_x\rangle=0$; hence
$\|y'_x\|^2=\langle y',y'_x\rangle=\langle y-\pgpr_{\mathcal Z}y,e_cy'_x\rangle=0$.
\end{proof}

\begin{lemma}[Taylor--Wiles places with root chains]\label{pri:TW}
In the setting of Theorem~\ref{pri:III} let $x_0$ be the system of a lift generator, with main
summand $R$. There are $J$, $\prity$, and for every $\nu>0$ places $w_j$ forming a triangle as in
Proposition~\ref{pri:triangle}, off $\pgS(\ell)\cup Q\cup\overline Q$ and the base level, and
polynomials $e_j\in\pgZb[X]$ with $e_j(1)=1$, such that every weak-lift base system $x$ whose
normalised characteristic polynomials are congruent modulo the maximal ideal to those of $x_0$ has,
at $w_j$, rescaled Hecke polynomial $\prod_{i<t'}(X-(\auNm w_j)^i)$ times a divisor of $e_j$.
\end{lemma}
\begin{proof}
Let $\mathcal W$ be the set of Frobenius data $(w,\phi)$ with $w$ off the bad places and the base
level, $\det(1-\psi(\phi)R(\phi))$ a unit, and such that for every weak-lift base system $x$ and every
$S$ with $S(1)$ a unit and $\det(X-\rho_x(\phi))\equiv S\cdot\prod_{i<t'}(X-\chi(\phi)^i)$, the rescaled
Hecke polynomial at $w$ is the chain times a divisor of an integral $e$ with $e(1)=1$. \emph{$\mathcal
W$ satisfies the hypothesis of Proposition~\ref{pri:triangle}}: by (M5) an inertial $\gamma$ above $p_0p_1$
acts on $R$ by a scalar $y$ with $y^{\ell^aj}-1$ a unit for $j\le4r'+2$, and is killed by $\psi$ and
$\pgcyc$; among $\gamma^i\delta^{\ell^\nu}g$, $i\le4r'+2$, one has $\det(1-\psi R)$ a unit on a
neighbourhood (for each eigenvalue $\mu$ of $\psi R(\delta^{\ell^\nu}g)$ at most one $i$ has $y^{i
\ell^\nu}\mu\equiv1$); the other conditions hold for almost all $w$ with Frobenius in a suitable open
set ((G1) in the ``most Frobenius'' form, and Lemma~\ref{pri:chain}: the chain divides the rescaled
Hecke polynomial $\iota^{-1}$-wise, with quotient $T$; then $\det(X-\rho_x(\phi))=\text{chain}\cdot
\iota T\equiv S\cdot\text{chain}$, monic factors cancel modulo the maximal ideal, so $\iota T\equiv S$,
$\iota T(1)$ is a unit and $e=T/T(1)$ works). Apply Proposition~\ref{pri:triangle}. For $x$ congruent
to $x_0$: $\det(X-\rho_{x_0}(\phi))=\det(X-\psi R(\phi))\cdot\prod_{i<t'}(X-\chi(\phi)^i)$ (equality of
traces of all powers, \eqref{pg:eq-trace}); so the condition in $\mathcal W$ applies with
$S=\det(X-\psi R(\phi_j))$. Take $e_j$ the product of the finitely many $e$ obtained.
\end{proof}

\begin{proof}[Proof of Theorem~\ref{pri:III}]
\emph{Step 1: one generator.} Let $v_0\ne0$ be a lift generator with Selmer bound $\mathrm{sel}=
e'_\ell$, of type $c$ and system $x_0$; let $\priX$ be the set of systems containing a nonzero lift
generator (with any Selmer bound), $Y$ the base systems. By Proposition~\ref{pri:BL} the $\rho_x$,
$x\in Y$, have integral characteristic polynomials and satisfy (c). Choose (Proposition~\ref{pri:triangle}
via Lemma~\ref{pri:TW}) $J$, and finitely many $g_k\in\pgG$ and nonzero integral $\mathrm{val}_k$ with
$\prod_k(\mathrm{val}_k-\det(1-\rho_x(g_k)))=0$ at every base system $x$ that is not a weak lift (for
each such $x$ pick $g$ with $\det(1-\rho_x(g))\ne0$). Let $\delta_0=\prod_k(\mathrm{val}_k-\det(1-
\rho(g_k)))\in\priT$, $n_0=\prod_k\mathrm{val}_k$; on $\priX$ each $\det(1-\rho_x(g))=0$ (the trivial
character is a summand), so $\delta_0=n_0$ there. Choose $\nu$ with $2n!(n_0N^2)^{\#J+1}\mid\ell^\nu$,
and places $w_j$ for this $\nu$ by Lemma~\ref{pri:TW}; let $A\in\priT$ be the idempotent of the
systems congruent to $x_0$ (it exists in the Hecke algebra: the product of differences of separating
coefficients). Then every base system with $A\delta_0\ne0$ is a
weak lift congruent to $x_0$, so has the root chain at each $w_j$ with $H_j=e_j$.

Now that the $w_j$ are fixed, replace $Y$ by the finite set of systems of eigenvalues of $\priT$
occurring in base level with principal level at the $w_j$; it contains $\priX$ (forms of base level
are of principal level at the $w_j$), and Proposition~\ref{pri:BL}(c) with $U=\{w_j\}$ applies to it,
as Proposition~\ref{pri:columns} requires. Each $x\in Y$ is a ring map $\priT\to\overline{\Q}_\ell$
(through $\iota$). Let $S$ be the image of $\priT[X_J]$ under $\epsilon\colon\priT[X_J]\to
\operatorname{Map}(Y,\overline{\Q}_\ell[X_J])$, $\epsilon(t)(x)=x(t)$ coefficientwise, and $\priJ=
\epsilon(\mathfrak a)$. $S$ contains the $X_j$ and the characteristic polynomial coefficients (they
lie in $\priT$, being limits of Hecke-polynomial coefficients), and $\priJ$ is an $S$-submodule
containing the families of the definition (the generators of $\mathfrak a$); so $\priJ$ is a trace
module. The map $\epsilon$ is injective: $\priT$ is a commutative algebra generated by normal
operators on a finite-dimensional space, hence simultaneously diagonalisable, so an element acting
by $0$ on every system is $0$. By Proposition~\ref{pri:columns} (with $\mathrm{ram}=e'_\ell$,
$\mathrm{sel}=e'_\ell$) there is a common Ribet element $f\in S$ at $x_0$ with value $d$,
$\pgv(N^2d)\le e'_\ell+2t'e'_\ell+r'\pgloc_\ell+e_\ell$. Write $f=\epsilon(\prid)$ with $\prid\in
\priT[X_J]$, unique by injectivity. For each $g$, $\det(1-\rho(g))\in\priT$ and $f\cdot\det(1-\rho(g))
=\epsilon(G)$ for some $G\in\mathfrak a$; injectivity gives $\prid\det(1-\rho(g))=G\in\mathfrak a$.
Hence
$(\delta_0-n_0)\prid\in\mathfrak a$ ($\delta_0-n_0$ lies in the ideal generated by the
$\det(1-\rho(g_k))$), $\prid_S=0$ on $\priX$ for $S\ne J$, and $\prid_J(x_0)=d$.

Let $y$ be integral and orthogonal to the false lifts, $l=e_cy$ and $y_1=e_c\pgpr_{\mathcal Z}y$. By
Proposition~\ref{pri:BL}(b), $Nl$ is integral; $y_1$ is supported on $\priX$, so $\delta_0y_1=n_0y_1$
and $\prid_Sy_1=0$ for $S\ne J$; and by Lemma~\ref{pri:weak}, $l-y_1$ is supported on non-weak-lift
systems, where $A\delta_0=0$, so $A\delta_0l=A\delta_0y_1=n_0Ay_1$. Proposition~\ref{pri:norm} (with
$c_0\mid N$, so that $2n!(n_0c_0^2)^{\#J+1}$ divides $2n!(n_0N^2)^{\#J+1}$ and hence $\ell^\nu$, and
$d_1=N$) shows that $Nc_0A\prid_Jy_1$ is integral, hence so is $N^2A\prid_Jy_1=(N/c_0)\cdot Nc_0A
\prid_Jy_1$, $N/c_0$ being an integer. Define $G(z)=N^2A\prid_J e_cz$; it acts
on each lift generator $v$ of type $c$ in system $x$ by the scalar $N^2A(x)\prid_J(x)$, on lift
generators killed by $e_c$ by $0$, and $G(\pgpr_{\mathcal Z}y)$ is integral; its eigenvalue at $v_0$
has $\pgv\le\mathrm{bound}:=e'_\ell+2t'e'_\ell+r'\pgloc_\ell+e_\ell$.

\emph{Step 2: all generators at once}. Choose a basis $b_1,\dots,b_m$ of $\mathcal Z$ among lift
generators. Step~1 for $v_0=b_p$ gives an operator $G_p$, diagonal in this basis, with eigenvalues
$l_{pk}$, $\pgv(l_{pp})\le\mathrm{bound}$, and with $G_p\pgpr_{\mathcal Z}y$ integral. The set of
operators with this integrality property is closed under integral linear combinations. For a square
array $(l_{pk})$ there are integral $d_p$ with $|\sum_pl_{pk}d_p|_\ell\ge|l_{kk}|_\ell$ for every $k$
(choose $d$ outside the finitely many hypersurfaces where some linear form $\sum_pl_{pk}X_p$ is smaller
than its largest coefficient; the residue field of $\pgZb$ is infinite). So $\Lambda_0=\sum_pd_pG_p$
is diagonal with eigenvalues $\mu_k$, $\pgv(\mu_k)\le\mathrm{bound}$, and preserves the Hecke
eigenspaces.

\emph{Step 3: the determinant.} Put $\Lambda=N\Lambda_0$. Its eigenvalues have $\pgv\le\mathrm{bound}+
\pgv(N)\le r'\beta_\ell$: indeed, with $r'\ge1$, $t'\le t$, $\mathrm{sel},\mathrm{ram}\le e'_\ell$,
\[
e'_\ell+2t'e'_\ell+r'\pgloc_\ell+e_\ell+\pgv(N)\le r'\bigl(\pgloc_\ell+(1+2t)e'_\ell+e_\ell+\pgv(N)
\bigr)=r'\beta_\ell .
\]
So $\pgv(\det\Lambda)\le\dim\mathcal Z\cdot r'\beta_\ell$. $\Lambda\pgpr_{\mathcal Z}y$ is integral at
$\iota$, and by Lemma~\ref{pri:forms}(b), $\langle y',\Lambda\pgpr_{\mathcal Z}y\rangle=N\langle y',
\Lambda_0\pgpr_{\mathcal Z}y\rangle\in\pgZb$ for $y'$ integral at every isomorphism. If $\mathcal Z=0$
there is nothing to prove.
\end{proof}

\subsection{The inputs of Theorem~\ref{pg:III}}\label{pri:ss-summary}

Theorem~\ref{pri:III}, hence Theorem~\ref{pg:III} (with the parity clause of Remark~\ref{pri:rem-compare}),
is proved from: (i) Proposition~\ref{pri:loss} (the local loss, Proposition~\ref{pg:loss}), proved in
\S\ref{pll:sec} with Theorem~\ref{pll:exponent} (\S\ref{pki:ss-deduce}); (ii) statements (G1)--(G5) of
\S\ref{pg:ss-galois}, proved in \S\ref{pga:sec} from published theorems; (iii) Chebotarev's theorem, the
Greenberg--Wiles formula and the other published results quoted in \S\S\ref{pod:sec}, \ref{pku:sec},
\ref{pax:sec}. The other inputs are proved as follows: Proposition~\ref{pri:ribet} in \S\ref{prm:ss-chars},
Proposition~\ref{pri:odd} as Proposition~\ref{pod:odd}, Proposition~\ref{pri:kummer} in \S\ref{pku:sec},
Proposition~\ref{pri:BL} in \S\ref{pax:ss-BL} and \S\ref{prm:ss-level}, Lemmas~\ref{pri:block},
\ref{pri:forms}, \ref{pri:chain} in \S\ref{pax:sec} and \S\ref{pay:sec}, Proposition~\ref{pri:dict} in
\S\ref{pax:ss-dict} and \S\ref{pay:ss-cons}, Proposition~\ref{pri:norm}(b) with Proposition~\ref{prm:normb},
and Lemma~\ref{pri:amitsur} as Proposition~\ref{pax:amitsur}.



%% ===== proofs/C-localloss.tex =====
% proofs/C-localloss.tex
% The local loss: Proposition pg:loss. Labels of this file start with pll:.
\providecommand{\auO}{\mathcal O}
\providecommand{\auQlb}{\overline{\Q}_\ell}
\providecommand{\panot}[1]{\par\smallskip{\raggedright\noindent\textbf{[Not proved here:} #1\textbf{]}\par}\smallskip}
\providecommand{\pgv}{v}
\providecommand{\pgloc}{\mathrm{loc}}
\providecommand{\pgcyc}{\varepsilon_\ell}
\providecommand{\pglen}{\operatorname{len}}
\providecommand{\pgHom}{\operatorname{Hom}}
% macros of this file (prefix \pll)
\providecommand{\pllZf}{Z_f}
\providecommand{\pllperp}{\perp}
\providecommand{\pllrk}{\operatorname{rk}}
\providecommand{\plldef}{\delta^{\vee}}
\providecommand{\pllB}{B_{\mathrm{cris}}}
\providecommand{\pllD}{\mathrm{dB}}

\section{The local loss}\label{pll:sec}

This section concerns Proposition~\ref{pg:loss}. We first fix exactly what is asserted
(\S\ref{pll:ss-statement}), then prove: the reduction, by local Tate duality, of the defect to a
\emph{dual defect} defined by cup products (\S\ref{pll:ss-duality}); the passage from an exponent
killing cup products to a bound for the dual defect (\S\ref{pll:ss-exponent}); a remark on another route for large
primes, through the theory of Fontaine and Laffaille (\S\ref{pll:ss-FL}); and the sum over places and
summands (\S\ref{pll:ss-sum}); the explicit exponents for the other primes
(\S\ref{pll:ss-gap}) are proved in \S\S\ref{pki:sec}--\ref{pkw:sec}.

The logical structure is as follows. Proposition~\ref{pll:dual} (which uses local Tate duality and
the results of Bloch and Kato), Lemma~\ref{pll:exp} and Lemma~\ref{pll:sum} reduce
Proposition~\ref{pg:loss}, for every prime $\ell$, to the explicit exponent bounds of
Theorem~\ref{pll:exponent}, which are proved in \S\ref{pki:sec} from
\S\S\ref{pkp:sec}--\ref{pkw:sec}; the deduction is carried out in Remark~\ref{pll:status}.

\emph{Conventions.} As in \S\ref{pg:sec}. Fix a prime $\ell$, a place $v\mid\ell$ of $E$, and write
$\Gamma=D_v$, identified with $\operatorname{Gal}(\auQlb/E_v)$, an open subgroup of
$G_{\Q_\ell}=\operatorname{Gal}(\auQlb/\Q_\ell)$ of index $[E_v:\Q_\ell]$. $\auO$ is a coefficient
ring with uniformiser $\varpi$; lengths are normalised ($\pglen\auO/\ell=1$), so that the length of
a finite $\auO$-module $X$ is $\log_\ell\#X/[\auO:\Z_\ell]$. A \emph{lattice} is a finite free $\auO$-module $T$
with a continuous $\auO$-linear action of $\Gamma$; $V=T[1/\ell]$; $T^*=\pgHom_\auO(T,\auO(1))$ is
its \emph{Tate dual}. The Hodge--Tate weights (``jumps'') of a potentially crystalline $V$ are
indexed by the embeddings $\gamma\colon\auO\to\C_\ell$, as in \S\ref{pg:sec}.

\subsection{The statement}\label{pll:ss-statement}

\begin{definition}\label{pll:def}
Let $T$ be a lattice and $\pllB$ Fontaine's crystalline period ring.
\begin{enumerate}
\item The \emph{crystalline cocycles} of $T$ are the continuous maps $c\colon\Gamma\to T$ for which some
$y\in\pllB\otimes_{\Z_\ell}T$ satisfies $1\otimes c(g)=gy-y$ for all $g\in\Gamma$; they form an
$\auO$-module $\pllZf(T)$, which contains the coboundaries $g\mapsto gx-x$, $x\in T$. (These are
the continuous cocycles whose class lies in $H^1_f(\Gamma,T)$, the preimage of the Bloch--Kato
subspace $H^1_f(\Gamma,V)=\ker(H^1(\Gamma,V)\to H^1(\Gamma,\pllB\otimes V))$.)
\item A map $s\colon\Gamma\to T/\ell^M$ \emph{obeys the crystalline relations} of $T$ if every
finite $\auO$-linear relation $\sum_g\lambda_g\,c(g)=0$ ($\lambda_g\in\operatorname{End}_\auO T$)
that holds for all $c\in\pllZf(T)$ also holds for $s$. The reductions modulo $\ell^M$ of the
elements of $\pllZf(T)$ obey them; we write $\overline{\pllZf}(T)_M$ for the module of these
reductions.
\item Let $\mathcal C$ be a class of lattices. $s\colon\Gamma\to T/\ell^M$ \emph{comes from
$\mathcal C$} if $s=\mathsf t\circ a$ with $Y\in\mathcal C$, $a\in\pllZf(Y)$ and $\mathsf t\colon
Y\to T/\ell^M$ $\auO$-linear and $\Gamma$-equivariant. For a submodule $N$ of maps
$\Gamma\to T/\ell^M$ all of whose elements come from $\mathcal C$, its \emph{defect} is the length
of $N/(N\cap\{\text{maps obeying the crystalline relations of }T\})$.
\item A \emph{pairing modulo $\ell^m$} of lattices $Y$, $L$ is a pair $(\psi,c_\psi)$ with
$\psi\in(Y\otimes_\auO L)^*$ and $c_\psi\in\pllZf((Y\otimes L)^*)$ such that $g\psi-\psi=
\ell^mc_\psi(g)$ for all $g$; with a fixed $\Z_\ell$-linear $\mathrm{f}\colon\auO\to\Z_\ell$ for
which $(x,y)\mapsto\mathrm f(xy)$ is a perfect pairing, it gives the $\Gamma$-equivariant pairing
$b_\psi(y,l)=\mathrm f(\psi(y\otimes l))\bmod\ell^m$ with values in $\Z/\ell^m(1)$.
\item The \emph{cup-orthogonal} $L^{\pllperp\mathcal C}$ of $\mathcal C$ in $L$ is the module of
maps $c\colon\Gamma\to L$ such that for every $Y\in\mathcal C$, every $m$, every pairing
$(\psi,c_\psi)$ modulo $\ell^m$ of $Y$ and $L$ and every $a\in\pllZf(Y)$, the $2$-cochain
$(g,h)\mapsto b_\psi(a(g),g\,c(h))$ is a coboundary in $C^2(\Gamma,\Z/\ell^m(1))$. The \emph{dual
defect} is
\[
\plldef_{\mathcal C}(L)=\pglen\bigl(\pllZf(L)/(\pllZf(L)\cap L^{\pllperp\mathcal C})\bigr),
\]
and $\mathcal C$ has \emph{cup exponent} $\kappa$ on $L$ if $\ell^\kappa\pllZf(L)\subseteq L^{\pllperp
\mathcal C}$.
\end{enumerate}
\end{definition}

The classes $\mathcal C_v$ of Proposition~\ref{pg:loss} are classes $\mathcal C(U)$ or
$\mathcal C(U,\mathrm{tame})$: for finite sets $U(\gamma)\subset\Z$, one for each embedding
$\gamma\colon\auO\to\C_\ell$, $\mathcal C(U)$ consists of the crystalline lattices whose jumps at
each $\gamma$ lie in $U(\gamma)$; if $\ell=p_i$, $\mathcal C(U,\mathrm{tame})$ consists of the
lattices that become crystalline over the tamely ramified extension of $E_v$ of degree
$e'=\operatorname{lcm}_{1\le a\le r}(\ell^a-1)$, have jumps in $U$, and such that for every inertial
$s$ both $X^{e'}-1$ and $q_s$ kill $s$ on $D_{\mathrm{pcris}}$ (Definition~\ref{pkw:def-pcris}). Here
$q_s=\prod_{y\in M_s}m_y\in\Z_\ell[X]$, with $m_y$ the minimal polynomial of $y$ over $\Q_\ell$ and
$M_s$ the multiset formed by the eigenvalues of $s$ on the inertial type that (G4)
(Proposition~\ref{pga:G45}) attaches to the label at the place above $\ell$, rescaled by the tame part
of the twist, together with eigenvalues $1$ up to the rank of the representation; $q_s$ is monic of degree at most
$\max(r+t,2r)\,r$ $(=(r+t)r$, as $t=5r)$ and depends only on $s$, not on the lattice
(Proposition~\ref{pkw:G4} and the discussion after it). (Becoming crystalline over the tame extension of
degree $e'$ then follows from the condition on $X^{e'}-1$, Theorem~\ref{pkw:T2}.) In Proposition~\ref{pg:loss}, $U(\gamma)$ is the union of the
jumps at $\gamma$ of the $t'+1$ summands $L_c$, $c\in\{\mathrm{main}\}\cup\{0,\dots,t'-1\}$: the
main summand $L_{\mathrm{main}}$ of rank $r'$ with the jumps $J_\sigma$ of (M1)--(M4) at the complex
embedding $\sigma=\sigma_\gamma$ attached to $\gamma$ through $\iota$ and $v$ (Convention~\ref{pkj:conv},
reading (R$^-$): $\sigma_\gamma=\iota^{-1}\circ g^{-1}\circ e_v$ for any $g$ extending $\gamma$), and the characters $L_e=\chi^e$ with the single jump $e$.

\begin{proposition}[Exact form of Proposition~\ref{pg:loss}]\label{pll:main}
In the situation of Proposition~\ref{pg:loss}, with $\mathcal C_v$ as above and $\auO$ large, let
$N_v(L_c)$ be the module of maps $\Gamma\to L_c/\ell^M$ that come from $\mathcal C_v$. Then its
defect $\delta_v(L_c)$ satisfies (here $N_v(L_c)$ is indeed a module: the classes
$\mathcal C(U)$ and $\mathcal C(U,\mathrm{tame})$ are stable under finite direct sums, since the
jump condition and the conditions on $X^{e'}-1$ and $q_s$ are checked summand by summand, and
$\mathsf t_1\circ a_1+\mathsf t_2\circ a_2=(\mathsf t_1,\mathsf t_2)\circ(a_1,a_2)$ with
$(a_1,a_2)\in\pllZf(Y_1\oplus Y_2)$)
\begin{equation}\label{pll:eq-place}
\delta_v(L_c)\le\plldef_{\mathcal C_v}(L_c^*)\le[E_v:\Q_\ell]\cdot\frac{\pllrk_\auO L_c}{x_c}
\cdot\pllD_\ell,\qquad
\pllD_\ell=10^6\,\frac mk\bigl(b\,\pgv(M)+\lambda_\ell\bigr),
\end{equation}
with $x_{\mathrm{main}}=1$, $x_e=t$ for the characters, and $\lambda_\ell$ as in
\S\ref{pg:ss-statement}. Consequently $\sum_{v\mid\ell}\sum_c\delta_v(L_c)\le r'\pgloc_\ell$.
\end{proposition}

The first inequality in \eqref{pll:eq-place} is Proposition~\ref{pll:dual}; the second is
Theorem~\ref{pll:exponent} with Lemma~\ref{pll:exp} (Remark~\ref{pll:status}); when $\ell\ge4b+2$ and
$\ell\nmid q_1dM$ see also Remark~\ref{pll:FL} for another route; the last assertion is Lemma~\ref{pll:sum}. Note that the statement concerns the
module of \emph{all} such cocycles, not only those that Ribet's method produces. The cocycles that
Ribet's method produces come from $\mathcal C_v$ by Lemma~\ref{pll:ribet} below, so they form a
submodule of $N_v(L_c)$, whose defect is no larger.

\subsection{Reduction to the dual defect}\label{pll:ss-duality}

We use three published facts, for a lattice $T$ over $\Gamma$ and $m\ge1$.
\begin{enumerate}
\item[(D1)] (Local Tate duality \cite[Theorem~7.2.6]{NSW}.) For a finite $\Z_\ell[\Gamma]$-module
$A$, cup product and the invariant map give a perfect pairing $H^1(\Gamma,A)\times
H^1(\Gamma,\pgHom(A,\mu_{\ell^\infty}))\to\Q_\ell/\Z_\ell$. We apply it to $A=T/\ell^m$, whose
dual is identified with $T^*/\ell^m$ through $\mathrm f$; the resulting pairing $H^1(\Gamma,
T/\ell^m)\times H^1(\Gamma,T^*/\ell^m)\to\Z/\ell^m$ is $\langle[s],[c]\rangle=
\operatorname{inv}[(g,h)\mapsto\mathrm f(c(h)\text{ evaluated at }s(g))]$, up to the sign
convention for $g$ acting on the second variable.
\item[(D2)] (Bloch--Kato \cite[Proposition~3.8]{BlochKato}, integral form.) The images of
$H^1_f(\Gamma,T)$ in $H^1(\Gamma,T/\ell^m)$ and of $H^1_f(\Gamma,T^*)$ in $H^1(\Gamma,T^*/\ell^m)$
are exact annihilators of each other under (D1), when $V$ is de Rham. (Bloch and Kato prove that
$H^1_f(V)$ and $H^1_f(V^*)$ are exact annihilators; the integral statement is deduced in
Lemma~\ref{pll:D2} below, after (D3). Compare \cite[Chapter~I, Proposition~4.3]{RubinES}, where the finite-level groups are defined by
propagation from $H^1_f(\Gamma,V)$.)
\item[(D3)] (Bloch--Kato \cite[Corollary~3.8.4]{BlochKato}.)
$\dim_{\Q_\ell}H^1_f(\Gamma,V)=\dim_{\Q_\ell}H^0(\Gamma,V)+\dim_{\Q_\ell}D_{\mathrm{dR}}(V)/
\mathrm{Fil}^0D_{\mathrm{dR}}(V)$ for de Rham $V$.
\end{enumerate}

\begin{lemma}\label{pll:D2}
Statement (D2) holds for every lattice $T$ with $V$ de Rham and every $m\ge1$.
\end{lemma}

\begin{proof}
All lengths in this proof are $\Z_\ell$-lengths; write $d=[E_v:\Q_\ell]$, $n=\pllrk_{\Z_\ell}T$,
$A\subseteq H^1(\Gamma,\bar T)$ and $A^*\subseteq H^1(\Gamma,\bar T^*)$ for the two images, where
$\bar T=T/\ell^m$, $\bar T^*=T^*/\ell^m$.

\emph{Orthogonality.} The pairing (D1) of the reductions of $x\in H^1(\Gamma,T)$ and $y\in
H^1(\Gamma,T^*)$ is the reduction modulo $\ell^m$ of the cup product $x\cup y\in H^2(\Gamma,
\Z_\ell(1))\cong\Z_\ell$, and $x\cup y$ is the image of $(x\cup y)\otimes1\in H^2(\Gamma,\Q_\ell(1))$
in which the rational classes are orthogonal by \cite[Proposition~3.8]{BlochKato}. As
$H^2(\Gamma,\Z_\ell(1))\cong\Z_\ell$ is torsion free, $x\cup y=0$ for $x\in H^1_f(T)$, $y\in
H^1_f(T^*)$. So $A^*\subseteq A^{\pllperp}$, and by (D1) it suffices to prove
$\pglen A+\pglen A^*=\pglen H^1(\Gamma,\bar T)$.

\emph{Length of $A$.} $H^1_f(T)$ is the preimage in the finitely generated module $H^1(\Gamma,T)$
of a subspace of $H^1(\Gamma,V)$; it contains the torsion $H^1(\Gamma,T)_{\rm tors}$, and it is
saturated, so $H^1_f(T)\cap\ell^mH^1(\Gamma,T)=\ell^mH^1_f(T)$. The sequence
$0\to H^1(\Gamma,T)/\ell^m\to H^1(\Gamma,\bar T)$ (from $0\to T\xrightarrow{\ell^m}T\to\bar T\to0$)
therefore gives $A\cong H^1_f(T)/\ell^m$, of length $m\,h_f+\pglen H^1(\Gamma,T)[\ell^m]$, where
$h_f=\dim_{\Q_\ell}H^1_f(\Gamma,V)$ (for a finite group $X$, $X/\ell^m$ and $X[\ell^m]$ have the same
order). The same sequence gives $0\to H^0(\Gamma,T)/\ell^m\to H^0(\Gamma,\bar T)\to
H^1(\Gamma,T)[\ell^m]\to0$, and $H^0(\Gamma,T)$ is free of rank $h_0=\dim H^0(\Gamma,V)$. Hence
$\pglen A=m(h_f-h_0)+\pglen H^0(\Gamma,\bar T)=m\dim_{\Q_\ell}D_{\rm dR}(V)/\mathrm{Fil}^0+
\pglen H^0(\Gamma,\bar T)$ by (D3). Likewise $\pglen A^*=m\dim_{\Q_\ell}D_{\rm dR}(V^*)/\mathrm{Fil}^0
+\pglen H^0(\Gamma,\bar T^*)$.

\emph{Comparison.} Since $V^*=\pgHom(V,\Q_\ell(1))$ (over $\auO$, equivalently over $\Q_\ell$ after
the trace form $\mathrm f$), $\mathrm{Fil}^0D_{\rm dR}(V^*)$ is the annihilator of
$\mathrm{Fil}^0D_{\rm dR}(V)$ under the perfect pairing $D_{\rm dR}(V)\times D_{\rm dR}(V^*)\to
D_{\rm dR}(\Q_\ell(1))$ (whose filtration jump is at $-1$), so the two quotient dimensions add up to
$\dim_{\Q_\ell}D_{\rm dR}(V)=dn$. By the local Euler characteristic formula of Tate \cite[Theorem~7.3.1]{NSW},
$\pglen H^1(\Gamma,\bar T)=\pglen H^0(\Gamma,\bar T)+\pglen H^2(\Gamma,\bar T)+dmn$, and by
\cite[Theorem~7.2.6]{NSW} $H^2(\Gamma,\bar T)$ is dual to $H^0(\Gamma,\bar T^*)$. Adding the two
lengths gives $\pglen A+\pglen A^*=dmn+\pglen H^0(\Gamma,\bar T)+\pglen H^0(\Gamma,\bar T^*)=
\pglen H^1(\Gamma,\bar T)$.
\end{proof}

\begin{proposition}\label{pll:dual}
Let $\mathcal C$ be a class of lattices containing $T$, all of whose members are potentially
crystalline, and let $N$ be a module of maps $\Gamma\to T/\ell^M$ that come from $\mathcal C$. Then
\[
\pglen\bigl(N/(N\cap\overline{\pllZf}(T)_M)\bigr)\le\plldef_{\mathcal C}(T^*).
\]
In particular the defect of $N$ is at most $\plldef_{\mathcal C}(T^*)$.
\end{proposition}

\begin{proof}
Write $m=M$, $\bar T=T/\ell^m$, $\bar T^*=T^*/\ell^m$.

\emph{Step 1: the elements of $N$ are cocycles, and $\overline{\pllZf}(T)_m$ is the module of
cocycles whose class lies in the image $A$ of $H^1_f(T)$.} If $s=\mathsf t\circ a$ then $s$ is a continuous cocycle, as $a$ is one and $\mathsf t$ is equivariant. A coboundary
$g\mapsto g\bar x-\bar x$ of $\bar T$ is the reduction of the coboundary of any lift $x\in T$,
which is crystalline. Hence $\overline{\pllZf}(T)_m$ is exactly the module of cocycles whose class
lies in $A$.

\emph{Step 2: $[N]$ is orthogonal to the image of $H:=\pllZf(T^*)\cap T^{*\pllperp\mathcal C}$.}
Let $s=\mathsf t\circ a\in N$, with $a\in\pllZf(Y)$, $Y\in\mathcal C$. As $Y$ is free, $\mathsf t$
lifts to an $\auO$-linear $q\colon Y\to T$, equivariant modulo $\ell^m$. Put
$\psi(y\otimes\phi)=\phi(q(y))$ on $Y\otimes T^*$. For $g\in\Gamma$,
$(g\psi-\psi)(y\otimes\phi)=\phi(g\,q(g^{-1}y)-q(y))\in\ell^m\auO$, so $g\psi-\psi=\ell^m
c_\psi(g)$ for a unique $c_\psi(g)\in(Y\otimes T^*)^*$; $\ell^mc_\psi$ is the coboundary of $\psi$,
hence crystalline, and $\pllZf$ is saturated in the maps $\Gamma\to(Y\otimes T^*)^*$ (if $\ell^mc$
is crystalline with period $y$, then $c$ is crystalline with period $\ell^{-m}y$ since
$\pllB\otimes$ is torsion-free over $\Z_\ell$ and $\ell$ is invertible in $\pllB$). Thus
$(\psi,c_\psi)$ is a pairing modulo $\ell^m$ of $Y$ and $T^*$, and
$b_\psi(a(g),g\,c(h))=\mathrm f\bigl((g\,c(h))(\mathsf t\,a(g))\bigr)=\mathrm f\bigl((g\,c(h))
(s(g))\bigr)\bmod\ell^m$. For $c\in H$ this cochain is a coboundary by definition of
$T^{*\pllperp\mathcal C}$; it is a cocycle representing the cup product of $[s]$ and $[\bar c]$
under the pairing of (D1). Hence $\langle[s],[\bar c]\rangle=0$.

\emph{Step 3: counting.} Let $A^*\subseteq H^1(\Gamma,\bar T^*)$ be the image of $H^1_f(T^*)$,
i.e.\ of $\pllZf(T^*)$, and $\bar H\subseteq A^*$ the image of $H$. By Step 2,
$[N]\subseteq\bar H^{\pllperp}$; by (D2), $A=A^{*\pllperp}$. By (D1) and the duality of finite
abelian groups, $\bar H^{\pllperp}/A^{*\pllperp}$ is dual to $A^*/\bar H$, hence has the same
order. By Step 1, $N/(N\cap\overline{\pllZf}(T)_m)$ injects into $\bar H^\pllperp/A$. Finally
$\pllZf(T^*)/H$ surjects onto $A^*/\bar H$. Taking normalised lengths (that is, $\log_\ell$ of orders divided by $[\auO:\Z_\ell]$, the same factor on both sides) gives the claim. The last
assertion holds because the elements of $\overline{\pllZf}(T)_m$ obey the crystalline relations.
\end{proof}

The next lemma shows that the cocycles of Ribet's method (cocycles obeying every $\auO$-linear
relation among the families $(u_x(g)-1)_x$) come from $\mathcal C_v$. Both classes
$\mathcal C(U)$ and $\mathcal C(U,\mathrm{tame})$ satisfy its hypotheses: they are stable under
finite direct sums (proof of Proposition~\ref{pll:main}); all members become crystalline over one
finite Galois extension $E'$ of $E_v$ ($E'=E_v$ for $\mathcal C(U)$; for $\mathcal C(U,\mathrm{tame})$,
the Galois closure of the tame extension of degree $e'$, since a representation crystalline over a
finite extension is crystalline over every larger finite extension); membership depends only on $V$; and for
a $\Gamma$-stable sublattice $S\subseteq Y$, $S[1/\ell]$ is a subrepresentation of $Y[1/\ell]$, so
it is crystalline (resp.\ crystalline over the same tame extension), its jumps form a
sub-multiset of those of $Y$, and $D_{\mathrm{pcris}}(S[1/\ell])$ is an $s$-stable subspace of
$D_{\mathrm{pcris}}(Y[1/\ell])$, on which $X^{e'}-1$ and $q_s$ still kill $s$.

\begin{lemma}\label{pll:ribet}
Let $\mathcal C$ be a class of lattices over $\Gamma$, stable under finite direct sums and under
passage to $\Gamma$-stable sublattices, all of whose members become crystalline over a fixed finite
Galois extension $E'/E_v$. Let $u_x$ ($x$ in a finite set $X$) be members of $\mathcal C$, of ranks
$n_x$, and let $z\colon\Gamma\to T/\ell^M$ be a continuous cocycle such that $\sum_i\lambda_i\,
z(g_i)=0$ whenever $\lambda_i\in\auO$, $g_i\in\Gamma$ and $\sum_i\lambda_i(u_x(g_i)-1)=0$ for every
$x$. Then $z$ comes from $\mathcal C$. The same holds if $\Gamma$ is replaced by a group $G$ in which
$\Gamma$ is a subgroup, the $u_x$ and $z$ are defined on $G$ and only the restrictions to $\Gamma$
are asserted to lie in $\mathcal C$ and to come from $\mathcal C$.
\end{lemma}

\begin{proof}
The last sentence follows from the first: relations among elements of $\Gamma$ are relations among
elements of $G$, so the restriction of $z$ to $\Gamma$ satisfies the hypothesis.
Let $Y=\bigoplus_x\operatorname{End}_\auO(u_x)$, with $g\in\Gamma$ acting on the $x$-th summand
by left multiplication by $u_x(g)$; then $Y\cong\bigoplus_xu_x^{\oplus n_x}\in\mathcal C$. Let
$1\in Y$ be the identity in each summand and $a(g)=(u_x(g)-1)_x=g\cdot1-1$, the coboundary of $1$,
so $a\in\pllZf(Y)$. Let $S\subseteq Y$ be the $\auO$-span of the values of $a$. Since
$g\,a(h)=a(gh)-a(g)$, $S$ is $\Gamma$-stable; it is finite free (a submodule of $Y$), so
$S\in\mathcal C$. By hypothesis $\Phi(\sum_i\lambda_ia(g_i)):=\sum_i\lambda_iz(g_i)$ is a
well-defined $\auO$-linear map $\Phi\colon S\to T/\ell^M$, and it is equivariant:
$\Phi(g\,a(h))=z(gh)-z(g)=g\,z(h)=g\,\Phi(a(h))$ by the cocycle relation. Thus $z=\Phi\circ a$,
and it remains to show $a\in\pllZf(S)$, i.e.\ to find $y\in\pllB\otimes_{\Z_\ell}S$ with
$gy-y=a(g)$.

Write $W_\Q=W\otimes_{\Z_\ell}\Q_\ell$, and recall $\pllB\otimes_{\Z_\ell}W=\pllB\otimes_{\Q_\ell}
W_\Q$ for a finite free $W$. The image $\bar1$ of $1$ in $(Y/S)_\Q$ is $\Gamma$-invariant, since
$g1-1\in S$. Put $\Gamma'=\operatorname{Gal}(\auQlb/E')$, a normal subgroup of finite index. The
sequence $0\to S_\Q\to Y_\Q\to(Y/S)_\Q\to0$ consists of representations crystalline over $E'$, and
$D_{\mathrm{cris},E'}=(\pllB\otimes-)^{\Gamma'}$ is exact on such representations (it is left exact
and its value has dimension $\dim_{\Q_\ell}$ of the representation over
$E'_0$), so $(\pllB\otimes Y_\Q)^{\Gamma'}\to(\pllB\otimes(Y/S)_\Q)^{\Gamma'}$ is onto. These are
$\Q_\ell$-vector spaces with an action of the finite group $\Gamma/\Gamma'$, so taking invariants
keeps the map onto. Hence some $y'\in(\pllB\otimes Y_\Q)^\Gamma$ maps to $1\otimes\bar1$. Then
$y=1\otimes1-y'$ maps to $0$ in $\pllB\otimes(Y/S)_\Q$, so $y\in\pllB\otimes S_\Q=\pllB\otimes S$,
and $gy-y=g\cdot1-1=a(g)$.
\end{proof}

\subsection{From an exponent to the dual defect}\label{pll:ss-exponent}

\begin{lemma}\label{pll:exp}
Let $L$ be a potentially crystalline lattice. Then $\pllZf(L)$ is a free $\auO$-module of rank at
most $2[E_v:\Q_\ell]\pllrk_\auO L$. Consequently, if $\mathcal C$ has cup exponent $\kappa$ on $L$, then
$\plldef_{\mathcal C}(L)\le2\kappa\,[E_v:\Q_\ell]\,\pllrk_\auO L$.
\end{lemma}

\begin{proof}
The coboundary map $L\to\pllZf(L)$ has kernel $L^\Gamma$, and $\pllZf(L)/\{\text{coboundaries}\}
=H^1_f(\Gamma,L)\subseteq H^1(\Gamma,L)$, a finitely generated $\Z_\ell$-module
(\cite[Theorem~7.1.8(iii) and Corollary~2.7.10]{NSW}). So $\pllZf(L)$ is finitely generated; it
is torsion free, being a module of $L$-valued maps; hence free. Over $\Z_\ell$ its rank is
$\pllrk L-\pllrk L^\Gamma+\dim H^1_f(\Gamma,V)$, and by (D3) this equals
$\pllrk_{\Z_\ell}L+\dim_{\Q_\ell}D_{\mathrm{dR}}(V)/\mathrm{Fil}^0\le\pllrk_{\Z_\ell}L+
[E_v:\Q_\ell]\pllrk_{\Z_\ell}L$, because $D_{\mathrm{dR}}(V)$ is free of rank $\dim_{\Q_\ell}V$
over $E_v$ (potentially crystalline implies de Rham). Dividing by $[\auO:\Z_\ell]$ gives the rank
bound. If $\ell^\kappa\pllZf(L)\subseteq L^{\pllperp\mathcal C}$, then $\pllZf(L)/(\pllZf(L)\cap
L^{\pllperp\mathcal C})$ is a quotient of $\pllZf(L)/\ell^\kappa\cong(\auO/\ell^\kappa)^{\pllrk}$, of
normalised length $\kappa\cdot\pllrk_\auO\pllZf(L)$.
\end{proof}

\subsection{Large primes: another route}\label{pll:ss-FL}

\begin{remark}\label{pll:FL}
Suppose $\ell\nmid q_1dM$ and that all jumps of all members of $\mathcal C_v$ and of the summands
$L_c$ lie in an interval $[w_0,w_0+\ell-2]$ containing $0$ (this holds if $\ell\ge4b+2$, since all
these jumps lie in $[-2b,2b]$). Then $E_v/\Q_\ell$ is unramified, $\mathcal C_v$ is a crystalline
class $\mathcal C(U)$, and after a twist by $\pgcyc^{-w_0}$ all jumps lie in $[0,\ell-2]$. In this
situation the conclusion $\delta_v(L_c)=0$ is given by Theorem~\ref{pll:exponent}(d) with
Proposition~\ref{pll:dual} (Remark~\ref{pll:status}). One could also try to argue through the
theory of Fontaine and Laffaille \cite{FontaineLaffaille} and \cite[\S4]{BlochKato}; this would
require integral statements on lattices and on subquotients of torsion objects that we do not find
in these sources in the needed form, and the paper does not use this route.
\end{remark}

\subsection{The remaining primes}\label{pll:ss-gap}

For the remaining primes we prove explicit cup exponents by an argument in integral $\ell$-adic
Hodge theory, written in \S\S\ref{pki:sec}--\ref{pkw:sec}. We state precisely what is needed.

\begin{theorem}\label{pll:exponent}
Let $\ell$ be any prime, $v\mid\ell$, $\auO$ large, $\mathcal C_v$ and the summands $L_c$ as in
Proposition~\ref{pll:main}, with $L_c\in\mathcal C_v$. Then $\mathcal C_v$ has, on $L_c^*$, a cup
exponent $k_c$ with $2k_c\le\pllD_\ell/x_c$; or, in the sub-case below where no exponent is
claimed, directly $\plldef_{\mathcal C_v}(L_c^*)\le[E_v:\Q_\ell]\pllrk(L_c)\pllD_\ell/x_c$. In
detail, with $h=r^8$, $Q=kt$ and $\lambda=\lambda_\ell$:
\begin{enumerate}
\item[(a)] (tame, $\ell=p_i$) exponent $3r^2(N+1)^2+150(h^2+4t+2)$, $N=\max(r+t,2r)r$;
\item[(b)] (crystalline, $\ell\le Q$) exponent $10^4\lfloor2\lambda/Q\rfloor(\lfloor\log_\ell
h^3\rfloor+1)$;
\item[(c)] (crystalline, $\ell>Q$, $\lambda\ge h^3$) exponent $10^4\lfloor2\lambda/Q\rfloor m$;
\item[(d)] (crystalline, $\ell>Q$, $\lambda<h^3$) $\plldef=0$ if $\lambda<\lfloor h/5\rfloor$ or
$L_c$ is a character, and in all cases $\plldef_{\mathcal C_v}(L_c^*)\le
[E_v:\Q_\ell]\cdot10^5\bigl(\pllrk(L_c)\lfloor2\lambda/Q\rfloor m+\mu\lambda\bigr)$ with
$\mu=\lfloor\pllrk(L_c)m/(kx_c)\rfloor$.
\end{enumerate}
\end{theorem}

Theorem~\ref{pll:exponent} is proved in \S\ref{pki:ss-deduce} (proof following
Theorem~\ref{pki:crys}) from Propositions~\ref{pki:generic}--\ref{pki:data}, which are proved in
\S\S\ref{pkc:sec}, \ref{pkd:sec} and \ref{pkj:sec} from \S\S\ref{pkp:sec}--\ref{pkw:sec}; the
dependency chain is recorded at the end of \S\ref{pki:ss-arch}. No published theorem is used in
place of this argument (Remark~\ref{pki:rem-published}).

\begin{remark}[Proof of \eqref{pll:eq-place} from Theorem~\ref{pll:exponent}]\label{pll:status}
Given Theorem~\ref{pll:exponent}, \eqref{pll:eq-place} follows: in cases (a)--(c) by
Lemma~\ref{pll:exp} applied to $L=L_c^*$ ($\pllrk L_c^*=\pllrk L_c$) and the inequality
$2k_c\le\pllD_\ell/x_c$; in case (a) this inequality uses $b\,\pgv(M)\ge b$ and the sizes
\eqref{eq:au-sizes} ($r^2(N+1)^2\le49r^6$ and $h^2$ are polynomial in $r$, far below $b\,m/(kt)$); in
cases (b)--(c), $\lfloor\log_\ell h^3\rfloor+1\le m$ because $2\log(4b)<m$, so $2k_c\le2\cdot10^4\cdot(2\lambda/(kt))\,m\le10^6(m/k)\lambda/x_c$ as $x_c\le t$. In case (d), $10^5(\rho\lfloor2\lambda/(kt)\rfloor m+\lfloor\rho
m/(kx)\rfloor\lambda)\le3\cdot10^5\rho m\lambda/(kx)\le\rho\,\pllD_\ell/x$ for $0<x\le t$.
For $\ell>2b$ one has $\ell\nmid(2b)!$ and $\ell>b+30h^3$ (as $30h^3<b$), hence $\lambda_\ell=0$, so if
moreover $\ell\nmid M$, case (d) gives $\plldef=0$: this is the ``no loss'' clause of
Proposition~\ref{pg:loss} (for $\ell>Q$, Theorem~\ref{pkj:chart} gives letters that read
thickness, and Definition~\ref{pki:def-reads} then forces ramification index $e=1$, so
$E_v/\Q_\ell$ is unramified and in particular $\ell\nmid q_1d$). For $\ell\ge4b+2$ prime to $q_1d$
see also Remark~\ref{pll:FL} for another route.
\end{remark}

\subsection{The sum over places and summands}\label{pll:ss-sum}

\begin{lemma}\label{pll:sum}
If \eqref{pll:eq-place} holds for all $v\mid\ell$ and all $c$, then $\sum_{v\mid\ell}\sum_c
\delta_v(L_c)\le r'\pgloc_\ell$.
\end{lemma}

\begin{proof}
$\sum_{v\mid\ell}[E_v:\Q_\ell]=[E:\Q]=2f$. Over the summands, $\sum_c\pllrk(L_c)/x_c=r'+t'/t\le
r'+1\le2r'$. Hence the double sum is at most $2f\cdot2r'\cdot\pllD_\ell=r'\cdot4f\cdot10^6
\frac mk(b\,\pgv(M)+\lambda_\ell)=r'\pgloc_\ell$.
\end{proof}



%% ===== proofs/C-galois.tex =====
% proofs/C-galois.tex
% The inputs of Theorem pg:I (main summands exist): the properties (G1)--(G5) of Galois
% representations, statements (L) and (C) on theta lifts, condition (M5), the two normalisations.
% Labels of this file start with pga:.
\providecommand{\auO}{\mathcal O}
\providecommand{\auQlb}{\overline{\Q}_\ell}
\providecommand{\auG}{\overline G}
\providecommand{\auNm}{\operatorname{Nm}}
\providecommand{\auU}{\mathrm U}
\providecommand{\auGL}{\mathrm{GL}}
\providecommand{\auA}{\mathbb A}
\providecommand{\auF}{\mathbb F}
\providecommand{\aubar}[1]{\overline{#1}}
\providecommand{\auf}[1]{\texttt{#1}}
\providecommand{\pgG}{G_E}
\providecommand{\pgtr}{\operatorname{tr}}
\providecommand{\pgFrob}{\operatorname{Frob}}
\providecommand{\pgcyc}{\varepsilon_\ell}
\providecommand{\pgNch}{N_{\mathrm{ch}}}
% macros of this file (prefix \pga)
\providecommand{\pgarec}{\operatorname{rec}}
\providecommand{\pgaWD}{\operatorname{WD}}
\providecommand{\pgaInd}{\operatorname{Ind}}
\providecommand{\pgaHT}{\operatorname{HT}}
\providecommand{\pgaord}{\operatorname{ord}}

\subsection{The inputs of Theorem~\ref{pg:I}}\label{pga:sec}

This subsection supplies the inputs of Theorem~\ref{pg:I} used in \S\ref{pg:ss-galois}:
(i) the properties (G1)--(G5) of Galois representations, deduced from published theorems whose
hypotheses we check (\S\S\ref{pga:ss-need}--\ref{pga:ss-local}); (ii) statement (L), the Hecke
polynomial of a theta lift, and statement (C), the central action on the lift from $\auU(1)$,
both proved here (\S\S\ref{pga:ss-L}, \ref{pga:ss-C}); (iii) condition (M5) and the two
normalisations of the proof of Theorem~\ref{pg:I} (\S\ref{pga:ss-M}). The final paragraph
(\S\ref{pga:ss-status}) lists the published results on which Theorem~\ref{pg:I} rests.

Notation is that of \S\ref{pg:sec}. Throughout, $\pgFrob_w$ is a \emph{geometric} Frobenius,
$\pgcyc$ the cyclotomic character and $\chi=\pgcyc^{-1}$, so that $\chi(\pgFrob_w)=\auNm w$. For an
algebraic Hecke character $\phi$ of $E$ we write $\phi_\ell$ for its $\ell$-adic Galois character
(through $\iota$), normalised so that $\phi_\ell(\pgFrob_w)=\iota\phi(\varpi_w)$ for $w\nmid\ell$
where $\phi$ is unramified, $\varpi_w$ a uniformiser at $w$ placed in the idele at $w$. Local
class field theory $\mathrm{Art}_w$ sends uniformisers to geometric Frobenius elements.

\subsubsection{What is needed, and in which form}\label{pga:ss-need}

Recall (G1)--(G5) of \S\ref{pg:ss-galois}. We make two of them precise.

\emph{(G1) and (G3), exact form.} Let $V$ be a positive definite standard hermitian space of rank
$n$ and $y$ an algebraic modular form on $\auU(V)$ of weight $\lambda$ that is an eigenform of the
$T_{w,i}$ at all split $w$ off a finite set. We use the Harish-Chandra parameter
$\varsigma=\lambda+\rho$ at each $\tau$, $\rho=(\frac{n-1}2,\frac{n-3}2,\dots,\frac{1-n}2)$, in the
normalisation in which the parameter of a character of $\auU(1)$ of weight $c$ is $c$. (G1) asks
for $\det(X-S_y(\pgFrob_w))=\iota H_w(y;X)$; (G3) asks that at $\sigma$ in the CM type the
Hodge--Tate weights be $\frac{n-1}2-\varsigma_{\sigma,i}$ and at $\aubar\sigma$ be
$\frac{n-1}2+\varsigma_{\sigma,i}$ ($1\le i\le n$), with the sign convention in which $\pgcyc$ has
weight $-1$, that is, $\chi$ has weight $1$. (With $\varsigma=\lambda+\rho$ these are the numbers
$i-1-\lambda_{\sigma,i}$ and $\lambda_{\sigma,i}+n-i$ of (G3).)

\emph{Where each property is used.} (G1): in the trace identity \eqref{pg:eq-trace} and in (M1).
(G2): in (M2) (unramified outside $\ell p_0p_1$, crystalline above $\ell\nmid M$) and in the
Selmer argument of Theorem~\ref{pg:III}. (G3): in (M2) (weights) and in the balance of the
Bloch--Kato conditions. (G4): in (M2) for $\ell=p_i$, (M3), (M4), (M5). (G5): only in
Theorem~\ref{pg:III}, at the places of $Q_\vartheta$.

\emph{Weak forms suffice.} Every use of (G1)--(G5) in Theorems~\ref{pg:I} and \ref{pg:III} is
through one of the following \emph{weak} statements: (G1$'$) $\pgtr S_y(\pgFrob_w)=\iota\,
(\text{coefficient})$ for a set of split $w$ of Dirichlet density one; (G2$'$), (G4$'$), (G5$'$)
``a given polynomial kills $S_y(g)$'' for $g$ a Frobenius, respectively inertial, element (for
instance $(S_y(s)-\psi_w(s))(S_y(s)-1)=0$); plus the crystalline and Hodge--Tate statements above
$\ell$. Indeed (G1$'$) gives \eqref{pg:eq-trace} by Chebotarev and continuity of traces, and the
characteristic polynomial of a semisimple representation is determined by its traces in
characteristic $0$; (M3)--(M5) use only that $R(\gamma_w)$ is killed by the polynomial whose roots
are the images of the roots of the rescaled label (which, $R(\gamma_w)$ having finite order, gives
that $R(\gamma_w)$ is semisimple with eigenvalues among these roots, and the trace identity then
gives the multiplicities). There is no short self-contained proof of the weak forms, since
already (G1$'$) contains the existence of the Galois representations. We therefore derive the
strong forms from the literature and note that they imply the weak ones.

\subsubsection{The global input}\label{pga:ss-global}

\begin{theorem}[Published results]\label{pga:lit}
Let $F$ be a CM field with maximal totally real subfield $F^+$, $\ell$ a prime and
$\iota\colon\auQlb\simeq\C$.
\begin{enumerate}
\item[(BC)] \emph{(Labesse~\cite{Labesse}, with M\oe glin--Waldspurger~\cite[Introduction, Th\'eor\`eme, p.~606]{MW89} for the
form of the discrete constituents.)} Assume $[F^+:\Q]\ge2$ and $F/F^+$ unramified at all finite
places. Let $G=\auU(V)$ for a positive definite
hermitian space $V$ of rank $n$ over $F$, quasi-split at all finite places, and $\pi$ an
automorphic representation of $G(\auA_{F^+})$ whose archimedean components are irreducible
algebraic representations (that is, $\pi_\infty$ is cohomological). Then there are a partition
$n=\sum_i n_im_i$ and cuspidal automorphic representations $\sigma_i$ of $\auGL_{n_i}(\auA_F)$ such
that $\Pi=\boxplus_i\operatorname{Speh}(\sigma_i,m_i)$ is an automorphic representation of
$\auGL_n(\auA_F)$ which is conjugate self-dual, cohomological with the infinitesimal character
determined by $\pi_\infty$, and is the local base change of $\pi_v$ at every place $v$ of $F^+$
split in $F$ and at every finite $v$ at which $G$ and $\pi_v$ are unramified; each
$\operatorname{Speh}(\sigma_i,m_i)$ is conjugate self-dual.
\item[(GR)] \emph{(Harris--Taylor~\cite{HarrisTaylor}, Shin~\cite{Shin}, Chenevier--Harris
\cite{ChenevierHarris}; collected in \cite[Thm.~2.1.1]{BLGGT}.)} Let $\Sigma$ be a regular
algebraic cuspidal automorphic representation of $\auGL_N(\auA_F)$ with $\Sigma^c\simeq\Sigma^\vee
\otimes(\phi\circ\auNm_{F/F^+}\circ\det)$ for an algebraic character $\phi$ of
$\auA_{F^+}^\times/F^{+\times}$ with $\phi_v(-1)$ independent of $v\mid\infty$. Then there is a
continuous semisimple $r_\iota(\Sigma)\colon G_F\to\auGL_N(\auQlb)$ with $r_\iota(\Sigma)^c\simeq
r_\iota(\Sigma)^\vee\pgcyc^{1-N}\phi_\ell$, de Rham above $\ell$ with Hodge--Tate weights given by
the infinitesimal character of $\Sigma_\infty$, and such that for every finite $v\nmid\ell$
\[
\pgaWD\bigl(r_\iota(\Sigma)|_{G_{F_v}}\bigr)^{\mathrm{F\text{-}ss}}\simeq\pgarec\bigl(\Sigma_v\otimes
|\det|_v^{(1-N)/2}\bigr)
\]
(after $\iota$). The equality \emph{including the monodromy operator} is \cite[Thm.~2.1.1(2)]{BLGGT}, stated
there for every regular algebraic cuspidal polarised pair, which rests on
Caraiani~\cite{CaraianiDuke}. For $v\mid\ell$, $r_\iota(\Sigma)|_{G_{F_v}}$ is de Rham, hence
potentially semistable, and for the Weil--Deligne representation of Fontaine's $D_{\mathrm{pst}}$
one has $\iota\pgaWD(r_\iota(\Sigma)|_{G_{F_v}})^{\mathrm{F\text{-}ss}}\prec\pgarec(\Sigma_v\otimes
|\det|_v^{(1-N)/2})$ in the sense of \cite[\S1]{BLGGTlp2}: the two sides have isomorphic
semisimplifications, and for each irreducible $W_{F_v}$-type the partial sums of the sizes of
the Jordan blocks of the monodromy on the left are bounded by those on the right
(Barnet-Lamb--Gee--Geraghty--Taylor \cite[Theorem~A]{BLGGTlp2}, stated there for pairs
$(\Sigma,\phi)$ as here, with no further hypothesis); with equality (``$\simeq$'') when the
weight is Shin-regular (loc.\ cit.), and in general when $\phi=1$ by
Caraiani~\cite[Theorem~1.1]{CaraianiANT}. In particular $r_\iota(\Sigma)$ is crystalline at
$v\mid\ell$ if $\Sigma_v$ is unramified \cite[Thm.~2.1.1(4)]{BLGGT}.
\end{enumerate}
\end{theorem}
These statements are quoted, not proved. In \cite{Labesse} the base change is stated
for a representation $\pi$ of the discrete spectrum with $\pi_\infty$ in the discrete series and
satisfying a sign condition on the archimedean traces of its near-equivalence class, which
holds automatically when $G(F^+\otimes\R)$ is compact (loc.\ cit.); the conclusion there is
$\Pi=\boxplus_i\Pi_i$ with $\Pi_i$ discrete automorphic and conjugate self-dual, and by
M\oe glin--Waldspurger~\cite[Introduction, Th\'eor\`eme, p.~606]{MW89} a discrete automorphic representation of $\auGL_N$ is
$\operatorname{Speh}(\sigma,m)$ with $\sigma$ cuspidal, $\sigma$ being conjugate self-dual when
$\operatorname{Speh}(\sigma,m)$ is. The assumption $[F^+:\Q]\ge2$ is made throughout
\cite{Labesse}. The sign condition is that of \cite[\S2.1]{BLGGT}
(polarized automorphic representations), and \cite[Theorem~2.1.1]{BLGGT} gives the identity at
$v\mid\ell$ only when $\Sigma_v$ has an Iwahori-fixed vector. Existence, the polarization,
de Rham-ness with the Hodge--Tate weights, the identity at $v\nmid\ell$ and crystallinity at
unramified $v\mid\ell$ are all contained in \cite[Theorem~2.1.1(1)--(4)]{BLGGT}; at ramified
$v\mid\ell$ the paper uses only the relation $\prec$ of \cite[Theorem~A]{BLGGTlp2}, see the proof
of Proposition~\ref{pga:G45}.

\begin{proposition}[(G1)--(G3), and no monodromy where the level is tame]\label{pga:G123}
Let $V$, $y$, $\lambda$ be as in \S\ref{pga:ss-need}, let $\pi$ be an irreducible constituent of
the automorphic representation generated by $y$ (on which the $T_{w,i}$ act by the eigenvalues
of $y$), and $\Pi=\boxplus_i\operatorname{Speh}(\sigma_i,m_i)$ as in (BC). Let $s_i\in\frac12\Z$
with $s_i\equiv(n_i-1-n+m_i)/2\pmod1$; then $\sigma_i|\det|^{s_i}$ is regular algebraic and satisfies
the hypothesis of (GR) (Proposition~\ref{pax:sign}). Put
\[
S_y=\bigoplus_i\bigoplus_{a=0}^{m_i-1}r_\iota\bigl(\sigma_i|\det|^{s_i}\bigr)\otimes\pgcyc^{\,e_{i,a}},
\quad e_{i,a}=\tfrac{m_i-1}2-a-s_i+\tfrac{n_i-n}2.
\]
Then $S_y$ satisfies (G1) at every split $w$ at which $y$ is spherical, (G2) and (G3); and for
every finite $v\nmid\ell$ the monodromy of $\pgaWD(S_y|_{G_{E_v}})$ vanishes as soon as the
monodromy of $\pgarec(\Pi_v)$ vanishes.
\end{proposition}
\begin{proof}
$\pi_\infty$ is the irreducible algebraic representation of weight $\lambda$ (definite group), so
(BC) applies; the hypothesis ``quasi-split at all finite places'' holds because $E/E^+$ is
unramified at all finite places and $V$ has a lattice unimodular at every finite place: the
stabiliser of such a lattice is hyperspecial, so $\auU(V)$ is unramified, in particular
quasi-split, at every finite place of $E^+$. Moreover $[E^+:\Q]\ge2$: every quadratic
extension of $\Q$ is ramified at some prime (Minkowski), while $E/E^+$ is unramified at all
finite places. No hypothesis at the places of $E^+$ ramified over $\Q$ is needed in
\cite{Labesse}. Since $\Pi$ is
regular (its infinitesimal character is $\varsigma$, regular), the constituents
$\operatorname{Speh}(\sigma_i,m_i)$ have pairwise disjoint infinitesimal characters, and each
$\sigma_i|\det|^{s_i}$ is regular algebraic (Proposition~\ref{pax:sign}). The Galois representation
of $\operatorname{Speh}(\sigma,m)$ is $\bigoplus_ar_\iota(\sigma|\det|^{s})\pgcyc^{e_a}$, $e_a=
\frac{m-1}2-a-s+\frac{N-n}2$ ($N$ the rank of $\sigma$), because the Satake parameters, respectively
the Langlands parameters, of $\operatorname{Speh}(\sigma,m)_v$ are those of $\sigma_v$ times
$|\cdot|^{(m-1)/2-a}$, $0\le a<m$ (Proposition~\ref{pax:sign}). Hence:

(G1). At a split $w$ at which $y$ is spherical, $\pi_v$ is unramified and $\Pi_w=\pi_v$ under
$\auU(V)(E^+_v)\simeq\auGL_n(E_w)$; the Satake parameters of $\Pi_w$ are the roots of
$H_w(y;X)$ rescaled by $(\auNm w)^{-(n-1)/2}$ (the normalisation of $H_w$: $\sum_i(-1)^i(\auNm
w)^{i(i-1)/2}T_{w,i}X^{n-i}=\prod_k(X-(\auNm w)^{(n-1)/2}\alpha_k)$ for the Satake parameters
$\alpha_k$ of $\Pi_w$, Satake's formula for $\auGL_n$), and $\pgarec(\Pi_w|\det|^{(1-n)/2})$ sends
$\pgFrob_w$ to $\operatorname{diag}((\auNm w)^{(n-1)/2}\alpha_k)$. By (GR) for each constituent
and the remark on Speh representations, $\det(X-S_y(\pgFrob_w))=\iota H_w(y;X)$ exactly, with no further twist: write
$q=\auNm w$, $\sigma=\sigma_i$, $m=m_i$, $s=s_i$, and let $\alpha_\sigma$ be a Satake parameter of
$\sigma_w$. Since $\pgarec(|\cdot|^u)(\pgFrob_w)=q^{-u}$ and $\pgcyc(\pgFrob_w)=q^{-1}$, the summand
$r_\iota(\sigma|\det|^{s})\otimes\pgcyc^{e_{i,a}}$ has the Frobenius eigenvalue
$\alpha_\sigma q^{-s}q^{(n_i-1)/2}q^{-e_{i,a}}=q^{(n-1)/2}\alpha_\sigma q^{a-(m-1)/2}$ by the
value of $e_{i,a}$ in the statement; as $a$ runs over $0,\dots,m-1$ these are $q^{(n-1)/2}$
times the Satake parameters $\alpha_\sigma q^{(m-1)/2-a}$ of $\operatorname{Speh}(\sigma,m)_w$.

(G2). If $y$ is fixed by a hyperspecial subgroup at $v$, $\pi_v$ and hence $\Pi_v$ is unramified
(compatibility of (BC) at unramified places), so each $\sigma_{i,v}$ is unramified (the cuspidal
support of an unramified representation is unramified); by (GR) the Frobenius-semisimplified
Weil--Deligne representation of $S_y$ at $v\nmid\ell$ is unramified with vanishing monodromy, so
$S_y$ is unramified at $v$; at $v\mid\ell$ each summand is crystalline.

(G3). By (GR) the Hodge--Tate weights of $r_\iota(\Sigma)$ at $\sigma$ are the entries of the
infinitesimal character of $\Sigma_\sigma$ shifted by $\frac{N-1}2$; the union over the summands
of $S_y$ gives, at $\sigma$ in the CM type, $\{\frac{n-1}2-\varsigma_{\sigma,i}\}$ and at
$\aubar\sigma$ the numbers $\frac{n-1}2+\varsigma_{\sigma,i}$ (the sign being fixed by the
convention of \S\ref{pga:ss-need}; the infinitesimal character of $\Pi_{\aubar\sigma}$ is the
negative of that of $\Pi_\sigma$ by conjugate self-duality).

Monodromy. $\pgarec(\operatorname{Speh}(\sigma,m)_v)=\bigoplus_a\pgarec(\sigma_v)|\cdot|^{(m-1)/2-a}$
with the monodromy of $\pgarec(\sigma_v)$ on each summand (the local component of a Speh
representation is the Langlands quotient, whose parameter carries no monodromy beyond that of
$\sigma_v$). By (GR), \cite[Thm.~2.1.1(2)]{BLGGT}, the monodromy of $S_y$ at $v$ is the direct sum of those of
the $\pgarec(\sigma_{i,v})$, hence it vanishes if that of $\pgarec(\Pi_v)$ does.
\end{proof}

\begin{remark}[The sign hypothesis]\label{pga:sign}
For $m_i>1$ the twist $|\det|^{(m_i-1)/2}$ need not make $\sigma_i$ algebraic; the correct twist is
$|\det|^{s_i}$ with $s_i\equiv(n_i-1-n+m_i)/2\pmod1$, and the polarisation character is then
$|\cdot|^{2s_i}$, whose sign is $+1$ at every infinite place. So the hypothesis of (GR) always holds;
see Proposition~\ref{pax:sign}.
\end{remark}

\subsubsection{Tame level: (G4) and (G5)}\label{pga:ss-local}

Let $K$ be a finite extension of $\Q_p$ with residue field $\auF_q$ and ring of integers $\auO_K$,
$\ell\ne p$ or $\ell=p$, and fix an injective character $\omega\colon\overline{\auF}_q^\times
\to\auQlb^\times$. For $d\ge1$ let $K_d$ be the unramified extension of degree $d$; the tame
inertia character $\theta_d\colon I_K\to\auF_{q^d}^\times$ is $s\mapsto s(\varpi^{1/(q^d-1)})/
\varpi^{1/(q^d-1)}\bmod\varpi$. A \emph{label} $c$ of degree $n$ over $\auF_q$ (monic, squarefree,
$c(0)\ne0$) factors as $c=\prod_kf_k$ with $f_k$ irreducible of degree $d_k$, roots $\xi_k^{q^a}$;
it names (Green) the irreducible character $\sigma_c$ of $\auGL_n(\auF_q)$ which is, up to sign,
the Deligne--Lusztig character $R_T(\theta)$ for the torus $T=\prod_k\operatorname{Res}\auF_{q^{
d_k}}^\times$ and the character $\theta$ of $T$ named by the $\xi_k$ through $\omega$; it is
Harish-Chandra induced from the cuspidal representation $\otimes_k\sigma_{f_k}$ of the Levi
$\prod_k\auGL_{d_k}(\auF_q)$.

\begin{lemma}[Depth zero]\label{pga:depth0}
Let $\Pi$ be an irreducible smooth representation of $\auGL_n(K)$ whose vectors fixed by the
principal congruence subgroup $K(1)=1+\varpi M_n(\auO_K)$ contain $\sigma_c$, inflated from
$\auGL_n(\auO_K)\to\auGL_n(\auF_q)$, for a label $c$. Then $\Pi\simeq\pi_1\times\dots\times\pi_s$
(irreducible normalised parabolic induction), where $\pi_k$ is a depth-zero supercuspidal
representation of $\auGL_{d_k}(K)$ containing $\sigma_{f_k}$; $\pgarec(\Pi)$ has monodromy $0$;
and $\pgarec(\Pi)|_{I_K}$ is semisimple, factors through the tame quotient, and, on an element
$s\in I_K$ with $\theta_d(s)$ a generator of $\auF_{q^d}^\times$, $d=\operatorname{lcm}(d_k)$, it
has the eigenvalues $\omega(\xi_k^{\,q^a})^{\epsilon}$, $0\le a<d_k$, where the roots $\xi_k$ are
normalised by $\xi_k=\theta_{d}(s)^{e_k}$ through the naming of $\sigma_c$, for a sign $\epsilon\in\{\pm1\}$ depending only on the normalisations of
$\pgarec$, $\mathrm{Art}$ and of Green's parametrisation, not on $c$ or $\Pi$.
\end{lemma}
\begin{proof}
\emph{Cuspidal support.} We prove the following statement, used also elsewhere with labels that
are not squarefree: \emph{if $\Pi^{K(1)}$ contains an irreducible representation $W$ of
$\auGL_n(\auF_q)$ with cuspidal support $(L,\tau)$, $L=\prod_{k=1}^s\auGL_{d_k}$,
$\tau=\otimes_k\sigma_{f_k}$ ($f_k$ irreducible, repetitions allowed), then $\Pi$ has depth zero and
is a subquotient of $\pi_1\times\dots\times\pi_s$, with $\pi_k$ a depth-zero supercuspidal
representation of $\auGL_{d_k}(K)$ containing $\sigma_{f_k}$; so its supercuspidal support is
$\{\pi_1,\dots,\pi_s\}$.} We use two papers of Moy and Prasad: \cite{MoyPrasad94} and
\cite{MoyPrasad}.

\emph{Step 1: $\Pi$ contains a depth-zero minimal $K$-type.} Let $Q=LU$ be the block upper
triangular parabolic subgroup of type $(d_1,\dots,d_s)$. Since $W$ has cuspidal support $(L,\tau)$, it
is a constituent of the Harish-Chandra induced representation $\operatorname{Ind}_{Q(\auF_q)}^{
\auGL_n(\auF_q)}\tau$ (inflated to $Q(\auF_q)$); by Harish-Chandra's theory this does not depend on
the parabolic subgroup with Levi factor $L$, \cite[Proposition~6.1(1)]{MoyPrasad}. Frobenius
reciprocity for the finite group gives $0\ne\operatorname{Hom}_{\auGL_n(\auF_q)}(W,\operatorname{Ind}
\tau)=\operatorname{Hom}_{Q(\auF_q)}(W,\tau)$. As $U(\auF_q)$ acts trivially on $\tau$, a non-zero
$Q(\auF_q)$-map $W\to\tau$ factors through the $U(\auF_q)$-coinvariants of $W$, which, the group being
finite and the coefficients of characteristic $0$, are identified $L(\auF_q)$-equivariantly with the
invariants $W^{U(\auF_q)}$. So $W^{U(\auF_q)}$ contains $\tau$ as an $L(\auF_q)$-subrepresentation
(complete reducibility). Let $\mathcal Q\subseteq\auGL_n(\auO_K)$ be the inverse image of
$Q(\auF_q)$ and $\mathcal Q^+$ the inverse image of $U(\auF_q)$. Then $\mathcal Q$ is a parahoric
subgroup $\mathcal G_x$ (for a point $x$ of the apartment of the diagonal torus), $\mathcal Q^+=
\mathcal G_x^+$ is its pro-unipotent radical, $\mathcal Q/\mathcal Q^+=L(\auF_q)$ is the group of
points of its reductive quotient, and $\Pi^{\mathcal Q^+}=(\Pi^{K(1)})^{U(\auF_q)}\supseteq
W^{U(\auF_q)}$. Hence the $(\mathcal Q,\tau)$-isotypic subspace of $\Pi$ is non-zero. Since $\tau$ is
cuspidal, $(\mathcal Q,\tau)$ is an unrefined minimal $K$-type of depth zero in the sense of
\cite[\S3.4, (i)]{MoyPrasad} (= \cite[Definition~5.1(i)]{MoyPrasad94}). Moreover $\Pi$ has depth zero:
$K(1)=P_{x_0,0+}$ for the hyperspecial vertex $x_0$ fixed by $\auGL_n(\auO_K)$, and $\Pi^{K(1)}\ne0$, so
the depth $\rho(\Pi)$, the smallest number $r\ge0$ with $\Pi^{P_{x,r+}}\neq0$ for some $x$
\cite[Theorem~5.2(1)]{MoyPrasad94} (\cite[Theorem~3.5(1)]{MoyPrasad} in any characteristic), is $0$.

\emph{Step 2: the supercuspidal support.} By \cite[Theorem~6.11(2)]{MoyPrasad}, an irreducible
admissible representation whose $(\mathcal G_x,\sigma)$-isotypic subspace is non-zero, for a minimal
$K$-type $(\mathcal G_x,\sigma)$ of depth zero, is a subquotient of the (unnormalised) induced
representation $I(\mathcal P,\theta)=\operatorname{Ind}_{\mathcal P}^{\mathcal G}\theta$, where
$\mathcal M$ is the Levi subgroup attached to $\mathcal G_x$ in \cite[6.3]{MoyPrasad}, $\mathcal P$ is
a parabolic subgroup with Levi factor $\mathcal M$, $\theta=\operatorname{c-Ind}_{\tilde{\mathcal
M}_x}^{\mathcal M}t$, $\tilde{\mathcal M}_x$ is the normaliser in $\mathcal M$ of the parahoric
$\mathcal M_x=\mathcal M\cap\mathcal G_x$, and $t$ is an irreducible representation of $\tilde{\mathcal
M}_x$ containing $\sigma$ (the set $\mathcal E(\sigma)$ of \cite[6.5]{MoyPrasad}); when $\mathcal G_x$ is
maximal, $\mathcal M=\mathcal G$ and $I(\mathcal P,\theta)=\theta$. We apply it to $(\mathcal
Q,\tau)$ (Step 1; $\Pi$ is admissible). The maximal $\auF_q$-split torus in the centre of $L$ is the
torus of block scalars; its lift $C$ to the diagonal torus has centraliser the block diagonal group,
so $\mathcal M=L(K)=\prod_k\auGL_{d_k}(K)$, $\mathcal M_x=\prod_k\auGL_{d_k}(\auO_K)$ (with
$\mathcal M_x/\mathcal M_x^+\cong\mathcal Q/\mathcal Q^+$, \cite[Proposition~6.4]{MoyPrasad}), and
$\tilde{\mathcal M}_x=\prod_kK^\times\auGL_{d_k}(\auO_K)$. An irreducible representation of this
product is $t=\otimes_kt_k$ with $t_k$ irreducible; $t$ contains $\tau=\otimes_k\sigma_{f_k}$, so each
$t_k$ contains $\sigma_{f_k}$, and $\theta=\otimes_k\theta_k$ with $\theta_k=\operatorname{c-Ind}_{
K^\times\auGL_{d_k}(\auO_K)}^{\auGL_{d_k}(K)}t_k$, an irreducible supercuspidal representation of
$\auGL_{d_k}(K)$ \cite[Proposition~6.6]{MoyPrasad} (applied to $\auGL_{d_k}$ and its maximal parahoric
$\auGL_{d_k}(\auO_K)$) which contains $\sigma_{f_k}$. Normalised and unnormalised induction differ by
$\delta_{\mathcal P}^{1/2}$, whose restriction to $\mathcal M$ is $\otimes_k|\det|^{a_k}$ for real
$a_k$, an unramified character; put $\pi_k=\theta_k\otimes|\det|^{a_k}=\operatorname{c-Ind}(t_k\otimes
|\det|^{a_k})$ (last assertion of \cite[Proposition~6.6]{MoyPrasad}), again a supercuspidal containing
$\sigma_{f_k}$, of depth zero by Step 1's criterion. Then $\Pi$ is a subquotient of
$\pi_1\times\dots\times\pi_s$ (normalised induction, in some order of the blocks; the order does not
change the supercuspidal support), as claimed.

\emph{Depth-zero supercuspidals with a given $\sigma_f$.} By the second assertion of
\cite[Proposition~6.6]{MoyPrasad}, an irreducible representation $\pi$ of $\auGL_d(K)$ containing the
(cuspidal) $\sigma_f$ is $\operatorname{c-Ind}_{K^\times\auGL_d(\auO_K)}^{\auGL_d(K)}\tilde\sigma_f$
for an irreducible $\tilde\sigma_f$ containing $\sigma_f$; as $\varpi$ is central, $\tilde\sigma_f$
restricts irreducibly to $\auGL_d(\auO_K)$, so it is an extension of $\sigma_f$, determined by the
scalar $\tilde\sigma_f(\varpi)$. Hence two such, $\pi$ and $\pi'$ with the same $\sigma_f$, differ by
twisting $\tilde\sigma_f$ by a character $\chi$ of $K^\times\auGL_d(\auO_K)/\auGL_d(\auO_K)\cong
\varpi^{\mathbb Z}$; choosing an unramified character $\mu$ of $K^\times$ with $\mu(\varpi)^d=
\chi(\varpi)$, $\chi$ is the restriction of $\mu\circ\det$, and $\pi'\simeq\pi\otimes(\mu\circ\det)$
by the last assertion of \cite[Proposition~6.6]{MoyPrasad}. Conversely, if $\pi$ contains $\sigma_f$
and $\pi\otimes(\mu\circ\det)$, $\mu$ unramified, contains $\sigma_{f'}$, then $\pi$ contains both
$(\auGL_d(\auO_K),\sigma_f)$ and $(\auGL_d(\auO_K),\sigma_{f'})$ ($\mu\circ\det$ is trivial on
$\auGL_d(\auO_K)$), which are therefore associates \cite[Theorem~5.2, last assertion]{MoyPrasad94}:
there is $g$ with $\auGL_d(\auO_K)\cap g\auGL_d(\auO_K)g^{-1}$ mapping onto both reductive quotients
and $\sigma_f\simeq\sigma_{f'}^g$ \cite[\S3.4, definition of associates, (i)]{MoyPrasad}. The image of
$\auGL_d(\auO_K)\cap g\auGL_d(\auO_K)g^{-1}$ in $\auGL_d(\auF_q)$ is the group of points of a
parabolic subgroup, proper unless $g$ fixes the vertex of $\auGL_d(\auO_K)$ (as in the proof of
\cite[Proposition~6.6]{MoyPrasad}); so $g\in K^\times\auGL_d(\auO_K)$, conjugation by $g$ induces an
inner automorphism of $\auGL_d(\auF_q)$, and $\sigma_f\simeq\sigma_{f'}$, that is $f=f'$. So two
depth-zero supercuspidals containing $\sigma_f$, $\sigma_{f'}$ are inertially equivalent (differ by an
unramified twist) if and only if $f=f'$.

\emph{Irreducibility.} Two supercuspidals $\pi_k$, $\pi_{k'}$ can be linked in Zelevinsky's sense
only if $\pi_{k'}\simeq\pi_k|\det|^{u}$, in particular only if they are inertially equivalent, so
only if $f_k=f_{k'}$, which does not happen because $c$ is squarefree. By Zelevinsky's
irreducibility criterion~\cite[Theorem~4.2]{Zelevinsky} (segments of length one, pairwise unlinked) the
induced representation $\pi_1\times\dots\times\pi_s$ is irreducible, and $\Pi$, which is a
subquotient of it, equals it. Its Langlands parameter is $\bigoplus_k\pgarec(\pi_k)$, with
monodromy $0$.

\emph{Inertia.} For a depth-zero supercuspidal $\pi$ of $\auGL_d(K)$ containing $\sigma_f$, $f$
irreducible with root $\xi$ generating $\auF_{q^d}$ over $\auF_q$ (that is, $\xi$ regular),
$\pgarec(\pi)\simeq\pgaInd_{W_{K_d}}^{W_K}\Lambda$ for a tamely ramified character $\Lambda$ of
$K_d^\times$ whose restriction to $\auO_{K_d}^\times$ is the inflation of the character
$\omega^{\pm1}$ of $\auF_{q^d}^\times$ named by $\xi$, the sign being fixed by the conventions
(the level zero case of the essentially tame correspondence: by
\cite[Corollary~3.3]{BHtame} the naive bijection is corrected by a tamely ramified character of
$K_d^\times$, and in level zero this character is the \emph{unramified} character
$\chi_2^{d-1}$, $\chi_2$ unramified of order $2$, \cite[\S7.5, (7.5.1)]{BHeffective}, which
does not change the restriction to inertia). Restricted to
$I_K=I_{K_d}$, the induced representation is $\bigoplus_{a<d}\Lambda^{q^a}\circ\mathrm{Art}^{-1}$
(Mackey), and $\Lambda\circ\mathrm{Art}^{-1}$ on $I_{K_d}$ factors through $\theta_d$; so the
eigenvalues at $s$ are the images under $\omega$ of the conjugates $\xi^{q^a}$, or uniformly of their
inverses. Summing over $k$ gives the
claim.
\end{proof}

\begin{lemma}[Parahoric level with a tame character]\label{pga:parahoric}
Let $n=n_1+n_2$, $\mathcal P\subseteq\auGL_n(\auO_K)$ the parahoric subgroup of the matrices that
are block lower triangular modulo $\varpi$ for the blocks $(n_1,n_2)$, $\psi$ a non-trivial
character of $\auF_q^\times$, and $\psi_{\mathcal P}(k)=\psi(\det k_{11}\bmod\varpi)$. If an
irreducible smooth $\Pi$ has a non-zero vector on which $\mathcal P$ acts by $\psi_{\mathcal P}$,
then $\Pi\simeq\Pi_1\times\Pi_2$ with $\Pi_1\otimes(\tilde\psi^{-1}\circ\det)$ unramified for an
extension $\tilde\psi$ of $\psi$ to a tame character of $K^\times$ and $\Pi_2$ unramified; the
monodromy of $\pgarec(\Pi)$ vanishes; and every $s\in I_K$ acts on $\pgarec(\Pi)$ by a matrix
$A$ with $(A-\tilde\psi(\mathrm{Art}^{-1}s)^{\epsilon})(A-1)=0$.
\end{lemma}
\begin{proof}
Let $P=MN$ be the upper block parabolic, $\bar N$ the opposite unipotent radical. $\mathcal P$ has
an Iwahori decomposition $\mathcal P=(\mathcal P\cap\bar N)(\mathcal P\cap M)(\mathcal P\cap N)$
with $\mathcal P\cap M=\auGL_{n_1}(\auO_K)\times\auGL_{n_2}(\auO_K)$ and $\psi_{\mathcal P}$ trivial on
the two unipotent parts. Let $\mathcal P_1$ be the pro-unipotent radical of $\mathcal P$ (the
kernel of reduction to $\auGL_{n_1}(\auF_q)\times\auGL_{n_2}(\auF_q)$) and $M_1=\mathcal P_1\cap M$.
We use the following statement (J), which is Moy--Prasad~\cite[Proposition~6.7, p.~116]{MoyPrasad}, and for
the Iwahori subgroup and trivial character is Borel's theorem:
\begin{enumerate}
\item[(J)] for every admissible $\Pi$ (in particular every irreducible smooth $\Pi$, irreducible smooth
representations of $\auGL_n(K)$ being admissible), the projection $\Pi\to\Pi_N$ restricts to an isomorphism
$\Pi^{\mathcal P_1}\to(\Pi_N)^{M_1}$, equivariant for $\mathcal P/\mathcal P_1=\mathcal P\cap M/M_1$.
\end{enumerate}
(The hypotheses of~\cite[6.3, Proposition~6.4, Proposition~6.7]{MoyPrasad}: $\mathcal P=\mathcal G_x$ is a
non-maximal parahoric ($n_1,n_2\ge1$) whose point $x$ lies in the apartment of the diagonal torus $S$; the
maximal $\auF_q$-split torus in the centre of its reductive quotient $\auGL_{n_1}\times\auGL_{n_2}$ is the
torus of block scalars, its lift $C\subset S$ has centraliser $M$, the block diagonal Levi subgroup, and
$P=MN$ may be either of the two opposite parabolics with Levi $M$; $\mathcal M_x=M\cap\mathcal G_x=
\mathcal P\cap M$, and $\mathcal M_x/\mathcal M_x^+\cong\mathcal G_x/\mathcal G_x^+$ by Proposition~6.4(2), so
$\mathcal M_x^+=M_1$ and $\mathcal G_x^+=\mathcal P_1$. Equivariance holds because the projection is
$M$-equivariant.)
(Casselman's notes \cite{CasselmanNotes} give the \emph{surjectivity} of $\Pi^{K_0}\to(\Pi_N)^{K_0\cap M}$
for open $K_0$ with an Iwahori factorisation (Theorem~3.3.3 there) and an isomorphism on the subspace
$\pi(K_0aK_0)\Pi^{K_0}$ (Proposition~4.1.4 there); injectivity on all of $\Pi^{\mathcal P_1}$ is the
additional content of (J).) Taking $(\mathcal P\cap M,\psi\circ\det\otimes1)$-isotypic parts in (J),
the isotypic vector of $\Pi$ gives a non-zero vector of $\Pi_N$ on which
$\auGL_{n_1}(\auO_K)\times\auGL_{n_2}(\auO_K)$ acts by $\psi\circ\det\otimes1$. As taking isotypic
parts for a compact open subgroup is exact and $\Pi_N$ has finite length, some irreducible
subquotient $\tau_1\otimes\tau_2$ of $\Pi_N$ has such a vector, so that $\tau_1\otimes(\tilde\psi^{-1}
\circ\det)$ and $\tau_2$ are spherical. The supercuspidal support of $\Pi$ is that of
$\tau_1\otimes\tau_2$ (every irreducible subquotient of the Jacquet module of an irreducible
representation has the supercuspidal support of that representation, Bernstein--Zelevinsky): $n_1$
characters with restriction $\psi$ to $\auO_K^\times$ and $n_2$ unramified characters. As $\psi\ne1$, no
segment of the first group is linked to one of the second, so by Zelevinsky's classification
$\Pi\simeq\Pi_1\times\Pi_2$ with $\Pi_1$, $\Pi_2$ irreducible with these two supports. By the
geometric lemma of Bernstein--Zelevinsky, the composition factors of the Jacquet module of
$\Pi_1\times\Pi_2$ for the parabolic of type $(n_1,n_2)$ are the representations
$(\alpha_1\times\alpha_2)\otimes(\beta_1\times\beta_2)$, with $\alpha_i\otimes\beta_i$ running over the
composition factors of Jacquet modules of $\Pi_i$; the first factor has supercuspidal support
consisting only of characters with restriction $\psi$ to $\auO_K^\times$ only when it is
$\Pi_1\otimes\Pi_2$ itself. Hence $\tau_1\otimes\tau_2\simeq\Pi_1\otimes\Pi_2$, so that
$\Pi_1\otimes(\tilde\psi^{-1}\circ\det)$ and $\Pi_2$ are spherical. Spherical
representations have parameters with monodromy $0$ and unramified; so $\pgarec(\Pi)=
\pgarec(\Pi_1)\oplus\pgarec(\Pi_2)$ has monodromy $0$, $I_K$ acts trivially on the second summand
and by $\tilde\psi\circ\mathrm{Art}^{-1}$ (to the power $\epsilon$) on the first.
\end{proof}

\begin{proposition}[(G4), (G5)]\label{pga:G45}
Let $y$, $S_y$ be as in Proposition~\ref{pga:G123} and $w$ a split
place, $w\mid p$, $v$ the place of $E^+$ below.
\begin{enumerate}
\item[(G4)] If $y$ is fixed by the principal congruence subgroup at $w$ and lies in the
$\sigma_c$-isotypic part for a label $c$, then $S_y|_{I_w}$ is tame and, if $w\nmid\ell$, of finite
order, semisimple, with the eigenvalues $\omega(\xi^{\epsilon})$, $\xi$ running over the roots of
$c$, on a generator $s$ of tame inertia with $\theta_n(s)$ a generator; if $w\mid\ell$, $S_y$ is
potentially crystalline at $w$, crystalline over $E_w(\varpi^{1/(q^d-1)})$, $d=\operatorname{lcm}_k(d_k)$,
with this inertial type. (Lemma~\ref{pga:depth0} is phrased with $\theta_d(s)$; since the norm
$\auF_{q^n}^\times\to\auF_{q^d}^\times$ is surjective, $\theta_n(s)$ a generator implies
$\theta_d(s)$ a generator.)
\item[(G5)] If $w\nmid\ell$ and $y$ is an eigenvector of the parahoric $\mathcal P$ of type
$(n_W,t)$ with character $\psi_{\mathcal P}$, $\psi=\psi_w\ne1$, then $(S_y(s)-\psi_w(s)^{
\epsilon})(S_y(s)-1)=0$ for $s\in I_w$, $\psi_w$ being viewed on $I_w$ through $\mathrm{Art}_w$.
\end{enumerate}
The sign $\epsilon$ is the same in (G4) and (G5) and for all $w$.
\end{proposition}
\begin{proof}
$\Pi_w=\pi_v$ ((BC) at split places). (G4): Lemma~\ref{pga:depth0} gives the parameter of
$\Pi_w\otimes|\det|^{(1-n)/2}$ (the twist is unramified); its monodromy is $0$, so by
Proposition~\ref{pga:G123} the monodromy of $S_y$ at $w$ is $0$ and $S_y|_{I_w}$ is the
restriction of the parameter, of finite order and semisimple, with the stated eigenvalues. For
$w\mid\ell$ we use only the relation $\prec$ of (GR) at $w$ (\cite[Theorem~A]{BLGGTlp2}, valid for
the polarized pairs $(\sigma_i|\det|^{s_i},|\cdot|^{2s_i})$ of Proposition~\ref{pga:G123}, so no
reduction to the case $\phi=1$ is needed). Let $\Sigma=\sigma_i|\det|^{s_i}$, $\rho=\iota\pgaWD
(r_\iota(\Sigma)|_{G_{E_w}})^{\mathrm{F\text{-}ss}}$ and $\rho'=\pgarec(\Sigma_w\otimes|\det|^{(1-
N)/2})$. Since $\Pi_w$ has a parameter with monodromy $0$, so has $\rho'$ (the parameter of $\Pi_w$
is the sum of those of the $\operatorname{Speh}(\sigma_{i,w},m_i)$, each containing that of
$\sigma_{i,w}$ up to an unramified twist, as in Proposition~\ref{pga:G123}). Thus in every
isotypic part $\rho'[\tau]$ all Jordan blocks have size $1$, and $\rho\prec\rho'$ gives first
$\rho^{\mathrm{ss}}\simeq\rho'^{\mathrm{ss}}$ (so $\rho[\tau]\ne0$ only if $\rho'[\tau]\ne0$) and then
that the largest block of $\rho[\tau]$ has size $\le1$; hence the monodromy of $\rho$ is $0$ and
$\rho\simeq\rho'$. Now $r_\iota(\Sigma)|_{G_{E_w}}$ is de Rham, hence potentially semistable
(Berger), and the monodromy of $D_{\mathrm{pst}}$ is that of $\pgaWD$, which is $0$; a potentially
semistable representation whose $D_{\mathrm{pst}}$ has $N=0$ is potentially crystalline, and it is
crystalline over a finite extension $L/E_w$ exactly when the Weil--Deligne representation is
unramified on $I_L$. The cyclotomic twists in $S_y$ are crystalline. So $S_y|_{G_{E_w}}$ is
potentially crystalline, and its inertial type is the restriction to $I_{E_w}$ of the parameter of
Lemma~\ref{pga:depth0}, which becomes trivial on $I_{E_w(\varpi^{1/(q^{d}-1)})}$. (For $n_i=1$,
$\Sigma$ is an algebraic Hecke character and the monodromy is $0$ for dimension reasons.) (G5):
Lemma~\ref{pga:parahoric} and Proposition~\ref{pga:G123} in the same way.
\end{proof}

\emph{The sign $\epsilon$ in Theorem~\ref{pg:I}.} Replacing every eigenvalue by its inverse does
not affect (M3)--(M5) nor Lemmas~\ref{pg:cyclic}, \ref{pg:burnside}: the inverse of a root of a
primitive label is a root of a primitive label, $\xi\mapsto\xi^{-1}$ preserves orders and the
property ``$\ne1$'', and $z\mapsto z^{-1}$ preserves $\auF_p^\times$. In Theorem~\ref{pg:III} the
sign must agree with the one in the definition of $\psi_\vartheta$ restricted to $I_w$; since
$\psi_\vartheta$ is itself defined through class field theory with the same $\mathrm{Art}_w$, and
(G5) involves the tame part of the same character $\mu_j^0\vartheta$ that defines the parahoric
character, the two signs coincide once $\epsilon$ is fixed, which we do by \emph{defining}
$\omega$ so that $\epsilon=1$.

\subsubsection{Statement (L): the Hecke polynomial of a lift}\label{pga:ss-L}

We state (L) for the Weil representation $\omega_\mu$ of the pair $(\auU(W),\auU(V))$ of
\S\ref{po:ss-set}: $(g,h)$ acts by
$\Psi(x)\mapsto a(\det h)\,b(\det g)^{-1}\cdots\Psi(h^{-1}xg)$ at a split place, where
$a=\chi_W=\upsilon^{r_-}$ is a power of the splitting character $\upsilon$ ($r_-$ the number of
negative signs of $W$) and $b=\chi_V=\mu\,\upsilon^{\lfloor n/2\rfloor}$ with $\mu=\mu^0_j\vartheta$.
Thus $\mu$ enters only through $\det g$, $g\in\auU(W)$, and $\auU(V)$ acts through a splitting
given by a power of $\upsilon$ alone; this is formula \eqref{po:eq-w5}. For an
automorphic character $c$ of $\auU(1)$ unramified at a split $w$ and $N\ge1$ put
\[
c^{(N)}_w=\bigl(c(\varpi_w/\aubar\varpi_w)\,\upsilon(\varpi_w)^{-N}\bigr)(\auNm w)^{N/2}
.
\]
Let $r=\operatorname{rk}W<n=\operatorname{rk}V$, $t=n-r$, $q=\auNm w$, $\alpha_w=\chi_{W,w}^{(r)}$,
$\beta_w=\chi_{V,w}^{(n)}$, and $P_w(X)=\sum_i(-X)^iq^{\binom i2}T_{w,i}=X^{\operatorname{rk}}
H_w(1/X)$ the reversed Hecke polynomial.

\begin{proposition}[(L)]\label{pga:L}
Let $w$ be split with $w,\aubar w$ off the level of the pair, let $y_0$ be a form on $\auU(W)$
spherical at $w$ with $P^W_w(X)y_0=p(X)y_0$, $p\in\C[X]$, and let $y=\Theta(\Phi)y_0$ be its theta
lift for a datum $\Phi$ which is the characteristic function of the integral matrices at $w$. Then
$y$ is spherical at $w$ and
\begin{equation}\label{pga:eq-L}
P_w^V(X)\,y=\prod_{l=0}^{t-1}(1-q^l\alpha_wX)\cdot p\bigl(q^t\alpha_w\beta_w^{-1}X\bigr)\,y .
\end{equation}
Equivalently, if $y\ne0$ and $\gamma_1,\dots,\gamma_r$ are the roots of $H^W_w(y_0;X)$, the roots
of $H^V_w(y;X)$ are $\alpha_wq^l$ ($0\le l<t$) and $q^t\alpha_w\beta_w^{-1}\gamma_k$; this is
\eqref{pg:eq-liftpoly} with $\lambda_w=\alpha_w$ and rescaling factor $q^t\alpha_w/\beta_w$.
Moreover the lift is $\auG_O$-equivariant as in Proposition~\ref{au:equiv}.
\end{proposition}
\begin{proof}
\emph{Coset representatives.} At the split place $w$, $\auU(X)(E^+_v)\simeq\auGL_N(E_w)$, $N=
\operatorname{rk}X$, and on vectors fixed by $K_w=\auGL_N(\auO_w)$ one has
\[
P_w(X)=\sum_{U}(-X)^{\dim U}q^{\binom{\dim U}2}\,R\bigl({}^t\gamma_U\bigr),\qquad
P_{\aubar w}(X)=\sum_U(-X)^{\dim U}q^{\binom{\dim U}2}\,R\bigl(\gamma_U^{-1}\bigr),
\]
the sums over the subspaces $U\subseteq\auF_q^N$, where $\gamma_U\in M_N(\auO_w)$ has
$\det\gamma_U$ of valuation $\dim U$ and the kernel of $\gamma_U\bmod\varpi$ is $U$, and
$\{\gamma_U:\dim U=i\}$ is a set of representatives of $K_w\backslash K_w
\operatorname{diag}(\varpi1_i,1_{N-i})K_w$ (the classical description of the Hecke operator
$T_i$ by lattices $\Lambda'$ with $\varpi\Lambda\subset\Lambda'\subset\Lambda$; the second formula is
the first transported by $g\mapsto{}^t\aubar g^{-1}$, which exchanges $w$ and $\aubar w$).

\emph{Translates of the vacuum.} Let $\mathbf 1$ be the characteristic function of
$M_{n\times r}(\auO_w)$. In the Schr\"odinger model at $w$ an element $h\in\auGL_n(E_w)$ acts by
$\mathbf 1\mapsto\chi_W(\det h)|\det h|^{-r/2}\cdots\mathbf 1(h^{-1}\,\cdot)$ and $g\in
\auGL_r(E_w)$ by $\mathbf 1\mapsto\chi_V(\det g)^{-1}|\det g|^{n/2}\cdots\mathbf 1(\cdot\,g)$ (the
character factors are those of the splitting: $a=\chi_W$ on the $V$ side, $b=\chi_V$ on the $W$
side). For $h={}^t\gamma_U$ this gives
\[
\omega({}^t\gamma_U)\mathbf 1(x)=\alpha_w^{\dim U}\,\mathbf 1\bigl({}^tx\,\gamma_U^{-1}\bigr)
\]
and for $g=\gamma_U^{-1}$ on the $W$ side, $\omega(\gamma_U^{-1})\mathbf 1(x)=\beta_w^{\dim U}
\mathbf 1(x\gamma_U^{-1})$; here $\alpha_w$, $\beta_w$ arise as $\chi(\varpi)q^{N/2}$ from the
character value at $\det\gamma_U$, of valuation $\dim U$, and from $|\det\gamma_U|^{-N/2}$.

\emph{The $q$-binomial identity.} For a matrix $z$ over $E_w$ with $J$ columns, $z\gamma_U^{-1}$ is
integral if and only if $z$ is integral and $U\subseteq\ker(z\bmod\varpi)$. By the $q$-binomial
theorem $\sum_{U\subseteq K}(-Y)^{\dim U}q^{\binom{\dim U}2}=\prod_{l<\dim K}(1-q^lY)$ for a
subspace $K$ of an $\auF_q$-space. Hence
\[
P^V_w(X/\alpha_w)\mathbf 1(x)=\mathbf 1(x)\prod_{l<n-\rho(x)}(1-q^lX),\qquad
P^W_{\aubar w}(Y/\beta_w)\mathbf 1(x)=\mathbf 1(x)\prod_{l<r-\rho(x)}(1-q^lY),
\]
$\rho(x)$ the rank of $x\bmod\varpi$ (for $P^V$ the kernel of ${}^tx\bmod\varpi$ has dimension
$n-\rho$, for $P^W$ that of $x\bmod\varpi$ has dimension $r-\rho$). As $n-\rho=t+(r-\rho)$,
$\prod_{l<n-\rho}(1-q^lX)=\prod_{l<t}(1-q^lX)\prod_{l<r-\rho}(1-q^{l+t}X)$, and therefore
\begin{equation}\label{pga:eq-vac}
P^V_w(X/\alpha_w)\,\mathbf 1=\prod_{l<t}(1-q^lX)\cdot P^W_{\aubar w}(q^tX/\beta_w)\,\mathbf 1 .
\end{equation}
(Both sides are also $K_w\times K_w$-fixed, the characters being unramified.)

\emph{Transfer to the lift.} By \eqref{po:eq-equiv}, $R_h\Theta(\Phi)y_0=\Theta(\omega(1,h)\Phi)
y_0$ and $\Theta(\omega(g,1)\Phi)y_0=\Theta(\Phi)(R_{g^{-1}}y_0)$. The first gives that $y$ is
$K_w$-fixed and $P^V_w(Z)y=\Theta(P^V_w(Z)\Phi)y_0$, with $P^V_w$ acting on the datum at $w$. The
second, summed with the coefficients of $P^W_{\aubar w}$, gives $\Theta(P^W_{\aubar w}(Z)\Phi)y_0
=\Theta(\Phi)(P^{W,*}_{\aubar w}(Z)y_0)$ where $P^{W,*}_{\aubar w}(Z)=\sum_U(\cdots)R_{\gamma_U}$
is the adjoint of $P^W_{\aubar w}(\bar Z)$ for the invariant pairing; on $K_w$-fixed vectors it is
$P^W_w(Z)$, because the adjoint of $T_{\aubar w,i}$ is $T_{w,i}$ (the double coset of
$\operatorname{diag}(\varpi1_i,1)^{-1}$ at $\aubar w$ is that of $\operatorname{diag}(\varpi 1_i,
1)$ at $w$ under the identification, and $\sum_{K\backslash KdK}R_{\gamma^{-1}}$ and
$\sum_{KdK/K}R_{\gamma}$ coincide on $K$-fixed vectors). Applying $\Theta(\cdot)y_0$ to \eqref{pga:eq-vac} with
$X=\alpha_wZ$,
\[
P^V_w(Z)y=\prod_{l<t}(1-q^l\alpha_wZ)\,\Theta(\Phi)\bigl(P^W_w(q^t\alpha_w\beta_w^{-1}Z)y_0\bigr)
=\prod_{l<t}(1-q^l\alpha_wZ)\,p(q^t\alpha_w\beta_w^{-1}Z)\,y .
\]
The statement about roots follows from $P(X)=X^{\deg}H(1/X)$. The $\auG_O$-equivariance is
(L2) of Lemma~\ref{po:facts}.
\end{proof}

\emph{Comparison with Lemma~\ref{po:facts}.} Lemma~\ref{po:facts}(L3) proves that the single
factor $1-q^{t-1}\alpha_wX$ kills the lift, which is enough for the false lifts; it does not give
the factor $p(q^t\alpha_w\beta_w^{-1}X)$, which is what the trace identity needs.
Lemma~\ref{po:facts}(L2) is the equivariance part of (L). Both statements, and the trace identity
below, are for the same representation $\omega_\mu$. (The lifts of the upper chain in
\S\ref{ps:ss-weil} are normalised differently: there the character twists both groups. That
convention is not used for the single lift $W_j\to V_j$.)

\subsubsection{Statement (C): the central action on the lift from $\auU(1)$}\label{pga:ss-C}

Let $O=O_{\nu,z}(\vartheta)$, $W$ of rank one, $V$ of rank $n=2z$, $J_1$ the level of $\auU(W)=
\auU(1)$ at the chosen places, free by Lemma~\ref{au:free}. Write $k(g,h)$ for the theta kernel
($g\in\auU(1)(\auA_f)$, $h\in\auU(V)(\auA_f)$, with values in the weight space), so that the lift
of the constant function is $L(h)=\sum_{c}k(s_c,h)$ over representatives $s_c$ of
$\auU(1)(E^+)\backslash\auU(1)(\auA_{E^+,f})/J_1$ (the level is free, so no stabiliser weights),
and the descent of $\Xi$ is $D\Xi(g)=\langle\Xi,k(g,\cdot)\rangle$ up to conjugation, as in
\eqref{po:eq-descent}. For $z\in\auU(1)(\auA_{E^+,f})$, viewed as the central element $z\cdot1_n$
of $\auU(V)(\auA_f)$, put $\varrho(z)\Xi=\chi_V(z)\,R_z\Xi$, where $\chi_V(z)$ is the value of the
character by which the Weil representation of the scalar pair $(z,z\cdot1_n)$ acts on the datum,
inverted.

\begin{proposition}[(C)]\label{pga:C}
\begin{enumerate}
\item[(a)] $\varrho(z)^{-1}k(g,\cdot)=k(gz,\cdot)$; hence $\varrho(z)L=L$.
\item[(b)] $\langle\varrho(z)\Xi,\Xi'\rangle=\langle\Xi,\varrho(z)^{-1}\Xi'\rangle$, and
$\varrho(z)$ commutes with every Hecke operator at a split place off the level and with the action
of $\auG_O$.
\item[(c)] If $z$ is integral above $\ell$ (and hence a unit there), $\varrho(z)^{\pm1}$ preserves
the forms integral at every $j\colon\C\to\C_\ell$.
\item[(d)] Every class in $\auU(1)(E^+)\backslash\auU(1)(\auA_{E^+,f})/J_1$ has a representative
integral above $\ell$.
\item[(e)] If $\Xi'$ is orthogonal to the false lifts in $\ker D$ and $s_1,\dots,s_{h_1}$ represent
the classes, then $\sum_c\varrho(s_c)^{-1}\Xi'$ is orthogonal to the false lifts in $(\C L)^\perp$.
\end{enumerate}
\end{proposition}
\begin{proof}
(a) The pair $(z,z1_n)\in\auU(W)\times\auU(V)$ acts on $\operatorname{Hom}(W,V)$ by $x\mapsto
z^{-1}xz=x$, so in the Weil representation it acts on every datum by the scalar given by the
splitting characters, which by definition is $\chi_V(z)^{-1}$: $\omega(z,z1_n)\Phi=\chi_V(z)^{-1}
\Phi$. Using the equivariance of the kernel, $k(gz,h)=$ (the kernel for the datum $\omega(z,1)\Phi$)
at $(g,h)$ $=\chi_V(z)^{-1}\cdot$(the kernel for $\omega(1,z^{-1}1_n)\Phi$) $=\chi_V(z)^{-1}k(g,hz^{-1})$,
i.e.\ $k(g z,\cdot)=\varrho(z^{-1})k(g,\cdot)=\varrho(z)^{-1}k(g,\cdot)$ . Summing over $g=s_c$: $\varrho(z)^{-1}L=\sum_ck(s_cz,\cdot)$,
and $c\mapsto$ class of $s_cz$ is a bijection of the class set (right translation by $z$ in the
abelian group $\auU(1)(\auA_f)$ permutes the double cosets; $k(\gamma g u,\cdot)=k(g,\cdot)$ for
$\gamma$ rational and $u\in J_1$ since the datum is $J_1$-invariant up to the character, which is
trivial on $J_1$ for the constant function); so $\varrho(z)^{-1}L=L$.

(b) $R_z$ preserves the inner product (the measure on $\Gamma\backslash G_f/J$ and the stabilisers
are translation invariant, and $z$ is central so $R_z$ preserves the level and the character), and
$|\chi_V(z)|=1$ ($\chi_V$ is unitary on $\auU(1)(\auA_f)$); so $\varrho(z)^*=\overline{\chi_V(z)}
R_{z^{-1}}=\varrho(z^{-1})$. Central translations commute with all translations, hence with Hecke
operators and with $\auG_O$, which acts by translations times characters of $\det$.

(c) $(\varrho(z^{-1})\Xi)(g)=\chi_V(z^{-1})\Xi(gz^{-1})$. If $\Xi$ is integral at $j$,
$j(\Xi(gz^{-1}))\in g_\ell z^{-1}_\ell W_\lambda(\auO_{\C_\ell})=g_\ell W_\lambda(\auO_{\C_\ell})$,
since $z_\ell^{\pm1}\cdot1_n\in\auGL_n(\auO_{E\otimes\Z_\ell})$ preserves $W_\lambda(\auO)$ (the
Schur matrices of integral matrices are integral). And $|j\chi_V(z)|=1$: $\chi_V$ is an algebraic
character of $\auU(1)$, so its value at a finite idele factors as $j\chi_V(z)=\chi_{V,\ell}(z)\cdot
\prod_{v\mid\ell}\theta_v(z_v)$, where $\chi_{V,\ell}$ is the $\ell$-adic avatar of $\chi_V$, a continuous
character of the compact group $\auU(1)(E^+)\backslash\auU(1)(\auA_{E^+,f})$ and hence unit-valued, and
each $\theta_v$ is an algebraic character of $\auU(1)(E^+_v)$, which takes unit values on the units
$z_v$. The same for $z^{-1}$.

(d) $\auU(1)(E^+)$ is dense in $\prod_{v\mid\ell}\auU(1)(E^+_v)$ (weak approximation for the torus
$\auU(1)$, which holds because $\auU(1)$ is a rational variety, being split by the quadratic
extension $E$; Cayley parametrisation). Given $g$, choose $\gamma\in\auU(1)(E^+)$ with $\gamma
g_v\in\auU(1)(\auO_{E^+_v})$ for all $v\mid\ell$ (an open subgroup); $\gamma g$ represents the same
class and is integral above $\ell$.

(e) Let $y$ be a false lift in $(\C L)^\perp$ at $w$: $y\in e\mathcal F(O)$, $H_w(\lambda_w)y=0$,
$y\perp L$. By (b), $\varrho(z)y\in e\mathcal F(O)\cap\ker H_w(\lambda_w)$ for every $z$. Put
$Y=\sum_c\varrho(s_c)y$. Then $D Y=0$: for every $g$, by (a) and (b),
\[
\langle Y,k(g,\cdot)\rangle=\Bigl\langle y,\sum_c\varrho(s_c)^{-1}k(g,\cdot)\Bigr\rangle=
\Bigl\langle y,\sum_ck(gs_c,\cdot)\Bigr\rangle=\langle y,\varrho(g)^{-1}L\rangle=\langle y,L\rangle=0,
\]
because $\sum_ck(gs_c,\cdot)=\varrho(g)^{-1}\sum_ck(s_c,\cdot)=\varrho(g)^{-1}L=L$ by (a) (the
group is abelian), and $y\perp L$. So $Y$ is a false lift in $\ker D$, and $\langle Y,\Xi'\rangle=0$ by
hypothesis; by (b) this is $\langle y,\sum_c\varrho(s_c)^{-1}\Xi'\rangle=0$. As $y$ was arbitrary,
the claim follows.
\end{proof}

This is statement (C) of Lemma~\ref{pg:line} exactly.

\subsubsection{Condition (M5) and the two normalisations}\label{pga:ss-M}

Recall from \S\ref{ss:au-notation} (data (D3), (D7)): $q_2=(q_0d)^{s+6}$ with $18480\mid q_0$,
$2^s=m^{22}$, $p_0,p_1\equiv1\pmod{q_2^2}$, $r\le3x^5$, $\sqrt x\le m\le2\sqrt x$; and $\chi_M$ has
order $q_2$ at every chosen place.

\begin{lemma}\label{pga:q2}
For $1\le j\le m$, $1\le u\le4r+2$, $b\ge0$ and every prime $\ell$: $q_2\nmid j\,u\,\ell^b$.
\end{lemma}
\begin{proof}
Let $p\in\{2,5\}$, $p\ne\ell$; $p\mid18480\mid q_0$, so $p^{s}\mid q_2$. If $q_2\mid ju\ell^b$ then
$p^s\mid ju$. But $ju\le m(4r+2)\le2\sqrt x(12x^5+2)<x^{11}\le m^{22}=2^s\le p^s$ (as $x\ge2^{20}$
by (S1)), a contradiction.
\end{proof}

\begin{proposition}[(M5)]\label{pga:M5}
In the proof of Theorem~\ref{pg:I}, (M5) holds for $R_{j,q}$ and for $R_{\nu,z}$.
\end{proposition}
\begin{proof}
Let $i$ with $p_i\ne\ell$, $p=p_i$, $w\mid p$ a chosen place whose label is primitive (case (a))
or $X-1$ (case (b)). By (G4) (Proposition~\ref{pga:G45}, $\epsilon=1$) and the twist by $\eta_j$,
$R(\gamma_w)$ is semisimple with eigenvalues $\omega(\xi')$, $\xi'$ running over the roots of the
rescaled label, that is, the roots $\xi$ of $c$ multiplied by the element $z\in\auF_p^\times$
naming the tame part of $\eta_j$ at $w$; this tame part is a power of $\chi_M$ (the twist being
trivial at the chosen places), so $z^{q_2}=1$.

(a) Let $\xi$ be a root of $c$; its order is $p^d-1$, $d$ the degree of the label. Put $e=(p^d-1)/(p-1)$ and $y_0=
\auNm_{\auF_{p^d}/\auF_p}(\xi z)=\xi^ez^d\in\auF_p^\times$. Every root $\xi'=(\xi z)^{p^a}$
satisfies $\xi'^e=y_0^{p^a}=y_0$. So $s=\gamma_w^e$ acts by the scalar $y=\omega(y_0)$. \emph{The
order of $y_0$ is divisible by $q_2^2$}: for a prime $p'\mid q_2$ let $p'^{\alpha}\|p-1$; then
$\alpha\ge2v_{p'}(q_2)$. The $p'$-primary component of $\xi^e$ (which has order $p-1$) has order
$p'^\alpha$, that of $z^d$ has order at most $p'^{v_{p'}(q_2)}<p'^\alpha$; in a cyclic $p'$-group
the product of an element of order $p'^\alpha$ and one of smaller order has order $p'^\alpha$. So
$q_2^2\mid\pgaord(y_0)$.

(b) The only root is $\xi'=z_w$, the element naming $\chi_M^{-j}$ at $w$ (the twist is trivial at
the chosen places); $s=\gamma_w$, $y=\omega(z_w)$, and $\pgaord(z_w)=q_2/\gcd(q_2,j)$.

In both cases: if $y_0^x=1$ then $q_2\mid jx$ (in (a) even $q_2^2\mid x$). Now
$y^{\ell^au}-1$ is a non-unit if and only if $y^{\ell^au}$ has $\ell$-power order, i.e.\ if and only
if $y_0^{\ell^{a+b}u}=1$ for some $b\ge0$; that would give $q_2\mid ju\ell^{a+b}$, which
Lemma~\ref{pga:q2} excludes for $1\le u\le4r+2$ ($\ge4n+2$ as $n\le r$). So $y^{\ell^au}-1$ is a
unit for all $a\ge0$, $1\le u\le4n+2$.
\end{proof}

Thus (M5) needs no further property of admissible data; it follows from $p_i\equiv1\pmod{q_2^2}$, $q_2=(q_0d)^{s+6}$
and $2^s=m^{22}$.

\begin{proposition}[First normalisation: the trace identity]\label{pga:trace}
With $\alpha_w,\beta_w$ as in Proposition~\ref{pga:L} for the pair $O_j(\vartheta)$, $n=r+t$:
\[
q^t\beta_w^{-1}=\psi_\vartheta(\pgFrob_w)\,\eta_j(\pgFrob_w),\qquad
\eta_j(y)=(\mu_j^0)^{-1}(y/\aubar y)\,|y|^{(r-t)/2},
\]
and $w\mapsto\alpha_w^{-1}$ is the Frobenius value of the $\ell$-adic character $\eta_O$ of the
algebraic Hecke character $y\mapsto(\chi_W(y/\aubar y)\upsilon(y)^{-r}|y|^{-r/2})^{-1}$. Hence
\eqref{pg:eq-trace} holds with this $\eta_O$ and $R_{j,q}=S_q\otimes\eta_j$.
\end{proposition}
\begin{proof}
$n=r+t$ is even ($r$, $t=5r$ odd). For even $N$ the space character of rank $N$ of the
character $\mu$ (the character $\chi_V$) is the
base change $y\mapsto\mu(y/\aubar y)$, the powers of $\upsilon$ cancelling; so $\beta_w=\mu(\varpi_w/\aubar\varpi_w)q^{n/2}$ with
$\mu=\mu_j^0\vartheta$. Hence $q^t\beta_w^{-1}=(\mu_j^0)^{-1}(\varpi_w/\aubar\varpi_w)\,
\vartheta^{-1}(\varpi_w/\aubar\varpi_w)\,q^{(t-r)/2}$, and $|\varpi_w|^{(r-t)/2}=q^{(t-r)/2}$:
this is $\eta_j(\pgFrob_w)\psi_\vartheta(\pgFrob_w)$ by definition of $\psi_\vartheta$
(\S\ref{pg:ss-reduction}) and of $\eta_j$. By (G1) and
Proposition~\ref{pga:L}, for split $w$ off the level the eigenvalues of $S_y(\pgFrob_w)$ are
$\alpha_wq^l$ and $q^t\alpha_w\beta_w^{-1}\gamma_k$, those of $S_q(\pgFrob_w)$ being the
$\gamma_k$; multiplying by $\eta_O(\pgFrob_w)=\alpha_w^{-1}$ gives $\chi(\pgFrob_w)^l$ ($l<t$) and
the eigenvalues of $(\psi_\vartheta\eta_jS_q)(\pgFrob_w)$. Traces agree on a set of Frobenius
elements of density one, hence everywhere (Chebotarev, continuity). The only properties of
$\eta_O$ used later are that it is the character of an algebraic Hecke character depending on the
pair only; it is unramified at every finite place not above $\ell$ (Lemma~\ref{pay:eta}).
\end{proof}

\begin{proposition}[Second normalisation: the weights]\label{pga:weights}
Let $\epsilon_\sigma=1$ for $\sigma$ in the CM type and $-1$ otherwise. The Hodge--Tate weights of
$R_{j,q}$ at $\sigma$ are
\[
\tfrac{t-1}2-\epsilon_\sigma\bigl(\varsigma_{\sigma,i}+\kappa_j(\sigma)\bigr),\qquad1\le i\le r,
\]
and (M2) holds. For $R_{\nu,z}$ the weight at $\sigma$ is $z-1-\epsilon_\sigma c(\sigma)$, $-c$ the
weight of $\nu$, and (M2) holds.
\end{proposition}
\begin{proof}
By (G3) for $y_0$ on the definite $\auU(W_m)$ (rank $r$, parameter $\varsigma$), $S_q$ has weights
$\frac{r-1}2-\epsilon_\sigma\varsigma_{\sigma,i}$. By (G3) in rank one (same convention: the
parameter of a character of $\auU(1)$ is its weight), the character $y\mapsto(\mu_j^0)^{-1}(y/\aubar
y)$, of $\auU(1)$-weight $+\kappa_j$, has weight $-\epsilon_\sigma\kappa_j(\sigma)$; and $|y|^{(r-t)/
2}$ has weight $\frac{t-r}2$ at every $\sigma$ (it is $\chi^{(t-r)/2}$ times a character of finite
order, up to the unitary normalisation). Adding, the weights of $R=S_q\otimes\eta_j$ are as
stated. (The two signs come from one and the same convention (G3), so a change of convention
$\sigma\leftrightarrow\aubar\sigma$ changes both, and the relative sign is convention-free.
Consistency check: the weights of $R$ at $\sigma$ and $\aubar\sigma$ pair to $t-1$, as (M1)
requires.)

Now $\varsigma_{\sigma,i}+\kappa_j=\kappa_j+(\frac k2-i)$ for $i<k$ and $=(j-i'+\frac12)\Delta$ for
the $i'$-th entry of the second group ($1\le i'\le m$). With $|\frac{t-1}2|<t$, $k\le r$ and
$t=5r$: the first $k-1$ weights $u$ satisfy $|u+\epsilon_\sigma\kappa_j|\le\frac k2+\frac{t-1}2\le
t$ (high, with the sign of (M2)); the other $m$ satisfy $\bigl||u|-|j-i'+\frac12|\Delta\bigr|\le
\frac{t-1}2<3r$, and $\frac12\Delta\le|j-i'+\frac12|\Delta\le(m-\frac12)\Delta$ with $\Delta=2r^8$, so
$\Delta/4\le|u|\le2m\Delta$ (low). Weights with $|u|\le2h^2=2r^{16}$: not the high ones, since
$\kappa_j\ge\kappa\ge xb/x_1>2r^{16}+t$ by (S1); all low ones, since $2m\Delta\le2r^{16}$; so there
are exactly $m\le rm/k$. Negative weights: at $\sigma$ in the CM type the $k-1$ high ones and the
low ones with $i'\le j$ ($j$ of them); at $\aubar\sigma$ the low ones with $i'>j$ ($m-j$); in total
$k-1+m=r$. For $R_{\nu,z}$ the same computation in rank one, with the character
$\nu^{-1}(y/\aubar y)|y|^{1-z}$, gives $z-1-\epsilon_\sigma c(\sigma)$; it is high since
$|c-\kappa_j|\le t$ and $z\le r$, and exactly one of the weights at $\sigma,\aubar\sigma$ is
negative.
\end{proof}

\emph{The convention.} By Lemma~\ref{pax:conv}, the automorphic representation attached to a form
of $\mathcal A_\lambda$ has $\pi_\infty\simeq\rho_\lambda\circ\tau$ at $\tau$, of parameter
$\varsigma=\lambda+\rho$, uniformly in the rank, including rank one; so the parameters of
$\mathcal A(W_m)$ and the weight of $\mu^0_j$ are measured in the same convention, and the relative
sign above is justified.

\subsubsection{The published inputs of Theorem~\ref{pg:I}}\label{pga:ss-status}

Theorem~\ref{pg:I} follows from the proof of \S\ref{pg:ss-galois} supplemented by
Propositions~\ref{pga:L} (statement (L)), \ref{pga:M5} ((M5)), \ref{pga:trace} and
\ref{pga:weights} (the two normalisations), Proposition~\ref{pax:sign} (the twist for Speh
constituents) and Lemma~\ref{pax:conv} (the convention), and the following published results:
\begin{enumerate}
\item[(1)] the published theorems quoted in Theorem~\ref{pga:lit} (base change of Labesse;
Galois representations and local--global compatibility of Harris--Taylor, Shin, Chenevier--Harris,
Caraiani, Barnet-Lamb--Gee--Geraghty--Taylor);
\item[(2)] the local results quoted in Lemmas~\ref{pga:depth0} and \ref{pga:parahoric}: Moy--Prasad
on depth-zero types and Jacquet functors, Zelevinsky's irreducibility criterion, the level-zero
case of the tame local Langlands correspondence (Bushnell--Henniart), Jacquet's lemma;
\end{enumerate}
Statement (C) (Proposition~\ref{pga:C}) is proved in full; it is used in Lemma~\ref{pg:line}, not
in Theorem~\ref{pg:I}. Proposition~\ref{au:mult} is used in the proof of Theorem~\ref{pg:main},
not in Theorem~\ref{pg:I}.


%% ===== proofs/B-discrete-series.tex =====
% proofs/B-discrete-series.tex
% The discrete series coefficient at a real place: Theorem pbe:ds of proofs/B-bessel-expansion.tex,
% the chamber of the discrete series, the bound for kappa_tau, and the two definitions of A_0.
% Labels of this file start with pds:.
\providecommand{\auO}{\mathcal O}
\providecommand{\auPl}{\Phi}
\providecommand{\auAz}{\mathcal A_0}
\providecommand{\auU}{\mathrm U}
\providecommand{\auA}{\mathbb A}
\providecommand{\aurd}{\mathrm{rd}_E}
\providecommand{\aubar}[1]{\overline{#1}}
\providecommand{\auf}[1]{\texttt{#1}}
\providecommand{\pbvol}{\operatorname{vol}}
\providecommand{\pbdeg}{\operatorname{deg}}
\providecommand{\pbip}[2]{\langle#1,#2\rangle}
\providecommand{\pbeq}{\mathsf q}
\providecommand{\pbW}{\mathbb W}
\providecommand{\pbel}{\ell}
% macros of this file (prefix \pds)
\providecommand{\pdsK}{\mathbf K}
\providecommand{\pdsg}{\mathfrak g}
\providecommand{\pdsk}{\mathfrak k}
\providecommand{\pdsp}{\mathfrak p}
\providecommand{\pdsm}{m}
\providecommand{\pdsD}{\Delta^{\mathrm{min}}}
\providecommand{\pdsP}{P}
\providecommand{\pdsHS}[1]{\|#1\|_{2}}

\section{The discrete series coefficient at a real place}\label{pds:sec}

This section treats Theorem~\ref{pbe:ds} of \S\ref{pbe:sec}, the four properties
\textup{(DS1)}--\textup{(DS4)} of the coefficient $c_\tau$, together with three points of the
set-up of \S\ref{pbe:ss-ds}: which discrete series is meant (\S\ref{pds:ss-chamber}), the
inequality \eqref{pbe:eq-kappa} for $\kappa_\tau$ (\S\ref{pds:ss-kappa}), and the agreement of two
definitions of $\auAz$ (\S\ref{pds:ss-A0}).

\emph{What is proved and what is not.}
\begin{itemize}
\item[(1)] \textup{(DS1)(i),(ii)} are Schur's orthogonality relations (published); see
\S\ref{pds:ss-ds1}.
\item[(2)] \textup{(DS1)(iii)} is \emph{equivalent} to Harish-Chandra's formula for the formal degree
in the normalisation by the compact dual (Proposition~\ref{pds:haar}); for parameters with gap at
least $500N^3+10N^5$ this formula is proved in Corollary~\ref{pbr:degree}, without Harish-Chandra's
theorem; all parameters of the paper have this gap (Proposition~\ref{pwf:haar}), so (DS1)(iii) is
proved in full.
\item[(3)] \textup{(DS2)} is proved in full (Theorem~\ref{pds:decay}) from the explicit
integral formula \eqref{pds:eq-FJ} for $c_\tau$, which is proved in Theorem~\ref{pco:FJ}
(the gap $h\ge2r^8$ exceeds $N^2+1$).
\item[(4)] \textup{(DS3)} is proved in full (Theorem~\ref{pds:tail}) from \textup{(DS2)}, the
definition of the measure in \textup{(DS1)(ii)}, and the value of the Casimir operator on a discrete
series (published). It does not use Harish-Chandra's formula.
\item[(5)] \textup{(DS4)} is proved in \S\ref{pbp:sec} (Theorem~\ref{pbp:main}); \S\ref{pds:ss-ds4}
only records how it is used.
\item[(6)] The chamber of \S\ref{pbe:ss-ds} is confirmed from the definition of the signs and of the
parameter (Lemma~\ref{pds:chamber}). The inequality \eqref{pbe:eq-kappa} holds for the generator
$\varpi$ of the different chosen in Lemma~\ref{pds:kappa}, not for an arbitrary purely imaginary
generator. The space $\auAz$ consists of the bounded solutions of the reproducing equation
(\S\ref{ss:au-spaces}); they are square-integrable forms of the required type
(Proposition~\ref{pds:A0}).
\end{itemize}

\subsection{Set-up}\label{pds:ss-setup}

Fix $\tau\in\auPl$ and drop it from the notation. Put $N=r=k+2\varrho+1$. As in
\S\ref{pbe:ss-real}, $Y_\tau$ has the normalised basis, in which the hermitian form has the matrix
$\mathsf J_1$, hermitian and unitary, and the norm $\|\cdot\|$ for which that basis is orthonormal.
Since $\mathsf J_1$ is hermitian and unitary, there is a unitary matrix $p$ with $p^{-1}\mathsf J_1p=
J:=\operatorname{diag}(\varepsilon_1,\dots,\varepsilon_N)$, $\varepsilon_i=\pm1$, with $k+\varrho$
signs $+$. Conjugation by $p$ preserves operator norms and carries $G_\tau$ onto
\[
G=\auU(J)=\{g\in\mathrm{GL}_N(\C):g^*Jg=J\},
\]
and $K_\tau=G_\tau\cap\auU(N)$ onto $\pdsK=G\cap\auU(N)$, the block diagonal elements of $G$ (an
element of $G$ is unitary if and only if it commutes with $J$). We work in $G$ from now on;
$\|g\|$ is the operator norm, and $\|g\|=\|g^{-1}\|\ge1$ because $g^{-1}=Jg^*J$. For a matrix $x$ let
\[
a(x)=\tfrac12(x+JxJ)
\]
be its \emph{block diagonal part} (the entries $x_{ij}$ with $\varepsilon_i=\varepsilon_j$). The Lie
algebra is $\pdsg=\{X:X^*J+JX=0\}=\pdsk\oplus\pdsp$ with $\pdsk=\pdsg\cap\{X=JXJ\}$ (block diagonal,
skew-hermitian) and $\pdsp=\pdsg\cap\{X=-JXJ\}$ (off the diagonal blocks, hermitian).

\subsubsection*{Parameters}
Let $\lambda=\varsigma_\tau\in\R^N$ be the Harish-Chandra parameter, its entries attached to the
coordinates as in \S\ref{pbe:ss-ds}; all entries are integers ($k$ is even and $\kappa(\tau)$,
$\Delta$ are integers), they are distinct, and $N-1$ is even, so $\lambda\in\frac{N-1}2+\Z^N$. The
positive roots are the $e_i-e_j$ with $\lambda_i>\lambda_j$. Put
\[
\rho_i=\tfrac12\bigl(\#\{j:\lambda_j<\lambda_i\}-\#\{j:\lambda_j>\lambda_i\}\bigr),\qquad
\mu=\lambda+\rho\in\Z^N,
\]
and let $\rho_{\mathrm c}$, $\rho_{\mathrm n}$ be the same half sums restricted to the $j$ with
$\varepsilon_j=\varepsilon_i$, respectively $\varepsilon_j\ne\varepsilon_i$. The highest weight of the
lowest $\pdsK$-type is $\Lambda=\lambda+\rho_{\mathrm n}-\rho_{\mathrm c}\in\Z^N$ (Blattner's
formula for the lowest $K$-type; e.g.\ \cite[Theorem~9.20(b), p.~310]{Knapp}, stated there for
linear connected semisimple groups; it applies to $\auU(p,q)=Z\cdot\mathrm{SU}(p,q)$, $Z\cong\auU(1)$
the compact centre, $Z\cap\mathrm{SU}(p,q)$ finite: a discrete series of $\auU(p,q)$ restricts
irreducibly to a discrete series of $\mathrm{SU}(p,q)$ with $Z$ acting by a character,
$\pdsK=Z\cdot(\pdsK\cap\mathrm{SU}(p,q))$, and $K$-types, their multiplicities and the lowest $K$-type
correspond under restriction of highest weights to $\mathfrak t\cap\mathfrak{su}(p,q)$; see the
paragraph ``On (D3)'' in \S\ref{par:ss-rigid}). The \emph{gap} is
\[
h=\min\{|\lambda_i-\lambda_j|:\varepsilon_i\ne\varepsilon_j\}\ \ge\ \Delta=2r^8,
\]
by \S\ref{pbe:ss-ds}. We put $L=\max_i|\lambda_i|$; then $L\le\kappa(\tau)+m\Delta\le2b$ (comparison
of sizes; $\kappa(\tau)\le b$), $|\rho_i|\le N/2$ and $|\Lambda_i|\le L+N$.

\subsubsection*{Products of minors}
For $x\in M_N(\C)$ and $S\subseteq\{1,\dots,N\}$ let $\Delta_S(x)=\det(x_{ij})_{i,j\in S}$
($\Delta_\emptyset=1$). For $w\in\Z^N$ and $t\in\Z$ let $S_t(w)=\{i:w_i\ge t\}$, and for invertible
$x$ put
\begin{equation}\label{pds:eq-m}
\pdsm_w(x)=\prod_{t\in\Z}\Delta_{S_t(w)}(x)\cdot(\det x)^{-[t\le0]},
\end{equation}
where $[t\le 0]$ is $1$ or $0$; almost all factors are $1$. Shifting $t$ shows
\begin{equation}\label{pds:eq-shift}
\pdsm_{w+c\mathbf 1}(x)=(\det x)^c\,\pdsm_w(x)\qquad(c\in\Z).
\end{equation}
For $B\in\N$ with $|w_i|<B$ for all $i$ we also use the finite product $\pdsP^B_w(x)=\prod_{-B<t\le B}
\Delta_{S_t(w)}(x)$; it equals $\pdsm_w(x)$ when $\det x=1$.

\subsection{The explicit formula}\label{pds:ss-FJ}

Every $g\in G$ is $g=kz$ with $k\in\pdsK$ and $z\in G$ positive definite: $z=(g^*g)^{1/2}$ lies in
$G$ because $J(g^*g)J=(g^*g)^{-1}$ forces $JzJ=z^{-1}$ (uniqueness of positive square roots), and
$k=gz^{-1}$ is unitary. Let $dk$ be the Haar measure of mass $1$ on $\pdsK$ and
$d_\Lambda=(\int_{\pdsK}|\pdsm_\Lambda|^2)^{-1}$.
\begin{equation}\label{pds:eq-FJ}
c_\tau(kz)=d_\Lambda\int_{\pdsK}\frac{\pdsm_\Lambda(ku)\,\aubar{\pdsm_\Lambda(u)}}{\pdsm_{-\mu}(u^{-1}zu)}
\,du\qquad(k\in\pdsK,\ z\in G\text{ positive definite}).
\end{equation}
(The denominator is positive by Lemma~\ref{pds:fischer} below.)

Formula \eqref{pds:eq-FJ} is Theorem~\ref{pco:FJ}, applied with $n=N$; its hypothesis $h\ge N^2+1$
holds since $h\ge2r^8$. (This is Flensted-Jensen's description of the coefficient of the lowest
$K$-type~\cite{FJ}; the proof in \S\ref{pco:sec} does not rely on~\cite{FJ}.) Only \textup{(DS2)} uses
\eqref{pds:eq-FJ}; \textup{(DS1)}, \textup{(DS3)} use only that $c_\tau$ is the coefficient of a unit
vector of the lowest $K$-type.

\subsection{The chamber}\label{pds:ss-chamber}

\begin{lemma}\label{pds:chamber}
The signs of \eqref{pm:eq-signs} on $Y_0$ are $\varepsilon_i=(-1)^{\max(0,\,i+1-k)}$ for the
coordinate numbered $i=0,\dots,r-1$, and the doubled parameter is twice
$\lambda_i=\frac k2-1-i$ for $i<k-1$ and $\lambda_i=-\kappa(\tau)-(i+2-k)\Delta$ for $i\ge k-1$.
Hence the entries $\frac k2-1,\dots,1-\frac k2$ and the entries $-\kappa(\tau)-j\Delta$ with $j$ odd
sit at coordinates of sign $+$, those with $j$ even at coordinates of sign $-$; there are $k+\varrho$
signs $+$ and $\varrho+1$ signs $-$. On the $k-1$ coordinates of the first group $\Lambda=0$, and on
the other $m=2\varrho+2$ coordinates, in the order of decreasing $\lambda$, the signs are $+,-,+,
\dots,-$.
\end{lemma}
\begin{proof}
The formulas restate the signs \eqref{pm:eq-signs} and the definition of $\varsigma_\tau$ in
\S\ref{ss:au-spaces} (the entry $i\ge k-1$ is $-\kappa(\tau)-j\Delta$ with $j=i+2-k$). With $j=i+2-k$
the sign is $(-1)^{j-1}$, which is the first claim;
the coordinates $i\le k-1$ give $k$ signs $+$ and the $m-1$ coordinates $i\ge k$ give $\varrho$
signs $+$ and $\varrho+1$ signs $-$. For $i<k-1$: the entries below $\lambda_i$ are the $k-2-i$
later entries of the first group and all $m$ entries of the second, the entries above are $i$; the
coordinates of sign $-$ are the $\varrho+1$ entries with $j$ even, all below. So $2\rho_{\mathrm n,i}
=\varrho+1$ and $2\rho_{\mathrm c,i}=(k-2-i)+(\varrho+1)-i$, and $2\Lambda_i=(k-2-2i)+(\varrho+1)-
(k-2-2i+\varrho+1)=0$. The last claim restates the first.
\end{proof}
This agrees with what \S\ref{pbe:ss-ds} states, and with the structure used for the
Laplace method of \textup{(DS4)} in \S\ref{pbp:sec} (weights $0$ on $k-1$ coordinates of sign $+$, then $2\varrho+2$
strictly decreasing weights in $[-3B,-B]$ with alternating signs starting with $+$, for
$B=\kappa(\tau)+\Delta/2$): indeed $\Lambda_i=\lambda_i+O(N)$ there, and $\kappa(\tau)+m\Delta+N\le
3B$ because $m\Delta\le\kappa(\tau)$ (comparison of sizes).

\subsection{(DS1)}\label{pds:ss-ds1}

\emph{(i)} $c_\tau(g)=\pbip{v_0}{\pi(g)v_0}$ is continuous because $\pi$ is strongly continuous;
$c_\tau(g^{-1})=\pbip{\pi(g)v_0}{v_0}=\aubar{c_\tau(g)}$; $|c_\tau|\le\|v_0\|^2=1$.

\emph{(ii)} $\pi$ is irreducible and square-integrable, so for any Haar measure $da$ there is
$d_\pi>0$ with $\int_G\pbip{u'}{\pi(a)v'}\allowbreak \aubar{\pbip{u}{\pi(a)v}}da=d_\pi^{-1}\pbip{u'}{u}
\pbip v{v'}$ for all $u,u',v,v'$ (Schur orthogonality, \cite[Proposition~9.6, p.~284]{Knapp}; written for inner products conjugate-linear in the first variable). With all four
vectors $v_0$, $d_\pi^{-1}=\int|c_\tau|^2$; so for the measure with $\int|c_\tau|^2=\pbdeg_\tau^{-1}$
we have $d_\pi=\pbdeg_\tau$, and with $u=v=v'=v_0$, $u'=\pi(g)^{-1}v_0$:
$\int_G\pbdeg_\tau\aubar{c_\tau(a)}\,c_\tau(ga)\,da=\pbip{v_0}{\pi(g)v_0}$, using $c_\tau(ga)=
\pbip{\pi(g)^{-1}v_0}{\pi(a)v_0}$. This is \textup{(DS1)(ii)}; the integrability of $c_\tau$ assumed
there is proved below (Theorem~\ref{pds:tail}).

\emph{(iii)} Let $\mathfrak u=\pdsk\oplus i\pdsp=\mathfrak u(N)$, the compact dual. A Haar measure on
$G$ (on $\auU(N)$) is determined by a Lebesgue measure on $\pdsg$ (on $\mathfrak u$) through the
exponential map near $0$; the measures are \emph{compatible} if the Lebesgue measures correspond
under $X+Y\mapsto X+iY$ ($X\in\pdsk$, $Y\in\pdsp$).

\begin{proposition}\label{pds:haar}
Let $dg$ be the Haar measure on $G$ compatible with the Haar measure of mass $1$ on $\auU(N)$. Then
$\int_G|c_\tau|^2dg=\pbdeg_\tau^{-1}$ with $\pbdeg_\tau=\prod_{i<j}(\lambda_i-\lambda_j)/(j-i)$
(entries in decreasing order). Thus the measures of \textup{(DS1)(ii)} and \textup{(DS1)(iii)} agree.
\end{proposition}
\begin{proof}
By (ii), $\int|c_\tau|^2dg=d_\pi^{-1}$ for the formal degree $d_\pi$ relative to $dg$, and
$d_\pi=\prod_{i<j}(\lambda_i-\lambda_j)/(j-i)$ by Corollary~\ref{pbr:degree}, which applies because
the parameter satisfies the gap condition \eqref{pbr:eq-gap} (Proposition~\ref{pwf:haar}). This is
Harish-Chandra's formula $d_\pi=\prod_{\alpha>0}|\pbip\lambda\alpha|/\pbip\rho\alpha$ in the
normalisation by the compact dual; \cite[Theorem~16, p.~96]{HC66} gives it as $c^{-1}[W_G]|\varpi(\lambda)|$ with a constant $c$ that is not
computed there for this measure, and the proof here does not use it.
\end{proof}

\begin{remark}\label{pds:cayley}
The Haar measure of Proposition~\ref{pds:haar} can also be described by the Cayley chart
$y\mapsto2(1-\frac12J^{-1}S_y)^{-1}-1$ with density $(|\det J|\cdot|\det(1-\frac12J^{-1}S_y)|^2)^{-N}$,
divided by the constant that makes the mass $1$ for $J=1$; $S_y$ runs over the skew-hermitian matrices. The differential of the chart at $0$ is $X\mapsto2X$ on $\pdsg$ written as
$X=\frac12J^{-1}S$, and $S=2JX$; the map $X\mapsto JX$ from $\pdsg$ to $\mathfrak u(N)$ and the map
$X+Y\mapsto X+iY$ both act on the off-diagonal blocks by an isometry of the real coordinates and on
$\pdsk$ by $\pm1$ on each block, so both have Jacobian $1$. Hence this measure is the one of
Proposition~\ref{pds:haar}, and for this measure the identity $\pbdeg=(\int|c|^2)^{-1}=\prod_{i<j}
(\lambda_i-\lambda_j)/(j-i)$ is the same statement.
\end{remark}

\subsection{(DS2): decay}\label{pds:ss-ds2}

\begin{lemma}\label{pds:blockdet}
For $g\in G$: $a(g)$ is invertible, $|\det a(g)|\ge(\|g\|/2)^2$, and $|\det a(kgk')|=|\det a(g)|$ for
$k,k'\in\pdsK$. If $g$ is positive definite, so is $a(g)=\frac12(g+g^{-1})$.
\end{lemma}
\begin{proof}
Elements of $\pdsK$ commute with $J$, so $a(kgk')=k\,a(g)k'$ with $|\det k|=|\det k'|=1$. Write
$g=kz$ as in \S\ref{pds:ss-FJ}; then $\|g\|=\|z\|$ and $a(z)=\frac12(z+JzJ)=\frac12(z+z^{-1})$.
The spectrum of $z$ is that of $z^{-1}=JzJ$, so with $x_1\ge\dots\ge x_N>0$ the eigenvalues,
$x_N=x_1^{-1}$ and $x_1=\|z\|$. Thus $\det a(z)=\prod_i\frac12(x_i+x_i^{-1})$, a product of factors
$\ge1$; if $N\ge2$ and $x_1\ne1$ the two factors for $x_1,x_N$ give $\det a(z)\ge\frac14(x_1+
x_1^{-1})^2\ge(\|z\|/2)^2$, and if $x_1=1$ then $\det a(z)=1\ge\frac14$. The last claim: $g^{-1}=
JgJ$ when $g=g^*$, and $g^{-1}$ is positive definite.
\end{proof}

\begin{lemma}\label{pds:fischer}
Let $w\in G$ be positive definite and $0\le c\le h$ an integer. Then $\pdsm_{-\mu}(w)=\pdsm_\mu(w)>0$
and $\det a(w)^c\le\pdsm_\mu(w)^2$.
\end{lemma}
\begin{proof}
We use three facts. (F1) \emph{Fischer's inequality}: for positive definite $x$ and disjoint $T,T'$,
$\Delta_{T\cup T'}(x)\le\Delta_T(x)\Delta_{T'}(x)$ \cite[Theorem~7.8.5]{HornJohnson}, applied to the
principal submatrix of $x$ on $T\cup T'$.
(F2) \emph{Jacobi}: $\Delta_S(x^{-1})=\Delta_{S^c}(x)/\det x$. (F3) $\det w=1$ (it is positive and
$|\det w|=1$), and $\Delta_S(w^{-1})=\Delta_S(JwJ)=\Delta_S(w)$; with (F2), $\Delta_S(w)=
\Delta_{S^c}(w)$ for all $S$.

By (F3), $\pdsm_{-\mu}(w)=\prod_t\Delta_{\{\mu_i\le-t\}}(w)=\prod_s\Delta_{S_s(\mu)^c}(w)=\prod_s
\Delta_{S_s(\mu)}(w)=\pdsm_\mu(w)$, and this is positive since principal minors of positive definite
matrices are positive. Let $v=\mu+c\mathbf 1$, choose $B$ with $|\mu_i|,|v_i|<B$, and write $a=a(w)$,
positive definite by Lemma~\ref{pds:blockdet}. Fix $t$ and put $S=S_t(v)\supseteq T=S_t(\mu)$ and
$D=S\setminus T=\{i:\mu_i<t\le\mu_i+c\}$. \emph{All $i\in D$ have the same sign}: for $i,j\in D$,
$|\mu_i-\mu_j|\le c-1$; but $\rho_i-\rho_j$ has the sign of $\lambda_i-\lambda_j$, so $|\mu_i-\mu_j|
\ge|\lambda_i-\lambda_j|\ge h\ge c$ if $\varepsilon_i\ne\varepsilon_j$. Hence $a$ and $w$ have the same
entries on $D\times D$, and
\[
\Delta_S(a)\overset{\textup{(F1)}}\le\Delta_T(a)\Delta_D(a)=\Delta_T(a)\Delta_D(w)
\overset{\textup{(F3)}}=\Delta_T(a)\Delta_{T\cup S^c}(w)\overset{\textup{(F1)}}\le\Delta_T(a)
\Delta_T(w)\Delta_{S^c}(w).
\]
Multiply over $-B<t\le B$. On the left, $\pdsP^B_v(a)=(\det a)^c\pdsP^B_\mu(a)$, as in
\eqref{pds:eq-shift} (the $c$ new factors have $S_t(\mu)$ equal to everything, the $c$ lost ones
have $S_t(\mu)=\emptyset$). On the right, $\prod_t\Delta_{S_t(v)^c}(w)=\prod_t\Delta_{S_t(v)}(w)=
\pdsm_v(w)=\pdsm_\mu(w)$ by (F3) and \eqref{pds:eq-shift}. So $(\det a)^c\pdsP^B_\mu(a)\le
\pdsP^B_\mu(a)\,\pdsm_\mu(w)^2$, and $\pdsP^B_\mu(a)>0$.
\end{proof}

\begin{theorem}[\textup{(DS2)}]\label{pds:decay}
For every $g\in G$ and every integer $0\le c\le h$ (formula \eqref{pds:eq-FJ} holds by Theorem~\ref{pco:FJ}),
\[
|c_\tau(g)|\le|\det a(g)|^{-c/2}\le\bigl(2/\|g\|\bigr)^c .
\]
In particular $|c_\tau(g)|\le(2/\|g\|)^{j+\pbeq}$ for $0\le j\le r^7$, $\pbeq=9r^2$, and for $g\in
G_\tau$ with the norm of \S\ref{pbe:ss-real}.
\end{theorem}
\begin{proof}
Write $g=kz$ and let $D=|\det a(g)|=\det a(z)$. For $u\in\pdsK$, $w_u=u^{-1}zu\in G$ is positive
definite and $\det a(w_u)=D$ (Lemma~\ref{pds:blockdet}), so $\pdsm_{-\mu}(w_u)\ge D^{c/2}$ by
Lemma~\ref{pds:fischer}. By \eqref{pds:eq-FJ} and $2|xy|\le|x|^2+|y|^2$,
\[
|c_\tau(g)|\le\frac{d_\Lambda}{D^{c/2}}\int_{\pdsK}\frac{|\pdsm_\Lambda(ku)|^2+|\pdsm_\Lambda(u)|^2}2du
=D^{-c/2},
\]
by the invariance of $du$. The second inequality is Lemma~\ref{pds:blockdet}. For the last
sentence, $j+9r^2\le r^7+9r^2\le\Delta\le h$, and conjugation by $p$ preserves the norm.
\end{proof}
The bound holds for all exponents up to the gap, which is more than \textup{(DS2)} asks; nothing
else in \S\ref{pbe:sec} needs to change.

\subsection{(DS3): the formal degree against the volume of balls}\label{pds:ss-ds3}

Let $dg$ be the Haar measure of \textup{(DS1)(ii)}, $\int_G|c_\tau|^2dg=\pbdeg_\tau^{-1}$. We use one
published fact. Let $E_{ab}$ be the matrix units, acting on the $\pdsK$-finite vectors of $\pi$
through the complexified differential $d\pi$, and $\Omega=\sum_{a,b}E_{ab}E_{ba}$,
$\Omega_{\pdsk}=\sum_{\varepsilon_a=\varepsilon_b}E_{ab}E_{ba}$. Then
\begin{equation}\label{pds:eq-casimir}
d\pi(\Omega)=|\lambda|^2-|\rho|^2,\qquad d\pi(\Omega_{\pdsk})v_0=\bigl(|\Lambda+\rho_{\mathrm c}|^2-
|\rho_{\mathrm c}|^2\bigr)v_0 .
\end{equation}
The first is the value of the Casimir operator on a representation with infinitesimal character
$\lambda$ (Harish-Chandra isomorphism; \cite[Theorem~8.18, (8.32)]{Knapp}); the second is the same computation for the irreducible representation of highest
weight $\Lambda$ of $\pdsK\cong\auU(k+\varrho)\times\auU(\varrho+1)$.

\begin{lemma}\label{pds:lip}
Put $C_0=3N(L+2N)$. If $g\in G$ and $\|g-1\|\le\delta:=1/(4C_0)$, then $|c_\tau(g)|\ge\frac12$.
\end{lemma}
\begin{proof}
\emph{Step 1: $\|d\pi(X)v_0\|\le C_0\|X\|$ for $X\in\pdsg$.} For $\varepsilon_a\ne\varepsilon_b$ the
matrices $E_{ab}+E_{ba}$ and $i(E_{ab}-E_{ba})$ lie in $\pdsg$, and $d\pi$ of an element of $\pdsg$ is
skew-adjoint; hence $d\pi(E_{ab})^*=-d\pi(E_{ba})$ on $\pdsK$-finite vectors, and
$\pbip{v_0}{d\pi(E_{ab}E_{ba})v_0}=-\|d\pi(E_{ba})v_0\|^2$. By \eqref{pds:eq-casimir},
\[
\sum_{\varepsilon_a\ne\varepsilon_b}\|d\pi(E_{ab})v_0\|^2=|\Lambda+\rho_{\mathrm c}|^2-
|\rho_{\mathrm c}|^2-|\lambda|^2+|\rho|^2\le N(L+2N)^2+N^3/4\le2N(L+2N)^2 .
\]
For $Y=\sum_{\varepsilon_a\ne\varepsilon_b}Y_{ab}E_{ab}\in\pdsp$, Cauchy--Schwarz gives $\|d\pi(Y)v_0\|
\le\pdsHS Y\,(2N)^{1/2}(L+2N)\le\sqrt2\,N(L+2N)\|Y\|$. For $X\in\pdsk$, $d\pi(X)v_0$ is the action of
$X$ in the $\pdsK$-type $\Lambda$; conjugating $X$ in $\pdsK$ to $\operatorname{diag}(i\theta)$,
$|\theta_a|\le\|X\|$, its eigenvalues there are $i\pbip\nu\theta$ for weights $\nu$ in the convex hull
of the Weyl group orbit of $\Lambda$, so its norm is at most $\sum_a|\Lambda_a|\,\|X\|\le N(L+N)\|X\|$.
A general $X\in\pdsg$ is $X_{\pdsk}+X_{\pdsp}$ with $X_{\pdsk}=a(X)$, and both parts have norm at
most $\|X\|$.

\emph{Step 2.} Let $\|g-1\|\le\delta\le\frac12$ and $X=\log g=\sum_{n\ge1}(-1)^{n+1}(g-1)^n/n$. Then
$\exp X=g$ and $\|X\|\le-\log(1-\delta)\le2\delta$. Also $X\in\pdsg$: $JX^*J=\log(Jg^*J)=\log(g^{-1})=
-X$, by the holomorphic functional calculus (the series for $\log(g^{-1})$ converges because
$\|g^{-1}-1\|=\|Jg^*J-1\|=\|g-1\|$). The vector $v_0$ is $\pdsK$-finite, hence smooth, so
$t\mapsto\pi(\exp tX)v_0$ has derivative $\pi(\exp tX)d\pi(X)v_0$ and
\[
|c_\tau(g)-1|=|\pbip{v_0}{\pi(g)v_0-v_0}|\le\int_0^1\|d\pi(X)v_0\|\,dt\le2C_0\delta=\tfrac12 .
\qedhere
\]
\end{proof}

\begin{lemma}\label{pds:balls}
For $T\ge1$ let $S_T=\{g\in G:\|g\|\le T\}$ and $B_\delta=\{g:\|g-1\|\le\delta\}$ with $\delta$ as in
Lemma~\ref{pds:lip}. Then $\pbvol(S_T)\le(5T^2/\delta)^{2N^2}\pbvol(B_\delta)\le4(5T^2/\delta)^{2N^2}
\pbdeg_\tau^{-1}$.
\end{lemma}
\begin{proof}
Choose a maximal family $g_1,\dots,g_\nu$ in $S_T$ with $\|g_ig_j^{-1}-1\|>\delta$ for $i\ne j$ (it is
finite by the estimate that follows, $S_T$ being compact). Since $g_i-g_j=(g_ig_j^{-1}-1)g_j$ and
$\|g_j^{-1}\|=\|g_j\|\le T$, $\|g_i-g_j\|>\delta/T$. The balls of radius $\delta/(2T)$ about the $g_i$
in $M_N(\C)$ (real dimension $2N^2$) are disjoint and lie in the ball of radius $2T$ about $0$, so
$\nu\le(4T^2/\delta)^{2N^2}$. By maximality every $g\in S_T$ has $gg_i^{-1}\in B_\delta$ for some $i$,
so $S_T\subseteq\bigcup_iB_\delta g_i$ and, $G$ being unimodular, $\pbvol(S_T)\le\nu\,\pbvol(B_\delta)$.
Finally $1\le4|c_\tau|^2$ on $B_\delta$ (Lemma~\ref{pds:lip}), so $\pbvol(B_\delta)\le4\int_G|c_\tau|^2
=4\pbdeg_\tau^{-1}$.
\end{proof}

\begin{theorem}[\textup{(DS3)}]\label{pds:tail}
For every real $\mathsf s>4N^2$, $\pbdeg_\tau\int_G\|g\|^{-\mathsf s}dg\le4(5/\delta)^{2N^2}\mathsf s/
(\mathsf s-4N^2)$. In particular $c_\tau$ is integrable, and
\[
\pbdeg_\tau\int_{G_\tau}\bigl(32/\|g\|\bigr)^{\pbeq}dg\le8\cdot32^{\pbeq}\bigl(60N(L+2N)\bigr)^{2N^2}
\le b^{99r^2}\qquad(\pbeq=9r^2).
\]
\end{theorem}
\begin{proof}
As $\|g\|\ge1$, $\|g\|^{-\mathsf s}=\int_1^\infty\mathsf sT^{-\mathsf s-1}[\|g\|\le T]\,dT$; by Fubini
and Lemma~\ref{pds:balls},
\[
\int_G\|g\|^{-\mathsf s}dg=\int_1^\infty\mathsf sT^{-\mathsf s-1}\pbvol(S_T)\,dT\le
\frac{4}{\pbdeg_\tau}\Bigl(\frac5\delta\Bigr)^{2N^2}\int_1^\infty\mathsf sT^{4N^2-\mathsf s-1}dT .
\]
Integrability of $c_\tau$ follows with Theorem~\ref{pds:decay} ($c=9N^2\le h$). For $\mathsf s=
\pbeq=9N^2$, $\mathsf s/(\mathsf s-4N^2)=9/5<2$ and $5/\delta=20C_0=60N(L+2N)$, which gives the
middle expression after multiplying by $32^{\pbeq}$. Lastly $N=r$, $L\le2b$ and $r\le b$ give
$60N(L+2N)\le b^3$, so the middle expression is at most $8\cdot2^{45r^2}b^{6r^2}\le b^{99r^2}$.
\end{proof}

\subsection{(DS4): the Bessel period at a stationary point}\label{pds:ss-ds4}

\smallskip\noindent\textup{(DS4)} is proved in \S\ref{pbp:sec} (Theorem~\ref{pbp:main}).\par\smallskip

\subsection{The bound for \texorpdfstring{$\kappa_\tau$}{kappa}}\label{pds:ss-kappa}

The inequality \eqref{pbe:eq-kappa} is not true for every purely imaginary generator $\varpi$ of the
different: $\varepsilon\varpi$, with $\varepsilon$ a unit of $E^+$ having one large conjugate,
violates it. It holds for the generator of Lemma~\ref{pds:kappa}, which is the one fixed in
\eqref{pb:eq-psi}. By Lemma~\ref{an:lem-sizes}(S7), $q_1=q_0t$ with every prime factor of $t$ dividing
$q_0$ (so that $\Phi_{q_1}(X)=\Phi_{q_0}(X^t)$ for the cyclotomic polynomials), and $(q_0/2)^{q_0}\le
q_0^{q_0}\le x$.

\begin{lemma}\label{pds:kappa}
Let $\zeta\in E$ be a primitive $q_1$-th root of unity and $\varpi=\gamma\,\Phi_{q_1}'(\zeta)\sqrt{-d}$,
where $\gamma\in\auO_E^\times$ satisfies $\gamma\aubar\gamma=1$ and $\aubar{\gamma\Phi'_{q_1}(\zeta)}=
\gamma\Phi'_{q_1}(\zeta)$ \textup{(}such $\gamma$ exist, e.g.\ $\gamma=\zeta^s$ with
$s=q_1/4+(q_1-\varphi(q_1))/2+1$; see the first paragraph of the proof\textup{)}. Then $\varpi$ is a purely imaginary generator of the different of $E$, and
for every $\tau\in\auPl$, $2\le|\tau(\varpi)|\le b$ and $1\le M|\tau(\varpi)|/2\pi\le b^{23}$.
\end{lemma}
\begin{proof}
\emph{Existence of $\gamma$.} Write $m=q_1$ and $u=\Phi_m'(\zeta)$, and let $g(X)=\prod_{e\mid m,\,e<m}
\Phi_e(X)$, so $X^m-1=\Phi_m(X)g(X)$, $g(\zeta)\ne0$, and differentiating at $\zeta$ gives
$u=m\zeta^{-1}/g(\zeta)$. Complex conjugation of the CM field $E$ sends $\zeta$ to $\zeta^{-1}$ and
fixes the rational coefficients, so $\aubar u=m\zeta/g(\zeta^{-1})$. Since $\Phi_1(1/X)=-X^{-1}\Phi_1(X)$
and $X^{\varphi(e)}\Phi_e(1/X)=\Phi_e(X)$ for $e\ge2$, we get $g(\zeta^{-1})=-\zeta^{-D}g(\zeta)$ with
$D=\deg g=m-\varphi(m)$, hence $\aubar u=-\zeta^{D+2}u=\zeta^{m/2+D+2}u=\zeta^{2s}u$ with
$s=m/4+D/2+1\in\mathbb Z$ (as $16\mid m$ and $\varphi(m)$ is even). For $\gamma=\zeta^s$:
$\gamma\aubar\gamma=1$ and $\aubar{\gamma u}=\zeta^{-s}\zeta^{2s}u=\gamma u$.

Since $d\equiv q_0/2-1\pmod{q_0}$ and $16\mid q_0$, $d\equiv-1$ modulo every prime factor of $q_0$ and
$d\equiv3\pmod 4$; so $d$ is prime to $q_1$, $\auO_K=\Z[\frac{1+\sqrt{-d}}2]$ has different
$(\sqrt{-d})$, and the discriminants of $K$ and $\Q(\zeta)$ are coprime. Hence $\auO_E=\auO_K[\zeta]$,
the different of $E/K$ is $(\Phi'_{q_1}(\zeta))$, and the different of $E/\Q$ is the product of those of
$E/K$ and $K/\Q$.

(\emph{$\auO_E=\auO_K[\zeta]$.} Let $F=\Q(\zeta)$; its ring of integers is $\Z[\zeta]$, with
$\Z$-basis $\beta_j=\zeta^{j}$, $0\le j<n=\varphi(q_1)$, and discriminant $d_F=\pm\auNm_{F/\Q}
\Phi'_{q_1}(\zeta)$~\cite[Chapter~I, \S2 and Proposition~(10.2)]{Neukirch}. Differentiating $X^{q_1}-1=\Phi_{q_1}(X)g(X)$ gives
$q_1\zeta^{q_1-1}=\Phi'_{q_1}(\zeta)g(\zeta)$, so $\Phi'_{q_1}(\zeta)$ divides $q_1$ and every prime
factor of $d_F$ divides $q_1$; as $d_K=-d$ is prime to $q_1$, $\gcd(d_K,d_F)=1$. Since $K$ is ramified
at $d$ and $F$ is not, $K\cap F=\Q$, so $[E:\Q]=2n$, $\Phi_{q_1}$ is irreducible over $K$, the
$\beta_j$ are a $K$-basis of $E$, and $\omega_1=1,\omega_2=\frac{1+\sqrt{-d}}2$ are an $F$-basis
of $E$. Let $x\in\auO_E$ and write $x=\sum_{i,j}c_{ij}\omega_i\beta_j$ with $c_{ij}\in\Q$. Put
$c_j=\sum_ic_{ij}\omega_i\in K$. The $n$ embeddings $\sigma$ of $E$ over $K$ restrict to the $n$
embeddings of $F$; solving $\sigma(x)=\sum_jc_j\sigma(\beta_j)$ by Cramer's rule gives
$c_j=\gamma_j/\delta$ with $\delta=\det(\sigma\beta_j)$, $\delta^2=d_F$, and $\gamma_j$ an algebraic
integer. So $d_Fc_j=\delta\gamma_j\in\auO_K$, i.e.\ $d_Fc_{ij}\in\Z$. The same argument with the
roles of $K$ and $F$ exchanged (embeddings of $E$ over $F$, basis $\omega_i$, $\det^2=d_K$) gives
$d_Kc_{ij}\in\Z$. Hence $c_{ij}\in\Z$, and $\auO_E=\auO_K[\zeta]$.

\emph{The different of $E/K$.} Let $A=\auO_K$, $B=\auO_E=A[\zeta]$, $f=\Phi_{q_1}$ (monic of degree
$n$, irreducible over $K$), and write $f(X)/(X-\zeta)=b_0+b_1X+\dots+b_{n-1}X^{n-1}$ with $b_k\in
B$. Lagrange interpolation at the $n$ distinct roots $\sigma(\zeta)$ gives, for $0\le i<n$,
$\sum_\sigma\frac{f(X)}{X-\sigma\zeta}\frac{\sigma(\zeta)^i}{f'(\sigma\zeta)}=X^i$, since both
sides have degree $<n$ and agree at every $\sigma(\zeta)$. Comparing coefficients of $X^k$:
$\mathrm{Tr}_{E/K}\bigl(\zeta^ib_k/f'(\zeta)\bigr)=\delta_{ik}$ (Euler's formula). So
$\{b_k/f'(\zeta)\}$ is the basis dual to $\{\zeta^i\}$ under the trace, and since $B$ is free over
$A$ on $\{\zeta^i\}$, $\{x:\mathrm{Tr}_{E/K}(xB)\subseteq A\}=\sum_kA\,b_k/f'(\zeta)$. Comparing
coefficients in $f(X)=(X-\zeta)\sum b_kX^k$ gives $b_{n-1}=1$ and $b_{k-1}=\zeta b_k+a_k$ ($a_k\in A$
the coefficients of $f$), so $b_{n-1-m}\in\zeta^m+\sum_{l<m}A\zeta^l$; the $b_k$ span $A[\zeta]=B$
over $A$. Hence $\mathfrak D_{E/K}^{-1}=f'(\zeta)^{-1}B$, i.e.\ $\mathfrak D_{E/K}=(\Phi'_{q_1}(\zeta))$.)

(Transitivity of the different: for a finite extension $B/A$ of rings of integers
of number fields, $\mathfrak D_{B/A}^{-1}=\{x:\mathrm{Tr}_{B/A}(xB)\subseteq A\}$, and an $A$-submodule $M$ of
$\mathrm{Frac}(A)$ satisfies $\mathrm{Tr}_{\mathrm{Frac}(A)/\Q}(M)\subseteq\Z$ iff $M\subseteq\mathfrak D_{A/\Z}^{-1}$. As
$\mathrm{Tr}_{E/\Q}=\mathrm{Tr}_{K/\Q}\circ\mathrm{Tr}_{E/K}$ and $\mathrm{Tr}_{E/K}(x\auO_E)$ is an
$\auO_K$-module, $x\in\mathfrak D_{E/\Q}^{-1}$ iff $\mathrm{Tr}_{E/K}(x\auO_E)\subseteq\mathfrak
D_{K/\Q}^{-1}$ iff $\mathrm{Tr}_{E/K}(x\mathfrak D_{K/\Q}\auO_E)\subseteq\auO_K$ iff $x\mathfrak
D_{K/\Q}\auO_E\subseteq\mathfrak D_{E/K}^{-1}$; so $\mathfrak D_{E/\Q}^{-1}=\mathfrak D_{E/K}^{-1}
\mathfrak D_{K/\Q}^{-1}\auO_E$, fractional ideals being invertible. See also \cite[Chapter~I, Proposition~(2.11), and Chapter~III, Propositions~(2.2) and~(2.4)]{Neukirch}.) As
$\gamma$ is a unit, $\varpi$ generates the different, and $\aubar\varpi=-\varpi$.

For $\tau\in\auPl$, $|\tau(\gamma)|=1$ and, by the chain rule, $|\tau\Phi_{q_1}'(\zeta)|=
t\,|\Phi_{q_0}'(\xi)|$ with $\xi=\tau(\zeta)^t$ a primitive $q_0$-th root of unity. Now
$|\Phi'_{q_0}(\xi)|=\prod_{\xi'}|\xi-\xi'|$ over the other primitive $q_0$-th roots, fewer than $q_0$
factors each in $[2\sin(\pi/q_0),2]\subseteq[4/q_0,2]$. So $|\tau(\varpi)|=A_\tau\,t\sqrt d$ with
$(4/q_0)^{q_0}\le A_\tau\le2^{q_0}$, and any two conjugates differ by a factor at most $(q_0/2)^{q_0}
\le x$. Since $\prod_{\tau\in\auPl}|\tau(\varpi)|^2=|\operatorname{N}_{E/\Q}\varpi|=|\operatorname{disc}
E|=\aurd^{2f}$, the geometric mean of the $|\tau(\varpi)|$ is $\aurd\in[x^{10},x^{12}]$
\eqref{eq:au-sizes}, so $x^9\le|\tau(\varpi)|\le x^{13}$, whence $2\le|\tau(\varpi)|\le b$. With $b^{20}
\le M=p_0p_1\le b^{22}$: $1\le M\cdot2/2\pi\le M|\tau(\varpi)|/2\pi\le b^{22}\cdot b/2\pi\le b^{23}$.
\end{proof}

\subsection{The two definitions of \texorpdfstring{$\auAz$}{A0}}\label{pds:ss-A0}

Let $e_\infty=\bigotimes_\tau e_\tau$, $e_\tau=\pbdeg_\tau\aubar{c_\tau}$, an integrable function on
$G_\infty$ (Theorem~\ref{pds:tail}), and $X=G_0(E^+)\backslash G_0(\auA_{E^+})/J_M$, of finite volume
\cite[Theorem~5.6(i)]{BorelFin}, see also \cite[Theorem~9.4]{BorelHC1962} (the centre of $G_0$ is anisotropic). Let $\auAz^{\mathrm b}$ be the space of bounded
continuous $y$ on $X$ with $\int_{G_\infty}e_\infty(a)y(ga)\,da=y(g)$ for all $g$ (this is $\auAz$, by
the definition of \S\ref{ss:au-spaces}), and $\auAz^2$ the space of $y\in L^2(X)$ with $\mathcal Ey=y$,
where $\mathcal Ey(g)=\int e_\infty(a)y(ga)\,da$ (the $L^2$ variant of the definition; only the
condition $\mathcal Ey=y$ is used below).

\begin{proposition}\label{pds:A0}
$\auAz^{\mathrm b}\subseteq\auAz^2$. Equality holds if every $y\in\auAz^2$ is (a.e.\ equal to) a bounded
continuous function.
\end{proposition}
\begin{proof}
A bounded $y$ is in $L^2(X)$ since $X$ has finite volume, and $\mathcal Ey=y$ pointwise. Conversely, if
$y\in\auAz^2$ is bounded and continuous, $\mathcal Ey$ is defined everywhere (as $e_\infty\in L^1$) and
continuous (by dominated convergence and the continuity of translation in $L^1(G_\infty)$), and
equals $y$ a.e., hence everywhere.
\end{proof}

\begin{remark}\label{pds:rem-A0}
The converse inclusion is not used in this paper, since $\auAz=\auAz^{\mathrm b}$ by definition.
For orientation: it amounts to the boundedness of every $y\in\auAz^2$. Elements of $\auAz^2$ are $K_\infty$-finite and
$Z(\pdsg)$-finite, hence smooth automorphic forms (Harish-Chandra). Boundedness would follow from
cuspidality, since cusp forms are rapidly decreasing; cuspidality of square-integrable forms whose
archimedean component is a discrete series with parameter far from the walls is a theorem of
Wallach~\cite{WallachCT}. Nothing in \S\S\ref{pb:sec}--\ref{pm:sec} uses it.
\end{remark}

\subsection{The inputs of Theorem~\ref{pbe:ds}}\label{pds:ss-summary}
Of Theorem~\ref{pbe:ds}: \textup{(DS1)(i),(ii)} follow from Schur orthogonality; \textup{(DS1)(iii)}
from Corollary~\ref{pbr:degree} \textup{(}Proposition~\ref{pds:haar}, Proposition~\ref{pwf:haar}\textup{)}; \textup{(DS2)} is proved in full from the explicit formula
\eqref{pds:eq-FJ}; \textup{(DS3)} is proved in full from \textup{(DS1)(ii)} and the Casimir values
\eqref{pds:eq-casimir}; \textup{(DS4)} is proved in \S\ref{pbp:sec}. Formula \eqref{pds:eq-FJ} is
Theorem~\ref{pco:FJ}. Published results used: Schur orthogonality and the value of the Casimir operator on a discrete series.


%% ===== proofs/B-pillars.tex =====
% proofs/B-pillars.tex
% Theorem pbe:pillars of proofs/B-bessel-expansion.tex: cell ranks (R), size (S), gain (G),
% count (C). Labels of this file start with ppi:.
\providecommand{\auO}{\mathcal O}
\providecommand{\auNm}{\operatorname{Nm}}
\providecommand{\auA}{\mathbb A}
\providecommand{\aubar}[1]{\overline{#1}}
\providecommand{\auf}[1]{\texttt{#1}}
\providecommand{\panot}[1]{\par\smallskip{\raggedright\noindent\textbf{[Not proved here:} #1\textbf{]}\par}\smallskip}
\providecommand{\pbR}{\overline R_0}
\providecommand{\pbDid}{\mathcal D_{\mathrm{id}}}
\providecommand{\pbDmain}{\mathcal D_{\mathrm{main}}}
\providecommand{\pbg}{\mathsf g}
\providecommand{\pbeN}{\mathcal N}
\providecommand{\pbet}{\beta_*}
\providecommand{\pbenp}{\mathsf n_{\mathrm p}}
\providecommand{\pbesz}{\operatorname{size}}
\providecommand{\pbeGs}{\Gamma^*}
\providecommand{\pbeAp}{A^{+}}
\providecommand{\pbeCL}{\mathrm{CL}}
% macros of this file (prefix \ppi)
\providecommand{\ppiF}{\mathsf F}
\providecommand{\ppiLc}{\Lambda^{X}_{c}}
\providecommand{\ppidisc}{\operatorname{disc}}

\section{Pairs of terms: cell ranks, size, gain and count}\label{ppi:sec}

This section concerns Theorem~\ref{pbe:pillars}. We use the notation of \S\S\ref{pb:sec}
and~\ref{pbe:sec}: the frame $f_0,\dots,f_\varrho$ with $f_a=e_{a+1}$ ($a<\varrho$) and $f_\varrho=z$,
$A=\langle f_0,\dots,f_{\varrho-1}\rangle$, $\pbeAp=A+Ez$, $Q=A^\perp/A=X_0\perp Ez$, the matrix $Z(
\gamma)$ with leading minors $\delta_l(\gamma)$, the discriminant $\pbg(\gamma)$, the link
$\Lambda_\gamma$, the isometry $\Upsilon_\gamma$ and the space $X_0^\gamma$ of
Lemma~\ref{pbe:link}, the set $\pbeGs$, and $\sigma$, $\psi_f$ as in Lemma~\ref{pbe:scalar}, so that
$\psi_f(u)=\psi_{E,f}(-\sigma(u)/(M\varpi))$ for $u\in U(E^+)$. For a coefficient vector $x=(x_0,
\dots,x_\varrho)\in E^n$ we write $fx=\sum_ax_af_a$ and $\gamma fx=\gamma(fx)$.

\emph{What is proved.}
\begin{itemize}
\item[(1)] The exact statement of Theorem~\ref{pbe:pillars} is fixed (Theorem~\ref{ppi:pillars}):
the set $\pbeN(c,c')$, the gauge, the discriminant and the size conditions under which it is
proved. Two size conditions on the data are needed that are not in \S\ref{pbe:sec}
(Remark~\ref{ppi:rem-sizes}).
\item[(2)] \textup{(R)}, and the existence of cell ranks with non-vanishing leading minor, are
\textbf{proved in full} (Proposition~\ref{ppi:cell}).
\item[(3)] The elementary parts of \textup{(G)} and \textup{(C)} are proved: the bound $1\le\mathsf
G\le b^{n^3f}$ for near terms (Lemma~\ref{ppi:gauge}), $\mathsf d\ge1$, the passage from the bound on
each common part to \textup{(G)} (Proposition~\ref{ppi:G}), and the reduction of \textup{(C)} to a
count of terms over one lattice and a sum over lattices (Proposition~\ref{ppi:C}).
The count over a lattice (C1) and the sum over lattices (C2) are proved in \S\ref{pbc:sec}. The bound for the common parts (G$'$) is proved in \S\ref{pbg:sec} (Theorem~\ref{pbg:main}). \textup{(S)} is proved in \S\ref{pbs:sec} given (O1) of \S\ref{pbo:sec}.
\item[(4)] Lemmas~\ref{pbe:rank}, \ref{pbe:pairs}, Proposition~\ref{pbe:two} and
Theorem~\ref{pbe:XL1} still follow from the corrected statement (\S\ref{ppi:ss-check}).
\end{itemize}

\subsection{The exact statement}\label{ppi:ss-statement}

\subsubsection*{Common vectors of two flags}
Let $\gamma\in\Gamma$, $0\le a\le\varrho$ and $\mu,\mu'\in E$. We say that \emph{$(\mu,\mu')$ occurs at
$a$} (for $\gamma$) if there are $x,y\in E^n$ with $fx=\gamma fy$, $x_b=y_b=0$ for $b>a$, and
$x_a=\mu$, $y_a=\mu'$. A common vector $fx=\gamma fy\ne0$ has a \emph{top index}: the largest $l$
with $(x_l,y_l)\ne(0,0)$; then $(x_l,y_l)$ occurs at $l$.

\begin{definition}[Cell rank]\label{ppi:def-cell}
Let $\gamma\in\Gamma$. If $\pbg(\gamma)\ne0$, the \emph{cell rank} of $\gamma$ is $j(\gamma)=\varrho+
1$. If $\pbg(\gamma)=0$, $j(\gamma)$ is the least $j\in\{0,\dots,\varrho+1\}$ such that, for one
$\zeta\in E$ with $\zeta\aubar\zeta=1$, $(\zeta,1)$ occurs at every $a$ with $j\le a\le\varrho$ (for
$j=\varrho+1$ the condition is empty, so $j(\gamma)$ exists).
\end{definition}

The cell rank is invariant under $\gamma\mapsto u\gamma u'$ with $u,u'$ in the group $U^0(E^+)$ of
the elements of $U(E^+)$ that fix every $f_a$ (if $fx=\gamma fy$ then $fx=ufx=u\gamma u'(u'^{-1}fy)=
u\gamma u'fy$). One may also require of a cell rank $\delta_{j}(\gamma)\ne0$
and $j\le\varrho\iff\pbg(\gamma)=0$; by Proposition~\ref{ppi:cell} below this holds for every
$\gamma\in\pbeGs$, so the two definitions agree on $\pbeGs$.

\subsubsection*{Near terms, labels, gauges, lattices}
Reduced and boxed points, and near elements, are defined in Definition~\ref{pbc:def-reduced}: $y$ is
reduced if the Gram--Schmidt coefficients of the vectors $y_\tau^{-1}\tilde f_a$, read in the frame
$f$, lie in a fixed Borel set of representatives of $\C^{\auPl}/M\auO_E$ with coordinates at most
$b^{26}$; boxed if moreover $\pbel_\tau(y_\tau)\le b^{6n+2}$ for all $\tau$; and $\gamma$ is near if
some boxed $y$ has $\gamma^{-1}y$ boxed. At a boxed point $\|y_\tau^{-1}\tilde f_a\|\le b^{37n}$
(Lemma~\ref{pbc:box}). For classes $c,c'$ let
$U_c^0=U_c\cap U^0$ and fix in each double coset $U_c^0\gamma U_{c'}^0$ one element, the \emph{chosen}
one.

For $\gamma\in\Gamma_{c,s,c'}$ write $k_\gamma=h_c^{-1}\gamma h_{c'}\in K$. Let $\Lambda=L$ be the
lattice of $Y_0$, $\ppiLc=h_cL\cap X_0$, and let $L_\gamma\subseteq\ppiLc$ be the set of the $v\in
\ppiLc$ such that some vector of $v+M(h_cL\cap A)$ is orthogonal to $\gamma f_0,\dots,\gamma
f_\varrho$. For $\gamma,\gamma'\in\Gamma_{c,s,c'}$ the \emph{common lattice} is $L_{\gamma\gamma'}=
L_\gamma\cap L_{\gamma'}\cap\pbeCL(\gamma,\gamma')$, and $\mathsf d(\gamma,\gamma')=\ppidisc
L_{\gamma\gamma'}$, where for a lattice $L'$ of rank $\rho$ in a subspace $V\subseteq X_0$ we put
$\ppidisc L'=[\ppiLc\cap V:L']\cdot\ppidisc(\ppiLc\cap V)$, and $\ppidisc(\ppiLc\cap V)$ is the square
root of the absolute norm of the ideal of $\auO_E$ generated by the Gram determinants of the $\rho$-tuples
of $\ppiLc\cap V$ (a generalised index).

Let $\gamma\in\Gamma_{c,s,c'}$ have cell rank $j$ and put $j'=\min(\varrho,j)$. The \emph{gauge} is
\begin{equation}\label{ppi:eq-gauge}
\mathsf G(\gamma)=M^{2j'f}\,\frac{\max\bigl(\auNm\pbg(\gamma),\ N^{f}\,\auNm\delta_j(\gamma)\,N^{-2(j-
\varrho)_+f}\bigr)}{i_\gamma\,i'_\gamma},
\end{equation}
where $\auNm x=\prod_{\tau\in\auPl}|\tau(x)|$ for $x\in E^+$ and $\auNm x=\prod_\tau|\tau(x)|^2$ for
$x\in E$, and $i_\gamma=[R_A:R_{<j'}]$, where $R_A$ (resp.\ $R_{<j'}$) is the $\auO_E$-span in $E^n$ of the
rows $a<\varrho$ (resp.\ $a<j'$) of $Z(\gamma)$, and $i'_\gamma$ is the same index for
$Z(\gamma^{-1})=Z(\gamma)^*$; both are finite and $\ge1$, and equal to $1$ if $j=\varrho+1$
(Remark~\ref{pbg:rem-gauge}). Equivalently, $\mathsf G(\gamma)=M^{2j'f}\max(\auNm\delta_*(\gamma)/
(i_\gamma i'_\gamma N^f),\auNm\pbg(\gamma))$ with the subframe minor $\delta_*$
(Lemma~\ref{pbg:subframe}).

Finally, $s\in\auG$ \emph{agrees} if $\psi_f(u)=\psi_f(u')$ for all integral $u,u'\in U$ whose
reductions satisfy $\bar us=s\bar u'$ (this notion, for $s$, is distinct from the agreement of an element $\gamma$ in \S\ref{pbe:ss-char}, which is the notion used in Lemma~\ref{ppi:climb} and Proposition~\ref{ppi:cell}). The \emph{label} of $\gamma\in\Gamma_{c,s,c'}$ is the component of $s$
in $\pbR$ for a fixed decomposition $\auG=\bigsqcup_\rho\rho\pbR$.

\begin{definition}\label{ppi:def-N}
$\pbeN(c,c')$ is the set of the $\gamma\in\pbeGs$ that are near and chosen and lie in
$\Gamma_{c,s,c'}$ for an agreeing $s$ (which $\gamma$ determines).
\end{definition}

Recall $\pbet=b^{kf/4}$, $\pbenp=\#\pbR\cdot\#\mathcal C/N^{kf}$, $\mathcal D(\gamma)$ and
$\pbesz(\gamma)=\pbet^2(\mathcal D(\gamma)/(\pbenp\mathsf G(\gamma)^k))^{1/2}$ from
\S\ref{pbe:ss-pillars}. Note that, for $\gamma\in\pbeGs$, $\mathcal D(\gamma)=\pbDmain$ if and only if $j(\gamma)=\varrho+1$ (Proposition~\ref{ppi:cell}).

\begin{theorem}[Size, gain and count; exact form]\label{ppi:pillars}
Assume the size conditions
\begin{equation}\label{ppi:eq-sizes}
\varrho\ge10^5,\qquad k\ge1000\,(2n)^9 .
\end{equation}
Let $W^\flat$ be as in Theorem~\ref{pbe:bessel}. With $\pbeN(c,c')$, $\mathsf G$ and $\mathsf d$ as
above, $\pbeN(c,c')$ contains at most one element of each double coset $U_c^0\gamma U_{c'}^0$,
$\mathsf G$ takes values in $(0,b^{n^3f}]$ on $\pbeN(c,c')$, $\mathsf d(\gamma,\gamma')=\mathsf d(\gamma',
\gamma)\ge1$, and:
\begin{enumerate}
\item[(R)] If $\gamma\in\pbeGs$ and $\pbg(\gamma)=0$, then $\pbeAp\cap\gamma\pbeAp\not\subseteq A$.
\item[(S)] For every stable $T\subseteq\pbeGs$ and every block $(c,s,c')$, $\|A^\flat_T(c,s,c')\|\le
\nu\sum_{\gamma\in\pbeN(c,c')\cap T\cap\Gamma_{c,s,c'}}\pbesz(\gamma)$. In particular $A^\flat_T(c,s,
c')=0$ if $s$ does not agree.
\item[(G)] For $\gamma,\gamma'\in\pbeN(c,c')\cap\Gamma_{c,s,c'}$ and $i=i(\gamma,\gamma')$,
$(\mathsf d(\gamma,\gamma')\,\mathsf x)^2\le(2/M)^{2if}\,\mathsf G(\gamma)\mathsf G(\gamma')$, with
$\mathsf x=\mathsf{mn}(s)$ if $i=0$ and $\pbg(\gamma)\ne0$, and $\mathsf x=1$ otherwise.
\item[(C)] For every integer $0\le l\le n^3$, $\sum_{c,c'}\sum_{(\gamma,\gamma')}\mathsf d(\gamma,
\gamma')^{-k}\le\pbet^{\,l+2}\pbenp$, the inner sum over the pairs of elements of $\pbeN(c,c')$ over a
common block with $k-\dim\pbeCL(\gamma,\gamma')=l$ and $\mathsf d(\gamma,\gamma')\le b^{n^3f}$.
\end{enumerate}
\end{theorem}

\begin{remark}[Changes against Theorem~\ref{pbe:pillars}]\label{ppi:rem-sizes}
(i) \eqref{ppi:eq-sizes} ($10^5\le\varrho$, $1000(2\varrho+2)^9\le k$) holds for the data by
Lemma~\ref{an:lem-sizes}(S7). It is used at least twice:
$167n+270n^2\le n^3$ in Lemma~\ref{ppi:gauge} (which needs $n\ge280$), and
$4000\,n^6(l+1)\le k/4$ for $l\le n^3$, i.e.\ $b^{4000n^6(l+1)f}\le\pbet$, in
Proposition~\ref{ppi:C} (which needs $k\ge32000\,n^9$). (ii) (R) is stated for all of $\pbeGs$, not only for $\pbeN(c,c')$. (iii) $\pbeN(c,c')$
is restricted to agreeing blocks, as (G) and (C) are proved only there; (S) is not
weakened by this, because the bound is proved with the sum restricted to agreeing $s$ (\S\ref{pbs:sec}). (iv)
The gauge, the common lattice and nearness are now defined. In (S), the constant in front of the sum
is $\nu=\pbvol(U_c\backslash U_\infty)$, the integral over $U_\infty$ of a cutoff for $U_c$, for the
Haar measure $du$ of convention (Cv3) of \S\ref{pb:ss-setting}; (S) as a statement does not depend on
the choice of $du$ (Remark~\ref{pbs:rem-nu}).
\end{remark}

\subsection{Cell ranks}\label{ppi:ss-cell}

In this subsection, as in \S\ref{pbg:sec}, hermitian forms are written linear in the second variable
and conjugate-linear in the first, the opposite of convention (Cv1) of \S\ref{pb:ss-setting}: we write
$h$ for the form $(x,y)\mapsto h(y,x)$. Isometries, orthogonality, isotropy, the matrix $Z(\gamma)$ up to
transposition and conjugation, the vanishing of the minors $\delta_l(\gamma)$ and of $\pbg(\gamma)$, and
$U(E^+)$ do not change under this replacement (\S\ref{pbv:ss-weak}(iv)), so the statements below hold
in either convention.

\begin{lemma}[Eichler elements]\label{ppi:eichler}
Let $0\le j<\varrho$, let $p\in\langle f_0,\dots,f_j\rangle$ and let $q\in Y_0$ be orthogonal to
$f_0,\dots,f_j$. Choose $\omega\in\auO_E$ with $\omega+\aubar\omega=1$ and put $\varsigma=\omega h(q,q)$.
Then
\[
\mathsf E_{p,q}(t)=t+h(q,t)\,p-h(p,t)\,(q+\varsigma p)
\]
lies in $U(E^+)$, and if $p=fx$ then $\sigma(\mathsf E_{p,q})=x_j\,h(q,f_{j+1})$. For $g\in G(E^+)$,
$g\mathsf E_{p,q}g^{-1}=\mathsf E_{gp,gq}$.
\end{lemma}
\begin{proof}
$p\in A$ is isotropic and $h(q,p)=0$. With these two facts and $\varsigma+\aubar\varsigma=h(q,q)$, a
direct expansion of $h(\mathsf E t,\mathsf Et')$ gives $h(t,t')$ (the Eichler transformation), so
$\mathsf E=\mathsf E_{p,q}$ is an isometry, defined over $E$. For $t\in A_i=\langle e_1,\dots,e_i
\rangle$ we have $h(p,t)=0$, and $h(q,t)=0$ if $i\le j+1$; so $\mathsf E$ fixes $A_{j+1}$ pointwise
and maps $t\in A_i$ to $t+h(q,t)p\in t+A_{j+1}$. For $t\in A^\perp$, $\mathsf Et=t+h(q,t)p\in t+A$.
So $\mathsf E$ preserves the flag of $A$, acts trivially on its graded pieces and on $A^\perp/A$:
$\mathsf E\in U(E^+)$. For $b<\varrho$, $\mathsf Ef_{b+1}-f_{b+1}=h(q,f_{b+1})p$, whose coefficient
at $f_b$ is $x_b\,h(q,f_{b+1})$; as $x_b=0$ for $b>j$ and $h(q,f_{b+1})=0$ for $b+1\le j$, only
$b=j$ remains in $\sigma(\mathsf E)$. The last claim holds because $g$ is an isometry.
\end{proof}

\begin{lemma}[Climbing]\label{ppi:climb}
Let $\gamma$ agree (in the sense of \S\ref{pbe:ss-char}) and let $(\mu,\mu')$ occur at $a\le\varrho$. If $a<\varrho$, then $(\mu,\mu')$
occurs at $a+1$. Hence it occurs at every $a'$ with $a\le a'\le\varrho$, and in all cases
$\mu\aubar\mu=\mu'\aubar{\mu'}$.
\end{lemma}
\begin{proof}
Let $fx=\gamma fy$ with $x,y$ supported on $\{0,\dots,a\}$, $x_a=\mu$, $y_a=\mu'$, and let $W_a=\langle
f_b,\gamma f_b:b\le a\rangle$. If $\mu f_{a+1}-\mu'\gamma f_{a+1}=\sum_{b\le a}(\alpha_bf_b-\beta_b
\gamma f_b)$, then $x'=\mu\,\epsilon_{a+1}-\sum_{b\le a}\alpha_b\epsilon_b$ and $y'=\mu'\epsilon_{a+1}-
\sum_{b\le a}\beta_b\epsilon_b$ (with $\epsilon_b$ the unit vectors) show that $(\mu,\mu')$ occurs at
$a+1$. So it suffices that every $E$-linear form $\lambda$ on $Y_0$ that vanishes on $W_a$ satisfies
$\Delta:=\mu\lambda(f_{a+1})-\mu'\lambda(\gamma f_{a+1})=0$. Write $\lambda=h(q,\cdot)$. Then $q\perp
f_b$ and $\gamma^{-1}q\perp f_b$ for $b\le a$. For $r\in E$ put $u=\mathsf E_{rfx,q}$ and $w=\mathsf
E_{rfy,\gamma^{-1}q}$ (Lemma~\ref{ppi:eichler}, with $h(\gamma^{-1}q,\gamma^{-1}q)=h(q,q)$, so the same
$\varsigma$). Both lie in $U(E^+)$, and $\gamma w\gamma^{-1}=\mathsf E_{r\gamma fy,q}=u$. So $u\in
U(E^+)\cap\gamma U(E^+)\gamma^{-1}$ and, $\gamma$ agreeing, $\psi_f(u)=\psi_f(w)$. By
Lemma~\ref{ppi:eichler}, $\sigma(u)=r\mu\lambda(f_{a+1})$ and $\sigma(w)=r\mu'h(\gamma^{-1}q,f_{a+1})=
r\mu'\lambda(\gamma f_{a+1})$. With $r=-M\varpi t$, $t\in E$, this gives $\psi_{E,f}(t\Delta)=1$ for
all $t\in E$. If $\Delta\ne0$, $\psi_{E,f}$ would be trivial on $E$, which is dense in $\auA_{E,f}$,
whereas the finite part of the standard character is not trivial. So $\Delta=0$.

By induction (trivially if $a=\varrho$) $(\mu,\mu')$ occurs at $\varrho$: $fx=\gamma fy$ with $x_\varrho=\mu$, $y_\varrho=\mu'$.
As $A$ is isotropic and orthogonal to $z$, $-N\mu\aubar\mu=h(fx,fx)=h(\gamma fy,\gamma fy)=h(fy,fy)=
-N\mu'\aubar{\mu'}$.
\end{proof}

\begin{proposition}[Cell ranks]\label{ppi:cell}
Let $\gamma\in\Gamma$ and $j=j(\gamma)$.
\begin{enumerate}
\item[(a)] If $\pbg(\gamma)\ne0$, then $\delta_n(\gamma)\ne0$.
\item[(b)] If $\gamma$ agrees and $\pbg(\gamma)=0$, then $j\le\varrho$, $\delta_j(\gamma)\ne0$ (with
$\delta_0=1$), and $\pbeAp\cap\gamma\pbeAp\not\subseteq A$.
\end{enumerate}
In particular \textup{(R)} holds, and every $\gamma\in\pbeGs$ has $\delta_{j(\gamma)}(\gamma)\ne0$ and
$j(\gamma)\le\varrho\iff\pbg(\gamma)=0$.
\end{proposition}
\begin{proof}
We use twice the following remark: \emph{if $y\in E^n$ and $\gamma fy$ is orthogonal to $f_0,\dots,
f_\varrho$, then $y_\varrho=0$ and $\gamma fy\in A$.} Indeed $(\pbeAp)^\perp=A\oplus X_0$, so $\gamma fy=
\alpha+\xi$ with $\alpha\in A$, $\xi\in X_0$, and $h(\xi,\xi)=h(\gamma fy,\gamma fy)=h(fy,fy)=-N|y_\varrho
|^2$; as $X_0$ is definite (positive at every real place), both sides vanish.

(a) If $\delta_n(\gamma)=0$ there is $y\ne0$ with $\gamma fy\perp f_a$ for all $a$; by the remark
$0\ne\gamma fy\in A\cap\gamma\pbeAp$, so $\pbg(\gamma)=0$ by Lemma~\ref{pbe:link}(a).

(b) By Lemma~\ref{pbe:link}(a) there is a common vector $fx=\gamma fy\ne0$. Let $l$ be its top index.
By Lemma~\ref{ppi:climb}, $|x_l|=|y_l|$, so both are non-zero; with $\zeta=x_l/y_l$, $(\zeta,1)$
occurs at $l$, hence at every $a\ge l$, and $\zeta\aubar\zeta=1$. So $j\le l\le\varrho$. At $a=\varrho\ge
j$ the definition of $j$ gives $fx=\gamma fy$ with $x_\varrho=\zeta\ne0$: this vector lies in $\pbeAp\cap
\gamma\pbeAp$ and not in $A$.

Let $j\ge1$ and suppose $\delta_j(\gamma)=0$. Then there is $y\ne0$, supported on $\{0,\dots,j-1\}$,
with $h(f_a,\gamma fy)=0$ for $a<j$. We show $h(f_a,\gamma fy)=0$ for $j\le a\le\varrho$ by induction
on $a$. By definition of $j$ there is $fx^{(a)}=\gamma fy^{(a)}$ with top index $a$, $x^{(a)}_a=\zeta$,
$y^{(a)}_a=1$. Then $h(fx^{(a)},\gamma fy)=h(fy^{(a)},fy)=-N\aubar{y^{(a)}_\varrho}y_\varrho=0$, as
$y_\varrho=0$; and $h(fx^{(a)},\gamma fy)=\sum_{b\le a}\aubar{x^{(a)}_b}h(f_b,\gamma fy)=\aubar\zeta\,
h(f_a,\gamma fy)$ by the choice of $y$ and the induction hypothesis. So $h(f_a,\gamma fy)=0$. By the
remark, $\gamma fy=fx\in A$ with $x\ne0$ (as $y\ne0$). Let $l$ be the top index of this common vector.
As above $x_l,y_l\ne0$ and $(x_l/y_l,1)$ occurs at every $a\ge l$ with $|x_l/y_l|=1$, so $l\ge j$ by
minimality; but $y_l\ne0$ forces $l<j$. This contradiction shows $\delta_j(\gamma)\ne0$.

The elements of $\pbeGs$ with $\pbg=0$ agree by definition, which gives the last sentence.
\end{proof}

\subsection{The gauge of a near term, and the gain}\label{ppi:ss-gain}

\begin{lemma}\label{ppi:gauge}
Let $\gamma\in\Gamma_{c,s,c'}$ be near. Then $|\tau(h(f_a,\gamma f_b))|\le b^{134n}$ for all $\tau,a,b$,
$\auNm\pbg(\gamma)\le b^{270n^2f}$, $\auNm\delta_l(\gamma)\le b^{270n^2f}$ for all $l$, and, if
$n\ge280$ and $\gamma\in\pbeGs$, $0<\mathsf G(\gamma)\le b^{n^3f}$.
\end{lemma}
\begin{proof}
Let $y$ and $\gamma^{-1}y$ be boxed; by Lemma~\ref{pbc:box}, $\|y_\tau^{-1}\tilde f_a\|,\|(\gamma^{-1}y)_\tau^{-1}\tilde f_a\|\le b^{37n}$. With the scales $s_a$ of \eqref{pbe:eq-scale}, $|s_a|\ge
N^{-1/2}$ (as $|\kappa_\tau|\ge1$), and by \eqref{pbe:eq-norm}
\[
|\tau(h(f_a,\gamma f_b))|=\frac{|h(y_\tau^{-1}\tilde f_a,(\gamma^{-1}y)_\tau^{-1}\tilde f_b)|}{|s_as_b|}
\le N\,b^{37n}\,b^{37n}\le b^{134n},
\]
using $N\le b^{60n}$. A minor of order $l\le n$ is a sum of $l!\le n^l\le b^l$ products of $l$
entries (as $n\le b$, Lemma~\ref{an:lem-sizes}(S7)(e)), so $|\tau(\delta_l)|^2\le b^{2l+268nl}\le b^{2n+268n^2}\le b^{270n^2}$. For the discriminant,
\[
|\tau(\pbg)|\le|\tau(\delta_n)|^2+N^2|\tau(\delta_{n-1})|^2\le b^{2n+268n^2}+b^{120n}\,b^{2(n-1)+268n(n-1)}
\le2\,b^{2n+268n^2}\le b^{270n^2},
\]
since $120n+2(n-1)+268n(n-1)=268n^2-146n-2\le2n+268n^2$, and $2\le b\le b^{2n^2-2n}$ as $2\le n\le b$
(Lemma~\ref{an:lem-sizes}(S7)(e) and \eqref{ppi:eq-sizes}). In particular (a weaker form used below and in \S\ref{pbs:sec}) $|\tau(\pbg)|\le b^{121n+270n^2}$. Taking products over the $f$ places,
and as $i_\gamma,i'_\gamma\ge1$, $N\le b^{60n}$ and $M\le b^{23}$,
\[
\mathsf G(\gamma)\le b^{46nf}\,b^{(121n+270n^2)f}\le b^{n^3f}\qquad(n\ge280).
\]
Positivity: by Proposition~\ref{ppi:cell}, $\delta_{j(\gamma)}(\gamma)\ne0$ for $\gamma\in\pbeGs$.
\end{proof}

\begin{proposition}[Gain, given the bound for the common parts]\label{ppi:G}
Assume the following statement \textup{(G$'$)}. Then \textup{(G)} holds, and $\mathsf d\ge1$.
\begin{enumerate}
\item[(G$'$)] Let $s$ agree, $\gamma,\gamma'\in\Gamma_{c,s,c'}\cap\pbeGs$, $C_\gamma=L_\gamma\cap
\pbeCL(\gamma,\gamma')$ and $t(\gamma)=[L\cap(A+k_\gamma A):L\cap A+k_\gamma(L\cap A)]=[h_cL\cap(A+\gamma A):h_cL\cap A+\gamma(h_cL\cap A)]$ (the
torsion index). Then $C_\gamma$ is a lattice of rank $\dim\pbeCL(\gamma,\gamma')$, $[C_\gamma:L_{\gamma
\gamma'}]$ divides $t(\gamma')$, and
\[
t(\gamma)\,\ppidisc C_\gamma\cdot t(\gamma')\,\ppidisc C_{\gamma'}\cdot\mathsf x^2\le(2/M)^{2if}\,
\mathsf G(\gamma)\mathsf G(\gamma').
\]
\end{enumerate}
\end{proposition}
\begin{proof}
$L_{\gamma\gamma'}\subseteq C_\gamma$ with index dividing $t(\gamma')$, so $L_{\gamma\gamma'}$ has the
same rank and $\mathsf d(\gamma,\gamma')=[C_\gamma:L_{\gamma\gamma'}]\ppidisc C_\gamma\le t(\gamma')
\ppidisc C_\gamma$; symmetrically ($L_{\gamma\gamma'}$ and $\pbeCL$ are symmetric in $\gamma,\gamma'$)
$\mathsf d\le t(\gamma)\ppidisc C_{\gamma'}$. Multiply and use (G$'$). The ideal generated by the
Gram determinants of $\ppiLc\cap V$ is integral, because $\ppiLc$ is an integral hermitian lattice;
so $\ppidisc(\ppiLc\cap V)\ge1$ and $\mathsf d\ge1$.
\end{proof}

Statement (G$'$) is Theorem~\ref{pbg:main}, proved in \S\ref{pbg:sec}.

\subsection{The count}\label{ppi:ss-count}

\begin{proposition}[Count, given two lattice estimates]\label{ppi:C}
Assume \eqref{ppi:eq-sizes}, \textup{(G$'$)} and the following two statements, for $0\le l\le n^3$,
$r=k-l$:
\begin{enumerate}
\item[(C1)] For a class $c$, a label $t\in\pbR$ and a lattice $L'\subseteq\ppiLc$ of rank $r$ with
$\ppidisc L'\le b^{n^3f}$, there are at most $b^{2000n^6(l+1)f}$ pairs $(c',\gamma)$ with $\gamma$ near,
chosen, of label $t$, in $\Gamma_{c,s,c'}\cap\pbeGs$ for some $s$, and with $L'\subseteq L_\gamma$.
\item[(C2)] $\displaystyle\sum_{c\in\mathcal C}\ \sum_{L'}\ (\ppidisc L')^{-k}\le\frac{\#\mathcal
C}{N^{kf}}\,\pbet^{\,l+1}$, the inner sum over the lattices $L'\subseteq\ppiLc$ of rank $r$ with
$\ppidisc L'\le b^{n^3f}$.
\end{enumerate}
Then \textup{(C)} holds.
\end{proposition}
\begin{proof}
Fix $c$. Send a pair $(\gamma,\gamma')$ counted in (C) for $c'$ to $(t,L_{\gamma\gamma'})$, where $t$
is the common label (both lie over $(c,s,c')$). By Proposition~\ref{ppi:G}, $L_{\gamma\gamma'}$ has
rank $\dim\pbeCL(\gamma,\gamma')=r$; it lies in $\ppiLc$, in $L_\gamma$ and in $L_{\gamma'}$, and
$\ppidisc L_{\gamma\gamma'}\le b^{n^3f}$. So both $(c',\gamma)$ and $(c',\gamma')$ lie in the set of
(C1) for $(t,L_{\gamma\gamma'})$, and the map $(c',\gamma,\gamma')\mapsto((c',\gamma),(c',\gamma'))$ is
injective. Hence each fibre has at most $b^{4000n^6(l+1)f}\le b^{8000n^9f}\le b^{kf/4}=\pbet$ elements, by
$l+1\le2n^3$ and \eqref{ppi:eq-sizes} ($32000\,n^9\le512000\,n^9\le k$). All pairs in a
fibre have the weight $(\ppidisc L_{\gamma\gamma'})^{-k}$, so
\[
\sum_{c'}\sum_{(\gamma,\gamma')}\mathsf d^{-k}\le\#\pbR\cdot\pbet\sum_{L'}(\ppidisc L')^{-k}.
\]
Summing over $c$ and using (C2), the total is at most $\#\pbR\cdot\pbet\cdot\#\mathcal CN^{-kf}
\pbet^{l+1}=\pbet^{l+2}\pbenp$.
\end{proof}

Statements (C1) and (C2) are proved in Theorems~\ref{pbc:C1} and~\ref{pbc:C2}; hence (C) holds, given (G$'$).

\subsection{Size}\label{ppi:ss-size}

\smallskip\noindent\textup{(S)} is Theorem~\ref{pbs:S}, proved in \S\ref{pbs:sec} from the open-cell bound \textup{(O1)} (Theorem~\ref{pbo:O1}).\par\smallskip

\subsection{The consequences still follow}\label{ppi:ss-check}

\begin{proposition}\label{ppi:check}
Lemmas~\ref{pbe:rank} and~\ref{pbe:pairs}, Proposition~\ref{pbe:two} and Theorem~\ref{pbe:XL1} hold
with Theorem~\ref{ppi:pillars} in place of Theorem~\ref{pbe:pillars}, with the same proofs.
\end{proposition}
\begin{proof}
The proofs in \S\ref{pbe:ss-XL1} use only the following, all contained in
Theorem~\ref{ppi:pillars}: (R) for $\gamma\in\pbeN(c,c')\subseteq\pbeGs$; $1\le\mathsf d$ and $\mathsf
G\le b^{n^3f}$; (G) with the same $\mathsf x$; (C) for every $l\le n^3$; and (S) for the stable sets
$T_\Lambda$, $T_0=\{\gamma\in\pbeGs:\pbg(\gamma)=0\}$ and $\pbeGs$, which are subsets of $\pbeGs$.
We check the points where the definitions enter. In Lemma~\ref{pbe:rank}(b), $\Upsilon_\gamma=
\Upsilon_{\gamma'}$ gives $X_0^\gamma=X_0^{\gamma'}$, and by Lemma~\ref{pbe:link}(c) and (R) $\dim
X_0^\gamma$ is $k-1$ or $k$ according as $\pbg\ne0$ or $\pbg=0$: unchanged. In
Lemma~\ref{pbe:rank}(c), $18i\le n^3$ gives $l\le i+1\le n^3$. In Lemma~\ref{pbe:pairs}, $\mathcal
D(\gamma)$ is $\pbDmain$ exactly at cell rank $\varrho+1$, so $\pbesz$ is the size used in Theorem~\ref{ppi:pillars}(S) (defined in \S\ref{pbe:ss-pillars}), and
$n^3+1\le b^4\le\pbet$. In Proposition~\ref{pbe:two}(a) with $i=0$, $\pbg(\gamma)=\pbg(\gamma')=0$ and
$\mathsf x=1$; with $i\ge1$, $\mathsf x\ge1$; in (b), $i=0$ and $\pbg\ne0$, so $\mathsf x=\mathsf{mn}(s)
\ge\mathsf z$. In Theorem~\ref{pbe:XL1}, a pair of elements of $\pbeN(c,c')\cap\Gamma_{c,s,c'}$ lies
over an agreeing block, as required by (G) and (C), and the sum in (S) runs over such elements only.
\end{proof}


%% ===== proofs/B-period.tex =====
% proofs/B-period.tex
% Proof of (DS4) of Theorem pbe:ds: the lower bound for the local Bessel period of the discrete
% series coefficient at a stationary point. Labels of this file start with pbp:.
\providecommand{\auU}{\mathrm U}
\providecommand{\aubar}[1]{\overline{#1}}
\providecommand{\pbvol}{\operatorname{vol}}
\providecommand{\pbip}[2]{\langle#1,#2\rangle}
\providecommand{\pbW}{\mathbb W}
\providecommand{\pbel}{\ell}
\providecommand{\pbeB}{\mathsf B}
\providecommand{\pdsK}{\mathbf K}
\providecommand{\pdsg}{\mathfrak g}
\providecommand{\pdsm}{m}
\providecommand{\pdsHS}[1]{\|#1\|_{2}}
\providecommand{\pcoD}{\mathsf D}
% macros of this file (prefix \pbp)
\providecommand{\pbpF}[1]{\|#1\|_{\mathrm F}}
\providecommand{\pbptr}{\operatorname{tr}}
\providecommand{\pbpct}{\operatorname{ct}}
\providecommand{\pbpCT}{\operatorname{CT}}
\providecommand{\pbpQ}{\mathcal Q}
\providecommand{\pbpLam}{\widehat\Lambda}
\providecommand{\pbpst}{\tilde\sigma}
\providecommand{\pbpI}{\mathsf I}

\section{The Bessel period at a stationary point}\label{pbp:sec}

This section proves \textup{(DS4)} of Theorem~\ref{pbe:ds}, stated in \S\ref{pds:ss-ds4}.
\emph{Dependencies.} Theorem~\ref{pbp:main} below, which is \textup{(DS4)}, uses:
the explicit formula \eqref{pds:eq-FJ} (Theorem~\ref{pco:FJ}) together with the independence of the
block diagonal factor (Proposition~\ref{pco:indep}) and the $K$-behaviour (Lemma~\ref{pco:K}) of
\S\ref{pco:sec}; the decay Theorem~\ref{pds:decay} and the lower bound of Lemma~\ref{pds:lip}; Witt's
theorem for hermitian forms (proved in the form used as Lemma~\ref{pbp:tb-witt}; see also \cite[Chapter~7, Theorem~9.1]{Scharlau}); Schwarz's lemma; and elementary linear algebra. The normalisation of
the character (Lemmas~\ref{pbp:psi}, \ref{pbp:sign}), the weight vector (Lemma~\ref{pbp:W}) and the
integrality of $\kappa(\tau)$ (Lemma~\ref{pbp:sizes}) are treated in \S\ref{pbp:ss-data}.

The argument is Laplace's method. A point $x$ is chosen (\S\S\ref{pbp:ss-pos}--\ref{pbp:ss-point})
so that, on the group $x^{-1}R_\tau x$, the phase of the character cancels the phases of the two
coefficients to first order (``stationary position''); near $1$ each coefficient is then its phase
times a number close to $[0,1]$ (\S\S\ref{pbp:ss-minor}--\ref{pbp:ss-gauss}, ``Gaussian shape'');
away from $1$ the integrand is either small or equal to a value taken near $1$
(\S\ref{pbp:ss-loc}); so the integrand lies in a cone about the positive axis up to an
exponentially small error, and a covering argument compares the integral with the volume of
$\pbeB_x$ (\S\ref{pbp:ss-cone}).

\subsection{Statement, conventions and sizes}\label{pbp:ss-stat}

Fix $\tau\in\auPl$; it is dropped from the notation when no confusion arises. A word on letters: $h$ denotes the gap below, and also the hermitian form in
\S\ref{pbp:ss-data} and an element of $H_\tau$ (or of $\auU(k)$) in Lemmas~\ref{pbp:R}, \ref{pbp:W}, \ref{pbp:compact} and in the
Reduction at the end of \S\ref{pbp:ss-data}, where the form itself does not occur; the image of
$s\in R_\tau$ in $H_\tau$ is written $h_s$. Likewise $k$
denotes an integer (a rank or a coordinate index) or an element of $\pdsK$; the meaning is fixed in
each lemma and the two never occur in the same formula. We use the set-up of
\S\ref{pds:ss-setup} (with rank $r$): $G=\auU(J)$ with $J=\operatorname{diag}(\varepsilon_1,\dots,
\varepsilon_r)$, $\pdsK=G\cap\auU(r)$, the parameter $\lambda$, $\rho$, $\mu=\lambda+\rho$, the
highest weight $\Lambda$ of the lowest $\pdsK$-type, the gap $h\ge\Delta=2r^8$, $L=\max_i|\lambda_i|$,
the products of minors $\pdsm_w$ of \eqref{pds:eq-m}, the block diagonal part $a(x)$, and the unitary
matrix $p$ with $p^{-1}\mathsf J_1p=J$, through which $G$ and $G_\tau$ are identified; thus
$c_\tau(g)=c(p^{-1}gp)$ for $g\in G_\tau$, where $c$ is the coefficient on $G$ given by
\eqref{pds:eq-FJ}. We regard $p$ as a map from $\C^r$, with the form $B(v,w)=v^*Jw$, to $Y_\tau$ with
the coordinates of the normalised basis; it is an isometry both for the hermitian forms and for
the Euclidean norms.

\emph{Norms.} For a matrix $X$, $\|X\|$ is the operator norm and $\pbpF X=(\pbptr X^*X)^{1/2}$ the
Frobenius norm. We use $\|X\|\le\pbpF X$, $\pbpF{XY}\le\pbpF X\pbpF Y$, $\pbpF{XY}\le\|X\|\pbpF Y$,
$|\pbptr XY|\le\pbpF X\pbpF Y$, $|X_{ij}|\le\pbpF X$, and that $\pbpF{\cdot}$ is invariant under
multiplication by unitary matrices; for an orthogonal projection $P$, $\pbpF{PX}\le\pbpF X$.

\emph{Coordinates of $G$.} By Lemma~\ref{pds:chamber} and the remark after it, the coordinates
$1,\dots,r$ of $\C^r$ split into a set $X$ of $k-1$ coordinates, on which $\varepsilon=+1$ and
$\Lambda=0$, and $2n$ further coordinates ($n=\varrho+1$, so $2n=2\varrho+2=m$), which we number
$0,1,\dots,2n-1$ in the order of decreasing $\Lambda$; on the coordinate numbered $i$ the sign is
$(-1)^i$, and the values $\mu_i:=\Lambda_i$ satisfy
\begin{equation}\label{pbp:eq-steep}
-3B\le\mu_{2n-1}<\dots<\mu_1<\mu_0\le-B,\qquad B:=\kappa(\tau)+\Delta/2 .
\end{equation}
(The values are integers; $|\Lambda_i-\lambda_i|<r$, the entries $\lambda_i=-\kappa(\tau)-j\Delta$,
$1\le j\le m$, decrease in steps of $\Delta>2r$, $r\le\Delta/2$ and $m\Delta+r\le2\kappa(\tau)+
\Delta/2$; the last two are comparisons of sizes, Lemma~\ref{pbp:sizes}.) To shorten we write
$\mu$ also for the diagonal matrix $\operatorname{diag}(\Lambda)$ when it acts on vectors; the
vector $\lambda+\rho$ of \S\ref{pds:ss-setup} will only appear as $\lambda+\rho$.

\emph{Coordinates of $H$.} $H_\tau=\auU(k)$ acts on the span of $x_1,\dots,x_k$; its coordinates are
$X'=\{1,\dots,k-1\}$ and the coordinate $k$. We fix a bijection between $X'$ and $X$ and use it to
identify them. We write $a_\tau$ for the number of columns of the diagram of $\lambda^X$ at $\tau$:
$a_\tau=a=xb$ if $\tau\in\auPa$ and $a_\tau=0$ otherwise (\S\ref{ss:au-spaces}), and
$\Lambda'=\operatorname{diag}(0,\dots,0,a_\tau)$.

\emph{Sizes.} Put
\begin{equation}\label{pbp:eq-T}
D=r+k,\qquad\beta=b^3,\qquad T=D^2\log\beta,\qquad h_0=\lceil10^{30}T^3\rceil,\qquad
\rho_0=(200\beta^5)^{-1}.
\end{equation}

\begin{lemma}\label{pbp:sizes}
\textup{(a)} $\kappa(\tau)$ is a positive integer, $1\le B$, $m\Delta+r\le\kappa(\tau)$,
$r\le\Delta/2$, and $3B+a_\tau\le b^2$. \textup{(b)} $L\le2b$, $\max_i|\Lambda_i|\le3b$, $8nr\le b$,
$r\le b$, $a_\tau\le b^2$. \textup{(c)} $\beta\ge10^{12}$, $\log\beta\ge1$, $r^2\le T$,
$18D^2\le h_0\le\Delta\le h$. \textup{(d)} $\rho_0\le(48b^2)^{-1}$ and $a_\tau\rho_0\le\frac12$.
\end{lemma}
\begin{proof}
(a) $\kappa(\tau)=b$ or $xb/x_1$, an integer because $q_1x_1$ divides $b$ (Definition~\ref{an:def-data},
(D4), (D6)); so $B\ge1$. By \eqref{eq:au-scales}, $x\le b^{10^{-12}}$, $r\le3x^5$
and $m\le2\sqrt x$, so $m\Delta+r\le4\sqrt x\,(3x^5)^8+3x^5\le x^{42}$, while $\kappa(\tau)\ge b/x_1
\ge b^{1-10^{-6}}$; and $r\le r^8=\Delta/2$. Finally $3B+a_\tau\le3b+3r^8+xb\le b^2$.
(b) $L=\kappa(\tau)+m\Delta\le2\kappa(\tau)\le2b$; $|\Lambda_i|\le L+r\le3b$; $8nr\le8r^2\le
72x^{10}\le b$; $a_\tau\le xb\le b^2$.
(c) $\beta=b^3$ is enormous; $T\ge r^2\log\beta\ge r^2$. Since $r/2\le k<r$, $D\le2r$ and $T\le
12r^2\log b$, so $10^{30}T^3+1\le1.8\cdot10^{33}r^6(\log b)^3\le2r^8$, because $r\ge x^5\ge
(\log b)^{500}$ and $\log b\ge128$ give $r^2\ge(\log b)^{1000}\ge10^{34}(\log b)^3$. Also $18D^2\le
72r^2\le h_0$, and $h\ge\Delta$ by \S\ref{pbe:ss-ds}.
(d) $48b^2\le200b^{15}$; $a_\tau\rho_0\le b^2/(200b^{15})$.
\end{proof}

\begin{theorem}[\textup{(DS4)}]\label{pbp:main}
There are $x_\tau\in G_\tau$ with $\pbel_\tau(x_\tau)\le b^{5n}$ and a unit vector $v_\tau\in
\pbW_\tau$ such that
\[
\pbvol(\pbeB_{x_\tau})\le b^{60(r+k)^2}\,\Bigl|\int_{R_\tau}c_\tau(x_\tau^{-1}s\,x_\tau)\,
\pbip{v_\tau}{\pi_R(s)^{-1}v_\tau}\,ds\Bigr| .
\]
\end{theorem}

This is \textup{(DS4)} as stated in Theorem~\ref{pbe:ds}. The proof is completed in
\S\ref{pbp:ss-cone}.

\subsection{The local data: the group, the character and the weight vector}\label{pbp:ss-data}

Write $A=\langle\tilde e_1,\dots,\tilde e_\varrho\rangle$, and recall the normalised frame
$(\tilde f_0,\dots,\tilde f_\varrho)=(\tilde e_1,\dots,\tilde e_\varrho,\tilde z)$ of
\S\ref{pbe:ss-real}; $A^\perp=A\oplus\langle x_1,\dots,x_k,\tilde z\rangle$. The form $h$ is linear
in the first variable; $h(\tilde e_a,\tilde e_b^*)=\delta_{ab}$, $h(x_i,x_j)=\delta_{ij}$,
$h(\tilde z,\tilde z)=-1$, all other products of basis vectors vanish.

\begin{lemma}\label{pbp:R}
$R_\tau$ is the set of $g\in G_\tau$ such that $g\tilde f_a\in\tilde f_a+\langle\tilde f_{a'}:a'<a
\rangle$ for $0\le a\le\varrho$. Every $s=uh\in R_\tau$ \textup($u\in U_\tau$, $h\in H_\tau$\textup)
satisfies $s\,x_i\in\sum_{i'}h_{i'i}x_{i'}+A$, where $(h_{i'i})$ is the matrix of $h$ in the basis
$(x_i)$. The map $s\mapsto F(s)$, where $F(s)$ is the unitriangular matrix of $s$ on the frame, is a
homomorphism; hence
\[
\tilde\sigma(s)=\sum_{a=1}^{\varrho}F(s)_{a-1,a}\qquad(\text{the coefficient of }\tilde f_{a-1}
\text{ in }s\tilde f_a,\text{ summed})
\]
is additive on $R_\tau$, and $\tilde\sigma(h)=0$ for $h\in H_\tau$.
\end{lemma}
\begin{proof}
Elements of $R_\tau$ have the stated form (\S\ref{pbe:ss-real}); and $u\in U_\tau$ preserves $A$, hence
$A^\perp$, and is trivial on $A^\perp/A$, so $ux_i\in x_i+A$, while $hx_i=\sum h_{i'i}x_{i'}$.
Conversely let $g$ be unitriangular on the frame. Then $g$ preserves each $\langle\tilde e_1,\dots,
\tilde e_a\rangle$, hence $A$ and $A^\perp$, and induces an isometry $\bar g$ of the nondegenerate
space $A^\perp/A=\langle\bar x_1,\dots,\bar x_k,\bar z\rangle$ with $\bar g\bar z=\bar z$. So $\bar g$
preserves $\bar z^\perp=\langle\bar x_i\rangle$ and acts there by a unitary matrix $h$; let $h\in
H_\tau$ also denote the element acting by $h$ on the $x_i$ and trivially on $\tilde e_a,\tilde z,
\tilde e_a^*$. Then $u=gh^{-1}$ preserves the flag $\langle\tilde e_1\rangle\subset\dots\subset A$, acts
trivially on its graded pieces (as $g$ is unitriangular and $h$ fixes the $\tilde e_a$) and trivially
on $A^\perp/A$; so $u\in U_\tau$ and $g=uh\in R_\tau$. The span of the frame is $A\oplus\langle\tilde z
\rangle$ and is preserved by $R_\tau$, so $F(s_1s_2)=F(s_1)F(s_2)$; the entries just above the
diagonal of a product of unitriangular matrices are the sums of those of the factors.
\end{proof}

\begin{lemma}[The character]\label{pbp:psi}
Let $\sigma$ be as in \eqref{pb:eq-psi}, $\kappa_\tau$ as in \eqref{pbe:eq-kappa}. For $u\in U_\tau$,
$\tilde\sigma(u)=\sigma(u)/\kappa_\tau$, and
\[
\aubar{\psi_\tau(u)}=\exp\bigl(\epsilon_0\,(\aubar{\tilde\sigma(u)}-\tilde\sigma(u))\bigr),
\]
where $\epsilon_0=+1$ if the component of $\psi_E$ at the place $\tau$ is $w\mapsto e^{-2\pi i(w+
\bar w)}$ and $\epsilon_0=-1$ if it is $w\mapsto e^{2\pi i(w+\bar w)}$.
\end{lemma}
\begin{proof}
Here $N$ is the level of \S\ref{pbe:ss-real}. From $\tilde e_a=\kappa_\tau^{\varrho-a+1}e_a/\sqrt N$, $\tilde e_a^*=\sqrt N e_a^*/\kappa_\tau^{\varrho-a+1}$
and $\tilde z=z/\sqrt N$ ($\kappa_\tau$ is real), $h(u\tilde e_{a+1},\tilde e_a^*)=h(ue_{a+1},
e_a^*)/\kappa_\tau$ and $h(u\tilde z,\tilde e_\varrho^*)=h(uz,e_\varrho^*)/\kappa_\tau$. Since $h(\cdot,
\tilde e_a^*)$ is the coefficient of $\tilde e_a$, these are the entries $F(u)_{a-1,a}$, and
$\tilde\sigma(u)=\sigma(u)/\kappa_\tau$. At $\tau$, $\psi_\tau(u)$ is the component of $\psi_E$ at
$w=-\sigma(u)/(M\tau(\varpi))$, where now $\sigma(u)$ is computed in $Y_\tau$. As $M\tau(\varpi)=
2\pi i\kappa_\tau$, $2\pi iw=-\sigma(u)/\kappa_\tau=-\tilde\sigma(u)$ and $2\pi i\bar w=
\aubar{\tilde\sigma(u)}$. So $\psi_\tau(u)=\exp(\mp(\aubar{\tilde\sigma}-\tilde\sigma))$, with the upper
sign for $e^{-2\pi i(w+\bar w)}$; the exponent is imaginary, and conjugation changes its sign.
\end{proof}

The paper uses the normalisation of $\psi_E$ only through \S\ref{pb:ss-layer2}; the next lemma shows
that the sign $\epsilon_0$ does not affect \textup{(DS4)}.

\begin{lemma}[The sign]\label{pbp:sign}
Let $\delta\in\operatorname{End}(Y_\tau)$ be $\delta\tilde e_a=(-1)^a\tilde e_a$, $\delta\tilde e_a^*=
(-1)^a\tilde e_a^*$, $\delta\tilde z=(-1)^{\varrho+1}\tilde z$, $\delta x_i=x_i$. Then $\delta\in G_\tau$,
$\delta$ is an isometry for $\|\cdot\|$, $\delta R_\tau\delta^{-1}=R_\tau$, $s\mapsto\delta s\delta^{-1}$
preserves the Haar measure of $R_\tau$, $\tilde\sigma(\delta s\delta^{-1})=-\tilde\sigma(s)$, the
images of $s$ and of $\delta s\delta^{-1}$ in $H_\tau$ coincide, $\pbel_\tau(\delta^{-1}y)=\pbel_\tau(y)$
and $\pbvol(\pbeB_{\delta^{-1}y})=\pbvol(\pbeB_y)$ for $y\in G_\tau$.
\end{lemma}
\begin{proof}
$\delta$ is diagonal with entries $\pm1$ in the orthonormal normalised basis and preserves the
products of the basis vectors, so it is in $G_\tau$ and isometric. We have $\delta\tilde f_a=
(-1)^{a+1}\tilde f_a$ for all $0\le a\le\varrho$; so $\delta s\delta^{-1}$ is unitriangular on the frame
with $F(\delta s\delta^{-1})_{a'a}=(-1)^{a+a'}F(s)_{a'a}$, hence lies in $R_\tau$
(Lemma~\ref{pbp:R}), and the entries just above the diagonal change sign. As $\delta$ is trivial on
the $x_i$ and preserves $A$, the actions of $s$ and $\delta s\delta^{-1}$ on $A^\perp/A$ agree on the
$\bar x_i$. The automorphism $s\mapsto\delta s\delta^{-1}$ has order $2$, so it multiplies the Haar
measure by a positive number whose square is $1$. Finally $\|s\delta^{-1}y\|=\|\delta^{-1}(\delta s
\delta^{-1})y\|=\|(\delta s\delta^{-1})y\|$, and $\pbeB_{\delta^{-1}y}=\delta^{-1}\pbeB_y\delta$.
\end{proof}

\begin{lemma}[The weight vector]\label{pbp:W}
As a representation of $H_\tau=\auU(k)$ in the basis $(x_i)$, $\pbW_\tau$ is isomorphic to the space
of homogeneous polynomials of degree $a_\tau$ in $k$ variables, with $(\tau_\tau(h)P)(t)=P(h^{\mathrm t}
t)$, i.e.\ $\tau_\tau(h)$ sends the monomial $t_j^{a_\tau}$ to $(\sum_ih_{ij}t_i)^{a_\tau}$; its
weights are the $\sum_i\alpha_ie_i$ with $\alpha_i\ge0$, $\sum\alpha_i=a_\tau$, and the highest weight
is $(a_\tau,0,\dots,0)$. For the unit vector $v_\tau=t_k^{a_\tau}/\|t_k^{a_\tau}\|$ \textup(any
invariant inner product\textup),
\begin{equation}\label{pbp:eq-cu}
c'(h):=\pbip{v_\tau}{\tau_\tau(h)v_\tau}=h_{kk}^{\,a_\tau}\qquad(h\in H_\tau).
\end{equation}
\end{lemma}
\begin{proof}
By \S\ref{ss:au-spaces}, $W_{\lambda^X}$ at $\tau$ is the dual $V_\tau$ of the space of forms of degree
$a_\tau$, in its model as the dual of the Schur module spanned by the products of minors ($1\times1$
minors here), on which $\mathrm{GL}_k(\C)$ acts polynomially by the Schur matrices; so as a
$\mathrm{GL}_k(\C)$-module $W_{\lambda^X,\tau}$ is the polynomial representation $\operatorname{Sym}^{a_\tau}$
of highest weight $(a_\tau,0,\dots,0)$. The group $H_\tau$ acts on it through $h\mapsto\tau(\bar h)=
{}^{\mathrm t}h^{-1}$ (\S\ref{ss:au-spaces}). The space $\pbW_\tau$ is the dual of $W_{\lambda^X,\tau}$ for the
bilinear pairing $\pbip w\xi=\sum w_i\xi_i$ with $\pbip{\tau_\infty(h)w}{h\xi}=\pbip w\xi$
(\S\ref{pb:ss-layer2}); so $\tau_\tau(h)$ is the transpose inverse of $\operatorname{Sym}^{a_\tau}
({}^{\mathrm t}h^{-1})$, which is $\operatorname{Sym}^{a_\tau}(h)$. In the model of polynomials this is the
stated action. The monomials are weight vectors and are orthogonal for any $\auU(k)$-invariant inner
product (distinct weights); the coefficient of $t_k^{a_\tau}$ in $(\sum_ih_{ik}t_i)^{a_\tau}$ is
$h_{kk}^{a_\tau}$, which gives \eqref{pbp:eq-cu}.
\end{proof}

\begin{remark}[What is used of the data]\label{pbp:rem-doubts}
We record: (1) Sign and scaling: by
Lemma~\ref{pbp:psi} the character in the normalised frame is $\exp(\pm(\aubar{\tilde\sigma}-
\tilde\sigma))$ with $\tilde\sigma=\sigma/\kappa_\tau$, i.e.\ the scaling by $\kappa_\tau$ is exactly
absorbed by the normalised frame; the sign depends on the normalisation of $\psi_E$ at infinity and is
irrelevant by Lemma~\ref{pbp:sign}. (2) $\pbW_\tau\cong\operatorname{Sym}^{a_\tau}$ has the extremal weight
vector $t_k^{a_\tau}$ of weight $a_\tau e_k$, a single nonnegative entry (Lemma~\ref{pbp:W}); this
uses the model of $W_\lambda$ in \S\ref{ss:au-spaces} (the Schur matrices are polynomial in $g$). (3) $\kappa(\tau)$ is an integer by (D4) and (D6) of
Definition~\ref{an:def-data} (Lemma~\ref{pbp:sizes}).
\end{remark}

\emph{Reduction.} Let $y\in G_\tau$ and $v_\tau$ be as in Lemma~\ref{pbp:W}. For $s=uh\in R_\tau$,
$\pi_R(s)^{-1}=\tau_\tau(h)^{-1}\psi_\tau(u)^{-1}$ and $\pbip{v_\tau}{\tau_\tau(h^{-1})v_\tau}=
\aubar{c'(h)}$, so by Lemmas~\ref{pbp:R} and \ref{pbp:psi}
\begin{equation}\label{pbp:eq-integrand}
c_\tau(y^{-1}sy)\pbip{v_\tau}{\pi_R(s)^{-1}v_\tau}=c_\tau(y^{-1}sy)\,\aubar{c'(h_s)}\,
\exp\bigl(\epsilon_0(\aubar{\tilde\sigma(s)}-\tilde\sigma(s))\bigr),
\end{equation}
$h_s$ the image of $s$ in $H_\tau$. By Lemma~\ref{pbp:sign}, the substitution $s\mapsto\delta s
\delta^{-1}$ turns the integral of \eqref{pbp:eq-integrand} for $\epsilon_0$ and $y$ into the integral
for $-\epsilon_0$ and $\delta^{-1}y$, without changing $\pbel_\tau$ and $\pbvol(\pbeB)$. \emph{It
therefore suffices to prove Theorem~\ref{pbp:main} with the integrand
\eqref{pbp:eq-integrand} for $\epsilon_0=+1$}, which we assume from now on.

\subsection{Positions and the two traces}\label{pbp:ss-pos}

For $l\in X$ let $e_l$ be the corresponding standard basis vector of $\C^r$.

\begin{definition}\label{pbp:def-pos}
A \emph{stationary position} is a bijective linear map $\iota\colon\C^r\to Y_\tau$ with
$h(\iota v,\iota w)=B(w,v)$ \textup(an isometry\textup) such that, for $\widehat\Lambda=\iota\mu
\iota^{-1}$,
\begin{enumerate}
\item[(S1)] $\widehat\Lambda\tilde f_a=\tilde f_{a+1}$ for $0\le a<\varrho$;
\item[(S2)] $h(\widehat\Lambda x_i,x_{i'})=\Lambda'_i\delta_{ii'}$ for $1\le i,i'\le k$;
\item[(S3)] $\iota e_l=x_l$ for $l\in X$.
\end{enumerate}
For $s\in R_\tau$ we then put $g_s=\iota^{-1}s\iota\in G$ and write $h_s\in\auU(k)$ for the image of
$s$ in $H_\tau$ (Lemma~\ref{pbp:R}), and
\[
\varphi(s)=\pbptr\bigl(\mu(g_s-1)\bigr)-\pbptr\bigl(\Lambda'(h_s-1)\bigr).
\]
\end{definition}

(The order of the arguments of $B$ only reflects that $h$ is linear in the first variable and $B$ in
the second.) Since $\mu$ is real diagonal and commutes with $J$, $\widehat\Lambda$ is self-adjoint for
$h$; and $h(sw,w')=h(w,s^{-1}w')$ for $s\in G_\tau$.

\begin{lemma}[Two traces]\label{pbp:traces}
Let $\iota$ be a stationary position and $s\in R_\tau$. Then
\begin{enumerate}
\item[(a)] $\pbptr(g_s-1)=\pbptr(h_s-1)$;
\item[(b)] $\varphi(s)=\tilde\sigma(s)-\aubar{\tilde\sigma(s)}$; in particular $\varphi$ is additive on
$R_\tau$ and purely imaginary;
\item[(c)] for every unitary $u$ on the coordinates $X$, the element $s_u\in H_\tau$ acting by $u$ on
$\langle x_l:l\in X\rangle$ and trivially on $x_k$, $\tilde e_a$, $\tilde z$, $\tilde e_a^*$ satisfies:
$g_{s_u}$ is $u$ on the coordinates $X$ and the identity on the others, $h_{s_u}=u\oplus1$, and
$\varphi(s_u)=0$.
\end{enumerate}
\end{lemma}
\begin{proof}
The coefficients of a vector $w$ on the basis $\tilde e_a,\tilde e_a^*,\tilde z,x_i$ are $h(w,\tilde
e_a^*)$, $h(w,\tilde e_a)$, $-h(w,\tilde z)$, $h(w,x_i)$, so for $T\in\operatorname{End}(Y_\tau)$
\begin{equation}\label{pbp:eq-trace}
\pbptr T=\sum_{a=1}^\varrho\bigl(h(T\tilde e_a,\tilde e_a^*)+h(T\tilde e_a^*,\tilde e_a)\bigr)
-h(T\tilde z,\tilde z)+\sum_{i=1}^kh(Tx_i,x_i).
\end{equation}
Let $Y=s-1$ and $Y^\dagger=s^{-1}-1$, so $h(Yw,w')=h(w,Y^\dagger w')$; $s^{-1}\in R_\tau$. By
Lemma~\ref{pbp:R}, $Y\tilde f_a\in\langle\tilde f_{a'}:a'<a\rangle$, $Y\tilde z\in A$ and $Yx_i\in
\sum_{i'}(h_s-1)_{i'i}x_{i'}+A$, and the same holds for $Y^\dagger$ with $h_s^{-1}$.

(b) Apply \eqref{pbp:eq-trace} to $T=\widehat\Lambda Y$; then $\pbptr T=\pbptr(\mu(g_s-1))$. First,
$Y\tilde e_a\in\langle\tilde e_1,\dots,\tilde e_{a-1}\rangle$ and $\widehat\Lambda\tilde e_b=\tilde e_{b+1}$
for $b<\varrho$, so $h(\widehat\Lambda Y\tilde e_a,\tilde e_a^*)$ is the coefficient of $\tilde e_{a-1}$ in
$Y\tilde e_a$, i.e.\ $F(s)_{a-2,a-1}$ ($2\le a\le\varrho$; zero for $a=1$). Second, write $Y\tilde z=
\sum_bc_b\tilde e_b$; $-h(\widehat\Lambda Y\tilde z,\tilde z)=-\sum_bc_bh(\widehat\Lambda\tilde e_b,
\tilde z)$, and $h(\widehat\Lambda\tilde e_b,\tilde z)$ is $h(\tilde e_{b+1},\tilde z)=0$ for $b<\varrho$
and $h(\tilde z,\tilde z)=-1$ for $b=\varrho$ (by (S1), $\widehat\Lambda\tilde f_{\varrho-1}=\tilde f_\varrho$);
so this term is $c_\varrho=F(s)_{\varrho-1,\varrho}$. The two give $\tilde\sigma(s)$. Third, $h(\widehat
\Lambda Yx_i,x_i)=h(Yx_i,\widehat\Lambda x_i)$; for $\alpha\in A$, $h(\alpha,\widehat\Lambda x_i)=
h(\widehat\Lambda\alpha,x_i)=0$ because $\widehat\Lambda A\subseteq\langle\tilde e_2,\dots,\tilde e_\varrho,
\tilde z\rangle\perp x_i$; and $h(x_{i'},\widehat\Lambda x_i)=\aubar{h(\widehat\Lambda x_i,x_{i'})}=
\Lambda'_i\delta_{ii'}$ by (S2). So the third sum is $\sum_i(h_s-1)_{ii}\Lambda'_i=\pbptr(\Lambda'
(h_s-1))$. Fourth, $h(\widehat\Lambda Y\tilde e_a^*,\tilde e_a)=h(\tilde e_a^*,Y^\dagger\widehat\Lambda
\tilde e_a)=\aubar{h(Y^\dagger\widehat\Lambda\tilde e_a,\tilde e_a^*)}$, and $\widehat\Lambda\tilde e_a=
\tilde f_a$; summing over $a$ gives $\aubar{\tilde\sigma(s^{-1})}=-\aubar{\tilde\sigma(s)}$
(Lemma~\ref{pbp:R}). Hence $\pbptr(\mu(g_s-1))=\tilde\sigma(s)+\pbptr(\Lambda'(h_s-1))-
\aubar{\tilde\sigma(s)}$.

(a) The same computation with $\widehat\Lambda$ replaced by the identity: $h(Y\tilde e_a,\tilde e_a^*)=0$,
$h(Y\tilde z,\tilde z)=0$ ($Y\tilde z\in A\perp\tilde z$), the third sum is $\pbptr(h_s-1)$, and
$h(Y\tilde e_a^*,\tilde e_a)=\aubar{h(Y^\dagger\tilde e_a,\tilde e_a^*)}=0$.

(c) By (S3), $g_{s_u}e_l=\iota^{-1}s_ux_l=\sum_{l'}u_{l'l}e_{l'}$ for $l\in X$. The isometry $\iota$ maps
the $B$-orthogonal complement of $\langle e_l:l\in X\rangle$, which is spanned by the other coordinate
vectors, onto the $h$-orthogonal complement of $\langle x_l:l\in X\rangle$, which $s_u$ fixes
pointwise. Finally $s_u\in H_\tau$ fixes the frame, so $\tilde\sigma(s_u)=0$ and $\varphi(s_u)=0$ by (b).
\end{proof}

\emph{The integrand.} If $\iota$ is a stationary position and $y=\iota p^{-1}\in G_\tau$, then
$c_\tau(y^{-1}sy)=c(p^{-1}y^{-1}syp)=c(g_s)$, and by Lemma~\ref{pbp:traces}(b) the integrand
\eqref{pbp:eq-integrand} (with $\epsilon_0=1$) is
\begin{equation}\label{pbp:eq-theta}
\theta(s)=c(g_s)\,\aubar{c'(h_s)}\,e^{-\varphi(s)},\qquad|\theta(s)|=|c(g_s)|\,|c'(h_s)|.
\end{equation}

\subsection{The stationary position}\label{pbp:ss-stationary}

Let $\mu_0>\dots>\mu_{2n-1}$ be as in \eqref{pbp:eq-steep}, $p(t)=\prod_{i=0}^{2n-1}(t-\mu_i)$ and
$q(t)=t-a_\tau$. As $a_\tau\ge0>\mu_i$, $q(\mu_i)<0$ and $p(a_\tau)>0$. For a polynomial $F$ let
$[F]\in\R^r\subset\C^r$ be the vector that vanishes on $X$ and has the entry
\[
[F]_i=F(\mu_i)\,\bigl|p'(\mu_i)q(\mu_i)\bigr|^{-1/2}
\]
on the coordinate numbered $i$. Then $\mu[F]=[tF]$. For real vectors $v,w$ vanishing on $X$,
$B(v,w)=\sum_i(-1)^iv_iw_i$.

\begin{lemma}[Residue sums]\label{pbp:residue}
Let $F$ be a polynomial of degree at most $2n-1$ and $c\in\R$ not a node. Then
$\sum_iF(\mu_i)/p'(\mu_i)$ is the coefficient of $t^{2n-1}$ in $F$, and $\sum_iF(\mu_i)/((\mu_i-c)
p'(\mu_i))=-F(c)/p(c)$.
\end{lemma}
\begin{proof}
By Lagrange interpolation $F(t)=\sum_iF(\mu_i)\prod_{i'\ne i}(t-\mu_{i'})/p'(\mu_i)$; compare the
coefficients of $t^{2n-1}$. For the second, $H(t)=(F(t)p(c)-F(c)p(t))/(t-c)$ is a polynomial of degree
at most $2n-1$ whose coefficient of $t^{2n-1}$ is $-F(c)$ ($p$ is monic of degree $2n$, $\deg F<2n$), and
$H(\mu_i)=F(\mu_i)p(c)/(\mu_i-c)$; apply the first statement to $H$.
\end{proof}

\begin{lemma}\label{pbp:gram}
\textup{(i)} For $a'+b'\le2\varrho$, $B([t^{a'}q],[t^{b'}q])=-1$ if $a'+b'=2\varrho$ and $0$ otherwise.
\textup{(ii)} $B([1],[t^bq])=0$ for $b\le2\varrho$. \textup{(iii)} $B([1],[1])=1/p(a_\tau)$.
\textup{(iv)} For $b\le2\varrho$ every entry of $[t^bq]$ has modulus at most $(3B+a_\tau)^{2n}$.
\end{lemma}
\begin{proof}
Among the $\mu_{i'}$, exactly $i$ exceed $\mu_i$, so $p'(\mu_i)=\prod_{i'\ne i}(\mu_i-\mu_{i'})$ has the
sign $(-1)^i$, and $|p'(\mu_i)q(\mu_i)|=(-1)^{i+1}p'(\mu_i)q(\mu_i)$. Hence
$B([F_1],[F_2])=-\sum_iF_1F_2(\mu_i)/(p'(\mu_i)q(\mu_i))$. (i) This is $-\sum_i\mu_i^{a'+b'}q(\mu_i)/
p'(\mu_i)$, minus the coefficient of $t^{2n-1}$ in $t^{a'+b'}q$, which has degree $a'+b'+1\le2n-1$.
(ii) It is $-\sum_i\mu_i^b/p'(\mu_i)=0$ since $b<2n-1$. (iii) It is $-\sum_i1/((\mu_i-a_\tau)
p'(\mu_i))=1/p(a_\tau)$ by Lemma~\ref{pbp:residue}. (iv) The $\mu_i$ are distinct integers, so
$|p'(\mu_i)|\ge1$, and $|[t^bq]_i|\le|\mu_i|^b|q(\mu_i)|^{1/2}\le(3B)^b(3B+a_\tau)^{1/2}$.
\end{proof}

\begin{lemma}[Witt's theorem in the form used]\label{pbp:tb-witt}
Let $F$ be a field of characteristic $\ne2$ with an involution $a\mapsto\aubar a$ \textup(possibly
trivial\textup), $F_0$ its fixed field, and $V,V'$ finite-dimensional $F$-spaces with non-degenerate
hermitian forms $h,h'$ \textup(linear in the first variable\textup) that are isometric.
\textup{(a)} If $e\in V$, $e'\in V$ and $h(e,e)=h(e',e')\ne0$, some isometry $g$ of $V$ has $ge=e'$.
\textup{(b)} Every injective linear isometry $\phi\colon S\to V'$ from a subspace $S\subseteq V$ extends to
a bijective isometry $V\to V'$.
\end{lemma}
\begin{proof}
For $u\in V$ with $H=h(u,u)\ne0$ and $\varepsilon\in F$ with $\varepsilon\aubar\varepsilon=1$, the map
$\sigma(x)=x-kh(x,u)u$, $k=(1-\varepsilon)/H$, satisfies $h(\sigma x,\sigma y)=h(x,y)+h(x,u)h(u,y)
(|k|^2H-k-\aubar k)=h(x,y)$, since $(1-\varepsilon)(1-\aubar\varepsilon)=(1-\varepsilon)+(1-\aubar
\varepsilon)$; it fixes $u^\perp$ and is injective, hence an isometry of $V$.

(a) Let $c=h(e,e)\in F_0$ and $b=h(e,e')$. For $u=e-e'$ put $z=h(e,u)=c-b$; then $h(u,u)=z+\aubar z$.
If $z+\aubar z\ne0$, take $\varepsilon=-\aubar z/z$; then $1-\varepsilon=(z+\aubar z)/z$, so
$\sigma(e)=e-u=e'$. Otherwise $b+\aubar b=2c$, and the same construction for the pair $(e,-e')$ applies,
because there $z+\aubar z=2c+b+\aubar b=4c\ne0$; it gives $g$ with $ge=-e'$, and $-g$ works.

(b) Composing with an isometry $V'\to V$ we may take $V'=V$. \emph{Cancellation.} A non-degenerate
space $X\ne0$ has an anisotropic vector (if $h(x,x)=0$ for all $x$ then $h(x,y)+h(y,x)=0$; replacing $y$
by $\lambda y$ with $\lambda\ne\aubar\lambda$, or using $h(x,y)=h(y,x)$ when the involution is trivial,
gives $h=0$), hence an orthogonal basis. If $X=Fe\perp U=Fe'\perp U'$ with $h(e,e)=h(e',e')\ne0$, then by
(a) $U'=e'^\perp=g(e^\perp)=g(U)\cong U$. By induction over an orthogonal basis: if $T,T'\subseteq V$ are
isometric non-degenerate subspaces, then $T^\perp\cong T'^\perp$. \emph{Non-degenerate $S$.} Then
$V=S\perp S^\perp=\phi S\perp(\phi S)^\perp$, and $\phi$ together with an isometry $S^\perp\to(\phi
S)^\perp$ is the extension. \emph{General $S$.} Let $R$ be the radical of $S$, with basis $r_1,\dots,
r_k$, and $S=R\oplus S_0$; $S_0$ is non-degenerate and $R\subseteq S_0^\perp$, which is non-degenerate.
So there are $y'_j\in S_0^\perp$ with $h(r_i,y'_j)=\delta_{ij}$; with $g_{ij}=h(y'_i,y'_j)$ the vectors
$y_j=y'_j-\tfrac12\sum_lg_{jl}r_l$ satisfy $h(r_i,y_j)=\delta_{ij}$ and $h(y_i,y_j)=g_{ij}-\tfrac12
g_{ij}-\tfrac12\aubar g_{ji}=0$. Then $\tilde S=S_0\oplus\langle r_i,y_i\rangle$ is non-degenerate.
Since $\phi(R)$ is the radical of $\phi(S)=\phi(R)\oplus\phi(S_0)$, the same construction gives $y''_j$
for $\phi(r_j)$ in $\phi(S_0)^\perp$, and $\phi$ extended by $y_j\mapsto y''_j$ is an injective isometry
of $\tilde S$; apply the non-degenerate case. (See also \cite[Chapter~7, Theorem~9.1]{Scharlau}.)
\end{proof}

\begin{proposition}\label{pbp:stationary}
There is a stationary position $\iota$ such that $\iota^{-1}\tilde f_a=\mu^av_1$ for $0\le a\le\varrho$,
where $v_1=[q]$; moreover, for $a,b\le2\varrho$, $\|\mu^av_1\|\le M_0:=\sqrt r\,(3B+a_\tau)^{2n}$ and
$B(\mu^av_1,\mu^bv_1)=-1$ if $a+b=2\varrho$, $0$ if $a+b<2\varrho$.
\end{proposition}
\begin{proof}
Put $v_j=[t^{j-1}q]=\mu^{j-1}v_1$ ($1\le j\le n$) and $w=p(a_\tau)^{1/2}[1]$. The vectors $v_1,\dots,
v_n,w$ are the evaluation vectors, scaled by nonzero factors, of the $n+1$ polynomials $t^{j-1}q$, $1$
of degree at most $n<2n$ at the $2n$ distinct nodes, so they are linearly independent; with the $e_l$,
$l\in X$, which have disjoint supports, they form a basis of a subspace $S\subseteq\C^r$. By
Lemma~\ref{pbp:gram} and disjointness of supports, the Gram matrix of $(v_1,\dots,v_n;e_l\,(l\in X);w)$
for $B$ is that of $(\tilde f_0,\dots,\tilde f_\varrho;x_l\,(l\in X);x_k)$ for $h$: all products vanish
except $B(v_n,v_n)=-1=h(\tilde z,\tilde z)$, $B(e_l,e_l)=1$ and $B(w,w)=1$ ((i) with $a'=b'=\varrho$,
(ii), (iii)). The second family is also linearly independent. By Witt's theorem (Lemma~\ref{pbp:tb-witt}(b),
with $V=\C^r$ and $V'=Y_\tau$, which are isometric via $p$), the linear map $S\to Y_\tau$ sending the
first family to the second extends to a bijective isometry $\iota\colon\C^r\to Y_\tau$. (S3) holds by
construction, and $\iota\mu^av_1=\iota v_{a+1}=\tilde f_a$. (S1): $\widehat\Lambda\tilde f_a=\iota\mu
v_{a+1}=\iota v_{a+2}=\tilde f_{a+1}$ for $a<\varrho$. (S2): $\iota^{-1}x_l=e_l$ and $\mu e_l=0$ for
$l\in X$; and $\mu w=p(a_\tau)^{1/2}[t]=p(a_\tau)^{1/2}[q]+a_\tau w=p(a_\tau)^{1/2}v_1+a_\tau w$. So
$h(\widehat\Lambda x_i,x_{i'})=B(\iota^{-1}x_{i'},\mu\iota^{-1}x_i)$ is $0$ for $i\in X$, and for $i=k$
it is $a_\tau B(\iota^{-1}x_{i'},w)+p(a_\tau)^{1/2}B(\iota^{-1}x_{i'},v_1)=a_\tau\delta_{i'k}$. The last
claims are Lemma~\ref{pbp:gram}(i),(iv) ($[t^bq]$ has at most $r$ entries).
\end{proof}

\subsection{The point and its length}\label{pbp:ss-point}

Let $\mathsf J_1$ also denote the operator on $Y_\tau$ with matrix $\mathsf J_1$ in the normalised basis,
so that $h(v,w)=\pbip{\mathsf J_1w}{v}$; it is unitary and self-adjoint, $\mathsf J_1\tilde e_a=\tilde
e_a^*$, $\mathsf J_1\tilde e_a^*=\tilde e_a$, $\mathsf J_1x_i=x_i$, $\mathsf J_1\tilde z=-\tilde z$, and
$g\mathsf J_1g^*=\mathsf J_1$ for $g\in G_\tau$ (because $g^{-1}=\mathsf J_1g^*\mathsf J_1$).

\begin{lemma}\label{pbp:frame}
For every unitriangular $(\varrho+1)\times(\varrho+1)$ matrix $T$ there is $s\in R_\tau$ with $F(s)=T$.
\end{lemma}
\begin{proof}
Define a linear map on $A^\perp=\langle\tilde f_0,\dots,\tilde f_\varrho,x_1,\dots,x_k\rangle$ by
$\tilde f_a\mapsto\sum_{a'}T_{a'a}\tilde f_{a'}$, $x_i\mapsto x_i$. It is injective and preserves $h$:
the images of $\tilde e_1,\dots,\tilde e_\varrho$ lie in $A$, which is isotropic and orthogonal to
$\tilde z$ and the $x_i$; the image of $\tilde z$ is $\tilde z+\alpha$ with $\alpha\in A$, of square
$-1$. By Witt's theorem (Lemma~\ref{pbp:tb-witt}(b)) it extends to some $s\in G_\tau$, which lies in $R_\tau$ by
Lemma~\ref{pbp:R}.
\end{proof}

\begin{proposition}[The point]\label{pbp:point}
Let $\iota$ be the stationary position of Proposition~\ref{pbp:stationary} and $y=\iota p^{-1}\in
G_\tau$. Then $\pbel_\tau(y)\le8(\varrho+1)M_0\le b^{5n}$.
\end{proposition}
\begin{proof}
\emph{Step 1.} Put $V_a=p\mu^av_1\in Y_\tau$ for $0\le a\le2\varrho$. Since $p$ is an isometry for both
structures, $\|V_a\|\le M_0$ and $h(V_a,V_b)=-1$ if $a+b=2\varrho$, $0$ if $a+b<2\varrho$
(Proposition~\ref{pbp:stationary}); and $W_a:=y^{-1}\tilde f_a=p\iota^{-1}\tilde f_a=V_a$ for
$a\le\varrho$.

\emph{Step 2.} Let $o_0,\dots,o_\varrho$ be the Gram--Schmidt orthogonalisation of $W_0,\dots,W_\varrho$
(Lemma~\ref{pbe:len}(c)), so that $o_a=W_a-\sum_{a'<a}t_{a'a}W_{a'}$ for some numbers $t_{a'a}$. Then
$\|o_a\|\le\|W_a\|\le M_0$, and $h(V_{2\varrho-a},o_a)=h(V_{2\varrho-a},V_a)-\sum_{a'<a}\aubar{t_{a'a}}
h(V_{2\varrho-a},V_{a'})=-1$; by \eqref{pbe:eq-norm}, $1\le\|V_{2\varrho-a}\|\,\|o_a\|$, so
$\|o_a\|\ge M_0^{-1}$.

\emph{Step 3.} By Lemma~\ref{pbp:frame} there is $s_1\in R_\tau$ with $s_1^{-1}\tilde f_a=\tilde f_a-
\sum_{a'<a}t_{a'a}\tilde f_{a'}$; then $y_1=s_1y$ has $y_1^{-1}\tilde f_a=o_a$.

\emph{Step 4.} For $1\le a\le\varrho$ put $b_a=\mathsf J_1o_{a-1}/\|o_{a-1}\|^2$, of norm $\|o_{a-1}\|^{-1}
\le M_0$. Then $h(o_c,b_a)=\pbip{o_{a-1}}{o_c}/\|o_{a-1}\|^2=\delta_{c,a-1}$, and $h(b_a,b_{a'})$ is a
multiple of $\aubar{h(o_{a'-1},o_{a-1})}=\aubar{h(\tilde e_{a'},\tilde e_a)}=0$. Hence $\beta_a=y_1b_a$
satisfies $h(\tilde e_c,\beta_a)=\delta_{ca}$, $h(\tilde z,\beta_a)=0$, $h(\beta_a,\beta_{a'})=0$. Define
$q\in\operatorname{End}(Y_\tau)$ by $q\tilde e_a=\tilde e_a$, $q\tilde z=\tilde z$, $q\tilde e_a^*=\beta_a$,
$qx_i=x_i-\sum_ah(x_i,\beta_a)\tilde e_a$. The images of the basis have the Gram matrix of the basis:
the only products to check are $h(qx_i,qx_j)=h(x_i,x_j)$ (the added vectors lie in $A$, which is
isotropic and orthogonal to the $x_j$) and $h(qx_i,\beta_b)=h(x_i,\beta_b)-\sum_ah(x_i,\beta_a)
\delta_{ab}=0$. So $q\in G_\tau$; it fixes the frame, so $q\in R_\tau$ (Lemma~\ref{pbp:R}).

\emph{Step 5.} Let $y_2=q^{-1}y_1\in R_\tau y$ and $\mathcal A=y_2^{-1}=y_1^{-1}q$. Then $\mathcal A\tilde
e_a=o_{a-1}$, $\mathcal A\tilde z=o_\varrho$, $\mathcal A\tilde e_a^*=b_a$, all of norm at most $M_0$. The
vectors $(1-\mathsf J_1)w$, $w$ in the normalised basis, are $\tilde e_a-\tilde e_a^*$, $\tilde e_a^*-
\tilde e_a$, $2\tilde z$ and $0$; so $\|\mathcal A(1-\mathsf J_1)\|\le\sum_w\|\mathcal A(1-\mathsf J_1)w\|
\le(4\varrho+2)M_0$. From $\mathcal A\mathsf J_1\mathcal A^*=\mathsf J_1$,
\[
\|\mathcal A\|^2=\|\mathcal A\mathcal A^*\|=\|\mathsf J_1+\mathcal A(1-\mathsf J_1)\mathcal A^*\|\le
1+(4\varrho+2)M_0\|\mathcal A\|,
\]
and as $\|\mathcal A\|\ge1$, $\|\mathcal A\|\le1+(4\varrho+2)M_0$. By Lemma~\ref{pbe:len}(a) and
\eqref{pbe:eq-norm}, $\pbel_\tau(y)=\pbel_\tau(y_2)\le\|y_2\|+1=\|\mathcal A\|+1\le8(\varrho+1)M_0$.
Finally $8nM_0=8n\sqrt r(3B+a_\tau)^{2n}\le8nr\,b^{4n}\le b^{5n}$ by Lemma~\ref{pbp:sizes}.
\end{proof}

From now on $\iota$ and $y$ are as in Propositions~\ref{pbp:stationary} and \ref{pbp:point}, $g_s=
\iota^{-1}s\iota$, and $\theta$ is \eqref{pbp:eq-theta}.

\subsection{One minor to second order}\label{pbp:ss-minor}

In this subsection all norms of matrices are Frobenius norms, $\|\cdot\|=\pbpF{\cdot}$, and the
dimension is arbitrary; no constant depends on it. Let $Q$ be an orthogonal projection. For matrices
$A,B$ put
\[
\pbpct_Q(A,B)=\pbptr\bigl(QA(1-Q)B\bigr)=\pbptr(QAB)-\pbptr(QAQB),\qquad c_Q(A)=QA(1-Q)+(1-Q)AQ,
\]
and $\pbpCT_Q(A_1,A_2,A_3)=\frac12\sum_i\pbpct_Q(A_i,A_i)+\sum_{i<j}\pbpct_Q(A_i,A_j)$. Let $U=2Q-1$;
then $U^2=1$, $UQ=QU$, $Uc_Q(A)U=-c_Q(A)$, $A-c_Q(A)=QAQ+(1-Q)A(1-Q)=\frac12(A+UAU)$ commutes
with $Q$, $\|A-c_Q(A)\|\le\|A\|$, and $c_Q(A)=U[Q,A]$, so $\|c_Q(A)\|=\|[Q,A]\|$.

\begin{lemma}[Third order]\label{pbp:third}
Let $\|A_i\|\le\rho\le\frac1{10}$ \textup($i=1,2,3$\textup). Then
\[
\det\bigl(1+Q(e^{A_1}e^{A_2}e^{A_3}-1)\bigr)=\exp\Bigl(\pbptr\bigl(Q(A_1+A_2+A_3)\bigr)+
\pbpCT_Q(A_1,A_2,A_3)+r\Bigr),\qquad|r|\le50\rho^3 .
\]
\end{lemma}
\begin{proof}
Expanding the three exponential series, $e^{A_1}e^{A_2}e^{A_3}-1=L+S+R$ with $L=\sum A_i$,
$S=\sum_i\frac12A_i^2+\sum_{i<j}A_iA_j$ and $R$ the sum of the monomials $A_1^{n_1}A_2^{n_2}A_3^{n_3}/
(n_1!n_2!n_3!)$ of degree $n_1+n_2+n_3\ge3$. Bounding each monomial by $\rho^{n_1+n_2+n_3}$ (the coefficients sum to those of
$e^{3\rho}$): $\|R\|\le e^{3\rho}-1-3\rho-\frac92\rho^2\le\frac{27}6\rho^3e^{0.3}\le6.1\rho^3$,
$\|Y-L\|\le e^{3\rho}-1-3\rho\le6.1\rho^2$ and $\|Y\|\le e^{3\rho}-1\le3.5\rho$, where $Y=L+S+R$; also
$\|L\|\le3\rho$. Each monomial of $R$ begins with some $A_i$, and $|\pbptr(QA_iW)|\le\|QA_i\|\,\|W\|\le
\rho\|W\|$; so $|\pbptr(QR)|\le6.1\rho^3$ as well. Let $M=QY$, $\|M\|\le0.35$. Then $\det(1+M)=
\exp\pbptr\log(1+M)$ with $\log(1+M)=\sum_{j\ge1}(-1)^{j+1}M^j/j$, and since $|\pbptr M^j|\le\|M\|^j$
for $j\ge2$,
\[
\pbptr\log(1+M)=\pbptr M-\tfrac12\pbptr M^2+r_1,\qquad|r_1|\le\frac{\|M\|^3}{3(1-\|M\|)}\le22\rho^3.
\]
Next $\pbptr M=\pbptr(QL)+\pbptr(QS)+\pbptr(QR)$ and $\pbptr M^2=\pbptr(QLQL)+r_2$ with $r_2=
\pbptr(Q(Y-L)QY)+\pbptr(QLQ(Y-L))$, $|r_2|\le6.1\rho^2(3.5\rho+3\rho)\le40\rho^3$. Finally, as
$\pbptr(QA_iQA_j)=\pbptr(QA_jQA_i)$,
\[
\pbptr(QS)-\tfrac12\pbptr(QLQL)=\sum_i\tfrac12\bigl(\pbptr(QA_i^2)-\pbptr(QA_iQA_i)\bigr)+\sum_{i<j}
\bigl(\pbptr(QA_iA_j)-\pbptr(QA_iQA_j)\bigr)=\pbpCT_Q(A_1,A_2,A_3).
\]
So $r=r_1+\pbptr(QR)-r_2/2$ and $|r|\le(22+6.1+20)\rho^3\le50\rho^3$.
\end{proof}

\begin{proposition}[Second order with crossing parts]\label{pbp:second}
Let $\|A_i\|\le\delta\le\frac1{20}$ and $\|[Q,A_i]\|\le\zeta$ \textup($i=1,2,3$\textup). Then the identity
of Lemma~\ref{pbp:third} holds with $|r|\le400\,\delta\zeta^2$.
\end{proposition}
\begin{proof}
If $\zeta\ge\delta$, Lemma~\ref{pbp:third} with $\rho=\delta$ gives $|r|\le50\delta^3\le50\delta\zeta^2$.
If $\zeta=0$, the $A_i$ commute with $Q$. Then $1+Q(e^{A}-1)=e^{QA}$ (as $(QA)^j=QA^j$), and
$1+Q(XY-1)=(1+Q(X-1))(1+Q(Y-1))$ when $X$ commutes with $Q$; so the determinant is
$\prod_ie^{\pbptr QA_i}$, while $\pbpct_Q(A_i,A_j)=\pbptr(A_iQ(1-Q)A_j)=0$; thus $r=0$.
Let $0<\zeta<\delta$ and put $B_i(t)=A_i-c_Q(A_i)+tc_Q(A_i)$ for $t\in\C$, and
\[
g(t)=\det\bigl(1+Q(e^{B_1(t)}e^{B_2(t)}e^{B_3(t)}-1)\bigr)\exp\Bigl(-\pbptr\bigl(Q\textstyle\sum_iB_i(t)
\bigr)-\pbpCT_Q(B_1(t),B_2(t),B_3(t))\Bigr),
\]
an entire function of $t$ with $g(1)=e^r$. For $|t|<R_0:=\delta/\zeta$, $\|B_i(t)\|\le\|A_i\|+|t|\,
\|[Q,A_i]\|\le2\delta\le\frac1{10}$, so by Lemma~\ref{pbp:third} $g(t)=e^{r(t)}$ with $|r(t)|\le
50(2\delta)^3=400\delta^3\le0.05$. Hence $f=\operatorname{Log}g$ (principal branch) is holomorphic on the
disc $|t|<R_0$, $f(t)=r(t)$ there, $|f|\le400\delta^3$, and $f(1)=r$. The $B_i(0)$ commute with $Q$, so
$f(0)=0$ by the case $\zeta=0$. Conjugation by $U$ sends $B_i(t)$ to $B_i(-t)$ and leaves the three
ingredients of $g$ unchanged: $\det(1+Q(UXU-1))=\det(U(1+Q(X-1))U)=\det(1+Q(X-1))$,
$\pbptr(QUBU)=\pbptr(QB)$, and $\pbpct_Q(UAU,UBU)=\pbptr(UQA(1-Q)BU)=\pbpct_Q(A,B)$. So $f$ is even,
$f(t)=\phi(t^2)$ with $\phi$ holomorphic on $|w|<R_0^2$, $\phi(0)=0$, $|\phi|\le400\delta^3$, and Schwarz's
lemma gives $|f(1)|=|\phi(1)|\le400\delta^3/R_0^2=400\delta\zeta^2$.
\end{proof}

\begin{lemma}[Size of one minor]\label{pbp:minorsize}
Let $k$ be unitary, $S$ a set of coordinates and $P$ the coordinate projection onto $S$. Then
$|\Delta_S(k)|^2\le\exp(-\|[P,k]\|^2/2)$.
\end{lemma}
\begin{proof}
Let $\mathcal M=(1-P)kP$ and $\mathcal Y=1-P+PkP$, so $\Delta_S(k)=\det\mathcal Y$. Using $k^*k=1$,
$\mathcal Y^*\mathcal Y=1-P+Pk^*PkP=1-\mathcal M^*\mathcal M$. The eigenvalues $\sigma_i\in[0,1]$ of
$\mathcal M^*\mathcal M$ give $|\det\mathcal Y|^2=\prod(1-\sigma_i)\le e^{-\sum\sigma_i}=e^{-\|\mathcal
M\|^2}$. Finally $\|[P,k]\|^2=\|Pk(1-P)\|^2+\|(1-P)kP\|^2$, and both terms equal $\pbptr P-
\pbptr(Pk^*PkP)$ ($k$ unitary), so $\|[P,k]\|^2=2\|\mathcal M\|^2$.
\end{proof}

\begin{proposition}[One minor near the identity]\label{pbp:minorconj}
Let $o$ be an injective real function on the coordinates, $S$ a set of coordinates with $o_i<o_j$
for $i\in S$, $j\notin S$, and $P$ the coordinate projection onto $S$. Let $Z\in\pdsg$ with
$\|Z\|\le\delta\le\frac1{20}$, and let $u_-$ \textup(resp.\ $u_+$\textup) keep the entries $Z_{ij}$ with
$\varepsilon_i=\varepsilon_j$ and $o_i<o_j$ \textup(resp.\ $o_i>o_j$\textup) and vanish elsewhere. Put
$\hat g=e^{-u_-}e^Ze^{-u_+}$. Then for every $k\in\pdsK$,
\[
\Delta_S(k^*\hat gk)=\exp\bigl(\pbptr(PZ)+\tfrac14\|[P,Z]\|^2+r\bigr),\qquad|r|\le10^3\delta\bigl(
\|[P,Z]\|^2+\|[P,k]\|^2\bigr),
\]
and $\det\hat g=e^{\pbptr Z}$.
\end{proposition}
\begin{proof}
$\det e^{-u_\pm}=e^{-\pbptr u_\pm}=1$. Let $Q=kPk^*$. As $\Delta_S(x)=\det(1+P(x-1))$ and $\det(1+
\mathcal A\mathcal B)=\det(1+\mathcal B\mathcal A)$, $\Delta_S(k^*\hat gk)=\det(1+Pk^*(\hat g-1)k)=\det(1+
Q(\hat g-1))$. Put $A_1=-u_-$, $A_2=Z$, $A_3=-u_+$, $x=\|[P,Z]\|$ and $\kappa=\|[P,k]\|$; then
$\|A_i\|\le\delta$, $Q-P=[k,P]k^*$ so $\|Q-P\|=\kappa$, and since $[P,\cdot]$ multiplies the entry
$(a,b)$ by $P_{aa}-P_{bb}$, $\|[P,A_i]\|\le x$. As $[Q,A]=[P,A]+[Q-P,A]$, $\|[Q,A_i]\|\le x+2\delta\kappa
=:\zeta\le x+\kappa$. Proposition~\ref{pbp:second} gives the formula with $\pbptr(Q(A_1+A_2+A_3))+
\pbpCT_Q(A_1,A_2,A_3)$ in place of $\pbptr(PZ)+x^2/4$ and an error at most $400\delta\zeta^2$. We compare.

(i) $W=A_1+Z+A_3$ has $W_{ab}=0$ for $a\ne b$ with $\varepsilon_a=\varepsilon_b$ (as $o$ is injective) and
$W_{aa}=Z_{aa}$, while $Q$ is block diagonal ($k\in\pdsK$ commutes with $J$). So $\pbptr(QW)-
\pbptr(PZ)=\sum_a(Q-P)_{aa}Z_{aa}$. Now $Q-P=-[P,k]k^*$, and $|[P,k]_{ab}k_{ab}|=|[P,k]_{ab}|^2$; so
$\sum_a|(Q-P)_{aa}|\le\kappa^2$ and $|\pbptr(QW)-\pbptr(PZ)|\le\delta\kappa^2$.

(ii) $\pbpCT_P(A_1,Z,A_3)=x^2/4$. Indeed $\pbpct_P(A,B)=\sum_{i\in S,j\notin S}A_{ij}B_{ji}$, and for
$i\in S$, $j\notin S$ we have $o_i<o_j$, so $(A_1)_{ij}=-Z_{ij}$ and $(A_3)_{ji}=-Z_{ji}$ if
$\varepsilon_i=\varepsilon_j$, while $(A_1)_{ji}=(A_3)_{ij}=0$, and $(A_1)_{ij}=(A_3)_{ji}=0$ if
$\varepsilon_i\ne\varepsilon_j$. As $Z\in\pdsg$, $Z_{ij}Z_{ji}=-|Z_{ji}|^2$ if $\varepsilon_i=
\varepsilon_j$ and $=|Z_{ji}|^2$ otherwise. With $C$ and $C'$ the sums of $|Z_{ji}|^2$ over $i\in S$,
$j\notin S$ with equal, resp.\ opposite, signs: $\pbpct_P(A_1,A_1)=\pbpct_P(A_3,A_3)=0$, $\pbpct_P(Z,Z)
=C'-C$, $\pbpct_P(A_1,Z)=C$, $\pbpct_P(A_1,A_3)=-C$, $\pbpct_P(Z,A_3)=C$; so $\pbpCT_P=\frac12(C+C')$.
And $x^2=\sum_{i\in S,j\notin S}(|Z_{ij}|^2+|Z_{ji}|^2)=2(C+C')$, since $|Z_{ij}|=|Z_{ji}|$.

(iii) For $R\in\{P,Q\}$, $\pbpct_R(A,B)=\pbptr(\mathcal X_R\mathcal Y_R)$ with $\mathcal X_R=RA(1-R)=
R[R,A](1-R)$ and $\mathcal Y_R=(1-R)BR$, so $\|\mathcal X_R\|\le\|[R,A]\|$, $\|\mathcal Y_R\|\le
\|[R,B]\|$; and $QA(1-Q)-PA(1-P)=(Q-P)A(1-Q)-PA(Q-P)$, so $\|\mathcal X_Q-\mathcal X_P\|\le2\delta\kappa$
for $\|A\|\le\delta$, and likewise for $\mathcal Y$. Hence $|\pbpct_Q(A_i,A_j)-\pbpct_P(A_i,A_j)|\le
2\delta\kappa\zeta+x\cdot2\delta\kappa\le4\delta\kappa\zeta$, and $|\pbpCT_Q-\pbpCT_P|\le\frac92\cdot
4\delta\kappa\zeta=18\delta\kappa\zeta$.

Altogether $|r|\le\delta\kappa^2+18\delta\kappa(x+\kappa)+400\delta(x+\kappa)^2\le10^3\delta(x^2+
\kappa^2)$.
\end{proof}

\subsection{The Gaussian shape of the coefficients}\label{pbp:ss-gauss}

Norms are again Frobenius norms unless stated otherwise. For a matrix $Y$ and an integer vector $w$
put $\pbpQ_w(Y)=\sum_{i,j}|w_i-w_j|\,|Y_{ij}|^2$. For $g\in G$ with $\|g-1\|\le\frac12$ let $\log g=
\sum_{j\ge1}(-1)^{j+1}(g-1)^j/j$; as in the proof of Lemma~\ref{pds:lip} (Step~2), $\log g\in\pdsg$,
$e^{\log g}=g$, and $\|\log g\|\le-\log(1-\|g-1\|)\le2\|g-1\|$. The same holds in $\auU(k)$, where
$\log h$ is skew-hermitian.

\begin{lemma}[Removing the unipotent factors]\label{pbp:reduced}
Let $Z\in\pdsg$ with $\|Z\|\le\frac1{20}$, and let $u_\pm$ and $\hat g$ be as in
Proposition~\ref{pbp:minorconj} for $o=\lambda$. Then
\[
c(e^Z)=d_\Lambda\int_{\pdsK}\frac{|\pdsm_\Lambda(k)|^2}{\pdsm_{-(\lambda+\rho)}(k^{-1}\hat gk)}\,dk .
\]
\end{lemma}
\begin{proof}
Let $x=e^Z$; $\|x-1\|\le e^{1/20}-1<1$ in operator norm, so $x$ is accretive and $1\in A_x$ in the
notation of \S\ref{pco:ss-indep}. By Theorem~\ref{pco:FJ} and Proposition~\ref{pco:indep}, $c(x)=
F(x)=d_\Lambda I(x;1)=d_\Lambda\mathcal J(1,1)$, with $\mathcal J$ and $U_x$ as in the proof of
Proposition~\ref{pco:indep}. By the second assertion of Proposition~\ref{pco:indep}, $\mathcal J$
is constant along every path in $U_x$. For $t\in[0,1]$,
$e^{tu_+}x^{-1}e^{tu_-}$ is within $e^{3/20}-1<1$ of $1$ in operator norm, hence accretive, so $(e^{tu_-},
e^{tu_+})\in U_x$, and
\[
\mathcal J(1,1)=\mathcal J(e^{u_-},e^{u_+})=\int_{\pdsK}\frac{\pdsm_\Lambda(e^{u_-}k)\,\pdsm_\Lambda(k^*
e^{u_+})}{\pdsm_{\lambda+\rho}(k^*e^{u_+}x^{-1}e^{u_-}k)}\,dk .
\]
The matrix $b=e^{u_-}$ has $b_{cd}=0$ unless $c=d$ or $\varepsilon_c=\varepsilon_d$ and $\lambda_c<
\lambda_d$ (the powers of $u_-$ have this shape), and $b_{cc}=1$; for such $(c,d)$, $\Lambda_c\le
\Lambda_d$ by Lemma~\ref{pco:minors}(c). So Lemma~\ref{pco:minors}(b) (left version; $b$ is triangular
for the order of decreasing $\lambda$) gives $\pdsm_\Lambda(bk)=\pdsm_\Lambda(k)$. Likewise $e^{u_+}$
has its off-diagonal entries at $(c,d)$ with $\lambda_c>\lambda_d$, $\Lambda_c\ge\Lambda_d$, and
$\pdsm_\Lambda(k^*e^{u_+})=\pdsm_\Lambda(k^*)=\aubar{\pdsm_\Lambda(k)}$. Finally
$(k^*e^{u_+}x^{-1}e^{u_-}k)^{-1}=k^{-1}\hat gk$, and $\pdsm_{\lambda+\rho}(y^{-1})=\pdsm_{-(\lambda+\rho)}
(y)$ (Lemma~\ref{pco:minors}(a)).
\end{proof}

\begin{lemma}\label{pbp:dLambda}
$d_\Lambda\le e^{2T}$.
\end{lemma}
\begin{proof}
Let $M_\Lambda=\max(1,\max_i|\Lambda_i|)\le3b$ and $\varepsilon_0=1/(6M_\Lambda r)$. By
\eqref{pds:eq-m}, $\pdsm_\Lambda(k)=(\det k)^{-M_\Lambda}\prod_{t=1-M_\Lambda}^{M_\Lambda}\Delta_{S_t
(\Lambda)}(k)$. If $k\in\pdsK$ and $\|k-1\|_{\mathrm{op}}\le\varepsilon_0$, then $|\det k|=1$ and each
$\Delta_S(k)=\det(1+(k-1)_{SS})$ is a product of at most $r$ factors $1+\xi$ with $|\xi|\le\varepsilon_0
\le\frac12$, so $|\Delta_S(k)|\ge e^{-2r\varepsilon_0}$ and $|\pdsm_\Lambda(k)|\ge e^{-4M_\Lambda r
\varepsilon_0}\ge\frac12$. Choose a maximal family $k_1,\dots,k_\nu\in\pdsK$ with $\|k_ik_j^{-1}-1\|_{
\mathrm{op}}>\varepsilon_0$ for $i\ne j$. Then $\|k_i-k_j\|_{\mathrm{op}}>\varepsilon_0$, the operator norm
balls of radius $\varepsilon_0/2$ about the $k_i$ are disjoint and lie in the ball of radius $2$ about
$0$ in $M_r(\C)$, of real dimension $2r^2$, so $\nu\le(4/\varepsilon_0)^{2r^2}$; and by maximality
$\pdsK$ is covered by the right translates $\{k:\|k-1\|_{\mathrm{op}}\le\varepsilon_0\}k_i$. As $dk$ has
mass $1$, $d_\Lambda^{-1}=\int|\pdsm_\Lambda|^2\ge\frac14(\varepsilon_0/4)^{2r^2}$, so $d_\Lambda\le
4(72b^2)^{2r^2}\le b^{6r^2}\le e^{2T}$ ($T\ge3r^2\log b$).
\end{proof}

\begin{proposition}[Gaussian shape, noncompact factor]\label{pbp:gauss}
Let $g\in G$ with $\|g-1\|\le\delta_1$, $\delta_1T\le5\cdot10^{-13}$, and $Z=\log g$. There is
$\gamma\in[0,1]$ with $|c(g)e^{-\pbptr(\mu Z)}-\gamma|\le\gamma/10+e^{-100T}$.
\end{proposition}
\begin{proof}
Put $\delta=2\delta_1$, so $\|Z\|\le\delta$, $\delta T\le10^{-12}$ and $\delta\le10^{-12}$. By
Lemma~\ref{pbp:reduced}, $c(g)=d_\Lambda\int|\pdsm_\Lambda(k)|^2/\allowbreak \pdsm_{-(\lambda+\rho)}(k^{-1}\hat gk)\,dk$.
Let $w=-(\lambda+\rho)$ and $M'=\max_i|w_i|$; then $\pdsm_w(y)=(\det y)^{-M'}\prod_{t=1-M'}^{M'}
\Delta_{S_t(w)}(y)$. For $i\in S_t(w)$, $j\notin S_t(w)$ we have $(\lambda+\rho)_i<(\lambda+\rho)_j$,
hence $\lambda_i<\lambda_j$ (Lemma~\ref{pco:minors}(c)), so Proposition~\ref{pbp:minorconj} applies to
each factor with $o=\lambda$. Since exactly $M'+w_i$ of the $t\in(-M',M']$ satisfy $t\le w_i$, and exactly
$|w_i-w_j|$ separate $i$ from $j$, multiplying gives
\[
\pdsm_w(k^{-1}\hat gk)=\exp\bigl(-\pbptr((\lambda+\rho)Z)+\tfrac14\pbpQ_w(Z)+r(k)\bigr),\qquad|r(k)|\le
10^3\delta\bigl(\pbpQ_w(Z)+\pbpQ_w(k)\bigr).
\]
Now $\lambda+\rho=\Lambda+2\rho_{\mathrm c}$ with $|2\rho_{\mathrm c,i}|\le r$, so $|\pbptr(2\rho_{\mathrm c}Z)|
\le r^2\delta$ and $|w_i-w_j|\le|\Lambda_i-\Lambda_j|+2r$; with $\sum_{ij}|k_{ij}|^2=r$ this gives
$\pbpQ_w(k)\le\pbpQ_\Lambda(k)+2r^2$. Put $q=\pbpQ_w(Z)/4$. Then $\pdsm_w(k^{-1}\hat gk)=\exp(q-
\pbptr(\mu Z)-r'(k))$ with $|r'(k)|\le10^4\delta\,s(k)$, $s(k)=q+\pbpQ_\Lambda(k)+r^2$, and since
$d_\Lambda\int|\pdsm_\Lambda|^2=1$,
\[
c(g)e^{-\pbptr(\mu Z)}-e^{-q}=d_\Lambda\int_{\pdsK}|\pdsm_\Lambda(k)|^2e^{-q}\bigl(e^{r'(k)}-1\bigr)\,dk .
\]
By Lemma~\ref{pbp:minorsize}, multiplied over $t$, $|\pdsm_\Lambda(k)|^2\le e^{-\pbpQ_\Lambda(k)/2}$. If
$s(k)\le5\cdot10^4T$, then $|r'(k)|\le5\cdot10^{-4}$ and $|e^{r'}-1|\le\frac1{10}$. Otherwise $|r'(k)|\le
s(k)/4$ and $|\pdsm_\Lambda(k)|^2e^{-q}|e^{r'}-1|\le e^{-\pbpQ_\Lambda(k)/2-q+s(k)/4}\le e^{-s(k)/4+r^2/2}
\le e^{-5100T}$, as $r^2\le T$. Integrating, $|c(g)e^{-\pbptr(\mu Z)}-e^{-q}|\le e^{-q}/10+d_\Lambda
e^{-5100T}$, and $d_\Lambda\le e^{2T}$ (Lemma~\ref{pbp:dLambda}). Take $\gamma=e^{-q}$.
\end{proof}

\begin{lemma}[Gaussian shape, compact factor]\label{pbp:compact}
Let $h\in\auU(k)$ with $\|h-1\|\le\delta_1$, $\delta_1T\le5\cdot10^{-13}$, and $Z'=\log h$. There is
$\gamma'\in[0,1]$ with $|c'(h)e^{-\pbptr(\Lambda'Z')}-\gamma'|\le\gamma'/10+e^{-100T}$.
\end{lemma}
\begin{proof}
Let $\delta=2\delta_1\ge\|Z'\|$, $P$ the projection onto the coordinate $k$, and $x=\|[P,Z']\|$. By
Proposition~\ref{pbp:second} with $Q=P$, $A_1=A_3=0$, $A_2=Z'$, $\zeta=x$: $h_{kk}=\det(1+P(e^{Z'}-1))=
\exp(Z'_{kk}+\frac12\pbpct_P(Z',Z')+r)$ with $|r|\le400\delta x^2$, and $\pbpct_P(Z',Z')=\sum_{i\ne k}
Z'_{ki}Z'_{ik}=-\sum_{i\ne k}|Z'_{ik}|^2=-x^2/2$. By \eqref{pbp:eq-cu}, $c'(h)e^{-\pbptr(\Lambda'Z')}=
h_{kk}^{a_\tau}e^{-a_\tau Z'_{kk}}=\gamma'e^{a_\tau r}$ with $\gamma'=e^{-a_\tau x^2/4}$. If $a_\tau x^2\le
5\cdot10^4T$, $|a_\tau r|\le2\cdot10^{-5}$ and $\gamma'|e^{a_\tau r}-1|\le\gamma'/10$. Otherwise
$\gamma'|e^{a_\tau r}-1|\le\gamma'e^{|a_\tau r|}\le e^{-a_\tau x^2/8}\le e^{-100T}$.
\end{proof}

\subsection{Localisation}\label{pbp:ss-loc}

Norms are Frobenius norms. For $s\in R_\tau$ write $x_s=g_s-1$, $x'_s=h_s-1$, and $y_s=g_s-a(g_s)$ for
the \emph{noncompact part} of $g_s$ (the entries at $(i,j)$ with $\varepsilon_i\ne\varepsilon_j$). Let
$P_{\mathrm{off}}$ be the coordinate projection onto the coordinates outside $X$.

\begin{lemma}[The phase at the logarithms]\label{pbp:phaselog}
If $\|x_s\|,\|x'_s\|\le\frac12$, then $\pbptr(\mu\log g_s)-\pbptr(\Lambda'\log h_s)=\varphi(s)$.
\end{lemma}
\begin{proof}
For pairs of matrices put $\Phi(X,X')=\pbptr(\mu X)-\pbptr(\Lambda'X')$ (linear), multiply pairs
componentwise, and let $\xi=(x_s,x'_s)$, $\xi_t=(x_t,x'_t)$ for $t\in R_\tau$. We show $\Phi(\xi^i\xi
(1+\xi_t))=0^i\varphi(s)$ for all $i\ge0$ and $t$. As $(1+\xi)(1+\xi_t)=1+\xi_{st}$, $\xi(1+\xi_t)=\xi_{st}
-\xi_t$, so for $i=0$ the left side is $\varphi(st)-\varphi(t)=\varphi(s)$ (Lemma~\ref{pbp:traces}(b)).
And $\xi^{i+1}\xi(1+\xi_t)=\xi^i\xi(\xi_{st}-\xi_t)=\xi^i\xi(1+\xi_{st})-\xi^i\xi(1+\xi_t)$, which gives the
induction step. With $t=1$, $\Phi(\xi^{i+1})=0^i\varphi(s)$; insert this in the series for $\log$.
\end{proof}

\begin{lemma}[Rows and columns]\label{pbp:rows}
For $s\in R_\tau$: $\|x_s\|^2=2\|y_s\|^2-2\Rel\pbptr x_s$, $\|x'_s\|^2=-2\Rel\pbptr x'_s$,
$\|P_{\mathrm{off}}x_s\|^2\le6\|y_s\|^2$ and $\|x_sP_{\mathrm{off}}\|^2\le6\|y_s\|^2$.
\end{lemma}
\begin{proof}
For $g\in G$, $gJg^*=J$; its diagonal entry $i$ reads $\sum_j\varepsilon_i\varepsilon_j|g_{ij}|^2=1$, i.e.\
$\sum_j|g_{ij}|^2-2\sum_j|y_{ij}|^2=1$ with $y=g-a(g)$; with $x=g-1$, $\sum_j|x_{ij}|^2=\sum_j|g_{ij}|^2-
2\Rel x_{ii}-1$, so
\[
2\Rel x_{ii}=2\textstyle\sum_j|y_{ij}|^2-\sum_j|x_{ij}|^2 ,
\]
and for $h\in\auU(k)$ likewise $2\Rel x'_{ii}=-\sum_j|x'_{ij}|^2$. Summing over $i$ gives the two
identities. Multiply the first row identity by $-\Lambda_i\ge0$ and sum: $\sum_i(-\Lambda_i)\sum_j
|x_{ij}|^2=2\sum_i(-\Lambda_i)\sum_j|y_{ij}|^2+2\Rel\pbptr(\mu x)$. By \eqref{pbp:eq-steep} the left side is
at least $B\|P_{\mathrm{off}}x\|^2$ and the first term on the right at most $6B\|y\|^2$. Multiplying the
second row identity by $\Lambda'_i\ge0$ gives $\Rel\pbptr(\Lambda'x'_s)\le0$, and $\Rel\pbptr(\mu x_s)=
\Rel\pbptr(\Lambda'x'_s)$ because $\varphi(s)$ is imaginary. Hence $B\|P_{\mathrm{off}}x_s\|^2\le6B
\|y_s\|^2$. For the columns apply this to $s^{-1}$: $x_{s^{-1}}=Jx_s^*J$, $y_{s^{-1}}=Jy_s^*J$, and
$P_{\mathrm{off}}$ commutes with $J$.
\end{proof}

\begin{lemma}\label{pbp:leftX}
If $k_0\in\pdsK$ is the identity on the coordinates outside $X$, then $c(k_0g)=c(g)$ for $g\in G$.
\end{lemma}
\begin{proof}
By \eqref{pco:eq-left} (applied with $k_1=k_0$ and with $k_1=1$, for the same $a\in A_g$) it suffices
that $\pdsm_\Lambda(k_0Y)=\pdsm_\Lambda(Y)$ for invertible $Y$. Let $S=S_t(\Lambda)$. If $t\ge1$, $S\cap X=
\emptyset$ ($\Lambda=0$ on $X$) and the rows of $k_0Y$ indexed by $S$ are those of $Y$, so $\Delta_S(k_0Y)
=\Delta_S(Y)$. If $t\le0$, $S\supseteq X$ and $(k_0)_{ij}=0$ for $i\in S$, $j\notin S$, so $(k_0Y)_{SS}=
(k_0)_{SS}Y_{SS}$ and $\Delta_S(k_0Y)=\det k_0\cdot\Delta_S(Y)$; this is compensated by the factor
$(\det k_0Y)^{-1}$ in \eqref{pds:eq-m}.
\end{proof}

\begin{lemma}[Far from $\pdsK$]\label{pbp:far}
For $g\in G$ with noncompact part $y$, $|c(g)|^4\le(1+\|y\|^2)^{-h_0}$.
\end{lemma}
\begin{proof}
From $g^*Jg=J$, with $J$ commuting with $a=a(g)$ and anticommuting with $y$, the block diagonal part
gives $a^*a-y^*y=1$. So $|\det a|^2=\det(1+y^*y)\ge1+\pbptr(y^*y)=1+\|y\|^2$, and
Theorem~\ref{pds:decay} with $c=h_0\le h$ (Lemma~\ref{pbp:sizes}) gives $|c(g)|^2\le|\det a|^{-h_0}$.
\end{proof}

\begin{lemma}[Localisation]\label{pbp:localise}
Let $0<\rho_1\le\frac16$ and $s\in R_\tau$ with $\|y_s\|\le\rho_1$. There is $t\in R_\tau$ with
$\theta(t)=\theta(s)$, $\|x_t\|\le4\rho_1$ and $\|x'_t\|\le4\rho_1$.
\end{lemma}
\begin{proof}
Let $\alpha=(g_s)_{XX}$ be the block of $g_s$ on $X$. The block on $X\times X$ of $g_s^*Jg_s=J$ is
$\alpha^*\alpha+\mathcal W^*J'\mathcal W=1$, where $\mathcal W$ consists of the rows outside $X$ and the
columns in $X$ of $g_s$, which are those of $x_s$, and $J'$ is the diagonal of signs outside $X$. So
$\|\alpha^*\alpha-1\|\le\|P_{\mathrm{off}}x_s\|^2\le6\rho_1^2\le\rho_1<1$ (Lemma~\ref{pbp:rows}); thus
$\alpha$ is invertible, $\alpha=k_1z_1$ with $k_1$ unitary and $z_1=(\alpha^*\alpha)^{1/2}$, and $\|z_1-1\|
\le\|z_1^2-1\|\le\rho_1$ (on eigenvalues $\nu\ge0$, $|\nu-1|\le|\nu^2-1|$). Let $u=k_1^{-1}$ and $t=s_us$
(Lemma~\ref{pbp:traces}(c)). Then $g_t=g_{s_u}g_s$ has the block $z_1$ on $X$; $c(g_t)=c(g_s)$ by
Lemma~\ref{pbp:leftX}; $(h_t)_{kk}=(h_s)_{kk}$ since $h_{s_u}=u\oplus1$, so $c'(h_t)=c'(h_s)$; and
$\varphi(t)=\varphi(s_u)+\varphi(s)=\varphi(s)$. So $\theta(t)=\theta(s)$. As $g_{s_u}$ is block diagonal and
unitary, $y_t=g_{s_u}y_s$ and $\|y_t\|=\|y_s\|$. Counting the entries of $x_t$ in the rows outside
$X$, in the columns outside $X$ and in the block on $X$, Lemma~\ref{pbp:rows} gives $\|x_t\|^2\le
6\rho_1^2+6\rho_1^2+\rho_1^2$, so $\|x_t\|\le4\rho_1$. Finally, by Lemma~\ref{pbp:rows} and
Lemma~\ref{pbp:traces}(a), $\|x'_t\|^2=-2\Rel\pbptr x'_t=-2\Rel\pbptr x_t=\|x_t\|^2-2\|y_t\|^2\le13\rho_1^2$.
\end{proof}

\begin{proposition}[The cone]\label{pbp:cone}
For every $s\in R_\tau$, $|\theta(s)|\le2\Rel\theta(s)+12e^{-100T}$.
\end{proposition}
\begin{proof}
Let $E=e^{-100T}$ and $\rho_1=10^{-13}/T\le\frac16$. If $\|y_s\|\ge\rho_1$, Lemma~\ref{pbp:far} gives
$|\theta(s)|^4\le|c(g_s)|^4\le(1+\rho_1^2)^{-h_0}$, and $h_0\log(1+\rho_1^2)\ge10^{30}T^3\rho_1^2/2=5000T$,
so $|\theta(s)|\le E\le2\Rel\theta(s)+12E$ (as $\Rel\theta\ge-|\theta|$). Otherwise take $t$ as in
Lemma~\ref{pbp:localise}, and $\delta_1=4\rho_1$, so $\delta_1T\le5\cdot10^{-13}$. Let $Z=\log g_t$,
$Z'=\log h_t$. The number $\pbptr(\Lambda'Z')$ is imaginary, so by Lemma~\ref{pbp:phaselog}
\[
\theta(t)=u_1\,u_2,\qquad u_1=c(g_t)e^{-\pbptr(\mu Z)},\quad u_2=\aubar{c'(h_t)e^{-\pbptr(\Lambda'Z')}} .
\]
By Proposition~\ref{pbp:gauss} and Lemma~\ref{pbp:compact} there are $\gamma,\gamma'\in[0,1]$ with
$|u_1-\gamma|\le\gamma/10+E$ and $|u_2-\gamma'|\le\gamma'/10+E$. Then, as $\gamma,\gamma',E\le1$,
\[
|u_1u_2-\gamma\gamma'|\le\gamma|u_2-\gamma'|+\gamma'|u_1-\gamma|+|u_1-\gamma||u_2-\gamma'|\le
0.21\gamma\gamma'+4E ,
\]
and for $z=\theta(s)=\theta(t)$ and $P=\gamma\gamma'\ge0$: $|z|-2\Rel z\le(P+|z-P|)-2(P-|z-P|)\le
-P+3(0.21P+4E)\le12E$.
\end{proof}

\subsection{Coverings and the proof of \texorpdfstring{\textup{(DS4)}}{(DS4)}}\label{pbp:ss-cone}

In this subsection $\|\cdot\|$ is the operator norm. Note $\|g_s\|=\|y^{-1}sy\|$ and $\|g_s-1\|=
\|y^{-1}sy-1\|$, since $p$ is unitary; so $\pbeB_y=\{s:\|g_s-1\|\le1\}$. The map $s\mapsto y^{-1}sy$ is a
homeomorphism of $R_\tau$ onto a closed subgroup of $G_\tau$, so the sets $\{s:\|g_s\|\le T'\}$ are
compact. Put
\[
\mathrm{Large}=\{s\in R_\tau:\|g_s\|\le\beta+1\},\qquad\mathrm{Small}=\{s\in R_\tau:\|g_s-1\|\le\rho_0,\
\|h_s-1\|\le\rho_0\}.
\]

\begin{lemma}\label{pbp:cover1}
For real $\mathsf c>4r^2$, $\int_{R_\tau}(2/\|g_s\|)^{\mathsf c}ds\le2^{\mathsf c}5^{2r^2}\frac{\mathsf c}{
\mathsf c-4r^2}\pbvol(\pbeB_y)$; for $\mathsf c=9D^2$ this is at most $16^{9D^2}\pbvol(\pbeB_y)$.
\end{lemma}
\begin{proof}
The proof of Lemma~\ref{pbe:cover} (with $x=y$ there) shows $\pbvol\{s:\|g_s\|\le T'\}\le5^{2r^2}
T'^{4r^2}\pbvol(\pbeB_y)$ for $T'\ge1$; integrate against $\mathsf cT'^{-\mathsf c-1}dT'$ as there. For
$\mathsf c=9D^2\ge9r^2$, $\frac{\mathsf c}{\mathsf c-4r^2}\le\frac95$ and $\frac95\cdot5^{2r^2}\le8^{9D^2}$.
\end{proof}

\begin{lemma}\label{pbp:cover2}
$\pbvol(\mathrm{Large})\le\beta^{16D^2}\pbvol(\mathrm{Small})$.
\end{lemma}
\begin{proof}
Let $\mathcal J(s)=g_s\oplus h_s\in\mathrm{GL}_D(\C)$; $\|\mathcal J(s)\|=\max(\|g_s\|,1)=\|g_s\|$,
$\|\mathcal J(s)^{-1}\|=\|\mathcal J(s)\|$, and $s\in\mathrm{Small}$ if and only if $\|\mathcal J(s)-1\|\le
\rho_0$. Let $T'=\beta+1$ and choose a maximal family $s_1,\dots,s_\nu$ in Large with $\|\mathcal J(s_i
s_j^{-1})-1\|>\rho_0$ for $i\ne j$. Since $\mathcal J(s_i)-\mathcal J(s_j)=(\mathcal J(s_is_j^{-1})-1)
\mathcal J(s_j)$, $\|\mathcal J(s_i)-\mathcal J(s_j)\|\ge\rho_0/T'$. The balls of radius $\rho_0/(2T')$ about
the $\mathcal J(s_i)$ in $M_D(\C)$ are disjoint and lie in the ball of radius $2T'$, so $\nu\le
(4T'^2/\rho_0)^{2D^2}=(800\beta^5(\beta+1)^2)^{2D^2}\le\beta^{16D^2}$. By maximality, Large is covered
by the sets $\mathrm{Small}\cdot s_i$, and the Haar measure of $R_\tau$ is right invariant.
\end{proof}

\begin{proof}[Proof of Theorem~\ref{pbp:main}]
By the reduction at the end of \S\ref{pbp:ss-data} we may assume $\epsilon_0=1$. Take $\iota$ and $y$
as in Propositions~\ref{pbp:stationary} and \ref{pbp:point}, $x_\tau=y$ and $v_\tau$ as in
Lemma~\ref{pbp:W}; then $\pbel_\tau(x_\tau)\le b^{5n}$, and the integrand is $\theta$ of
\eqref{pbp:eq-theta}. Put $G(s)=(2/\|g_s\|)^{9D^2}$ and $E=\beta^{-50D^2}$, so $E^2=e^{-100T}$.

\emph{Decay.} By Theorem~\ref{pds:decay} with the exponents $9D^2$ and $18D^2$ (both at most $h$,
Lemma~\ref{pbp:sizes}), $|\theta(s)|\le|c(g_s)|\le G(s)$ and $|\theta(s)|\le G(s)^2$. By
Lemma~\ref{pbp:cover1}, $G$ and hence the continuous function $\theta$ are integrable, and $\int G\le
16^{9D^2}\pbvol(\pbeB_y)$.

\emph{Cone with decaying error.} By Proposition~\ref{pbp:cone}, $\omega:=|\theta|-2\Rel\theta\le12E^2$;
also $\omega\le3|\theta|\le3G^2$. Hence $\omega\le6EG$: otherwise $36E^2G^2<\omega^2\le12E^2\cdot3G^2$.

\emph{Small ball.} For $s\in\mathrm{Small}$, $|c(g_s)|\ge\frac12$ by Lemma~\ref{pds:lip}, because
$\rho_0\le(48b^2)^{-1}\le(4C_0)^{-1}$ with $C_0=3r(L+2r)\le12b^2$ (Lemma~\ref{pbp:sizes}); and
$|c'(h_s)|=|(h_s)_{kk}|^{a_\tau}\ge(1-\rho_0)^{a_\tau}\ge1-a_\tau\rho_0\ge\frac12$. So $|\theta|\ge\frac14$
on Small.

\emph{Integration.} With $I=|\int_{R_\tau}\theta\,ds|$, $V=\pbvol(\mathrm{Large})$ and $W=\pbvol(\mathrm{Small})$,
and using $\pbeB_y\subseteq\mathrm{Large}$ ($\|g\|\le1+\|g-1\|\le2$),
\[
\tfrac14W\le\int|\theta|\le2\Rel\int\theta+6E\int G\le2I+6E\,16^{9D^2}\,V .
\]
By Lemma~\ref{pbp:cover2}, $V\le\beta^{16D^2}W\le8\beta^{16D^2}I+24\cdot16^{9D^2}\beta^{-34D^2}V$, and
$24\cdot16^{9D^2}\beta^{-34D^2}\le\frac12$ as $\beta\ge10^{12}$. As $V<\infty$, $V\le16\beta^{16D^2}I\le
\beta^{20D^2}I$. Finally $\pbvol(\pbeB_y)\le V$ and $\beta^{20D^2}=b^{60(r+k)^2}$.
\end{proof}

\subsection{Conclusion: the statement (DS4)}\label{pbp:ss-summary}
\textup{(DS4)} of Theorem~\ref{pbe:ds} holds (Theorem~\ref{pbp:main}), with $x_\tau$ the point of
Proposition~\ref{pbp:point} (composed with $\delta^{-1}$ of Lemma~\ref{pbp:sign} if the component of
$\psi_E$ at infinity is $e^{2\pi i(w+\bar w)}$) and $v_\tau$ the normalised monomial $t_k^{a_\tau}$ of
Lemma~\ref{pbp:W}. Published results used: Schwarz's lemma (Witt's theorem is proved as Lemma~\ref{pbp:tb-witt}; see also \cite[Chapter~7, Theorem~9.1]{Scharlau});
the rest is proved here or in \S\S\ref{pco:sec}, \ref{pds:sec}.


%% ===== proofs/B-gain.tex =====
% proofs/B-gain.tex
% Statement (G') of proofs/B-pillars.tex (Proposition ppi:G): torsion index times discriminant of
% the common parts. Labels of this file start with pbg:.
\providecommand{\auO}{\mathcal O}
\providecommand{\auNm}{\operatorname{Nm}}
\providecommand{\auA}{\mathbb A}
\providecommand{\aubar}[1]{\overline{#1}}
\providecommand{\pbg}{\mathsf g}
\providecommand{\pbeAp}{A^{+}}
\providecommand{\pbeCL}{\mathrm{CL}}
\providecommand{\ppiLc}{\Lambda^{X}_{c}}
\providecommand{\ppidisc}{\operatorname{disc}}
% macros of this file (prefix \pgn)
\providecommand{\pgnD}{\mathfrak d}
\providecommand{\pgnd}{\operatorname{d}}
\providecommand{\pgnN}{\mathrm N}
\providecommand{\pgnsrc}{\operatorname{src}}
\providecommand{\pgnlm}[2]{\langle #1,#2\rangle_{\Lambda}}
\providecommand{\pgnlf}[2]{\langle\!\langle #1,#2\rangle\!\rangle}

\section{The gain for the common parts}\label{pbg:sec}

This section proves statement (G$'$) of Proposition~\ref{ppi:G}, from which that proposition
deduces the gain (G) of Theorem~\ref{ppi:pillars}. We use the notation of \S\S\ref{pbe:sec}
and~\ref{ppi:sec}: the frame $f_0,\dots,f_\varrho$ ($f_a=e_{a+1}$ for $a<\varrho$, $f_\varrho=z$),
$A=\langle f_0,\dots,f_{\varrho-1}\rangle$, $Q=A^\perp/A=X_0\perp Ez$, $Q^+=X_0\perp Ez^+$, the matrix
$Z(\gamma)=(h(f_a,\gamma f_b))_{a,b}$ with leading minors $\delta_l(\gamma)$ ($\delta_0=1$), the
discriminant $\pbg(\gamma)=\delta_n\aubar{\delta_n}-N^2\delta_{n-1}\aubar{\delta_{n-1}}$, the link
$\Lambda_\gamma$ and the isometry $\Upsilon_\gamma$ of Lemma~\ref{pbe:link}, the cell rank
(Definition~\ref{ppi:def-cell}) and the gauge \eqref{ppi:eq-gauge}. Hermitian forms are written here
linear in the second variable, the opposite of convention (Cv1) of \S\ref{pb:ss-setting}; the statements
of this section concern lattices, indices and absolute norms of Gram determinants, which do not change
when $h$ is replaced by $\aubar h$ (\S\ref{pbv:ss-weak}(iv)). For $x\in E$ we write $\auNm x=\prod_{\tau}|\tau(x)|^2=|\pgnN_{E/\Q}(x)|$, and for $x\in
E^+$, $\auNm x=\prod_\tau|\tau(x)|$, as in \S\ref{ppi:sec}. Gram determinants, the constant $c$ of
Lemma~\ref{pbg:jacobi} and the integer $N$ are always normed as elements of $E$, $\auNm x=|\pgnN_{E/\Q}(x)|$;
only the discriminant $\pbg(\gamma)$ is normed as an element of $E^+$.

The lattice of a link is computed exactly through a bordered
Gram determinant (Proposition~\ref{pbg:L1}); passing to the level costs $M^{2j'f}$
(Proposition~\ref{pbg:Lplus}); the link lattice is cut along the intersection with a second link
(Proposition~\ref{pbg:cut}); the two complements are paired integrally after division by $M$
(Proposition~\ref{pbg:congr}); and the common lattice has index dividing a torsion index in the
common part (Proposition~\ref{pbg:index}).

\subsection{Lattices and discriminants}\label{pbg:ss-lattices}

Write $\auO=\auO_E$. A \emph{lattice} of rank $\rho$ in an $E$-subspace $V$ of dimension $\rho$ of a
hermitian space is a finitely generated $\auO$-submodule $L\subseteq V$ with $EL=V$. For a hermitian
form $q$ on $V$ (not necessarily the restriction of $h$) let $\pgnD_q(L)$ be the fractional ideal of
$\auO$ generated by the Gram determinants $\det(q(\ell_a,\ell_b))_{a,b\le\rho}$ of the
$\rho$-tuples of elements of $L$, and
\[
\pgnd_q(L)=\pgnN(\pgnD_q(L))^{1/2}\in[0,\infty),
\]
$\pgnN$ the absolute norm of fractional ideals (with $\pgnN(0)=0$). If $q$ is nondegenerate on $V$,
$\pgnd_q(L)>0$. We drop $q$ when it is $h$. For a lattice $L'$ in $V\subseteq X_0$ contained in
$\ppiLc$, $\ppidisc L'$ of \S\ref{ppi:ss-statement} equals $\pgnd(L')$, by Lemma~\ref{pbg:basic}(a)
applied to $L'\subseteq\ppiLc\cap V$. Every lattice has a \emph{pseudo-basis}: an $E$-basis $\ell=
(\ell_1,\dots,\ell_\rho)$ of $V$ and fractional ideals $\mathfrak a_l$ with $L=\bigoplus_l\mathfrak a_l
\ell_l$ (Steinitz). Writing $q(\ell,\ell')$ for the matrix $(q(\ell_a,\ell'_b))_{a,b}$ and $\mathfrak a(L)=
\prod_l\mathfrak a_l$, one has $\pgnD_q(L)=\mathfrak a(L)\aubar{\mathfrak a(L)}\det q(\ell,\ell)$, since
the determinants of the $\rho\times\rho$ matrices with columns in $\bigoplus\mathfrak a_l$ generate
$\mathfrak a(L)$. In particular $\pgnd_q(L)^2=\pgnN(\mathfrak a(L))^2\auNm\det q(\ell,\ell)$.

\begin{lemma}\label{pbg:basic}
Let $q$ be a hermitian form on $V$.
\begin{enumerate}
\item[(a)] If $L'\subseteq L$ are lattices of the same rank, then $\pgnd_q(L')=[L:L']\,\pgnd_q(L)$.
\item[(b)] Let $V=V_1\perp V_2$ for $q$, $L$ a lattice in $V$, $L_1=L\cap V_1$ (a lattice in $V_1$) and
$L_2$ the image of $L$ under the orthogonal projection $\pi_2$ to $V_2$. Then $L_2$ is a lattice in
$V_2$ and $\pgnD_q(L)=\pgnD_q(L_1)\pgnD_q(L_2)$.
\item[(c)] \textup{(Cauchy--Schwarz for Gram determinants.)} Let $q$ be positive definite at every complex place, and let $L$,
$L'$ be lattices of rank $r$ in subspaces of $V$. Let $\mathfrak p$ be the fractional ideal generated
by the determinants $\det q(\ell,\ell')$, $\ell$, $\ell'$ running over the $r$-tuples of $L$ and $L'$.
If $\mathfrak p\subseteq\auO$, $\mathfrak p\ne0$ and $x\le\pgnN\mathfrak p$, then $x^2\le\pgnd_q(L)^2\,
\pgnd_q(L')^2$.
\item[(d)] Let $q$ be positive definite at every complex place, $L'$ a lattice of rank $i$ and $u$ an
isometry of $V$ with $1-u$ injective on $EL'$. Then $\pgnd_q((1-u)L')\le4^{if}\pgnd_q(L')$.
\end{enumerate}
\end{lemma}
\begin{proof}
(a) Localising at the primes of $E^+$ (where $\auO$ is a principal ideal ring), $L'=L C$ for bases,
$\det q(\ell',\ell')=\aubar{\det C}\det C\det q(\ell,\ell)$, and $\pgnN(\det C)$ is the local index
(elementary divisors); multiply over the primes.

(b) $L/L_1\cong L_2$ is torsion free, hence projective; so $L=L_1\oplus L_2'$ with $\pi_2:L_2'\to L_2$
bijective. Take pseudo-bases $(\ell_1,\mathfrak a_1)$ of $L_1$ and $(\ell'_2,\mathfrak a_2)$ of $L'_2$;
then $(\ell_2,\mathfrak a_2)$ with $\ell_2=\pi_2\ell'_2$ is one of $L_2$, and $(\ell_1,\ell'_2)$ with
the ideals $\mathfrak a_1,\mathfrak a_2$ is one of $L$. Now $\ell'_2=\ell_2+\ell_1Y$ for a matrix $Y$
over $E$ (as $\ell'_2-\pi_2\ell'_2\in V_1$), and column operations give $\det q((\ell_1,\ell'_2),(\ell_1,
\ell'_2))=\det q((\ell_1,\ell_2),(\ell_1,\ell_2))=\det q(\ell_1,\ell_1)\det q(\ell_2,\ell_2)$, as
$\ell_1\perp\ell_2$.

(c) Take pseudo-bases $(\ell,\mathfrak a_l)$, $(\ell',\mathfrak a'_l)$; then $\mathfrak p=\aubar{\mathfrak
a(L)}\mathfrak a(L')\det q(\ell,\ell')$. The Gram matrix of $(\ell,\ell')$ is positive semidefinite at
every complex place $\tau$, so $|\tau\det q(\ell,\ell')|^2\le\tau\det q(\ell,\ell)\cdot\tau\det q(\ell',
\ell')$ (for a positive semidefinite $\bigl(\begin{smallmatrix}P&R\\R^*&S\end{smallmatrix}\bigr)$ one has
$|\det R|^2\le\det P\det S$: if $P>0$ this follows from $S-R^*P^{-1}R\ge0$, and in general by
continuity). The product over $\tau$ gives $\pgnN(\mathfrak p)^2\le\pgnd_q(L)^2\pgnd_q(L')^2$.

(d) With a pseudo-basis $(b,\mathfrak a_l)$ of $L'$, $((1-u)b,\mathfrak a_l)$ is one of $(1-u)L'$; put
$b'=ub$. At each complex place $G_\pm=q(b\pm b',b\pm b')$ are positive semidefinite and $G_-+G_+=
2q(b,b)+2q(b',b')=4q(b,b)$. The determinant is monotone on positive semidefinite matrices, so
$\tau\det G_-\le4^i\tau\det q(b,b)$. Take the product of the squares over $\tau$.
\end{proof}

\begin{lemma}[Jacobi's identity with a border]\label{pbg:jacobi}
Let $(\bar Q,h)$ be a unimodular hermitian lattice of rank $m$. Let $F=(F_1,\dots,F_\sigma)$ be
linearly independent vectors of $\bar Q$, let $p=(p_1,p_2)$ be vectors of $\bar Q$, and let $0\ne
c\in\auO$. Let $P^\perp=\{x\in\bar Q:h(F_a,x)=0\ \forall a\}$, of rank $m-\sigma$, and let $\iota=
[F^{\mathrm{sat}}:\textstyle\sum\auO F_a]$, where $F^{\mathrm{sat}}=\bar Q\cap\sum EF_a$. On $\bar
Q\otimes E$ put $\tilde h(x,y)=h(x,y)+c^{-1}\sum_{i=1,2}\aubar{h(p_i,x)}h(p_i,y)$. Then, with
\[
\mathsf G=\begin{pmatrix}h(F,F)&h(F,p)\\h(p,F)&h(p,p)+c\cdot1_2\end{pmatrix},
\qquad
\pgnd_{\tilde h}(P^\perp)^2\,\iota^2\,\auNm(c)^2=\auNm\det\mathsf G .
\]
\end{lemma}
\begin{proof}
Localise. Let $g$ be a basis of $F^{\mathrm{sat}}$, so that $F=gC$ with $\pgnN(\det C)=\iota$;
$F^{\mathrm{sat}}$ is a direct summand of $\bar Q$, so $g$ extends to a basis $R=(g,e)$ of $\bar Q$.
Let $\mathfrak G=h(R,R)$; $\det\mathfrak G$ is a unit. The dual basis $R^\#=R\mathfrak G^{-1}$
(with $h(R,R^\#)=1$) is a basis of $\bar Q^\#=\bar Q$; a vector $x=\sum_kR^\#_kh(R_k,x)$ of $\bar Q$
lies in $P^\perp$ if and only if its coefficients at $g$ vanish, so the last $m-\sigma$ vectors $n$ of
$R^\#$ form a basis of $P^\perp$, and $h(n,n)=(\mathfrak G^{-1})_{22}$ (the block of the indices of
$e$). Write $p=R\pi$; then $h(R^\#,p)=\pi$, so $h(n,p)=\pi_2$ (the rows of $\pi$ at $e$) and
$\tilde h(n,n)=(\mathfrak G^{-1})_{22}+c^{-1}\pi_2\pi_2^*$. Consider
\[
\tilde{\mathfrak G}=\begin{pmatrix}\mathfrak G&\mathfrak G\pi\\\pi^*\mathfrak G&\pi^*\mathfrak
G\pi+c\end{pmatrix}=\begin{pmatrix}1&0\\\pi^*&1\end{pmatrix}\begin{pmatrix}\mathfrak G&0\\0&c
\end{pmatrix}\begin{pmatrix}1&\pi\\0&1\end{pmatrix},\quad
\tilde{\mathfrak G}^{-1}=\begin{pmatrix}\mathfrak G^{-1}+c^{-1}\pi\pi^*&-c^{-1}\pi\\-c^{-1}\pi^*&c^{-1}
\end{pmatrix},
\]
with $\det\tilde{\mathfrak G}=c^2\det\mathfrak G$. Let $S_1$ be the indices of $g$ and of the last two
rows, $S_2$ those of $e$. Jacobi's theorem on complementary minors, $\det(\tilde{\mathfrak
G}^{-1})_{S_2S_2}=\det\tilde{\mathfrak G}_{S_1S_1}/\det\tilde{\mathfrak G}$, gives $\det\tilde
h(n,n)=\det\hat{\mathsf G}/(c^2\det\mathfrak G)$ with $\hat{\mathsf G}=\tilde{\mathfrak G}_{S_1S_1}$,
which is $\mathsf G$ with $g$ in place of $F$. Finally $\det\mathsf G=\aubar{\det C}\det C\det\hat{
\mathsf G}$. Take norms.
\end{proof}

\begin{lemma}\label{pbg:indexmap}
Let $D\subseteq B'$ and $K$, $K'$ be lattices, and $\phi$ a linear map with $\phi(D)\subseteq K$ and
$K'\subseteq\phi(B')+K$. If $[B':D]$ is finite, then $[K':K\cap K']$ divides $[B':D]$.
\end{lemma}
\begin{proof}
$K'/(K\cap K')$ embeds into $(\phi(B')+K)/K$, a quotient of $B'/D$ (as $\phi(D)\subseteq K$); the
order of a subgroup of a quotient of a finite group divides the order of the group.
\end{proof}

\subsection{The subframe minor and the gauge}\label{pbg:ss-minor}

Fix classes $c,c'$ and $\gamma\in\Gamma_{c,s,c'}$, and put $\Lambda=h_cL$ (this differs from the $\Lambda$ of \S\ref{ppi:sec}). As $h_c\in H(\auA_f)$ fixes
the $e_a$, $e_a^*$ and $z$, and $L=\bigoplus_a(\auO e_a\oplus\auO e_a^*)\perp L_{\mathrm{mid}}$ with
$L_{\mathrm{mid}}$ unimodular and $z\in L_{\mathrm{mid}}$ (proof of Lemma~\ref{pb:KR}), the lattice
$\Lambda$ is unimodular, $\Lambda\cap A=\bigoplus_a\auO e_a$, a vector of $\Lambda$ orthogonal to $A$
is $a+m$ with $a\in\Lambda\cap A$, $m\in\Lambda\cap(X_0\oplus Ez)$, and $z\in\Lambda$. As $k_\gamma=
h_c^{-1}\gamma h_{c'}\in K$, $\gamma h_{c'}L=\Lambda$; in particular $\gamma f_b\in\Lambda$ for all $b$.

Let $j=j(\gamma)$ be the cell rank and $j'=\min(\varrho,j)$, and assume $\delta_j(\gamma)\ne0$ and
$j\le\varrho\iff\pbg(\gamma)=0$ (by Proposition~\ref{ppi:cell} this holds for $\gamma\in\Gamma^*$). Put
\[
\mu(\gamma)=\auNm\delta_j(\gamma)\,N^{-2(j-\varrho)_+f},
\]
so that \eqref{ppi:eq-gauge} reads $\mathsf G(\gamma)=M^{2j'f}\max(\auNm\pbg(\gamma),N^f\mu(\gamma))/(
i_\gamma i'_\gamma)$. Let $S_j=\{0,\dots,j'-1\}\cup\{\varrho\}$ and let $\delta_*(\gamma)=\det Z(\gamma)_{
S_jS_j}$ be the \emph{subframe minor}.

\begin{lemma}[Shared vectors]\label{pbg:shared}
Let $I=A\cap\gamma A$.
\begin{enumerate}
\item[(a)] If $j\le\varrho$, there is a common vector $f\hat x=\gamma f\hat y$ with $\hat y_\varrho=1$
and $\hat y_a=0$ for $j\le a<\varrho$; it satisfies $\hat x_\varrho\aubar{\hat x_\varrho}=1$ and
$(Z(\gamma)\hat y)_b=0$ for $b<\varrho$.
\item[(b)] $I\cap\langle f_0,\dots,f_{j'-1}\rangle=0$, and $I$ has a basis $u_b=f_b+\sum_{a<j'}c_{ab}
f_a$, $j'\le b<\varrho$. So $\dim I=\varrho-j'$.
\item[(c)] A vector $a\in A$ with $h(\gamma f_b,a)=0$ for $0\le b\le\varrho$ lies in $I$.
\end{enumerate}
\end{lemma}
\begin{proof}
(a) By the definition of $j$ there is $\zeta$,
$\zeta\aubar\zeta=1$, such that for each $a$ with $j\le a\le\varrho$ there is a common vector $fx^{(a)}=
\gamma fy^{(a)}$ with top index $a$ and $(x^{(a)}_a,y^{(a)}_a)=(\zeta,1)$. Starting from $a=\varrho$ and
subtracting suitable multiples of the vectors with $a=\varrho-1,\dots,j$, we get a common vector
$f\hat x=\gamma f\hat y$ with $\hat y_\varrho=1$ and $\hat y_a=0$ for $j\le a<\varrho$. As $A$ is
isotropic and orthogonal to $z$, $-N\hat x_\varrho\aubar{\hat x_\varrho}=h(f\hat x,f\hat x)=h(f\hat y,f
\hat y)=-N$; and for $b<\varrho$, $(Z\hat y)_b=h(f_b,\gamma f\hat y)=h(f_b,f\hat x)=0$.

(b) Let $fx=\gamma fy\in I$ with $x$ supported on $\{0,\dots,j'-1\}$; then $y_\varrho=0$. For every
$b$, $h(\gamma f_b,fx)=h(f_b,fy)=0$ ($fy\in A$). For $b<j'$ this says that the transpose conjugate of
the leading block of order $j'$ of $Z$ kills $x$. If $j\le\varrho$ this block has determinant
$\delta_j\ne0$; if $j=\varrho+1$ use all $b\le\varrho$ and $\delta_n\ne0$ (with $x$ extended by $0$).
So $x=0$. For $j'\le a<\varrho$ (so $j\le\varrho$) the common vector $fx^{(a)}=\gamma fy^{(a)}$ of (a)
has top index $a<\varrho$, so it lies in $A\cap\gamma A=I$ and has coefficient $\zeta\ne0$ at $f_a$.
Triangular elimination gives the $u_b$; by the first claim they span $I$.

(c) Subtract from $a$ the combination of the $u_b$ with the same coefficients at $f_b$, $b\ge j'$;
the difference $a''$ lies in $\langle f_0,\dots,f_{j'-1}\rangle$ and is still orthogonal to all
$\gamma f_b$ (as $I\subseteq\gamma A$ is orthogonal to $\gamma A+E\gamma z$). The argument of (b) gives
$a''=0$.
\end{proof}

\begin{lemma}[The subframe minor]\label{pbg:subframe}
$\auNm\delta_*(\gamma)=N^{2f}\mu(\gamma)$.
\end{lemma}
\begin{proof}
If $j=\varrho+1$, $S_j=\{0,\dots,\varrho\}$ and $\delta_*=\delta_n=\delta_j$, $\mu=\auNm\delta_n/N^{2f}$.
Let $j\le\varrho$, and let $P$ be the matrix whose rows $a<j$ are those of $Z=Z(\gamma)$ and whose other
rows are those of the identity; $\det P=\delta_j\ne0$. By Lemma~\ref{pbg:shared}(a), $P\hat y=
\epsilon_\varrho$ (the last unit vector). Expanding along the identity rows, $\delta_*$ is the
determinant of $P$ with its last row replaced by the last row $Z_\varrho$ of $Z$; by Cramer's rule
this is $(Z_\varrho\hat y)\det P$. And $Z_\varrho\hat y=h(z,\gamma f\hat y)=h(z,f\hat x)=-N\hat x_\varrho$.
So $\delta_*=-N\hat x_\varrho\delta_j$ and, as $|\tau(\hat x_\varrho)|=1$ at every $\tau$, $\auNm\delta_*
=N^{2f}\auNm\delta_j$.
\end{proof}

\begin{remark}[The gauge]\label{pbg:rem-gauge}
(i) The gauge may also be written as $M^{2j'f}\max(\mathsf g^\sharp,\auNm\pbg)$ with $\mathsf g^\sharp=
\auNm\delta_*/(i_\gamma i'_\gamma N^f)$. By Lemma~\ref{pbg:subframe}, $\mathsf g^\sharp=N^f\mu/(i_\gamma
i'_\gamma)$; and if $\pbg(\gamma)\ne0$ then $j=\varrho+1$ and $i_\gamma=i'_\gamma=1$ (see (ii)), so the
expression with $\mathsf g^\sharp$ is \eqref{ppi:eq-gauge}. (ii) The indices are, exactly: $i_\gamma=[R_A:R_{<j'}]$, where
$R_A$ (resp.\ $R_{<j'}$) is the $\auO$-span in $E^n$ of the rows $a<\varrho$ (resp.\ $a<j'$) of
$Z(\gamma)$, and $i'_\gamma$ is the same for $Z(\gamma^{-1})=Z(\gamma)^*$. (In lattice coordinates the
rows are the products of $e_a$ with the frame $k_\gamma f$, as $h(e_a,k_\gamma f_b)=h(e_a,\gamma f_b)$.)
They are finite: by Lemma~\ref{pbg:shared}(b),(c), the kernel of $a\mapsto(h(\gamma f_b,a))_b$ on $A$ is
$I$, of codimension $j'$, so rows $a<j'$ span the row space over $E$. For $j=\varrho+1$, $j'=\varrho$
and $i_\gamma=i'_\gamma=1$.
\end{remark}

\subsection{The lattice of a link}\label{pbg:ss-link}

For a pair $P=(y_0+\lambda'z,\ y_0'+\lambda z)\in Q^2$ ($y_0,y_0'\in X_0$) put $\pgnsrc P=y_0+\lambda z^+
\in Q^+$. On $Q^2$ we use two forms: the \emph{link metric}
$\pgnlm PR=\tfrac12\bigl(h^+(\pgnsrc P,\pgnsrc R)+h^+(\pgnsrc\bar P,\pgnsrc\bar R)\bigr)$, where $\bar
P=(y_0'+\lambda z,y_0+\lambda'z)$ is the swapped pair, which is positive definite at every complex
place; and the \emph{link form} $\pgnlf{(p,p')}{(r,r')}=h(p,r)-h(p',r')$. By Lemma~\ref{pbe:link}(b),
for $P,R\in\Lambda_\gamma$ one has $\pgnsrc\bar P=\Upsilon_\gamma\pgnsrc P$, so $\pgnlm PR=h^+(\pgnsrc
P,\pgnsrc R)$, and $\pgnsrc:\Lambda_\gamma\to Q^+$ is an isometric isomorphism.

\begin{lemma}\label{pbg:linkforms}
Let $\gamma,\gamma'\in\Gamma$.
\begin{enumerate}
\item[(a)] For $t,t'\in(A+\gamma A)^\perp$ with images $P,R\in\Lambda_\gamma$, $\pgnlf PR=0$, and
$\Lambda_\gamma$ is its own orthogonal for $\pgnlf\cdot\cdot$ in $Q^2$.
\item[(b)] For $t\in(A+\gamma A)^\perp$, $t'\in(A+\gamma'A)^\perp$ with images $P\in\Lambda_\gamma$,
$R\in\Lambda_{\gamma'}$, $\pgnlf PR=h(t,t')-h(t,\gamma\gamma'^{-1}t')$.
\item[(c)] For $P\in\Lambda_\gamma$, $R\in\Lambda_{\gamma'}$, $\pgnlf PR=h^+(a,b-ub)$, where $a=\pgnsrc P$,
$b=\pgnsrc R$ and $u=\Upsilon_\gamma^{-1}\Upsilon_{\gamma'}$.
\item[(d)] For $t\in(A+\gamma A)^\perp$ with image $P$, $\pgnlm PP=h(t,t)+N^{-1}(|h(z,t)|^2+|h(\gamma z,
t)|^2)$; and the two linear forms $\zeta_1(p,p')=h(z,p)$, $\zeta_2(p,p')=h(z,p')$ take at $P$ the values
$h(z,t)$, $h(\gamma z,t)$. On $\Lambda_\gamma$, $\zeta_1=\pgnlm{\varepsilon_1}\cdot$ and $\zeta_2=
\pgnlm{\varepsilon_2}\cdot$, where $\varepsilon_1,\varepsilon_2\in\Lambda_\gamma$ have sources
$-\Upsilon_\gamma^{-1}z^+$ and $-z^+$.
\end{enumerate}
\end{lemma}
\begin{proof}
Classes in $Q=A^\perp/A$ carry the induced form. (b) $h([t],[t'])=h(t,t')$ as $t,t'\perp A$, and
$h([\gamma^{-1}t],[\gamma'^{-1}t'])=h(\gamma^{-1}t,\gamma'^{-1}t')=h(t,\gamma\gamma'^{-1}t')$ as
$\gamma^{-1}t,\gamma'^{-1}t'\perp A$. (a) is (b) with $\gamma'=\gamma$; as $\dim\Lambda_\gamma=k+1$ is
half of $\dim Q^2$ and $\pgnlf\cdot\cdot$ is nondegenerate, $\Lambda_\gamma$ is its own orthogonal. (c) Write
$P=(y_0+\lambda'z,y_0'+\lambda z)$, $R=(w_0+\mu'z,w_0'+\mu z)$. Then $h(p,r)-h(p',r')=(h(y_0,w_0)-
N\aubar{\lambda'}\mu')-(h(y_0',w_0')-N\aubar\lambda\mu)=h^+(a,b)-h^+(\Upsilon_\gamma a,\Upsilon_{\gamma'}b)$,
and $h^+(\Upsilon_\gamma a,\Upsilon_{\gamma'}b)=h^+(a,ub)$. (d) With $t=\alpha+y_0+\lambda'z$, $\alpha\in A$,
and $\gamma^{-1}t=\alpha'+y_0'+\lambda z$: $h(t,t)=|y_0|^2-N|\lambda'|^2$, $h(z,t)=-N\lambda'$, $h(\gamma
z,t)=h(z,\gamma^{-1}t)=-N\lambda$, and $\pgnlm PP=|y_0|^2+N|\lambda|^2$. For the last claim,
$\pgnlm{\varepsilon_1}P=-h^+(\Upsilon_\gamma^{-1}z^+,a)=-h^+(z^+,\Upsilon_\gamma a)=-N\lambda'=h(z,p)$, and
$\pgnlm{\varepsilon_2}P=-h^+(z^+,a)=-N\lambda=h(z,p')$.
\end{proof}

Let $\gamma\in\Gamma_{c,s,c'}$ be as in \S\ref{pbg:ss-minor}. Define the lattices of $\Lambda_\gamma$
(for the link metric)
\[
L^1_\gamma=\{([t],[\gamma^{-1}t]):t\in\Lambda\cap(A+\gamma A)^\perp\},\qquad
L^+_\gamma=\{([t],[\gamma^{-1}t]):t\in\Lambda\cap(A+\gamma A)^\perp,\ t-m\in M(\Lambda\cap A)\},
\]
where in the second $t=\alpha+m$ with $\alpha\in A$, $m\in X_0\oplus Ez$. They are lattices of rank $k+1$
($t\mapsto([t],[\gamma^{-1}t])$ maps $(A+\gamma A)^\perp$ onto $\Lambda_\gamma$ with kernel $I$, and
$L^+_\gamma\supseteq ML^1_\gamma$). We write $\pgnd$ for $\pgnd_{\pgnlm\cdot\cdot}$ on $\Lambda_\gamma$.
Let $t(\gamma)=[\Lambda\cap(A+\gamma A):\Lambda\cap A+\gamma(\Lambda\cap A)]$, the torsion index of
Proposition~\ref{ppi:G} written in the coordinates of $\Lambda$.

\begin{proposition}[The lattice $L^1_\gamma$]\label{pbg:L1}
$\pgnd(L^1_\gamma)\,i_\gamma\,i'_\gamma\,t(\gamma)\le\mu(\gamma)$.
\end{proposition}
\begin{proof}
\emph{The quotient.} $I=A\cap\gamma A$ is isotropic and $\Lambda\cap I$ is primitive in $\Lambda$. Let
$\bar Q=(\Lambda\cap I^\perp)/(\Lambda\cap I)$, a lattice in $I^\perp/I$ with the induced form. It is
unimodular: a linear form $\phi$ on $\bar Q$ lifts to $\Lambda\cap I^\perp$, and extends to $\Lambda$
(as $\Lambda\cap I^\perp$ is primitive); since $\Lambda$ is unimodular, $\phi=h(w,\cdot)$ with $w\in
\Lambda$, and $w\perp\Lambda\cap I$, so $w\in\Lambda\cap I^\perp$ (as $\Lambda\cap I$ spans $I$) and
$\phi=h([w],\cdot)$; and $\bar Q$ is integral. With $n_{\mathrm{tot}}=\dim Y_0$, its rank is $\dim I^\perp/I=n_{\mathrm{tot}}-2(\varrho-j')=2j'+k+1$.

\emph{The two lists.} The vectors $f_a,\gamma f_a$ ($a<\varrho$) lie in $\Lambda\cap I^\perp$ (as $A$,
$\gamma A$ are isotropic and contain $I$), and so do $z$ and $\gamma z$ ($z\perp A$, $\gamma z\perp\gamma
A$). Let $F$ be the classes in $\bar Q$ of $f_a,\gamma f_a$ for $a<j'$, and $p=([z],[\gamma z])$. By
Lemma~\ref{pbg:shared}(b), the class of $f_b$ for $b\ge j'$ is that of $-\sum_{a<j'}c_{ab}f_a$, and
likewise for $\gamma f_b$ (apply (b) to $\gamma^{-1}$, which has the same cell rank, the same $\pbg$
and $\delta_l(\gamma^{-1})=\aubar{\delta_l(\gamma)}$, since $Z(\gamma^{-1})=Z(\gamma)^*$, and use $A\cap
\gamma^{-1}A=\gamma^{-1}I$); hence the span $W$ of $F$ is the image of
$A+\gamma A$, which has dimension $(\varrho+j')-(\varrho-j')=2j'$, and $F$ is linearly independent.

\emph{The complement.} A vector $t\in\Lambda\cap I^\perp$ is orthogonal to $A+\gamma A$ if and only if
$[t]\perp W$. So $t\mapsto[t]$ maps $\Lambda\cap(A+\gamma A)^\perp$ onto $P^\perp=\{x\in\bar Q:x\perp
F\}$ (onto, because $\Lambda\cap I^\perp\to\bar Q$ is onto), with kernel $\Lambda\cap I$; composing
with the isomorphism induced by $t\mapsto([t],[\gamma^{-1}t])$, whose kernel is also $I$, we get an
isomorphism of $P^\perp$ onto $L^1_\gamma$, which by Lemma~\ref{pbg:linkforms}(d) is isometric from
$\tilde h=h+N^{-1}(|h(p_1,\cdot)|^2+|h(p_2,\cdot)|^2)$ to the link metric. Lemma~\ref{pbg:jacobi} with
$c=N$ gives
\[
\pgnd(L^1_\gamma)^2\,\iota^2\,N^{4f}=\auNm\det\mathsf G,\qquad\iota=[F^{\mathrm{sat}}:\textstyle\sum
\auO F_a].
\]

\emph{The bordered determinant.} $\mathsf G$ is the Gram matrix of the two lists $(f_a\ (a<j'),\ z+w_1)$
and $(\gamma f_a\ (a<j'),\ \gamma z+w_2)$ in $I^\perp/I\perp Ew_1\perp Ew_2$, with $h(w_i,w_i)=N$. Each
list is isotropic ($h(z+w_1,z+w_1)=-N+N=0$), and the products of the first with the second form
the matrix $Z(\gamma)_{S_jS_j}$. So $\det\mathsf G=\pm\delta_*\aubar{\delta_*}$ and, by
Lemma~\ref{pbg:subframe}, $\auNm\det\mathsf G=(N^{2f}\mu)^2$. Hence $\pgnd(L^1_\gamma)\iota=\mu$.

\emph{The index.} It remains to show that $i_\gamma i'_\gamma t(\gamma)$ divides $\iota$. Let $L_1$ be
the $\auO$-span of $F$, $L_2$ that of the classes of $f_a$ ($a<\varrho$) and $\gamma f_a$ ($a<j'$), $L_3$
that of the classes of all $f_a,\gamma f_a$ ($a<\varrho$), and $L_4=\bar Q\cap W=F^{\mathrm{sat}}$; then
$L_1\subseteq L_2\subseteq L_3\subseteq L_4$ and $\iota=[L_4:L_1]$. (i) The map $\phi(x)=(h(\gamma f_b,
x))_{b\le\varrho}$ is defined on $\bar Q$ (as $\gamma f_b\perp I$); it sends the class of $f_a$ to the
conjugate of row $a$ of $Z$, and that of $\gamma f_a$ to $(h(f_b,f_a))_b=0$. By
Lemma~\ref{pbg:indexmap} (with $B'=L_2$, $D=L_1$, $K=\phi(L_1)$, $K'=\phi(L_2)$), $i_\gamma$ divides
$[L_2:L_1]$. (ii) In the same way, with $\phi'(x)=(h(f_b,x))_b$, which kills the classes of the
$f_a$ and sends that of $\gamma f_a$ to the conjugate of row $a$ of $Z(\gamma^{-1})$, $i'_\gamma$
divides $[L_3:L_2]$. (iii) Choose a section $\beta:\bar Q\to\Lambda\cap I^\perp$, $\auO$-linear, of the
projection. For $w\in\Lambda\cap(A+\gamma A)$, $w-\beta([w])\in\Lambda\cap I\subseteq\Lambda\cap A$; so
$\beta$ maps $L_3$ into $T_0=\Lambda\cap A+\gamma(\Lambda\cap A)$ and $L_4$ into $\Lambda\cap(A+\gamma A)$,
and $\Lambda\cap(A+\gamma A)\subseteq\beta(L_4)+T_0$. Lemma~\ref{pbg:indexmap} gives $t(\gamma)\mid[L_4:
L_3]$. Multiply.
\end{proof}

\begin{proposition}[The lattice $L^+_\gamma$]\label{pbg:Lplus}
$t(\gamma)\,\pgnd(L^+_\gamma)\le M^{2j'f}\mu(\gamma)/(i_\gamma i'_\gamma)$.
\end{proposition}
\begin{proof}
Let $\Phi(\alpha)=(h(\gamma f_b,\alpha))_{b\le\varrho}$ for $\alpha\in A$, and $R=\Phi(\Lambda\cap A)$, a
lattice of rank $\varrho-\dim I=j'$ (Lemma~\ref{pbg:shared}(c)). For $P=([t],[\gamma^{-1}t])\in
L^1_\gamma$ with $t=\alpha+m$, the class of $\Phi(\alpha)$ in $R/MR$ depends only on $P$ ($t$ is
determined modulo $\Lambda\cap I$, and $\Phi(I)=0$). If it vanishes, $\Phi(\alpha)=M\Phi(\omega)$ with
$\omega\in\Lambda\cap A$, so $\alpha-M\omega\in\Lambda\cap I$ by Lemma~\ref{pbg:shared}(c), and $t'=t-(
\alpha-M\omega)$ lies in $\Lambda\cap(A+\gamma A)^\perp$, has the same image $P$ and $A$-part $M\omega$:
$P\in L^+_\gamma$. So $L^+_\gamma$ contains the kernel of $L^1_\gamma\to R/MR$, and $[L^1_\gamma:
L^+_\gamma]\le\#R/MR=\pgnN(M\auO)^{j'}=M^{2j'f}$. By Lemma~\ref{pbg:basic}(a) and
Proposition~\ref{pbg:L1}, $t(\gamma)\pgnd(L^+_\gamma)\le M^{2j'f}t(\gamma)\pgnd(L^1_\gamma)\le M^{2j'f}
\mu/(i_\gamma i'_\gamma)$.
\end{proof}

\subsection{The cut along a second link}\label{pbg:ss-cut}

\begin{lemma}[The vectors $z^+$ and $\Upsilon_\gamma^{-1}z^+$]\label{pbg:linkvec}
Let $\varkappa_\gamma$ be the $z^+$-coordinate of $\Upsilon_\gamma^{-1}z^+$, and $\mathsf M(\gamma)=\max(
\auNm\pbg(\gamma),N^f\mu(\gamma))$. The Gram matrix of $(z^+,\Upsilon_\gamma^{-1}z^+)$ for $h^+$ is
$\bigl(\begin{smallmatrix}N&N\varkappa_\gamma\\N\aubar{\varkappa_\gamma}&N\end{smallmatrix}\bigr)$. If $j=\varrho+1$,
then $\pbg(\gamma)=\delta_n\aubar{\delta_n}(1-\varkappa_\gamma\aubar{\varkappa_\gamma})$. Consequently, for every
principal minor $\mathsf m$ of order $r\le1+(j-\varrho)_+$ of this Gram matrix, $\auNm\mathsf m\le(\mathsf
M(\gamma)/\mu(\gamma))^2$.
\end{lemma}
\begin{proof}
The diagonal entries are $h^+(z^+,z^+)=N$ and $h^+(\Upsilon^{-1}z^+,\Upsilon^{-1}z^+)=N$; the
off-diagonal one is $h^+(z^+,\Upsilon^{-1}z^+)=N\varkappa_\gamma$. Let $j=\varrho+1$, so $\delta=\delta_n\ne0$;
let $y=\operatorname{adj}(Z)\epsilon_\varrho$, so that $Zy=\delta\epsilon_\varrho$ and $y_\varrho=
\delta_{n-1}$. Put $t=\gamma fy$. For $b<\varrho$, $h(f_b,t)=(Zy)_b=0$ and $h(\gamma f_b,t)=h(f_b,fy)=0$; so
$t\perp A+\gamma A$. Now $[\gamma^{-1}t]=[fy]=\delta_{n-1}z$, and $h(z,t)=(Zy)_\varrho=\delta$, so the
$z$-coefficient of $[t]$ is $-\delta/N$. By Lemma~\ref{pbe:link}(b), $\Upsilon_\gamma(y_0+\delta_{n-1}z^+)=
-(\delta/N)z^+$, where $y_0$ is the $X_0$-part of $t$; so $\Upsilon_\gamma^{-1}z^+=-(N/\delta)(y_0+
\delta_{n-1}z^+)$ and $\varkappa_\gamma=-N\delta_{n-1}/\delta$. Hence $1-\varkappa_\gamma\aubar{\varkappa_\gamma}=(\delta
\aubar\delta-N^2\delta_{n-1}\aubar{\delta_{n-1}})/\delta\aubar\delta=\pbg/\delta\aubar\delta$.

For the minors: $r=0$ gives $1\le(N^f\mu/\mu)^2$; $r=1$ gives $\auNm N=N^{2f}=(N^f\mu/\mu)^2$; $r=2$
occurs only if $j=\varrho+1$, and then $\mathsf m=N^2(1-\varkappa_\gamma\aubar{\varkappa_\gamma})=N^2\pbg/\delta_n
\aubar{\delta_n}$, so $\auNm\mathsf m=N^{4f}(\auNm\pbg)^2/(\auNm\delta_n)^2=(\auNm\pbg/\mu)^2$, as $\mu=
\auNm\delta_n/N^{2f}$ and $\pbg\in E^+$.
\end{proof}

The factor $\mathsf{mn}(s)$ rests on the following residual form of Lemma~\ref{ppi:climb}. Recall
(\S\ref{pb:sec}) that $s\in\auG$ meets at a place $\mathfrak P\in\{w,\aubar w\}$ above a chosen $w$ if
the spans of the frame and of its image under $s$ modulo $\mathfrak P$ meet nontrivially, and
$\mathsf{mn}(s)=\prod\auNm\mathfrak P$ over these places.

\begin{lemma}[Climbing modulo a level place]\label{pbg:climbp}
Let $s$ agree, let $\kappa\in K$ reduce to $s$, and let $s$ meet at $\mathfrak P$. Then there are $x,y\in
\auO^n$ with $fx-\kappa fy\in\mathfrak PL$ and $(x_\varrho,y_\varrho)\not\equiv(0,0)\bmod\mathfrak P$.
\end{lemma}

\begin{proof}
Write $\mathfrak P$ as a pair $(w,\epsilon)$: $\mathfrak P=w$ or $\mathfrak P=\aubar w$. An element of
$K$ is determined modulo $J$ by its reductions modulo the places $(w',\epsilon)$, $w'$ chosen, with the
\emph{same} $\epsilon$ (the reduction modulo $\aubar{w'}$ is the inverse adjoint of that modulo $w'$);
so for $u,u'\in K$, $[u]\,s=s\,[u']$ in $\auG$ (where $[u]$ denotes the image of $u$ in $\auG$) if and only if $u\kappa\equiv \kappa u'$ modulo every place
$(w',\epsilon)$. Write $\tilde\cdot$ for reduction modulo $\mathfrak P$, $\mathbb F=\auO/\mathfrak P$, and
$\tilde\kappa$ for the matrix of $\kappa$ modulo $\mathfrak P$.

\emph{Lifts.} By the Chinese remainder theorem, every $\theta\in\mathbb F$ has a lift $\hat\theta\in
\widehat\auO_E$ with $\hat\theta\equiv0$ modulo every place above $M$ other than $\mathfrak P$ (in
particular modulo $\aubar{\mathfrak P}$), and $\hat0=0$. Then $\aubar{\hat\theta}\equiv0$ modulo every
place above $M$ other than $\aubar{\mathfrak P}$, in particular modulo $\mathfrak P$ and modulo all
$(w',\epsilon)$.

\emph{One step.} Let $\tilde f x=\tilde\kappa\tilde f y$ with $x,y\in\mathbb F^n$ vanishing above $a<\varrho$,
$(x_a,y_a)=(\xi,\xi')$. As in the proof of Lemma~\ref{ppi:climb}, $(\xi,\xi')$ occurs at $a+1$ modulo
$\mathfrak P$ as soon as every $\mathbb F$-linear form $\lambda$ on $L/\mathfrak PL$ that vanishes on
$\tilde f_c,\tilde\kappa\tilde f_c$ ($c\le a$) satisfies $\xi\lambda(\tilde f_{a+1})=\xi'\lambda(\tilde\kappa\tilde f_{a+1})$.
The frame vectors $f_c$ belong to a basis $\beta$ of $L$. Define $\widehat\auO$-linear forms $m,m'$ on
$L\otimes\widehat\auO$ by $m(\beta_l)=\widehat{\lambda(\tilde\beta_l)}$ and $m'(\beta_l)=\widehat{\lambda(\tilde
\kappa\tilde\beta_l)}$ (lifts as above). Then $m\equiv\lambda$, $m'\equiv\lambda\circ\tilde\kappa$ modulo $\mathfrak P$,
both vanish modulo every place above $M$ other than $\mathfrak P$, and $m(f_c)=m'(f_c)=0$ exactly for
$c\le a$ (as $\hat0=0$). As $L$ is unimodular, $m=h(q,\cdot)$ and $m'=h(q',\cdot)$ with $q,q'\in L\otimes
\widehat\auO$, orthogonal to $f_0,\dots,f_a$. Let $\hat x,\hat y$ be the coordinatewise lifts of $x,y$,
$p=f\hat x$, $p'=f\hat y\in\langle f_0,\dots,f_a\rangle$, fix $\omega$ as in Lemma~\ref{ppi:eichler} and
$r\in\widehat\auO_E$, and let
\[
u=\mathsf E_{rp,q},\qquad u'=\mathsf E_{rp',q'}
\]
be the Eichler elements of Lemma~\ref{ppi:eichler} (with $j=a<\varrho$; the lemma is pure algebra and
holds over $\auA_f$). They are integral, so $u,u'\in U(\auA_f)\cap K$, and $\sigma(u)=r\hat x_am(f_{a+1})$,
$\sigma(u')=r\hat y_am'(f_{a+1})$. Now
\[
u(t)=t+h(q,t)\,rp-\aubar r\,h(p,t)\,(q+\varsigma rp),\qquad h(p,t)=\textstyle\sum_b\aubar{\hat x_b}\,h(f_b,t).
\]
\emph{This is where both $w$ and $\aubar w$ enter:} the last term carries the conjugates
$\aubar{\hat x_b}$, which vanish modulo $\mathfrak P$ because $\hat x_b\equiv0$ modulo $\aubar{\mathfrak
P}$, whereas the second term carries $\hat x_b$ and $m$. Hence modulo $\mathfrak P$, $\tilde u=1+\tilde r
\,(\tilde fx)\,\lambda$, and modulo every place $(w',\epsilon)\ne\mathfrak P$ both terms vanish ($\hat x$,
$\aubar{\hat x}$ and $m$ vanish there), so $u\equiv1$. Likewise $\tilde u'=1+\tilde r(\tilde fy)(\lambda\circ
\tilde\kappa)$ and $u'\equiv1$ at the other places $(w',\epsilon)$. Modulo $\mathfrak P$,
\[
\tilde u\tilde\kappa=\tilde\kappa+\tilde r(\tilde fx)(\lambda\circ\tilde\kappa)=\tilde\kappa+\tilde r(\tilde\kappa\tilde fy)(\lambda\circ\tilde\kappa)=
\tilde\kappa\tilde u' ,
\]
and $u\kappa\equiv \kappa\equiv \kappa u'$ at the other places. So $[u]\,s=s\,[u']$ in $\auG$ and, $s$ agreeing, $\psi_f(u)=
\psi_f(u')$, that is $\psi_{E,f}(r(\Theta-\Theta')/(M\varpi))=1$ with $\Theta=\hat x_am(f_{a+1})$,
$\Theta'=\hat y_am'(f_{a+1})$, for every $r\in\widehat\auO_E$. Since $\varpi$ generates the different of
$E$, the $\alpha\in\auA_{E,f}$ with $\psi_{E,f}(r\alpha)=1$ for all $r\in\widehat\auO_E$ form $\varpi^{-1}
\widehat\auO_E$ (the dual of $\widehat\auO_E$ for $\psi_{E,f}$ is the inverse different, (Cv2) of \S\ref{pb:ss-setting}). Hence $\Theta-\Theta'\in
M\widehat\auO_E$, and modulo $\mathfrak P$: $\xi\lambda(\tilde f_{a+1})=\xi'\lambda(\tilde\kappa\tilde f_{a+1})$.

\emph{Climbing.} As $s$ meets at $\mathfrak P$, there are $x,y\in\mathbb F^n$, $x\ne0$, with $\tilde fx=
\tilde\kappa\tilde fy$. Let $a$ be the top index (the largest index with $(x_a,y_a)\ne(0,0)$; it exists as
$x\ne0$). If $a<\varrho$, by the step applied $\varrho-a$ times, $(x_a,y_a)$ occurs at $a+1,\dots,\varrho$. Lift the resulting $x,y$ to $\auO^n$.
\end{proof}

\begin{lemma}[Meeting places]\label{pbg:meet}
Let $s$ agree, $\gamma\in\Gamma_{c,s,c'}$, and $t_1,t_2\in\Lambda\cap(A+
\gamma A)^\perp$. Then $\Delta=\det\bigl(\begin{smallmatrix}h(z,t_1)&h(z,t_2)\\h(\gamma z,t_1)&h(\gamma
z,t_2)\end{smallmatrix}\bigr)\in\auO$ lies in $\aubar{\mathfrak P}$ for every place $\mathfrak P$ at
which $s$ meets; so every ideal generated by such determinants has norm at least $\mathsf{mn}(s)$ if it
is not $0$.
\end{lemma}
\begin{proof}
Put $\kappa=k_\gamma$ and $\tilde t_l=h_c^{-1}t_l\in L$. As $h_c$, $h_{c'}$ fix the frame and $h_c$ is an
isometry, $h(z,t_l)=h(z,\tilde t_l)$, $h(\gamma z,t_l)=h(\kappa z,\tilde t_l)$, and $\tilde t_l\perp A+\kappa A$. Take
$x,y$ as in Lemma~\ref{pbg:climbp}. Since $\tilde t_l\perp A$ and $\tilde t_l\perp \kappa A$, $h(fx,\tilde t_l)=
\aubar{x_\varrho}h(z,\tilde t_l)$ and $h(\kappa fy,\tilde t_l)=\aubar{y_\varrho}h(\kappa z,\tilde t_l)$, and $h(fx-\kappa fy,
\tilde t_l)\in h(\mathfrak PL,L)\subseteq\aubar{\mathfrak P}$. So $(\aubar{x_\varrho},-\aubar{y_\varrho})$, which
is nonzero modulo $\aubar{\mathfrak P}$, is in the left kernel of the matrix of $\Delta$ modulo
$\aubar{\mathfrak P}$, and $\Delta\in\aubar{\mathfrak P}$. The entries are integral as $z,\gamma z,t_l
\in\Lambda$. The places $\aubar{\mathfrak P}$ are distinct primes with the norms $\auNm\mathfrak P$.
\end{proof}

\begin{proposition}[The cut]\label{pbg:cut}
Let $s$ agree, $\gamma\in\Gamma_{c,s,c'}\cap\Gamma^*$ and $\gamma'\in\Gamma$. Let $\Lambda_\cap=
\Lambda_\gamma\cap\Lambda_{\gamma'}$, $i=k+1-\dim\Lambda_\cap$, $V_y$ the orthogonal complement of
$\Lambda_\cap$ in $\Lambda_\gamma$ for the link metric, $\pi_y$ the orthogonal projection onto $V_y$, and
$L_y=\pi_y(L^+_\gamma)$, a lattice of rank $i$ in $V_y$. Let $C_\gamma=L_\gamma\cap\pbeCL(\gamma,\gamma')$.
Let $x$ be an integer with $1\le x\le\mathsf{mn}(s)$, and $x=1$ unless $i=0$ and $\pbg(\gamma)\ne0$. Then $C_\gamma$ is a lattice of rank $\dim\pbeCL(\gamma,\gamma')$
and
\[
t(\gamma)\,\pgnd(C_\gamma)\,x\,\pgnd(L_y)\le\mathsf G(\gamma).
\]
\end{proposition}
\begin{proof}
\emph{Three pieces.} Let $K_z=\ker\zeta_1\cap\ker\zeta_2$, the pairs of vectors of $X_0$; $V_m=
\Lambda_\cap\cap K_z$; $V_u$ the orthogonal complement of $V_m$ in $\Lambda_\cap$, of dimension $r$. The
first projection maps $V_m$ isometrically (link metric to $h$) onto $\pbeCL(\gamma,\gamma')$, by the
definition of $\pbeCL$ (a pair in $\Lambda_\gamma$ is determined by its source, which for a pair in
$K_z$ is its first member). The space $\Lambda_\gamma\cap K_z$ is the graph of $\Upsilon_\gamma$ on
$X_0^\gamma$, of dimension $k-1$ if $\pbg(\gamma)\ne0$ and $k$ otherwise (Lemma~\ref{pbe:link}(c) and
Proposition~\ref{ppi:cell}), that is, $k-(j-\varrho)_+$. As $\Lambda_\gamma\cap K_z+\Lambda_\cap\subseteq
\Lambda_\gamma$, $\dim V_m\ge k-(j-\varrho)_++\dim\Lambda_\cap-(k+1)$, so $r\le1+(j-\varrho)_+$; and if
$i=0$ and $j=\varrho+1$, then $\Lambda_\cap=\Lambda_\gamma$, $V_m=\Lambda_\gamma\cap K_z$ and $r=2$. The
forms $\zeta=(\zeta_1,\zeta_2)$ are injective on $V_u$, as $V_u\cap K_z\subseteq V_u\cap V_m=0$.

Apply Lemma~\ref{pbg:basic}(b) to $\Lambda_\gamma=\Lambda_\cap\perp V_y$ and $L^+_\gamma$, and then to
$\Lambda_\cap=V_m\perp V_u$ and $L^+_\gamma\cap\Lambda_\cap$: with $L_m=L^+_\gamma\cap V_m$ and $L_u$ the
projection of $L^+_\gamma\cap\Lambda_\cap$ to $V_u$,
\[
\pgnd(L^+_\gamma)=\pgnd(L_m)\,\pgnd(L_u)\,\pgnd(L_y).
\]

\emph{The common part.} Let $P=([t],[\gamma^{-1}t])\in L_m$, $t=\alpha+m\in\Lambda\cap(A+\gamma A)^\perp$
with $\alpha\in M(\Lambda\cap A)$. As $\zeta_1(P)=h(z,t)=0$, $m\in X_0$, and $m=t-\alpha\in\Lambda$, so
$m\in\ppiLc$; it is the first member of $P$, so $m\in\pbeCL(\gamma,\gamma')$. As $\zeta_2(P)=h(\gamma
z,t)=0$ and $t\perp\gamma A$, the vector $m+\alpha\in m+M(\Lambda\cap A)$ is orthogonal to $\gamma f_0,
\dots,\gamma f_\varrho$, so $m\in L_\gamma$. Hence the first projection maps $L_m$ isometrically into
$C_\gamma\subseteq\ppiLc\cap\pbeCL(\gamma,\gamma')$; its image has rank $\dim\pbeCL$. So $C_\gamma$ is a
lattice of that rank and, by Lemma~\ref{pbg:basic}(a), $\pgnd(C_\gamma)\le\pgnd(L_m)$.

\emph{The part seen by the $z$-forms.} Choose $\rho\subseteq\{1,2\}$ with $\#\rho=r$ such that
$\zeta_\rho$ is an isomorphism of $V_u$ onto $E^r$ ($\rho=\{1,2\}$ if $r=2$). By
Lemma~\ref{pbg:linkforms}(d), $\zeta_\rho=\pgnlm{\varepsilon_\rho}\cdot$ on $\Lambda_\gamma$. For $P\in
L_u$, $P=P'-v$ with $P'\in L^+_\gamma\cap\Lambda_\cap$ and $v\in V_m\subseteq K_z$, so $\zeta(P)=
\zeta(P')=(h(z,t),h(\gamma z,t))$ for a lift $t\in\Lambda$ of $P'$, which is integral. So the ideal
$\mathfrak p$ generated by the $\det\zeta_\rho(\ell)$, $\ell$ an $r$-tuple of $L_u$, is integral and
nonzero, and $\pgnN\mathfrak p\ge x$: trivially if $x=1$, and by Lemma~\ref{pbg:meet} if $x\ne1$ (then
$r=2$). Apply Lemma~\ref{pbg:basic}(c) to the free lattice $\sum_{l\in\rho}\auO\varepsilon_l$ (whose
$r$-tuples give the same ideal $\mathfrak p$) and $L_u$; the Gram matrix of $\varepsilon_\rho$ is a
principal minor of that of $(\Upsilon_\gamma^{-1}z^+,z^+)$, so by Lemma~\ref{pbg:linkvec} (for $r=0$ all
is trivial)
\[
x\,\mu(\gamma)\le\mathsf M(\gamma)\,\pgnd(L_u).
\]

\emph{Conclusion.} With Proposition~\ref{pbg:Lplus},
\[
t(\gamma)\pgnd(C_\gamma)\,x\,\pgnd(L_y)\le t(\gamma)\pgnd(L_m)\pgnd(L_u)\pgnd(L_y)\frac{\mathsf M}{\mu}=
t(\gamma)\pgnd(L^+_\gamma)\frac{\mathsf M}{\mu}\le\frac{M^{2j'f}\mathsf M(\gamma)}{i_\gamma i'_\gamma}=
\mathsf G(\gamma). \qedhere
\]
\end{proof}

\subsection{The congruence and the index}\label{pbg:ss-congr}

\begin{proposition}[The congruence]\label{pbg:congr}
Let $\gamma,\gamma'\in\Gamma_{c,s,c'}\cap\Gamma^*$, and let $L_y\subseteq V_y\subseteq\Lambda_\gamma$ and
$L'_y\subseteq V'_y\subseteq\Lambda_{\gamma'}$ be as in Proposition~\ref{pbg:cut}, for $(\gamma,\gamma')$
and for $(\gamma',\gamma)$, both of rank $i$. Then $1\le(2/M)^{2if}\,\pgnd(L_y)\,\pgnd(L'_y)$.
\end{proposition}
\begin{proof}
If $i=0$ both sides are $1$. \emph{Nondegeneracy.} Let $R\in V'_y$ with $\pgnlf PR=0$ for all $P\in V_y$.
As $\Lambda_\cap\subseteq\Lambda_{\gamma'}$, $\pgnlf PR=0$ also for $P\in\Lambda_\cap$
(Lemma~\ref{pbg:linkforms}(a)), so $R$ is orthogonal to $\Lambda_\gamma$, hence $R\in\Lambda_\gamma$ by
Lemma~\ref{pbg:linkforms}(a); so $R\in\Lambda_\cap\cap V'_y=0$. Thus $\pgnlf\cdot\cdot$ pairs $V_y$ and
$V'_y$ nondegenerately.

\emph{Integrality.} $\delta=\gamma\gamma'^{-1}$ satisfies $h_c^{-1}\delta h_c=k_\gamma k_{\gamma'}^{-1}\in K$
with reduction $ss^{-1}=1$, so $h_c^{-1}\delta h_c\in J=J_M$ and $(\delta-1)\Lambda\subseteq M\Lambda$.
For $P\in L^1_\gamma$, $R\in L^1_{\gamma'}$ with lifts $t,t'\in\Lambda$, Lemma~\ref{pbg:linkforms}(b)
gives $\pgnlf PR=h(t,(1-\delta)t')\in M\auO$. An element of $L_y$ is $P-C$ with $P\in L^+_\gamma\subseteq
L^1_\gamma$ and $C\in\Lambda_\cap$, and an element of $L'_y$ is $R-C'$ likewise; as $C,C'\in\Lambda_
\gamma\cap\Lambda_{\gamma'}$, Lemma~\ref{pbg:linkforms}(a) gives $\pgnlf{P-C}{R-C'}=\pgnlf PR\in M\auO$.

\emph{The norm bound.} Let $u=\Upsilon_\gamma^{-1}\Upsilon_{\gamma'}$, $S=M^{-1}\pgnsrc(L_y)$ and $S'=(1-u)
\pgnsrc(L'_y)$, lattices in $Q^+$. By Lemma~\ref{pbg:linkforms}(c), $\pgnlf PR=h^+(\pgnsrc P,(1-u)\pgnsrc
R)$ for $P\in V_y$, $R\in V'_y$. So $1-u$ is injective on $\pgnsrc V'_y$ (by the nondegeneracy), and the
ideal generated by the determinants $\det h^+(\ell,\ell')$, $\ell$, $\ell'$ running over the $i$-tuples
of $S$, $S'$, is $M^{-i}$ times the ideal generated by the $\det\pgnlf{\ell_1}{\ell'_1}$ over the
$i$-tuples of $L_y$, $L'_y$; it is integral and nonzero. Lemma~\ref{pbg:basic}(c) with $q=h^+$ and $x=1$,
then Lemma~\ref{pbg:basic}(d) and $\pgnd_{h^+}(\pgnsrc L)=\pgnd(L)$, give
\[
1\le\pgnd_{h^+}(S)^2\,\pgnd_{h^+}(S')^2\le M^{-4if}\pgnd(L_y)^2\cdot4^{2if}\pgnd(L'_y)^2 ,
\]
as scaling by $M^{-1}$ multiplies the Gram determinants of $i$-tuples by $M^{-2i}$. Take square roots.
\end{proof}

\begin{proposition}[The index]\label{pbg:index}
Let $\gamma,\gamma'\in\Gamma_{c,s,c'}$, and let $C_\gamma=L_\gamma\cap\pbeCL(\gamma,\gamma')$. Then
$[C_\gamma:C_\gamma\cap L_{\gamma'}]$ divides $t(\gamma')$; note $C_\gamma\cap L_{\gamma'}=L_{\gamma\gamma'}$.
\end{proposition}
\begin{proof}
$L_{\gamma\gamma'}=L_\gamma\cap L_{\gamma'}\cap\pbeCL=C_\gamma\cap L_{\gamma'}$. Let $\phi(v)=(h(\gamma'f_b,v)
)_{b\le\varrho}$ and $R'=\phi(\Lambda\cap A)$. A vector $v\in\ppiLc$ lies in $L_{\gamma'}$ if and only if
$\phi(v)\in MR'$ ($v+M\omega\perp\gamma'f$ means $\phi(v)=-M\phi(\omega)$). So $[C_\gamma:C_\gamma\cap
L_{\gamma'}]=[\phi(C_\gamma):\phi(C_\gamma)\cap MR']$.

\emph{Claim: $\phi(C_\gamma)\subseteq M\phi(\Lambda\cap(A+\gamma'A))+MR'$.} Let $v\in C_\gamma$, and
$\omega\in\Lambda\cap A$ with $y=v-M\omega\perp\gamma f_0,\dots,\gamma f_\varrho$. Then $y\in\Lambda$,
$y\perp A$ (as $v\in X_0$), so $y\in(A+\gamma A)^\perp$ and $y\perp\gamma z$: the pair $([y],[\gamma^{-1}
y])=(v,v')$ lies in $\Lambda_\gamma$ with $v'\in X_0$. As $v\in\pbeCL(\gamma,\gamma')$, some pair $(v,v'')\in
\Lambda_\gamma\cap\Lambda_{\gamma'}$ has $v''\in X_0$; both pairs lie in $\Lambda_\gamma$ and have the
source $v$, so $v''=v'$ and $(v,v')\in\Lambda_{\gamma'}$: there is $t'\perp A+\gamma'A$ with $[t']=v$,
$[\gamma'^{-1}t']=v'$. Then $t'-y\in A$ and $\gamma^{-1}y-\gamma'^{-1}t'\in A$, and with $m=\gamma'\gamma^{-1}$,
\[
(m-1)y=(t'-y)+\gamma'(\gamma^{-1}y-\gamma'^{-1}t')\in A+\gamma'A .
\]
As in Proposition~\ref{pbg:congr}, $(m-1)\Lambda\subseteq M\Lambda$, so $w=-M^{-1}(m-1)y\in\Lambda\cap(A+
\gamma'A)$. Since $my\perp m\gamma f=\gamma'f$, $\phi(y)=\phi(my)-\phi((m-1)y)=M\phi(w)$; and $\phi(y)=
\phi(v)-M\phi(\omega)$. So $\phi(v)=M\phi(w)+M\phi(\omega)$.

\emph{Conclusion.} $\phi(\gamma'(\Lambda\cap A))=0$ ($h(\gamma'f_b,\gamma'e_a)=h(f_b,e_a)=0$), so $M\phi$
maps $T'_0=\Lambda\cap A+\gamma'(\Lambda\cap A)$ into $MR'$; by the claim, $\phi(C_\gamma)\subseteq
M\phi(\Lambda\cap(A+\gamma'A))+MR'$. Lemma~\ref{pbg:indexmap} (with $B'=\Lambda\cap(A+\gamma'A)$,
$D=T'_0$, $K=MR'$, $K'=\phi(C_\gamma)$) shows that $[\phi(C_\gamma):\phi(C_\gamma)\cap MR']$ divides $[B':
T'_0]=t(\gamma')$, which is finite as $T'_0$ spans $A+\gamma'A$.
\end{proof}

\subsection{Statement (G$'$)}\label{pbg:ss-main}

\begin{theorem}\label{pbg:main}
Let $s$ agree, $\gamma,\gamma'\in\Gamma_{c,s,c'}\cap\Gamma^*$, $C_\gamma=L_\gamma\cap\pbeCL(\gamma,\gamma')$,
$C_{\gamma'}=L_{\gamma'}\cap\pbeCL(\gamma',\gamma)$, $i=i(\gamma,\gamma')$, and $\mathsf x=\mathsf{mn}(s)$ if
$i=0$ and $\pbg(\gamma)\ne0$, $\mathsf x=1$ otherwise. Then $C_\gamma$ is a lattice of rank $\dim\pbeCL(
\gamma,\gamma')$, $[C_\gamma:L_{\gamma\gamma'}]$ divides $t(\gamma')$, and
\[
t(\gamma)\,\ppidisc C_\gamma\cdot t(\gamma')\,\ppidisc C_{\gamma'}\cdot\mathsf x^2\le(2/M)^{2if}\,\mathsf
G(\gamma)\,\mathsf G(\gamma').
\]
That is, statement \textup{(G$'$)} of Proposition~\ref{ppi:G} holds.
\end{theorem}
\begin{proof}
The first two claims are Propositions~\ref{pbg:cut} and~\ref{pbg:index}. $\pbeCL$ and $i$ are symmetric
in $\gamma,\gamma'$. If $i=0$ then $\Lambda_\gamma=\Lambda_{\gamma'}$, so $X_0^\gamma=X_0^{\gamma'}$, and by
Lemma~\ref{pbe:link}(c) and Proposition~\ref{ppi:cell} $\pbg(\gamma)\ne0$ if and only if $\pbg(\gamma')\ne
0$. So $x=\mathsf x$ is admissible in Proposition~\ref{pbg:cut} for $(\gamma,\gamma')$ and for $(\gamma',
\gamma)$, and multiplying the two bounds,
\[
t(\gamma)\ppidisc C_\gamma\,t(\gamma')\ppidisc C_{\gamma'}\,\mathsf x^2\,\pgnd(L_y)\pgnd(L'_y)\le\mathsf G(
\gamma)\mathsf G(\gamma').
\]
Multiply by $(2/M)^{2if}$ and use Proposition~\ref{pbg:congr}. Here $\ppidisc=\pgnd$ on lattices of
$\ppiLc$, and $t(\gamma)$ is the torsion index of (G$'$) as defined in
Proposition~\ref{ppi:G}; the two expressions given there agree because $h_c,h_{c'}$ fix $A$ pointwise,
so $\gamma A=h_ck_\gamma A$ and $\gamma(h_cL\cap A)=h_ck_\gamma(L\cap A)$.
\end{proof}


%% ===== proofs/B-count.tex =====
% proofs/B-count.tex
% Statements (C1) and (C2) of Proposition ppi:C (proofs/B-pillars.tex): the count of near terms over
% one lattice, and the sum over the lattices of a given corank. Labels of this file start with pbc:.
\providecommand{\auO}{\mathcal O}
\providecommand{\auNm}{\operatorname{Nm}}
\providecommand{\auA}{\mathbb A}
\providecommand{\aubar}[1]{\overline{#1}}
\providecommand{\auf}[1]{\texttt{#1}}
\providecommand{\pbR}{\overline R_0}
\providecommand{\pbet}{\beta_*}
\providecommand{\pbeGs}{\Gamma^*}
\providecommand{\pbeAp}{A^{+}}
\providecommand{\pbel}{\ell}
\providecommand{\ppiLc}{\Lambda^{X}_{c}}
\providecommand{\ppidisc}{\operatorname{disc}}
% macros of this file (prefix \pbc)
\providecommand{\pbcD}{\mathfrak D}
\providecommand{\pbcB}{\mathcal B}
\providecommand{\pbcadj}{\operatorname{adj}}
\providecommand{\pbcspan}{\operatorname{span}}
\providecommand{\pbcrd}{\mathrm{rd}_E}
\providecommand{\pbcHom}{\operatorname{Hom}}

\section{The count over one lattice and the sum over lattices}\label{pbc:sec}

This section proves the two statements (C1) and (C2) of Proposition~\ref{ppi:C}, from which that
proposition deduces the count (C) of Theorem~\ref{ppi:pillars}. We use the notation of
\S\S\ref{pbe:sec} and~\ref{ppi:sec}: the frame $f_0,\dots,f_\varrho$ ($f_a=e_{a+1}$ for $a<\varrho$,
$f_\varrho=z$), $n=\varrho+1$, $A=\langle e_1,\dots,e_\varrho\rangle$, the decomposition $Y_0=A\oplus
X_0\oplus Ez\oplus A^*$ with $h(e_a,e_b^*)=\delta_{ab}$, the lattice $L$ of $Y_0$, the classes
$c\in\mathcal C$ with representatives $h_c\in H(\auA_f)$, the sets $\Gamma_{c,s,c'}$, the groups $U_c$,
$U_c^0=U_c\cap U^0$, the level $M$, the prime $N=\pi_N\aubar{\pi_N}$, the number $f$ of places
$\tau\in\auPl$ (so $[E:\Q]=2f$), $\pbet=b^{kf/4}$, $\ppiLc=h_cL\cap X_0$, the lattices $L_\gamma$, the
discriminant $\ppidisc$, the labels, and the cell rank $j(\gamma)$ with the leading minors
$\delta_l(\gamma)$ of the matrix $Z(\gamma)=(h(f_a,\gamma f_b))_{0\le a,b\le\varrho}$. For $x\in E$,
$\auNm x=\prod_{\tau}|\tau(x)|^2=|N_{E/\Q}(x)|$, and for a nonzero ideal $\mathfrak a$ of $\auO_E$,
$\auNm\mathfrak a=\#(\auO_E/\mathfrak a)$. Lattices are finitely generated $\auO_E$-submodules; a
lattice in $X_0$ given by local data (such as $\ppiLc$, or $u\auO_E^k$ for $u\in H(\auA_f)$) is the
$\auO_E$-lattice with these completions.

In this section, as in \S\ref{pbg:sec} and \S\ref{ppi:ss-cell}, hermitian forms are written linear in the second variable
and conjugate-linear in the first, the opposite of convention (Cv1) of \S\ref{pb:ss-setting}: we write $h$ for the form
$(x,y)\mapsto h(y,x)$. This is the reading under which the map $u$ of Lemma~\ref{pbc:moves} is linear, $\beta_{ab}=h(e_a,\gamma f_b)$
is the $e_a^*$-coordinate of $\gamma f_b$, the map $v\mapsto\alpha$ in Step~2 of the proof of Theorem~\ref{pbc:C1} is
$\auO_E$-linear, and \eqref{pbc:eq-lift} holds as written. Orthogonality, integrality, lattices and their duals, indices,
discriminants, isometries, $U_c^0$, and the norms $\auNm$ of the entries and minors of $Z(\gamma)$ (which is replaced by its
entrywise conjugate) do not change under this replacement (\S\ref{pbv:ss-weak}(iv)); so (C1) and (C2), which are counts,
and Lemmas~\ref{pbc:box}, \ref{pbc:indep}, which concern norms, hold in either convention, while the explicit formulas
inside the section (such as the map $u$ of Lemma~\ref{pbc:moves}) are to be read in the convention just stated.

We define the reduced and boxed points (Definition~\ref{pbc:def-reduced}) and prove the bound
$\|y_\tau^{-1}\tilde f_a\|\le b^{37n}$ at boxed points used in Lemma~\ref{ppi:gauge} (Lemma~\ref{pbc:box}).
We then prove (C1) (Theorem~\ref{pbc:C1}), under $\varrho\ge280$ and the sizes of Lemma~\ref{an:lem-sizes},
and (C2) (Theorem~\ref{pbc:C2}), using Minkowski's second theorem for lattices in
$\R^D$~\cite[Chapter~VIII, \S4.3, Theorem~V]{Cassels} (see also \cite[Theorem~1.3]{HenkMinima}), under the
size conditions of Remark~\ref{pbc:rem-sizes}, which are deduced there from Lemma~\ref{an:lem-sizes}.

\subsection{Reduced and boxed points}\label{pbc:ss-box}

At a place $\tau$ and for $y\in G_\tau$ let $w_a=y^{-1}\tilde f_a$ and let $o_a$, $t_{a'a}$ be as in
Lemma~\ref{pbe:len}(c): $w_a=o_a+\sum_{a'<a}t_{a'a}(y)\,o_{a'}$. For $y=(y_\tau)\in G_\infty$ and
$a'<a$ put
\[
T_{a'a}(y)=\Bigl(t_{a'a}(y_\tau)\,\frac{s_{a'}}{s_a}\Bigr)_{\tau\in\auPl}\in\C^{\auPl},
\]
with the scales $s_a=s_{a,\tau}$ of \eqref{pbe:eq-scale}. (These are the coefficients of the
Gram--Schmidt process read in the frame $f$ instead of the normalised frame $\tilde f$; Step~1 of
Proposition~\ref{pbe:count} shows how they change under $U_{c'}$.) Let $M\auO_E\subset\C^{\auPl}$ be the
lattice of the vectors $(\tau(Mx))_\tau$, $x\in\auO_E$.

\begin{lemma}\label{pbc:domain}
There is a Borel set $\pbcD\subset\C^{\auPl}$ that contains exactly one point of each coset of
$M\auO_E$ and whose points have all coordinates of absolute value at most $b^{26}$.
\end{lemma}
\begin{proof}
Let $Q=\{w\in\C^{\auPl}:|w_\tau|\le b^{26}\ \forall\tau\}$. By Step~1 of Proposition~\ref{pbe:count}
every coset of $M\auO_E$ meets $Q$; it meets $Q$ in a finite set, because $M\auO_E$ is discrete. Let
$\lambda_1,\lambda_2,\dots$ be an enumeration of $M\auO_E$, let $P$ be a half-open fundamental
parallelepiped of $M\auO_E$ (a Borel set of representatives), and for $x\in P$ let $i(x)$ be the least
$i$ with $x+\lambda_i\in Q$. Then $\pbcD=\{x+\lambda_{i(x)}:x\in P\}$ is a set of representatives in
$Q$, and it is Borel: $\pbcD=\bigcup_i\bigl(((P\cap(Q-\lambda_i))\setminus\bigcup_{i'<i}(Q-
\lambda_{i'}))+\lambda_i\bigr)$.
\end{proof}

\begin{definition}[Reduced and boxed points]\label{pbc:def-reduced}
Fix $\pbcD$ as in Lemma~\ref{pbc:domain}. A point $y\in G_\infty$ is \emph{reduced} if $T_{a'a}(y)\in
\pbcD$ for all $0\le a'<a\le\varrho$, and \emph{boxed} if moreover $\pbel_\tau(y_\tau)\le b^{6n+2}$ for
every $\tau$. An element $\gamma\in\Gamma$ is \emph{near} if there is a boxed $y$ with $\gamma^{-1}y$
boxed (here $\gamma$ acts through its image in $G_\infty$).
\end{definition}

This is the set $\mathfrak D$ of reduced points used in \S\ref{ppi:ss-statement}. The set of boxed points
is Borel (the $t_{a'a}$ and $\pbel_\tau$ are continuous on $G_\tau$), and both sets are invariant under
left multiplication by every $s\in\prod_\tau G_\tau$ with $s_\tau^{-1}f_a=f_a$ for all $a$ and
$\tau$, and lying in $\prod_\tau R_\tau$: such $s$ do not change the $w_a$, nor $\pbel_\tau$ by
Lemma~\ref{pbe:len}(a). In particular they are invariant under $U^0(E^+)$.

\begin{lemma}[Box]\label{pbc:box}
Let $y$ be boxed. Then $|t_{a'a}(y_\tau)|\le b^{26}$ for $a'<a$, and $\|y_\tau^{-1}\tilde f_a\|\le
b^{27}\pbel_\tau(y_\tau)\le b^{37n}$ for all $\tau$ and $a$.
\end{lemma}
\begin{proof}
By \eqref{pbe:eq-scale}, $|s_a/s_{a'}|\le1$ for $a'<a$, so $|t_{a'a}|=|\tau\text{-coordinate of }T_{a'a}|
\cdot|s_a/s_{a'}|\le b^{26}$. By Lemma~\ref{pbe:len}(c), $\|o_{a'}\|<\pbel_\tau(y_\tau)$, hence $\|w_a\|\le
\sum_{a'\le a}|t_{a'a}|\,\|o_{a'}\|\le n\,b^{26}\pbel_\tau(y_\tau)\le b^{27}b^{6n+2}\le b^{37n}$, using
$n\le b$ and $t_{aa}=1$.
\end{proof}

\begin{remark}\label{pbc:rem-box}
Lemma~\ref{pbc:box} is the input of Lemma~\ref{ppi:gauge}. Thus, for near $\gamma$, $|\tau(Z(\gamma)_{
ab})|\le b^{134n}$ and $\auNm\delta_l(\gamma)\le b^{270n^2f}$ for all $\tau,a,b,l$; these two bounds are
the only properties of nearness used in (C1), besides Lemma~\ref{pbc:box} itself.
\end{remark}

\subsection{Tools}\label{pbc:ss-tools}

\subsubsection*{Lattices and discriminants}
For a lattice $L'$ of rank $\rho$ in a subspace $V\subseteq X_0$ with $h(L',L')\subseteq\auO_E$, let
$\mathfrak d(L')$ be the ideal of $\auO_E$ generated by the Gram determinants $\det(h(v_i,v_{i'}))_{i,
i'\le\rho}$ of the $\rho$-tuples of $L'$, let $L'^\vee=\pbcHom_{\auO_E}(L',\auO_E)$, and let $L'^h
\subseteq L'^\vee$ be the image of $L'$ under $x\mapsto h(x,\cdot)|_{L'}$ (injective, as $h$ is definite).

\begin{lemma}\label{pbc:disc}
Let $L'\subseteq\Lambda$ be lattices in $X_0$, $\Lambda$ integral, $V=EL'$ of dimension $\rho$, $\Lambda_V=
\Lambda\cap V$.
\begin{enumerate}
\item[(a)] $\ppidisc L'=\auNm\mathfrak d(L')^{1/2}$; in particular $\ppidisc L'=[\Lambda_V:L']\,\ppidisc
\Lambda_V$, $\ppidisc L'$ does not depend on $\Lambda$, $\ppidisc(\eta L')=\ppidisc L'$ for $\eta\in
H(E^+)$, and $\ppidisc L'\ge1$.
\item[(b)] $[L'^\vee:L'^h]=(\ppidisc L')^2$, and $[L'^\vee:\Lambda_V^\vee]=[\Lambda_V:L']$ (restriction
$\Lambda_V^\vee\hookrightarrow L'^\vee$).
\item[(c)] If $\Lambda$ is unimodular, then $\ppidisc(\Lambda\cap V)=\ppidisc(\Lambda\cap V^\perp)$, where
$V^\perp$ is the orthogonal of $V$ in $X_0$.
\end{enumerate}
\end{lemma}
\begin{proof}
All indices are products of local indices, so we may localise at a prime $\mathfrak p$ of $E^+$; then
$\auO=\auO_{E,\mathfrak p}$ is a principal ideal ring, all lattices are free, and for an ideal $\mathfrak
a$, $\auNm\mathfrak a$ is the order of the quotient. (a) If $v$ is a basis of $L'$ with Gram matrix $G$,
the Gram determinant of the tuple $vC$ ($C$ a $\rho\times\rho$ matrix over $\auO$) is $\det(C)^-\det(C)
\det G$, so $\mathfrak d(L')=(\det G)$. If $v'$ is a basis of $\Lambda_V$ and $v=v'C$, then $[\Lambda_V:L']=
\auNm(\det C)$ and $\auNm\mathfrak d(L')=\auNm(\det C)^2\auNm\mathfrak d(\Lambda_V)$; this is the first
formula, given the definition of $\ppidisc$ in \S\ref{ppi:ss-statement}. An isometry preserves Gram
determinants, and $\mathfrak d(L')$ is integral and nonzero. (b) In the basis $v$ and the dual basis
of $L'^\vee$, the map $L'\to L'^\vee$ has the matrix ${}^tG$ (up to the conjugation, which preserves
norms), so its cokernel has order $\auNm(\det G)$. The second claim: in the bases $v'$ and $v=v'C$
the restriction has the matrix ${}^tC$. (c) Put $\Lambda_1=\Lambda\cap V$ and $\Lambda_2=\Lambda\cap
V^\perp$. As $\Lambda/\Lambda_1$ is torsion free, $\Lambda_1$ is a direct summand of $\Lambda$, and as
$\Lambda$ is unimodular, $x\mapsto h(x,\cdot)|_{\Lambda_1}$ maps $\Lambda$ onto $\Lambda_1^\vee$, with
kernel $\Lambda_2$. Its composite with $\Lambda_1^\vee\to\Lambda_1^\vee/\Lambda_1^h$ has kernel $\Lambda_1
\oplus\Lambda_2$ (if $h(x,\cdot)=h(x_1,\cdot)$ on $\Lambda_1$ with $x_1\in\Lambda_1$, then $x-x_1\in
\Lambda_2$). So $[\Lambda:\Lambda_1\oplus\Lambda_2]=(\ppidisc\Lambda_1)^2$ by (b), and by symmetry ($V=
(V^\perp)^\perp$, $\Lambda_2$ is also a direct summand) it equals $(\ppidisc\Lambda_2)^2$.
\end{proof}

\begin{lemma}[Packing]\label{pbc:pack}
\textup{(a)} Let $m\ge1$ and $R\ge0$. At most $(1+2R)^{2mf}$ vectors $v\in\auO_E^m$ satisfy $|\tau
(v_i)|\le R$ for all $i$ and $\tau$.
\textup{(b)} Let $\Sigma\subset\R^D$ be a set with $\|x\|\ge1$ and $\|x-x'\|\ge1$ for all $x\ne x'$ in
$\Sigma$ (Euclidean norm). Then $\sum_{x\in\Sigma}\|x\|^{-s}\le5^s$ for every real $s\ge D+2$.
\end{lemma}
\begin{proof}
(a) If $v\ne v'$, some $v_i-v'_i$ is a nonzero algebraic integer, so $\auNm(v_i-v'_i)\ge1$ and
$|\tau(v_i-v'_i)|\ge1$ for some $\tau$. Hence the polydiscs of radius $\frac12$ in $\C^{m\auPl}$ around
the points $(\tau(v_i))_{i,\tau}$ are disjoint; they lie in the polydisc of radius $R+\frac12$, and
comparing volumes gives the bound. (b) The balls of radius $\frac12$ around the points of $\Sigma$
are disjoint, so for $T\ge1$ the number $\nu(T)$ of $x\in\Sigma$ with $\|x\|\le T$ is at most
$(2T+1)^D\le(3T)^D$. Hence $\sum\|x\|^{-s}=\int_1^\infty sT^{-s-1}\nu(T)\,dT\le s3^D\int_1^\infty
T^{D-s-1}dT=s3^D/(s-D)\le\frac s2\,3^s\le5^s$.
\end{proof}

\subsubsection*{Moves}
The following elements of $U_c^0$ are used to move frames.

\begin{lemma}[Moves]\label{pbc:moves}
Let $c\in\mathcal C$, let $x_1,\dots,x_\varrho\in\ppiLc$, and let $S=(S_{ba})$ be a $\varrho\times\varrho$
matrix over $\auO_E$ with ${}^t\aubar S=-S$. Let $G=(h(x_b,x_a))_{b,a}$, let $\omega\in\auO_E$ be as in
Lemma~\ref{ppi:eichler} ($\omega+\aubar\omega=1$), put $W=M\omega G-S$, and
\[
u(t)=t+M\sum_{a=1}^\varrho h(x_a,t)\,e_a-M\sum_{a=1}^\varrho h(e_a,t)\Bigl(x_a+\sum_{b=1}^\varrho W_{ba}
e_b\Bigr)\qquad(t\in Y_0).
\]
Then $u\in U_c^0$. For $\gamma\in\Gamma$ write $\gamma f_b=\sum_a\alpha_{ab}e_a+\xi_b+\zeta_bz+\sum_a
\beta_{ab}e_a^*$ with $\xi_b\in X_0$, and let $\alpha(\gamma)=(\alpha_{ab})$, $\xi(\gamma)=(\xi_0,\dots,
\xi_\varrho)$, $\beta(\gamma)=(\beta_{ab})$ ($1\le a\le\varrho$, $0\le b\le\varrho$). Then $Z(u\gamma)=Z(
\gamma)$, $\beta(u\gamma)=\beta(\gamma)$, $\zeta(u\gamma)=\zeta(\gamma)$, and
\[
\xi(u\gamma)=\xi(\gamma)-M\,X\beta(\gamma),\qquad\text{and, if }X=0,\qquad\alpha(u\gamma)=\alpha(\gamma)+
M\,S\beta(\gamma),
\]
where $X\beta$ is the row $(\sum_ax_a\beta_{ab})_b$.
\end{lemma}
\begin{proof}
Write $u=1+M\nu$. \emph{Unitarity.} Put $\mathsf a_a=h(x_a,t)$, $\mathsf b_a=h(e_a,t)$ and similarly
$\mathsf a'_a,\mathsf b'_a$ for $t'$. Using that $A$ is isotropic and orthogonal to $X_0$, and $h(e_a,
e_b^*)=\delta_{ab}$, one finds $h(\nu t,t')+h(t,\nu t')=-\sum_{a,b}\aubar{\mathsf b_a}\mathsf b'_b(W+
{}^t\aubar W)_{ab}$ (the terms with $\mathsf a$ cancel in pairs) and $h(\nu t,\nu t')=\sum_{a,b}
\aubar{\mathsf b_a}\mathsf b'_bG_{ab}$. As $W+{}^t\aubar W=M(\omega+\aubar\omega)G=MG$, $h(ut,ut')=h(t,
t')+M\bigl(h(\nu t,t')+h(t,\nu t')+Mh(\nu t,\nu t')\bigr)=h(t,t')$. \emph{Unipotence.} $\nu$ vanishes on
$A$, maps $A^\perp=A\oplus X_0\oplus Ez$ into $A$ and $Y_0$ into $A\oplus X_0\subseteq A^\perp$. So $u$
is the identity on $A$, preserves the flag of $A$, acts trivially on $A^\perp/A$ and on $Y_0/A^\perp$;
as in the proof of Lemma~\ref{ppi:eichler}, $u\in U(E^+)$; also $\nu^3=0$. \emph{The frame.} $h(x_a,
f_b)=0$ and $h(e_a,f_b)=0$, so $\nu f_b=0$ and $u\in U^0(E^+)$. \emph{Level.} Let $L^\sharp=h_cL$; as
in the proof of Lemma~\ref{pbe:eichler}, $L^\sharp=\bigoplus_a(\widehat\auO e_a\oplus\widehat\auO
e_a^*)\oplus L^\sharp_{\rm mid}$ is unimodular, so $h(x_a,t),h(e_a,t)\in\widehat\auO$ for $t\in
L^\sharp$, and $x_a,e_a\in L^\sharp$, $W$ is integral; hence $\nu L^\sharp\subseteq L^\sharp$ and $(u-1)
L^\sharp\subseteq ML^\sharp$, $(u^{-1}-1)L^\sharp=(-M\nu+M^2\nu^2)L^\sharp\subseteq ML^\sharp$. By the
description of $U_c$ in the proof of Lemma~\ref{pbe:eichler}, $u\in U_c$.
\emph{Action.} $h(f_a,u\gamma f_b)=h(u^{-1}f_a,\gamma f_b)=h(f_a,\gamma f_b)$, and $\beta_{ab}=h(e_a,
\gamma f_b)$, $-N\zeta_b=h(z,\gamma f_b)$ are entries of $Z$. Finally $\nu(\gamma f_b)=\sum_a h(x_a,
\xi_b)e_a-\sum_a\beta_{ab}(x_a+\sum_bW_{ba}e_b)$ has $X_0$-part $-\sum_a\beta_{ab}x_a$ and, if $X=0$
(so $W=-S$), $A$-part $\sum_{a,b'}S_{b'a}\beta_{ab}e_{b'}$.
\end{proof}

\subsubsection*{Uniform independence at boxed points, and shared vectors}

\begin{lemma}\label{pbc:indep}
Let $y\in G_\tau$ with $\|y^{-1}\tilde f_a\|\le b^{37n}$ for all $a$ and $\pbel_\tau(y)\le b^{7n}$ (by
Lemma~\ref{pbc:box} this holds for every component of a boxed point). Then for every $c\in\C^n$ and every $a$,
\[
|c_a|\le b^{46n^2}\Bigl\|\sum_{a'}c_{a'}\,y^{-1}\tilde f_{a'}\Bigr\| .
\]
\end{lemma}
\begin{proof}
Let $v=\sum c_{a'}y^{-1}\tilde f_{a'}$. As $\pbel_\tau(y)\le b^{7n}$ there is $s\in R_\tau$ with $\|sy\|\le
b^{7n}$. Let $T$ be the matrix of $s$ on the normalised frame: $s\tilde f_{a'}=\sum_{a''}T_{a''a'}
\tilde f_{a''}$; it is unitriangular (\S\ref{pbe:ss-real}), so $\det T=1$ and $c=\pbcadj(T)\,Tc$. The
normalised basis is orthonormal, so $|(Tc)_{a''}|\le\|\sum_{a''}(Tc)_{a''}\tilde f_{a''}\|=\|sy\,v\|\le
b^{7n}\|v\|$, and likewise $|T_{a''a'}|\le\|sy\|\,\|y^{-1}\tilde f_{a'}\|\le b^{44n}$. The entries of
$\pbcadj T$ are minors of order $\varrho$, at most $\varrho!\,b^{44n\varrho}\le b^{(1+44n)\varrho}$ in
absolute value. Hence $|c_a|\le n\,b^{(1+44n)\varrho}b^{7n}\|v\|\le b^{46n^2}\|v\|$, as $1+(1+44n)
\varrho+7n\le46n^2$.
\end{proof}

For $\gamma\in\Gamma$ of cell rank $j=j(\gamma)$ with $\delta_j(\gamma)\ne0$, let $Y=Y_j(\gamma)$ be the
$n\times n$ matrix whose rows $a<j$ are those of $Z(\gamma)$ and whose rows $a\ge j$ are those of the
identity matrix; $\det Y=\delta_j(\gamma)$. For a coefficient vector $x\in E^n$ we write $fx=\sum x_b f_b$
and $\gamma fx=\gamma(fx)$, as in \S\ref{ppi:sec}.

\begin{lemma}[Shared vectors]\label{pbc:shared}
Let $\gamma\in\Gamma$ have cell rank $j$ with $\delta_j(\gamma)\ne0$, and $Y=Y_j(\gamma)$.
\begin{enumerate}
\item[(a)] For $j\le k\le\varrho$, $\gamma f(Y^{-1}\epsilon_k)$ lies in $\langle f_0,\dots,f_\varrho\rangle$
($\epsilon_k$ the unit vectors).
\item[(b)] If $j=\varrho+1$, the $X_0$-part $\xi$ of $\gamma f(Y^{-1}\epsilon_\varrho)$ satisfies
$\tau(h(\xi,\xi))\le1/N$ for every $\tau$.
\item[(c)] For $u\in U^0(E^+)$ the element $u\gamma$ has the same cell rank, $Z(u\gamma)=Z(\gamma)$,
$Y_j(u\gamma)=Y$, and \textup{(a)}, \textup{(b)} hold for $u\gamma$.
\end{enumerate}
\end{lemma}
\begin{proof}
(a) By definition of the cell rank there are $\zeta$ and, for each $a\ge j$, vectors $x^{(a)},y^{(a)}$
supported on $\{0,\dots,a\}$ with $fx^{(a)}=\gamma fy^{(a)}$ and $y^{(a)}_a=1$. Subtracting from
$y^{(k)}$ suitable multiples of the $y^{(a)}$, $j\le a<k$, by descending induction on the index, we
get $y$ with $y_a=\delta_{ak}$ for $a\ge j$ and $\gamma fy=fx\in\langle f\rangle$ for some $x$. For
$a<j$ we have $a<\varrho$, so $(Zy)_a=h(f_a,fx)=0$ ($A$ is isotropic and orthogonal to $z$); and $(Yy)_a=
y_a=\delta_{ak}$ for $a\ge j$. So $Yy=\epsilon_k$.
(b) Now $Y=Z$. Let $x=Z^{-1}\epsilon_\varrho$ and $t=\gamma fx$. Then $h(e_a,t)=0$ for all $a$, so $t$
has no $A^*$-part, and $h(z,t)=1$, so its $z$-coordinate is $-1/N$. Hence $h(t,t)=h(\xi,\xi)-N\cdot N^{-2}$,
while $h(t,t)=h(fx,fx)=-N|x_\varrho|^2$. So $h(\xi,\xi)=N^{-1}-N|x_\varrho|^2$, and $\tau(h(\xi,\xi))\le
1/N$.
(c) The cell rank is $U^0(E^+)$-bi-invariant (\S\ref{ppi:ss-statement}) and $Z(u\gamma)=Z(\gamma)$ as
in Lemma~\ref{pbc:moves}.
\end{proof}

\subsection{The count over one lattice: (C1)}\label{pbc:ss-C1}

\begin{theorem}[(C1)]\label{pbc:C1}
Assume $\varrho\ge280$, the sizes $M\le b^{23}$, $N\le b^{60n}$, $n\le b$ of Lemma~\ref{an:lem-sizes}\textup{(S5)}, \textup{(S7)},
and $1\le|\kappa_\tau|\le b^{23}$ \eqref{pbe:eq-kappa} (Lemma~\ref{ppi:gauge} uses only $|\kappa_\tau|\ge1$). Let $c\in
\mathcal C$, $t\in\pbR$, $0\le l\le k$, $r=k-l$, and let $L'\subseteq\ppiLc$ be a lattice of rank $r$ with
$\ppidisc L'\le b^{n^3f}$. Let $\pbcB$ be the set of the pairs $(c',\gamma)$ with $c'\in\mathcal C$ and
$\gamma\in\Gamma_{c,s,c'}\cap\pbeGs$ for some $s$, such that $\gamma$ is near, chosen, of label $t$, and
$L'\subseteq L_\gamma$. Then $\#\pbcB\le b^{2000n^6(l+1)f}$.
\end{theorem}

The proof occupies the rest of this subsection. For $(c',\gamma)\in\pbcB$ we use the coordinates
$\alpha(\gamma),\xi(\gamma),\zeta(\gamma),\beta(\gamma)$ of Lemma~\ref{pbc:moves}; the \emph{frame} of
$\gamma$ is the tuple $\gamma f=(\gamma f_0,\dots,\gamma f_\varrho)$. Two facts about the lattice
$L^\sharp=h_cL=\bigoplus_a(\widehat\auO e_a\oplus\widehat\auO e_a^*)\oplus L^\sharp_{\rm mid}$ (proof of
Lemma~\ref{pbe:eichler}) are used: since $h_c^{-1}\gamma h_{c'}\in K$ and $h_{c'}$ fixes the $f_b$,
$\gamma f_b\in L^\sharp$, so $\alpha(\gamma)$ and $\beta(\gamma)$ are integral; and $L^\sharp_{\rm mid}\cap
X_0=L^\sharp\cap X_0=\ppiLc$ (by the proof of Lemma~\ref{pb:KR} the lattice of $X_0\perp Ez$ is
spanned by $x_1,\dots,x_{k-1},x_k/\aubar{\pi_N}+z/N,z$, and $h_c$ acts on $X_0$ only).

\emph{Step 0: frames modulo moves separate.} Let $(c',\gamma),(c'',\gamma')\in\pbcB$ and $u,u'\in U_c^0$
with $u\gamma f_b=u'\gamma'f_b$ for all $b$. Then $(c',\gamma)=(c'',\gamma')$.

Indeed, $h_c^{-1}uh_c\in J$ has reduction $1$, so $u\gamma\in\Gamma_{c,s,c'}$ and $u'\gamma'\in\Gamma_{c,
s',c''}$ where $s,s'$ are the reductions of $\gamma,\gamma'$. The element $g=(u\gamma)^{-1}u'\gamma'$ fixes
the frame, so (Step~2 of Proposition~\ref{pbe:count}) $g\in R_0(E^+)$, and $h_{c'}^{-1}gh_{c''}=(h_c^{-1}
u\gamma h_{c'})^{-1}(h_c^{-1}u'\gamma'h_{c''})\in K\cap R_0(\auA_f)=K_R$ has reduction $s^{-1}s'\in\pbR$. So
$s\pbR=s'\pbR$; as the labels (the components in $\pbR$ for the fixed decomposition $\auG=\bigsqcup_\rho
\rho\pbR$) agree, $s=s'$. Now Step~2 of Proposition~\ref{pbe:count}, applied to $u\gamma,u'\gamma'\in
\Gamma_{c,s,\cdot}$, gives $c''=c'$ and $g\in U_{c'}$; $g$ fixes the frame, so $g\in U_{c'}^0$. Then $\gamma'
\in U_c^0\gamma U_{c'}^0$, and both are chosen, so $\gamma'=\gamma$.

\emph{Step 1: cell rank and products.} By Proposition~\ref{ppi:cell}, every $(c',\gamma)\in\pbcB$ has a
cell rank $j(\gamma)\in\{0,\dots,\varrho+1\}$ with $\delta_{j(\gamma)}(\gamma)\ne0$; there are $\varrho+2\le
b^2$ values. The entries of $Z(\gamma)$ are integral (they are products of vectors of $L^\sharp$) and
$|\tau(Z_{ab})|\le b^{134n}$ (Remark~\ref{pbc:rem-box}); by Lemma~\ref{pbc:pack}(a) with $m=n^2$ there are
at most $(1+2b^{134n})^{2n^2f}\le b^{270n^3f}$ values of $Z(\gamma)$. We now fix $j$ and $Z$ with $\delta:=
\delta_j(Z)\ne0$, and bound the number of elements of the set $\pbcB_{j,Z}$ of $(c',\gamma)\in\pbcB$ with
$j(\gamma)=j$, $Z(\gamma)=Z$. Let $Y=Y_j(Z)$, $P=\beta(\gamma)$ (rows $0,\dots,\varrho-1$ of $Z$; it does not
depend on $\gamma$), $d=M\delta$ and $d'=M\delta\aubar\delta=M\det(Y\,{}^t\aubar Y)$. By
Remark~\ref{pbc:rem-box}, $\auNm\delta\le b^{270n^2f}$, so with $M\le b^{23}$ and $n\ge280$:
\begin{equation}\label{pbc:eq-moduli}
\auNm d=M^{2f}\auNm\delta\le b^{2n^3f},\qquad\auNm d'=M^{2f}(\auNm\delta)^2\le b^{2n^3f}.
\end{equation}

\emph{Step 2: lifts of $L'$.} Let $(c',\gamma)\in\pbcB_{j,Z}$ and let $K_\gamma\subseteq\auO_E^\varrho$ be the
set of $\alpha$ with $\sum_a\alpha_ae_a\perp\gamma f_b$ for all $b$. For $v\in L'\subseteq L_\gamma$ there is
$\alpha\in\auO_E^\varrho$ ($h_cL\cap A=\bigoplus\auO_Ee_a$) with $v+M\sum_a\alpha_ae_a\perp\gamma f_b$ for all
$b$, and $\alpha$ is unique modulo $K_\gamma$. So $v\mapsto\alpha\bmod K_\gamma$ is an $\auO_E$-linear map
$L'\to\auO_E^\varrho/K_\gamma$; as $L'$ is projective, it lifts to a linear map $\varphi_\gamma=(\varphi_{\gamma,
1},\dots,\varphi_{\gamma,\varrho})\in(L'^\vee)^\varrho$ with
\begin{equation}\label{pbc:eq-lift}
h(v,\xi_b(\gamma))+M\sum_a\aubar{\varphi_{\gamma,a}(v)}\,P_{ab}=0\qquad(v\in L',\ 0\le b\le\varrho),
\end{equation}
since $h(v,\gamma f_b)=h(v,\xi_b)$ and $h(e_a,\gamma f_b)=P_{ab}$. Colour $\pbcB_{j,Z}$ by the class of
$\varphi_\gamma$ in $(L'^\vee/L'^h)^\varrho$: by Lemma~\ref{pbc:disc}(b) there are at most $(\ppidisc
L')^{2\varrho}\le b^{2n^3\varrho f}$ colours. Fix one colour $\mathcal C_1$ and $\gamma_1\in\mathcal C_1$. For
$\gamma\in\mathcal C_1$ there are $\upsilon_a\in L'$ with $\varphi_{\gamma,a}-\varphi_{\gamma_1,a}=h(\upsilon_a,\cdot)$; let
$u_\gamma\in U_c^0$ be the move of Lemma~\ref{pbc:moves} with $x_a=-\upsilon_a$ ($\in\ppiLc$) and $S=0$. Then
$\xi(u_\gamma\gamma)=\xi(\gamma)+Y'P$ with $Y'=(M\upsilon_a)_a$, and by \eqref{pbc:eq-lift}, for $v\in L'$,
\[
h(v,\xi_b(u_\gamma\gamma))=-M\sum_a\bigl(\aubar{\varphi_{\gamma,a}(v)}-h(v,\upsilon_a)\bigr)P_{ab}=-M\sum_a
\aubar{\varphi_{\gamma_1,a}(v)}P_{ab},
\]
as $\aubar{h(\upsilon_a,v)}=h(v,\upsilon_a)$. So $h(v,\xi_b(u_\gamma\gamma))$ depends only on $v,b$ and the colour.

\emph{Step 3: the part in $X_0$.} Let $\gamma_0\in\mathcal C_1$ and put $\tilde\gamma=u_\gamma\gamma$ for
$\gamma\in\mathcal C_1$. By Step~2 the columns of $\xi(\tilde\gamma)-\xi(\tilde\gamma_0)$ are orthogonal to
$L'$; they lie in $\ppiLc$, because $\tilde\gamma f_b,\tilde\gamma_0f_b\in L^\sharp$ have the same
$z$-coordinate $\zeta_b$ (Lemma~\ref{pbc:moves}), so their difference of middle parts lies in
$L^\sharp_{\rm mid}\cap X_0=\ppiLc$. So they lie in $P'=\ppiLc\cap(EL')^\perp$, a lattice of rank $l$.
Colour $\mathcal C_1$ by the class of $\xi(\tilde\gamma)-\xi(\tilde\gamma_0)$ in $(P'/dP')^n$: there are at
most $(\auNm d)^{ln}\le b^{2n^4lf}$ colours. Fix one, $\mathcal C_2$, and $\gamma_2\in\mathcal C_2$; for
$\gamma\in\mathcal C_2$ the columns of $v=\xi(\tilde\gamma)-\xi(\tilde\gamma_2)$ lie in $d\,\ppiLc$.

We claim that $v=M\mathsf H P$ for a row $\mathsf H=(\eta_1,\dots,\eta_\varrho)$ of vectors of $\ppiLc$. Put $\mathsf H'=d^{-1}
v\,\pbcadj(Y)$, whose columns lie in $\ppiLc$ ($\pbcadj Y$ is integral), so that $vY^{-1}=M\mathsf H'$. Let $i$
satisfy $i\ge j$ or $i=\varrho$. The column $i$ of $vY^{-1}$ is the difference of the $X_0$-parts of
$\tilde\gamma f(Y^{-1}\epsilon_i)$ and $\tilde\gamma_2f(Y^{-1}\epsilon_i)$. By Lemma~\ref{pbc:shared}
(applied to $\tilde\gamma,\tilde\gamma_2$, which have cell rank $j$ and products $Z$), each of these
$X_0$-parts is $0$ (if $i\ge j$) or has $\tau(h(\cdot,\cdot))\le1/N$ (if $i=\varrho<j$); so the column
$x\in\ppiLc$ satisfies $\tau(h(x,x))\le4/N<1$ for all $\tau$. Since $h(x,x)$ is a totally nonnegative
algebraic integer, $h(x,x)=0$ and $x=0$. Hence the columns $i\ge j$ and the column $\varrho$ of $\mathsf H'$
vanish, and $v=M\mathsf H'Y=M\sum_{i<j,\,i<\varrho}\mathsf H'_i\,(\text{row }i\text{ of }Y)=M\mathsf HP$, with $\mathsf H$ the first
$\varrho$ columns of $\mathsf H'$ (rows $i<j$ of $Y$ are rows of $P$; the columns $j\le i<\varrho$ of $\mathsf H$ vanish).

Let $u'_\gamma\in U_c^0$ be the move with $x_a=\eta_a$ and $S=0$. Then $\xi(u'_\gamma\tilde\gamma)=\xi(
\tilde\gamma)-M\mathsf HP=\xi(\tilde\gamma_2)$. Put $\hat\gamma=u'_\gamma\tilde\gamma=u'_\gamma u_\gamma\gamma$ for
$\gamma\in\mathcal C_2$ (and $u'_{\gamma_2}=1$). All $\hat\gamma$ have the products $Z$, the cell rank $j$,
the same $\zeta$ and $\beta=P$, and the same $\xi$.

\emph{Step 4: shared vectors.} For $\gamma\in\mathcal C_2$ and $j\le i\le\varrho$ let $\sigma_i(\gamma)=
\alpha(\gamma)\,\pbcadj(Y)\epsilon_i\in\auO_E^\varrho$, the $A$-part of $\gamma f(\pbcadj(Y)\epsilon_i)=\delta\,
\gamma f(Y^{-1}\epsilon_i)$; put $\sigma_i=0$ for $i<j$. By Lemma~\ref{pbc:shared}(a) this vector lies in
$\langle f\rangle$, so it is fixed by every move, and $\sigma(\hat\gamma)=\sigma(\gamma)$. We bound
$\sigma(\gamma)$. Let $x=\pbcadj(Y)\epsilon_i$; the entries of $Y$ are entries of $Z$ or $0,1$, so $|\tau(
x_b)|\le\varrho!\,b^{134n\varrho}\le b^{(1+134n)\varrho}$. Write $\gamma fx=fw$ ($w\in E^n$; $\sigma_i=(w_0,
\dots,w_{\varrho-1})$). Let $y$ and $\gamma^{-1}y$ be boxed. At $\tau$, as $\tilde f_b=s_bf_b$,
\[
\sum_b\frac{\tau(w_b)}{s_b}\,y_\tau^{-1}\tilde f_b=y_\tau^{-1}(fw)=y_\tau^{-1}\gamma fx=\sum_b
\frac{\tau(x_b)}{s_b}\,(\gamma^{-1}y)_\tau^{-1}\tilde f_b .
\]
By Lemma~\ref{pbc:box} at $\gamma^{-1}y$, the right side has norm at most $n\,b^{(1+134n)\varrho}b^{37n}
\max_b|s_b|^{-1}$; by Lemma~\ref{pbc:indep} at $y$, $|\tau(w_a)|/|s_a|\le b^{46n^2}$ times this norm; and
$|s_a/s_b|\le|\kappa_\tau|^\varrho\le b^{23\varrho}$ by \eqref{pbe:eq-kappa}. So $|\tau(w_a)|\le b^{46n^2+1+(1+134n)\varrho+37n+
23\varrho}\le b^{182n^2}$ (with $\varrho=n-1$ the exponent is $180n^2-73n-23$). The entries of $\sigma$ are integral ($\gamma f_b\in L^\sharp$, $\pbcadj Y$ integral),
so by Lemma~\ref{pbc:pack}(a) with $m=\varrho n$ there are at most $(1+2b^{182n^2})^{2\varrho nf}\le
b^{366n^4f}$ values of $\sigma$. Colour $\mathcal C_2$ by $\sigma$ and fix a colour $\mathcal C_3$.

\emph{Step 5: the part in $A$.} Colour $\mathcal C_3$ by the class of $\alpha(\hat\gamma)$ modulo $d'$ (an
integral $\varrho\times n$ matrix): at most $(\auNm d')^{\varrho n}\le b^{2n^4\varrho f}$ colours. We claim
that each colour has at most one element; with Steps~1--4 and \eqref{pbc:eq-moduli} this gives
\[
\#\pbcB\le b^{2}\,b^{270n^3f}\,b^{2n^3\varrho f}\,b^{2n^4lf}\,b^{366n^4f}\,b^{2n^4\varrho f}\le
b^{2000n^6(l+1)f}.
\]
Let $\gamma,\gamma'$ have the same colour and put $\Xi=\alpha(\hat\gamma')-\alpha(\hat\gamma)=d'\Theta$, $\Theta$
integral. The frames $\hat\gamma f$ and $\hat\gamma'f$ differ only in their $A$-parts. (i) \emph{${}^t\aubar
\Theta P$ is skew-hermitian.} Both frames have the Gram matrix of $f$; expanding $h(\hat\gamma'f_b,\hat\gamma'
f_{b'})$ and using that $A$ is isotropic and $h(e_a,\hat\gamma f_{b'})=P_{ab'}$ gives ${}^t\aubar\Xi P+
{}^t\aubar P\Xi=0$; divide by $d'=\aubar{d'}$. (ii) Let $\alpha^\circ$ be the $n\times\varrho$ matrix of the
first $\varrho$ columns of $\pbcadj Y$. \emph{The columns $i\ge j$ and the column $\varrho$ of $\Theta\pbcadj Y$
vanish}: for $i\ge j$, $d'\Theta\pbcadj(Y)\epsilon_i=\sigma_i(\gamma')-\sigma_i(\gamma)=0$. If $j=\varrho+1$ (so
$Y=Z$), the rows $a<\varrho$ of $Y\pbcadj Y=\delta\cdot1$ give $P\pbcadj(Y)\epsilon_\varrho=0$ and $P\alpha^\circ
=\delta\cdot1$; then ${}^t\aubar P\Theta\pbcadj(Y)\epsilon_\varrho=-{}^t\aubar \Theta P\pbcadj(Y)\epsilon_\varrho=0$ and
$\aubar\delta \Theta\pbcadj(Y)\epsilon_\varrho={}^t\aubar{\alpha^\circ}\,{}^t\aubar P\Theta\pbcadj(Y)\epsilon_\varrho=0$.
(iii) Hence $\delta \Theta=\Theta\pbcadj(Y)Y=\Theta\alpha^\circ P$, since only the columns $i<\min(j,\varrho)$ of
$\Theta\pbcadj Y$ survive and the rows $i<j$ of $Y$ are rows of $P$. (iv) Put
\[
S=\aubar\delta\,\Theta\alpha^\circ-\delta\,{}^t\aubar{(\Theta\alpha^\circ)}-{}^t\aubar{\alpha^\circ}\,({}^t\aubar
P\Theta)\,\alpha^\circ .
\]
By (i), ${}^t\aubar S=-S$; $S$ is integral; and by (i), (iii),
\[
SP=\aubar\delta\delta \Theta-\delta\,{}^t\aubar{\alpha^\circ}\,{}^t\aubar\Theta P-{}^t\aubar{\alpha^\circ}\,{}^t
\aubar P\,\delta \Theta=\delta\aubar\delta \Theta-\delta\,{}^t\aubar{\alpha^\circ}({}^t\aubar \Theta P+{}^t\aubar P\Theta)=\delta
\aubar\delta \Theta .
\]
The move $u''\in U_c^0$ of Lemma~\ref{pbc:moves} with $x=0$ and this $S$ changes only the $A$-part, by
$MSP=d'\Theta=\Xi$. So $u''\hat\gamma f=\hat\gamma'f$, that is, $u''u'_\gamma u_\gamma\gamma f=u'_{\gamma'}
u_{\gamma'}\gamma'f$, and Step~0 gives $(c',\gamma)=(c'',\gamma')$. \hfill$\square$

\begin{remark}\label{pbc:rem-C1}
The exponent: $2+270n^3+2n^3\varrho+366n^4+2n^4\varrho+2n^4l\le2000n^6(l+1)$. The theorem is the
statement (C1) of Proposition~\ref{ppi:C}: the elements of $\pbeGs$ have a cell rank with non-vanishing
leading minor by Proposition~\ref{ppi:cell}, and no other property of $\pbeGs$ is used.
\end{remark}

\subsection{The sum over lattices of a given corank: (C2)}\label{pbc:ss-C2}

Let $H=\auU(X_0)$, $K^\circ\subseteq H(\auA_f)$ the stabiliser of $\auO_E^k$, $K_H\subseteq K^\circ$ the group of
Lemma~\ref{pb:KR}, and $L_1=L\cap X_0=\auO_Ex_1+\dots+\auO_Ex_{k-1}+\pi_N\auO_Ex_k$, so that $K_H$ stabilises
$L_1$ and $J_X=K_H\cap J\subseteq K_H$. Recall that $\mathcal C=H(E^+)\backslash H(\auA_f)/J_X$, that $c=
H(E^+)h_cJ_X$, and that the level is free: $H(E^+)\cap hJ_Xh^{-1}=1$ for every $h\in H(\auA_f)$
(\S\ref{pb:sec}). Then $\ppiLc=h_cL_1$. For $g\in H(\auA_f)$ write $[g]\in\mathcal C$ for its class.

\begin{lemma}[Averaging]\label{pbc:average}
Let $\mathcal A$ be a finite nonempty set, $a\colon\mathcal A\to H(\auA_f)$, $F\colon\mathcal C\to[0,
\infty)$ and $B\ge0$. If $\sum_{\phi\in\mathcal A}F([ga(\phi)])\le\#\mathcal A\cdot B$ for every $g\in
H(\auA_f)$, then $\sum_{c\in\mathcal C}F(c)\le\#\mathcal C\cdot B$.
\end{lemma}
\begin{proof}
Let $J'=J_X\cap\bigcap_\phi a(\phi)J_Xa(\phi)^{-1}$, an open compact subgroup, and $\mathcal X=H(E^+)
\backslash H(\auA_f)/J'$, a finite set ($H(E^+)\backslash H(\auA_f)$ is compact, $H$ being anisotropic).
The function $g\mapsto F([ga(\phi)])$ is right $J'$-invariant, so the hypothesis summed over $\mathcal
X$ gives $\sum_\phi\sum_{x\in\mathcal X}F([xa(\phi)])\le\#\mathcal X\#\mathcal A B$. For a fixed $\phi$,
$x\mapsto xa(\phi)$ is a bijection of $\mathcal X$ onto $H(E^+)\backslash H(\auA_f)/J''$, $J''=a(\phi)^{-1}
J'a(\phi)\subseteq J_X$, and the fibre of $H(E^+)\backslash H(\auA_f)/J''\to\mathcal C$ over $c$ is
$(H(E^+)\cap h_cJ_Xh_c^{-1})\backslash h_cJ_X/J''\cong J_X/J''$, of order $\pbvol(J_X)/\pbvol(J'')=
\pbvol(J_X)/\pbvol(J')=:\iota$ (the level is free, and $H(\auA_f)$ is unimodular). So $\sum_{x\in\mathcal
X}F([xa(\phi)])=\iota\sum_cF(c)$; likewise $\#\mathcal X=\iota\#\mathcal C$. Dividing by $\iota\#\mathcal A$
gives the claim.
\end{proof}

\begin{lemma}[Cosets]\label{pbc:index}
Let $\mathcal A=K^\circ/K_H$. For $y\in K^\circ$ let $\lambda_y$ be the last row of $y^{-1}$. Then
$yK_H\mapsto\lambda_y\bmod\pi_N$ is well defined and injective, $yL_1=\{v\in\auO_E^k:\lambda_yv\in\pi_N
\widehat\auO\}$, and $\#\mathcal A\ge N^{kf}/2$ if $4f\le5^k$.
\end{lemma}
\begin{proof}
By Lemma~\ref{pb:KR}, $y^{-1}y'\in K_H$ if and only if $\lambda_y y'\equiv(0,\dots,0,1)=\lambda_{y'}y'$
modulo $\pi_N$, that is, $\lambda_y\equiv\lambda_{y'}$. The description of $yL_1$ is that of $L_1$
transported by $y$. As in the proof of Lemma~\ref{pb:KR}, at each prime $w\mid\pi_N$ the group $H$ is
$\auGL_k(E_w)$, with $K^\circ_w=\auGL_k(\auO_w)$, and the congruence modulo $\pi_N$ is the congruence
modulo every $w\mid\pi_N$; as $K^\circ$ is a product of local groups, the image contains all rows that are
nonzero modulo every $w$. With $q_w=\auNm w\ge N\ge5$, $\prod_wq_w=\auNm(\pi_N)=N^f$ and at most $2f$
primes $w$, $\#\mathcal A\ge\prod_w(q_w^k-1)\ge N^{kf}(1-2f\cdot5^{-k})\ge N^{kf}/2$.
\end{proof}

\begin{lemma}[Forms vanishing on a lattice]\label{pbc:forms}
Let $\Lambda$ be a unimodular lattice of $X_0$, $L''\subseteq\Lambda$ a lattice of rank $r=k-l$, $V=EL''$
and $i=[\Lambda\cap V:L'']$. Then at most $i\,N^{fl}$ classes $\lambda\in\Lambda^\vee/\pi_N\Lambda^\vee$
satisfy $\lambda(L'')\subseteq\pi_N\auO_E$.
\end{lemma}
\begin{proof}
$\Lambda_V=\Lambda\cap V$ is a direct summand: $\Lambda=\Lambda_V\oplus\Lambda_2$, $\Lambda_2$ of rank $l$, and
$\Lambda^\vee/\pi_N\Lambda^\vee=\Lambda_V^\vee/\pi_N\Lambda_V^\vee\times\Lambda_2^\vee/\pi_N\Lambda_2^\vee$. The
second factor has $\auNm(\pi_N)^l=N^{fl}$ elements. In the first, the admissible $\lambda_1$ lie in $\pi_N
L''^\vee\cap\Lambda_V^\vee$ (restriction to $L''$), which contains $\pi_N\Lambda_V^\vee$ with index at most
$[\pi_NL''^\vee:\pi_N\Lambda_V^\vee]=[L''^\vee:\Lambda_V^\vee]=i$ by Lemma~\ref{pbc:disc}(b).
\end{proof}

\begin{lemma}[Sublattices]\label{pbc:zeta}
Let $\Lambda'$ be a lattice of rank $\rho$ and $s>\rho$. Then $\sum_{M'}[\Lambda':M']^{-s}=\prod_{i=0}^{\rho-1}
\zeta_E(s-i)$, the sum over the sublattices $M'$ of finite index. If $s\ge\rho+2\varrho+4$ and $2f\le
2^{2\varrho}$, the sum is at most $2^\rho$.
\end{lemma}
\begin{proof}
A sublattice of finite index is the same as a family $(M'_{\mathfrak p})$ of sublattices of finite index
of the completions, almost all equal to $\Lambda'_{\mathfrak p}$, and the index is the product of the
local indices; so the sum is the product over the primes $\mathfrak p$ of $\auO_E$ of the local sums.
Let $\varpi$ be a uniformiser and $q=\auNm\mathfrak p$; $\Lambda'_{\mathfrak p}\cong\auO_{\mathfrak p}^\rho$.
Each $M'_{\mathfrak p}$ has a unique basis matrix in Hermite form: upper triangular with diagonal
$\varpi^{a_1},\dots,\varpi^{a_\rho}$ and with the entries to the right of the diagonal in row $i$ taken in
a fixed set of representatives of $\auO_{\mathfrak p}/\varpi^{a_i}$. There are $\prod_iq^{a_i(\rho-i)}$ such
matrices for given $a_i$, and the index is $q^{\sum a_i}$; so the local sum is $\prod_{i=1}^\rho(1-q^{
\rho-i-s})^{-1}$, and the product over $\mathfrak p$ is $\prod_{i=0}^{\rho-1}\zeta_E(s-i)$. For the bound,
each factor is at most $\zeta_E(\sigma)$ with $\sigma=2\varrho+5$. Above each rational prime $p$ there are
at most $2f$ primes, of norm at least $p$, so $\zeta_E(\sigma)\le\zeta(\sigma)^{2f}$; and $\zeta(\sigma)\le
1+2^{-\sigma}+\int_2^\infty x^{-\sigma}dx\le1+2^{1-\sigma}$. Hence $\zeta_E(\sigma)\le\exp(2f\,2^{1-\sigma})
=\exp(2f\,2^{-2\varrho-4})\le\exp(1/16)\le2$.
\end{proof}

For $x\in X_0$ put $\|x\|^2=\sum_{\tau}\tau(h(x,x))$, the Euclidean norm of the image of $x$ in $X_{0,
\infty}=\prod_\tau X_0\otimes_{E,\tau}\C\cong\C^{kf}\cong\R^{2kf}$ (coordinates in the orthonormal
basis $x_1,\dots,x_k$). If $x\ne0$ lies in an integral lattice, $h(x,x)$ is a totally positive integer of
$E^+$, so $\|x\|^2\ge f\,(N_{E^+/\Q}h(x,x))^{1/f}\ge1$. Let $D_E$ be the discriminant of $E$ and
$\pbcrd=|D_E|^{1/2f}$ its root discriminant.

\begin{lemma}[Covolume of $\auO_E$]\label{pbc:tb-covol}
Let $\auPl$ be a set of representatives of the $f$ pairs of complex conjugate embeddings of the CM
field $E$ ($[E:\Q]=2f$). The image of $\auO_E$ under $a\mapsto(\tau(a))_{\tau\in\auPl}\in\C^f\cong\R^{2f}$
is a lattice of covolume $2^{-f}|D_E|^{1/2}$ for the Lebesgue measure $\prod_\tau d\Rel\,d\Iml$.
\end{lemma}
\begin{proof}
Let $\alpha_1,\dots,\alpha_{2f}$ be a $\Z$-basis of $\auO_E$ and $A$ the $2f\times2f$ matrix with rows
$(\sigma(\alpha_i))_\sigma$, $\sigma$ running over the $2f$ embeddings, ordered as $\tau_1,\aubar\tau_1,\dots,
\tau_f,\aubar\tau_f$; by definition $D_E=\det(A)^2$. The image lattice has the basis whose $i$-th vector has
coordinates $(\Rel\tau(\alpha_i),\Iml\tau(\alpha_i))_\tau$, so its covolume is $|\det B|$ for the matrix
$B$ of these rows. $B$ is obtained from $A$ by replacing, for each $\tau$, the pair of columns
$(\tau,\aubar\tau)$ by $(\tfrac12(\tau+\aubar\tau),\tfrac1{2i}(\tau-\aubar\tau))$, i.e.\ $B=AM$ with $M$
block diagonal with blocks $\bigl(\begin{smallmatrix}1/2&1/(2i)\\1/2&-1/(2i)\end{smallmatrix}\bigr)$
of determinant $-1/(2i)$. Hence $|\det B|=2^{-f}|\det A|=2^{-f}|D_E|^{1/2}$; in particular $\det B\ne0$
and the image is a lattice. (See also \cite[Chapter~I, (5.1) and Proposition~(5.2)]{Neukirch}.)
\end{proof}

\begin{lemma}[Covolume]\label{pbc:covol}
Let $\Lambda'$ be an integral lattice ($h(\Lambda',\Lambda')\subseteq\auO_E$) of rank $\rho$, spanning a subspace $W\subseteq X_0$. Its image in $W_\infty=\prod_\tau
W\otimes_{E,\tau}\C\cong\R^{2\rho f}$ (with the norm above) is a lattice of covolume $2^{-\rho f}|D_E|^{\rho/2}
\ppidisc\Lambda'$.
\end{lemma}
\begin{proof}
For $\auO_E\subset\C^f$ (embedded by the $\tau\in\auPl$) the covolume for the Lebesgue measure is
$2^{-f}|D_E|^{1/2}$ (Lemma~\ref{pbc:tb-covol}). Let $w_1,\dots,w_\rho\in\Lambda'$ be an $E$-basis of $W$ with Gram
matrix $G_w$ and $F=\bigoplus\auO_Ew_i\subseteq\Lambda'$, an integral lattice of finite index in $\Lambda'$. The map $\C^{\rho f}\to W_\infty$ given at $\tau$ by $\tau(w)$ (in
orthonormal coordinates) has real Jacobian $\prod_\tau|\det\tau(w)|^2=\prod_\tau\tau(\det G_w)=\auNm(\det
G_w)^{1/2}=\ppidisc F$ (Lemma~\ref{pbc:disc}(a); $\det G_w\in E^+$ is totally positive). So the claim holds
for $F$. Since $F\subseteq\Lambda'$ has finite index, $\operatorname{covol}(F)=[\Lambda':F]\operatorname{covol}(\Lambda')$
and, by Lemma~\ref{pbc:disc}(a), $\ppidisc F=[\Lambda':F]\ppidisc\Lambda'$; so the claim holds for $\Lambda'$.
\end{proof}

\begin{lemma}[Siegel's lemma]\label{pbc:siegel}
Let $\Lambda$ be a unimodular lattice of $X_0$ and $V\subseteq X_0$ a subspace of codimension $l\ge1$. There
are $E$-linearly independent $X_1,\dots,X_l\in\Lambda\cap V^\perp$ with $V=\{v:h(X_i,v)=0\ \forall i\}$ and
\[
\prod_{i=1}^l\|X_i\|^{2f}\le(lf\,\pbcrd)^{lf}\,\ppidisc(\Lambda\cap V).
\]
\end{lemma}
\begin{proof}
Let $\Lambda'=\Lambda\cap V^\perp$, of rank $l$ and integral as $\Lambda$ is, viewed as a lattice of rank $D=2lf$ over $\Z$ in $V^\perp_\infty
\cong\R^D$. By Lemmas~\ref{pbc:covol} and~\ref{pbc:disc}(c) its covolume is $2^{-lf}|D_E|^{l/2}\ppidisc
(\Lambda\cap V)$. Let $\lambda_1\le\dots\le\lambda_D$ be its successive minima for the Euclidean unit ball, of
volume $\pi^{lf}/(lf)!$. Minkowski's second theorem (\cite[Chapter~VIII, \S4.3, Theorem~V]{Cassels}; see also \cite[Theorem~1.3]{HenkMinima}) gives
\[
\lambda_1\cdots\lambda_D\le\frac{2^D\,2^{-lf}|D_E|^{l/2}\ppidisc(\Lambda\cap V)}{\pi^{lf}/(lf)!}\le(lf)^{lf}
\pbcrd^{\,lf}\ppidisc(\Lambda\cap V),
\]
as $|D_E|^{1/2}=\pbcrd^{\,f}$ and $(2/\pi)^{lf}\le1$. Let $v_1,\dots,v_D$ be $\Z$-independent with $\|v_j\|=
\lambda_j$. An $E$-subspace of dimension $e$ contains at most $2fe$ $\Z$-independent vectors, so
$v_1,\dots,v_{2f(i-1)+1}$ span an $E$-subspace of dimension at least $i$; choose inductively
$E$-independent $X_i$ among them with $\|X_i\|\le\lambda_{2f(i-1)+1}$. Then $\|X_i\|^{2f}\le\prod_{t=1}^{2f}
\lambda_{2f(i-1)+t}$, and the product over $i$ is at most $\lambda_1\cdots\lambda_D$. The $X_i$ span
$V^\perp$, so their orthogonal is $V$.
\end{proof}

\begin{lemma}[Height zeta function]\label{pbc:height}
Let $\Lambda$ be a unimodular lattice of $X_0$ and $0\le l\le k$. Then $\sum_V\ppidisc(\Lambda\cap V)^{-(k+1)}
\le(25\,lf\,\pbcrd)^{f(k+1)l}$ (read as $1$ for $l=0$), the sum over the subspaces $V\subseteq X_0$ of
codimension $l$.
\end{lemma}
\begin{proof}
For $l=0$ the only term is $V=X_0$, with $\ppidisc\Lambda=1$. Let $l\ge1$. Attach to each $V$ the tuple
$(X_1,\dots,X_l)$ of Lemma~\ref{pbc:siegel}; $V$ is determined by it. Raising the inequality of that
lemma to the power $k+1$,
\[
\ppidisc(\Lambda\cap V)^{-(k+1)}\le(lf\pbcrd)^{lf(k+1)}\prod_i\|X_i\|^{-2f(k+1)} .
\]
Summing over all tuples of nonzero vectors of $\Lambda$, the total is at most $(lf\pbcrd)^{lf(k+1)}
\bigl(\sum_{0\ne X\in\Lambda}\|X\|^{-2f(k+1)}\bigr)^l$. The image of $\Lambda\smallsetminus0$ in $\R^{2kf}$
satisfies the hypotheses of Lemma~\ref{pbc:pack}(b) (differences of lattice vectors are lattice vectors),
and $2f(k+1)\ge2kf+2$; so the inner sum is at most $5^{2f(k+1)}$.
\end{proof}

\begin{theorem}[(C2)]\label{pbc:C2}
Assume the size conditions of Remark~\ref{pbc:rem-sizes}. Let $0\le l\le n^3$ and $r=k-l$. Then
\[
\sum_{c\in\mathcal C}\ \sum_{L'}\ (\ppidisc L')^{-k}\le\frac{\#\mathcal C}{N^{kf}}\,\pbet^{\,l+1},
\]
the inner sum over the lattices $L'\subseteq\ppiLc$ of rank $r$ with $\ppidisc L'\le b^{n^3f}$.
\end{theorem}
\begin{proof}
Put $\mathsf D=b^{n^3f}$ and let $F(c)$ be the inner sum. Let $\mathcal A=K^\circ/K_H$ and fix
representatives $y_\phi\in K^\circ$ of the cosets $\phi$. By Lemma~\ref{pbc:average} it suffices to show,
for every $g\in H(\auA_f)$,
\begin{equation}\label{pbc:eq-goal}
\sum_{\phi\in\mathcal A}F([gy_\phi])\le\frac{\pbet^{\,l+1}}2\le\#\mathcal A\cdot\frac{\pbet^{\,l+1}}{N^{kf}},
\end{equation}
the second inequality being Lemma~\ref{pbc:index} ($4f\le2^{2\varrho+1}\le5^k$). Fix $g$ and put $\Lambda=g
\auO_E^k$, a unimodular lattice of $X_0$. For $\phi\in\mathcal A$ let $c=[gy_\phi]$; then $h_c=\eta g y_\phi j$
with $\eta\in H(E^+)$ and $j\in J_X\subseteq K_H$, so $\ppiLc=h_cL_1=\eta gy_\phi L_1$. By
Lemma~\ref{pbc:index}, $\eta^{-1}\ppiLc=\{v\in\Lambda:\lambda_\phi(v)\in\pi_N\widehat\auO\}$, where
$\lambda_\phi=\lambda_{y_\phi}\circ g^{-1}\in\Lambda^\vee$, and the classes of the $\lambda_\phi$ modulo $\pi_N
\Lambda^\vee$ are distinct. As $L'\mapsto\eta^{-1}L'$ preserves rank and discriminant
(Lemma~\ref{pbc:disc}(a)),
\[
\sum_\phi F([gy_\phi])\le\sum_{L''}\#\{\phi:\lambda_\phi(L'')\subseteq\pi_N\auO_E\}\,(\ppidisc L'')^{-k},
\]
the sum over the lattices $L''\subseteq\Lambda$ of rank $r$ with $\ppidisc L''\le\mathsf D$. For such $L''$ let
$V=EL''$, $d_V=\ppidisc(\Lambda\cap V)\ge1$ and $i=[\Lambda\cap V:L'']\ge1$, so that $\ppidisc L''=d_Vi\le\mathsf
D$. By Lemma~\ref{pbc:forms} the count is at most $iN^{fl}$. With $\mu=2\varrho+4$,
\[
i\,(d_Vi)^{-k}=i(d_Vi)^{\mu+1}d_V^{-(k+\mu+1)}i^{-(k+\mu+1)}\le\mathsf D^{\mu+2}\,d_V^{-(k+1)}\,i^{-(r+\mu)},
\]
using $i\le\mathsf D$, $d_Vi\le\mathsf D$, $d_V,i\ge1$ and $r\le k$. Group the $L''$ by $V$ (a subspace of
codimension $l$) and sum over the sublattices of finite index of $\Lambda\cap V$ by Lemma~\ref{pbc:zeta}
($2f\le2^{2\varrho}$), then over $V$ by Lemma~\ref{pbc:height}:
\[
\sum_\phi F([gy_\phi])\le N^{fl}\mathsf D^{2\varrho+6}\,2^r\,(25\,lf\,\pbcrd)^{f(k+1)l}
\le\bigl(2^{k}\mathsf D^{2\varrho+6}\bigr)\bigl(N^f(25n^3f\pbcrd)^{f(k+1)}\bigr)^l .
\]
By the size conditions, $2^{k+1}\mathsf D^{2\varrho+6}\le b^{(k+1)/8+n^3(2n+4)f}\le\pbet$ and $N^f(25n^3f
\pbcrd)^{f(k+1)}\le b^{60nf+f(k+1)/1000}\le\pbet$. This gives \eqref{pbc:eq-goal}.
\end{proof}

\begin{remark}[Size conditions]\label{pbc:rem-sizes}
Theorem~\ref{pbc:C2} uses: (Z1) $\log b\ge8\log2$, $N\ge5$; (Z2) $k\ge8n^3(2n+4)+1$, $k\ge2\varrho$,
and $60n+(k+1)/1000\le k/4$; (Z3) $N\le b^{60n}$; (Z4) $2f\le2^{2\varrho}$; (Z5) $(25n^3f\pbcrd)^{1000}\le
b$; and the description of $K_H$ at the primes $w\mid\pi_N$ in the proof of Lemma~\ref{pb:KR}.
Theorem~\ref{pbc:C1} uses $\varrho\ge280$, $M\le b^{23}$, $N\le b^{60n}$, $1\le|\kappa_\tau|\le b^{23}$, $n\le b$.
All follow from Lemma~\ref{an:lem-sizes}: (Z1) from (S1) and $\log N\ge m\log b$ (S5); (Z2), (Z3),
$\varrho\ge280$ from (S7) ($k\ge512000n^9$, $N\le b^{60n}$, $\varrho\ge10^5$); $M\le b^{22}$ from (S5);
the bound on $\kappa_\tau$ is \eqref{pbe:eq-kappa}; $n\le b$ from $18\varrho f\le b^8$ (S7) or (S1). For (Z4) and (Z5): by the
proof of (S2), $\sqrt x\le m\le2\sqrt x$, so $n=m/2\le\sqrt x$ and $\varrho\ge\sqrt x/2-1$; by (S3),
(S4), $f\le x^{12}$ and $\pbcrd\le x^{12}$. Since $\log x\ge10^{14}$ (S1), $2^{2\varrho}\ge2^{\sqrt x-2}\ge2x^{12}
\ge2f$, and $25n^3f\pbcrd\le x^{26}$, while $x^{26000}\le x^{10^6}\le x_1\le b$ (S1). (Here $f$ is the
number of places $\tau$, as in the rest of the paper.)
\end{remark}


%% ===== proofs/B-size.tex =====
% proofs/B-size.tex
% Statement (S) (size) of Theorem ppi:pillars of proofs/B-pillars.tex: the bound of one block of
% amplitudes by the sizes of its near orbits. Labels of this file start with pbs:.
\providecommand{\auO}{\mathcal O}
\providecommand{\auNm}{\operatorname{Nm}}
\providecommand{\auA}{\mathbb A}
\providecommand{\aubar}[1]{\overline{#1}}
\providecommand{\auf}[1]{\texttt{#1}}
\providecommand{\panot}[1]{\par\smallskip{\raggedright\noindent\textbf{[Not proved here:} #1\textbf{]}\par}\smallskip}
\providecommand{\pbR}{\overline R_0}
\providecommand{\pbDid}{\mathcal D_{\mathrm{id}}}
\providecommand{\pbDmain}{\mathcal D_{\mathrm{main}}}
\providecommand{\pbvol}{\operatorname{vol}}
\providecommand{\pbdeg}{\operatorname{deg}}
\providecommand{\pbg}{\mathsf g}
\providecommand{\pbeN}{\mathcal N}
\providecommand{\pbet}{\beta_*}
\providecommand{\pbenp}{\mathsf n_{\mathrm p}}
\providecommand{\pbesz}{\operatorname{size}}
\providecommand{\pbeGs}{\Gamma^*}
\providecommand{\pbeAp}{A^{+}}
\providecommand{\pbel}{\ell}
\providecommand{\pbeL}{\mathsf L}
% macros of this file (prefix \pbs)
\providecommand{\pbsv}{\mathsf v}
\providecommand{\pbsK}{\mathsf K}
\providecommand{\pbsF}{\mathcal F}
\providecommand{\pbsO}{\mathcal O}
\providecommand{\pbsrb}{\rho_\square}
\providecommand{\pbsth}{\theta_\square}
\providecommand{\pbsp}{\mathsf p}
\providecommand{\pbsi}{\mathsf i}
\providecommand{\pbsmn}{\mathsf m}

\section{The size of a block}\label{pbs:sec}

This section proves statement \textup{(S)} of Theorem~\ref{ppi:pillars}. We use the notation of
\S\S\ref{pb:sec}, \ref{pbe:sec} and~\ref{ppi:sec}: $n=\varrho+1$, the frame $f_0,\dots,f_\varrho$
($f_a=e_{a+1}$ for $a<\varrho$, $f_\varrho=z$), $A=\langle f_0,\dots,f_{\varrho-1}\rangle$, the matrix
$Z(\gamma)=(h(f_a,\gamma f_b))_{a,b}$ with leading minors $\delta_l(\gamma)$, the discriminant
$\pbg(\gamma)=\delta_n\aubar{\delta_n}-N^2\delta_\varrho\aubar{\delta_\varrho}$, the groups $U_c$, $U^0$,
$U_c^0=U_c\cap U^0$, the cell rank $j(\gamma)$ (Definition~\ref{ppi:def-cell}), the gauge $\mathsf
G(\gamma)$ \eqref{ppi:eq-gauge}, near and chosen elements, agreeing reductions, the set $\pbeN(c,c')$
(Definition~\ref{ppi:def-N}) and $\pbesz(\gamma)=\pbet^2(\mathcal D(\gamma)/(\pbenp\mathsf
G(\gamma)^k))^{1/2}$. At a real place $\tau$ we use the normalised basis, the norm $\|\cdot\|$, the
scales $s_a$ \eqref{pbe:eq-scale}, the length $\pbel_\tau$ and the cut of \S\ref{pbe:ss-real}.
In this section, as in \S\S\ref{pbg:sec}, \ref{ppi:ss-cell} and \ref{pbc:sec}, hermitian forms are
written conjugate-linear in the first variable and linear in the second, the opposite of convention
(Cv1); $Z(\gamma)$ is the matrix of \S\ref{ppi:ss-statement} in this reading.
Put
\[
\pbsrb=b^{37n},\qquad\pbsth=b^{6n+2},\qquad\mathsf j'=\min(\varrho,j),
\]
the radius of the boxes and the radius of the cut.

\emph{What is proved.}
\begin{itemize}
\item[(1)] The exact statement is Theorem~\ref{pbs:S}. Its right-hand side is that of
Theorem~\ref{ppi:pillars}(S). The constant in front of the sum is $\nu=\pbvol(U_c\backslash
U_\infty)$ for the Haar measure $du$ of convention (Cv3) of \S\ref{pb:ss-setting}; its explicit value
$\nu_{\mathrm{src}}=\nu_*$ is Lemma~\ref{pb:covol}, used in Lemma~\ref{pbs:open-arith}
(Remark~\ref{pbs:rem-nu}).
\item[(2)] The reduction of (S) to one orbit of one cell rank, the vanishing for non-agreeing $s$,
and the choice of representatives are \textbf{proved in full} (\S\ref{pbs:ss-reduction}).
\item[(3)] The bound for the orbits that are summed by Schur's test (all orbits except those on the
open cell with large discriminant) is \textbf{proved in full} (\S\S\ref{pbs:ss-schur}--%
\ref{pbs:ss-orbit}): the weighted Cauchy--Schwarz inequality, the computation of the weight at one
real place by a triangular basis, the Ostrowski--Taussky inequality and Hua's integral, the
colouring of an orbit by a lattice index, packing, and the final comparison of sizes.
\item[(4)] For the orbits on the open cell the bound is deduced (\S\ref{pbs:ss-open}) from two
statements: the open-cell pairing bound (O1) for the Bessel functions, stated exactly and
proved in \S\ref{pbo:sec}, and the comparison (O2) of the Siegel constant with the main second
moment, \textbf{proved here} (Lemma~\ref{pbs:O2}) from the mass formula of
Lemma~\ref{pb:massratio}.
\item[(5)] The size conditions on the data that are used are listed in \S\ref{pbs:ss-sizes};
all are in Lemma~\ref{an:lem-sizes}.
\end{itemize}

\subsection{The statement}\label{pbs:ss-statement}

For a class $c$ let $\Lambda_c=h_cL\cap X_0$, a lattice of the positive definite space $X_0$
($\dim X_0=k$). For $\gamma\in\Gamma_{c,s,c'}$ let $k_\gamma=h_c^{-1}\gamma h_{c'}$, and for $0\le
j\le n$ let $i_\gamma=\pbsi_j(k_\gamma)$, $i'_\gamma=\pbsi_j(k_\gamma^{-1})$ be the indices of
\eqref{ppi:eq-gauge}: $\pbsi_j(k)$ is the index, in the $\widehat\auO$-module spanned by the
vectors $\xi_a(k)=(h(f_a,kf_b))_{0\le b\le\varrho}\in\auA_{E,f}^{\,n}$, $0\le a<\varrho$, of the
submodule spanned by the $\xi_a(k)$ with $a<\mathsf j'$. (These modules are finitely generated
$\widehat\auO$-modules of the same rank by Lemma~\ref{pbs:indep} below, so the index is finite; it
is $1$ for $j=n$.) For a cell rank $j$ put
\[
\pbsmn_j(\gamma)=\frac{|\auNm\delta_j(\gamma)|}{N^{2(j-\varrho)_+f}},\qquad
\mathsf G_j(\gamma)=\frac{\max\bigl(\auNm\pbg(\gamma),\ N^f\pbsmn_j(\gamma)\bigr)}{i_\gamma
i'_\gamma},
\]
where $\auNm x=\prod_{\tau\in\auPl}|\tau(x)|$ for $x\in E^+$ and $\auNm x=\prod_\tau|\tau(x)|^2$ for
$x\in E$. Then the gauge \eqref{ppi:eq-gauge} is $\mathsf G(\gamma)=M^{2\mathsf j'f}\mathsf
G_{j(\gamma)}(\gamma)$, and for $\gamma$ of cell rank $j$
\[
\pbesz(\gamma)=\mathsf S_j\bigl(\mathsf G_j(\gamma)\bigr),\qquad\mathsf S_j(x):=\pbet^2\Bigl(
\frac{\mathcal D_j}{\pbenp\,(M^{2\mathsf j'f}x)^k}\Bigr)^{1/2},\qquad\mathcal D_j=\begin{cases}
\pbDmain&j=n,\\ \pbet^{-16}\pbDid&j\le\varrho,\end{cases}
\]
since $\mathcal D(\gamma)=\pbDmain$ exactly when $\pbg(\gamma)\ne0$, that is when $j(\gamma)=n$
(Proposition~\ref{ppi:cell}).

\begin{theorem}[Size]\label{pbs:S}
Assume the size conditions of \S\ref{pbs:ss-sizes}, Theorem~\ref{pbe:ds}, and the statement
\textup{(O1)} of \S\ref{pbs:ss-open} (Theorem~\ref{pbo:O1}). Let $W^\flat$ be as in
Theorem~\ref{pbe:bessel}. For every stable $T\subseteq\pbeGs$ and every block $(c,s,c')$,
\[
\bigl\|A^\flat_T(c,s,c')\bigr\|\le\nu\sum_{\gamma\in\pbeN(c,c')\cap T\cap\Gamma_{c,s,c'}}
\pbesz(\gamma),
\]
and $A^\flat_T(c,s,c')=0$ if $s$ does not agree. This is \textup{(S)} of
Theorem~\ref{ppi:pillars}.
\end{theorem}

The proof occupies the rest of the section; it is assembled in \S\ref{pbs:ss-proof}.

\subsection{Measures, cutoffs and reduced points}\label{pbs:ss-measures}

\subsubsection*{Measures} On $G_\infty=\prod_\tau G_\tau$ we use the product $dg$ of the Haar
measures of \S\ref{pbe:ss-real}. At each $\tau$ the group $U_\tau$ is the group of real points of a
split unipotent group; on it $du_\tau$ is the Haar measure of convention (Cv3) of
\S\ref{pb:ss-setting}, on $U_\infty$ the product $du$, and on $U_\tau\backslash G_\tau$ the quotient
measure (Remark~\ref{pbs:rem-nu}). On the subgroup $U^0_\tau$ of the elements of $U_\tau$ that fix
every $\tilde f_a$, also a unipotent group, we fix a Haar measure $du^0$. Unipotent groups are
unimodular. Recall $\nu=\pbvol(U_c\backslash U_\infty)$, finite and independent of $c$
(Lemma~\ref{pb:ortho}).

A \emph{cutoff} for a discrete group $\Delta$ acting on the left on $G_\infty$ is a measurable
$p\colon G_\infty\to[0,\infty)$ with $\sum_{\delta\in\Delta}p(\delta y)=1$ for all $y$. For a
discrete subgroup of $G_\infty$ the indicator of a measurable set of representatives of
$\Delta\backslash G_\infty$ is one. If $p$ is a cutoff for $\Delta$ and $F\ge0$ is measurable and
left invariant under $\Delta$, then $\int_{G_\infty}pF\,dg=\int_{\Delta\backslash G_\infty}F$; in
particular the left side does not depend on the cutoff.

\subsubsection*{Reduced points}
Let $\mathsf D=\pbcD\subseteq\prod_\tau\C$ be the Borel set of Lemma~\ref{pbc:domain}: it contains
exactly one representative of each coset of the lattice $M\auO_E$ (embedded by $x\mapsto(\tau(x))_\tau$),
and all coordinates of its points have absolute value at most $b^{26}$. For $y\in G_\infty$ and each $\tau$ let $w_a=y_\tau^{-1}\tilde f_a$,
let $o_a$ be their Gram--Schmidt orthogonalisation and $w_a=o_a+\sum_{a'<a}t_{a'a}o_{a'}$ as in
Lemma~\ref{pbe:len}(c). The point $y$ is \emph{reduced} if for all $a'<a$ the point
$(t_{a'a}(y_\tau)s_{a'}/s_a)_\tau$ lies in $\mathsf D$; let $\omega$ be the indicator of the set of
reduced points. A point is \emph{boxed} if it is reduced and $\pbel_\tau(y_\tau)\le\pbsth$ for all
$\tau$. (These are the reduced and boxed points of Definition~\ref{pbc:def-reduced}.)

\begin{lemma}[Reduced points]\label{pbs:reduced}
\begin{enumerate}
\item[(a)] $\omega(uy)=\omega(y)$ for every $u=(u_\tau)$ with $u_\tau$ in the stabiliser $F_\tau$
of $\tilde f_0,\dots,\tilde f_\varrho$ in $G_\tau$.
\item[(b)] For every $y$, the set of $x\in U_c$ with $xy$ reduced is a single coset $U_c^0x_0$.
Hence $\sum_{x\in U_c}a(x)\omega(xy)=1$ for every function $a\colon U_c\to[0,\infty]$ whose sum
over each coset $U^0_cx$ is $1$.
\item[(c)] If $y$ is reduced and $\pbel_\tau(y_\tau)\le\pbsth R$ with $R\ge1$, then
$\|y_\tau^{-1}\tilde f_a\|\le\pbsrb R$ for all $a$.
\end{enumerate}
\end{lemma}
\begin{proof}
(a) $(u_\tau y_\tau)^{-1}\tilde f_a=y_\tau^{-1}\tilde f_a$. (b) For $x\in U_c$ write $xf_a=f_a+\sum_{a'
<a}C_{a'a}f_{a'}$; then $C_{a'a}\in M\auO_E$, because $U_c$ preserves $h_cL$, acts trivially on it
modulo $M$, and the $f_a$ are part of a basis of $h_cL$ adapted to the frame (proof of
Lemma~\ref{pbe:eichler}). As in Step~1 of the proof of Proposition~\ref{pbe:count}, replacing $y$
by $xy$ does not change the $o_a$ and replaces $t_{a'a}s_{a'}/s_a$ by $t_{a'a}s_{a'}/s_a+\tau(
C_{a'a})+\sum_{a'<a''<a}t_{a'a''}(s_{a'}/s_{a''})\tau(C_{a''a})$ (for $x^{-1}$ in place of $x$, which
is the same statement). By Lemma~\ref{pbe:eichler} every strictly upper triangular matrix over
$M\auO_E$ occurs as $(C_{a'a})$; choosing its entries by induction on $a-a'$ we reach a reduced
point, and at each step the choice is unique because $\mathsf D$ is a set of representatives. So
the $x$ with $xy$ reduced are those with one prescribed matrix $(C_{a'a})$, and two elements of
$U_c$ with the same matrix differ by an element of $U_c$ fixing the frame, that is of $U_c^0$. For
the second claim, the sum is the sum of $a$ over the coset $U^0_cx_0$. (c) The coefficients of a
reduced point satisfy $|t_{a'a}|\le|t_{a'a}s_{a'}/s_a|\le b^{26}$ by \eqref{pbe:eq-scale}, and
$\|o_a\|<\pbel_\tau(y_\tau)$ by Lemma~\ref{pbe:len}(c), so $\|w_a\|\le n\,b^{26}\pbsth R\le\pbsrb R$.
\end{proof}

\subsubsection*{Class cutoffs} For each class $c$ fix a cutoff $p_c$ for $U_c^0$ and put
$\beta_c=p_c\,\omega$. By Lemma~\ref{pbs:reduced}(b) applied to $a(x)=p_c(xy)$, $\beta_c$ is a
cutoff for $U_c$:
\begin{equation}\label{pbs:eq-cutoff}
\sum_{x\in U_c}\beta_c(xy)=1\qquad(y\in G_\infty).
\end{equation}
Hence, for $F\ge0$ measurable and left invariant under $U_\infty$,
\begin{equation}\label{pbs:eq-nu}
\int_{G_\infty}\beta_c\,F\,dg=\int_{U_c\backslash G_\infty}F=\nu\int_{U_\infty\backslash G_\infty}F .
\end{equation}

\begin{remark}[The constant $\nu$]\label{pbs:rem-nu}
Put $\nu_{\mathrm{src}}=(M^{2d+\varrho}/N^k)^f\pbvol(\widehat\auO)^{-d}|D_E|^{\varrho/4}/\Lambda$
with $d=\sum_{a<\varrho}(2a+k+1)=\varrho(\varrho+k)$ and $\Lambda=\prod_\tau\prod_{a\le\varrho}
(\kappa_\tau^{2(\varrho-a)}/N)^{k+2(\varrho-a)}$; this is the constant $\nu_*$ of Lemma~\ref{pb:covol}
($M^{f(2d+\varrho)}=\#\pbUb$, $|D_E|^{\varrho/4}=\aurd^{f\varrho/2}$), so $\nu=\nu_{\mathrm{src}}$ for the
measure (Cv3) by that lemma. The statement of Theorem~\ref{pbs:S} itself does not depend on the
choice of $du$: if $du$ is multiplied by $\lambda>0$, the quotient measure on $U_\infty\backslash G_\infty$ is
divided by $\lambda$, the normalisation $\int_{U_\tau\backslash G_\tau}\|W_\tau\|^2=d_\tau$ multiplies
$W$ and $W^\flat$ by $\lambda^{1/2}$ (one factor $\lambda^{1/2f}$ at each place), so $A^\flat_T$ is
multiplied by $\lambda$, and so is $\nu$, while $\pbesz(\gamma)$ does not change. Its proof, however,
uses (O1), which is proved for the measure (Cv3), and the explicit value of $\nu$
(Lemmas~\ref{pbs:O2}, \ref{pbs:open-arith}).
\end{remark}

\subsubsection*{Amplitudes with weights}
For $\gamma\in\Gamma$ and $y\in G_\infty$ put $\pbsK^\flat_\gamma(y)=W^\flat(y)W^\flat(\gamma^{-1}y)
^*$ and $\pbsK_\gamma(y)=W(y)W(\gamma^{-1}y)^*$, matrices indexed by $I\times I$ with the Euclidean
norm; $\|\pbsK^\flat_\gamma(y)\|=\|W^\flat(y)\|\,\|W^\flat(\gamma^{-1}y)\|$. For $u\in U_c$, $u'\in
U_{c'}$, $\pbsK^\flat_{u\gamma u'^{-1}}(uy)=\pbsK^\flat_\gamma(y)$, because $W^\flat(uy)=\psi_\infty(u)
W^\flat(y)$ and $\psi_\infty$ is trivial on $U_c$ and on $U_{c'}$ (\S\ref{pb:ss-layer2}).

\begin{lemma}[Change of weight]\label{pbs:weight}
Let $T$ be stable, $X=T\cap\Gamma_{c,s,c'}$, and let $\phi\colon X\times G_\infty\to[0,\infty)$ be
measurable in $y$ with
\[
\sum_{(u,u')\in U_c\times U_{c'}}\phi(u\gamma u'^{-1},uy)=1\qquad(\gamma\in X,\ y\in G_\infty).
\]
Then $\sum_{\gamma\in X}\int_{G_\infty}\phi(\gamma,y)\|\pbsK_\gamma(y)\|\,dy<\infty$,
\[
A^\flat_T(c,s,c')=\sum_{\gamma\in X}\int_{G_\infty}\phi(\gamma,y)\,\pbsK^\flat_\gamma(y)\,dy,\qquad
\sum_{\gamma\in X}\int_{G_\infty}\phi(\gamma,y)\|\pbsK_\gamma(y)-\pbsK^\flat_\gamma(y)\|\,dy\le
\frac\nu2\,b^{-k^3f}.
\]
\end{lemma}
\begin{proof}
The group $\Delta=U_c\times U_{c'}$ acts on $X$ by $(u,u')\gamma=u\gamma u'^{-1}$ ($X$ is stable
under it because $T$ is stable and $\Gamma_{c,s,c'}$ is stable under $U_c\times U_{c'}$). The
stabiliser of $\gamma$ is $\{(u,\gamma^{-1}u\gamma):u\in\Lambda_\gamma\}$ with $\Lambda_\gamma=U_c\cap
\gamma U_{c'}\gamma^{-1}$. Let $F\colon X\times G_\infty\to[0,\infty]$ be measurable in $y$ with
$F(u\gamma u'^{-1},uy)=F(\gamma,y)$. For an orbit $\Delta\gamma$,
\[
\sum_{\delta\in\Delta\gamma}\int\phi(\delta,y)F(\delta,y)\,dy=\sum_{(u,u')\in\Delta/\mathrm{Stab}}
\int\phi(u\gamma u'^{-1},uy)F(\gamma,y)\,dy=\int_{\Lambda_\gamma\backslash G_\infty}F(\gamma,y)\,dy,
\]
by the substitution $y\mapsto uy$ and the cutoff property (the function $y\mapsto\sum_{\Delta/
\mathrm{Stab}}\phi(u\gamma u'^{-1},uy)$ is a cutoff for $\Lambda_\gamma$, since summing it over
$\Lambda_\gamma$ gives the full sum over $\Delta$). So the sum over $X$ of $\int\phi F$ is the same
for all admissible $\phi$, namely $\sum_{\Delta\gamma}\int_{\Lambda_\gamma\backslash G_\infty}F$. With
$F=\|\pbsK_\gamma\|$ this is $\sum_{\Delta\gamma}\int_{\Lambda_\gamma\backslash G_\infty}\|W(y)W(\gamma^{-1}y)^*\|\,dy$,
the sum of the integrals of the norms in \eqref{pb:eq-amp} (with $W^\flat=W$), which is finite because
Proposition~\ref{pb:unfold} states that the sums \eqref{pb:eq-amp} for $W$ converge absolutely for every
stable $T$; and then the same computation without absolute values (justified by Fubini) gives the
expression \eqref{pb:eq-amp} for $A^\flat_T$, using $\|W^\flat\|\le\|W\|$. For the last claim take
$\phi(\gamma,y)=\beta_c(y)\beta_{c'}(\gamma^{-1}y)$, which is admissible by \eqref{pbs:eq-cutoff};
then the sum is at most the left side of the first display in the proof of
Proposition~\ref{pbe:cut} (for the block $(c,s,c')$ only), and the proof there gives the bound
$\nu\,b^{-2r^3f+12n^3rf+400r^2f+2kf+1}\le\frac\nu2b^{-k^3f}$; the last inequality is
$12n^3r+400r^2+2k+2+k^3\le2r^3$. Indeed $2r^3-k^3\ge r^3\ge(12n^3+500)r^2\ge12n^3r+400r^2+2k+2$,
as $k<r$ and $12n^3+500\le k<r$ (Lemma~\ref{an:lem-sizes}(S7)(a)).
\end{proof}

\begin{lemma}[Non-agreeing reductions]\label{pbs:agree}
If $T$ is stable and $s$ does not agree, then $A^\flat_T(c,s,c')=0$.
\end{lemma}
\begin{proof}
Let $u,u'$ be integral elements of $U(\auA_f)$ (in the stabiliser of the lattice) with reductions
$\bar u s=s\bar u'$ and $\psi_f(u)\ne\psi_f(u')$. By strong approximation for the split unipotent
group $U$ there are $x,x'\in U(E^+)$ with $h_c^{-1}xh_c\in u\,(K_U\cap J)$ and $h_{c'}^{-1}x'h_{c'}\in
u'(K_U\cap J)$, where $K_U$ is the integral subgroup of $U(\auA_f)$; as $\psi_f$ is trivial on $U(\auA_f)
\cap J$ and invariant under conjugation by $h_c$, $\psi_f(x)=\psi_f(u)$ and $\psi_f(x')=\psi_f(u')$.
The map $\gamma\mapsto x\gamma x'^{-1}$ is a bijection of $T\cap\Gamma_{c,s,c'}$ onto $T\cap\Gamma_{c,
\bar us\bar u'^{-1},c'}=T\cap\Gamma_{c,s,c'}$ ($T$ is stable). By the substitution $y\mapsto xy$ in
Lemma~\ref{pbs:weight} with the weight $\phi(\delta,y)=\beta_c(xy)\beta_{c'}(x'\delta^{-1}
y)$ for $\delta=x^{-1}\gamma x'$ (admissible: $x$ normalises $U_c=U(E^+)\cap h_cJh_c^{-1}$ because
$h_c^{-1}xh_c\in K$ and $J$ is normal in $K$, so $\sum_{u\in U_c}\beta_c(xuy)=\sum_{u\in U_c}\beta_c(u\,
xy)=1$, and likewise for $x'$), and $W^\flat(xy)=\psi_\infty(x)W^\flat(y)$,
we get $A^\flat_T(c,s,c')=\psi_\infty(x)\aubar{\psi_\infty(x')}A^\flat_T(c,s,c')=\aubar{\psi_f(u)}
\psi_f(u')A^\flat_T(c,s,c')$, using $\psi_\infty\psi_f=1$ on $U(E^+)$. As $\psi_f(u)\ne\psi_f(u')$, the
block vanishes.
\end{proof}

\subsection{Reduction to one orbit}\label{pbs:ss-reduction}

For $\gamma\in\Gamma$ and classes $c,c'$ the \emph{amplitude majorant} is
\[
\mathsf a_{c,c'}(\gamma)=\int_{G_\infty}\beta_c(y)\,\beta_{c'}(\gamma^{-1}y)\,\|W^\flat(y)\|\,
\|W^\flat(\gamma^{-1}y)\|\,dy\in[0,\infty].
\]
By Lemma~\ref{pbs:weight} with $\phi(\gamma,y)=\beta_c(y)\beta_{c'}(\gamma^{-1}y)$,
\begin{equation}\label{pbs:eq-majorant}
\|A^\flat_T(c,s,c')\|\le\sum_{\gamma\in T\cap\Gamma_{c,s,c'}}\mathsf a_{c,c'}(\gamma).
\end{equation}

We say that $\gamma$ \emph{has cell rank $j$} ($0\le j\le n$) if (i) there is $\zeta\in E$ with
$\zeta\aubar\zeta=1$ such that $(\zeta,1)$ occurs at every $a$ with $j\le a\le\varrho$
(\S\ref{ppi:ss-statement}); (ii) $\delta_j(\gamma)\ne0$ ($\delta_0=1$); (iii) $j=n$ if and only if
$\pbg(\gamma)\ne0$. We never need that such a $j$ is unique; we always use $j=j(\gamma)$ of
Definition~\ref{ppi:def-cell}, by the following consequence of Proposition~\ref{ppi:cell}.

\begin{lemma}[Cell ranks on $\pbeGs$]\label{pbs:cells}
Every $\gamma\in\pbeGs$ has cell rank $j(\gamma)$, and $j(\gamma)\ge1$. The cell rank, the minors
$\delta_l$ and $\pbg$ are invariant under $\gamma\mapsto u\gamma u'$ for $u,u'\in U(E^+)$, and the
matrix $Z(\gamma)$ is invariant under $\gamma\mapsto u\gamma u'$ for $u,u'\in U^0(E^+)$. The sets
$\{\gamma\in T:j(\gamma)=j\}$ and $\{\gamma\in T:\auNm\pbg(\gamma)>N^f\pbsmn_j(\gamma)\}$ are stable
when $T$ is.
\end{lemma}
\begin{proof}
For $\gamma\in\pbeGs$, (i) holds at $j(\gamma)$ by definition, and (ii), (iii) by
Proposition~\ref{ppi:cell}. If $j(\gamma)=0$, then $\pbg(\gamma)=0$ and, for each $a$, there is a
common vector $fx^{(a)}=\gamma fy^{(a)}$ with $x^{(a)},y^{(a)}$ of top index $a$ and $y^{(a)}_a=1$;
by induction on $a$, $\langle f_0,\dots,f_a\rangle=\gamma\langle f_0,\dots,f_a\rangle$ (both sides have
dimension $a+1$ and the left contains $fx^{(a)}$ and the right side for $a-1$). For $a=\varrho-1$
this is $\gamma A=A$, which is excluded in $\pbeGs$. For $u\in U(E^+)$ the vectors $uf_a$ are $f_a$
plus a combination of the $f_{a'}$, $a'<a$; so $Z(u\gamma u')=C_{u^{-1}}^*Z(\gamma)C_{u'}$ with
unitriangular $C$, which preserves the leading minors, hence $\pbg$, $\pbsmn_j$ and $\auNm\pbg$; and
it preserves the occurrence of $(\zeta,1)$ at $a$ ($fx=\gamma fy$ becomes $u fx=u\gamma u'(u'^{-1}fy)$,
and $u$, $u'^{-1}$ act unitriangularly on coefficient vectors). For $u,u'\in U^0$ the matrices $C$
are $1$. The last claim follows.
\end{proof}

\begin{lemma}[Near orbits]\label{pbs:near}
If $\mathsf a_{c,c'}(\gamma)\ne0$ then $\gamma$ is near. If $\gamma$ is near, so is $u\gamma u'$ for
$u,u'\in U^0(E^+)$.
\end{lemma}
\begin{proof}
If $\mathsf a_{c,c'}(\gamma)\ne0$, some $y$ has $\beta_c(y)\beta_{c'}(\gamma^{-1}y)W^\flat(y)\otimes
W^\flat(\gamma^{-1}y)\ne0$. Then $y$ and $\gamma^{-1}y$ are reduced, and at every place $W^\flat_\tau
\ne0$ at their components, so $\pbel_\tau\le b^{6n+2}$ there (Proposition~\ref{pbe:local}(e)): both
are boxed. If $y$ and $\gamma^{-1}y$ are boxed, so are $uy$ and $(u\gamma u')^{-1}uy=u'^{-1}\gamma^{-1}
y$, by Lemma~\ref{pbs:reduced}(a) and Lemma~\ref{pbe:len}(a) ($U^0\subseteq F_\tau\cap R_\tau$).
\end{proof}

The heart of the proof is the following bound for one orbit, proved in
\S\S\ref{pbs:ss-schur}--\ref{pbs:ss-orbit}.

\begin{theorem}[One orbit]\label{pbs:orbit}
Let $\gamma\in\Gamma_{c,s,c'}$ have cell rank $j\ge1$ and satisfy $\auNm\pbg(\gamma)\le N^f
\pbsmn_j(\gamma)$, and let $\pbsO=U^0_c\gamma U^0_{c'}$. Then
\[
\sum_{\delta\in\pbsO}\mathsf a_{c,c'}(\delta)\le\nu\,\mathsf S_j\bigl(\mathsf G_j(\gamma)\bigr).
\]
\end{theorem}

For the orbits on the open cell we use the following bound, deduced in \S\ref{pbs:ss-open} from
(O1) and Lemma~\ref{pbs:O2}.

\begin{proposition}[Open cell]\label{pbs:open-block}
Let $T$ be stable with $\auNm\pbg(\gamma)>N^f\pbsmn_n(\gamma)$ on $T$. Then $\|A^\flat_T(c,s,c')\|\le
\nu\sum_\gamma\mathsf S_n(\auNm\pbg(\gamma))$, the sum over the near and chosen
$\gamma\in T\cap\Gamma_{c,s,c'}$.
\end{proposition}

\begin{proposition}[One cell rank]\label{pbs:rank}
Let $T\subseteq\pbeGs$ be stable, all of whose elements have cell rank $j$. Then
\[
\|A^\flat_T(c,s,c')\|\le\nu\sum_{\gamma}\mathsf S_j\bigl(\mathsf G_j(\gamma)\bigr),
\]
the sum over the near and chosen $\gamma\in T\cap\Gamma_{c,s,c'}$.
\end{proposition}
\begin{proof}
Let $\mathcal K=\{\gamma:\auNm\pbg(\gamma)>N^f\pbsmn_j(\gamma)\}$; $T\cap\mathcal K$ and $T\smallsetminus
\mathcal K$ are stable (Lemma~\ref{pbs:cells}) and $A^\flat_T=A^\flat_{T\cap\mathcal K}+A^\flat_{T
\smallsetminus\mathcal K}$ (the sums converge absolutely). On $\mathcal K$, $\pbg\ne0$, so $j=n$, the
indices are $1$ and $\mathsf G_n(\gamma)=\auNm\pbg(\gamma)$; Proposition~\ref{pbs:open-block} bounds
the first block. For the second, by \eqref{pbs:eq-majorant} it suffices to bound $\sum\mathsf a_{c,c'
}(\gamma)$ over $X=(T\smallsetminus\mathcal K)\cap\Gamma_{c,s,c'}$. Group $X$ into the orbits $\pbsO$
of $U^0_c\times U^0_{c'}$ (they lie in $X$: $U^0_c\subseteq U_c$ and $T$ is stable); each orbit has
one chosen element $\gamma_\pbsO$. If $\sum_{\delta\in\pbsO}\mathsf a_{c,c'}(\delta)\ne0$, some $\delta
\in\pbsO$ is near (Lemma~\ref{pbs:near}), hence so is $\gamma_\pbsO$, and by Theorem~\ref{pbs:orbit}
applied to $\gamma_\pbsO$ (which has cell rank $j\ge1$ and lies outside $\mathcal K$) the orbit
contributes at most $\nu\mathsf S_j(\mathsf G_j(\gamma_\pbsO))$. Adding the two bounds gives the
claim, as the near chosen elements of $T\cap\mathcal K$ and of $T\smallsetminus\mathcal K$ are
disjoint sets.
\end{proof}

\subsection{The weight of Schur's test}\label{pbs:ss-schur}

Fix a real place $\tau$. In the normalised basis the hermitian form is $h(v,w)=\langle\mathsf J_1v,
w\rangle$ (the form of \S\ref{pbe:ss-real} with its variables exchanged). Let $V_-$ be the $-1$-eigenspace of $\mathsf J_1$, spanned by the
$\tilde e_a-\tilde e_a^*$ and $\tilde z$; it has dimension $n$ and $h(v,v)=-\|v\|^2$ on it. For $y\in
G_\tau$ the space $P_y=yV_-$ is a maximal negative definite subspace, and $h(yv,yv)=-\|v\|^2$ for
$v\in V_-$. Let $N_0=\langle\tilde f_0,\dots,\tilde f_\varrho\rangle=A_\tau\oplus\C\tilde z$; then
$h(\sum_ax_a\tilde f_a,\sum_ax_a\tilde f_a)=-|x_\varrho|^2\le0$. Let $\pi_y$ be the $h$-orthogonal
projection onto $P_y$ (the form is non-degenerate on $P_y$). Let $F_\tau$ be the stabiliser of
$\tilde f_0,\dots,\tilde f_\varrho$ in $G_\tau$, and $U^0_\tau=U_\tau\cap F_\tau$.

\begin{lemma}\label{pbs:proj}
$\pi_y$ is injective on $N_0$, hence an isomorphism $N_0\to P_y$.
\end{lemma}
\begin{proof}
If $w\in N_0$ and $\pi_yw=0$, then $w\in P_y^\perp$, on which $h$ is positive definite (the signature
is $(k+\varrho,n)$ and $P_y$ is negative definite of dimension $n$); as $h(w,w)\le0$, $w=0$.
\end{proof}

\begin{definition}[Volume ratio]\label{pbs:def-v}
For $g,y\in G_\tau$ let $T_{g,y}=(\pi_y|_{N_0})^{-1}\circ\pi_y\circ g|_{N_0}\in\operatorname{End}(N_0)$ and
$\pbsv_y(g)=|\det T_{g,y}|$. (Up to the common factor $|\det(\pi_y|_{N_0})|$ this is the ratio of the
volumes of the projections of $gN_0$ and of $N_0$ to $P_y$.)
\end{definition}

\begin{lemma}\label{pbs:v-props}
Let $g,y\in G_\tau$.
\begin{enumerate}
\item[(a)] $0<\pbsv_y(g)<\infty$, and $y\mapsto\pbsv_y(g)$ is continuous.
\item[(b)] $\pbsv_{g^{-1}y}(g^{-1})=\pbsv_y(g)^{-1}$.
\item[(c)] For $u\in F_\tau$, $\pbsv_{uy}(g)=\pbsv_y(u^{-1}g)$ and $\pbsv_y(gu)=\pbsv_y(g)$.
\end{enumerate}
\end{lemma}
\begin{proof}
(a) $\pi_y g|_{N_0}$ is injective by Lemma~\ref{pbs:proj} applied to $gN_0$ in place of $N_0$ (the
form is $\le0$ on $gN_0$ too); the continuity follows from that of $y\mapsto\pi_y$. (b) $\pi_{g^{-1}y}
=g^{-1}\pi_yg$, so $T_{g^{-1},g^{-1}y}=(g^{-1}\pi_yg|_{N_0})^{-1}g^{-1}\pi_y|_{N_0}=(\pi_yg|_{N_0})^{-1}
\pi_y|_{N_0}=T_{g,y}^{-1}$. (c) $\pi_{uy}=u\pi_yu^{-1}$ and $u|_{N_0}=1$, so $T_{g,uy}=(\pi_yu^{-1}|_{N_0}
)^{-1}\pi_yu^{-1}g|_{N_0}=T_{u^{-1}g,y}$; and $gu|_{N_0}=g|_{N_0}$.
\end{proof}

For $y\in G_\infty$ and $\gamma\in\Gamma$ put $\pbsv_y(\gamma)=\prod_\tau\pbsv_{y_\tau}(\gamma_\tau)$. Let
$\beta^\square_c=\beta_c\cdot\mathbf 1_{\text{boxed}}$. For an integer $i\ge0$ the \emph{Schur
integral} is
\[
I^{(i)}_{c,c'}(\gamma)=\int_{G_\infty}\beta^\square_c(y)\,\|W^\flat(y)\|^2\,\beta^\square_{c'}(
\gamma^{-1}y)\,\pbsv_y(\gamma)^{-(k+i)}\,dy .
\]

\begin{lemma}[Cauchy--Schwarz with a weight]\label{pbs:cs}
For every $i\ge0$, $\mathsf a_{c,c'}(\gamma)^2\le I^{(i)}_{c,c'}(\gamma)\,I^{(i)}_{c',c}(\gamma^{-1})$.
\end{lemma}
\begin{proof}
Where $\beta_c(y)\beta_{c'}(\gamma^{-1}y)W^\flat(y)\otimes W^\flat(\gamma^{-1}y)\ne0$, both points are
boxed (proof of Lemma~\ref{pbs:near}), so $\beta_c=\beta_c^\square$ there. Write $e=k+i$ and
$\pbsv=\pbsv_y(\gamma)$. The integrand of $\mathsf a_{c,c'}(\gamma)$ is the product of $(\beta^\square_c
(y)\|W^\flat(y)\|^2\beta^\square_{c'}(\gamma^{-1}y)\pbsv^{-e})^{1/2}$ and $(\beta^\square_{c'}(\gamma^{-1}y)
\|W^\flat(\gamma^{-1}y)\|^2\beta^\square_c(y)\pbsv^{e})^{1/2}$. By the Cauchy--Schwarz inequality,
$\mathsf a_{c,c'}(\gamma)^2$ is at most $I^{(i)}_{c,c'}(\gamma)$ times the integral of the square of
the second factor. In that integral substitute $y=\gamma y'$; by Lemma~\ref{pbs:v-props}(b) (with
$g=\gamma^{-1}$, applied at each place), $\pbsv_{\gamma y'}(\gamma)=\pbsv_{y'}(\gamma^{-1})^{-1}$, so the
integral is $I^{(i)}_{c',c}(\gamma^{-1})$.
\end{proof}

\begin{lemma}[Integration in stages]\label{pbs:stages}
Let $\Phi\colon G_\infty\to[0,\infty]$ be measurable and left invariant under $U^0_c$. Let $(X,m)$
be a $\sigma$-finite measure space, $t\colon X\to U^0_\infty$ measurable, $B\subseteq X$ with $0<
m(B)<\infty$, and $C\ge0$ such that
\[
\int_B\Phi(u\,t(x)\,y)\,dm(x)\le m(B)\,C\qquad\text{for all }u\in U^0_\infty\text{ and all boxed }y.
\]
Then $\int_{G_\infty}\beta^\square_c(y)\|W^\flat(y)\|^2\Phi(y)\,dy\le C\,\nu\,d_\infty$, where $d_\infty=
\dim\pbW=\prod_\tau d_\tau$.
\end{lemma}
\begin{proof}
Disintegrate $dg$ along the orbits of $U^0_\infty$: $\int_{G_\infty}F=\int_{U^0_\infty\backslash G_\infty}
\int_{U^0_\infty}F(ua)\,du^0\,da$ for $F\ge0$. On an orbit $U^0_\infty a$ the functions $\omega$,
$\mathbf 1_{\text{boxed}}$ and $\|W^\flat\|$ are constant (Lemma~\ref{pbs:reduced}(a),
Lemma~\ref{pbe:len}(a), and $W^\flat(uy)=\psi_\infty(u)W^\flat(y)$), so $\beta^\square_c(ua)=p_c(ua)\,
\omega(a)\mathbf 1_{\text{boxed}}(a)$. It therefore suffices to show, for boxed $a$,
\begin{equation}\label{pbs:eq-stage}
\int_{U^0_\infty}p(u)\,\Phi_a(u)\,du^0\le C\int_{U^0_\infty}p(u)\,du^0,\qquad p(u)=p_c(ua),\ \Phi_a(u)=
\Phi(ua);
\end{equation}
for then the left side of the lemma is at most $C\int\beta_c\|W^\flat\|^2=C\nu\int_{U_\infty\backslash
G_\infty}\|W^\flat\|^2$ by \eqref{pbs:eq-nu}, and $\int_{U_\tau\backslash G_\tau}\|W^\flat_\tau\|^2\le
\int\|W_\tau\|^2=d_\tau$ (Proposition~\ref{pbe:local}(e)). Now $p$ is a cutoff for the left action of
the discrete group $U^0_c$ on $U^0_\infty$, and $\Phi_a$ is invariant under it. For $s\in U^0_\infty$,
$\int p(u)\Phi_a(us)\,du^0=\int_{U^0_c\backslash U^0_\infty}\Phi_a(us)=\int_{U^0_c\backslash U^0_\infty}
\Phi_a(u)=\int p\,\Phi_a$, because right translation by $s$ preserves the invariant measure of
$U^0_c\backslash U^0_\infty$ ($U^0_\infty$ is unimodular). Integrating over $x\in B$ with $s=t(x)$ and
using Tonelli,
\[
m(B)\int p\,\Phi_a=\int_{U^0_\infty}p(u)\int_B\Phi(u\,t(x)\,a)\,dm(x)\,du^0\le m(B)\,C\int p ,
\]
which is \eqref{pbs:eq-stage}.
\end{proof}

\subsection{The weight at one real place}\label{pbs:ss-oneplace}

Fix $\tau$ and drop it from the notation. For $g\in G_\tau$ let $\tilde Z(g)=(h(\tilde f_a,g\tilde
f_b))_{a,b}$, and let $\Phi(g)=(\Phi(g)_b)_{0\le b\le\varrho}$ with $\Phi(g)_b\in X_{0,\tau}$ the
component in $X_{0,\tau}$ of $g\tilde f_b$ for the decomposition $Y_\tau=A_\tau\oplus X_{0,\tau}\oplus
\C\tilde z\oplus A^*_\tau$, $A^*_\tau=\langle\tilde e^*_a\rangle$. We say that $(\zeta,1)$ \emph{occurs at
$a$ for $g$} if there are $x,x'\in\C^n$ supported on $\{0,\dots,a\}$ with $\sum x_b\tilde f_b=g\sum
x'_b\tilde f_b$, $x_a=\zeta$, $x'_a=1$, and that \emph{norm-one pairs occur from $j$} if for one $\zeta$
with $|\zeta|=1$ this holds at every $a$ with $j\le a\le\varrho$. If $\gamma\in\Gamma$ has cell rank
$j$, norm-one pairs occur from $j$ for $\gamma_\tau$ (rescale the common vectors of
Definition~\ref{ppi:def-cell}: $f_b=\tilde f_b/s_b$, and multiply by $s_a$; the leading coefficients
become $\tau(\zeta)$ and $1$), and $\tilde Z(\gamma_\tau)_{ab}=s_as_b\,\tau(Z(\gamma)_{ab})$.

Let $0\le j\le n$, $i=\min(\varrho,j)$. The \emph{subframe} is $\mathsf s=(\tilde f_0,\dots,\tilde
f_{i-1},\tilde f_\varrho)$; for a matrix $Z$ indexed by the frame, $B(Z)$ is its $(i+1)\times(i+1)$
block on the indices of the subframe, and $B_<(Z)$ the first $i$ rows of $B(Z)$. For $g\in G_\tau$
write $g\mathsf s$ for $(g\tilde f_0,\dots,g\tilde f_{i-1},g\tilde f_\varrho)$ and $\Phi_{\mathsf s}(g)$ for
the corresponding columns of $\Phi(g)$.

\begin{lemma}[Common vectors]\label{pbs:common}
Let norm-one pairs occur from $j$ for $g$, $Z=\tilde Z(g)$.
\begin{enumerate}
\item[(a)] For each $b$ with $j\le b<\varrho$ there are $x^{(b)},x'^{(b)}$, supported on $\{0,\dots,
b\}$, with $x'^{(b)}_b=1$, $|x^{(b)}_b|=1$ and $c_b:=\sum_ax^{(b)}_a\tilde f_a=g\sum_ax'^{(b)}_a\tilde
f_a=:gc'_b$; and $c_b\in A_\tau\cap gA_\tau$ is $h$-orthogonal to $gN_0$, and $x^{(b)*}Z=0$.
\item[(b)] If $j\le\varrho$ there are $x\in\C^n$ with $|x_\varrho|=1$ and $\mathsf c\in\C^{i+1}$ with
$\mathsf c_{i}=1$ (last entry) such that $g\mathsf s\,\mathsf c=\sum_ax_a\tilde f_a$, and $B(Z)\,\mathsf c
=(0,\dots,0,-x_\varrho)^{\mathsf T}$.
\item[(c)] $|\det B(Z)|=|\det(Z_{ab})_{a,b<j}|$ (for $j=n$ this is $\det Z$).
\end{enumerate}
\end{lemma}
\begin{proof}
(a) The vectors exist by definition. As $c_b=gc'_b$ and $g$ is an isometry, $|x^{(b)}_\varrho|=|x'^{(b)
}_\varrho|$ ($h(c,c)=-|x_\varrho|^2$ on $N_0$), and both vanish as $b<\varrho$. So $c_b\in A_\tau$, $c'_b\in
A_\tau$. For $w\in N_0$, $h(c_b,gw)=h(c'_b,w)=0$ because $A_\tau\perp N_0$; this is $x^{(b)*}Z=0$.
(b) Take the common vector at $a=\varrho$: $\sum x_a\tilde f_a=g\sum x'_a\tilde f_a$ with $x'_\varrho=1$,
$|x_\varrho|=1$. Subtracting the common vectors of (a) (multiplied by $x'_b$, for $b=\varrho-1,\dots,j$
in turn) we may assume $x'_b=0$ for $i\le b<\varrho$; then $x'$ is supported on the subframe and
$\mathsf c$ is its restriction (the norm equality still gives $|x_\varrho|=1$). The products of the
subframe with $\sum x_a\tilde f_a$ are $h(\tilde f_a,\sum x_b\tilde f_b)=-x_\varrho\delta_{a\varrho}$.
(c) If $j=n$ the subframe is the whole frame. If $j\le\varrho$, $i=j$; by Cramer's rule, replacing the
last column of $B(Z)$ by $B(Z)\mathsf c$ multiplies the determinant by $\mathsf c_i=1$; expanding the new
matrix along its last column, whose only non-zero entry is $-x_\varrho$ in the last row, gives $\pm
x_\varrho\det(Z_{ab})_{a,b<j}$.
\end{proof}

\begin{lemma}[The weight modulo common vectors]\label{pbs:mod}
Let norm-one pairs occur from $j$ for $g$, let $y\in G_\tau$, and let $\omega_0,\dots,\omega_i\in P_y$
be orthogonal to $A_\tau\cap gA_\tau$, with $\det(h(\mathsf s_a,\omega_l))_{a,l}\ne0$. Then
\[
\pbsv_y(g)=\frac{|\det(h(g\mathsf s_a,\omega_l))_{a,l}|}{|\det(h(\mathsf s_a,\omega_l))_{a,l}|}.
\]
\end{lemma}
\begin{proof}
Let $\lambda\colon Y_\tau\to\C^{i+1}$, $\lambda(w)=(h(w,\omega_l))_l$. As $\omega_l\in P_y$,
$\lambda\circ\pi_y=\lambda$, hence $\lambda\circ T_{g,y}=\lambda\circ g$ on $N_0$. The map $\lambda$ is
an isomorphism on $\langle\mathsf s\rangle$ and vanishes on the vectors $c_b$ ($j\le b<\varrho$, and
$i\le b$) of Lemma~\ref{pbs:common}(a). Two bases of $N_0$: $\mathcal B'=(\mathsf s,(c'_b)_{i\le b<
\varrho})$ and $\mathcal B=(\mathsf s,(c_b)_{i\le b<\varrho})$. In the frame order, the matrix of
$\mathcal B'$ (resp.\ $\mathcal B$) in the basis $\tilde f$ is triangular with diagonal entries $1$
(resp.\ $1$ or $x^{(b)}_b$, of modulus $1$); so $|\det T_{g,y}|$ is the modulus of the determinant of
the matrix of $T_{g,y}$ from $\mathcal B'$ to $\mathcal B$. As $T_{g,y}c'_b=c_b$ (because $\pi_ygc'_b=
\pi_yc_b$), this matrix is block triangular $\bigl(\begin{smallmatrix}\mathsf X&0\\\mathsf Y&1
\end{smallmatrix}\bigr)$, with $T_{g,y}\mathsf s_l=\sum_a\mathsf X_{al}\mathsf s_a+\sum_b\mathsf Y_{bl}c_b$.
Applying $\lambda$: $\lambda(g\mathsf s_l)=\lambda(T_{g,y}\mathsf s_l)=\sum_a\mathsf X_{al}\lambda(\mathsf
s_a)$. Taking determinants, $\det(\lambda(g\mathsf s))=\det\mathsf X\cdot\det(\lambda(\mathsf s))$ and
$\pbsv_y(g)=|\det\mathsf X|$.
\end{proof}

\begin{lemma}[Triangular bases]\label{pbs:tri}
Let $Z$ be a matrix indexed by the frame with $\delta:=\det B(Z)\ne0$, and $y\in G_\tau$. Put $D=Z_{
\cdot,\mathsf s}B(Z)^{-1}$ ($n\times(i+1)$). There are $\omega_0,\dots,\omega_i\in P_y$ and a lower
triangular invertible $L$ with $h(\omega_l,\omega_m)=-\delta_{lm}$ and $(h(\tilde f_a,\omega_l))_{a,l}=
DL$. If $g$ has $\tilde Z(g)=Z$ and norm-one pairs occur from $j$ for $g$, then:
\begin{enumerate}
\item[(a)] the $\omega_l$ are orthogonal to $A_\tau\cap gA_\tau$ and $(h(\mathsf s_a,\omega_l))=L$; hence
$\pbsv_y(g)=|\det(h(g\mathsf s_a,\omega_l))|/|\det L|$;
\item[(b)] with $\omega'=g\mathsf s\,B(Z)^{-1}L$ and $v=(v_0,\dots,v_i)$, $v_l\in X_{0,\tau}$ the
$X_0$-component of $\omega'_l-\omega_l$:
\[
|\delta|\det(1+vv^*)\le2^{i+1}|\det L|^2\,\pbsv_y(g);
\]
\item[(c)] if $\|y^{-1}\tilde f_a\|\le\rho$ for all $a$, then $|\det L|\le\rho^{i+1}$;
\item[(d)] if $\pbel_\tau(y)\le\theta$, then $|L_{ii}|^2\le\theta^{2(j-i)}(1+\|v_i\|^2)$.
\end{enumerate}
Here $vv^*$ is the endomorphism $x\mapsto\sum_lv_l\langle x,v_l\rangle$ of $X_{0,\tau}$.
\end{lemma}
\begin{proof}
\emph{Construction.} By Lemma~\ref{pbs:proj}, $w\mapsto(h(\tilde f_a,w))_a$ is an isomorphism $P_y\to
\C^n$ ($P_y\cong N_0^*$ through $h$, as $\pi_y|_{N_0}$ is an isomorphism and $h$ is non-degenerate on
$P_y$). Let $\omega^0_l\in P_y$ have products $D_{\cdot l}$; they are linearly independent because
$D$ has the left inverse $(\text{rows of }\mathsf s)$: the rows of $D$ at the subframe indices form
$B(Z)B(Z)^{-1}=1$. The matrix $\mathsf G_0=(-h(\omega^0_l,\omega^0_m))_{l,m}$ is positive definite, since $-h$ is
positive definite on $P_y$; by the Cholesky decomposition $\mathsf G_0^{-1}=LL^*$ with $L$ lower
triangular and invertible, i.e.\ $L^*\mathsf G_0L=1$. Put $\omega_m=\sum_l\omega^0_lL_{lm}$; then
$h(\omega_l,\omega_m)=-\delta_{lm}$. Then the products with the frame are $DL$.

(a) Let $c=\sum x_a\tilde f_a\in A_\tau\cap gA_\tau$; as in Lemma~\ref{pbs:common}(a), $x^*Z=0$
(write $c=gc'$ with $c'\in A_\tau$; then $h(c,gw)=h(c',w)=0$ for $w\in N_0$). Then $h(c,\omega_l)=(x^*DL)_l=(x^*Z_{\cdot,\mathsf s}B^{-1}L)_l=0$. The
products of the subframe with $\omega$ are the subframe rows of $DL$, that is $L$. Lemma~\ref{pbs:mod}
gives the formula.

(b) The products of the frame with $\omega'$ are $Z_{\cdot,\mathsf s}B^{-1}L=DL$, the same as with
$\omega$. So $\omega'_l-\omega_l\perp N_0$, i.e.\ $\omega'_l-\omega_l\in(A_\tau\oplus\C\tilde z)^\perp=
A_\tau\oplus X_{0,\tau}$, and as $A_\tau$ is isotropic and orthogonal to $A_\tau\oplus X_{0,\tau}$,
$h(\omega'_l-\omega_l,\omega'_m-\omega_m)=\langle v_m,v_l\rangle$ (with $h=\langle\ ,\ \rangle$ on
$X_{0,\tau}$). Let $\mathsf K=(h(\omega_l,\omega'_m))_{l,m}$ and $\mathsf G'=(h(\omega'_l,\omega'_m))$.
Expanding the Gram matrix of $\omega'-\omega$ gives $v^\dagger v=\mathsf G'-\mathsf K-\mathsf K^*+(-1)$,
where $v^\dagger v=(\langle v_m,v_l\rangle)_{l,m}$; so $(-\mathsf K)+(-\mathsf K)^*-(1+v^\dagger v)=-\mathsf G'$,
which is positive semidefinite since $\omega'_l\in gN_0$ and $h\le0$ on $gN_0$. By the
Ostrowski--Taussky inequality (Lemma~\ref{pbs:OT} below), $\det(1+v^\dagger v)\le2^{i+1}|\det\mathsf K|$.
Now $\mathsf K=(h(\omega_l,g\mathsf s_m))\,B^{-1}L$, and $(h(\omega_l,g\mathsf s_m))_{l,m}$ is the
conjugate transpose of $(h(g\mathsf s_m,\omega_l))_{m,l}$, so by (a) $|\det\mathsf K|=\pbsv_y(g)|\det L|
^2/|\delta|$. Finally $\det(1+v^\dagger v)=\det(1+vv^*)$.

(c) $L_{al}=h(\mathsf s_a,\omega_l)=h(y^{-1}\mathsf s_a,y^{-1}\omega_l)$, and the $y^{-1}\omega_l\in V_-$
are orthonormal for $\|\cdot\|$ (since $h=-\|\cdot\|^2$ on $V_-$). By \eqref{pbe:eq-norm}, $|L_{ll}|\le
\|y^{-1}\mathsf s_l\|\le\rho$; and $\det L=\prod_lL_{ll}$.

(d) \emph{Case $j\le\varrho$} ($i=j$). Let $x,\mathsf c$ be as in Lemma~\ref{pbs:common}(b). The last
column of $L$ is $L_{ii}e_i$ ($L$ is lower triangular), and $B^{-1}e_i=-\mathsf c/x_\varrho$, so $\omega'_i=
L_{ii}\,g\mathsf sB^{-1}e_i=t\sum_ax_a\tilde f_a$ with $t=-L_{ii}/x_\varrho$. Thus $\omega'_i\in N_0$,
hence $\omega'_i\perp\omega'_i-\omega_i$, and $h(\omega'_i,\omega'_i)=-|tx_\varrho|^2=-|L_{ii}|^2$. So
$-1=h(\omega_i,\omega_i)=h(\omega'_i,\omega'_i)+h(\omega'_i-\omega_i,\omega'_i-\omega_i)=-|L_{ii}|^2+\|v_i
\|^2$, i.e.\ $|L_{ii}|^2=1+\|v_i\|^2$. \emph{Case $j=n$} ($i=\varrho$, $j-i=1$). Choose $s\in R_\tau$
with $\|sy\|\le\theta$ (possible since $\pbel_\tau(y)=\inf_{s\in R_\tau}(\|sy\|+1)\le\theta$). Then $(sy)^{-1}\tilde z=y^{-1}s^{-1}\tilde z=y^{-1}(\tilde z+\sum_{a<\varrho}
\mu_a\tilde f_a)$ (elements of $R_\tau$ are unitriangular on the frame), and $h(\tilde f_a,\omega_\varrho)
=L_{a\varrho}=0$ for $a<\varrho$. Hence $L_{\varrho\varrho}=h((sy)^{-1}\tilde z,y^{-1}\omega_\varrho)$, of
modulus at most $\|(sy)^{-1}\|=\|sy\|\le\theta$; and $1\le1+\|v_\varrho\|^2$.
\end{proof}

\begin{lemma}[Ostrowski--Taussky]\label{pbs:OT}
Let $H$ and $\Theta+\Theta^*-H$ be positive semidefinite $m\times m$ complex matrices. Then $|\det H|\le
2^m|\det\Theta|$.
\end{lemma}
\begin{proof}
If $\det H=0$ there is nothing to prove; otherwise $H$ is positive definite. Let $\ell$ be an
eigenvalue of $2H^{-1}\Theta$ with eigenvector $v\ne0$: $2\Theta v=\ell Hv$. Then $\ell\,v^*Hv=2v^*\Theta
v$, and taking real parts, $(\operatorname{Re}\ell)\,v^*Hv=v^*(\Theta+\Theta^*)v\ge v^*Hv>0$; so
$\operatorname{Re}\ell\ge1$ and $|\ell|\ge1$. Hence $|\det(2H^{-1}\Theta)|=2^m|\det\Theta|/\det H\ge1$.
\end{proof}

\begin{lemma}[Hua's integral]\label{pbs:hua}
Let $k\ge7$, $\mathsf G$ a positive definite $k\times k$ matrix, $\mathsf S$ an invertible $i\times i$
matrix and $\mathsf c$ an $i\times k$ matrix. With Lebesgue measure on $\C^{i\times k}$,
\[
\int_{\C^{i\times k}}\det\bigl(\mathsf G+(\mathsf S\mathsf A+\mathsf c)^*(\mathsf S\mathsf A+\mathsf
c)\bigr)^{-(k+i)}d\mathsf A\le|\det\mathsf S|^{-2k}(\det\mathsf G)^{-k}.
\]
\end{lemma}
\begin{proof}
The real-linear map $\mathsf A\mapsto\mathsf S\mathsf A$ has determinant $|\det\mathsf S|^{2k}$, so after
the affine substitution it suffices to treat $\mathsf S=1$, $\mathsf c=0$. Then $\mathsf A^*\mathsf A=
\sum_l w_l^*w_l$ over the rows $w_l\in\C^k$ (row vectors). We show by induction on the number $i$ of
rows: for $\mathsf M>0$ and $r\ge k$,
\[
\int\det\Bigl(\mathsf M+\sum_{l<i}w_l^*w_l\Bigr)^{-(i+r)}dw\le\Bigl(\frac{\pi^k}{k!}\Bigr)^i(\det
\mathsf M)^{-r}.
\]
For one row: write $\mathsf M=\mathsf C^*\mathsf C$ and $w=u\mathsf C$; then $\det(\mathsf M+w^*w)=
\det\mathsf M\,(1+\|u\|^2)$, $dw=\det\mathsf M\,du$, and $\int_{\C^k}(1+\|u\|^2)^{-(r+1)}du\le\int(1+\|u\|^2)
^{-(k+1)}du=\pi^k/k!$; so the integral is at most $(\pi^k/k!)(\det\mathsf M)^{-r}$. For $i$ rows,
integrate first over the rows $l\ge1$ with $\mathsf M'=\mathsf M+w_0^*w_0$ and $r+1$ in place of $r$
(induction hypothesis), then over $w_0$ with the one-row case. Finally $\pi^k/k!\le1$ for $k\ge7$.
\end{proof}

\begin{theorem}[One real place]\label{pbs:oneplace}
Let $k\ge7$, $y\in G_\tau$ with $\|y^{-1}\tilde f_a\|\le\rho$ for all $a$ and $\pbel_\tau(y)\le\theta$,
where $\rho,\theta\ge1$. Let $1\le j\le n$, $i=\min(\varrho,j)$, let $g_0\in G_\tau$ be such that
norm-one pairs occur from $j$ for $g_0$, and let $Z=\tilde Z(g_0)$ with $\delta=\det B(Z)\ne0$. There is
a measurable $\varphi\colon X_{0,\tau}^{\,i}\to[0,\infty)$ with
\[
\int_{X_{0,\tau}^i}\varphi\le\frac{(2\rho^2)^{(i+1)i}\bigl(2^{i+1}\theta^{2(j-i)}\bigr)^k}{|\delta|^{k+i}}
\]
such that $\pbsv_y(g)^{-(k+i)}\le\varphi(a_0,\dots,a_{i-1})$ for every $g\in G_\tau$ with $\tilde
Z(g)=Z$, norm-one pairs occurring from $j$, and $\Phi_{\mathsf s}(g)=\Phi_{\mathsf s}(g_0)+\sum_{l<i}a_l
\,B_<(Z)_{l,\cdot}$ (that is, $\Phi_{\mathsf s}(g)_m=\Phi_{\mathsf s}(g_0)_m+\sum_{l<i}B(Z)_{lm}a_l$).
\end{theorem}
\begin{proof}
Take $\omega,L$ from Lemma~\ref{pbs:tri} for $Z$ and $y$; they do not depend on $g$. Let $v^{(0)}$ be
the vectors of Lemma~\ref{pbs:tri}(b) for $g_0$. For $g$ as in the statement, $\omega'(g)=g\mathsf s\,
B^{-1}L$, so the $X_0$-components satisfy $v(g)=\Phi_{\mathsf s}(g)B^{-1}L-(\omega)_{X_0}$ (in the
obvious matrix notation, columns indexed by the subframe), whence
\[
v(g)=v^{(0)}+\Bigl(\sum_{l<i}a_lB_{l\cdot}\Bigr)B^{-1}L=v^{(0)}+\Bigl(\sum_{l<i}a_lL_{l\cdot}\Bigr),
\]
because the first $i$ rows of $B\,B^{-1}L$ are those of $L$. As $L$ is lower triangular, $L_{li}=0$ for
$l<i$: the last vector $v_i(g)=v^{(0)}_i$ does not depend on $g$, and the others are $v_m(g)=v^{(0)}_m+
\sum_{l<i}L_{lm}a_l$ ($m<i$). Put $\mathsf G=1+v^{(0)}_iv^{(0)*}_i$, $L_1=(L_{lm})_{l,m<i}$ and, in
coordinates of an orthonormal basis of $X_{0,\tau}$, $\mathsf w(\mathsf a)=L_1^{\mathsf T}\mathsf a+
\mathsf c$ ($i\times k$; rows $w_m=v_m(g)$). Then $1+v(g)v(g)^*=\mathsf G+\mathsf w^*\mathsf w$, and by
Lemma~\ref{pbs:tri}(a),(b)
\[
\pbsv_y(g)^{-(k+i)}\le\varphi(\mathsf a):=\Bigl(\frac{2^{i+1}|\det L|^2}{|\delta|}\Bigr)^{k+i}\det(\mathsf
G+\mathsf w(\mathsf a)^*\mathsf w(\mathsf a))^{-(k+i)} .
\]
By Lemma~\ref{pbs:hua} ($\det L_1^{\mathsf T}=\det L_1\ne0$, $|\det L|=|\det L_1||L_{ii}|$, $\det\mathsf G=1+
\|v^{(0)}_i\|^2$),
\[
\int\varphi\le\frac{(2^{i+1}|\det L|^2)^{k+i}}{|\delta|^{k+i}|\det L_1|^{2k}(\det\mathsf G)^k}=\frac{(2^{i+1}
|\det L|^2)^i\,\bigl(2^{i+1}|L_{ii}|^2/(1+\|v^{(0)}_i\|^2)\bigr)^k}{|\delta|^{k+i}}.
\]
Lemma~\ref{pbs:tri}(c),(d) (for $g_0$) bound this by $(2^{i+1}\rho^{2(i+1)})^i(2^{i+1}\theta^{2(j-i)})^k/
|\delta|^{k+i}$, which is the claim.
\end{proof}

\subsection{An orbit: colouring, shifts and packing}\label{pbs:ss-orbit}

Fix $\gamma\in\Gamma_{c,s,c'}$ of cell rank $j\ge1$, $i=\mathsf j'=\min(\varrho,j)\ge1$, and its orbit
$\pbsO=U^0_c\gamma U^0_{c'}$. For $\delta\in\Gamma$ let $\Phi(\delta)=(\Phi(\delta)_b)_{b\le\varrho}\in
X_0^n$, $\Phi(\delta)_b$ the $X_0$-component of $\delta f_b$ for $Y_0=A\oplus X_0\oplus Ez\oplus A^*$.

\begin{lemma}\label{pbs:indep}
Let $\mathfrak R$ (resp.\ $\mathfrak R'$) be the $\auO_E$-module spanned by the rows $Z(\gamma)_{a,\cdot}
\in E^n$ with $a<\varrho$ (resp.\ $a<i$). The rows with $a<i$ are linearly independent and span the
same $E$-space as those with $a<\varrho$; so $[\mathfrak R:\mathfrak R']$ is finite, and it equals
$i_\gamma$ of \S\ref{pbs:ss-statement}.
\end{lemma}
\begin{proof}
Independence: the minor $\delta_j(\gamma)$ (for $j\le\varrho$), resp.\ $\det Z(\gamma)$ (for $j=n$), is
non-zero. If $j\le a<\varrho$, the common vector $fx=\gamma fy$ with top index $a$ and $x_a=\zeta\ne0$
has $y_\varrho=0$, so $\sum_{a'\le a}\aubar{x_{a'}}Z_{a'b}=h(fx,\gamma f_b)=h(fy,f_b)=0$ for all $b$:
the row $a$ is a combination of the rows $a'<a$. For the index: $h_c$ and $h_{c'}$ fix the $f_a$, so
$h(f_a,k_\gamma f_b)=h(f_a,\gamma f_b)$ and the vectors $\xi_a(k_\gamma)$ of \S\ref{pbs:ss-statement}
are the rows of $Z(\gamma)$; the index of the completions $\widehat\auO\otimes\mathfrak R'\subseteq
\widehat\auO\otimes\mathfrak R$ is the index $[\mathfrak R:\mathfrak R']$.
\end{proof}

\begin{lemma}[Moving inside an orbit]\label{pbs:move}
Let $\delta=u\gamma u'$ with $u\in U^0_c$, $u'\in U^0_{c'}$. Then $Z(\delta)=Z(\gamma)$, $\delta$ has cell
rank $j$, $\delta\in\Gamma_{c,s,c'}$, and $\delta f_b-\gamma f_b\in A\oplus X_0$ with
\[
\Phi(\delta)_b=\Phi(\gamma)_b+\sum_{a<\varrho}Z(\gamma)_{ab}\,q_a(u),\qquad q_a(u)\in M\Lambda_c,
\]
where $q_a(u)$ is the $X_0$-component of $(u-1)e^*_{a+1}$.
\end{lemma}
\begin{proof}
The first three claims are Lemma~\ref{pbs:cells} and the stability of $\Gamma_{c,s,c'}$. Now $\delta f_b
=u\gamma f_b$. The element $u\in U^0(E^+)$ fixes $A$ and $z$; it maps $X_0$ into $X_0+A$ (it is trivial
on $A^\perp/A$); and $ue^*_a=e^*_a+q+\alpha$ with $q\in X_0$, $\alpha\in A$ (its products with $z$ and the
$e_b$ are those of $e^*_a$). For $w\in Y_0$ the $A^*$-component of $w$ is $\sum_ah(e_a,w)e^*_a$. Hence
$uw-w\in A\oplus X_0$ with $X_0$-component $\sum_ah(e_a,w)q_a$; for $w=\gamma f_b$ this is the formula.
For integrality: $u\in U_c$ means $(u-1)h_cL\subseteq Mh_cL$ (at every finite place), and $e^*_a\in h_cL$
($h_c$ fixes $e^*_a$). Write $h_cL=L_{\mathrm{hyp}}\oplus L'_{\mathrm{mid}}$ as in the proof of
Lemma~\ref{pbe:eichler}; the vector $(u-1)e^*_a=q_a+\alpha$ has $L_{\mathrm{hyp}}$-component $\alpha$ and
$L'_{\mathrm{mid}}$-component $q_a$, so $q_a\in ML'_{\mathrm{mid}}\cap X_0=M\Lambda_c$ locally everywhere,
hence $q_a\in M\Lambda_c$.
\end{proof}

\begin{lemma}[Colouring]\label{pbs:colour}
The orbit $\pbsO$ is the union of at most $i_\gamma^{\,k}$ subsets $\mathcal B$ with the following
property: for $p_0\in\mathcal B$ and every $p\in\mathcal B$ there is $l_p=(l_{p,0},\dots,l_{p,i-1})\in
\Lambda_c^{\,i}$ with
\[
\Phi(p)_b=\Phi(p_0)_b+M\sum_{a<i}Z(\gamma)_{ab}\,l_{p,a}\qquad(0\le b\le\varrho).
\]
\end{lemma}
\begin{proof}
Let $\mathfrak M=\Lambda_c\otimes_{\auO_E}\mathfrak R$ and $\mathfrak M'=\Lambda_c\otimes\mathfrak R'$,
regarded inside $X_0\otimes_EE^n=X_0^n$ via $l\otimes\xi\mapsto(\xi_bl)_b$ (injective, as $\Lambda_c$ is
projective, hence flat). Then $[\mathfrak M:\mathfrak M']=|\Lambda_c\otimes\mathfrak R/\mathfrak R'|=
[\mathfrak R:\mathfrak R']^k=i_\gamma^k$ (locally $\Lambda_c\cong\auO^k$). By Lemma~\ref{pbs:move},
$\Phi(\delta)-\Phi(\gamma)\in M\mathfrak M$ for $\delta\in\pbsO$. Colour $\delta$ by the class of $M^{-1}
(\Phi(\delta)-\Phi(\gamma))$ in $\mathfrak M/\mathfrak M'$. Within a colour, $\Phi(p)-\Phi(p_0)\in
M\mathfrak M'$, which is the claim.
\end{proof}

\subsubsection*{Shifts}
At a place $\tau$, every $u\in U^0_\tau$ satisfies, as in Lemma~\ref{pbs:move}, $\Phi(u^{-1}g)_b=\Phi(
g)_b+\sum_{a<\varrho}\tilde Z(g)_{ab}\,t_a(u)$ with $t_a(u)$ the $X_{0,\tau}$-component of $u^{-1}\tilde
e^*_{a+1}$; conversely, for every $t=(t_a)_{a<\varrho}\in X_{0,\tau}^\varrho$ there is $u_t\in U^0_\tau$ with
$t_a(u_t)=t_a$, depending continuously on $t$: put $u_t^{-1}=1$ on $N_0$, $u_t^{-1}x=x-\sum_ah(t_a,x)\tilde
e_{a+1}$ on $X_{0,\tau}$, and $u_t^{-1}\tilde e^*_{a+1}=\tilde e^*_{a+1}+t_a-\frac12\sum_{a'}h(t_{a'},t_a)
\tilde e_{a'+1}$; a direct check shows that $u_t^{-1}$ preserves $h$, and it lies in $U^0_\tau$. Let
$V_{\mathrm{col}}=\prod_\tau X_{0,\tau}^{\,i}$, with Lebesgue measure for orthonormal coordinates
(the basis $x_1,\dots,x_k$ of $X_0$ is orthonormal at every $\tau$), and for $\mathsf x=(\mathsf x_{\tau,
a})\in V_{\mathrm{col}}$ let $u_{\mathsf x}\in U^0_\infty$ have components $u_{t(\tau)}$ with $t(\tau)_a=
(M/s_a)\mathsf x_{\tau,a}$ for $a<i$ and $t(\tau)_a=0$ for $a\ge i$. Let $L_{\mathrm{col}}\subseteq
V_{\mathrm{col}}$ be the image of $\Lambda_c^{\,i}$ under $l\mapsto(\tau(l_a))_{\tau,a}$.

\begin{lemma}\label{pbs:col}
$L_{\mathrm{col}}$ is a lattice in $V_{\mathrm{col}}$ with a measurable fundamental domain $\mathsf D$
with $1\le\pbvol(\mathsf D)<\infty$.
\end{lemma}
\begin{proof}
$\Lambda_c^{\,i}$ is a projective $\auO_E$-module of rank $ik$, and $l\mapsto(\tau(l))_\tau$ maps it onto a
lattice of $V_{\mathrm{col}}$ (real dimension $2ikf=[E:\Q]\,ik$). For $l,l'\in\Lambda_c^{\,i}$ the
Euclidean product of their images is $\sum_\tau\operatorname{Re}\tau(\alpha)$ with $\alpha=\sum_ah(l_a,
l'_a)\in\auO_E$ ($h_cL$ is an integral hermitian lattice). Since $E$ contains the $5040$-th roots of
unity, $i=\sqrt{-1}\in E$, and $\sum_\tau\operatorname{Re}\tau(\alpha)=\frac12\operatorname{Tr}_{E/\Q}
\alpha=\frac12\operatorname{Tr}_{\Q(i)/\Q}(\operatorname{Tr}_{E/\Q(i)}\alpha)=\operatorname{Re}\beta\in
\Z$ for $\beta=\operatorname{Tr}_{E/\Q(i)}\alpha\in\Z[i]$. So the Gram matrix of a $\Z$-basis of
$L_{\mathrm{col}}$ is an integral positive definite matrix, its determinant is $\ge1$, and the
covolume, its square root, is $\ge1$.
\end{proof}

\begin{lemma}[Shifted points]\label{pbs:shifted}
Let $\mathcal B$, $p_0$, $l_p$ be as in Lemma~\ref{pbs:colour}, $u\in U^0_\infty$, $\mathsf x\in V_{\mathrm{
col}}$. Then for every $p\in\mathcal B$ and every $\tau$: $\tilde Z((uu_{\mathsf x})^{-1}_\tau p_\tau)=
\tilde Z(\gamma_\tau)$, norm-one pairs occur from $j$ for $(uu_{\mathsf x})_\tau^{-1}p_\tau$, and
\[
\Phi\bigl((uu_{\mathsf x})^{-1}_\tau p_\tau\bigr)_b=\Phi(u_\tau^{-1}p_{0,\tau})_b+\sum_{a<i}\tilde Z(\gamma_
\tau)_{ab}\,\frac M{s_a}\bigl(\tau(l_{p,a})+\mathsf x_{\tau,a}\bigr).
\]
\end{lemma}
\begin{proof}
Elements of $F_\tau$ do not change $\tilde Z$ ($\tilde Z(w^{-1}g)_{ab}=h(w\tilde f_a,g\tilde f_b)$) nor the
occurrence of norm-one pairs; and $\tilde Z(p_\tau)=\tilde Z(\gamma_\tau)$ by Lemma~\ref{pbs:move}. The
normalised components are $\Phi(p_\tau)_b=s_b\,\tau(\Phi(p)_b)$, and $\tilde Z_{ab}=s_as_b\tau(Z_{ab})$, so
by Lemma~\ref{pbs:colour} $\Phi(p_\tau)_b=\Phi(p_{0,\tau})_b+\sum_{a<i}\tilde Z_{ab}(M/s_a)\tau(l_{p,a})$. Apply
the shift formula for $u^{-1}$ (same $t_a(u)$ for $p$ and $p_0$, same $\tilde Z$) and then for $u_{\mathsf
x}^{-1}$.
\end{proof}

\begin{lemma}[Packing]\label{pbs:packing}
Let $Y\in G_\infty$ and let $\mathcal P\subseteq\Gamma_{c,s,c'}$ be a set such that $p^{-1}Y$ is boxed for
every $p\in\mathcal P$ and $pf_b-p'f_b\in A$ for all $p,p'\in\mathcal P$ and all $b$. Then $\sum_{p\in
\mathcal P}\beta_{c'}(p^{-1}Y)\le b^{12n^3\varrho f}\le\pbet$.
\end{lemma}
\begin{proof}
Group $\mathcal P$ by the frame $(pf_b)_b$. If $pf_b=p'f_b$ for all $b$, then by Step~2 of the proof of
Proposition~\ref{pbe:count} $p^{-1}p'\in U_{c'}$; so the elements with a given frame are of the form
$pu$ with distinct $u\in U_{c'}$, and their contribution is at most $\sum_{u\in U_{c'}}\beta_{c'}(u^{-1}
p^{-1}Y)=1$. To count the frames attach to $p$ the point $\mathsf F(p)=(Y_\tau^{-1}p\tilde f_b)_{\tau,b}$ of
$(\C^r)^{nf}$ with the norm $\max_{\tau,b}\|\cdot\|$. By Lemma~\ref{pbs:reduced}(c) (with $R=1$) the
points lie in the ball of radius $\pbsrb$, and by hypothesis in one translate of the real subspace
$\mathsf V=\bigoplus_\tau(Y_\tau^{-1}A_\tau)^n$, of real dimension at most $2\varrho nf$. If two
elements have different frames, their points are at distance at least $(N\pbsrb)^{-1}$: this is
Step~3 of the proof of Proposition~\ref{pbe:count}, with $y=Y$, $\rho=\pbsrb$, and $\pbel_\tau\le\pbsth$
in place of \eqref{pbe:eq-lenR}; the last inequality there becomes $b^{23\varrho+6n+2}(Nb^{37n})^{-1}<1$,
which holds. Comparing volumes in $\mathsf V$, the number of frames is at most $(1+4N\pbsrb^2)^{2\varrho
nf}\le(b^{6n^2})^{2\varrho nf}$, as $1+4N\pbsrb^2\le b^{134n+1}\le b^{6n^2}$. Finally $12n^3\varrho f\le
kf/4$.
\end{proof}

\subsubsection*{The kernel bound}
For a place $\tau$ let $\theta=\pbsth$, $\rho=\pbsrb$, $\delta_\tau=|\tau(\delta_j(\gamma))|$, and
\[
\kappa_\gamma(\tau)=\prod_{a<i}\Bigl|\frac{s_a}M\Bigr|^{2k}\cdot\frac{(2\rho^2)^{(i+1)i}(2^{i+1}\theta^{2(j-
i)})^k}{|\det B(\tilde Z(\gamma_\tau))|^{k+i}} .
\]

\begin{lemma}\label{pbs:kernel}
$\prod_\tau\kappa_\gamma(\tau)\le\pbet\,(b^{13}/M^2)^{ikf}\,\pbsmn_j(\gamma)^{-k/2}$.
\end{lemma}
\begin{proof}
By Lemma~\ref{pbs:common}(c) and $\tilde Z_{ab}=s_as_b\tau(Z_{ab})$, $|\det B(\tilde Z)|=w^2N^{-j}\delta_\tau$
with $w=|\kappa_\tau|^{\sum_{a<j}(\varrho-a)}\ge1$; and $\prod_{a<i}|s_a/M|=w/(M\sqrt N)^i$ (the index
$a=\varrho$, present when $j=n$, contributes $\varrho-a=0$). Hence
\[
\kappa_\gamma(\tau)=\Bigl(\frac{(2\rho^2)^{i+1}N^j}{w^2\delta_\tau}\Bigr)^i\Bigl(\frac{2^{i+1}\theta^{2(j-i)}
N^{j-i}}{M^{2i}\delta_\tau}\Bigr)^k\le\Bigl(\frac{(2\rho^2)^{i+1}N^j}{\delta_\tau}\Bigr)^i\Bigl(\frac{(b^{13}/
M^2)^iN^{(j-\varrho)_+}}{\delta_\tau}\Bigr)^k,
\]
using $w\ge1$, $j-i=(j-\varrho)_+$ and $2^{i+1}\theta^{2(j-i)}\le b^{13i}$ (for $j\le\varrho$ this is $2^{i+1}
\le b^{13i}$; for $j=n$, $i=\varrho$ and $2^{\varrho+1}b^{12n+4}\le b^{13\varrho}$). Let $P=\prod_\tau
\delta_\tau$. The entries of $Z(\gamma)$ are algebraic integers ($h(f_a,\gamma f_b)=h(f_a,k_\gamma f_b)$ with
$k_\gamma\in K$), so $\delta_j(\gamma)$ is a non-zero algebraic integer and $P\ge1$; also $P^2=|\auNm\delta_j|$,
so $P/N^{(j-\varrho)_+f}=\pbsmn_j(\gamma)^{1/2}$. The product of the first factors is at most $((2\rho^2)^{
i+1}N^j)^{fi}\le b^{135n^2f\varrho}\le\pbet$ (as $2\rho^2\le b^{74n+1}$, $N\le b^{60n}$, $135n^2\varrho\le k/4$),
and the product of the second factors is $((b^{13}/M^2)^{if}N^{(j-\varrho)_+f}/P)^k=(b^{13}/M^2)^{ikf}
\pbsmn_j^{-k/2}$.
\end{proof}

\begin{proposition}[Sum of Schur integrals over an orbit]\label{pbs:orbit-sum}
$\displaystyle\sum_{\delta\in\pbsO}I^{(i)}_{c,c'}(\delta)\le\nu\,d_\infty\,\pbet^{\,2}\,(b^{13}/M^2)^{ikf}
\bigl(i_\gamma\,\pbsmn_j(\gamma)^{-1/2}\bigr)^k$.
\end{proposition}
\begin{proof}
Let $\Phi_\pbsO(y)=\sum_{p\in\pbsO}\beta^\square_{c'}(p^{-1}y)\pbsv_y(p)^{-(k+i)}$, so that the left side is
$\int\beta^\square_c\|W^\flat\|^2\Phi_\pbsO$. It is left invariant under $U^0_c$: for $w\in U^0_c$,
$\pbsv_{wy}(p)=\pbsv_y(w^{-1}p)$ by Lemma~\ref{pbs:v-props}(c), $(w^{-1}p)^{-1}y=p^{-1}wy$, and $p\mapsto
w^{-1}p$ permutes $\pbsO$. We apply Lemma~\ref{pbs:stages} with $X=V_{\mathrm{col}}$, $t(\mathsf x)=u_{
\mathsf x}$, $B=\mathsf D$ (Lemma~\ref{pbs:col}) and $C=i_\gamma^k\,\pbet\prod_\tau\kappa_\gamma(\tau)$; by
Lemma~\ref{pbs:kernel} this gives the proposition. Fix $u\in U^0_\infty$ and a boxed $y$; we must show
$\int_{\mathsf D}\Phi_\pbsO(uu_{\mathsf x}y)\,d\mathsf x\le\pbvol(\mathsf D)\,C$; as $\pbvol(\mathsf D)\ge1$ it
suffices that the integral is at most $C$, and by Lemma~\ref{pbs:colour} that $\int_{\mathsf D}\sum_{p\in
\mathcal B}(\cdots)\le\pbet\prod_\tau\kappa_\gamma(\tau)$ for each colour class $\mathcal B$ (with base point
$p_0$ and the $l_p$).

At each $\tau$ apply Theorem~\ref{pbs:oneplace} to the point $y_\tau$ (boxed: $\|y_\tau^{-1}\tilde f_a\|\le
\rho$ by Lemma~\ref{pbs:reduced}(c), $\pbel_\tau\le\theta$), $g_0=u_\tau^{-1}p_{0,\tau}$ and $Z=\tilde Z(
\gamma_\tau)$ ($\det B(Z)\ne0$ by Lemma~\ref{pbs:common}(c) and $\delta_j\ne0$; $k\ge7$). Let $\varphi_\tau$ be
the function obtained, and $\varphi(\mathsf x)=\prod_\tau\varphi_\tau\bigl(((M/s_a)\mathsf x_{\tau,a})_a\bigr)$
on $V_{\mathrm{col}}$; by the substitution $a_a=(M/s_a)\mathsf x_{\tau,a}$, $\int_{V_{\mathrm{col}}}\varphi\le
\prod_\tau\kappa_\gamma(\tau)$. By Lemma~\ref{pbs:v-props}(c), $\pbsv_{uu_{\mathsf x}y}(p)=\pbsv_y((uu_{\mathsf
x})^{-1}p)$, and by Lemma~\ref{pbs:shifted} and Theorem~\ref{pbs:oneplace},
\[
\pbsv_{uu_{\mathsf x}y}(p)^{-(k+i)}\le\varphi(\lambda(l_p)+\mathsf x),\qquad\lambda(l)=(\tau(l_a))_{\tau,a}
\in L_{\mathrm{col}} .
\]
Group $\mathcal B$ by $l_p$. For fixed $l$, the $p$ with $l_p=l$ have equal $\Phi(p)$, equal $Z(p)$ and
(Lemma~\ref{pbs:move}) the same components along $Ez$ and $A^*$ (these are read off from $Z$), so
$pf_b-p'f_b\in A$; by Lemma~\ref{pbs:packing} with $Y=uu_{\mathsf x}y$, $\sum_{l_p=l}\beta^\square_{c'}(p^{-1}
Y)\le\pbet$. Hence $\sum_{p\in\mathcal B}\beta^\square_{c'}(p^{-1}Y)\pbsv_Y(p)^{-(k+i)}\le\pbet\sum_{w\in
L_{\mathrm{col}}}\varphi(w+\mathsf x)$, and $\int_{\mathsf D}\sum_w\varphi(w+\mathsf x)\,d\mathsf x=\int_{V_{
\mathrm{col}}}\varphi$.
\end{proof}

\begin{proof}[Proof of Theorem~\ref{pbs:orbit}]
By Lemma~\ref{pbs:cs} and the Cauchy--Schwarz inequality for sums,
\[
\Bigl(\sum_{\delta\in\pbsO}\mathsf a_{c,c'}(\delta)\Bigr)^2\le\sum_{\delta\in\pbsO}I^{(i)}_{c,c'}(\delta)\cdot
\sum_{\delta\in\pbsO}I^{(i)}_{c',c}(\delta^{-1}).
\]
The elements $\delta^{-1}$ form the orbit $U^0_{c'}\gamma^{-1}U^0_c$ of $\gamma^{-1}\in\Gamma_{c',s^{-1},c}$.
Since $Z(\gamma^{-1})=Z(\gamma)^*$, $\gamma^{-1}$ has the same $\pbsmn_j$; it has cell rank $j$ (a common
vector $fx=\gamma fy$ gives $fy=\gamma^{-1}fx$, and $(1,\zeta)$ is proportional to $(\aubar\zeta,1)$); and
$i_{\gamma^{-1}}=i'_\gamma$ as $k_{\gamma^{-1}}=k_\gamma^{-1}$. Proposition~\ref{pbs:orbit-sum} for both orbits
gives
\[
\Bigl(\sum_{\delta\in\pbsO}\mathsf a_{c,c'}(\delta)\Bigr)^2\le\bigl(\nu d_\infty\pbet^2(b^{13}/M^2)^{ikf}
\bigr)^2\Bigl(\frac{i_\gamma i'_\gamma}{\pbsmn_j(\gamma)}\Bigr)^k=\bigl(\nu d_\infty\pbet^2(b^{13}/M^2)^{ikf}
\bigr)^2\Bigl(\frac{N^f}{\mathsf G_j(\gamma)}\Bigr)^k,
\]
because $\auNm\pbg\le N^f\pbsmn_j$ gives $\mathsf G_j(\gamma)=N^f\pbsmn_j/(i_\gamma i'_\gamma)$. It remains to
show that the right side is at most $\nu^2\mathsf S_j(\mathsf G_j)^2=\nu^2\pbet^4\mathcal D_j/(\pbenp M^{2ifk}
\mathsf G_j^k)$, i.e.
\[
d_\infty^2\,(b^{13}/M^2)^{2ikf}\,N^{fk}\le\mathcal D_j/(\pbenp M^{2ifk}).
\]
We use: $d_\infty\le b^{2kf}=\pbet^8$; $b^{13}/M^2\le(b^4M)^{-1}$ (as $b^{17}\le M$); $\mathcal D_j\ge\pbet^{8-
32i}\pbDid$ (for $j\le\varrho$, $\mathcal D_j=\pbet^{-16}\pbDid$ and $i\ge1$; for $j=n$, $\pbDmain\ge b^{11kf}
\pbDid=\pbet^{44}\pbDid$); $\pbDid=\#\pbR\cdot\#\mathcal C\cdot d_\infty$ and $\pbenp=\#\pbR\cdot\#\mathcal C\,
N^{-kf}$. Then the left side is at most $d_\infty\pbet^8b^{-8ikf}M^{-2ikf}N^{fk}$ and the right side at
least $\pbet^{8-32i}d_\infty N^{kf}M^{-2ifk}$; these are equal since $b^{8ikf}=\pbet^{32i}$.
\end{proof}

\subsection{The open cell}\label{pbs:ss-open}

For $\gamma\in\Gamma$ and a place $\tau$ let $\sigma_\tau=\prod_{a\le\varrho}s_a^2$ and $\pbg_\tau(\gamma)=
|\det\tilde Z(\gamma_\tau)|^2-|\tilde D_\varrho|^2$, $\tilde D_\varrho$ the leading minor of order $\varrho$ of
$\tilde Z(\gamma_\tau)$. As $\tilde Z_{ab}=s_as_b\tau(Z_{ab})$ and $s_\varrho^2=1/N$, $\pbg_\tau(\gamma)=\sigma_
\tau^2\,\tau(\pbg(\gamma))$. Let $\mathrm{vol}\,\auU(l)=\prod_{i<l}2\pi^{i+1}/i!$ and
\[
c_G=\frac{\mathrm{vol}\,\auU(k)}{\mathrm{vol}\,\auU(r)}\cdot\frac{2\pi^k}{(k-1)!}=\frac{2\pi^k}{(k-1)!}
\prod_{i=k}^{r-1}\frac{i!}{2\pi^{i+1}}
\]
and $\pbdeg^{\mathrm B}_\tau:=\pbdeg_\tau$, the formal
degree of \S\ref{pbe:ss-ds}. Put
\[
\pbsp_\tau(\gamma)=\Bigl(\frac{\pbsrb^{9n}}{\pbg_\tau(\gamma)}\Bigr)^{10^5n^4}\frac{d_\tau^2}{\pbdeg^{
\mathrm B}_\tau\,c_G\,(\pbg_\tau(\gamma)/2)^{k-1}} .
\]
We use the following statement, proved in \S\ref{pbo:sec} (Theorem~\ref{pbo:O1}).

\begin{itemize}
\item[(O1)] \emph{(Open-cell pairing.)} Let $\gamma\in\Gamma$ with $\pbg_\tau(\gamma)>0$ for all $\tau$, and
suppose some $y_0$ has $y_0$ and $\gamma^{-1}y_0$ boxed. Let $\phi\colon G_\infty\to[0,1]$ be measurable,
left invariant under $\prod_\tau F_\tau$, and such that $\phi(y)\ne0$ and $\pbel_\tau(y_\tau)\le\pbsth$
imply $\|y_\tau^{-1}\tilde f_a\|\le\pbsrb$ for all $a$. Then, for $W,W^\flat$ of Theorem~\ref{pbe:bessel},
\[
\Bigl\|\int_{G_\infty}\phi\,\pbsK^\flat_\gamma\Bigr\|\le\Bigl(\prod_\tau\pbsp_\tau(\gamma)\Bigr)^{1/2}+\int_{
G_\infty}\phi\,\|\pbsK_\gamma-\pbsK^\flat_\gamma\| .
\]
\end{itemize}

The second input is proved here. Let $\pbvol(\widehat\auO)=(2/\mathrm{rd}_E)^f$ (the normalisation of
Remark~\ref{pbs:rem-nu}), $c_S=(\pi^k/((k-1)!N^k))^f|D_E|^{1/4}\pbvol(\widehat\auO)^k$, $\#\pbUb=M^{f
\varrho(2k+2\varrho+1)}$, $a_*=\#\pbR\cdot\#\pbUb\cdot\#\mathcal C\cdot c_S$, $\pbdeg^{\mathrm B}=\prod_\tau
\pbdeg^{\mathrm B}_\tau=\pbdeg$, and $\Lambda$, $\nu_{\mathrm{src}}$ as in Remark~\ref{pbs:rem-nu}.

\begin{lemma}[(O2): the Siegel constant]\label{pbs:O2}
$a_*\le(1+b^{-9})\,\pbDmain\,\pbdeg\,c_G^{\,f}\,\pbvol(\widehat\auO)^{-n}\,\bigl(\nu_{\mathrm{src}}\Lambda/
d_\infty\bigr)^2$.
\end{lemma}
\begin{proof}
Write $\mathrm{rd}=\mathrm{rd}_E$, so $|D_E|=\mathrm{rd}^{2f}$ and $\pbvol(\widehat\auO)=(2/\mathrm{rd})^f$.
\emph{Step 1 (frame mass).} Put $\mathsf{fr}=N^{kf}\prod_{i=k}^{r-1}\bigl((2\pi)^{i+1}/(\mathrm{rd}^{\,i+1/2}
\,i!)\bigr)^f$. We claim
\begin{equation}\label{pbs:eq-fr}
c_S=\frac{\mathsf{fr}\cdot c_G^{\,f}\,|D_E|^{\varrho/2}}{N^{2kf}\,\pbvol(\widehat\auO)^{\varrho+1+2\varrho(
\varrho+k)}} .
\end{equation}
Both sides are positive; compare logarithms divided by $f$. With $\sum_{i=k}^{r-1}i=(2\varrho+1)(k+\varrho)$
and $\sum_{i=k}^{r-1}(i+\frac12)=(2\varrho+1)(k+\varrho+\frac12)$, the logarithm of $\mathsf{fr}\,c_G^f
|D_E|^{\varrho/2}$ divided by $f$ is
\[
k\log N-\log(k-1)!+k\log\pi+\bigl(1+(2\varrho+1)(k+\varrho)\bigr)\log2-\bigl((2\varrho+1)(k+\varrho+\tfrac12)-\varrho
\bigr)\log\mathrm{rd}
\]
(the factorials $i!$ and the powers $\pi^{i+1}$ cancel between $\mathsf{fr}$ and $c_G$), and that of
$c_SN^{2kf}\pbvol(\widehat\auO)^{\varrho+1+2\varrho(\varrho+k)}$ is
\[
k\log N+k\log\pi-\log(k-1)!+(k+\varrho+1+2\varrho(\varrho+k))\log2-(k-\tfrac12+\varrho+1+2\varrho(\varrho+k))
\log\mathrm{rd};
\]
the coefficients agree.

\emph{Step 2 (covolume).} From the definition, with $d'=\varrho(\varrho+k)$ and $M^{(2d'+\varrho)f}=\#\pbUb$,
$\nu_{\mathrm{src}}\Lambda=\#\pbUb\,\allowbreak N^{-kf}\pbvol(\widehat\auO)^{-d'}|D_E|^{\varrho/4}$. With
\eqref{pbs:eq-fr} and $\#\pbR=\#\pbUb\cdot\#\pbHb$ \eqref{pb:eq-orders},
\[
a_*=\#\pbHb\,\#\pbUb^{\,2}\,\#\mathcal C\,c_S=\#\pbHb\cdot\#\mathcal C\cdot\mathsf{fr}\cdot c_G^{\,f}\,
\pbvol(\widehat\auO)^{-n}\,(\nu_{\mathrm{src}}\Lambda)^2 ,
\]
as $n=\varrho+1$. Since $\pbDmain=(\#\mathcal C\,d_\infty)^2/(\pbdeg\,\pbmass_r)$ (\eqref{pb:eq-D},
\eqref{pb:eq-DX}), the lemma is equivalent to $\#\pbHb\cdot\mathsf{fr}\le(1+b^{-9})\,\#\mathcal C/\pbmass_r$, and
by $\#\mathcal C=\#\pbHb\,(N^k-1)^f\pbmass_k$ (proof of Lemma~\ref{pb:sizes}) to
\begin{equation}\label{pbs:eq-fr2}
\mathsf{fr}\cdot\pbmass_r\le(1+b^{-9})\,(N^k-1)^f\,\pbmass_k .
\end{equation}

\emph{Step 3 (masses).} By Lemma~\ref{pb:massratio}, for $k\le i<r$, $\pbmass_{i+1}/\pbmass_i=\bigl(\mathrm
{rd}^{\,i+1/2}i!/(2\pi)^{i+1}\bigr)^fL_i$ with $|\log L_i|\le f2^{1-i}$. Telescoping, $\mathsf{fr}\,\pbmass_r=
N^{kf}\pbmass_k\prod_iL_i$, and $\sum_i|\log L_i|\le f2^{2-k}$. Also $\log(N^k/(N^k-1))\le1/(N^k-1)\le2N^{-k}
\le2^{1-k}$. So \eqref{pbs:eq-fr2} follows from $f(2^{2-k}+2^{1-k})\le\log(1+b^{-9})$, and $\log(1+b^{-9})\ge
b^{-9}/2$. By Lemma~\ref{an:lem-sizes}, $f\le x^{12}$ (S3), $k\ge r/2$ (S2) and $r\ge x^5$ (Definition~\ref{an:def-data}), so $k\ge x^5/2$, and $\log b\le x^{1/100}$
(S1, as $(\log b)^{100}\le x$); hence $f\,2^{3-k}\le8x^{12}2^{-x^5/2}\le\frac12e^{-9x^{1/100}}\le b^{-9}/2$.
\end{proof}

\begin{lemma}[Arithmetic of the open cell]\label{pbs:open-arith}
Assume $\nu=\nu_{\mathrm{src}}$ (Lemma~\ref{pb:covol}). Let $\gamma\in\Gamma_{c,s,c'}$ be near with $\pbg(\gamma)
\ne0$. Then $\auNm\pbg(\gamma)\ge1$, $\pbg_\tau(\gamma)>0$ for all $\tau$,
\[
\Bigl(\prod_\tau\pbsp_\tau(\gamma)\Bigr)^{1/2}\le\frac\nu2\,\mathsf S_n\bigl(\auNm\pbg(\gamma)\bigr),
\qquad\frac\nu2\,b^{-k^3f}\le\frac\nu2\,\mathsf S_n\bigl(\auNm\pbg(\gamma)\bigr).
\]
\end{lemma}
\begin{proof}
$\pbg(\gamma)\in E^+$ is a non-zero algebraic integer (the entries of $Z(\gamma)$ are integral), so
$\auNm\pbg=|\mathrm N_{E^+/\Q}\pbg|\ge1$. Fix $\tau$. By the proof of Lemma~\ref{pbe:link}(a), the Gram
matrix of the $2n$ vectors $f_0,\dots,f_\varrho,\gamma f_0,\dots,\gamma f_\varrho$ is $\bigl(\begin{smallmatrix}
\mathsf D&Z^*\\Z&\mathsf D\end{smallmatrix}\bigr)$, and exchanging the two block columns (sign $(-1)^{n^2}
=(-1)^n$) and applying the determinant formula there shows that its determinant is $(-1)^n\pbg(\gamma)$.
As $\pbg\ne0$ the $2n$ vectors span a non-degenerate subspace $\mathsf V_\tau$ of $Y_\tau$ of dimension $2n$.
It contains $N_0$, of dimension $n$, on which $h\le0$; in a non-degenerate space a subspace on which
$h\le0$ has dimension at most the negative index, so the negative index of $\mathsf V_\tau$ is $\ge n$,
hence $=n$ (that of $Y_\tau$). The determinant of a hermitian Gram matrix of a basis of $\mathsf V_\tau$
has the sign $(-1)^n$. Hence $\tau(\pbg(\gamma))>0$ and $\pbg_\tau(\gamma)=\sigma_\tau^2\tau(\pbg)>0$.

\emph{Second inequality.} $\mathsf S_n(x)=\pbet^2(\pbDmain/(\pbenp M^{2\varrho fk}x^k))^{1/2}\ge(M^{2\varrho fk}x^k)
^{-1/2}$, using $\pbet\ge1$ and $\pbenp\le\pbDid\le\pbDmain$. By Lemma~\ref{ppi:gauge}, $x=\auNm\pbg\le
b^{(121n+270n^2)f}\le b^{271n^2f}$, and $M\le b^{22}$; so $M^{2\varrho fk}x^k\le b^{(44\varrho+271n^2)kf}\le
b^{2k^3f}$.

\emph{First inequality.} We have $\prod_\tau\pbsp_\tau=C_{\mathrm{open}}/(\auNm\pbg)^{10^5n^4+k-1}\le C_{
\mathrm{open}}/(\auNm\pbg)^k$ with $C_{\mathrm{open}}=\prod_\tau(\pbsrb^{9n}/\sigma_\tau^2)^{10^5n^4}\allowbreak d_\tau^2/
(\pbdeg^{\mathrm B}_\tau c_G(\sigma_\tau^2/2)^{k-1})$. Each $s_a^2=\kappa_\tau^{2(\varrho-a)}/N$ lies in $[b^{-60n},
b^{46n}]$, so $b^{-60n^2}\le\sigma_\tau\le b^{46n^2}$ and $\Lambda_\tau:=\prod_as_a^{2(k+2(\varrho-a))}\le\sigma_\tau^k
b^{92n^3}$. Hence the factor of $C_{\mathrm{open}}$ at $\tau$ is at most
\[
b^{453n^2\cdot10^5n^4}\cdot2^k\sigma_\tau^2\frac{d_\tau^2}{\pbdeg^{\mathrm B}_\tau c_G\sigma_\tau^{2k}}\le2^kb^{
453\cdot10^5n^6+92n^2+184n^3}\frac{d_\tau^2}{\pbdeg^{\mathrm B}_\tau c_G\Lambda_\tau^2},
\]
and, as $453\cdot10^5n^6+92n^2+184n^3\le10^5n^9\le k/4$, $C_{\mathrm{open}}\le2^{kf}\pbet\,d_\infty^2/(\pbdeg^{
\mathrm B}c_G^f\Lambda^2)$. By Lemma~\ref{pbs:O2}, $d_\infty^2/(\pbdeg^{\mathrm B}c_G^f\Lambda^2)\le(1+b^{-9})\pbDmain\nu^2/(a_*
\pbvol(\widehat\auO)^n)$. Next, $2^{kf}\pbenp M^{2\varrho fk}\le\pbet\,a_*\pbvol(\widehat\auO)^n$: the factors
$\#\pbR\#\mathcal C$ and $N^{-kf}$ cancel, $M^{2\varrho fk}\le M^{f\varrho(2k+2\varrho+1)}$, and with
$x_0=200fk\,\mathrm{rd}_E$ we have $2^{kf}\le x_0^{kf}$, $(k-1)!\le x_0^k$, $\pi^k\ge1$, $|D_E|\ge1$ and
$\pbvol(\widehat\auO)=(2/\mathrm{rd}_E)^f\ge x_0^{-f}$, so it suffices that $x_0^{kf}\cdot x_0^{kf}\cdot
x_0^{(k+n)f}\le\pbet$, which follows from $x_0^{1000}\le b$ (size condition (Z3)). Altogether
\[
C_{\mathrm{open}}\le\pbet^2(1+b^{-9})\nu^2\frac{\pbDmain}{\pbenp M^{2\varrho fk}}\le\Bigl(\frac{\nu\pbet^2}2
\Bigr)^2\frac{\pbDmain}{\pbenp M^{2\varrho fk}},
\]
as $4(1+b^{-9})\le\pbet^2$; dividing by $(\auNm\pbg)^k$ gives $(\nu\mathsf S_n(\auNm\pbg)/2)^2$.
\end{proof}

\begin{lemma}[Trivial stabilisers]\label{pbs:stab}
If $\pbg(\gamma)\ne0$, $w\in U(E^+)$, $p\in U^0(E^+)$ and $p\gamma=\gamma w$, then $w=1$ (and $p=1$).
\end{lemma}
\begin{proof}
$w$ fixes the vectors $\gamma^{-1}f_b$, since $w\gamma^{-1}f_b=\gamma^{-1}pf_b=\gamma^{-1}f_b$. By
Proposition~\ref{ppi:cell}(a), $\det Z(\gamma)\ne0$, so the $\varrho\times n$ matrix $(h(e_a,\gamma^{-1}f_b))$
has rank $\varrho$ and there are $v_1,\dots,v_\varrho$ in the span of the $\gamma^{-1}f_b$ with $h(e_a,v_b)=
\delta_{ab}$; $w$ fixes them. Now $w\in U$ acts on $A$ by a unitriangular matrix $C$ and on $Y_0/A^\perp
\cong A^*$ by $(C^*)^{-1}$; the $A^*$-components of $wv_b$ and $v_b$ agree, so $C=1$ and $w$ fixes $A$.
Then $we^*_b=e^*_b+q_b+\alpha_b$ ($q_b\in X_0\oplus Ez$, $\alpha_b\in A$) and $w-1$ maps $A^\perp$ into $A$;
comparing $X_0\oplus Ez$-components of $wv_b=v_b$ gives $q_b=0$. For $x\in X_0\oplus Ez$, $wx=x+\sum_a
\lambda_a(x)e_a$ and $0=h(x,e^*_b)=h(wx,we^*_b)=h(x+\sum\lambda_a(x)e_a,e^*_b+\alpha_b)$, which equals
$\aubar{\lambda_b(x)}$ since $h(x,\alpha_b)=h(e_a,\alpha_b)=0$ and $h(e_a,e^*_b)=\delta_{ab}$; so $w$ is trivial on $X_0\oplus Ez$, and then
$wv_b=v_b$ gives $\alpha_b=0$. So $w=1$, and $p=\gamma w\gamma^{-1}=1$.
\end{proof}

\begin{proof}[Proof of Proposition~\ref{pbs:open-block}]
Let $X=T\cap\Gamma_{c,s,c'}$; $\pbg\ne0$ on $X$. In each double coset $U^0_c\delta U_{c'}$ ($\delta\in X$)
choose a representative $\mathsf r(\delta)$, near and chosen if the double coset contains a near and chosen
element. Let $\phi(\delta,y)=\mathbf 1[\delta=\mathsf r(\delta)]\,\omega(y)$. \emph{$\phi$ is admissible in
Lemma~\ref{pbs:weight}.} For $\delta\in X$ and $x\in U_c$ let $\mathsf n(x)=\#\{u'\in U_{c'}:x\delta u'^{-1}=
\mathsf r(x\delta)\}$. For each coset $U^0_cx$, $\sum_{v\in U^0_c}\mathsf n(vx)$ is the number of $(v,u')\in
U^0_c\times U_{c'}$ with $vx\delta u'^{-1}=\mathsf r(x\delta)$; there is at least one, and if $(v,u')$,
$(v_1,u'_1)$ are two, then $(v_1^{-1}v)(x\delta)=(x\delta)(u_1'^{-1}u')$, so $u'=u'_1$ and $v=v_1$ by
Lemma~\ref{pbs:stab} ($\pbg(x\delta)=\pbg(\delta)\ne0$). By Lemma~\ref{pbs:reduced}(b),
$\sum_{(x,u')}\phi(x\delta u'^{-1},xy)=\sum_x\mathsf n(x)\omega(xy)=1$.

Hence $A^\flat_T(c,s,c')=\sum_{\delta\in X,\ \delta=\mathsf r(\delta)}\int\omega\,\pbsK^\flat_\delta$. If $\delta=
\mathsf r(\delta)$ is not near and chosen, then $\omega\pbsK^\flat_\delta=0$ identically: otherwise some reduced
$y$ has $W^\flat(y)\ne0\ne W^\flat(\delta^{-1}y)$, so $y$ is boxed and $\pbel_\tau\le\pbsth$ at $\delta^{-1}y$;
choose $u\in U_{c'}$ with $\beta_{c'}(u\delta^{-1}y)\ne0$ \eqref{pbs:eq-cutoff}; then $u\delta^{-1}y$ is boxed,
$\delta u^{-1}$ is near, and the chosen element of $U^0_c\delta u^{-1}U^0_{c'}$ is near
(Lemma~\ref{pbs:near}) and lies in $U^0_c\delta U_{c'}$, contradicting the choice of $\mathsf r(\delta)$. If
$\delta$ is near and chosen, (O1) applies with $\phi=\omega$ (Lemma~\ref{pbs:reduced}(a),(c) and
Lemma~\ref{pbs:open-arith}). Summing, with Lemma~\ref{pbs:weight} for the error terms and
Lemma~\ref{pbs:open-arith} for the main terms,
\[
\|A^\flat_T(c,s,c')\|\le\sum_{\delta}\frac\nu2\mathsf S_n(\auNm\pbg(\delta))+\frac\nu2b^{-k^3f},
\]
the sum over the near and chosen $\delta\in X$. If there is no such $\delta$, the block is $0$; otherwise
the last term is at most $\frac\nu2\mathsf S_n(\auNm\pbg(\delta_0))$ for any one of them, and the
proposition follows.
\end{proof}

\subsection{Size conditions}\label{pbs:ss-sizes}

The proof uses the following conditions on the data, besides those of \S\ref{pbe:sec}:
\begin{itemize}
\item[(Z1)] $\varrho\ge10^5$ and $k\ge1000(2n)^9$ \eqref{ppi:eq-sizes}. Used for: $k\ge7$;
$12n^3\varrho\le k/4$ (Lemma~\ref{pbs:packing}); $135n^2\varrho\le k/4$ (Lemma~\ref{pbs:kernel});
$10^5n^9\le k/4$ (Lemma~\ref{pbs:open-arith}); $44\varrho+271n^2\le2k^2$; $2^{\varrho+1}b^{12n+4}\le
b^{13\varrho}$.
\item[(Z2)] From \S\ref{pbe:sec} and Lemma~\ref{pb:sizes}: $b\ge1000$; $b^{20}\le M\le b^{22}$ (we use
$b^{17}\le M$); $N\le b^{60n}$; $1\le|\kappa_\tau|\le b^{23}$; $d_\infty\le b^{2kf}$; $b^{11kf}\pbDid\le
\pbDmain$; and $\pbDid=\#\pbR\cdot\#\mathcal C\cdot d_\infty$ (the dimension of the space of source forms
is $\#\mathcal C\cdot d_\infty$).
\item[(Z3)] $(200\,fk\,\mathrm{rd}_E)^{1000}\le b$ (Lemma~\ref{pbs:open-arith}). It holds for the data: this is
Lemma~\ref{an:lem-sizes}(S7)(d).
\end{itemize}

\subsection{Proof of Theorem~\ref{pbs:S}}\label{pbs:ss-proof}

If $s$ does not agree, $A^\flat_T(c,s,c')=0$ by Lemma~\ref{pbs:agree}, and the sum is empty by
Definition~\ref{ppi:def-N}. Let $s$ agree. By Lemma~\ref{pbs:cells}, $T$ is the disjoint union of the
stable sets $T_j=\{\gamma\in T:j(\gamma)=j\}$, $1\le j\le n$, all of whose elements have cell rank $j$, and
$A^\flat_T=\sum_jA^\flat_{T_j}$. By Proposition~\ref{pbs:rank},
\[
\|A^\flat_T(c,s,c')\|\le\sum_j\nu\sum_{\gamma}\mathsf S_j(\mathsf G_j(\gamma))=\nu\sum_\gamma\pbesz(\gamma),
\]
the inner sums over the near and chosen $\gamma\in T_j\cap\Gamma_{c,s,c'}$; together these are the elements
of $\pbeN(c,c')\cap T\cap\Gamma_{c,s,c'}$ (they lie in $\pbeGs$, are near and chosen, and $s$ agrees), and
$\mathsf S_{j(\gamma)}(\mathsf G_{j(\gamma)}(\gamma))=\pbesz(\gamma)$ by \S\ref{pbs:ss-statement}. \qed


%% ===== proofs/B-opencell.tex =====
% proofs/B-opencell.tex
% Statement (O1) of proofs/B-size.tex (the open-cell pairing bound for the cut Bessel functions).
% Labels of this file start with pbo:.
\providecommand{\auO}{\mathcal O}
\providecommand{\aubar}[1]{\overline{#1}}
\providecommand{\pbeq}{\mathsf q}
\providecommand{\pbvol}{\operatorname{vol}}
\providecommand{\pbdeg}{\operatorname{deg}}
\providecommand{\pbg}{\mathsf g}
\providecommand{\pbel}{\ell}
\providecommand{\pbsrb}{\rho_\square}
\providecommand{\pbsth}{\theta_\square}
\providecommand{\pbsK}{\mathsf K}
\providecommand{\pbsp}{\mathsf p}
% macros of this file (prefix \pbo)
\providecommand{\pboZ}{\mathsf Z}
\providecommand{\pboF}{\mathsf F}
\providecommand{\pboFd}{\mathsf F'}
\providecommand{\pboO}{\Omega}
\providecommand{\pbonu}{\nu}
\providecommand{\pboU}{U^0}
\providecommand{\pboR}{R^0}
\providecommand{\pbocG}{c_G}
\providecommand{\pboV}{\operatorname{vol}\mathrm U}
\providecommand{\pboq}{\mathsf q}
\providecommand{\pboS}{\mathcal S}
\providecommand{\pboB}{\mathsf B}
\providecommand{\pbod}{\mathsf d}

\section{The open-cell pairing}\label{pbo:sec}

This section proves statement \textup{(O1)} of \S\ref{pbs:ss-open}, the bound for the pairing of
the cut Bessel functions over an element $\gamma$ whose frame matrix lies in the open cell. We use
the notation of \S\S\ref{pbe:sec} and~\ref{pbs:sec}: at a real place $\tau$ the group $G_\tau$, the
normalised basis $\tilde e_a,x_i,\tilde z,\tilde e^*_a$ with Gram matrix $\mathsf J_1$, the
Euclidean norm for which it is orthonormal, the operator norm $\|g\|$, the normalised frame
$\tilde f_0,\dots,\tilde f_\varrho$, the matrix $\tilde Z(g)=(h(\tilde f_a,g\tilde f_b))_{a,b}$, the
groups $U_\tau$, $H_\tau=\mathrm U(X_{0,\tau})$, $R_\tau=U_\tau H_\tau$, the length $\pbel_\tau$, the
Bessel functions $W_\tau$, $W^\flat_\tau=\chi_\tau W_\tau$ of Proposition~\ref{pbe:local}, the formal
degree $\pbdeg_\tau$ and $d_\tau=\dim\mathbb W_\tau$. As in \S\ref{pbs:sec}, $n=\varrho+1$,
$r=k+2\varrho+1$, $\pbsrb=b^{37n}$, $\pbsth=b^{6n+2}$.

The exact statement is Theorem~\ref{pbo:O1}; it is \textup{(O1)} of \S\ref{pbs:ss-open} with the two
constants made explicit (the degree $\pbdeg^{\mathrm B}_\tau$ is the formal degree $\pbdeg_\tau$ of
\S\ref{pbe:ss-ds}, and $c_G$ is the number \eqref{pbo:eq-cG}), and with $\phi$ allowed complex valued,
$|\phi|\le1$. It is deduced from properties \textup{(BS1)--(BS5)} (Definition~\ref{pbo:def-system}), which
the Bessel functions satisfy by Theorem~\ref{pbo:system}; the latter uses Schmid's theorem on the lowest
$K$-type, Harish-Chandra's value of the Casimir operator, and the existence of $\beta\in C^\infty_c(G)$
with $\pi(\beta)v_0=v_0$. The normalisation of Haar measure as a Gram fibre measure (Lemma~\ref{pbo:GF})
uses Remark~\ref{pds:cayley}.

\subsection{The statement}\label{pbo:ss-statement}

Fix a real place and drop $\tau$ from the notation until \S\ref{pbo:ss-proof}. Let $G=G_\tau$, and
let $\pboF$, $\pboFd$ be the $r\times n$ matrices whose columns are $(\tilde e_1,\dots,\tilde
e_\varrho,\tilde z)$ and $(\tilde e^*_1,\dots,\tilde e^*_\varrho,x_{k_0})$, for a fixed index $k_0$.
For an $n\times n$ complex matrix $Z$ let $D_\varrho(Z)$ be its leading minor of order $\varrho$ and
\[
\pbg(Z)=|\det Z|^2-|D_\varrho(Z)|^2,\qquad \xi(Z)=(Z^{-1})_{\varrho\varrho},\qquad
\pbonu(Z)^2=1-|\xi(Z)|^2 .
\]
The \emph{open cell} is $\pboO=\{Z:\det Z\ne0,\ |\xi(Z)|<1\}$. By Cramer's rule $\xi=D_\varrho/\det Z$,
so for invertible $Z$
\begin{equation}\label{pbo:eq-gnu}
\pbg(Z)=|\det Z|^2\,\pbonu(Z)^2,\qquad\text{hence}\qquad\pboO=\{Z:\pbg(Z)>0\}.
\end{equation}
For $g\in G$, $\tilde Z(g)=\pboF^*\mathsf J_1g\pboF$; we write $\pbg(g)=\pbg(\tilde Z(g))$, which is
$\pbg_\tau(\gamma)$ of \S\ref{pbs:ss-open} for $g=\gamma_\tau$. Let $\pboU=U^0_\tau$ be the elements of
$U_\tau$ fixing every $\tilde f_a$ and $\pboR=\pboU H$ the fixer of the frame in $R_\tau$ (it is the
stabiliser $F_\tau\cap R_\tau$; in fact $F_\tau=\pboR$, Lemma~\ref{pbo:witt}). Let
\[
\pboV(j)=\prod_{i=1}^{j}\frac{2\pi^i}{(i-1)!},\qquad
\pbocG=\frac{\pboV(k)}{\pboV(r)}\cdot\frac{2\pi^k}{(k-1)!}.
\tag{\theequation}\label{pbo:eq-cG}
\]
\stepcounter{equation}
For a matrix $Z$ with $\pbg(Z)>0$, numbers $\rho\ge1$, $D>0$, $d\ge1$ put
\[
\pbsp(Z;D,\rho,d)=\Bigl(\frac{\rho^{9n}}{\pbg(Z)}\Bigr)^{10^5n^4}\frac{d^2}{D\,\pbocG\,(\pbg(Z)/2)^{k-1}}.
\]

\begin{theorem}[Open-cell pairing]\label{pbo:O1}
Assume the size conditions of \S\ref{pbo:ss-sizes}, Theorem~\ref{pbe:ds}. Let $W=\bigotimes_\tau W_\tau$ and $W^\flat=\bigotimes_\tau
W^\flat_\tau$ be as in Theorem~\ref{pbe:bessel}, and for $\gamma\in G_\infty$ let $\pbsK_\gamma(y)=
W(y)W(\gamma^{-1}y)^*$, $\pbsK^\flat_\gamma(y)=W^\flat(y)W^\flat(\gamma^{-1}y)^*$ (Hilbert--Schmidt
norm). Let $\gamma=(\gamma_\tau)\in G_\infty$ with $\pbg(\gamma_\tau)>0$ for all $\tau$, and let $\rho_\tau\ge
\pbsth$ be numbers such that some $y_0$ has $\|y_{0,\tau}^{-1}\tilde f_a\|\le\rho_\tau$ and
$\|y_{0,\tau}^{-1}\gamma_\tau\tilde f_a\|\le\rho_\tau$ for all $\tau$ and $a$. Let $\phi\colon G_\infty\to
\C$ be measurable with $|\phi|\le1$, with $\phi(uy)=\phi(y)$ for $u\in\prod_\tau\pboR_\tau$, and such
that $\phi(y)\pbsK^\flat_\gamma(y)\ne0$ implies $\|y_\tau^{-1}\tilde f_a\|\le\rho_\tau$ for all $\tau,a$.
Then $\pbsK^\flat_\gamma$ is continuous, $\phi\pbsK^\flat_\gamma$ is bounded, measurable and vanishes outside a compact set (so it is integrable), and
\[
\Bigl\|\int_{G_\infty}\phi\,\pbsK^\flat_\gamma\Bigr\|\le\Bigl(\prod_\tau\pbsp\bigl(\tilde Z(\gamma_\tau);
\pbdeg_\tau,\rho_\tau,d_\tau\bigr)\Bigr)^{1/2}+\int_{G_\infty}|\phi|\,\|\pbsK_\gamma-\pbsK^\flat_\gamma\| .
\]
\end{theorem}

\begin{remark}[Relation with \textup{(O1)} of \S\ref{pbs:ss-open}]\label{pbo:rem-O1}
\textup{(O1)} is the case $\rho_\tau=\pbsrb$, $\phi$ with values in $[0,1]$, with $\pbdeg^{\mathrm B}_\tau=
\pbdeg_\tau$ and $c_G$ as in \eqref{pbo:eq-cG}. Its hypotheses imply ours: if $y_0$ and $\gamma^{-1}y_0$
are boxed, Lemma~\ref{pbs:reduced}(c) with $R=1$ gives $\|y_{0,\tau}^{-1}\tilde f_a\|\le\pbsrb$ and
$\|(\gamma_\tau^{-1}y_{0,\tau})^{-1}\tilde f_a\|\le\pbsrb$; and if $\phi(y)\pbsK^\flat_\gamma(y)\ne0$ then
$W^\flat_\tau(y_\tau)\ne0$, so $\pbel_\tau(y_\tau)\le\pbsth$ (Proposition~\ref{pbe:local}(e)) and the
hypothesis of \textup{(O1)} on $\phi$ gives $\|y_\tau^{-1}\tilde f_a\|\le\pbsrb$. (Invariance under
$\prod F_\tau$ implies invariance under $\prod\pboR_\tau$.) The exponent $9n$ and the constants are
those of \S\ref{pbs:ss-open}, so Lemma~\ref{pbs:open-arith} applies unchanged.
\end{remark}

The proof occupies \S\S\ref{pbo:ss-system}--\ref{pbo:ss-proof}. Its structure: at one place, every
$g$ with $\pbg(g)>0$ is $s\gamma_Zt^{-1}$ with $s,t$ in the fixer of the frame and $\gamma_Z$ an
explicit smooth section of $g\mapsto\tilde Z(g)$ (\S\ref{pbo:ss-section}); Haar measure in the
coordinates $(s',Z,s)\mapsto s'\gamma_Zs$ has density at least $\pbocG\pbg(Z)^{k-1}$
(\S\ref{pbo:ss-GF}); a Sobolev inequality on a small cube of matrices near $\tilde Z(g)$, together
with Schur's orthogonality bound for $W$ and its derivatives, bounds the mean square over shifts
$t\in\pboR$ of the pairing (\S\S\ref{pbo:ss-sobolev}--\ref{pbo:ss-onplace}); the bounds multiply over
the places, one shift is no worse than the mean, and the cut is compared with $W$
(\S\ref{pbo:ss-proof}).

\subsection{Bessel pairs}\label{pbo:ss-system}

For a word $w=(\ell_1,\dots,\ell_m)$ in the alphabet $\mathcal E$ of pairs $(a,b)$ of basis indices
let $|w|=m$ and $(a,b)w$ the word with the letter $(a,b)$ prepended. A function $V\colon G\to\mathbb W$
into a finite-dimensional unitary representation $\mathbb W$ of $H$ is a \emph{Bessel function} if it
is continuous and bounded and $V(uhy)=\psi(u)\,h\cdot V(y)$ for $u\in U_\tau$, $h\in H$; then
$\|V(sy)\|=\|V(y)\|$ for $s\in R_\tau$, and more generally
\begin{equation}\label{pbo:eq-inv}
\Bigl\|\int_G\beta(y)V(sgy)\,dy\Bigr\|=\Bigl\|\int_G\beta(y)V(gy)\,dy\Bigr\|\qquad(s\in R_\tau,\ \beta\in
L^1(G)),
\end{equation}
because left translation by $s$ acts on the values by the unitary operator $\psi(u)h$, which
commutes with the integral. We write $L^2(U\backslash G)$ for the quotient measure of
\S\ref{pbe:ss-real}.

\begin{definition}[Bessel pair]\label{pbo:def-system}
Let $D>0$ and $b\ge1$. A \emph{Bessel pair of degree $D$} is a pair $(W,W^\flat)$ of functions $G\to
\mathbb W$, $d=\dim\mathbb W$, such that:
\begin{itemize}
\item[(BS1)] $W$ and $W^\flat$ are Bessel functions;
\item[(BS2)] $\int_{U\backslash G}\|W\|^2=d$;
\item[(BS3)] $W^\flat(y)\ne0$ implies $\pbel(y)\le b^{6n+2}$;
\item[(BS4)] (\emph{Schur bound with derivatives.}) for every $P\ge b^5$ there is a family
$(V_w)$ of Bessel functions indexed by the words, with $V_\varnothing=W$, such that (i) for every
$C^1$ curve $c$ in $G$ with $c'(t)=c(t)X$, $(V_w\circ c)'(t)=\sum_{a,b}X_{ab}V_{(a,b)w}(c(t))$; (ii)
for $|w|\le2n^2$ and every $\varphi\in L^1(G)$,
\[
\int_{U\backslash G}\Bigl\|\int_G\varphi(y)V_w(gy)\,dy\Bigr\|^2dg\le P^{2|w|}\,\frac dD\int_G|\varphi|^2 ;
\]
\item[(BS5)] $b\ge8r$.
\end{itemize}
\end{definition}


\begin{theorem}[The Bessel functions form Bessel pairs]\label{pbo:system}
Assume Theorem~\ref{pbe:ds} and the size conditions of \S\ref{pbo:ss-sizes}. For
every real place $\tau$, $(W_\tau,W^\flat_\tau)$ of Proposition~\ref{pbe:local} is a Bessel pair of
degree $\pbdeg_\tau$, with values in $\mathbb W_\tau$, $d=d_\tau$.
\end{theorem}

The proof is in \S\ref{pbo:ss-pairs}. Theorem~\ref{pbo:O1} will be deduced from the following local
statement, proved in \S\ref{pbo:ss-onplace}.

\begin{proposition}[Open-cell bound at one place]\label{pbo:local}
Let $(W,W^\flat)$ be a Bessel pair of degree $D$, $g\in G$ with $\pbg(g)>0$, $\rho\ge b$, and let some
$y$ have $\|y^{-1}\tilde f_a\|\le\rho$ and $\|y^{-1}g\tilde f_a\|\le\rho$ for all $a$. There are a compact
set $A\subseteq G$, a probability space $(T,m)$, a measurable map $t\colon T\to\pboR$ and a finite
measure $\mu$ on $G$ with $\mu(G)\le\pbsp(\tilde Z(g);D,\rho,d)$, such that
\begin{enumerate}
\item[(a)] for every bounded measurable $\varphi$ on $G$,
$\displaystyle\int_T\sum_{i,j}\Bigl|\int_A\varphi(z)\,W_i(gt(x)z)\,\aubar{W_j(t(x)z)}\,dz\Bigr|^2dm(x)\le
\int_G|\varphi|^2d\mu$, the coordinates being taken in an orthonormal basis of $\mathbb W$;
\item[(b)] for $m$-almost every $x$, $A$ contains every $z$ with $\|z^{-1}\tilde f_a\|\le\rho$ for all $a$
and $\pbel(gt(x)z)\le\rho$.
\end{enumerate}
\end{proposition}

\subsection{The section of the frame matrix over the open cell}\label{pbo:ss-section}

Norms of matrices are operator norms. For $x\in\C$ let $\Delta(x)=\operatorname{diag}(1,\dots,1,x)$
($n\times n$). The Gram matrix of the frame is $\pboF^*\mathsf J_1\pboF=\Delta(0)-1$, and
$\pboF^*\mathsf J_1\pboFd=\Delta(0)$; the columns of $(\pboF\,|\,\pboFd)$ are orthonormal.

\begin{definition}[The section]\label{pbo:def-section}
For $Z\in\pboO$ let $\zeta$ be the last row of $Z^{-1}$, $\pbonu=\pbonu(Z)>0$, and let $Y(Z)$ be the
$n\times n$ matrix with last row $(0,\dots,0,-1)$ and, for $a<\varrho$, entries $Y_{ab}=-c_b\,\aubar{
\zeta_a}\zeta_b$ with $c_b=\frac12$ for $b<\varrho$ and $c_\varrho=1$. Put $F^\vee(Z)=\pboF Y(Z)+\pboFd
\Delta(\pbonu)$, an $r\times n$ matrix. Let $\gamma_Z$ be the linear map that is the identity on the
orthogonal complement of the columns of $(\pboF\,|\,\pboFd)$ (spanned by the $x_i$, $i\ne k_0$) and
maps the columns of $\pboF Z^{-1}$ to those of $F^\vee(Z)$ and the columns of $F^\vee(Z^*)$ to those
of $\pboF(Z^*)^{-1}$.
\end{definition}

\begin{lemma}[The section]\label{pbo:section}
For $Z\in\pboO$: (a) $\gamma_Z$ is well defined and lies in $G$; (b) $\gamma_Z\pboF=F^\vee(Z)Z$, and
$\tilde Z(\gamma_Z)=Z$; (c) the matrix of $\gamma_Z$ is a real-analytic function of $Z$ on $\pboO$; it
is the restriction to $W=Z^*$ of a function $M(Z,W)$ holomorphic on $\pboO\times\pboO$, given by
rational expressions in the entries of $Z,Z^{-1},W^{-1}$ and in $\sigma=(1-\xi(Z)\xi(W))^{1/2}$
(principal branch; $|\xi(Z)\xi(W)|<1$), with $\sigma(Z,Z^*)=\pbonu(Z)$.
\end{lemma}
\begin{proof}
The vectors $\pboF Z^{-1}$ and $F^\vee(Z^*)$ (together $2n$ vectors) lie in the span $S$ of the columns
of $(\pboF|\pboFd)$, and so do $F^\vee(Z)$ and $\pboF(Z^*)^{-1}$. In the basis $(\pboF|\pboFd)$ of $S$
the two families are the columns of the block matrices $x=\bigl(\begin{smallmatrix}Z^{-1}&Y(Z^*)
\Delta(\pbonu)^{-1}\\0&1\end{smallmatrix}\bigr)\cdot\operatorname{diag}(1,\Delta(\pbonu))$ and $y=
\bigl(\begin{smallmatrix}Y(Z)&(Z^*)^{-1}\\\Delta(\pbonu)&0\end{smallmatrix}\bigr)$ up to the same
right factor, and $x$ is invertible (block triangular with invertible diagonal blocks), so $\gamma_Z|_S
=yx^{-1}$ is well defined. On $S$ the form is given by $h(\tilde e_a,\tilde e^*_b)=\delta_{ab}$, $h(\tilde z,\tilde z)=-1$, $h(x_{k_0},x_{k_0})=1$, all other products of basis vectors
of $S$ zero. One checks entry by entry that $y^*J_Sy=x^*J_Sx$: both are polynomial identities in the
entries of $Z^{-1}$ once $\pbonu^2=1-|\zeta_\varrho|^2$ is used. (For instance the Gram matrix of
$F^\vee(Z)$ is $Y^*(\Delta(0)-1)Y+Y^*\Delta(0)\Delta(\pbonu)+\Delta(\pbonu)\Delta(0)Y+\Delta(\pbonu)^2
\cdot e_{\varrho\varrho}$, and with the explicit $Y$ this is $-\zeta^*\zeta$, which is the Gram matrix
$(Z^{-1})^*(\Delta(0)-1)Z^{-1}$ of $\pboF Z^{-1}$; the products of $F^\vee(Z)$ with $\pboF$ are
$(\Delta(0)-1)Y+\Delta(0)\Delta(\pbonu)=1$, those of $\pboF Z^{-1}$ with $F^\vee(Z^*)$ are
$(Z^{-1})^*\cdot1$, and so on.) Hence $yx^{-1}$ is an isometry of $S$, and extended by the identity
on $S^\perp$ it is an isometry of $Y_\tau$: (a). (b) $\gamma_Z\pboF=\gamma_Z(\pboF Z^{-1})Z=F^\vee(Z)Z$,
and $\tilde Z(\gamma_Z)=\pboF^*\mathsf J_1F^\vee(Z)Z=((\Delta(0)-1)Y+\Delta(0)\Delta(\pbonu))Z=Z$, as
$(\Delta(0)-1)Y$ keeps the negated last row $(0,\dots,0,1)$ of $Y$ and $\Delta(0)\Delta(\pbonu)=
\Delta(0)$. (c) Replacing $Z^*$ by an independent $W\in\pboO$ and $\pbonu$ by $\sigma$ in the formula
$yx^{-1}$ gives $M(Z,W)$; $1-\xi(Z)\xi(W)$ lies in the disc of centre $1$ and radius $1$, where the
principal square root is holomorphic, and $\sigma(Z,Z^*)=(1-|\xi|^2)^{1/2}=\pbonu$.
\end{proof}

\begin{lemma}[Double cosets over the open cell]\label{pbo:witt}
The stabiliser $F_\tau$ of $\tilde f_0,\dots,\tilde f_\varrho$ in $G$ is $\pboR=\pboU H$. If $g\in G$ and
$Z=\tilde Z(g)\in\pboO$, there are $s,s'\in\pboR$ with $g=s\gamma_Zs'^{-1}$.
\end{lemma}
\begin{proof}
An element fixing the frame preserves $N_0^\perp$, where $N_0=\langle\tilde f_a\rangle$, hence acts on
$N_0^\perp/(N_0\cap N_0^\perp)=N_0^\perp/A_\tau\cong X_{0,\tau}$ by some $h\in H$; then $h^{-1}$ times it
fixes the frame and acts trivially on $N_0^\perp/A_\tau$, and such an element is unipotent and lies in
$U_\tau$, hence in $\pboU$. For the second claim, consider the $2n$ vectors $(\pboF|g\pboF)$ and
$(\pboF|\gamma_Z\pboF)$. Their Gram matrices are both $\bigl(\begin{smallmatrix}\Delta(0)-1&Z\\Z^*&\Delta
(0)-1\end{smallmatrix}\bigr)$ (as $g,\gamma_Z$ are isometries and $\tilde Z(g)=\tilde Z(\gamma_Z)=Z$), and
its determinant is $\pm\pbg(Z)\ne0$ (proof of Lemma~\ref{pbs:open-arith}, with $N=1$ in the normalised
basis). So both families are bases of non-degenerate subspaces, and by Witt's theorem there is
$x\in G$ with $x\pboF=\pboF$ and $xg\pboF=\gamma_Z\pboF$. Then $x\in\pboR$ and $\gamma_Z^{-1}xg\in\pboR$, and
$g=x^{-1}\gamma_Z(\gamma_Z^{-1}xg)$.
\end{proof}

\begin{lemma}[Sizes at and near a point of the open cell]\label{pbo:sizes}
Say that $Z$ is \emph{$P$-bounded} if $Z\in\pboO$ and $\|Z\|,\|Z^{-1}\|,\pbonu(Z)^{-1}\le P$.
\begin{enumerate}
\item[(a)] If $n\le M$, $|Z_{ab}|\le M^2$ for all $a,b$ and $\pbg(Z)\ge M^{-1}$, then $Z$ is
$T$-bounded for every $T\ge M^{8n}$.
\item[(a$'$)] If $n\le M$ and $|Z_{ab}|\le M^2$ for all $a,b$, then $\pbg(Z)\le|\det Z|^2\le M^{8n}$.
\item[(b)] If $Z$ is $T$-bounded with $T\ge4$ and $\|Z'-Z\|<T^{-6}$, then $Z'$ is $2T$-bounded; if also
$T\ge\max(8,n)$, then $\pbg(Z')\ge\pbg(Z)/2$.
\item[(c)] If $Z$ is $T$-bounded, $T\ge4$, then $\|\gamma_Z\|\le T^{18}$ and the real derivative of
$Z\mapsto\gamma_Z$ at $Z$ has norm at most $T^{25}$.
\item[(d)] If $s\in U_\tau$ and $\tilde Z(y)$, $\tilde Z(sy)$ are $P$-bounded, the entries of the matrix
$m(s)$ of $s$ on the frame ($s\tilde f_b=\sum_am(s)_{ab}\tilde f_a$) have modulus at most $P^2$.
\end{enumerate}
\end{lemma}
\begin{proof}
(a) Every square submatrix $B$ of $Z$ of size $l\le n$ has $\|B\|\le lM^2\le M^4$ (bound the operator
norm by $l$ times the largest entry), so $|\det B|\le M^{4l}$. From $M^{-1}\le\pbg=|\det Z|^2\pbonu^2$
and $\pbonu\le1$, $|\det Z|\ge M^{-1/2}$ and $\pbonu\ge M^{-1/2}|\det Z|^{-1}\ge M^{-1/2-4n}$; so
$\pbonu^{-1}\le M^{4n+1}$. By Cramer's rule $\|Z^{-1}\|\le|\det Z|^{-1}\cdot n\max|\text{minors}|\le
M\cdot M^2M^{4\varrho}=M^{4n-1}$. And $\|Z\|\le M^4$. (a$'$) $\pbg\le|\det Z|^2\le M^{8n}$ by the bound on $|\det B|$ with $B=Z$,
which uses only $n\le M$ and $|Z_{ab}|\le M^2$. (b) $Z'=Z(1+Z^{-1}
(Z'-Z))$ with $\|Z^{-1}(Z'-Z)\|\le T^{-5}\le\frac12$, so $Z'$ is invertible, $\|Z'^{-1}\|\le2T$ and
$\|Z'^{-1}-Z^{-1}\|=\|Z'^{-1}(Z-Z')Z^{-1}\|\le2T^{-4}$. Hence $|\xi(Z')|\le|\xi(Z)|+2T^{-4}$, and as
$\pbonu(Z)^2\ge T^{-2}$ one gets $\pbonu(Z')^2\ge(1-8T^{-2})\pbonu(Z)^2\ge\frac14T^{-2}$ (for $T\ge4$),
i.e.\ $\pbonu(Z')^{-1}\le2T$; and $\|Z'\|\le T+1\le2T$. If $T\ge\max(8,n)$, then $\pbonu(Z')^2\ge\frac78
\pbonu(Z)^2$, and from $Z=(1-(Z'-Z)Z'^{-1})Z'$, $|\det Z|\le(1+2T^{-5})^n|\det Z'|\le\frac98|\det Z'|$;
by \eqref{pbo:eq-gnu}, $\pbg(Z')\ge\frac78(\frac89)^2\pbg(Z)\ge\frac12\pbg(Z)$. (c) Let $Z$ be $P$-bounded,
$P\ge8$, and $W\in\pboO$ with $\|W^{-1}\|\le P$. Write $Y(Z,W)=\Delta(0)-1-\Delta(0)W^{-1}(1-\Delta(\frac12))
Z^{-1}\Delta(2)$, so that $Y(Z)=Y(Z,Z^*)$ (compare entries with Definition~\ref{pbo:def-section}), and
$F^\vee(Z,W)=\pboF Y(Z,W)+\pboFd\Delta(\sigma)$, $\sigma=(1-\xi(Z)\xi(W))^{1/2}$. Unwinding $yx^{-1}$ in the proof of
Lemma~\ref{pbo:section}, the matrix of $\gamma_Z$ is $M(Z,Z^*)$ with
\[
M(Z,W)=1+(F^\vee(Z,W)Z-\pboF)\pboF^*+\bigl((\pboF W^{-1}-F^\vee(Z,W)Z\,Y(W,Z))\Delta(\sigma)^{-1}-\pboFd\bigr)\pboFd^* .
\]
Now $\|\pboF\|=\|\pboFd\|=1$, $\|\Delta(x)\|=\max(1,|x|)$, $\|\Delta(0)-1\|=1$, $\|1-\Delta(\frac12)\|\le1$, $\|\Delta(2)\|=2$, so
$\|Y(Z,W)\|,\|Y(W,Z)\|\le1+2P^2$; $|\sigma|^2\le1+|\xi(Z)\xi(W)|\le2$; and $|1-\xi(Z)\xi(W)|\ge1-|\xi(Z)|\ge\frac12\pbonu(Z)^2$,
so $|\sigma^{-1}|\le\sqrt2\,\pbonu(Z)^{-1}\le2P$. Hence $\|F^\vee(Z,W)Z\|\le(1+2P^2+2)P\le3P^3$ (as $P\ge8$),
$\|(\pboF W^{-1}-F^\vee ZY(W,Z))\Delta(\sigma)^{-1}\|\le(P+3P^3(1+2P^2))\cdot2P\le2P\cdot7P^5=14P^6$, and
\[
\|M(Z,W)\|\le1+(3P^3+1)+(14P^6+1)\le15P^6\le P^{12}\qquad(P\ge8).
\] By (b), for $Z$ $T$-bounded and $T\ge4$,
every $Z'$ with $\|Z'-Z\|<T^{-6}$ is $2T$-bounded, and the same holds for $W'$ near $Z^*$; so $M$ is
holomorphic and bounded by $(2T)^{12}$ on the product of the operator-norm balls of radius $T^{-6}$ about $Z$ and $Z^*$ (Cauchy's estimate in one variable along each direction; \S\ref{pbv:ss-weak}(iii)). The Cauchy
estimate gives a derivative of norm at most $2(2T)^{12}T^6\le T^{25}$, and $\|\gamma_Z\|\le(2T)^{12}\le
T^{18}$. (d) $\tilde Z(y)=m(s)^*\tilde Z(sy)$ ($s$ maps the frame to itself by $m(s)$ and is an
isometry; this is the computation of Lemma~\ref{pbs:cells}), so $m(s)^*=\tilde Z(y)\tilde Z(sy)^{-1}$
has norm at most $P^2$.
\end{proof}

\begin{lemma}[Bounds from the length]\label{pbo:norm-z}
Let $8n\le M$, $M^{9n}\le R$, $Z\in\pboO$ with $\pbg(Z)\ge M^{-1}$ and $\|Z^{-1}\|\le R$, $s\in R_\tau$,
$v\in\pboU$, $h\in H$, and let $z\in G$ satisfy $\|z^{-1}\tilde f_a\|\le M$ for all $a$ and $\pbel(s\gamma_Z
(vh)^{-1}z)\le M$. Then $\|v^{-1}z\|\le R^5$.
\end{lemma}
\begin{proof}
Put $z_1=(vh)^{-1}z$; it has the same frame vectors $z_1^{-1}\tilde f_a=z^{-1}\tilde f_a$ (as $vh$ fixes
the frame), $\pbel(\gamma_Zz_1)\le M$ (Lemma~\ref{pbe:len}(a)) and $\|v^{-1}z\|=\|hz_1\|=\|z_1\|$. So we
may assume $s=v=h=1$. Choose $w\in U_\tau$ with $\|w\gamma_Zz\|\le M$ (the infimum defining $\pbel$ over
$R_\tau=U_\tau H$ may be taken over $U_\tau$, as $H$ preserves norms) and put $y=w\gamma_Zz$, $Z'=\tilde Z(
w\gamma_Z)=m(w)^{-*}Z$ (Lemma~\ref{pbo:sizes}(d)); $\pbg(Z')=\pbg(Z)$, as $m(w)$ is unitriangular. The
entries $Z'_{ab}=h(y^{-1}\tilde f_a,y^{-1}w\gamma_Z\tilde f_b)=h(y^{-1}\tilde f_a,z^{-1}\tilde f_b)$ have
modulus $\le M^2$, so by Lemma~\ref{pbo:sizes}(a) $\|Z'^{-1}\|\le M^{8n}\le R$. Reading off the
column at $\tilde e^*_a$ of the identity $z^{-1}=(y^{-1}w)\gamma_Z$, expanded with the formula for
$\gamma_Z$ (only the second block column of $yx^{-1}$ contributes, and $m(w)$ cancels against $Z$),
\[
z^{-1}\tilde e^*_a=\sum_{i\le\varrho}\bigl((Z'^*)^{-1}\bigr)_{ia}\,y^{-1}\tilde f_i+\sum_{i<\varrho}\tfrac12
(Z^{-1})_{i\varrho}\bigl((Z^*)^{-1}\bigr)_{\varrho a}\,z^{-1}\tilde f_i .
\]
Hence $\|z^{-1}\tilde e^*_a\|\le nR\cdot M+\varrho R^2M/2\le R^4$, and $\|z^{-1}\tilde f_a\|\le M\le R^4$.
Finally, for $A=z^{-1}$, the relation $A\mathsf J_1A^*=\mathsf J_1$ gives $\|A\|^2=\|\mathsf J_1+A(1-\mathsf
J_1)A^*\|\le1+\|A(1-\mathsf J_1)\|\,\|A\|$, and the columns of $1-\mathsf J_1$ are $0$, $2\tilde z$ and
$\pm(\tilde e_a-\tilde e^*_a)$, so $\|A(1-\mathsf J_1)\|\le(4\varrho+2)R^4$; as $\|A\|\ge1$, $\|z\|=\|A\|
\le1+(4\varrho+2)R^4\le R^5$.
\end{proof}

\subsection{Haar measure in the coordinates of the open cell}\label{pbo:ss-GF}

On complex matrices we use Lebesgue measure for the real and imaginary parts of the entries, and on
hermitian $m\times m$ matrices $T$ Lebesgue measure $dT$ in the real coordinates $T_{aa}$, $\Rel T_{ab}$,
$\Iml T_{ab}$ ($a<b$). For a tuple $E$ of $m$ vectors of $Y_\tau$ (an $r\times m$ matrix) let $\Phi(E)=
E^*\mathsf J_1E$. On the open set of tuples with $\Phi$ a submersion (e.g.\ $E$ of rank $m$ spanning a
non-degenerate subspace) Lebesgue measure disintegrates as $dE=\int\lambda_T\,dT$ with fibre measures
$\lambda_T$ on $\{\Phi=T\}$ (coarea formula). For a measure $\mu$ on $G$ and a tuple $E_0$ with
$\Phi(E_0)=T_0$, write $\mu\preceq\lambda_{E_0}$ if the image of $\mu$ under $g\mapsto gE_0$ is at most
$\lambda_{T_0}$, in the weak sense: $\int K(gE_0)\,d\mu(g)\le C$ whenever $K\ge0$ is lower semicontinuous
and $\int_{\Phi(E)\in S}K\,dE\le C\,\pbvol(S)$ for all open $S$ in a neighbourhood of $T_0$.

\begin{lemma}[Gram fibre normalisation \textup{(GF)}]\label{pbo:GF}
Let $dg$ be the Haar measure of \textup{(DS1)(iii)} and $\sigma$ a real-analytic section of $\Phi$ near
$\mathsf J_1$ with $\sigma(\mathsf J_1)=1$ and invertible values (e.g.\ $\sigma(T)=(\mathsf J_1T)^{1/2}$, the
power series square root near $1$). Then $(g,T)\mapsto g\sigma(T)$ carries $\frac{\pboV(r)}{2^r}\,dg\,dT$ to
Lebesgue measure $dY$ on a neighbourhood of $G=\{\Phi=\mathsf J_1\}$. Consequently every measure
$\mu\preceq\lambda_1$ satisfies $\mu\le\frac{\pboV(r)}{2^r}dg$.
\end{lemma}
\begin{proof}
\emph{Reduction to a constant.} The map $(g,T)\mapsto g\sigma(T)$ is a local diffeomorphism near
$G\times\{\mathsf J_1\}$ (its differential is computed below and is bijective) and commutes with
$Y\mapsto g_0Y$, $g_0\in G$, which preserves $dY$ ($|\det g_0|=1$) and $dg\,dT$. Hence the image of
$dg\,dT$ is $c(T)^{-1}dY$ for a function $c$ of $T$ alone. It does not depend on $T$ either: for
$A\in GL_r(\C)$ near $1$ and $T$ near $\mathsf J_1$ put $g'(T)=\sigma(T)A\,\sigma(A^*TA)^{-1}$; since
$\Phi(\sigma(T)A)=A^*TA=\Phi(\sigma(A^*TA))$, we get $\Phi(g'(T))=\mathsf J_1$, i.e.\ $g'(T)\in G$, and
$g\sigma(T)A=gg'(T)\,\sigma(A^*TA)$. Thus $Y\mapsto YA$ corresponds to $(g,T)\mapsto(gg'(T),A^*TA)$. The
first map multiplies $dY$ by $|\det A|^{2r}$; the second multiplies $dg\,dT$ by $|\det A|^{2r}$, because
it is triangular, right translation by $g'(T)$ preserves $dg$ ($G$ is unimodular) and $T\mapsto A^*TA$
multiplies $dT$ by $|\det A|^{2r}$. Hence $c(T)=c(A^*TA)$, and every $T$ near $\mathsf J_1$ is
$A^*\mathsf J_1A$ with $A=\sigma(T)$ near $1$. So $c$ is a constant $C$, the Jacobian at $(1,\mathsf J_1)$. Writing $T=\mathsf J_1+H$,
the differential there is $(X,H)\mapsto X+\frac12\mathsf J_1H$ on $\mathfrak g_0\times\mathrm{Herm}_r$, because
$\Phi(1+X+\frac12\mathsf J_1H)=\mathsf J_1+H+O(2)$ ($X^*\mathsf J_1+\mathsf J_1X=0$, $\mathsf J_1^2=1$). With
$X_0=\mathsf J_1X\in\mathfrak u(r)$ it is $\mathsf J_1\circ\bigl((X_0,H)\mapsto X_0+\frac12H\bigr)$, and $\mathsf J_1$ is
unitary. By Remark~\ref{pds:cayley}, the density of $dg$ at $1$ in the coordinate $X_0$ equals that of the
Haar probability measure of $\mathrm U(r)$ at $1$ in the coordinate $X_0$ (the Cayley density at $0$ divided
by the same normalising constant, and $X\mapsto\mathsf J_1X$ has Jacobian $1$). So $C$ is the constant of the
case $\mathsf J_1=1$, $G=\mathrm U(r)$, $dg$ of mass one.

\emph{The compact case.} For $\mathsf J_1=1$ the same argument applies at every point of
$\mathrm U(r)\times\{T>0\}$ with the global diffeomorphism $(g,T)\mapsto gT^{1/2}$ onto the invertible
matrices (polar decomposition). For $A\in GL_r(\C)$, $Y\mapsto YA$ multiplies $dY$ by $|\det A|^{2r}$ and
$T\mapsto A^*TA$ multiplies $dT$ by $|\det A|^{2r}$; and $gT^{1/2}A=gu_0(A^*TA)^{1/2}$ with $u_0=T^{1/2}A(A^*TA)^{-1/2}$
unitary, so the fibre over $T$ is carried to the fibre over $A^*TA$ by $g\mapsto gu_0$, which preserves $dg$.
Hence the density of $dY$ against $dg\,dT$ is the same at every $T>0$, namely $C$. So $dY=C\,dg\,dT$ globally. Integrate $e^{-\operatorname{tr}Y^*Y}$: $\pi^{r^2}=C\int_{T>0}
e^{-\operatorname{tr}T}dT=C\,\pi^{r(r-1)/2}\prod_{j=1}^r(j-1)!$ (complex Siegel gamma integral at $s=r$:
$\int_{T>0}e^{-\operatorname{tr}T}(\det T)^{s-r}dT=\pi^{r(r-1)/2}\prod_{j=1}^r\Gamma(s-j+1)$, proved by induction on
$r$ with $T=\bigl(\begin{smallmatrix}T'&v\\v^*&t\end{smallmatrix}\bigr)$, $t=t_1+v^*T'^{-1}v$, $\det T=t_1\det T'$, and
the Gaussian integral in $v$ giving $\pi^{r-1}\det T'$). Hence $C=\prod_{j=1}^r\pi^j/(j-1)!=\pboV(r)/2^r$.

\emph{The consequence.} If $\mu\preceq\lambda_1$ and $O\subseteq G$ is open, test with $K$ the indicator
of the open set $\{g\sigma(T):g\in O,T\in S_0\}$ for a small open $S_0\ni\mathsf J_1$; for open $S\subseteq S_0$,
$\int_{\Phi\in S}K\,dY=C\pbvol(O)\pbvol(S)$, so $\mu(O)\le C\pbvol(O)$; outer regularity of $dg$ gives
$\mu\le C\,dg$.
\end{proof}


\begin{lemma}[Appending vectors]\label{pbo:append}
\begin{enumerate}
\item[(a)] Let $E=(E_1|E_2)$ be a tuple split into two parts, let $\mu\preceq\lambda_{E_1}$, and let
$\nu$ be a measure on the stabiliser $G_{E_1}$ of $E_1$ such that the image of $\nu$ under $u\mapsto
uE_2$ is at most the fibre measure, over $(E_1^*\mathsf J_1E_2,\Phi(E_2))$, of the map $X\mapsto(E_1^*
\mathsf J_1X,\Phi(X))$ on tuples $X$ of the size of $E_2$. Then $\mu*\nu\preceq\lambda_E$.
\item[(b)] Let $\mu\preceq\lambda_{\pboF}$ and let $ds=du^0\,dh$ be the measure on $\pboR$, $du^0$ the
Haar measure on $\pboU$ that is the image of Lebesgue measure on the parameters $q=(\beta,\sigma)$
($\beta\in M_{\varrho\times k}(\C)$, $\sigma$ real skew-hermitian coordinates)  $q$ of
$u_q$ in \eqref{pbo:eq-uq} below, and $dh$ of mass one. Then $\mu*ds\le\frac{\pboV(r)}{2^{\varrho+1}\pboV(k)}\,dg$.
\end{enumerate}
\end{lemma}
\begin{proof}
(a) For $K$ as in the definition, disintegrate $dE=dE_1\,dE_2$ first along $\Phi(E_1)$ and then, for
fixed $E_1$, along $(E_1^*\mathsf J_1E_2,\Phi(E_2))$; the Gram matrix of $E$ is determined by these
blocks. Moving $E_1$ to $gE_1$ by $g$ (in the support of $\mu$) and using $\nu$ on its stabiliser,
$\int K(gE)\,d(\mu*\nu)(g)=\int\int K(guE)\,d\nu(u)\,d\mu(g)$, and the inner integral is bounded through
the hypothesis on $\nu$ by the integral of $K(g\,\cdot)$ over the fibre in the second variable, then
the outer one through $\mu\preceq\lambda_{E_1}$. (b) Append to $\pboF$ first the vectors $\tilde
e^*_1,\dots,\tilde e^*_\varrho$, then $x_1,\dots,x_k$. The tuples $X=(\tilde e^*_a)+\pboF'' p$, with $\pboF''$
the columns $\tilde e_a,x_i$ (orthonormal, $h$-orthogonal to $\pboF$), are exactly those with $\pboF^*
\mathsf J_1X=\pboF^*\mathsf J_1(\tilde e^*_a)$; writing such a tuple as $u_q(\tilde e^*+\tilde e\,T/2)$
with $T$ its Gram matrix, Lebesgue measure is $dp=2^{-\varrho}dT\,dq$ (translate the $\tilde e$-part
by the $q$-dependent quadratic term and split a complex matrix into hermitian and skew-hermitian parts,
which costs $2^{-\varrho}$ in the diagonal; this is the computation \eqref{pbo:eq-dual} below with
$k=0$). So (a) applies with $\nu=2^{-\varrho}du^0$. For the last $k$ vectors the stabiliser is $H$,
and the image of $\frac{\pboV(k)}{2^k}dh$ under $h\mapsto(hx_i)$ is the fibre measure of $X\mapsto
(X^*X)$ on $k$-tuples of $X_0$ over $1$ (\textup{(GF)} for $\mathrm U(k)$); tuples of $Y_\tau$
orthogonal to the first $2\varrho+1$ vectors are tuples in $X_0$, and the linear part of the map costs
nothing (orthonormal coordinates). By (a) twice, $\mu*(2^{-\varrho}du^0)*(\frac{\pboV(k)}{2^k}dh)\preceq
\lambda_1$, so by \textup{(GF)} $\mu*ds\le2^\varrho\frac{2^k}{\pboV(k)}\frac{\pboV(r)}{2^r}dg$, and $r=k+
2\varrho+1$.
\end{proof}

We make the parametrisation of $\pboU$ explicit. For $\beta\in M_{\varrho\times k}(\C)$ and real
skew-hermitian $\sigma\in M_\varrho(\C)$ let $u_q$ ($q=(\beta,\sigma)$) fix $\tilde e_a$ and $\tilde z$, map
$x\mapsto x+\tilde e\,\beta$ (that is, $x_i\mapsto x_i+\sum_a\beta_{ai}\tilde e_a$), and
\begin{equation}\label{pbo:eq-uq}
u_q(\tilde e^*+\tilde e\,a)=\tilde e^*+\tilde e\,(\sigma+a-\tfrac12\beta\beta^*)-x\,\beta^*,
\end{equation}
in matrix notation for tuples ($\tilde e$, $x$, $\tilde e^*$ the rows of basis vectors). One checks directly
that $u_q$ preserves $h$ (the only non-trivial products are $h(u_qx_i,u_q\tilde e^*_b)=\aubar{\beta_{bi}}-\aubar{\beta_{bi}}=0$
and the Gram matrix $(\sigma+\sigma^*)-\beta\beta^*+\beta\beta^*=0$ of $u_q\tilde e^*$). Every element of $\pboU$ is a unique $u_q$, and $q\mapsto u_q$ carries
Lebesgue measure to a Haar measure $du^0$ (it is a two-step nilpotent group in exponential-type
coordinates). We fix this $du^0$; nothing below depends on the choice except through quantities in
which it cancels.

\begin{proposition}[Mean value inequality]\label{pbo:meanvalue}
Assume \textup{(GF)}. For measurable $f\ge0$ on $G$,
\[
\int_{\pboO}\pbocG\,\pbg(Z)^{k-1}\int_{\pboR}\int_{\pboR}f(s'\gamma_Zs)\,ds\,ds'\,dZ\le\int_Gf(g)\,dg .
\]
\end{proposition}
\begin{proof}
\emph{Step 1: tuples dual to the frame.} An $n$-tuple $E$ with $\pboF^*\mathsf J_1E=Z$ invertible is
$E=E'Z$ with $\pboF^*\mathsf J_1E'=1$; such $E'$ are $E'=\mathsf J_1\pboF+\pboF''p$ with $p$ a
$(\varrho+k)\times n$ matrix, and $T'=\Phi(E')=(Z^{-1})^*\Phi(E)Z^{-1}$. For a hermitian $T'$ with
$1+T'_{\varrho\varrho}>0$ let $E'(T')$ be the tuple with columns $\tilde e^*_b+\frac12\sum_aT'_{ab}\tilde
e_a$ ($b<\varrho$) and $-\tilde z+\sum_aT'_{a\varrho}\tilde e_a+(1+T'_{\varrho\varrho})^{1/2}x_{k_0}$; it is dual
to the frame with Gram matrix $T'$. Every dual tuple is $s\,E'(T')$ with $s\in\pboR$: the first $\varrho$
vectors are moved by $\pboU$ (as in Lemma~\ref{pbo:append}(b)), and the $X_0$-part of the last vector,
of squared length $1+T'_{\varrho\varrho}$, by $H$ (polar coordinates in $X_0\cong\C^k$). In these
coordinates
\begin{equation}\label{pbo:eq-dual}
dp=2^{-\varrho}\,\mathsf m_k(1+T'_{\varrho\varrho})\,dT'\,ds,\qquad\mathsf m_k(t)=\frac{\pi^kt^{k-1}}{(k-1)!}
\ (t>0),
\end{equation}
where $\mathsf m_k(t)$ is the mass of the sphere $\{\|x\|^2=t\}$ of $\C^k$ for $dx/d\|x\|^2$ (polar
coordinates: $dx=\mathsf m_k(t)\,dt\,dh$ with $dh$ of mass one), and the factor $2^{-\varrho}$ comes from
the first $\varrho$ vectors. The substitution $E=E'Z$ multiplies $dp$ by $|\det Z|^{2(\varrho+k)}$ and
$T=Z^*T'Z$ multiplies $dT'$ by $|\det Z|^{-2n}$; so on the tuples over $Z$,
\[
dp_E=|\det Z|^{2(k-1)}2^{-\varrho}\mathsf m_k\bigl(1+((Z^{-1})^*TZ^{-1})_{\varrho\varrho}\bigr)\,dT\,ds .
\]
\emph{Step 2: the fibre over the frame.} At $T=\Phi(\pboF)=\Delta(0)-1$ we have $1+((Z^{-1})^*T Z^{-1})_{
\varrho\varrho}=1-|\xi(Z)|^2=\pbonu(Z)^2$, and the tuple $E'((Z^{-1})^*TZ^{-1})Z$ is $\gamma_Z\pboF$ (compare
coefficients with Definition~\ref{pbo:def-section}: the coefficients of $\tilde e_a$, $x_{k_0}$, $\tilde
z$, $\tilde e^*_a$ are those of $Y(Z)$ and $\Delta(\pbonu)$). By \eqref{pbo:eq-gnu},
$|\det Z|^{2(k-1)}2^{-\varrho}\mathsf m_k(\pbonu^2)=2^{-\varrho-1}\frac{2\pi^k}{(k-1)!}\pbg(Z)^{k-1}$. As the
tuples $\mathsf J_1\pboF Z+\pboF''p$ ($Z\in\pboO$) are a subset of all tuples and $(Z,p)\mapsto\mathsf J_1
\pboF Z+\pboF''p$ preserves Lebesgue measure, the measure $\mu_0$, image of $2^{-\varrho-1}\frac{2\pi^k}{(
k-1)!}\pbg(Z)^{k-1}\mathbf 1_{\pboO}(Z)\,dZ\,ds'$ under $(Z,s')\mapsto s'\gamma_Z$, satisfies $\mu_0\preceq
\lambda_{\pboF}$ (the coordinates $(T,Z,s')$ are continuous at $T_0$ and the density is continuous
there, which is what the weak definition needs).

\emph{Step 3.} By Lemma~\ref{pbo:append}(b), $\mu_0*ds\le\frac{\pboV(r)}{2^{\varrho+1}\pboV(k)}dg$.
Integrating $f$ and dividing by the constant gives the claim, since $\frac{2^{\varrho+1}\pboV(k)}{\pboV(r)}
\cdot2^{-\varrho-1}\frac{2\pi^k}{(k-1)!}=\pbocG$.
\end{proof}

\subsection{Unfolding, Sobolev's inequality, and the bound at the section}\label{pbo:ss-sobolev}

\begin{lemma}[Unfolding along $U_\tau$]\label{pbo:unfold}
Let $w,\varphi\ge0$ be measurable on $G$, $w$ left invariant under $R_\tau$, and $K,R\ge1$ such that
$\int_{\pboU}\varphi(uy)\,du^0\le K$ for all $y$, and such that $\varphi(y)\ne0\ne\varphi(sy)$ with $s\in U_\tau$
implies $|m(s)_{ab}|\le R$ for all $a,b$. Then $\int_G\varphi w\,dg\le K(2R)^{2n^2}\int_{U\backslash G}w$,
for the Haar measure $du=du^0\,dx$ on $U_\tau$, $dx$ Lebesgue measure in the strictly upper triangular
entries $x$ of $m(s)$.
\end{lemma}
\begin{proof}
The map $s\mapsto m(s)$ is a homomorphism from $U_\tau$ onto the unitriangular group of the frame with
kernel $\pboU$, so every $s\in U_\tau$ is $s=u\,\ell(x)$ with $u\in\pboU$ and $\ell$ a measurable section
with $m(\ell(x))=1+x$. The measure $du^0dx$ is the measure $du$ of convention \textup{(Cv3)}: the
\textup{(Cv3)} coordinates of $s$ in the frame columns are the entries of $m(s)-1$, and on $\pboU$ the
other \textup{(Cv3)} coordinates are $\beta$, the entries $(\sigma-\frac12\beta\beta^*)_{ab}$ ($a<b$) and
$\Iml\sigma_{aa}$, which differ from $q$ by a translation depending on $\beta$ (so Lebesgue measure is
unchanged); by Weil's formula for the normal subgroup $\pboU$ of $U_\tau$ the two Haar measures
coincide (the measure of (Cv3) is a Haar measure by Lemma~\ref{pbv:haar}, and the two have the same
density at $1$ in these coordinates). By unfolding, $\int_G\varphi w=\int_{U\backslash G}w(y)\int_{U_\tau}\varphi(sy)\,ds\,dy$.
Fix $y$ and some $s_0$ with $\varphi(s_0y)\ne0$ (if none, the inner integral is $0$); substituting $s\mapsto
ss_0$, the inner integral is $\int dx\int_{\pboU}\varphi(u\ell(x)s_0y)\,du^0$, whose inner integral is $\le K$
and vanishes unless $|x_{ab}|\le R$ (apply the hypothesis to $s_0y$ and $u\ell(x)$, whose frame matrix is
$1+x$). The box of strictly upper triangular complex $x$ with $|x_{ab}|\le R$ has volume $\le(\pi R^2)^{
n^2/2}\le(2R)^{2n^2}$.
\end{proof}

\begin{lemma}[Integral over the section]\label{pbo:section-int}
Assume \textup{(GF)}. Let $P\ge1$, $c>0$, $Q\subseteq\pboO$ measurable with $\pbg\ge c$ and every
$Z\in Q$ $P$-bounded, and $w\ge0$ measurable and left invariant under $R_\tau$. Then
\[
\int_Q\int_{H\times\pboU}w\bigl(\gamma_Z(vh)^{-1}\bigr)\,\pbocG c^{k-1}\,dh\,dv\,dZ\le(2P^2)^{2n^2}\int_{U
\backslash G}w .
\]
\end{lemma}
\begin{proof}
Let $\chi\ge0$ be a cutoff for $\pboU$ ($\int_{\pboU}\chi(uy)du^0=1$ for all $y$) and $\chi_Q=\chi\cdot\mathbf
1[\tilde Z(\cdot)\in Q]$. By Lemma~\ref{pbo:unfold} with $K=1$, $R=P^2$ and Lemma~\ref{pbo:sizes}(d),
$\int_G\chi_Qw\le(2P^2)^{2n^2}\int_{U\backslash G}w$. By Proposition~\ref{pbo:meanvalue} applied to $\chi_Qw$
and restricted to $Q$, $\int_G\chi_Qw\ge\int_Q\pbocG\pbg^{k-1}\int\int\chi_Q(s'\gamma_Zs)w(s'\gamma_Zs)\,ds\,
ds'\,dZ$. For $s,s'\in\pboR$, $\tilde Z(s'\gamma_Zs)=Z$ and $w(s'\gamma_Zs)=w(\gamma_Zs)$; integrating $\chi$
over $s'=uh$ gives $1$ (the measure $ds$ is invariant under inversion, as $(uh)^{-1}=(h^{-1}u^{-1}h)h^{-1}$
and conjugation by $H$ preserves $du^0$). So the right side is $\int_Q\pbocG\pbg^{k-1}\int w(\gamma_Z(vh)^{-1})
\,dh\,dv\,dZ$, and $\pbg\ge c$ on $Q$.
\end{proof}

\begin{lemma}[Sobolev's inequality on a cube]\label{pbo:sobolev}
Let $f(x,\omega)$ be functions of $x\in[0,\varepsilon]^N$, $\varepsilon\le1$, and of $\omega$ in a measure space,
with values in a finite-dimensional Hilbert space, indexed by words $w$ in a finite alphabet $L$, such
that for each coordinate $j$, each $w$, $x$ and $\omega$, $t\mapsto f_w(x+te_j,\omega)$ is differentiable on
the cube with derivative of norm $\le\Lambda\sum_{l\in L}\|f_{lw}(x+te_j,\omega)\|$. Then
\[
\int\|f_\varnothing(0,\omega)\|^2d\omega\le a^N\max_{|w|\le N}\int_{[0,\varepsilon]^N}\int\|f_w(x,\omega)\|^2\,d\omega\,
dx,\qquad a=\frac2\varepsilon\bigl(1+(|L|\Lambda)^2\bigr).
\]
\end{lemma}
\begin{proof}
For a function $g$ on $[0,\varepsilon]$ with values in a Hilbert space, differentiable, $g(0)=g(t)-\int_0^tg'$ gives
$\|g(0)\|^2\le2\|g(t)\|^2+2\varepsilon\int_0^\varepsilon\|g'\|^2$, and averaging over $t$, $\|g(0)\|^2\le\frac2\varepsilon
\int_0^\varepsilon\|g\|^2+2\varepsilon\int_0^\varepsilon\|g'\|^2$. If $\|g'\|\le\Lambda\sum_{l\in L}\|g_l\|$, then $\|g'\|^2\le
\Lambda^2|L|\sum_l\|g_l\|^2$. For $0\le j\le N$ and a word $w$ let $I_j(w)$ be the integral of $\|f_w(x,\omega)\|^2$
over $\omega$ and over $x$ with $x_1,\dots,x_j\in[0,\varepsilon]$ and $x_{j+1}=\dots=x_N=0$. Applying the one-variable
inequality in $x_{j+1}$ with the other variables fixed and integrating,
\[
I_j(w)\le\frac2\varepsilon I_{j+1}(w)+2\varepsilon\Lambda^2|L|\sum_{l\in L}I_{j+1}(lw).
\]
Iterating from $I_0(\varnothing)$, each of the $N$ steps either keeps the word (factor $2/\varepsilon$) or prepends
one of $|L|$ letters (factor $2\varepsilon\Lambda^2|L|$ each); all words reached have length $\le N$, so with
$K=\max_{|w|\le N}I_N(w)$,
\[
I_0(\varnothing)\le\sum_{m=0}^N\binom Nm\Bigl(\frac2\varepsilon\Bigr)^{N-m}(2\varepsilon\Lambda^2|L|^2)^mK=\Bigl(\frac2\varepsilon+
2\varepsilon(|L|\Lambda)^2\Bigr)^NK\le a^NK,
\]
as $\varepsilon\le1$.
\end{proof}

\begin{theorem}[$L^2$ bound at the section]\label{pbo:schur-section}
Assume \textup{(GF)}. Let $(W,W^\flat)$ be a Bessel pair of degree $D$, $P\ge\max(8r,b^5)$, $c>0$, and
$Z$ such that every $Z'$ with $\|Z'-Z\|<P^{-6}$ is $P$-bounded with $\pbg(Z')\ge c$. Let $\varphi\in
L^1(G)$ vanish where $\|y\|>P^{20}$. Then, with $\mathcal B=\{(h,v)\in H\times\pboU:\|v\|\le P^{20}\}$,
\[
\int_{\mathcal B}\Bigl\|\int_G\varphi(y)W\bigl(\gamma_Z(vh)^{-1}y\bigr)dy\Bigr\|^2\pbocG c^{k-1}dh\,dv\le
P^{2000n^2}\frac dD\int_G|\varphi|^2 .
\]
\end{theorem}
\begin{proof}
Take $(V_w)$ from (BS4) for this $P$. For words $i$ in the alphabet $\mathcal E^{(2)}$ of letters
$((a,b),c,d)$ define weights $\psi_\varnothing=\varphi$, $\psi_{((a,b),c,d)i}(y)=\psi_i(y)\,(y^{-1})_{ad}\,y_{cb}/
P^{40}$; then $|\psi_i|\le|\varphi|$ since $|y_{cb}|,|(y^{-1})_{ad}|\le\|y\|\le P^{20}$ on the support. Put
$\mathcal W_i(g)=\int_G\psi_i(y)V_{\bar i}(gy)\,dy$, $\bar i$ the word of first components. If $c$ is a $C^1$
curve with $c'=cX$, then $(c(t)y)'=c(t)y\,(y^{-1}Xy)$, so by (BS4)(i) and differentiation under the
integral (the $V_w$ are bounded, $\psi_i$ integrable),
\[
\frac{d}{dt}\mathcal W_i(c(t))=\sum_{a,b,c,d}P^{40}X_{dc}\,\mathcal W_{((a,b),c,d)i}(c(t)),
\]
since $(y^{-1}Xy)_{ab}=\sum_{c,d}(y^{-1})_{ad}X_{dc}y_{cb}$. Hence $\|\frac d{dt}\mathcal W_i(c(t))\|\le P^{40}
\|X\|\sum_\ell\|\mathcal W_{\ell i}\|$. Choose real linear coordinates $x\mapsto Z+e(x)$ on $M_n(\C)$
($N=2n^2$ coordinates, the real and imaginary parts of the entries), measure preserving, with $\|e(x)\|
\le2n^2\|x\|_\infty$; with $\varepsilon=P^{-8}$ the cube maps into the ball of radius $P^{-6}/8$ about $Z$, and
$4\|e(e_j)\|\le P^2$. For $(h,v)\in\mathcal B$ and $x$ in the cube, $Z'=Z+e(x)$ is $P$-bounded, so by
Lemma~\ref{pbo:sizes}(c) (with $T=P$; the ball hypothesis gives the bounds at every point of the segment)
the curve $t\mapsto\gamma_{Z+e(x+te_j)}(vh)^{-1}$ has logarithmic derivative of norm $\le\|\gamma_{Z'}^{-1}\|\,
\|(vh)^{-1}\|\,P^{25}\cdot\frac{P^2}4\cdot\|vh\|\le P^{18+20+25+2+20}=P^{85}$. Lemma~\ref{pbo:sobolev} with
$f_i(x,(h,v))=\mathcal W_i(\gamma_{Z+e(x)}(vh)^{-1})$, $\Lambda=P^{40}P^{85}$, $|L|=r^4\le P^4$, gives the
factor $a^N$ with $a(2P^2)\le\frac2{P^{-8}}(1+P^{258})2P^2\le P^{275}$. For each word $i$ with $|i|\le2n^2$,
by Lemma~\ref{pbo:section-int} with $Q$ the ball of radius $P^{-6}$ about $Z$ and $w=\|\mathcal W_i\|^2$
(left invariant under $R_\tau$ by \eqref{pbo:eq-inv}), and by (BS4)(ii),
\[
\int_{\rm cube}\int_{\mathcal B}\|f_i\|^2\pbocG c^{k-1}\le(2P^2)^{2n^2}\int_{U\backslash G}\|\mathcal W_i\|^2\le
(2P^2)^{2n^2}P^{4n^2}\frac dD\int|\varphi|^2 .
\]
Altogether the left side of the theorem is at most $(P^{275})^{2n^2}P^{4n^2}\frac dD\int|\varphi|^2\le P^{2000n^2}
\frac dD\int|\varphi|^2$.
\end{proof}

\subsection{The bound at one place}\label{pbo:ss-onplace}

\begin{lemma}[Tubes]\label{pbo:tube}
For $T\ge8r$ there are a compact $A\subseteq G$ and a compact $\Omega_0\subseteq\pboU$ of positive measure
with: $\|y\|\le T^{10}$ on $A$; $vy\in A$ whenever $v\in\Omega_0$ and $\|y\|\le T$; and $\pbvol\{v\in\pboU:vy\in
A\}\le3\pbvol(\Omega_0)$ for every $y$.
\end{lemma}
\begin{proof}
\emph{A gauge on $\pboU$.} For $u=u_q$ let $A(u)=\sigma-\frac12\beta\beta^*$ be the $\tilde e$-coefficient matrix in
\eqref{pbo:eq-uq} and $N(u)=\max(\|\beta\|^2,\|A(u)\|)$. From \eqref{pbo:eq-uq} and $u_qx=x+\tilde e\beta$,
$u_qu_{q'}$ has parameters $\beta+\beta'$ and $A=A(u)+A(u')-\beta\beta'^*$; hence $\sqrt{N(uu')}\le\sqrt{N(u)}+
\sqrt{N(u')}$. Also $A(u^{-1})=A(u)^*$ (from $A+A^*=-\beta\beta^*$), so $N(u^{-1})=N(u)$. As $\beta$ and $A$ are
blocks of $u-1$, $N(u)\le\max(\|u-1\|^2,\|u-1\|)\le(\|u\|+1)^2$; and $\|u\|\le1+2\|\beta\|+\|A\|\le(1+\sqrt{N(u)})^2$.
Conjugation by the element $a_t$ of $G$ with $a_t\tilde e_a=t\tilde e_a$, $a_t\tilde e^*_a=t^{-1}\tilde e^*_a$, fixing
$x_i,\tilde z$, maps $u_{(\beta,\sigma)}$ to $u_{(t\beta,t^2\sigma)}$; it multiplies $N$ by $t^2$ and $du^0$ by $t^Q$,
$Q=2\varrho k+\varrho^2$ (real dimension of the parameters). So the balls $B(\rho)=\{N\le\rho\}$ satisfy
$\pbvol B(\rho')=(\rho'/\rho)^{Q/2}\pbvol B(\rho)$, and they are compact with non-empty interior.

\emph{The tube.} Let $\rho=Q^2(T^2+1)^2$, $\Omega_0=B(\rho)$, $A=\Omega_0\cdot\{y:\|y\|\le T\}$, compact. On $A$,
$\|vy\|\le(1+\sqrt\rho)^2T\le(1+Q(T^2+1))^2T\le T^9\le T^{10}$, using $Q\le r^2/2\le T^2/128$. Fix $y$ and $v_0$ with
$v_0y=v_1y_1\in A$ ($v_1\in\Omega_0$, $\|y_1\|\le T$); if also $vy=v_2y_2\in A$, then $w:=v_2^{-1}vv_0^{-1}v_1=y_2y_1^{-1}$
lies in $\pboU$ and has $\|w\|\le T^2$, so $\sqrt{N(w)}\le T^2+1$, and $\sqrt{N(vv_0^{-1}v_1)}=\sqrt{N(v_2w)}\le\sqrt\rho+
T^2+1$. Thus $\{v:vy\in A\}\subseteq B(\rho')g^{-1}$ with $g=v_0^{-1}v_1$, $\sqrt{\rho'}=\sqrt\rho+T^2+1$. By right
invariance of $du^0$, its measure is at most $(\rho'/\rho)^{Q/2}\pbvol(\Omega_0)=(1+1/Q)^Q\pbvol(\Omega_0)\le
e\,\pbvol(\Omega_0)\le3\pbvol(\Omega_0)$.
\end{proof}

\begin{proof}[Proof of Proposition~\ref{pbo:local}]
Let $Z_0=\tilde Z(g)$. As $|Z_{0,ab}|=|h(y^{-1}\tilde f_a,y^{-1}g\tilde f_b)|\le\rho^2$ and $n\le\rho$, put
$M=\rho^{9n}/\pbg(Z_0)$; by Lemma~\ref{pbo:sizes}(a$'$) with $M=\rho$, $\pbg(Z_0)\le\rho^{8n}$, so $M\ge\rho$, and $M^{-1}\le
\pbg(Z_0)$, $|Z_{0,ab}|\le M^2$. Let $M_0=M^{9n}$, $P=M_0^{5n}$; then $P\ge M\ge\rho\ge b\ge8r$ and $P\ge
b^5$. By Lemma~\ref{pbo:witt} write $g=s\gamma_{Z_0}s'^{-1}$ with $s,s'\in\pboR$, and for $x=(h,v)\in H\times
\pboU$ put $t(x)=s'h^{-1}v^{-1}$, so that $gt(x)=s\gamma_{Z_0}(vh)^{-1}$. By Lemma~\ref{pbo:sizes}(a),(b)
(with $T=P/2\ge M^{8n}$), every $Z'$ with $\|Z'-Z_0\|<P^{-6}$ is $P$-bounded with $\pbg(Z')\ge c:=\pbg(Z_0)/2$.
Let $A,\Omega_0$ be a tube for $T=M_0^5$ (Lemma~\ref{pbo:tube}); then $\|y\|\le P^{20}$ on $A$ and $\Omega_0
\subseteq A$. Let $\mu_c=\pbocG c^{k-1}dh\,dv$, $m=\mu_c(\cdot\,|\,H\times\Omega_0)$, a probability measure, and
\[
\mu=\frac{P^{2000n^2}d}{D\,\mu_c(H\times\Omega_0)}\,\mathbf 1_A\|W\|^2\,dg .
\]
(a) Fix $\varphi$ bounded measurable. By \eqref{pbo:eq-inv} and the Bessel property, for fixed $x$ the
Hilbert--Schmidt norm squared of $\int_A\varphi(z)W(gt(x)z)W(t(x)z)^*dz$ equals $\sum_j\|\int_G\beta_j(z)
W(\gamma_{Z_0}(vh)^{-1}z)dz\|^2$ with $\beta_j=\mathbf 1_A\varphi\,\aubar{W_j}$: expand in the coordinates $j$
of the second factor, by \textup{(BS1)}, $W(t(x)z)=\pi_R(t(x))W(z)$ and $W(gt(x)z)=\pi_R(s)W(\gamma_{Z_0}(vh)^{-1}z)$ with
$\pi_R(t(x))$, $\pi_R(s)$ unitary (where $\pi_R(uh)=\psi(u)h$), and the Hilbert--Schmidt norm is invariant
under unitary factors on either side. As $\beta_j$ vanishes off $A$, Theorem~\ref{pbo:schur-section} applies (with
$\mathcal B\supseteq H\times\Omega_0$, since $\|v\|\le T^{10}\le P^{20}$ on $\Omega_0$):
\[
\int_T\cdots dm\le\frac1{\mu_c(H\times\Omega_0)}P^{2000n^2}\frac dD\sum_j\int|\beta_j|^2=\int|\varphi|^2d\mu .
\]
(b) Let $x=(h,v)$ with $v\in\Omega_0$, and $z$ with $\|z^{-1}\tilde f_a\|\le\rho\le M$ and $\pbel(gt(x)z)\le\rho
\le M$. By Lemma~\ref{pbo:norm-z} (with $R=M_0^{\,}\ge M^{9n}$; $\|Z_0^{-1}\|\le M^{8n}\le M_0$) $\|v^{-1}z\|\le
M_0^5=T$, so $z=v(v^{-1}z)\in A$.

\emph{The mass of $\mu$.} By Lemma~\ref{pbo:unfold} with $\varphi=\mathbf 1_A$, $w=\|W\|^2$ (left invariant), $K=3
\pbvol(\Omega_0)$ (tube) and $R=P^{40}$ (if $y,sy\in A$ then $\|s\|\le P^{40}$, so $|m(s)_{ab}|\le P^{40}$), and
(BS2), $\int_A\|W\|^2\le3\pbvol(\Omega_0)(2P^{40})^{2n^2}d$. Since $\mu_c(H\times\Omega_0)=\pbocG c^{k-1}\pbvol(
\Omega_0)$, the volume of $\Omega_0$ cancels:
\[
\mu(G)\le\frac{3(2P^{40})^{2n^2}P^{2000n^2}d^2}{D\,\pbocG\,(\pbg(Z_0)/2)^{k-1}}\le M^{10^5n^4}\frac{d^2}{D\pbocG(
\pbg(Z_0)/2)^{k-1}}=\pbsp(Z_0;D,\rho,d),
\]
because $3(2P^{40})^{2n^2}P^{2000n^2}\le M_0^5(M P^{40})^{2n^2}P^{2000n^2}=M^{93600n^4+2n^2+45n}\le M^{10^5n^4}$
(with $P=M^{45n^2}$, $M_0^5=M^{45n}$, $3\le M_0^5$, $2\le M$). Finally $A$ is compact, $t$ is measurable
and $m$ is a probability measure.
\end{proof}

\subsection{Proof of Theorem~\ref{pbo:O1}}\label{pbo:ss-proof}

By Theorem~\ref{pbo:system}, $(W_\tau,W^\flat_\tau)$ is a Bessel pair of degree $\pbdeg_\tau$ at every $\tau$,
and $\rho_\tau\ge\pbsth\ge b$. The kernel $\pbsK^\flat_\gamma(y)$ is the adjoint of $\pbsK^{\flat\prime}(y):=
W^\flat(\gamma^{-1}y)W^\flat(y)^*$, and the adjoint is an isometry of the Hilbert--Schmidt norm commuting
with integration; and $\pbg(\gamma_\tau^{-1})=\pbg(\gamma_\tau)$ since $\tilde Z(g^{-1})=\tilde Z(g)^*$. So it
suffices to prove the bound for $\int\bar\phi\,\pbsK^{\flat\prime}$ with $g_\tau=\gamma^{-1}_\tau$ and the point
$y_1=\gamma^{-1}y_0$ in place of $y_0$ (the hypotheses are symmetric).

Apply Proposition~\ref{pbo:local} at each $\tau$ to $g_\tau$, $\rho_\tau$, obtaining $A_\tau,(T_\tau,m_\tau),t_\tau,
\mu_\tau$. For $\mathbf x=(x_\tau)\in\prod T_\tau$ let $t(\mathbf x)=(t_\tau(x_\tau))$ and
\[
\mathcal N(\mathbf x)=\Bigl\|\int_{\prod A_\tau}\bar\phi(z)\,W(g\,t(\mathbf x)z)\,W(t(\mathbf x)z)^*\,dz\Bigr\|^2 .
\]
\emph{Products of contractions.} If kernels $k_\tau(x_\tau,z_\tau)$ (matrix valued) satisfy
$\int\|\int\varphi k_\tau\|^2dm_\tau\le\int|\varphi|^2d\mu_\tau$ for all bounded $\varphi$, then the tensor product
kernel satisfies $\int\|\int\varphi\bigotimes k_\tau\|^2\prod dm_\tau\le\int|\varphi|^2\prod d\mu_\tau$ for bounded
$\varphi$ on the product: by induction on the number of factors, integrate first over $z_{\tau_1}$ against the
kernel of $\tau_1$ with the other variables frozen as a vector-valued function ($\varphi$ replaced by the
partial integral, which is a bounded function of $z_{\tau_1}$ with values in a Hilbert space; the
contraction property for scalar $\varphi$ extends to Hilbert-space-valued $\varphi$ coordinatewise), then apply
the inductive hypothesis. Hence $\int\mathcal N\,d\prod m_\tau\le\int|\phi|^2d\prod\mu_\tau\le\prod_\tau\mu_\tau(
G_\tau)$.

\emph{One good shift.} The set of $\mathbf x$ for which (b) fails at some $\tau$ is null; as $\prod m_\tau$ is a
probability measure, there is $\mathbf x$ outside it with $\mathcal N(\mathbf x)\le\prod_\tau\mu_\tau(G_\tau)$.
Let $t=t(\mathbf x)\in\prod\pboR_\tau$.

\emph{Support.} If $\phi(tz)W^\flat(gtz)W^\flat(tz)^*\ne0$, then (hypothesis on $\phi$, via $y=tz$ after the
symmetrisation) $\|(tz_\tau)^{-1}\tilde f_a\|=\|z_\tau^{-1}\tilde f_a\|\le\rho_\tau$, and $W^\flat_\tau(g_\tau t_\tau
z_\tau)\ne0$ so $\pbel_\tau(g_\tau t_\tau z_\tau)\le b^{6n+2}\le\rho_\tau$ by (BS3); so $z\in\prod A_\tau$ by (b).
Hence $z\mapsto\bar\phi(tz)\pbsK^{\flat\prime}(tz)$ vanishes off the compact set $\prod A_\tau$, it is continuous
there up to the measurable factor $\phi$, and $\int_{G_\infty}\bar\phi\,\pbsK^{\flat\prime}=\int_{\prod A_\tau}\bar\phi(tz)
\pbsK^{\flat\prime}(tz)\,dz$ by left invariance of Haar measure; also $\phi(tz)=\phi(z)$.

\emph{Comparison with the uncut kernel.} Writing $\pbsK^{\flat\prime}=\pbsK'-(\pbsK'-\pbsK^{\flat\prime})$ on
$\prod A_\tau$,
\[
\Bigl\|\int\bar\phi\,\pbsK^{\flat\prime}\Bigr\|\le\mathcal N(\mathbf x)^{1/2}+\int_{G_\infty}|\phi|\,\|\pbsK'-
\pbsK^{\flat\prime}\|\le\Bigl(\prod_\tau\pbsp(\tilde Z(g_\tau);\pbdeg_\tau,\rho_\tau,d_\tau)\Bigr)^{1/2}+\int|\phi|\,
\|\pbsK_\gamma-\pbsK^\flat_\gamma\| ,
\]
using $\pbsp(\tilde Z(g^{-1}))=\pbsp(\tilde Z(g))$ ($\pbg$ is unchanged) and that the adjoint preserves the norm of
the difference. That $\phi\pbsK^\flat_\gamma$ vanishes outside a compact set follows from the support statement (translated back by $t$ and by the symmetrisation); it is bounded since $W$, $W^\flat$ are bounded, and measurable as a product of measurable functions; $\pbsK^\flat_\gamma$ is continuous by \textup{(BS1)}. \qed

\subsection{The Bessel functions form Bessel pairs}\label{pbo:ss-pairs}

Fix $\tau$; $c=c_\tau$, $\pbdeg=\pbdeg_\tau$, $\pi=\pi_\tau$, $v_0$ as in \S\ref{pbe:ss-ds}, $da=dg$ the Haar measure
of \textup{(DS1)}. Let $\mathfrak g_0$ be the Lie algebra of $G$ (skew-hermitian for $\mathsf J_1$), with the
operator norm. For $X\in\mathfrak g_0$ write $dX$ for the right derivative $(dXf)(g)=\frac d{dt}f(ge^{tX})|_{t=0}$.
Let $L_\tau=\max_i|2\varsigma_{\tau,i}|$.

\begin{lemma}[Derivatives of $K$-finite vectors]\label{pbo:K}
Let $\mathcal H^s$ be the span of the vectors $d\pi(X_1)\cdots d\pi(X_j)v_0$, $j\le s$, $X_i\in\mathfrak g_0$
(these are $K$-finite smooth vectors). For $x\in\mathcal H^s$ and $X\in\mathfrak g_0$,
\[
\|d\pi(X)x\|\le(r+1)^4\,(L_\tau+r+1+s)\,\|X\|\,\|x\|,\qquad d\pi(X)x\in\mathcal H^{s+1}.
\]
\end{lemma}
\begin{proof}
$G=\mathrm U(k+\varrho,\varrho+1)$ is connected, $K=\mathrm U(k+\varrho)\times\mathrm U(\varrho+1)$ is connected and
contains the diagonal torus $T_c$, a compact Cartan subgroup; $\pi$ is a discrete series representation with
Harish-Chandra parameter $\varsigma$ (regular, entries in $\frac12\Z$; the chamber is fixed in \S\ref{pbe:ss-ds}).
Use the trace form $B(X,Y)=\operatorname{tr}(XY)$ on $\mathfrak g$, the Cartan decomposition $\mathfrak g_0=
\mathfrak k_0\oplus\mathfrak p_0$, and coordinates on $\mathfrak t^*$ in which the roots are $\epsilon_i-\epsilon_j$.
We use two published results.

(i) \emph{Lowest $K$-type} (Schmid \cite[Lemma~(2.22), Lemma~(3.8), Corollary~(3.9), Corollary~(4.19) and the proof of Theorem~(1.3) on p.~560]{Schmid75}; see also \cite[Theorem~(1.3), p.~130]{HechtSchmid75}; for groups outside their standing hypotheses these apply to $\auU(p,q)=Z\cdot\mathrm{SU}(p,q)$, $Z$ the compact centre acting by a character, through $\mathrm{SU}(p,q)$, see the paragraph ``On (D3)'' in \S\ref{par:ss-rigid}): for $G$ connected semisimple, or reductive with compact centre component as here, with a
compact Cartan subgroup, the discrete series with Harish-Chandra parameter $\varsigma$ contains the $K$-type of
highest weight $\mu_0=\varsigma+\rho_n-\rho_c$ with multiplicity one, $\rho_n,\rho_c$ the half sums of the
noncompact and compact roots positive for $\varsigma$. Its hypotheses hold by the above. As $v_0$ is a highest
weight vector of this $K$-type and the entries of $\rho_n$, $\rho_c$ have absolute value $\le r/2$, $|\mu_{0,i}|\le
L_\tau/2+r$.

(ii) \emph{Casimir} (Harish-Chandra; see \cite[Theorem~8.18, (8.32) and the Example on p.~225]{Knapp}): on the $(\mathfrak g,K)$-module of $\pi$ the Casimir element $\Omega_G=\sum
X_iX^i$ (dual bases for $B$) acts by $|\varsigma|^2-|\rho|^2$; on an irreducible representation of the
connected group $K$ of highest weight $\mu$, $\Omega_K$ acts by $|\mu+\rho_K|^2-|\rho_K|^2$, with $\rho_K=\rho_c$ for
the compact positive roots. (Norms of the standard inner product in the coordinates $\epsilon_i$.)

\emph{$K$-types of $\mathcal H^s$.} The adjoint representation of $K$ on $\mathfrak p_\C$ has weights $\pm(\epsilon_i-
\epsilon_j)$ (one index in each block), with entries in $\{0,\pm1\}$. If $x$ lies in a sum of $K$-types of highest
weights $\mu$, then $d\pi(X)x$ ($X\in\mathfrak p$) lies in the image of $\mathfrak p_\C\otimes V_\mu$, whose
irreducible constituents have highest weights $\mu+\beta$ with $\beta$ a weight of $\mathfrak p_\C$ (a constituent's
highest weight is a weight of the tensor product, hence of the form (weight of $V_\mu$)$+\beta$, and dominance
bounds it entrywise by $\max|\mu_i|+1$); $d\pi(X)$, $X\in\mathfrak k$, preserves isotypic components. So every
$K$-type $\mu$ of $\mathcal H^s$ has $\max_i|\mu_i|\le L_\tau/2+r+s=:\Lambda_s$.

\emph{$\mathfrak k$-part.} For $X\in\mathfrak k_0$, conjugate by $K$ to $X=\operatorname{Ad}(k)H$, $H\in\mathfrak t_0$ with
$\|H\|=\|X\|$; on $V_\mu$, $d\pi(H)$ is diagonal in a weight basis with eigenvalues $\langle\nu,H\rangle$, $|\langle\nu,
H\rangle|\le\sum_i|\nu_i||H_i|\le r\Lambda_s\|X\|$ (the weights $\nu$ of $V_\mu$ satisfy $|\nu_i|\le\max|\mu_j|$). Hence $\|d\pi(X)
x\|\le r\Lambda_s\|X\|\|x\|$.

\emph{$\mathfrak p$-part.} Let $(Y_j)$ be a basis of $\mathfrak p_0$ orthonormal for $B$ (positive on $\mathfrak p_0$) and
$(Z_i)$ a basis of $\mathfrak k_0$ with $B(Z_i,Z_{i'})=-\delta_{ii'}$; then $\Omega_G=\sum Y_j^2-\sum Z_i^2$ and $\Omega_K=-\sum Z_i^2$, so
$-\sum_jd\pi(Y_j)^2=\Omega_K-\Omega_G$ on $K$-finite vectors. For $X=\sum a_jY_j\in\mathfrak p_0$, $(\sum a_j^2)^{1/2}=\|X\|_{HS}\le
\sqrt r\|X\|$, and by skew-adjointness of $d\pi(Y_j)$ on smooth vectors,
\[
\|d\pi(X)x\|^2\le\|X\|_{HS}^2\sum_j\|d\pi(Y_j)x\|^2=\|X\|_{HS}^2\langle(\Omega_K-\Omega_G)x,x\rangle\le r\|X\|^2(|\mu+\rho_K|^2_{\max}+|\rho|^2)\|x\|^2,
\]
where $|\mu+\rho_K|^2\le r(\Lambda_s+r)^2$ and $|\rho|^2\le r\cdot r^2$; so $\|d\pi(X)x\|\le\sqrt2\,r(\Lambda_s+r)\|X\|\|x\|$.

\emph{Total.} For $X\in\mathfrak g_0$, the components $X_{\mathfrak k}=\frac12(X+\theta X)$, $X_{\mathfrak p}=\frac12(X-\theta X)$ have norm $\le\|X\|$, so
$\|d\pi(X)x\|\le(r+\sqrt2r)(\Lambda_s+r)\|X\|\|x\|\le3r\cdot(L_\tau+r+1+s)\cdot2\cdot\|X\|\|x\|$, and $6r\le(r+1)^4$. The same
holds for the complex matrix units $E_{ab}=X+iX'$ with $\|X\|,\|X'\|\le1$, at the cost of a factor $2$, still within
$(r+1)^4$.
\end{proof}

\begin{lemma}[Schur orthogonality for translates]\label{pbo:translates}
Let $\psi=\pbdeg\cdot c$. Then $\psi(g^{-1})=\aubar{\psi(g)}$, $\psi(1)=\pbdeg$, and for all $t,g\in G$
\begin{equation}\label{pbo:eq-schur}
\int_G\psi(s)\,\psi(t^{-1}s^{-1}g)\,ds=\frac{\psi(t^{-1})}{\pbdeg}\,\psi(g);
\end{equation}
in particular ($t=1$) $\psi*\psi=\psi$. For a Bessel function $F$ with $\int_{U\backslash G}\|F\|<\infty$ and $F*\psi=F$, where
$(F*k)(y)=\int k(a^{-1})F(ya)\,da$, the translates satisfy
\[
\int_{U\backslash G}\langle F(ya),F(yb)\rangle\,dy=\frac{\psi(a^{-1}b)}{\pbdeg}\int_{U\backslash G}\|F\|^2 .
\]
\end{lemma}
\begin{proof}
The first two identities are \textup{(DS1)(i)}. \eqref{pbo:eq-schur} is Schur's orthogonality relation
$\int_G\langle u,\pi(s)v\rangle\,\aubar{\langle u',\pi(s)v'\rangle}\,ds=\pbdeg^{-1}\langle u,u'\rangle\aubar{\langle v,v'\rangle}$
(conjugate-linear first variable; $\pi$ square integrable and irreducible, $\int|c|^2=\pbdeg^{-1}$), applied to
$u=v_0$, $v=v_0$, $u'=\pi(t)v_0$, $v'=\pi(g)v_0$ after writing $\psi(t^{-1}s^{-1}g)=\pbdeg\,\aubar{\langle\pi(g)v_0,
\pi(s)\pi(t)v_0\rangle}$; it contains \textup{(DS1)(ii)}.
Let $k(g)=\int_{U\backslash G}\langle F(y),F(yg)\rangle dy$; by the transformation law the integrand is left
$U_\tau$-invariant, and the quotient measure is right invariant, so the left side is $k(a^{-1}b)$; $k$ is
bounded, $k(g^{-1})=\aubar{k(g)}$, and $F*\psi=F$ gives $k(g)=\int\psi(t^{-1})k(gt)\,dt$ and $k(g)=\int\psi(t)
k(t^{-1}g)\,dt$. Combining, $k(g)=\int\int\psi(s)\psi(t^{-1})k(s^{-1}gt)\,ds\,dt$, and (substitute $s^{-1}gt=u$ and use \eqref{pbo:eq-schur} in the form $\int\psi(s)\psi(u^{-1}s^{-1}g)\,ds=\psi(u^{-1})\psi(g)/
\pbdeg$) yields $k(g)=\frac{\psi(g)}{\pbdeg}\int\psi(u^{-1})k(u)\,du=\frac{\psi(g)}{\pbdeg}k(1)$.
\end{proof}

\begin{proof}[Proof of Theorem~\ref{pbo:system}]
(BS1) is Proposition~\ref{pbe:local}(a) and Lemma~\ref{pbe:len}(a) (for $W^\flat=\chi W$, $\chi$ continuous and
$R_\tau$-invariant). (BS2) is the normalisation of $W_\tau$. (BS3) is Proposition~\ref{pbe:local}(e).
(BS5) is a size condition (\S\ref{pbo:ss-sizes}).

(BS4). Let $\psi=\pbdeg\,c$. By Theorem~\ref{pbe:bessel}(B7), $\int e(a)W(ya)\,da=W(y)$ with $e=\pbdeg\,\bar c$,
i.e.\ $W*\psi=W$, since $\psi(a^{-1})=\pbdeg\,\aubar{c(a)}=e(a)$. Choose $\beta\in C^\infty_c(G)$ with $\pi(\beta)
v_0=v_0$, where $\pi(\beta)=\int\beta(g)\pi(g)\,dg$ (it exists because $v_0$ is $K$-finite and
$Z(\mathfrak g)$-finite: by Harish-Chandra's theorem \cite[Theorem~1, p.~18]{HC66}, used in this form in
\cite[Ch.~I, \S4, proof of Lemma~14, p.~15]{HarishChandraAF},
a smooth function $f$ on $G$ that is left $K$-finite and $Z(\mathfrak g)$-finite satisfies $f=\alpha*f$ for some
$\alpha\in C^\infty_c(G)$, $(\alpha*f)(x)=\int\alpha(y)f(y^{-1}x)\,dy$ (the published statement is $f^\vee*\alpha'=f^\vee$
for the right $K$-finite function $f^\vee(x)=f(x^{-1})$; put $\alpha(y)=\alpha'(y^{-1})$, as in (D5) of
\S\ref{ptq:ss-H}); apply this to $f(x)=\langle u,\pi(x^{-1})v_0\rangle$ for a fixed $u\ne0$, which is
linear in $v_0$ (the inner product is conjugate-linear in the first variable), smooth, left $K$-finite
($f(kx)=\langle u,\pi(x^{-1})\pi(k^{-1})v_0\rangle$) and $Z(\mathfrak g)$-finite (the left regular action of
$z\in Z(\mathfrak g)$ on $f$ is the action of an element of $Z(\mathfrak g)$ on $v_0$, $Z(\mathfrak g)$ being stable
under the principal antiautomorphism), because $v_0$ is smooth, $K$-finite and $Z(\mathfrak g)$-finite. Then
$(\alpha*f)(x)=\int\alpha(y)\langle u,\pi(x^{-1})\pi(y)v_0\rangle dy=\langle u,\pi(x^{-1})\pi(\alpha)v_0\rangle$, so
$\pi(\alpha)v_0-v_0$ is orthogonal to $\pi(G)u$, which spans a dense subspace as $\pi$ is unitary and
irreducible; take $\beta=\alpha$). For a word $w=(\ell_1,\dots,\ell_m)$ of matrix units $E_{ab}$ let $x_w=d\pi(E_{\ell_1})\cdots
d\pi(E_{\ell_m})v_0$ (complexified action) and $X_w(g)=\pbdeg\,\langle v_0,\pi(g)x_w\rangle$; $X_\varnothing=\psi$.
Then $X_w=\psi*X_w$ (as $x_w$ lies in the irreducible $\pi$: Schur orthogonality), $X_w$ is integrable (it is
$\psi*(\text{derivative of }\check\beta)$, a convolution of $L^1$ functions), and by Schur orthogonality and
Lemma~\ref{pbo:K},
\[
\frac1{\pbdeg}\int_G|X_w|^2=\|x_w\|^2\le\bigl((r+1)^4(L_\tau+r+1+|w|)\bigr)^{2|w|}\le P^{2|w|}
\]
for $|w|\le2n^2$ and $P\ge b^5$, by the size condition (Z4). Put $V_w=W*X_w$. These are Bessel functions
(left translation commutes with right convolution; bounded as $W$ is bounded and $X_w\in L^1$). For a curve
with $c'=cX$, $X=\sum X_{ab}E_{ab}$, $V_w(c(t))=\int W(a)X_w(a^{-1}c(t))\,da$ and the right derivative of
$X_w$ along $E_{ab}$ is $X_{(a,b)w}$ (as $\frac d{dt}\pi(ge^{tE})x=\pi(g)d\pi(E)x$); differentiation under the
integral is justified as follows. From $v_0=\pi(\beta)v_0$ we get $x_w=\pi(\beta_w)v_0$ with $\beta_w\in
C^\infty_c(G)$ (left derivatives of $\beta$), so $X_w(h)=\int\beta_w(a)\psi(ha)\,da$. For $c(t)$ in a compact
set $C$, $\sup_t|X_{(a,b)w}(a'^{-1}c(t))|\le\|\beta_{(a,b)w}\|_1\sup\{|\psi(a'^{-1}cb')|:c\in C,\ b'\in
\operatorname{supp}\beta_{(a,b)w}\}\le C'\bigl(2/\|a'\|\bigr)^{\pbeq}$ by \textup{(DS2)} with $j=0$ and
$\|a'^{-1}cb'\|\ge\|a'\|/(\|c\|\,\|b'\|)$, for a constant $C'$ depending on $C$ and $w$; this bound is
integrable in $a'$ by \textup{(DS3)}, and $W$ is bounded, so dominated convergence applies. This is
(BS4)(i).

For (BS4)(ii): by Lemma~\ref{pbo:translates} the vectors $\omega_a=F(\cdot\,a)$ of $L^2(U\backslash G)$, $F=W$, have
$\langle\omega_a,\omega_b\rangle=\frac{\psi(a^{-1}b)}{\pbdeg}d$. Let $\mathcal Z=X_w*\check\varphi$; then $g\mapsto\int\varphi(y)
V_w(gy)dy=W*\mathcal Z$, and in $L^2(U\backslash G)$ this function $\Phi$ satisfies $\|\Phi\|^2=\langle\Phi,\int
\mathcal Z(a^{-1})\omega_a\,da\rangle\le\|\Phi\|\,\|\int\mathcal Z(a^{-1})\omega_a\,da\|$. The vector $V=\int
\mathcal Z(a^{-1})\omega_a\,da=\int\varphi(g)\,\Big(\int X_w(a^{-1})\omega_{ga}\,da\Big)dg$ lies in the closed span of the
$\omega_a$; for $u=\sum\lambda_i\omega_{a_i}$, $\langle u,\int X_w(a^{-1})\omega_{ga}da\rangle=\frac d{\pbdeg}\sum\bar\lambda_iX_w(a_i^{-1}g)$
(by $\psi*X_w=X_w$), whose $L^2(dg)$-norm squared is $\frac{d^2}{\pbdeg^2}\cdot\frac{\pbdeg\|u\|^2}{d}\cdot\frac{\int|X_w|^2}{\pbdeg}$
by Schur orthogonality ($\int X_w(a^{-1}g)\aubar{X_w(b^{-1}g)}dg=\frac{\psi(a^{-1}b)}{\pbdeg}\int|X_w|^2$). By
Cauchy--Schwarz in $g$, $|\langle u,V\rangle|^2\le\int|\varphi|^2\cdot\frac{d}{\pbdeg}\frac{\int|X_w|^2}{\pbdeg}\|u\|^2$, and by
density $\|V\|^2\le\frac{\int|X_w|^2}{\pbdeg}\cdot\frac d{\pbdeg}\int|\varphi|^2\le P^{2|w|}\frac d{\pbdeg}\int|\varphi|^2$.
This is (BS4)(ii) with $D=\pbdeg$.
\end{proof}

\subsection{Size conditions}\label{pbo:ss-sizes}

Besides those of \S\S\ref{pbe:sec}, \ref{pbs:ss-sizes}, the proof uses:
\begin{itemize}
\item[(Z4)] $8r\le b$ (BS5), and $(r+1)^4(L_\tau+r+1+2n^2)\le b^5$ for every $\tau$, where $L_\tau$ is the
largest absolute value of an entry of $2\varsigma_\tau$. The first is in \S\ref{pbe:sec} ($r\le b$ is listed there;
$8r\le b$ holds since $r\le3x^5$ and $b\ge x^{10^{12}}$ by Lemma~\ref{an:lem-sizes}). The second holds as soon
as $L_\tau\le b^4$, by the same comparison; and $L_\tau\le b^2$ holds: the entries of $\varsigma_\tau$ are $k-1$ numbers of absolute value $<k\le3x^5$ and
the numbers $-\kappa(\tau)-a\Delta$ with $|\kappa(\tau)|\le b$, $1\le a\le m\le2\sqrt x$, $\Delta=2r^8\le2\cdot3^8x^{40}$
(\S\ref{ss:au-notation}, Lemma~\ref{an:lem-sizes}(S2),(S5)), so $L_\tau\le2(b+m\Delta)\le b^2$.
\end{itemize}


%% ===== proofs/B-frames-real.tex =====
% proofs/B-frames-real.tex
% The local theory of two frames (shared by the real and the finite places) and the sieve system at
% a real place: Theorem pbm:inputs (R).
% Labels of this file start with pfr:.
\providecommand{\panot}[1]{\par\smallskip{\raggedright\noindent\textbf{[Not proved here:} #1\textbf{]}\par}\smallskip}
\providecommand{\pbdeg}{\operatorname{deg}}
\providecommand{\pbvol}{\operatorname{vol}}
\providecommand{\pbmTh}{\Theta}
% macros of this file (prefix \pfr)
\providecommand{\pfrG}{\mathsf G}
\providecommand{\pfrH}{\mathsf H}
\providecommand{\pfrU}{\mathsf U}
\providecommand{\pfrR}{\mathsf R}
\providecommand{\pfrM}{\mathsf M}
\providecommand{\pfrm}{\mathsf m}
\providecommand{\pfrK}{\mathsf K}
\providecommand{\pfrP}{\mathcal P}
\providecommand{\pfrCl}{\mathcal C}
\providecommand{\pfrI}{\mathcal I}
\providecommand{\pfrF}{\mathsf F}
\providecommand{\pfrB}{\mathsf B}
\providecommand{\pfrmod}{\operatorname{mod}}
\providecommand{\pfrdep}{\operatorname{dp}}
\providecommand{\pfrwt}{\wp}
\providecommand{\pfrkap}{\kappa}
\providecommand{\pfrS}{\mathsf S}
\providecommand{\pfrL}{\Lambda}
\providecommand{\pfrN}{\mathsf N}
\providecommand{\pfrX}{\Xi}
\providecommand{\pfrip}[2]{\langle#1,#2\rangle}
\providecommand{\pfrhw}{\mathsf w}
\providecommand{\pfrell}{\ell}

\section{Two frames: the local theory, and the sieve system at a real place}\label{pfr:sec}

This section concerns part \textup{(R)} of Theorem~\ref{pbm:inputs}, and the local theory of
\emph{pairs of frames} that parts \textup{(R)} and \textup{(F)} of that theorem share.

\emph{What is proved.}
\begin{itemize}
\item[(I)] \emph{The exact statement of \textup{(R)}.} Theorem~\ref{pfr:R} states \textup{(R)} with
its hypotheses. Its constants are \emph{exactly} those of Theorem~\ref{pbm:inputs}\,(R) as used in
Theorem~\ref{pbm:sieve}, Lemma~\ref{pbm:const} and Theorem~\ref{pbm:L2}: mass $(1+b^{-10})/\pbdeg_\tau$,
weights $\Lambda^\infty$ with the radius $b^{60n^2}$ and the exponent $n^6$, decay $\beta^\infty$ with
the constant $\pbmTh\,b^{2000n^9}$ and the exponent $k-n^7$; the top form is the Gram form of the
normalised orbital integrals described in \S\ref{pbm:ss-remains}. The comparison is made term by term
in Remark~\ref{pfr:rem-match}. Two points have to be added to the statement printed in
Theorem~\ref{pbm:inputs}: the identification \textup{(R-b)} of the top form, which Siegel's step uses,
and the hypotheses on the data under which \textup{(R)} holds (Remark~\ref{pfr:rem-match}).
\item[(II)] \emph{The reduction of \textup{(R)} to four local estimates is proved in full}
(Theorem~\ref{pfr:R}, \S\ref{pfr:ss-Rproof}): positivity, continuity and descent of the forms, the
mass, the decay in the frequency and the decay for unequal entries.
\item[(III)] \emph{The mass is proved} (Proposition~\ref{pfr:mass}), from the bound for the remote
positions (Proposition~\ref{pfr:remote}) and the comparison of the floor measure with Haar measure
(Lemma~\ref{pfr:haar}), both proved in \S\ref{prp:sec}.
\item[(IV)] \emph{The shared theory.} The objects (place data, floors, correlations, forms,
majorants, depths, gauges, shells) are defined in \S\ref{pfr:ss-setting}--\S\ref{pfr:ss-gauge}. The
sum over the shells (Theorem~\ref{pfr:shells}) is proved in full from the one-frame bound
(Theorem~\ref{pfr:oneframe}) and the gains from a deep vector. Theorems~\ref{pfr:nogain} and~\ref{pfr:deep}
are proved (\S\ref{pfr:ss-gainsproof}), and so is Lemma~\ref{pfr:depth-rare}, with all the facts on
floors, quadrics and sections they use (\S\ref{pfr:ss-floors}, \S\ref{pfr:ss-section}); Lemma~\ref{pfr:descent}
is proved (\S\ref{pfr:ss-descent}). Theorems~\ref{pfr:oneframe}, \ref{pfr:gain-char} and \ref{pfr:gain-heads}
are proved in \S\ref{pfs:sec}, which also supplies the facts on sections and on the skew Haar measure used below.
\item[(V)] \emph{The real place.} Lemmas~\ref{pfr:realgauge}, \ref{pfr:deepreal} and
Propositions~\ref{pfr:conv}, \ref{pfr:top}, \ref{pfr:formbound}, \ref{pfr:remote} are proved in
\S\ref{prp:sec}, and so is Lemma~\ref{pfr:haar} (using Lemmas~\ref{pbo:GF}, \ref{pbo:append}).
\end{itemize}

\subsection{Places, floors and place data}\label{pfr:ss-setting}

\subsubsection*{Places}
A \emph{place} is a locally compact commutative ring $A$ with an involution $c\mapsto\bar c$, a Haar
measure $dc$, a continuous unitary character $\psi$ of $(A,+)$, a \emph{modulus}
$\pfrmod\colon A\to[0,\infty)$ (the factor by which multiplication by $c$ changes $dc$) and a
constant $P\ge1$ (the constant of Plancherel's inequality: axiom~(P) of \S\ref{pfs:ss-planch}). The two places used
are:
\begin{itemize}
\item the \emph{real place} $A=\C$, $\bar c$ complex conjugation, $dc$ Lebesgue measure,
$\psi(c)=e^{c-\bar c}$, $\pfrmod(c)=|c|^2$, $P=\pi^2$;
\item the \emph{finite place} $A=E_T=\prod_{w\in T}E_w$ of \S\ref{pbm:ss-sieve}, with the measure
giving $\auO_T$ volume $1$, $\psi=\psi_T$, $\pfrmod(c)=|c|_T$; it is treated in
\textup{B-frames-finite}.
\end{itemize}
Each place also carries an element $\tau_0$ with $\tau_0+\bar\tau_0=1$ and a skew unit $\delta$
($\bar\delta=-\delta$): $\tau_0=\frac12$, $\delta=\frac i2$ at the real place, and the elements of
Lemma~\ref{pfg:psi}(b) at the finite place. Put $A^+=\{c:\bar c=c\}$ and $A_{\rm sk}=\{c:\bar c=-c\}$. We fix
the Haar measure $d\beta_{\rm sk}$ of $A_{\rm sk}$ of Lemma~\ref{pfs:skewhaar} (the one for which
$(\beta_1,\beta_2)\mapsto\tau_0\delta^{-1}\beta_1+\beta_2$ carries $d\beta_{\rm sk}\otimes d\beta_{\rm sk}$ to $dc$), and the
Haar measure $dh$ of $A^+$: Lebesgue measure of $\R$ at the real place, the measure giving $A^+\cap\auO_T$ volume $1$
at the finite place. The \emph{skew constant} $\kappa_A$ of the place is defined by: the isomorphism
$A^+\times A_{\rm sk}\to A$, $(h,b)\mapsto\tau_0h+b$, carries $dh\otimes d\beta_{\rm sk}$ to $\kappa_A\,dc$. At the
real place, $\tau_0\delta^{-1}=-i$, so $(it_1,it_2)\mapsto t_1+it_2$ is measure preserving and $d\beta_{\rm sk}$ is the
Lebesgue measure $dt$ of $\beta=it$; then $(h,it)\mapsto\frac h2+it$ has Jacobian $\frac12$, and $\kappa_\C=2$. At the
finite place $\kappa_{E_T}=1$, by Lemma~\ref{pfg:duality}(d) (whose proof shows that $\beta\mapsto\delta\beta$ carries
the measure of $A^+$ to that of $A_{\rm sk}$). At both places multiplication by $c\bar c$ ($c$ a unit) multiplies $dh$
by $\pfrmod(c)$ (at the finite place by Lemma~\ref{pfg:duality}(d), as $|c\bar c|_T^{1/2}=|c|_T$).

Hermitian forms follow convention (Cv1) of \S\ref{pb:ss-setting} throughout this section: $h$ is linear in the
first variable, so the coordinate of a vector $x$ along $e_a^*$ is $h(x,e_a)$ and along $e_a$ it is $h(x,e_a^*)$.

\subsubsection*{The space and its floors}
Let $\varrho\ge1$, $n=\varrho+1$, and let $V$ be a free $A$-module with a basis
\[
e_0,\dots,e_{\varrho-1};\quad z;\quad (x_i)_{i\in\iota};\quad e^*_{\varrho-1},\dots,e^*_0
\]
and a hermitian form $h$ for which the $(e_a,e^*_a)$ are hyperbolic pairs
($h(e_a,e^*_b)=\delta_{ab}$, $h(e_a,e_b)=h(e_a^*,e_b^*)=0$), orthogonal to $z$ and to the $x_i$,
$h(z,z)=-1$ at the real place, and the $x_i$ are an orthonormal family orthogonal to $z$ (at the
finite place the middle block is the given unimodular lattice; see \textup{B-frames-finite}). Put
$k=\#\iota$ and $r=\operatorname{rank}V=2\varrho+k+1$. The \emph{frame} is $f_a=e_a$ $(a<\varrho)$,
$f_\varrho=z$. Let $\pfrG=\auU(V,h)$.

The \emph{floor} of $e_a$ and of $e_a^*$ is $a$, the floor of $z$ is $\varrho$ and the floor of the
$x_i$ is $n=\varrho+1$. For $l\ge0$ let $\pfrG^{(l)}$ be the subgroup of the $g\in\pfrG$ that fix the
basis vectors of floor $<l$; so $\pfrG^{(0)}=\pfrG$, $\pfrG^{(\varrho)}=\pfrG^\flat$ fixes the
$e_a,e_a^*$, and $\pfrH=\pfrG^{(n)}$ fixes them and $z$. Let $\pfrR^{(l)}$ be the group of the $g\in\pfrG$
that change each $f_a$, $a<l$, by a combination of the $f_b$, $b<a$ (that is, $gf_a\in f_a+\langle
f_b:b<a\rangle$), let $\pfrU=\pfrU_{(<\varrho+1)}$ be the Bessel unipotent group (the block upper
unitriangular elements for the flag of the frame, with middle block $1$), and for $u\in\pfrU$ let
$\sigma(u)$ be the sum of the entries of $u$ just above the diagonal of its block on the frame, as in
\eqref{pb:eq-psi}. For $l\le\varrho$ the \emph{Heisenberg group of floor $l$}, $\pfrN_l\subseteq
\pfrG^{(l)}\cap\pfrU$, is the unipotent radical of the stabiliser in $\pfrG^{(l)}$ of the line
$Af_l$. For $l<\varrho$ its elements are the $n(\omega,\beta)$, with $\omega$ in the span $V_{l+1}$ of the
basis vectors of floor $>l$ and $\beta\in A$ with $\beta+\bar\beta=-h(\omega,\omega)$, acting by $f_l\mapsto f_l$,
$v\mapsto v-h(v,\omega)f_l$ ($v\in V_{l+1}$), $f_l^*\mapsto f_l^*+\omega+\beta f_l$, and as the identity on the basis
vectors of floor $<l$ (\S\ref{pfr:ss-floors}). Writing $\beta=b-\tau_0h(\omega,\omega)$ with $b\in A_{\rm sk}$, the
measure $\eta_l$ is the image of $d\omega\otimes d\beta_{\rm sk}(b)$ under $(\omega,b)\mapsto n(\omega,b-\tau_0h(\omega,\omega))$,
where $d\omega$ is the Haar measure of $V_{l+1}\cong A^{d_l}$ in the coordinates of the basis ($d_l=\operatorname{rank}
V_{l+1}$; Lebesgue measure at the real place, the measure giving the lattice volume $1$ at the finite place); it is a
Haar measure of $\pfrN_l$ ((N2) of \S\ref{pfr:ss-floors}). At the real place, with $b=it$, this is the image of
$d\omega\,dt$. For $l\ge\varrho$ we put $\pfrN_l=1$, $\eta_l=\delta_1$.

For a sequence $p=(p_0,p_1,\dots)$ of units of $A$ and $j\le\varrho$ the \emph{Weyl element}
$w_p=w_{j,p}$ is $e_a\mapsto p_ae^*_a$, $e_a^*\mapsto\bar p_a^{\,-1}e_a$ $(a<j)$, and the identity on the
other basis vectors; and $w_{q}^{-1}w_p$ is the \emph{Weyl ratio}, the diagonal element with entries
$p_a/q_a$ at $e_a$ and $\overline{q_a/p_a}$ at $e_a^*$ $(a<j)$.

\subsubsection*{Place data}
A \emph{place datum} $\pfrF$ consists of
\begin{itemize}
\item a measure $\theta_\varrho$ on $\pfrG^{(\varrho)}$ (the \emph{top floor measure}) whose image under
$\sigma\mapsto\sigma^{-1}z$ is $c_\varrho\mu_\varrho$ for a constant $0<c_\varrho\le2$, where $\mu_\varrho$ is the
quadric measure of $z$ defined below;
\item a non-zero measure $\mu_\pfrH$ carried by $\pfrH$ and invariant under right translation by
$\pfrH$;
\item a unit $M$ of $A$ with $\bar M=M$ (the \emph{level});
\item an open set $\pfrCl$ of pairs of vectors containing $0$, with compact closure (the \emph{close
pairs});
\item a compact set $\pfrK\subseteq\pfrG$ (the \emph{box}) and a measurable set $\pfrP\subseteq\pfrG$
(the \emph{positions});
\end{itemize}
such that $\mu_\pfrH\{h:\exists u\in\pfrU,\ uhg\in\pfrK\}\le1$ for every $g\in\pfrG$. Its
\emph{floor measures} are
\[
\theta_l=\lambda_l*\eta_l\ (l<\varrho),\qquad\theta_\varrho,\qquad\theta_{n}=\mu_\pfrH,\qquad
\theta_l=\delta_1\ (l>n),
\]
where $\lambda_l$ is the image of the measure $\pfrmod(c)^{d_l}\,dc$ on $A^\times$ under $c\mapsto w_{l,c}^{-1}$
(the law of the inverse Weyl elements of floor $l$; $w_{l,c}$ is the Weyl factor $e_l\mapsto ce_l^*$,
$e_l^*\mapsto\bar c^{-1}e_l$, the identity on the other basis vectors), so that
\[
\int F\,d\theta_l=\iint F\bigl(w_{l,c}^{-1}\,n\bigr)\,\pfrmod(c)^{d_l}\,dc\,d\eta_l(n)\qquad(l<\varrho,\ F\ge0\text{ Borel}).
\]
\emph{Quadric measures.} For $l\le\varrho$ let $V_l$ be the span of the basis vectors of floor $\ge l$, with the Haar
measure $dy$ in the coordinates of the basis (Lebesgue measure at the real place; the measure giving the lattice
spanned by these basis vectors volume $1$ at the finite place). The \emph{quadric measure} $\mu_l$ of $f_l$ is the
fibre, at the value $h(f_l,f_l)$, of the disintegration of $dy$ along $y\mapsto h(y,y)$ (with $dh$ on $A^+$):
\begin{equation}\label{pfr:eq-slab}
\int\varphi\,d\mu_l=\lim_{B\to0}\frac1{\pbvol(B)}\int_{\{y\in V_l:\ h(y,y)-h(f_l,f_l)\in B\}}\varphi(y)\,dy,
\end{equation}
for $\varphi$ continuous with compact support (locally constant with compact support at the finite place), where
$B=(-\varepsilon,\varepsilon)$, $\varepsilon\to0$, at the real place and $B=\varpi_1^{\ell}(A^+\cap\auO_T)$, $\ell\to\infty$, at the
finite place ($\varpi_1$ of Lemma~\ref{pfg:psi}(c)). The limit exists and defines a positive Radon measure carried by
the quadric $\{h(y,y)=h(f_l,f_l)\}$: at the finite place this is Proposition~\ref{pfg:fibres}(a),(b) (the lattice of
$V_l$ is unimodular of rank $\ge3$); for $l<\varrho$ it also follows from Lemma~\ref{pfr:floorquad} below, and for
$l=\varrho$ at the real place from the computation in its proof. As $dy$ and the sets in
\eqref{pfr:eq-slab} are invariant under $\auU(V_l)$, \emph{$\mu_l$ is invariant under $\auU(V_l)\supseteq\pfrG^{(l)}$}
(restricted to $V_l$). This is the measure called the disintegration of $dy$ in Lemma~\ref{pfr:block}.
The constants of the floor measures are $c_l$ for $l\le\varrho$: for $l<\varrho$, $c_l=\kappa_A$ (Lemma~\ref{pfr:floorquad});
$c_\varrho$ is the constant of the top floor measure. At the real place $c_l=2$ for all $l\le\varrho$
(\S\ref{pfr:ss-realforms}); at the finite place $c_l=1$ for all $l\le\varrho$ (\S\ref{pff:ss-large}). In all cases
$0<c_l\le2$, and the image of $\theta_l$ under $\sigma\mapsto\sigma^{-1}f_l$ is $c_l\mu_l$ ($l\le\varrho$).

For measures $\nu_a$ on $\pfrG$ write
$\nu_{[l,l+i)}=\nu_{l+i-1}*\dots*\nu_l$ (convolution on $\pfrG$). Put
\[
\pfrM=\bigl((\eta*\theta)\bigr)_{[0,\varrho+2)},\qquad\pfrm_j=\bigl((\eta*\theta)\bigr)_{[0,j)},\qquad
\theta_{\ge l}=\theta_{[l,\varrho+2)},
\]
where $(\eta*\theta)_a=\eta_a*\theta_a$: the \emph{group measure} $\pfrM$, the measure $\pfrm_j$ of the
floors below $j$, and the measure $\theta_{\ge l}$ that represents the quotient
$(\pfrU\cap\pfrG^{(l)})\backslash\pfrG^{(l)}$. We always assume that $\pfrM$ is a Haar measure of
$\pfrG$ (at the real place this is Lemma~\ref{pfr:haar}; at the finite place see \S\ref{pff:ss-large}).

\begin{lemma}[Floor measures and quadrics]\label{pfr:floorquad}
Let $l<\varrho$. For every $\varphi$ as in \eqref{pfr:eq-slab} the limit \eqref{pfr:eq-slab} exists, and
\[
\int\varphi(\sigma^{-1}f_l)\,d\theta_l(\sigma)=\kappa_A\int\varphi\,d\mu_l .
\]
That is, the image of $\theta_l$ under $\sigma\mapsto\sigma^{-1}f_l$ is $c_l\mu_l$ with $c_l=\kappa_A$: $c_l=2$ at the
real place and $c_l=1$ at the finite place.
\end{lemma}
\begin{proof}
\emph{Coordinates.} In the basis of $V_l$ (the vectors $f_l$, $f_l^*$ and the basis of $V_{l+1}$) write
$y=sf_l+\omega'+tf_l^*$ with $s,t\in A$,
$\omega'\in V_{l+1}$; then $dy=ds\,d\omega'\,dt$ and, by (Cv1), $h(y,y)=s\bar t+\bar st+h(\omega',\omega')$. The non-units
$t$ form a null set (axiom (P) of \S\ref{pfs:ss-planch}); for a unit $t$ write $s=t(\tau_0h+b)$ with $h\in A^+$,
$b\in A_{\rm sk}$, and $h'=t\bar t\,h$. Then $s\bar t+\bar st=t\bar t(\tau_0h+b+\bar\tau_0h-b)=h'$, and, for fixed
$t,\omega'$, the map $(h',b)\mapsto s$ carries $dh'\otimes d\beta_{\rm sk}$ to $\kappa_A\,ds$: indeed $(h,b)\mapsto\tau_0h+b$
carries $dh\otimes d\beta_{\rm sk}$ to $\kappa_A\,du$, $u\mapsto tu$ multiplies $du$ by $\pfrmod(t)$, and $h\mapsto t\bar t\,h$
multiplies $dh$ by $\pfrmod(t)$ (\S\ref{pfr:ss-setting}). Write $Y(t,\omega',b,h')=sf_l+\omega'+tf_l^*$ for the
point with these coordinates, $s=\tau_0h'\bar t^{-1}+tb$. Hence, with $c_0=h(f_l,f_l)=0$,
\[
\begin{gathered}\frac1{\pbvol(B)}\int_{\{h(y,y)\in B\}}\varphi\,dy=\frac1{\kappa_A}\iint F_B(t,\omega')\,dt\,d\omega',\\
F_B(t,\omega')=\frac1{\pbvol(B)}\iint\mathbf 1\bigl[h'+h(\omega',\omega')\in B\bigr]\,\varphi\bigl(Y(t,\omega',b,h')\bigr)\,
db\,dh'.\end{gathered}
\]
\emph{A bound.} Let $\varphi$ vanish outside the set of the $y$ whose coordinates lie in a compact set $R_0$
($R_0=\{|c|\le R\}$, resp.\ $R_0=\varpi_1^{-a}\auO_T$), and let $B\subseteq B_0$, $B_0=(-1,1)$ resp.\ $A^+\cap\auO_T$.
For fixed $t,\omega',h'$ the set of the $b$ with $Y(t,\omega',b,h')$ in the support of $\varphi$ has
$d\beta_{\rm sk}$-measure at most $C_0\pfrmod(t)^{-1/2}$ with $C_0$ depending only on $R_0$: at the real place $s=
\frac{h'}{2\bar t}+tb$ with $b=ix$, and $\frac{h'}{2\bar t}=\frac{h't}{2|t|^2}$ is orthogonal to $tix$ in $\C=\R^2$, so
$|s|\le R$ forces $|x|\le R/|t|$; at the finite place the set is empty or a coset of
$\{b:tb\in\varpi_1^{-a}\auO_T\}$, of measure at most $|\varpi_1^at|_T^{-1/2}$ by Lemma~\ref{pfg:duality}(d). As the
$h'$ with $h'+h(\omega',\omega')\in B$ form a set of measure $\pbvol(B)$,
\[
|F_B(t,\omega')|\le\sup|\varphi|\cdot C_0\,\pfrmod(t)^{-1/2}\,\mathbf 1[t\in R_0,\ \omega'\in R_0^{d_l}],
\]
and $\int_{R_0}\pfrmod(t)^{-1/2}dt<\infty$ (at the real place $\int_{|t|\le R}|t|^{-1}dt=2\pi R$; at the finite place
the integral is a finite product over $w\in T$ of the sums $\sum_{j\ge-a'}(1-q_w^{-1})q_w^{-j}q_w^{j/2}$).
\emph{The limit.} Fix $t,\omega'$ and put $h'_0=-h(\omega',\omega')$ and $G(h')=\int\varphi(Y(t,\omega',b,h'))\,db$. The
function $G$ is continuous at $h'_0$: at the real place by dominated convergence (the $b$ with non-zero integrand lie in
a fixed bounded set, by the bound above); at the finite place $\varphi$ is invariant under translation by a lattice
$\varpi_1^{m}L$ ($L$ the lattice of $V_l$), and for $h'$ so close to $h'_0$ that $\tau_0(h'-h'_0)\bar t^{-1}\in\varpi_1^m\auO_T$
the points $Y(t,\omega',b,h')$ and $Y(t,\omega',b,h'_0)$ differ by $\tau_0(h'-h'_0)\bar t^{-1}f_l\in\varpi_1^mL$ for every $b$, so
$G(h')=G(h'_0)$. As $F_B(t,\omega')$ is the mean of $G$ over
$h'_0+B$, $F_B(t,\omega')\to G(h'_0)$. By dominated convergence the limit \eqref{pfr:eq-slab} exists and
\[
\kappa_A\int\varphi\,d\mu_l=\iiint\varphi\bigl(Y(t,\omega',b,-h(\omega',\omega'))\bigr)\,db\,d\omega'\,dt .
\]
\emph{The image of $\theta_l$.} For $\sigma=w_{l,c}^{-1}n$ with $n=n(\omega,\beta)$, $\beta=b-\tau_0h(\omega,\omega)$, we have
$n^{-1}=n(-\omega,\bar\beta)$ (as $\beta+\bar\beta=-h(\omega,\omega)$) and $\sigma^{-1}f_l=n^{-1}(cf_l^*)=c(f_l^*-\omega+\bar\beta f_l)$. Its
coordinates are $t=c$, $\omega'=-c\omega$, $s=c\bar\beta$; with $h'=-h(\omega',\omega')=-c\bar c\,h(\omega,\omega)$ we get
$s=\tau_0h'\bar c^{-1}+cb'$ for $b'=\bar\beta+\tau_0h(\omega,\omega)=-b+(\tau_0-\bar\tau_0)h(\omega,\omega)\in A_{\rm sk}$. So
$\sigma^{-1}f_l=Y(c,-c\omega,b',-h(\omega',\omega'))$. For fixed $c$, $\omega\mapsto-c\omega$ multiplies $d\omega$ by
$\pfrmod(c)^{d_l}$, and for fixed $c,\omega$, $b\mapsto b'$ preserves $d\beta_{\rm sk}$. By the definition of $\theta_l$,
\[
\int\varphi(\sigma^{-1}f_l)\,d\theta_l(\sigma)=\iiint\varphi\bigl(Y(c,-c\omega,b',\cdot)\bigr)\pfrmod(c)^{d_l}\,dc\,d\omega\,db=
\iiint\varphi\bigl(Y(t,\omega',b',-h(\omega',\omega'))\bigr)\,db'\,d\omega'\,dt,
\]
which is $\kappa_A\int\varphi\,d\mu_l$. The values of $\kappa_A$ are those of \S\ref{pfr:ss-setting}.
\end{proof}

A point $g$ is \emph{near at $j$} if $\pfrR^{(j)}g$ meets $\pfrK$ (for $j\ge\varrho+2$ we put $\pfrR^{(j)}=\pfrU$). A function $W\colon\pfrG\to\C^I$ ($I$
a finite set) is a \emph{cut Bessel function} for $\pfrF$ if it is bounded and measurable,
$W(ug)=\psi(M^{-1}\sigma(u))W(g)$ for $u\in\pfrU$, and $W(g)\ne0$ only if $ug\in\pfrK$ for some
$u\in\pfrU$.

\subsubsection*{Correlations and forms}
For a kernel $\mathsf k$ on $\pfrG\times\pfrG$ the \emph{correlation from floor $l$ on} is
\[
\mathsf k_l(g,y)=\int_\pfrG\mathsf k(xg,xy)\,d\theta_{\ge l}(x).
\]
For a cut Bessel function $W$ and $d\in\pfrG$ put $\mathsf k^{d}(g,y)=\pfrip{W(g)}{W(dy)}$. Fix $j\le
\varrho$ and sequences of units $p,q$; write $w=w_{j,p}$ and $d=w_{j,q}^{-1}w_{j,p}$. The
\emph{head density} is
\[
\pfrkap_j(p,q)=\prod_{a<j}\pfrmod(p_aq_a)^{d_a/2},\qquad d_a=k+1+2(\varrho-a-1).
\]
For a position $x$, a point $g$ and a bounded measurable weight $\omega\colon\pfrG\to\C$ put
\[
\Phi_j(x;\omega;g)=\iint\omega(\sigma\tau^{-1})\,\mathsf k^1_{j+1}(\sigma g,\sigma gx)\,
\overline{\mathsf k^{d}_{j+1}(\tau w^{-1}g,\tau w^{-1}gx)}\,d\theta_j(\tau)\,d\theta_j(\sigma)
\]
The
\emph{form of floor $j$} is
\[
\pfrI_j(p,q;\omega)=\pfrkap_j(p,q)\int_{\pfrG\times\pfrG}\Phi_j(x;\omega;g)\,d\pfrM(x)\,d\pfrm_j(g),
\]
whenever the integral converges absolutely. Its \emph{integrand norm} and \emph{majorant} at the
position $x$ are
\[
\pfrF_j(x;\omega)=\int|\Phi_j(x;\omega;g)|\,d\pfrm_j(g),\qquad
\pfrF_j^+(x)=\int\pfrB^1_j(g,gx)\,\pfrB^{d}_j(w^{-1}g,w^{-1}gx)\,d\pfrm_j(g),
\]
where $\pfrB^{d}_j(g,y)=\int|\mathsf k^{d}_{j+1}(\sigma g,\sigma y)|\,d\theta_j(\sigma)$ is the
\emph{pivot norm}.

\begin{lemma}\label{pfr:L1}
Let $W$ be a cut Bessel function, $j\le\varrho$, $p,q$ sequences of units and $|\omega|\le1$. Then
$\pfrF_j(x;\omega)\le\pfrF^+_j(x)$ for every $x$, and, whenever $\int\pfrF_j^+\,d\pfrM<\infty$, the
form $\pfrI_j(p,q;\omega)$ is defined and
$|\pfrI_j(p,q;\omega)|\le\pfrkap_j(p,q)\int\pfrF_j(x;\omega)\,d\pfrM(x)$.
\end{lemma}
\begin{proof}
For fixed $g$, $|\Phi_j(x;\omega;g)|\le\iint|\mathsf k^1_{j+1}(\sigma g,\sigma gx)|\,
|\mathsf k^{d}_{j+1}(\tau w^{-1}g,\tau w^{-1}gx)|\,d\theta_j(\tau)\,d\theta_j(\sigma)$, and by
Tonelli's theorem the right side is $\pfrB^1_j(g,gx)\pfrB^{d}_j(w^{-1}g,w^{-1}gx)$. Integrating in $g$
gives the first claim. The second is the triangle inequality for the integral defining $\pfrI_j$, whose
integrand is then absolutely integrable by the first claim and Tonelli's theorem.
\end{proof}

\subsection{Depth, weights and gauges}\label{pfr:ss-gauge}

\subsubsection*{Stable vectors and depth}
Let $Y^{(l)}$ be the module of the vectors orthogonal to $e_0,\dots,e_{l-1}$ modulo their span. A
vector $\omega$ is \emph{stable at $(g,g')$ for $l$} if it is orthogonal to the $e_b$, $b<l$, and
$(g^{-1}\omega,g'^{-1}\omega')\in\pfrCl$ for some $\omega'\equiv\omega$ in $Y^{(l)}$. For $a<\varrho$ and
points $g$, $gx$ let
\[
\mathfrak d_a(g,x)=\{h(\omega,e_a):\ \omega\text{ stable at }(g,gx)
\text{ for }a\}\subseteq A\ \text{(the $e_a^*$-coordinates)},\qquad\pfrdep_a(g,x)=\pbvol(\mathfrak d_a(g,x))^{-1},
\]
the \emph{depth} of $e_a$ at the two points. At the real place $\mathfrak d_a$ is a disc and
$\pfrdep_a(g,x)=\pi^{-1}\inf_\zeta(|\zeta|_g^2+|\zeta|_{gx}^2)$, the infimum over the duals $\zeta$ of
$e_a$ in $Y^{(a)}$ (Lemma~\ref{pfr:realgauge}); at the finite places it is the index of an ideal.
Let $X'$ be the span of the orthonormal vectors $x_i$ left after two blocks of $2m+2$ of them are set
aside ($m$ the number of further hyperbolic pairs of the place datum; $m=0$ at the real place), and, at
the real place, also the vector $x_{i_0}$ of the boosts (\S\ref{pfr:ss-realdata}), and
$s=\operatorname{rank}X'-(3n+1+m)$, the \emph{volume exponent}. The \emph{depth weight} of length $j$ is
$\pfrwt_j(g,x)=\prod_{a<j}\pfrdep_a(g,x)^{-s}$.

\subsubsection*{Second moments}
For a kernel $\mathsf k$, a position $x$ and $j\le n$ let
\[
\mathsf a_j(g)=\pfrwt_j(g,x)\,|\mathsf k_{j+1}(g,gx)|^2,\qquad
\mathsf b_j(g)=\int\bigl(\mathsf a_j+\mathsf b_{j+1}\bigr)(\sigma g)\,d\theta_j(\sigma)\ (j<n),\quad
\mathsf b_j=0\ (j\ge n),
\]
and $\mathsf S_j(\mathsf k;x)=\int\mathsf b_j\,d\theta_{[0,j)}$, the \emph{weighted second moment} from
floor $j$ on.

\subsubsection*{Gauges}
A \emph{gauge} $\Gamma$ for the place datum $\pfrF$ is a filtration of $A$ by compact sets
$\Gamma(0)\subseteq\Gamma(1)\subseteq\dots$ (``the numbers of size $i$'': $|c|\le L^i$ at the real place
for a scale $L$, $L^ic$ integral at the finite places), together with a
number $\Lambda=\Lambda_\Gamma\in[2,\infty)$, the \emph{loss}, such that:
\begin{enumerate}
\item[(G0)] (\emph{size rules}) $\Gamma(i)\Gamma(i')\subseteq\Gamma(i+i')$; a sum of at most $r+1$ elements of
$\Gamma(i)$ lies in $\Gamma(i+1)$; $\Gamma(i)$ is stable under $c\mapsto-c$ and $c\mapsto\bar c$ and contains $0,1$;
the entries of the Gram matrix $J$ of the basis and of $J^{-1}$ lie in $\Gamma(1)$. A vector or matrix
\emph{has size $i$} if its entries lie in $\Gamma(i)$. Consequences used below: if $X,Y$ have sizes $i,i'$,
then $XY$ has size $i+i'+1$; an isometry $k$ of size $i$ has $k^{-1}=J^{-1}k^*J$ of size $i+4$; an
$l\times l$ determinant of entries of size $i$ ($l\le r+1$) has size $l(i+1)$ (Laplace expansion and
induction on $l$);
\item[(G1)] $m\le9n$, $P\ge1$, $M\in\Gamma(1)$, $\tau_0\delta^{-1}\in\Gamma(1)$ for a fixed element $\tau_0$ of trace one and a skew unit $\delta$, the basis and the elements of $\pfrK$ have entries in
$\Gamma(1)$;
\item[(G2)] for a close pair $(\omega_1,\omega_2)$ the vectors $\omega_1$ and $M^{-1}(\omega_1-\omega_2)$
have entries in $\Gamma(1)$; and if $M^{-1}\in\Gamma(1)$, for every $i$ there is a unit $c$ with
$c^{-1}\in\Gamma(i+2)$ such that $(c\omega_1,c\omega_2)$ is close for all vectors $\omega_1,\omega_2$ with
entries in $\Gamma(i)$;
\item[(G3)] $\pfrmod(u)\le\Lambda^i$ for units $u\in\Gamma(i)$, and $\pbvol\Gamma(i)\le\Lambda^{i+1}$;
\item[(G4)] (\emph{quadrics of a block}) Lebesgue measure on $A^{2m+2}$ disintegrates along the
standard hermitian form into measures each of which gives mass at most $\Lambda^{(2m+2)(i+1)}$ to the
vectors with entries in $\Gamma(i)$;
\item[(G5)] (\emph{common vectors}) for $g\in\pfrK$, a position $x\in\pfrP$ and $\varrho$ linear forms
$v$ on $X'$, let $\mathcal W$ be the set of the $\omega\in X'$ killed by $v$ with $(g^{-1}\omega,
(gx)^{-1}\omega)\in\pfrCl$, and let $\mathcal D\subseteq A$ be the set of the differences
$M^{-1}\bigl(h(\omega_1,\omega_1')-h(\omega_2,\omega_2')\bigr)$ for close pairs $(\omega_1,\omega_2)$,
$(\omega_1',\omega_2')$; there is a measurable $\phi\colon X'\to[0,\infty]$ such that for every $u\in X'$
\[
\pbvol\{h(\omega,u):\omega\in\mathcal W\}^{-1}\le\phi(u),\qquad
\pbvol\{c\in A:h(\omega,cu)\in\mathcal D\ \forall\omega\in\mathcal W\}\le\phi(u),
\]
and $\int_{\Gamma(i)^{X'}}\phi^s\le\Lambda^{99n(i+1)}$, the integral over the $u\in X'$ whose coordinates in the
orthonormal basis of $X'$ lie in $\Gamma(i)$, for the Haar measure in these coordinates.
\end{enumerate}
In (G5) the form $h$ is that of (Cv1): for $u=\sum_iu_ix_i$ and $\omega=\sum_i\omega_ix_i$ in $X'$,
$h(\omega,u)=\sum_i\omega_i\bar u_i=u^*\omega$, and $h(\omega,cu)=\bar c\,h(\omega,u)$. The set $\mathcal D$ is stable
under conjugation ($\bar M=M$, $\overline{h(\omega,\omega')}=h(\omega',\omega)$, and the roles of the two pairs of close
pairs may be exchanged), so the second set in (G5) is also $\{c:h(cu,\omega)\in\mathcal D\ \forall\omega\in\mathcal W\}$.
For a vector $y$ we write $y_{X'}=\sum_ih(y,x_i)x_i$ ($x_i$ running over the orthonormal basis of $X'$) for its
orthogonal projection on $X'$; its coordinates $h(y,x_i)$ are coordinates of $y$, and $h(\omega,y)=h(\omega,y_{X'})$
for $\omega\in X'$.
A sequence $p$ is a \emph{head of length $j$} for $\Gamma$ if $(p_0\cdots p_{l-1})^{-1}\in\Gamma(1)$ for
$l\le j$. The \emph{position ball} of $\Gamma$ is the set of the positions $x\in\pfrP$ such that
$M^{-1}(x-1)$ has entries in $\Gamma(1)$. Put $C_\Gamma=\Lambda^{10^6n^3}$ and
\[
\mathsf Q_\Gamma(x)=\mathsf Q_{\Gamma,j,p,q,W}(x)=\Bigl(\prod_{a<j}\mathsf c_a(p_a^{-1})\Bigr)^2C_\Gamma^3\,
\mathsf S_0(\mathsf k^1;x)\,\mathsf S_j(\mathsf k^{d};x),\qquad\mathsf c_l(g)=\pfrmod(g)^{d_l+1/2},
\]
the \emph{weighted square product}; we say that a number $f\ge0$ \emph{has the gain $e$ at $x$} if
$f^2\le e^2\mathsf Q_\Gamma(x)$. Put also
\[
\mathsf N(W)=\Bigl(\prod_{l<\varrho}\mathsf c_l(M)\Bigr)P^{\,\varrho\operatorname{rank}X'}\iint
|\pfrip{W(g)}{W(y)}|^2\,d\theta_{\ge0}(g)\,d\theta_{\ge0}(y),
\]
the \emph{shell normalisation}. A point $g$ is \emph{active} at $x$ (for $j,p,q,W$) if
$\pfrip{W(g)}{W(gx)}\ne0$ and $w_p^{-1}g$, $w_q^{-1}gx$ are near at $j$; the position $x$ \emph{has a
deep vector of depth $D$} if every active point $g$ has an $a$ with $j\le a<\varrho$ and
$\pfrdep_a(g,x)\ge D$. Finally, for $V\ge0$, \emph{the mass of floor $j$ is at most $V$} if
$\theta_j\{\sigma:\sigma g\text{ near at }j+1\}\le V$ for every $g$.

\subsection{The shared estimates}\label{pfr:ss-shared}

Throughout this subsection $\pfrF$ is a place datum whose group measure $\pfrM$ is a Haar measure,
$\Gamma$ a gauge with loss $\Lambda$, $W$ a cut Bessel function, $j\le\varrho$, $p,q$ sequences of units,
and $\mathsf Q=\mathsf Q_{\Gamma,j,p,q,W}$, $\pfrF^+_j$, $\mathsf N(W)$ are as in \S\ref{pfr:ss-gauge}.
These are the statements that \textup{B-frames-finite} also uses.

\begin{lemma}[Descent]\label{pfr:descent}
Let $i<\varrho$ and assume $\int\pfrF^+_i\,d\pfrM<\infty$ (for the given $p,q$). For $c\in A^\times$ let
$p^c$, $q^c$ be $p$, $q$ with the entry of index $i$ replaced by $c$. Then for every $\xi\in A$ the
function $c\mapsto\psi(\xi c)\pfrI_{i+1}(p^c,q^c;1)$ is integrable on $A$ (the complement of $A^\times$
has measure $0$), and
\[
\int_A\psi(\xi c)\,\pfrI_{i+1}(p^c,q^c;1)\,dc=\pfrI_i\bigl(p,q;\ g\mapsto\psi(\xi\,h(gf_i,f_i))\bigr).
\]
\end{lemma}
(Proof in \S\ref{pfr:ss-descent}.)

\begin{theorem}[One frame]\label{pfr:oneframe}
Let $q$ be a head of length $j$ for $\Gamma$, $\Gamma'$ a second gauge, $S$ a subset of the position
ball of $\Gamma'$, and $f\ge0$ measurable with the gain $e$ at every $x\in S$. Then
\[
\pfrkap_j(p,q)\int_Sf\,d\pfrM\le e\,\Lambda^{1500000n^3+1}\,\Lambda_{\Gamma'}^{25000n^2}\,\mathsf N(W).
\]
\end{theorem}
(Proof in \S\ref{pfs:ss-oneframe}.)

\begin{theorem}[One position, no gain]\label{pfr:nogain}
Let $p$ be a head of length $j$ for $\Gamma$, $x\in\pfrP$, and let the mass of floor $j$ be at most $V$.
Then $\pfrF^+_j(x)^2\le V^2\,\mathsf Q(x)$.
\end{theorem}

\begin{theorem}[The gain from a deep vector]\label{pfr:deep}
Let $p$ be a head of length $j$ for $\Gamma$, $x\in\pfrP$, let the mass of floor $j$ be at most $V\ge1$,
and let $0<\varsigma$ be such that $x$ has a deep vector of depth $\varsigma^{-1}$. Then
$\pfrF^+_j(x)^2\le\varsigma^{\,s}\,V^2\,\mathsf Q(x)$. If $j=0$, then moreover
$\pfrF^+_0(x)^2\le\varsigma^{\,2s}\,\mathsf Q(x)$ (at floor $0$ the two frames coincide and both pivots gain).
\end{theorem}
(Proofs in \S\ref{pfr:ss-gainsproof}.)

\begin{theorem}[The gain from a frequency]\label{pfr:gain-char}
Let $p$ be a head of length $j$ for $\Gamma$, $x\in\pfrP$, $\xi$ a unit and $\eta\in A$ with
$\psi(\eta c)=\psi(\xi c)$ for all $c\in\Gamma(200n^2)$. Then
\[
\pfrF_j\bigl(x;\ g\mapsto\psi(\eta\,h(gf_j,f_j))\bigr)^2\le\Bigl(\frac{P}{\pfrmod(\xi)}\Bigr)^{\operatorname{rank}X'}
\mathsf Q(x).
\]
\end{theorem}

\begin{theorem}[The gain from two heads]\label{pfr:gain-heads}
Let $j=i+1$, $p$ a head of length $j$ for $\Gamma$, $x\in\pfrP$, and assume that
$\xi=M^{-1}(p_i^{-1}-q_i^{-1})$ is a unit. Then
$\pfrF_j(x;1)^2\le(P/\pfrmod(\xi))^{\operatorname{rank}X'}\,\mathsf Q(x)$.
\end{theorem}
(Proofs in \S\ref{pfs:ss-gains}; the constant $P$ is that of axiom~(P) of \S\ref{pfs:ss-planch}.)

\begin{lemma}[Deep vectors are rare]\label{pfr:depth-rare}
For $x\in\pfrP$, $a<\varrho$ and every $g$,
$\int_{\{\sigma:\ \sigma g\text{ near at }a+1\}}\pfrdep_a(\sigma g,x)^s\,d\theta_a(\sigma)\le\Lambda^{10^4n}$.
\end{lemma}
We use two facts on the floor measures, stated exactly.

\begin{lemma}[Charts of a floor]\label{pfr:chart}
Let $a\le\varrho$ and let $c\ge0$ be measurable with $c(\xi g)=c(g)$ for $\xi\in\pfrR^{(a+1)}$. Suppose that for
every $g\in\pfrK$ there is a measurable $\Psi_g\ge0$ on vectors with $c(\sigma g)\le\Psi_g(\sigma^{-1}f_a)$ for
$\sigma\in\pfrG^{(a)}$ and $\int_{\Gamma(7)}\Psi_g\,d\mu_a\le C$, where $\mu_a$ is the measure of the quadric
of $f_a$ and $\Gamma(7)$ the vectors with entries of size $7$. Then
$\int_{\{\sigma:\sigma g\text{ near at }a+1\}}c(\sigma g)\,d\theta_a(\sigma)\le c_aC\le2C$ for every $g$ (so
$\le C$ at the finite place, where $c_a=1$).
\end{lemma}

\begin{lemma}[From the quadric to $X'$]\label{pfr:quadricX}
Let $a\le\varrho$. For measurable $\phi\ge0$ on $X'$ with $\int_{\Gamma(11)^{X'}}\phi^s\le B$, the function
$\Psi(v)=\phi(v_{X'})^s$, $v_{X'}$ the orthogonal projection of $v$ on $X'$ (\S\ref{pfr:ss-gauge}), satisfies
$\int_{\Gamma(7)}\Psi\,d\mu_a\le\Lambda^{99n\cdot13}B$.
\end{lemma}
(Lemma~\ref{pfr:chart} is proved in \S\ref{pfr:ss-floors}.)
(Lemma~\ref{pfr:quadricX} is proved in \S\ref{pfr:ss-section}.)

\begin{proof}[Proof of Lemma~\ref{pfr:depth-rare}]
Let $c(g)=\pfrdep_a(g,x)^s$; it is $\pfrR^{(a+1)}$-invariant by (C3) of \S\ref{pfr:ss-gainsproof}. Let
$g\in\pfrK$ and let $\phi$ be the function of gauge axiom (G5) for $g$, $x$ and $v=0$, and $\mathcal W$ its set
of common vectors. For $\sigma\in\pfrG^{(a)}$ and $\omega\in\mathcal W$, the vector $\sigma\omega$ is orthogonal to the
$e_b$, $b<a$ ($\omega\in X'$ is orthogonal to them and $\sigma$ fixes them), and $(\,(\sigma g)^{-1}\sigma\omega,
(\sigma gx)^{-1}\sigma\omega)=(g^{-1}\omega,(gx)^{-1}\omega)$ is close; so $\sigma\omega$ is stable at $(\sigma g,\sigma gx)$,
and its $e_a^*$-coordinate is $h(\sigma\omega,f_a)=h(\omega,\sigma^{-1}f_a)=h(\omega,u)$ with $u=(\sigma^{-1}f_a)_{X'}$. Hence
$\mathfrak d_a(\sigma g,x)\supseteq\{h(\omega,u):\omega\in\mathcal W\}$, and $\pfrdep_a(\sigma g,x)\le\phi(u)$ by the first
inequality of (G5). By Lemma~\ref{pfr:quadricX} with $B=\Lambda^{99n\cdot12}$ (axiom (G5) at size $11$) and
Lemma~\ref{pfr:chart} with $\Psi_g(v)=\phi(v_{X'})^s$, the integral is at most
$c_a\Lambda^{99n\cdot25}\le2\Lambda^{2475n}\le\Lambda^{10^4n}$ (as $\Lambda\ge2$).
\end{proof}

\subsubsection*{Shells}
\begin{definition}\label{pfr:shell}
Let $H\in[1,\infty)$. A set $S\subseteq\pfrG$ is a \emph{shell of height $H$} (for $\Gamma,j,p,q,W$)
if
\begin{enumerate}
\item[(a)] $S$ is covered by the position balls of a family of at most $2H^{9n}$ gauges $\Gamma'$ with
$\Lambda_{\Gamma'}\le\Lambda H^2$;
\item[(b)] for every $\varsigma\in(0,1)$ with $\Lambda^{1200n^2}\le\varsigma^{400n^2}H$, every $x\in S$
has a deep vector of depth $\varsigma^{-2}$.
\end{enumerate}
\end{definition}
By (a) and the definition of the position balls, a shell lies in $\pfrP$.

\begin{lemma}[One shell]\label{pfr:oneshell}
Let $S$ be a shell of height $H$, $q$ a head of length $j$, and $f\ge0$ measurable with the gain $e$
at every point of $S$. Then
\[
\pfrkap_j(p,q)\int_Sf\,d\pfrM\le e\,\Lambda^{2\cdot10^6n^3}H^{10^5n^2}\,\mathsf N(W).
\]
\end{lemma}
\begin{proof}
Let $\mathcal G$ be the family of (a). By Theorem~\ref{pfr:oneframe} applied to $S\cap(\text{position
ball of }\Gamma')$ for each $\Gamma'\in\mathcal G$, and $\Lambda_{\Gamma'}\le\Lambda H^2$,
\[
\pfrkap_j\int_{S}f\,d\pfrM\le\sum_{\Gamma'\in\mathcal G}e\,\Lambda^{1500000n^3+1}(\Lambda H^2)^{25000n^2}
\mathsf N(W)\le2H^{9n}\,e\,\Lambda^{1500000n^3+25000n^2+1}H^{50000n^2}\mathsf N(W).
\]
As $2\le\Lambda$ and $H\ge1$, the right side is at most $e\,\Lambda^{2\cdot10^6n^3}H^{10^5n^2}\mathsf N(W)$.
\end{proof}

\begin{lemma}[A remote shell]\label{pfr:remoteshell}
Let $S$ be a shell of height $H>\Lambda^{2400n^2}$, let $p,q$ be heads of length $j$, let the mass of
floor $j$ be at most $V\ge1$, and let $\tau\in(0,1)$ with $\tau^{800n^2}=H^{-1}$. Then
\[
\pfrkap_j(p,q)\int_S\pfrF^+_j\,d\pfrM\le V\,\tau^{\,s}\,\Lambda^{2\cdot10^6n^3}H^{10^5n^2}\,\mathsf N(W).
\]
\end{lemma}
\begin{proof}
$\tau<1$ because $H>1$. Moreover $\Lambda^{2400n^2}\le H=\tau^{800n^2}H^2$, so
$\Lambda^{1200n^2}\le\tau^{400n^2}H$, and by (b) every $x\in S$ has a deep vector of depth $\tau^{-2}$.
By Theorem~\ref{pfr:deep} with $\varsigma=\tau^2$, $\pfrF^+_j(x)$ has the gain $V\tau^{s}$ at every
$x\in S$; Lemma~\ref{pfr:oneshell} concludes.
\end{proof}

\begin{lemma}[A few factors of the gain pay for the height]\label{pfr:payheight}
Let $S$ be a shell of height $H$, $p,q$ heads of length $j$, the mass of floor $j$ at most $V\ge1$,
$e\le1$, $r'\ge0$ with $9\cdot10^7n^4+r'\le s$, and $f\ge0$ measurable with $f\le\pfrF_j^+$ everywhere
and with the gain $Ve^{s}$ on $S$. Then
\[
\pfrkap_j(p,q)\int_Sf\,d\pfrM\le H^{-1}\,V\,e^{r'}\,\Lambda^{3\cdot10^8n^4}\,\mathsf N(W).
\]
\end{lemma}
\begin{proof}
It suffices to find $\mathsf g\le e$ with
\begin{equation}\label{pfr:eq-pay}
\mathsf g^{9\cdot10^7n^4}H^{1+10^5n^2}\le\Lambda^{29\cdot10^7n^4},\qquad
\pfrkap_j\int_Sf\,d\pfrM\le V\mathsf g^{s}\Lambda^{2\cdot10^6n^3}H^{10^5n^2}\mathsf N(W):
\end{equation}
indeed $\mathsf g^s\le\mathsf g^{9\cdot10^7n^4}e^{r'}$ (as $\mathsf g\le e\le1$), and then the right side
of the second inequality is at most $H^{-1}Ve^{r'}\Lambda^{29\cdot10^7n^4+2\cdot10^6n^3}\mathsf N(W)$, and
$29\cdot10^7n^4+2\cdot10^6n^3\le3\cdot10^8n^4$.

If $H\le\Lambda^{2400n^2}$ take $\mathsf g=e$: the first inequality holds since $e\le1$ and
$H^{1+10^5n^2}\le\Lambda^{2400n^2(1+10^5n^2)}\le\Lambda^{29\cdot10^7n^4}$, and the second is
Lemma~\ref{pfr:oneshell}. If $H>\Lambda^{2400n^2}$, let $\tau^{800n^2}=H^{-1}$ and
$\mathsf g=\min(e,\tau)$. Then $\mathsf g^{9\cdot10^7n^4}H^{1+10^5n^2}\le(\tau^{800n^2})^{112500n^2}
H^{112500n^2}=1$. If $e\le\tau$ the second inequality is Lemma~\ref{pfr:oneshell}; if $\tau<e$ it is
Lemma~\ref{pfr:remoteshell}, since $f\le\pfrF^+_j$.
\end{proof}

\begin{theorem}[The sum over the shells]\label{pfr:shells}
Let $S_i$ $(i\ge0)$ be shells of heights $2^i$ that cover the support of a measurable $f\ge0$. Let $p,q$
be heads of length $j\le\varrho$ for $\Gamma$, let the mass of floor $j$ be at most $V\ge1$, let $f\le
\pfrF_j^+$ everywhere, and let $e\ge0$ and $r'$ with $9\cdot10^7n^4+r'\le s$ be such that, if $e<1$, $f$
has the gain $Ve^{s}$ at every $x\in\pfrP$. Then
\[
\pfrkap_j(p,q)\int_\pfrG f\,d\pfrM\le V\,\min(1,e)^{r'}\,\Lambda^{4\cdot10^8n^4}\,\mathsf N(W).
\]
\end{theorem}
\begin{proof}
Let $e'=\min(1,e)$. On each $S_i\subseteq\pfrP$, $f$ has the gain $Ve'^{\,s}$: if $e<1$ by hypothesis,
and if $e\ge1$ by $f\le\pfrF^+_j$ and Theorem~\ref{pfr:nogain} (as $V\ge1$). By
Lemma~\ref{pfr:payheight}, $\pfrkap_j\int_{S_i}f\le2^{-i}Ve'^{\,r'}\Lambda^{3\cdot10^8n^4}\mathsf N(W)$. Summing
over $i$ (the $S_i$ cover the support of $f$) gives the factor $\sum_i2^{-i}=2\le\Lambda^{10^8n^4}$.
\end{proof}

\subsection{The real place: data and the statement of \textup{(R)}}\label{pfr:ss-realdata}

\subsubsection*{The group}
Fix $\tau\in\auPl$. As in \S\ref{pbm:ss-remains} we work in the normalised basis of
\S\ref{pbe:ss-real}: $\pfrG=G_\tau$ is the unitary group of the hermitian form on $\C^r$ with $\varrho$
hyperbolic pairs $e_a,e_a^*$, an orthonormal family $(x_i)_{i\in\iota}$, $\#\iota=k$, and $z$ with
$h(z,z)=-1$; the frame is $(e_0,\dots,e_{\varrho-1},z)$, $\pfrU$ is the Bessel unipotent group and
$\pfrH=\auU(k)$ the fixer of the $e_a,e_a^*,z$. For $\gamma\in\pfrG$ let $Z(\gamma)=(h(\gamma f_b,f_a))_{a,b}$
be its \emph{frame matrix}, with diagonal entries $Z_{aa}(\gamma)=h(\gamma f_a,f_a)$. For a sequence
$x=(x_0,\dots,x_\varrho)$ of non-zero complex numbers with $|x_\varrho|>1$ let $w_x=w_{\varrho,x}$ and let
$\sigma_x\in\pfrG^{(\varrho)}$ be the boost in the plane of $z$ and $x_{i_0}$ ($i_0\in\iota$ fixed) with
$h(\sigma_xz,z)=x_\varrho$ (for $|x_\varrho|\le1$ we let $\sigma_x$ be the boost of parameter
$x_\varrho/\min(1,|x_\varrho|)$; this case does not occur in the sieve, as the last entries of the
open terms satisfy $|\tau(t)|>N$, that is $|x_\varrho|>1$ after normalisation).

\subsubsection*{The cut Bessel function and its orbital integrals}
Let $\mathbb W$ be a finite-dimensional unitary representation of $\pfrH$, of dimension $\dim\mathbb W=\#I$,
irreducible, and let $W,W^\flat\colon\pfrG\to\mathbb W$ be continuous Bessel functions of type $\mathbb W$: they transform on the left by the
character $u\mapsto e^{\sigma(u)-\overline{\sigma(u)}}$ under $\pfrU$ and by the representation under $\pfrH$, as in
Definition~\ref{pb:bessel}. Write $\pfrip{A}{R(x)B}=\int_{\pfrU\backslash\pfrG}\pfrip{A(y)}{B(yx)}\,dy$ for
the matrix coefficients in $L^2(\pfrU\backslash\pfrG)$ (the invariant measure normalised as in
\S\ref{pbe:ss-real}). For $\sigma\in\pfrG^{(\varrho)}$ and a sequence $x$ as above the \emph{normalised
orbital integral} is
\[
\pfrN_\sigma(x)=\frac{\sqrt{c_G}\,\bigl((|x_\varrho|^2-1)^+\bigr)^{(k-1)/2}}{\#I}\prod_{a<\varrho}|x_a|^{k+1+2(\varrho-a-1)}
\int_\pfrG W^\flat(y)\otimes\overline{W^\flat\bigl((w_x\sigma)^{-1}y\bigr)}\,dy\ \in\mathbb W\otimes\overline{\mathbb W},
\]
with $c_G=\pbvol\auU(k)^2/(\pbvol\auU(r)\pbvol\auU(k-1))$ as in \S\ref{pbm:ss-inputs}, and
$\pfrX(x)=\pfrN_{\sigma_x}(x)$. (The integrand has compact support, by Proposition~\ref{pbm:cell}
and the support of $W^\flat$, as in \S\ref{pbm:ss-cell}.)

\subsubsection*{Hypotheses}
Let $b\ge2$, $\rho=b^{6n+2}$, $f_W\in\N$ and $\pbdeg_\tau>0$. We assume:
\begin{enumerate}
\item[(H0)] (\emph{sizes}) $8r\le\rho$; $k=s_0+3n+6$ with $s_0>10^9n^3$; $n^7\le k$; $\varrho\ge10^8$;
\item[(H1)] (\emph{cut pair}) $\int_{\pfrU\backslash\pfrG}\|W\|^2=\#I$, $\int_{\pfrU\backslash\pfrG}\|W^\flat\|^2\le\#I$,
$\int_{\pfrU\backslash\pfrG}\|W^\flat-W\|^2\le b^{-2r^3f_W}\#I$, and $W^\flat(y)\ne0$ only if
$\|uy\|\le\rho$ for some $u\in\pfrU$ (operator norm);
\item[(H2)] (\emph{formal degree}) $\int_\pfrG|\pfrip{W}{R(x)W}|^2\,dx\le(\#I)^2/\pbdeg_\tau$;
\item[(H3)] $\pbdeg_\tau\,\pbmTh\le b^{4r^2}$, with $\pbmTh=\pi^{(2\varrho+1)k}\pbvol\auU(k)/\pbvol\auU(r)$;
\item[(H4)] (\emph{balls}) $\pbdeg_\tau\cdot\pbvol\{g\in\pfrG:\|g\|\le T\}\le T^{9r^2}b^{99r^2}$ for $T\ge1$;
\item[(H5)] $f_W\ge4000n$.
\end{enumerate}
The Haar measure $dx$ of $\pfrG$ in (H2), (H4) and in $\pfrN_\sigma$ is the one fixed in
\S\ref{pbe:ss-real}; (H2) says that $\pbdeg_\tau$ is at most the formal degree of the discrete series
generated by $W$ for that measure.

\begin{theorem}[Theorem~\ref{pbm:inputs}\,\textup{(R)}]\label{pfr:R}
Assume \textup{(H0)}--\textup{(H5)}. There are functions $J_i=J^\tau_i$, $0\le i\le n$, on pairs of
sequences of non-zero complex numbers such that:
\begin{enumerate}
\item[(R-a)] $(J_i)_{0\le i\le n}$ is a sieve system (Definition~\ref{pbm:system}) at the place
$(\C,\text{Lebesgue},c\mapsto e^{c-\bar c},w\mapsto w^{-1})$ with mass $(1+b^{-10})/\pbdeg_\tau$,
weights $\Lambda^\infty_i$ and decay $\beta^\infty$ of \eqref{pbm:eq-beta}:
\[
\Lambda^\infty_i(p)=\prod_{l=1}^i\bigl[|\delta_l|\le b^{60n^2}\bigr]\max(1,|\delta_l|^{-1})^{n^6},\quad
\delta_l=p_0\cdots p_{l-1};\qquad\beta^\infty(\xi)=\pbmTh\,b^{2000n^9}\,\pfrell_{k-n^7}(|\xi|^{-1}),
\]
$\pfrell_K(t)=t^K\min(1,t)^6$;
\item[(R-b)] $J_n(x,x')=\pfrip{\pfrX(x)}{\pfrX(x')}$; and for every $\sigma\in\pfrG^{(\varrho)}$ and all $x,x'$
with $x_\varrho=x'_\varrho=h(\sigma z,z)$, $J_n(x,x')=\pfrip{\pfrN_\sigma(x)}{\pfrN_\sigma(x')}$.
\end{enumerate}
\end{theorem}

\begin{remark}[Comparison with Theorem~\ref{pbm:inputs}]\label{pfr:rem-match}
(a) \emph{The constants agree exactly.} The mass $(1+b^{-10})/\pbdeg_\tau$, the weights
$\Lambda^\infty$ (radius $b^{60n^2}$, exponent $n^6$, product over the partial products
$\delta_l=p_0\cdots p_{l-1}$, $1\le l\le i$) and the decay $\beta^\infty$ (constant $\pbmTh\,b^{2000n^9}$,
the same $\pbmTh$ as in \S\ref{pbm:ss-inputs}, exponent $k-n^7$, the factor $\min(1,t)^6$) are those of
\eqref{pbm:eq-beta} and Theorem~\ref{pbm:inputs}\,(R), which is all that Theorem~\ref{pbm:sieve},
Lemma~\ref{pbm:const} and Theorem~\ref{pbm:L2} use from \textup{(R)}. The constant $c_G$ in
$\pfrN_\sigma$ is that of \S\ref{pbm:ss-inputs}: $c_G=\frac{\pbvol\auU(k)}{\pbvol\auU(r)}\cdot
\frac{\pbvol\auU(k)}{\pbvol\auU(k-1)}=\frac{\pbvol\auU(k)}{\pbvol\auU(r)}\cdot\frac{2\pi^k}{(k-1)!}$. The
normalisation of $\pfrN_\sigma$ is that of \S\ref{pbm:ss-remains}. The indexing of
Definition~\ref{pbm:system} matches: (S1) is required for $i\le\varrho$ (integration over the entry of
index $i$ of the forms $J_{i+1}$), (S5) for $1\le i\le\varrho$ with $\xi=p_{i-1}^{-1}-q_{i-1}^{-1}$.

(b) \emph{What must be added to Theorem~\ref{pbm:inputs}\,(R).} First, the identification (R-b) of the
top form with the Gram form of the normalised orbital integrals: Siegel's step \textup{(S)} is about
the forms $J_n^\tau$, and it is through (R-b) that they are related to the orbit sums of
Proposition~\ref{pbm:orbit}. Second, the hypotheses (H0)--(H5) on the data and the Bessel functions;
they are properties of the Bessel functions of Theorem~\ref{pbe:bessel} and of the data, and have to be
checked there. Among them, (H0) asks $\varrho\ge10^8$ and $k\ge n^7$, $k>10^9n^3+3n+6$, which is
stronger than the condition $n\ge100$ listed after Theorem~\ref{pbm:L2}, but holds for the data: by
Lemma~\ref{an:lem-sizes} (S1), (S2), (S7), $\varrho=m/2-1\ge\sqrt x/2-1$ with $x\ge128^{100}$, $k\ge x^5/2$ and
$n\le\sqrt x$, so $\varrho\ge10^8$, $k\ge x^5/2\ge n^{10}/2\ge n^7$ and $k-3n-6\ge n^{10}/2-3n-6>10^9n^3$ (these three conditions are also stated directly in (S7)(a)). The remaining condition of (H0), $8r\le\rho$, holds since $8r\le8nr\le b\le b^{6n+2}=\rho$ by Lemma~\ref{an:lem-sizes}\,(S7)(b).
(c) The functions $J_i$, $i\le\varrho$, are the normalised forms of floor $i$:
$J_i(p,q)=\mathsf c\cdot\pfrI_i(p,q;1)$ for the place datum of \S\ref{pfr:ss-realforms} and the constant
$\mathsf c$ of \eqref{pfr:eq-c}; this is the description of \S\ref{pbm:ss-remains} (``$J_i$ is $J_n$
integrated over the last $n-i$ entries'').
\end{remark}

\subsection{The real place: the place datum, the forms and the local estimates}\label{pfr:ss-realforms}

\subsubsection*{The place datum at $\C$}
The \emph{place datum at $\tau$} is: level $M=1$; close pairs the $(\omega_1,\omega_2)$ with
$|\omega_1|^2+|\omega_2|^2<1$ (Euclidean norm of the coordinates); box $\pfrK=\{g:\|g\|\le\rho\}$; all
of $\pfrG$ as positions; $\mu_\pfrH$ the Haar probability measure of $\pfrH=\auU(k)$; $\lambda_l$ ($l<\varrho$)
the image of the measure $|c|^{2d_l}\,dc$ on $\C^\times$ (the \emph{cell density} of floor $l$;
$2d_l=2(k-1)+4(\varrho-l)$) under $c\mapsto w_{l,c}^{-1}$, as in \S\ref{pfr:ss-setting} ($\pfrmod(c)^{d_l}=|c|^{2d_l}$);
and $\theta_\varrho=\lambda_\varrho*\mu_\pfrH$, where $\lambda_\varrho$ is the image of the measure
$\vartheta(c)\,dc$, $\vartheta(c)=\frac{2\pi^k}{(k-1)!}\bigl((|c|^2-1)^+\bigr)^{k-1}$, under $c\mapsto\sigma_c^{-1}$ (the
boosts of \S\ref{pfr:ss-realdata}); that is, $\int F\,d\theta_\varrho=\iint F(\sigma_c^{-1}k)\,\vartheta(c)\,dc\,d\mu_\pfrH(k)$.
The constant $\frac{2\pi^k}{(k-1)!}=\frac{\pbvol\auU(k)}{\pbvol\auU(k-1)}$ is the one in $c_G$
(Remark~\ref{pfr:rem-match}(a)), and it is forced by the identity $N(p^c)N(q^c)=\mathsf c\lambda^2\vartheta(c)\pfrkap_\varrho(p,q)$
of the proof of Proposition~\ref{pfr:top}. The mass condition holds since $\mu_\pfrH$ is a probability measure.

\emph{The constants $c_l$ at $\C$.} For $l<\varrho$, $c_l=\kappa_\C=2$ by Lemma~\ref{pfr:floorquad}. For $l=\varrho$,
$c_\varrho=2$: write $y=y_zz+y_X\in V_\varrho$ with $y_X$ in the span $X_\iota\cong\C^k$ of the $x_i$, so $h(y,y)=|y_X|^2-|y_z|^2$.
For fixed $y_z$, polar coordinates give $dy_X=\mathsf m_k(R)\,dR\,dS$, with $R=|y_X|^2$,
$\mathsf m_k(R)=\pi^kR^{k-1}/(k-1)!$ (the derivative of the volume $\pi^kR^k/k!$ of the ball) and $dS$ the
$\auU(k)$-invariant probability measure on the sphere of radius $\sqrt R$. As $R\mapsto\mathsf m_k(R)\int\varphi\,dS$ is
continuous, \eqref{pfr:eq-slab} gives $\int\varphi\,d\mu_\varrho=\int_{|y_z|>1}\mathsf m_k(|y_z|^2-1)\int\varphi(y_zz+y_X)\,
dS(y_X)\,dy_z$, the inner integral over the sphere of radius $(|y_z|^2-1)^{1/2}$ (dominated convergence; $\varphi$ has
compact support). On the other side, for $\sigma=\sigma_c^{-1}k$, $\sigma^{-1}z=k^{-1}\sigma_cz=-cz+k^{-1}v_c$ with $|v_c|^2=
|c|^2-1$ (proof of Proposition~\ref{pfr:top}), and $k^{-1}v_c$ is distributed by $dS$ when $k$ is distributed by
$\mu_\pfrH$. So the image of $\theta_\varrho$ is $\int\vartheta(c)\int\varphi(-cz+y_X)\,dS\,dc=2\int\varphi\,d\mu_\varrho$
(substitute $y_z=-c$, and $\vartheta(c)=2\mathsf m_k(|c|^2-1)$ for $|c|>1$). Thus $c_l=2$ for every $l\le\varrho$ at the real place. The gauge of scale
$L\ge\rho$, $\Gamma_L$, has $\Gamma_L(i)=\{c:|c|\le L^i\}$ and loss $\Lambda=L^2$. All of this is the
specialisation of \S\ref{pfr:ss-setting}--\S\ref{pfr:ss-gauge}; that the gauge axioms (G0)--(G5)
hold for $\Gamma_L$ under (H0) is Lemma~\ref{pfr:realgauge}.

\begin{lemma}[Haar measure]\label{pfr:haar}
The group measure $\pfrM$ of the place datum at $\tau$ is $\lambda$ times the Haar measure $dx$ of
$\pfrG$ for a number $0<\lambda<\infty$ with $\lambda\,\pbvol\auU(k)/\pbvol\auU(r)\le1$; and for every
measurable $\phi\ge0$ on $\pfrG$ invariant under left translation by $\pfrU$,
$\int\phi\,d\theta_{\ge0}=\lambda\int_{\pfrU\backslash\pfrG}\phi$.
\end{lemma}
(Proof in \S\ref{prp:ss-haar}.)

\begin{lemma}[The gauge at $\C$]\label{pfr:realgauge}
Under \textup{(H0)}, for every $L\ge\rho$ the filtration $\Gamma_L$ with loss $L^2$ is a gauge for the
place datum at $\tau$, and its volume exponent is $s=k-(3n+1)-2\cdot2-1=s_0$ ($x_{i_0}$ and two blocks of rank
$2m+2=2$ are set aside, $m=0$); by (H0), $s\ge10^{10}n^4$. The depth of $e_a$ at
$(g,gx)$ is $\pi^{-1}\inf_\zeta(|g^{-1}\zeta|^2+|(gx)^{-1}\zeta|^2)$ over the duals $\zeta$ of $e_a$ in
$Y^{(a)}$.
\end{lemma}
(Proof in \S\ref{prp:ss-gauge}.)

\subsubsection*{The forms}
Put
\begin{equation}\label{pfr:eq-c}
\mathsf c=\frac{\pbvol\auU(k)}{\pbvol\auU(r)}\cdot\frac1{(\lambda\,\#I)^2},\qquad
J_i(p,q)=\mathsf c\,\pfrI_i(p,q;1)\ (0\le i\le\varrho),\qquad J_n(p,q)=\pfrip{\pfrX(p)}{\pfrX(q)},
\end{equation}
where $\pfrI_i$ is the form of floor $i$ of the place datum at $\tau$ for the cut Bessel function
$W^\flat$ (that $W^\flat$ is a cut Bessel function for this datum is (H1) with $M=1$).

\begin{proposition}[Absolute convergence and support]\label{pfr:conv}
Under \textup{(H0)}--\textup{(H1)}, for $j\le\varrho$ and all $p,q$, $\int\pfrF^+_j\,d\pfrM<\infty$.
If $|p_0\cdots p_{l-1}|>\rho^{9n}$ or $|q_0\cdots q_{l-1}|>\rho^{9n}$ for some $l\le j$, then
$\pfrF^+_j=0$, so $\pfrI_j(p,q;\omega)=0$ for $|\omega|\le1$.
\end{proposition}

\begin{proposition}[The top floor]\label{pfr:top}
Under \textup{(H0)}--\textup{(H1)}: $\pfrX$ is continuous; for $\sigma\in\pfrG^{(\varrho)}$ and
$x_\varrho=x'_\varrho=h(\sigma z,z)$, $\pfrip{\pfrX(x)}{\pfrX(x')}=\pfrip{\pfrN_\sigma(x)}{\pfrN_\sigma(x')}$; and
for $\xi\in\C$ the function $c\mapsto e^{\xi c-\overline{\xi c}}\pfrip{\pfrX(p^c)}{\pfrX(q^c)}$ ($p^c,q^c$:
the entry of index $\varrho$ replaced by $c$) is integrable, with integral
$\mathsf c\,\pfrI_\varrho\bigl(p,q;\ g\mapsto e^{\xi Z_{zz}(g)-\overline{\xi Z_{zz}(g)}}\bigr)$.
\end{proposition}
(Proposition~\ref{pfr:conv} is proved in \S\ref{prp:ss-sum}, Proposition~\ref{pfr:top} in \S\ref{prp:ss-top}.)

\begin{proposition}[The bound for a form]\label{pfr:formbound}
Under \textup{(H0)}--\textup{(H1)}, let $i\le\varrho$, $p,q$ sequences, $\nu\in\C^\times$, and let
$\omega$ be the weight $g\mapsto e^{\nu Z_{ii}(g)-\overline{\nu Z_{ii}(g)}}$, or, if $i\ge1$ and
$\nu=p_{i-1}^{-1}-q_{i-1}^{-1}$, the weight $\omega=1$. Then
\[
\#I\cdot\mathsf c\,|\pfrI_i(p,q;\omega)|\le\pbmTh\,b^{2000n^9}\,\pfrell_{k-n^7}(|\nu|^{-1})\,
\Lambda^\infty_i(p)\Lambda^\infty_i(q).
\]
\end{proposition}
(Proof in \S\ref{prp:ss-sum}.)

\begin{proposition}[Remote positions]\label{pfr:remote}
Under \textup{(H0)}--\textup{(H1)}, let $\mathcal B$ be the set of the $x\in\pfrG$ all of whose entries
of $x-1$ have modulus at most $(b^{6r}\rho^{11n})^{99n}$ (the \emph{position ball}). Then
\[
\mathsf c\,\pfrkap_0\int_{\pfrG\smallsetminus\mathcal B}\pfrF^+_0\,d\pfrM\le\frac{\pbmTh\,\rho^{10^9n^4}\rho^{24k}
(b^{6r})^{10^8n^4}}{(b^{6r})^{s_0}}.
\]
\end{proposition}
(Proof in \S\ref{prp:ss-sum}.)

\begin{lemma}[Far means deep]\label{pfr:deepreal}
Let $p,q$ be heads of length $j\le\varrho$ for $\Gamma_L$, $L\ge\rho$, let $T\ge L^{11n}$ and
$T^{-1}\le\tau'/2$ with $0<\tau'\le1$. If some entry of $x-1$ has modulus $>T^{99n}$, then $x$ has a deep
vector of depth $\tau'^{-2}$.
\end{lemma}
(Proof in \S\ref{prp:ss-deep}.)

\subsection{The mass}\label{pfr:ss-mass}

\begin{proposition}[The mass at a real place]\label{pfr:mass}
Assume \textup{(H0)}--\textup{(H5)}. Then $|J_0(p,q)|\le(1+b^{-10})/\pbdeg_\tau$ for all $p,q$.
\end{proposition}

We first collect three lemmas. Write $\mathsf C(x)=\pfrip{W^\flat}{R(x)W^\flat}$ and
$\mathsf C_0(x)=\pfrip{W}{R(x)W}$ for the matrix coefficients.

\begin{lemma}\label{pfr:F0}
For every $x$ and all $p,q$, $\pfrF_0(x;1)=\lambda^2|\mathsf C(x)|^2$, where $\lambda$ is the number of
Lemma~\ref{pfr:haar}.
\end{lemma}
\begin{proof}
At floor $0$ the Weyl elements and the Weyl ratio are $1$, $\pfrm_0=\delta_1$, the weight is $1$, and the
two factors of $\Phi_0(x;1;1)$ are conjugate: $\Phi_0(x;1;1)=|\mathsf k^1_0(1,x)|^2$, where $\mathsf
k^1_0(1,x)=\int\pfrip{W^\flat(y)}{W^\flat(yx)}\,d\theta_{\ge0}(y)$ (the integration over $\theta_0$ in $\Phi_0$ and
the correlation from floor $1$ on combine to the correlation from floor $0$ on). The function
$y\mapsto\pfrip{W^\flat(y)}{W^\flat(yx)}$ is invariant under left translation by $\pfrU$, because the two
factors $\psi(\sigma(u))$ cancel. By Lemma~\ref{pfr:haar} (applied to the positive and negative parts of
the real and imaginary parts, each integrable by (H1) and Cauchy--Schwarz), $\mathsf k^1_0(1,x)=\lambda\,
\mathsf C(x)$.
\end{proof}

\begin{lemma}\label{pfr:coeffdiff}
For every $x$, $|\mathsf C(x)-\mathsf C_0(x)|\le2b^{-r^3f_W}\,\#I$.
\end{lemma}
\begin{proof}
$\mathsf C(x)-\mathsf C_0(x)=\pfrip{W^\flat-W}{R(x)W^\flat}+\pfrip{W}{R(x)(W^\flat-W)}$. The difference of two
Bessel functions of the same type is a Bessel function, right translation preserves the norm of
$L^2(\pfrU\backslash\pfrG)$, and by Cauchy--Schwarz and (H1) each term has modulus at most
$(b^{-2r^3f_W}\#I)^{1/2}(\#I)^{1/2}$.
\end{proof}

\begin{lemma}\label{pfr:ballnorm}
Let $\mathcal B$ be the position ball of Proposition~\ref{pfr:remote}. Under \textup{(H0)},
$\|x\|\le b^{700nr}$ for $x\in\mathcal B$, and
$\int_{\mathcal B}|\mathsf C(x)|^2dx\le(1+b^{-12})^2(\#I)^2/\pbdeg_\tau$ under \textup{(H1)}, \textup{(H2)},
\textup{(H4)}, \textup{(H5)}.
\end{lemma}
\begin{proof}
Let $R_0=(b^{6r}\rho^{11n})^{99n}$. For $x\in\mathcal B$ the entries of $x-1$ are at most $R_0$, so
$\|x\|\le1+r^2R_0\le\rho^2R_0$ (as $r^2+1\le\rho^2$ by $8r\le\rho$). With $\rho=b^{6n+2}\le b^{7n}$ (as $n\ge2$),
$\rho^2R_0\le b^{14n+99n(6r+77n^2)}\le b^{700nr}$, because $14n+7623n^3\le106nr$ when $r\ge10^9n^3$ (which
follows from (H0): $r=2\varrho+k+1\ge k>10^9n^3$).

For the second claim, let $\alpha=(\#I)\pbdeg_\tau^{-1/2}$, so that $\int_{\mathcal B}|\mathsf C_0|^2\le\alpha^2$ by
(H2). Let $T=b^{700nr}$ and $V_0=T^{9r^2}b^{99r^2}$; by the first claim and (H4),
$\pbvol(\mathcal B)\le V_0/\pbdeg_\tau$. By Lemma~\ref{pfr:coeffdiff},
\[
\int_{\mathcal B}(2b^{-r^3f_W}\#I)^2\,dx\le(2b^{-r^3f_W})^2V_0\cdot\alpha^2\le(b^{-12}\alpha)^2,
\]
since $(2b^{-r^3f_W})^2V_0\le b^{2+6300nr^3+99r^2-2r^3f_W}\le b^{-24}$: indeed $2+6300nr^3+99r^2+24\le
2r^3\cdot4000n\le2r^3f_W$ by (H5) and $r\ge1$, $n\ge1$. By Minkowski's inequality in $L^2(\mathcal B)$
applied to $|\mathsf C|\le|\mathsf C_0|+2b^{-r^3f_W}\#I$,
$\bigl(\int_{\mathcal B}|\mathsf C|^2\bigr)^{1/2}\le\alpha+b^{-12}\alpha$.
\end{proof}

\begin{proof}[Proof of Proposition~\ref{pfr:mass}]
By Proposition~\ref{pfr:conv}, Lemma~\ref{pfr:L1} and $\pfrkap_0=1$,
$|J_0(p,q)|\le\mathsf c\int\pfrF_0(x;1)\,d\pfrM(x)$. Split the integral into $\mathcal B$ and its
complement. On $\mathcal B$, by Lemmas~\ref{pfr:haar}, \ref{pfr:F0} and~\ref{pfr:ballnorm},
\[
\mathsf c\int_{\mathcal B}\pfrF_0\,d\pfrM=\mathsf c\,\lambda\cdot\lambda^2\int_{\mathcal B}|\mathsf C|^2dx\le
\frac{\pbvol\auU(k)}{\pbvol\auU(r)}\frac{\lambda^3}{\lambda^2(\#I)^2}\cdot\frac{(1+b^{-12})^2(\#I)^2}{\pbdeg_\tau}
\le\frac{(1+b^{-12})^2}{\pbdeg_\tau},
\]
using $\lambda\,\pbvol\auU(k)/\pbvol\auU(r)\le1$. On the complement, $\pfrF_0\le\pfrF_0^+$
(Lemma~\ref{pfr:L1}), and by Proposition~\ref{pfr:remote} and (H3) the contribution is at most
\[
\pbmTh\,\rho^{10^9n^4}\rho^{24k}(b^{6r})^{10^8n^4-s_0}\le\frac{b^{4r^2}}{\pbdeg_\tau}\,b^{7n(10^9n^4+24k)+6r\cdot
10^8n^4-6rs_0}\le\frac{b^{-12}}{\pbdeg_\tau}.
\]
For the last step write $r=s_0+5n+5$ and $k=s_0+3n+6$ (by (H0) and $r=2\varrho+k+1$), and note
$s_0\ge10^{10}n^4$: by (H0), $s_0=k-3n-6\ge n^7-3n-6$ and $n\ge10^8$. We need $E+12\le6rs_0$ with
$E=4r^2+7\cdot10^9n^5+168nk+6\cdot10^8n^4r$. Now $6rs_0-4r^2=2r(s_0-10n-10)\ge1.9\,rs_0$, while
$6\cdot10^8n^4r\le0.06\,rs_0$, $7\cdot10^9n^5\le0.7\,n\cdot10^{10}n^4\le0.7\,rs_0$ (as $r\ge s_0\ge n$),
and $168nk+12\le169\,nr\le0.01\,rs_0$. Altogether
\[
|J_0(p,q)|\le\frac{(1+b^{-12})^2+b^{-12}}{\pbdeg_\tau}\le\frac{1+b^{-10}}{\pbdeg_\tau},
\]
since $(1+u^{12})^2+u^{12}=1+(3+u^{12})u^2\cdot u^{10}\le1+u^{10}$ for $0<u\le\frac12$.
\end{proof}

\subsection{Proof of Theorem~\ref{pfr:R}}\label{pfr:ss-Rproof}

\begin{proof}[Proof of Theorem~\ref{pfr:R} from Propositions~\ref{pfr:conv}--\ref{pfr:formbound}, Lemma~\ref{pfr:descent} and Proposition~\ref{pfr:mass}]
Let $J_i$ be as in \eqref{pfr:eq-c}. Heads are sequences of non-zero complex numbers; in (S1) the entry
$c$ ranges over $\C$, and $c=0$ (where the head is not defined) is a null set.

(S1) Let $i\le\varrho$, $p,q$ heads, $\xi\in\C$. If $i<\varrho$, by Proposition~\ref{pfr:conv} and
Lemma~\ref{pfr:descent} (with $\psi(c)=e^{c-\bar c}$), $c\mapsto e^{\xi c-\overline{\xi c}}J_{i+1}(p^c,q^c)$ is
integrable and its integral is $\mathsf c\,\pfrI_i(p,q;\omega_\xi)$ with $\omega_\xi(g)=e^{\xi Z_{ii}(g)-
\overline{\xi Z_{ii}(g)}}$. If $i=\varrho$ the same holds by Proposition~\ref{pfr:top}. For $\xi=0$,
$\omega_0=1$ and the integral is $J_i(p,q)$. For $\xi\ne0$, by Proposition~\ref{pfr:formbound} and
$\#I\ge1$,
$|\mathsf c\,\pfrI_i(p,q;\omega_\xi)|\le\pbmTh\,b^{2000n^9}\pfrell_{k-n^7}(|\xi|^{-1})\Lambda^\infty_i(p)\Lambda^\infty_i(q)
=\beta^\infty(\xi)\Lambda^\infty_i(p)\Lambda^\infty_i(q)$.

(S2) For heads $p^{(1)},\dots,p^{(\kappa)}$ (with any common last entry) the matrix
$(\pfrip{\pfrX(p^{(a)})}{\pfrX(p^{(b)})})_{a,b}$ is a Gram matrix, hence hermitian and positive
semidefinite.

(S3) $J_n$ is continuous because $\pfrX$ is (Proposition~\ref{pfr:top}).

(S4) is Proposition~\ref{pfr:mass}.

(S5) Let $1\le i\le\varrho$ and $\xi=p_{i-1}^{-1}-q_{i-1}^{-1}\ne0$. By Proposition~\ref{pfr:formbound}
with $\omega=1$, $|J_i(p,q)|\le\beta^\infty(\xi)\Lambda^\infty_i(p)\Lambda^\infty_i(q)$.

(R-b) is the definition of $J_n$ and the second claim of Proposition~\ref{pfr:top}.
\end{proof}

\subsection{Proofs of Theorems~\ref{pfr:nogain} and~\ref{pfr:deep}}\label{pfr:ss-gainsproof}

We keep the assumptions of \S\ref{pfr:ss-shared}. Two facts on the floors are used without proof in
this subsection; they are stated exactly.

\begin{lemma}[Weyl invariance of the lower floors]\label{pfr:weylinv}
Let $j\le\varrho$, $p$ a sequence of units and $\Phi\ge0$ measurable and invariant under left
translation by $\pfrG^{(j)}$. Then $\int\Phi(w_{j,p}^{-1}g)\,d\pfrm_j(g)=\int\Phi(g)\,d\pfrm_j(g)$.
\end{lemma}

\begin{lemma}[The inverse weight on the floors below a head]\label{pfr:invweight}
Let $p$ be a head of length $j\le\varrho$ for $\Gamma$, $x\in\pfrP$, $j\le b'$, and let $\Phi\ge0$ be
measurable, invariant under left translation by $\pfrR^{(b')}$ and zero at the points not near at $b'$.
Then
\[
\int\Phi(g)\,\pfrwt^{-}_j(w_{j,p}g)\,d\pfrm_j(g)\le\Bigl(\prod_{a<j}\mathsf c_a(p_a^{-1})\Bigr)C_\Gamma
\int\Phi\,d\theta_{[0,j)},
\]
where $\pfrwt^-_j(g)=\pfrwt_j(g,x)^{-1}$ if $g$ is near at $j$, and $0$ otherwise.
\end{lemma}
(Proofs in \S\ref{pfr:ss-floors}.)

\subsubsection*{Elementary facts on correlations}
Let $\mathsf k=\mathsf k^{d'}$ for a cut Bessel function $W$ and an element $d'$ commuting with
$\pfrG^{(j)}$ (for us $d'=1$ or the Weyl ratio $w_q^{-1}w_p$, which is diagonal on the hyperbolic
pairs of index $<j$ and the identity elsewhere). We use:
\begin{enumerate}
\item[(C1)] $\mathsf k_l(g,y)=\int\mathsf k_{l+1}(\sigma g,\sigma y)\,d\theta_l(\sigma)$, by the
definition $\theta_{\ge l}=\theta_{\ge l+1}*\theta_l$ and Fubini's theorem (the integrals converge
absolutely: $W$ is bounded and, by the cut condition and the compactness of $\pfrK$, the integrand is
supported in a set of finite measure); hence $|\mathsf k_l(g,y)|\le\pfrB_l(g,y)$, the pivot norm
$\int|\mathsf k_{l+1}(\sigma g,\sigma y)|d\theta_l(\sigma)$.
\item[(C2)] If $\mathsf k_l(g,y)\ne0$ then $g$ is near at $l$: some $x'g$, $x'\in\pfrG^{(l)}$, has
$W(x'g)\ne0$, hence $ux'g\in\pfrK$ for some $u\in\pfrU$; and the element $ux'$ lies in $\pfrR^{(l)}$
after the factor of $u$ that moves the $f_a$, $a\ge l$, is absorbed (precisely: $\pfrR^{(l)}\supseteq
\pfrU\pfrG^{(l)}$).
\item[(C3)] The depths $\pfrdep_a(g,x)$, $a<l$, and therefore $\pfrwt_l(g,x)$, are invariant under
$g\mapsto\xi g$ for $\xi\in\pfrR^{(l)}$, in particular for $\xi\in\pfrG^{(l)}$: the stable vectors for
$a$ at $(\xi g,\xi gx)$ are the $\xi\omega$ with $\omega$ stable at $(g,gx)$, and $\xi$ preserves
$\langle e_b:b\le a\rangle$ and the coordinate on $e^*_a$ modulo that span.
\item[(C4)] $|\mathsf k_{l+1}(ng,n'y)|=|\mathsf k_{l+1}(g,y)|$ for $n,n'\in\pfrU_{(<l)}$ (the equivariance of
the correlation by a character, (K1) of Lemma~\ref{pfs:kernels}); the function $\mathsf a_l$ is invariant under
$\pfrR^{(l+1)}$ ((K3)) and $\mathsf b_l$ under $\pfrR^{(l)}$ (Lemma~\ref{pfr:floorright}, descending induction) on the
left, and both vanish off the points near at $l+1$, resp.\ $l$, by (C2).
\end{enumerate}

\begin{lemma}[A weighted square on one floor]\label{pfr:floorsq}
For $l\le\varrho$ and all $g$:
\[
\int|\mathsf k_{l+1}(\sigma g,\sigma gx)|^2\,d\theta_l(\sigma)\le\mathsf b_l(g)\,\pfrwt_l(g,x)^{-1}.
\]
If moreover the mass of floor $l$ is at most $V$, then $\pfrB_l(g,gx)^2\le V\int|\mathsf k_{l+1}(\sigma g,\sigma
gx)|^2\,d\theta_l(\sigma)$.
\end{lemma}
\begin{proof}
$\theta_l$ is carried by $\pfrG^{(l)}$, on which $\pfrwt_l(\sigma g,x)=\pfrwt_l(g,x)$ by (C3); so the left
side is $\pfrwt_l(g,x)^{-1}\int\mathsf a_l(\sigma g)\,d\theta_l(\sigma)\le\pfrwt_l(g,x)^{-1}\mathsf b_l(g)$ (the
weight is finite and non-zero, as each depth is). For the second claim, by (C2) the integrand of
$\pfrB_l(g,gx)$ vanishes unless $\sigma g$ is near at $l+1$, and by Cauchy--Schwarz on that set, of
$\theta_l$-measure at most $V$, $\pfrB_l^2\le V\int|\mathsf k_{l+1}|^2$.
\end{proof}

\begin{lemma}[Separating the frames]\label{pfr:separate}
Let $p$ be a head of length $j\le\varrho$, $x\in\pfrP$, $K\ge0$ and $f\ge0$ measurable on $\pfrG$ with
\[
f(g)^2\le K\,C_\Gamma\,\frac{\mathsf b^{(1)}_j(g)}{\pfrwt_j(g,x)}\cdot\frac{\mathsf b^{(d)}_j(w_p^{-1}g)}{\pfrwt_j(w_p^{-1}g,x)}
\quad\text{wherever }f(g)\ne0,
\]
where $\mathsf b^{(1)}$, $\mathsf b^{(d)}$ are the functions $\mathsf b$ of the kernels $\mathsf k^1$,
$\mathsf k^{d}$. Then $\bigl(\int f\,d\pfrm_j\bigr)^2\le K\,\mathsf Q(x)$.
\end{lemma}
\begin{proof}
Put $B=\mathsf b^{(1)}_j$, $B'=\mathsf b^{(d)}_j$; by (C4) both are $\pfrR^{(j)}$-invariant and vanish
off the points near at $j$, so we may replace $B(g)/\pfrwt_j(g,x)$ by $B(g)\pfrwt^-_j(g)$, and likewise
for $B'$. Crossing the weights, $f(g)^2\le\bigl(B'(w_p^{-1}g)\pfrwt^-_j(g)\bigr)\cdot KC_\Gamma\bigl(B(g)
\pfrwt^-_j(w_p^{-1}g)\bigr)$, and by Cauchy--Schwarz
\[
\Bigl(\int f\,d\pfrm_j\Bigr)^2\le\int B'(w_p^{-1}g)\pfrwt^-_j(g)\,d\pfrm_j\cdot KC_\Gamma\int B(g)\pfrwt^-_j(w_p^{-1}g)\,
d\pfrm_j .
\]
In the first integral substitute $g\mapsto w_pg$ by Lemma~\ref{pfr:weylinv}: the function
$g\mapsto B'(g)\pfrwt^-_j(w_pg)$ is $\pfrG^{(j)}$-invariant (as $w_p$ commutes with $\pfrG^{(j)}$ and by
(C3), (C4)); then Lemma~\ref{pfr:invweight} bounds it by $D\int B'\,d\theta_{[0,j)}$ with
$D=\prod_{a<j}\mathsf c_a(p_a^{-1})C_\Gamma$. In the second, $w_p^{-1}=w_{\bar p}$ (indeed $w_p^{-1}e_a=\bar p_ae_a^*$ and $w_p^{-1}e_a^*=p_a^{-1}e_a$),
$\bar p$ is again a head for $\Gamma$ (the filtration is stable under the involution) and
$\mathsf c_a(\bar p_a^{-1})=\mathsf c_a(p_a^{-1})$; so Lemma~\ref{pfr:invweight} for the head $\bar p$
gives the bound $D\int B\,d\theta_{[0,j)}$. Hence
$(\int f)^2\le KC_\Gamma D^2\,\mathsf S_j(\mathsf k^1;x)\,\mathsf S_j(\mathsf k^d;x)$. Finally
$\mathsf S_j(\mathsf k^1;x)\le\mathsf S_0(\mathsf k^1;x)$ (unfolding $\mathsf b_0$ down to floor $j$ gives
$\mathsf S_0=\mathsf S_j+\sum_{l<j}\int\mathsf a_l\,d\theta_{[0,l+1)}\ge\mathsf S_j$), and $KC_\Gamma D^2\mathsf S_0
\mathsf S_j=K\mathsf Q(x)$ by definition.
\end{proof}

\begin{proof}[Proof of Theorem~\ref{pfr:nogain}]
Let $y$ be a point and $z=w_p^{-1}y$. By Lemma~\ref{pfr:floorsq} for both kernels,
\[
\bigl(\pfrB^1_j(y,yx)\pfrB^{d}_j(z,zx)\bigr)^2\le V^2\frac{\mathsf b^{(1)}_j(y)}{\pfrwt_j(y,x)}\cdot\frac{\mathsf
b^{(d)}_j(z)}{\pfrwt_j(z,x)},
\]
and $C_\Gamma\ge1$. Lemma~\ref{pfr:separate} with $K=V^2$ and $f(y)=\pfrB^1_j(y,yx)\pfrB^d_j(z,zx)$ gives the
claim, since $\pfrF^+_j(x)=\int f\,d\pfrm_j$.
\end{proof}

\subsubsection*{The cut at the first deep vector}
For a kernel $\mathsf k$, a position $x$, $D>0$ and $a<\varrho$ let $\mathsf d_a(y)=|\mathsf k_{a+1}(y,
yx)|$ if $\pfrdep_a(y,x)\ge D$, and $0$ otherwise.

\begin{lemma}\label{pfr:cut}
Let $l\le\varrho$ and $g$ be such that for every $\sigma\in\pfrG^{(l)}$ with $\mathsf k(\sigma g,\sigma gx)\ne0$
there is $a$ with $l\le a<\varrho$ and $\pfrdep_a(\sigma g,x)\ge D$. Then
\[
\pfrB_l(g,gx)\le\sum_{l\le a<\varrho}\int\mathsf d_a(\sigma g)\,d\theta_{[l,a+1)}(\sigma).
\]
\end{lemma}
\begin{proof}
Descending induction on $l$. If $l=\varrho$ no $a$ is available, so $\mathsf k(\sigma g,\sigma gx)=0$ for
all $\sigma\in\pfrG^{(l)}$, hence $\mathsf k_{l+1}(\sigma g,\sigma gx)=0$ for $\sigma\in\pfrG^{(l)}$ and
$\pfrB_l=0$ ($\theta_l$ is carried by $\pfrG^{(l)}$). Let $l<\varrho$ and $\sigma\in\pfrG^{(l)}$. If
$\pfrdep_l(\sigma g,x)\ge D$, then $|\mathsf k_{l+1}(\sigma g,\sigma gx)|=\mathsf d_l(\sigma g)$, which is the
term $a=l$. Otherwise, for $\tau\in\pfrG^{(l+1)}$, $\pfrdep_l(\tau\sigma g,x)=\pfrdep_l(\sigma g,x)<D$ by (C3),
so the hypothesis for $\tau\sigma\in\pfrG^{(l)}$ gives a deep vector with $a\ge l+1$; by (C1) and the
induction hypothesis at $(l+1,\sigma g)$,
$|\mathsf k_{l+1}(\sigma g,\sigma gx)|\le\pfrB_{l+1}(\sigma g,\sigma gx)\le\sum_{l<a<\varrho}\int\mathsf
d_a(\tau\sigma g)\,d\theta_{[l+1,a+1)}(\tau)$. Integrating over $\sigma$ against $\theta_l$ and using
$\theta_{[l,a+1)}=\theta_{[l+1,a+1)}*\theta_l$ gives the claim.
\end{proof}

\begin{lemma}\label{pfr:invweight-floors}
For $x\in\pfrP$, $l\le a<\varrho$ and every $g$,
$\int\pfrwt^-_{a+1}(\sigma g)\,d\theta_{[l,a+1)}(\sigma)\le\Lambda^{10^4n(a-l+1)}\,\pfrwt^-_l(g)$.
\end{lemma}
\begin{proof}
By induction it suffices to show $\int\pfrwt^-_{b+1}(\sigma g)\,d\theta_b(\sigma)\le\Lambda^{10^4n}\pfrwt^-_b(g)$
for every $g$ and $b<\varrho$. For $\sigma\in\pfrG^{(b)}$ with $\sigma g$ near at $b+1$, $g$ is near at $b$
(as $\pfrR^{(b+1)}\pfrG^{(b)}\subseteq\pfrR^{(b)}$) and $\pfrwt^-_{b+1}(\sigma g)=\pfrdep_b(\sigma
g,x)^{s}\pfrwt^-_b(g)$ by (C3); if $\sigma g$ is not near at $b+1$ the integrand is $0$. Lemma~\ref{pfr:depth-rare}
concludes.
\end{proof}

\begin{lemma}[The gain on one pivot]\label{pfr:pivotdeep}
Let $x\in\pfrP$, $l\le\varrho$, $\varepsilon>0$, and let $g$ satisfy the hypothesis of
Lemma~\ref{pfr:cut} with $D=\varepsilon^{-1}$. Then
$\pfrB_l(g,gx)^2\le\varepsilon^{s}\,C_\Gamma\,\mathsf b_l(g)\,\pfrwt_l(g,x)^{-1}$.
\end{lemma}
\begin{proof}
Fix $a$ with $l\le a<\varrho$. Where $\mathsf d_a(y)\ne0$ we have $1\le\varepsilon\,\pfrdep_a(y,x)$, $y$ near
at $a+1$ by (C2), and so $\mathsf d_a(y)^2\le\varepsilon^s\pfrdep_a(y,x)^s|\mathsf k_{a+1}(y,yx)|^2=
\varepsilon^s\pfrwt^-_{a+1}(y)\,\mathsf a_a(y)$. By Cauchy--Schwarz, Lemma~\ref{pfr:invweight-floors} and
$\int\mathsf a_a(\sigma g)\,d\theta_{[l,a+1)}(\sigma)\le\mathsf b_l(g)$ (unfold the recursion of $\mathsf b$),
\[
\Bigl(\int\mathsf d_a(\sigma g)\,d\theta_{[l,a+1)}\Bigr)^2\le\varepsilon^s\Lambda^{10^4n\varrho}\,\pfrwt_l(g,x)^{-1}\mathsf b_l(g).
\]
By Lemma~\ref{pfr:cut} and the triangle inequality, $\pfrB_l^2\le\varrho^2\cdot\varepsilon^s\Lambda^{10^4n\varrho}
\pfrwt_l^{-1}\mathsf b_l$, and $\varrho^2\Lambda^{10^4n\varrho}\le\Lambda^{2\varrho+10^4n\varrho}\le C_\Gamma$.
\end{proof}

\begin{proof}[Proof of Theorem~\ref{pfr:deep}]
Let $y$ be a point with $\pfrB^1_j(y,yx)\pfrB^d_j(z,zx)\ne0$, $z=w_p^{-1}y$. The second factor is
bounded by Lemma~\ref{pfr:floorsq}: $\pfrB^d_j(z,zx)^2\le V\mathsf b^{(d)}_j(z)/\pfrwt_j(z,x)$. As it is not
zero, there is $\tau\in\pfrG^{(j)}$ with $W(\tau z)\ne0$ and $W(\tau dzx)\ne0$, $d=w_q^{-1}w_p$; by the cut
condition and since $\tau$ commutes with the Weyl elements, $w_p^{-1}(\rho y)$ and $w_q^{-1}(\rho yx)$ are near
at $j$ for every $\rho\in\pfrG^{(j)}$. Hence for $\rho\in\pfrG^{(j)}$ with $\mathsf k^1(\rho y,\rho yx)\ne0$, the
point $\rho y$ is active at $x$, and by hypothesis it has a deep vector of depth $\varsigma^{-1}$ with index
$\ge j$. Lemma~\ref{pfr:pivotdeep} with $l=j$, $\varepsilon=\varsigma$ gives $\pfrB^1_j(y,yx)^2\le\varsigma^sC_\Gamma
\mathsf b^{(1)}_j(y)/\pfrwt_j(y,x)$. Multiply, use $V\le V^2$, and apply Lemma~\ref{pfr:separate} with
$K=\varsigma^sV^2$.

\emph{The case $j=0$.} Here $\pfrm_0=\delta_1$, $w_p=1$, $d=1$ and the two kernels coincide, so
$\pfrF^+_0(x)=\pfrB^1_0(1,x)^2$. Every point is near at $0$ (as $\pfrR^{(0)}=\pfrG$), so every $\rho\in\pfrG$ with
$\mathsf k^1(\rho,\rho x)\ne0$ is active at $x$ and has a deep vector of depth $\varsigma^{-1}$.
Lemma~\ref{pfr:pivotdeep} with $l=0$, $g=1$, $\varepsilon=\varsigma$ gives
$\pfrB^1_0(1,x)^2\le\varsigma^sC_\Gamma\,\mathsf b^{(1)}_0(1)$ (the weight $\pfrwt_0$ is $1$), and
$\mathsf b^{(1)}_0(1)=\mathsf S_0(\mathsf k^1;x)$ since $\theta_{[0,0)}=\delta_1$. For $j=0$,
$\mathsf Q(x)=C_\Gamma^3\,\mathsf S_0(\mathsf k^1;x)^2$; hence
$\pfrF^+_0(x)^2\le\varsigma^{2s}C_\Gamma^2\,\mathsf S_0(\mathsf k^1;x)^2\le\varsigma^{2s}\,\mathsf Q(x)$, as $C_\Gamma\ge1$.
\end{proof}

\subsection{Facts on the floors: proofs of Lemmas~\ref{pfr:weylinv}, \ref{pfr:invweight}, \ref{pfr:chart}}\label{pfr:ss-floors}

Three structural facts on place data are used; they are stated exactly.

\begin{lemma}[The upper floors]\label{pfr:upperinv}
For $0\le j\le\varrho+2$ let $\rho_j=\bigl((\eta*\theta)\bigr)_{[j,\varrho+2)}$. Then $\pfrM=\rho_j*\pfrm_j$,
$\rho_j$ is carried by $\pfrG^{(j)}$, and $\rho_j*\delta_t=\rho_j$ for every $t\in\pfrG^{(j)}$.
\end{lemma}

\begin{lemma}[One Heisenberg floor after a head]\label{pfr:section}
Let $a<\varrho$, $p$ a head of length $a+1$ for $\Gamma$, $x\in\pfrP$, $y\in\pfrK$ and $\nu\in\pfrU_{(<a)}$. Then
\[
\int\mathbf 1[w_{a+1,p}n\nu y\text{ near at }a+1]\;\pfrdep_a(w_{a+1,p}n\nu y,x)^s\,d\eta_a(n)\le\mathsf c_a(p_a^{-1})
\,\Lambda^{10^6n^2}.
\]
\end{lemma}

\begin{lemma}[Near points on a floor]\label{pfr:nearsize}
Let $b\le\varrho$, $g\in\pfrK$ and $\sigma\in\pfrG^{(b)}$ with $\sigma g$ near at $b+1$. Then $\sigma^{-1}f_b$ has
entries of size $7$ for $\Gamma$.
\end{lemma}
(Lemmas~\ref{pfr:upperinv} and~\ref{pfr:nearsize} are proved below; Lemma~\ref{pfr:section} in \S\ref{pfr:ss-section}.)

\begin{proof}[Proof of Lemma~\ref{pfr:weylinv}]
Write $w=w_{j,p}^{-1}$, $\rho=\rho_j$, $\nu=\pfrm_j$. Since $w$ commutes with $\pfrG^{(j)}$, which carries $\rho$,
$\rho*\delta_w=\delta_w*\rho$. Choose $f>0$ measurable with $\int f\,d\pfrM<\infty$ ($\pfrM$ is a Haar measure,
hence $\sigma$-finite), and put $\mathrm{av}(g)=\int f(yg)\,d\rho(y)$; by Lemma~\ref{pfr:upperinv},
$\mathrm{av}(tg)=\mathrm{av}(g)$ for $t\in\pfrG^{(j)}$, and $\mathrm{av}>0$ everywhere as $\rho\ne0$ ($\pfrM\ne0$).
Let $\mu$ be any $\sigma$-finite measure with $\pfrM=\rho*\mu$. Then $\mathrm{av}<\infty$ $\mu$-a.e., and for
$\Phi\ge0$ measurable and $\pfrG^{(j)}$-invariant,
\[
\int\frac{\Phi}{\mathrm{av}}f\,d\pfrM=\int\!\!\int\frac{\Phi(yg)}{\mathrm{av}(yg)}f(yg)\,d\rho(y)\,d\mu(g)=
\int\frac{\Phi(g)}{\mathrm{av}(g)}\,\mathrm{av}(g)\,d\mu(g)=\int\Phi\,d\mu,
\]
since $\rho$ is carried by $\pfrG^{(j)}$. The left side does not depend on $\mu$. Apply this to $\mu=\nu$ and to
$\mu=\delta_w*\nu$, which is admissible because $\rho*\delta_w*\nu=\delta_w*\rho*\nu=\delta_w*\pfrM=\pfrM$
(left invariance of the Haar measure $\pfrM$). Hence $\int\Phi\,d(\delta_w*\nu)=\int\Phi\,d\nu$, that is
$\int\Phi(wg)\,d\nu(g)=\int\Phi\,d\nu$.
\end{proof}

\begin{proof}[Proof of Lemma~\ref{pfr:invweight}]
\emph{Step 1: a pointwise bound.} We show: for $g$ near at $j$,
\begin{equation}\label{pfr:eq-iw}
\int\pfrwt^-_j(w_{j,p}\nu g)\,d\eta_{[0,j)}(\nu)\le\prod_{a<j}\mathsf c_a(p_a^{-1})\cdot C_\Gamma,
\end{equation}
where $\eta_{[0,j)}=\eta_{j-1}*\dots*\eta_0$ is carried by $\pfrU_{(<j)}$. Let $\xi g\in\pfrK$ with
$\xi\in\pfrR^{(j)}$ and write $\xi=\gamma n_0$ with $\gamma\in\pfrG^{(j)}$, $n_0\in\pfrU_{(<j)}$. The measure
$\eta_{[0,j)}$ is invariant under right translation by $n_0$ and under conjugation by $\gamma$ (which
normalises each Heisenberg group of floor $<j$ and preserves its Haar measure), $\gamma$ commutes with
$w_{j,p}$, and $\pfrwt^-_j$ is $\pfrG^{(j)}$-invariant by (C3); so we may replace $g$ by $y=\xi g\in\pfrK$. (The right
invariance: by the last step of the proof of Lemma~\ref{pfr:rearr}, applied repeatedly, $\eta_{[0,j)}=
\eta_0*\dots*\eta_{j-1}$. As $\pfrU_{(<j)}=\pfrN_0\cdots\pfrN_{j-1}$, it suffices to treat $n\in\pfrN_a$, $a<j$.
Write $\eta_{[0,j)}=\eta_0*\dots*\eta_{a-1}*\eta_a*A$ with $A=\eta_{a+1}*\dots*\eta_{j-1}$, carried by
$\pfrG^{(a+1)}$ since $\pfrN_b\subseteq\pfrG^{(b)}\subseteq\pfrG^{(a+1)}$ for $b>a$. For $x\in\pfrG^{(a+1)}$,
$n_x:=xnx^{-1}\in\pfrN_a$ by (N2), so for $\Phi\ge0$, $\iint\Phi(zxn)\,d\eta_a(z)\,dA(x)=\iint\Phi(zn_xx)\,d\eta_a(z)\,dA(x)
=\iint\Phi(zx)\,d\eta_a(z)\,dA(x)$ by the invariance $\eta_a*\delta_{n_x}=\eta_a$ of (N2). Thus
$\eta_a*A*\delta_n=\eta_a*A$, and so $\eta_{[0,j)}*\delta_n=\eta_{[0,j)}$.) We
prove by induction on $a\le j$ that $I_a:=\int\pfrwt^-_a(w_{a,p}\nu y)\,d\eta_{[0,a)}(\nu)\le\prod_{l<a}\mathsf
c_l(p_l^{-1})\cdot\Lambda^{10^6n^2a}$. For $a=0$ the integrand is $\le1$ ($\pfrwt_0=1$). For the step, write
$\eta_{[0,a+1)}=\eta_a*\eta_{[0,a)}$ and $\nu'=n\nu$. For $n$ in the Heisenberg group of floor $a$, the
weight of length $a$ at $w_{a+1,p}n\nu y$ equals that at $w_{a,p}\nu y$ ($w_{a+1,p}=w_{a,p}\cdot$(the factor of
floor $a$), and both the factor of floor $a$ and $n$ lie in $\pfrG^{(a)}$, under which the depths of index
$<a$ are invariant); if the point is near at $a+1$ it is near at $a$, and
$\pfrwt^-_{a+1}=\pfrwt^-_a\cdot\pfrdep_a^{s}$. Hence, by Lemma~\ref{pfr:section},
$\int\pfrwt^-_{a+1}(w_{a+1,p}n\nu y)\,d\eta_a(n)\le\pfrwt^-_a(w_{a,p}\nu y)\,\mathsf c_a(p_a^{-1})\Lambda^{10^6n^2}$, and
integrating over $\nu$ gives the step. Finally $\Lambda^{10^6n^2j}\le\Lambda^{10^6n^3}=C_\Gamma$.

\emph{Step 2.} $\pfrm_j=\eta_{[0,j)}*\theta_{[0,j)}$ up to the order of the factors; more precisely, since the
Heisenberg groups of floors $<j$ are normalised by the floor measures of the lower floors, we
write $\int\Psi\,d\pfrm_j=\int\!\int\Psi(\nu y)\,d\eta_{[0,j)}(\nu)\,d\theta_{[0,j)}(y)$ (the ``ascending''
convolution of the Heisenberg measures followed by that of the floors). For fixed $y$,
$\Phi(\nu y)=\Phi(y)$ for $\nu\in\pfrU_{(<j)}\subseteq\pfrR^{(b')}$; if $y$ is not near at $b'$, $\Phi(y)=0$;
otherwise $y$ is near at $j$ ($j\le b'$) and \eqref{pfr:eq-iw} applies. Integrating over $y$ gives the claim.
\end{proof}
(The rearrangement is Lemma~\ref{pfr:rearr}.)

\begin{lemma}[Right translation on one floor]\label{pfr:floorright}
Let $b\le\varrho+1$, $\xi\in\pfrR^{(b)}$ and $\Psi\ge0$ measurable on $\pfrG$. Assume either
(a) $\Psi(\zeta\sigma)=\Psi(\sigma)$ for all $\zeta\in\pfrR^{(b+1)}$ and $\sigma\in\pfrG$; or
(b) $\xi\in\pfrG^{(b)}$ and $\Psi(\zeta\sigma)=\Psi(\sigma)$ for all $\zeta\in\pfrR^{(b+1)}\cap\pfrG^{(b)}$ and
$\sigma\in\pfrG^{(b)}$.
Then $\int\Psi(\sigma\xi)\,d\theta_b(\sigma)=\int\Psi(\sigma)\,d\theta_b(\sigma)$.
\end{lemma}
\begin{proof}[Proof of Lemma~\ref{pfr:chart}]
Write $b=a$ and $\Phi=c\cdot\mathbf 1[\text{near at }b+1]$, which is $\pfrR^{(b+1)}$-invariant on the left. If
$g$ is not near at $b$, then no $\sigma g$, $\sigma\in\pfrG^{(b)}$, is near at $b+1$ (as
$\pfrR^{(b+1)}\pfrG^{(b)}\subseteq\pfrR^{(b)}$), and the integral is $0$ since $\theta_b$ is carried by $\pfrG^{(b)}$.
Otherwise let $\xi\in\pfrR^{(b)}$ with $g'=\xi g\in\pfrK$. By Lemma~\ref{pfr:floorright}(a) applied to
$\Psi(\sigma)=\Phi(\sigma g)$, $\int\Phi(\sigma g)\,d\theta_b(\sigma)=\int\Phi(\sigma\xi g)\,d\theta_b(\sigma)=\int\Phi(\sigma g')\,d\theta_b(\sigma)$. For
$\sigma\in\pfrG^{(b)}$ with $\sigma g'$ near at $b+1$, $\sigma^{-1}f_b\in\Gamma(7)$ by Lemma~\ref{pfr:nearsize}, and
$\Phi(\sigma g')\le\Psi_{g'}(\sigma^{-1}f_b)$. As $\theta_b$ maps to $c_b\mu_b$ under $\sigma\mapsto\sigma^{-1}f_b$, the
integral is at most $c_b\int_{\Gamma(7)}\Psi_{g'}\,d\mu_b\le c_bC$, with $c_b\le2$ (\S\ref{pfr:ss-setting}).
\end{proof}

\subsubsection*{The Heisenberg groups and the floors}
We record the elementary group structure behind \S\ref{pfr:ss-setting}. Let $l\le\varrho$, $V_l$ the span of the
basis vectors of floor $\ge l$, and $\pfrN_l$ the Heisenberg group of floor $l$: the isometries of $V_l$
(extended by the identity) of the form $n_{\omega,\beta}\colon f_l\mapsto f_l$, $v\mapsto v-h(v,\omega)f_l$ for $v\perp
f_l,f_l^*$ in $V_{l+1}$-part, $f^*_l\mapsto f_l^*+\omega+\beta f_l$, with $\beta+\bar\beta=-h(\omega,\omega)$ (for $l=\varrho$
we have $f_\varrho=z$ and $\pfrN_\varrho=1$ in our normalisation, the top floor being handled by $\theta_\varrho$ alone).
\begin{enumerate}
\item[(N1)] \emph{$\pfrR^{(l+1)}\cap\pfrG^{(l)}=\pfrG^{(l+1)}\pfrN_l$.} An element $t$ of the left side fixes the basis
vectors of floor $<l$, and $tf_l\in f_l+\langle f_b:b<l\rangle$ is orthogonal to the $e_b^*$, $b<l$, hence $tf_l=f_l$.
So $t$ stabilises $f_l$, hence $f_l^\perp\cap V_l$ and the flag $Af_l\subset f_l^\perp\cap V_l\subset V_l$; the
stabiliser of this flag in $\auU(V_l)$ fixing $f_l$ is $\pfrN_l\rtimes\auU(V_{l+1})$, and $\auU(V_{l+1})=\pfrG^{(l+1)}$.
\item[(N2)] \emph{$\eta_l$ is a Haar measure of $\pfrN_l$}, so $\eta_l*\delta_n=\delta_n*\eta_l=\eta_l$ for $n\in\pfrN_l$;
and for $u\in\pfrG^{(l+1)}$, $u\,n_{\omega,\beta}u^{-1}=n_{u\omega,\beta}$, a linear change of the coordinate
$\omega$ of determinant of modulus $1$, which preserves $\eta_l$; hence $\delta_u*\eta_l=\eta_l*\delta_u$.
\item[(N3)] $\pfrN_a$ for $a<l$ is normalised by $\pfrG^{(l)}$ with the same formula, so the ascending convolution
$\eta_0*\eta_1*\dots*\eta_{l-1}$ commutes with $\delta_u$ for $u\in\pfrG^{(l)}$.
\end{enumerate}

\begin{lemma}\label{pfr:rearr}
$\pfrm_j=(\eta_0*\dots*\eta_{j-1})*\theta_{[0,j)}$.
\end{lemma}
\begin{proof}
Induction on $j$; $j=0$ is trivial. As $\theta_j$ is carried by $\pfrG^{(j)}$, (N3) gives $\theta_j*(\eta_0*\dots*\eta_{j-1})=
(\eta_0*\dots*\eta_{j-1})*\theta_j$ (both sides integrate $\Phi$ to $\iint\Phi(us)\,d\eta(u)\,d\theta_j(s)$ after
the substitution $u\mapsto sus^{-1}$). Hence $\pfrm_{j+1}=\eta_j*\theta_j*\pfrm_j=\eta_j*\theta_j*(\eta_0*\dots*\eta_{j-1})*
\theta_{[0,j)}=\eta_j*(\eta_0*\dots*\eta_{j-1})*\theta_j*\theta_{[0,j)}$. Finally $\eta_j*(\eta_0*\dots*\eta_{j-1})=
\eta_0*\dots*\eta_{j}$: by (N3), conjugation by $n\in\pfrN_j\subseteq\pfrG^{(j)}$ preserves each $\eta_a$, $a<j$
($\omega\mapsto n\omega$ has modulus $|\det n|=1$), hence preserves $\nu'=\eta_0*\dots*\eta_{j-1}$; so $\iint\Phi(nx)\,d\nu'(x)\,
d\eta_j(n)=\iint\Phi(xn)\,d\nu'(x)\,d\eta_j(n)$, that is $\eta_j*\nu'=\nu'*\eta_j$ (details in \S\ref{pbv:ss-weak}(i)).
\end{proof}

\begin{proof}[Proof of Lemma~\ref{pfr:floorright}]
In case (b) put $\gamma=\xi$. In case (a) write $\xi=\gamma\nu$ with $\gamma\in\pfrG^{(b)}$ and
$\nu\in\pfrU_{(<b)}$: an element of $\pfrR^{(b)}$ preserves the flag $\langle e_0\rangle\subset\dots\subset\langle
e_0,\dots,e_{b-1}\rangle$ and acts unitriangularly on it, so it lies in the parabolic subgroup of this flag with trivial
torus part, which is $\pfrU_{(<b)}\rtimes\auU(V_b)$; its Levi component $\gamma$ lies in $\auU(V_b)=\pfrG^{(b)}$. For
$\sigma\in\pfrG^{(b)}$, $\sigma\gamma\nu=(\sigma\gamma\nu(\sigma\gamma)^{-1})\,\sigma\gamma$ and $\sigma\gamma\nu(\sigma\gamma)^{-1}\in\pfrU_{(<b)}
\subseteq\pfrR^{(b+1)}$ (conjugation by $\pfrG^{(b)}$ preserves $\pfrU_{(<b)}$), so $\Psi(\sigma\xi)=\Psi(\sigma\gamma)$.
In both cases $\Psi$ restricted to $\pfrG^{(b)}$ is invariant on the left under $\pfrR^{(b+1)}\cap\pfrG^{(b)}$, and
only this is used below.
If $b<n$: $\Psi$ restricted to $\pfrG^{(b)}$ is invariant under $\pfrG^{(b+1)}\pfrN_b=\pfrR^{(b+1)}\cap\pfrG^{(b)}$, which
is the stabiliser of $f_b$ in $\pfrG^{(b)}$ by (N1); by Lemma~\ref{pfs:section}, $\Psi(\sigma)=\Psi_0(\sigma^{-1}f_b)$ on
$\pfrG^{(b)}$ for a Borel function $\Psi_0$ on $V_b$, and $\mu_b$ is carried by the orbit of $f_b$ (Remark~\ref{pfs:witt}:
no transitivity on the whole quadric is needed). Then $\Psi(\sigma\gamma)=\Psi_0(\gamma^{-1}\sigma^{-1}f_b)$, and since
$\theta_b$ maps to $c_b\mu_b$, a multiple of the quadric measure $\mu_b$, which is $\pfrG^{(b)}$-invariant,
$\int\Psi(\sigma\gamma)\,d\theta_b=c_b\int\Psi_0(\gamma^{-1}v)\,d\mu_b(v)=c_b\int\Psi_0\,d\mu_b=\int\Psi\,d\theta_b$. If $b=n$, $\theta_n=\mu_\pfrH$
is right invariant under $\pfrH=\pfrG^{(n)}$. If $b>n$, $\pfrG^{(b)}=1$.
\end{proof}

\begin{proof}[Proof of Lemma~\ref{pfr:upperinv}]
The factorisation is associativity of convolution, and the support claim holds factor by factor ($\eta_a$,
$\theta_a$ are carried by $\pfrG^{(a)}\subseteq\pfrG^{(j)}$ for $a\ge j$). Invariance by descending induction on
$j$; for $j\ge\varrho+2$, $\pfrG^{(j)}=1$. Let $\rho_j=\rho'*\theta_j$ with $\rho'=\rho_{j+1}*\eta_j$. For
$t=un\in\pfrR^{(j+1)}\cap\pfrG^{(j)}$ ((N1)), $\rho'*\delta_t=\rho_{j+1}*\delta_u*\eta_j*\delta_n=\rho_{j+1}*\eta_j=\rho'$
by (N2) and the induction hypothesis. Let $f\ge0$ and $\mathrm{av}(y)=\int f(xy)\,d\rho'(x)$; then $\mathrm{av}(ty)=
\mathrm{av}(y)$ for such $t$, and for $s\in\pfrG^{(j)}$, by Lemma~\ref{pfr:floorright}(b) (with $\xi=s$),
$\int f\,d(\rho_j*\delta_s)=\int\mathrm{av}(ys)\,d\theta_j(y)=\int\mathrm{av}(y)\,d\theta_j(y)=\int f\,d\rho_j$.
\end{proof}

\begin{proof}[Proof of Lemma~\ref{pfr:nearsize}]
Let $\xi\in\pfrR^{(b+1)}$ with $k:=\xi\sigma g\in\pfrK$. Then $\sigma^{-1}\xi^{-1}=g\,k^{-1}$, a product of an element of
$\pfrK$ and the inverse of one; by (G1) and the closure properties of the filtration it has entries of
size $\le7$ ($k^{-1}=J^{-1}k^*J$ has size $5$ by (G0)).
Now $\xi^{-1}f_b=f_b+\sum_{c<b}\lambda_cf_c$ and $\sigma^{-1}f_c=f_c$ for $c<b$; so the coordinates of
$\sigma^{-1}f_b$ of floor $\ge b$ are those of the column $\sigma^{-1}\xi^{-1}f_b$, and those of floor $<b$ vanish
($\sigma^{-1}$ preserves $V_b$). Hence $\sigma^{-1}f_b\in\Gamma(7)$.
\end{proof}

\subsection{Quadrics, sections and the proof of Lemmas~\ref{pfr:quadricX}, \ref{pfr:section}}\label{pfr:ss-section}

\subsubsection*{Coordinates}
For $b\le\varrho$ let $V_b$ be the span of the basis vectors of floor $\ge b$ and let its \emph{floor basis}
consist of $f_b,f_b^*$, the remaining hyperbolic pairs of $V_b$, and the orthonormal vectors $x_i$. In
these coordinates $h$ is the standard form; the coordinates of $v$ on the $x_i$ of $X'$ are the $h(v,x_i)$, the
coordinates of $v_{X'}$ (\S\ref{pfr:ss-gauge}). A vector of size $j$ has floor coordinates of size $j+4$ (they are entries of
$B^*Jv$ with $B$ the basis matrix of size $1$, by (G0)). Recall (G3): $\pbvol\Gamma(i)\le\Lambda^{i+1}$, and
(G4): Lebesgue measure on $A^{2m+2}$ disintegrates along the standard form, $dy=\int\kappa_\beta\,d\beta$,
with $\kappa_\beta(\Gamma(i)^{2m+2})\le\Lambda^{(2m+2)(i+1)}$ for every value $\beta$.

\begin{lemma}[A solved block]\label{pfr:block}
Let $\varrho_b$ be Lebesgue measure on $V_b$ or the quadric measure $\mu_b$ (the disintegration of
Lebesgue measure of $V_b$ along $y\mapsto h(y,y)$, at the value $h(f_b,f_b)$, given by \eqref{pfr:eq-slab}). Fix a block $\beta$ of $2m+2$ of
the orthonormal coordinates, disjoint from those of $X'$ (possible by $2(2m+2)\le\#\iota$). Then, in floor
coordinates $y=(y',y_\beta)$, $\varrho_b=\int\kappa'_{y'}\,dy'$ with Lebesgue measure $dy'$ in the coordinates off
the block and measures $\kappa'_{y'}$ on the block with $\kappa'_{y'}(\Gamma(i)^{2m+2})\le\Lambda^{(2m+2)(i+2)}$ for
every $i$ and every $y'$.
\end{lemma}
\begin{proof}
For Lebesgue measure take $\kappa'_{y'}$ Lebesgue on the block, of mass $\pbvol\Gamma(i)^{2m+2}\le\Lambda^{(2m+2)(i+1)}$ on
$\Gamma(i)^{2m+2}$, and $\Lambda^{(2m+2)(i+1)}\le\Lambda^{(2m+2)(i+2)}$. For $\mu_b$ put $L_i=\Lambda^{(2m+2)(i+2)}$ and let
$U_i$ be the interior of $\Gamma(i+1)^{2m+2}$; then $\Gamma(i)^{2m+2}\subseteq U_i$ (at the real place
$\Gamma(i)=\{|c|\le L^i\}$ with $L\ge r+1\ge2$, so $\Gamma(i)\subseteq\{|c|<L^{i+1}\}$; at the finite places $\Gamma(i)$,
defined by an integrality condition, is open) and $\bigcup_iU_i=A^{2m+2}$. Let $t_0=h(f_b,f_b)$.

\emph{Step 1: $\mu_b(E'\times U_i)\le L_i\,\pbvol(E')$ for every Borel set $E'$ of the coordinates off the block.} Let
$\varphi'\ge0$ and $\chi$ with $0\le\chi\le\mathbf 1_{U_i}$ be test functions as in \eqref{pfr:eq-slab} in the two groups
of coordinates. We have $h(y,y)=h(y',y')+h(y_\beta,y_\beta)$ with the standard form on the block, and by (G4)
$dy_\beta=\int\kappa_s\,ds$. By Tonelli, for each $y'$ the substitution $s=t-h(y',y')$ (a translation, preserving $dh$)
and $\kappa_s(\chi)\le\kappa_s(\Gamma(i+1)^{2m+2})\le L_i$,
\[
\int_{\{h(y,y)-t_0\in B\}}\varphi'(y')\chi(y_\beta)\,dy=\int\varphi'(y')\int_{\{s:\ s+h(y',y')-t_0\in B\}}\kappa_s(\chi)\,ds\,dy'
\le\pbvol(B)\,L_i\int\varphi'\,dy'.
\]
Dividing by $\pbvol(B)$ and letting $B\to0$, \eqref{pfr:eq-slab} gives $\int\varphi'\otimes\chi\,d\mu_b\le L_i\int\varphi'dy'$.
Taking the supremum over $\chi\uparrow\mathbf 1_{U_i}$ ($U_i$ is open; monotone convergence), then the infimum over
$\varphi'\downarrow\mathbf 1_{E'}$ for $E'$ compact (dominated convergence, $\mu_b$ being Radon), gives the claim for
compact $E'$. The Borel measure $\pi_i(E')=\mu_b(E'\times U_i)$ is thus finite on compact sets, hence inner regular,
and the claim follows for every Borel $E'$.

\emph{Step 2: disintegration.} Let $w(y)=\sum_i2^{-i}(1+L_i)^{-1}\mathbf 1_{U_i}(y_\beta)$, a Borel function with
$0<w\le1$. By Step 1 the image $\pi_w$ of $w\mu_b$ under $y\mapsto y'$ satisfies $\pi_w(E')\le\sum_i2^{-i}\pbvol(E')
\le2\pbvol(E')$, so $\pi_w=g\,dy'$ with $0\le g\le2$. By the disintegration theorem (all spaces are locally compact
and second countable), $w\mu_b=\int\lambda_{y'}\,d\pi_w(y')$ with probability measures $\lambda_{y'}$ on the fibre
$\{y'\}\times A^{2m+2}$, identified with the block. Put $\kappa'_{y'}=g(y')\,w^{-1}\lambda_{y'}$; then
$\mu_b=\int\kappa'_{y'}\,dy'$, and for every Borel $E'$, $\int_{E'}\kappa'_{y'}(U_i)\,dy'=\pi_i(E')\le L_i\pbvol(E')$, so
$\kappa'_{y'}(U_i)\le L_i$ for almost every $y'$. Replacing $\kappa'_{y'}$ by $0$ on the null set of the $y'$ where
this fails for some $i$ (countably many conditions), $\kappa'_{y'}(\Gamma(i)^{2m+2})\le\kappa'_{y'}(U_i)\le L_i$ for
every $i$ and $y'$.
\end{proof}

\begin{proof}[Proof of Lemma~\ref{pfr:quadricX}]
More generally, for $\varrho_b$ as in Lemma~\ref{pfr:block} and $\int_{\Gamma(j+4)^{X'}}\phi\le B$ we show
$\int_{\Gamma(j)}\phi(v_{X'})\,\allowbreak d\varrho_b(v)\le\Lambda^{99n(j+6)}B$; the case $j=7$, $\phi\mapsto\phi^s$ is the lemma. A
vector of size $j$ has floor coordinates in $\Gamma(j+4)$, so the integral is at most
$\int_{\Gamma(j+4)}\phi(y_{X'})\,d\kappa'_{y'}(y_\beta)\,dy'$ over $y'$ in $\Gamma(j+4)$. Integrate first over the block
($\le\Lambda^{(2m+2)(j+6)}$), then over the coordinates of $X'$ ($\le B$), then over the remaining $d$ coordinates
($\le(\Lambda^{j+5})^d$). The number of coordinates of $V_b$ off $X'$ is at most $2+2(m+\varrho)+4(m+1)\le99n-(2m+2)$
(by $m\le9n$); so the product is at most $\Lambda^{(2m+2)(j+6)+(99n-(2m+2))(j+5)}B\le\Lambda^{99n(j+6)}B$.
\end{proof}

\subsubsection*{The Heisenberg group of floor $a$ in coordinates}
For $a<\varrho$ the elements of $\pfrN_a$ are $n=n(\omega,\beta)$ with $\omega\in V_{a+1}$ ($r_a:=d_a=\dim V_{a+1}$
coordinates) and $\beta=b-\tau_0h(\omega,\omega)$, $b\in A_{\rm sk}$; as defined in \S\ref{pfr:ss-setting}, $\eta_a$ is
the image of $d\omega\otimes d\beta_{\rm sk}(b)$, where $d\beta_{\rm sk}$ is the Haar measure of the skew elements for which
$(\beta_1,\beta_2)\mapsto\tau_0\delta^{-1}\beta_1+\beta_2$ is a measure-preserving isomorphism of $A_{\rm sk}^2$ onto $A$
(Lemma~\ref{pfs:skewhaar}; $\tau_0\delta^{-1}\in\Gamma(1)$ by (G1)). As $n^{-1}=n(-\omega,\bar\beta)$,
$n^{-1}e_a^*=e_a^*-\omega+\mu e_a$ with $\mu=\bar\beta=-b-\bar\tau_0h(\omega,\omega)$; in the decomposition $\mu=\tau_0(\mu+\bar\mu)
+(\mu-\tau_0(\mu+\bar\mu))$ its skew part is $-b+(\tau_0-\bar\tau_0)h(\omega,\omega)$, that is $-b$ up to a term depending only
on $\omega$. Below, ``the $\beta$-integral'' is the integral in $b$ against $d\beta_{\rm sk}$.

\begin{lemma}[Skew sets]\label{pfr:skew}
If $S$ is a set of skew elements and $u$ a unit with $u^{-1}(\beta-\beta')\in\Gamma(j+1)$ for $\beta,\beta'\in S$, then
$\pbvol_{\rm sk}(S)\le\pfrmod(u)^{1/2}\Lambda^{j+2}$.
\end{lemma}
\begin{proof}
The image of $S\times S$ under $(\beta_1,\beta_2)\mapsto\tau_0\delta^{-1}\beta_1+\beta_2$ has measure $\pbvol_{\rm sk}(S)^2$, and,
after fixing a point $(\beta^0_1,\beta^0_2)$ of $S\times S$, all its points $c$ satisfy
$u^{-1}(c-c^0)\in\Gamma(1+(j+1)+1)$ by (G0). Hence $\pbvol_{\rm sk}(S)^2\le\pfrmod(u)\pbvol\Gamma(j+3)\le\pfrmod(u)\Lambda^{j+4}
\le(\pfrmod(u)^{1/2}\Lambda^{j+2})^2$.
\end{proof}

\begin{lemma}[Integration over $\pfrN_a$]\label{pfr:heis}
Let $u$ be a unit, $\Psi\ge0$ measurable on $V_{a+1}$ with $\int_{\Gamma(j)}\Psi\le I$, $C\ge0$, and $\Phi\ge0$
measurable on $\pfrG$ such that whenever $\Phi(n(\omega,\beta))\ne0$: $u^{-1}(n^{-1}-1)e_a^*$ is congruent modulo
$\langle e_b:b<a\rangle$ to a vector of size $j$, and $\Phi(n)\le C\Psi(u^{-1}\omega)$. Then
$\int\Phi\,d\eta_a\le C\,\mathsf c_a(u)\,\Lambda^{j+2}I$, where $\mathsf c_a(u)=\pfrmod(u)^{d_a+1/2}$.
\end{lemma}
\begin{proof}
Substitute $\omega=uv$ (Jacobian $\pfrmod(u)^{d_a}$). For fixed $v$, on the set $S$ of the skew coordinates $b$
(with $n=n(\omega,b-\tau_0h(\omega,\omega))$) with $\Phi(n)\ne0$, the vector $u^{-1}(n^{-1}-1)e_a^*=u^{-1}(-\omega+\mu e_a)$ has,
modulo the $e_{b'}$ ($b'<a$), size $j$: its $\omega$-part gives $v\in\Gamma(j)$ and its $e_a$-coordinate gives
$u^{-1}\mu=-u^{-1}(b+c_v)\in\Gamma(j)$ for the number $c_v=\bar\tau_0h(\omega,\omega)$ depending only on $v$; so differences
satisfy $u^{-1}(b-b')\in\Gamma(j+1)$, and by Lemma~\ref{pfr:skew} the $b$-integral is at most $\pfrmod(u)^{1/2}\Lambda^{j+2}C\Psi(v)\mathbf 1[v\in\Gamma(j)]$. Integrating over $v$ gives the
claim.
\end{proof}

For a unit $c$ the \emph{section} of the quadric of $f_a$ ($h(y,y)=0$ in $V_a$) at $h(y,e_a)=c$ (the
$e_a^*$-coordinate) is parametrised by
$n\mapsto y_c(n)=c\,n^{-1}e_a^*$, $n\in\pfrN_a$; its \emph{section measure} is $\pfrmod(c)^{d_a}$ times the image of
$\eta_a$. By the definition of $\theta_a=\lambda_a*\eta_a$ ($\lambda_a$ the image of $\pfrmod(c)^{d_a}dc$ under
$c\mapsto w_{a,c}^{-1}$, and $(w_{a,c}^{-1}n)^{-1}e_a=c\,n^{-1}e_a^*$), the image of $\theta_a$ under $\sigma\mapsto\sigma^{-1}e_a$ is
$\int(\text{section measure at }c)\,dc$, which is $c_a\mu_a$ by Lemma~\ref{pfr:floorquad}.

\begin{lemma}[Integration over a section]\label{pfr:sectionint}
If $\int_{\Gamma(j+5)^{X'}}\phi\le B$, then $\int_{\Gamma(j)}\phi(y_{X'})\,d(\text{section at }c)(y)\le
\pfrmod(c^{-1})^{1/2}\Lambda^{j+3}\Lambda^{99n(j+7)}B$.
\end{lemma}
\begin{proof}
$y_c(n)=c(e_a^*-\omega+\mu e_a)$, so $y_c(n)_{X'}=(-c\,\omega)_{X'}$ and $c(n^{-1}-1)e_a^*=y_c(n)-ce_a^*$, which has size
$j+1$ when $y_c(n)$ has size $j$ (then $c=h(e_a,y)$-coordinate has size $j$ and (G0)); hence $(-c)(n^{-1}-1)e_a^*$ has size $j+1$ too ((G0): $\Gamma(i)=-\Gamma(i)$). Apply Lemma~\ref{pfr:heis} with
$u=(-c)^{-1}$, $j+1$ in place of $j$, $C=1$, $\Psi(v)=\phi(v_{X'})$ on $V_{a+1}$, and $I=\Lambda^{99n(j+7)}B$ from the
general form of Lemma~\ref{pfr:quadricX} for Lebesgue measure on $V_{a+1}$ (with $j+1$ in place of $j$). The result is
$\mathsf c_a(c^{-1})\Lambda^{j+3}\Lambda^{99n(j+7)}B$; multiplying by $\pfrmod(c)^{d_a}$ leaves $\pfrmod(c^{-1})^{1/2}$.
\end{proof}

\subsubsection*{Sizes at a head}
\begin{lemma}[Entries at a head]\label{pfr:headentries}
Let $j\le\varrho$, $p$ a head of length $j$, $g\in\pfrG$, $s\in\pfrR^{(j)}$ with $sg\in\pfrK$ and $\eta\in\pfrR^{(j)}$ with
$\eta w_p^{-1}g\in\pfrK$, and $a<j$. Then $p_a$, $p_a^{-1}$ and the entries of the column $se_a^*$ have size $10n$.
\end{lemma}
\begin{proof}
Let $k_1=sg$, $k_2=\eta w_p^{-1}g$, of size $1$ by (G1); then $\gamma:=k_1k_2^{-1}=sw_p\eta^{-1}$ has size $1+5+1=7$ by
(G0). \emph{Products.} For $u\in\pfrR^{(j)}$, $uf_c\in f_c+\langle f_{c'}:c'<c\rangle$ for $c<j$, and the same holds for
$u^{-1}$. Hence, for $l\le j$, the leading $l\times l$ block of the frame matrix of $\gamma$ is, since
$h(sw_p\eta^{-1}f_c,f_{c'})=h(w_p\eta^{-1}f_c,s^{-1}f_{c'})$, a product $\mathsf L\,D_l\,\mathsf R$ with unitriangular
$\mathsf L,\mathsf R$ and $D_l=(h(w_pf_c,f_{c'}))_{c,c'<l}=\operatorname{diag}(p_c)$ up to conjugation of the entries
(as $w_pe_c=p_ce_c^*$); so its determinant is $P_l=p_0\cdots p_{l-1}$ or $\bar P_l$. The entries of $Z(\gamma)$ are entries of $J\gamma$, of size
$9$; so $P_l$ (or $\bar P_l$, the same size) has size $10l\le10\varrho$. As $P_a^{-1},P_{a+1}^{-1}\in\Gamma(1)$ (head),
$p_a=P_a^{-1}P_{a+1}$ and $p_a^{-1}=P_{a+1}^{-1}P_a$ have size $10\varrho+1\le10n$.
\emph{Columns} (Cramer's rule). Fix a basis index $i$. Let $Z'$ be the $(a+1)\times(a+1)$ block of $\gamma$ with rows
$e_0^*,\dots,e_{a-1}^*,i$ and columns $e_0,\dots,e_a$; it has size $7$. Since $\eta^{-1}e_d\in e_d+\langle e_c:c<d\rangle$,
$Z'=\mathsf L'\mathsf C$ with $\mathsf C$ the unitriangular block of $\eta^{-1}$ on $e_0,\dots,e_a$ and $\mathsf L'$ the block
of $sw_p$ with the same rows and columns $e_0,\dots,e_a$; as $w_pe_d=p_de_d^*$, $\mathsf L'_{bd}=s_{R(b),e_d^*}p_d$. For
$b\le d\le a$ with $b<a$, $s_{e_b^*,e_d^*}=h(se_d^*,e_b)=h(e_d^*,s^{-1}e_b)=\delta_{bd}$ (as $s^{-1}e_b\in e_b+\langle e_c:c<b\rangle$); so $\mathsf
L'$ is lower triangular with diagonal $p_0,\dots,p_{a-1},s_{i,e_a^*}p_a$. Hence $\det Z'=P_{a+1}s_{i,e^*_a}$, of size
$(a+1)\cdot8\le8\varrho$, and $s_{i,e_a^*}=P_{a+1}^{-1}\det Z'$ has size $8\varrho+1\le10n$.
\end{proof}

\begin{proof}[Proof of Lemma~\ref{pfr:section}]
Write $w=w_{a+1,p}$ and let $\phi$ be the function of (G5) for the point $y\in\pfrK$, the position $x$, and the
$\varrho$ linear forms $\omega\mapsto(\nu\omega)_{e_b}$ on $X'$ (the $e_b$-coordinates, $b<\varrho$), with common vectors
$\mathcal W$.

\emph{The depth is bounded by $\phi$.} Let $n\in\pfrN_a$ and $\omega\in\mathcal W$. The vector $\omega'=wn\nu\omega$ is stable at
$(wn\nu y,wn\nu yx)$ for $a$: the pair $((wn\nu y)^{-1}\omega',(wn\nu yx)^{-1}\omega')=(y^{-1}\omega,(yx)^{-1}\omega)$ is close,
and for $b<a$ the $e_b^*$-coordinate of $\omega'$ is $p_b(n\nu\omega)_{e_b}=p_b(\nu\omega)_{e_b}=0$ ($n$ fixes the coordinates
of floor $<a$; $\omega$ is killed by the forms), i.e. $\omega'\perp e_b$. Its $e_a^*$-coordinate is $h(\omega',e_a)=
h(\omega,(wn\nu)^{-1}e_a)$, and $(wn\nu)^{-1}e_a=\nu^{-1}n^{-1}\bar p_ae_a^*$ is congruent modulo $\langle e_b:b<a\rangle$ to
the section point $y_{\bar p_a}(n)$ (as $\nu^{-1}\in\pfrU_{(<a)}$ and $y_{\bar p_a}(n)$ has no $e^*_b$-coordinates, $b<a$);
congruent vectors have the same projection on $X'$ (the $e_b$ are orthogonal to $X'$). So, with $u=y_{\bar p_a}(n)_{X'}$,
the $e_a^*$-coordinate of $\omega'$ is $h(\omega,u)$, $\mathfrak d_a\supseteq\{h(\omega,u):\omega\in\mathcal W\}$, and
$\pfrdep_a(wn\nu y,x)\le\phi(u)$ by the first inequality of (G5).

\emph{Size.} If $wn\nu y$ is near at $a+1$, apply Lemma~\ref{pfr:headentries} with $j=a+1$, the head $\bar p$
(for which $w_{\bar p}^{-1}=w$, \S\ref{pfr:ss-gainsproof}), $g=n\nu y$, $s=(n\nu)^{-1}\in\pfrR^{(a+1)}$ ($sg=y\in\pfrK$) and
the $\eta$ given by nearness: $\bar p_a$ and the column $(n\nu)^{-1}e_a^*$ have size $10n$, so
$\bar p_a(n\nu)^{-1}e_a^*$ has size $20n$, and so has $y_{\bar p_a}(n)$, which agrees with it in the coordinates of
floor $\ge a$ and vanishes in the others.

\emph{Integration.} Hence the integrand is at most $\Psi(y_{\bar p_a}(n))$ with $\Psi(v)=\mathbf 1[v\in\Gamma(20n)]\,
\phi(v_{X'})^s$. By the definition of the section measure and Lemma~\ref{pfr:sectionint} (with $j=20n$ and
$B=\Lambda^{99n(20n+6)}$ from (G5)),
\[
\int\Psi(y_{\bar p_a}(n))\,d\eta_a(n)=\pfrmod(p_a^{-1})^{d_a}\int\Psi\,d(\text{section at }\bar p_a)\le\pfrmod(p_a^{-1})^{d_a+1/2}
\Lambda^{20n+3+99n(20n+7)+99n(20n+6)},
\]
and the exponent is at most $10^6n^2$.
\end{proof}

\subsection{Proof of Lemma~\ref{pfr:descent}}\label{pfr:ss-descent}

Write $a=i$, $w=w_{a,p}$, $d=w_{a,q}^{-1}w_{a,p}$, and for $c\in A^\times$ let $w_c=w_{a+1,c}$ restricted to the pair
$(e_a,e_a^*)$: $w_ce_a=ce_a^*$, $w_ce_a^*=\bar c^{-1}e_a$. For points $h,x$ and an element $v$ put
\[
\Pi(x;v;h)=\mathsf k^1_{a+1}(h,hx)\,\overline{\mathsf k^{d}_{a+1}(v^{-1}h,v^{-1}hx)} .
\]
\emph{Facts.} (i) $w_{a+1,p^c}=w\,w_c$, the Weyl ratio of $(p^c,q^c)$ at length $a+1$ is $d$ (the entries $c/c$
are $1$), and $\pfrkap_{a+1}(p^c,q^c)=\pfrkap_a(p,q)\pfrmod(c)^{d_a}$. (ii) By (C1) of \S\ref{pfr:ss-gainsproof},
applied to both factors (the elements $w,w_c$ commute with $\pfrG^{(a+1)}$, which carries $\theta_{a+1}$),
$\Phi_{a+1}(x;1;g')=\Pi(x;ww_c;g')$ for the pair $(p^c,q^c)$, and by definition
$\Phi_a(x;\omega;g)=\iint\omega(\sigma\tau^{-1})\,\Pi(x;\sigma w\tau^{-1};\sigma g)\,d\theta_a(\tau)\,d\theta_a(\sigma)$. (iii) For
$t$ in $\pfrR^{(a+1)}\cap\pfrG^{(a)}$, $tf_a=f_a$ ((N1)), so $Z_{aa}(tyt')=Z_{aa}(y)$; and $\mathsf k^{1}_{a+1}(th,thx)=
\mathsf k^{1}_{a+1}(h,hx)$, $\mathsf k^{d}_{a+1}(th,thx)=\mathsf k^{d}_{a+1}(h,hx)$ by (K2), second claim, of Lemma~\ref{pfs:kernels} (applied with
$l=a$, $\zeta=t$; for the factor coming from $\pfrN_a$ one also sees directly that the character factors
$\psi(M^{-1}(\sigma(dtd^{-1})-\sigma(t)))$ are $1$ since $dtd^{-1}=t$, $d$ being trivial on the floors $\ge a$ where $t$
acts). (iv) $\theta_a$ is the image of $\pfrmod(c)^{d_a}dc\times\eta_a$ under $(c,n)\mapsto w_c^{-1}n$
(the definition of the floor measures in \S\ref{pfr:ss-setting}, at every place), and $Z_{aa}(w_c)=c$.

\begin{proof}[Proof of Lemma~\ref{pfr:descent}]
Let $\omega=\varphi\circ Z_{aa}$ with $\varphi(c)=\psi(\xi c)$, $|\varphi|\le1$. Put
$F(x,g,\sigma,\tau)=\omega(\sigma\tau^{-1})\Pi(x;\sigma w\tau^{-1};\sigma g)$ on $\pfrG^4$ with the measure
$\pfrM\otimes\pfrm_a\otimes\theta_a\otimes\theta_a$, and $E(x,g',c)=\varphi(c)\Pi(x;ww_c;g')$ with the measure
$\pfrM\otimes\pfrm_{a+1}\otimes\pfrmod(c)^{d_a}dc$.

\emph{$|F|$ is integrable.} By Tonelli and $|\omega|\le1$, $\int|F|\le\int\pfrF_a^+\,d\pfrM<\infty$: the integral over
$(\sigma,\tau)$ of $|\Pi|$ is the product of the two pivot norms, as in Lemma~\ref{pfr:L1}.

\emph{$F$ and $E$ have the same law.} Fix $x,g$ and $\sigma\in\pfrG^{(a)}$. The function $\tau\mapsto F(x,g,\sigma,\tau)$
is invariant under $\tau\mapsto t\tau$ for $t\in\pfrR^{(a+1)}\cap\pfrG^{(a)}$, by (iii) (on the first factor
$\omega(\sigma\tau^{-1}t^{-1})$, and on the second because $(\sigma w\tau^{-1}t^{-1})^{-1}\sigma g=t\tau w^{-1}g$). By
Lemma~\ref{pfr:floorright}(b) (with $\xi=\sigma\in\pfrG^{(a)}$) and (iv), for every measurable $f\ge0$,
\[
\int f(F(x,g,\sigma,\tau))\,d\theta_a(\tau)=\int f(F(x,g,\sigma,\tau\sigma))\,d\theta_a(\tau)=\iint
f\bigl(F(x,g,\sigma,w_c^{-1}n\sigma)\bigr)\,d\eta_a(n)\,\pfrmod(c)^{d_a}dc .
\]
Now $\sigma w(w_c^{-1}n\sigma)^{-1}=\sigma w\sigma^{-1}n^{-1}w_c=wn^{-1}w_c$ ($\sigma$ commutes with $w$), so
$\omega(\sigma\tau^{-1})=\varphi(Z_{aa}(n^{-1}w_c))=\varphi(c)$ by (iii), (iv); and
$\Pi(x;wn^{-1}w_c;\sigma g)=\Pi(x;ww_c;n\sigma g)$ by (iii) applied to both correlations ($n$ commutes with $w$).
Hence $F(x,g,\sigma,w_c^{-1}n\sigma)=E(x,n\sigma g,c)$. Integrating over $\sigma\sim\theta_a$, $g\sim\pfrm_a$ and $x\sim\pfrM$,
and using $\pfrm_{a+1}=(\eta_a*\theta_a)*\pfrm_a$, we get $\int f\circ F=\int f\circ E$ for all $f\ge0$.

\emph{Conclusion.} Hence $E$ is integrable, $\int E=\int F$, and by Fubini's theorem (in the order $c$ last)
$c\mapsto\pfrmod(c)^{d_a}\varphi(c)\iint\Pi(x;ww_c;g')\,d\pfrM\,d\pfrm_{a+1}$ is integrable with integral $\int F$. By
(i), (ii), the integrand is $\psi(\xi c)\pfrI_{a+1}(p^c,q^c;1)/\pfrkap_a(p,q)$, and $\pfrkap_a\int F=\pfrI_a(p,q;\omega)$ by
definition.
\end{proof}


%% ===== proofs/B-frames-shared.tex =====
% proofs/B-frames-shared.tex
% The shared estimates of the local theory of two frames: the one-frame bound (Theorem pfr:oneframe)
% and the gains from a frequency and from two heads (Theorems pfr:gain-char, pfr:gain-heads).
% Labels of this file start with pfs:. It continues proofs/B-frames-real.tex (labels pfr:).
\providecommand{\panot}[1]{\par\smallskip{\raggedright\noindent\textbf{[Not proved here:} #1\textbf{]}\par}\smallskip}
\providecommand{\pbvol}{\operatorname{vol}}
\providecommand{\pfrG}{\mathsf G}
\providecommand{\pfrH}{\mathsf H}
\providecommand{\pfrU}{\mathsf U}
\providecommand{\pfrR}{\mathsf R}
\providecommand{\pfrM}{\mathsf M}
\providecommand{\pfrm}{\mathsf m}
\providecommand{\pfrK}{\mathsf K}
\providecommand{\pfrP}{\mathcal P}
\providecommand{\pfrI}{\mathcal I}
\providecommand{\pfrF}{\mathsf F}
\providecommand{\pfrB}{\mathsf B}
\providecommand{\pfrmod}{\operatorname{mod}}
\providecommand{\pfrdep}{\operatorname{dp}}
\providecommand{\pfrwt}{\wp}
\providecommand{\pfrkap}{\kappa}
\providecommand{\pfrN}{\mathsf N}
\providecommand{\pfrip}[2]{\langle#1,#2\rangle}
% macros of this file (prefix \pfs)
\providecommand{\pfsA}{\mathsf A}
\providecommand{\pfsBall}{\mathcal B}
\providecommand{\pfscm}{\operatorname{cm}}
\providecommand{\pfscut}{\operatorname{cut}}
\providecommand{\pfsCB}{\mathcal C\!\mathcal B}
\providecommand{\pfsc}{\mathsf c}
\providecommand{\pfsrk}{\operatorname{rk}}
\providecommand{\pfsfc}{\mathsf f}
\providecommand{\pfsO}{\mathcal O}
\providecommand{\pfsst}{\operatorname{St}}

\section{Two frames: the shared estimates}\label{pfs:sec}

This section proves the three statements of \S\ref{pfr:ss-shared} stated in
\S\ref{pfr:sec}: the one-frame bound (Theorem~\ref{pfr:oneframe}), the gain from a frequency
(Theorem~\ref{pfr:gain-char}) and the gain from two heads (Theorem~\ref{pfr:gain-heads}). It also
supplies the three facts that \S\ref{pfr:sec} used as standard (the action on the quadric, measurable
sections, the Haar measure on the skew elements), in \S\ref{pfs:ss-conv}.

\emph{What is proved.} All three theorems are proved in full, from the statements of \S\ref{pfr:sec}
that are proved there and the gauge axioms (G0)--(G5), together with one property of the place that
\S\ref{pfr:sec} names (the constant $P$) but does not state; it is stated here as axiom~(P)
(\S\ref{pfs:ss-planch}). Fact (C4) of \S\ref{pfr:ss-gainsproof} is proved here as (K1)--(K3), and we
recall the convention $\pfrR^{(j)}=\pfrU$ ($j\ge\varrho+2$) in the definition of nearness
(Remark~\ref{pfs:rem-corr}).

Throughout, the assumptions of \S\ref{pfr:ss-shared} are in force: $\pfrF$ is a place datum whose group
measure $\pfrM$ is a Haar measure, $\Gamma$ is a gauge with loss $\Lambda$, $W$ is a cut Bessel function,
$n=\varrho+1$, and $\pfsrk:=\operatorname{rank}X'$.

\subsection{Conventions, kernels, and three standard facts}\label{pfs:ss-conv}

\subsubsection*{Conventions}
\begin{enumerate}
\item[(i)] The hermitian form $h$ is linear in its first and conjugate linear in its second variable
(convention (Cv1) of \S\ref{pb:ss-setting});
this is the reading under which the maps $v\mapsto v-h(v,\omega)f_l$ of \S\ref{pfr:ss-floors} are
linear. The coordinate of a vector $x$ along $e_a^*$ is $h(x,e_a)$, and along $e_a$ it is
$h(x,e_a^*)$. The scalar product $\pfrip{\cdot}{\cdot}$ on $\C^I$ is conjugate linear in its first
variable. Since $\psi$ is unitary, $\overline{\psi(c)}=\psi(-c)$.
\item[(ii)] For $l\ge0$ put $\pfrU_{(<l)}=\pfrN_0\pfrN_1\cdots\pfrN_{l-1}$ and
$\pfsA_l=\eta_0*\eta_1*\dots*\eta_{l-1}$ (the ascending convolution of \S\ref{pfr:ss-floors};
$\pfsA_0=\delta_1$, and $\pfsA_{l}=\pfsA_{\varrho}$ for $l\ge\varrho$ since $\eta_l=\delta_1$ for $l\ge\varrho$).
By Lemma~\ref{pfr:rearr} and its proof, $\pfsA_l$ is a Haar measure of the nilpotent (hence unimodular)
group $\pfrU_{(<l)}$, so $\delta_\nu*\pfsA_l=\pfsA_l*\delta_\nu=\pfsA_l$ for $\nu\in\pfrU_{(<l)}$; by (N3) of
\S\ref{pfr:ss-floors}, $\delta_\sigma*\pfsA_l=\pfsA_l*\delta_\sigma$ for $\sigma\in\pfrG^{(l)}$; and
$\pfsA_{l+1}=\pfsA_l*\eta_l=\eta_l*\pfsA_l$ (associativity, and (N3) for $\eta_l$, which is carried by
$\pfrN_l\subseteq\pfrG^{(l)}$).
\item[(iii)] The facts on $\pfrU_{(<l)}$ used in \S\ref{pfr:ss-floors}: it fixes the $e_b$, $b<l$; it acts on
$V_l$ by $v\mapsto v+\sum_{b<l}\lambda_b(v)e_b$ (each $\pfrN_b$, $b<l$, acts on $V_{b+1}\supseteq V_l$ by
$v\mapsto v-h(v,\omega)e_b$ and fixes the $e_c$, $c\le b$, and moves the $e_c$, $c>b$, by multiples of $e_b$);
it is normalised by $\pfrG^{(l)}$; and $\pfrR^{(l)}=\pfrG^{(l)}\pfrU_{(<l)}$, by (N1) applied on each floor
$b<l$. In particular, $\pfrU_{(<l)}$ acts trivially on $Y^{(l)}$.
\item[(iv)] \emph{The Heisenberg group of floor $a<\varrho$.} Write $n(\omega,\beta)$ for its elements as in
\S\ref{pfr:ss-floors}, with $\omega\in V_{a+1}$ orthogonal to $f_a,f_a^*$. Directly from the formulas:
$n(\omega,\beta)n(\omega',\beta')=n(\omega+\omega',\beta'')$ for some $\beta''$ (the vector part is additive);
$\gamma n(\omega,\beta)\gamma^{-1}=n(\gamma\omega,\beta)$ for $\gamma\in\pfrG^{(a+1)}$ ((N2)); and the entry of
$n(\omega,\beta)$ just above the diagonal of its frame block, at $(a,a+1)$, is the coefficient of $f_a$ in
$n(\omega,\beta)f_{a+1}=f_{a+1}-h(f_{a+1},\omega)f_a$, so that
\begin{equation}\label{pfs:eq-sigma}
\sigma(n(\omega,\beta))=-h(f_{a+1},\omega),\qquad
\sigma\bigl(\gamma n(\omega,\beta)\gamma^{-1}\bigr)=-h(\gamma^{-1}f_{a+1},\omega)\quad(\gamma\in\pfrG^{(a+1)}).
\end{equation}
\item[(v)] The size rules (G0) are used without comment; ``size $i$'' always refers to $\Gamma$ unless
another gauge is named. A vector of $V_b$ ``of size $i$'' is one whose entries lie in $\Gamma(i)$.
\end{enumerate}

\subsubsection*{Kernels}
A \emph{kernel at floor $j$} is $\mathsf k=\mathsf k^{d'}$, $\mathsf k^{d'}(g,y)=\pfrip{W(g)}{W(d'y)}$, where
$d'$ is diagonal in the basis, of the form $e_a\mapsto\lambda_ae_a$, $e_a^*\mapsto\bar\lambda_a^{-1}e_a^*$ ($a<j$),
and the identity on the basis vectors of floor $\ge j$; such a $d'$ is an isometry, normalises $\pfrU$, and
commutes with $\pfrG^{(j)}$. Examples: $d'=1$, and the Weyl ratio $w_{j,q}^{-1}w_{j,p}$ of two sequences of units.
A kernel at floor $j$ is a kernel at every floor $j'\ge j$. The correlations $\mathsf k_l$ are those of \S\ref{pfr:ss-setting};
facts (C1) and (C2) of \S\ref{pfr:ss-gainsproof} hold for them. We use three further facts, in the form in
which the argument needs them.

\begin{lemma}\label{pfs:kernels}
Let $\mathsf k$ be a kernel at floor $j$.
\begin{enumerate}
\item[(K1)] \textup{(Equivariance.)} For $l\le n$, $\nu,\nu'\in\pfrU_{(<l)}$ and all $g,y$,
\[
\mathsf k_{l+1}(\nu g,\nu'y)=\psi\bigl(M^{-1}(\sigma(d'\nu'd'^{-1})-\sigma(\nu))\bigr)\,\mathsf k_{l+1}(g,y).
\]
\item[(K2)] \textup{(Invariance.)} For $j\le l$, $\gamma\in\pfrG^{(l)}$ and all $h,h'$:
$\mathsf k_l(\gamma h,\gamma h')=\mathsf k_l(h,h')$. For $j\le l\le\varrho$ and $\zeta\in\pfrR^{(l+1)}\cap\pfrG^{(l)}$:
$\mathsf k_{l+1}(\zeta h,\zeta h')=\mathsf k_{l+1}(h,h')$.
\item[(K3)] For $l\le\varrho$ with $j\le l$, and every position $x$, the function
$\mathsf a_l(g)=\pfrwt_l(g,x)|\mathsf k_{l+1}(g,gx)|^2$ of \S\ref{pfr:ss-gauge} satisfies $\mathsf a_l(\xi g)=\mathsf
a_l(g)$ for $\xi\in\pfrR^{(l+1)}$; for every $l$, $|\mathsf k_{l+1}(g,\nu y)|=|\mathsf k_{l+1}(g,y)|$ for
$\nu\in\pfrU_{(<l)}$.
\end{enumerate}
\end{lemma}
\begin{proof}
(K1) The measure $\theta_{\ge l+1}$ is carried by $\pfrG^{(l+1)}$. For $x\in\pfrG^{(l+1)}$, $x\nu x^{-1}$ lies in
$\pfrU_{(<l)}\subseteq\pfrU$ and has the same entries as $\nu$ just above the diagonal of the frame block: these are
the entries $(a,a+1)$ with $a<l$, so $a+1\le l$, that is $h(\nu f_{a+1},f^\vee_a)$ for the dual
vector $f_a^\vee$ of $f_a$, and $x$ fixes $f_{a+1}$ and $f_a^\vee$ (both of floor $\le l<l+1$). So $\sigma(x\nu x^{-1})=\sigma(\nu)$, and
$W(x\nu g)=W((x\nu x^{-1})xg)=\psi(M^{-1}\sigma(\nu))W(xg)$. Likewise $d'x\nu'x^{-1}d'^{-1}\in\pfrU$ ($d'$ is diagonal
on the frame, so it normalises $\pfrU$ and multiplies each entry above the diagonal by a scalar that does not
depend on $x$), $\sigma(d'x\nu'x^{-1}d'^{-1})=\sigma(d'\nu'd'^{-1})$, and $W(d'x\nu'y)=\psi(M^{-1}\sigma(d'\nu'd'^{-1}))
W(d'xy)$. Insert both in $\mathsf k(x\nu g,x\nu'y)$, use convention (i) and integrate over $x$.

(K2) Descending induction on $l$. For $l>n$, $\pfrG^{(l)}=1$. For $l=n$, $\gamma\in\pfrH$ and
$\theta_{\ge n}=\theta_{\ge n+1}*\theta_n$ with $\theta_{\ge n+1}=\delta_1$ and $\theta_n=\mu_\pfrH$; so
$\mathsf k_n(\gamma h,\gamma h')=\int\mathsf k(x\gamma h,x\gamma h')\,d\mu_\pfrH(x)=\mathsf k_n(h,h')$ by the right
invariance of $\mu_\pfrH$. Let $j\le l\le\varrho$ and assume the first claim at $l+1$. \emph{Second claim at
$l$.} By (N1), $\zeta=\gamma'n_0$ with $\gamma'\in\pfrG^{(l+1)}$ and $n_0\in\pfrN_l$. By the first claim at $l+1$,
$\gamma'$ may be dropped. For $n_0$: by (C1), $\mathsf k_{l+1}(n_0h,n_0h')=\int\mathsf k_{l+2}(\sigma n_0h,\sigma
n_0h')\,d\theta_{l+1}(\sigma)$, and for $\sigma\in\pfrG^{(l+1)}$ the element $s=\sigma n_0\sigma^{-1}$ lies in
$\pfrN_l\subseteq\pfrU_{(<l+1)}$; by (K1) at level $l+2$ (note $l+1\le n$), $\mathsf k_{l+2}(s\sigma h,s\sigma h')=
\psi(M^{-1}(\sigma(d'sd'^{-1})-\sigma(s)))\mathsf k_{l+2}(\sigma h,\sigma h')$, and $d'sd'^{-1}=s$ as $s\in\pfrG^{(l)}
\subseteq\pfrG^{(j)}$. So the factor is $1$. \emph{First claim at $l$.} By (C1), $\mathsf k_l(\gamma h,\gamma h')=
\int\Psi(\sigma\gamma)\,d\theta_l(\sigma)$ with $\Psi(\sigma)=\mathsf k_{l+1}(\sigma h,\sigma h')$, which is integrable
for $\theta_l$ and, by the second claim, satisfies $\Psi(\zeta\sigma)=\Psi(\sigma)$ for $\zeta\in\pfrR^{(l+1)}\cap
\pfrG^{(l)}$. Lemma~\ref{pfr:floorright}(b) with $\xi=\gamma\in\pfrG^{(l)}$, applied to $\varphi\circ\Psi$ for every
Borel $\varphi\ge0$ on $\C$, shows that $\sigma\mapsto\Psi(\sigma\gamma)$ and $\Psi$ have the same law under
$\theta_l$; hence the same integral.

(K3) The second claim is (K1) with $\nu=1$ and $|\psi|=1$. For the first, write $\xi=\gamma\nu$ with
$\gamma\in\pfrG^{(l)}$, $\nu\in\pfrU_{(<l)}$ (convention (iii)); as $\nu\in\pfrR^{(l+1)}$, also $\gamma\in\pfrR^{(l+1)}
\cap\pfrG^{(l)}$. By (K2), $\mathsf k_{l+1}(\gamma\nu g,\gamma\nu gx)=\mathsf k_{l+1}(\nu g,\nu gx)$, and by (K1) its
modulus is $|\mathsf k_{l+1}(g,gx)|$. The weight $\pfrwt_l$ is invariant under $\pfrR^{(l)}\supseteq\pfrR^{(l+1)}$ by
(C3).
\end{proof}

\begin{remark}[Two points of \S\ref{pfr:sec} recalled]\label{pfs:rem-corr}
(a) Fact (C4) of \S\ref{pfr:ss-gainsproof} is stated in the form that holds and that the proofs in
\S\ref{pfr:ss-gainsproof} use, namely (K1)--(K3): invariance of $|\mathsf k_{l+1}|$ under $\pfrU_{(<l)}$, and of
$\mathsf a_l$ under $\pfrR^{(l+1)}$ (the functions $\mathsf b_l$ are then $\pfrR^{(l)}$-invariant by
Lemma~\ref{pfr:floorright}, by descending induction). (b) A kernel
$\mathsf k_{\varrho+2}=\mathsf k$ is non-zero only at points $g$ with $\pfrU g\cap\pfrK\ne\emptyset$; so in (C2) for
$l=\varrho+2$, and wherever \S\ref{pfr:sec} speaks of points ``near at $\varrho+2$'', the group $\pfrR^{(\varrho+2)}$ is
to be read as $\pfrU$ (this convention is part of the definition of nearness in \S\ref{pfr:ss-setting}). We use this only in Lemma~\ref{pfs:climb} below, where
we argue directly.
\end{remark}

\subsubsection*{Three standard facts}
\begin{lemma}[The Haar measure on the skew elements]\label{pfs:skewhaar}
Let $\tau_0\in A$ with $\tau_0+\bar\tau_0=1$ and let $\delta$ be a skew unit. The map
$\Sigma(\beta_1,\beta_2)=\tau_0\delta^{-1}\beta_1+\beta_2$ is an isomorphism of topological groups from
$A_{\rm sk}^2$ onto $A$ ($A_{\rm sk}=\{\beta:\beta+\bar\beta=0\}$), and there is exactly one Haar measure
$d\beta_{\rm sk}$ on $A_{\rm sk}$ such that $\Sigma$ carries $d\beta_{\rm sk}\otimes d\beta_{\rm sk}$ to $dc$.
\end{lemma}
\begin{proof}
For skew $\beta_1$, $r=\delta^{-1}\beta_1$ is hermitian ($\bar r=(-\delta^{-1})(-\beta_1)$). If
$c=\tau_0r+\beta_2$, then $c+\bar c=(\tau_0+\bar\tau_0)r=r$. Hence $\Sigma$ has the continuous inverse
$c\mapsto(\delta(c+\bar c),\,c-\tau_0(c+\bar c))$ (the second entry is skew: its sum with its conjugate is
$(c+\bar c)(1-\tau_0-\bar\tau_0)=0$). Let $\nu_0$ be a Haar measure of $A_{\rm sk}$. The image of $\nu_0\otimes\nu_0$
under the isomorphism $\Sigma$ is a Haar measure of $A$, so it is $\lambda\,dc$ with $\lambda>0$; then
$\nu=\lambda^{-1/2}\nu_0$ is the required measure, and it is unique since any other is $c\nu$ with $c^2=1$,
$c>0$.
\end{proof}
At the real place ($\tau_0=\frac12$, $\delta=\frac i2$, as fixed in \S\ref{pfr:ss-setting}), $\Sigma(it_1,it_2)=t_1+it_2$, so
$d\beta_{\rm sk}$ is the Lebesgue measure $dt$ of $\beta=it$. At the finite place $d\beta_{\rm sk}$ gives the integral skew elements
volume $1$ (Lemma~\ref{pfg:duality}(d)).

\begin{lemma}[Orbits and measurable sections]\label{pfs:section}
Let $b\le\varrho$, let $\pfsst_b=\{\sigma\in\pfrG^{(b)}:\sigma f_b=f_b\}$ and $\pfsO_b=\{\sigma^{-1}f_b:\sigma\in\pfrG^{(b)}\}
\subseteq V_b$. Then $\pfsst_b=\pfrR^{(b+1)}\cap\pfrG^{(b)}$; $\pfsO_b$ is a Borel set; and there is a Borel map
$s\colon\pfsO_b\to\pfrG^{(b)}$ with $s(v)^{-1}f_b=v$. Consequently, if $\Psi$ is a Borel function on $\pfrG^{(b)}$
(with values in $\C$ or $[0,\infty]$) with $\Psi(\zeta\sigma)=\Psi(\sigma)$ for $\zeta\in\pfsst_b$, then
$\Psi(\sigma)=\Psi_0(\sigma^{-1}f_b)$ for the Borel function $\Psi_0=\Psi\circ s$ on $\pfsO_b$, extended by $0$
to $V_b$. Moreover the quadric measure $\mu_b$ is carried by $\pfsO_b$.
\end{lemma}
\begin{proof}
The equality $\pfsst_b=\pfrR^{(b+1)}\cap\pfrG^{(b)}$ is shown in (N1) of \S\ref{pfr:ss-floors} (one inclusion is
proved there; conversely an element of $\pfrG^{(b)}$ fixing $f_b$ fixes all $f_c$, $c\le b$, hence lies in
$\pfrR^{(b+1)}$). The group $\pfrG^{(b)}$ is a closed subgroup of $\mathrm{GL}_r(A)$, a Polish group ($A$ is $\C$ or
a finite product of local fields), and $\pfsst_b$ is a closed subgroup. By the theorem on Borel transversals of
cosets of a closed subgroup of a Polish group (Kechris, \emph{Classical descriptive set theory}, Theorem~12.17,
applied to the left cosets of $\pfsst_b$ after the inversion $\sigma\mapsto\sigma^{-1}$) there is a Borel set $T\subseteq\pfrG^{(b)}$ meeting each coset $\pfsst_b\sigma$ in exactly one
point. The continuous map $o(\sigma)=\sigma^{-1}f_b$ is constant on these cosets ($o(\zeta\sigma)=\sigma^{-1}\zeta^{-1}
f_b=o(\sigma)$) and separates them ($o(\sigma)=o(\sigma')$ gives $\sigma'\sigma^{-1}\in\pfsst_b$); so $o|_T$ is a
continuous injection with image $\pfsO_b$. By the Lusin--Souslin theorem (Kechris, op.\ cit., Theorem~15.1) the
image $\pfsO_b$ is Borel and $s=(o|_T)^{-1}$ is Borel. For $\Psi$ as stated, $s(\sigma^{-1}f_b)\in\pfsst_b\sigma$,
so $\Psi(s(\sigma^{-1}f_b))=\Psi(\sigma)$. Finally, the image of $\theta_b$ under $o$ is $c_b\mu_b$ with $c_b>0$
(definition of a place datum, \S\ref{pfr:ss-setting}), and $\theta_b$ is carried by $\pfrG^{(b)}$.
\end{proof}

\begin{remark}[Witt's theorem]\label{pfs:witt}
The proof of Lemma~\ref{pfr:floorright} (and of Lemma~\ref{pfr:upperinv}, which uses it) invokes the transitivity of
$\pfrG^{(b)}$ on the quadric of $f_b$ in order to write a $\pfsst_b$-invariant function as a function of $\sigma^{-1}f_b$.
Transitivity is not needed: Lemma~\ref{pfs:section} gives $\Psi=\Psi_0(\sigma^{-1}f_b)$ on $\pfrG^{(b)}$ with
$\Psi_0$ Borel, and the integrals in that proof are integrals against $\mu_b$, which is carried by the orbit
$\pfsO_b$; the remaining step, the $\pfrG^{(b)}$-invariance of $\mu_b$, is used there as stated. For the record:
at the real place $\pfsO_b$ is the whole quadric minus $\{0\}$, by Witt's extension theorem for hermitian forms
over $\C$ (e.g.\ Scharlau, \emph{Quadratic and Hermitian Forms}), applied to
the map $f_b\mapsto v$ between the lines $Af_b$ and $Av$ for $v\ne0$ with $h(v,v)=h(f_b,f_b)$.
\end{remark}

\subsection{Plancherel's inequality and the quadrics}\label{pfs:ss-planch}

\subsubsection*{The constant of a place}
The constant $P$ of a place (\S\ref{pfr:ss-setting}) is the constant of the following axiom, which
we add to the definition of a place, together with two properties of the modulus that every place of
\S\ref{pfr:ss-setting} has.
\begin{enumerate}
\item[(P)] There is a continuous integrable function $G\colon A\to[0,\infty)$ with $G(0)=1$ whose
Fourier transform $\widehat G(z)=\int_AG(v)\psi(\bar vz)\,dv$ satisfies $\int_A|\widehat G(z)|\,dz\le P$.
Moreover $\pfrmod(\bar c)=\pfrmod(c)$, $\pfrmod$ is multiplicative, $d(cv)=\pfrmod(c)\,dv$ for units $c$, and
the non-units of $A$ form a null set.
\end{enumerate}
For $d\ge1$ and $v,w\in A^d$ write $\pfrip vw=\sum_i\bar v_iw_i$.

\begin{proposition}[Plancherel's inequality]\label{pfs:planch}
Let $\eta$ be a unit and $f\in L^1(A^d)$. Then
\[
\int_{A^d}\Bigl|\int_{A^d}\psi(\eta\pfrip vw)f(w)\,dw\Bigr|^2dv\le\Bigl(\frac P{\pfrmod(\eta)}\Bigr)^d\int_{A^d}|f|^2 .
\]
\end{proposition}
\begin{proof}
We may assume $f\in L^2$. Write $\widehat f(v)=\int\psi(\eta\pfrip vw)f(w)\,dw$ and, for $x\in A$,
$I(x)=\int|\widehat f(v)|^2\prod_iG(\bar xv_i)\,dv$. \emph{Expansion.} The function $(v,w,w')\mapsto\prod_iG(\bar x
v_i)|f(w)||f(w')|$ is integrable, so by Fubini's theorem, with $z=w'-w$,
\[
I(x)=\iint\overline{f(w)}f(w')\prod_i\widehat G_x(z_i)\,dw\,dw',\qquad\widehat G_x(t)=\int_AG(\bar xv)\psi(\eta\bar vt)\,dv .
\]
For a unit $x$ substitute $v=\bar x^{-1}v'$: $dv=\pfrmod(x)^{-1}dv'$ and $\eta\bar vt=\eta x^{-1}\bar v't$, so
$\widehat G_x(t)=\pfrmod(x)^{-1}\widehat G(\eta x^{-1}t)$ and $\int|\widehat G_x|\le\pfrmod(x)^{-1}\pfrmod(\eta
x^{-1})^{-1}P=P/\pfrmod(\eta)$. By Cauchy--Schwarz, $\int|f(w)||f(w+z)|\,dw\le\int|f|^2$ for every $z$; hence,
substituting $w'=w+z$,
\[
I(x)\le\int\prod_i|\widehat G_x(z_i)|\int|f(w)||f(w+z)|\,dw\,dz\le\Bigl(\frac P{\pfrmod(\eta)}\Bigr)^d\int|f|^2 .
\]
\emph{Limit.} The units are conull, hence dense, so there are units $x_m\to0$. As $G$ is continuous with
$G(0)=1$, $\prod_iG(\bar x_mv_i)\to1$ for every $v$, and Fatou's lemma gives $\int|\widehat f|^2\le\liminf
I(x_m)$.
\end{proof}

\subsubsection*{The Fourier kernel between two quadrics}
Fix $b\le\varrho$ and use the floor coordinates of $V_b$ (\S\ref{pfr:ss-section}): in them $h$ is the standard
form, a vector of size $i$ has floor coordinates in $\Gamma(i+4)$, and for $y,v\in V_b$ the product $h(y,v)$ is
$\sum_{c}(y_{c}\bar v_{c^*}+y_{c^*}\bar v_c)+\sum_{t}y_t\bar v_t$, the first sum over the hyperbolic pairs of $V_b$,
the second over the orthonormal coordinates. Fix two disjoint blocks $\beta,\beta'$ of $2m+2$ orthonormal
coordinates each, disjoint from those of $X'$ (by the definition of $X'$ in \S\ref{pfr:ss-gauge}). Recall that
$V_b$ has at most $99n$ floor coordinates off $X'$ (proof of Lemma~\ref{pfr:quadricX}).

\begin{proposition}[Plancherel on a quadric]\label{pfs:quadplanch}
Let $\rho$ be Lebesgue measure on $V_b$ or the quadric measure $\mu_b$, let $i,j\ge0$ and let $\xi$ be a unit.
For every Borel $\varphi\in L^1(\mu_b)$ vanishing off the vectors of size $j$,
\[
\int_{\{v\text{ of size }i\}}\Bigl|\int\psi(\xi h(y,v))\varphi(y)\,d\mu_b(y)\Bigr|^2d\rho(v)\le
\Lambda^{99n(i+6)}\Lambda^{99n(j+6)}\Bigl(\frac P{\pfrmod(\xi)}\Bigr)^{\pfsrk}\int|\varphi|^2\,d\mu_b .
\]
\end{proposition}
\begin{proof}
By Lemma~\ref{pfr:block}, applied to $\mu_b$ with the block $\beta$ and to $\rho$ with the block $\beta'$, write
$y=(y_o,y_X,y_\beta)$ and $v=(v_o,v_X,v_{\beta'})$, where the subscript $X$ denotes the coordinates of $X'$ and
$o$ the remaining (outer) ones; then $d\mu_b(y)=dy_o\,dy_X\,\kappa_{(y_o,y_X)}(dy_\beta)$ and
$d\rho(v)=dv_o\,dv_X\,\kappa'_{(v_o,v_X)}(dv_{\beta'})$ with Lebesgue measures in the outer and $X'$ coordinates,
and fibre measures giving mass at most $\Lambda^{(2m+2)(l+2)}$ to the vectors with entries in $\Gamma(l)$. Let
$R=\Gamma(j+4)$, $R'=\Gamma(i+4)$, and let $N_o$, $N_o'$ be the numbers of outer coordinates of $y$ and of $v$.

\emph{The phase.} The coordinates of $\beta'$ are outer coordinates of $y$ (as $\beta'$ is disjoint from $\beta$ and
from $X'$). Hence $h(y,v)=\pfrip{v_X}{y_X}+B_v(y_o)+a_{v_o}(y)$, where $\pfrip{v_X}{y_X}=\sum_{t\in X'}\bar v_ty_t$ is
the pairing of Proposition~\ref{pfs:planch}, $B_v(y_o)=\sum_{t\in\beta'}y_t\bar v_t$
depends on $y$ only through $y_o$, and $a_{v_o}(y)$ collects the terms containing a coordinate of $v_o$. So
$\psi(\xi h(y,v))=\psi(\xi B_v(y_o))\,\psi(\xi\pfrip{v_X}{y_X})\,\psi(\xi a_{v_o}(y))$.

\emph{Cauchy--Schwarz in $y_o$.} Put $F_{v_o}(y_o,y_X)=\int\psi(\xi a_{v_o}(y))\varphi(y)\,\kappa_{(y_o,y_X)}(dy_\beta)$.
As $\varphi$ vanishes unless $y_o\in R^{N_o}$, and $|\psi(\xi B_v)|=1$,
\[
T(v):=\Bigl|\int\psi(\xi h(y,v))\varphi\,d\mu_b\Bigr|^2\le\pbvol(R)^{N_o}\,H(v_o,v_X),\quad
H(v_o,v_X)=\int\Bigl|\int\psi(\xi\pfrip{v_X}{y_X})F_{v_o}(y_o,y_X)\,dy_X\Bigr|^2dy_o .
\]
\emph{Plancherel in $y_X$, Cauchy--Schwarz in the block.} By Proposition~\ref{pfs:planch} for each $y_o$ (the
function $y_X\mapsto F_{v_o}(y_o,y_X)$ is integrable for almost every $y_o$), and then
$|F_{v_o}|^2\le\kappa_{(y_o,y_X)}(R^{2m+2})\int|\varphi|^2d\kappa\le\Lambda^{(2m+2)(j+6)}\int|\varphi|^2d\kappa$ (as
$\varphi$ vanishes unless $y_\beta\in R^{2m+2}$),
\[
\int_{A^{X'}}H(v_o,v_X)\,dv_X\le\Bigl(\frac P{\pfrmod(\xi)}\Bigr)^{\pfsrk}\Lambda^{(2m+2)(j+6)}\int|\varphi|^2\,d\mu_b
\quad\text{for every }v_o .
\]
\emph{Integration in $v$.} A vector $v$ of size $i$ has floor coordinates in $R'$. As $T(v)\le\pbvol(R)^{N_o}H(v_o,v_X)$
does not depend on $v_{\beta'}$, integrating first over $v_{\beta'}$ (mass $\le\Lambda^{(2m+2)(i+6)}$), then over
$v_X\in A^{X'}$, then over $v_o\in R'^{N'_o}$ gives the bound
\[
\pbvol(R')^{N'_o}\Lambda^{(2m+2)(i+6)}\cdot\pbvol(R)^{N_o}\Lambda^{(2m+2)(j+6)}\Bigl(\frac P{\pfrmod(\xi)}\Bigr)^{\pfsrk}\int|\varphi|^2d\mu_b .
\]
By (G3), $\pbvol(R)\le\Lambda^{j+5}$, $\pbvol(R')\le\Lambda^{i+5}$, and $N_o+2m+2$, $N_o'+2m+2$ are at most the number
$99n$ of coordinates off $X'$; as $\Lambda\ge1$, $(i+5)N_o'+(i+6)(2m+2)\le99n(i+6)$ and likewise for $j$, so the
constant is at most $\Lambda^{99n(i+6)}\Lambda^{99n(j+6)}$.
\end{proof}

\subsubsection*{The Fourier transform of a correlation along a quadric}
\begin{lemma}[Near points, general size]\label{pfs:nearsize}
Let $b\le\varrho$, let $h'\in\pfrG$ have size $i$, and let $\rho\in\pfrG^{(b)}$ with $\rho h'$ near at $b+1$. Then
$\rho^{-1}f_b$ has size $i+6$.
\end{lemma}
\begin{proof}
As in the proof of Lemma~\ref{pfr:nearsize}: let $\xi\in\pfrR^{(b+1)}$ with $k=\xi\rho h'\in\pfrK$; then
$\rho^{-1}\xi^{-1}=h'k^{-1}$ has size $i+5+1$ ($k^{-1}$ has size $5$ by (G0), (G1)); $\xi^{-1}f_b=f_b+\sum_{c<b}\lambda_cf_c$
and $\rho^{-1}$ fixes the $f_c$, $c<b$, and preserves $V_b$; so the coordinates of $\rho^{-1}f_b$ are those of
floor $\ge b$ of the column $h'k^{-1}f_b$, and vanish on the floors $<b$.
\end{proof}

\begin{proposition}[Plancherel for a correlation]\label{pfs:corrplanch}
Let $b\le\varrho$, let $\mathsf k$ be a kernel at floor $j_0\le b$, let $t\in\pfrG^{(b)}$ and $q'\in\pfrG$ with
$t^{-1}q'\in\pfrU_{(<b)}$, let $g,y\in\pfrG$ with $q'g$ of size $i$, let $\xi$ be a unit, $v_0\in V_b$, $i'\ge0$, and
let $\rho$ be Lebesgue measure on $V_b$ or $\mu_b$. Put
$\mathcal F(v)=\int\psi(\xi h(t\sigma^{-1}f_b,v+v_0))\,\mathsf k_{b+1}(\sigma g,\sigma y)\,d\theta_b(\sigma)$. Then
\[
\int_{\{v\text{ of size }i'\}}|\mathcal F(v)|^2\,d\rho(v)\le2\Lambda^{99n(i'+6)}\Lambda^{99n(i+12)}\Bigl(\frac
P{\pfrmod(\xi)}\Bigr)^{\pfsrk}\int|\mathsf k_{b+1}(\sigma g,\sigma y)|^2\,d\theta_b(\sigma).
\]
\end{proposition}
\begin{proof}
Let $\kappa(\sigma)=\mathsf k_{b+1}(\sigma g,\sigma y)$; it is bounded, Borel, and integrable for $\theta_b$ by
(C1). \emph{Support.} Let $\sigma\in\pfrG^{(b)}$ with $\kappa(\sigma)\ne0$. By (C2), $\sigma g$ is near at $b+1$. The
element $\sigma(t^{-1}q')\sigma^{-1}$ lies in $\pfrU_{(<b)}\subseteq\pfrR^{(b+1)}$, so
$(\sigma t^{-1})(q'g)=\sigma(t^{-1}q')\sigma^{-1}\cdot\sigma g$ is near at $b+1$, and by Lemma~\ref{pfs:nearsize} (with
$\rho=\sigma t^{-1}$, $h'=q'g$), $t\sigma^{-1}f_b$ has size $i+6$. \emph{Factorisation.} By (K2), $\kappa(\zeta\sigma)=
\kappa(\sigma)$ for $\zeta\in\pfrR^{(b+1)}\cap\pfrG^{(b)}=\pfsst_b$, so by Lemma~\ref{pfs:section}, $\kappa(\sigma)=\kappa_0(
\sigma^{-1}f_b)$ on $\pfrG^{(b)}$ with $\kappa_0$ Borel. Put $\varphi(w)=\mathbf 1[w\text{ of size }i+6]\,\kappa_0(t^{-1}w)\,
\psi(\xi h(w,v_0))$. The map $\sigma\mapsto t\sigma^{-1}f_b=(\sigma t^{-1})^{-1}f_b$ carries $\theta_b$ to $c_b\mu_b$ (the
image of $\theta_b$ under $\sigma\mapsto\sigma^{-1}f_b$ is $c_b\mu_b$ with $0<c_b\le2$, \S\ref{pfr:ss-setting} and
Lemma~\ref{pfr:floorquad}, and $\mu_b$ is $\pfrG^{(b)}$-invariant). Hence, by the
support property, for $\theta_b$-almost every $\sigma$, $\kappa(\sigma)\psi(\xi h(t\sigma^{-1}f_b,v_0))=\varphi(t\sigma^{-1}
f_b)$, so $\mathcal F(v)=c_b\int\psi(\xi h(w,v))\varphi(w)\,d\mu_b(w)$ (additivity of $h$ in the second variable), and
$c_b\int|\varphi|^2d\mu_b=\int|\kappa|^2d\theta_b$, $\varphi\in L^1(\mu_b)$. Proposition~\ref{pfs:quadplanch} with $j=i+6$
bounds the left side by $c_b^2\Lambda^{99n(i'+6)}\Lambda^{99n(i+12)}(P/\pfrmod(\xi))^{\pfsrk}\int|\varphi|^2d\mu_b$, and
$c_b^2\int|\varphi|^2d\mu_b=c_b\int|\kappa|^2d\theta_b\le2\int|\kappa|^2d\theta_b$.
\end{proof}

\subsubsection*{Correlations along a Heisenberg group}
\begin{lemma}\label{pfs:heiscorr}
Let $a<\varrho$, let $\mathsf k$ be a kernel at floor $a+1$, $\varepsilon\in\{0,1\}$ and $\xi\in A$ with
\begin{equation}\label{pfs:eq-twist}
M^{-1}\bigl(\sigma(d'sd'^{-1})-\varepsilon\,\sigma(s)\bigr)=\xi\,\sigma(s)\qquad(s\in\pfrN_a).
\end{equation}
Then for $n=n(\omega,\beta)\in\pfrN_a$ and all $g,y$,
\[
\mathsf k_{a+1}(n^\varepsilon g,ny)=\int\psi\bigl(-\xi\,h(\sigma^{-1}f_{a+1},\omega)\bigr)\,\mathsf k_{a+2}(\sigma g,\sigma y)\,d\theta_{a+1}(\sigma).
\]
\end{lemma}
\begin{proof}
By (C1), $\mathsf k_{a+1}(n^\varepsilon g,ny)=\int\mathsf k_{a+2}(\sigma n^\varepsilon g,\sigma ny)\,d\theta_{a+1}(\sigma)$. For
$\sigma\in\pfrG^{(a+1)}$ put $s=\sigma n\sigma^{-1}\in\pfrN_a\subseteq\pfrU_{(<a+1)}$; then $\sigma n^\varepsilon g=s^\varepsilon\sigma g$
and $\sigma ny=s\sigma y$. By (K1) at level $a+2$ (as $a+1\le n$),
$\mathsf k_{a+2}(s^\varepsilon\sigma g,s\sigma y)=\psi(M^{-1}(\sigma(d'sd'^{-1})-\sigma(s^\varepsilon)))\mathsf k_{a+2}(\sigma g,\sigma y)$, and
$\sigma(s^\varepsilon)=\varepsilon\sigma(s)$. By \eqref{pfs:eq-twist} and \eqref{pfs:eq-sigma} the phase is
$\psi(\xi\sigma(s))=\psi(-\xi h(\sigma^{-1}f_{a+1},\omega))$.
\end{proof}

\begin{proposition}[Parseval along a Heisenberg group]\label{pfs:parseval}
Let $a<\varrho$, $\mathsf k$ a kernel at floor $a+1$, $\varepsilon\in\{0,1\}$, $\xi$ and $u$ units satisfying
\eqref{pfs:eq-twist}, $C\ge0$, $i,j\ge0$, $g,y\in\pfrG$, $S\subseteq\pfrN_a$, and $\Phi\ge0$ Borel on $\pfrG$ such
that:
\begin{enumerate}
\item[(a)] if $n\in\pfrN_a$ and $\Phi(n)\ne0$, then $n\in S$ and $\Phi(n)\le C|\mathsf k_{a+1}(n^\varepsilon g,ny)|^2$;
\item[(b)] for every $n_0\in S$ there are $t_0\in\pfrR^{(a+1)}\cap\pfrG^{(a)}$ and $q\in\pfrU_{(<a+1)}$ such that
$t_0qg$ has size $i$ and, for every $n\in S$, some vector $v$ congruent to $(t_0n_0(t_0n)^{-1}-1)e_a^*$ modulo
$\langle e_c:c<a\rangle$ has $u^{-1}v$ of size $j$.
\end{enumerate}
Then
\[
\int\Phi\,d\eta_a\le C\Bigl(\frac P{\pfrmod(\xi u)}\Bigr)^{\pfsrk}\pfsc_a(u)\,\Lambda^{100n(i+j+31)}
\int|\mathsf k_{a+2}(\sigma g,\sigma y)|^2\,d\theta_{a+1}(\sigma),
\]
with $\pfsc_a(u)=\pfrmod(u)^{d_a+1/2}$ as in \S\ref{pfr:ss-gauge}.
\end{proposition}
\begin{proof}
If $\int\Phi\,d\eta_a=0$ there is nothing to prove; otherwise, as $\eta_a$ is carried by $\pfrN_a$, there is
$n_0\in\pfrN_a$ with $\Phi(n_0)\ne0$, so $n_0\in S$. Take $t_0,q$ as in (b) and write $t_0=tr$ with
$t\in\pfrG^{(a+1)}$, $r\in\pfrN_a$ ((N1)). Let $z$ be the vector part of $n_0$, and for $n'=n(\omega',\beta')\in\pfrN_a$ put
\[
N(n')=r^{-1}\,(t^{-1}n't)\,r\,n_0\in\pfrN_a ,
\]
whose vector part is $t^{-1}\omega'+z$ (convention (iv)). \emph{Substitution.} The measure $\eta_a$ is invariant
under left and right translation by elements of $\pfrN_a$ and under conjugation by $t\in\pfrG^{(a+1)}$ ((N2)); so
$\int\Phi\,d\eta_a=\int\Phi(N(n'))\,d\eta_a(n')$.

\emph{The correlation at $N(n')$.} By Lemma~\ref{pfs:heiscorr} and $h(\sigma^{-1}f_{a+1},t^{-1}\omega'+z)=h(t\sigma^{-1}f_{a+1},
\omega'+tz)$,
\[
\mathsf k_{a+1}(N(n')^\varepsilon g,N(n')y)=\int\psi\bigl(\xi'\,h(t\sigma^{-1}f_{a+1},u^{-1}\omega'+v_0)\bigr)\,
\mathsf k_{a+2}(\sigma g,\sigma y)\,d\theta_{a+1}(\sigma)=:\mathcal F(u^{-1}\omega'),
\]
with $\xi'=-\xi\bar u$ and $v_0=u^{-1}tz$ (as $h(w,\omega')=\bar u\,h(w,u^{-1}\omega')$ by convention (i)).

\emph{The lift.} If $\Phi(N(n'))\ne0$, then $N(n')\in S$, and (b) with $n=N(n')$ gives a vector $v\equiv(t_0n_0(t_0
N(n'))^{-1}-1)e_a^*$ with $u^{-1}v$ of size $j$; and $t_0n_0(t_0N(n'))^{-1}=trn_0\,n_0^{-1}r^{-1}t^{-1}n'^{-1}tr\,r^{-1}t^{-1}=
n'^{-1}$. Also $\Phi(N(n'))\le C|\mathcal F(u^{-1}\omega')|^2$ by (a).

\emph{Integration.} Apply Lemma~\ref{pfr:heis} to $n'\mapsto\Phi(N(n'))$ with the unit $u$, the size $j$, the
constant $C$ and $\Psi=|\mathcal F|^2$ on $V_{a+1}$; the integral $I$ of $\Psi$ over the vectors of size $j$
is bounded by Proposition~\ref{pfs:corrplanch} (with $b=a+1$, the kernel $\mathsf k$ at floor $a+1$, $t$, $q'=t_0q$, so that
$t^{-1}q'=rq\in\pfrU_{(<a+1)}$ and $q'g=t_0qg$ has size $i$, the unit $\xi'$, $v_0$, $i'=j$ and Lebesgue measure).
As $\pfrmod(\xi')=\pfrmod(\xi u)$, the result is
\[
\int\Phi\,d\eta_a\le 2C\,\pfsc_a(u)\Lambda^{j+2}\Lambda^{99n(j+6)}\Lambda^{99n(i+12)}\Bigl(\frac P{\pfrmod(\xi u)}\Bigr)^{\pfsrk}
\int|\mathsf k_{a+2}(\sigma g,\sigma y)|^2\,d\theta_{a+1}(\sigma),
\]
and, as $2\le\Lambda$, $j+3+99n(j+6)+99n(i+12)\le100n(i+j+31)$.
\end{proof}

\subsection{Depth along the Heisenberg groups}\label{pfs:ss-depth}

Let $\pfsBall=\pfsBall_\Gamma$ be the position ball of $\Gamma$ (the $x\in\pfrP$ with $M^{-1}(x-1)$ of size $1$),
and for $a\le\varrho$ and $g\in\pfrG$ let
\[
\pfsCB_a(g)=\{z\in\pfrG:\ g^{-1}\nu z\in\pfsBall\text{ for some }\nu\in\pfrU_{(<a)}\}.
\]

\begin{lemma}[Two points of the position ball]\label{pfs:twopos}
If $g\in\pfrK$ and $x,y\in\pfsBall$, then $M^{-1}(gx(gy)^{-1}-1)$ has size $14$.
\end{lemma}
\begin{proof}
$y=M\cdot M^{-1}(y-1)+1$ has size $3$ ((G1): $M\in\Gamma(1)$), so $gy$ has size $5$ and $(gy)^{-1}$ size $9$. The
matrix $g(M^{-1}(x-1)-M^{-1}(y-1))$ has size $1+2+1=4$, and $M^{-1}(gx(gy)^{-1}-1)=g\,M^{-1}(x-y)\,(gy)^{-1}$ has size
$4+9+1$.
\end{proof}

\begin{lemma}[The lift of a Heisenberg element]\label{pfs:lift}
Let $a<\varrho$, $g\in\pfrK$, $y\in\pfrG$ with $g^{-1}y\in\pfsBall$, and $n\in\pfrN_a$ with $ny\in\pfsCB_a(g)$. Then
some vector $v$ congruent to $(n^{-1}-1)e_a^*$ modulo $\langle e_c:c<a\rangle$ has $M^{-1}v$ of size $17$.
\end{lemma}
\begin{proof}
Let $\nu\in\pfrU_{(<a)}$ with $g^{-1}\nu ny\in\pfsBall$, and $\mu=n^{-1}\nu^{-1}n\in\pfrU_{(<a)}$ (as $n\in\pfrG^{(a)}$
normalises $\pfrU_{(<a)}$). Then $g(g^{-1}y)\bigl(g(g^{-1}\nu ny)\bigr)^{-1}=n^{-1}\nu^{-1}=\mu n^{-1}$, so by
Lemma~\ref{pfs:twopos} the matrix $M^{-1}(\mu n^{-1}-1)$ has size $14$, and so has its column $M^{-1}v$ at $e_a^*$, where
$v=(\mu n^{-1}-1)e_a^*$. As $n^{-1}e_a^*\in V_a$ and $\mu$ acts on
$V_a$ by adding combinations of the $e_c$, $c<a$ (convention (iii)), $v\equiv(n^{-1}-1)e_a^*$. Size $14$ implies
size $17$.
\end{proof}

\begin{lemma}[A Heisenberg identity]\label{pfs:heisid}
Let $a<\varrho$, $n=n(\omega,\beta)\in\pfrN_a$, let $\zeta\equiv\zeta'$ modulo $\langle e_c:c<a\rangle$, and let $o\in X'$.
Then $h(\zeta,o)-h(n^{-1}\zeta',o)=h(\zeta,e_a)\,h(\omega,o)$.
\end{lemma}
\begin{proof}
$o\in V_{a+1}$ is orthogonal to $f_a,f_a^*$, so $no=o-h(o,\omega)e_a$ and $h(n^{-1}\zeta',o)=h(\zeta',no)=h(\zeta',o)-
\overline{h(o,\omega)}\,h(\zeta',e_a)$. Here $h(\zeta',o)=h(\zeta,o)$ and $h(\zeta',e_a)=h(\zeta,e_a)$, as $\zeta-\zeta'$ is a
combination of the $e_c$, $c<a$, which are orthogonal to $o$ and to $e_a$; and $\overline{h(o,\omega)}=h(\omega,o)$.
\end{proof}

\begin{lemma}[The inverse depth on a Heisenberg group]\label{pfs:invdepthN}
Let $a<\varrho$, $g\in\pfrK$ and $y\in\pfrG$ with $g^{-1}y\in\pfsBall$. Then
\[
\int\mathbf 1[ny\in\pfsCB_a(g)]\;\pfrdep_a(g,g^{-1}ny)^{-s}\,d\eta_a(n)\le\pfsc_a(M)\,\Lambda^{4600n}.
\]
\end{lemma}
\begin{proof}
\emph{The inverse depth is bounded by a common-vector function.} Let $\phi$ be the function of (G5) for the point
$g\in\pfrK$, the position $g^{-1}y\in\pfrP$ and $v=0$, with set of common vectors $\mathcal W$ (the $o\in X'$ with
$(g^{-1}o,y^{-1}o)$ close) and set of differences $\mathcal D$. Let $n=n(\omega,\beta)$ and $u=(M^{-1}\omega)_{X'}$. We claim
$\pfrdep_a(g,g^{-1}ny)^{-1}=\pbvol\mathfrak d_a(g,g^{-1}ny)\le\phi(u)$. Let $c\in\mathfrak d_a(g,g^{-1}ny)$: $c=h(\zeta,e_a)$
(the $e_a^*$-coordinate) for a vector $\zeta$ stable at $(g,ny)$ for $a$, with $(g^{-1}\zeta,(ny)^{-1}\zeta')$ close for some
$\zeta'\equiv\zeta$ modulo $\langle e_c:c<a\rangle$. For $o\in\mathcal W$ the pair $(g^{-1}o,y^{-1}o)$ is close, so
\[
\mathcal D\ni M^{-1}\bigl(h(g^{-1}\zeta,g^{-1}o)-h((ny)^{-1}\zeta',y^{-1}o)\bigr)=M^{-1}\bigl(h(\zeta,o)-h(n^{-1}\zeta',o)\bigr)=
c\,h(M^{-1}\omega,o)=h(cu,o),
\]
by Lemma~\ref{pfs:heisid}, $\bar M=M$, and $h(\omega',o)=h(\omega'_{X'},o)$ for $o\in X'$ (here $\omega'_{X'}$ is read as the
vector $\omega'_{X'}=\sum_th(\omega',x_t)x_t$ of $X'$, \S\ref{pfr:ss-gauge}). As $\mathcal D$ is stable under conjugation
(\S\ref{pfr:ss-gauge}) and $\overline{h(cu,o)}=h(o,cu)$, also $h(o,cu)\in\mathcal D$ for every $o\in\mathcal W$. So $c$ lies in the
set $\{c:h(o,cu)\in\mathcal D\ \forall o\in\mathcal W\}$ of (G5), whose volume is at most $\phi(u)$.

\emph{Integration.} If the integrand is non-zero at $n$, Lemma~\ref{pfs:lift} gives a vector $v\equiv(n^{-1}-1)e_a^*$
with $M^{-1}v$ of size $17$, and the integrand is at most $\Psi(M^{-1}\omega)$ with $\Psi(w)=\phi(w_{X'})^s$ on $V_{a+1}$.
By the general form of Lemma~\ref{pfr:quadricX} (proved there for Lebesgue measure) and (G5) at size $21$,
$\int_{\{w\text{ of size }17\}}\Psi\le\Lambda^{99n\cdot23}\Lambda^{99n\cdot22}$. Lemma~\ref{pfr:heis} with $u=M$, $j=17$, $C=1$ bounds the
integral by $\pfsc_a(M)\Lambda^{19}\Lambda^{99n\cdot45}\le\pfsc_a(M)\Lambda^{4600n}$.
\end{proof}

\begin{lemma}[The depth weight over the unipotent group]\label{pfs:depthweight}
Let $l\le\varrho$, and let $g\in\pfrG$ be near at $l$. Then for every $y\in\pfrG$,
\[
\int\mathbf 1[g^{-1}\nu y\in\pfsBall]\;\pfrwt_l(g,g^{-1}\nu y)\,d\pfsA_l(\nu)\le\prod_{a<l}\bigl(\pfsc_a(M)\Lambda^{4600n}\bigr).
\]
\end{lemma}
\begin{proof}
\emph{Reduction to $g\in\pfrK$.} Let $\xi g\in\pfrK$ with $\xi\in\pfrR^{(l)}$, $\xi=\gamma\nu_0$ with $\gamma\in\pfrG^{(l)}$,
$\nu_0\in\pfrU_{(<l)}$. Substitute $\nu\mapsto\nu_0^{-1}\gamma^{-1}\nu\gamma$, which preserves $\pfsA_l$ (convention (ii)). Then
$g^{-1}\nu_0^{-1}\gamma^{-1}\nu\gamma y=(\xi g)^{-1}\nu(\gamma y)$, and for every $x$, $\pfrwt_l(g,x)=\pfrwt_l(\xi g,x)$ by (C3)
(the second point being $gx$, resp.\ $\xi gx$). So the integral equals the one for $(\xi g,\gamma y)$; we may assume
$g\in\pfrK$.

\emph{Induction on the length.} Let $g\in\pfrK$ and $J_a(y)$ the integral with $a$ in place of $l$, $a\le l$. We show
$J_a(y)\le C_a:=\prod_{b<a}\bigl(\pfsc_b(M)\Lambda^{4600n}\bigr)$ for all $y$. For $a=0$, $\pfsA_0=\delta_1$ and $\pfrwt_0=1$. Let $a<l$ and
suppose $J_a\le C_a$ everywhere. If $J_{a+1}(y)\ne0$, there is $\tau\in\pfrU_{(<a+1)}$ with $g^{-1}\tau y\in\pfsBall$
($\pfsA_{a+1}$ is carried by $\pfrU_{(<a+1)}$). By the right invariance of $\pfsA_{a+1}$ under $\tau$ and
$\pfsA_{a+1}=\pfsA_a*\eta_a$ (convention (ii)),
\[
J_{a+1}(y)=\iint\mathbf 1[g^{-1}\nu nz\in\pfsBall]\,\pfrwt_{a+1}(g,g^{-1}\nu nz)\,d\pfsA_a(\nu)\,d\eta_a(n),\qquad z=\tau y .
\]
For fixed $n$ and $z'=nz$: $\pfrwt_{a+1}(g,g^{-1}\nu z')=\pfrwt_a(g,g^{-1}\nu z')\,\pfrdep_a(g,g^{-1}\nu z')^{-s}$, and
$\pfrdep_a(g,g^{-1}\nu z')=\pfrdep_a(g,g^{-1}z')$ for $\nu\in\pfrU_{(<a)}$: the second point is $\nu z'$, and the vectors
stable at $(g,\nu z')$ and at $(g,z')$ for $a$ are the same, because $(\nu z')^{-1}\zeta'=z'^{-1}(\nu^{-1}\zeta')$ and
$\nu^{-1}\zeta'\equiv\zeta'$ modulo $\langle e_c:c<a\rangle$ for $\zeta'$ orthogonal to the $e_c$, $c<a$ (convention (iii)). Moreover
the inner integral vanishes unless $z'\in\pfsCB_a(g)$. Hence the inner integral is at most
$J_a(z')\,\mathbf 1[z'\in\pfsCB_a(g)]\,\pfrdep_a(g,g^{-1}z')^{-s}\le C_a\,\mathbf 1[nz\in\pfsCB_a(g)]\,\pfrdep_a(g,g^{-1}nz)^{-s}$,
and Lemma~\ref{pfs:invdepthN} (for the point $z=\tau y$, with $g^{-1}z\in\pfsBall$) gives $J_{a+1}(y)\le C_a\pfsc_a(M)
\Lambda^{4600n}=C_{a+1}$.
\end{proof}

\subsection{One frame: proof of Theorem~\ref{pfr:oneframe}}\label{pfs:ss-oneframe}

Throughout this subsection $\Gamma$ is any gauge for $\pfrF$ (it will be applied to the gauge $\Gamma'$ of
Theorem~\ref{pfr:oneframe}), $\pfsBall$ its position ball, and $\mathsf k$ a kernel at floor $j\le\varrho$. Put
\[
\pfsfc_l=\pfsc_l(M)\,P^{\pfsrk}\Lambda^{4900n}\ (l<\varrho),\qquad\pfsfc_\varrho=1,\qquad
\|\mathsf k\|^2=\iint|\mathsf k(g,y)|^2\,d\theta_{\ge0}(y)\,d\theta_{\ge0}(g).
\]
For $l\le\varrho+1$ and $g,y\in\pfrG$ let
\[
\pfscm_l(g,y)=\pfsA_l\{\nu:\ g^{-1}\nu y\in\pfsBall\},\qquad
\pfscut_l(g,y)=\begin{cases}|\mathsf k_{l+1}(g,y)|^2&\text{if }\pfscm_l(g,y)>0,\\0&\text{otherwise}.\end{cases}
\]
Both are Borel functions of $(g,y)$ (Tonelli). If $\pfscm_l(g,y)>0$ then $y\in\pfsCB_l(g)$, since $\pfsA_l$ is carried by
$\pfrU_{(<l)}$.

\begin{lemma}[Translation rules]\label{pfs:cmrules}
For $l\le\varrho+1$: (a) $\pfscm_l(\sigma g,\sigma y)=\pfscm_l(g,y)$ for $\sigma\in\pfrG^{(l)}$; (b) $\pfscm_l(g,\nu y)=
\pfscm_l(g,y)$ for $\nu\in\pfrU_{(<l)}$; (c) $\int\pfscm_l(g,ny)\,d\eta_l(n)=\pfscm_{l+1}(g,y)$ for $l\le\varrho$.
\end{lemma}
\begin{proof}
(a) $\pfscm_l(\sigma g,\sigma y)=\int\mathbf 1_\pfsBall(g^{-1}\sigma^{-1}\nu\sigma y)\,d\pfsA_l(\nu)$, and $\nu\mapsto\sigma^{-1}\nu\sigma$
preserves $\pfsA_l$ (convention (ii)). (b) Right invariance of $\pfsA_l$ under $\nu$. (c) $\pfsA_{l+1}=\pfsA_l*\eta_l$ and Tonelli.
\end{proof}

\begin{lemma}[A factorisation of the Haar measure]\label{pfs:factor}
For $l\le\varrho+1$ let $\rho_l=\eta_l*\eta_{l+1}*\dots*\eta_\varrho*\theta_{\ge0}$ ($\rho_{\varrho+1}=\theta_{\ge0}$). Then
$\rho_l$ is s-finite, $\rho_l=\eta_l*\rho_{l+1}$ for $l\le\varrho$, and $\pfrM=\pfsA_l*\rho_l$. If $\Phi\ge0$ is Borel with
$\Phi(\nu y)=\Phi(y)$ for all $\nu\in\pfrU_{(<l)}$ and $y$, then $\int\Phi(\sigma y)\,d\rho_l(y)=\int\Phi\,d\rho_l$ for every
$\sigma\in\pfrG^{(l)}$.
\end{lemma}
\begin{proof}
By Lemma~\ref{pfr:rearr} with $j=\varrho+2$ (and $\eta_\varrho=\eta_{\varrho+1}=\delta_1$), $\pfrM=\pfsA_{\varrho+1}*\theta_{[0,\varrho+2)}=
\pfsA_{\varrho+1}*\theta_{\ge0}$; and $\pfsA_{\varrho+1}=\pfsA_l*\eta_l*\dots*\eta_\varrho$, which gives $\pfrM=\pfsA_l*\rho_l$ by associativity.
Convolutions of finitely many s-finite measures are s-finite. For the invariance we use the averaging argument of
the proof of Lemma~\ref{pfr:weylinv}. Choose a Borel $f>0$ with $\int f\,d\pfrM<\infty$ and put
$\mathrm{av}(y)=\int f(\nu y)\,d\pfsA_l(\nu)\in(0,\infty]$; $\mathrm{av}(\nu y)=\mathrm{av}(y)$ for $\nu\in\pfrU_{(<l)}$
by right invariance of $\pfsA_l$. For every s-finite $\mu$ with $\pfrM=\pfsA_l*\mu$ we have $\int\mathrm{av}\,d\mu=\int f\,d\pfrM<\infty$,
so $\mathrm{av}<\infty$ $\mu$-a.e., and
\[
\int\frac{\Phi f}{\mathrm{av}}\,d\pfrM=\iint\frac{\Phi(\nu y)f(\nu y)}{\mathrm{av}(\nu y)}\,d\pfsA_l(\nu)\,d\mu(y)=\int\Phi\,d\mu .
\]
The left side does not depend on $\mu$. For $\sigma\in\pfrG^{(l)}$, $\mu=\delta_\sigma*\rho_l$ is admissible:
$\pfsA_l*\delta_\sigma*\rho_l=\delta_\sigma*\pfsA_l*\rho_l=\delta_\sigma*\pfrM=\pfrM$ by convention (ii) and left invariance. Hence
$\int\Phi(\sigma y)\,d\rho_l(y)=\int\Phi\,d(\delta_\sigma*\rho_l)=\int\Phi\,d\rho_l$.
\end{proof}

\begin{lemma}[One step up]\label{pfs:climb}
Let $j\le l\le\varrho$ and $g,y\in\pfrG$. Then
\[
\int\pfscut_l(g,ny)\,d\eta_l(n)\le\pfsfc_l\int\pfscut_{l+1}(\sigma g,\sigma y)\,d\theta_{l+1}(\sigma).
\]
\end{lemma}
\begin{proof}
If $\pfscm_{l+1}(g,y)=0$, then by Lemma~\ref{pfs:cmrules}(c) $\pfscm_l(g,ny)=0$ for $\eta_l$-almost every $n$, and the left
side is $0$. Otherwise, by Lemma~\ref{pfs:cmrules}(a), $\pfscm_{l+1}(\sigma g,\sigma y)=\pfscm_{l+1}(g,y)>0$ for
$\sigma\in\pfrG^{(l+1)}$, which carries $\theta_{l+1}$; so the right side is $\pfsfc_l\int|\mathsf k_{l+2}(\sigma g,\sigma y)|^2\,
d\theta_{l+1}(\sigma)$, and $\pfscut_l(g,ny)\le|\mathsf k_{l+1}(g,ny)|^2$.

\emph{The top floor, $l=\varrho$.} Here $\eta_\varrho=\delta_1$ and $\pfsfc_\varrho=1$. As $\theta_{\ge\varrho+2}=\delta_1$ and
$\theta_{\varrho+1}=\mu_\pfrH$, (C1) gives $\mathsf k_{\varrho+1}(g,y)=\int\pfrip{W(\sigma g)}{W(d'\sigma y)}\,d\mu_\pfrH(\sigma)$; the integrand vanishes
unless $W(\sigma g)\ne0$, hence unless $u\sigma g\in\pfrK$ for some $u\in\pfrU$, a set of $\mu_\pfrH$-mass at most $1$ by the
definition of a place datum. By Cauchy--Schwarz, $|\mathsf k_{\varrho+1}(g,y)|^2\le\int|\mathsf k_{\varrho+2}(\sigma g,\sigma y)|^2\,
d\theta_{\varrho+1}(\sigma)$.

\emph{Below the top floor, $l<\varrho$.} We apply Proposition~\ref{pfs:parseval} with $a=l$, the kernel $\mathsf k$ (a kernel at
floor $l+1$), $\varepsilon=0$, $\xi=M^{-1}$, $u=M$, $C=1$, $i=1$, $j=17$, the points $g$ and $y$, and
\[
S=\{n\in\pfrN_l:\ ny\in\pfsCB_l(g)\text{ and }g\text{ is near at }l+1\},\qquad\Phi(n)=\pfscut_l(g,ny).
\]
\eqref{pfs:eq-twist}: for $s\in\pfrN_l\subseteq\pfrG^{(j)}$, $d'sd'^{-1}=s$, so $M^{-1}\sigma(d'sd'^{-1})=M^{-1}\sigma(s)$. (a): if
$\Phi(n)\ne0$, then $\pfscm_l(g,ny)>0$, so $ny\in\pfsCB_l(g)$, and $\mathsf k_{l+1}(g,ny)\ne0$, so $g$ is near at $l+1$ by (C2);
and $\Phi(n)\le|\mathsf k_{l+1}(n^0g,ny)|^2$. (b): let $n_0\in S$ and $\nu_0\in\pfrU_{(<l)}$ with $g^{-1}\nu_0n_0y\in\pfsBall$. As
$g$ is near at $l+1$, $\xi_0g\in\pfrK$ for some $\xi_0\in\pfrR^{(l+1)}\subseteq\pfrR^{(l)}$; write $\xi_0=t\nu$ with $t\in\pfrG^{(l)}$,
$\nu\in\pfrU_{(<l)}$ (convention (iii)). Then $t=\xi_0\nu^{-1}\in\pfrR^{(l+1)}\cap\pfrG^{(l)}$, $q=\nu\in\pfrU_{(<l+1)}$, and
$tqg=\xi_0g\in\pfrK$ has size $1$. Let $n\in S$, with $\mu\in\pfrU_{(<l)}$ such that $g^{-1}\mu ny\in\pfsBall$. Put $\delta=nn_0^{-1}
\in\pfrN_l$, $g'=t\nu g\in\pfrK$, $y'=t\nu\nu_0n_0y$ and $n'=t\delta t^{-1}$. Then $n'\in\pfrN_l$ (write $t=s'r'$ with
$s'\in\pfrG^{(l+1)}$, $r'\in\pfrN_l$ by (N1), and use (N2)); $g'^{-1}y'=g^{-1}\nu_0n_0y\in\pfsBall$; and
$n'y'\in\pfsCB_l(g')$, because $\kappa=t\nu\mu\,\delta(\nu\nu_0)^{-1}\delta^{-1}t^{-1}$ lies in $\pfrU_{(<l)}$ ($\delta$ and $t$ lie in
$\pfrG^{(l)}$, which normalises $\pfrU_{(<l)}$) and $g'^{-1}\kappa n'y'=g^{-1}\mu ny\in\pfsBall$. By Lemma~\ref{pfs:lift} (for
$g'$, $y'$, $n'$) some $v\equiv(n'^{-1}-1)e_l^*$ has $M^{-1}v$ of size $17$; and $n'^{-1}=tn_0n^{-1}t^{-1}=tn_0(tn)^{-1}$.

Proposition~\ref{pfs:parseval} gives the bound $(P/\pfrmod(M^{-1}M))^{\pfsrk}\pfsc_l(M)\Lambda^{100n\cdot49}\int|\mathsf k_{l+2}
(\sigma g,\sigma y)|^2d\theta_{l+1}(\sigma)$, and $100n\cdot49=4900n$.
\end{proof}

\begin{lemma}[All floors above $l$]\label{pfs:allfloors}
For $j\le l\le\varrho+1$,
\[
\Bigl(\prod_{a<l}\pfsfc_a\Bigr)\iint\pfscut_l(g,y)\,d\rho_l(y)\,d\theta_{[0,l+1)}(g)\le\Bigl(\prod_{a<\varrho}\pfsfc_a\Bigr)\|\mathsf k\|^2 .
\]
\end{lemma}
\begin{proof}
Descending induction on $l$. For $l=\varrho+1$: $\rho_{\varrho+1}=\theta_{[0,\varrho+2)}=\theta_{\ge0}$, $\pfscut_{\varrho+1}\le|\mathsf k_{\varrho+2}|^2=|\mathsf k|^2$
($\theta_{\ge\varrho+2}=\delta_1$), and $\prod_{a\le\varrho}\pfsfc_a=\prod_{a<\varrho}\pfsfc_a$. Let $j\le l\le\varrho$. For fixed $g$, by
$\rho_l=\eta_l*\rho_{l+1}$, Lemma~\ref{pfs:climb} and Tonelli,
\[
\int\pfscut_l(g,y)\,d\rho_l(y)=\iint\pfscut_l(g,ny)\,d\eta_l(n)\,d\rho_{l+1}(y)\le\pfsfc_l\iint\pfscut_{l+1}(\sigma g,\sigma y)\,
d\rho_{l+1}(y)\,d\theta_{l+1}(\sigma).
\]
For $\sigma\in\pfrG^{(l+1)}$ the function $\Phi(y)=\pfscut_{l+1}(\sigma g,y)$ satisfies $\Phi(\nu y)=\Phi(y)$ for $\nu\in\pfrU_{(<l+1)}$,
by Lemma~\ref{pfs:cmrules}(b) and (K3); so by Lemma~\ref{pfs:factor} the inner integral is $\int\pfscut_{l+1}(\sigma g,y)\,d\rho_{l+1}(y)$.
Integrating over $g$ against $\theta_{[0,l+1)}$ and using $\theta_{[0,l+2)}=\theta_{l+1}*\theta_{[0,l+1)}$, the left side of the
claim at $l$ is at most the left side at $l+1$.
\end{proof}

\begin{proposition}[One frame, one pivot]\label{pfs:pivot}
For $j\le l\le\varrho$ (and $\mathsf k$ a kernel at floor $j$),
\[
\int_{\pfsBall}\int\mathsf a_l(g;x)\,d\theta_{[0,l+1)}(g)\,d\pfrM(x)\le\Bigl(\prod_{a<\varrho}\pfsfc_a\Bigr)\|\mathsf k\|^2,
\qquad\mathsf a_l(g;x)=\pfrwt_l(g,x)\,|\mathsf k_{l+1}(g,gx)|^2 .
\]
\end{proposition}
\begin{proof}
By Tonelli it suffices to bound, for fixed $g$, $\int\mathbf 1_\pfsBall(x)\mathsf a_l(g;x)\,d\pfrM(x)$ by
$(\prod_{a<l}\pfsfc_a)\int\pfscut_l(g,y)\,d\rho_l(y)$, and then to apply Lemma~\ref{pfs:allfloors}. By the left invariance of
$\pfrM$ under $g^{-1}$ and $\pfrM=\pfsA_l*\rho_l$,
\[
\int\mathbf 1_\pfsBall(x)\mathsf a_l(g;x)\,d\pfrM(x)=\iint\mathbf 1_\pfsBall(g^{-1}\nu y)\,\pfrwt_l(g,g^{-1}\nu y)\,|\mathsf k_{l+1}(g,\nu y)|^2
\,d\pfsA_l(\nu)\,d\rho_l(y).
\]
By (K3), $|\mathsf k_{l+1}(g,\nu y)|=|\mathsf k_{l+1}(g,y)|$ for $\nu\in\pfrU_{(<l)}$. Fix $y$. If $\mathsf k_{l+1}(g,y)=0$, the inner
integral vanishes. If $\pfscm_l(g,y)=0$, the indicator vanishes for $\pfsA_l$-almost every $\nu$. Otherwise $g$ is near at $l+1$,
hence at $l$ (C2), $\pfscut_l(g,y)=|\mathsf k_{l+1}(g,y)|^2$, and by Lemma~\ref{pfs:depthweight} the inner integral is at most
$\prod_{a<l}\bigl(\pfsc_a(M)\Lambda^{4600n}\bigr)\cdot\pfscut_l(g,y)\le\prod_{a<l}\pfsfc_a\cdot\pfscut_l(g,y)$, as $P\ge1$.
\end{proof}

\begin{proposition}[One frame]\label{pfs:frame}
For a kernel $\mathsf k$ at floor $j\le\varrho$,
$\int_{\pfsBall}\mathsf S_j(\mathsf k;x)\,d\pfrM(x)\le n\,(\prod_{a<\varrho}\pfsfc_a)\,\|\mathsf k\|^2$.
\end{proposition}
\begin{proof}
Unfolding the recursion of \S\ref{pfr:ss-gauge} with $\theta_{[0,l+1)}=\theta_l*\theta_{[0,l)}$ gives $\int\mathsf b_l\,d\theta_{[0,l)}=
\int\mathsf a_l\,d\theta_{[0,l+1)}+\int\mathsf b_{l+1}\,d\theta_{[0,l+1)}$ for $l\le\varrho$, and $\mathsf b_n=0$; so
$\mathsf S_j(\mathsf k;x)=\sum_{l=j}^{\varrho}\int\mathsf a_l(g;x)\,d\theta_{[0,l+1)}(g)$, a sum of at most $n$ terms, each bounded
by Proposition~\ref{pfs:pivot}.
\end{proof}

\begin{lemma}[Translation by a Weyl ratio]\label{pfs:weylratio}
Let $j\le\varrho$, $p,q$ sequences of units, $d=w_{j,q}^{-1}w_{j,p}$, and $\Phi\ge0$ Borel on $\pfrG$. Then
\[
\int\Phi(dy)\,d\theta_{\ge0}(y)=\prod_{a<j}\pfrmod(p_a/q_a)^{d_a+1}\int\Phi\,d\theta_{\ge0}.
\]
\end{lemma}
\begin{proof}
Induction on $j$; $j=0$ is trivial. Let $j<\varrho$, write $d_l=w_{l,q}^{-1}w_{l,p}$, and for a unit $c$ let $w^{(j)}_c$ be the
Weyl factor of floor $j$: $e_j\mapsto ce_j^*$, $e_j^*\mapsto\bar c^{-1}e_j$, the identity elsewhere. By (iv) of
\S\ref{pfr:ss-descent}, $\theta_j=\lambda_j*\eta_j$ with $\lambda_j$ the image of $\pfrmod(c)^{d_j}dc$ under $c\mapsto(w^{(j)}_c)^{-1}$;
these elements commute with $\pfrG^{(j+1)}$, which carries $\theta_{\ge j+1}$. Hence $\theta_{\ge0}=\lambda_j*\mu$ with
$\mu=\theta_{\ge j+1}*\eta_j*\theta_{[0,j)}$, and for Borel $\Psi\ge0$,
$\int\Psi\,d\theta_{\ge0}=\int\pfrmod(c)^{d_j}\int\Psi((w^{(j)}_c)^{-1}y)\,d\mu(y)\,dc$. With $r=p_j/q_j$ one checks on $e_j,e_j^*$
that $d_{j+1}(w^{(j)}_{rc})^{-1}=d_j(w^{(j)}_c)^{-1}$ (both send $e_j^*\mapsto c^{-1}e_j$ and $e_j\mapsto\bar ce_j^*$, and agree with
$d_j$ elsewhere). Substituting $c=rc'$ ($dc=\pfrmod(r)\,dc'$; the non-units are null),
\[
\int\Phi(d_{j+1}y)\,d\theta_{\ge0}=\pfrmod(r)^{d_j+1}\int\pfrmod(c')^{d_j}\int\Phi(d_j(w^{(j)}_{c'})^{-1}y)\,d\mu\,dc'=
\pfrmod(r)^{d_j+1}\int\Phi(d_jy)\,d\theta_{\ge0},
\]
and the induction hypothesis concludes.
\end{proof}

\begin{proof}[Proof of Theorem~\ref{pfr:oneframe}]
Let $\pfsBall'$ be the position ball of $\Gamma'$, $\Lambda'=\Lambda_{\Gamma'}$, and $\pfsfc'_l$ the constants above for $\Gamma'$. Let
$d=w_{j,q}^{-1}w_{j,p}$, $D=(\prod_{a<j}\pfsc_a(p_a^{-1}))^2$, $J=\prod_{a<j}\pfrmod(p_a/q_a)^{d_a+1}$ and
$E_1=\int_{\pfsBall'}\mathsf S_0(\mathsf k^1;x)\,d\pfrM$, $E_2=\int_{\pfsBall'}\mathsf S_j(\mathsf k^d;x)\,d\pfrM$. Let
$\|W\|^2=\iint|\pfrip{W(g)}{W(y)}|^2\,d\theta_{\ge0}(g)\,d\theta_{\ge0}(y)$, so that $\mathsf N(W)=\prod_{l<\varrho}\pfsc_l(M)
P^{\varrho\pfsrk}\|W\|^2$.

\emph{Cauchy--Schwarz.} On $S\subseteq\pfsBall'$, $f^2\le e^2\mathsf Q(x)=e^2DC_\Gamma^3\mathsf S_0(\mathsf k^1;x)\mathsf S_j(\mathsf k^d;x)$, so
$(\int_Sf\,d\pfrM)^2\le e^2DC_\Gamma^3E_1E_2$. \emph{The two frames.} By Proposition~\ref{pfs:frame} for $\Gamma'$ and the kernels $\mathsf k^1$
(at floor $0$) and $\mathsf k^d$ (at floor $j$), $E_1\le nV\|\mathsf k^1\|^2$ and $E_2\le nV\|\mathsf k^d\|^2$ with $V=\prod_{l<\varrho}\pfsfc'_l$;
here $\|\mathsf k^1\|^2=\|W\|^2$ and, by Lemma~\ref{pfs:weylratio} applied to $y\mapsto|\pfrip{W(g)}{W(y)}|^2$ for each $g$,
$\|\mathsf k^d\|^2=J\|W\|^2$. \emph{The head.} For each $a<j$,
$\pfrmod(p_aq_a)^{d_a}\cdot\pfrmod(p_a)^{-(2d_a+1)}\cdot\pfrmod(p_a/q_a)^{d_a+1}=\pfrmod(q_a)^{-1}$; hence
$\pfrkap_j(p,q)^2DJ=\pfrmod((q_0\cdots q_{j-1})^{-1})\le\Lambda$ by the head condition on $q$ and (G3). Altogether
\[
\pfrkap_j(p,q)^2\Bigl(\int_Sf\,d\pfrM\Bigr)^2\le\Lambda\,e^2\,C_\Gamma^3\,(nV)^2\,\|W\|^4 .
\]
\emph{Constants.} $n\le2^n\le\Lambda'^n$, and $V=\prod_{l<\varrho}\pfsc_l(M)P^{\varrho\pfsrk}\Lambda'^{4900n\varrho}$; as $n+4900n\varrho\le25000n^2$,
$nV\|W\|^2\le\Lambda'^{25000n^2}\mathsf N(W)$. Taking square roots, with $C_\Gamma^{3/2}=\Lambda^{1500000n^3}$ and $\Lambda^{1/2}\le\Lambda$,
\[
\pfrkap_j(p,q)\int_Sf\,d\pfrM\le e\,\Lambda^{1500000n^3+1}\Lambda'^{25000n^2}\mathsf N(W).\qedhere
\]
\end{proof}

\subsection{The gains from a frequency and from two heads}\label{pfs:ss-gains}

\subsubsection*{Sizes at two near points}
\begin{lemma}[An element of the unipotent group]\label{pfs:unisize}
Let $j\le\varrho$, $d\ge1$, and $s\in\pfrU_{(<j)}$ such that the columns $se_a^*$, $a<j$, have size $d$. Then $s$ has size
$j(d+1)+4$ and $s^{-1}$ has size $j(d+1)+8$.
\end{lemma}
\begin{proof}
Let $C_{ab}=h(se_a^*,e_b)$ ($a,b<j$), the $e_b^*$-coordinate of $se_a^*$; it has size $d$. As $s^{-1}e_b\in e_b+\langle e_c:c<b\rangle$
(convention (iii)), $C_{ab}=h(e_a^*,s^{-1}e_b)=\delta_{ab}$ for $b\le a$; so the matrix $\bar C=(\bar C_{ab})$ is unitriangular,
and the entries of $\bar C^{-1}$ are $(j-1)\times(j-1)$ minors, of size $(j-1)(d+1)$ by (G0). Note $h(e_c,se_a^*)=\bar C_{ac}$.
\emph{Columns at $e_b$, $b<j$.} $se_b=e_b+\sum_{c<b}\mu_ce_c$, and $h(se_b,se_a^*)=h(e_b,e_a^*)=\delta_{ab}$ for $a<j$ gives
$\bar C\mu'=\mathbf e_b$ with $\mu'=(\mu_0,\dots,\mu_{b-1},1,0,\dots)$; so the $\mu_c$ are entries of $\bar C^{-1}$. \emph{Columns at basis
vectors $v$ of floor $\ge j$.} $sv=v+\sum_{c<j}\lambda_ce_c$ (convention (iii)), and $h(sv,se_a^*)=h(v,e_a^*)=0$ gives
$\bar C\lambda=-(h(v,se_a^*))_{a<j}$. Each $h(v,se_a^*)$ has size $d+3$ (products with the Gram matrix, of size $1$, and a sum of
at most $r$ terms), so each $\lambda_c$ has size $(j-1)(d+1)+(d+3)+1+1=j(d+1)+4$. \emph{Columns at $e_a^*$, $a<j$}: size $d$.
So $s$ has size $j(d+1)+4$, and $s^{-1}=J^{-1}s^*J$ has size $j(d+1)+8$ by (G0).
\end{proof}

\begin{lemma}[One translate for two near points]\label{pfs:twosizes}
Let $j\le\varrho$, $p$ a head of length $j$ for $\Gamma$, and $g\in\pfrG$ such that $g$ and $w_p^{-1}g$ are near at $j$. Then there
is $t\in\pfrG^{(j)}$ such that $tg$ has size $60n^2$ and $tw_p^{-1}g$ has size $99n^2$.
\end{lemma}
\begin{proof}
Let $\xi g\in\pfrK$ with $\xi\in\pfrR^{(j)}$, $\xi=t\nu$ with $t\in\pfrG^{(j)}$, $\nu\in\pfrU_{(<j)}$, and $s=t\nu t^{-1}\in\pfrU_{(<j)}
\subseteq\pfrR^{(j)}$; then $s(tg)=\xi g\in\pfrK$. The element $t$ commutes with $w_p$ ($w_p$ acts on the pairs of index $<j$, which $t$
fixes, and is the identity on $V_j$, which $t$ preserves), so $w_p^{-1}(tg)=t\,w_p^{-1}g$ is near at $j$ ($t\in\pfrR^{(j)}$). By
Lemma~\ref{pfr:headentries} at the point $tg$, $p_a$, $p_a^{-1}$ and the columns $se_a^*$ ($a<j$) have size $10n$. By
Lemma~\ref{pfs:unisize}, $s^{-1}$ has size $j(10n+1)+8\le20n^2$ (as $j<n$ and $n\ge2$), so $tg=s^{-1}(\xi g)$ has size $20n^2+2\le60n^2$.
The entries of $w_p^{-1}$ are $0$, $1$, $\bar p_a$ and $p_a^{-1}$, of size $10n$, so $tw_p^{-1}g=w_p^{-1}(tg)$ has size
$10n+60n^2+1\le99n^2$.
\end{proof}

\subsubsection*{The gain from a frequency}
\begin{proof}[Proof of Theorem~\ref{pfr:gain-char}]
Write $w=w_{j,p}$, $d=w_{j,q}^{-1}w_{j,p}$, $\omega(\gamma)=\psi(\eta h(\gamma f_j,f_j))$, $K=(P/\pfrmod(\xi))^{\pfsrk}$, and for
$\sigma\in\pfrG^{(j)}$ write $y_\sigma=\sigma^{-1}f_j$. Fix $g$ and put
\[
f(\sigma)=\mathsf k^1_{j+1}(\sigma g,\sigma gx),\qquad f'(\tau)=\mathsf k^d_{j+1}(\tau w^{-1}g,\tau w^{-1}gx).
\]
As $\sigma$ is an isometry, $h(\sigma\tau^{-1}f_j,f_j)=h(y_\tau,y_\sigma)$, so $\Phi_j(x;\omega;g)=\iint\psi(\eta h(y_\tau,y_\sigma))f(\sigma)
\overline{f'(\tau)}\,d\theta_j(\tau)\,d\theta_j(\sigma)$. By Lemma~\ref{pfr:separate} (with $f$ there the function $|\Phi_j(x;\omega;\cdot)|$,
whose $\pfrm_j$-integral is $\pfrF_j(x;\omega)$) and the first claim of Lemma~\ref{pfr:floorsq} for the two kernels (which bounds
$\int|f|^2d\theta_j$ by $\mathsf b^{(1)}_j(g)/\pfrwt_j(g,x)$ and $\int|f'|^2d\theta_j$ by $\mathsf b^{(d)}_j(w^{-1}g)/\pfrwt_j(w^{-1}g,x)$), it suffices to
show, for every $g$ with $\Phi_j(x;\omega;g)\ne0$,
\begin{equation}\label{pfs:eq-char}
|\Phi_j(x;\omega;g)|^2\le K\,C_\Gamma\int|f|^2\,d\theta_j\int|f'|^2\,d\theta_j .
\end{equation}
\emph{A common translate.} As $\Phi_j(x;\omega;g)\ne0$ and $\theta_j$ is carried by $\pfrG^{(j)}$, there are $\sigma_0,\tau_0\in\pfrG^{(j)}$
with $f(\sigma_0)\ne0\ne f'(\tau_0)$; by (C2), $\sigma_0g$ and $\tau_0w^{-1}g$ are near at $j+1$, so $g$ and $w^{-1}g$ are near at $j$
($\pfrR^{(j+1)}\pfrG^{(j)}\subseteq\pfrR^{(j)}$). Let $t$ be given by Lemma~\ref{pfs:twosizes}. \emph{Supports.} If $f(\sigma)\ne0$ ($\sigma\in
\pfrG^{(j)}$), then $\sigma g=(\sigma t^{-1})(tg)$ is near at $j+1$, and by Lemma~\ref{pfs:nearsize}, $ty_\sigma=(\sigma t^{-1})^{-1}f_j$ has size
$60n^2+6$. Likewise $f'(\tau)\ne0$ implies that $ty_\tau$ has size $99n^2+6$. Let $S$ and $S'$ be the sets of these $\sigma$, resp.\ $\tau$.
\emph{The character.} For $\sigma\in S$, $\tau\in S'$, $h(y_\tau,y_\sigma)=h(ty_\tau,ty_\sigma)$ has size $(60n^2+6)+(99n^2+6)+5\le200n^2$ (by
(G0); $n\ge2$), so $\psi(\eta h(y_\tau,y_\sigma))=\psi(\xi h(ty_\tau,ty_\sigma))$. Hence, with
\[
G(v)=\int\psi\bigl(-\xi\,h(ty_\tau,v)\bigr)f'(\tau)\,d\theta_j(\tau),
\]
the inner integral equals $f(\sigma)\overline{G(ty_\sigma)}$ for $\theta_j$-almost every $\sigma$ (both sides vanish for $\sigma\notin S$, and
$\overline{\psi(c)}=\psi(-c)$), and $\Phi_j(x;\omega;g)=\int_Sf(\sigma)\overline{G(ty_\sigma)}\,d\theta_j(\sigma)$. \emph{Cauchy--Schwarz and
Plancherel.} The map $\sigma\mapsto ty_\sigma$ carries $\theta_j$ to $c_j\mu_j$, $c_j\le2$; so
\[
|\Phi_j(x;\omega;g)|^2\le2\int|f|^2\,d\theta_j\cdot\int_{\{v\text{ of size }60n^2+6\}}|G(v)|^2\,d\mu_j(v),
\]
and Proposition~\ref{pfs:corrplanch} (with $b=j$, the kernel $\mathsf k^d$ at floor $j$, $t$, $q'=t$, the points $w^{-1}g$ and $w^{-1}gx$,
$i=99n^2$, the unit $-\xi$, $v_0=0$, $i'=60n^2+6$, $\rho=\mu_j$) bounds the last integral by
$2\Lambda^{99n(60n^2+12)}\Lambda^{99n(99n^2+12)}K\int|f'|^2d\theta_j$. As $4\le\Lambda^2$ and $99n(159n^2+24)+2\le10^6n^3$, \eqref{pfs:eq-char} follows.
\end{proof}

\subsubsection*{The gain from two heads}
\begin{lemma}\label{pfs:col}
Let $i<\varrho$, $t,t'\in\pfrG^{(i+1)}$, $n_0,n\in\pfrG$ and $h\in\pfrG$ such that $tn_0h$ and $t'nh$ have size $s$. Then the column
$(tn_0(tn)^{-1}-1)e_i^*$ has size $2s+6$.
\end{lemma}
\begin{proof}
$t^{-1}$ and $t'^{-1}$ fix $e_i^*$, so $tn_0(tn)^{-1}e_i^*=tn_0n^{-1}e_i^*=(tn_0h)(t'nh)^{-1}e_i^*$; the matrix $(tn_0h)(t'nh)^{-1}-1$ has size
$s+(s+4)+1+1$ by (G0).
\end{proof}

\begin{lemma}[Parseval along the last vector of two heads]\label{pfs:secbound}
Let $j=i+1\le\varrho$, $p$ a head of length $j$ for $\Gamma$, $x$ a position, $\xi=M^{-1}(p_i^{-1}-q_i^{-1})$ a unit, $w=w_{j,p}$,
$d=w_{j,q}^{-1}w_{j,p}$. Then for every $h\in\pfrG$,
\[
\int\mathbf 1[wnh\text{ near at }j]\,|\mathsf k^d_{i+1}(nh,nhx)|^2\,d\eta_i(n)\le\Bigl(\frac P{\pfrmod(\xi)}\Bigr)^{\pfsrk}\pfsc_i(p_i^{-1})\,C_\Gamma
\int|\mathsf k^d_{i+2}(\sigma h,\sigma hx)|^2\,d\theta_{i+1}(\sigma).
\]
\end{lemma}
\begin{proof}
Apply Proposition~\ref{pfs:parseval} with $a=i$, the kernel $\mathsf k^d$ (at floor $i+1$), $\varepsilon=1$, the units $-\xi p_i$ and
$u=p_i^{-1}$, $C=1$, sizes $99n^2$ and $999n^2$, the points $h$, $hx$, the function $\Phi(n)=\mathbf 1[wnh\text{ near at }j]|\mathsf k^d_{i+1}(nh,
nhx)|^2$ and $S=\{n\in\pfrN_i:\ wnh\text{ and }nh\text{ are near at }j\}$.
\eqref{pfs:eq-twist}: $d$ multiplies $f_i$ by $p_i/q_i$ and fixes $f_{i+1}$, so for $s=n(\omega,\beta)\in\pfrN_i$,
$dsd^{-1}f_{i+1}=f_{i+1}-(p_i/q_i)h(f_{i+1},\omega)f_i$, that is $\sigma(dsd^{-1})=(p_i/q_i)\sigma(s)$, and
$M^{-1}(p_i/q_i-1)\sigma(s)=-\xi p_i\,\sigma(s)$. (a): if $\Phi(n)\ne0$, then $wnh$ is near at $j$ and, by (C2), so is $nh$; and
$\Phi(n)\le|\mathsf k^d_{i+1}(n^1h,nhx)|^2$. (b): let $n_0\in S$. Lemma~\ref{pfs:twosizes} at the point $wn_0h$ (near at $j$, with
$w^{-1}(wn_0h)=n_0h$ near at $j$) gives $t\in\pfrG^{(j)}$ with $tn_0h$ of size $99n^2$; take $t_0=t\in\pfrR^{(i+1)}\cap\pfrG^{(i)}$ and
$q=n_0\in\pfrU_{(<i+1)}$. For $n\in S$, Lemma~\ref{pfs:twosizes} likewise gives $t'\in\pfrG^{(j)}$ with $t'nh$ of size $99n^2$, and by
Lemma~\ref{pfs:col} the vector $v=(tn_0(tn)^{-1}-1)e_i^*$ has size $198n^2+6$. By Lemma~\ref{pfr:headentries} at the point $wnh$ ($wnh$ and
$w^{-1}(wnh)=nh$ are near at $j$), $p_i$ has size $10n$, so $u^{-1}v=p_iv$ has size $10n+198n^2+7\le999n^2$.

Proposition~\ref{pfs:parseval} gives the bound with the constant $(P/\pfrmod(-\xi p_ip_i^{-1}))^{\pfsrk}\pfsc_i(p_i^{-1})\Lambda^{100n(1098n^2+31)}$,
and $100n(1098n^2+31)\le10^6n^3$.
\end{proof}

\begin{proof}[Proof of Theorem~\ref{pfr:gain-heads}]
Let $j=i+1$, $w=w_{j,p}$, $d=w_{j,q}^{-1}w_{j,p}$, $K=(P/\pfrmod(\xi))^{\pfsrk}$, $\varphi_1(g)=|\mathsf k^1_{j}(g,gx)|$, $\varphi_2(g)=|\mathsf k^d_{j}(g,
gx)|$, and $\pfrwt^-_l$ as in Lemma~\ref{pfr:invweight}. For $\omega=1$ the double integral defining $\Phi_j$ factorises, and by (C1)
$|\Phi_j(x;1;g)|=\varphi_1(g)\varphi_2(w^{-1}g)$; so $\pfrF_j(x;1)=\int\varphi_1(g)\varphi_2(w^{-1}g)\,d\pfrm_j(g)$.

\emph{Crossing the weights.} Put $u(h)=\mathbf 1[wh\text{ near at }j]\,\pfrwt_{j}(h,x)\,\pfrwt^-_i(w_{i,p}h)\,\varphi_2(h)^2$. We claim
\[
\bigl(\varphi_1(g)\varphi_2(w^{-1}g)\bigr)^2\le\bigl(\pfrwt_i(g,x)\varphi_1(g)^2\bigr)\,\pfrwt^-_j(w^{-1}g)\cdot u(w^{-1}g).
\]
The left side vanishes unless $g$ and $w^{-1}g$ are near at $j$ ((C2)). Then $\pfrwt^-_j(w^{-1}g)=\pfrwt_j(w^{-1}g,x)^{-1}$; and
$w_{i,p}w^{-1}$ is the inverse of the Weyl factor of floor $i$, an element of $\pfrG^{(i)}\subseteq\pfrR^{(i)}$, under which
$\pfrwt^-_i$ is invariant ((C3) and the invariance of nearness), so $\pfrwt^-_i(w_{i,p}w^{-1}g)=\pfrwt^-_i(g)=\pfrwt_i(g,x)^{-1}$. The
weights are finite and non-zero, and the right side is the left side.

\emph{Cauchy--Schwarz.} $\pfrF_j(x;1)^2\le I_1I_2$ with $I_1=\int\pfrwt_i(g,x)\varphi_1(g)^2\pfrwt^-_j(w^{-1}g)\,d\pfrm_j(g)$ and
$I_2=\int u(w^{-1}g)\,d\pfrm_j(g)$.

\emph{The first integral.} $w^{-1}=w_{\bar p}$ with the head $\bar p$ (proof of Lemma~\ref{pfr:separate}), and the function
$\pfrwt_i(g,x)\varphi_1(g)^2=\mathsf a^{(1)}_i(g)$ is invariant under $\pfrR^{(j)}$ by (K3) and vanishes off the points near at $j$ by (C2).
Lemma~\ref{pfr:invweight} (head $\bar p$ of length $j$, $b'=j$) gives $I_1\le\prod_{a<j}\pfsc_a(p_a^{-1})\,C_\Gamma\int\mathsf a^{(1)}_i\,d\theta_{[0,j)}$.

\emph{The second integral.} The function $u$ is invariant under $\pfrG^{(j)}$: for $\sigma\in\pfrG^{(j)}$, $w\sigma h=\sigma wh$, $\pfrwt_j(\sigma h,x)=
\pfrwt_j(h,x)$ (C3), $\pfrwt^-_i(w_{i,p}\sigma h)=\pfrwt^-_i(\sigma w_{i,p}h)=\pfrwt^-_i(w_{i,p}h)$, and $\varphi_2(\sigma h)=\varphi_2(h)$ by (K2) for the
kernel $\mathsf k^d$ at floor $j$. By Lemma~\ref{pfr:weylinv}, $I_2=\int u\,d\pfrm_j$. By Lemma~\ref{pfr:rearr} and convention (ii),
$\pfrm_j=\pfsA_j*\theta_{[0,j)}=\eta_i*\nu$ with $\nu=\pfsA_i*\theta_{[0,j)}$. For $n\in\pfrN_i$: $\pfrwt_j(nh,x)=\pfrwt_j(h,x)$ ($n\in\pfrR^{(j)}$),
and $\pfrwt^-_i(w_{i,p}nh)=\pfrwt^-_i(nw_{i,p}h)=\pfrwt^-_i(w_{i,p}h)$ ($n\in\pfrG^{(i)}$ commutes with $w_{i,p}$ and lies in $\pfrR^{(i)}$). Hence, by
Lemma~\ref{pfs:secbound} and the first claim of Lemma~\ref{pfr:floorsq} (at floor $j$, for $\mathsf k^d$),
\[
\int u(nh)\,d\eta_i(n)\le K\,\pfsc_i(p_i^{-1})\,C_\Gamma\,\pfrwt^-_i(w_{i,p}h)\,\mathsf b^{(d)}_j(h).
\]
The function $\mathsf b^{(d)}_j$ is invariant under $\pfrR^{(j)}$ and vanishes off the points near at $j$ (Remark~\ref{pfs:rem-corr}(a)). The
proof of Lemma~\ref{pfr:invweight} for the head $p$ of length $i$ and $b'=j$, with $\theta_{[0,j)}$ in place of $\theta_{[0,i)}$ in its Step~2
(Step~2 uses only that $\Phi(\nu y)=\Phi(y)$ for $\nu\in\pfrU_{(<i)}$ and the pointwise bound of Step~1), gives
$\int\mathsf b^{(d)}_j(h)\pfrwt^-_i(w_{i,p}h)\,d\nu(h)\le\prod_{a<i}\pfsc_a(p_a^{-1})\,C_\Gamma\,\mathsf S_j(\mathsf k^d;x)$. So
$I_2\le K\prod_{a<j}\pfsc_a(p_a^{-1})\,C_\Gamma^2\,\mathsf S_j(\mathsf k^d;x)$.

\emph{Conclusion.} $\int\mathsf a^{(1)}_i\,d\theta_{[0,j)}\le\mathsf S_0(\mathsf k^1;x)$ (unfold the recursion, as in the proof of
Proposition~\ref{pfs:frame}). Hence $\pfrF_j(x;1)^2\le K\bigl(\prod_{a<j}\pfsc_a(p_a^{-1})\bigr)^2C_\Gamma^3\,\mathsf S_0(\mathsf k^1;x)\,\mathsf S_j(\mathsf k^d;x)=
K\,\mathsf Q(x)$.
\end{proof}




%% ===== proofs/B-frames-realplace.tex =====
% proofs/B-frames-realplace.tex
% The seven real-place statements of B-frames-real (pfr:haar, pfr:realgauge, pfr:conv, pfr:top,
% pfr:formbound, pfr:remote, pfr:deepreal), and the check of the hypotheses (H1)--(H5).
% Labels of this file start with prp:.
\providecommand{\panot}[1]{\par\smallskip{\raggedright\noindent\textbf{[Not proved here:} #1\textbf{]}\par}\smallskip}
\providecommand{\pbdeg}{\operatorname{deg}}
\providecommand{\pbvol}{\operatorname{vol}}
\providecommand{\pbmTh}{\Theta}
\providecommand{\pfrG}{\mathsf G}
\providecommand{\pfrH}{\mathsf H}
\providecommand{\pfrU}{\mathsf U}
\providecommand{\pfrR}{\mathsf R}
\providecommand{\pfrM}{\mathsf M}
\providecommand{\pfrm}{\mathsf m}
\providecommand{\pfrK}{\mathsf K}
\providecommand{\pfrP}{\mathcal P}
\providecommand{\pfrCl}{\mathcal C}
\providecommand{\pfrI}{\mathcal I}
\providecommand{\pfrF}{\mathsf F}
\providecommand{\pfrB}{\mathsf B}
\providecommand{\pfrdep}{\operatorname{dp}}
\providecommand{\pfrkap}{\kappa}
\providecommand{\pfrN}{\mathsf N}
\providecommand{\pfrX}{\Xi}
\providecommand{\pfrip}[2]{\langle#1,#2\rangle}
\providecommand{\pfrell}{\ell}
% macros of this file (prefix \prp)
\providecommand{\prpD}{\mathsf D}
\providecommand{\prpvolU}[1]{\pbvol\auU(#1)}
\providecommand{\prpdens}{\mathsf h}

\section{The real place: the local statements}\label{prp:sec}

This section proves the statements of the real place that \S\ref{pfr:ss-realforms} states without
proof: Lemmas~\ref{pfr:haar}, \ref{pfr:realgauge}, \ref{pfr:deepreal} and
Propositions~\ref{pfr:conv}, \ref{pfr:top}, \ref{pfr:formbound}, \ref{pfr:remote}. We use the notation
of \S\ref{pfr:sec} throughout, and its results with their proofs; of the statements of
\S\ref{pfr:ss-shared} that have no proof there, we use Theorems~\ref{pfr:oneframe}, \ref{pfr:gain-char}
and~\ref{pfr:gain-heads} as stated (they are shared with the finite places and are proved elsewhere).

\emph{What is proved and what is not.}
\begin{itemize}
\item Lemma~\ref{pfr:realgauge} (the gauge axioms (G0)--(G5) and the formula for the depth),
Lemma~\ref{pfr:deepreal}, Propositions~\ref{pfr:conv}, \ref{pfr:top}, \ref{pfr:formbound}
and~\ref{pfr:remote} are proved in full (from the shared theorems just named).
\item Lemma~\ref{pfr:haar} is proved in full; the comparison with Haar measure (Lemma~\ref{prp:fibre})
uses the Gram fibre normalisation Lemma~\ref{pbo:GF} and Lemma~\ref{pbo:append} of \S\ref{pbo:ss-GF}.
\item Consequently Theorem~\ref{pfr:R} depends only on Theorems~\ref{pfr:oneframe}, \ref{pfr:gain-char},
\ref{pfr:gain-heads}, on results proved elsewhere in the paper, and on the hypotheses (H0)--(H5)
(\S\ref{prp:ss-R}).
\S\ref{prp:ss-H} says on what each of (H1)--(H5) rests for the Bessel functions of
Theorem~\ref{pbe:bessel}.
\end{itemize}

Throughout, $\pfrG=G_\tau$ in the normalised basis, which is orthonormal for the Euclidean norm
$|\cdot|$; $\|\cdot\|$ is the operator norm. The Gram matrix $J$ of the basis is a signed permutation
matrix ($J=J^{-1}=J^*$), so $g^{-1}=Jg^*J$ and $\|g^{-1}\|=\|g\|$ for $g\in\pfrG$; elements of $\pfrH$ are
Euclidean isometries. The place datum at $\tau$ is that of \S\ref{pfr:ss-realforms}, with box
$\pfrK=\{\|g\|\le\rho\}$, close pairs $|\omega_1|^2+|\omega_2|^2<1$, level $M=1$, and $\Gamma_L$ is the
filtration $\Gamma_L(i)=\{c:|c|\le L^i\}$ with loss $L^2$, for $L\ge\rho$. We write
$P_l(p)=p_0\cdots p_{l-1}$. By (H0), $\rho\ge8r\ge8$, $k-5=s_0+3n+1$, and
$r=2\varrho+k+1=s_0+5n+5$.

\subsection{The Haar ratio: proof of Lemma~\ref{pfr:haar}}\label{prp:ss-haar}

\subsubsection*{The Heisenberg measures}
At the real place the Heisenberg group $\pfrN_a$ of floor $a<\varrho$ consists of the elements
$n(v,\beta)$, $v\in V_{a+1}$ (the span of the basis vectors of floor $>a$), $\beta\in i\R$, with
$n(v,\beta)e_a^*=e_a^*+v+(\beta-\tfrac12h(v,v))e_a$, $n(v,\beta)e_a=e_a$, $n(v,\beta)w=w-h(w,v)e_a$ for
$w\in V_{a+1}$, and the identity on the basis vectors of floor $<a$; $\eta_a$ is the image of
Lebesgue measure $dv\,dt$ under $(v,t)\mapsto n(v,it)$. Put $\pfrU_{(<l)}=\pfrN_0\cdots\pfrN_{l-1}$
(so $\pfrU_{(<\varrho)}=\pfrU$, Lemma~\ref{pfr:rearr}) and $\eta_\pfrU=\eta_0*\eta_1*\dots*\eta_{\varrho-1}$,
that is $\int f\,d\eta_\pfrU=\int f(n_0n_1\cdots n_{\varrho-1})\,d\eta_0(n_0)\cdots d\eta_{\varrho-1}(n_{\varrho-1})$.
For $l\le\varrho$ and $g\in\pfrG$ let
\[
G_l(g)=\exp\Bigl(-\pi\sum_{b<l}\Bigl(\sum_{i:\ \text{floor}(i)>b}|g_{i,e_b^*}|^2+(\Iml g_{e_b,e_b^*})^2\Bigr)\Bigr),
\]
the Gaussian in the \emph{free entries} of the columns $e_b^*$, $b<l$, of $g$.

\begin{lemma}\label{prp:heis}
\textup{(a)} For $\nu\in\pfrU_{(<a)}$, $a<\varrho$, and $n=n(v,it)\in\pfrN_a$:
$G_{a+1}(\nu n)=e^{-\pi(|v|^2+t^2)}\,G_a(\nu)$.
\textup{(b)} $\int G_\varrho\,d\eta_\pfrU=1$.
\textup{(c)} The map $(n_0,\dots,n_{\varrho-1})\mapsto n_0\cdots n_{\varrho-1}$ is a homeomorphism of
$\pfrN_0\times\dots\times\pfrN_{\varrho-1}$ onto $\pfrU$, and $\eta_\pfrU$ is a Haar measure of $\pfrU$.
\end{lemma}
\begin{proof}
(a) Every element of $\pfrN_m$, $m<a$, maps a vector $w\in V_a$ to $w$ plus a multiple of $e_m$ (as
$V_a\subseteq V_{m+1}$); hence $\nu w-w\in\langle e_b:b<a\rangle$ for $w\in V_a$. The column $e_a^*$ of $\nu n$ is
$\nu(ne_a^*)$ with $ne_a^*\in V_a$; so its entries in the rows of floor $\ge a$ are those of $ne_a^*$: the
vector $v$ in the rows of floor $>a$, and $\beta-\frac12h(v,v)$, of imaginary part $t$ ($h(v,v)$ is
real), in the row $e_a$. For $b<a$, $n$ fixes $e_b^*$, so the column $e_b^*$ of $\nu n$ is that of $\nu$.

(b) Induction on $l\le\varrho$ for $\int G_l\,d(\eta_0*\dots*\eta_{l-1})=1$; for $l=0$ both sides are
$1$. By (a), Fubini's theorem and $\int_{\C^d\times\R}e^{-\pi(|v|^2+t^2)}dv\,dt=1$, the integral for $l=a+1$
equals the integral for $l=a$ (the measure $\eta_0*\dots*\eta_{a-1}$ is carried by $\pfrU_{(<a)}$).

(c) By the proof of (a) (applied to $\nu=n_0\cdots n_{a-1}$, and since $n_m$, $m>a$, fixes $e_a^*$), the free
entries of the column $e_a^*$ of $u=n_0\cdots n_{\varrho-1}$ are the coordinates $(v_a,t_a)$ of $n_a$. So
the product map is injective, with the continuous inverse $u\mapsto(n(v_a,it_a))_a$ read off the columns
of $u$; it is onto $\pfrU$ by Lemma~\ref{pfr:rearr}. Hence $\eta_\pfrU$ is the image of a Lebesgue measure under
a homeomorphism onto $\pfrU$: it is a nonzero Radon measure on $\pfrU$. It is right invariant under each
$\pfrN_l$: for $n\in\pfrN_l$ write $n_0\cdots n_{\varrho-1}n=n_0\cdots n_l\,n'\,n_{l+1}\cdots n_{\varrho-1}$ with
$n'=gng^{-1}$, $g=n_{l+1}\cdots n_{\varrho-1}\in\pfrG^{(l+1)}$; by (N2) of \S\ref{pfr:ss-floors}, $n'\in\pfrN_l$,
and integrating first in $n_l$ (the measure $\eta_l$ is right invariant under $\pfrN_l$) removes $n'$. As
the $\pfrN_l$ generate $\pfrU$, $\eta_\pfrU$ is a right Haar measure of $\pfrU$; $\pfrU$ is nilpotent, hence
unimodular, so $\eta_\pfrU$ is a Haar measure.
\end{proof}

\emph{Convention on $du$.} We take for the Haar measure $du$ of $\pfrU$ in \S\ref{pbe:ss-real} the measure
$\eta_\pfrU$; equivalently (by Lemma~\ref{prp:heis}(b) and the uniqueness of Haar measure, $\pfrU$ being
unimodular so that $\int G_\varrho(u^{-1})\,du=\int G_\varrho\,du$), $du$ is the Haar measure with
$\int_\pfrU G_\varrho(u^{-1})\,du=1$. Since $u^{-1}=Ju^*J$, $G_\varrho(u^{-1})$ is the standard Gaussian in the
free entries of the rows $e_a$ of $u$ (the complex entries above the diagonal block that are not fixed by
unitarity, and the imaginary parts of the entries $(e_a,e_a^*)$); so $du$ is Lebesgue measure in these
entries (Lemma~\ref{pbv:haar}). This is the measure $du$ fixed in \S\ref{pb:ss-setting}. The invariant measure on $\pfrU\backslash\pfrG$ is
normalised by $\int_\pfrG f\,dx=\int_{\pfrU\backslash\pfrG}\int_\pfrU f(uy)\,du\,dy$ ($f\ge0$ measurable).

\subsubsection*{The comparison with Haar measure}

\begin{lemma}[Fibre measures of frames]\label{prp:fibre}
For every Borel set $B\subseteq\pfrG$, $\pfrM(\{g:g^{-1}\in B\})\le\dfrac{\prpvolU r}{\prpvolU k}\,\pbvol(B)$,
where $\pbvol$ is the Haar measure $dx$ of \S\ref{pbe:ss-real} and $\prpvolU N$ is the volume of
$\auU(N)$ for the measure compatible with the compact dual \textup{(}Remark~\ref{pfr:rem-match}\textup{)}.
\end{lemma}
\begin{proof}
We use the fibre measures of \S\ref{pbo:ss-GF}: for a tuple $E_0$ the relation $\mu\preceq\lambda_{E_0}$,
Lemma~\ref{pbo:append}(a) (appending vectors) and Lemma~\ref{pbo:GF} (if $\mu\preceq\lambda_1$, $1$ the tuple of
the basis, then $\mu\le\frac{\prpvolU r}{2^r}dg$; the order of the basis vectors in the tuple is irrelevant, a
permutation of the columns preserving Lebesgue measure and $dT$). Write $\nu^{\rm inv}$ for the image of $\nu$
under $g\mapsto g^{-1}$. Since $\pfrM=(\eta_{\varrho+1}*\theta_{\varrho+1})*\dots*(\eta_0*\theta_0)$ with
$\theta_{\varrho+1}=\mu_\pfrH$, $\eta_\varrho=\eta_{\varrho+1}=\delta_1$, and $\eta_a^{\rm inv}=\eta_a$
($n(v,it)^{-1}=n(-v,-it)$), $\mu_\pfrH^{\rm inv}=\mu_\pfrH$,
\[
\pfrM^{\rm inv}=\theta_0^{\rm inv}*\eta_0*\theta_1^{\rm inv}*\eta_1*\dots*\theta_{\varrho-1}^{\rm inv}*\eta_{\varrho-1}*
\theta_\varrho^{\rm inv}*\mu_\pfrH .
\]
Let $E_a=(e_0,e_0^*,\dots,e_{a-1},e_{a-1}^*)$ and $\mu'_a=\theta_0^{\rm inv}*\eta_0*\dots*\theta_{a-1}^{\rm inv}*\eta_{a-1}$.
We show $4^{-a}\mu'_a\preceq\lambda_{E_a}$ by induction ($a=0$: $\delta_1$ and the empty tuple). Write $V_a$ for the
span of the basis vectors of floor $\ge a$; the conditions $E_a^*\mathsf J_1X=0$ say $X\in V_a$, and they are $2a$
orthonormal complex coordinates of $X$, so the fibre measure of $X\mapsto(E_a^*\mathsf J_1X,\Phi(X))$ over
$(0,t)$ is the disintegration of Lebesgue measure of $V_a$ along $h(X,X)$ at $t$.

\emph{Appending $e_a$.} $\nu=\frac12\theta_a^{\rm inv}$ is carried by $\pfrG^{(a)}$, which fixes $E_a$; the image of
$\theta_a^{\rm inv}$ under $u\mapsto ue_a$ is the image of $\theta_a$ under $\sigma\mapsto\sigma^{-1}e_a$, which is
$c_a\mu_a=2\mu_a$ with $\mu_a$ the quadric measure of $e_a$ in $V_a$ (Lemma~\ref{pfr:floorquad}, $\kappa_\C=2$), and $\mu_a$
is the fibre measure just described at $t=0$. So the image of $\nu$ is that fibre measure, and by
Lemma~\ref{pbo:append}(a), $4^{-a}\cdot\frac12\,\mu'_a*\theta_a^{\rm inv}\preceq\lambda_{(E_a,e_a)}$.

\emph{Appending $e_a^*$.} $\nu=\frac12\eta_a$ is carried by $\pfrN_a$, which fixes $E_a$ and $e_a$. The tuples $X$
with $h(X,e_b)=h(X,e_b^*)=0$ ($b<a$), $h(X,e_a)=1$ are $X=e_a^*+w+\mu e_a$, $w\in V_{a+1}$, and
$\Phi(X)=h(w,w)+2\Rel\mu$; in the coordinates $(w,\Rel\mu,\Iml\mu)$ Lebesgue measure disintegrates along $\Phi$ with
fibre measure $\frac12\,dw\,d\Iml\mu$ (as $d\Phi=2\,d\Rel\mu$). The image of $\eta_a$ under $n\mapsto ne_a^*=
e_a^*+v+(it-\frac12h(v,v))e_a$ is $dv\,dt$ in these coordinates. So the image of $\frac12\eta_a$ is the fibre measure,
and Lemma~\ref{pbo:append}(a) gives $4^{-a-1}\mu'_{a+1}\preceq\lambda_{E_{a+1}}$.

\emph{Appending $z$, then the $x_i$.} $\theta_\varrho^{\rm inv}$ is carried by $\pfrG^{(\varrho)}$, fixing $E_\varrho$, and
its image under $u\mapsto uz$ is the image of $\theta_\varrho$ under $\sigma\mapsto\sigma^{-1}z$, which is $c_\varrho\mu_\varrho=
2\mu_\varrho$ with $\mu_\varrho$ the quadric measure of $z$ in $V_\varrho$ (at $h=-1$; \S\ref{pfr:ss-realforms}), the fibre
measure for appending $z$; so $4^{-\varrho}\cdot\frac12\,\mu'_\varrho*\theta_\varrho^{\rm inv}\preceq\lambda_{(E_\varrho,z)}$. Finally,
as in the proof of Lemma~\ref{pbo:append}(b), the image of $\frac{\prpvolU k}{2^k}\mu_\pfrH$ under $h\mapsto(hx_i)_i$ is the fibre
measure of $X\mapsto X^*X$ on $k$-tuples of $X_0$ over $1$ (Lemma~\ref{pbo:GF} for $\auU(k)$), and the tuples
orthogonal to $(E_\varrho,z)$ are the tuples of $X_0$, in orthonormal coordinates. Hence
$2^{-2\varrho-1}\frac{\prpvolU k}{2^k}\pfrM^{\rm inv}\preceq\lambda_1$, and by Lemma~\ref{pbo:GF}
\[
\pfrM^{\rm inv}\le2^{2\varrho+1+k}\frac{\prpvolU r}{2^r\prpvolU k}\,dg=\frac{\prpvolU r}{\prpvolU k}\,dg,
\]
since $r=k+2\varrho+1$. Here $dg$ is the Haar measure of \textup{(DS1)(iii)}, the measure $dx$ of
\S\ref{pbe:ss-real}.
\end{proof}

\begin{proof}[Proof of Lemma~\ref{pfr:haar}]
\emph{$\pfrM$ is right invariant.} This is Lemma~\ref{pfr:upperinv} for $j=0$: $\rho_0=\pfrM$ and
$\pfrM*\delta_t=\pfrM$ for $t\in\pfrG^{(0)}=\pfrG$. (Its proof uses only the structure of the floors, not that
$\pfrM$ is a Haar measure.)

\emph{$\pfrM=\lambda\,dx$ with $0<\lambda<\infty$.} $\pfrM\neq0$, since each factor $\eta_a$, $\theta_a$ is a
nonzero measure carried by a group. By Lemma~\ref{prp:fibre} applied to $B^{-1}$, $\pfrM(B)\le C\pbvol(B^{-1})
=C\pbvol(B)$ with $C=\prpvolU r/\prpvolU k$ ($\pfrG$ is unimodular, so $dx$ is invariant under inversion);
hence $\pfrM$ is finite on compact sets, and, as $\pfrG$ is second countable and locally compact, it is a
right invariant Radon measure. By the uniqueness of right Haar measure, $\pfrM=\lambda\,dx$ (the right Haar
measures of the unimodular group $\pfrG$ are the multiples of $dx$), with $\lambda>0$ as $\pfrM\ne0$ and
$\lambda\le C<\infty$ by the inequality just written for a compact $B$ with $0<\pbvol(B)$. In particular
$\lambda\,\prpvolU k/\prpvolU r\le1$.

\emph{The quotient formula.} By Lemma~\ref{pfr:rearr} for $j=\varrho+2$ (the factors $\eta_l$, $l\ge\varrho$, are
$\delta_1$), $\int f\,d\pfrM=\iint f(uy)\,d\eta_\pfrU(u)\,d\theta_{\ge0}(y)$ for $f\ge0$ measurable, and
$\eta_\pfrU=du$ by our convention. Let $b(g)=e^{-\|g\|^2}$. For $y\in\pfrG$, $\|uy\|\ge\|u\|/\|y^{-1}\|$, and
$\|u\|$ is at least the modulus of each coordinate $v_a,t_a$ of the chart of Lemma~\ref{prp:heis}(c) (these
are entries of $u$); so $B(y):=\int_\pfrU b(uy)\,du$ is finite, positive, measurable in $y$, and
$B(u'y)=B(y)$ for $u'\in\pfrU$. Put $\beta=b/B$, so that $\int_\pfrU\beta(uy)\,du=1$ for all $y$. Let $\phi\ge0$ be measurable and
left invariant under $\pfrU$. Then
\[
\int\phi\,d\theta_{\ge0}=\iint\beta(uy)\phi(uy)\,du\,d\theta_{\ge0}(y)=\int\beta\phi\,d\pfrM=\lambda\int_\pfrG\beta\phi\,dx
=\lambda\int_{\pfrU\backslash\pfrG}\phi(y)\int_\pfrU\beta(uy)\,du\,dy=\lambda\int_{\pfrU\backslash\pfrG}\phi.\qedhere
\]
\end{proof}

\subsection{The gauge at $\C$: proof of Lemma~\ref{pfr:realgauge}}\label{prp:ss-gauge}

Let $L\ge\rho$. Then $L\ge8r\ge8$ and $L\ge r+1$. The space $X'$ is spanned by the $x_i$ other than $x_{i_0}$
and two blocks of two of them, so $\operatorname{rank}X'=k-5=s_0+3n+1$, $m=0$, and the volume exponent is
$s=\operatorname{rank}X'-(3n+1)=s_0$; by (H0), $s_0=k-3n-6\ge n^7-3n-6\ge10^{10}n^4$ (as $n>10^8$).

\begin{lemma}[A disc of products]\label{prp:disc}
Let $g,y\in\pfrG$ and let $v$ be a family of at most $n$ linear forms on $X'$. There are orthonormal
vectors $w_1,\dots,w_{s+1}$ of $X'$, killed by $v$, such that $|g^{-1}w|=|y^{-1}w|=|w|$ for every $w$ in their
span $O$. For $u\in X'$ put $q(u)=\sum_i|h(w_i,u)|^2$ \textup{(}on $X'$, $h(w,u)=u^*w$ in the coordinates of
the orthonormal basis, by (Cv1)\textup{)}. Then every $t\in\C$ with $|t|\le\sqrt{q(u)}/2$ is $h(w,u)$ for some
$w\in O$ with $|w|\le1/2$.
\end{lemma}
\begin{proof}
The vectors $w$ with $Jw=w$ are those with $w_{e_a}=w_{e_a^*}$ for all $a$ and $w_z=0$, the kernel of
$\varrho+1$ linear forms $N$; for them $h(w,w)=w^*Jw=|w|^2$. Let $C$ be the family of linear forms on $X'$
made of $v$, $N\circ g^{-1}$ and $N\circ y^{-1}$: at most $n+2(\varrho+1)\le3n$ forms. Its kernel has
dimension at least $\operatorname{rank}X'-3n=s+1$; let $w_1,\dots,w_{s+1}$ be orthonormal in it. For $w\in O$,
$vw=0$, and $g^{-1}w$ is fixed by $J$, so $|g^{-1}w|^2=h(g^{-1}w,g^{-1}w)=h(w,w)=|w|^2$ ($w\in X'$, on which $h$
is the Euclidean form); likewise for $y^{-1}$. Given $t$ with $|t|\le\sqrt{q}/2$, $q=q(u)>0$, put
$w=(t/q)\sum_ih(u,w_i)\,w_i$. Then, $h$ being linear in the first variable, $h(w,u)=(t/q)\sum_ih(u,w_i)h(w_i,u)=
(t/q)\sum_i|h(w_i,u)|^2=t$, and $|w|^2=|t|^2\sum_i|h(u,w_i)|^2/q^2=|t|^2/q\le\frac14$. (If
$q=0$ only $t=0$ occurs, with $w=0$.)
\end{proof}

\begin{proof}[Proof of Lemma~\ref{pfr:realgauge}]
\emph{(G0).} $|cc'|\le L^{i+i'}$; a sum of at most $r+1$ numbers of modulus $\le L^i$ has modulus
$\le(r+1)L^i\le L^{i+1}$; $\Gamma_L(i)$ is stable under $c\mapsto-c,\bar c$ and contains $0,1$; $J=J^{-1}$ has
entries $0,\pm1$. The consequences listed in (G0) follow from these rules, as stated there.

\emph{(G1).} $m=0\le9n$; $P=\pi^2\ge1$; $M=1$; at $\C$ the fixed element of trace one is $\tau_0=\frac12$ and
the skew unit is $\delta=\frac i2$, so $\tau_0\delta^{-1}=-i$; the basis vectors have entries of modulus at
most $2\le L$; an element of $\pfrK$ has entries of modulus at most $\|g\|\le\rho\le L$.

\emph{(G2).} If $|\omega_1|^2+|\omega_2|^2<1$, the entries of $\omega_1$ have modulus $\le1$ and those of
$\omega_1-\omega_2$ modulus $\le2\le L$. Given $i$, let $c=L^{-(i+2)}$, so $c^{-1}\in\Gamma_L(i+2)$; if
$\omega_1,\omega_2$ have entries of modulus $\le L^i$, then $|c\omega_j|^2\le rL^{2i}L^{-2i-4}$, and
$|c\omega_1|^2+|c\omega_2|^2\le2r/L^4<1$.

\emph{(G3).} For $|u|\le L^i$, $\operatorname{mod}(u)=|u|^2\le L^{2i}=\Lambda^i$; $\pbvol\Gamma_L(i)=\pi L^{2i}\le
L^{2i+2}=\Lambda^{i+1}$.

\emph{(G4).} Here $2m+2=2$, and the standard form on the block is $|y_1|^2+|y_2|^2$. In polar coordinates,
Lebesgue measure on $\C^2$ is $\int_0^\infty\kappa_t\,dt$, where $\kappa_t$ is the rotation-invariant
measure on the sphere $|y|^2=t$ of total mass $\pi^2t$ (the derivative of $\pbvol\{|y|^2\le t\}=\pi^2t^2/2$);
put $\kappa_t=0$ for $t\le0$. If $\kappa_t$ charges $\Gamma_L(i)^2$ then $t\le2L^{2i}$, so
$\kappa_t(\Gamma_L(i)^2)\le2\pi^2L^{2i}\le L^4L^{4i}=\Lambda^{2(i+1)}$.

\emph{(G5).} Let $g\in\pfrK$, $x\in\pfrG$, $v$ a family of $\varrho$ linear forms on $X'$, and $O$, $q$ as in
Lemma~\ref{prp:disc} for $g$ and $y=gx$. Put $\phi(u)=36/q(u)$ ($=\infty$ if $q(u)=0$); $\phi$ is measurable.
Each $w\in O$ with $|w|\le\frac12$ has $vw=0$ and $|g^{-1}w|^2+|(gx)^{-1}w|^2=2|w|^2<1$, so it is in
$\mathcal W$. Hence, by Lemma~\ref{prp:disc}, $\{h(\omega,u):\omega\in\mathcal W\}$ contains the disc of radius
$\sqrt q/2$, of area $\pi q/4\ge q/36$: the first condition. For the second, a product of two close pairs
satisfies $|h(\omega_1,\omega_1')|\le|\omega_1||\omega_1'|\le\frac12(|\omega_1|^2+|\omega_1'|^2)$, so (as $M=1$) every
element $h(\omega_1,\omega_1')-h(\omega_2,\omega_2')$ of $\mathcal D$ has modulus
$<\frac12(|\omega_1|^2+|\omega_2|^2)+\frac12(|\omega'_1|^2+|\omega'_2|^2)<1$. If $q>0$, take $\omega\in\mathcal W$ with
$h(\omega,u)=\sqrt q/2$; if $h(\omega,cu)=\bar c\,h(\omega,u)\in\mathcal D$ then $|c|\sqrt q/2<1$. So the set of these $c$ lies in
the disc of radius $2/\sqrt q$, of area $4\pi/q\le36/q$.

For the integral, complete $(w_i)$ to an orthonormal basis of $X'$ and let $Q$ be the unitary change of
coordinates; $u\mapsto Q^*u$ preserves Lebesgue measure, $q(u)=|y'|^2$ where $y'\in\C^{s+1}$ are the first
coordinates of $Q^*u$, and the other $3n$ coordinates $y''$ form the rest. If the entries of $u$ have
modulus $\le L^i$, those of $Q^*u$ have modulus $\le(\operatorname{rank}X')L^i\le L^{i+1}$; so $|y'|^2\le R:=
L\cdot L^{2i+2}$ and $y''$ lies in a polydisc of volume $(\pi L^{2i+2})^{3n}$. The law of $|y'|^2$ under
Lebesgue measure on $\C^{s+1}$ has density $\pi^{s+1}t^s/s!$ on $(0,\infty)$; hence
\[
\int_{\Gamma_L(i)^{X'}}\phi^s\le\int_0^R\Bigl(\frac{36}t\Bigr)^s\frac{\pi^{s+1}t^s}{s!}\,dt\cdot(\pi
L^{2i+2})^{3n}=\pi\frac{(36\pi)^s}{s!}R\,(\pi L^{2i+2})^{3n}\le4e^{144}R\,(\pi L^{2i+2})^{3n},
\]
using $\pi\le4$ and $(144)^s/s!\le e^{144}$. As $4,e,\pi\le L$, this is at most
$L^{1+144+(2i+3)+(2i+3)3n}=L^{145+(2i+3)(3n+1)}\le L^{2\cdot99n(i+1)}=\Lambda^{99n(i+1)}$, since
$198n(i+1)-145-(2i+3)(3n+1)=192ni-2i+189n-148\ge0$ for $n\ge1$.

\emph{The depth.} Let $a<\varrho$ and $D=\inf(|g^{-1}\zeta|^2+|(gx)^{-1}\zeta'|^2)$ over the $\zeta\perp e_b$
($b<a$) with $h(\zeta,e_a)=1$ and the $\zeta'\equiv\zeta$ modulo $\langle e_b:b<a\rangle$, that is, the
infimum of $|\zeta|_g^2+|\zeta|_{gx}^2$ over the duals of $e_a$ in $Y^{(a)}$ for the induced lengths. We
show $\mathfrak d_a(g,x)=\{c:|c|^2D<1\}$. The $e_a^*$-coordinate of $\omega$ is $h(\omega,e_a)$. If $c=0$,
take $\omega=0$. If $c\ne0$ and $c=h(\omega,e_a)$ with $\omega$ stable, $(g^{-1}\omega,(gx)^{-1}\omega')$ close, then
$\zeta=c^{-1}\omega$, $\zeta'=c^{-1}\omega'$ give $D\le|c|^{-2}(|g^{-1}\omega|^2+|(gx)^{-1}\omega'|^2)<|c|^{-2}$.
Conversely, if $|c|^2D<1$ choose $\zeta,\zeta'$ with $|c|^2(|g^{-1}\zeta|^2+|(gx)^{-1}\zeta'|^2)<1$; then
$\omega=c\zeta$ is stable with $h(\omega,e_a)=c$. So $\mathfrak d_a$ is the open disc of radius $D^{-1/2}$ (all of
$\C$ if $D=0$), of area $\pi/D$, and $\pfrdep_a(g,x)=D/\pi$.
\end{proof}

\begin{lemma}[The place axiom (P) at $\C$]\label{prp:planch}
The real place $(\C,\ dc\text{ Lebesgue},\ \psi(c)=e^{c-\bar c},\ \operatorname{mod}(c)=|c|^2)$ satisfies the axiom
\textup{(P)} of \S\ref{pfs:ss-planch} with $P=\pi^2$, with equality $\int_\C|\widehat G|=\pi^2$ for $G(v)=e^{-\pi|v|^2}$.
\end{lemma}
\begin{proof}
$G$ is continuous, integrable, $G(0)=1$. For $v=v_1+iv_2$, $z=z_1+iz_2$: $\bar vz-v\bar z=2i\Iml(\bar vz)=
2i(v_1z_2-v_2z_1)$, so $\widehat G(z)=\int_{\R^2}e^{-\pi|v|^2}e^{i\,k\cdot v}dv$ with $k=2(z_2,-z_1)$, $|k|=2|z|$; the
Gaussian integral gives $\widehat G(z)=e^{-|k|^2/(4\pi)}=e^{-|z|^2/\pi}>0$, and $\int_\C e^{-|z|^2/\pi}dz=\pi\cdot\pi=\pi^2$.
Moreover $\operatorname{mod}(\bar c)=|c|^2=\operatorname{mod}(c)$, $\operatorname{mod}$ is multiplicative, $d(cv)=|c|^2dv$
(the real Jacobian of $v\mapsto cv$ on $\R^2$), and the non-units $\{0\}$ form a null set.
\end{proof}

\emph{Conventions.} The convention of \S\ref{pfs:sec} ($h$ linear in its first and conjugate linear in its second
variable; $\pfrip vw=\sum\bar v_iw_i$) is the one used in this section: the $e_a^*$-coordinate of $\omega$ is
$h(\omega,e_a)$ and the $e_a$-coordinate of $v\in V_a$ is $h(v,e_a^*)$, both linear in $\omega$, $v$; the products
$h(w,u)=u^*w$ ($u,w\in X'$) of Lemma~\ref{prp:disc} and of (G5) are linear in $w$ and conjugate linear in $u$, and the
verification of (G5) above is written in this convention; the other uses of $h$ (isometries,
$h(\zeta,u^{-1}e_a)$, $s_w=h(w,v^\#)$ in Lemma~\ref{prp:unip}) hold in either convention. Only moduli of the
leading minors are used (Lemmas~\ref{prp:minor}, \ref{prp:headsupport}), which do not depend on it.

\subsection{Minors, supports, and four auxiliary bounds}\label{prp:ss-aux}

For $\gamma\in\pfrG$ let $\delta_l(\gamma)$ be the leading $l\times l$ minor of the frame matrix
$Z(\gamma)=(h(\gamma f_b,f_a))_{a,b}$.

\begin{lemma}[Minors of near points]\label{prp:minor}
Let $l\le n$, $\gamma,y\in\pfrG$, and suppose $y$ and $\gamma^{-1}y$ are near at $l$. Then
$|\delta_j(\gamma)|\le\rho^{2j}\le\rho^{9n}$ for $j\le l$.
\end{lemma}
\begin{proof}
Let $\xi y=k_1$, $\xi'\gamma^{-1}y=k_2$ with $\xi,\xi'\in\pfrR^{(l)}$ and $k_1,k_2\in\pfrK$. Then
$\gamma=\xi^{-1}k\,\xi'$ with $k=k_1k_2^{-1}$, $\|k\|\le\rho^2$. For $b<l$, $\xi'f_b\in f_b+\langle f_c:c<b\rangle$
and $\xi f_a\in f_a+\langle f_c:c<a\rangle$ (also for the inverses), so
$h(\gamma f_b,f_a)=h(k\xi'f_b,\xi f_a)$ and the leading $j\times j$ block of $Z(\gamma)$ is $\mathsf L^*Z_j(k)\mathsf R$
with unitriangular $\mathsf L,\mathsf R$, where $Z_j(k)$ is the leading block of $Z(k)$. Its entries are
entries of $Jk$, a $j\times j$ submatrix of $Jk$, of operator norm $\le\rho^2$; so
$|\delta_j(\gamma)|=|\det Z_j(k)|\le\rho^{2j}$.
\end{proof}

\begin{lemma}[Heads of near points]\label{prp:headsupport}
Let $j\le\varrho$, $W$ a cut Bessel function for the place datum at $\tau$, and $p,q$ sequences of units.
If $\pfrF^+_j(x)\ne0$ for some $x$, then $|P_l(p)|\le\rho^{9n}$ and $|P_l(q)|\le\rho^{9n}$ for all $l\le j$.
\end{lemma}
\begin{proof}
Let $w=w_{j,p}$, $d=w_{j,q}^{-1}w$, and $g$ with $\pfrB^1_j(g,gx)\pfrB^d_j(w^{-1}g,w^{-1}gx)\ne0$. As $\theta_j$
is carried by $\pfrG^{(j)}$ and $\theta_{\ge j+1}$ by $\pfrG^{(j+1)}$, the first factor gives $\varsigma\in\pfrG^{(j)}$
with $W(\varsigma g)\ne0\ne W(\varsigma gx)$, and the second gives $\tau\in\pfrG^{(j)}$ with $W(\tau w^{-1}g)\ne0$
and $W(d\tau w^{-1}gx)=W(\tau w_{j,q}^{-1}gx)\ne0$ ($d$ commutes with $\pfrG^{(j)}$). A point $y$ with $W(y)\ne0$
is near at every $l$ (cut condition, and $\pfrU\subseteq\pfrR^{(l)}$), and nearness at $j$ is invariant under left
translation by $\pfrG^{(j)}\subseteq\pfrR^{(j)}$. Hence $y_1=\varsigma g$ and $w^{-1}y_1=\varsigma\tau^{-1}(\tau
w^{-1}g)$ are near at $j$; and $y_2=\varsigma gx$ and $w_{j,q}^{-1}y_2$ are near at $j$. The frame matrix of
$w_{j,p}$ has leading block $\operatorname{diag}(p_0,\dots,p_{j-1})$ up to complex conjugation of the
entries ($w_pe_b=p_be_b^*$ and $h(e_b^*,e_a)=\delta_{ab}$), so $|\delta_l(w_{j,p})|=|P_l(p)|$, and
Lemma~\ref{prp:minor} applies to $(w_{j,p},y_1)$ and $(w_{j,q},y_2)$.
\end{proof}

\begin{lemma}[Heads for one gauge]\label{prp:exgauge}
Let $j\le\varrho$, $K\ge0$ and $p,q$ sequences of units. Let $L=\max\bigl(\rho,\max_{l\le j}|P_l(p)|^{-1},
\max_{l\le j}|P_l(q)|^{-1}\bigr)$. Then $L\ge\rho$, $p$ and $q$ are heads of length $j$ for $\Gamma_L$, and
$L^K\le\rho^K\prod_{l=1}^j\max(1,|P_l(p)|^{-1})^K\max(1,|P_l(q)|^{-1})^K$. In particular, if
$|P_l(p)|,|P_l(q)|\le b^{60n^2}$ for $l\le j$, then $L^K\le\rho^K\Lambda^\infty_j(p)\Lambda^\infty_j(q)$ for the
weights of Theorem~\ref{pfr:R} with exponent $K$ in place of $n^6$.
\end{lemma}
\begin{proof}
A head of length $j$ for $\Gamma_L$ means $|P_l(p)|^{-1}\le L$ for $l\le j$. The bound holds since the maximum of
numbers $\ge1$ is at most their product, $\rho\ge1$ and $P_0=1$.
\end{proof}

\begin{lemma}[Schur's bound on the quotient]\label{prp:schur}
Let $W$ be a measurable Bessel function of irreducible type $\mathbb W$ \textup{(}as in
\S\ref{pfr:ss-realdata}\textup{)} with $\int_{\pfrU\backslash\pfrG}\|W\|^2\le\#I$. Then
$\int_{\pfrU\backslash\pfrG}|\pfrip v{W(y)}|^2dy\le|v|^2$ for every $v\in\mathbb W$.
\end{lemma}
\begin{proof}
The function $y\mapsto W(y)W(y)^*$ is left invariant under $\pfrU$ (the character factors cancel), and
integrable on $\pfrU\backslash\pfrG$; let $A=\int_{\pfrU\backslash\pfrG}W(y)W(y)^*dy$, a positive endomorphism of
$\mathbb W$ with $\operatorname{tr}A=\int\|W\|^2\le\#I$. For $h\in\pfrH$, $W(hy)=\pi(h)W(y)$ with $\pi$ the
representation, and $hy$ runs over $\pfrU\backslash\pfrG$ with the invariant measure when $y$ does: $h$
normalises $\pfrU$, conjugation by $h$ preserves $du$ (the modulus of the compact group $\pfrH$ acting on
$\pfrU$ is $1$), and $dx$ is left invariant. Hence $\pi(h)A\pi(h)^*=A$, and by Schur's lemma $A=c\cdot1$ with
$c\,\#I=\operatorname{tr}A\le\#I$. So $\int|\pfrip v{W}|^2=\pfrip v{Av}=c|v|^2\le|v|^2$.
\end{proof}

\begin{lemma}[Two bounds for the cut Bessel function]\label{prp:normW}
Under \textup{(H0)}--\textup{(H1)}, with $\lambda$ the ratio of Lemma~\ref{pfr:haar}:
\textup{(a)} $\mathsf N(W^\flat)\le\pi^{2(k-5)\varrho}\lambda^2\,\#I$;
\textup{(b)} for $j\le\varrho$ and $L\ge\rho$, the mass of floor $j$ is at most
$V:=(\pi\rho^{24})^k\rho^{9999n}$, and $V\ge1$.
\end{lemma}
\begin{proof}
(a) At $\C$, $M=1$, $P=\pi^2$, $\operatorname{rank}X'=k-5$, so $\mathsf N(W^\flat)=\pi^{2\varrho(k-5)}\iint
|\pfrip{W^\flat(g)}{W^\flat(y)}|^2d\theta_{\ge0}(g)\,d\theta_{\ge0}(y)$. For fixed $g$ the function
$y\mapsto|\pfrip{W^\flat(g)}{W^\flat(y)}|^2$ is left invariant under $\pfrU$, so by Lemma~\ref{pfr:haar} and
Lemma~\ref{prp:schur} (with (H1)) its $\theta_{\ge0}$-integral is $\lambda\int_{\pfrU\backslash\pfrG}\le\lambda
\|W^\flat(g)\|^2$; and $\int\|W^\flat\|^2d\theta_{\ge0}=\lambda\int_{\pfrU\backslash\pfrG}\|W^\flat\|^2\le\lambda\,\#I$.

(b) Lemmas~\ref{pfr:chart} and~\ref{pfr:quadricX} hold for every $a\le\varrho$, with the constant $c_a=2$ of the
real place (\S\ref{pfr:ss-realforms}). Apply them, for the gauge $\Gamma_\rho$ (loss $\Lambda=\rho^2$), with $c\equiv1$
(invariant), $\Psi_g\equiv1$ and $\phi\equiv1$, for which $\int_{\Gamma_\rho(11)^{X'}}1=(\pi\rho^{22})^{k-5}\le
\Lambda^{12(k-5)}$: the mass of floor $j$ is at most $2\Lambda^{99n\cdot13+12(k-5)}=2\rho^{2574n}\rho^{24(k-5)}\le V$
(as $2\le\pi^k$ and $\rho\ge1$). The mass does not depend on the gauge. $V\ge1$ is clear.
\end{proof}

\subsection{Far means deep: proof of Lemma~\ref{pfr:deepreal}}\label{prp:ss-deep}

For $a\le\varrho$ let $V_a$ be the span of the basis vectors of floor $\ge a$ (so $V_0=\C^r$, and $V_\varrho$ is the
span of $z$ and the $x_i$), $E_a=\langle e_b:b<a\rangle$, and $\pi_a$ the coordinate projection onto $V_a$ along
$E_a\oplus\langle e_b^*:b<a\rangle$; on $E_a^\perp=V_a\oplus E_a$ it is the projection along $E_a$, and
$|\pi_aw|\le|w|$.

\begin{lemma}[Bounded duals bound a unipotent element]\label{prp:unip}
Let $m\in\pfrU$, $n=m^{-1}$, and $R\ge1$. Suppose that for every $a<\varrho$ there are vectors
$\alpha_a,\beta_a\perp E_a$ with $h(\alpha_a,e_a)=h(\beta_a,e_a)=1$, $|\alpha_a|,|\beta_a|\le R$,
$m\alpha_a\equiv\beta_a$ and $n\beta_a\equiv\alpha_a$ modulo $E_a$. Then $\|m\|,\|n\|\le(10R^2)^\varrho$.
\end{lemma}
\begin{proof}
An element of $\pfrU$ preserves each $E_a$, hence $E_a^\perp$; so $\bar m_a=\pi_a\circ m|_{V_a}$ is well defined on
$V_a$, and $\pi_a(mw)=\bar m_a(\pi_aw)$ for $w\in E_a^\perp$; likewise $\bar n_a$, and $\bar n_a=\bar m_a^{-1}$. As $E_a$
is isotropic and orthogonal to $V_a$, $\bar m_a$ preserves $h$ on $V_a$. Moreover: $\bar m_a e_a=e_a$
($me_a\in e_a+E_a$); $\bar m_\varrho=1$ (the middle block of an element of $\pfrU$ is $1$); and for $w\in V_{a+1}$,
$\pi_{a+1}(\bar m_aw)=\bar m_{a+1}w$, so $\bar m_aw=\bar m_{a+1}w+s_w e_a$ with $s_w=h(\bar m_aw,e_a^*)$. Let
$D_a=\max(\|\bar m_a-1\|,\|\bar n_a-1\|)$ (operator norms on $V_a$); $D_\varrho=0$. We show
$1+D_a\le10R^2(1+D_{a+1})$; then $\|m\|=\|\bar m_0\|\le1+D_0\le(10R^2)^\varrho$, and likewise for $n$.

Write $\pi_a\alpha_a=e_a^*+w_\alpha+t_\alpha e_a$ and $\pi_a\beta_a=e_a^*+w_\beta+t_\beta e_a$ with
$w_\alpha,w_\beta\in V_{a+1}$ (the $e_a^*$-coordinate is $h(\cdot,e_a)=1$); all these pieces have norm $\le R$. From
$m\alpha_a\equiv\beta_a$: $\bar m_a(\pi_a\alpha_a)=\pi_a\beta_a$, and similarly $\bar n_a(\pi_a\beta_a)=\pi_a\alpha_a$.
Write $\bar m_ae_a^*=e_a^*+w^\#+s^\#e_a$ and $\bar n_ae_a^*=e_a^*+v^\#+t^\#e_a$ ($w^\#,v^\#\in V_{a+1}$; the
$e_a^*$-coordinate is $h(\bar m_ae_a^*,e_a)=h(e_a^*,\bar n_ae_a)=1$).

\emph{The coefficients $s_w$.} For $w\in V_{a+1}$, $s_w=h(\bar m_aw,e_a^*)=h(w,\bar n_ae_a^*)=h(w,v^\#)$, so
$|s_w|\le|w|\,|v^\#|$. Comparing the $V_{a+1}$-components of $\bar n_a(\pi_a\beta_a)=\pi_a\alpha_a$ gives
$v^\#+\bar n_{a+1}w_\beta=w_\alpha$, so $v^\#=(w_\alpha-w_\beta)-(\bar n_{a+1}-1)w_\beta$ and
$|v^\#|\le(2+D_{a+1})R$. Symmetrically $|w^\#|\le(2+D_{a+1})R$, and the coefficient of $e_a$ in $\bar n_aw$ has
modulus $\le|w||w^\#|$.

\emph{The column $e_a^*$.} Comparing the $e_a$-components of $\bar m_a(\pi_a\alpha_a)=\pi_a\beta_a$:
$s^\#+s_{w_\alpha}+t_\alpha=t_\beta$, so $|s^\#|\le2R+R(2+D_{a+1})R$. So
$|(\bar m_a-1)e_a^*|\le|w^\#|+|s^\#|\le(2+D_{a+1})(R+R^2)+2R$.

\emph{Conclusion.} On $V_a=\C e_a\oplus V_{a+1}\oplus\C e_a^*$ (orthogonal for the Euclidean norm),
$(\bar m_a-1)e_a=0$, $|(\bar m_a-1)w|\le|(\bar m_{a+1}-1)w|+|s_w|\le(D_{a+1}+(2+D_{a+1})R)|w|$, and the bound on
$e_a^*$ above; so $\|\bar m_a-1\|$ is at most the sum of the two coefficients, $\le(1+D_{a+1})(1+6R+2R^2)$. The
same holds for $\bar n_a$, and $1+D_a\le(1+D_{a+1})(2+6R+2R^2)\le10R^2(1+D_{a+1})$.
\end{proof}

\begin{proof}[Proof of Lemma~\ref{pfr:deepreal}]
Let $y$ be active at $x$: $\pfrip{W^\flat(y)}{W^\flat(yx)}\ne0$, and $w_p^{-1}y$, $w_q^{-1}yx$ are near at $j$.
Suppose that $\pfrdep_a(y,x)<\tau'^{-2}$ for every $a$ with $j\le a<\varrho$; we show that all entries of $x-1$
have modulus $\le T^{99n}$, which proves the lemma. By the cut condition (H1) there are $u,v\in\pfrU$ with
$k_1=uy\in\pfrK$ and $k_2=vyx\in\pfrK$. Put $m=vu^{-1}\in\pfrU$, $n=m^{-1}$. Then $x=y^{-1}(yx)=k_1^{-1}nk_2$, so
$\|x\|\le\rho^2\|n\|$.

\emph{Duals for $a\ge j$.} By Lemma~\ref{pfr:realgauge}, the infimum $D$ defining the depth is $\pi\,\pfrdep_a(y,x)<\pi\tau'^{-2}\le\pi T^2/4<T^2$
(as $\tau'\ge2/T$); this gives $\zeta\perp E_a$ with $h(\zeta,e_a)=1$ and $\zeta'\equiv\zeta$ modulo $E_a$ with
$|y^{-1}\zeta|^2+|(yx)^{-1}\zeta'|^2<T^2$.

\emph{Duals for $a<j$.} Take $\zeta=\zeta'=e_a^*$. By Lemma~\ref{pfr:headentries} (for $\Gamma_L$, the head $p$,
the point $y$, $s=u\in\pfrU\subseteq\pfrR^{(j)}$ and the element $\eta\in\pfrR^{(j)}$ given by the nearness of
$w_p^{-1}y$), the column $ue_a^*$ has entries of modulus $\le L^{10n}$, so $|y^{-1}e_a^*|=|k_1^{-1}ue_a^*|\le
\rho\sqrt rL^{10n}\le L^{10n+2}$; likewise, with $q$, $yx$ and $v$, $|(yx)^{-1}e_a^*|\le L^{10n+2}\le T$.

In both cases put $\alpha_a=u\zeta$ and $\beta_a=v\zeta'$. They are orthogonal to $E_a$ (as $u,v$ preserve $E_a$
and $E_a^\perp$), $h(\alpha_a,e_a)=h(\zeta,u^{-1}e_a)=h(\zeta,e_a)=1$ and likewise $h(\beta_a,e_a)=1$;
$|\alpha_a|=|k_1y^{-1}\zeta|\le\rho T$ and $|\beta_a|=|k_2(yx)^{-1}\zeta'|\le\rho T$; and modulo $E_a$,
$m\alpha_a=v\zeta\equiv v\zeta'=\beta_a$ and $n\beta_a=u\zeta'\equiv u\zeta=\alpha_a$. By Lemma~\ref{prp:unip} with
$R=\rho T$, $\|n\|\le(10\rho^2T^2)^\varrho$. Hence every entry of $x-1$ has modulus at most
$1+\rho^2(10\rho^2T^2)^\varrho\le2\cdot10^\varrho\rho^{2\varrho+2}T^{2\varrho}\le T^{1+\varrho+(2\varrho+2)+2\varrho}=T^{5\varrho+3}\le
T^{99n}$, using $2,10,\rho\le L\le T$.
\end{proof}

\begin{lemma}[Dyadic shells]\label{prp:dyadic}
Let $p,q$ be heads of length $j\le\varrho$ for $\Gamma_L$, $L\ge\rho$, and $W=W^\flat$. For $i\ge0$ the set
$S_i$ of the $x$ for which $2^i$ is the least power of $2$ bounding the moduli of all entries of $x-1$ is a
shell of height $2^i$ for $\Gamma_L,j,p,q,W^\flat$ (Definition~\ref{pfr:shell}); the $S_i$ cover $\pfrG$.
\end{lemma}
\begin{proof}
Every $x$ lies in exactly one $S_i$. (a) Let $H=2^i$ and $L'=L2^i\ge\rho$. The position ball of $\Gamma_{L'}$
($M=1$, all positions allowed) contains $S_i$, and $\Lambda_{\Gamma_{L'}}=L^24^i=\Lambda H^2$; one gauge
suffices, and $1\le2H^{9n}$. (b) Let $\varsigma\in(0,1)$ with $\Lambda^{1200n^2}\le\varsigma^{400n^2}H$; then $i\ge1$
(otherwise $L^{2400n^2}\le\varsigma^{400n^2}<1$). Put $T=L^{11n+1}/\varsigma\ge L^{11n}$; then
$T^{-1}\le\varsigma/L\le\varsigma/2$. For $x\in S_i$ some entry of $x-1$ exceeds $2^{i-1}=H/2$, and
\[
\varsigma^{400n^2}\cdot2T^{99n}\le\varsigma^{99n}\,L\,(L^{11n+1}/\varsigma)^{99n}=L^{1089n^2+99n+1}\le L^{2400n^2}
\le\varsigma^{400n^2}H,
\]
so $2T^{99n}\le H$ and some entry of $x-1$ exceeds $T^{99n}$. By Lemma~\ref{pfr:deepreal} with $\tau'=\varsigma$,
$x$ has a deep vector of depth $\varsigma^{-2}$.
\end{proof}

\subsection{The sum over the dyadic shells; Propositions~\ref{pfr:formbound}, \ref{pfr:conv}, \ref{pfr:remote}}\label{prp:ss-sum}

Recall $\pbmTh=\pi^{(2\varrho+1)k}\allowbreak \prpvolU k/\prpvolU r$ and $\mathsf c=\frac{\prpvolU k}{\prpvolU r}(\lambda\,\#I)^{-2}$
\eqref{pfr:eq-c}.

\begin{proposition}[The sum over the shells at $\C$]\label{prp:sum}
Assume \textup{(H0)}--\textup{(H1)}. Let $j\le\varrho$, $L\ge\rho$, $p,q$ heads of length $j$ for $\Gamma_L$,
$t\ge1$, and $f\ge0$ measurable with $f\le\pfrF^+_j$ everywhere \textup{(}for $W=W^\flat$\textup{)} and, if
$t>1$, $f(x)^2\le(V/t^{s_0})^2\mathsf Q_{\Gamma_L}(x)$ for all $x$, with $V$ as in Lemma~\ref{prp:normW}. Then
\[
\#I\cdot\mathsf c\,\pfrkap_j(p,q)\int f\,d\pfrM\le\pbmTh\,L^{10^9n^4}\rho^{24k}\,t^{10^8n^4}/t^{s_0}.
\]
\end{proposition}
\begin{proof}
Apply Theorem~\ref{pfr:shells} to the dyadic shells $S_i$ of Lemma~\ref{prp:dyadic}, with the mass bound $V\ge1$
of Lemma~\ref{prp:normW}(b), $e=t^{-1}$, and $r'=s_0-10^8n^4$; then $9\cdot10^7n^4+r'\le s_0=s$, and if $e<1$ the
gain $Ve^{s}=V/t^{s_0}$ is our hypothesis. So
$\pfrkap_j\int f\,d\pfrM\le V\,t^{-r'}(L^2)^{4\cdot10^8n^4}\mathsf N(W^\flat)$. By Lemma~\ref{prp:normW}(a),
$\#I\cdot\mathsf c\cdot\mathsf N(W^\flat)\le\frac{\prpvolU k}{\prpvolU r}\pi^{2(k-5)\varrho}$. Hence the left side is at most
\[
\frac{\prpvolU k}{\prpvolU r}\,\pi^{k+2(k-5)\varrho}\rho^{24k}\rho^{9999n}L^{8\cdot10^8n^4}\frac{t^{10^8n^4}}{t^{s_0}}
\le\pbmTh\,L^{10^9n^4}\rho^{24k}\frac{t^{10^8n^4}}{t^{s_0}},
\]
since $k+2(k-5)\varrho\le(2\varrho+1)k$, $\rho\le L$ and $9999n+8\cdot10^8n^4\le10^9n^4$.
\end{proof}

\begin{lemma}\label{prp:arith}
Assume \textup{(H0)}. \textup{(a)} If $0<\nu$, $1\le\rho^{24}\nu$, then $(\pi^2/\nu^2)^{k-5}\le\bigl(V/(\rho^{24}\nu)^{s_0}\bigr)^2$.
\textup{(b)} Let $\nu>0$, $L\ge\rho$, $H\ge0$ with $L^{n^6}\le\rho^{n^6}H$, and $t=\max(1,\rho^{24}\nu)$. Then
$\pbmTh L^{10^9n^4}\rho^{24k}t^{10^8n^4}/t^{s_0}\le\pbmTh\,b^{2000n^9}\pfrell_{k-n^7}(\nu^{-1})\,H$.
\end{lemma}
\begin{proof}
(a) $(\pi/\nu)^{k-5}=(\pi\rho^{24})^{k-5}/(\rho^{24}\nu)^{k-5}\le(\pi\rho^{24})^k\rho^{9999n}/(\rho^{24}\nu)^{s_0}$, as
$\pi\rho^{24}\ge1$, $\rho^{24}\nu\ge1$ and $s_0\le k-5$.

(b) Write $k=n^7+K$ ($K\ge0$ by (H0)). As $n>10^8$, $10^9n^4\le n^6$, so $L^{10^9n^4}\le L^{n^6}\le\rho^{n^6}H$; and
$s_0=k-3n-6\ge10^8n^4+K+6$ (as $n^7-3n-12\ge10^8n^4$), so $t^{10^8n^4}/t^{s_0}\le t^{-K-6}$. Hence the left side
is at most $\pbmTh\,\rho^{n^6+24n^7}H\,(\rho^{24}/t)^Kt^{-6}$. Now $\rho^{24}/t\le\nu^{-1}$ and $t^{-1}\le\min(1,\nu^{-1})$,
so $(\rho^{24}/t)^Kt^{-6}\le\pfrell_K(\nu^{-1})$; and $\rho^{n^6+24n^7}=b^{(6n+2)(n^6+24n^7)}\le b^{200n^8}\le b^{2000n^9}$.
\end{proof}

\begin{proof}[Proof of Proposition~\ref{pfr:formbound}]
Let $\omega$ be one of the two weights; $|\omega|\le1$. If $|P_l(p)|>b^{60n^2}$ or $|P_l(q)|>b^{60n^2}$ for some
$l\le i$, then, as $\rho^{9n}=b^{(6n+2)9n}\le b^{60n^2}$, Lemma~\ref{prp:headsupport} gives $\pfrF^+_i=0$, so
$\pfrI_i(p,q;\omega)=0$ by Lemma~\ref{pfr:L1}. Otherwise let $L$ be given by Lemma~\ref{prp:exgauge} with $K=n^6$, so that
$p,q$ are heads of length $i$ for $\Gamma_L$ and $L^{n^6}\le\rho^{n^6}H$ with $H=\Lambda^\infty_i(p)\Lambda^\infty_i(q)$.
Let $f(x)=\pfrF_i(x;\omega)\le\pfrF^+_i(x)$ (Lemma~\ref{pfr:L1}) and $t=\max(1,\rho^{24}|\nu|)$. If $t>1$, then
$t=\rho^{24}|\nu|$, and Theorem~\ref{pfr:gain-char} (for the character weight, with $\xi=\eta=\nu$) or
Theorem~\ref{pfr:gain-heads} (for $\omega=1$, $\nu=p_{i-1}^{-1}-q_{i-1}^{-1}$, $j=i$) gives, at $\C$ where
$P/\operatorname{mod}(\nu)=\pi^2/|\nu|^2$ and $\operatorname{rank}X'=k-5$,
$f(x)^2\le(\pi^2/|\nu|^2)^{k-5}\mathsf Q_{\Gamma_L}(x)\le(V/t^{s_0})^2\mathsf Q_{\Gamma_L}(x)$ by
Lemma~\ref{prp:arith}(a). By Lemma~\ref{pfr:L1} (the integrability is Proposition~\ref{pfr:conv}, proved below
without using this proposition), Proposition~\ref{prp:sum} and Lemma~\ref{prp:arith}(b),
\[
\#I\cdot\mathsf c\,|\pfrI_i(p,q;\omega)|\le\#I\cdot\mathsf c\,\pfrkap_i\int f\,d\pfrM\le\pbmTh\,b^{2000n^9}
\pfrell_{k-n^7}(|\nu|^{-1})\,\Lambda^\infty_i(p)\Lambda^\infty_i(q).\qedhere
\]
\end{proof}

\begin{proof}[Proof of Proposition~\ref{pfr:conv}]
The support claim is Lemma~\ref{prp:headsupport} (and Lemma~\ref{pfr:L1} for $\pfrI_j$). For the finiteness,
let $L$ be as in Lemma~\ref{prp:exgauge} (any $K$), and apply Proposition~\ref{prp:sum} with $t=1$ and $f=\pfrF^+_j$
(the gain hypothesis is void): $\#I\cdot\mathsf c\,\pfrkap_j(p,q)\int\pfrF^+_j\,d\pfrM<\infty$, and $\mathsf c>0$,
$\pfrkap_j(p,q)>0$.
\end{proof}

\begin{proof}[Proof of Proposition~\ref{pfr:remote}]
Let $L=\rho$, $t=b^{6r}\ge1$, $T=t\rho^{11n}$, so that $\mathcal B$ is the set of the $x$ all of whose entries of
$x-1$ have modulus $\le T^{99n}$. At floor $0$ all sequences are heads of length $0$. Let
$f=\mathbf 1_{\pfrG\smallsetminus\mathcal B}\,\pfrF^+_0\le\pfrF^+_0$. For $x\notin\mathcal B$, Lemma~\ref{pfr:deepreal}
with $\tau'=t^{-1}$ (indeed $T\ge\rho^{11n}$ and $T^{-1}\le t^{-1}/2$) shows that $x$ has a deep vector of depth
$t^2$, and Theorem~\ref{pfr:deep} with $\varsigma=t^{-2}$ and the mass bound $V$ gives
$f(x)^2\le t^{-2s_0}V^2\mathsf Q_{\Gamma_\rho}(x)$; for $x\in\mathcal B$, $f(x)=0$. Proposition~\ref{prp:sum} and
$\#I\ge1$, $\pfrkap_0=1$ give
$\mathsf c\int_{\pfrG\smallsetminus\mathcal B}\pfrF^+_0\,d\pfrM\le\pbmTh\rho^{10^9n^4}\rho^{24k}t^{10^8n^4}/t^{s_0}$, which
is the claim.
\end{proof}

\subsection{The top floor: proof of Proposition~\ref{pfr:top}}\label{prp:ss-top}

Write $J(\gamma)=\int_\pfrG W^\flat(y)\otimes\overline{W^\flat(\gamma^{-1}y)}\,dy$, so that
$\pfrN_\sigma(x)=N(x)J(w_x\sigma)$ with $N(x)=\frac{\sqrt{c_G}}{\#I}\bigl((|x_\varrho|^2-1)^+\bigr)^{(k-1)/2}\prod_{a<\varrho}
|x_a|^{d_a}$, $d_a=k+1+2(\varrho-a-1)$. (For $|x_\varrho|\le1$ the factor $(|x_\varrho|^2-1)^{(k-1)/2}$ of
\S\ref{pfr:ss-realdata} is $0$.) Let $\prpdens(c)=\frac{2\pi^k}{(k-1)!}((|c|^2-1)^+)^{k-1}$, the
density of $\theta_\varrho$, and recall $c_G=\frac{\prpvolU k}{\prpvolU r}\frac{2\pi^k}{(k-1)!}$
(Remark~\ref{pfr:rem-match}).

\begin{lemma}[Compact supports]\label{prp:support}
Let $x^0$ be a sequence of non-zero complex numbers. There are a neighbourhood $\mathcal N$ of $x^0$ and a
compact set $C\subseteq\pfrG$ such that for $x\in\mathcal N$, $\sigma\in\pfrG^{(\varrho)}$ and $y\in\pfrG$:
if $W^\flat(y)\ne0$ and $W^\flat((w_x\sigma)^{-1}y)\ne0$, then $y\in C$.
\end{lemma}
\begin{proof}
Let $\varepsilon=\frac12\min_{l\le\varrho}|P_l(x^0)|$ and $\mathcal N$ the set of the $x$ with $|P_l(x)|>\varepsilon$ for
$l\le\varrho$. Let $\gamma=w_x\sigma$; the leading block of order $\varrho$ of $Z(\gamma)$ is that of $Z(w_x)$ ($\sigma$ fixes
$e_0,\dots,e_{\varrho-1}$), so $|\delta_l(\gamma)|=|P_l(x)|>\varepsilon$ for $l\le\varrho$. By (H1) there are $u,u'\in\pfrU$
with $k_1=uy\in\pfrK$, $k_2=u'\gamma^{-1}y\in\pfrK$; then $\kappa:=k_1k_2^{-1}=u\,w_x\sigma\,u'^{-1}$ lies in the compact
set $C_0=\{\kappa\in\pfrK\pfrK^{-1}:|\delta_l(\kappa)|\ge\varepsilon,\ l\le\varrho\}$ ($\delta_l$ is invariant under
$\gamma\mapsto u\gamma u'^{-1}$, Lemma~\ref{prp:minor}'s proof). $C_0$ is contained in the open set where the $\delta_l$,
$l\le\varrho$, are invertible, on which Proposition~\ref{pbm:cell} decomposes every element uniquely as
$u(\kappa)w_{p}\sigma'u'(\kappa)$ with $u(\cdot)$ continuous. By uniqueness $u=u(\kappa)$, which lies in the compact
set $u(C_0)$; so $y=u^{-1}k_1\in C:=u(C_0)^{-1}\pfrK$.
\end{proof}

\begin{lemma}[Descent on the top floor]\label{prp:topdescent}
Let $\varphi\colon\C\to\C$ be measurable with $|\varphi|\le1$, let $p,q$ be sequences of units with
$\int\pfrF^+_\varrho\,d\pfrM<\infty$, $w=w_{\varrho,p}$, and $\Pi$ as in \S\ref{pfr:ss-descent} with $a=\varrho$. Then
$c\mapsto\prpdens(c)\varphi(Z_{zz}(\sigma_c))\allowbreak \iint\Pi(x;w\sigma_c;g')\,d\pfrM(x)\,d\pfrM(g')$ is integrable on $\C$, and
$\pfrkap_\varrho(p,q)$ times its integral is $\pfrI_\varrho(p,q;\varphi\circ Z_{zz})$.
\end{lemma}
\begin{proof}
This is the proof of Lemma~\ref{pfr:descent} (\S\ref{pfr:ss-descent}) with $a=\varrho$ and the following changes. Fact
(iv) is replaced by the definition of $\theta_\varrho$: it is the image of $\prpdens(c)\,dc\times\mu_\pfrH$ under
$(c,h)\mapsto\sigma_c^{-1}h$, with $\mu_\pfrH$ carried by $\pfrH=\pfrR^{(\varrho+1)}\cap\pfrG^{(\varrho)}$ (by (N1), as
$\pfrN_\varrho=1$). Fact (iii) holds for $t\in\pfrH$: $tz=z$, so $Z_{zz}(tyt')=Z_{zz}(y)$, and the correlations from floor
$\varrho+1$ are averages over $\mu_\pfrH$, which is right invariant under $\pfrH$, and $d=w_{\varrho,q}^{-1}w$ commutes with
$\pfrH$. With $F(x,g,\sigma,\tau)=\varphi(Z_{zz}(\sigma\tau^{-1}))\Pi(x;\sigma w\tau^{-1};\sigma g)$ on
$\pfrM\otimes\pfrm_\varrho\otimes\theta_\varrho\otimes\theta_\varrho$, $\int|F|\le\int\pfrF^+_\varrho\,d\pfrM<\infty$, and
$\pfrkap_\varrho\int F=\pfrI_\varrho(p,q;\varphi\circ Z_{zz})$. For $\sigma\in\pfrG^{(\varrho)}$, Lemma~\ref{pfr:floorright}
with $\xi=\sigma$ and the description of $\theta_\varrho$ give $\int f(F(x,g,\sigma,\tau))\,d\theta_\varrho(\tau)=\iint
f(F(x,g,\sigma,\sigma_c^{-1}h\sigma))\,\prpdens(c)\,dc\,d\mu_\pfrH(h)$; and $\sigma w(\sigma_c^{-1}h\sigma)^{-1}=wh^{-1}\sigma_c$,
$Z_{zz}(h^{-1}\sigma_c)=Z_{zz}(\sigma_c)$, $\Pi(x;wh^{-1}\sigma_c;\sigma g)=\Pi(x;w\sigma_c;h\sigma g)$. So
$F(x,g,\sigma,\sigma_c^{-1}h\sigma)=E(x,h\sigma g,c)$ with $E(x,g',c)=\varphi(Z_{zz}(\sigma_c))\Pi(x;w\sigma_c;g')$. When
$\sigma\sim\theta_\varrho$, $g\sim\pfrm_\varrho$ and $h\sim\mu_\pfrH$, the point $h\sigma g$ has law
$\mu_\pfrH*\theta_\varrho*\pfrm_\varrho=\pfrM$ (the factors $\eta_\varrho,\eta_{\varrho+1}$ are trivial and
$\theta_{\varrho+1}=\mu_\pfrH$). Hence $F$ and $E$ (the latter on $\pfrM\otimes\pfrM\otimes\prpdens(c)dc$) have the same
law, and Fubini's theorem concludes as in Lemma~\ref{pfr:descent}.
\end{proof}

\begin{proof}[Proof of Proposition~\ref{pfr:top}]
\emph{Continuity.} Let $x^0$, $\mathcal N$, $C$ be as in Lemma~\ref{prp:support}. On $\mathcal N$, $x\mapsto N(x)$,
$x\mapsto w_x$ and $x\mapsto\sigma_x$ are continuous (the boost depends continuously on its parameter
$x_\varrho/\min(1,|x_\varrho|)$), and $\pfrX(x)=N(x)\int_CW^\flat(y)\otimes\overline{W^\flat((w_x\sigma_x)^{-1}y)}\,dy$, with a
jointly continuous integrand bounded by $(\sup\|W^\flat\|)^2$ on the compact set $C$. By dominated convergence $\pfrX$ is
continuous at $x^0$.

\emph{The equality.} Let $\sigma\in\pfrG^{(\varrho)}$ and $c=h(\sigma z,z)=x_\varrho=x'_\varrho$. Write $\sigma z=-cz+v$ with $v$ in the
span $X_0$ of the $x_i$ (the $z$-coefficient of a vector $w$ is $-h(w,z)$, since $h(z,z)=-1$); $h(\sigma z,\sigma z)=-1$ gives $|v|^2=|c|^2-1$. If $|c|\le1$ both sides vanish ($N(x)=N(x')=0$).
If $|c|>1$, also $\sigma_cz=-cz+v_c$ with $|v_c|=|v|>0$, and there is $h_1\in\pfrH=\auU(X_0)$ with $h_1v_c=v$; then
$h_1\sigma_cz=\sigma z$, so $h_2:=\sigma_c^{-1}h_1^{-1}\sigma$ fixes $z$ and the $e_a,e_a^*$, i.e. $h_2\in\pfrH$, and
$\sigma=h_1\sigma_ch_2$. Since $h_1$ commutes with $w_x$, $W^\flat(hy)=\pi(h)W^\flat(y)$ for $h\in\pfrH$, and $dy$ is left
invariant, the substitution $y\mapsto h_1y$ gives $J(w_x\sigma)=(\pi(h_1)\otimes\overline{\pi(h_2)^{-1}})\,J(w_x\sigma_c)$,
with a unitary operator that does not depend on $x$. Hence
$\pfrip{\pfrN_\sigma(x)}{\pfrN_\sigma(x')}=\pfrip{\pfrN_{\sigma_c}(x)}{\pfrN_{\sigma_c}(x')}=\pfrip{\pfrX(x)}{\pfrX(x')}$.

\emph{The descent.} Let $\varphi(c)=e^{\xi c-\overline{\xi c}}$. By Proposition~\ref{pfr:conv},
$\int\pfrF^+_\varrho\,d\pfrM<\infty$, so Lemma~\ref{prp:topdescent} applies. We identify its integrand. The correlations
from floor $\varrho+1$ are averages over $\mu_\pfrH$, and $\pfrH$ acts unitarily on the values of $W^\flat$ and commutes
with $d$; so $\mathsf k^1_{\varrho+1}(g,y)=\pfrip{W^\flat(g)}{W^\flat(y)}$ and $\mathsf k^d_{\varrho+1}(g,y)=
\pfrip{W^\flat(g)}{W^\flat(dy)}$, and, as $\sigma_c$ commutes with $w$ and $d\sigma_c^{-1}w^{-1}=\sigma_c^{-1}w_{\varrho,q}^{-1}$,
\[
\Pi(x;w\sigma_c;g')=\pfrip{W^\flat(g')}{W^\flat(g'x)}\,\overline{\pfrip{W^\flat((w_{\varrho,p}\sigma_c)^{-1}g')}{W^\flat((w_{\varrho,q}
\sigma_c)^{-1}g'x)}}.
\]
With $\pfrM=\lambda\,dx$ and the measure-preserving map $(g',x)\mapsto(g',g'x)$ of $\pfrG\times\pfrG$, and by
Lemma~\ref{prp:support} (compact supports, so everything is integrable), expanding the inner products coordinatewise,
$\iint\Pi\,d\pfrM\,d\pfrM=\lambda^2\pfrip{J(w_{\varrho,p}\sigma_c)}{J(w_{\varrho,q}\sigma_c)}$. For $|c|>1$,
$Z_{zz}(\sigma_c)=c$ and
\[
N(p^c)N(q^c)=\frac{c_G}{(\#I)^2}\bigl((|c|^2-1)^+\bigr)^{k-1}\pfrkap_\varrho(p,q)=\mathsf c\,\lambda^2\,\prpdens(c)\,\pfrkap_\varrho(p,q);
\]
for $|c|\le1$ both sides are $0$. Hence $\varphi(c)\pfrip{\pfrX(p^c)}{\pfrX(q^c)}=\mathsf c\,\pfrkap_\varrho(p,q)\,
\prpdens(c)\varphi(Z_{zz}(\sigma_c))\iint\Pi(x;w\sigma_c;g')\,d\pfrM\,d\pfrM$ for all $c\ne0$, and
Lemma~\ref{prp:topdescent} gives the integrability and the value $\mathsf c\,\pfrI_\varrho(p,q;\varphi\circ Z_{zz})$.
\end{proof}

\subsection{The hypotheses (H1)--(H5) for the Bessel functions of Theorem~\ref{pbe:bessel}}\label{prp:ss-H}

Let $W=W_\tau$, $W^\flat=W_\tau^\flat$ be the functions of Proposition~\ref{pbe:local} and Theorem~\ref{pbe:bessel},
$\mathbb W=\pbW_\tau$, $\#I=d_\tau=\dim\pbW_\tau$, $f_W=f$ (the exponent of Proposition~\ref{pbe:local}(e)). They are
continuous Bessel functions of type $\mathbb W$ by (B1), and $\mathbb W$ is irreducible under $\pfrH$
(\S\ref{pbe:ss-real}).

\emph{(H1).} $\int\|W\|^2=d_\tau$ is the normalisation of Proposition~\ref{pbe:local}; $\int\|W^\flat\|^2\le d_\tau$
and $\int\|W^\flat-W\|^2\le b^{-2r^3f}d_\tau$ are Proposition~\ref{pbe:local}(e). If $W^\flat(y)\ne0$ then
$\inf_{s\in R_\tau}(\|sy\|+1)\le\rho$ (Proposition~\ref{pbe:local}(e)), so some $s=u'h\in R_\tau=\pfrU\pfrH$ has
$\|sy\|<\rho$; as $h$ is a Euclidean isometry normalising $\pfrU$, $\|uy\|=\|sy\|$ with $u=h^{-1}u'h\in\pfrU$.

\emph{(H2), with equality.} Let $\mathsf k(g)=\pfrip W{R(g)W}=\int_{\pfrU\backslash\pfrG}\pfrip{W(y)}{W(yg)}\,dy$. As
$\|W\|$ is bounded (Proposition~\ref{pbe:local}(a)) and integrable on $\pfrU\backslash\pfrG$ (B2), $\mathsf k$ is
bounded and measurable; by the right invariance of the measure of $\pfrU\backslash\pfrG$,
$\int\pfrip{W(ys)}{W(yt)}dy=\mathsf k(s^{-1}t)$, and $\mathsf k(g^{-1})=\overline{\mathsf k(g)}$. By the reproduction
property (B7), $W(y)=\int e(a)W(ya)\,da$ with $e=e_\tau=\pbdeg_\tau\overline{c_\tau}$ integrable; inserting it in the
second factor of $\mathsf k(g)=\int\pfrip{W(yg^{-1})}{W(y)}dy$ (Fubini: $e\in L^1$, $W$ bounded, $\|W\|$ integrable)
gives $\mathsf k(g)=\int e(a)\mathsf k(ga)\,da$; conjugating, $\mathsf k(g)=\int\overline{e(a)}\mathsf k(a^{-1}g)\,da$.
Hence, substituting $t'=s^{-1}gt$,
\[
\mathsf k(g)=\iint\overline{e(s)}e(t)\mathsf k(s^{-1}gt)\,dt\,ds=\int\mathsf k(t')\Bigl(\int\overline{e(s)}\,e(g^{-1}st')\,ds\Bigr)dt'.
\]
The inner integral is $\pbdeg_\tau^2\int c_\tau(s)\overline{c_\tau(g^{-1}st')}\,ds$; with $a=s^{-1}$ and
$c_\tau(t'^{-1}ag)=\pfrip{\pi(t')v_0}{\pi(a)\pi(g)v_0}$ it equals $\pbdeg_\tau^2\int\overline{c_\tau(a)}\,c_\tau(t'^{-1}ag)\,da$,
and Schur's orthogonality relations (\S\ref{pds:ss-ds1}(ii), for the four vectors $\pi(t')v_0,\pi(g)v_0,v_0,v_0$)
give $\pbdeg_\tau\,c_\tau(t'^{-1})\,c_\tau(g)=\pbdeg_\tau\overline{c_\tau(t')}c_\tau(g)$. So
$\mathsf k(g)=c_\tau(g)\int e(t')\mathsf k(t')\,dt'=c_\tau(g)\mathsf k(1)=d_\tau\,c_\tau(g)$, and
$\int_\pfrG|\mathsf k|^2=d_\tau^2\int|c_\tau|^2=d_\tau^2/\pbdeg_\tau$ by \textup{(DS1)(ii)}. This is (H2) with equality.

\emph{(H3).} The entries of $\varsigma_\tau$ are $\pm(\frac k2-1),\dots$ (modulus $\le r\le b$) and
$-\kappa(\tau)-j\Delta$ with $j\Delta\le m\Delta\le\kappa(\tau)\le b$ (Lemma~\ref{an:lem-sizes}\,(S5) and (S7)(g)); so $|\varsigma_{\tau,i}|\le2b$, every factor $(\varsigma_{\tau,i}-\varsigma_{\tau,j})/(j-i)$ is
$\le4b\le b^3$, and $\pbdeg_\tau\le b^{3r^2/2}$. Next, $\prpvolU N=\prod_{j=1}^N\frac{2\pi^j}{(j-1)!}$, so
$\prpvolU k/\prpvolU r=\prod_{j=k+1}^r\frac{(j-1)!}{2\pi^j}\le(r^r)^{r-k}$; with $r-k=2\varrho+1\le r/2$ (as $k\ge r/2$,
\S\ref{pbe:ss-real} sizes) and $r\le b$, this is $\le b^{r^2/2}$; and $\pi^{(2\varrho+1)k}\le4^{(r-k)k}\le4^{r^2/4}\le
b^{r^2/2}$. Hence $\pbdeg_\tau\pbmTh\le b^{3r^2/2+r^2/2+r^2/2}\le b^{4r^2}$.

\emph{(H4)} is \textup{(DS3)} for the Haar measure of \textup{(DS1)}: since $\|g\|\le T$ implies $(32/\|g\|)^{9r^2}\ge
(32/T)^{9r^2}$, $\pbdeg_\tau\pbvol\{\|g\|\le T\}\le(T/32)^{9r^2}b^{99r^2}\le T^{9r^2}b^{99r^2}$.

\emph{(H5).} $f\ge4000\,n$, Lemma~\ref{an:lem-sizes}\,(S7)(b). \emph{(H0)} is
Remark~\ref{pfr:rem-match}\,(b).

\subsection{What Theorem~\ref{pfr:R} now rests on}\label{prp:ss-R}

With the proofs of this section, Theorem~\ref{pfr:R} (through \S\ref{pfr:ss-Rproof}, Proposition~\ref{pfr:mass} and
the propositions of \S\ref{pfr:ss-realforms}) depends only on: the shared Theorems~\ref{pfr:oneframe},
\ref{pfr:gain-char} and~\ref{pfr:gain-heads}; the results of \S\ref{pfr:sec} with their proofs (and the standard
facts used there: Witt's theorem, measurable sections of orbit maps); Proposition~\ref{pbm:cell} (in
Lemma~\ref{prp:support}); Lemmas~\ref{pbo:GF} and~\ref{pbo:append} (in Lemma~\ref{prp:fibre}); and the hypotheses
(H0)--(H5). For the Bessel functions of Theorem~\ref{pbe:bessel}, (H0)--(H5) hold by \S\ref{prp:ss-H}, using
Proposition~\ref{pbe:local} and Theorem~\ref{pbe:bessel}\,(B1), (B2), (B7) (resting on \textup{(DS1)}--\textup{(DS4)},
Theorem~\ref{pbe:ds}, with \textup{(DS4)} = Theorem~\ref{pbp:main}), the Schur relations of \S\ref{pds:ss-ds1},
\textup{(DS3)}, and Lemma~\ref{an:lem-sizes}. Schur's lemma, the uniqueness of Haar measure and the polar formula are
used as standard facts.


%% ===== proofs/B-frames-finite.tex =====
% proofs/B-frames-finite.tex
% The local systems at the finite places and Siegel's step: parts (F) and (S) of
% Theorem pbm:inputs. Labels of this file start with pff:.
\providecommand{\auO}{\mathcal O}
\providecommand{\auPl}{\Phi}
\providecommand{\auU}{\mathrm U}
\providecommand{\auGL}{\mathrm{GL}}
\providecommand{\auA}{\mathbb A}
\providecommand{\auF}{\mathbb F}
\providecommand{\aurd}{\mathrm{rd}_E}
\providecommand{\aubar}[1]{\overline{#1}}
\providecommand{\panot}[1]{\par\smallskip{\raggedright\noindent\textbf{[Not proved here:} #1\textbf{]}\par}\smallskip}
\providecommand{\pbR}{\overline R_0}
\providecommand{\pbUb}{\overline U}
\providecommand{\pbHb}{\overline H}
\providecommand{\pbvol}{\operatorname{vol}}
\providecommand{\pbFix}{\operatorname{Fix}}
\providecommand{\pbDmain}{\mathcal D_{\mathrm{main}}}
\providecommand{\pbmGf}{G^{\flat}}
\providecommand{\pbmXp}{X^{+}}
\providecommand{\auG}{\overline G}
\providecommand{\pbmOrb}{\mathcal O}
\providecommand{\pbmT}{\mathcal T}
\providecommand{\pbmRT}{R_T}
\providecommand{\pbmE}{\mathsf E}
\providecommand{\pbmn}{\mathsf n}
\providecommand{\pbmden}{\mathsf{den}}
\providecommand{\pbmvol}{\mathsf v}
\providecommand{\pbmaS}{\mathsf a_*}
\providecommand{\pbmrS}{\mathsf r_*}
\providecommand{\pbmcG}{c_G}
% macros of this file (prefix \pff)
\providecommand{\pffET}{E_T}
\providecommand{\pffOT}{\auO_T}
\providecommand{\pfgOh}{\auO_{\mathrm h}}
\providecommand{\pfgQ}{\mathsf Q_N}
\providecommand{\pffG}{G_T}
\providecommand{\pffK}{K_T}
\providecommand{\pffKM}{K_T(M)}
\providecommand{\pffch}{\chi_M}
\providecommand{\pffW}{W_Z}
\providecommand{\pffm}{\varkappa}
\providecommand{\pffr}{\rho}
\providecommand{\pffl}{\mathsf f}
\providecommand{\pffI}{\mathcal I}
\providecommand{\pffZ}{\mathsf Z}
\providecommand{\pfrF}{\mathsf F}
\providecommand{\pfrB}{\mathsf B}
\providecommand{\pfrM}{\mathsf M}
\providecommand{\pfrm}{\mathsf m}
\providecommand{\pfrN}{\mathsf N}
\providecommand{\pfrkap}{\kappa}
\providecommand{\pfrH}{\mathsf H}
\providecommand{\pbmKl}{\mathsf{Kl}}
\providecommand{\pbmass}{\operatorname{mass}}
\providecommand{\pfrmod}{\operatorname{mod}}
\providecommand{\pfrI}{\mathcal I}
\providecommand{\pfrip}[2]{\langle#1,#2\rangle}

\section{The local systems at the finite places and Siegel's step}\label{pff:sec}

This section concerns the parts \textup{(F)} and \textup{(S)} of Theorem~\ref{pbm:inputs}. The
part \textup{(R)}, and the theory of forms on pairs of heads that is common to the real and the
finite places (floor measures, correlations, gauges and sizes, depth of frame vectors, the sum
over dyadic shells of positions), are treated in \S\ref{pfr:sec}; from there we quote, with
their labels, the statements that we need, and we do not prove them again.

\emph{The results of this section.} We list the results by the label of the statement.
\begin{itemize}
\item[(i)] \emph{The exact statements} (\S\ref{pff:ss-statements}, Theorems~\ref{pff:F}
and~\ref{pff:S}). We fix the forms $J^T_i$ at the finite places, and we check that the
constants of \textup{(F)} and \textup{(S)} of Theorem~\ref{pbm:inputs} are exactly those that
Theorem~\ref{pbm:L2} and Lemma~\ref{pbm:const} use: the mass $(1+b^{-6})b^{-6k\#\mathsf A}$, the
weights $\Lambda^T$, the decay $\beta^T$, the place $(E_T,dx,\psi_T,w\mapsto M^{-1}w^{-1})$, and in
\textup{(S)} the factor $(1+b^{-6})\pbmaS$ against $\pbmrS^{\,2}$. One reading has to be made precise in
Theorem~\ref{pbm:inputs}: in \textup{(F)} the forms depend on a group $\pffZ$ that is
chosen together with them, and \textup{(S)} is stated with the same group. The requirement that $T$
contain the places at which the chosen basis of the lattice is not integral is already part of the
hypothesis of \textup{(F)}, as these places lie above $N$. This does not affect the
deduction of Theorem~\ref{pbm:L2} (Remark~\ref{pff:rem-L2}).
\item[(ii)] \emph{Theorem~\ref{pff:F} (part \textup{(F)}) is proved}, using the statements of the
common theory of \S\ref{pfr:ss-shared} (Lemmas~\ref{pfr:L1}, \ref{pfr:descent},
Theorems~\ref{pfr:deep}, \ref{pfr:gain-char}, \ref{pfr:gain-heads}, \ref{pfr:shells}; proved in
\S\S\ref{pfr:sec}, \ref{pfs:sec}) and Lemma~\ref{pff:twogauges}, proved in \S\ref{pfg:ss-two}. What is specific to the finite
places is written here: the basis at $T$, the gauges, the depth ideals, the floor masses and the volume of
$\pffK$, the dyadic shells, the mass (with the gain $b^{-6k}$ for each fixed vector,
Theorem~\ref{pff:fixed}), the decay (Theorem~\ref{pff:decay}), and the density at the last entry
(Proposition~\ref{pff:top}).
\item[(iii)] \emph{Theorem~\ref{pff:S} (Siegel's step) is proved} (\S\ref{pff:ss-siegel}), modulo Weil's Siegel--Weil theorem, from which Theorem~\ref{pff:siegelformula} is derived in \S\ref{pfg:ss-siegel}, and the top forms at the real places of \S\ref{pfr:sec}.
The lemmas used here and proved in \S\ref{pfg:sec} are listed in \S\ref{pff:ss-pending}.
\end{itemize}

\subsection{The setting at a finite set of places}\label{pff:ss-setting}

We use the notation of \S\ref{pb:ss-setting} and \S\ref{pbm:sec}. Hermitian forms follow convention (Cv1) of
\S\ref{pb:ss-setting} in this section, as in \S\ref{pfr:sec} (linear in the first variable); \S\ref{pfg:sec} writes them
in the opposite convention, as declared there. Let $T$ be a finite set of
finite places of $E$, stable under conjugation, containing the places above $MN$. Let
$\pffET=\prod_{w\in T}E_w$, with ring of integers $\pffOT=\prod_{w\in T}\auO_w$, the Haar measure
$dx$ for which $\pffOT$ has volume $1$, and the module $|\xi|_T=\prod_{w\in T}|\xi|_w$; for a unit
$c$ of $E_T$ the map $x\mapsto cx$ multiplies $dx$ by $|c|_T$. The involution of $E/E^+$ acts on
$\pffET$, and $\psi_T(x)=\psi_{E,f}(-x/\varpi)$ is the character of \S\ref{pbm:ss-sieve}; it is
trivial on $\pffOT$ (as $\varpi$ generates the different), and $\psi_T(x+\aubar x)=1$ for all
$x\in\pffET$: with $y=-x/\varpi$ and $\aubar\varpi=-\varpi$ we have $-(x+\aubar x)/\varpi=y-\aubar
y$, whose trace to $E^+_T$ vanishes, and $\psi_E$ factors through the trace.

\emph{The lattice.} Let $Y_0$ be the hermitian space of \S\ref{pb:ss-setting} (dimension
$r=2\varrho+k+1$, with the hyperbolic pairs $e_a,e_a^*$ and the middle space $\pbmXp=X_0\perp
Ez$), and $L\subset Y_0$ the lattice of \S\ref{pb:ss-setting}, whose stabiliser is $K$; it is
unimodular, and we fix a basis of it over $\auO_E$ (the \emph{lattice basis}; it differs from the
basis $e_a,x_i,z,e_a^*$ of $Y_0$ only in the plane of $x_{k_0}$ and $z$, by a matrix that is integral
with integral inverse at the places not above $N$). Let $V_T=Y_0\otimes_EE_T$,
$L_T=L\otimes\pffOT$, $\pffG=\auU(V_T)$, and
\[
\pffK=\{g\in\pffG:gL_T=L_T\},\qquad\pffKM=\{g\in\pffK:(g-1)L_T\subseteq ML_T\}.
\]
Both are compact open subgroups, $\pffKM$ is normal in $\pffK$, and $\pffKM$ is the component at $T$
of the group $J$ (\S\ref{pb:ss-setting}; $J$ is the kernel of the reduction of $K$ modulo $M$),
as $T$ contains the places above $M$; the component of $K$ at $T$ is $\pffK$. We write $U_T$,
$H_T$, $\pbmGf_T$ for the groups $U$, $H$, $\pbmGf$ over $\pffET$, and $\sigma(u)$ for the sum of
the entries just above the diagonal blocks of $u\in U_T$ (in the frame
$f_0,\dots,f_\varrho$ of \S\ref{pbm:ss-cell}), so that $u\mapsto\psi_T(M^{-1}\sigma(u))$ is the
component at $T$ of the character $\psi_f$ of \S\ref{pbm:ss-cell}.

\emph{The function $\pffch$.} For $u\in U_T\cap\pffKM$ we have $\sigma(u)\in M\pffOT$, so
$\psi_T(M^{-1}\sigma(u))=1$. Hence
\begin{equation}\label{pff:eq-chi}
\pffch(uk)=\psi_T\bigl(M^{-1}\sigma(u)\bigr)\quad(u\in U_T,\ k\in\pffKM),\qquad\pffch=0\text{ off }U_T\pffKM,
\end{equation}
is well defined; $\pffch$ is locally constant, $|\pffch|\le1$, $\pffch(ug)=\psi_T(M^{-1}\sigma(u))
\pffch(g)$ for $u\in U_T$, and $\pffch(gk)=\pffch(g)$ for $k\in\pffKM$. It is the function $\chi_M$
of \S\ref{pbm:ss-remains}.

\emph{Intermediate groups.} An \emph{intermediate group} is a subgroup $\pffZ$ with $\pffKM\subseteq
\pffZ\subseteq\pffK$; it is compact open, normalises $\pffKM$, and has finitely many cosets
$\pffKM a$. Choosing a representative $a_c\in\pffZ$ of each coset $c$, the \emph{Bessel vector} of
$\pffZ$ is
\[
\pffW\colon\pffG\to\C^{\pffKM\backslash\pffZ},\qquad\pffW(y)_c=\pffch(ya_c^{-1}).
\]
By the right invariance of $\pffch$ under $\pffKM$ and the normality of $\pffKM$ in $\pffK$, a
change of representatives multiplies no coordinate, so $\pffW$ does not depend on the choices. It
satisfies $\pffW(uy)=\psi_T(M^{-1}\sigma(u))\pffW(y)$ for $u\in U_T$, and $\pffW(y)\ne0$ only if
$uy\in\pffK$ for some $u\in U_T$: it is a \emph{cut Bessel function} in the sense of
\S\ref{pfr:ss-setting}, for the box $\pffK$ and the level $M$.

\emph{Measures.} We fix on $\pffG$ the Haar measure $\pbvol=\pfrM$ of \S\ref{pfr:ss-setting}
(the \emph{group measure}, built floor by floor from the frame), for the place $(\pffET,dx,\psi_T)$
and the lattice $L_T$. The group $\pffG$ is unimodular, so $\pbvol$ is left and right invariant.
We use the following properties of it, all from \S\ref{pfr:sec}.
\begin{enumerate}
\item[(M1)] \emph{(Unfolding along $U_T$.)} There is a Haar measure $du$ on $U_T$ with
$\int_{U_T\cap\pffK}du=1$, and a measure $\theta$ on $\pffG$ (standing for the quotient
$U_T\backslash\pffG$), such that $\int_{\pffG}F\,d\pbvol=\int_{\pffG}\int_{U_T}F(uy)\,du\,
d\theta(y)$ for $F\ge0$ measurable. (Here $\theta=\theta_{\ge0}$ and $du=\eta_{[0,\varrho+2)}$, which is
carried by $U_T$; see \S\ref{pff:ss-large}.)
\item[(M2)] \emph{(Cosets of $\pffKM$ in $U_T$.)} For $y\in\pffG$, the set $\{u\in U_T:uy\in
\pffKM\}$ is empty or a coset of $U_T\cap\pffKM$; we write
\[
\pffl=\int_{U_T\cap\pffKM}du ,
\]
the \emph{unipotent mass}. (This is the product of the floor moduli of $M$ in
\S\ref{pfr:sec}.)
\end{enumerate}
We write $\pffm(\pffZ)$ for the \emph{mass integral} and $\pffr(\pffZ)$ for the
\emph{normalised mass} of an intermediate group:
\begin{equation}\label{pff:eq-mass}
\pffm(\pffZ)=\int_{\pffG}\Bigl|\int_{\pffZ}\pffch(yxy^{-1})\,dy\Bigr|^2dx,\qquad
\pffr(\pffZ)=\frac{\pbvol(\pffZ)\,\pbvol(\pffKM)}{\pffl^{\,2}},
\end{equation}
both integrals for $\pbvol$. The function $x\mapsto\int_{\pffZ}\pffch(yxy^{-1})dy$ is the
function $\psi^{\natural}$ of Lemma~\ref{pbm:pos}(c) for the $p$-adic group, up to the factor
$\pbvol(\pffZ)$; compare Remark~\ref{pbm:rem-E4}(c).

\subsection{The Bessel vector and the mass integral}\label{pff:ss-bessel}

Throughout, $\pffZ$ is an intermediate group and $\pffW$ its Bessel vector.

\begin{lemma}[Unfolding the Bessel vector]\label{pff:unfold}
For $x,y\in\pffG$,
\[
\pffl\cdot\langle\pffW(y),\pffW(yx)\rangle=\int_{U_T}\mathbf 1_{\pffZ}(uy)\,\pffch\bigl(uy\,x\,(uy)^{-1}
\bigr)\,du ,
\]
where $\langle a,b\rangle=\sum_c\aubar{a_c}\,b_c$.
\end{lemma}
\begin{proof}
For a coset $c$ let $S_c=\{u\in U_T:uya_c^{-1}\in\pffKM\}$, and $\alpha_c=\aubar{\pffW(y)_c}\,
\pffW(yx)_c$. We claim that for every $u\in U_T$
\begin{equation}\label{pff:eq-unf}
\mathbf 1_{\pffZ}(uy)\,\pffch\bigl(uyx(uy)^{-1}\bigr)=\sum_c\mathbf 1_{S_c}(u)\,\alpha_c .
\end{equation}
If $uy\notin\pffZ$, then $u\notin S_c$ for all $c$, since $uya_c^{-1}\in\pffKM\subseteq\pffZ$ would
give $uy\in\pffZ$. If $uy\in\pffZ$, then $uy$ lies in exactly one coset $\pffKM a_d$, so $u\in S_c$
exactly for $c=d$. Put $k=uya_d^{-1}\in\pffKM$. Then $\pffW(y)_d=\pffch(u^{-1}k)=\psi_T(M^{-1}
\sigma(u^{-1}))$, and
\[
\pffW(yx)_d=\pffch(yxa_d^{-1})=\pffch\bigl(u^{-1}\cdot uyx(uy)^{-1}\cdot k\bigr)=\psi_T(M^{-1}\sigma(
u^{-1}))\,\pffch\bigl(uyx(uy)^{-1}\bigr),
\]
by the two invariance properties of $\pffch$. As $|\psi_T|=1$, $\alpha_d=\pffch(uyx(uy)^{-1})$, which
proves \eqref{pff:eq-unf}. By (M2), $S_c$ is empty or a coset of $U_T\cap\pffKM$, of measure
$\pffl$; and if $S_c$ is empty then $\pffW(y)_c=0$ (for $\pffW(y)_c\ne0$ means $ya_c^{-1}\in
U_T\pffKM$, that is $S_c\ne\emptyset$), so $\alpha_c=0$. Integrating \eqref{pff:eq-unf} over
$U_T$ gives $\sum_c\pffl\,\alpha_c=\pffl\langle\pffW(y),\pffW(yx)\rangle$.
\end{proof}

\begin{corollary}\label{pff:k0}
Put $k_0(x)=\int\langle\pffW(y),\pffW(yx)\rangle\,d\theta(y)$ (the \emph{correlation from floor
$0$}, an absolutely convergent integral). Then for all $x\in\pffG$
\[
\int_{\pffZ}\pffch(yxy^{-1})\,dy=\pffl\cdot k_0(x),\qquad\text{hence}\qquad\pffm(\pffZ)=\pffl^{\,2}
\int_{\pffG}|k_0(x)|^2\,dx .
\]
Moreover $\pffl\,\|\pffW(y)\|^2\le1$ for all $y$, and $\pffl\int\|\pffW(y)\|^2d\theta(y)=
\pbvol(\pffZ)$.
\end{corollary}
\begin{proof}
The function $F(z)=\mathbf 1_{\pffZ}(z)\pffch(zxz^{-1})$ is bounded with compact support. By (M1)
and Fubini's theorem, $\int_{\pffZ}\pffch(yxy^{-1})dy=\int\int_{U_T}F(uy)\,du\,d\theta(y)$, and the
inner integral is $\pffl\langle\pffW(y),\pffW(yx)\rangle$ by Lemma~\ref{pff:unfold}. For $x=1$,
$\pffch(1)=1$ and Lemma~\ref{pff:unfold} gives $\pffl\|\pffW(y)\|^2=\int_{U_T}\mathbf 1_{\pffZ}(uy)\,du$.
This is at most $1$: the set $\{u:uy\in\pffZ\}$ lies in $\{u:uy\in\pffK\}$, which is empty or a coset
of $U_T\cap\pffK$, of measure $1$ by (M1). Integrating over $\theta$ and using (M1) for
$F=\mathbf 1_{\pffZ}$ gives the last identity; it also shows that the integral defining $k_0$
converges absolutely, by the Cauchy--Schwarz inequality.
\end{proof}

\begin{lemma}[The positions in $\pffKM$]\label{pff:inKM}
$\displaystyle\int_{\pffKM}\Bigl|\int_{\pffZ}\pffch(yxy^{-1})\,dy\Bigr|^2dx=\pbvol(\pffZ)^2\,
\pbvol(\pffKM)$.
\end{lemma}
\begin{proof}
For $x\in\pffKM$ and $y\in\pffZ$ we have $yxy^{-1}\in\pffKM$, as $\pffZ\subseteq\pffK$ normalises
$\pffKM$; so $\pffch(yxy^{-1})=\pffch(1)=1$ and the inner integral is $\pbvol(\pffZ)$.
\end{proof}

\begin{lemma}[Volumes]\label{pff:vol}
Assume the conditions \textup{(L1)}--\textup{(L3)} of \S\ref{pff:ss-large} below. Then:
\begin{enumerate}
\item[(a)] $\pbvol(\pffK)\le1+2^{-n}$;
\item[(b)] if the non-units of $\pffOT$ have volume at most $\frac12$, then $\pbvol(\pffK)\ge
2^{-n}$;
\item[(c)] for an intermediate group $\pffZ$, $\pbvol(\pffZ)=[\pffZ:\pffKM]\,\pbvol(\pffKM)$.
\end{enumerate}
\end{lemma}
\begin{proof}
(c) is the invariance of $\pbvol$. (a) and (b) are proved in \S\ref{pff:ss-floors}
(Lemma~\ref{pff:floormass}); in outline: in the unfolding of $\pbvol$ floor by floor, the frame vector of floor
$l$ ranges over a quadric of $\pffOT$-points of rank $2(m'+1)+\#\iota$, whose volume is at most the
ideal zeta sum $\zeta_{\pffOT}(m')=\sum_{I}[\pffOT:I]^{-m'}$ (sum over the ideals of finite index)
and at least the volume of $\pffOT^\times$; and $\zeta_{\pffOT}(m')^n\le1+2^{-n}$ is (L2).
\end{proof}

\subsection{The datum at $T$; positions outside $\pffKM$}\label{pff:ss-large}

The group measure and the common theory of \S\ref{pfr:sec} are applied at $T$ to a basis of $L_T$
adapted to the frame. We record the conditions under which the statements of \S\ref{pfr:sec} are
quoted (\emph{large lattice datum}); $m'=3\varrho+5$ and $\iota=k-6\varrho-13$.
\begin{enumerate}
\item[(L1)] $L_T$ has an $\pffOT$-basis consisting of the frame and dual frame
$e_1,\dots,e_\varrho$, the hyperbolic pair $(u,u')=(-v,z-N\tau v)$ spanning the plane $\langle v,z\rangle$
of Lemma~\ref{pff:L1} (so the frame vector $z=u'-N\tau u$ lies in this plane and has coordinate $1$ on $u'$),
$m'+1$ further hyperbolic pairs, $\iota$ orthonormal vectors, and $e_1^*,\dots,e^*_\varrho$; so
$k+1=\iota+2(m'+1)+2$ and $r=2\varrho+2+2(m'+1)+\iota$.
\item[(L2)] $\zeta_{\pffOT}(m')^n\le1+2^{-n}$, where $\zeta_{\pffOT}(s)=\sum_I[\pffOT:I]^{-s}=
\prod_{w\in T}(1-q_w^{-s})^{-1}$, the sum over the ideals of finite index; and $m'<9n$, $53n\le
\iota$.
\item[(L3)] $M\in\pffOT$, $\aubar M=M$, $\psi_T$ is trivial on $\pffOT$ and on the elements
$x+\aubar x$, and the hermitian form is unimodular on $L_T$.
\end{enumerate}
For the data, (L3) was seen in \S\ref{pff:ss-setting}. In (L2), $m'<9n$ and $53n\le\iota$ hold as
$k\ge x^5/2$ and $n\le2\sqrt x$; and, as $\zeta(s)\le1+2^{2-s}$ for $s\ge2$ and the number of places
of $E$ above a rational prime is at most $2f$,
\[
\log\zeta_{\pffOT}(m')^n\le n\cdot2f\cdot2^{2-m'}\le2^{-\varrho}\cdot fn\,2^{-2\varrho-1}\le
\log(1+2^{-n}),
\]
because $fn\le x^{13}\le2^{2\varrho-1}$ for the data ($\varrho\ge\sqrt x/2-1$ and $(\log b)^{100}\le
x$). Condition (L1) is Lemma~\ref{pff:L1}.

\subsubsection*{The finite place as a place datum}
We now place ourselves in the framework of \S\ref{pfr:ss-setting}, with the finite place
$A=\pffET$ ($\pfrmod(c)=|c|_T$, Plancherel constant $P=1$: axiom (P) of \S\ref{pfs:ss-planch} holds with $P=1$, Lemma~\ref{pfg:P}), the basis of (L1), and the place datum
$\mathfrak F_T$ whose box is $\pffK$, whose close pairs are the pairs $(\omega_1,\omega_2)$ with $\omega_1\in
L_T$ and $M^{-1}(\omega_1-\omega_2)\in L_T$, whose positions are the elements $y^{-1}uyk$
($y\in\pffG$, $u\in U_T$, $k\in\pffKM$), whose level is $M$, and whose measures are as follows.
\begin{itemize}
\item $\mu_H$ is the right Haar measure of the closed subgroup $H_T$ normalised by $\mu_H(H_T\cap\pffK)=1$. For every $g$
the $h\in H_T$ with $uhg\in\pffK$ for some $u\in U_T$ lie in one right coset of $H_T\cap\pffK$: if $u_1h_1g,u_2h_2g\in\pffK$, then
$(u_1h_1g)(u_2h_2g)^{-1}=u'h_1h_2^{-1}\in\pffK$ with $u'=u_1(h_1h_2^{-1}u_2^{-1}h_2h_1^{-1})\in U_T$ ($H_T$ normalises $U_T$), and in the
basis (L1) an element $u'h$ ($u'\in U_T$, $h\in H_T$) is block upper triangular with the diagonal blocks of $h$, as is its
inverse $h^{-1}u'^{-1}$ with the diagonal blocks of $h^{-1}$; so $u'h\in\pffK$ forces $h,h^{-1}$ integral, that is $h\in\pffK$; so the condition $\mu_H\{h:\exists u,\ uhg\in\pffK\}\le1$ holds.
\item The floor measures $\theta_l$, $l<\varrho$, are those of \S\ref{pfr:ss-setting}: $\theta_l=\lambda_l*\eta_l$ with
$\lambda_l$ the image of $|c|_T^{d_l}dc$ under $c\mapsto w_{l,c}^{-1}$, and $\eta_l$ the image of $d\omega\otimes d\beta_{\rm sk}$,
where $d\omega$ gives the lattice of $V_{l+1}$ volume $1$ and $d\beta_{\rm sk}$ gives the integral skew elements volume $1$
(Lemma~\ref{pfg:duality}(d)). By Lemma~\ref{pfr:floorquad} ($\kappa_{E_T}=1$) the image of $\theta_l$ under
$\sigma\mapsto\sigma^{-1}f_l$ is exactly the quadric measure $\mu_l$: $c_l=1$. Here $\mu_l$ is the fibre measure at $0$ of
Proposition~\ref{pfg:fibres} for the unimodular lattice of $V_l$ spanned by the basis vectors of floor $\ge l$.
\item The quadric measure of the frame vector $z$ is $\mu_\varrho=\nu^+_{-N}$, the fibre measure at $h(z,z)=-N$ of
\S\ref{pfg:ss-fibre} for the lattice $L^+$ of $V_\varrho=\pbmXp_T$; it is carried by the quadric $\pfgQ$, which is the orbit
$\{\sigma^{-1}z:\sigma\in\pffG^{(\varrho)}\}$ (Lemma~\ref{pfg:orbit}). Let $s\colon\pfgQ\to\pffG^{(\varrho)}$ be a Borel map
with $s(e)^{-1}z=e$ (Lemma~\ref{pfs:section} for $b=\varrho$). The \emph{top floor measure} is
\[
\theta_\varrho=s_*\mu_\varrho,\qquad\int F\,d\theta_\varrho=\int_{\pfgQ} F(s(e))\,d\mu_\varrho(e)\quad(F\ge0\text{ Borel}).
\]
It is finite on compact sets ($s^{-1}(C)\subseteq\{\sigma^{-1}z:\sigma\in C\}$, a compact set), and its image under
$\sigma\mapsto\sigma^{-1}z$ is $\mu_\varrho$: $c_\varrho=1$. The stabiliser of $z$ in $\pffG^{(\varrho)}$ is $H_T=\pfrR^{(\varrho+1)}\cap
\pffG^{(\varrho)}$, and $s(\sigma^{-1}z)\in H_T\sigma$. Every function integrated against $\theta_\varrho$ below is invariant under
left translation by $H_T$ (the correlations from floor $\varrho+1$ are $\mu_H$-averages, $\mu_H$ being right invariant
under $H_T$; nearness at $\varrho+1$ and the weights are $\pfrR^{(\varrho+1)}$-invariant), and for such $F$,
$\int F\,d\theta_\varrho=\int F_0\,d\mu_\varrho$ with $F(\sigma)=F_0(\sigma^{-1}z)$; so these integrals do not depend on the
choice of $s$.
\end{itemize}
\emph{The group measure is a Haar measure.} The group measure $\pfrM$ of $\mathfrak F_T$ is right invariant by
Lemma~\ref{pfr:upperinv} for $j=0$ (whose proof uses only the structure of the floors and the invariance of the $\mu_l$);
$0<\pfrM(\pffK)<\infty$ by Lemma~\ref{pff:floormass}(b),(c), whose proofs do not use that $\pfrM$ is a Haar measure; a compact set
is covered by finitely many right cosets $\pffK g$, each of measure $\pfrM(\pffK)$; so $\pfrM$ is a non-zero right invariant
Radon measure, hence a right Haar measure, hence a Haar measure ($\pffG$ is unimodular). It is the measure $\pbvol$ of
\S\ref{pff:ss-setting}. (M1) is the decomposition $\pfrM=\pfrm\ast\theta_{\ge0}$
with $\pfrm=\eta_{[0,\varrho+2)}$ carried by $U_T$, and (M2) holds with $\pffl=\prod_{l<\varrho}
\eta_l(\pfrN_l\cap\pffKM)$ (each Heisenberg coordinate lies in $M\pffOT$ or $\pffOT$).
The Bessel vector $\pffW$ is a cut Bessel function for $\mathfrak F_T$ (\S\ref{pff:ss-setting}).

\begin{lemma}[Gauges at $T$]\label{pff:gauge}
Let $\alpha\in\pffOT$ be a unit of $\pffET$ with $\aubar\alpha=\alpha$. Put $\Gamma_\alpha(i)=\{c\in
\pffET:\alpha^ic\in\pffOT\}$ and $\Lambda_\alpha=2/|\alpha|_T$. Under (L1)--(L3), $\Gamma_\alpha$ with
the loss $\Lambda_\alpha$ is a gauge for $\mathfrak F_T$ (\S\ref{pfr:ss-gauge}, axioms (G0)--(G5)). For
$\alpha=1$ the sets $\Gamma_1(i)$ are all $\pffOT$ and the loss is $2$.
\end{lemma}
\begin{proof}
This is Lemma~\ref{pfg:gauge}, where (G0)--(G5) are verified in their present form, including the clause
$\tau_0\delta^{-1}\in\Gamma(1)$ of (G1) (with $\tau_0$, $\delta$ of Lemma~\ref{pfg:psi}), the count (G4) from the
Gauss sum bound of Lemma~\ref{pfg:gauss} and Proposition~\ref{pfg:fibres}, and the count (G5) from
Lemmas~\ref{pfg:common}, \ref{pfg:Dint}, \ref{pfg:content}.
\end{proof}

\begin{lemma}[Depth at $T$]\label{pff:dideal}
Let $a<\varrho$, $x,y\in\pffG$.
\begin{enumerate}
\item[(a)] $\mathfrak d_a(y,x)=c\,\pffOT$ for a unit $c$ of $\pffET$; if $y$ is near at $a+1$, then
$c\in\pffOT$ and $\operatorname{dp}_a(y,x)=|c|_T^{-1}=[\pffOT:c\pffOT]$.
\item[(b)] If $\langle\pffW(y),\pffW(z)\rangle\ne0$ for some $z$, then $y$ is near at every floor and
$\mathfrak d_a(y,x)$ contains a $c\in\pffOT$ with $\pbvol\mathfrak d_a(y,x)=|c|_T$.
\end{enumerate}
\end{lemma}
\begin{proof}
(a) The pairs $(w,d)$ with $w$ stable at $(y,yx)$ for $a$ and $w+d$ its second lift are the solutions
of $X(w,d)\in L_T\oplus L_T$ with $w\perp e_b$, $d\in\langle e_b:b<a\rangle$ ($b<a$), for the invertible
matrix $X=\bigl(\begin{smallmatrix}y^{-1}&0\\M^{-1}(y^{-1}-(yx)^{-1})&-M^{-1}(yx)^{-1}\end{smallmatrix}\bigr)$.
The set of the values of one coordinate (here the $e_a^*$-coordinate of $w$) on the solutions of an
integrality condition $X v\in\pffOT^{2r}$ with $X$ invertible, subject to the vanishing of the
coordinates of larger weight, is an $\pffOT$-module that is open and compact, i.e.\ $c\pffOT$ for a unit
$c$ ($\pffOT$ is a finite product of discrete valuation rings, hence a principal ideal ring). If
$\xi y\in\pffK$ with $\xi\in R^{(a+1)}$, then $h(\xi y\cdot y^{-1}w,e_a)$-type coordinates are
integral: the $e_a^*$-coordinate of $w$ equals that of $\xi w$ ($\xi$ moves $e_b$, $b\le a$, by
combinations of the $e_{b'}$, $b'<b$, and $w\perp e_b$ for $b<a$), and $\xi w=(\xi y)(y^{-1}w)$ is
integral. So $c\in\pffOT$, and the volume of $c\pffOT$ is $|c|_T$.
(b) By the definition of $\pffch$, $\pffW(y)\ne0$ gives $u\in U_T$ with $uy\in\pffK$, so $y$ is near
at every floor ($U_T\subseteq R^{(l)}$); (a) applies.
\end{proof}

\begin{lemma}[Two gauges]\label{pff:twogauges}
Let
$\Gamma$, $\Gamma'$ be gauges with $\Gamma(10n)\subseteq\Gamma'(1)$ and $\pffOT\subseteq\Gamma'(2)$ (for instance any gauges $\Gamma_\alpha$ of Lemma~\ref{pff:gauge}), $p,q$ heads of length $j$ for
$\Gamma$, $M^{-1}\in\Gamma(1)$ if $j\ne0$. Let $y$ be active at $x$ with $\langle\pffW(y),
\pffW(yx)\rangle\ne0$, and suppose that for every $a$ with $j\le a<\varrho$ the ideal $\mathfrak
d_a(y,x)$ contains a unit $d_a$ with $d_a^{-1}\in\Gamma'(1)$. Then $M^{-1}(x-1)$ has entries in
$\Gamma'(99n)$.
\end{lemma}
(Proved in Corollary~\ref{pfg:twogaugesT}, from Theorem~\ref{pfg:twogauges}.)

\begin{lemma}[Dyadic shells]\label{pff:dyadic}
For $x$ a position, let $\mathfrak J(x)\subseteq\pffOT$ be the ideal of the $\varepsilon\in\pffOT$ with
$\varepsilon M^{-1}(x-1)$ integral, and for $i\ge0$ let $S_i$ be the set of positions with
$\lceil\log_2[\pffOT:\mathfrak J(x)]\rceil=i$. Let $\Gamma=\Gamma_\alpha$, $p,q$ heads of length $j\le\varrho$
for $\Gamma$, and $M^{-1}\in\Gamma(1)$ if $j\ne0$. Then $S_i$ is a shell of height $2^i$ for
$\Gamma,j,p,q,\pffW$ (Definition~\ref{pfr:shell}).
\end{lemma}
\begin{proof}
(a) By (L2) the number of ideals of index at most $H=2^i$ is at most $\zeta_{\pffOT}(m')H^{m'}\le2H^{9n}$.
For each such ideal $\mathfrak J=c\pffOT$, the gauge $\Gamma_{\alpha\aubar cc}$ has loss
$\Lambda|c|_T^{-2}\le\Lambda H^2$, and contains in its position ball the positions with $\mathfrak J(x)=
\mathfrak J$, since $\aubar c\,c\,M^{-1}(x-1)$ is integral. (b) Let $x\in S_i$, $\mathfrak J(x)=c\pffOT$,
$[\pffOT:\mathfrak J(x)]\ge H/2$, and $\varsigma<1$ with $\Lambda^{1200n^2}\le\varsigma^{400n^2}H$; then,
as $2\le\Lambda$, $|c|_T\Lambda^{990n^2}\le\varsigma^{400n^2}$. Let $y$ be active at $x$ and suppose that
$\pbvol\mathfrak d_a(y,x)>\varsigma^2$ for all $a\ge j$. By Lemma~\ref{pff:dideal}(b) choose $d_a\in\pffOT
\cap\mathfrak d_a(y,x)$ with $|d_a|_T\ge\varsigma^2$, and let $C=\prod_{a\ge j}d_a\aubar d_a$, so that
$|C|_T\ge\varsigma^{4n}$, $\aubar C=C$ and $C/d_a\in\pffOT$. For $\Gamma'=\Gamma_{\alpha^{10n}C}$ we have
$\Gamma(10n)\subseteq\Gamma'(1)$ and $d_a^{-1}\in\Gamma'(1)$, so by Lemma~\ref{pff:twogauges}
$(\alpha^{10n}C)^{99n}$ lies in $\mathfrak J(x)=c\pffOT$. Taking modules, $|C|_T^{99n}\le|c|_T|\alpha|_T^{-990n^2}\le|c|_T\Lambda^{990n^2}\le\varsigma^{400n^2}<
\varsigma^{396n^2}\le|C|_T^{99n}$, a contradiction. So some $a\ge j$ has depth $\ge\varsigma^{-2}$.
\end{proof}

\begin{lemma}[A deep vector outside $\pffKM$]\label{pff:deep}
Let $Q\ge2$ be such that $|c|_T\le Q^{-1}$ for every $c\in\pffOT\cap\pffET^\times$ with
$c^{-1}\notin\pffOT$. If $x\notin\pffKM$, then for the heads of length $0$ every point $y$ active at $x$
with $\langle\pffW(y),\pffW(yx)\rangle\ne0$ has a frame vector of depth at least $Q$.
\end{lemma}
\begin{proof}
Otherwise, by Lemma~\ref{pff:dideal}, each $\mathfrak d_a(y,x)$ contains $c_a\in\pffOT$ with
$|c_a|_T=\pbvol\mathfrak d_a>Q^{-1}$, hence $c_a^{-1}\in\pffOT$. Apply Lemma~\ref{pff:twogauges} with
$\Gamma=\Gamma'=\Gamma_1$ (all sizes mean integral): $M^{-1}(x-1)$ is integral. As $x$ is unitary and the form
unimodular, $x^{-1}=\mathsf J^{-1}{}^t\aubar x\mathsf J$ is integral, so $x\in\pffKM$.
\end{proof}
Active points with $\langle\pffW(y),\pffW(yx)\rangle=0$ contribute nothing to $\pfrF_0^+(x)$, whose
integrand involves this pairing; so only the points covered by Lemma~\ref{pff:deep} matter below.

\begin{proposition}[The positions outside $\pffKM$]\label{pff:outside}
Let $Q$ be as in Lemma~\ref{pff:deep} and $t\ge0$ with $t+2\cdot10^9n^4\le k$. Then
\[
\int_{\pffG\smallsetminus\pffKM}\Bigl|\int_{\pffZ}\pffch(yxy^{-1})\,dy\Bigr|^2dx\le Q^{-t}\,\pbvol(
\pffZ)\,\pbvol(\pffKM).
\]
\end{proposition}
\begin{proof}
At floor $0$ with the heads $1,1$ the majorant is $\pfrF_0^+(x)=\pfrB^1_0(1,x)^2\ge|k_0(x)|^2$
(Lemma~\ref{pfr:L1} and the definition of $k_0$ in Corollary~\ref{pff:k0}). Let $f=\mathbf
1_{\pffG\smallsetminus\pffKM}\pfrF^+_0$. In the gauge $\Gamma_1$ (loss $2$) the mass of floor $0$ is at most
$V=2^{1287n}$ (Lemma~\ref{pff:floormass}), and by Lemma~\ref{pff:deep} and the last assertion of Theorem~\ref{pfr:deep}
(the case $j=0$) with $\varsigma=Q^{-1}$, $\pfrF^+_0(x)^2\le Q^{-2s}\mathsf Q(x)$ for $x\notin\pffKM$, so that $f$ has the
gain $Q^{-s}\le VQ^{-s}$ at every position, where $s\ge9\cdot10^7n^4+r'$ for $r'=t+10^9n^4$ (as $s=\iota-2(2m'+2)-
(3n+1+m')$ and $k$ is large). By Lemma~\ref{pff:dyadic} and Theorem~\ref{pfr:shells}
\[
\int f\,d\pbvol\le V\,Q^{-t-10^9n^4}\,2^{4\cdot10^8n^4}\,\mathsf N(\pffW)\le Q^{-t}\pffr(\pffZ),
\]
using $Q\ge2$ and $\mathsf N(\pffW)\le\pffr(\pffZ)$; the last inequality follows from
$\pffl\iint|\langle\pffW(g),\pffW(y)\rangle|^2d\theta\,d\theta\le\pffr(\pffZ)$ (for fixed $g$, the $y$
with a non-zero pairing lie in a set of $\theta$-mass $\pbvol(\pffKM)/\pffl$, and
$\pffl|\langle\pffW(g),\pffW(y)\rangle|^2\le\|\pffW(g)\|^2$ by Corollary~\ref{pff:k0}; then integrate in $g$)
and $\prod_{l<\varrho}\mathsf c_l(M)=\pffl$ at a finite place (the floor moduli). Multiply by
$\pffl^{\,2}$ and use Corollary~\ref{pff:k0}.
\end{proof}

\begin{theorem}[The mass integral]\label{pff:massZ}
Let $Q$, $t$ be as in Proposition~\ref{pff:outside}. For every intermediate group $\pffZ$,
\[
\pbvol(\pffZ)^2\,\pbvol(\pffKM)\ \le\ \pffm(\pffZ)\ \le\ \pbvol(\pffZ)\,\pbvol(\pffKM)\,\bigl(\pbvol(
\pffZ)+Q^{-t}\bigr).
\]
\end{theorem}
\begin{proof}
Split the outer integral of $\pffm(\pffZ)$ into $\pffKM$ and its complement, and apply
Lemma~\ref{pff:inKM} and Proposition~\ref{pff:outside}; the lower bound is Lemma~\ref{pff:inKM}
alone.
\end{proof}

\begin{corollary}\label{pff:massK}
If $b\ge1000$ and $2^{-\varrho}\le b^{-6}$, then $\pffm(\pffK)\le(1+b^{-6})\,\pbvol(\pffK)\,
\pbvol(\pffKM)$.
\end{corollary}
\begin{proof}
Take $Q=2$: a $c\in\pffOT$ with $c^{-1}\notin\pffOT$ generates a proper ideal, of index
$|c|_T^{-1}\ge2$. Take $t=n$; then $t+2\cdot10^9n^4\le k$. By Theorem~\ref{pff:massZ} and
Lemma~\ref{pff:vol}(a), $\pffm(\pffK)\le\pbvol(\pffK)\pbvol(\pffKM)(1+2^{-n}+2^{-n})$, and
$1+2^{1-n}=1+2^{-\varrho}\le1+b^{-6}$.
\end{proof}

For the data, $2^{-\varrho}\le b^{-6}$ holds with room: $\varrho\ge\sqrt x/2-1$ and $(\log b)^{100}
\le x$ (Lemma~\ref{an:lem-sizes}).

\subsection{The basis at $T$, floor masses and volumes}\label{pff:ss-floors}

\begin{lemma}[The basis (L1)]\label{pff:L1}
For the data, $L_T$ has a basis as in \textup{(L1)}, consisting of integral vectors with integral
dual basis.
\end{lemma}
\begin{proof}
In the lattice basis the middle block $X_0\perp Ez$ has the vectors $x_i$ $(i\ne k_0)$, orthonormal,
and $v=\aubar\pi_N^{-1}x_{k_0}+N^{-1}z$, $z$, with $h(v,v)=0$, $h(v,z)=-1$, $h(z,z)=-N$. The field
$E/E^+$ is unramified at every finite place (\S\ref{pb:ss-setting}, data of
\S\ref{ss:au-notation}), so at every $w\in T$ the norm map $\auO_w^\times\to\auO_{E^+,v}^\times$ and the
trace $\auO_w\to\auO_{E^+,v}$ are onto (for a split place both are clear); hence there are
$c,\tau\in\pffOT$ with $\aubar cc=-1$ and $\tau+\aubar\tau=1$. Then $-v$ and $z-N\tau v$ form a
hyperbolic pair ($h(-v,z-N\tau v)=1$, $h(z-N\tau v,z-N\tau v)=-N+N\tau+N\aubar\tau=0$), and the second
vector has coordinate $1$ at $z$; for two orthonormal vectors $x,x'$, the vectors $p=x+cx'$ and
$q=\aubar\tau x-c\tau\,x'$ form a hyperbolic pair: $h(p,p)=1+\aubar cc=0$, $h(q,q)=(1+c\aubar c)\tau\aubar\tau=0$,
$h(p,q)=\tau+\aubar\tau=1$, and the change of basis $(x,x')\mapsto(p,q)$ has the matrix
$\bigl(\begin{smallmatrix}1&\aubar\tau\\c&-c\tau\end{smallmatrix}\bigr)$ of determinant $-c$ (the explicit basis and its
Gram matrix are written out in \S\ref{pbv:ss-weak}(ii)). Using $2(m'+1)$ of the $x_i$ this way and keeping
$\iota$ of them, the change of basis is integral with determinant a unit, as $c,\tau\in\pffOT$ and
$c^{-1}=-\aubar c$; and $k+1=\iota+2(m'+1)+2$.
\end{proof}

The places where the lattice basis is not integral with integral inverse: the matrix from the basis
$e_a,x_i,z,e_a^*$ to the lattice basis has the entries $\aubar\pi_N^{-1}$ and $(\pi_N\aubar\pi_N)^{-1}=
N^{-1}$, its inverse the entries $\aubar\pi_N$ and $-\pi_N^{-1}$; so both are integral at the places
not dividing $N$. This verifies the claim of \S\ref{pff:ss-setting}.

\begin{lemma}[Floor masses and the volume of $\pffK$]\label{pff:floormass}
Assume (L1)--(L3) and let $\Gamma=\Gamma_1$.
\begin{enumerate}
\item[(a)] For $l\le\varrho$ and every $g$, $\theta_l\{\sigma:\sigma g\text{ near at }l+1\}\le\zeta_{\pffOT}(m')$;
in the notation of \S\ref{pfr:ss-gauge}, the mass of every floor is at most $V=2\le2^{1287n}$ (the mass involves only
the box $\pffK$, not the gauge, so this holds for every gauge $\Gamma_\alpha$).
\item[(b)] $\pbvol(\pffK)\le\zeta_{\pffOT}(m')^n\le1+2^{-n}$.
\item[(c)] $\pbvol(\pffK)\ge\pbvol(\pffOT^\times)^n$; in particular $\pbvol(\pffK)\ge2^{-n}$ if the
non-units have volume at most $\frac12$.
\end{enumerate}
\end{lemma}
\begin{proof}
(a) Apply Lemma~\ref{pfr:chart} with $a=l$, the gauge $\Gamma_1$ (for which $\Gamma_1(7)$ is the set of integral
vectors), $c\equiv1$ and $\Psi_g\equiv1$: the mass is at most $c_l\,\mu_l(L_l)$, where $L_l$ is the lattice of the integral
vectors of $V_l$. Here $c_l=1$ (\S\ref{pff:ss-large}), and by Proposition~\ref{pfg:fibres}(e) for the unimodular lattice
$L_l$ (of rank $\kappa\ge k+1\ge2m'+1$), with $\mu=m'$, $g=1$, $x_0=0$: $\mu_l(L_l)\le\zeta_{\pffOT}(m')$. Finally
$\zeta_{\pffOT}(m')\le2$ by (L2).
(b) If $g\in\pffK$ then $g$ is near at every floor. Unfolding $\pbvol=\pfrM$ floor by floor
($\pfrM=\eta_{[0,\varrho+2)}\ast\theta_{\ge0}$ and the measures $\eta_l$ give $\pfrN_l\cap\pffK$ mass $1$),
$\pbvol(\pffK)\le\theta_{\ge0}\{g:g\text{ near at }\varrho+2\}\le\prod_{l<n}\zeta_{\pffOT}(m')$ by (a)
applied at each floor, and $\zeta_{\pffOT}(m')^n\le1+2^{-n}$ by (L2).
(c) $\pffK$ is a group, so for measures $\nu,\nu'$ on $\pffG$, $(\nu\ast\nu')(\pffK)\ge\nu(\pffK)\nu'(\pffK)$. As
$\eta_\varrho=\eta_{\varrho+1}=\delta_1$ and $\theta_{\varrho+1}=\mu_H$,
\[
\pfrM=(\mu_H\ast\theta_\varrho)\ast\eta_{\varrho-1}\ast\theta_{\varrho-1}\ast\dots\ast\eta_0\ast\theta_0,\qquad
\pbvol(\pffK)\ge(\mu_H\ast\theta_\varrho)(\pffK)\prod_{l<\varrho}\eta_l(\pffK)\,\theta_l(\pffK).
\]
\emph{Lower floors.} $\eta_l(\pffK)=\eta_l(\pfrN_l\cap\pffK)=1$: the element $n(\omega,b-\tau_0h(\omega,\omega))$ and its inverse
$n(-\omega,\bar\beta)$ are integral if and only if $\omega$ and $b$ are integral ($\tau_0h(\omega,\omega)$ is then integral, and the
lattice of $V_{l+1}$ is unimodular), a set of $d\omega\otimes d\beta_{\rm sk}$-measure $1$. For a unit $c$ of $\pffOT$,
$w_{l,c}^{-1}\in\pffK$; so $\theta_l(\pffK)\ge\int_{\pffOT^\times}|c|_T^{d_l}\,\eta_l(\pffK)\,dc=\pbvol(\pffOT^\times)$.

\emph{The top floor.} $(\mu_H\ast\theta_\varrho)(\pffK)=\int_{\pfgQ}\mu_H\{k\in H_T:ks(e)\in\pffK\}\,d\mu_\varrho(e)$. Use the hyperbolic
pair $(u,u')$ of (L1), with $z=u'-N\tau u$, and the orthogonal decomposition $L^+=\pffOT u\oplus\pffOT u'\oplus L'$, $L'$ the
span of the remaining basis vectors of $L^+$ (unimodular). Let $E_1\subseteq L^+$ be the compact open set of the
$e=au+bu'+y'$ ($a\in\pffOT$, $b\in\pffOT^\times$, $y'\in L'$). (i) \emph{For $e\in E_1\cap\pfgQ$ there is $\sigma_e\in\pffK\cap
\pffG^{(\varrho)}$ with $\sigma_e^{-1}z=e$.} For $\omega\in L'$ and integral $\beta$ with $\beta+\bar\beta=-h(\omega,\omega)$ let
$n_{\omega,\beta}$ be the isometry of $\pbmXp_T$ fixing $u$, sending $x\mapsto x-h(x,\omega)u$ on $L'\otimes\pffET$ and
$u'\mapsto u'+\omega+\beta u$; for a unit $c$ of $\pffOT$ let $d_c$ fix $L'$ and send $u\mapsto cu$, $u'\mapsto\bar c^{-1}u'$. These
are isometries (as in \S\ref{pfr:ss-floors}), preserve $L^+$ together with their inverses, and extended by the identity on
the $e_a,e_a^*$ they lie in $\pffK\cap\pffG^{(\varrho)}$ (by (L1), $L_T$ is the orthogonal sum of the planes $\langle e_a,e_a^*\rangle$
and $L^+$). With $\omega=-b^{-1}y'$ and $\beta=-\tau_0h(\omega,\omega)$, $n_{\omega,\beta}e=a_1u+bu'$ for some $a_1\in\pffOT$; then
$d_{\bar b}(a_1u+bu')=a_2u+u'$ with $a_2=a_1\bar b$; as these maps are isometries, $a_2+\bar a_2=h(e,e)=-N$. With the skew
integral $\beta'=-N\tau-a_2$ ($\beta'+\bar\beta'=-N+N=0$), $n_{0,\beta'}(a_2u+u')=-N\tau u+u'=z$. So $\gamma=n_{0,\beta'}d_{\bar b}
n_{\omega,\beta}\in\pffK\cap\pffG^{(\varrho)}$ has $\gamma e=z$, and $\sigma_e=\gamma$ works. (ii) For such $e$, $s(e)\in H_T\sigma_e$, say
$s(e)=k_0\sigma_e$; then $\{k:ks(e)\in\pffK\}=(H_T\cap\pffK\sigma_e^{-1})k_0^{-1}=(H_T\cap\pffK)k_0^{-1}$, of $\mu_H$-measure $1$ by
right invariance. Hence $(\mu_H\ast\theta_\varrho)(\pffK)\ge\mu_\varrho(E_1)$. (iii) $\mu_\varrho(E_1)=\pbvol(\pffOT^\times)$: by
Proposition~\ref{pfg:fibres}(a), $\mu_\varrho(E_1)=\lim_\ell A_\ell(\mathbf 1_{E_1};-N)$, and in the coordinates $(a,b,y')$, in which
$dy=da\,db\,dy'$ and $h(e,e)=a\bar b+\bar ab+h(y',y')$,
\[
A_\ell(\mathbf 1_{E_1};-N)=\frac1{\pbvol(\varpi_1^\ell\pfgOh)}\int_{\pffOT^\times}\int_{L'}\pbvol\{a\in\pffOT:\
a\bar b+\overline{a\bar b}\in-N-h(y',y')+\varpi_1^\ell\pfgOh\}\,dy'\,db=\pbvol(\pffOT^\times),
\]
because $a\mapsto a\bar b$ preserves the Haar measure of $\pffOT$ ($b$ a unit of $\pffOT$), the trace $x\mapsto x+\bar x$ maps
$\pffOT$ onto $\pfgOh$ (Lemma~\ref{pfg:duality}(a)) and so carries the Haar probability measure of $\pffOT$ to that of $\pfgOh$, and
$-N-h(y',y')\in\pfgOh$. So $(\mu_H\ast\theta_\varrho)(\pffK)\ge\pbvol(\pffOT^\times)$.

Altogether $\pbvol(\pffK)\ge\pbvol(\pffOT^\times)^{\varrho+1}=\pbvol(\pffOT^\times)^n$. If the non-units have volume $\le\frac12$,
$\pbvol(\pffOT^\times)\ge\frac12$.
\end{proof}

\subsubsection*{The density at the last entry}
For $l\in\pbmGf_T$ and a sequence of units $p$ let
\[
\mathcal O_p(l)=\pfrkap_\varrho(p,1)\int_{\pffG}\pffW(y)\otimes\aubar{\pffW\bigl((w_pl^{-1})^{-1}y\bigr)}
\,dy\ \in\C^{C\times C}\qquad(C=\pffKM\backslash\pffZ),
\]
the \emph{scaled orbital integral} at the element $w_pl^{-1}$ (head $p$, link $l^{-1}$), and
\[
\Pi_{p,q}(l^{-1}z)=\int_{H_T}\langle\mathcal O_p(kl),\mathcal O_q(kl)\rangle\,d\mu_H(k),
\]
the \emph{fibre pairing}; it depends only on $e=l^{-1}z$, as two $l$ with the same $e$ differ on the left
by an element of $H_T$ and $\mu_H$ is right invariant. Put $\Pi_{p,q}=0$ off the orbit of $z$.

\begin{definition}\label{pff:lastentry}
Let $\pi_N$ be the component at $T$ of the idele with $\pi_N\aubar\pi_N=N$. For heads $p,q$ of length
$n$ with $p_\varrho=q_\varrho=t$, the \emph{density at the last entry} is
\[
\mathcal D(p,q)=|\pi_N|_T^{-1}\int_{X_{0,T}}\Pi_{p,q}\bigl(x-N^{-1}t\,z\bigr)\,d\nu_{\beta(t)}(x),
\qquad\beta(t)=\frac{t\aubar t}N-N,
\]
where $\nu_\beta$ is the measure $dx/d\langle x,x\rangle$ on the sphere $\{x\in X_{0,T}:h(x,x)=\beta\}$ of
$X_{0,T}=X_0\otimes\pffET$ (the fibre measures of the disintegration of $dx$ along $x\mapsto h(x,x)$, constructed in Proposition~\ref{pfg:fibres}),
and the vector $x-N^{-1}tz$ is read in the lattice basis.
\end{definition}

The vector $e=x-N^{-1}tz$ satisfies $h(e,e)=h(x,x)-t\aubar t/N=-N$ and, by (Cv1), $h(e,z)=-N^{-1}t\,h(z,z)=t$: it runs over the
points of the quadric $h(e,e)=-N$ of $\pbmXp\otimes\pffET$ with last entry $t$.

\begin{lemma}[Properties of the fibre measures and of $\Pi$]\label{pff:fibre}
Assume (L1)--(L3).
\begin{enumerate}
\item[(a)] (\emph{Coarea.}) For $g$ continuous with compact support on $\pbmXp\otimes\pffET$,
$\int g\,d\theta_\varrho^{\sharp}=\int_{\pffET}|\pi_N|_T^{-1}\int g(x-N^{-1}tz)\,d\nu_{\beta(t)}(x)\,dt$, where
$\theta_\varrho^\sharp$ is the measure of the top quadric (image of $\theta_\varrho$ under $\sigma\mapsto
\sigma^{-1}z$); and $t\mapsto\int g(x-N^{-1}tz)d\nu_{\beta(t)}(x)$ is continuous.
\item[(b)] (\emph{Step function.}) There is a function $\Phi_{p,q}$ on $\pbmXp\otimes\pffET$, locally constant with
compact support and invariant under $H_T\cap\pffKM$, that equals $\Pi_{p,q}$ on the quadric $h(e,e)=-N$; and
$\Pi_{p',q'}=\Pi_{p,q}$ for all $(p',q')$ in a neighbourhood of $(p,q)$ ($\Pi$ does not depend on the entries of
index $\varrho$).
\end{enumerate}
\end{lemma}
\begin{proof}
(a) is Lemma~\ref{pfg:coarea}, with the fibre measures of Proposition~\ref{pfg:fibres} (Remark~\ref{pfg:rem-fibres});
(b) is Proposition~\ref{pfg:fibrepairing}(c),(d).
\end{proof}

\begin{proposition}[Top floor]\label{pff:top}
Assume (L1)--(L3), $k\ge10^{10}n^4$. For heads $p,q$ and $\xi\in\pffET$, writing $p^t,q^t$ for $p,q$ with
the last entry replaced by $t$:
\begin{enumerate}
\item[(a)] $t\mapsto\psi_T(\xi t)\mathcal D(p^t,q^t)$ is integrable and
$\int\psi_T(\xi t)\mathcal D(p^t,q^t)\,dt=\pfrI_\varrho\bigl(p,q;\ g\mapsto\psi_T(\xi h(gz,z))\bigr)$;
\item[(b)] for heads $p^{(1)},\dots,p^{(\kappa)}$ with a common last entry, the matrix
$(\mathcal D(p^{(a)},p^{(b)}))_{a,b}$ is hermitian and positive semidefinite;
\item[(c)] $(p,q)\mapsto\mathcal D(p,q)$ is continuous.
\end{enumerate}
\end{proposition}
\begin{proof}
(a) Write $\omega(g)=\psi_T(\xi h(gz,z))$, $\nu=\mu_H\ast\theta_\varrho$, $w=w_p$, $d=w_q^{-1}w_p$. Three facts are used.
(i) $\omega$ is invariant under left and right translation by $H_T$: $H_Tz=z$ and $h(kgz,z)=h(gz,k^{-1}z)$ for $k\in H_T$.
(ii) $\nu$ is invariant under right translation by $\pffG^{(\varrho)}$: for Borel $F\ge0$ put $F^H(\sigma)=\int F(h\sigma)\,
d\mu_H(h)$; it is left $H_T$-invariant ($\mu_H$ is right invariant), so $F^H(\sigma)=F_0(\sigma^{-1}z)$ with $F_0$ on $\pfgQ$, and
for $g\in\pffG^{(\varrho)}$, $\int F(lg)\,d\nu(l)=\int F_0(g^{-1}s(e)^{-1}z)\,d\mu_\varrho(e)=\int F_0(g^{-1}e)\,d\mu_\varrho(e)=\int F_0\,d\mu_\varrho
=\int F\,d\nu$, by the $\pffG^{(\varrho)}$-invariance of $\mu_\varrho$ (\S\ref{pfr:ss-setting}). Moreover $\nu$ is finite on compact
sets $C$: $\nu(C)=\int\mu_H(H_T\cap Cs(e)^{-1})\,d\mu_\varrho(e)$, the integrand vanishes unless $e\in C^{-1}z$ (if $hs(e)\in C$ then
$e=(hs(e))^{-1}z$), and $H_T\cap Cs(e)^{-1}\subseteq(H_T\cap CC^{-1})h_0$ for any of its elements $h_0$, so
$\nu(C)\le\mu_H(H_T\cap CC^{-1})\,\mu_\varrho(C^{-1}z)<\infty$.
(iii) $\pfrM=\nu\ast\pfrm_\varrho$, that is $\int F\,d\pfrM=\iint F(sg)\,d\nu(s)\,d\pfrm_\varrho(g)$: this is the definition of
$\pfrM$ and $\pfrm_\varrho$ (\S\ref{pfr:ss-setting}), as $\eta_\varrho=\eta_{\varrho+1}=\delta_1$ and $\theta_{\varrho+1}=\mu_H$; and
$\theta_{\ge\varrho+1}=\mu_H$, so $\mathsf k^{d'}_{\varrho+1}(g,y)=\int\pfrip{\pffW(hg)}{\pffW(d'hy)}\,d\mu_H(h)$.
Also $\theta_\varrho$, $\mu_H$ are carried by $\pffG^{(\varrho)}$, whose elements commute with $w_p$, $w_q$ and $d$.

By the definitions of \S\ref{pfr:ss-setting} with $j=\varrho$, $\pfrI_\varrho(p,q;\omega)$ is $\pfrkap_\varrho(p,q)$ times the integral over
$g$ ($\pfrm_\varrho$), $x$ ($\pfrM$), $\sigma,\tau$ ($\theta_\varrho$) and $h,h'$ ($\mu_H$) of
\[
\omega(\sigma\tau^{-1})\,\pfrip{\pffW(h\sigma g)}{\pffW(h\sigma gx)}\,\overline{\pfrip{\pffW(h'\tau w^{-1}g)}{\pffW(dh'\tau w^{-1}gx)}}.
\]
We first transform the integral of the absolute value of this function by Tonelli's theorem; the same steps, applied to the
function itself, are then justified by Fubini's theorem. Put $s=h\sigma$, $L=h'\tau s^{-1}\in\pffG^{(\varrho)}$. By (i),
$\omega(\sigma\tau^{-1})=\omega(s(h'\tau)^{-1})=\omega(L^{-1})$. Substitute $y=sgx$ in the $x$-integral ($\pfrM$ is a Haar measure of
the unimodular group $\pffG$): then $h'\tau w^{-1}g=Lw^{-1}sg$ and $dh'\tau w^{-1}gx=dLw^{-1}y=Lw_q^{-1}y$. Next, for fixed $s$,
substitute $l=(h'\tau)s^{-1}$ in the $(h',\tau)$-integral, which is an integral against $\nu$; by (ii) it preserves $\nu$. Finally
integrate over $(h,\sigma,g)$ with $u=sg$, which by (iii) is an integral against $\pfrM$. The result is
\[
\pfrI_\varrho(p,q;\omega)=\pfrkap_\varrho(p,q)\int\omega(l^{-1})\iint\pfrip{\pffW(u)}{\pffW(y)}\,
\overline{\pfrip{\pffW(lw_p^{-1}u)}{\pffW(lw_q^{-1}y)}}\,du\,dy\,d\nu(l),
\]
and the same with absolute values inside for the absolute integral. The latter is finite: by Lemma~\ref{pfg:Op}(a),(b) (for $p$
and for $q$) the inner integrand vanishes unless $u\in S_p$, $y\in S_q$ and $l\in C_p$, compact sets, it is bounded by (W3),
and $\nu(C_p)<\infty$ by (ii). So $\pfrI_\varrho(p,q;\omega)$ is defined, and the identity holds. With the pairing
$\langle a\otimes\bar b,c\otimes\bar d\rangle=\pfrip ac\,\overline{\pfrip bd}$ on $\C^{C\times C}$, the inner double integral is
$\langle\int\pffW(u)\otimes\overline{\pffW(lw_p^{-1}u)}\,du,\int\pffW(y)\otimes\overline{\pffW(lw_q^{-1}y)}\,dy\rangle$, and
$\pfrkap_\varrho(p,q)=\pfrkap_\varrho(p,1)\pfrkap_\varrho(q,1)$; hence
\[
\pfrI_\varrho(p,q;\omega)=\int\omega(l^{-1})\,\langle\mathcal O_p(l),\mathcal O_q(l)\rangle\,d\nu(l).
\]
Thus the argument of the weight is $l^{-1}$, where $l$ is the argument of the scaled orbital integrals: no conjugation of $\xi$
occurs. Writing $l=hs(e)$ ($h\in H_T$, $e\in\pfgQ$), so that $l^{-1}z=e$ and $\omega(l^{-1})=\psi_T(\xi h(e,z))$, and integrating
first over $h$ gives $\int\psi_T(\xi h(e,z))\,\Pi_{p,q}(e)\,d\mu_\varrho(e)$, where $\mu_\varrho=\theta_\varrho^\sharp$ ($c_\varrho=1$).
Lemma~\ref{pff:fibre}(a) with $g(e)=\psi_T(\xi h(e,z))\Phi_{p,q}(e)$ (locally constant with compact support by
Lemma~\ref{pff:fibre}(b); on the quadric $h(x-N^{-1}tz,z)=t$) gives $\int\psi_T(\xi t)\mathcal D(p^t,q^t)\,dt$, the function
$t\mapsto\psi_T(\xi t)\mathcal D(p^t,q^t)$ being continuous with compact support, hence integrable ($\Pi$ does not depend on the
last entry of the heads).
(b) For $u\in\C^\kappa$, $\sum_{a,b}\aubar u_au_b\Pi_{p^{(a)},p^{(b)}}(e)=\int_{H_T}\|\sum_bu_b\mathcal
O_{p^{(b)}}(kl)\|^2d\mu_H\ge0$, and hermitian symmetry is clear; integrate against the positive
measure $\nu_\beta$.
(c) Near $(p,q)$, by Lemma~\ref{pff:fibre}(b), $\mathcal D(p',q')=|\pi_N|_T^{-1}\int\Phi_{p,q}(x-N^{-1}
p'_\varrho z)\,d\nu_{\beta(p'_\varrho)}(x)$, continuous in $p'_\varrho$ by Lemma~\ref{pff:fibre}(a).
\end{proof}

\subsubsection*{The decay at $T$}
For a sequence $r$ and $K\ge0$ let $\Lambda_K(r)=\prod_{l\le i}[\delta_l\in\pffOT]\,\pbmn(\delta_l)^{K/2}$
($\delta_l=r_0\cdots r_{l-1}$), so that $\Lambda^T_i=\Lambda_{n^6}$ for heads of length $i$.

\begin{lemma}[Heads of a non-zero form]\label{pff:heads}
If the majorant $\pfrF^+_j$ of $p,q,\pffW$ is not identically zero, then $\delta_l(p),\delta_l(q)\in
\pffOT$ for $l\le j$, and there is $\alpha\in\pffOT$, $\aubar\alpha=\alpha$, such that in $\Gamma_\alpha$:
$p,q$ are heads of length $j$; $M^{-1}\in\Gamma_\alpha(1)$ if $j\ne0$; $p_i^{-1}-q_i^{-1}\in
\Gamma_\alpha(1)$ for $i<j$; and $\Lambda_\alpha^E\le(2/|M|_T)^K\Lambda_K(p)\Lambda_K(q)$ whenever
$4E\le K$.
\end{lemma}
\begin{proof}
Where the integrand of the majorant is non-zero there are points $y$, $yx$ with
$w_p^{-1}y$, $w_q^{-1}yx$ near at $j$ (\S\ref{pfr:ss-gauge}); in the gauge $\Gamma_1$ the leading
minors of order $l\le j$ of the frame matrix of an element moved into $\pffK$ by $R^{(j)}$ are integral, and
these minors are the $\delta_l$. Let $\beta=\prod_{l\le j}\delta_l(p)$, $\gamma=\prod_{l\le j}\delta_l(q)$;
they are integral multiples of each $\delta_l$ and of each $p_i$, resp.\ $q_i$ ($p_i=\delta_{i+1}/\delta_i$),
and $\Lambda_K(p)=\pbmn(\beta)^{K/2}$, $\Lambda_K(q)=\pbmn(\gamma)^{K/2}$. Put $\alpha=M\,\aubar{\beta\gamma}\,
\beta\gamma$. Then $\alpha\delta_l(p)^{-1}$, $\alpha\delta_l(q)^{-1}$, $\alpha M^{-1}$, $\alpha(p_i^{-1}-q_i^{-1})$
are integral, and $|\alpha|_T=|M|_T(|\beta|_T|\gamma|_T)^2$, so $\Lambda_\alpha=2/|\alpha|_T=(2/|M|_T)(\pbmn
(\beta)\pbmn(\gamma))^2$; as each factor is $\ge1$ and $4E\le K$, $\Lambda_\alpha^E\le(2/|M|_T)^K\pbmn(\beta)^{K/2}
\pbmn(\gamma)^{K/2}$.
\end{proof}

\begin{theorem}[Decay at $T$]\label{pff:decay}
Assume (L1)--(L3), $j\le\varrho$, $t\ge0$ with $t+10^{10}n^4\le k$, and $K\ge10^{10}n^4$. Let $\xi$ be a
unit and either (i) $\omega(g)=\psi_T(\xi\,h(gf_j,f_j))$, or (ii) $\omega=1$, $j=i+1$ and
$\xi=M^{-1}(p_i^{-1}-q_i^{-1})$. Then
\[
|\pfrI_j(p,q;\omega)|\le\pffr(\pffZ)\,\pbmn(\xi)^{t/2}\,\pbmden(\xi)^{-6}\,(2/|M|_T)^K\,\Lambda_K(p)\Lambda_K(q).
\]
\end{theorem}
\begin{proof}
If $\pfrF^+_j\equiv0$ the form vanishes (Lemma~\ref{pfr:L1}). Otherwise take $\alpha$ of
Lemma~\ref{pff:heads}, $\Gamma=\Gamma_\alpha$, $\Lambda=\Lambda_\alpha$, and $\beta'=\alpha^{200n^2}$.

\emph{A unit $\eta$ dividing $\xi$.} In case (i), as $\pffOT$ is a finite product of discrete valuation
rings, there is a unit $\eta=\xi+c\beta'$ of $\pffET$, $c\in\pffOT$, with $\xi\eta^{-1}$ and $\eta^{-1}\beta'$
integral (at each component take $\eta$ to be $\xi$ or $\xi+\beta'$, whichever generates the ideal
$(\xi,\beta')$ there); then $\psi_T(\xi c')=\psi_T(\eta c')$ for $c'\in\Gamma(200n^2)$, because
$\psi_T$ is trivial on $\pffOT$ and $c\beta'c'\in\pffOT$. Put $w=\eta^{-1}\beta'$. In case (ii) put
$\eta=\xi$ and $w=\beta'$; here $\xi\in\Gamma(2)\subseteq\Gamma(200n^2)$ by Lemma~\ref{pff:heads}, so
$\xi w\in\pffOT$. In both cases $w$, $\xi w$, $\xi\eta^{-1}$ are integral, and $|\eta|_T^{-1}=|w|_T
|\beta'|_T^{-1}$ in case (i), while in case (ii) $|w|_T|\beta'|_T^{-1}=1$.

\emph{The gain.} Let $e=|\eta|_T^{-1/2}$. By Theorem~\ref{pfr:gain-char} (case (i), with $\eta$ in place
of the character's argument) or Theorem~\ref{pfr:gain-heads} (case (ii)), with $P=1$,
$\pfrF_j(x;\omega)^2\le e^{2\operatorname{rank}X'}\mathsf Q(x)$ at every position; if $e<1$, as the volume
exponent satisfies $s\le\operatorname{rank}X'$, $\pfrF_j(x;\omega)$ has the gain $e^{s}\le Ve^s$.

\emph{Sum over the shells.} By Lemma~\ref{pfr:L1}, Lemma~\ref{pff:dyadic} (the dyadic shells cover the
positions, outside of which $\pfrF_j^+$ vanishes since the integrand involves $\langle\pffW(y),\pffW(yx)
\rangle\ne0$, which forces $x$ to be a position), Lemma~\ref{pff:floormass} ($V=2$) and
Theorem~\ref{pfr:shells} with $r'=t+12$ (admissible as $s\ge9\cdot10^7n^4+t+12$ for $t+10^{10}n^4\le k$),
\[
|\pfrI_j(p,q;\omega)|\le\pfrkap_j\int\pfrF_j(x;\omega)\,d\pbvol\le2\min(1,e)^{t+12}\Lambda^{4\cdot10^8n^4}
\mathsf N(\pffW),\qquad\mathsf N(\pffW)\le\pffr(\pffZ)
\]
(the last inequality as in the proof of Proposition~\ref{pff:outside}).

\emph{From $\eta$ to $\xi$.} As $\xi\eta^{-1}$ is integral, $\pbmn(\xi)=\pbmn(\xi\eta^{-1})\pbmn(\eta)\ge
\pbmn(\eta)=e^2$, so $\min(1,e)^t\le\pbmn(\xi)^{t/2}$. As $w\pffOT\subseteq\pffOT\cap\xi^{-1}\pffOT$,
$|w|_T\le\pbmden(\xi)^{-1}$; and $|\beta'|_T^{-1}\le\Lambda^{200n^2}$. Moreover $\min(1,e)^2\le|w|_T|\beta'|_T^{-1}$: in case (i) because $e^2=|\eta|_T^{-1}=|w|_T|\beta'|_T^{-1}$, in case (ii) because the right side is $1$. Hence $\min(1,e)^{12}\le(\min(1,e)^2)^6
\le(|w|_T|\beta'|_T^{-1})^6\le\pbmden(\xi)^{-6}\Lambda^{1200n^2}$. Altogether
\[
|\pfrI_j(p,q;\omega)|\le\pffr(\pffZ)\,\pbmn(\xi)^{t/2}\pbmden(\xi)^{-6}\,\Lambda^{E},\qquad E=1+1200n^2+4\cdot10^8
n^4,
\]
and $4E\le10^{10}n^4\le K$, so Lemma~\ref{pff:heads} gives $\Lambda^E\le(2/|M|_T)^K\Lambda_K(p)\Lambda_K(q)$.
\end{proof}


\subsection{Fixed vectors: the gain $b^{-6k}$}\label{pff:ss-fixed}

Let $\mathsf A$ be a set of chosen places above $p_0$ and $\mathsf d=(\mathsf d_w)$ a family of
frame vectors (\S\ref{pb:ss-setting}). Let $K^{\mathrm{ad}}\subset G_f$ be the stabiliser of the
lattice and $\mathrm{red}\colon K^{\mathrm{ad}}\to\auG$ the reduction modulo the chosen places; it
is onto (strong approximation, \S\ref{pb:ss-setting}). The residue field at a chosen place above
$p_0$ is $\auF_{p_0}$.

\begin{lemma}[Few non-units]\label{pff:nonunits}
Let $T'$ be a set of places each of which is a chosen place or the conjugate of one, with
$\#T'\le b$. Then the non-units of $\auO_{T'}$ have volume at most $b^{-9}$.
\end{lemma}
\begin{proof}
An element $y\in\auO_{T'}$ that is not a unit has a component in the maximal ideal $\mathfrak m_w$
for some $w\in T'$. The set of such $y$, for fixed $w$, has volume $q_w^{-1}\le b^{-10}$, as
$q_w\in\{p_0,p_1\}$ and $p_0,p_1\ge b^{10}$. Sum over the at most $b$ places.
\end{proof}

For the data $\#T'\le4f\le b$ (the chosen places above $p_0$ and $p_1$ and their conjugates).

\begin{lemma}\label{pff:index}
Let $w\in\mathsf A$ and $H_w=\{\kappa\in K^{\mathrm{ad}}:\mathrm{red}(\kappa)\in\pbFix(\mathsf d_w)\}$.
Then $[K^{\mathrm{ad}}:H_w]=p_0^{\,r}-1$.
\end{lemma}
\begin{proof}
As $\mathrm{red}$ is onto, $[K^{\mathrm{ad}}:H_w]=[\auG:\pbFix(\mathsf d_w)]$. The map $s\mapsto
s_w$ (or $s\mapsto s_w^\vee$, if $\mathsf d_w\in F_{\aubar w}$) maps $\auG$ onto $\auGL_r(\auF_{p_0})$,
and $\pbFix(\mathsf d_w)$ is the preimage of the stabiliser of the non-zero vector $\mathsf d_w$; as
$\auGL_r(\auF_{p_0})$ acts transitively on the $p_0^r-1$ non-zero vectors, the index is
$p_0^r-1$.
\end{proof}

\begin{proposition}[One place]\label{pff:oneplace}
Let $w\in\mathsf A$, let $T'=T\cap\{w,\aubar w\}$, and use the notation of \S\ref{pff:ss-setting}
with $T'$ in place of $T$. Assume (L1)--(L3) at $T'$, $b\ge1000$, $k\ge3\cdot10^{10}n^4$ and $r\ge
k+2$. Let $\pffZ_w$ be the image in $G_{T'}$ of $H_w$. Then $\pffZ_w$ is an intermediate group, it
contains the components at $T'$ of all $\kappa\in K^{\mathrm{ad}}$ with $\mathrm{red}(\kappa)\in
\pbFix(\mathsf A,\mathsf d)$, and
\[
\pffm(\pffZ_w)\,\pbvol(K_{T'})\,b^{6k}\le\pffm(K_{T'})\,\pbvol(\pffZ_w).
\]
\end{proposition}
\begin{proof}
Let $p\colon K^{\mathrm{ad}}\to G_{T'}$ be the projection; its image is $K_{T'}$ (strong
approximation for the lattice: a point of $K_{T'}$ is the component of an element of $G_f$ that is
integral elsewhere, which lies in $K^{\mathrm{ad}}$). If $p(\kappa)\in K_{T'}(M)$, then $\kappa-1$
vanishes modulo $w$, so $\mathrm{red}(\kappa)_w=1$ and $\kappa\in H_w$. Hence $p^{-1}(K_{T'}(M))
\subseteq H_w$, so $K_{T'}(M)\subseteq p(H_w)=\pffZ_w\subseteq K_{T'}$, and $\ker p\subseteq H_w$
gives $[K_{T'}:\pffZ_w]=[K^{\mathrm{ad}}:H_w]=p_0^r-1$ (Lemma~\ref{pff:index}). The second claim
is clear, as $\pbFix(\mathsf A,\mathsf d)\subseteq\pbFix(\mathsf d_w)$.

Put $B=b^9$; by Lemma~\ref{pff:nonunits}, the non-units of $\auO_{T'}$ have volume at most
$B^{-1}\le\frac12$. Let $t=k-10^{10}n^4$; then $t+1+n+2\cdot10^9n^4\le k$, $6k\le9t$ and $t+1<r$.

\emph{Small volume.} As $B\le b^{10}\le p_0$ and $t+1<r$, $B^{t+1}\le p_0^{t+1}\le p_0^r-1=
[K_{T'}:\pffZ_w]$, so by Lemma~\ref{pff:vol}(c)
\begin{equation}\label{pff:eq-rare}
B^{t+1}\,\pbvol(\pffZ_w)\le\pbvol(K_{T'}).
\end{equation}
\emph{Upper bound for $\pffZ_w$.} A $c\in\auO_{T'}$ with $c^{-1}\notin\auO_{T'}$ generates an ideal
of non-units, so $|c|_{T'}=\pbvol(c\auO_{T'})\le B^{-1}$. By Theorem~\ref{pff:massZ} with $Q=B$
and the exponent $t+1+n$, by $B^{-n}\le2^{-n}$ and by Lemma~\ref{pff:vol}(b),
\[
\pffm(\pffZ_w)\le\pbvol(\pffZ_w)\pbvol(K_{T'}(M))\bigl(\pbvol(\pffZ_w)+B^{-t-1}\,\pbvol(K_{T'})\bigr).
\]
\emph{Lower bound for $K_{T'}$.} By Theorem~\ref{pff:massZ}, $\pbvol(K_{T'})^2\pbvol(K_{T'}(M))\le
\pffm(K_{T'})$.

\emph{Conclusion.} By \eqref{pff:eq-rare} and $B\ge2$, $\bigl(\pbvol(\pffZ_w)+B^{-t-1}\pbvol(K_{T'})
\bigr)B^t\le2B^{-1}\pbvol(K_{T'})\le\pbvol(K_{T'})$ (each term times $B^{t+1}$ is at most
$\pbvol(K_{T'})$). As $b^{6k}\le b^{9t}=B^t$,
\[
\pffm(\pffZ_w)\pbvol(K_{T'})b^{6k}\le\pbvol(\pffZ_w)\pbvol(K_{T'}(M))\pbvol(K_{T'})^2\le\pffm(K_{T'})
\pbvol(\pffZ_w).\qedhere
\]
\end{proof}

\begin{lemma}[Products]\label{pff:products}
Let $T=T_1\sqcup\dots\sqcup T_s$ be a partition into sets stable under conjugation, and assume
(L1)--(L3) at $T$ and at each $T_i$. Then $G_T=\prod_iG_{T_i}$, and there is a constant $\lambda>0$
such that for all closed subgroups $Y_i\subseteq K_{T_i}$ containing $K_{T_i}(M)$ and $Y=\prod Y_i$
\[
\pbvol(Y)=\lambda\prod_i\pbvol(Y_i),\qquad\pffm(Y)=\lambda^3\prod_i\pffm(Y_i).
\]
\end{lemma}
\begin{proof}
The decomposition of $E_T$ as $\prod E_{T_i}$ gives $G_T=\prod_iG_{T_i}$, with $K_T=\prod K_{T_i}$,
$K_T(M)=\prod K_{T_i}(M)$, $U_T=\prod U_{T_i}$, and $\psi_T=\prod\psi_{T_i}$, $\sigma$ being additive;
hence $\chi_{M}^{(T)}(g)=\prod_i\chi_M^{(T_i)}(g_i)$. Both $\pbvol_T$ and $\prod_i\pbvol_{T_i}$ are Haar
measures on $G_T$, so $\pbvol_T=\lambda\prod_i\pbvol_{T_i}$. Then $\int_Y\chi_M(yxy^{-1})dy=\lambda
\prod_i\int_{Y_i}\chi_M^{(T_i)}(y_ix_iy_i^{-1})dy_i$, and integrating the square of its absolute
value over $x$ gives another factor $\lambda$ by Fubini's theorem.
\end{proof}

\begin{theorem}[The gain for fixed vectors]\label{pff:fixed}
Assume (L1)--(L3) at $T$ and at its subsets stable under conjugation, $b\ge1000$, $4f\le b$,
$k\ge3\cdot10^{10}n^4$, and that $T$ contains the places above $M$. For every set $\mathsf A$ of
chosen places above $p_0$ and every family $\mathsf d$ of frame vectors there is an intermediate
group $\pffZ=\pffZ(\mathsf A,\mathsf d)$ such that
\begin{enumerate}
\item[(a)] $\pffZ$ contains the component at $T$ of every $\kappa\in K^{\mathrm{ad}}$ with
$\mathrm{red}(\kappa)\in\pbFix(\mathsf A,\mathsf d)$;
\item[(b)] $\pffm(\pffZ)\,\pbvol(\pffK)\,b^{6k\#\mathsf A}\le\pffm(\pffK)\,\pbvol(\pffZ)$.
\end{enumerate}
\end{theorem}
\begin{proof}
A place of $T$ lies in $\{w,\aubar w\}$ for at most one $w\in\mathsf A$ (the conjugate of a chosen
place is not chosen). Partition $T$ into the sets $T_w=T\cap\{w,\aubar w\}$, $w\in\mathsf A$, which
are non-empty as $T$ contains the places above $M$, and the rest $T_0$. On $T_w$ take $\pffZ_w$ of
Proposition~\ref{pff:oneplace}, on $T_0$ take $K_{T_0}$ (for which (b) holds with the factor $1$),
and let $\pffZ=\prod_w\pffZ_w\times K_{T_0}$, an intermediate group. (a) holds componentwise. By
Lemma~\ref{pff:products} applied to $Y=\pffZ$ and to $Y=\pffK$, the inequality (b) is the product of
the inequalities of Proposition~\ref{pff:oneplace}, the constants $\lambda^3\cdot\lambda$ being the
same on both sides.
\end{proof}

\subsection{The forms at $T$ and the exact statements of \textup{(F)} and \textup{(S)}}\label{pff:ss-statements}

\subsubsection*{The forms}
Let $\pffZ$ be an intermediate group and $\pffW$ its Bessel vector. For heads $p,q$ of length $i$
(tuples of units of $\pffET$) and $0\le i\le\varrho$, let $\pffI_i(p,q)=\pffI_i(p,q;\pffW,1)$ be the
\emph{form of floor $i$} of the cut Bessel function $\pffW$, with weight $1$
(\S\ref{pfr:ss-setting}); for heads of length $n$ with a common last entry, let
$\mathcal D(p,q)=\mathcal D_{\pi_N}(p,q;\pffZ)$ be the \emph{density at the last entry}
(Definition~\ref{pff:lastentry}): the integral, over the points of the quadric $\{y\in\pbmXp\otimes
\pffET:h(y,y)=-N\}$ with last entry $h(y,z)=p_\varrho$, of the pairing of the scaled orbital
integrals of $\pffW$ at $w_{p'}\sigma_y$ and $w_{q'}\sigma_y$ ($p',q'$ the first $\varrho$ entries),
for the unit $\pi_N$ of $\pffET$ with $\pi_N\aubar\pi_N=N$. The \emph{forms at $T$} are
\begin{equation}\label{pff:eq-forms}
J^T_i(p,q)=\frac{\pffI_i(p,q)}{\pffr(\pffZ)}\quad(0\le i\le\varrho),\qquad J^T_n(p,q)=
\frac{\mathcal D(p,q)}{\pffr(\pffZ)} .
\end{equation}
They depend on $T$ and on $\pffZ$; in \textup{(F)} and \textup{(S)} the group is
$\pffZ=\pffZ(\mathsf A,\mathsf d)$ of Theorem~\ref{pff:fixed}. This is the description of
\S\ref{pbm:ss-remains}, item \textup{(F)}: the vector of translates of $\chi_M$ by the cosets of
$\pffKM$ in $\pffZ$, divided by $\pbvol(\pffZ)\pbvol(\pffKM)/\pffl^{\,2}$, and $\pffl^{-1}=\#\pbUb$
for the data (the volume of $U_T\cap\pffKM$ for the measure that gives $U_T\cap\pffK$ volume $1$ is
$[U_T\cap\pffK:U_T\cap\pffKM]^{-1}$, and reduction identifies this quotient with $\pbUb$).

\begin{lemma}[The form of floor $0$]\label{pff:J0}
$\pffI_0(p,q)=\int_{\pffG}|k_0(x)|^2dx=\pffm(\pffZ)/\pffl^{\,2}$ for all $p,q$; hence
$0\le J^T_0=\pffm(\pffZ)/(\pbvol(\pffZ)\pbvol(\pffKM))$.
\end{lemma}
\begin{proof}
At floor $0$ the head density is $1$, $\pfrm_0=\delta_1$, and $\theta_{\ge0}=\theta_{\ge1}\ast
\theta_0$, so $\Phi_0(x;1;1)=\bigl|\int\mathsf k^1_1(\sigma,\sigma x)\,d\theta_0(\sigma)\bigr|^2=|k_0(x)|^2$
and $\pfrI_0(p,q)=\int|k_0|^2\,d\pbvol$; the second is
Corollary~\ref{pff:k0}; the third follows from \eqref{pff:eq-mass}.
\end{proof}

\begin{theorem}[The mass in \textup{(F)}]\label{pff:mass}
Under the hypotheses of Theorem~\ref{pff:fixed} and $2^{-\varrho}\le b^{-6}$, for
$\pffZ=\pffZ(\mathsf A,\mathsf d)$ and all heads $p,q$,
\[
0\le J^T_0(p,q)\le(1+b^{-6})\,b^{-6k\#\mathsf A}.
\]
\end{theorem}
\begin{proof}
By Lemma~\ref{pff:J0}, Theorem~\ref{pff:fixed}(b) and Corollary~\ref{pff:massK},
\[
J^T_0=\frac{\pffm(\pffZ)}{\pbvol(\pffZ)\pbvol(\pffKM)}\le\frac{b^{-6k\#\mathsf A}\,\pffm(\pffK)}{
\pbvol(\pffK)\pbvol(\pffKM)}\le(1+b^{-6})\,b^{-6k\#\mathsf A}.\qedhere
\]
\end{proof}

\begin{remark}\label{pff:rem-E4}
For $\mathsf A=\emptyset$ the theorem says that the mass of $\pffK$ exceeds its main part
$\pbvol(\pffK)^2\pbvol(\pffKM)$, which comes from the positions in $\pffKM$, by at most the factor
$1+b^{-6}$ (up to $\pbvol(\pffK)\le1+2^{-n}$). The positions in $\pffK\smallsetminus\pffKM$ are among
those treated by Proposition~\ref{pff:outside}; on them the statement is of the same nature as
Theorem~\ref{pbm:E4}, which is however not used here. The gain for fixed vectors does not come from
a bound for the sum of $f_{\mathsf D}$ over a stabiliser as suggested in Remark~\ref{pbm:rem-E4}(b),
but from the index $p_0^r-1$ of the stabiliser of a vector, which makes $\pbvol(\pffZ)$ small, while
the mass integral of $\pffZ$ is controlled by its main part $\pbvol(\pffZ)^2\pbvol(\pffKM)$ and an
error that is small relative to $\pbvol(\pffK)$.
\end{remark}

\subsubsection*{The exact statements}
Recall $\Lambda^T$, $\beta^T$ of \eqref{pbm:eq-beta}, and the place $\mathfrak P_T=(\pffET,dx,\psi_T,
w\mapsto M^{-1}w^{-1})$ of Theorem~\ref{pbm:sieve}. A head of length $i$ is a tuple $p=(p_0,\dots,
p_{i-1})$ of units of $\pffET$, and $(p,c)$ is $p$ followed by $c$.

\begin{theorem}[Local system at $T$; Theorem~\ref{pbm:inputs}\,\textup{(F)}]\label{pff:F}
Assume the hypotheses on the data of Theorem~\ref{pbm:L2}. Let $T$ be a finite set of finite places
of $E$, stable under conjugation, containing the places above $MN$. Let $\mathsf A$ be a set of
chosen places above $p_0$ and $\mathsf d$ a family of frame vectors, and let $\pffZ=\pffZ(\mathsf A,
\mathsf d)$ be the intermediate group of Theorem~\ref{pff:fixed}. Then the forms
\eqref{pff:eq-forms} for $\pffZ$ satisfy:
\begin{enumerate}
\item[(a)] $\pffZ$ contains the component at $T$ of every $\kappa\in K^{\mathrm{ad}}$ with
$\mathrm{red}(\kappa)\in\pbFix(\mathsf A,\mathsf d)$;
\item[(b)] $(J^T_i)_{0\le i\le n}$ is a sieve system at $\mathfrak P_T$
(Definition~\ref{pbm:system}) with mass $(1+b^{-6})\,b^{-6k\#\mathsf A}$, weights $\Lambda^T$ and
decay $\beta^T$.
\end{enumerate}
\end{theorem}

\begin{proof}[Proof of Theorem~\ref{pff:F}]
(a) is Theorem~\ref{pff:fixed}(a). For (b): \emph{(S4)} is Theorem~\ref{pff:mass}.
\emph{(S2), (S3)} are Proposition~\ref{pff:top}(b), (c), divided by $\pffr(\pffZ)>0$.
\emph{(S1).} For $i<\varrho$ the identity and the integrability are Lemma~\ref{pfr:descent}, which
applies as $\int\pfrF^+_{i+1}d\pbvol<\infty$ (Theorem~\ref{pfr:shells} with $e=1$, the right side being
finite); for $i=\varrho$ they are Proposition~\ref{pff:top}(a). In both cases the twisted integral is the
form of floor $i$ with the weight $\psi_T(\xi h(gf_i,f_i))$, and for a unit $\xi$
Theorem~\ref{pff:decay}(i) with $t=k-n^7$, $K=n^6$ (admissible: $10^{10}n^4\le n^6$ and $n^7+10^{10}n^4\le
k$ for the data) bounds it by $\pffr(\pffZ)\pbmn(\xi)^{(k-n^7)/2}\pbmden(\xi)^{-6}(2/|M|_T)^{n^6}
\Lambda^T_i(p)\Lambda^T_i(q)$. \emph{(S5)} is the same with Theorem~\ref{pff:decay}(ii), $j=i$ (the floor
$i$ form for heads of length $i$ with $\xi=M^{-1}(p_{i-1}^{-1}-q_{i-1}^{-1})$). Finally $\pbmden(\xi)^{-6}\le
\pbmden(\xi)^{-3}$ and $(2/|M|_T)^{n^6}=(2M^{2f})^{n^6}\le(2b^{44f})^{n^6}\le b^{n^9f}$ (as $M\le b^{22}$,
$2\le b$ and $44f+1\le n^3f$), which gives the decay $\beta^T$; dividing by $\pffr(\pffZ)$ gives
\eqref{pff:eq-forms}.
\end{proof}

\begin{theorem}[Siegel's step; Theorem~\ref{pbm:inputs}\,\textup{(S)}]\label{pff:S}
Assume the hypotheses of Theorem~\ref{pff:F}, and in addition that $T$ contains the places at which
an element of the set $\mathsf B$ of \S\ref{pbm:ss-cell} is not a unit. Let $J^\tau_n$ be the top
forms at the real places of Theorem~\ref{pfr:R}. Let $(c_j,c'_j,s_j,\sigma_j)_{j\in\mathsf I}$ be a
finite family of pairwise distinct quadruples of classes $c_j,c'_j$, elements $s_j\in\pbFix(\mathsf
A,\mathsf d)$ and $\sigma_j\in\pbmGf(E^+)$ such that $t_j=t(\sigma_j)\in\pbmT$, $|\tau(t_j)|>N$
for all $\tau$, and $h_{c_j}^{-1}\sigma_jh_{c'_j}$ stabilises the lattice at every finite place
outside $T$. Then
\[
\pbmrS^{\,2}\sum_{j\in\mathsf I}\Bigl\|\sum_{p\in\pbmT^\varrho}\pbmOrb_{c_j,s_j,c'_j}(w_p\sigma_j)
\Bigr\|^2\le(1+b^{-6})\,\pbmaS\,\Rel\sum_{t\in\{t_j\}}\ \sum_{p,q\in\pbmT^\varrho}\prod_\tau J^\tau_n
\bigl([p,t]_\tau,[q,t]_\tau\bigr)\,J^T_n\bigl((p,t),(q,t)\bigr),
\]
with $J^T_n$ the top form \eqref{pff:eq-forms} for the group $\pffZ(\mathsf A,\mathsf d)$ of
Theorem~\ref{pff:F}.
\end{theorem}

\begin{remark}[Comparison with Theorem~\ref{pbm:inputs}, and the constants]\label{pff:rem-L2}
(1) \emph{The constants of \textup{(F)}.} The mass $(1+b^{-6})b^{-6k\#\mathsf A}$, the weights
$\Lambda^T_i(p)=\prod_{l\le i}[\delta_l\in\pffOT]\,\pbmn(\delta_l)^{n^6/2}$, the decay
$\beta^T(\xi)=\pbmn(\xi)^{(k-n^7)/2}\pbmden(\xi)^{-3}b^{n^9f}$, and the place $\mathfrak P_T$ (with
$\iota(w)=M^{-1}w^{-1}$, so that the condition (S5) concerns $\xi=M^{-1}(p_{i-1}^{-1}-q_{i-1}^{-1})$)
are exactly those of Theorem~\ref{pbm:inputs}\,\textup{(F)} and of Theorem~\ref{pbm:sieve}, which
Theorem~\ref{pbm:L2} applies. Here $\pbmn(\xi)=|\xi|_T^{-1}$ is the generalised index $[\pffOT:\xi
\pffOT]$ and $\pbmden(\xi)=[\pffOT:\pffOT\cap\xi^{-1}\pffOT]$. The decay of Theorem~\ref{pff:decay} carries the
factor $(2/|M|_T)^{n^6}=(2M^{2f})^{n^6}$ (as $|M|_T=M^{-2f}$), which is bounded by $b^{n^9f}$ because
$M\le b^{22}$ and $44f+1\le n^3f$; this is the factor $b^{n^9f}$ of $\beta^T$, and the power
$\pbmden(\xi)^{-6}$ of Theorem~\ref{pff:decay} is weakened to $\pbmden(\xi)^{-3}$.

(2) \emph{The constant of \textup{(S)}.} With $\#\pbR=\#\pbHb\,\#\pbUb$ and $\#\pbUb=M^{f\varrho(2k+
2\varrho+1)}$ (the unipotent radical has dimension $\varrho(2k+2\varrho+1)$, and reduction modulo $M$ gives
$[U_T\cap\pffK:U_T\cap\pffKM]=|\pffOT^+/M\pffOT^+|^{\dim U}$ with $|\pffOT^+/M\pffOT^+|=M^f$, $\pffOT^+$ the
integers of $E^+_T$),
\[
\pbmaS=\#\pbR\cdot\#\pbUb\cdot\#\mathcal C\cdot\Bigl(\frac{\pi^k}{(k-1)!\,N^k}\Bigr)^f\aurd^{\,f/2}
\pbmvol^{\,k},\qquad\aurd^{\,f/2}=|\operatorname{disc}E|^{1/4},\qquad\pbmvol=(2/\aurd)^f=\pbvol(
\widehat\auO_E),
\]
where $\pbvol(\widehat\auO_E)$ is the volume of $\widehat\auO_E$ for the Haar measure of $\auA_{E,f}$ which,
together with Lebesgue measure on each $E_\tau=\C$, gives $E\backslash\auA_E$ volume $1$
(Lemma~\ref{pfg:tamagawa}(a)); up to the factor $1+b^{-6}$ it is
(number of classes) $\times\#\pbR\#\pbUb\times$ (Siegel's constant). The
normalisation is $\pbmrS=d^{-1}\prod_\tau\bigl(\sqrt{\pbmcG}\prod_as_{\tau,a}^{k-1+2(\varrho-a)}\bigr)$. Lemma~\ref{pbm:const} uses exactly $\pbmaS$, $\pbmrS$ and the mass
$(1+b^{-6})b^{-6k\#\mathsf A}$; so the constants transcribed in Theorem~\ref{pbm:inputs} are right.

(3) \emph{Two additions.} (a) In Theorem~\ref{pbm:inputs}, the functions $J^T_i$ are said to
depend on $T,\mathsf A,\mathsf d$; they do so through the group $\pffZ(\mathsf A,\mathsf d)$, and
\textup{(S)} holds for the top form of \emph{that} group: the elements $s_j\in\pbFix(\mathsf A,
\mathsf d)$ are reductions of elements $\kappa_j\in K^{\mathrm{ad}}$, whose components at $T$ lie in
$\pffZ$ by (a) of Theorem~\ref{pff:F}, and this is what Siegel's step uses. (b) $T$ must contain
the places at which the lattice basis is not integral with integral inverse; these lie above $N$
(\S\ref{pff:ss-setting}), so the hypothesis of Theorem~\ref{pbm:inputs}\,\textup{(F)} already
contains it. With these readings the proof of Theorem~\ref{pbm:L2} applies verbatim: it uses
\textup{(F)} and \textup{(S)} for the same $T,\mathsf A,\mathsf d$.
\end{remark}

\subsection{Siegel's step: proof of Theorem~\ref{pff:S}}\label{pff:ss-siegel}

Throughout, $T$, $\mathsf A$, $\mathsf d$, $\pffZ=\pffZ(\mathsf A,\mathsf d)$ are as in
Theorem~\ref{pff:S}, $C=\pffKM\backslash\pffZ$ is the set of cosets, and for a class $c$ we write
$h_c\in H(\auA_f)$ for its representative. For a rational $\gamma$ we write $\gamma_T$ for its component
at $T$ and put $g(c,c',\sigma)=(h_c^{-1}\sigma h_{c'})_T\in\pbmGf_T$.

\subsubsection*{Siegel's formula and the density at the finite places}
Let $\mathcal C=H(E^+)\backslash H(\auA_f)/J_X$ be the set of classes (\S\ref{pb:ss-setting}); by
Lemma~\ref{au:free} the stabilisers $H(E^+)\cap hJ_Xh^{-1}$ are trivial. On $\auA_{E,f}^k$ we use the
Haar measure for which $\widehat\auO_E$ has volume $\pbvol(\widehat\auO_E)=\pbmvol=(2/\aurd)^f$ in each
coordinate (the measure of Lemma~\ref{pfg:tamagawa}(a)), and on the self-adjoint finite adeles $\mathrm{Herm}_f$
the Haar measure for which the codifferent $\{z:\psi_{E,f}(z\widehat\auO_E)=1\}\cap\mathrm{Herm}_f$ has
volume $1$. For $\Phi\ge0$ locally constant with compact support on $\auA_{E,f}^k$ and $\alpha\in
\auA_{E,f}$ self-adjoint put
\[
\delta(\Phi,\alpha)=\int_{\mathrm{Herm}_f}\int_{\auA_{E,f}^k}\Phi(x)\,\psi_{E,f}\Bigl(\tfrac12z\,(x^*x-\alpha)
\Bigr)\,dx\,dz ,
\]
whenever the outer integral converges absolutely, and $0$ otherwise.

\begin{theorem}[Siegel's formula]\label{pff:siegelformula}
Let $k\ge3$, let $\Phi\ge0$ be locally constant with compact support on $\auA_{E,f}^k$ and invariant under
$J_X$, and let $\alpha\in E^+$ be totally positive. Then
\[
\sum_{c\in\mathcal C}\ \sum_{v\in E^k,\ v^*v=\alpha}\Phi(h_c^{-1}v)=\#\mathcal C\cdot|\operatorname{disc}E|^{1/4}
\prod_{\tau\in\auPl}\frac{\pi^k\,\tau(\alpha)^{k-1}}{(k-1)!}\cdot\delta(\Phi,\alpha),
\]
the inner sums being finite.
\end{theorem}
This is Siegel's formula for the representations of a totally positive number by the genus of $x^*x$; it is
derived, with every constant, from Weil's Siegel--Weil theorem (Theorem~\ref{pfg:weil}) in
Proposition~\ref{pfg:siegel}: the measures are Tamagawa measures (Lemma~\ref{pfg:tamagawa}), the archimedean
density is computed, and the $z$-integral defining $\delta(\Phi,\alpha)$ converges absolutely (classically: for quadratic forms over $\Q$
\cite[\S7]{pff-Siegel35}, for hermitian forms \cite{pff-Braun41}).

\begin{lemma}[The density off $T$ is at most $(1+b^{-6})\pbmvol^k$]\label{pff:offT}
Let $\Psi\ge0$ be locally constant with compact support on $X_{0,T}=\pffET^k$, let $\Phi=\Psi\circ
\mathrm{pr}_T\cdot\mathbf 1_{\prod_{v\notin T}\auO_v^k}$, and let $\alpha\in E^+$ (identified with its image in
$\auA_{E,f}$). Then, for the data,
\[
|\delta(\Phi,\alpha)|\le(1+b^{-6})\,\pbmvol^{\,k}\int\Psi\,d\nu_{\alpha,T},
\]
where $\nu_{\alpha,T}$ is the fibre measure at $\alpha$ of $x\mapsto x^*x$ on $\pffET^k$
(Definition~\ref{pff:lastentry}).
\end{lemma}
\begin{proof}
\emph{Factorisation.} Write $\auA_{E,f}=\pffET\times\auA^T$, $x=(x_T,x^T)$, $z=(z_T,z^T)$. Because
$\Phi(x)$ depends on $x_T$ and on the condition $x^T$ integral, and because $\psi_{E,f}(\frac12z^T(x^{T*}x^T-
\alpha))$ is $1$ when $z^T$ is in the codifferent and $x^T$ integral, an exact computation (translate $x$
by the integral vectors off $T$ and $z$ by the self-adjoint elements that are in the codifferent off $T$;
the measures are translation invariant and the codifferent has volume $1$) gives
\[
\delta(\Phi,\alpha)=\Bigl(\int_{\mathrm{Herm}^{(T)}}\ \int_{\widehat\auO_E^k}\psi_{E,f}\bigl(\tfrac12z(w^*w-\alpha^{(T)})
\bigr)\,dw\,dz\Bigr)\cdot\int_{\mathrm{Skew}_T}\int_{\pffET^k}\Psi(y)\,\psi_T\bigl(\tau b(y^*y-\alpha)\bigr)\,dy\,db,
\]
where $\mathrm{Herm}^{(T)}$ is the set of self-adjoint $z$ whose components at $T$, multiplied by $-\varpi$,
are integral, $\alpha^{(T)}$ is $\alpha$ with its components at $T$ removed, $\tau\in\pffOT$ has
$\tau+\aubar\tau=1$, and $db$ is the Haar measure of the skew elements of $\pffET$ that gives the skew
integers volume $1$ (the substitution $z_T=-\varpi^{-1}b$ turns $\psi_{E,f}(\frac12 z_T\cdot)$ into
$\psi_T(\tau b\,\cdot)$ on the relevant set).

\emph{The Gauss sums.} For self-adjoint $z$ let $I_z\subseteq\widehat\auO_E$ be the largest ideal on which
$v\mapsto\psi_{E,f}(zv)$ is trivial, and $G(z)=\int_{\widehat\auO_E^k}\psi_{E,f}(\frac12z(w^*w-\alpha^{(T)}))dw$. By
Weyl differencing, $|G(z)|^2=\int_{\widehat\auO^k}\int_{\widehat\auO^k}\aubar{F(w)}F(w+d)\,dw\,dd$ with $F$ the
integrand; the inner integral vanishes unless every $d_i$ lies in $K_z=\{r:\psi_{E,f}(zr\widehat\auO)=1\}$
(translate $w$ by $\aubar v e_i$, $v\in\widehat\auO$, which multiplies it by $\psi_{E,f}(zd_iv)$), so
$|G(z)|^2\le(\pbvol(K_z\cap\widehat\auO)\pbmvol)^k=([\widehat\auO:I_z]^{-1}\pbmvol^2)^k$. The $z$ with $I_z=I$
lie in at most $[\widehat\auO:I]$ cosets of the codifferent (one for each character of $\widehat\auO/I$), each
of volume $1$. Hence, for $\mu$ with $2(\mu+1)\le k$,
\[
\int|G(z)|\,dz\le\sum_I[\widehat\auO:I]\cdot[\widehat\auO:I]^{-k/2}\pbmvol^k\le\zeta_{\widehat\auO}(\mu)\,
\pbmvol^k,\qquad\zeta_{\widehat\auO}(\mu)=\sum_I[\widehat\auO:I]^{-\mu}.
\]
Take $\mu=3\varrho+4$ ($2(\mu+1)\le k$ for the data). As $\zeta_{\widehat\auO}(\mu)\le\zeta_E(\mu)$ and
$\log\zeta_E(\mu)\le2f\cdot2^{2-\mu}$, we get $\zeta_{\widehat\auO}(\mu)\le1+2^{-\varrho}\le1+b^{-6}$ for
the data (as in \S\ref{pff:ss-large}).

\emph{The density at $T$.} For $\Psi\ge0$ locally constant with compact support,
$\bigl|\int\!\int\Psi(y)\psi_T(\tau b(y^*y-\alpha))\,dy\,db\bigr|\le\int\Psi\,d\nu_{\alpha,T}$. Indeed, if the
outer integral does not converge absolutely the left side is $0$ by convention. Otherwise, cut the
frequency to $b\in u\aubar u\,\mathcal O$ for a unit $u$ ($\mathcal O$ the skew integers): by orthogonality of
characters the cut integral equals $\pbvol(u\aubar u\mathcal O)$ times the mass, for $dy$, of the slab
$\{y:y^*y-\alpha\in(u\aubar u)^{-1}\mathcal O_{E^+,T}\}$ weighted by $\Psi$, where $\mathcal O_{E^+,T}$ is the lattice of self-adjoint integers, dual to $\mathcal O$ under
$(b,X)\mapsto\psi_T(\tau bX)$ since $E/E^+$ is unramified (so the two volumes multiply to $1$); this mass, by the disintegration, is the
mean of $\int\Psi\,d\nu_{\alpha+\beta,T}$ over $\beta$ in that slab. As $u\to0$ the cut integrals tend
to the full integral (dominated convergence) and the means tend to $\int\Psi\,d\nu_{\alpha,T}$
(continuity of $\beta\mapsto\int\Psi\,d\nu_{\beta,T}$, Lemma~\ref{pff:fibre}(a)); so the integral equals
$\int\Psi\,d\nu_{\alpha,T}\ge0$.

Multiplying the two bounds gives the lemma.
\end{proof}

\begin{proposition}[Siegel's formula as an upper bound]\label{pff:siegelsum}
Let $t\in E$ with $|\tau(t)|>N$ for all $\tau$, $\alpha=(t\aubar t-N^2)/N$, and let $\Psi\ge0$ on
$V_T=Y_0\otimes\pffET$ be locally constant with compact support and invariant under $H_T\cap\pffKM$.
For a class $c$ and $v\in E^k$ let $x(c,v)=h_c^{-1}v$. Then for every finite set $\mathcal S$ of pairs
$(c,v)$ with $Nv^*v=t\aubar t-N^2$ and $x(c,v)$ integral outside $T$,
\[
\sum_{(c,v)\in\mathcal S}\Psi\bigl(x(c,v)_T-N^{-1}tz\bigr)\le(1+b^{-6})\,\#\mathcal C\,\mathsf c_S\,G_t\cdot
|\pi_N|_T^{-1}\int\Psi\bigl(y-N^{-1}tz\bigr)\,d\nu_{\alpha,T}(y),
\]
with $\mathsf c_S=(\pi^k/((k-1)!N^k))^f|\operatorname{disc}E|^{1/4}\pbmvol^k$ and $G_t=\prod_\tau(|\tau(t)|^2-
N^2)^{k-1}$.
\end{proposition}
\begin{proof}
Let $\Phi(x)=\Psi(x_T-N^{-1}tz)\,\mathbf 1[x\text{ integral off }T]$ on $\auA_{E,f}^k$. It is locally
constant with compact support; it is invariant under $J_X$: an element $\kappa\in J_X$ is integral off $T$
(so preserves the integrality condition), and at $T$ it acts on the vectors $x-N^{-1}tz$ as an element of
$H_T\cap\pffKM$ (Lemma~\ref{pb:KR}: $K_H\cap J=J_X$), under which $\Psi$ is invariant. The pairs of
$\mathcal S$ are among the terms of the left side of Theorem~\ref{pff:siegelformula} (for each class, the
$v$ with $v^*v=\alpha$ and $\Phi(h_c^{-1}v)=\Psi(\cdots)$), and all terms are $\ge0$; so the left side is at
most $\#\mathcal C|\operatorname{disc}E|^{1/4}\prod_\tau\pi^k\tau(\alpha)^{k-1}/(k-1)!\cdot|\delta(\Phi,\alpha)|$.
By Lemma~\ref{pff:offT} (with $\Psi(\cdot-N^{-1}tz)$), $|\delta|\le(1+b^{-6})\pbmvol^k\int\Psi(y-N^{-1}tz)
d\nu_{\alpha,T}$. Finally $\prod_\tau\tau(\alpha)^{k-1}=N^{-(k-1)f}G_t$, and the product formula for the
$T$-unit $N$ ($N^{-1}$ is integral off $T$) gives $|\pi_N|_T^{-1}=|N|_T^{-1/2}=N^f$; so
$N^{-(k-1)f}=N^{-kf}\cdot|\pi_N|_T^{-1}$, which is the stated constant.
\end{proof}

\subsubsection*{From class sums to the fibre pairing}

\begin{lemma}[The real places]\label{pff:topreal}
Let $\tau\in\auPl$, $\sigma\in\pbmGf_\tau$ and $x$ a head with $h(\sigma z,z)=x_\varrho$ (in the normalised
basis). Then $\pfrip{\pfrN_\sigma(x)}{\pfrN_\sigma(x')}=J^\tau_n(x,x')$ whenever $x_\varrho=x'_\varrho$.
\end{lemma}
\begin{proof}
$\sigma^{-1}z$ and $\sigma_x^{-1}z$ have the same norm and the same pairing with $z$, so their components in
the span of the $x_i$ have the same norm, and by Witt's theorem there is $h\in\pfrH$ with $\sigma^{-1}z=
h\sigma_x^{-1}z$. Then $h_2=\sigma_xh^{-1}\sigma^{-1}$ fixes $z$ and the $e_a,e_a^*$, so $h_2\in\pfrH$ and
$\sigma=h_2^{-1}\sigma_xh^{-1}$. As $w_x$ commutes with $\pfrH$ and $W^\flat(h'y)=\pi(h')W^\flat(y)$ for
$h'\in\pfrH$, the substitution $y\mapsto h_2^{-1}y$ gives $\pfrN_\sigma(x)=(\pi(h_2)^{-1}\otimes\aubar{\pi(h)})
\pfrN_{\sigma_x}(x)$, with a unitary operator that does not depend on $x$. Take inner products and use
the definition of $J^\tau_n$ in \eqref{pfr:eq-c}.
\end{proof}

\begin{lemma}[Orbital integrals at $T$ are Kloosterman sums]\label{pff:locKl}
Let $\sigma\in\pbmGf(E^+)$, $c,c'$ classes, $p$ a head of units of $\pbmRT$ (entries and inverses integral
off $T$), $\gamma=w_p\sigma$ with $h_c^{-1}\gamma h_{c'}$ integral off $T$, $l=g(c,c',\sigma)$, and let
$\kappa\in K$ and cosets $a,b\in C$ with $\kappa_T^{-1}a_{\vphantom{1}}b^{-1}\in\pffKM$. Then
\[
\mathcal O_p(l^{-1})_{(a,b)}\cdot O(\gamma)=\pfrkap_\varrho(p,1)\,\pbvol(\pffKM)\,\pbmOrb_{c,\bar\kappa,c'}(\gamma),
\]
where $\mathcal O_p$ is the scaled orbital integral of \S\ref{pff:ss-floors}, $O(\gamma)$ the orbital integral
at the real places (\S\ref{pbm:ss-cell}), and $\bar\kappa$ the reduction of $\kappa$.
\end{lemma}
\begin{proof}
By Proposition~\ref{pbm:orbit}, $\pbmOrb_{c,\bar\kappa,c'}(\gamma)=\pbmKl_{c,\bar\kappa,c'}(\gamma)O(\gamma)$, and by
the remark after it $\pbmKl$ is the sum of $\psi_f(u)\aubar{\psi_f(u')}$ over $(u,u')\in(U(\auA_f)/U(\auA_f)\cap
h_cJh_c^{-1})\times(\cdots)$ with $h_c^{-1}u^{-1}\gamma u'h_{c'}\in K$ of reduction $\bar\kappa$. Off $T$ all
data are integral and the characters trivial, so the sum is over $U_T/(U_T\cap\pffKM)$ in both variables,
with the condition $u^{-1}(w_pl)u'\in\kappa_T\pffKM$ at $T$ (the classes being represented in $H$, which
normalises $U$). On the other hand, by definition,
\[
\int_{\pffG}\pffW(y)_a\aubar{\pffW((w_pl)^{-1}y)_b}\,dy=\int_{\pffG}\pffch(ya^{-1})\aubar{\pffch(l^{-1}w_p^{-1}yb^{-1})}dy .
\]
The first factor vanishes unless $ya^{-1}=uk$ ($u\in U_T$, $k\in\pffKM$); write $y=uka$ and integrate over
$u\in U_T/(U_T\cap\pffKM)$ and $k\in\pffKM$. The second factor is non-zero iff $l^{-1}w_p^{-1}uka b^{-1}\in
u'\pffKM$, i.e.\ $u'^{-1}(w_pl)^{-1}u\in\pffKM\,ba^{-1}=(\kappa_T\pffKM)^{-1}$ (using $\kappa_T^{-1}ab^{-1}\in\pffKM$
and normality), and then the integrand is $\psi_T(M^{-1}\sigma(u))\aubar{\psi_T(M^{-1}\sigma(u'))}$. The
$k$-integral gives $\pbvol(\pffKM)$; so the integral is $\pbvol(\pffKM)$ times the Kloosterman sum above (the
inversion $(u,u')\leftrightarrow$ the pairs of the global sum being the identity on the condition). Multiply by
$\pfrkap_\varrho(p,1)$.
\end{proof}

\begin{lemma}[One position]\label{pff:headsum}
Let $c,c'$ be classes and $\sigma\in\pbmGf(E^+)$ with $h_c^{-1}\sigma h_{c'}$ integral off $T$, let $\mathcal P$ be
a finite set of heads of units of $\pbmRT$, $\rho_p>0$ with $\rho_p\pfrkap_\varrho(p,1)=w$ for all $p$, and let
$R$ be the Gram matrix of the vectors $\rho_pO(w_p\sigma)$. Let $\kappa_i\in K$ ($i\in I$) have distinct
reductions $s_i$ and components $\kappa_{i,T}\in\pffZ$. Then, with $l=g(c,c',\sigma)$,
\[
\pbvol(\pffKM)^2\,\#C\sum_{i\in I}w^2\Bigl\|\sum_{p\in\mathcal P}\pbmOrb_{c,s_i,c'}(w_p\sigma)\Bigr\|^2\le
\Rel\sum_{p,q}R_{pq}\pfrip{\mathcal O_p(l^{-1})}{\mathcal O_q(l^{-1})}=:\mathcal H(l).
\]
\end{lemma}
\begin{proof}
$\mathcal H(l)=\sum_{\mathbf x\in C\times C}v_{\mathbf x}^*Rv_{\mathbf x}$ with $v_{\mathbf x}=(\mathcal
O_p(l^{-1})_{\mathbf x})_p$, a sum of non-negative terms ($R$ is a Gram matrix). For $i\in I$ and $b\in C$ let
$A_i(b)$ be the coset with $\kappa_{i,T}^{-1}bA_i(b)^{-1}\in\pffKM$ (it exists as $\kappa_{i,T}\in\pffZ$). For
$\mathbf x=(b,A_i(b))$, Lemma~\ref{pff:locKl} gives $\sum_pv_{\mathbf x,p}\rho_pO(w_p\sigma)=w\,\pbvol(\pffKM)\sum_p
\pbmOrb_{c,s_i,c'}(w_p\sigma)$, so $v_{\mathbf x}^*Rv_{\mathbf x}=w^2\pbvol(\pffKM)^2\|\sum_p\pbmOrb\|^2$. The pairs
$(b,A_i(b))$ are distinct for distinct $(i,b)$: if $A_i(b)=A_{i'}(b)$ then $\kappa_{i,T}^{-1}\kappa_{i',T}\in
\pffKM$, so (as both are integral off $T$, hence in $K$) $\kappa_i^{-1}\kappa_{i'}\in J$ and $s_i=s_{i'}$. Summing
over $i$ and the $\#C$ cosets $b$ gives the claim.
\end{proof}

\begin{lemma}[Separation]\label{pff:sep}
Let $(c,c_1',\sigma_1)$, $(c,c_2',\sigma_2)$ have the same left class $c$, with $\sigma_i\in\pbmGf(E^+)$,
$h(\sigma_1z,z)=h(\sigma_2z,z)$, the same $X_0$-component of $\sigma_iz$, $h_c^{-1}\sigma_ih_{c_i'}$ integral off
$T$, and $g(c,c_1',\sigma_1)^{-1}g(c,c_2',\sigma_2)\in\pffKM$. Then $c_1'=c_2'$ and $\sigma_1=\sigma_2$.
\end{lemma}
\begin{proof}
$\sigma_1z=\sigma_2z$ (same $X_0$-part and same pairing with $z$; the $z$-coordinate is determined by the
latter), so $\nu=\sigma_1^{-1}\sigma_2\in\pbmGf(E^+)$ fixes $z$: $\nu\in H(E^+)$. Moreover $h_{c_1'}^{-1}\nu
h_{c_2'}=(h_c^{-1}\sigma_1h_{c_1'})^{-1}(h_c^{-1}\sigma_2h_{c_2'})$ is integral off $T$ and in $\pffKM$ at $T$, so in
$J\cap H(\auA_f)=J_X$ (Lemma~\ref{pb:KR}). Hence $c_1'=c_2'$ in $\mathcal C$, and then $\nu\in H(E^+)\cap
h_{c_1'}J_Xh_{c_1'}^{-1}=1$ by Lemma~\ref{au:free}.
\end{proof}

\begin{lemma}[Fibre pairing]\label{pff:fibresum}
Let $R$ be positive semidefinite, $\mathcal P$ a finite set of heads, and $l_\theta\in\pbmGf_T$ ($\theta\in\Theta$) with
the same $e=l_\theta^{-1}z$ and in distinct right cosets of $\pffKM$ ($l_\theta l_{\theta'}^{-1}\in\pffKM\Rightarrow
\theta=\theta'$). Then $\mu_H(H_T\cap\pffKM)\sum_\theta\mathcal H(l_\theta^{-1})^{\vphantom{1}}\le\Rel\sum_{p,q}
R_{pq}\Pi_{p,q}(e)$, with $\mathcal H$ formed with the $\mathcal O_p(\cdot)$ as in Lemma~\ref{pff:headsum}.
\end{lemma}
\begin{proof}
Fix $l_0$ with $l_0^{-1}z=e$; then $l_\theta=k_\theta l_0$ with $k_\theta\in H_T$. By definition
$\sum R_{pq}\Pi_{p,q}(e)=\int_{H_T}\phi(k)\,d\mu_H(k)$ with $\phi(k)=\sum_{p,q}R_{pq}\pfrip{\mathcal O_p(kl_0)}
{\mathcal O_q(kl_0)}\ge0$. For $k'\in H_T\cap\pffKM$, $\mathcal O_p(k'kl_0)=\mathcal O_p(kl_0)$: by
definition $\mathcal O_p(k'kl_0)=\pfrkap_\varrho(p,1)\int_{\pffG}\pffW(y)\otimes\overline{\pffW(k'kl_0w_p^{-1}y)}\,dy$, and
$\pffW(k'x)=\pffW(x)$ for $k'\in H_T\cap\pffKM$ by \textup{(W2)} in the setting of Proposition~\ref{pfg:fibrepairing}
($H_T$ normalises $U_T$ and fixes the frame). So $\phi$ is constant on the cosets
$(H_T\cap\pffKM)k_\theta$, where it equals $\mathcal H(l_\theta^{-1})$. These cosets are disjoint, and each has
measure $\mu_H(H_T\cap\pffKM)$ by right invariance.
\end{proof}

\begin{lemma}[The normalised mass]\label{pff:rhoD}
Put $D=\#\pbHb\cdot\#\pbUb^{\,2}$. Then $\pffr(\pffZ)\le\mu_H(H_T\cap\pffKM)\,\pbvol(\pffKM)^2\,\#C\cdot D$, and
$\pbmaS=D\cdot\#\mathcal C\cdot\mathsf c_S$.
\end{lemma}
\begin{proof}
$\pbvol(\pffZ)=\#C\,\pbvol(\pffKM)$ (Lemma~\ref{pff:vol}(c)); $\pffl^{-1}=\#\pbUb$ (\S\ref{pff:ss-statements}); and
$\mu_H(H_T\cap\pffK)=1$ with $[H_T\cap\pffK:H_T\cap\pffKM]=\#\pbHb$ (reduction, Lemma~\ref{pb:KR}), so
$\mu_H(H_T\cap\pffKM)\#\pbHb=1$. Hence $\pffr(\pffZ)=\#C\,\pbvol(\pffKM)^2\#\pbUb^2=\mu_H(H_T\cap\pffKM)\pbvol(\pffKM)^2\#C
\cdot D$. The second identity is $\#\pbR=\#\pbHb\#\pbUb$ and Remark~\ref{pff:rem-L2}(2).
\end{proof}

\begin{lemma}[The weights]\label{pff:weights}
For a head $p\in\pbmT^\varrho$ with last entry $t$, let
\[
\rho_p=\prod_\tau\frac{\sqrt{\pbmcG}(|x_{\tau,\varrho}|^2-1)^{(k-1)/2}}
{d_\tau}\prod_{a<\varrho}|x_{\tau,a}|^{k+1+2(\varrho-a-1)}
\]
with $x_\tau=[p,t]_\tau$ and $d=\prod_\tau d_\tau$. Then
$\rho_p\,\pfrkap_\varrho(p,1)=\pbmrS\sqrt{G_t}$.
\end{lemma}
\begin{proof}
Put $d_a=k+1+2(\varrho-a-1)=k-1+2(\varrho-a)$. Then $\pfrkap_\varrho(p,1)=\prod_{a<\varrho}|p_a|_T^{d_a/2}$, and $|x_{\tau,a}|=
s_{\tau,a}|\tau(p_a)|$. The product formula for the $T$-units $p_a$ gives $|p_a|_T\prod_\tau|\tau(p_a)|^2=1$, so
$\prod_\tau\prod_{a<\varrho}|x_{\tau,a}|^{d_a}\cdot\pfrkap_\varrho(p,1)=\prod_\tau\prod_{a<\varrho}s_{\tau,a}^{d_a}$. And $s_{\tau,\varrho}=
N^{-1}$, so $|x_{\tau,\varrho}|^2-1=(|\tau(t)|^2-N^2)/N^2$ and $\prod_\tau(|x_{\tau,\varrho}|^2-1)^{(k-1)/2}=\sqrt{G_t}
\prod_\tau s_{\tau,\varrho}^{k-1}$. Comparing with the definition of $\pbmrS$ in \eqref{pbm:eq-consts} gives the
claim.
\end{proof}

\begin{proof}[Proof of Theorem~\ref{pff:S}]
Group the family by the last entry $t\in\{t_j\}$; it suffices to prove the inequality for the subfamily
$\mathsf I_t$ with last entry $t$, the right side being additive in $t$. Fix $t$, let $\mathcal P=\pbmT^\varrho$
and, for $p\in\mathcal P$, let $\tilde p=(p,t)$; put $R_{pq}=\prod_\tau J^\tau_n([p,t]_\tau,[q,t]_\tau)$. By
Lemma~\ref{pff:topreal} and the multiplicativity of inner products of tensors over $\tau$, for every $j\in
\mathsf I_t$ the matrix $R$ is the Gram matrix of the vectors $\rho_pO(w_p\sigma_j)$ ($O$ the orbital integral
of the cut Bessel function at the real places, with $\rho_p$ of Lemma~\ref{pff:weights}); in particular $R$ is
positive semidefinite. For $j\in\mathsf I_t$ choose $\kappa_j\in K$ with reduction $s_j$; by
Theorem~\ref{pff:F}(a), $\kappa_{j,T}\in\pffZ$.

Let $v_j\in E^k$ be the $X_0$-component of $\sigma_jz$; then $Nv_j^*v_j=t\aubar t-N^2$ \eqref{pbm:eq-last}, and
$g_j=g(c_j,c'_j,\sigma_j)$ satisfies $g_jz=x(c_j,v_j)_T-N^{-1}tz=:e_j$, read in the lattice basis; $x(c_j,
v_j)$ is integral off $T$ because $h_{c_j}^{-1}\sigma_jh_{c_j'}$ and the lattice basis are. Group $\mathsf I_t$ by
the pair $(c_j,v_j)$. For a fixed pair, group further by $(c'_j,\sigma_j)$: by Lemma~\ref{pff:sep} the elements
$l=g^{-1}$ for distinct $(c',\sigma)$ lie in distinct right cosets of $\pffKM$ ($l_1l_2^{-1}=g_1^{-1}g_2\in\pffKM$
forces equality), and all satisfy $l^{-1}z=e$; and for fixed $(c,c',\sigma)$ the $s_j$ are distinct.
By Lemma~\ref{pff:headsum} with $w=\pbmrS\sqrt{G_t}$ (Lemma~\ref{pff:weights}) and Lemma~\ref{pff:fibresum},
\[
\mu_H(H_T\cap\pffKM)\pbvol(\pffKM)^2\#C\sum_{j:\,(c_j,v_j)=(c,v)}\pbmrS^{\,2}G_t\Bigl\|\sum_p\pbmOrb_j\Bigr\|^2\le\Rel\sum_{p,q}
R_{pq}\Pi_{\tilde p,\tilde q}(e)=:\Psi_t(e),
\]
where $\pbmOrb_j=\pbmOrb_{c_j,s_j,c_j'}(w_p\sigma_j)$. By Lemma~\ref{pff:fibre}(b), $\Psi_t$ agrees on the quadric with a
function $\Psi\ge0$ (the positive part of the real part of $\sum R_{pq}\Phi_{\tilde p,\tilde q}$) that is locally
constant with compact support and invariant under $H_T\cap\pffKM$. By Lemma~\ref{pff:rhoD}, the left side is at
least $\pffr(\pffZ)D^{-1}\pbmrS^{\,2}G_t\sum_{j}\|\sum_p\pbmOrb_j\|^2$. Summing over the pairs $(c,v)$ and applying
Proposition~\ref{pff:siegelsum},
\[
\begin{aligned}\frac{\pffr(\pffZ)}D\pbmrS^{\,2}G_t\sum_{j\in\mathsf I_t}\Bigl\|\sum_p\pbmOrb_j\Bigr\|^2&\le(1+b^{-6})\#\mathcal C\mathsf c_SG_t\,
|\pi_N|_T^{-1}\int\Psi(y-N^{-1}tz)\,d\nu_{\alpha,T}(y)\\&=(1+b^{-6})\#\mathcal C\mathsf c_SG_t\,\Rel\sum_{p,q}R_{pq}\mathcal
D(\tilde p,\tilde q),\end{aligned}
\]
the last equality because $\Rel\sum R_{pq}\Pi_{\tilde p\tilde q}(e)\ge0$ on the quadric (the entrywise product of the positive semidefinite matrices $R$ and $(\Pi_{pq}(e))$ is positive semidefinite, as in the proof of Proposition~\ref{pff:top}(b)), so that there $\Psi$ is this function itself, and because almost every point $y-N^{-1}tz$ of the sphere lies on the quadric
(Definition~\ref{pff:lastentry}). Divide by $G_t\pffr(\pffZ)/D$ and use $D\#\mathcal C\mathsf c_S=\pbmaS$
(Lemma~\ref{pff:rhoD}) and $J^T_n=\mathcal D/\pffr(\pffZ)$.
\end{proof}
\subsection{The inputs of \textup{(F)} and \textup{(S)} proved in \S\ref{pfg:sec}}\label{pff:ss-pending}

The points used above that are proved in \S\ref{pfg:sec} are: Lemma~\ref{pff:twogauges}
(Corollary~\ref{pfg:twogaugesT}), the axioms (G0)--(G5) at $T$ (Lemma~\ref{pfg:gauge}), Lemma~\ref{pff:fibre}
(Lemma~\ref{pfg:coarea}, Proposition~\ref{pfg:fibrepairing}), and Theorem~\ref{pff:siegelformula}, derived from Weil's
theorem (Proposition~\ref{pfg:siegel}). The published results used are Theorem~\ref{pfg:weil}, Witt's theorem for hermitian forms
(in the form of Lemma~\ref{pbp:tb-witt}) and the Smith normal form.


%% ===== proofs/B-frames-finite2.tex =====
% proofs/B-frames-finite2.tex
% The finite places, second part: the open points of B-frames-finite (two gauges, the counts (G4),
% (G5), the fibre measures and the step function, Siegel's formula). Labels of this file start with pfg:.
\providecommand{\auO}{\mathcal O}
\providecommand{\auA}{\mathbb A}
\providecommand{\auU}{\mathrm U}
\providecommand{\auGL}{\mathrm{GL}}
\providecommand{\auPl}{\Phi}
\providecommand{\aurd}{\mathrm{rd}_E}
\providecommand{\aubar}[1]{\overline{#1}}
\providecommand{\panot}[1]{\par\smallskip{\raggedright\noindent\textbf{[Not proved here:} #1\textbf{]}\par}\smallskip}
\providecommand{\pbvol}{\operatorname{vol}}
\providecommand{\pbmvol}{\mathsf v}
\providecommand{\pffET}{E_T}
\providecommand{\pffOT}{\auO_T}
\providecommand{\pffG}{G_T}
\providecommand{\pffK}{K_T}
\providecommand{\pffKM}{K_T(M)}
\providecommand{\pffch}{\chi_M}
\providecommand{\pffW}{W_Z}
\providecommand{\pffZ}{\mathsf Z}
\providecommand{\pfrK}{\mathsf K}
\providecommand{\pfrG}{\mathsf G}
\providecommand{\pfrU}{\mathsf U}
\providecommand{\pfrR}{\mathsf R}
\providecommand{\pfrCl}{\mathcal C}
\providecommand{\pfrkap}{\kappa}
\providecommand{\pfrdep}{\operatorname{dp}}
\providecommand{\pfrip}[2]{\langle#1,#2\rangle}
\providecommand{\pbmXp}{X^{+}}
% macros of this file (prefix \pfg)
\providecommand{\pfgH}{\mathrm{Herm}_T}
\providecommand{\pfgS}{\mathrm{Skew}_T}
\providecommand{\pfgOh}{\auO_{\mathrm h}}
\providecommand{\pfgOs}{\auO_{\mathrm s}}
\providecommand{\pfgE}{\mathsf e}
\providecommand{\pfgD}{\mathsf D}
\providecommand{\pfgEb}{\mathcal E}
\providecommand{\pfgQ}{\mathsf Q_N}
\providecommand{\pfgz}{\zeta_{\pffOT}}

\section{The finite places: gauges at the top, fibres and Siegel's formula}\label{pfg:sec}

This section proves the lemmas used in \S\ref{pff:sec} and listed in \S\ref{pff:ss-pending}. We use the
notation of \S\ref{pff:sec} and of \S\ref{pfr:sec} (place data, gauges, the axioms (G0)--(G5) of
\S\ref{pfr:ss-gauge} in their present form, heads, depth ideals $\mathfrak d_a$, nearness), and cite their
statements by label.

\emph{Contents of this section.}
\begin{itemize}
\item[(i)] \emph{Two gauges} (\S\ref{pfg:ss-two}): Theorem~\ref{pfg:twogauges}, a statement for every place
datum and every pair of gauges, and its case at $T$, Corollary~\ref{pfg:twogaugesT}, which is
Lemma~\ref{pff:twogauges} exactly as stated there. The proof is complete: the comparison of lengths is
carried down from $Y^{(\varrho)}$ to the whole space through short duals, and below the length $j$ of the
heads the short duals come from units of $\mathfrak d_a$ produced by Cramer's rule
(Lemmas~\ref{pfg:cramer}, \ref{pfg:headunits}). The off-diagonal block $h(se_b^*,e_c)$ is not needed.
\item[(ii)] \emph{Fourier analysis at $T$} (\S\ref{pfg:ss-fourier}): the duality of self-adjoint and skew
elements, a bound for Gauss sums of residue classes (Lemma~\ref{pfg:gauss}), and the construction of
the fibre measures of a unimodular hermitian lattice, with their continuity, the disintegration and a
bound for the mass of residue classes (Proposition~\ref{pfg:fibres}).
\item[(iii)] \emph{The gauges at $T$} (\S\ref{pfg:ss-gauge}): Lemma~\ref{pfg:gauge}, a complete proof of
Lemma~\ref{pff:gauge}, including the axioms (G0) and the clause $\tau_0\delta^{-1}\in\Gamma(1)$ of (G1),
and the counts (G4) and (G5) written out in full.
\item[(iv)] \emph{The fibre measures at the last entry and the step function} (\S\ref{pfg:ss-fibre}):
Lemma~\ref{pfg:coarea} (Lemma~\ref{pff:fibre}(a)) and Proposition~\ref{pfg:fibrepairing} (Lemma~\ref{pff:fibre}(b)); what
Proposition~\ref{pff:top}(c) uses is that $\Pi_{p',q'}=\Pi_{p,q}$ for \emph{all} $(p',q')$ near $(p,q)$.
\item[(v)] \emph{Siegel's formula} (\S\ref{pfg:ss-siegel}): we state Weil's Siegel--Weil theorem for the
dual pair $(\auU(1,1),\auU(V))$ in the range of absolute convergence (Theorem~\ref{pfg:weil}, cited) and
derive from it Theorem~\ref{pff:siegelformula} with every constant (Proposition~\ref{pfg:siegel}): the
measures are checked to be Tamagawa measures, and the archimedean density is computed. So the
normalisation of Theorem~\ref{pff:siegelformula} is verified directly.
\end{itemize}
The only results cited without proof are Theorem~\ref{pfg:weil} and the Smith normal form over a principal
ideal domain; Witt's theorem is used in the form of Lemma~\ref{pbp:tb-witt}, proved in the paper.

\subsection{Conventions}\label{pfg:ss-conv}
The hermitian form is written $h(u,v)=u^*\mathsf Jv$, conjugate linear in $u$ (the opposite of
convention (Cv1) of \S\ref{pb:ss-setting}; by Remark~\ref{pfg:rem-conv} the statements of this section
hold for either), with Gram matrix $\mathsf J$ of the
basis of \S\ref{pfr:ss-setting}; vectors are columns of coordinates in that basis, and matrices act on them.
For a vector $\omega$ we write $\omega_{e^*_a}$ for its coordinate on $e^*_a$ and $\omega_{e_a}$ for its
coordinate on $e_a$; since $e_a$ is orthogonal to every basis vector except $e^*_a$ and $h(e_a,e^*_a)=1$,
\[
\omega_{e^*_a}=h(e_a,\omega).
\]
Put $\pfgEb_{<l}=\langle e_b:b<l\rangle$. A vector $\omega$ is \emph{orthogonal below $l$} if $h(e_b,\omega)=0$
for $b<l$, i.e.\ $\omega_{e_b^*}=0$ for $b<l$; and $\omega\equiv\omega'\ (\mathrm{mod}\ \pfgEb_{<l})$ means
$\omega-\omega'\in\pfgEb_{<l}$. The module $Y^{(l)}$ of \S\ref{pfr:ss-gauge} is the quotient of the vectors
orthogonal below $l$ by $\pfgEb_{<l}$. Sizes refer to a gauge (\S\ref{pfr:ss-gauge}); we use the
consequences of (G0) listed there: a sum of at most $r+1$ terms of size $i$ has size $i+1$, a product of
numbers of sizes $i,i'$ has size $i+i'$, a product of matrices (or a matrix and a vector) of sizes $i,i'$ has
size $i+i'+1$, $h(u,v)$ has size $i+i'+3$ for $u,v$ of sizes $i,i'$ (as $\mathsf J$ has size $1$), the
inverse of an isometry of size $i$ has size $i+4$, and an $l\times l$ determinant ($l\le r+1$) of entries of
size $i$ has size $l(i+1)$. Sizes are monotone: size $i$ implies size $i'$ for $i'\ge i$ (as $1\in\Gamma(i'-i)$).

We use two properties of the groups of \S\ref{pfr:ss-setting}.
\begin{enumerate}
\item[(R)] An element $\xi\in\pfrR^{(l)}$ satisfies $\xi e_a\in e_a+\pfgEb_{<a}$ for $a<l$, and so does
$\xi^{-1}$. Hence $\xi$ preserves each $\pfgEb_{<a}$ ($a\le l$), maps vectors orthogonal below $a$ to
vectors orthogonal below $a$ for $a\le l$ (as $h(e_b,\xi\omega)=h(\xi^{-1}e_b,\omega)$), and does not
change the coordinate on $e^*_a$ of a vector orthogonal below $a$, for $a<l$.
\item[(U)] The Bessel unipotent group $\pfrU$ is contained in $\pfrR^{(\varrho)}$, and every $u\in\pfrU$ acts
trivially on $Y^{(\varrho)}$: $u\omega\equiv\omega\ (\mathrm{mod}\ \pfgEb_{<\varrho})$ for every $\omega$
orthogonal below $\varrho$. (The group $\pfrU$ is block unitriangular for the flag $\pfgEb_{<1}\subset\dots
\subset\pfgEb_{<\varrho}\subset\pfgEb_{<\varrho}^\perp\subset\dots$, with middle block $1$.) In coordinates:
for a vector $w$ orthogonal below $\varrho$, $(u-1)w$ has non-zero coordinates only on the $e_a$, $a<\varrho$.
\end{enumerate}

\subsection{Two gauges}\label{pfg:ss-two}

Throughout, $\mathfrak F$ is a place datum (\S\ref{pfr:ss-setting}) with level $M$, box $\pfrK$ and close pairs
$\pfrCl$, and sizes refer to a gauge $\Gamma$ for $\mathfrak F$ unless another gauge is named.

\begin{definition}\label{pfg:def}
Let $g,g'\in\pfrG$, let $r$ be a unit with $\bar r=r$ and $r\in\Gamma(1)$, $l\le\varrho$, $a<\varrho$ and
$t,c\ge0$.
\begin{enumerate}
\item[(a)] \emph{The lengths at $g$ and $g'$ differ by $t$ modulo $r$ on $Y^{(l)}$}, written $\pfgD_r(g,g';l,t)$,
if for every $i\ge0$ and every $\omega$ orthogonal below $l$ with $g^{-1}\omega$ of size $i$ there is
$\omega'\equiv\omega\ (\mathrm{mod}\ \pfgEb_{<l})$ such that $r^{-1}(g^{-1}\omega-g'^{-1}\omega')$ has size
$i+t$.
\item[(b)] $e_a$ has a \emph{short lift of size $c$ at $g$} if some $f\equiv e_a\ (\mathrm{mod}\
\pfgEb_{<a})$ has $g^{-1}f$ of size $c$.
\item[(c)] $e_a$ has a \emph{short dual of size $c$ for $(g,g')$ modulo $r$} if there are $\zeta$ orthogonal
below $a$ with $\zeta_{e^*_a}=1$ and $\zeta'\equiv\zeta\ (\mathrm{mod}\ \pfgEb_{<a})$ such that $g^{-1}\zeta$ and
$r^{-1}(g^{-1}\zeta-g'^{-1}\zeta')$ have size $c$.
\end{enumerate}
\end{definition}

\begin{lemma}[The pairing with a dual]\label{pfg:pairing}
Let $\zeta$ be orthogonal below $a$ with $\zeta_{e^*_a}=1$, $\zeta'\equiv\zeta\ (\mathrm{mod}\ \pfgEb_{<a})$, and
let $\omega$ be orthogonal below $a+1$ and $\omega'\equiv\omega\ (\mathrm{mod}\ \pfgEb_{<a+1})$. Then
$h(\zeta',\omega')-h(\zeta,\omega)=\mu$, the coordinate of $\omega'-\omega$ on $e_a$.
\end{lemma}
\begin{proof}
$\omega'-\mu e_a\equiv\omega\ (\mathrm{mod}\ \pfgEb_{<a})$. The vector $\zeta'$ is orthogonal below $a$ (as
$\zeta$ is and the $e_b$ are isotropic and mutually orthogonal), so $h(\zeta',\omega'-\mu e_a)=h(\zeta',\omega)$;
and $h(\zeta',\omega)=h(\zeta,\omega)$ as $\zeta'-\zeta\in\pfgEb_{<a}$ and $\omega$ is orthogonal below $a$.
Finally $h(\zeta',e_a)=h(\zeta,e_a)=\overline{\zeta_{e^*_a}}=1$ (again because $\zeta'-\zeta\in\pfgEb_{<a}$).
Hence $h(\zeta',\omega')-\mu=h(\zeta,\omega)$.
\end{proof}

\begin{lemma}[One floor down]\label{pfg:step}
Let $a<\varrho$, $r$ as in Definition~\ref{pfg:def}, and suppose that $e_a$ has short lifts of size $c$ at
$g$ and at $g'$ and a short dual of size $c$ for $(g,g')$ modulo $r$. Then
$\pfgD_r(g,g';a+1,t)$ implies $\pfgD_r(g,g';a,t+4c+12)$.
\end{lemma}
\begin{proof}
Let $f,f'$ be the short lifts at $g,g'$ and $\zeta,\zeta'$ the short dual. Let $\omega$ be orthogonal below $a$
with $g^{-1}\omega$ of size $i$. Put $\alpha=\omega_{e^*_a}$. As $f\equiv e_a\ (\mathrm{mod}\ \pfgEb_{<a})$ and
$\omega$ is orthogonal below $a$, $\alpha=h(e_a,\omega)=h(f,\omega)=h(g^{-1}f,g^{-1}\omega)$, of size $c+i+3$.
Put $\omega_0=\omega-\alpha\zeta$; it is orthogonal below $a+1$ (below $a$ since $\omega,\zeta$ are, and
$(\omega_0)_{e^*_a}=\alpha-\alpha=0$), and $g^{-1}\omega_0=g^{-1}\omega-\alpha g^{-1}\zeta$ has size
$i_0:=2c+i+4$. By $\pfgD_r(g,g';a+1,t)$ there is $\omega_0'\equiv\omega_0\ (\mathrm{mod}\ \pfgEb_{<a+1})$ with
$r^{-1}(g^{-1}\omega_0-g'^{-1}\omega_0')$ of size $i_0+t$. Let $\mu$ be the coordinate of $\omega_0'-\omega_0$ on
$e_a$. By Lemma~\ref{pfg:pairing} and the invariance of $h$ under $g,g'$,
\[
\mu=h(g'^{-1}\zeta',g'^{-1}\omega'_0)-h(g^{-1}\zeta,g^{-1}\omega_0).
\]
Write $v=g^{-1}\zeta$, $v'=g'^{-1}\zeta'$, $w=g^{-1}\omega_0$, $w'=g'^{-1}\omega_0'$; then $v$, $r^{-1}(v-v')$ have
size $c$ and $w$, $r^{-1}(w-w')$ have size $i_0+t$. Since $\bar r=r$,
$r^{-1}(h(v',w')-h(v,w))=-h(v',r^{-1}(w-w'))-h(r^{-1}(v-v'),w)$, and $v'=v-r\cdot r^{-1}(v-v')$ has size $c+2$; so
$r^{-1}\mu$ has size $c+i_0+t+6$. Put
\[
\omega'=\alpha\zeta'+\omega_0'-\mu f' .
\]
Then $\omega'\equiv\omega\ (\mathrm{mod}\ \pfgEb_{<a})$: indeed $\omega_0'-\mu f'\equiv\omega_0'-\mu e_a\equiv\omega_0$
(the coordinates of $\omega_0'-\omega_0\in\pfgEb_{<a+1}$ on $e_b$, $b<a$, do not matter, and the one on $e_a$ is
$\mu$), and $\zeta'\equiv\zeta$. Finally
\[
r^{-1}(g^{-1}\omega-g'^{-1}\omega')=r^{-1}(g^{-1}\omega_0-g'^{-1}\omega'_0)+\alpha\,r^{-1}(g^{-1}\zeta-g'^{-1}\zeta')
+(r^{-1}\mu)\,g'^{-1}f' ,
\]
a sum of three vectors of sizes $i_0+t$, $(c+i+3)+c$ and $(c+i_0+t+6)+c$, each at most
$i+t+4c+10$; so the sum has size $i+t+4c+11\le i+t+(4c+12)$.
\end{proof}

\begin{corollary}[Down to the whole space]\label{pfg:whole}
Let $j\le\varrho$, and suppose that for every $a<j$ the vector $e_a$ has short lifts of size $c$ at $g$ and at
$g'$ and a short dual of size $c$ for $(g,g')$ modulo $r$. Then $\pfgD_r(g,g';j,t)$ implies
$\pfgD_r(g,g';0,t+(4c+12)j)$.
\end{corollary}
\begin{proof}
Apply Lemma~\ref{pfg:step} for $a=j-1,j-2,\dots,0$.
\end{proof}

\begin{lemma}[From lengths to the matrix]\label{pfg:matrix}
If $\pfgD_r(g,gx;0,i)$, then $r^{-1}(x-1)$ has size $i+4$.
\end{lemma}
\begin{proof}
For a basis vector $e$ put $\omega=ge$; then $g^{-1}\omega=e$ has size $0$. On $Y^{(0)}$ the only lift of $\omega$
is $\omega$ itself, so $r^{-1}(e-x^{-1}e)$ has size $i$: the matrix $N=r^{-1}(x^{-1}-1)$ has size $i$. As $x$ is
an isometry, $x^{-1}=\mathsf J^{-1}x^*\mathsf J$, and applying this to $x^{-1}$ gives
$x=\mathsf J^{-1}(x^{-1})^*\mathsf J$; with $\bar r=r$ we get $r^{-1}(x-1)=\mathsf J^{-1}N^*\mathsf J$. As
$\mathsf J,\mathsf J^{-1}$ have size $1$ (G0), $N^*\mathsf J$ has size $i+2$ and $\mathsf J^{-1}N^*\mathsf J$ size
$i+4$.
\end{proof}

\begin{lemma}[The last quotient]\label{pfg:top}
Let $u,v\in\pfrU$ with $uy\in\pfrK$, $vyx\in\pfrK$, and $M^{-1}(vyx-uy)$ of size $2$. Then
$\pfgD_M(y,yx;\varrho,9)$.
\end{lemma}
\begin{proof}
$M$ is a unit with $\bar M=M$ and $M\in\Gamma(1)$ (\S\ref{pfr:ss-setting}, (G1)). Let $\omega$ be orthogonal below
$\varrho$ with $y^{-1}\omega$ of size $i$, and put $\omega'=v^{-1}u\,\omega$; by (U),
$\omega'\equiv\omega\ (\mathrm{mod}\ \pfgEb_{<\varrho})$. Then
\[
M^{-1}\bigl(y^{-1}\omega-(yx)^{-1}\omega'\bigr)=(vyx)^{-1}\cdot M^{-1}(vyx-uy)\cdot y^{-1}\omega ,
\]
because $(vyx)^{-1}(vyx-uy)y^{-1}=y^{-1}-(yx)^{-1}v^{-1}u$. Here $(vyx)^{-1}$ has size $1+4=5$ by (G1), so the
right side has size $(5+2+1)+i+1=i+9$.
\end{proof}

\begin{lemma}[Short lifts at a near point]\label{pfg:lift}
If $g$ is near at $j$, then for every $a<j$ the vector $e_a$ has a short lift of size $5$ at $g$.
\end{lemma}
\begin{proof}
Let $\xi\in\pfrR^{(j)}$ with $\xi g\in\pfrK$, and $f=\xi^{-1}e_a$; by (R), $f\equiv e_a\ (\mathrm{mod}\
\pfgEb_{<a})$. And $g^{-1}f=(\xi g)^{-1}e_a$ is a column of $(\xi g)^{-1}$, of size $5$ by (G1), (G0).
\end{proof}

\begin{lemma}[A unit of the depth ideal gives a short dual]\label{pfg:dualunit}
Let $c$ be a unit with $c\in\mathfrak d_a(g,x)$ and $c^{-1}\in\Gamma(1)$. Then $e_a$ has a short dual of size
$2$ for $(g,gx)$ modulo $M$.
\end{lemma}
\begin{proof}
By the definition of $\mathfrak d_a$ there is $\omega$ orthogonal below $a$ with $\omega_{e^*_a}=c$, and
$\omega'\equiv\omega\ (\mathrm{mod}\ \pfgEb_{<a})$ with $(g^{-1}\omega,(gx)^{-1}\omega')\in\pfrCl$. By (G2),
$g^{-1}\omega$ and $M^{-1}(g^{-1}\omega-(gx)^{-1}\omega')$ have size $1$. Put $\zeta=c^{-1}\omega$, $\zeta'=c^{-1}\omega'$:
$\zeta$ is orthogonal below $a$, $\zeta_{e^*_a}=1$, $\zeta'\equiv\zeta$, and $g^{-1}\zeta$,
$M^{-1}(g^{-1}\zeta-(gx)^{-1}\zeta')$ have size $1+1=2$.
\end{proof}

\begin{lemma}[Cramer's rule at a head]\label{pfg:cramer}
Let $j\le\varrho$, $r$ a head of length $j$ for $\Gamma$, $z\in\pfrG$, and $s,\eta\in\pfrR^{(j)}$ with $sz\in\pfrK$
and $\eta w_r^{-1}z\in\pfrK$. Then for $a<j$ the column $se^*_a$ has size $8\varrho+1$.
\end{lemma}
\begin{proof}
This is the column part of Lemma~\ref{pfr:headentries}, with the size read off its proof; we repeat the
argument. $\gamma:=(sz)(\eta w_r^{-1}z)^{-1}=sw_r\eta^{-1}$ has size $1+5+1=7$. Fix a basis index $i$ and let
$Z'$ be the $(a+1)\times(a+1)$ block of $\gamma$ with rows $e_0^*,\dots,e^*_{a-1},i$ and columns
$e_0,\dots,e_a$. As $\eta^{-1}e_d\in e_d+\pfgEb_{<d}$ (R), $Z'=\mathsf L'\mathsf C$ with $\mathsf C$ the block
of $\eta^{-1}$ on $e_0,\dots,e_a$, unitriangular, and $\mathsf L'$ the block of $sw_r$ with the same rows and
columns $e_0,\dots,e_a$; as $w_re_d=r_de_d^*$, $\mathsf L'_{bd}=s_{R(b),e^*_d}r_d$ ($R(b)$ the $b$-th row index).
For $b\le d$, $s_{e^*_b,e^*_d}=h(e_b,se^*_d)=h(s^{-1}e_b,e^*_d)=\delta_{bd}$ by (R). So $\mathsf L'$ is lower
triangular with diagonal $r_0,\dots,r_{a-1},s_{i,e^*_a}r_a$, and $\det Z'=P_{a+1}s_{i,e^*_a}$ with
$P_{a+1}=r_0\cdots r_a$. By (G0), $\det Z'$ has size $8(a+1)$; as $P_{a+1}^{-1}\in\Gamma(1)$ (head),
$s_{i,e^*_a}=P_{a+1}^{-1}\det Z'$ has size $8(a+1)+1\le8\varrho+1$.
\end{proof}

\begin{lemma}[Units of the depth ideals below a head]\label{pfg:headunits}
Let $j\le\varrho$, $j\ne0$, $p,q$ heads of length $j$ for $\Gamma$, $M^{-1}\in\Gamma(1)$, and let $y$, $yx$,
$w_p^{-1}y$, $w_q^{-1}yx$ be near at $j$. Then for every $a<j$ the ideal $\mathfrak d_a(y,x)$ contains a unit $c$
with $c^{-1}\in\Gamma(8\varrho+9)\subseteq\Gamma(10n)$.
\end{lemma}
\begin{proof}
Let $z\in\{y,yx\}$ and $r\in\{p,q\}$ the corresponding head. Let $s\in\pfrR^{(j)}$ with $sz\in\pfrK$; by
Lemma~\ref{pfg:cramer} (with $\eta$ given by the nearness of $w_r^{-1}z$) the column $se^*_a$ has size
$8\varrho+1$, so $z^{-1}e^*_a=(sz)^{-1}(se^*_a)$ has size $5+(8\varrho+1)+1=8\varrho+7$. By (G2) (as
$M^{-1}\in\Gamma(1)$) there is a unit $c$ with $c^{-1}\in\Gamma(8\varrho+9)$ such that $(c\omega_1,c\omega_2)\in\pfrCl$
for all vectors $\omega_1,\omega_2$ of size $8\varrho+7$. Put $\omega=ce^*_a$. It is orthogonal below $a$
($h(e_b,e^*_a)=0$ for $b\ne a$), and $(y^{-1}\omega,(yx)^{-1}\omega)=(c\,y^{-1}e^*_a,c\,(yx)^{-1}e^*_a)\in\pfrCl$; so
$\omega$ is stable at $(y,yx)$ for $a$ (with $\omega'=\omega$), and $c=\omega_{e^*_a}\in\mathfrak d_a(y,x)$. Finally
$8\varrho+9\le10\varrho+10=10n$.
\end{proof}

\begin{theorem}[Two gauges]\label{pfg:twogauges}
Let $\Gamma,\Gamma'$ be gauges for $\mathfrak F$, $j\le\varrho$, $p,q$ heads of length $j$ for $\Gamma$, and if
$j\ne0$ assume $M^{-1}\in\Gamma(1)$ and $\Gamma(10n)\subseteq\Gamma'(1)$. Let $y,x\in\pfrG$ and suppose:
\begin{enumerate}
\item[(a)] there are $u,v\in\pfrU$ with $uy\in\pfrK$, $vyx\in\pfrK$ and $M^{-1}(vyx-uy)$ of size $2$ for $\Gamma'$;
\item[(b)] $w_p^{-1}y$ and $w_q^{-1}yx$ are near at $j$;
\item[(c)] for every $a$ with $j\le a<\varrho$ the ideal $\mathfrak d_a(y,x)$ contains a unit $d_a$ with
$d_a^{-1}\in\Gamma'(1)$.
\end{enumerate}
Then $M^{-1}(x-1)$ has entries in $\Gamma'(99n)$.
\end{theorem}
\begin{proof}
All sizes are for $\Gamma'$ except where $\Gamma$ is named. \emph{Units for every $a<\varrho$.} For $a\ge j$ take
$d_a$ of (c). For $a<j$ (so $j\ne0$): $y$ and $yx$ are near at $j$ by (a), since $\pfrU\subseteq\pfrR^{(\varrho)}
\subseteq\pfrR^{(j)}$; with (b), Lemma~\ref{pfg:headunits} for $\Gamma$ gives a unit $d_a\in\mathfrak d_a(y,x)$ with
$d_a^{-1}\in\Gamma(10n)\subseteq\Gamma'(1)$. \emph{Short lifts and duals.} By Lemma~\ref{pfg:dualunit} (for
$\Gamma'$), each $e_a$, $a<\varrho$, has a short dual of size $2\le5$ for $(y,yx)$ modulo $M$; by (a) $y$ and $yx$
are near at $\varrho$, so by Lemma~\ref{pfg:lift} each $e_a$ has short lifts of size $5$ at $y$ and at $yx$.
\emph{Descent.} By Lemma~\ref{pfg:top}, $\pfgD_M(y,yx;\varrho,9)$; by Corollary~\ref{pfg:whole} with $c=5$,
$\pfgD_M(y,yx;0,9+32\varrho)$; by Lemma~\ref{pfg:matrix}, $M^{-1}(x-1)$ has size $13+32\varrho\le99n$.
\end{proof}

\begin{corollary}[Lemma~\ref{pff:twogauges}]\label{pfg:twogaugesT}
At the finite places, with the place datum $\mathfrak F_T$ of \S\ref{pff:ss-large}, the statement of
Lemma~\ref{pff:twogauges} holds for all gauges $\Gamma,\Gamma'$ for $\mathfrak F_T$ with $\pffOT\subseteq\Gamma'(2)$
(in particular for all gauges $\Gamma_\alpha$ of Lemma~\ref{pff:gauge}, whose sets $\Gamma_\alpha(0)$ contain
$\pffOT$).
\end{corollary}
\begin{proof}
We verify (a)--(c) of Theorem~\ref{pfg:twogauges}. (b) is part of ``$y$ active at $x$''; (c) is assumed. For
(a): $\langle\pffW(y),\pffW(yx)\rangle\ne0$ gives a coset $\pffKM a_c$ with $\pffch(ya_c^{-1})\ne0$ and
$\pffch(yxa_c^{-1})\ne0$, so $ya_c^{-1}=u^{-1}k_1$ and $yxa_c^{-1}=v^{-1}k_2$ with $u,v\in U_T$ and
$k_1,k_2\in\pffKM$ \eqref{pff:eq-chi}. Then $uy=k_1a_c$ and $vyx=k_2a_c$ lie in $\pffK$, and
$M^{-1}(vyx-uy)=k_1\cdot M^{-1}(k_1^{-1}k_2-1)\cdot a_c$ is integral, because $k_1^{-1}k_2\in\pffKM$ and
$k_1,a_c\in\pffK$ are integral. So its entries lie in $\pffOT\subseteq\Gamma'(2)$.
\end{proof}

We also record the bound for the unipotent factor and its consequence, used in \S\ref{pfg:ss-fibre}.

\begin{lemma}[Inverse of a unipotent factor]\label{pfg:unipinv}
Let $s\in\pfrU$ and suppose that the columns $se^*_a$, $a<\varrho$, have size $d$. Then $s^{-1}$ has size
$(4d+16)\varrho+5$.
\end{lemma}
\begin{proof}
Take $g=1$, $g'=s^{-1}$, $r=1$ (a unit with $\bar r=r$, $1\in\Gamma(1)$). \emph{Top.} $\pfgD_1(1,s^{-1};\varrho,0)$:
for $\omega$ orthogonal below $\varrho$, $\omega':=s^{-1}\omega\equiv\omega$ by (U), and $\omega-(s^{-1})^{-1}\omega'=0$.
\emph{Lifts.} $f=e_a$ at $1$, and $f'=s^{-1}e_a\equiv e_a\ (\mathrm{mod}\ \pfgEb_{<a})$ (R) at $s^{-1}$, with
$(s^{-1})^{-1}f'=e_a$; both of size $0$. \emph{Duals.} $\zeta=\zeta'=e^*_a$: orthogonal below $a$, coordinate $1$
on $e^*_a$, and $e^*_a-se^*_a$ has size $d+1$. By Corollary~\ref{pfg:whole} with $c=d+1$ and $j=\varrho$,
$\pfgD_1(1,s^{-1};0,(4d+16)\varrho)$, and by Lemma~\ref{pfg:matrix} (with $x=s^{-1}$) $s^{-1}-1$ has size
$(4d+16)\varrho+4$; so $s^{-1}=1+(s^{-1}-1)$ has size $(4d+16)\varrho+5$.
\end{proof}

\begin{corollary}[Size of a point under two cut conditions]\label{pfg:sizehead}
Let $W$ be a cut Bessel function for $\mathfrak F$, $r$ a head of length $\varrho$ for $\Gamma$, $l\in\pfrG^{(\varrho)}$,
and $y\in\pfrG$ with $W(y)\ne0$ and $W(lw_r^{-1}y)\ne0$. Then $y$ has size $60n^2$.
\end{corollary}
\begin{proof}
By the cut condition there are $s,u'\in\pfrU$ with $sy\in\pfrK$ and $u'lw_r^{-1}y\in\pfrK$; $\eta:=u'l$ lies in
$\pfrR^{(\varrho)}$ (which contains $\pfrU$ by (U) and $\pfrG^{(\varrho)}$, whose elements fix the $e_a$). By
Lemma~\ref{pfg:cramer} with $j=\varrho$, the columns $se^*_a$ have size $8\varrho+1$; by Lemma~\ref{pfg:unipinv},
$s^{-1}$ has size $(32\varrho+20)\varrho+5$; so $y=s^{-1}(sy)$ has size $32\varrho^2+20\varrho+7\le60n^2$.
\end{proof}

\subsection{Fourier analysis on hermitian lattices at $T$}\label{pfg:ss-fourier}

We keep the setting of \S\ref{pff:ss-setting}: $\pffET=\prod_{w\in T}E_w$, $T$ stable under conjugation,
$E/E^+$ unramified at every finite place, $|\cdot|_T$ the module, $dx$ the Haar measure with
$\pbvol(\pffOT)=1$, $\psi_T(x)=\psi_{E,f}(-x/\varpi)$.

\begin{lemma}[The character and two elements]\label{pfg:psi}
\begin{enumerate}
\item[(a)] $\psi_T(y+\bar y)=1$ for all $y$, and for $c\in\pffET$: $\psi_T(c\pffOT)=1$ if and only if $c\in\pffOT$.
\item[(b)] There are $\tau_0\in\pffOT$ with $\tau_0+\bar\tau_0=1$ and a unit $\delta$ of $\pffET$ with
$\bar\delta=-\delta$ and $\delta,\delta^{-1}\in\pffOT$; hence $\tau_0\delta^{-1}\in\pffOT$.
\item[(c)] There is $\varpi_1\in\pffOT$ with $\bar\varpi_1=\varpi_1$, a unit of $\pffET$, with $|\varpi_1|_w<1$ for
every $w\in T$.
\end{enumerate}
\end{lemma}
\begin{proof}
(a) The first claim is in \S\ref{pff:ss-setting}. The standard character $\psi_{E,f}=\psi_{\Q,f}\circ\mathrm{Tr}$ has
$\psi_{E,f}(y\widehat\auO_E)=1$ exactly for $y$ in the inverse different; at $T$ the different is generated by
$\varpi$, so $\psi_T(c\pffOT)=\psi_{E,f}(-c\varpi^{-1}\pffOT)=1$ iff $c\varpi^{-1}\in\varpi^{-1}\pffOT$, i.e.
$c\in\pffOT$. (b) $\tau_0$ is the element $\tau$ of Lemma~\ref{pff:L1}. For $\delta$ we work at a pair
$\{w,\bar w\}\subseteq T$. If $w\ne\bar w$, the components at $w,\bar w$ are exchanged by the involution, and
$\delta=(1,-1)$ works. If $w=\bar w$, $E_w/E^+_v$ is an unramified quadratic extension, the involution induces
the non-trivial automorphism of the residue extension, and for a unit $\epsilon$ of $\auO_w$ whose residue is
not in the residue field of $E^+_v$, $\delta=\epsilon-\bar\epsilon$ is skew and has non-zero residue, so it is a
unit of $\auO_w$. (c) At a pair $w\ne\bar w$ take $(\varpi_v,\varpi_v)$ for a uniformiser $\varpi_v$ of $E^+_v$; at
$w=\bar w$ a uniformiser of $E^+_v$, which is one of $E_w$ (unramified).
\end{proof}

The trace-one element $\tau_0$ and the skew unit $\delta$ are fixed from now on; they are the elements of axiom
(G1). Put $\pfgH=\{\beta\in\pffET:\bar\beta=\beta\}$, $\pfgS=\{b:\bar b=-b\}$, $\pfgOh=\pfgH\cap\pffOT$,
$\pfgOs=\pfgS\cap\pffOT$, with the Haar measures $d\beta$, $db$ giving $\pfgOh$, $\pfgOs$ volume $1$, and the
pairing
\[
\pfgE(b,\beta)=\psi_T(\tau_0b\beta)\qquad(b\in\pfgS,\ \beta\in\pfgH).
\]

\begin{lemma}[Duality]\label{pfg:duality}
\begin{enumerate}
\item[(a)] Every $x\in\pffOT$ can be written $x=\tau_0h+s$ with $h=x+\bar x\in\pfgOh$, $s\in\pfgOs$, and
$x=\tau_0\delta h_1+h_2$ with $h_1=\delta^{-1}(x-\bar x)\in\pfgOh$, $h_2\in\pfgOh$. $\psi_T(\gamma)=1$ for
$\gamma\in\pfgH$.
\item[(b)] $\pfgE$ is a continuous bicharacter, trivial on $\pfgOs\times\pfgOh$; for $b\in\pfgS$, $\pfgE(b,\cdot)\equiv1$
on $\pfgOh$ implies $b\in\pfgOs$; for $\beta\in\pfgH$, $\pfgE(\cdot,\beta)\equiv1$ on $\pfgOs$ implies $\beta\in
\pfgOh$. For $c\in\pfgH$, $\pfgE(b,c\beta)=\pfgE(cb,\beta)$.
\item[(c)] For $b\in\pfgS$ and $z\in\pffET$: $\pfgE(b,z+\bar z)=\psi_T(bz)$.
\item[(d)] The map $\pfgS\times\pfgS\to\pffET$, $(b_1,b_2)\mapsto\tau_0\delta^{-1}b_1+b_2$, is an isomorphism of
topological groups carrying $\pfgOs\times\pfgOs$ onto $\pffOT$, and so carries $db\,db$ to $dx$. Consequently, for
a unit $u$ of $\pffET$, $\pbvol\{b\in\pfgS:ub\in\pffOT\}^2\le|u|_T^{-1}$; and for a unit $c\in\pfgH$,
multiplication by $c$ multiplies $db$ and $d\beta$ by $|c|_T^{1/2}$.
\end{enumerate}
\end{lemma}
\begin{proof}
(a) $s=x-\tau_0(x+\bar x)=\bar\tau_0x-\tau_0\bar x$ is skew and integral. $x-\bar x=\tau_0\delta h_1-
\overline{\tau_0\delta h_1}=(\tau_0+\bar\tau_0)\delta h_1=\delta h_1$, and $h_2=x-\tau_0(x-\bar x)=\bar\tau_0x+\tau_0\bar x$
is self-adjoint and integral. For $\gamma\in\pfgH$, $\gamma=\tau_0\gamma+\overline{\tau_0\gamma}$ is a trace, so
$\psi_T(\gamma)=1$. (b) Continuity and multiplicativity are clear, and $\tau_0b\beta\in\pffOT$ for integral $b,\beta$.
If $\pfgE(b,\cdot)\equiv1$ on $\pfgOh$, then for $x=\tau_0h+s\in\pffOT$ as in (a), $\psi_T(bx)=\psi_T(\tau_0bh)
\psi_T(bs)=1$, as $bs\in\pfgH$; so $b\in\pffOT$ by Lemma~\ref{pfg:psi}(a). If $\pfgE(\cdot,\beta)\equiv1$ on $\pfgOs$, then
for $x=\tau_0\delta h_1+h_2$, $\psi_T(x\beta)=\pfgE(\delta h_1,\beta)\psi_T(h_2\beta)=1$ (as $\delta h_1\in\pfgOs$ and
$h_2\beta\in\pfgH$); so $\beta\in\pffOT$. The last identity is clear. (c) $\overline{\tau_0b\bar z}=-\bar\tau_0bz$, and
$\psi_T(y)=\psi_T(-\bar y)$ by (a) of Lemma~\ref{pfg:psi}; so $\psi_T(\tau_0b\bar z)=\psi_T(\bar\tau_0bz)$ and
$\pfgE(b,z+\bar z)=\psi_T((\tau_0+\bar\tau_0)bz)=\psi_T(bz)$. (d) The inverse is $x\mapsto(\delta(x+\bar x),x-\tau_0(x+
\bar x))$: indeed $x+\bar x=(\tau_0+\bar\tau_0)\delta^{-1}b_1=\delta^{-1}b_1$ for $x=\tau_0\delta^{-1}b_1+b_2$. Both maps
preserve integrality, so the compact open subgroups correspond and the Haar measures, both normalised to give them
volume $1$, correspond. If $S=\{b:ub\in\pffOT\}$, the image of $S\times S$ lies in $\{x:ux\in\pffOT\}$ (as
$\tau_0\delta^{-1}$ is integral), of volume $|u|_T^{-1}$. Multiplication by $c\in\pfgH$ on $\pfgS\times\pfgS$
corresponds to multiplication by $c$ on $\pffET$, of module $|c|_T$, and acts the same way on both factors; and
$\beta\mapsto\delta\beta$ is a measure preserving isomorphism $\pfgH\to\pfgS$ commuting with $c$.
\end{proof}

\begin{lemma}[Orthogonality]\label{pfg:orth}
Let $\varpi_1$ be as in Lemma~\ref{pfg:psi}(c), $\ell\ge0$, $B_\ell=\varpi_1^{-\ell}\pfgOs$. For $y\in\pfgH$,
\[
\int_{B_\ell}\pfgE(b,y)\,db=\pbvol(B_\ell)\,\mathbf 1[y\in\varpi_1^{\ell}\pfgOh],\qquad
\pbvol(B_\ell)=|\varpi_1|_T^{-\ell/2}=\pbvol(\varpi_1^\ell\pfgOh)^{-1}.
\]
\end{lemma}
\begin{proof}
$b\mapsto\pfgE(b,y)$ is a character of the compact group $B_\ell$; it is trivial iff $\pfgE(b',\varpi_1^{-\ell}y)=1$
for all $b'\in\pfgOs$, i.e.\ (Lemma~\ref{pfg:duality}(b)) iff $\varpi_1^{-\ell}y\in\pfgOh$. The volumes follow from
Lemma~\ref{pfg:duality}(d).
\end{proof}

Let $\pfgz(\mu)=\sum_I[\pffOT:I]^{-\mu}=\prod_{w\in T}(1-q_w^{-\mu})^{-1}$, the sum over the ideals of finite index
(\S\ref{pff:ss-large}); it is finite for every $\mu>0$. Every ideal of finite index is $w\pffOT$ for some
$w\in\pffOT\cap\pffET^\times$, with $[\pffOT:w\pffOT]=|w|_T^{-1}$ ($\pffOT$ is a finite product of discrete valuation
rings).

\emph{Setting for the rest of this subsection.} $V$ is a free $\pffET$-module of rank $\kappa$ with a hermitian form
$h$ and a lattice $L$ (a free $\pffOT$-module spanning $V$) on which $h$ is \emph{unimodular}: in a basis of $L$ the
Gram matrix $\mathsf J$ satisfies $\mathsf J^*=\mathsf J$, $\mathsf J,\mathsf J^{-1}$ integral. Coordinates and the
Haar measure $dx$ of $V$ ($\pbvol(L)=1$) are taken in that basis; $h(x,x)\in\pfgH$.

\begin{lemma}[Gauss sums of residue classes]\label{pfg:gauss}
Let $g$ be a unit of $\pffET$, $x_0\in V$, $\beta\in\pfgH$, and for $b\in\pfgS$
\[
G(b)=\int_{x_0+gL}\pfgE\bigl(b,h(x,x)-\beta\bigr)\,dx .
\]
Then $|G(b)|^2\le|g|_T^{2\kappa}[\pffOT:I_b]^{-\kappa}$, where $I_b=\{c\in\pffOT:bg\bar gc\in\pffOT\}$; and for every
$\mu>0$ with $2\mu+1\le\kappa$,
\[
\int_{\pfgS}|G(b)|\,db\le\pfgz(\mu)\,|g|_T^{\kappa-1}.
\]
\end{lemma}
\begin{proof}
\emph{Weyl differencing.} Writing the second point of $C=x_0+gL$ as $x+d$, $d\in gL$,
\[
|G(b)|^2=\int_C\int_{gL}\overline{\pfgE(b,h(x,x)-\beta)}\,\pfgE(b,h(x+d,x+d)-\beta)\,dd\,dx=\int_{gL}\pfgE(b,h(d,d))
\int_C\psi_T\bigl(b\,h(x,d)\bigr)\,dx\,dd ,
\]
since $h(x+d,x+d)-h(x,x)=h(d,d)+z+\bar z$ with $z=h(x,d)$, and Lemma~\ref{pfg:duality}(c). In the inner integral put
$x=x_0+gm$, $m\in L$: it is $\psi_T(bh(x_0,d))$ times the integral over $gL$ of the character $gm\mapsto\psi_T(b\bar g
h(m,d))$, which is $\pbvol(gL)=|g|_T^\kappa$ if the character is trivial and $0$ otherwise. The values $h(m,d)=m^*\mathsf
Jd$, $m\in L$, form the $\pffOT$-module generated by the coordinates of $\mathsf Jd$; by Lemma~\ref{pfg:psi}(a) the
character is trivial iff $b\bar g\,\mathsf Jd$ is integral, i.e.\ (unimodularity) iff $b\bar gd\in L$. Writing
$d=gm'$, this says $bg\bar gm'\in L$, i.e.\ every coordinate of $m'$ lies in $I_b$. The ideal $I_b$ has finite index
(at each $w\in T$ it is $\auO_w$ or $(bg\bar g)_w^{-1}\auO_w$), so $I_b=w_b\pffOT$ and the set of such $d$ is $gw_bL$,
of volume $|g|_T^\kappa|w_b|_T^\kappa$. Hence $|G(b)|^2\le|g|_T^{2\kappa}|w_b|_T^\kappa$.

\emph{Sum over the ideals.} For an ideal $I=w\pffOT$ of finite index let $S_I=\{b:I_b=I\}$. As $w\in I_b$ for $b\in S_I$,
$S_I\subseteq\{b:bg\bar gw\in\pffOT\}$, and by Lemma~\ref{pfg:duality}(d) $\pbvol(S_I)\le|g\bar gw|_T^{-1/2}=|g|_T^{-1}
[\pffOT:I]^{1/2}$. The sets $S_I$ cover $\pfgS$ and are measurable: $S_I$ is the set $\{b:bg\bar gw\in\pffOT\}$ minus the union of the
countably many sets $\{b:bg\bar gw'\in\pffOT\}$ with $w'\pffOT\supsetneq I$. So
\[
\int|G|\le\sum_I\pbvol(S_I)\,|g|_T^\kappa[\pffOT:I]^{-\kappa/2}\le|g|_T^{\kappa-1}\sum_I[\pffOT:I]^{-(\kappa-1)/2}\le
|g|_T^{\kappa-1}\pfgz(\mu),
\]
as $(\kappa-1)/2\ge\mu$ and every index is $\ge1$.
\end{proof}

\begin{proposition}[Fibre measures]\label{pfg:fibres}
Assume $\kappa\ge3$. For $\phi\colon V\to\C$ locally constant with compact support and $\beta\in\pfgH$ put
$F_\phi(b;\beta)=\int_V\phi(x)\,\pfgE(b,h(x,x)-\beta)\,dx$ and
\[
A_\ell(\phi;\beta)=\pbvol(\varpi_1^\ell\pfgOh)^{-1}\int_{\{x:\,h(x,x)-\beta\in\varpi_1^\ell\pfgOh\}}\phi(x)\,dx\qquad(\ell\ge0).
\]
\begin{enumerate}
\item[(a)] $b\mapsto F_\phi(b;\beta)$ is integrable on $\pfgS$, with $\int|F_\phi(b;\beta)|db=:C_\phi$ independent of
$\beta$, and
\[
\nu_\beta(\phi):=\int_{\pfgS}F_\phi(b;\beta)\,db=\lim_{\ell\to\infty}A_\ell(\phi;\beta),\qquad|A_\ell(\phi;\beta)|\le C_\phi .
\]
\item[(b)] If $\phi\ge0$ then $\nu_\beta(\phi)\ge0$. Hence $\nu_\beta$ is a positive Radon measure on $V$ (the
\emph{fibre measure} at $\beta$), and it is carried by the sphere $S_\beta=\{x:h(x,x)=\beta\}$.
\item[(c)] $\beta\mapsto\nu_\beta(\phi)$ is continuous.
\item[(d)] (\emph{Disintegration.}) $\int_{\pfgH}\nu_\beta(\phi)\,d\beta=\int_V\phi\,dx$; hence $\int_Vf\,dx=\int_{\pfgH}
\int f\,d\nu_\beta\,d\beta$ for every Borel $f\ge0$, $\beta\mapsto\int f\,d\nu_\beta$ being measurable.
\item[(e)] For every unit $g$, every $x_0$ and every $\mu>0$ with $2\mu+1\le\kappa$:
$\nu_\beta(x_0+gL)\le\pfgz(\mu)|g|_T^{\kappa-1}$.
\item[(f)] (\emph{Uniqueness.}) If $(\nu'_\beta)$ is a family of Radon measures, $\nu'_\beta$ carried by $S_\beta$, with
$\int\nu'_\beta(\phi)\,d\beta=\int\phi\,dx$ and $\beta\mapsto\nu'_\beta(\phi)$ continuous for every such $\phi$, then
$\nu'_\beta=\nu_\beta$ for every $\beta$.
\end{enumerate}
\end{proposition}
\begin{proof}
Choose a unit $g$ such that $\phi$ is constant on the classes $x_0+gL$ and vanishes off finitely many of them; then
$\phi$ is a finite combination of indicators of such classes.

(a) For an indicator, $F$ is the Gauss sum $G$ of Lemma~\ref{pfg:gauss}, integrable (take $\mu=1$; $\kappa\ge3$).
As $\pfgE(b,h(x,x)-\beta)=\pfgE(b,h(x,x))\pfgE(b,-\beta)$, $|F_\phi(b;\beta)|=|F_\phi(b;0)|$. By Fubini's theorem on the
compact set $B_\ell$ and Lemma~\ref{pfg:orth}, $\int_{B_\ell}F_\phi(b;\beta)\,db=\pbvol(B_\ell)\int\phi(x)\mathbf 1[h(x,x)-\beta
\in\varpi_1^\ell\pfgOh]\,dx=A_\ell(\phi;\beta)$. So $|A_\ell|\le C_\phi$, and as the $B_\ell$ increase to $\pfgS$ (as
$|\varpi_1|_w<1$ at every $w$), dominated convergence gives the limit.

(b) $A_\ell(\phi;\beta)\ge0$ for $\phi\ge0$. A positive linear functional on the locally constant compactly supported
functions extends uniquely to $C_c(V)$ (these functions are dense in $C_c(V)$ for the uniform norm with supports in a
fixed compact open set, and $|\nu_\beta(\phi)|\le\nu_\beta(\mathbf 1_U)\sup|\phi|$ for $\phi$ supported in $U$), hence is a
Radon measure (Riesz). If $\operatorname{supp}\phi\cap S_\beta=\emptyset$, the compact set $\{h(x,x)-\beta:x\in
\operatorname{supp}\phi\}$ does not contain $0$, so it misses $\varpi_1^\ell\pfgOh$ for large $\ell$, and
$A_\ell(\phi;\beta)=0$; hence $\nu_\beta(V\smallsetminus S_\beta)=0$.

(c) $\nu_\beta(\phi)=\int F_\phi(b;0)\pfgE(b,-\beta)\,db$ with $F_\phi(\cdot;0)$ integrable; dominated convergence.

(d) By Fubini, $\int_{\pfgH}A_\ell(\phi;\beta)\,d\beta=\int\phi(x)\,\pbvol(\varpi_1^\ell\pfgOh)^{-1}\pbvol\{\beta:h(x,x)-\beta\in
\varpi_1^\ell\pfgOh\}\,dx=\int\phi\,dx$. For $\ell\ge0$, $A_\ell(\phi;\beta)=0$ unless $\beta$ lies in the compact set
$\{h(x,x):x\in\operatorname{supp}\phi\}+\pfgOh$, and $|A_\ell|\le C_\phi$; dominated convergence gives the identity. Both $\phi\mapsto\int\phi\,dx$ and
$\phi\mapsto\int\nu_\beta(\phi)d\beta$ are then Radon measures agreeing on a dense class, so they are equal; the
extension to Borel $f\ge0$ and the measurability follow by the monotone class theorem.

(e) For $\phi=\mathbf 1_{x_0+gL}$, $\nu_\beta(\phi)=\int G\le\int|G|$ and Lemma~\ref{pfg:gauss}.

(f) Let $\phi_\ell(x)=\phi(x)\mathbf 1[h(x,x)-\beta\in\varpi_1^\ell\pfgOh]$, again locally constant with compact support.
On $S_{\beta'}$ it equals $\phi\cdot\mathbf 1[\beta'-\beta\in\varpi_1^\ell\pfgOh]$, so, $\nu'_{\beta'}$ being carried by
$S_{\beta'}$, the disintegration gives $A_\ell(\phi;\beta)=\pbvol(\varpi_1^\ell\pfgOh)^{-1}\int_{\beta+\varpi_1^\ell\pfgOh}
\nu'_{\beta'}(\phi)\,d\beta'$, a mean of a continuous function over a shrinking neighbourhood of $\beta$. Hence
$\nu'_\beta(\phi)=\lim_\ell A_\ell(\phi;\beta)=\nu_\beta(\phi)$ by (a).
\end{proof}

\begin{remark}\label{pfg:rem-fibres}
By (f), the fibre measures of any continuous disintegration (for instance those built from a flat chart of the quadric) coincide with these. In particular the measures $\nu_\beta$ of Definition~\ref{pff:lastentry}
(``$dx/d\langle x,x\rangle$'' on $X_{0,T}$) and the measure of the top quadric are those of
Proposition~\ref{pfg:fibres}, and the formula $\nu_\alpha(\Psi)=\int_{\pfgS}\int\Psi(y)\psi_T(\tau_0b(y^*y-\alpha))\,dy\,db$
used at the end of the proof of Lemma~\ref{pff:offT} is Proposition~\ref{pfg:fibres}(a) (the integral over $b$
converges absolutely, so the convention ``$0$ otherwise'' there is never invoked).
\end{remark}

\subsection{The gauges at $T$: the axioms in full}\label{pfg:ss-gauge}

We use the place datum $\mathfrak F_T$ of \S\ref{pff:ss-large} under (L1)--(L3), with the parameter $m$ of
\S\ref{pfr:ss-gauge} taken to be $m'$ (it enters only through (G1), the rank $2m+2$ of the blocks of (G4) and the
exponent $s$), and $X'$ as in
\S\ref{pfr:ss-gauge}: the span of the orthonormal vectors of (L1) left after two blocks of $2m'+2$ of them are set
aside; $d':=\operatorname{rank}X'$ and the volume exponent is $s=d'-(3n+1+m')$. Coordinates are taken in the basis of
(L1), which is a basis of $L_T$ with unimodular Gram matrix. For $u\in\pffET^{X'}$ (a vector of coordinates on the
orthonormal basis $x'_1,\dots,x'_{d'}$ of $X'$) and $\omega\in X'$ put $\langle u,\omega\rangle=\sum_i\bar u_i\omega_i$,
$\omega_i$ the coordinates of $\omega$.

\begin{lemma}[Products of close pairs are integral]\label{pfg:Dint}
For close pairs $(\omega_1,\omega_2)$, $(\omega'_1,\omega'_2)$ of $\mathfrak F_T$, the number
$M^{-1}(h(\omega_1,\omega'_1)-h(\omega_2,\omega'_2))$ is integral.
\end{lemma}
\begin{proof}
Write $\omega_2=\omega_1-Mb$, $\omega'_2=\omega'_1-Md$ with $b,d\in L_T$. As $\bar M=M$,
$M^{-1}(h(\omega_1,\omega'_1)-h(\omega_2,\omega'_2))=h(\omega_1,d)+h(b,\omega'_1)-M\,h(b,d)$, and $h$ is integral on
$L_T$ (unimodular), $M\in\pffOT$.
\end{proof}

\begin{lemma}[Common vectors]\label{pfg:common}
Let $g\in\pffK$, $x=y^{-1}uyk$ a position ($y\in\pffG$, $u\in U_T$, $k\in\pffKM$), and $v$ a $\varrho\times d'$ matrix
(the coefficients of $\varrho$ linear forms on $X'$). Put $d=d'-3\varrho$. There is $V\in\auGL_{d'}(\pffOT)$ such that for
every $t\in\pffOT^{d}$ the vector $\omega=\sum_i\tau_ix'_i$, $\tau=V\binom0t$, lies in the set $\mathcal W$ of (G5):
$v\tau=0$ and $(g^{-1}\omega,(gx)^{-1}\omega)\in\pfrCl$.
\end{lemma}
\begin{proof}
Let $B$ be the $r\times d'$ matrix whose columns are the coordinates of the $x'_i$ (integral, by (L1)), $N=yg^{-1}B$,
and $C$ the $3\varrho\times d'$ matrix whose rows are: the $\varrho$ rows of $v$; the rows of $N$ of index $e^*_a$
($a<\varrho$); the rows of $(u^{-1}-1)N$ of index $e_a$ ($a<\varrho$). \emph{Column reduction.} At each $w\in T$ (the
components of $\pffET$ are the fields $E_w$), after multiplying $C_w$ by a power of a uniformiser to make it integral,
the Smith normal form over the principal ideal domain $\auO_w$ gives $C_w=P_wD_wQ_w$ with $P_w\in\auGL_{3\varrho}(\auO_w)$,
$Q_w\in\auGL_{d'}(\auO_w)$ and $D_w$ diagonal of size $3\varrho\times d'$, whose columns of index $>3\varrho$ vanish. Let
$V=(Q_w^{-1})_w\in\auGL_{d'}(\pffOT)$; then the last $d$ columns of $CV$ vanish. \emph{Verification.} For
$\tau=V\binom0t$ we have $C\tau=0$: so $v\tau=0$; $N\tau$ has zero coordinates on the $e^*_a$, i.e.\ is orthogonal
below $\varrho$; and $(u^{-1}-1)N\tau$ has zero coordinates on the $e_a$. By (U) applied to $u^{-1}\in U_T$ and the
vector $N\tau$, the other coordinates of $(u^{-1}-1)N\tau$ vanish too; so $u^{-1}N\tau=N\tau$. As $N\tau=yg^{-1}\omega$,
\[
(gx)^{-1}\omega=k^{-1}y^{-1}u^{-1}yg^{-1}\omega=k^{-1}y^{-1}u^{-1}N\tau=k^{-1}y^{-1}N\tau=k^{-1}g^{-1}\omega .
\]
Now $g^{-1}\omega=g^{-1}B\tau$ is integral ($g\in\pffK$, $B$, $\tau$ integral), and $M^{-1}(g^{-1}\omega-k^{-1}g^{-1}\omega)
=M^{-1}(1-k^{-1})g^{-1}\omega$ is integral as $k^{-1}\in\pffKM$. So the pair is close.
\end{proof}

\begin{lemma}[Integral of a content index]\label{pfg:content}
Let $d'=a+d$, $V\in\auGL_{d'}(\pffOT)$, and for $u\in\pffET^{d'}$ let $w(u)\in\pffET^d$ be the last $d$ coordinates of
$V^*u$ and $\phi(u)=\pbvol\{c\in\pffET:c\,w(u)_j\in\pffOT\ \forall j\}\in(0,\infty]$. Let $\beta\in\pffOT$ be a unit of
$\pffET$ and $0\le s'<d$. Then
\[
\int_{\{u:\,\beta u\in\pffOT^{d'}\}}\phi(u)^{s'}\,du\le\pfgz(d-s')\,|\beta|_T^{-a-(d-s')} .
\]
\end{lemma}
\begin{proof}
$u\mapsto V^*u$ preserves $du$ and the set $\{\beta u\in\pffOT^{d'}\}$ ($V^*\in\auGL_{d'}(\pffOT)$). So the integral is
$\pbvol(\beta^{-1}\pffOT)^a\allowbreak \int_{\{\beta w\in\pffOT^d\}}\phi_0(w)^{s'}dw=|\beta|_T^{-a}\int\phi_0^{s'}$, with
$\phi_0(w)=\pbvol\{c:cw\in\pffOT^d\}$. For almost every $w$, at every $w'\in T$ some coordinate of $w$ has non-zero
component, and the content $\mathfrak c(w)=\sum_j\pffOT w_j$ is $t\pffOT$ for a unit $t$; then $\phi_0(w)=|t|_T^{-1}$, and if
$\beta w$ is integral, $I(w):=\beta t\pffOT$ is an ideal of finite index with $[\pffOT:I(w)]=|\beta t|_T^{-1}$. Given $I$,
$t$ is determined up to a unit of $\pffOT$, and the $w$ with $I(w)=I$ lie in $(t\pffOT)^d$, of volume $|t|_T^d$. Hence
\[
\int\phi_0^{s'}\le\sum_I|t|_T^{-s'}|t|_T^{d}=\sum_I\bigl(|\beta|_T[\pffOT:I]\bigr)^{-(d-s')}=|\beta|_T^{-(d-s')}\pfgz(d-s').\qedhere
\]
\end{proof}

\begin{lemma}[Gauges at $T$; Lemma~\ref{pff:gauge}]\label{pfg:gauge}
Let $\alpha\in\pffOT$ be a unit of $\pffET$ with $\bar\alpha=\alpha$, $\Gamma_\alpha(i)=\{c:\alpha^ic\in\pffOT\}$ and
$\Lambda_\alpha=2/|\alpha|_T$. Under (L1)--(L3), $\Gamma_\alpha$ with the loss $\Lambda_\alpha$ is a gauge for $\mathfrak F_T$:
it satisfies (G0)--(G5) of \S\ref{pfr:ss-gauge} in their present form, with $\tau_0,\delta$ of Lemma~\ref{pfg:psi}(b).
\end{lemma}
\begin{proof}
$\Gamma_\alpha(i)=\alpha^{-i}\pffOT$ is compact and increasing in $i$ ($\alpha$ integral), and $\Lambda_\alpha\ge2$ as
$|\alpha|_T\le1$.

(G0) $\alpha^{i+i'}cc'=(\alpha^ic)(\alpha^{i'}c')$; for a sum $\sigma$ of elements of $\Gamma_\alpha(i)$, $\alpha^i\sigma$ is
integral, so $\alpha^{i+1}\sigma$ is; $\pffOT$ is stable under $-$ and conjugation and $\bar\alpha=\alpha$; $0,1\in
\Gamma_\alpha(i)$; the Gram matrix of the basis (L1) and its inverse are integral (unimodularity, (L3)), hence in
$\Gamma_\alpha(0)\subseteq\Gamma_\alpha(1)$.

(G1) $m'\le9n$ by (L2); $P=1$; $M\in\pffOT$; $\tau_0\delta^{-1}\in\pffOT$ (Lemma~\ref{pfg:psi}(b)); the basis vectors are
integral (Lemma~\ref{pff:L1}) and the elements of $\pffK$ are integral matrices in that basis. All these lie in
$\Gamma_\alpha(1)$.

(G2) For a close pair, $\omega_1$ and $M^{-1}(\omega_1-\omega_2)$ are integral. If $M^{-1}\in\Gamma_\alpha(1)$ and $i\ge0$,
take $c=\alpha^{i+1}$: $\alpha^{i+2}c^{-1}=\alpha$ is integral, so $c^{-1}\in\Gamma_\alpha(i+2)$; and for $\omega_1,\omega_2$ of
size $i$, $c\omega_1=\alpha(\alpha^i\omega_1)$ is integral and $M^{-1}(c\omega_1-c\omega_2)=(\alpha M^{-1})\alpha^i(\omega_1-\omega_2)$
is integral; so $(c\omega_1,c\omega_2)$ is close.

(G3) For a unit $u$ with $\alpha^iu\in\pffOT$, $|\alpha|_T^i|u|_T\le1$, so $|u|_T\le|\alpha|_T^{-i}\le\Lambda_\alpha^i$; and
$\pbvol\Gamma_\alpha(i)=|\alpha|_T^{-i}\le\Lambda_\alpha^{i+1}$.

(G4) A block of $2m'+2$ of the orthonormal vectors of (L1) spans a free module of rank $\kappa=2m'+2\ge3$ on which
$h$ is the standard form $\sum_i\bar y_iy_i$, with the unimodular lattice $\pffOT^{2m'+2}$, and its Haar measure is the one of \S\ref{pfg:ss-fourier}. By Proposition~\ref{pfg:fibres} its
measure disintegrates along the form into the fibre measures $\nu_\beta$, and by (e) there with $\mu=m'$
($2m'+1\le\kappa$), $g=\alpha^{-i}$, $x_0=0$:
\[
\nu_\beta\bigl(\Gamma_\alpha(i)^{2m'+2}\bigr)\le\pfgz(m')\,|\alpha|_T^{-i(2m'+1)}\le2\Bigl(\frac{\Lambda_\alpha}2\Bigr)^{i(2m'+1)}
\le\Lambda_\alpha^{(2m'+2)(i+1)},
\]
using $\pfgz(m')\le2$ ((L2): $\pfgz(m')^n\le1+2^{-n}$).

(G5) Let $g\in\pffK$, $x$ a position, $v$ as there, and $V$ as in Lemma~\ref{pfg:common}; put $d=d'-3\varrho$,
$w(u)$ and $\phi(u)$ as in Lemma~\ref{pfg:content} (with $a=3\varrho$). $\phi$ is measurable. Axiom (G5) is stated with
the form of convention (Cv1), for which, for $u,\omega\in X'$ with coordinates $u_i,\omega_i$ on the orthonormal basis
$x'_i$, the product is $h_{\rm (Cv1)}(\omega,u)=\sum_i\omega_i\bar u_i=u^*\omega=\langle u,\omega\rangle$ (in the convention of
\S\ref{pfg:ss-conv} this is $h(u,\omega)$; the number is the same), and $h_{\rm (Cv1)}(\omega,cu)=\bar c\langle u,\omega\rangle$. We
write the verification with $\langle u,\omega\rangle=u^*\omega$. \emph{First condition.} By Lemma~\ref{pfg:common}, $\mathcal W$
contains the vectors $\omega$ with coordinates $\tau=V\binom0t$, $t\in\pffOT^d$, and for them
$\langle u,\omega\rangle=u^*V\binom0t=\sum_j\overline{w(u)_j}\,t_j$. These values fill the $\pffOT$-module generated by the
$\overline{w(u)_j}$, which is $\overline{\mathfrak c(w(u))}$; if $\mathfrak c(w(u))=t_0\pffOT$ with $t_0$ a unit, this is $\bar t_0\pffOT$,
of volume $|\bar t_0|_T=|t_0|_T=\phi(u)^{-1}$ ($T$ is stable under conjugation), and otherwise $\phi(u)=\infty$. So
$\pbvol\{h_{\rm (Cv1)}(\omega,u):\omega\in\mathcal W\}^{-1}\le\phi(u)$. \emph{Second condition.} The set $\mathcal D$ of (G5) consists of
the numbers $M^{-1}(h_{\rm (Cv1)}(\omega_1,\omega_1')-h_{\rm (Cv1)}(\omega_2,\omega_2'))$, the conjugates of the numbers of
Lemma~\ref{pfg:Dint}, so $\mathcal D\subseteq\pffOT$ by that lemma ($\pffOT$ is stable under conjugation). If
$h_{\rm (Cv1)}(\omega,cu)\in\mathcal D$ for all $\omega\in\mathcal W$, then $\bar c\,\overline{w(u)_j}=h_{\rm (Cv1)}(\omega_j,cu)\in\pffOT$ for the
vectors $\omega_j$ with $t=e_j$; conjugating, $c\,w(u)_j\in\pffOT$ for all $j$, and $c$ lies in the set whose volume is $\phi(u)$. \emph{The integral.} By
Lemma~\ref{pfg:content} with $\beta=\alpha^i$, $a=3\varrho$, $s'=s$ and $d-s=d'-3\varrho-s=m'+4$:
\[
\int_{\Gamma_\alpha(i)^{X'}}\phi^s\le\pfgz(m'+4)\,|\alpha|_T^{-i(3\varrho+m'+4)}\le2\Bigl(\frac{\Lambda_\alpha}2\Bigr)^{i(3\varrho+m'+4)}
\le\Lambda_\alpha^{99n(i+1)},
\]
as $\pfgz(m'+4)\le\pfgz(m')\le2$ and $3\varrho+m'+4\le3n+9n\le99n$.
\end{proof}

For $\alpha=1$ all the sets $\Gamma_1(i)$ are $\pffOT$ and the loss is $2$, as stated in Lemma~\ref{pff:gauge}.

\subsection{The fibre measures at the last entry, and the step function}\label{pfg:ss-fibre}

\emph{Notation.} $\pbmXp_T=X_{0,T}\perp\pffET z$; its \emph{lattice basis} is $x_i$ ($i\ne k_0$), $v=\bar\pi_N^{-1}x_{k_0}+N^{-1}z$,
$z$ (\S\ref{pff:ss-floors}), with Gram matrix unimodular ($x_i$ orthonormal and orthogonal to $v,z$; $h(v,v)=0$,
$h(v,z)=-1$, $h(z,z)=-N$). Let $dy$ be the Haar measure of $\pbmXp_T$ giving the $\pffOT$-span $L^+$ of the lattice basis
volume $1$, $dx$ that of $X_{0,T}$ giving $\bigoplus\pffOT x_i$ volume $1$, and $\nu^+_\beta$, $\nu^0_\beta$ the fibre
measures of Proposition~\ref{pfg:fibres} for these two unimodular lattices (ranks $k+1$ and $k$, both $\ge3$). As in
Remark~\ref{pfg:rem-fibres}, $\nu^0_\beta$ is the measure $\nu_\beta$ of Definition~\ref{pff:lastentry}, and the measure
$\theta^\sharp_\varrho$ of the top quadric (the image of the top floor measure $\theta_\varrho$ of $\mathfrak F_T$ under
$\sigma\mapsto\sigma^{-1}z$) is normalised to be $\nu^+_{-N}$, with $dy$ as just fixed (the continuous disintegration of $dy$ along $y\mapsto h(y,y)$, at
the value $h(z,z)=-N$). Put $\pfgQ=\{e\in\pbmXp_T:h(e,e)=-N\}$.

\begin{lemma}[Coarea; Lemma~\ref{pff:fibre}(a)]\label{pfg:coarea}
For $\phi$ locally constant with compact support on $\pbmXp_T$ and $\beta\in\pfgH$, put $\phi_t(x)=\phi(x-N^{-1}tz)$ on
$X_{0,T}$. Then $t\mapsto\nu^0_{\beta+t\bar t/N}(\phi_t)$ is continuous with compact support, and
\[
\nu^+_\beta(\phi)=|\pi_N|_T^{-1}\int_{\pffET}\nu^0_{\beta+t\bar t/N}(\phi_t)\,dt .
\]
For $\beta=-N$ this is Lemma~\ref{pff:fibre}(a), with $\beta(t)=t\bar t/N-N$; the case of a continuous $g$ with compact
support follows by uniform approximation by locally constant functions with supports in a fixed compact open set (both
sides are finite measures on it).
\end{lemma}
\begin{proof}
\emph{Measures.} Let $L_0=\bigoplus\pffOT x_i\oplus\pffOT z$ and $dy_0$ the Haar measure with $\pbvol(L_0)=1$. The matrix taking
the basis $(x_1,\dots,x_k,z)$ to the lattice basis has determinant $\bar\pi_N^{-1}$, so $\pbvol_{y_0}(L^+)=|\pi_N|_T^{-1}$ and
$dy=|\pi_N|_T\,dy_0$. The map $\Theta(x,t)=x-N^{-1}tz$ carries $dx\,dt$ to $|N|_T\,dy_0=|\pi_N|_T\,dy$ (the coordinate on $z$
is $-N^{-1}t$, and $|N|_T=|\pi_N|_T^2$). So $\int\phi\,dy=|\pi_N|_T^{-1}\iint\phi_t(x)\,dx\,dt$. Moreover
$h(\Theta(x,t),\Theta(x,t))=h(x,x)-t\bar t/N$, as $h(z,z)=-N$, $\bar N=N$, $x\perp z$.

\emph{Slabs.} With $A_\ell$ as in Proposition~\ref{pfg:fibres}, these two facts give
$A^+_\ell(\phi;\beta)=|\pi_N|_T^{-1}\int A^0_\ell(\phi_t;\beta+t\bar t/N)\,dt$. Choose $a$ with $\phi$ invariant under translation
by $\varpi_1^aL^+$; as $z\in L^+$, $\phi_t$ depends only on $t$ modulo $\varpi_1^aN\pffOT$. As $t=h(z,y)$ for
$y=\Theta(x,t)$, $\phi_t=0$ unless $t$ lies in the compact set $h(z,\operatorname{supp}\phi)$. So only finitely many
functions $\phi_t$ occur, and $|A^0_\ell(\phi_t;\cdot)|\le C:=\max_tC_{\phi_t}$. By Proposition~\ref{pfg:fibres}(a) for both
lattices and dominated convergence in $t$, letting $\ell\to\infty$ gives the formula. Continuity of the integrand: $\phi_t$
is locally constant in $t$, and $\beta'\mapsto\nu^0_{\beta'}(\phi_t)$ is continuous (Proposition~\ref{pfg:fibres}(c)).
\end{proof}

\begin{lemma}[The top orbit]\label{pfg:orbit}
$\pffG^{(\varrho)}=\auU(\pbmXp_T)$ (the elements fixing the $e_a,e^*_a$) acts transitively on $\pfgQ$, and for $e_0\in\pfgQ$ the
map $\sigma\mapsto\sigma^{-1}e_0$ from $\pffG^{(\varrho)}$ to $\pfgQ$ is open.
\end{lemma}
\begin{proof}
\emph{Transitivity.} It suffices to work at one pair $\{w,\bar w\}\subseteq T$. Let $e,e'\in\pfgQ$. As $h(e,e)=-N$ is a unit,
$\pbmXp=\pffET e\perp e^\perp=\pffET e'\perp e'^\perp$. If $w=\bar w$, $E_w$ is a field with non-trivial involution, and by
Lemma~\ref{pbp:tb-witt}(a) (characteristic $0$, $h(e,e)=h(e',e')\ne0$; see also \cite[Chapter~7, Theorem~9.1]{Scharlau}) there is an element of
$\auU(\pbmXp_w)$ taking $e$ to $e'$. If $w\ne\bar w$, the components at $w,\bar w$ are $X_1,X_2$ (vector spaces over
$E_w\cong E^+_v$), $h$ is given by a non-degenerate bilinear pairing $B\colon X_1\times X_2\to E_w$, the unitary group is
$\auGL(X_1)$ acting on $X_2$ contragrediently, and $e=(e_1,e_2)$, $e'=(e'_1,e'_2)$ with $B(e_1,e_2)=B(e'_1,e'_2)=c\ne0$. An
element $\gamma\in\auGL(X_1)$ with $\gamma e_1=e'_1$ and $\gamma(\ker B(\cdot,e_2))=\ker B(\cdot,e'_2)$ exists (both kernels are
complements of the lines), and its contragredient takes $e_2$ to $e'_2$ (both are characterised by their values on
$e'_1$ and on $\ker B(\cdot,e'_2)$).

\emph{Openness.} $\pffG^{(\varrho)}$ is locally compact and second countable, $\pfgQ$ is closed in $\pbmXp_T$, hence locally
compact. Let $U$ be a neighbourhood of $1$ and $W$ a compact neighbourhood of $1$ with $W^{-1}W\subseteq U$. Countably many
translates $g_nW$ cover $\pffG^{(\varrho)}$, so $\pfgQ=\bigcup_ng_nWe_0$ with compact $g_nWe_0$; by Baire's theorem some $g_nWe_0$,
hence $We_0$, has an interior point $we_0$ ($w\in W$), and $e_0$ is interior to $w^{-1}We_0\subseteq Ue_0$. So $\sigma\mapsto
\sigma e_0$ is open at $1$, hence everywhere, and so is $\sigma\mapsto\sigma^{-1}e_0$.
\end{proof}

We use three properties of the Bessel vector $\pffW$ of an intermediate group $\pffZ$ (\S\ref{pff:ss-setting}):
\begin{enumerate}
\item[(W1)] $\pffW(gk)=\pffW(g)$ for $k\in\pffKM$: $\pffch(gka_c^{-1})=\pffch(ga_c^{-1}\cdot a_cka_c^{-1})$ and $a_cka_c^{-1}\in\pffKM$;
\item[(W2)] $\pffW(kg)=\pffW(g)$ for $k\in H_T\cap\pffKM$: $H_T$ normalises $U_T$ and fixes the frame, so $\sigma(kuk^{-1})=\sigma(u)$;
if $ga_c^{-1}=uk'$ ($u\in U_T$, $k'\in\pffKM$) then $kga_c^{-1}=(kuk^{-1})(kk')$ and $\pffch(kga_c^{-1})=\pffch(ga_c^{-1})$, and if
$ga_c^{-1}\notin U_T\pffKM$ the same holds for $kga_c^{-1}$;
\item[(W3)] $\pffW$ is a cut Bessel function with $|\pffW(g)_c|\le1$.
\end{enumerate}
Recall $\mathcal O_p(l)=\pfrkap_\varrho(p,1)\int_{\pffG}\pffW(y)\otimes\overline{\pffW(lw_p^{-1}y)}\,dy$ for $l\in\pffG^{(\varrho)}$
(as $(w_pl^{-1})^{-1}=lw_p^{-1}$); it involves only the entries $p_a$, $a<\varrho$.

\begin{lemma}[A gauge for heads near $p$]\label{pfg:headgauge}
Let $p$ be a sequence of units and $\mathcal N'_p=\{r:r_bp_b^{-1}\in\pffOT^\times\ (b<\varrho)\}$, an open neighbourhood of $p$.
There is $\alpha$ as in Lemma~\ref{pfg:gauge} such that every $r\in\mathcal N'_p$ and its conjugate $\bar r$ are heads of length
$\varrho$ for $\Gamma_\alpha$; and $\pfrkap_\varrho(r,1)=\pfrkap_\varrho(p,1)$ for $r\in\mathcal N'_p$.
\end{lemma}
\begin{proof}
Take $\alpha=\varpi_1^a$ with $a$ so large that $\varpi_1^a(p_0\cdots p_{l-1})^{-1}\in\pffOT$ for $l\le\varrho$. For $r\in\mathcal N'_p$,
$(r_0\cdots r_{l-1})^{-1}$ is $(p_0\cdots p_{l-1})^{-1}$ times a unit of $\pffOT$, so has size $1$, and so has its conjugate
($\bar\alpha=\alpha$). And $|r_b|_T=|p_b|_T$, so the head densities agree.
\end{proof}

\begin{lemma}[Support and local constancy of $\mathcal O_p$]\label{pfg:Op}
Let $p$ be a sequence of units, $\alpha$ as in Lemma~\ref{pfg:headgauge}, and $S=\{y\in\pffG:y\text{ has size }60n^2\text{ for }
\Gamma_\alpha\}$, a compact set.
\begin{enumerate}
\item[(a)] For $r\in\mathcal N'_p$ and $l\in\pffG^{(\varrho)}$, the integrand of $\mathcal O_r(l)$ vanishes off $S$; so
$\mathcal O_r(l)$ converges absolutely and $|\mathcal O_r(l)|\le\pfrkap_\varrho(p,1)\,\#C\,\pbvol(S)$.
\item[(b)] If $\mathcal O_p(l)\ne0$ ($l\in\pffG^{(\varrho)}$), then $l\in C_p:=\{y'y^{-1}w_p:y,y'\in S\}$, a compact set.
\item[(c)] There is a neighbourhood $\mathcal N_p$ of $(p,1)$ in $\mathcal N'_p\times\pffG^{(\varrho)}$ such that
$\mathcal O_r(ls)=\mathcal O_p(l)$ for all $(r,s)\in\mathcal N_p$ and all $l\in\pffG^{(\varrho)}$.
\end{enumerate}
\end{lemma}
\begin{proof}
(a) If $\pffW(y)\ne0$ and $\pffW(lw_r^{-1}y)\ne0$, then $y\in S$ by Corollary~\ref{pfg:sizehead} ($r$ a head for $\Gamma_\alpha$,
$l\in\pffG^{(\varrho)}$). The bound follows from (W3). (b) Some $y$ has $\pffW(y)\ne0$ and $\pffW(y')\ne0$ for $y'=lw_p^{-1}y$; so
$y\in S$. As $l$ commutes with $w_p$ and $w_p^{-1}=w_{\bar p}$ (\S\ref{pfr:ss-gainsproof}), $l^{-1}w_{\bar p}^{-1}y'=
l^{-1}w_plw_p^{-1}y=y$; Corollary~\ref{pfg:sizehead} for the head $\bar p$ and $l^{-1}\in\pffG^{(\varrho)}$ gives $y'\in S$. And
$l=y'y^{-1}w_p$. (c) The map $(r,s,y)\mapsto y^{-1}w_psw_r^{-1}y$ is continuous and equals $1$ at $(p,1,y)$; as $\pffKM$ is open
and $S$ compact, there is a neighbourhood $\mathcal N_p\subseteq\mathcal N'_p\times\pffG^{(\varrho)}$ of $(p,1)$ with
$y^{-1}w_psw_r^{-1}y\in\pffKM$ for $(r,s)\in\mathcal N_p$, $y\in S$. Let $(r,s)\in\mathcal N_p$, $l\in\pffG^{(\varrho)}$, and
$y$ with $\pffW(y)\ne0$. If $\pffW(lsw_r^{-1}y)\ne0$ or $\pffW(lw_p^{-1}y)\ne0$, then $y\in S$ (by (a) for $r$ and $ls$, or for $p$
and $l$), and $lsw_r^{-1}y=lw_p^{-1}y\cdot(y^{-1}w_psw_r^{-1}y)$, so the two values agree by (W1). Hence the integrands of
$\mathcal O_r(ls)$ and $\mathcal O_p(l)$ coincide; and $\pfrkap_\varrho(r,1)=\pfrkap_\varrho(p,1)$.
\end{proof}

\begin{proposition}[The fibre pairing; Lemma~\ref{pff:fibre}(b)]\label{pfg:fibrepairing}
Assume (L1)--(L3). Let $p,q$ be heads. Then $\Pi_{p,q}$ is defined on all of $\pfgQ$, depends only on the entries of index
$<\varrho$ of $p,q$, and:
\begin{enumerate}
\item[(a)] the integral defining $\Pi_{p,q}(e)$ converges absolutely, and $\Pi_{p,q}(ke)=\Pi_{p,q}(e)$ for $k\in H_T\cap\pffKM$;
\item[(b)] $\Pi_{p,q}$ vanishes outside a compact subset of $\pfgQ$, and is locally constant on $\pfgQ$;
\item[(c)] there is a neighbourhood $\mathcal N$ of $(p,q)$ such that $\Pi_{p',q'}=\Pi_{p,q}$ on $\pfgQ$ for $(p',q')\in\mathcal N$;
\item[(d)] there is a function $\Phi_{p,q}$ on $\pbmXp_T$, locally constant with compact support and invariant under $H_T\cap
\pffKM$, with $\Phi_{p,q}=\Pi_{p,q}$ on $\pfgQ$.
\end{enumerate}
\end{proposition}
\begin{proof}
By Lemma~\ref{pfg:orbit}, $\pfgQ$ is the orbit of $z$, so $\Pi_{p,q}$ is defined everywhere on it. Fix $e=l^{-1}z$,
$l\in\pffG^{(\varrho)}$.

(a) By Lemma~\ref{pfg:Op}(a),(b), $k\mapsto\langle\mathcal O_p(kl),\mathcal O_q(kl)\rangle$ is bounded and vanishes unless
$k\in H_T\cap C_pl^{-1}$, a compact set, of finite $\mu_H$-measure. For $k\in H_T\cap\pffKM$, $ke=(lk^{-1})^{-1}z$, and
$\mathcal O_p(glk^{-1})=\mathcal O_p(gl)$: substitute $y\mapsto ky$ (Haar measure is invariant), use $\pffW(ky)=\pffW(y)$ (W2) and
$k^{-1}w_p^{-1}k=w_p^{-1}$ ($k\in\pffG^{(\varrho)}$ commutes with $w_p$).

(b) If $\Pi_{p,q}(e)\ne0$, some $k\in H_T$ has $\mathcal O_p(kl)\ne0$, so $kl\in C_p$ and $e=(kl)^{-1}z$ (as $k$ fixes $z$) lies in
the compact set $\{c^{-1}z:c\in C_p\cap\pffG^{(\varrho)}\}$. Local constancy at $e_0=l_0^{-1}z$: by Lemma~\ref{pfg:Op}(c) there
is a neighbourhood $U$ of $1$ in $\pffG^{(\varrho)}$ with $(p,\sigma)\in\mathcal N_p$ and $(q,\sigma)\in\mathcal N_q$ for $\sigma\in U$. By
Lemma~\ref{pfg:orbit}, every $e\in\pfgQ$ near $e_0$ is $\sigma^{-1}e_0=(l_0\sigma)^{-1}z$ with $\sigma\in U$; then
$\mathcal O_p(kl_0\sigma)=\mathcal O_p(kl_0)$ and likewise for $q$, for all $k\in H_T$, so $\Pi_{p,q}(e)=\Pi_{p,q}(e_0)$.

(c) Let $\mathcal N=\{(p',q'):(p',1)\in\mathcal N_p,\ (q',1)\in\mathcal N_q\}$. By Lemma~\ref{pfg:Op}(c) with $s=1$,
$\mathcal O_{p'}=\mathcal O_p$ and $\mathcal O_{q'}=\mathcal O_q$ on $\pffG^{(\varrho)}$.

(d) Let $C\subseteq\pfgQ$ be the compact set of (b) and $L^\sharp$ the lattice spanned by a basis of $\pbmXp_T$ in which
$H_T\cap\pffKM$ acts by integral matrices (the lattice basis; $\pffKM\subseteq\pffK$ is integral, and $H_T$ preserves $\pbmXp_T$).
\emph{A uniform modulus.} For each $y\in C$ choose an integral unit $u_y$ with $\Pi(x)=\Pi(y)$ for $x\in\pfgQ\cap(y+u_yL^\sharp)$
(by (b)); cover $C$ by finitely many $y_i+\varpi_1u_{y_i}L^\sharp$, and put $u=\varpi_1\prod_iu_{y_i}$. If $y\in C$, $x\in\pfgQ$ and
$x-y\in uL^\sharp$, then for $i$ with $y\in y_i+\varpi_1u_{y_i}L^\sharp$ both $x,y$ lie in $y_i+u_{y_i}L^\sharp$, so
$\Pi(x)=\Pi(y_i)=\Pi(y)$. \emph{The extension.} Put $\Phi(x)=\Pi(y)$ if there is $y\in\pfgQ$ with $\Pi(y)\ne0$ and
$x-y\in uL^\sharp$, and $\Phi(x)=0$ otherwise. It is well defined (two such $y,y'$ lie in $C$ and $y'-y\in uL^\sharp$, so
$\Pi(y')=\Pi(y)$), constant on the classes $x+uL^\sharp$, and supported in the compact set $C+uL^\sharp$. On $\pfgQ$:
if $\Pi(x)\ne0$ take $y=x$; if $\Pi(x)=0$ and some $y$ as above existed, then $\Pi(x)=\Pi(y)\ne0$, a contradiction; so
$\Phi=\Pi$ on $\pfgQ$. For $k\in H_T\cap\pffKM$: $k$ preserves $\pfgQ$, $uL^\sharp$ and the values of $\Pi$ by (a), so it maps the
set of admissible $y$ for $x$ onto that for $kx$, and $\Phi(kx)=\Phi(x)$.
\end{proof}

\begin{remark}\label{pfg:rem-top}
Proposition~\ref{pfg:fibrepairing}(c),(d) is what the proof of Proposition~\ref{pff:top}(c) uses (for $(p',q')$ near $(p,q)$,
$\mathcal D(p',q')=|\pi_N|_T^{-1}\int\Phi_{p,q}(x-N^{-1}p'_\varrho z)\,d\nu_{\beta(p'_\varrho)}(x)$, continuous in $p'_\varrho$ by
Lemma~\ref{pfg:coarea}); (d) and Lemma~\ref{pfg:coarea} are what the proofs of Proposition~\ref{pff:top}(a) and of
Theorem~\ref{pff:S} use (there $\Psi\ge0$ is the positive part of $\Rel\sum R_{pq}\Phi_{\tilde p,\tilde q}$, again locally
constant, compactly supported and $H_T\cap\pffKM$-invariant).
\end{remark}

\subsection{The place axiom (P) at $T$, and conventions of the shared file}\label{pfg:ss-P}

\begin{lemma}[Axiom (P) at $T$]\label{pfg:P}
The place $(\pffET,dx,\psi_T)$ (with modulus $|\cdot|_T$) of Theorem~\ref{pff:F} satisfies axiom (P) of
\S\ref{pfs:ss-planch} with $P=1$. (The map $w\mapsto M^{-1}w^{-1}$ of the place $\mathfrak P_T$ plays no role in (P).)
\end{lemma}
\begin{proof}
Take $G=\mathbf 1_{\pffOT}$: continuous (as $\pffOT$ is open and closed), integrable, $G(0)=1$. Then
$\widehat G(z)=\int_{\pffOT}\psi_T(\bar vz)\,dv$ is the integral of a character of the compact group $\pffOT$
($\bar v$ runs over $\pffOT$), equal to $1$ if $\psi_T(z\pffOT)=1$ and $0$ otherwise; by Lemma~\ref{pfg:psi}(a),
$\widehat G=\mathbf 1_{\pffOT}$, so $\int|\widehat G|=1$. The module $|c|_T=\prod_{w\in T}|c_w|_w$ is multiplicative,
$|\bar c|_T=|c|_T$ because conjugation permutes the places of the conjugation-stable set $T$ and maps $E_w$
isometrically onto $E_{\bar w}$, and $d(cv)=|c|_Tdv$ by definition of the module. A non-unit of $\pffET$ has some
component $0$, and $\{x:x_w=0\}$ is null for each of the finitely many $w$.
\end{proof}

\begin{remark}[Conventions of \S\ref{pfs:ss-conv}]\label{pfg:rem-conv}
(1) The shared proofs take $h$ linear in its first variable (convention (Cv1) of \S\ref{pb:ss-setting}), so that the coordinate on $e^*_a$ is $h(\omega,e_a)$;
\S\ref{pfg:ss-conv} wrote $h(u,v)=u^*\mathsf Jv$, linear in the second variable. The two readings differ by
replacing $h(u,v)$ by $h(v,u)=\overline{h(u,v)}$. Every statement of this section is about sizes, integrality,
orthogonality, volumes, or characters evaluated on traces, and each of these is unchanged under conjugation of the
values of $h$ (the sets $\Gamma(i)$, $\pffOT$, $\pfgOh$, $\pfgOs$ are stable under conjugation, conjugation preserves
$dx$, and $\psi_T(\bar y)=\psi_T(-y)$); in the proofs, $h(e_a,\omega)$ becomes $h(\omega,e_a)$, $h(gm,d)=\bar g\,h(m,d)$
becomes $g\,h(m,d)$, with the same conclusions; the verification of (G5) in Lemma~\ref{pfg:gauge} is written directly
with the form $h_{\rm (Cv1)}$ in which (G5) is stated. In particular
Lemma~\ref{pff:gauge}, in the form Lemma~\ref{pfg:gauge}, holds verbatim under the convention of \S\ref{pfs:ss-conv}.
(2) Fact (C4) is used in \S\ref{pff:sec} and here only through Theorems~\ref{pfr:deep}, \ref{pfr:shells}, i.e.\ in the
indexing recalled in Remark~\ref{pfs:rem-corr}; nothing in Lemma~\ref{pff:gauge} or in this
section depends on the indexing. (3) ``Near at $\varrho+2$'' with $\pfrR^{(\varrho+2)}=\pfrU$: in the proof of
Lemma~\ref{pff:floormass}(b) an element $g\in\pffK$ is near at $\varrho+2$ (take $\xi=1\in\pfrU$), which is all that is
used; and $\pfrU\subseteq\pfrR^{(j)}$ for $j\le\varrho$ is property (U) of \S\ref{pfg:ss-conv}. So the axioms verified in
Lemma~\ref{pfg:gauge} and the use of nearness in \S\ref{pff:sec} are consistent with (C4) and with the definition of nearness in \S\ref{pfr:sec} (Remark~\ref{pfs:rem-corr}).
\end{remark}

\subsection{Siegel's formula: the published theorem and the normalisation}\label{pfg:ss-siegel}

\emph{Data.} $E^+$ is totally real of degree $f$, $E$ is CM, unramified over $E^+$ at every finite place, $\psi_E=
\psi_\Q\circ\mathrm{Tr}_{E/\Q}$ the standard character of $\auA_E/E$ (its components at the infinite places are
$x\mapsto e^{\mp2\pi i(x+\bar x)}$), and $\psi_{E^+}(y):=\psi_E(\tfrac12y)=\psi_\Q(\mathrm{Tr}_{E^+/\Q}y)$ for
$y\in\auA_{E^+}$, the standard character of $\auA_{E^+}/E^+$. $V=E^k$ with $h(x,y)=x^*y$, $H=\auU(V)$, and $\mathcal C$,
$h_c$, $J_X$ as in \S\ref{pff:ss-siegel}. $G'=\auU(1,1)$ is the unitary group of the split skew-hermitian plane over $E$,
$P=MN$ its Borel subgroup, $n(b)\in N(\auA)$ ($b\in\auA_{E^+}$) and $w$ the Weyl element.

\emph{The Weil representation.} Fix a splitting character $\chi$ of $\auA_E^\times/E^\times$ (with
$\chi|_{\auA_{E^+}^\times}=\varepsilon_{E/E^+}^k$). The Weil representation $\omega=\omega_{\psi,\chi}$ of
$G'(\auA)\times H(\auA)$ on the Schwartz--Bruhat space $\mathcal S(V(\auA_E))$ (Schr\"odinger model) satisfies
\[
\omega(1,h)\Phi(x)=\Phi(h^{-1}x),\quad\omega(n(b))\Phi(x)=\psi_{E^+}(b\,h(x,x))\Phi(x),\quad
\omega(w)\Phi(0)=\gamma\int_{V(\auA_E)}\Phi(x)\,dx ,
\]
where $dx$ is the self-dual measure for the pairing $\psi_E(h(x,y))$, i.e.\ the Haar measure for which $V(E)\backslash
V(\auA_E)$ has volume $1$, and $\gamma=\prod_v\gamma_v=1$ is the product of the local Weil indices; and
$\omega(m(a))\Phi(0)=\chi(a)|a|^{k/2}\Phi(0)$, which is $\Phi(0)$ for $a\in E^\times$. Put $\theta(g,h;\Phi)=\sum_{\xi\in V(E)}
\omega(g,h)\Phi(\xi)$ and $E(g,\Phi)=\sum_{\gamma\in P(E^+)\backslash G'(E^+)}\omega(\gamma g)\Phi(0)$.

\begin{theorem}[Siegel--Weil formula, convergent anisotropic case; Weil]\label{pfg:weil}
Let $V$ be anisotropic of dimension $k\ge3$. Then $\sum_\gamma|\omega(\gamma g)\Phi(0)|$ converges and is locally bounded in $g$,
and for every $\Phi\in\mathcal S(V(\auA_E))$ and $g\in G'(\auA)$
\[
\int_{H(E^+)\backslash H(\auA)}\theta(g,h;\Phi)\,dh=E(g,\Phi),
\]
$dh$ the Haar measure of total mass $1$ on the compact quotient.
\end{theorem}
This is Weil's theorem \cite[Th\'eor\`eme~5]{WeilSiegel} for the dual pair $(\auU(1,1),\auU(V))$; the condition $k>2$ is
the range of absolute convergence of the Eisenstein series on $\auU(1,1)$ at the point given by $V$ ($k>2n$ for
$\auU(n,n)$), in which Weil's hypotheses hold for anisotropic $V$ (condition (B) of \cite[no.~38]{WeilSiegel} in the hermitian case
$\varepsilon=\frac12$ is $k>2n$ with $n=1$, and the theta integral converges since $[\auU(V)]$ is compact, \cite[Proposition~8]{WeilSiegel},
$r=0$); the statement in this notation for unitary groups is recalled in \cite[Introduction, p.~242]{IchinoSW}. We use
it as stated.

\begin{lemma}[The measures are Tamagawa measures]\label{pfg:tamagawa}
\begin{enumerate}
\item[(a)] The product of Lebesgue measure $dx\,dy$ on each $E_\tau=\C$ and the Haar measure on $\auA_{E,f}$ giving
$\widehat\auO_E$ volume $\pbmvol=(2/\aurd)^f$ gives $E\backslash\auA_E$ volume $1$.
\item[(b)] The product of $|\operatorname{disc}E|^{1/4}$, Lebesgue measure on each $E^+_\tau=\R$ and the Haar measure $dz$
on $\mathrm{Herm}_f=\auA_{E^+,f}$ of \S\ref{pff:ss-siegel} (codifferent of volume $1$) gives $E^+\backslash\auA_{E^+}$
volume $1$.
\end{enumerate}
\end{lemma}
\begin{proof}
(a) $\auA_E=E+(E_\infty\times\widehat\auO_E)$ with $E\cap(E_\infty\times\widehat\auO_E)=\auO_E$, so the volume is
$\operatorname{covol}(\auO_E\subset\C^f)\cdot\pbmvol=2^{-f}|\operatorname{disc}E|^{1/2}\cdot2^f\aurd^{-f}=1$, as
$\aurd^{f}=|\operatorname{disc}E|^{1/2}$. (b) For $z\in\auA_{E^+,f}$, $\psi_{E,f}(z\widehat\auO_E)=\psi_{E^+,f}(z\,
\mathrm{Tr}_{E/E^+}\widehat\auO_E)$ and $\mathrm{Tr}_{E/E^+}\widehat\auO_E=\widehat\auO_{E^+}$ (unramified), so the
self-adjoint codifferent is the inverse different $\widehat{\mathfrak d}^{-1}_{E^+}$, and $dz(\widehat\auO_{E^+})=
N(\mathfrak d_{E^+})^{-1}=|\operatorname{disc}E^+|^{-1}$. The volume is $|\operatorname{disc}E|^{1/4}\cdot
|\operatorname{disc}E^+|^{1/2}\cdot|\operatorname{disc}E^+|^{-1}=1$, since $|\operatorname{disc}E|=|\operatorname{disc}E^+|^2$
($E/E^+$ unramified at all finite places).
\end{proof}

\begin{proposition}[Theorem~\ref{pff:siegelformula}]\label{pfg:siegel}
Theorem~\ref{pff:siegelformula} holds as stated, with the measures and the normalisation of \S\ref{pff:ss-siegel}; moreover
the $z$-integral defining $\delta(\Phi,\alpha)$ converges absolutely.
\end{proposition}
\begin{proof}
$V=E^k$ with $h(x,y)=x^*y$ is anisotropic: for $0\ne x\in E^k$ and every embedding $\tau$,
$\tau(x^*x)=\sum_i|\tau(x_i)|^2>0$; so Theorem~\ref{pfg:weil} applies.
Let $\Phi_f$ be the function $\Phi$ of Theorem~\ref{pff:siegelformula}, $\Phi_\infty(x)=\exp(-\sum_\tau x_\tau^*x_\tau)$, and
$\Phi=\Phi_\infty\otimes\Phi_f$. Measures: $dx$ as in Lemma~\ref{pfg:tamagawa}(a) on each coordinate (the self-dual
measure: covolume $1$), and $db$ the measure of Lemma~\ref{pfg:tamagawa}(b) on $\auA_{E^+}$.

\emph{Coefficient of the theta side.} For $h\in H(\auA)$, $\theta(n(b),h;\Phi)=\sum_\xi\psi_{E^+}(b\,\xi^*\xi)\Phi(h^{-1}\xi)$,
absolutely and uniformly convergent on compacta. Since $E^+\backslash\auA_{E^+}$ has volume $1$ and $\psi_{E^+}$ is
non-trivial and trivial on $E^+$, $\int_{E^+\backslash\auA_{E^+}}\theta(n(b),h;\Phi)\psi_{E^+}(-\alpha b)\,db=\sum_{\xi^*\xi=\alpha}
\Phi(h^{-1}\xi)$. For such $\xi$, $\Phi_\infty(h_\infty^{-1}\xi)=e^{-\sum_\tau\tau(\alpha)}$. Integrating over the compact
$H(E^+)\backslash H(\auA)$ (Fubini), the $\alpha$-coefficient of the left side of Theorem~\ref{pfg:weil} at $g=1$ is
$e^{-\sum\tau(\alpha)}\int\Sigma(h_f)\,dh$ with $\Sigma(h_f)=\sum_{\xi^*\xi=\alpha}\Phi_f(h_f^{-1}\xi)$ (a finite sum: $\Phi_f$ has compact
support and the $\xi$ lie in a lattice and on a compact sphere at infinity). $H(\auA)=\bigsqcup_cH(E^+)h_cH(\R)J_X$, the image of
each piece in the quotient is $(H(\R)\times J_X)$ divided by $H(E^+)\cap h_cJ_Xh_c^{-1}=1$ (Lemma~\ref{au:free}), so each class
has mass $1/\#\mathcal C$; $\Sigma$ is right $J_X$-invariant. Hence the coefficient is
$e^{-\sum\tau(\alpha)}\#\mathcal C^{-1}\sum_c\Sigma(h_c)$.

\emph{Coefficient of the Eisenstein side.} $P(E^+)\backslash G'(E^+)=\{1\}\sqcup\{wn(\beta):\beta\in E^+\}$ (Bruhat). The
identity term contributes $\Phi(0)\int\psi_{E^+}(-\alpha b)db=0$. By the absolute convergence of Theorem~\ref{pfg:weil} and
Tonelli, $\int_{E^+\backslash\auA}\sum_\beta|\omega(wn(\beta+b))\Phi(0)|\,db=\int_{\auA_{E^+}}|\omega(wn(b))\Phi(0)|\,db<\infty$, so
the coefficient is
\[
\mathcal E_\alpha=\int_{\auA_{E^+}}\mathcal W(b)\,db,\qquad\mathcal W(b)=\int_{V(\auA_E)}\Phi(x)\,\psi_{E^+}\bigl(b(x^*x-\alpha)\bigr)\,dx ,
\]
with $\gamma=1$. The integrand factorises: $\mathcal W(b)=\prod_\tau\mathcal W_\tau(b_\tau)\cdot\mathcal W_f(b_f)$, where
$\mathcal W_f(z)=\int\Phi_f(x)\psi_{E,f}(\tfrac12z(x^*x-\alpha))dx$ and, with $a=\tau(\alpha)>0$ and $t=b_\tau\in\R$,
\[
\mathcal W_\tau(t)=\int_{\C^k}e^{-|x|^2}e^{\mp2\pi it(|x|^2-a)}\,dx=e^{\pm2\pi ita}\Bigl(\frac\pi{1\pm2\pi it}\Bigr)^k .
\]
$\mathcal W_\tau$ is integrable ($k\ge2$) and not identically $0$, and $\int|\mathcal W|<\infty$; so $\mathcal W_f$ is integrable for
$dz$ (this is the absolute convergence of $\delta$), and $\mathcal E_\alpha=|\operatorname{disc}E|^{1/4}\prod_\tau\int_\R
\mathcal W_\tau\cdot\delta(\Phi_f,\alpha)$ by Lemma~\ref{pfg:tamagawa}(b). \emph{The archimedean density.} The image of Lebesgue
measure of $\C^k=\R^{2k}$ under $x\mapsto|x|^2$ is $\sigma(r)dr$ with $\sigma(r)=\pi^kr^{k-1}/(k-1)!$ (the ball of radius
$\sqrt r$ has volume $\pi^kr^k/k!$). So $\mathcal W_\tau(t)=\int_0^\infty F(r)e^{\mp2\pi it(r-a)}dr$ with $F(r)=e^{-r}\sigma(r)$,
continuous at $a$ and integrable, with integrable Fourier transform; by Fourier inversion
$\int_\R\mathcal W_\tau(t)\,dt=F(a)=e^{-\tau(\alpha)}\pi^k\tau(\alpha)^{k-1}/(k-1)!$.

\emph{Conclusion.} Theorem~\ref{pfg:weil} equates the two coefficients:
$e^{-\sum\tau(\alpha)}\#\mathcal C^{-1}\sum_c\Sigma(h_c)=|\operatorname{disc}E|^{1/4}\allowbreak \prod_\tau e^{-\tau(\alpha)}\allowbreak \pi^k\tau(\alpha)^{k-1}
/(k-1)!\cdot\delta(\Phi_f,\alpha)$. Cancelling $e^{-\sum\tau(\alpha)}$ and multiplying by $\#\mathcal C$ gives
Theorem~\ref{pff:siegelformula}.
\end{proof}

\begin{remark}\label{pfg:rem-siegel}
Every constant of Theorem~\ref{pff:siegelformula} is now accounted for: $\#\mathcal C$ from the class decomposition and the
freeness of $J_X$; $|\operatorname{disc}E|^{1/4}$ from Lemma~\ref{pfg:tamagawa}(b); $\pi^k\tau(\alpha)^{k-1}/(k-1)!$ from the
archimedean density; $\pbmvol^k$ enters through the measure $dx$ in $\delta$, fixed by Lemma~\ref{pfg:tamagawa}(a). The check
through Lemma~\ref{pb:massratio} and \cite{GHY} in \S\ref{pff:ss-siegel} is not needed.
\end{remark}


%% ===== proofs/C-exchange.tex =====
% proofs/C-exchange.tex
% The exchange identity (Theorem pc:X of C-determinant.tex): statement in the convention of
% \S pc:ss-data, dictionary of splittings, non-vanishing of the constant, architecture of the proof,
% and its proof. Labels of this file start with pcx:.
% Macros: those of C-determinant.tex (\auU, \auA, \auNm, \pcm, \pcR, \pcMat, \psvac, ...).
\providecommand{\pcxH}{\mathbb H}
\providecommand{\pcxstr}{\mathrm{r}}
\providecommand{\pcxc}{\mathfrak c}
\providecommand{\pctcone}{\operatorname{cone}}

\section{The exchange identity}\label{pcx:sec}

This section states the exchange identity and proves it, together with
\S\S\ref{pcu:sec}, \ref{pct:sec}, \ref{pcn:sec} (see \S\ref{pcx:ss-status}); it is
Theorem~\ref{pc:X}, on which Theorems~\ref{pc:tri} and~\ref{pc:auxcore} and
Proposition~\ref{po:exchange} rest. Notation is that of \S\ref{pc:ss-data}.

\subsection{Statement, splitting and non-vanishing}\label{pcx:ss-statement}

\subsubsection*{Conventions}
For spaces $S,T$ and automorphic characters $\alpha,\beta$ of $\auU(1)$ we use, as in
\S\ref{pc:ss-data},
\[
\omega_{\alpha,\beta}(g,h)=\alpha(\det g)^{-1}\beta(\det h)\,\omega_{\auHom(S,T)}(g,h),\qquad
\theta^{S,T}_{\alpha,\beta}(g,h;f)=\theta\bigl(\omega_{\alpha,\beta}(g,h)f\bigr),
\]
$g\in\auU(S)(\auA)$, $h\in\auU(T)(\auA)$, where $\omega_{\auHom(S,T)}(g,h)$ is the Weil representation
of $\auU(\auHom(S,T))(\auA)$, defined with $\upsilon$, composed with $x\mapsto hxg^{-1}$
(\S\ref{ps:ss-weil}). \emph{The first index is the character of the source group $\auU(S)$, entering
through its inverse; the second is the character of the target group $\auU(T)$.} We write
$\upsilon$ also for $\upsilon|_{\auU(1)}$.

Let $\pcm$ be a modulus and $S=S_\pcm$ the set of primes $w$ of $E$ dividing $\pcm$. A
\emph{tame character for $\pcm$ of weight $\mu$} is an automorphic character $\nu$ of $\auU(1)$ of
weight $\mu=(\mu_\tau)$ with all $\mu_\tau\ge1$, unramified outside $\pcm$, tamely ramified with
non-trivial tame part at every $w\in S$. Put $F_S=\prod_{w\in S}E_w$, $\auO_S=\prod_{w\in
S}\auO_{E_w}$, and let $\chi_S\colon F_S^\times\to\C^\times$, $t\mapsto\nu(t/\bar t)$ ($t$ read as an
idele supported on $S$), extended by $0$ to the non-invertible elements; on $\auO_S^\times$ it
factors through $\pcR_\pcm^\times$ and gives the character $\chi$ of the residue ring given by the
tame parts. Write $q_w=\#\auF_w$ and, for $e\in\Z_{\ge0}$,
\begin{equation}\label{pcx:eq-Lstar}
L_*(\pcm,\nu,e)=\Bigl(\prod_{w\in S}q_w\Bigr)^{e^2}\prod_{z=1}^e|L^S(z,\nu)|^2,\qquad
L_*(\pcm,\nu,-e)=L_*(\pcm,\nu,e)^{-1},
\end{equation}
where $L^S(s,\nu)$ is the Hecke $L$-function of $x\mapsto\nu(x/\bar x)$ with the Euler factors at
$S\cup\bar S$ removed. For squarefree $\pcm$ (e.g.\ $\pcm_\vartheta$; in general only the primes of $\pcm$ enter) $\prod_{w\in S}q_w=
\auNm\pcm$, and \eqref{pcx:eq-Lstar} is the $L_*$ of \S\ref{pc:ss-data}.

\begin{definition}[Matched functions]\label{pcx:def-matched}
Let $R$ be a commutative ring, $\chi$ a character of $R^\times$, $N=b+a$, the indices of $R^N$
ordered so that the first $b$ come first. For $k\in\auGL_N(R)$ let $k_{[b]}\in\pcMat_{N\times b}(R)$
be its first $b$ columns and $k^{[a]}\in\pcMat_{a\times N}(R)$ the last $a$ rows of $k^{-1}$.
Functions $\Phi$ on $\pcMat_{a\times N}(R)$ and $\Psi$ on $\pcMat_{N\times b}(R)$ are \emph{matched
for $\chi$} if $\Phi$ vanishes off $\{k^{[a]}\}$, $\Psi$ vanishes off $\{k_{[b]}\}$, and
$\Phi(k^{[a]})=\chi(\det k)\Psi(k_{[b]})$ for all $k\in\auGL_N(R)$. Over $F_S$ the same definition is
used with $k\in\auGL_N(\auO_S)$.
\end{definition}

\begin{theorem}[Exchange identity]\label{pcx:X}
Let $W,V,V'$ be spaces of ranks $N=a+b$, $a$, $b$, with $a=b+2e$, $e\ge0$ (so $N$ is even). Let
$q,q_V,q_{V'}$ be their numbers of negative signs; assume $q=q_V+q_{V'}$ and $2q+2\le N$. Let
$\mu=(\mu_\tau)$ be integers with $N\le2\mu_\tau+2q$ for all $\tau$, and let $v$ be a vector of the
Fock space of $\auHom(V',W)\otimes\R$ with finite moments. Then there are $C>0$ and a vector $x$ of
the Fock space of $\auHom(W,V)\otimes\R$ with finite moments, such that: for every modulus
$\pcm\ne(1)$ and every tame character $\nu$ for $\pcm$ of weight $\mu$ there is $c\in\C$ with
\[
|c|^2=C\cdot L_*(\pcm,\nu,e)\ne0,
\]
such that for all automorphic characters $\alpha,\alpha'$ of $\auU(1)$ with
\[
\alpha\alpha'=\gamma:=\nu^{-1}\upsilon^{N/2},
\]
all functions $\Phi$ on $\pcMat_{a\times N}(\pcR_\pcm)$ and $\Psi$ on $\pcMat_{N\times b}(\pcR_\pcm)$
matched for $\chi$, and all $g\in\auU(W)(\auA)$,
\begin{equation}\label{pcx:eq-X}
\int_{[\auU(V)]}\theta^{W,V}_{\alpha,\gamma}(g,h;x,\Phi)\,dh=c\int_{[\auU(V')]}
\theta^{V',W}_{\gamma,\alpha'}(h',g;v,\Psi)\,dh' .
\end{equation}
Here $\theta(\cdot\,;x,\Phi)$ is the kernel of the datum $[x,\Phi]$ of \S\ref{pc:ss-data} (matrices
in the frames), and $C$, $x$ do not depend on $\pcm,\nu,\alpha,\alpha',\Phi,\Psi,g$.
\end{theorem}
Thus on the left $\auU(W)$ acts through $\alpha(\det g)^{-1}$ and $\auU(V)$ through $\gamma(\det h)$;
on the right $\auU(V')$ acts through $\gamma(\det h')^{-1}$ and $\auU(W)$ through $\alpha'(\det g)$.
Both integrals converge absolutely (Corollary~\ref{pcx:conv}).

\subsubsection*{The dictionary with the other splitting}
The other convention in use (\S\ref{po:ss-set}, \S\ref{pga:ss-L}, Remark~\ref{pc:rem-norm}) writes, for characters $a,b$,
\[
\omega'_{a,b}(g,h)=a(\det h)\,b(\det g)^{-1}\,\omega_{\auHom(S,T)}(g,h),\qquad
\theta'^{S,T}_{a,b}=\theta\circ\omega'_{a,b}.
\]
The underlying operator $\omega_{\auHom(S,T)}(g,h)$ is the same in both (the Weil representation
of the unitary group of $\auHom(S,T)$ defined with $\upsilon$, composed with $x\mapsto hxg^{-1}$), so
\begin{equation}\label{pcx:eq-dict}
\omega'_{a,b}=\omega_{b,a},\qquad\theta'^{S,T}_{a,b}=\theta^{S,T}_{b,a}
\end{equation}
\emph{exactly}: the dictionary is the transposition of the two indices, with no factor. In the
primed convention \eqref{pcx:eq-X} reads
\[
\int_{[\auU(V)]}\theta'^{W,V}_{\gamma,\alpha}(g,h;x,\Phi)\,dh=c\int_{[\auU(V')]}
\theta'^{V',W}_{\alpha',\gamma}(h',g;v,\Psi)\,dh',
\]
with the \emph{same} $c$, $C$, $x$. Theorem~\ref{pc:X} is printed in the convention of
\S\ref{pc:ss-data}, with the indices of \eqref{pcx:eq-X}. The representation $\omega_\mu$ of \S\ref{ps:ss-weil} is
$\omega_{\mu\upsilon_{N'},\mu\upsilon_N}=\omega'_{\mu\upsilon_N,\mu\upsilon_{N'}}$.

Changing the splitting characters multiplies kernels by characters of the determinants:
\begin{equation}\label{pcx:eq-twist}
\theta^{S,T}_{\alpha\lambda,\beta\lambda'}(g,h;f)=\lambda(\det g)^{-1}\lambda'(\det h)\,
\theta^{S,T}_{\alpha,\beta}(g,h;f),
\end{equation}
immediately from the definition, since $\theta$ is linear. In particular the passage from $\Theta'$
to $\Theta$ of Remark~\ref{pc:rem-norm} (multiplication by $\chi(\det h)$ on the target group,
$|\chi|=1$) does not change $|c|$: if \eqref{pcx:eq-X} is used with both sides multiplied by the
same unitary character of $\det g$ or integrated against forms twisted by unitary characters, the
constant $c$ is unchanged, and if the twist is put on one side only it is not an instance of
\eqref{pcx:eq-X}.

\subsubsection*{Non-vanishing of the constant}

\begin{lemma}\label{pcx:Lpos}
Let $\nu$ be a tame character for $\pcm$ of weight $\mu\ne0$. Then $L^S(z,\nu)\ne0$ for every
real $z\ge1$; hence $L_*(\pcm,\nu,e)>0$ for every $e\in\Z$.
\end{lemma}
\begin{proof}
The Hecke character $\lambda(x)=\nu(x/\bar x)$ of $\auA_E^\times/E^\times$ is unitary, and it is not
trivial: its component at an infinite place $\tau$ is $x\mapsto(x/\bar x)^{\mu_\tau}$ (up to the
normalisation of the weight), which is non-trivial since $\mu_\tau\ne0$. For real $z>1$ the Euler
product of $L^S(z,\lambda)$ converges absolutely and each factor $(1-\lambda(\varpi_w)q_w^{-z})^{-1}$
is non-zero, so $L^S(z,\nu)\ne0$. For $z=1$, $L(s,\lambda)$ has no zero on $\Rel s=1$ for a
non-trivial unitary Hecke character (Hecke; see e.g.\ Lang, \emph{Algebraic Number Theory}), and removing finitely many Euler factors, each finite and non-zero
at $s=1$, keeps this. Thus every factor of \eqref{pcx:eq-Lstar} is positive.
\end{proof}

\begin{corollary}[The constant is non-zero for every $e$]\label{pcx:cne0}
In Theorem~\ref{pcx:X}, $c\ne0$ for every $e\ge0$; for $e=0$, $|c|^2=C$.
\end{corollary}
\begin{proof}
$C>0$ and $L_*(\pcm,\nu,e)>0$ by Lemma~\ref{pcx:Lpos} (the weight is non-zero because
$\mu_\tau\ge(N-2q)/2\ge1$).
\end{proof}

\begin{remark}[The exponents $e$ used downstream]\label{pcx:rem-e}
Theorem~\ref{pcx:X} is applied: (i) in the commutation of one square of the triangle
(Theorem~\ref{pc:tri}(b), \S\ref{pc:ss-exchange}), with $(W,V,V')=(\pcG_{i,n+1}\oplus\pcG_{i+1,n+1},
\pcG_{i,n+2},\pcG_{i+1,n})$ and $e=t$; (ii) in the adjoint law (Theorem~\ref{pc:tri}(c),
\S\ref{pc:ss-exchange}(v)) with $e=0$; (iii) for $V$ positive definite with explicit determinant
data at the infinite places, which is the form used for Proposition~\ref{po:exchange} and for
Theorem~\ref{pc:auxcore}, with $e=d$ where $a=b+2d$ is dictated by the ranks of the spaces of that
application, an arbitrary integer $d\ge0$ as far as the statement is concerned. Since
Corollary~\ref{pcx:cne0} holds for all $e\ge0$, no restriction on $e$ is needed in any of these
applications; the only input is Lemma~\ref{pcx:Lpos}, valid because the characters used have
non-zero weight. (The non-vanishing can also be read off from $C_1|c|^2=L_*(\pcm,\nu,e)\,C_0$ (proof of Theorem~\ref{pcx:X}, $C_1$ the constant
of Proposition~\ref{pcx:P-matches}) with $C_0\neq0$ the equal-rank constant of Proposition~\ref{pcx:P-equal} and $L_*>0$.)
\end{remark}

\subsection{Architecture of the proof}\label{pcx:ss-arch}

Throughout this subsection $W$ is a space of even rank $N$ with $q$ negative signs, $\nu$ runs
through the tame characters of a fixed weight $\mu$ (for all moduli), $S$ is the set of primes of
the modulus of $\nu$, $\chi_S$ is as in \S\ref{pcx:ss-statement}, and
\[
\gamma=\gamma_\nu=\nu^{-1}\upsilon^{N/2}.
\]
Data at $S$ are functions on matrices over $F_S$: for spaces $S_1,S_2$ of ranks $n_1,n_2$, a Fock
vector $y$ of $\auHom(S_1,S_2)\otimes\R$ and a locally constant compactly supported $\Phi$ on
$\pcMat_{n_2\times n_1}(F_S)$, $[y,\Phi]_S$ denotes the datum which is $y$ at the infinite places,
$\Phi$ at the places of $S$ (in the model (W5) at $w\in S$) and the vacuum $\psvac$ of the standard
lattice elsewhere. $[y,\Phi]_{\bar S}$ is the same with $\Phi$ put at the places of $\bar S$.

\begin{definition}[The two theta integrals]\label{pcx:def-large}
For spaces $W$ (rank $N$) and $T$ (rank $r$), a Fock vector $v$ of $\auHom(W,T)\otimes\R$ and
$\Phi$ on $\pcMat_{r\times N}(F_S)$, $\Psi$ on $\pcMat_{N\times r}(F_S)$, put
\[
\Theta^{\mathrm L}_{W,T}(v,\Phi)(g)=\int_{[\auU(T)]}\theta^{W,T}_{\gamma,\gamma}(g,h;[v,\Phi]_S)\,dh,
\qquad
\Theta^{\mathrm S}_{W,T}(v,\Psi)(g)=\int_{[\auU(T)]}\theta^{W,T}_{\upsilon^{r},\,\nu\upsilon^{N/2}}
(g,h;[v,\Psi^{\mathsf t}]_{\bar S})\,dh,
\]
with $\Psi^{\mathsf t}(y)=\Psi(y^{\mathsf t})$. (In the primed convention, these are the
integrals with characters $(\gamma,\gamma)$, resp.\ $(\nu\upsilon^{N/2},\upsilon^r)$.)
\end{definition}

\begin{definition}[Standard data]\label{pcx:def-std}
Let $N=b+a$. For $z\in\pcMat_{b\times N}(\auO_S)$ define on $\pcMat_{a\times N}(F_S)$ and
$\pcMat_{N\times b}(F_S)$
\[
\Phi^{\mathrm{row}}_z(x)=\chi_S\bigl(\det\tbinom zx\bigr)^{-1},\qquad
\Psi^{\mathrm{col}}_z(y)=\chi_S\bigl(\det(zy)\bigr)^{-1},
\]
where $\Phi^{\mathrm{row}}_z(x)=0$ unless $\binom zx\in\auGL_N(\auO_S)$, and
$\Psi^{\mathrm{col}}_z(y)=0$ unless $y$ is integral and $zy\in\auGL_b(\auO_S)$.
\end{definition}

\begin{definition}[Equal-rank ratio]\label{pcx:def-ratio}
For a space $X$ of even rank $2k$ and a Fock vector $v$ of $\auHom(X,X)\otimes\R$ with finite
moments let
\[
\pcxstr_X(\nu;v)=\frac{\bigl|\Theta^{\mathrm L}_{X,X}(v,\Psi^{\mathrm{col}}_1)(1)\bigr|^2}
{L_*(\pcm,\nu,k)},
\]
where $\Psi^{\mathrm{col}}_1=\chi_S^{-1}\circ\det$ on $\auGL_{2k}(\auO_S)$ and $0$ elsewhere.
\end{definition}

\begin{definition}[Exchange on standard data]\label{pcx:def-RC}
Let $X$, $T$, $T'$ be spaces of ranks $N=b+a$, $a$, $b$, and $\ell$ a function of the tame
characters of weight $\mu$. Vectors $v$ (for $(X,T)$) and $y$ (for $(X,T')$) with finite moments
satisfy $\mathrm{RC}(X,T,T';v,y;\ell)$ if there is a real $C$ such that every $\nu$ has a $c\in\C$
with
\begin{enumerate}
\item[(a)] $C|c|^2=L_*(\pcm,\nu,a-N/2)\,\ell(\nu)$;
\item[(b)] $\Theta^{\mathrm L}_{X,T}(v,\Phi^{\mathrm{row}}_z)=c\,\Theta^{\mathrm S}_{X,T'}
(y,\Psi^{\mathrm{col}}_z)$ for every $z\in\pcMat_{b\times N}(\auO_S)$.
\end{enumerate}
\end{definition}

The proof of Theorem~\ref{pcx:X} uses the following intermediate statements. In each, the
quantifiers over $\nu$ are as above.

\begin{proposition}[Exchange on standard data in even rank]\label{pcx:P-std}
Let $W$ have even rank $N=b+r$, let $U$ (rank $b$) and $V$ (rank $r$) be spaces with
$q_U+q=b+q_V$, and let $W_0=\pcxH\oplus W$, $\pcxH$ a hyperbolic plane. For Fock vectors $z$ of
$(W_0,W_0)$ and $y$ of $(W,U)$ with finite moments there is $v$ for $(W,V)$ with
$\mathrm{RC}(W,V,U;v,y;\nu\mapsto\pcxstr_{W_0}(\nu;z))$.
\end{proposition}

This is the heart of the matter. Its proof has three parts.

\begin{proposition}[Large triple]\label{pcx:P-top}
Let $W_0=\pcxH\oplus W$ be of even rank, $U$ of rank $b$, $X=\pcxH^b\oplus W_0$, and
$T=U\oplus W_0$ (no sign condition is needed; $W_0$ splits off an adapted plane). For $z$ of $(W_0,W_0)$ and
$w$ of $(W_0,U)$ there are $x$ for $(X,T)$ and $x'$ for $(X,U)$, with $x'$ having zero term $w$
along the $b$ planes (for the splitting with the two isotropic frames exchanged, \S\ref{pcu:ss-mixed}),
such that $\mathrm{RC}(X,T,U;x,x';\nu\mapsto\pcxstr_{W_0}(\nu;z))$.
\end{proposition}

\begin{proposition}[Real vectors with prescribed zero term]\label{pcx:P-slice}
Let $W'=W\oplus\pcxH^p$ be an adapted splitting. For every Fock vector $y$ of $(W,U)$ with finite
moments there is a Fock vector $w$ of $(W',U)$ with finite moments whose zero term (restriction of
the real Schr\"odinger model to the slice of matrices vanishing on the isotropic part) is $y$.
\end{proposition}

\begin{proposition}[Removing planes]\label{pcx:P-descent}
Let $W'=W\oplus\pcxH$ be an adapted splitting, $N=\operatorname{rk}W$ even, $T'$ of rank $b$,
$T$ of rank $r'$, and $T_0$ any space of rank $r'-2$ with $q_{T_0}=q_T-1$ (zero terms taken with the
orientation of Remark~\ref{pcu:rem-descent}). If
$\mathrm{RC}(W',T,T';x,x';\ell)$ holds, then there is $v$ with $\mathrm{RC}(W,T_0,T';v,y;\ell)$,
where $y$ is the zero term of $x'$. (Here $y=Z_Ax'$ is the zero term of $x'$ along the given adapted
splitting, $\ell$ is kept and the constant changes. With $d=r'-2-N/2$, $v$ is the cone integral
$\pctcone_{d+1}(x)$ if $d\ge0$ and the dual cone integral $\pctcone^\vee_{-d}(x)$ if $d<0$; the ratio of
the constant terms has modulus $C\,L^\sharp(d+1)$, resp.\ $C\,L^\sharp(-d)$, $L^\sharp(z)=|L^S(z,\nu)|
\prod_{w\in S}q_w^{z-1/2}$; only values of $L^S$ at points $\ge1$ occur and no functional equation is
used. Proof: Proposition~\ref{pct:descent}.)
\end{proposition}

\begin{proposition}[From standard data to matched functions]\label{pcx:P-match}
Let $W$ have even rank $N=b+r$. If $v$, $y$ have finite moments and
$\Theta^{\mathrm L}_{W,V}(v,\Phi^{\mathrm{row}}_z)=c\,\Theta^{\mathrm S}_{W,U}(y,\Psi^{\mathrm{col}}_z)$
for all integral $z$, then for all $\Phi$, $\Psi$ matched for $\chi_S$ over $\auO_S$ with $\Psi$
locally constant, $\Psi$ is a Schwartz--Bruhat function supported on left-invertible matrices and
$\Theta^{\mathrm L}_{W,V}(v,\Phi)=c\,\Theta^{\mathrm S}_{W,U}(y,\Psi)$.
\end{proposition}

\begin{proposition}[Passing to adjoints]\label{pcx:P-flip}
Let $W$ have even rank $N=b+r$, $U$ of rank $b$, $U^-$ the space with the negated form, and $y$ a
Fock vector of $(U,W)$. There is a Fock vector $y^{\mathrm{ad}}$ of $(W,U^-)$, with the same
moments, a number $d_0>0$, and for each $\nu$ a number $d$ with $|d|=d_0$, such that for all $\Phi,\Psi$ matched for
$\chi_S$ with $\Psi$ Schwartz--Bruhat on left-invertible matrices and all $g$,
\[
\gamma(\det g)\,\Theta^{\mathrm S}_{W,U^-}(y^{\mathrm{ad}},\Psi)(g)
=d\int_{[\auU(U)]}\theta^{U,W}_{\gamma,\gamma}(h,g;[y,\Psi]_S)\,dh .
\]
\end{proposition}


\begin{proposition}[Value of the equal-rank ratio]\label{pcx:P-equal}
Let $W$ have rank $2k+2q$ with $q$ negative signs, and $k\le\mu_\tau$ for all $\tau$. There are a
Fock vector $z$ of $(\pcxH\oplus W,\pcxH\oplus W)$ with finite moments and $C_0\ne0$ such that
$\pcxstr_{\pcxH\oplus W}(\nu;z)=C_0$ for every tame character $\nu$ of weight $\mu$.
\end{proposition}
Its proof uses Propositions~\ref{pcx:P-matches} (hence \ref{pcx:P-std}--\ref{pcx:P-flip}) and three
further statements:

\begin{proposition}[Non-vanishing for one character]\label{pcx:P-nonvan}
Let $J$ be the positive definite standard space of rank $2k$, $k\le\mu_\tau$, $Y=\pcxH\oplus J$.
For each $\nu$ there is $z$ with $\pcxstr_Y(\nu;z)\ne0$.
\end{proposition}

\begin{proposition}[Adding a plane]\label{pcx:P-plane}
Let $W'=W\oplus\pcxH$ be an adapted splitting, $\operatorname{rk}W$ even. For every $v$ of $(W,W)$
with finite moments there are $v'$ of $(W',W')$ with finite moments and $C>0$ with
$\pcxstr_{W'}(\nu;v')=C\,\pcxstr_W(\nu;v)$ for all $\nu$.
\end{proposition}

\begin{proposition}[Only the signs matter]\label{pcx:P-signs}
If $Y$, $Y_1$ are spaces of the same even rank and the same signs, and $Y$ contains a hyperbolic
plane split off adaptedly, then $\pcxstr_{Y_1}(\nu;v)=\pcxstr_Y(\nu;v)$ for all $\nu$, $v$.
\end{proposition}

Finally two elementary statements, proved in \S\ref{pcx:ss-elem}:
Lemma~\ref{pcx:std-matched} (standard data are matched) and Lemma~\ref{pcx:mod} (data modulo
$\pcm$ come from matched data on $F_S$).

\begin{proposition}[Exchange for matched functions]\label{pcx:P-matches}
Let $W,U,V$ have ranks $N=b+r$ (even), $b$, $r$, with $q=q_V+q_U$. For $z$ of $(\pcxH\oplus W,
\pcxH\oplus W)$ and $y$ of $(U,W)$ with finite moments there are $C\in\R$ and $v$ of $(W,V)$ with
finite moments such that every $\nu$ has $c$ with
\begin{enumerate}
\item[(a)] $C|c|^2=L_*(\pcm,\nu,r-N/2)\;\pcxstr_{\pcxH\oplus W}(\nu;z)$;
\item[(b)] for all $\Phi,\Psi$ matched for $\chi_S$ over $\auO_S$, $\Psi$ locally constant, and all
$g$:
$\displaystyle\int_{[\auU(V)]}\theta^{W,V}_{1,\gamma}(g,h;[v,\Phi]_S)\,dh=c\int_{[\auU(U)]}
\theta^{U,W}_{\gamma,\gamma}(h,g;[y,\Psi]_S)\,dh$.
\end{enumerate}
\end{proposition}

The logical structure is:
\[
\begin{gathered}\ref{pcx:P-slice}+\ref{pcx:P-top}+\ref{pcx:P-descent}\Rightarrow\ref{pcx:P-std};\\
\ref{pcx:P-std}+\ref{pcx:P-match}+\ref{pcx:P-flip}\Rightarrow\ref{pcx:P-matches};\\
\ref{pcx:P-matches}+\ref{pcx:P-nonvan}+\ref{pcx:P-plane}+\ref{pcx:P-signs}\Rightarrow
\ref{pcx:P-equal};\\
\ref{pcx:P-matches}+\ref{pcx:P-equal}+\ref{pcx:mod}+\eqref{pcx:eq-twist}\Rightarrow
\ref{pcx:X}.\end{gathered}
\]

\begin{lemma}[Adapted planes exist]\label{pcx:adapted}
Let $X$ be a space of rank $n$ with $q_X$ negative signs. For every sign vector of rank $n+2$ with
$q_X+1$ negative signs there are a space $Y$ with these signs and an adapted splitting
$Y=X\oplus\pcxH$ (frames orthonormal at the infinite places and integral, with unimodular sum, at
the finite places).
\end{lemma}
\begin{proof}
Let $\pcxH$ be the plane with basis $e,f$ and Gram matrix $\bigl(\begin{smallmatrix}0&1\\1&0
\end{smallmatrix}\bigr)$; $\auO_Ee\oplus\auO_Ef$ is unimodular. Let $Y=X\oplus\pcxH$ with the
lattice $L_X\oplus\widehat\auO_Ee\oplus\widehat\auO_Ef$, unimodular, and at each infinite place the
orthonormal basis consisting of that of $X$ and of $(e+f)/\sqrt2$, $(e-f)/\sqrt2$ (signs $+1,-1$).
The signature of $Y$ at every infinite place is that of $X$ plus $(1,1)$, so $Y$ has rank $n+2$ and
$q_X+1$ negative signs; ordering the frame so that the signs come in the prescribed order gives
the required sign vector. The splitting $Y=X\oplus\pcxH$, with isotropic frame $e,f$, is adapted
by construction.
\end{proof}

\begin{lemma}[Spaces with the same signs]\label{pcx:landherr}
Let $Y,Y_1$ be spaces of the same rank and signs. There are an $E$-isometry $\phi\colon Y_1(E)\to Y(E)$
and $k\in\auU(Y)(\auA)$ such that $\phi$ carries the frame of $Y_1$ (its lattice at the finite places
and its orthonormal bases at the infinite places) to $k$ applied to the frame of $Y$.
\end{lemma}
\begin{proof}
At each finite place both lattices are unimodular, hence the local hermitian spaces have
determinant a unit, which is a local norm since $E/E^+$ is unramified at finite places; at the
infinite places the signatures agree. By Landherr's theorem (hermitian spaces over $E/E^+$ are
classified by rank, the class of the determinant in $E^{+\times}/\auNm E^\times$, and the
signatures;~\cite{Landherr}) together with the Hasse norm theorem for
the cyclic extension $E/E^+$, the determinants agree in $E^{+\times}/\auNm E^\times$ and $Y_1\cong Y$.
Transport the frame of $Y_1$ by $\phi$: at each finite place it is a unimodular lattice of $Y_v$,
hence of the form $k_vL_Y$ with $k_v\in\auU(Y)(E^+_v)$, and $k_v=1$ for almost all $v$ (both lattices
are the standard ones almost everywhere) (over an unramified extension a unimodular hermitian lattice has an orthonormal basis, proof: the reduction of the lattice is a non-degenerate hermitian space over the quadratic extension of residue fields, so it has a vector $\bar x$ with $\bar h(\bar x,\bar x)\ne0$ (otherwise polarisation and the surjectivity of the residue trace give $\bar h=0$); a lift $x$ has $h(x,x)$ a unit of the base ring, hence the norm of a unit since the extension is unramified, so after scaling $h(x,x)=1$; then the lattice is $\mathcal O x\perp(x^\perp\cap\text{lattice})$ with unimodular second summand, and one inducts on the rank; see also Jacobowitz~\cite{Jacobowitz}; at a split place every
unimodular lattice is standard); at the infinite places two orthonormal bases of the same signs
differ by an element of the real unitary group.
\end{proof}


\subsection{Elementary lemmas}\label{pcx:ss-elem}

\begin{lemma}[Standard data are matched]\label{pcx:std-matched}
For $z\in\pcMat_{b\times N}(\auO_S)$, $\Phi^{\mathrm{row}}_z$ and $\Psi^{\mathrm{col}}_z$ are matched
for $\chi_S$ over $\auO_S$, and $\Psi^{\mathrm{col}}_z$ is locally constant.
\end{lemma}
\begin{proof}
Let $k\in\auGL_N(\auO_S)$. Then
\[
\tbinom{z}{k^{[a]}}\,k=\begin{pmatrix}zk_{[b]}&z k_{(a)}\\0&1_a\end{pmatrix},
\]
where $k_{(a)}$ are the last $a$ columns of $k$, because the last $a$ rows of $k^{-1}k$ are
$(0\ 1_a)$. Taking determinants, $\det\binom z{k^{[a]}}\det k=\det(zk_{[b]})$. The matrix
$\binom z{k^{[a]}}$ is integral, and it lies in $\auGL_N(\auO_S)$ if and only if $\det(zk_{[b]})$ is a
unit; $k_{[b]}$ is integral. Hence both sides vanish together, and otherwise
$\Phi^{\mathrm{row}}_z(k^{[a]})=\chi_S(\det(zk_{[b]}))^{-1}\chi_S(\det k)=\chi_S(\det k)
\Psi^{\mathrm{col}}_z(k_{[b]})$. Supports: if $\Phi^{\mathrm{row}}_z(x)\ne0$ then $M=\binom zx\in
\auGL_N(\auO_S)$ and $x=k^{[a]}$ for $k=M^{-1}$. If $\Psi^{\mathrm{col}}_z(y)\ne0$ then $y$ is integral
and $(zy)^{-1}z$ is an integral left inverse of $y$; over the product of complete discrete
valuation rings $\auO_S$ this means that the reduction of $y$ modulo each prime has rank $b$, so the
columns of $y$ extend to a basis of $\auO_S^N$: there is $k\in\auGL_N(\auO_S)$ with $k_{[b]}=y$.
Finally $\Psi^{\mathrm{col}}_z$ is locally constant: the set of integral $y$ with $zy$ invertible is
open and closed, and on it $\chi_S^{-1}(\det(zy))$ is locally constant because $\chi_S$ is tame,
hence trivial on $1+\varpi\auO_S$.
\end{proof}

\begin{lemma}[Data modulo $\pcm$]\label{pcx:mod}
Let $\nu$ be a tame character for $\pcm$, $S$ the primes of $\pcm$, and let $\Phi$, $\Psi$ be
functions on $\pcMat_{a\times N}(\pcR_\pcm)$, $\pcMat_{N\times b}(\pcR_\pcm)$ matched for $\chi$. Let
$\Phi'(x)=\Phi(x\bmod\pcm)$ for integral $x\in\pcMat_{a\times N}(F_S)$ and $\Phi'(x)=0$ otherwise,
and define $\Psi'$ in the same way. Then $\Phi'$, $\Psi'$ are matched for $\chi_S$ over $\auO_S$,
$\Psi'$ is locally constant, and $[x,\Phi]=[x,\Phi']_S$, $[v,\Psi]=[v,\Psi']_S$ for all Fock
vectors $x$, $v$ (the left sides being the data of \S\ref{pc:ss-data}).
\end{lemma}
\begin{proof}
The reduction $\pi\colon\auO_S\to\pcR_\pcm=\prod_{w\in S}\auO_{E_w}/w^{n_w}$ is a surjective ring
homomorphism, and $u\in\auO_S$ is a unit if and only if $\pi(u)$ is (both mean: non-zero residue at
every $w\in S$). On units $\chi_S(u)=\prod_{w\in S}\chi_w(u_w\bmod w)=\chi(\pi(u))$, by the
definition of the tame parts. The equality of data is the definition of $\Phi^\natural$ in
\S\ref{pc:ss-data}; $\Psi'$ is locally constant since it factors through reduction modulo $\pcm$
on the open and closed set of integral matrices.

Supports. If $\Phi'(x)\ne0$, then $x$ is integral and $\pi(x)=\bar k^{[a]}$ for some $\bar k\in
\auGL_N(\pcR_\pcm)$. Lift the first $b$ rows of $\bar k^{-1}$ to an integral $z$ and put
$M=\binom zx$; $\pi(M)=\bar k^{-1}$, so $\det M$ is a unit and $M\in\auGL_N(\auO_S)$; then $x=k^{[a]}$
for $k=M^{-1}$. If $\Psi'(y)\ne0$, then $\pi(y)=\bar k_{[b]}$; lift the last $a$ columns of $\bar k$ to
an integral $u$ and put $k=(y\ u)$, whose reduction is $\bar k$; so $k\in\auGL_N(\auO_S)$ and
$y=k_{[b]}$. Values. For $k\in\auGL_N(\auO_S)$, $\pi(k^{[a]})=\pi(k)^{[a]}$ and $\pi(k_{[b]})=
\pi(k)_{[b]}$, so $\Phi'(k^{[a]})=\Phi(\pi(k)^{[a]})=\chi(\det\pi(k))\Psi(\pi(k)_{[b]})=
\chi_S(\det k)\Psi'(k_{[b]})$.
\end{proof}

\noindent\emph{Coordinates at $S$.} The data of \S\ref{pc:ss-data} are those of
Definition~\ref{pgx:conv}: at $w\in S$, $\Phi^\natural(x)=\Phi(a_w^{-1}x\bmod\pcm)$ on $a_w\pcMat(\auO_w)$,
$a_w$ a generator of $\alpha_w\auO_w$. Throughout \S\S\ref{pcx:sec}--\ref{pcn:sec}, matrices at $w\in S$
are written in the coordinates $\tilde x$ of Lemma~\ref{pgx:dict}: $\tilde x=a_w^{-1}x$ on the
$w$-component and $a_{\aubar w}^{-1}x$ on the $\aubar w$-component, in which both components of the
self-dual lattice are the integral matrices and the Heisenberg pairing is
$\lambda_w(\operatorname{tr}({}^t\tilde y\tilde x))^{\pm1}$ with $\lambda_w$ of conductor $\auO_w$. As $a_w$ is a scalar this
commutes with the actions $x\mapsto k'^{-1}xk$ of (W5) and with the measure normalisations; the vacuum
lattice $a_w\pcMat(\auO_w)$ (Lemma~\ref{pgx:vac}) becomes $\pcMat(\auO_w)$, and $\Phi^\natural$ becomes
the function $\Phi'$ of the lemma. So Lemma~\ref{pcx:mod} holds with the data of
Definition~\ref{pgx:conv}, $[x,\Phi']_S$ being read in these coordinates; where $\alpha_w$ is a unit,
$a_w=1$ and nothing changes. (See \S\ref{pgx:ss-other}.)

\subsection{Proof of Theorem~\ref{pcx:X} from the intermediate propositions}\label{pcx:ss-assembly}

\begin{proof}[Proof of Proposition~\ref{pcx:P-matches} from Propositions~\ref{pcx:P-std},
\ref{pcx:P-match}, \ref{pcx:P-flip}]
The space $U^-$ has $b-q_U$ negative signs, so $q_{U^-}+q=b+q_V$ is equivalent to $q=q_U+q_V$.
Let $y^{\mathrm{ad}}$ be the vector of Proposition~\ref{pcx:P-flip}; it has finite moments.
Proposition~\ref{pcx:P-std} for $(W,U^-,V)$, $z$ and $y^{\mathrm{ad}}$ gives $v$ and $C$, and for each
$\nu$ a number $c_0$ with $C|c_0|^2=L_*(\pcm,\nu,r-N/2)\pcxstr_{\pcxH\oplus W}(\nu;z)$ and
$\Theta^{\mathrm L}_{W,V}(v,\Phi^{\mathrm{row}}_z)=c_0\Theta^{\mathrm S}_{W,U^-}(y^{\mathrm{ad}},
\Psi^{\mathrm{col}}_z)$ for integral $z$. By Proposition~\ref{pcx:P-match}, for $\Phi,\Psi$ matched
with $\Psi$ locally constant, $\Psi$ is Schwartz--Bruhat on left-invertible matrices and
$\Theta^{\mathrm L}_{W,V}(v,\Phi)=c_0\Theta^{\mathrm S}_{W,U^-}(y^{\mathrm{ad}},\Psi)$. By
\eqref{pcx:eq-twist}, $\theta^{W,V}_{1,\gamma}(g,h)=\gamma(\det g)\theta^{W,V}_{\gamma,\gamma}(g,h)$, so
\[
\begin{aligned}&\int_{[\auU(V)]}\theta^{W,V}_{1,\gamma}(g,h;[v,\Phi]_S)\,dh=\gamma(\det g)\Theta^{\mathrm L}_{W,V}
(v,\Phi)(g)\\&\quad=c_0\,\gamma(\det g)\Theta^{\mathrm S}_{W,U^-}(y^{\mathrm{ad}},\Psi)(g)
\\&\quad=c_0d\int_{[\auU(U)]}\theta^{U,W}_{\gamma,\gamma}(h,g;[y,\Psi]_S)\,dh\end{aligned}
\]
by Proposition~\ref{pcx:P-flip}. Put $c=c_0d$; since $|d|=d_0$ does not depend on $\nu$,
$Cd_0^{-2}|c|^2=C|c_0|^2$ gives (a) with the constant $Cd_0^{-2}$.
\end{proof}

\begin{proof}[Proof of Theorem~\ref{pcx:X} from Propositions~\ref{pcx:P-matches},
\ref{pcx:P-equal} and Lemmas~\ref{pcx:Lpos}, \ref{pcx:mod}]
Put $k=N/2-q$. From $N\le2\mu_\tau+2q$ we get $k\le\mu_\tau$, and $\operatorname{rk}W=2k+2q$.
Proposition~\ref{pcx:P-equal} gives $z$ and $C_0\ne0$ with $\pcxstr_{\pcxH\oplus W}(\nu;z)=C_0$ for
all tame $\nu$ of weight $\mu$. Proposition~\ref{pcx:P-matches} for $(W,U,V)=(W,V',V)$, this $z$
and $y=v$ gives $C_1\in\R$ and $x$ (the vector of the theorem). Let $\nu$ be tame for $\pcm$ of weight
$\mu$ (note $\mu_\tau\ge k\ge1$ since $2q+2\le N$). The number $c$ of
Proposition~\ref{pcx:P-matches} satisfies, as $a-N/2=e$,
\[
C_1|c|^2=L_*(\pcm,\nu,e)\,C_0 .
\]
The right side is non-zero by Lemma~\ref{pcx:Lpos}; hence $C_1\ne0$, $c\ne0$, and
$|c|^2=C\,L_*(\pcm,\nu,e)$ with $C=C_0/C_1$, which is then positive. ($C$ does not depend on
$\pcm,\nu$; if no tame character of weight $\mu$ exists the theorem is empty and we take $C=1$.)

Let $\Phi,\Psi$ be matched for $\chi$ modulo $\pcm$, and let $\Phi',\Psi'$ be as in
Lemma~\ref{pcx:mod}. Proposition~\ref{pcx:P-matches}(b) for $\Phi',\Psi'$ is \eqref{pcx:eq-X} for
$\alpha=1$, $\alpha'=\gamma$. For general $\alpha\alpha'=\gamma$, \eqref{pcx:eq-twist} gives
$\theta^{W,V}_{\alpha,\gamma}(g,h)=\alpha(\det g)^{-1}\theta^{W,V}_{1,\gamma}(g,h)$ and
$\theta^{V',W}_{\gamma,\alpha'}(h',g)=(\alpha'\gamma^{-1})(\det g)\,\theta^{V',W}_{\gamma,\gamma}(h',g)
=\alpha(\det g)^{-1}\theta^{V',W}_{\gamma,\gamma}(h',g)$: both sides of \eqref{pcx:eq-X} are
multiplied by $\alpha(\det g)^{-1}$, and the identity persists with the same $c$.
\end{proof}

\subsection{The equal-rank ratio: proof of Proposition~\ref{pcx:P-equal}}\label{pcx:ss-equal}

Let $J$ be positive definite of rank $2k$, $Y=\pcxH\oplus J$, and $v$ a Fock vector of $(Y,Y)$ with
finite moments. Write $\mathrm{vol}_0$ for the volume of $[\auU(0)]$ (a point) and $\mathrm{vol}_J$
for that of the compact space $[\auU(J)]$; both are positive. For a datum $f$ of $(J,J)$ the kernel
$\theta^{J,J}_{\gamma,\gamma}(h,g;f)$ is continuous on the compact space $[\auU(J)]^2$, so Fubini
applies.

\begin{lemma}[The identity used twice]\label{pcx:twice}
Assume Proposition~\ref{pcx:P-matches}. There is $D\ge0$ with $\pcxstr_Y(\nu;v)^2=D$ for all tame
$\nu$ of weight $\mu$.
\end{lemma}
\begin{proof}
\emph{The triple $(J,J,0)$.} Apply Proposition~\ref{pcx:P-matches} with $(W,U,V)=(J,0,J)$ (ranks
$2k$, $0$, $2k$, sign condition $0=0+0$), $z=v$ and $y$ the vacuum of the zero Fock space
$\auHom(0,J)\otimes\R=0$. It gives $B$ and $y_J$ for $(J,J)$, and for each $\nu$ a number $c_1$ with
$B|c_1|^2=L_*(\pcm,\nu,k)\pcxstr_Y(\nu;v)$. Take $\Phi=\Phi^{\mathrm{row}}_\emptyset=\chi_S^{-1}\circ
\det$ (on $\auGL_{2k}(\auO_S)$, zero elsewhere) and $\Psi=\Psi^{\mathrm{col}}_\emptyset$, the
function $1$ on the one-point space $\pcMat_{2k\times0}$; they are matched by
Lemma~\ref{pcx:std-matched} (with $b=0$). On $\auHom(0,J)=0$ the theta functional of the datum
$[\mathrm{vac},1]$ is $1$, so the right side of (b) is $c_1\gamma(\det g)\mathrm{vol}_0$, while the left
side is $\gamma(\det g)\Theta^{\mathrm L}_{J,J}(y_J,\Phi)(g)$ by \eqref{pcx:eq-twist}. With
$f=[y_J,\chi_S^{-1}\circ\det]_S$:
\begin{equation}\label{pcx:eq-A}
\int_{[\auU(J)]}\theta^{J,J}_{\gamma,\gamma}(g,h;f)\,dh=c_1\,\mathrm{vol}_0\qquad(g\in\auU(J)(\auA)).
\end{equation}
\emph{The triple $(J,0,J)$.} Apply Proposition~\ref{pcx:P-matches} with $(W,U,V)=(J,J,0)$, $z=v$,
$y=y_J$. It gives $C$ and a vector $x$ of the zero Fock space $\auHom(J,0)\otimes\R$, and for each
$\nu$ a number $c_2$ with $C|c_2|^2=L_*(\pcm,\nu,-k)\pcxstr_Y(\nu;v)$. Take $\Phi=1$ on
$\pcMat_{0\times2k}$ and $\Psi=\chi_S^{-1}\circ\det$, matched by Lemma~\ref{pcx:std-matched} with
$z=1_{2k}$. On $\auHom(J,0)=0$ the kernel $\theta^{J,0}_{1,\gamma}(g,h;[x,1]_S)$ is the number
$\xi=\theta([x,1]_S)$, which does not depend on $g$, on $h\in\auU(0)$ or on $\nu$ (the datum is $x$
tensored with the constant $1$ at every finite place). Hence
\begin{equation}\label{pcx:eq-B}
\mathrm{vol}_0\,\xi=c_2\int_{[\auU(J)]}\theta^{J,J}_{\gamma,\gamma}(h,g;f)\,dh\qquad(g\in\auU(J)(\auA)).
\end{equation}
\emph{Comparison.} Integrate \eqref{pcx:eq-B} over $g\in[\auU(J)]$ and use Fubini and
\eqref{pcx:eq-A} (with the roles of the variables renamed):
$\mathrm{vol}_J\mathrm{vol}_0\,\xi=c_2\int_{[\auU(J)]}\int_{[\auU(J)]}\theta^{J,J}_{\gamma,\gamma}(h,g;f)
\,dg\,dh=c_2c_1\mathrm{vol}_0\mathrm{vol}_J$, so $\xi=c_1c_2$. Since $L_*(k)L_*(-k)=1$,
\[
\pcxstr_Y(\nu;v)^2=\bigl(L_*(\pcm,\nu,k)\pcxstr_Y(\nu;v)\bigr)\bigl(L_*(\pcm,\nu,-k)\pcxstr_Y(\nu;v)
\bigr)=BC|c_1c_2|^2=BC|\xi|^2=:D,
\]
which does not depend on $\nu$.
\end{proof}

\begin{proof}[Proof of Proposition~\ref{pcx:P-equal} from Propositions~\ref{pcx:P-matches},
\ref{pcx:P-nonvan}, \ref{pcx:P-plane}, \ref{pcx:P-signs} and Lemma~\ref{pcx:adapted}]
If there is no tame character of weight $\mu$ there is nothing to prove. Otherwise fix one, $\nu_0$.

\emph{First floor: $W=J$ the standard positive definite space of rank $2k$.} By
Proposition~\ref{pcx:P-nonvan} there is $z$ with $\pcxstr_Y(\nu_0;z)\ne0$. By
Lemma~\ref{pcx:twice}, $\pcxstr_Y(\nu;z)^2=\pcxstr_Y(\nu_0;z)^2$ for all $\nu$; as
$\pcxstr_Y\ge0$ by definition (Lemma~\ref{pcx:Lpos}), $\pcxstr_Y(\nu;z)=\pcxstr_Y(\nu_0;z)=:C_0\ne0$.

\emph{Induction on $q$.} For $q=0$ the space $\pcxH\oplus W$ has the signs of $Y=\pcxH\oplus J$, and
$Y$ contains the adapted plane $\pcxH$; Proposition~\ref{pcx:P-signs} gives
$\pcxstr_{\pcxH\oplus W}(\nu;z)=\pcxstr_Y(\nu;z)=C_0$. For $q+1$: let $W_0=\pcxH^q\oplus J$, of rank
$2k+2q$ with $q$ negative signs; by induction there are $z_0$ and $C\ne0$ with
$\pcxstr_{\pcxH\oplus W_0}(\nu;z_0)=C$ for all $\nu$. By Lemma~\ref{pcx:adapted} there is a space $Y'$
with the signs of $\pcxH\oplus W$ and an adapted splitting $Y'=(\pcxH\oplus W_0)\oplus\pcxH$
(the ranks differ by $2$ and the numbers of negative signs by $1$). Proposition~\ref{pcx:P-plane}
gives $z'$ and $C'>0$ with $\pcxstr_{Y'}(\nu;z')=C'C$, and Proposition~\ref{pcx:P-signs} gives
$\pcxstr_{\pcxH\oplus W}(\nu;z')=\pcxstr_{Y'}(\nu;z')=C'C\ne0$.
\end{proof}

\subsection{Convergence}\label{pcx:ss-conv}

For a space $S_1$ of rank $n_1$, a space $S_2$ of rank $n_2$ and a set $P$ of primes of $E$ with
$P\cap\bar P=\emptyset$, consider data $[y,\Phi]_P$ of $\auHom(S_1,S_2)$, $\Phi$ locally constant
with compact support on $\pcMat_{n_2\times n_1}(F_P)$. Write $\mathrm{ht}$ for the height on the
adelic unitary groups (\S\ref{ss:au-heights}).

\begin{proposition}[Decay under a rank condition]\label{pcx:decay}
\textup{(a)} Suppose that every $x$ with $\Phi(x)\ne0$ is surjective at every place of $P$ (each
component has a right inverse). For every $N_0$ there are $B,M,C$ such that for all
Fock vectors $y$ with finite moment $\|y\|_M=\int(1+|z|)^M|y(z)|\,dz$, all characters
$\alpha,\beta$ and all $(g,h)$,
\[
|\theta^{S_1,S_2}_{\alpha,\beta}(g,h;[y,\Phi]_P)|\,\mathrm{ht}(h)^{N_0}\le C\,\|y\|_M\,
\mathrm{ht}(g)^B .
\]
\textup{(b)} If every $x$ with $\Phi(x)\ne0$ is injective at every place of $P$, the same holds with
the roles of $g$ and $h$ exchanged.
\end{proposition}
\noindent The proof occupies the rest of this subsection. Heights are those of
Definition~\ref{pcx:def-ht} below (not the heights of families of \S\ref{ss:au-heights}).

\begin{definition}[Lattice and height at a point]\label{pcx:def-ht}
Let $S_1$ be a space of rank $n$ with frame giving $S_1(\auA)\cong\auA_E^n$. For $g\in\auU(S_1)(\auA)$
the \emph{lattice at $g$} is $\Lambda_g=\{\gamma\in S_1(E): g^{-1}\gamma\text{ is integral at all
finite places}\}$, with the length $|\gamma|_g=\|(g^{-1}\gamma)_\infty\|$ (the euclidean norm of the
archimedean coordinates). Let $\lambda_1(g)=\min\{|\gamma|_g:0\ne\gamma\in\Lambda_g\}$ and
$\mathrm{ht}(g)=\max(1,\lambda_1(g)^{-1})$.
\end{definition}

We use three classical facts.
\begin{enumerate}
\item[(R1)] (Reduction theory.) For $N_0$ large, $\int_{[\auU(S_1)]}\mathrm{ht}(g)^{-N_0}dg<\infty$; and
there is $Z$ such that $\Lambda_g$ has an $E$-basis $c_1,\dots,c_n$ with $|c_j|_g\le Z\,\mathrm{ht}(g)^n$
(Minkowski's second theorem: $\Lambda_g$ is a $\Z$-lattice of rank $n[E:\Q]$ and covolume independent
of $g$, so the product of its $n[E:\Q]$ successive minima is bounded above and below. It is an
$\auO_E$-module and $|\omega\gamma|_g\le C_E|\gamma|_g$ for $\omega$ in a fixed $\Z$-basis of $\auO_E$,
so if $\mu_1\le\dots\le\mu_n$ are the successive minima over $E$ (the least $\mu$ such that the
vectors of length $\le\mu$ span an $E$-subspace of dimension $j$), the $\Z$-minima are, up to
constants depending only on $E$, the $\mu_j$ each repeated $[E:\Q]$ times. Hence
$\prod_j\mu_j^{[E:\Q]}\asymp1$, and with $\mu_j\ge\lambda_1(g)\ge\mathrm{ht}(g)^{-1}$ this gives
$\mu_n\ll\mathrm{ht}(g)^{n-1}\le\mathrm{ht}(g)^n$; a basis realising the $\mu_j$ is an $E$-basis). (Borel, \emph{Introduction aux groupes arithm\'etiques}, reduction theory.)
\item[(R2)] (Moderate growth of theta functions; Lemma~\ref{pcx:thetagrowth} below.) Let
$X=\auHom(S_1,S_2)$, $2D=\dim_\R(X\otimes\R)$, $s=D+1$, and let $\|\cdot\|_s$ and $\pzF_s$ be the norm and
the Banach space of Lemma~\ref{pz:growth} for the archimedean Fock space of $X$. A \emph{datum of class
$s$} is a finite sum of tensors $\Psi\otimes f$ with $\Psi\in\pzF_s$ and $f$ a finite datum; $\theta$ is
defined on these and continuous in the first factor (Lemma~\ref{pz:growth}(c)), and the Weil
representation preserves them (Lemma~\ref{pz:growth}(a) at infinity). For a linear functional
$\Lambda$ on the data of class $s$ and a finite datum $f_0$ put
\[
\|\Lambda\|_{f_0}=\sup\Bigl\{\frac{|\Lambda(\psrho_\infty(z)\Psi\otimes f_0)|}{\|\Psi\|_s}:
z\in X\otimes\R,\ 0\ne\Psi\in\pzF_s\Bigr\}\in[0,\infty]
\]
($\psrho_\infty(z)$ preserves $\pzF_s$ by Lemma~\ref{pcx:superpos}(a)). As $\psrho_\infty(y)\psrho_\infty(z)$
is a scalar of absolute value $1$ times $\psrho_\infty(y+z)$, $\|\Lambda\circ\psrho_\infty(y)\|_{f_0}=
\|\Lambda\|_{f_0}$ for $y\in X\otimes\R$. Lemma~\ref{pcx:thetagrowth} says: for every finite datum $f_0$
there are $Y,A$ with
\[
\|\theta\circ\omega_{\alpha,\beta}(g,h)\|_{f_0}\le Y(\mathrm{ht}(g)\mathrm{ht}(h))^A
\]
for all characters $\alpha,\beta$ and all $g,h$. Since $\mathrm{ht}\ge1$ we may and do take $A\ge2s$.
\item[(R3)] (Short vectors are isotropic.) If $\gamma,\gamma'\in\Lambda_g$ and $|\gamma|_g|\gamma'|_g<1$
then $\langle\gamma,\gamma'\rangle=0$: indeed $\langle\gamma,\gamma'\rangle\in\auO_E$ (the lattice at $g$
is unimodular, $g$ being unitary at the finite places), and each archimedean absolute value is at
most $|\gamma|_g|\gamma'|_g<1$ (Cauchy--Schwarz in the coordinates at $g$, where the form is
$\mathrm{diag}(\pm1)$); an algebraic integer all of whose conjugates have absolute value $<1$ is $0$.
\end{enumerate}
In particular, if $\mathrm{ht}(h)>2$ then $\lambda_1(h)<1/2$ and there is $0\ne d\in\Lambda_h$ with
$|d|_h<2\lambda_1(h)\le 2/\mathrm{ht}(h)<1$, isotropic by (R3).

\begin{lemma}[Reduction to a Siegel set]\label{pcx:reduce}
Let $S_1$ be a space of rank $n$ with a frame, $G=\auU(S_1)$, viewed as a connected reductive group over
$E^+$, and $d=[E^+:\Q]$. There are a compact set $C_f\subset G(\auA_f)$ and a constant $C$ such that every
$g\in G(\auA)$ can be written $g=\gamma^{-1}x$ with $\gamma\in G(E^+)$, $x=x_\infty x_f$, $x_f,x_f^{-1}\in C_f$
and $\|x_\infty\|,\|x_\infty^{-1}\|\le C\,\mathrm{ht}(g)$ (operator norms in the
coordinates of the frame).
\end{lemma}
\begin{proof}
If $S_1$ is anisotropic, $G(\auA)/G(E^+)$ is compact~\cite[Theorem~5.6(ii)]{BorelFin}, so $G(\auA)=
C\,G(E^+)$ with $C$ compact; apply this to $g^{-1}$. Otherwise let $r\ge1$ be the Witt index and
$e_1,\dots,e_r,e'_1,\dots,e'_r$ rational vectors with $\langle e_i,e'_j\rangle=\delta_{ij}$, $\langle e_i,e_j
\rangle=\langle e'_i,e'_j\rangle=0$, spanning a subspace with anisotropic orthogonal complement $S_0$. The
torus $T$ acting by $e_i\mapsto a_ie_i$, $e'_i\mapsto a_i^{-1}e'_i$ ($a_i\in\mathrm{GL}_1$ over $E^+$) and
trivially on $S_0$ is a maximal $E^+$-split torus, and the stabiliser $P$ of the flag $Ee_1\subset Ee_1+
Ee_2\subset\cdots$ is a minimal parabolic subgroup containing it; the roots of $T$ in the Lie algebra of
the unipotent radical of $P$ are $a_i/a_j$ $(i<j)$, $a_ia_j$, $a_i^2$ and, if $S_0\ne0$, $a_i$, so the simple
roots are $a_i/a_{i+1}$ $(i<r)$ and $a_r$ or $a_r^2$. By Godement's form of Borel's
theorem~\cite[Th\'eor\`eme~7]{Godement63} there are a compact $M\subset G(\auA)$ with $G(\auA)=MP(\auA)$, a
relatively compact $F\subset P(\auA)^0$ with $P(\auA)^0=F\,P(E^+)$ (here $P(\auA)^0$ is the common kernel of
the $|\chi|$, $\chi$ a rational character of $P$; it is normal in $P(\auA)$ and contains $P(E^+)$ by the
product formula) and $c\ge1$ such that $G(\auA)=M\,T_\auA(c)\,F\,G(E^+)$, where $T_\auA(c)$ is the set of $t\in
T(\auA)$ with $|\alpha(t)|<c$ for the simple roots $\alpha$ (idelic absolute value).

Write $g^{-1}=mtf\gamma_0$ accordingly. \emph{Normalising $t$.} Let $X_f\subset\auA_{E^+,f}^\times$ be a
finite set of representatives of the ideal classes of $E^+$. Every idele $a$ of $E^+$ is $\xi b$ with
$\xi\in E^{+\times}$, $b_f\in X_f\widehat\auO_{E^+}^\times$, and, by Dirichlet's unit theorem (the logarithms
of the units form a lattice in the trace-zero hyperplane), $C_1^{-1}|a|^{1/d}\le|b_\tau|\le C_1|a|^{1/d}$ at
every real place $\tau$ ($b\mapsto\xi\varepsilon b$, $\varepsilon$ a unit, $|a|$ the idelic absolute value, which
equals $|b|$). Doing this for each $a_i(t)$ gives $t=\xi t'$ with $\xi\in T(E^+)$, $t'_f$ in a fixed
compact set, and $|a_i(t'_\tau)|\in[C_1^{-1},C_1]\,|a_i(t)|^{1/d}$; and $|\alpha(t')|=|\alpha(t)|<c$. Then
$g^{-1}=mt'(\xi f\xi^{-1})\xi\gamma_0$ and $\xi f\xi^{-1}\in P(\auA)^0=FP(E^+)$, say $\xi f\xi^{-1}=f'p$.
So $g^{-1}=mt'f'\gamma'$ with $\gamma'=p\xi\gamma_0\in G(E^+)$, $f'\in F$, $t'\in T_\auA(c)$ normalised.

\emph{Bounds.} Put $\gamma=\gamma'$ and $x=\gamma g=f'^{-1}t'^{-1}m^{-1}$. Its finite part and that of its
inverse $mt'f'$ lie in a fixed compact set $C_f$ ($M$, $\bar F$ and the finite parts of the normalised
$t'$ do). From $|a_i/a_{i+1}|<c$ and $|a_r|<c$ (or $|a_r|^2<c$), with $a_i=a_i(t)$: $|a_i|<c^r$ and
$|a_i|^{-1}<c^{r}|a_1|^{-1}$ for all $i$. The archimedean eigenvalues of $t'$ in the basis
$(e_i,e'_i,\text{basis of }S_0)$ are $a_i(t'_\tau)^{\pm1}$ and $1$, so, the change of basis to the frame being
fixed, $\|x_\infty\|,\|x_\infty^{-1}\|\le C_2\max_i\max(|a_i|^{1/d},|a_i|^{-1/d})\le C_3\max(1,|a_1|^{-1/d})$.
Finally the vector $\eta=N\gamma'^{-1}e_1\in S_1(E)$ satisfies $g^{-1}\eta=Nmt'f'e_1=N\,\chi(f')a_1(t')\,me_1$,
where $f'e_1=\chi(f')e_1$ ($P$ stabilises $Ee_1$) with $\chi(f')$ in a relatively compact subset of
$\auA_E^\times$. The finite parts of $\chi(f')a_1(t')me_1$ lie in a compact subset of $S_1(\auA_f)$, hence in
$N^{-1}$ times the integral points for a fixed integer $N\ge1$; so $\eta\in\Lambda_g$, and
$|\eta|_g\le C_4|a_1|^{1/d}$. Thus $\lambda_1(g)\le C_4|a_1|^{1/d}$, i.e.\ $|a_1|^{-1/d}\le C_4\,\mathrm{ht}(g)$,
so $\|x_\infty\|,\|x_\infty^{-1}\|\le C_3\max(1,C_4\,\mathrm{ht}(g))\le C\,\mathrm{ht}(g)$, since
$\mathrm{ht}(g)\ge1$.
\end{proof}

\begin{lemma}[Growth of theta functions]\label{pcx:thetagrowth}
Let $S_1,S_2$ be spaces, $X=\auHom(S_1,S_2)$, $D=\frac12\dim_\R(X\otimes\R)$, $s=D+1$, $\pzF_s$ and
$\|\cdot\|_s$ as in Lemma~\ref{pz:growth} for $W=X$, and $f_0$ a finite datum. There are $Y,A$ such that for
all characters $\alpha,\beta$, $g\in\auU(S_1)(\auA)$, $h\in\auU(S_2)(\auA)$, $z\in X\otimes\R$ and
$\Psi\in\pzF_s$,
\[
|\theta(\omega_{\alpha,\beta}(g,h)(\psrho_\infty(z)\Psi\otimes f_0))|\le Y\,\|\Psi\|_s\,
(\mathrm{ht}(g)\mathrm{ht}(h))^{A};\qquad A=2s\text{ will do.}
\]
\end{lemma}
\begin{proof}
\emph{Lemma~\ref{pz:growth} applies to $X$.} Its statement and proof concern $W=\auHom(Y_j,Y_{j+1})$ but
use only: the Fock realisation of the archimedean Heisenberg representation of $W$ for a positive complex
structure commuting with the maximal compact subgroups, the commutation of $\omega_\infty$ with the
Heisenberg operators (W1), the continuity of the Weil representation on smooth vectors, the existence of
a non-zero theta functional continuous in the archimedean factor (W0), and the action of a lattice on
finite data by characters (W4); none of the properties of the chain $Y_j$ enters. All of these hold
for any pair of spaces, so (a), (b), (c) hold for $X$ with $D$ as above.

\emph{Uniformity in $z$.} In the proof of (c), replacing $\Psi$ by $\psrho_\infty(z)\Psi$ replaces
$v=\gamma^\natural-w_0$ by $v+z$ (up to a phase), so $|\theta(\psrho_\infty(z)\Psi\otimes f)|\le C(D+1)^{2s}\|f\|
\|f'\|\,\|\Psi\|_s\sup_{u}\sum_{\lambda}(1+|\lambda+u|)^{-2s}$, the sum over a coset of a fixed lattice of
$X\otimes\R\cong\R^{2D}$ and the supremum over $u\in\R^{2D}$; this supremum is finite as $2s>2D$. The
lattice depends on $f$ only through the integer $\mathsf N$ of that proof. Hence for every finite set
$\mathcal S$ of finite data there is $C_{\mathcal S}$ with $|\theta(\psrho_\infty(z)\Psi\otimes f)|\le
C_{\mathcal S}\|\Psi\|_s$ for $f\in\mathcal S$, all $z$, all $\Psi$.

\emph{Reduction.} By Lemma~\ref{pcx:reduce} write $g=\gamma^{-1}x$, $h=\delta^{-1}y$ with rational
$\gamma,\delta$. The theta functional is invariant under the rational points ((W1) of Lemma~\ref{ps:weil})
and the characters are trivial on them, so the left side is
$|\theta(\omega_{\alpha,\beta}(x,y)(\psrho_\infty(z)\Psi\otimes f_0))|=|\theta(\psrho_\infty(z')\,\omega_\infty
(x_\infty,y_\infty)\Psi\otimes\omega_f(x_f,y_f)f_0)|$ with $z'=y_\infty zx_\infty^{-1}$, the characters
contributing factors of absolute value $1$. As $f_0$ is fixed by an open subgroup $K'$ of the finite
group and $(x_f,y_f)$ runs over a compact set, covered by finitely many cosets of $K'$, the vectors
$\omega_f(x_f,y_f)f_0$ lie, up to factors of absolute value $1$, in a finite set $\mathcal S$. By the
previous step and Lemma~\ref{pz:growth}(a),
the left side is at most $C_{\mathcal S}(2D\|x_\infty\|\|y_\infty\|)^{2s}\|\Psi\|_s\le Y(\mathrm{ht}(g)
\mathrm{ht}(h))^{2s}\|\Psi\|_s$ by the bounds of Lemma~\ref{pcx:reduce}.
\end{proof}

\begin{lemma}[Fock vectors are superpositions of translated vacua]\label{pcx:superpos}
Let $\mathcal F$ be the Fock space of a real hermitian space $\C^m$ (real dimension $2m$), realised as in
\S\ref{psp:ss-fock} as the space of functions $v\in L^1\cap L^2$ on $\C^m$ with $v(x)=\int v(y)\beta(y,x)
e(x-y)\,dy$, where $e(x)=e^{-|x|^2/2}$ is the vacuum, $|\beta|=1$, and the Heisenberg translations are
$(\rho(y)v)(x)=\beta(y,x)v(x-y)$; the measure is the normalised Lebesgue measure. For $N\ge0$ put
$\|v\|_N=\int(1+|x|)^N|v(x)|\,dx\in[0,\infty]$. Let $s\ge0$ be an integer, and let $\|\cdot\|_s$ and
$\mathcal F_s\subseteq\mathcal F$ be the norm and the Banach space of Lemma~\ref{pz:growth} (with $D=m$),
transported by the unitary isomorphism with the model \eqref{pk:eq-fock} that intertwines the
Heisenberg translations and carries vacuum to vacuum; the operators $\partial_w$, being the derivatives
of $\rho$, correspond. Let $c_\rho$ be a constant with
$\rho(w)\rho(y)=e^{\kappa(w,y)}\rho(y)\rho(w)$, $\kappa(w,y)\in i\R$ real bilinear,
$|\kappa(w,y)|\le c_\rho|w||y|$ (in the model \eqref{pk:eq-fock}, $\kappa(w,y)=2i\Iml\operatorname{tr}(y^*w)$ by the
formula for $\psrho_\infty(x)\psrho_\infty(y)$ recalled after it, so $c_\rho=2$ will do).
\begin{enumerate}
\item[\textup{(a)}] For $y\in\C^m$, $\rho(y)$ maps $\mathcal F_s$ into itself, with
$\|\rho(y)\Psi\|_s\le(1+c_\rho|y|)^{2s}\|\Psi\|_s$, and $y\mapsto\rho(y)\Psi$ is continuous into
$\mathcal F_s$ for every $\Psi\in\mathcal F_s$.
\item[\textup{(b)}] If $\|v\|_N<\infty$ then $v=\int v(y)\rho(y)e\,dy$, a Bochner integral converging for
the norm $\|\cdot\|_N$, because $\|\rho(y)e\|_N\le(1+|y|)^N\|e\|_N$. If $\|v\|_{2s}<\infty$, then
$v\in\mathcal F_s$, the same integral converges in $\mathcal F_s$, and $\|v\|_s\le c_1\|v\|_{2s}$ with
$c_1=\max(1,c_\rho)^{2s}\|e\|_s$.
\item[\textup{(c)}] Let $\ell$ be a continuous linear functional on $\mathcal F_s$ with $|\ell(\rho(z)e)|\le
S(1+|z|)^N$ for all $z$. Then $|\ell(v)|\le S\|v\|_N$ for every $v$ with $\|v\|_{\max(N,2s)}<\infty$.
\end{enumerate}
\end{lemma}
\begin{proof}
(a) Let $u$ be a smooth vector and $w\in\mathsf B$ (the real basis defining $\|\cdot\|_s$).
Differentiating $\rho(tw)\rho(y)u=e^{t\kappa(w,y)}\rho(y)\rho(tw)u$ at $t=0$ gives $\partial_w\rho(y)u=
\rho(y)(\partial_w+\kappa(w,y))u$; by induction,
\[
\partial_{w_1}\cdots\partial_{w_n}\rho(y)u=\rho(y)\,(\partial_{w_1}+\kappa_1)\cdots(\partial_{w_n}+\kappa_n)u,
\qquad\kappa_i=\kappa(w_i,y).
\]
Expanding the product, the right side is $\rho(y)$ applied to a sum over the subsets $\Sigma\subseteq
\{1,\dots,n\}$ of $\prod_{i\in\Sigma}\kappa_i$ times $\partial_{w_{j_1}}\cdots\partial_{w_{j_r}}u$, where
$j_1<\dots<j_r$ are the indices not in $\Sigma$. For $n\le2s$ each such word has norm at most
$\|u\|_s$, $|\prod_{i\in\Sigma}\kappa_i|\le(c_\rho|y|)^{|\Sigma|}$, and $\rho(y)$ is unitary; so
$\|\partial_{w_1}\cdots\partial_{w_n}\rho(y)u\|\le(1+c_\rho|y|)^n\|u\|_s$, and $\|\rho(y)u\|_s\le
(1+c_\rho|y|)^{2s}\|u\|_s$. The same identity shows that $y\mapsto\partial_{w_1}\cdots\partial_{w_n}
\rho(y)u$ is continuous into $\mathcal F$ ($\rho$ is strongly continuous on $\mathcal F$ and the $\kappa_i$
are continuous in $y$), i.e.\ $y\mapsto\rho(y)u$ is continuous into $\mathcal F_s$. Since the smooth
vectors are dense in $\mathcal F_s$ and $\rho(y)$ is bounded on them by a constant that is locally
uniform in $y$, both statements extend to $\mathcal F_s$ (consistently with $\rho(y)$ on $\mathcal F$, as
$\|\cdot\|\le\|\cdot\|_s$).

(b) The reproducing formula says $v(x)=\int v(y)(\rho(y)e)(x)\,dy$ for every $x$. From $1+|x|\le
(1+|y|)(1+|x-y|)$ and $|(\rho(y)e)(x)|=e(x-y)$ we get $\|\rho(y)e\|_N\le(1+|y|)^N\|e\|_N$; the map
$y\mapsto\rho(y)e$ is continuous for $\|\cdot\|_N$ and $\int|v(y)|\|\rho(y)e\|_N\,dy\le\|e\|_N\|v\|_N$, so
the Bochner integral exists in the space of functions with finite $\|\cdot\|_N$; evaluating it
pointwise (which is legitimate in $L^1_{\mathrm{loc}}$) gives $v$. If $\|v\|_{2s}<\infty$, the map
$y\mapsto v(y)\rho(y)e$ is continuous into $\mathcal F_s$ by (a), with norm at most
$|v(y)|(1+c_\rho|y|)^{2s}\|e\|_s\le c_1(1+|y|)^{2s}|v(y)|$, which is integrable; so the Bochner integral
exists in $\mathcal F_s$, with norm at most $c_1\|v\|_{2s}$. Its image in $\mathcal F$ is the $\mathcal
F$-valued integral, which is $v$: pairing with $\rho(x)e$ and using $c(x)=\langle\rho(x)e,c\rangle$ for
$c\in\mathcal F$ (\S\ref{psp:ss-fock}) gives $\int v(y)(\rho(y)e)(x)\,dy=v(x)$.

(c) By (b) with the index $\max(N,2s)$, $v=\int v(y)\rho(y)e\,dy$ in $\mathcal F_s$, and $\ell$, being
continuous, commutes with the Bochner integral: $|\ell(v)|\le\int|v(y)||\ell(\rho(y)e)|\,dy\le
S\int(1+|y|)^N|v(y)|\,dy=S\|v\|_N$.
\end{proof}

\begin{lemma}[Short differences gain powers]\label{pcx:diffgain}
Keep the notation of Lemma~\textup{\ref{pcx:superpos}}. For a linear functional $\ell$ on $\mathcal F_s$
put
\[
\|\ell\|^*=\sup\bigl\{|\ell(\rho(z)\Psi)|/\|\Psi\|_s:\ z\in\C^m,\ 0\ne\Psi\in\mathcal F_s\bigr\}\in[0,\infty].
\]
For each $N\ge0$ there is $C$, depending only on $N$, $m$, $s$ and $c_\rho$, with the following
property. Let $t\ge0$ and let $\ell$ be a linear functional on $\mathcal F_s$ with $\|\ell\|^*<\infty$
(so $\ell$ is continuous) such that, if $t<1$: for every $f\in\mathcal F_s$ and every $B\ge0$, if
$|\ell(f)-\ell(\rho(y)f)|\le B$ for all $|y|\le t$, then $|\ell(f)|\le B$. Then
\[
|\ell(v)|\le C\,\|\ell\|^*\,t^N\,\|v\|_N\qquad\text{for every }v\in\mathcal F\text{ with }\|v\|_{\max(N,2s)}<\infty .
\]
\end{lemma}
\begin{proof}
If $t\ge1$: $|\ell(\rho(z)e)|\le\|\ell\|^*\|e\|_s\le\|\ell\|^*\|e\|_st^N$ for all $z$, and
Lemma~\ref{pcx:superpos}(c) with the index $0$ gives $|\ell(v)|\le\|\ell\|^*\|e\|_st^N\|v\|_0\le
\|\ell\|^*\|e\|_st^N\|v\|_N$.

Let $t<1$. For a polynomial $P$ on $\C^m$ let $v_P$ be the vector of $\mathcal F$ that corresponds to $P$
in the model \eqref{pk:eq-fock} (a smooth vector, $v_1=e$). For $w\in\C^m$, $\partial_wv_P=v_{P_w}$
with $P_w(z)=(z\cdot w)P(z)-(\mathsf D_{\bar w}P)(z)$ (proof of Lemma~\ref{pz:growth}(b)), a polynomial
depending $\R$-linearly on $w$; write $\mathsf B=(b_1,\dots,b_{2m})$, so that $\partial_y=\sum_ky_k
\partial_{b_k}$ for $y=\sum_ky_kb_k$, $|y_k|\le|y|$. We show by induction on $N$:

\smallskip\noindent$(H_N)$ \emph{for every polynomial $P$ there is $C_{N,P}$, depending only on
$N,P,m,s,c_\rho$, with $|\ell(\rho(z)v_P)|\le C_{N,P}\|\ell\|^*t^N(1+|z|)^N$ for all $z\in\C^m$.}

\smallskip\noindent For $N=0$ take $C_{0,P}=\|v_P\|_s$. Assume $(H_N)$, fix $P$, $z$ and $y$ with
$|y|\le t$, and put $f(\sigma)=\ell(\rho(\sigma y)\rho(z)v_P)$ for $\sigma\in[0,1]$. With
$\kappa=\kappa(y,z)$ we have $\rho(\sigma y)\rho(z)=e^{\sigma\kappa}\rho(z)\rho(\sigma y)$; the orbit map
$\sigma\mapsto\rho(\sigma y)v_P$ of the smooth vector $v_P$ is $C^1$ for the topology of smooth vectors,
hence into $\mathcal F_s$, with derivative $\rho(\sigma y)\partial_yv_P$ ($\rho$ is additive on the line
$\R y$); $\rho(z)$ is bounded on $\mathcal F_s$ and $\ell$ is continuous. Hence
\[
f'(\sigma)=e^{\sigma\kappa}\,\ell\bigl(\rho(z)\rho(\sigma y)(\kappa v_P+\partial_yv_P)\bigr),
\]
and $\rho(z)\rho(\sigma y)$ is a scalar of absolute value $1$ times $\rho(z+\sigma y)$. Since
$|\kappa|\le c_\rho|y||z|\le c_\rho t(1+|z|)$, $|y_k|\le t$ and $1+|z+\sigma y|\le2(1+|z|)$, $(H_N)$
for $P$ and for the $P_{b_k}$ gives
\[
|f'(\sigma)|\le2^N\Bigl(c_\rho C_{N,P}+\sum_kC_{N,P_{b_k}}\Bigr)\|\ell\|^*t^{N+1}(1+|z|)^{N+1}
=:C_{N+1,P}\|\ell\|^*t^{N+1}(1+|z|)^{N+1}.
\]
By the mean value inequality, $|\ell(\rho(z)v_P)-\ell(\rho(y)\rho(z)v_P)|=|f(0)-f(1)|$ is at most this
bound, for every $|y|\le t$. The hypothesis on $\ell$, applied to $f=\rho(z)v_P\in\mathcal F_s$, gives
$|\ell(\rho(z)v_P)|\le C_{N+1,P}\|\ell\|^*t^{N+1}(1+|z|)^{N+1}$, which is $(H_{N+1})$.

Finally $(H_N)$ for $P=1$ and Lemma~\ref{pcx:superpos}(c) with $S=C_{N,1}\|\ell\|^*t^N$ give the lemma
with $C=\max(C_{N,1},\|e\|_s)$.
\end{proof}

\begin{lemma}[Comparison with the mean over an isotropic subspace]\label{pcx:mean}
Let $X$ be a space (here $X=\auHom(S_1,S_2)$), $s$, the data of class $s$ and $\|\cdot\|_{f_0}$ as in
\textup{(R2)}, $f_0$ a finite datum (a finite combination of finite Heisenberg translates of the
vacuum), $l\ge0$ and $\iota$ a finite set. There is $C$ such that: if $k\in\auU(X)(\auA)$ (acting
through the Weil representation), $\delta\ge0$, $z_i\in\Lambda_k$ $(i\in\iota)$ with $|z_i|_k\le\delta$
span an isotropic subspace $Z\subseteq X(E)$, $\Lambda$ is a linear functional on the data of class $s$,
invariant under $\rho(Z(E))$, such that $x\mapsto\Lambda(\rho(x)F)$ is continuous on $Z(\auA)$ for every
datum $F$ of class $s$ and $\|\Lambda\circ\omega(k)\|_{f_0}<\infty$, and
$\Lambda_0(F)=\int_{Z(\auA)/Z(E)}\Lambda(\rho(x)F)\,dx$ (probability Haar measure on the compact group
$Z(\auA)/Z(E)$), then for every Fock vector $v$ of $X\otimes\R$ with $\|v\|_{\max(l,2s)}<\infty$,
\[
|(\Lambda-\Lambda_0)(\omega(k)(v\otimes f_0))|\le C\,\|\Lambda\circ\omega(k)\|_{f_0}\,\delta^l\,\|v\|_l .
\]
(Here $v\in\pzF_s$ by Lemma~\textup{\ref{pcx:superpos}(b)}, so $v\otimes f_0$ is a datum of class $s$.)
\end{lemma}
\begin{proof}
Since $Z$ is isotropic, $\rho$ is additive on $Z(\auA)$ (no cocycle), so for $F$ of class $s$ the
function $x\mapsto\Lambda(\rho(x)F)$ is $Z(E)$-periodic and continuous; $\Lambda_0$ is a linear
functional on the data of class $s$, invariant under $\rho(Z(\auA))$ (invariance of Haar measure), in
particular under $\rho(Z(E))$. Put $\Lambda'=\Lambda-\Lambda_0$ and $\ell(u)=\Lambda'(\omega(k)(u\otimes
f_0))$ for $u\in\pzF_s$. Then $\Lambda'$ is invariant under $\rho(Z(E))$ and has mean zero:
$\int_{Z(\auA)/Z(E)}\Lambda'(\rho(x)F)\,dx=\Lambda_0(F)-\Lambda_0(F)=0$.

\emph{Rational translates are real translations.} There is $M\ge1$ such that $\rho(x_f)f_0=f_0$ for
every finite vector $x_f$ in $M$ times the integral points (finite data are locally constant under
finite translations, and the integral points form a compact group). There is $R$ such that every
adele is congruent modulo $E$ to $M$ times an adele with integral finite part and archimedean
components of size $\le R$ (a compact fundamental domain for $E\backslash\auA_E$). Given $x=\sum
c_iz_i\in Z(\auA)$, write $c_i+a_i=Mw_i$, $a_i\in E$, and $\gamma=\sum a_iz_i\in Z(E)$; then
$k^{-1}(\gamma+x)=M\sum w_ik^{-1}z_i=:M\mathsf q$ has integral finite part and archimedean coordinates of
size $\le M\#\iota R\delta=:C_0\delta$. As $\rho$ is additive on the isotropic space $k^{-1}Z(\auA)$ and
$\rho(M\mathsf q)=\rho_\infty(M\mathsf q_\infty)\otimes\rho_f(M\mathsf q_f)$ with $\rho_f(M\mathsf q_f)f_0=f_0$, we get, for every
$u\in\pzF_s$, $\rho(k^{-1}(\gamma+x))(u\otimes f_0)=(\rho_\infty(y_x)u)\otimes f_0$ with
$y_x=M\mathsf q_\infty$, $|y_x|\le C_0\delta$. Let $\Lambda_1$ be $\Lambda$ or $\Lambda_0$ (both are invariant under
$\rho(Z(E))$). Since $\rho(\gamma)\rho(x)=\rho(\gamma+x)$ and $\rho(\gamma+x)\omega(k)=\omega(k)
\rho(k^{-1}(\gamma+x))$ by (W1),
\begin{equation}\label{pcx:eq-meantrans}
\Lambda_1\bigl(\rho(x)\omega(k)(u\otimes f_0)\bigr)=\Lambda_1\bigl(\omega(k)((\psrho_\infty(y_x)u)\otimes
f_0)\bigr)\qquad(x\in Z(\auA),\ u\in\pzF_s).
\end{equation}

\emph{The norm of $\ell$.} Take $u=\psrho_\infty(z)\Psi$ in \eqref{pcx:eq-meantrans} with $\Lambda_1=\Lambda$;
as $\psrho_\infty(y_x)\psrho_\infty(z)$ is a scalar of absolute value $1$ times $\psrho_\infty(y_x+z)$,
$|\Lambda(\rho(x)\omega(k)(\psrho_\infty(z)\Psi\otimes f_0))|\le\|\Lambda\circ\omega(k)\|_{f_0}\|\Psi\|_s$
for all $x$; integrating over $x$, the same bound holds for $\Lambda_0$. Hence, in the notation of
Lemma~\ref{pcx:diffgain} for the Fock space of $X\otimes\R$, $\|\ell\|^*\le2\|\Lambda\circ\omega(k)\|_{f_0}$.

\emph{The hypothesis of Lemma~\ref{pcx:diffgain} with $t=C_0\delta$.} Let $u\in\pzF_s$ and
$F=\omega(k)(u\otimes f_0)$. By the mean zero property, $\ell(u)=\int_{Z(\auA)/Z(E)}(\ell(u)-
\Lambda'(\rho(x)F))\,dx$, and by \eqref{pcx:eq-meantrans} for $\Lambda$ and for $\Lambda_0$,
$\Lambda'(\rho(x)F)=\ell(\psrho_\infty(y_x)u)$ with $|y_x|\le t$. So if $|\ell(u)-\ell(\psrho_\infty(y)u)|
\le B$ for all $|y|\le t$, then $|\ell(u)|\le B$.

\emph{Conclusion.} Lemma~\ref{pcx:diffgain} with $N=l$ gives, for $\|v\|_{\max(l,2s)}<\infty$,
$|\ell(v)|\le C_l\|\ell\|^*(C_0\delta)^l\|v\|_l\le2C_lC_0^l\,\|\Lambda\circ\omega(k)\|_{f_0}\,\delta^l\,
\|v\|_l$.
\end{proof}

\begin{lemma}[The mean kills data supported off the orthogonal]\label{pcx:meanzero}
Let $\Lambda_0$ be a bounded functional on the data of $X=\auHom(S_1,S_2)$ invariant under $\rho(Z(\auA))$
for an isotropic $Z\subseteq X(E)$, and $\Phi$ locally constant compactly supported at $P$. If for
every $x$ with $\Phi(x)\ne0$ there is $\zeta\in Z(\auA)$ whose finite part at $\bar P$ pairs
non-trivially with $x$ (in the pairing of \eqref{pl:eq-heis}), then $\Lambda_0([y,\Phi]_P\otimes f)=0$.
\end{lemma}
\begin{proof}
Here \emph{bounded} is used only through the following: $\Lambda_0$ is linear and
$\Lambda_0(\rho(\zeta)F)=\Lambda_0(F)$ for $\zeta\in Z(\auA)$ and all data $F=[y,\Phi_i]_P\otimes f$
with $\Phi_i$ the indicator of a coset (times $\Phi$). Since $Z(\auA)=Z\otimes_E\auA_E$, with
$\zeta\in Z(\auA)$ also every $a\zeta$, $a\in\auA_E$, lies in $Z(\auA)$, and we may take $a$
supported at the places of $\bar P$ only, so that $a\zeta$ is an element $a'$ of $X(\auA)$ supported
at $\bar P$. By \eqref{pl:eq-heis}, $\rho(a')$ acts on functions of $x$ at $P$ by multiplication by
$\lambda(\mathrm{tr}({}^ta'x))$, where $\lambda=\prod\lambda_w$ with $\lambda_w$ trivial on $\auO_w$
and non-trivial on $\varpi_w^{-1}\auO_w$. Fix $x_0$ with $\Phi(x_0)\ne0$ and $\zeta$ as in the
hypothesis, so $\mathrm{tr}({}^t\zeta x_0)\ne0$ at some place $w$ of $\bar P$. Choosing $a=t$ at $w$
(and $0$ elsewhere) with $t\in E_w$ such that $\lambda_w(t\,\mathrm{tr}({}^t\zeta x_0))\ne1$, the
character $x\mapsto\lambda(\mathrm{tr}({}^t(a\zeta)x))$ is constant, with value $c\ne1$, on a small
coset $x_0+\varpi^{n}\pcMat(O_P)$ (it is trivial on $\varpi^n\pcMat(O_P)$ for $n$ large). Then
$\rho(a\zeta)$ multiplies the datum with the indicator of that coset by $c$, and invariance gives
$\Lambda_0=c\Lambda_0$ on it, so $\Lambda_0$ kills it. By compactness finitely many such cosets cover
the support of $\Phi$ (shrinking them to be disjoint and refining so that $\Phi$ is constant on each);
by linearity $\Lambda_0([y,\Phi]_P\otimes f)=0$.
\end{proof}

\begin{proof}[Proof of Proposition~\ref{pcx:decay}]
(a) Let $A\ge2s$ be as in (R2) for the finite datum $f_0=\Phi\otimes\psvac$ (so $F=[y,\Phi]_P=
y\otimes f_0$), put $l=N_0+A\ge2s$, and take $M=l$; let $y$ have $\|y\|_M<\infty$, so that $y\in\pzF_s$ and
$\|y\|_s\le c_1\|y\|_{2s}\le c_1\|y\|_M$ (Lemma~\ref{pcx:superpos}(b)). If $\mathrm{ht}(h)\le2$, (R2) (with
$z=0$) gives $|\theta(\omega(k)F)|\le\|\theta\circ\omega(k)\|_{f_0}\|y\|_s\le Yc_12^A\|y\|_M\,
\mathrm{ht}(g)^A$, and $\mathrm{ht}(h)^{N_0}\le2^{N_0}$; this is the claim. Fix $g,h$ with
$\mathrm{ht}(h)>2$. Let $d\in\Lambda_h$ be
isotropic with $|d|_h\le2/\mathrm{ht}(h)$, and $c_1,\dots,c_{n_1}$ a basis of $S_1(E)$ in $\Lambda_g$
with $|c_j|_g\le Z\mathrm{ht}(g)^{n_1}$ (R1). The rank-one maps $z_j=d\langle c_j,\cdot\rangle\in X(E)$
lie in the lattice at $k=(g,h)$, have length $\le|d|_h|c_j|_g\le2Z\mathrm{ht}(g)^{n_1}/\mathrm{ht}(h)
=:\delta$, and span the isotropic subspace $Z=\{d\lambda:\lambda\in S_1(E)^*\}$ ($\langle d\lambda,
d\lambda'\rangle$ is a multiple of $\langle d,d\rangle=0$). $\theta$ is invariant under $\rho(Z(E))$.
Lemma~\ref{pcx:mean} with $\Lambda=\theta$ and $l$ as above (its hypotheses hold: $\theta$ is invariant
under $\rho(Z(E))$, $x\mapsto\theta(\rho(x)F')$ is continuous for $F'$ of class $s$ by
Lemmas~\ref{pz:growth}(c) and~\ref{pcx:superpos}(a) and the local constancy of finite data,
$\|\theta\circ\omega(k)\|_{f_0}<\infty$ by (R2), and $\|y\|_{\max(l,2s)}=\|y\|_l<\infty$) gives
$|\theta(\omega(k)F)-\Lambda_0(\omega(k)F)|\le C\|\theta\circ\omega(k)\|_{f_0}\delta^l\|y\|_l$. Now
$\Lambda_0(\omega(k)F)=\Lambda_0'(F)$ with $\Lambda_0'=\Lambda_0\circ\omega(k)$ invariant under
$\rho(k^{-1}Z(\auA))$, and the pairing condition of Lemma~\ref{pcx:meanzero} holds: in coordinates at
$k$, the elements of $k^{-1}Z(\auA)$ are $d'\,\lambda'$ with $d'=h^{-1}d\ne0$ and $\lambda'$ running over
all of the dual (because $g^{-1}c$ is an adelic basis), and the pairing of $x$ with $\zeta=d'\lambda'$ is
$\mathrm{tr}({}^t\zeta\,x)=\lambda'({}^td'x)$ (with ${}^td'x$ the row vector of the linear form
$\lambda'$ is evaluated on), which is non-zero for some $\lambda'$ iff ${}^td'x\ne0$; this holds
because $x$ is surjective at $P$ and the component of $d'$ at $\bar P$ is non-zero ($d$ is a
non-zero rational vector). So $\Lambda_0(\omega(k)F)=0$. With (R2),
$\|\theta\circ\omega(k)\|_{f_0}\le Y(\mathrm{ht}(g)\mathrm{ht}(h))^A$, so
\[
|\theta(\omega(k)F)|\le C'\|y\|_l\,\mathrm{ht}(g)^{A+n_1l}\,\mathrm{ht}(h)^{A-l},
\]
which is (a) since $l-A=N_0$. (b) is the same with the roles of $S_1$, $S_2$ exchanged: $d$ a short
isotropic vector of $S_1$ at $g$, $c$ a short basis of $S_2$ at $h$, and the condition $x\,d'\ne0$ from
injectivity.
\end{proof}


\begin{corollary}[Absolute convergence in the exchange identity]\label{pcx:conv}
Assume Proposition~\ref{pcx:decay}. \textup{(i)} If $\Phi$ on $\pcMat_{r\times N}(F_S)$ is supported
on matrices of rank $r$ (at every place of $S$) then $h\mapsto\theta^{W,V}_{\alpha,\beta}(g,h;[v,\Phi]_S)$
is integrable on $[\auU(V)]$, uniformly for $g$ in compact sets, and the integral has moderate
growth in $g$. \textup{(ii)} The same holds for $h'\mapsto\theta^{V',W}_{\alpha,\beta}(h',g;[v,\Psi]_S)$
when $\Psi$ on $\pcMat_{N\times b}(F_S)$ is supported on matrices of rank $b$. \textup{(iii)} In
particular, for $\Phi,\Psi$ matched for $\chi$ (over $\pcR_\pcm$, or over $\auO_S$), both integrals of
\eqref{pcx:eq-X} converge absolutely.
\end{corollary}
\begin{proof}
(i) Proposition~\ref{pcx:decay}(a) with $N_0$ large and the finiteness of $\int_{[\auU(V)]}
\mathrm{ht}(h)^{-N_0}dh$ for large $N_0$ (reduction theory; for $V$ anisotropic $[\auU(V)]$ is
compact). (ii) is (b) for the pair $(V',W)$, the integration being over the source group. (iii)
Matched functions are supported on $\{k^{[a]}\}$, resp.\ $\{k_{[b]}\}$, $k$ invertible, i.e.\ on
matrices of full rank $a$, resp.\ $b$ (Lemma~\ref{pcx:mod} reduces data modulo $\pcm$ to this case).
\end{proof}

\begin{remark}[Which integrals are used downstream]\label{pcx:rem-conv}
In all applications (Theorem~\ref{pc:tri}(b) with $e=t$, Theorem~\ref{pc:tri}(c) with $e=0$,
Proposition~\ref{po:exchange}, Theorem~\ref{pc:auxcore}) the theta integrals of
Theorem~\ref{pcx:X} are formed with matched finite data, or with point data $\delta_\xi$ at
$\xi$ of full rank; Corollary~\ref{pcx:conv} covers all of them, and no other integrals occur. So
Theorem~\ref{pcx:X} holds with absolutely convergent integrals, as stated in Theorem~\ref{pc:X}.
\end{remark}

\subsection{From standard data to matched functions: proof of Proposition~\ref{pcx:P-match}}
\label{pcx:ss-match}

Throughout, $F=F_S$, $O=\auO_S$, $\varpi\in O$ is the product of uniformisers of the factors, and
Haar measures on $\auGL_m(F)$ give $\auGL_m(O)$ volume $1$. A matrix $y\in\pcMat_{N\times b}(F)$ is
\emph{primitive} if it is integral and has an integral left inverse; equivalently (each factor of
$O$ being a discrete valuation ring) $y$ is integral and its reduction modulo every prime of $S$ has
rank $b$, equivalently $y=k_{[b]}$ for some $k\in\auGL_N(O)$. Write $\mathrm{Pr}_b\subset
\pcMat_{N\times b}(F)$ for the primitive matrices and $\mathrm{Pr}^a=\{x\in\pcMat_{a\times N}(F):
x^{\mathsf t}\text{ primitive}\}=\{k^{[a]}:k\in\auGL_N(O)\}$ (apply the first description to
$k^{-1}$). Both sets are open and compact. $\mathrm{SB}(X)$ denotes the locally constant compactly
supported functions with support in $X$; $\mathrm{Inj}_b$ is the set of matrices in
$\pcMat_{N\times b}(F)$ of rank $b$, $\mathrm{Surj}^a$ that of matrices in $\pcMat_{a\times N}(F)$ of
rank $a$ (at every factor).

Let $\xi(t)=\chi_S(t)|t|^{N/2}$ on $F^\times$; then $\xi=\chi_S$ on $O^\times$.

\begin{lemma}[Cosets]\label{pcx:cosets}
\textup{(a)} For $y\in\mathrm{Inj}_b$ the set of $u\in\auGL_b(F)$ with $yu$ primitive is one coset
$u_0\auGL_b(O)$. \textup{(b)} For $x\in\mathrm{Surj}^a$ the set of $u\in\auGL_a(F)$ with $ux\in
\mathrm{Pr}^a$ is one coset $\auGL_a(O)u_0$. \textup{(c)} If $y$ is primitive and $yg$ is primitive,
$g\in\auGL_b(F)$, then $g\in\auGL_b(O)$.
\end{lemma}
\begin{proof}
(a) $\Lambda=\{c\in F^b:yc\in O^N\}$ is an $O$-lattice in $F^b$ ($y$ is injective with a left
inverse $L$, so $\Lambda\subseteq L(O^N)$ is bounded, and it contains a small multiple of $O^b$). The
matrix $yu$ is integral iff $uO^b\subseteq\Lambda$, and then primitive iff $yuO^b=yF^b\cap O^N$, i.e.\
iff $uO^b=\Lambda$; the $u$ with $uO^b=\Lambda$ form one coset of $\auGL_b(O)$. (b) is (a) for
$x^{\mathsf t}$: $ux\in\mathrm{Pr}^a$ iff $x^{\mathsf t}u^{\mathsf t}$ is primitive, so the
$u^{\mathsf t}$ form a coset $u_0'\auGL_a(O)$ and the $u$ form $\auGL_a(O)u_0'^{\mathsf t}$ (a left
coset, as it must be, $\mathrm{Pr}^a$ being stable under left multiplication by $\auGL_a(O)$). (c) is (a) applied with $y$, since $u=1$ and $u=g$ both qualify.
\end{proof}

\begin{lemma}[Right invariance of matched functions]\label{pcx:rightinv}
Let $\Phi,\Psi$ be matched for $\chi_S$ over $O$. Then $\Psi(ym)=\chi_S(\det m)^{-1}\Psi(y)$ and
$\Phi(m'x)=\chi_S(\det m')^{-1}\Phi(x)$ for $m\in\auGL_b(O)$, $m'\in\auGL_a(O)$. Moreover, if $j\in
\auGL_N(F)$ is such that $j_{[b]}$ is primitive and $j^{[a]}\in\mathrm{Pr}^a$, then
$\Phi(j^{[a]})=\xi(\det j)\Psi(j_{[b]})$.
\end{lemma}
\begin{proof}
For $k\in\auGL_N(O)$ let $D=\mathrm{diag}(m,1_a)$; then $(kD)_{[b]}=k_{[b]}m$, $(kD)^{[a]}=k^{[a]}$,
$\det kD=\det k\det m$, and the matching relation for $k$ and $kD$ gives the first claim on the
support $\{k_{[b]}\}$, both sides vanishing elsewhere (the support is stable under right
multiplication by $\auGL_b(O)$). The second claim is the same with $D=\mathrm{diag}(1_b,m'^{-1})$.
For the last claim let $i\in\auGL_N(O)$ with $i^{[a]}=j^{[a]}$. The columns of $i_{[b]}$ and of
$j_{[b]}$ both form a basis of the kernel of $j^{[a]}$, so $j_{[b]}g=i_{[b]}$ with $g\in\auGL_b(F)$,
and $g\in\auGL_b(O)$ by Lemma~\ref{pcx:cosets}(c). The matrix $j^{-1}i$ has last rows $j^{[a]}i=
i^{[a]}i=(0\ 1_a)$ and first columns $j^{-1}j_{[b]}g=\binom g0$, so $\det i=\det j\det g$. Hence
$\Phi(j^{[a]})=\Phi(i^{[a]})=\chi_S(\det i)\Psi(j_{[b]}g)=\xi(\det j)\xi(\det g)\chi_S(\det g)^{-1}
\Psi(j_{[b]})=\xi(\det j)\Psi(j_{[b]})$, using $\xi=\chi_S$ on units and $|\det i|=1$.
\end{proof}

\begin{definition}[Sections and same-section pairs]\label{pcx:def-section}
For $\Phi\in\mathrm{SB}(\mathrm{Surj}^a)$, $x\in\mathrm{Surj}^a$, and $\Psi\in\mathrm{SB}(\mathrm{Inj}_b)$,
$y\in\mathrm{Inj}_b$, put
\[
\mathsf S_\Phi(x)=\int_{\auGL_a(F)}\xi(\det r)\Phi(rx)\,dr,\qquad
\mathsf T_\Psi(y)=\int_{\auGL_b(F)}\xi(\det r)\Psi(yr)\,dr
\]
(the integrands have compact support, since $r\mapsto rx$ and $r\mapsto yr$ are proper). A pair
$(\Phi,\Psi)\in\mathrm{SB}(\mathrm{Surj}^a)\times\mathrm{SB}(\mathrm{Inj}_b)$ is a
\emph{same-section pair} if $\mathsf S_\Phi(k^{[a]})=\xi(\det k)\,\mathsf T_\Psi(k_{[b]})$ for all
$k\in\auGL_N(F)$. They form a vector space $\mathcal P$; let $\mathcal M$ be its image under
$(\Phi,\Psi)\mapsto\Psi$.
\end{definition}

\begin{lemma}\label{pcx:section-laws}
\textup{(a)} If $(\Phi,\Psi)\in\mathcal P$ and $g\in\auGL_N(F)$ then $(\xi(\det g)\Phi(\cdot\,g),
\Psi(g^{-1}\cdot))\in\mathcal P$. In particular $\mathcal M$ is stable under $\Psi\mapsto\Psi(g\,\cdot)$.
\textup{(b)} If $\Phi,\Psi$ are matched for $\chi_S$ over $O$ and $\Psi$ is locally constant, then
$(\Phi,\Psi)\in\mathcal P$.
\end{lemma}
\begin{proof}
(a) The support conditions are preserved. Since $k^{[a]}g=(g^{-1}k)^{[a]}$ and $g^{-1}k_{[b]}=
(g^{-1}k)_{[b]}$,
$\mathsf S_{\xi(\det g)\Phi(\cdot g)}(k^{[a]})=\xi(\det g)\mathsf S_\Phi((g^{-1}k)^{[a]})=
\xi(\det g)\xi(\det g^{-1}k)\mathsf T_\Psi((g^{-1}k)_{[b]})=\xi(\det k)\mathsf T_{\Psi(g^{-1}\cdot)}
(k_{[b]})$.

(b) $\Psi$ is supported on $\mathrm{Pr}_b$, compact open, so $\Psi\in\mathrm{SB}(\mathrm{Inj}_b)$.
$\Phi$ is supported on $\mathrm{Pr}^a$; it is locally constant because the map $\auGL_N(O)\to
\mathrm{Pr}^a$, $k\mapsto k^{[a]}$, is continuous, open and surjective, and $\Phi(k^{[a]})=
\chi_S(\det k)\Psi(k_{[b]})$ is a locally constant function of $k$; so
$\Phi\in\mathrm{SB}(\mathrm{Surj}^a)$. Let $k\in\auGL_N(F)$. By Lemma~\ref{pcx:cosets} choose $u$ with
$k_{[b]}u$ primitive and $u'$ with $u'k^{[a]}\in\mathrm{Pr}^a$. By Lemma~\ref{pcx:rightinv} the
integrand of $\mathsf T_\Psi(k_{[b]})$ vanishes off $u\auGL_b(O)$ and is constant on it, so
$\mathsf T_\Psi(k_{[b]})=\xi(\det u)\Psi(k_{[b]}u)$; likewise the integrand of $\mathsf S_\Phi(k^{[a]})$ vanishes off the left coset
$\auGL_a(O)u'$ and is constant on it (by $\Phi(m'x)=\chi_S(\det m')^{-1}\Phi(x)$), so
$\mathsf S_\Phi(k^{[a]})=\xi(\det u')\Phi(u'k^{[a]})$. Let $j=k\,\mathrm{diag}(u,u'^{-1})$: $j_{[b]}=k_{[b]}u$, $j^{[a]}=u'k^{[a]}$,
$\det j=\det k\det u/\det u'$. By the last part of Lemma~\ref{pcx:rightinv},
$\mathsf S_\Phi(k^{[a]})=\xi(\det u')\xi(\det j)\Psi(j_{[b]})=\xi(\det k)\mathsf T_\Psi(k_{[b]})$.
\end{proof}

\begin{lemma}[Linear factor]\label{pcx:linfactor}
Let $X$ be a vector space and $\Lambda\colon\mathrm{SB}(\mathrm{Surj}^a)\to X$ linear with
$\xi(\det g)\Lambda(\Phi(g\,\cdot))=\Lambda(\Phi)$ for all $g\in\auGL_a(F)$. There is a linear map
$\mathcal T\colon\mathrm{SB}(\mathrm{Inj}_b)\to X$ with $\Lambda(\Phi)=\mathcal T(\Psi)$ for all
$(\Phi,\Psi)\in\mathcal P$.
\end{lemma}
\begin{proof}
Let $\Phi\in\mathrm{SB}(\mathrm{Surj}^a)$. By Lemma~\ref{pcx:cosets}(b), $\mathrm{Surj}^a$ is the
disjoint union of the open sets $r\,\mathrm{Pr}^a$, $r$ running through representatives of
$\auGL_a(F)/\auGL_a(O)$; finitely many of them, $r_1,\dots,r_m$, cover the support of $\Phi$. Put
$\Phi_i=\Phi\cdot\mathbf 1_{r_i\mathrm{Pr}^a}\in\mathrm{SB}(\mathrm{Surj}^a)$. For $k\in\auGL_a(O)$ the
hypothesis gives $\Lambda(\Phi_i)=\xi(\det r_ik)\Lambda(\Phi_i(r_ik\,\cdot))$. The function
$k\mapsto\xi(\det r_ik)\Phi_i(r_ik\,\cdot)$ is locally constant on the compact group $\auGL_a(O)$ with
values in a finite-dimensional space, so averaging over $k$ is a finite sum and commutes with
$\Lambda$:
\[
\Lambda(\Phi)=\sum_i\Lambda(\Phi_i)=\Lambda(\Phi^0),\qquad\Phi^0(x)=\sum_i\int_{\auGL_a(O)}
\xi(\det r_ik)\Phi_i(r_ikx)\,dk .
\]
For $x\in\mathrm{Pr}^a$ we have $r_ikx\in r_i\mathrm{Pr}^a$, so $\Phi_i(r_ikx)=\Phi(r_ikx)$, and
$\Phi(rx)=0$ unless $r$ lies in some $r_i\auGL_a(O)$; hence $\Phi^0(x)=\mathsf S_\Phi(x)$. For
$x\notin\mathrm{Pr}^a$, $r_ikx\notin r_i\mathrm{Pr}^a$ and $\Phi^0(x)=0$. So $\Lambda(\Phi)=
\Lambda(\mathbf 1_{\mathrm{Pr}^a}\mathsf S_\Phi)$. For each $x\in\mathrm{Pr}^a$ fix $k_x\in
\auGL_N(O)$ with $k_x^{[a]}=x$; let $\tilde\Lambda$ be a linear extension of $\Lambda$ to all
functions on $\pcMat_{a\times N}(F)$, and put $\mathcal T(\Psi)=\tilde\Lambda\bigl(x\mapsto\mathbf
1_{\mathrm{Pr}^a}(x)\,\xi(\det k_x)\mathsf T_\Psi((k_x)_{[b]})\bigr)$, which is linear in $\Psi$.
For $(\Phi,\Psi)\in\mathcal P$
the function is $\mathbf 1_{\mathrm{Pr}^a}\mathsf S_\Phi$, so $\mathcal T(\Psi)=\Lambda(\Phi)$.
\end{proof}

\begin{lemma}[Generation by one standard datum]\label{pcx:generate}
Let $\mathcal Q\subseteq\mathrm{SB}(\mathrm{Inj}_b)$ be a subspace stable under $\Psi\mapsto
\Psi(g\,\cdot)$ for $g\in\auGL_N(F)$, and suppose $\Psi^{\mathrm{col}}_s\in\mathcal Q$ for an integral
$s\in\pcMat_{b\times N}(O)$ with $sx_0=1$ for some integral $x_0$. Then every locally constant $\Psi$
having a matched $\Phi$ (for $\chi_S$ over $O$) lies in $\mathcal Q$.
\end{lemma}
\begin{proof}
Write $p=\Psi^{\mathrm{col}}_s$. As $\Psi$ is locally constant with compact support there is $c\ge0$
with $\Psi(\varpi^ct+y)=\Psi(y)$ for all integral $t$ and all $y$. Let $e=x_0s$, an idempotent, and
$d=\varpi^{-c}(1-e)+e\in\auGL_N(F)$ (inverse $\varpi^c(1-e)+e$). Note $sd=s$ and $dx_0=x_0$.

\emph{Claim: for $k\in\auGL_N(O)$ and all $y$, $\Psi(k^{-1}x_0)\,p(dky)=|p(dky)|\,\Psi(y)$.} If
$p(dky)=0$ both sides vanish. Otherwise $z=dky$ is integral and $g=sz=sky\in\auGL_b(O)$, and $p(z)=
\chi_S(\det g)^{-1}$. Since $z-x_0g=(d-e)ky=\varpi^{-c}(1-e)ky$, we have
$y=k^{-1}\bigl(\varpi^c(z-x_0g)+x_0g\bigr)$ with $t=k^{-1}(z-x_0g)$ integral, so
$\Psi(y)=\Psi(k^{-1}x_0g)=\chi_S(\det g)^{-1}\Psi(k^{-1}x_0)$ by Lemma~\ref{pcx:rightinv}. Thus
$\Psi(k^{-1}x_0)p(z)=\Psi(y)$, and $|p(z)|=1$.

By stability, $y\mapsto\Psi(k^{-1}x_0)p(dky)$ lies in $\mathcal Q$, hence so does $y\mapsto|p(dky)|
\Psi(y)$, for every $k\in\auGL_N(O)$; this function is locally constant in $k$ (uniformly in $y$), so
its average $I(y)\Psi(y)$, $I(y)=\int_{\auGL_N(O)}|p(dky)|\,dk$, is a finite combination and lies in
$\mathcal Q$. If $\Psi(y)\ne0$ then $y=l_{[b]}$ with $l\in\auGL_N(O)$, and $x_0=j_{[b]}$ with
$j\in\auGL_N(O)$ ($x_0$ is primitive since $sx_0=1$), so $y=lj^{-1}x_0$ and $I(y)=I(x_0)$ by right
invariance of $dk$. Finally $p(dx_0)=p(x_0)=\chi_S(\det sx_0)^{-1}=1$, so $I(x_0)>0$ by continuity.
Hence $I(x_0)\Psi\in\mathcal Q$ and $\Psi\in\mathcal Q$.
\end{proof}

\begin{lemma}[Action at $S$]\label{pcx:Sact}
Let $W$ (rank $N$, even), $V$ (rank $r$), $U$ (rank $b=N-r$) and $\eta(t)=\upsilon_w(t)|t|^{1/2}$ on
$F^\times$, where at $w\in S$ the local component $\upsilon_w$ of $\upsilon$ is evaluated on
$t_w\in E_w^\times$. For $g\in\auGL_N(F)$, $h\in\auGL_r(F)$, read via the frames as elements of
$\auU(W)(\auA)$, $\auU(V)(\auA)$ supported on $S\cup\bar S$, and for $\Phi\in\mathrm{SB}(\mathrm{Surj}^r)$,
$\Psi\in\mathrm{SB}(\mathrm{Inj}_b)$:
\begin{enumerate}
\item[(a)] $\Theta^{\mathrm L}_{W,V}(v,\xi(\det h)\Phi(h\,\cdot))=\Theta^{\mathrm L}_{W,V}(v,\Phi)$;
\item[(b)] $\Theta^{\mathrm L}_{W,V}(v,\xi(\det g)^{-1}\Phi(\cdot\,g^{-1}))(x g)=\eta(\det g)^{-b}
\Theta^{\mathrm L}_{W,V}(v,\Phi)(x)$;
\item[(c)] $\Theta^{\mathrm S}_{W,U}(y,\Psi(g\,\cdot))(xg)=\eta(\det g)^{-b}\Theta^{\mathrm S}_{W,U}(y,
\Psi)(x)$.
\end{enumerate}
\end{lemma}
\begin{proof}
By Corollary~\ref{pcx:conv} all integrals converge absolutely. All scalars are computed from one
formula. Let $v=w\aubar w$ with $w\in S$. By Proposition~\ref{pwf:W5}, in the model at $w$ of
$\auHom(S_1,S_2)$ (ranks $n_1,n_2$; source group element $g$, target group element $h$) the untwisted
operator is
\[
\Psi\mapsto\upsilon_w(\det h_w)^{-n_1}\upsilon_w(\det g_w)^{n_2}|\det h_w|_w^{-n_1/2}
|\det g_w|_w^{n_2/2}\,\Psi(h_w^{-1}xg_w).
\]
At $v$, $\auU(1)(E^+_v)=\{(t,t^{-1})\}\cong E_w^\times$ through the $w$-component, and, as in the
proof of Proposition~\ref{pwf:W5}, one has
$\upsilon(t,t^{-1})=\upsilon_w(t)^2$ and $\nu(t,t^{-1})=\chi_S(t)$; in particular, for $z\in
E_w^\times$, $\upsilon_w(z)^{n}=\upsilon(z)^{n/2}$ for even $n$, where $\upsilon(z)$ means
$\upsilon|_{\auU(1)}$ at $(z,z^{-1})$. Thus $\gamma(z)=\chi_S(z)^{-1}\upsilon_w(z)^{N}$, and the
untwisted operator equals $\upsilon(\det h)^{-n_1/2}\upsilon(\det g)^{n_2/2}|\det h|^{-n_1/2}
|\det g|^{n_2/2}\Psi(h^{-1}xg)$ when $n_1,n_2$ are even, and in general with $\upsilon_w^{-n_1}$,
$\upsilon_w^{n_2}$ in place of $\upsilon^{-n_1/2}$, $\upsilon^{n_2/2}$.

(a) For $\auHom(W,V)$, $n_1=N$, $n_2=r$. The operator $\omega_{\gamma,\gamma}(1,h^{-1})$ is
$\gamma(\det h)^{-1}$ times the untwisted operator for $h^{-1}$, i.e.\ it sends $\Phi$ to
\[
\chi_S(\det h)\upsilon_w(\det h)^{-N}\cdot\upsilon_w(\det h)^{N}|\det h|^{N/2}\,\Phi(h\,\cdot)
=\xi(\det h)\Phi(h\,\cdot):
\]
the power $\upsilon_w^{N}$ of the model cancels the one carried by $\gamma=\nu^{-1}\upsilon^{N/2}$.
So $\omega_{\gamma,\gamma}(1,h^{-1})[v,\Phi]_S=[v,\xi(\det h)\Phi(h\,\cdot)]_S$, using
\eqref{pcx:eq-twist} for the twist and the invariance of the infinite and unramified components
under the element $h$, which is supported on $S\cup\bar S$. Since
$\int_{[\auU(V)]}\theta(g,h h_0)\,dh=\int_{[\auU(V)]}\theta(g,h)\,dh$, (a) follows.

(b) $\Theta^{\mathrm L}(v,\Phi)(xg)$ is $\Theta^{\mathrm L}$ of the datum translated by
$\omega_{\gamma,\gamma}(g,1)$, evaluated at $x$, and $\omega_{\gamma,\gamma}(g,1)$ sends $\Phi$ to
\[
\gamma(\det g)^{-1}\upsilon_w(\det g)^{r}|\det g|^{r/2}\Phi(\cdot\,g)
=\chi_S(\det g)\upsilon_w(\det g)^{r-N}|\det g|^{r/2}\Phi(\cdot\,g).
\]
Applied to $\xi(\det g)^{-1}\Phi(\cdot\,g^{-1})$, with $\xi(\det g)^{-1}=\chi_S(\det g)^{-1}
|\det g|^{-N/2}$, this gives $\upsilon_w(\det g)^{-b}|\det g|^{-b/2}\Phi=\eta(\det g)^{-b}\Phi$,
since $r-N=-b$.

(c) For $\Theta^{\mathrm S}_{W,U}$ the characters are $(\upsilon^{b},\nu\upsilon^{N/2})$ and the
datum $\Psi^{\mathsf t}$ sits at $\bar w$, in the model at $\bar w$ of $\auHom(W,U)$ ($n_1=N$,
$n_2=b$). Apply the formula above with the roles of $w,\bar w$ exchanged. The component of $g$ at
$\bar w$ is $g_{\bar w}=\bar g^{-\mathsf t}$, so, in the identification of $E_{\bar w}$ with $E_w$ by
conjugation, $\det g_{\bar w}=(\det g)^{-1}$, and the element $\det g$ of $\auU(1)(E^+_v)$ is
$((\det g),(\det g)^{-1})$ with $\bar w$-component $(\det g)^{-1}$. The untwisted operator therefore
multiplies by $\upsilon_{\bar w}(\det g_{\bar w})^{b}|\det g_{\bar w}|^{b/2}=\upsilon_{\bar w}
((\det g)^{-1})^{b}|\det g|^{-b/2}=\upsilon_w(\det g)^{b}|\det g|^{-b/2}$, because
$\upsilon_{\bar w}=\upsilon_w^{-1}$ on $E^{+\times}_v$ (this holds for $b$ of either parity; reading
$(t,t^{-1})$ through its $\bar w$-component $t^{-1}$ gives the same value $\upsilon_{\bar
w}(t^{-1})^2=\upsilon_w(t)^2$ of $\upsilon|_{\auU(1)}$), and acts on $\Psi^{\mathsf t}$
through $X\mapsto X g_{\bar w}$, i.e.\ on $\Psi$ through $y\mapsto g^{-1}y$ up to the
identification by conjugation. The source character contributes $\upsilon^{b}(\det g)^{-1}=
\upsilon_w(\det g)^{-2b}$. Hence $\Theta^{\mathrm S}(y,\Psi)(xg)=\upsilon_w(\det g)^{-b}
|\det g|^{-b/2}\,\Theta^{\mathrm S}(y,\Psi(g^{-1}\,\cdot))(x)$; replacing $\Psi$ by $\Psi(g\,\cdot)$
gives (c), with the scalar $\eta(\det g)^{-b}$, the same as in (b).
\end{proof}

\begin{proof}[Proof of Proposition~\ref{pcx:P-match}]
Let $\Lambda(\Phi)=\Theta^{\mathrm L}_{W,V}(v,\Phi)$ on $\mathrm{SB}(\mathrm{Surj}^r)$ (convergent by
Corollary~\ref{pcx:conv}); by Lemma~\ref{pcx:Sact}(a) it satisfies the hypothesis of
Lemma~\ref{pcx:linfactor}, which gives $\mathcal T$. Put $\mathcal D(\Psi)=\mathcal T(\Psi)-c\,
\Theta^{\mathrm S}_{W,U}(y,\Psi)$ and $\mathcal Q=\{\Psi\in\mathcal M:\mathcal D(\Psi)=0\}$.

$\mathcal Q$ is stable under $\Psi\mapsto\Psi(g\,\cdot)$: if $(\Phi,\Psi)\in\mathcal P$ then
$(\Phi',\Psi(g\,\cdot))\in\mathcal P$ with $\Phi'=\xi(\det g)^{-1}\Phi(\cdot\,g^{-1})$
(Lemma~\ref{pcx:section-laws}(a) for $g^{-1}$), and by Lemma~\ref{pcx:Sact}(b),(c),
$\mathcal D(\Psi(g\,\cdot))(xg)=\eta(\det g)^{-b}\mathcal D(\Psi)(x)$ (the two parts of
$\mathcal D$ are multiplied by the same scalar: in (b) the twist $\gamma(\det g)^{-1}=
\chi_S(\det g)\upsilon_w(\det g)^{-N}$ and the model factor $\upsilon_w(\det g)^{r}|\det g|^{r/2}$
combine with $\xi(\det g)^{-1}$ to $\upsilon_w(\det g)^{-b}|\det g|^{-b/2}$, and in (c) the twist
$\upsilon^b(\det g)^{-1}=\upsilon_w(\det g)^{-2b}$ and the model factor
$\upsilon_w(\det g)^{b}|\det g|^{-b/2}$ at $\bar w$ give the same product), which vanishes if
$\mathcal D(\Psi)=0$.

Take $s$ the first $b$ rows of $1_N$ and $x_0$ the first $b$ columns. By Lemma~\ref{pcx:std-matched}
and Lemma~\ref{pcx:section-laws}(b), $(\Phi^{\mathrm{row}}_s,\Psi^{\mathrm{col}}_s)\in\mathcal P$, so
$\mathcal D(\Psi^{\mathrm{col}}_s)=\Theta^{\mathrm L}(v,\Phi^{\mathrm{row}}_s)-c\Theta^{\mathrm S}(y,
\Psi^{\mathrm{col}}_s)=0$ by hypothesis, i.e.\ $\Psi^{\mathrm{col}}_s\in\mathcal Q$. By
Lemma~\ref{pcx:generate}, every locally constant $\Psi$ with a matched $\Phi$ lies in $\mathcal Q$;
then $(\Phi,\Psi)\in\mathcal P$ (Lemma~\ref{pcx:section-laws}(b)), $\Psi\in\mathrm{SB}(\mathrm{Inj}_b)$,
and $\Theta^{\mathrm L}(v,\Phi)=\mathcal T(\Psi)=c\Theta^{\mathrm S}(y,\Psi)$.
\end{proof}

\subsection{Passing to adjoints: proof of Proposition~\ref{pcx:P-flip}}\label{pcx:ss-flip}

Let $U^-$ be $U$ with the negated form; $\auU(U^-)=\auU(U)$ as groups, with the same rational points,
the same lattice $L_U$ and hence the same Haar measure, so $[\auU(U^-)]=[\auU(U)]$ as measure spaces,
and $\det$ is the same. Proposition~\ref{pf:flip} (fact (W6)) gives, for the pair $(U,W)$ of ranks
$b,N$, an isomorphism $\mathrm{Fl}\colon\psHS(\auHom(U,W))\to\psHS(\auHom(W,U^-))$ induced by $w\mapsto
w^\dagger$, a constant $c_\theta\ne0$ with $\theta'\circ\mathrm{Fl}=c_\theta\theta$, and
\[
\mathrm{Fl}\circ\omega_{\auHom(U,W)}(h,g)=\upsilon(\det h)^{N}\upsilon(\det g)^{-b}\,
\omega_{\auHom(W,U^-)}(g,h)\circ\mathrm{Fl}\qquad(h\in\auU(U)(\auA),\ g\in\auU(W)(\auA)),
\]
(the factor is $\upsilon(\det\nolimits_{\auHom(U,W)}(h,g))^{-1}$, the determinant of $(h,g)$ on the
$E$-module $\auHom(U,W)$ being $\det(g)^b\det(h)^{-N}$). We
normalise $\mathrm{Fl}$ as a tensor product of local isomorphisms: at the infinite places the map of
Fock spaces induced by $w\mapsto w^\dagger$ on real points, which defines $y\mapsto y^{\mathrm{ad}}$ and
preserves the moments $\|\cdot\|_M$ (it is induced by a measure-preserving isometry of the real
vector spaces); at the finite places the map induced by $w\mapsto w^\dagger$, normalised to be unitary.
By Lemma~\ref{pgx:dict}(f) the latter carries the self-dual lattice $\alpha\auHom(L_U,L_W)$ to
$u\,\alpha\auHom(L_W,L_U)$, $u=\aubar a/a$, a central idele of $\auU(U)=\auU(U^-)$ with $u\aubar u=1$ and
$u_v=1$ at the places not above $q_1$. Put $\mathrm{Fl}^\circ=\omega(1,u)^{-1}\circ\mathrm{Fl}$ ($u$
acting on $\auHom(W,U^-)$ through $\auU(U^-)$). Then: at finite places outside $S\cup\bar S$,
$\mathrm{Fl}^\circ$ sends the vacuum of the standard lattice $\alpha\auHom(L_U,L_W)$ to a multiple of
modulus one of the vacuum of $\alpha\auHom(L_W,L_U)$ (Lemma~\ref{pgx:dict}(e)); at $w\in S$, in the
models (W5) and the coordinates $\tilde x$ of Lemma~\ref{pgx:dict}, $\mathrm{Fl}^\circ$ is, up to a
constant of modulus one, the map $\Psi\mapsto(X\mapsto\Psi(\epsilon X^{\mathsf t}))$ from functions on
$\pcMat_{N\times b}(E_w)$ to functions on $\pcMat_{b\times N}(E_{\bar w})$ (identified by conjugation),
where $\epsilon=-1$: in the frames $w^\dagger=-\bar w^{\mathsf t}$ (the Gram matrices being unimodular
at $S$), the sign coming from the negated form; indeed by Lemma~\ref{pgx:dict}(f)
$\mathrm{Fl}$ is $\tilde X=-\nu^{-1}\,{}^t\tilde x$ with $\nu=\aubar{u_{\bar w}}^{-1}$, and
$\omega(1,u)^{-1}$ multiplies the variable $\tilde X$ by $\nu^{-1}$ ((W5): $h=u^{-1}$ acts by
$\Psi\mapsto\Psi(u_{\bar w}\,\cdot)$). The tensor product $\mathrm{Fl}$ of the local maps intertwines
the Heisenberg operators, so it is an admissible $\mathrm{Fl}$, and $c_\theta$ is a constant depending
only on $U$, $W$ and the frames. Since $u$ is central, Lemma~\ref{pgx:dict}(e) gives, for the theta
integrals over $[\auU(U)]$ used below,
$\int\theta'(\omega'(g,h)\mathrm{Fl}f)\,dh=\zeta_u\int\theta'(\omega'(g,h)\mathrm{Fl}^\circ f)\,dh$ with
$\zeta_u$ a value at $\det u$ of the target character, of modulus one; so in these integrals
$\mathrm{Fl}$ may be replaced by $\mathrm{Fl}^\circ$ at the cost of a constant of modulus one.

\begin{proof}[Proof of Proposition~\ref{pcx:P-flip}]
\emph{The datum.} By the normalisation, $\mathrm{Fl}^\circ[y,\Psi]_S=a_1[y^{\mathrm{ad}},\Psi_\epsilon^{\mathsf
t}]_{\bar S}$ with $\Psi_\epsilon(X)=\Psi(\epsilon X)$ and $|a_1|=1$. For $\Psi$ matched with some $\Phi$,
Lemma~\ref{pcx:rightinv} with $m=\epsilon 1_b$ gives $\Psi(\epsilon X)=\chi_S(\epsilon)^{-b}\Psi(X)$.
So $\mathrm{Fl}^\circ[y,\Psi]_S=a\,[y^{\mathrm{ad}},\Psi^{\mathsf t}]_{\bar S}$ with $a=a_1\chi_S(\epsilon)^{-b}$,
$|a|=1$.

\emph{The kernels.} By definition and the displayed law,
\[
\theta^{U,W}_{\gamma,\gamma}(h,g;f)=\gamma(\det h)^{-1}\gamma(\det g)\,\theta(\omega(h,g)f)
=c_\theta^{-1}\gamma(\det h)^{-1}\upsilon(\det h)^{N}\gamma(\det g)\upsilon(\det g)^{-b}\,
\theta'(\omega'(g,h)\mathrm{Fl}f).
\]
In the convention of \S\ref{pcx:ss-statement} for the pair $(W,U^-)$ this is
$c_\theta^{-1}\theta^{W,U^-}_{\gamma^{-1}\upsilon^b,\ \gamma^{-1}\upsilon^N}(g,h;\mathrm{Fl}f)$, and
$\gamma^{-1}\upsilon^N=\nu\upsilon^{N/2}$; by \eqref{pcx:eq-twist},
$\theta^{W,U^-}_{\gamma^{-1}\upsilon^b,\nu\upsilon^{N/2}}=\gamma(\det g)\,\theta^{W,U^-}_{\upsilon^b,
\nu\upsilon^{N/2}}$.

\emph{Integration.} With $f=[y,\Psi]_S$, integrating over $h\in[\auU(U)]=[\auU(U^-)]$ (absolutely
convergent by Corollary~\ref{pcx:conv}, $\Psi$ being supported on injective matrices), and
replacing $\mathrm{Fl}$ by $\mathrm{Fl}^\circ$ in the integral at the cost of the constant $\zeta_u$ of
modulus one (end of \S\ref{pcx:ss-flip}; absorbed in $a$),
\[
\int_{[\auU(U)]}\theta^{U,W}_{\gamma,\gamma}(h,g;[y,\Psi]_S)\,dh=c_\theta^{-1}a\,\gamma(\det g)\,
\Theta^{\mathrm S}_{W,U^-}(y^{\mathrm{ad}},\Psi)(g).
\]
So the proposition holds with $d=c_\theta a^{-1}$, $|d|=|c_\theta|=:d_0$, independent of $\nu$.
\end{proof}

\begin{remark}
One can normalise so that $|d|=1$; only the independence of $|d|$ from $\nu$ is needed, since constants independent of $\nu$ are absorbed in $C$.
\end{remark}

\subsection{Only the signs matter: proof of Proposition~\ref{pcx:P-signs}}\label{pcx:ss-signs}

\begin{lemma}[Change of space]\label{pcx:move}
Let $W,W_1$ have the same rank and signs, and $V,V_1$ likewise. There are $k\in\auU(W)(\auA)$ and,
for all characters $\alpha,\beta$, a number $u$ of modulus $1$ (independent of the data and of $g$)
such that for all data $[v,\Phi]_P$ and all $g$,
$\int_{[\auU(V_1)]}\theta^{W_1,V_1}_{\alpha,\beta}(g,h;[v,\Phi]_P)\,dh=u\int_{[\auU(V)]}
\theta^{W,V}_{\alpha,\beta}(gk,h;[v,\Phi]_P)\,dh$ (absolutely convergent integrals), where on the
right $g$ stands for its transport $\phi g\phi^{-1}$ by the isometry $\phi$ of the proof.
\end{lemma}
\begin{proof}
Apply Lemma~\ref{pcx:landherr} to $(W_1,W)$ and $(V_1,V)$, giving $\phi,k$ and $\phi',k'$. A datum is a
vector of $\psHS(\auHom)$ written in the frames, and the theta functional of $\auHom(W_1,V_1)$ is
that of $\auHom(W,V)$ transported by $(\phi,\phi')$ (it only depends on the rational points).
Transporting, the frame-datum $[v,\Phi]_P$ for $(W_1,V_1)$ becomes $\omega(k,k')$-applied to the same
frame-datum for $(W,V)$, up to the characters; hence
$\theta^{W_1,V_1}_{\alpha,\beta}(g,h;\cdot)=\alpha(\det k)\beta(\det k')^{-1}\theta^{W,V}_{\alpha,\beta}
(gk\,,hk';\cdot)$, with $g,h$ transported by $\phi,\phi'$ (indeed
$\theta(\omega(g,h)\omega(k,k')D)=\theta(\omega(gk,hk')D)$). Finally $h\mapsto hk'$ carries the quotient
by $\auU(V_1)(E^+)=\phi'^{-1}\auU(V)(E^+)\phi'$ measure-preservingly onto that by $\auU(V)(E^+)$
(conjugation by the transport, right translation by $k'$, both preserving Haar measure). So the
lemma holds with $u=\alpha(\det k)\beta(\det k')^{-1}$.
\end{proof}

\begin{lemma}[Equal rank: the theta integral is constant]\label{pcx:const}
Let $Y$ have even rank $2k$ and contain an adapted hyperbolic plane. For every Fock vector $v$ of
$(Y,Y)$ with finite moments and every tame $\nu$, the function
$g\mapsto\Theta^{\mathrm L}_{Y,Y}(v,\Psi^{\mathrm{col}}_1)(g)$ is constant on $\auU(Y)(\auA)$.
\end{lemma}
\begin{proof}
See \S\ref{pct:ss-const}: the hypotheses of Lemma~\ref{pct:SA} are verified there for
$D=\Theta^{\mathrm L}_{Y,Y}(v,\Psi^{\mathrm{col}}_1)$, the torus condition coming from
Theorem~\ref{pct:plane} (independence of $u$ in equal rank).
\end{proof}

\begin{proof}[Proof of Proposition~\ref{pcx:P-signs}]
By Lemma~\ref{pcx:move} with $(W,V)=(Y,Y)$, $(W_1,V_1)=(Y_1,Y_1)$, $\alpha=\beta=\gamma$, $P=S$,
$\Phi=\Psi^{\mathrm{col}}_1$: $\Theta^{\mathrm L}_{Y_1,Y_1}(v,\Phi)(1)=u\,\Theta^{\mathrm L}_{Y,Y}(v,\Phi)(k)
=u\,\Theta^{\mathrm L}_{Y,Y}(v,\Phi)(1)$ by Lemma~\ref{pcx:const}. Taking absolute values and dividing
by $L_*$ gives $\pcxstr_{Y_1}(\nu;v)=\pcxstr_Y(\nu;v)$.
\end{proof}

\subsection{The inputs of the exchange identity}\label{pcx:ss-status}

\emph{(A) Proved here} (modulo the published facts cited: Landherr's theorem, the Hasse norm
theorem, Jacobowitz, reduction theory (R1), Godement's form of Borel's reduction theorem and Dirichlet's unit theorem in
Lemma~\ref{pcx:reduce}, Hecke's
non-vanishing; (R2) is Lemma~\ref{pcx:thetagrowth}): the dictionary \eqref{pcx:eq-dict} and the twist law \eqref{pcx:eq-twist};
Lemma~\ref{pcx:Lpos}, Corollary~\ref{pcx:cne0}; Lemmas~\ref{pcx:adapted}, \ref{pcx:landherr},
\ref{pcx:std-matched}, \ref{pcx:mod}, \ref{pcx:reduce}, \ref{pcx:thetagrowth}, \ref{pcx:superpos}, \ref{pcx:diffgain}, \ref{pcx:mean},
\ref{pcx:meanzero}, \ref{pcx:cosets}, \ref{pcx:rightinv}, \ref{pcx:section-laws},
\ref{pcx:linfactor}, \ref{pcx:generate}, \ref{pcx:Sact}, \ref{pcx:move};
Proposition~\ref{pcx:decay} and Corollary~\ref{pcx:conv} (absolute convergence);
Propositions~\ref{pcx:P-match} and~\ref{pcx:P-flip} (the latter with Proposition~\ref{pf:flip});
the deductions
$\ref{pcx:P-std}+\ref{pcx:P-match}+\ref{pcx:P-flip}\Rightarrow\ref{pcx:P-matches}$,
$\ref{pcx:P-matches}+\ref{pcx:P-equal}\Rightarrow\ref{pcx:X}$,
$\ref{pcx:P-matches}+\ref{pcx:P-nonvan}+\ref{pcx:P-plane}+\ref{pcx:P-signs}\Rightarrow
\ref{pcx:P-equal}$ (Lemma~\ref{pcx:twice}); Proposition~\ref{pcx:P-signs}.

\emph{(B) Proved in the following sections:} Proposition~\ref{pcx:P-slice} in \S\ref{pcu:ss-slice};
Proposition~\ref{pcx:P-top} in \S\S\ref{pcu:ss-witt}--\ref{pcu:ss-std-data} (with (U1), (U2) in
\S\ref{pcu:ss-open}); the deduction of Proposition~\ref{pcx:P-std} in \S\ref{pcu:ss-std};
Lemma~\ref{pcx:const} and Propositions~\ref{pcx:P-descent}, \ref{pcx:P-plane}, \ref{pcx:P-nonvan} in
\S\ref{pct:sec}, the last with Lemma~\ref{pct:negcoeff}, proved in \S\ref{pcn:sec}. The full list of
dependencies of Theorem~\ref{pcx:X} is \S\ref{pcu:ss-final}.


%% ===== proofs/C-exchange-top.tex =====
% proofs/C-exchange-top.tex
% Continuation of C-exchange.tex: real vectors with prescribed zero term (Proposition pcx:P-slice),
% the large triple (Proposition pcx:P-top), the deduction of pcx:P-std, and the list of what
% Theorem pcx:X depends on. Labels of this file start with pcu:. Macros: those of C-exchange.tex.
\providecommand{\pcuS}{\mathsf S}
\providecommand{\pcuF}{\mathcal F}

\section{The large triple and the assembly of the exchange identity}\label{pcu:sec}

This section continues \S\ref{pcx:sec}; notation is that of \S\ref{pcx:ss-arch}. It proves
Proposition~\ref{pcx:P-slice}, proves Proposition~\ref{pcx:P-top} (the unfolding statements (U1),
(U2) being proved in \S\ref{pcu:ss-open}), deduces Proposition~\ref{pcx:P-std}, and lists in \S\ref{pcu:ss-final} every
statement on which Theorem~\ref{pcx:X} depends.

\subsection{Splittings and the mixed model at the real places}\label{pcu:ss-mixed}

Let $X$ be a space and $X=(\ell\oplus\ell')\perp X_1$ an adapted splitting
(Lemma~\ref{pcx:adapted}): $\ell=\langle e_1,\dots,e_p\rangle$, $\ell'=\langle f_1,\dots,f_p\rangle$
are isotropic, $\langle e_i,f_j\rangle=\delta_{ij}$, $X_1=(\ell\oplus\ell')^\perp$, and the frame
$(e,f,\text{frame of }X_1)$ is orthonormal (after the change $e\pm f$) at the infinite places and
integral with unimodular sum at the finite places. We write $D$ for the splitting and $D^{\mathrm
s}$ for the one with $e$ and $f$ exchanged. For a space $T$ of rank $a$ the real points of
$\auHom(X,T)$ decompose as $\auHom(\ell',T)_\R\oplus\auHom(\ell,T)_\R\oplus\auHom(X_1,T)_\R$, and
$\auHom(\ell',T)_\R\oplus\auHom(\ell,T)_\R$ is a polarised space whose Lagrangians are the two
summands. Identify $\auHom(\ell,T)_\R$ with $\prod_\tau\pcMat_{a\times p}(\C)$ through the frames.

\emph{The mixed model.} The Fock space of $\auHom(X,T)_\R$ is isomorphic, compatibly with the
Heisenberg group, to the space of functions $x\mapsto\Lambda(x)$ on $\prod_\tau\pcMat_{a\times
p}(\C)$ with values in the Fock space of $\auHom(X_1,T)_\R$ (Schr\"odinger model in the
$\auHom(\ell,T)$ variable, Fock model in the rest; this is the real form of the models (W2), (W2$'$)
of \S\ref{ps:ss-weil}, obtained by applying the Bargmann transform to the $\auHom(\ell\oplus\ell',T)$
factor of the Fock space). The value $\Lambda(x)$ is the \emph{slice} of
$\Lambda$ at $x$ along $D$, written $\mathrm{sl}_D(x)\Lambda$; for an adelic matrix $x$,
$\mathrm{sl}_D(x)$ means the slice at the archimedean components of $x$. The \emph{zero term}
of $\Lambda$ along $D$ is $\mathrm{sl}_D(0)\Lambda$; it is the restriction of the real
Schr\"odinger model to the matrices vanishing on the isotropic part, as in
Proposition~\ref{pcx:P-slice}. We use four standard properties.
\begin{enumerate}
\item[(M1)] \emph{Prescribed slices} (proved in \S\ref{pcu:ss-M1}). If $\tau$ is a smooth compactly supported function on
$\prod_\tau\pcMat_{a\times p}(\C)$, $u$ a smooth map from there to the real unitary group
$\prod_\tau\auU(X_1\otimes_\tau\C)\times\prod_\tau\auU(T\otimes_\tau\C)$, and $v$ a Fock vector of
$\auHom(X_1,T)_\R$ with finite moments, there is a Fock vector $\Lambda$ of $\auHom(X,T)_\R$ with
finite moments and $\mathrm{sl}_D(x)\Lambda=\tau(x)\,\omega_\infty(u(x))v$ for all $x$. (The unitary
groups here need not be compact; the proof in \S\ref{pcu:ss-M1} uses only that $\omega_\infty(g)$ carries the
vacuum to a Gaussian vector depending smoothly on $g$.)
\item[(M2)] \emph{Change of polarisation.} There is a constant $c_D>0$, depending only on the
normalisation of measures, such that for every $\Lambda$ with finite moments
$\mathrm{sl}_{D^{\mathrm s}}(0)\Lambda=c_D\int\mathrm{sl}_D(x)\Lambda\,dx$. \emph{Proof:} the
intertwiner between the Schr\"odinger models in $\auHom(\ell,T)$ and in $\auHom(\ell',T)$ is the
partial Fourier transform $\phi\mapsto c_D\int\phi(x)\psi_\infty(\langle x,x'\rangle)dx$ (Weil's
intertwining operator for two transversal Lagrangians, unique up to a scalar by Stone--von~Neumann);
at $x'=0$ it is $c_D\int\phi$, the integral converging absolutely by Lemma~\ref{pcu:slicedecay}:
$\int\|\mathrm{sl}_D(x)\Lambda\|_0dx\le C\|\Lambda\|_{2p+1}$.
\item[(M3)] \emph{Equivariance.} The real unitary group of $T$ acts on slices by
$\mathrm{sl}_D(x)\bigl(\omega_\infty(1,h)\Lambda\bigr)=\omega_\infty(1,h)\,\mathrm{sl}_D(h^{-1}x)\Lambda$,
and $\auU(X_1)_\R\subset\auU(X)_\R$ acts on the slices pointwise; with the characters of the
representation $\omega_{\alpha,\beta}$ (\S\ref{pcx:ss-statement}) the factor $\alpha^{-1}\beta$ of
the determinants is carried along.
\item[(M4)] \emph{Stages.} If $X_1=(\ell_1\oplus\ell'_1)\perp X_2$ is adapted, then $X=((\ell\oplus\ell_1)
\oplus(\ell'\oplus\ell'_1))\perp X_2$ is adapted, and the zero term along it is the zero term along
$X_1=(\ell_1\oplus\ell_1')\perp X_2$ of the zero term along $D$.
\end{enumerate}
(M3) and (M4) are immediate from the definition of the mixed model as a tensor product.
(M2)--(M4) are proved as indicated, (M1) in \S\ref{pcu:ss-M1}.

\begin{lemma}[The mixed model is unitary up to a constant]\label{pcu:mixedmodel}
Let $V_\ell=\auHom(\ell,T)_\R$ with a Haar measure $dx$, $\mathcal F'$, $\mathcal F$ the Fock spaces of
$\auHom(X,T)_\R$ and $\auHom(X_1,T)_\R$, and $\mathsf H=L^2(V_\ell,dx;\mathcal F)$. There is a bounded operator
$\mathsf M\colon\mathcal F'\to\mathsf H$ and a constant $c_{\mathsf M}>0$, depending only on $dx$, such that
\textup{(a)} $\|\mathsf Mv\|_{\mathsf H}^2=c_{\mathsf M}\|v\|^2$ for all $v\in\mathcal F'$; \textup{(b)} $\mathsf M$
intertwines the Heisenberg operators of $\mathcal F'$ with the operators of the mixed model (translation of the
variable by $\auHom(\ell,T)_\R$, multiplication by characters for $\auHom(\ell',T)_\R$, action on the values for
$\auHom(X_1,T)_\R$), in the form \eqref{pcu:eq-mixed}; \textup{(c)} for every $v$ with finite moments,
$(\mathsf Mv)(x)=\mathrm{sl}_D(x)v$ for all $x$. The family $\mathrm{sl}_D$ is determined on vectors with finite moments by
\eqref{pcu:eq-mixed}.
\end{lemma}
\begin{proof}
\emph{Uniqueness.} Suppose two families of bounded maps $\mathrm{sl}(x)$, $\mathrm{sl}'(x)$ satisfy
\eqref{pcu:eq-mixed}. Both send the vacuum $e$ to $e^{-|x|^2}e$, and \eqref{pcu:eq-mixed} expresses
$\mathrm{sl}(x)\psrho(y)e$ through $\mathrm{sl}(\cdot)e$ alone; so $\mathrm{sl}(x)\psrho(y)e=\mathrm{sl}'(x)\psrho(y)e$
for all $x,y$. A vector with finite moments is the absolutely convergent superposition
$v=\int v(y)\psrho(y)e\,dy$ (Lemma~\ref{pcx:superpos}), and $\mathrm{sl}(x)$, $\mathrm{sl}'(x)$ are bounded, so
$\mathrm{sl}(x)v=\mathrm{sl}'(x)v$.

\emph{Existence and unitarity.} Write $A=\auHom(\ell\oplus\ell',T)_\R$ and $B=\auHom(X_1,T)_\R$, so that
$\auHom(X,T)_\R=A\oplus B$ orthogonally, and use the Fock model \eqref{pk:eq-fock}: entire functions $F$ with
$\int|F|^2e^{-|z|^2}<\infty$, in coordinates for which the frames are orthonormal. (i) \emph{Tensor factorisation.}
In these coordinates $z=(z_A,z_B)$, $|z|^2=|z_A|^2+|z_B|^2$, and the Gaussian measure is a product; the monomials
$z_A^\alpha z_B^{\alpha'}$ are orthogonal with $\|z_A^\alpha z_B^{\alpha'}\|^2=\|z_A^\alpha\|^2\|z_B^{\alpha'}\|^2$, and
the polynomials are dense. Hence $F_A\otimes F_B\mapsto F_A(z_A)F_B(z_B)$ extends to a unitary isomorphism
$\mathcal F(A)\widehat\otimes\mathcal F(B)\to\mathcal F'$ (Hilbert tensor product). By \eqref{pk:eq-fock},
$e^{z\cdot y-|y|^2/2}F(z-\bar y)$ factorises for $y=y_A+y_B$, so the isomorphism carries
$\psrho_A(y_A)\otimes\psrho_B(y_B)$ to $\psrho(y_A+y_B)$, and the vacuum $1\otimes1$ to the vacuum.
(ii) \emph{Bargmann transform.} $A=\auHom(\ell,T)_\R\oplus\auHom(\ell',T)_\R$ is a sum of two transversal
Lagrangians for the alternating form of $\beta$, and the frame of $D$ at the infinite places is
orthonormal after the change $e\pm f$ (\S\ref{pcu:ss-mixed}), so the complex structure of the Fock model of $A$
exchanges the two summands. The Bargmann transform is a unitary operator, up to a positive constant fixed by
the normalisation of $dx$, from the Schr\"odinger model $L^2(V_\ell,dx)$ of $A$ onto $\mathcal F(A)$, which
intertwines the Schr\"odinger operators (translations by $\auHom(\ell,T)_\R$, multiplications by the characters
given by $\auHom(\ell',T)_\R$) with the Fock operators and carries a Gaussian $e^{-|x|^2}$, in suitable linear
coordinates $x$ on $V_\ell$, to a non-zero multiple of the vacuum (Folland~\cite[Chapter~1, \S6, pp.~40--44, (1.71)--(1.72)]{FollandPS}, where the Gaussian is $e^{-\boldsymbol\pi|x|^2}$). Let $\mathsf B\colon\mathcal F(A)\to L^2(V_\ell,dx)$ be its inverse, multiplied by the positive
constant for which $\mathsf B1=e^{-|x|^2}$. Then $\mathsf M=(\mathsf B\otimes1)\circ(\text{isomorphism of
(i)})^{-1}$ maps $\mathcal F'$ onto $\mathsf H=L^2(V_\ell)\widehat\otimes\mathcal F$, multiplies norms by a constant
$c_{\mathsf M}^{1/2}>0$, intertwines the Heisenberg operators as stated in (b) (this defines the linear maps $L$, $K$,
$I$ of \eqref{pcu:eq-mixed}: $L$ and $K$ are the inclusions of $\auHom(\ell,T)_\R$, $\auHom(\ell',T)_\R$ read in the
coordinates $x$, $I$ that of $B$), and $\mathsf Me=e^{-|x|^2}e$. This is (a) and (b).
(iii) \emph{Pointwise values.} For $v=\psrho(y)e$ the function $\mathsf Mv$ is, by (b), the continuous function
$x\mapsto(\text{phase})\,e^{-|x+Ay|^2}\psrho(Wy)e$ (notation of \S\ref{pcu:ss-M1}); so $x\mapsto(\mathsf M\psrho(y)e)(x)$
satisfies \eqref{pcu:eq-mixed} pointwise, and by the uniqueness above it equals $\mathrm{sl}_D(x)\psrho(y)e$. For $v$
with finite moments, $\mathsf Mv=\int v(y)\mathsf M\psrho(y)e\,dy$ in $\mathsf H$ ($\mathsf M$ bounded), and the
integral also converges pointwise in $\mathcal F$, uniformly in $x$ (the integrand has norm $\le|v(y)|$); the two
limits agree almost everywhere, and the pointwise one is $\mathrm{sl}_D(x)v$. This is (c).
\end{proof}

\subsection{Real vectors with prescribed zero term: proof of Proposition~\ref{pcx:P-slice}}
\label{pcu:ss-slice}

\begin{lemma}[Prescribed constant slices]\label{pcu:bump}
Let $D$ be an adapted splitting $X=(\ell\oplus\ell')\perp X_1$, $T$ a space, $\tau\ge0$ smooth with
compact support on $\prod_\tau\pcMat_{a\times p}(\C)$, $\tau\not\equiv0$, and $v$ a Fock vector of
$\auHom(X_1,T)_\R$ with finite moments. There are $\kappa\in\C$ and a Fock vector $\Lambda$ of
$\auHom(X,T)_\R$ with finite moments such that $\mathrm{sl}_D(x)\Lambda=\tau(x)\kappa\,v$ for all $x$ and
$\mathrm{sl}_{D^{\mathrm s}}(0)\Lambda=v$.
\end{lemma}
\begin{proof}
$\int\tau>0$. Put $\kappa=(c_D\int\tau)^{-1}$ and take $\Lambda$ from (M1) with the section
$\tau\kappa$ and $u\equiv1$. By (M2), $\mathrm{sl}_{D^{\mathrm s}}(0)\Lambda=c_D\kappa(\int\tau)v=v$.
\end{proof}

\begin{proof}[Proof of Proposition~\ref{pcx:P-slice}]
Apply Lemma~\ref{pcu:bump} to $D^{\mathrm s}$ (the splitting $W'=W\oplus\pcxH^p$ with the two
isotropic frames exchanged), $T=U$, any bump $\tau$ and $v=y$. The vector $w=\Lambda$ has finite
moments and its zero term along $(D^{\mathrm s})^{\mathrm s}=D$ is $y$.
\end{proof}

\begin{remark}
The zero term in Proposition~\ref{pcx:P-slice} is taken along the given splitting; by the lemma the
same holds along the splitting with the frames exchanged. Proposition~\ref{pcx:P-top} below produces
$x'$ whose zero term along $D^{\mathrm s}$ is $w$; Proposition~\ref{pcx:P-descent} must use the
same orientation (see \S\ref{pcu:ss-std}).
\end{remark}

\subsection{The large triple: setting}\label{pcu:ss-top-setting}

Fix the data of Proposition~\ref{pcx:P-top}: $W_0=\pcxH\oplus W$ of even rank $n=2j$ with the plane
split off adaptedly (splitting $D_1$), $U$ of rank $b$ with signs $t'$, $X=\pcxH^b\oplus W_0$ with
an adapted splitting $D\colon X=(\ell\oplus\ell')\perp W_0$, $\dim\ell=b$ (Lemma~\ref{pcx:adapted}
applied $b$ times, together with (M4)), and $T=U\oplus W_0$, of rank $b+n$ with signs $t'\oplus s$
($s$ those of $W_0$). Let $J=\mathrm{diag}(t'\oplus s)$ at each infinite place and
$x_0=\binom{1_b}{0}$, the frame of $U$ in $T$; its Gram matrix is $\mathrm{diag}(t')$, and its
stabiliser in $\auU(T)$ is $\auU(W_0)$. Write $\mathrm{Fr}=\prod_\tau\pcMat_{(b+n)\times b}(\C)$ for
families of frames, one for each infinite place. Note $\operatorname{rk}T-\operatorname{rk}X/2=
b+n-(2b+n)/2=j$, so in $\mathrm{RC}(X,T,U;\cdot)$ the factor is $L_*(\pcm,\nu,j)$, and
$\pcxstr_{W_0}$ also has $L_*(\pcm,\nu,j)$ in its denominator.

The parabolic $P_D$ of $\auU(X)$ stabilising $\ell$ consists of the products $n(c,z)\,m(q,g)$, with
$m(q,g)$ acting by $q\in\auGL_b$ on $\ell$ (in the frame $e$), by $q^{*-1}$ on $\ell'$ and by $g\in
\auU(W_0)$ on $W_0$, and $n(c,z)$ in the unipotent radical. Precisely, with $J_0$ the Gram matrix of the
frame of $W_0$, $c\in\pcMat_{n\times b}$ and $z\in\pcMat_b$ with
\[
z+z^*=-c^*J_0c ,
\]
$n(c,z)$ fixes $e$, maps $w\in W_0$ to $w+e\,c^*J_0w$ and $f$ to $f-c+e\,z$ (in the frames: the $k$-th vector
$f_k$ goes to $f_k-\sum_lc_{lk}w_l+\sum_ie_iz_{ik}$, $(w_l)$ the frame of $W_0$). The relation $z+z^*=-c^*J_0c$ is
exactly the condition that $n(c,z)$ is an isometry ($\langle e_i,f_j\rangle=\delta_{ij}$, the forms being
conjugate-linear in the first variable), and $n(c,z)n(c',z')=n(c+c',z+z'-c^*J_0c')$, so
$n(c,z)^{-1}=n(-c,z^*)$. The centre is $Z_D=\{n(0,z):z+z^*=0\}$: \emph{its elements are parametrised by the
skew-hermitian $b\times b$ matrices $z$}, and $z\mapsto n(0,z)$ is a homomorphism from the additive group.
For an automorphic function $f$ on $\auU(X)$ and a hermitian $b\times b$ matrix $\beta$ over $E$ the
\emph{centre coefficient} is
\begin{equation}\label{pcu:eq-psibeta}
\pcuF_\beta f(p)=\int_{Z_D(E)\backslash Z_D(\auA)}f(n(0,z)p)\,\psi_\beta(z)^{-1}\,dz ,\qquad
\psi_\beta(z)=\psi\bigl(-\tfrac12\jmath\operatorname{tr}(z\beta)\bigr),
\end{equation}
probability measure, $\psi$ the standard character and $\jmath=-1/\sqrt{-d}$ the constant of the Weil data
(\S\ref{ps:ss-weil}). For $z$ skew-hermitian and $\beta$ hermitian, $\operatorname{tr}(z\beta)$ is purely imaginary,
so $-\frac12\jmath\operatorname{tr}(z\beta)$ lies in $\auA_{E^+}$. We use the same formula for $\psi_\beta(z)$ when
$z+z^*=-c^*J_0c$: since $\psi$ is trivial on $\jmath\auA_{E^+}$ ($\operatorname{Tr}_{E/E^+}\jmath=0$), only the
skew-hermitian part $z+\frac12c^*J_0c$ of $z$ contributes. Writing $z=\jmath h$ with $h$ hermitian,
$\psi_\beta(z)=\psi(\operatorname{tr}(h\beta)/2d)$; the trace pairing on hermitian matrices being non-degenerate, the
characters of $Z_D(\auA)$ trivial on $Z_D(E)$ are exactly the $\psi_\beta$, $\beta$ hermitian over $E$. This is the
character by which the centre acts on the terms (Lemma~\ref{pcu:laws}): the term of $x$ is multiplied by
$\psi_{x^*J_Vx}$, and $\pcuF_\beta$ keeps the $x$ with Gram matrix $x^*J_Vx=\beta$.

A pair $(\beta,q)$, $\beta$ hermitian over $E$, $q\in\auGL_b(\auA_E)$, is a \emph{test pair for $U$} if
either (a) $\beta=J_U$ is the form of $U$ and $q=s_0\,x_U\,Y$, where $s_0$ is the finite idele with
$s_0\widehat\auO_E=\alpha\widehat\auO_E$ fixed in \S\ref{ps:ss-weil} (denoted $a$ there; $\alpha$ is the idele by
which the Weil data scale $\psi$; for the set of test pairs only the ideal $s_0\widehat\auO_E$ matters: replacing $s_0$
by $s_0\varepsilon$ with $\varepsilon$ a unit idele replaces $Y$ by $\varepsilon Y$, and $\varepsilon Yk=Y(\varepsilon k)$
with $\varepsilon k\in\auGL_b(\widehat\auO_E)$), $x_U$ the adelic frame of $U$, and the Gram matrix of the columns
of $x_UY_f$ for the form of $U$ lies in $\auGL_b(\widehat\auO_E)$, i.e.\ these columns span a lattice self-dual
for $J_U$ (a condition invariant under $Y\mapsto Yk$, $k\in\auGL_b(\widehat\auO_E)$, and $Y\mapsto\varepsilon Y$); or (b) $\beta$ is invertible and not of the form $a^*\mathrm{diag}(t')a$ at some infinite
place.

\emph{Coordinates.} At the finite places, matrices of $\auHom(X,V)$ are written in the coordinates
$\tilde x=s_0^{-1}x$ of Lemma~\ref{pgx:dict} (on both components at a split place). In them the self-dual
lattice $\alpha\auHom(L_X,L_V)$, whose vacuum is the datum at the finite places outside $S\cup\bar S$
(the \emph{standard vacuum}), is the lattice of integral matrices, at $w\in S$ the data of
Definition~\ref{pgx:conv} are functions of $\tilde x$ (the paragraph after Lemma~\ref{pcx:mod}),
and ``integral'' always refers to these coordinates. By Lemma~\ref{pgx:dict}(d), a point
$x=s_0y$ of $\auHom(\ell,V)(\auA_f)$ has coordinate $\tilde x=y$.

Write $z_0\in\pcMat_{b\times(2b+n)}(\auO_S)$ for the block $f^*$ of the inverse of the frame of $D$
read at $S$ (the matrix whose rows are the coordinate functionals of the $\ell'$-part), so that
$\Phi^{\mathrm{row}}_{z_0}$ and $\Psi^{\mathrm{col}}_{z_0}$ are the standard data of the frame of $D$.

\subsection{Smooth Witt data}\label{pcu:ss-witt}

\begin{lemma}[Smooth local Witt theorem for spans]\label{pcu:witt}
Let $J$ be a hermitian $m\times m$ complex matrix and $x$ an $m\times b$ matrix with $\det(x^*Jx)\ne
0$. There is a smooth map $k$ from $\pcMat_{m\times b}(\C)$ to the unitary group $\auU(J)$ such that
for all $y$ near $x$ there is $c\in\pcMat_b(\C)$ with $k(y)\,x\,c=y$.
\end{lemma}
\begin{proof}
Put $w=(x^*Jx)^{-1}x^*$, so $ywJx=y$ for every $y$. For $a\in\pcMat_m(\C)$ let $L(a)=(axw-xwa^*)J$.
Then $L(a)^*J+JL(a)=0$, i.e.\ $L(a)$ is in the Lie algebra of $\auU(J)$, and $(a-L(a))x=x(wa^*Jx)$
because $L(a)x=axwJx-xwa^*Jx=ax-xwa^*Jx$. Let $T(a)=\exp(L(a))(a-L(a))$. Since $L(1)=0$, $T(1)=1$,
and the derivative of $T$ at $1$ is $(\mathrm{id}-L)+L=\mathrm{id}$. By the inverse function
theorem there is a smooth $A$ on $\pcMat_m(\C)$ with $T(A(z))=z$ for $z$ near $1$ (extend a local
inverse smoothly). For $y$ put $z(y)=1-(x-y)wJ$, which is near $1$ for $y$ near $x$, and
$k(y)=\exp(L(A(z(y))))\in\auU(J)$, smooth in $y$. For $y$ near $x$, with $a=A(z(y))$,
\[
k(y)\,x\,(wa^*Jx)=\exp(L(a))(a-L(a))x=T(a)x=z(y)x=x-(x-y)wJx=y. \qedhere
\]
\end{proof}

\begin{lemma}[Witt data with a bump]\label{pcu:wittbump}
For $J=\mathrm{diag}(t'\oplus s)$ and $x_0$ as above there are: a compact neighbourhood $N$ of
$\mathrm{diag}(t')$ in $\auGL_b(\C)$ and a smooth $\iota$ on $\pcMat_b(\C)$ with $\iota(A)=A^{-1}$ for $A\in N$; a smooth $k\colon\pcMat_{(b+n)\times b}(\C)\to\auU(J)$; and
a smooth compactly supported $\sigma\ge0$ on $\mathrm{Fr}$ with $\sigma(x_0,\dots,x_0)=1$, such that
whenever $\sigma(y)\ne0$, for every $\tau$: $y_\tau^*Jy_\tau\in N$ (so
$\iota(y_\tau^*Jy_\tau)=(y_\tau^*Jy_\tau)^{-1}$), and $k(y_\tau)x_0c=y_\tau$ for some $c$.
\end{lemma}
\begin{proof}
Let $N$ be a compact neighbourhood of $\mathrm{diag}(t')$ in the open set $\auGL_b(\C)$, $\chi$ a smooth
function with compact support in $\auGL_b(\C)$ equal to $1$ on $N$, and $\iota(A)=\chi(A)\operatorname{adj}(A)/\det A$
($0$ off the support of $\chi$); it is smooth on $\pcMat_b(\C)$ and inverts on $N$. Lemma~\ref{pcu:witt} gives $k$.
The conditions $y^*Jy\in N$ and $k(y)x_0c=y$ hold on a neighbourhood of $x_0$ (the first because
$y\mapsto y^*Jy$ is continuous and $x_0^*Jx_0=\mathrm{diag}(t')$ is interior to $N$), and $\sigma$ is a
product over $\tau$ of a bump supported in it.
\end{proof}

For $x\in\pcMat_{(b+n)\times b}(\C)$ let $\mathrm{pr}(x)=x\,\iota(x^*Jx)\,x^*Jx_0$ (where $\iota$ inverts,
the $J$-orthogonal projection of $x_0$ to the span of $x$); $\mathrm{pr}$ is smooth on all of
$\pcMat_{(b+n)\times b}(\C)$. If $C$ is invertible and $\iota$ inverts the Gram matrices of $x$ and $xC$, then
$\mathrm{pr}(xC)=xC\,C^{-1}(x^*Jx)^{-1}C^{*-1}C^*x^*Jx_0=\mathrm{pr}(x)$.

\begin{lemma}\label{pcu:frame}
If $\sigma(\mathrm{pr}\circ x)\ne0$ then for every $\tau$, $x_\tau=k(\mathrm{pr}(x_\tau))x_0c$ with
$c^*\mathrm{diag}(t')c=x_\tau^*Jx_\tau$; in particular $x_\tau^*Jx_\tau$ has the form
$c^*\mathrm{diag}(t')c$.
\end{lemma}
\begin{proof}
By Lemma~\ref{pcu:wittbump} applied to $y=\mathrm{pr}\circ x$, $y_\tau=k(y_\tau)x_0c_1$ and $y_\tau$ has
invertible Gram matrix; since $y_\tau=x_\tau A$ with $A=\iota(\cdot)x_\tau^*Jx_0$ and
$\det(y_\tau^*Jy_\tau)\ne0$, $A$ is invertible and $x_\tau=k(y_\tau)x_0c_1A^{-1}$. Put $c=c_1A^{-1}$;
as $k(y_\tau)$ is unitary, $x_\tau^*Jx_\tau=c^*x_0^*Jx_0c=c^*\mathrm{diag}(t')c$.
\end{proof}

\subsection{The two real vectors and the identity of periods}\label{pcu:ss-period}

Let $\mu$ be the weight ($\mu_\tau\ge1$). Let $w$ be the Fock vector of $(W_0,U)$ and $z$ that of
$(W_0,W_0)$ given in Proposition~\ref{pcx:P-top}. Let $\omega^\flat$ be the real Weil representation
of $\auU(W_0)_\R\times\auU(T)_\R$ on the Fock space of $\auHom(W_0,T)_\R$ twisted by
$\prod_\tau\det(h_\tau)^{j-\mu_\tau}$ on $\auU(T)_\R$, and $w\boxtimes z$ the external product (a vector
of $\auHom(W_0,U\oplus W_0)_\R$ with finite moments, its moments being bounded by products of those of
$w$, $z$ and the vacuum, since it is an iterated integral of translated vacua,
Lemma~\ref{pcx:superpos}).

\emph{The vector $\Lambda$ for $(X,U)$.} By Lemma~\ref{pcu:bump} with
$\tau(B)=\sigma((B_\tau;0)_\tau)\,\bigl|\prod_\tau\det(B_\tau)^{j+\mu_\tau}\bigr|^2$ ($\tau(1)=1$) there are
$\kappa$ and $\Lambda$ with finite moments, zero term $w$ along $D^{\mathrm s}$, and
$\mathrm{sl}_D(B)\Lambda=\sigma((B;0))\,\bigl|\prod_\tau\det B_\tau^{j+\mu_\tau}\bigr|^2\kappa\,w$.

\emph{The vector $\Phi$ for $(X,T)$.} For $y,x\in\mathrm{Fr}$ put $x'_\tau=k(y_\tau)^{-1}x_\tau$,
$(x'_\tau)_1$ its upper $b\times b$ block, and
\[
\pcuS(y,x)=\sigma(y)\sigma(x')\,\kappa\prod_\tau\overline{\det(x'_\tau)_1}^{\,2\mu_\tau}\;
\omega^\flat(1,k\circ y)(w\boxtimes z).
\]
The function $x\mapsto\pcuS(\mathrm{pr}\circ x,x)$ is smooth and compactly supported (on its support
$x_\tau=k(y_\tau)x'_\tau$ with $y$, $x'$ in the support of $\sigma$), of the form
$\tau(x)\omega^\flat(u(x))v$ with $\tau$, $u$ smooth (the twist is a polynomial in the entries of
$k(y_\tau)$ and their conjugates). By (M1) there is $\Phi$ with finite moments and
$\mathrm{sl}_D(x)\Phi=\pcuS(\mathrm{pr}\circ x,x)$.

\begin{lemma}[Support of the slices]\label{pcu:support}
If $\mathrm{sl}_D(x)\Phi\ne0$ then $x_\tau^*Jx_\tau$ has the form $c^*\mathrm{diag}(t')c$ for every
$\tau$.
\end{lemma}
\begin{proof}
Then $\sigma(\mathrm{pr}\circ x)\ne0$; apply Lemma~\ref{pcu:frame}.
\end{proof}

For a finite datum $F$ of a pair $(W_0,V)$, an adelic matrix $x$, a point $(g,h)$ and characters
$\alpha,\beta$, write $\vartheta^{\alpha,\beta}_{F,x,(g,h)}(\phi)=\theta\bigl(\psrho(x)\omega_{\alpha,
\beta}(g,h)(\phi\otimes F)\bigr)$ for the theta functional on real vectors $\phi$.

\begin{lemma}[Product formula]\label{pcu:product}
For a finite datum $f$ of $(W_0,U)$, a datum $F$ of $(W_0,W_0)$, real vectors $\phi$, $v_1$ and
characters $\alpha,\beta$,
\[
\vartheta^{\alpha,\alpha\beta}_{f\boxtimes F,(x;0),(g,y\oplus l)}(\phi\boxtimes v_1)=
\vartheta^{\alpha,\beta}_{f,x,(g,y)}(\phi)\;\theta^{W_0,W_0}_{\alpha,\alpha}(g,l;v_1\otimes F).
\]
\end{lemma}
\begin{proof}
$\auHom(W_0,U\oplus W_0)=\auHom(W_0,U)\oplus\auHom(W_0,W_0)$ orthogonally; the Weil representation,
the Heisenberg operators and $\theta$ are tensor products along this decomposition (the character
$\alpha\beta$ on $\auU(U\oplus W_0)$ restricted to $\auU(U)\times\auU(W_0)$ is $\alpha\beta\otimes
\alpha\beta$, and the factor $\beta$ is assigned to the first factor, $\alpha$ to the second, by the
convention \eqref{pcx:eq-twist}), and the translation is $0$ on the second factor.
\end{proof}

Let $\gamma=\nu^{-1}\upsilon^{j}$ (the character of \S\ref{pcx:ss-arch} for the rank $n$ of $W_0$),
$I_0(\nu)=\Theta^{\mathrm L}_{W_0,W_0}(z,\Psi^{\mathrm{col}}_1)(1)$, and for $h\in\auU(T)_\R$ let
$\rho(h)=\sigma\bigl(\mathrm{pr}\circ(h_\tau^{-1}x_0)_\tau\bigr)\ge0$; $\rho$ is continuous and
$\rho(1)=1$. For $B\in\auGL_b(\auA_E)$ let $\zeta_B$ be the idele equal to $\det B$ at the infinite
places and $1$ at the finite ones, and $\mathrm N(\zeta)=\prod_\tau|\zeta_\tau|^2$.

\begin{lemma}[Identity of periods]\label{pcu:period}
For every $h\in\auU(T)_\R$ there is $p\in\auU(U)_\R\times\auU(W_0)_\R$ such that for all
$B\in\auGL_b(\auA_E)$, characters $\beta$, finite data $f$ of $(W_0,U)$, adelic $x$, and $g$, $y$,
\[
\int_{[\auU(W_0)]}\vartheta^{\gamma,\gamma\beta}_{f\boxtimes F_1,(x;0),(g,y\oplus u)}
\Bigl(\omega^\flat(1,ph)\,\mathrm{sl}_D\bigl((ph)^{-1}(B;0)\bigr)\Phi\Bigr)du
=\frac{\rho(h)\,I_0(\nu)}{\nu_{\mathrm{bc}}(\zeta_B)\,\mathrm N(\zeta_B)^{j}}\;
\vartheta^{\gamma,\beta}_{f,x,(g,y)}\bigl(\mathrm{sl}_D(B)\Lambda\bigr),
\]
where $F_1$ is the finite datum $\Psi^{\mathrm{col}}_1=\chi_S^{-1}\circ\det$ at $S$ and the vacuum
elsewhere, and $\nu_{\mathrm{bc}}(t)=\nu(t/\bar t)$.
\end{lemma}
\begin{proof}
Let $x_{1,\tau}=h_\tau^{-1}x_0$, so $x_{1,\tau}^*Jx_{1,\tau}=\mathrm{diag}(t')$ and $\iota$ inverts it;
let $y=\mathrm{pr}\circ x_1$, so $\rho(h)=\sigma(y)$.

If $\sigma(y)\ne0$, Lemma~\ref{pcu:frame} gives $x_{1,\tau}=k(y_\tau)x_0u_\tau$ with $u_\tau^*\mathrm{diag}
(t')u_\tau=\mathrm{diag}(t')$, i.e.\ $u\in\auU(U)_\R$. Then $h_\tau k(y_\tau)\mathrm{diag}(u_\tau,1)$ fixes
$x_0$, hence (being unitary) is $\mathrm{diag}(1,q_\tau)$ with $q_\tau\in\auU(W_0)_\R$. Put
$p=(u,q^{-1})$, so $p h=k\circ y^{-1}$ pointwise. If $\sigma(y)=0$ take $p=1$.

At each $\tau$, $(ph)^{-1}(B;0)=h^{-1}x_0u^{-1}B_\tau=x_{1,\tau}C_\tau$ with $C_\tau=u_\tau^{-1}B_\tau$
invertible, and the slice of $\Phi$ there is $\pcuS(y',x_1C)$ with $y'=\mathrm{pr}\circ(x_1C)$.
\emph{If this slice is non-zero, then $y'=y$.} Indeed then $\sigma(x'')\ne0$ for $x''_\tau=k(y'_\tau)^{-1}x_{1,\tau}C_\tau$,
so by Lemma~\ref{pcu:wittbump} the Gram matrix $x''^*_\tau Jx''_\tau=C_\tau^*\mathrm{diag}(t')C_\tau$ ($k(y'_\tau)$ is
unitary) lies in $N$; as $\mathrm{diag}(t')\in N$ too, $\iota$ inverts the Gram matrices of $x_1$ and of $x_1C$,
and $y'=\mathrm{pr}\circ(x_1C)=\mathrm{pr}\circ x_1=y$.
If $\sigma(y)=0$ the slice therefore vanishes (it contains the factor $\sigma(y')=\sigma(y)$ when non-zero), and so does the right
side ($\rho(h)=0$). Otherwise $k(y_\tau)^{-1}x_{1,\tau}C_\tau=x_0u_\tau C_\tau=(B_\tau;0)$. If $\sigma((B;0))=0$, the right
side vanishes ($\mathrm{sl}_D(B)\Lambda$ contains the factor $\sigma((B;0))$), and so does the left: a non-zero slice
would have $y'=y$, hence contain the factor $\sigma(k(y)^{-1}x_1C)=\sigma((B;0))=0$. If $\sigma((B;0))\ne0$, then
$C_\tau^*\mathrm{diag}(t')C_\tau=(B_\tau;0)^*J(B_\tau;0)\in N$ by Lemma~\ref{pcu:wittbump}, so again $y'=y$ and the
slice is $\pcuS(y,x_1C)$. It remains to treat this last case, $\sigma(y)\sigma((B;0))\ne0$, where
\[
\omega^\flat(1,ph)\pcuS(y,x_1C)=\sigma(y)\,\sigma((B;0))\,\kappa\prod_\tau\overline{\det B_\tau}^{\,2\mu_\tau}
\,(w\boxtimes z),
\]
the operators $\omega^\flat(1,ph)$ and $\omega^\flat(1,k\circ y)$ cancelling. By
Lemma~\ref{pcu:product} with $\alpha=\gamma$, the integrand is this scalar times
$\vartheta^{\gamma,\beta}_{f,x,(g,y)}(w)\,\theta^{W_0,W_0}_{\gamma,\gamma}(g,u;z\otimes F_1)$; the integral
over $u$ of the last factor is $\Theta^{\mathrm L}_{W_0,W_0}(z,\Psi^{\mathrm{col}}_1)(g)=I_0(\nu)$ by
Lemma~\ref{pcx:const}. On the right, $\mathrm{sl}_D(B)\Lambda=\sigma((B;0))|\prod\det B_\tau^{j+\mu_\tau}|^2
\kappa w$. It remains to see $\prod_\tau\bar\zeta_\tau^{2\mu_\tau}=|\prod_\tau\zeta_\tau^{j+\mu_\tau}|^2/
(\nu_{\mathrm{bc}}(\zeta)\mathrm N(\zeta)^j)$ for $\zeta=\zeta_B$. With $P=\prod\zeta_\tau$ and
$Q=\prod\zeta_\tau^{\mu_\tau}$ one has $\nu_{\mathrm{bc}}(\zeta)=Q^2/|Q|^2$ (weight $\mu$ at the
infinite places, \S\ref{pcx:ss-statement}), $\mathrm N(\zeta)^j=|P|^{2j}$, and
$|P^jQ|^2|Q|^2/(Q^2|P|^{2j})=|Q|^4/Q^2=\bar Q^2$.
\end{proof}

\subsection{Centre coefficients at test pairs}\label{pcu:ss-coeff}

Put $G=\auU(T)(\auA)$, $H=\auU(U)(\auA)\times\auU(W_0)(\auA)\subset G$ (block diagonal), and
$\Gamma_2=\auU(W_0)(E)$. The unfolding of the theta kernel along $D$ uses three facts, stated here;
(U1), (U2) are proved in \S\ref{pcu:ss-open} and (U3) below.

\begin{enumerate}
\item[(U1)] \emph{Unfolding.} Let $V$ be a space of rank $a$, $F=[\phi,\Phi]$ a datum of
$\auHom(X,V)$ with $\Phi$ one of the standard data $\Phi^{\mathrm{row}}_{z_0}$ (for $V=T$) or
$(\Psi^{\mathrm{col}}_{z_0})^{\mathsf t}$ at $\bar S$ (for $V=U$), and $p=n(c,z)m(q,g)$. The kernel
$\theta(\omega(p,h)F)$ is $\sum_{x\in\pcMat_{a\times b}(E)}\mathrm{term}_x(p,h)$, where
$\mathrm{term}_x(p,h)$ is the partial theta sum over the $\xi\in\auHom(X,V)(E)$ whose restriction to
$\ell$ is $x$. Let $\beta$ be invertible. If there is no $x$ with $x^*J_Vx=\beta$, or if
$\mathrm{term}_x(p,\cdot)\equiv0$ for all such $x$, then $\pcuF_\beta\int_{[\auU(V)]}\theta(\omega(\cdot,h)F)dh$
vanishes at $p$. Otherwise the $x$ with $x^*J_Vx=\beta$ form one $\auU(V)(E)$-orbit, and if
$\mathrm{term}_{x}(p,\cdot)$ is absolutely integrable on $\mathrm{Stab}(x)(E)\backslash\auU(V)(\auA)$ (condition
\eqref{pcu:eq-unfoldabs}) then $\pcuF_\beta\int_{[\auU(V)]}\theta(\omega(\cdot,h)F)dh$ at $p$ equals
$\int_{\mathrm{Stab}(x)(E)\backslash\auU(V)(\auA)}\mathrm{term}_x(p,h)\,dh$. Condition \eqref{pcu:eq-unfoldabs}
holds in the two situations of Lemma~\ref{pcu:integrable} (the orbit of $x_0$ on the large side and the
orbit of $1$ on the small side, at a test pair of type (a)), which are the ones used in
Lemma~\ref{pcu:coeff} and in type (a) of the proof of Proposition~\ref{pct:vantop}.
\item[(U2)] \emph{Term formula.} With $x'=x_V^{-1}x$ (standard coordinates) and $x'q=s_0Y$,
\[
\mathrm{term}_x(p,h)=\psi_{x^*J_Vx}(z)\,\chi_V(\det q)^{-1}
|\det q|^{a/2}\,\mathbf 1_{\mathrm{int}}(y)\,\varepsilon_V\;
\vartheta_{F^{(y)},\,x'c^*J_{W_0},\,(g,h)}\bigl(\mathrm{sl}_D(h^{-1}Y)\phi\bigr),
\]
where $y$ is the finite part of $h^{-1}Y$, $\mathbf 1_{\mathrm{int}}(y)$ is $1$ if $y$ is integral and
$0$ otherwise, $\varepsilon_V\ne0$ is the constant in the slices of the vacuum, $F^{(y)}$ is the
finite datum of $(W_0,V)$ equal at $S$ (or $\bar S$) to the slice $y'\mapsto\Phi((y\ y'))$ of $\Phi$
at $y$ and to the vacuum elsewhere, $\chi_V(u)=\alpha_V(u/\bar u)\upsilon(u)^{-a}$ for the source
character $\alpha_V$, and the characters of $\vartheta$ are those of the pair $(X,V)$ with the target
character divided by $\upsilon^b$.
\item[(U3)] \emph{Stages.} Let $\Theta\colon\Gamma_2\backslash G\to\C$ be integrable, and suppose
there are a measurable $\rho'\ge0$ on $G$, left $H$-invariant, non-zero on a neighbourhood of $1$, and
$c_0\in\C$ with $\int_{\Gamma_2\backslash H}\Theta(\xi g)\,d\xi=c_0\rho'(g)$ for almost all $g$. Then
$\int_{\Gamma_2\backslash G}\Theta=R\,c_0$ with $R=\int_{H\backslash G}\rho'\in(0,\infty]$, where if
$R=\infty$ then $c_0=0$ and the identity reads $\int\Theta=0$.
\end{enumerate}
\begin{proof}[Proof of (U3)]
By the quotient integral formula, $\int_{\Gamma_2\backslash G}\Theta=\int_{H\backslash G}
\int_{\Gamma_2\backslash H}\Theta(\xi g)d\xi\,dg=c_0\int_{H\backslash G}\rho'$. $R>0$ because
$\rho'$ is non-zero near $1$ and $H\backslash G$ carries a quasi-invariant measure positive on open
sets; if $R=\infty$, integrability of $\Theta$ forces $c_0=0$.
\end{proof}

\begin{lemma}[Slices of a datum at a split place]\label{pcu:W5slice}
Let $v=w\bar w$ be a place of $E^+$ split in $E$, $V$ a space of rank $a$, and realise $\psHS_v(\auHom(X,V))$ on
$\auS(\auHom(X_w,V_w))$ as in \textup{(W5)}: $F$ is a function of $a\times(2b+n)$ matrices
$\xi=(\xi_\ell\ \xi_{\ell'}\ \xi_{W_0})$, written in the frame $(e,f,\text{frame of }W_0)$ of $D$, with
$\psrho(a_w,0)$ acting by translation and $\psrho(0,a_{\bar w})$ by multiplication by $[\xi,a_{\bar w}]$. Then the
slice of $F$ along $D$ at a point $u=(u_w,u_{\bar w})\in\auHom(\ell,V)_v=\pcMat_{a\times b}(E_w)\times\pcMat_{a\times
b}(E_{\bar w})$ is, up to a constant $c_v>0$ depending only on Haar measures, the function of
$\xi_{W_0}\in\auHom(W_{0,w},V_w)$ (the \textup{(W5)} model of $\auHom(W_0,V)_v$)
\[
F[u](\xi_{W_0})=c_v\int_{\pcMat_{a\times b}(E_w)}F(u_w\ t\ \xi_{W_0})\,[t,u_{\bar w}]\,dt ,
\]
where $t$ runs through the $\ell'$-columns and $[t,u_{\bar w}]$ is the pairing $[\,,]$ between $\auHom(\ell'_w,V_w)$
and $\auHom(\ell_{\bar w},V_{\bar w})$. If $F$ is the indicator of the integral $\ell'$-columns times a function
$F_1(\xi_\ell,\xi_{W_0})$, and the integral matrices of $\auHom(\ell'_w,V_w)$ and $\auHom(\ell_{\bar w},V_{\bar w})$
are dual lattices for $[\,,]$, then $F[u](\xi_{W_0})=c'_v\,\mathbf 1_{\mathrm{int}}(u_{\bar w})\,F_1(u_w,\xi_{W_0})$,
$c'_v>0$.
\end{lemma}
\begin{proof}
In the mixed model of $D$ at $v$ the Schr\"odinger variable is $u\in\auHom(\ell,V)_v=I^*_v$ (maps vanishing on
$\ell'\oplus W_0$), $I_v$ (maps vanishing on $\ell\oplus W_0$) acts by multiplication, and $\auHom(W_0,V)_v$ on the
values (\S\ref{pcu:ss-open}). Compare with the \textup{(W5)} model, which is the Schr\"odinger model of the
polarisation $\auHom(X_w,V_w)\oplus\auHom(X_{\bar w},V_{\bar w})$. The two polarisations agree on
$\auHom(\ell_w,V_w)$ (translated in both), on $\auHom(\ell'_{\bar w},V_{\bar w})$ (multiplication in both), and on
$\auHom(W_0,V)_v$, which is the orthogonal summand on which both are the \textup{(W5)} model of $(W_0,V)$; they differ
exactly on $\auHom(\ell'_w,V_w)$ (translated in \textup{(W5)}, in $I_v$ for $D$) and $\auHom(\ell_{\bar w},V_{\bar w})$
(multiplication in \textup{(W5)}, in $I^*_v$ for $D$). These two are in duality for $[\,,]$, since
$\langle e_i,f_j\rangle=\delta_{ij}$ and the form pairs the $w$-component with the $\bar w$-component. The
intertwiner between the Schr\"odinger models of the two transversal Lagrangians of the symplectic space
$\auHom(\ell'_w,V_w)\oplus\auHom(\ell_{\bar w},V_{\bar w})$ is the partial Fourier transform
$\varphi\mapsto c_v\int\varphi(t)[t,\cdot\,]\,dt$ (Weil's intertwining operator, unique up to a scalar by (P1); one
checks directly that it carries translation by $t_0$ to multiplication by $[\cdot,t_0]=[t_0,\cdot\,]^{-1}$ and multiplication by
$[\cdot,y_0]$ to translation by $y_0$, which are the operators of (P2) in the new model). Tensoring with the identity on the other factors gives
the intertwiner from the \textup{(W5)} model to the mixed model, and the slice at $u$ is its value at $u$; this is the
formula. For the second statement, $\int_{\mathrm{int}}[t,u_{\bar w}]\,dt$ is the volume of the integral matrices if
$u_{\bar w}$ lies in the dual lattice, which is the lattice of integral matrices by hypothesis, and $0$ otherwise.
\end{proof}

\begin{lemma}[Finite slices at test pairs of type (a)]\label{pcu:finslice}
Let $(J_U,q)$ be a test pair of type (a), $q=s_0x_UY$, and $p=n(c,z)m(q,g)$; let $F=[\phi,\Phi^{\mathrm{row}}_{z_0}]_S$
($V=T$), $K_f$ the stabiliser in $\auU(T)(\auA_f)$ of the standard lattice $L_T=L_U\perp L_{W_0}$, $H_f=\auU(U)(\auA_f)
\times\auU(W_0)(\auA_f)$, and let $M_0\subset U(\auA_f)$ be the lattice spanned by the columns of $x_0Y_f$ (self-dual for the
form of $U$, by the definition of a test pair). The slice of $\Phi=\Phi^{\mathrm{row}}_{z_0}=\epsilon\cdot(\text{standard datum of
the frame of }D)$ entering (U2) is computed as follows. In the frame of $D$ at $S$ the rows of $z_0$ are the coordinate functionals of
$\ell'$, so for $\xi=(\xi_\ell\ \xi_{\ell'}\ \xi_{W_0})$ the matrix $\binom{z_0}{\xi}$ is $\begin{pmatrix}0&1_b&0\\\xi_\ell&\xi_{\ell'}&\xi_{W_0}
\end{pmatrix}$; it lies in $\auGL_{2b+n}(\auO_S)$ if and only if $\xi_{\ell'}$ is integral and $(\xi_\ell\ \xi_{W_0})\in\auGL_{b+n}(\auO_S)$, and
then its determinant is $\pm\det(\xi_\ell\ \xi_{W_0})$ (expansion along the rows of $z_0$), with $\chi_S(\pm1)=\nu(1)=1$. Hence
$\Phi(\xi)=\epsilon\,\mathbf 1_{\mathrm{int}}(\xi_{\ell'})\,\Phi_1(\xi_\ell,\xi_{W_0})$ with $\Phi_1(y,y')=\chi_S^{-1}(\det(y\ y'))$ if the square
matrix $(y\ y')$ (columns of $\ell\oplus W_0$) lies in $\auGL_{b+n}(\auO_S)$ and $0$ otherwise. At each place $v=w\bar w$ with $w\in S$
the frame of $D$ is integral with unimodular Gram matrix, so $\alpha\auHom(L_X,L_T)_v$ is the sum of its intersections
with the summands of the decomposition along $\ell,\ell',W_0$; hence (Lemma~\ref{pgx:dict}(b)) the integral $\ell'_w$-columns
and the integral matrices of $\auHom(\ell_{\bar w},T_{\bar w})$, both in the coordinates $\tilde x$, are dual lattices, and Lemma~\ref{pcu:W5slice} gives the slice at $u=(u_w,u_{\bar w})$:
$c'_v\,\mathbf 1_{\mathrm{int}}(u_{\bar w})\,\epsilon\,\Phi_1(u_w,\cdot)$. The factor $\mathbf 1_{\mathrm{int}}(u_{\bar w})$ is the integrality
condition at $\bar S$ recorded in $\mathbf 1_{\mathrm{int}}(y)$ in (U2), and $c'_v$ is part of $\varepsilon_T$; so in (U2) the datum
$F^{(y)}$ at $S$ is $y'\mapsto\Phi((y\ y'))=\epsilon\,\Phi_1(y,y')$, $y$ the $S$-component of the slice point.
\begin{enumerate}
\item[\textup{(a)}] For $\kappa\in K_f$, $\omega_{\gamma_X,\gamma_X}(1,\kappa)F=F$; hence $\mathrm{term}_{x_0}(p,h\kappa)=\mathrm{term}_{x_0}(p,h)$.
\item[\textup{(b)}] If $\mathrm{term}_{x_0}(p,h)\ne0$, then the finite part of $h^{-1}x_0Y$ is integral and $h_f\in H_fK_f$.
\item[\textup{(c)}] Let $h_f=(y_f,y'_f)\in H_f$ and $x_f=y_f^{-1}Y_f$. If $\mathrm{term}_{x_0}(p,h)\ne0$ then $x_f$ is integral,
$\det x_f\in\auO_S^\times$, and the finite datum $F^{(y)}$ of (U2) at $S$ is
$\epsilon\,\chi_S^{-1}(\det x_f)\,f_U\otimes F_1$, where $f_U$ is the vacuum datum (indicator of the integral matrices) of
$\auHom(W_0,U)$ at $S$ and $F_1=\Psi^{\mathrm{col}}_{1_n}$ is the datum of $\auHom(W_0,W_0)$ at $S$ of
Definition~\ref{pcx:def-std} (with $z=1_n$), supported on $\auGL_n(\auO_S)$.
\end{enumerate}
\end{lemma}
\begin{proof}
(a) $\mathrm{term}_{x_0}(p,h\kappa)=\theta_{W_0}((\omega(p,h)\omega(1,\kappa)F)[x_0])$, so the second claim follows from the
first. Here $\omega=\omega_{\gamma_X,\gamma_X}$ with $\gamma_X=\nu^{-1}\upsilon^{(2b+n)/2}=\gamma\,\upsilon^{b}$ (the character attached to the rank $2b+n$ of $X$; $\gamma=\nu^{-1}\upsilon^j$ as before Lemma~\ref{pcu:period}, the shift by $\upsilon^b$ being the one in (U2)), so $\omega(1,\kappa)=\gamma_X(\det\kappa)\,
\omega_{\auHom(X,T)}(1,\kappa)$, and $\det\kappa=(\det\kappa_v)_v$ with $\det\kappa_v$ a norm-one unit at every finite place $v$ of
$E^+$ ($\kappa_v$ stabilises the lattice $L_{T,v}$). Two facts on the characters are used: $\upsilon$ is unramified at every finite
place, i.e.\ trivial on $\widehat\auO_E^\times$ (\S\ref{ps:ss-weil}); and $\nu$ is unramified outside $\pcm$, i.e.\ trivial on
$\auU(1)(\auO_{E^+_v})$ for $v$ not below a prime of $S$ (\S\ref{pcx:ss-statement}). We compute $\omega(1,\kappa_v)$ place by place.

\emph{$v$ not below $S$.} The datum at $v$ is the vacuum of the self-dual lattice $\alpha\auHom(L_X,L_T)_v$ (the integral
matrices in the coordinates $\tilde x$), which $\kappa_v$
stabilises; by the last assertion of (W4) (for the splitting by $\upsilon$ of Definition~\ref{pwf:def}),
$\omega_{\auHom(X,T),v}(1,\kappa_v)$ fixes it. The character gives $(\gamma_X)_v(\det\kappa_v)=\nu_v(\det\kappa_v)^{-1}
\upsilon_v(\det\kappa_v)^{(2b+n)/2}=1$, both factors being trivial on norm-one units by the two facts above.

\emph{$v=w\bar w$ with $w\in S$.} By Proposition~\ref{pwf:W5} (first formula, $g=1$, source rank $2b+n$),
$(\omega_{\auHom(X,T),v}(1,\kappa_v)\Phi)(\xi)=\upsilon_w(\det\kappa_w)^{-(2b+n)}|\det\kappa_w|_w^{-(2b+n)/2}\Phi(\kappa_w^{-1}\xi)$,
where $\kappa_w\in\auGL_{b+n}(\auO_{E_w})$ is the component of $\kappa_v$ on $T_w$; the first two factors are $1$
($\upsilon$ unramified, $\det\kappa_w$ a unit). Since $\binom{z_0}{\kappa_w^{-1}\xi}=\mathrm{diag}(1,\kappa_w^{-1})\binom{z_0}\xi$,
$\Phi(\kappa_w^{-1}\xi)=\chi_S(\det\kappa_w)\Phi(\xi)$ (the support condition is unchanged). The character: under
$\auU(1)(E^+_v)=\{(t,t^{-1})\}\cong E_w^\times$ the element $\det\kappa_v$ is $(t,t^{-1})$ with $t=\det\kappa_w$, and
$\nu(t,t^{-1})=\nu(t/\bar t)=\chi_S(t)$ (for $t\in E_w^\times$, $t/\bar t=(t,t^{-1})$, as in the proof of
Proposition~\ref{pwf:W5}), while $\upsilon_v(\det\kappa_v)=\upsilon_w(t)\upsilon_{\bar w}(t^{-1})=1$. So $(\gamma_X)_v(\det\kappa_v)=
\chi_S(\det\kappa_w)^{-1}$, and the total factor at $v$ is $\chi_S(\det\kappa_w)^{-1}\chi_S(\det\kappa_w)=1$.

The real component of $F$ is not acted on ($\kappa$ is finite). Hence $\omega(1,\kappa)F=F$.

(b) By (U2), $\mathrm{term}_{x_0}(p,h)\ne0$ forces $\mathbf 1_{\mathrm{int}}(y)=1$ for the finite part $y=h_f^{-1}x_0Y_f$; its
columns span $h_f^{-1}M_0$, so $h_f^{-1}M_0\subset L_T$, and $h_f^{-1}M_0$ is unimodular ($h_f$ is an isometry). A
unimodular sublattice $M$ of a unimodular lattice $L$ is an orthogonal summand: for $l\in L$ the form
$m\mapsto\langle m,l\rangle$ is integral on $M$, hence equals $\langle\cdot,m_l\rangle$ for some $m_l\in M$, and
$l-m_l\perp M$; so $L=M\perp(L\cap M^\perp)$ with $L\cap M^\perp$ unimodular. Thus $L_T=h_f^{-1}M_0\perp N_1$ and
$h_fL_T=M_0\perp h_fN_1$, with $h_fN_1$ a unimodular lattice of $M_0^\perp=W_0(\auA_f)$. Self-dual hermitian lattices of
the same rank in the same space are isometric ($E/E^+$ is unramified at every finite place, so they have orthonormal
bases, Jacobowitz~\cite{Jacobowitz}, as in Step~2 of the proof of Proposition~\ref{pcu:selfdual}); so there are
$u_1\in\auU(U)(\auA_f)$ with $u_1M_0=L_U$ and $u_2\in\auU(W_0)(\auA_f)$ with $u_2h_fN_1=L_{W_0}$. Then
$\kappa=(u_1,u_2)h_f$ stabilises $L_T$, i.e.\ $\kappa\in K_f$, and $h_f=(u_1,u_2)^{-1}\kappa\in H_fK_f$.

(c) Here $y=h_f^{-1}x_0Y_f=x_0x_f=\binom{x_f}0$. By (b) $x_f$ is integral. The Gram matrix of the columns of $x_0x_f$
equals that of the columns of $x_0Y_f$ ($y_f$ is an isometry), which lies in $\auGL_b(\widehat\auO_E)$ by the definition
of a test pair; it is $x_f^*G_0x_f$ with $G_0$ the Gram matrix of the integral basis, a unit matrix. Hence
$\det x_f\cdot\overline{\det x_f}$ is a unit at every finite place of $E^+$; as $\det x_f$ is integral (on both components at a
split place), $\det x_f$ is a unit. For $y'=\binom{y'_1}{y'_2}$ ($y'_1$ of size $b\times n$, $y'_2$ of size $n\times n$),
$(y\ y')=\begin{pmatrix}x_f&y'_1\\0&y'_2\end{pmatrix}$ has determinant $\det x_f\det y'_2$, and lies in $\auGL_{b+n}(\auO_S)$
if and only if $y'_1$, $y'_2$ are integral and $\det y'_2$ is a unit (a product of two integers is a unit only if both
are). Hence $F^{(y)}(y')=\epsilon\chi_S^{-1}(\det x_f)\,\mathbf 1_{\mathrm{int}}(y'_1)\,\chi_S^{-1}(\det y'_2)
\mathbf 1_{\auGL_n(\auO_S)}(y'_2)$, which is the claim.
\end{proof}

\begin{lemma}[Coefficients at test pairs]\label{pcu:coeff}
There is $K>0$ such that for every tame $\nu$ of weight $\mu$ there is $C_\nu\in\C$ with
$|C_\nu|=K|I_0(\nu)|$ and, for every test pair $(\beta,q)$ for $U$, all $c,z,g$ and $p=n(c,z)m(q,g)$,
\[
\pcuF_\beta\,\Theta^{\mathrm L}_{X,T}(\Phi,\Phi^{\mathrm{row}}_{z_0})(p)=C_\nu\,
\pcuF_\beta\,\Theta^{\mathrm S}_{X,U}(\Lambda,\Psi^{\mathrm{col}}_{z_0})(p).
\]
\end{lemma}
\begin{proof}
\emph{Case (b).} On the small side ($V=U$, $a=b$) an $x$ with $x^*J_Ux=\beta$ is invertible and
exhibits $\beta$ as $a^*\mathrm{diag}(t')a$ at every infinite place; there is none, and the coefficient
is $0$ by (U1). On the large side, by (U2) the term of $x$ is a multiple of the slice of $\Phi$ at
$h^{-1}Y$, whose Gram matrix is $q'^*\beta q'$ ($q=s_0q'$); by Lemma~\ref{pcu:support} it vanishes
unless $\beta$ is represented by $t'$ at every infinite place, which it is not. Both sides are $0$.

\emph{Case (a): reduction to one term.} Now $\beta=J_U$. On the large side the $x\in\pcMat_{(b+n)\times
b}(E)$ with $x^*J_Tx=J_U$ are the isometric embeddings of $U$ in $T=U\oplus W_0$; by Witt's theorem
they form one orbit, that of $x_0$, with stabiliser $\Gamma_2=\auU(W_0)(E)$. On the small side they
are the elements of $\auU(U)(E)$, one orbit with trivial stabiliser. By (U1), whose condition
\eqref{pcu:eq-unfoldabs} holds on both sides by Lemma~\ref{pcu:integrable},
\[
\pcuF_\beta\Theta^{\mathrm L}(p)=\int_{\Gamma_2\backslash G}\mathrm{term}^{\mathrm L}_{x_0}(p,h)dh,\qquad
\pcuF_\beta\Theta^{\mathrm S}(p)=\int_{\auU(U)(\auA)}\mathrm{term}^{\mathrm S}_1(p,y)dy .
\]

\emph{Inner integral.} Write $q=s_0x_UY$ with the columns of $x_UY_f$ spanning a self-dual lattice; let
$K_f$ be the stabiliser of the standard lattice of $T$ in $\auU(T)(\auA_f)$, and for $g'\in G$ let
$\rho'(g')=\rho(g'_\infty)$ if $g'_f\in H_fK_f$ and $0$ otherwise; $\rho'$ is left $H$-invariant
(Lemma~\ref{pcu:period} gives $\rho$ as a function of the span of $h^{-1}x_0$, which is unchanged by
left multiplication by $H_\infty$) and non-zero near $1$. We claim: for $g'=(h_\infty,u_f)$ there is
$k'\in H$ with
\[
\int_{\auU(U)(\auA)}\int_{[\auU(W_0)]}\mathrm{term}^{\mathrm L}_{x_0}(p,(y,y')k'g')\,dy'\,dy=
\rho'(g')\,\kappa_0\,\epsilon\,\nu_{\mathrm{bc}}(d_0)\,I_0(\nu)\int_{\auU(U)(\auA)}
\mathrm{term}^{\mathrm S}_1(p,y)dy,
\]
with $d_0=\det(s_0x_U)$, $\kappa_0=\varepsilon_T\varepsilon_U^{-1}\mathrm N((d_0)_\infty)^{j}$ (with $\mathrm N$ as defined before Lemma~\ref{pcu:period}, applied to the archimedean components), and $\epsilon$ the
constant of modulus $1$ with $\Phi^{\mathrm{row}}_{z_0}=\epsilon\cdot(\text{standard datum of the frame
of }D)$. If $u_f\notin H_fK_f$, then $(y,y')g'$ has finite part in $(y_f,y'_f)u_f\notin H_fK_f$, so the
term vanishes by Lemma~\ref{pcu:finslice}(b), and both sides are $0$ ($k'=1$). If $u_f=(u_1,u_2)\kappa$ with $\kappa\in K_f$, let $(\delta_1,\delta_2)=p$ be the element
given by Lemma~\ref{pcu:period} for $h_\infty$ and $k'=((\delta_1,u_1^{-1}),(\delta_2,u_2^{-1}))$
(adjusted so that $k'g'=(\delta h_\infty,\kappa)$). The term is right $K_f$-invariant (Lemma~\ref{pcu:finslice}(a)), so $\kappa$ drops. Fix $y\in\auU(U)(\auA)$ and put
$B=y^{-1}Y$, $x_f$ its finite part. By Lemma~\ref{pcu:finslice}(c) on the large side (and (U2)
on the small side), both terms vanish unless $x_f$ is integral, and then $\det x_f$ is a unit and the large finite
slice is $\epsilon\chi_S^{-1}(\det x_f)$ times the vacuum of $(W_0,U)$ tensored with $F_1$; on the small side the constant $\chi_S^{-1}(\det\bar x_f)$ times
the vacuum. The archimedean point $\delta h_\infty$ is moved into the real slice by (M3), and the
integral over $[\auU(W_0)]$ is the left side of Lemma~\ref{pcu:period}; replacing it by the right side
and comparing scalars gives the claim. The scalar comparison uses: $\det B=\zeta_B\cdot\det x_f$ with
$|\det y_\tau|=1$ at every infinite place and $\det x_f$ a unit, so $\mathrm N(\zeta_B)=\prod_\tau|\det Y_\tau|^2$; $\chi_S^{-1}(\bar\epsilon')=
\chi_S^{-1}(\epsilon')\nu_{\mathrm{bc}}(\epsilon')$ for units $\epsilon'$ at $S$ (definition of
$\chi_S$); and the factor $\nu(\det y)^2$ between the characters $\nu\upsilon^j$ and $\nu^{-1}
\upsilon^j$ of $\auU(U)$ on the two sides; $j=n/2$ makes the powers of $\upsilon$ agree.

\emph{Conclusion.} By (U3) applied to $\Theta(h)=\mathrm{term}^{\mathrm L}_{x_0}(p,h)$ (integrable by
Lemma~\ref{pcu:integrable}(a)), the large coefficient is $R\,\kappa_0\epsilon\nu_{\mathrm{bc}}(d_0)I_0(\nu)$ times the small one,
where $R=\int_{H\backslash G}\rho'$ depends only on $\sigma$ and the measures. If $R<\infty$ put
$K=R|\kappa_0|$ and $C_\nu=R\kappa_0\epsilon\nu_{\mathrm{bc}}(d_0)I_0(\nu)$ ($|\epsilon|=|\nu_{\mathrm{bc}}|=1$).
If $R=\infty$ then (U3) gives $0$ on the left and $c_0=0$, i.e.\ $I_0(\nu)\cdot(\text{small coefficient})
=0$, so the identity holds with $K=|\kappa_0|$ and $C_\nu=\kappa_0\epsilon\nu_{\mathrm{bc}}(d_0)I_0(\nu)$.
\end{proof}
\subsection{From test pairs to all standard data}\label{pcu:ss-std-data}

Fix $\nu$ and $C=C_\nu$. The difference
$\mathcal D(\Psi)=\mathcal T(\Psi)-C\,\Theta^{\mathrm S}_{X,U}(\Lambda,\Psi)$, where $\mathcal T$ is the
linear map of Lemma~\ref{pcx:linfactor} attached to $\Phi'\mapsto\Theta^{\mathrm L}_{X,T}(\Phi,\Phi')$
(Lemma~\ref{pcx:Sact}(a)), satisfies, by Lemma~\ref{pcx:section-laws} and Lemma~\ref{pcx:Sact}:
(E1) on the space $\mathcal M$ of \S\ref{pcx:ss-match} it is linear with values automorphic functions
of moderate growth (Corollary~\ref{pcx:conv}); (E2) $\mathcal D(\Psi(g\,\cdot))(xg)=\eta(\det g)^{-b}
\mathcal D(\Psi)(x)$ for $g$ at $S$; (E3) its values are right invariant under the stabiliser of the
standard lattice away from $S$; (E4) $\mathcal D(\Psi^{\mathrm{col}}_{z'})=\Theta^{\mathrm L}(\Phi,
\Phi^{\mathrm{row}}_{z'})-C\,\Theta^{\mathrm S}(\Lambda,\Psi^{\mathrm{col}}_{z'})$ for integral $z'$
(Lemmas~\ref{pcx:std-matched}, \ref{pcx:section-laws}(b)). We call it the \emph{equivariant lift}.

\begin{proposition}[Self-dual lattices suffice]\label{pcu:selfdual}
Let $\mathcal D$ satisfy (E1)--(E4) and suppose $\pcuF_\beta\mathcal D(\Psi^{\mathrm{col}}_{z_0})
(n(c,z)m(q,g))=0$ for every test pair $(\beta,q)$ for $U$ and all $c,z,g$. Then $\mathcal D(\Psi)=0$ for
every locally constant $\Psi$ having a matched $\Phi$ for $\chi_S$; in particular
$\mathcal D(\Psi^{\mathrm{col}}_{z'})=0$ for all integral $z'$.
\end{proposition}
\begin{proof}
\emph{Step 1 (reduction to coefficients at unimodular indices).} By Lemma~\ref{pcu:gauss} below it
suffices that $\pcuF_\beta\mathcal D(\Psi^{\mathrm{col}}_{z_0})((1,g_f))=0$ for all $g_f$
trivial at $S$ and all hermitian $\beta$ with $s_0^{-1}\bar s_0^{-1}
\beta$ (the normalised index) in $\auGL_b(\auO_S)$.

\emph{Step 2 (from unimodular indices to test pairs).} For such $\beta$ the statement of
Lemma~\ref{pcu:unimod} reduces the vanishing at $(1,g_f)$ to the vanishing at points
$n(l,z)m(s_0q,g)$ of $P_D$ at which the Gram matrix $q^*\beta q$ is unimodular at all finite places,
i.e.\ $q k$ has orthonormal columns for $\beta$ at the finite places, $(qk)^*\beta(qk)=1$, for some $k\in
\auGL_b(\widehat\auO_E)$ (a unimodular hermitian lattice over $\widehat\auO_E$ has an orthonormal
basis, because $E/E^+$ is unramified at every finite place (\S\ref{sec:automorphic}), so that every unit of
$\widehat\auO_{E^+}$ is a local norm; proof: the reduction of the lattice is a non-degenerate hermitian space over the quadratic extension of residue fields, so it has a vector $\bar x$ with $\bar h(\bar x,\bar x)\ne0$ (otherwise polarisation and the surjectivity of the residue trace give $\bar h=0$); a lift $x$ has $h(x,x)$ a unit of the base ring, hence the norm of a unit since the extension is unramified, so after scaling $h(x,x)=1$; then the lattice is $\mathcal O x\perp(x^\perp\cap\text{lattice})$ with unimodular second summand, and one inducts on the rank; this is the argument at a place inert in $E$; at a place split in $E$, $\auO_{E,v}=\auO_v\times\auO_v$ and a unimodular hermitian lattice is a pair of lattices $(L',L'')$ in perfect duality under the induced bilinear pairing, so a basis of $L'$ together with its dual basis of $L''$ is an orthonormal basis; see also Jacobowitz~\cite{Jacobowitz}).

\emph{Step 3 (Landherr).} Let $(\beta,q,k)$ be as in Step 2. If $\beta$ is not represented by $t'$ at
some infinite place, $(\beta,s_0q)$ is a test pair of type (b). Otherwise $\beta$ has signature $t'$
at every infinite place and an orthonormal lattice at every finite place, so by
Lemma~\ref{pcx:landherr} there are $c\in\auGL_b(E)$ and $w\in\auU(U)(\auA)$ with $\beta=c^*J_Uc$ and
$x_\beta=c^{-1}x_Uw$, where $x_\beta=qk$ (at the finite places). The coefficient is unchanged by the
rational point $m(c,1)\in P_D(E)$: $\pcuF_{c^*J_Uc}f(p)=\pcuF_{J_U}f(m(c,1)p)$ for automorphic $f$, and
$m(c,1)n(l,z)m(s_0q,g)=n(l',z')m(s_0cq,g)$. Now $cq=x_UY$ with $Y=x_U^{-1}cq$, and the Gram matrix of the columns of $x_UY_f=cq$ for $J_U$ is
$(cq)^*J_U(cq)=k^{*-1}\bigl((cqk)^*J_U(cqk)\bigr)k^{-1}=k^{*-1}k^{-1}$ (since
$(cqk)^*J_U(cqk)=(qk)^*\beta(qk)=1$, as $\beta=c^*J_Uc$), which lies in $\auGL_b(\widehat\auO_E)$ as $k$ does; so $(J_U,s_0cq)$ is a test pair of type (a). In both cases the hypothesis gives $0$.
\end{proof}

\begin{lemma}[Gauss sums and weak approximation]\label{pcu:gauss}
Under (E1)--(E4), if $\pcuF_\beta\mathcal D(\Psi^{\mathrm{col}}_{z_0})(1,g_f)=0$ for every $g_f$ trivial at
$S$ and every $\beta$ with unimodular normalised index, then $\mathcal D(\Psi)=0$ for every locally
constant $\Psi$ with a matched function.
\end{lemma}
\begin{proof}
Let $\mathcal Q=\mathcal M\cap\ker\mathcal D$; by (E2) it is stable under $\Psi\mapsto\Psi(g\,\cdot)$. Let
$\Psi_0=\Psi^{\mathrm{col}}_{z_0}$ and, for $w\in\pcMat_b(F_S)$, $\Psi_w(y)=\Psi_0(y+L\,w\,(z_0y))$ with $L$
the $\ell$-block of the frame of $D$ at $S$ (the shear by $n(0,w)$ at $S$). For $f\in\mathcal M$: an
automorphic function vanishes if it vanishes at all $(1,g_f)$ with $g_f$ trivial at $S$ (weak
approximation at $S$ and (E3), since $\mathcal D f$ is continuous and left $\auU(X)(E)$-invariant), and
a function on $Z_D(E)\backslash Z_D(\auA)$ vanishes if all its Fourier coefficients do, the characters of
this compact group being the $\psi_\beta$, $\beta$ hermitian over $E$ (\S\ref{pcu:ss-top-setting}). For fixed
$x=(1,g_f)$ and $\beta$, the functional $\Lambda_{\beta,x}\colon f\mapsto\pcuF_\beta\mathcal Df(x)$ is linear
on $\mathcal M$ and satisfies $\Lambda_{\beta,x}(\Psi_w)=\psi_S(\operatorname{tr}(\tau_\beta w))
\Lambda_{\beta,x}(\Psi_0)$ by (E2) (the shear is the action of $n(0,w)$ at $S$, which multiplies the
coefficient by the character), where here and below $\psi_S(t)$ stands for $\psi_S(\jmath s_0\bar s_0t)$, a character of
$F_S$ trivial exactly on $\auO_S$ (as $s_0\widehat\auO_E=\alpha\widehat\auO_E$ and $\jmath\alpha\bar\alpha\widehat\auO_E$ is the dual of
$\widehat\auO_E$ for $\psi$, \S\ref{ps:ss-weil}), and $\tau_\beta=-s_0^{-1}\bar s_0^{-1}\beta$ is, up to sign, the
normalised index (at $v=w\bar w$, $w\in S$, $Z_D(E^+_v)\cong\pcMat_b(E_w)$ through the $w$-component, and
$\psi_{\beta,v}(z)=\psi_w(-\jmath\operatorname{tr}(z_w\beta))$ by \eqref{pcu:eq-psibeta}; only whether $\det\tau_\beta$ is a
unit matters below). So: every $f\in\mathcal M$ killed
by all functionals $\Lambda$ with $\Lambda(\Psi_w)=\psi_S(\operatorname{tr}(\tau w))\Lambda(\Psi_0)$ for all
$w$ and $\Lambda(\Psi_0)=0$ when $\tau\in\auGL_b(\auO_S)$ (hypothesis) lies in $\mathcal Q$. Let $\pi\colon
\auO_S\to\pcR$ be reduction modulo the product of the primes of $S$ and $\eta$ the character of
$\pcR^\times$ with $\eta\circ\pi=\chi_S^{-1}$ on units, extended by $0$. Consider
$f=\sum_{r\in\pcMat_b(\pcR)}\eta(\det r)\Psi_{\varpi^{-1}\tilde r}$, $\tilde r$ a lift of $r$. For a
functional as above, $\Lambda(\Psi_w)=\Lambda(\Psi_0)$ for integral $w$ (as $\Psi_w=\Psi_0$ then), so
$\tau$ is integral if $\Lambda(\Psi_0)\ne0$; and
$\Lambda(f)=\Lambda(\Psi_0)\sum_r\eta(\det r)\psi_S(\varpi^{-1}\operatorname{tr}(\tau\tilde r))$, a
Gauss sum on $\pcMat_b(\pcR)$ which vanishes when $\det\tau$ is not a unit ($\eta$ primitive, as
$\chi_S$ is tame with non-trivial tame part at each $w\in S$); when $\det\tau$ is a unit,
$\Lambda(\Psi_0)=0$. So $f\in\mathcal Q$, hence $f(g\,\cdot)\in\mathcal Q$ for $g=-\varpi^{-1}1_N$. A
direct computation shows that every non-zero term of
$f(g\,\cdot)$ at $y$ equals $\Psi^{\mathrm{col}}_{z'_0}(y)$ for the complementary block $z'_0$ of the
inverse frame, so $f(g\,\cdot)=n(\cdot)\Psi^{\mathrm{col}}_{z'_0}$ with $n(y)\in\Z_{\ge0}$ the number of
non-vanishing terms; $n$ is locally constant, and $n>0$ at the point $x_0$ of
Lemma~\ref{pcx:generate} for $s=z'_0$ (the term $r=1$ is $1$ there). The proof of
Lemma~\ref{pcx:generate} goes through verbatim with $p=\Psi^{\mathrm{col}}_{s}$ replaced by $np$:
the claim becomes $\Psi(k^{-1}x_0)\,n(dky)p(dky)=n(dky)|p(dky)|\Psi(y)$, the average
$I(y)=\int n(dky)|p(dky)|\,dk$ is still constant on the support of $\Psi$ by right invariance of
$dk$, and $I(x_0)>0$ by continuity. Hence every locally constant $\Psi$ with a matched function lies
in $\mathcal Q$.
\end{proof}

\begin{lemma}[Unimodular Gram points suffice]\label{pcu:unimod}
Let $f=\mathcal D(\Psi^{\mathrm{col}}_{z_0})$ and $\beta$ with unimodular normalised index. If
$\pcuF_\beta f$ vanishes at every point $n(l,z)m(s_0q,g)\in P_D(\auA)$ whose Gram matrix $q^*\beta q$
is unimodular at all finite places, then $\pcuF_\beta f(1,g_f)=0$ for all $g_f$ trivial at $S$.
\end{lemma}
\noindent The proof is in \S\ref{pcu:ss-unimod}.

\begin{proof}[Proof of Proposition~\ref{pcx:P-top}]
Let $\Phi$, $\Lambda$ be the vectors of \S\ref{pcu:ss-period}, $x=\Phi$, $x'=\Lambda$; they have finite
moments and the zero term of $x'$ along $D^{\mathrm s}$ (the $b$ planes) is $w$. Let $K$, $C_\nu$ be as
in Lemma~\ref{pcu:coeff}. For each $\nu$, the difference $\mathcal D$ built with $C_\nu$ satisfies the
hypothesis of Proposition~\ref{pcu:selfdual}: by (E4) and the linearity of $\pcuF_\beta$ (both
integrals are automorphic of moderate growth), $\pcuF_\beta\mathcal D(\Psi^{\mathrm{col}}_{z_0})=
\pcuF_\beta\Theta^{\mathrm L}(\Phi,\Phi^{\mathrm{row}}_{z_0})-C_\nu\pcuF_\beta\Theta^{\mathrm S}(\Lambda,
\Psi^{\mathrm{col}}_{z_0})$, which vanishes at test pairs by Lemma~\ref{pcu:coeff}. Hence
$\Theta^{\mathrm L}_{X,T}(\Phi,\Phi^{\mathrm{row}}_{z'})=C_\nu\Theta^{\mathrm S}_{X,U}(\Lambda,
\Psi^{\mathrm{col}}_{z'})$ for every integral $z'$: this is (b) of $\mathrm{RC}$ with $c=C_\nu$. For (a):
$|c|^2=K^2|I_0(\nu)|^2=K^2L_*(\pcm,\nu,j)\,\pcxstr_{W_0}(\nu;z)$ by Definition~\ref{pcx:def-ratio}
($\operatorname{rk}W_0=2j$), and $L_*(\pcm,\nu,\operatorname{rk}T-\operatorname{rk}X/2)=L_*(\pcm,\nu,j)$;
so (a) holds with the constant $K^{-2}$.
\end{proof}

\begin{remark}[Hypotheses of Proposition~\ref{pcx:P-top}]\label{pcu:rem-top}
No sign condition is used: the slices of $\Phi$ vanish outside the real class of $U$ by construction
(Lemma~\ref{pcu:support}), which is what case (b) of Lemma~\ref{pcu:coeff} needs. The only hypotheses
are that $W_0$ has even rank and splits off an adapted plane (for Lemma~\ref{pcx:const}), and
$\mu_\tau\ge1$.
\end{remark}

\subsection{Unfolding along $P_D$: proofs of (U1) and (U2)}\label{pcu:ss-open}

\emph{Adelic slices.} For a datum $f$ of $\auHom(X,V)$ (all places) and $x\in\pcMat_{a\times b}(\auA_E)$
let $f[x]$ be the datum of $\auHom(W_0,V)$ obtained by restricting the mixed model of $D$ at every
place (Schr\"odinger in the $\auHom(\ell,V)$ variable at every place, as in (W2), (W2$'$) of
\S\ref{ps:ss-weil}) to the fibre over $x$: $f[x]=(\psrho(x\,e^*)f)[0]$, $[0]$ the zero term. The
theta distribution of $\auHom(X,V)$ is $\theta(f)=\sum_{x\in\pcMat_{a\times b}(E)}\theta_{W_0}(f[x])$:
in the mixed model the rational points of $\auHom(X,V)$ decompose as $\auHom(\ell,V)(E)\oplus
\auHom(W_0,V)(E)\oplus\auHom(\ell',V)(E)$, the Schr\"odinger variable is summed over the first and
the remaining sum is $\theta_{W_0}$ by Poisson summation in $\auHom(\ell',V)$ (absolutely convergent
for Schwartz data). Put $\mathrm{term}_x(p,h)=\theta_{W_0}((\omega(p,h)f)[x])$.

\begin{lemma}[Levi and radical laws]\label{pcu:laws}
Let $\chi_V(u)=\alpha(u/\bar u)\upsilon(u)^{-a}$ for the source character $\alpha$. Then
\begin{align*}
(\omega_{\alpha,\beta\upsilon^b}(m(q,g),h)f)[x]&=\chi_V(\det q)^{-1}|\det q|^{a/2}\,
\omega_{\alpha,\beta}(g,h)\bigl(f[h^{-1}xq]\bigr),\\
(\omega(n(c,z),1)f)[x]&=\psi_{x^*J_Vx}(z)\,\psrho(x c^*J_{W_0})
\bigl(f[x]\bigr),\qquad\psi_{x^*J_Vx}(z)=\psi\bigl(-\tfrac12\jmath\operatorname{tr}(z\,x^*J_Vx)\bigr),
\end{align*}
with $\psi_\beta$ as in \eqref{pcu:eq-psibeta}.
\end{lemma}
\begin{proof}
These are the Schr\"odinger-model formulas (W2), (W2$'$) for the Siegel-type parabolic of the
polarisation $\auHom(\ell,V)\oplus\auHom(\ell',V)$, tensored with the action on $\auHom(W_0,V)$. For
$m(q,g)$: together with $h$ it acts on the Schr\"odinger variable $x\in\auHom(\ell,V)$ by
$x\mapsto hxq^{-1}$, so the new function at $x$ is the old one at $h^{-1}xq$, with the factor of (W2) for $\mathrm{GL}(\auHom(\ell,V))$, which is
$\upsilon$ of its determinant $\det(q)^{a}\det(h)^{-b}$ times $|\det q|^{a/2}$; the factor
$\upsilon(\det h)^{-b}$ is absorbed by the shift $\beta\upsilon^b\to\beta$ in the target character
and the factor $\alpha(\det q/\overline{\det q})$ comes from the $\auU(X)$-character restricted along
$\ell\oplus\ell'$, giving $\chi_V(\det q)^{-1}$; $g$ acts on $\auHom(W_0,V)$.

\emph{The radical.} Write $\mathsf W=\auHom(X,V)=I\oplus I^*\oplus\mathsf W'$ with $I^*=\auHom(\ell,V)$ (maps vanishing on
$\ell'\oplus W_0$), $I$ the maps vanishing on $\ell\oplus W_0$, $\mathsf W'=\auHom(W_0,V)$; the form of $\mathsf W$ is
$\langle w_1,w_2\rangle=\operatorname{tr}(w_1^\dagger w_2)$, and in the frames $w^\dagger=G_X^{-1}M^*J_V$ for a map with matrix $M$, $G_X$
the Gram matrix of the frame of $X$. The mixed model is $\auS(I^*;\psHS(\mathsf W'))$ with $\psrho(u_0)$ ($u_0\in I^*$) the
translation, $\psrho(i)$ ($i\in I$) the multiplication by $[u,i]$, and $\mathsf W'$ acting on the values, as in (P2) and
Lemma~\ref{pwf:dist}. The element $(g,1)$ acts on $\mathsf W$ by $w\mapsto wg^{-1}$. Let $u\in I^*$ with $u(e)=x$. Since
$n(c,z)^{-1}=n(-c,z^*)$, $u\,n(c,z)^{-1}=u+d(u)+b(u)$ with $d(u)\in\mathsf W'$, $d(u)(w)=-x\,c^*J_0w$, and $b(u)\in I$,
$b(u)(f)=x\,z^*$; also $\mathsf W'\ni w'\mapsto w'n(c,z)^{-1}=w'+C(w')$ with $C(w')\in I$, and $I$ is fixed.

(i) \emph{The centre, $c=0$.} Then $d=0$, $z^*=-z$, and $n(0,z)$ acts as the element $n(b)$ of the Siegel parabolic of
$I\oplus I^*$, tensored with the identity on $\mathsf W'$ (Lemma~\ref{pwf:sum}). By Lemma~\ref{pwf:siegel},
$\omega(n(b))\varphi(u)=\beta(u,bu)^{-1}\varphi(u)$, where $\beta(u,bu)=\psi(\frac12\jmath\langle bu,u\rangle)$. In the frames,
$G_X^{-1}$ exchanges the $e$- and $f$-blocks, the matrix of $u$ is $(x\ 0\ 0)$ and that of $bu$ is $(0\ {-xz}\ 0)$, so
$\langle bu,u\rangle=\operatorname{tr}(G_X^{-1}(bu)^*J_Vu)=\operatorname{tr}(-z^*x^*J_Vx)=\operatorname{tr}(z\,x^*J_Vx)$. Hence
$\omega(n(0,z),1)$ is multiplication by $\psi(-\frac12\jmath\operatorname{tr}(z\,x^*J_Vx))=\psi_{x^*J_Vx}(z)$.

(ii) \emph{The element $n(c,z_2)$, $z_2=-\frac12c^*J_0c$.} Here $z_2$ is hermitian. We claim that $A\colon F\mapsto
(u\mapsto\psrho(-d(u))F(u))=(u\mapsto\psrho(xc^*J_0)F(u))$ satisfies $A\psrho(y)A^{-1}=\psrho(n(c,z_2)y)$ for all $y\in\mathsf W$; we
use $\psrho(y)^{-1}=\psrho(-y)$ and $\beta(y,y')=\psi(\frac12\jmath\langle y',y\rangle)$, and that $\psi(\jmath r)=1$ for $r\in\auA_{E^+}$. For
$y=i\in I$ both sides are multiplication by $[u,i]$. For $y=w'\in\mathsf W'$, $A\psrho(w')A^{-1}F(u)=[-d(u),w']\psrho(w')F(u)$, while
$\psrho(w'+C(w'))=\psrho(C(w'))\psrho(w')$ ($\langle w',C(w')\rangle=0$) multiplies by $[u,C(w')]$ and then applies $\psrho(w')$; and
$[u,C(w')]=[d(u),w']^{-1}$, because $n(c,z_2)$ preserves $[\,,]$ and $u\perp w'$ (expand $[u+d(u)+b(u),w'+C(w')]=1$ using
$I\perp\mathsf W'$, $I$ isotropic). For $y=u_0\in I^*$, with $d_0=d(u_0)$, $b_0=b(u_0)$: on one side
$A\psrho(u_0)A^{-1}F(u)=\psrho(-d(u))\psrho(d(u)+d_0)F(u+u_0)=\beta(d(u),d_0)^{-1}\psrho(d_0)F(u+u_0)$; on the other,
$\psrho(u_0+d_0+b_0)F(u)=\beta(b_0,u_0)^{-1}[u,b_0]\psrho(d_0)F(u+u_0)$. Now $\beta(b_0,u_0)=\psi(\frac12\jmath\operatorname{tr}(x_0^*J_Vx_0z_2))=1$
($x_0=u_0(e)$, and the trace of a product of two hermitian matrices is real); with $P=x^*J_Vx_0$, $M=c^*J_0c$ one finds
$\langle u,b_0\rangle=\operatorname{tr}(Pz_2)=-\frac12\operatorname{tr}(PM)$, so $[u,b_0]=\psi(\frac12\jmath(\overline{\langle u,b_0\rangle}-\langle u,b_0\rangle))
=\psi(-\jmath\langle u,b_0\rangle)=\psi(\frac12\jmath\operatorname{tr}(PM))$, and $\langle d_0,d(u)\rangle=\operatorname{tr}(x_0^*J_Vx\,M)=
\operatorname{tr}(P^*M)=\overline{\operatorname{tr}(PM)}$, so $\beta(d(u),d_0)^{-1}=\psi(-\frac12\jmath\,\overline{\operatorname{tr}(PM)})=
\psi(\frac12\jmath\operatorname{tr}(PM))$. The two sides agree, which proves the claim. By (P1), $\omega(n(c,z_2),1)$ is a scalar
multiple of $A$; the scalar is $1$, because by (W2) (with $\chi_I=1$ on unipotent elements, and $n(c,z_2)$ acting trivially on
$I^\perp/I=\mathsf W'$) $\omega(n(c,z_2),1)$ commutes with evaluation at $u=0$, as $A$ does.

(iii) \emph{General $n(c,z)$.} $n(c,z)=n(0,z_1)\,n(c,z_2)$ with $z_1=z+\frac12c^*J_0c$ skew-hermitian, by the composition law. By
(i) and (ii), $\omega(n(c,z),1)F(u)=\psi_{x^*J_Vx}(z_1)\,\psrho(xc^*J_0)F(u)$, and $\psi_{x^*J_Vx}(z_1)=\psi_{x^*J_Vx}(z)$ since $\psi$ is
trivial on $\jmath\auA_{E^+}$ (\S\ref{pcu:ss-top-setting}). The characters $\alpha$, $\beta$ are trivial on unipotent elements. This is
the second formula.
\end{proof}

\begin{corollary}[Local form]\label{pcu:laws-local}
The same formulas hold at each place $v$ of $E^+$ separately: for a datum $f_v$ of $\auHom(X,V)_v$ in the mixed model of $D$
at $v$ and $x\in\pcMat_{a\times b}(E_v)$,
$(\omega_v(m(q,g),h)f_v)[x]=\chi_{V,v}(\det q)^{-1}|\det q|_v^{a/2}\,\omega_v(g,h)(f_v[h^{-1}xq])$ and
$(\omega_v(n(c,z),1)f_v)[x]=\psi_v(-\frac12\jmath\operatorname{tr}(z\,x^*J_Vx))\,\psrho_v(xc^*J_{W_0})(f_v[x])$, with the local components of the
characters. At an infinite place $\tau$, $f_\tau[x]=\mathrm{sl}_D(x)f_\tau$.
\end{corollary}
\begin{proof}
The proof of Lemma~\ref{pcu:laws} is place by place: (W2), (W2$'$) are identities between local operators, the adelic
mixed model, the slices $f[x]$ and the Weil representation are restricted tensor products of the local ones, and the
characters are products of their local components. Each step of that proof is therefore an identity at the place $v$,
and the adelic formulas are the tensor products of the local ones. At $\tau$ the mixed model is that of
\S\ref{pcu:ss-mixed} and the slice is $\mathrm{sl}_D$ by definition.
\end{proof}

\begin{lemma}[Surjectivity of the standard data]\label{pcu:rowsurj}
\textup{(a)} If $\Phi^{\mathrm{row}}_z(x)\ne0$ ($z\in\pcMat_{b\times N}(\auO_S)$, $x\in\pcMat_{a\times N}(F_S)$, $N=a+b$), then at
every place of $S$ the matrix $x$ has rank $a$, i.e.\ $x\colon F_w^N\to F_w^a$ is surjective. \textup{(b)} The same holds
for $\Psi^{\mathrm{col}}_{1_n}$ on $\pcMat_{n\times n}(F_S)$: it is supported on $\auGL_n(\auO_S)$. \textup{(c)} If
$(\Psi^{\mathrm{col}}_z)^{\mathsf t}(x)=\Psi^{\mathrm{col}}_z(x^{\mathsf t})\ne0$ ($x\in\pcMat_{b\times N}$), then $x$ has rank $b$ at
every place, i.e.\ is surjective. Hence, as $S\cap\bar S=\emptyset$ (beginning of \S\ref{pct:sec}),
Proposition~\ref{pcx:decay}(a) applies with $P=S$ to the data $[y,\Phi^{\mathrm{row}}_z]_S$ of $\auHom(X,T)$ and
$[y,\Psi^{\mathrm{col}}_{1_n}]_S$ of $\auHom(W_0,W_0)$, and with $P=\bar S$ to $[y,(\Psi^{\mathrm{col}}_z)^{\mathsf t}]_{\bar S}$ of
$\auHom(X,U)$.
\end{lemma}
\begin{proof}
(a) By Definition~\ref{pcx:def-std}, $\binom zx\in\auGL_N(\auO_S)$; the $a$ rows of $x$ are rows of an invertible matrix,
hence linearly independent over each $E_w$, $w\in S$. (b) By Definition~\ref{pcx:def-std} with $z=1_n$. (c) By
Definition~\ref{pcx:def-std}, $zx^{\mathsf t}\in\auGL_b(\auO_S)$, so $x^{\mathsf t}$, hence $x$, has rank $b$.
\end{proof}

\begin{proof}[Proof of (U2)]
Apply Lemma~\ref{pcu:laws} to $p=n(c,z)m(q,g)$: \[\mathrm{term}_x(p,h)=\psi_{x^*J_Vx}(z)\chi_V(\det q)^{-1}|\det q|^{a/2}\theta_{W_0}(\psrho(x'c^*J_{W_0})\omega(g,h)f[h^{-1}xq]).\]
With $f=[\phi,\Phi]$ and $x'q=s_0Y$, the slice $f[h^{-1}s_0Y]$ is the tensor product of the real
slice $\mathrm{sl}_D(h^{-1}Y)\phi$ (the factor $s_0$ is trivial at infinity), the slice of $\Phi$ at
the $S$-component $y$ of $h^{-1}Y$ (by definition of the Schr\"odinger model at $S$), and at the
other finite places the slice of the standard vacuum, which is $\varepsilon_V\mathbf 1_{\mathrm{int}}$
times the vacuum of $\auHom(W_0,V)$ (the vacuum of the self-dual lattice $\alpha\auHom(L_X,L_V)$ is the
indicator of the lattice, which factors along the unimodular splitting; $\varepsilon_V$ is the
product of the local volume normalisations, independent of $x$). Here the slice point is the finite part of
$h^{-1}s_0Y$, whose coordinate is $y=h^{-1}Y$ by Lemma~\ref{pgx:dict}(d); so its integrality, and the
slice of $\Phi$ at $S$, are read on $y$, as in (U2). This is (U2).
\end{proof}

\begin{proof}[Proof of (U1)]
Fix $p=n(c,z)m(q,g)$. Let $I\subset\auHom(X,V)$ be the isotropic subspace of the maps vanishing on
$\ell^\perp=\ell\oplus W_0$, so that $I^\perp$ is the space of maps vanishing on $\ell$,
$I^\perp/I=\auHom(W_0,V)$, and $w\mapsto w|_\ell$ identifies $\auHom(X,V)(E)/I^\perp(E)$ with
$\pcMat_{a\times b}(E)$. We claim $\theta_{W_0}(F[x])=c_x(F)$ in the notation of Lemma~\ref{pz:absgrowth},
i.e.\ $F[x]=\psZ_I\psrho(w)F$ for any $w$ with $w|_\ell=x$. Let $I^*$ be the maps vanishing on
$\ell'\oplus W_0$; it is isotropic, in duality with $I$, and $\auHom(X,V)=(I\oplus I^*)\perp\auHom(W_0,V)$
with $I^*=\auHom(\ell,V)$. The mixed model of $D$ is the model $\psHS(I^*)\hat\otimes\psHS(\auHom(W_0,V))$
in which $\psrho(I^*)$ translates the variable $u\in I^*(\auA)$, $\psrho(I)$ multiplies by the characters
$u\mapsto\psi(\langle u,i\rangle)$, and $\psrho(\auHom(W_0,V))$ acts on the values. The map
$F\mapsto F[0]$ (evaluation at $u=0$) satisfies $(\psrho(i)F)[0]=F[0]$ for $i\in I(\auA)$ and
$(\psrho(w)F)[0]=\psrho(w)(F[0])$ for $w\in\auHom(W_0,V)(\auA)$, i.e.\ the intertwining properties of
$\psZ_I$ in (W2). These determine the map up to a scalar: a continuous map with them factors through the
quotient of $\psHS(I^*)$ by the span of the $(\psi(\langle\cdot,i\rangle)-1)\varphi$, whose dual consists of
the distributions $T$ with $(\psi(\langle\cdot,i\rangle)-1)T=0$ for all $i$; such $T$ is supported at $u=0$, and
at the archimedean places no derivative of $\delta_0$ survives (if $T=\sum_{|\alpha|\le m}c_\alpha
\partial^\alpha\delta_0$ with some $c_\alpha\ne0$, $|\alpha|=m\ge1$, multiplying by $\psi(\langle\cdot,i\rangle)-1$,
which vanishes to first order at $0$ with arbitrary linear part, gives a non-zero distribution of order
$m-1$), so $T$ is a multiple of $\delta_0$. Hence $F[0]=c\,\psZ_IF$, and $c=1$ because both satisfy the
expansion of $\theta$: $\theta(F)=\sum_{u\in I^*(E)}\theta_{W_0}((\psrho(u)F)[0])$ (Poisson summation in the
mixed model, as recalled before Lemma~\ref{pcu:laws}) and (W3) with $W(E)/I^\perp(E)=I^*(E)$, and
$\theta\ne0$. As $F[x]=(\psrho(xe^*)F)[0]$ with $xe^*\in I^*(\auA)$ restricting to $x$ on $\ell$, this is the
claim; the expansion $\theta=\sum_xc_x$ is (W3).

\emph{Step 1: absolute convergence and moderate growth of the absolute series.} For every finite datum
$f_0$ there is $Y$, depending on $f_0$ and on $p$, with
\begin{equation}\label{pcu:eq-absmod}
\sum_{x}\bigl|c_x\bigl(\omega(p,h)(\psrho_\infty(z')\Psi\otimes f_0)\bigr)\bigr|\le Y\,\mathrm{ht}(h)^{4s}\,
\|\Psi\|_{2s}\qquad(h\in\auU(V)(\auA),\ z'\in\auHom(X,V)_\infty,\ \Psi\in\pzF_{2s}),
\end{equation}
$s$ as in (R2) for $\auHom(X,V)$. Indeed, for $\gamma\in\auU(V)(E)$, $c_x\circ\omega(1,\gamma)=c_{\gamma^{-1}x}$
up to a factor of absolute value $1$ ((W1), (W2): $\gamma$ stabilises $I$, and the characters are trivial on
rational points), so the left side is unchanged when $h$ is replaced by $\gamma h$; by
Lemma~\ref{pcx:reduce} we may assume $h_f$ in a fixed compact set and $\|h_\infty\|\le C\,\mathrm{ht}(h)$.
Then $\omega(p,h)(\psrho_\infty(z')\Psi\otimes f_0)$ is a phase times $\psrho_\infty(z'')\omega_\infty
(p_\infty,h_\infty)\Psi\otimes\omega_f(p_f,h_f)f_0$ ((W1)), the finite vector lies up to phases in a finite
set, and Lemma~\ref{pz:absgrowth}(c) followed by Lemma~\ref{pz:growth}(a) with index $2s$ (``Other
spaces'' paragraph before Lemma~\ref{pz:absgrowth}) gives \eqref{pcu:eq-absmod} with
$Y=C'(2DC\|p_\infty\|)^{4s}$. In particular $\sum_x|\mathrm{term}_x(p,h)|<\infty$ for every $h$, and
$|\mathrm{term}_x(n(0,z)p,h)|=|\mathrm{term}_x(p,h)|$ by Lemma~\ref{pcu:laws}.

\emph{Step 2: the coefficient as an integral of the $\beta$-part.} By Proposition~\ref{pcx:decay}(a) (its
hypothesis holds for the standard data, Lemma~\ref{pcu:rowsurj}) $|\theta(\omega(n(0,z)p,h)F)|\le
C\,\mathrm{ht}(h)^{-N_0}$ uniformly for $z$ in a compact fundamental set of $Z_D(E)\backslash Z_D(\auA)$,
integrable on $[\auU(V)]$ by (R1); so Fubini applies to the double integral over $Z_D(E)\backslash
Z_D(\auA)\times[\auU(V)]$. For fixed $h$, Step~1 and Lemma~\ref{pcu:laws} with $c=0$, $q=1$
($\mathrm{term}_x(n(0,z)p,h)=\psi_{x^*J_Vx}(z)\mathrm{term}_x(p,h)$, $\psi_\beta$ as in \eqref{pcu:eq-psibeta}) allow integrating term by term against $\psi_\beta(z)^{-1}$ over the compact
group, which keeps exactly the $x$ with $x^*J_Vx=\beta$. Hence
\[
\pcuF_\beta\int_{[\auU(V)]}\theta(\omega(\cdot,h)F)\,dh\,(p)=\int_{[\auU(V)]}\theta_\beta(h)\,dh,\qquad
\theta_\beta(h)=\sum_{x^*J_Vx=\beta}\mathrm{term}_x(p,h),
\]
the integral on the right absolutely convergent ($|\theta_\beta(h)|\le\sup_z|\theta(\omega(n(0,z)p,h)F)|$).

\emph{Step 3: unfolding.} $\mathrm{term}_{\gamma x}(p,\gamma h)=\mathrm{term}_x(p,h)$ for $\gamma\in
\auU(V)(E)$. If there is no $x$ with $x^*J_Vx=\beta$, or all their terms vanish identically, then
$\theta_\beta\equiv0$ and the coefficient is $0$ by Step~2. Otherwise, $\beta$ being invertible, the $x$
with $x^*J_Vx=\beta$ form one $\auU(V)(E)$-orbit (Witt's theorem: an isometric embedding of the
nondegenerate space $(E^b,\beta)$ into $V$ is unique up to $\auU(V)(E)$). Suppose that
\begin{equation}\label{pcu:eq-unfoldabs}
\int_{\mathrm{Stab}(x_0)(E)\backslash\auU(V)(\auA)}|\mathrm{term}_{x_0}(p,h)|\,dh<\infty
\qquad\Bigl(\text{equivalently }\int_{[\auU(V)]}\sum_{x^*J_Vx=\beta}|\mathrm{term}_x(p,h)|\,dh<\infty\Bigr)
\end{equation}
for one (hence every) $x_0$ of the orbit; the two integrals are equal by Tonelli (non-negative terms,
$\sum_{x}|\mathrm{term}_x(p,h)|=\sum_{\gamma\in\mathrm{Stab}(x_0)(E)\backslash\auU(V)(E)}|\mathrm{term}_{x_0}(p,\gamma h)|$).
Then Fubini gives $\int_{[\auU(V)]}\theta_\beta=\int_{\mathrm{Stab}(x_0)(E)\backslash\auU(V)(\auA)}
\mathrm{term}_{x_0}(p,h)\,dh$, which is (U1).
\end{proof}

\begin{remark}[What (U1) does not assert]\label{pcu:rem-U1}
Condition \eqref{pcu:eq-unfoldabs} is not automatic. For the orbit of $x_0$ on the large side the
stabiliser is $\auU(W_0)$, and by Lemma~\ref{pcn:slicechar} the inner integral over $[\auU(W_0)]$ is a theta
integral for the pair $(W_0,W_0)$ of equal ranks $n$; $W_0\supset\pcxH$ is isotropic, so Weil's
criterion for absolute convergence of such integrals for arbitrary Schwartz--Bruhat data (which needs the
rank of the source to be smaller than the rank minus the Witt index of the target) fails. What makes
\eqref{pcu:eq-unfoldabs} true below is the finite datum $\Phi^{\mathrm{row}}_{z_0}$: its slices along $x_0$
contain the factor $\chi_S^{-1}\circ\det$ supported on invertible matrices, to which
Proposition~\ref{pcx:decay}(a) applies. Moderate growth of $\sum_x|\mathrm{term}_x|$ (Step~1) and decay of the
signed sum $\theta_\beta$ (Step~2) do not by themselves give \eqref{pcu:eq-unfoldabs}.
\end{remark}

\begin{lemma}[Polynomial volume of the orbits]\label{pcu:orbitvol}
Let $m\ge b$, $J$ a hermitian $m\times m$ matrix $\mathrm{diag}(\pm1)$ and $\beta_0$ an invertible hermitian
$b\times b$ matrix, at each infinite place $\tau$. Let $\mathrm{Fr}=\prod_\tau\pcMat_{m\times b}(\C)$ with
Lebesgue measure and euclidean norm $|\cdot|$, $Q(x)=(x_\tau^*Jx_\tau)_\tau$, and $\Omega=Q^{-1}(\beta_0)$.
Suppose $\Omega\ne\emptyset$. Then $\Omega$ is a closed submanifold and a single orbit of $G=\prod_\tau\auU(J)$
acting by left multiplication; the orbit map identifies $G/G_{x_1}$ ($x_1\in\Omega$) with $\Omega$, and for
every $K>\dim_\R\mathrm{Fr}$ the $G$-invariant measure $\mu$ on $\Omega$ satisfies
$\int_\Omega(1+|x|)^{-K}d\mu(x)<\infty$.
\end{lemma}
\begin{proof}
On the open set $\mathrm{Fr}^\circ$ where every $Q(x)_\tau$ is invertible, $Q$ is a submersion onto an open
subset of $\mathcal H=\prod_\tau\mathrm{Herm}_b(\C)$: $dQ_x(\xi)=\xi^*Jx+x^*J\xi$, and $\xi=\frac12x\,Q(x)^{-1}\eta$
gives $\eta$. So $\Omega$ is a closed submanifold, and for each $\beta\in Q(\mathrm{Fr}^\circ)$ the fibre carries
the Leray measure $\mu_\beta=dx/dQ$ (the quotient of the Lebesgue measures), with
$\int_{\mathrm{Fr}^\circ}f\,dx=\int_{\mathcal H}\int_{Q^{-1}(\beta)}f\,d\mu_\beta\,d\beta$ for measurable $f\ge0$.
Left multiplication by $h\in G$ preserves $Q$ and $dx$ ($|\det h_\tau|=1$), hence each $\mu_\beta$. By Witt's
theorem for hermitian forms over $\C$, $\Omega$ is one $G$-orbit; a transitive smooth action of a Lie group
with countably many components identifies $G/G_{x_1}$ with $\Omega$, and $G$, $G_{x_1}$ (a product of
unitary groups) are unimodular, so the invariant measure on $\Omega$ is unique up to a scalar and $\mu$ is a
multiple of $\mu_{\beta_0}$. For $c\in\prod_\tau\auGL_b(\C)$, right multiplication $x\mapsto xc$ maps
$Q^{-1}(\beta)$ to $Q^{-1}(c^*\beta c)$ and multiplies $dx$ and $d\beta$ by positive constants depending
continuously on $c$, so $\int_{Q^{-1}(c^*\beta_0c)}f\,d\mu_{c^*\beta_0c}=J(c)\int_\Omega f(xc)\,d\mu_{\beta_0}(x)$ with
$J$ continuous and positive. The map $c\mapsto c^*\beta_0c$ is a submersion at $c=1$ (derivative
$\eta\mapsto\eta^*\beta_0+\beta_0\eta$, onto with $\eta=\frac12\beta_0^{-1}\zeta$), so there are an open
neighbourhood $N$ of $\beta_0$ and a continuous $c\colon N\to\prod\auGL_b(\C)$ with $c(\beta)^*\beta_0c(\beta)=\beta$
and $\|c(\beta)\|\le2$. Put $I=\int_\Omega(1+|x|)^{-K}d\mu_{\beta_0}$. As $1+|xc|\le2(1+|x|)$ for $\|c\|\le2$,
\[
\infty>\int_{\mathrm{Fr}}(1+|x|)^{-K}dx\ge\int_N J(c(\beta))\int_\Omega(1+|xc(\beta)|)^{-K}d\mu_{\beta_0}(x)\,d\beta
\ge2^{-K}I\int_NJ(c(\beta))\,d\beta,
\]
and the last integral is positive; so $I<\infty$.
\end{proof}

\begin{lemma}[Short representatives]\label{pcu:section}
In the setting of \S\ref{pcu:ss-top-setting} ($J=\mathrm{diag}(t'\oplus s)$, $x_0=\binom{1_b}0$, $J_U=\mathrm{diag}(t')$),
for every $x\in\mathrm{Fr}$ with $x_\tau^*Jx_\tau=J_U$ for all $\tau$ there is $h\in\auU(T)_\R$ with
$h^{-1}x_0=x$ and $\|h\|\le(1+\sqrt n)(1+|x|)$ (operator norm, maximum over $\tau$).
\end{lemma}
\begin{proof}
Fix $\tau$ and drop it. Let $P=xJ_Ux^*J$ and $Q=1-P$; then $x^*JQ=x^*J-J_U^2x^*J=0$, so $Q$ maps into
$x^\perp=\ker(x^*J)$, $Q$ is the identity on $x^\perp$, and $\|Q\|\le1+|x|^2$. Let $z_1$ be an $(n+b)\times n$
matrix whose columns are a euclidean orthonormal basis of $x^\perp$ and $G=z_1^*Jz_1$. For $w\in x^\perp$
put $w'=QJw\in x^\perp$; since $Q^*=1-JxJ_Ux^*$ and $x^*Jw=0$, $w'^*Jw=w^*JQ^*Jw=w^*w$, and $|w'|\le\|Q\||w|$.
So for every $\xi\in\C^n$ there is $\xi'$ with $|\xi'^*G\xi|\ge|\xi||\xi'|/\|Q\|$, i.e.\ $\|G^{-1}\|\le\|Q\|$. Write
$G=u^*Du$ with $u$ unitary and $D$ real diagonal; $T=\mathrm{span}(x)\perp x^\perp$ and Sylvester's law give
$x^\perp$ the signs of $W_0$, so after permuting, $z=z_1u^*|D|^{-1/2}$ satisfies $z^*Jz=\mathrm{diag}(s)$ and
$|z|\le\sqrt n\,\|Q\|^{1/2}\le\sqrt n(1+|x|)$. The matrix $k=(x\ z)$ has $k^*Jk=J$, so $k\in\auU(J)$, $kx_0=x$;
put $h=k^{-1}=Jk^*J$, so $\|h\|=\|k\|\le|x|+|z|\le(1+\sqrt n)(1+|x|)$.
\end{proof}

\begin{lemma}[Moments under the Weil representation]\label{pcu:moments}
Let $\mathcal F$ be the Fock space of a real hermitian space $\C^m$, realised as in
Lemma~\ref{pcx:superpos}, and let $k$ be a unitary operator on $\mathcal F$ with $k\psrho(w)k^{-1}=
\psrho(\sigma w)$ for a real symplectic map $\sigma$ of $\C^m$ (for instance $\omega_\infty(g,h)$, by (W1)).
For every $N\ge0$ there is $C_N$, depending only on $N$, $m$ and the normalisations, with
$\|kv\|_N\le C_N\|\sigma\|^{2N+2m}\|v\|_N$ for every $v$ with $\|v\|_N<\infty$.
\end{lemma}
\begin{proof}
By Stone--von Neumann a unitary operator with this intertwining property is determined by $\sigma$ up to a
scalar of absolute value $1$; write $M(\sigma)$ for any choice. Let $E=k^{-1}e$ and $E(w)=\langle\psrho(w)e,E\rangle$
its function in the realisation of Lemma~\ref{pcx:superpos}. \emph{Reduction.} We first show $|E(w)|\le
\exp(-|w|^2/(1+\|\sigma\|^2))$. Write the
polar decomposition $\sigma^{-1}=\kappa P$ with $P=(\sigma^{-\mathsf t}\sigma^{-1})^{1/2}$ positive symmetric and
$\kappa$ orthogonal; both are symplectic, so $\kappa$ commutes with the complex structure. The vacuum is, up
to scalars, the only vector killed by the annihilation operators of the complex structure, a family that
$M(\kappa)$ permutes; so $M(\kappa)e\in\C e$, $M(\kappa)^{-1}\psrho(w)e\in\C\,\psrho(\kappa^{-1}w)e$, and
$|E(w)|=|\langle\psrho(\kappa^{-1}w)e,M(P)e\rangle|$, with $|\kappa^{-1}w|=|w|$. $P=\exp X$ with $X$ symmetric and
anticommuting with the complex structure $i$, so $i$ exchanges the eigenspaces of $X$ for $\mu$ and $-\mu$ and
there is a euclidean orthonormal basis $(e_j,ie_j)$ of $\C^m$ with $Pe_j=\lambda_je_j$, $P(ie_j)=\lambda_j^{-1}ie_j$,
$\lambda_j\ge1$, $\lambda_j\le\|P\|=\|\sigma^{-1}\|=\|\sigma\|$ (for a symplectic $\sigma$, $\sigma^{-1}=
\mathsf J^{-1}\sigma^{\mathsf t}\mathsf J$ with $\mathsf J$ the orthogonal matrix of the symplectic form in a basis $(e_j,ie_j)$). A unitary change of basis is again of the form $M(\kappa')$ and
does not change the moduli, so we may assume this basis is the standard one. Then $\mathcal F$ is the tensor
product of the Fock spaces of the lines $\C e_j$, $e$ and $M(P)$ factor (by the same uniqueness), and in the
Schr\"odinger model of each line, with $\psrho(p,q)f(t)=e^{2\pi iqt}f(t-p)$ up to a phase, the vacuum is the Gaussian
$\gamma(t)=2^{1/4}e^{-\pi t^2}$ and $M(P_j)f(t)=\lambda_j^{-1/2}f(t/\lambda_j)$. \emph{The normalisation.} The commutation
relation $\psrho(p,0)\psrho(0,q)=e^{-2\pi ipq}\psrho(0,q)\psrho(p,0)$ of this model, compared with $\psrho(w)\psrho(y)=
e^{\kappa(w,y)}\psrho(y)\psrho(w)$, $\kappa(w,y)=2i\Iml\operatorname{tr}(y^*w)$ (Lemma~\ref{pcx:superpos}), forces the
coordinate of $(p,q)$ on the line $\C e_j$ to be $\sqrt\pi(p+iq)$, so $|w|^2=\pi(p^2+q^2)$. For a normalised Gaussian
$\gamma_a(t)=(2a)^{1/4}e^{-\pi at^2}$ ($a>0$), completing the square gives $|\langle\psrho(p,q)\gamma_a,\gamma_a\rangle|=
\exp(-\frac\pi2(ap^2+q^2/a))$, while in the model of Lemma~\ref{pcx:superpos} $|\langle\psrho(w)e,e\rangle|=e(w)=
e^{-|w|^2/2}=\exp(-\frac\pi2(p^2+q^2))$; the unitary isomorphism carrying $e$ to $\gamma_a$ (up to a phase) preserves these
moduli, so $a=1$. Then a Gaussian integral gives
\[
|\langle\psrho(p,q)\gamma,M(P_j)\gamma\rangle|=\Bigl(\frac{2\lambda_j}{1+\lambda_j^2}\Bigr)^{1/2}
\exp\Bigl(-\pi\,\frac{p^2+\lambda_j^2q^2}{1+\lambda_j^2}\Bigr)\le\exp\Bigl(-\pi\,\frac{p^2+q^2}{1+\lambda_j^2}\Bigr)
=\exp\Bigl(-\frac{|w_j|^2}{1+\lambda_j^2}\Bigr)
\]
($w_j$ the component of $w$ on $\C e_j$; the exponent is obtained by writing $(t-p)^2+t^2/\lambda_j^2=(1+\lambda_j^{-2})(t-t_0)^2+
p^2/(1+\lambda_j^2)$, and the prefactor is $\le1$). Taking the product over $j$, with $\lambda_j\le\|\sigma\|$, gives the bound on
$|E(w)|$. Hence $\|E\|_N\le\int(1+|w|)^Ne^{-|w|^2/(1+\|\sigma\|^2)}dw\le C'_N\|\sigma\|^{N+2m}$ ($\|\sigma\|\ge1$).
\emph{The estimate.} As $\|E\|_0<\infty$, $E=\int E(w)\psrho(w)e\,dw$ in $\mathcal F$ (Lemma~\ref{pcx:superpos}(b)).
Then $(kv)(x)=\langle\psrho(x)e,kv\rangle=\langle k^{-1}\psrho(x)e,v\rangle=\langle\psrho(\sigma^{-1}x)E,v\rangle$, and
$\psrho(\sigma^{-1}x)\psrho(w)$ is a scalar of absolute value $1$ times $\psrho(\sigma^{-1}x+w)$; so
$|(kv)(x)|\le\int|E(w)||v(\sigma^{-1}x+w)|dw$. Substituting $u=\sigma^{-1}x+w$ ($|\det\sigma|=1$) and using
$1+|\sigma(u-w)|\le\|\sigma\|(1+|u|)(1+|w|)$,
$\|kv\|_N\le\|\sigma\|^N\|E\|_N\|v\|_N\le C'_N\|\sigma\|^{2N+2m}\|v\|_N$.
\end{proof}

\begin{lemma}[Integrability at test pairs of type (a)]\label{pcu:integrable}
Let $(J_U,q)$ be a test pair of type (a) for $U$ and $p=n(c,z)m(q,g)$.
\begin{enumerate}
\item[\textup{(a)}] Let $V=T$ and $F=[\phi,\Phi^{\mathrm{row}}_{z_0}]_S$ with $\phi$ a Fock vector of $\auHom(X,T)_\R$
with finite moments. Then $\int_{\Gamma_2\backslash G}|\mathrm{term}_{x_0}(p,h)|\,dh<\infty$ ($G=\auU(T)(\auA)$,
$\Gamma_2=\auU(W_0)(E)$, the stabiliser of $x_0$).
\item[\textup{(b)}] Let $V=U$ and $F=[\lambda,(\Psi^{\mathrm{col}}_{z_0})^{\mathsf t}]_{\bar S}$ with $\lambda$ a Fock
vector of $\auHom(X,U)_\R$ with finite moments. Then $\int_{\auU(U)(\auA)}|\mathrm{term}_1(p,y)|\,dy<\infty$.
\end{enumerate}
In particular \eqref{pcu:eq-unfoldabs} holds for $\beta=J_U$ in Lemma~\ref{pcu:coeff} (with $\phi=\Phi$,
$\lambda=\Lambda$ of \S\ref{pcu:ss-period}) and in type (a) of the proof of Proposition~\ref{pct:vantop}
(with $\phi=v$).
\end{lemma}
\begin{proof}
Write $q=s_0x_UY$ as in the proof of Lemma~\ref{pcu:coeff}, so that, in the notation of (U2), $Y$ is the
$b\times b$ block and $h^{-1}(Y;0)=h^{-1}x_0Y$. Constants $C$ may depend on $p$, $\phi$, $\lambda$.

(a) \emph{Integration in stages.} With $H_2=\auU(W_0)(\auA)$, Tonelli and the quotient integral formula
give $\int_{\Gamma_2\backslash G}|\mathrm{term}_{x_0}(p,h)|dh=\int_{H_2\backslash G}j(h)\,dh$ with
$j(h)=\int_{[\auU(W_0)]}|\mathrm{term}_{x_0}(p,s'h)|ds'$, a measurable function of $H_2h$ with values in
$[0,\infty]$. By Lemma~\ref{pcu:finslice}(a), (b), $|\mathrm{term}_{x_0}(p,\cdot)|$ is
right invariant under $K_f$, and $\mathrm{term}_{x_0}(p,\cdot)$ vanishes at $h=(h_\infty,u_f)$ unless $u_f\in H_fK_f$
($H=\auU(U)(\auA)\times H_2$). If $u_f=(u_1,u_2)\kappa$ with $\kappa\in K_f$, left invariance of $j$ under $H_2$
and right invariance under $K_f$ give $j(h)=j(h_0)$ with $h_0=(k_\infty,(y_f,1))$, $y_f=u_1$, for any
$k_\infty\in H_{2,\infty}h_\infty$. In
$H_2\backslash G=(H_{2,\infty}\backslash G_\infty)\times(H_{2,f}\backslash G_f)$ the archimedean coset of $h_0$
is determined by $x=k_\infty^{-1}x_0\in\Omega=\{x:x_\tau^*Jx_\tau=J_U\}$, and by Lemma~\ref{pcu:section} we may
take $\|k_\infty\|\le C(1+|x|)$.

\emph{The term at $s'h_0$.} Let $s'\in H_2$. The finite part of $(s'h_0)^{-1}x_0Y$ is $x_0y_f^{-1}Y_f$, and
its archimedean part is $xY_\infty$ ($s'$ fixes $x_0$). By (U2) and Lemma~\ref{pcu:finslice}(c) (with $h_f=(y_f,s'_f)$; the finite slice does not depend on $s'_f$),
$\mathrm{term}_{x_0}(p,s'h_0)=0$ unless $x_f=y_f^{-1}Y_f$ is integral, and then
\[
|\mathrm{term}_{x_0}(p,s'h_0)|=C_0\,\bigl|\theta\bigl(\psrho(t)\,\omega(g,s'h_0)(\mathrm{sl}_D(xY_\infty)\phi\otimes
f_U\otimes F_1)\bigr)\bigr|,
\]
with $C_0$ independent of $s'$, $x$, $y_f$, $t=x'c^*J_{W_0}$ fixed, $f_U$ the vacuum datum of $(W_0,U)$ and $F_1$
the datum $\Psi^{\mathrm{col}}_1$ of $(W_0,W_0)$ at $S$ (vacuum elsewhere). The set $\mathcal Y$ of $y_f\in\auU(U)(\auA_f)$
with $y_f^{-1}Y_f$ integral is compact ($y_f^{-1}\in\pcMat_b(\widehat\auO_E)Y_f^{-1}$ is bounded, and unitary
groups are closed). Put $A_x=\omega_\infty(1,k_\infty)\,\mathrm{sl}_D(xY_\infty)\phi$, a vector of
$\auHom(W_0,T)_\R$. By Lemma~\ref{pcu:moments} and Lemma~\ref{pcu:slicedecay}, for every $M'$ and $K$,
\[
\|A_x\|_{M'}\le C\|k_\infty\|^{2M'+2m}\|\mathrm{sl}_D(xY_\infty)\phi\|_{M'}\le C(1+|x|)^{2M'+2m}(1+|xY_\infty|)^{-K}
\|\phi\|_{M'+K}\le C'(1+|x|)^{2M'+2m-K}
\]
($m$ the complex dimension of $\auHom(W_0,T)_\R$; $|xY_\infty|\ge|x|/\|Y_\infty^{-1}\|$). Since $s'$ and
$(1,(y_f,1))$ are block diagonal for $T=U\oplus W_0$, $\omega(g,s'h_0)(\mathrm{sl}_D(xY_\infty)\phi\otimes f_U\otimes
F_1)=\omega(g,s')\bigl(A_x\otimes\omega_f(1,y_f)f_U\otimes F_1\bigr)$. Expand $A_x=\int A_x(w)\psrho(w_1)e\boxtimes
\psrho(w_2)e\,dw$ along $\auHom(W_0,T)=\auHom(W_0,U)\oplus\auHom(W_0,W_0)$ (Lemma~\ref{pcx:superpos}(b); the
integral converges in the space of data of class $s$, on which $\theta\circ\psrho(t)\circ\omega(g,s')$ is
continuous) and apply Lemma~\ref{pcu:product} to each coherent state:
\[
|\mathrm{term}_{x_0}(p,s'h_0)|\le C_0\int|A_x(w)|\,\bigl|\theta\bigl(\psrho(t)\omega(g,1)(\psrho(w_1)e\otimes
\omega_f(1,y_f)f_U)\bigr)\bigr|\,\bigl|\theta^{W_0,W_0}(g,s';\psrho(w_2)e\otimes F_1)\bigr|\,dw .
\]
The first theta factor is bounded uniformly in $w_1$ and $y_f\in\mathcal Y$: write $\psrho(t)\omega(g,1)=
\omega(g,1)\psrho(g^{-1}t)$; the finite data $\psrho_f((g^{-1}t)_f)\omega_f(1,y_f)f_U$ lie, up to factors of absolute
value $1$, in a finite set (finite data are fixed by an open subgroup, and $\mathcal Y$ is compact), and
Lemma~\ref{pcx:thetagrowth} is uniform in the archimedean translation. For the second, $F_1$ is supported on
$\auGL_n(\auO_S)$, i.e.\ on surjective matrices (Lemma~\ref{pcu:rowsurj}(b)), so Proposition~\ref{pcx:decay}(a) gives, for $N_0$ with
$\int_{[\auU(W_0)]}\mathrm{ht}^{-N_0}<\infty$ (R1), constants $B,M$ with $|\theta^{W_0,W_0}(g,s';\psrho(w_2)e\otimes
F_1)|\le C\|\psrho(w_2)e\|_M\mathrm{ht}(s')^{-N_0}\mathrm{ht}(g)^B\le C(1+|w_2|)^M\mathrm{ht}(s')^{-N_0}$. So, with
$M'=\max(M,2s)$,
\[
j(h_0)\le C\,\mathbf 1_{\mathcal Y}(y_f)\,\|A_x\|_{M'}\int_{[\auU(W_0)]}\mathrm{ht}(s')^{-N_0}ds'\le C''\,
\mathbf 1_{\mathcal Y}(y_f)\,(1+|x|)^{2M'+2m-K}.
\]

\emph{Conclusion.} The finite cosets with $j\ne0$ lie in the image of the compact set
$\{(y_f,1):y_f\in\mathcal Y\}K_f$, of finite measure. On the archimedean factor, $H_{2,\infty}\backslash G_\infty$ is identified with $\Omega$ by
$h\mapsto h^{-1}x_0$ (Lemma~\ref{pcu:orbitvol} with $m=b+n$, $\beta_0=J_U$), the quotient measure being a
multiple of $\mu$; choosing $K>2M'+2m+\dim_\R\mathrm{Fr}$, Lemma~\ref{pcu:orbitvol} gives
$\int_\Omega(1+|x|)^{2M'+2m-K}d\mu<\infty$. Hence $\int_{H_2\backslash G}j<\infty$.

(b) The stabiliser of $1$ is trivial. By (U2) (factor $\mathbf 1_{\mathrm{int}}$), $\mathrm{term}_1(p,y)=0$ unless $x_f=y_f^{-1}Y_f$ is
integral, i.e.\ $y_f\in\mathcal Y$, and then
$|\mathrm{term}_1(p,y)|=C_0|\theta(\psrho(t)\omega(g,y)(\mathrm{sl}_D(y_\infty^{-1}Y_\infty)\lambda\otimes f^{(x_f)}))|$, where
$f^{(x_f)}$ is the finite datum of (U2); as the datum at $\bar S$ is locally constant with compact support and $x_f$
runs over a compact set, $f^{(x_f)}$ takes only finitely many values. Writing $\psrho(t)\omega(g,y)=\omega(g,y)\psrho((g,y)^{-1}t)$ and using
Lemma~\ref{pcx:thetagrowth} (uniform in the archimedean translation; the finite data form a finite set up to
phases since $y_f\in\mathcal Y$), $|\mathrm{term}_1(p,y)|\le C\|\mathrm{sl}_D(y_\infty^{-1}Y_\infty)\lambda\|_s
\,\mathrm{ht}(y)^{A}$. Here $\|\cdot\|_s\le c_1\|\cdot\|_{2s}$ (Lemma~\ref{pcx:superpos}(b)) and
Lemma~\ref{pcu:slicedecay} give $\|\mathrm{sl}_D(y_\infty^{-1}Y_\infty)\lambda\|_s\le C(1+|y_\infty^{-1}|)^{-K}
\|\lambda\|_{2s+K}$; and $\mathrm{ht}(y)\le C(1+\|y_\infty\|)$ for $y_f\in\mathcal Y$ (the lattice at $y$ lies in a
fixed lattice $N^{-1}\auO_E^b$, whose non-zero vectors have archimedean length $\ge c>0$, and
$|(y^{-1}\gamma)_\infty|\ge|\gamma_\infty|/\|y_\infty\|$), with $\|y_\infty\|=\|y_\infty^{-1}\|\le|y_\infty^{-1}|$
($y_\infty^{-1}=J_Uy_\infty^*J_U$). So $\int_{\auU(U)(\auA)}|\mathrm{term}_1|\le C\,\mathrm{vol}(\mathcal Y)
\int_{\auU(U)_\R}(1+|y|)^{A-K}dy$, finite for $K$ large by Lemma~\ref{pcu:orbitvol} with $m=b$,
$\beta_0=J_U$, $\Omega=\auU(U)_\R$ (inversion preserves Haar measure).
\end{proof}

\subsection{Singular forms are governed by unimodular coefficients: proof of
Lemma~\ref{pcu:unimod}}\label{pcu:ss-unimod}

This is an adelic unitary form of a theorem of Freitag (\emph{Singular modular forms and theta
relations}, Lecture Notes in Math.\ 1487, Springer 1991; see also J.-S.~Li, Amer.\ J.\ Math.\ 119
(1997)); we give the proof. Throughout, $\Sigma$ is a set of finite places, $\Sigma^c=\Sigma\cup\bar
\Sigma$, $\mathcal O_\Sigma$ the adeles integral at $\Sigma^c$, $K_\Sigma$ the integral points of
$\auU(X)$ supported on $\Sigma^c$, $\delta=s_0\bar s_0$ (a finite idele of $E^+$), and for $\gamma\in\pcMat_b(\auA_E)$
\emph{Siegel's bracket} $\beta[\gamma]=\delta^{-1}\gamma^*\beta\gamma$. At a finite place $v$ the form $\psi$ of
\S\ref{ps:ss-weil} is $\psi_\Q\circ\operatorname{Tr}_{E/\Q}$, and $\jmath\alpha\bar\alpha\auO_{E,v}$ is the dual $\mathfrak d_{E,v}^{-1}$ of
$\auO_{E,v}$ for it (this is the choice of $\alpha$ in \S\ref{ps:ss-weil}: $\widehat\auO_E$ is its own dual for
$(z,u)\mapsto\psi(\jmath\alpha\bar\alpha\bar zu)$; it holds at every finite place, including those above $2$, $d$
and $q_1$, so no further fractional ideal enters; note that on $\auA_{E^+}$ one has
$\psi(t)=\psi_\Q(\operatorname{Tr}_{E^+/\Q}(2t))$, which cancels the factor $\frac12$ in \eqref{pcu:eq-psibeta}; and
$s_0\bar s_0/\alpha\bar\alpha$ is a unit idele, so $\delta\widehat\auO_{E^+}=\alpha\bar\alpha\widehat\auO_{E^+}$); as $E/E^+$ is unramified at $v$, there is a unit $\theta_v\in\auO_{E,v}^\times$ with
$\bar\theta_v=-\theta_v$ ($\theta_v=\zeta-\bar\zeta$ for a unit $\zeta$ generating the residue extension if $v$ is inert,
$\theta_v=(1,-1)$ if $v$ splits), so the integral skew-hermitian matrices at $v$ are $\theta_v$ times the integral
hermitian ones. A point $g$ \emph{acts on the
frame of $D$ by $\gamma$} if $g\,e=e\,\gamma$ ($e$ the frame of $\ell$, adelically).

\emph{Setting.} $D_1$ gives an adapted splitting $X=(\langle u\rangle\oplus\langle u'\rangle)\perp X'$
(one plane, the splitting $A$), and $X'=(\ell\oplus\ell')\perp W_2$ (the splitting $B$) with
$\ell$ the frame of $D$; in Lemma~\ref{pcu:unimod} this is $X=\pcxH\oplus(\pcxH^b\oplus W)$ with
the plane of $D_1$ (by (M4) and Lemma~\ref{pcx:adapted} the frames can be chosen so). For a row
$a\in\pcMat_{1\times b}(\auA_E)$ let $n_+(a)\in\auU(X)(\auA)$ be the unipotent element $u\mapsto u-e\,a^*$
(with the forced action on $u'$ and $\ell'$, identity on $W_2$), and $n_-(a)$ the same with $u,u'$
exchanged. For $f$ automorphic and $\beta$ with $\det\beta\ne0$ put
\[
f^+(g)=\int_{\pcMat_{1\times b}(E)\backslash\pcMat_{1\times b}(\auA_E)}\pcuF_\beta f(n_-(b)g)\,db,\qquad
\varphi^+_g(a)=f^+(n_+(a)g),
\]
and $f^-,\varphi^-_g$ the same with $u,u'$ exchanged.

\begin{lemma}[Partial means]\label{pcu:pm}
Let $f$ be automorphic and $\det\beta\ne0$, and put $\chi_\beta(a,b)=\psi_\beta(b^*a-a^*b)=\psi(-\jmath\,a\beta b^*)$.
\textup{(a)} $n_+(a)n_-(b)=n(0,b^*a-a^*b)\,n_-(b)n_+(a)$, so
$\pcuF_\beta f(n_-(b)n_+(a)g)=\chi_\beta(a,b)^{-1}\,\pcuF_\beta f(n_+(a)n_-(b)g)$. \textup{(b)} $f^+$ is continuous and
$f^+(n_-(b)g)=f^+(g)$; $\varphi^+_{n_+(a)g}(x)=\varphi^+_g(x+a)$; $\varphi^-_{n_+(a)g}(b)=\chi_\beta(a,b)^{-1}
\varphi^-_g(b)$. \textup{(c)} If $\varphi^+_{g'}=c\,\varphi^+_g$ then $\pcuF_\beta f(n_-(b)n_+(a)g')=c\,
\pcuF_\beta f(n_-(b)n_+(a)g)$ for all $a,b$, and $\varphi^-_{g'}=c\,\varphi^-_g$. In particular
$\varphi^+_g=0$ implies $\pcuF_\beta f(g)=0$ and $\varphi^-_g=0$.
\end{lemma}
\begin{proof}
(a) The isometry conditions force $n_+(a)$: $u\mapsto u-ea^*$, $f_k\mapsto f_k+a_ku'$, and the
identity on $u'$, $e$, $W_2$; and $n_-(b)$: $u'\mapsto u'-eb^*$, $f_k\mapsto f_k+b_ku$, the identity on $u$, $e$, $W_2$.
Both products agree on $u,u',e,W_2$, and $n_+(a)n_-(b)f_k-n_-(b)n_+(a)f_k=\sum_ie_i(\bar b_ia_k-\bar a_ib_k)$, so
they differ by the element $n(0,z)$, $z=b^*a-a^*b$ (skew-hermitian), of the centre $Z_D$; it commutes with $n_\pm$, and
$\pcuF_\beta f(n(0,z)p)=\psi_\beta(z)\pcuF_\beta f(p)$ (substitute $z'+z$ in \eqref{pcu:eq-psibeta}). For the
formula for $\chi_\beta$: $\operatorname{tr}(z\beta)=s-\bar s$ with $s=a\beta b^*$, $\bar\jmath=-\jmath$, so
$-\frac12\jmath(s-\bar s)=-\frac12(\jmath s+\overline{\jmath s})$, and $\psi(\bar t)=\psi(t)$. (b) $b\mapsto\pcuF_\beta f(n_-(b)g)$ is continuous
and invariant under rational $b$ ($n_-(b)$ rational fixes the frame of $D$, so is in $P_D(E)$ with
trivial Levi part on $\ell$ and $\pcuF_\beta f$ is left invariant under it); the mean over the compact
quotient is translation invariant. The second identity is the definition; the third follows from
(a) and the first. (c) For rational $a_0$, by (a) and the left invariance under $n_+(a_0)$,
$\varphi^+_g(a_0+a)$ is the Fourier coefficient at the character $b\mapsto\chi_\beta(a_0,b)$ of the
periodic function $b\mapsto\pcuF_\beta f(n_-(b)n_+(a)g)$. As $\beta$ is invertible, $a_0\beta$ runs through
$\pcMat_{1\times b}(E)$, so $a_0\mapsto\chi_\beta(a_0,\cdot)=\psi(-\jmath\,a_0\beta\,\overline{\cdot}^{\mathsf t})$ runs through all characters of $\pcMat_{1\times b}(E)\backslash
\pcMat_{1\times b}(\auA_E)$, so a continuous periodic function is determined by these coefficients
(Fourier uniqueness on the compact group). This gives the first claim; exchanging the order of the
shifts by (a) (the factor does not depend on the point) and averaging over $a$ gives the second.
\end{proof}

\begin{lemma}[Singularity]\label{pcu:sing}
Let $\mathcal U$ be a non-empty set of finite places. If $f^+(\varsigma g)=f^+(g)$ for every
$\varsigma\in N_A(E_{\mathcal U})$ (the shears $\varsigma=1+uc$, $cu=0$, at the places of $\mathcal U$) and every $g$ fixing the frame of $D$ off
$\Sigma^c$, then $f^+(rg)=f^+(g)$ for every $r$ in the radical $N_A(\auA)$ of the stabiliser of $u$
whose components off $\Sigma^c$ are trivial, for such $g$.
\end{lemma}
\begin{proof}
Call $g$ \emph{admissible} if it fixes the frame of $D$ off $\Sigma^c$. $N_A$ is a two-step unipotent
group: its centre is $Z_A\cong\{z\in\auA_E:z+\bar z=0\}$ (the elements $u'\mapsto u'+zu$) and
$N_A/Z_A\cong X'$ (the elements $n_A(c)$, $c\in X'$). The function $f^+$ itself need not be left invariant
under $N_A(E)$; we pass to its mean over $Z_A$.

\emph{The group $Z'$.} Let $I_1=\langle u\rangle\oplus\ell$ and let $Z'$ be the group of the isometries
$1+\nu$ of $X$ with $\nu=0$ on $I_1^\perp=I_1\oplus W_2$ and $\nu(X)\subseteq I_1$. The product of two such
$\nu$ is $0$, so $Z'$ is abelian and $(1+\nu)(1+\nu')=1+\nu+\nu'$; it contains $Z_D$, the elements $n_-(b)$
and $Z_A$, and it is the direct product of these three subgroups (write $\nu(f_k)=\sum_ie_iz_{ik}+b_ku$,
$\nu(u')=-eb^*+\zeta u$). Let $\psi'(1+\nu)=\psi_\beta(z)$, a character of $Z'(\auA)$ trivial on
$Z'(E)$, on the $n_-(b)$ and on $Z_A$, and put
\[
f^0(g)=\int_{Z'(E)\backslash Z'(\auA)}f(z'g)\,\psi'(z')^{-1}\,dz'=\int_{Z_A(E)\backslash Z_A(\auA)}f^+(\zeta g)\,d\zeta
\]
(probability measures; the two expressions agree by the definitions of $\pcuF_\beta$ and $f^+$ and Fubini's
theorem on the compact quotient). It is continuous. In particular $f^+$ is left invariant under $Z_A(E)$
(the elements of $Z_A$ commute with $Z_D$ and with the $n_-(b)$).

\emph{$f^0$ is left invariant under $N_A(E)$.} Let $\gamma\in N_A(E)$. It stabilises $I_1$ (it fixes $u$
and maps $e_k$ to $e_k-\langle c,e_k\rangle u$), hence $I_1^\perp$, so it normalises $Z'$; and
$\gamma^{-1}(1+\nu)\gamma=1+\gamma^{-1}\nu\gamma$ with $(\gamma^{-1}\nu\gamma)(f_k)=\gamma^{-1}\nu(f_k)$ (as
$\gamma f_k-f_k\in\langle u\rangle\subseteq\ker\nu$), whose components along $e$ are those of $\nu(f_k)$ (as
$\gamma^{-1}$ changes elements of $I_1$ by multiples of $u$). So $\psi'(\gamma^{-1}z'\gamma)=\psi'(z')$, and the
substitution $z'\mapsto\gamma z'\gamma^{-1}$, which preserves $Z'(E)$ and the measure, gives
$f^0(\gamma g)=\int f(\gamma\,\gamma^{-1}z'\gamma\,g)\psi'(z')^{-1}dz'=\int f(z'g)\psi'(z')^{-1}dz'=f^0(g)$, by the
left invariance of $f$ under $\gamma$.

\emph{Consequences of the hypothesis.} Let $g$ be admissible. For $\zeta\in Z_A(\auA)$ the point $\zeta g$ is
admissible ($Z_A$ fixes $e$), and $Z_A(E_{\mathcal U})\subseteq N_A(E_{\mathcal U})$; so the continuous function
$\zeta\mapsto f^+(\zeta g)$ on $Z_A(\auA)$ is invariant under $Z_A(E)$ and, by hypothesis, under
$Z_A(E_{\mathcal U})$. By strong approximation for $\mathbb G_a$ ($\mathcal U\ne\emptyset$, so $E$ is dense in the
adeles away from $\mathcal U$ and $E+E_{\mathcal U}$ is dense in $\auA_E$; the same holds for the vector spaces $Z_A$
and $X'$) it is constant. Hence
\[
f^+(g)=f^0(g)\qquad\text{for every admissible }g .
\]
Moreover, for $\varsigma\in N_A(E_{\mathcal U})$, as $Z_A$ is central in $N_A$ and $\zeta g$ is admissible,
$f^0(\varsigma g)=\int f^+(\varsigma\zeta g)\,d\zeta=\int f^+(\zeta g)\,d\zeta=f^0(g)$.

\emph{Conclusion.} Let $r\in N_A(\auA)$ be trivial off $\Sigma^c$ and $g$ admissible; then $rg$ is admissible.
The set $N_A(E)\cdot N_A(E_{\mathcal U})$ is dense in $N_A(\auA)$: for $c\in X'(E)$, $c_{\mathcal U}\in X'(E_{\mathcal U})$,
$n_A(c)n_A(c_{\mathcal U})=n_A(c+c_{\mathcal U})z$ with $z\in Z_A(\auA)$, and one approximates first the image of $r$
in $X'(\auA)$ and then, multiplying by elements of $Z_A(E)$ and $Z_A(E_{\mathcal U})$, the central component. Choose
$\gamma_i\in N_A(E)$, $\varsigma_i\in N_A(E_{\mathcal U})$ with $\gamma_i\varsigma_i\to r$. By the two invariances of
$f^0$ and its continuity, $f^0(rg)=\lim f^0(\gamma_i\varsigma_ig)=\lim f^0(\varsigma_ig)=f^0(g)$. As $g$ and $rg$
are admissible, $f^+(rg)=f^0(rg)=f^0(g)=f^+(g)$.
\end{proof}

\begin{definition}[Straight points and dual index]\label{pcu:def-straight}
A \emph{straight point} is $P=(g,g_0,t,t',\gamma)$ with $g$ fixing the frame of $D$ off $\Sigma^c$ and
equal to $g_0$ on $\Sigma^c$, where $g_0$ acts on $(u,e)$ by $t\oplus\gamma$ and on $(u',e)$ by
$t'\oplus\gamma$, with $t'^*t=1$. Its \emph{dual index} $\mathrm{di}(P)$ is the index in
$\widehat\auO_{E,\Sigma^c}$ of $\{x:\ x\det(\beta[\gamma])^{-1}\in\widehat\auO_{E,\Sigma^c}\}$ (the
denominator on $\Sigma^c$ of $\det\beta[\gamma]^{-1}$), a positive integer when $\beta[\gamma]$ is
integral on $\Sigma^c$.
\end{definition}

\begin{lemma}\label{pcu:bracket-int}
Let $f$ be automorphic and right $K_\Sigma$-invariant. If $\pcuF_\beta f(Y)\ne0$ for a $Y$ equal on
$\Sigma^c$ to a point acting on the frame of $D$ by $\gamma$, then $\beta[\gamma]$ is integral on
$\Sigma^c$. If moreover $\varphi^+_P\ne0$ for a straight point $P$, then $\beta[\gamma_P]$ is integral on
$\Sigma^c$.
\end{lemma}
\begin{proof}
For integral $M$, the element $n(0,M)$ taken on $\Sigma^c$ is in $K_\Sigma$; moving it across $Y$ turns it
into $n(0,\gamma M\gamma^*)$ on $\Sigma^c$, which multiplies $\pcuF_\beta f(Y)$ by $\psi_\beta(\gamma M\gamma^*)=
\psi(-\frac12\jmath\operatorname{tr}(Mh))$ restricted to $\Sigma^c$, $h=\gamma^*\beta\gamma$ (see \eqref{pcu:eq-psibeta}). Right
invariance makes this $1$ for all integral skew-hermitian $M$. Write $M=\theta_vH$ at $v\in\Sigma^c$, $H$ integral
hermitian, and $c=\jmath\theta_v\in E^+_v$. The diagonal entries give
$\psi(-\frac12cH_{ii}h_{ii})=\psi_\Q(\operatorname{Tr}_{E^+_v/\Q_p}(-cH_{ii}h_{ii}))=1$ for all $H_{ii}\in\auO_{E^+,v}$, and the
pairs of off-diagonal entries give $\psi(-\frac12c\operatorname{Tr}_{E/E^+}(H_{ij}h_{ji}))=\psi(-cH_{ij}h_{ji})=1$ for all
$H_{ij}\in\auO_{E,v}$. Hence $c\,h_{ji}\in\mathfrak d_{E,v}^{-1}$ for all $i,j$ (for $i=j$: $ch_{ii}\in
\mathfrak d_{E^+,v}^{-1}\subseteq\mathfrak d_{E,v}^{-1}$), i.e.\ $h\in c^{-1}\jmath\alpha\bar\alpha\,
\pcMat_b(\auO_{E,v})=\delta\,\pcMat_b(\auO_{E,v})$ ($\theta_v$ a unit): $\beta[\gamma]=\delta^{-1}h$ is integral on $\Sigma^c$.
For the second claim, $\varphi^+_P(a)\ne0$ gives by Lemma~\ref{pcu:pm}(b) and the definition some $b$
with $\pcuF_\beta f(n_-(b)n_+(a)g)\ne0$, and $n_\pm$ fix the frame of $D$, so the point acts by
$\gamma_P$ on $\Sigma^c$.
\end{proof}

\begin{lemma}[Iwasawa decomposition for the flag]\label{pcu:iwasawa}
\textup{(a)} For every $g\in\auU(X)(\auA)$ there is $\kappa$, integral at all finite places, such that
$g\kappa$ stabilises the flag $\langle u\rangle\subset\langle u\rangle\oplus\ell$, i.e.\ acts on $(u,e)$ by a
block upper triangular matrix $\binom{\alpha\ \beta'}{0\ \gamma}$. \textup{(b)} If $g'$ acts on $(u,e)$ by
such a matrix, then $g'=r\,m$ with $r\in N_A(\auA)$ (radical of the stabiliser of $\langle u\rangle$)
and $m$ in the Levi $\auGL_1\times\auU(X')$ acting on $u$ by $\alpha$, on $u'$ by $\bar\alpha^{-1}$, and on $e$
by $\gamma$.
\end{lemma}
\begin{proof}
(a) At each finite place $v$, let $\Lambda$ be the unimodular lattice of $X_v$ and
$Y_1=g_v^{-1}\langle u\rangle\subset Y_2=g_v^{-1}(\langle u\rangle\oplus\ell)$ (isotropic). The
sublattices $Y_i\cap\Lambda$ are primitive and isotropic; in a unimodular hermitian lattice over the
local ring $\auO_{E,v}$ every primitive isotropic flag extends to a hyperbolic basis adapted to it
(lift a basis of $Y_1\cap\Lambda$, then of $Y_2\cap\Lambda$ modulo it, pick dual vectors in $\Lambda$ by
unimodularity, and orthogonalise; Jacobowitz~\cite{Jacobowitz}, also at split places; here again
$E/E^+$ is unramified at $v$), so there is
$\kappa_v\in\auU(\Lambda)$ carrying the standard flag lattices $\auO u\subset\auO u\oplus\auO e$ to
$Y_i\cap\Lambda$. For almost all $v$, $g_v\in\auU(\Lambda)$ and $\kappa_v=g_v^{-1}$ works up to the standard
flag; put $\kappa_\infty=1$ (at the infinite places there is no condition: take $\kappa_\infty$ with
$g_\infty\kappa_\infty$ in the parabolic, by the real Iwasawa decomposition, the maximal compact acting
transitively on isotropic flags of the given type, and note $\kappa_\infty$ is integral vacuously).
Then $g\kappa$ maps $\langle u\rangle$ to itself and $\langle u\rangle\oplus\ell$ to itself.
(b) $g'$ stabilises $\langle u\rangle$, hence lies in $P_A=M_AN_A$ with $M_A$ the stabiliser of $u$ and
$u'$ lines and of $X'$. Let $m\in M_A$ act by $\alpha$ on $u$ and by the component $g'_{X'}$ (the map induced
on $\langle u\rangle^\perp/\langle u\rangle\cong X'$) on $X'$; then $r=g'm^{-1}$ acts trivially on
$\langle u\rangle$ and on $\langle u\rangle^\perp/\langle u\rangle$, so $r\in N_A$. Since $g'$ maps $e$ into
$e\gamma+u\beta'$, the induced map on $X'$ sends $e$ to $e\gamma$.
\end{proof}

\begin{lemma}[Straightening]\label{pcu:straighten}
Under the hypotheses of Lemma~\ref{pcu:unimod}, let $X_0$ fix the frame of $D$ off $\Sigma^c$, equal
$g$ on $\Sigma^c$, with $f^+(X_0)\ne0$, and let $g\kappa$ ($\kappa\in K_\Sigma$-integral) act on the
frame $(u,e)$ by $\binom{\alpha\ \beta'}{0\ \gamma}$ (such $\kappa$ exists by Lemma~\ref{pcu:iwasawa}). Then there is a straight point $P$ with $\gamma_P=\gamma$
and $\varphi^+_P\ne0$.
\end{lemma}
\begin{proof}
By Lemma~\ref{pcu:iwasawa}(b) write $g\kappa=r\,m$ with $r\in N_A$ and $m$ in the Levi of $A$ ($m$ acts by $\alpha$ on $u$ and
preserves $X'$, acting on $e$ by $\gamma$). Put $P_0=r^{-1}_{\Sigma^c}X_0\kappa_{\Sigma^c}$ (parts on
$\Sigma^c$): it equals $m$ on $\Sigma^c$, fixes the frame off $\Sigma^c$, and $f^+(P_0)=f^+(X_0\kappa_{\Sigma^c})
=f^+(X_0)\ne0$ by Lemma~\ref{pcu:sing} and right $K_\Sigma$-invariance. Then $P=(P_0,m,\alpha,
\bar\alpha^{-1},\gamma)$ is straight and $\varphi^+_P(0)=f^+(P_0)\ne0$.
\end{proof}

\begin{lemma}[Index decreases]\label{pcu:index}
Let $h$ act on $(u,e)$ by $\binom{t\ 0}{x\ \gamma}$, with $\gamma$ and $\beta[\gamma]$ invertible and
$\gamma^{-1}x$ not integral on $\Sigma^c$, and let $h\kappa$ ($\kappa$ integral) act by $\binom{\alpha\ \beta'}
{0\ \gamma'}$ with $\beta[\gamma']$ integral on $\Sigma^c$. Then $\mathrm{di}(\gamma')<\mathrm{di}(\gamma)$.
\end{lemma}
\begin{proof}
$\kappa^{-1}$ acts on the frame by an integral matrix $a$ (the frame is integral and unimodular), and
comparing, $x=\gamma'a_{21}$, $\gamma=\gamma'a_{22}$; hence $\gamma^{-1}x=a_{22}^{-1}a_{21}$ with $a_{21}$
integral, so $\det a_{22}$ is not a unit on $\Sigma^c$. Since $\beta[\gamma]=a_{22}^*\beta[\gamma']a_{22}$ and
$\beta[\gamma']$ is integral, $\det\beta[\gamma]^{-1}=\det\beta[\gamma']^{-1}\,\mathrm N(\det a_{22})^{-1}$ with
$\mathrm N(\det a_{22})=\det a_{22}\overline{\det a_{22}}$ integral and not a unit on $\Sigma^c$; as
$\det\beta[\gamma']$ is integral, the denominator of $\det\beta[\gamma]^{-1}$ is that of
$\det\beta[\gamma']^{-1}$ multiplied by the index of $\mathrm N(\det a_{22})\widehat\auO_{E,\Sigma^c}$, which is
$>1$.
\end{proof}

\begin{proof}[Proof of Lemma~\ref{pcu:unimod}]
(We prove it in the generality above, with $\Sigma=$ the places outside $S\cup\bar S$ used in
Proposition~\ref{pcu:selfdual}: there $f=\mathcal D(\Psi^{\mathrm{col}}_{z_0})$ is right invariant under
$K_\Sigma$ by (E3), and the shear invariance of $f^\pm$ at a non-empty set of places $\mathcal U$ (namely $\bar S$)
is the lemma proved after this one.)

Suppose $\pcuF_\beta f(g)\ne0$ for some $g$ fixing the frame off $\Sigma^c$. By Lemma~\ref{pcu:pm}(c),
$\varphi^+_g\ne0$, i.e.\ $f^+(n_+(a)g)\ne0$ for some $a$; by Lemma~\ref{pcu:straighten} there is a straight
point with $\varphi^+\ne0$. Choose one, $P$, with least $\mathrm{di}(P)$ (finite by
Lemma~\ref{pcu:bracket-int}); write $\gamma=\gamma_P$, $t,t'$ its scalars.

\emph{Supports.} If $\varphi^+_P(a)\ne0$ then $\gamma^{-1}a^*t$ is integral on $\Sigma^c$: otherwise the
point $n_+(a)g_0$ acts on $(u,e)$ by $\binom{t\ \ 0}{-a^*t\ \gamma}$, straightening $n_+(a)P$
(Lemma~\ref{pcu:straighten}, legitimate since $f^+(n_+(a)P)=\varphi^+_P(a)\ne0$) gives $Q$ with
$\varphi^+_Q\ne0$, and Lemma~\ref{pcu:index} with Lemma~\ref{pcu:bracket-int} gives
$\mathrm{di}(Q)<\mathrm{di}(P)$, a contradiction. The same holds for $\varphi^-$ with $t'$ (apply the
argument to the straight point with $u,u'$ exchanged, which has the same $\gamma$, and note that
$\varphi^+_{P}=0$ would force $\varphi^-_P=0$ by Lemma~\ref{pcu:pm}(c)).

\emph{Periods.} Let $x$ be a column with $\beta[\gamma]x$ integral on $\Sigma^c$ and
$a=\mathbf 1_{\Sigma^c}(\gamma x t'^*)^*$. For $b$ in the support of $\varphi^-_P$, $\chi_\beta(a,b)=1$:
indeed, with $c'=\gamma^{-1}b^*t'$ (integral on $\Sigma^c$ by Supports) and $t'^*t=1$, on $\Sigma^c$
$a\beta b^*=t'x^*\gamma^*\beta\gamma c't'^{-1}=\delta\,(\beta[\gamma]x)^*c'\in\delta\widehat\auO_E$, so
$\jmath a\beta b^*\in\jmath\alpha\bar\alpha\widehat\auO_E$, on which $\psi$ is trivial, and $a=0$ off $\Sigma^c$. By Lemma~\ref{pcu:pm}(b),
$\varphi^-_{n_+(a)P}=\varphi^-_P$, hence by Lemma~\ref{pcu:pm}(c) (with $u,u'$ exchanged)
$\varphi^+_{n_+(a)P}=\varphi^+_P$, i.e.\ $a$ is a period of $\varphi^+_P$. If $\varphi^+_P(y)\ne0$ then
$\varphi^+_P(y+a)\ne0$, and by Supports $\gamma^{-1}(y+a)^*t$ and $\gamma^{-1}y^*t$ are integral, so
$\gamma^{-1}a^*t=x$ (using $t'^*t=1$) is integral on $\Sigma^c$. Thus $\beta[\gamma]x$ integral implies $x$
integral: $\beta[\gamma]^{-1}$ is integral on $\Sigma^c$. With Lemma~\ref{pcu:bracket-int}, $\beta[\gamma]$ is
unimodular on $\Sigma^c$.

\emph{Contradiction.} By Lemma~\ref{pcu:pm}(c) and $\varphi^+_P\ne0$ there is a point $Y=n_-(b)n_+(a)P$,
equal on $\Sigma^c$ to a point acting on the frame of $D$ by $\gamma$ and fixing it elsewhere, with
$\pcuF_\beta f(Y)\ne0$. This contradicts the hypothesis of Lemma~\ref{pcu:unimod} (in the form: vanishing
at points acting by $\gamma$ on $\Sigma^c$, trivially elsewhere, with $\beta[\gamma]$ unimodular on $\Sigma^c$;
these are points $n(l,z)m(s_0q,g)$ of the statement: $Y$ stabilises $\ell$, so lies in $P_D(\auA)$, and its Levi
component on $\ell$ is $s_0q$ with $q=s_0^{-1}\gamma$ on $\Sigma^c$ and $q=s_0^{-1}$ at the finite places of $S\cup\bar S$;
then $q^*\beta q=\beta[\gamma]$ on $\Sigma^c$ and $q^*\beta q=s_0^{-1}\bar s_0^{-1}\beta$, the normalised index, at $S\cup\bar S$,
both unimodular).
\end{proof}

\begin{proof}[Shear invariance of $f^\pm$ for $f=\mathcal D(\Psi^{\mathrm{col}}_{z_0})$]
Let $\varsigma$ be a shear at $\bar S$ ($1+uc$ there, $cu=0$, trivial off $S\cup\bar S$) and $g$ fixing the frame
of $D$ off $\Sigma^c$. By (E2), $\mathcal D(\Psi')(xg)$ with $\Psi'=\Psi_0\circ(1+k\,c'\,\cdot)$
($k$, $c'$ the $S$-components of $(\varsigma g)^{-1}u$ and $c\,\varsigma g$) equals $\mathcal D(\Psi_0)$ at $\varsigma^{-1}$-translated
points, so $f^+(\varsigma g)$ is the value at $\Psi'$ of the linear form $\Lambda=$ ``$f^+$ at $\varsigma g$'' (linear on
$\mathcal M$ by Lemma~\ref{pcu:pm}(b) and (E1)), while $f^+(g)=\Lambda(\Psi_0)$ after moving $\varsigma$. The form
$\Lambda$ is invariant under the shears $y\mapsto y+k\,w\,z_0y$ ($w$ any row), because these come from the
shifts $n_-(\cdot)$ at $S$, which fix $f^+$ (Lemma~\ref{pcu:pm}(b)); here $z_0 k=0$ and $c'k=0$. It remains:
\emph{if a linear form $\Lambda$ on a space of Schwartz--Bruhat functions is invariant under all
$y\mapsto y+k\,w\,z_0y$, then $\Lambda(\Psi_0\circ(1+kr\,\cdot))=\Lambda(\Psi_0)$ for every row $r$ with
$rk=0$.} Average over integral $t$: $\Lambda$ is unchanged by replacing $\Phi$ with
$\bar\Phi(y)=\int\Phi(y+k\,r\,t\,z_0y)dt$ (a finite combination, by local constancy), and
$\bar{\Psi_0\circ(1+kr\cdot)}=\bar\Psi_0$: if $\Psi_0(y+k\rho)\ne0$ for some $\rho$, then $z_0y$ has an
integral left inverse $g'$ and $(1+kr)(y+krtz_0y)=y+kr(t+t_0)z_0y$ with $t_0=yg'$ integral, so the two
averages agree by translation invariance of $dt$; otherwise both vanish.
\end{proof}

\subsection{Prescribed slices: proof of (M1), and the decay of slices}\label{pcu:ss-M1}

Write the real Fock spaces as $\mathcal F'$ (of $\auHom(X,T)_\R\cong\C^{\iota'}$) and $\mathcal F$ (of
$\auHom(X_1,T)_\R\cong\C^{\iota}$), with moments $\|v\|_N=\int(1+|y|)^N|v(y)|dy$, Heisenberg operators
$\psrho$ and vacua $e$. The mixed model of $D$ is characterised by: there are linear maps $L,K\colon
\C^p\to\C^{\iota'}$, $I\colon\C^\iota\to\C^{\iota'}$ (here $p$ is the complex dimension of
$\auHom(\ell,T)_\R$) with $\C^{\iota'}=L\C^p\oplus I\C^\iota\oplus K\C^p$, and bounded maps
$\mathrm{sl}(x)\colon\mathcal F'\to\mathcal F$, such that
\begin{equation}\label{pcu:eq-mixed}
\mathrm{sl}(x)\,\psrho(La'+Iw+Kb)v=\beta(2x+a',b)\,\psrho(w)\,\mathrm{sl}(x+a')v,\qquad
\mathrm{sl}(x)e=e^{-|x|^2}e,
\end{equation}
$\beta(x,b)=\exp(i\,\Im\,b^*x)$ (the Schr\"odinger variable is translated by $L$, multiplied by a
character by $K$, and the rest acts on the values). Let $y\mapsto(A y,Wy,By)$ be the inverse of
$(a,w,b)\mapsto La+Iw+Kb$.

\begin{lemma}[Slices decay]\label{pcu:slicedecay}
For all $N,M$ there is $C$ with $(1+|x|)^M\|\mathrm{sl}(x)v\|_N\le C\|v\|_{N+M}$ for all $x$, $v$.
In particular $\int\|\mathrm{sl}(x)v\|_0\,dx\le C'\|v\|_{2p+1}$.
\end{lemma}
\begin{proof}
For $v=\psrho(y)e$, \eqref{pcu:eq-mixed} gives $\mathrm{sl}(x)v=(\text{phase})\,e^{-|x+Ay|^2}\psrho(Wy)e$.
From $1+|x|\le(1+|x+Ay|)(1+|Ay|)$ and $(1+t)^Me^{-t^2}\le e^{M^2}$,
$(1+|x|)^Me^{-|x+Ay|^2}\le e^{M^2}(1+\|A\||y|)^M$; and $\|\psrho(w)e\|_N\le(1+|w|)^N\|e\|_N$ (translate
the weight). Hence $(1+|x|)^M\|\mathrm{sl}(x)\psrho(y)e\|_N\le C(1+|y|)^{N+M}$, and
Lemma~\ref{pcx:superpos} ($v=\int v(y)\psrho(y)e\,dy$, $\mathrm{sl}(x)$ bounded) gives the bound. The
second statement follows with $M=2p+1$, since $\int(1+|x|)^{-2p-1}dx<\infty$ on $\C^p=\R^{2p}$.
\end{proof}
This is the estimate used in (M2): the integral $\int\mathrm{sl}_D(x)\Lambda\,dx$ converges
absolutely in $\mathcal F$ as soon as $\|\Lambda\|_{2p+1}<\infty$.

\begin{proof}[Proof of (M1)]
Let $\tau$, $u$, $v$ be as in (M1); write $\xi=v$, and for $g$ in the real unitary group of (M1) let $S_g$ be the real
symplectic map of $\auHom(X_1,T)_\R=\C^\iota$ by which $g$ acts, so that $\omega(g)\psrho(w)\omega(g)^{-1}=\psrho(S_gw)$
((W1)). Write $S_x=S_{u(x)}$; on the compact set $\mathrm{supp}\,\tau$, $x\mapsto S_x$ is smooth and
$B=\sup\|S_x\|$, $B'=\sup\|S_x^{-1}\|$ are finite. No compactness of the unitary groups is used. Functions of vectors
are taken in the realisation of Lemma~\ref{pcx:superpos}: $F(w)=\langle\psrho(w)e,F\rangle$.

\emph{Step 1: Gaussian vectors.} Put $E_g=\omega(g)e$. (i) $|E_g(w)|\le\exp(-|w|^2/(1+\|S_g^{-1}\|^2))$: this is the
bound proved in the first part (``Reduction'') of the proof of Lemma~\ref{pcu:moments}, applied to the unitary operator
$k=\omega(g)^{-1}$, for which $k\psrho(w)k^{-1}=\psrho(S_g^{-1}w)$ and $k^{-1}e=E_g$. (ii) In a Schr\"odinger model of
$\C^\iota$ in which $e$ is the Gaussian $\gamma_{Z_0}(t)=\exp(\boldsymbol\pi i\,t^{\mathsf t}Z_0t)$, $Z_0=i1$, the operator
$\omega(g)$ is, up to a scalar, the metaplectic operator of $S_g$, and the metaplectic operators carry $\gamma_{Z_0}$ to a
multiple of $\gamma_{Z}$ with $Z$ the image of $Z_0$ under the linear fractional action on the Siegel upper half space of
$S_g$ or of $(S_g^{\mathsf t})^{-1}$, according to the normalisation (Folland~\cite[Chapter~4, Thm.~4.65]{FollandPS}: the metaplectic operator of $A$ carries $\gamma_Z$ to a multiple of $\gamma_{\alpha((A^{\mathsf t})^{-1})Z}$). So $E_g=c(g)\gamma_{Z(g)}$ with $Z(g)$ smooth in $g$; and $c(g)=\langle
E_g,e\rangle/\langle\gamma_{Z(g)},\gamma_{Z_0}\rangle$ is smooth, because the denominator is an explicit Gaussian integral,
smooth and nowhere zero in $Z$, and $g\mapsto\langle\omega(g)e,e\rangle$ is smooth, $e$ being a Schwartz function and hence a
smooth vector (Folland~\cite[Chapter~4, Thm.~4.45]{FollandPS}). Computing the Gaussian integral
$\langle\psrho(w)\gamma_{Z_0},\gamma_Z\rangle$ gives $E_g(w)=c'(g)\exp Q_g(w)$ with $Q_g$ a quadratic form in the real
coordinates of $w$, with complex coefficients, and $c'$, $Q_g$ smooth in $g$. Hence every derivative of order $|\alpha|$ in
$g$ of $E_g(w)$, and in $w$ of order $|\alpha'|$, is $E_g(w)$ times a polynomial of degree $\le2|\alpha|+|\alpha'|$ in $w$
whose coefficients are smooth in $g$.

\emph{Step 2: the functions of the prescribed slices.} By Lemma~\ref{pcx:superpos} and boundedness of $\omega(g)$,
$\omega(g)\xi=\int\xi(y)\psrho(S_gy)E_g\,dy$. The Heisenberg relations give
$\langle\psrho(w)e,\psrho(a)E\rangle=e^{i\vartheta(a,w)}E(w-a)$ with $\vartheta$ a real quadratic polynomial in $(a,w)$
(the cocycle of the Heisenberg group). So the function of $G(x)=\tau(x)\omega(u(x))\xi$ is
\[
g_x(w)=\tau(x)\int\xi(y)\,e^{i\vartheta(S_xy,w)}\,E_{u(x)}(w-S_xy)\,dy .
\]
Put $\varrho=1+|y|+|w|$. By Step 1 and the chain rule, each $x$-derivative of the integrand brings a factor $O(\varrho^2)$,
uniformly on $\mathrm{supp}\,\tau$, so with $c=(1+B'^2)^{-1}$
\[
\bigl|\partial_x^\alpha\bigl[\tau(x)e^{i\vartheta(S_xy,w)}E_{u(x)}(w-S_xy)\bigr]\bigr|\le C_\alpha\,\varrho^{2|\alpha|}e^{-c|w-S_xy|^2}.
\]
Since $\varrho\le(1+|w-S_xy|)(1+B)(1+|y|)$, this is $\le C'_\alpha(1+|y|)^{2|\alpha|}(1+|w-S_xy|)^{2|\alpha|}e^{-c|w-S_xy|^2}$,
which is integrable against $|\xi(y)|$; so differentiation under the integral is legitimate, and, using
$1+|w|\le(1+|w-S_xy|)(1+B|y|)$ and integrating first in $w$,
\begin{equation}\label{pcu:eq-M1bound}
\int\!\!\int(1+|w|)^N\,|\partial_x^\alpha g_x(w)|\,dw\,dx\le C_{N,\alpha}\,\|\xi\|_{N+2|\alpha|}<\infty
\end{equation}
for all $N$, $\alpha$ (the $x$-integral runs over the compact set $\mathrm{supp}\,\tau$).

\emph{Step 3: construction.} Put $h_w(x)=e^{|x|^2}g_x(w)$, smooth with compact support in $x$;
$\hat h_w(b)=\int h_w(x)e^{-2\pi i\langle x,b\rangle}dx$ (Fourier transform in $x\in\C^p=\R^{2p}$); and
$k(b)=K(-\pi i\,b)$, so that by \eqref{pcu:eq-mixed} $\mathrm{sl}(x)\psrho(k(b))=e^{2\pi i\langle x,b\rangle}\mathrm{sl}(x)$.
Define
\[
\Lambda=\int\!\!\int\hat h_w(b)\,\psrho(k(b))\,\psrho(Iw)e\,db\,dw .
\]
\emph{Moments.} $\|\psrho(k(b))\psrho(Iw)e\|_N\le(1+\|k\||b|)^N(1+\|I\||w|)^N\|e\|_N$ (translation of the weight).
Integrating by parts in $x$, $(1+|b|)^M|\hat h_w(b)|\le C_M\sum_{|\alpha|\le M}\int|\partial_x^\alpha h_w(x)|dx\le
C'_M\sum_{|\alpha|\le M}\int|\partial_x^\alpha g_x(w)|dx$ ($e^{|x|^2}$ and its derivatives are bounded on
$\mathrm{supp}\,\tau$). Taking $M=N+2p+1$, $\int(1+\|k\||b|)^N(1+|b|)^{-M}db<\infty$, and \eqref{pcu:eq-M1bound} gives
$\iint|\hat h_w(b)|\,\|\psrho(k(b))\psrho(Iw)e\|_N\,db\,dw<\infty$. So the integral converges absolutely for every
$\|\cdot\|_N$, and $\|\Lambda\|_N<\infty$ for all $N$.

\emph{Slices.} The Hilbert norm is $\le\|\cdot\|_0$ (Lemma~\ref{pcx:superpos}), so the integral converges absolutely in
$\mathcal F'$ and the bounded operator $\mathrm{sl}(x)$ passes inside. By \eqref{pcu:eq-mixed} with $a'=0$, $b=0$,
$\mathrm{sl}(x)\psrho(Iw)e=\psrho(w)\mathrm{sl}(x)e=e^{-|x|^2}\psrho(w)e$; and by Fourier inversion ($h_w$ smooth with
compact support), $\int\hat h_w(b)e^{2\pi i\langle x,b\rangle}db=h_w(x)$. Hence
\[
\mathrm{sl}(x)\Lambda=\int h_w(x)e^{-|x|^2}\psrho(w)e\,dw=\int g_x(w)\psrho(w)e\,dw=G(x)=\tau(x)\,\omega(u(x))\xi,
\]
the last but one equality by Lemma~\ref{pcx:superpos}, as $\|G(x)\|_0<\infty$ by \eqref{pcu:eq-M1bound} with $\alpha=0$
(pointwise in $x$ by the same bound before integrating in $x$).
\end{proof}

\subsection{Exchange on standard data: proof of Proposition~\ref{pcx:P-std}}\label{pcu:ss-std}

\begin{proof}[Proof of Proposition~\ref{pcx:P-std} from Propositions~\ref{pcx:P-slice},
\ref{pcx:P-top} and~\ref{pcx:P-descent}]
Let $W$ (rank $N=b+r$, even, $q$ negative signs), $U$ (rank $b$), $V$ (rank $r$) with
$q_U+q=b+q_V$, $W_0=\pcxH\oplus W$ (adapted, Lemma~\ref{pcx:adapted}), $z$ for $(W_0,W_0)$ and $y$ for
$(W,U)$ be given.

\emph{Step 1.} Proposition~\ref{pcx:P-slice} for the adapted splitting $W_0=W\oplus\pcxH$ gives $w$ for
$(W_0,U)$ with finite moments and zero term $y$.

\emph{Step 2.} Proposition~\ref{pcx:P-top} for $W_0$, $U$, $z$, $w$ gives $X=\pcxH^b\oplus W_0=
\pcxH^{b+1}\oplus W$, $T=U\oplus W_0$, and $x$, $x'$ with zero term $w$ along the $b$ planes and
$\mathrm{RC}(X,T,U;x,x';\nu\mapsto\pcxstr_{W_0}(\nu;z))$. Here $\operatorname{rk}T=b+N+2$ and $T$ has
$q_U+q+1$ negative signs.

\emph{Step 3: the tower.} For $i=0,\dots,b+1$ let $X_i=\pcxH^{i}\oplus W$ (so $X_{b+1}=X$, $X_0=W$,
each $X_{i+1}=X_i\oplus\pcxH$ adapted and of even rank). Let $T_{b+1}=T$ and, for $0\le i\le b$,
$T_i=V\oplus\pcxH^i$ (Lemma~\ref{pcx:adapted}), of rank $r+2i$ with $q_V+i$ negative signs. Then
$\operatorname{rk}T_{i+1}=\operatorname{rk}T_i+2$ and $q_{T_{i+1}}=q_{T_i}+1$ for all $i$: for $i<b$
this is clear, and for $i=b$, $\operatorname{rk}T=b+N+2=r+2b+2$ and $q_T=q_U+q+1=q_V+b+1$ by the
hypothesis $q_U+q=b+q_V$.

\emph{Step 4: descent.} Apply Proposition~\ref{pcx:P-descent} $b+1$ times, to $W'=X_{i+1}=X_i\oplus
\pcxH$, $T=T_{i+1}$, $T_0=T_i$, $T'=U$, $i=b,b-1,\dots,0$: each application keeps $\ell$ and replaces
the second vector by its zero term along the plane removed. After $b+1$ steps one gets $v$ for
$(W,V)$ with $\mathrm{RC}(W,V,U;v,y';\ell)$, where $y'$ is the iterated zero term of $x'$ along the
$b+1$ planes. By (M4) this is the zero term of $x'$ along $X=\pcxH^b\oplus W_0$ followed by the zero
term along $W_0=\pcxH\oplus W$, i.e.\ the zero term of $w$, which is $y$ (Steps 1--2). So
$\mathrm{RC}(W,V,U;v,y;\nu\mapsto\pcxstr_{W_0}(\nu;z))$, which is Proposition~\ref{pcx:P-std}.
\end{proof}

\begin{remark}[Consistency conditions on Proposition~\ref{pcx:P-descent}]\label{pcu:rem-descent}
The proof needs from Proposition~\ref{pcx:P-descent}, as stated in \S\ref{pcx:ss-arch}: (i) it is
applied with the zero term taken along the splitting with the same orientation as the one in which
Proposition~\ref{pcx:P-top} and Proposition~\ref{pcx:P-slice} give zero terms (along $D^{\mathrm s}$
for the $b$ planes); (ii) the sign condition ``with the signs
required'' means $q_{T_0}=q_T-1$; (iii) the exponent in $\mathrm{RC}$(a) changes from
$\operatorname{rk}T-\operatorname{rk}W'/2=r-N/2+i+1$ to $r-N/2+i$, and the constant $C$ is allowed
to change while $\ell$ is kept; (iv) $T_0$ is an arbitrary space of rank $\operatorname{rk}T-2$ with
$q_T-1$ negative signs (not necessarily $T$ minus a plane); removal of $p\ge1$ planes follows by
induction from one plane, as in Step 4. Proposition~\ref{pcx:P-descent} is proved for $T=T_0\oplus
\pcxH$ (adapted), as in the set-up \S\ref{pct:ss-setup}, and Step~3 works with this form: take $T_b=U\oplus W$ (rank
$b+N=r+2b$, $q_U+q=q_V+b$ negative signs), so $T=U\oplus W_0=T_b\oplus\pcxH$; descend once to
$(X_b,T_b)$; then $T_b$ and $V\oplus\pcxH^b$ have the same rank and signs, and Lemma~\ref{pcx:move}
(with $W_1=W=X_b$, $k=1$) transfers $\mathrm{RC}(X_b,T_b,U;\cdot)$ to
$\mathrm{RC}(X_b,V\oplus\pcxH^b,U;\cdot)$, multiplying $c$ by a number of modulus $1$ independent of
the data and keeping $C$; the remaining $b$ descents are of the form $T_{i+1}=T_i\oplus\pcxH$.
\end{remark}

\subsection{Dependencies of Theorem~\ref{pcx:X}}\label{pcu:ss-final}

In the deductions of \S\ref{pcx:ss-assembly} and \S\ref{pcx:ss-equal} the statements of \S\ref{pcx:ss-arch} are applied
as follows: (1) in the proof of Proposition~\ref{pcx:P-matches},
Proposition~\ref{pcx:P-std} is applied to $(W,U^-,V)$, whose sign condition $q_{U^-}+q=b+q_V$ is
equivalent to $q=q_U+q_V$; the moments of $y^{\mathrm{ad}}$ are those of $y$; (2) in the proof of
Theorem~\ref{pcx:X}, Proposition~\ref{pcx:P-matches} is applied to $(W,V',V)$ with $y=v$ a vector of
$(V',W)$, $N=a+b=2b+2e$ is even, $q=q_V+q_{V'}$, and $a-N/2=e$; (3) Lemma~\ref{pcx:twice} applies
Proposition~\ref{pcx:P-matches} to $(J,0,J)$ and $(J,J,0)$, which satisfy its hypotheses (for $b=0$
Proposition~\ref{pcx:P-top} is the identity on $X=T=W_0$ and the descent is from $(W_0,W_0)$ to
$(J,J)$). By Corollary~\ref{pcx:conv}(iii) both integrals of \eqref{pcx:eq-X} converge absolutely. The
statements on which Theorem~\ref{pcx:X} depends:

\medskip
{\small\noindent\begin{tabular}{@{}p{0.36\textwidth}p{0.6\textwidth}@{}}
\toprule
Statement & Where proved\\ \midrule
Theorem~\ref{pcx:X} & \S\ref{pcx:ss-assembly}, from \ref{pcx:P-matches}, \ref{pcx:P-equal},
\ref{pcx:Lpos}, \ref{pcx:mod}, \eqref{pcx:eq-twist}\\
Convergence (Cor.~\ref{pcx:conv}) & \S\ref{pcx:ss-conv}, from Prop.~\ref{pcx:decay}, (R1), (R2), (R3)\\
Lemma~\ref{pcx:Lpos}, Cor.~\ref{pcx:cne0} & \S\ref{pcx:ss-statement}; Hecke's non-vanishing on
$\Rel s=1$\\
Lemmas~\ref{pcx:std-matched}, \ref{pcx:mod} & \S\ref{pcx:ss-elem}\\
Prop.~\ref{pcx:P-matches} & \S\ref{pcx:ss-assembly}, from \ref{pcx:P-std}, \ref{pcx:P-match},
\ref{pcx:P-flip}\\
Prop.~\ref{pcx:P-match} & \S\ref{pcx:ss-match} (with Cor.~\ref{pcx:conv})\\
Prop.~\ref{pcx:P-flip} & \S\ref{pcx:ss-flip}, from Prop.~\ref{pf:flip}\\
Prop.~\ref{pcx:P-std} & \S\ref{pcu:ss-std}, from \ref{pcx:P-slice}, \ref{pcx:P-top},
\ref{pcx:P-descent}, \ref{pcx:adapted}, (M4)\\
Prop.~\ref{pcx:P-slice} & \S\ref{pcu:ss-slice}, from (M1), (M2)\\
Prop.~\ref{pcx:P-top} & \S\S\ref{pcu:ss-witt}--\ref{pcu:ss-std-data}, from \ref{pcx:const},
\ref{pcx:landherr}, \ref{pcx:generate}, \ref{pcx:linfactor}, \ref{pcx:section-laws}, \ref{pcx:Sact},
(M1)--(M3), (U1)--(U3), Lemmas~\ref{pcu:unimod} (\S\ref{pcu:ss-unimod}), \ref{pcu:integrable}\\
Prop.~\ref{pcx:P-descent} & \S\ref{pct:ss-descent} (Proposition~\ref{pct:descent})\\
Prop.~\ref{pcx:P-equal} & \S\ref{pcx:ss-equal}, from \ref{pcx:P-matches}, \ref{pcx:P-nonvan},
\ref{pcx:P-plane}, \ref{pcx:P-signs}, \ref{pcx:adapted}\\
Prop.~\ref{pcx:P-nonvan} & \S\ref{pct:ss-nonvan}, with Lemma~\ref{pct:negcoeff} $=$ Theorem~\ref{pcn:negcoeff} (local vanishing, Lemma~\ref{pcn:local}; no unfolding)\\
Prop.~\ref{pcx:P-plane}, Lemma~\ref{pcx:const} & \S\S\ref{pct:ss-plane}, \ref{pct:ss-const}\\
Prop.~\ref{pcx:P-signs} & \S\ref{pcx:ss-signs}, from \ref{pcx:move}, \ref{pcx:const}\\
Lemmas~\ref{pcx:adapted}, \ref{pcx:landherr}, \ref{pcx:move} & \S\S\ref{pcx:ss-arch},
\ref{pcx:ss-signs}; Landherr, Hasse norm theorem, Jacobowitz\\
Prop.~\ref{pcx:decay} & \S\ref{pcx:ss-conv}; (R1) reduction theory, (R2) moderate growth\\
(M1)--(M4), Lemma~\ref{pcu:mixedmodel} & \S\ref{pcu:ss-mixed}, \S\ref{pcu:ss-M1}; Bargmann transform (Folland)\\
\bottomrule
\end{tabular}}
\medskip

Thus Theorem~\ref{pcx:X}, with absolutely convergent integrals, rests on the proofs listed and on
the published results listed in \S\S\ref{pct:ss-status}, \ref{pcn:ss-dep}.


%% ===== proofs/C-exchange-tower.tex =====
% proofs/C-exchange-tower.tex
% The plane constant term of theta integrals on standard data and what rests on it:
% Lemma pcx:const, Propositions pcx:P-plane, pcx:P-descent, pcx:P-nonvan of C-exchange.tex.
% Labels of this file start with pct:. Macros: those of C-exchange.tex and C-determinant.tex.
\providecommand{\pctN}{\mathsf N}
\providecommand{\pctP}{\mathsf P}
\providecommand{\pctL}{\mathcal L}
\providecommand{\pctZ}{\mathcal Z}
\providecommand{\pctK}{\mathsf K}
\providecommand{\pctsl}{\operatorname{sl}}
\providecommand{\pctcone}{\operatorname{cone}}

\section{The plane constant term}\label{pct:sec}

This section proves Lemma~\ref{pcx:const} and Propositions~\ref{pcx:P-plane},
\ref{pcx:P-descent} and (modulo one input from the large triple, see \S\ref{pct:ss-nonvan})
\ref{pcx:P-nonvan}. All four rest on one computation, the constant term of a theta integral
$\Theta^{\mathrm L}$ along the isotropic line of an added hyperbolic plane
(Theorem~\ref{pct:plane}). \emph{Convention:} in this section a Fock vector $v$ \emph{has finite
moments} if $\|v\|_N=\int(1+|y|)^N|v(y)|\,dy<\infty$ for every $N\ge0$ (moments of all orders; the norms
of Lemma~\ref{pcx:superpos}); by Lemma~\ref{pcx:superpos}(b) such a $v$ lies in $\pzF_{s'}$ for every
integer $s'\ge0$. Notation is that of \S\ref{pcx:ss-arch}: $\nu$ is a tame character of
weight $\mu$ (all $\mu_\tau\ge1$) for a modulus $\pcm\ne(1)$, $S\ne\emptyset$ is the set of primes
of $\pcm$ (so $S\cap\bar S=\emptyset$), $F=F_S$, $O=\auO_S$, $\chi=\chi_S$, $Q=\prod_{w\in S}q_w$,
and for a space $W$ of even rank $N$
\[
\gamma_W=\nu^{-1}\upsilon^{N/2},\qquad
\Theta^{\mathrm L}_{W,T}(v,\Phi)(g)=\int_{[\auU(T)]}\theta^{W,T}_{\gamma_W,\gamma_W}(g,h;[v,\Phi]_S)\,dh .
\]
We write $\tilde\nu(x)=\nu(x/\bar x)$ for the base change of $\nu$, a unitary Hecke character of
$\auA_E^\times/E^\times$, $|\cdot|$ for the idele norm of $E$, and
\[
L^\sharp(x)=|L^S(x,\nu)|\,Q^{x-1/2}\qquad(x\ge1).
\]
Then $L_*(\pcm,\nu,e+1)=L^\sharp(x_e)^2L_*(\pcm,\nu,e)$ for every $e\in\Z$, where $x_e=e+1$ if
$e\ge0$ and $x_e=-e$ if $e<0$: for $e\ge0$ this is $Q^{(e+1)^2-e^2}|L^S(e+1,\nu)|^2$, and for
$e=-(m+1)$ with $m\ge0$, inverting the case $e=m$ through $L_*(-e')=L_*(e')^{-1}$ gives
$L_*(-m)=L^\sharp(m+1)^2L_*(-m-1)$, which is the claim since $x_e=m+1$. By Lemma~\ref{pcx:Lpos}, $L^\sharp(x)>0$ for $x\ge1$. Facts
(W1)--(W5) are those of Lemma~\ref{ps:weil}.

The base change $\tilde\nu$ has archimedean component $\tilde\nu_\infty(a)=\prod_\tau(a_\tau/|a_\tau|)^{2\mu_\tau}$
on $(E\otimes\R)^\times=\prod_\tau\C^\times$; at a place $v=w\bar w$ with $w\in S$, $\tilde\nu_w$ and
$\tilde\nu_{\bar w}$ are ramified (their restrictions to units are $\chi_w$ and $\chi_w^{-1}$ composed
with conjugation), and $\tilde\nu$ is unramified outside $S\cup\bar S\cup\infty$. Hence the Euler
factors of $L(s,\tilde\nu)$ at $S\cup\bar S$ are $1$ and $L^S(s,\nu)=L(s,\tilde\nu)$.

\subsection{Set-up: the plane, its parabolic and the line terms}\label{pct:ss-setup}

Let $W'=W\oplus\pcxH_A$ and $V'=V\oplus\pcxH_X$ be adapted splittings (Lemma~\ref{pcx:adapted}):
$\pcxH_A$ has the isotropic basis $e_A,f_A$ with $\langle e_A,f_A\rangle=1$, the lattice of $W'$ is
$L_W\oplus\widehat\auO_Ee_A\oplus\widehat\auO_Ef_A$, and likewise $e_X,f_X$ for $V'$. Put
$N=\operatorname{rk}W$ (always even below), $r=\operatorname{rk}V$, $b=N-r$ and
\[
d=r-N/2\in\Z .
\]
Then $\operatorname{rk}W'=N+2$, $\operatorname{rk}V'=r+2$, $(N+2)-(r+2)=b$, and
$\gamma_{W'}=\gamma_W\upsilon$.

\emph{Parabolics.} $\pctP_A\subset\auU(W')$ is the stabiliser of the line $Ee_A$, $\pctN_A$ its
unipotent radical, and for $u\in\auA_E^\times$, $g\in\auU(W)(\auA)$ let $m_A(u,g)$ be the Levi
element acting by $g$ on $W$, by $e_A\mapsto ue_A$ and $f_A\mapsto\bar u^{-1}f_A$. For $V'$ we use the
parametrisation $\ell_X(a,h)$: $h$ on $V$, $e_X\mapsto a^{-1}e_X$, $f_X\mapsto\bar af_X$. The group
$\pctN_A(E)\backslash\pctN_A(\auA)$ is compact; $dn$ is its probability measure. For a continuous
function $F$ on $\auU(W')(\auA)$, left invariant under $\pctN_A(E)$, its \emph{constant term} is
$F_A(x)=\int_{\pctN_A(E)\backslash\pctN_A(\auA)}F(nx)\,dn$. The group $\pctN_X(\auA)$ has
coordinates $(c,t)$, $c\in V(\auA)$, $t$ in the $\auA_{E^+}$-line $\{t\in\auA_E:t+\bar t=0\}$ (a
central extension of $V(\auA)$ by this line), and $\ell_X(a,h)$ acts by conjugation through
$c\mapsto hc\,\bar a^{-1}$, $t\mapsto(a\bar a)^{-1}t$; as $|\det h|=1$, this gives
$\int f(\ell_X(a,h)n\ell_X(a,h)^{-1})\,dn=|a|^{r+1}\int f(n)\,dn$ for a Haar measure $dn$ of
$\pctN_X(\auA)$.

\emph{The plane datum.} Let $\mathrm{pr}\in\pcMat_{N\times(N+2)}(O)$ be the matrix, in the frames, of
the orthogonal projection $W'\to W$; it is integral because the splitting is adapted. For
$z\in\pcMat_{b\times N}(O)$ put
\[
\Phi'_z=\Phi^{\mathrm{row}}_{z\,\mathrm{pr}}\qquad\text{on }\pcMat_{(r+2)\times(N+2)}(F)
\]
(Definition~\ref{pcx:def-std} for $(W',V')$; $z\,\mathrm{pr}\in\pcMat_{b\times(N+2)}(O)$). For $b=0$,
$\Phi'_\emptyset=\Phi^{\mathrm{row}}_\emptyset=\chi^{-1}\circ\det$ on $\auGL_{N+2}(O)$, which is
$\Psi^{\mathrm{col}}_1$ of Definition~\ref{pcx:def-ratio}.

\emph{Coordinates on $\auHom(W',V')$.} Write a map $\phi\colon W'\to V'$ as a $3\times3$ block matrix
with rows indexed by $(V,e_X,f_X)$ and columns by $(W,e_A,f_A)$. In the form
$\langle\phi,\psi\rangle=\operatorname{tr}(\psi^*\phi)$ the block $(y,a)$ pairs non-trivially only
with the block $(y^\vee,a^\vee)$, where $e^\vee=f$, $f^\vee=e$, $V^\vee=V$, $W^\vee=W$. Put
\[
\lambda_{\mathrm t}=E_{e_X,e_A},\qquad\lambda_{\mathrm f}=E_{f_X,f_A},\qquad
I_0=(V,f_A)\oplus(e_X,W)\oplus(e_X,f_A),
\]
where $E_{y,a}$ is the map $a\mapsto y$ (zero on the other basis vectors and on $W$), and $(y,a)$
denotes the corresponding block. Then $I_0$ is isotropic, $\Pi=E\lambda_{\mathrm t}\oplus
E\lambda_{\mathrm f}$ is a hyperbolic plane orthogonal to $I_0$, and
$I_0^\perp/I_0=\Pi\oplus\auHom(W,V)$. Also $I_A=(V',f_A)=(V,f_A)\oplus(e_X,f_A)\oplus(f_X,f_A)$
(the maps vanishing on $W\oplus Ee_A$) is isotropic and $I_A^\perp/I_A=\auHom(W,V')$, and
$I_A\cong V'$ by $x\mapsto(w'\mapsto\langle w',e_A\rangle x)$.

\begin{definition}[Slices and line terms]\label{pct:def-line}
For $x\in V'(\auA)$ let $x\otimes f_A^\vee\in\auHom(W',V')(\auA)$ be $w'\mapsto\langle w',f_A\rangle x$
(this is $x$ placed in the column $e_A$), and define the \emph{slice}
$\pctsl_A(x)=\psZ_{I_A}\circ\psrho(x\otimes f_A^\vee)\colon\psHS(\auHom(W',V'))\to
\psHS(\auHom(W,V'))$ (zero terms as in (W2)). For $\alpha\in\auA_E$ the two \emph{line terms}
$\pctL_{\mathrm t}(\alpha),\pctL_{\mathrm f}(\alpha)\colon\psHS(\auHom(W',V'))\to\psHS(\auHom(W,V))$
are
\[
\pctL_{\mathrm t}(\alpha)=\psZ_{I_0+E\lambda_{\mathrm f}}\circ\psrho(\alpha\lambda_{\mathrm t}),\qquad
\pctL_{\mathrm f}(\alpha)=\psZ_{I_0+E\lambda_{\mathrm t}}\circ\psrho(\alpha\lambda_{\mathrm f}).
\]
By the transitivity of zero terms in (W2), $\pctL_{\mathrm t}(\alpha)=\psZ_{(e_X,W)}\circ
\pctsl_A(\alpha e_X)$ (first the zero term along $I_A$, then along $\auHom(W,Ee_X)\subset
\auHom(W,V')$), and $\pctL_{\mathrm f}(\alpha)$ is the zero term along $\auHom(W',Ee_X)$ followed by
the zero term along $(V,f_A)$, after $\psrho(\alpha\lambda_{\mathrm f})$. All these operators are
tensor products of local operators; the archimedean factors, on Fock spaces, are written
$\pctsl_{A,\infty}(x)$ and $\pctL^\infty_\beta(a)$ ($\beta\in\{\mathrm t,\mathrm f\}$,
$a\in(E\otimes\R)$).
\end{definition}

\begin{lemma}[Poisson's formula along the line]\label{pct:poisson}
For every datum $f$ of $\auHom(W',V')$,
$\sum_{\xi\in E}\theta_{W,V}(\pctL_{\mathrm t}(\xi)f)=\sum_{\xi\in E}\theta_{W,V}(\pctL_{\mathrm f}(\xi)f)$,
both series converging absolutely.
\end{lemma}
\begin{proof}
Put $F_0=\psZ_{I_0}f\in\psHS(\Pi\oplus\auHom(W,V))$. Since $\lambda_{\mathrm t},\lambda_{\mathrm f}\in
I_0^\perp$, (W2) gives $\pctL_{\mathrm t}(\xi)f=\psZ_{E\lambda_{\mathrm f}}\psrho(\xi\lambda_{\mathrm t})F_0$
and $\pctL_{\mathrm f}(\xi)f=\psZ_{E\lambda_{\mathrm t}}\psrho(\xi\lambda_{\mathrm f})F_0$, the zero
terms being taken in $\Pi\oplus\auHom(W,V)$. Apply (W3) in this space, once with the isotropic line
$I=E\lambda_{\mathrm f}$ (then $W/I^\perp\cong E\lambda_{\mathrm t}$) and once with $I=E\lambda_{\mathrm t}$:
both sums equal $\theta_{\Pi\oplus\auHom(W,V)}(F_0)$, with absolute convergence.
\end{proof}

\begin{lemma}[Levi elements]\label{pct:levi}
Let $u,c\in\auA_E^\times$, $g\in\auU(W)(\auA)$, $h\in\auU(V)(\auA)$, and assume $u=1$ or $r=N$.
Put $x=cu$, $P_{\mathrm t}(x)=x$, $P_{\mathrm f}(x)=\bar x^{-1}$. Then for every datum $f$ and
$\alpha\in\auA_E$,
\[
\pctL_\beta(\alpha)\,\omega_{\gamma_{W'},\gamma_{W'}}\bigl(m_A(u,g),\ell_X(c,h)\bigr)f
=\frac{|u|^{r+1}\,\tilde\nu(x)}{|x|^{N/2+[\beta=\mathrm f]}}\;
\omega_{\gamma_W,\gamma_W}(g,h)\,\pctL_\beta\bigl(\alpha P_\beta(x)\bigr)f .
\]
\end{lemma}
\begin{proof}
We use (W2) in the following explicit form for a line split off one of the two spaces. For
characters $a,b$ put $\chi^{b}_k(x)=b(x/\bar x)\upsilon(x)^{-k}$ ($x\in\auA_E^\times$; the character of
a space of rank $k$). (\emph{Target line}) If $T$ has rank $k$, $V'=V\oplus\pcxH_X$ and $\psZ$ is the
zero term along $\auHom(T,Ee_X)$, then for $q\in\auU(T)(\auA)$
\[
\psZ\,\omega_{a,\,b\upsilon}(q,\ell_X(c,h))=\chi^{a}_k(c)^{-1}|c|^{-k/2}\,\omega_{a,b}(q,h)\,\psZ ;
\]
(\emph{source line}) if $T'$ has rank $k$ and $\psZ$ is the zero term along the maps $W'\to T'$
vanishing on $W\oplus Ee_A$, then $\psZ\,\omega_{a\upsilon,\,b}(m_A(u,g),q)=\chi^{b}_k(u)^{-1}|u|^{k/2}\,
\omega_{a,b}(g,q)\,\psZ$. Both are (W2): the element stabilises the isotropic subspace $I$ and induces
$(q,h)$, resp.\ $(g,q)$, on $I^\perp/I$; its determinant on $I(\auA)$ is $c^{-k}$ times a determinant of
$q$, resp.\ $\bar u^{k}$ times a determinant of $q$, and $\chi_I=\upsilon(\det)^{-1}|\det|^{1/2}$
combines with the change of the splitting characters ($\det\ell_X(c,h)=(\bar c/c)\det h$,
$\det m_A(u,g)=(u/\bar u)\det g$, the extra $\upsilon$ being evaluated there) to the stated scalar
(the first character of the splitting is that of the target group). Moreover, by (W1),
$\psrho(\alpha\lambda_\beta)\omega(p)=\omega(p)\psrho(p^{-1}\alpha\lambda_\beta)$ with
$p^{-1}\lambda_{\mathrm t}=cu\,\lambda_{\mathrm t}$ and $p^{-1}\lambda_{\mathrm f}=\overline{cu}^{-1}
\lambda_{\mathrm f}$, which turns $\alpha$ into $\alpha P_\beta(x)$.

$\beta=\mathrm t$: $\pctL_{\mathrm t}=\psZ_{(e_X,W)}\circ\pctsl_A$; first the source line on
$\auHom(W',V')$ ($k=r+2$, $a=\gamma_W$, $b=\gamma_{W'}$), then the target line on $\auHom(W,V')$ ($k=N$,
$a=\gamma_W$, $b=\gamma_W$). The scalar is $\chi^{\gamma_{W'}}_{r+2}(u)^{-1}|u|^{(r+2)/2}\cdot
\chi^{\gamma_W}_N(c)^{-1}|c|^{-N/2}$. Now $\chi^{\gamma_W}_N(x)=\nu(x/\bar x)^{-1}\upsilon(x/\bar x)^{N/2}
\upsilon(x)^{-N}=\tilde\nu(x)^{-1}$, because the character $\upsilon$ of $\auU(1)$ is normalised by
$\upsilon(x/\bar x)=\upsilon(x)^2$ (well defined since $\upsilon^2$ is trivial on $\auA_{E^+}^\times$;
$N$ is even); likewise
$\chi^{\gamma_{W'}}_{N+2}=\tilde\nu^{-1}$. If $u=1$ the first factor is $1$; if $r=N$ it is
$\tilde\nu(u)|u|^{(N+2)/2}=\tilde\nu(u)|u|^{r+1-N/2}$. In both cases the scalar is
$\tilde\nu(cu)|u|^{r+1}|cu|^{-N/2}$.

$\beta=\mathrm f$: first the target line on $\auHom(W',V')$ ($k=N+2$, $a=\gamma_{W'}$), then the source
line on $\auHom(W',V)\to\auHom(W,V)$ ($k=r$); the scalar is $\tilde\nu(c)|c|^{-(N+2)/2}\cdot
\chi^{\gamma_W}_r(u)^{-1}|u|^{r/2}$, which for $u=1$ or $r=N$ equals $\tilde\nu(cu)|u|^{r+1}
|cu|^{-N/2-1}$ ($r/2=r+1-N/2-1$ when $r=N$).

The hypothesis is used exactly to identify $\chi_{r+2}$, resp.\ $\chi_r$, at $u$ with $\tilde\nu^{-1}$ and
the power of $|u|$; for $r\ne N$, $u\ne1$ the formula is false.
\end{proof}

\begin{remark}
The restriction $u=1$ or $r=N$ is essential: for $r\ne N$ the constant term at $m_A(u,g)$ is not
independent of $u$. Only the cases used below occur: $u=1$ (Propositions~\ref{pcx:P-plane},
\ref{pcx:P-descent}, \ref{pcx:P-nonvan}) and $r=N$ (Lemma~\ref{pcx:const}).
\end{remark}

\subsection{Cone integrals and the statement}\label{pct:ss-statement}

Let $\pctK'=\pctK'_\infty\pctK'_f$ be the maximal compact subgroup of $\auU(V')(\auA)$ ($\pctK'_f$ the
stabiliser of the lattice, $\pctK'_\infty$ maximal compact in $\auU(V')(E^+\otimes\R)$), with
probability Haar measure $d\kappa$, and let $da$ be a Haar measure of $(E\otimes\R)^\times=
\prod_\tau\C^\times$. Write $\omega'_\infty$ for the archimedean component of
$\omega_{\gamma_{W'},\gamma_{W'}}$ and $\|\cdot\|_M$ for the moments of \S\ref{pcx:ss-conv}.

\begin{lemma}[Slices decay]\label{pct:slicedecay}
For all $N_0,M$ there is $C$ such that for every Fock vector $w$ of $\auHom(W',V')\otimes\R$,
$a\in E\otimes\R$, $x\in V'(E\otimes\R)$, $\kappa\in\pctK'_\infty$ and $\beta\in\{\mathrm t,\mathrm f\}$:
$(1+|a|_\infty)^M\|\pctL^\infty_\beta(a)w\|_{N_0}\le C\|w\|_{N_0+M}$,
$\|\pctsl_{A,\infty}(x)w\|_{N_0}\le C\|w\|_{N_0}$ and $\|\omega'_\infty(1,\kappa)w\|_{M}=\|w\|_M$, where
$|a|_\infty=\max_\tau|a_\tau|$.
\end{lemma}
\begin{proof}
By Lemma~\ref{pcx:superpos}, $w=\int w(y)\psrho_\infty(y)e\,dy$, so it suffices to bound the operators
on the coherent states $\psrho_\infty(y)e$ by $C(1+|y|)^{N_0+M}$, resp.\ $C(1+|y|)^{N_0}$, in the norm
$\|\cdot\|_{N_0}$. Let $I$ be an isotropic subspace of a real hermitian space $Y=I\oplus I'\oplus
Y_I$ ($I'$ dual to $I$). In the mixed model $\psHS_\infty(Y)=\mathcal S(I')\otimes\psHS_\infty(Y_I)$
the vacuum is a Gaussian $\varphi_0\otimes e_I$, $\psrho(y_1)$ for $y_1\in I^\perp$ acts by a phase
times $\psrho(\bar y_1)$ on the second factor, and $\psrho(y_2)$, $y_2\in I'$, translates the first
factor; $\psZ_I$ is evaluation of the first factor at $0$. Writing $y=y_1+y_2$, one gets
$\psZ_I\psrho(y)e=\epsilon\,\varphi_0(y_2)\psrho(\bar y_1)e_I$ with $|\epsilon|=1$, and
$\|\psrho(\bar y_1)e_I\|_{N_0}\le C(1+|y_1|)^{N_0}$. Applied to $\pctL^\infty_\beta(a)=
\psZ_{I_\beta}\psrho(a\lambda_\beta)$, the $I'$-component of $y+a\lambda_\beta$ contains $a$ in the
coordinate dual to $\lambda_{\beta^\vee}$, so $|\varphi_0(y_2)|\le Ce^{-c|a-y'|^2}$ with $y'$ a
coordinate of $y$, and $e^{-c|a-y'|^2}\le C_M(1+|a|)^{-M}(1+|y|)^M$. The slice is the case
$I=I_A$ without the factor in $a$. Finally $\pctK'_\infty$ fixes the line of $e$ (it lies in the
maximal compact subgroup of the symplectic group of $\auHom(W',V')\otimes\R$, which acts on $e$
by a character), so $\omega'_\infty(1,\kappa)\psrho(y)e=\epsilon'\psrho(\kappa y)e$ with $|\epsilon'|=1$
and $|\kappa y|=|y|$; this gives the equality of moments.
\end{proof}

\begin{definition}[Cone integrals]\label{pct:def-cone}
For a Fock vector $v$ of $\auHom(W',V')\otimes\R$ with finite moments and $x>0$ put
\[
\pctcone_x(v)=\int_{\pctK'_\infty\times(E\otimes\R)^\times}|a|^x\,\tilde\nu_\infty(a)\,
\pctL^\infty_{\mathrm t}(a)\,\omega'_\infty(1,\kappa)v\;d\kappa\,da,
\]
and $\pctcone^\vee_x(v)$ the same with $\pctL^\infty_{\mathrm f}$. (Here $\tilde\nu_\infty(a)=
\prod_\tau(a_\tau/|a_\tau|)^{2\mu_\tau}$ depends only on $\mu$.)
\end{definition}

By Lemma~\ref{pct:slicedecay} the integrand has $\|\cdot\|_{N_0}$-norm at most
$C|a|^x(1+|a|_\infty)^{-M}\|v\|_{N_0+M}$, and $\int|a|^x(1+|a|_\infty)^{-M}da<\infty$ for $x>0$ and
$M>2x[E^+:\Q]$ (in polar coordinates $|a|^x=\prod_\tau r_\tau^{2x}$ and $da=\prod_\tau
r_\tau^{-1}dr_\tau\,d\theta_\tau$). So the cone integrals converge absolutely in every norm
$\|\cdot\|_{N_0}$ and have finite moments.

For a datum $f$ of $\auHom(W',V')$, $g\in\auU(W)(\auA)$, $h\in\auU(V)(\auA)$, $\kappa\in\pctK'$
and $\alpha\in\auA_E$ put
\[
\Lambda_\beta(f;\alpha;g,h,\kappa)=\theta_{W,V}\bigl(\omega_{\gamma_W,\gamma_W}(g,h)\,\pctL_\beta(\alpha)\,
\omega_{\gamma_{W'},\gamma_{W'}}(1,\kappa)f\bigr),
\]
and, for $u\in\auA_E^\times$ and real $x$, the \emph{zeta integral}
\[
\pctZ(f,g,u;x)=\int_{E^\times\backslash\auA_E^\times}\int_{[\auU(V)]}\int_{\pctK'}\tilde\nu(au)|au|^x
\sum_{\xi\in E^\times}\Lambda_{\mathrm t}(f;\xi au;g,h,\kappa)\,d\kappa\,dh\,da,
\]
whenever it converges absolutely. The substitution $a\mapsto au^{-1}$ shows
$\pctZ(f,g,u;x)=\pctZ(f,g,1;x)$.

\begin{theorem}[The plane constant term]\label{pct:plane}
Let $N$ be even and $d=r-N/2$. There are constants $C_{\mathrm t},C_{\mathrm f}>0$ and an idele
$\alpha_0$, depending only on the splittings, on $\mu$ and on the Weil data, such that for every
tame $\nu$ of weight $\mu$ there are $\varepsilon_{\mathrm t},\varepsilon_{\mathrm f}\in\C$ with
$|\varepsilon_\beta|=C_\beta$ and the following holds for every Fock vector $v$ of
$\auHom(W',V')\otimes\R$ with finite moments, every $z\in\pcMat_{b\times N}(O)$, every
$g\in\auU(W)(\auA)$ and every $u\in\auA_E^\times$ with $u=1$ or $r=N$. Let
$f=[v,\Phi'_z]_S$ and $P(x)=|\alpha_0|^xQ^{x-1/2}$.
\begin{enumerate}
\item[(a)] There is $c_X\ne0$, depending only on Haar measures, such that the function
$\pctZ_{v,z,g}(x)=c_X\pctZ(f,g,u;x)$ is defined by an absolutely convergent integral for every
real $x$, does not depend on $u$, and is real-analytic on $\R$.
\item[(b)] $\bigl(\Theta^{\mathrm L}_{W',V'}(v,\Phi'_z)\bigr)_A(m_A(u,g))=\pctZ_{v,z,g}(d+1)$.
\item[(c)] For $x\ge1$:
$\pctZ_{v,z,g}(x)=\varepsilon_{\mathrm t}P(x)L^S(x,\nu)\,\allowbreak \Theta^{\mathrm L}_{W,V}(\pctcone_xv,
\Phi^{\mathrm{row}}_z)(g)$ and
$\pctZ_{v,z,g}(1-x)=\varepsilon_{\mathrm f}P(x)L^S(x,\nu)\,\allowbreak \Theta^{\mathrm L}_{W,V}(\pctcone^\vee_xv,
\Phi^{\mathrm{row}}_z)(g)$.
\end{enumerate}
In particular there are $C>0$ and, for each $\nu$, $\kappa_\nu\in\C$ with $|\kappa_\nu|=C\,
L^\sharp(d+1)$ if $d\ge0$, $|\kappa_\nu|=C\,L^\sharp(-d)$ if $d<0$, such that, for all $v,z,g$ and
$u$ as above,
\begin{equation}\label{pct:eq-plane}
\bigl(\Theta^{\mathrm L}_{W',V'}(v,\Phi'_z)\bigr)_A(m_A(u,g))=\begin{cases}
\kappa_\nu\,\Theta^{\mathrm L}_{W,V}(\pctcone_{d+1}v,\Phi^{\mathrm{row}}_z)(g),&d\ge0,\\
\kappa_\nu\,\Theta^{\mathrm L}_{W,V}(\pctcone^\vee_{-d}v,\Phi^{\mathrm{row}}_z)(g),&d<0.
\end{cases}
\end{equation}
If $r=N$, the left side does not depend on $u$.
\end{theorem}

The last two assertions follow from (a)--(c): for $d\ge0$ take $x=d+1\ge1$ in the first formula
of (c), for $d<0$ take $x=-d\ge1$ in the second, so that $1-x=d+1$; then $\kappa_\nu=
\varepsilon_\beta P(x)L^S(x,\nu)$ has modulus $C_\beta|\alpha_0|^xL^\sharp(x)$, and $x$ is fixed by
the spaces. Independence of $u$ for $r=N$ is (a)+(b). Note that only values of $L^S$ at points
$\ge1$ enter, and no functional equation is used. The rest of this subsection and the next prove
(a)--(c).

\subsection{Proof of Theorem~\ref{pct:plane}}\label{pct:ss-proof}

\emph{Step 1: expansion of the constant term.} For $y\in V'(E)$ and a datum $F$ of
$\auHom(W',V')$ put $T_y(F)=\int_{\pctN_A(E)\backslash\pctN_A(\auA)}\theta_{W,V'}(\pctsl_A(y)
\omega(n,1)F)\,dn$, with $\omega=\omega_{\gamma_{W'},\gamma_{W'}}$.

\begin{lemma}\label{pct:expand}
Let $x'\in\auU(W')(\auA)$ and $f=[v,\Phi'_z]_S$ as in Theorem~\ref{pct:plane}, where the Fock vector $v$
has moments of all orders ($\|v\|_N<\infty$ for every $N$). Then
$\bigl(\Theta^{\mathrm L}_{W',V'}(v,\Phi'_z)\bigr)_A(x')=\int_{[\auU(V')]}\sum_{y\in V'(E)}
T_y(\omega(x',h)f)\,dh$, and $T_y=0$ unless $\langle y,y\rangle=0$.
\end{lemma}
\begin{proof}
By (W3) for the isotropic subspace $I_A$ (whose dual is the column $e_A$, identified with $V'$ by
$y\mapsto y\otimes f_A^\vee$), $\theta(F)=\sum_{y\in V'(E)}\theta_{W,V'}(\pctsl_A(y)F)$, the series
converging absolutely. Let $\mathsf N_c\subset\pctN_A(\auA)$ be a compact set of representatives of
$\pctN_A(E)\backslash\pctN_A(\auA)$. For fixed $x',h$, Lemma~\ref{pz:absgrowth}(b) (the elements
$(nx',h)$, $n\in\mathsf N_c$, have finite parts in a compact set and bounded archimedean parts; the
archimedean vector $v$ has moments of all orders by hypothesis, so it lies in every $\pzF_{s'}$
(Lemma~\ref{pcx:superpos}(b)), hence is a smooth vector (Lemma~\ref{pcn:smooth}), and
Lemma~\ref{pz:absgrowth}(b) applies to it directly) bounds
$\sum_y|\theta_{W,V'}(\pctsl_A(y)\omega(nx',h)f)|$ uniformly for $n\in\mathsf N_c$, and Fubini's theorem
gives
$\int\theta(\omega(nx',h)f)\,dn=\sum_yT_y(\omega(x',h)f)$ for each $h$. The kernel
$\theta^{W',V'}_{\gamma_{W'},\gamma_{W'}}(nx',h;f)$ is bounded by $C\,\mathrm{ht}(nx')^B
\mathrm{ht}(h)^{-N_0}$ (Proposition~\ref{pcx:decay}(a): $\Phi'_z$ is supported on matrices of full
rank $r+2$ by Lemma~\ref{pcx:std-matched}); $\mathrm{ht}$ is invariant under rational points and
$\pctN_A(E)\backslash\pctN_A(\auA)x'$ is compact, so Fubini applies on
$(\pctN_A(E)\backslash\pctN_A(\auA))\times[\auU(V')]$ by (R1). This gives the formula. The centre of
$\pctN_A(\auA)$ consists of the elements $\phi\mapsto\phi+t\langle\cdot,e_A\rangle e_A$-type
transvections ($t+\bar t=0$), which act on $\pctsl_A(y)$ by the character
$t\mapsto\psi(t\langle y,y\rangle)$ ((W1) and the commutation relation of the Heisenberg group);
its mean over $E\backslash\auA$ vanishes unless $\langle y,y\rangle=0$.
\end{proof}

\emph{Step 2: the term $y=0$ vanishes.} The element $m_A(u,g)$ stabilises $I_A$, so by (W2)
$T_0(\omega(m_A(u,g),h)f)$ is a multiple of $\theta_{W,V'}(\omega(g,h)\pctsl_A(0)f)$ (the
radical acts trivially on the slice at $0$). At a place $w\in S$, in the model (W5), the zero term
along $I_A$ is the composite of an average over translations by $I_{A,w}$ and of the average over
the characters $x\mapsto[x,a]$, $a\in I_{A,\bar w}$; the latter restricts a function to the
annihilator of $I_{A,\bar w}$, which consists of the matrices $x$ with $xe_A=0$. Now
$\Phi'_z(x)\ne0$ forces $\binom{z\,\mathrm{pr}}{x}\in\auGL_{N+2}(O)$, while $z\,\mathrm{pr}\,e_A=0$ and
$xe_A=0$; so $\Phi'_z$ vanishes on that annihilator and $\pctsl_A(0)f=0$.

\emph{Step 3: unfolding along the isotropic lines.}
\begin{lemma}\label{pct:unfold}
Let $x'=m_A(u,g)$ with $u=1$ or $r=N$, and put $\Phi_{x'}(h)=\sum_{\xi\in E^\times}T_{\xi e_X}(\omega(x',h)f)$.
\begin{enumerate}
\item[(i)] For $\xi\ne0$, $T_{\xi e_X}(F)=\theta_{W,V}(\pctL_{\mathrm t}(\xi)F)$.
\item[(ii)] $\Phi_{x'}$ is left invariant under $\pctN_X(\auA)$ and under the rational points
$\pctP_X(E)$ of the stabiliser of $Ee_X$, and $\sum_{y\ne0,\ \langle y,y\rangle=0}T_y(\omega(x',h)f)=
\sum_{\delta\in\pctP_X(E)\backslash\auU(V')(E)}\Phi_{x'}(\delta h)$.
\item[(iii)] If the integral $\pctZ(f,g,u;d+1)$ converges absolutely, then
$\bigl(\Theta^{\mathrm L}_{W',V'}(v,\Phi'_z)\bigr)_A(x')=c_X\pctZ(f,g,u;d+1)$, with $c_X\ne0$
depending only on Haar measures.
\end{enumerate}
\end{lemma}
\begin{proof}
(i) $\pctN_A(\auA)$ modulo its centre is $W(\auA)$, acting on $\pctsl_A(\xi e_X)F$ through
$\psrho$ of the subspace $\auHom(W,E\xi e_X)=(e_X,W)$ (by (W1)); the centre acts trivially as
$\langle\xi e_X,\xi e_X\rangle=0$. By the first formula of (W2), the mean over
$(e_X,W)(E)\backslash(e_X,W)(\auA)$ of $\theta_{W,V'}\circ\psrho$ is $\theta_{W,V}\circ\psZ_{(e_X,W)}$,
and $\psZ_{(e_X,W)}\pctsl_A(\xi e_X)=\pctL_{\mathrm t}(\xi)$ (Definition~\ref{pct:def-line}).

(ii) An element of $\pctN_X(\auA)$ fixes $e_X$, stabilises $I_0+E\lambda_{\mathrm f}$ and induces the
identity on $\auHom(W,V)$; by (W1), (W2) it commutes with $\pctL_{\mathrm t}(\xi)$ (the character
$\chi_I$ being trivial on unipotent elements). For $\delta\in\auU(V')(E)$, (W1) and the
invariance $\theta\circ\omega(\delta)=\theta$ give $T_{\delta y}(\omega(1,\delta)F)=T_y(F)$; an element
of $\pctP_X(E)$ permutes the vectors $\xi e_X$, $\xi\in E^\times$. By Witt's theorem the nonzero
isotropic vectors of $V'(E)$ are the $\delta^{-1}\xi e_X$, $\delta\in\auU(V')(E)$, $\xi\in E^\times$,
the pair $(\pctP_X(E)\delta,\xi)$ being determined up to the action of $\pctP_X(E)$ on $\xi$; grouping
the vectors by the line they span gives the identity.

(iii) By Lemma~\ref{pct:levi} with $c=a$ (and (W1) to move $\omega(1,\kappa)$),
$|a|^{r+1}\Phi_{x'}(\ell_X(a,h)\kappa)=\sum_{\xi}\tilde\nu(au)|au|^{d+1}\allowbreak \Lambda_{\mathrm t}(f;\xi au;g,h,\kappa)$.
By the Iwasawa decomposition $\auU(V')(\auA)=\pctN_X(\auA)\,\ell_X(\auA_E^\times\times\auU(V)(\auA))
\,\pctK'$ and the modulus computed at the end of \S\ref{pct:ss-setup}, for a measurable function $\Phi$ on
$\auU(V')(\auA)$, left invariant under $\pctN_X(\auA)$ and $\pctP_X(E)$,
\[
\int_{\pctP_X(E)\backslash\auU(V')(\auA)}\Phi=c_X\int_{E^\times\backslash\auA_E^\times}\int_{[\auU(V)]}
\int_{\pctK'}|a|^{r+1}\Phi(\ell_X(a,h)\kappa)\,d\kappa\,dh\,da,
\]
whenever the right side converges absolutely; here $\pctP_X(E)=\pctN_X(E)\ell_X(E^\times\times
\auU(V)(E))$ and $\pctN_X(E)\backslash\pctN_X(\auA)$ has volume $1$ (the standard integration
formula for a parabolic subgroup; e.g.\ Moeglin--Waldspurger, \emph{Spectral decomposition and
Eisenstein series}). Applied to $|\Phi_{x'}|$ (whose right side is
bounded by the absolutely convergent $\pctZ$ integrand) and then to $\Phi_{x'}$, and combined with
(ii) and the unfolding $\int_{[\auU(V')]}\sum_{\delta\in\pctP_X(E)\backslash\auU(V')(E)}\Phi(\delta h)\,dh=
\int_{\pctP_X(E)\backslash\auU(V')(\auA)}\Phi$ (legitimate since $|\Phi|$ is integrable on the right),
Lemma~\ref{pct:expand} and Step~2 give the claim.
\end{proof}

\emph{Step 4: the line terms of the plane datum factor.} Let $\varpi\in O$ be the product of
uniformisers, $\psi$ the standard additive character, and $\alpha_0$ the finite idele denoted $\alpha$ in
\S\ref{ps:ss-weil} (we keep the name $\alpha_0$, as $\alpha$ is a variable idele below). This is the choice the theorem requires: the data $f=[v,\Phi'_z]_S$ are data of
\S\ref{pc:ss-data}, whose components off $S$ are the vacua of the lattices $\alpha_0\auHom(L_{W'},L_{V'})_v$
(self-dual because $\widehat\auO_E$ is self-dual for $\psi(\jmath\alpha_0\aubar\alpha_0\,\aubar zu)$,
Lemma~\ref{pgx:vac}), and whose components at $w\in S$ are supported on $a_w\pcMat(\auO_w)$ with
$a_w\auO_w=\alpha_{0,w}\auO_w$ (Definition~\ref{pgx:conv}). (Self-duality alone would fix only the products
$\alpha_{0,w}\overline{\alpha_{0,\bar w}}$; what is used below is that the scaling idele is the one defining the
data, and this holds by definition.) The components of $\alpha_0$ are units outside the primes dividing the
relative different of $E/K$ (proof of Proposition~\ref{pgv:alpha}(a)), and $\alpha_0$ depends only on the
Weil data. At $w\in S$ we use the coordinates of Definition~\ref{pgx:conv}
described after Lemma~\ref{pcx:mod}, in which the vacuum lattice is $\pcMat(\auO_w)$ and the
component of $\alpha_0$ at $w$ becomes the unit $\alpha_{0,w}/a_w$, its \emph{unit component} (equal to
$\alpha_{0,w}$ where $\alpha_{0,w}$ is a unit, as then $a_w=1$). For $\alpha\in\auA_E$ write $\alpha_S=(\alpha_w)_{w\in S}\in F$
and $\alpha_{\bar S}^\flat=(\overline{\alpha_{\bar w}})_{w\in S}\in F$, and put
\[
\lambda_\varpi(\alpha)=\begin{cases}\chi(\alpha_S)^{-1}\chi(\varpi\alpha^\flat_{\bar S}),&\alpha_f\text{ integral
off }S\cup\bar S,\ \alpha_S\in O^\times,\ \varpi\alpha^\flat_{\bar S}\in O^\times,\\0,&\text{otherwise},
\end{cases}
\]
and $\lambda_1$ the same with $\varpi$ replaced by $1$.

\begin{lemma}[Gauss sums]\label{pct:gauss}
Let $w\in S$, $\chi_w$ the tame part (a character of $O_w^\times$ trivial on $1+\varpi_wO_w$,
non-trivial), and $\psi_w$ any character of $E_w$ of conductor $O_w$ (it will be applied to
$\lambda_w^{\pm1}$ of Lemma~\ref{pgx:dict}(b)). There is $g_w\in\C$ with
$|g_w|^2=q_w^{-1}$ such that for all $c\in E_w$,
$\int_{O_w^\times}\chi_w^{-1}(t)\psi_w(ct)\,dt=g_w\,\chi_w(\varpi_wc)$ if $\varpi_wc\in O_w^\times$,
and $0$ otherwise ($dt$ of total mass $1$ on $O_w$).
\end{lemma}
\begin{proof}
If $c\in O_w$ the integrand's $\psi_w$-factor is $1$ and $\int_{O_w^\times}\chi_w^{-1}=0$. If
$v(c)\le-2$, replacing $t$ by $t(1+\varpi_ws)$, $s\in O_w$, does not change $\chi_w^{-1}(t)$ and
multiplies $\psi_w(ct)$ by $\psi_w(c\varpi_wst)$; averaging over $s$ gives $0$ as $c\varpi_wt\notin
O_w$. If $c=\varpi_w^{-1}c_0$, $c_0\in O_w^\times$, the integral is $q_w^{-1}\sum_{t\in\auF_w^\times}
\chi_w^{-1}(t)\psi_w(\varpi_w^{-1}c_0t)=q_w^{-1}\chi_w(c_0)G$ with the Gauss sum $G$ of $\chi_w^{-1}$
and the non-trivial character $t\mapsto\psi_w(\varpi_w^{-1}t)$ of $\auF_w$, $|G|^2=q_w$. Put
$g_w=q_w^{-1}G$.
\end{proof}

\begin{lemma}[Factorisation]\label{pct:factor}
There are $l_{\mathrm t}=l_{\mathrm t}(\nu),l_{\mathrm f}=l_{\mathrm f}(\nu)\in\C^\times$ whose moduli
$|l_{\mathrm t}|,|l_{\mathrm f}|$ depend only on the splittings and the Weil data, and for each $\nu$ a number $g_S$ with $|g_S|^2=Q^{-1}$, such that for all $v$, all
$z\in\pcMat_{b\times N}(O)$ and all $\alpha\in\auA_E$,
\[
\pctL_\beta(\alpha_0\alpha)[v,\Phi'_z]_S=l_\beta\,g_S\,\lambda_\varpi(\alpha)\,
[\pctL^\infty_\beta(\alpha_\infty)v,\Phi^{\mathrm{row}}_z]_S .
\]
Moreover the finite points $\pctK'_f$ fix $[v,\Phi'_z]_S$ in $\omega_{\gamma_{W'},\gamma_{W'}}$.
\end{lemma}
\begin{proof}
All operators are tensor products over the places, and so is the datum. At $\infty$ the factor is
$\pctL^\infty_\beta(\alpha_\infty)v$ by definition.

\emph{Places $v\notin S\cup\bar S\cup\infty$.} The datum is the vacuum $\psvac_L$ of the self-dual
lattice $L=\alpha_0\auHom(L_{W'},L_{V'})_v$, for which $\lambda_{\mathrm t},\lambda_{\mathrm f}$ and the blocks
of $I_\beta$ are spanned by vectors of $\auHom(L_{W'},L_{V'})_v$ (the splittings are adapted). Let $I=I_\beta$ and
$y=\alpha_0\alpha\lambda_\beta$. For $i\in I\cap L$, $\psZ_I\psrho(i)=\psZ_I$ and $\psrho(i)\psvac_L=
\psvac_L$ (with the normalisation $\psrho^\circ$ of (W4)), while
$\psrho(i)\psrho(y)=\psi([i,y])\psrho(y)\psrho(i)$; hence $\psZ_I\psrho(y)\psvac_L=0$ unless
$\psi([i,y])=1$ for all $i\in I\cap L$, i.e.\ unless $y\in L+I^\perp$, i.e.\ (as $\lambda_\beta$ pairs
with the primitive lattice vector $\lambda_{\beta^\vee}\in I$, and $\alpha_0$ makes $L$ self-dual)
unless $\alpha_v$ is integral. If it is, $\psZ_I\psrho(y)\psvac_L$ is a root of unity times
$\psZ_I\psvac_L$, the vacuum of the self-dual lattice of $I^\perp/I$; the root of unity depends only
on the Weil data and goes into $l_\beta$.

\emph{Places $w\in S$} (and their conjugates, which carry the same local factor $v=w\bar w$). In
the model (W5) data are functions on $\pcMat_{(r+2)\times(N+2)}(E_w)$ (matrices in the frames of
$W'$, $V'$). \emph{Change of coordinates.} Let $P=(P_W\,|\,P_{e}\,|\,P_f)\in\auGL_{N+2}(O_w)$ be the
matrix whose columns are the coordinates of the frame of $W$, $e_A$, $f_A$ in the frame of $W'$, and
$R=P^{-1}$, with block rows $R_W=\mathrm{pr}$, $R_e$, $R_f$; similarly $G\in\auGL_{r+2}(O_w)$ for
$(V,e_X,f_X)$ in $V'$. Both are integral with integral inverse because the splittings are adapted.
Write $x=G\tilde xR$; then $\tilde x$ is the block matrix with rows $(V,e_X,f_X)$ and columns
$(W,e_A,f_A)$ used above, the substitution preserves the measure (unimodular), and
\[
\begin{pmatrix}z\,\mathrm{pr}\\ x\end{pmatrix}=\begin{pmatrix}1&0\\0&G\end{pmatrix}
\begin{pmatrix}z\,\mathrm{pr}\,P\\ \tilde x\end{pmatrix}R,\qquad z\,\mathrm{pr}\,P=(z\,|\,0\,|\,0),
\]
so $\Phi'_z(x)=\epsilon_w\,\chi_w^{-1}\bigl(\det\binom{(z\,|\,0\,|\,0)}{\tilde x}\bigr)$ (with value $0$
when the matrix is not in $\auGL_{N+2}(O_w)$), where $\epsilon_w=\chi_w^{-1}(\det G\det R)$ is a constant of
modulus one (it depends on the frames and, through $\chi=\chi_S$, on $\nu$); it goes into $l_\beta$. The Heisenberg operators and the
isotropic subspaces of Definition~\ref{pct:def-line} are, in the coordinates $\tilde x$, the
translations and characters in the blocks indicated. As in Step~2, the zero term along an isotropic $I$ is $\Phi\mapsto(x\mapsto
\int_{I_w}\Phi(x+i)\,di)$ on the annihilator of $I_{\bar w}$, which is the set of matrices vanishing
in the blocks dual to those of $I$; $\psrho(y)$ is translation by $y_w$ times the character
$[\,\cdot\,,y_{\bar w}]$. For $\beta=\mathrm t$, $I=I_{\mathrm t}=(V,f_A)\oplus(e_X,W)\oplus(e_X,f_A)
\oplus(f_X,f_A)$, the dual blocks are $(V,e_A),(f_X,W),(f_X,e_A),(e_X,e_A)$, and the translation puts
$\alpha_w$ (times the unit component of $\alpha_0$) in the block $(e_X,e_A)$ and multiplies by
$\lambda_w(\epsilon'_w\overline{\alpha_{\bar w}}\,t)^{\pm1}$, $t$ the entry $(f_X,f_A)$, where
$\lambda_w(t)=\psi_w(\jmath(a\bar a)_wt)$ is the character of conductor $O_w$ of
Lemma~\ref{pgx:dict}(b) (the Heisenberg pairing in the coordinates $\tilde x=a^{-1}x$) and
$\epsilon'_w$ is the conjugate of the unit component of $\alpha_0$ at $\bar w$; replacing $c$ by
$\pm\epsilon'_wc$ in Lemma~\ref{pct:gauss} (applied to $\lambda_w^{\pm1}$) multiplies the result by
$\chi_w(\pm\epsilon'_w)$, of modulus one, which goes into $l_\beta$; the $\chi_w$-values of these unit components, which arise when they are removed from the
determinant, have modulus one and go into $l_\beta$. The resulting matrix is
\[
M=\begin{pmatrix}z&0&0\\ y&0&p\\ q&\alpha_w&s\\0&0&t\end{pmatrix}
\]
with $y\in\pcMat_{r\times N}$ the variable of $\auHom(W,V)$ and $p,q,s,t$ integrated over $O_w$
(the measure of $O_w$ being $1$). Expanding along the column $e_A$, $\det M=\pm\alpha_wt\det\binom zy$,
and $M\in\auGL_{N+2}(O_w)$ iff all entries are integral and $\alpha_w,t,\det\binom zy$ are units. So
the integrand is $\chi_w^{-1}(\pm1)\chi_w^{-1}(\alpha_w)\mathbf 1_{O_w^\times}(\alpha_w)\cdot
\Phi^{\mathrm{row}}_z(y)\cdot\chi_w^{-1}(t)\mathbf 1_{O_w^\times}(t)\lambda_w(\epsilon'_w\overline{\alpha_{\bar w}}t)^{\pm1}\cdot
\mathbf 1_{O_w}(p,q,s)$, and Lemma~\ref{pct:gauss} gives the factor
$g_w\chi_w(\alpha_w)^{-1}\chi_w(\varpi_w\overline{\alpha_{\bar w}})$ times $\Phi^{\mathrm{row}}_z(y)$ when
$\alpha_w$ and $\varpi_w\overline{\alpha_{\bar w}}$ are units, and $0$ otherwise. For $\beta=\mathrm f$
the same computation with $I_{\mathrm f}$ puts $t$ in the corner $(e_X,e_A)$ and $\alpha_w$ in the
corner $(f_X,f_A)$; the determinant and the result are the same. Put $g_S=\prod_{w\in S}g_w$.

\emph{Invariance under $\pctK'_f$.} At $w\in S$, $k\in\auGL_{r+2}(O_w)$ acts by $\Phi\mapsto
\gamma_{W'}(\det k)\,\Phi(k^{-1}\cdot)$ (W5), and $\Phi'_z(k^{-1}x)=\chi_w(\det k)\Phi'_z(x)$ (the
row datum is matched, Lemma~\ref{pcx:std-matched}, and Lemma~\ref{pcx:rightinv}), while
$\gamma_{W'}(\det k)=\chi_w(\det k)^{-1}$ on units ($\upsilon$ being trivial on the determinants of
integral points); elsewhere $\pctK'_f$ fixes the vacuum by (W4).
\end{proof}

\emph{Step 5: Tate's integral.} The function $\lambda_\varpi$ satisfies $\lambda_\varpi(\alpha)=
\lambda_1(\alpha c)$ for the finite idele $c$ equal to $\bar\varpi_w$ at the places $\bar w\in\bar S$
and to $1$ elsewhere, and $|c^{-1}|=Q$. On the units $\widehat\auO_E^\times$ of the finite ideles,
$\lambda_1(u)\tilde\nu_f(u)=1$ (at $w\in S$: $\tilde\nu_w(u_w)\tilde\nu_{\bar w}(u_{\bar w})=
\chi_w(u_w)\chi_w(\overline{u_{\bar w}})^{-1}$; elsewhere $\tilde\nu$ is unramified), and $\lambda_1(a)=0$
unless $a$ is integral and a unit at $S\cup\bar S$.

\begin{lemma}\label{pct:fintate}
Let $\rho$ be a Haar measure on $\auA_{E,f}^\times$. For $x>1$,
$\int_{\auA_{E,f}^\times}\lambda_1(a)\tilde\nu_f(a)|a|^x\,d\rho(a)=L^S(x,\nu)\,\rho(\widehat\auO_E^\times)$,
absolutely convergent.
\end{lemma}
\begin{proof}
Write $a=u\varpi^m$ with $u$ a unit and $\varpi^m=\prod_{v\notin S\cup\bar S}\varpi_v^{m_v}$, $m_v\ge0$
almost all $0$. The integrand is $\tilde\nu(\varpi^m)|\varpi^m|^x$, so the integral is
$\rho(\widehat\auO_E^\times)\prod_{v\notin S\cup\bar S}(1-\tilde\nu(\varpi_v)q_v^{-x})^{-1}$.
\end{proof}

\begin{lemma}[Mellin transforms]\label{pct:mellin}
Let $T$, $\mathsf n>0$ be measurable on a measure space with $\int|T|\mathsf n^R<\infty$ and
$\int|T|\mathsf n^{1-R}<\infty$ for every $R>1$. Then $\int|T|\mathsf n^x<\infty$ for every real $x$
and $x\mapsto\int T\mathsf n^x$ is real-analytic on $\R$.
\end{lemma}
\begin{proof}
$\mathsf n^x\le\mathsf n^{1-R}+\mathsf n^R$ for $x\in[1-R,R]$. Near $x_0$, with $0<\delta<1$,
$|\log\mathsf n|^k\mathsf n^{x_0}\le k!\delta^{-k}(\mathsf n^{x_0+\delta}+\mathsf n^{x_0-\delta})$, so the
series $\sum_k\frac{(x-x_0)^k}{k!}\int T(\log\mathsf n)^k\mathsf n^{x_0}$ converges absolutely for
$|x-x_0|<\delta$ and equals $\int T\mathsf n^x$.
\end{proof}

\begin{lemma}[Poisson on the line, in the variable $a$]\label{pct:poisson-a}
For $f=[v,\Phi'_z]_S$, $a\in\auA_E^\times$ and $(g,h,\kappa)$:
$|a|\sum_{\xi\in E^\times}\Lambda_{\mathrm t}(f;\xi a;g,h,\kappa)=\sum_{\xi\in E^\times}
\Lambda_{\mathrm f}(f;\xi\bar a^{-1};g,h,\kappa)$.
\end{lemma}
\begin{proof}
By Lemma~\ref{pct:factor}, $\Lambda_\beta(f;0;\cdot)=0$, since $\lambda_\varpi(0)=0$. Apply
Lemma~\ref{pct:poisson} to $F=\omega(m_A(1,g),\ell_X(a,h))\omega(1,\kappa)f$ and use
Lemma~\ref{pct:levi} with $u=1$, $c=a$: the terms become $\tilde\nu(a)|a|^{-N/2}\Lambda_{\mathrm t}(\xi a)$
and $\tilde\nu(a)|a|^{-N/2-1}\Lambda_{\mathrm f}(\xi\bar a^{-1})$; cancel $\tilde\nu(a)|a|^{-N/2-1}$.
\end{proof}

\begin{proof}[Proof of Theorem~\ref{pct:plane}]
Let $K_0=c^{-1}\alpha_0\in\auA_E^\times$ and $\mathcal K(g,h;w)=\theta^{W,V}_{\gamma_W,\gamma_W}(g,h;
[w,\Phi^{\mathrm{row}}_z]_S)$. By Lemma~\ref{pct:factor} (the finite part of $\kappa$ acting trivially),
for $\alpha\in\auA_E^\times$,
\begin{multline*}
\tilde\nu(\alpha K_0)|\alpha K_0|^x\Lambda_\beta(f;\alpha K_0;g,h,\kappa)\\
=\Bigl[\tilde\nu_\infty(\alpha_\infty)|\alpha_\infty|^x\mathcal K\bigl(g,h;\pctL^\infty_\beta(\alpha_\infty)
\omega'_\infty(1,\kappa_\infty)v\bigr)\Bigr]\cdot\Bigl[l_\beta g_S\tilde\nu(K_0)|K_0|^x\,
\lambda_1(\alpha_f)\tilde\nu_f(\alpha_f)|\alpha_f|^x\Bigr].
\end{multline*}
\emph{Convergence for $x>1$, side $\mathrm t$.} In $\pctZ(f,g,u;x)$ substitute $\xi au=\alpha K_0$;
the sum over $E^\times$ and the integral over $E^\times\backslash\auA_E^\times$ combine into an integral
over $\auA_E^\times=(E\otimes\R)^\times\times\auA_{E,f}^\times$. By Tonelli the integral of the absolute
values is the product of $\int_{\auA_{E,f}^\times}|\lambda_1\tilde\nu_f||\cdot|^x\,d\rho<\infty$
(Lemma~\ref{pct:fintate}) and of $\int|\alpha_\infty|^x|\mathcal K(g,h;\pctL^\infty_{\mathrm t}(\alpha_\infty)
\omega'_\infty(1,\kappa)v)|\,d\alpha_\infty\,dh\,d\kappa$, which is finite: by
Proposition~\ref{pcx:decay}(a) ($\Phi^{\mathrm{row}}_z$ is supported on matrices of rank $r$),
$|\mathcal K(g,h;w)|\le C\,\mathrm{ht}(g)^B\mathrm{ht}(h)^{-N_0}\|w\|_{M'}$, and Lemma~\ref{pct:slicedecay}
and (R1) bound the rest. Hence the integral converges absolutely and, by Fubini and
Lemma~\ref{pct:fintate},
\[
\pctZ(f,g,u;x)=l_{\mathrm t}g_S\tilde\nu(K_0)|K_0|^x\rho(\widehat\auO_E^\times)L^S(x,\nu)\int_{[\auU(V)]}
\mathcal K(g,h;\pctcone_xv)\,dh,
\]
where the cone integral was taken inside $\mathcal K$ (a continuous linear functional of $w$ for
the norm $\|\cdot\|_{M'}$, Proposition~\ref{pcx:decay}(a)). The last integral is
$\Theta^{\mathrm L}_{W,V}(\pctcone_xv,\Phi^{\mathrm{row}}_z)(g)$.

\emph{Side $\mathrm f$.} By Lemma~\ref{pct:poisson-a}, the integrand of $\pctZ(f,g,u;1-x)$ is
$\tilde\nu(au)|au|^{-x}\sum_\xi\Lambda_{\mathrm f}(f;\xi P_{\mathrm f}(au);\cdot)$, and
$\tilde\nu(P_{\mathrm f}(y))=\tilde\nu(y)$, $|P_{\mathrm f}(y)|=|y|^{-1}$. The substitution
$y'=P_{\mathrm f}(au)$ preserves the Haar measure of $E^\times\backslash\auA_E^\times$, and the same
computation with $\beta=\mathrm f$ shows absolute convergence for $x>1$ and
$\pctZ(f,g,u;1-x)=l_{\mathrm f}g_S\tilde\nu(K_0)|K_0|^x\rho(\widehat\auO_E^\times)L^S(x,\nu)\,
\Theta^{\mathrm L}_{W,V}(\pctcone^\vee_xv,\Phi^{\mathrm{row}}_z)(g)$.

\emph{(a).} The integrand of $\pctZ$ is $T\cdot|au|^x$ with $T$ independent of $x$; by the two
previous paragraphs Lemma~\ref{pct:mellin} applies. Independence of $u$ was noted before the
theorem.

\emph{(b).} Lemma~\ref{pct:unfold}(iii), with the absolute convergence at $d+1$ from (a).

\emph{(c).} For $x>1$ these are the two formulas just proved, with
$\varepsilon_\beta=c_Xl_\beta g_S\sqrt Q\,\tilde\nu(K_0)\rho(\widehat\auO_E^\times)$: indeed
$|K_0|^x=Q^x|\alpha_0|^x$ and $|g_S|=Q^{-1/2}$, so $|\varepsilon_\beta|=|c_Xl_\beta|\rho(\widehat
\auO_E^\times)=:C_\beta$ does not depend on $\nu$, and $c_X\,l_\beta g_S\tilde\nu(K_0)|K_0|^x=
\varepsilon_\beta P(x)$. At $x=1$ both sides are continuous on $[1,\infty)$: the left side by (a),
$L^S(x,\nu)=L(x,\tilde\nu)$ because $L(s,\tilde\nu)$ is entire ($\tilde\nu$ is a non-trivial unitary
Hecke character, Lemma~\ref{pcx:Lpos}; Hecke, see Tate's thesis, Main Theorem~4.4.1), and the
theta integral of the cone integral because it is, by the computation above, a Mellin transform
$\int\mathcal K(g,h;\cdots)|\alpha_\infty|^x$ converging absolutely for all $x>0$,
hence continuous on $(0,\infty)$ (the proof of Lemma~\ref{pct:mellin} applies on every compact
subinterval).
\end{proof}

\subsection{Strong approximation and the constancy in equal rank}\label{pct:ss-const}

Let $Y=Y_0\oplus\pcxH$ be an adapted splitting, $\pcxH=Ee\oplus Ef$, $\langle e,f\rangle=1$, the form
linear in the first variable. For $c\in Y_0$ and $\beta\in E$ with $\beta+\bar\beta=-\langle c,c\rangle$
(over $E$ or over $\auA_E$) let $n(c,\beta)$ and $n^-(c,\beta)$ be the isometries
\[
n(c,\beta)\colon e\mapsto e,\ y\mapsto y-\langle y,c\rangle e,\ f\mapsto f+c+\beta e;\qquad
n^-(c,\beta)\colon f\mapsto f,\ y\mapsto y-\langle y,c\rangle f,\ e\mapsto e+c+\beta f
\]
($y\in Y_0$); one checks directly that they preserve the form. The $n(c,\beta)$ form the unipotent
radical $\pctN$ of the stabiliser $\pctP$ of $Ee$, the $n^-(c,\beta)$ that of the opposite parabolic,
and $m(t,\delta)$ ($t\in\auA_E^\times$, $\delta\in\auU(Y_0)$) acts by $e\mapsto te$, $f\mapsto\bar t^{-1}f$,
$y\mapsto\delta y$; every element of $\pctP$ is $n\,m(t,\delta)$, uniquely. For a set $R$ of finite
places of $E$ and $x\in\auU(Y)(\auA)$ let $\ell_R(x)$ be the element with the components of $x$ at the
places of $E^+$ below $R\cup\bar R$ and $1$ elsewhere (also at $\infty$); $\ell_R$ is a homomorphism,
and $\auU(Y)(F_R)=\ell_R(\auU(Y)(\auA))$. Let $\pctK=\prod_v\pctK_v$ be the standard maximal compact
subgroup (lattice stabilisers at finite $v$) and $\pctK^S=\{k\in\prod_{v<\infty}\pctK_v:\ell_S(k)=1\}$.

\begin{lemma}[Quasi-reflections from radicals]\label{pct:qref}
Let $v\in Y_0(E)$, $\sigma\in E^\times$ with $\sigma+\bar\sigma=\langle v,v\rangle$, and
$\rho(y)=y-v\sigma^{-1}\langle y,v\rangle$ the quasi-reflection. Let $j\in E^\times$ with $\bar j=-j$,
$\beta_1=-\sigma-j^{-1}$, $\lambda=(j\bar\sigma)^{-1}$. Then
\[
u=n^-(v,\beta_1)\,n(0,j)\,n^-(-\lambda v,\lambda\bar\beta_1)\in\pctP(E),\qquad u=n\,m(\lambda,\rho)\
\text{ for some }n\in\pctN(E).
\]
\end{lemma}
\begin{proof}
The parameters are admissible: $\beta_1+\bar\beta_1=-(\sigma+\bar\sigma)-(j^{-1}+\bar j^{-1})=
-\langle v,v\rangle$, $j+\bar j=0$, and for $\beta_3=\lambda\bar\beta_1$, $c_3=-\lambda v$ one needs
$\beta_3+\bar\beta_3=-\langle c_3,c_3\rangle$, i.e.\
$\lambda\bar\beta_1+\bar\lambda\beta_1=\lambda\bar\lambda(\beta_1+\bar\beta_1)$, i.e.\
$\bar\beta_1\bar\lambda^{-1}+\beta_1\lambda^{-1}=\beta_1+\bar\beta_1$. Here $\lambda^{-1}=j\bar\sigma=
1-j\bar\beta_1$ (as $1-j\bar\beta_1=1+j\bar\sigma+j\bar j^{-1}$ and $j\bar j^{-1}=-1$) and
$\bar\lambda^{-1}=1-\bar j\beta_1=1+j\beta_1$; so the left side is $\bar\beta_1(1+j\beta_1)+\beta_1(1-j\bar\beta_1)=\beta_1+\bar\beta_1$. Now compute. With
$c_3=-\lambda v$, $\beta_3=\lambda\bar\beta_1$: $n^-(c_3,\beta_3)e=e+c_3+\beta_3f$; then $n(0,j)$ gives
$(1+j\beta_3)e+c_3+\beta_3f$, and $1+j\beta_3=\lambda$ (as $\lambda(1-j\bar\beta_1)=1$); then
$n^-(v,\beta_1)$ gives $\lambda(e+v+\beta_1f)+c_3-\langle c_3,v\rangle f+\beta_3f=\lambda e+(\lambda\beta_1+
\lambda\langle v,v\rangle+\lambda\bar\beta_1)f=\lambda e$. So $ue=\lambda e$ and $u\in\pctP(E)$. For
$y\in Y_0$: $n^-(c_3,\beta_3)y=y+\bar\lambda\langle y,v\rangle f$, then $n(0,j)$ gives
$y+\bar\lambda\langle y,v\rangle(f+je)$, and the $Y_0$-component after $n^-(v,\beta_1)$ is
$y+j\bar\lambda\langle y,v\rangle v$. Since $j\bar\lambda=j(\bar j\sigma)^{-1}=-\sigma^{-1}$, this is
$\rho(y)$. Hence the Levi component of $u$ is $m(\lambda,\rho)$.
\end{proof}

\begin{lemma}[Constancy by strong approximation]\label{pct:SA}
Let $S$ be a non-empty finite set of finite places of $E$, and $D\colon\auU(Y)(\auA)\to\C$ continuous
with: (i) $D(\gamma x)=D(x)$ for $\gamma\in\auU(Y)(E)$; (ii) $D(xs)=D(x)$ for $s\in\auU(Y)(F_S)$;
(iii) $D(xk)=D(x)$ for $k\in\pctK^S$; (iv) if $D(nx)=D(x)$ for all $n\in\pctN(\auA)$ and all $x$,
then $D(m(t,y))$ depends only on $y$. Then $D$ is constant.
\end{lemma}
\begin{proof}
\emph{(1)} If $a,b\in\auU(Y)(\auA)$ agree at all places not below $S\cup\bar S$, then $D(a)=D(b)$:
$a=b\,\ell_S(b^{-1}a)$ and (ii).

\emph{(2)} $D$ is left invariant under $\pctN(\auA)$ and $\pctN^-(\auA)$. The parameters $(c,\beta)$
form a closed subset of $Y_0(\auA)\times\auA_E$, and by strong approximation for the additive
group ($E+\prod_{w\in S\cup\bar S}E_w$ is dense in $\auA_E$, as $S\ne\emptyset$; Cassels--Fr\"ohlich,
\emph{Algebraic Number Theory}, strong approximation) applied to the coordinates of $c$ and to
the solutions $\beta$, the admissible parameters of the form
$(c_E+c_S,\beta_E+\beta_S)$, with $(c_E,\beta_E)$ rational and $(c_S,\beta_S)$ supported at $S\cup\bar S$,
are dense. Indeed, let $(c,\beta)$ be admissible, $\beta+\bar\beta=-\langle c,c\rangle$. Choose $c_E\in Y_0(E)$
and $c_S$ supported at $S\cup\bar S$ with $c'=c_E+c_S$ close to $c$. Then
$\delta=\langle c',c'\rangle-\langle c_E,c_E\rangle=\langle c_E,c_S\rangle+\langle c_S,c_E\rangle+\langle c_S,c_S\rangle\in\auA_{E^+}$
is supported at the places below $S$; put $\beta'_E=-\frac12\langle c_E,c_E\rangle\in E^+$ and
$\beta'_S=-\frac12\delta$, so that $(c',\beta'_E+\beta'_S)$ is admissible. The difference
$\beta-\beta'_E-\beta'_S$ is $\tau+\eta$ with $\tau=-\frac12(\langle c,c\rangle-\langle c',c'\rangle)$, small,
and $\eta$ of trace zero. The trace-zero elements form the line $j\auA_{E^+}$ for a fixed $j\in E^\times$
with $\bar j=-j$, and by strong approximation for $\auA_{E^+}$ (the set of places below $S$ is non-empty)
there are $\eta_E\in jE^+$ and $\eta_S\in j\auA_{E^+}$ supported below $S$ with $\eta-\eta_E-\eta_S$ small.
Then $(c',\beta'_E+\eta_E+\beta'_S+\eta_S)$ is admissible, of the required form, and close to $(c,\beta)$. For those, $n(c_E+c_S,\beta_E+\beta_S)x$ agrees with $n(c_E,\beta_E)x$ off $S\cup\bar S$, so
(1) and (i) give $D(n(\cdot)x)=D(x)$; continuity extends this to all of $\pctN(\auA)$. The same for
$\pctN^-$.

\emph{(3)} By (2) and (iv) there is $D_0$ with $D(m(t,y))=D_0(y)$.

\emph{(4)} Let $R$ be a set of finite places. The set $\Gamma_R$ of $\delta\in\auU(Y_0)(E)$ for which
there is $t\in\auA_E^\times$ with $D(\ell_R(m(t,\delta))x)=D(x)$ for all $x$ is a subgroup ($\ell_R$
and $m$ are homomorphisms). It contains the quasi-reflections: with $u=n_1n_2n_3=n\,m(\lambda,\rho)$
from Lemma~\ref{pct:qref}, $\ell_R(m(\lambda,\rho))=\ell_R(n)^{-1}\ell_R(n_1)\ell_R(n_2)\ell_R(n_3)$, a
product of elements of $\pctN(\auA)$ and $\pctN^-(\auA)$ (cutting the parameters to $R$ gives
admissible adelic parameters), under which $D$ is left invariant by (2). As $\auU(Y_0)(E)$ is
generated by quasi-reflections (Witt's theorem for hermitian forms, proved by quasi-reflections;
see e.g.\ Hahn--O'Meara, \emph{The Classical Groups and K-Theory}),
$\Gamma_R=\auU(Y_0)(E)$.

\emph{(5)} For finite $R$, all $t$ and all $y\in\auU(Y_0)(\auA)$: $D(\ell_R(m(t,y)))=D(1)$. By weak
approximation for $\auU(Y_0)$ (a rational variety by the Cayley transform; Platonov--Rapinchuk,
\emph{Algebraic Groups and Number Theory}, strong approximation) the elements $\gamma k$ with
$\gamma\in\auU(Y_0)(E)$ and $\ell_R(k)=1$ are dense, and $y\mapsto D(\ell_R(m(t,y)))$ is continuous, so
it suffices to take $y=\gamma k$, for which $\ell_R(m(t,y))=\ell_R(m(t,\gamma))=m(\ell_R(t),\ell_R(\gamma))$.
By (3) its value is $D_0(\ell_R(\gamma))=D(\ell_R(m(a,\gamma)))$ for the $a$ given by (4), which is
$D(1)$ (take $x=1$ in (4)).

\emph{(6)} By weak approximation for $\auU(Y)$ at the infinite places, $\auU(Y)(E)\,\auU(Y)(\auA_f)$ is
dense; by (i) and continuity it suffices to show $D(k)=D(1)$ for $k\in\auU(Y)(\auA_f)$. Let $T$ be the
cofinite set of finite places with $k_v\in\pctK_v$, and use the Iwasawa decomposition
$\auU(Y)(\auA)=\pctP(\auA)\pctK$ (valid for the lattice stabilisers since the splitting is adapted:
the lattice contains the hyperbolic pair $e,f$): $k\kappa=n\,m(t,y)$ with $\kappa\in\pctK$. The element
$b=k^{-1}\ell_{T^c}(k\kappa)$ has components $\kappa_v$ at $v\in T^c$, $k_v^{-1}\in\pctK_v$ at $v\in T$ and
$1$ at $\infty$; so $b=s\,k'$ with $s=\ell_S(b)\in\auU(Y)(F_S)$ and $k'\in\pctK^S$. By (ii), (iii),
$D(k)=D(kb)=D(\ell_{T^c}(k\kappa))=D(\ell_{T^c}(n)\,\ell_{T^c}(m(t,y)))=D(\ell_{T^c}(m(t,y)))=D(1)$, by
(2) and (5) with $R=T^c$ finite.
\end{proof}

\begin{proof}[Proof of Lemma~\ref{pcx:const}]
Let $Y$ have rank $2k$, $Y=Y_0\oplus\pcxH$ adapted, $v$ with finite moments, and
$F(x)=\Theta^{\mathrm L}_{Y,Y}(v,\Phi)(x)$ with $\Phi=\Psi^{\mathrm{col}}_1=\chi^{-1}\circ\det$ on
$\auGL_{2k}(O)$ (zero elsewhere), which is the row datum $\Phi^{\mathrm{row}}_\emptyset$ for $b=0$. We
check the hypotheses of Lemma~\ref{pct:SA} for $D=F$.

(i) and continuity: the kernel is left invariant in $x$ under $\auU(Y)(E)$ by (W1) ($\gamma_Y$ is
trivial on rational determinants), and the integral over $[\auU(Y)]$ converges uniformly for $x$ in
compact sets (Corollary~\ref{pcx:conv}(i): $\Phi$ is supported on invertible matrices).

(ii): Let $g\in\auGL_{2k}(F)=\auU(Y)(F_S)$ (read in the frames, $S$ consisting of split places).
In the notation of \S\ref{pcx:ss-match} with $b=0$, $a=2k$: $\mathrm{Inj}_0$ is a point, and
$\mathsf T_1=1$. For $\Phi_g=\xi(\det g)^{-1}\Phi(\cdot\,g^{-1})$ one computes, substituting
$r'=rxg^{-1}\in\auGL_{2k}(O)$, $\mathsf S_{\Phi}(x)=\xi(\det x)^{-1}\int_{\auGL_{2k}(O)}\xi(\det r')
\chi^{-1}(\det r')\,dr'=\xi(\det x)^{-1}$ and $\mathsf S_{\Phi_g}(x)=\xi(\det g)^{-1}\mathsf S_\Phi(xg^{-1})=
\xi(\det x)^{-1}$ for $x\in\auGL_{2k}(F)$ (using $\xi=\chi$ on $O^\times$). Since $k^{[a]}=k^{-1}$ for
$b=0$, both $(\Phi,1)$ and $(\Phi_g,1)$ are same-section pairs (Definition~\ref{pcx:def-section}).
Lemma~\ref{pcx:linfactor}, applied to $\Lambda=\Theta^{\mathrm L}_{Y,Y}(v,\cdot)$ (whose hypothesis is
Lemma~\ref{pcx:Sact}(a)), gives $\Lambda(\Phi_g)=\mathcal T(1)=\Lambda(\Phi)$. With
Lemma~\ref{pcx:Sact}(b) for $b=0$: $F(xg)=\Lambda(\Phi_g)(xg)=\Lambda(\Phi)(x)=F(x)$.

(iii): For $k\in\pctK^S$, $\omega_{\gamma_Y,\gamma_Y}(k,1)[v,\Phi]_S=[v,\Phi]_S$: the components of $k$
at $S\cup\bar S$ and $\infty$ are $1$, the vacuum is fixed by (W4) at the other finite places, and
$\gamma_Y(\det k)=1$ ($\nu$ and $\upsilon$ are trivial on the determinants of integral points off
$S\cup\bar S$). Hence $F(xk)=F(x)$.

(iv): If $F$ is left invariant under $\pctN(\auA)$, then $F=F_A$ (the constant term along the line
$Ee$), and by the last assertion of Theorem~\ref{pct:plane} applied with $W'=V'=Y$, $W=V=Y_0$,
$b=0$ ($r=N=2k-2$, $d=k-1\ge0$), $F(m(t,y))=F_A(m(t,y))$ does not depend on $t$.

By Lemma~\ref{pct:SA} ($S\ne\emptyset$ as $\pcm\ne(1)$), $F$ is constant.
\end{proof}

\subsection{The image of the cone integral}\label{pct:ss-cone}

Keep the set-up of \S\ref{pct:ss-setup}. At the infinite places write $Z_X$ for the archimedean
zero term along $(e_X,W)=\auHom(W,Ee_X)\subset\auHom(W,V')$, from Fock vectors of
$\auHom(W,V')\otimes\R$ to those of $\auHom(W,V)\otimes\R$; thus $\pctL^\infty_{\mathrm t}(a)=
Z_X\circ\pctsl_{A,\infty}(ae_X)$. The following two facts about the real mixed models are used.

\begin{lemma}[Real mixed models]\label{pct:realmixed}
\textup{(a)} For every Fock vector $v$ of $\auHom(W,V)\otimes\R$ with finite moments there is a Fock
vector $v_B$ of $\auHom(W,V')\otimes\R$ with finite moments and $Z_Xv_B=v$.
\textup{(b)} Let $\beta$ be a smooth compactly supported function on $V'(E\otimes\R)$, $x\mapsto k(x)$ a
smooth map to $\auU(V')(E^+\otimes\R)$, and $w$ a Fock vector of $\auHom(W,V')\otimes\R$ with finite
moments. There is a Fock vector $v'$ of $\auHom(W',V')\otimes\R$ with finite moments such that
$\pctsl_{A,\infty}(x)v'=\beta(x)\,\omega_\infty(1,k(x))w$ for all $x$.
\end{lemma}
\begin{proof}
We use the mixed model and its properties (M1)--(M4) of \S\ref{pcu:ss-mixed} (with (M1) proved in
\S\ref{pcu:ss-M1}, and Lemma~\ref{pcu:slicedecay}).

(b) This is (M1) for the adapted splitting $A$ of the source space $W'$ (with $p=1$, $T=V'$): the
slice $\mathrm{sl}_A(x)$ there is the value of the mixed model at the value $x$ on $e_A$, i.e.\
$\pctsl_{A,\infty}(x)$ of Definition~\ref{pct:def-line} (translation by $x$ in the column $e_A$
followed by the zero term along the column $f_A$, which is evaluation of the Schr\"odinger variable,
\S\ref{pcu:ss-mixed}); take $\tau=\beta$ and $u(x)=(1,k(x))$.

(a) $Z_X$ is a zero term along a line of the \emph{target}. The flip $\mathrm{Fl}$ of
\S\ref{pcx:ss-flip} (fact (W6)), at the infinite places, is induced by the measure-preserving
isometry $w\mapsto w^\dagger$ of $\auHom(W,V')_\R$ onto $\auHom(V'^-,W)_\R$, so it preserves the moments
$\|\cdot\|_M$ and intertwines the Heisenberg operators; it carries $(e_X,W)=\auHom(W,Ee_X)$ onto the
maps $V'^-\to W$ vanishing on $(Ee_X)^{\perp}$, i.e.\ the isotropic subspace attached to the adapted
splitting $V'^-=V^-\oplus\pcxH_X$ of the \emph{source} $V'^-$, and hence carries $Z_X$ to the zero term
$\mathrm{sl}(0)$ along that splitting (both are characterised by (W2) through the same isotropic
subspace). So it suffices to show that $\mathrm{sl}(0)$ is onto on vectors with finite moments, for a
source splitting. Given $v$ with finite moments, let $\tau$ be smooth with compact support and
$\tau(0)=1$; by (M1) with $u\equiv1$ there is $\Lambda$ with finite moments and
$\mathrm{sl}(x)\Lambda=\tau(x)v$; then $\mathrm{sl}(0)\Lambda=v$. Put $v_B=\mathrm{Fl}^{-1}\Lambda$.
\end{proof}

\begin{lemma}[Equivariance of the real zero term]\label{pct:zeroterm-equiv}
Let $p\in\auU(V')(E^+\otimes\R)$ with $pe_X=ae_X$, $a\in(E\otimes\R)^\times$. There is
$h\in\auU(V)(E^+\otimes\R)$ such that $Z_X\,\omega_\infty(1,p)=\tilde\nu_\infty(a)^{-1}|a|^{N/2}\,
\omega_\infty(1,h)\,Z_X$ on Fock vectors with finite moments, the representations being the
archimedean components of $\omega_{\gamma_W,\gamma_{W'}}$ on $\auHom(W,V')$ and of
$\omega_{\gamma_W,\gamma_W}$ on $\auHom(W,V)$.
\end{lemma}
\begin{proof}
$p$ stabilises $Ee_X$, hence $(e_X,W)$, and induces on $(e_X,W)^\perp/(e_X,W)=\auHom(W,V)$ the
action of an element $h$ of $\auU(V)$ (the image of $p$ in the Levi quotient of the stabiliser of
$Ee_X$). The archimedean form of (W2) gives $Z_X\omega(1,p)=\chi(p)\omega(1,h)Z_X$ with
$|\chi(p)|=|\det p|_{(e_X,W)}|^{1/2}=|a|^{N/2}$ ($p$ acts on $\auHom(W,Ee_X)$ by the scalar $a$). The
unitary part of $\chi(p)$, together with the change of the archimedean splitting characters from
$\gamma_{W'}$ to $\gamma_W$ on the target, is the character $\prod_\tau(a_\tau/|a_\tau|)^{-2\mu_\tau}=
\tilde\nu_\infty(a)^{-1}$ of the weight (the determinant of $p$ is $(a/\bar a)\det h$).
\end{proof}

\begin{lemma}[Smooth Witt]\label{pct:smoothwitt}
There are a neighbourhood $\mathcal U$ of $e_X$ in $V'(E\otimes\R)$ and a smooth map
$k\colon\mathcal U\to\auU(V')(E^+\otimes\R)$ with $k(x)e_X=x$ for every isotropic $x\in\mathcal U$.
\end{lemma}
\begin{proof}
Place by place. The set $\mathcal C$ of non-zero isotropic vectors of $V'_\tau\cong\C^{p,q}$ is a
smooth submanifold, and $\auU(p,q)$ acts transitively on it (Witt's theorem over $\C/\R$). The orbit
map $\pi\colon k\mapsto ke_X$ is a smooth surjective equivariant map of homogeneous spaces, hence a
submersion, and has a smooth local section $s$ near $e_X$ with $s(e_X)=1$. Let
$P\colon\mathcal U\to\mathcal C$ be a smooth retraction of a neighbourhood of $e_X$ onto $\mathcal C$
(a tubular neighbourhood), and put $k=s\circ P$; for isotropic $x$, $P(x)=x$.
\end{proof}

\begin{proposition}[The image of the cone integral]\label{pct:coneimage}
Let $N$ be even. There is a function $B>0$ on $\R$, real-analytic,
such that every Fock vector $v$ of $\auHom(W,V)\otimes\R$ with finite moments has a Fock vector $v'$
of $\auHom(W',V')\otimes\R$ with finite moments satisfying, for all $x>0$, all tame $\nu$ of weight
$\mu$, all $z\in\pcMat_{b\times N}(O)$ and all $g$,
\[
\Theta^{\mathrm L}_{W,V}(\pctcone_xv',\Phi^{\mathrm{row}}_z)(g)=B(x)\,\Theta^{\mathrm L}_{W,V}(v,
\Phi^{\mathrm{row}}_z)(g).
\]
In particular, for fixed $x>0$ there is $v'$ with $\Theta^{\mathrm L}_{W,V}(\pctcone_xv',\cdot)=
\Theta^{\mathrm L}_{W,V}(v,\cdot)$ on all these data (replace $v$ by $B(x)^{-1}v$).
\end{proposition}
\begin{proof}
Let $\mathcal U$, $k$ be as in Lemma~\ref{pct:smoothwitt}, and $\beta\ge0$ smooth with compact support in
$\mathcal U$, $\beta(e_X)=1$, vanishing unless $\frac12|e_{X,\tau}|^2\le|x_\tau|^2\le2|e_{X,\tau}|^2$ for
every $\tau$. Lemma~\ref{pct:realmixed}(b) needs $k$ on all of $V'(E\otimes\R)$; we extend it. Since
$k(e_X)=1$, after shrinking $\mathcal U$ we may assume $k(\mathcal U)\subset\exp(\mathcal B)$, where $\mathcal B$
is a ball about $0$ in the Lie algebra $\mathfrak g$ of $\auU(V')(E^+\otimes\R)$ on which $\exp$ is a
diffeomorphism onto an open set, so $k=\exp\circ X$ with $X\colon\mathcal U\to\mathcal B$ smooth. Take
$\beta$ as above with support in this smaller $\mathcal U$, a smooth $\phi$ with compact support in $\mathcal U$
and $\phi=1$ on $\mathrm{supp}\,\beta$, and replace $k$ by $\exp(\phi X)$ (extended by $1$ outside
$\mathcal U$): it is smooth on $V'(E\otimes\R)$ and equal to $k$ on $\mathrm{supp}\,\beta$, which is all
that enters below. By Lemma~\ref{pct:realmixed}(a) write $v=Z_Xv_B$, and by (b) choose $v'$ with
$\pctsl_{A,\infty}(y)v'=\beta(y)\omega_\infty(1,k(y))v_B$. Let $\kappa\in\pctK'_\infty$, $a\in(E\otimes\R)^\times$,
$y=\kappa^{-1}ae_X$ (isotropic). By (W1), (W2) at $\infty$ (the slice is equivariant for the action
of $\auU(V')$ on the target), $\pctsl_{A,\infty}(ae_X)\omega'_\infty(1,\kappa)v'=\omega_\infty(1,\kappa)
\pctsl_{A,\infty}(y)v'=\beta(y)\,\omega_\infty(1,\kappa k(y))v_B$. If $\beta(y)\ne0$, then
$\kappa k(y)e_X=\kappa y=ae_X$, and Lemma~\ref{pct:zeroterm-equiv} gives $h=h(\kappa,a)$ with
\[
|a|^x\tilde\nu_\infty(a)\pctL^\infty_{\mathrm t}(a)\omega'_\infty(1,\kappa)v'=\varrho(\kappa,a)|a|^x\,
\omega_\infty(1,h)v,\qquad\varrho(\kappa,a)=\beta(\kappa^{-1}ae_X)|a|^{N/2}\ge0 .
\]
The functional $w\mapsto\Theta^{\mathrm L}_{W,V}(w,\Phi^{\mathrm{row}}_z)(g)$ is continuous for some
moment norm (Proposition~\ref{pcx:decay}(a)) and invariant under $w\mapsto\omega_\infty(1,h)w$, $h\in
\auU(V)(E^+\otimes\R)$ (right translation of the variable of integration over $[\auU(V)]$). Hence
$\Theta^{\mathrm L}_{W,V}(\pctcone_xv',\Phi^{\mathrm{row}}_z)(g)=B(x)\Theta^{\mathrm L}_{W,V}(v,
\Phi^{\mathrm{row}}_z)(g)$ with $B(x)=\int\varrho(\kappa,a)|a|^x\,d\kappa\,da$. On the support of
$\varrho$, $|a_\tau|$ lies in $[2^{-1/2},2^{1/2}]$ (as $\kappa$ preserves lengths), so $\varrho$ is continuous
with compact support in $\pctK'_\infty\times(E\otimes\R)^\times$; thus $B$ is real-analytic
(differentiate under the integral) and $B(x)>0$ since $\varrho\ge0$ and $\varrho(1,1)=\beta(e_X)=1>0$.
\end{proof}

\subsection{Adding a plane: proof of Proposition~\ref{pcx:P-plane}}\label{pct:ss-plane}

\begin{proof}[Proof of Proposition~\ref{pcx:P-plane}]
Let $W'=W\oplus\pcxH$ be adapted, $\operatorname{rk}W=N=2d$, and $v$ of $(W,W)$ with finite moments.
Apply \S\ref{pct:ss-setup} with $V'=W'$, $V=W$, $X=A$, $b=0$: then $r=N$ and the integer of
Theorem~\ref{pct:plane} is $r-N/2=d\ge0$. Proposition~\ref{pct:coneimage} with $x=d+1$ gives $v'$
with finite moments and $\Theta^{\mathrm L}_{W,W}(\pctcone_{d+1}v',\Phi^{\mathrm{row}}_\emptyset)=
\Theta^{\mathrm L}_{W,W}(v,\Phi^{\mathrm{row}}_\emptyset)$ for all $\nu$. Let $F=\Theta^{\mathrm L}_{W',W'}
(v',\Psi^{\mathrm{col}}_1)$; $\Psi^{\mathrm{col}}_1=\Phi'_\emptyset$. By Lemma~\ref{pcx:const} (proved in
\S\ref{pct:ss-const}), $F$ is constant, so $F(1)=F_A(m_A(1,1))$, and by \eqref{pct:eq-plane}
\[
F(1)=\kappa_\nu\,\Theta^{\mathrm L}_{W,W}(\pctcone_{d+1}v',\Psi^{\mathrm{col}}_1)(1)=\kappa_\nu\,
\Theta^{\mathrm L}_{W,W}(v,\Psi^{\mathrm{col}}_1)(1),\qquad|\kappa_\nu|=C\,L^\sharp(d+1).
\]
Since $L_*(\pcm,\nu,d+1)=L^\sharp(d+1)^2L_*(\pcm,\nu,d)$,
$\pcxstr_{W'}(\nu;v')=|F(1)|^2/L_*(\pcm,\nu,d+1)=C^2\,|\Theta^{\mathrm L}_{W,W}(v,\Psi^{\mathrm{col}}_1)(1)|^2/\allowbreak 
L_*(\pcm,\nu,d)=C^2\,\pcxstr_W(\nu;v)$, with $C^2>0$ independent of $\nu$.
\end{proof}

\subsection{Removing a plane: proof of Proposition~\ref{pcx:P-descent}}\label{pct:ss-descent}

The precise statement proved (the parenthetical remark of Proposition~\ref{pcx:P-descent} summarises
it) is:

\begin{proposition}\label{pct:descent}
Let $W'=W\oplus\pcxH_A$ be an adapted splitting, $N=\operatorname{rk}W$ even, $T'$ a space of rank
$b$, $T_0$ of rank $r$ and $T$ of rank $r+2$, with $N=b+r$ and $q_T=q_{T_0}+1$, and $d=r-N/2$. If
$\mathrm{RC}(W',T,T';x,x';\ell)$ holds, then $\mathrm{RC}(W,T_0,T';v,y;\ell)$ holds with $y=Z_Ax'$
the zero term of $x'$ along the line $Ee_A$ (defined below) and $v=\pctcone_{d+1}(x)$ if $d\ge0$,
$v=\pctcone^\vee_{-d}(x)$ if $d<0$, the cone integrals being formed with an adapted splitting
$V'=T_0\oplus\pcxH_X$ having the signs of $T$. Only values $L^S(z,\nu)$ with $z\ge1$ enter; no
functional equation is used.
\end{proposition}

For a space $U$ of rank $b$ let $I^U_A\subset\auHom(W',U)$ be the maps vanishing on $W\oplus Ee_A$
(isotropic, $\cong U$, with $I^{U\perp}_A/I^U_A=\auHom(W,U)$), $Z_A$ the archimedean zero term along
it, and recall $\Theta^{\mathrm S}_{X,U}(y,\Psi)=\int_{[\auU(U)]}\theta^{X,U}_{\upsilon^b,\nu\upsilon^{n/2}}
(\cdot,h;[y,\Psi^{\mathsf t}]_{\bar S})\,dh$ for $X$ of rank $n$ (Definition~\ref{pcx:def-large}).

\begin{lemma}[The constant term on the small side]\label{pct:small}
There is $\kappa_A\ne0$, depending only on $A$ and $U$, such that for all $y'$ of $(W',U)$ with
finite moments, $z\in\pcMat_{b\times N}(O)$ and $g\in\auU(W)(\auA)$,
\[
\bigl(\Theta^{\mathrm S}_{W',U}(y',\Psi^{\mathrm{col}}_{z\,\mathrm{pr}})\bigr)_A(m_A(1,g))=\kappa_A\,
\Theta^{\mathrm S}_{W,U}(Z_Ay',\Psi^{\mathrm{col}}_z)(g).
\]
\end{lemma}
\begin{proof}
As in Lemma~\ref{pct:expand} (with Proposition~\ref{pcx:decay}(b) for the bound, the finite datum
being supported on injective matrices), the constant term is $\int_{[\auU(U)]}\sum_{u\in U(E)}
T_u\,dh$, with $T_u$ the mean over $\pctN_A$ of the term of $u$ in the expansion (W3) along $I^U_A$,
and $T_u=0$ for anisotropic $u$.

\emph{$u=0$.} As $m_A(1,g)$ fixes $f_A$ and stabilises $I^U_A$, (W2) gives
$T_0=\theta_{W,U}(\omega(g,h)\psZ_{I^U_A}f')$ for $f'=[y',\Psi^{\mathrm{col}\,\mathsf t}_{z\,\mathrm{pr}}]_{\bar S}$,
the characters $(\upsilon^b,\nu\upsilon^{N/2+1})$ of $(W',U)$ becoming those $(\upsilon^b,\nu\upsilon^{N/2})$
of $(W,U)$ because $\chi_{I^U_A}(m_A(1,g),h)=\upsilon(\det h)^{-1}$. At $\bar w\in\bar S$, in the model (W5),
$\Psi^{\mathrm{col}}_{z\,\mathrm{pr}}(Y)=\chi^{-1}(\det(z\,\mathrm{pr}\,Y))$ (for integral $Y$) depends only on
the rows of $Y$ along $W$; the zero term restricts to the matrices vanishing in the row dual to $I^U_A$
and integrates the row of $I^U_A$ over $O$, which gives $\Psi^{\mathrm{col}}_z$ times $1$. At the other
finite places the zero term of the vacuum is a non-zero multiple of the vacuum (W4), and at $\infty$
it is $Z_A$. So $\psZ_{I^U_A}f'=\kappa_A[Z_Ay',\Psi^{\mathrm{col}\,\mathsf t}_z]_{\bar S}$.

\emph{Isotropic $u\ne0$.} As in Lemma~\ref{pct:unfold}(i), $T_u$ is $\theta$ of the zero term along
$\auHom(W,Eu)$ of the slice at $u$ of $\omega(m_A(1,g),h)f'$; the slice is $\omega(g,h)$ applied to a
datum whose component at $\bar S$ is supported, like $\Psi^{\mathrm{col}\,\mathsf t}_z$, on matrices of full
rank $b$. The functional $\theta\circ\psZ_{\auHom(W,Eu)}\circ\omega(g,h)$ is invariant under the
translations by $\auHom(W,E\,h^{-1}u)(\auA)$, an isotropic subspace whose elements pair
non-trivially at $\bar S$ with every full-rank matrix (as $h^{-1}u\ne0$), so it kills such data by
Lemma~\ref{pcx:meanzero}. Hence $T_u=0$, and the lemma follows.
\end{proof}

\begin{proof}[Proof of Proposition~\ref{pct:descent} (hence of Proposition~\ref{pcx:P-descent})]
By Lemma~\ref{pcx:adapted} there is an adapted splitting $V'=T_0\oplus\pcxH_X$ with the signs of
$T$; by Lemma~\ref{pcx:move} (with $W_1=W=W'$, where $k=1$ may be taken) and the definition of RC,
$\mathrm{RC}(W',V',T';x,x';\ell)$ holds with $c$ replaced by $u^{-1}c$, $|u|=1$. Let $C_0$ and $c$
be as in Definition~\ref{pcx:def-RC} for this identity:
$C_0|c|^2=L_*(\pcm,\nu,d+1)\ell(\nu)$ (as $(r+2)-(N+2)/2=d+1$). For $z\in\pcMat_{b\times N}(O)$
apply RC(b) to the integral matrix $z\,\mathrm{pr}$ ($\Phi^{\mathrm{row}}_{z\,\mathrm{pr}}=\Phi'_z$) and take
constant terms along $A$ at $m_A(1,g)$. By Theorem~\ref{pct:plane} with $u=1$ and Lemma~\ref{pct:small},
\[
\kappa_\nu\,\Theta^{\mathrm L}_{W,T_0}(v,\Phi^{\mathrm{row}}_z)(g)=c\,\kappa_A\,\Theta^{\mathrm S}_{W,T'}
(Z_Ax',\Psi^{\mathrm{col}}_z)(g),
\]
with $v$ as in the statement and $|\kappa_\nu|=C\,L^\sharp(x_d)$, $x_d=d+1$ or $-d$; $\kappa_\nu\ne0$ by
Lemma~\ref{pcx:Lpos}. Put $c_0=c\kappa_A/\kappa_\nu$ and $C'=C_0C^2/|\kappa_A|^2$. Then
$C'|c_0|^2=C_0|c|^2C^2/|\kappa_\nu|^2=L_*(\pcm,\nu,d+1)\ell(\nu)/L^\sharp(x_d)^2=L_*(\pcm,\nu,d)\ell(\nu)$, by
$L_*(e+1)=L^\sharp(e+1)^2L_*(e)$ ($e=d\ge0$) and its form $L_*(-z)=L^\sharp(z+1)^2L_*(-z-1)$ for
$d=-(z+1)<0$. The vectors $v$ and $y=Z_Ax'$ have finite moments (\S\ref{pct:ss-statement} and
Lemma~\ref{pct:slicedecay}). This is $\mathrm{RC}(W,T_0,T';v,y;\ell)$.
\end{proof}

\begin{remark}[Orientation of the zero term]\label{pct:rem-orient}
Proposition~\ref{pct:descent} holds for \emph{every} adapted splitting $A$ of a plane off $W'$, the
zero term $Z_A$ being taken along $I^U_A$ (maps vanishing on $W\oplus Ee_A$), i.e.\ in the orientation of
the isotropic frame $(e_A,f_A)$ of $A$; the space $T_0$ is arbitrary of rank $\operatorname{rk}T-2$ with
$q_{T_0}=q_T-1$, and the exponent of RC(a) passes from $\operatorname{rk}T-(N+2)/2=d+1$ to $d$, with
$\ell$ kept and the constant changed. Applied to the splitting $D^{\mathrm s}$ with the isotropic frames
exchanged (as in Remark~\ref{pcu:rem-descent}) it gives $y=\mathrm{sl}_{D^{\mathrm s}}(0)x'$:
Proposition~\ref{pct:descent} is stated for an arbitrary adapted splitting and is applied with that
splitting, so no conversion is needed.
\end{remark}

\subsection{Non-vanishing: proof of Proposition~\ref{pcx:P-nonvan}}\label{pct:ss-nonvan}

For spaces $W$ (even rank $N$), $V$ (rank $r$), $b=N-r$, say that $\mathrm{Van}_\nu(W,V)$ holds if
$\Theta^{\mathrm L}_{W,V}(v,\Phi^{\mathrm{row}}_z)(g)=0$ for all $v$ with finite moments, all
$z\in\pcMat_{b\times N}(O)$ and all $g$.

\begin{proposition}[Vanishing descends]\label{pct:vandesc}
Let $k\in\Z_{\ge0}$ with $k\le\mu_\tau$ for all $\tau$, let $W'=W\oplus
\pcxH_A$, $V'=V\oplus\pcxH_X$ be adapted, $N$ even and $d=r-N/2\ge-k$. Then
$\mathrm{Van}_\nu(W',V')\Rightarrow\mathrm{Van}_\nu(W,V)$. The same holds for $W'$, $V'$ of the
ranks and signs of $W\oplus\pcxH^j$, $V\oplus\pcxH^j$, $j\ge1$.
\end{proposition}
\begin{proof}
Let $v,z,g$ be given and $I=\Theta^{\mathrm L}_{W,V}(v,\Phi^{\mathrm{row}}_z)(g)$. Take $B$ and $v'$ of
Proposition~\ref{pct:coneimage}, and $\pctZ=\pctZ_{v',z,g}$ of Theorem~\ref{pct:plane} (with $u=1$).
By (c), for $x\ge1$, $\pctZ(x)=\varepsilon_{\mathrm t}P(x)L(x,\tilde\nu)B(x)I$. The function
$L(s,\tilde\nu)$ is entire, and $L(j,\tilde\nu)\ne0$ for every integer $j>-k$: for $j\ge1$ by
Lemma~\ref{pcx:Lpos}; for $-k<j\le0$ by the functional equation $\Lambda(s,\tilde\nu)=
\epsilon\Lambda(1-s,\tilde\nu^{-1})$ of the completed $L$-function $\Lambda(s,\tilde\nu)=\prod_\tau
\Gamma_\C(s+\mu_\tau)L(s,\tilde\nu)$ (Tate's thesis, Main Theorem~4.4.1; the archimedean factor of
$a\mapsto(a/\bar a)^{\mu_\tau}$ is $\Gamma_\C(s+\mu_\tau)$), since $\Gamma_\C(j+\mu_\tau)$ and
$\Gamma_\C(1-j+\mu_\tau)$ are finite and non-zero and $L(1-j,\tilde\nu^{-1})\ne0$ ($1-j\ge1$). Both
$\pctZ$ and $\varepsilon_{\mathrm t}PL(\cdot,\tilde\nu)BI$ are real-analytic on $\R$ and agree on
$[1,\infty)$, hence everywhere. At $x=d+1>-k$, Theorem~\ref{pct:plane}(b) and the hypothesis give
$\pctZ(d+1)=(\Theta^{\mathrm L}_{W',V'}(v',\Phi'_z))_A(m_A(1,g))=0$, while
$\varepsilon_{\mathrm t}P(d+1)L(d+1,\tilde\nu)B(d+1)\ne0$. So $I=0$. For $j$ planes, induct, replacing
$W'$, $V'$ by spaces with the same signs (Lemma~\ref{pcx:move}, $\mathrm{Van}$ depending only on
signs) and using that $d$ only increases going up.
\end{proof}

The second input is a special case (period $0$) of the large-triple identity of
Proposition~\ref{pcx:P-top}:

\begin{proposition}[Vanishing at the top]\label{pct:vantop}
Let $J$ be positive definite of rank $2k$, $Y=\pcxH\oplus J$, $X=\pcxH^{2k}\oplus Y$ (adapted),
$T=J^-\oplus Y$. If $\Theta^{\mathrm L}_{Y,Y}(w,\Psi^{\mathrm{col}}_1)\equiv0$ for all $w$ with finite
moments, then $\mathrm{Van}_\nu(X,T)$.
\end{proposition}
In the notation of \S\S\ref{pcu:ss-top-setting}--\ref{pcu:ss-std-data} take $W_0=Y$ (rank $n=2j$,
$j=k+1$, one negative sign, the plane split off adaptedly), $U=J^-$ ($b=2k$, signs $t'=-1$), so that
$X=\pcxH^{b}\oplus W_0$ and $T=U\oplus W_0$ are those of Proposition~\ref{pct:vantop}. The
proof of Proposition~\ref{pcx:P-top} in \S\ref{pcu:sec} is written for the particular vector $\Phi$ of
\S\ref{pcu:ss-period}, whose slices are supported on the real class of $U$ (Lemma~\ref{pcu:support});
for an arbitrary $v$ that support property is replaced, in case (b) of Lemma~\ref{pcu:coeff}, by the
following vanishing, which uses the signs.

\begin{lemma}[Coefficients at non-negative indices]\label{pct:negcoeff}
In the situation just described, let $v$ be any Fock vector of $\auHom(X,T)_\R$ with finite moments
and $\beta$ an invertible hermitian $b\times b$ matrix over $E$ that is not negative definite at some
infinite place. Then $\pcuF_\beta\Theta^{\mathrm L}_{X,T}(v,\Phi^{\mathrm{row}}_{z_0})\equiv0$ on
$P_D(\auA)$.
\end{lemma}
\noindent Lemma~\ref{pct:negcoeff} is Theorem~\ref{pcn:negcoeff}, proved in \S\ref{pcn:sec}. It uses
only that the signs of $U$ are $-1$, that $W_0$ has a positive sign, and that
$\mu_\tau+p_0-\frac n2\ge1$ ($p_0=2k+1$ positive signs of $W_0$, $n=2k+2$), i.e.\
$k+1<2k+1+\mu_\tau$, which holds. The proof does not unfold: it regards the coefficient as a continuous
functional of the archimedean datum at a place $\tau$ where $\beta$ is not negative definite, invariant under
$\auU(T_\tau)$ and $\beta$-equivariant under the centre of the unipotent radical; such a functional is carried by
the orbit of index $\beta$ and descends to a functional on the fibre invariant under the stabiliser of a point, and a
compact subgroup of that stabiliser admits no non-zero invariant vector, by computing the action of the rotations of
a negative line orthogonal to the point on polynomial Fock vectors (Lemmas~\ref{pcn:rot}--\ref{pcn:nofix},
\ref{pcn:local}).

\begin{proof}[Proof of Proposition~\ref{pct:vantop}]
Fix $\nu$ and $v$ with finite moments, and let $\mathcal D=\mathcal T$ be the linear map of
Lemma~\ref{pcx:linfactor} attached to $\Phi'\mapsto\Theta^{\mathrm L}_{X,T}(v,\Phi')$ (this is the
difference of \S\ref{pcu:ss-std-data} with $C=0$). It satisfies (E1)--(E4) of \S\ref{pcu:ss-std-data}
(their proofs do not use the particular vector), with
$\mathcal D(\Psi^{\mathrm{col}}_{z'})=\Theta^{\mathrm L}_{X,T}(v,\Phi^{\mathrm{row}}_{z'})$. By
Proposition~\ref{pcu:selfdual} it suffices to show $\pcuF_\beta\Theta^{\mathrm L}_{X,T}(v,
\Phi^{\mathrm{row}}_{z_0})(p)=0$ at every test pair $(\beta,q)$ for $U$.

Type (b): $\beta$ is not of the form $a^*\mathrm{diag}(t')a=-a^*a$ at some infinite place, i.e.\ not
negative definite there; Lemma~\ref{pct:negcoeff}.

Type (a): $\beta=J_U$. By (U1) and Witt's theorem the coefficient is
$\int_{\auU(W_0)(E)\backslash\auU(T)(\auA)}\mathrm{term}_{x_0}(p,h)\,dh$ (the stabiliser of $x_0$ is
$\auU(W_0)$), absolutely convergent by Lemma~\ref{pcu:integrable}(a) (with $\phi=v$, which has finite
moments of all orders). By the quotient integral formula it is
$\int_{\auU(W_0)(\auA)\backslash\auU(T)(\auA)}\int_{[\auU(W_0)]}\mathrm{term}_{x_0}(p,lh)\,dl\,dh$, and we show
the inner integral is $0$ for every $h$. By (U2), $\mathrm{term}_{x_0}(p,lh)$ is a scalar (independent
of $l$, since $l$ fixes $x_0$ and $|\det|=1$) times $\vartheta_{F^{(y)},x'c^*J_{W_0},(g,lh)}(\phi)$ with
$\phi=\mathrm{sl}_D(h^{-1}Y)v$, a vector of $\auHom(W_0,T)_\R=\auHom(W_0,U)_\R\oplus\auHom(W_0,W_0)_\R$
with finite moments, and $F^{(y)}=f\boxtimes F_1$, $F_1=[\,\cdot\,,\Psi^{\mathrm{col}}_1]_S$, the slice
of $\Phi^{\mathrm{row}}_{z_0}$ being $\chi^{-1}\circ\det$ in the $W_0$-block (as in the proof of
Lemma~\ref{pcu:period}). Expand $\phi=\int\phi(y)\psrho_\infty(y)e\,dy$ (Lemma~\ref{pcx:superpos}); a
coherent state is $\psrho(y_1)e\boxtimes\psrho(y_2)e$ along the orthogonal decomposition, so by
Lemma~\ref{pcu:product} the integrand is
$\int\phi(y)\,\vartheta_{f,\dots}(\psrho(y_1)e)\,\theta^{W_0,W_0}_{\gamma,\gamma}(g,l;\psrho(y_2)e\otimes
F_1)\,dy$. The bound of Proposition~\ref{pcx:decay}(a) for the second factor
($\le C(1+|y_2|)^M\mathrm{ht}(l)^{-N_0}\mathrm{ht}(g)^B$) and $\|\phi\|_M<\infty$ allow Fubini in
$(y,l)$; the integral over $l$ of the second factor is
$\Theta^{\mathrm L}_{W_0,W_0}(\psrho(y_2)e,\Psi^{\mathrm{col}}_1)(g)=0$ by hypothesis (coherent states
have finite moments). So the coefficient vanishes, and Proposition~\ref{pcu:selfdual} gives
$\Theta^{\mathrm L}_{X,T}(v,\Phi^{\mathrm{row}}_{z'})=0$ for all integral $z'$, i.e.\
$\mathrm{Van}_\nu(X,T)$.
\end{proof}

\begin{proof}[Proof of Proposition~\ref{pcx:P-nonvan}]
Suppose $\pcxstr_Y(\nu;z)=0$ for all $z$. By Lemma~\ref{pcx:const}, $\Theta^{\mathrm L}_{Y,Y}(z,
\Psi^{\mathrm{col}}_1)\equiv0$, so $\mathrm{Van}_\nu(X,T)$ by Proposition~\ref{pct:vantop}. Here $X$ has
rank $6k+2$ and $2k+1$ negative signs, $T$ rank $4k+2$ and $2k+1$ negative signs; they have the ranks
and signs of $J\oplus\pcxH^{2k+1}$ and $0\oplus\pcxH^{2k+1}$. Proposition~\ref{pct:vandesc} with
$j=2k+1$ (the integer $d$ is $\ge-k$ at every level, $-k$ at the bottom) gives
$\mathrm{Van}_\nu(J,0)$. But for $(J,0)$, $b=2k$, $z=1_{2k}$, $v$ the vacuum and $g=1$, the kernel
on $\auHom(J,0)=0$ is $\Phi^{\mathrm{row}}_{1}(\emptyset)=\chi^{-1}(\det1)=1$, so the integral is the
volume of $[\auU(0)]$, non-zero. Contradiction.
\end{proof}

\subsection{The inputs of Theorem~\ref{pcx:X}}\label{pct:ss-status}
This section proves Theorem~\ref{pct:plane}; Lemmas~\ref{pct:poisson}--\ref{pct:poisson-a},
\ref{pct:qref}, \ref{pct:SA}, \ref{pct:realmixed}, \ref{pct:zeroterm-equiv}, \ref{pct:smoothwitt},
\ref{pct:small}; Lemma~\ref{pcx:const}; Propositions~\ref{pct:coneimage}, \ref{pcx:P-plane},
\ref{pct:descent} (= \ref{pcx:P-descent}), \ref{pct:vandesc}, \ref{pct:vantop}, hence
Proposition~\ref{pcx:P-nonvan}. Lemma~\ref{pct:negcoeff} is proved in \S\ref{pcn:sec}
(Theorem~\ref{pcn:negcoeff}), without using (U1); type (a) of Proposition~\ref{pct:vantop} uses Lemma~\ref{pcu:integrable}.

With \S\ref{pcu:ss-final} and \S\ref{pcn:ss-dep}, Theorem~\ref{pcx:X} therefore rests only on the
following published results (with their hypotheses verified where used): the facts (W1)--(W6) on the Weil
representation (Weil, Kudla, Harris--Kudla--Sweet; Lemma~\ref{ps:weil}, Proposition~\ref{pf:flip});
Stone--von Neumann and Weil's intertwining operator (M2); reduction theory (R1) (Borel) and moderate
growth of theta functions (R2) (Lemma~\ref{pcx:thetagrowth}, from Godement's form of Borel's reduction theorem); Landherr's theorem, the Hasse norm theorem and Jacobowitz's
classification of unimodular hermitian lattices (Lemmas~\ref{pcx:landherr}, \ref{pcu:selfdual}); Witt's
theorem and generation of unitary groups by quasi-reflections (Hahn--O'Meara); strong approximation
for the additive group (Cassels--Fr\"ohlich); weak approximation for unitary groups
(Platonov--Rapinchuk); the Iwasawa decomposition and the integration formula for parabolic
subgroups (Moeglin--Waldspurger); Hecke's theorems for the non-trivial unitary character
$\tilde\nu$ --- $L(s,\tilde\nu)$ entire, non-vanishing on $\Rel s=1$, and Tate's functional equation
(Tate's thesis, Main Theorem~4.4.1); classical Gauss sums (Lemma~\ref{pct:gauss}, proved); the existence
of smooth local sections of submersions (Lemma~\ref{pct:smoothwitt}); Witt's theorem over $\C$, the
Fock-model action of the maximal compact subgroup (Folland~\cite[Prop.~4.39]{FollandPS},
Kashiwara--Vergne~\cite{KashiwaraVergne}), the unitarity of the Bargmann transform (Lemma~\ref{pcu:mixedmodel}), Stone's theorem and the inverse function theorem. The density of
polynomial (finite Hermite) vectors among smooth vectors, the equivalence of the norms $\|\cdot\|_s$ with
$\|(\mathsf N+1)^s\cdot\|$, the bound of moments by these norms, the smoothness of the action on smooth vectors, the
mixed model on smooth vectors, invariant distributions and adapted charts are proved in
\S\S\ref{pcn:ss-smooth}, \ref{pcn:ss-local} (Lemmas~\ref{pcn:smooth}, \ref{pcn:smoothaction}, \ref{pcn:mixed},
\ref{pcn:local}; see also Reed--Simon~\cite[Appendix to \S V.3, Lemma~2 and Thm.~V.13]{ReedSimonI} for Hermite expansions).


%% ===== proofs/C-exchange-neg.tex =====
% proofs/C-exchange-neg.tex
% Proof of Lemma pct:negcoeff of C-exchange-tower.tex: centre coefficients of the theta integral
% vanish at indices that are not negative definite at some real place.
% Labels of this file start with pcn:. Macros: those of C-exchange.tex, C-exchange-top.tex and
% C-exchange-tower.tex.
\providecommand{\pcnK}{\mathsf K}
\providecommand{\pcnC}{\mathsf C}
\providecommand{\pcnS}{\mathsf S}

\section{Centre coefficients at indices that are not negative definite}\label{pcn:sec}

This section proves Lemma~\ref{pct:negcoeff}, the last input of Proposition~\ref{pct:vantop} and hence of
Proposition~\ref{pcx:P-nonvan}. The argument: the coefficient, as a function of the archimedean
component of the datum at a real place $\tau$ where the index is not negative definite, is a continuous linear
functional which is invariant under $\auU(T_\tau)$ (right invariance of the measure on $[\auU(T)]$) and transforms by
the character of index $\beta$ under the centre of the unipotent radical. Such a functional is carried by the orbit
$\{x:x^*Jx=\beta\}$ at $\tau$ and descends to a functional on the fibre invariant under the stabiliser of a point
$x$; a compact subgroup of that stabiliser contains, for every point of the Fock space, a circle of rotations of a
negative line; these act on polynomial Fock vectors by a substitution times a positive power of the rotation
parameter, which forces every invariant functional to vanish. No unfolding of the global integral is needed. We first
fix the statement and the conventions, then do the computation at one real place (\S\S\ref{pcn:ss-fock},
\ref{pcn:ss-local}), and then the adelic argument (\S\ref{pcn:ss-proof}).

\subsection{Statement and conventions}\label{pcn:ss-statement}

We are in the situation of Proposition~\ref{pct:vantop}, in the notation of
\S\S\ref{pcu:ss-top-setting}--\ref{pcu:ss-coeff}: $J$ is positive definite of rank $2k$ ($k\ge0$),
$W_0=Y=\pcxH\oplus J$ has rank $n=2k+2$ and signature $(p_0,q_0)=(2k+1,1)$ at every infinite place,
$U=J^-$ has rank $b=2k$ and signs $t'=-1$, $X=\pcxH^b\oplus W_0$ with the adapted splitting
$D\colon X=(\ell\oplus\ell')\perp W_0$, and $T=U\oplus W_0$, of rank $b+n$ and signature
\[
(p',q')=(p_0,\,b+q_0)=(2k+1,\,2k+1)
\]
at every infinite place. The weight $\mu$ of $\nu$ has $\mu_\tau\ge1$ for all $\tau$. For a space $W$ of
even rank $N$, $\gamma_W=\nu^{-1}\upsilon^{N/2}$; the theta integral is
$\Theta^{\mathrm L}_{X,T}(v,\Phi)(g)=\int_{[\auU(T)]}\theta^{X,T}_{\gamma_X,\gamma_X}(g,h;[v,\Phi]_S)\,dh$.
The centre coefficient $\pcuF_\beta$ is that of \S\ref{pcu:ss-top-setting}: $z$ runs through the skew-hermitian
matrices, $\beta$ is hermitian, and $\psi_\beta(z)=\psi(-\frac12\jmath\operatorname{tr}(z\beta))$ \eqref{pcu:eq-psibeta}; the
coefficient $\pcuF_\beta$ keeps the terms of the $x$ with Gram matrix $x^*J_Tx=\beta$.

\begin{theorem}[$=$ Lemma~\ref{pct:negcoeff}]\label{pcn:negcoeff}
Let $v$ be a Fock vector of $\auHom(X,T)_\R$ with finite moments, and let $\beta$ be an invertible
hermitian $b\times b$ matrix over $E$ such that, for some infinite place $\tau$, $\tau(\beta)$ is not
negative definite. Then $\pcuF_\beta\Theta^{\mathrm L}_{X,T}(v,\Phi^{\mathrm{row}}_{z_0})(p)=0$ for every
$p\in P_D(\auA)$.
\end{theorem}

The only properties of the situation that the proof uses are:
\begin{enumerate}
\item[(H1)] every sign of $U$ is $-1$ at $\tau$, so that $T_\tau$ has exactly $b$ more negative signs
than $W_{0,\tau}$;
\item[(H2)] $W_0$ has a positive sign at $\tau$ ($p_0\ge1$);
\item[(H3)] the integer $N_\tau=\mu_\tau+p_0-\frac n2$ is $\ge1$. Here $N_\tau=\mu_\tau+k\ge1$.
\end{enumerate}
In terms of the weight $\mu^\flat_\tau=\frac n2-\mu_\tau$ of the character $\gamma_{W_0}$ by which
$\auU(T)$ acts on the slices (Lemma~\ref{pcn:slicechar}), (H3) reads $\mu^\flat_\tau<p_0$, i.e.\
$k+1<2k+1+\mu_\tau$. No condition of the form $2q_{W_0}\le n$ is needed (it holds anyway: $q_0=1$).

\begin{lemma}[The character on the slices]\label{pcn:slicechar}
For $x\in\pcMat_{(b+n)\times b}(E)$ let $\mathrm{term}_x(p,h)=\theta_{W_0}((\omega(p,h)F)[x])$ be the
term of (U1), $F=[v,\Phi^{\mathrm{row}}_{z_0}]_S$, $\omega=\omega_{\gamma_X,\gamma_X}$. Let
$\pcnS_x\subset\auU(T)(\auA)$ be the pointwise stabiliser of the columns of $x$. For $s\in\pcnS_x$,
\[
\mathrm{term}_x(p,sh)=\theta_{W_0}\bigl(\omega^\flat(1,s)\,\Xi_h\bigr),\qquad\Xi_h=(\omega(p,h)F)[x],
\]
where $\omega^\flat=\omega_{\gamma_X\upsilon^{-b},\,\gamma_X\upsilon^{-b}}$ is the Weil representation
of the pair $(W_0,T)$ whose $\auU(T)$-character is $\gamma_X\upsilon^{-b}=\gamma_{W_0}$.
\end{lemma}
\begin{proof}
$\omega(p,1)$ and $\omega(1,s)$ commute, so $(\omega(p,sh)F)[x]=(\omega(1,s)\omega(p,h)F)[x]$. The Levi law
of Lemma~\ref{pcu:laws} with $q=1$, $g=1$ gives $(\omega(1,s)f)[x]=\omega^\flat(1,s)(f[s^{-1}x])$, and
$s^{-1}x=x$. The character: in the Schr\"odinger model of the variable $\auHom(\ell,T)\cong T^b$, the
group $\auU(T)$ acts linearly with determinant $\det(h)^b$, and the factor $\upsilon(\det h)^{b}$ of
(W2) for this action is absorbed by the \emph{target} character, which becomes $\gamma_X\upsilon^{-b}$.
(This is the target-character shift of Lemma~\ref{pcu:laws}. Since $\Theta^{\mathrm L}$ uses the
same character $\gamma_X$ on both groups, the scalar $\chi_V$ of Lemma~\ref{pcu:laws} is the same with
either character.) Finally
$\gamma_X\upsilon^{-b}=\nu^{-1}\upsilon^{(2b+n)/2-b}=\nu^{-1}\upsilon^{n/2}=\gamma_{W_0}$.
\end{proof}

\emph{Convention on $\upsilon$.} All characters above are evaluated on determinants, which lie in the
norm-one ideles $\auU(1)(\auA)$; $\upsilon$ there means the restriction $\upsilon|_{\auU(1)}$, as in
$\upsilon_k=(\upsilon|_{\auU(1)})^{\lfloor k/2\rfloor}$ of \S\ref{ps:ss-weil} (for even rank $N$,
$\upsilon_N=\upsilon^{N/2}$ on $\auU(1)$, which is the $\upsilon^{N/2}$ in $\gamma_W$). The identity
$\upsilon(x/\bar x)=\upsilon(x)^2$ ($x\in\auA_E^\times$), used in \S\ref{pct:sec} to pass between a
character of $\auU(1)$ and one of $\auA_E^\times$, is exact: $\upsilon(\bar x)=\upsilon(x\bar x)
\upsilon(x)^{-1}$, and $\upsilon(x\bar x)=\omega_{E/E^+}(\mathrm N x)=1$ because the quadratic character
of $E/E^+$ is trivial on norms. So no factor is affected by this normalisation. At an infinite place
$\upsilon_\tau|_{\auU(1)}(z)=z$ (Theorem~\ref{pco:V}) and $\nu_\tau|_{\auU(1)}(z)=z^{\mu_\tau}$ (the
base change $\tilde\nu(x)=\nu(x/\bar x)$ has $\tilde\nu_\tau(a)=(a/|a|)^{2\mu_\tau}$), so
$\gamma_{W_0,\tau}(z)=z^{n/2-\mu_\tau}=z^{\mu^\flat_\tau}$.

\subsection{Smooth vectors of Fock spaces}\label{pcn:ss-smooth}

This subsection proves the facts on smooth vectors used below (and in \S\ref{pct:sec}); everything is
elementary once the Fock model is written in monomials. Let $W_\infty$ be a real symplectic space of
dimension $2D$ with a positive complex structure, realised as in the paragraph preceding
Lemma~\ref{pz:growth} (which applies to every pair of spaces, see the first paragraph of the proof of
Lemma~\ref{pcx:thetagrowth}): $\pzF$ is the Fock space of entire functions on $\C^D$ with
$\|F\|^2=\boldsymbol\pi^{-D}\int|F(z)|^2e^{-|z|^2}dz$, $\psrho=\psrho_\infty$ is given by \eqref{pk:eq-fock},
$\mathsf B=(\mathsf e_k,i\mathsf e_k)_{k\le D}$ is the real orthonormal basis used there, $z_k=z\cdot\mathsf e_k$,
and for $w\in W_\infty$, $\partial_w$ is the generator of the strongly continuous unitary one-parameter group
$s\mapsto\psrho(sw)$ (a closed skew-adjoint operator, by Stone's theorem). On a polynomial $P$, by
\eqref{pk:eq-fock} (proof of Lemma~\ref{pz:growth}(b)),
\begin{equation}\label{pcn:eq-dw}
\partial_wP=(z\cdot w)P-\mathsf D_{\bar w}P,\qquad a^+_kP:=z_kP=\tfrac12(\partial_{\mathsf e_k}-i\partial_{i\mathsf e_k})P,
\qquad a^-_kP:=\mathsf D_{\bar{\mathsf e}_k}P=-\tfrac12(\partial_{\mathsf e_k}+i\partial_{i\mathsf e_k})P,
\end{equation}
so on polynomials $\partial_w=\sum_k(c_k(w)a^+_k+c'_k(w)a^-_k)$ with $|c_k(w)|,|c'_k(w)|\le|w|$, real linear in $w$.
The inner product of $\pzF$ is the Fischer inner product (\S\ref{par:ss-fock}): $a^-_k=\partial/\partial z_k$
is adjoint to $a^+_k$ on polynomials, the monomials $z^\alpha$ are orthogonal with $\|z^\alpha\|^2=\alpha!$,
and they form an orthogonal basis of $\pzF$ (the Taylor series of $F\in\pzF$ converges to $F$ in $\pzF$).
Write $v=\sum_\alpha c_\alpha z^\alpha$ and put
\[
p_k(v)^2=\sum_\alpha(1+|\alpha|)^{2k}\alpha!\,|c_\alpha|^2\in[0,\infty],\qquad
\mathcal S(\pzF)=\{v\in\pzF:\ p_k(v)<\infty\text{ for all }k\ge0\},
\]
a Fr\'echet space (an intersection of weighted $\ell^2$ spaces). Note $p_k(v)=\|(\mathsf N+1)^kv\|$ for the
number operator $\mathsf N=\sum_ka^+_ka^-_k$, $\mathsf Nz^\alpha=|\alpha|z^\alpha$. Recall
$\|\Psi\|_s=\max\|\partial_{w_1}\cdots\partial_{w_j}\Psi\|$ over $j\le2s$, $w_i\in\mathsf B$
(before Lemma~\ref{pz:growth}). If $\pzF=\pzF'\hat\otimes\pzF''$ is the Fock space of an orthogonal sum
$W'\oplus W''$ (monomials in the two sets of variables), then $p_k(v'\otimes v'')\le p_k(v')p_k(v'')$, since
$1+|\alpha'|+|\alpha''|\le(1+|\alpha'|)(1+|\alpha''|)$.

\begin{lemma}[Smooth vectors]\label{pcn:smooth}
\textup{(a)} Every $v\in\mathcal S(\pzF)$ lies in the domain of every product $\partial_{w_1}\cdots\partial_{w_j}$
(each factor applied to a vector of the domain of the next), the product acts on the Taylor series
term by term through \eqref{pcn:eq-dw}, $w\mapsto\partial_wv$ is real linear, $\partial_w\mathcal S(\pzF)\subset
\mathcal S(\pzF)$, and
\[
\|\partial_{w_1}\cdots\partial_{w_j}v\|\le(2D)^jj^{j/2}|w_1|\cdots|w_j|\;p_{\lceil j/2\rceil}(v).
\]
\textup{(b)} Let $\mathsf B'$ be any real basis of $W_\infty$ and $v\in\pzF$ a vector in the domain of every
product $\partial_{u_1}\cdots\partial_{u_j}$ with $u_i\in\mathsf B'$, $j\le2s$. Then $p_s(v)\le C(\mathsf B',s)\max
\|\partial_{u_1}\cdots\partial_{u_j}v\|$ (maximum over these products). In particular $\mathcal S(\pzF)$ is the
space $\pzF^{\mathrm{sm}}$ of vectors in the domains of all products of the $\partial_w$, it equals
$\bigcap_s\pzF_s$, the norms $\|\cdot\|_s$ define its topology
($(D+1)^{-s}p_s\le\|\cdot\|_s\le(2D)^{2s}(2s)^{s}p_s$), and the polynomials are dense in it.
\textup{(c)} In the realisation of Lemma~\ref{pcx:superpos}, $(1+|x|)^{2s}|v(x)|\le(D+1)^{2s}\|v\|_s$ for all $x$,
hence $\|v\|_M\le C_{M,s}\|v\|_s$ whenever $2s>M+2D$, with $C_{M,s}=(D+1)^{2s}\int(1+|x|)^{M-2s}dx$.
\textup{(d)} Let $k$ be a unitary operator on $\pzF$ with $k\psrho(w)k^{-1}=\psrho(\sigma w)$ for a real
symplectic map $\sigma$ (for instance $\omega_\infty(g,h)$, by (W1)). Then $k\,\mathcal S(\pzF)=\mathcal S(\pzF)$,
$\partial_wkv=k\,\partial_{\sigma^{-1}w}v$ for $v\in\mathcal S(\pzF)$, and $\|kv\|_s\le(2D\|\sigma^{-1}\|)^{2s}\|v\|_s$.
The Heisenberg operators $\psrho(y)$ preserve $\mathcal S(\pzF)$, with $\|\psrho(y)v\|_s\le(1+c_\rho|y|)^{2s}\|v\|_s$.
\end{lemma}
\begin{proof}
(a) A product of $j$ operators $a^\pm_{k}$ maps $z^\alpha$ to a multiple of $z^{\alpha+\delta}$, where $\delta$
depends only on the product, and $\|a^+_kz^\alpha\|^2=(\alpha_k+1)\alpha!$, $\|a^-_kz^\alpha\|^2=\alpha_k\alpha!$;
so it multiplies the norm of $z^\alpha$ by at most $\prod_{i<j}(|\alpha|+i+1)^{1/2}\le j^{j/2}(1+|\alpha|)^{j/2}$, and
distinct $\alpha$ go to orthogonal vectors. Expanding $\partial_{w_1}\cdots\partial_{w_j}$ on polynomials into
$(2D)^j$ such products with coefficients of modulus $\le|w_1|\cdots|w_j|$, the termwise operator $T$ given by
\eqref{pcn:eq-dw} satisfies $\|Tv\|\le(2D)^jj^{j/2}|w_1|\cdots|w_j|p_{\lceil j/2\rceil}(v)$ and, similarly,
$p_k(Tv)\le C\,p_{k+j}(v)$. For $j=1$: if $v_N$ is the truncation of the Taylor series of $v\in\mathcal S(\pzF)$ to
degrees $\le N$, then $v_N\to v$ and $\partial_wv_N=Tv_N\to Tv$ in $\pzF$; as $\partial_w$ is closed, $v$ is in
its domain and $\partial_wv=Tv\in\mathcal S(\pzF)$. Induction on $j$ gives the rest; linearity in $w$ holds
because $T$ is linear in $w$.

(b) For a polynomial $q$ and $v$ as in (b), skew-adjointness gives $\langle v,\partial_{u_1}\cdots\partial_{u_j}q
\rangle=(-1)^j\langle\partial_{u_j}\cdots\partial_{u_1}v,q\rangle$ (one factor at a time; polynomials lie in every
domain by (a)). Expand $(\mathsf N+1)^s=(1+\sum_ka^+_ka^-_k)^s$ into $(D+1)^s$ products of at most $2s$
operators $a^\pm_k$, each a combination, through \eqref{pcn:eq-dw} and the linearity of $w\mapsto\partial_w$ on
polynomials, of products of at most $2s$ operators $\partial_u$, $u\in\mathsf B'$, with coefficients bounded in
terms of $\mathsf B'$ and $s$ (for $\mathsf B'=\mathsf B$ the coefficients have total modulus $\le1$ for each
product of $a^\pm$'s). Hence $\langle v,(\mathsf N+1)^sq\rangle=\langle\tilde v,q\rangle$ for a vector $\tilde v$
with $\|\tilde v\|\le C(\mathsf B',s)\max\|\partial_{u_1}\cdots\partial_{u_j}v\|$ (for $\mathsf B'=\mathsf B$,
$C=(D+1)^s$). With $q=z^\alpha$: $(1+|\alpha|)^s\alpha!\,c_\alpha=\langle\tilde v,z^\alpha\rangle$, and Parseval
gives $p_s(v)^2=\sum_\alpha|\langle\tilde v,z^\alpha\rangle|^2/\alpha!=\|\tilde v\|^2$. This is the first claim.
With (a) it gives $\pzF^{\mathrm{sm}}=\mathcal S(\pzF)$ and the two inequalities between $\|\cdot\|_s$ and
$p_s$. A $\|\cdot\|_s$-Cauchy sequence of smooth vectors is $p_s$-Cauchy, so its limit has $p_s<\infty$; hence
$\bigcap_s\pzF_s=\mathcal S(\pzF)$. The truncations $v_N$ converge to $v$ in every $p_k$.

(c) The isomorphism with the model of Lemma~\ref{pcx:superpos} intertwines the Heisenberg operators and
carries vacuum to vacuum, so it preserves the operators $\partial_w$ and the norms $\|\cdot\|_s$, and $v(x)=
\langle\psrho(x)e,v\rangle$ corresponds to $\langle\psrho(x)1,\Psi\rangle$, which by \eqref{pk:eq-fock} (as in the
proof of Lemma~\ref{pz:growth}(c)) has absolute value $e^{-|x|^2/2}|\Psi(\bar x)|$. By \eqref{pcn:eq-dw},
$z^\beta\Psi$ is a product of $|\beta|$ operators $a^+_k$ applied to $\Psi$, each a combination of two
$\partial_w$ with coefficients $\frac12$, so $\|z^\beta\Psi\|\le\|\Psi\|_s$ for $|\beta|\le2s$; with $|F(z)|\le
\|F\|e^{|z|^2/2}$ (reproducing kernel) and the multinomial expansion of $(1+\sum_k|z_k|)^{2s}\ge(1+|z|)^{2s}$,
whose coefficients add up to $(D+1)^{2s}$, this gives the pointwise bound (this is the proof of
Lemma~\ref{pz:growth}(b), which works for every $s$), and integrating gives the moment bound.

(d) Conjugating $s\mapsto\psrho(s\sigma^{-1}w)$ by $k$ gives $s\mapsto\psrho(sw)$, so $k$ carries the domain of
$\partial_{\sigma^{-1}w}$ onto that of $\partial_w$ and $\partial_wk=k\,\partial_{\sigma^{-1}w}$ there. For $v\in\mathcal
S(\pzF)$, induction with (a) shows that $kv$ is in the domain of every product of the $\partial_w$, with
$\partial_{w_1}\cdots\partial_{w_j}kv=k\,\partial_{\sigma^{-1}w_1}\cdots\partial_{\sigma^{-1}w_j}v$; by (b), $kv\in\mathcal
S(\pzF)$, and writing $\sigma^{-1}w_i$ in the basis $\mathsf B$ (coordinates of modulus $\le\|\sigma^{-1}\|$) and
using linearity gives the bound, as in the proof of Lemma~\ref{pz:growth}(a). The same applies to $k^{-1}$. The
statement on $\psrho(y)$ is Lemma~\ref{pcx:superpos}(a) (with $\bigcap_s\pzF_s=\mathcal S(\pzF)$).
\end{proof}

\begin{lemma}[Smooth action of a Lie group]\label{pcn:smoothaction}
Let $G$ be a Lie group, $\omega$ a strongly continuous unitary representation of $G$ on $\pzF$ with
$\omega(g)\psrho(w)\omega(g)^{-1}=\psrho(\sigma_gw)$ for a continuous homomorphism $g\mapsto\sigma_g$ into the real
symplectic group, and suppose that for every $X$ in the Lie algebra the generator $A_X$ of $t\mapsto\omega(\exp tX)$
has the polynomials in its domain and acts on them by a polynomial $\mathsf Q_X$ of degree $\le2$ in the
$a^\pm_k$ (for the Weil representation of a unitary group this is Lemma~\ref{par:vac}, which holds by
Theorem~\ref{pco:V}). Then:
\begin{enumerate}
\item[\textup{(i)}] $\mathcal S(\pzF)$ lies in the domain of $A_X$, and $A_X=\mathsf Q_X$ there (term by term), a
continuous operator of $\mathcal S(\pzF)$; the map $G\times\mathcal S(\pzF)\to\mathcal S(\pzF)$, $(g,v)\mapsto
\omega(g)v$, is continuous, and the operators $\omega(g)$, $g$ in a compact set, are equicontinuous.
\item[\textup{(ii)}] If $\psi\colon G\to\mathcal S(\pzF)$ is smooth, so is $g\mapsto\omega(g)\psi(g)$, and for the
left-invariant vector field $\tilde X$ of $X$, $\tilde X(\omega\psi)(g)=\omega(g)\bigl(\mathsf Q_X\psi(g)+(\tilde X\psi)(g)
\bigr)$. Consequently, for every compact $K\subset G$, $k$ and $r$ there are $C,k'$ with
$\sup_K p_k(\tilde X_1\cdots\tilde X_{r'}(\omega\psi))\le C\max\sup_K p_{k'}(\tilde X_{i_1}\cdots\tilde X_{i_{r''}}\psi)$
($r'\le r$, the maximum over $r''\le r'$ and over a fixed basis of the Lie algebra).
\end{enumerate}
\end{lemma}
\begin{proof}
(i) $\mathsf Q_X$ acts termwise continuously on $\mathcal S(\pzF)$ ($p_k(\mathsf Q_Xv)\le Cp_{k+1}(v)$, as in
Lemma~\ref{pcn:smooth}(a)); since $A_X$ is closed, truncation as in Lemma~\ref{pcn:smooth}(a) gives the first claim.
By Lemma~\ref{pcn:smooth}(d), $\partial_{w_1}\cdots\partial_{w_j}\omega(g)v=\omega(g)\sum_Vc_V(g)Vv$, a sum over
products $V$ of $j$ operators $\partial_u$, $u\in\mathsf B$, with coefficients $c_V$ polynomial in the entries of
$\sigma_g^{-1}$, hence continuous in $g$ and bounded on compact sets; with the strong continuity of $\omega$ on
$\pzF$ this gives the continuity and the equicontinuity.

(ii) First, for $u\in\mathcal S(\pzF)$, $t^{-1}(\omega(e^{tX})u-u)\to\mathsf Q_Xu$ in $\mathcal S(\pzF)$. Indeed, for a
product $W$ of $j$ operators $\partial_w$, write as above $W\omega(e^{tX})=\omega(e^{tX})W_t$ on $\mathcal S(\pzF)$,
$W_t=\sum_Vc_V(t)V$ with $c_V$ smooth and $W_0=W$; then
\[
W\,\frac{\omega(e^{tX})u-u}t=\sum_V\frac{c_V(t)-c_V(0)}t\,\omega(e^{tX})Vu+\frac{\omega(e^{tX})Wu-Wu}t
\longrightarrow W'u+\mathsf Q_XWu\quad\text{in }\pzF,
\]
$W'=\sum_Vc'_V(0)V$. On polynomials $P,q$: $\langle W\omega(e^{tX})P,q\rangle=\langle\omega(e^{tX})P,W^\dagger q\rangle$
($W^\dagger$ the reversed product times $(-1)^j$), whose derivative at $0$ is $\langle\mathsf Q_XP,W^\dagger q\rangle=
\langle W\mathsf Q_XP,q\rangle$, while the derivative of $\langle\omega(e^{tX})W_tP,q\rangle$ is $\langle\mathsf Q_XWP+
W'P,q\rangle$. So $W\mathsf Q_X=\mathsf Q_XW+W'$ on polynomials, hence on $\mathcal S(\pzF)$ (both sides are
continuous $\mathcal S(\pzF)\to\pzF$, Lemma~\ref{pcn:smooth}(a), and polynomials are dense). Thus the limit is
$W\mathsf Q_Xu$, for every $W$, which is the claim by Lemma~\ref{pcn:smooth}(b). Now
\[
\frac{\omega(ge^{tX})\psi(ge^{tX})-\omega(g)\psi(g)}t=\omega(g)\Bigl[\omega(e^{tX})\frac{\psi(ge^{tX})-\psi(g)}t+
\frac{\omega(e^{tX})\psi(g)-\psi(g)}t\Bigr];
\]
the second term tends to $\mathsf Q_X\psi(g)$, and the first to $(\tilde X\psi)(g)$ by the equicontinuity and
continuity in (i). So $\tilde X(\omega\psi)=\omega\cdot\psi_1$ with $\psi_1=\mathsf Q_X\psi+\tilde X\psi$ continuous, of the
same form; by induction all iterated derivatives along the left-invariant fields of a basis exist, are given by
the stated formula and are continuous. The left-invariant fields form a global frame, so in local coordinates all
partial derivatives exist and are continuous, which (with the mean value inequality for each seminorm) means that
$\omega\psi$ is smooth. The estimate follows from the formula, the continuity of $\mathsf Q_X$ and (i).
\end{proof}

\subsection{Rotations of a negative line in the Fock model}\label{pcn:ss-fock}

Fix an infinite place $\tau$. Choose orthonormal bases of $T_\tau=T\otimes_{E,\tau}\C$ and
$W_{0,\tau}$, with signs $s'_a$ ($a$ running through the $b+n$ rows) and $\varepsilon_j$ ($j$ running
through the $n$ columns); the first $b$ rows are those of $U$, with $s'_a=-1$ (H1). An element of
$\auHom(W_0,T)_\tau$ is a matrix $X=(X_{aj})\in\pcMat_{(b+n)\times n}(\C)$, the form is
$\operatorname{tr}(X^\dagger X)=\sum_{a,j}s'_a\varepsilon_j|X_{aj}|^2$, and $h\in\auU(T_\tau)$ acts by
$X\mapsto hX$. As in \S\ref{par:ss-fock}, the Fock space is realised on functions of $X$, the
holomorphic coordinates being
\[
Z_{aj}=\overline{X_{aj}}\ \text{ if }s'_a\varepsilon_j=1,\qquad Z_{aj}=X_{aj}\ \text{ if }s'_a\varepsilon_j=-1,
\]
the vacuum $e$ is a Gaussian function of $\sum_{a,j}|X_{aj}|^2$, and the polynomial vectors are the
functions $P(Z)e$, $P$ a polynomial. The Fock space of $\auHom(W_0,T)_\R$ is the tensor product of
those at the infinite places; its polynomial vectors are $P(Z)e$ with $P$ a polynomial in the
variables $Z_{w,aj}$ of all infinite places $w$. A polynomial vector determines its polynomial (the map
$P\mapsto P(Z)e$ is injective, $e$ having no zeros and $Z$ being a real-linear coordinate system).
In this realisation the polynomial vectors $P(Z)e$ are the polynomial vectors (polynomials) of the model of
\S\ref{pcn:ss-smooth} (\S\ref{par:ss-fock}), and the smooth vectors are the space $\mathcal S(\pzF)$ of
Lemma~\ref{pcn:smooth}. Write $\pcnK_\tau=\{c\in\auU(T_\tau):
c\,\mathrm{diag}(s')=\mathrm{diag}(s')c\}\cong\auU(p')\times\auU(q')$ for the maximal compact
subgroup; $c\in\pcnK_\tau$ is unitary for the standard form, and $c=\mathrm{diag}(c_+,c_-)$ along the
positive and negative rows.

\begin{lemma}[Compact action on Fock functions]\label{pcn:Kaction}
Let $\omega^\flat_\tau(1,\cdot)$ be the action of $\auU(T_\tau)$ (placed at $\tau$) in the
representation $\omega^\flat$ of Lemma~\ref{pcn:slicechar}. For $c\in\pcnK_\tau$ and every
polynomial $P$,
\[
\omega^\flat_\tau(1,c)\bigl(P(Z)e\bigr)=\chi_\tau(c)\,P\bigl(Z(c^{-1}X)\bigr)\,e,\qquad
\chi_\tau(c)=\det(c_+)^{\mu^\flat_\tau-q_0}\det(c_-)^{\mu^\flat_\tau-p_0},
\]
where $Z(c^{-1}X)$ means the coordinates $Z$ evaluated at $c^{-1}X$ (the variables of the other
infinite places are unchanged). In particular the space $\mathcal P_d$ of polynomial vectors $P(Z)e$
with $P$ homogeneous of degree $d$ is finite-dimensional and $\pcnK_\tau$-stable.
\end{lemma}
\begin{proof}
An element $c\in\pcnK_\tau$ preserves the form and the standard euclidean structure of
$\auHom(W_0,T)_\tau$, hence the complex structure defining the Fock model; such a symplectic
transformation acts on polynomial vectors by linear substitution of the variables times a
character: this is the statement ``$K\times K'$ acts by linear substitutions'' in the proof of
Lemma~\ref{par:vac}, applicable since (V) holds by Theorem~\ref{pco:V} (for the standard form of the
action of $\auU(m)\subset\mathrm{Sp}(2m,\R)$ in the Fock model see Folland~\cite[Ch.~4, Prop.~4.39, p.~184]{FollandPS},
and Kashiwara--Vergne~\cite{KashiwaraVergne}). As $\pcnK_\tau$ is connected and polynomial vectors are
dense (Lemma~\ref{pcn:smooth}(b)), this determines the action: $f\mapsto\chi(c)f(c^{-1}\cdot)$ for a
character $\chi$ of $\pcnK_\tau$. (The proof of Lemma~\ref{par:vac} uses about the pair
$(Y_j,Y_{j+1})$ only (V) and the shape of the twisting characters, so it applies verbatim to the
pair $(W_0,T)$ with the character $\gamma_{W_0}$ on $\auU(T)$; the linear-substitution form is the
statement that $\pcnK_\tau$ preserves the positive complex structure, so that the metaplectic
action of $c$ is $f\mapsto f(c^{-1}\cdot)$ up to a scalar character of the metaplectic cover of
$\auU(p')\times\auU(q')$.) As $e$ is a
function of $|X|^2$ and $|c^{-1}X|=|X|$, $f=P(Z)e$ goes to $\chi(c)P(Z(c^{-1}X))e$. The map
$X\mapsto c^{-1}X$ is complex linear in the coordinates $Z$: for $s'_a\varepsilon_j=1$,
$Z_{aj}(c^{-1}X)=\sum_{a'}\overline{(c^{-1})_{aa'}}Z_{a'j}$, and the sum runs over $a'$ with
$s'_{a'}=s'_a$ because $c^{-1}$ is block diagonal; likewise without conjugation if
$s'_a\varepsilon_j=-1$. So homogeneous polynomials of degree $d$ go to homogeneous polynomials of
degree $d$.

The character $\chi$ is the eigencharacter of the vacuum, $\chi(c)=\det(c_+)^{a_+}\det(c_-)^{a_-}$
for integers $a_\pm$, which we read off from the derived action on $e$ ($P=1$). By (V) of
Lemma~\ref{par:vac}, which holds by Theorem~\ref{pco:V}, the untwisted operator
$\omega_{\auHom(W_0,T)}(1,h)$ at $\tau$ is the oscillator action times $\det(h)^{-n/2}$, and the
Weyl-ordered oscillator operator of the diagonal generator of the row $c$ is
$-\sum_ja^-_{jc}a^+_{jc}+\frac n2$ in the notation of \eqref{par:eq-osc} (see the proof of
Lemma~\ref{par:vac}). On $e=1$ the operator $a^-_{jc}a^+_{jc}$ gives $1$ if $\eta_{jc}=
-\varepsilon_js'_c=1$ ($a^+=z$, $a^-=\partial$) and $0$ otherwise ($a^+=-\partial$). The twist of
$\omega^\flat$ on $\auU(T)$ is $\gamma_{W_0,\tau}(\det h)=\det(h)^{\mu^\flat_\tau}$. So the weight of
$e$ on the row $c$ is $-\#\{j:\varepsilon_j\ne s'_c\}+\frac n2-\frac n2+\mu^\flat_\tau$, that is
$\mu^\flat_\tau-q_0$ for a positive row and $\mu^\flat_\tau-p_0$ for a negative row.
\end{proof}

\begin{lemma}[Moving a positive vector into the span]\label{pcn:witt}
Let $x\in\pcMat_{(b+n)\times b}(\C)$ have Gram matrix $x^*\,\mathrm{diag}(s')\,x$ which is not negative
semidefinite, and let $i$ be a row with $s'_i=1$ (it exists by (H2)). There are $u\in\auU(T_\tau)$
and $w\in\C^b$ with $uxw=e_i$, the $i$-th basis vector.
\end{lemma}
\begin{proof}
There is $w$ with $w^*x^*\mathrm{diag}(s')xw>0$; rescaling, the vector $xw$ has length $1$, as does
$e_i$. Both span non-degenerate lines of the same length, so by Witt's theorem for the hermitian space
$T_\tau$ there is an isometry $u$ with $u(xw)=e_i$.
\end{proof}

\begin{lemma}[A negative line orthogonal to the frame]\label{pcn:line}
Let $x''\in\pcMat_{(b+n)\times b}(\C)$ and $w\in\C^b$ with $x''w=e_i$, $s'_i=1$, and let $y_j\in\C^{b+n}$
be given for the $q_0$ columns $j$ with $\varepsilon_j=-1$. There is $r\in\C^{b+n}$ with $r_a=0$ whenever
$s'_a=1$, $r^{\mathsf t}x''=0$, $r^{\mathsf t}y_j=0$ for all $j$ with $\varepsilon_j=-1$, and
$r^{\mathsf t}\bar r=1$.
\end{lemma}
\begin{proof}
Consider the linear map $f\colon\C^{b+n}\to\C^b\times\C^{q_0}\times\C^{p'}$,
$f(r)=(r^{\mathsf t}x'',(r^{\mathsf t}y_j)_j,(r_a)_{s'_a=1})$. By (H1), $q'=b+q_0$, so source and
target both have dimension $b+q_0+p'=b+n$. If $\ker f=0$, $f$ is bijective, and there is $r$ with
$f(r)=(0,0,(1)_a)$; then $r_i=1$, but $r_i=r^{\mathsf t}e_i=r^{\mathsf t}x''w=0$, a contradiction. So
there is $r\ne0$ in $\ker f$; normalise $r^{\mathsf t}\bar r=\sum|r_a|^2=1$.
\end{proof}

\begin{lemma}[Rotations of the line]\label{pcn:rot}
Let $r$ be as in Lemma~\ref{pcn:line}, $\Pi=\bar r\,r^{\mathsf t}$ and, for $|\eta|=1$,
$c_\eta=1+(\bar\eta-1)\Pi$. Then $c_\eta\in\pcnK_\tau$, $c_\eta x''=x''$, and for every polynomial $P$
\[
\omega^\flat_\tau(1,c_\eta)\bigl(P(Z)e\bigr)=\eta^{N_\tau}\,(P\circ\sigma_\eta)(Z)\,e ,
\]
where $N_\tau=\mu_\tau+p_0-\frac n2$ and $\sigma_\eta$ is the linear substitution fixing the variables
of the other places and the variables of the positive rows, and sending, for a negative row $a$,
\[
Z_{aj}\mapsto Z_{aj}+(\bar\eta-1)\,r_a\sum_{a'}\bar r_{a'}Z_{a'j}\quad(\varepsilon_j=-1),\qquad
Z_{aj}\mapsto Z_{aj}+(\eta-1)\,\bar r_a\sum_{a'}r_{a'}Z_{a'j}\quad(\varepsilon_j=1),
\]
the sums running over the negative rows $a'$.
\end{lemma}
\begin{proof}
$\Pi$ is hermitian, $\Pi^2=\Pi$ (as $r^{\mathsf t}\bar r=1$), and $\mathrm{diag}(s')\Pi=\Pi\,
\mathrm{diag}(s')=-\Pi$ ($r$ is supported on negative rows). For a hermitian $M$ with $M\Pi=\Pi M=
\epsilon\Pi$, $c_\eta^*Mc_\eta=M+\epsilon\bigl((\eta-1)+(\bar\eta-1)+(\eta-1)(\bar\eta-1)\bigr)\Pi=M$,
since $\eta\bar\eta=1$. With $M=1$ and $M=\mathrm{diag}(s')$, $c_\eta$ is unitary for both forms, so
$c_\eta\in\pcnK_\tau$; and $c_\eta x''=x''+(\bar\eta-1)\bar r(r^{\mathsf t}x'')=x''$. Its positive block
is $1$ and its negative block is $1+(\bar\eta-1)\bar r_-r_-^{\mathsf t}$, of determinant
$1+(\bar\eta-1)r^{\mathsf t}\bar r=\bar\eta$. By Lemma~\ref{pcn:Kaction},
$\chi_\tau(c_\eta)=\bar\eta^{\,\mu^\flat_\tau-p_0}=\eta^{p_0-n/2+\mu_\tau}=\eta^{N_\tau}$. Finally
$c_\eta^{-1}=c_{\bar\eta}$, so $(c_\eta^{-1}X)_{aj}=X_{aj}+(\eta-1)\bar r_a\sum_{a'}r_{a'}X_{a'j}$;
this changes only negative rows $a$, and involves only negative rows $a'$. If $\varepsilon_j=1$ then
$Z_{a'j}=X_{a'j}$ for all negative $a'$ and this is the second formula; if $\varepsilon_j=-1$ then
$Z_{a'j}=\overline{X_{a'j}}$ and conjugating gives the first.
\end{proof}

\begin{lemma}[Fixed polynomial vectors vanish on a hyperplane section]\label{pcn:fixed}
Assume (H3). Let $r$ be as in Lemma~\ref{pcn:line} and $P$ a polynomial such that
$\omega^\flat_\tau(1,c_\eta)(P(Z)e)=P(Z)e$ for all $|\eta|=1$. Then $P(z)=0$ at every point $z$ (a value
of all the variables) with $\sum_{a'}\bar r_{a'}z_{a'j}=0$ for every $j$ with $\varepsilon_j=-1$.
\end{lemma}
\begin{proof}
By Lemma~\ref{pcn:rot} and injectivity of $P\mapsto P(Z)e$, $\eta^{N_\tau}P\circ\sigma_\eta=P$ as
polynomials, for $|\eta|=1$. At such a point $z$, the first formula of Lemma~\ref{pcn:rot} leaves the
variables $z_{aj}$ ($\varepsilon_j=-1$) unchanged, so $\sigma_\eta(z)=z+(\eta-1)z'$, where $z'_{aj}=
\bar r_a\sum_{a'}r_{a'}z_{a'j}$ for negative rows $a$ and $\varepsilon_j=1$, and $z'=0$ otherwise. The
function $g(\eta)=\eta^{N_\tau}P(z+(\eta-1)z')$ is a polynomial in $\eta$, equal to $P(z)$ on the unit
circle, hence identically $P(z)$; and $g(0)=0$ because $N_\tau\ge1$. So $P(z)=0$.
\end{proof}

\begin{lemma}[No fixed polynomial vectors]\label{pcn:nofix}
Assume (H1)--(H3). Let $x''$, $w$, $i$ be as in Lemma~\ref{pcn:line} and
$\pcnK'=\{c\in\pcnK_\tau:cx''=x''\}$. If a polynomial vector $P(Z)e$ is fixed by
$\omega^\flat_\tau(1,c)$ for all $c\in\pcnK'$, then $P=0$.
\end{lemma}
\begin{proof}
Fix a point $z$. Apply Lemma~\ref{pcn:line} with $y_j=(\overline{z_{\tau,aj}})_a$; then $r^{\mathsf t}y_j=0$
gives, after conjugation, $\sum_a\bar r_az_{\tau,aj}=0$ for $\varepsilon_j=-1$. By
Lemma~\ref{pcn:rot} all $c_\eta$ lie in $\pcnK'$, so they fix $P(Z)e$, and Lemma~\ref{pcn:fixed} gives
$P(z)=0$. As $z$ is arbitrary, $P=0$.
\end{proof}

\begin{lemma}[No invariant functionals]\label{pcn:nofunct}
Assume (H1)--(H3) and let $\pcnK'$ be as in Lemma~\ref{pcn:nofix}. Let $\mathcal S$ be either the space
$\mathcal S(\pzF)$ of Lemma~\ref{pcn:smooth} for the Fock space of $\auHom(W_0,T)_\R$ (all infinite places), or the
same space for the Fock space of $\auHom(W_0,T)_\tau$ alone. Every continuous linear functional
$l$ on $\mathcal S$ with $l\circ\omega^\flat_\tau(1,c)=l$ for all $c\in\pcnK'$ is zero.
\end{lemma}
\begin{proof}
$\pcnK'$ is a closed subgroup of the compact group $\pcnK_\tau$, hence compact; let $dc$ be its
probability Haar measure. Let $f\in\mathcal S$ be a polynomial vector, homogeneous of degree $d$ in the
variables of $\tau$ (any polynomial vector is a finite sum of such vectors times polynomials in the variables of
the other places, on which $\omega^\flat_\tau(1,c)$ acts trivially). By Lemma~\ref{pcn:Kaction}, $c\mapsto
\omega^\flat_\tau(1,c)f$ is a continuous (indeed polynomial in the entries of $c$, $\bar c$) map into
a finite-dimensional space of such vectors, so $m=\int_{\pcnK'}\omega^\flat_\tau(1,c)f\,dc$ lies in it,
and $m$ is fixed by $\pcnK'$ (invariance of $dc$). By Lemma~\ref{pcn:nofix} (applied to $m$, resp.\ to $m\otimes
e^\tau$ with $e^\tau$ the vacuum of the other infinite places in the local case), $m=0$. As $l$ restricted
to a finite-dimensional space is continuous, $l(m)=\int l(\omega^\flat_\tau(1,c)f)
\,dc=l(f)$. So $l$ vanishes on every polynomial vector. Polynomial vectors are dense in $\mathcal S$
(Lemma~\ref{pcn:smooth}(b)). By continuity, $l=0$.
\end{proof}

\subsection{A local vanishing statement at the place $\tau$}\label{pcn:ss-local}

Keep the place $\tau$ of \S\ref{pcn:ss-fock}. Write $G_\tau=\auU(T_\tau)$, $V_\ell=\auHom(\ell,T)_\tau\cong
\pcMat_{(b+n)\times b}(\C)\cong\C^{p}$ ($p=(b+n)b$; a real vector space of dimension $2p$ with euclidean norm
$|\cdot|$), and $Q(x)=x^*J_Tx\in\mathrm{Herm}_b(\C)$ for $x\in V_\ell$. Let $\pzF_\tau$ and $\pzF^\flat_\tau$ be the
Fock spaces of $\auHom(X,T)_\tau$ and $\auHom(W_0,T)_\tau$, and $\mathcal S_\tau=\mathcal S(\pzF_\tau)$, $\mathcal
S_\flat=\mathcal S(\pzF^\flat_\tau)$ their spaces of smooth vectors (Lemma~\ref{pcn:smooth}), with the seminorms
$p_k$. On $\pzF^\flat_\tau$, $G_\tau$ acts by $\omega^\flat_\tau(1,\cdot)$ (Lemma~\ref{pcn:slicechar}); by
Lemma~\ref{pcn:smoothaction} (with Lemma~\ref{par:vac} and Theorem~\ref{pco:V}) this action on $\mathcal S_\flat$ is
continuous and smooth in the sense of that lemma. Let $\mathfrak z_\tau$ be the centre $Z_D$ at $\tau$, written
additively: the space of skew-hermitian $b\times b$ complex matrices (\S\ref{pcu:ss-top-setting}). By
Lemma~\ref{pcu:laws}, $z\in\mathfrak z_\tau$ acts on the slice at $x$ through the character
$e(z,Q(x))$, where $e(z,y)=\psi_\tau(-\frac12\jmath\operatorname{tr}(zy))$ for $y\in\mathrm{Herm}_b(\C)$ (with $\jmath$ read
through $\tau$, $\tau(\jmath)=i/\sqrt d$). Writing $z=iH$ with $H$ hermitian,
$-\frac12\jmath\operatorname{tr}(zy)=\operatorname{tr}(Hy)/(2\sqrt d)$ is real, and $\psi_\tau(t)=e^{-4\pi it}$ for real $t$
($\psi_\tau=\psi_{\Q,\infty}\circ\operatorname{Tr}_{\C/\R}$). As $(H,y)\mapsto\operatorname{tr}(Hy)$ is a positive definite real
inner product on $\mathrm{Herm}_b(\C)$, the pairing $(z,y)\mapsto e(z,y)$ is non-degenerate: every unitary character of
$\mathrm{Herm}_b(\C)$ is $e(z,\cdot)$ for exactly one $z\in\mathfrak z_\tau$, and the real-linear forms
$y\mapsto-\frac12\jmath\operatorname{tr}(zy)$, $z\in\mathfrak z_\tau$, span the dual of $\mathrm{Herm}_b(\C)$.

\emph{The mixed model at $\tau$.} By \S\ref{pcu:ss-mixed} at the single place $\tau$, there is an isomorphism
$\mathsf M\colon\pzF_\tau\to L^2(V_\ell;\pzF^\flat_\tau)$, $(\mathsf Mv)(x)=\mathrm{sl}_D(x)v$, compatible with the
Heisenberg operators as in \eqref{pcu:eq-mixed}; it is unitary up to a constant $c_{\mathsf M}>0$,
$\|\mathsf Mv\|^2=c_{\mathsf M}\|v\|^2$, by Lemma~\ref{pcu:mixedmodel} (at the single place $\tau$). In the notation of \eqref{pcu:eq-mixed}, taking two of
$a',w,b$ equal to $0$,
\begin{equation}\label{pcn:eq-heis}
\mathsf M\psrho(La')\mathsf M^{-1}f(x)=f(x+a'),\quad\mathsf M\psrho(Kb)\mathsf M^{-1}f(x)=e^{2i\,\Im\,b^*x}f(x),\quad
\mathsf M\psrho(Iw)\mathsf M^{-1}f(x)=\psrho^\flat(w)f(x),
\end{equation}
and $\mathsf Me=e^{-|x|^2}e^\flat$. By the local form of Lemma~\ref{pcu:laws} (Corollary~\ref{pcu:laws-local}) and
by Lemma~\ref{pcn:slicechar} together with (M3) for the character, for $v\in\pzF_\tau$, $z\in\mathfrak z_\tau$, $h\in G_\tau$:
\begin{equation}\label{pcn:eq-mixedlaws}
\mathsf M(\omega_\tau(n(0,z),1)v)(x)=e(z,Q(x))\,(\mathsf Mv)(x),\qquad
\mathsf M(\omega_\tau(1,h)v)(x)=\omega^\flat_\tau(1,h)\bigl((\mathsf Mv)(h^{-1}x)\bigr).
\end{equation}
Let $\mathcal T$ be the space of functions $V_\ell\to\mathcal S_\flat$ which are finite sums of elements of
$C_c^\infty(V_\ell;\mathcal S_\flat)$ and of functions $g\otimes u$, $g$ a Schwartz function on $V_\ell$, $u\in\mathcal
S_\flat$, with the seminorms $q_r(f)=\max_{|\beta|\le r}\sup_x(1+|x|)^{r+2p}p_r(\partial^\beta f(x))$
($\partial^\beta$ the real partial derivatives on $V_\ell$). For a compact $B\subset V_\ell$, $C_c^\infty(B;\mathcal S_\flat)$
(functions supported in $B$) is a Fr\'echet space for the $q_r$, and $C_c^\infty(V_\ell;\mathcal S_\flat)$ is their union.

\begin{lemma}[Smooth vectors in the mixed model]\label{pcn:mixed}
\textup{(a)} For $f\in\mathcal T$, $\iota(f)=\mathsf M^{-1}f$ lies in $\mathcal S_\tau$, and for every $k$ there are
$C,r$ with $p_k(\iota f)\le C\,q_r(f)$.
\textup{(b)} $\iota(C_c^\infty(V_\ell;\mathcal S_\flat))$ is dense in $\mathcal S_\tau$.
\textup{(c)} For $f\in C_c^\infty(V_\ell;\mathcal S_\flat)$, $z\in\mathfrak z_\tau$, $h\in G_\tau$:
$\omega_\tau(n(0,z),1)\iota(f)=\iota(e'_zf)$ with $e'_z(x)=e(z,Q(x))$, and
$\omega_\tau(1,h)\iota(f)=\iota(h\cdot f)$ with $(h\cdot f)(x)=\omega^\flat_\tau(1,h)f(h^{-1}x)$; both $e'_zf$ and
$h\cdot f$ lie in $C_c^\infty(V_\ell;\mathcal S_\flat)$.
\end{lemma}
\begin{proof}
(a) Let $D_{a'}$ be the derivative in the direction $a'$ and $m_b$ the multiplication by $2i\,\Im\,b^*x$. For
$f\in\mathcal T$, $t^{-1}(f(\cdot+ta')-f)\to D_{a'}f$ in $L^2(V_\ell;\pzF^\flat_\tau)$ (mean value inequality and
dominated convergence, the difference quotients being bounded by $\sup_{|s|\le1}\|D_{a'}f(x+sa')\|\le
C(1+|x|)^{-2p}q_1(f)$); $t^{-1}(e^{2it\Im b^*x}-1)f\to m_bf$ likewise; and $t^{-1}(\psrho^\flat(tw)f(x)-f(x))\to
\partial^\flat_wf(x)$ for each $x$, with $\|t^{-1}(\psrho^\flat(tw)u-u)\|\le\|\partial^\flat_wu\|$ ($u\in\mathcal S_\flat$;
write $\psrho^\flat(tw)u-u=\int_0^t\psrho^\flat(sw)\partial^\flat_wu\,ds$), so again in $L^2$. By \eqref{pcn:eq-heis},
$\iota(f)$ is in the domains of $\partial_{La'}$, $\partial_{Kb}$, $\partial_{Iw}$, and $\mathsf M\partial_{La'}\iota(f)=
D_{a'}f$, $\mathsf M\partial_{Kb}\iota(f)=m_bf$, $\mathsf M\partial_{Iw}\iota(f)=\partial^\flat_wf$ (applied to the values;
Lemma~\ref{pcn:smooth}(a)). These three operators map $\mathcal T$ into itself. Let $\mathsf B'$ be a real basis of
$\auHom(X,T)_\tau$ made of vectors $La'$, $Kb$, $Iw$. By induction, $\iota(f)$ is in the domain of every product
$\partial_{u_1}\cdots\partial_{u_j}$, $u_i\in\mathsf B'$, and $\mathsf M\partial_{u_1}\cdots\partial_{u_j}\iota(f)=P f$ for
a product $P$ of $j$ of the operators $D_{a'}$, $m_b$, $\partial^\flat_w$. By Lemma~\ref{pcn:smooth}(a) and the
definition of $q_r$, $\|Pf\|_{L^2}^2\le\int C(1+|x|)^{2j}\,p_j(\text{derivatives of order}\le j\text{ of }f(x))^2dx
\le C'q_{2j}(f)^2$. Lemma~\ref{pcn:smooth}(b) for the basis $\mathsf B'$ gives $\iota(f)\in\mathcal S_\tau$ and the bound.

(b) By Lemma~\ref{pcn:smooth}(b) the polynomial vectors are dense in $\mathcal S_\tau$; each is a combination of
products of operators $a^+_k$ applied to the vacuum $e$, and each $a^+_k$ is a combination of operators
$\partial_u$, $u\in\mathsf B'$ (by \eqref{pcn:eq-dw} and linearity). As $\mathsf Me=e^{-|x|^2}e^\flat$ lies in
$\mathcal S(V_\ell)\otimes\mathcal S_\flat$, which the three operators of (a) preserve, $\mathsf M$ maps every polynomial
vector into the algebraic tensor product $\mathcal S(V_\ell)\otimes\mathcal S_\flat$. For $g\in\mathcal S(V_\ell)$,
$u\in\mathcal S_\flat$ and a cut-off $\chi$ equal to $1$ near $0$, $q_r((1-\chi(\cdot/R))g\otimes u)\to0$ as
$R\to\infty$, so by (a) $\iota(\chi(\cdot/R)g\otimes u)\to\iota(g\otimes u)$ in $\mathcal S_\tau$. Hence the closure of
$\iota(C_c^\infty(V_\ell;\mathcal S_\flat))$ contains every polynomial vector.

(c) The identities follow from \eqref{pcn:eq-mixedlaws} and the injectivity of $\mathsf M$. $e'_zf$ is smooth with
the same support; $h\cdot f$ is smooth with compact support because $(x,u)\mapsto\omega^\flat_\tau(1,h)u$ is
continuous linear in $u$ (Lemma~\ref{pcn:smooth}(d)).
\end{proof}

\begin{lemma}[Local vanishing]\label{pcn:local}
Assume (H1)--(H3). Let $\beta_\tau\in\mathrm{Herm}_b(\C)$ be invertible and not negative definite, and let
$\lambda$ be a continuous linear functional on $\mathcal S_\tau$ such that, for all $v\in\mathcal S_\tau$,
$z\in\mathfrak z_\tau$ and $h\in G_\tau$,
\[
\text{(i)}\ \ \lambda(\omega_\tau(n(0,z),1)v)=e(z,\beta_\tau)\,\lambda(v),
\qquad\text{(ii)}\ \ \lambda(\omega_\tau(1,h)v)=\lambda(v).
\]
Then $\lambda=0$.
\end{lemma}
\begin{proof}
Put $\Lambda=\lambda\circ\iota$ on $C_c^\infty(V_\ell;\mathcal S_\flat)$; by Lemma~\ref{pcn:mixed}(a) its restriction to
each $C_c^\infty(B;\mathcal S_\flat)$ is continuous, and by Lemma~\ref{pcn:mixed}(b) it suffices to show $\Lambda=0$.
Write $e_z(x)=e(z,Q(x)-\beta_\tau)=e(z,Q(x))\,e(z,\beta_\tau)^{-1}$, a smooth function of modulus $1$ on
$V_\ell$; by (i) and Lemma~\ref{pcn:mixed}(c), $\Lambda((e_z-1)g)=0$ for all $g\in C_c^\infty(V_\ell;\mathcal S_\flat)$ and
all $z$. Let $\Omega=Q^{-1}(\beta_\tau)\subset V_\ell$.

\emph{Step A: $\Lambda$ kills the functions vanishing on $\Omega$.} (1) Let $x_0\notin\Omega$. By the
non-degeneracy of the pairing there is $z$ with $e_z(x_0)\ne1$, hence $e_z\ne1$ on a neighbourhood $U_{x_0}$. If
$g\in C^\infty_c(V_\ell;\mathcal S_\flat)$ has support in $U_{x_0}$, then $g=(e_z-1)g'$ with $g'=g/(e_z-1)\in
C_c^\infty(V_\ell;\mathcal S_\flat)$, so $\Lambda(g)=0$. (2) For $g\in C_c^\infty(B;\mathcal S_\flat)$ and $z$ fixed,
$s^{-1}(e_{sz}-1)g\to\vartheta_zg$ in $C_c^\infty(B;\mathcal S_\flat)$ as $s\to0$, where $\vartheta_z(x)=2\pi i\,\jmath
\operatorname{tr}(z(Q(x)-\beta_\tau))$ (from $e(z,y)=\exp(2\pi i\jmath\operatorname{tr}(zy))$; every derivative of $s^{-1}(e^{s\vartheta}-1)-
\vartheta$, $\vartheta$ a polynomial, is bounded on $B$ by $|s|$ times a constant). By continuity
$\Lambda(\vartheta_zg)=0$; the real-linear forms $y\mapsto-\frac12\jmath\operatorname{tr}(zy)$ span the dual
of $\mathrm{Herm}_b(\C)$, so $\Lambda(t_kg)=0$ for each real coordinate $t_k$ ($1\le k\le b^2$) of $Q-\beta_\tau$ in a
fixed basis of $\mathrm{Herm}_b(\C)$. (3) \emph{Charts adapted to $\Omega$.} Let $x_1\in\Omega$. The differential
$dQ_{x_1}(\xi)=\xi^*J_Tx_1+x_1^*J_T\xi$ is onto $\mathrm{Herm}_b(\C)$: given $\eta$, $\xi=\frac12x_1\beta_\tau^{-1}\eta$
gives $\eta$ (as $x_1^*J_Tx_1=\beta_\tau$). Let $K_1=\ker dQ_{x_1}$ (real dimension $2p-b^2$) with an orthogonal
projection $\pi_1\colon V_\ell\to K_1$. The map $\Xi(x)=(\pi_1(x-x_1),Q(x)-\beta_\tau)\in K_1\times\mathrm{Herm}_b(\C)$
has bijective differential at $x_1$, so by the inverse function theorem it is a diffeomorphism from an open
neighbourhood $N_{x_1}$ of $x_1$ onto a product $B_1\times B_2$ of open balls centred at $0$; in the coordinates
$(y,t)=\Xi(x)$, $\Omega\cap N_{x_1}=\{t=0\}$. If $g\in C_c^\infty(V_\ell;\mathcal S_\flat)$ has support in $N_{x_1}$ and
$g|_\Omega=0$, write $\tilde g=g\circ\Xi^{-1}$; for $(y,t)\in B_1\times B_2$, $\tilde g(y,t)=\tilde g(y,t)-\tilde g(y,0)=
\sum_kt_k\tilde g_k(y,t)$ with $\tilde g_k(y,t)=\int_0^1\partial_{t_k}\tilde g(y,st)\,ds$ (Hadamard; $B_2$ is convex,
and differentiation under the integral sign is legitimate for the continuous seminorms $p_r$). With a cut-off
$\chi\in C_c^\infty(N_{x_1})$ equal to $1$ on the support of $g$, $g=\sum_kt_k\,g_k$ with $g_k=\chi\cdot(\tilde g_k\circ\Xi)
\in C_c^\infty(V_\ell;\mathcal S_\flat)$, so $\Lambda(g)=0$ by (2). (4) A general $g\in C_c^\infty(V_\ell;\mathcal S_\flat)$ with
$g|_\Omega=0$ is $\sum_i\rho_ig$ for a finite smooth partition of unity on its support subordinate to sets $U_{x_0}$
and $N_{x_1}$; each $\rho_ig$ vanishes on $\Omega$, so $\Lambda(g)=0$. In particular, if $\Omega=\emptyset$ then
$\Lambda=0$ and we are done. Assume henceforth $\Omega\ne\emptyset$.

\emph{Step B: the induced functional on $\Omega$.} By Lemma~\ref{pcu:orbitvol} (at the single place $\tau$, with
$m=b+n$, $\beta_0=\beta_\tau$), $\Omega$ is a closed submanifold and one orbit of $G_\tau$; fix $x_1\in\Omega$ and let
$S_1=\mathrm{Stab}_{G_\tau}(x_1)$, the unitary group of the orthogonal complement of the span of $x_1$, so $G_\tau$
and $S_1$ are unimodular and the orbit map identifies $G_\tau/S_1$ with $\Omega$ (so it is a submersion with smooth
local sections). Cover $\Omega$ by charts $N_{x_i}$ as in Step~A(3), with coordinates $(y,t)$, and let $(\rho_i)$ be a
locally finite smooth partition of unity on $\Omega$ with $\mathrm{supp}\,\rho_i\cap\Omega$ compact in
$N_{x_i}\cap\Omega$; fix $\chi_i\in C_c^\infty(B_{2,i})$ equal to $1$ near $0$. For $f\in C_c^\infty(\Omega;\mathcal S_\flat)$ put
\[
\mathrm{ext}(f)=\sum_i\bigl((y,t)\mapsto\chi_i(t)\,(\rho_if)(\Xi_i^{-1}(y,0))\bigr)\circ\Xi_i ,
\]
each term extended by $0$ outside $N_{x_i}$ (it has compact support in $N_{x_i}$: $\mathrm{supp}(\rho_i)\cap\Omega$ is a
compact subset of $\Xi_i^{-1}(B_{1,i}\times\{0\})$, so its $y$-projection is compact in $B_{1,i}$). The sum is finite on each compact set, so
$\mathrm{ext}$ is a continuous linear map $C_c^\infty(A;\mathcal S_\flat)\to C_c^\infty(V_\ell;\mathcal S_\flat)$ for every
compact $A\subset\Omega$, and $\mathrm{ext}(f)|_\Omega=\sum_i\rho_if=f$. Put $E(f)=\Lambda(\mathrm{ext}(f))$. By Step~A,
$\Lambda(g)=E(g|_\Omega)$ for every $g\in C_c^\infty(V_\ell;\mathcal S_\flat)$ ($g-\mathrm{ext}(g|_\Omega)$ vanishes on
$\Omega$), and $E(f)=\Lambda(g)$ for every compactly supported smooth extension $g$ of $f$. The group $G_\tau$ acts on
$C_c^\infty(\Omega;\mathcal S_\flat)$ by $(h\cdot f)(x)=\omega^\flat_\tau(1,h)f(h^{-1}x)$; $h\cdot\mathrm{ext}(f)$ (the action of
Lemma~\ref{pcn:mixed}(c)) is a compactly supported smooth extension of $h\cdot f$, so by (ii) and
Lemma~\ref{pcn:mixed}(c), $E(h\cdot f)=\Lambda(h\cdot\mathrm{ext}(f))=\lambda(\omega_\tau(1,h)\iota(\mathrm{ext}f))=E(f)$.

\emph{Step C: descent to the fibre.} Let $dg$, $ds$ be Haar measures on $G_\tau$, $S_1$. For $F\in C_c^\infty(G_\tau;\mathcal
S_\flat)$ put $f_F(gx_1)=\int_{S_1}\omega^\flat_\tau(1,gs)F(gs)\,ds$. By Lemma~\ref{pcn:smoothaction}(ii), $g\mapsto
\omega^\flat_\tau(1,g)F(g)$ is smooth with compact support, so the integral (over a compact part of $S_1$ for $g$ in a
compact set) defines a smooth right $S_1$-invariant function of $g$, hence (local sections) an element of
$C_c^\infty(\Omega;\mathcal S_\flat)$, and by the estimate of Lemma~\ref{pcn:smoothaction}(ii) $F\mapsto f_F$ is continuous
from $C_c^\infty(K;\mathcal S_\flat)$ to $C_c^\infty(Kx_1;\mathcal S_\flat)$ for each compact $K\subset G_\tau$ (topologies of
uniform convergence of all derivatives in all $p_k$). It is onto: given $f$, let $K_1$ be a compact neighbourhood of
$\mathrm{supp}f$; choose $a\in C_c^\infty(G_\tau)$, $a\ge0$, with $A(g)=\int_{S_1}a(gs)\,ds>0$ for $gx_1\in K_1$ (cover
the compact set $K_1$ by finitely many open sets $U_jx_1$, $U_j\subset G_\tau$ open relatively compact, and take
$a=\sum a_j$ with $a_j\ge0$ positive on $U_j$), and put $F(g)=a(g)A(g)^{-1}\omega^\flat_\tau(1,g^{-1})f(gx_1)$ where
$gx_1\in K_1$, $F=0$ elsewhere; $F$ is smooth with compact support (Lemma~\ref{pcn:smoothaction}(ii) applied after
the inversion $g\mapsto g^{-1}$) and $f_F=f$, as $A(gs)=A(g)$. Let $\tilde E(F)=E(f_F)$, a linear functional continuous on
each $C_c^\infty(K;\mathcal S_\flat)$. For $(L_hF)(g)=F(h^{-1}g)$, the substitution $g\mapsto hg$ gives $f_{L_hF}=h\cdot
f_F$, so $\tilde E\circ L_h=\tilde E$ by Step~B. For $s\in S_1$ and $(R_sF)(g)=\omega^\flat_\tau(1,s)F(gs)$, right
invariance of $ds$ gives $f_{R_sF}=f_F$, so $\tilde E\circ R_s=\tilde E$.

\emph{Haar measure.} Fix $a\in C_c^\infty(G_\tau)$ with $\int a=1$. For $F\in C_c^\infty(G_\tau;\mathcal S_\flat)$ the map
$h\mapsto L_hF$ is continuous into $C_c^\infty(K';\mathcal S_\flat)$ for a compact $K'$ (left translation commutes with
left-invariant vector fields, and compactly supported continuous functions are uniformly continuous), so
$a*F=\int a(h)L_hF\,dh$ exists in that Fr\'echet space, $\tilde E(a*F)=\int a(h)\tilde E(L_hF)dh=\tilde E(F)$, and
evaluating at $g$, $(a*F)(g)=\int a(h)F(h^{-1}g)dh=\int a(gk^{-1})F(k)\,dk$ ($k=h^{-1}g$; $G_\tau$ is unimodular). The
map $k\mapsto a(\cdot\,k^{-1})\otimes F(k)$ is likewise continuous into some $C_c^\infty(K'';\mathcal S_\flat)$, and its
integral has the same values, so equals $a*F$. Hence
\begin{equation}\label{pcn:eq-haar}
\tilde E(F)=\int_{G_\tau}\tilde E\bigl(a(\cdot\,k^{-1})\otimes F(k)\bigr)\,dk .
\end{equation}
For $\varphi\in\mathcal S_\flat$, $T_\varphi(b)=\tilde E(b\otimes\varphi)$ is a continuous linear functional on $C_c^\infty(G_\tau)$
with $T_\varphi(b(h^{-1}\cdot))=T_\varphi(b)$. Applying \eqref{pcn:eq-haar} to $F=b\otimes\varphi$ gives $T_\varphi(b)=
\int b(k)\phi(k)\,dk$ with $\phi(k)=T_\varphi(a(\cdot\,k^{-1}))$, a continuous function. Left invariance gives
$\int b(k)(\phi(hk)-\phi(k))\,dk=0$ for all $b$, so $\phi(hk)=\phi(k)$ and $\phi$ is a constant $\epsilon(\varphi)$:
$T_\varphi(b)=\epsilon(\varphi)\int b$ (a left-invariant distribution is a multiple of Haar measure). $\epsilon$ is linear,
and continuous ($\epsilon(\varphi)=\tilde E(a\otimes\varphi)$). Then \eqref{pcn:eq-haar} and $\int a(gk^{-1})dg=\int a$ give
$\tilde E(F)=\int_{G_\tau}\epsilon(F(k))\,dk$. Finally $R_s(b\otimes\varphi)=b(\cdot\,s)\otimes\omega^\flat_\tau(1,s)\varphi$ and
$\int b(gs)\,dg=\int b$, so $\tilde E\circ R_s=\tilde E$ gives
\[
\epsilon\circ\omega^\flat_\tau(1,s)=\epsilon\qquad(s\in S_1).
\]

\emph{Step D: no invariant functional on the fibre.} The Gram matrix $\beta_\tau$ of $x_1$ is invertible and not
negative definite, hence not negative semidefinite. Let $u$, $w$, $i$ be as in Lemma~\ref{pcn:witt} for $x_1$,
$x''=ux_1$ and $\pcnK'=\{c\in\pcnK_\tau:cx''=x''\}$; then $u^{-1}\pcnK'u\subset S_1$. The functional
$\epsilon'=\epsilon\circ\omega^\flat_\tau(1,u)^{-1}$ is continuous on $\mathcal S_\flat$ (Lemma~\ref{pcn:smooth}(d)), and for
$c\in\pcnK'$, $\epsilon'\circ\omega^\flat_\tau(1,c)=\epsilon\circ\omega^\flat_\tau(1,u^{-1}cu)\circ\omega^\flat_\tau(1,u)^{-1}=\epsilon'$.
By Lemma~\ref{pcn:nofunct} (the local case), $\epsilon'=0$, so $\epsilon=0$.

\emph{Conclusion.} $\tilde E(F)=\int\epsilon(F(k))dk=0$ for all $F$. As $F\mapsto f_F$ is onto, $E=0$; by Step~B,
$\Lambda(g)=E(g|_\Omega)=0$ for every $g\in C_c^\infty(V_\ell;\mathcal S_\flat)$; so $\lambda=0$.
\end{proof}

\subsection{Proof of Theorem~\ref{pcn:negcoeff}}\label{pcn:ss-proof}

\emph{Step 0: reduction to coherent states.} Let $\Phi=\Phi^{\mathrm{row}}_{z_0}$. For a Fock vector
$w$ with finite moments, Proposition~\ref{pcx:decay}(a), applied with $P=S$ to the pair $(X,T)$ (its
hypothesis holds: by Lemma~\ref{pcu:rowsurj}(a) every $x$ with $\Phi(x)\ne0$ is surjective at every place of $S$,
and $S\cap\bar S=\emptyset$), gives $|\theta(\omega(p,h)[w,\Phi]_S)|\le C\|w\|_M\,\mathrm{ht}(h)^{-N_0}$, with $C$
uniform for $p$ in a compact set. By Lemma~\ref{pcx:superpos},
$\theta(\omega(p,h)[w,\Phi]_S)=\int w(y)\,\theta(\omega(p,h)[\psrho(y)e,\Phi]_S)\,dy$, and since
$\|\psrho(y)e\|_M\le(1+|y|)^M\|e\|_M$ and $\int|w(y)|(1+|y|)^Mdy=\|w\|_M<\infty$, Fubini (with (R1) for the
integral over $[\auU(T)]$ and compactness of $Z_D(E)\backslash Z_D(\auA)$) gives
\[
\pcuF_\beta\Theta^{\mathrm L}_{X,T}(w,\Phi)(p)=\int w(y)\,\pcuF_\beta\Theta^{\mathrm L}_{X,T}(\psrho(y)e,\Phi)(p)\,dy .
\]
So it suffices to prove the theorem for $v=\psrho(y)e$. Up to a factor of modulus $1$ this vector is a product
$v=v_\tau\otimes v^\tau$, $v^\tau=\bigotimes_{w\ne\tau}v_w$, of the vectors $v_w=\psrho_w(y_w)e_w$, which lie in the spaces
$\mathcal S(\pzF_w)$ of Lemma~\ref{pcn:smooth} (the vacuum is a polynomial vector, and Lemma~\ref{pcn:smooth}(d)).

\emph{Step 1: the coefficient as a continuous functional at $\tau$.} Let $\pzF_\infty$ be the archimedean
Fock space of $\auHom(X,T)_\R$, of complex dimension $D_\infty$, $\pzF_\tau$ its factor at $\tau$, and $\mathcal S_\tau=
\mathcal S(\pzF_\tau)$. Write $p=p_\infty p_f$ and $p_\infty=(p_w)_w$, and put $u^\tau=\bigotimes_{w\ne\tau}\omega_w(p_w,1)v_w$,
a vector of $\mathcal S(\bigotimes_{w\ne\tau}\pzF_w)$ (Lemma~\ref{pcn:smooth}(d)). For $\varphi\in\mathcal S_\tau$ let
\[
\lambda(\varphi)=\pcuF_\beta\Theta^{\mathrm L}_{X,T}\bigl(w_\varphi,\Phi\bigr)(p),\qquad
w_\varphi=\omega_\infty(p_\infty,1)^{-1}(\varphi\otimes u^\tau),
\]
so that $\lambda(\omega_\tau(p_\tau,1)v_\tau)$ is the coefficient to be computed ($w_\varphi=v$ for this $\varphi$). Fix $M$ as
in Step~0 and an integer $s$ with $2s>M+2D_\infty$. By Lemma~\ref{pcn:smooth}(c), (d) and the tensor product
inequality for $p_s$,
\[
\|w_\varphi\|_M\le C_{M,s}\|w_\varphi\|_s\le C\|\varphi\otimes u^\tau\|_s\le C'p_s(\varphi\otimes u^\tau)\le C'p_s(u^\tau)\,p_s(\varphi),
\]
with $C$ depending on $\|\sigma_{p_\infty}\|$ only. By Step~0's bound,
\[
|\lambda(\varphi)|\le\sup_{z}\int_{[\auU(T)]}\bigl|\theta\bigl(\omega(n(0,z)p,h)[w_\varphi,\Phi]_S\bigr)\bigr|\,dh
\le C''\|w_\varphi\|_M\le C'''p_s(\varphi)
\]
($z$ in a compact fundamental domain of $Z_D(E)\backslash Z_D(\auA)$, so $n(0,z)p$ stays in a compact set; (R1) for
$\int\mathrm{ht}^{-N_0}<\infty$; probability measure on the centre). So $\lambda$ is a well-defined continuous linear
functional on $\mathcal S_\tau$, every integral involved converging absolutely.

\emph{Step 2: the two invariance properties.} (ii) Let $h_0\in G_\tau$, placed at $\tau$. As $\omega_\tau(1,h_0)$
commutes with $\omega_\infty(p_\infty,1)$ and with the other places,
$\theta(\omega(n(0,z)p,h)[w_{\omega_\tau(1,h_0)\varphi},\Phi]_S)=\theta(\omega(n(0,z)p,hh_0)[w_\varphi,\Phi]_S)$. The integral
over $[\auU(T)]$ of an absolutely integrable function is invariant under right translation by $h_0$ (the quotient
measure of the unimodular group $\auU(T)(\auA)$ by the discrete subgroup $\auU(T)(E)$ is right invariant), so,
integrating then over $z$ (Fubini, by the bounds of Steps~0--1), $\lambda(\omega_\tau(1,h_0)\varphi)=\lambda(\varphi)$.
(i) Let $z_\tau\in\mathfrak z_\tau$ (the centre at $\tau$, \S\ref{pcn:ss-local}) and $n_\tau=n(0,z_\tau)$, placed at $\tau$. Since
$\omega(p,h)[w,\Phi]_S=\omega(1,h)[\omega_\infty(p_\infty,1)w,\omega_f(p_f)\Phi]_S$, we get, with
$\Phi'=\omega_f(p_f)\Phi$,
(here $[\cdot,\Phi']$ denotes the datum whose finite part is $\omega_f(p_f)$ applied to the finite part of $[\cdot,\Phi]_S$)
\[
\theta\bigl(\omega(n(0,z)p,h)[w_{\omega_\tau(n_\tau,1)\varphi},\Phi]_S\bigr)
=\theta\bigl(\omega(n(0,z)n_\tau,h)[\varphi\otimes u^\tau,\Phi']\bigr)
=\theta\bigl(\omega(n(0,z+z_\tau)p,h)[w_\varphi,\Phi]_S\bigr),
\]
because $n(0,\cdot)$ is a homomorphism from the additive group into $Z_D(\auA)$ and $n_\tau$ commutes with the
finite components. Integrating over $h$ and then against $\psi_\beta(z)^{-1}dz$ over the compact group
$Z_D(E)\backslash Z_D(\auA)$ and substituting $z\mapsto z-z_\tau$ (invariance of its Haar measure; the coset of $z+z_\tau$
is well defined since $z_\tau\in Z_D(\auA)$) gives
$\lambda(\omega_\tau(n_\tau,1)\varphi)=\psi_{\beta,\tau}(z_\tau)\lambda(\varphi)$ with $\psi_{\beta,\tau}(z_\tau)=
\psi_\tau(-\frac12\jmath\operatorname{tr}(z_\tau\tau(\beta)))=e(z_\tau,\tau(\beta))$.

\emph{Step 3: conclusion.} Let $\tau$ be the infinite place of the theorem, where $\tau(\beta)$ is invertible and not
negative definite; (H1)--(H3) hold (\S\ref{pcn:ss-statement}). By Steps~1--2, $\lambda$ satisfies the hypotheses of
Lemma~\ref{pcn:local} with $\beta_\tau=\tau(\beta)$ (on $\mathcal S_\tau$ itself; no change of model is needed), so $\lambda=0$; in particular
$\pcuF_\beta\Theta^{\mathrm L}_{X,T}(v,\Phi)(p)=\lambda(\omega_\tau(p_\tau,1)v_\tau)=0$. With Step~0 this proves the theorem
for every $v$ with finite moments. \qed

\begin{remark}[What the argument does not use]\label{pcn:rem}
The proof uses neither the unfolding (U1) nor any integrability of a term over a stabiliser quotient: the only
global inputs are the absolute convergence of the theta integral and of its coefficient
(Proposition~\ref{pcx:decay}(a), (R1)) and the right invariance of the measure on $[\auU(T)]$. The orbit
$\{x:x^*J_Tx=\beta\}$ and its stabiliser enter only locally at $\tau$, through Lemma~\ref{pcn:local}. The
finite part of $\Phi$ plays no role beyond the hypothesis of Proposition~\ref{pcx:decay}(a), and nothing is assumed
at the infinite places other than $\tau$.
\end{remark}

\subsection{Dependencies}\label{pcn:ss-dep}
With Theorem~\ref{pcn:negcoeff}, Lemma~\ref{pct:negcoeff} is proved, and with it
Proposition~\ref{pct:vantop} and Proposition~\ref{pcx:P-nonvan}. The proof uses results proved
in the paper (Lemma~\ref{pcu:laws} in its local form, Corollary~\ref{pcu:laws-local}; Lemmas~\ref{pcu:rowsurj},
\ref{pcu:orbitvol}, \ref{pcx:superpos}, \ref{pz:growth}; Proposition~\ref{pcx:decay}; Lemma~\ref{par:vac} with
Theorem~\ref{pco:V}; the mixed model of \S\ref{pcu:ss-mixed}), the facts on smooth vectors, mixed models, invariant
distributions and adapted charts proved in \S\S\ref{pcn:ss-smooth}, \ref{pcn:ss-local}, and the following published
results: Witt's theorem for hermitian spaces over $\C$; the Fock-model formula for the action of the
maximal compact subgroup (background reference: Folland~\cite[Prop.~4.39]{FollandPS}, Kashiwara--Vergne~\cite{KashiwaraVergne}; the proof rests on Lemma~\ref{par:vac}); Stone's theorem on
one-parameter unitary groups; the inverse function theorem; reduction theory (R1). It does not use (U1). Together with
\S\ref{pct:ss-status} and \S\ref{pcu:ss-final}, Theorem~\ref{pcx:X} depends only on published results (those
listed in \S\ref{pct:ss-status}, and the ones just listed).


%% ===== proofs/C-triangle.tex =====
% proofs/C-triangle.tex
% Archimedean types at the nodes of the triangle (Lemma pc:arch for the pairs of the triangle) and
% the analytic steps (ii), (iii) of Theorem pc:tri. Labels of this file start with ptr:.
% Only amsmath/amssymb/amsthm are used. Conventions and macros: proofs/C-determinant.tex.
\providecommand{\auU}{\mathrm U}
\providecommand{\auPl}{\Phi}
\providecommand{\auHom}{\operatorname{Hom}}
\providecommand{\auNm}{\operatorname{Nm}}
\providecommand{\auGL}{\mathrm{GL}}
\providecommand{\auA}{\mathbb A}
\providecommand{\aubar}[1]{\overline{#1}}
\providecommand{\psrho}{\varrho}
\providecommand{\pzF}{\mathfrak F}
\providecommand{\pcF}{\mathcal F}
\providecommand{\pcG}{\mathsf G}
\providecommand{\pcm}{\mathfrak m}
\providecommand{\pcR}{R}
\providecommand{\pcMat}{\mathrm M}
\providecommand{\pdsm}{m}
\providecommand{\pdsK}{K}
\providecommand{\ptrnot}[1]{\par\smallskip{\raggedright\noindent\textsc{Not proved here.} [#1]\par}\smallskip}
\providecommand{\ptrwt}{\operatorname{wt}}
\providecommand{\ptrht}{\operatorname{ht}}
\providecommand{\ptrsgn}{\operatorname{sgn}}

\section{The archimedean types at the nodes of the triangle}\label{ptr:sec}

This section treats the archimedean part of Lemma~\ref{pc:arch} for the pairs of the triangle of
\S\ref{pc:ss-triangle}, and the steps (ii) (convergence) and (iii) (archimedean multiplicity one)
of the deduction of Theorem~\ref{pc:tri} described at the end of \S\ref{pc:ss-exchange}. The
notation is that of \S\ref{pc:ss-triangle}, of \S\ref{par:sec} and of \S\ref{pz:ss-types}; inner
products are conjugate-linear in the first variable.

\emph{Contents of this section.}
\begin{enumerate}
\item[(1)] The Harish-Chandra parameters of the types $e_{i,n}$ at the nodes, as defined in
\S\ref{pc:ss-triangle}, \emph{match} (in the sense of \eqref{par:eq-match}) along every
ascending pair $(\pcG_{i,n},\pcG_{i,n+1})$ and along every descending pair
$(\pcG_{i+1,n},\pcG_{i,n+1})$, for the characters of the lifts $a_{i,n}$ and for the characters of
the descents determined in Definition~\ref{ptr:desc-char}; the latter are the characters forced by
the hypotheses of the exchange identity (Lemma~\ref{ptr:chars}). This is checked from the
parameters themselves (Proposition~\ref{ptr:match}).
\item[(2)] Consequently (A1) (integrability, cuspidality, rapid decay) and (A2) (the lifts and the
transposed lifts with the joint harmonic preserve the types: the \emph{edge} property) hold at
all nodes and along all these pairs, and the joint harmonic exists and is unique
(Theorem~\ref{ptr:edge}).
\item[(3)] Convergence of all integrals that occur in the squares (Proposition~\ref{ptr:conv}).
\item[(4)] $Z_\infty\neq0$ is reduced to the positivity (F) of the Fock coefficients of a termwise
harmonic polynomial at positive definite elements (Proposition~\ref{ptr:Z}); (F) is
Theorem~\ref{ptq:F}.
\item[(5)] Archimedean multiplicity one (step (iii)) is reduced to the local smoothing statement
(H) (Theorem~\ref{ptr:square}); (H) is Theorem~\ref{ptq:H}.
\end{enumerate}
Step (v) of \S\ref{pc:ss-exchange} (the adjoint law, Theorem~\ref{pc:tri}(c)) is treated in
\S\ref{ptv:sec}.

\subsection{Parameters at the nodes}\label{ptr:ss-param}

\subsubsection*{Notation}
Recall from \S\ref{pc:ss-triangle}: $\pcG_{i,n}$ $(i,n\ge0,\ i+n\le m)$ has rank $N_{i,n}=r+nt$, a
stem of rank $r$ and $n$ positive definite blocks of rank $t$ numbered $i+1,\dots,i+n$. The stem
coordinates are numbered $p=1,\dots,r$; we put $a(p)=(p-k)^+$. The sign of the stem coordinate $p$
is $\epsilon_p^{(i,n)}=(-1)^{E_{i,n}(a(p))}$ with
\begin{equation}\label{ptr:eq-E}
E_{i,n}(a)=i+a+\min\bigl((a-i)^+,n\bigr)=\begin{cases}i+a,&a\le i,\\ 2a,&i<a\le i+n,\\
i+a+n,&a>i+n.\end{cases}
\end{equation}
The stem of the parameter consists of the numbers of $\varsigma_\tau$ (the case $j=0$ of
\eqref{pz:eq-param}): the $k-1$ numbers $\tfrac k2-1,\dots,1-\tfrac k2$ (these are the coordinates
$p<k$) and the numbers $-\kappa(\tau)-i'\Delta$, $1\le i'\le m$ (the coordinate $p=k-1+i'$, so that
$a(p)=i'-1$). The parameter $\Lambda^{(i,n)}_\tau$ of $e_{i,n}$ is: the stem numbers with the signs
$\epsilon^{(i,n)}$, and for each block $i+l$ $(1\le l\le n)$ the $t$ numbers
$-\kappa_{i+l}(\tau)+\tfrac{t-1}2-b$ $(0\le b<t)$ with sign $+$; all numbers lowered by
$\tfrac12$ if $n$ is odd. Here $\kappa_l=\kappa+(l+\tfrac12)\Delta$.

\begin{lemma}[Consistency with the lower chain and with the first column]\label{ptr:cons}
\textup{(a)} $\Lambda^{(0,n)}_\tau=\Lambda^{(n)}_\tau$ of \eqref{pz:eq-param}, with the same signs.
\textup{(b)} $\epsilon^{(i,0)}=(-1)^i\epsilon^{(0,0)}$, and $\Lambda^{(i,0)}_\tau$ and $\Lambda^{(0,0)}_\tau$
define the same discrete series of the common group $\auU(\pcG_{i,0})=\auU(Y_0)$.
\textup{(c)} The number of negative signs of $\pcG_{i,n}$ is at most $r$; the ranks $N_{i,n}$ are
odd exactly when $n$ is even (so $r$ and $t$ are odd).
\end{lemma}
\begin{proof}
(a) For $i=0$, \eqref{ptr:eq-E} gives the sign $+$ when $a\le n$ and $(-1)^{a+n}$ when $a>n$. For
the number $-\kappa-i'\Delta$ ($a=i'-1$) the sign is $-$ exactly when $i'\ge n+2$ and
$i'\equiv n\pmod2$, which is the rule of \eqref{pz:eq-param} for $j=n$; the small numbers ($a=0$)
have sign $+$; the blocks and the half-shift are the same. (b) For $n=0$, $E_{i,0}(a)=i+a$ in all
cases, while $E_{0,0}(a)=a$. The group of a hermitian form and of its negative is the same, the
compact Cartan subgroup and the root system are the same, and the compact roots are those
$e_a-e_b$ with equal signs, which does not change when all signs change; the discrete series with
Harish-Chandra parameter $\lambda$ is characterised by $\lambda$ and the set of compact roots
(Harish-Chandra; \cite[Theorem~9.20, p.~310]{Knapp}, stated for linear connected semisimple groups; it
applies to $\auU(p,q)=Z\cdot\mathrm{SU}(p,q)$, $Z\cong\auU(1)$ the compact centre,
$Z\cap\mathrm{SU}(p,q)$ finite, since a discrete series of $\auU(p,q)$ restricts irreducibly to one of
$\mathrm{SU}(p,q)$ and $Z$ acts by a character, which is determined by $\lambda$ as $Z$ is connected;
see the paragraph ``On (D3)'' in \S\ref{par:ss-rigid}). (c) The negative signs occur only on
the stem. For (c), second part: $N_j=r+jt$ is odd exactly when $j$ is even (proof of
Lemma~\ref{par:vac}); hence $r$ is odd and $t$ is odd.
\end{proof}

\subsubsection*{The derived constants for arbitrary characters}
For automorphic characters $\alpha,\beta$ of $\auU(1)$ write $\ptrwt\alpha=(\ptrwt_\tau\alpha)_\tau$
for their weights, normalised as in \S\ref{par:sec} ($\upsilon_\tau$ has weight $1$,
$\upsilon_k=\upsilon^{\lfloor k/2\rfloor}$, $\mu_j$ has weight $-\kappa_j$).

\begin{lemma}\label{ptr:const}
Let $S,T$ be spaces of ranks $n<n'$ and $\alpha,\beta$ characters. Under \textup{(V)}
(Lemma~\ref{par:vac}, proved in Theorem~\ref{pco:V}) the derived actions of
$\omega_{\alpha,\beta,\infty}(\cdot,1)$ and $\omega_{\alpha,\beta,\infty}(1,\cdot)$ at $\tau$ on
polynomial vectors are \eqref{par:eq-osc} with $\gamma=\tau_0+\tfrac{n'}2$, $\gamma'=\tau'_0+\tfrac n2$,
\[
\tau_0=\tfrac{n'}2-\ptrwt_\tau\alpha,\qquad\tau'_0=\ptrwt_\tau\beta-\tfrac n2 .
\]
\end{lemma}
\begin{proof}
The proof of Lemma~\ref{par:vac} uses only (V) and the weights of the two characters: the factor
$\alpha(\det g)^{-1}\beta(\det h)$ contributes $\det(g)^{-\ptrwt\alpha}\det(h)^{\ptrwt\beta}$ at
$\tau$, (V) contributes $\det(g)^{n'/2}\det(h)^{-n/2}$, and the Weyl-ordered compact operators carry
the constants $\tfrac{n'}2$, $\tfrac n2$. (For $\omega_{\mu_{j+1}}=\omega_{\mu_{j+1}\upsilon_{n'},
\mu_{j+1}\upsilon_n}$ this gives back $\tau_0=\kappa_{j+1}+\tfrac{n'}2-\lfloor\tfrac{n'}2\rfloor$,
$\tau'_0=-\kappa_{j+1}+\lfloor\tfrac n2\rfloor-\tfrac n2$.)
\end{proof}

By the computation in the proof of Lemma~\ref{par:match} (which uses only the definitions of
$\xi,\xi'$), for a row parameter $\lambda$ and a column parameter $\lambda'$ the conditions
\eqref{par:eq-match} are equivalent to the following \emph{block rule}: there is a set $B$ of $t$
columns and a bijection $r\mapsto c(r)$ of the rows onto the other columns such that
\begin{equation}\label{ptr:eq-rule}
\begin{gathered}
\lambda'_{c(r)}=-\lambda_r+\tau_0+\tau'_0,\qquad|\lambda'_{c(r)}-\tau'_0|>\tfrac t2,\qquad
\varepsilon'_{c(r)}=\varepsilon_r\,\ptrsgn(\lambda'_{c(r)}-\tau'_0),\\
\{\lambda'_c:c\in B\}=\{\tau'_0+\tfrac{t-1}2-b:0\le b<t\},\qquad\varepsilon'\text{ constant on }B.
\end{gathered}
\end{equation}
Indeed $\xi_r=\lambda_r-\tau_0$, $\xi'_c=\lambda'_c-\tau'_0$, and $\eta_{rc(r)}=-\varepsilon_r
\varepsilon'_{c(r)}=\ptrsgn\xi_r=-\ptrsgn\xi'_{c(r)}$. In words: the column parameter is the
negative of the row parameter shifted by $\tau_0+\tau'_0$, plus a string of $t$ numbers of one
sign centred at $\tau'_0$; the signs are kept above the centre and changed below it.

\subsubsection*{The two kinds of pairs}
An \emph{ascending pair} is $(\pcG_{i,n},\pcG_{i,n+1})$ with the characters of $a_{i,n}$:
$\alpha=\mu_{i+n+1}\upsilon_{r+(n+1)t}$, $\beta=\mu_{i+n+1}\upsilon_{r+nt}$; rows are the
coordinates of $\pcG_{i,n}$, with $\lambda=-\Lambda^{(i,n)}_\tau$; columns those of $\pcG_{i,n+1}$,
with $\lambda'=\Lambda^{(i,n+1)}_\tau$. Thus, as in \S\ref{par:sec}, the lift with the joint
harmonic is to take forms of type $e_{i,n}$ to forms of type $e_{i,n+1}$.

A \emph{descending pair} is $(\pcG_{i+1,n},\pcG_{i,n+1})$, $i+n+1\le m$, rows the coordinates of
$\pcG_{i+1,n}$ with $\lambda=\Lambda^{(i+1,n)}_\tau$, columns those of $\pcG_{i,n+1}$ with
$\lambda'=-\Lambda^{(i,n+1)}_\tau$. Then (as in Proposition~\ref{pz:types}(c)) the transposed lift
with the joint harmonic is to take forms of type $e_{i,n+1}$ (the conjugate of the type of the
column parameter $-\Lambda^{(i,n+1)}$) to forms of type $e_{i+1,n}$ (the type of the row parameter):
this is the shape of the descent $d_{i,n}\colon\pcF_{i,n+1}\to\pcF_{i+1,n}$.

\begin{definition}[Characters of the descents]\label{ptr:desc-char}
For $i+n+1\le m$ let $\alpha^\downarrow_{i,n},\beta^\downarrow_{i,n}$ be automorphic characters of
$\auU(1)$ with
\[
\ptrwt_\tau\alpha^\downarrow_{i,n}=\kappa_{i+1}(\tau)+\Bigl\lceil\frac{r+(n+1)t}2\Bigr\rceil,\qquad
\ptrwt_\tau\beta^\downarrow_{i,n}=\kappa_{i+1}(\tau)+\Bigl\lceil\frac{r+nt}2\Bigr\rceil
\]
(their finite parts are fixed in Lemma~\ref{ptr:chars}). The descent kernel is
$\theta_{\alpha^\downarrow_{i,n},\beta^\downarrow_{i,n}}$ for $\auHom(\pcG_{i+1,n},\pcG_{i,n+1})$.
\end{definition}

\begin{proposition}[Matching at the nodes]\label{ptr:match}
Under \textup{(V)}, for every $\tau$:
\textup{(a)} every ascending pair satisfies the block rule \eqref{ptr:eq-rule}, with
$\tau_0=\kappa_{i+n+1}+\tfrac12[n\text{ odd}]$, $\tau'_0=-\kappa_{i+n+1}-\tfrac12[n\text{ even}]$, the
string $B$ being block $i+n+1$;
\textup{(b)} every descending pair, with the characters of Definition~\ref{ptr:desc-char},
satisfies \eqref{ptr:eq-rule}, with $\tau_0=-\kappa_{i+1}-\tfrac12[n\text{ odd}]$,
$\tau'_0=\kappa_{i+1}+\tfrac12[n\text{ even}]$, the string $B$ being block $i+1$;
\textup{(c)} conversely, the weights of Definition~\ref{ptr:desc-char} are the only ones for which
\textup{(b)} holds with the string $B$ equal to block $i+1$ and the old coordinates paired with the
same coordinates.
\end{proposition}
\begin{proof}
Throughout, $[P]$ is $1$ if $P$ holds and $0$ otherwise, and we use Lemma~\ref{ptr:cons}(c):
$N_{i,n}=r+nt$ is odd iff $n$ is even, so $\tfrac{N_{i,n}}2-\lfloor\tfrac{N_{i,n}}2\rfloor=\tfrac12
[n\text{ even}]$ and $\lceil\tfrac{N_{i,n}}2\rceil-\tfrac{N_{i,n}}2=\tfrac12[n\text{ even}]$.

(a) \emph{Constants.} By Lemma~\ref{ptr:const} with $n'=N_{i,n+1}$, $\ptrwt\alpha=-\kappa_{i+n+1}+
\lfloor n'/2\rfloor$: $\tau_0=\kappa_{i+n+1}+\tfrac12[n+1\text{ even}]$; with $\ptrwt\beta=
-\kappa_{i+n+1}+\lfloor N_{i,n}/2\rfloor$: $\tau'_0=-\kappa_{i+n+1}-\tfrac12[n\text{ even}]$. So
$\tau_0+\tau'_0=\tfrac12[n\text{ odd}]-\tfrac12[n\text{ even}]$.
\emph{Old numbers.} With $\lambda=-\Lambda^{(i,n)}$ the first condition of \eqref{ptr:eq-rule}
reads $\Lambda^{(i,n+1)}_{c(r)}=\Lambda^{(i,n)}_r+\tau_0+\tau'_0$. Pair each stem coordinate and
each coordinate of the blocks $i+1,\dots,i+n$ of $\pcG_{i,n}$ with the same coordinate of
$\pcG_{i,n+1}$. Both carry the same unshifted number, lowered by $\tfrac12$ according to the parity
of $n$, resp.\ $n+1$; the difference is $-\tfrac12$ if $n$ is even and $+\tfrac12$ if $n$ is odd,
which is $\tau_0+\tau'_0$.
\emph{New string.} Block $i+n+1$ of $\pcG_{i,n+1}$ carries $-\kappa_{i+n+1}+\tfrac{t-1}2-b$
lowered by $\tfrac12[n+1\text{ odd}]=\tfrac12[n\text{ even}]$: the string centred at $\tau'_0$, with
the constant sign $+$.
\emph{Distances.} The block $i+l$ $(l\le n)$ is centred at $-\kappa_{i+l}\pm\tfrac12$, at distance at
least $(n+1-l)\Delta-1\ge\Delta-1$ from $\tau'_0$, so its numbers are at distance $\ge\Delta-\tfrac{t+1}2>
\tfrac t2$. The stem numbers are $\tfrac k2-1,\dots$ (at distance $>\kappa_{i+n+1}-k>\tfrac t2$) and
$-\kappa-i'\Delta$, at distance $|(i+n+\tfrac32-i')\Delta|-\tfrac12\ge\tfrac12\Delta-\tfrac12>\tfrac t2$,
since $\Delta-t\ge r^8+2$ (\S\ref{pz:ss-types}).
\emph{Signs.} The blocks lie above the centre (their centres $-\kappa_{i+l}$ exceed
$-\kappa_{i+n+1}$) and keep the sign $+$; the small stem numbers keep their sign $(-1)^i$. The number $-\kappa-i'\Delta$,
$a=i'-1$, lies above the centre iff $i'<i+n+\tfrac32$, i.e.\ $a\le i+n$. By \eqref{ptr:eq-E}:
for $a\le i$, $E_{i,n}(a)=E_{i,n+1}(a)=i+a$; for $i<a\le i+n$ both are $2a$; for $a=i+n+1$,
$E_{i,n}(a)=2i+2n+1$ is odd and $E_{i,n+1}(a)=2a$ is even; for $a>i+n+1$, $E_{i,n+1}(a)=E_{i,n}(a)+1$.
So the sign is kept exactly when $a\le i+n$, which is the rule. (The small numbers have $a=0\le i$.)

(b) \emph{Constants.} Here $n'=N_{i,n+1}$ (columns) and the rows have rank $N_{i+1,n}$. By
Lemma~\ref{ptr:const}, $\tau_0=\tfrac{n'}2-\kappa_{i+1}-\lceil\tfrac{n'}2\rceil=-\kappa_{i+1}-\tfrac12
[n+1\text{ even}]=-\kappa_{i+1}-\tfrac12[n\text{ odd}]$, and $\tau'_0=\kappa_{i+1}+\lceil
\tfrac{N_{i+1,n}}2\rceil-\tfrac{N_{i+1,n}}2=\kappa_{i+1}+\tfrac12[n\text{ even}]$. So
$\tau_0+\tau'_0=\tfrac12[n\text{ even}]-\tfrac12[n\text{ odd}]$.
\emph{Old numbers.} With $\lambda=\Lambda^{(i+1,n)}$, $\lambda'=-\Lambda^{(i,n+1)}$ the first condition
reads $\Lambda^{(i,n+1)}_{c(r)}=\Lambda^{(i+1,n)}_r-(\tau_0+\tau'_0)$. Pair the stem coordinates, and
the blocks $i+2,\dots,i+n+1$ of $\pcG_{i+1,n}$ with the same blocks of $\pcG_{i,n+1}$. Both carry
the same unshifted number, lowered according to $n$, resp.\ $n+1$; the difference
$\Lambda^{(i+1,n)}-\Lambda^{(i,n+1)}$ is $-\tfrac12[n\text{ odd}]+\tfrac12[n\text{ even}]=\tau_0+\tau'_0$.
\emph{New string.} Block $i+1$ of $\pcG_{i,n+1}$ carries $-\kappa_{i+1}+\tfrac{t-1}2-b-\tfrac12
[n\text{ even}]$; its negatives, the entries of $\lambda'$, are $\kappa_{i+1}+\tfrac12[n\text{ even}]
-\tfrac{t-1}2+b$, the string centred at $\tau'_0$, with the constant sign $+$.
\emph{Distances.} The entries of $\lambda'$ on old coordinates are the negatives of those of
$\Lambda^{(i,n+1)}$: the blocks $i+l$, $l\ge2$, centred at $\kappa_{i+l}\pm\tfrac12$, at distance
$\ge\Delta-1$ from $\tau'_0$; the small stem numbers, at distance $>\kappa_{i+1}-k$; the numbers
$\kappa+i'\Delta\pm\tfrac12$, at distance $\ge\tfrac12\Delta-\tfrac12$. All exceed $\tfrac t2$.
\emph{Signs.} The blocks $i+l$, $l\ge2$, lie above $\tau'_0$ and keep $+$. The small stem
numbers (negated) lie below and change sign: indeed $E_{i+1,n}(0)=i+1$, $E_{i,n+1}(0)=i$. The number
$\kappa+i'\Delta$ ($a=i'-1$) lies above $\kappa_{i+1}+\tfrac12[\,]$ iff $i'\ge i+2$, i.e.\ $a\ge i+1$.
By \eqref{ptr:eq-E}: for $a\le i$, $E_{i+1,n}(a)=i+1+a$ and $E_{i,n+1}(a)=i+a$ differ by one; for
$a=i+1$, $E_{i+1,n}(a)=2i+2$ and $E_{i,n+1}(a)=2a$ are both even; for $i+1<a\le i+n+1$ both are $2a$;
for $a>i+n+1$ both are $i+a+n+1$. So the sign is changed exactly when $a\le i$, i.e.\ below the
centre, which is the rule.

(c) The block rule determines $\tau'_0$ (the centre of the new string) and $\tau_0+\tau'_0$ (the
shift), hence $\tau_0,\tau'_0$, hence by Lemma~\ref{ptr:const} the weights.
\end{proof}

\begin{remark}\label{ptr:rem-desc}
For $n=0$ the descending pair is $(\pcG_{i+1,0},\pcG_{i,1})=(\pcG_{i,0}^-,\pcG_{i,1})$, and by (W7)
(Lemma~\ref{pc:SW}) the complex conjugate of the kernel of $a_{i,0}$ is a kernel for this pair,
with parameters negated: $(-\Lambda^{(i,0)},\Lambda^{(i,1)})\mapsto(\Lambda^{(i,0)},-\Lambda^{(i,1)})$,
and $\Lambda^{(i,0)}$, $\Lambda^{(i+1,0)}$ define the same representation (Lemma~\ref{ptr:cons}(b)).
This is the descending pair of Proposition~\ref{ptr:match}(b) for $n=0$, consistent with
$d_{i,0}=a_{i,0}^*$ (Theorem~\ref{pc:tri}(a)); by (c) the weights of the conjugate characters are
those of Definition~\ref{ptr:desc-char}.
\end{remark}

\begin{lemma}[Gap and regularity]\label{ptr:gap}
Every $\Lambda^{(i,n)}_\tau$ consists of $N_{i,n}$ distinct numbers in $\tfrac{N_{i,n}-1}2+\Z$, two
numbers with different signs differ by at least $\tfrac12(\Delta-t)-1\ge\tfrac12r^8$, and
$N_{i,n}\le N_m$. Hence the conclusions of Theorem~\ref{pco:FJ}, Lemma~\ref{pco:decay},
Theorem~\ref{par:A1} and Corollary~\ref{par:A2} hold for the types $e_{i,n}$.
\end{lemma}
\begin{proof}
The numbers are those of the stem of $\varsigma$ and of $n$ strings, all shifted by the same
$-\tfrac12[n\text{ odd}]$. The parity: for $n=0$ they lie in $\tfrac{r-1}2+\Z$ by
\eqref{pz:eq-param} ($j=0$); every string adds $t$ (odd) to the rank and the shift by $\tfrac12$
when $n$ changes parity restores $\tfrac{N-1}2+\Z$. Distinctness and the gap: by \eqref{ptr:eq-E}
negative signs occur only on stem numbers (the small numbers when $i$ is odd, and some of the
numbers $-\kappa-i'\Delta$, shifted); a small number is at distance $\ge\kappa-k$ from every
non-small number, and a number $-\kappa-i'\Delta$ is at distance
$\ge\min(\tfrac12\Delta-\tfrac{t-1}2-\tfrac12,\ \Delta)$ from the other non-small numbers; both
bounds are at least $\tfrac12(\Delta-t)-1$ as in \S\ref{pz:ss-types}; distinct strings are $\Delta$ apart, and within
a string the numbers are consecutive. As $N_{i,n}\le r+mt=N_m$, the numerical hypotheses of the four
cited statements (gap $\ge N^2+1$, resp.\ $\ge B+N^2+1$, $\ge5N^3$) hold by the inequality
$\tfrac12r^8\ge500(N_j+N_{j+1})^3+10N_j^5+x^{25}$ of \S\ref{par:sec}; these statements use the
parameter only through its regularity, the gap and the rank (Remark~\ref{par:FJ}, proof of
Lemma~\ref{pco:decay}).
\end{proof}

\subsection{Joint harmonics and the edge property at the nodes}\label{ptr:ss-edge}

For a pair $(S,T)$ (ascending or descending) with row parameter $\lambda$ and column parameter
$\lambda'$ at every $\tau$, let $e_S=d\,\overline{c}$ be the type on $S$ whose conjugate has parameter
$\lambda$ (so $e_S=e_{i,n}$ for an ascending pair and $e_S=\overline{e_{i+1,n}}$ for a descending one)
and $e_T$ the type with parameter $\lambda'$ ($e_{i,n+1}$, resp.\ $\overline{e_{i,n+1}}$). Put
$\mathsf P=\int e_S(x)\,\omega_\infty(x^{-1},1)\,dx$ and $\mathsf P'=\int e_T(x)\,\omega_\infty(1,x)\,dx$
on the Fock space, as in Theorem~\ref{par:A3}.

\begin{theorem}[Edges]\label{ptr:edge}
Assume \textup{(V)} and \textup{(D1)--(D3)} of \S\ref{par:ss-rigid}. For every ascending and every
descending pair of the triangle:
\begin{enumerate}
\item[(a)] a joint harmonic $P_{S,T}=\bigotimes_\tau\Phi_\tau$ exists, is unique up to a scalar, and
no monomial of it contains two variables of one column from rows of opposite signs, or of one row
from columns of opposite signs;
\item[(b)] $\mathsf P P_{S,T}=\mathsf P'P_{S,T}$;
\item[(c)] every form of type $e_S$ or $\overline{e_S}$, $e_T$ or $\overline{e_T}$ is a cusp form,
rapidly decreasing on Siegel sets, and the forms of a given level form a finite-dimensional space;
\item[(d)] for every finite datum $f$ and every form $y$ of type $e_S$, the lift
$h\mapsto\int_{[\auU(S)]}\theta(\omega(g,h)(P_{S,T}\otimes f))\,y(g)\,dg$ converges absolutely and is a
form of type $e_T$; for every form $z$ of type $\overline{e_T}$ the transposed lift
$g\mapsto\int_{[\auU(T)]}\theta(\omega(g,h)(P_{S,T}\otimes f))\,z(h)\,dh$ converges absolutely and is a
form of type $\overline{e_S}$.
\end{enumerate}
Here $\omega$ is $\omega_{\alpha,\beta}$ with the characters of $a_{i,n}$, resp.\ of
Definition~\ref{ptr:desc-char}. In particular \textup{(A1)} and \textup{(A2)} of Lemma~\ref{pc:arch}
hold for the pairs $(\pcG_{i,0},\pcG_{i,1})$, and the maps $a_{i,n}$, $d_{i,n}$ of
Theorem~\ref{pc:tri} (the latter formed with any finite datum) are well defined between the spaces
$\pcF_{i,n}$.
\end{theorem}
\begin{proof}
(a) By Proposition~\ref{ptr:match} the parameters match at every $\tau$; Propositions~\ref{par:exist}
and~\ref{par:unique} are stated for an arbitrary matching pair with distinct entries, and the
passage to the tensor product over $\tau$ is as in the proof of Theorem~\ref{par:A3}.

(b) The proof of Theorem~\ref{par:A3} uses: Lemma~\ref{par:vac} (here Lemma~\ref{ptr:const}), the
matching (Proposition~\ref{ptr:match}), the conclusion of Theorem~\ref{par:A1} for the two types
(Lemma~\ref{ptr:gap}), Lemma~\ref{par:rigid} (stated for any regular $\lambda$ with integrable
$e_\lambda$), Lemma~\ref{par:casimir} (stated for any matching pair), and (D1)--(D3). It applies
verbatim. (For a descending pair the row and column representations are $\pi_{\Lambda^{(i+1,n)}}$
and $\pi_{-\Lambda^{(i,n+1)}}$, and $\mathsf P$, $\mathsf P'$ are the projectors $\sigma(e_\lambda)$,
$\sigma'(e_{\lambda'})$ exactly as there.)

(c) Proposition~\ref{pz:types}(a) for $G_j$ uses only (A1), (A2) for the type $e_j$, i.e.\ the
conclusions of Theorem~\ref{par:A1} and Corollary~\ref{par:A2}; these hold for every $e_{i,n}$ by
Lemma~\ref{ptr:gap}. The groups $\auU(\pcG_{i,n})$ with $n=m-i$ may be compact (no negative sign);
then the argument is the one given there for $j=m$.

(d) The proof of Proposition~\ref{pz:types}(b), (c) uses: the growth Lemma~\ref{pz:growth} (which is
stated for an arbitrary pair of spaces and characters: the norm $\|\cdot\|_s$ and the bound
$\|\omega_\infty(g,h)\|_s\le(2D\|g\|\|h\|)^{2s}$ depend only on $D=fNN'$), reduction theory,
statement (c) above, the conclusion of Theorem~\ref{par:A1}, and the identity (b). It applies
verbatim, with $\Phi_j$ replaced by $P_{S,T}$, and gives the two assertions. The last sentence:
$a_{i,n}$ is the lift of (d) for the ascending pair with $f=\delta_\iota$ at $\pcm_\vartheta$ and
the vacuum elsewhere, so it maps $\pcF_{i,n}=\pcF(\pcG_{i,n},e_{i,n})$ into forms of type
$e_{i,n+1}$; and the descent along a descending pair is a transposed lift of (d) with
$e_S=\overline{e_{i+1,n}}$, $e_T=\overline{e_{i,n+1}}$, so it maps forms of type
$\overline{e_T}=e_{i,n+1}$ to forms of type $\overline{e_S}=e_{i+1,n}$. Rapid decay in (A1) of
Lemma~\ref{pc:arch} is (c), and in the form used below it reads: on a Siegel set
$\mathfrak S_\infty\mathsf cK'$ every such form is bounded by $C_N\|g_\infty\|^{-N}$ for every $N$.
(The steps of the proofs of Theorem~\ref{par:A3} and Proposition~\ref{pz:types} referred to in (a)--(d)
are gone through one by one, for an arbitrary pair of standard spaces satisfying the block rule, in
Steps 0--4 of the proof of Theorem~\ref{pgv:G3}; for the present pairs the excess is $t$ and the
properties used there are Proposition~\ref{ptr:match} and Lemma~\ref{ptr:gap}.)
\end{proof}

\begin{remark}[What the previous text assumed]\label{ptr:rem-prev}
In \S\ref{pc:ss-triangle} the parameter at the node $(i,n)$ is defined directly, and
Lemma~\ref{pc:arch} asserts the edge property for $(\pcG_{i,0},\pcG_{i,1})$ only. What the squares
need (step (iii) of \S\ref{pc:ss-exchange}; Theorem~\ref{ptr:square} below) is the edge property
for all four pairs of every square: for the square at $(i,n)$, with $U_1=\pcG_{i,n+1}$,
$U_2=\pcG_{i+1,n+1}$, $A=\pcG_{i,n+2}$, $A'=\pcG_{i+1,n}$, the pairs $(U_1,A)$ and $(A',U_2)$ are
ascending (at $(i,n+1)$ and $(i+1,n)$) and $(U_2,A)$, $(A',U_1)$ are descending (at $(i,n+1)$ and
$(i,n)$). All are covered by Theorem~\ref{ptr:edge}. The sign identity used in
\S\ref{pc:ss-exchange} for the numbers of negative signs is a consequence of \eqref{ptr:eq-E} and
was checked there.
\end{remark}

\subsection{The characters of a square}\label{ptr:ss-chars}

Fix a square at $(i,n)$, $i+n+2\le m$: $U_1=\pcG_{i,n+1}$, $U_2=\pcG_{i+1,n+1}$, $A=\pcG_{i,n+2}$,
$A'=\pcG_{i+1,n}$, $W=U_1\oplus U_2$ of rank $N=2M$, $M=r+(n+1)t$. For an $m$-fold twist $\vartheta$
let $\nu=(\nu_{n+1}\vartheta^{\,n+1})^{-1}$, of weight $(n+1)\Delta$ (\S\ref{pc:ss-exchange}), and
$\gamma=\nu^{-1}\upsilon^{M}$, so that $\ptrwt_\tau\gamma=M-(n+1)\Delta$.

In step (i) of \S\ref{pc:ss-exchange} the kernel of $(W,A)$ restricted to $\auU(U_1)\times\auU(U_2)$
is the product of the kernels of $(U_1,A)$ and $(U_2,A)$, and the kernel of $(A',W)$ is the product
of those of $(A',U_1)$ and $(A',U_2)$; for this the target characters on $\auU(A)$ of the first two
must multiply to $\gamma$, and the source characters on $\auU(A')$ of the last two must multiply to
$\gamma$ (Theorem~\ref{pcx:X}: $\auU(V)$ acts through $\gamma(\det h)$ on the left side and $\auU(V')$
through $\gamma(\det h')^{-1}$ on the right side), the Weil representations being multiplicative for
orthogonal sums. The pair $(U_1,A)$ is ascending at $(i,n+1)$, with target character
$\mu_{i+n+2}\upsilon_M$; $(U_2,A)$ is descending at $(i,n+1)$; $(A',U_2)$ is ascending at $(i+1,n)$; $(A',U_1)$ is descending at $(i,n)$. In the conventions of
\S\ref{pc:ss-data} the target character of $a_{i,n+1}$ and the source character of $a_{i+1,n}$ are
both $\mu_{i+n+2}\upsilon_{r+(n+1)t}=\mu_{i+n+2}\upsilon_M$.

\begin{lemma}[The descents' characters are those of the squares]\label{ptr:chars}
Put $\beta^\downarrow_{i,n+1}=\gamma\,(\mu_{i+n+2}\upsilon_M)^{-1}$ and
$\alpha^\downarrow_{i,n}=\gamma\,(\mu_{i+n+2}\upsilon_M)^{-1}$ (for the twist $\vartheta$ of the
square). Their weights are those of Definition~\ref{ptr:desc-char}. Consequently the characters
required by the exchange identity for the four kernels of a square make the four pairs edges (Theorem~\ref{ptr:edge}).
\end{lemma}
\begin{proof}
$\ptrwt(\mu_{i+n+2}\upsilon_M)=-\kappa_{i+n+2}+\lfloor M/2\rfloor$, so both characters have weight
$M-(n+1)\Delta+\kappa_{i+n+2}-\lfloor M/2\rfloor=\kappa_{i+1}+\lceil M/2\rceil$, because
$\kappa_{i+n+2}-(n+1)\Delta=\kappa_{i+1}$. Definition~\ref{ptr:desc-char} asks
$\ptrwt\beta^\downarrow_{i,n+1}=\kappa_{i+1}+\lceil\tfrac{r+(n+1)t}2\rceil$ and
$\ptrwt\alpha^\downarrow_{i,n}=\kappa_{i+1}+\lceil\tfrac{r+(n+1)t}2\rceil$: the same. The last
sentence is Proposition~\ref{ptr:match}(c).
\end{proof}
The characters of $d_{i,n}$ are fixed by the squares, consistently, and $\beta^\downarrow_{i,0}$
agrees with the character given by $d_{i,0}=a_{i,0}^*$ (Remark~\ref{ptr:rem-desc},
Proposition~\ref{ptv:chars}). The square at $(i,n)$ also fixes, by (E2) at $U_1$ and $U_2$,
$\beta^\downarrow_{i,n}=\gamma\,(\text{source character of }a_{i,n+1})^{-1}$ and
$\alpha^\downarrow_{i,n+1}=\gamma\,(\text{target character of }a_{i+1,n})^{-1}$ (weights checked in
Corollary~\ref{ptr:tri-b}); in particular $\alpha^\downarrow_{i,m-i-1}$ is fixed by the square at
$(i,m-i-2)$ when $m-i\ge2$, and the characters of $d_{m-1,0}$ are those of (W7)
(Proposition~\ref{ptv:chars}). They depend on $\vartheta$ through $\gamma$; the tame
part of $\gamma$ at the primes of $\pcm_\vartheta$ is the character $\chi$ of the square, as required
in step (i).

\subsection{Convergence in the squares}\label{ptr:ss-conv}

\begin{proposition}[Step (ii)]\label{ptr:conv}
Let a square be as in \S\ref{ptr:ss-chars}, $\pcm=\pcm_\vartheta$, $y_1$ a form of type
$e_{i,n+1}$ on $U_1$ and $y_2$ a form of type $\overline{e_{i+1,n+1}}$ on $U_2$, $\xi\in\pcMat_{a\times
N}(\pcR_\pcm)$ of rank $a=\dim A$ at every prime of $\pcm$ and $\eta\in\pcMat_{N\times b}(\pcR_\pcm)$ of rank
$b=\dim A'$, and let $\chi_2$ be a unitary character of $\auU(1)$. There are $M_0,C$ such that
for every Fock vector $x$ of $\auHom(W,A)\otimes\R$ and $v$ of $\auHom(A',W)\otimes\R$ with finite
moments of order $M_0$:
\begin{enumerate}
\item[(a)] the function $(g_1,g_2,h)\mapsto y_1(g_1)y_2(g_2)\chi_2(\det g_2)\,\theta^{W,A}(g_1\oplus g_2,h;
[x,\delta_\xi])$ is integrable on $[\auU(U_1)]\times[\auU(U_2)]\times[\auU(A)]$, with integral of absolute
value at most $C\|x\|_{M_0}$;
\item[(b)] the function $(g_1,g_2,h')\mapsto y_1(g_1)y_2(g_2)\chi_2(\det g_2)\,\theta^{A',W}(h',g_1\oplus
g_2;[v,\delta_\eta])$ is integrable on $[\auU(U_1)]\times[\auU(U_2)]\times[\auU(A')]$, with the same bound
in $\|v\|_{M_0}$.
\end{enumerate}
(The characters of the kernels are any; $\|\cdot\|_M$ is the moment of Proposition~\ref{pcx:decay}.)
In particular all iterated integrals in steps (i), (iii) of \S\ref{pc:ss-exchange} converge
absolutely and may be computed in any order.
\end{proposition}
\begin{proof}
\emph{Heights of block matrices.} With the heights of Definition~\ref{pcx:def-ht}, for $g_l\in
\auU(U_l)(\auA)$ the lattice at $g=g_1\oplus g_2$ is $\Lambda_{g_1}\oplus\Lambda_{g_2}$ (the frame of
$W$ is the union of the frames) and $|(\gamma_1,\gamma_2)|_g^2=|\gamma_1|_{g_1}^2+|\gamma_2|_{g_2}^2$; so
$\lambda_1(g)=\min(\lambda_1(g_1),\lambda_1(g_2))$ and $\ptrht(g)=\max(\ptrht g_1,\ptrht g_2)\le
\ptrht(g_1)\ptrht(g_2)$.

\emph{Heights on Siegel sets.} Let $\mathfrak S_\infty\mathsf cK'$ be as in the proof of
Proposition~\ref{pz:types}(b) for $U_l$. For $g=g_\infty c k'$ in it and $0\ne\gamma\in\Lambda_g$,
$\gamma$ lies in the lattice $(ck')L$ of $S(E)$, which belongs to a finite set of lattices as $c$
runs over the finite set $\mathsf c$ and $k'$ over $K'$; so $\|\gamma\|\ge c_0$ for a fixed euclidean
norm and some $0<c_0\le1$. Since $\|\gamma\|=\|g_\infty g_\infty^{-1}\gamma\|\le\|g_\infty\|\,
\|g_\infty^{-1}\gamma\|$ (operator norm; the norm of $g_\infty=(g_\tau)_\tau$ is the maximum over the
archimedean factors), $|\gamma|_g=\|g_\infty^{-1}\gamma\|\ge c_0\|g_\infty\|^{-1}$. Hence
$\ptrht(g)=\max(1,\lambda_1(g)^{-1})\le c_0^{-1}\max(1,\|g_\infty\|)$ on the Siegel set. By
Theorem~\ref{ptr:edge}(c), $|y_l(g)|\le C_{N'}\|g_\infty\|^{-N'}$ there for every $N'\ge0$; using
$N'=0$ where $\|g_\infty\|\le1$ and $N'=B$ where $\|g_\infty\|\ge1$,
$|y_l(g)|\ptrht(g)^B\le c_0^{-B}\max(C_0,C_B)$ on the Siegel set, so it is bounded there, hence on $\auU(U_l)(\auA)$ (left
$\auU(U_l)(E^+)$-invariance), for every $B$. As $[\auU(U_l)]$ has finite volume,
$|y_l|\ptrht^B$ is integrable.

(a) By Proposition~\ref{pcx:decay}(a) for $(S_1,S_2)=(W,A)$ and $P$ the primes of $\pcm$ (the point
datum $\delta_\xi$ is supported on a matrix of rank $a$, i.e.\ surjective, at every place of $P$),
for $N_0$ large there are $B,M_0,C$ with $|\theta^{W,A}(g,h;[x,\delta_\xi])|\le C\|x\|_{M_0}
\ptrht(g)^B\ptrht(h)^{-N_0}$. With the two preceding paragraphs the integrand is bounded by
$C\|x\|_{M_0}\,|y_1|\ptrht^B(g_1)\cdot|y_2|\ptrht^B(g_2)\cdot\ptrht(h)^{-N_0}$, a product of integrable
functions by (R1) of \S\ref{pcx:ss-conv}. (b) The same with Proposition~\ref{pcx:decay}(b) for
$(S_1,S_2)=(A',W)$: $\delta_\eta$ is supported on a matrix of rank $b$, i.e.\ injective. The last
sentence follows by Fubini's theorem; the path integrals of steps (i), (iii) (lift to $A$, then
pairing on $[\auU(A)]$; transposed lift to $A'$, then lift to $U_2$) are iterated integrals of these
functions, by the factorisation of the kernels in step (i).
\end{proof}

\subsection{Archimedean multiplicity one}\label{ptr:ss-mult}

Let $(S,T)$ be a pair with matching parameters as in Theorem~\ref{ptr:edge} (or any pair for which
the conclusions (a)--(d) of that theorem hold), $P=P_{S,T}$ its joint harmonic, $\mathsf P$,
$\mathsf P'$ as there, and $\pzF_s$ the Banach space of \S\ref{pz:ss-types}. We use the following
local statement.

\begin{enumerate}
\item[(H)] For every polynomial vector $Q$ of the Fock space of $\auHom(S,T)\otimes\R$ there is a
continuous compactly supported function $\alpha_Q$ on $\auU(T)_\infty$ with
$\mathsf PQ=\int\alpha_Q(x)\,\omega_\infty(1,x)\,\mathsf PP\,dx$ (Bochner integral in $\pzF_s$).
\end{enumerate}
\noindent (H) holds for every ascending and descending pair of the triangle: this is Theorem~\ref{ptq:H}
(proved by Howe's argument, Proposition~\ref{ptq:howe}, modulo (V) and the published results
(D1)--(D5)).

\begin{lemma}\label{ptr:conv-lift}
Let $(S,T)$ be a pair for which \textup{(H)} holds (for instance any pair of the triangle, by
Theorem~\ref{ptq:H}). For every polynomial $Q$, every finite datum $f$ and every form $y$ of type
$e_S$, the lift $F_Q$ of $y$ with $Q\otimes f$ (absolutely convergent as in
Theorem~\ref{ptr:edge}(d)) is $F_Q(h)=\int\alpha_Q(x)F_P(hx)\,dx$, with $\alpha_Q$ independent of
$f$, $y$, the characters and the spaces of the same signatures.
\end{lemma}
\begin{proof}
In the proof of Proposition~\ref{pz:types}(b), $F_\Psi(h)=L_h(\Psi)$ with $L_h$ continuous on
$\pzF_s$, $L_h(\Psi)=L_h(\mathsf P\Psi)$ and $L_h(\omega_\infty(1,x)\Psi)=L_{hx}(\Psi)$. Apply $L_h$ to
(H): $F_Q(h)=L_h(\mathsf PQ)=\int\alpha_Q(x)L_{hx}(\mathsf PP)\,dx=\int\alpha_Q(x)F_P(hx)\,dx$.
\end{proof}

\begin{lemma}\label{ptr:avg}
Let $e$ be an integrable type on a space $T$ with parameter of the class of Lemma~\ref{ptr:gap},
$\alpha_1,\alpha_2\in C_c(\auU(T)_\infty)$. There is $r=r(\alpha_1,\alpha_2)\in\C$ such that for
every form $F_1$ of type $e$ and every form $F_2$ of type $\overline e$ (both rapidly decreasing),
$\int_{[\auU(T)]}(\alpha_1*F_1)(\alpha_2*F_2)=r\int_{[\auU(T)]}F_1F_2$, where $(\alpha*F)(h)=\int\alpha(x)
F(hx)dx$.
\end{lemma}
\begin{proof}
$\overline{F_2}$ is of type $e$. By \S\ref{pc:ss-data} (consequence of Schur orthogonality),
$F_1=\Phi_1(v_0)$, $\overline{F_2}=\Phi_2(v_0)$ for intertwiners $\Phi_l\colon\pi\to L^2$, and
$\langle\Phi_2(u),\Phi_1(v)\rangle=\langle\overline{F_2},F_1\rangle\langle u,v\rangle$. Hence
$B(x,x'):=\int F_1(hx)F_2(hx')dh=\langle\rho(x')\overline{F_2},\rho(x)F_1\rangle=\langle\pi(x')v_0,
\pi(x)v_0\rangle\int F_1F_2$, and by Fubini (the integrand is bounded by
$|\alpha_1(x)\alpha_2(x')|\,|F_1(hx)F_2(hx')|$, integrable since the forms are bounded and in $L^2$)
the claim holds with $r=\iint\alpha_1(x)\alpha_2(x')\langle\pi(x')v_0,\pi(x)v_0\rangle dx\,dx'$.
\end{proof}

\begin{lemma}[Density of polynomials for moment norms]\label{ptr:dense}
In the realisation of Lemma~\ref{pcx:superpos}, every Fock vector $x$ with $\|x\|_{M'}<\infty$ for
all $M'$ is, for each $M$, a limit in $\|\cdot\|_M$ of polynomial vectors. Here the polynomial vectors are the images of
those of \eqref{pk:eq-fock} under the isomorphism of Lemma~\ref{pcx:superpos}; it carries the vacuum
to the vacuum and intertwines the Heisenberg action, so they are the vectors $Pe$ of the realisation.
\end{lemma}
\begin{proof}
By Lemma~\ref{pcx:superpos}, $x=\int x(z)\rho(z)e\,dz$; since $\|\rho(z)e\|_M\le C(1+|z|)^M$ and
$z\mapsto\rho(z)e$ is continuous for $\|\cdot\|_M$, the truncations to $|z|\le R$ converge to $x$
(since $\|\rho(z)e\|_M\le C(1+|z|)^M$, the tail is at most $C\int_{|z|>R}(1+|z|)^M|x(z)|\,dz$, which
tends to $0$ because $\|x\|_M<\infty$, Lemma~\ref{pcx:superpos}(b)), and each truncation is a limit of finite
combinations of vectors $\rho(z)e$. Each $\rho(z)e$ is $\exp(\ell_z)\,e$ with $\ell_z$ affine in the
coordinates in which the polynomial vectors are the $Pe$ (the vacuum is a Gaussian $e$); the Taylor
truncations $\sum_{k\le K}\ell_z^k/k!\cdot e$ are polynomial vectors, bounded by
$e^{|\ell_z|}|e|$, and $(1+|w|)^Me^{|\ell_z(w)|}|e(w)|$ is integrable; so they converge in
$\|\cdot\|_M$ by dominated convergence.
\end{proof}

\subsection{The square law}\label{ptr:ss-square}

We state the passage from the exchange identity to one square in a form that serves both the
triangle (Theorem~\ref{pc:tri}(b)) and the grids (Proposition~\ref{pgr:square}).

\begin{definition}[Edge square]\label{ptr:def-square}
An \emph{edge square} of excess $d\ge0$ and weight $\mu$ consists of spaces $U_1,U_2,A,A'$ of
ranks $n_1,n_2,a,b$ with $a=b+2d$, $N=n_1+n_2=a+b$, $W=U_1\oplus U_2$, such that, with $q$, $q_A$,
$q_{A'}$ the numbers of negative signs of $W,A,A'$: $q=q_A+q_{A'}$, $2q+2\le N\le2\mu_\tau+2q$; of
types $e_1,e_2,e,e'$ on $U_1,U_2,A,A'$; and of a family $I$ of data $(\pcm,\nu,\chi,\text{characters})$,
$\pcm\ne(1)$ a modulus, $\nu$ a tame character for $\pcm$ of weight $\mu$ with tame character $\chi$,
$\gamma=\nu^{-1}\upsilon^{N/2}$, and characters for the four kernels such that:
\begin{enumerate}
\item[(E1)] (edges) the four pairs are edges in the sense of Theorem~\ref{ptr:edge}(a)--(d) with
polynomials $P_1,P_2,P'_1,P'_2$: $(U_1,A)$ from $e_1$ to $e$, $(U_2,A)$ from $\overline{e_2}$ to
$\overline e$, $(A',U_1)$ from $\overline{e'}$ to $\overline{e_1}$, $(A',U_2)$ from $e'$ to $e_2$;
the polynomials and types do not depend on the member of $I$, only the characters do, with fixed
weights; and \textup{(H)} holds for $(U_1,A)$ and $(U_2,A)$;
\item[(E2)] (characters) the target characters on $A$ of the kernels of $(U_1,A)$, $(U_2,A)$
multiply to $\gamma$; the source characters on $A'$ of $(A',U_1)$, $(A',U_2)$ multiply to $\gamma$;
and for $l=1,2$ the source character on $U_l$ of $(U_l,A)$ times the target character on $U_l$ of
$(A',U_l)$ is $\gamma$.
\end{enumerate}
For $\xi\in\pcMat_{a\times N}(\pcR_\pcm)$ with blocks $\xi_l$ (columns of $U_l$) and
$\eta\in\pcMat_{N\times b}(\pcR_\pcm)$ with blocks $\eta_l$ (rows of $U_l$), forms $y_1$ of type $e_1$
and $y_2$ of type $\overline{e_2}$, put
\[
L(\xi)=\int_{[\auU(A)]}\Lambda^{(1)}_{\xi_1}y_1\cdot\Lambda^{(2)}_{\xi_2}y_2,\qquad
L'(\eta)=\int_{[\auU(U_2)]}y_2\cdot\Lambda'^{(2)}_{\eta_2}T'^{(1)}_{\eta_1}y_1,
\]
where $\Lambda^{(l)}_\zeta$ is the lift along $(U_l,A)$ with $P_l$ and the point datum $\delta_\zeta$,
$T'^{(1)}_\zeta$ the transposed lift along $(A',U_1)$ and $\Lambda'^{(2)}_\zeta$ the lift along
$(A',U_2)$. (By (E1) and Theorem~\ref{ptr:edge}(d) these are forms of the right types, and the
integrals converge.)
\end{definition}

\begin{theorem}[Square law]\label{ptr:square}
For an edge square there are $\kappa\in\C$ and $C>0$, depending only on the archimedean data, and
for each member of $I$ a number $c$ with $|c|^2=C\,L_*(\pcm,\nu,d)\ne0$, such that for all forms
$y_1,y_2$ the functions $L$ (restricted to the $\xi$ of rank $a$) and $L'$ (restricted to the $\eta$
of rank $b$) satisfy
\begin{equation}\label{ptr:eq-square}
\sum_\eta\Psi(\eta)L'(\eta)=\frac{\kappa}{c}\sum_\xi\Phi(\xi)L(\xi)
\end{equation}
for all $\Phi,\Psi$ matched for $\chi$; moreover $L(h\xi)=\chi(\det h)\,L(\xi)$ and
$L'(\eta h')=\chi(\det h')\,L'(\eta)$ for $h\in\auGL_a(\pcR_\pcm)$, $h'\in\auGL_b(\pcR_\pcm)$. Thus
$L,L'$ satisfy the square law of \S\ref{pc:ss-exchange} with constant $\kappa/c$ (matched functions
vanish off full-rank matrices).
\end{theorem}
\begin{proof}
\emph{Transformation laws.} For $h\in\auGL_a(\pcR_\pcm)$ choose $k\in K_A$ with reduction $h$. By
Lemma~\ref{pc:level}, $\omega(1,k^{-1})$ sends the datum $[P_l,\delta_{h\xi_l}]$ to the target character
of $(U_l,A)$ at $\det k$, inverted, times $[P_l,\delta_{\xi_l}]$; translating the lifted forms by
$k^{-1}$ (the pairing on $[\auU(A)]$ is invariant) gives $L(h\xi)=c_1(\det k)^{-1}c_2(\det k)^{-1}L(\xi)$
with $c_1c_2=\gamma$ the product of the two target characters (E2), and $\gamma^{-1}$ restricted to
$K_A$ at the primes of $\pcm$ is $\chi\circ\det$ by the definition of $\chi$ (the factor
$\upsilon^{N/2}$ is unramified). The law of $L'$ is the same with $K_{A'}$ and the source characters.

\emph{Kernels multiply.} $\psrho$, $\theta$ and the Weil representations of
$\auHom(W,A)=\auHom(U_1,A)\oplus\auHom(U_2,A)$ are tensor products, and
$[x_1x_2,\delta_\xi]=[x_1,\delta_{\xi_1}]\otimes[x_2,\delta_{\xi_2}]$ (finite data multiply blockwise;
the vacuum of an orthogonal sum of lattices is the tensor product). So for Fock vectors $x$ of
$\auHom(W,A)\otimes\R$ the function
\[
L_x(\xi)=\int_{[\auU(U_1)]\times[\auU(U_2)]\times[\auU(A)]}y_1(g_1)\,y_2(g_2)\,\lambda(\det g_2)\,
\theta^{W,A}_{\alpha,\gamma}(g_1\oplus g_2,h;[x,\delta_\xi])
\]
($\alpha=\alpha_1$, where $\alpha_l$ is the source character of $(U_l,A)$, and
$\lambda=\alpha_1\alpha_2^{-1}$, unitary; by (E2) also $\lambda=\alpha'_2\alpha'^{-1}_1$ for the target
characters $\alpha'_l=\gamma\alpha_l^{-1}$ of $(A',U_l)$, so that $y_1y_2\lambda(\det g_2)
\theta^{W,A}_{\alpha_1,\gamma}$ and $y_1y_2\lambda(\det g_2)\theta^{A',W}_{\gamma,\alpha'_1}$ are the
products of the two kernels on each side) equals $L(\xi)$ for $x=P_1P_2$, by Fubini; the triple integral converges
absolutely by Proposition~\ref{ptr:conv}(a) (whose proof uses only (E1) and the full rank of $\xi$),
and $|L_x(\xi)|\le C_\xi\|x\|_{M_0}$. In the same way $L'(\eta)$ is the triple integral
$L'_{v}(\eta)$ over $[\auU(U_1)]\times[\auU(U_2)]\times[\auU(A')]$ of $y_1y_2\lambda(\det g_2)
\theta^{A',W}_{\gamma,\alpha'}(h',g_1\oplus g_2;[v,\delta_\eta])$ with $v=P'_1P'_2$ and
$\alpha'=\gamma\alpha^{-1}$ by (E2), absolutely convergent by Proposition~\ref{ptr:conv}(b) (here the
lift along $(A',U_2)$ is moved to the other side of the pairing on $U_2$ by Fubini).

\emph{Multiplicity one.} Fix $y_1,y_2$. For a polynomial $Q=Q_1Q_2$ with $Q_l$ on
$\auHom(U_l,A)\otimes\R$, Lemma~\ref{ptr:conv-lift} and Lemma~\ref{ptr:avg} (with $F_1=\Lambda^{(1)}
_{\xi_1}y_1$ of type $e$ and $F_2=\Lambda^{(2)}_{\xi_2}y_2$ of type $\overline e$) give
$L_{Q}(\xi)=r(\alpha_{Q_1},\alpha_{Q_2})L(\xi)$, with $r$ independent of $\xi$, $y_1,y_2$ and the
member of $I$. Every polynomial on $\auHom(W,A)\otimes\R$ is a sum of such products; so for every
polynomial $Q$ there is $\kappa(Q)$ with $L_Q=\kappa(Q)L$ on all cases (all $\xi$ of rank $a$, all
$y_1,y_2$, all members of $I$). The map $x\mapsto(L_x(\text{case}))_{\text{cases}}$ is, case by case,
continuous for $\|\cdot\|_{M_0}$, and the set of functions on the cases that are multiples of $L$ is a
line, closed for pointwise convergence. By Lemma~\ref{ptr:dense}, for every $x$ with all moments
finite there is $\kappa(x)$ with $L_x=\kappa(x)L$.

\emph{Exchange.} Apply Theorem~\ref{pcx:X} to $(W,V,V')=(W,A,A')$ with the vector $v=P'_1P'_2$ (a
polynomial, so all its moments are finite): it gives $C>0$ and $x$ with finite moments, and for each
member of $I$ a number $c$ with $|c|^2=CL_*(\pcm,\nu,d)\ne0$ (Corollary~\ref{pcx:cne0}; the excess called $e$ there is $d$ here), such that
for matched $\Phi,\Psi$ and all $g$,
$\int_{[\auU(A)]}\theta^{W,A}_{\alpha,\gamma}(g,h;[x,\Phi])dh=c\int_{[\auU(A')]}\theta^{A',W}_{\gamma,
\alpha'}(h',g;[v,\Psi])dh'$. Write $[x,\Phi]=\sum_\xi\Phi(\xi)[x,\delta_\xi]$ (a finite sum, $\Phi$
vanishing off rank $a$), likewise for $\Psi$; multiply by $y_1(g_1)y_2(g_2)\lambda(\det g_2)$, put
$g=g_1\oplus g_2$, and integrate over $[\auU(U_1)]\times[\auU(U_2)]$; by the absolute convergence just
proved the integrals may be exchanged with the finite sums, and
$\sum_\xi\Phi(\xi)L_x(\xi)=c\sum_\eta\Psi(\eta)L'_v(\eta)$. With $L_x=\kappa L$, $\kappa=\kappa(x)$,
and $L'_v=L'$ this is \eqref{ptr:eq-square}.
\end{proof}

\begin{corollary}[Theorem~\ref{pc:tri}(b)]\label{ptr:tri-b}
Assume \textup{(V)} (Theorem~\ref{pco:V}) and \textup{(D1)--(D5)}. Then the squares of the triangle at $(i,n)$, $i+n+2\le m$, with the family of $m$-fold
twists, are edge squares of excess $t$ and weight $(n+1)\Delta$, and Theorem~\ref{pc:tri}(b) holds,
with $\kappa^\sharp_{i,n}=\kappa$, $c_{i,n}(\vartheta)=c$.
\end{corollary}
\begin{proof}
The rank and sign conditions are verified in \S\ref{pc:ss-exchange}; (E1) is
Theorem~\ref{ptr:edge} with Remark~\ref{ptr:rem-prev}; (E2) holds for the weights by
Lemma~\ref{ptr:chars} and the analogous computation at $U_1$, $U_2$: the weights of the source
character of $a_{i,n+1}$ and of $\beta^\downarrow_{i,n}$ add up to
$-\kappa_{i+n+2}+\lfloor\tfrac{M+t}2\rfloor+\kappa_{i+1}+\lceil\tfrac{M-t}2\rceil=M-(n+1)\Delta=\ptrwt
\gamma$, and those of $\alpha^\downarrow_{i,n+1}$ and of the target character of $a_{i+1,n}$ add up to
$\kappa_{i+1}+\lceil\tfrac{M+t}2\rceil-\kappa_{i+n+2}+\lfloor\tfrac{M-t}2\rfloor=\ptrwt\gamma$
($M\pm t$ have the same parity). The descent characters are then fixed by (E2), consistently in all
squares (Proposition~\ref{ptv:chars}); (H) for the ascending and descending pairs is
Theorem~\ref{ptq:H}. Theorem~\ref{ptr:square} gives the square law for $L,L'$ with constant
$\kappa/c$, for the splitting of Definition~\ref{pcx:def-matched} (the first $b$ indices of the frame of
$W=U_1\oplus U_2$). The splitting $(\mathsf B,\mathsf A)=(\mathsf N_0,\mathsf N_1\sqcup\mathsf T')$ of
Lemma~\ref{pc:next} is obtained from it by the permutation that moves the $r+nt$ indices of $\mathsf N_0$
in front of the $r+(n+1)t$ indices of $\mathsf N_1$, of sign $(-1)^{(r+nt)(r+(n+1)t)}=1$ because $r+nt$ and
$r+(n+1)t$ have opposite parities ($t$ is odd, Lemma~\ref{ptr:cons}(c)); so by Lemma~\ref{ptw:split} the law holds
for this splitting with the same constant, and step (iv) of \S\ref{pc:ss-exchange}
(Lemma~\ref{pc:next}) turns it into (b).
\end{proof}

\begin{remark}[The grid squares]\label{ptr:rem-grid}
Proposition~\ref{pgr:square} is Theorem~\ref{ptr:square} for the square $(V',V;W_1,W_2)$ of a grid,
$U_l=W_l$, $A=V$, $A'=V'$, with excess $d=t$ and $\nu=(\mu^0_{a+1}\mu^0_{b+1}\vartheta^n)^{-1}$; its
functions $L_{f,l}$, $L'_{f,l}$ are obtained from $L,L'$ above by replacing the pairing with $y_2$ by
the linear functional $l$, which changes nothing in the proof (Fubini is applied with
$l$ applied after the integrals, $l$ being a linear functional on a finite-dimensional space of forms).
The rank, sign and weight conditions are verified in Lemma~\ref{pgv:esq}; (E1) requires the edge
property along the four edges of a grid square; it is checked in Proposition~\ref{pgv:match} and
Theorem~\ref{pgv:G3}, clause by clause in Lemma~\ref{pgv:E1g} (for edges read with the groups
exchanged use Theorem~\ref{ptw:Hc} and Proposition~\ref{ptw:revedge}); (E2) is the relation
$\chi_{f_1}=\chi_{e_2}$, $\chi_{f_2}=\chi_{e_1}$ of the edge characters there together with the
product condition at $V$, $V'$. The transformation laws of Theorem~\ref{ptr:square} are the
invariance used in Proposition~\ref{pgr:turned}, and the identity \eqref{ptr:eq-square} is the input
of Lemma~\ref{pgr:sumturn}.
\end{remark}

\subsection{The zeta integral does not vanish}\label{ptr:ss-Z}

Let $(S,T)$ be one of the pairs above, $\tau$ fixed, $G=\auU(S)_\tau$ with row parameter $\lambda$,
$\sigma=\omega_\infty(\cdot,1)$ at $\tau$, $\Phi=\Phi_\tau$ the local joint harmonic, $c_\lambda$ the
coefficient of a unit highest weight vector of the lowest $K$-type $\Lambda$ of $\pi_\lambda$,
$e_\lambda=d_\lambda\overline{c_\lambda}$ and $\mathsf P=\sigma(e_\lambda)$.

\begin{enumerate}
\item[(F)] For every positive definite $w\in G$, $\langle\Phi,\sigma(w)\Phi\rangle>0$.
\end{enumerate}
\noindent (F) holds: this is Theorem~\ref{ptq:F} (proved in full; the vacuum coefficient is Lemma~\ref{ptq:vac}).

\begin{proposition}\label{ptr:Z}
Then $Z_\infty\neq0$ for the pair $(S,T)$; precisely
$Z_\infty=\prod_\tau d_{\lambda_\tau}^{-1}\langle\Phi_\tau,\mathsf P_\tau\Phi_\tau\rangle$ and each factor is
positive.
\end{proposition}
\begin{proof}
$Z_\infty$ is the integral of $\langle P,\sigma(g)P\rangle\langle v_0,\pi_S(g)v_0\rangle$; as
$\overline{\pi_S}\cong\pi_\lambda$ with $\overline{v_0}$ the highest vector (proof of
Theorem~\ref{par:A3}), $\langle v_0,\pi_S(g)v_0\rangle=\overline{c_\lambda(g)}$, and the integral is a
product over $\tau$ of $\int\overline{c_\lambda}\langle\Phi,\sigma(g)\Phi\rangle=d_\lambda^{-1}
\langle\Phi,\mathsf P\Phi\rangle$ (Fubini; $c_\lambda\in L^1$ by Lemma~\ref{ptr:gap}). Fix $\tau$.

\emph{Two averages over $K$.} $\Phi$ is a highest weight vector of weight $\Lambda$ for $K$
(Definition~\ref{par:def-harm}, Lemma~\ref{par:relw}(iii)), so it generates an irreducible $K$-module
$V$ of highest weight $\Lambda$. The function $k\mapsto\langle\Phi,\sigma(k)\Phi\rangle/\|\Phi\|^2$
extends holomorphically to $K_\C$, equals $1$ at $1$, and is multiplied by $\Lambda(l)\Lambda(u)$
under $x\mapsto lxu$ ($l$ lower, $u$ upper triangular in each block, for the order of decreasing
$\lambda$); $\pdsm_\Lambda$ of \eqref{pds:eq-m} has the same properties, the Gauss decomposition is
dense, so $\langle\Phi,\sigma(k)\Phi\rangle=\|\Phi\|^2\pdsm_\Lambda(k)$, and likewise
$\langle v,kv\rangle=\pdsm_\Lambda(k)$ for the unit highest vector $v$ of the lowest $K$-type. By Schur
orthogonality on $K$ (mass $1$), with $d_\Lambda=\dim V=(\int|\pdsm_\Lambda|^2)^{-1}$, for $u\in K$:
\begin{equation}\label{ptr:eq-Kavg}
\int_K\pdsm_\Lambda(k^{-1}u)\,\sigma(k)\Phi\,dk=d_\Lambda^{-1}\sigma(u)\Phi .
\end{equation}
(Identify $V$ with the lowest $K$-type by $\Phi\mapsto\|\Phi\|v$; then $\pdsm_\Lambda(k^{-1}u)=
\langle kv,uv\rangle$ and $\int\langle kv,uv\rangle kv\,dk=d_\Lambda^{-1}uv$.)

\emph{Positivity of the averaged integrand.} Since $e_\lambda\in L^1$ and
$|\langle\Phi,\sigma(g)\Phi\rangle|\le\|\Phi\|^2$, $A(g)=\overline{c_\lambda(g)}\langle\Phi,\sigma(g)\Phi
\rangle$ is integrable and, by bi-invariance of Haar measure and Fubini,
$\int_GA=\int_G\bar A(g)\,dg$, $\bar A(g)=\iint_{K\times K}A(k_1gk_2)\,dk_1dk_2$. Fix $g=\kappa z$
($\kappa\in K$, $z$ positive definite, \S\ref{pds:ss-FJ}). Then $k_1\kappa zk_2=(k_1\kappa k_2)
(k_2^{-1}zk_2)$ with $k_2^{-1}zk_2$ positive definite, and \eqref{pds:eq-FJ} (Theorem~\ref{pco:FJ},
applicable by Lemma~\ref{ptr:gap}) with the substitution $u\mapsto k_2^{-1}u$ gives
\[
\overline{c_\lambda(k_1\kappa zk_2)}=d_\Lambda\int_K\frac{\overline{\pdsm_\Lambda(k_1\kappa u)}\,
\pdsm_\Lambda(k_2^{-1}u)}{\pdsm_{-\mu}(u^{-1}zu)}\,du ,
\]
{\emergencystretch=6em the denominator being positive (Lemma~\ref{pds:fischer}). Insert this in $\bar A(g)$, with
$\langle\Phi,\sigma(k_1\kappa zk_2)\Phi\rangle=\langle\sigma(k_1^{-1})\Phi,\allowbreak \sigma(\kappa z)\sigma(k_2)
\Phi\rangle$; all integrands are continuous on compact sets, so Fubini applies. By
\eqref{ptr:eq-Kavg}, $\int\pdsm_\Lambda(k_2^{-1}u)\sigma(k_2)\Phi\,dk_2=d_\Lambda^{-1}\sigma(u)\Phi$ and,
with $k=k_1^{-1}$, $\int\overline{\pdsm_\Lambda(k_1\kappa u)}\langle\sigma(k_1^{-1})\Phi,\cdot\rangle
dk_1=\langle d_\Lambda^{-1}\sigma(\kappa u)\Phi,\cdot\rangle$. Hence
\[
\bar A(\kappa z)=d_\Lambda^{-1}\int_K\frac{\langle\sigma(\kappa u)\Phi,\sigma(\kappa z)\sigma(u)\Phi
\rangle}{\pdsm_{-\mu}(u^{-1}zu)}du=d_\Lambda^{-1}\int_K\frac{\langle\Phi,\sigma(u^{-1}zu)\Phi\rangle}
{\pdsm_{-\mu}(u^{-1}zu)}\,du>0
\]
by (F) (Theorem~\ref{ptq:F}). So $\langle\Phi,\mathsf P\Phi\rangle=d_\lambda\int_G\bar A>0$.\par}
\end{proof}



%% ===== proofs/C-triangle2.tex =====
% proofs/C-triangle2.tex
% Continuation of proofs/C-triangle.tex (labels ptr:): the positivity (F) of Fock coefficients and
% the smoothing lemma (H) (Howe's argument). Labels of this file start with ptq:.
% Only amsmath/amssymb/amsthm are used. Conventions and macros: proofs/C-triangle.tex.
\providecommand{\auU}{\mathrm U}
\providecommand{\auHom}{\operatorname{Hom}}
\providecommand{\pzF}{\mathfrak F}
\providecommand{\pcG}{\mathsf G}
\providecommand{\ptqsgn}{\operatorname{sgn}}

\section{Positivity of Fock coefficients and the smoothing lemma}\label{ptq:sec}

This section continues \S\ref{ptr:sec} and proves the two local statements used there: the
positivity (F) of \S\ref{ptr:ss-Z} (Theorem~\ref{ptq:F}) and the smoothing statement (H) of
\S\ref{ptr:ss-mult} (Theorem~\ref{ptq:H}). Notation is that of \S\ref{ptr:sec} and of
\S\ref{par:ss-fock}--\S\ref{par:ss-A3}; inner products are conjugate-linear in the first variable.
Throughout, $(S,T)$ is one of the ascending or descending pairs of the triangle
(\S\ref{ptr:ss-param}), $\tau$ is fixed, $R$ indexes the rows (an orthogonal basis of
$S\otimes_{E,\tau}\C$, signs $\varepsilon_r$), $C$ the columns (signs $\varepsilon'_c$),
$\lambda,\lambda'$ are the row and column parameters, which match in the sense of
\eqref{ptr:eq-rule} (Proposition~\ref{ptr:match}), $G=\auU(S)_\tau$, $G'=\auU(T)_\tau$,
$\sigma=\omega_\infty(\cdot,1)$, $\sigma'=\omega_\infty(1,\cdot)$ at $\tau$, with derived actions on
polynomial vectors given by \eqref{par:eq-osc} (Lemma~\ref{ptr:const}, using Theorem~\ref{pco:V}),
and $\mathsf P=\sigma(e_\lambda)$, $\mathsf P'=\sigma'(e_{\lambda'})$, where for $f\in L^1$ we write
$\sigma(f)=\int f(x)\sigma(x)\,dx$. By Lemma~\ref{ptr:gap} and Theorem~\ref{par:A1}, $e_\lambda$,
$e_{\lambda'}$ and all their left and right derivatives are integrable.

We prove (F) (the vacuum coefficient is treated in Lemma~\ref{ptq:vac}) and (H). Besides
(D1)--(D3) (already used in \S\ref{ptr:ss-edge}), the proof of (H) uses two published results: the theorem of Harish-Chandra
that a $K$-finite, $\mathfrak Z$-finite function is reproduced by convolution with a smooth
compactly supported function, and Schmid's theorem on the $K$-types of discrete series.

\subsection{Matrix coefficients between harmonic polynomials}\label{ptq:ss-coef}

In this subsection $I$ is a finite set, $\eta\colon I\to\{\pm1\}$, $\pzF$ the Fock space in
variables $z_i$ ($i\in I$) with the Fischer inner product ($z_i$ adjoint to $\partial_i$, $\langle
1,1\rangle=1$), $a^\pm_i$ as in \S\ref{par:ss-fock} ($a^+_i=z_i$, $a^-_i=\partial_i$ if $\eta_i=1$;
$a^+_i=-\partial_i$, $a^-_i=z_i$ if $\eta_i=-1$), so that $(a^+_i)^*=\eta_ia^-_i$ on polynomials.
The space $W=\C^I$ carries the hermitian form $\sum_i\eta_i\bar w_iw'_i$ and its unitary group
$\auU(\eta)$; $\varrho$ is the Heisenberg representation on $\pzF$, normalised so that on
polynomial vectors
\begin{equation}\label{ptq:eq-heis}
d\varrho(w)=c_0\sum_i\bigl(w_i\,a^+_i-\eta_i\,\bar w_i\,a^-_i\bigr)\qquad(w\in W)
\end{equation}
for a nonzero real constant $c_0$ (this is the Fock model of \S\ref{par:ss-fock}: $z_i$ is the
coordinate or its conjugate according to $\eta_i$; see Remark~\ref{ptq:rem-conv} for the other
normalisation). Let $u\in\auU(\eta)$ and let $\omega(u)$ be a unitary operator on $\pzF$ with
\begin{equation}\label{ptq:eq-cov}
\omega(u)\varrho(w)\omega(u)^{-1}=\varrho(uw)\qquad(w\in W),
\end{equation}
which maps smooth vectors of $\varrho$ to smooth vectors. For polynomials $Q,P$ write
$\langle Q,\omega(u)P\rangle$ for the coefficient between the polynomial vectors.

\begin{lemma}[Moving the ladder operators]\label{ptq:ladder}
For polynomials $P,Q$ and $j\in I$:
\[
\langle Q,\omega(u)a^+_jP\rangle=\sum_iu_{ij}\eta_i\langle a^-_iQ,\omega(u)P\rangle,\qquad
\eta_j\langle Q,\omega(u)a^-_jP\rangle=\sum_i\overline{u_{ij}}\,\langle a^+_iQ,\omega(u)P\rangle .
\]
\end{lemma}
\begin{proof}
By \eqref{ptq:eq-cov}, $\omega(u)d\varrho(w)P=d\varrho(uw)\omega(u)P$, and $d\varrho(uw)$ is
skew-adjoint (it is $c_0\sum_i((uw)_ia^+_i-((uw)_ia^+_i)^*)$). Hence
$\langle Q,\omega(u)d\varrho(w)P\rangle=-\langle d\varrho(uw)Q,\omega(u)P\rangle$. Expanding both sides
with \eqref{ptq:eq-heis} and using conjugate-linearity in the first variable,
\[
\sum_jw_j\langle Q,\omega(u)a^+_jP\rangle-\sum_j\eta_j\bar w_j\langle Q,\omega(u)a^-_jP\rangle
=-\sum_i\overline{(uw)_i}\langle a^+_iQ,\omega(u)P\rangle+\sum_i\eta_i(uw)_i\langle a^-_iQ,\omega(u)P\rangle .
\]
Both sides are sums of a complex-linear and a complex-antilinear function of $w\in\C^I$; such a
decomposition is unique. Comparing the linear parts gives the first identity, the antilinear parts
the second.
\end{proof}

Let $P_\pm$ be the projections onto the coordinates with $\eta_i=\pm1$, $\bar u$ the entrywise
conjugate, and
\[
C(u)=P_+uP_++P_-\bar uP_- .
\]
Its diagonal blocks are the block $u_{++}$ of $u$ and the conjugate of the block $u_{--}$; both are
invertible for $u\in\auU(\eta)$ (from $u^*\mathrm{diag}(\eta)u=\mathrm{diag}(\eta)$,
$u_{++}^*u_{++}=1+u_{-+}^*u_{-+}$ and $u_{--}^*u_{--}=1+u_{+-}^*u_{+-}$), so $C(u)$ is invertible.
Fix a map $d$ from $I$ to a set of \emph{labels}, and assume that $u$ is \emph{block diagonal}:
$u_{ij}=0$ when $d(i)\ne d(j)$. Call a polynomial $P$ \emph{harmonic} if $\partial_i\partial_jP=0$
whenever $d(i)=d(j)$ and $\eta_i\ne\eta_j$.

\begin{lemma}\label{ptq:pderiv}
\textup{(a)} If $\partial_kQ=0$ for all $k$ with $d(k)=d(i)$, $\eta_k\ne\eta_i$, then
$\langle Q,\omega(u)\partial_iP\rangle=\sum_k\overline{C(u)_{ki}}\,\langle z_kQ,\omega(u)P\rangle$.

\textup{(b)} If $\partial_iQ=0$ for all $i$ with $d(i)=d(j)$, $\eta_i\ne\eta_j$, then
$\langle z_jQ,\omega(u)P\rangle=\sum_i\overline{A_{ij}}\,\langle Q,\omega(u)\partial_iP\rangle$ with
$A=C(u)^{-1}$.
\end{lemma}
\begin{proof}
(a) If $\eta_i=1$, then $\partial_i=a^-_i$ and the second identity of Lemma~\ref{ptq:ladder} gives
$\sum_k\overline{u_{ki}}\langle a^+_kQ,\omega(u)P\rangle$. For $\eta_k=1$, $a^+_k=z_k$ and
$C(u)_{ki}=u_{ki}$. For $\eta_k=-1$, $a^+_kQ=-\partial_kQ$, and the term vanishes: either
$d(k)\ne d(i)$ and $u_{ki}=0$, or $d(k)=d(i)$, $\eta_k\ne\eta_i$ and $\partial_kQ=0$; and
$C(u)_{ki}=0$. If $\eta_i=-1$, then $\partial_i=-a^+_i$ and the first identity gives
$-\sum_ku_{ki}\eta_k\langle a^-_kQ,\omega(u)P\rangle$. For $\eta_k=-1$, $a^-_k=z_k$, the coefficient is
$u_{ki}=\overline{C(u)_{ki}}$; for $\eta_k=1$, $a^-_kQ=\partial_kQ$ and the term vanishes as before,
while $C(u)_{ki}=0$.

(b) Put $w_k=\langle z_kQ,\omega(u)P\rangle$ and let $b=\{i:d(i)=d(j),\ \eta_i=\eta_j\}$. For $i\in b$
the hypothesis of (a) holds (if $d(k)=d(i)$ and $\eta_k\ne\eta_i$ then $d(k)=d(j)$, $\eta_k\ne\eta_j$),
so $(C(u)^*w)_i=\langle Q,\omega(u)\partial_iP\rangle$. The matrix $C(u)$ is block diagonal for the
partition of $I$ by the pair $(d,\eta)$ ($C(u)_{ki}=0$ if $\eta_k\ne\eta_i$, and $u_{ki}=0$ if
$d(k)\ne d(i)$), and so are $C(u)^*$ and $(C(u)^*)^{-1}=A^*$; $b$ is one of the blocks. Hence
$w_j=(A^*C(u)^*w)_j=\sum_{i\in b}(A^*)_{ji}(C(u)^*w)_i$, and $(A^*)_{ji}=\overline{A_{ij}}$ vanishes for
$i\notin b$.
\end{proof}

\begin{lemma}[Harmonic induction]\label{ptq:induct}
Let $\mathcal S$ be a space of polynomials containing $1$ and containing $z_jP$ whenever $P\in\mathcal S$
and $\partial_iP=0$ for all $i$ with $d(i)=d(j)$, $\eta_i\ne\eta_j$. Then $\mathcal S$ contains every
harmonic polynomial.
\end{lemma}
\begin{proof}
The operators $\partial_i\partial_j$ are homogeneous, so the homogeneous components of a harmonic
polynomial are harmonic; it suffices to treat a harmonic $P$ homogeneous of degree $n$, by induction
on $n$. For $n=0$, $P\in\C\cdot1\subset\mathcal S$. For $n\ge1$, Euler's identity gives
$nP=\sum_jz_j\,\partial_jP$; each $\partial_jP$ is harmonic, homogeneous of degree $n-1$, hence in
$\mathcal S$, and $\partial_i\partial_jP=0$ whenever $d(i)=d(j)$, $\eta_i\ne\eta_j$. So
$z_j\partial_jP\in\mathcal S$ and $P\in\mathcal S$.
\end{proof}

For a matrix $X$ let $L_X$ be the substitution $(L_XQ)(z)=Q(zX)$, that is,
$z_j\mapsto\sum_iX_{ij}z_i$. Then $L_XL_Y=L_{XY}$, and $\langle L_XQ,P\rangle=\langle Q,L_{X^*}P\rangle$:
the Fischer space of polynomials of degree $n$ is the $n$-th symmetric power of the space of linear
forms with the induced inner product, $L_X$ is the symmetric power of its action on linear forms,
and on linear forms $\langle L_Xz_j,z_k\rangle=\overline{X_{kj}}=\langle z_j,L_{X^*}z_k\rangle$.

\begin{proposition}[Coefficients between harmonic polynomials]\label{ptq:coef}
Let $u\in\auU(\eta)$ be block diagonal and $c_u=\langle1,\omega(u)1\rangle$. For harmonic $P,Q$,
\[
\langle Q,\omega(u)P\rangle=c_u\,\langle L_{A}Q,P\rangle,\qquad A=C(u)^{-1}.
\]
\end{proposition}
\begin{proof}
\emph{Step 1: $Q=1$.} Let $\mathcal S=\{P:\langle1,\omega(u)P\rangle=c_u\langle1,P\rangle\}$; it
contains $1$. Let $R\in\mathcal S$ with $\partial_iR=0$ for $d(i)=d(j)$, $\eta_i\ne\eta_j$. The element
$u^{-1}=\mathrm{diag}(\eta)u^*\mathrm{diag}(\eta)$ is block diagonal, and $\omega(u)^*=\omega(u)^{-1}$
satisfies \eqref{ptq:eq-cov} for $u^{-1}$. By Lemma~\ref{ptq:pderiv}(b) for $u^{-1}$ and $P=1$,
$\langle z_jR,\omega(u)^{-1}1\rangle=\sum_i\overline{A'_{ij}}\langle R,\omega(u)^{-1}\partial_i1\rangle=0$. Hence
$\langle1,\omega(u)z_jR\rangle=\overline{\langle z_jR,\omega(u)^{-1}1\rangle}=0=c_u\langle\partial_j1,R\rangle
=c_u\langle1,z_jR\rangle$, so $z_jR\in\mathcal S$. By Lemma~\ref{ptq:induct}, every harmonic $P$
lies in $\mathcal S$.

\emph{Step 2.} Let $\mathcal S'$ be the space of $Q$ such that the identity holds for all harmonic
$P$; $1\in\mathcal S'$ by Step 1. Let $R\in\mathcal S'$ with $\partial_iR=0$ for $d(i)=d(j)$,
$\eta_i\ne\eta_j$, and $P$ harmonic. The derivatives $\partial_iP$ are harmonic. By
Lemma~\ref{ptq:pderiv}(b) and the induction hypothesis,
\[
\langle z_jR,\omega(u)P\rangle=\sum_i\overline{A_{ij}}\,c_u\langle L_AR,\partial_iP\rangle
=c_u\sum_i\overline{A_{ij}}\langle z_iL_AR,P\rangle=c_u\Bigl\langle\sum_iA_{ij}z_i\,L_AR,P\Bigr\rangle
=c_u\langle L_A(z_jR),P\rangle .
\]
So $z_jR\in\mathcal S'$, and Lemma~\ref{ptq:induct} gives the claim for all harmonic $Q$.
\end{proof}

\begin{remark}\label{ptq:rem-conv}
If the Fock model is normalised with $\bar w_ia^+_i-\eta_iw_ia^-_i$ in \eqref{ptq:eq-heis}, then
Lemma~\ref{ptq:ladder} and Proposition~\ref{ptq:coef} hold with $u$ replaced by $\bar u$ in
the coefficients (the proof is the same, with linear and antilinear parts exchanged). Below $u$ is
hermitian positive definite, and so is $\bar u$; the conclusions do not change.
\end{remark}

\subsection{The positivity (F)}\label{ptq:ss-F}

We use the following lemma on the vacuum coefficient.

\begin{lemma}[The vacuum coefficient does not vanish]\label{ptq:vac}
In the setting of \S\ref{ptq:ss-coef}, let $u\in\auU(\eta)$ be block diagonal and $\omega(u)$ as
there. Then $c_u=\langle1,\omega(u)1\rangle\ne0$.
\end{lemma}
\begin{proof}
Let $I_\pm=\{i:\eta_i=\pm1\}$, $W_\pm=\{w\in W:w_i=0\text{ for }i\notin I_\pm\}$, $\mathcal P_\pm$ the
polynomials in the variables $z_i$, $i\in I_\pm$, and $H_\pm$ the closure of $\mathcal P_\pm$ in $\pzF$.
Recall that $\varrho$ is a Heisenberg representation: $\varrho(x)\varrho(y)=\beta(x,y)\varrho(x+y)$ with
$|\beta|=1$, and $\varrho$ is irreducible (\S\ref{pwf:ss-def}).

\emph{Step 1: $\varrho(w)1\in H_\pm$ for $w\in W_\pm$.} By \eqref{ptq:eq-heis}, for $w\in W_+$,
$d\varrho(w)=c_0\sum_{i\in I_+}(w_iz_i-\bar w_i\partial_i)$, and for $w\in W_-$,
$d\varrho(w)=c_0\sum_{i\in I_-}(\bar w_iz_i-w_i\partial_i)$; so $d\varrho(w)$ maps $\mathcal P_\pm$ into
itself. In the Fischer norm, for $P$ of degree $\le d$, $\|z_iP\|^2=\|P\|^2+\|\partial_iP\|^2$ and
$\sum_i\|\partial_iP\|^2\le d\|P\|^2$ (both by $[\partial_i,z_i]=1$, $z_i^*=\partial_i$ and Euler's identity
on homogeneous components, which are orthogonal); hence $\|z_iP\|,\|\partial_iP\|\le\sqrt{d+1}\,\|P\|$ and
$\|d\varrho(w)P\|\le M\sqrt{d+1}\,\|P\|$ with $M=2|c_0|\sum_i|w_i|$. Thus $\|d\varrho(w)^n1\|\le
M^n\sqrt{n!}$. Let $A$ be the skew-adjoint generator of the unitary group $s\mapsto\varrho(sw)$ (a group because $\beta(sw,s'w)=1$, $\beta$ being the exponential of an alternating form); it extends
$d\varrho(w)$ on polynomial vectors, which are smooth, so $1\in\bigcap_nD(A^n)$ with $A^n1=d\varrho(w)^n1$.
Since $\sum_n|s|^nM^n\sqrt{n!}/n!<\infty$ for every $s$, the spectral theorem gives
$\varrho(sw)1=\sum_n s^nA^n1/n!$, a norm-convergent series of elements of $\mathcal P_\pm$. Take $s=1$.

\emph{Step 2.} Suppose $c_u=0$. Polynomials in $\mathcal P_+$ and in $\mathcal P_-$ are harmonic
(each involves only variables with one value of $\eta$), so by Proposition~\ref{ptq:coef}
$\langle Q,\omega(u)P\rangle=0$ for $Q\in\mathcal P_-$, $P\in\mathcal P_+$, and by continuity
$\langle x',\omega(u)x\rangle=0$ for $x'\in H_-$, $x\in H_+$. By Step 1, for $w\in W_+$, $w'\in W_-$,
\[
0=\langle\varrho(w')1,\omega(u)\varrho(w)1\rangle=\langle1,\varrho(w')^{-1}\omega(u)\varrho(w)1\rangle
=\langle1,\omega(u)\varrho(u^{-1}w')^{-1}\varrho(w)1\rangle ,
\]
by \eqref{ptq:eq-cov}; and $\varrho(u^{-1}w')^{-1}\varrho(w)$ is a scalar of modulus one times
$\varrho(w-u^{-1}w')$ (Heisenberg relations; $\varrho(y)^{-1}$ is a scalar times $\varrho(-y)$). Hence
$\langle\omega(u)^{-1}1,\varrho(y)1\rangle=0$ for every $y\in W_++u^{-1}W_-$.

\emph{Step 3: $W_++u^{-1}W_-=W$.} The element $u^{-1}$ lies in $\auU(\eta)$, so its block
$(u^{-1})_{--}$ is invertible (as shown before Lemma~\ref{ptq:pderiv}). If $w'\in W_-$ and
$u^{-1}w'\in W_+$, then $(u^{-1})_{--}w'=0$, so $w'=0$. Thus $W_+\cap u^{-1}W_-=0$, and the complex
dimensions $|I_+|+|I_-|=|I|$ add up.

\emph{Step 4.} By Steps 2 and 3, $\omega(u)^{-1}1$ is orthogonal to $\varrho(y)1$ for all $y\in W$. The
closed span of the $\varrho(y)1$ is $\varrho(W)$-stable (Heisenberg relations) and non-zero, hence is
$\pzF$ by irreducibility. So $\omega(u)^{-1}1=0$, which is absurd since $\omega(u)$ is unitary. Hence
$c_u\ne0$.
\end{proof}

\begin{theorem}[(F)]\label{ptq:F}
Let $(S,T)$ be an ascending or descending pair of the triangle, $\Phi=\Phi_\tau$ the local joint
harmonic of Theorem~\ref{ptr:edge}, and $w\in G$ positive definite. Then
$\langle\Phi,\sigma(w)\Phi\rangle>0$. More generally $\langle P,\sigma(w)P\rangle>0$ for every
polynomial $P\ne0$ none of whose monomials contains two variables of one column from rows of
opposite signs.
\end{theorem}
\begin{proof}
Write $w=\exp X$ with $X$ in the $-1$ eigenspace $\mathfrak p$ of the Cartan involution (hermitian,
off-diagonal for the signs); $\operatorname{tr}X=0$, so $\det w=1$, and $w_t=\exp tX$ ($0\le t\le1$) are
positive definite.

Let $u_t=\iota(w_t,1)$, $x\mapsto xw_t^{-1}$ on $W=\auHom(S,T)_\tau$.

\emph{The element $u=u_1$.} In the coordinates $x_{rc}$, $u$ acts on each column separately (it
mixes the row index only), so it is block diagonal for the labels $d(r,c)=c$; it is hermitian
positive definite for the positive majorant (the right multiplication by the positive definite
$w^{-1}$), and lies in $\auU(\eta)$, $\eta_{rc}=-\varepsilon_r\varepsilon'_c$. A polynomial is harmonic
for $d$ if $\partial_{rc}\partial_{r'c}P=0$ when $\eta_{rc}\ne\eta_{r'c}$, i.e.\ $\varepsilon_r\ne
\varepsilon_{r'}$; this holds if no monomial of $P$ contains $z_{rc}z_{r'c}$ with $\varepsilon_r\ne
\varepsilon_{r'}$, in particular for $\Phi$ (Theorem~\ref{ptr:edge}(a)). The operator $\sigma(w)$
satisfies \eqref{ptq:eq-cov} for $u$ by (W1$_\infty$) of \S\ref{pco:ss-V}.

\emph{The constant.} $c(t)=\langle1,\sigma(w_t)1\rangle$ is continuous (strong continuity) and $c(0)=1$.
It is real: let $J=\mathrm{diag}(\varepsilon)\in K$; since $w_t$ is hermitian and $w_t^*Jw_t=J$,
$w_t^{-1}=Jw_tJ$. The group $K$ acts on polynomial vectors by linear substitutions times a character
(proof of Lemma~\ref{par:vac}), so $\sigma(J)1=\chi1$ with $|\chi|=1$, and
$\overline{c(t)}=\langle1,\sigma(w_t^{-1})1\rangle=\langle\sigma(J)^{-1}1,\sigma(w_t)\sigma(J)^{-1}1\rangle=
|\chi|^{-2}c(t)=c(t)$. It never vanishes: everything said above about $u$ holds for $u_t$ (block
diagonal for the labels $d(r,c)=c$, in $\auU(\eta)$, and \eqref{ptq:eq-cov} holds for $\sigma(w_t)$ and
$u_t$ by (W1$_\infty$)), so Lemma~\ref{ptq:vac} gives $c(t)\ne0$. Hence $c(1)=c_u>0$. Moreover
$C(u)$ is positive definite: its blocks are the principal submatrix $u_{++}$ of the positive definite
$u$ and the conjugate of $u_{--}$.

\emph{Positivity.} $C=C(u)$ is positive definite, hence so is $A=C^{-1}$; write $A=BB^*$ with $B$
invertible. By Proposition~\ref{ptq:coef} with $Q=P$,
\[
\langle P,\sigma(w)P\rangle=c_u\langle L_BL_{B^*}P,P\rangle=c_u\langle L_{B^*}P,L_{B^*}P\rangle=c_u\,
\|L_{B^*}P\|^2>0,
\]
since $L_{B^*}$ is invertible (with inverse $L_{(B^*)^{-1}}$) and $P\ne0$.
\end{proof}

\begin{remark}\label{ptq:rem-F}
Theorem~\ref{ptq:F} is (F) of \S\ref{ptr:ss-Z}. Hence Proposition~\ref{ptr:Z} holds
unconditionally: $Z_\infty\ne0$ for every pair of the triangle, with each local factor positive.
The same argument gives $\langle P,\sigma'(w)P\rangle>0$ for positive definite $w\in G'$ and $P\ne0$
none of whose monomials contains two variables of one row from columns of opposite signs (now
$u=\iota(1,w)$ is block diagonal for the labels $d(r,c)=r$).
\end{remark}

\subsection{The smoothing lemma (H)}\label{ptq:ss-H}

We use, besides (D1)--(D3) of \S\ref{par:ss-rigid}, the following published results.
\begin{enumerate}
\item[(D4)] \emph{(Schmid \cite[Lemma~(2.22), Lemma~(3.8), Corollary~(3.9), Corollary~(4.19) and the proof of Theorem~(1.3) on p.~560]{Schmid75}. Also a consequence of Blattner's
formula, proved by Hecht and Schmid. Schmid's positive system is $\{\alpha:\langle\alpha,\lambda\rangle>0\}$, as here; the discrete series concerned are not holomorphic in general, see the paragraph ``On (D3)'' in \S\ref{par:ss-rigid}.)} Every $K$-type of $\pi_\lambda$ has highest weight of the form
$\Lambda^K(\lambda)+\sum_\beta n_\beta\beta$, $n_\beta\in\Z_{\ge0}$, the sum over the positive
noncompact roots $\beta$ (positive meaning $\langle\beta,\lambda\rangle>0$).
\item[(D5)] \emph{(Harish-Chandra \cite[Theorem~1, p.~18]{HC66}; used in \cite[Chapter~I, proof of Lemma~14, p.~15]{HarishChandraAF}; see also \cite[\S2.2, p.~193]{BorelJacquet}.)} Let $f$ be a smooth
function on a connected reductive group, $K$-finite on the left and on the right and
$\mathfrak Z(\mathfrak g)$-finite. Then there is $\beta\in C_c^\infty$ with $\beta*f=f$, where
$(\beta*f)(g)=\int\beta(y)f(y^{-1}g)\,dy$. The published statement is the right-convolution
form: for $f$ smooth, $\mathfrak Z(\mathfrak g)$-finite and right $K$-finite there is
$\alpha\in C_c^\infty$ with $f*\alpha=f$, $(f*\alpha)(g)=\int f(gy^{-1})\alpha(y)\,dy$. The left
form follows: $f^\vee(g)=f(g^{-1})$ is right $K$-finite (as $f$ is left $K$-finite) and
$\mathfrak Z(\mathfrak g)$-finite (as $\mathfrak Z(\mathfrak g)$ is stable under the principal
antiautomorphism of $U(\mathfrak g)$, which intertwines the right and left regular actions under
$f\mapsto f^\vee$). Take $\alpha$ with $f^\vee*\alpha=f^\vee$ and put $\beta(y)=\alpha(y^{-1})$.
Substituting $y\mapsto y^{-1}$ (the group is unimodular),
$(\beta*f)(g)=\int\alpha(y)f(yg)\,dy=\int f^\vee(g^{-1}y^{-1})\alpha(y)\,dy=(f^\vee*\alpha)(g^{-1})=f(g)$.
\end{enumerate}
(D5) applies to $f=e_{\lambda'}$ on $G'$: $e_{\lambda'}=d_{\lambda'}\overline{c_{\lambda'}}$ is a
coefficient of the $K'$-finite vector $v'$ of the irreducible unitary $\pi_{\lambda'}$, so it is
smooth, bi-$K'$-finite, and $\mathfrak Z(\mathfrak g')$ acts on it through the infinitesimal
character. Let $H_\lambda\subset\pzF_\tau$ be the closed span of the images of the $G$-equivariant
isometries $\pi_\lambda\to\sigma$; it is $G$-stable, its orthogonal projection $p$ commutes with
$\sigma(G)$, and $H_\lambda\cong\pi_\lambda\hat\otimes M$ for a Hilbert space $M$, so the $K$-types of
$H_\lambda$ are those of $\pi_\lambda$.

\subsubsection*{Four properties of $\mathsf P$}

\begin{lemma}\label{ptq:P}
\textup{(a)} \emph{(Integration by parts.)} For a polynomial $Q$ and $Y$ in $\mathfrak g_\C$,
$\mathsf P(d\sigma(Y)Q)=-\sigma(R_{Y}e_\lambda)Q$, where $R_Y$ is the right derivative extended
complex-linearly. Consequently $\mathsf P(E_{ba}Q)=0$ if $\varepsilon_a=\varepsilon_b$ and
$\lambda_a>\lambda_b$, and $\mathsf P(E_{aa}Q)=\Lambda_a\mathsf PQ$, $\Lambda=\Lambda^K(\lambda)$.

\textup{(b)} \emph{(Image.)} For every $x\in\pzF_\tau$, $\mathsf Px\in H_\lambda$; and $\mathsf P$ is
self-adjoint.

\textup{(c)} \emph{(Off the cone.)} Let $Q$ be a polynomial with $E_{aa}Q=\nu_aQ$ for all $a$ and
$E_{ab}Q=0$ whenever $\varepsilon_a=\varepsilon_b$, $\lambda_a>\lambda_b$, where $\nu-\Lambda$ is not a
sum of roots $e_a-e_b$ with $\lambda_a>\lambda_b$ with nonnegative coefficients. Then
$\mathsf P\,\sigma(y)Q=0$ for all $y\in G$, and $\mathsf P(E(X)Q)=0$ for every $X\in\mathfrak{gl}_n(\C)$,
$E(X)$ the row action \eqref{par:eq-osc}.

\textup{(d)} \emph{(Fixing.)} There is $\beta_0\in C_c^\infty(G')$ with
$\sigma'(\beta_0)\mathsf P\Phi=\mathsf P\Phi$.
\end{lemma}
\begin{proof}
(a) For real $Y$: since $Q$ is a smooth vector with derivative $d\sigma(Y)Q$,
$\mathsf Pd\sigma(Y)Q=\frac d{dt}\big|_0\mathsf P\sigma(\exp tY)Q$, and
$\mathsf P\sigma(\exp tY)Q=\int e_\lambda(g\exp(-tY))\sigma(g)Q\,dg$ (right invariance of Haar
measure). Now $e_\lambda(g\exp(-tY))-e_\lambda(g)=-\int_0^t(R_Ye_\lambda)(g\exp(-sY))ds$, so by Fubini
(as $R_Ye_\lambda\in L^1$ and $\sigma$ is unitary)
$\mathsf P\sigma(\exp tY)Q-\mathsf PQ=-\int_0^t\int(R_Ye_\lambda)(g)\sigma(g\exp sY)Q\,dg\,ds$; the inner
integral is continuous in $s$ by dominated convergence, which gives the formula at $t=0$. Both sides
are complex-linear in $Y$. Now $e_\lambda(g)=d_\lambda\langle\pi_\lambda(g)v,v\rangle$ ($v$ the unit highest
vector of the lowest $K$-type, $c_\lambda(g)=\langle v,\pi_\lambda(g)v\rangle$), so for real $Y$,
$R_Ye_\lambda(g)=d_\lambda\langle\pi_\lambda(g)d\pi_\lambda(Y)v,v\rangle$, and since this is antilinear in
$d\pi_\lambda(Y)v$, for complex $Y$ it is $d_\lambda\langle\pi_\lambda(g)d\pi_\lambda(\bar Y)v,v\rangle$, where
$\bar Y=-JY^*J$ ($J=\mathrm{diag}(\varepsilon)$) is the conjugation of $\mathfrak g_\C$ with respect to
$\mathfrak g$. For $\varepsilon_a=\varepsilon_b$, $\bar E_{ba}=-E_{ab}$ and $\bar E_{aa}=-E_{aa}$. As
$d\pi_\lambda(E_{ab})v=0$ for $\lambda_a>\lambda_b$ and $d\pi_\lambda(E_{aa})v=\Lambda_av$ (with $\Lambda_a$ real),
$R_{E_{ba}}e_\lambda=0$ and $R_{E_{aa}}e_\lambda=-\Lambda_ae_\lambda$; with $d\sigma(E_{ab})=E_{ab}$ on
polynomials this gives the two consequences.

(b) $e_\lambda(x^{-1})=\overline{e_\lambda(x)}$ (proof of Theorem~\ref{par:A3}), so $\sigma(e_\lambda)^*=
\sigma(e_\lambda)$. Let $y=\mathsf Px$. For $Y\in\mathfrak g$, $\sigma(\exp tY)y=\int e_\lambda(\exp(-tY)g)
\sigma(g)x\,dg$ and $e_\lambda(\exp(-tY)g)=d_\lambda\langle\pi_\lambda(g)v,\pi_\lambda(\exp tY)v\rangle$. As the left
derivatives of $e_\lambda$ are integrable, $y$ is a smooth vector and, for $Y,Y'\in\mathfrak g_\C$,
$d\sigma(Y)d\sigma(Y')y=d_\lambda\int\langle\pi_\lambda(g)v,d\pi_\lambda(Y)d\pi_\lambda(Y')v\rangle\sigma(g)x\,dg$
(the expression is linear in the second slot, so complex-linear in $Y,Y'$). Hence $y$ has
$T$-weight $\Lambda$, is killed by the $E_{ab}$ with $\varepsilon_a=\varepsilon_b$, $\lambda_a>\lambda_b$,
and $d\sigma(\Omega)y=\kappa_\lambda y$, because $v$ has these properties in $\pi_\lambda$. By the proof of
Lemma~\ref{par:rigid} (which uses only these three properties, $e_\lambda\in L^1$ and (D1)--(D3)), the
closed span of $\sigma(G)y$ is a direct integral of copies of $\pi_\lambda$, hence is contained in
$H_\lambda$.

(c) The polynomial vectors of bounded degree form a finite-dimensional space stable under $K$
(which acts by linear substitutions, proof of Lemma~\ref{par:vac}), so $Q$ is $K$-finite; being a
highest weight vector of weight $\nu$ for $\mathfrak k_\C$ in a finite-dimensional unitary
representation, it generates an irreducible $K$-module of highest weight $\nu$. The projection $p$
commutes with $K$ and with the derived action on smooth vectors, so $pQ$ is zero or generates a
$K$-type of $H_\lambda$ of highest weight $\nu$. By (D4) the highest weights of the $K$-types of
$H_\lambda$ lie in $\Lambda+\sum\Z_{\ge0}\beta$, contained in $\Lambda$ plus the cone of the hypothesis;
so $pQ=0$. Let $y\in G$ and $x=\mathsf P\sigma(y)Q$. By (b), $x\in H_\lambda$ and $\mathsf Px\in H_\lambda$,
and
\[
\|x\|^2=\langle\mathsf P\sigma(y)Q,x\rangle=\langle\sigma(y)Q,\mathsf Px\rangle=\langle p\,\sigma(y)Q,\mathsf
Px\rangle=\langle\sigma(y)pQ,\mathsf Px\rangle=0 .
\]
Differentiating $t\mapsto\mathsf P\sigma(\exp tY)Q=0$ at $0$ gives $\mathsf P(E(Y)Q)=0$ for
$Y\in\mathfrak g$, hence for $Y\in\mathfrak g_\C=\mathfrak{gl}_n(\C)$ by linearity.

(d) By (D5) choose $\beta_0\in C_c^\infty(G')$ with $\beta_0*e_{\lambda'}=e_{\lambda'}$. Since
$\overline{c_{\lambda'}(y^{-1}g)}=c_{\lambda'}(g^{-1}y)$, this reads $\int\beta_0(y)c_{\lambda'}(hy)\,dy=
c_{\lambda'}(h)$ for all $h$. Let $x=\mathsf P\Phi$. In the proof of Theorem~\ref{par:A3} (valid for the
present pair, Theorem~\ref{ptr:edge}(b)) it is shown that $x$ is a smooth vector for $\sigma'$ of
$T'$-weight $\Lambda'$, killed by the compact raising operators of the columns, with
$d\sigma'(\Omega')x=\kappa_{\lambda'}x$; so Lemma~\ref{par:rigid} gives $\langle x,\sigma'(h)x\rangle=
\|x\|^2c_{\lambda'}(h)$. Then for $g\in G'$,
\[
\langle\sigma'(g)x,\sigma'(\beta_0)x\rangle=\int\beta_0(y)\langle x,\sigma'(g^{-1}y)x\rangle dy=\|x\|^2\!\int
\beta_0(y)c_{\lambda'}(g^{-1}y)dy=\|x\|^2c_{\lambda'}(g^{-1})=\langle\sigma'(g)x,x\rangle .
\]
So $\sigma'(\beta_0)x-x$, which lies in the closed span of $\sigma'(G')x$, is orthogonal to it, and
is $0$.
\end{proof}

\subsubsection*{Fischer's argument}

For a polynomial space we write $\mathcal D_j$ for the polynomials of degree $\le j$. The operators
$E_{ab}$, $E'_{cd}$ with equal signs preserve degree, and on polynomials $E_{ab}^*=E_{ba}$,
$E'^*_{cd}=E'_{dc}$ for the Fischer product (from $(a^+_i)^*=\eta_ia^-_i$, $s_{ab}=1$; the shifts are
real). For $\varepsilon_a\ne\varepsilon_b$, $\Delta_{ab}=\sum_{c:\eta_{ac}=-1}\partial_{ac}\partial_{bc}$ as in
\S\ref{par:ss-fock}; note that $\eta_{bc}=-\eta_{ac}$.

\begin{lemma}\label{ptq:fischer}
Let $\mathcal G$ be a space of polynomials containing $\mathcal D_{j-1}$, and $Q\in\mathcal D_j$
orthogonal to $\mathcal G\cap\mathcal D_j$. Let $\varepsilon_a\ne\varepsilon_b$ and suppose that
$E_{ba}\Delta_{ab}Q\in\mathcal G$. Then $\Delta_{ab}Q=0$. The same holds for columns, with
$\Delta'_{cd}$ and $E'_{dc}$.
\end{lemma}
\begin{proof}
Put $Q''=\Delta_{ab}Q\in\mathcal D_{j-2}$. By \eqref{par:eq-osc}, $\pm E_{ba}=\sum_ca^+_{bc}a^-_{ac}$; for
$\eta_{ac}=-1$ (so $\eta_{bc}=1$) the term is $z_{bc}z_{ac}$, and for $\eta_{ac}=1$ it is
$-\partial_{bc}\partial_{ac}$. Hence $\sum_{c:\eta_{ac}=-1}z_{ac}z_{bc}Q''=\pm E_{ba}Q''+
\sum_{c:\eta_{ac}=1}\partial_{ac}\partial_{bc}Q''$. The first term is in $\mathcal G\cap\mathcal D_j$; the
second is in $\mathcal D_{j-4}\subset\mathcal G$. By adjointness,
\[
\|Q''\|^2=\langle\Delta_{ab}Q,Q''\rangle=\Bigl\langle Q,\sum_{c:\eta_{ac}=-1}z_{ac}z_{bc}Q''\Bigr\rangle=0 .
\]
For columns the computation is the same with rows and columns exchanged.
\end{proof}

\begin{lemma}\label{ptq:tensor}
Let $Q$ satisfy $E_{aa}Q=\Lambda_aQ$ for all $a$ and $E_{ab}Q=0$ for $\varepsilon_a=\varepsilon_b$,
$\lambda_a>\lambda_b$. Suppose that, for $\varepsilon_a\ne\varepsilon_b$, $\Delta_{ab}Q=0$ whenever
$\Delta_{ab}Q$ has these two properties for the weight $\Lambda+e_a-e_b$. Then $\Delta_{ab}Q=0$ for all
$\varepsilon_a\ne\varepsilon_b$.
\end{lemma}
\begin{proof}
From the canonical commutation relations, for $\varepsilon_x=\varepsilon_z$ and $\varepsilon_a\ne\varepsilon_b$,
$[E_{xz},\Delta_{ab}]=\delta_{za}\Delta_{xb}-\delta_{bx}\Delta_{az}$ (the $\Delta_{ab}$ transform like
the matrix units). Argue by induction on $\lambda_b-\lambda_a$. $\Delta_{ab}Q$ has $T$-weight
$\Lambda+e_a-e_b$ by the case $x=z$. Let $\varepsilon_x=\varepsilon_z$, $\lambda_x>\lambda_z$. Then
$E_{xz}\Delta_{ab}Q=\delta_{za}\Delta_{xb}Q-\delta_{bx}\Delta_{az}Q$; in the first term $\lambda_x>\lambda_a$,
so $\lambda_b-\lambda_x<\lambda_b-\lambda_a$ (and $\varepsilon_x=\varepsilon_a\ne\varepsilon_b$); in the second
$\lambda_z<\lambda_b$, so $\lambda_z-\lambda_a<\lambda_b-\lambda_a$. By induction both vanish, so
$\Delta_{ab}Q$ has the two properties, and is $0$ by hypothesis.
\end{proof}

\subsubsection*{Howe's argument}

\begin{proposition}[Howe's argument]\label{ptq:howe}
Let $\Phi$ be the joint harmonic of the matching pair $(\lambda,\lambda')$ (\S\ref{par:ss-harm}) and
$\mathcal G$ a space of polynomials such that: \textup{(i)} $E'_{cd}\mathcal G\subset\mathcal G$ for all
$c,d$; \textup{(ii)} $\Phi\in\mathcal G$; \textup{(iii)} $E(X)Q\in\mathcal G$ for every $X$ and every $Q$ as in
Lemma~\ref{ptq:P}(c); \textup{(iv)} $E_{ba}\mathcal G'\subset\mathcal G$ for all polynomials $\mathcal G'$,
when $\varepsilon_a=\varepsilon_b$, $\lambda_a>\lambda_b$; \textup{(v)} $(E_{aa}-\Lambda_a)Q\in\mathcal G$ for
every polynomial $Q$. Then $\mathcal G$ contains all polynomials.
\end{proposition}
\begin{proof}
We show $\mathcal D_j\subset\mathcal G$ by induction on $j$ ($\mathcal D_{-1}=0$). Suppose
$\mathcal D_{j-1}\subset\mathcal G$ and $\mathcal D_j\not\subset\mathcal G$. Let $\mathcal N$ be the
orthogonal complement of $\mathcal G\cap\mathcal D_j$ in $\mathcal D_j$; $\mathcal N\ne0$.

\emph{$\mathcal N$ is stable under the compact column operators.} If $\varepsilon'_c=\varepsilon'_d$, then
$E'_{dc}$ preserves $\mathcal G\cap\mathcal D_j$ by (i) and degree, so its adjoint $E'_{cd}$ preserves
$\mathcal N$. The $E'_{cc}$ are self-adjoint and commute, so $\mathcal N$ is a finite-dimensional
module for the compact column algebra on which the Cartan subalgebra acts semisimply; it contains
a nonzero vector $Q$ that is an $E'_{cc}$-eigenvector for all $c$ and is killed by every $E'_{cd}$ with
$\varepsilon'_c=\varepsilon'_d$, $\lambda'_c>\lambda'_d$ (a weight vector whose weight is maximal for a linear
order compatible with these roots).

\emph{Row weight and raising.} By (v), $\mathcal G$ is $E_{aa}$-stable, so $\mathcal N$ is, and
$(E_{aa}-\Lambda_a)Q\in\mathcal N\cap\mathcal G=0$. If $\varepsilon_a=\varepsilon_b$, $\lambda_a>\lambda_b$, then
by (iv) $E_{ba}\mathcal D_j\subset\mathcal G\cap\mathcal D_j$, so $E_{ab}Q$, which lies in $\mathcal D_j$,
is orthogonal to $\mathcal D_j$, and $E_{ab}Q=0$.

\emph{Column Laplacians.} For $\varepsilon'_c\ne\varepsilon'_d$, $E'_{dc}\Delta'_{cd}Q\in\mathcal G$ by (i)
(as $\Delta'_{cd}Q\in\mathcal D_{j-2}\subset\mathcal G$), so $\Delta'_{cd}Q=0$ by
Lemma~\ref{ptq:fischer}.

\emph{Row Laplacians.} We apply Lemma~\ref{ptq:tensor}: let $\varepsilon_a\ne\varepsilon_b$ and suppose
$\Delta_{ab}Q$ has $T$-weight $\Lambda+e_a-e_b$ and is killed by the compact raising operators. If
$e_a-e_b$ is not a nonnegative combination of roots $e_x-e_z$ with $\lambda_x>\lambda_z$, then (iii)
gives $E_{ba}\Delta_{ab}Q\in\mathcal G$, and $\Delta_{ab}Q=0$ by Lemma~\ref{ptq:fischer}. Otherwise,
pairing with $\lambda$ shows $\lambda_a>\lambda_b$. All row parameters satisfy $|\xi_r|>\tfrac t2$
(every row is matched, \eqref{par:eq-match}), and $\xi_r-\lambda_r$ is constant; so either
$\xi_a>\tfrac t2$, or $\xi_b<\xi_a<-\tfrac t2$. By Lemma~\ref{par:relw}(i), $s\,w_r\ge0$ for
$s=\ptqsgn\xi_r$; the polynomial $Q$ is killed by the row raising operators and by the $\Delta'_{cd}$
and has $E_{aa}Q=\Lambda_aQ$, so Lemma~\ref{par:polar}(c) gives: in the first case
$\partial_{ac}Q=0$ whenever $\eta_{ac}=-1$; in the second $\partial_{bc}Q=0$ whenever $\eta_{bc}=1$. Every
term $\partial_{ac}\partial_{bc}Q$ of $\Delta_{ab}Q$ has $\eta_{ac}=-1$, $\eta_{bc}=1$, so
$\Delta_{ab}Q=0$. By Lemma~\ref{ptq:tensor}, all row Laplacians kill $Q$.

\emph{Conclusion.} $Q$ is a pre-joint harmonic (Definition~\ref{par:def-harm}), so $Q=c\Phi$ by
Proposition~\ref{par:unique}. By (ii), $Q\in\mathcal G\cap\mathcal N=0$, a contradiction.
\end{proof}

\begin{theorem}[(H)]\label{ptq:H}
Let $(S,T)$ be an ascending or descending pair of the triangle. For every polynomial $Q$ at $\tau$
there is $\beta_Q\in C_c^\infty(G')$ with $\mathsf PQ=\sigma'(\beta_Q)\mathsf P\Phi$. Consequently
\textup{(H)} of \S\ref{ptr:ss-mult} holds: for every polynomial vector $Q$ of the Fock space of
$\auHom(S,T)\otimes\R$ there is $\alpha_Q\in C_c(\auU(T)_\infty)$ with $\mathsf PQ=\int\alpha_Q(x)
\omega_\infty(1,x)\mathsf PP_{S,T}\,dx$.
\end{theorem}
\begin{proof}
Let $\mathcal G$ be the space of polynomials $Q$ such that $\mathsf PQ=\sigma'(\beta)\mathsf P\Phi$ for some
$\beta\in C_c^\infty(G')$. We check the hypotheses of Proposition~\ref{ptq:howe}.

(i) Let $Q\in\mathcal G$, $\mathsf PQ=\sigma'(\beta)\mathsf P\Phi$, and $Y\in\mathfrak g'$. As $\mathsf P$ is
bounded and commutes with $\sigma'$, and $Q$ is a smooth vector for $\sigma'$ with derivative
$d\sigma'(Y)Q=E'(Y)Q$, we get $\mathsf PE'(Y)Q=\frac d{dt}|_0\sigma'(\exp tY)\mathsf PQ$. And
$\sigma'(\exp tY)\sigma'(\beta)\mathsf P\Phi=\sigma'(\beta(\exp(-tY)\,\cdot))\mathsf P\Phi$, whose derivative at
$0$ is $\sigma'(\beta_Y)\mathsf P\Phi$ with $\beta_Y(x)=\frac d{dt}|_0\beta(\exp(-tY)x)\in C_c^\infty$
(differentiation under the integral, $\beta$ being smooth with compact support and $\sigma'$ unitary).
So $E'(Y)Q\in\mathcal G$ for $Y\in\mathfrak g'$, hence for $Y\in\mathfrak g'_\C$ by linearity.

(ii) Lemma~\ref{ptq:P}(d).

(iii), (iv), (v): for such polynomials $\mathsf P$ of them is $0=\sigma'(0)\mathsf P\Phi$, by
Lemma~\ref{ptq:P}(c) and (a) respectively.

So $\mathcal G$ contains all polynomials. For the second assertion: $\mathsf P=\bigotimes_\tau\mathsf
P_\tau$, $P_{S,T}=\bigotimes_\tau\Phi_\tau$, $\omega_\infty(1,x)=\bigotimes_\tau\sigma'_\tau(x_\tau)$, and every
polynomial vector is a finite sum of products $\bigotimes_\tau Q_\tau$; for a product take
$\alpha_Q(x)=\prod_\tau\beta_{Q_\tau}(x_\tau)$, and sum. The identity holds in the Hilbert space; both sides
lie in $\pzF_s$: $\mathsf P$ maps $\pzF_s$ into itself (the integral defining it converges absolutely in
$\pzF_s$ by Lemma~\ref{pz:growth}(a) and (A1), as in the proof of Proposition~\ref{pz:types}(b)), so
$\mathsf PQ,\mathsf PP_{S,T}\in\pzF_s$, and the Bochner integral of the continuous compactly supported
$\pzF_s$-valued function $x\mapsto\alpha_Q(x)\omega_\infty(1,x)\mathsf PP_{S,T}$ converges in $\pzF_s$ by
Lemma~\ref{pz:growth}(a); and $\pzF_s$ injects continuously into the Hilbert space, so the identity
holds in $\pzF_s$.
\end{proof}

\begin{remark}\label{ptq:rem-H}
Theorem~\ref{ptq:H} is (H) of \S\ref{ptr:ss-mult}, for all ascending and descending pairs, hence for
the pairs $(U_1,A)$, $(U_2,A)$ of every square of the triangle (Remark~\ref{ptr:rem-prev}). So
Lemma~\ref{ptr:conv-lift}, Theorem~\ref{ptr:square} and Corollary~\ref{ptr:tri-b} hold without the
hypothesis (H). The proof uses only that the parameters match with distinct entries, the
integrability of $e_\lambda$, $e_{\lambda'}$ and their derivatives, (V) and (D1)--(D5); it applies
verbatim to any pair satisfying these (for instance the edges of the grids,
Theorem~\ref{pgv:G3}(e)); for edges read with the groups exchanged use Theorem~\ref{ptw:Hc} and
Proposition~\ref{ptw:revedge}.
\end{remark}

\subsection{Dependencies of Propositions~\ref{au:inj} and~\ref{au:factor}}\label{ptq:ss-status}

Theorem~\ref{ptq:F} uses Lemma~\ref{ptq:vac} (proved here), (W1$_\infty$) and Theorem~\ref{pco:V}; Lemma~\ref{ptq:P}, Proposition~\ref{ptq:howe} and
Theorem~\ref{ptq:H} use (V) and (D1)--(D5); Theorem~\ref{ptr:square} and Corollary~\ref{ptr:tri-b} use
(D1)--(D5) and Theorem~\ref{pcx:X}. The finite parts of the descent characters and the comparison with
\S\ref{pc:ss-triangle} are treated in \S\ref{ptv:sec} (Propositions~\ref{ptv:chars}, \ref{ptv:arch},
\ref{ptv:tri}); the adjoint law is proved there from (RS), which is Theorem~\ref{ptw:RS}; the edge
hypothesis (E1) for the grid squares (Remark~\ref{ptr:rem-grid}) is Lemma~\ref{pgv:E1g}. Hence
Propositions~\ref{au:inj} and~\ref{au:factor} rest, besides the proofs written here, on Theorem~\ref{pcx:X}
and the published results (D1)--(D5), (W1$_\infty$), (W2$_\infty$), (M), (S).



%% ===== proofs/C-triangle3.tex =====
% proofs/C-triangle3.tex
% Continuation of proofs/C-triangle.tex (ptr:) and proofs/C-triangle2.tex (ptq:): the embedding of
% F_s in the Fock space, the characters of the descents, the reduction of the adjoint law
% (Theorem pc:tri(c)) to the square law of the reversed squares, and the final check of
% Lemma pc:arch and Theorem pc:tri. Labels of this file start with ptv:.
% Only amsmath/amssymb/amsthm are used. Conventions and macros: proofs/C-triangle.tex.
\providecommand{\auU}{\mathrm U}
\providecommand{\auHom}{\operatorname{Hom}}
\providecommand{\auGL}{\mathrm{GL}}
\providecommand{\auNm}{\operatorname{Nm}}
\providecommand{\auA}{\mathbb A}
\providecommand{\auO}{\mathcal O}
\providecommand{\auF}{\mathbb F}
\providecommand{\aubar}[1]{\overline{#1}}
\providecommand{\psrho}{\varrho}
\providecommand{\psHS}{\mathcal H}
\providecommand{\psvac}{\mathbf v}
\providecommand{\pzF}{\mathfrak F}
\providecommand{\pcF}{\mathcal F}
\providecommand{\pcG}{\mathsf G}
\providecommand{\pcm}{\mathfrak m}
\providecommand{\pcR}{R}
\providecommand{\pcMat}{\mathrm M}
\providecommand{\ptrwt}{\operatorname{wt}}
\providecommand{\ptvtr}{\operatorname{tr}}
\providecommand{\ptvF}{\mathsf F}

\section{The adjoint law and the clauses of the triangle}\label{ptv:sec}

This section continues \S\ref{ptr:sec} and \S\ref{ptq:sec}. Notation is that of
\S\ref{pc:ss-data}, \S\ref{pc:ss-triangle}, \S\ref{pc:ss-exchange} and \S\ref{ptr:sec}; inner
products are conjugate-linear in the first variable; for a function $f$ on a group, $\aubar f$ is
its complex conjugate.

\emph{Outline.}
\begin{enumerate}
\item[(1)] The Banach space $\pzF_s$ of \S\ref{pz:ss-types} injects continuously into the Fock
space (Lemma~\ref{ptv:embed}); this is the fact used in the last step of the proof of
Theorem~\ref{ptq:H}.
\item[(2)] The characters of the descents are fixed consistently by the squares and by
Theorem~\ref{pc:tri}(a) (Proposition~\ref{ptv:chars}): the compatibility of finite parts
stated after Lemma~\ref{ptr:chars} holds, as an identity of automorphic characters.
\item[(3)] The adjoint law, Theorem~\ref{pc:tri}(c), is reduced to the square law for the
\emph{reversed squares} (statement (RS) of \S\ref{ptv:ss-RS}), by complete arguments at the finite
places: reversal of a theta kernel modulo an ideal (Proposition~\ref{ptv:rev}), translation by
central elements (Lemma~\ref{ptv:central}), and a Gauss sum identity for the extended data
(Proposition~\ref{ptv:fourier}); (RS) itself is Theorem~\ref{ptw:RS}.
\item[(4)] Lemma~\ref{pc:arch} is proved for all its pairs, and every clause of
Theorem~\ref{pc:tri} is proved, clause (c) through (RS) (\S\ref{ptv:ss-check}).
\end{enumerate}
The edge hypothesis (E1) for the grid squares (Remark~\ref{ptr:rem-grid}) is Lemma~\ref{pgv:E1g}.

\subsection{The space $\pzF_s$ in the Fock space}\label{ptv:ss-embed}

Recall from \S\ref{pz:ss-types}: $\pzF$ is the Fock space (a Hilbert space), $\pzF^{\mathrm{sm}}$ its
space of smooth vectors, $\partial_w$ ($w\in W_\infty$) the derivative at $s=0$ of $\psrho_\infty(sw)$
on $\pzF^{\mathrm{sm}}$, $\|\Psi\|_s$ the maximum of the norms $\|\partial_{w_1}\cdots\partial_{w_n}\Psi\|$
over $0\le n\le 2s$ and $w_i$ in a fixed real basis $\mathsf B$, and $\pzF_s$ the completion of
$\pzF^{\mathrm{sm}}$ for $\|\cdot\|_s$. The definition asserts that this completion is ``contained
in $\pzF$''; this is the content of the following lemma.

\begin{lemma}\label{ptv:embed}
The identity of $\pzF^{\mathrm{sm}}$ extends to a continuous linear map $\iota\colon\pzF_s\to\pzF$ with
$\|\iota x\|\le\|x\|_s$, and $\iota$ is injective. Consequently a Bochner integral in $\pzF_s$ of a
continuous compactly supported (or $\|\cdot\|_s$-integrable) function is mapped by $\iota$ to the
Bochner integral in $\pzF$ of the same function, and two elements of $\pzF_s$ with the same image
in $\pzF$ are equal.
\end{lemma}
\begin{proof}
The term $n=0$ of the maximum gives $\|\Psi\|\le\|\Psi\|_s$ on $\pzF^{\mathrm{sm}}$, so the identity is
uniformly continuous and extends to $\iota$ with the stated bound. Let $x\in\pzF_s$ with $\iota x=0$
and let $x_k\in\pzF^{\mathrm{sm}}$ converge to $x$ in $\|\cdot\|_s$; so $x_k\to0$ in $\pzF$. Fix a word
$D=\partial_{w_1}\cdots\partial_{w_n}$ with $n\le2s$, $w_i\in\mathsf B$. The sequence $Dx_k$ is Cauchy in
$\pzF$, with a limit $y$. For $w\in W_\infty$ the operators $\psrho_\infty(sw)$, $s\in\R$, form a
strongly continuous one-parameter unitary group (on a line the cocycle of the Heisenberg group is
trivial), so $\psrho_\infty(sw)^*=\psrho_\infty(-sw)$ and, for $a,b\in\pzF^{\mathrm{sm}}$,
$\langle\partial_wa,b\rangle=-\langle a,\partial_wb\rangle$ (differentiate
$\langle\psrho_\infty(sw)a,b\rangle=\langle a,\psrho_\infty(-sw)b\rangle$ at $s=0$); and
$\pzF^{\mathrm{sm}}$ is stable under the $\partial_w$. Hence for every $b\in\pzF^{\mathrm{sm}}$
\[
\langle y,b\rangle=\lim_k\langle Dx_k,b\rangle=\lim_k(-1)^n\langle x_k,\partial_{w_n}\cdots\partial_{w_1}b
\rangle=0 .
\]
As $\pzF^{\mathrm{sm}}$ is dense in $\pzF$, $y=0$. So $\|Dx_k\|\to0$ for each of the finitely many
words $D$, i.e.\ $\|x_k\|_s\to0$, and $x=\lim x_k=0$. The statements on Bochner integrals follow
because a continuous linear map commutes with Bochner integrals.
\end{proof}

\begin{remark}\label{ptv:rem-embed}
This is exactly what the last step of the proof of Theorem~\ref{ptq:H} uses: the two sides of the
identity there are elements of $\pzF_s$ (the Bochner integral of $x\mapsto\alpha_Q(x)\omega_\infty(1,x)
\mathsf PP_{S,T}$ converges in $\pzF_s$ by Lemma~\ref{pz:growth}(a); $\mathsf PQ\in\pzF_s$ for the
same reason, $\mathsf P$ being a Bochner integral of $e_S(x)\omega_\infty(x^{-1},1)Q$), their images in
$\pzF$ coincide, and Lemma~\ref{ptv:embed} gives equality in $\pzF_s$.
\end{remark}

\subsection{The characters of the descents}\label{ptv:ss-chars}

Recall (\S\ref{pc:ss-data}) that $\theta_{\alpha,\beta}$ for $\auHom(S,T)$ carries
$\alpha(\det g)^{-1}\beta(\det h)$, $g\in\auU(S)$, $h\in\auU(T)$: $\alpha$ is the \emph{source}
character and $\beta$ the \emph{target} character. Write $M_k=r+kt$ for the rank of the nodes in
column $k$. The lift $a_{i,n}$ has source character $\mu_{i+n+1}\upsilon_{M_{n+1}}$ and target
character $\mu_{i+n+1}\upsilon_{M_n}$. The descent $d_{i,n}$ is a transposed lift for
$\auHom(\pcG_{i+1,n},\pcG_{i,n+1})$ with source character $\alpha^\downarrow_{i,n}$ (on
$\pcG_{i+1,n}$) and target character $\beta^\downarrow_{i,n}$ (on $\pcG_{i,n+1}$). For a square at
$(i,n)$ and a twist $\vartheta$ put $\gamma_n=\nu_{n+1}\vartheta^{\,n+1}\upsilon^{M_{n+1}}$; this is
$\gamma=\nu^{-1}\upsilon^{N/2}$ of \S\ref{ptr:ss-chars} ($N=2M_{n+1}$). Recall
$\mu_j=\mu^\flat\nu_1^{\,j}\vartheta^j$ (\S\ref{pc:ss-triangle}), so
$\mu_{j+1}=\mu_j\,\nu_1\vartheta$, and $\upsilon_k=\upsilon^{\lfloor k/2\rfloor}$.

In the square at $(i,n)$ (with $U_1=\pcG_{i,n+1}$, $U_2=\pcG_{i+1,n+1}$, $A=\pcG_{i,n+2}$,
$A'=\pcG_{i+1,n}$), condition (E2) of Definition~\ref{ptr:def-square} consists of four equations:
\begin{align*}
&\text{at }A:&\mu_{i+n+2}\upsilon_{M_{n+1}}\cdot\beta^\downarrow_{i,n+1}&=\gamma_n,&
&\text{at }A':&\alpha^\downarrow_{i,n}\cdot\mu_{i+n+2}\upsilon_{M_{n+1}}&=\gamma_n,\\
&\text{at }U_1:&\mu_{i+n+2}\upsilon_{M_{n+2}}\cdot\beta^\downarrow_{i,n}&=\gamma_n,&
&\text{at }U_2:&\alpha^\downarrow_{i,n+1}\cdot\mu_{i+n+2}\upsilon_{M_n}&=\gamma_n .
\end{align*}
(At $A$: target characters of $a_{i,n+1}$ and of $d_{i,n+1}$; at $A'$: source characters of
$d_{i,n}$ and of $a_{i+1,n}$; at $U_1$: source character of $a_{i,n+1}$ times target character of
$d_{i,n}$; at $U_2$: source character of $d_{i,n+1}$ times target character of $a_{i+1,n}$.) So each
square fixes all four characters of $d_{i,n}$ and $d_{i,n+1}$, and the characters of $d_{i,n+1}$
($n+1\ge1$) are fixed twice, by the squares at $(i,n)$ and at $(i,n+1)$; those of $d_{i,0}$ are
fixed by the square at $(i,0)$ and, for $\beta^\downarrow_{i,0}$, also by
Theorem~\ref{pc:tri}(a).

\begin{proposition}\label{ptv:chars}
\textup{(a)} For $i+n+2\le m$ define $\alpha^\downarrow_{i,n},\beta^\downarrow_{i,n},
\alpha^\downarrow_{i,n+1},\beta^\downarrow_{i,n+1}$ by the four equations of the square at $(i,n)$.
If also $i+n+3\le m$, the characters of $d_{i,n+1}$ so obtained coincide with those defined by the
square at $(i,n+1)$:
\[
\begin{gathered}\beta^\downarrow_{i,n+1}=\gamma_n\,(\mu_{i+n+2}\upsilon_{M_{n+1}})^{-1}=\gamma_{n+1}\,
(\mu_{i+n+3}\upsilon_{M_{n+3}})^{-1},\\
\alpha^\downarrow_{i,n+1}=\gamma_n\,(\mu_{i+n+2}\upsilon_{M_n})^{-1}=\gamma_{n+1}\,
(\mu_{i+n+3}\upsilon_{M_{n+2}})^{-1}.\end{gathered}
\]
So for each $\vartheta$ there is one family of characters $(\alpha^\downarrow_{i,n},
\beta^\downarrow_{i,n})_{i+n+1\le m}$ satisfying \textup{(E2)} in every square of the triangle (the
characters of $d_{i,m-i-1}$, which lies in no square as the lower edge, are fixed by the square at
$(i,m-i-2)$ when $m-i\ge2$; for $i=m-1$ the characters of $d_{m-1,0}$ are those given by
\textup{(W7)}, Remark~\ref{ptw:rem-W7}, and \textup{(b)} identifies its target character). Their weights are those of Definition~\ref{ptr:desc-char}, and their
tame parts at the primes of $\pcm_\vartheta$ are those of $\mu_{i+1}^{-1}$.

\textup{(b)} The target character of the descent $d_{i,0}$ defined by Theorem~\ref{pc:tri}(a)
(through \textup{(W7)}, Remark~\ref{ptr:rem-desc}) is $\beta^\downarrow_{i,0}$ of \textup{(a)}.
\end{proposition}
\begin{proof}
(a) In the quotient of the two expressions for $\beta^\downarrow_{i,n+1}$,
\[
\frac{\gamma_{n+1}}{\gamma_n}\cdot\frac{\mu_{i+n+2}}{\mu_{i+n+3}}\cdot\frac{\upsilon_{M_{n+1}}}
{\upsilon_{M_{n+3}}}=\nu_1\vartheta\,\upsilon^{t}\cdot(\nu_1\vartheta)^{-1}\cdot
\upsilon^{\lfloor M_{n+1}/2\rfloor-\lfloor M_{n+3}/2\rfloor}=\upsilon^{t-t}=1,
\]
because $\gamma_{n+1}/\gamma_n=\nu_1\vartheta\,\upsilon^{M_{n+2}-M_{n+1}}$, $\nu_{n+2}=\nu_{n+1}\nu_1$,
and $M_{n+3}=M_{n+1}+2t$, so $\lfloor M_{n+3}/2\rfloor=\lfloor M_{n+1}/2\rfloor+t$. The same computation
with $M_n,M_{n+2}$ in place of $M_{n+1},M_{n+3}$ gives the second identity. The weights:
$\ptrwt\mu_j=-\kappa_j$, $\ptrwt\nu_k=-k\Delta$, $\vartheta$ has weight $0$ and $\kappa_{j}+\Delta=
\kappa_{j+1}$; so $\ptrwt\beta^\downarrow_{i,n}=-(n+1)\Delta+M_{n+1}+\kappa_{i+n+2}-\lfloor M_{n+2}/2
\rfloor=\kappa_{i+1}+M_{n+1}-\lfloor(M_{n+1}+t)/2\rfloor$, and as $M_{n+1}+t$ and $M_n=M_{n+1}-t$ have
the same parity, $M_{n+1}-\lfloor(M_{n+1}+t)/2\rfloor=\lceil(M_{n+1}-t)/2\rceil=\lceil M_n/2\rceil$, which
is Definition~\ref{ptr:desc-char}; for $\alpha^\downarrow_{i,n}$ the computation is that of
Lemma~\ref{ptr:chars}. The tame parts: $\upsilon$ is unramified at all finite places, and
$\gamma_n\mu_{i+n+2}^{-1}=\nu_{n+1}\vartheta^{n+1}\mu_{i+n+2}^{-1}\upsilon^{M_{n+1}}=\mu_{i+1}^{-1}
\upsilon^{M_{n+1}}$, because $\mu_{i+n+2}=\mu_{i+1}(\nu_1\vartheta)^{n+1}$; so every character above
is $\mu_{i+1}^{-1}$ times a power of $\upsilon$.

(b) By (a), $\beta^\downarrow_{i,0}=\gamma_0(\mu_{i+2}\upsilon_{M_2})^{-1}=\mu_{i+1}^{-1}\upsilon^{k_0}$
for an integer $k_0$. On the other hand, by (W7) (Lemma~\ref{pc:SW}) the complex conjugate of the
kernel of $a_{i,0}$ is a kernel for $\auHom(\pcG_{i,0}^-,\pcG_{i,1})=\auHom(\pcG_{i+1,0},\pcG_{i,1})$
whose target character is the conjugate $\mu_{i+1}^{-1}\upsilon_{M_0}^{-1}$ of the target character
of $a_{i,0}$ times a power of $\upsilon$ (the character $\lambda$ of (W7) is a product of powers of
$\upsilon\circ\det$); so it is $\mu_{i+1}^{-1}\upsilon^{k_1}$ for an integer $k_1$ (explicitly $k_1=\lceil r/2\rceil$,
Remark~\ref{ptw:rem-W7}). Its weight is that of
Definition~\ref{ptr:desc-char} by Remark~\ref{ptr:rem-desc} (Proposition~\ref{ptr:match}(c)), and
so is the weight of $\beta^\downarrow_{i,0}$ by (a). As $\upsilon$ has weight $1$ at every
$\tau$, $k_0=k_1$, and the two characters are equal.
\end{proof}

\begin{remark}\label{ptv:rem-chars}
A priori the characters on $U_1$ and $U_2$ of the two kernels at $A$ need not agree, and the
discrepancy $\nu_0$ is carried by a factor $\nu_0(\det g_2)$ in the integrand; in
Theorem~\ref{ptr:square} this is the factor $\lambda(\det g_2)$. Proposition~\ref{ptv:chars} shows
that in the triangle no choice is involved: the characters of every descent are determined by any
one of the conditions that involve it, and all these determinations agree. With it,
Corollary~\ref{ptr:tri-b} needs no further input concerning finite parts.
\end{remark}

\subsection{Adjoints of lifts and translations by central elements}\label{ptv:ss-adj}

Throughout this subsection $(S,T)$ is a pair of spaces of ranks $N\le N'$, $\pcm$ a modulus,
$\alpha,\beta$ characters, $P$ a Fock vector with finite moments of $\auHom(S,T)\otimes\R$ and $\Psi$ a
function on $\pcMat_{N'\times N}(\pcR_\pcm)$ \emph{supported on injective matrices} (at every prime of
$\pcm$). We assume that $(S,T)$ with $(\alpha,\beta,P)$ is an edge in the sense of
Theorem~\ref{ptr:edge}(c), (d), from a type $e_S$ to a type $e_T$: the lift
$\Theta=\Theta_{\alpha,\beta}(P,\Psi)$ maps $\pcF(S,e_S)$ to $\pcF(T,e_T)$, the transposed lift
$\Theta^{\mathsf t}=\Theta^{\mathsf t}_{\alpha,\beta}(P,\Psi)$ maps forms of type $\aubar{e_T}$ to forms
of type $\aubar{e_S}$, and all these forms are bounded and rapidly decreasing.

\begin{lemma}[Adjoints]\label{ptv:adj}
For $z\in\pcF(T,e_T)$ put $\Theta^*z(g)=\int_{[\auU(T)]}\aubar{\theta_{\alpha,\beta}(g,h;P,\Psi)}\,z(h)
\,dh$. Then $\Theta^*z=\aubar{\Theta^{\mathsf t}\aubar z}\in\pcF(S,e_S)$, and
$\langle\Theta y,z\rangle=\langle y,\Theta^*z\rangle$ for all $y\in\pcF(S,e_S)$, the double integral
being absolutely convergent.
\end{lemma}
\begin{proof}
$\aubar z$ is a form of type $\aubar{e_T}$, so $\Theta^{\mathsf t}\aubar z$ is a form of type
$\aubar{e_S}$ and its conjugate, which is $\Theta^*z$, lies in $\pcF(S,e_S)$ (it is bounded,
continuous, invariant under an open subgroup because the datum is, and reproduced by $e_S$). By
Proposition~\ref{pcx:decay}(b) ($\Psi$ is supported on injective matrices), for $N_0$ given there
are $B,C$ with $|\theta_{\alpha,\beta}(g,h;P,\Psi)|\le C\,\mathrm{ht}(g)^{-N_0}\mathrm{ht}(h)^{B}$. As in
the proof of Proposition~\ref{ptr:conv}, $|z|\,\mathrm{ht}^B$ is integrable on $[\auU(T)]$, $y$ is
bounded and $\mathrm{ht}^{-N_0}$ is integrable on $[\auU(S)]$ for $N_0$ large ((R1) of
\S\ref{pcx:ss-conv}). So $\iint|\theta(g,h)\,y(g)\,z(h)|\,dg\,dh<\infty$, and Fubini's theorem gives
$\langle\Theta y,z\rangle=\iint\aubar{\theta(g,h)}\,\aubar{y(g)}\,z(h)\,dg\,dh=\langle y,\Theta^*z\rangle$.
\end{proof}

The same computation applies to a transposed lift: if $\Theta^{\mathsf t}$ maps $\pcF(T,\aubar{e_T}')$
to $\pcF(S,\aubar{e_S}')$ for types $e'$ and its datum is supported on injective matrices, its adjoint
is $(\Theta^{\mathsf t})^*y(h)=\int_{[\auU(S)]}\aubar{\theta(g,h)}\,y(g)\,dg$, with
$\langle\Theta^{\mathsf t}x,y\rangle=\langle x,(\Theta^{\mathsf t})^*y\rangle$.

\begin{lemma}[Central translations]\label{ptv:central}
Let $\zeta\in\auU(1)(\auA_{E^+,f})$ and write $\zeta$ also for the scalar element of $\auU(S)(\auA_f)$
and of $\auU(T)(\auA_f)$; it is central. There is $\eta\in\C$, $|\eta|=1$, depending on $\zeta$,
$\alpha$, $\beta$, $S$ and $T$ but not on $y$, with $\Theta(\psrho(\zeta)y)=\eta\,\psrho(\zeta)\Theta(y)$ for all
$y$, where $\psrho$ is right translation; the same holds for transposed lifts.
\end{lemma}
\begin{proof}
The pair $(\zeta,\zeta)$ acts on $\auHom(S,T)(\auA)$ by $x\mapsto\zeta x\zeta^{-1}=x$. By (W1) the operator
$\omega_{\auHom(S,T)}(\zeta,\zeta)$ commutes with all $\psrho(x)$, hence is a scalar, of modulus one as
$\omega$ is unitary; so $\omega_{\alpha,\beta}(\zeta,\zeta)=\eta_0$ with $|\eta_0|=1$, and as
$(\zeta,\zeta)$ is central, $\theta_{\alpha,\beta}(g\zeta,h\zeta;P,\Psi)=\eta_0\,\theta_{\alpha,\beta}
(g,h;P,\Psi)$. Since $\zeta$ is central, $g\mapsto g\zeta^{-1}$ preserves $[\auU(S)]$ and its measure,
and
\[
\Theta(\psrho(\zeta)y)(h)=\int\theta(g,h)\,y(g\zeta)\,dg=\int\theta(g\zeta^{-1},h)\,y(g)\,dg=
\aubar{\eta_0}\int\theta(g,h\zeta)\,y(g)\,dg=\aubar{\eta_0}\,\Theta(y)(h\zeta),
\]
using $\theta(g\zeta^{-1},h)=\aubar{\eta_0}\theta(g,h\zeta)$. Put $\eta=\aubar{\eta_0}$. For transposed
lifts exchange the roles of $g$ and $h$.
\end{proof}

\subsection{Reversal of a theta kernel modulo an ideal}\label{ptv:ss-rev}

The adjoint of a descent is a lift whose kernel is the complex conjugate of a kernel of
$\auHom(\pcG_{i+1,n},\pcG_{i,n+1})$, viewed as a function on $\auU(\pcG_{i,n+1})\times
\auU(\pcG_{i+1,n})$. We write it as a kernel of the reversed space $\auHom(\pcG_{i,n+1},\pcG_{i+1,n})$.

\begin{lemma}[Invariant pairings]\label{ptv:pairing}
Let $F$ be a finite field, $p,q\ge1$, and $e\colon\pcMat_{p\times q}(F)\times\pcMat_{q\times p}(F)\to
\C^\times$ bi-additive, non-degenerate (for each $\rho\ne0$, $e(\rho,\cdot)\ne1$, and for each $v\ne0$,
$e(\cdot,v)\ne1$) and invariant: $e(g\rho h^{-1},hvg^{-1})=e(\rho,v)$ for $g\in\auGL_p(F)$,
$h\in\auGL_q(F)$. Then there is a non-trivial character $\psi$ of $F$ with
$e(\rho,v)=\psi(\ptvtr(\rho v))$.
\end{lemma}
\begin{proof}
Fix a non-trivial character $\psi_1$ of $F$. Every character of $\pcMat_{p\times q}(F)$ is
$\rho\mapsto\psi_1(\ptvtr(\rho X))$ for a unique $X\in\pcMat_{q\times p}(F)$ (the pairing
$(\rho,X)\mapsto\psi_1(\ptvtr\rho X)$ is non-degenerate, and both groups have the same order). So
$e(\rho,v)=\psi_1(\ptvtr(\rho\,\mu(v)))$ for a map $\mu$ of $\pcMat_{q\times p}(F)$ to itself, additive by
bi-additivity, injective by non-degeneracy. Invariance gives $\mu(hvg^{-1})=h\mu(v)g^{-1}$. The
invertible matrices span $\pcMat_k(F)$ additively for every $k\ge1$ and every $F$ (for $k=1$ because
$1$ generates $\mathbb F_p\subset F$ and $F^\times$ spans $F$; for $k\ge2$ every elementary matrix
$E_{ij}$ is a difference of two invertible matrices); so $\mu(XvY)=X\mu(v)Y$ for all $X\in\pcMat_q(F)$,
$Y\in\pcMat_p(F)$. With $v=E_{ij}$ and $X=E_{ii}$, $Y=E_{jj}$, $\mu(E_{ij})=E_{ii}\mu(E_{ij})E_{jj}$ is a
multiple $c_{ij}E_{ij}$; with $X=E_{ki}$, $Y=E_{jl}$, $c_{kl}=c_{ij}=:c$, and with $X=a1$ ($a\in F$),
$\mu(av)=a\mu(v)$; so $\mu(v)=cv$, $c\ne0$. Put $\psi(x)=\psi_1(cx)$.
\end{proof}

For spaces $S,T$ let $S^-$ be $S$ with the negated form. By (W7) there is a conjugate-linear isometry
$\mathsf C\colon\psHS(\auHom(S,T))\to\psHS(\auHom(S^-,T))$ with $\mathsf C\psrho(x)=\psrho^-(x)\mathsf C$
and $\theta^-\circ\mathsf C=\aubar\theta$; by Proposition~\ref{pf:flip}, applied to
$\auHom(S^-,T)$, the map $w\mapsto w^\dagger$ identifies $\auHom(S^-,T)'$ with $\auHom(T,S)$ and gives
$\mathrm{Fl}$ with $\mathrm{Fl}\,\psrho^-(x)=\psrho'(x^\dagger)\mathrm{Fl}$ and
$\theta'\circ\mathrm{Fl}=c_\theta\theta^-$. Put $\mathcal A=\mathrm{Fl}\circ\mathsf C$ and
$\phi(x)=x^\dagger$, where $\dagger$ is the adjoint for the forms of $S$ and $T$ (the adjoint for
$S^-$ and $T$ followed by the identification is $x\mapsto-x^\dagger$; the sign is absorbed in
$\phi$ and changes nothing below). Then $\mathcal A$ is conjugate-linear, isometric (the local maps are
normalised to be isometric below; the normalisation constant of $\mathrm{Fl}$ is absorbed in
$c_\theta$), and
\begin{equation}\label{ptv:eq-A}
\mathcal A\,\psrho(x)=\psrho'(\phi(x))\,\mathcal A,\qquad\theta'\circ\mathcal A=c_\theta\,\aubar\theta,
\qquad\mathcal A\,\omega_{\auHom(S,T)}(g,h)=\lambda_S(\det g)\lambda_T(\det h)^{-1}\,
\omega_{\auHom(T,S)}(h,g)\,\mathcal A,
\end{equation}
where $\lambda_S,\lambda_T$ are the products of the character of (W7) with the character
$\upsilon(\det g)^{N'}\upsilon(\det h)^{-N}$ of Proposition~\ref{pf:flip} ($\det$ the
determinants on $S$ and $T$); they are powers of $\upsilon$ depending only on $N,N'$, and in fact both
are trivial by Lemma~\ref{ptw:W7}. We keep them in the formulas below and use $\lambda_S=\lambda_T=1$
from \S\ref{ptv:ss-RS} on. As in \S\ref{pcx:ss-flip} we take
$\mathrm{Fl}$, $\mathsf C$, hence $\mathcal A$, as tensor products of local maps $\mathcal A_v$, each
conjugate-linear, isometric and intertwining the local Heisenberg operators through $\phi$, with
the normalisation at the finite places fixed below; then $c_\theta\ne0$ depends only on $S$, $T$
and the frames.

Let $\zeta_\alpha=\aubar\alpha/\alpha\in\auU(1)(\auA_{E^+,f})$, with $\alpha$ the finite idele of
\S\ref{pc:ss-data} (the vacuum of the datum is that of the lattice $\alpha\auHom(L_S,L_T)$).

\begin{proposition}[Reversal]\label{ptv:rev}
For every modulus $\pcm$ there are a character $\psi_\pcm$ of the additive group of $\pcR_\pcm$,
non-trivial on every factor $\auF_w$, a scalar $\zeta_\pcm\in\auU(1)(\auA_{E^+,f})$ equal to
$\zeta_\alpha$ outside the primes of $\pcm$, both independent of $S,T$, and for all $\alpha,\beta$ a
number $c=c(S,T,\pcm,\alpha)$ with $|c|^2=C_{S,T}\,(\#\pcR_\pcm)^{-NN'}$, $C_{S,T}>0$ depending only on
$S,T$, such that for every Fock vector $P$, every $\Psi$ on $\pcMat_{N'\times N}(\pcR_\pcm)$ and all
$g\in\auU(S)(\auA)$, $h\in\auU(T)(\auA)$
\[
\aubar{\theta_{\alpha,\beta}(g,h;P,\Psi)}=c\,\theta_{\beta\lambda_T,\,\alpha\lambda_S}(h,g\zeta_\pcm;
P^\flat,\Psi^\flat),\qquad\Psi^\flat(\varrho)=\sum_{v\in\pcMat_{N'\times N}(\pcR_\pcm)}\aubar{\Psi(v)}\,
\psi_\pcm(\ptvtr(\varrho v)),
\]
where the right side is a kernel of $\auHom(T,S)$, $\varrho\in\pcMat_{N\times N'}(\pcR_\pcm)$, and
$P^\flat=\mathcal A_\infty P$, a Fock vector with the same moments as $P$.
\end{proposition}
\begin{proof}
\emph{Places outside $\pcm$.} Let $v$ be a finite place, $L=\alpha\auHom(L_S,L_T)_v$. By (W4) the
vacuum $\psvac_L$ spans the line of vectors on which $\psrho(l)$, $l\in L$, acts by $\epsilon_L$; by
\eqref{ptv:eq-A}, $\mathcal A_v\psvac_L$ spans the line on which $\psrho'(l')$, $l'\in\phi(L)$, acts by
$\aubar{\epsilon_L}\circ\phi^{-1}$. Now $\phi(\alpha x)=\aubar\alpha\phi(x)$ and
$\auHom(L_S,L_T)^\dagger=\auHom(L_T,L_S)$ (the lattices are unimodular), so $\phi(L)=\zeta_\alpha L'$,
$L'=\alpha\auHom(L_T,L_S)_v$, and the scalar $\zeta_\alpha$ acts on $\auHom(T,S)$ as the element
$(1,\zeta_\alpha)$ of $\auU(T)\times\auU(S)$ (by $x'\mapsto\zeta_\alpha x'$). By (W4), $\omega'(1,\zeta_\alpha)
\psvac_{L'}$ spans the line of $\zeta_\alpha L'$; we normalise $\mathcal A_v$ by
$\mathcal A_v\psvac_L=\omega'_v(1,\zeta_{\alpha,v})\psvac_{L'}$ (both are unit vectors). At almost all
$v$, $\zeta_{\alpha,v}$ is a unit and $\omega'_v(1,\zeta_{\alpha,v})\psvac_{L'}=\psvac_{L'}$ (W4), so the
local maps form a restricted tensor product.

\emph{Places of $\pcm$.} Let $v=w\aubar w$, $w\mid\pcm$, $q=\#\auF_w$, $\varpi$ a uniformiser. In the
model (W5), $\psHS_v(\auHom(S,T))$ is the space of Schwartz--Bruhat functions on
$X=\pcMat_{N'\times N}(E_w)$; $\psrho(a)$, $a\in X$, is translation by $a$, and $\psrho(a')$, $a'$ in
the $\aubar w$-component $X^*$, is multiplication by $x\mapsto[x,a']$, a character which, for $a'$
integral, is trivial on $M_w=a_w\pcMat(\auO_w)$ exactly when $a'\in M_{\aubar w}$, the dual lattice
(Lemma~\ref{pgx:vac}; $M_w=\pcMat(\auO_w)$ and $M_{\aubar w}=\pcMat(\auO_{\aubar w})$ when $\alpha$ is a unit at
$w,\aubar w$; in general, with the data of Definition~\ref{pgx:conv}, read $\pcMat(\auO_w)$ as $M_w$ and
$\pcMat(\auO_{\aubar w})$ as $M_{\aubar w}$ in the rest of this paragraph), and which identifies
$\varpi^{-1}\pcMat(\auO_{\aubar w})/\pcMat(\auO_{\aubar w})$ with the dual of
$\pcMat(\auO_w)/\varpi\pcMat(\auO_w)$ (\eqref{pl:eq-heis}: $\lambda_w$ is trivial on $\auO_w$ and not on
$\varpi^{-1}\auO_w$). So the space $D$ of the functions $\Psi^\natural$ ($\Psi$ on
$\pcMat_{N'\times N}(\auF_w)$) is the space of vectors fixed by $\psrho(\Lambda)$,
$\Lambda=\varpi\pcMat(\auO_w)\oplus\pcMat(\auO_{\aubar w})$, and the operators $\psrho(a)$,
$a\in\pcMat(\auO_w)$, and $\psrho(a')$, $a'\in\varpi^{-1}\pcMat(\auO_{\aubar w})$, act on $D$ by
translations $\tau_a$ and by all the characters $m_b$ of $\pcMat_{N'\times N}(\auF_w)$. The adjoint
exchanges the components at $w$ and $\aubar w$ (it involves the conjugation of $E_v$), so $\phi$ maps
$X$ onto the $\aubar w$-component $X'^*$ of $\auHom(T,S)_v$ and $X^*$ onto $X'=\pcMat_{N\times N'}(E_w)$,
and it is equivariant: $\phi(hxg^{-1})=g\phi(x)h^{-1}$ for unitary $g,h$. The only lattices of $X'$ (of
$X'^*$) stable under $\auGL_N(\auO_w)\times\auGL_{N'}(\auO_w)$ are the $\varpi^k\pcMat(\auO)$ (such a
lattice is an $\auO$-module stable under left and right multiplication by matrix rings, hence a
module over $\pcMat_N(\auO)\otimes\pcMat_{N'}(\auO)^{\mathrm{op}}$ inside $\pcMat_{N\times N'}(E_w)$); and
$\phi$ is an isometry for $[\,,]$. Hence $\phi(\Lambda)=\varpi^{k}\pcMat(\auO_w)\oplus\varpi^{1-k}
\pcMat(\auO_{\aubar w})$ for an integer $k$ (depending only on $\alpha_w$, the frames being unimodular),
and $\mathcal A_vD$ is the space of functions on $X'$ that are $\varpi^k\pcMat(\auO_w)$-periodic and
supported on $\varpi^{k-1}\pcMat(\auO_w)$. Let $\zeta_{\pcm,v}\in\auU(1)(E^+_v)$ be the scalar whose
$w$-component moves $\varpi^{k-1}\pcMat(\auO_w)$ to $\pcMat(\auO_w)$ under the action of $(1,\zeta)$ in
(W5); then $T_v=\omega'_v(1,\zeta_{\pcm,v})^{-1}\mathcal A_v$ maps $D$ bijectively onto the space $D'$ of the
functions $\Psi'^\natural$, $\Psi'$ on $\pcMat_{N\times N'}(\auF_w)$ (as $(1,\zeta)$ normalises the
Heisenberg group, $T_v$ intertwines $\psrho$ and $\psrho'\circ\phi_\zeta$ with $\phi_\zeta$ the
composite of $\phi$ and the action of $(1,\zeta)^{-1}$). Through $T_v$ the translations $\tau_a$ on $D$
become multiplications by characters $\rho\mapsto e(\rho,a)$ on $D'$, and the multiplications $m_b$
become translations; $e$ is bi-additive, non-degenerate ($T_v$ is bijective) and invariant under
$\auGL_N(\auF_w)\times\auGL_{N'}(\auF_w)$ (by the equivariance of $\phi$, of the reduction maps of
$K_S$, $K_T$ (Lemma~\ref{pc:level}) and the commutation of $(1,\zeta)$ with $\omega'$). By
Lemma~\ref{ptv:pairing}, $e(\rho,a)=\psi_w(\ptvtr(\rho a))^{-1}$ for a non-trivial character $\psi_w$
of $\auF_w$. The map $\ptvF\colon\Psi\mapsto\Psi^\flat$ of the statement (at $w$) satisfies
$\ptvF(\tau_a\Psi)(\rho)=\psi_w(\ptvtr\rho a)^{-1}\ptvF(\Psi)(\rho)$, the same rule. For the
multiplications: if $T_vm_b=\tau_{b'}T_v$, the commutation relation $\tau_am_b=m_b(a)m_b\tau_a$ and the
conjugate-linearity of $T_v$ give $e(b',a)=m_b(a)$ for all $a$, which determines $b'$; the same
computation for $\ptvF$ gives the same $b'$. So $U=\ptvF^{-1}T_v$ is a linear map of $D$ commuting with
all $\tau_a$ and all $m_b$. The operators commuting with all $m_b$ are multiplications by functions,
and such a multiplication commutes with all translations only if it is constant; so $U$ is a scalar
$\varepsilon_v$, i.e.\ $T_v(\Psi^\natural)=\varepsilon_v(\Psi^\flat)^\natural$. Norms: in the
normalisation of \S\ref{pc:ss-data}, $\|\Psi^\natural\|^2=q^{-NN'}\sum|\Psi|^2$, and by finite
Plancherel $\sum_\rho|\Psi^\flat(\rho)|^2=q^{NN'}\sum_v|\Psi(v)|^2$; $\mathcal A_v$ and $\omega'_v$ preserve
norms; so $|\varepsilon_v|^2=q^{-NN'}$. We normalise $\mathcal A_v$ at $w$ by $\mathcal A_v\psvac=
\omega'_v(1,\zeta_{\pcm,v})u'$, where $u'$ is the unit vector proportional to
$(\mathbf 1^\flat)^\natural=q^{NN'}\delta_0^\natural$ (the characteristic function of $\varpi\pcMat(\auO_w)$, i.e.\ the vacuum of
$\varpi\pcMat(\auO_w)\oplus\varpi^{-1}\pcMat(\auO_{\aubar w})$): the vacuum is $\mathbf 1^\natural$ up to its norm, its image
lies on the line of $\omega'_v(1,\zeta_{\pcm,v})(\mathbf 1^\flat)^\natural$ because $T_v(\mathbf 1^\natural)=\varepsilon_v(\mathbf 1^\flat)^\natural$, so this is possible; this
changes $\varepsilon_v$ by a number of modulus one. Put $\psi_\pcm=\prod_w\psi_w$ on
$\pcR_\pcm=\prod_w\auF_w$; it is non-trivial on every factor.

\emph{Globally.} Put $\zeta_\pcm=\prod_v\zeta_{\pcm,v}$ with $\zeta_{\pcm,v}=\zeta_{\alpha,v}$ for $v$ not
below $\pcm$. By the local statements, $\mathcal A[P,\Psi]=\varepsilon\,\omega'(1,\zeta_\pcm)
[P^\flat,\Psi^\flat]$ with $\varepsilon=\prod_{w\mid\pcm}\varepsilon_w$, $|\varepsilon|^2=(\#\pcR_\pcm)^{-NN'}$.
By \eqref{ptv:eq-A},
\begin{align*}
\aubar{\theta_{\alpha,\beta}(g,h;P,\Psi)}&=\alpha(\det g)\beta(\det h)^{-1}\,\aubar{\theta(\omega(g,h)
[P,\Psi])}=c_\theta^{-1}\alpha(\det g)\beta(\det h)^{-1}\theta'(\mathcal A\,\omega(g,h)[P,\Psi])\\
&=c_\theta^{-1}(\alpha\lambda_S)(\det g)(\beta\lambda_T)(\det h)^{-1}\,\varepsilon\,\theta'\bigl(
\omega'(h,g)\,\omega'(1,\zeta_\pcm)[P^\flat,\Psi^\flat]\bigr)\\
&=c_\theta^{-1}\varepsilon\,(\alpha\lambda_S)(\det\zeta_\pcm)^{-1}\,\theta_{\beta\lambda_T,\alpha\lambda_S}
(h,g\zeta_\pcm;P^\flat,\Psi^\flat).
\end{align*}
So $c=c_\theta^{-1}\varepsilon\,(\alpha\lambda_S)(\det\zeta_\pcm)^{-1}$ and $|c|^2=|c_\theta|^{-2}
(\#\pcR_\pcm)^{-NN'}$. The archimedean map $\mathcal A_\infty$ is induced by the real-linear isometry
$\phi$ followed by complex conjugation, which preserve the moments (\S\ref{pcx:ss-flip}). The
integer $k$ and the scalar $\zeta_{\pcm}$ depend only on $\alpha$ and $\lambda_w$, not on $S$, $T$.
Neither does $\psi_w$: for $a\in X$ the operator $\mathcal A_v\psrho(a)\mathcal A_v^{-1}=\psrho'(\phi(a))$
is the multiplication by $x'\mapsto[x',\phi(a)]'$ on the functions on $X'$, and by the definition of the
bracket (\S\ref{ps:ss-weil}: $[x,y]=\psi(\tfrac12\jmath(\langle y,x\rangle-\langle x,y\rangle))$ with
$\langle w_1,w_2\rangle=\ptvtr(w_1^\dagger w_2)$, for every pair of spaces) and $\phi(a)^\dagger=a$,
$[x',\phi(a)]'=\psi_v\bigl(\tfrac12\jmath(\ptvtr(ax')-\aubar{\ptvtr(ax')})\bigr)$ is the value at
$\ptvtr(ax')\in E_w$ of a character of $E_w$ that depends only on $\psi_v$ and $\jmath$; conjugation by
$\omega'_v(1,\zeta_{\pcm,v})$ replaces $x'$ by a fixed scalar multiple, and the reductions modulo $\pcm$ are
taken in the coordinates $a_w^{-1}x$, $a_w^{-1}x'$ for all spaces. So $e(\rho,a)$ is the value at
$\ptvtr(\rho a)$ of one character of $\auF_w$, the same for all $S$, $T$.
\end{proof}

\begin{corollary}[Adjoints of descents]\label{ptv:desc-adj}
Let $(S,T)=(\pcG_{i+1,n},\pcG_{i,n+1})$ and let $d=\Theta^{\mathsf t}_{\alpha,\beta}(P,\varepsilon)\colon
\pcF_{i,n+1}\to\pcF_{i+1,n}$ be a descent, with $\varepsilon$ supported on injective matrices. Its
adjoint $d^*\colon\pcF_{i+1,n}\to\pcF_{i,n+1}$, $d^*y(h)=\int_{[\auU(S)]}\aubar{\theta(g,h)}y(g)\,dg$, is
\[
d^*=c\cdot R\circ\psrho(\zeta_\pcm^{-1}),\qquad R=\Theta^{\mathsf t}_{\beta\lambda_T,\alpha\lambda_S}
(P^\flat,\varepsilon^\flat),
\]
where $R$, the \emph{reversed lift}, is the transposed lift of the space $\auHom(T,S)$ (integration
over $[\auU(S)]$), and $c$ is the number of Proposition~\ref{ptv:rev}.
\end{corollary}
\begin{proof}
By Proposition~\ref{ptv:rev}, $d^*y(h)=c\int\theta_{\beta\lambda_T,\alpha\lambda_S}(h,g\zeta_\pcm)y(g)\,dg$,
and the substitution $g\mapsto g\zeta_\pcm^{-1}$ (central, measure preserving) gives
$c\int\theta(h,g)\,y(g\zeta_\pcm^{-1})\,dg$. That $d^*$ is the adjoint of $d$ and takes values in
$\pcF_{i,n+1}$ is the remark after Lemma~\ref{ptv:adj} with Theorem~\ref{ptr:edge}(d) for the
descending pair (the conjugate of the lift of $\aubar y$ along it).
\end{proof}

\subsection{A Gauss sum identity for extended data}\label{ptv:ss-fourier}

In this subsection $\pcR$ is a finite product of finite fields $\auF_w$, $\psi$ a character of
$(\pcR,+)$ non-trivial on every factor, and $\chi$ a character of $\pcR^\times$ non-trivial on every
factor $\auF_w^\times$, extended by $0$ to non-units. For finite sets $\mathsf I,\mathsf J$ and a
function $f$ on $\pcMat_{\mathsf I\times\mathsf J}(\pcR)$ put
$\ptvF f(\rho)=\sum_vf(v)\psi(\ptvtr(\rho v))$, $\rho\in\pcMat_{\mathsf J\times\mathsf I}(\pcR)$. Since
$(\rho,v)\mapsto\psi(\ptvtr\rho v)$ is a non-degenerate pairing, $\sum_s\psi(\ptvtr(Cs))=
\#\pcMat\cdot[C=0]$ and $\sum_\rho|\ptvF f(\rho)|^2=\#\pcMat_{\mathsf I\times\mathsf J}(\pcR)\sum_v|f(v)|^2$.

\begin{lemma}[Gauss sum of $\chi\circ\det$]\label{ptv:gauss}
Let $\mathsf T$ be a finite set and $G=\sum_{d\in\auGL_{\mathsf T}(\pcR)}\chi(\det d)\psi(\ptvtr d)$. Then
$|G|^2=(\#\pcR)^{|\mathsf T|^2}$, and every function $S$ on $\pcMat_{\mathsf T}(\pcR)$ with
$S(gd)=\chi(\det g)S(d)$ for all $g\in\auGL_{\mathsf T}(\pcR)$ satisfies
$\ptvF S(y)=G\,\chi(\det y)^{-1}S(1)$ for all $y$ (with $\chi(\det y)^{-1}=0$ if $\det y$ is not a
unit).
\end{lemma}
\begin{proof}
If $d$ is not invertible, some component $d_w$ is singular; choose $u\ne0$ with ${}^tu\,d_w=0$ and,
for $t\in\auF_w^\times$, $v$ with ${}^tuv=t-1$; then $g_w=1+v\,{}^tu$ has $g_wd_w=d_w$ and $\det g_w=t$.
With $g=g_w$ at $w$ and $1$ elsewhere, $S(d)=\chi_w(t)S(d)$ for all $t$, so $S(d)=0$. For invertible $d$,
$S(d)=\chi(\det d)S(1)$. Hence $\ptvF S(y)=S(1)\sum_{d\in\auGL}\chi(\det d)\psi(\ptvtr(yd))$. If $y$ is
invertible, the substitution $d=y^{-1}d'$ gives $\chi(\det y)^{-1}G$. If $y_w$ is singular, choose
$g_w$ as above with $y_wg_w=y_w$ (right kernel) and $\det g_w=t$; the substitution $d\mapsto gd$
multiplies the sum by $\chi_w(t)$, so the sum vanishes. For $|G|$: take $S_0=\chi\circ\det$ on
$\auGL$, $0$ elsewhere; Plancherel gives $|G|^2\#\auGL=(\#\pcR)^{|\mathsf T|^2}\#\auGL$.
\end{proof}

\subsubsection*{The extension of a datum}
Let $\mathsf N$, $\mathsf T$ be finite sets, $\mathsf M=\mathsf N\sqcup\mathsf T$ (in the triangle:
$\mathsf N$ indexes the nodes of column $n$, $\mathsf M$ those of column $n+1$, $\mathsf T$ the new
block). For $p$ on $\pcMat_{\mathsf M\times\mathsf N}(\pcR)$ its extension $p^+$ is the function on
$\pcMat_{(\mathsf M\sqcup\mathsf T)\times(\mathsf N\sqcup\mathsf T)}(\pcR)$ with
\[
p^+\begin{pmatrix}\eta&s\\0&1\end{pmatrix}=\frac{p(-\eta)}{\#\pcMat_{\mathsf M\times\mathsf T}(\pcR)}
\qquad(\eta\in\pcMat_{\mathsf M\times\mathsf N},\ s\in\pcMat_{\mathsf M\times\mathsf T}),
\]
and $p^+=0$ at all other matrices; this is the function $p^+$ of \S\ref{pc:ss-exchange} (Lemma~\ref{pc:next}),
for $\mathsf N_1=\mathsf M$, $\mathsf N_0=\mathsf N$, $\mathsf T'=\mathsf T$. Note
$\aubar{p^+}=(\aubar p)^+$. If $p$ is supported on injective matrices, so is $p^+$.

\begin{lemma}\label{ptv:F-ext}
For $\rho=\begin{pmatrix}A&B\\C&D\end{pmatrix}$ with $A\in\pcMat_{\mathsf N\times\mathsf M}$,
$B\in\pcMat_{\mathsf N\times\mathsf T}$, $C\in\pcMat_{\mathsf T\times\mathsf M}$, $D\in\pcMat_{\mathsf T}$:
$\ptvF(p^+)(\rho)=[C=0]\,\ptvF p(-A)\,\psi(\ptvtr D)$.
\end{lemma}
\begin{proof}
$\ptvtr\bigl(\rho\bigl(\begin{smallmatrix}\eta&s\\0&1\end{smallmatrix}\bigr)\bigr)=\ptvtr(A\eta)+
\ptvtr(Cs)+\ptvtr D$; sum over $s$ (giving $\#\pcMat\cdot[C=0]$) and over $\eta$ (giving
$\sum_\eta p(-\eta)\psi(\ptvtr A\eta)=\ptvF p(-A)$).
\end{proof}

\subsubsection*{The reversed square law}
Let $\mathsf N,\mathsf T,\mathsf M$ be as above, $\mathsf T'$ a copy of $\mathsf T$, $\mathsf M'=\mathsf M
\sqcup\mathsf T'$. Put $\mathsf U=\mathsf N\sqcup\mathsf M'=\mathsf N\sqcup\mathsf M\sqcup\mathsf T'$,
$\mathsf B=\mathsf M$ (the middle block) and $\mathsf A=\mathsf N\sqcup\mathsf T'$, so that
$\mathsf U=\mathsf B\sqcup\mathsf A$ and $|\mathsf A|=|\mathsf B|$; we identify $\mathsf A$ with
$\mathsf M=\mathsf N\sqcup\mathsf T$ through $\mathsf T'\cong\mathsf T$. Let $L$ on
$\pcMat_{\mathsf M\times\mathsf U}(\pcR)$ and $L'$ on $\pcMat_{\mathsf U\times\mathsf M}(\pcR)$ satisfy the
square law of \S\ref{pc:ss-exchange} for $(\mathsf B,\mathsf A)$ with constant $c$, and assume that
$L(kx)=\chi(\det k)L(x)$ for all $k\in\auGL_{\mathsf M}(\pcR)$ and \emph{all} $x$. Let
$\iota\in\pcMat_{\mathsf M\times\mathsf N}$ and $\iota'\in\pcMat_{\mathsf M'\times\mathsf M}$ be the
inclusion matrices; for $y\in\pcMat_{\mathsf N\times\mathsf M}$ write $\binom y{\iota'}\in
\pcMat_{\mathsf U\times\mathsf M}$, and for $x\in\pcMat_{\mathsf M\times\mathsf M'}$ write $(\iota\ x)\in
\pcMat_{\mathsf M\times\mathsf U}$.

\begin{proposition}\label{ptv:fourier}
There is $\gamma$ with $|\gamma|^2=(\#\pcR)^{|\mathsf T|(|\mathsf T|+2|\mathsf N|)}$, independent of
$L,L',c,p$, such that for every $p$ on $\pcMat_{\mathsf M\times\mathsf N}(\pcR)$
\[
\gamma\sum_{y\in\pcMat_{\mathsf N\times\mathsf M}}\ptvF(\aubar p)(y)\,L'\tbinom y{\iota'}=c\sum_{x\in
\pcMat_{\mathsf M\times\mathsf M'}}\ptvF(\aubar{p^+})(x)\,L(\iota\ x).
\]
\end{proposition}
\begin{proof}
\emph{Step 1.} For $q$ on $\pcMat_{\mathsf N\times\mathsf M}$ we show
$\sum_yq(y)L'\binom y{\iota'}=c\sum_xq^+(x)L(\iota\ x)$, where $q^+$ is the extension for the sets
$(\mathsf N,\mathsf M,\mathsf T)$ in the roles of $(\mathsf M,\mathsf N,\mathsf T)$:
$q^+\bigl(\begin{smallmatrix}\eta&s\\0&1\end{smallmatrix}\bigr)=q(-\eta)/\#\pcMat_{\mathsf N\times
\mathsf T}$, $\eta\in\pcMat_{\mathsf N\times\mathsf M}$, $s\in\pcMat_{\mathsf N\times\mathsf T'}$, rows
$\mathsf N\sqcup\mathsf T$ ($=\mathsf M$), columns $\mathsf M\sqcup\mathsf T'$. For $y\in\pcMat_{\mathsf N
\times\mathsf M}$ let $w\in\pcMat_{\mathsf A\times\mathsf B}$ have rows $-y$ (on $\mathsf N$) and $0$ (on
$\mathsf T'$). In the decomposition $\mathsf U=\mathsf B\sqcup\mathsf A$, $\binom y{\iota'}$ has the
block $1$ on the rows $\mathsf B$ and $-w$ on the rows $\mathsf A$, i.e.\ it is $\binom1{-w}$; and
$(w\ 1)\in\pcMat_{\mathsf A\times\mathsf U}$, read with rows $\mathsf M$, is $(\iota\ x_y)$ with
$x_y=\bigl(\begin{smallmatrix}-y&0\\0&1\end{smallmatrix}\bigr)$. By Lemma~\ref{pc:pointlaw}, as
$\#\auGL_{\mathsf A}=\#\auGL_{\mathsf B}$, $L'\binom y{\iota'}=c\,L(\iota\ x_y)$. For
$s\in\pcMat_{\mathsf N\times\mathsf T'}$ the matrix $u_s=\bigl(\begin{smallmatrix}1&s\\0&1\end{smallmatrix}
\bigr)\in\auGL_{\mathsf M}$ (rows and columns $\mathsf N,\mathsf T$) has determinant $1$, $u_s\iota=\iota$
and $u_s x_y=\bigl(\begin{smallmatrix}-y&s\\0&1\end{smallmatrix}\bigr)$; so $L(\iota\ x)$ is the same for
all $x=\bigl(\begin{smallmatrix}-y&s\\0&1\end{smallmatrix}\bigr)$, and summing against $q^+$ gives the
claim.

\emph{Step 2.} Apply Step 1 to $q=\ptvF(\aubar p)$. It remains to show
$\sum_x\ptvF(\aubar{p^+})(x)L(\iota\ x)=\gamma\sum_xq^+(x)L(\iota\ x)$ with $\gamma=G\cdot\#\pcMat_{\mathsf N
\times\mathsf T}(\pcR)$, $G$ the Gauss sum of Lemma~\ref{ptv:gauss} for $\mathsf T$; then $|\gamma|^2=
(\#\pcR)^{|\mathsf T|^2+2|\mathsf N||\mathsf T|}$. By Lemma~\ref{ptv:F-ext} applied to $\aubar p$ (note
$\aubar{p^+}=(\aubar p)^+$; here $x$ plays the role of $\rho$, with rows $\mathsf N\sqcup\mathsf T$ and
columns $\mathsf M\sqcup\mathsf T'$),
\[
\sum_x\ptvF(\aubar{p^+})(x)L(\iota\ x)=\sum_{A,B}q(-A)\sum_D\psi(\ptvtr D)\,K(A,B,D),\qquad
K(A,B,D)=L\Bigl(\iota\ \begin{pmatrix}A&B\\0&D\end{pmatrix}\Bigr).
\]
For $g\in\auGL_{\mathsf T}$, $k=\mathrm{diag}(1,g)\in\auGL_{\mathsf M}$ fixes $\iota$ and multiplies the
last block row by $g$, so $K(A,B,gD)=\chi(\det g)K(A,B,D)$; by Lemma~\ref{ptv:gauss} at $y=1$,
$\sum_D\psi(\ptvtr D)K(A,B,D)=G\,K(A,B,1)$. Hence the sum is $G\sum_{A,B}q(-A)K(A,B,1)=G\,\#\pcMat_{\mathsf
N\times\mathsf T}\sum_xq^+(x)L(\iota\ x)$.
\end{proof}

\subsection{The reversed squares and the adjoint law}\label{ptv:ss-RS}

Fix $i+n+2\le m$ and a twist $\vartheta$, with $\pcm=\pcm_\vartheta$, $\pcR=\pcR_\vartheta$,
$\psi=\psi_\pcm$, $\zeta=\zeta_\pcm$ of Proposition~\ref{ptv:rev}, and $\chi=\chi_\vartheta$ the tame
character of $(\nu_{n+1}\vartheta^{n+1})^{-1}$, non-trivial at every prime of $\pcm$
(\S\ref{pc:ss-exchange}). The nodes of column $k$ share an index set; let $\mathsf N$, $\mathsf M=
\mathsf N\sqcup\mathsf T$, $\mathsf M'=\mathsf M\sqcup\mathsf T'$ be those of columns $n$, $n+1$, $n+2$,
the last block being $\mathsf T$, resp.\ $\mathsf T'$; the inclusion matrices $\iota$, $\iota'$ of
\S\ref{ptv:ss-fourier} are the point data of $a_{i+1,n}$ and $a_{i,n+1}$. For a matrix $\xi$ write
$a[\xi]$ for the lift along the same ascending pair as $a$ with the point datum $\delta_\xi$, so
$a_{i+1,n}=a_{i+1,n}[\iota]$, $a_{i,n+1}=a_{i,n+1}[\iota']$.

Let $d_{i,k}=\Theta^{\mathsf t}(P_{i,k},\varepsilon_{i,k})$ ($k=n,n+1$) be the descents of
\S\ref{pc:ss-exchange}(iv): $P_{i,k}$ is the joint harmonic of the descending pair at $(i,k)$,
the characters are those of Proposition~\ref{ptv:chars}, $\varepsilon_{i,0}$ is the datum of
Theorem~\ref{pc:tri}(a) and $\varepsilon_{i,k+1}=\varepsilon_{i,k}^+$. By induction from
$\varepsilon_{i,0}$ (which is $\delta_\iota$ up to a constant: by (W7) in the model (W5) the map
$\mathsf C$ is complex conjugation of functions at the primes of $\pcm$, up to a constant), all
$\varepsilon_{i,k}$ are supported on injective matrices. Let $R_{k}[\varrho]$ be the reversed lift of
Corollary~\ref{ptv:desc-adj} for $d_{i,k}$ with the point datum $\delta_\varrho$ in place of
$\varepsilon_{i,k}^\flat$, and $R_k[f]=\sum_\varrho f(\varrho)R_k[\varrho]$. For $y\in\pcF_{i+1,n}$ and
$z\in\pcF_{i,n+2}$ define, for $\xi=(\xi_1\ \xi_2)\in\pcMat_{\mathsf M\times(\mathsf N\sqcup\mathsf M')}
(\pcR)$ and $\eta=\binom{\eta_1}{\eta_2}\in\pcMat_{(\mathsf N\sqcup\mathsf M')\times\mathsf M}(\pcR)$,
\[
\mathfrak L(\xi)=\bigl\langle z,\,R_{n+1}[\xi_2]\,a_{i+1,n}[\xi_1]\,y\bigr\rangle,\qquad
\mathfrak L'(\eta)=\bigl\langle z,\,a_{i,n+1}[\eta_2]\,R_n[\eta_1]\,y\bigr\rangle .
\]
These are the two composites around the \emph{reversed square} with corners
$U_1=\pcG_{i+1,n}$ (source), $U_2=\pcG_{i,n+2}$ (target), $A=\pcG_{i+1,n+1}$, $A'=\pcG_{i,n+1}$, all
maps being well defined between the spaces $\pcF$ by Theorem~\ref{ptr:edge}(d) and
Corollary~\ref{ptv:desc-adj}.

\begin{enumerate}
\item[(RS)] There is a function $\kappa^{\mathrm{rev}}_{i,n}$ on $m$-fold twists with modulus $1$ up to norm
powers (Definition~\ref{pc:modulus}) such that for every $\vartheta$ and all $y,z$:
$\mathfrak L(kx)=\chi(\det k)\mathfrak L(x)$ for $k\in\auGL_{\mathsf M}(\pcR)$ and all $x$, and
$\mathfrak L,\mathfrak L'$ satisfy the square law of \S\ref{pc:ss-exchange} for
$(\mathsf B,\mathsf A)=(\mathsf M,\mathsf N\sqcup\mathsf T')$ with constant $\kappa^{\mathrm{rev}}_{i,n}
(\vartheta)$.
\end{enumerate}

\noindent\emph{(RS) is proved in Theorem~\ref{ptw:RS}}: the reversed square is an edge square of excess
$e=0$ and weight $(n+1)\Delta$ (so $L_*(\pcm,\nu,0)=1$ and $|c|^2=C$), with $y_1=y$, $y_2=\aubar z$.
The rank and sign conditions are those of the square at $(i,n)$ rearranged; (E1) for the reversed
pairs $(A',U_1)$, $(U_2,A)$, including (H) for $(U_2,A)$, is Proposition~\ref{ptw:revedge} (through
Howe's argument with the groups exchanged, Theorem~\ref{ptw:Hc}, and the preservation of polynomial
vectors by $\mathcal A_\infty$, Lemma~\ref{ptw:poly}); (E2) holds because $\lambda_S=\lambda_T=1$
(Lemma~\ref{ptw:W7}), so the four conditions are those of Proposition~\ref{ptv:chars}.

\begin{theorem}[Theorem~\ref{pc:tri}(c)]\label{ptv:adjoint}
The adjoint law \textup{(T4)} holds for the triangle
of \S\ref{pc:ss-triangle}, with functions $\kappa_{i,n}$ of modulus $1$ up to norm powers.
\end{theorem}
\begin{proof}
We use \textup{(RS)} for all $i+n+2\le m$, which is Theorem~\ref{ptw:RS}.
\emph{The adjoint of $a_{i,n+1}$.} By Lemma~\ref{ptv:adj} for the ascending pair at $(i,n+1)$ (the
point datum $\delta_{\iota'}$ is supported on an injective matrix), every $z\in\pcF_{i,n+2}$ has
$x=a_{i,n+1}^*z\in\pcF_{i,n+1}$ with $\langle a_{i,n+1}v,z\rangle=\langle v,x\rangle$ for all $v$.

\emph{Reduction.} For $v'\in\pcF_{i+1,n}$, with $d^*$ as in Corollary~\ref{ptv:desc-adj},
$\langle d_{i,n}x,v'\rangle=\langle x,d^*_{i,n}v'\rangle=\langle z,a_{i,n+1}d^*_{i,n}v'\rangle$ and
$\langle d_{i,n+1}z,a_{i+1,n}v'\rangle=\langle z,d^*_{i,n+1}a_{i+1,n}v'\rangle$. So (T4) with constant
$\kappa_{i,n}$ follows from
\begin{equation}\label{ptv:eq-adj}
a_{i,n+1}\,d^*_{i,n}=\kappa_{i,n}(\vartheta)\,d^*_{i,n+1}\,a_{i+1,n}\qquad\text{on }\pcF_{i+1,n}.
\end{equation}

\emph{The reversed commutation.} Let $p=\varepsilon_{i,n}$, a function on $\pcMat_{\mathsf M\times
\mathsf N}(\pcR)$, so $\varepsilon_{i,n+1}=p^+$, $\varepsilon_{i,n}^\flat=\ptvF(\aubar p)$ and
$\varepsilon_{i,n+1}^\flat=\ptvF(\aubar{p^+})$ (Proposition~\ref{ptv:rev}, with $\psi=\psi_\pcm$). By
linearity of lifts in the datum,
\[
\langle z,a_{i,n+1}R_n[\ptvF\aubar p]y\rangle=\sum_{y'}\ptvF(\aubar p)(y')\,\mathfrak L'\tbinom{y'}{\iota'},
\qquad\langle z,R_{n+1}[\ptvF\aubar{p^+}]a_{i+1,n}y\rangle=\sum_x\ptvF(\aubar{p^+})(x)\,\mathfrak L(\iota\ x).
\]
By (RS) and Proposition~\ref{ptv:fourier}, $\gamma\langle z,a_{i,n+1}R_n[\ptvF\aubar p]y\rangle=
\kappa^{\mathrm{rev}}_{i,n}\langle z,R_{n+1}[\ptvF\aubar{p^+}]a_{i+1,n}y\rangle$ for all $z\in\pcF_{i,n+2}$.
Both vectors lie in $\pcF_{i,n+2}$, on which the Petersson product is definite; so
\[
a_{i,n+1}R_n=\frac{\kappa^{\mathrm{rev}}_{i,n}}{\gamma}\,R_{n+1}\,a_{i+1,n}\qquad\text{on }\pcF_{i+1,n},
\]
$R_k$ now denoting the reversed lift with the datum $\varepsilon_{i,k}^\flat$.

\emph{Conclusion.} By Corollary~\ref{ptv:desc-adj}, $d^*_{i,k}=c_k\,R_k\,\psrho(\zeta^{-1})$, with the same
central $\zeta$ for $k=n,n+1$ (it does not depend on the spaces). By Lemma~\ref{ptv:central},
$a_{i+1,n}\psrho(\zeta^{-1})=\eta\,\psrho(\zeta^{-1})a_{i+1,n}$ with $|\eta|=1$. Hence
\[
a_{i,n+1}d^*_{i,n}=c_n\frac{\kappa^{\mathrm{rev}}_{i,n}}{\gamma}R_{n+1}a_{i+1,n}\psrho(\zeta^{-1})=
\frac{c_n\,\eta\,\kappa^{\mathrm{rev}}_{i,n}}{\gamma\,c_{n+1}}\,d^*_{i,n+1}a_{i+1,n},
\]
which is \eqref{ptv:eq-adj} with $\kappa_{i,n}=c_n\eta\kappa^{\mathrm{rev}}_{i,n}/(\gamma c_{n+1})$. By
Proposition~\ref{ptv:rev}, $|c_n|^2=C\,(\#\pcR_\vartheta)^{-M_nM_{n+1}}$ and $|c_{n+1}|^2=C'(\#\pcR_\vartheta)^{
-M_{n+1}M_{n+2}}$; by Proposition~\ref{ptv:fourier}, $|\gamma|^2=(\#\pcR_\vartheta)^{t(t+2M_n)}$; and
$\#\pcR_\vartheta=M^f\auNm Q_\vartheta$ (\S\ref{pc:ss-data}). So $|\kappa_{i,n}|^2$ is a constant times
$|\kappa^{\mathrm{rev}}_{i,n}|^2$ times an integral power of $\auNm Q_\vartheta$, and $\kappa_{i,n}$ has
modulus $1$ up to norm powers by Lemma~\ref{pc:mod}(a).
\end{proof}

\subsection{Final check of Lemma~\ref{pc:arch} and Theorem~\ref{pc:tri}}\label{ptv:ss-check}

\begin{proposition}[Lemma~\ref{pc:arch}]\label{ptv:arch}
Lemma~\ref{pc:arch} holds for all its pairs, modulo the inputs of Theorems~\ref{ptr:edge},
\ref{ptq:F} and Proposition~\ref{plf:L4} \textup{((V), (D1)--(D3))}.
\end{proposition}
\begin{proof}
\emph{The pairs $(\pcG_{i,0},\pcG_{i,1})$.} (A1) and (A2) are Theorem~\ref{ptr:edge}(c), (d) (the last
paragraph of its statement). (A3): Proposition~\ref{ptr:Z} with Theorem~\ref{ptq:F}
(Remark~\ref{ptq:rem-F}) gives $Z_\infty=\prod_\tau d_{\lambda_\tau}^{-1}\langle\Phi_\tau,\mathsf P_\tau
\Phi_\tau\rangle$ with every factor positive, and the integral converges absolutely because
$|\langle P,\omega_\infty(g,1)P\rangle|\le\|P\|^2$ and the coefficient is integrable
(Lemma~\ref{ptr:gap}); it involves only archimedean data, at which the twists are trivial.

\emph{The pairs $(W_j,V_j)$.} (A1): for $W_j$ this is Lemma~\ref{plf:A12}; $[\auU(V_j)]$ is compact
(the group $\auU(V_j)_\infty$ is compact), so the type $e_{V_j}$ is integrable and its forms are bounded
on a compact quotient, which is the decay statement there. (A3): write $\sigma=\omega_\infty(\cdot,1)$
and $\mathsf P=\int e_{W_j}(x)\sigma(x^{-1})\,dx$. As in the first paragraph of the proof of
Proposition~\ref{ptr:Z} (with $\lambda=-\varsigma_\tau$, $\langle v_0,\pi_{W_j}(g)v_0\rangle=
\aubar{c_\lambda(g)}$, proof of Proposition~\ref{plf:L4}, Step 1),
$Z_\infty=\prod_\tau d_{\lambda_\tau}^{-1}\langle\Phi_\tau,\sigma_\tau(e_{\lambda_\tau})\Phi_\tau\rangle$, the
integral being absolutely convergent by Lemma~\ref{plf:A12}. Step 1 of the proof of
Proposition~\ref{plf:L4} shows $\sigma_\tau(e_{\lambda_\tau})\Phi_\tau=\Phi_\tau$ for the joint harmonic
$\Phi_\tau\ne0$. Hence $Z_\infty=\prod_\tau d_{\lambda_\tau}^{-1}\|\Phi_\tau\|^2>0$, independent of the
twist. (A2): the lift $h\mapsto\int\theta(g,h;P,\Psi)\varphi(g)\,dg$ converges absolutely and defines a
bounded function right invariant under an open subgroup, by the argument of the proof of
Proposition~\ref{pz:types}(b) (Lemma~\ref{pz:growth} and rapid decay, Lemma~\ref{plf:A12}), which uses
the row side only through $\mathsf PP=P$ (Step 1 of the proof of Proposition~\ref{plf:L4}). It is
reproduced by $e_{V_j}$: writing the lift as $L_h(P)$ with $L_h$ continuous on $\pzF_s$ and
$L_h(\omega_\infty(1,x)\Psi)=L_{hx}(\Psi)$ (proof of Lemma~\ref{ptr:conv-lift}), right convolution by
$e_{V_j}$ is $L_h(\int e_{V_j}(x)\omega_\infty(1,x)P\,dx)$, and $\int e_{V_j}(x)\omega_\infty(1,x)P\,dx=P$
by Schur orthogonality on the compact group $\auU(V_j)_\infty$, because $P$ is a highest weight vector
of the irreducible representation of highest weight $\Lambda'=\lambda'-\rho'$ (Step 1 of the proof of
Proposition~\ref{plf:L4}) and $e_{V_j}$ is, by definition, the type of that representation
(Proposition~\ref{ptw:A2} and Remark~\ref{ptw:rem-A2}).
\end{proof}

\begin{proposition}[The clauses of Theorem~\ref{pc:tri}]\label{ptv:tri}
\textup{(i)} The maps $a^\vartheta_{i,n}$ are well defined between the spaces $\pcF_{i,n}$
(Theorem~\ref{ptr:edge}(d)) and have adjoints (Lemma~\ref{ptv:adj}).
\textup{(ii)} With $d_{i,0}$ the transposed lift along $(\pcG_{i+1,0},\pcG_{i,1})$ whose kernel is
the complex conjugate of that of $a_{i,0}$ (it is a kernel for this pair by \textup{(W7)}, with the
characters of Proposition~\ref{ptv:chars}(b)), $d_{i,0}=a_{i,0}^*$ under $\pcF_{i+1,0}=\pcF_{i,0}$: this
is Lemma~\ref{ptv:adj} for $a_{i,0}$, because $\pcG_{i+1,0}=\pcG_{i,0}^-$ has the same group, quotient
and measure. This is \textup{(a)}.
\textup{(iii)} \textup{(b)} is Corollary~\ref{ptr:tri-b} (with Theorem~\ref{ptq:H} for (H) and
Proposition~\ref{ptv:chars} for the finite parts of the characters), modulo the published results
(D1)--(D5) and Theorem~\ref{pcx:X}.
\textup{(iv)} \textup{(c)} is Theorem~\ref{ptv:adjoint} (which uses \textup{(RS)}, Theorem~\ref{ptw:RS}).
\end{proposition}
\begin{proof}
Each item is the cited statement; for (ii), $a_{i,0}^*z(g)=\int\aubar{\theta(g,h)}z(h)\,dh$ and by
(W7) $\aubar{\theta^{\pcG_{i,0},\pcG_{i,1}}_{\alpha,\beta}(g,h;F)}=\theta^{\pcG_{i,0}^-,\pcG_{i,1}}
_{\alpha',\beta'}(g,h;\mathsf CF)$ with $\alpha',\beta'$ the conjugate characters times the powers of
$\upsilon$ of (W7), and $\mathsf C[P,\delta_\iota]$ is a datum $[P^-,c\,\delta_\iota]$, $|c|=1$ (at the
primes of the modulus, in the model (W5), $\mathsf C$ is complex conjugation up to a constant, since
$[x,a']^-=\aubar{[x,a']}$; elsewhere it maps the vacuum to a multiple of the vacuum, (W4)).
\end{proof}



%% ===== proofs/C-triangle4.tex =====
% proofs/C-triangle4.tex
% Continuation of proofs/C-triangle.tex (ptr:), C-triangle2.tex (ptq:), C-triangle3.tex (ptv:):
% Howe's argument with rows and columns exchanged, the smoothing lemma (H) for the reversed pairs,
% the characters of the reversed kernels, and the reversed square law (RS) of \S\ref{ptv:ss-RS};
% then the edge hypothesis (E1) for the grid squares and the check of (A2) for the pairs (W_j,V_j).
% Labels of this file start with ptw:. Only amsmath/amssymb/amsthm are used.
% Conventions and macros: proofs/C-triangle.tex, C-triangle3.tex.
\providecommand{\auU}{\mathrm U}
\providecommand{\auHom}{\operatorname{Hom}}
\providecommand{\auGL}{\mathrm{GL}}
\providecommand{\auNm}{\operatorname{Nm}}
\providecommand{\auA}{\mathbb A}
\providecommand{\auO}{\mathcal O}
\providecommand{\aubar}[1]{\overline{#1}}
\providecommand{\psrho}{\varrho}
\providecommand{\psHS}{\mathcal H}
\providecommand{\pzF}{\mathfrak F}
\providecommand{\pcF}{\mathcal F}
\providecommand{\pcG}{\mathsf G}
\providecommand{\pcm}{\mathfrak m}
\providecommand{\pcR}{R}
\providecommand{\pcMat}{\mathrm M}
\providecommand{\pop}{\mathsf p}
\providecommand{\ptrwt}{\operatorname{wt}}
\providecommand{\ptwsgn}{\operatorname{sgn}}

\section{The reversed square law, the grid edges, and the pairs $(W_j,V_j)$}\label{ptw:sec}

This section continues \S\ref{ptr:sec}, \S\ref{ptq:sec} and \S\ref{ptv:sec}, whose notation it
uses; inner products are conjugate-linear in the first variable.

\emph{Contents of this section.}
\begin{enumerate}
\item[(1)] Howe's argument and the uniqueness of the joint harmonic hold with the roles of rows and
columns exchanged (Propositions~\ref{ptw:unique} and~\ref{ptw:howe}); the only point that does not
transpose word for word, the string $B$ of columns that are not matched, is handled by the
observation that a column of $B$ carries no variable of any polynomial that is highest of the
lowest weight for the columns and harmonic for the rows (Lemma~\ref{ptw:B}). Consequently the
smoothing statement holds with the groups exchanged, for every ascending and descending pair
(Theorem~\ref{ptw:Hc}).
\item[(2)] At the infinite places the reversal map $\mathcal A_\infty$ of \S\ref{ptv:ss-rev} maps
polynomial vectors onto polynomial vectors (Lemma~\ref{ptw:poly}), and the characters $\lambda_S$,
$\lambda_T$ of \eqref{ptv:eq-A} are trivial; equivalently, the character of (W7) is
$\upsilon(\det g)^{-N'}\upsilon(\det h)^{N}$ (Lemma~\ref{ptw:W7}).
\item[(3)] Hence the reversed pairs are edges in the sense used by Theorem~\ref{ptr:square},
including (H) (Proposition~\ref{ptw:revedge}), the reversed squares satisfy (E2), and the reversed
square law (RS) of \S\ref{ptv:ss-RS} holds (Theorem~\ref{ptw:RS}). With Theorem~\ref{ptv:adjoint}
this proves the adjoint law, Theorem~\ref{pc:tri}(c), modulo the same published results
(D1)--(D5) as Theorem~\ref{pc:tri}(b).
\end{enumerate}
Finally, \S\ref{ptw:ss-A2} identifies the type on $V_j$ used in the check of (A2) for the pairs
$(W_j,V_j)$ in Proposition~\ref{ptv:arch}, and \S\ref{ptw:ss-grid} records where the edge hypothesis
(E1) for the grid squares is proved.

\subsection{Rows and columns exchanged}\label{ptw:ss-sym}

We work at one place $\tau$, in the polynomial model of \S\ref{par:ss-fock}, for a matching pair
$(\lambda,\lambda')$ with distinct entries in the sense of \eqref{par:eq-match}: rows $R$ ($n$ of
them), columns $C$ ($n'=n+t$), every row $r$ matched with a column $c(r)$, and the string
$B=C\smallsetminus c(R)$ of $t$ columns with $\{\xi'_c:c\in B\}=\{\tfrac{t-1}2-b:0\le b<t\}$ and
$\varepsilon'$ constant on $B$. (By Proposition~\ref{ptr:match} and Lemma~\ref{ptr:gap} this is the
situation at every $\tau$ for every ascending and every descending pair of the triangle, the
derived actions being \eqref{par:eq-osc} by Lemma~\ref{ptr:const}.) Notation is that of
\S\ref{par:ss-harm}: $\Lambda=\Lambda^K(\lambda)$, $\Lambda'=\Lambda^K(\lambda')$, $v_a$, $v'_c$ the
weights of the vacuum, $w_r=\Lambda_r-v_r$, $w'_c=\Lambda'_c-v'_c$.

\subsubsection*{The column operators in the form of the row operators}
Put $\eta'_{rc}=-\eta_{rc}=\varepsilon_r\varepsilon'_c$. For $\varepsilon'_c=\varepsilon'_d$ we have
$\eta_{rc}=\eta_{rd}$ for all $r$, and \eqref{par:eq-osc} gives, using $\partial_{rc}z_{rd}=
z_{rd}\partial_{rc}+\delta_{cd}$,
\[
E'_{cd}=-\sum_{r:\eta_{rc}=1}\partial_{rc}z_{rd}+\sum_{r:\eta_{rc}=-1}z_{rc}\partial_{rd}+\gamma'\delta_{cd}
=D'^{+}_{c\leftarrow d}-D'^{-}_{d\leftarrow c}+\delta_{cd}v'_c,\qquad
D'^{s}_{c\leftarrow d}=\sum_{r:\eta'_{rd}=s}z_{rc}\partial_{rd},
\]
with $v'_c=\gamma'-\#\{r:\eta'_{rc}=-1\}$. For $\varepsilon_a=\varepsilon_b$, the proof of
Proposition~\ref{par:exist} writes $E_{ab}=D^+_{a\leftarrow b}-D^-_{b\leftarrow a}+\delta_{ab}v_a$,
$D^s_{a\leftarrow b}=\sum_{c:\eta_{bc}=s}z_{ac}\partial_{bc}$, $v_a=\gamma-\#\{c:\eta_{ac}=-1\}$. For
$\varepsilon'_c\ne\varepsilon'_d$ the degree-lowering part of $E'_{cd}$ is
$-\sum_{r:\eta'_{rc}=-1}\partial_{rc}\partial_{rd}$, so up to sign
$\Delta'_{cd}=\sum_{r:\eta'_{rc}=-1}\partial_{rc}\partial_{rd}$; and for rows $a,b$ and a column $c$ with
$\eta'_{ac}=1=-\eta'_{bc}$ (so $\varepsilon_a\ne\varepsilon_b$),
$\sum_{d:\varepsilon'_d=\varepsilon'_c}\partial_{ad}\partial_{bd}=\Delta_{ab}$, because $\eta_{ad}=-1$ means
$\varepsilon'_d=\varepsilon_a=\varepsilon'_c$. Thus the substitution
\[
(R,C,\varepsilon,\varepsilon',\eta,\lambda,\lambda',E,\Delta,D,v)\longmapsto
(C,R,\varepsilon',\varepsilon,\eta',\lambda',\lambda,E',\Delta',D',v')
\]
carries the formulas for the row operators into those for the column operators and conversely,
and the Fischer product is unchanged. Lemma~\ref{par:polar}, Lemma~\ref{ptq:fischer} and
Lemma~\ref{ptq:tensor} are statements about these formulas and the Fischer product only (their
proofs use nothing else: in Lemma~\ref{par:polar} the identities $E_{ab}=D^+-D^-+\delta v$, the
adjointness of $z_i$ and $\partial_i$, and the vanishing of $\sum_b\partial_{bc}\partial_{bd}$ through
$\Delta'_{cd}$; in Lemma~\ref{ptq:tensor} the commutation relations of the $E$ with the $\Delta$). So
they hold with rows and columns exchanged. We record the column form of Lemma~\ref{par:polar}.

\begin{lemma}[Column polarisations]\label{ptw:polar}
Let $P$ be a polynomial killed by the column raising operators $E'_{cd}$ ($\varepsilon'_c=\varepsilon'_d$,
$\lambda'_c>\lambda'_d$) and by the row Laplacians $\Delta_{ab}$ ($\varepsilon_a\ne\varepsilon_b$). Let
$\varepsilon'_c=\varepsilon'_d$ and $s=\pm1$.
\textup{(b)} If $s\lambda'_c>s\lambda'_d$, then $D'^s_{c\leftarrow d}P=0$.
\textup{(c)} If $E'_{cc}P=\Lambda'_cP$ and $k=s\,w'_c\ge0$, then $D'^s_{c\leftarrow c}P=kP$ and
$\partial_{rc}P=0$ for every $r$ with $\eta'_{rc}=-s$.
\end{lemma}
\begin{proof}
Lemma~\ref{par:polar}(b), (c) after the substitution above.
\end{proof}

\begin{lemma}[Column weights]\label{ptw:relw}
For a matched column $c=c(r)$, $w'_c\,\ptwsgn\xi'_c\ge0$; for $c\in B$, $w'_c=0$.
\end{lemma}
\begin{proof}
By Lemma~\ref{par:relw}(ii), $w'_{c(r)}=-w_r$, and $\xi'_{c(r)}=-\xi_r$ by \eqref{par:eq-match}; so
$w'_{c(r)}\ptwsgn\xi'_{c(r)}=w_r\ptwsgn\xi_r\ge0$ by Lemma~\ref{par:relw}(i). The second assertion is
Lemma~\ref{par:relw}(ii).
\end{proof}

\begin{lemma}[The string carries no variable]\label{ptw:B}
Let $P$ be killed by the column raising operators and by the row Laplacians, with
$E'_{cc}P=\Lambda'_cP$ for a column $c\in B$. Then $\partial_{rc}P=0$ for every row $r$: no monomial of
$P$ contains a variable of the column $c$.
\end{lemma}
\begin{proof}
By Lemma~\ref{ptw:relw}, $w'_c=0$, so $k=s\,w'_c=0\ge0$ for both $s=1$ and $s=-1$.
Lemma~\ref{ptw:polar}(c) with $s=1$ gives $\partial_{rc}P=0$ when $\eta'_{rc}=-1$, and with $s=-1$ when
$\eta'_{rc}=1$.
\end{proof}

The uniqueness of Proposition~\ref{par:unique} omits the column weight condition. Howe's argument
with the groups exchanged produces polynomials for which the \emph{row} weight is not known, so we
need the uniqueness with that condition omitted instead. Its proof is that of
Proposition~\ref{par:unique} with rows and columns exchanged, except that the columns of $B$ have
to be treated separately.

\begin{proposition}[Uniqueness, rows and columns exchanged]\label{ptw:unique}
Call a polynomial $F\ne0$ a \emph{column pre-joint harmonic} if it satisfies all conditions of
Definition~\ref{par:def-harm} except possibly $E_{aa}F=\Lambda_aF$. Two column pre-joint harmonics
are proportional. In particular every column pre-joint harmonic is a multiple of the joint harmonic.
\end{proposition}
\begin{proof}
Let $x_0$ be the matrix with $z_{rc}(x_0)=1$ if $c=c(r)$ and $0$ otherwise, as in the proof of
Proposition~\ref{par:unique}. As there, it suffices to show that a column pre-joint harmonic $F$ has
$F(x_0)\ne0$: for two of them $P,Q$, the polynomial $P-\frac{P(x_0)}{Q(x_0)}Q$ satisfies the same
linear conditions and vanishes at $x_0$, so it is $0$.

Suppose $F(x_0)=0$. Write $F(x_0+u)=\sum_mc_mu^m$, choose a monomial $u^m$ of least positive degree
$\delta$ with $c_m\ne0$ and a variable $i$ occurring in it; then $\partial_iF\ne0$. Let $V=\sum_k
z_{a_k}\partial_{b_k}$ be a vector field such that exactly one term, $z_{a_1}\partial_{b_1}$ with
$b_1=i$, has $z_{a_1}(x_0)\ne0$. In $VF$ the term $z_{a_1}(x_0)\partial_iF$ contributes a non-zero
coefficient of $u^{m-e_i}$, of degree $\delta-1$, and every other term $z_{a_k}\partial_{b_k}F$ has
only terms of degree $\ge\delta$ ($z_{a_k}$ vanishes at $x_0$ and $\partial_{b_k}F$ has no terms of
degree $<\delta-1$). As $F$ has no term of degree $\le\delta-1$, $VF$ is not a multiple of $F$. We
reach a contradiction by producing such a $V$ with $VF\in\C F$.

Let $i=(r,c)$ and $s=\eta'_{rc}$.

\emph{$c\in B$.} By Lemma~\ref{ptw:B} (applicable: $F$ is killed by the column raising operators and
the row Laplacians and $E'_{cc}F=\Lambda'_cF$), $\partial_{rc}F=0$, a contradiction.

\emph{$c=c(r_1)$ is matched.} Then $|\xi'_c|>\tfrac t2$. If $-s\,\xi'_c>\tfrac t2$, then
$\ptwsgn\xi'_c=-s$, so $(-s)w'_c\ge0$ by Lemma~\ref{ptw:relw}, and Lemma~\ref{ptw:polar}(c) with $-s$
gives $\partial_{r'c}F=0$ whenever $\eta'_{r'c}=s$, in particular $\partial_iF=0$: a contradiction. So
$s\,\xi'_c>\tfrac t2$. Then $\ptwsgn\xi_{r_1}=-\ptwsgn\xi'_c=-s$, so by \eqref{par:eq-match}
$\eta_{r_1c}=-s$, i.e.\ $\eta'_{r_1c}=s=\eta'_{rc}$, and hence $\varepsilon_{r_1}=\varepsilon_r$. Three
cases.

($\alpha$) $r=r_1$. By Lemma~\ref{ptw:polar}(c) with $s$ ($s\,w'_c\ge0$ by
Lemma~\ref{ptw:relw}), $F$ is an eigenvector of $D'^s_{c\leftarrow c}=\sum_{r':\eta'_{r'c}=s}
z_{r'c}\partial_{r'c}$. At $x_0$, $z_{r'c}(x_0)\ne0$ only for $c=c(r')$, i.e.\ $r'=r_1=r$; so the only
term not vanishing at $x_0$ is $z_{rc}\partial_{rc}=z_{rc}\partial_i$. Contradiction.

($\beta$) $r\ne r_1$ and $s\lambda_{r_1}<s\lambda_r$ (the entries of $\lambda$ are distinct). Then
$(-s)\lambda_{r_1}>(-s)\lambda_r$ and $\varepsilon_{r_1}=\varepsilon_r$, so by Lemma~\ref{par:polar}(b) (the
row raising operators and the column Laplacians kill $F$) the field
$D^{-s}_{r_1\leftarrow r}=\sum_{d:\eta_{rd}=-s}z_{r_1d}\partial_{rd}$ kills $F$. Here $z_{r_1d}(x_0)\ne0$
only for $d=c(r_1)=c$, and the term $d=c$ occurs because $\eta_{rc}=-\eta'_{rc}=-s$; it is
$z_{r_1c}\partial_i$. Contradiction.

($\gamma$) $r\ne r_1$ and $s\lambda_{r_1}>s\lambda_r$. Since $\xi-\lambda$ is constant on the rows,
$s\xi_r<s\xi_{r_1}=-s\xi'_c<-\tfrac t2$; so $\ptwsgn\xi_r=-s$, and with $c_2=c(r)\ne c$,
\eqref{par:eq-match} gives $\eta_{rc_2}=-s$, i.e.\ $\eta'_{rc_2}=s=\eta'_{rc}$, so
$\varepsilon'_{c_2}=\varepsilon'_c$. Moreover $s\xi'_{c_2}=-s\xi_r>-s\xi_{r_1}=s\xi'_c$, so
$s\lambda'_{c_2}>s\lambda'_c$. By Lemma~\ref{ptw:polar}(b), $D'^s_{c_2\leftarrow c}=\sum_{r':\eta'_{r'c}=s}
z_{r'c_2}\partial_{r'c}$ kills $F$. Here $z_{r'c_2}(x_0)\ne0$ only for $c_2=c(r')$, i.e.\ $r'=r$, and
that term occurs ($\eta'_{rc}=s$); it is $z_{rc_2}\partial_i$. Contradiction.

So $F(x_0)\ne0$. The joint harmonic of Proposition~\ref{par:exist} is a column pre-joint harmonic,
whence the last assertion.
\end{proof}

\begin{proposition}[Howe's argument, groups exchanged]\label{ptw:howe}
Let $\Phi$ be the joint harmonic of $(\lambda,\lambda')$ and $\mathcal G$ a space of polynomials such
that: \textup{(i)} $E_{ab}\mathcal G\subset\mathcal G$ for all rows $a,b$; \textup{(ii)}
$\Phi\in\mathcal G$; \textup{(iii)} $E'(X)Q\in\mathcal G$ for every $X\in\mathfrak{gl}_{n'}(\C)$ and
every polynomial $Q$ with $E'_{cc}Q=\nu'_cQ$ for all $c$ and $E'_{cd}Q=0$ whenever
$\varepsilon'_c=\varepsilon'_d$, $\lambda'_c>\lambda'_d$, where $\nu'-\Lambda'$ is not a non-negative
combination of roots $e_c-e_d$ with $\lambda'_c>\lambda'_d$; here $E'(X)$ is the column action;
\textup{(iv)} $E'_{dc}Q\in\mathcal G$ for every polynomial $Q$, when $\varepsilon'_c=\varepsilon'_d$,
$\lambda'_c>\lambda'_d$; \textup{(v)} $(E'_{cc}-\Lambda'_c)Q\in\mathcal G$ for every polynomial $Q$
and every $c$. Then $\mathcal G$ contains all polynomials.
\end{proposition}
\begin{proof}
We show $\mathcal D_j\subset\mathcal G$ by induction on $j$, $\mathcal D_{-1}=0$. Suppose
$\mathcal D_{j-1}\subset\mathcal G$, $\mathcal D_j\not\subset\mathcal G$, and let $\mathcal N\ne0$ be the
orthogonal complement of $\mathcal G\cap\mathcal D_j$ in $\mathcal D_j$ for the Fischer product.

\emph{A highest vector for the rows.} The row operators preserve degree when
$\varepsilon_a=\varepsilon_b$, and $E_{ab}^*=E_{ba}$ for the Fischer product (\S\ref{ptq:ss-H}). If
$\varepsilon_a=\varepsilon_b$, $E_{ba}$ preserves $\mathcal G\cap\mathcal D_j$ by (i), so $E_{ab}$ preserves
$\mathcal N$. The $E_{aa}$ are self-adjoint and commute; so $\mathcal N$ contains $Q\ne0$ with
$E_{aa}Q=\nu_aQ$ for all $a$ (some $\nu$) and $E_{ab}Q=0$ whenever $\varepsilon_a=\varepsilon_b$,
$\lambda_a>\lambda_b$ (a weight vector of maximal weight for a linear order compatible with these
roots).

\emph{Column weight and column raising.} By (v), $\mathcal G$ is $E'_{cc}$-stable, hence so is
$\mathcal N$ ($E'_{cc}$ is self-adjoint and preserves degree), and $(E'_{cc}-\Lambda'_c)Q\in\mathcal N\cap
\mathcal G=0$. If $\varepsilon'_c=\varepsilon'_d$ and $\lambda'_c>\lambda'_d$, then by (iv)
$E'_{dc}\mathcal D_j\subset\mathcal G\cap\mathcal D_j$, so $E'_{cd}Q\in\mathcal D_j$ is orthogonal to
$\mathcal D_j$, and $E'_{cd}Q=0$.

\emph{Row Laplacians.} For $\varepsilon_a\ne\varepsilon_b$, $\Delta_{ab}Q\in\mathcal D_{j-2}\subset\mathcal G$,
so $E_{ba}\Delta_{ab}Q\in\mathcal G$ by (i), and $\Delta_{ab}Q=0$ by Lemma~\ref{ptq:fischer}.

\emph{Column Laplacians.} We apply Lemma~\ref{ptq:tensor} for the columns (valid by
\S\ref{ptw:ss-sym}): let $\varepsilon'_c\ne\varepsilon'_d$ and suppose that $\Delta'_{cd}Q$ has column
weight $\Lambda'+e_c-e_d$ and is killed by the column raising operators. If $e_c-e_d$ is not a
non-negative combination of roots $e_x-e_z$ with $\lambda'_x>\lambda'_z$, then (iii) gives
$E'_{dc}\Delta'_{cd}Q\in\mathcal G$, and $\Delta'_{cd}Q=0$ by Lemma~\ref{ptq:fischer} for columns.
Otherwise, pairing with $\lambda'$ gives $\lambda'_c>\lambda'_d$. Now $Q$ is killed by the column raising
operators and by the row Laplacians, and $E'_{xx}Q=\Lambda'_xQ$ for all $x$. If $c\in B$ or $d\in B$,
Lemma~\ref{ptw:B} shows that $Q$ contains no variable of that column, so every term
$\partial_{rc}\partial_{rd}Q$ of $\Delta'_{cd}Q$ vanishes. Otherwise $c,d$ are matched, $|\xi'_c|,|\xi'_d|>
\tfrac t2$, and $\xi'_c>\xi'_d$ ($\xi'-\lambda'$ is constant); so either $\xi'_c>\tfrac t2$, or
$\xi'_d<\xi'_c<-\tfrac t2$. In the first case $w'_c\ge0$ (Lemma~\ref{ptw:relw}) and
Lemma~\ref{ptw:polar}(c) with $s=1$ gives $\partial_{rc}Q=0$ whenever $\eta'_{rc}=-1$; in the second
$-w'_d\ge0$ and Lemma~\ref{ptw:polar}(c) with $s=-1$ gives $\partial_{rd}Q=0$ whenever $\eta'_{rd}=1$.
Every term of $\Delta'_{cd}=\sum_{r:\eta'_{rc}=-1}\partial_{rc}\partial_{rd}$ has $\eta'_{rc}=-1$ and
$\eta'_{rd}=1$ ($\varepsilon'_c\ne\varepsilon'_d$), so $\Delta'_{cd}Q=0$. By Lemma~\ref{ptq:tensor} for
the columns, all column Laplacians kill $Q$.

\emph{Conclusion.} $Q$ is a column pre-joint harmonic, so $Q=c\Phi$ by
Proposition~\ref{ptw:unique}. By (ii), $Q\in\mathcal G\cap\mathcal N=0$, a contradiction.
\end{proof}

\begin{remark}\label{ptw:rem-howe}
In the proof of Proposition~\ref{ptq:howe} the step ``row Laplacians'' uses that every row is
matched; in the exchanged argument the corresponding step meets the columns of $B$, and
Lemma~\ref{ptw:B} disposes of them. What Howe's argument really needs from the parameters, in either
direction, is: for two indices $x,y$ of the projected group with opposite signs and $\lambda_x>\lambda_y$
(resp.\ $\lambda'_x>\lambda'_y$), either $x$ carries no variable of one of the two kinds or $y$ carries
no variable of the other kind. For matched indices this follows from Lemma~\ref{par:polar}(c) and
the sign of $w$; for the unmatched string it is the statement $w'=0$ of Lemma~\ref{par:relw}(ii),
through Lemma~\ref{ptw:B}.
\end{remark}

\subsection{The smoothing lemma with the groups exchanged}\label{ptw:ss-H}

Let $(S,T)$ be an ascending or descending pair of the triangle, $\tau$ fixed, $G=\auU(S)_\tau$,
$G'=\auU(T)_\tau$, $\sigma=\omega_\infty(\cdot,1)$, $\sigma'=\omega_\infty(1,\cdot)$ at $\tau$, $\Phi$ the
joint harmonic, and $\mathsf P=\sigma(e_\lambda)$, $\mathsf P'=\sigma'(e_{\lambda'})$ the two projections
(proof of Theorem~\ref{par:A3}, valid for the pair by Theorem~\ref{ptr:edge}(b)). Let
$H_{\lambda'}\subset\pzF_\tau$ be the closed span of the images of the $G'$-equivariant isometries
$\pi_{\lambda'}\to\sigma'$. For $\beta\in C_c^\infty(G)$ write $\sigma(\beta)=\int\beta(x)\sigma(x)\,dx$.

\begin{lemma}\label{ptw:P}
\textup{(a)} For a polynomial $Q$, $\mathsf P'(E'_{dc}Q)=0$ if $\varepsilon'_c=\varepsilon'_d$ and
$\lambda'_c>\lambda'_d$, and $\mathsf P'(E'_{cc}Q)=\Lambda'_c\mathsf P'Q$.
\textup{(b)} $\mathsf P'$ is self-adjoint and $\mathsf P'x\in H_{\lambda'}$ for every $x$.
\textup{(c)} If $Q$ is a polynomial with $E'_{cc}Q=\nu'_cQ$ for all $c$ and $E'_{cd}Q=0$ whenever
$\varepsilon'_c=\varepsilon'_d$, $\lambda'_c>\lambda'_d$, where $\nu'-\Lambda'$ is not a non-negative
combination of roots $e_c-e_d$ with $\lambda'_c>\lambda'_d$, then $\mathsf P'(E'(X)Q)=0$ for every
$X\in\mathfrak{gl}_{n'}(\C)$.
\textup{(d)} There is $\beta_0\in C_c^\infty(G)$ with $\sigma(\beta_0)\mathsf P'\Phi=\mathsf P'\Phi$.
\end{lemma}
\begin{proof}
(a)--(c) are Lemma~\ref{ptq:P}(a)--(c) for $(G',\sigma',\lambda',\mathsf P')$ in place of
$(G,\sigma,\lambda,\mathsf P)$. Their proofs use only: that the operator is $\sigma(e_\lambda)$ for a
unitary representation $\sigma$ whose derived action on polynomial vectors is the row action
\eqref{par:eq-osc} (Lemma~\ref{ptr:const}), with $e_\lambda=d_\lambda\aubar{c_\lambda}$ the type of the
discrete series $\pi_\lambda$ and $c_\lambda$ the coefficient of the highest vector of its lowest
$K$-type; the integrability of $e_\lambda$ and of its derivatives (Lemma~\ref{ptr:gap}); Fischer
adjointness; that polynomial vectors of bounded degree form a finite-dimensional space stable
under $K$; Lemma~\ref{par:rigid} and (D1)--(D4). All of these hold for $(G',\sigma',\lambda')$ with the
column action (the parameter $\lambda'$ is of the class of Lemma~\ref{ptr:gap}, and
Lemma~\ref{par:rigid} and (D1)--(D4) are stated for an arbitrary unitary group and regular parameter).

(d) By (D5) applied to $e_\lambda$ on $G$ (smooth, bi-$K$-finite and $\mathfrak Z(\mathfrak g)$-finite,
as in \S\ref{ptq:ss-H}) choose $\beta_0\in C_c^\infty(G)$ with $\beta_0*e_\lambda=e_\lambda$; as
$\aubar{c_\lambda(y^{-1}g)}=c_\lambda(g^{-1}y)$, this reads $\int\beta_0(y)c_\lambda(hy)\,dy=c_\lambda(h)$ for
all $h$. Let $x=\mathsf P'\Phi$; by Theorem~\ref{ptr:edge}(b), $x=\mathsf P\Phi$. In the proof of
Lemma~\ref{ptq:P}(b) it is shown that $\mathsf P\Phi$ is a smooth vector for $\sigma$ of $T$-weight
$\Lambda$, killed by the compact raising operators of the rows, with $d\sigma(\Omega)x=\kappa_\lambda x$.
So Lemma~\ref{par:rigid} gives $\langle x,\sigma(h)x\rangle=\|x\|^2c_\lambda(h)$, and for $g\in G$
\[
\langle\sigma(g)x,\sigma(\beta_0)x\rangle=\int\beta_0(y)\langle x,\sigma(g^{-1}y)x\rangle dy=\|x\|^2\!\int
\beta_0(y)c_\lambda(g^{-1}y)dy=\|x\|^2c_\lambda(g^{-1})=\langle\sigma(g)x,x\rangle .
\]
So $\sigma(\beta_0)x-x$ lies in the closed span of $\sigma(G)x$ and is orthogonal to it; it is $0$.
\end{proof}

\begin{theorem}[(H) with the groups exchanged]\label{ptw:Hc}
Let $(S,T)$ be an ascending or descending pair of the triangle. For every polynomial $Q$ at $\tau$
there is $\beta_Q\in C_c^\infty(G)$ with $\mathsf P'Q=\sigma(\beta_Q)\mathsf P'\Phi$. Consequently, for
every polynomial vector $Q$ of the Fock space of $\auHom(S,T)\otimes\R$ there is a continuous
compactly supported $\beta_Q$ on $\auU(S)_\infty$ with
\begin{equation}\label{ptw:eq-Hc}
\mathsf P'Q=\int_{\auU(S)_\infty}\beta_Q(x)\,\omega_\infty(x,1)\,\mathsf P'P_{S,T}\,dx,
\end{equation}
the integral being a Bochner integral in $\pzF_s$; here $\mathsf P'=\int e_T(x)\omega_\infty(1,x)\,dx$
is the operator of \S\ref{ptr:ss-edge}.
\end{theorem}
\begin{proof}
Let $\mathcal G$ be the space of polynomials $Q$ with $\mathsf P'Q=\sigma(\beta)\mathsf P'\Phi$ for some
$\beta\in C_c^\infty(G)$. We check the hypotheses of Proposition~\ref{ptw:howe}.

(i) Let $Q\in\mathcal G$, $\mathsf P'Q=\sigma(\beta)\mathsf P'\Phi$, and $Y\in\mathfrak g$. As $\mathsf P'$ is
bounded and commutes with $\sigma$, and $Q$ is a smooth vector for $\sigma$ with derivative
$d\sigma(Y)Q=E(Y)Q$ (Lemma~\ref{ptr:const}), $\mathsf P'E(Y)Q=\frac d{dt}|_0\sigma(\exp tY)\mathsf P'Q$; and
$\sigma(\exp tY)\sigma(\beta)\mathsf P'\Phi=\sigma(\beta(\exp(-tY)\,\cdot))\mathsf P'\Phi$, whose derivative
at $0$ is $\sigma(\beta_Y)\mathsf P'\Phi$ with $\beta_Y(x)=\frac d{dt}|_0\beta(\exp(-tY)x)\in C_c^\infty(G)$.
So $E(Y)Q\in\mathcal G$ for $Y\in\mathfrak g$, hence for $Y\in\mathfrak g_\C=\mathfrak{gl}_n(\C)$, in
particular for the $E_{ab}$.

(ii) Lemma~\ref{ptw:P}(d). (iii) Lemma~\ref{ptw:P}(c): $\mathsf P'$ of such a polynomial is
$0=\sigma(0)\mathsf P'\Phi$. (iv), (v) Lemma~\ref{ptw:P}(a), in the same way.

So $\mathcal G$ contains all polynomials. The global statement follows from the local one exactly
as in the proof of Theorem~\ref{ptq:H}: $\mathsf P'=\bigotimes_\tau\mathsf P'_\tau$, $P_{S,T}=
\bigotimes_\tau\Phi_\tau$, $\omega_\infty(x,1)=\bigotimes_\tau\sigma_\tau(x_\tau)$, and every polynomial
vector is a finite sum of products $\bigotimes_\tau Q_\tau$, for which $\beta_Q=\prod_\tau
\beta_{Q_\tau}$. The identity holds in the Hilbert space; both sides lie in $\pzF_s$ (the right side
is a Bochner integral of a continuous compactly supported $\pzF_s$-valued function, by
Lemma~\ref{pz:growth}(a); the left side because $\mathsf P'$ is the Bochner integral of
$e_T(x)\omega_\infty(1,x)Q$, convergent in $\pzF_s$ by Lemma~\ref{pz:growth}(a) and the integrability
of $e_T$ times a power of $\|x\|$, Lemma~\ref{ptr:gap}); so it holds in $\pzF_s$ by
Lemma~\ref{ptv:embed}.
\end{proof}

\subsection{The reversal at the infinite places and the characters of (W7)}\label{ptw:ss-rev}

Let $(S,T)$ be spaces of ranks $N\le N'$, and $\mathcal A=\mathrm{Fl}\circ\mathsf C$, $\phi(x)=x^\dagger$ as in
\S\ref{ptv:ss-rev}, with $\mathcal A=\bigotimes_v\mathcal A_v$. Write $W_\infty=\auHom(S,T)\otimes\R$,
$W'_\infty=\auHom(T,S)\otimes\R$, $\pzF$, $\pzF'$ their Fock spaces (\S\ref{pz:ss-types}), and
$K\times K'$ the maximal compact subgroup of $\auU(S)_\infty\times\auU(T)_\infty$ used there.

\subsubsection*{The complex structures}
In the realisation \eqref{pk:eq-fock} of \S\ref{pz:ss-types} there are a real-linear isomorphism
$u\colon W_\infty\to\C^D$ and the standard hermitian form $h(z,z')=\sum_kz_k\aubar{z'_k}$ with
$(\psrho_\infty(y)F)(z)=e^{h(z,u(y))-|u(y)|^2/2}F(z-u(y))$ (in \eqref{pk:eq-fock}, $u(y)=\aubar y$ and
$z\cdot y=h(z,\aubar y)$); the complex structure $J=u^{-1}\circ i\circ u$ of $W_\infty$ is the positive
complex structure commuting with $K\times K'$ of \S\ref{pz:ss-types}. From the formula,
$\psrho_\infty(x)\psrho_\infty(y)=e^{iA(x,y)}\psrho_\infty(y)\psrho_\infty(x)$ with
$A(x,y)=2\Iml h(u(y),u(x))$ (the exponents of the two products differ by $-h(u(x),u(y))+h(u(y),u(x))$),
so $A(w,Jw)=2|u(w)|^2>0$ for $w\ne0$. The same holds for $W'_\infty$ with $u',J',A'$.

\begin{lemma}\label{ptw:J}
\textup{(a)} A complex structure $J_1$ on $W_\infty$ commuting with $K\times K'$ and with
$A(w,J_1w)>0$ for all $w\ne0$ equals $J$.
\textup{(b)} $A'(\phi x,\phi y)=-A(x,y)$, and $\phi\circ J=-J'\circ\phi$.
\end{lemma}
\begin{proof}
(a) At each $\tau$, in orthogonal bases of $S_\tau$ and $T_\tau$, $W_\tau=\auHom_\C(S_\tau,T_\tau)$ is
the sum of the four blocks $\auHom(S_{\tau,\pm},T_{\tau,\pm})$ ($S_{\tau,\pm}$ the spans of the basis
vectors of sign $\pm$), stable under $K_\tau\times K'_\tau=\auU(S_{\tau,+})\times\auU(S_{\tau,-})\times
\auU(T_{\tau,+})\times\auU(T_{\tau,-})$, which acts on a non-zero block $\auHom(S_{\tau,a},T_{\tau,b})$
through $x\mapsto h_bxg_a^{-1}$. This is an irreducible complex representation, and it is not
isomorphic to its complex conjugate (the centres of $\auU(S_{\tau,a})$ and $\auU(T_{\tau,b})$ act by
$\zeta^{-1}\zeta'$, resp.\ $\zeta\zeta'^{-1}$); so it is irreducible as a real representation and its
real commutant is $\C$ (the complex structure $i$ of the block and the scalars). Two distinct
non-zero blocks, also for different $\tau$, are non-isomorphic real representations (some central
factor acts non-trivially on one and trivially on the other). Hence $J_1$ preserves every block and
acts on it by an element $a+bi$ of the commutant with $(a+bi)^2=-1$, i.e.\ by $\epsilon i$,
$\epsilon=\pm1$. The same holds for $J$, with signs $\epsilon_0$. On a block, $A(w,J_1w)=\epsilon
A(w,iw)$ and $A(w,Jw)=\epsilon_0A(w,iw)$ are both positive for $w\ne0$; so $\epsilon=\epsilon_0$ and
$J_1=J$.

(b) By \eqref{ptv:eq-A}, $\mathcal A_\infty$ is conjugate-linear with $\mathcal A_\infty\psrho_\infty(x)=
\psrho'_\infty(\phi x)\mathcal A_\infty$. Applying $\mathcal A_\infty$ to $\psrho(x)\psrho(y)=e^{iA(x,y)}
\psrho(y)\psrho(x)$ gives $\psrho'(\phi x)\psrho'(\phi y)=e^{-iA(x,y)}\psrho'(\phi y)\psrho'(\phi x)$, so
$e^{iA'(\phi x,\phi y)}=e^{-iA(x,y)}$ for all real $x,y$; replacing $x$ by $sx$, $s\in\R$, gives
$A'(\phi x,\phi y)=-A(x,y)$. Put $J_1=-\phi^{-1}J'\phi$. It is a complex structure; it commutes with
$K\times K'$, because $\phi(hxg^{-1})=g\phi(x)h^{-1}$ for unitary $g,h$ (\S\ref{ptv:ss-rev}) and $J'$
commutes with $K'\times K$; and $A(w,J_1w)=-A(w,\phi^{-1}J'\phi w)=A'(\phi w,J'\phi w)>0$ for $w\ne0$.
By (a), $J_1=J$, i.e.\ $\phi J=-J'\phi$.
\end{proof}

\begin{lemma}[$\mathcal A_\infty$ preserves polynomial vectors]\label{ptw:poly}
$\mathcal A_\infty$ maps smooth vectors to smooth vectors, with $\mathcal A_\infty\partial_y=
\partial'_{\phi y}\mathcal A_\infty$, maps the vacuum to a multiple of the vacuum, and maps the space of
polynomial vectors of degree $\le d$ onto the corresponding space of $\pzF'$, for every $d$. The
same holds for $\mathcal A_\infty^{-1}$.
\end{lemma}
\begin{proof}
For a smooth vector $v$, $s\mapsto\psrho'(s\phi y)\mathcal A_\infty v=\mathcal A_\infty\psrho(sy)v$ is
differentiable with derivative $\mathcal A_\infty\partial_yv$ at $0$ ($\mathcal A_\infty$ is a bounded
real-linear map); this gives the first assertion. From the formula for $\psrho_\infty$,
$\partial_yF=h(z,u(y))F-D_{u(y)}F$, where $D_{u}$ is the derivative in the direction $u$ (complex-linear
in $u$, $F$ being entire), and $u(Jy)=iu(y)$; hence
\[
\partial_y+i\partial_{Jy}=2\,h(z,u(y)),\qquad\partial_y-i\partial_{Jy}=-2D_{u(y)}
\]
(multiplication by a linear form, resp.\ a derivation), and likewise on $\pzF'$. By
conjugate-linearity and Lemma~\ref{ptw:J}(b),
\[
\mathcal A_\infty(\partial_y\pm i\partial_{Jy})=(\partial'_{\phi y}\mp i\partial'_{\phi Jy})\mathcal A_\infty
=(\partial'_{\phi y}\pm i\partial'_{J'\phi y})\mathcal A_\infty .
\]
So $\mathcal A_\infty$ carries the operators $D_{u(y)}$ to the operators $D'_{u'(\phi y)}$ and the
multiplications by the linear forms $h(z,u(y))$ to the multiplications by $h(z',u'(\phi y))$; as $y$
runs over $W_\infty$ these run over all directional derivatives and all linear forms. The vacuum $1$
is killed by all $D_u$; so $\mathcal A_\infty1$ is a smooth vector (an entire function) killed by all
$D'_{u'}$, i.e.\ a constant $c$, and $c\ne0$ since $\mathcal A_\infty$ is injective. The polynomial
vectors of degree $\le d$ are the span of the products of at most $d$ multiplications by linear forms
applied to $1$; so $\mathcal A_\infty$ maps them into the polynomial vectors of degree $\le d$ of
$\pzF'$. The same argument for $\mathcal A_\infty^{-1}$ (which intertwines $\psrho'$ and
$\psrho\circ\phi^{-1}$) gives the reverse inclusion.
\end{proof}

\begin{lemma}[The exponents in (W7)]\label{ptw:W7}
Let $S,T$ have ranks $N,N'$. The character $\lambda$ of \textup{(W7)} (Lemma~\ref{pc:SW}) is
$\lambda(g,h)=\upsilon(\det g)^{-N'}\upsilon(\det h)^{N}$. Consequently the characters $\lambda_S,\lambda_T$
of \eqref{ptv:eq-A} are trivial: $\mathcal A\,\omega_{\auHom(S,T)}(g,h)=\omega_{\auHom(T,S)}(h,g)\,\mathcal A$,
and Proposition~\ref{ptv:rev} and Corollary~\ref{ptv:desc-adj} hold with $\lambda_S=\lambda_T=1$:
\[
\aubar{\theta_{\alpha,\beta}(g,h;P,\Psi)}=c\,\theta_{\beta,\alpha}(h,g\zeta_\pcm;P^\flat,\Psi^\flat).
\]
\end{lemma}
\begin{proof}
By (W7), $\lambda(g,h)=\upsilon(\det g)^a\upsilon(\det h)^b$ for integers $a,b$; we determine them at an
infinite place $\tau$. Let $W=\auHom(S,T)$ and $W^-=\auHom(S^-,T)$; at $\tau$ both have the same
underlying real space, symplectic forms of opposite signs, the same symplectic group and hence the
same metaplectic double cover $\mathrm{Mp}$, and the same image of $\auU(S)_\tau\times\auU(T)_\tau$. Let
$\omega^{\mathrm{Mp}}$, $\omega^{\mathrm{Mp},-}$ be the metaplectic representations on the Fock spaces
of $\psrho_\tau$, $\psrho^-_\tau$ (\S\ref{pco:ss-V}). The local component $\mathsf C_\tau$ of $\mathsf C$
is conjugate-linear with $\mathsf C_\tau\psrho_\tau(x)=\psrho^-_\tau(x)\mathsf C_\tau$, so
$\tilde u\mapsto\mathsf C_\tau\omega^{\mathrm{Mp}}(\tilde u)\mathsf C_\tau^{-1}$ is a strongly continuous
unitary representation of $\mathrm{Mp}$ with $\mathsf C_\tau\omega^{\mathrm{Mp}}(\tilde u)
\mathsf C_\tau^{-1}\psrho^-(w)(\cdots)^{-1}=\psrho^-(uw)$. By the irreducibility of $\psrho^-_\tau$ it
differs from $\omega^{\mathrm{Mp},-}$ by a continuous character of $\mathrm{Mp}$; as $\mathrm{Mp}$ is
connected with perfect Lie algebra $\mathfrak{sp}$, this character is trivial. By (V) for $W$ and
for $W^-$ (Theorem~\ref{pco:V}, whose proof uses only that $\omega$ is the Weil representation of
$\auU(W)$ defined with $\upsilon$, together with (W1$_\infty$), (W2$_\infty$), (M), and (S) when $W$
is definite), for $u=\iota(g,h)$ with a lift $\tilde u$,
$\omega_\tau(u)=d(\tilde u)\,\omega^{\mathrm{Mp}}(\tilde u)$ and $\omega^-_\tau(u)=d(\tilde u)\,
\omega^{\mathrm{Mp},-}(\tilde u)$ with the same genuine function $d(\tilde u)=\det(g)^{N'/2}
\det(h)^{-N/2}$, $|d|=1$. Hence
\[
\mathsf C_\tau\omega_\tau(u)\mathsf C_\tau^{-1}=\aubar{d(\tilde u)}\,\omega^{\mathrm{Mp},-}(\tilde u)=
d(\tilde u)^{-2}\omega^-_\tau(u)=\det(g)^{-N'}\det(h)^{N}\,\omega^-_\tau(u),
\]
and since $\upsilon_\tau(z)=z$ on $\auU(1)$ (Theorem~\ref{pco:V}), $a=-N'$, $b=N$ ($\det$ maps
$\auU(S)_\tau$ onto $\auU(1)_\tau$). By \eqref{ptv:eq-A}, $\lambda_S(\det g)\lambda_T(\det h)^{-1}$ is
$\lambda(g,h)$ times the factor $\upsilon(\det g)^{N'}\upsilon(\det h)^{-N}$ of Proposition~\ref{pf:flip}
(\S\ref{pcx:ss-flip}), i.e.\ $1$; as $\det$ maps $\auU(S)(\auA)$ and $\auU(T)(\auA)$ onto
$\auU(1)(\auA_{E^+})$, $\lambda_S=\lambda_T=1$. The rest is Proposition~\ref{ptv:rev} and
Corollary~\ref{ptv:desc-adj} with this value.
\end{proof}

\begin{remark}\label{ptw:rem-W7}
For the descent $d_{i,0}$ this gives the exponent needed in Proposition~\ref{ptv:chars}(b): the
conjugate of the kernel of $a_{i,0}$ (characters $\alpha=\mu_{i+1}\upsilon_{M_1}$,
$\beta=\mu_{i+1}\upsilon_{M_0}$) is a kernel of $\auHom(\pcG_{i,0}^-,\pcG_{i,1})$ with characters
$\alpha^{-1}\upsilon^{M_1}$ and $\beta^{-1}\upsilon^{M_0}$, since $\aubar{\theta_{\alpha,\beta}(g,h;F)}=
\alpha(\det g)\beta(\det h)^{-1}\lambda(g,h)\theta^-(\omega^-(g,h)\mathsf CF)$. Its target character is
$\mu_{i+1}^{-1}\upsilon^{M_0-\lfloor M_0/2\rfloor}=\mu_{i+1}^{-1}\upsilon^{\lceil r/2\rceil}$, of weight
$\kappa_{i+1}+\lceil r/2\rceil$, as Definition~\ref{ptr:desc-char} requires.
\end{remark}

\subsubsection*{The reversed pairs}
Let $(S,T)=(\pcG_{i+1,k},\pcG_{i,k+1})$ be a descending pair, $i+k+1\le m$, with the characters
$\alpha=\alpha^\downarrow_{i,k}$, $\beta=\beta^\downarrow_{i,k}$ of Proposition~\ref{ptv:chars} (for a
twist $\vartheta$), the types $e_S=\aubar{e_{i+1,k}}$, $e_T=\aubar{e_{i,k+1}}$ and the joint harmonic
$P=P_{S,T}$ of \S\ref{ptr:ss-edge}. The \emph{reversed pair} is $(T,S)$, i.e.\ the space
$\auHom(T,S)$, with source character $\beta$ on $\auU(T)$, target character $\alpha$ on $\auU(S)$, the
Weil representation $\omega^{\mathrm{rev}}=\omega_{\beta,\alpha}$ of $\auHom(T,S)$, the polynomial
$P^\flat=\mathcal A_\infty P$ (a polynomial vector by Lemma~\ref{ptw:poly}), source type $e_T$ and
target type $e_S$. Put $\mathsf P^{\mathrm{rev}}=\int_{\auU(T)_\infty}e_T(x)\,\omega^{\mathrm{rev}}_\infty
(x^{-1},1)\,dx$, the operator $\mathsf P$ of \S\ref{ptr:ss-edge} for this pair.

\begin{proposition}[Reversed pairs are edges]\label{ptw:revedge}
\textup{(c)} Every form of type $e_T$, $\aubar{e_T}$, $e_S$ or $\aubar{e_S}$ is a cusp form, rapidly
decreasing on Siegel sets.
\textup{(d)} For every finite datum $f$ and every form $y$ of type $e_T$, the lift
$g\mapsto\int_{[\auU(T)]}\theta^{\mathrm{rev}}(\omega^{\mathrm{rev}}(h,g)(P^\flat\otimes f))\,y(h)\,dh$
converges absolutely and is a form of type $e_S$; for every form $z$ of type $\aubar{e_S}$ the
transposed lift $h\mapsto\int_{[\auU(S)]}\theta^{\mathrm{rev}}(\omega^{\mathrm{rev}}(h,g)(P^\flat\otimes
f))\,z(g)\,dg$ converges absolutely and is a form of type $\aubar{e_T}$.
\textup{(H)} For every polynomial vector $Q$ of the Fock space of $\auHom(T,S)\otimes\R$ there is
$\alpha_Q\in C_c(\auU(S)_\infty)$ with $\mathsf P^{\mathrm{rev}}Q=\int\alpha_Q(x)\,
\omega^{\mathrm{rev}}_\infty(1,x)\,\mathsf P^{\mathrm{rev}}P^\flat\,dx$ (Bochner integral in $\pzF'_s$).
\textup{(L)} Lemma~\ref{ptr:conv-lift} holds for the reversed pair: for every polynomial vector $Q$,
finite datum $f$ and form $y$ of type $e_T$, the lift $F_Q$ of $y$ with $Q\otimes f$ satisfies
$F_Q(g)=\int\alpha_Q(x)F_{P^\flat}(gx)\,dx$.
\end{proposition}
\begin{proof}
(c) is Theorem~\ref{ptr:edge}(c). By Lemma~\ref{ptw:W7} and the first lines of the computation
``Globally'' in the proof of Proposition~\ref{ptv:rev} (valid for every datum $F$), for every
$F\in\psHS(\auHom(S,T))$
\[
\aubar{\theta_{\alpha,\beta}(g,h;F)}=c_\theta^{-1}\,\theta_{\beta,\alpha}(h,g;\mathcal AF),
\]
and $\mathcal A$ is a bijection. (d) Let $F'=P^\flat\otimes f$ and $F=\mathcal A^{-1}F'=P\otimes f_1$
($f_1=\mathcal A_f^{-1}f$). Then $\int\theta_{\beta,\alpha}(h,g;F')y(h)\,dh=c_\theta\,\aubar{\int
\theta_{\alpha,\beta}(g,h;F)\,\aubar{y(h)}\,dh}$. As $\aubar y$ has type $\aubar{e_T}$, the integral on
the right is the transposed lift along the descending pair, absolutely convergent and a form of type
$\aubar{e_S}$ by Theorem~\ref{ptr:edge}(d); its conjugate is a form of type $e_S$. Likewise the
transposed reversed lift of $z$ is $c_\theta$ times the conjugate of the lift of $\aubar z$ (type
$e_S$) along the descending pair, a form of type $e_T$, whose conjugate has type $\aubar{e_T}$.

(H) At the infinite places, Lemma~\ref{ptw:W7} gives $\mathcal A_\infty\omega_\infty(g,h)=
\omega^{\mathrm{rev}}_\infty(h,g)\mathcal A_\infty$ (the archimedean factors of the characters move
through the conjugate-linear map as in the proof of Proposition~\ref{ptv:rev}). With
$e_T(x^{-1})=\aubar{e_T(x)}$ and the unimodularity of $\auU(T)_\infty$,
\[
\mathcal A_\infty\mathsf P'v=\int\aubar{e_T(x)}\,\omega^{\mathrm{rev}}_\infty(x,1)\mathcal A_\infty v\,dx
=\int e_T(x)\,\omega^{\mathrm{rev}}_\infty(x^{-1},1)\mathcal A_\infty v\,dx=\mathsf P^{\mathrm{rev}}
\mathcal A_\infty v,
\]
$\mathsf P'=\int e_T(x)\omega_\infty(1,x)dx$ being the column projection of the descending pair
($\mathcal A_\infty$ is continuous and real-linear, so it commutes with Bochner integrals up to
conjugation of the scalar factor). Given $Q$, the vector $Q_1=\mathcal A_\infty^{-1}Q$ is a polynomial
vector (Lemma~\ref{ptw:poly}); by Theorem~\ref{ptw:Hc}, $\mathsf P'Q_1=\int\beta(x)\omega_\infty(x,1)
\mathsf P'P\,dx$. Applying $\mathcal A_\infty$,
\[
\mathsf P^{\mathrm{rev}}Q=\int\aubar{\beta(x)}\,\omega^{\mathrm{rev}}_\infty(1,x)\,\mathsf P^{\mathrm{rev}}
P^\flat\,dx,
\]
an identity in the Hilbert space; both sides lie in $\pzF'_s$ (as in the proof of
Theorem~\ref{ptw:Hc}), so it holds there by Lemma~\ref{ptv:embed}. Take $\alpha_Q=\aubar\beta$.

(L) The proof of Lemma~\ref{ptr:conv-lift} uses the functional $L_g(\Psi)=\int_{[\auU(T)]}
\theta^{\mathrm{rev}}(\omega^{\mathrm{rev}}(h,g)(\Psi\otimes f))y(h)\,dh$, continuous on $\pzF'_s$ with
$L_g(\omega^{\mathrm{rev}}_\infty(1,x)\Psi)=L_{gx}(\Psi)$ and $L_g(\mathsf P^{\mathrm{rev}}\Psi)=L_g(\Psi)$.
These are established in the proof of Proposition~\ref{pz:types}(b) from Lemma~\ref{pz:growth}
(stated for an arbitrary pair of spaces), reduction theory, the rapid decay of $y$ (by (c)) and
$y=R(e_T)y$; no property of the parameters is used. With (H) the proof of
Lemma~\ref{ptr:conv-lift} applies verbatim.
\end{proof}

\subsection{The reversed square law}\label{ptw:ss-RS}

\begin{lemma}[Change of splitting]\label{ptw:split}
In the square law of \S\ref{pc:ss-exchange} let $\mathsf N=\mathsf B\sqcup\mathsf A=\mathsf B'\sqcup
\mathsf A'$ be two splittings with $|\mathsf B|=|\mathsf B'|$. If $L,L'$ satisfy the square law with
constant $c$ for $(\mathsf B,\mathsf A)$, they satisfy it with constant $\chi(\epsilon)c$ for
$(\mathsf B',\mathsf A')$, for a sign $\epsilon\in\{\pm1\}\subset\pcR^\times$ depending only on the
two splittings (the functions on matrices being read through bijections of $\mathsf A$ with
$\mathsf A'$ and of $\mathsf B$ with $\mathsf B'$); in particular $|\chi(\epsilon)c|=|c|$.
\end{lemma}
\begin{proof}
Let $\sigma$ be the permutation matrix of $\mathsf N$ carrying $\mathsf B$ onto $\mathsf B'$ and $\mathsf A$
onto $\mathsf A'$ through the given bijections, $\epsilon=\det\sigma$. For $k\in\auGL_{\mathsf N}$ put
$\tilde k=k\sigma$: the columns $\mathsf B$ of $\tilde k$ are the columns $\mathsf B'$ of $k$, and the rows
$\mathsf A$ of $\tilde k^{-1}=\sigma^{-1}k^{-1}$ are the rows $\mathsf A'$ of $k^{-1}$. So if $(\Phi,\Psi)$ is
matched for $(\mathsf B',\mathsf A')$, then $\Phi(x_{\tilde k})=\chi(\det k)\Psi(y_{\tilde k})=
\chi(\epsilon)^{-1}\chi(\det\tilde k)\Psi(y_{\tilde k})$, i.e.\ $(\Phi,\chi(\epsilon)^{-1}\Psi)$ is matched for
$(\mathsf B,\mathsf A)$, and $\chi(\epsilon)^{-1}\sum\Psi L'=c\sum\Phi L$. The transformation laws do not
involve the splitting. As $\epsilon^2=1$, $\chi(\epsilon)=\pm1$.
\end{proof}

\begin{theorem}[(RS)]\label{ptw:RS}
Let $i+n+2\le m$. The reversed square with $U_1=\pcG_{i+1,n}$, $U_2=\pcG_{i,n+2}$, $A=\pcG_{i+1,n+1}$,
$A'=\pcG_{i,n+1}$, types $e_1=e_{i+1,n}$, $e_2=e_{i,n+2}$, $e=e_{i+1,n+1}$, $e'=e_{i,n+1}$, the
ascending pairs $(U_1,A)$, $(A',U_2)$ (the lifts $a_{i+1,n}$, $a_{i,n+1}$ with their joint harmonics),
the reversed pairs $(U_2,A)$, $(A',U_1)$ of the descending pairs at $(i,n+1)$ and $(i,n)$ (with
$P^\flat_{i,n+1}$, $P^\flat_{i,n}$ and the characters of Proposition~\ref{ptv:chars}), and the family of
$m$-fold twists $\vartheta$ with $\pcm=\pcm_\vartheta$, $\nu=(\nu_{n+1}\vartheta^{n+1})^{-1}$,
$\chi=\chi_\vartheta$, has excess $0$ and weight $(n+1)\Delta$ and satisfies Definition~\ref{ptr:def-square} with (E1) for the
reversed pairs replaced by (c), (d), (H) and (L) of Proposition~\ref{ptw:revedge}. The conclusion of
Theorem~\ref{ptr:square} holds for it (see the proof). Consequently
\textup{(RS)} of \S\ref{ptv:ss-RS} holds, with $\kappa^{\mathrm{rev}}_{i,n}(\vartheta)=\pm\kappa/c(\vartheta)$,
$|c(\vartheta)|^2=C$; so $|\kappa^{\mathrm{rev}}_{i,n}|^2=|\kappa|^2/C$ is constant and
$\kappa^{\mathrm{rev}}_{i,n}$ has modulus $1$ up to norm powers.
\end{theorem}
\begin{proof}
\emph{Ranks and signs.} $a=\operatorname{rk}A=M_{n+1}=\operatorname{rk}A'=b$, so $e=0$, and
$N=M_n+M_{n+2}=2M_{n+1}=a+b$. Write $q(X)$ for the number of negative signs. The sign identity of
\S\ref{pc:ss-exchange} for the square at $(i,n)$ is $q(\pcG_{i,n+1})+q(\pcG_{i+1,n+1})=
q(\pcG_{i,n+2})+q(\pcG_{i+1,n})$; rearranged it is $q(U_1)+q(U_2)=q(A)+q(A')$, the identity required
here. Moreover the number $q=q(U_1)+q(U_2)$ equals the number $q(\pcG_{i,n+1})+q(\pcG_{i+1,n+1})$ of
the square at $(i,n)$, and $N$ is the same; so the inequalities $2q+2\le N\le2\mu_\tau+2q$, with
$\mu=(n+1)\Delta$ the weight of $\nu$, are those verified in \S\ref{pc:ss-exchange} for that square
(Corollary~\ref{ptr:tri-b}).

\emph{(E1).} The pairs $(U_1,A)=(\pcG_{i+1,n},\pcG_{i+1,n+1})$ and $(A',U_2)=(\pcG_{i,n+1},
\pcG_{i,n+2})$ are ascending, from $e_1$ to $e$ and from $e'$ to $e_2$, and are edges by
Theorem~\ref{ptr:edge}; (H) for $(U_1,A)$ is Theorem~\ref{ptq:H}. The pair $(U_2,A)$ is the reversed
pair of the descending pair $(\pcG_{i+1,n+1},\pcG_{i,n+2})$, for which $e_T=\aubar{e_{i,n+2}}=
\aubar{e_2}$, $e_S=\aubar{e_{i+1,n+1}}=\aubar e$: it goes from $\aubar{e_2}$ to $\aubar e$, as
Definition~\ref{ptr:def-square} requires; and $(A',U_1)$ is the reversed pair of
$(\pcG_{i+1,n},\pcG_{i,n+1})$, from $\aubar{e_{i,n+1}}=\aubar{e'}$ to $\aubar{e_{i+1,n}}=\aubar{e_1}$.
By Proposition~\ref{ptw:revedge}, (c), (d) hold for them and (H) for $(U_2,A)$. The polynomials and
types do not depend on $\vartheta$; the characters do, with fixed weights.

\emph{What the proof of Theorem~\ref{ptr:square} uses of (E1).} The transformation laws use only
Lemma~\ref{pc:level} and (E2); the factorisation of the kernels uses only the tensor structure of
the Weil representations; Proposition~\ref{ptr:conv} uses the rapid decay of $y_1,y_2$ ((c)) and
finite moments of the Fock vectors; the multiplicity-one step uses Lemma~\ref{ptr:conv-lift} for
$(U_1,A)$ and $(U_2,A)$, Lemma~\ref{ptr:avg} for the forms $\Lambda^{(1)}_{\xi_1}y_1$ of type $e$ and
$\Lambda^{(2)}_{\xi_2}y_2$ of type $\aubar e$ ((d); $e=e_{i+1,n+1}$ is integrable of the class of
Lemma~\ref{ptr:gap}), and Lemma~\ref{ptr:dense}; the exchange step uses Theorem~\ref{pcx:X} with the
polynomial $v=P'_1P'_2$. Parts (a), (b) of Theorem~\ref{ptr:edge} enter only through (d) and
Lemma~\ref{ptr:conv-lift}, which for the reversed pairs are Proposition~\ref{ptw:revedge}(d), (L).
So the proof applies.

\emph{(E2).} By Lemma~\ref{ptw:W7} the kernel of the reversed pair of the descent $d_{i,k}$ has
source character $\beta^\downarrow_{i,k}$ (on $\pcG_{i,k+1}$) and target character
$\alpha^\downarrow_{i,k}$ (on $\pcG_{i+1,k}$). The four conditions, with $\gamma=\nu^{-1}\upsilon^{N/2}=
\gamma_n$ of \S\ref{ptv:ss-chars}, are:
\begin{align*}
&\text{at }A=\pcG_{i+1,n+1}:&\mu_{i+n+2}\upsilon_{M_n}\cdot\alpha^\downarrow_{i,n+1}&=\gamma_n,&
&\text{at }A'=\pcG_{i,n+1}:&\beta^\downarrow_{i,n}\cdot\mu_{i+n+2}\upsilon_{M_{n+2}}&=\gamma_n,\\
&\text{at }U_1=\pcG_{i+1,n}:&\mu_{i+n+2}\upsilon_{M_{n+1}}\cdot\alpha^\downarrow_{i,n}&=\gamma_n,&
&\text{at }U_2=\pcG_{i,n+2}:&\beta^\downarrow_{i,n+1}\cdot\mu_{i+n+2}\upsilon_{M_{n+1}}&=\gamma_n
\end{align*}
(at $A$: target characters of $a_{i+1,n}$ and of the reversed $d_{i,n+1}$; at $A'$: source
characters of the reversed $d_{i,n}$ and of $a_{i,n+1}$; at $U_1$: source character of $a_{i+1,n}$
times target character of the reversed $d_{i,n}$; at $U_2$: source character of the reversed
$d_{i,n+1}$ times target character of $a_{i,n+1}$). These are exactly the four equations of the
square at $(i,n)$ in \S\ref{ptv:ss-chars} (at $U_2$, $U_1$, $A'$, $A$ there, respectively), which hold
by Proposition~\ref{ptv:chars}.

\emph{Conclusion.} Theorem~\ref{ptr:square} gives $\kappa$, $C>0$ and, for each $\vartheta$, $c$ with
$|c|^2=C\,L_*(\pcm,\nu,0)=C$ (Corollary~\ref{pcx:cne0}), such that $L,L'$ satisfy the square law with
constant $\kappa/c$, for the splitting used in Theorem~\ref{pcx:X}. With $y_1=y$, $y_2=\aubar z$ we
have $L(\xi)=\mathfrak L(\xi)$ for $\xi$ of rank $a$ and $L'(\eta)=\mathfrak L'(\eta)$ for $\eta$ of rank
$b$, the only matrices on which matched functions are supported: indeed $\Lambda^{(1)}_{\xi_1}=a_{i+1,n}[\xi_1]$,
$R_{n+1}[\xi_2]$ is the transposed lift along the reversed pair $(U_2,A)$, so by Fubini (absolute
convergence, Proposition~\ref{ptr:conv}(a))
\[
\mathfrak L(\xi)=\int_{[\auU(U_2)]}\aubar{z(h)}\int_{[\auU(A)]}\theta^{\mathrm{rev}}(h,g;\xi_2)\,
(a_{i+1,n}[\xi_1]y)(g)\,dg\,dh=\int_{[\auU(A)]}\Lambda^{(1)}_{\xi_1}y\cdot\Lambda^{(2)}_{\xi_2}\aubar z
=L(\xi),
\]
and $\mathfrak L'(\eta)=\int_{[\auU(U_2)]}\aubar z\cdot a_{i,n+1}[\eta_2]R_n[\eta_1]y=L'(\eta)$ with
$T'^{(1)}_{\eta_1}=R_n[\eta_1]$, $\Lambda'^{(2)}_{\eta_2}=a_{i,n+1}[\eta_2]$ (Fubini, by
Proposition~\ref{ptr:conv}(b)). The law $\mathfrak L(h\xi)=\chi(\det h)\mathfrak L(\xi)$ for
$h\in\auGL_{\mathsf M}(\pcR)$ and \emph{all} $\xi$ is proved directly as the transformation law in the
proof of Theorem~\ref{ptr:square}: choose $k\in K_A$ reducing to $h$; by Lemma~\ref{pc:level},
$\omega(1,k^{-1})$ on $\auHom(U_1,A)$ and on $\auHom(U_2,A)$ send the point data
at $h\xi_1$, $h\xi_2$ to the target characters at $\det k$, inverted, times the data at $\xi_1,\xi_2$;
so $a_{i+1,n}[h\xi_1]y=c_1(\det k)^{-1}\psrho(k^{-1})a_{i+1,n}[\xi_1]y$, and the substitution
$g\mapsto gk$ in the integral over $[\auU(A)]$ defining $R_{n+1}[h\xi_2]$ gives
$R_{n+1}[h\xi_2]\psrho(k^{-1})=c_2(\det k)^{-1}R_{n+1}[\xi_2]$; with $c_1c_2=\gamma$ ((E2) at $A$) and
$\gamma^{-1}|_{K_A}=\chi\circ\det$ at the primes of $\pcm$, $\mathfrak L(h\xi)=\chi(\det h)\mathfrak L(\xi)$.
By Lemma~\ref{ptw:split} the
law holds for the splitting $(\mathsf B,\mathsf A)=(\mathsf M,\mathsf N\sqcup\mathsf T')$ of (RS) with
constant $\pm\kappa/c$. This is (RS).
\end{proof}

\begin{corollary}[The adjoint law]\label{ptw:adjoint}
Theorem~\ref{pc:tri}\textup{(c)} holds: the adjoint law \textup{(T4)} holds for the triangle, with
functions $\kappa_{i,n}$ of modulus $1$ up to norm powers. It depends on the same inputs as
Theorem~\ref{pc:tri}\textup{(b)}: (V) (Theorem~\ref{pco:V}), Theorem~\ref{pcx:X}, the published results
(D1)--(D5), and (W7) (Lemmas~\ref{pc:SW} and~\ref{ptw:W7}).
\end{corollary}
\begin{proof}
Theorem~\ref{ptv:adjoint} with Theorem~\ref{ptw:RS}. In the constant of Theorem~\ref{ptv:adjoint},
$|c_n|^2$, $|c_{n+1}|^2$, $|\gamma|^2$ are constants times integral powers of $\auNm Q_\vartheta$ and do not
vanish, and $|\kappa^{\mathrm{rev}}_{i,n}|^2$ is constant; Lemma~\ref{pc:mod}(a).
\end{proof}

\subsection{The pairs $(W_j,V_j)$: the type on $V_j$}\label{ptw:ss-A2}

The proof of (A2) for $(W_j,V_j)$ in Proposition~\ref{ptv:arch} uses that $e_{V_j}$ is the type of
the representation of highest weight $\Lambda'$ generated by the joint harmonic $P_j=P^W_j$, as in
the definition of $\Theta_{j,\vartheta}$ in \S\ref{pc:sec}. This subsection proves the facts about
that type which (A2) requires.

Let $\sigma'=\omega_{\alpha,\beta,\infty}(1,\cdot)$ on the Fock space of $\auHom(W_j,V_j)\otimes\R$, with
the characters of the pair, $G'=\auU(V_j)_\infty$ (compact, connected) with Haar measure of total
mass $1$, and let $\rho_j$ be the representation of $G'$ on the span of $\sigma'(G')P_j$. For a
finite-dimensional irreducible unitary representation $\pi$ of $G'$ with a unit highest weight
vector $v$ put $e_\pi=d_\pi\,\aubar{c_\pi}$, $c_\pi(x)=\langle v,\pi(x)v\rangle$ (the convention
$e=d\,\aubar c$ of \S\ref{pc:ss-data}).

\begin{proposition}\label{ptw:A2}
\textup{(a)} $\rho_j$ is irreducible, of highest weight $\Lambda'=\lambda'-\rho'$ at every $\tau$, and
$P_j$ is a highest weight vector of it.
\textup{(b)} $\int_{G'}e_{\rho_j}(x)\,\sigma'(x)P_j\,dx=P_j$.
\textup{(c)} For every $\pcm$, $\Psi$, characters as in \textup{(A2)} and form $y$ of type $e_{W_j}$,
the lift $F(h)=\int_{[\auU(W_j)]}\theta(g,h;P_j,\Psi)\,y(g)\,dg$ converges absolutely, is bounded,
right invariant under an open subgroup of $\auU(V_j)(\auA_f)$, and $R(e_{\rho_j})F=F$.
Hence \textup{(A2)} of Lemma~\ref{pc:arch} holds for $(W_j,V_j)$ with
$e_{V_j}=e_{\rho_j}$, and with no other type of the form $e_\pi$ unless $e_\pi=e_{\rho_j}$.
\end{proposition}
\begin{proof}
(a) By the proof of Lemma~\ref{par:vac}, $K'=G'$ acts on polynomial vectors by linear substitutions
times a character, so $\sigma'(G')P_j$ spans a finite-dimensional unitary representation. By Step~1
of the proof of Proposition~\ref{plf:L4}, $P_j$ (the tensor product of the $\pop_{\xi_0,\tau}$,
Lemma~\ref{plf:harm}) has $T'$-weight $\Lambda'$ and is killed by the $E'_{cd}$ with
$\lambda'_c>\lambda'_d$ (all roots are compact), i.e.\ it is a highest weight vector. A
finite-dimensional representation of a compact connected group generated by a highest weight vector
$v$ of weight $\Lambda'$ is irreducible: decompose $v=\sum_iv_i$ along a decomposition into
irreducibles; each non-zero $v_i$ is a highest weight vector of weight $\Lambda'$, so the summands
with $v_i\ne0$ are isomorphic to the irreducible representation $\pi_{\Lambda'}$; inside the isotypic
part $\pi_{\Lambda'}\otimes M$ the highest weight vectors are $v^+\otimes M$, so $v=v^+\otimes m$ and $v$
generates $\pi_{\Lambda'}\otimes m$.

(b) Put $\hat P=P_j/\|P_j\|$. By Schur orthogonality, for every $u$ in the space of $\rho_j$,
$\langle u,\int e_{\rho_j}(x)\rho_j(x)\hat P\,dx\rangle=d\int\aubar{\langle\hat P,\rho_j(x)\hat P\rangle}
\langle u,\rho_j(x)\hat P\rangle\,dx=\langle u,\hat P\rangle$; both vectors lie in the space of $\rho_j$, so
they are equal.

(c) Convergence, boundedness and invariance are shown in Proposition~\ref{ptv:arch}. With $L_h$ as
in the proof of Lemma~\ref{ptr:conv-lift} ($L_h$ continuous on $\pzF_s$,
$L_h(\sigma'(x)\Psi)=L_{hx}(\Psi)$), $R(e_{\rho_j})F(h)=\int e_{\rho_j}(x)L_h(\sigma'(x)P_j)\,dx=
L_h(\int e_{\rho_j}(x)\sigma'(x)P_j\,dx)=L_h(P_j)=F(h)$ by (b). The last assertion: if $R(e_\pi)F=F$ for
all such $F$ and some $F\ne0$, then $\pi\cong\rho_j$ (the $G'$-span of the right translates of $F$ is
a quotient of $\rho_j$ by the formula $F(hx)=L_h(\sigma'(x)P_j)$), and the highest weight vectors
agree up to a scalar, so $e_\pi=e_{\rho_j}$.
\end{proof}

\begin{remark}[The scalar functions on $V_j$]\label{ptw:rem-A2}
The scalar functions attached to $\mathcal A(V_j;\vartheta)$ in \S\ref{pc:sec} are
$h\mapsto\langle\xi_0,\pi^\natural(h_\infty)^{-1}\varphi(h_f)\rangle$, with $\xi_0$ as in
Lemma~\ref{plf:harm} ($\pop_{\xi_0}=P_j$). This is consistent with Proposition~\ref{ptw:A2}: by
\eqref{po:eq-lift} with $\xi=\pi^\natural(h_\infty)\xi_0$ and $\pop_{\pi^\natural(h_\infty)\xi_0}=
\sigma'(h_\infty)\pop_{\xi_0}$, the scalar lift with $P_j$ is
$h\mapsto\langle\pi^\natural(h_\infty)\xi_0,(\Theta'y)(h_f)\rangle$ (up to the character $\chi\circ\det$ of
Remark~\ref{pc:rem-norm}). Whether $\xi_0$ is the highest or the lowest weight vector of
$W_{\lambda_V}$ for the standard action depends on whether $\pi^\natural$ is the standard action or its
composite with complex conjugation (compare Remark~\ref{par:Ym}); nothing depends on it, since
Lemma~\ref{pc:algpet} holds for every unit vector in place of $v^+$, and (A2), (A3) and the norms
use only Proposition~\ref{ptw:A2}.
\end{remark}

\subsection{The edge hypothesis for the grid squares}\label{ptw:ss-grid}

The edge hypothesis (E1) for the squares of the grids (Remark~\ref{ptr:rem-grid},
Proposition~\ref{pgr:square}) is proved in Lemma~\ref{pgv:E1g} from the block rule of
Proposition~\ref{pgv:match}; every edge of a grid square is read in the direction of
Theorem~\ref{pgv:G3}(d), so Theorems~\ref{ptr:edge} and~\ref{ptq:H} suffice there and the exchanged
forms of this section are not needed.


%% ===== proofs/C-grid.tex =====
% proofs/C-grid.tex
% Theorem pc:auxcore (integral and short lifts through the second grid): architecture, the
% combinatorics of paths through two grids, transfer of periods along turns, and the deduction of
% pc:auxcore (a), (b), (c) from the intermediate propositions. Labels start with pgr:.
% Conventions and macros: those of proofs/C-determinant.tex.
\providecommand{\auI}{\mathcal I}
\providecommand{\auAut}{\operatorname{Aut}}
\providecommand{\auNm}{\operatorname{Nm}}
\providecommand{\auU}{\mathrm U}
\providecommand{\auGL}{\mathrm{GL}}
\providecommand{\auAz}{\mathcal A_0}
\providecommand{\auAm}{\mathcal A_m}
\providecommand{\auO}{\mathcal O}
\providecommand{\pcFF}{\mathfrak F}
\providecommand{\pcm}{\mathfrak m}
\providecommand{\pcb}{\mathfrak b}
\providecommand{\pgrY}{\widetilde Y}
\providecommand{\pgupt}{\mathsf u}
\providecommand{\pcMat}{\mathrm M}
\providecommand{\aubar}[1]{\overline{#1}}
\providecommand{\pgrG}{\Gamma}
\providecommand{\pgrfirst}{\Gamma_{\mathrm{first}}}
\providecommand{\pgrlast}{\Gamma_{\mathrm{last}}}
\providecommand{\pgrsq}{\operatorname{Sq}}
\providecommand{\pgrE}{\operatorname{Ed}}
\providecommand{\pgrTr}{\mathrel{\vdash}}

\section{Lifts through the second grid: Theorem~\ref{pc:auxcore}}\label{pgr:sec}

This section concerns Theorem~\ref{pc:auxcore}, from which Proposition~\ref{au:aux} is deduced in
\S\ref{pc:ss-aux}. The notation is that of \S\ref{pc:sec}. \S\ref{pgr:ss-arch} fixes the objects
and states the intermediate propositions exactly; \S\ref{pgr:ss-paths} proves the combinatorics of
paths through the two grids; \S\ref{pgr:ss-transfer} proves that the turns compose to a transfer
from the first path to the last; \S\ref{pgr:ss-deduce} deduces Theorem~\ref{pc:auxcore} from the
intermediate propositions.

\subsection{Architecture}\label{pgr:ss-arch}

\subsubsection*{The two grids}
For $0\le a,b\le m$ there are two spaces with frames (spaces in the sense of \S\ref{pc:ss-data}),
defined completely in \S\ref{pgv:ss-def} (Definition~\ref{pgv:def-spaces}):
\begin{itemize}
\item $Y^a_b$, \emph{the first grid}: a stem of rank $r$, $a$ blocks of rank $t$ of the first kind and
$b$ blocks of rank $t$ of the second kind; $Y^0_b=Y_b$ with the lattice of \S\ref{ps:ss-Y};
\item $\pgrY^b_a$, \emph{the second grid}: a stem of rank $r+1$, $a$ blocks of rank $t+1$ and $b$
blocks of rank $t-1$; $\pgrY^0_0=Y_0\oplus E\tilde e$ with $\langle\tilde e,\tilde e\rangle=1$.
\end{itemize}
The signs are \eqref{pgv:eq-sign}: they depend only on the second index of the node $(a,b)$.
$Y^m_m$, $\pgrY^m_m$ and $\pgrY^m_0$ are positive definite, $\pgrY^l_l=Y^l_l\oplus E\tilde e_l$ with
$\langle\tilde e_l,\tilde e_l\rangle=(-1)^l$ as spaces with frames, and $\pgrY:=\pgrY^m_0$ is the space
of \S\ref{pc:ss-aux} (rank $\tilde n=r+1+m(t-1)$). A half $h_1$ of a two-grid path runs through the
nodes $(a,b)\mapsto\pgrY^b_a$. Each node carries a type $e_p$ (Harish-Chandra parameter $-\lambda_p$,
\eqref{pgv:eq-param}); the forms moved at a node are the forms of type $e_p$: in the second grid these
are forms, in the first grid \emph{conjugate forms}.

For an $m$-fold twist $\vartheta$ let $\pcm=\pcm(\vartheta)=\pcb\prod_{w\in Q_\vartheta}w$, with
$\pcb=(\pi_M)^b$, and $R_\pcm=\auO_E/\pcm$; here the exponent $b$ is the global parameter of \S\ref{ss:au-spaces}, not the second grid index. This modulus is distinct from the modulus
$\pcm_\vartheta=(\pi_M)\prod_{w\in Q_\vartheta}w$ of \S\ref{pc:ss-aux}, which appears below only
through the finite ring $R_{\pcm_\vartheta}=\auO_E/\pcm_\vartheta$ (a product of finite fields) in
the index $[K_m:J^\circ_\vartheta]=\#\auGL_{r+mt}(R_{\pcm_\vartheta})$ of the principal congruence subgroup
$J^\circ_\vartheta=J_{Y_m}(\pcm_\vartheta)$ (Lemma~\ref{pgr:jcirc}). An \emph{edge} is a pair of adjacent nodes in one grid,
oriented towards the larger index; it is \emph{vertical} if it changes the second index in the first
grid or the first index in the second grid, and \emph{horizontal} otherwise. Its characters
$\alpha_e,\beta_e$, its edge character $\chi_e$ of $R_\pcm^\times$, its joint harmonic $P_e$ and its
\emph{moves} (the lift $\Lambda_{S_e}$ with the point datum $\delta_{S_e}$ when the first index grows,
the transposed lift $T_{S_e}$ when the second index grows, $S_e\in\pcMat_{q\times p}(R_\pcm)$) are
those of Definition~\ref{pgv:def-edges}; $\chi_e$ depends only on the kind and the index of $e$.

\subsubsection*{Paths and turns}
A \emph{grid path} from $(a_0,b_0)$ to $(a,b)$ is a finite sequence of nodes starting at
$(a_0,b_0)$ and ending at $(a,b)$, each step being \emph{down}, $(a',b'+1)\to(a',b')$, or
\emph{right}, $(a',b')\to(a'+1,b')$. A \emph{two-grid path} $\pgrG=(l,h_0,h_1)$ consists of an index
$0\le l\le m$ (the \emph{diagonal index}) and two grid paths $h_0,h_1$ from $(0,m)$ to $(l,l)$, the
\emph{halves}: $h_0$ lives in the first grid (node $(a,b)$ is $Y^a_b$) and $h_1$ in the second
(node $(a,b)$ is $\pgrY^b_a$). Read as one path from $\pgrY^m_0$ to $Y^0_m$, $\pgrG$ goes along $h_1$
backwards in the second grid to $\pgrY^l_l$, crosses to $Y^l_l$, and goes along $h_0$ backwards to
$Y^0_m$. The \emph{first path} $\pgrfirst=(0,c,c)$ has both halves equal to the column
$c\colon(0,m)\to(0,m-1)\to\dots\to(0,0)$; the \emph{last path} $\pgrlast=(m,\rho,\rho)$ has both halves
equal to the row $\rho\colon(0,m)\to(1,m)\to\dots\to(m,m)$.

A \emph{corner} of a grid path is a node $(a,b)$ reached by a step down from $(a,b+1)$ and left by a
step right to $(a+1,b)$. For two-grid paths $\pgrG,\pgrG'$ and $s=(i,j)$ we write
$\pgrG\pgrTr_s\pgrG'$ (\emph{a turn at $s$}) in two cases:
\begin{itemize}
\item[(off)] for some $z\in\{0,1\}$ the half $h_z$ of $\pgrG$ has a corner at $(a,b)$, $\pgrG'$
equals $\pgrG$ except that in $h_z$ the two steps $(a,b+1)\to(a,b)\to(a+1,b)$ are replaced by
$(a,b+1)\to(a+1,b+1)\to(a+1,b)$, and $s=(a,b)$ if $z=0$, $s=(b,a)$ if $z=1$;
\item[(diag)] $\pgrG=(j,h_0,h_1)$ where both halves end with the step $(j,j+1)\to(j,j)$, and
$\pgrG'=(j+1,h'_0,h'_1)$ where $h'_z$ is $h_z$ with this last step replaced by
$(j,j+1)\to(j+1,j+1)$; here $s=(j,j)$.
\end{itemize}
The set of \emph{edges} $\pgrE(\pgrG)$ is the set of edges of $h_0$ in the first grid together with
those of $h_1$ in the second.

\subsubsection*{Path data and periods}
An \emph{edge tuple} for $\vartheta$ is a family $S=(S_e)$ of matrices over $R_\pcm$, one for each
edge. A \emph{node family} is a family $g=(g_p)$ with $g_p\in\auGL_{p}(R_\pcm)$; it is
\emph{admissible} if $g_{Y^0_m}=1$, $g_{\pgrY^m_0}\equiv1\bmod\pcb$, and $g$ at $\pgrY^a_a$ is
$g_{Y^a_a}\oplus1$ for every $a$ (\S\ref{pgv:ss-def}). It acts by
$(gS)_e=g_q^{-1}S_eg_p$ for $e=(p,q)$, and we put $\chi_w(g)=\prod_{e\in w}\chi_e(\det g_q/\det g_p)$
for a set $w$ of edges. A complex number is a \emph{sum of at most $N$ roots of unity} if it is
$\sum_{k<N'}\zeta_k$ with $N'\le N$ and roots of unity $\zeta_k$.

\begin{definition}[Path data]\label{pgr:def-datum}
Let $w$ be a finite set of edges and $N\ge0$. A function $\Psi$ of edge tuples (for $\vartheta$) is a
\emph{path datum on $w$ with bound $N$} if: it depends only on $(S_e)_{e\in w}$;
$\chi_w(g)\Psi(gS)=\Psi(S)$ for every admissible $g$ and every $S$; $\Psi(S)\ne0$ implies that
$S_e$ has a left inverse for every vertical $e\in w$; and every value $\Psi(S)$ is a sum of at most
$N$ roots of unity.
\end{definition}

For a two-grid path $\pgrG$ with diagonal index $l$, a half $h$ in the grid $z$ and an edge tuple
$S$, let $h_*(S)$ be the composite of the moves of $S$ along $h$; it maps conjugate forms on $Y_m$
(if $z=0$), respectively forms on $\pgrY^m_0$ (if $z=1$), to conjugate forms on $Y^l_l$,
respectively forms on $\pgrY^l_l$. A \emph{period triple} is $\xi=(\vartheta,y',\varphi)$ with
$\vartheta$ an $m$-fold twist, $y'$ a conjugate form on $Y_m$ and $\varphi$ a form on $\pgrY^m_0$.
For a real point $g$ of $\auU(\pgrY^l_l)$ the \emph{diagonal period} of a conjugate form $y''$ on
$Y^l_l$ and a form $\varphi''$ on $\pgrY^l_l$ is $\mathsf P_{l,g}(y'',\varphi'')=\int_{[\auU(Y^l_l)]}
y''(x)\,\varphi''((x\oplus\pgupt^l)g)\,dx$, with $\pgupt=\pgupt_\pcm=\aubar u_0/u_0$, $u_0=\alpha
\varpi_\pcm$ (Definition~\ref{pgv:def-period}). For a family $\Psi=(\Psi_\vartheta)$
the \emph{period of $\Psi$ on $\pgrG$ at $g$} is
\[
I_{\pgrG}(\Psi;g)(\xi)=\sum_{S\in\mathcal S_\vartheta(\pgrG)}\Psi_\vartheta(S)\,\mathsf P_{l,g}\bigl(h_{0*}(S)y',
h_{1*}(S)\varphi\bigr),
\]
where $\mathcal S_\vartheta(\pgrG)$ is the finite set of edge tuples that vanish off $\pgrE(\pgrG)$.

\begin{definition}[Proportionality, transfer]\label{pgr:def-transfer}
Let $\Xi$ be a complex function of period triples. $\Psi$ is \emph{proportional to $\Xi$ on
$\pgrG$}, written $\Psi\propto_{\pgrG}\Xi$, if for every real point $g$ of $\auU(\pgrY^l_l)$ there is
$\beta(g)\in\C$ with $I_{\pgrG}(\Psi;g)(\xi)=\beta(g)\Xi(\xi)$ for all period triples $\xi$. For two-grid
paths $\pgrG',\pgrG$, a function $k\ge1$ with integer values on the $m$-fold twists and a function
$c$ with values in $\C^\times$, we say that \emph{$(\pgrG',\pgrG)$ transfers periods with bound $k$
and constant $c$} if for every function $N$ and every family $\Psi'$ such that each
$\Psi'_\vartheta$ is a path datum on $\pgrE(\pgrG')$ with bound $N(\vartheta)$, there is a family
$\Psi$ such that each $\Psi_\vartheta$ is a path datum on $\pgrE(\pgrG)$ with bound
$k(\vartheta)N(\vartheta)$, and $\Psi\propto_{\pgrG}\Xi$ implies
$\Psi'\propto_{\pgrG'}\Xi/c$ (the function $(\vartheta,y',\varphi)\mapsto\Xi(\vartheta,y',\varphi)/
c(\vartheta)$) for every $\Xi$.
\end{definition}

Put $V(\vartheta)=(\#R_{\pcm(\vartheta)})^{2n_m}$ (a bound for the number of column vectors over
$R_\pcm$ at a node) and $\mathcal E(\vartheta)=(\#R_{\pcm(\vartheta)})^{(2n_m)^2}$ (a bound for the
number of matrices of an edge). For $0\le i,j<m$ let $\rho'_{ij}$ be the ramified ratio
\eqref{pc:eq-rhoprime}.

\subsubsection*{The intermediate statements}
The proof of Theorem~\ref{pc:auxcore} consists of the following statements. Those of
\S\S\ref{pgr:ss-paths}--\ref{pgr:ss-deduce} are proved below; the others are stated here exactly and
proved in the later sections cited after each of them (see \S\ref{pgr:ss-status}).

\begin{proposition}[First path]\label{pgr:first}
There are a family $\Psi^0=(\Psi^0_\vartheta)$, each $\Psi^0_\vartheta$ a path datum on
$\pgrE(\pgrfirst)$ with bound $\mathcal E(\vartheta)^{2m}$, a real point $g_0$ of
$\auU(\pgrY^0_0)$, a number $\alpha_{\mathrm{first}}\in\C$ and a surjective linear map
$\mathsf r\colon\pcFF\to\auAz$, $\pcFF$ the space of forms of type $e_{(1,0,m)}$ on $\pgrY^m_0=\pgrY$
fixed by $J(\pcb)$,
such that for every period triple $(\vartheta,y',\varphi)$ with $\varphi\in\pcFF$
\[
\int_{[\auU(Y_m)]}y'\cdot X_{0,m}^\vartheta\,(\mathsf r\varphi)=\alpha_{\mathrm{first}}\cdot
I_{\pgrfirst}(\Psi^0;g_0)(\vartheta,y',\varphi),
\]
where $X^\vartheta_{0,m}=w^\vartheta_{m-1}a_{0,m-1}\circ\dots\circ w^\vartheta_0a_{0,0}$ is the chain of
\S\ref{pc:ss-triangle}, so that $X^\vartheta_{0,m}|_{\auAz}=\iota_\vartheta$ by \eqref{pc:eq-iota}.
\end{proposition}
\noindent Proposition~\ref{pgr:first} is proved in \S\ref{pgw:ss-first} (Theorem~\ref{pgw:first}) from
three inputs stated there: (I) the injectivity of the untwisted chain of the first column of the
second grid, which is Theorem~\ref{pgy:I}; (V) the deep-level vanishing, which is
Theorem~\ref{pgq:V}; and (P) the compatibility of data under the change of modulus $\pcb\to\pcm$,
which is Proposition~\ref{pgy:P}.

\begin{proposition}[One turn]\label{pgr:turn}
For $0\le i,j<m$ there is $c_{ij}$, with values in $\C^\times$ on the $m$-fold twists and of modulus
$\rho'_{ij}$ up to norm powers (Definition~\ref{pc:modulus}), such that every pair
$(\pgrG',\pgrG)$ with $\pgrG'\pgrTr_{(i,j)}\pgrG$ transfers periods with bound $V$ and constant
$c_{ij}$.
\end{proposition}
\noindent Proposition~\ref{pgr:turn} is proved in \S\ref{pgr:ss-off} from
Propositions~\ref{pgr:square} (the square law on a square of a grid) and~\ref{pgr:turndiag} (the
turn on the diagonal), which are stated there.

\begin{proposition}[Last path]\label{pgr:last}
Let $\Xi_{\mathrm{last}}(\Psi)(\vartheta,y',\varphi)=\Lambda(\Psi_\vartheta;y',\varphi)/
\#\auGL_{r+mt}(R_{\pcm_\vartheta})$, where $\Lambda$ is the sum \eqref{pgw:eq-Lambda}: over
$S\in\mathcal S_\vartheta(\pgrlast)$, $\Psi_\vartheta(S)$ times the class integral over
$\Gamma\backslash\auU(Y^m_m)(\auA_f)$, for the measure giving the level of $Y^m_m$ the volume $1$, of
the branching pairing of the algebraic lifts (each normalised by the level of its source node) of
$y'$ and $\varphi$ along the two halves of $\pgrlast$ with the point data $S$, the second evaluated
at $h\oplus\pgupt_\pcm^m$. Then $\Psi\propto_{\pgrlast}\Xi_{\mathrm{last}}(\Psi)$ for every family
$\Psi$.
\end{proposition}
\noindent Proposition~\ref{pgr:last} is proved in \S\ref{pgw:ss-last} (Theorem~\ref{pgw:last}).

\begin{proposition}[Denominators and conjugates on the last path]\label{pgr:den}
There is an integer $N_0\ge1$ with $N_0\le\exp(x^{99}b^2)$ such that, with
$L(\vartheta)=V(\vartheta)^{m^2}\mathcal E(\vartheta)^{2m}$:
\begin{enumerate}
\item[(a)] if $\Psi_\vartheta$ is a path datum on $\pgrE(\pgrlast)$, the pair $(y',\varphi)$ is
fixed by the levels of the two outer nodes, $\iota\in\auI$ with $\ell(\iota)\nmid\auNm Q_\vartheta$,
and both $y'$ and $\varphi$, read as algebraic forms, have gauge at most $1$ at $\iota$ at every
point, then $N_0\Lambda(\Psi_\vartheta;y',\varphi)$ is integral at $\iota$;
\item[(b)] if $\vartheta=1$, $\Psi_1$ is a path datum on $\pgrE(\pgrlast)$ with bound $L(1)$, the
pair is fixed by the outer levels and $\|\sigma(y'(g))\|^2\le B_0$, $\|\sigma(\varphi(g))\|^2\le
B_1$ for all $g$, then $|\sigma(\Lambda(\Psi_1;y',\varphi))|\le\exp(x^{200}b^2)\sqrt{B_0B_1}$.
\end{enumerate}
\end{proposition}
\noindent Here $\Lambda$ is the sum \eqref{pgw:eq-Lambda}, formed with the integral choices of the
kernels and of the branching form of Definition~\ref{pgz:def-choice} (Theorem~\ref{pgw:last} holds for
them, with another constant $\beta(g)$). Part (a) is proved in \S\ref{pgz:ss-a}
(Theorem~\ref{pgz:a}), with $N_0=\mathsf M\,\mathsf d^{2m}$; the class sums in $\Lambda$ carry no
denominators (Proposition~\ref{pgz:expand}, Remark~\ref{pgz:rem-den}). Part (b) is proved in
Theorem~\ref{pgz:b} from the size bounds (B1)--(B3) of \S\ref{pgz:ss-b}, which are proved in
\S\ref{pgt:sec}.

\begin{proposition}[Values read by delta forms]\label{pgr:delta}
Let $\lambda'$ be the weight of the conjugate forms on $Y_m$ dual to the weight of
$\auAm(\vartheta)$, $\langle\ ,\ \rangle_{\mathrm{du}}$ the integral duality pairing between them,
and $h\ge1$ its hook number. For a finite adelic point $x_f$ of $\auU(Y_m)$ and a vector $v$ of the
dual weight let $\delta_{x_f,v}$ be the conjugate delta form at $x_f$ with value $v$. Then:
\begin{enumerate}
\item[(a)] $[K_m:J^\circ_\vartheta]\int_{[\auU(Y_m)]}\delta_{x_f,v}\cdot X^\vartheta_{0,m}y=
\langle v,\iota_\vartheta y(x_f)\rangle_{\mathrm{du}}/\dim\lambda'$ for $y\in\auAz$;
\item[(b)] $\delta_{x_f,v}$ is fixed by the outer level; if $x_f$ is integral above $\ell(\iota)$ (in the sense of Corollary~\ref{pl:wacor}) and $v$
is integral, $\delta_{x_f,v}$ has gauge at most $1$ at $\iota$; and
$\|\sigma(\delta_{x_f,v}(g))\|^2\le\|\sigma(v)\|^2$;
\item[(c)] a vector $u$ with $\langle w,u\rangle_{\mathrm{du}}\in h\cdot\mathcal O_\iota$ for all
integral $w$ is integral at $\iota$; and if $|\sigma\langle v,u\rangle_{\mathrm{du}}|\le K\|\sigma v\|$
for all $v$ then $\|\sigma u\|^2\le C_{\mathrm{du}}K^2$, with $h\,C_{\mathrm{du}}^{1/2}\exp(x^{99}b^2)
\le\exp(x^{201}b^2-x^{200}b^2)$.
\end{enumerate}
\end{proposition}
\noindent Proposition~\ref{pgr:delta} is proved in \S\ref{pgs:ss-delta} (Proposition~\ref{pgs:delta}),
with $C_{\mathrm{du}}=1$.

\begin{proposition}[Short forms of deep level]\label{pgr:short}
$\pcFF$ is spanned by forms that are integral at every $\iota\in\auI$ and whose values have all
conjugates of squared norm at most $\exp(x^{99}b^2)$.
\end{proposition}

\subsection{Paths through two grids}\label{pgr:ss-paths}

For a grid path $h$ let $\pgrsq(h)$ be the set of unit squares, named by their corner with least
indices, between the first column and the steps down of $h$: a step down from $(a,b+1)$ to $(a,b)$
contributes the squares $(i,b)$, $i<a$. For a two-grid path $\pgrG=(l,h_0,h_1)$ put
\[
\pgrsq(\pgrG)=\{(i,j):i,j<l\}\cup\pgrsq(h_0)\cup\{(j,i):(i,j)\in\pgrsq(h_1)\}.
\]
Then $\pgrsq(\pgrfirst)=\emptyset$ (no square lies left of the first column) and
$\pgrsq(\pgrlast)=\{0,\dots,m-1\}^2$ (the halves have no step down and $l=m$).

\begin{lemma}\label{pgr:sqturn}
If $\pgrG\pgrTr_s\pgrG'$, then $s\notin\pgrsq(\pgrG)$, $\pgrsq(\pgrG')=\pgrsq(\pgrG)\cup\{s\}$, and
$s\in\{0,\dots,m-1\}^2$.
\end{lemma}
\begin{proof}
(off) Let $h_z$ have the corner $(a,b)$, so $h_z=q\cdot(a,b+1)\to(a,b)\to(a+1,b)\cdot r$ with $q$ a
path from $(0,m)$ to $(a,b+1)$ and $r$ one from $(a+1,b)$ to $(l,l)$. As $r$ only goes down and
right, $a+1\le l\le b$. A step down from $(a',b'+1)$ to $(a',b')$ contributes squares in the row
$b'$ only, and each row is descended at most once by a path; so the squares of $h_z$ in row $b$
are exactly the $(i,b)$, $i<a$, contributed by the step $(a,b+1)\to(a,b)$, and
$(a,b)\notin\pgrsq(h_z)$. After the turn the step down in row $b$ is $(a+1,b+1)\to(a+1,b)$,
contributing $(i,b)$, $i\le a$, and the other steps down are unchanged: exactly one new square,
$(a,b)$. Next, every square $(i,j)$ of a half of $\pgrG$ satisfies $j\ge l$ (a path from $(0,m)$ to
$(l,l)$ descends only rows $\ge l$) and $i<l$ (the first index along the path is at most $l$, and
a square has first index less than that of its step). Since $a<l$, the pair $(b,a)$ is not a square
of either half; since $b\ge l$, neither $(a,b)$ nor $(b,a)$ lies in $\{(i,j):i,j<l\}$. This gives
$s\notin\pgrsq(\pgrG)$ and $\pgrsq(\pgrG')=\pgrsq(\pgrG)\cup\{s\}$ in both cases $z=0,1$ (for $z=1$
after the transposition $(a,b)\mapsto(b,a)$ in the definition of $\pgrsq(\pgrG)$). Finally $a<l\le
b<m$ since the step from $(a,b+1)$ exists and $b+1\le m$.

(diag) Here $\pgrG=(j,h_0,h_1)$ with both halves ending with $(j,j+1)\to(j,j)$. The square $(j,j)$
is not in $\{i,i'<j\}$, and not a square of a half nor its mirror image, since the squares of a half
of $\pgrG$ have first index $<j$. After the turn, $l=j+1$, so $\{i,i'<j+1\}$ gains the squares
$(j,j)$, $(i,j)$ and $(j,i)$ for $i<j$; on the other hand the last step down of each half, which
contributed the squares $(i,j)$, $i<j$, is removed. So $\pgrsq(h'_0)=\pgrsq(h_0)\smallsetminus\{(i,j):
i<j\}$, $\pgrsq(h'_1)=\pgrsq(h_1)\smallsetminus\{(i,j):i<j\}$, and
\[
\pgrsq(\pgrG')=\{i,i'<j\}\cup\{(j,j)\}\cup\{(i,j),(j,i):i<j\}\cup\bigl(\pgrsq(h_0)\cup\pgrsq(h_1)^{\mathrm t}
\smallsetminus\{(i,j),(j,i):i<j\}\bigr)=\pgrsq(\pgrG)\cup\{(j,j)\},
\]
where $^{\mathrm t}$ denotes the mirror image. As the step from $(j,j+1)$ exists, $j+1\le m$.
\end{proof}

\begin{lemma}\label{pgr:exturn}
Every two-grid path $\pgrG\ne\pgrlast$ admits a turn $\pgrG\pgrTr_s\pgrG'$.
\end{lemma}
\begin{proof}
A grid path from $(a_0,b_0)$ to $(a,b)$ either has a corner, or ends with a step down, or has no
step down (then $b=b_0$): by induction on the length, appending a step right to a path ending with a
step down creates a corner, and appending a step right to a path without step down keeps it so. If
a half of $\pgrG=(l,h_0,h_1)$ has a corner, (off) applies. Otherwise, if $l=m$, both halves go from
$(0,m)$ to $(m,m)$ without a step down, so they are the row and $\pgrG=\pgrlast$. If $l<m$, each half
must descend from row $m$ to row $l$, so it is not without step down and, having no corner, ends with
a step down $(l,l+1)\to(l,l)$; then (diag) applies.
\end{proof}

\begin{proposition}[Turn induction]\label{pgr:induction}
Let $\mathsf M(\pgrG,T)$ be a property of two-grid paths $\pgrG$ and finite sets $T$ of pairs. If
$\mathsf M(\pgrlast,\emptyset)$ holds, and $\mathsf M(\pgrG',T)$ implies
$\mathsf M(\pgrG,T\cup\{s\})$ whenever $\pgrG\pgrTr_s\pgrG'$ and $s\notin T$, then
$\mathsf M(\pgrfirst,\{0,\dots,m-1\}^2)$ holds.
\end{proposition}
\begin{proof}
Put $D(\pgrG)=\{0,\dots,m-1\}^2\smallsetminus\pgrsq(\pgrG)$. We show $\mathsf M(\pgrG,D(\pgrG))$
for every $\pgrG$ by induction on $\#D(\pgrG)$; for $\pgrfirst$ it is the claim. If
$\pgrG=\pgrlast$, $D(\pgrG)=\emptyset$. Otherwise take a turn $\pgrG\pgrTr_s\pgrG'$
(Lemma~\ref{pgr:exturn}). By Lemma~\ref{pgr:sqturn}, $s\in\{0,\dots,m-1\}^2\smallsetminus
\pgrsq(\pgrG)$ and $D(\pgrG)=D(\pgrG')\cup\{s\}$ with $s\notin D(\pgrG')$, so
$\#D(\pgrG')<\#D(\pgrG)$, $\mathsf M(\pgrG',D(\pgrG'))$ holds by induction, and the hypothesis gives
$\mathsf M(\pgrG,D(\pgrG))$.
\end{proof}

\subsection{Transfer of periods from the first path to the last}\label{pgr:ss-transfer}

\begin{lemma}\label{pgr:transfer-alg}
\begin{enumerate}
\item[(a)] $(\pgrG,\pgrG)$ transfers periods with bound $1$ and constant $1$.
\item[(b)] If $(\pgrG',\pgrG)$ transfers periods with bound $k$ and constant $c$, it does so with
every bound $k'\ge k$.
\item[(c)] If $(\pgrG_1,\pgrG_2)$ transfers periods with bound $k$ and constant $c$, and
$(\pgrG_2,\pgrG_3)$ with bound $k'$ and constant $c'$, then $(\pgrG_1,\pgrG_3)$ transfers periods
with bound $kk'$ and constant $cc'$.
\end{enumerate}
\end{lemma}
\begin{proof}
(a) Take $\Psi=\Psi'$. (b) A path datum with bound $kN$ has bound $k'N$. (c) Let $\Psi^1$ be path
data on $\pgrE(\pgrG_1)$ with bounds $N$. The first transfer gives $\Psi^2$ on $\pgrE(\pgrG_2)$ with
bounds $kN$; the second, applied with the bounds $kN$, gives $\Psi^3$ on $\pgrE(\pgrG_3)$ with
bounds $k'kN$. If $\Psi^3\propto_{\pgrG_3}\Xi$, then $\Psi^2\propto_{\pgrG_2}\Xi/c'$, hence
$\Psi^1\propto_{\pgrG_1}(\Xi/c')/c=\Xi/(cc')$.
\end{proof}

\begin{proposition}[First to last]\label{pgr:firstlast}
With the constants $c_{ij}$ of Proposition~\ref{pgr:turn}, $c=\prod_{0\le i,j<m}c_{ij}$ has modulus
$\prod_{0\le i,j<m}\rho'_{ij}$ up to norm powers, and $(\pgrfirst,\pgrlast)$ transfers periods with
bound $V^{m^2}$ and constant $c$.
\end{proposition}
\begin{proof}
The modulus follows from Lemma~\ref{pc:mod}(a). Apply Proposition~\ref{pgr:induction} to the
property ``$(\pgrG,\pgrlast)$ transfers periods with bound $V^{\#T}$ and constant
$\prod_{s\in T}c_s$''. For $(\pgrlast,\emptyset)$ it is Lemma~\ref{pgr:transfer-alg}(a). If
$\pgrG\pgrTr_s\pgrG'$, $s\notin T$, and the property holds for $(\pgrG',T)$, then $s\in\{0,\dots,
m-1\}^2$ by Lemma~\ref{pgr:sqturn}, $(\pgrG,\pgrG')$ transfers with bound $V$ and constant $c_s$ by
Proposition~\ref{pgr:turn}, and Lemma~\ref{pgr:transfer-alg}(c) gives the property for
$(\pgrG,T\cup\{s\})$, with bound $V\cdot V^{\#T}=V^{\#(T\cup\{s\})}$.
\end{proof}

\subsection{Proof of Theorem~\ref{pc:auxcore} from the intermediate statements}\label{pgr:ss-deduce}

Theorem~\ref{pc:auxcore} is to be read with $\pcFF$ the space of forms of level $J(\pcb)$ on
$\pgrY$, of the weight $\tilde\lambda$ of the outer node of the second grid, and with ``$\varphi$
integral at $\iota$'' meaning that $\varphi$, read as an algebraic form, has gauge at most $1$ at
$\iota$ at every point.

\begin{lemma}[The two levels on $Y_m$]\label{pgr:jcirc}
Let $\vartheta$ be an $m$-fold twist, $J_\vartheta$ and $\chi_\vartheta$ the level and the character of
$\auAm(\vartheta)$ (\S\ref{ss:au-spaces}, Lemma~\ref{pz:lastlevel}), and
$J^\circ_\vartheta=J_{Y_m}(\pcm_\vartheta)$ the kernel of the reduction $K_m\to\auGL_{r+mt}
(R_{\pcm_\vartheta})$ (\S\ref{pc:ss-data}).
\begin{enumerate}
\item[(a)] $[K_m:J^\circ_\vartheta]=\#\auGL_{r+mt}(R_{\pcm_\vartheta})$.
\item[(b)] $J^\circ_\vartheta\subseteq J_\vartheta$, $\chi_\vartheta$ is trivial on $J^\circ_\vartheta$,
and every element of $\auAm(\vartheta)$, in particular $\iota_\vartheta y$ for $y\in\auAz$, is right
invariant under $J^\circ_\vartheta$. The group $J^\circ_\vartheta$ is free.
\item[(c)] At $w\in Q_\vartheta$ the reduction of $J_\vartheta$ is the block upper triangular
subgroup of $\auGL_{r+mt}(\auF_w)$ with diagonal blocks of sizes $r,t,\dots,t$, so that
$[J_\vartheta:J^\circ_\vartheta]=C\prod_{w\in Q_\vartheta}g_{w,r}\,g_{w,t}^{\,m}\,q_w^{e_m}$ with
$g_{w,N}=\#\auGL_N(\auF_w)$, $e_m=mrt+\frac12m(m-1)t^2$ and $C\ge1$ an integer independent of
$\vartheta$. In particular $J_\vartheta\ne J^\circ_\vartheta$ as soon as $Q_\vartheta\ne\emptyset$.
\end{enumerate}
\end{lemma}
\begin{proof}
(a) The reduction $K_m\to\auGL_{r+mt}(R_{\pcm_\vartheta})$ is onto (\S\ref{pc:ss-data}), with kernel
$J^\circ_\vartheta$. (b) An element of $J^\circ_\vartheta$ is $\equiv1$ modulo $\pi_M$, hence at the
chosen places, and its reduction at $w\in Q_\vartheta$ is the identity, which is block upper
triangular with trivial diagonal blocks; so it lies in $J_\vartheta$, and $\chi_\vartheta$, a product
of tame characters of the determinants of the diagonal blocks (Lemma~\ref{pz:lastlevel}), is $1$ on
it. A form $\psi$ of level $(J_\vartheta,\chi_\vartheta)$ satisfies $\psi(gk)=\chi_\vartheta(k)\psi(g)$
for $k\in J_\vartheta$, hence $\psi(gk)=\psi(g)$ for $k\in J^\circ_\vartheta$. A subgroup of a free level is free.
(c) The level $J_\vartheta$ consists of the elements of $K_m$ that are $\equiv1$ at the chosen places
and block upper triangular modulo $Q_\vartheta$ (\S\ref{ss:au-spaces}); Lemma~\ref{pz:lastlevel}
shows that $\iota_\vartheta z$ is equivariant under this group: a $k\in K_m$ whose reduction
$\bar k_w$ at $w\in Q_\vartheta$ preserves the flag of frames of $Y_0\subset Y_1\subset\dots\subset
Y_m$ satisfies $\bar k_{i+1}\iota_i=\iota_i\bar k_i$ for its upper left blocks, which is the
condition of Lemma~\ref{ps:level}. By (a) and the Chinese remainder theorem, $[J_\vartheta:
J^\circ_\vartheta]$ is the order of the reduction of $J_\vartheta$ in $\prod_{w\mid\pcm_\vartheta}
\auGL_{r+mt}(\auF_w)$; the factor at the chosen places does not depend on $\vartheta$, and at
$w\in Q_\vartheta$ it is the order $g_{w,r}g_{w,t}^{\,m}q_w^{e_m}$ of the parabolic subgroup, $e_m$
being the dimension of its unipotent radical.
\end{proof}
\noindent Lemma~\ref{pgr:jcirc}(c) is the index used in Lemma~\ref{pc:index}, where the inner product
of $\auAm(\vartheta)$ (normalised by $J_\vartheta$) enters the chain determinant. Below, the delta forms
and the integrals over $[\auU(Y_m)]$ are normalised by $J^\circ_\vartheta$; since
$[K_m:J_\vartheta]/[K_m:J^\circ_\vartheta]=C^{-1}\prod_w(g_{w,r}g_{w,t}^mq_w^{e_m})^{-1}$ is not a norm
power, the two normalisations must not be interchanged.

\begin{proposition}\label{pgr:lastchain}
There are a function $c$
with values in $\C^\times$, of modulus $\prod_{i,j<m}\rho'_{ij}$ up to norm powers, a family $\Psi$
with $\Psi_\vartheta$ a path datum on $\pgrE(\pgrlast)$ with bound $L(\vartheta)$, a number
$\beta\in\C$ and the surjection $\mathsf r$ of Proposition~\ref{pgr:first}, such that for every
period triple with $\varphi\in\pcFF$
\[
[K_m:J^\circ_\vartheta]\int_{[\auU(Y_m)]}y'\cdot X^\vartheta_{0,m}(\mathsf r\varphi)=\beta\,
\frac{\Lambda(\Psi_\vartheta;y',\varphi)}{c(\vartheta)} .
\]
\end{proposition}
\begin{proof}
Let $\Psi^0,g_0,\alpha_{\mathrm{first}},\mathsf r$ be as in Proposition~\ref{pgr:first}, with bounds
$\mathcal E^{2m}$. By Proposition~\ref{pgr:firstlast} there is $\Psi$ on $\pgrE(\pgrlast)$ with
bounds $V^{m^2}\mathcal E^{2m}=L$, and, since $\Psi\propto_{\pgrlast}\Xi_{\mathrm{last}}(\Psi)$
(Proposition~\ref{pgr:last}), $\Psi^0\propto_{\pgrfirst}\Xi_{\mathrm{last}}(\Psi)/c$. At $g_0$ this
gives $\beta_0$ with $I_{\pgrfirst}(\Psi^0;g_0)(\xi)=\beta_0\Lambda(\Psi_\vartheta;y',\varphi)/
(\#\auGL_{r+mt}(R_{\pcm_\vartheta})\,c(\vartheta))$, and $[K_m:J^\circ_\vartheta]=
\#\auGL_{r+mt}(R_{\pcm_\vartheta})$ by Lemma~\ref{pgr:jcirc}(a). With
Proposition~\ref{pgr:first}, $\beta=\alpha_{\mathrm{first}}\beta_0$ serves.
\end{proof}

\begin{proof}[Proof of Theorem~\ref{pc:auxcore}\textup{(a)}, \textup{(b)} from
Propositions~\ref{pgr:first}--\ref{pgr:delta}]
Let $c,\Psi,\beta,\mathsf r$ be as in Proposition~\ref{pgr:lastchain}, and $N_0$ as in
Proposition~\ref{pgr:den}. By Proposition~\ref{pgr:delta}(a) applied to $y=\mathsf r\varphi$, for
every $x_f$ and $v$,
\begin{equation}\label{pgr:eq-value}
\Bigl\langle v,\frac{c(\vartheta)}{c(1)}\iota_\vartheta(\mathsf r\varphi)(x_f)\Bigr\rangle_{\mathrm{du}}=
\beta'\,\Lambda(\Psi_\vartheta;\delta_{x_f,v},\varphi),\qquad\beta'=\frac{\dim\lambda'\cdot\beta}{c(1)}.
\end{equation}
Choose $s\in\C^\times$ and $\varepsilon\in\{0,1\}$ with $s\beta'=h\,\varepsilon N_0$: if
$\beta'=0$ take $s=1$, $\varepsilon=0$; otherwise $s=hN_0/\beta'$, $\varepsilon=1$. Then
\[
\Bigl\langle v,\frac{c(\vartheta)}{c(1)}\iota_\vartheta(s\,\mathsf r\varphi)(x_f)\Bigr\rangle_{\mathrm{du}}
=h\,\varepsilon\,N_0\Lambda(\Psi_\vartheta;\delta_{x_f,v},\varphi).
\]

(a) Let $\ell(\iota)\nmid\auNm Q_\vartheta$ and $\varphi$ be integral at $\iota$. Integrality of an
algebraic form at $\iota$ is, by definition (\S\ref{pl:ss-norm}, (N2)), the bound $\le1$ on the
gauge at $(\iota,x_f)$ of its value at $x_f$, for every $x_f\in\auU(Y_m)(\auA_f)$. It suffices to check
this at points $x_f$ integral above $\ell=\ell(\iota)$: for any $x_f$, Corollary~\ref{pl:wacor}(a)
(with $d=1$) gives $\gamma\in\Gamma$ with $\gamma x_f$ integral above $\ell$, and for an algebraic
form $\psi$ one has $\psi(\gamma x_f)=\gamma\psi(x_f)$, so by (G1) of \S\ref{pl:ss-norm} the gauge
of $\psi(x_f)$ at $(\iota,x_f)$ equals the gauge of $\psi(\gamma x_f)$ at $(\iota,\gamma x_f)$. For such
$x_f$ and integral $v$, $\delta_{x_f,v}$ is fixed by the outer level and has gauge at most $1$
(Proposition~\ref{pgr:delta}(b)); $\varphi\in\pcFF$ is fixed by the level of $\pgrY^m_0$. So
$N_0\Lambda(\Psi_\vartheta;\delta_{x_f,v},\varphi)$ is integral at $\iota$ by
Proposition~\ref{pgr:den}(a), the right side lies in $h\mathcal O_\iota$, and
Proposition~\ref{pgr:delta}(c) shows that $(c(\vartheta)/c(1))\iota_\vartheta(s\mathsf r\varphi)(x_f)$
is integral at $\iota$.

(b) Let $\vartheta=1$ and $\|\sigma(\varphi(g))\|^2\le B$ for all $g$. By
Proposition~\ref{pgr:delta}(b) the delta form satisfies $\|\sigma\delta_{x_f,v}(g)\|^2\le
\|\sigma v\|^2$, and by Proposition~\ref{pgr:den}(b)
$|\sigma\Lambda(\Psi_1;\delta_{x_f,v},\varphi)|\le\exp(x^{200}b^2)\|\sigma v\|\sqrt B$. As
$|\sigma\varepsilon|\le1$, $N_0\le\exp(x^{99}b^2)$ and $h$ is a positive integer,
\[
\bigl|\sigma\langle v,\iota_1(s\mathsf r\varphi)(x_f)\rangle_{\mathrm{du}}\bigr|\le h\exp(x^{99}b^2)
\exp(x^{200}b^2)\sqrt B\,\|\sigma v\| .
\]
By Proposition~\ref{pgr:delta}(c) with $K=h\exp((x^{99}+x^{200})b^2)\sqrt B$,
\[
\|\sigma(\iota_1(s\mathsf r\varphi)(x_f))\|^2\le C_{\mathrm{du}}h^2\exp(2(x^{99}+x^{200})b^2)B\le
\exp(2x^{201}b^2)\,B .
\]
Since $x_f$ was an arbitrary finite adelic point and algebraic forms are functions on
$\auU(Y_m)(\auA_f)$, this is the bound of Theorem~\ref{pc:auxcore}(b) for all $g$. Finally $c$ has modulus $\prod\rho'_{ij}$ up to norm powers by Proposition~\ref{pgr:lastchain}, and
$\mathsf r$ is onto.
\end{proof}

\begin{proof}[Proof of Theorem~\ref{pc:auxcore}\textup{(c)}]
This is Proposition~\ref{pgr:short}, proved in \S\ref{pgs:ss-short} (Proposition~\ref{pgs:short}).
\end{proof}

\subsection{Turns off the diagonal}\label{pgr:ss-off}

Throughout, $\vartheta$ is an $m$-fold twist, $R=R_{\pcm(\vartheta)}$, and for finite sets
$\mathsf X,\mathsf Y$ we write $\mathrm M_{\mathsf X\times\mathsf Y}$ for the matrices over $R$ with
rows indexed by $\mathsf X$ and columns by $\mathsf Y$.

\subsubsection*{Matched functions}
Let $\mathsf N$ be a finite set with a bijection $\varepsilon\colon\mathsf B\sqcup\mathsf A\to
\mathsf N$, and $\chi$ a character of $R^\times$ (extended by $0$). For $\gamma\in\auGL_{\mathsf N}
(R)$ let $y_\gamma\in\mathrm M_{\mathsf N\times\mathsf B}$ be its columns $\varepsilon(\mathsf B)$
and $x_\gamma\in\mathrm M_{\mathsf A\times\mathsf N}$ the rows $\varepsilon(\mathsf A)$ of
$\gamma^{-1}$; matched functions and the square law are as in \S\ref{pc:ss-exchange}.

\begin{lemma}\label{pgr:extend}
A matrix $y\in\mathrm M_{\mathsf N\times\mathsf B}$ with a left inverse is $y_\gamma$ for some
$\gamma\in\auGL_{\mathsf N}(R)$.
\end{lemma}
\begin{proof}
$R=\prod_w\auO_E/w^{k_w}$ is a finite product of local rings, and it suffices to treat one factor
$R_w$. If $Ly=1$, then $R_w^{\mathsf N}=y(R_w^{\mathsf B})\oplus\ker L$, and $\ker L$ is a direct
summand of a free module over a local ring, hence free, of rank $\#\mathsf N-\#\mathsf B=\#\mathsf A$
(compare ranks modulo the maximal ideal). A basis of $\ker L$, indexed by $\mathsf A$, together with
the columns of $y$ gives $\gamma$.
\end{proof}

\begin{lemma}\label{pgr:match}
\begin{enumerate}
\item[(a)] If $x_\gamma=x_{\gamma'}$, there is $h\in\auGL_{\mathsf B}(R)$ with $y_{\gamma'}=y_\gamma h$
and $\det\gamma'=\det\gamma\det h$.
\item[(b)] Let $\Psi$ on $\mathrm M_{\mathsf N\times\mathsf B}$ vanish off the matrices $y_\gamma$ and
satisfy $\chi(\det h)\Psi(yh)=\Psi(y)$ for $h\in\auGL_{\mathsf B}(R)$. Then
$\Phi=\mathrm{match}_\chi\Psi$, defined by $\Phi(x_\gamma)=\chi(\det\gamma)\Psi(y_\gamma)$ and $\Phi=0$
off the $x_\gamma$, is well defined, and $\Phi,\Psi$ are matched. Conversely, if $\Phi,\Psi$ are
matched then $\Phi=\mathrm{match}_\chi\Psi$.
\item[(c)] If $\Phi,\Psi$ are matched, $G\in\auGL_{\mathsf N}(R)$ and $k\in\auGL_{\mathsf A}(R)$,
then $x\mapsto\chi(\det G)\Phi(xG)$ and $y\mapsto\Psi(G^{-1}y)$ are matched, and
$\Phi(kx)=\chi(\det k)^{-1}\Phi(x)$.
\end{enumerate}
\end{lemma}
\begin{proof}
(a) Order $\mathsf N$ as $\mathsf B\sqcup\mathsf A$ by $\varepsilon$ and put $u=\gamma^{-1}\gamma'$.
The rows $\mathsf A$ of $u^{-1}=\gamma'^{-1}\gamma$ are $x_{\gamma'}\gamma=x_\gamma\gamma=(0\ 1)$, so
$u^{-1}$, hence $u$, is block upper triangular with lower right block $1$:
$u=\bigl(\begin{smallmatrix}h&*\\0&1\end{smallmatrix}\bigr)$, $h$ invertible, $\det u=\det h$. Then
$y_{\gamma'}=\gamma u\binom10=\gamma\binom h0=y_\gamma h$. (b) Well defined by (a) and the law of
$\Psi$. The support conditions and the identity $\Phi(x_\gamma)=\chi(\det\gamma)\Psi(y_\gamma)$ are
the definition. A function matched with $\Psi$ vanishes off the $x_\gamma$ and has the same values
there. (c) $x_\gamma G=x_{G^{-1}\gamma}$, $y_{G^{-1}\gamma}=G^{-1}y_\gamma$ and
$\det(G^{-1}\gamma)=\det\gamma/\det G$; the support conditions transform alike. For the last claim
let $\gamma'=\gamma\,\mathrm{diag}(1,k^{-1})$: then $x_{\gamma'}=kx_\gamma$, $y_{\gamma'}=y_\gamma$
and $\det\gamma'=\det\gamma/\det k$.
\end{proof}

\begin{lemma}\label{pgr:compl}
Let $\mathsf N=\mathsf N_1\sqcup\mathsf N_2$ and $\gamma\in\auGL_{\mathsf N}(R)$. If the rows
$\mathsf N_2$ of $y_\gamma$ form a matrix with a left inverse, then the columns $\mathsf N_1$ of
$x_\gamma$ form a matrix with a left inverse; likewise with $\mathsf N_1,\mathsf N_2$ exchanged.
\end{lemma}
\begin{proof}
Let $L\in\mathrm M_{\mathsf B\times\mathsf N_2}$ with $L\,(y_\gamma)_{\mathsf N_2}=1$, and let $P\in
\mathrm M_{(\mathsf B\sqcup\mathsf A)\times\mathsf N}$ have the rows $(0\ L)$ (columns $\mathsf N_1$,
$\mathsf N_2$) and $x_\gamma$. Then $P\gamma=\bigl(\begin{smallmatrix}1&*\\0&1\end{smallmatrix}
\bigr)$ (as $x_\gamma y_\gamma=0$ and $x_\gamma\gamma_{\mathsf A}=1$), so $P$ is invertible. The
block $\mathsf N_1\times\mathsf N_1$ of $P^{-1}P=1$ reads $P^{-1}_{\mathsf N_1,\mathsf A}\,
(x_\gamma)_{\mathsf N_1}=1$, because the columns $\mathsf N_1$ of the rows $\mathsf B$ of $P$
vanish.
\end{proof}

\subsubsection*{A square and its turn}
Fix $z\in\{0,1\}$ and $a\le b<m$. In the grid $z$ let $V'=(a,b)$, $W_1=(a,b+1)$, $W_2=(a+1,b)$,
$V=(a+1,b+1)$; the \emph{old edges} are $e_1=(V',W_1)$, $e_2=(V',W_2)$ and the \emph{new edges}
$f_1=(W_1,V)$, $f_2=(W_2,V)$. Let $\mathsf B,\mathsf A,\mathsf N_1,\mathsf N_2$ be the frames of
$V',V,W_1,W_2$ and $\mathsf N=\mathsf N_1\sqcup\mathsf N_2$; the ranks satisfy
$\#\mathsf B+\#\mathsf A=\#\mathsf N$ (both are twice the rank of $V'$ plus $2t$ in the first grid,
and likewise in the second), and we fix a bijection $\varepsilon\colon\mathsf B\sqcup\mathsf A\to
\mathsf N$. For an edge tuple $S$ put $y(S)=\binom{S_{e_1}}{S_{e_2}}\in\mathrm M_{\mathsf N\times
\mathsf B}$, $x(S)=(S_{f_1}\ S_{f_2})\in\mathrm M_{\mathsf A\times\mathsf N}$, and let $S[y]$,
$S[x]$ be $S$ with the old, respectively new, edges replaced by the blocks of $y$, respectively $x$.
Put $\chi=(\chi_{e_1}\chi_{e_2})^{-1}$.

The edge characters satisfy $\chi_{f_1}=\chi_{e_2}$ and $\chi_{f_2}=\chi_{e_1}$: $e_1$ and $f_2$
change the second index at $b$, $e_2$ and $f_1$ change the first index at $a$, and the character of
an edge depends only on whether it is vertical and on the index that changes. In the grid $z=0$ the
edges $e_1,f_2$ are vertical and $e_2,f_1$ horizontal; in the grid $z=1$ it is the reverse. If
$a<b$, neither $V'$ nor $V$ is an outer node or a diagonal node, so admissible families are
unconstrained at $V'$ and $V$.

From here to Lemma~\ref{pgr:sumturn} assume $a<b$.
Let $w$ be a set of edges none of which has $V'$ as an end or $V$ as its target.

\begin{proposition}[Turning a path datum]\label{pgr:turned}
Let $\Psi'$ be a path datum on $w\cup\{e_1,e_2\}$ with bound $N$. For an edge tuple $S$ let
$\Psi'_S(y)=\Psi'(S[y])$. Then $\Psi(S)=(\mathrm{match}_\chi\Psi'_S)(x(S))$ defines a path datum on
$w\cup\{f_1,f_2\}$ with bound $N$, and for every $S$ the functions $x\mapsto\Psi(S[x])$ and
$\Psi'_S$ are matched for $\chi$.
\end{proposition}
\begin{proof}
\emph{The match exists.} If $\Psi'_S(y)\ne0$, the block of $y$ on the vertical old edge has a left
inverse, hence so has $y$, and $y=y_\gamma$ by Lemma~\ref{pgr:extend}. For
$k\in\auGL_{\mathsf B}(R)$ let $g'$ be the node family equal to $k$ at $V'$ and to $1$ elsewhere; it
is admissible, it maps $S[y]$ to $S[yk]$ (the edges of $w$ do not meet $V'$), and
$\chi_{w\cup\{e_1,e_2\}}(g')=\chi_{e_1}(\det k)^{-1}\chi_{e_2}(\det k)^{-1}=\chi(\det k)$. The
invariance of $\Psi'$ gives $\chi(\det k)\Psi'_S(yk)=\Psi'_S(y)$. By Lemma~\ref{pgr:match}(b),
$\mathrm{match}_\chi\Psi'_S$ is defined and matched with $\Psi'_S$; since $\Psi'_S$ depends on $S$
only through $(S_e)_{e\in w}$, so does the match, and $\Psi(S[x])=(\mathrm{match}_\chi\Psi'_S)(x)$.
In particular $\Psi$ depends only on the new edges and on $w$.

\emph{Invariance.} Let $g$ be admissible, $G=\mathrm{diag}(g_{W_1},g_{W_2})\in\auGL_{\mathsf N}(R)$,
and let $g''$ be $g$ with $1$ at $V'$ (admissible). Then $g''(S[y])=(gS)[G^{-1}y]$, and the
invariance of $\Psi'$ gives $\Psi'_S(y)=\mathsf a\,\Psi'_{gS}(G^{-1}y)$ with
$\mathsf a=\chi_{e_1}(\det g_{W_1})\chi_{e_2}(\det g_{W_2})\chi_w(g)$. By Lemma~\ref{pgr:match}(c) and
the uniqueness in (b), $(\mathrm{match}_\chi\Psi'_S)(x)=\mathsf a\,\chi(\det G)\,(\mathrm{match}_\chi
\Psi'_{gS})(xG)$. Moreover $x(gS)=g_V^{-1}x(S)G$, so by the last claim of
Lemma~\ref{pgr:match}(c), $\Psi(gS)=\chi(\det g_V)(\mathrm{match}_\chi\Psi'_{gS})(x(S)G)$. Hence
$\Psi(S)=\mathsf a\,\chi(\det G)\chi(\det g_V)^{-1}\Psi(gS)$, and
\[
\mathsf a\,\chi(\det G)\chi(\det g_V)^{-1}=\chi_{e_2}(\det g_{W_1})^{-1}\chi_{e_1}(\det g_{W_2})^{-1}
\chi_{e_1}\chi_{e_2}(\det g_V)\,\chi_w(g)=\chi_{w\cup\{f_1,f_2\}}(g)
\]
by $\chi_{f_1}=\chi_{e_2}$, $\chi_{f_2}=\chi_{e_1}$.

\emph{Left inverses.} Let $\Psi(S)\ne0$. Then $x(S)=x_\gamma$ with $\Psi'_S(y_\gamma)\ne0$. The
edges of $w$ carry the same matrices in $S$ and $S[y_\gamma]$, so their vertical ones are split
injective. If $f_1$ is vertical, so is $e_2$; then the rows $\mathsf N_2$ of $y_\gamma$ are split
injective, and the columns $\mathsf N_1$ of $x_\gamma$, which form $S_{f_1}$, are split injective
by Lemma~\ref{pgr:compl}. The case of $f_2$ is symmetric.

\emph{Values.} $\Psi(S)=\chi(\det\gamma)\Psi'_S(y_\gamma)$ is a root of unity times a sum of at most
$N$ roots of unity.
\end{proof}

\begin{lemma}[Sums over a turned square]\label{pgr:sumturn}
In the situation of Proposition~\ref{pgr:turned}, let $J$ be a set, $C\in\C$, and for $i\in J$ let
$L_i$ on $\mathrm M_{\mathsf A\times\mathsf N}$ and $L'_i$ on $\mathrm M_{\mathsf N\times\mathsf B}$
satisfy the square law for $\chi$ with constant $C$. Let $i(S)\in J$ depend only on
$(S_e)_{e\in w}$. Then
\[
\sum_{S\in\mathcal S(w\cup\{e_1,e_2\})}\Psi'(S)\,L'_{i(S)}(y(S))=C\sum_{S\in\mathcal S(w\cup\{f_1,
f_2\})}\Psi(S)\,L_{i(S)}(x(S)),
\]
$\mathcal S(u)$ being the set of edge tuples vanishing off $u$.
\end{lemma}
\begin{proof}
Since $e_1,e_2\notin w$, $(S,y)\mapsto S[y]$ is a bijection from $\mathcal S(w)\times\mathrm
M_{\mathsf N\times\mathsf B}$ onto $\mathcal S(w\cup\{e_1,e_2\})$, and likewise for the new edges.
For fixed $S\in\mathcal S(w)$ the square law applied to the matched pair of
Proposition~\ref{pgr:turned} gives $\sum_y\Psi'_S(y)L'_{i(S)}(y)=C\sum_x\Psi(S[x])L_{i(S)}(x)$; sum
over $S$.
\end{proof}

\subsubsection*{The square law on a square of a grid}
For $f$ a form at $W_1$ (a conjugate form if $z=0$) and $l$ a linear functional on the forms at
$W_2$, let
\[
L_{f,l}(x)=l\bigl(T_{x_{\mathsf N_2}}\Lambda_{x_{\mathsf N_1}}f\bigr),\qquad
L'_{f,l}(y)=l\bigl(\Lambda_{y_{\mathsf N_2}}T_{y_{\mathsf N_1}}f\bigr),
\]
where $\Lambda_M$ and $T_M$ denote the moves (lift and transposed lift) with the point datum $M$
along the edge concerned: $L_{f,l}$ goes through $V$, $L'_{f,l}$ through $V'$.

\begin{proposition}[Square law on a grid square]\label{pgr:square}
For $z\in\{0,1\}$ and $a,b<m$ there are $\kappa\in\C$ and a function $c^z_{ab}$ with values in
$\C^\times$, of modulus $\rho'_{ab}$ if $z=0$ and $\rho'_{ba}$ if $z=1$ up to norm powers, such that
for every $\vartheta$, every $f$ and every $l$ the functions $L_{f,l}$, $L'_{f,l}$ satisfy the square
law for $\chi$ with constant $\kappa/c^z_{ab}(\vartheta)$.
\end{proposition}
\noindent Proposition~\ref{pgr:square} is proved in \S\ref{pgv:ss-square} (Theorem~\ref{pgv:square})
from Theorem~\ref{ptr:square}; the hypothesis (E1) of an edge square holds for the four pairs of a
grid square by Lemma~\ref{pgv:E1g}. The character is $\nu=(\mu^0_{a+1}\mu^0_{b+1}\vartheta^n)^{-1}$,
$n=b+1$ if $z=0$ and $n=a+1$ if $z=1$.

\begin{proposition}[A turn off the diagonal]\label{pgr:turnoff}
Let $z\in\{0,1\}$, $a<b<m$, and let $(\pgrG',\pgrG)$ be a turn
of type (off) at the corner $(a,b)$ of the half $h_z$. Then $(\pgrG',\pgrG)$ transfers periods with
bound $1$ and constant $c^z_{ab}$.
\end{proposition}
\begin{proof}
Write $h_z=q\cdot(V'$-corner$)\cdot r$ as in the definition of (off), and let $w$ be the set of the
edges of $q$, of $r$ and of the other half. The edges of $q$ have both ends in the rectangle
$\{a'\le a,\ b'\ge b+1\}$, those of $r$ in $\{a'\ge a+1,\ b'\le b\}$, and those of the other half in
the other grid; so no edge of $w$ meets $V'$ or ends at $V$, and $\pgrE(\pgrG')=w\cup\{e_1,e_2\}$,
$\pgrE(\pgrG)=w\cup\{f_1,f_2\}$. Let $\Psi'$ be path data on $\pgrE(\pgrG')$ with bounds $N$, and let
$\Psi$ be given by Proposition~\ref{pgr:turned}; it has bounds $N$.

Let $(\vartheta,y',\varphi)$ be a period triple and $g$ a real point. Let $F$ be the starting form of
the half $h_z$ ($y'$ if $z=0$, $\varphi$ if $z=1$). For $S$ put $f_S=q_*(S)F$, a form at $W_1$, and
let $l_S$ be the functional $u\mapsto\mathsf P_{l,g}(r_*(S)u,\,h_{1*}(S)\varphi)$ if $z=0$, and
$u\mapsto\mathsf P_{l,g}(h_{0*}(S)y',\,r_*(S)u)$ if $z=1$. Both depend on $S$ only through $w$. As
$h_{z*}$ for $\pgrG'$ is $r_*\circ\Lambda_{S_{e_2}}\circ T_{S_{e_1}}\circ q_*$ and for $\pgrG$ it is
$r_*\circ T_{S_{f_2}}\circ\Lambda_{S_{f_1}}\circ q_*$,
\[
I_{\pgrG'}(\Psi';g)(\xi)=\sum_S\Psi'_\vartheta(S)\,L'_{f_S,l_S}(y(S)),\qquad
I_{\pgrG}(\Psi;g)(\xi)=\sum_S\Psi_\vartheta(S)\,L_{f_S,l_S}(x(S)).
\]
By Proposition~\ref{pgr:square} and Lemma~\ref{pgr:sumturn} (with $J$ the pairs $(f,l)$),
$I_{\pgrG'}(\Psi';g)(\xi)=\kappa\,c^z_{ab}(\vartheta)^{-1}I_{\pgrG}(\Psi;g)(\xi)$. If
$I_{\pgrG}(\Psi;g)=\beta(g)\Xi$, then $I_{\pgrG'}(\Psi';g)=\kappa\beta(g)\,\Xi/c^z_{ab}$.
\end{proof}

\begin{proposition}[A turn on the diagonal]\label{pgr:turndiag}
For $0\le j<m$ there is a function $c_{jj}$ with values in $\C^\times$, of modulus $\rho'_{jj}$ up to
norm powers, such that every turn $\pgrG'\pgrTr_{(j,j)}\pgrG$ transfers periods with bound $V$ and
constant $c_{jj}$.
\end{proposition}
\noindent Proposition~\ref{pgr:turndiag} is proved in \S\ref{pgu:sec} (Theorem~\ref{pgu:turndiag})
from Proposition~\ref{pgr:square} and the properties (G1)--(G4) of the grid data listed in
\S\ref{pgu:ss-data}, which are proved in \S\ref{pgv:sec}, with the line reversal of
Lemma~\ref{pgx:line}; $c_{jj}=c^0_{jj}/c'_{\mathrm{line}}$.

\begin{proof}[Proof of Proposition~\ref{pgr:turn} from Propositions~\ref{pgr:square}
and~\ref{pgr:turndiag}]
Let $\pgrG'\pgrTr_{(i,j)}\pgrG$. In case (off) with $z=0$ at the corner $(a,b)$ we have $(i,j)=(a,b)$
and $a<b$ (proof of Lemma~\ref{pgr:sqturn}); with $z=1$, $(i,j)=(b,a)$ and $a<b$; in case (diag),
$i=j$. So for $i<j$ the turn is (off) with $z=0$ at $(i,j)$, and we take $c_{ij}=c^0_{ij}$, of
modulus $\rho'_{ij}$; for $i>j$ it is (off) with $z=1$ at $(j,i)$, and $c_{ij}=c^1_{ji}$, of modulus
$\rho'_{ij}$; for $i=j$ we take $c_{jj}$ of Proposition~\ref{pgr:turndiag}. By
Proposition~\ref{pgr:turnoff} and Lemma~\ref{pgr:transfer-alg}(b) (as $V\ge1$) the off-diagonal
turns transfer periods with bound $V$.
\end{proof}

\subsection{Where the propositions of this section are proved}\label{pgr:ss-status}
Theorem~\ref{pc:auxcore}(a),(b) is deduced above from Propositions~\ref{pgr:first}--\ref{pgr:delta}, and
Theorem~\ref{pc:auxcore}(c) is Proposition~\ref{pgr:short}. These propositions are proved as follows:
Proposition~\ref{pgr:first} in \S\ref{pgw:ss-first} (Theorem~\ref{pgw:first}), with its inputs in
\S\S\ref{pgy:sec}, \ref{pgq:sec}--\ref{pgm:sec}; Proposition~\ref{pgr:square} in \S\ref{pgv:ss-square}
(Theorem~\ref{pgv:square}, with Lemma~\ref{pgv:E1g}); Proposition~\ref{pgr:turndiag} in \S\ref{pgu:sec}
(Theorem~\ref{pgu:turndiag}), using \S\ref{pgv:sec} and Lemma~\ref{pgx:line}; Proposition~\ref{pgr:last} in
\S\ref{pgw:ss-last} (Theorem~\ref{pgw:last}); Proposition~\ref{pgr:den} in \S\S\ref{pgz:ss-a},
\ref{pgz:ss-b}, \ref{pgt:sec} (Theorems~\ref{pgz:a}, \ref{pgz:b}); Propositions~\ref{pgr:delta} and
\ref{pgr:short} in \S\ref{pgs:sec}. Besides these proofs, Theorem~\ref{pc:auxcore} uses Theorem~\ref{pcx:X}
and the published results named in \S\S\ref{ptr:sec}--\ref{ptw:sec} and \S\S\ref{pgq:sec}--\ref{pgm:sec}.



%% ===== proofs/C-grid2.tex =====
% proofs/C-grid2.tex
% Continuation of proofs/C-grid.tex (labels pgr:): proofs of the open propositions pgr:delta and
% pgr:short (= pc:auxcore(c)). Labels start with pgs:. Conventions and macros: those of
% proofs/C-grid.tex and proofs/C-determinant.tex.
\providecommand{\auI}{\mathcal I}
\providecommand{\auAut}{\operatorname{Aut}}
\providecommand{\auGL}{\mathrm{GL}}
\providecommand{\auU}{\mathrm U}
\providecommand{\auO}{\mathcal O}
\providecommand{\pcFF}{\mathfrak F}
\providecommand{\pcb}{\mathfrak b}
\providecommand{\auA}{\mathbb A}
\providecommand{\pgsT}{\mathbb T}
\providecommand{\pgsC}{\mathsf C}
\providecommand{\pgsev}{\mathrm{ev}}
\providecommand{\pgsdu}{\mathrm{du}}

\section{Lifts through the second grid: Weyl modules, delta forms and short forms}\label{pgs:sec}

This section continues \S\ref{pgr:sec} and uses its notation. It proves
Proposition~\ref{pgr:delta} (values read by delta forms; \S\ref{pgs:ss-delta}), after the facts on
integral Weyl modules that it needs (\S\ref{pgs:ss-weyl}), and Proposition~\ref{pgr:short}, that
is Theorem~\ref{pc:auxcore}(c) (\S\ref{pgs:ss-short}).

\subsection{Integral Weyl modules and their duality}\label{pgs:ss-weyl}

All statements are multiplicative over the infinite places $\tau$, and we write them for one
group $\auGL_N$; for a weight $\lambda=(\kappa,e)$ of $\prod_\tau\auGL_N$ one takes tensor products
over $\tau$ (and $\det^{e_\tau}$ at $\tau$). Let $\mu$ be a diagram with columns $\mathfrak c$ of
heights $1\le h_{\mathfrak c}<N$, given by the numbers $\kappa(h)$ of columns of height $h$, and
$|\mu|=\sum_{\mathfrak c}h_{\mathfrak c}$ its number of boxes. We use the model of
\S\ref{pk:ss-branch}: the tensor space $\pgsT_\mu(R)=\bigotimes_{\mathfrak c}\Lambda^{h_{\mathfrak
c}}R^N$ with the orthonormal basis $(e_U)$ indexed by the fillings $U$, the evaluation tensors
$\pgsev_y=\bigotimes_{\mathfrak c}y_1\wedge\dots\wedge y_{h_{\mathfrak c}}$ ($y_i$ the columns of
$y\in\mathrm M_N(R)$), their span $\pgsC_\mu\subseteq\pgsT_\mu(\C)$, the products of minors
$b_U(y)=(e_U,\pgsev_y)$, and the isomorphism $\mathsf m\colon\pgsC_\mu\to W_\mu$,
$\mathsf m(u)(y)=(u,\pgsev_y)$, onto the span $W_\mu$ of the products of minors; $\Sigma_\mu$ is the
$\Z$-span of the $b_U$, with the $\Z$-basis $(b_T)$ of standard products of minors, and
$\Sigma_\mu^\#$ its dual lattice for the invariant form.

\subsubsection*{The coordinates of the paper}
The weight $\lambda$ is realised (\S\ref{ss:au-spaces}) on $W_\lambda$ with the basis $(e_T)$ dual to
the standard products of minors, with integral lattice $W_\lambda(\Z)=\bigoplus\Z e_T$, and with an
invariant hermitian form whose Gram matrix $G_\lambda$ in this basis is a positive definite integer
matrix. We identify $W_\lambda$ with $W_\mu$ (twisted by $\det^{e}$) through the invariant form, so
that $e_T$ is the element of $W_\mu$ with $\langle e_T,b_{T'}\rangle=\delta_{TT'}$. By
\eqref{pk:eq-dual}, $\langle\mathsf m u,b_{U}\rangle=(u,e_U)$ for $u\in\pgsC_\mu$; hence
\begin{equation}\label{pgs:eq-lattice}
W_\lambda(\Z)=\Sigma_\mu^\#=\mathsf m\bigl(\pgsC_\mu\cap\pgsT_\mu(\Z)\bigr),\qquad
\pgsev(y):=\sum_Tb_T(y)\,e_T=\mathsf m(\pgsev_y),
\end{equation}
the second because $\langle\mathsf m\,\pgsev_y,b_T\rangle=(\pgsev_y,e_T)=b_T(y)$. The invariant form
of $W_\lambda$ is the standard form of $\pgsT_\mu$ restricted to $\pgsC_\mu$
(\S\ref{pk:ss-branch}); in particular
\begin{equation}\label{pgs:eq-evnorm}
\|\pgsev(y)\|^2=\prod_{\mathfrak c}\det\bigl(y_{[h_{\mathfrak c}]}^*\,y_{[h_{\mathfrak c}]}\bigr)
\le\prod_{\mathfrak c}\ \prod_{i\le h_{\mathfrak c}}\|y_i\|^2 ,
\end{equation}
where $y_{[h]}$ consists of the first $h$ columns of $y$; the inequality is Hadamard's. A vector of
$W_\lambda$ is \emph{integral at $\iota$} if its coordinates in $(e_T)$ lie in $\auO_\iota$.

\begin{lemma}[Gram matrix and Fischer form]\label{pgs:gram}
Let $F_\mu=(\langle b_T,b_{T'}\rangle_F)_{T,T'}$ be the matrix of the Fischer form of
Lemma~\ref{pk:hook} on the standard products of minors, and $H_\mu$ the hook number of that lemma.
Then $F_\mu$ has integer entries and $G_\lambda F_\mu=F_\mu G_\lambda=H_\mu\cdot1$.
\end{lemma}
\begin{proof}
The entries are integers since the $b_T$ have integer coefficients. By Lemma~\ref{pk:hook}, the
invariant form of $W_\mu$ is $H_\mu^{-1}\langle\ ,\,\rangle_F$, so the Gram matrix of $(b_T)$ is
$H_\mu^{-1}F_\mu$. The basis $(e_T)$ is dual to $(b_T)$ for this real symmetric form, so its Gram
matrix is the inverse, $G_\lambda=H_\mu F_\mu^{-1}$.
\end{proof}

\begin{lemma}[Duality by complementary minors]\label{pgs:star}
Let $\kappa'$ be the diagram with $\kappa'(h)=\kappa(N-h)$ for $1\le h<N$, $\mu'$ the corresponding
diagram, and $d=\sum_{1\le h<N}\kappa(h)$. There is an integer matrix $J$ from coordinates of
weight $\kappa$ to coordinates of weight $\kappa'$ (in the bases $(e_T)$) which has an integral
inverse, satisfies ${}^tJ\,G_{\mu'}J=G_\mu$, and $S_{\kappa'}({}^tg^{-1})\,J=\det(g)^{-d}\,J\,
S_\kappa(g)$ for $g\in\auGL_N(\C)$, where $S$ denotes the Schur matrices.
\end{lemma}
\begin{proof}
For a subset $U\subseteq[N]$ of $h$ elements let $\star e_U=\epsilon_Ue_{U^c}$ with the sign
$\epsilon_U$ of the permutation $(U,U^c)$ of $[N]$; then $\star\colon\Lambda^hR^N\to\Lambda^{N-h}R^N$
is a signed bijection of bases, characterised by $(\star\omega,\beta)\,e_{[N]}=\omega\wedge\beta$
for $\omega\in\Lambda^h$, $\beta\in\Lambda^{N-h}$ (as $e_U\wedge e_{U^c}=\epsilon_Ue_{[N]}$). For
invertible $g$, $(\star\Lambda^hg\,\omega,\beta)e_{[N]}=g\omega\wedge g(g^{-1}\beta)=\det(g)\,
\omega\wedge g^{-1}\beta=\det(g)\,(\Lambda^{N-h}({}^tg^{-1})\star\omega,\beta)e_{[N]}$, that is
(Jacobi's identity for complementary minors)
$\star\circ\Lambda^hg=\det(g)\,\Lambda^{N-h}({}^tg^{-1})\circ\star$. Let $P=\bigotimes_{\mathfrak c}
\star\colon\pgsT_\mu\to\pgsT_{\mu'}$, the columns of height $h$ going to columns of height $N-h$; $P$
is a signed permutation of the orthonormal bases, so it is orthogonal and preserves the integral
tensors, and $P\circ(\bigotimes\Lambda^{h_{\mathfrak c}}g)=\det(g)^{d}(\bigotimes\Lambda^{N-h_{\mathfrak
c}}{}^tg^{-1})\circ P$. The subspace $\pgsC_\mu$ is the $\auGL_N$-submodule generated by
$\pgsev_1$ (as $g\,\pgsev_y=\pgsev_{gy}$), and $P\pgsev_1=\pm\pgsev_{w_0}$ with $w_0$ the matrix
of the reversal of $[N]$ (each $\star e_{[h]}=\pm e_{\{h+1,\dots,N\}}$, the exterior product of the first $N-h$ columns $e_N,\dots,e_{h+1}$ of $w_0$); so $P$ maps $\pgsC_\mu$ onto
the submodule generated by $\pgsev_{w_0}$, which is $\pgsC_{\mu'}$ ($\dim\pgsC_\mu=\dim\pgsC_{\mu'}$,
both irreducible). As $P$ and $P^{-1}={}^tP$ preserve integral tensors, $P$ maps
$\pgsC_\mu\cap\pgsT_\mu(\Z)$ onto $\pgsC_{\mu'}\cap\pgsT_{\mu'}(\Z)$. Let $J$ be the matrix of
$\mathsf m'\circ P\circ\mathsf m^{-1}$ in the bases $(e_T)$, $(e_{T'})$. By \eqref{pgs:eq-lattice}
it maps $W_\mu(\Z)$ onto $W_{\mu'}(\Z)$, so it and its inverse are integral; it is isometric because
$P$ is orthogonal and the invariant forms are the restrictions of the standard forms; and the
intertwining relation is that of $P$ (for the action of $g$ on $W_\mu$ by $S_\kappa(g)$, the
transpose conventions of \S\ref{pl:ss-norm} being the same on both sides).
\end{proof}

\begin{lemma}[The duality pairing]\label{pgs:pairing}
Let $\lambda=(\kappa,e)$ and $\lambda'=(\kappa',e')$ be weights of $\prod_\tau\auGL_N$ whose doubled
Harish-Chandra parameters are $l$ and $-l$. Then $\kappa'(h)=\kappa(N-h)$ for $1\le h<N$ and
$e+e'=-d$ at every $\tau$. Put $A=G_{\lambda'}J$, with $J$ the tensor product over $\tau$ of the
matrices of Lemma~\ref{pgs:star}, and $\langle v',v\rangle_{\pgsdu}={}^tv'Av$ for coordinate vectors
$v'$ of weight $\lambda'$ and $v$ of weight $\lambda$. Then:
\begin{enumerate}
\item[(P1)] $\langle\ ,\,\rangle_\pgsdu$ is invariant: $\langle\rho'(g)v',\rho(g)v\rangle_\pgsdu=
\langle v',v\rangle_\pgsdu$ for $g\in\prod_\tau\auGL_N(\C)$;
\item[(P2)] there is an integer matrix $C$ with $CA=h\cdot1$, where $h=H_{\lambda'}=\prod_\tau
H_{\mu'_\tau}$ is the hook number of $\lambda'$;
\item[(P3)] for every $v$ of weight $\lambda$ the vector $v^\flat=\overline{Jv}$ of weight $\lambda'$
satisfies $\|v^\flat\|=\|v\|$ and $\langle a,v\rangle_\pgsdu=\langle v^\flat,a\rangle$ for all $a$,
where $\langle p,q\rangle={}^t\bar pG_{\lambda'}q$ is the invariant hermitian form of weight
$\lambda'$.
\end{enumerate}
\end{lemma}
\begin{proof}
With $\mu_p=e+\sum_{h>p}\kappa(h)$ the highest weight, the doubled parameter is
$(2\mu_p+N-1-2p)_{p<N}$; its negative, reordered decreasingly, is the doubled parameter of the
highest weight $\mu'_p=-\mu_{N-1-p}$. Since $\mu_p-\mu_{p+1}=\kappa(p+1)$, this gives
$\kappa'(h)=\mu'_{h-1}-\mu'_h=\mu_{N-h-1}-\mu_{N-h}=\kappa(N-h)$ and $e'=\mu'_{N-1}=-\mu_0=-e-d$.

(P1) The invariance of the hermitian form $G=G_{\lambda'}$ says ${}^t\overline{\rho'(g)}G\rho'(g)=G$
for $g$ in the compact group $\prod_\tau\auU(N)$; since $\rho'$ has integral (hence real)
coefficients in the entries of $g$ and in $\det g^{\pm1}$, $\overline{\rho'(g)}=\rho'(\bar g)=
\rho'({}^tg^{-1})$, so ${}^t\rho'(g)\,G\,\rho'({}^tg^{-1})=G$ on $\prod\auU(N)$, hence on
$\prod\auGL_N(\C)$ (Zariski density of the unitary group). By Lemma~\ref{pgs:star} and $e+e'=-d$,
$J\rho(g)=J\det(g)^eS_\kappa(g)=\det(g)^{e+d}S_{\kappa'}({}^tg^{-1})J=\rho'({}^tg^{-1})J$. So
${}^t(\rho'(g)v')GJ\rho(g)v={}^tv'\,{}^t\rho'(g)G\rho'({}^tg^{-1})Jv={}^tv'GJv$.

(P2) By Lemma~\ref{pgs:gram} for $\lambda'$, $F_{\lambda'}G_{\lambda'}=h\cdot1$; with $J^{-1}$
integral (Lemma~\ref{pgs:star}), $C=J^{-1}F_{\lambda'}$ gives $CA=J^{-1}F_{\lambda'}G_{\lambda'}J=h$.

(P3) $\langle a,v\rangle_\pgsdu={}^taGJv={}^t(Jv)Ga$ ($G=G_{\lambda'}$ symmetric), which is
$\langle\overline{Jv},a\rangle$. And $\|\overline{Jv}\|^2={}^t\overline v\,{}^tJG_{\lambda'}J\,v=
\|v\|^2$ by the isometry of $J$ (Lemma~\ref{pgs:star}; $J$ is real).
\end{proof}

\begin{remark}
The pairing of Proposition~\ref{pgr:delta} is this $A=G_{\lambda'}J$, and it
measures integrality in the coordinates $(e_T)$; by \eqref{pgs:eq-lattice} this is the lattice
$W_\lambda(\Z)=\Sigma_\mu^\#$, as the plan of \S\ref{pgr:sec} required. Lemma~\ref{pgs:gram} is the
precise form of the inclusion $H_\mu\Sigma_\mu\subseteq\Sigma_\mu^\#$ of Lemma~\ref{pk:hook} that is
used: in the coordinates $(e_T)$ the vector $b_{T'}$ has coordinates $G_\lambda^{-1}$ times the
unit vector, and $H_\mu G_\lambda^{-1}=F_\mu$ is integral.
\end{remark}

\subsection{Values read by delta forms: proof of Proposition~\ref{pgr:delta}}\label{pgs:ss-delta}

\subsubsection*{The weights}
Let $\lambda=(\kappa,e)$ be the weight of $\auAm(\vartheta)$ on $Y_m$ (of rank $N=r+mt$), with
doubled Harish-Chandra parameter $l$, and $\lambda'=(\kappa',e')$ the weight of the conjugate forms
on $Y_m$, the weight of the outer node $Y^0_m$ of the first grid; its doubled parameter is $-l$
(the type of a conjugate form is the negated type). So Lemma~\ref{pgs:pairing} applies, and we let
$\langle\ ,\,\rangle_\pgsdu$ be the pairing $A=G_{\lambda'}J$ of that lemma and $h=H_{\lambda'}$.
The forms on $Y_m$ of weight $\lambda$ and the conjugate forms of weight $\lambda'$ are read from
algebraic forms $\psi$, $\eta$ (of weights $\lambda$, $\lambda'$) as in Lemma~\ref{pc:algpet}:
\[
f_\psi(u,g)=\langle v'^{+},u^{-1}\psi(g)\rangle_\pgsdu,\qquad f_\eta(u,g)=\langle u^{-1}\eta(g),v^{+}
\rangle_\pgsdu\qquad(u\in\auU(Y_m)_\infty,\ g\in G_f),
\]
with $v^+$, $v'^+$ the highest weight vectors of the parameters $l$, $-l$, normalised by
$\langle v'^+,v^+\rangle_\pgsdu=1$. (The reading of conjugate forms is defined through the pairing in
this way, and $J$ maps the evaluation vector at the frame ordered by $l$, which is $v^+$, to the one at
the reversed frame, which is $v'^+$.) For a finite adelic point $x$
and a vector $v$ of weight $\lambda'$, the \emph{delta form} $\delta_{x,v}$ is the algebraic form of
weight $\lambda'$ and level $J^\circ_\vartheta=J_{Y_m}(\pcm_\vartheta)$ (the principal congruence
subgroup of Lemma~\ref{pgr:jcirc}, not the level $J_\vartheta$ of $\auAm(\vartheta)$) supported on
$\Gamma xJ^\circ_\vartheta$ with $\delta_{x,v}(\gamma xk)=\gamma v$; it is well defined because
$J^\circ_\vartheta$ is free (Lemma~\ref{au:free}: it consists of elements $\equiv1$ modulo the
chosen places, and the rank is $\le2n_m$). We use the same name for the conjugate form $f_{\delta_{x,v}}$.

\begin{lemma}[Size of the hook number]\label{pgs:hooksize}
$h\le L^{3L}$, where $L=\sum_\tau|\mu'_\tau|\le40f\,n_mb\le x^{19}b$, and
$h\exp(x^{99}b^2)\exp(x^{200}b^2)\le\exp(x^{201}b^2)$.
\end{lemma}
\begin{proof}
$h=\prod_\tau H_{\mu'_\tau}\le\prod_\tau L_\tau^{3L_\tau}\le L^{3L}$ by Lemma~\ref{pk:hook}. The
bound for $L$ is Lemma~\ref{pgs:boxes} below (with $N\le2n_m$), together with $f\le x^{12}$ and
$n_m\le40x^6$ (Lemma~\ref{pz:sizes}(b)) and $1600\le x$. Then $\log L\le\log b+19\log x\le2\log b$,
so $\log h\le3L\log L\le6x^{19}b\log b\le x^{20}b^2$, and $x^{20}+x^{99}+x^{200}\le x^{201}$.
\end{proof}

\begin{proposition}[Proposition~\ref{pgr:delta}]\label{pgs:delta}
Proposition~\ref{pgr:delta} holds, with $C_{\mathrm{du}}=1$. More precisely:
\begin{enumerate}
\item[(a)] for $y\in\auAz$, $[K_m:J^\circ_\vartheta]\int_{[\auU(Y_m)]}\delta_{x,v}\cdot X^\vartheta_{0,m}y=
\langle v,\iota_\vartheta y(x)\rangle_\pgsdu/\dim W_{\lambda}$;
\item[(b)] $\delta_{x,v}$ is fixed by the principal congruence subgroup modulo $\pcm(\vartheta)$ (the
level of the outer node $Y^0_m$); if $x$ is integral above $\ell(\iota)$ and $v$ is integral at
$\iota$, then $\delta_{x,v}$ has gauge at most $1$ at $\iota$ at every point; and
$\|\sigma(\delta_{x,v}(g))\|^2\le\|\sigma(v)\|^2$ for all $g$ and $\sigma\in\auAut(\C)$;
\item[(c)] if $u$ is a vector of weight $\lambda$ with $\langle w,u\rangle_\pgsdu\in h\auO_\iota$ for
all integral $w$, then $u$ is integral at $\iota$; if $\sigma\in\auAut(\C)$, $K\ge0$ and
$|\sigma\langle v,u\rangle_\pgsdu|\le K\|\sigma v\|$ for all $v$, then $\|\sigma u\|^2\le K^2$; and
$h\exp(x^{99}b^2)\le\exp(x^{201}b^2-x^{200}b^2)$.
\end{enumerate}
\end{proposition}
\begin{proof}
(a) Let $\psi=\iota_\vartheta y$, the algebraic form of weight $\lambda$ whose form is
$X^\vartheta_{0,m}y$ \eqref{pc:eq-iota}; it is right invariant under $J^\circ_\vartheta$
(Lemma~\ref{pgr:jcirc}(b)). Let $\eta=\delta_{x,v}$. As in the proof of
Lemma~\ref{pc:algpet}, $[\auU(Y_m)]$ is the disjoint union over the classes
$x'\in\Gamma\backslash G_f/J^\circ_\vartheta$ of the images of $\auU(Y_m)_\infty\times x'J^\circ_\vartheta$, of
volume $[K_m:J^\circ_\vartheta]^{-1}$ (the group $J^\circ_\vartheta$ is free), and the integrand is right
$J^\circ_\vartheta$-invariant. On the class of $x'$ the integral of $f_\eta f_\psi$ is
$[K_m:J^\circ_\vartheta]^{-1}\int_{\auU(Y_m)_\infty}\langle u^{-1}\eta(x'),v^+\rangle_\pgsdu\langle v'^+,
u^{-1}\psi(x')\rangle_\pgsdu\,du$. By Schur orthogonality for the irreducible representation
$W_\lambda$ of the compact group $\auU(Y_m)_\infty$ of mass $1$ and its dual $W_{\lambda'}$, paired
invariantly by (P1),
\[
\int\langle u^{-1}a,v^+\rangle_\pgsdu\langle v'^+,u^{-1}b\rangle_\pgsdu\,du=
\frac{\langle a,b\rangle_\pgsdu\langle v'^+,v^+\rangle_\pgsdu}{\dim W_\lambda}=
\frac{\langle a,b\rangle_\pgsdu}{\dim W_\lambda}.
\]
Only the class of $x$ contributes, and there $\eta(x)=v$.

(b) The group $J^\circ_\vartheta$ is the principal congruence subgroup of $K_m$
modulo $\pcm_\vartheta$, which contains the principal congruence subgroup modulo
$\pcm(\vartheta)=\pcb\prod_{w\in Q}w$, since $\pi_M\mid\pcb$. At the point $\gamma xk$ the value is
$\gamma v$, and by (G1), (G2) of \S\ref{pl:ss-norm} its gauge at $(\iota,\gamma xk)$ is the gauge of
$v$ at $(\iota,x)$, that is the largest absolute value of a coordinate of $\rho(x_\ell)^{-1}\iota(v)$;
if $x$ is integral above $\ell$, $\rho(x_\ell)^{-1}$ has integral entries (G2), so this is $\le1$
for integral $v$. Finally $\sigma(\gamma v)=\rho(\gamma^\sigma)\sigma(v)$, where $\rho(\gamma^\sigma)$
is the action of $\gamma$ through the embeddings $\sigma\circ\tau\circ c$, again an element of the
compact group $\prod_\tau\auU(N)$, which preserves the invariant form; so $\|\sigma(\gamma v)\|=
\|\sigma v\|$. (Off the class of $x$ the value is $0$.)

(c) Let $C$ be as in (P2) and $c_T$ its rows, integral. Then $h\,u_T=\langle c_T,u\rangle_\pgsdu\in
h\auO_\iota$, so $u_T\in\auO_\iota$. For the second claim let $\zeta=\sigma(u)$ (coordinatewise),
$a=\zeta^\flat$ as in (P3) and $v=\sigma^{-1}(a)$. As $A$ has integer entries,
$\sigma\langle v,u\rangle_\pgsdu=\langle a,\zeta\rangle_\pgsdu=\langle\zeta^\flat,\zeta^\flat\rangle=
\|\zeta\|^2$, while $\|\sigma v\|=\|a\|=\|\zeta\|$. So $\|\zeta\|^2\le K\|\zeta\|$, whence
$\|\zeta\|^2\le K^2$. The last claim is Lemma~\ref{pgs:hooksize}.
\end{proof}

\begin{remark}
In Proposition~\ref{pgr:delta}(c) the constant $C_{\mathrm{du}}$ may therefore be taken equal to $1$,
and the hypothesis $hC_{\mathrm{du}}^{1/2}\exp(x^{99}b^2)\le\exp(x^{201}b^2-x^{200}b^2)$ used in the
proof of Theorem~\ref{pc:auxcore}(b) holds. The weight $\lambda'$ is the weight of the conjugate
forms on $Y_m$ of the last path, and the bound on its boxes is Lemma~\ref{pgs:boxes} applied to the outer node $Y^0_m$ of the first grid.
\end{remark}

\subsection{Short forms of deep level: proof of Proposition~\ref{pgr:short}}\label{pgs:ss-short}

\subsubsection*{The parameters of the grids}
At a node $(a,b)$ of either grid the type is given (\S\ref{pgr:ss-arch}, extending
\S\ref{pc:ss-triangle}) by its doubled Harish-Chandra parameter at each $\tau$, which consists of:
on the stem, the numbers $2\varsigma'_{\tau,p}$, $p<r$ (and, in the second grid, $p\le r$ for the
extra line), where $\varsigma'_{\tau,p}=k/2-1-p$ for $p<k-1$ and $\varsigma'_{\tau,p}=\Delta/2-
\kappa_{p+2-k}(\tau)$ otherwise; on the $l$-th block of the first kind ($l<a$, rank $t$ or $t+1$),
the numbers $2\kappa_{l+1}(\tau)+t-1-2b'$; on the $l$-th block of the second kind ($l<b$, rank $t$ or
$t-1$), the numbers $t-1-2(\kappa_{l+1}(\tau)+z)-2b'$; here $z\in\{0,1\}$ is the grid, $b'$ runs over
the rank of the block, and all entries are then raised by $1$ if $z=1$, and lowered by $1$ if $z=0$
and $a+b$ is odd; the weight of the forms moved at the node has the negative of this parameter in
the second grid and in the first grid (where conjugate forms are moved). The weight $\tilde\lambda$ of
Theorem~\ref{pc:auxcore} is the weight at the outer node $\pgrY^m_0=\pgrY$ (in the internal
indexing: no block of the first kind, $m$ blocks of rank $t-1$ of the second kind).

\begin{lemma}[Sizes of the weights on the grids]\label{pgs:boxes}
Let $\lambda=(\kappa,e)$ be the weight at a node of either grid (indices $\le m$), of rank $N$.
Then $N\le2n_m$; every entry of its doubled parameter has absolute value at most $20b$; $|e_\tau|\le
10b$; and its diagram at $\tau$ has at most $20Nb$ boxes, so at most $40n_mb$, and at most $40fn_mb
\le x^{19}b$ in total.
\end{lemma}
\begin{proof}
Recall the data used: $t=5r$ and $n_m=k+m(t+1)$ (proof of Lemma~\ref{pz:sizes}(b)); $r=k+2\varrho+1$
with $\varrho=m/2-1$ (Lemma~\ref{an:lem-sizes}(S7)), so $r+1=k+m$; and $\kappa_j(\tau)=\kappa(\tau)+(j+\frac12)\Delta$
with $\Delta=2r^8$, the parameters of the grids. The rank is at most $r+1+a(t+1)+b(t-1)\le r+1+2mt=k+m+2mt\le2(k+m(t+1))=2n_m$.
For $j\le2m$, $\kappa_j(\tau)=\kappa(\tau)+(j+\frac12)\Delta\le b+(2m+1)\cdot2r^8\le b+x^{42}\le2b$, by
$\kappa(\tau)\le b$ (Lemma~\ref{an:lem-sizes}(S5)), $m\le2\sqrt x$, $r\le3x^5$, $x\ge(\log
b)^{100}\ge128^{100}$ and $x^{42}\le b$ (S1). The indices occurring are $p+2-k\le r+2-k=m+1$ and
$l+1\le m$. Hence the stem entries satisfy $|2\varsigma'_{\tau,p}|\le\max(k,2\kappa_j(\tau))\le4b$
(as $k\le r\le b$, and, for the entries of the form $\Delta-2\kappa_j(\tau)$, $|\Delta-2\kappa_j(\tau)|\le2\kappa_j(\tau)$
because $\kappa_j(\tau)\ge\Delta/2$, as $\kappa(\tau)\ge1$ by (S5) and $j\ge0$), the entries of the blocks are at most $2\kappa_{l+1}(\tau)+2t+2\le4b+10r+2\le5b$
in absolute value, and the final shift adds at most $1$: all entries are at most $6b\le20b$.

Let $l_0>l_1>\dots>l_{N-1}$ be the entries, so $l_p=2\mu_p+N-1-2p$ with the highest weight
$\mu_0\ge\dots\ge\mu_{N-1}=e_\tau$. The row lengths of the diagram are
$\mu_p-\mu_{N-1}=\frac12(l_p-l_{N-1})-(N-1-p)\le\frac12(|l_p|+|l_{N-1}|)\le20b$, so the diagram has at
most $20Nb$ boxes. Also $2e_\tau+N-1\le2\mu_0+N-1=l_0\le20b$ and $2e_\tau-(N-1)=l_{N-1}\ge-20b$, so
$|e_\tau|\le10b$. Finally $40fn_mb\le40\cdot x^{12}\cdot40x^6\,b\le x^{19}b$ by
Lemma~\ref{pz:sizes}(b) and Lemma~\ref{an:lem-sizes}(S3).
\end{proof}

\subsubsection*{Short forms span at a free level}
The following lemma is Step~1 of the proof of Proposition~\ref{au:short}, for an arbitrary weight;
the only new point is the determinant part of the weight, which is absorbed by short bases of the
lattices of the powers $g^{e_\tau}$.

\begin{lemma}[Short forms span]\label{pgs:shortspan}
Let $S$ be a positive definite standard space of rank $n\ge1$ (so $\auU(S)(\auA_{E^+,f})=
\auU(E^n)(\auA_{E^+,f})$ via the frame), $\lambda=(\kappa,e)$ a weight, $|\kappa|=\sum_\tau|\mu_\tau|$
its total number of boxes, $J\subseteq K_S$ a free level, $c_0$ a positive integer in the different
of $E$, $C=6c_0(fn)^2$ the constant of Lemma~\ref{au:shortbasis}, and $B\ge C^{6(nf+|\kappa|)}$.
Then $\mathcal A_\lambda(J,1)$ is spanned by forms $\varphi$ that are integral at every
$j\colon\C\to\C_\ell$, for every prime $\ell$, and satisfy $\|\sigma(\varphi(g))\|^2\le B$ for all
$g\in G_f$ and $\sigma\in\auAut(\C)$.
\end{lemma}
\begin{proof}
As in Step~1 of the proof of Proposition~\ref{au:short}, since $J$ is free,
$\mathcal A_\lambda(J,1)$ is the direct sum over the classes $x\in\Gamma\backslash G_f/J$ of the
spaces $\{\delta_{x,\xi}:\xi\in W_\lambda\}$, where $\delta_{x,\xi}(\gamma xk)=\gamma\xi$; the form
$\delta_{x,\xi}$ is integral at $j$ if and only if $\xi$ has gauge at most $1$ at $(j,x)$ (by (G1),
(G2)), and $\|\sigma(\delta_{x,\xi}(g))\|=\|\sigma\xi\|$ or $0$ (proof of
Proposition~\ref{pgs:delta}(b)). So it suffices, for fixed $x$, to span $W_\lambda$ by vectors $\xi$
of gauge at most $1$ at $(j,x)$ for all $j$ and with $\|\sigma\xi\|^2\le B$ for all $\sigma$.

By Lemma~\ref{au:shortbasis} choose $Y_0\in\auGL_n(E)$ with $x^{-1}Y_0$ integral and, for each $\tau$,
$Y_\tau\in\auGL_n(E)$ with $(x^{e_\tau})^{-1}Y_\tau$ integral (Lemma~\ref{au:shortbasis} applied to the
element $x^{e_\tau}$ of $G_f$; $e_\tau$ may be negative), all entries of all of them having all
conjugates of absolute value at most $C$. For a family $\zeta=(\zeta_\tau)$ of $n\times n$ matrices
whose entries are roots of unity put
\[
\xi(\zeta)=\prod_\tau\tau'(\det Y_\tau)\cdot\rho_0(Y_0)\,\pgsev(\zeta),
\]
where $\rho_0$ is the action of weight $(\kappa,0)$ through the embeddings, $\tau'=\tau\circ c$ is the
embedding through which the factor at $\tau$ is acted on, and $\pgsev(\zeta)=\bigotimes_\tau
\pgsev(\zeta_\tau)$ as in \eqref{pgs:eq-lattice}.

\emph{They span.} $\rho_0(Y_0)$ and the scalar are invertible, and the $\pgsev(\zeta)$ span
$W_\lambda$: a linear form vanishing on them is $\zeta\mapsto\sum_Tc_Tb_T(\zeta)$, a polynomial
vanishing at all families with root of unity entries, hence $0$, and the $b_T$ are linearly
independent, so $c=0$.

\emph{Integrality.} The action of $x_\ell^{-1}$ in weight $\lambda$ is $\prod_\tau\tau'(\det
x^{-1})^{e_\tau}$ times $\rho_0(x^{-1})$, and $\rho_0(x^{-1})\rho_0(Y_0)=\rho_0(x^{-1}Y_0)$, $\det(
x^{-1})^{e_\tau}\det Y_\tau=\det((x^{e_\tau})^{-1}Y_\tau)$, all components at the primes above $\ell$
taken through $j$. So $\rho(x_\ell)^{-1}j(\xi(\zeta))$ is a product of determinants of integral
matrices times $\rho_0(x^{-1}Y_0)$, a matrix whose entries are polynomials with integer coefficients
in the entries of the integral matrix $x^{-1}Y_0$ (Schur matrices; \S\ref{pl:ss-norm} (N1)), applied
to $\pgsev(\zeta)$, whose coordinates $b_T(\zeta)$ are integer polynomials in roots of unity. All
coordinates are integral: the gauge is at most $1$.

\emph{Size.} Fix $\sigma$. The coordinates of $\sigma(\xi(\zeta))$ are those of
$\prod_\tau\det(Y'_\tau)\cdot\rho_0(Y'_0)\pgsev(\zeta')$, where $Y'_\tau=\sigma\tau'(Y_\tau)$, $Y'_0$ is
the family $\sigma\tau'(Y_0)$, and $\zeta'=\sigma(\zeta)$ has root of unity entries. By (N1),
$\rho_0(Y'_0)\pgsev(\zeta')=\bigotimes_\tau\pgsev(z_\tau)$ with $z_\tau=\tilde Y_\tau\zeta'_\tau$,
$\tilde Y_\tau$ being $\sigma\tau'(Y_0)$ or a transpose of it. Each entry of $z_\tau$ is a sum of $n$
products of a number of absolute value $\le C$ and a root of unity, so each column of $z_\tau$ has
squared norm at most $n(nC)^2\le C^5$ (as $n\le C$). By \eqref{pgs:eq-evnorm},
$\|\pgsev(z_\tau)\|^2\le C^{5|\mu_\tau|}$. By Hadamard's inequality $|\det Y'_\tau|^2\le(nC^2)^n\le
C^{3n}$. As the invariant form is the tensor product over $\tau$,
\[
\|\sigma\xi(\zeta)\|^2\le\prod_\tau C^{3n}\cdot\prod_\tau C^{5|\mu_\tau|}\le C^{3nf+5|\kappa|}\le B.
\qedhere
\]
\end{proof}

\begin{proposition}[Proposition~\ref{pgr:short}; Theorem~\ref{pc:auxcore}(c)]\label{pgs:short}
$\pcFF=\mathcal A_{\tilde\lambda}(J(\pcb),1)$ is spanned by forms that are integral at every
$j\colon\C\to\C_\ell$ for every prime $\ell$ (in particular at every $\iota\in\auI$) and whose values
have all conjugates of squared norm at most $\exp(x^{99}b^2)$.
\end{proposition}
\begin{proof}
Apply Lemma~\ref{pgs:shortspan} with $S=\pgrY$, $n=\tilde n=n_m-m\le n_m$, $\lambda=\tilde\lambda$,
$J=J(\pcb)$ (free by Lemma~\ref{au:free}, as recalled in \S\ref{pc:ss-aux}) and $c_0=q_1d$ (which
lies in the different of $E\subseteq\Q(\zeta_{q_1d})$, as in the proof of Proposition~\ref{au:short}).
By Lemma~\ref{pz:sizes}(b), $C=6q_1d(f\tilde n)^2\le x^{102}$; by Lemma~\ref{pgs:boxes},
$|\tilde\kappa|\le x^{19}b$; and $\tilde nf\le40x^{18}$. Hence
\[
\log C^{6(\tilde nf+|\tilde\kappa|)}\le6\cdot102\,(40x^{18}+x^{19}b)\log x\le x^{21}b\log b\le
x^{99}b^2,
\]
and $B=\exp(x^{99}b^2)$ is admissible.
\end{proof}

\begin{remark}
In the printed Theorem~\ref{pc:auxcore}, ``integral at $\iota$'' for $\varphi\in\pcFF$ means gauge at
most $1$ at $\iota$ at every point, which is the notion used here. For the number of boxes of
$\tilde\lambda$ at each place we use the explicit bound $x^{19}b$ (Lemma~\ref{pgs:boxes}).
\end{remark}



%% ===== proofs/C-grid3.tex =====
% proofs/C-grid3.tex
% Continuation of proofs/C-grid.tex (pgr:) and proofs/C-grid2.tex (pgs:): the turn on the
% diagonal, Proposition pgr:turndiag. Labels start with pgu:. Conventions and macros: those of
% proofs/C-grid.tex and proofs/C-determinant.tex.
\providecommand{\auI}{\mathcal I}
\providecommand{\auGL}{\mathrm{GL}}
\providecommand{\auU}{\mathrm U}
\providecommand{\auO}{\mathcal O}
\providecommand{\auA}{\mathbb A}
\providecommand{\auNm}{\operatorname{Nm}}
\providecommand{\auHom}{\operatorname{Hom}}
\providecommand{\aubar}[1]{\overline{#1}}
\providecommand{\pcm}{\mathfrak m}
\providecommand{\pcMat}{\mathrm M}
\providecommand{\pctr}{\operatorname{tr}}
\providecommand{\psHS}{\mathcal H}
\providecommand{\psrho}{\varrho}
\providecommand{\psvac}{\mathbf v}
\providecommand{\psZ}{Z}
\providecommand{\pgrY}{\widetilde Y}
\providecommand{\pgrG}{\Gamma}
\providecommand{\pgrE}{\operatorname{Ed}}
\providecommand{\pgrTr}{\mathrel{\vdash}}
\providecommand{\pcb}{\mathfrak b}
\providecommand{\pguK}{\mathsf K}
\providecommand{\pguL}{\mathsf L}
\providecommand{\pgupt}{\mathsf u}

\section{The turn on the diagonal}\label{pgu:sec}

This section continues \S\ref{pgr:sec} and uses its notation. It proves
Proposition~\ref{pgr:turndiag} from Proposition~\ref{pgr:square} (the square law on a square of
a grid, at the square $(j,j)$ of the first grid), in the form of Theorem~\ref{pgu:turndiag}
below. The argument has five ingredients: the theta kernel of an edge ending at a diagonal node of
the second grid splits as a product of a kernel of the first grid and a kernel of a line
(\S\ref{pgu:ss-sum}); the kernel of a line can be reversed, at the cost of a finite Fourier
transform (\S\ref{pgu:ss-line}); two seesaw identities move periods between the diagonal nodes and
the node $S$ common to the two paths (\S\ref{pgu:ss-seesaw}); a path datum is turned across the
square after a Fourier transform in the column of the line (\S\ref{pgu:ss-turned}); and these are
assembled in \S\ref{pgu:ss-proof}.

\subsection{The data at a diagonal turn}\label{pgu:ss-data}

Fix $0\le j<m$, an $m$-fold twist $\vartheta$, $\pcm=\pcm(\vartheta)$ and $R=R_\pcm$. We use the
following nodes.
\begin{itemize}
\item In the first grid: $V'=Y^j_j$ (removed by the turn), $W_1$ the node $(j,j+1)$ (kept),
$V=Y^{j+1}_{j+1}$ (added), and $S_1$ the node $(j+1,j)$ (not on either path).
\item In the second grid: $\widetilde V'=\pgrY^j_j=V'\oplus L$ (removed), $S$ the node at the place
$(j,j+1)$ of the half $h_1$ (kept), and $\widetilde V=\pgrY^{j+1}_{j+1}=V\oplus L'$ (added). Here
$L=E\tilde e_j$ and $L'=E\tilde e_{j+1}$ are lines with $\langle\tilde e_l,\tilde e_l\rangle=
\epsilon_l$, $\epsilon_{l}=-\epsilon_{l+1}$, and the frames of $\widetilde V'$, $\widetilde V$ are
those of $V'$, $V$ followed by $\tilde e_j$, $\tilde e_{j+1}$.
\end{itemize}
We write $\mathsf B,\mathsf A,\mathsf N_1,\mathsf N_2$ for the frames of $V',V,W_1,S$, $\bullet$ for
the index of the line, and $\mathsf N=\mathsf N_1\sqcup\mathsf N_2$.

\subsubsection*{Facts on the grid data}
The proof uses the following properties of the grid data. They are proved in \S\ref{pgv:sec} from
the complete definition of the grids given there: (G1) is Proposition~\ref{pgv:G1}, (G2) is
Proposition~\ref{pgv:G2} (with $\alpha_L=\alpha_{V'S_1}$, $\beta^{\mathrm c}=1$,
$\alpha^{\mathrm r}=\upsilon$, $\beta_{L'}=\beta_{S_1V}$), (G3) is Theorem~\ref{pgv:G3}, (G4) is
Proposition~\ref{pgv:G4}.
\begin{enumerate}
\item[(G1)] $S$ and $S_1$ have the same rank and signs, and the frames identify them as spaces
(\S\ref{pc:ss-data}); we identify them, so that $\mathsf N_2$ is also the frame of $S_1$, and
$\#\mathsf B+\#\mathsf A=\#\mathsf N$ (both sides equal $2(r+(2j+1)t)$). (The types at $S$ and $S_1$
differ; only spaces, groups, lattices and frames are identified, which is all that is used:
Proposition~\ref{pgx:G1}.)
\item[(G2)] (Splitting characters.) Write the kernel of an edge $(p,q)$ with point datum $X$ and
Fock vector $P$ as $\theta_{(p,q)}(g_p,g_q;P,\delta_X)=\theta^{p,q}_{\alpha_{pq},\beta_{pq}}(g_p,
g_q;P,\delta_X)$ (notation of \S\ref{pc:ss-data}). There are automorphic characters
$\alpha_L,\beta^{\mathrm c}$ (column line) and $\alpha^{\mathrm r},\beta_{L'}$ (row line) of
$\auU(1)$, unramified outside $\pcm$ and tamely ramified at $\pcm$, such that
\[
\alpha_{\widetilde V'S}(\det(h\oplus k))=\alpha_{V'S_1}(\det h)\,\alpha_L(k),\qquad
\beta_{\widetilde V'S}=\beta_{V'S_1}\,\beta^{\mathrm c},
\]
\[
\alpha_{S\widetilde V}=\alpha_{S_1V}\,\alpha^{\mathrm r},\qquad
\beta_{S\widetilde V}(\det(h\oplus k))=\beta_{S_1V}(\det h)\,\beta_{L'}(k),
\]
and $\alpha^{\mathrm r}=(\beta^{\mathrm c})^{-1}\upsilon$. The tame characters of the edges satisfy
$\chi_{(\widetilde V',S)}=\chi_{(V',S_1)}$ and $\chi_{(S,\widetilde V)}=\chi_{(S_1,V)}$; moreover
$\chi_{(W_1,V)}=\chi_{(V',S_1)}$ and $\chi_{(S_1,V)}=\chi_{(V',W_1)}$ (opposite edges of a square
have equal characters, as in \S\ref{pgr:ss-off}).
\item[(G3)] (Edges.) The pairs of the grids are edges in the sense of
Theorem~\ref{ptr:edge}\textup{(a)--(d)}, with joint harmonics $P_{pq}$, and the local statement
\textup{(H)} of \S\ref{ptr:ss-mult} holds for them. This is the input that
Proposition~\ref{pgr:square} also uses (step (iii) of \S\ref{pc:ss-exchange}).
\item[(G4)] (Diagonal period.) There is $\pgupt=\pgupt_\pcm\in\auU(1)(\auA_{E^+,f})$, depending only
on $\pcm$, such that the diagonal period on a path with diagonal index $l$ is
\[
\mathsf P_{l,g}(y'',\varphi'')=\int_{[\auU(Y^l_l)]}y''(x)\,\varphi''\bigl((x\oplus\pgupt^l)g\bigr)\,dx ,
\]
$g$ a real point of $\auU(\pgrY^l_l)$, $x\oplus\pgupt^l$ acting on $Y^l_l\oplus E\tilde e_l$. The
element $\pgupt$ is $\aubar u_0/u_0$, $u_0=\alpha\varpi_\pcm$ (Definition~\ref{pgv:def-period}); it is
exactly the element of Lemma~\ref{pgx:line}.
\end{enumerate}
All forms on the grids are rapidly decreasing on the automorphic quotients (Theorem~\ref{ptr:edge}(c)
for the types of (G3)).

The edges of the turn are $e_1=(V',W_1)$, $e_2=(\widetilde V',S)$ (old) and $f_1=(W_1,V)$,
$f_2=(S,\widetilde V)$ (new). In the first grid $e_1$ is vertical and $f_1$ horizontal; in the
second grid $e_2$ is horizontal and $f_2$ vertical. For an edge tuple $S$ (we write $\mathbf S$ for
edge tuples in this section, to avoid a clash with the node $S$) we cut the matrices of the edges
at the line:
\[
\mathbf S_{e_2}=(y_2\ |\ v),\quad y_2\in\pcMat_{\mathsf N_2\times\mathsf B},\ v\in
\pcMat_{\mathsf N_2\times\bullet};\qquad
\mathbf S_{f_2}=\binom{x_2}{v'},\quad x_2\in\pcMat_{\mathsf A\times\mathsf N_2},\ v'\in
\pcMat_{\bullet\times\mathsf N_2},
\]
and put $y(\mathbf S)=\binom{\mathbf S_{e_1}}{y_2}\in\pcMat_{\mathsf N\times\mathsf B}$,
$v(\mathbf S)=v$, $x(\mathbf S)=(\mathbf S_{f_1}\ x_2)\in\pcMat_{\mathsf A\times\mathsf N}$,
$v'(\mathbf S)=v'$. Conversely $\mathbf S[y,v]$ and $\mathbf S[x,v']$ denote $\mathbf S$ with the
old, respectively new, edges replaced by the blocks of $(y,v)$, respectively $(x,v')$.

\subsection{Orthogonal sums and the splitting of kernels}\label{pgu:ss-sum}

\begin{lemma}[Orthogonal sums]\label{pgu:sum}
Let $W=W_1\perp W_2$ be hermitian spaces with self-dual lattices $L_1,L_2$ (for $[\,,]$), and
identify $\psHS(W)=\psHS(W_1)\otimes\psHS(W_2)$ (restricted tensor products; $\psrho(w_1+w_2)=
\psrho(w_1)\otimes\psrho(w_2)$). Then $\omega_W(g_1,g_2)=\omega_{W_1}(g_1)\otimes\omega_{W_2}(g_2)$
for $g_l\in\auU(W_l)(\auA)$, and $\theta_W=c\,\theta_{W_1}\otimes\theta_{W_2}$ with $c\ne0$.
\end{lemma}
\begin{proof}
Since $\beta_W(x_1+x_2,y_1+y_2)=\beta_{W_1}(x_1,y_1)\beta_{W_2}(x_2,y_2)$, the tensor product is a
model of $\psrho_W$. The functional $\theta_{W_1}\otimes\theta_{W_2}$ is invariant under
$\psrho(W(E))$, so it is $c^{-1}\theta_W$ by (W1). By (W1) for $W,W_1,W_2$ the operator
$\kappa(g)=(\omega_{W_1}(g_1)\otimes\omega_{W_2}(g_2))^{-1}\omega_W(g_1,g_2)$ commutes with all
$\psrho(w)$, hence is a scalar, and $\kappa$ is a continuous character of $\auU(W_1)(\auA)\times
\auU(W_2)(\auA)$. It suffices to show $\kappa(g_1,1)=1$ (and symmetrically). We follow the proof of
Proposition~\ref{pf:flip}, with $\kappa_0=\kappa(\cdot,1)$ in place of $\kappa_0$ there.
\emph{Rational points:} $\theta_W\circ\omega_W(\gamma)=\theta_W$ and the same for $W_1$ give
$\kappa_0(\gamma)=1$. \emph{Stabiliser of $L_1$:} $L=L_1\oplus L_2$ is self-dual, its vacuum is
$\psvac_{L_1}\otimes\psvac_{L_2}$ (by (W4), the line of the character is unique), and both
$\omega_W(k,1)$ and $\omega_{W_1}(k)\otimes1$ fix it for $k$ in the stabiliser of $L_1$ (last
assertion of (W4)); so $\kappa_0(k)=1$. \emph{Parabolic elements:} let $I\subset W_{1,v}$ be
isotropic and $p\in P_I(E^+_v)$ acting trivially on $I^\perp/I$ (in $W_1$). Then $I$ is isotropic in
$W_v$, $I^{\perp_W}/I=(I^{\perp_{W_1}}/I)\perp W_2$, and $(p,1)$ acts trivially on it; by the uniqueness
of the zero term (\S\ref{pf:ss-flip}), $\psZ^W_I=c'\,\psZ^{W_1}_I\otimes1$, and (W2$'$) for $W$ and $W_1$
gives $\chi_I(p)\psZ^W_I=\psZ^W_I\omega_W(p,1)=\kappa_0(p)\chi_I(p)\psZ^W_I$, the character $\chi_I(p)$
being the same for $W$ and $W_1$ (it depends on $\det p|_I$ only). So $\kappa_0(p)=1$. The remaining
steps (all finite places, by the Iwasawa decomposition; the infinite places, by density of the
rational points) are those of the proof of Proposition~\ref{pf:flip}, word for word.
\end{proof}

\begin{corollary}[Splitting at the diagonal nodes]\label{pgu:split}
Let $\pguK^{\mathrm c}_{v,B}(k,g)=\theta^{L,S}_{\alpha_L,\beta^{\mathrm c}}(k,g;B,\delta_v)$ for
$v\in\pcMat_{\mathsf N_2\times\bullet}$ and a polynomial $B$ on $\auHom(L,S)\otimes\R$, and
$\pguK^{\mathrm r}_{v',B'}(g,k)=\theta^{S,L'}_{\alpha^{\mathrm r},\beta_{L'}}(g,k;B',\delta_{v'})$ for
$v'\in\pcMat_{\bullet\times\mathsf N_2}$ and a polynomial $B'$ on $\auHom(S,L')\otimes\R$. Then
there are constants $c_1,c_2\ne0$, independent of $\vartheta$, such that for polynomials $A$ on
$\auHom(V',S)\otimes\R$, $A'$ on $\auHom(S,V)\otimes\R$, and all $h,k,g$,
\begin{align*}
\theta_{e_2}\bigl(h\oplus k,g;A\otimes B,\delta_{(y_2|v)}\bigr)&=c_1\,\theta_{(V',S_1)}(h,g;A,\delta_{y_2})\,
\pguK^{\mathrm c}_{v,B}(k,g),\\
\theta_{f_2}\bigl(g,h\oplus k;A'\otimes B',\delta_{\binom{x_2}{v'}}\bigr)&=c_2\,\theta_{(S_1,V)}(g,h;A',\delta_{x_2})\,
\pguK^{\mathrm r}_{v',B'}(g,k).
\end{align*}
\end{corollary}
\begin{proof}
$\auHom(\widetilde V',S)=\auHom(V',S)\perp\auHom(L,S)$, with the self-dual lattices
$\alpha\auHom(L_{V'},L_S)$ and $\alpha\auHom(L_L,L_S)$, and $(h\oplus k,g)$ acts as $(h,g)$ on the
first summand and $(k,g)$ on the second. By Lemma~\ref{pgu:sum}, $\omega(h\oplus k,g)=\omega(h,g)
\otimes\omega(k,g)$ and $\theta=c\,\theta\otimes\theta$. The datum factors:
$A\otimes B$ at the infinite places, the vacuum $\psvac^\pcm=\psvac\otimes\psvac$ outside
$S(\pcm)$, and at $v\in S(\pcm)$ the function $\delta^\natural_{(y_2|v)}(x_1|x_2)=
\delta^\natural_{y_2}(x_1)\delta^\natural_v(x_2)$ (the norms at $S(\pcm)$ are $L^2$ norms of product
measures, so the normalisations factor up to a constant independent of $\vartheta$). The characters
factor by (G2). This gives the first identity; the second is the same with
$\auHom(S,\widetilde V)=\auHom(S,V)\perp\auHom(S,L')$.
\end{proof}

\subsection{Reversing a line}\label{pgu:ss-line}

This subsection treats the case where $\alpha$ is a unit at the primes of $\pcm$. The general case is
Lemma~\ref{pgx:line} (\S\ref{pgx:ss-line}), which has the statement below with the data of
Definition~\ref{pgx:conv} and with $\pgupt$ exactly the element of Definition~\ref{pgv:def-period}; it is
the form used in the rest of the paper, and its proof follows the one given here. In this subsection we
assume that $\alpha$ is a unit at the primes of $\pcm$, so that at $v\in S(\pcm)$ the self-dual
lattices $\alpha\auHom(L_S,L_T)_v$ are the lattices of integral matrices. Let $n=\#\mathsf
N_2$ be the rank of $S$, and $L^-$ the line $L$ with the opposite form; it is the line $L'$ (same
frame vector, sign $-\epsilon_j=\epsilon_{j+1}$), and we identify them.

\begin{lemma}[Line reversal when $\alpha$ is a unit at the primes of $\pcm$; the general case is Lemma~\ref{pgx:line}]\label{pgu:line}
Assume that $\alpha$ is a unit at the primes of $\pcm$. There are an element $\pgupt=\pgupt_\pcm$ of $\auU(1)(\auA_{E^+,f})$ depending only on $\pcm$, a
linear bijection $\tau$ from the polynomials on $\auHom(L,S)\otimes\R$ to those on
$\auHom(S,L')\otimes\R$, independent of $\vartheta$, a primitive additive character $\psi_\pcm$ of
$R$, and $c_{\mathrm{line}}=c_{\mathrm{line}}(\vartheta)\in\C^\times$ with
\[
|c_{\mathrm{line}}|^2=C_0\,(\auNm\pcm)^{-n},\qquad C_0>0\text{ independent of }\vartheta,
\]
such that for every $k\in\auU(L)(\auA_f)$, $g\in\auU(S)(\auA)$, $v$ and $B$
\[
\pguK^{\mathrm c}_{v,B}(k,g)=c_{\mathrm{line}}\,\zeta(k)
\sum_{v'\in\pcMat_{\bullet\times\mathsf N_2}(R)}\psi_\pcm(v'v)\,\pguK^{\mathrm r}_{v',\tau B}(g,k\pgupt)
\]
where $\zeta(k)$ is a number of modulus one depending only on $k$, on the characters and on $\pcm$.
Here $\pgupt$ is the element constructed in the proof (it is $k_0$ outside $S(\pcm)$ and has
$w$-component $\varpi^{-k_w}$ at $v\in S(\pcm)$); it is in general not the element $\pgupt$ of
(G4), and the case of (G4) is treated by Lemma~\ref{pgx:line}.
\end{lemma}

\begin{proof}
\emph{The flip.} Let $\mathrm{Fl}\colon\psHS(\auHom(L,S))\to\psHS(\auHom(S,L^-))$ be the isomorphism
of Proposition~\ref{pf:flip} ((W6)): $\mathrm{Fl}\,\psrho(w)=\psrho(w^\dagger)\mathrm{Fl}$,
$\theta\circ\mathrm{Fl}=c_\theta\theta$, and $\mathrm{Fl}\,\omega(k,g)=\upsilon(\det k)^n\upsilon(\det
g)^{-1}\omega(g,k)\mathrm{Fl}$. With the twists of \S\ref{pc:ss-data},
\[
\mathrm{Fl}\,\omega^{L,S}_{\alpha_L,\beta^{\mathrm c}}(k,g)=\omega^{S,L^-}_{(\beta^{\mathrm c})^{-1}
\upsilon,\ \alpha_L^{-1}\upsilon^n}(g,k)\,\mathrm{Fl}
=\alpha_L(k)^{-1}\upsilon(k)^n\beta_{L'}(k)^{-1}\,
\omega^{S,L'}_{\alpha^{\mathrm r},\beta_{L'}}(g,k)\,\mathrm{Fl},
\]
using $\alpha^{\mathrm r}=(\beta^{\mathrm c})^{-1}\upsilon$ (G2), $\omega_{a,b}(g,h)=a(\det g)^{-1}b(\det h)\,\omega(g,h)$, and
$\beta^{\mathrm c}(g)\upsilon(g)^{-1}=\alpha^{\mathrm r}(g)^{-1}$. (Write
$\zeta_1(k)=\alpha_L(k)^{-1}\upsilon(k)^n\beta_{L'}(k)^{-1}$, as in Lemma~\ref{pgx:line}; it is a character of $\auU(1)$; at the fixed points $k$ used below it is a number of
modulus one; this is the factor $\zeta(k)$ of the statement, together with the value at
$\pgupt$ in the last step.) Hence
\[
\pguK^{\mathrm c}_{v,B}(k,g)=c_\theta^{-1}\theta\bigl(\mathrm{Fl}\,\omega^{L,S}(k,g)[B,\delta_v]\bigr)
=\zeta_1(k)\,c_\theta^{-1}\,\theta\bigl(\omega^{S,L'}_{\alpha^{\mathrm r},\beta_{L'}}(g,k)\,
\mathrm{Fl}[B,\delta_v]\bigr),
\]
and it remains to compute $\mathrm{Fl}[B,\delta_v]$.

\emph{Local factors.} The Heisenberg representations of $\auHom(L,S)_v$ and $\auHom(S,L^-)_v$ are
irreducible, and $w\mapsto w^\dagger$ preserves $\beta$ locally; so there are unitary local
intertwiners $\mathrm{Fl}_v$, unique up to factors of modulus one, and by the uniqueness of
$\mathrm{Fl}$, $\mathrm{Fl}=c_0\bigotimes_v\mathrm{Fl}_v$, where $|c_0|$ does not depend on the
choice of the unitary $\mathrm{Fl}_v$, hence not on $\vartheta$.

\emph{Infinite places.} In the Fock model, $\psHS_\infty(W)$ consists of the functions on
$W\otimes\R$ holomorphic for the complex structure of $W$ (multiplication by $i$ on the coordinates
of positive sign, by $-i$ on those of negative sign) and square integrable for the Gaussian; the
vacuum is $1$, and $\psrho_\infty$ is defined from $\beta$ and the complex structure. The map
$\dagger$ is $\R$-linear, preserves $\beta$, and carries the complex structure of $\auHom(L,S)$ to
that of $\auHom(S,L^-)$: it is antilinear, and the signs of the coordinates are negated (the form
of $L^-$ is that of $L$ negated). Hence $F\mapsto F\circ\dagger^{-1}$ is a unitary intertwiner, and
we take it as $\mathrm{Fl}_\infty$. It maps polynomials of degree $d$ to polynomials of degree $d$;
put $\tau B=B\circ\dagger^{-1}$ (in coordinates, a transposition of the variables with signs). It
does not depend on $\vartheta$.

\emph{Finite places outside $S(\pcm)$.} $\mathrm{Fl}_v$ maps the vacuum of the self-dual lattice
$M_v=\alpha\auHom(L_L,L_S)_v$ to a vector on which $\psrho(l^\dagger)$, $l\in M_v$, acts by
$\epsilon_{M_v}(l)$; as in the proof of Proposition~\ref{pf:flip} (stabiliser of $L$), this is the
character of the self-dual lattice $M_v^\dagger=\aubar\alpha\auHom(L_S,L_L)_v$, so by (W4)
$\mathrm{Fl}_v\psvac_{M_v}=c_v\psvac_{M^\dagger_v}$ with $|c_v|=1$. Let $k_0\in\auU(L')(\auA_f)$ be
the scalar $\aubar\alpha/\alpha$ on the line (it has norm one). Then $M^\dagger_v=k_{0,v}
\,\alpha\auHom(L_S,L_{L'})_v$, and by (W4) $\psvac_{M^\dagger_v}=c'_v\,\omega(1,k_{0,v})
\psvac_{\alpha\auHom(L_S,L_{L'})_v}$ with $|c'_v|=1$ ($\omega$ is unitary). At the places where
$\alpha$ is a unit both constants are $1$ for a suitable choice of $\mathrm{Fl}_v$.

\emph{Places of $S(\pcm)$.} Let $v=w\aubar w\in S(\pcm)$ with $w\mid\pcm$, $\varpi$ a uniformiser at
$w$ and $k_w=\operatorname{ord}_w\pcm$. In the model (W5), $\psHS_v(\auHom(L,S))$ is the space of
Schwartz--Bruhat functions of the columns $x\in E_w^n$, translations by $\auHom(L_w,S_w)$ act by
translation and those by $\auHom(L_{\aubar w},S_{\aubar w})$ by multiplication with characters;
similarly $\psHS_v(\auHom(S,L'))$ consists of functions of the rows $\rho\in E_w^n$. The map
$\dagger$ is antilinear, so it exchanges the components at $w$ and $\aubar w$: it carries
translations of the first model into multiplications of the second and conversely. An operator with
these commutation relations is, up to a scalar (irreducibility), a Fourier transform: there is
$\kappa_w\in E_w^\times$ with
\[
\mathrm{Fl}_v\Psi(\rho)=\int_{E_w^n}\Psi(x)\,\psi_w(\kappa_w\rho x)\,dx ,
\]
$\psi_w$ the local component of $\psi$ and $dx$ the measure giving $\auO_w^n$ volume $1$, up to a
factor of modulus one. As $\mathrm{Fl}_v$ maps the vacuum $1_{\auO_w^n}$ of the integral lattice to
a multiple of the vacuum $1_{\auO_w^n}$ of the target (as in the previous paragraph, $\alpha$ being a
unit at $w$), and the transform of $1_{\auO_w^n}$ is the indicator of the dual lattice of
$\auO_w^n$ for $(\rho,x)\mapsto\psi_w(\kappa_w\rho x)$, this dual lattice is $\auO_w^n$; so $dx$ is
self-dual, the formula is unitary (Plancherel), and the factor has modulus one. Now let
$\tilde v\in\auO_w^n$ lift $v$. Then
\[
\mathrm{Fl}_v\delta_v^\natural(\rho)=\int_{\tilde v+\varpi^{k_w}\auO_w^n}\psi_w(\kappa_w\rho x)\,dx
=q_w^{-k_wn}\,\psi_w(\kappa_w\rho\tilde v)\,1_{\varpi^{-k_w}\auO_w^n}(\rho),
\]
$q_w$ the order of the residue field. Let $\pgupt_v\in\auU(L')(E^+_v)\cong E_w^\times$ have
$w$-component $\varpi^{k_w}$. By (W5) and unitarity, $(\omega(1,\pgupt_v)F)(\rho)=\eta\,
q_w^{k_wn/2}F(\varpi^{-k_w}\rho)$ with $|\eta|=1$ (the factor $|\det|^{-n/2}$ makes the operator
unitary). Hence
\[
\omega(1,\pgupt_v)\,\mathrm{Fl}_v\delta_v^\natural(\rho)=\eta\,q_w^{-k_wn/2}\,\psi_w(\kappa_w
\varpi^{-k_w}\rho\tilde v)\,1_{\auO_w^n}(\rho)=\eta\,q_w^{-k_wn/2}\sum_{v'}\psi_{w}^{\pcm}(v'v)\,
\delta^\natural_{v'}(\rho),
\]
where $\psi^\pcm_w(r)=\psi_w(\kappa_w\varpi^{-k_w}\tilde r)$ is a well defined character of
$\auO_w/\varpi^{k_w}$ ($\kappa_w$ is a unit for $\psi_w$ in the sense that $\auO_w$ is self-dual for
$\psi_w(\kappa_w\,\cdot\,)$), primitive. Put $\psi_\pcm=\prod_{w\mid\pcm}\psi^\pcm_w$ on
$R=\prod_w\auO_w/\varpi^{k_w}$.

\emph{Assembling.} Let $\pgupt$ be $k_0^{-1}$ outside $S(\pcm)$ (at the places where $\alpha$ is
not a unit) and $\pgupt_v$ at $v\in S(\pcm)$; it depends only on $\pcm$. The local results
give
\[
\mathrm{Fl}[B,\delta_v]=c\;\omega^{S,L'}(1,\pgupt^{-1})\Bigl[\tau B,\sum_{v'}\psi_\pcm(v'v)\delta_{v'}\Bigr],
\qquad|c|=|c_0|\,\prod_{w\mid\pcm}q_w^{-k_wn/2}=|c_0|\,(\auNm\pcm)^{-n/2},
\]
the operators $\omega(1,\cdot)$ being taken without the determinant twists (which only contribute
a factor of modulus one, depending on $\pcm$ and the characters). Inserting this in the formula of
the first step and using $\omega(g,k)\omega(1,\pgupt^{-1})=\omega(g,k\pgupt^{-1})$, we get the
identity with $k\pgupt^{-1}$ in place of $k\pgupt$; replacing $\pgupt$ by $\pgupt^{-1}$ (both depend
only on $\pcm$) gives the statement, with $c_{\mathrm{line}}=c/c_\theta$ and
$C_0=|c_0/c_\theta|^2$.
\end{proof}

\subsection{The seesaw}\label{pgu:ss-seesaw}

A \emph{moderate weight} from a space $S_0$ to a space $T'$ is a pair $(N,j)$ of a continuous
function $N$ on $\auU(S_0)(\auA)$, left invariant under $\auU(S_0)(E^+)$ and bounded by a power of
the height (Definition~\ref{pcx:def-ht}), and a continuous map $j\colon\auU(S_0)(\auA)\to
\auU(T')(\auA)$ compatible with the rational points and with $\mathrm{ht}(j(x))\le C\,\mathrm{ht}(x)^C$.
For forms $y$ on $S_0$, $F$ on $T'$ and a real point $r$ of $\auU(T')$ put
\[
\langle y,F\rangle_{N,j,r}=\int_{[\auU(S_0)]}N(x)\,y(x)\,F(j(x)r)\,dx ;
\]
it converges absolutely when $y$ is rapidly decreasing and $F$ is bounded.

\begin{lemma}[Seesaw]\label{pgu:seesaw}
Let $(S_0,T)$ and $(T',S')$ be pairs of spaces that are edges (\textup{(G3)}) with joint harmonics
$\Phi$ and $\Phi'$ and fixed characters, for a family of moduli. Let $X$ be a set, and for each
$x\in X$ let $W_x$ be a finite set and, for $i\in W_x$, let $u_{x,i}\in\C$, $y_{x,i}$ a form on
$S_0$ of the source type of $\Phi$, $f_{x,i}$ a form on $S'$ of the conjugate of the target type of
$\Phi'$, $(N_{x,i},j_{x,i})$ a moderate weight from $S_0$ to $T'$, and $\Psi'_{x,i}$ a finite
datum for $(T',S')$ modulo $\pcm_x$. Let $\Xi$ be a function on $X$. Call a polynomial $p$ on
$\auHom(T',S')\otimes\R$ \emph{matched} if there are a polynomial $Q$ on $\auHom(S_0,T)\otimes\R$,
finite data $\Psi_{x,i}$ for $(S_0,T)$ modulo $\pcm_x$ and moderate weights $(M_{x,i},j'_{x,i})$ from
$S'$ to $T$ such that
\begin{enumerate}
\item[(a)] for every real point $r'$ of $\auU(T)$ there is $\beta(r')$ with
$\sum_{i\in W_x}u_{x,i}\langle f_{x,i},\Theta(\Phi,\Psi_{x,i})y_{x,i}\rangle_{M_{x,i},j'_{x,i},r'}=
\beta(r')\,\Xi(x)$ for all $x$;
\item[(b)] for all $x,i$ and all $a\in\auU(S_0)(\auA)$, $b\in\auU(S')(\auA)$,
$N_{x,i}(a)\,\allowbreak \theta^{T',S'}(j_{x,i}(a),b;p,\Psi'_{x,i})=M_{x,i}(b)\,\allowbreak \theta^{S_0,T}(a,j'_{x,i}(b);Q,
\Psi_{x,i})$.
\end{enumerate}
If the matched polynomials span all polynomials, then for every real point $r$ of $\auU(T')$ there
is $\beta\in\C$ with
\[
\sum_{i\in W_x}u_{x,i}\,\bigl\langle y_{x,i},\Theta^{\mathsf t}(\Phi',\Psi'_{x,i})f_{x,i}
\bigr\rangle_{N_{x,i},j_{x,i},r}=\beta\,\Xi(x)\qquad\text{for all }x\in X .
\]
\end{lemma}

\begin{proof}
For a Fock vector $\mathsf v$ of $\auHom(T',S')\otimes\R$ with finite moments put
\[
\Lambda_{x,i}(\mathsf v)=\int_{[\auU(S_0)]}\int_{[\auU(S')]}N_{x,i}(a)\,y_{x,i}(a)\,
\theta^{T',S'}(j_{x,i}(a),b;\mathsf v,\Psi'_{x,i})\,f_{x,i}(b)\,db\,da ,\qquad
\Lambda_x=\sum_{i\in W_x}u_{x,i}\Lambda_{x,i}.
\]
\emph{Convergence and continuity.} By (R2) of \S\ref{pcx:ss-conv} there are $M,A,Y$ with
$|\theta^{T',S'}(a',b;\mathsf v,\Psi')|\le Y\|\mathsf v\|_M(\mathrm{ht}(a')\mathrm{ht}(b))^A$, where
$\|\cdot\|_M$ is the moment norm of Proposition~\ref{pcx:decay}. The forms $y_{x,i}$ and $f_{x,i}$
are rapidly decreasing and $N_{x,i}$, $j_{x,i}$ are moderate, so the integrand is bounded by
$\|\mathsf v\|_M$ times an integrable function ((R1) of \S\ref{pcx:ss-conv}), and $\Lambda_{x,i}$ is a
continuous linear functional for $\|\cdot\|_M$; the integrals may be taken in any order.

\emph{Matched polynomials.} Let $p$ be matched, with $Q,\Psi,M$. By Lemma~\ref{ptr:conv-lift}
(this uses (H), included in (G3)) there is $\alpha_Q\in C_c(\auU(T)_\infty)$, independent of $x,i$,
with $\Theta(Q,\Psi_{x,i})y_{x,i}(h)=\int\alpha_Q(r')\,\Theta(\Phi,\Psi_{x,i})y_{x,i}(hr')\,dr'$. By
Fubini, (b), and the definition of the lift,
\[
\Lambda_{x,i}(p)=\int_{[\auU(S')]}M_{x,i}(b)f_{x,i}(b)\,\Theta(Q,\Psi_{x,i})y_{x,i}(j'_{x,i}(b))\,db=
\int\alpha_Q(r')\,\langle f_{x,i},\Theta(\Phi,\Psi_{x,i})y_{x,i}\rangle_{M_{x,i},j'_{x,i},r'}\,dr',
\]
the exchange of integrals being justified by the boundedness of the lifted forms and the rapid
decrease of $f_{x,i}$. Multiplying by $u_{x,i}$, summing and using (a),
$\Lambda_x(p)=\bigl(\int\alpha_Q\beta\bigr)\,\Xi(x)$ for all $x$. So the vector $(\Lambda_x(p))_{x\in
X}\in\C^X$ lies on the line $\C\,\Xi$.

\emph{All Fock vectors.} The set of $\mathsf v$ with $(\Lambda_x(\mathsf v))_x\in\C\,\Xi$ is a
linear subspace, closed for $\|\cdot\|_M$ (each $\Lambda_x$ is continuous, and a line is closed in
$\C^X$ with the product topology). It contains the matched polynomials, hence all polynomials, hence
(Lemma~\ref{ptr:dense}) every Fock vector with finite moments of all orders.

\emph{Descent.} Let $r$ be a real point of $\auU(T')$ and $\mathsf v_r=\omega_\infty(r,1)\Phi'$
(without determinant twists), which has finite moments of all orders. Translating the transposed
lift by $r$ moves the Fock vector: $\Theta^{\mathsf t}(\Phi',\Psi')f(a'r)=\int\theta^{T',S'}(a'r,b;
\Phi',\Psi')f(b)\,db=\zeta_r\int\theta^{T',S'}(a',b;\mathsf v_r,\Psi')f(b)\,db$, with $\zeta_r$ the
value at $r$ of the archimedean component of the determinant twist (of modulus one, independent of
$x,i$, since all characters of the family have the same weights). Hence
$\langle y_{x,i},\Theta^{\mathsf t}(\Phi',\Psi'_{x,i})f_{x,i}\rangle_{N_{x,i},j_{x,i},r}=\zeta_r
\Lambda_{x,i}(\mathsf v_r)$, and the claim follows from the previous step applied to $\mathsf v_r$.
\end{proof}

\subsection{Turning a path datum with a line}\label{pgu:ss-turned}

We keep the notation of \S\ref{pgu:ss-data}. Let $w$ be a set of edges none of which has $V'$ or
$\widetilde V'$ as an end or $V$ or $\widetilde V$ as its target, let $\chi=(\chi_{e_1}
\chi_{e_2})^{-1}$, and fix a bijection $\varepsilon\colon\mathsf B\sqcup\mathsf A\to\mathsf N$ and an
additive character $\psi$ of $R$. Matched functions are those of \S\ref{pc:ss-exchange} for $\chi$
and $\varepsilon$ (as in \S\ref{pgr:ss-off}). For a path datum $\Psi'$ on $w\cup\{e_1,e_2\}$ and an
edge tuple $\mathbf S$ put
\[
\widehat\Psi'_{\mathbf S}(v',y)=\sum_{v\in\pcMat_{\mathsf N_2\times\bullet}(R)}\psi(v'v)\,
\Psi'(\mathbf S[y,v])\qquad(v'\in\pcMat_{\bullet\times\mathsf N_2},\ y\in\pcMat_{\mathsf N\times
\mathsf B}),
\]
the Fourier transform in the column of the line; it depends on $\mathbf S$ only through
$(\mathbf S_e)_{e\in w}$.

\begin{proposition}[Turning with a line]\label{pgu:turned}
Let $\Psi'$ be a path datum on $w\cup\{e_1,e_2\}$ with bound $N$. Then
\[
\Psi(\mathbf S)=\bigl(\mathrm{match}_\chi\widehat\Psi'_{\mathbf S}(v'(\mathbf S),\cdot\,)\bigr)(x(\mathbf S))
\]
defines a path datum on $w\cup\{f_1,f_2\}$ with bound $\#\pcMat_{\mathsf N_2\times\bullet}(R)\cdot N
\le V(\vartheta)N$, and for every $\mathbf S$ and $v'$ the functions $x\mapsto\Psi(\mathbf S[x,v'])$
and $\widehat\Psi'_{\mathbf S}(v',\cdot)$ are matched.
\end{proposition}

\begin{proof}
\emph{Transformation law of the transform.} Let $g$ be an admissible node family and
$k\in\auGL_{\mathsf B}(R)$. Let $g'$ be $g$ with $k$ at $V'$ and $k\oplus1$ at $\widetilde V'$; it
is admissible (these nodes are neither outer nor touched by the other condition except through
$k\oplus1$). Put $G=\mathrm{diag}(g_{W_1},g_S)\in\auGL_{\mathsf N}(R)$. Then
$g'(\mathbf S[y,v])=(g\mathbf S)[G^{-1}yk,\ g_S^{-1}v]$, the edges of $w$ not meeting $V',\widetilde
V'$. The invariance of $\Psi'$ under $g'$ and the substitution $v\mapsto g_Sv$ in the sum give
\begin{equation}\label{pgu:eq-law}
\widehat\Psi'_{\mathbf S}(v',y)=\mathsf a(g)\,\chi(\det k)\,\widehat\Psi'_{g\mathbf S}(v'g_S,\,
G^{-1}yk),\qquad\mathsf a(g)=\chi_{e_1}(\det g_{W_1})\chi_{e_2}(\det g_S)\chi_w(g),
\end{equation}
since $\psi(v'g_S\cdot g_S^{-1}v)=\psi(v'v)$ and $\chi_{w\cup\{e_1,e_2\}}(g')=\chi(\det k)\mathsf
a(g)$ ($\det(k\oplus1)=\det k$).

\emph{The match exists.} If $\widehat\Psi'_{\mathbf S}(v',y)\ne0$ then $\Psi'(\mathbf S[y,v])\ne0$
for some $v$, so $\mathbf S_{e_1}$ (the rows $\mathsf N_1$ of $y$; $e_1$ is vertical) has a left
inverse, hence so has $y$, and $y=y_\gamma$ by Lemma~\ref{pgr:extend}; and \eqref{pgu:eq-law} with
$g=1$ gives $\chi(\det k)\widehat\Psi'_{\mathbf S}(v',yk)=\widehat\Psi'_{\mathbf S}(v',y)$. By
Lemma~\ref{pgr:match}(b) the match exists and is matched with $\widehat\Psi'_{\mathbf S}(v',\cdot)$.
Since $v'(\mathbf S[x,v'])=v'$, $x(\mathbf S[x,v'])=x$ and the transform depends only on
$(\mathbf S_e)_{e\in w}$, the last assertion follows, and $\Psi$ depends only on the edges of
$w\cup\{f_1,f_2\}$.

\emph{Invariance.} For admissible $g$: $v'(g\mathbf S)=v'(\mathbf S)g_S$, because $g$ is
$g_V\oplus1$ at $\widetilde V$, and $x(g\mathbf S)=g_V^{-1}x(\mathbf S)G$. Let
$H=\widehat\Psi'_{g\mathbf S}(v'(\mathbf S)g_S,\cdot)$ and $\Phi=\mathrm{match}_\chi H$. By
\eqref{pgu:eq-law} with $k=1$, $\widehat\Psi'_{\mathbf S}(v'(\mathbf S),y)=\mathsf a(g)H(G^{-1}y)$, so by
Lemma~\ref{pgr:match}(c) and uniqueness, $\Psi(\mathbf S)=\mathsf a(g)\chi(\det G)\Phi(x(\mathbf S)G)$.
On the other hand $\Psi(g\mathbf S)=\Phi(g_V^{-1}x(\mathbf S)G)=\chi(\det g_V)\Phi(x(\mathbf S)G)$
(last claim of Lemma~\ref{pgr:match}(c)). Hence $\Psi(\mathbf S)=\mathsf a(g)\chi(\det G)\chi(\det
g_V)^{-1}\Psi(g\mathbf S)$, and exactly as at the end of the invariance step of the proof of
Proposition~\ref{pgr:turned}, using $\chi_{f_1}=\chi_{e_2}$ and $\chi_{f_2}=\chi_{e_1}$ (G2) and
$\det(g_V\oplus1)=\det g_V$, the factor is $\chi_{w\cup\{f_1,f_2\}}(g)$.

\emph{Left inverses.} Let $\Psi(\mathbf S)\ne0$. Then $x(\mathbf S)=x_\gamma$ with
$\widehat\Psi'_{\mathbf S}(v'(\mathbf S),y_\gamma)\ne0$, so $\Psi'(\mathbf S[y_\gamma,v])\ne0$ for some $v$:
the vertical edges of $w$ carry split injective matrices, and so does $e_2$ if it is vertical
(it is not: $e_2$ is horizontal). The new vertical edge is $f_2$; its opposite old edge $e_1$ is
vertical, so the rows $\mathsf N_1$ of $y_\gamma$ form a split injective matrix and, by
Lemma~\ref{pgr:compl}, so do the columns $\mathsf N_2$ of $x_\gamma$, i.e.\ $x_2$. If $Lx_2=1$, then
$(L\ 0)\binom{x_2}{v'}=1$, so $\mathbf S_{f_2}$ is split injective.

\emph{Values.} $\Psi(\mathbf S)=\chi(\det\gamma)\sum_v\psi(v'v)\Psi'(\mathbf S[y_\gamma,v])$, a root of
unity times a sum of $\#\pcMat_{\mathsf N_2\times\bullet}(R)$ terms, each a root of unity times a
sum of at most $N$ roots of unity. Finally $\#\pcMat_{\mathsf N_2\times\bullet}(R)=(\#R)^n\le
(\#R)^{2n_m}=V(\vartheta)$.
\end{proof}

\begin{lemma}[Sums over the turned square with a line]\label{pgu:sumturn}
In the situation of Proposition~\ref{pgu:turned}, let $J$ be a set, $C\in\C$, and for $i\in J$ let
$L_i$ on $\pcMat_{\mathsf A\times\mathsf N}$ and $L'_i$ on $\pcMat_{\mathsf N\times\mathsf B}$
satisfy the square law for $\chi$ with constant $C$. Let $i(\mathbf S,v')\in J$ depend only on
$(\mathbf S_e)_{e\in w}$ and $v'$. Then
\[
\sum_{\mathbf S\in\mathcal S(w\cup\{e_1,e_2\})}\Psi'(\mathbf S)\sum_{v'}\psi(v'v(\mathbf S))\,
L'_{i(\mathbf S,v')}(y(\mathbf S))=C\sum_{\mathbf S\in\mathcal S(w\cup\{f_1,f_2\})}\Psi(\mathbf S)\,
L_{i(\mathbf S,v'(\mathbf S))}(x(\mathbf S)).
\]
\end{lemma}

\begin{proof}
$(\mathbf S,y,v)\mapsto\mathbf S[y,v]$ is a bijection from $\mathcal S(w)\times\pcMat_{\mathsf
N\times\mathsf B}\times\pcMat_{\mathsf N_2\times\bullet}$ onto $\mathcal S(w\cup\{e_1,e_2\})$, and
$(\mathbf S,x,v')\mapsto\mathbf S[x,v']$ one onto $\mathcal S(w\cup\{f_1,f_2\})$. For fixed
$\mathbf S\in\mathcal S(w)$ the left side is $\sum_{v'}\sum_y\widehat\Psi'_{\mathbf S}(v',y)
L'_{i(\mathbf S,v')}(y)$, which by the square law applied to the matched pair of
Proposition~\ref{pgu:turned} is $C\sum_{v'}\sum_x\Psi(\mathbf S[x,v'])L_{i(\mathbf S,v')}(x)$; sum
over $\mathbf S$.
\end{proof}

\subsection{Proof of Proposition~\ref{pgr:turndiag}}\label{pgu:ss-proof}

Let $\pgrG'\pgrTr_{(j,j)}\pgrG$ be a turn of type (diag): $\pgrG'=(j,h'_0,h'_1)$ where $h'_z=q_z\cdot
\bigl((j,j+1)\to(j,j)\bigr)$, and $\pgrG=(j+1,h_0,h_1)$ with $h_z=q_z\cdot\bigl((j,j+1)\to(j+1,j+1)
\bigr)$, $q_z$ a grid path from $(0,m)$ to $(j,j+1)$. Let $w$ be the set of edges of $q_0$ (first
grid) and $q_1$ (second grid). The nodes of $q_z$ lie in $\{(a,b):a\le j,\ b\ge j+1\}$, so no edge of
$w$ meets $V',\widetilde V'$ or ends at $V,\widetilde V$, and $\pgrE(\pgrG')=w\cup\{e_1,e_2\}$,
$\pgrE(\pgrG)=w\cup\{f_1,f_2\}$. For a period triple $x=(\vartheta,y',\varphi)$ and an edge tuple
$\mathbf S$ let $f_{\mathbf S}=q_{0*}(\mathbf S)y'$ (a conjugate form at $W_1$) and
$\varphi_{\mathbf S}=q_{1*}(\mathbf S)\varphi$ (a form at $S$); they depend on $\mathbf S$ only through
$w$. We write $T_M$, $\Lambda_M$ for the transposed lift and the lift with the point datum $\delta_M$
and the joint harmonic of the edge concerned.

\subsubsection*{Column and row periods}
For a real point $z$ of $\auU(S)$, a polynomial $B$ on $\auHom(L,S)\otimes\R$, a polynomial $B'$ on
$\auHom(S,L')\otimes\R$ and families $\Psi'$, $\Psi$ of path data on $\pgrE(\pgrG')$,
$\pgrE(\pgrG)$, put
\begin{align*}
I^{\mathrm c}(\Psi';z,B)(x)&=\sum_{\mathbf S\in\mathcal S_\vartheta(\pgrG')}\Psi'_\vartheta(\mathbf S)
\int_{[\auU(S)]}\pguK^{\mathrm c}_{v(\mathbf S),B}(\pgupt^j,b)\,\varphi_{\mathbf S}(b)\,
\bigl(\Lambda_{y_2}T_{\mathbf S_{e_1}}f_{\mathbf S}\bigr)(bz)\,db,\\
I^{\mathrm r}(\Psi;z,B')(x)&=\sum_{\mathbf S\in\mathcal S_\vartheta(\pgrG)}\Psi_\vartheta(\mathbf S)
\int_{[\auU(S)]}\pguK^{\mathrm r}_{v'(\mathbf S),B'}(b,\pgupt^{j+1})\,\varphi_{\mathbf S}(b)\,
\bigl(T_{x_2}\Lambda_{\mathbf S_{f_1}}f_{\mathbf S}\bigr)(bz)\,db,
\end{align*}
where $\Lambda_{y_2}$ is the lift along $(V',S_1)$, $T_{x_2}$ the transposed lift along $(S_1,V)$,
and the forms on $S_1$ are read on $S$ by (G1). The integrals converge absolutely (rapid decrease of
$\varphi_{\mathbf S}$, boundedness of the other form, (R2) of \S\ref{pcx:ss-conv} for the line
kernel).

\begin{lemma}[First seesaw]\label{pgu:see1}
Let $\Xi$ be a function of period triples. If for every $z$ and $B$ there is $\beta(z,B)$ with
$I^{\mathrm c}(\Psi';z,B)=\beta(z,B)\,\Xi$, then $\Psi'\propto_{\pgrG'}\Xi$.
\end{lemma}
\begin{proof}
Apply Lemma~\ref{pgu:seesaw} with $(S_0,T)=(V',S_1)$, $(T',S')=(\widetilde V',S)$, $W_x=\mathcal
S_\vartheta(\pgrG')$, $u_{x,\mathbf S}=\Psi'_\vartheta(\mathbf S)$, $y_{x,\mathbf S}=T_{\mathbf
S_{e_1}}f_{\mathbf S}$, $f_{x,\mathbf S}=\varphi_{\mathbf S}$, $N=1$, $j(a)=a\oplus\pgupt^j$,
$\Psi'_{x,\mathbf S}=\delta_{\mathbf S_{e_2}}$, $\Phi'$ the harmonic of $e_2$. The conclusion of the
lemma at a real point $r$ of $\auU(\widetilde V')$ is $I_{\pgrG'}(\Psi';r)=\beta\,\Xi$ by (G4) and the
definition of $T_{\mathbf S_{e_2}}$ (the half $h'_1$ ends with the transposed lift along $e_2$). The
polynomials on $\auHom(\widetilde V',S)\otimes\R=\auHom(V',S)\otimes\R\oplus\auHom(L,S)\otimes\R$
are spanned by the products $A\otimes B$. Each product is matched: take $Q=A$, $\Psi_{x,\mathbf
S}=\delta_{y_2(\mathbf S)}$, $M_{x,\mathbf S}(b)=c_1\pguK^{\mathrm c}_{v(\mathbf S),B}(\pgupt^j,b)$ and
$j'$ the identification of (G1); then (b) is Corollary~\ref{pgu:split}, and the sum in (a) at
$r'=z$ is $c_1I^{\mathrm c}(\Psi';z,B)$, a multiple of $\Xi$ by hypothesis.
\end{proof}

\begin{lemma}[Second seesaw]\label{pgu:see2}
If $\Psi\propto_{\pgrG}\Xi$, then for every $z$ and $B'$ there is $\beta$ with $I^{\mathrm
r}(\Psi;z,B')=\beta\,\Xi$.
\end{lemma}
\begin{proof}
Apply Lemma~\ref{pgu:seesaw} with $(S_0,T)=(S,\widetilde V)$, $(T',S')=(S_1,V)$, $W_x=\mathcal
S_\vartheta(\pgrG)$, $u_{x,\mathbf S}=\Psi_\vartheta(\mathbf S)$, $y_{x,\mathbf S}=\varphi_{\mathbf S}$,
$f_{x,\mathbf S}=\Lambda_{\mathbf S_{f_1}}f_{\mathbf S}$, $N_{x,\mathbf S}(b)=c_2\pguK^{\mathrm
r}_{v'(\mathbf S),B'}(b,\pgupt^{j+1})$, $j$ the identification of (G1), $\Psi'_{x,\mathbf
S}=\delta_{x_2(\mathbf S)}$, $\Phi'$ the harmonic of $(S_1,V)$, and $r=z$; the conclusion is the
claim (up to the factor $c_2$). Every polynomial $A'$ on $\auHom(S_1,V)\otimes\R$ is matched: take
$Q=A'\otimes B'$, $\Psi_{x,\mathbf S}=\delta_{\mathbf S_{f_2}}$, $M=1$ and $j'(a)=a\oplus\pgupt^{j+1}$.
Then (b) is Corollary~\ref{pgu:split}, and the sum in (a) at a real point $r'$ of
$\auU(\widetilde V)$ is $I_{\pgrG}(\Psi;r')$ by (G4), a multiple of $\Xi$ by hypothesis.
\end{proof}

\begin{theorem}[Proposition~\ref{pgr:turndiag}]\label{pgu:turndiag}
For $0\le j<m$ put
$c_{jj}=c^0_{jj}/c'_{\mathrm{line}}$, with $c^0_{jj}$ of Proposition~\ref{pgr:square} ($z=0$,
$a=b=j$) and $c'_{\mathrm{line}}$ of Lemma~\ref{pgx:line}. Then $c_{jj}$ has modulus $\rho'_{jj}$ up
to norm powers, and every turn $\pgrG'\pgrTr_{(j,j)}\pgrG$ transfers periods with bound $V$ and
constant $c_{jj}$.
\end{theorem}

\begin{proof}
We use Proposition~\ref{pgr:square}, which is Theorem~\ref{pgv:square}, and the facts (G1)--(G4), which are
Propositions~\ref{pgv:G1}, \ref{pgv:G2}, \ref{pgv:G4} and Theorem~\ref{pgv:G3}.
\emph{Modulus.} $|c^0_{jj}|^2=C(\auNm Q_\vartheta)^e\rho'^2_{jj}$ (with $e$ the norm-power exponent of Definition~\ref{pc:modulus}) and $|c'_{\mathrm{line}}|^{-2}=
C_0^{-1}(\auNm\pcm)^n$ with $n=r+(2j+1)t$; as $\auNm\pcm=\auNm\pcb\,\auNm Q_\vartheta$ with $\pcb=(\pi_M)^b$,
$|c_{jj}|^2=C\,C_0^{-1}(\auNm\pcb)^{n}(\auNm Q_\vartheta)^{e+n}\rho'^2_{jj}$.

\emph{The datum.} Let $\Psi'$ be a family of path data on $\pgrE(\pgrG')$ with bounds $N$. For each
$\vartheta$ let $\Psi_\vartheta$ be the datum of Proposition~\ref{pgu:turned} (with $w$ as above,
$\psi=\psi_\pcm$ of Lemma~\ref{pgx:line}, and the bijection $\varepsilon$ of the square law of
Proposition~\ref{pgr:square}); it is a path datum on $\pgrE(\pgrG)$ with bound $V(\vartheta)N(\vartheta)$.

\emph{The identity.} We show $I^{\mathrm c}(\Psi';z,B)(x)=c'_{\mathrm{line}}(\vartheta)\,\kappa\,
c^0_{jj}(\vartheta)^{-1}\,I^{\mathrm r}(\Psi;z,\tau B)(x)$. By Lemma~\ref{pgx:line} at $k=\pgupt^j$,
\[
\begin{gathered}I^{\mathrm c}(\Psi';z,B)(x)=c'_{\mathrm{line}}\sum_{\mathbf S}\Psi'(\mathbf S)\sum_{v'}\psi_\pcm(v'v(\mathbf
S))\,l_{\mathbf S,v'}\bigl(\Lambda_{y_2}T_{\mathbf S_{e_1}}f_{\mathbf S}\bigr),\\
l_{\mathbf S,v'}(F)=\int_{[\auU(S)]}\pguK^{\mathrm r}_{v',\tau B}(b,\pgupt^{j+1})\varphi_{\mathbf S}(b)F(bz)\,db\end{gathered}
\]
(a finite sum of absolutely convergent integrals). Here $l_{\mathbf S,v'}$ is a linear functional on
the conjugate forms at $S_1$, depending on $\mathbf S$ only through $w$, and
$l(\Lambda_{y_2}T_{\mathbf S_{e_1}}f)=L'_{f,l}(y(\mathbf S))$ in the notation of \S\ref{pgr:ss-off}
for the square $V'=(j,j)$, $W_1=(j,j+1)$, $W_2=S_1=(j+1,j)$, $V=(j+1,j+1)$ of the first grid.
By Proposition~\ref{pgr:square} ($z=0$, $a=b=j$), the functions $L_{f,l},L'_{f,l}$ satisfy the square
law for $(\chi_{e_1}\chi_{(V',S_1)})^{-1}=\chi$ (G2) with constant $\kappa/c^0_{jj}(\vartheta)$. By
Lemma~\ref{pgu:sumturn} with $J$ the pairs $(f,l)$ and $i(\mathbf S,v')=(f_{\mathbf S},l_{\mathbf
S,v'})$,
\[
I^{\mathrm c}(\Psi';z,B)(x)=c'_{\mathrm{line}}\frac{\kappa}{c^0_{jj}}\sum_{\mathbf S\in\mathcal
S_\vartheta(\pgrG)}\Psi(\mathbf S)\,L_{f_{\mathbf S},l_{\mathbf S,v'(\mathbf S)}}(x(\mathbf S)),
\]
and $L_{f,l}(x(\mathbf S))=l(T_{x_2}\Lambda_{\mathbf S_{f_1}}f)$, so the right side is
$c'_{\mathrm{line}}\kappa\,c^0_{jj}{}^{-1}I^{\mathrm r}(\Psi;z,\tau B)(x)$.

\emph{Transfer.} Let $\Psi\propto_{\pgrG}\Xi$. By Lemma~\ref{pgu:see2}, $I^{\mathrm r}(\Psi;z,\tau B)
=\beta(z,B)\Xi$ for every $z,B$. By the identity, $I^{\mathrm c}(\Psi';z,B)=\kappa\beta(z,B)\,
\Xi/c_{jj}$, where $\Xi/c_{jj}$ is the function $x\mapsto\Xi(x)/c_{jj}(\vartheta)$ and $\kappa\beta$
does not depend on $x$. By Lemma~\ref{pgu:see1} applied to $\Xi/c_{jj}$, $\Psi'\propto_{\pgrG'}
\Xi/c_{jj}$.
\end{proof}



%% ===== proofs/C-grid4.tex =====
% proofs/C-grid4.tex
% Continuation of proofs/C-grid.tex (pgr:), C-grid2.tex (pgs:), C-grid3.tex (pgu:): the complete
% definition of the two grids, the facts (G1)--(G4) of \S pgu:ss-data, and Proposition pgr:square.
% Labels start with pgv:. Conventions and macros: those of proofs/C-grid.tex and C-determinant.tex.
\providecommand{\auI}{\mathcal I}
\providecommand{\auGL}{\mathrm{GL}}
\providecommand{\auU}{\mathrm U}
\providecommand{\auO}{\mathcal O}
\providecommand{\auA}{\mathbb A}
\providecommand{\auNm}{\operatorname{Nm}}
\providecommand{\aubar}[1]{\overline{#1}}
\providecommand{\pcm}{\mathfrak m}
\providecommand{\pcb}{\mathfrak b}
\providecommand{\pcMat}{\mathrm M}
\providecommand{\pgrY}{\widetilde Y}
\providecommand{\pcnot}[1]{\par\smallskip{\raggedright\noindent\textsc{Not proved here.} [#1]\par}\smallskip}
\providecommand{\pgupt}{\mathsf u}
\providecommand{\pgvsgn}{\operatorname{sgn}}
\providecommand{\pgvwt}{\operatorname{wt}}
\providecommand{\pzF}{\mathfrak F}
\providecommand{\auPl}{\Phi}

\section{The grids: definition, the facts (G1)--(G4), and the square law}\label{pgv:sec}

This section continues \S\S\ref{pgr:sec}, \ref{pgs:sec}, \ref{pgu:sec} and uses their notation.
\S\ref{pgr:ss-arch} describes the two grids briefly, and \S\ref{pgu:ss-data} isolates four
properties (G1)--(G4) of them on which the diagonal turn (Theorem~\ref{pgu:turndiag}) rests. Here
\S\ref{pgv:ss-def} gives the complete definition (spaces, frames, signs, parameters, characters of
every edge, the moves, the levels and the diagonal period), \S\ref{pgv:ss-G} proves (G1), (G2), (G4)
and discusses the unit property of the idele $\alpha$ at the primes of $\pcm$ assumed in
\S\ref{pgu:ss-line}, \S\ref{pgv:ss-G3} proves (G3) (every pair of the grids is an edge, and (H)
holds), and \S\ref{pgv:ss-square} deduces Proposition~\ref{pgr:square} from Theorem~\ref{ptr:square}.


Throughout, $r$, $k=r+1-m$, $t=5r$, $\Delta=2r^8$, $\kappa$, $\kappa_l=\kappa+(l+\frac12)\Delta$ and
the characters $\mu^0_l$ are those of the data (\S\ref{ss:au-spaces}); $r$ and $t$ are odd, $m$ is
even, so $k$ is even; $\kappa_l(\tau)$ is an integer ($\Delta$ is even). We write $\upsilon$ also for
the restriction of $\upsilon$ to $\auU(1)$, of weight $1$ and unramified at every finite place, and
$\pgvwt$ for weights of characters of $\auU(1)$, normalised as in \S\ref{ptr:ss-param}, so that
$\pgvwt\mu^0_l=-\kappa_l$ and $\pgvwt\vartheta=0$ for a twist. For an $m$-fold twist $\vartheta$,
$\pcm=\pcm(\vartheta)=\pcb\prod_{w\in Q_\vartheta}w$ with $\pcb=(\pi_M)^b$, and $R=\auO_E/\pcm$, as
in \S\ref{pgr:ss-arch}. Since the primes of $\pi_M$ are prime to their conjugates and to $Q_\vartheta$
(\S\ref{ss:au-spaces}: $Q\cap\aubar Q=\emptyset$, $\auNm Q$ prime to $M$), $\pcm$ is a modulus in the
sense of \S\ref{pc:ss-data}.

\subsection{Definition of the grids}\label{pgv:ss-def}

\subsubsection*{Nodes, frames, ranks and signs}
A \emph{node} is a triple $p=(z,a,b)$ with $z\in\{0,1\}$ (the grid) and $0\le a,b\le m$. Its
\emph{frame} $\mathsf F_p$ is the disjoint union of
\begin{itemize}
\item the \emph{stem}, with coordinates numbered $0\le\varrho<r+z$;
\item $a$ \emph{blocks of the first kind}, numbered $0\le l<a$, each with $t+z$ coordinates
(positions $0\le\beta<t+z$);
\item $b$ \emph{blocks of the second kind}, numbered $0\le l<b$, each with $t-z$ coordinates
(positions $0\le\beta<t-z$).
\end{itemize}
So the rank is
\begin{equation}\label{pgv:eq-rank}
N_p=N^z_{a,b}=r+z+a(t+z)+b(t-z).
\end{equation}
For a stem coordinate $\varrho$ put $a(\varrho)=\max(\varrho+1-k,0)$; as $r+1=k+m$, $a(\varrho)=0$ for
$\varrho\le k-1$ and $a(\varrho)=\varrho+1-k$ runs through $1,\dots,m-1+z$ for $k\le\varrho<r+z$. The
\emph{signs} at $p$ are
\begin{equation}\label{pgv:eq-sign}
\epsilon_p(\varrho)=(-1)^{a(\varrho)+\min(a(\varrho),b)}\ \ (\varrho\text{ in the stem}),\qquad
\epsilon_p=+1\text{ on the blocks}.
\end{equation}
Thus $\epsilon_p(\varrho)=+1$ if $a(\varrho)\le b$ and $\epsilon_p(\varrho)=(-1)^{a(\varrho)+b}$ if
$a(\varrho)>b$; the signs depend on $z$ and $b$ only (and $z$ only through the length of the stem).
For $z=0$, $a=0$ these are the signs $\epsilon^{(0,b)}$ of $Y_b$ (\eqref{ptr:eq-E} with $i=0$, the
stem being numbered there from $1$).

\begin{definition}[The spaces]\label{pgv:def-spaces}
The space at the node $p$ is the following space with frame $\mathsf F_p$ (spaces in the sense of
\S\ref{pc:ss-data}):
\begin{itemize}
\item $p=(0,0,b)$: the space $Y_b$ with its frame and lattice (\S\ref{ps:ss-Y});
\item $p=(1,0,0)$: $Y_0\oplus E\tilde e$ with $\langle\tilde e,\tilde e\rangle=1$, the frame of $Y_0$
followed by $\tilde e$ as the stem coordinate $\varrho=r$;
\item every other node: $E^{N_p}$ with the standard lattice and the diagonal hermitian form with the
entries $\epsilon_p$ in the frame.
\end{itemize}
The node $(0,a,b)$ is written $Y^a_b$ (\emph{first grid}) and the node $(1,a,b)$ is written
$\pgrY^b_a$ (\emph{second grid}): \emph{$\pgrY^b_a$ has $a$ blocks of rank $t+1$ and $b$ blocks of
rank $t-1$}. So $Y^0_b=Y_b$, $\pgrY:=\pgrY^m_0$ (no block of rank $t+1$, $m$ blocks of rank $t-1$,
rank $\tilde n=r+1+m(t-1)$) is the space of \S\ref{pc:ss-aux}, and a half $h_1$ of a two-grid path
runs through the nodes $(1,a,b)=\pgrY^b_a$, as in \S\ref{pgr:ss-arch}.
\end{definition}

\begin{remark}[Indexing]\label{pgv:rem-index}
In the second grid the upper index counts the blocks of rank $t-1$ and the lower index those of rank
$t+1$: $\pgrY^b_a$ has $a$ blocks of rank $t+1$ and $b$ of rank $t-1$. With the opposite convention
$\pgrY^m_0$ would have rank $r+1+m(t+1)$, not $\tilde n$, and the node $(j,j+1)$ of $h_1$ would not
have the rank of $Y^{j+1}_j$ (compare Proposition~\ref{pgv:G1}).
\end{remark}

\subsubsection*{Parameters and types}
For a node $p=(z,a,b)$ and $\tau\in\auPl$ let $s_p=\frac12$ if $z=0$ and $a+b$ is odd, $s_p=0$
otherwise, and let $\lambda_{p,\tau}$ be the vector on $\mathsf F_p$ with entries
\begin{equation}\label{pgv:eq-param}
\begin{aligned}
&\varsigma'_{\tau,\varrho}+\tfrac z2-s_p&&\text{(stem coordinate }\varrho),\\
&\kappa_{l+1}(\tau)+\tfrac{t+z-1}2-\beta-s_p&&\text{(position }\beta\text{ of the block }l\text{ of the first kind}),\\
&-\kappa_{l+1}(\tau)+\tfrac{t-z-1}2-\beta-s_p&&\text{(position }\beta\text{ of the block }l\text{ of the second kind}),
\end{aligned}
\end{equation}
where $\varsigma'_{\tau,\varrho}=\frac k2-1-\varrho$ for $\varrho\le k-2$ and
$\varsigma'_{\tau,\varrho}=-\kappa(\tau)-(\varrho+2-k)\Delta=\frac\Delta2-\kappa_{\varrho+2-k}(\tau)$
for $\varrho\ge k-1$. (This is the description of \S\ref{pgs:ss-short}, written without doubling.)
The \emph{type at $p$}, $e_p$, is the type of $\auU(\mathsf F_p)_\infty$ whose Harish-Chandra
parameter at $\tau$ is $-\lambda_{p,\tau}$, with the signs $\epsilon_p$. The forms \emph{moved} at
$p$ are the forms of type $e_p$, of any level: in the first grid these are the conjugates of forms
of the type with parameter $\lambda_p$ (the \emph{conjugate forms} of \S\ref{pgr:ss-arch}); in the
second grid the parameter of the space $\pgrY^b_a$ is $-\lambda_{(1,a,b)}$ and they are forms. For
$p=(0,0,n)$, $\lambda_p=\Lambda^{(0,n)}$ of \S\ref{ptr:ss-param}, so $e_{(0,0,n)}=\overline{e_{0,n}}$.

\subsubsection*{Edges and their characters}
An \emph{edge} is a pair $(p,q)$ of nodes of one grid $z$ with $q=(z,a+1,b)$ (a \emph{right} edge;
its \emph{index} is $l=a$) or $q=(z,a,b+1)$ (an \emph{up} edge; index $l=b$), where $p=(z,a,b)$. Its
\emph{kind} is $c=1$ (\emph{vertical}) if it is an up edge of the first grid or a right edge of the
second grid, and $c=0$ (\emph{horizontal}) otherwise; this agrees with \S\ref{pgr:ss-arch}. Note
$N_q-N_p=d_{pq}$ with $d_{pq}=t+z$ for right edges and $t-z$ for up edges.

\begin{definition}[Characters, kernels, moves]\label{pgv:def-edges}
Let $\vartheta$ be an $m$-fold twist and $(p,q)$ an edge of kind $c$ and index $l$. Put
$\nu_{pq}=\mu^0_{l+1}\vartheta^{c(l+1)}$ and
\[
\alpha_{pq}=\nu_{pq}\,\upsilon^{\lfloor(N_q+1-c)/2\rfloor},\qquad
\beta_{pq}=\nu_{pq}\,\upsilon^{\lfloor(N_p+1-c)/2\rfloor}
\]
(so the power of $\upsilon$ is $\lfloor N/2\rfloor$ for vertical and $\lceil N/2\rceil$ for horizontal
edges). The \emph{kernel} of the edge with Fock vector $P$ and finite datum $\Psi$ on
$\pcMat_{q\times p}(R)$ (rows indexed by $\mathsf F_q$, columns by $\mathsf F_p$) is
$\theta_{\alpha_{pq},\beta_{pq}}(g_p,g_q;P,\Psi)$ for $\auHom(p,q)$, in the notation of
\S\ref{pc:ss-data}, with the modulus $\pcm(\vartheta)$; recall $\omega_{\alpha,\beta}(g,h)=
\alpha(\det g)^{-1}\beta(\det h)\omega(g,h)$, $g$ on the source $p$, $h$ on the target $q$. The
\emph{edge character} $\chi_{pq}$ is the character $x\mapsto\prod_{w\mid\pcm}\psi_w(x\bmod w)$ of
$R^\times$, $\psi_w$ the tame part of $\nu_{pq}$ at $w$ (trivial at the $w\mid\pcm$ where $\nu_{pq}$
is unramified). Let $P_{pq}$ be the joint harmonic of the edge (Theorem~\ref{pgv:G3}(a)). For
$X\in\pcMat_{q\times p}(R)$ the \emph{move} along the edge is
\[
\Lambda_X=\Theta_{\alpha_{pq},\beta_{pq}}(P_{pq},\delta_X)\colon\{\text{forms of type }e_p\}\to\{\text{forms
of type }e_q\}\quad\text{if }(p,q)\text{ is a right edge},
\]
\[
T_X=\Theta^{\mathsf t}_{\alpha_{pq},\beta_{pq}}(P_{pq},\delta_X)\colon\{\text{forms of type }e_q\}\to
\{\text{forms of type }e_p\}\quad\text{if }(p,q)\text{ is an up edge}.
\]
That these maps are defined and take values as stated is Theorem~\ref{pgv:G3}(d).
\end{definition}

\noindent So $\pgvwt\alpha_{pq}=-\kappa_{l+1}+\lfloor\frac{N_q+1-c}2\rfloor$ and $\pgvwt\beta_{pq}=
-\kappa_{l+1}+\lfloor\frac{N_p+1-c}2\rfloor$. For the edge $(0,0,n)\to(0,0,n+1)$ (vertical, index $n$)
the characters are $\mu_{n+1}\upsilon_{r+(n+1)t}$, $\mu_{n+1}\upsilon_{r+nt}$, those of $a_{0,n}$ in
\S\ref{pc:ss-triangle}. The edge character depends only on the kind and the index, and $\chi_{pq}$ is
also the character of $R^\times$ through which $\alpha_{pq}$ and $\beta_{pq}$ act on $\det$ of
elements of $K$ (Lemma~\ref{pc:level}; $\upsilon$ is unramified).

\subsubsection*{Levels, admissible families and the diagonal period}
The \emph{level} of the node $p$ is the subgroup of $\auGL_{\mathsf F_p}(R)$ equal to $\{1\}$ if
$p=(0,0,b)$ and to the kernel of $\auGL_{\mathsf F_p}(R)\to\auGL_{\mathsf F_p}(\auO_E/\pcb)$
otherwise. For $0\le l\le m$ we identify $\pgrY^l_l$ with $Y^l_l\oplus E\tilde e_l$,
$\langle\tilde e_l,\tilde e_l\rangle=(-1)^l$, as spaces with frames: for $l=0$ this is
Definition~\ref{pgv:def-spaces}, and for $l\ge1$ we fix a bijection
$\mathsf F_{(0,l,l)}\sqcup\{\bullet\}\to\mathsf F_{(1,l,l)}$ preserving signs, $\bullet$ having the
sign $(-1)^l$ (it exists by Lemma~\ref{pgv:count}(c)); both spaces are then standard diagonal spaces
with the same signs, and the bijection is an isometry carrying frame to frame and lattice to lattice.
A node family $g=(g_p)$ is \emph{admissible} if $g_{(0,0,m)}=1$, $g_{(1,0,m)}\equiv1\bmod\pcb$, and
$g_{(1,a,a)}=g_{(0,a,a)}\oplus1$ for all $a$ under this identification. (This makes precise the
definition in \S\ref{pgr:ss-arch}: the level of $Y^0_m=Y_m$ reduces to $\{1\}$ modulo $\pcm$, that of
$\pgrY^m_0$ is $J(\pcb)$.)

\begin{definition}[The diagonal period]\label{pgv:def-period}
Let $\varpi_\pcm\in\auA_{E,f}^\times$ be a finite idele that is $1$ at the primes not dividing
$\pcm$ and has order $\operatorname{ord}_w\pcm$ at each $w\mid\pcm$ (it generates $\pcm\widehat
\auO_E$), and put $u_0=\alpha\varpi_\pcm$ and $\pgupt=\pgupt_\pcm=\aubar u_0/u_0$, an element of
$\auU(1)(\auA_{E^+,f})$ depending only on $\pcm$ (and the fixed idele $\alpha$ of
\S\ref{ps:ss-weil}). For $0\le l\le m$, a real point $g$ of $\auU(\pgrY^l_l)$, a form $y''$ of type
$e_{(0,l,l)}$ and a form $\varphi''$ of type $e_{(1,l,l)}$, the diagonal period is
\[
\mathsf P_{l,g}(y'',\varphi'')=\int_{[\auU(Y^l_l)]}y''(x)\,\varphi''\bigl((x\oplus\pgupt^l)\,g\bigr)\,dx ,
\]
$x\oplus\pgupt^l$ acting on $Y^l_l\oplus E\tilde e_l$ (the line component has trivial archimedean
part).
\end{definition}

\noindent With these definitions the periods $I_\Gamma(\Psi;g)$, path data and transfers of
\S\ref{pgr:ss-arch} are defined without further choice. The integral converges absolutely: $y''$ and
$\varphi''$ are bounded (Theorem~\ref{pgv:G3}(c)) and $[\auU(Y^l_l)]$ has finite volume.

\subsection{The facts (G1), (G2), (G4); the idele \texorpdfstring{$\alpha$}{alpha} at \texorpdfstring{$\pcm$}{m}}\label{pgv:ss-G}

\begin{lemma}[Ranks and signs]\label{pgv:count}
Let $p=(z,a,b)$ be a node.
\begin{enumerate}
\item[(a)] $N^1_{a,b}$ is even; $N^0_{a,b}$ is odd if and only if $a+b$ is even.
\item[(b)] The number of negative signs at $p$ is $\lfloor(m-b+z)/2\rfloor$.
\item[(c)] For $0\le l\le m$: $N^1_{l,l}=N^0_{l,l}+1$, and the number of negative signs at $(1,l,l)$
is that at $(0,l,l)$ plus $1$ if $l$ is odd and plus $0$ if $l$ is even.
\item[(d)] For $0\le j<m$: $N^1_{j,j+1}=N^0_{j+1,j}=r+(2j+1)t$, and $(1,j,j+1)$, $(0,j+1,j)$ have the
same number $\lfloor(m-j)/2\rfloor$ of negative signs.
\end{enumerate}
\end{lemma}
\begin{proof}
(a) $N^1_{a,b}=r+1+a(t+1)+b(t-1)$ with $r+1$, $t\pm1$ even; $N^0_{a,b}=r+(a+b)t$ with $r,t$ odd.
(b) Only stem coordinates can be negative, and the coordinate with $a(\varrho)=a'$ is negative iff
$a'>b$ and $a'+b$ is odd. The values $a'=a(\varrho)\ge1$ are $1,\dots,m-1+z$, each once; so the count
is the number of $a'\in(b,m-1+z]$ with $a'\equiv b+1\pmod 2$, which is $\lfloor(m-1+z-b+1)/2\rfloor$.
(c) The ranks: $r+1+l(t+1)+l(t-1)=r+2lt+1$. By (b) the counts are $\lfloor(m-l+1)/2\rfloor$ and
$\lfloor(m-l)/2\rfloor$; as $m$ is even they are equal if $l$ is even and differ by $1$ if $l$ is
odd. (d) $r+1+j(t+1)+(j+1)(t-1)=r+(2j+1)t=r+(j+1)t+jt$; by (b) the counts are
$\lfloor(m-(j+1)+1)/2\rfloor$ and $\lfloor(m-j)/2\rfloor$.
\end{proof}

\noindent Recall the nodes of a diagonal turn at $(j,j)$ (\S\ref{pgu:ss-data}): $V'=Y^j_j=(0,j,j)$,
$W_1=(0,j,j+1)$, $V=Y^{j+1}_{j+1}$, $S_1=(0,j+1,j)$ in the first grid; $\widetilde V'=\pgrY^j_j=
(1,j,j)$, $S=(1,j,j+1)$ (the node at the place $(j,j+1)$ of $h_1$, i.e.\ $\pgrY^{j+1}_j$),
$\widetilde V=\pgrY^{j+1}_{j+1}$ in the second; $L=E\tilde e_j$, $L'=E\tilde e_{j+1}$.

\begin{proposition}[(G1)]\label{pgv:G1}
For $0\le j<m$ there is a bijection $\sigma\colon\mathsf F_S\to\mathsf F_{S_1}$ preserving the signs,
and for any such $\sigma$ the map $E^{\mathsf F_S}\to E^{\mathsf F_{S_1}}$ it induces is an isometry
of spaces with frames (frame to frame, standard lattice to standard lattice). Moreover
$N_{V'}+N_V=N_{W_1}+N_S=2(r+(2j+1)t)$.
\end{proposition}
\begin{proof}
$S$ and $S_1$ are not among the exceptional nodes of Definition~\ref{pgv:def-spaces} (the first index
of $S_1$ and the second index of $S$ are $\ge1$), so both are standard diagonal spaces. By
Lemma~\ref{pgv:count}(d) they have the same rank and the same number of negative signs, so a
sign-preserving bijection of the frames exists, and it carries the diagonal form with entries
$\epsilon_S$ to the one with entries $\epsilon_{S_1}$ and the standard lattice to the standard
lattice. The ranks: $N_{V'}+N_V=(r+2jt)+(r+2(j+1)t)$ and $N_{W_1}=N_S=r+(2j+1)t$.
\end{proof}

\noindent The types at $S$ and $S_1$ are \emph{not} identified by $\sigma$ (for instance the first
kind blocks of $S$ have $t+1$ entries and those of $S_1$ have $t$); (G1) identifies spaces, groups,
lattices and frames only, which is all that \S\ref{pgu:sec} uses.

\begin{proposition}[(G2)]\label{pgv:G2}
For $0\le j<m$ and every $m$-fold twist, with the characters of Definition~\ref{pgv:def-edges},
\[
\alpha_{\widetilde V'S}=\alpha_{V'S_1},\quad\beta_{\widetilde V'S}=\beta_{V'S_1},\quad
\alpha_{S\widetilde V}=\alpha_{S_1V}\,\upsilon,\quad\beta_{S\widetilde V}=\beta_{S_1V}.
\]
Hence (G2) of \S\ref{pgu:ss-data} holds with $\alpha_L=\alpha_{V'S_1}=\mu^0_{j+1}\upsilon^{N_S/2}$,
$\beta^{\mathrm c}=1$, $\alpha^{\mathrm r}=\upsilon$ and $\beta_{L'}=\beta_{S_1V}=\mu^0_{j+1}
\vartheta^{j+1}\upsilon^{\lfloor N_S/2\rfloor}$, all unramified outside $\pcm$ and tamely ramified at
$\pcm$; in particular $\alpha^{\mathrm r}=(\beta^{\mathrm c})^{-1}\upsilon$. Moreover
$\chi_{(\widetilde V',S)}=\chi_{(V',S_1)}=\chi_{(W_1,V)}$ and $\chi_{(S,\widetilde V)}=\chi_{(S_1,V)}=
\chi_{(V',W_1)}$; and for every square $V'=(z,a,b)$, $W_1=(z,a,b+1)$, $W_2=(z,a+1,b)$,
$V=(z,a+1,b+1)$ of a grid, $\chi_{(W_1,V)}=\chi_{(V',W_2)}$ and $\chi_{(W_2,V)}=\chi_{(V',W_1)}$.
\end{proposition}
\begin{proof}
The edges $(\widetilde V',S)=((1,j,j),(1,j,j+1))$ and $(V',S_1)=((0,j,j),(0,j+1,j))$ are horizontal
(an up edge of the second grid, a right edge of the first) with index $j$; so $\nu=\mu^0_{j+1}$ for
both, and the powers of $\upsilon$ are $\lceil N/2\rceil$. For $\alpha$ the relevant ranks are those
of the targets, $N_S=N_{S_1}$ (Lemma~\ref{pgv:count}(d)): equal. For $\beta$ they are
$N_{\widetilde V'}=N_{V'}+1$ and $N_{V'}=r+2jt$, which is odd; so $\lceil(N_{V'}+1)/2\rceil=
\lceil N_{V'}/2\rceil$: equal. The edges $(S,\widetilde V)=((1,j,j+1),(1,j+1,j+1))$ and
$(S_1,V)=((0,j+1,j),(0,j+1,j+1))$ are vertical with index $j$, so $\nu=\mu^0_{j+1}\vartheta^{j+1}$ and
the powers are $\lfloor N/2\rfloor$. For $\beta$ the ranks of the sources agree. For $\alpha$ they are
$N_{\widetilde V}=N_V+1$ and $N_V=r+2(j+1)t$, odd, so $\lfloor(N_V+1)/2\rfloor=\lfloor N_V/2\rfloor+1$:
$\alpha_{S\widetilde V}=\alpha_{S_1V}\upsilon$.

Now $\det(h\oplus k)=\det h\cdot k$ for $h\in\auU(V')$, $k\in\auU(L)=\auU(1)$, so
$\alpha_{\widetilde V'S}(\det(h\oplus k))=\alpha_{V'S_1}(\det h)\alpha_L(k)$ with $\alpha_L=
\alpha_{V'S_1}$, and $\beta_{\widetilde V'S}=\beta_{V'S_1}\beta^{\mathrm c}$ with $\beta^{\mathrm c}=1$;
similarly $\alpha_{S\widetilde V}=\alpha_{S_1V}\alpha^{\mathrm r}$ with $\alpha^{\mathrm r}=\upsilon$,
and $\beta_{S\widetilde V}(\det(h\oplus k))=\beta_{S_1V}(\det h)\beta_{L'}(k)$ with $\beta_{L'}=
\beta_{S_1V}$. These are the identities of (G2). $N_S$ is even ($r$ and $(2j+1)t$ are odd). The
characters $\mu^0_{j+1}$ are tamely ramified at the chosen places and unramified elsewhere,
$\vartheta$ is tamely ramified at $Q_\vartheta$ and unramified elsewhere, $\upsilon$ is unramified at
all finite places: so all four characters are unramified outside $\pcm$ and tamely ramified at
$\pcm$.

The edge character depends only on $\nu_{pq}$, i.e.\ on the kind and the index. $(\widetilde V',S)$,
$(V',S_1)$ and $(W_1,V)=((0,j,j+1),(0,j+1,j+1))$ are horizontal of index $j$; $(S,\widetilde V)$,
$(S_1,V)$ and $(V',W_1)=((0,j,j),(0,j,j+1))$ are vertical of index $j$. In a square of the grid $z$,
$(W_1,V)$ and $(V',W_2)$ are right edges of index $a$, and $(W_2,V)$, $(V',W_1)$ are up edges of
index $b$, of the same kind within each pair.
\end{proof}

\begin{proposition}[(G4)]\label{pgv:G4}
The element $\pgupt_\pcm$ of Definition~\ref{pgv:def-period} lies in $\auU(1)(\auA_{E^+,f})$,
depends only on $\pcm$, and the diagonal period is the integral displayed in (G4) of
\S\ref{pgu:ss-data}. At $v\in S(\pcm)$, $v=w\aubar w$ with $w\mid\pcm$, under
$\auU(1)(E^+_v)\cong E_w^\times$ (projection to the $w$-component) $\pgupt_v$ is
$(\alpha_w\varpi_w^{k_w})^{-1}\aubar{\alpha}_{w}$, where $\varpi_w^{k_w}$ is the $w$-component of
$\varpi_\pcm$, $k_w=\operatorname{ord}_w\pcm$, and $\aubar\alpha_w$ is the $w$-component of
$\aubar\alpha$; outside $S(\pcm)$ it is $\aubar\alpha/\alpha$.
\end{proposition}
\begin{proof}
$\aubar u_0/u_0$ has norm $1$, so it is a finite adelic point of $\auU(1)$; it is defined from $\pcm$
and the fixed $\alpha$ only. (G4) is Definition~\ref{pgv:def-period}. At $v=w\aubar w\in S(\pcm)$ the
component of $\varpi_\pcm$ at $\aubar w$ is $1$, because $\aubar w\nmid\pcm$ ($\pcm$ is prime to
$\aubar\pcm$); the $w$-component of $\aubar u_0$ is $\aubar\alpha_w\cdot\aubar{\varpi_\pcm}{}_{,w}=
\aubar\alpha_w$ (the $w$-component of the conjugate is the conjugate of the $\aubar w$-component).
Outside $S(\pcm)$, $\varpi_\pcm=1$.
\end{proof}

\noindent By Lemma~\ref{pgx:line}, $\pgupt_\pcm$ is exactly the element of the line reversal.

\subsubsection*{The idele $\alpha$ at the primes of $\pcm$}
\S\ref{pgu:ss-line} assumes that $\alpha$ is a unit at the primes of $\pcm$. By \S\ref{ps:ss-weil},
$\alpha=1/P'(\zeta_{q_1})$ at the primes dividing $\omega=(1+\sqrt{-d})/2$ and $\alpha=1$ elsewhere,
$P$ the minimal polynomial of $\zeta_{q_1}$ over $K$; and $P'(\zeta_{q_1})$ generates the relative
different $\mathfrak d_{E/K}$, which divides $q_1$ (proof of Lemma~\ref{pf:frame}).

\begin{proposition}\label{pgv:alpha}
\begin{enumerate}
\item[(a)] $\alpha$ is a unit at every prime $w$ of $E$ with $w\nmid q_1$; in particular at every
chosen place, and at every prime of $\pcm(\vartheta)$ if $\auNm Q_\vartheta$ is prime to $q_1$.
\item[(b)] Let $p\mid q_1$ be a prime that ramifies in $\Q(\zeta_{q_1})$ (every odd $p\mid q_1$, and
$p=2$ if $4\mid q_1$). Then $p$ splits in $K$, and every prime $w$ of $E$ above $p$ is split over $E^+$
($w\ne\aubar w$) and ramified over $K$; $\alpha$ is not a unit at the primes $w\mid p$ that divide
$\omega$.
\end{enumerate}
\end{proposition}
\begin{proof}
(a) At $w\nmid\omega$, $\alpha_w=1$. At $w\mid\omega$ with $w\nmid q_1$, $w\nmid\mathfrak d_{E/K}$, so
$P'(\zeta_{q_1})$ is a unit at $w$. A chosen place lies above $p_0$ or $p_1$, and $p_i\equiv1\pmod{q_1}$
(\S\ref{ss:au-spaces}: $p_i\equiv1\pmod{q_2^2}$ and $q_1\mid q_2$), so $p_i\nmid q_1$.

(b) If $p$ is odd: $\omega\aubar\omega=(1+d)/4$ is divisible by every prime dividing $q_1$
(\S\ref{pf:sec}), so $d\equiv-1\pmod p$, $p\nmid d$, $-d$ is a non-zero square modulo $p$, and $p$
splits in $K$. The prime $2$ splits in $K$ (\S\ref{pf:ss-adm}). Let $D\subseteq\mathrm{Gal}(E/\Q)$ be
the decomposition group of $w$; its image in $\mathrm{Gal}(K/\Q)$ is the decomposition group of $p$ in
$K$, which is trivial. Complex conjugation restricts to the non-trivial element of $\mathrm{Gal}(K/\Q)$,
so it is not in $D$, i.e.\ $w\ne\aubar w$. As $E\supseteq\Q(\zeta_{q_1})$ is ramified at $p$ and $K$ is
not, $w$ is ramified over $K$, so $w\mid\mathfrak d_{E/K}$ and $P'(\zeta_{q_1})$ is not a unit at $w$.
\end{proof}

\noindent So the assumption of \S\ref{pgu:ss-line} holds at the chosen places, but it does \emph{not}
follow from the definitions for the primes of $Q_\vartheta$ above $q_1$: the definition of a twist
(\S\ref{ss:au-spaces}) requires only $\auNm Q$ prime to $MN$ and $Q\cap\aubar Q=\emptyset$, and by
(b) the primes above $q_1$ meet the second condition. With the data of Definition~\ref{pgx:conv} the line reversal holds without the assumption
(Lemma~\ref{pgx:line}), so the definition of a twist is not changed. The rest of this section does
not use the assumption.

\subsection{(G3): the pairs of the grids are edges, and (H) holds}\label{pgv:ss-G3}

\subsubsection*{The block rule for an arbitrary excess}
Let $(S,T)$ be spaces of ranks $n<n'$, $d=n'-n$, with characters $\alpha,\beta$ of integral
weights, and at $\tau$ let $\lambda$ (on the frame $\mathsf R$ of $S$, signs $\varepsilon$) and
$\lambda'$ (on the frame $\mathsf C$ of $T$, signs $\varepsilon'$) be parameters with distinct
entries, the entries of $\lambda$ in $\frac{n-1}2+\Z$ and those of $\lambda'$ in $\frac{n'-1}2+\Z$. Put
$\tau_0=\frac{n'}2-\pgvwt_\tau\alpha$ and $\tau'_0=\pgvwt_\tau\beta-\frac n2$ (Lemma~\ref{ptr:const}).
We say that $(\lambda,\lambda')$ satisfies the \emph{block rule of excess $d$} if there are a set
$\mathsf B\subseteq\mathsf C$ of $d$ columns and a bijection $\varrho\mapsto c(\varrho)$ of $\mathsf R$ onto
$\mathsf C\smallsetminus\mathsf B$ with
\begin{equation}\label{pgv:eq-rule}
\begin{gathered}
\lambda'_{c(\varrho)}=-\lambda_\varrho+\tau_0+\tau'_0,\qquad|\lambda'_{c(\varrho)}-\tau'_0|>\tfrac d2,\qquad
\varepsilon'_{c(\varrho)}=\varepsilon_\varrho\,\pgvsgn(\lambda'_{c(\varrho)}-\tau'_0),\\
\{\lambda'_c:c\in\mathsf B\}=\{\tau'_0+\tfrac{d-1}2-\beta:0\le\beta<d\},\qquad\varepsilon'\text{ constant on }
\mathsf B.
\end{gathered}
\end{equation}
For $d=t$ this is \eqref{ptr:eq-rule}.

\begin{lemma}\label{pgv:rule}
In this situation put $\xi_\varrho=\lambda_\varrho-\tau_0$, $\xi'_c=\lambda'_c-\tau'_0$. Then all
$\xi_\varrho,\xi'_c$ lie in $\frac{d-1}2+\Z$, and the block rule of excess $d$ is the matching
condition \eqref{par:eq-match} with $t$ replaced by $d$. The statements and proofs of
Lemma~\ref{par:relw}, Propositions~\ref{par:exist}, \ref{par:unique}, Lemmas~\ref{par:polar},
\ref{par:casimir} and of the local identity in the proof of Theorem~\ref{par:A3} hold for such
$(\lambda,\lambda')$ with $t$ replaced by $d$.
\end{lemma}
\begin{proof}
The weights are integers, so $\tau_0\in\frac{n'}2+\Z$, $\tau'_0\in\frac n2+\Z$, $\xi_\varrho\in
\frac{n-1-n'}2+\Z=\frac{d-1}2+\Z$ (as $-d-1\equiv d-1\pmod 2$) and $\xi'_c\in\frac{n'-1-n}2+\Z=\frac{d-1}2+\Z$.
The equivalence with \eqref{par:eq-match} is the computation recalled after Lemma~\ref{ptr:const}:
$\eta_{\varrho c(\varrho)}=-\varepsilon_\varrho\varepsilon'_{c(\varrho)}$, and
$\eta_{\varrho c(\varrho)}=\pgvsgn\xi_\varrho=-\pgvsgn\xi'_{c(\varrho)}$ is the sign condition. In the cited
statements and proofs $t$ occurs only as the number $n'-n$ of columns of $T$ not matched with rows:
as the size of $\mathsf B$ and of its string $\{\frac{t-1}2-\beta\}$ (Lemma~\ref{par:relw}(ii),
Proposition~\ref{par:exist}: columns in $\mathsf B$ carry no variable), as the half-integer class
$\frac{t-1}2+\Z$ of the $\xi$ and the bound $|\xi|\ge\frac{t+1}2$ that it implies with $|\xi|>\frac t2$
(Lemma~\ref{par:relw}(i), Proposition~\ref{par:unique}), and in Lemma~\ref{par:casimir} through
$|\lambda'|^2=|\lambda|^2-2s_0\sum\lambda+ns_0^2+t\xi_0^2+\frac{t(t^2-1)}{12}$ and $n'=n+t$, i.e.\ again
as $|\mathsf B|=n'-n$. With the first assertion, each of these uses holds with $d$ in place of $t$; no
other property of $t$ is used (the parity of $t$ enters only through the class of the $\xi$).
\end{proof}

\subsubsection*{The rule at every edge of the grids}
For an edge $(p,q)$ of the grid $z$ at $\tau$ we take as row parameter (on $\mathsf F_p$) and column
parameter (on $\mathsf F_q$)
\[
(\lambda,\lambda')=(\lambda_{p,\tau},-\lambda_{q,\tau})\ \text{for a right edge},\qquad
(\lambda,\lambda')=(-\lambda_{p,\tau},\lambda_{q,\tau})\ \text{for an up edge},
\]
with the signs $\epsilon_p,\epsilon_q$ and the characters $\alpha_{pq},\beta_{pq}$. In the language of
\S\ref{ptr:ss-edge} (type $e_S$ on $S=p$ with conjugate of parameter $\lambda$, type $e_T$ on $T=q$ with
parameter $\lambda'$): for a right edge $e_S=e_p$, $e_T=e_q$, so the lift goes from forms of type $e_p$
to forms of type $e_q$; for an up edge $e_S=\overline{e_p}$, $e_T=\overline{e_q}$, so the transposed
lift goes from forms of type $\overline{e_T}=e_q$ to forms of type $\overline{e_S}=e_p$. These are the
directions of the moves of Definition~\ref{pgv:def-edges}.

\begin{proposition}[Matching along the edges]\label{pgv:match}
For every edge $(p,q)$ of either grid, every $\tau$ and every $m$-fold twist, the pair
$(\lambda,\lambda')$ just defined satisfies the block rule of excess $d_{pq}=N_q-N_p$, with $c$ the
identity on the stem and on the blocks of $p$ and $\mathsf B$ the new block of $q$ (the first kind
block $a$ for a right edge, the second kind block $b$ for an up edge). The constants are:
\begin{center}
\begin{tabular}{lllll}
\hline
edge&$d$&$\tau_0$&$\tau'_0$&$\tau_0+\tau'_0$\\
\hline
right, $z=0$&$t$&$\kappa_{a+1}-\frac12[a+b\text{ odd}]$&$-\kappa_{a+1}+\frac12[a+b\text{ even}]$&$s_q-s_p$\\
up, $z=0$&$t$&$\kappa_{b+1}+\frac12[a+b\text{ odd}]$&$-\kappa_{b+1}-\frac12[a+b\text{ even}]$&$s_p-s_q$\\
right, $z=1$&$t+1$&$\kappa_{a+1}$&$-\kappa_{a+1}$&$0$\\
up, $z=1$&$t-1$&$\kappa_{b+1}$&$-\kappa_{b+1}$&$0$\\
\hline
\end{tabular}
\end{center}
\end{proposition}
\begin{proof}
Write $\lambda^\circ_p$ for $\lambda_p+s_p$ (the entries of \eqref{pgv:eq-param} without the shift);
on the common coordinates of $p$ and $q$ (stem and old blocks) $\lambda^\circ_q=\lambda^\circ_p$, since the
entries \eqref{pgv:eq-param} do not depend on $a,b$ except through $s$. We use $\kappa_{l+1}\ge
\frac32\Delta$, $\kappa_{l'+1}-\kappa_{l+1}\ge\Delta$ for $l'>l$, $k\le r$, and $\Delta=2r^8>4t+4r$.

\emph{Constants.} By Lemma~\ref{pgv:count}(a) and the weights after Definition~\ref{pgv:def-edges},
$\tau_0=\kappa_{l+1}+\frac{N_q}2-\lfloor\frac{N_q+1-c}2\rfloor$ and $\tau'_0=-\kappa_{l+1}+
\lfloor\frac{N_p+1-c}2\rfloor-\frac{N_p}2$. For $z=1$ all ranks are even and these are $\pm\kappa_{l+1}$.
For $z=0$, a right edge is horizontal ($c=0$, ceilings) and an up edge vertical ($c=1$, floors), and
$N^0_{a,b}$ is odd iff $a+b$ is even; this gives the table. In the last column we used
$s_p=\frac12[a+b\text{ odd}]$, $s_q=\frac12[a+b\text{ even}]$ for $z=0$.

\emph{First condition.} Right edge: it reads $-\lambda_q=-\lambda_p+\tau_0+\tau'_0$ on the common
coordinates, i.e.\ $\lambda^\circ_q-s_q=\lambda^\circ_p-s_p-(s_q-s_p)$: true. Up edge: it reads
$\lambda_q=\lambda_p+\tau_0+\tau'_0$, i.e.\ $\lambda^\circ_q-s_q=\lambda^\circ_p-s_p+(s_p-s_q)$: true. (For
$z=1$, $s=0$.)

\emph{The new block.} Right edge, $z=0$: the block $a$ of the first kind of $q$ carries
$\kappa_{a+1}+\frac{t-1}2-\beta-s_q$, $0\le\beta<t$; the entries of $\lambda'=-\lambda_q$ there are
$-\kappa_{a+1}+s_q-\frac{t-1}2+\beta$, i.e.\ the set $\{\tau'_0+\frac{t-1}2-\beta\}$ as $\tau'_0=
-\kappa_{a+1}+s_q$. Right edge, $z=1$: the block has $t+1$ entries $\kappa_{a+1}+\frac t2-\beta$; their
negatives form $\{\tau'_0+\frac t2-\beta:0\le\beta\le t\}$, the string of excess $t+1$. Up edge, $z=0$:
the block $b$ of the second kind of $q$ carries $-\kappa_{b+1}+\frac{t-1}2-\beta-s_q=\tau'_0+\frac{t-1}2-
\beta$. Up edge, $z=1$: it carries the $t-1$ entries $-\kappa_{b+1}+\frac{t-2}2-\beta=\tau'_0+\frac{(t-1)-1}2-
\beta$. In all cases the signs on the new block are $+1$.

\emph{Distances and signs, right edges.} On a common coordinate $\varrho$,
$\lambda'_{c(\varrho)}-\tau'_0=-\lambda_q(\varrho)+\kappa_{a+1}-s_q=\kappa_{a+1}-\lambda^\circ_p(\varrho)$
(for both $z$). Every entry of $\lambda^\circ_p$ is at most $\kappa_{a+1}-\Delta+\frac t2$: the stem
entries are $\le\frac k2$, the first kind blocks $l<a$ are $\le\kappa_a+\frac t2$, the second kind
blocks are negative. So $\lambda'_{c(\varrho)}-\tau'_0\ge\Delta-\frac t2>\frac d2$ and its sign is $+$:
the rule asks $\epsilon_q(\varrho)=\epsilon_p(\varrho)$. This holds: the stem signs \eqref{pgv:eq-sign}
depend only on $z$ and on the second index, which a right edge does not change, and the block signs
are $+1$.

\emph{Distances and signs, up edges.} Here $\lambda'_{c(\varrho)}-\tau'_0=\lambda_q(\varrho)+\kappa_{b+1}+
s_q=\lambda^\circ_p(\varrho)+\kappa_{b+1}$ (for both $z$). For a first kind block this is
$\ge2\kappa_1-\frac t2>0$, far from $0$. For a second kind block $l<b$ it is
$\kappa_{b+1}-\kappa_{l+1}+\frac{t-z-1}2-\beta\ge\Delta-\frac t2$. For a stem coordinate with
$\varrho\le k-2$ it is $\frac k2-1-\varrho+\frac z2+\kappa_{b+1}\ge\frac32\Delta-r$. For a stem coordinate
with $\varrho\ge k-1$ put $i=\varrho+2-k\ge1$, so that $a(\varrho)=i-1$; then it is
$-\kappa-i\Delta+\frac z2+\kappa+(b+\frac32)\Delta=(b+\frac32-i)\Delta+\frac z2$, of absolute value
$\ge\frac\Delta2-\frac12>\frac d2$, positive iff $i\le b+1$, i.e.\ iff $a(\varrho)\le b$. So the rule asks:
the blocks keep their sign $+1$, and a stem coordinate keeps its sign iff $a(\varrho)\le b$. By
\eqref{pgv:eq-sign}, with $a'=a(\varrho)$: if $a'\le b$ both exponents are $2a'$; if $a'\ge b+1$ the
exponent at $q$ is $a'+b+1$, one more than at $p$. This is the rule.
\end{proof}

\begin{lemma}[Gap and regularity]\label{pgv:gap}
For every node $p$ and $\tau$, $\lambda_{p,\tau}$ consists of $N_p$ distinct numbers in
$\frac{N_p-1}2+\Z$; two of them at coordinates of different signs differ by at least $\Delta-t\ge
\frac12(\Delta-t)-1$; and $N_p\le2n_m$. Hence the conclusions of Theorem~\ref{pco:FJ},
Lemma~\ref{pco:decay}, Theorem~\ref{par:A1} and Corollary~\ref{par:A2} hold for the types $e_p$ and
$\overline{e_p}$.
\end{lemma}
\begin{proof}
\emph{Classes.} For $z=0$ the unshifted entries are integers ($k$ even, $t$ odd), and $N^0_{a,b}-1$ is
even iff $a+b$ is even, which is when $s_p=0$. For $z=1$ every entry lies in $\frac12+\Z$ (the stem
$+\frac12$; blocks of even lengths $t\pm1$), and $N^1_{a,b}$ is even.

\emph{Distinctness.} The stem entries are strictly decreasing in $\varrho$; each block is a string of
consecutive numbers; the centres $\pm\kappa_{l+1}$ of the blocks are $\ge\Delta$ apart from each other
and $\ge\frac32\Delta-r$ from the small stem entries ($|\cdot|\le\frac k2$); and the large stem entries
$-\kappa-i\Delta$ are at distance $|i-l-\frac32|\Delta\ge\frac\Delta2$ from the centres $-\kappa_{l+1}$ and
$>\Delta$ from the centres $+\kappa_{l+1}$. All these exceed the half-lengths $\le\frac{t+1}2$ plus $1$.

\emph{Gap.} A negative coordinate is a stem coordinate with $a(\varrho)=i-1>b$, $i=\varrho+2-k$; so its
entry is $-\kappa-i\Delta$ (shifted) with $i\ge b+2$. Its distance to another stem entry is $\ge\Delta$
(large ones) or $\ge\kappa+\Delta$ (small ones); to the second kind blocks $l<b$ it is
$\ge(i-l-\frac32)\Delta-\frac t2\ge\frac32\Delta-\frac t2$; to the first kind blocks it is
$>2\kappa+\Delta$. So any two entries of different signs differ by at least $\Delta-t$.

\emph{Rank.} $N_p\le r+1+m(t+1)+m(t-1)\le2(k+m(t+1))=2n_m$, as $r+1=k+m$ (Lemma~\ref{pgs:boxes}).

\emph{Consequences.} As in the proof of Lemma~\ref{ptr:gap}, the four cited statements use the
parameter only through its regularity, the gap and the rank, and their numerical hypotheses hold if
$\frac12r^8\ge500(N+N')^3+10N^5+x^{25}$ for the ranks $N,N'\le2n_m$ involved. Now $n_m=k+m(t+1)\le
r+2\sqrt x(5r+1)\le13\sqrt x\,r\le13r^{11/10}$ (as $m\le2\sqrt x$ and $x^5\le r$), so the right side is
at most $500\cdot64\cdot13^3r^{33/10}+320\cdot13^5r^{11/2}+r^5\le2\cdot10^8r^{11/2}\le\frac12r^8$, since
$r\ge x^5$ and $x\ge128^{100}$ (\S\ref{pgs:ss-short}, proof of Lemma~\ref{pgs:boxes}).
\end{proof}

\begin{theorem}[(G3)]\label{pgv:G3}
Assume \textup{(V)} and \textup{(D1)--(D5)} (\S\ref{par:ss-rigid}, \S\ref{ptq:ss-H}). For every edge
$(p,q)$ of the grids, with the parameters above and the characters $\alpha_{pq},\beta_{pq}$ of any
$m$-fold twist:
\begin{enumerate}
\item[(a)] a joint harmonic $P_{pq}=\bigotimes_\tau\Phi_\tau$ exists, is unique up to a scalar,
does not depend on the twist, and no monomial of it contains two variables of one column from rows
of opposite signs, or of one row from columns of opposite signs;
\item[(b)] $\mathsf PP_{pq}=\mathsf P'P_{pq}$, with $\mathsf P,\mathsf P'$ as in \S\ref{ptr:ss-edge};
\item[(c)] every form of type $e_p$, $\overline{e_p}$, $e_q$ or $\overline{e_q}$ is a cusp form, rapidly
decreasing on Siegel sets, and the forms of a given level form a finite-dimensional space;
\item[(d)] for every finite datum the lift and the transposed lift of Theorem~\ref{ptr:edge}(d)
converge absolutely; in particular the moves of Definition~\ref{pgv:def-edges} are defined, $\Lambda_X$
maps forms of type $e_p$ to forms of type $e_q$ and $T_X$ maps forms of type $e_q$ to forms of type
$e_p$;
\item[(e)] the statement \textup{(H)} of \S\ref{ptr:ss-mult} holds for $(p,q)$.
\end{enumerate}
This is (G3) of \S\ref{pgu:ss-data}.
\end{theorem}
\begin{proof}
Fix an edge $(p,q)$ and write $S=E^{\mathsf F_p}$, $T=E^{\mathsf F_q}$, of ranks $n=N_p<n'=N_q$,
$d=d_{pq}=n'-n$, $\alpha=\alpha_{pq}$, $\beta=\beta_{pq}$, and $(\lambda,\lambda')$ for the row and column
parameters at $\tau$ fixed before Proposition~\ref{pgv:match}; $e_S$, $e_T$ are the types of
\S\ref{ptr:ss-edge} ($e_S=e_p$, $e_T=e_q$ for a right edge, $e_S=\overline{e_p}$, $e_T=\overline{e_q}$
for an up edge: in both cases the conjugate of $e_S$ has parameter $\lambda$ and $e_T$ has parameter
$\lambda'$), and $\sigma=\omega_{\alpha,\beta,\infty}(\cdot,1)$, $\sigma'=\omega_{\alpha,\beta,\infty}(1,\cdot)$
at $\tau$. The proof follows that of Theorem~\ref{ptr:edge}, which is the proof of
Theorem~\ref{par:A3} and of Proposition~\ref{pz:types} for the pair $(Y_j,Y_{j+1})$; we go through its
steps and say for each what it uses.

\emph{Step 0: the derived action.} Theorem~\ref{pco:V} is stated for $W=\auHom(Y_j,Y_{j+1})$ at $\tau$,
but Steps 1--3 of its proof concern $W=\auHom(V,V')$ for arbitrary nondegenerate hermitian spaces
$V,V'$ over $E_\tau$: they use only (W1$_\infty$), (W2$_\infty$), (M) and (S), which are stated for every
nondegenerate hermitian space, the identity $\det_W\iota(g,h)=\det(h)^n\det(g)^{-n'}$, and that
$W\oplus W'$ is indefinite. So (V) holds for $(S,T)$. The computation of $d\omega^{\mathrm{Mp}}$ in the
proof of Lemma~\ref{par:vac} is made for arbitrary signs $\varepsilon,\varepsilon'$ and ranks (it uses
only $\eta_{rc}\eta_{r'c}=s_{rr'}$ and its analogue for columns). Hence Lemma~\ref{ptr:const} holds for
$(S,T)$: the derived actions are \eqref{par:eq-osc} with $\gamma=\tau_0+\frac{n'}2$,
$\gamma'=\tau'_0+\frac n2$, and $\tau_0,\tau'_0$ are the numbers of Proposition~\ref{pgv:match}. They
depend on $\alpha,\beta$ only through the weights, which do not depend on the twist
($\pgvwt\vartheta=0$).

\emph{Step 1: proof of} (a). By Proposition~\ref{pgv:match} the pair $(\lambda,\lambda')$ satisfies the
block rule of excess $d$, and by Lemma~\ref{pgv:gap} its entries are distinct and lie in
$\frac{n-1}2+\Z$, $\frac{n'-1}2+\Z$. By Lemma~\ref{pgv:rule} this is the matching condition
\eqref{par:eq-match} with $d$ in place of $t$, and Lemma~\ref{par:relw} holds with $d$ in place of $t$:
in its proof, $\xi\in\frac{d-1}2+\Z$ and $|\xi_k|>\frac d2$ give $|\xi_k|\ge\frac{d+1}2$, at most
$\xi_r-\frac{d+1}2$ indices have $0<\xi_k<\xi_r$, $|\sum_{c\in\mathsf B}\varepsilon'_c|=d$, so
$2w_r\ge2\xi_r-1-(2\xi_r-d-1)-d=0$; and for $c\in\mathsf B$, $2\xi'_c+\sum_{c'\in\mathsf B}\pgvsgn(\xi'_{c'}-
\xi'_c)=0$ because $\mathsf B$ is a string of $d$ consecutive numbers centred at $0$. The statement and
the proof of Proposition~\ref{par:exist} do not contain $t$; they use the matching and
Lemma~\ref{par:relw} only, and give a joint harmonic $\Phi_\tau$ with the property of the monomials.
Lemma~\ref{par:polar} does not involve the matching. In the proof of Proposition~\ref{par:unique},
$t$ occurs in the dichotomy $-s\xi_r>\frac t2$ or $s\xi_r>\frac t2$ and in the third case, where
$s\lambda'_{c(r)}>s\lambda'_c$ gives $s\xi'_c<s\xi'_{c(r)}=-s\xi_r<-\frac t2$, so that $c\notin\mathsf B$;
both hold with $d$, since $|\xi'_c|\le\frac{d-1}2$ on $\mathsf B$. So $\Phi_\tau$ is unique up to a
scalar. The passage to $P_{pq}=\bigotimes_\tau\Phi_\tau$ is the first paragraph of the proof of
Theorem~\ref{par:A3}, which uses only that the conditions of Definition~\ref{par:def-harm} are
conditions on the separate factors. The operators \eqref{par:eq-osc}, hence the joint harmonic, depend
only on the signs, on $\lambda,\lambda'$ and on $\gamma,\gamma'$; by Step~0 they do not depend on the
twist.

\emph{Step 2: proof of} (b). As in the proof of Theorem~\ref{par:A3} both sides are tensor products over
$\tau$, and it suffices to prove $\mathsf P\Phi=\mathsf P'\Phi$ at one $\tau$, where
$\mathsf P=\int e_S(x)\sigma(x^{-1})\,dx=\sigma(e_\lambda)$ and $\mathsf P'=\sigma'(e_{\lambda'})$ by the
description of $e_S$, $e_T$ above and $e(x^{-1})=\overline{e(x)}$. By Lemma~\ref{pgv:gap} the
conclusion of Theorem~\ref{par:A1} holds for $e_\lambda$ and $e_{\lambda'}$: they and their derivatives
are integrable, so $\mathsf P$, $\mathsf P'$ are bounded commuting operators. Let $u=\mathsf P\Phi$.
Since $\mathsf P$ is bounded and commutes with $\sigma'$, $u$ is smooth for $\sigma'$, has $T'$-weight
$\Lambda^K(\lambda')$ and is killed by the compact raising operators of the columns. The first identity
of Lemma~\ref{par:casimir} is an identity between the operators \eqref{par:eq-osc} for arbitrary
$n,n',\gamma,\gamma'$; its second assertion uses the matching through
$\lambda'=\{-\lambda_\varrho+s_0\}\cup\{\xi_0+\frac{d-1}2-\beta:0\le\beta<d\}$, so that
$|\lambda'|^2=|\lambda|^2-2s_0\sum\lambda+ns_0^2+d\xi_0^2+\frac{d(d^2-1)}{12}$, and $n'=n+d$; the final
reduction is an identity of polynomials in $n$, $d$, $s_0$, $\xi_0$, valid for every $d$. So
$d\sigma'(\Omega')u=\mathsf P\,d\sigma(\Omega)\Phi+c\,\mathsf P\Phi=\kappa_{\lambda'}u$, the integration
by parts being justified by the integrability of the derivatives of $e_\lambda$. Lemma~\ref{par:rigid}
is stated for every group $\auU(p,q)$ (the compact case included, Lemma~\ref{par:D3}), every regular
parameter and every unitary representation; applied to $\sigma'$, $\lambda'$ and $u$ it gives
$\mathsf P'u=u$. Symmetrically $\mathsf P\mathsf P'\Phi=\mathsf P'\Phi$, and
$\mathsf P\Phi=\mathsf P'\mathsf P\Phi=\mathsf P\mathsf P'\Phi=\mathsf P'\Phi$.

\emph{Step 3: proof of} (c). This is the proof of Proposition~\ref{pz:types}(a) with $G_j$ replaced by
$\auU(S)$ or $\auU(T)$ and $e_j$ by the type concerned. It uses of the type only (A1) and (A2), that is,
the conclusions of Theorem~\ref{par:A1} and Corollary~\ref{par:A2}, which hold by Lemma~\ref{pgv:gap},
and the general theorems quoted there. If the node has no negative sign (the nodes with $b\ge m-1$ in
the first grid and $b=m$ in the second, Lemma~\ref{pgv:count}(b)) the archimedean group is compact, and
the argument is the one given there for $j=m$; if the group over $E^+$ has no proper parabolic subgroup,
the cusp condition is empty and the rest is unchanged.

\emph{Step 4: proof of} (d). This is the proof of Proposition~\ref{pz:types}(b), (c) with
$(G_j,G_{j+1},\Phi_j,\omega_{\mu_{j+1}},e_j,e_{j+1})$ replaced by $(\auU(S),\auU(T),P_{pq},
\omega_{\alpha,\beta},e_S,e_T)$. It uses: Lemma~\ref{pz:growth}, whose statement and proof involve the
pair of spaces only through $D=fnn'$ and the facts (W0), (W1), (W4); the inequality $2s=2D+2\le B=x^{25}$,
which holds here because $f\le x^{12}$ and $n,n'\le2n_m\le26r^{11/10}$ (proof of Lemma~\ref{pgv:gap}), so
$D\le676\,x^{12}r^{11/5}\le x^{24}$; reduction theory for $\auU(S)$, resp.\ $\auU(T)$; the bound of
(c) on Siegel sets; (A1) for $e_S$ and $e_T$ with the weight $\|x\|^B$ (Lemma~\ref{pgv:gap}); the
identity $\mathsf PP_{pq}=\mathsf P'P_{pq}$ of (b); and the finite-dimensionality in (c) for the
truncation argument. Each of these has been established for $(S,T)$, and the argument is otherwise
word for word the same. It gives: the lift maps forms of type $e_S$ to forms of type $e_T$, and the
transposed lift maps forms of type $\overline{e_T}$ to forms of type $\overline{e_S}$. For a right edge
this is $e_p\to e_q$ for the lift; for an up edge the transposed lift maps forms of type $e_q$ to forms
of type $e_p$. These are the moves of Definition~\ref{pgv:def-edges}.

\emph{Step 5: proof of} (e). This is the proof of Theorem~\ref{ptq:H} for $(S,T)$. Its ingredients are:
Lemma~\ref{ptq:P}(a)--(d), whose proofs use the integrability of $e_\lambda$, $e_{\lambda'}$ and of
their derivatives (Lemma~\ref{pgv:gap}), Lemma~\ref{par:rigid}, the facts established for $u=\mathsf
P\Phi$ in Step~2, (D4) and (D5); Lemmas~\ref{ptq:fischer} and~\ref{ptq:tensor}, which are identities
between the operators \eqref{par:eq-osc}; and Proposition~\ref{ptq:howe}, in whose proof the excess
occurs once, in the step on the row Laplacians: there $\lambda_a>\lambda_b$ gives $\xi_a>\xi_b$, and as
all $|\xi_\varrho|>\frac d2$, either $\xi_a>\frac d2$ or $\xi_b<\xi_a<-\frac d2$; the rest of that step
uses Lemma~\ref{par:relw}(i) and Lemma~\ref{par:polar}(c), and the conclusion uses
Proposition~\ref{par:unique}, all available by Step~1. The passage from one $\tau$ to the product and to
$\pzF_s$ at the end of the proof of Theorem~\ref{ptq:H} uses Lemma~\ref{pz:growth}(a) and (A1), as in
Step~4.
\end{proof}

\subsection{The square law on a grid square: Proposition~\ref{pgr:square}}\label{pgv:ss-square}

Fix $z\in\{0,1\}$ and $0\le a,b<m$, and the square $V'=(z,a,b)$, $W_1=(z,a,b+1)$, $W_2=(z,a+1,b)$,
$V=(z,a+1,b+1)$ of \S\ref{pgr:ss-off}, with $e_1=(V',W_1)$, $e_2=(V',W_2)$, $f_1=(W_1,V)$,
$f_2=(W_2,V)$. For an $m$-fold twist $\vartheta$ put $n=b+1$ if $z=0$, $n=a+1$ if $z=1$, and
\[
\nu=\nu^z_{ab}(\vartheta)=\bigl(\mu^0_{a+1}\mu^0_{b+1}\vartheta^{\,n}\bigr)^{-1},\qquad
\chi=(\chi_{e_1}\chi_{e_2})^{-1},\qquad M_0=\tfrac12\bigl(N_{W_1}+N_{W_2}\bigr),\qquad\gamma=\nu^{-1}\upsilon^{M_0}.
\]
By \eqref{pgv:eq-rank}, $N_{W_2}-N_{W_1}=2z$: in the first grid $N_{W_1}=N_{W_2}=M_0$, and in the second
grid $N_{W_2}=N_{W_1}+2$, all ranks are even (Lemma~\ref{pgv:count}(a)) and $M_0=N_{W_1}+1$ is an integer.
\begin{lemma}[(E1) for the grid squares]\label{pgv:E1g}
Assume \textup{(V)} and \textup{(D1)--(D5)}, as in Theorem~\ref{pgv:G3}. For every such square and every
$m$-fold twist, the four pairs $(U_1,A)=(W_1,V)$, $(U_2,A)=(W_2,V)$, $(A',U_1)=(V',W_1)$,
$(A',U_2)=(V',W_2)$, with the characters of Definition~\ref{pgv:def-edges} and the types
$e_1=e_{W_1}$, $e_2=e_{W_2}$, $e=e_V$, $e'=e_{V'}$, satisfy \textup{(E1)} of
Definition~\ref{ptr:def-square}, with the polynomials $P_{f_1},P_{f_2},P_{e_1},P_{e_2}$. We refer to
this statement as \textup{(E1g)}.
\end{lemma}
\begin{proof}
We check the clauses of (E1) one by one. Recall (\S\ref{pgr:ss-paths}) that a step which raises the
first index is a right edge and one which raises the second index is an up edge, and that for a right
edge $(p,q)$ Theorem~\ref{pgv:G3} is applied with $e_S=e_p$, $e_T=e_q$, for an up edge with
$e_S=\overline{e_p}$, $e_T=\overline{e_q}$ (the convention fixed before Proposition~\ref{pgv:match}).

\emph{The four pairs are edges, in the required directions.} $V'=(z,a,b)$, $W_1=(z,a,b+1)$,
$W_2=(z,a+1,b)$, $V=(z,a+1,b+1)$. So $f_1=(W_1,V)$ and $e_2=(V',W_2)$ raise the first index (right
edges), and $f_2=(W_2,V)$ and $e_1=(V',W_1)$ raise the second (up edges). By Theorem~\ref{pgv:G3}(a)--(d)
each of them satisfies the conclusions (a)--(d) of Theorem~\ref{ptr:edge} with $(S,T)$ the pair and the
types just recalled. Hence: $(U_1,A)=f_1$ is an edge from $e_S=e_{W_1}=e_1$ to $e_T=e_V=e$;
$(U_2,A)=f_2$ is an edge from $\overline{e_{W_2}}=\overline{e_2}$ to $\overline{e_V}=\overline e$;
$(A',U_1)=e_1$ is an edge from $\overline{e_{V'}}=\overline{e'}$ to $\overline{e_{W_1}}=\overline{e_1}$; and
$(A',U_2)=e_2$ is an edge from $e_{V'}=e'$ to $e_{W_2}=e_2$. These are exactly the four directions
required in (E1). In each case Theorem~\ref{ptr:edge}(d) gives both the lift from $e_S$ to $e_T$ and the
transposed lift from $\overline{e_T}$ to $\overline{e_S}$, which is what Definition~\ref{ptr:def-square}
uses.

\emph{Independence of the member of $I$.} The members of $I$ are indexed by the $m$-fold twists. The
types $e_p$ are defined by the node parameters \eqref{pgv:eq-param}, which do not involve $\vartheta$;
the joint harmonics $P_{pq}$ do not depend on the twist (Theorem~\ref{pgv:G3}(a)); the characters
$\alpha_{pq},\beta_{pq}$ of Definition~\ref{pgv:def-edges} depend on $\vartheta$ only through the factor
$\vartheta^{c(l+1)}$, of weight $0$, so their weights (computed after Definition~\ref{pgv:def-edges}) are
fixed.

\emph{(H) for $(U_1,A)$ and $(U_2,A)$.} This is Theorem~\ref{pgv:G3}(e) for the edges $f_1$ and $f_2$,
read with the same parameters (and hence the same projections $\mathsf P$, $\mathsf P'$) as above.

Every clause of (E1) is thus a conclusion of Theorem~\ref{pgv:G3}, which assumes only (V) and
(D1)--(D5). No edge needs to be read with the groups exchanged.
\end{proof}

\begin{lemma}[Grid squares are edge squares]\label{pgv:esq}
Assume \textup{(V)} and \textup{(D1)--(D5)}. The square, with the family $I$ of the data $(\pcm(\vartheta),\nu^z_{ab}(\vartheta),
\chi,\text{the characters of Definition~\ref{pgv:def-edges}})$ over the $m$-fold twists $\vartheta$, is
an edge square of excess $t$ and weight $\kappa_{a+1}+\kappa_{b+1}$ (Definition~\ref{ptr:def-square}),
and $\chi$ is the tame character of $\nu$.
\end{lemma}
\begin{proof}
\emph{Ranks.} $N_V=N_{V'}+2t$ (both indices grow by one; $(t+z)+(t-z)=2t$), and
$N_{W_1}+N_{W_2}=N_V+N_{V'}=2M_0$, so $\dim A=\dim A'+2t$ and $N=2M_0$.
\emph{Signs.} By Lemma~\ref{pgv:count}(b), with $\nu_-(b')=\lfloor(m-b'+z)/2\rfloor$: $q=\nu_-(b+1)+
\nu_-(b)$ (second indices of $W_1$, $W_2$) and $q_A+q_{A'}=\nu_-(b+1)+\nu_-(b)$ (second indices of $V$,
$V'$). Moreover $q\le m+1$, so $2q+2\le2m+4\le2r\le N$; and $N\le4n_m\le52r^{11/10}<3\Delta\le
2(\kappa_{a+1}+\kappa_{b+1})$ (proof of Lemma~\ref{pgv:gap}).
\emph{The data.} $\pcm(\vartheta)\ne(1)$ is a modulus; $\mu^0_l$ is tamely ramified at the chosen places
and unramified elsewhere, $\vartheta$ tamely ramified at $Q_\vartheta$ and unramified elsewhere, so $\nu$
is a tame character for $\pcm(\vartheta)$, of weight $\kappa_{a+1}+\kappa_{b+1}$. In the grid $z=0$,
$e_1$ is an up edge of index $b$, vertical, so $\nu_{e_1}=\mu^0_{b+1}\vartheta^{b+1}$, and $e_2$ is a
right edge of index $a$, horizontal, $\nu_{e_2}=\mu^0_{a+1}$; in the grid $z=1$, $\nu_{e_1}=\mu^0_{b+1}$
and $\nu_{e_2}=\mu^0_{a+1}\vartheta^{a+1}$. In both cases $\nu_{e_1}\nu_{e_2}=\nu^{-1}$, so the tame
character of $\nu$ is $(\chi_{e_1}\chi_{e_2})^{-1}=\chi$. By Proposition~\ref{pgv:G2} (opposite edges)
also $\nu_{f_1}\nu_{f_2}=\nu^{-1}$.
\emph{(E1)} is Lemma~\ref{pgv:E1g}.
\emph{(E2).} In each of the three products of (E2) the two edges are one right and one up edge of
the same grid, hence of the two kinds $c,c'$ with $\{c,c'\}=\{0,1\}$, and their $\nu$'s multiply to
$\nu^{-1}$ (the pairs are $\{f_1,f_2\}$ at $A$, $\{e_1,e_2\}$ at $A'$, and $\{f_1,e_1\}$ at $U_1$,
$\{f_2,e_2\}$ at $U_2$, each consisting of an edge of index $a$ and of the corresponding edge of index
$b$). It remains to add the powers of $\upsilon$. At $A=V$ the target characters are $\beta_{f_1}$,
$\beta_{f_2}$, with powers computed from the ranks of the sources $W_1$, $W_2$: for $z=0$,
$N_{W_1}=N_{W_2}=M_0$ and $\lfloor\frac{M_0+1-c}2\rfloor+
\lfloor\frac{M_0+1-c'}2\rfloor=\lceil\frac{M_0}2\rceil+\lfloor\frac{M_0}2\rfloor=M_0$; for $z=1$ both ranks
are even and the sum is $\frac{N_{W_1}}2+\frac{N_{W_2}}2=M_0$. At $A'=V'$ the source
characters are $\alpha_{e_1},\alpha_{e_2}$, with powers computed from the targets $W_1$, $W_2$: again
$M_0$. At $U_l=W_l$ the characters are $\alpha_{f_l}$ (power from $N_V$) and $\beta_{e_l}$ (power from
$N_{V'}$), of the two kinds; as $N_V=N_{V'}+2t$ have the same parity and $N_V+N_{V'}=2M_0$, the sum of a
floor and a ceiling (or of two exact halves) of $\frac{N_V}2$, $\frac{N_{V'}}2$ is $M_0$. So each product
is $\nu^{-1}\upsilon^{M_0}=\gamma$.
\end{proof}

\begin{theorem}[Proposition~\ref{pgr:square}]\label{pgv:square}
Assume \textup{(V)} and \textup{(D1)--(D5)} (the inputs of Theorems~\ref{ptr:square} and~\ref{pgv:G3}).
For $z\in\{0,1\}$ and $0\le a,b<m$ there are $\kappa_{\square}\in\C$ and a function
$c^z_{ab}$ on the $m$-fold twists with values in $\C^\times$, of modulus $\rho'_{ab}$ if $z=0$ and
$\rho'_{ba}$ if $z=1$ up to norm powers, such that for every $\vartheta$, every form $f$ of type
$e_{W_1}$ and every linear functional $l$ on the forms of type $e_{W_2}$, the functions $L_{f,l}$,
$L'_{f,l}$ of \S\ref{pgr:ss-off} satisfy the square law for $\chi$ with constant
$\kappa_{\square}/c^z_{ab}(\vartheta)$.
\end{theorem}
\begin{proof}
By Lemma~\ref{pgv:esq} and Theorem~\ref{ptr:square} there are $\kappa_{\square}\in\C$ and $C>0$ depending only on
the archimedean data, and for each $\vartheta$ a number $c=c^z_{ab}(\vartheta)$ with
$|c|^2=C\,L_*(\pcm(\vartheta),\nu,t)\ne0$, such that for all forms $y_1$ of type $e_1=e_{W_1}$ and $y_2$
of type $\overline{e_{W_2}}$ the functions $L,L'$ of Definition~\ref{ptr:def-square} satisfy the square
law for $\chi$ with constant $\kappa_{\square}/c$.

\emph{From $y_2$ to $l$.} Let $\mathcal F_0$ be the space of forms of type $e_{W_2}$ fixed by
$J_{W_2}(\pcm)$. For $X$ a matrix over $R$ and $k'\in J_T(\pcm)$, Lemma~\ref{pc:level} gives
$\omega_{\alpha,\beta}(1,k')[P,\delta_X]=[P,\delta_X]$ (the reduction of $k'$ is $1$), so every lift
with a point datum modulo $\pcm$ is fixed by $J(\pcm)$ of its target, and likewise every transposed lift
(with $k\in J_S(\pcm)$). Hence $T_{x_{\mathsf N_2}}\Lambda_{x_{\mathsf N_1}}f$ and
$\Lambda_{y_{\mathsf N_2}}T_{y_{\mathsf N_1}}f$ lie in $\mathcal F_0$ for all $x,y$. By
Theorem~\ref{pgv:G3}(c), $\mathcal F_0$ is finite-dimensional, and the Petersson product is positive
definite on it; so there is $w\in\mathcal F_0$ with $l(u)=\langle w,u\rangle=\int_{[\auU(W_2)]}
\aubar w\,u$ for $u\in\mathcal F_0$, and $y_2=\aubar w$ is a form of type $\overline{e_{W_2}}$. Since only
the restriction of $l$ to $\mathcal F_0$ enters $L_{f,l},L'_{f,l}$, we may replace $l$ by
$u\mapsto\int y_2u$.

\emph{Identification.} With $y_1=f$: $L'_{f,l}(y)=\int_{[\auU(W_2)]}y_2\cdot\Lambda_{y_{\mathsf N_2}}
T_{y_{\mathsf N_1}}f$ is $L'(\eta)$ for $\eta=y$ (the lift along $(A',U_2)=e_2$ after the transposed lift
along $(A',U_1)=e_1$). And $L_{f,l}(x)=\int_{[\auU(W_2)]}y_2\cdot T_{x_{\mathsf N_2}}F$ with
$F=\Lambda_{x_{\mathsf N_1}}f$ on $V$; the transposed lift along $f_2$ and the lift of $y_2$ along $f_2$
have the same kernel $\theta_{f_2}(g,h)$, and the double integral of
$y_2(g)\theta_{f_2}(g,h)F(h)$ over $[\auU(W_2)]\times[\auU(V)]$ converges absolutely
(Proposition~\ref{ptr:conv}, whose proof uses only (E1), Lemma~\ref{pgv:E1g}), so by Fubini $L_{f,l}(x)=
\int_{[\auU(V)]}\Lambda_{x_{\mathsf N_1}}f\cdot\Lambda_{x_{\mathsf N_2}}y_2=L(\xi)$ for $\xi=x$ (columns
$\mathsf N_1$ for $U_1$, $\mathsf N_2$ for $U_2$). So $L_{f,l},L'_{f,l}$ satisfy the square law for $\chi$
with constant $\kappa_{\square}/c$.

\emph{Modulus.} By definition (after \eqref{pc:eq-rho}) $L_*(\pcm,\nu,t)=(\prod_{w\mid\pcm}q_w)^{t^2}\prod_{z'=1}^t|L^\pcm(z',\nu)|^2$, with
$\prod_{w\mid\pcm(\vartheta)}q_w=\auNm(\pi_M)\,\auNm Q_\vartheta=M^f\auNm Q_\vartheta$ (the ideal $(\pi_M)$ is prime, and only the primes dividing $\pcm(\vartheta)$ count, not their multiplicity $b$). The primes of $\pcm(\vartheta)$ are
those of $\pcm_\vartheta$, so $L^{\pcm(\vartheta)}$ is the $L$-function of \eqref{pc:eq-rhoprime} with the
convention recalled after \eqref{pc:eq-rho}. For $z=0$, $\nu=(\mu^0_{a+1}\mu^0_{b+1}\vartheta^{b+1})^{-1}$ and
\[
|c^0_{ab}(\vartheta)|^2=C\,M^{ft^2}\prod_{z'=1}^t|L(z',(\mu^0_{a+1}\mu^0_{b+1})^{-1})|^2\cdot
(\auNm Q_\vartheta)^{t^2}\,\rho'_{ab}(\vartheta)^2 ;
\]
the product of the $L$-values does not depend on $\vartheta$ and is not zero, because $c^0_{ab}(1)\ne0$.
So $c^0_{ab}$ has modulus $\rho'_{ab}$ up to norm powers (Definition~\ref{pc:modulus}, with exponent $t^2$). For
$z=1$, $\nu=(\mu^0_{b+1}\mu^0_{a+1}\vartheta^{a+1})^{-1}$, and the same computation gives the modulus
$\rho'_{ba}$.
\end{proof}

\begin{remark}[(E1g) and Theorem~\ref{pgv:G3}]\label{pgv:rem-E1}
Lemma~\ref{pgv:E1g} shows that Theorem~\ref{pgv:G3} supplies every clause of (E1) for the grid squares:
$(W_1,V)$ is a right edge, lifting $e_{W_1}$ to $e_V$; $(W_2,V)$ is an up edge, lifting
$\overline{e_{W_2}}$ to $\overline{e_V}$; $(V',W_1)$ is an up edge, lifting $\overline{e_{V'}}$ to
$\overline{e_{W_1}}$; $(V',W_2)$ is a right edge, lifting $e_{V'}$ to $e_{W_2}$. The reduction of (E1) to
the block rule recorded in \S\ref{ptw:ss-grid} is thereby completed: the block rule is
Proposition~\ref{pgv:match}.
\end{remark}



%% ===== proofs/C-grid5.tex =====
% proofs/C-grid5.tex
% Continuation of proofs/C-grid.tex (pgr:), C-grid2.tex (pgs:), C-grid3.tex (pgu:), C-grid4.tex
% (pgv:): the first path (Proposition pgr:first) and the last path (Proposition pgr:last).
% Labels start with pgw:. Conventions and macros: those of proofs/C-grid.tex and C-determinant.tex.
\providecommand{\auI}{\mathcal I}
\providecommand{\auGL}{\mathrm{GL}}
\providecommand{\auU}{\mathrm U}
\providecommand{\auO}{\mathcal O}
\providecommand{\auA}{\mathbb A}
\providecommand{\auNm}{\operatorname{Nm}}
\providecommand{\auHom}{\operatorname{Hom}}
\providecommand{\aubar}[1]{\overline{#1}}
\providecommand{\pcm}{\mathfrak m}
\providecommand{\pcb}{\mathfrak b}
\providecommand{\pcMat}{\mathrm M}
\providecommand{\pcFF}{\mathfrak F}
\providecommand{\pcF}{\mathcal F}
\providecommand{\auPl}{\Phi}
\providecommand{\psvac}{\mathbf v}
\providecommand{\pgrY}{\widetilde Y}
\providecommand{\pgrG}{\Gamma}
\providecommand{\pgrfirst}{\Gamma_{\mathrm{first}}}
\providecommand{\pgrlast}{\Gamma_{\mathrm{last}}}
\providecommand{\pgrE}{\operatorname{Ed}}
\providecommand{\pgupt}{\mathsf u}
\providecommand{\pgwK}{\mathcal K}
\providecommand{\pgwT}{\mathcal T}
\providecommand{\pgwP}{\mathcal P}
\providecommand{\pgwev}{\operatorname{ev}}
\providecommand{\pgwint}[1]{\int^{#1}}

\section{The first and the last path}\label{pgw:sec}

This section continues \S\S\ref{pgr:sec}, \ref{pgs:sec}, \ref{pgu:sec}, \ref{pgv:sec} and uses their
notation; the grids, their characters, moves, levels and the diagonal period are those of
\S\ref{pgv:ss-def}. \S\ref{pgw:ss-first} reduces Proposition~\ref{pgr:first} (Theorem~\ref{pgw:first})
to three precisely stated inputs: the injectivity of the untwisted chain of the second grid, a
deep-level vanishing theorem, and a compatibility of data under change of modulus.
\S\ref{pgw:ss-last} proves Proposition~\ref{pgr:last} in full, in the form of
Theorem~\ref{pgw:last}, which makes the normalisation of $\Lambda$ explicit.

\subsection{The first path}\label{pgw:ss-first}

\subsubsection*{The columns}
The first path $\pgrfirst$ runs down the first columns of the two grids: the nodes $(0,0,b)=Y_b$
and $(1,0,b)=\pgrY^b_0$, $0\le b\le m$, joined by the up edges $e^z_b=((z,0,b),(z,0,b+1))$, $b<m$.
In the first grid these edges are vertical, with the characters of $a_{0,b}$
(\S\ref{pgv:ss-def}); in the second grid they are horizontal, so their characters do not involve
$\vartheta$, and their tame parts are supported at the chosen places. The moves along them are
transposed lifts $T_X$ (Definition~\ref{pgv:def-edges}). For a node family $g$ and a set $w$ of
edges we use $\chi_w(g)$ and the action $(gS)_e=g_q^{-1}S_eg_p$ of \S\ref{pgr:ss-arch}. Let
$\pcm_0=\pcm_\vartheta=(\pi_M)\prod_{w\in Q_\vartheta}w$; it has the same primes as $\pcm=\pcb
\prod_{w\in Q_\vartheta}w$. For $z\in\{0,1\}$ let $\pcm^z$ be $\pcm_0$ if $z=0$ and $\pcb$ if
$z=1$, and let $\iota$ be the inclusion matrix of frames.

\emph{The chains.} By uniqueness of the joint harmonic (Theorem~\ref{ptr:edge}(a), (G3)), the joint
harmonic $P_{e^0_b}$ is a non-zero multiple $c_bP_{0,b}$ of the harmonic of $a_{0,b}$, with $c_b$
independent of $\vartheta$; replacing $P_{e^0_b}$ by $P_{0,b}$ changes the moves by these constants,
and we absorb them below into the constant $\alpha_{\mathrm{first}}$ (they do not depend on $\vartheta$, $S$ or the
forms). Let $a^{\mathsf t}_{0,b}$ be the transposed lift with the datum of $a_{0,b}$ (modulus $\pcm_0$,
point datum $\delta_\iota$), and
\[
X^{\vartheta,\mathsf t}_{0,m}=w_0a^{\mathsf t}_{0,0}\circ w_1a^{\mathsf t}_{0,1}\circ\dots\circ
w_{m-1}a^{\mathsf t}_{0,m-1}
\]
(conjugate forms on $Y_m$ to conjugate forms on $Y_0$). In the second grid let $T^\flat_b$ be the
transposed lift along $e^1_b$ with the modulus $\pcb$ and the point datum $\delta_\iota$ on
$\pcMat(\auO_E/\pcb)$, and $\Theta^\flat=T^\flat_0\circ\dots\circ T^\flat_{m-1}$; it does not depend on
$\vartheta$. Let $a^\flat_b$ be the lift with the same datum, so that $a^\flat=a^\flat_{m-1}\circ\dots
\circ a^\flat_0$ maps forms of type $\overline{e_{(1,0,0)}}$ on $\pgrY^0_0$ to forms of type
$\overline{e_{(1,0,m)}}$ on $\pgrY^m_0$ (Theorem~\ref{ptr:edge}(d) for the pairs of the grids, (G3)).

\emph{Levels.} By Lemma~\ref{pc:level}: a pair $(k,k')\in K_p\times K_q$ with $k\equiv1$ and
$k'\equiv1$ modulo $\pcm^z$ fixes the data $[P,\delta_\iota]$ of the edge $e^z_b$ for the modulus
$\pcm^z$ (the tame parts of the characters are trivial on such elements, and $\bar k'^{-1}\iota\bar
k=\iota$). Hence $X^{\vartheta,\mathsf t}_{0,m}y'$ is fixed by $J_{Y_0}(\pcm_0)$, $\Theta^\flat\varphi$
by $J_{\pgrY^0_0}(\pcb)$, and $a^\flat y$ by $J_{\pgrY^m_0}(\pcb)$, for all forms $y',\varphi,y$.

\emph{Transposition.} For an edge $(p,q)$ with datum $F$, a form $w$ at $p$ and $x$ at $q$ of the
conjugate types,
\begin{equation}\label{pgw:eq-transp}
\int_{[G_q]}x(h)\,\bigl(\Theta(F)w\bigr)(h)\,dh=\int_{[G_p]}\bigl(\Theta^{\mathsf t}(F)x\bigr)(g)\,w(g)\,dg,
\end{equation}
both sides being $\iint\theta(\omega(g,h)F)w(g)x(h)\,dg\,dh$ by Fubini: theta functions of Schwartz
data have moderate growth, and the forms are rapidly decreasing (Theorem~\ref{ptr:edge}(c), (G3)),
so the double integral converges absolutely.

\subsubsection*{Three inputs}
\begin{enumerate}
\item[(I)] (Injectivity of the second column.) $a^\flat$ is injective on the forms of type
$\overline{e_{(1,0,0)}}$ on $\pgrY^0_0$ fixed by $J_{\pgrY^0_0}(\pcb)$.
\item[(V)] (Deep-level vanishing.) There is a real point $g_0$ of $\auU(\pgrY^0_0)$ such that every
$y\in\auAz$ with $\int_{[\auU(Y_0)]}\aubar{y(x)}\,\psi((x\oplus1)g_0)\,dx=0$ for all forms $\psi$ of
type $e_{(1,0,0)}$ on $\pgrY^0_0$ fixed by $J_{\pgrY^0_0}(\pcb)$ is zero.
\item[(P)] (Change of modulus.) For an edge of the second column and $\Psi$ on
$\pcMat(\auO_E/\pcb)$, the data $[P,\Psi]$ for the modulus $\pcb$ and $[P,\Psi\circ\mathrm{red}]$ for
the modulus $\pcm$ coincide, $\mathrm{red}$ being reduction modulo $\pcb$.
\end{enumerate}
\noindent (I) is Theorem~\ref{pgy:I} and (P) is Proposition~\ref{pgy:P} (for the data of
Definition~\ref{pgx:conv}, without any assumption on $\alpha$). (V) is Theorem~\ref{pgq:V}, with $g_0=\omega$
the point of input (S) there; its inputs (C), (Q), (S) are Theorems~\ref{pgk:C}, \ref{pgk:Q}
and~\ref{pgl:S}.

\subsubsection*{The restriction map}
Let $\pcFF$ be the forms of type $e_{(1,0,m)}$ on $\pgrY=\pgrY^m_0$ fixed by $J(\pcb)$, and $\pcFF_0$
those of type $e_{(1,0,0)}$ on $\pgrY^0_0$ fixed by $J(\pcb)$; both are finite-dimensional
(Theorem~\ref{ptr:edge}(c)). For a form $y$ on $Y_0$ of type $e_0$ and $\psi\in\pcFF_0$ the diagonal
period at $l=0$ is $\mathsf P_{0,g}(\aubar y,\psi)=\int_{[\auU(Y_0)]}\aubar{y(x)}\psi((x\oplus1)g)\,dx$
(Definition~\ref{pgv:def-period}, $\pgupt^0=1$).

\begin{lemma}\label{pgw:onto}
 $\Theta^\flat(\pcFF)=\pcFF_0$.
\end{lemma}
\begin{proof}
$\Theta^\flat(\pcFF)\subseteq\pcFF_0$: the levels give the invariance under $J(\pcb)$, and the type of a transposed lift is given by Theorem~\ref{ptr:edge}(d). Let $\psi\in\pcFF_0$ be orthogonal to
$\Theta^\flat(\pcFF)$ for the Petersson product. Put $U=a^\flat(\aubar\psi)$; it is fixed by
$J(\pcb)$, so $\aubar U\in\pcFF$. By \eqref{pgw:eq-transp} applied $m$ times,
$\langle U,U\rangle=\int\aubar U\,U=\int\Theta^\flat(\aubar U)\cdot\aubar\psi=\aubar{\langle\Theta^\flat
(\aubar U),\psi\rangle}=0$, with $\langle\varphi',\varphi\rangle=\int\aubar{\varphi'}\varphi$. So $U=0$ and
$\psi=0$ by (I) (Theorem~\ref{pgy:I}). As $\pcFF_0$ is finite-dimensional, $\Theta^\flat(\pcFF)=\pcFF_0$.
\end{proof}

\begin{lemma}\label{pgw:res}
Let $g$ be a real point of $\auU(\pgrY^0_0)$. There is a linear map $\mathsf r_g\colon\pcFF\to\auAz$
such that for every $m$-fold twist $\vartheta$, every $\varphi\in\pcFF$ and every form $y$ of type
$e_0$ on $Y_0$ fixed by $J_{Y_0}(\pcm_\vartheta)$,
$\langle y,\mathsf r_g\varphi\rangle=\mathsf P_{0,g}(\aubar y,\Theta^\flat\varphi)$. For the point
$g=g_0$ of \textup{(V)}, $\mathsf r_{g_0}$ is onto.
\end{lemma}
\begin{proof}
$\auAz=\pcF(Y_0,e_0)^{J_M}$ is finite-dimensional, so there is a unique $\mathsf r_g\varphi\in\auAz$
with $\langle y,\mathsf r_g\varphi\rangle=\mathsf P_{0,g}(\aubar y,\Theta^\flat\varphi)$ for $y\in\auAz$
(the right side is conjugate linear in $y$). Let now $y$ be fixed by $J(\pcm_\vartheta)$ and $u\in
J_M$. Let $v\in J_M$ be $u_w^{-1}$ at the chosen places $w$ and $1$ elsewhere. Then $v$ is trivial at
the primes of $Q_\vartheta$ and at the chosen places lies in $J_M$, so $v\in J(\pcm_\vartheta)$ fixes
$y$; and $uv$ is $1$ at the chosen places, so $uv\in J(\pcb)$ and $(uv)\oplus1\in J_{\pgrY^0_0}(\pcb)$
fixes $\Theta^\flat\varphi$. Since $g$ is real it commutes with finite points, and the substitution
$x\mapsto x(uv)^{-1}$ gives
\[
\mathsf P_{0,g}(\overline{\rho(u)y},\Theta^\flat\varphi)=\mathsf P_{0,g}(\overline{\rho(uv)y},
\rho((uv)\oplus1)\Theta^\flat\varphi)=\mathsf P_{0,g}(\aubar y,\Theta^\flat\varphi);
\]
also $\langle\rho(u)y,\mathsf r_g\varphi\rangle=\langle y,\mathsf r_g\varphi\rangle$ since
$\mathsf r_g\varphi$ is fixed by $J_M$. Both sides of the claimed identity are therefore unchanged when
$y$ is replaced by its mean $\hat y=\int_{J_M}\rho(u)y\,du\in\auAz$ (volume $1$), and for $\hat y$ the
identity is the definition. If $y\in\auAz$ is orthogonal to the image of
$\mathsf r_{g_0}$, then $\mathsf P_{0,g_0}(\aubar y,\Theta^\flat\varphi)=0$ for all $\varphi\in\pcFF$,
hence $\mathsf P_{0,g_0}(\aubar y,\psi)=0$ for all $\psi\in\pcFF_0$ by Lemma~\ref{pgw:onto}, and $y=0$
by (V) (Theorem~\ref{pgq:V}).
\end{proof}

\begin{lemma}\label{pgw:chain}
For every real point $g$, every $\vartheta$, every conjugate form $y'$ on $Y_m$ and every
$\varphi\in\pcFF$,
$\int_{[\auU(Y_m)]}y'\cdot X^\vartheta_{0,m}(\mathsf r_g\varphi)=\mathsf P_{0,g}\bigl(X^{\vartheta,
\mathsf t}_{0,m}y',\Theta^\flat\varphi\bigr)$.
\end{lemma}
\begin{proof}
By \eqref{pgw:eq-transp} applied to the $m$ edges, the left side is $\int_{[\auU(Y_0)]}(X^{\vartheta,
\mathsf t}_{0,m}y')\cdot\mathsf r_g\varphi$. The form $y=\overline{X^{\vartheta,\mathsf t}_{0,m}y'}$ is
of type $e_0$ and fixed by $J(\pcm_\vartheta)$, so this is $\langle y,\mathsf r_g\varphi\rangle=
\mathsf P_{0,g}(\aubar y,\Theta^\flat\varphi)$ by Lemma~\ref{pgw:res}.
\end{proof}

\subsubsection*{The path datum}
For an edge tuple $S$ let
\[
\Psi^1_\vartheta(S)=\prod_{z\in\{0,1\}}\prod_{b<m}\mathbf 1\bigl[S_{e^z_b}\equiv\iota\bmod\pcm^z\bigr]
\in\{0,1\}.
\]
Let $\mathcal C$ be the set of families $x=(x_{z,b})_{z\in\{0,1\},\,b<m}$ of left cosets
$x_{z,b}\in\auGL_{\mathsf F_{(z,0,b)}}(R)/H_{z,b}$, $H_{z,b}$ the level of the node $(z,0,b)$
(\S\ref{pgv:ss-def}: $H_{0,b}=1$, $H_{1,b}=\ker(\to\auGL(\auO_E/\pcb))$). Fix representatives
$\gamma_{z,b}\in x_{z,b}$ and let $g_x$ be the node family equal to $\gamma_{z,b}$ at $(z,0,b)$ for
$b<m$, except $g_x=\gamma_{0,0}\oplus1$ at $(1,0,0)$, and to $1$ at all other nodes. Put
\[
\Psi^0_\vartheta(S)=\sum_{x\in\mathcal C}\chi_{\pgrE(\pgrfirst)}(g_x)\,\Psi^1_\vartheta(g_xS).
\]

\begin{lemma}\label{pgw:datum}
\begin{enumerate}
\item[(a)] If $h$ is a node family with $h_{(z,0,b)}\in H_{z,b}$ for all $b\le m$ \textup{(}with
$H_{z,m}$ the levels of the outer nodes\textup{)}, then $\chi_{\pgrE(\pgrfirst)}(h)=1$ and
$\Psi^1_\vartheta(hS)=\Psi^1_\vartheta(S)$.
\item[(b)] $\Psi^0_\vartheta$ is a path datum on $\pgrE(\pgrfirst)$ with bound $\#\mathcal C\le
\mathcal E(\vartheta)^{2m}$.
\end{enumerate}
\end{lemma}
\begin{proof}
(a) At the nodes of the first column $h=1$. At $(1,0,b)$, $h\equiv1\bmod\pcb$, and the characters of
the second column are trivial on such determinants (their tame parts live at the chosen places).
Moreover $h_q^{-1}S_eh_p\equiv S_e$ modulo $\pcm^z$.

(b) $\Psi^0_\vartheta(S)$ depends only on the $S_e$, $e\in\pgrE(\pgrfirst)$, and is a sum of
$\#\mathcal C$ terms each $0$ or a root of unity; $\#\mathcal C\le\prod_{z,b}\#\auGL_{N_{(z,0,b)}}(R)
\le\mathcal E^{2m}$. \emph{Left inverses.} The vertical edges of $\pgrfirst$ are the $e^0_b$. If
$\Psi^0_\vartheta(S)\ne0$, some $\Psi^1_\vartheta(g_xS)\ne0$, so $Y=(g_xS)_{e^0_b}\equiv\iota\bmod
\pcm_0$; then $\iota^{\mathrm t}Y\equiv1$ modulo $\pcm_0$, hence is invertible modulo $\pcm$ (same
primes), $Y$ has a left inverse, and so has $S_{e^0_b}=g_{x,q}Yg_{x,p}^{-1}$. \emph{Invariance.} Let
$g$ be admissible. Left multiplication by $g_{(z,0,b)}$ permutes the cosets; let $\sigma(x)$ be the
image of $x$. Then $(g_xgS)$ is $((gg_x)S)$ with $(gg_x)_p=g_pg_{x,p}$, and $gg_x=g_{\sigma x}h$ with
$h_p=g_{\sigma x,p}^{-1}g_pg_{x,p}\in H_p$ at the nodes $(z,0,b)$, $b<m$; at the outer nodes
$h_p=g_p\in H_p$ (admissibility); and at $(1,0,0)$, as $H_{0,0}=1$ forces $g_{\sigma x,(0,0,0)}=
g_{(0,0,0)}g_{x,(0,0,0)}$, admissibility ($g_{(1,0,0)}=g_{(0,0,0)}\oplus1$) gives $h_{(1,0,0)}=1$. Since
$((g_{\sigma x}h)S)=h(g_{\sigma x}S)$ and $\chi$ is multiplicative, (a) gives
$\chi(g)\chi(g_x)\Psi^1((gg_x)S)=\chi(g_{\sigma x})\Psi^1(g_{\sigma x}S)$. Summing over $x$:
$\chi_{\pgrE(\pgrfirst)}(g)\Psi^0_\vartheta(gS)=\Psi^0_\vartheta(S)$.
\end{proof}

\begin{lemma}\label{pgw:transl}
Let $g$ be a node family with $g=1$ at $(0,0,m)$ and $(1,0,m)$ and $g_{(1,0,0)}=g_{(0,0,0)}\oplus1$.
For every real point $g_\infty$, every conjugate form $y'$ on $Y_m$, every form $\varphi$ on
$\pgrY^m_0$ and every $S$, with $h_z=$ the column of the grid $z$,
\[
\mathsf P_{0,g_\infty}\bigl(h_{0*}(gS)y',h_{1*}(gS)\varphi\bigr)=\chi_{\pgrE(\pgrfirst)}(g)\,
\mathsf P_{0,g_\infty}\bigl(h_{0*}(S)y',h_{1*}(S)\varphi\bigr).
\]
\end{lemma}
\begin{proof}
Choose $k_p\in K_p$ reducing to $g_p$ at each node $(z,0,b)$, with $k=1$ at the outer nodes and
$k_{(1,0,0)}=k_{(0,0,0)}\oplus1$. For an edge $e=(p,q)$ of a column, $X=S_e$ and $Y=(gS)_e=
g_q^{-1}Xg_p$, Lemma~\ref{pc:level} gives $\omega(k_p,k_q)[P,\delta_Y]=\chi_e(\det g_q/\det g_p)
[P,\delta_X]$ (as $\bar k_qY\bar k_p^{-1}=X$, and both characters of the edge act through
$\chi_e$ on determinants of elements of $K$), so $\theta(gk_p,hk_q;P,\delta_Y)=\chi_e(\det g_q/\det
g_p)\theta(g,h;P,\delta_X)$ and, substituting $h\mapsto hk_q$ in the transposed lift,
$\rho(k_p)T_Y(F)=\chi_e(\det g_q/\det g_p)\,T_X(\rho(k_q)F)$. By induction down the column of the grid
$z$, $\rho(k_{(z,0,0)})h_{z*}(gS)F=\chi_{w_z}(g)h_{z*}(S)F$, $w_z$ the edges of the column. Finally
$\mathsf P_{0,g_\infty}(\rho(k)y'',\rho(k\oplus1)\varphi'')=\mathsf P_{0,g_\infty}(y'',\varphi'')$ for
finite $k$ (substitute $x\mapsto xk^{-1}$; $g_\infty$ commutes with $k\oplus1$).
\end{proof}

\begin{theorem}[Proposition~\ref{pgr:first}]\label{pgw:first}
With $g_0$ as in \textup{(V)}, $\mathsf r=
\mathsf r_{g_0}$ and $\Psi^0=(\Psi^0_\vartheta)$, there is $\alpha_{\mathrm{first}}\in\C$ such that for every period
triple $(\vartheta,y',\varphi)$ with $\varphi\in\pcFF$
\[
\int_{[\auU(Y_m)]}y'\cdot X^\vartheta_{0,m}(\mathsf r\varphi)=\alpha_{\mathrm{first}}\,I_{\pgrfirst}(\Psi^0;g_0)
(\vartheta,y',\varphi);
\]
$\Psi^0_\vartheta$ is a path datum on $\pgrE(\pgrfirst)$ with bound $\mathcal E(\vartheta)^{2m}$, and
$\mathsf r$ is onto $\auAz$.
\end{theorem}
\begin{proof}
The inputs (I), (V), (P) hold: Theorem~\ref{pgy:I}, Theorem~\ref{pgq:V} (with $g_0=\omega$) and
Proposition~\ref{pgy:P}. $\mathsf r$ is onto by Lemma~\ref{pgw:res}, and $\Psi^0_\vartheta$ is a path datum by
Lemma~\ref{pgw:datum}. By Lemma~\ref{pgw:chain} the left side is $\mathsf P_{0,g_0}(X^{\vartheta,
\mathsf t}_{0,m}y',\Theta^\flat\varphi)$. \emph{Expansion.} For the edge $e^0_b$ the datum
$\delta_\iota$ modulo $\pcm_0$ and the function $X\mapsto\mathbf 1[X\equiv\iota\bmod\pcm_0]$ modulo
$\pcm$ give the same function $\Psi^\natural$ (Proposition~\ref{pgy:P} with $\pcm'=\pcm\subseteq\pcm_0$, as $\pcb\subseteq(\pi_M)$, the generators $a_w$ being the same), so
$a^{\mathsf t}_{0,b}=\sum_{X\equiv\iota\,(\pcm_0)}T_X$ up to the constant $c_b$; for $e^1_b$, (P) gives
$T^\flat_b=\sum_{X\equiv\iota\,(\pcb)}T_X$, the sums over $X\in\pcMat(R)$. Expanding the composites
(edge tuples on $\pgrE(\pgrfirst)$ are the pairs of tuples on the two columns),
\[
\mathsf P_{0,g_0}(X^{\vartheta,\mathsf t}_{0,m}y',\Theta^\flat\varphi)=W(\vartheta)\sum_{S\in\mathcal S_\vartheta
(\pgrfirst)}\Psi^1_\vartheta(S)\,\mathsf P_{0,g_0}\bigl(h_{0*}(S)y',h_{1*}(S)\varphi\bigr),\qquad
W(\vartheta)=\prod_{b<m}\frac{w_b^\vartheta}{c_b}.
\]
\emph{Symmetrisation.} Each $g_x$ satisfies the hypotheses of Lemma~\ref{pgw:transl}, and $S\mapsto
g_xS$ is a bijection of $\mathcal S_\vartheta(\pgrfirst)$; so
$\sum_S\chi(g_x)\Psi^1(g_xS)\mathsf P(S)=\sum_S\Psi^1(g_xS)\mathsf P(g_xS)=\sum_S\Psi^1(S)\mathsf P(S)$,
writing $\mathsf P(S)$ for the period in the display. Summing over $x$,
$I_{\pgrfirst}(\Psi^0;g_0)(\vartheta,y',\varphi)=\#\mathcal C\cdot\sum_S\Psi^1_\vartheta(S)\mathsf P(S)$.
\emph{Constants.} $\#\mathcal C=\prod_{b<m}\#\auGL_{N_{(0,0,b)}}(R)\,\#\auGL_{N_{(1,0,b)}}(\auO_E/\pcb)$
and $w^\vartheta_b=\#\auGL_{N_{(0,0,b)}}(\auO_E/\pcm_0)$. As in the proof of Theorem~\ref{pgw:last},
the Chinese remainder theorem gives $\#\auGL_N(R)=\kappa_N\#\auGL_N(\auO_E/\pcm_0)$ with $\kappa_N=
\#\auGL_N(\auO_E/\pcb)/\#\auGL_N(\auO_E/(\pi_M))$ independent of $\vartheta$. Hence
$W(\vartheta)/\#\mathcal C=\alpha_{\mathrm{first}}$ does not depend on $\vartheta$, and the theorem follows.
\end{proof}

\subsection{The last path}\label{pgw:ss-last}

Throughout this subsection $\vartheta$ is an $m$-fold twist, $\pcm=\pcm(\vartheta)$, and $p,q$ denote
nodes $(z,a,m)$ of the row $b=m$ of either grid. By \eqref{pgv:eq-sign}, $\epsilon_p(\varrho)=+1$
whenever $a(\varrho)\le b$; for $b=m$ this holds for every stem coordinate, since
$a(\varrho)\le m-1+z\le m$. Hence \emph{every node of the row $m$ is a positive definite space},
and so is every space on the last path $\pgrlast$, whose halves run through the nodes
$(0,a,m)=Y^a_m$ and $(1,a,m)=\pgrY^m_a$, $0\le a\le m$, along right edges.

For a definite node $p$ write $G_p=\auU(p)$, $G_{p,\infty}$ for its group of real points (a compact
group, of total mass $1$), $G_{p,f}$ for its finite adelic points, $\Gamma_p=G_p(E^+)$, and $K_p$
for the stabiliser of the lattice. Since $G_{p,\infty}$ is compact, $\Gamma_p$ projects injectively
onto a discrete subgroup of $G_{p,f}$ (an element with trivial finite components is $1$), which we
denote again by $\Gamma_p$; the stabilisers $\Gamma_p\cap xJx^{-1}$ of open compact subgroups $J$ are
finite. For an open compact subgroup $J\subseteq G_{p,f}$ and a function $F$ on
$\Gamma_p\backslash G_{p,f}$ that is right invariant under an open subgroup $J'\subseteq J$ put
\[
\pgwint{J}_{\Gamma_p\backslash G_{p,f}}F(x)\,dx=\sum_{x\in\Gamma_p\backslash G_{p,f}/J'}\frac{F(x)}{[J:J']\,
\#(\Gamma_p\cap xJ'x^{-1})},
\]
the integral for the Haar measure giving $J$ the volume $1$ (a finite sum, as $\Gamma_p\backslash
G_{p,f}$ is compact). It does not depend on $J'$, and $\int^{J}=[K_p:J]\int^{K_p}$ for $J\subseteq
K_p$. With the measures of \S\ref{pc:ss-data} (mass $1$ on $G_{p,\infty}$, volume $1$ on $K_p$),
Fubini's theorem gives, for continuous $F$ on $[G_p]=\Gamma_p\backslash(G_{p,\infty}\times G_{p,f})$ right invariant under an open subgroup $J'$ of $G_{p,f}$ (the group $\Gamma_p\cap xJ'x^{-1}$ acts freely on $G_{p,\infty}\times xJ'$, so the fibre over $x$ has mass $\mathrm{vol}(J')/\#(\Gamma_p\cap xJ'x^{-1})$),
\begin{equation}\label{pgw:eq-unfold}
\int_{[G_p]}F=\pgwint{K_p}_{\Gamma_p\backslash G_{p,f}}\int_{G_{p,\infty}}F(ux)\,du\,dx .
\end{equation}
The \emph{level} $J_p\subseteq K_p$ of the node $p$ is the preimage of the level of $p$ in
$\auGL_{\mathsf F_p}(R)$ (\S\ref{pgv:ss-def}): $J_p=J_p(\pcm)$ for $p=(0,0,m)$ (the node $Y_m$) and
$J_p=J_p(\pcb)$ otherwise. Thus $[K_p:J_p]=\#\auGL_{N_p}(\auO_E/\pcb)$ does not depend on
$\vartheta$ for $p\ne(0,0,m)$.

\subsubsection*{Vector forms}
Let $p$ be definite and $(W,\sigma)$ a finite-dimensional complex representation of $G_{p,\infty}$;
an element $\gamma\in\Gamma_p$ acts by $\sigma(\gamma)=\sigma(\gamma_\infty)$. A \emph{$W$-valued form}
is a function $\psi\colon G_{p,f}\to W$, right invariant under an open subgroup, with
$\psi(\gamma x)=\sigma(\gamma)\psi(x)$ for $\gamma\in\Gamma_p$. For $\mu\in W^*$ put
\[
f_{\mu,\psi}(g)=\mu\bigl(\sigma(g_\infty)^{-1}\psi(g_f)\bigr)\qquad(g\in G_p(\auA)),
\]
a function on $[G_p]$: $\sigma(\gamma_\infty g_\infty)^{-1}\psi(\gamma g_f)=\sigma(g_\infty)^{-1}\psi(g_f)$.
(For a weight $\lambda$ and $\mu=\langle v^+,\cdot\,\rangle$ this is the function $f_\varphi$ of
Lemma~\ref{pc:algpet}; we use an abstract representation and a linear form so as not to depend on
the conventions for $W_\lambda$.)

\begin{lemma}[Forms of a type on a definite space]\label{pgw:vec}
Let $p$ be definite, $e$ a type at $p$ with representation $\pi$ and unit vector $v_0$, and let
$W=\pi^*$ with the contragredient action $\sigma(u)\lambda=\lambda\circ\pi(u)^{-1}$ and
$\mu_0=\pgwev_{v_0}\in W^*$. Then $\psi\mapsto f_{\mu_0,\psi}$ is a bijection from the $W$-valued
forms onto $\pcF(p,e)$, and $f_{\mu_0,\psi}$ is fixed by an open subgroup $J$ if and only if $\psi$ is.
\end{lemma}
\begin{proof}
By \S\ref{pc:ss-data} every $\varphi\in\pcF(p,e)$ is $\Phi(v_0)$ for a map $\Phi\colon\pi\to
L^2([G_p])$ intertwining $\pi$ with right translation; replacing $\Phi$ by its composite with the
projection to the forms fixed by the open subgroup $J$ fixing $\varphi$ (which commutes with right
translation by $G_{p,\infty}$) we may assume that the $\Phi(v)$ are fixed by $J$, and, as $\pi$ is
finite-dimensional and $\varphi$ continuous, that they are continuous (they are linear combinations
of right translates of $\varphi$ by elements of $G_{p,\infty}$, by irreducibility). Put
$\Psi(g)(v)=\Phi(v)(g)$, so $\Psi(g)\in\pi^*$, $\Psi(gu)=\Psi(g)\circ\pi(u)$ for $u\in
G_{p,\infty}$ and $\Psi(\gamma g)=\Psi(g)$. Let $\psi(x)=\Psi(x)$ for $x\in G_{p,f}$ (viewed in
$G_p(\auA)$ with trivial archimedean component). Then $\Psi(g)=\Psi(g_fg_\infty)=\sigma(g_\infty)^{-1}
\psi(g_f)$, hence $\varphi(g)=\Psi(g)(v_0)=f_{\mu_0,\psi}(g)$. For $\gamma\in\Gamma_p$, written
$\gamma=\gamma_\infty\gamma_f$ in $G_p(\auA)$, and $x\in G_{p,f}$: $\Psi(\gamma_fx)=\Psi(\gamma x\,
\gamma_\infty^{-1})=\Psi(\gamma x)\circ\pi(\gamma_\infty)^{-1}=\Psi(x)\circ\pi(\gamma_\infty)^{-1}$, that
is, $\psi(\gamma x)=\sigma(\gamma)\psi(x)$ in the notation above; and $\psi$ is fixed by $J$. Uniqueness: $f_{\mu_0,\psi}(ux)=\psi(x)(\pi(u)v_0)$
determines $\psi(x)$ on the span of the $\pi(u)v_0$, which is $\pi$. Conversely, for a $W$-valued
form $\psi$ the map $\Phi(v)(g)=\psi(g_f)(\pi(g_\infty)v)$ intertwines $\pi$ with right translation,
its values are bounded continuous functions fixed by the subgroup fixing $\psi$, and
$f_{\mu_0,\psi}=\Phi(v_0)$; by \S\ref{pc:ss-data} (right convolution with $e$ is the orthogonal
projection onto the image of $\auHom_{G_{p,\infty}}(\pi,L^2)\otimes v_0$) it lies in $\pcF(p,e)$.
\end{proof}

\subsubsection*{One lift between definite spaces}
Let $(p,q)$ be a right edge of the row $m$ (of either grid), with characters
$\alpha=\alpha_{pq}$, $\beta=\beta_{pq}$ and joint harmonic $P=P_{pq}$ (Definition~\ref{pgv:def-edges}).
Then $\omega_{\alpha,\beta}=\omega_{\alpha,\beta,\infty}\otimes\omega_{\alpha,\beta,f}$, and the
archimedean factor does not depend on $\vartheta$, which is trivial at the infinite places. Let
$\pgwP$ be the space of polynomials in the Fock model of $\auHom(p,q)\otimes\R$; it is stable under
$\omega_{\alpha,\beta,\infty}(G_{p,\infty}\times G_{q,\infty})$, which preserves the degree, and
$P\in\pgwP$. For a finite datum $\Phi_f$ (a vector of the restricted tensor product of the local
Schwartz spaces at the finite places) and $Q\in\pgwP$ the theta functional $\theta(Q\otimes\Phi_f)$ is
defined and linear in $Q$; for a finite-dimensional subspace of $\pgwP$ it is a linear form, and
$\theta(\omega_{\alpha,\beta}(\gamma,\gamma')F)=\theta(F)$ for rational $\gamma,\gamma'$, because
$\alpha,\beta$ are automorphic.

\begin{lemma}[Compact Howe duality]\label{pgw:howe}
$\pgwP=\bigoplus_\nu\pgwP_\nu$, where each $\pgwP_\nu$ is stable under $G_{p,\infty}\times
G_{q,\infty}$ and isomorphic to $A_\nu\otimes B_\nu$ with $A_\nu$, $B_\nu$ irreducible
finite-dimensional representations of $G_{p,\infty}$, $G_{q,\infty}$; the $A_\nu$ are pairwise
inequivalent, and so are the $B_\nu$.
\end{lemma}
\begin{proof}
In the Fock model, $\omega_{\alpha,\beta,\infty}(g,h)$ acts on $\pgwP$ by $f(z)\mapsto c(g,h)\,f(h^{-1}zg)$ on $\prod_\tau\pcMat_{N_q\times N_p}(\C)$, with $c$ a character of $G_{p,\infty}\times G_{q,\infty}$ (the product of $\alpha_\infty(\det g)^{-1}\beta_\infty(\det h)$ and the determinant factor of the Fock model; the splitting characters make this a genuine representation). Up to the twist by $c$, this is the decomposition of the polynomials on $\pcMat_{N_q\times N_p}(\C)$
under $\auGL_{N_q}\times\auGL_{N_p}$ into the modules $W_\nu(\C^{N_q})^*\otimes W_\nu(\C^{N_p})^*$,
each once, $\nu$ running over the diagrams with at most $\min(N_p,N_q)$ rows, applied at each
$\tau\in\auPl$ and restricted to the unitary groups, as in Lemma~\ref{pk:cauchy}
(\cite{HoweTheta}, \cite[Theorem~5.6.7]{GoodmanWallach}, \cite{KashiwaraVergne}); distinct families
$(\nu_\tau)_\tau$ give inequivalent representations of $\prod_\tau\auU(N_p)$ and of
$\prod_\tau\auU(N_q)$, because a diagram with at most $\min(N_p,N_q)$ rows is determined by the
highest weight of $W_\nu(\C^N)$ for $N\ge\min(N_p,N_q)$. Twisting by the character $c$ changes $A_\nu,B_\nu$ by fixed characters and preserves
irreducibility and pairwise inequivalence.
\end{proof}

\begin{lemma}[One lift]\label{pgw:one}
Let $(W,\sigma)$ be an irreducible representation of $G_{p,\infty}$, or $W=0$, and $\mu\in W^*$ (if $W=0$ take $W'=0$). There are an
irreducible representation $(W',\sigma')$ of $G_{q,\infty}$ (possibly $W'=0$), a vector $\mu'\in
W'^*$, and for every finite datum $\Phi_f$ a linear map $\pgwK(\Phi_f)\colon W\to W'$, linear in
$\Phi_f$, all depending only on $(p,q)$, $W$, $\mu$ and not on $\vartheta$, such that for every
$m$-fold twist $\vartheta$ (with $\omega=\omega_{\alpha,\beta}$ formed with $\vartheta$):
\begin{enumerate}
\item[(a)] $\pgwK(\omega_f(\gamma x,\gamma'h)\Phi_f)=\sigma'(\gamma')\,\pgwK(\omega_f(x,h)\Phi_f)\,
\sigma(\gamma)^{-1}$ for $x\in G_{p,f}$, $h\in G_{q,f}$, $\gamma\in\Gamma_p$, $\gamma'\in\Gamma_q$;
\item[(b)] for every $W$-valued form $\psi$ and every finite datum $\Phi_f$,
\[
\int_{[G_p]}\theta\bigl(\omega(g,h)(P\otimes\Phi_f)\bigr)\,f_{\mu,\psi}(g)\,dg=f_{\mu',\psi'}(h),\qquad
\psi'(h)=\pgwint{K_p}_{\Gamma_p\backslash G_{p,f}}\pgwK\bigl(\omega_f(x,h)\Phi_f\bigr)\psi(x)\,dx,
\]
and $\psi'$ is a $W'$-valued form.
\end{enumerate}
\end{lemma}
\begin{proof}
\emph{The data.} If no $A_\nu$ of Lemma~\ref{pgw:howe} is isomorphic to $W$, put $W'=0$. Otherwise
let $\nu_0$ be the unique $\nu$ with $A_\nu\cong W$, $T=B_{\nu_0}$ with its action $\tau$, and fix an
isomorphism $j\colon W\otimes T\to\pgwP_{\nu_0}$ of $G_{p,\infty}\times G_{q,\infty}$-modules. Let
$W'=T^*$ with $\sigma'(h)\lambda=\lambda\circ\tau(h)^{-1}$; it is irreducible. For $Q\in\pgwP$ and
$w\in W$ put
\[
Q_w=\int_{G_{p,\infty}}\mu\bigl(\sigma(u)^{-1}w\bigr)\,\omega_\infty(u,1)Q\,du ,
\]
an integral in the finite-dimensional space of polynomials of degree at most $\deg Q$. The
substitution $u\mapsto u_0u$ gives $Q_{\sigma(u_0)w}=\omega_\infty(u_0,1)Q_w$, so $w\mapsto Q_w$ is
$G_{p,\infty}$-equivariant. Writing $Q=\sum_\nu Q^\nu$ with $Q^\nu\in\pgwP_\nu$, the map $w\mapsto
(Q^\nu)_w$ is an equivariant map $W\to\pgwP_\nu\cong A_\nu^{\oplus\dim B_\nu}$, hence zero unless
$A_\nu\cong W$. So $Q_w=0$ for all $w$ if $W'=0$; otherwise $w\mapsto Q_w$ is an equivariant map
$W\to W\otimes T$, and by Schur's lemma there is a unique $t(Q)\in T$ with $Q_w=j(w\otimes t(Q))$ for
all $w$; $t$ is linear. Since $\omega_\infty(1,h)$ commutes with $\omega_\infty(u,1)$ and $j$ is
$G_{q,\infty}$-equivariant, $t(\omega_\infty(1,h)Q)=\tau(h)t(Q)$. Put $\mu'=\pgwev_{t(P)}\in W'^*$
and define $\pgwK(\Phi_f)\in\auHom(W,W')$ by
\[
\bigl(\pgwK(\Phi_f)w\bigr)(s)=\theta\bigl(j(w\otimes s)\otimes\Phi_f\bigr)\qquad(w\in W,\ s\in T).
\]
None of these depends on $\vartheta$.

(a) For rational $\gamma,\gamma'$, $\omega(\gamma,\gamma')=\omega_\infty(\gamma,\gamma')\otimes
\omega_f(\gamma,\gamma')$ and $\theta\circ\omega(\gamma,\gamma')=\theta$, so
\[
\theta\bigl(j(w\otimes s)\otimes\omega_f(\gamma x,\gamma'h)\Phi_f\bigr)=\theta\bigl(\omega_\infty(\gamma,
\gamma')^{-1}j(w\otimes s)\otimes\omega_f(x,h)\Phi_f\bigr)=\theta\bigl(j(\sigma(\gamma)^{-1}w\otimes
\tau(\gamma')^{-1}s)\otimes\omega_f(x,h)\Phi_f\bigr),
\]
which is $\bigl(\sigma'(\gamma')\pgwK(\omega_f(x,h)\Phi_f)\sigma(\gamma)^{-1}w\bigr)(s)$.

(b) By (a) with $\gamma'=1$, $x\mapsto\pgwK(\omega_f(x,h)\Phi_f)\psi(x)$ is left
$\Gamma_p$-invariant; it is right invariant under the open subgroup fixing $\psi$ and fixing
$\Phi_f$ under $\omega_f(\cdot,1)$, so $\psi'$ is defined, and by (a) with $\gamma=1$,
$\psi'(\gamma'h)=\sigma'(\gamma')\psi'(h)$; $\psi'$ is fixed by the open subgroup of $G_{q,f}$ fixing
$\Phi_f$ under $\omega_f(1,\cdot)$. By \eqref{pgw:eq-unfold} the left side of (b) is
\[
\pgwint{K_p}_{\Gamma_p\backslash G_{p,f}}\int_{G_{p,\infty}}\theta\bigl(\omega_\infty(u,1)Q\otimes
\omega_f(x,h_f)\Phi_f\bigr)\,\mu\bigl(\sigma(u)^{-1}\psi(x)\bigr)\,du\,dx,\qquad Q=\omega_\infty(1,
h_\infty)P,
\]
and the inner integral is $\theta(Q_{\psi(x)}\otimes\omega_f(x,h_f)\Phi_f)$ by linearity of $\theta$
on a finite-dimensional space. If $W'=0$ this vanishes, and so does $f_{\mu',\psi'}$. Otherwise
$Q_{\psi(x)}=j(\psi(x)\otimes\tau(h_\infty)t(P))$, and the inner integral is
\[
\bigl(\pgwK(\omega_f(x,h_f)\Phi_f)\psi(x)\bigr)\bigl(\tau(h_\infty)t(P)\bigr)=\mu'\bigl(\sigma'(h_\infty)^{-1}
\pgwK(\omega_f(x,h_f)\Phi_f)\psi(x)\bigr).
\]
Integrating over $x$ (a finite sum) gives $f_{\mu',\psi'}(h)$.
\end{proof}

\begin{remark}\label{pgw:rem-jets}
$\pgwK$ is the kernel of jets of \eqref{ps:eq-kernel}: by Lemma~\ref{pk:cauchy} one may take
$j(w\otimes s)$ to be the polynomial $z\mapsto s(S_{\nu_0}(z)w)$, $S_{\nu_0}$ the Schur functor, and
then by \eqref{ps:eq-bargmann} $\pgwK(\Phi_f)$ is the matrix of the Taylor coefficients of the theta
function of $\Phi_f$ along the entries of $S_{\nu_0}$. A different choice of $j$ multiplies
$\pgwK$ by a non-zero constant and $\mu'$ by its inverse (Schur's lemma for the irreducible module
$W\otimes T$). Only the existence of $\pgwK$ with (a), (b) is used below; the integrality of the
kernels (needed for Proposition~\ref{pgr:den}) concerns the explicit choice of
Definition~\ref{pgz:def-choice}(a).
\end{remark}

\subsubsection*{The branching pairing}
Let $Y=Y^m_m$ and $\pgrY^m_m=Y\oplus E\tilde e_m$, identified as spaces with frames
(\S\ref{pgv:ss-def}); both are definite, of ranks $n$ and $n+1$. For $x\in G_Y(\auA)$ and
$u\in\auU(1)(\auA_{E^+})$ we write $x\oplus u\in G_{\pgrY^m_m}(\auA)$.

\begin{lemma}[Branching]\label{pgw:branch}
Let $(W_0,\sigma_0)$, $(W_1,\sigma_1)$ be irreducible representations of $G_{Y,\infty}$ and
$G_{\pgrY^m_m,\infty}$, and $\mu_z\in W_z^*$.
\begin{enumerate}
\item[(a)] The bilinear forms $B$ on $W_0\times W_1$ with $B(\sigma_0(x)w_0,\sigma_1(x\oplus1)w_1)=
B(w_0,w_1)$ for $x\in G_{Y,\infty}$ form a space of dimension at most $1$.
\item[(b)] Fix such a form $B$, non-zero if one exists. For every real point $g$ of
$G_{\pgrY^m_m}$ there is $c(g)\in\C$ such that for all $W_0$-valued forms $\eta$ on $G_Y$, all
$W_1$-valued forms $\varphi$ on $G_{\pgrY^m_m}$ and all $u\in\auU(1)(\auA_{E^+,f})$
\[
\int_{[G_Y]}f_{\mu_0,\eta}(x)\,f_{\mu_1,\varphi}\bigl((x\oplus u)g\bigr)\,dx=c(g)\,
\pgwint{K_Y}_{\Gamma_Y\backslash G_{Y,f}}B\bigl(\eta(h),\varphi(h\oplus u)\bigr)\,dh .
\]
\end{enumerate}
\end{lemma}
\begin{proof}
(a) $G_{\pgrY^m_m,\infty}=\prod_\tau\auU(n+1)$ and $W_1=\bigotimes_\tau W_{1,\tau}$. The restriction
of an irreducible representation of $\auU(n+1)$ to $\auU(n)$ (embedded as $x\mapsto x\oplus1$) is
multiplicity free (the branching law of Weyl, \cite[Theorem~8.1.1]{GoodmanWallach}). Hence the restriction of $W_1$ to $\prod_\tau\auU(n)$ is a direct sum of pairwise
inequivalent irreducible representations $\bigotimes_\tau V_\tau$, and an invariant form $B$ is an
element of $\auHom_{G_{Y,\infty}}(W_1,W_0^*)$, a space of dimension at most $1$ since $W_0^*$ is
irreducible.

(b) Put $b_g(w_0,w_1)=\int_{G_{Y,\infty}}\mu_0(\sigma_0(x)^{-1}w_0)\,\mu_1\bigl(\sigma_1(g)^{-1}
\sigma_1(x\oplus1)^{-1}w_1\bigr)dx$. The substitution $x\mapsto yx$ shows that $b_g$ is invariant in
the sense of (a), so $b_g=c(g)B$ for some $c(g)$ (and $b_g=0$ if no non-zero invariant form exists).
The function $h\mapsto B(\eta(h),\varphi(h\oplus u))$ is left $\Gamma_Y$-invariant, because
$\varphi((\gamma h)\oplus u)=\varphi((\gamma\oplus1)(h\oplus u))=\sigma_1(\gamma\oplus1)\varphi(h\oplus u)$
and $(\gamma\oplus1)_\infty=\gamma_\infty\oplus1$; it is right invariant under an open subgroup. For
$x=x_\infty x_f$ the point $(x\oplus u)g$ has archimedean component $(x_\infty\oplus1)g$ and finite
component $x_f\oplus u$, so by \eqref{pgw:eq-unfold} the left side is
\[
\pgwint{K_Y}_{\Gamma_Y\backslash G_{Y,f}}\int_{G_{Y,\infty}}\mu_0\bigl(\sigma_0(x)^{-1}\eta(h)\bigr)\,
\mu_1\bigl(\sigma_1(g)^{-1}\sigma_1(x\oplus1)^{-1}\varphi(h\oplus u)\bigr)\,dx\,dh=\pgwint{K_Y}c(g)\,
B\bigl(\eta(h),\varphi(h\oplus u)\bigr)\,dh. \qedhere
\]
\end{proof}

\subsubsection*{The algebraic lifts along the last path and the sum $\Lambda$}
For $z\in\{0,1\}$ and $1\le a\le m$ let $e^z_a$ be the edge $((z,a-1,m),(z,a,m))$; these are the
edges of the half $h_z$ of $\pgrlast$. Let $\pi^z_0$ be the representation of the type at the outer
node $(z,0,m)$, with unit vector $v^z_0$, $W^z_0=(\pi^z_0)^*$ and $\mu^z_0=\pgwev_{v^z_0}$
(Lemma~\ref{pgw:vec}); applying Lemma~\ref{pgw:one} successively to $e^z_1,\dots,e^z_m$ defines
irreducible representations $W^z_a$ of $G_{(z,a,m),\infty}$, vectors $\mu^z_a\in(W^z_a)^*$ and
kernels $\pgwK^z_a$ for $e^z_a$. For a matrix $X\in\pcMat_{q\times p}(R)$ on the edge $e^z_a=(p,q)$ let
$[\delta_X]_f=\delta_X^\natural\otimes\psvac^\pcm$ be the finite part of the datum of
\S\ref{pc:ss-data}, and define the \emph{algebraic lift}
\[
(\pgwT^z_{a,X}\psi)(h)=\pgwint{J_p}_{\Gamma_p\backslash G_{p,f}}\pgwK^z_a\bigl(\omega_f(x,h)[\delta_X]_f
\bigr)\psi(x)\,dx ,
\]
for the measure giving the level $J_p$ of the node $p$ the volume $1$. For an edge tuple $S$ put
$\pgwT^z_S=\pgwT^z_{m,S_{e^z_m}}\circ\dots\circ\pgwT^z_{1,S_{e^z_1}}$. Fix a form $B$ as in
Lemma~\ref{pgw:branch}(a) for $(W^0_m,W^1_m)$, non-zero if possible. For a function $\Psi$ of edge
tuples, a form $y'$ of type $e_{(0,0,m)}$ (a conjugate form on $Y_m$) and a form $\varphi$ of type
$e_{(1,0,m)}$ on $\pgrY^m_0$, with $y'=f_{\mu^0_0,\psi_{y'}}$ and $\varphi=f_{\mu^1_0,\psi_\varphi}$
(Lemma~\ref{pgw:vec}), put
\begin{equation}\label{pgw:eq-Lambda}
\Lambda(\Psi;y',\varphi)=\sum_{S\in\mathcal S_\vartheta(\pgrlast)}\Psi(S)\,
\pgwint{J_Y}_{\Gamma_Y\backslash G_{Y,f}}B\bigl((\pgwT^0_S\psi_{y'})(h),\,(\pgwT^1_S\psi_\varphi)
(h\oplus\pgupt_\pcm^m)\bigr)\,dh ,
\end{equation}
where $Y=Y^m_m$, $J_Y=J_Y(\pcb)$ is its level and $\pgupt_\pcm$ is the element of
Definition~\ref{pgv:def-period}. This is the precise form of the sum $\Lambda$ of
Proposition~\ref{pgr:last}: the pairing $B$ is the branching pairing, the $\pgwT$ are the algebraic
lifts (with the kernels of jets, Remark~\ref{pgw:rem-jets}), and each class sum is normalised by the
level of the node over which it is taken. Replacing $B$ or the kernels $\pgwK^z_a$ by non-zero
multiples multiplies $\Lambda$ by a constant independent of $\vartheta$, $\Psi$, $y'$, $\varphi$.

\begin{theorem}[Proposition~\ref{pgr:last}]\label{pgw:last}
Let $\Xi_{\mathrm{last}}(\Psi)(\vartheta,y',\varphi)=\Lambda(\Psi_\vartheta;y',\varphi)/
\#\auGL_{r+mt}(R_{\pcm_\vartheta})$, with $\Lambda$ as in \eqref{pgw:eq-Lambda} and
$\pcm_\vartheta=(\pi_M)\prod_{w\in Q_\vartheta}w$. Then $\Psi\propto_{\pgrlast}\Xi_{\mathrm{last}}(\Psi)$
for every family $\Psi=(\Psi_\vartheta)$ of functions of edge tuples.
\end{theorem}
\begin{proof}
Fix a real point $g$ of $G_{\pgrY^m_m}$. Let $\vartheta$, $y'$, $\varphi$, $S$ be given and write
$X_a=S_{e^z_a}$ when the half is clear. For $p=(z,a-1,m)$ the algebraic lift with the measure
normalised by $K_p$ is $[K_p:J_p]^{-1}\pgwT^z_{a,X}$. By Lemma~\ref{pgw:one}(b), applied $m$ times
in each half (the move along $e^z_a$ is the lift with the datum $P_{e^z_a}\otimes[\delta_{X_a}]_f$,
Definition~\ref{pgv:def-edges}),
\[
h_{z*}(S)F=f_{\mu^z_m,\,\psi^z_S},\qquad\psi^z_S=\Bigl(\prod_{a=0}^{m-1}[K_{(z,a,m)}:J_{(z,a,m)}]
\Bigr)^{-1}\pgwT^z_S\psi_F ,
\]
for $F=y'$ ($z=0$) and $F=\varphi$ ($z=1$). By Definition~\ref{pgv:def-period} (that is, (G4)) and
Lemma~\ref{pgw:branch}(b) with $u=\pgupt_\pcm^m$,
\[
\begin{aligned}\mathsf P_{m,g}\bigl(h_{0*}(S)y',h_{1*}(S)\varphi\bigr)&=c(g)\pgwint{K_Y}_{\Gamma_Y\backslash G_{Y,f}}
B\bigl(\psi^0_S(h),\psi^1_S(h\oplus\pgupt_\pcm^m)\bigr)\,dh\\&=\frac{c(g)}{D(\vartheta)}\pgwint{J_Y}
B\bigl((\pgwT^0_S\psi_{y'})(h),(\pgwT^1_S\psi_\varphi)(h\oplus\pgupt_\pcm^m)\bigr)\,dh,\end{aligned}
\]
with $D(\vartheta)=[K_Y:J_Y]\prod_{z}\prod_{a<m}[K_{(z,a,m)}:J_{(z,a,m)}]$. Multiplying by
$\Psi_\vartheta(S)$ and summing over $S\in\mathcal S_\vartheta(\pgrlast)$ gives
$I_{\pgrlast}(\Psi;g)(\vartheta,y',\varphi)=c(g)\Lambda(\Psi_\vartheta;y',\varphi)/D(\vartheta)$.

All the indices in $D(\vartheta)$ are those of $J(\pcb)$, independent of $\vartheta$, except
$[K_{Y_m}:J_{Y_m}(\pcm)]=\#\auGL_N(\auO_E/\pcm)$ for the node $(0,0,m)=Y_m$, $N=r+mt$. The primes of
$\pi_M$ are distinct from those of $Q_\vartheta$ and $(\pi_M)$ is squarefree (\S\ref{pc:ss-data}:
$R_\vartheta$ is a product of fields), so by the Chinese remainder theorem
$\#\auGL_N(\auO_E/\pcm)=\#\auGL_N(\auO_E/\pcb)\prod_{w\in Q_\vartheta}\#\auGL_N(\auO_E/w)=\kappa_0\,
\#\auGL_N(R_{\pcm_\vartheta})$ with $\kappa_0=\#\auGL_N(\auO_E/\pcb)/\#\auGL_N(\auO_E/(\pi_M))$. Hence
$D(\vartheta)=D_0\,\#\auGL_N(R_{\pcm_\vartheta})$ with $D_0>0$ independent of $\vartheta$, and
$I_{\pgrlast}(\Psi;g)(x)=\beta(g)\,\Xi_{\mathrm{last}}(\Psi)(x)$ for all period triples $x$, with
$\beta(g)=c(g)/D_0$.
\end{proof}

\noindent The proof uses only the definitions of \S\ref{pgv:ss-def} (spaces, moves, levels, the
diagonal period (G4)), the compact Howe duality and Weyl's branching law; in particular it does not
use (G3) beyond the definition of the joint harmonic $P_{pq}$ as a polynomial. (That the moves take
forms of type $e_p$ to forms of type $e_q$, which (G3) asserts, is recovered here: $W^z_{a}$ is then
the representation of the type at $(z,a,m)$, but this is not needed.)



%% ===== proofs/C-grid6.tex =====
% proofs/C-grid6.tex
% Continuation of proofs/C-grid.tex (pgr:), C-grid2.tex (pgs:), C-grid3.tex (pgu:), C-grid4.tex
% (pgv:): the line reversal without the unit assumption on alpha, the caveat on (G1), and
% Proposition pgr:den. Labels start with pgx:. Conventions and macros: those of proofs/C-grid3.tex.
\providecommand{\auI}{\mathcal I}
\providecommand{\auGL}{\mathrm{GL}}
\providecommand{\auU}{\mathrm U}
\providecommand{\auO}{\mathcal O}
\providecommand{\auA}{\mathbb A}
\providecommand{\auNm}{\operatorname{Nm}}
\providecommand{\auHom}{\operatorname{Hom}}
\providecommand{\aubar}[1]{\overline{#1}}
\providecommand{\pcm}{\mathfrak m}
\providecommand{\pcMat}{\mathrm M}
\providecommand{\pcR}{\mathsf R}
\providecommand{\psHS}{\mathcal H}
\providecommand{\psrho}{\varrho}
\providecommand{\psvac}{\mathbf v}
\providecommand{\pgrY}{\widetilde Y}
\providecommand{\pgrTr}{\mathrel{\vdash}}
\providecommand{\pcnot}[1]{\par\smallskip{\raggedright\noindent\textsc{Not proved here.} [#1]\par}\smallskip}
\providecommand{\pguK}{\mathsf K}
\providecommand{\pgupt}{\mathsf u}

\section{The line reversal at all primes of the modulus; the identification (G1)}\label{pgx:sec}

This section continues \S\S\ref{pgu:sec}, \ref{pgv:sec} and uses their notation. Lemma~\ref{pgu:line}
assumes that the finite idele $\alpha$ of \S\ref{ps:ss-weil} is a unit at the primes of $\pcm$, and
Proposition~\ref{pgv:alpha}(b) shows that this can fail at the primes of $Q_\vartheta$ above $q_1$.
Here no such assumption is made: the data at the primes of $\pcm$ are placed on the self-dual lattice
$\alpha\auHom$ (Definition~\ref{pgx:conv}), and the reversal of a line is proved with the element
$\pgupt=\pgupt_\pcm=\aubar u_0/u_0$, $u_0=\alpha\varpi_\pcm$, of Definition~\ref{pgv:def-period}
(Lemma~\ref{pgx:line}, which contains Lemma~\ref{pgu:line}). \S\ref{pgx:ss-later} lists the uses of
the line reversal in \S\ref{pgu:sec}, \S\ref{pgx:ss-other} the other places of the paper that use the
data at the primes of a modulus, and \S\ref{pgx:ss-G1} shows that the seesaw steps use (G1) only as an
identification of spaces.

\subsection{The datum at the primes of the modulus}\label{pgx:ss-conv}

Let $\pcm$ be a modulus (\S\ref{pc:ss-data}), $v=w\aubar w\in S(\pcm)$ with $w\mid\pcm$,
$k_w=\operatorname{ord}_w\pcm$, and let $S,T$ be spaces of ranks $N,N'$. By (W5),
$\psHS_v(\auHom(S,T))$ is the space of Schwartz--Bruhat functions on $X_w=\auHom(S_w,T_w)=
\pcMat_{N'\times N}(E_w)$ (matrices in the frames), $\psrho(a)$ is translation by $a$ for $a$ in the
$w$-component and multiplication by $x\mapsto[x,a']$ for $a'$ in the $\aubar w$-component $X_{\aubar w}$.

\begin{lemma}[The vacuum at a split place]\label{pgx:vac}
The self-dual lattice $M_v=\alpha\auHom(L_S,L_T)_v$ is $M_w\oplus M_{\aubar w}$ with
$M_w=\alpha_w\pcMat_{N'\times N}(\auO_w)$ and $M_{\aubar w}$ the dual of $M_w$ for the pairing
$(x,a')\mapsto[x,a']$ of $X_w\times X_{\aubar w}$. In the model (W5) the vacuum $\psvac_{M_v}$ is,
up to a scalar, the indicator function $\mathbf 1_{M_w}$.
\end{lemma}
\begin{proof}
The lattice $L_S\otimes\auO_v$ is $L_{S,w}\oplus L_{S,\aubar w}$, and similarly for $T$, so $M_v$ is
the direct sum of its components $M_w=\alpha_w\auHom(L_{S,w},L_{T,w})$, which is
$\alpha_w\pcMat_{N'\times N}(\auO_w)$ in the frames, and $M_{\aubar w}$. The hermitian form of
$\auHom(S,T)$ pairs the $w$-component with the $\aubar w$-component and vanishes on each of them; so
$[\,,]$ is trivial on $X_w\times X_w$ and on $X_{\aubar w}\times X_{\aubar w}$, and the self-duality
of $M_v$ says exactly that $M_{\aubar w}$ is the dual of $M_w$ under $[\,,]$. The function
$\mathbf 1_{M_w}$ is invariant under the translations by $M_w$ and under the multiplications by
$[\,\cdot\,,a']$, $a'\in M_{\aubar w}$ (these are trivial on $M_w$). On $M_v$ the character
$\epsilon_{M_v}$ of (W4) is trivial on $M_w$ and on $M_{\aubar w}$ (since $\beta$ is trivial on each
component), so by (W4) $\psvac_{M_v}$ is the unique line with these invariances, which is
$\C\mathbf 1_{M_w}$.
\end{proof}

\begin{definition}[Data at the primes of a modulus]\label{pgx:conv}
For every $w\mid\pcm$ let $a_w$ be the generator of the ideal $\alpha_w\auO_w$ fixed in
\S\ref{ps:ss-weil} ($a_w=\alpha_w$ if $\alpha_w$ is not a unit, $a_w=1$ otherwise). For a function $\Psi$ on
$\pcMat_{N'\times N}(\pcR_\pcm)$ let $\Psi^\natural$ be the function on $\prod_{w\mid\pcm}
\pcMat_{N'\times N}(E_w)$ equal to $\Psi(a^{-1}x\bmod\pcm)$ at $x\in a\,\pcMat_{N'\times N}
(\widehat\auO)$ (that is, $x_w\in a_w\pcMat(\auO_w)$ for all $w\mid\pcm$) and to $0$ elsewhere, and
use at $v\in S(\pcm)$ the $L^2$-norm for the measure giving $M_w=a_w\pcMat(\auO_w)$ the volume $1$.
The datum $[P,\Psi]$, the kernels, the lifts and the transposed lifts are then defined by the
formulas of \S\ref{pc:ss-data}.
\end{definition}

\noindent At the primes where $\alpha_w$ is a unit $a_w=1$; then Definition~\ref{pgx:conv}
is literally the definition of \S\ref{pc:ss-data}. The two definitions differ only at the primes of
$\pcm$ where $\alpha$ is not a unit, which by Proposition~\ref{pgv:alpha}(a) lie above $q_1$; at
such a prime a function supported on the integral matrices would not be the vacuum
(Lemma~\ref{pgx:vac}); for this reason \S\ref{pc:ss-data} uses Definition~\ref{pgx:conv}. With it,
$\mathbf 1^\natural=\mathbf 1_{M_w}$ is the vacuum at every $w\mid\pcm$, and
$\|\Psi^\natural\|^2=(\#\pcMat_{N'\times N}(\pcR_\pcm))^{-1}\sum|\Psi|^2$ exactly as before. Another
choice of the generators, $a'_w=a_w\varepsilon_w$ with units $\varepsilon_w$, replaces $\Psi^\natural$ by
$(\Psi\circ\bar\varepsilon^{-1})^\natural$, $\bar\varepsilon$ the image of $\varepsilon$ in
$\pcR_\pcm^\times$ acting by multiplication. Lemma~\ref{pc:level} holds without change: for $k\in K_S$,
$k'\in K_T$ the matrix $k'^{-1}xk$ lies in $a\pcMat(\widehat\auO)$ iff $x$ does, and
$a^{-1}k'^{-1}xk=k'^{-1}(a^{-1}x)k$ since $a$ is a scalar; the proof of that lemma then applies word
for word. The same remark shows that Corollary~\ref{pgu:split} holds with Definition~\ref{pgx:conv}:
$\delta^\natural_{(y_2|v)}(x_1|x_2)=\delta^\natural_{y_2}(x_1)\delta^\natural_v(x_2)$, the lattices
$a_w\pcMat(\auO_w)$ being products of the corresponding lattices of the two summands and the
measures product measures.

\subsubsection*{The coordinates at split places}
The idele $\alpha$ enters the data at a split place in three ways: through the frames (coordinates
$x$), through the coordinates $T$ of (N5) (\S\ref{pl:ss-norm}), and through the coordinates $a^{-1}x$
of Definition~\ref{pgx:conv}. The following lemma fixes one coordinate and gives the dictionary; all
data at split places in the proof of Theorem~C are written in it.

\begin{lemma}[The normalised coordinates]\label{pgx:dict}
Let $a$ be the idele of \S\ref{ps:ss-weil}, $S,T$ spaces of ranks $N,N'$, $v=w\aubar w$ a place of
$E^+$ split in $E$, and $X_w,X_{\aubar w}$ the components of $\auHom(S,T)_v$, written as matrices $x$
in the frames. Put $\tilde x=a_w^{-1}x$ on $X_w$ and $\tilde x=a_{\aubar w}^{-1}x$ on $X_{\aubar w}$, and
more generally $\tilde x=a^{-1}x$ on $\auHom(S,T)(\auA_f)$.
\begin{enumerate}
\item[(a)] $\alpha_w$ is not a unit if and only if $w$ lies above a prime $p\mid q_1$ and $w\mid\omega$;
then $w\ne\aubar w$, $\aubar w\nmid\omega$ and $\alpha_{\aubar w}=a_{\aubar w}=1$. Every place of $E^+$
above $p\mid q_1$ is $w\aubar w$ with exactly one of $w,\aubar w$ dividing $\omega$. No chosen place
and no prime above $N$ lies above $q_1$. A prime of $Q_\vartheta$ may lie above $q_1$ and divide
$\omega$; such primes are the only primes of a modulus $\pcm_\vartheta$, of $S(\pcm_\vartheta)$, or of
the moduli of the grids at which $a_w\ne1$.
\item[(b)] The self-dual lattice $\alpha\auHom(L_S,L_T)$ is $a\auHom(L_S,L_T)$; at $v$ its components
are $M_w=a_w\pcMat(\auO_w)$ and $M_{\aubar w}=a_{\aubar w}\pcMat(\auO_{\aubar w})$, i.e.\ the integral
matrices in the coordinates $\tilde x$, and they are dual to each other for $[\,,]$. For $x\in X_w$,
$y\in X_{\aubar w}$ one has $[x,y]=\lambda_w(\operatorname{tr}({}^t\tilde y\tilde x))^{\pm1}$ with
$\lambda_w(t)=\psi_w(\jmath(a\aubar a)_wt)$, trivial on $\auO_w$ and not on $\varpi_w^{-1}\auO_w$
($E_{\aubar w}$ identified with $E_w$ by conjugation). In the frame coordinates the dual of
$\pcMat(\auO_w)$ is $\aubar{a_w}\,a_{\aubar w}\pcMat(\auO_{\aubar w})$, which is not
$\pcMat(\auO_{\aubar w})$ when $a_w\ne1$ or $a_{\aubar w}\ne1$.
\item[(c)] On $X_w$ the coordinate $T$ of (N5) and the coordinate $a_w^{-1}x$ of
Definition~\ref{pgx:conv} are both $\tilde x$. Hence for $w\mid\pcm$:
$\Psi^\natural=\Psi(\tilde x\bmod\pcm)\,\mathbf 1[\tilde x\ \text{integral}]$; the datum
$\delta_\iota^\natural$ is $\mathbf 1_{a_w(\iota+\varpi\pcMat(\auO_w))}$ in the frames, which is the
datum of \S\ref{ps:ss-Y} at a level place; and a datum defined as a function of $T$ by (N5) (as in
\S\S\ref{po:ss-glue}, \ref{plf:ss-KS}, \ref{pte:sec}) is the same function of $\tilde x$, so it is
of the form $\Psi^\natural$ exactly when, as a function of $T$, it is supported on the integral
matrices and depends only on $T$ modulo $\pcm$.
\item[(d)] For $u\in\auHom(S,T)(\auA_f)$, $u\in\alpha\auHom(L_S,L_T)$ if and only if $\tilde u$ is
integral. If $u=s_0y$ with $s_0=a$, then $\tilde u=y$; for $h\in\auU(T)(\auA_f)$,
$\widetilde{h^{-1}u}=h^{-1}y$.
\item[(e)] Let $u$ be a finite idele of $E$ with $u\aubar u=1$, acting as a central element of
$\auU(T)(\auA_f)$. Then $\omega(1,u)$ carries the vacuum of a self-dual lattice $\Lambda$ of
$\auHom(S,T)(\auA_f)$ to a multiple of modulus one of the vacuum of $u\Lambda$, and for every datum $f$,
all characters $\alpha',\beta'$ and $g\in\auU(S)(\auA)$,
\[
\int_{[\auU(T)]}\theta\bigl(\omega_{\alpha',\beta'}(g,h)\,\omega(1,u)f\bigr)\,dh=\beta'(\det u)^{-1}
\int_{[\auU(T)]}\theta\bigl(\omega_{\alpha',\beta'}(g,h)f\bigr)\,dh
\]
whenever either side converges absolutely.
\item[(f)] The map $w\mapsto w^\dagger$ carries $\alpha\auHom(L_S,L_T)$ onto $\aubar\alpha\auHom(L_T,L_S)=
u\,\alpha\auHom(L_T,L_S)$, $u=\aubar a/a$ (a finite idele with $u\aubar u=1$, central in $\auU(S)$,
equal to $1$ at the places not above $q_1$). At $v$, in the coordinates $\tilde x$ of $\auHom(S,T)$ and
$\tilde X$ of $\auHom(T,S^-)$ (frames of $S$ and $T$), it is $\tilde X=-\nu^{-1}\,{}^t\tilde x$ after
identifying $E_{\aubar w}$ with $E_w$ by conjugation, where $\nu=\aubar{a_{\aubar w}}/a_w$ is the
conjugate of $u_{\aubar w}^{-1}$.
\end{enumerate}
\end{lemma}
\begin{proof}
(a) If $w\nmid\omega$ then $\alpha_w=1$. If $w\mid\omega$ and $w\nmid q_1$, $\alpha_w$ is a unit
(Proposition~\ref{pgv:alpha}(a)). If $w\mid\omega$ and $w$ lies above $p\mid q_1$: every prime $p\mid q_1$
ramifies in $\Q(\zeta_{q_1})$ ($p$ odd, or $p=2$ and $4\mid q_1$, \S\ref{an:ssec-param}), so $\alpha_w$
is not a unit and $w\ne\aubar w$ (Proposition~\ref{pgv:alpha}(b)). Since $\omega+\aubar\omega=1$, $\omega$
and $\aubar\omega$ are coprime; $p$ divides $\omega\aubar\omega=(1+d)/4$ and splits in $K$ as
$\mathfrak p\aubar{\mathfrak p}$, so exactly one of $\mathfrak p,\aubar{\mathfrak p}$ divides $\omega$,
the primes of $E$ above it divide $\omega$, and their conjugates (above the other one) do not. Hence
$\aubar w\nmid\omega$ and $\alpha_{\aubar w}=1$. The chosen places lie above $p_0,p_1$ and the primes
dividing $\pi_N$ above $N$; these are $\equiv1\pmod{q_2^2}$ and $q_1\mid q_2$, so they are prime to
$q_1$. The definition of a twist (\S\ref{ss:au-spaces}) asks only that $\auNm Q$ be prime to $MN$ and
$Q\cap\aubar Q=\emptyset$, which such primes satisfy (Proposition~\ref{pgv:alpha}(b)). The moduli in
question are divisible only by chosen places and primes of $Q_\vartheta$, which gives the last claim.

(b) $a/\alpha$ is a unit idele, so $\alpha\auHom(L_S,L_T)=a\auHom(L_S,L_T)$, whose components at $v$
are $a_w\pcMat(\auO_w)$ and $a_{\aubar w}\pcMat(\auO_{\aubar w})$; their duality is Lemma~\ref{pgx:vac}.
For the pairing: $[x,y]=\psi(\jmath\operatorname{tr}(y^\dagger x))^{\pm1}$ (the sign depending on the
order of the two components), and with $x=a_w\tilde x$, $y=a_{\aubar w}\tilde y$,
$\operatorname{tr}(y^\dagger x)=a_w\aubar{a_{\aubar w}}\operatorname{tr}(\tilde y^\dagger\tilde x)$, the
frames having unimodular Gram matrices at $v$ (which are absorbed in the coordinates, as in
(N5)); so $[x,y]=\lambda_w(\operatorname{tr}({}^t\tilde y\tilde x))^{\pm1}$. The character
$t\mapsto\psi_w(\jmath\alpha_w\aubar{\alpha_{\aubar w}}t)$ is trivial exactly on $\auO_w$ by the
self-duality of $\widehat\auO_E$ for $\psi(\jmath\alpha\aubar\alpha\,\cdot)$ (\S\ref{ps:ss-weil}), and
$a_w\aubar{a_{\aubar w}}/\alpha_w\aubar{\alpha_{\aubar w}}$ is a unit. For the last assertion, for
$c\in E_w^\times$ one has $[cx,y]=[x,\aubar cy]$ ($\aubar c\in E_{\aubar w}$), so the dual of
$\pcMat(\auO_w)=a_w^{-1}M_w$ is $\aubar{a_w}M_{\aubar w}=\aubar{a_w}a_{\aubar w}\pcMat(\auO_{\aubar w})$,
and $\aubar{a_w}a_{\aubar w}$ is a unit only if $a_w=a_{\aubar w}=1$ (by (a) at most one of them is
$\ne1$, and that one is not a unit).

(c) By definition (N5) uses $T=a_w^{-1}x_w$, and Definition~\ref{pgx:conv} uses $a_w^{-1}x$ with the
same $a_w$. The formula for $\Psi^\natural$ is Definition~\ref{pgx:conv} rewritten;
$\delta_\iota^\natural(x)=1$ if and only if $a_w^{-1}x\in\iota+\varpi\pcMat(\auO_w)$. The last
assertion is a restatement.

(d) $\alpha\auHom(L_S,L_T)=a\auHom(L_S,L_T)$ by (b), and $a$ is a scalar commuting with the action of
$\auU(T)(\auA_f)$.

(e) $\omega(1,u)\psrho(l)\omega(1,u)^{-1}=\psrho(ul)$ (W1), so $\omega(1,u)\psvac_\Lambda$ is fixed by
$\psrho(u\Lambda)$ up to the character $\epsilon_{u\Lambda}$, hence is a multiple of
$\psvac_{u\Lambda}$ by (W4), of modulus one as $\omega(1,u)$ is unitary. Since $u$ is central,
$\omega_{\alpha',\beta'}(g,h)\omega(1,u)=\beta'(\det u)^{-1}\omega_{\alpha',\beta'}(g,hu)$, and the
measure on $[\auU(T)]$ is invariant under right translation by $u$.

(f) $(cx)^\dagger=\aubar cx^\dagger$ for $c\in\auA_{E,f}$, and $w\mapsto w^\dagger$ carries
$\auHom(L_S,L_T)$ onto $\auHom(L_T,L_S)$ (unimodular lattices); so $\alpha\auHom(L_S,L_T)$ goes to
$\aubar\alpha\auHom(L_T,L_S)=\aubar a\auHom(L_T,L_S)=u\,a\auHom(L_T,L_S)$, and $u$ acts on
$\auHom(T,S^-)$ through the centre of $\auU(S^-)=\auU(S)$. By (a), $u_v=1$ unless $v$ lies above
$q_1$. In the frames $x^\dagger=-{}^t\aubar x$ at $v$ (as in \S\ref{pcx:ss-flip}, the Gram matrices
being absorbed), and $x=a_w\tilde x$, $x^\dagger=a_{\aubar w}\tilde X$ give
$\tilde X=-a_{\aubar w}^{-1}\aubar{a_w}\,{}^t\aubar{\tilde x}$, which is the formula after the
identification by conjugation.
\end{proof}

\subsection{The line reversal}\label{pgx:ss-line}

We keep the notation of \S\ref{pgu:ss-line}: $n=\#\mathsf N_2$ is the rank of $S$, $L=E\tilde e_j$,
$L'=E\tilde e_{j+1}$ is $L$ with the opposite form, and $\pguK^{\mathrm c}_{v,B}$,
$\pguK^{\mathrm r}_{v',B'}$ are the line kernels of Corollary~\ref{pgu:split}, now with the data of
Definition~\ref{pgx:conv}. No assumption on $\alpha$ at the primes of $\pcm$ is made.

\begin{lemma}[Line reversal]\label{pgx:line}
Let $\pgupt=\pgupt_\pcm=\aubar u_0/u_0\in\auU(1)(\auA_{E^+,f})$, $u_0=\alpha\varpi_\pcm$, be the
element of Definition~\ref{pgv:def-period}. There are a linear bijection $\tau$ from the polynomials
on $\auHom(L,S)\otimes\R$ to those on $\auHom(S,L')\otimes\R$, independent of $\vartheta$ and
preserving degrees, a primitive additive character $\psi_\pcm$ of $R=\auO_E/\pcm$ depending only on
$\pcm$ (and the fixed choices of $\alpha$, $\varpi_\pcm$, $a$), and
$c_{\mathrm{line}}=c_{\mathrm{line}}(\vartheta)\in\C^\times$ with
\[
|c_{\mathrm{line}}|^2=C_0\,(\auNm\pcm)^{-n},\qquad C_0>0\text{ independent of }\vartheta,
\]
such that for every $k\in\auU(L)(\auA_f)$, $g\in\auU(S)(\auA)$, $v\in\pcMat_{\mathsf N_2\times
\bullet}(R)$ and every polynomial $B$
\[
\pguK^{\mathrm c}_{v,B}(k,g)=c_{\mathrm{line}}\,\zeta(k)\sum_{v'\in\pcMat_{\bullet\times\mathsf N_2}(R)}
\psi_\pcm(v'v)\,\pguK^{\mathrm r}_{v',\tau B}(g,k\pgupt),
\]
where $\zeta(k)$ is a number of modulus one depending only on $k$, on $\pcm$ and on the characters of
the two line kernels. In particular, for $0\le j<m$,
$\pguK^{\mathrm c}_{v,B}(\pgupt^j,g)=c'_{\mathrm{line}}\sum_{v'}\psi_\pcm(v'v)\,
\pguK^{\mathrm r}_{v',\tau B}(g,\pgupt^{j+1})$ with $c'_{\mathrm{line}}=c_{\mathrm{line}}
\zeta(\pgupt^j)$, $|c'_{\mathrm{line}}|=|c_{\mathrm{line}}|$.
\end{lemma}

\begin{proof}
\emph{The flip.} Let $\mathrm{Fl}\colon\psHS(\auHom(L,S))\to\psHS(\auHom(S,L'))$ be the
isomorphism of Proposition~\ref{pf:flip} ((W6)) for the pair $(L,S)$, $L'=L^-$:
$\mathrm{Fl}\,\psrho(w)=\psrho(w^\dagger)\mathrm{Fl}$, $\theta\circ\mathrm{Fl}=c_\theta\theta$ with
$c_\theta\ne0$, and $\mathrm{Fl}\,\omega(k,g)=\upsilon(\det k)^n\upsilon(\det g)^{-1}\omega(g,k)
\mathrm{Fl}$. Exactly as in the first step of the proof of Lemma~\ref{pgu:line} (this step does not
involve $\alpha$), using $\alpha^{\mathrm r}=(\beta^{\mathrm c})^{-1}\upsilon$ (G2),
\[
\mathrm{Fl}\,\omega^{L,S}_{\alpha_L,\beta^{\mathrm c}}(k,g)=\zeta_1(k)\,
\omega^{S,L'}_{\alpha^{\mathrm r},\beta_{L'}}(g,k)\,\mathrm{Fl},\qquad
\zeta_1(k)=\alpha_L(k)^{-1}\upsilon(k)^n\beta_{L'}(k)^{-1},
\]
so that
\begin{equation}\label{pgx:eq-flip}
\pguK^{\mathrm c}_{v,B}(k,g)=c_\theta^{-1}\,\zeta_1(k)\,\theta\bigl(\omega^{S,L'}_{\alpha^{\mathrm r},
\beta_{L'}}(g,k)\,\mathrm{Fl}[B,\delta_v]\bigr).
\end{equation}
It remains to compute $\mathrm{Fl}[B,\delta_v]$. Write $M=\alpha\auHom(L_L,L_S)$ and
$M'=\alpha\auHom(L_S,L_{L'})$ for the self-dual lattices of the source and the target, and
$M^\dagger=\{x^\dagger:x\in M\}$. Since $(\alpha x)^\dagger=\aubar\alpha x^\dagger$ and
$\auHom(L_L,L_S)^\dagger=\auHom(L_S,L_{L'})$ (the lattices are unimodular),
$M^\dagger=\aubar\alpha\auHom(L_S,L_{L'})=\pgupt^\circ M'$, where $\pgupt^\circ=\aubar\alpha/\alpha$
acts on $\auHom(S,L')$ through the line, $x\mapsto\pgupt^\circ x$, i.e.\ as the element
$(1,\pgupt^\circ)$ of $\auU(S)\times\auU(L')$.

\emph{Local factors.} The local Heisenberg representations are irreducible and $w\mapsto
w^\dagger$ preserves $\beta$ locally, so there are unitary local intertwiners $\mathrm{Fl}_v$, unique
up to factors of modulus one. At every finite place $v$, $\mathrm{Fl}_v\psvac_{M_v}$ spans the line
on which $\psrho(l^\dagger)$, $l\in M_v$, acts by $\epsilon_{M_v}(l)$, i.e.\ (W4) the line of
$\psvac_{M^\dagger_v}$; and by (W4), $\omega_v(1,h)\psvac_{M'_v}$ spans the line of
$\psvac_{hM'_v}$ for $h\in\auU(L')(E^+_v)$ ($\omega_v$ the local Weil representation without
determinant twists). We normalise $\mathrm{Fl}_v$, for finite $v\notin S(\pcm)$, by
\begin{equation}\label{pgx:eq-norm}
\mathrm{Fl}_v\psvac_{M_v}=\omega_v(1,\pgupt_v)\,\psvac_{M'_v}\qquad(\text{unit vectors}),
\end{equation}
which is possible because there $\pgupt_v=\pgupt^\circ_v$ (Proposition~\ref{pgv:G4}), so that
$\pgupt_vM'_v=M^\dagger_v$; at $v\in S(\pcm)$ we only require $\mathrm{Fl}_v$ to be unitary (it is
made explicit below). At almost all $v$ (those outside
$S(\pcm)$ where $\alpha_v$ and $\aubar\alpha_v$ are units) $\pgupt_v$ stabilises $M'_v$, and
$\omega_v(1,\pgupt_v)\psvac_{M'_v}=\psvac_{M'_v}$ (last assertion of (W4)); so the $\mathrm{Fl}_v$ form a
restricted tensor product. At the infinite places $\mathrm{Fl}_\infty$ is the map $F\mapsto F\circ
\dagger^{-1}$ of the Fock models, and $\tau B=B\circ\dagger^{-1}$, exactly as in the proof of
Lemma~\ref{pgu:line} (that paragraph does not involve $\alpha$): $\tau$ is a signed transposition of
the variables, independent of $\vartheta$ and degree-preserving. By the uniqueness of $\mathrm{Fl}$
up to a scalar, $\mathrm{Fl}=c_0\bigotimes_v\mathrm{Fl}_v$; as $\bigotimes_v\mathrm{Fl}_v$ is unitary,
$|c_0|$ is the operator norm of $\mathrm{Fl}$ and does not depend on $\pcm$ or $\vartheta$.

\emph{Places of $S(\pcm)$.} Let $v=w\aubar w\in S(\pcm)$, $w\mid\pcm$, $\varpi$ a
uniformiser at $w$ with $\varpi^{k_w}$ the $w$-component of $\varpi_\pcm$, $q$ the order of the residue
field. By Lemma~\ref{pgx:vac}, in the models (W5) the source space $\psHS_v(\auHom(L,S))$ consists of
the Schwartz--Bruhat functions of the columns $x\in E_w^n$, with vacuum $\mathbf 1_{A}$,
$A=M_w=a_w\auO_w^n$, and the target space $\psHS_v(\auHom(S,L'))$ of the functions of the rows
$\rho\in E_w^n$, with vacuum $\mathbf 1_{A'}$, $A'=M'_w=a_w\auO_w^n$. The map $\dagger$ involves the
conjugation of $E_v$, so it carries the $w$-component of the source onto the $\aubar w$-component of
the target and the $\aubar w$-component of the source onto the $w$-component of the target; hence
$\mathrm{Fl}_v$ carries the translations of the source model into multiplications by characters of
the target model and the multiplications into translations. By the irreducibility of the
Heisenberg representation, an operator with these commutation relations is a multiple of a Fourier
transform: there are $\kappa\in E_w^\times$ and $\gamma\in\C^\times$ with
\[
\mathrm{Fl}_v\Phi(\rho)=\gamma\int_{E_w^n}\Phi(x)\,\psi_w(\kappa\rho x)\,dx ,
\]
$dx$ the measure giving $A$ the volume $1$ and $\psi_w$ the component of the additive character $\psi$.
In detail: the relation $\mathrm{Fl}_v\psrho(l)=\psrho(l^\dagger)\mathrm{Fl}_v$ says that
$\mathrm{Fl}_v$ intertwines translation by $x_0$ in the source with multiplication by the character
$\rho\mapsto[x_0^\dagger,\rho]$ of the target (the pairing $[\,,]$ of the target model between its
$\aubar w$- and $w$-components), and multiplication by a character of the source with a translation of
the target. The pairing $(x_0,\rho)\mapsto[x_0^\dagger,\rho]$ is a non-degenerate bi-additive pairing
$E_w^n\times E_w^n\to\C^\times$ of the form $\psi_w(b(x_0,\rho))$ with $b$ $E_w$-bilinear in the
coordinates (the map $\dagger$ is $E$-semilinear on matrices and the conjugation is absorbed in the
identification of $X_{\aubar w}$ with the matrices over $E_w$). Since $x^\dagger$ is the adjoint of $x$ and the pairing of the target is defined through the
adjoint for the forms of $S$ and $L'=L^-$, one has $(x_0^\dagger)^\dagger=-x_0$, the Gram matrices of
$S$ cancelling; hence $b(x_0,\rho)=\kappa\,\rho x_0$ for some $\kappa\in E_w^\times$. The displayed
transform satisfies the same commutation relations, and by the Stone--von~Neumann theorem (irreducibility
of the Heisenberg representation and Schur's lemma) an intertwiner is unique up to a scalar; hence
$\mathrm{Fl}_v$ is a multiple of it.
The $w$-component of $M^\dagger_v$ is $B'=\aubar\alpha_w\auO_w^n$ (the $w$-component of
$\aubar\alpha$ is the conjugate of $\alpha_{\aubar w}$, and that of $\auHom(L_S,L_{L'})$ is
$\auO_w^n$). As $\mathrm{Fl}_v\mathbf 1_A$ is a multiple of $\psvac_{M^\dagger_v}=\mathbf 1_{B'}$
(Lemma~\ref{pgx:vac} for $M^\dagger_v$) and the transform of $\mathbf 1_A$ is the indicator of the
annihilator $A^\perp=\{\rho:\psi_w(\kappa\rho x)=1\ \forall x\in A\}$, we get
\begin{equation}\label{pgx:eq-dual}
A^\perp=B',\qquad\text{i.e.}\quad\psi_w(\kappa a_w\aubar\alpha_w\,\cdot\,)\text{ is trivial on }\auO_w
\text{ and on no larger fractional ideal}.
\end{equation}
(In particular $\psi_w(\kappa\rho x)=1$ for $x\in A$ and $\rho\in B'$.)
Now $\pgupt_v$ has $w$-component $\pgupt_w=\aubar\alpha_w\alpha_w^{-1}\varpi^{-k_w}$
(Proposition~\ref{pgv:G4}), so $\pgupt_wA'=\aubar\alpha_w\varpi^{-k_w}\auO_w^n=\varpi^{-k_w}B'$. We fix $\gamma$ so that $\mathrm{Fl}_v$ is unitary (the transform maps $\mathbf 1_A$ to
$\mathrm{vol}(A)\mathbf 1_{B'}$, and by Plancherel a positive multiple of it is unitary); a change of
$\gamma$ by a number of modulus one only changes $c_0$ by a number of modulus one.

Let $v\in\pcMat_{\mathsf N_2\times\bullet}(R)$ and let $\tilde v\in\auO_w^n$ lift its $w$-component.
By Definition~\ref{pgx:conv}, $\delta^\natural_v$ at $w$ is the indicator of $a_w\tilde v+a_w
\varpi^{k_w}\auO_w^n$, and
\[
\mathrm{Fl}_v\delta^\natural_v(\rho)=\gamma\int_{a_w\tilde v+a_w\varpi^{k_w}\auO_w^n}\psi_w(\kappa\rho x)
\,dx=\gamma\,q^{-k_wn}\,\psi_w(\kappa a_w\rho\tilde v)\,\mathbf 1_{\varpi^{-k_w}B'}(\rho),
\]
since the annihilator of $a_w\varpi^{k_w}\auO_w^n$ is $\varpi^{-k_w}A^\perp=\varpi^{-k_w}B'$. Substitute
$\rho=\pgupt_w\rho_1$. As $\pgupt_wA'=\varpi^{-k_w}B'$, the indicator becomes $\mathbf 1_{A'}(\rho_1)$.
Write $\rho_1=a_w(\tilde v'+\varpi^{k_w}y)$ with $\tilde v'\in\auO_w^n$ (a row) and $y\in\auO_w^n$. Then
\[
\kappa a_w\pgupt_w\rho_1\tilde v=\kappa a_w^2\pgupt_w\,\tilde v'\tilde v+\kappa a_w\aubar\alpha_w
\,\varepsilon_w\,y\tilde v,\qquad\varepsilon_w=a_w/\alpha_w\in\auO_w^\times ,
\]
because $a_w^2\pgupt_w\varpi^{k_w}=a_w\aubar\alpha_w\varepsilon_w$. By \eqref{pgx:eq-dual} the second
term lies in the kernel of $\psi_w$. Put $\psi^\pcm_w(r)=\psi_w(\kappa a_w^2\pgupt_w\tilde r)$ for
$r\in\auO_w/\varpi^{k_w}$ with a lift $\tilde r$; since $\kappa a_w^2\pgupt_w=\varpi^{-k_w}\cdot\kappa a_w
\aubar\alpha_w\varepsilon_w$ and $\psi_w(\kappa a_w\aubar\alpha_w\varepsilon_w\,\cdot\,)$ is trivial on
$\auO_w$ and on no larger fractional ideal (\eqref{pgx:eq-dual}; $\varepsilon_w$ is a unit),
$\psi^\pcm_w$ is a well defined character of $\auO_w/\varpi^{k_w}$, primitive (non-trivial on the ideal
$\varpi^{k_w-1}\auO_w/\varpi^{k_w}$). Also the value does not depend on the lift $\tilde v$. Hence
\[
\mathrm{Fl}_v\delta^\natural_v(\pgupt_w\rho_1)=\gamma\,q^{-k_wn}\sum_{v'}\psi^\pcm_w(v'_wv_w)\,
\delta^\natural_{v'}(\rho_1),
\]
the sum over the rows $v'$ modulo $\varpi^{k_w}$, with $\delta^\natural_{v'}$ the target datum of
Definition~\ref{pgx:conv} (the indicator of $a_w\tilde v'+a_w\varpi^{k_w}\auO_w^n$). By (W5), for
$h\in\auU(L')(E^+_v)$ with $w$-component $h_w$, $\omega_v(1,h)F(\rho)=\eta\,|h_w|_w^{-n/2}F(h_w^{-1}
\rho)$ with $|\eta|=1$ (the factor makes the operator unitary). Therefore
\[
\mathrm{Fl}_v\delta^\natural_v=\gamma_v\,\omega_v(1,\pgupt_v)\Bigl[\sum_{v'}\psi^\pcm_w(v'_wv_w)\,
\delta^\natural_{v'}\Bigr]
\]
for a constant $\gamma_v$ independent of the rows $v$ and $v'$. Its modulus follows from the norms:
$\|\delta^\natural_v\|^2=q^{-k_wn}$ (measure giving $A$ volume $1$), and the bracket has squared norm
$\sum_{v'}\|\delta^\natural_{v'}\|^2=q^{k_wn}\cdot q^{-k_wn}=1$ (the $\delta^\natural_{v'}$ have disjoint
supports; measure giving $A'$ volume $1$); $\mathrm{Fl}_v$ and $\omega_v(1,\pgupt_v)$ are unitary. So
$|\gamma_v|=q^{-k_wn/2}$. Put $\psi_\pcm=\prod_{w\mid\pcm}\psi^\pcm_w$ on $R=\prod_w\auO_w/\varpi^{k_w}$; it
is primitive, and depends only on $\pcm$ (through $\kappa$, $a_w$, $\alpha$, $\varpi_\pcm$).

\emph{Assembling.} The datum $[B,\delta_v]$ is $B\otimes\bigotimes_{w\mid\pcm}\delta^\natural_{v,w}
\otimes\psvac^\pcm$. Outside $S(\pcm)$, by \eqref{pgx:eq-norm}, $\mathrm{Fl}_v$ maps the unit vector
$\psvac_{M_v}$ to $\omega_v(1,\pgupt_v)\psvac_{M'_v}$. At the infinite places $\mathrm{Fl}_\infty B=
\tau B$. At $S(\pcm)$ we have just computed the image. Hence
\[
\mathrm{Fl}[B,\delta_v]=c\;\omega^{S,L'}(1,\pgupt)\Bigl[\tau B,\sum_{v'\in\pcMat_{\bullet\times
\mathsf N_2}(R)}\psi_\pcm(v'v)\,\delta_{v'}\Bigr],\qquad |c|=|c_0|\prod_{w\mid\pcm}q_w^{-k_wn/2}=
|c_0|\,(\auNm\pcm)^{-n/2},
\]
with $\omega^{S,L'}$ the Weil representation without determinant twists ($\pgupt$ has trivial
archimedean component). Inserting this into \eqref{pgx:eq-flip} and using
$\omega^{S,L'}_{\alpha^{\mathrm r},\beta_{L'}}(g,k)\,\omega^{S,L'}(1,\pgupt)=\beta_{L'}(\pgupt)^{-1}
\omega^{S,L'}_{\alpha^{\mathrm r},\beta_{L'}}(g,k\pgupt)$ (the twisted and untwisted operators differ by
the character $\beta_{L'}(\det\cdot)$ on the target group), we obtain the identity of the statement
with $c_{\mathrm{line}}=c/c_\theta$, $\zeta(k)=\zeta_1(k)\beta_{L'}(\pgupt)^{-1}$, and
$C_0=|c_0/c_\theta|^2$, which is independent of $\vartheta$. The values of the characters at $k$ and at
$\pgupt$ have modulus one, and so does $\zeta(k)$. The last assertion is the case $k=\pgupt^j$.
\end{proof}

\begin{remark}\label{pgx:rem-line}
(a) Lemma~\ref{pgx:line} has the statement of Lemma~\ref{pgu:line}, with two differences: the data are
those of Definition~\ref{pgx:conv} (which agree with those of \S\ref{pc:ss-data} where $\alpha$ is a
unit at the primes of $\pcm$), and the element is \emph{exactly} the element $\pgupt_\pcm$ of
Definition~\ref{pgv:def-period}. In both lemmas
the characters of $k$ enter only through the factor $\zeta(k)$ of modulus one.

(b) The role of the lattices at $v\in S(\pcm)$: when $\alpha$ is a unit at $w$, the vacuum of the
source and of the target is the indicator of $\auO_w^n$, the self-dual measure gives $\auO_w^n$
volume $1$ and the dilation needed is by $\varpi^{k_w}$. In general the vacuum lattices are
$a_w\auO_w^n$ (Lemma~\ref{pgx:vac}), the image of
the source vacuum is the indicator of $\aubar\alpha_w\auO_w^n$, and the dilation needed is by
$\pgupt_w=\aubar\alpha_w/(\alpha_w\varpi^{k_w})$, which carries $a_w\auO_w^n$ to
$\varpi^{-k_w}\aubar\alpha_w\auO_w^n$. With the data of \S\ref{pc:ss-data} (support on the integral
matrices) the computation also closes, but with the dilation by $\varpi^{-k_w}\alpha_w\aubar\alpha_w$
instead of $\pgupt_w$; this element differs from $\pgupt_w$ by $\alpha_w^2$, which is not a unit at
the primes of Proposition~\ref{pgv:alpha}(b), and it would not match the diagonal period of
Definition~\ref{pgv:def-period}. For this reason the data are placed on the lattices $a_w\auO_w^n$
and the period is that of Definition~\ref{pgv:def-period}.
\end{remark}

\subsection{The line reversal in \S\ref{pgu:sec} for general $\alpha$}\label{pgx:ss-later}

The line reversal and its constant enter \S\ref{pgu:sec} at the following places. In each, Lemma~\ref{pgx:line}
is used, with the data of Definition~\ref{pgx:conv}; it contains Lemma~\ref{pgu:line} as the case where
$\alpha$ is a unit at the primes of $\pcm$.
\begin{itemize}
\item \emph{(G4) and the remark after Proposition~\ref{pgv:G4}.} (G4) asks for the element of the line
reversal; Lemma~\ref{pgx:line} produces exactly $\pgupt_\pcm$ of Definition~\ref{pgv:def-period}, so
(G4) holds as stated.
\item \emph{Corollary~\ref{pgu:split}.} Holds with Definition~\ref{pgx:conv}: the factorisation
of the point data and of the norms at $S(\pcm)$ holds because $a_w$ is a scalar (end of
\S\ref{pgx:ss-conv}); the rest of the proof does not involve the primes of $\pcm$.
\item \emph{Lemma~\ref{pgu:see1}.} Holds with Definition~\ref{pgx:conv}. It uses the line kernel only through the function
$M_{x,\mathbf S}(b)=c_1\pguK^{\mathrm c}_{v(\mathbf S),B}(\pgupt^j,b)$, a continuous function
of $b$, left invariant under the rational points and of moderate growth by (R2) of
\S\ref{pcx:ss-conv}, and through Corollary~\ref{pgu:split}; neither the identity of the line
reversal nor $c_{\mathrm{line}}$ occurs. The point $\pgupt^j$ is that of (G4).
\item \emph{Lemma~\ref{pgu:see2}.} Holds, for the same reason, with
$N_{x,\mathbf S}(b)=c_2\pguK^{\mathrm r}_{v'(\mathbf S),B'}(b,\pgupt^{j+1})$ and $j'(a)=a\oplus
\pgupt^{j+1}$.
\item \emph{Proposition~\ref{pgu:turned} and Lemma~\ref{pgu:sumturn}.} They are statements
about functions on matrices over $R$ for an arbitrary additive character $\psi$ of $R$, and do not
involve kernels.
\item \emph{Theorem~\ref{pgu:turndiag}.} Statement and proof hold with $c'_{\mathrm{line}}$ and
$\psi_\pcm$ those of Lemma~\ref{pgx:line}. The proof uses: the identity at $k=\pgupt^j$ (the last
assertion of Lemma~\ref{pgx:line}); the primitivity of
$\psi_\pcm$ is not even needed (Proposition~\ref{pgu:turned} allows any $\psi$); the modulus
$|c'_{\mathrm{line}}|^2=C_0(\auNm\pcm)^{-n}$, $n=r+(2j+1)t$; and that
$\tau$ is a bijection of polynomials independent of $\vartheta$ (in the transfer step, $B$ runs over
all polynomials, so $\tau B$ does too).
\end{itemize}
So Theorem~\ref{pgu:turndiag} $=$ Proposition~\ref{pgr:turndiag} holds without any condition on
$\alpha$ at the primes of $\pcm$.

\subsection{Other uses of the integral matrices at the primes of a modulus}\label{pgx:ss-other}

Definition~\ref{pgx:conv} changes the data only at primes $w\mid\pcm$ with $\alpha_w$ not a unit; by
Proposition~\ref{pgv:alpha} these lie above $q_1$, and they occur only for moduli containing primes of
$Q_\vartheta$ above $q_1$ (the moduli of the grids, and $\pcm_\vartheta$ of \S\ref{pc:ss-data}). We
list the places where the proof of Theorem~C uses the data at such primes.
\begin{enumerate}
\item \S\ref{pc:ss-data}, definition of $\Psi^\natural$ and of the norm: this is
Definition~\ref{pgx:conv}. With data supported on the integral matrices the identity $\mathbf
1^\natural=$ vacuum, used whenever a datum with $\Psi=\mathbf 1$ is identified with the vacuum, would be
false at the primes of Proposition~\ref{pgv:alpha}(b) (Lemma~\ref{pgx:vac}).
\item Lemma~\ref{pc:level}: holds verbatim with Definition~\ref{pgx:conv} (\S\ref{pgx:ss-conv}).
\item Proposition~\ref{pc:rallis}, step ``At $\pcm$'': uses only the support condition for the right
action of $\auGL_N(E_w)$, Lemma~\ref{pc:level} and $\|\Psi^\natural\|^2=\|\Psi\|^2/\#\pcMat(\pcR_\pcm)$;
all three hold with Definition~\ref{pgx:conv} (the lattice $a_w\pcMat(\auO_w)$ is stable under the
right action exactly when $\pcMat(\auO_w)$ is, $a_w$ being a scalar).
\item Lemma~\ref{pc:clean}: uses only that the component at one split place is supported on
injective matrices, which is independent of the lattice.
\item Proposition~\ref{ptv:rev}: its proof at $w\mid\pcm$ describes the data as
the vectors fixed by $\psrho(\Lambda)$, $\Lambda=\varpi\pcMat(\auO_w)\oplus\pcMat(\auO_{\aubar w})$, and
asserts that $\pcMat(\auO_{\aubar w})$ is the dual of $\pcMat(\auO_w)$. By Lemma~\ref{pgx:dict}(b) the dual
of $\pcMat(\auO_w)$ is $\aubar{a_w}a_{\aubar w}\pcMat(\auO_{\aubar w})$, so that sentence is correct
only when $\alpha$ is a unit at $w$ and $\aubar w$. With Definition~\ref{pgx:conv} the proof applies
verbatim after replacing $\pcMat(\auO_w)$ by $M_w=a_w\pcMat(\auO_w)$ and $\pcMat(\auO_{\aubar w})$ by its
dual $M_{\aubar w}$: the argument uses only that these two lattices are mutually dual and stable under
$\auGL_N(\auO_w)\times\auGL_{N'}(\auO_w)$, and the scalar $\zeta_\pcm$ there is constructed from them.
(For the squarefree
moduli of \S\ref{pc:ss-data} this concerns the primes of $Q_\vartheta$ above $q_1$.)
\item The exchange identity (\S\S\ref{pcx:sec}--\ref{pcn:sec}). Its proof works at $w\in S$ in the
coordinates $\tilde x$ of Lemma~\ref{pgx:dict} ($a_w^{-1}x$ on the $w$-component, $a_{\aubar w}^{-1}x$ on
the $\aubar w$-component), in which the vacuum lattice $M_w$ and its dual $M_{\aubar w}$ both become the
integral matrices and $[\,,]$ becomes $\lambda_w(\operatorname{tr}({}^t\tilde y\tilde x))^{\pm1}$
(Lemma~\ref{pgx:dict}(b)). Since $a_w$ is a scalar, this
change of coordinates commutes with the actions $x\mapsto k'^{-1}xk$ of the model (W5) and preserves
the measure normalisations of Definition~\ref{pgx:conv}; it carries the datum $\Psi^\natural$ of
Definition~\ref{pgx:conv} to the function $\Psi(\tilde x\bmod\pcm)$ on the integral matrices, and
matched pairs to matched pairs. Hence Lemma~\ref{pcx:mod} (``$[x,\Phi]=[x,\Phi']_S$'') holds with the
data of Definition~\ref{pgx:conv}, the right side being read in these coordinates; this is stated
after that lemma. In the proof of Lemma~\ref{pct:factor} (Step~4 of the proof of
Theorem~\ref{pct:plane}), the idele $\alpha_0$ by which the Weil data scale $\psi$ is the idele
$\alpha$ of \S\ref{ps:ss-weil} (the scaling that makes the standard lattices self-dual), and in the
coordinates $\tilde x$ its component at $w\in S$ becomes the unit $a_w^{-1}(\alpha_0)_w$; this is the
``unit component of $\alpha_0$'' that appears there, and the additive character to which the Gauss sum
Lemma~\ref{pct:gauss} is applied has conductor $\auO_w$ in these coordinates, by the duality
\eqref{pgx:eq-dual} of the two lattices. The map $w\mapsto w^\dagger$ of \S\ref{pcx:ss-flip} carries
the self-dual lattice $\alpha\auHom$ to $\aubar\alpha\auHom$, which differs from $\alpha\auHom$ at the places
above $q_1$; the comparison is made by the central translation of Lemma~\ref{pgx:dict}(e),(f), as
explained in \S\ref{pcx:ss-flip}. These statements hold with the data of Definition~\ref{pgx:conv}.
\item \S\ref{pgu:ss-line}, Lemma~\ref{pgu:line}: the general case is Lemma~\ref{pgx:line}.
\item (G4), the remark after Proposition~\ref{pgv:G4}, and the discussion after
Proposition~\ref{pgv:alpha}: the data are those of Definition~\ref{pgx:conv}, and the diagonal period is
that of Definition~\ref{pgv:def-period}.
\end{enumerate}
No other step of the proof of Theorem~C uses the integral matrices at a prime of a modulus as a vacuum
lattice.

\subsection{The identification (G1) in the seesaw steps}\label{pgx:ss-G1}

\begin{proposition}\label{pgx:G1}
The proofs of Lemmas~\ref{pgu:see1} and~\ref{pgu:see2}, of Corollary~\ref{pgu:split} and of
Theorem~\ref{pgu:turndiag} use the bijection $\sigma$ of Proposition~\ref{pgv:G1} only through the
induced isometry $S\to S_1$ of spaces with frames: the identification of $\auU(S)$ with
$\auU(S_1)$ (adelic points, rational points, the groups $K$ and the reductions modulo $\pcm$), of
the frames $\mathsf N_2$, and of $\auHom(V',S)$ with $\auHom(V',S_1)$ and $\auHom(S,V)$ with
$\auHom(S_1,V)$ (with lattices, Weil representations, theta functionals and Fock models). No type of
$S$ is identified with a type of $S_1$.
\end{proposition}
\begin{proof}
\emph{Corollary~\ref{pgu:split}.} The isometry $\sigma_*\colon S\to S_1$ is $E$-linear, carries the frame
to the frame and $L_S$ to $L_{S_1}$, and preserves the signs; so it induces isometries
$\auHom(V',S)\cong\auHom(V',S_1)$, $\auHom(S,V)\cong\auHom(S_1,V)$ carrying the self-dual lattices
$\alpha\auHom$ to each other, the complex structures at the infinite places to each other, and
intertwining the Weil representations along $\auU(S)\cong\auU(S_1)$. By (W1) the theta functional is
carried to a non-zero multiple of the theta functional (the isometry is rational), and the multiple is
absorbed in $c_1,c_2$. The datum is carried to the datum: the matrices are written in the frames, which
$\sigma$ identifies, and the polynomials $A$, $A'$ are transported. The characters of (G2) are characters
of $\det$, and $\det(\sigma_*g\sigma_*^{-1})=\det g$. Nothing else is used.

\emph{Lemma~\ref{pgu:see1}.} In the application of Lemma~\ref{pgu:seesaw} with $(S_0,T)=(V',S_1)$ and
$(T',S')=(\widetilde V',S)$: the form $f_{x,\mathbf S}=\varphi_{\mathbf S}$ lives on $S'=S$ and must be
of the conjugate of the target type of $\Phi'=P_{e_2}$, which involves only the type at $S$; the form
$y_{x,\mathbf S}$ lives on $V'$ and has the source type of $\Phi=P_{(V',S_1)}$; the lifts
$\Theta(\Phi,\Psi)y$ are forms on $S_1$ of the type at $S_1$ (Theorem~\ref{pgv:G3}(d)), and they enter
only through the pairing $\langle f,F\rangle_{M,j',r'}=\int_{[\auU(S)]}M(b)f(b)F(j'(b)r')\,db$, in which
$F$ is merely a bounded continuous function on $[\auU(S_1)]$ composed with the group isomorphism
$j'$ of (G1); the moderate-weight condition on $(M,j')$ requires $j'$ to be continuous, compatible with
rational points and to satisfy $\mathrm{ht}(j'(x))\le C\,\mathrm{ht}(x)^C$, and $j'$, induced by a
rational isometry of spaces with frames, preserves the heights (computed in the frames). Condition (b)
is Corollary~\ref{pgu:split}. The rapid decrease of $\varphi_{\mathbf S}$ and the boundedness of the
lifts are properties of the types at $S$ and at $S_1$ separately (Theorem~\ref{pgv:G3}(c),(d)).

\emph{Lemma~\ref{pgu:see2}.} Symmetric: $(S_0,T)=(S,\widetilde V)$, $(T',S')=(S_1,V)$; $y=\varphi_{\mathbf
S}$ lives on $S$ and has the source type of $P_{(S,\widetilde V)}$, the transposed lift
$\Theta^{\mathsf t}(\Phi',\Psi')f$ lives on $S_1$ with the type at $S_1$, and they are paired by
$\int_{[\auU(S)]}N(b)\,y(b)\,F(j(b)r)\,db$ with $j$ the group isomorphism of (G1). Again only the groups
and Corollary~\ref{pgu:split} are used.

\emph{Theorem~\ref{pgu:turndiag}.} The functionals $l_{\mathbf S,v'}$ are defined on functions on
$[\auU(S_1)]$ by integration over $[\auU(S)]$ through the group isomorphism, and the square law
(Proposition~\ref{pgr:square}) is applied in the first grid to forms at $S_1$ of the type at $S_1$; the
forms at $S$ (of the type at $S$) enter only through $\varphi_{\mathbf S}$ inside $l_{\mathbf S,v'}$.
\end{proof}

\noindent Thus the proofs of \S\ref{pgu:sec} use (G1) only in the form described in the remark after
Proposition~\ref{pgv:G1}.



%% ===== proofs/C-grid7.tex =====
% proofs/C-grid7.tex
% Continuation of proofs/C-grid.tex ... C-grid6.tex: the inputs (I) and (P) of Theorem pgw:first.
% Labels start with pgy:. Conventions and macros: those of proofs/C-grid5.tex and C-determinant.tex.
\providecommand{\auGL}{\mathrm{GL}}
\providecommand{\auU}{\mathrm U}
\providecommand{\auO}{\mathcal O}
\providecommand{\auHom}{\operatorname{Hom}}
\providecommand{\aubar}[1]{\overline{#1}}
\providecommand{\pcm}{\mathfrak m}
\providecommand{\pcb}{\mathfrak b}
\providecommand{\pcMat}{\mathrm M}
\providecommand{\pcF}{\mathcal F}
\providecommand{\pcR}{\mathsf R}
\providecommand{\psvac}{\mathbf v}
\providecommand{\pgrY}{\widetilde Y}
\providecommand{\pcnot}[1]{\par\smallskip{\raggedright\noindent\textsc{Not proved here.} [#1]\par}\smallskip}
\providecommand{\pgvwt}{\operatorname{wt}}
\providecommand{\pgyG}{\mathsf G}
\providecommand{\pcRal}{\mathcal R}
\providecommand{\pcmc}{c}
\providecommand{\pcW}{\mathcal W}

\section{Two inputs of the first path: injectivity (I) and change of modulus (P)}\label{pgy:sec}

This section continues \S\S\ref{pgv:sec}, \ref{pgw:sec}, \ref{pgx:sec} and uses their notation.
Theorem~\ref{pgw:first} was proved from three inputs (I), (V), (P) stated in \S\ref{pgw:ss-first}.
\S\ref{pgy:ss-tri} constructs the triangle of the second grid (in the sense of
Definition~\ref{pc:triangle}), verifies (T1)--(T3) for it, and deduces (I) from
Proposition~\ref{pc:tri-inj} (Theorem~\ref{pgy:I}). \S\ref{pgy:ss-P} proves (P) for the data of
Definition~\ref{pgx:conv}, with no assumption on the idele $\alpha$ (Proposition~\ref{pgy:P}).
Input (V) is Theorem~\ref{pgq:V}.

Throughout, $r,k,t=5r,\Delta,\kappa_l,\mu^0_l$ are as in \S\ref{pgv:sec}; $\mu^0_j=\mu^\flat\nu_1^{\,j}$
with $\nu_1=\upsilon^{-\Delta}\chi_M$ as recalled in \S\ref{pc:ss-triangle}, so $\pgvwt\mu^0_j=-\kappa_j$
and the tame parts of $\nu_1$ at the chosen places have order $>2m$; $\pcb=(\pi_M)^b$.
Concerning hypotheses: the archimedean statement called (V) in Lemma~\ref{ptr:const} and
Theorem~\ref{pgv:G3} (Lemma~\ref{par:vac}, proved in Theorem~\ref{pco:V}) is unrelated to the input
(V) of \S\ref{pgw:ss-first}; to avoid confusion we write it $(\mathrm V_\infty)$. Like
Theorem~\ref{pgv:G3}, everything below assumes $(\mathrm V_\infty)$ and (D1)--(D5).

\subsection{The triangle of the second grid and the input (I)}\label{pgy:ss-tri}

\subsubsection*{Spaces, parameters, types}
For a stem coordinate $\varrho$ ($0\le\varrho\le r$) put $a(\varrho)=\max(\varrho+1-k,0)$, as in
\eqref{pgv:eq-sign}, and for integers $i,n\ge0$
\[
E_{i,n}(a)=i+a+\min\bigl((a-i)^+,n\bigr)\qquad\text{(the exponent of \eqref{ptr:eq-E})}.
\]
For $i+n\le m$ let $\pgyG_{i,n}$ be the following space with a frame consisting of a stem of $r+1$
coordinates and $n$ blocks of $t-1$ coordinates, numbered $i+1,\dots,i+n$; its rank is
$N_n=r+1+n(t-1)$, which is even.
\begin{itemize}
\item $\pgyG_{i,0}$ is $\pgrY^0_0=Y_0\oplus E\tilde e$ (Definition~\ref{pgv:def-spaces}) with its
form multiplied by $(-1)^i$;
\item for $n\ge1$, $\pgyG_{i,n}=E^{N_n}$ with the standard lattice and the diagonal form with entry
$(-1)^{E_{i,n}(a(\varrho))}$ at the stem coordinate $\varrho$ and $+1$ on the blocks.
\end{itemize}
The signs of $\pgyG_{i,0}$ are $(-1)^{i+a(\varrho)}=(-1)^{E_{i,0}(a(\varrho))}$, since those of
$\pgrY^0_0$ are $(-1)^{a(\varrho)}$ by \eqref{pgv:eq-sign}. For $i=0$, $E_{0,n}(a)=a+\min(a,n)$, so by
\eqref{pgv:eq-sign} and Definition~\ref{pgv:def-spaces}
\begin{equation}\label{pgy:eq-node}
\pgyG_{0,n}=\pgrY^n_0\quad\text{(the node $(1,0,n)$ of the second grid)},\qquad0\le n\le m .
\end{equation}
At $\tau\in\auPl$ let $\lambda^{(i,n)}_\tau$ be the vector with entries
$\varsigma'_{\tau,\varrho}+\frac12$ at the stem coordinate $\varrho$ ($\varsigma'$ as in
\eqref{pgv:eq-param}) and $-\kappa_{i+l}(\tau)+\frac{t-2}2-\beta$ at the position $0\le\beta<t-1$ of
the block $i+l$. For $i=0$ this is $\lambda_{(1,0,n),\tau}$ of \eqref{pgv:eq-param} ($z=1$, $s=0$). Let
$f_{i,n}$ be the type of $\auU(\pgyG_{i,n})_\infty$ whose Harish-Chandra parameter is
$\lambda^{(i,n)}$ (with the signs of $\pgyG_{i,n}$); for $i=0$, $f_{0,n}=\overline{e_{(1,0,n)}}$, since
$e_{(1,0,n)}$ has parameter $-\lambda_{(1,0,n)}$. For $n=0$ the parameter does not depend on $i$,
and $\pgyG_{i,0},\pgyG_{0,0}$ have the same group; as in Lemma~\ref{ptr:cons}(b) the types
$f_{i,0}$ and $f_{0,0}$ coincide (a change of all signs does not change the compact roots). Put
\[
\pcF_{i,n}=\pcF(\pgyG_{i,n},f_{i,n}),\qquad V=\pcF_{0,0}^{J(\pcb)},\qquad j_i\colon V\hookrightarrow
\pcF_{i,0}=\pcF_{0,0}
\]
the inclusion, $J(\pcb)=J_{\pgrY^0_0}(\pcb)$. So $V$ is the space of forms of type
$\overline{e_{(1,0,0)}}$ on $\pgrY^0_0$ fixed by $J(\pcb)$ of the input (I).

\subsubsection*{Lifts and descents}
For $i+n<m$ let
\[
\alpha_{i,n}=\mu^0_{i+n+1}\upsilon^{N_{n+1}/2},\quad\beta_{i,n}=\mu^0_{i+n+1}\upsilon^{N_n/2},\qquad
a_{i,n}=\Theta_{\alpha_{i,n},\beta_{i,n}}(P_{i,n},\delta_\iota)\colon\pcF_{i,n}\to\pcF_{i,n+1},
\]
with the modulus $\pcb$, the point datum at the inclusion matrix $\iota$ (stem and blocks
$i+1,\dots,i+n$ of $\pgyG_{i,n}$ to the same coordinates of $\pgyG_{i,n+1}$) and the joint harmonic
$P_{i,n}$ of the \emph{ascending pair} $(\pgyG_{i,n},\pgyG_{i,n+1})$, rows with
$\lambda=-\lambda^{(i,n)}$, columns with $\lambda'=\lambda^{(i,n+1)}$. Let $w_n=\#\auGL_{N_n}(\auO_E/\pcb)$.
For $i=0$ these are the characters of the up edge $e^1_n=((1,0,n),(1,0,n+1))$ of
Definition~\ref{pgv:def-edges} (horizontal, index $n$, $\nu=\mu^0_{n+1}$, powers
$\lceil N/2\rceil=N/2$), the joint harmonic is $P_{e^1_n}$ (both are joint harmonics for the same
parameters, unique up to a scalar, Theorem~\ref{pgv:G3}(a)), and so, after normalising
$P_{0,n}=P_{e^1_n}$,
\begin{equation}\label{pgy:eq-chain}
a_{0,n}=a^\flat_n,\qquad X_{0,m}=w_{m-1}\cdots w_0\;a^\flat_{m-1}\circ\dots\circ a^\flat_0=
w_{m-1}\cdots w_0\;a^\flat ,
\end{equation}
with $a^\flat_n$, $a^\flat$ of \S\ref{pgw:ss-first} and $X_{i,n}$ the chains of
Definition~\ref{pc:triangle}.

For $i+n+1\le m$ the \emph{descending pair} is $(\pgyG_{i+1,n},\pgyG_{i,n+1})$, rows with
$\lambda=\lambda^{(i+1,n)}$, columns with $\lambda'=-\lambda^{(i,n+1)}$, and characters
\begin{equation}\label{pgy:eq-desc}
\alpha^\downarrow_{i,n}=(\mu^0_{i+1})^{-1}\upsilon^{N_{n+1}/2+c_{i}},\qquad
\beta^\downarrow_{i,n}=(\mu^0_{i+1})^{-1}\upsilon^{N_n/2+c_i}
\end{equation}
for an integer $c_i$ fixed in Lemma~\ref{pgy:chars}; their weights are $\kappa_{i+1}+N_{n+1}/2+c_i$ and
$\kappa_{i+1}+N_n/2+c_i$. The descent $d_{i,n}\colon\pcF_{i,n+1}\to\pcF_{i+1,n}$ is the transposed lift
along the descending pair with its joint harmonic and a finite datum $\varepsilon_{i,n}$ modulo
$\pcb$: $\varepsilon_{i,0}$ is the datum for which $d_{i,0}=a_{i,0}^*$ (Lemma~\ref{pgy:T2}), and
$\varepsilon_{i,n+1}=\varepsilon_{i,n}^+$ (the extension of Lemma~\ref{pc:next}), as in step (iv) of
\S\ref{pc:ss-exchange}. The triangle is over the one-point set $I=\{1\}$ (nothing depends on a twist).

\subsubsection*{Matching}
\begin{lemma}\label{pgy:match}
Let $c_i=0$ in \eqref{pgy:eq-desc}. At every $\tau$, every ascending pair satisfies the block rule
of excess $t-1$ \eqref{pgv:eq-rule}, with $\tau_0=\kappa_{i+n+1}$, $\tau'_0=-\kappa_{i+n+1}$, $c$ the
identity on the stem and the blocks of $\pgyG_{i,n}$, and $\mathsf B$ the block $i+n+1$; every
descending pair satisfies it with $\tau_0=-\kappa_{i+1}$, $\tau'_0=\kappa_{i+1}$, $c$ the identity on
the stem and on the blocks $i+2,\dots,i+n+1$, and $\mathsf B$ the block $i+1$ of $\pgyG_{i,n+1}$.
Conversely these are the only weights of the descent characters for which the descending pair
satisfies the block rule; hence $c_i=0$ is forced.
\end{lemma}
\begin{proof}
All ranks are even. \emph{Constants} (Lemma~\ref{ptr:const}, valid for any ranks and characters):
ascending, $\tau_0=\frac{N_{n+1}}2-\pgvwt\alpha_{i,n}=\kappa_{i+n+1}$, $\tau'_0=\pgvwt\beta_{i,n}-
\frac{N_n}2=-\kappa_{i+n+1}$; descending, $\tau_0=\frac{N_{n+1}}2-\pgvwt\alpha^\downarrow_{i,n}=
-\kappa_{i+1}-c_i$, $\tau'_0=\pgvwt\beta^\downarrow_{i,n}-\frac{N_n}2=\kappa_{i+1}+c_i$. In both cases
$\tau_0+\tau'_0=0$.

\emph{Old coordinates.} The stem entries of $\lambda^{(i,n)}$ depend neither on $i$ nor on $n$; the
block $i+l$ carries the same entries in every $\pgyG_{i',n'}$ that contains it. Ascending: the
condition $\lambda'_{c(\varrho)}=-\lambda_\varrho$ reads $\lambda^{(i,n+1)}=\lambda^{(i,n)}$ on the common
coordinates: true. Descending: it reads $-\lambda^{(i,n+1)}=-\lambda^{(i+1,n)}$ on the common
coordinates (stem, blocks $i+2,\dots,i+n+1$): true.

\emph{New block.} Ascending: block $i+n+1$ of $\pgyG_{i,n+1}$ carries $-\kappa_{i+n+1}+\frac{(t-1)-1}2-
\beta=\tau'_0+\frac{d-1}2-\beta$, $d=t-1$. Descending: block $i+1$ of $\pgyG_{i,n+1}$ carries
$-\kappa_{i+1}+\frac{t-2}2-\beta$; the entries of $\lambda'$ there are $\kappa_{i+1}-\frac{t-2}2+\beta$,
the string $\{\tau'_0+\frac{d-1}2-\beta'\}$ exactly when $c_i=0$. The signs on blocks are $+1$.
This gives the converse: the rule determines $\tau'_0$ (centre of the new string) and
$\tau_0+\tau'_0$, hence the weights.

\emph{Distances and signs, ascending.} On a common coordinate $\lambda'_{c(\varrho)}-\tau'_0=
\lambda^{(i,n)}_\varrho+\kappa_{i+n+1}$. Blocks $i+l$, $l\le n$: $\kappa_{i+n+1}-\kappa_{i+l}+\frac{t-2}2-
\beta\ge\Delta-\frac t2>\frac d2$, positive. Stem, $\varrho\le k-2$ (then $a(\varrho)=0$):
$\frac k2-\frac12-\varrho+\kappa_{i+n+1}\ge\frac32\Delta-r>0$. Stem, $\varrho\ge k-1$, $a=a(\varrho)=
\varrho+1-k$: $-\kappa-(a+1)\Delta+\frac12+\kappa+(i+n+\frac32)\Delta=(i+n+\frac12-a)\Delta+\frac12$, of
absolute value $\ge\frac\Delta2-\frac12>\frac d2$, positive iff $a\le i+n$. So the rule asks: a stem
coordinate keeps its sign iff $a\le i+n$ (for $\varrho\le k-2$, $a=0\le i+n$). By the definition of
$E_{i,n}$: for $a\le i$, $E_{i,n}(a)=E_{i,n+1}(a)=i+a$; for $i<a\le i+n$ both are $2a$; for $a>i+n$,
$E_{i,n+1}(a)=E_{i,n}(a)+1$. This is the rule.

\emph{Distances and signs, descending} ($c_i=0$). On a common coordinate $\lambda'_{c(\varrho)}-
\tau'_0=-\lambda^{(i,n+1)}_\varrho-\kappa_{i+1}$. Blocks $i+l$, $l\ge2$: $\kappa_{i+l}-\kappa_{i+1}-
\frac{t-2}2+\beta\ge\Delta-\frac t2$, positive. Stem, $\varrho\le k-2$: $-\frac k2+\frac12+\varrho-\kappa_{i+1}
<-\frac32\Delta+r$, negative. Stem, $\varrho\ge k-1$: $\kappa+(a+1)\Delta-\frac12-\kappa-(i+\frac32)\Delta=
(a-i-\frac12)\Delta-\frac12$, of absolute value $>\frac d2$, positive iff $a\ge i+1$. So the rule asks: a
stem coordinate changes sign iff $a\le i$. Compare $E_{i+1,n}$ (rows) with $E_{i,n+1}$ (columns): for
$a\le i$ they are $i+1+a$ and $i+a$; for $a=i+1$ both are $2i+2$; for $i+1<a\le i+n+1$ both are $2a$;
for $a>i+n+1$ both are $i+a+n+1$. This is the rule.
\end{proof}

\begin{lemma}[Gap and regularity]\label{pgy:gap}
Every $\lambda^{(i,n)}_\tau$ consists of $N_n$ distinct numbers in $\frac{N_n-1}2+\Z$, two entries at
coordinates of opposite signs differ by at least $\frac12(\Delta-t)-1$, and $N_n\le2n_m$. Hence the
conclusions of Theorem~\ref{pco:FJ}, Lemma~\ref{pco:decay}, Theorem~\ref{par:A1} and
Corollary~\ref{par:A2} hold for the types $f_{i,n}$ and $\overline{f_{i,n}}$.
\end{lemma}
\begin{proof}
The entries are in $\frac12+\Z$ ($\varsigma'$ is integral, $k$ even; blocks of even length $t-1$), and
$N_n$ is even. Group the entries into: the small stem entries ($\varrho\le k-2$, all with
$a(\varrho)=0$, hence all of the sign $(-1)^{i}$, of absolute value $\le\frac k2$), each large stem entry
$-\kappa-(a+1)\Delta+\frac12$ ($0\le a\le m$), and each block (entries $-\kappa_{i+l}+\frac{t-2}2-\beta$,
$0\le\beta\le t-2$: centre $-\kappa_{i+l}$, half-length $\frac{t-2}2<\frac t2$). Within a group the sign is constant and the entries are
distinct. Between groups: small stem to anything else $\ge\kappa-\frac k2-\frac t2$; two large stem
entries $\ge\Delta$; a large stem entry to a block centre $-\kappa-(i+l+\frac12)\Delta$ at distance
$|a+1-i-l-\frac12|\Delta\ge\frac\Delta2$, so to the block entries $\ge\frac\Delta2-\frac t2-\frac12$; two
blocks $\ge\Delta-t$. All these exceed $\frac12(\Delta-t)-1$ (and $1$), giving distinctness and the gap.
$N_n\le r+1+m(t-1)\le2n_m$ as in Lemma~\ref{pgv:gap}. The consequences follow exactly as in the proof
of Lemma~\ref{pgv:gap} (the cited statements use the parameter only through regularity, gap and rank;
the gap bound $\frac12(\Delta-t)-1$ is the one used there and in Lemma~\ref{ptr:gap}).
\end{proof}

\begin{proposition}[Edges]\label{pgy:edge}
For every ascending and every descending pair (with $c_i=0$), the conclusions (a)--(d) of
Theorem~\ref{ptr:edge} hold, with $e_S=f_{i,n}$, $e_T=f_{i,n+1}$ for an ascending pair and
$e_S=\overline{f_{i+1,n}}$, $e_T=\overline{f_{i,n+1}}$ for a descending one; \textup{(H)} of
\S\ref{ptr:ss-mult} holds for them; and \textup{(A3)} of Lemma~\ref{pc:arch} holds for
$(\pgyG_{i,0},\pgyG_{i,1})$. In particular $a_{i,n}$ maps $\pcF_{i,n}$ to $\pcF_{i,n+1}$ and $d_{i,n}$ maps
$\pcF_{i,n+1}$ to $\pcF_{i+1,n}$, all forms involved are cusp forms, rapidly decreasing, and those of a
given level form finite-dimensional spaces.
\end{proposition}
\begin{proof}
Let $(S,T)$ be one of the pairs, of ranks $n<n'$ with $n'-n=d=t-1$, with its characters
$\alpha,\beta$ and the row and column parameters $(\lambda,\lambda')$ of Lemma~\ref{pgy:match}. Steps
0--5 of the proof of Theorem~\ref{pgv:G3} are written for a pair of standard spaces $(S,T)$ and use the
following properties only, each of which holds here.
\begin{itemize}
\item[(i)] The characters have integral weights, and $\tau_0=\frac{n'}2-\pgvwt_\tau\alpha$,
$\tau'_0=\pgvwt_\tau\beta-\frac n2$ (Step 0: (V) and Lemma~\ref{ptr:const} hold for every pair of
nondegenerate spaces). Here all ranks are even and the constants are computed in the proof of
Lemma~\ref{pgy:match}.
\item[(ii)] $(\lambda,\lambda')$ satisfies the block rule of excess $d$ (Steps 1, 2, 5, through
Lemma~\ref{pgv:rule}). This is Lemma~\ref{pgy:match}, with $d=t-1$.
\item[(iii)] The entries of $\lambda$ and of $\lambda'$ are distinct and lie in $\frac{n-1}2+\Z$,
$\frac{n'-1}2+\Z$, and the conclusions of Theorem~\ref{pco:FJ}, Lemma~\ref{pco:decay},
Theorem~\ref{par:A1} and Corollary~\ref{par:A2} hold for the types of the two nodes and their conjugates
(Steps 1--5). This is Lemma~\ref{pgy:gap}.
\item[(iv)] $n,n'\le2n_m$, so that $2fnn'+2\le x^{25}$ (Step 4). This is Lemma~\ref{pgy:gap}.
\end{itemize}
No other property of the spaces enters; in particular the position of the stem relative to the blocks
is not used. So Steps 0--4 give (a)--(d) of Theorem~\ref{ptr:edge} for $(S,T)$, and Step 5 gives (H).
In Step 3 the nodes without negative sign (if any) have compact archimedean group and are treated as
$j=m$ in Proposition~\ref{pz:types}. The directions in Step 4: for an ascending pair the lift goes from
forms of type $e_S=f_{i,n}$ to forms of type $e_T=f_{i,n+1}$; for a descending pair the transposed lift
goes from forms of type $\overline{e_T}=f_{i,n+1}$ to forms of type $\overline{e_S}=f_{i+1,n}$.

For (A3) of Lemma~\ref{pc:arch} we go through the proofs of Theorem~\ref{ptq:F} and
Proposition~\ref{ptr:Z} for $(S,T)=(\pgyG_{i,0},\pgyG_{i,1})$. Theorem~\ref{ptq:F} is proved for every
polynomial $P\ne0$ none of whose monomials contains two variables of one column from rows of opposite
signs; its proof uses the Fock model of $\auHom(S,T)$ at $\tau$ with the signs $\eta_{rc}=-\varepsilon_r
\varepsilon'_c$, (W1$_\infty$), the strong continuity of $\sigma$ and the fact that $K$ acts on polynomial
vectors by linear substitutions times a character (Step 0), and nothing else about the pair. The joint
harmonic has the property of the monomials by (a). So (F) holds. The proof of
Proposition~\ref{ptr:Z} uses: $c_\lambda\in L^1$ and formula \eqref{pds:eq-FJ} (here by (iii)); that
$\Phi$ is a highest weight vector of weight $\Lambda^K(\lambda)$ for $K$ (Definition~\ref{par:def-harm} and
Lemma~\ref{par:relw}(iii), which does not involve the excess); Schur orthogonality on $K$; and (F).
Hence $Z_\infty=\prod_\tau d_{\lambda_\tau}^{-1}\langle\Phi_\tau,\mathsf P_\tau\Phi_\tau\rangle$ with every
factor positive, which is (A3), as in the proof of Proposition~\ref{ptv:arch}.
\end{proof}

\subsubsection*{Characters of the squares}
\begin{lemma}\label{pgy:chars}
With $c_i=0$, for every square at $(i,n)$, $i+n+2\le m$ \textup{(}$U_1=\pgyG_{i,n+1}$,
$U_2=\pgyG_{i+1,n+1}$, $A=\pgyG_{i,n+2}$, $A'=\pgyG_{i+1,n}$, $M=N_{n+1}$\textup{)}, put
$\nu=\nu_1^{-(n+1)}$ and $\gamma=\nu^{-1}\upsilon^{M}$. Then the four conditions of \textup{(E2)} hold (at $A$, $A'$, $U_1$, $U_2$ respectively, with
the source and target characters as listed in Proposition~\ref{ptv:chars}):
\[
\beta_{i,n+1}\beta^\downarrow_{i,n+1}=\alpha^\downarrow_{i,n}\alpha_{i+1,n}=
\alpha_{i,n+1}\beta^\downarrow_{i,n}=\alpha^\downarrow_{i,n+1}\beta_{i+1,n}=\gamma .
\]
Moreover the target character of the
descent $d_{i,0}=a_{i,0}^*$ obtained from \textup{(W7)} is $\beta^\downarrow_{i,0}$.
\end{lemma}
\begin{proof}
Write $\mu^0_{i+n+2}=\mu^0_{i+1}\nu_1^{\,n+1}$, so $\gamma(\mu^0_{i+n+2})^{-1}=(\mu^0_{i+1})^{-1}
\upsilon^M$. At $A$: $\beta_{i,n+1}\beta^\downarrow_{i,n+1}=\mu^0_{i+n+2}\upsilon^{M/2}\allowbreak (\mu^0_{i+1})^{-1}
\upsilon^{M/2}=\nu_1^{n+1}\upsilon^M=\gamma$. At $A'$: $\alpha^\downarrow_{i,n}\alpha_{i+1,n}=(\mu^0_{i+1})^{-1}\upsilon^{N_{n+1}/2}
\allowbreak \mu^0_{i+n+2}\upsilon^{N_{n+1}/2}=\gamma$. At $U_1$:
$\alpha_{i,n+1}\beta^\downarrow_{i,n}=\mu^0_{i+n+2}\upsilon^{N_{n+2}/2}\allowbreak (\mu^0_{i+1})^{-1}\upsilon^{N_n/2}=
\gamma$, since $N_n+N_{n+2}=2M$. At $U_2$: the source character $\alpha^\downarrow_{i,n+1}=
(\mu^0_{i+1})^{-1}\upsilon^{N_{n+2}/2}$ of $d_{i,n+1}$ times the target character
$\beta_{i+1,n}=\mu^0_{i+n+2}\upsilon^{N_n/2}$ of $a_{i+1,n}$ is again $\gamma$. The last assertion is the argument of Proposition~\ref{ptv:chars}(b): by (W7) the target character of
the conjugate kernel is $(\mu^0_{i+1})^{-1}$ times an integral power of $\upsilon$ (the target
character $\beta_{i,0}=\mu^0_{i+1}\upsilon^{N_0/2}$ conjugated, times powers of $\upsilon$), its weight is
forced by the block rule (Lemma~\ref{pgy:match}, converse part, for $n=0$, which applies because the
conjugate kernel is a kernel of the descending pair with negated parameters, as in
Remark~\ref{ptr:rem-desc}); $\beta^\downarrow_{i,0}$ has the same form and the same weight, and
$\upsilon$ has weight $1$ at every $\tau$, so the two are equal.
\end{proof}

\subsubsection*{(T1): the squares}
\begin{proposition}[(T1)]\label{pgy:T1}
For $i+n+2\le m$ there is $\sigma_{i,n}\in\C$ with $w_n\,a_{i+1,n}\circ d_{i,n}=\sigma_{i,n}\,w_{n+2}\,
d_{i,n+1}\circ a_{i,n+1}$ on $\pcF_{i,n+1}$.
\end{proposition}
\begin{proof}
We check that the square at $(i,n)$, with the one-member family $(\pcb,\nu,\chi,\text{the
characters above})$, $\chi$ the tame character of $\nu$, is an edge square of excess $e=t-1$ and
weight $\mu_\tau=(n+1)\Delta$ (Definition~\ref{ptr:def-square}). \emph{Ranks}:
$\dim A=N_{n+2}=N_n+2(t-1)$ and $N_{n+1}+N_{n+1}=N_{n+2}+N_n=N=2M$. \emph{Signs}: all negative signs are
on stems, and comparing the exponents $E$ coordinate by coordinate as in \S\ref{pc:ss-exchange}
(for $a\le i$ both sides have the exponents $\{i+a,i+a+1\}$; for $i<a\le i+n+1$ all four are $2a$;
for $a\ge i+n+2$ both sides have $\{i+a+n+1,i+a+n+2\}$) gives $q=q_A+q_{A'}$. As each stem has $r+1$
coordinates, $q\le2r+2$, so $2q+2\le4r+6\le2(r+1)+2(t-1)\le N$. And $N\le4n_m\le52r^{11/10}<2\Delta\le
2\mu_\tau$ (proof of Lemma~\ref{pgv:gap}). \emph{Data}: $\pcb\ne(1)$ is a modulus; $\nu=\nu_1^{-(n+1)}$
has weight $(n+1)\Delta$, is unramified outside the chosen places and tamely ramified there with tame
part $\chi_M^{-(n+1)}$, non-trivial since $n+1\le m$ and the tame parts of $\chi_M$ have order $>2m$.
\emph{(E1)} is Proposition~\ref{pgy:edge} for the four pairs $(U_1,A)$, $(A',U_2)$ (ascending at
$(i,n+1)$, $(i+1,n)$) and $(U_2,A)$, $(A',U_1)$ (descending at $(i,n+1)$, $(i,n)$), whose polynomials
and types do not depend on anything else. \emph{(E2)} is Lemma~\ref{pgy:chars} (the factor
$\upsilon^{N/2}=\upsilon^M$ is unramified, so the tame part of $\gamma^{-1}$ at $\pcb$ is $\chi$).
Theorem~\ref{ptr:square} (which rests on Theorem~\ref{pcx:X} and Proposition~\ref{ptr:conv}; the
latter uses only (E1) and full rank) gives the square law of \S\ref{pc:ss-exchange} for the functions
$L,L'$ of the square with constant $\kappa/c$. Step (iv) of \S\ref{pc:ss-exchange}, i.e.\
Lemma~\ref{pc:next} with $\mathsf N_0$ the frame of $A'$ (so $\#\auGL_{\mathsf N_0}(\auO_E/\pcb)=w_n$) and
$\mathsf A$ the frame of $A$ ($\#\auGL_{\mathsf A}=w_{n+2}$), with $p=\varepsilon_{i,n}$ and
$p^+=\varepsilon_{i,n+1}$, turns it into the identity, by pairing both sides with arbitrary forms of
$\pcF_{i+1,n}$ (the pairings separate, the spaces of a given level being finite-dimensional and the
Petersson product non-degenerate). So $\sigma_{i,n}=\kappa/c$.
\end{proof}

\subsubsection*{(T2), (T3): the first lifts}
\begin{lemma}[(T2)]\label{pgy:T2}
There is a finite datum $\varepsilon_{i,0}$ modulo $\pcb$, with $|\varepsilon_{i,0}|=|\delta_{\iota}|$ up to a
constant of modulus one, such that the transposed lift $d_{i,0}$ along $(\pgyG_{i+1,0},\pgyG_{i,1})$ with it
is the adjoint of $a_{i,0}$ under $\pcF_{i+1,0}=\pcF_{i,0}$. In particular
$\langle z,a_{i,0}j_iy\rangle=\langle d_{i,0}z,j_{i+1}y\rangle$ for $y\in V$, $z\in\pcF_{i,1}$.
\end{lemma}
\begin{proof}
This is item (ii) of Proposition~\ref{ptv:tri}, whose proof applies verbatim: Lemma~\ref{ptv:adj}
(its hypotheses: $\delta_\iota$ is supported on injective matrices, and the edge property is
Proposition~\ref{pgy:edge}) gives $a_{i,0}^*z(g)=\int\aubar{\theta(g,h)}z(h)\,dh$; by (W7) the conjugate
kernel is the kernel for $(\pgyG_{i,0}^-,\pgyG_{i,1})=(\pgyG_{i+1,0},\pgyG_{i,1})$ of the datum
$\mathsf C[P_{i,0},\delta_\iota]=[\aubar{P_{i,0}},c\,\delta_{\iota}]$, $|c|=1$ (at the primes of $\pcb$
the idele $\alpha$ is a unit by Proposition~\ref{pgv:alpha}(a), so $a_w=1$ in
Definition~\ref{pgx:conv}, and in the model (W5) $\mathsf C$ is complex conjugation of functions up to
a constant of modulus one, since $[x,a']^-=\aubar{[x,a']}$, as in the proof of
Proposition~\ref{ptv:tri}; outside $S(\pcb)$ it maps the vacuum to a multiple of the vacuum by (W4));
$\aubar{P_{i,0}}$ is the joint harmonic of the descending pair for
$n=0$ (Remark~\ref{ptr:rem-desc}), and its characters are those of Lemma~\ref{pgy:chars}.
\end{proof}

\begin{lemma}[The series in rank $r+1$]\label{pgy:series}
Let $S$ be a space of rank $n=r+1$, $M'=r+t$, $\pcm$ a modulus and $\alpha$ a character unramified
outside $\pcm$. Then $\sum_{\gamma\in G^\pcm/K^\pcm,\ \gamma\notin K^\pcm}|\pcmc_{M',\alpha}(\gamma)|<\frac12$,
and the conclusions of Lemma~\ref{pc:series}\textup{(c)} hold for $\pcRal_{M',\alpha}$ on $S$.
\end{lemma}
\begin{proof}
Here $h:=M'/2=3r=3n-3$, whereas Lemma~\ref{pc:series}(b) assumed $M'\ge6n$; we redo the count of its
proof with the same notation ($q=q_w\ge5$). \emph{Split $v$}: as there,
$1+\varepsilon_v\le\prod_{i<n}(1-q^{i-h})^{-2}$; here $i-h\le-2n+2$, $\sum_{i<n}q^{i-h}\le\frac54q^{-2n+2}\le
\frac15$, so $1+\varepsilon_v\le\exp(5q^{-2n+2})$ and $\varepsilon_v\le10q^{-2n+2}\le6250\cdot25^{-n}q^{-2}\le
2(4/5)^nq^{-2}$, as $3125\le20^n$ for $n\ge3$. \emph{Non-split $v$}: as there, $\varepsilon_v\le\sum_{k\ge1}
s_k^2q^{-kh}$ with $s_k\le2^{k+n-1}q^{k(n-1)}$, so $\varepsilon_v\le4^{n-1}\sum_{k\ge1}(4q^{2n-2-h})^k=
4^{n-1}\sum_k(4q^{-n+1})^k\le2\cdot4^nq^{-n+1}$ (as $4q^{-n+1}\le\frac12$), and $2\cdot4^nq^{-n+1}=
2q^{-2}(4/5)^n(5/q)^nq^3\le250(4/5)^nq^{-2}$ (as $(5/q)^nq^3\le125$ for $n\ge3$, $q\ge5$).
\emph{All places}: as there, $\sum_v\varepsilon_v\le250\cdot2f(4/5)^n\cdot\frac12=250f(4/5)^n$. With
$n>r\ge x^5$, $f\le x^{12}$ (the sizes used in the proof of Lemma~\ref{pc:series}(c)) and
$x\ge128^{100}$, $250f(4/5)^n<\frac13$, so the sum is at most $e^{1/3}-1<\frac12$. The proof of
Lemma~\ref{pc:series}(c) then applies unchanged.
\end{proof}

\begin{proposition}[First lifts]\label{pgy:first}
For $0\le i<m$ there is $C\ne0$ with
$\langle a_{i,0}y',a_{i,0}y\rangle=C\langle y',\pcRal\,y\rangle$ for all $y,y'\in V$, where
$\pcRal=\pcRal_{r+t,\alpha_{i,0}}$ with the modulus $\pcb$; and $a_{i,0}$ is injective on $V$. Moreover for
$u\in J(\pcb)$ there is $g\in K_{\pgyG_{i,1}}$ with $a_{i,0}\rho(u)=\rho(g)a_{i,0}$.
\end{proposition}
\begin{proof}
Proposition~\ref{pc:rallis} is stated for ranks $(r,r+t)$; its proof uses the ranks only through
(SW) of Lemma~\ref{pc:SW} and through Lemma~\ref{pc:series}. For the ranks $(r+1,r+t)$ of
$(\pgyG_{i,0},\pgyG_{i,1})$ we use Lemma~\ref{pgy:series} in place of Lemma~\ref{pc:series}, and (SW) for
the doubled space $\pcW=\pgyG_{i,0}\oplus\pgyG_{i,0}^-$ of rank $2(r+1)$ and $T=\pgyG_{i,1}$ of rank
$M'=6r$: as recalled in the comments to Lemma~\ref{pc:SW}, Weil's theorem needs the convergence of the
Eisenstein series, $M'>2(r+1)$, and of the theta integral, $M'$ minus the Witt index of $T$ greater
than $r+1$; the Witt index is at most $r+1$ (all negative signs are on the stem), and $6r-(r+1)>r+1$.
(These are condition (B) and Proposition~8 of~\cite{WeilSiegel} in the hermitian case, the hypotheses of \cite[Th\'eor\`eme~5]{WeilSiegel}; here $r\ge1$, so $6r>2(r+1)$.)
Lemma~\ref{pc:clean} is rank-free. (A1), (A2), (A3) hold by Proposition~\ref{pgy:edge}. The point
datum $\delta_\iota$ on $\pcMat_{N_1\times N_0}(\auO_E/\pcb)$ is sharp for $J=J(\pcb)$ (with $A$ the projection
to the first $N_0$ rows, $A\iota=1\in\bar J=\{1\}$, and $\alpha_{i,0}$ trivial on $\det J(\pcb)$, its tame
parts living at the chosen places, i.e.\ at the primes of $\pcb$, and being characters of $\auO_E/\pcb$
that are trivial on $1$). So Proposition~\ref{pc:rallis}, with $\#\pcMat_{(r+t)\times(r+1)}$ in place of
$\#\pcMat_{(r+t)\times r}$, gives the formula and the injectivity on $\pcF(\pgyG_{i,0},f_{i,0})^{J(\pcb)}=V$.
For the last assertion choose $g$ with reduction $\operatorname{diag}(\bar u,1)=1$ modulo $\pcb$, e.g.\
$g=1$ at the primes of $\pcb$, and argue as in the proof of Proposition~\ref{pc:first}: here $\bar u=1$, so
by Lemma~\ref{pc:level} the pair $(u^{-1},1)$ fixes $[P,\delta_\iota]$ and $a_{i,0}\rho(u)=a_{i,0}$, i.e.\
the claim with $g=1$.
\end{proof}

\begin{lemma}[(T3)]\label{pgy:T3}
For $i<m$, $d_{i,0}a_{i,0}j_i(V)\subseteq j_{i+1}(V)$, and $d_{i,0}a_{i,0}j_i$ is injective.
\end{lemma}
\begin{proof}
As in the proof of Corollary~\ref{pc:istriangle}. Let $x\in V$ and $z=d_{i,0}a_{i,0}x\in\pcF_{i,0}$. By
Lemma~\ref{pgy:T2}, $\langle z,f\rangle=\langle a_{i,0}x,a_{i,0}f\rangle$ for all $f\in\pcF_{i,0}$. For
$u\in J(\pcb)$, by Proposition~\ref{pgy:first}, $a_{i,0}\rho(u^{-1})=a_{i,0}$, hence
$\langle\rho(u)z,f\rangle=\langle z,\rho(u^{-1})f\rangle=\langle a_{i,0}x,a_{i,0}f\rangle=\langle z,f\rangle$; so
$z$ is $J(\pcb)$-invariant, $z\in V$ (it lies in $\pcF_{i+1,0}=\pcF_{0,0}$ by Proposition~\ref{pgy:edge}).
If $z=0$ then $\langle a_{i,0}x,a_{i,0}x\rangle=0$ and $x=0$ by Proposition~\ref{pgy:first}.
\end{proof}

\begin{theorem}[Input (I)]\label{pgy:I}
The data $(V,\pcF_{i,n},j_i,w_n,a_{i,n},d_{i,n})$ form a triangle of size $m$ over $\{1\}$
(Definition~\ref{pc:triangle}). Consequently $a^\flat=a^\flat_{m-1}\circ\dots\circ a^\flat_0$ is injective
on the forms of type $\overline{e_{(1,0,0)}}$ on $\pgrY^0_0$ fixed by $J_{\pgrY^0_0}(\pcb)$: input (I) of
\S\ref{pgw:ss-first} holds.
\end{theorem}
\begin{proof}
The $\pcF_{i,n}$ are inner product spaces (Petersson products of bounded forms on finite-volume
quotients), $j_i$ is an isometry, $w_n>0$. (T1) is Proposition~\ref{pgy:T1}, (T2) Lemma~\ref{pgy:T2},
(T3) Lemma~\ref{pgy:T3}. By Proposition~\ref{pc:tri-inj} with $(i,n)=(0,m)$, $X_{0,m}\circ j_0$ is
injective on $V$, and $X_{0,m}=w_{m-1}\cdots w_0\,a^\flat$ by \eqref{pgy:eq-node} and \eqref{pgy:eq-chain}.
\end{proof}

\begin{remark}\label{pgy:rem-I}
Only the existence of the constants $\sigma_{i,n}$ is used, not their size or non-vanishing (the
latter follows from Proposition~\ref{pc:tri-inj} if $V\ne0$). The triangle is untwisted, so no
modulus of a twist and no $L$-value estimate enters.
\end{remark}

\subsection{The change of modulus (P)}\label{pgy:ss-P}

We recall the data of \S\ref{pc:ss-data} (Definition~\ref{pgx:conv}) at the primes of a modulus: generators $a_w$ of
$\alpha_w\auO_w$, fixed once for every prime $w$ (independently of the modulus), $a_w=1$ where $\alpha_w$
is a unit, and the vacuum at $v\in S(\pcm)$ normalised as the function $\mathbf 1_{M_w}$,
$M_w=a_w\pcMat(\auO_w)$, of norm $1$ for the measure of Definition~\ref{pgx:conv} (a vacuum by
Lemma~\ref{pgx:vac}). Outside $S(\pcm)$ the vacuum is that of $\alpha\auHom(L_S,L_T)$ as in
\S\ref{pc:ss-data}, with the same normalisation at split places.

\begin{proposition}[Input (P)]\label{pgy:P}
Let $\pcm'\subseteq\pcm$ be moduli, $\mathrm{red}\colon\auO_E/\pcm'\to\auO_E/\pcm$ the reduction, $S,T$
spaces, $P$ a Fock vector and $\Psi$ a function on $\pcMat_{N'\times N}(\auO_E/\pcm)$. Then the datum
$[P,\Psi\circ\mathrm{red}]$ for the modulus $\pcm'$ equals the datum $[P,\Psi]$ for the modulus $\pcm$.
In particular, for every edge of the second column and every $\Psi$ on $\pcMat(\auO_E/\pcb)$, the data
$[P,\Psi]$ for $\pcb$ and $[P,\Psi\circ\mathrm{red}]$ for $\pcm(\vartheta)=\pcb\prod_{w\in Q_\vartheta}w$ coincide, with
no assumption on $\alpha$; so the kernels, lifts and transposed lifts coincide.
\end{proposition}
\begin{proof}
Both data are $P$ at the infinite places and the vacuum of $\alpha\auHom(L_S,L_T)$ at the finite places
outside $S(\pcm')$. At $v\in S(\pcm)$, $v=w\aubar w$, $w\mid\pcm$: the $w$-component of
$(\Psi\circ\mathrm{red})^\natural$ is $x\mapsto\Psi(\mathrm{red}(a^{-1}x\bmod\pcm'))=\Psi(a^{-1}x\bmod\pcm)$ on
$a\pcMat(\widehat\auO)$, which is the $w$-component of $\Psi^\natural$; the measures used for the norms are
the same. (The decomposition into components is the factorisation of $\Psi^\natural$ over the primes of
$\pcm$ given by the Chinese remainder theorem, $\auO_E/\pcm=\prod_{w\mid\pcm}\auO_E/w^{k_w}$; for
non-factorisable $\Psi$ the identity is the same computation on $\prod_{w\mid\pcm}\pcMat(E_w)$.) At
$v\in S(\pcm')\smallsetminus S(\pcm)$, $w\mid\pcm'$, $w\nmid\pcm$: the value of $\Psi\circ\mathrm{red}$
does not depend on the $w$-component, so $(\Psi\circ\mathrm{red})^\natural$ is, in the variable $x_w$, the
indicator of $a_w\pcMat(\auO_w)=M_w$, which is the normalised vacuum at $v$ by Lemma~\ref{pgx:vac} and our
normalisation; this is the component of $[P,\Psi]$ at $v$, where the modulus $\pcm$ puts the vacuum. For
the second-column edges take $\pcm'=\pcm(\vartheta)$, $\pcm=\pcb$ (moduli: \S\ref{pgv:sec}).
\end{proof}

\begin{remark}\label{pgy:rem-P}
\S\ref{pc:ss-data} uses the normalised coordinates $a^{-1}x$ of Definition~\ref{pgx:conv}. Had
$\Psi^\natural$ been supported on the integral matrices in the frame coordinates, (P) would still hold at
the primes where $\alpha_w$ is a unit, but it would fail at a prime $w\mid Q_\vartheta$ where $\alpha_w$ is
not a unit, because $\mathbf 1_{\pcMat(\auO_w)}$ is not the vacuum there (Lemma~\ref{pgx:vac}). The proof
of Theorem~\ref{pgw:first} uses (P) only to write $T^\flat_b=\sum_{X\equiv\iota\,(\pcb)}T_X$, and, for the
first column, to identify $\delta_\iota$ modulo $\pcm_0$ with $\mathbf 1[X\equiv\iota\bmod\pcm_0]$ modulo
$\pcm$; both are cases of Proposition~\ref{pgy:P}.
\end{remark}



%% ===== proofs/C-grid8.tex =====
% proofs/C-grid8.tex
% Continuation of proofs/C-grid.tex (pgr:), C-grid2.tex (pgs:), C-grid4.tex (pgv:), C-grid5.tex
% (pgw:), C-grid6.tex (pgx:): Proposition pgr:den (denominators and conjugates on the last path).
% Labels start with pgz:. Conventions and macros: those of proofs/C-grid5.tex.
\providecommand{\auI}{\mathcal I}
\providecommand{\auGL}{\mathrm{GL}}
\providecommand{\auU}{\mathrm U}
\providecommand{\auO}{\mathcal O}
\providecommand{\auA}{\mathbb A}
\providecommand{\auNm}{\operatorname{Nm}}
\providecommand{\auHom}{\operatorname{Hom}}
\providecommand{\aubar}[1]{\overline{#1}}
\providecommand{\pcm}{\mathfrak m}
\providecommand{\pcb}{\mathfrak b}
\providecommand{\pcMat}{\mathrm M}
\providecommand{\pcR}{\mathsf R}
\providecommand{\pcF}{\mathcal F}
\providecommand{\auPl}{\Phi}
\providecommand{\psvac}{\mathbf v}
\providecommand{\psrho}{\varrho}
\providecommand{\pgrY}{\widetilde Y}
\providecommand{\pgrlast}{\Gamma_{\mathrm{last}}}
\providecommand{\pgrE}{\operatorname{Ed}}
\providecommand{\pcnot}[1]{\par\smallskip{\raggedright\noindent\textsc{Not proved here.} [#1]\par}\smallskip}
\providecommand{\pgupt}{\mathsf u}
\providecommand{\pgwK}{\mathcal K}
\providecommand{\pgwT}{\mathcal T}
\providecommand{\pgwint}[1]{\int^{#1}}
\providecommand{\pgzd}{\mathsf d}
\providecommand{\pgzM}{\mathsf M}

\section{Denominators and conjugates on the last path: Proposition~\ref{pgr:den}}\label{pgz:sec}

This section proves part (a) of Proposition~\ref{pgr:den} in full, for the sum $\Lambda$ of
\eqref{pgw:eq-Lambda} with the integral choices of the branching form and of the kernels fixed in
\S\ref{pgz:ss-choice}, and reduces part (b) to three size bounds stated in \S\ref{pgz:ss-b}. The
notation is that of \S\S\ref{pgv:sec}, \ref{pgw:sec}, \ref{pgx:sec}. Throughout, $\vartheta$ is an
$m$-fold twist, $\pcm=\pcm(\vartheta)=\pcb\prod_{w\in Q_\vartheta}w$ with $\pcb=(\pi_M)^b$, and the
data at the primes of $\pcm$ are those of Definition~\ref{pgx:conv}. For $z\in\{0,1\}$ and $0\le a\le
m$ let $p^z_a=(z,a,m)$ be the nodes of the half $h_z$ of $\pgrlast$, $G^z_a=G_{p^z_a}$,
$\Gamma^z_a=\Gamma_{p^z_a}$, $J^z_a=J_{p^z_a}$ its level (\S\ref{pgw:ss-last}: $J^0_0=J(\pcm)$, all
others $J(\pcb)$), $N^z_a$ its rank, and $e^z_a=(p^z_{a-1},p^z_a)$ for $1\le a\le m$. Put
$Y=Y^m_m$, so $p^0_m$ is the node $Y$ and $p^1_m$ the node $\pgrY^m_m=Y\oplus E\tilde e_m$.

The structure of the argument is the following. (i) All levels along $\pgrlast$ are free, and the
invariance of a path datum makes the integrand of $\Lambda$ invariant under these levels in every
variable separately; so $\Lambda$ is a plain sum, over the tuples of classes of $2m+1$ levels and
over the edge tuples, of $\Psi_\vartheta(S)$ times branching pairings of vectors reached by $m$
lifts from each end (Proposition~\ref{pgz:expand}). \emph{No index $[J:J']$ and no order of a
stabiliser occurs}, so these contribute nothing to $N_0$. (ii) One lift with an edge datum, multiplied
by the integer $\pgzd$ of \eqref{pgz:eq-d}, does not increase gauges at $\iota$ when
$\ell(\iota)\nmid\auNm Q_\vartheta$ (Lemma~\ref{pgz:kernel}). (iii) The branching pairing of two
vectors of gauge at most $1$ at $h$ and at $h\oplus\pgupt^m$ is integral after multiplication by the
integer $\pgzM$ of \eqref{pgz:eq-M} (Lemma~\ref{pgz:pair}). Then $N_0=\pgzM\pgzd^{2m}$.

\subsection{Integral structures and the integral choices}\label{pgz:ss-choice}

\subsubsection*{The representations along the last path}
Let $W^z_a$, $\mu^z_a$, $\pgwK^z_a$ be as in \S\ref{pgw:ss-last} (the paragraph before
\eqref{pgw:eq-Lambda}). If some $W^z_a$ is $0$, then $\pgwT^z_S=0$ for all $S$ and $\Lambda=0$, and
there is nothing to prove; we assume all $W^z_a\ne0$.

\emph{$W^z_a$ is the representation of the type at $p^z_a$}, more precisely $W^z_a\cong\pi_{p^z_a}^*$ with
$\pi_p$ the representation of the type $e_p$, as $W^z_0=(\pi^z_0)^*$ in Lemma~\ref{pgw:vec}. By induction on
$a$; let $(p,q)=e^z_a$, a right edge (the first index grows along the row), and $W^z_{a-1}\cong\pi_p^*$. In
the convention fixed before Proposition~\ref{pgv:match}, $e_p$ is the type whose parameter is the conjugate
of the row parameter $\lambda$, and $e_q$ the type of the column parameter $\lambda'$; so $\pi_p^*\cong
\pi_\lambda$ and $\pi_q\cong\pi_{\lambda'}$. By Theorem~\ref{pgv:G3}(a),(b) the joint harmonic $P=P_{pq}\ne0$
satisfies $\mathsf PP=\mathsf P'P=P$ with $\mathsf P=\omega_\infty(e_\lambda,1)$, $\mathsf P'=\omega_\infty(1,
e_{\lambda'})$, $e_\lambda=d_\lambda\overline{c_\lambda}$; by Schur orthogonality $\mathsf P$ maps into the
$\pi_\lambda$-isotypic part for $G_{p,\infty}$ and $\mathsf P'$ into the $\pi_{\lambda'}$-isotypic part for
$G_{q,\infty}$. In the decomposition of Lemma~\ref{pgw:howe}, $P$ therefore lies in $\pgwP_{\nu_1}$ for the
unique $\nu_1$ with $A_{\nu_1}\cong\pi_\lambda$ (the $A_\nu$ are pairwise inequivalent), and
$B_{\nu_1}\cong\pi_{\lambda'}$. As $W=W^z_{a-1}\cong\pi_p^*\cong\pi_\lambda=A_{\nu_1}$, the index $\nu_0$ of
the proof of Lemma~\ref{pgw:one} is $\nu_1$, and $W^z_a=B_{\nu_1}^*\cong\pi_{\lambda'}^*\cong\pi_q^*$. The bounds
used below (the extreme entries of the highest weight, Lemmas~\ref{pgz:deg}, \ref{pgz:char}) are the
same for $\pi_q$ and $\pi_q^*$. We realise $W^z_a$ as $W_{\lambda^z_a}$ with the basis $(e_T)$ and the lattice
$W_{\lambda^z_a}(\Z)=\bigoplus\Z e_T$ of \S\ref{pgs:ss-weyl}, $\lambda^z_a$ the corresponding weight of
$\prod_\tau\auGL_{N^z_a}$. Its matrix coefficients in the basis $(e_T)$ are polynomials with integer
coefficients in the entries of $g$ and in $\det g^{-1}$. For a prime $\ell$ and $j\colon\C\to\C_\ell$,
the \emph{gauge} at $(j,g)$, $g\in G^z_{a,f}$, of $\xi\in W^z_a$ is the largest absolute value of a
coordinate of $g_\ell^{-1}\cdot j(\xi)$, the action of $g_\ell$ being through the embeddings of
\S\ref{ss:au-spaces} (as in the definition of integral forms there). A $W^z_a$-valued form is
\emph{integral at $\iota$} if each value $\psi(g)$ has gauge at most $1$ at $(\iota,g)$; this is the
hypothesis of Proposition~\ref{pgr:den}(a) on $y'=f_{\mu^0_0,\psi_{y'}}$ and
$\varphi=f_{\mu^1_0,\psi_\varphi}$.

Two facts on the gauge are used. (G-a) For $k\in G^z_{a,f}$ integral with integral inverse at the
places above $\ell$, $k_\ell$ acts on $W_{\lambda^z_a}(\auO_{\C_\ell})$ by a matrix invertible over
$\auO_{\C_\ell}$ (its entries and those of its inverse are integral polynomials in entries that are
integral). (G-b) For a rational point $\gamma\in\Gamma^z_a$, the archimedean action of $\gamma$ on
$W^z_a$ followed by $j$ equals the action of $\gamma_\ell$ on $W^z_a\otimes_j\C_\ell$: both are given
by the same integral polynomials evaluated at the images of the entries of $\gamma$ under the
embeddings concerned (\S\ref{ss:au-spaces}). Hence the gauge of $\gamma\xi$ at $(j,\gamma g)$ equals
the gauge of $\xi$ at $(j,g)$.

\begin{definition}[Integral choices]\label{pgz:def-choice}
(a) \emph{Kernels.} For the edge $e^z_a=(p,q)$ let $S_\nu(z)$ be the Schur matrix of the diagram
$\nu=\nu_0$ of the edge (Lemma~\ref{pgw:one}, Remark~\ref{pgw:rem-jets}) in the bases $(e_T)$, an
integer matrix of polynomials in the variables $z_{\tau,c,a'}$, homogeneous of total degree
$e(\nu)=\sum_\tau|\nu_\tau|$, and let $\Theta_f$ be the function of \eqref{pk:eq-theta} attached to a
finite datum $f$ on $\auHom(p,q)$, with the scale $s_0$ of \eqref{pk:eq-s0}. We take
\[
\pgwK^z_a(f)=c_{\mathrm{vac}}^{\,N_pN_q}\,S_\nu(\partial)\bigl[\Theta_f(s_0\,\cdot)\bigr](0),
\]
$c_{\mathrm{vac}}$ the constant of \eqref{pk:eq-vac}. By Remark~\ref{pgw:rem-jets} this is the kernel
of Lemma~\ref{pgw:one} for one choice of $j$, multiplied by a non-zero constant independent of
$\vartheta$.

(b) \emph{Branching form.} By Lemma~\ref{pgw:branch}(a) the invariant forms $B$ on $W^0_m\times
W^1_m$ form a space of dimension at most $1$. In the bases $(e_T)$ the conditions of invariance are
linear equations with rational coefficients (the representation matrices have rational coefficients
and the invariance under the compact group $G_{Y,\infty}$ is equivalent, by Zariski density, to the
polynomial identities $B(\rho_0(g)w_0,\rho_1(g\oplus1)w_1)=B(w_0,w_1)$ on $\prod_\tau\auGL_n$). So
if a non-zero invariant form exists, there is one whose matrix in the bases $(e_T)$ has coprime
rational integer entries; we take $B$ to be this matrix (it is unique up to sign), and $B=0$ if
there is no non-zero invariant form.
\end{definition}

\noindent\emph{The bases $(e_T)$ and the Schur matrices.} The Schur matrix $S_\nu$ of (a) is the matrix,
in the bases $(e_T)$ of \S\ref{pgs:ss-weyl}, of the Schur functor $Z\mapsto S(Z)$ of
\eqref{ps:eq-schurnorm}; these are the coordinates in which \S\ref{ss:au-spaces} writes the matrix of $g$
as the Schur matrix, so $S_\nu$ is the matrix of \eqref{ps:eq-kernel} for the edge. Its entries are
polynomials with integer coefficients in the entries of $Z$: in the basis $(b_T)$ of standard products of
minors, $S(Z)b_{T'}(x)=b_{T'}(Z^{\mathsf t}x)$ is, by the Cauchy--Binet formula and straightening, an integer
combination of the $b_T$ with coefficients integer polynomials in $Z$; and $(e_T)$ is the dual basis for
the invariant form, for which $S(Z)^*=S(Z^*)$, so the matrix of $S(Z)$ in $(e_T)$ is the conjugate
transpose of that of $S(Z^*)$ in $(b_T)$, again with integer polynomial entries in $Z$. The operator norm
\eqref{ps:eq-schurnorm} and the degree $e(\nu)$ do not depend on the basis. These are the only properties
of $S_\nu$ used in Lemmas~\ref{pgz:deg}, \ref{pgz:char} and~\ref{pgt:jet}.

\noindent By the last sentence of the paragraph containing \eqref{pgw:eq-Lambda}, Theorem~\ref{pgw:last}
holds for these choices, with another constant $\beta(g)$. The choice of the matrix $B$ only affects
(b) of Proposition~\ref{pgr:den}; see \S\ref{pgz:ss-b}.

\begin{lemma}[The degree of the kernels]\label{pgz:deg}
For every edge $e^z_a$, $e(\nu)\le40fn_mb\le x^{19}b$.
\end{lemma}
\begin{proof}
By Lemma~\ref{pgw:howe} and the proof of Lemma~\ref{pgw:one}, $W'=W^z_a$ is the dual of
$B_{\nu}$, which is the restriction to $\prod_\tau\auU(N_q)$ of $\bigotimes_\tau W_{\nu_\tau}(
\C^{N_q})^*$ twisted by a character; by Lemma~\ref{pgw:howe}, $\nu_\tau$ has at most $\min(N_p,N_q)$
rows. Along the half $h_z$ the node $p^z_a$ has $a$ blocks of the first kind and $p^z_{a-1}$ has
$a-1$, with the same stem and the same blocks of the second kind (\S\ref{pgw:ss-last}), so
$N_q=N_p+t'$ with $t'\in\{t,t+1\}$ the rank of a block of the first kind; thus $N_p<N_q$ and
$\nu_\tau$ has at most $N_q-1$ rows. So the highest weight of $W^z_a$ at $\tau$ is, up to adding a
constant vector (the twisting character, which is the same for all entries), $\pm\nu_\tau$ padded by
at least one zero, reordered. Adding a constant vector does not change differences of entries, and
since a padded zero entry occurs, the largest minus the smallest entry of $\pm\nu_\tau$ padded is
exactly $\nu_{\tau,1}-0$; hence the longest row $\nu_{\tau,1}$ equals the difference between the
largest and the smallest entry of the highest weight. By the proof of Lemma~\ref{pgs:boxes}, for the weight at the node $q$ this difference
is at most $20b$. Hence $|\nu_\tau|\le N_q\nu_{\tau,1}\le2n_m\cdot20b$, and summing over the $f$
places $\tau$ gives $e(\nu)\le40fn_mb\le x^{19}b$ (Lemma~\ref{pgs:boxes}).
\end{proof}

\subsection{The levels are free, and $\Lambda$ is a sum without denominators}\label{pgz:ss-free}

\begin{lemma}[Free levels on the last path]\label{pgz:free}
For every $z$, $a\le m$ and $g\in G^z_{a,f}$, $\Gamma^z_a\cap gJ^z_ag^{-1}=1$. Consequently, for every
open subgroup $J'\subseteq J^z_a$ and every function $F$ on $G^z_{a,f}$ that is left
$\Gamma^z_a$-invariant and right $J'$-invariant,
\[
\pgwint{J^z_a}_{\Gamma^z_a\backslash G^z_{a,f}}F(x)\,dx=\sum_{x\in\Gamma^z_a\backslash G^z_{a,f}/J^z_a}
\int_{J^z_a}F(xk)\,dk ,
\]
$dk$ the Haar probability measure of $J^z_a$; if $F$ is right $J^z_a$-invariant this is
$\sum_{x}F(x)$, a sum over representatives of the classes with coefficient $1$.
\end{lemma}
\begin{proof}
The rank is $N^z_a\le2n_m$ (Lemma~\ref{pgs:boxes}). Both $J(\pcb)$ and $J(\pcm)$ are contained in
$J(\pi_M)$, whose elements are $\equiv1$ modulo every chosen place, because the chosen places divide
$\pi_M$ (proof of Lemma~\ref{pz:sizes}(a)). The space at $p^z_a$ is a standard positive definite
space (Definition~\ref{pgv:def-spaces}: for $a=0$, $z=0$ it is $Y_m$, definite by
\S\ref{pgr:ss-arch}; the other nodes of row $m$ are diagonal with all signs $+1$, by
Lemma~\ref{pgv:count}), so Lemma~\ref{au:free} applies and gives the first assertion. For the
formula: the class $\Gamma xJ^z_a$ is the disjoint union of the classes $\Gamma xkJ'$, $k$ running
over $J^z_a/J'$, which are pairwise distinct since $\Gamma\cap xJ^z_ax^{-1}=1$; each has measure
$[J^z_a:J']^{-1}$ for the measure giving $J^z_a$ the volume $1$ (its stabiliser in $\Gamma$ being
trivial). Hence the integral over $\Gamma\backslash\Gamma xJ^z_a$ is
$[J^z_a:J']^{-1}\sum_{k\in J^z_a/J'}F(xk)=\int_{J^z_a}F(xk)\,dk$.
\end{proof}

\begin{lemma}[Moving the data by the levels]\label{pgz:move}
Let $(p,q)=e^z_a$, $X\in\pcMat_{q\times p}(\pcR_\pcm)$, $k\in K_p$, $k'\in K_q$ with reductions
$\bar k,\bar k'$ modulo $\pcm$, and $x\in G_{p,f}$, $h\in G_{q,f}$. Then
\[
\pgwK^z_a\bigl(\omega_f(xk,hk')[\delta_X]_f\bigr)=\chi_{pq}\bigl(\det\bar k'/\det\bar k\bigr)\,
\pgwK^z_a\bigl(\omega_f(x,h)[\delta_{\bar k'X\bar k^{-1}}]_f\bigr).
\]
\end{lemma}
\begin{proof}
$\omega_f(xk,hk')=\omega_f(x,h)\omega_f(k,k')$. By Lemma~\ref{pc:level}, which holds for the data of
Definition~\ref{pgx:conv} (remark after that definition), $\omega_f(k,k')[\delta_X]_f$ is the datum
of the function $y\mapsto\alpha_\pcm(\det\bar k)^{-1}\beta_\pcm(\det\bar k')\delta_X(\bar k'^{-1}y\bar
k)=\alpha_\pcm(\det\bar k)^{-1}\beta_\pcm(\det\bar k')\delta_{\bar k'X\bar k^{-1}}(y)$; and
$\alpha_{pq}$, $\beta_{pq}$ act on determinants of elements of $K$ through the same character
$\chi_{pq}$ (remark after Definition~\ref{pgv:def-edges}). Apply the linear map $\pgwK^z_a\circ
\omega_f(x,h)$.
\end{proof}

For an edge tuple $S$, points $\mathbf x=(x^z_a)_{z\in\{0,1\},\,0\le a<m}$ with $x^z_a\in G^z_{a,f}$,
and $h\in G_{Y,f}$, put $x^0_m=h$, $x^1_m=h\oplus\pgupt^m$ ($\pgupt=\pgupt_\pcm$ of
Definition~\ref{pgv:def-period}) and define the vectors \emph{reached along $h_z$}
\begin{equation}\label{pgz:eq-reached}
v^z_S(\mathbf x,h)=\pgwK^z_m\bigl(\omega_f(x^z_{m-1},x^z_m)[\delta_{S_{e^z_m}}]_f\bigr)\cdots
\pgwK^z_1\bigl(\omega_f(x^z_0,x^z_1)[\delta_{S_{e^z_1}}]_f\bigr)\,\psi^z(x^z_0)\in W^z_m,
\end{equation}
with $\psi^0=\psi_{y'}$, $\psi^1=\psi_\varphi$.

\begin{proposition}[$\Lambda$ as a sum over classes]\label{pgz:expand}
Let $\Psi_\vartheta$ be a path datum on $\pgrE(\pgrlast)$ with bound $N$, and let $y'$, $\varphi$ be
fixed by the levels $J^0_0=J(\pcm)$ and $J^1_0=J(\pcb)$ of the two outer nodes. Choose representatives
$\mathcal R^z_a$ of $\Gamma^z_a\backslash G^z_{a,f}/J^z_a$ ($a<m$) and $\mathcal R_Y$ of
$\Gamma_Y\backslash G_{Y,f}/J_Y$. Then
\[
\Lambda(\Psi_\vartheta;y',\varphi)=\sum_{\mathbf x\in\prod_{z,a<m}\mathcal R^z_a}\ \sum_{h\in\mathcal
R_Y}\ \sum_{S\in\mathcal S_\vartheta(\pgrlast)}\Psi_\vartheta(S)\,B\bigl(v^0_S(\mathbf x,h),
v^1_S(\mathbf x,h)\bigr).
\]
\end{proposition}
\begin{proof}
\emph{Unfolding.} For fixed $S$, the composite $\pgwT^z_S\psi^z$ evaluated at $x^z_m$ is, by the
definition of the algebraic lifts in \S\ref{pgw:ss-last}, the iterated integral
$\pgwint{J^z_{m-1}}\cdots\pgwint{J^z_0}$ of \eqref{pgz:eq-reached} over $x^z_{m-1},\dots,x^z_0$ (each
integral is a finite sum, so the order is immaterial). As $B$ is bilinear and all sums are finite,
\[
\Lambda=\pgwint{J_Y}_{\Gamma_Y\backslash G_{Y,f}}\ \pgwint{J^\bullet}\ F(\mathbf x,h)\,d\mathbf x\,dh,
\qquad F(\mathbf x,h)=\sum_{S\in\mathcal S_\vartheta(\pgrlast)}\Psi_\vartheta(S)\,B\bigl(v^0_S(\mathbf
x,h),v^1_S(\mathbf x,h)\bigr),
\]
where $\pgwint{J^\bullet}$ is the iterated class integral over the $2m$ variables $x^z_a$ ($a<m$), each
for the measure giving $J^z_a$ the volume $1$. Each function of one variable obtained by fixing the
others is left invariant under the corresponding $\Gamma$ (Lemma~\ref{pgw:one}(a), the $\Gamma$-equivariance
of $\psi^z$, and for $h$ also Lemma~\ref{pgw:branch}(b): $B(\eta(\gamma h),\varphi((\gamma
h)\oplus\pgupt^m))=B(\eta(h),\varphi(h\oplus\pgupt^m))$ by the invariance of $B$ and
$(\gamma h)\oplus\pgupt^m=(\gamma\oplus1)(h\oplus\pgupt^m)$), and right invariant under an open
subgroup.

\emph{Invariance.} Let $k^z_a\in J^z_a$ for $a<m$, $k_Y\in J_Y$, and put $k^0_m=k_Y$,
$k^1_m=k_Y\oplus1$ (an element of $J^1_m=J(\pcb)$ of $\pgrY^m_m$). We claim
\begin{equation}\label{pgz:eq-inv}
F\bigl((x^z_ak^z_a),hk_Y\bigr)=F(\mathbf x,h).
\end{equation}
Note $x^0_mk^0_m=hk_Y$ and $x^1_mk^1_m=(h\oplus\pgupt^m)(k_Y\oplus1)=hk_Y\oplus\pgupt^m$, so the left
side is formed with the points $x^z_ak^z_a$ for all $a\le m$. Let $g$ be the node family with
$g_{p^z_a}=\bar k^z_a$ (reduction modulo $\pcm$) for $a\le m$ and $g_p=1$ at the other nodes. It is
admissible (\S\ref{pgv:ss-def}): $g_{(0,0,m)}=1$ because $k^0_0\in J(\pcm)$; $g_{(1,0,m)}\equiv1
\bmod\pcb$ because $k^1_0\in J(\pcb)$; $g_{(1,m,m)}=g_{(0,m,m)}\oplus1$; and at the nodes $(z,a,a)$
with $a<m$, which are not on $\pgrlast$, both sides are $1$. By the outer invariance
$\psi^z(x^z_0k^z_0)=\psi^z(x^z_0)$, and by Lemma~\ref{pgz:move} applied to each factor of
\eqref{pgz:eq-reached},
\[
v^z_S\bigl((x^{z'}_ak^{z'}_a),hk_Y\bigr)=\prod_{a=1}^m\chi_{e^z_a}\Bigl(\frac{\det g_{p^z_a}}{\det
g_{p^z_{a-1}}}\Bigr)\cdot v^z_{g^{-1}S}(\mathbf x,h),
\]
where $(g^{-1}S)_e=g_qS_eg_p^{-1}$ for $e=(p,q)$. The product of the character factors over both
halves is $\chi_{\pgrE(\pgrlast)}(g)$. So the left side of \eqref{pgz:eq-inv} is
$\sum_S\Psi_\vartheta(S)\chi_{\pgrE(\pgrlast)}(g)B(v^0_{g^{-1}S},v^1_{g^{-1}S})(\mathbf x,h)$. The
map $S\mapsto g^{-1}S$ is a bijection of $\mathcal S_\vartheta(\pgrlast)$; substituting $S=gS'$ and
using $\chi_{\pgrE(\pgrlast)}(g)\Psi_\vartheta(gS')=\Psi_\vartheta(S')$ (Definition~\ref{pgr:def-datum})
gives $F(\mathbf x,h)$.

\emph{Conclusion.} By \eqref{pgz:eq-inv} with all but one of the $k$ equal to $1$, $F$ is right
invariant under $J^z_a$ in the variable $x^z_a$ and under $J_Y$ in $h$, the other variables being
fixed. Applying Lemma~\ref{pgz:free} to each of the $2m+1$ class integrals in turn replaces it by the
sum over the representatives, with coefficient $1$.
\end{proof}

\begin{remark}[The denominators of the class sums]\label{pgz:rem-den}
Proposition~\ref{pgz:expand} settles the first point of Proposition~\ref{pgr:den}: in $\Lambda$ every
class sum is normalised by the level over which it is taken, the levels are free
(Lemma~\ref{pgz:free}), and the joint invariance \eqref{pgz:eq-inv} makes each class sum a sum with
coefficients $1$. The indices $[J:J']$ of auxiliary subgroups fixing the individual terms, and the
orders of stabilisers, cancel or are $1$; they contribute nothing to $N_0$, and no condition relating
them to $\ell(\iota)$ is needed. The indices $[K_p:J_p]$ occur only in the comparison of $\Lambda$ with
the periods (the number $D(\vartheta)$ in the proof of Theorem~\ref{pgw:last}), not in $\Lambda$.
\end{remark}

\begin{lemma}[Counting the terms]\label{pgz:count}
The number of triples $(\mathbf x,h,S)$ in Proposition~\ref{pgz:expand} is
$\#\mathcal S_\vartheta(\pgrlast)\cdot\#(\Gamma_Y\backslash G_{Y,f}/J_Y)\prod_{z,\,a<m}
\#(\Gamma^z_a\backslash G^z_{a,f}/J^z_a)$, and $\#\mathcal S_\vartheta(\pgrlast)=\prod_{z,a}
\#\pcMat_{N^z_a\times N^z_{a-1}}(\pcR_\pcm)$.
\end{lemma}
\begin{proof}
An edge tuple vanishing off $\pgrE(\pgrlast)$ is a choice of one matrix for each of the $2m$ edges
$e^z_a$.
\end{proof}

\subsection{One lift does not increase gauges after multiplication by $\pgzd$}\label{pgz:ss-lift}

Put
\begin{equation}\label{pgz:eq-d}
\pgzd=M^{4bfn_m^2}\,(dM^b)^{40fn_mb},
\end{equation}
$d$ the discriminant of $K$ (a positive integer; $40fn_mb$ is an integer).

\begin{lemma}[The edge data are standard]\label{pgz:std}
Let $\iota\in\auI$ with $\ell=\ell(\iota)\nmid\auNm Q_\vartheta$, $(p,q)=e^z_a$, $X\in\pcMat_{q\times
p}(\pcR_\pcm)$, and let $x\in G_{p,f}$, $h\in G_{q,f}$ be integral with integral inverses at the places
above $\ell$. Then $c_{\mathrm{vac}}^{\,N_pN_q}M^{bfN_pN_q}\,\omega_f(x,h)[\delta_X]_f$ is an
$\iota$-unit times a standard datum of level $M^b$ at $\ell$ (\S\ref{pf:ss-frames}).
\end{lemma}
\begin{proof}
\emph{The datum $[\delta_X]_f$.} Outside $S(\pcm)$ it is the vacuum. At $v=w\aubar w\in S(\pcm)$,
$w\mid\pcm$, it is $\delta_X^\natural$, the indicator of $a_w(\tilde X+\pcm M_w^\circ)$ with
$M_w^\circ=\pcMat_{q\times p}(\auO_w)$ and $\tilde X$ an integral lift of $X$ (Definition~\ref{pgx:conv}).
By Lemma~\ref{pgx:vac} the vacuum at $v$ is the indicator of $M_w=a_wM^\circ_w$. As in the proof of
Lemma~\ref{plf:KS} (paragraph ``the datum as a combination of translates of the vacuum'', with the
lattice $M_w$ in place of the integral matrices; only the self-duality of $M_v=M_w\oplus M_{\aubar w}$ is
used there), the indicator of $a_w\tilde X+\pcm M_w$ is a root of unity times $\psrho^\circ(a_w\tilde
X)\mathbf 1_{\pcm M_w}$, and $\mathbf 1_{\pcm M_w}=[M_w:\pcm M_w]^{-1}\sum_y\zeta_y\psrho^\circ(y)
\mathbf 1_{M_w}$, with $[M_w:\pcm M_w]$ terms, roots of unity $\zeta_y$, and $y$ in the
$\aubar w$-component with $y\in\aubar\pcm^{-1}M_{\aubar w}$. Taking the product over $v\in S(\pcm)$ and
writing $\psvac=c_{\mathrm{vac}}^{-N_pN_q}\psvac^\circ$ \eqref{pk:eq-vac},
\[
c_{\mathrm{vac}}^{\,N_pN_q}\,\auNm(\pcm)^{N_pN_q}\,[\delta_X]_f=\sum_{x'}\zeta_{x'}\,\psrho^\circ(x')
\psvac^\circ ,
\]
a finite sum with roots of unity $\zeta_{x'}$ and translations $x'$ that vanish outside $S(\pcm)$ and
at $v\in S(\pcm)$ have components $a_w\tilde X\in M_w$ and $y\in\aubar\pcm^{-1}M_{\aubar w}$.

\emph{Level $M^b$ at $\ell$.} Let $w\mid\pcm$ lie above $\ell$. As $\ell\nmid\auNm Q_\vartheta$, $w$ is
not in $Q_\vartheta$, so $w\mid\pi_M$; $\ell\nmid q_1$ (since $\ell\mid M$ and $p_0,p_1\equiv1\pmod
{q_1}$), so $a_w=1$ (Proposition~\ref{pgv:alpha}(a) and the convention after Definition~\ref{pgx:conv}),
$M_w=M^\circ_w$, and $M_{\aubar w}$ is the integral lattice at $\aubar w$. Then $M^b\in\pcm_w\cap
\aubar\pcm_{\aubar w}$ ($\pcm_w=\pi_M^b\auO_w$ and $\aubar{\pi_M}^b\mid M^b$), so $M^bx'$ is integral at
the places above $\ell$. Hence the right side is a standard datum of level $M^b$ at $\ell$, with
coefficients roots of unity.

\emph{Norms.} $\auNm(\pcm)=\auNm(\pcb)\auNm Q_\vartheta$ (the primes of $\pi_M$ are not in $Q_\vartheta$),
$\auNm Q_\vartheta$ is prime to $\ell$, and $\auNm(\pcb)=\auNm_{E/\Q}(\pi_M)^b=M^{bf}$ ($\pi_M\in K$,
$[E:K]=f$). So $c_{\mathrm{vac}}^{N_pN_q}M^{bfN_pN_q}[\delta_X]_f$ is an $\iota$-unit times a standard
datum of level $M^b$.

\emph{The action of $(x,h)$.} By the remark preceding Theorem~\ref{pf:std}, $\omega(x,h)$ preserves
the standard data of level $M^b$ at $\ell$; $\omega_f=\omega_{\alpha,\beta,f}$ differs from it by
$\alpha(\det x)^{-1}\beta(\det h)$ and the unramified factors of \eqref{po:eq-w5}, which are
$\iota$-units for $x,h$ integral with integral inverses above $\ell$ (values of algebraic Hecke
characters at ideles that are units above $\ell$, and powers of residue characteristics different from
$\ell$), exactly as in the proof of Lemma~\ref{plf:KS}(KS1).
\end{proof}

\begin{lemma}[One lift]\label{pgz:kernel}
Let $\iota\in\auI$ with $\ell(\iota)\nmid\auNm Q_\vartheta$, $(p,q)=e^z_a$, $X\in\pcMat_{q\times
p}(\pcR_\pcm)$, $x\in G_{p,f}$ and $h\in G_{q,f}$. If $\xi\in W^z_{a-1}$ has gauge at most $1$ at
$(\iota,x)$, then $\pgzd\cdot\pgwK^z_a(\omega_f(x,h)[\delta_X]_f)\,\xi$ has gauge at most $1$ at
$(\iota,h)$.
\end{lemma}
\begin{proof}
\emph{Reduction to integral points.} Unitary groups are rational varieties and satisfy weak
approximation~\cite[\S7.1, Propositions~7.3 and~7.4]{PlatonovRapinchuk}, so there are rational points $\gamma\in\Gamma_p$,
$\gamma'\in\Gamma_q$ such that $k=\gamma^{-1}x$ and $k'=\gamma'^{-1}h$ are integral with integral
inverses at the places above $\ell$ (as in the proof of Proposition~\ref{pk:K12}). By
Lemma~\ref{pgw:one}(a), $\pgwK^z_a(\omega_f(x,h)\Phi)=\gamma'\circ\pgwK^z_a(\omega_f(k,k')\Phi)\circ
\gamma^{-1}$, the rational points acting on $W^z_{a-1}$, $W^z_a$ by the archimedean action. By (G-b)
of \S\ref{pgz:ss-choice}, with $\eta=\pgwK^z_a(\omega_f(x,h)\Phi)\xi$,
\[
h_\ell^{-1}\iota(\eta)=k'^{-1}_\ell\,\iota\bigl(\pgwK^z_a(\omega_f(k,k')\Phi)\bigr)\,k_\ell\,
\bigl(x_\ell^{-1}\iota(\xi)\bigr).
\]
By (G-a), $k_\ell$ and $k'^{-1}_\ell$ are matrices invertible over $\auO_{\C_\ell}$, and
$x_\ell^{-1}\iota(\xi)$ has integral coordinates. So it suffices to show that the matrix of
$\pgzd\,\pgwK^z_a(\omega_f(k,k')[\delta_X]_f)$ in the bases $(e_T)$ has entries in $\auO_\iota$.

\emph{Integrality.} By Lemma~\ref{pgz:std}, $c_{\mathrm{vac}}^{N_pN_q}\omega_f(k,k')[\delta_X]_f=
u\,M^{-bfN_pN_q}f'$ with an $\iota$-unit $u$ and a standard datum $f'$ of level $M^b$ at $\ell$. By
Definition~\ref{pgz:def-choice}(a) and the homogeneity of the entries of $S_\nu$ (of total degree
$e=e(\nu)$),
\[
\pgwK^z_a(\omega_f(k,k')[\delta_X]_f)=u\,M^{-bfN_pN_q}(dM^b)^{-e}\,S_\nu(\partial)\bigl[\Theta_{f'}
(dM^bs_0\,\cdot)\bigr](0).
\]
The series $\Theta_{f'}(dM^bs_0\,\cdot)$ is $\partial$-integral at $\iota$: for $\ell\ne d$ by the last
assertion of Theorem~\ref{pf:std} with $\delta=M^b$ ($\aubar\delta=\delta$) and $c=d$; for $\ell=d$ by
the first assertion of Theorem~\ref{pf:std} and the argument of the proof of Lemma~\ref{plf:KS}(KS1)
for $\ell=d$ (there written for $\delta=M^2$; it uses only the product formula of
Theorem~\ref{pf:std}, the $\partial$-integrality of $y$ at every prime, and that of $\theta(d\,\cdot)$).
The entries of $S_\nu$ are polynomials with integer coefficients, so the jets are in $\auO_\iota$.
Finally $N_pN_q\le4n_m^2$ (Lemma~\ref{pgs:boxes}) and $e\le40fn_mb$ (Lemma~\ref{pgz:deg}), so
$M^{bfN_pN_q}(dM^b)^e$ divides $\pgzd$.
\end{proof}

\begin{corollary}[Vectors reached along a half]\label{pgz:reach}
Let $\iota$ be as in Lemma~\ref{pgz:kernel} and let $\psi^z$ be integral at $\iota$. Then for every
$S$, $\mathbf x$ and $h$, $\pgzd^mv^0_S(\mathbf x,h)$ has gauge at most $1$ at $(\iota,h)$ and
$\pgzd^mv^1_S(\mathbf x,h)$ has gauge at most $1$ at $(\iota,h\oplus\pgupt^m)$.
\end{corollary}
\begin{proof}
Induction on the number of factors in \eqref{pgz:eq-reached}: $\psi^z(x^z_0)$ has gauge at most $1$
at $x^z_0$, and Lemma~\ref{pgz:kernel} (gauges scale with the absolute value of the scalar) passes from
$x^z_{a-1}$ to $x^z_a$, the last point being $x^z_m$.
\end{proof}

\subsection{The branching pairing at the point $h\oplus\pgupt^m$}\label{pgz:ss-pair}

\begin{lemma}[The idele $\alpha$ and the flip scalar]\label{pgz:flip}
(a) $\alpha^{-1}$ and $q_1\alpha$ are integral finite ideles of $E$.
(b) $\mathsf a=\auNm Q_\vartheta\cdot M^b\cdot q_1\cdot\pgupt$ and $\aubar{\mathsf a}=\auNm Q_\vartheta
\cdot M^b\cdot q_1\cdot\pgupt^{-1}$ are integral finite ideles of $E$.
\end{lemma}
\begin{proof}
(a) Let $P\in\auO_K[x]$ be the minimal polynomial of $\zeta=\zeta_{q_1}$ over $K$ (\S\ref{pgv:ss-G},
before Proposition~\ref{pgv:alpha}). It is a monic divisor of $x^{q_1}-1$, so $x^{q_1}-1=PG$ with $G$
monic; $G$ has coefficients in $\auO_K$, being a monic factor of a monic polynomial with integral
coefficients. Differentiating at $\zeta$: $q_1\zeta^{q_1-1}=P'(\zeta)G(\zeta)$, so $q_1/P'(\zeta)=
\zeta^{-1}G(\zeta)\in\auO_E$. Since $\alpha$ is $1/P'(\zeta)$ or $1$ at every prime, $\alpha^{-1}$ and
$q_1\alpha$ are integral.

(b) $\pgupt=\aubar u_0/u_0$ with $u_0=\alpha\varpi_\pcm$, so $\mathsf a=(\auNm Q_\vartheta M^b
\varpi_\pcm^{-1})\cdot\aubar{\varpi_\pcm}\cdot\aubar{q_1\alpha}\cdot\alpha^{-1}$. The integer
$\auNm Q_\vartheta M^b$ lies in $\pcm$ ($\pi_M^b\mid M^b$, and $\auNm w\in w$ for each $w\in
Q_\vartheta$), and $\varpi_\pcm$ generates $\pcm\widehat\auO_E$; so the first factor is integral. The
others are integral by (a) and since complex conjugation preserves integrality. As $\pgupt\aubar\pgupt
=1$, $\aubar{\mathsf a}=\auNm Q_\vartheta M^bq_1\pgupt^{-1}$, integral as the conjugate of $\mathsf a$.
\end{proof}

\begin{lemma}[The character of the branching form]\label{pgz:char}
Suppose $B\ne0$. There are integers $L_\tau$ $(\tau\in\auPl)$ with $|L_\tau|\le30b$ such that, as
polynomial identities in $g\in\prod_\tau\auGL_n$ and $t\in\prod_\tau\auGL_1$ (over any field of
characteristic $0$, through the integral representation matrices),
\[
B\bigl(\rho_0(g)w_0,\rho_1(g\oplus t)w_1\bigr)=\prod_\tau t_\tau^{L_\tau}\,B(w_0,w_1),
\]
$\rho_0,\rho_1$ the representations of $\prod_\tau\auGL_n$, $\prod_\tau\auGL_{n+1}$ on $W^0_m$,
$W^1_m$ in the bases $(e_T)$.
\end{lemma}
\begin{proof}
For $t$ in the compact group $\prod_\tau\auU(1)$ the form $(w_0,w_1)\mapsto B(w_0,\rho_1(1\oplus t)w_1)$
is invariant in the sense of Lemma~\ref{pgw:branch}(a), since $1\oplus t$ commutes with $g\oplus1$; by
the uniqueness in that lemma it is $\chi(t)B$ for a continuous character $\chi$ of
$\prod_\tau\auU(1)$, so $\chi(t)=\prod_\tau t_\tau^{L_\tau}$ with $L_\tau\in\Z$. Together with the
invariance, $B(\rho_0(g)w_0,\rho_1(g\oplus t)w_1)=\chi(t)B(w_0,w_1)$ on $\prod\auU(n)\times\prod\auU(1)$;
both sides are polynomials in the entries of $g$, $t$ and in $\det g^{-1},t^{-1}$, with rational
coefficients, so by Zariski density the identity holds on $\prod\auGL_n\times\prod\auGL_1$ and then over
every field of characteristic $0$. For the bound, decompose $W^1_m$ into the eigenspaces of
$\rho_1(1\oplus t)$, $t\in\prod\auGL_1$; on each the action is by $\prod_\tau t_\tau^{k_\tau}$, with
$k_\tau$ the last coordinate at $\tau$ of a weight of $W^1_m$. Since $B\ne0$, it is non-zero on some
eigenspace, and there $\chi(t)=\prod t_\tau^{k_\tau}$, so $L_\tau=k_\tau$. The coordinates of the
weights of an irreducible representation lie between the smallest and the largest entry of its highest
weight; by the proof of Lemma~\ref{pgs:boxes} these are at most $|e_\tau|+20b\le30b$ in absolute value.
\end{proof}

Put
\begin{equation}\label{pgz:eq-M}
\pgzM=\prod_{\tau\in\auPl}(M^bq_1)^{m|L_\tau|}\qquad(\pgzM=1\text{ if }B=0).
\end{equation}

\begin{lemma}[Clearing the branching pairing]\label{pgz:pair}
Let $\iota\in\auI$ with $\ell=\ell(\iota)\nmid\auNm Q_\vartheta$, $h\in G_{Y,f}$, and let $v'\in W^0_m$
have gauge at most $1$ at $(\iota,h)$ and $v\in W^1_m$ gauge at most $1$ at $(\iota,h\oplus\pgupt^m)$.
Then $\pgzM\,B(v',v)\in\auO_\iota$.
\end{lemma}
\begin{proof}
We may assume $B\ne0$. Let $w'=h_\ell^{-1}\iota(v')$ and $w=(h\oplus\pgupt^m)_\ell^{-1}\iota(v)$, vectors
with coordinates in $\auO_{\C_\ell}$. Here $(h\oplus\pgupt^m)_\ell=(h_\ell\oplus1)(1\oplus\pgupt^m_\ell)$,
and $\pgupt_\ell$ acts through its images $t_\tau\in\C_\ell^\times$ at the primes $w_\tau$ of
\S\ref{ss:au-spaces}. By Lemma~\ref{pgz:char} over $\C_\ell$ (applied with $g=h_\ell$, $t=\pgupt^m_\ell$),
\[
\iota(B(v',v))=B\bigl(h_\ell w',(h_\ell\oplus\pgupt^m_\ell)w\bigr)=\prod_\tau t_\tau^{mL_\tau}\cdot
B(w',w).
\]
$B$ has integer entries, so $B(w',w)\in\auO_{\C_\ell}$. For each $\tau$, $(M^bq_1)^{m|L_\tau|}
t_\tau^{mL_\tau}$ is $(\auNm Q_\vartheta)^{-m|L_\tau|}$ times the image of $\mathsf a^{m|L_\tau|}$ or of
$\aubar{\mathsf a}^{m|L_\tau|}$ (according to the sign of $L_\tau$), which is integral by
Lemma~\ref{pgz:flip}(b); and $\auNm Q_\vartheta$ is an $\ell$-adic unit. Multiplying over $\tau$ gives
the claim.
\end{proof}

\subsection{Proof of Proposition~\ref{pgr:den}(a)}\label{pgz:ss-a}

\begin{theorem}[Proposition~\ref{pgr:den}(a)]\label{pgz:a}
Let $N_0=\pgzM\,\pgzd^{2m}$, with $\pgzd$, $\pgzM$ of \eqref{pgz:eq-d}, \eqref{pgz:eq-M}. Then
$1\le N_0\le\exp(x^{99}b^2)$, $N_0$ does not depend on $\vartheta$, $\Psi$, $y'$, $\varphi$, and the
assertion (a) of Proposition~\ref{pgr:den} holds for $\Lambda$ of \eqref{pgw:eq-Lambda} with the
choices of Definition~\ref{pgz:def-choice}.
\end{theorem}
\begin{proof}
\emph{Integrality.} Let $\Psi_\vartheta$, $(y',\varphi)$ and $\iota$ be as in (a). By
Proposition~\ref{pgz:expand}, $\Lambda$ is a finite sum of terms $\Psi_\vartheta(S)B(v^0_S,v^1_S)$.
$\Psi_\vartheta(S)$ is a sum of roots of unity, hence in $\auO_\iota$. By Corollary~\ref{pgz:reach},
$\pgzd^mv^0_S$ and $\pgzd^mv^1_S$ have gauge at most $1$ at $h$ and $h\oplus\pgupt^m$; by
Lemma~\ref{pgz:pair}, $\pgzM\pgzd^{2m}B(v^0_S,v^1_S)=\pgzM B(\pgzd^mv^0_S,\pgzd^mv^1_S)\in\auO_\iota$.

\emph{Independence.} $\pgzd$ depends only on $M,b,f,n_m,d$. The integers $L_\tau$ depend only on the
weights at $p^0_m$, $p^1_m$ (Lemma~\ref{pgz:char}), not on $\vartheta$.

\emph{Size.} By Lemma~\ref{an:lem-sizes}: $\log M\le22\log b$ (S5), $f\le x^{12}$ (S3), $d\le\log x$
((S1)), $q_1\le x^{12}$ (S3), $m\le2\sqrt x$ (D3), and $n_m\le40x^6$ (Lemma~\ref{pz:sizes}(b)). So
\[
\log\pgzd=4bfn_m^2\log M+40fn_mb(\log d+b\log M)\le4\cdot1600\,x^{24}\cdot22\,b\log b+1600\,x^{18}b
(\log x+22\,b\log b)\le x^{26}b^2,
\]
using $\log b\le x^{1/100}$ (S1). By Lemma~\ref{pgz:char}, $\log\pgzM\le f\cdot m\cdot30b\cdot(b\log M+
\log q_1)\le x^{12}\cdot2\sqrt x\cdot30b\cdot23\,b\log b\le x^{14}b^2$. Hence $\log N_0\le
x^{14}b^2+4\sqrt x\,x^{26}b^2\le x^{99}b^2$.
\end{proof}

\subsection{Proposition~\ref{pgr:den}(b): reduction to three size bounds}\label{pgz:ss-b}

Let $\vartheta=1$, so $\pcm=\pcb$, $\#\pcR_\pcb=\auNm\pcb=M^{bf}$, and put $G_0=\exp(x^{99}b^2)$.
Part (b) follows from Proposition~\ref{pgz:expand} and the following three bounds, which are proved
in \S\ref{pgt:sec}: (B1) is Proposition~\ref{pgt:B1}, (B2) is Proposition~\ref{pgt:B2} and (B3) is
Corollary~\ref{pgt:B3}.

\begin{enumerate}
\item[(B1)] (One lift.) For every edge $e^z_a=(p,q)$, $X\in\pcMat_{q\times p}(\pcR_\pcb)$, $x\in
G_{p,f}$, $h\in G_{q,f}$, $\sigma\in\auAut(\C)$ and $\xi\in W^z_{a-1}$:
$\|\sigma(\pgwK^z_a(\omega_f(x,h)[\delta_X]_f)\xi)\|\le G_0\|\sigma(\xi)\|$, with the kernels of
Definition~\ref{pgz:def-choice}(a), $\sigma$ acting on coordinates in the bases $(e_T)$ and
$\|\ \|$ the invariant norms.
\item[(B2)] (The branching form.) If a non-zero invariant form exists, there is a non-zero invariant integer matrix $B'$ (in the bases
$(e_T)$) with $|B'(v',v)|\le C\|v'\|\,\|v\|$ for all $v',v$, where $C=(L^L)^3$ and $L\ge1$ is an integer
bound for the number of boxes of the weight at $p^1_m$; by Lemma~\ref{pgs:boxes} one may take
$L=\lfloor x^{19}b\rfloor$.
\item[(B3)] (Classes.) For $z\in\{0,1\}$ and $a\le m$, $\#(\Gamma^z_a\backslash G^z_{a,f}/J^z_a)\le G_0$
when $\vartheta=1$.
\end{enumerate}

\noindent Given (B2), the matrix $B$ of Definition~\ref{pgz:def-choice}(b) satisfies the same bound:
by uniqueness $B'=cB$ with $c$ rational, and as $B$ has coprime integer entries and $B'$ integer entries,
$c$ is a non-zero integer, so $|B(v',v)|\le|B'(v',v)|$.

\begin{theorem}[Proposition~\ref{pgr:den}(b)]\label{pgz:b}
Let $\Psi_1$ be a path datum on $\pgrE(\pgrlast)$ with bound $L(1)$, let
$(y',\varphi)$ be fixed by the outer levels, and $\|\sigma(\psi_{y'}(g))\|^2\le B_0$,
$\|\sigma(\psi_\varphi(g))\|^2\le B_1$ for all $g$. Then $|\sigma(\Lambda(\Psi_1;y',\varphi))|\le
\exp(x^{200}b^2)\sqrt{B_0B_1}$.
\end{theorem}
\begin{proof}
The bounds (B1), (B2), (B3) hold by Propositions~\ref{pgt:B1}, \ref{pgt:B2} (if no non-zero invariant
form exists, $B=0$ and the bound is trivial, Remark~\ref{pgt:rem-B2}) and Corollary~\ref{pgt:B3}.
By Proposition~\ref{pgz:expand}, $\sigma(\Lambda)=\sum\sigma(\Psi_1(S))\,B(\sigma v^0_S,\sigma v^1_S)$, as
$B$ has rational entries. By (B1) and induction along \eqref{pgz:eq-reached},
$\|\sigma v^0_S\|\le G_0^m\sqrt{B_0}$ and $\|\sigma v^1_S\|\le G_0^m\sqrt{B_1}$; by (B2) each pairing is
at most $C\,G_0^{2m}\sqrt{B_0B_1}$; $|\sigma\Psi_1(S)|\le L(1)$. By Lemma~\ref{pgz:count} and (B3) the
number of terms is at most $\#\mathcal S_1(\pgrlast)\,G_0^{2m+1}$, and $\#\mathcal S_1(\pgrlast)\le
M^{bf\cdot2m(2n_m)^2}$ (Lemma~\ref{pgs:boxes}: ranks $\le2n_m$). Also $L(1)=V(1)^{m^2}\mathcal
E(1)^{2m}=M^{bf(2n_mm^2+2m(2n_m)^2)}$. With the bounds of Lemma~\ref{an:lem-sizes} used in the proof of
Theorem~\ref{pgz:a},
\[
\log\bigl(L(1)\,\#\mathcal S_1(\pgrlast)\bigr)\le bf\log M\,(2n_mm^2+4m(2n_m)^2)\le x^{12}\cdot22\,b\log b
\cdot(320x^7+51200\,x^{12.5})\le x^{28}b^2,
\]
$\log C=3L\log L\le x^{20}b^2$, and $(4m+1)x^{99}b^2\le x^{100}b^2$. The total is at most
$x^{28}b^2+x^{20}b^2+x^{100}b^2\le x^{200}b^2$.
\end{proof}





%% ===== proofs/C-grid9.tex =====
% proofs/C-grid9.tex
% Input (V) of proofs/C-grid5.tex (deep-level vanishing by Poincare series).
% Labels start with pgq:. Conventions and macros: those of proofs/C-grid5.tex and C-grid4.tex.
\providecommand{\auU}{\mathrm U}
\providecommand{\auO}{\mathcal O}
\providecommand{\auA}{\mathbb A}
\providecommand{\auAz}{\mathcal A_0}
\providecommand{\auPl}{\Phi}
\providecommand{\aubar}[1]{\overline{#1}}
\providecommand{\pcb}{\mathfrak b}
\providecommand{\pgrY}{\widetilde Y}
\providecommand{\pcnot}[1]{\par\smallskip{\raggedright\noindent\textsc{Not proved here.} [#1]\par}\smallskip}
\providecommand{\pgqvol}{\operatorname{vol}}
\providecommand{\pgqcol}{\operatorname{col}}
\providecommand{\pgqSl}{\mathsf{Sl}}

\section{Deep-level vanishing}\label{pgq:sec}

This section proves input (V) of \S\ref{pgw:ss-first} (Theorem~\ref{pgq:V}), from three
inputs (C), (Q), (S) stated in \S\ref{pgq:ss-inputs} and proved in \S\S\ref{pgk:sec}, \ref{pgl:sec}. The argument halves every bound of
$|y|$: high in the cusp by the cusp property of $y$ and the slow variation of its reproducing kernel
(\S\ref{pgq:ss-cusp}); elsewhere by a Poincar\'e series of deep level whose kernel is a translated
matrix coefficient of the type of $\pgrY^0_0$. Its period against $y$ vanishes by the hypothesis of
(V), and by unfolding it equals a smoothing of $y$ up to an error that is small because the level
$\pcb=(\pi_M)^b$ is deep (\S\ref{pgq:ss-poinc}); averaging the smoothing against the kernel of $y$
returns $y$ times a product of local pairings, which is not small at the chosen point $g_0$
(\S\ref{pgq:ss-bulk}).

\subsection{Setting and statement}\label{pgq:ss-stat}

\emph{Spaces and groups.} Let $V=Y_0$, of rank $n=r$, and $Z=\pgrY^0_0=V\oplus E\tilde e$ with
$\langle\tilde e,\tilde e\rangle=1$ (Definition~\ref{pgv:def-spaces}); the lattice of $Z$ is
$L_Z=L_V\oplus\auO_E\tilde e$. Put $G=\auU(V)$, $H=\auU(Z)$, embed $G\to H$ by $h\mapsto h\oplus1$,
and write $\Gamma_G=G(E^+)$, $\Gamma_H=H(E^+)$, $[G]=\Gamma_G\backslash G(\auA)$. Let $f=\#\auPl$.
For $\tau\in\auPl$ we use on $V_\tau=V\otimes_{E,\tau}\C$ the normalised basis of
\S\ref{pds:ss-setup}, in which the hermitian form has a hermitian unitary matrix $\mathsf J_\tau$, and
on $Z_\tau$ that basis followed by $\tilde e$ (matrix $\mathsf J_\tau\oplus1$). Norms of vectors are
the Euclidean norms in these bases and $\|g\|$ is the operator norm; thus $\|g\|=\|g^{-1}\|\ge1$ for
$g$ in $G_\tau$ or $H_\tau$ (as $g^{-1}=\mathsf Jg^*\mathsf J$), and $\|r\oplus1\|=\|r\|$. For
$g\in G_\infty=\prod_\tau G_\tau$ put $\|g\|=\prod_\tau\|g_\tau\|$ and $\operatorname{dist}(g,1)=
\max_\tau\|g_\tau-1\|$; likewise on $H_\infty$. We write $G_\infty\subset G(\auA)$ and use that real
points commute with finite ones.

\emph{Measures.} $G_\tau$ carries the Haar measure of \textup{(DS1)(ii)} (Proposition~\ref{pds:haar}),
$G_\infty$ the product measure, $G(\auA_f)$ any Haar measure, $G(\auA)$ the product, and $[G]$ the
quotient measure, so that
\begin{equation}\label{pgq:eq-unfold}
\int_{[G]}\sum_{\delta\in\Gamma_G}F(\delta h)\,dh=\int_{G(\auA)}F(h)\,dh
\end{equation}
for $F\ge0$ measurable, and for $F\in L^1(G(\auA))$ (the inner sum then converging absolutely for
almost all $h$). Replacing the measure of $[G]$ by a positive multiple changes neither the hypothesis
nor the conclusion of (V). Put $V_\tau=\pgqvol\{r\in G_\tau:\|r-1\|\le1\}>0$.

\emph{Coordinates and the first minimum.} For $x\in G(\auA)$ and $v\in V(E)$ write $x^{-1}v\in
V(\auA)$ in the frame; $v$ lies in the lattice $L_V(x)$ if the finite coordinates of $x^{-1}v$ are
integral, and $|v|_x$ is the Euclidean norm of its archimedean coordinates in $\prod_\tau\C^n$.
Likewise $L_Z(x')$, $|\cdot|_{x'}$ for $x'\in H(\auA)$. The \emph{first minimum} is
\[
\lambda(x)=\min\bigl(1,\ \min\{|v|_x:0\ne v\in L_V(x)\}\bigr)\in(0,1].
\]

\emph{Types.} Let $c=\prod_\tau c_\tau$ and $e=d\,\aubar c$ ($d=\prod_\tau d_\tau$) be the coefficient
and the type of $\auAz$ (\S\ref{ss:au-spaces}), $c_\tau=\langle v_\tau,\pi_\tau(\cdot)v_\tau\rangle$; and
let $c'=\prod_\tau c'_\tau$, $e'=d'\aubar{c'}$ be those of the type $e_{(1,0,0)}$ of $Z$
(\S\ref{pgv:ss-def}: Harish-Chandra parameter $-\lambda_{(1,0,0),\tau}$, signs $\epsilon_{(1,0,0)}$). The
forms of type $e$ on $G$ (Definition~\ref{pz:def-type}) are bounded, continuous, cuspidal
(Proposition~\ref{pz:types}(a)). Let $J_G(\mathfrak a)$, $J_H(\mathfrak a)$ be the principal
congruence subgroups of level $\mathfrak a$ in the stabilisers of $\widehat L_V$, $\widehat L_Z$. Then
$J_M=J_G(\pi_M)\supseteq J_G(\pcb)$ and $J_M\supseteq J_G(M)$, as $\pi_M$ divides $\pcb=\pi_M^b$ and
$M=\pi_M\aubar{\pi_M}$ (\S\ref{ss:au-spaces}); and $h\oplus1\in J_H(\mathfrak a)$ if and only if
$h\in J_G(\mathfrak a)$. Since $\pi_M\in\auO_K$, $|\tau(\pi_M)|^2=M$ for every $\tau$.

\begin{theorem}[Input (V)]\label{pgq:V}
There is a real point
$g_0=\omega$ of $H$ such that every $y\in\auAz$ with
\[
\int_{[G]}\aubar{y(x)}\,\psi\bigl((x\oplus1)\omega\bigr)\,dx=0
\]
for all forms $\psi$ of type $e'$ on $H$ fixed by $J_H(\pcb)$ is zero.
\end{theorem}

\subsection{Three inputs}\label{pgq:ss-inputs}

Let $N_d=q_1d$, a natural number in the different of $E$ (\S\ref{psp:ss-sparse}), and
$\kappa_{\mathrm c}=8fn\,\mathrm{rd}_E+48n^6f^2N_d+1$.

\begin{enumerate}
\item[(C)] (\emph{Cusp estimate.}) Let $F$ be a form of type $e$ on $G$ fixed by $J_G(M)$, let
$0<\delta\le1$, and let $x\in G(\auA)$ satisfy $\lambda(x)<(\delta/(M\kappa_{\mathrm c}^2))^n$. If
$|F(xu)-F(x)|\le C$ for all $u\in G_\infty$ with $\operatorname{dist}(u,1)\le\delta$, then
$|F(x)|\le C$.
\item[(Q)] (\emph{Kernel forms.}) Let $K\subseteq H(\auA_f)$ be an open subgroup of the stabiliser of
$\widehat L_Z$ and $z\in H(\auA)$. The series $\sum_{\gamma\in\Gamma_H}e'((z\gamma g)_\infty)
\mathbf 1_K((z\gamma g)_f)$ converges absolutely for every $g\in H(\auA)$, and there is a form $Q$ of
type $e'$ on $H$ fixed by $K$ with $\aubar{Q(g)}$ equal to its sum.
\item[(S)] (\emph{Restriction pairing.}) Let $\beta_{\mathsf S}=b^{3m}$ and $N_{\mathsf S}=20(2n+1)^2$.
For every $\tau$ there is $\omega_\tau\in H_\tau$ with $\|\omega_\tau\|^2\le\beta_{\mathsf S}$ and
\[
V_\tau\le\beta_{\mathsf S}^{N_{\mathsf S}}\Bigl|\int_{G_\tau}c_\tau(r)\,\aubar{c'_\tau\bigl(
\omega_\tau^{-1}(r\oplus1)\omega_\tau\bigr)}\,dr\Bigr| .
\]
\end{enumerate}
\noindent Inputs (C) and (Q) are proved in \S\ref{pgk:sec} (Theorems~\ref{pgk:C} and~\ref{pgk:Q}), and
input (S) in \S\ref{pgl:sec} (Theorem~\ref{pgl:S}).

\begin{lemma}[Steepness]\label{pgq:steep}
Let $B=\kappa(\tau)+\Delta/2$ and $\Lambda$, $\Lambda'$ the highest weights of the lowest $K$-types for
the parameters $\varsigma_\tau$ of $V$ and $\lambda'=\lambda_{(1,0,0),\tau}$ of $Z$ \textup{(}signs
$\epsilon_{(0,0,0)}$, $\epsilon_{(1,0,0)}$\textup{)}. On the first $k-1$ coordinates both signs are $+$
and $\Lambda=\Lambda'=0$. On the remaining $m$, resp.\ $m+1$, coordinates, in the order of decreasing
parameter, the signs alternate starting with $+$, and $\Lambda$, resp.\ $\Lambda'$, is strictly
decreasing with values in $[-3B,-B]$.
\end{lemma}
\begin{proof}
For $V$ this is Lemma~\ref{pds:chamber} and the remark after it. For $Z$: by \eqref{pgv:eq-param}
with $z=1$, $a=b=0$, $s_p=0$, the stem coordinate $\varrho\in\{0,\dots,r\}$ carries
$\lambda'_\varrho=\varsigma'_{\tau,\varrho}+\frac12$, that is $\frac k2-\frac12-\varrho$ for $\varrho\le
k-2$ and $-\kappa(\tau)-j\Delta+\frac12$ with $j=\varrho+2-k\in\{1,\dots,m+1\}$ otherwise; by
\eqref{pgv:eq-sign} its sign is $(-1)^{a(\varrho)}=(-1)^{j-1}$ for $\varrho\ge k-1$ and $+$ before.
So the far coordinates, in decreasing order, have signs $+,-,\dots,+$, and $m/2$ of them are
negative. Use $\Lambda'=\lambda'+\rho_{\mathrm n}-\rho_{\mathrm c}$ (\S\ref{pds:ss-setup}). For $i\le k-2$
there are $i$ entries above $\lambda'_i$, all of sign $+$, and below it $k-2-i$ entries of sign $+$ and
the $m+1$ far entries, $m/2$ of sign $-$. Hence $2\rho_{\mathrm n,i}=m/2$, $2\rho_{\mathrm c,i}=
(k-2-i)+(m/2+1)-i$, and $2\Lambda'_i=(k-1-2i)+m/2-(k-2-2i)-(m/2+1)=0$. On the far coordinates
$|\Lambda'_j-\lambda'_j|\le n+1$, where $\lambda'_j=-\kappa(\tau)-j\Delta+\frac12$; the $\lambda'_j$ decrease in steps $\Delta>2n+2$, so $\Lambda'$ is
strictly decreasing, and $\Lambda'_j\ge-\kappa-(m+1)\Delta-n-\frac12\ge-3B$, $\Lambda'_j\le-\kappa-\Delta+n+\frac32\le-B$ because $m\Delta+r\le
\kappa(\tau)$ (Lemma~\ref{an:lem-sizes}(S7)(g)), and $n+2\le r^8=\Delta/2$ as $\Delta=2r^8$, $r\ge100$ (Lemma~\ref{an:lem-sizes}(S7)(b)); this also gives $\Delta>2n+2$.
\end{proof}

\subsection{Local estimates}\label{pgq:ss-local}

Let $h_V\ge\Delta$ and $h_Z\ge\Delta-t$ be the gaps of $\varsigma_\tau$ and of $-\lambda_{(1,0,0),\tau}$
(\S\ref{pds:ss-setup}, Lemma~\ref{pgv:gap}), $L_\tau=\max_i|\varsigma_{\tau,i}|\le\kappa(\tau)+m\Delta\le2b$ (this is the $L$ of \S\ref{pds:ss-setup}, bounded there; it is not the quantity $\max_i|2\varsigma_{\tau,i}|$ of Lemma~\ref{an:lem-sizes}), $C_{0,\tau}=
3n(L_\tau+2n)$ and $\delta_\tau=1/(4C_{0,\tau})$.

\begin{lemma}\label{pgq:local}
Let $\tau\in\auPl$, $g\in G_\tau$, $g'\in H_\tau$.
\begin{enumerate}
\item[(a)] $|c_\tau(g)|\le(2/\|g\|)^j$ for $0\le j\le h_V$, and $|c'_\tau(g')|\le(2/\|g'\|)^j$ for
$0\le j\le h_Z$.
\item[(b)] For real $s\ge8n^2$, $d_\tau\int_{G_\tau}\|g\|^{-s}dg\le8(5/\delta_\tau)^{2n^2}$.
\item[(c)] For $u\in G_\tau$ with $\|u-1\|\le\frac12$, $|c_\tau(gu)-c_\tau(g)|\le2C_{0,\tau}\|u-1\|$.
\item[(d)] $\aubar{c_\tau(g)}=c_\tau(g^{-1})$, $|c_\tau|\le1$, $|c'_\tau|\le1$.
\end{enumerate}
\end{lemma}
\begin{proof}
(a) Theorem~\ref{pds:decay}; for $c'_\tau$, Lemma~\ref{pco:decay} with Theorem~\ref{pco:FJ}, which apply
by Lemma~\ref{pgv:gap}. (b) Theorem~\ref{pds:tail}, as $s/(s-4n^2)\le2$. (c) Steps~1 and 2 of the
proof of Lemma~\ref{pds:lip} give $\|\pi_\tau(u)v_\tau-v_\tau\|\le C_{0,\tau}\|\log u\|\le2C_{0,\tau}
\|u-1\|$; and $c_\tau(gu)-c_\tau(g)=\langle v_\tau,\pi_\tau(g)(\pi_\tau(u)v_\tau-v_\tau)\rangle$.
(d) $c_\tau(g^{-1})=\langle\pi_\tau(g)v_\tau,v_\tau\rangle$; unit vectors.
\end{proof}

\begin{lemma}[Balls]\label{pgq:balls}
For real $s\ge8n^2$, $\int_{G_\tau}\|r\|^{-s}dr\le2\cdot16^{n^2}V_\tau$.
\end{lemma}
\begin{proof}
For $T\ge1$ let $S_T=\{r:\|r\|\le T\}$, a compact set. Let $r_1,\dots,r_\nu\in S_T$ with
$\|r_ir_j^{-1}-1\|>1$ for $i\ne j$. As $\|r_ir_j^{-1}-1\|\le\|r_i-r_j\|\,\|r_j^{-1}\|$, $\|r_i-r_j\|>1/T$,
so the balls of radius $1/(2T)$ about the $r_i$ in $\mathrm M_n(\C)\cong\R^{2n^2}$ are disjoint and
lie in the ball of radius $2T$ about $0$; hence $\nu\le(4T^2)^{2n^2}$. Take such a family maximal.
Then every $r\in S_T$ has $\|rr_i^{-1}-1\|\le1$ for some $i$, so $S_T\subseteq\bigcup_iB_1r_i$ with
$B_1=\{\|r-1\|\le1\}$, and $\pgqvol(S_T)\le(4T^2)^{2n^2}V_\tau$ ($G_\tau$ is unimodular). Since
$\|r\|^{-s}=\int_1^\infty sT^{-s-1}\mathbf 1[\|r\|\le T]\,dT$, Fubini gives $\int\|r\|^{-s}dr\le
16^{n^2}V_\tau\int_1^\infty sT^{4n^2-s-1}dT=16^{n^2}V_\tau\,s/(s-4n^2)$.
\end{proof}

\begin{lemma}[Slices]\label{pgq:slice}
For $g\in H_\tau$ and real $s\ge16n^2$, $\int_{G_\tau}\|g(r\oplus1)\|^{-s}dr\le2\cdot16^{n^2}V_\tau$.
\end{lemma}
\begin{proof}
The continuous function $r\mapsto\|g(r\oplus1)\|\ge\|r\|/\|g^{-1}\|$ tends to infinity, so it attains
its minimum $t_0$ at some $r_0$. For every $r$,
\[
\|rr_0^{-1}\|=\|(rr_0^{-1})\oplus1\|=\bigl\|(g(r_0\oplus1))^{-1}g(r\oplus1)\bigr\|\le t_0\,\|g(r\oplus1)\|
\le\|g(r\oplus1)\|^2 .
\]
Hence $\|g(r\oplus1)\|^{-s}\le\|rr_0^{-1}\|^{-s/2}$; integrate, use right invariance and
Lemma~\ref{pgq:balls} with $s/2$.
\end{proof}

\begin{lemma}[Averaging]\label{pgq:schur}
Let $F$ be a bounded measurable function on $G_\infty$ with $F(g)=\int_{G_\infty}e(z)F(gz)\,dz$ for all
$g$. Then $\int_{G_\infty}e(u)F(ur)\,du=c(r)F(1)$ for every $r\in G_\infty$.
\end{lemma}
\begin{proof}
The function $e$ is integrable (Lemma~\ref{pgq:local}(a),(b)), so all integrals below converge
absolutely. By the hypothesis at the point $ur$ and the substitution $z=r^{-1}u^{-1}w$,
\[
\int e(u)F(ur)\,du=\int F(w)\Bigl(\int e(u)\,e(r^{-1}u^{-1}w)\,du\Bigr)dw .
\]
Now $e(u)=d\langle\pi(u)v,v\rangle$ and $e(r^{-1}u^{-1}w)=d\,\aubar{\langle\pi(u)\pi(r)v,\pi(w)v\rangle}=
d\langle\pi(w)v,\pi(u)\pi(r)v\rangle$ ($\pi=\bigotimes\pi_\tau$, $v=\bigotimes v_\tau$). Schur's
orthogonality relations for the square-integrable irreducible representation $\pi$ of formal degree
$d$ (conjugate-linear in the first variable: $\int\langle a,\pi(u)b\rangle\aubar{\langle a',\pi(u)b'
\rangle}du=d^{-1}\langle a,a'\rangle\langle b',b\rangle$; see \eqref{par:eq-coef}) give
$\int e(u)e(r^{-1}u^{-1}w)\,du=d\langle v,\pi(r)v\rangle\langle\pi(w)v,v\rangle=c(r)e(w)$. So the left
side is $c(r)\int e(w)F(w)\,dw=c(r)F(1)$.
\end{proof}

\begin{lemma}[Variation of the kernel]\label{pgq:var}
Let $\mathsf D(g)=\prod_\tau(2/\|g_\tau\|)^{2A}$ on $G_\infty$, with $4A\le h_V$, $2A\ge8n^2$, and
$T_\tau=d_\tau\int_{G_\tau}(2/\|g\|)^{2A}dg$. Then $|e|\le d\,\mathsf D^2$,
$\int d\,\mathsf D=\prod_\tau T_\tau$, and for $u\in G_\infty$ with $\operatorname{dist}(u,1)\le\delta
\le\frac12$,
\[
\int_{G_\infty}|e(u^{-1}z)-e(z)|\,dz\le2\bigl(2fC_0\delta\bigr)^{1/2}\prod_\tau T_\tau,\qquad C_0=
\max_\tau C_{0,\tau}.
\]
\end{lemma}
\begin{proof}
$|e|\le d\mathsf D^2$ is Lemma~\ref{pgq:local}(a) with $j=4A$, and the integral of $d\mathsf D$ is a
product by Fubini. Let $a_\tau=\aubar{c_\tau(u_\tau^{-1}z_\tau)}$, $b_\tau=\aubar{c_\tau(z_\tau)}$. By
Lemma~\ref{pgq:local}(d),(c), $|a_\tau-b_\tau|=|c_\tau(z_\tau^{-1}u_\tau)-c_\tau(z_\tau^{-1})|\le
2C_0\delta$, and as all $|a_\tau|,|b_\tau|\le1$, telescoping gives $|\prod a_\tau-\prod b_\tau|\le
2fC_0\delta$. Also $|\prod a_\tau|\le\mathsf D(u^{-1}z)^2$ and $|\prod b_\tau|\le\mathsf D(z)^2$. Hence
$|\prod a-\prod b|^2\le2fC_0\delta\,(\mathsf D(u^{-1}z)+\mathsf D(z))^2$, and
$|e(u^{-1}z)-e(z)|\le d(2fC_0\delta)^{1/2}(\mathsf D(u^{-1}z)+\mathsf D(z))$. Integrate, using left
invariance.
\end{proof}

\subsection{Sizes}\label{pgq:ss-sizes}

Put
\[
A=4f(n+1)^2,\qquad\mathsf D_0=b^{x^{36}},\qquad\Pi=\prod_\tau\Bigl(T_\tau\,(4\beta_{\mathsf S})^{2A}\,
2\cdot16^{n^2}\,\beta_{\mathsf S}^{N_{\mathsf S}}\Bigr) .
\]

\begin{lemma}\label{pgq:sizes}
\begin{enumerate}
\item[(a)] $4A\le h_V$, $2A\le h_Z$, $A\ge16(n+1)^2$, $A\ge2f(n+1)+2$.
\item[(b)] $\mathsf D_0\ge\max\bigl(10,\ f,\ 20\,\Pi,\ 4(2fC_0)^{1/2}\prod_\tau T_\tau\bigr)$ and
$\mathsf D_0^{2n+2}(M\kappa_{\mathrm c}^2)^n\le M^{b/2}$.
\item[(c)] $\beta_{\mathsf S}\ge10^{12}$, $\beta_{\mathsf S}\ge2n+1+2b$, $\beta_{\mathsf S}\ge(n+1)^4
(6B)^{2m+2}$, and $\min(h_V,h_Z)\ge500(n+1)^3+10^{30}\bigl((2n+1)^2\log\beta_{\mathsf S}\bigr)^3$.
\end{enumerate}
\end{lemma}
\begin{proof}
We use Lemma~\ref{an:lem-sizes} (here $n=r$; the $n$ of (S7) is $m/2$): $x^5/2\le k\le r=n\le3x^5$, $m\le2\sqrt x$, $x^{11}\le f\le x^{12}$,
$\mathrm{rd}_E\le x^{12}$, $q_1\le x^{12}$, $d\le\log x$, $(\log b)^{100}\le x$, $x\ge128^{100}$,
$b\ge x^{10^{12}}$, $b^{20}\le M\le b^{22}$, $\Delta=2r^8$, $t=5r$, $B\le2b$; moreover $L_\tau\le2b$ (\S\ref{pds:ss-setup}).
(a) $4A\le16x^{12}(4x^5)^2\le x^{40}/2^{9}\le\Delta-t\le\min(h_Z,h_V)$ (as $r\ge x^5/2$, $h_Z\ge\Delta-t$ by Lemma~\ref{pgv:gap} and $h_V\ge\Delta$); the other two
hold as $f\ge4$. (b) $5/\delta_\tau=60n(L_\tau+2n)\le b^2$, so $T_\tau\le2^{2A}\cdot8b^{4n^2}\le
b^{2A+4n^2+3}$ by Lemma~\ref{pgq:local}(b). With $4\beta_{\mathsf S}\le b^{3m+2}$ and
$2\cdot16^{n^2}\le b^{n^2+1}$,
\[
\log_b\Pi\le f\bigl(2A(3m+3)+5n^2+4+60m(2n+1)^2\bigr)\le f\,m\bigl(48f(n+1)^2+244(n+1)^2\bigr)\le
50f^2m(n+1)^2\le x^{35},
\]
since $50x^{24}\cdot2x^{1/2}\cdot16x^{10}\le x^{35}$. So $20\Pi\le b^{x^{35}+1}\le\mathsf D_0$; the
bound for $4(2fC_0)^{1/2}\prod T_\tau\le b^2\prod T_\tau$ is the same computation, and $10,f\le
\mathsf D_0$. Next $\kappa_{\mathrm c}\le8x^{12}\cdot3x^5\cdot x^{12}+48(3x^5)^6x^{24}x^{13}+1\le x^{70}
\le b$, so $\log_b(\mathsf D_0^{2n+2}(M\kappa_{\mathrm c}^2)^n)\le(2n+2)x^{36}+24n\le x^{42}\le10b\le
\log_b M^{b/2}$. (c) $\log_b\beta_{\mathsf S}=3m$ and $b\ge x^{10^{12}}$ give the first two;
$(n+1)^4(12b)^{2m+2}\le b\cdot b^{1.01(2m+2)}\le b^{3m}$ as $m\ge200$ (Lemma~\ref{an:lem-sizes}(S7)).
Finally $10^{30}((2n+1)^2\,3m\log b)^3\le10^{35}(37x^{10}\cdot6x^{1/2}x^{1/100})^3\le x^{33}$ and
$500(n+1)^3\le x^{17}$, while $h_Z\ge\Delta-t\ge x^{40}/2^8-15x^5$.
\end{proof}

\noindent From now on $A$, $\mathsf D_0$ are as above, $\omega=(\omega_\tau)$ is the real point of $H$
given by (S), and for a real point $\alpha$ of $H$ we put
\[
\phi_\alpha(g)=\prod_\tau\aubar{c'_\tau\bigl(\alpha_\tau^{-1}g_\tau\omega_\tau\bigr)}=\frac{e'(\alpha^{-1}g
\omega)}{d'}\quad(g\in H_\infty),\qquad\pgqSl(\alpha)=\prod_\tau\bigl(2\|\alpha_\tau\|\|\omega_\tau\|\bigr)^{2A}
\,2\cdot16^{n^2}V_\tau .
\]

\begin{lemma}[Slices of the translated coefficient]\label{pgq:slicecoef}
Let $g\in H_\infty$ and $t>0$ with $t\le\prod_\tau\|g_\tau(r_\tau\oplus1)\|$ for all $r\in G_\infty$.
Then $r\mapsto\phi_\alpha(g(r\oplus1))$ is integrable on $G_\infty$ and
$\int_{G_\infty}|\phi_\alpha(g(r\oplus1))|\,dr\le\pgqSl(\alpha)\,t^{-A}$.
\end{lemma}
\begin{proof}
Write $X_\tau=g_\tau(r_\tau\oplus1)$. As $\|X_\tau\|\le\|\alpha_\tau\|\,\|\alpha_\tau^{-1}X_\tau\omega_\tau\|\,
\|\omega_\tau\|$ ($\|\alpha_\tau\|=\|\alpha_\tau^{-1}\|$, $\|\omega_\tau^{-1}\|=\|\omega_\tau\|$),
Lemma~\ref{pgq:local}(a) with $j=2A\le h_Z$ gives $|c'_\tau(\alpha_\tau^{-1}X_\tau\omega_\tau)|\le
(2\|\alpha_\tau\|\|\omega_\tau\|)^{2A}\|X_\tau\|^{-A}\|X_\tau\|^{-A}$. The product over $\tau$ of the
first factors $\|X_\tau\|^{-A}$ is at most $t^{-A}$; integrate the rest by Fubini and
Lemma~\ref{pgq:slice} ($A\ge16n^2$). The function is continuous, hence measurable.
\end{proof}

\subsection{The cusp}\label{pgq:ss-cusp}

\begin{lemma}\label{pgq:cusp}
Let $y\in\auAz$ and $S\ge\sup|y|$. If $\lambda(x)<\theta:=(\mathsf D_0^{-2}/(M
\kappa_{\mathrm c}^2))^n$, then $|y(x)|\le S/2$.
\end{lemma}
\begin{proof}
Let $\delta=\mathsf D_0^{-2}\le\frac12$ and $\operatorname{dist}(u,1)\le\delta$. By the reproducing
property at $xu$ and at $x$, and the substitution $z\mapsto u^{-1}z$,
$y(xu)-y(x)=\int(e(u^{-1}z)-e(z))y(xz)\,dz$; by Lemma~\ref{pgq:var} (its hypotheses hold by
Lemma~\ref{pgq:sizes}(a)) and Lemma~\ref{pgq:sizes}(b),
$|y(xu)-y(x)|\le S\cdot2(2fC_0)^{1/2}\mathsf D_0^{-1}\prod_\tau T_\tau\le S/2$. As $y$ is a form of
type $e$ fixed by $J_M\supseteq J_G(M)$, (C) gives $|y(x)|\le S/2$.
\end{proof}

\subsection{Poincar\'e series of deep level}\label{pgq:ss-poinc}

Fix $x\in G(\auA)$ and put $\hat x=x\oplus1$. Let $\mathcal R\subset\Gamma_H$ be a set of
representatives of $\Gamma_H/(\Gamma_G\oplus1)$ containing $1$; then $(\sigma,\delta)\mapsto
\sigma(\delta\oplus1)$ is a bijection $\mathcal R\times\Gamma_G\to\Gamma_H$. For $\sigma\in\Gamma_H$
let $\pgqcol(\sigma)\in\prod_\tau\C^{n+1}\cong\R^{2f(n+1)}$ be the archimedean coordinates of
$\hat x^{-1}\sigma\tilde e$. Let $\mathcal R^\flat$ be the set of $\sigma\in\mathcal R$ for which some
$k\in G(\auA_f)$ has $(\hat x^{-1}\sigma(k\oplus1))_f\in J_H(\pcb)$.

\begin{lemma}[Columns]\label{pgq:cols}
\begin{enumerate}
\item[(a)] For $\sigma\in\mathcal R$ and $h\in G(\auA)$, $g=\hat x^{-1}\sigma(h\oplus1)$ satisfies
$\prod_\tau\|g_\tau\|\ge|\pgqcol(\sigma)|/\sqrt f$; and $|\pgqcol(1)|=\sqrt f$.
\item[(b)] For distinct $\sigma,\sigma'\in\mathcal R^\flat$,
$|\pgqcol(\sigma)-\pgqcol(\sigma')|\ge M^{b/2}\lambda(x)$.
\end{enumerate}
\end{lemma}
\begin{proof}
(a) $g_\infty\tilde e=\hat x_\infty^{-1}\sigma\tilde e$, as $h\oplus1$ fixes $\tilde e$; so
$|\pgqcol(\sigma)|^2=\sum_\tau|g_\tau\tilde e|^2\le\sum_\tau\|g_\tau\|^2\le f\prod_\tau\|g_\tau\|^2$ (all
factors are $\ge1$). For $\sigma=1$ the coordinates are those of $\tilde e$, of norm $1$ at each
$\tau$.

(b) Let $v=\sigma\tilde e$, $v'=\sigma'\tilde e\in Z(E)$. If $v=v'$, then $\sigma^{-1}\sigma'$ fixes
$\tilde e$, hence preserves $V=\tilde e^\perp$, so $\sigma^{-1}\sigma'\in\Gamma_G\oplus1$ and
$\sigma=\sigma'$; thus $v\ne v'$. Choose $k,k'$ as in the definition of $\mathcal R^\flat$ and put
$g=(\hat x^{-1}\sigma(k\oplus1))_f\in J_H(\pcb)$. Then $\hat x_f^{-1}v=g\tilde e\equiv\tilde e$ modulo
$\pcb\widehat L_Z$, and likewise for $v'$, so $\hat x_f^{-1}(v-v')\in\pi_M^b\widehat L_Z$ and
$w=\pi_M^{-b}(v-v')$ is a non-zero element of $L_Z(\hat x)$. Since $|\tau(\pi_M)|=M^{1/2}$,
$|\pgqcol(\sigma)-\pgqcol(\sigma')|=|v-v'|_{\hat x}=M^{b/2}|w|_{\hat x}$. Finally $L_Z(\hat x)=
L_V(x)\oplus\auO_E\tilde e$ orthogonally in the coordinates, so for $w=w_0+a\tilde e\ne0$: if
$w_0\ne0$ then $|w|_{\hat x}\ge|w_0|_x\ge\lambda(x)$, and if $w_0=0$ then $a\in\auO_E\smallsetminus0$
and $\max_\tau|\tau(a)|\ge1$, as $\prod_\tau|\tau(a)|^2=|\mathrm N_{E/\Q}(a)|\ge1$.
\end{proof}

Let $y$ be a bounded measurable function on $[G]$ fixed on the right by $J_G(\pcb)$, $S\ge\sup|y|$,
$\pgqvol=\pgqvol(J_G(\pcb))$, and for a real point $\alpha$ of $H$ and $g\in H(\auA)$ let
$\Phi_\alpha(g)=\phi_\alpha(g_\infty)\mathbf 1_{J_H(\pcb)}(g_f)$. Put
\[
\mathsf P_x(\alpha)=\int_{[G]}\sum_{\gamma\in\Gamma_H}\Phi_\alpha\bigl(\hat x^{-1}\gamma(h\oplus1)\bigr)
\,y(h)\,dh,\qquad\mathsf c_x(\alpha)=\int_{G_\infty}\phi_\alpha(r\oplus1)\,y(xr)\,dr .
\]

\begin{lemma}[Unfolding]\label{pgq:unfold}
Let $\mathsf D\ge\max(10,\sqrt f)$ with $\mathsf D^2\le M^{b/2}\lambda(x)$. Then the series defining
$\mathsf P_x(\alpha)$ converges absolutely for almost all $h$, $\mathsf c_x(\alpha)$ converges
absolutely, and
\[
\bigl|\mathsf P_x(\alpha)-\pgqvol\cdot\mathsf c_x(\alpha)\bigr|\le\pgqvol\cdot\pgqSl(\alpha)\,S\,
\frac{10}{\mathsf D}.
\]
\end{lemma}
\begin{proof}
For $\sigma\in\mathcal R$ put $T_\sigma(h)=\Phi_\alpha(\hat x^{-1}\sigma(h\oplus1))y(h)$ on $G(\auA)$.
\emph{Support.} If $T_\sigma(h)\ne0$ then $(\hat x^{-1}\sigma(h\oplus1))_f\in J_H(\pcb)$, so
$\sigma\in\mathcal R^\flat$. \emph{One coset.} Let $\sigma\in\mathcal R^\flat$, with $k$ as in the
definition, and $g_0=\hat x^{-1}\sigma(k\oplus1)$. For $h\in G(\auA)$, $(g_0(h\oplus1))_f\in J_H(\pcb)$ if
and only if $h_f\in J_G(\pcb)$, as $g_{0,f}\in J_H(\pcb)$. Substituting $h\mapsto kh$,
\[
\int_{G(\auA)}|T_\sigma|\le S\,\pgqvol\int_{G_\infty}|\phi_\alpha(g_{0,\infty}(r\oplus1))|\,dr\le S\,\pgqvol\,
\pgqSl(\alpha)\,t_\sigma^{-A}
\]
by Lemma~\ref{pgq:slicecoef}, where $t_\sigma=\max(1,|\pgqcol(\sigma)|/\mathsf D)$ is admissible by
Lemma~\ref{pgq:cols}(a) ($g_0(r\oplus1)=\hat x^{-1}\sigma((kr)\oplus1)$, and $\sqrt f\le\mathsf D$).
\emph{The sum over cosets.} By Lemma~\ref{pgq:cols}, the points $\pgqcol(\sigma)$, $\sigma\in
\mathcal R^\flat$, are $\mathsf D^2$-separated and $|\pgqcol(1)|\le\mathsf D$; so $|\pgqcol(\sigma)|\ge
\mathsf D^2-\mathsf D\ge\mathsf D^2/2$ for $\sigma\ne1$, and the points $p_\sigma=\pgqcol(\sigma)/(\mathsf
D^2/2)$, $\sigma\in\mathcal R^\flat\smallsetminus1$, satisfy $|p_\sigma|\ge1$, $|p_\sigma-p_{\sigma'}|\ge1$.
By Lemma~\ref{pbc:pack}(b) in $\R^{2f(n+1)}$ with exponent $A\ge2f(n+1)+2$,
\[
\sum_{\sigma\in\mathcal R^\flat\smallsetminus1}\Bigl(\frac{\mathsf D}{|\pgqcol(\sigma)|}\Bigr)^A=\Bigl(
\frac2{\mathsf D}\Bigr)^A\sum|p_\sigma|^{-A}\le\Bigl(\frac{10}{\mathsf D}\Bigr)^A\le\frac{10}{\mathsf D}.
\]
Hence $\sum_{\sigma\in\mathcal R}\int|T_\sigma|<\infty$. \emph{Unfolding.} As $y(\delta h)=y(h)$, the
bijection $\mathcal R\times\Gamma_G\to\Gamma_H$ gives $\sum_\gamma\Phi_\alpha(\hat x^{-1}\gamma(h\oplus1))
y(h)=\sum_{\sigma}\sum_{\delta}T_\sigma(\delta h)$. By \eqref{pgq:eq-unfold} applied to $|T_\sigma|$ and
summed over $\sigma$, the double series converges absolutely for almost all $h$ and is integrable on
$[G]$; by dominated convergence and \eqref{pgq:eq-unfold} again,
$\mathsf P_x(\alpha)=\sum_{\sigma\in\mathcal R^\flat}\int_{G(\auA)}T_\sigma$. \emph{Main term.} For
$\sigma=1$ substitute $h\mapsto xh$: $\int T_1=\int_{G(\auA)}\phi_\alpha(h_\infty\oplus1)
\mathbf 1_{J_G(\pcb)}(h_f)y(xh)\,dh=\pgqvol\cdot\mathsf c_x(\alpha)$, as $y(xh_\infty h_f)=y(xh_\infty)$
for $h_f\in J_G(\pcb)$; $\mathsf c_x(\alpha)$ converges absolutely by Lemma~\ref{pgq:slicecoef} with
$g=1$, $t=1$. The other terms are bounded by $S\,\pgqvol\,\pgqSl(\alpha)(\mathsf D/|\pgqcol(\sigma)|)^A$,
whose sum is at most $S\,\pgqvol\,\pgqSl(\alpha)\,10/\mathsf D$.
\end{proof}

\begin{lemma}[The period vanishes]\label{pgq:zero}
Let $y$ be as above, and orthogonal in the sense of Theorem~\ref{pgq:V} at the
point $\omega$. Then $\mathsf P_x(\alpha)=0$ for every real point $\alpha$ of $H$.
\end{lemma}
\begin{proof}
Let $z=(\alpha^{-1},1)\hat x^{-1}\in H(\auA)$. For $h\in G(\auA)$ and $\gamma\in\Gamma_H$, as $\alpha$
and $\omega$ are real, $z\gamma(h\oplus1)(\omega,1)$ has archimedean part $\alpha^{-1}(\hat x^{-1}\gamma
(h\oplus1))_\infty\omega$ and finite part $(\hat x^{-1}\gamma(h\oplus1))_f$. So
$\Phi_\alpha(\hat x^{-1}\gamma(h\oplus1))=d'^{-1}e'((z\gamma g)_\infty)\mathbf 1_{J_H(\pcb)}((z\gamma g)_f)$
with $g=(h\oplus1)\omega$. By (Q) with $K=J_H(\pcb)$ there is a form $Q$ of type $e'$ fixed by
$J_H(\pcb)$ with $\sum_\gamma\Phi_\alpha(\hat x^{-1}\gamma(h\oplus1))=d'^{-1}\aubar{Q((h\oplus1)\omega)}$.
Hence $\mathsf P_x(\alpha)=d'^{-1}\,\aubar{\int_{[G]}\aubar{y(h)}Q((h\oplus1)\omega)\,dh}=0$.
\end{proof}

\subsection{Away from the cusp}\label{pgq:ss-bulk}

\begin{lemma}\label{pgq:bulk}
Let $y\in\auAz$ be orthogonal as in Theorem~\ref{pgq:V} at $\omega$,
$S\ge\sup|y|$, and $\lambda(x)\ge\theta$ \textup{(}Lemma~\ref{pgq:cusp}\textup{)}. Then $|y(x)|\le S/2$.
\end{lemma}
\begin{proof}
$y$ is fixed by $J_M\supseteq J_G(\pcb)$. By Lemma~\ref{pgq:sizes}(b), $\mathsf D=\mathsf D_0$ satisfies
$\mathsf D\ge\max(10,\sqrt f)$ and $\mathsf D^2\le M^{b/2}\theta\le M^{b/2}\lambda(x)$. So for every real
$\alpha$, Lemmas~\ref{pgq:unfold} and \ref{pgq:zero} give
\begin{equation}\label{pgq:eq-conv}
|\mathsf c_x(\alpha)|\le\pgqSl(\alpha)\,S\,\frac{10}{\mathsf D_0}.
\end{equation}
\emph{Translation.} Put $F(g)=y(xg)$ on $G_\infty$; it is bounded and $F(g)=\int e(z)F(gz)\,dz$. For
$u\in G_\infty$, $\phi_{(u\oplus1)\omega}(r\oplus1)=\phi_\omega((u^{-1}r)\oplus1)$, so the substitution
$r\mapsto ur$ gives $\mathsf c_x((u\oplus1)\omega)=\int_{G_\infty}\phi_\omega(r\oplus1)F(ur)\,dr$.
\emph{Averaging.} By
Fubini (the integrand is bounded by $S|e(u)||\phi_\omega(r\oplus1)|$, integrable) and
Lemma~\ref{pgq:schur},
\[
\int_{G_\infty}e(u)\,\mathsf c_x((u\oplus1)\omega)\,du=\int\phi_\omega(r\oplus1)\int e(u)F(ur)\,du\,dr
=\eta\,y(x),\qquad\eta=\int_{G_\infty}\phi_\omega(r\oplus1)c(r)\,dr .
\]
By Fubini, $\eta=\prod_\tau\int_{G_\tau}c_\tau(r)\aubar{c'_\tau(\omega_\tau^{-1}(r\oplus1)\omega_\tau)}dr$,
so $|\eta|\ge\prod_\tau V_\tau\beta_{\mathsf S}^{-N_{\mathsf S}}$ by (S). \emph{Bound.} As
$\|(u_\tau\oplus1)\omega_\tau\|\le\|u_\tau\|\beta_{\mathsf S}^{1/2}$ and $\|\omega_\tau\|\le
\beta_{\mathsf S}^{1/2}$, $\pgqSl((u\oplus1)\omega)\le\prod_\tau(2\|u_\tau\|\beta_{\mathsf S})^{2A}
2\cdot16^{n^2}V_\tau$; with $|e(u)|\le d\prod_\tau(2/\|u_\tau\|)^{4A}$ (Lemma~\ref{pgq:var}) and
\eqref{pgq:eq-conv},
\[
|e(u)\,\mathsf c_x((u\oplus1)\omega)|\le S\frac{10}{\mathsf D_0}\prod_\tau d_\tau\Bigl(\frac2{\|u_\tau\|}
\Bigr)^{2A}(4\beta_{\mathsf S})^{2A}\,2\cdot16^{n^2}V_\tau .
\]
Integrating over $u$, $|\eta|\,|y(x)|\le S\frac{10}{\mathsf D_0}\prod_\tau T_\tau(4\beta_{\mathsf S})^{2A}
2\cdot16^{n^2}V_\tau$, and dividing by the lower bound for $|\eta|$ (the $V_\tau$ cancel),
$|y(x)|\le S\cdot10\,\Pi/\mathsf D_0\le S/2$ by Lemma~\ref{pgq:sizes}(b).
\end{proof}

\subsection{Proof of Theorem~\ref{pgq:V}}\label{pgq:ss-proof}

The inputs (C), (Q), (S) hold by Theorems~\ref{pgk:C}, \ref{pgk:Q}, \ref{pgl:S}. Let $\omega$ be
given by (S) and let $y\in\auAz$ be orthogonal as in the theorem. The form $y$ is
bounded; let $S=\sup|y|$. For every $x\in G(\auA)$, $|y(x)|\le S/2$: by Lemma~\ref{pgq:cusp} if
$\lambda(x)<\theta$ and by Lemma~\ref{pgq:bulk} otherwise. Hence $S\le S/2$, $S=0$ and $y=0$.
\qed



%% ===== proofs/C-grid10.tex =====
% proofs/C-grid10.tex
% Continuation of proofs/C-grid8.tex (pgz:): the size bounds (B1)--(B3) of \S\ref{pgz:ss-b}.
% Labels start with pgt:. Conventions and macros: those of proofs/C-grid8.tex.
\providecommand{\auI}{\mathcal I}
\providecommand{\auGL}{\mathrm{GL}}
\providecommand{\auU}{\mathrm U}
\providecommand{\auO}{\mathcal O}
\providecommand{\auNm}{\operatorname{Nm}}
\providecommand{\pcb}{\mathfrak b}
\providecommand{\auPl}{\Phi}
\providecommand{\pkm}{\mathsf m}
\providecommand{\pkev}{\operatorname{ev}}
\providecommand{\pgsdu}{\mathrm{du}}
\providecommand{\aneu}{\mathrm{eu}_d}
\providecommand{\anrd}{\mathrm{rd}}
\providecommand{\pcnot}[1]{\par\smallskip{\raggedright\noindent\textsc{Not proved here.} [#1]\par}\smallskip}

\section{The size bounds on the last path}\label{pgt:sec}

This section proves the bounds of \S\ref{pgz:ss-b} that part (b) of Proposition~\ref{pgr:den} was
reduced to in Theorem~\ref{pgz:b}. The notation is that of \S\ref{pgz:sec}; $\vartheta=1$, so
$\pcm=\pcb=(\pi_M)^b$, and $G_0=\exp(x^{99}b^2)$.

\subsection{(B3): the number of classes for ranks up to $2n_m$}\label{pgt:ss-B3}

The proof of Proposition~\ref{pf:classes} uses the hypothesis $n\le n_m$ only in its last paragraph
(``Sizes''): the counting argument (the two elementary facts, the invariant $\mathsf I(g)$, and the
passage to levels) holds for every rank $n\ge1$, since Lemma~\ref{au:shortbasis}, which it uses with
$c_0=q_1d$, holds for every $n\ge1$. We record the statement for all ranks and redo the size
comparison for $n\le2n_m$.

\begin{proposition}[Classes in rank $\le 2n_m$]\label{pgt:classes}
Let $X=E^n$ be the positive definite standard space of rank $n\ge1$ (form $\sum x_i\aubar{x_i}$,
lattice $\widehat\auO_E^{\,n}$), $G=\auU(X)$, $\Gamma=G(E^+)$, $K^\circ$ the stabiliser of
$\widehat\auO_E^{\,n}$ in $G_f$, and for an integer $\mathsf G\ge1$ let $K(\mathsf G)=\{k\in K^\circ:
k\equiv1\bmod\mathsf G\}$. Then, with $C_{\mathrm{sb}}=6q_1d(fn)^2$ and $\mathsf B=n^2C_{\mathrm{sb}}^2
\cdot2f\sqrt d$,
\begin{equation}\label{pgt:eq-classes}
\#\bigl(\Gamma\backslash G_f/J\bigr)\le(2\mathsf B+1)^{4fn^2}\,\mathsf G^{2fn^2}
\end{equation}
for every subgroup $J\supseteq K(\mathsf G)$ of $G_f$. If $n\le2n_m$ and $\mathsf G=M^b$, the right
side is at most $\exp(x^{26}b^2)\le G_0$.
\end{proposition}
\begin{proof}
\emph{The bound \eqref{pgt:eq-classes}.} This is the proof of Proposition~\ref{pf:classes}, which
we check uses no restriction on $n$. Fact (i) ($\auA_{E,f}=E+\widehat\auO_E$ with a representative of
conjugates at most $\mathsf B_0=2f\sqrt d$) and fact (ii) (at most $(2B+1)^{2f}$ elements of $\auO_E$
with all conjugates at most $B$) do not involve $n$. For $g\in G_f$, Lemma~\ref{au:shortbasis} (for
the given $n$, with $c_0=q_1d$, which lies in the different of $E$ as in the proof of
Proposition~\ref{au:short}) gives $Y\in\auGL_n(E)$ with $A=g^{-1}Y$ integral and entries with
conjugates at most $C_{\mathrm{sb}}$; fact (i) gives $X'\in M_n(E)$ with $A^{-1}-X'$ integral and
entries bounded by $\mathsf B_0$. As $g^*g=1$, $Y^*Y=A^*A$ and $Y^*YX'=A^*(1-A(A^{-1}-X'))$ are in
$M_n(\auO_E)$, with conjugates at most $nC_{\mathrm{sb}}^2$ and $n^2C_{\mathrm{sb}}^2\mathsf B_0=\mathsf B$;
by fact (ii) the pair $\mathsf I(g)=(Y^*Y,Y^*YX')$ takes at most $(2\mathsf B+1)^{4fn^2}$ values. If
$\mathsf I(g)=\mathsf I(g')$ (with $X'_g$, $X'_{g'}$ the matrices chosen for $g$, $g'$), then $Y^*Y=Y'^*Y'$ is invertible and the second components give $X'_g=X'_{g'}=:X'$; so $A'X'=1-A'(A'^{-1}-X')$ and $A'(A^{-1}-X')$ are integral. Moreover $\gamma=YY'^{-1}$ satisfies $\gamma^*\gamma=1$, so $\gamma\in\Gamma$,
and $k=g'^{-1}\gamma^{-1}g=A'A^{-1}=A'(A^{-1}-X')+A'X'$ is integral and unitary, so $k\in K^\circ$ and
$g=\gamma g'k$. Hence $\#(\Gamma\backslash G_f/K^\circ)\le(2\mathsf B+1)^{4fn^2}$. Finally
$\Gamma\backslash G_f/(J\cap K^\circ)\to\Gamma\backslash G_f/J$ is onto and the fibres of
$\Gamma\backslash G_f/(J\cap K^\circ)\to\Gamma\backslash G_f/K^\circ$ have at most $[K^\circ:K(\mathsf
G)]\le\#M_n(\auO_E/\mathsf G)=\mathsf G^{2fn^2}$ elements, since $k\mapsto k\bmod\mathsf G$ embeds
$K^\circ/K(\mathsf G)$ in $M_n(\auO_E/\mathsf G)$ and $[E:\Q]=2f$.

\emph{Sizes for $n\le2n_m$.} By Lemma~\ref{an:lem-sizes}: $f\le q_1\le x^{12}$ (S3), $d\le\log x$
((S1): $\log x\ge10^{14}\cdot2d$), $\log M\le22\log b$ (S5), $\log b\le x^{1/100}$ (S1); and
$n_m\le40x^6$ (Lemma~\ref{pz:sizes}(b)). So $n\le80x^6$ and
\[
C_{\mathrm{sb}}=6q_1d(fn)^2\le6\cdot x^{12}\cdot\log x\cdot x^{24}\cdot6400x^{12}\le x^{50}.
\]
Since $n\le C_{\mathrm{sb}}$ and $\mathsf B_0=2f\sqrt d\le C_{\mathrm{sb}}$, we have $\mathsf B\le
C_{\mathrm{sb}}^5$ and $2\mathsf B+1\le C_{\mathrm{sb}}^6$, so $\log(2\mathsf B+1)\le300\log x$. Also
$fn^2\le x^{12}\cdot6400x^{12}=6400x^{24}$. Hence
\[
4fn^2\log(2\mathsf B+1)\le25600\,x^{24}\cdot300\log x\le x^{26},\qquad
2fn^2\log(M^b)\le12800\,x^{24}\cdot22\,b\log b\le x^{25}b^2,
\]
and the sum is at most $x^{26}b^2\le x^{99}b^2$.
\end{proof}

\begin{corollary}[(B3)]\label{pgt:B3}
Let $\vartheta=1$. For $z\in\{0,1\}$ and $a\le m$, $\#(\Gamma^z_a\backslash G^z_{a,f}/J^z_a)\le
\exp(x^{26}b^2)\le G_0$; likewise $\#(\Gamma_Y\backslash G_{Y,f}/J_Y)\le G_0$.
\end{corollary}
\begin{proof}
The node $p^z_a=(z,a,m)$ has rank $N^z_a\le2n_m$ (Lemma~\ref{pgs:boxes}). Its space is a positive
definite standard space: for $(z,a)\ne(0,0)$ it is $E^{N^z_a}$ with the standard lattice and the
diagonal form with the signs $\epsilon_{p^z_a}$ (Definition~\ref{pgv:def-spaces}; the node $(1,0,0)$
is not on $\pgrlast$ since $m\ge1$), and all these signs are $+1$, since the number of negative signs
at $(z,a,m)$ is $\lfloor z/2\rfloor=0$ (Lemma~\ref{pgv:count}(b)); for $(z,a)=(0,0)$ it is $Y_m$,
which via its frame is the positive definite standard space of rank $n_m-1$ with the lattice
$\widehat\auO_E^{\,n_m-1}$ (\S\ref{ps:ss-Y}, and Lemma~\ref{pz:sizes}(a) for the same identification).
So Proposition~\ref{pgt:classes} applies with $G=G^z_a$, $\Gamma=\Gamma^z_a$, $K^\circ=K_{p^z_a}$.
For $\vartheta=1$ every level on $\pgrlast$ is $J(\pcb)$, the principal congruence subgroup of
$K_{p^z_a}$ of level $\pcb$ (\S\ref{pgw:ss-last}, \S\ref{pgz:sec}). As $\pcb=(\pi_M)^b$ and
$\pi_M\aubar{\pi_M}=M$, $M^b\in\pcb$, so $K(M^b)\subseteq J(\pcb)$, and \eqref{pgt:eq-classes} with
$\mathsf G=M^b$ gives the bound. The node $Y=Y^m_m$ (not the space $Y_m$ above) is $p^0_m$, so the last assertion is the case $a=m$,
$z=0$.
\end{proof}

\begin{remark}\label{pgt:rem-B3}
Corollary~\ref{pgt:B3} is (B3) of \S\ref{pgz:ss-b}. The statement of Proposition~\ref{pf:classes}
may be widened from $n\le n_m$ to all $n$ for the bound \eqref{pgt:eq-classes}, without change of
proof; only its final size estimate $\exp(x^{30})$ uses $n\le n_m$ and $\mathsf G=MN$.
\end{remark}

\subsection{(B2): an integral branching matrix for an arbitrary interlacing pair}\label{pgt:ss-B2}

We use the tensor model of \S\ref{pk:ss-branch} and the coordinates of \S\ref{pgs:ss-weyl}. As
there, all statements are multiplicative over $\tau\in\auPl$; we write the constructions for one
$\tau$ and take tensor products over $\tau$ at the end. Put $N=n+1$ and regard $\auGL_n\subset
\auGL_N$ as $g\mapsto g\oplus1$, the new coordinate being the last one, and $\auGL_n\times\auGL_1
\subset\auGL_N$ as $(g,s)\mapsto g\oplus s$. In the model of \S\ref{pk:ss-branch} columns of height
equal to the rank are allowed; for $\auGL_n$ a column of height $n$ contributes the factor
$\Lambda^n R^n=Re_{[n]}$, on which $g$ acts by $\det g$, and the vector $e_{[n]}$ has norm $1$. So for a
diagram $\nu$ with $c$ columns of height $n$ and $\nu^-$ the diagram of its other columns,
$\mathbb T_\nu=\mathbb T_{\nu^-}\otimes(Re_{[n]})^{\otimes c}$ identifies $\mathsf C_\nu$ with
$\mathsf C_{\nu^-}\otimes\det^c$ isometrically, preserving integral vectors and products of minors
(a product of minors for $\nu$ is one for $\nu^-$ times $\det^c$). Hence the weight $(\nu^-,e+c)$ of
\S\ref{pgs:ss-weyl} may be realised on $\mathsf C_\nu$ twisted by $\det^e$, with its lattice
\eqref{pgs:eq-lattice} and its invariant form.

\subsubsection*{Interlacing}
Let $\lambda=(\kappa,e)$ be a weight of $\auGL_N$ at $\tau$, $\mu$ its diagram (columns of heights
$<N$), so the highest weight is $\mu+e$ with $\mu$ padded to $N$ entries. Let $\lambda_0$ be a weight
of $\auGL_n$ whose highest weight $\mu_0'$ interlaces it: $\mu_1+e\ge\mu'_{0,1}\ge\mu_2+e\ge\dots\ge
\mu'_{0,n}\ge\mu_N+e=e$. Then $\nu=\mu_0'-e$ is a diagram with at most $n$ rows and $\mu_{i+1}\le
\nu_i\le\mu_i$, i.e.\ $\mu/\nu$ is a horizontal strip: every column of $\mu$, of height $h$, either
\emph{keeps} its height in $\nu$ or is \emph{lost} and has height $h-1$ in $\nu$ (columns of height $1$
that are lost disappear); the columns of $\nu$ are exactly these. Let $\ell$ be the number of lost
columns. Then $\lambda_0=(\nu^-,e+c)$ with $c$ the number of kept columns of height $n$. Note
$|\nu^-|\le|\nu|\le|\mu|$.

\begin{lemma}[The branching map]\label{pgt:br}
In this situation let $\varepsilon\colon\mathbb T_\nu(\Z^n)\to\mathbb T_\mu(\Z^N)$ send $e_U$ to the
tensor with factor $e_{U_{\mathfrak c}}$ in a kept column and $e_N\wedge e_{U_{\mathfrak c}}$ in a lost
one, and put $i=\pkm\circ\varepsilon\circ(\pkm^\nu|_{\mathsf C_\nu})^{-1}\colon W_\nu\to W_\mu$ and
$\mathrm{br}=H_{\nu}\,i^*\colon W_\mu\to W_\nu$, $H_\nu$ the hook number of Lemma~\ref{pk:hook} and
$i^*$ the adjoint for the invariant forms. Then:
\begin{enumerate}
\item[(a)] $\mathrm{br}\circ\rho_\mu(g\oplus s)=s^{\ell}\,\rho_\nu(g)\circ\mathrm{br}$ for $g\in\auGL_n$,
$s\in\auGL_1$, over every field of characteristic $0$ (through the integral representation matrices);
\item[(b)] $\|\mathrm{br}\,v\|\le H_\nu\|v\|$ for the invariant norms;
\item[(c)] $\mathrm{br}$ maps $W_\mu(\Z)$ into $W_\nu(\Z)$, i.e.\ its matrix in the bases $(e_T)$ is
integral;
\item[(d)] $\mathrm{br}\ne0$, and $H_\nu\le|\nu|^{3|\nu|}$.
\end{enumerate}
\end{lemma}
\begin{proof}
This is the proof of Proposition~\ref{pk:branching}, in which the specific structure of the pair
$(\lambda(X_m),\lambda^\circ)$ was used only to take a smaller denominator; we take the full hook
number $H_\nu$ instead, and then no structure is needed.

\emph{The map $i$.} $\varepsilon$ maps the orthonormal basis $(e_U)$ injectively to $\pm$ elements of
the orthonormal basis of $\mathbb T_\mu$ (the sets $U_{\mathfrak c}\subseteq[n]$ do not contain $N$),
so it is isometric and integral, with $\varepsilon^t$ integral of norm $\le1$. For $x\in M_n(R)$,
whose columns are regarded in $R^N$ with last coordinate $0$, and $(g,s)$, one has
$\varepsilon((\Lambda g)u)=s^{-\ell}(\Lambda(g\oplus s))\varepsilon(u)$ factor by factor (a lost column
gets $s e_N\wedge$), so $\varepsilon$ is $\auGL_n$-equivariant; as in \S\ref{pk:ss-branch}, $\pkm$ factors through $\pi_\mu$, and under the
identifications \eqref{pk:eq-dual} $i=\pi_\mu\varepsilon|_{\mathsf C_\nu}$ and
$i^*=\pi_\nu\varepsilon^t|_{\mathsf C_\mu}$. The projections $\pi_\mu,\pi_\nu$ commute with the
actions of the unitary groups (the subspaces $\mathsf C$ are invariant and the forms are), so
$i^*\rho_\mu(g\oplus s)=s^\ell\rho_\nu(g)i^*$ for unitary $g,s$; both sides are polynomial in the
entries of $g,s$ and $\det g^{-1},s^{-1}$ with rational coefficients, so the identity holds on
$\auGL_n\times\auGL_1$ by Zariski density, and then over every field of characteristic $0$. This is
(a). (b): $\|i^*\|\le\|\pi_\nu\|\,\|\varepsilon^t\|\le1$, the invariant forms being the restrictions of
the standard forms (\S\ref{pk:ss-branch}).

\emph{The kernel.} For $x\in M_n(R)$ and $y\in M_N(R)$ let $\mathcal K(x,y)=(\varepsilon\,\pkev^\nu_x,
\pkev_y)$. By \eqref{pk:eq-dual}, as in \eqref{pk:eq-kernelbr},
$\mathcal K(x,y)=(\pkev^\nu_x,\pi_\nu\varepsilon^t\pkev_y)$ and $(\pi_\nu\varepsilon^t\pkev_y,e_{T'})=
i(b^\nu_{T'})(y)$ for every filling $T'$ of $\nu$. All columns of $\mu$ have height $\le n$, so the
last column of $y$ does not occur in $\pkev_y$. Let first the last row of $y$ be $(\mathsf
t,0,\dots,0)$ and $g\in M_n(R)$ the matrix of the first $n$ rows and columns of $y$. In a kept column
of height $h$ the factor of $\mathcal K$ is $\det(x_a\cdot g_b)_{a,b\le h}=\Delta_{[h],[h]}(g^tx)$; in a
lost column of height $h$ it is the determinant with first row $(\mathsf t,0,\dots,0)$ and rows
$(x_a\cdot g_b)_{b\le h}$, $a<h$, that is $\mathsf t\,\Delta_{\{2,\dots,h\},[h-1]}(g^tx)$. So
$\mathcal K(x,y)=\mathsf t^{\,\ell}\mathsf P(g^tx)$ with $\mathsf P=b^\nu_U$ the product of minors of a
filling $U$ of $\nu$. By Lemma~\ref{pk:hook} and \eqref{pk:eq-dual}, $H_\nu\mathsf P=\pkm^\nu(w_0)$ with
$w_0\in\mathsf C_\nu\cap\mathbb T_\nu(\Z)$, so $H_\nu\mathcal K(x,y)=\mathsf t^{\,\ell}(w_0,
\pkev^\nu_{g^tx})=(\mathsf t^{\,\ell}gw_0,\pkev^\nu_x)$. For general $y$, over $R=\Z[y_{ab}]$, let $y'$
have the columns $y'_1=y_1$, $y'_b=y_{N1}y_b-y_{Nb}y_1$ $(b\ge2)$; its last row is
$(y_{N1},0,\dots,0)$ and $\pkev_{y'}=y_{N1}^{\,\mathsf M}\pkev_y$ with $\mathsf M=\sum_{\mathfrak
c}(h_{\mathfrak c}-1)$. Exactly as in the proof of Proposition~\ref{pk:branching},
$y_{N1}^{\,\mathsf M}H_\nu\pi_\nu\varepsilon^t\pkev_y=y_{N1}^{\,\ell}g'w_0$ ($g'$ the first $n$ rows and
columns of $y'$), both being vectors of $\mathsf C_\nu\otimes\Q(y)$ with the same pairings against all
$\pkev^\nu_x$; the right side has coordinates in $R$. Hence $y_{N1}^{\,\mathsf M}H_\nu
i(b^\nu_{T'})\in\Z[y_{ab}]$, and since $y_{N1}$ is prime in $\Z[y_{ab}]$ and $H_\nu i(b^\nu_{T'})\in
\Q[y_{ab}]$, $H_\nu i(b^\nu_{T'})\in\Z[y_{ab}]$; by \eqref{pk:eq-sat}, $H_\nu i(b^\nu_{T'})\in\Sigma_\mu$.

\emph{(c).} For $u\in W_\mu(\Z)=\Sigma_\mu^\#$ and every $T'$, $\langle\mathrm{br}\,u,b^\nu_{T'}\rangle=
\langle u,H_\nu i(b^\nu_{T'})\rangle\in\Z$; the coordinate of $\mathrm{br}\,u$ on $e_{T'}$ is this
number (the basis $(e_T)$ being dual to $(b_T)$, \S\ref{pgs:ss-weyl}), so $\mathrm{br}\,u\in
W_\nu(\Z)$. (The products of minors $b^\nu_{T'}$ for $\nu$ and for $\nu^-$ differ by $\det^c$, which
does not affect this.)

\emph{(d).} Let $y\in M_N(\C)$ have first $n$ rows and columns $1_n$ and last row $(1,0,\dots,0)$, so
$\mathsf t=1$, $g=1$ and $\mathcal K(x,y)=\mathsf P(x)$. As $\mathsf P$ is a product of minors of a
generic matrix, it is a non-zero polynomial, so $\mathsf P(x)\ne0$ for some $x$. Then
$(\pkev^\nu_x,i^*\pkev_y)=\mathcal K(x,y)\ne0$ (by the paragraph on the kernel, since $\pkev_y\in\mathsf C_\mu$ by the definition of $\mathsf C_\mu$ and $i^*=\pi_\nu\varepsilon^t|_{\mathsf C_\mu}$), so $i^*\ne0$ and $\mathrm{br}\ne0$. The bound on
$H_\nu$ is Lemma~\ref{pk:hook}.
\end{proof}

\begin{proposition}[(B2)]\label{pgt:B2}
Let $\lambda^0$, $\lambda^1$ be the weights of $W^0_m$ (rank $n$) and $W^1_m$ (rank $N=n+1$), and
suppose there is a non-zero invariant form on $W^0_m\times W^1_m$ (in the sense of
Lemma~\ref{pgw:branch}(a)). Then there is a non-zero integer matrix $B'$ in the bases $(e_T)$ with
$B'(\rho_0(g)w_0,\rho_1(g\oplus1)w_1)=B'(w_0,w_1)$ for $g\in\prod_\tau\auGL_n$ and
\[
|B'(v',v)|\le L^{3L}\,\|v'\|\,\|v\|\qquad\text{for all }v'\in W^0_m,\ v\in W^1_m\ (\text{over }\C),
\]
for every integer $L$ at least the total number of boxes $\sum_\tau|\mu^1_\tau|$ of $\lambda^1$; in
particular for $L=\lfloor x^{19}b\rfloor$ (Lemma~\ref{pgs:boxes}). The same bound holds after
applying any $\sigma\in\auAut(\C)$ to the coordinates of $v'$, $v$, and it holds for the matrix $B$ of
Definition~\ref{pgz:def-choice}(b). As $L^{3L}=(L^L)^3$, this is (B2) of \S\ref{pgz:ss-b}.
\end{proposition}
\begin{proof}
\emph{Interlacing.} A non-zero invariant form is a non-zero element of
$\auHom_{\prod_\tau\auU(n)}(W^1_m,(W^0_m)^*)$, so $(W^0_m)^*=\bigotimes_\tau(W^0_{m,\tau})^*$ occurs in
the restriction of $W^1_m=\bigotimes_\tau W^1_{m,\tau}$ to $\prod_\tau\auU(n)$; hence, at each $\tau$,
$(W^0_{m,\tau})^*$ occurs in the restriction of $W^1_{m,\tau}$ to $\auU(n)$, and by the branching law
of Weyl (\cite[Theorem~8.1.1]{GoodmanWallach}) its highest weight interlaces that of
$W^1_{m,\tau}$. Let $\lambda_0$ be the weight of $(W^0_m)^*$; by the proof of
Lemma~\ref{pgs:pairing} its doubled parameter is the negative of that of $\lambda^0$. At each $\tau$
we are in the situation of Lemma~\ref{pgt:br} with $\lambda=\lambda^1_\tau$ and $\lambda_0$ at $\tau$;
let $\mathrm{br}=\bigotimes_\tau\mathrm{br}_\tau$, an integral matrix $W^1_m\to W_{\lambda_0}$
(Lemma~\ref{pgt:br}(c), tensor products of integral matrices in the tensor product bases). With the
twists by $\det^{e}$ of $\lambda^1$ and $\det^{e+c}$ of $\lambda_0$ at each $\tau$, Lemma~\ref{pgt:br}(a)
with $s=1$ gives $\mathrm{br}\,\rho_1(g\oplus1)=\rho_{\lambda_0}(g)\,\mathrm{br}$, since $\det(g\oplus1)
=\det g$ and the kept columns of height $n$ contribute $\det(g)^c$.

\emph{The matrix.} Let $A=G_{\lambda^0}J$ be the duality pairing of Lemma~\ref{pgs:pairing} between
weight $\lambda^0$ (in the role of $\lambda'$) and $\lambda_0$ (in the role of $\lambda$), and put
$B'(v',v)=\langle v',\mathrm{br}\,v\rangle_\pgsdu={}^tv'A\,\mathrm{br}\,v$. It is an integer matrix.
By (P1), $B'(\rho_0(g)v',\rho_1(g\oplus1)v)=\langle\rho_0(g)v',\rho_{\lambda_0}(g)\mathrm{br}\,v
\rangle_\pgsdu=B'(v',v)$. It is non-zero since $\mathrm{br}\ne0$ (Lemma~\ref{pgt:br}(d)) and $A$ is
invertible ($J$ is, and $G_{\lambda^0}$ is positive definite).

\emph{The bound.} By (P3), $\langle v',u\rangle_\pgsdu=\langle u^\flat,v'\rangle$ with $\|u^\flat\|=
\|u\|$, so by Cauchy--Schwarz for the invariant hermitian form and Lemma~\ref{pgt:br}(b),
\[
|B'(v',v)|\le\|v'\|\,\|\mathrm{br}\,v\|\le\Bigl(\prod_\tau H_{\nu_\tau}\Bigr)\|v'\|\,\|v\|\le
\prod_\tau L_\tau^{3L_\tau}\,\|v'\|\,\|v\|\le L^{3L}\|v'\|\,\|v\|,
\]
with $L_\tau=|\nu_\tau|\le|\mu^1_\tau|$ and $\sum_\tau L_\tau\le L$ (the norm of a tensor product of
operators being the product of the norms). For $\sigma\in\auAut(\C)$: $B'$ has integer entries, so
$\sigma(B'(v',v))=B'(\sigma v',\sigma v)$ and the bound applies to $\sigma v',\sigma v$. For $B$: by
Lemma~\ref{pgw:branch}(a) $B'=cB$ with $c$ rational, and as $B$ is primitive and $B'$ integral, $c$ is
a non-zero integer; so $|B(v',v)|\le|B'(v',v)|$. Finally $\sum_\tau|\mu^1_\tau|\le40fn_mb\le x^{19}b$ by
Lemma~\ref{pgs:boxes}.
\end{proof}

\begin{remark}\label{pgt:rem-B2}
If there is no non-zero invariant form, $B=0$ by Definition~\ref{pgz:def-choice}(b) and the bound
used in the proof of Theorem~\ref{pgz:b} holds trivially; so Proposition~\ref{pgt:B2} is all that
proof needs. The constant obtained is $L^{3L}=(L^L)^3$, the constant asked for in (B2).
\end{remark}

\subsection{(B1): the operator bound for one jet kernel}\label{pgt:ss-B1}

Proposition~\ref{au:kernel} is stated for the kernels of the chain $X_j\to X_{j+1}$, with their row
lengths, and Lemma~\ref{plf:KS}(KS2) needs $M^2\in\pcm$, which fails for $\pcm=\pcb=(\pi_M)^b$ when
$b>2$. Their proofs, however, give an operator bound for any jet kernel between standard definite
spaces of increasing ranks with a translate of the vacuum as datum, and in the archimedean bound the
modulus plays no role: an indicator datum is an average of translates of the vacuum with coefficients
of absolute value $1$ (after normalisation). We write this out.

\begin{lemma}[Jets of translates of the vacuum]\label{pgt:jet}
Let $X=E^n$, $X'=E^{n'}$ be positive definite standard spaces, $n<n'$, $W=\auHom(X,X')$,
$D=fnn'$. Let $\nu=(\nu_\tau)$ be a diagram with at most $n$ rows at each $\tau$, $e=\sum_\tau|\nu_\tau|$,
and $S_\nu(z)$, $z\in\prod_\tau M_{n'\times n}(\C)$, the matrix of the Schur functor
$W_\nu(\C^n)\to W_\nu(\C^{n'})$ of \S\ref{ps:ss-kernel} (tensor product over $\tau$), in the bases
$(e_T)$. Let $\phi=\psrho^\circ(x)\omega(k)\psvac^\circ$ with $x\in W(\auA_f)$, $k\in\auU(W)(\auA_f)$. Then
for all $\sigma\in\auAut(\C)$ and all $\xi\in W_\nu(\C^n)$,
\[
\bigl\|\sigma\bigl(S_\nu(\partial)[\Theta_\phi(s_0\,\cdot)](0)\,\xi\bigr)\bigr\|\le G_1\,\|\sigma\xi\|,\qquad
G_1=\Omega^{-e/2}(2e)^{e/2}\,4^{D}\exp\bigl(C_0(q_1nn'+1)^2\bigr),
\]
$C_0=C_0(q_0,d)$ the constant of Lemma~\ref{ps:conj}{\rm(ii)} and $(2e)^{e/2}=1$ if $e=0$.
\end{lemma}
\begin{proof}
This is Steps 2--4 of the proof of Proposition~\ref{au:kernel} at the scale $s_0$ instead of
$\pi_Ms_0$, with a cruder treatment of the Schur matrix. As there we may assume $\sigma|_K=1$: if
$\sigma'=c\circ\sigma$ ($c$ complex conjugation), then $|\sigma'w|=|\sigma w|$ for numbers and
$\|\sigma'v\|=\|\sigma v\|$ for vectors, the Gram matrices being real.

Let $\eta$ be a covector on $W_\nu(\C^{n'})$ with $\|\sigma\eta\|\le1$ for the dual norm, and
$Q(z)=\eta\cdot S_\nu(z)\xi$, a polynomial homogeneous of degree $e$. By Lemma~\ref{ps:conj}(i) and
\eqref{ps:eq-bargmann} for $\Lambda_\sigma\phi$ at the scale $t_\sigma$ (Step~2 of the proof of
Proposition~\ref{au:kernel}),
\[
\sigma\bigl(Q(\partial)[\Theta_\phi(s_0\,\cdot)](0)\bigr)=\int Q^\sigma(t_\sigma\bar y)\,e^{-|y|^2/2}\,
\theta\bigl(\psrho_\infty(y)(1\otimes\Lambda_\sigma\phi)\bigr)\,d\mu(y),\qquad d\mu(y)=\pi^{-D}dy,
\]
$Q^\sigma$ obtained by applying $\sigma$ to the coefficients and $\varpi_\sigma$ to the variables. The
entries of $S_\nu$ have integer coefficients, so $Q^\sigma(z)=\sigma\eta\cdot S_\nu(z')\sigma\xi$ with
$z'$ a permutation of the components of $z$, and by \eqref{ps:eq-schurnorm} (at each $\tau$, and the
operator norm of a tensor product being the product of the norms)
\[
|Q^\sigma(z)|\le\|\sigma\xi\|\prod_{\tau,p}s_{\tau,p}^{\,l_{\tau,p}}\le\|\sigma\xi\|\,|z|^{e},
\]
where $s_{\tau,p}\le|z_\tau|_{\mathrm{HS}}\le|z|$ are the singular values and $l_{\tau,p}$ the rows of
$\nu_\tau$, $\sum_{\tau,p}l_{\tau,p}=e$. For $z=t_\sigma\bar y$, $|z|^2=\Omega^{-1}|y|^2$
(Lemma~\ref{ps:conj}). By Lemma~\ref{ps:conj}(ii) the theta value is at most $\exp(C_0(q_1nn'+1)^2)$.
Using $u^{e/2}e^{-u/4}\le(2e)^{e/2}e^{-e/2}\le(2e)^{e/2}$ for $u\ge0$ (maximum at $u=2e$) with $u=|y|^2$,
\[
\int|y|^{e}e^{-|y|^2/2}\,d\mu(y)\le(2e)^{e/2}\int e^{-|y|^2/4}\,d\mu(y)=(2e)^{e/2}\,4^{D}.
\]
Hence $|\sigma(\eta\cdot S_\nu(\partial)[\Theta_\phi(s_0\,\cdot)](0)\xi)|\le G_1\|\sigma\xi\|$; since $S_\nu$ has
integer coefficients, the left side is $|\sigma\eta\cdot\sigma(S_\nu(\partial)[\cdots](0)\xi)|$, and
taking the supremum over $\eta$ gives the claim.
\end{proof}

\begin{proposition}[(B1)]\label{pgt:B1}
Let $\vartheta=1$, $e^z_a=(p,q)$ an edge of $\pgrlast$, $X\in\pcMat_{q\times p}(\pcR_\pcb)$, $x\in
G_{p,f}$, $h\in G_{q,f}$, $\sigma\in\auAut(\C)$ and $\xi\in W^z_{a-1}$. With the kernels of
Definition~\ref{pgz:def-choice}(a),
\[
\bigl\|\sigma\bigl(\pgwK^z_a(\omega_f(x,h)[\delta_X]_f)\,\xi\bigr)\bigr\|\le\exp(x^{52}b\log b)\,
\|\sigma\xi\|\le G_0\,\|\sigma\xi\| .
\]
This is (B1) of \S\ref{pgz:ss-b}.
\end{proposition}
\begin{proof}
\emph{The spaces.} The nodes $p=p^z_{a-1}$, $q=p^z_a$ carry positive definite standard spaces
(proof of Corollary~\ref{pgt:B3}) of ranks $N_p<N_q\le2n_m$ (Lemma~\ref{pgs:boxes}; the rank grows by
$t+1$ along $h_z$). By Definition~\ref{pgz:def-choice}(a), $\pgwK^z_a(\Phi)=c_{\mathrm{vac}}^{\,N_pN_q}
S_\nu(\partial)[\Theta_\Phi(s_0\,\cdot)](0)$ with $S_\nu$ the Schur matrix of the diagram $\nu$ of the
edge, of degree $e=e(\nu)\le40fn_mb\le x^{19}b$ (Lemma~\ref{pgz:deg}); $\nu_\tau$ has at most $N_p$
rows (Lemma~\ref{pgw:howe} and the proof of Lemma~\ref{pgw:one}: $W^z_{a-1}\cong A_{\nu}$ is, up to a
character, $W_\nu$ for $\auGL_{N_p}$).

\emph{The datum.} As $\vartheta=1$, $Q_\vartheta=\emptyset$ and $S(\pcb)$ consists of the places
$v=w\aubar w$ with $w\mid\pi_M$, where $a_w=1$ (proof of Lemma~\ref{pgz:std}). By the first paragraph of
the proof of Lemma~\ref{pgz:std},
\[
c_{\mathrm{vac}}^{\,N_pN_q}\,[\delta_X]_f=\auNm(\pcb)^{-N_pN_q}\sum_{x'}\zeta_{x'}\,\psrho^\circ(x')
\psvac^\circ ,
\]
a sum of $\prod_{w\mid\pcb}[M_w:\pcb M_w]=\auNm(\pcb)^{N_pN_q}$ terms with roots of unity $\zeta_{x'}$.
By Step~1 of the proof of Proposition~\ref{au:kernel}, in the form used in the proof of
Lemma~\ref{plf:KS}(KS2) for twisted Weil representations $\omega_{a,b}$ with unitary Hecke characters
of algebraic weights (which is the case of $\omega_f=\omega_{\alpha,\beta,f}$, as in the proof of
Lemma~\ref{pgz:std}), $\omega_f(x,h)\psrho^\circ(x')\psvac^\circ=\epsilon_{x'}\psrho^\circ(x'')\omega(k)
\psvac^\circ$ with $k\in\auU(W)(\auA_f)$ the image of $(x,h)$, some $x''\in W(\auA_f)$, and
$|\sigma\epsilon_{x'}|=1$. (Here $\epsilon_{x'}$ is a root of unity times the values at $\det x$, $\det h$
of the splitting characters $\alpha,\beta$ of Definition~\ref{pgv:def-edges}, which are products of the
characters $\mu^0_l$ and of powers of $\upsilon|_{\auU(1)}$, all of algebraic weight. The unramified
factors $\eta_w$ of the model at a split place, \eqref{po:eq-w5}, do not enter $\epsilon_{x'}$ with their
non-unitary part: $|\eta_w(1,h)|=|\det h|_w^{-r/2}$ is the normalisation that makes the substitution
$\Psi\mapsto\Psi(h^{-1}Tg)$ unitary, i.e.\ it belongs to the operator $\omega(k)$, and the remaining
unitary part of $\eta_w$ is a value of a power of $\upsilon$. Finally, for a character $\lambda$ of
$\auU(1)(\auA)$ trivial on $\auU(1)(E^+)$ with algebraic weight and $a\in\auU(1)(\auA_f)$: as
$\auU(1)$ is anisotropic, $a^h=\gamma\kappa$ for some $h\ge1$, $\gamma\in\auU(1)(E^+)$ and $\kappa$ in a
compact open subgroup on which $\lambda$ has finite order, so $\lambda(a)^h$ is a root of unity times
$\prod_\tau\tau(\gamma)^{-k_\tau}$ with $\gamma\aubar\gamma=1$; all conjugates of this number have modulus
$1$, $E$ being a CM field, hence so have those of $\lambda(a)$. This is the argument of Step~1 of the
proof of Proposition~\ref{au:kernel}.) Hence
\[
\pgwK^z_a(\omega_f(x,h)[\delta_X]_f)=\auNm(\pcb)^{-N_pN_q}\sum_{x'}\zeta_{x'}\epsilon_{x'}\,
S_\nu(\partial)\bigl[\Theta_{\phi_{x'}}(s_0\,\cdot)\bigr](0),\qquad \phi_{x'}=\psrho^\circ(x'')\omega(k)\psvac^\circ ,
\]
and as the matrices $S_\nu(\partial)[\cdots](0)$ are applied to $\xi$ and then conjugated by $\sigma$
(which commutes with the finite sum and the rational scalar $\auNm(\pcb)^{-N_pN_q}$), the triangle
inequality and Lemma~\ref{pgt:jet} give the bound $G_1\|\sigma\xi\|$, an average of
$\auNm(\pcb)^{N_pN_q}$ terms each at most $G_1\|\sigma\xi\|$.

\emph{Sizes.} By Lemma~\ref{an:lem-sizes}(S6) and (S4), $\Omega^{-1}=e^{-\aneu}\anrd_E/2\pi\le
x^{12}e^{|\aneu|}\le x^{13}$ ((S1): $|\aneu|\le\log x$). With $e\le x^{19}b$ and $\log(2e)\le
1+19\log x+\log b\le2\log b$ ((S1)), $\tfrac e2\log\Omega^{-1}+\tfrac e2\log(2e)\le x^{19}b(7\log x+\log
b)\le2x^{19}b\log b$. Next $D=fN_pN_q\le x^{12}(80x^6)^2\le x^{25}$, so $D\log4\le x^{26}$; and
$C_0\le\log\log b\le x$ (Lemma~\ref{pz:C0}(c), (S1)), $q_1N_pN_q+1\le x^{12}\cdot6400x^{12}+1\le x^{25}$,
so $C_0(q_1N_pN_q+1)^2\le x^{51}$. In all, $\log G_1\le2x^{19}b\log b+x^{26}+x^{51}\le x^{52}b\log b\le
x^{99}b^2$.
\end{proof}

\begin{remark}\label{pgt:rem-B1}
The bound does not depend on $\pcm$, $X$ or the points $x,h$; in particular it holds verbatim for
every twist $\vartheta$ (with $\pcm(\vartheta)$ in place of $\pcb$, using the full statement of
Lemma~\ref{pgz:std}), though only $\vartheta=1$ is used. The row lengths of $\nu$ enter only through
their sum $e$; a sharper bound, with $\sum l\log l$ in place of $e\log e$, follows from
Lemma~\ref{ps:monomial} as in Step~3 of the proof of Proposition~\ref{au:kernel}, but is not needed.
\end{remark}



%% ===== proofs/C-grid11.tex =====
% proofs/C-grid11.tex
% Inputs (C) and (Q) of proofs/C-grid9.tex (Section pgq:ss-inputs). Labels start with pgk:.
% Conventions and macros: those of proofs/C-grid9.tex.
\providecommand{\auU}{\mathrm U}
\providecommand{\auO}{\mathcal O}
\providecommand{\auA}{\mathbb A}
\providecommand{\auAz}{\mathcal A_0}
\providecommand{\auPl}{\Phi}
\providecommand{\aubar}[1]{\overline{#1}}
\providecommand{\pgrY}{\widetilde Y}
\providecommand{\pgkvol}{\operatorname{vol}}
\providecommand{\pgkP}{\mathsf P}

\section{The cusp estimate and the kernel forms}\label{pgk:sec}

This section proves the inputs (Q) and (C) of \S\ref{pgq:ss-inputs} (Theorems~\ref{pgk:Q} and
\ref{pgk:C}). Notation is that of \S\ref{pgq:ss-stat}: $V=Y_0$ of rank $n$, $Z=V\oplus E\tilde e$,
$G=\auU(V)$, $H=\auU(Z)$, $\Gamma_G=G(E^+)$, $\Gamma_H=H(E^+)$, $f=\#\auPl$, the normalised bases at
the places $\tau\in\auPl$, operator norms $\|\cdot\|$, $\operatorname{dist}(u,1)=\max_\tau\|u_\tau-1\|$,
the coefficients $c,c'$ and the types $e=d\aubar c$ of $\auAz$ and $e'=d'\aubar{c'}$ of $e_{(1,0,0)}$.
Forms of a type are meant in the sense of Definition~\ref{pz:def-type} (for $H$, with $H$ in place of
$G_j$): bounded continuous functions, left invariant under the rational points, right invariant under
an open subgroup of the finite adelic points, with $R(e)y=y$, $(R(e)y)(g)=\int e(x)y(gx)\,dx$. For
$y\in H_\infty$ (or $G_\infty$) put $\pgkP(y)=\prod_\tau\|y_\tau\|\ge1$; then $\pgkP(y^{-1})=\pgkP(y)$
and $\pgkP(yy')\le\pgkP(y)\pgkP(y')$.

\subsection{Kernel forms of a level: input (Q)}\label{pgk:ss-Q}

\begin{remark}[What is already in the paper]\label{pgk:rem-Q}
Lemma~\ref{pm:siegel}(c) bounds, uniformly in $g$, the sum over $\gamma$ of $\pgkP((g^{-1}\gamma x)_\infty)
^{-B'}\mathbf 1_{\Omega_f}((g^{-1}\gamma x)_f)$, with a constant depending on the reduction of $x$, and
Proposition~\ref{pm:bound}(a) gives the convergence of kernels $\sum_\gamma(e\otimes\varphi')(x^{-1}
\gamma y)$; both are written for a group $\auU(V)$ with the norms adapted to a Witt basis. We need the
bound only for one fixed point $z$ and uniformly in $g$, for the group $H$ with the norms of the
normalised basis; the disjointness argument of the proof of Lemma~\ref{pm:siegel}(c) gives it directly,
and we write it out in Lemma~\ref{pgk:sum}, which avoids Siegel sets and the change of norms. What
remains, that the sum is a form of type $e'$, is the reproducing property $e'*e'=e'$.
\end{remark}

Fix $B'=2(n+1)^2$. By Lemma~\ref{pgq:local}(a) with $j=B'\le h_Z$ (Lemma~\ref{pgq:sizes}(a): $h_Z\ge2A
\ge B'$), for $y\in H_\infty$
\begin{equation}\label{pgk:eq-decay}
|e'(y)|=d'\prod_\tau|c'_\tau(y_\tau)|\le d'\,2^{fB'}\,\pgkP(y)^{-B'} .
\end{equation}
By Lemma~\ref{pm:kernel}(c), applied to the unitary group $H_\tau$ of rank $n+1$ (its proof uses only
the Cartan decomposition $H_\tau=K_\tau A^+K_\tau$, and the operator norm of the normalised basis is
bi-invariant under the maximal compact subgroup $H_\tau\cap\auU(n+1)$, since the matrix of the form is
hermitian and unitary, \S\ref{pds:ss-setup}), $I_\tau=\int_{H_\tau}\|h\|^{-B'}dh<\infty$, as
$B'>(n+1)^2$.

\begin{lemma}[A uniform sum]\label{pgk:sum}
Let $K\subseteq H(\auA_f)$ be a compact open subgroup and $x\in H(\auA)$. Put $\mathsf M(y)=
\pgkP(y_\infty)^{-B'}\mathbf 1_K(y_f)$ on $H(\auA)$. There is $C_x<\infty$ with
$\sum_{\gamma\in\Gamma_H}\mathsf M(g^{-1}\gamma x)\le C_x$ for all $g\in H(\auA)$.
\end{lemma}
\begin{proof}
$\Gamma_H$ is discrete in $H(\auA)$, so there is an open neighbourhood $W$ of $1$ with $W\cap\Gamma_H=
\{1\}$. Choose $0<\rho\le\frac14$ and an open subgroup $K''\subseteq K$ such that $Z_{4\rho}\times K''
\subseteq x^{-1}Wx$, where $Z_r=\{u\in H_\infty:\operatorname{dist}(u,1)\le r\}$; this is possible since
these sets form a basis of neighbourhoods of $1$. Let $U=Z_\rho\times K''$. For $u,u'\in Z_\rho$,
$\|u_\tau^{-1}\|=\|u_\tau\|\le2$ and $u'_\tau u_\tau^{-1}-1=(u'_\tau-u_\tau)u_\tau^{-1}$, so
$u'u^{-1}\in Z_{4\rho}$; hence $UU^{-1}\subseteq x^{-1}Wx$.

\emph{Disjointness.} If $g^{-1}\gamma xu=g^{-1}\gamma'xu'$ with $u,u'\in U$, then $\gamma'^{-1}\gamma=
xu'u^{-1}x^{-1}\in W\cap\Gamma_H$, so $\gamma=\gamma'$. Thus the sets $g^{-1}\gamma xU$ are disjoint.

\emph{Comparison.} For $y\in H(\auA)$, $u\in U$: $\pgkP((yu)_\infty)\le2^f\pgkP(y_\infty)$ and
$\mathbf 1_K(y_fu_f)=\mathbf 1_K(y_f)$, so $\mathsf M(y)\le2^{fB'}\mathsf M(yu)$. Averaging over $U$ and
using left invariance of Haar measure,
\[
\pgkvol(U)\sum_\gamma\mathsf M(g^{-1}\gamma x)\le2^{fB'}\sum_\gamma\int_{g^{-1}\gamma xU}\mathsf M\le
2^{fB'}\int_{H(\auA)}\mathsf M=2^{fB'}\pgkvol(K)\prod_\tau I_\tau .
\]
So $C_x=2^{fB'}\pgkvol(U)^{-1}\pgkvol(K)\prod_\tau I_\tau$ works ($\pgkvol(U)>0$, $U$ having non-empty
interior).
\end{proof}

\begin{theorem}[Input (Q)]\label{pgk:Q}
Let $K\subseteq H(\auA_f)$ be an open subgroup of the stabiliser of $\widehat L_Z$ and $z\in H(\auA)$.
For $g\in H(\auA)$ the series
\[
\mathsf F(g)=\sum_{\gamma\in\Gamma_H}e'\bigl((z\gamma g)_\infty\bigr)\,\mathbf 1_K\bigl((z\gamma g)_f\bigr)
\]
converges absolutely, $\sup_g\sum_\gamma|\cdots|<\infty$, and $Q=\aubar{\mathsf F}$ is a form of type
$e'$ on $H$ fixed by $K$. This is \textup{(Q)} of \S\ref{pgq:ss-inputs}.
\end{theorem}
\begin{proof}
$K$ is compact (closed in a compact group) and open. Write $\Phi(y)=e'(y_\infty)\mathbf 1_K(y_f)$.
By \eqref{pgk:eq-decay}, $\pgkP(y^{-1})=\pgkP(y)$ and $K=K^{-1}$,
$|\Phi(y)|\le d'2^{fB'}\mathsf M(y^{-1})$. With $y=z\gamma g$, $y^{-1}=g^{-1}\gamma^{-1}z^{-1}$, and
Lemma~\ref{pgk:sum} with $x=z^{-1}$ (sum over $\gamma^{-1}$) gives
\begin{equation}\label{pgk:eq-unif}
\sum_\gamma|\Phi(z\gamma g)|\le C:=d'2^{fB'}C_{z^{-1}}\qquad\text{for all }g\in H(\auA).
\end{equation}
So $\mathsf F$ is defined, $|\mathsf F|\le C$, and $\mathsf F$ is measurable (a pointwise limit of
finite sums of measurable functions). \emph{Invariance.} For $\delta\in\Gamma_H$, $\mathsf F(\delta g)=
\mathsf F(g)$ (substitute $\gamma\mapsto\gamma\delta^{-1}$); for $k\in K$, $(z\gamma gk)_f=(z\gamma g)_fk$
and $(z\gamma gk)_\infty=(z\gamma g)_\infty$, so $\mathsf F(gk)=\mathsf F(g)$.

\emph{Reproducing property.} Since $e'(h^{-1})=\aubar{e'(h)}$ ($c'_\tau(h^{-1})=\aubar{c'_\tau(h)}$ as in
Lemma~\ref{pgq:local}(d)), $R(e')Q=Q$ is equivalent to
\[
\mathsf F(g)=\int_{H_\infty}e'(h^{-1})\,\mathsf F(gh)\,dh\qquad(g\in H(\auA)).
\]
For each $\gamma$, with $y=(z\gamma g)_\infty$, $\int e'(h^{-1})\Phi(z\gamma gh)\,dh=\mathbf 1_K((z\gamma
g)_f)\int e'(yh)e'(h^{-1})\,dh=\Phi(z\gamma g)$, because $\int e'(yh)e'(h^{-1})dh=(e'*e')(y)=e'(y)$. The
last identity is Schur's orthogonality relation for the discrete series $\pi'=\bigotimes_\tau\pi'_\tau$
of formal degree $d'$ and unit vector $v'$ (with the convention recalled in the proof of
Lemma~\ref{pgq:schur}): $e'(y)=d'\langle\pi'(y)v',v'\rangle$, and
$\int\langle v',\pi'(h)v'\rangle\aubar{\langle\pi'(y^{-1})v',\pi'(h)v'\rangle}dh=d'^{-1}\langle v',\pi'(y^{-1})
v'\rangle$, which is the integral $d'^{-2}\int e'(yh)e'(h^{-1})dh$ written out. The exchange of sum and
integral is justified by Tonelli: $\int|e'(h^{-1})|\sum_\gamma|\Phi(z\gamma gh)|\,dh\le C\|e'\|_{L^1}
<\infty$ by \eqref{pgk:eq-unif} and the integrability of $e'$ \eqref{pgk:eq-decay}.

\emph{Continuity.} By the reproducing property, $\mathsf F(g)=\int e'(h^{-1}g_\infty)\mathsf F(h,g_f)\,dh$
after the substitution $h\mapsto g_\infty^{-1}h$. As $g_\infty\mapsto e'(\,\cdot\,g_\infty)$ is
continuous into $L^1(H_\infty)$ (continuity of translation in $L^1$) and $\mathsf F$ is bounded,
$\mathsf F$ is continuous in $g_\infty$, uniformly in $g_f$; and it is locally constant in $g_f$ by
right $K$-invariance. Hence $\mathsf F$ and $Q$ are continuous, and $Q$ is a form of type $e'$ fixed
by $K$ with $\aubar{Q(g)}=\mathsf F(g)$.
\end{proof}

\subsection{Coordinates, lattices and two arithmetic facts}\label{pgk:ss-coord}

\emph{Coordinates.} $V$ is a space in the sense of \S\ref{pc:ss-data}: it has a frame, an $E$-basis
whose $\widehat\auO_E$-span $\widehat L_V$ is unimodular. We identify $V(E)=E^n$ by the frame; the
form has a hermitian Gram matrix $\mathsf H\in M_n(E)$, integral with unit determinant at every finite
place. For $\tau\in\auPl$ let $\mathsf P_\tau\in\mathrm{GL}_n(\C)$ be the matrix whose columns are the
normalised basis in frame coordinates, so $\mathsf P_\tau^*\tau(\mathsf H)\mathsf P_\tau=\mathsf J_\tau$.
For $x\in G(\auA)$ and $v\in V(\auA)$ the \emph{coordinates at $x$} are $\xi_x(v)=x^{-1}v$, read at
$\tau$ in the normalised basis (that is, $\mathsf P_\tau^{-1}x_\tau^{-1}\tau(v)$) and at the finite
places in the frame; this is the reading of \S\ref{pgq:ss-stat}, so that $L_V(x)=\{v\in E^n:(x^{-1}v)_f
\text{ integral}\}$, $|v|_x^2=\sum_\tau|\xi_x(v)_\tau|^2$, and $\lambda(x)$ is the first minimum defined
there. In these coordinates the form is $\mathsf S$, with $\mathsf S_\tau=\mathsf J_\tau$ (hermitian and
unitary, so $\|\mathsf S_\tau\|=1$) and $\mathsf S_f=x_f^*\mathsf Hx_f=\mathsf H$; for rational vectors
$v,w$, $\langle v,w\rangle=\xi_x(v)^*\mathsf S\,\xi_x(w)$ at every place. We extend $\xi_x$ to
matrices whose columns are vectors.

\begin{lemma}[Lattice vectors]\label{pgk:lat}
Let $x\in G(\auA)$ and $v,w\in L_V(x)$.
\begin{enumerate}
\item[(a)] $\langle v,w\rangle\in\auO_E$ and $|\tau\langle v,w\rangle|\le|v|_x|w|_x$ for every $\tau$.
If $|v|_x|w|_x<1$ then $\langle v,w\rangle=0$.
\item[(b)] If $v\ne0$ and $y_1,\dots,y_n\in L_V(x)$ span $V$ with $|y_j|_x\le R$, then $|v|_xR\ge1$.
\item[(c)] $L_V(x)$ is a lattice of rank $2fn$ in $\prod_\tau\C^n\cong\R^{2fn}$ \textup{(}embedded by
the archimedean coordinates\textup{)} of covolume $(2^{-f}|D_E|^{1/2})^n$, $D_E$ the discriminant of
$E$.
\end{enumerate}
\end{lemma}
\begin{proof}
(a) At the finite places $\langle v,w\rangle=\xi_x(v)^*\mathsf H\xi_x(w)$ with integral entries, so
$\langle v,w\rangle\in E\cap\widehat\auO_E=\auO_E$. At $\tau$, $|\tau\langle v,w\rangle|=|\xi_x(v)_\tau^*
\mathsf J_\tau\xi_x(w)_\tau|\le|\xi_x(v)_\tau||\xi_x(w)_\tau|\le|v|_x|w|_x$. A non-zero $a\in\auO_E$
has $\prod_{\tau\in\auPl}|\tau(a)|^2=|\mathrm N_{E/\Q}(a)|\ge1$, so $\max_\tau|\tau(a)|\ge1$.
(b) Otherwise $|v|_x|y_j|_x<1$ for all $j$, so $v$ is orthogonal to all $y_j$ by (a), hence to $V$, and
$v=0$ as the form is non-degenerate.
(c) $L_V(x)=E^n\cap x_f\widehat\auO_E^{\,n}$. Since $x_f$ preserves the unimodular form,
$\det x_f$ is a unit at every finite place, so $x_f\widehat\auO_E^{\,n}$ and $\widehat\auO_E^{\,n}$ have the same
volume in $\auA_{E,f}^n$; hence the covolumes of $L_V(x)$ and of $\auO_E^n$ in $E_\infty^n$ (embedded by
$v\mapsto(\tau(v))_\tau$) agree, both being the covolume of $E^n$ in $\auA_E^n$ divided by that volume. The latter is $(2^{-f}
|D_E|^{1/2})^n$ (Lemma~\ref{pbc:tb-covol}; see also \cite[Chapter~I, (5.1) and Proposition~(5.2)]{Neukirch}). The archimedean coordinate map is $v\mapsto(\mathsf
P_\tau^{-1}x_\tau^{-1}\tau(v))_\tau$, of real Jacobian $\prod_\tau|\det\mathsf P_\tau|^{-2}|\det x_\tau|
^{-2}$. Here $|\det x_\tau|=1$ ($x_\tau$ is unitary for a non-degenerate form), and $|\det\mathsf P_\tau|
^{-2}=|\tau(\det\mathsf H)|$ because $\det\mathsf J_\tau=\pm1$; as $\det\mathsf H\in E^+$ is a unit at
every finite place, $\prod_\tau|\tau(\det\mathsf H)|=|\mathrm N_{E^+/\Q}\det\mathsf H|=1$.
\end{proof}

\begin{lemma}[Rounding]\label{pgk:round}
Let $N_d=q_1d$, a natural number in the different $\mathfrak d_E$ of $E$ (\S\ref{psp:ss-sparse}). For
every $c\in\auA_E$ there is $\varrho\in E$ with $c_f-\varrho\in\widehat\auO_E$ and $|c_\tau-\tau(\varrho)|
\le2f^2N_d$ for every $\tau\in\auPl$.
\end{lemma}
\begin{proof}
As $E+\widehat\auO_E=\auA_{E,f}$ (approximation theorem), there is $\varrho_0\in E$ with $c_f-\varrho_0
\in\widehat\auO_E$. Embed $E_\infty=\prod_\tau\C\cong\R^{2f}$ with the inner product $(a,b)=\Rel\sum_\tau
\tau(a)\aubar{\tau(b)}=\frac12\mathrm{Tr}_{E/\Q}(a\bar b)$ on $E$. The dual of $\auO_E$ is
$\{b:\frac12\bar b\in\mathfrak d_E^{-1}\}=2\mathfrak d_E^{-1}$. For $0\ne b\in2\mathfrak d_E^{-1}$,
$N_db/2\in N_d\mathfrak d_E^{-1}\subseteq\auO_E$ is non-zero, so $|N_db/2|^2=\sum_\tau|\tau(N_db/2)|^2
\ge f\ge1$ (arithmetic and geometric means); hence the dual has no non-zero vector shorter than $2/N_d$.
By Lemma~\ref{ps:transfer} with $m=2f$, every point of $E_\infty$ is within $(2f)^2N_d/2$ of $\auO_E$:
there is $o\in\auO_E$ with $|(c_\tau-\tau(\varrho_0)-\tau(o))_\tau|\le2f^2N_d$. Take $\varrho=\varrho_0+o$.
\end{proof}

\subsection{A short vector gives a cusp}\label{pgk:ss-cusp}

\begin{lemma}[Minkowski basis]\label{pgk:mink}
Let $x\in G(\auA)$ and $\kappa\ge(fn\,\mathrm{rd}_E)^{1/2}$, where $\mathrm{rd}_E=|D_E|^{1/2f}$. There
is an $E$-basis $y_1,\dots,y_n$ of $V$ in $L_V(x)$ with $\alpha_j=|y_j|_x$ non-decreasing and
$\prod_j\alpha_j\le\kappa^n$.
\end{lemma}
\begin{proof}
Let $D=2fn$ and $\lambda_1\le\dots\le\lambda_D$ the successive minima of $L_V(x)\subset\R^D$ for the
Euclidean unit ball, of volume $\pi^{fn}/(fn)!$. By Minkowski's second theorem
(\cite[Chapter~VIII, \S4.3, Theorem~V]{Cassels}; see also \cite[Theorem~1.3]{HenkMinima}) and Lemma~\ref{pgk:lat}(c),
\[
\lambda_1\cdots\lambda_D\le\frac{2^D\,2^{-fn}|D_E|^{n/2}}{\pi^{fn}/(fn)!}\le\Bigl(\frac2\pi\Bigr)^{fn}(fn)^{fn}
\mathrm{rd}_E^{\,fn}\le(fn\,\mathrm{rd}_E)^{fn},
\]
as $|D_E|^{1/2}=\mathrm{rd}_E^{\,f}$. Exactly as in the proof of Lemma~\ref{pbc:siegel}, choose
$\Z$-independent $v_1,\dots,v_D$ with $|v_j|_x=\lambda_j$; an $E$-subspace of dimension $e$ contains at
most $2fe$ $\Z$-independent vectors, so we may pick $E$-independent $y_i$ among $v_1,\dots,
v_{2f(i-1)+1}$ with $|y_i|_x\le\lambda_{2f(i-1)+1}$, $1\le i\le n$. Then $|y_i|_x^{2f}\le\prod_{t=1}^{2f}
\lambda_{2f(i-1)+t}$, so $\prod_i|y_i|_x^{2f}\le(fn\,\mathrm{rd}_E)^{fn}$ and $\prod_i|y_i|_x\le(fn\,
\mathrm{rd}_E)^{n/2}\le\kappa^n$. Reorder by length.
\end{proof}

\begin{lemma}[Cusp frame]\label{pgk:frame}
Let $\kappa\ge\max(1,(fn\,\mathrm{rd}_E)^{1/2})$, $\theta>0$ with $\theta\kappa\le1$, and $x\in G(\auA)$
with $\lambda(x)<\theta^n$. There are $1\le i\le n/2$, a non-zero isotropic subspace $A\subset V$ of
dimension $i$, and lattice vectors in $L_V(x)$: $y_1,\dots,y_i$ spanning $A$, of length at most
$\ell_X$, and $y_1,\dots,y_{n-i}$ spanning $A^\perp$, of length at most $\ell_Y$, with $\ell_X\le\ell_Y$
and $\ell_X\ell_Y<\theta\kappa$.
\end{lemma}
\begin{proof}
As $\lambda(x)<\theta^n\le1$, there is $0\ne\xi\in L_V(x)$ with $|\xi|_x<\theta^n$. Take $y_j$,
$\alpha_j$ from Lemma~\ref{pgk:mink}. By Lemma~\ref{pgk:lat}(b) with $R=\alpha_n$, $|\xi|_x\alpha_n\ge1$, so
\[
\prod_{j<n}\alpha_j\le|\xi|_x\prod_{j\le n}\alpha_j<\theta^n\kappa^n .
\]
If $n=1$ this reads $1<\theta\kappa$, which is false; so $n\ge2$. Put $k=n-1$ and $c=\theta\kappa\le1$.
Then $\prod_{j=1}^k\alpha_j\alpha_{k+1-j}=(\prod_{j\le k}\alpha_j)^2<c^{2n}\le c^k$, so some factor has
$\alpha_j\alpha_{k+1-j}<c$; replacing $j$ by $k+1-j$ if needed we get $i$ with $i\le k+1-i=n-i$ (so
$\alpha_i\le\alpha_{n-i}$ by monotonicity) and $\alpha_i\alpha_{n-i}<c$. Put $\ell_X=\alpha_i$,
$\ell_Y=\alpha_{n-i}$, $A=\langle y_1,\dots,y_i\rangle$, $B=\langle y_1,\dots,y_{n-i}\rangle$. For
$a\le i$ and $b\le n-i$, $|y_a|_x|y_b|_x\le\ell_X\ell_Y<1$, so $\langle y_a,y_b\rangle=0$ by
Lemma~\ref{pgk:lat}(a). Thus $A\subseteq B\subseteq A^\perp$; $A$ is isotropic; and $\dim B=n-i=
\dim A^\perp$, so $B=A^\perp$.
\end{proof}

\subsection{The unipotent radical and its reduction}\label{pgk:ss-rad}

Let $A\subset V$ be a non-zero isotropic subspace. Its \emph{unipotent radical} $N_A$ is the algebraic
subgroup of $G$ of the $u$ that fix $A$ pointwise and induce the identity on $A^\perp/A$. As for $G$, we write $N_A(R)$ for the elements of $N_A$ read as matrices over $R=E,\auA_E,\dots$; thus $N_A(E)=N_A(E^+)\subseteq\Gamma_G$ and $N_A(\auA)=N_A(\auA_{E^+})$. In the
coordinates at a point $x$, let $\mathsf a$ (an $n\times i$ matrix) and $\mathsf y$ (an $n\times(n-i)$
matrix) be the coordinates of the columns $X=(y_1,\dots,y_i)$ and $Y=(y_1,\dots,y_{n-i})$ of
Lemma~\ref{pgk:frame}; so $\mathsf a^*\mathsf S\mathsf a=0$, $\mathsf a^*\mathsf S\mathsf y=0$. Let $R$ be
$E$, $\auA_E$ or one of its completions, $\mathsf S$ the form (with $\mathsf a$, $\mathsf y$) read
over $R$. For an $n\times i$ matrix $w$ over $R$ with $\mathsf a^*\mathsf Sw=0$ put
\[
P_w=(w\mathsf a^*-\mathsf aw^*)\mathsf S,\qquad\mathsf E(w)=1+2P_w+2P_w^2=1+2\bigl(w\mathsf a^*-\mathsf a
w^*-\mathsf a\,w^*\mathsf Sw\,\mathsf a^*\bigr)\mathsf S .
\]

\begin{lemma}\label{pgk:param}
\begin{enumerate}
\item[(a)] $P_w$ is skew-adjoint for $\mathsf S$, $P_w^2=-\mathsf a(w^*\mathsf Sw)\mathsf a^*\mathsf S$,
$P_w^3=0$, and $\mathsf E(w)=\exp(2P_w)$ lies in $N_A(R)$ \textup{(}read in coordinates\textup{)}.
\item[(b)] $\mathsf E(w)\mathsf E(w')=\mathsf E\bigl(w+w'-\mathsf a(w^*\mathsf Sw')\bigr)$.
\item[(c)] Every $u\in N_A(R)$ is $\mathsf E(w)$ for some $w=\mathsf yC$, $C$ an $(n-i)\times i$ matrix
over $R$.
\end{enumerate}
\end{lemma}
\begin{proof}
Note $w^*\mathsf S\mathsf a=(\mathsf a^*\mathsf Sw)^*=0$. (a) $\mathsf SP_w=\mathsf S(w\mathsf a^*-\mathsf a
w^*)\mathsf S$ is skew-hermitian. Expanding $P_w^2$, the terms containing $\mathsf a^*\mathsf S\mathsf a$,
$\mathsf a^*\mathsf Sw$ or $w^*\mathsf S\mathsf a$ vanish, leaving $-\mathsf aw^*\mathsf Sw\mathsf a^*\mathsf
S$; and $P_wP_w^2=0$ for the same reason. So $\exp(2P_w)=1+2P_w+2P_w^2$ is unitary for $\mathsf S$ (the
exponential of an $\mathsf S$-skew-adjoint nilpotent matrix). It fixes the columns of $\mathsf a$
($P_w\mathsf a=0$), and for $v$ with $\mathsf a^*\mathsf Sv=0$, $P_wv=-\mathsf a(w^*\mathsf Sv)\in A$ and
$P_w^2v\in\mathsf a\cdot R^i$; so it is the identity on $A^\perp/A$.
(b) For two-step nilpotent matrices, $\exp X\exp Y=\exp(X+Y+\frac12[X,Y])$. Here, expanding as in (a),
$P_wP_{w'}=-\mathsf a(w^*\mathsf Sw')\mathsf a^*\mathsf S$, so $[P_w,P_{w'}]=-\mathsf a(k-k^*)\mathsf a^*
\mathsf S$ with $k=w^*\mathsf Sw'$, and this commutes with $P_w,P_{w'}$. Also $P_{-\mathsf ak}=-\mathsf a
(k-k^*)\mathsf a^*\mathsf S$ (and $\mathsf a^*\mathsf S\mathsf ak=0$). Hence $2P_w+2P_{w'}+2[P_w,P_{w'}]=
2P_{w+w'-\mathsf ak}$.
(c) Let $u\in N_A(R)$ and $X=\log u=(u-1)-\frac12(u-1)^2$; since $(u-1)V\subseteq A^\perp$,
$(u-1)A^\perp\subseteq A$, $(u-1)A=0$, we have $(u-1)^3=0$, $u=\exp X$, and $X$ is $\mathsf S$-skew-adjoint
with the same three properties. Choose $\mathsf b$ ($n\times i$ over $R$) with $\mathsf a^*\mathsf S
\mathsf b=1_i$ (possible: $\mathsf a$ has independent columns and $\mathsf S$ is invertible); then
$R^n=A^\perp\oplus\mathsf bR^i$. Put $w_0=X\mathsf b$, so $\mathsf a^*\mathsf Sw_0=-(X\mathsf a)^*\mathsf
S\mathsf b=0$, and $c=w_0^*\mathsf S\mathsf b=-\mathsf b^*\mathsf SX\mathsf b$, which is skew-hermitian.
We claim $X=P_{w_0}+\mathsf ac\mathsf a^*\mathsf S$. On $\mathsf b$: the right side gives $w_0-\mathsf a
(w_0^*\mathsf S\mathsf b)+\mathsf ac=w_0=X\mathsf b$. On $v\in A^\perp$: write $Xv=\mathsf at$; then
$t=\mathsf b^*\mathsf S\mathsf at=\mathsf b^*\mathsf SXv$, and the right side gives $-\mathsf aw_0^*\mathsf
Sv=-\mathsf a\mathsf b^*X^*\mathsf Sv=\mathsf a\mathsf b^*\mathsf SXv=Xv$. Now $\mathsf ac\mathsf a^*
\mathsf S=P_{\mathsf ac/2}$ (as $c^*=-c$), so $X=P_w$ with $w=w_0+\mathsf ac/2$, and $\mathsf a^*\mathsf S
w=0$. Hence $u=\exp(2P_{w/2})=\mathsf E(w/2)$. The columns of $w/2$ lie in $A^\perp$, which is spanned
by the columns of $\mathsf y$, so $w/2=\mathsf yC$.
\end{proof}

\begin{proposition}[Reduction in the radical]\label{pgk:cover}
Let $M\ge1$ be an integer, $0<\delta\le1$, $\kappa\ge48n^6f^2N_d$, $x\in G(\auA)$, and let $A$,
$\ell_X$, $\ell_Y$ be as in Lemma~\ref{pgk:frame} with $\ell_X\ell_Y\le\delta/(M\kappa)$. For every
$u\in N_A(\auA)$ there is $\gamma\in N_A(E)$ such that $k=x^{-1}\gamma ux$ satisfies
$\operatorname{dist}(k_\infty,1)\le\delta$ and $k_f\in J_G(M)$.
\end{proposition}
\begin{proof}
Work in the coordinates at $x$: $\tilde u=x^{-1}ux\in N_A$ read in coordinates. By
Lemma~\ref{pgk:param}(c), $\tilde u=\mathsf E(M\mathsf yc)$ for an adelic $(n-i)\times i$ matrix $c$.
Let $\varrho(\cdot)$ be the entrywise rounding of Lemma~\ref{pgk:round}, put $d=-M\varrho(c)^*(\mathsf y^*
\mathsf S\mathsf y)c$ (an adelic $i\times i$ matrix) and $Z=-M(Y\varrho(c)-X\varrho(d))$, a matrix over
$E$ whose columns lie in $A^\perp$. Let $\gamma=\mathsf E_{\mathsf H}(Z)$ be the element defined by the
same formula in frame coordinates (with $X$, $\mathsf H$ in place of $\mathsf a$, $\mathsf S$); by
Lemma~\ref{pgk:param}(a) over $R=E$, $\gamma\in N_A(E)$. Since the coordinates at $x$ are obtained by
applying $x^{-1}$ (with $\mathsf S=x^*\mathsf Hx$, $\mathsf a=x^{-1}X$, $\mathsf y=x^{-1}Y$), $x^{-1}\gamma
x=\mathsf E(z)$ with $z=x^{-1}Z=-M(\mathsf y\varrho(c)-\mathsf a\varrho(d))$. By Lemma~\ref{pgk:param}(b),
$k=\mathsf E(z)\mathsf E(M\mathsf yc)=\mathsf E(z+M\mathsf yc-\mathsf az^*\mathsf SM\mathsf yc)$. Using
$\mathsf a^*\mathsf S\mathsf y=0$, $z^*\mathsf SM\mathsf yc=-M^2\varrho(c)^*\mathsf y^*\mathsf S\mathsf yc=
Md$, so
\[
k=\mathsf E(MT),\qquad T=\mathsf y\,(c-\varrho(c))-\mathsf a\,(d-\varrho(d)).
\]
\emph{Finite places.} The entries of $c-\varrho(c)$, $d-\varrho(d)$, $\mathsf a_f$, $\mathsf y_f$ (lattice
vectors) and $\mathsf S_f=\mathsf H$ are integral, so $k_f=1+M\mathsf Z$ with $\mathsf Z$ integral. As
$k_f$ preserves the unimodular form, $\det k_f$ is a unit, so $k_f$ lies in the stabiliser of
$\widehat L_V$ and is $\equiv1$ modulo $M$: $k_f\in J_G(M)$.
\emph{Real places.} A matrix with $p$ columns each of length at most $\ell$ has operator norm at most
$\sqrt p\,\ell$, and a $p\times q$ matrix with entries at most $\beta$ has operator norm at most
$\sqrt{pq}\,\beta$. So $\|\mathsf a_\tau\|\le\sqrt n\ell_X$, $\|\mathsf y_\tau\|\le\sqrt n\ell_Y$,
$\|(c-\varrho(c))_\tau\|,\|(d-\varrho(d))_\tau\|\le2nf^2N_d$, and $\|T_\tau\|\le4n^{3/2}f^2N_d\ell_Y$ (as
$\ell_X\le\ell_Y$). With $q=\|T_\tau\|\|\mathsf a_\tau\|\le4n^2f^2N_d\ell_X\ell_Y$ and $\|\mathsf S_\tau
\|=1$,
\[
\|k_\tau-1\|\le4Mq+2M^2q^2,\qquad Mq\le\frac{4n^2f^2N_d\,\delta}{\kappa}\le\frac\delta{12},
\]
so $\|k_\tau-1\|\le\frac\delta3+\frac{\delta^2}{72}\le\delta$.
\end{proof}

\subsection{The cusp estimate: input (C)}\label{pgk:ss-C}

\begin{theorem}[Input (C)]\label{pgk:C}
Let $\kappa_{\mathrm c}=8fn\,\mathrm{rd}_E+48n^6f^2N_d+1$. Let $F$ be a form of type $e$ on $G$ fixed by
$J_G(M)$, let $0<\delta\le1$, and let $x\in G(\auA)$ with $\lambda(x)<(\delta/(M\kappa_{\mathrm c}^2))^n$.
If $|F(xu)-F(x)|\le C$ for all $u\in G_\infty$ with $\operatorname{dist}(u,1)\le\delta$, then $|F(x)|\le
C$. This is \textup{(C)} of \S\ref{pgq:ss-inputs}.
\end{theorem}
\begin{proof}
\emph{The cusp.} Let $\kappa=\kappa_{\mathrm c}$ and $\theta=\delta/(M\kappa^2)$. Then $\kappa\ge1$,
$\kappa\ge8fn\,\mathrm{rd}_E\ge(fn\,\mathrm{rd}_E)^{1/2}$ ($\mathrm{rd}_E\ge1$), $\kappa\ge48n^6f^2N_d$,
and $\theta\kappa=\delta/(M\kappa)\le1$. By Lemma~\ref{pgk:frame} there is a non-zero isotropic $A$ with
$\ell_X\ell_Y<\theta\kappa=\delta/(M\kappa)$, and Proposition~\ref{pgk:cover} applies: every $u\in
N_A(\auA)$ has $\gamma_u\in N_A(E)$ with $x^{-1}\gamma_uux\in\mathcal B:=Z_\delta\times J_G(M)$, where
$Z_\delta=\{z\in G_\infty:\operatorname{dist}(z,1)\le\delta\}$.

\emph{Compactness.} Put $\mathcal K=N_A(\auA)\cap x\mathcal Bx^{-1}$, compact ($N_A(\auA)$ is closed and
$\mathcal B$ compact). Every $u\in N_A(\auA)$ lies in $N_A(E)\mathcal K$. So $N_A(E)\backslash N_A(\auA)$
is compact; give it the quotient of a Haar measure of $N_A(\auA)$, of finite mass $\mu>0$.

\emph{Oscillation.} For $u\in N_A(\auA)$ write $x^{-1}\gamma_uux=k=(k_\infty,k_f)$. As $F$ is left
invariant under $\gamma_u\in\Gamma_G$, right invariant under $J_G(M)$, and real and finite points commute,
$F(ux)=F(\gamma_uux)=F(xk_\infty k_f)=F(xk_\infty)$, so $|F(ux)-F(x)|\le C$.

\emph{The constant term vanishes.} The function $\mathsf F=F$ is bounded, continuous, left invariant
under $N_A(E)\subseteq\Gamma_G$ and satisfies $\mathsf F(g)=\int_{G_\infty}\mathsf F(gh)e(h)\,dh$; and the
real points of $N_A$ at $\tau$ form the unipotent radical of the stabiliser in $G_\tau$ of the non-zero
isotropic subspace $\tau(A)$, a proper parabolic subgroup. The argument of the first paragraph of the
proof of Lemma~\ref{pm:decay} (``the constant term vanishes'') therefore gives
$\int_{N_A(E)\backslash N_A(\auA)}F(ug)\,du=0$ for all $g$. That argument uses Lemma~\ref{pm:cusp} for
$e_\tau=d_\tau\aubar{c_\tau}$; its hypotheses (the bound and smoothness of Lemma~\ref{pm:kernel}(a),(b)
with $B=5n^3$) hold here by Lemma~\ref{pgq:local}(a) with $j=5n^3\le h_V$ (Lemma~\ref{pgq:sizes}(c)
gives $h_V\ge500(n+1)^3\ge5n^3$) and the proof of Lemma~\ref{pm:kernel}(b), which uses only that $\pi_\tau$ is
irreducible unitary and admissible (Remark~\ref{pm:rem-other}(1),(3)).

\emph{Conclusion.} The function $u\mapsto F(ux)$ is continuous and left $N_A(E)$-invariant, so
\[
\mu\,|F(x)|=\Bigl|\int_{N_A(E)\backslash N_A(\auA)}\bigl(F(x)-F(ux)\bigr)du\Bigr|\le\mu\,C .
\qedhere
\]
\end{proof}



%% ===== proofs/C-grid12.tex =====
% proofs/C-grid12.tex
% Input (S) of proofs/C-grid9.tex (the restriction pairing), proved by Laplace's method.
% Labels start with pgl:. Conventions and macros: those of proofs/C-grid9.tex and B-period.tex.
\providecommand{\auU}{\mathrm U}
\providecommand{\aubar}[1]{\overline{#1}}
\providecommand{\pgqvol}{\operatorname{vol}}
\providecommand{\pdsK}{\mathbf K}
\providecommand{\pdsg}{\mathfrak g}
\providecommand{\pdsm}{m}
\providecommand{\pbptr}{\operatorname{tr}}
\providecommand{\pglF}{\mathsf F}
\providecommand{\pglLarge}{\mathrm{Large}}
\providecommand{\pglSmall}{\mathrm{Small}}

\section{The restriction pairing}\label{pgl:sec}

This section proves input (S) of \S\ref{pgq:ss-inputs} (Theorem~\ref{pgl:S}). The argument is the
Laplace method of \S\ref{pbp:sec}, with two changes: the compact factor of \S\ref{pbp:sec} is replaced
by a second noncompact group, and the stationary position has a single vector. Lemmas of
\S\S\ref{pbp:sec}, \ref{pds:sec}, \ref{pco:sec} are reused in the general form (arbitrary rank and
parameter) stated and proved in \S\ref{pgm:sec}; what differs is written in full. \emph{Dependencies:} Theorem~\ref{pgl:S} is proved in full, modulo the
published results already used in \S\S\ref{pbp:sec}, \ref{pds:sec}, \ref{pco:sec} (Schwarz's lemma,
the Casimir value \eqref{pds:eq-casimir} through Lemma~\ref{pds:lip}).

\subsection{Statement and reductions}\label{pgl:ss-stat}

We use the notation of \S\ref{pgq:ss-stat} and \S\ref{pgq:ss-local}: $n=r$, $m=2\varrho+2$, $k$,
$B=\kappa(\tau)+\Delta/2$, $\beta_{\mathsf S}=b^{3m}$, $N_{\mathsf S}=20(2n+1)^2$, $V_\tau=\pgqvol\{r\in
G_\tau:\|r-1\|\le1\}$, and the coefficients $c_\tau$ (type $e_0$, Harish-Chandra parameter
$\varsigma_\tau$, \S\ref{pds:ss-setup}) and $c'_\tau$ (type $e_{(1,0,0)}$, Harish-Chandra parameter
$-\lambda'$ with $\lambda'=\lambda_{(1,0,0),\tau}$).

\begin{theorem}[Input (S)]\label{pgl:S}
For every $\tau\in\auPl$ there is $\omega_\tau\in H_\tau$ with $\|\omega_\tau\|^2\le\beta_{\mathsf S}$ and
\[
V_\tau\le\beta_{\mathsf S}^{N_{\mathsf S}}\Bigl|\int_{G_\tau}c_\tau(r)\,\aubar{c'_\tau\bigl(\omega_\tau^{-1}
(r\oplus1)\omega_\tau\bigr)}\,dr\Bigr| .
\]
\end{theorem}

Fix $\tau$ and drop it. \emph{Matrices.} Let $p$ be the unitary matrix of \S\ref{pds:ss-setup} with
$p^{-1}\mathsf J_\tau p=J:=\operatorname{diag}(\varepsilon_i)$ (signs on $V$), and put $p_Z=p\oplus1$,
$J_Z=J\oplus1$. Conjugation by $p$, $p_Z$ identifies $G_\tau$, $H_\tau$ with $G=\auU(J)$, $H=\auU(J_Z)$,
preserves operator norms, carries $r\oplus1$ to $r\oplus1$, and carries the Haar measure to a Haar
measure. We work in $G$, $H$ from now on; the ball $\{\|r-1\|\le1\}$ is preserved. By
Lemma~\ref{pgm:coords} (with Lemmas~\ref{pds:chamber} and \ref{pgq:steep}), the coordinates of $\C^n$ split as $X\sqcup\{1,\dots,m\}$
and those of $\C^{n+1}$ as $X\sqcup\{0,\dots,m\}$ (the same set $X$ of $k-1$ coordinates, of sign $+$,
the extra coordinate of $\C^{n+1}$ lying among the far ones), numbered so that: on $X$ both weights
vanish; for $V$, $\Lambda_1>\dots>\Lambda_m$ lie in $[-3B,-B]$ and the coordinate $j$ has sign
$(-1)^{j-1}$; for $Z$, $\Lambda'_0>\dots>\Lambda'_m$ lie in $[-3B,-B]$ and the coordinate $i$ has sign
$(-1)^i$. Here $\Lambda=\Lambda(\varsigma)$ and $\Lambda'=\Lambda(\lambda')$ are the highest weights of the
lowest $K$-types (\S\ref{pds:ss-setup}). All weights are integers, and $m$ is even.

\emph{The functions $\pglF_\lambda$.} For a unitary group $\auU(J_0)$ of rank $N$ and a regular
$\lambda\in(\frac{N-1}2+\Z)^N$ let $\pglF_\lambda$ be the function $F$ of \eqref{pco:eq-F}, with its
$\Lambda=\Lambda(\lambda)$, $\mu=\lambda+\rho$. By Theorem~\ref{pco:FJ} (gap $\ge N^2+1$, which holds here
by Lemma~\ref{pgq:sizes}(c)), $c=\pglF_\varsigma$ on $G$ and $c'=\pglF_{-\lambda'}$ on $H$.

\begin{lemma}[Negation]\label{pgl:neg}
$\Lambda(-\lambda)=-\Lambda(\lambda)$, $\mu(-\lambda)=-\mu(\lambda)$, $-\lambda$ has the same gap as
$\lambda$, and $\pglF_{-\lambda}=\aubar{\pglF_\lambda}$.
\end{lemma}
\begin{proof}
Replacing $\lambda$ by $-\lambda$ exchanges the sets $\{j:\lambda_j<\lambda_i\}$ and $\{j:\lambda_j>
\lambda_i\}$, so $\rho$, $\rho_{\mathrm c}$, $\rho_{\mathrm n}$ change sign; this gives the first two
claims, and the third is clear. For unitary $u$, Lemma~\ref{pco:minors}(a) gives $\pdsm_{-w}(u)=
\pdsm_w(u^{-1})=\pdsm_w(u^*)=\aubar{\pdsm_w(u)}$. Hence $|\pdsm_{-\Lambda}|=|\pdsm_\Lambda|$ on $K$,
$d_{-\Lambda}=d_\Lambda$, and the numerator of \eqref{pco:eq-F} for $-\lambda$ is
$\aubar{\pdsm_\Lambda(ku)}\pdsm_\Lambda(u)$, the conjugate of that for $\lambda$. The denominators are
$\pdsm_\mu(w)$ and $\pdsm_{-\mu}(w)$ for positive definite $w=u^{-1}zu\in\auU(J_0)$; they are equal and
positive by Lemma~\ref{pds:fischer} (its statement $\pdsm_{-\mu}(w)=\pdsm_\mu(w)>0$ does not use the
integer $c$).
\end{proof}

So $\aubar{c'}=\pglF_{\lambda'}$ on $H$ and $c=\aubar{\pglF_{-\varsigma}}$ on $G$, and with
\[
\mu_Z:=\operatorname{diag}(\Lambda')\ (\text{weight of }\lambda'),\qquad\mu_V:=\operatorname{diag}(-\Lambda)\
(\text{weight of }-\varsigma),
\]
the integrand of Theorem~\ref{pgl:S} is $\pglF_{\lambda'}(\omega^{-1}(r\oplus1)\omega)\,
\aubar{\pglF_{-\varsigma}(r)}$. The weight $\mu_Z$ is \emph{steep} in the sense of
\eqref{pbp:eq-steep} (zero on $X$, strictly decreasing in $[-3B,-B]$ with alternating signs elsewhere),
and $\mu_V$ is $0$ on $X$ and lies in $[B,3B]$ elsewhere.

\subsection{The stationary position}\label{pgl:ss-pos}

Write $\mu_i=\Lambda'_i$ ($0\le i\le m$) and $\nu_j=-\Lambda_j$ ($1\le j\le m$): integers with
$\mu_0>\dots>\mu_m$, $\nu_1<\dots<\nu_m$, $|\mu_i|\le3B$, $B\le\nu_j\le3B$, so $1\le2B\le\nu_j-\mu_i\le
6B$. Let $p(t)=\prod_i(t-\mu_i)$ (degree $m+1$), $q(t)=\prod_j(t-\nu_j)$ (degree $m$) and $q_j=q/(t-\nu_j)$.
For a real polynomial $F$ let $[F]\in\R^{n+1}$ vanish on $X$ and have the entry
$[F]_i=F(\mu_i)\,|p'(\mu_i)q(\mu_i)|^{-1/2}$ on the coordinate $i$; then $\mu_Z[F]=[tF]$. For real
vectors vanishing on $X$ put $\mathcal B(x,y)=\sum_i(-1)^ix_iy_i$, the form $x^*J_Zy$.

\begin{lemma}\label{pgl:gram}
\textup{(i)} $\mathcal B([F_1],[F_2])=\sum_iF_1F_2(\mu_i)/(p'(\mu_i)q(\mu_i))$.
\textup{(ii)} $\mathcal B([q],[q])=1$, $\mathcal B([q_j],[q])=0$, and $\mathcal B([q_j],[q_l])=
(-1)^{j-1}\delta_{jl}\,|q'(\nu_j)|/p(\nu_j)$, with $p(\nu_j)>0$.
\textup{(iii)} $|[q]_i|\le(6B)^{m/2}$, and with $s_j=(p(\nu_j)/|q'(\nu_j)|)^{1/2}$, $|s_j[q_j]_i|\le
(6B)^m$.
\end{lemma}
\begin{proof}
(i) Exactly $i$ of the $\mu$'s exceed $\mu_i$, so $p'(\mu_i)$ has sign $(-1)^i$; $q(\mu_i)=
\prod_j(\mu_i-\nu_j)$ has sign $(-1)^m=+1$. So $|p'(\mu_i)q(\mu_i)|=(-1)^ip'(\mu_i)q(\mu_i)$.
(ii) By (i) and Lemma~\ref{pgm:residue} (with $s=m+1$ nodes): $\mathcal B([q],[q])=\sum_iq(\mu_i)/
p'(\mu_i)$ is the coefficient of $t^m$ in $q$, which is $1$; $\mathcal B([q_j],[q])=\sum_iq_j(\mu_i)/
p'(\mu_i)$ is the coefficient of $t^m$ in $q_j$ (degree $m-1$), which is $0$; and $\mathcal B([q_j],
[q_l])=\sum_iq_j(\mu_i)/((\mu_i-\nu_l)p'(\mu_i))=-q_j(\nu_l)/p(\nu_l)$, as $\deg q_j\le m$ and $\nu_l$ is
not a node. This vanishes for $j\ne l$, and for $j=l$ it is $-q'(\nu_j)/p(\nu_j)$. Now $p(\nu_j)>0$ as
$\nu_j>\mu_i$ for all $i$, and $q'(\nu_j)=\prod_{l\ne j}(\nu_j-\nu_l)$ has sign $(-1)^{m-j}=(-1)^j$.
(iii) The nodes are distinct integers, so $|p'(\mu_i)|\ge1$, $|q'(\nu_j)|\ge1$, and $|q_j(\mu_i)|\le
|q(\mu_i)|$ as $|\mu_i-\nu_j|\ge1$. Hence $[q]_i^2\le|q(\mu_i)|\le(6B)^m$, and $(s_j[q_j]_i)^2\le
p(\nu_j)\,q_j(\mu_i)^2/|q(\mu_i)|\le p(\nu_j)|q_j(\mu_i)|\le(6B)^{m+1}(6B)^{m-1}$.
\end{proof}

\begin{proposition}[Stationary position]\label{pgl:stat}
Let $\iota$ be the $(n+1)\times n$ matrix that is the identity from the coordinates $X$ of $\C^n$ to
those of $\C^{n+1}$ and sends the far coordinate vector $e_j$ of $\C^n$ to $w_j:=s_j[q_j]$, and let
$v=[q]$. Then $w:=(\iota\mid v)\in H$, $\|w\|^2\le\beta_{\mathsf S}$, $\iota^\dagger\mu_Z\iota=\mu_V$
where $\iota^\dagger=J\iota^*J_Z$, and $\iota^\dagger\iota=1$, $\iota^\dagger v=0$.
\end{proposition}
\begin{proof}
By Lemma~\ref{pgl:gram}(ii) and disjointness of supports, the Gram matrix $w^*J_Zw$ is
$J\oplus1=J_Z$ (the coordinate $j$ of $V$ has sign $(-1)^{j-1}$); so $w\in H$, and $w^{-1}=J_Zw^*J_Z$
has the rows $\iota^\dagger$ and $v^*J_Z$, which gives $\iota^\dagger\iota=1$, $\iota^\dagger v=0$.
Since $t\,q_j=\nu_jq_j+q$, $\mu_Zw_j=\nu_jw_j+s_jv$, and $\mu_Z$ vanishes on the $X$-coordinates; so
$\iota^\dagger\mu_Z\iota e_j=\nu_je_j$ and $\iota^\dagger\mu_Z\iota e_l=0$ ($l\in X$), i.e.\
$\iota^\dagger\mu_Z\iota=\mu_V$. All entries of $w$ have modulus at most $(6B)^m$
(Lemma~\ref{pgl:gram}(iii)), so $\|w\|^2\le\|w\|_{\mathrm F}^2\le(n+1)^2(6B)^{2m}\le\beta_{\mathsf S}$ by
Lemma~\ref{pgq:sizes}(c).
\end{proof}

We put $\omega=w^{-1}$ (so $\|\omega\|=\|w\|$ and $\omega^{-1}(r\oplus1)\omega=w(r\oplus1)w^{-1}$) and,
for $r\in G$, $g_r=w(r\oplus1)w^{-1}\in H$.

\begin{lemma}[The group of the position]\label{pgl:group}
For $r\in G$: \textup{(a)} $g_r-1=w((r-1)\oplus0)w^{-1}$, $g_r\iota=\iota r$, $g_rv=v$, $r=\iota^\dagger
g_r\iota$; \textup{(b)} $\pbptr(g_r-1)=\pbptr(r-1)$ and $\pbptr(\mu_Z(g_r-1))=\pbptr(\mu_V(r-1))$;
\textup{(c)} if $\|r-1\|\le\frac12$, then $\|g_r-1\|\le\frac12\beta_{\mathsf S}$ and, when also
$\|g_r-1\|\le\frac12$, $\log g_r=w((\log r)\oplus0)w^{-1}$ and $\pbptr(\mu_Z\log g_r)=\pbptr(\mu_V\log r)$;
\textup{(d)} if $u$ is unitary on the coordinates $X$ and $r_u=u\oplus1$, then $g_{r_u}=u\oplus1$;
\textup{(e)} $r\mapsto g_r$ is a homeomorphism of $G$ onto a closed subgroup of $H$, and
$\|g_r\|\le\|w\|^2\|r\|$.
\end{lemma}
\begin{proof}
(a) Immediate from $w(r\oplus1)=(\iota r\mid v)$ and $\iota^\dagger\iota=1$. (b) $\pbptr(D\,w(Y\oplus0)
w^{-1})=\pbptr((Y\oplus0)w^{-1}Dw)$, and the upper left $n\times n$ block of $w^{-1}Dw$ is
$\iota^\dagger D\iota$; take $D=1$, resp.\ $D=\mu_Z$, and use Proposition~\ref{pgl:stat}. (c) The
first bound is $\|w\|\|w^{-1}\|\|r-1\|$. The logarithm series commutes with conjugation by $w$ and
with $Y\mapsto Y\oplus0$ termwise (both series converge, $\|(r-1)^j\oplus0\|=\|(r-1)^j\|$), so
$\log g_r=w((\log r)\oplus0)w^{-1}$ (as $(1+Y)\oplus1$, conjugated, is $g_r$); then argue as in (b).
(d) $\iota$ and $v$ are the identity on $X$ and vanish there off it: $w=1_X\oplus w'$. (e) $g_r$
is continuous in $r$, $r=\iota^\dagger g_r\iota$, and the image is the closed subgroup of $h\in H$
fixing $v$ (an $h$ fixing $v$ preserves $\iota(\C^n)=v^\perp$, and $w^{-1}hw=r\oplus1$ with
$r=\iota^\dagger h\iota\in G$).
\end{proof}

\emph{The integrand.} By \S\ref{pgl:ss-stat}, the integrand of Theorem~\ref{pgl:S} is
\begin{equation}\label{pgl:eq-theta}
\theta(r)=\pglF_{\lambda'}(g_r)\,\aubar{\pglF_{-\varsigma}(r)},\qquad|\theta(r)|=|c'(g_r)|\,|c(r)| .
\end{equation}
The phase difference of \S\ref{pbp:ss-pos} is $\varphi(r)=\pbptr(\mu_Z(g_r-1))-\pbptr(\mu_V(r-1))=0$
by Lemma~\ref{pgl:group}(b): the position is exactly stationary and there is no character.

\subsection{Sizes and the coefficient lemmas of \S\ref{pbp:sec}}\label{pgl:ss-sizes}

Put $D=2n+1$, $\beta=\beta_{\mathsf S}$, $T=D^2\log\beta$, $h_0=\lceil10^{30}T^3\rceil$,
$\rho_0=(200\beta^5)^{-1}$.

\begin{lemma}\label{pgl:sizes}
\textup{(a)} $\beta\ge10^{12}$, $(n+1)^2\le T$, $18D^2\le h_0\le\min(h_V,h_Z)$.
\textup{(b)} For $\pglF=\pglF_{\lambda'}$ on $H$ and $\pglF=\pglF_{-\varsigma}$ on $G$ \textup{(}rank
$N\le n+1$, weight $\Lambda_\pglF$\textup{)}: $\max_i|\Lambda_{\pglF,i}|\le6b$ and $d_{\Lambda_\pglF}\le
e^{2T}$. \textup{(c)} $\rho_0\le1/(4C)$ for $C=3(n+1)(2b+2n+2)$.
\end{lemma}
\begin{proof}
(a) Lemma~\ref{pgq:sizes}(c) gives $\beta\ge10^{12}$ and $\min(h_V,h_Z)\ge500(n+1)^3+10^{30}T^3\ge h_0$;
$\log\beta\ge1$ gives $(n+1)^2\le T$, and $18D^2\le h_0$. (b) $|\Lambda_i|\le3B\le6b$
(Lemma~\ref{pgq:steep}, $B\le2b$ by Lemma~\ref{an:lem-sizes}). Lemma~\ref{pgm:dLambda} gives
$d_\Lambda\le4(24M_\Lambda N)^{2N^2}$ with $M_\Lambda\le6b$: so $d_\Lambda\le4(144b(n+1))^{2(n+1)^2}\le
b^{4(n+1)^2+1}$ (as $144(n+1)\le b$), and $(4(n+1)^2+1)\log b\le2D^2\cdot3m\log b=2T$.
(c) $4C\le48b^2\le200b^{15m}$.
\end{proof}

\begin{lemma}[Transfer]\label{pgl:transfer}
For each of the two functions $\pglF$ of Lemma~\ref{pgl:sizes}(b), with $\pdsK$, $J_0$, the weight
$\Lambda_\pglF$ and the parameter $\lambda_\pglF$ of its group, the following hold.
\begin{enumerate}
\item[(a)] \textup{(Decay)} $|\pglF(g)|\le(2/\|g\|)^{c}$ for integers $0\le c\le h_0$, and
$|\pglF(g)|^4\le(1+\|y\|_{\mathrm F}^2)^{-h_0}$, $y=g-a(g)$ the noncompact part.
\item[(b)] \textup{(Gaussian shape)} If $\|g-1\|_{\mathrm F}\le\delta_1$ with $\delta_1T\le5\cdot10^{-13}$
and $Z=\log g$, there is $\varrho\in[0,1]$ with $|\pglF(g)e^{-\pbptr(\Lambda_\pglF Z)}-\varrho|\le\varrho/10
+e^{-100T}$.
\item[(c)] \textup{(Left $X$-invariance)} If $k_0\in\pdsK$ is the identity off $X$, $\pglF(k_0g)=\pglF(g)$.
\item[(d)] \textup{(Near $1$)} If $\|g-1\|\le\rho_0$, $|\pglF(g)|\ge\frac12$.
\end{enumerate}
\end{lemma}
\begin{proof}
Both functions are $\pglF_\lambda$ for a regular $\lambda\in(\frac{N-1}2+\Z)^N$ ($N=n$, $\lambda=-\varsigma$;
$N=n+1$, $\lambda=\lambda'$), so the general lemmas of \S\ref{pgm:sec} apply (\S\ref{pgm:ss-effect}).
(a) Lemma~\ref{pgm:decay} with $c\le h_0\le\min(h_V,h_Z)$ (Lemma~\ref{pgl:sizes}(a); Lemma~\ref{pgl:neg}
for the gap of $-\varsigma$); the second bound uses $|\pglF|\le|\det a(g)|^{-h_0/2}$, the intermediate
bound of the proof of Lemma~\ref{pco:decay}.
(b) Proposition~\ref{pgm:gauss}, whose hypotheses $T\ge N^2$ and $d_{\Lambda_\pglF}\le e^{2T}$ are
Lemma~\ref{pgl:sizes}(a),(b).
(c) Lemma~\ref{pgm:leftX}: $X$ has sign $+$ and $\Lambda'=0$, $-\Lambda=0$ there (Lemma~\ref{pgm:coords}).
(d) $|\pglF|=|c'|$, resp.\ $|c|$ (Theorem~\ref{pco:FJ}, gaps $\ge N^2+1$). Lemma~\ref{pgm:lip} with
$C=3N(L+2N)$: for $c$, $N=n$, $L=L_\tau\le2b$; for $c'$, $N=n+1$ and $L=\max_i|\lambda'_i|\le2b$ (the entries
are $\frac k2-\frac12-i$ ($0\le i\le k-2$) or $-\kappa(\tau)-j\Delta+\frac12$ ($1\le j\le m+1$), Lemma~\ref{pgm:coords}; the first
kind has $|\frac k2-\frac12-i|\le\frac k2\le b$, and the second has
$|\lambda'_i|\le\kappa(\tau)+(m+1)\Delta\le b+(2m+1)\cdot2r^8\le b+x^{42}\le2b$ by $\kappa(\tau)\le b$
(Lemma~\ref{an:lem-sizes}(S5)), $\Delta=2r^8$, $m\le2\sqrt x$, $r\le3x^5$ and $x^{42}\le b$ (S1)). In
both cases $C\le3(n+1)(2b+2n+2)$, and $\rho_0\le1/(4C)$ by Lemma~\ref{pgl:sizes}(c).
\end{proof}

\subsection{Rows, localisation and the cone}\label{pgl:ss-cone}

Norms are Frobenius norms in this subsection. For $r\in G$ let $x=g_r-1$, $x'=r-1$, $y$, $y'$ the
noncompact parts of $g_r$, $r$, and $P_{\mathrm{off}}$ the coordinate projection of $\C^{n+1}$ onto the
coordinates outside $X$.

\begin{lemma}[Rows]\label{pgl:rows}
$\|x\|^2=2\|y\|^2-2\Rel\pbptr x$, $\|x'\|^2=2\|y'\|^2-2\Rel\pbptr x'$, and
$\|P_{\mathrm{off}}x\|^2\le6(\|y\|^2+\|y'\|^2)$, $\|xP_{\mathrm{off}}\|^2\le6(\|y\|^2+\|y'\|^2)$.
\end{lemma}
\begin{proof}
The row identity $2\Rel x_{ii}=2\sum_j|y_{ij}|^2-\sum_j|x_{ij}|^2$ of the proof of
Lemma~\ref{pbp:rows} holds in any $\auU(J_0)$; summing gives the two identities. Multiply the row
identity of $g_r$ by $-\mu_{Z,i}\ge0$ and sum: $\sum_i(-\mu_{Z,i})\sum_j|x_{ij}|^2=2\sum_i(-\mu_{Z,i})
\sum_j|y_{ij}|^2+2\Rel\pbptr(\mu_Zx)$; the left side is $\ge B\|P_{\mathrm{off}}x\|^2$ and the first
term on the right $\le6B\|y\|^2$. Multiply that of $r$ by $\mu_{V,i}\in[0,3B]$:
$2\Rel\pbptr(\mu_Vx')\le6B\|y'\|^2$. As $\pbptr(\mu_Zx)=\pbptr(\mu_Vx')$ (Lemma~\ref{pgl:group}(b)),
$B\|P_{\mathrm{off}}x\|^2\le6B(\|y\|^2+\|y'\|^2)$. For the columns apply this to $r^{-1}$:
$g_{r^{-1}}-1=J_Zx^*J_Z$, and noncompact parts of inverses have the same norms.
\end{proof}

\begin{lemma}[Localisation]\label{pgl:loc}
Let $0<\rho_1\le\frac16$ and $\|y\|^2+\|y'\|^2\le\rho_1^2$. There is $t\in G$ with $\theta(t)=
\theta(r)$, $\|g_t-1\|\le4\rho_1$ and $\|t-1\|\le4\rho_1$.
\end{lemma}
\begin{proof}
As in Lemma~\ref{pbp:localise}: the block $\alpha$ of $g_r$ on $X\times X$ satisfies $\|\alpha^*\alpha-1\|
\le\|P_{\mathrm{off}}x\|^2\le6\rho_1^2\le\rho_1$ (Lemma~\ref{pgl:rows}), so $\alpha=k_1z_1$, $k_1$
unitary, $\|z_1-1\|\le\rho_1$. Let $u=k_1^{-1}$ and $t=r_ur$ (Lemma~\ref{pgl:group}(d)); then
$g_t=(u\oplus1)g_r$ has the block $z_1$ on $X$, $\pglF_{\lambda'}(g_t)=\pglF_{\lambda'}(g_r)$ and
$\pglF_{-\varsigma}(t)=\pglF_{-\varsigma}(r)$ by Lemma~\ref{pgl:transfer}(c), so $\theta(t)=\theta(r)$;
left multiplication by the block diagonal unitaries $u\oplus1$ keeps $\|y\|$, $\|y'\|$. Counting rows
off $X$, columns off $X$ and the block on $X$, $\|g_t-1\|^2\le13\rho_1^2$. By the two identities of
Lemma~\ref{pgl:rows} and $\pbptr(g_t-1)=\pbptr(t-1)$, $\|t-1\|^2=\|g_t-1\|^2-2\|y_t\|^2+2\|y'_t\|^2\le
15\rho_1^2$.
\end{proof}

\begin{proposition}[Cone]\label{pgl:cone}
For every $r\in G$, $|\theta(r)|\le2\Rel\theta(r)+12e^{-100T}$.
\end{proposition}
\begin{proof}
Let $E=e^{-100T}$, $\rho_1=10^{-13}/T$. If $\|y\|^2+\|y'\|^2\ge\rho_1^2$, Lemma~\ref{pgl:transfer}(a)
gives $|\theta|^4\le(1+\|y\|^2)^{-h_0}(1+\|y'\|^2)^{-h_0}\le(1+\rho_1^2)^{-h_0}\le E^4$ (as in
Proposition~\ref{pbp:cone}), so $|\theta|\le2\Rel\theta+12E$. Otherwise take $t$ by
Lemma~\ref{pgl:loc}, $\delta_1=4\rho_1$, $Z=\log g_t$, $Z'=\log t$. By Lemma~\ref{pgl:group}(c),
$\pbptr(\mu_ZZ)=\pbptr(\mu_VZ')$, and this is imaginary ($Z'\in\pdsg$ has imaginary diagonal). Hence
$\theta(t)=u_1u_2$ with $u_1=\pglF_{\lambda'}(g_t)e^{-\pbptr(\mu_ZZ)}$, $u_2=\aubar{\pglF_{-\varsigma}(t)
e^{-\pbptr(\mu_VZ')}}$. By Lemma~\ref{pgl:transfer}(b) there are $\varrho,\varrho'\in[0,1]$ with
$|u_1-\varrho|\le\varrho/10+E$, $|u_2-\varrho'|\le\varrho'/10+E$, and Lemma~\ref{pgm:conearith} gives the
claim.
\end{proof}

\subsection{Proof of Theorem~\ref{pgl:S}}\label{pgl:ss-proof}

Operator norms again. Let $\pglLarge=\{r:\|g_r\|\le\beta+1,\ \|r\|\le\beta+1\}$ and
$\pglSmall=\{r:\|g_r-1\|\le\rho_0,\ \|r-1\|\le\rho_0\}$; they are compact (Lemma~\ref{pgl:group}(e)).

\begin{lemma}\label{pgl:cover}
$\pgqvol(\pglLarge)\le\beta^{16D^2}\pgqvol(\pglSmall)$, and $\{\|r-1\|\le1\}\subseteq\pglLarge$.
\end{lemma}
\begin{proof}
Lemma~\ref{pgm:cover} with $\Gamma=G$, $\mathcal J(r)=g_r\oplus r\in\mathrm{GL}_D(\C)$, $T'=\beta+1$, $\rho=\rho_0$
(Remark~\ref{pgm:rem-apply}): $\|\mathcal J(r)\|=\max(\|g_r\|,\|r\|)=\|\mathcal J(r)^{-1}\|$, $r\in\pglSmall$
iff $\|\mathcal J(r)-1\|\le\rho_0$, and $(8(\beta+1)^2/\rho_0)^{2D^2}\le\beta^{16D^2}$ (as $1/\rho_0\le200\beta^5$ and $\beta\ge10^{12}$). If $\|r-1\|\le1$ then $\|r\|\le2$ and $\|g_r\|\le1+\|w\|^2\|r-1\|\le1+\beta$
(Lemma~\ref{pgl:group}(a), Proposition~\ref{pgl:stat}).
\end{proof}

\begin{proof}[Proof of Theorem~\ref{pgl:S}]
Take $\omega$ as above (and $p_Z\omega p_Z^{-1}$ in $H_\tau$); $\|\omega\|^2\le\beta$. Let
$\mathcal G(r)=(2/\|r\|)^{9D^2}$ and $E=\beta^{-50D^2}$, so $E^2=e^{-100T}$. \emph{Decay.} By
Lemma~\ref{pgl:transfer}(a) with $c=9D^2,18D^2\le h_0$, $|\theta|\le|c(r)|\le\mathcal G$ and
$|\theta|\le\mathcal G^2$; by Lemma~\ref{pgq:balls} ($9D^2\ge8n^2$), $\int_G\mathcal G\le2^{9D^2}\cdot
2\cdot16^{n^2}V_\tau\le16^{9D^2}V_\tau$, so $\theta$ is integrable. \emph{Cone.} $\Omega:=|\theta|-
2\Rel\theta\le12E^2$ (Proposition~\ref{pgl:cone}) and $\Omega\le3\mathcal G^2$, so $\Omega\le6E\mathcal G$.
\emph{Small.} On $\pglSmall$, $|\theta|\ge\frac14$ by Lemma~\ref{pgl:transfer}(d). \emph{Integration.}
With $I=|\int_G\theta|$, $V=\pgqvol(\pglLarge)\ge V_\tau$, $W=\pgqvol(\pglSmall)$:
$\frac14W\le\int|\theta|\le2I+6E\cdot16^{9D^2}V_\tau\le2I+6E\cdot16^{9D^2}V$. By Lemma~\ref{pgl:cover},
$V\le8\beta^{16D^2}I+24\cdot16^{9D^2}\beta^{-34D^2}V$, and $24\cdot16^{9D^2}\beta^{-34D^2}\le\frac12$; as
$V<\infty$, $V_\tau\le V\le16\beta^{16D^2}I\le\beta^{20D^2}I=\beta_{\mathsf S}^{N_{\mathsf S}}I$.
\end{proof}



%% ===== proofs/C-grid13.tex =====
% proofs/C-grid13.tex
% The coefficient lemmas of B-period.tex / B-discrete-series.tex in arbitrary rank, as used by
% proofs/C-grid12.tex (pgl:), and the two coordinate facts taken from pgq:steep.
% Labels start with pgm:. Conventions and macros: those of proofs/C-grid12.tex and B-period.tex.
\providecommand{\auU}{\mathrm U}
\providecommand{\aubar}[1]{\overline{#1}}
\providecommand{\pdsK}{\mathbf K}
\providecommand{\pdsg}{\mathfrak g}
\providecommand{\pdsm}{m}
\providecommand{\pbptr}{\operatorname{tr}}
\providecommand{\pbpF}[1]{\|#1\|_{\mathrm F}}
\providecommand{\pbpQ}{\mathcal Q}
\providecommand{\pglF}{\mathsf F}
\providecommand{\pgmvol}{\operatorname{vol}}
\providecommand{\pdsk}{\mathfrak k}
\providecommand{\pdsp}{\mathfrak p}
\providecommand{\pcoH}{\mathsf H}

\section{The coefficient lemmas in arbitrary rank}\label{pgm:sec}

Section~\ref{pgl:sec} applies several lemmas of \S\S\ref{pbp:sec}, \ref{pds:sec} to functions other
than the coefficient $c$ for which they are stated: to $\pglF_{-\varsigma}$ on $G$ (rank $n$) and to
$\pglF_{\lambda'}$ on $H$ (rank $n+1$). This section states each of these lemmas for an arbitrary
unitary group $\auU(J_0)$ of rank $N$ and an arbitrary parameter, with the hypotheses that are
actually needed and explicit constants, and proves it by going through its proof and
recording every place where the rank, the parameter or the group enters. It also proves the two
facts on coordinates that \S\ref{pgl:ss-stat} takes from Lemma~\ref{pgq:steep}, and re-checks the
numerical inputs of Theorem~\ref{pgl:S}. Two inputs of \S\ref{pgl:sec} are used in the form proved here: the bound needed from
Lemma~\ref{pco:decay} is the intermediate bound of its proof (Lemma~\ref{pgm:decay}), and the constant of
Lemma~\ref{pbp:dLambda} is the one for the rank at hand (Lemma~\ref{pgm:dLambda}); with them
Theorem~\ref{pgl:S} holds as stated (\S\ref{pgm:ss-effect}).

\subsection{Setting}\label{pgm:ss-set}

Throughout, $J_0=\operatorname{diag}(\varepsilon_1,\dots,\varepsilon_N)$ with $\varepsilon_i=\pm1$,
$G_0=\auU(J_0)$, $\pdsK_0=G_0\cap\auU(N)$ (the block diagonal elements), $\pdsg_0$ the Lie algebra,
and $\lambda\in(\frac{N-1}2+\Z)^N$ has distinct entries (\emph{regular}). From $\lambda$ we form, as in
\S\ref{pco:ss-stat}, $\rho,\rho_{\mathrm c},\rho_{\mathrm n}$, $\Lambda=\Lambda(\lambda)=\lambda+
\rho_{\mathrm n}-\rho_{\mathrm c}\in\Z^N$, $\mu=\lambda+\rho=\Lambda+2\rho_{\mathrm c}\in\Z^N$, the gap
$h=\min\{|\lambda_i-\lambda_j|:\varepsilon_i\ne\varepsilon_j\}$ (put $h=\infty$ if all signs agree),
$L=\max_i|\lambda_i|$, $M_\Lambda=\max(1,\max_i|\Lambda_i|)$, $d_\Lambda=(\int_{\pdsK_0}
|\pdsm_\Lambda|^2)^{-1}$, and $\pglF_\lambda$, the function $F$ of \eqref{pco:eq-F} on $G_0$. Norms:
$\|\cdot\|$ operator norm, $\pbpF\cdot$ Frobenius norm, $a(x)=\frac12(x+J_0xJ_0)$, and the noncompact
part of $g\in G_0$ is $y=g-a(g)$. We record the elementary bounds
\begin{equation}\label{pgm:eq-rho}
|\rho_i|\le\tfrac{N-1}2,\quad|2\rho_{\mathrm c,i}|\le N-1,\quad|\rho_{\mathrm n,i}|+|\rho_{\mathrm c,i}|
\le\tfrac{N-1}2,\quad|\Lambda_i|\le L+N,
\end{equation}
each half sum running over at most $N-1$ indices $j\ne i$.

\begin{lemma}[What is stated in general]\label{pgm:general}
The following statements are stated, and proved, for an arbitrary $G_0$ and an arbitrary regular
$\lambda\in(\frac{N-1}2+\Z)^N$, with no condition on the gap: Lemma~\ref{pco:minors},
Corollary~\ref{pco:equiv}, Lemma~\ref{pco:density}, Proposition~\ref{pco:indep}, Lemmas~\ref{pco:local}
and~\ref{pco:K} \textup{(}in particular \eqref{pco:eq-left}\textup{)}, Lemma~\ref{pds:blockdet}, the
first assertion $\pdsm_{-\mu}(w)=\pdsm_\mu(w)>0$ of Lemma~\ref{pds:fischer}, and, for integers
$0\le c\le h$, its second assertion. Lemmas~\ref{pbp:third}, \ref{pbp:minorsize},
Propositions~\ref{pbp:second} and~\ref{pbp:minorconj} are stated in arbitrary dimension, with
constants independent of it.
\end{lemma}
\begin{proof}
For \S\ref{pco:sec} this is the convention of \S\ref{pco:ss-stat} ($G=\auU(p,q)$ of arbitrary rank $n$,
arbitrary regular $\lambda$); the gap enters \S\ref{pco:sec} only in Lemma~\ref{pco:decay} and
Theorem~\ref{pco:FJ}. Lemma~\ref{pds:blockdet} and the first assertion of Lemma~\ref{pds:fischer}
use only $g^*J_0g=J_0$ (facts (F2), (F3) of that proof use $w\in G_0$ positive definite and nothing
else); the second assertion uses the gap only through ``$|\lambda_i-\lambda_j|\ge h\ge c$ if
$\varepsilon_i\ne\varepsilon_j$'' and the sign of $\rho_i-\rho_j$, valid for any regular $\lambda$.
The proofs in \S\ref{pbp:ss-minor} begin with ``the dimension is arbitrary; no constant depends on
it'', and indeed use only norm inequalities valid in every dimension.
\end{proof}

\subsection{Decay and the bound far from \texorpdfstring{$\pdsK_0$}{K}}\label{pgm:ss-decay}

\begin{lemma}[Decay, general form of Lemma~\ref{pco:decay} and Lemma~\ref{pbp:far}]\label{pgm:decay}
For $g\in G_0$ and every integer $0\le c\le h$,
\[
|\pglF_\lambda(g)|\le|\det a(g)|^{-c/2}\le(2/\|g\|)^c,\qquad|\pglF_\lambda(g)|^4\le(1+\pbpF y^2)^{-c},
\]
where $y$ is the noncompact part of $g$.
\end{lemma}
\begin{proof}
This is the proof of Theorem~\ref{pds:decay}, with $\pglF_\lambda$ in place of $c_\tau$: write $g=kz$,
$D=\det a(z)=|\det a(g)|$; for $u\in\pdsK_0$, $w_u=u^{-1}zu$ is positive definite in $G_0$ with $\det
a(w_u)=D$ (Lemma~\ref{pds:blockdet}), so $\pdsm_{-\mu}(w_u)=\pdsm_\mu(w_u)\ge D^{c/2}$ by
Lemma~\ref{pds:fischer} (valid by Lemma~\ref{pgm:general}). Then $|\pglF_\lambda(g)|\le d_\Lambda
D^{-c/2}\int\frac12(|\pdsm_\Lambda(ku)|^2+|\pdsm_\Lambda(u)|^2)du=D^{-c/2}$, by invariance of $du$ and the
definition of $d_\Lambda$. The second inequality is Lemma~\ref{pds:blockdet}. For the last, as in
Lemma~\ref{pbp:far}: the block diagonal part of $g^*J_0g=J_0$ is $a^*a-y^*y=1$ ($J_0$ commutes with
$a=a(g)$ and anticommutes with $y$), so $|\det a|^2=\det(1+y^*y)\ge1+\pbptr(y^*y)=1+\pbpF y^2$, and
$|\pglF_\lambda|^4\le|\det a|^{-2c}\le(1+\pbpF y^2)^{-c}$. Neither $N$ nor $\lambda$ enters beyond
Lemma~\ref{pgm:general}.
\end{proof}

Lemma~\ref{pbp:far} and Lemma~\ref{pgl:transfer}(a) use the bound $|\det a(g)|^{-c/2}$ of
Lemma~\ref{pgm:decay}.

\subsection{The Gaussian shape}\label{pgm:ss-gauss}

\begin{lemma}[General form of Lemma~\ref{pbp:reduced}]\label{pgm:reduced}
Let $Z\in\pdsg_0$ with $\pbpF Z\le\frac1{20}$; let $u_-$ \textup(resp.\ $u_+$\textup) keep the entries
$Z_{ij}$ with $\varepsilon_i=\varepsilon_j$ and $\lambda_i<\lambda_j$ \textup(resp.\ $\lambda_i>
\lambda_j$\textup) and vanish elsewhere, and $\hat g=e^{-u_-}e^Ze^{-u_+}$. Then
\[
\pglF_\lambda(e^Z)=d_\Lambda\int_{\pdsK_0}\frac{|\pdsm_\Lambda(k)|^2}{\pdsm_{-\mu}(k^{-1}\hat gk)}\,dk .
\]
\end{lemma}
\begin{proof}
We go through the proof of Lemma~\ref{pbp:reduced}. (1) $x=e^Z$ has $\|x-1\|\le e^{1/20}-1<1$, so
$x+x^*=2+(x-1)+(x-1)^*$ is positive definite and $1\in A_x$; with the polar factor $k_x\in A_x$,
Proposition~\ref{pco:indep} gives $\pglF_\lambda(x)=d_\Lambda I(x;k_x)=d_\Lambda I(x;1)=d_\Lambda
\mathcal J(1,1)$. The proof invokes Theorem~\ref{pco:FJ} at this point only to replace $c$ by
$F$; for $\pglF_\lambda$ this step is the definition, and no gap is needed. (2) Step (iv) of the proof
of Proposition~\ref{pco:indep} gives both $\partial_1\mathcal J=0$ and $\partial_2\mathcal J=0$ on $U_x$,
for any $\pcoH$ (rank-free). (3) $\|u_\pm\|\le\pbpF Z$, so $e^{tu_+}x^{-1}e^{tu_-}$ is within
$e^{3/20}-1<1$ of $1$ for $t\in[0,1]$, hence accretive; $e^{tu_\pm}$ are block diagonal (only entries
with $\varepsilon_i=\varepsilon_j$), so the path lies in $U_x$ and $\mathcal J(1,1)=\mathcal J(e^{u_-},
e^{u_+})$. (4) $b=e^{u_-}$ has unit diagonal and $b_{cd}=0$ unless $c=d$ or ($\varepsilon_c=
\varepsilon_d$, $\lambda_c<\lambda_d$); for such pairs $\Lambda_c\le\Lambda_d$ by
Lemma~\ref{pco:minors}(c), which holds for every regular $\lambda$. Lemma~\ref{pco:minors}(b) (left
version; $b$ is triangular for the order of decreasing $\lambda$, which is a total order since
$\lambda$ is regular) gives $\pdsm_\Lambda(bk)=\pdsm_\Lambda(k)$; likewise $\pdsm_\Lambda(k^*e^{u_+})=
\aubar{\pdsm_\Lambda(k)}$. (5) $(k^*e^{u_+}x^{-1}e^{u_-}k)^{-1}=k^{-1}\hat gk$ and
Lemma~\ref{pco:minors}(a). The rank enters nowhere; the parameter enters only through its
regularity (total order, Lemma~\ref{pco:minors}(c)).
\end{proof}

\begin{lemma}[General form of Lemma~\ref{pbp:dLambda}]\label{pgm:dLambda}
$d_\Lambda\le4\,(24M_\Lambda N)^{2N^2}$.
\end{lemma}
\begin{proof}
We go through the proof of Lemma~\ref{pbp:dLambda}, with $r$ replaced by $N$ and $3b$ by $M_\Lambda$.
Let $\varepsilon_0=1/(6M_\Lambda N)\le\frac16$. By \eqref{pds:eq-m}, $\pdsm_\Lambda(k)=(\det k)^{-M_\Lambda}
\prod_{t=1-M_\Lambda}^{M_\Lambda}\Delta_{S_t(\Lambda)}(k)$ (factors with $t\le-M_\Lambda$ are $\det k\cdot
(\det k)^{-1}$, those with $t>M_\Lambda$ are $1$): $2M_\Lambda$ minors. If $k\in\pdsK_0$ and
$\|k-1\|\le\varepsilon_0$, then $|\det k|=1$, and $\Delta_S(k)=\det(1+(k-1)_{SS})$ is a product of at
most $N$ factors $1+\xi$ with $|\xi|\le\|(k-1)_{SS}\|\le\varepsilon_0\le\frac12$ (eigenvalues), so
$|\Delta_S(k)|\ge e^{-2N\varepsilon_0}$ and $|\pdsm_\Lambda(k)|\ge e^{-4M_\Lambda N\varepsilon_0}=
e^{-2/3}\ge\frac12$. A maximal family $k_1,\dots,k_\nu\in\pdsK_0$ with $\|k_ik_j^{-1}-1\|>
\varepsilon_0$ ($i\ne j$) has $\|k_i-k_j\|>\varepsilon_0$; the balls of radius $\varepsilon_0/2$ about the
$k_i$ in $M_N(\C)$ (real dimension $2N^2$) are disjoint and lie in the ball of radius $2$, so
$\nu\le(4/\varepsilon_0)^{2N^2}$. By maximality $\pdsK_0$ is covered by the $\nu$ right translates
$\{\|k-1\|\le\varepsilon_0\}k_i$, each of $dk$-mass equal to that of $\{\|k-1\|\le\varepsilon_0\}$; so
this ball has mass $\ge1/\nu$, and $d_\Lambda^{-1}\ge\frac14\nu^{-1}$, i.e.\ $d_\Lambda\le4(4/
\varepsilon_0)^{2N^2}$. The rank enters through the number of factors of a minor ($N$), the
dimension $2N^2$, and $\varepsilon_0$; the parameter only through $M_\Lambda$.
\end{proof}

\begin{proposition}[General form of Proposition~\ref{pbp:gauss}]\label{pgm:gauss}
Let $T\ge N^2$ with $d_\Lambda\le e^{2T}$. Let $g\in G_0$ with $\pbpF{g-1}\le\delta_1$, $\delta_1T\le
5\cdot10^{-13}$, and $Z=\log g$. There is $\varrho\in[0,1]$ with
\[
\bigl|\pglF_\lambda(g)\,e^{-\pbptr(\Lambda Z)}-\varrho\bigr|\le\varrho/10+e^{-100T}.
\]
Here $\Lambda$ is regarded as the diagonal matrix $\operatorname{diag}(\Lambda)$; no condition on the
gap is needed.
\end{proposition}
\begin{proof}
We go through the proof of Proposition~\ref{pbp:gauss}, naming each place where $N$ or $\lambda$
enters. (1) $\delta=2\delta_1$; as $\pbpF{g-1}\le\frac12$, $\pbpF Z\le-\log(1-\pbpF{g-1})\le\delta$
(Frobenius is submultiplicative), $Z\in\pdsg_0$ by the argument of Step~2 of Lemma~\ref{pds:lip}
(rank-free), and $\delta T\le10^{-12}$; as $T\ge N^2\ge1$, $\delta\le10^{-12}$. (2)
Lemma~\ref{pgm:reduced}. (3) With $w=-\mu$, $M'=\max|w_i|$: $\pdsm_w(y)=(\det y)^{-M'}\prod_{t=1-M'}^{M'}
\Delta_{S_t(w)}(y)$. For $i\in S_t(w)$, $j\notin S_t(w)$, $\mu_i<\mu_j$, hence $\lambda_i<\lambda_j$
(Lemma~\ref{pco:minors}(c)); so Proposition~\ref{pbp:minorconj} (arbitrary dimension) applies to each
factor with $o=\lambda$ (injective: regularity). Counting $t$ as stated,
$\pdsm_w(k^{-1}\hat gk)=\exp(-\pbptr(\mu Z)+\frac14\pbpQ_w(Z)+r(k))$ with $|r(k)|\le10^3\delta
(\pbpQ_w(Z)+\pbpQ_w(k))$, where $\pbpQ_w(Y)=\sum_{i,j}|w_i-w_j||Y_{ij}|^2$. (4) \emph{Rank enters.} By
\eqref{pgm:eq-rho}, $|\pbptr(2\rho_{\mathrm c}Z)|\le(N-1)\sum_i|Z_{ii}|\le(N-1)\sqrt N\pbpF Z\le N^2\delta$;
$|w_i-w_j|\le|\Lambda_i-\Lambda_j|+2(N-1)$, and $\sum_{ij}|k_{ij}|^2=N$, so $\pbpQ_w(k)\le\pbpQ_\Lambda(k)+
2N^2$. With $q=\pbpQ_w(Z)/4$: $\pdsm_w(k^{-1}\hat gk)=\exp(q-\pbptr(\Lambda Z)-r'(k))$, $|r'(k)|\le N^2
\delta+10^3\delta(4q+\pbpQ_\Lambda(k)+2N^2)\le10^4\delta\,s(k)$, $s(k)=q+\pbpQ_\Lambda(k)+N^2$. (5) As
$d_\Lambda\int|\pdsm_\Lambda|^2=1$, $\pglF_\lambda(g)e^{-\pbptr(\Lambda Z)}-e^{-q}=d_\Lambda\int
|\pdsm_\Lambda(k)|^2e^{-q}(e^{r'(k)}-1)dk$. (6) Lemma~\ref{pbp:minorsize} (any dimension), multiplied
over $t$, gives $|\pdsm_\Lambda(k)|^2\le e^{-\pbpQ_\Lambda(k)/2}$. (7) If $s(k)\le5\cdot10^4T$, then
$|r'|\le5\cdot10^8\delta T\le5\cdot10^{-4}$ and $|e^{r'}-1|\le\frac1{10}$. Otherwise $|r'|\le s/4$ (as
$10^4\delta\le\frac14$) and $|\pdsm_\Lambda|^2e^{-q}|e^{r'}-1|\le 2e^{-\pbpQ_\Lambda/2-q+s/4}\le
2e^{-s/4+N^2/2}\le 2e^{-12500T+T/2}\le e^{-5100T}$ (using $|e^{r'}-1|\le e^{|r'|}+1\le2e^{s/4}$); \emph{this is the place where $T\ge N^2$ is
used}. (8) Integrating, the difference is at most $e^{-q}/10+d_\Lambda e^{-5100T}\le e^{-q}/10+
e^{-5098T}$ by the hypothesis $d_\Lambda\le e^{2T}$; take $\varrho=e^{-q}$. The constants are those of
the proof of Proposition~\ref{pbp:gauss}, which uses $r^2\le T$ and $d_\Lambda\le e^{2T}$, and these are replaced by the hypotheses.
\end{proof}

\subsection{Left invariance and the lower bound near \texorpdfstring{$1$}{1}}\label{pgm:ss-near}

\begin{lemma}[General form of Lemma~\ref{pbp:leftX}]\label{pgm:leftX}
Let $X$ be a set of coordinates on which $\varepsilon\equiv+1$ and $\Lambda=0$. If $k_0\in\pdsK_0$ is the
identity on the coordinates outside $X$, then $\pglF_\lambda(k_0g)=\pglF_\lambda(g)$ for $g\in G_0$.
\end{lemma}
\begin{proof}
The proof uses \eqref{pco:eq-left} (general, Lemma~\ref{pgm:general}), $\Lambda=0$ on $X$, and
that $k_0$ is block diagonal, which holds as $X$ lies in one sign class. For $t\ge1$, $S_t(\Lambda)\cap
X=\emptyset$ and the rows of $k_0Y$ in $S$ are those of $Y$; for $t\le0$, $S_t(\Lambda)\supseteq X$, and
$(k_0)_{ij}=0$ for $i\in S$, $j\notin S$ (as $j\notin X$), so $\Delta_S(k_0Y)=\det k_0\,\Delta_S(Y)$, and
this factor cancels against $(\det k_0Y)^{-1}$ in \eqref{pds:eq-m}. Rank and parameter do not enter.
\end{proof}

\begin{lemma}[General form of Lemma~\ref{pds:lip}]\label{pgm:lip}
Let $\pi$ be a discrete series of $G_0$ with Harish-Chandra parameter $\lambda$, $v_0$ a unit vector of
its lowest $\pdsK_0$-type, of highest weight $\Lambda$, and $c_0=\langle v_0,\pi(\cdot)v_0\rangle$. Put
$C=3N(L+2N)$. If $\|g-1\|\le1/(4C)$, then $|c_0(g)|\ge\frac12$. If $h\ge N^2+1$, the same holds for
$\pglF_\lambda$.
\end{lemma}
\begin{proof}
Step~1 of the proof of Lemma~\ref{pds:lip} uses \eqref{pds:eq-casimir} for $\pi$ (published, valid for
any discrete series of any $\auU(p,q)$: the Casimir value depends only on the infinitesimal character
$\lambda$, and the $\pdsk$-Casimir value only on the $\pdsK_0$-type $\Lambda$, a product of two unitary
groups), and the bounds
$|\Lambda_i+\rho_{\mathrm c,i}|=|\lambda_i+\rho_{\mathrm n,i}|\le L+N$, $|\rho|^2\le N^3/4$,
$|\Lambda_a|\le L+N$ (\eqref{pgm:eq-rho}): they give $\sum_{\varepsilon_a\ne\varepsilon_b}
\|d\pi(E_{ab})v_0\|^2\le N(L+2N)^2+N^3/4\le2N(L+2N)^2$, hence $\|d\pi(Y)v_0\|\le\sqrt2N(L+2N)\|Y\|$
on $\pdsp_0$ (Cauchy--Schwarz gives $\|d\pi(Y)v_0\|\le\pbpF Y\,(2N)^{1/2}(L+2N)$, and $\pbpF Y\le
\sqrt N\|Y\|$), and $\|d\pi(X)v_0\|\le N(L+N)\|X\|$ on $\pdsk_0$. So
$\|d\pi(X)v_0\|\le C\|X\|$ on $\pdsg_0$. Step~2 is rank-free: $\|\log g\|\le2\|g-1\|$, $\log g\in\pdsg_0$,
and $|c_0(g)-1|\le\int_0^1\|d\pi(\log g)v_0\|dt\le2C\|g-1\|\le\frac12$. The rank enters through $C$
only, the parameter through $L$ only. If $h\ge N^2+1$, Theorem~\ref{pco:FJ} gives $\pglF_\lambda=
c_0$ for the discrete series with parameter $\lambda$ (whose lowest $\pdsK_0$-type is $\Lambda$, by the
formula for the lowest $K$-type quoted in \S\ref{pds:ss-setup}).
\end{proof}

\subsection{The lemmas for two noncompact groups}\label{pgm:ss-two}

Lemmas~\ref{pgl:rows}, \ref{pgl:loc} and Proposition~\ref{pgl:cone} are proved in \S\ref{pgl:ss-cone}
with complete arguments except for two references, ``the last lines of the proof of
Proposition~\ref{pbp:cone}'' and ``the proof of Lemma~\ref{pbp:cover2}'', and Lemma~\ref{pgl:gram}
refers to Lemma~\ref{pbp:residue} with $2n$ nodes. We state and prove the three general facts behind
these references; they involve neither the rank nor the parameter.

\begin{lemma}[Residue sums, any number of nodes]\label{pgm:residue}
Let $\mu_0,\dots,\mu_{s-1}$ be distinct reals, $p(t)=\prod_i(t-\mu_i)$, $F$ a polynomial of degree
$\le s-1$ and $c$ not a node. Then $\sum_iF(\mu_i)/p'(\mu_i)$ is the coefficient of $t^{s-1}$ in $F$, and
$\sum_iF(\mu_i)/((\mu_i-c)p'(\mu_i))=-F(c)/p(c)$.
\end{lemma}
\begin{proof}
The proof of Lemma~\ref{pbp:residue} with $2n$ replaced by $s$: Lagrange interpolation at the $s$
nodes, and $H(t)=(F(t)p(c)-F(c)p(t))/(t-c)$, of degree $\le s-1$ with coefficient $-F(c)$ of $t^{s-1}$.
\end{proof}

\begin{lemma}[The cone inequality]\label{pgm:conearith}
Let $0\le E\le1$, $\varrho,\varrho'\in[0,1]$ and $u_1,u_2\in\C$ with $|u_1-\varrho|\le\varrho/10+E$,
$|u_2-\varrho'|\le\varrho'/10+E$. Then $z=u_1u_2$ satisfies $|z|\le2\Rel z+12E$.
\end{lemma}
\begin{proof}
Put $P=\varrho\varrho'\ge0$. Then $|z-P|\le\varrho|u_2-\varrho'|+\varrho'|u_1-\varrho|+|u_1-\varrho|
|u_2-\varrho'|\le(0.2P+2E)+(0.01P+0.2E+E^2)\le0.21P+3.2E$, using $\varrho,\varrho',E\le1$. Hence
$|z|-2\Rel z\le(P+|z-P|)-2(P-|z-P|)=-P+3|z-P|\le-0.37P+9.6E\le12E$.
\end{proof}

\begin{lemma}[Covering]\label{pgm:cover}
Let $\Gamma$ be a unimodular locally compact group with a Haar measure, $D\ge1$, and $\mathcal J\colon
\Gamma\to\mathrm{GL}_D(\C)$ a continuous injective homomorphism with $\|\mathcal J(x)^{-1}\|=\|\mathcal
J(x)\|$, such that $\{x:\|\mathcal J(x)\|\le T'\}$ is compact for every $T'$. For $T'\ge1$ and
$0<\rho\le1$,
\[
\pgmvol\{x:\|\mathcal J(x)\|\le T'\}\le(8T'^2/\rho)^{2D^2}\,\pgmvol\{x:\|\mathcal J(x)-1\|\le\rho\}.
\]
\end{lemma}
\begin{proof}
Let $x_1,\dots,x_\nu$ lie in $L=\{\|\mathcal J\|\le T'\}$ with $\|\mathcal J(x_ix_j^{-1})-1\|>\rho$ for
$i\ne j$. As $\mathcal J(x_i)-\mathcal J(x_j)=(\mathcal J(x_ix_j^{-1})-1)\mathcal J(x_j)$ and
$\|\mathcal J(x_j)^{-1}\|\le T'$, $\|\mathcal J(x_i)-\mathcal J(x_j)\|>\rho/T'$; the balls of radius
$\rho/(2T')$ about the $\mathcal J(x_i)$ in $M_D(\C)\cong\R^{2D^2}$ are disjoint and lie in the ball of
radius $2T'$, so $\nu\le(4T'^2/\rho)^{2D^2}$. Take such a family maximal and put $S:=\{\|\mathcal
J-1\|\le\rho\}$. If $x\in L$ is not among the $x_i$, it cannot be added, so for some $i$ either
$\|\mathcal J(xx_i^{-1})-1\|\le\rho$ or $\|\mathcal J(x_ix^{-1})-1\|\le\rho$, i.e.\ $xx_i^{-1}\in S$ or
$xx_i^{-1}\in S^{-1}$ (the separation condition need not be symmetric in $i,j$). Hence $L\subseteq
\bigcup_i(S\cup S^{-1})x_i$, and by right invariance and unimodularity ($\pgmvol(S^{-1})=\pgmvol(S)$)
$\pgmvol(L)\le2\nu\pgmvol(S)\le2(4T'^2/\rho)^{2D^2}\pgmvol(S)\le(8T'^2/\rho)^{2D^2}\pgmvol(S)$.
\end{proof}

\begin{remark}[Application in \S\ref{pgl:sec}]\label{pgm:rem-apply}
Lemma~\ref{pgl:cover} is Lemma~\ref{pgm:cover} with $\Gamma=G$, $\mathcal J(r)=g_r\oplus r$, $D=2n+1$,
$T'=\beta+1$, $\rho=\rho_0$: $\mathcal J$ is a continuous injective homomorphism
(Lemma~\ref{pgl:group}(e)), $\|\mathcal J(r)^{\pm1}\|=\max(\|g_r\|,\|r\|)$ as $\|h^{-1}\|=\|h\|$ in $G$
and $H$, $\{\|\mathcal J\|\le T'\}$ is closed in the compact set $\{\|r\|\le T'\}$, and
$(8(\beta+1)^2\cdot200\beta^5)^{2D^2}\le\beta^{16D^2}$ as $\beta\ge10^{12}$. Proposition~\ref{pgl:cone}
ends with Lemma~\ref{pgm:conearith} applied to $u_1,u_2$ (with $E=e^{-100T}$). The proofs of
Lemmas~\ref{pgl:rows} and~\ref{pgl:loc} are complete as stated:
in Lemma~\ref{pgl:rows} the weights used are $-\mu_{Z,i}\in\{0\}\cup[B,3B]$ and $\mu_{V,i}\in\{0\}\cup
[B,3B]$, with zero exactly on $X$ (Lemma~\ref{pgm:coords}); in Lemma~\ref{pgl:loc} the element $r_u=
u\oplus1$ lies in $G$ because $X$ has sign $+$, $g_{r_u}=u\oplus1$ needs that $X$ is the same set in
$\C^n$ and $\C^{n+1}$ (Lemma~\ref{pgm:coords}), and the invariance $\theta(r_ur)=\theta(r)$ is
Lemma~\ref{pgm:leftX} on $H$ and on $G$.
\end{remark}

\subsection{The coordinates of \texorpdfstring{$V$}{V} and \texorpdfstring{$Z$}{Z}}\label{pgm:ss-coords}

Fix $\tau$. By \S\ref{pgq:ss-stat}, $V_\tau$ carries the normalised basis of $Y_0$, which after
conjugation by $p$ (\S\ref{pds:ss-setup}) gives the coordinates $\varrho=0,\dots,r-1$ of $\C^n$ ($n=r$),
attached to the stem coordinates of the frame of $Y_0$; and $Z_\tau$ carries that basis followed by
$\tilde e$, so that after conjugation by $p_Z=p\oplus1$ the coordinates of $\C^{n+1}$ are the same
$\varrho=0,\dots,r-1$ followed by $\varrho=r$, the coordinate of $\tilde e$
(Definition~\ref{pgv:def-spaces}: the frame of $Z=\pgrY^0_0$ is that of $Y_0$ followed by $\tilde e$ as
the stem coordinate $\varrho=r$). The inclusion $r\mapsto r\oplus1$ is the identity on the
coordinates $\varrho<r$. Put $a(\varrho)=\max(\varrho+1-k,0)$.

\begin{lemma}[The two coordinate facts]\label{pgm:coords}
Let $X_V=\{\varrho<r:\Lambda_\varrho=0\}$ and $X_Z=\{\varrho\le r:\Lambda'_\varrho=0\}$, where
$\Lambda=\Lambda(\varsigma_\tau)$ \textup(signs of $V$\textup) and $\Lambda'=\Lambda(\lambda')$,
$\lambda'=\lambda_{(1,0,0),\tau}$ \textup(signs $\epsilon_{(1,0,0)}$\textup). Then:
\begin{enumerate}
\item[(a)] $X_V=X_Z=\{0,1,\dots,k-2\}$, and on these coordinates the signs of $V$ and of $Z$ are $+$;
\item[(b)] at every coordinate $\varrho<r$ the signs of $V$ and $Z$ agree, namely $(-1)^{a(\varrho)}$,
and $\lambda'_\varrho=\varsigma_{\tau,\varrho}+\frac12$;
\item[(c)] the extra coordinate $\varrho=r$ of $Z$ has sign $+$, is not in $X_Z$, and carries the
smallest entry $\lambda'_r=-\kappa(\tau)-(m+1)\Delta+\frac12$ of $\lambda'$ and the smallest value of
$\Lambda'$; in the numbering $0,\dots,m$ of the far coordinates of $Z$ by decreasing $\Lambda'$ used in
\S\ref{pgl:ss-stat}, it is the coordinate $m$. The far coordinates are $\varrho=k-1+i$ \textup($0\le
i\le m$\textup) for $Z$, of sign $(-1)^i$, and $\varrho=k-2+j$ \textup($1\le j\le m$\textup) for $V$, of
sign $(-1)^{j-1}$.
\end{enumerate}
\end{lemma}
\begin{proof}
\emph{Signs.} By \eqref{pgv:eq-sign} for $p=(z,0,0)$, $\epsilon_p(\varrho)=(-1)^{a(\varrho)+\min(a(\varrho),0)}
=(-1)^{a(\varrho)}$ for every stem coordinate, for $z=0$ (the space $Y_0=Y^0_0$, whose signs are
those of Lemma~\ref{pds:chamber}: $(-1)^{\max(0,\varrho+1-k)}$) and for $z=1$ alike. So the signs agree
at $\varrho<r$; they are $+$ for $\varrho\le k-1$; and at $\varrho=r$, $a(r)=r+1-k=m$ (as $r+1=k+m$) is
even, so the sign is $+$.

\emph{Parameters.} By \eqref{pgv:eq-param} with $z=1$, $a=b=0$, $s_p=0$, $\lambda'_\varrho=
\varsigma'_{\tau,\varrho}+\frac12$, and $\varsigma'_{\tau,\varrho}$ is the entry $\varsigma_{\tau,\varrho}$ of
Lemma~\ref{pds:chamber} ($\frac k2-1-\varrho$ for $\varrho\le k-2$, $-\kappa(\tau)-(\varrho+2-k)\Delta$ for
$\varrho\ge k-1$); this gives (b), and $\lambda'_r=-\kappa(\tau)-(m+1)\Delta+\frac12$, the smallest
entry since the entries decrease along $\varrho$ (by $1$ on $\varrho\le k-2$, then from $\lambda'_{k-2}=\frac32-\frac k2$ to
$\lambda'_{k-1}=-\kappa(\tau)-\Delta+\frac12<\frac32-\frac k2$ as $\kappa(\tau)+\Delta>k$, then by $\Delta$).

\emph{Zero sets.} By Lemma~\ref{pds:chamber}, $\Lambda_\varrho=0$ for $\varrho\le k-2$, and by the remark
after it (and Lemma~\ref{pgq:steep}) $\Lambda_\varrho\in[-3B,-B]$, in particular $\ne0$ as $B\ge1$, for
$\varrho\ge k-1$; so $X_V=\{0,\dots,k-2\}$. The computation in the proof of Lemma~\ref{pgq:steep}
(which uses only the signs and parameters just determined) gives $\Lambda'_\varrho=0$ for $\varrho\le k-2$
and $\Lambda'_\varrho\in[-3B,-B]$ for $k-1\le\varrho\le r$; so $X_Z=\{0,\dots,k-2\}=X_V$ as sets of
coordinates of $\C^{n+1}\supset\C^n$, and $r\notin X_Z$. This is (a).

\emph{Order.} On $\varrho\ge k-1$, $\lambda'$ decreases strictly with $\varrho$, in steps $\Delta$, and
$|\Lambda'_\varrho-\lambda'_\varrho|\le|\rho_{\mathrm n,\varrho}|+|\rho_{\mathrm c,\varrho}|\le n/2$
\eqref{pgm:eq-rho}; as $\Delta=2r^8>n+1$, $\Lambda'$ also decreases strictly
with $\varrho$ there, so the order by decreasing $\Lambda'$ is the order of $\varrho$, $\varrho=k-1+i$,
and $\varrho=r$ is $i=m$; its sign is $(-1)^{a(\varrho)}=(-1)^{i}$. The same argument for $V$ (steps
$\Delta$, $|\Lambda_\varrho-\varsigma_{\tau,\varrho}|\le(n-1)/2$) gives $\varrho=k-2+j$, sign
$(-1)^{\varrho+1-k}=(-1)^{j-1}$. This is (c).
\end{proof}

So the two facts that \S\ref{pgl:ss-stat} attributes to Lemma~\ref{pgq:steep} (``$X$ is the same set
of coordinates in $V$ and $Z$'', ``the extra coordinate of $Z$ lies among the far ones'') follow
from the definitions (``on the first $k-1$ coordinates'' in Lemma~\ref{pgq:steep}, the frame being
that of Definition~\ref{pgv:def-spaces}).

\subsection{Effect on Theorem~\ref{pgl:S}}\label{pgm:ss-effect}

We list the parameters with which the lemmas above are applied in \S\ref{pgl:sec} and check every
hypothesis. Notation of \S\ref{pgl:ss-sizes}: $D=2n+1$, $\beta=b^{3m}$, $T=D^2\log\beta$,
$h_0=\lceil10^{30}T^3\rceil$, $\rho_0=(200\beta^5)^{-1}$.

\begin{enumerate}
\item[(1)] \emph{The two functions.} $\pglF_{-\varsigma}$ on $G=\auU(J)$: $N=n$, parameter $-\varsigma$,
weight $-\Lambda$; $\pglF_{\lambda'}$ on $H=\auU(J_Z)$: $N=n+1$, parameter $\lambda'$, weight $\Lambda'$.
Both parameters are regular and lie in $(\frac{N-1}2+\Z)^N$: the entries of $\varsigma$ are integers and
$n=r$ is odd; those of $\lambda'$ are in $\frac12+\Z$ ($k$ even) and $(N-1)/2=r/2\in\frac12+\Z$.
Lemma~\ref{pgl:neg} gives the weights and the gaps.
\item[(2)] \emph{Lemma~\ref{pgm:decay}} with $c\le h_0$ requires $h_0\le\min(h_V,h_Z)$:
Lemma~\ref{pgq:sizes}(c) gives $\min(h_V,h_Z)\ge500(n+1)^3+10^{30}T^3\ge h_0$.
\item[(3)] \emph{Proposition~\ref{pgm:gauss}} requires $T\ge N^2$ and $d_\Lambda\le e^{2T}$. First,
$N^2\le(n+1)^2\le D^2\le T$ as $\log\beta\ge1$. Second, $M_\Lambda\le3B\le6b$ (Lemma~\ref{pgm:coords},
Lemma~\ref{pgq:steep}; $B=\kappa(\tau)+\Delta/2\le b+r^8\le2b$ by Lemma~\ref{an:lem-sizes}(S5), and
$r^8\le b$ as $r\le3x^5$, $b\ge x^{10^{12}}$), so Lemma~\ref{pgm:dLambda} gives $d_\Lambda\le4(144b
(n+1))^{2(n+1)^2}\le b^{4(n+1)^2+1}$, using $144(n+1)\le b$ ((S7)(b): $8\cdot\frac m2\cdot r\le b$ with
$\frac m2\ge100$); and $(4(n+1)^2+1)\log b\le2D^2\cdot3m\log b=2T$. This is the constant of
Lemma~\ref{pgl:sizes}(b), now justified by Lemma~\ref{pgm:dLambda} instead of Lemma~\ref{pbp:dLambda}.
\item[(4)] \emph{Lemma~\ref{pgm:leftX}} with $X$ of Lemma~\ref{pgm:coords}: sign $+$, and $\Lambda'=0$,
$-\Lambda=0$ there.
\item[(5)] \emph{Lemma~\ref{pgm:lip}} for $c$ on $G$ ($N=n$, $L=L_\tau\le2b$) and for $c'$ on $H$
($N=n+1$, parameter $-\lambda'$, $L'=\max_i|\lambda'_i|\le\kappa(\tau)+(m+1)\Delta\le b+x^{42}\le2b$,
using $m\le2\sqrt x$, $\Delta=2r^8\le2\cdot3^8x^{40}$). The gaps exceed $N^2+1$ (item (2)), so
$|\pglF|=|c|$, $|c'|$ and the lemma applies to both $\pglF$. Since $C\le3(n+1)(2b+2n+2)$ and
$4C\le48b^2\le200\beta^5$, $\rho_0\le1/(4C)$.
\item[(6)] \emph{Stationary position.} Lemma~\ref{pgl:gram} uses Lemma~\ref{pgm:residue} with $s=m+1$;
Proposition~\ref{pgl:stat} uses $\|w\|^2\le(n+1)^2(6B)^{2m}\le\beta$, which follows from
Lemma~\ref{pgq:sizes}(c) ($\beta\ge(n+1)^4(6B)^{2m+2}$).
\item[(7)] \emph{Final count.} The proof of Theorem~\ref{pgl:S} uses $9D^2,18D^2\le h_0$, $9D^2\ge8n^2$
(Lemma~\ref{pgq:balls}), $\beta\ge10^{12}$ (Lemma~\ref{pgq:sizes}(c)), Lemma~\ref{pgm:cover} with the
constant $\beta^{16D^2}$, and $16\beta^{16D^2}\le\beta^{20D^2}$, all of which hold here.
\end{enumerate}

Hence every lemma used by \S\ref{pgl:sec} holds with the constants used there, and
Theorem~\ref{pgl:S} holds as stated, with $\beta_{\mathsf S}=b^{3m}$ and $N_{\mathsf S}=20(2n+1)^2$.



%% ===== proofs/C-exchanges3.tex =====
% proofs/C-exchanges3.tex
% Proposition po:rallis (the norm of the lift of a character from U(1) to U(2z) is an L-value, for
% every tame character of modulus Q), statement (IM) of C-selmer.tex for z >= 2 and for z = 1. To be \input after C-determinant.tex,
% C-exchange.tex and C-onelift.tex. Labels of this file start with pex:. Only amsmath/amssymb/amsthm.
\providecommand{\auO}{\mathcal O}
\providecommand{\auU}{\mathrm U}
\providecommand{\auGL}{\mathrm{GL}}
\providecommand{\auHom}{\operatorname{Hom}}
\providecommand{\auNm}{\operatorname{Nm}}
\providecommand{\auA}{\mathbb A}
\providecommand{\auF}{\mathbb F}
\providecommand{\auPl}{\Phi}
\providecommand{\aubar}[1]{\overline{#1}}
\providecommand{\plone}{\mathbf 1}
\providecommand{\auf}[1]{\texttt{#1}}
\providecommand{\pop}{\mathsf p}
\providecommand{\pcR}{R}
\providecommand{\pcW}{\mathsf W}
\providecommand{\psHS}{\mathcal S}
\providecommand{\pao}{\mathfrak o}
\providecommand{\psvac}{{\phi^{\mathrm{vac}}}}
% macros of this file (prefix \pex)
\providecommand{\pexK}{K^{(z)}}
\providecommand{\pexmc}{\mathfrak c}
\providecommand{\pexMat}{\mathrm M}

\subsection{The norm of the lift of a character, and integral measures}\label{pex:sec}

This subsection proves Proposition~\ref{po:rallis}, in the slightly stronger form needed below
(Proposition~\ref{pex:rallis}: all tame characters of modulus $Q$, not only twists, with the
constant made explicit), and deduces from it statement (IM) of \S\ref{pse:ss-stick} for the
critical pairs with $z\ge2$ (Proposition~\ref{pex:IM}) and, with Theorem~\ref{pcx:X}, for $z=1$
(Proposition~\ref{pex:IM1}). Proposition~\ref{po:exchange} is proved in \S\ref{pte:sec} and \S\ref{ptf:sec}, with the explicit form of
Theorem~\ref{pcx:X} given in Corollary~\ref{pte:explicit}.

\subsubsection{Setting}\label{pex:ss-set}

We use the notation of \S\ref{po:ss-hecke} and \S\ref{pse:ss-stick}. A \emph{tame character of
modulus $Q$} is as in \S\ref{pse:ss-stick}: a character $\vartheta$ of $\auU(1)(E^+)\backslash
\auU(1)(\auA_{E^+})$, trivial at the infinite places, unramified outside $Q$ and tamely ramified
(possibly unramified) at the places of $Q$, where $Q$ is a finite set of places of $E$, prime to
$MN$, with $Q\cap\aubar Q=\emptyset$ and $Q\cap\mathsf M=\emptyset$; the twists of
\S\ref{po:ss-set} are the tame characters ramified at every place of their modulus. The trivial
character has modulus $\emptyset$. For $w\in Q$ let $\psi_w$ be the tame part of $\vartheta$ at
$w$, a character of $\auF_w^\times$ (trivial if $\vartheta$ is unramified at $w$), extended by
$\psi_w(0)=0$. Put $\mathsf S=\mathsf M\cup Q$, $\auNm Q=\prod_{w\in Q}q_w$, $q_w=\#\auF_w$.
Every place of $\mathsf S$ lies above a place of $E^+$ that splits in $E$, and $E/E^+$ is
unramified at all finite places (\S\ref{pc:ss-data}).

Fix a critical pair $p=(\nu,z)$ with $z\ge2$; $\nu$ has weight $-c$, $z\le c(\tau)$, and its tame
part is supported on $\mathsf M$. Put $l_\tau=c(\tau)-z\ge0$ and $\mu=\nu\vartheta$. The objects
of \S\ref{po:ss-hecke}, \emph{the lift of a character}, are defined verbatim for a tame character
$\vartheta$ of modulus $Q$ (the definition in \S\ref{po:ss-hecke} uses of the twist only $Q$ and the
tame parts $\psi_w$): the definite line $E$ in place of $W$, the standard space $V^{(z)}=E^{2z}$, the
block datum
\[
\Phi^{(1)}_{\vartheta}=\bigotimes_{w\in\mathsf M}\plone[x\equiv e_1\bmod w]\otimes
\bigotimes_{w\in Q}\plone[x\text{ integral}]\,\psi_w(x_1\bmod w)\otimes(\text{vacuum elsewhere})
\]
on $\auHom(E,V^{(z)})(\auA_f)=E^{2z}(\auA_f)$, columns written in the frame and read in the
coordinates $T=a_w^{-1}x_w$ of (N5) (Lemma~\ref{pgx:dict}; so ``integral'' means $x_w\in a_w\auO_w^{2z}=\alpha_w\auO_w^{2z}$), $x_1$ the first
coordinate; the level $\pexK_\vartheta\subseteq K_{V^{(z)}}$ of the $k$ whose reduction fixes $e_1$
at the places of $\mathsf M$ and preserves the span of $e_2,\dots,e_{2z}$ at the places of $Q$; the
space $\mathcal A(V^{(z)};p,\vartheta)$ of algebraic forms of the weight with one row of length
$l_\tau$ at $\tau$, level $\pexK_\vartheta$ and the character by which $\pexK_\vartheta$ acts on
$\Phi^{(1)}_\vartheta$; the lift $\Theta^{(1)}_\vartheta$ of \eqref{po:eq-lift}, from functions on
$\auU(1)(E^+)\backslash\auU(1)(\auA_{E^+})/K^1_M$ (with $K^1_M$ the units $\equiv1$ at
$\mathsf M$, and all units at the other finite places); and the normalised lift
\begin{equation}\label{pex:eq-phi}
\varphi_{p,\vartheta}=\kappa_p\,[K^1:K^1_M]\,\Theta^{(1)}_\vartheta(\plone),\qquad\kappa_p=
c_{\mathrm{vac}}^{\,2z}\,(\pi_Ms_0)^{\sum_\tau l_\tau},
\end{equation}
where $K^1$ is the group of all integral points of $\auU(1)(\auA_{E^+,f})$ and $\kappa_p\ne0$ does not
depend on $\vartheta$. For a twist this is the $\varphi_{p,\vartheta}$ of \S\ref{po:ss-hecke}. The
norm $\|\cdot\|$ on $\mathcal A(V^{(z)};p,\vartheta)$ is that of \S\ref{ss:au-spaces}.

\emph{The two conventions.} By \eqref{po:eq-w5} and Remark~\ref{pc:rem-norm}, the representation
used to define $\Theta^{(1)}_\vartheta$ is $\omega'_{a,b}$ with $a=\upsilon^{r_-}=1$ (the line is
positive definite) and $b=\mu\,\upsilon^{z}$; by \eqref{pcx:eq-dict} it is $\omega_{\alpha,\beta}$ of
\S\ref{pc:ss-data} with
\[
\alpha=\nu\vartheta\,\upsilon^{z}\ \text{ on }\auU(1),\qquad\beta=1\ \text{ on }\auU(V^{(z)}).
\]
Let $v^+$ be the highest weight vector of the weight of $\mathcal A(V^{(z)};p,\vartheta)$ and
$P=\pop_{v^+}$ the polynomial \eqref{po:eq-poly} (a polynomial of degree $l_\tau$ at $\tau$ in the
$2z$ holomorphic variables of $\auHom(E,V^{(z)})\otimes_\tau\R=\C^{2z}$). By \eqref{po:eq-lift}, for
a function $y$ on $[\auU(1)]$ and $h\in\auU(V^{(z)})(\auA)$,
\begin{equation}\label{pex:eq-f}
f_{\Theta^{(1)}_\vartheta y}(h)=\bigl\langle v^+,h_\infty^{-1}(\Theta^{(1)}_\vartheta y)(h_f)\bigr\rangle
=\int_{[\auU(1)]}y(g)\,\theta\bigl(\omega_{\alpha,\beta}(g,h)[P,\Psi_\vartheta]\bigr)\,dg
=\bigl(\Theta_{\alpha,\beta}(P,\Psi_\vartheta)y\bigr)(h),
\end{equation}
in the notation of Lemma~\ref{pc:algpet} and \S\ref{pc:ss-data}, where $\Psi_\vartheta$ is the
function on $\pexMat_{2z\times1}(\pcR_{\mathfrak m})$, $\mathfrak m=(\pi_M)\prod_{w\in Q}w$, of which
$\Phi^{(1)}_\vartheta$ is the datum $\Psi^\natural_\vartheta\otimes\psvac^{\mathfrak m}$ of
Definition~\ref{pgx:conv} (both are supported on $a_w\auO_w^{2z}$ at $w\mid\mathfrak m$ and are the same
function of $a_w^{-1}x$, the coordinate $T$ of (N5), by Lemma~\ref{pgx:dict}(c)): the first
equality is the definition of $f$, the second is \eqref{po:eq-lift} with $\xi=v^+$ and the
$\auU(V^{(z)})(\R)$-equivariance of $\xi\mapsto\pop_\xi$.

\subsubsection{Local computations}\label{pex:ss-local}

\begin{lemma}[Sharpness and the finite masses]\label{pex:sharp}
Let $J\subseteq K^1$ be the group of the integral points $\equiv1$ at the places of $\mathsf M$
(no condition at the places of $Q$).
\begin{enumerate}
\item[(a)] $\Psi_\vartheta$ is sharp for $J$ (Definition~\ref{pc:sharp}, for the modulus
$\mathfrak m$ and the character $\alpha$).
\item[(b)] $\|\Psi_\vartheta\|^2/\#\pexMat_{2z\times1}(\pcR_{\mathfrak m})=\prod_{w\in\mathsf M}q_w^{-2z}
\cdot\prod_{w\in Q}(q_w-1)/q_w$, and $[K^1:J]=\prod_{w\in\mathsf M}(q_w-1)$.
\item[(c)] $[K_{V^{(z)}}:\pexK_\vartheta]=c_{\mathsf M}\prod_{w\in Q}\dfrac{q_w^{2z}-1}{q_w-1}$, with
$c_{\mathsf M}=\prod_{w\in\mathsf M}(q_w^{2z}-1)$ independent of $\vartheta$.
\end{enumerate}
\end{lemma}
\begin{proof}
(a) $J$ contains $J_{E}(\mathfrak m)$ and the full components of $K^1$ outside $\mathsf S$. At
$w\in\mathsf M$ the reduction of $J$ is trivial and there is nothing to check for the first condition
of sharpness. At $w\in Q$ the tame part of $\alpha$ is that of $\vartheta$, namely $\psi_w$ ($\nu$
and $\upsilon$ are unramified at $w$), and for $\bar k\in\auF_w^\times$,
$\psi_w(\bar k)^{-1}\Psi_{\vartheta,w}(x\bar k)=\psi_w(\bar k)^{-1}\psi_w(x_1\bar k)=\Psi_{\vartheta,w}
(x)$. For the second condition take $A(x)=$ the first coordinate of $x$ at every place of
$\mathfrak m$: on the support of $\Psi_\vartheta$ it is $1$ at $w\in\mathsf M$ and a unit at $w\in Q$,
hence lies in the reduction $\bar J=\prod_{w\in\mathsf M}\{1\}\times\prod_{w\in Q}\auF_w^\times$.
(b) At $w\in\mathsf M$ the support is one vector out of $q_w^{2z}$; at $w\in Q$ it consists of the
$(q_w-1)q_w^{2z-1}$ vectors with $x_1\ne0$, where $|\Psi|=1$. The index is the order of
$\prod_{w\in\mathsf M}\auF_w^\times$. (c) At $w\in\mathsf M$ the reduction of $\pexK_\vartheta$ is
the stabiliser of the vector $e_1$ in $\auGL_{2z}(\auF_w)$, of index the number $q_w^{2z}-1$ of
non-zero vectors; at $w\in Q$ it is the stabiliser of a hyperplane, of index the number of
hyperplanes $(q_w^{2z}-1)/(q_w-1)$.
\end{proof}

\begin{lemma}[The unramified sum]\label{pex:vac}
Let $\pexmc=\pexmc_{2z,\alpha}$ be the coefficient of the vacuum of \S\ref{pc:ss-data} for $n=1$,
$M'=2z$ and the modulus $\mathfrak m$, on $G^{\mathfrak m}=\auU(1)(\auA^{\mathsf S}_{E^+,f})$ (the finite
points with trivial components at the places below $\mathsf S$). Then $\sum_{\gamma\in G^{\mathfrak
m}/K^{\mathfrak m}}|\pexmc(\gamma)|<\infty$ and
\[
\sum_{\gamma\in G^{\mathfrak m}/K^{\mathfrak m}}\pexmc(\gamma)=\frac{L^{\mathsf S}\bigl(z,(\nu\vartheta)^{-1}
\bigr)}{\zeta^{\mathsf S}_{E^+}(2z)},
\]
where $L^{\mathsf S}$ is as in \S\ref{po:ss-hecke} and $\zeta^{\mathsf S}_{E^+}$ is the Dedekind zeta
function of $E^+$ without the Euler factors at the places below $\mathsf S$.
\end{lemma}
\begin{proof}
$G^{\mathfrak m}/K^{\mathfrak m}$ is the restricted product of the $\auU(1)(E^+_v)/K_v$ over the
places $v$ of $E^+$ not below $\mathsf S$, and $\pexmc$ is the product of local coefficients
$\pexmc_v(\gamma_v)=\alpha_v(\gamma_v)^{-1}\langle\psvac_v,\omega_v(\gamma_v,1)\psvac_v\rangle$, equal to
$1$ on $K_v$ (by (W4); $\alpha$ is unramified at $v$). So it suffices to compute, for each $v$, the
local sum $\Sigma_v=\sum_{\gamma\in\auU(1)(E^+_v)/K_v}\pexmc_v(\gamma)$, and to check
$\sum_v(\sum|\pexmc_v|-1)<\infty$.

\emph{$v$ inert} (it is unramified). $\auU(1)(E^+_v)=\{x\in E_v^\times:x\aubar x=1\}$ consists of
units, so it equals $K_v$ and $\Sigma_v=1$. The Euler factor at the place $w=v$ of $E$ of
$L(s,(\nu\vartheta)^{-1})$ (the $L$-function of $x\mapsto(\nu\vartheta)^{-1}(x/\aubar x)$) is
$(1-q_w^{-s})^{-1}=(1-q_v^{-2s})^{-1}$, because a uniformiser $\varpi$ of $E^+_v$ is one of $E_w$
and $\varpi/\aubar\varpi=1$. So $\Sigma_v=1=L_w(z,(\nu\vartheta)^{-1})/\zeta_{E^+,v}(2z)$.

\emph{$v=w\aubar w$ split.} $\auU(1)(E^+_v)\cong E_w^\times$ by $x\mapsto x_w$ (the component at
$\aubar w$ is then $\aubar{x_w}^{-1}$), $K_v\cong\auO_w^\times$, and the classes are represented by
$\varpi_w^k$, $k\in\Z$. In the model (W5) at $v$, the space is that of Schwartz--Bruhat functions
on $E_w^{2z}$, $\psvac_v=\plone_{M_w}$ with $M_w=a_w\auO_w^{2z}$ (Lemma~\ref{pgx:vac}); in the
coordinate $\tilde x=a_w^{-1}x$ of Lemma~\ref{pgx:dict} this is $\plone[\tilde x\in\auO_w^{2z}]$, and
the integral below is computed in $\tilde x$, with the measure giving $\auO_w^{2z}$ volume $1$ (the
substitution $x=a_w\tilde x$ does not change its value). By \eqref{po:eq-w5} (with $n=2z$, the character
$\mu=\nu\vartheta$ on the source group, and $\alpha=\mu\upsilon^z$, whose $\upsilon$-part is the
splitting) $\pexmc_v(\varpi_w^k)=\chi_w(\varpi_w)^{-k}|\varpi_w^k|_w^{z}\int_{E_w^{2z}}\plone[\auO_w^{2z}](x)
\,\plone[\auO_w^{2z}](x\varpi_w^k)\,dx$, where $\chi_w$ is the (unramified) base change of
$\nu\vartheta$ at $w$; the modulus $|\pexmc_v(\varpi^k)|=q_w^{-|k|z}$ agrees with
Lemma~\ref{pc:series}(a). The integral is $1$ for $k\ge0$ and $q_w^{2zk}$ for $k<0$. With
$a=\chi_w(\varpi_w)^{-1}$, $|a|=1$, and $u=q_w^{-z}$,
\[
\Sigma_v=\sum_{k\ge0}(au)^k+\sum_{k\ge1}(a^{-1}u)^k=\frac{1-u^2}{(1-au)(1-a^{-1}u)} .
\]
Since $\chi_{\aubar w}(\varpi_{\aubar w})=\chi_w(\varpi_w)^{-1}$ (the character $x\mapsto\mu(x/\aubar x)$
is inverted by complex conjugation), the denominator is the inverse of the product of the Euler
factors at $w$ and $\aubar w$ of $L(z,(\nu\vartheta)^{-1})$, and $1-u^2=\zeta_{E^+,v}(2z)^{-1}$. The
sum of $|\pexmc_v|$ is $(1+u)/(1-u)\le1+3q_w^{-2}$ ($z\ge2$), and $\sum_wq_w^{-2}<\infty$.
Multiplying over $v$ gives the assertion; the products converge absolutely because $z\ge2$.
\end{proof}

\begin{lemma}[The archimedean factor]\label{pex:arch}
For $g\in\auU(1)(E^+\otimes\R)$, $\omega_{\alpha,\beta,\infty}(g,1)P=P$. Hence, in
Proposition~\ref{pc:rallis} for the pair $(E,V^{(z)})$, the constant function on $\auU(1)_\infty$ as
the type of the source, and the polynomial $P$, the archimedean integral is
$Z_\infty=\langle P,P\rangle_{\mathrm{Fock}}>0$; it does not depend on $\vartheta$.
\end{lemma}
\begin{proof}
At $\tau$, $\auU(1)(\R)$ is the circle, acting on the Fock space of $\auHom(E,V^{(z)})\otimes_\tau\R=
\C^{2z}$ (all variables holomorphic, both spaces being definite) by the Weil representation, which in
the Fock model multiplies a homogeneous polynomial of degree $d$ in the holomorphic variables by
$\epsilon^{-d}$ times the value at $\epsilon$ of the archimedean splitting character
($\epsilon\in\auU(1)(\R)$; (N1)--(N5) of \S\ref{pl:ss-norm}, \cite{KashiwaraVergne}). The factor
$\alpha_\tau(\epsilon)^{-1}$ in $\omega_{\alpha,\beta}$ cancels the splitting character (the
$\upsilon^z$-part of $\alpha$) and multiplies by $\epsilon^{c(\tau)-z}=\epsilon^{l_\tau}$, the
character $\nu_\tau^{-1}$ ($\nu$ has weight $-c$, $\vartheta_\tau=1$). Since $P$ is homogeneous of
degree $l_\tau$ at $\tau$, the product is trivial. Then the integrand of $Z_\infty$ in
Lemma~\ref{pc:arch}(A3) is the constant $\langle P,P\rangle\cdot1$ on a group of total mass $1$. The
vector $P$ is non-zero because $v^+\ne0$ and $\xi\mapsto\pop_\xi$ is injective on the weight
(\S\ref{po:ss-set}). (This is the condition that the constant function lifts into a space of the
weight of \S\ref{po:ss-hecke}, i.e.\ (A2) of Lemma~\ref{pc:arch} for this pair; (A1) holds because
$[\auU(1)]$ and $[\auU(V^{(z)})]$ are compact.)
\end{proof}

\subsubsection{The norm of the lift}\label{pex:ss-rallis}

\begin{proposition}[Norm of the lift of a character]\label{pex:rallis}
Let $p=(\nu,z)$ be a critical pair with $z\ge2$. There is a real number $A_p>0$, independent of
$\vartheta$, such that for every tame character $\vartheta$ of modulus $Q$
\[
\|\varphi_{p,\vartheta}\|^2=A_p\,(\auNm Q)^{2z-1}\,L^{\mathsf M\cup Q}\bigl(z,(\nu\vartheta)^{-1}\bigr),
\]
and $L^{\mathsf M\cup Q}(z,(\nu\vartheta)^{-1})>0$. Explicitly
\[
A_p=|\kappa_p|^2\,\prod_{w\in\mathsf M}(q_w-1)\cdot\dim W_{\lambda}\cdot\|v^+\|^{-2}\cdot c'\,
\langle P,P\rangle_{\mathrm{Fock}}\,\mathrm{vol}([\auU(V^{(z)})])\,\mathrm{vol}([\auU(1)])\,
\zeta_{E^+}(2z)^{-1},
\]
with $c'\ne0$ the constant of Lemma~\ref{pc:clean}(b) (all factors of $A_p$ are non-zero, and $A_p$ is real and positive, see the proof). In particular, for twists $\vartheta$, $\rho_p(\vartheta)=\|\varphi_{p,\vartheta}\|^2/
\bigl((\auNm Q_\vartheta)^{2z-1}\|\varphi_{p,1}\|^2\bigr)$.
\end{proposition}

\begin{proof}
\emph{Step 1: the inner product formula.} We apply Proposition~\ref{pc:rallis} to $S=E$ (rank $n=1$,
definite), $T=V^{(z)}$ (rank $M'=2z$, definite), the characters $\alpha,\beta$ of
\S\ref{pex:ss-set}, the modulus $\mathfrak m$, the polynomial $P$, the function $\Psi_\vartheta$,
which is sharp for $J$ (Lemma~\ref{pex:sharp}(a)), and $y=y'=\plone$, which is $J$-invariant. We
check that its proof applies. Step~1 there uses (W7), (SW) and Lemma~\ref{pc:clean}: the component of
$[P,\Psi_\vartheta]$ at a place of $\mathsf M$ is supported on injective $2z\times1$ matrices
(columns $\equiv e_1$, hence non-zero), as Lemma~\ref{pc:clean} requires; (SW) is used for
$\pcW=E\oplus E^-$ of rank $2$ and $T$ of rank $2z$. Both conditions of (SW) in Lemma~\ref{pc:SW} hold: (i) the theta
integral converges absolutely because $T$ is definite at every infinite place, hence anisotropic; (ii)
$M'=2z>2=2n$. So (SW) applies directly (for $z=2$ this is the case $M'=4$, $n=1$; Weil's
identity~\cite[Th\'eor\`eme~5]{WeilSiegel} in its range of absolute convergence); the height bounds of (SW) are not needed, since $g,g'$ run over the compact space $[\auU(1)]$. Step~2 uses only that $y,y'$
are bounded on the compact space $[\auU(1)]$. In Step~3 the series is $\sum_\gamma\pexmc(\gamma)\rho
(\gamma)$, which converges absolutely by Lemma~\ref{pex:vac} (Lemma~\ref{pc:series}(b),(c) is stated
for $n\ge2$ and is replaced by Lemma~\ref{pex:vac}); applied to $\plone$ it gives the constant
$\Sigma\cdot\plone$, $\Sigma=\sum_\gamma\pexmc(\gamma)$. The archimedean factor is $Z_\infty$ of
Lemma~\ref{pex:arch}. Hence, with \eqref{pex:eq-f},
\[
\langle f_{\Theta^{(1)}_\vartheta\plone},f_{\Theta^{(1)}_\vartheta\plone}\rangle_{[\auU(V^{(z)})]}
=c'Z_\infty\,\mathrm{vol}([\auU(V^{(z)})])\cdot\frac{\|\Psi_\vartheta\|^2}{\#\pexMat_{2z\times1}
(\pcR_{\mathfrak m})\,[K^1:J]}\cdot\mathrm{vol}([\auU(1)])\,\Sigma .
\]

\emph{Step 2: the algebraic norm.} By Lemma~\ref{pc:algpet},
$\|\Theta^{(1)}_\vartheta\plone\|^2=[K_{V^{(z)}}:\pexK_\vartheta]\cdot\dim W_\lambda\cdot\|v^+\|^{-2}\cdot
\langle f,f\rangle_{[\auU(V^{(z)})]}$, $f=f_{\Theta^{(1)}_\vartheta\plone}$.

\emph{Step 3: collecting the factors.} By Lemma~\ref{pex:vac},
$\Sigma=L^{\mathsf S}(z,(\nu\vartheta)^{-1})\,\zeta_{E^+}(2z)^{-1}\prod_{w\in\mathsf S}(1-q_w^{-2z})^{-1}$
(the places of $\mathsf S$ lie above distinct split places $v$ of $E^+$ with $q_v=q_w$, because
$\mathsf S\cap\aubar{\mathsf S}=\emptyset$). Insert Lemma~\ref{pex:sharp}(b),(c). The factors
depending on a place $w$ of $\mathsf S$ are:
\[
w\in\mathsf M:\ \frac{q_w^{-2z}}{q_w-1}\cdot(q_w^{2z}-1)\cdot(1-q_w^{-2z})^{-1}=\frac1{q_w-1};\qquad
w\in Q:\ \frac{q_w-1}{q_w}\cdot\frac{q_w^{2z}-1}{q_w-1}\cdot(1-q_w^{-2z})^{-1}=q_w^{2z-1}.
\]
So $\|\Theta^{(1)}_\vartheta\plone\|^2=B\,(\auNm Q)^{2z-1}L^{\mathsf S}(z,(\nu\vartheta)^{-1})$ with
\[B=\prod_{w\in\mathsf M}(q_w-1)^{-1}\cdot\dim W_\lambda\,\|v^+\|^{-2}\,c'Z_\infty\,\mathrm{vol}([\auU(V^{(z)})])
\,\mathrm{vol}([\auU(1)])\,\zeta_{E^+}(2z)^{-1},\]
which is independent of $\vartheta$. Multiplying by
$|\kappa_p|^2[K^1:K^1_M]^2=|\kappa_p|^2\prod_{w\in\mathsf M}(q_w-1)^2$ gives the formula with the stated
$A_p$.

\emph{Positivity.} $L^{\mathsf S}(z,(\nu\vartheta)^{-1})$ is a convergent Euler product ($z\ge2$) of
non-zero factors, and it is real and positive: the factors at the places $w\ne\aubar w$ not in
$\mathsf S\cup\aubar{\mathsf S}$ come in pairs $(1-aq^{-z})^{-1}(1-\aubar aq^{-z})^{-1}=|1-aq^{-z}|^{-2}$,
$|a|=1$, and the factors at $w=\aubar w$ are $(1-q_w^{-z})^{-1}>0$ (proof of Lemma~\ref{pex:vac}).
All factors of $A_p$ are non-zero: $c'\ne0$ by Lemma~\ref{pc:clean}(b),
$Z_\infty$, hence $\langle P,P\rangle_{\mathrm{Fock}}$, is non-zero by Lemma~\ref{pex:arch}, the volumes,
$\dim W_\lambda$, $\|v^+\|^{-2}$ and $\zeta_{E^+}(2z)^{-1}$ are positive, and $\kappa_p\ne0$; so $A_p\ne0$.
Since $\|\varphi_{p,\vartheta}\|^2\ge0$ and $(\auNm Q)^{2z-1}L^{\mathsf M\cup Q}(z,(\nu\vartheta)^{-1})>0$,
the formula shows that $A_p$ is real and $\ge0$; hence $A_p>0$, and $\|\varphi_{p,\vartheta}\|^2>0$, so $\varphi_{p,\vartheta}
\ne0$. The last assertion follows by dividing the formula for $\vartheta$ by that for $\vartheta=1$
(modulus $\emptyset$) and recalling $\rho_p(\vartheta)=|L^{\mathsf M\cup Q}(z,(\nu\vartheta)^{-1})|/
|L^{\mathsf M}(z,\nu^{-1})|$.
\end{proof}

\begin{remark}\label{pex:rem-rallis}
(i) Proposition~\ref{po:rallis} is the case of twists; the extension to all
tame characters of modulus $Q$ (in particular to a twist of modulus $Q$ regarded as a tame
character of a larger modulus) is what (IM)(c) needs. (ii) The proof uses the following inputs of
other sections: Proposition~\ref{pc:rallis} and Lemmas~\ref{pc:SW}, \ref{pc:clean}, \ref{pc:algpet};
(SW) is used for an anisotropic target
with $M'=2z>2n$, inside the range of Lemma~\ref{pc:SW} (absolute convergence of Weil's theorem). Lemma~\ref{pc:arch} is not used: its
content for this pair is Lemma~\ref{pex:arch}. (iii) The constant $A_p$ contains the volumes of the
two automorphic quotients and $c'$, which we do not compute; only $A_p>0$ and its independence of
$\vartheta$ are used.
\end{remark}

\subsubsection{Integral measures: statement {\rm(IM)} for $z\ge2$}\label{pex:ss-IM}

Let $\mathsf N_z=\#\auGL_{4fz}(\auF_3)$. By the proof of Lemma~\ref{po:naut}, the order of every
group $\Gamma\cap h\,\pexK h^{-1}$, $h\in\auU(V^{(z)})(\auA_{E^+,f})$, $\pexK\subseteq K_{V^{(z)}}$, divides
$\mathsf N_z$, and for forms $\Xi,\Xi'\in\mathcal A(V^{(z)};p,\vartheta)$ integral at $\aubar\iota$
and at $\iota$, $\mathsf N_z\langle\Xi,\Xi'\rangle\in\auO_\iota$ (same proof, with
$2f\cdot2z$ in place of $2fn$).

\emph{Kernel integrality.} We use the following consequence of Theorem~\ref{pf:std}, proved exactly
as Proposition~\ref{pk:K12} and Corollary~\ref{pf:K1}(a) (the kernel is the $\auU(1)$, $\auU(2z)$
instance of \eqref{ps:eq-kernel}): \emph{there is an integer $n_p\ge1$ such that for every tame
$\vartheta$ of modulus $Q$, every prime $\ell\ne d$ with $\ell\nmid\auNm Q$, every
$\jmath\colon\C\to\C_\ell$, every $x\in\auU(1)(\auA_{E^+,f})$ and $h\in\auU(V^{(z)})(\auA_{E^+,f})$, the
vector $n_p\,h_\ell^{-1}k^0_Q(x,h)$ has integral coordinates at $\jmath$}, where $k^0_Q(x,h)\in
W_\lambda$ is the kernel $k_{\Phi}(x,h;\cdot)$ of \S\ref{po:ss-set} (scaled by $\kappa_p$) for the
character $\nu$ and the \emph{point datum} $\Phi^0_Q=\bigotimes_{w\in\mathsf M}\plone[x\equiv e_1]
\otimes\bigotimes_{w\in Q}\plone[x\text{ integral},\,x_1\equiv1\bmod w]\otimes(\text{vacuum})$.
Indeed, by finite Fourier analysis at the places of $\mathsf M\cup Q$ (the indicator of a coset of a
sublattice of index $q$ of $\auO_w^{2z}$ is $q^{-1}$ times a sum of $q$ Heisenberg translates of the
local vacuum, with roots of unity as coefficients; the index is $q_w^{2z}$ at $w\in\mathsf M$ and
$q_w$ at $w\in Q$, and $\prod_{w\in\mathsf M}q_w^{2z}=M^{2fz}$), the function $M^{2fz}\auNm Q\cdot
c_{\mathrm{vac}}^{2z}\,\Phi^0_Q$ is a standard datum of level $\aubar{\pi_M}$ at every prime $\ell$;
$\auNm Q$ is a unit at $\ell$, and the values of $\nu$ and $\upsilon$ at ideles
that are units above $\ell$ are units. So $n_p$ may be taken to be $M^{2fz}$ times the constant of
Proposition~\ref{pk:K12} for this instance; it depends only on $p$
(not on $\vartheta$ or $Q$), and its value is not used (for flat pairs, $l_\tau=0$, one may take
$n_p=M^{2fz}$ at every $\ell\nmid\auNm Q$: see the proof of Lemma~\ref{psu:I}(b)).
The prime $\ell=N$ is covered by Corollary~\ref{pf:K1}(b) in the same way.

\begin{lemma}[Kernel integrality at $\ell=d$]\label{pex:dint}
Let $\ell=d$, $\iota\colon\C\to\C_d$, $\delta\in\auO_K\smallsetminus\{0\}$ and $f$ a standard datum of
level $\delta$ at $d$. Then $\Theta_f(d\,\aubar\delta s_0\,\cdot)$ is $\partial$-integral; hence for every
$k$ and every homogeneous polynomial $Q$ of degree $k$ with coefficients in $\pao$, $d^kQ(\partial)
\Theta_f(\aubar\delta s_0\,\cdot)(0)\in\pao$.
\end{lemma}
\begin{proof}
\emph{The curve at $d$.} Normalise the valuation by $v(d)=1$. In the proof of
Proposition~\ref{pk:curve}(C2), $x$, $y$ are $\partial$-integral at every prime, $g^{(n)}(0)=x^{(n-2)}(0)$
for $n\ge3$, and $-2d\,g''(0)=\sum_{u}x(u-t)-\tfrac{d+1}{12}a_1^2$, the sum over the non-zero points of
order $d$. A point of order $d$ either has integral coordinates or lies in the group $\mathcal E_1$ of
points reducing to the origin of the integral model, which is the group of the formal group over the
integers $\mathfrak o$ of a finite extension of $\Q_d$ \cite[Chapter~VII, Proposition~2.2]{SilvermanAEC} (that proposition
assumes a minimal Weierstrass equation; the model $y^2+a_1xy=x^3+x$ of \eqref{pa:eq-curve} is
minimal at $d$, since $a_1$ is integral and its discriminant $a_1^4-64$ is a unit at the odd prime $d$,
Lemma~\ref{pa:model}); a torsion point of
exact order $d$ there has parameter $\mathsf z$ with $v(\mathsf z)\le v(d)/(d-1)$
\cite[Chapter~IV, Theorem~6.1]{SilvermanAEC}, and $x=\mathsf z^{-2}(1+\mathsf z\,\mathfrak o[[\mathsf z]])$,
so $v(x)\ge-2/(d-1)$. Since $2\ne d$ and $12\mid d+1$ (Proposition~\ref{pk:curve}, proof of (C2)),
$v(g''(0))\ge-1-2/(d-1)>-2$ ($d\ge5$), so $d^2g''(0)\in\pao$. Hence $g(dw)=\sum_{n\ge2}g^{(n)}(0)d^nw^n/n!$
is $\partial$-integral with zero constant term, and $\theta(dw)=\exp(-g(dw))$ and $y(dw)$ are
$\partial$-integral (the rules recalled before Proposition~\ref{pk:curve}).

\emph{The datum.} Theorem~\ref{pf:std} holds at every prime: $\Theta_f(\aubar\delta s_0z)=(\prod_j\theta
(\mu_j(z))^{m_j})^{\auNm\delta}\mathsf H(z)$ with $\mathsf H$ an integral power series in the
$y(\mu_j(z))$ and $\mu_j$ linear forms with integral coefficients (Lemma~\ref{pf:scale},
Corollary~\ref{pf:CS}, both valid at every prime). Replace $z$ by $dz$: $\mu_j(dz)=d\mu_j(z)$. A
one-variable $\partial$-integral series composed with a linear form with integral coefficients is
$\partial$-integral in several variables (expand $(\sum_ia_iz_i)^n/n!=\sum\prod_ia_i^{k_i}z_i^{k_i}/k_i!$),
so $\theta(d\mu_j(z))$ and $y(d\mu_j(z))$ are $\partial$-integral, the latter with zero constant term;
products and $\mathsf H(\dots)$ of them are $\partial$-integral. Finally a jet of degree $k$ of
$\Theta_f(\aubar\delta s_0\,\cdot)$ is $d^{-k}$ times the same jet of $\Theta_f(d\aubar\delta s_0\,\cdot)$.
\end{proof}

Consequently kernel integrality holds at $\ell=d$ (if $d\nmid N\auNm Q$) with $n_p$ replaced by
$n_pd^{\sum_\tau l_\tau}$: the reduction of Proposition~\ref{pk:K12} is independent of the prime, and the
jets are of degree $\sum_\tau l_\tau$. From now on $n_p$ includes it, so kernel integrality
holds at all $\ell\nmid\auNm Q$.

\begin{lemma}[The lift as a sum over points]\label{pex:points}
Let $Q$ be a finite set of places as above and $\mathcal P_Q=\bigl(\auU(1)(E^+)\backslash
\auU(1)(\auA_{E^+,f})/K^1_M\bigr)\times\prod_{w\in Q}\auF_w^\times$. Choose representatives $x_c$ of
the classes and, for $a=(a_w)\in\prod_{w\in Q}\auF_w^\times$, an integral $u_a\in\auU(1)(\auA_{E^+,f})$
equal to $1$ outside $Q\cup\aubar Q$ whose component at $w\in Q$ reduces to $a_w^{-1}$. Put
$x_{c,a}=x_cu_a$. Then for every tame $\vartheta$ of modulus $Q$ and every $h$,
\[
\varphi_{p,\vartheta}(h)=\sum_{(c,a)\in\mathcal P_Q}\vartheta(x_{c,a})^{-1}\,k^0_Q(x_{c,a},h).
\]
\end{lemma}
\begin{proof}
In \eqref{po:eq-w5} the character $\vartheta$ enters $\omega_\mu$, $\mu=\nu\vartheta$, only through
the scalar $\vartheta(\det g)^{-1}$ on $\auU(1)$ (it is trivial at the infinite places); so
$k_{\Phi}(g,h;\xi)$ for $\nu\vartheta$ is $\vartheta(g)^{-1}$ times the kernel for $\nu$. At $w\in Q$,
$\Phi^{(1)}_{\vartheta,w}=\sum_{a\in\auF_w^\times}\psi_w(a)\,\plone[x\text{ integral},x_1\equiv a]$, and
$\plone[x\text{ integral},x_1\equiv a_w]=\omega_\nu(u_a,1)\plone[x\text{ integral},x_1\equiv1]$ by
\eqref{po:eq-w5}, $\nu$ being unramified at $w$ (the unramified scalar is $1$ on units). Since the
component of $\vartheta$ at $w$ is $\psi_w$ on units, $\psi_w(a_w)=\vartheta_w(u_a)^{-1}$. Hence the
integrand of $\Theta^{(1)}_\vartheta(\plone)$ is $F(g)=\sum_a\vartheta(gu_a)^{-1}k^0_Q(gu_a,h)/\kappa_p$.
This function of $g\in[\auU(1)]$ is right invariant under $K^1_M$: for $k\in K^1_M$ the elements $ku_a$
run over representatives of the $u_a$ modulo the units $\equiv1$ at $Q$, under which the summand is
invariant ($\Phi^0_Q$ is invariant under them, by \eqref{po:eq-w5}). So $\int_{[\auU(1)]}F=
\mathrm{vol}(K^1_M)\sum_cF(x_c)$, the stabilisers $\auU(1)(E^+)\cap x_cK^1_Mx_c^{-1}$ being trivial
(the group $K^1_M$ is abelian, so such an element is a root of unity $\zeta\in E$ with $\zeta\equiv1$
modulo every $w\in\mathsf M$; if $\zeta$ has exact order $m>1$, then $\zeta-1$ is a unit when $m$ is not a
prime power and divides $r$ when $m$ is a power of a prime $r$, so $\zeta-1$ is a unit at every place
not above $r$; as $\mathsf M$ contains places above each of the two distinct primes $p_0,p_1$, this is
impossible, so $\zeta=1$), and
the volume of $[\auU(1)]$-measure is normalised by $\mathrm{vol}(K^1)=1$ (\S\ref{pc:ss-data}), so
$\mathrm{vol}(K^1_M)=[K^1:K^1_M]^{-1}$. Multiply by $\kappa_p[K^1:K^1_M]$.
\end{proof}

\begin{proposition}[Statement (IM) for $z\ge2$]\label{pex:IM}
Let $p=(\nu,z)$ be a ramified critical pair with $z\ge2$. Put $I=\{z\}$, $m=0$ and
\[
F(\vartheta)=\mathsf N_z\,n_p^2\,\frac{\|\varphi_{p,\vartheta}\|^2}{(\auNm Q)^{2z-1}}\qquad
(\vartheta\text{ tame of modulus }Q).
\]
Then: (a) $F(1)\ne0$; (b) for every $Q$ there are finitely many $x_a\in\auU(1)(\auA_{E^+,f})$ and
$c_a\in\C$, with $c_a$ integral at every prime $\ell\nmid\auNm Q$ under every embedding $\C\to\C_\ell$,
such that $F(\vartheta)=\sum_ac_a\vartheta(x_a)$ for every tame $\vartheta$ of modulus $Q$; (c) if $v\notin Q$ and $\vartheta'$ is
$\vartheta$ regarded as of modulus $Q\cup\{v\}$, then $F(\vartheta')=F(\vartheta)\prod_\pm(1-\alpha_v^{\pm1}
q_v^{-z})$, $\alpha_v$ the value at a uniformiser of the base change of $\nu\vartheta$ at $v$; (d) for
every twist $\vartheta$ of modulus $Q$, $F(\vartheta)=F(1)\,\rho_p(\vartheta)$.
\end{proposition}
\begin{proof}
By Proposition~\ref{pex:rallis}, $F(\vartheta)=\mathsf N_zn_p^2A_p\,L^{\mathsf M\cup Q}(z,(\nu\vartheta)^{-1})$
with $\mathsf N_zn_p^2A_p>0$ independent of $\vartheta$. (a) $L^{\mathsf M}(z,\nu^{-1})>0$. (d) is the
quotient of the formulas for $\vartheta$ and $1$, the $L$-values being positive. (c) The two
$L$-functions differ by the Euler factors at $v$ and $\aubar v$ (we may assume $v$, $\aubar v\notin
\mathsf M\cup Q\cup\aubar Q$, as for an auxiliary modulus; then $\vartheta$ is unramified at $v$ and the
place below $v$ splits): $L^{\mathsf S\cup\{v\}}=L^{\mathsf S}\cdot(1-aq_v^{-z})(1-a^{-1}q_v^{-z})$ with
$a$ the value at $\varpi_v$ of the base change of $(\nu\vartheta)^{-1}$, using
$\chi_{\aubar v}(\varpi_{\aubar v})=\chi_v(\varpi_v)^{-1}$ as in Lemma~\ref{pex:vac}; the product is
symmetric in $a\leftrightarrow a^{-1}=\alpha_v$.

(b) Fix $Q$ and the points $x_{c,a}$ of Lemma~\ref{pex:points}, which depend only on $Q$. The level
$\pexK_\vartheta$ and the class representatives $h$ of $\Gamma\backslash\auU(V^{(z)})(\auA_f)/\pexK_\vartheta$
depend only on $Q$; let $e_h=\#(\Gamma\cap h\pexK_\vartheta h^{-1})$. By Lemma~\ref{pex:points} and the
definition of the norm (\S\ref{ss:au-spaces}),
\[
\|\varphi_{p,\vartheta}\|^2=\sum_h\frac1{e_h}\bigl\langle\varphi_{p,\vartheta}(h),\varphi_{p,\vartheta}(h)
\bigr\rangle=\sum_{\mathsf p,\mathsf p'\in\mathcal P_Q}\vartheta\bigl(x_{\mathsf p}x_{\mathsf p'}^{-1}\bigr)
\,G_{\mathsf p\mathsf p'},\qquad G_{\mathsf p\mathsf p'}=\sum_h\frac1{e_h}\bigl\langle k^0_Q(x_{\mathsf
p},h),k^0_Q(x_{\mathsf p'},h)\bigr\rangle,
\]
using $\aubar{\vartheta(x)^{-1}}=\vartheta(x)$ and the sesquilinearity of the form (conjugate-linear in
the first variable, the convention of \S\ref{pc:ss-data}); $G_{\mathsf p\mathsf p'}$ depends only on $Q$. Put $x_{(\mathsf p,\mathsf
p')}=x_{\mathsf p}x_{\mathsf p'}^{-1}$ and $c_{(\mathsf p,\mathsf p')}=\mathsf N_zn_p^2G_{\mathsf p\mathsf
p'}(\auNm Q)^{1-2z}$. Then $F(\vartheta)=\sum c_{(\mathsf p,\mathsf p')}\vartheta(x_{(\mathsf p,\mathsf
p')})$ for every tame $\vartheta$ of modulus $Q$. Let $\ell\nmid\auNm Q$ and $\jmath\colon\C\to\C_\ell$.
The form on $W_\lambda$ is invariant under $h_\ell$ and the Gram matrix of the basis of $W_\lambda(\Z)$ is integral (\S\ref{ss:au-spaces}), so
$\langle k^0_Q(x,h),k^0_Q(x',h)\rangle=\langle h_\ell^{-1}k^0_Q(x,h),h_\ell^{-1}k^0_Q(x',h)\rangle$ is,
after multiplication by $n_p^2$, a sum of products of a coordinate integral at $\jmath\circ c$ (complex
conjugation; the conjugate of a coordinate integral at $\jmath\circ c$ is integral at $\jmath$) and one
integral at $\jmath$; so $n_p^2\langle\cdot,\cdot\rangle$ is integral at $\jmath$, by kernel
integrality applied at $\jmath$ and at $\jmath\circ c$. As $e_h\mid\mathsf N_z$, $\mathsf N_zn_p^2G_{\mathsf
p\mathsf p'}$ is integral at $\jmath$, and $\auNm Q$ is a unit at $\ell$.
\end{proof}

\begin{remark}[On the form of (IM)]\label{pex:rem-IM}
(i) In (d) the ratio is $\rho_p(\vartheta)=|L^{\mathsf M\cup Q}(z,(\nu\vartheta)^{-1})|/|L^{\mathsf M}(z,
\nu^{-1})|$, with the Euler factors at $\mathsf M$ removed in both and at $Q$ in the numerator; this is
the convention of $r(\vartheta)$ in \S\ref{pse:ss-stick} (``conventions of
Theorem~\ref{pg:main}(b)'') if those conventions are read as in \S\ref{po:ss-hecke}; at the places of $Q$
the twisted character is ramified, so for a twist nothing is removed there. (ii) In (c) the Euler
factor is stated for $v$ with $v,\aubar v\notin\mathsf M\cup Q\cup\aubar Q$, which is the case used in
\eqref{pse:eq-norm} ($Q\cup\{v\}$ auxiliary). (iii) In (b) the points $x_a$ are finite ideles of
$\auU(1)$ and $\vartheta(x_a)$ is the value of $\vartheta$ there; the use in \S\ref{pse:ss-stick} of
$[x_a^{-\delta}]$ in the class group is unaffected. (iv) $\ell=d$ is covered by Lemma~\ref{pex:dint}.
\end{remark}

\subsubsection{Statement {\rm(IM)} for $z=1$}\label{pex:ss-IM1}

Let $p=(\nu,1)$ be a ramified critical pair; $\nu$ has weight $-c$, $c(\tau)\ge1$. Let $I_2=E^2$ be the
standard positive definite plane, $\vartheta$ a tame character of modulus $Q$, $\mathfrak m_\vartheta=
(\pi_M)\prod_{w\in Q}w$, $\chi_\vartheta$ the character of $\pcR_{\mathfrak m_\vartheta}^\times$ given by the
tame parts of $\nu\vartheta$ (trivial at the places of $Q$ where $\vartheta$ is unramified), extended by
$0$, and $\lambda_\vartheta=\nu\vartheta\upsilon$ (weight $-(c-1)$). Put
$P_c=\prod_\tau\det(y_\tau)^{c(\tau)-1}$, a polynomial on $\auHom(I_2,I_2)\otimes\R=\prod_\tau M_2(\C)$,
and define the \emph{period}
\[
\Lambda(\vartheta)=\sum_{h\in\Gamma\backslash\auU(I_2)(\auA_f)/K_{I_2}}\frac1{e_h}\,\theta^{I_2,I_2}_{
\lambda_\vartheta,\lambda_\vartheta}\bigl(1,h;P_c,\chi_\vartheta\circ\det\bigr),\qquad e_h=\#(\Gamma\cap
hK_{I_2}h^{-1}),
\]
in the notation of \S\ref{pc:ss-data} (the datum $[P_c,\chi_\vartheta\circ\det]$ modulo
$\mathfrak m_\vartheta$). The summand is right $K_{I_2}$-invariant in $h$ by Lemma~\ref{pc:level}: for
$k\in K_{I_2}$ the factor $\lambda_\vartheta(\det k)$ is $\chi_\vartheta(\det\bar k)$ and cancels against
$\chi_\vartheta(\det(\bar k^{-1}x))$. (One may equally take the modulus $(\pi_M)^2\prod w$; as
$\chi_\vartheta$ is tame this changes nothing.)

\begin{lemma}[The isotypic vectors]\label{pex:iso}
Let $k=(k_\tau)\in\Z_{\ge0}^{\auPl}$. A polynomial $R$ on $\prod_\tau M_2(\C)$ with
$R((h_\tau^{-1}y_\tau g_\tau)_\tau)=\prod_\tau\det(g_\tau)^{k_\tau}\det(h_\tau)^{-k_\tau}R(y)$ for all
$g,h\in\prod_\tau\auU(2)$ is a multiple of $\prod_\tau\det(y_\tau)^{k_\tau}$. A polynomial with this
property for some $k\in\Z^{\auPl}$ with some $k_\tau<0$ is $0$.
\end{lemma}
\begin{proof}
Both sides are polynomial in the entries of $g_\tau,h_\tau$ and their inverses, and $\auU(2)$ is
Zariski dense in $\auGL_2(\C)$ (its Lie algebra spans $\mathfrak{gl}_2(\C)$ over $\C$), so the identity
holds for $g,h\in\prod\auGL_2(\C)$. Take $h=1$, $y=1$ and $g=y'$ invertible: $R(y')=\prod_\tau\det
(y'_\tau)^{k_\tau}R(1)$ on a dense set. If some $k_\tau<0$ the right side has a pole unless $R(1)=0$.
\end{proof}

\begin{proposition}[The equal-rank period; (E1)]\label{pex:E1}
There is $C_1>0$ such that for $\vartheta=1$ and for every twist $\vartheta$ of modulus $Q$,
\[
C_1\,|\Lambda(\vartheta)|^2=L_*\bigl(\mathfrak m_\vartheta,(\nu\vartheta)^{-1},1\bigr)=\auNm(\mathsf M\cup Q)
\,\bigl|L^{\mathsf M\cup Q}\bigl(1,(\nu\vartheta)^{-1}\bigr)\bigr|^2\ne0 .
\]
\end{proposition}
\begin{proof}
Apply Theorem~\ref{pcx:X} with $W=V=I_2$, $V'=0$: $N=a=2$, $b=0$, $e=1$, $q=q_V=q_{V'}=0$, $2q+2\le N$;
$\mu=c$, and $N\le2\mu_\tau$ since $c(\tau)\ge1$; $v$ is the (unique up to a scalar) vector $1$ of the
Fock space of the zero space. This gives $C>0$ and a Fock vector $x$ of $\auHom(I_2,I_2)\otimes\R$,
independent of the modulus and the character. The character $\nu'=(\nu\vartheta)^{-1}$ is a tame
character for $\mathfrak m_\vartheta$ of weight $c$ (its tame part is non-trivial at the places of
$\mathsf M$ because $p$ is ramified, and at those of $Q$ because $\vartheta$ is a twist and $\nu$ is
unramified there; for $\vartheta=1$, $\mathfrak m=(\pi_M)$). Then $\gamma=\nu'^{-1}\upsilon=\lambda_\vartheta$;
take $\alpha=\lambda_\vartheta$, $\alpha'=1$. Matched functions for $b=0$: $\Psi$ on $\pexMat_{2\times0}$
is the constant $1$, and $\Phi$ is matched iff $\Phi(k^{-1})=\chi(\det k)$ for $k\in\auGL_2(\pcR)$ and
$\Phi=0$ off $\auGL_2$, where $\chi$ is the character of Definition~\ref{pcx:def-matched} attached to
$\nu'$, i.e.\ $\chi=\chi_\vartheta^{-1}$; so $\Phi=\chi_\vartheta\circ\det$. The right side of
\eqref{pcx:eq-X} is $c$ times the kernel of $\auHom(0,I_2)=0$ with character $\alpha'=1$, i.e.\ $c$. So
\begin{equation}\label{pex:eq-E1}
\int_{[\auU(I_2)]}\theta^{I_2,I_2}_{\lambda_\vartheta,\lambda_\vartheta}(g,h;x,\chi_\vartheta\circ\det)\,dh=c,
\qquad|c|^2=C\,L_*(\mathfrak m_\vartheta,\nu',1),
\end{equation}
for all $g\in\auU(I_2)(\auA)$, and $c\ne0$ (Corollary~\ref{pcx:cne0}).

\emph{Replacing $x$ by $P_c$.} The compact group $K_\infty=\auU(I_2)_\infty\times\auU(I_2)_\infty$ acts on
the Fock space by $\omega_{\lambda,\lambda,\infty}$, and $\theta(g k_1,hk_2;x)=\theta(g,h;\omega_\infty
(k_1,k_2)x)$. Integrate \eqref{pex:eq-E1}, with $g$ replaced by $gk_1$, against $dk_1$ over
$\auU(I_2)_\infty$ (total mass $1$), and use the invariance of $dh$ under $h\mapsto hk_2$: \eqref{pex:eq-E1}
holds with $x$ replaced by its projection $x'$ onto the $K_\infty$-invariant vectors (the integral over
$k_2$ is built into $\int dh$, the one over $k_1$ fixes the constant $c$). The $K_\infty$-invariant
vectors of the Fock space are polynomials (the $K_\infty$-types of the Fock space of a pair of definite
spaces are finite-dimensional spaces of polynomials). In the Fock model $\omega_\infty(k_1,k_2)R(y)$ is
$R(k_2^{-1}yk_1)$ times the splitting factor and $\lambda_\infty(\det k_1)^{-1}\lambda_\infty(\det k_2)$
(\S\ref{pc:ss-data}); since $\lambda=\nu\vartheta\upsilon$ has weight $-(c-1)$ and the splitting by
$\upsilon$ is cancelled as in Lemma~\ref{pex:arch}, invariance means $R(k_2^{-1}yk_1)=\prod_\tau\det
(k_{1,\tau})^{c_\tau-1}\det(k_{2,\tau})^{-(c_\tau-1)}R(y)$ (this is the condition that characterises $P_c$). By Lemma~\ref{pex:iso}, $x'=\kappa P_c$, and $\kappa\ne0$
because $c\ne0$. $\kappa$ depends only on $x$ and the archimedean characters, not on $\vartheta$.

\emph{From the integral to the class sum.} By Lemma~\ref{pc:level} and the invariance under
$\auU(I_2)_\infty$ just used, the integrand of \eqref{pex:eq-E1} at $g=1$ with $x'$ is a function of the
class of $h$ in $\Gamma\backslash\auU(I_2)(\auA_f)/K_{I_2}$; with $\mathrm{vol}(K_{I_2})=1$ and
$\mathrm{vol}(\auU(I_2)_\infty)=1$, $\int_{[\auU(I_2)]}=\sum_h e_h^{-1}(\cdot)(h)$. So
$\kappa\Lambda(\vartheta)=c$, and $C_1=|\kappa|^2/C$.
\end{proof}

\begin{proposition}[Statement (IM) for $z=1$]\label{pex:IM1}
Let $p=(\nu,1)$ be a ramified critical pair, $I=\{0,1\}$, $m=1$, and $F(\vartheta)=|\mathsf N_1n_1
\Lambda(\vartheta)|^2$, with $\mathsf N_1=\#\auGL_{4f}(\auF_3)$ and $n_1\ge1$ an integer, depending only on $p$, making the kernel values below integral at all
$\ell\nmid\auNm Q$ (it exists by Proposition~\ref{pk:K12}, Corollary~\ref{pf:K1} and Lemma~\ref{pex:dint}). Then {\rm(a)}--{\rm(d)}
of {\rm(IM)} hold: $F(1)\ne0$; $F$ is an integral measure at every modulus; for $v,\aubar v\notin\mathsf M
\cup Q\cup\aubar Q$, $F(\vartheta')=F(\vartheta)\prod_{e\in\{0,1\}}\prod_\pm(1-\alpha_v^{\pm1}q_v^{-e})$; and
$F(\vartheta)=F(1)\,\auNm Q\,\rho_p(\vartheta)^2$ for every twist.
\end{proposition}
\begin{proof}
(a) and (d): Proposition~\ref{pex:E1}, since $L_*(\mathfrak m_\vartheta,\nu',1)/L_*((\pi_M),\nu^{-1},1)=
\auNm Q\cdot\rho_p(\vartheta)^2$.

(b) Fix $Q$; the classes $h$, the numbers $e_h$ (which divide $\mathsf N_1$, as in Lemma~\ref{po:naut})
and the ring $\pcR=\pcR_{\mathfrak m_\vartheta}$ depend only on $Q$. Write $\chi_\vartheta\circ\det=\sum_{y\in
\auGL_2(\pcR)}\chi_\vartheta(\det y)\delta_y$ and $\delta_y(x)=\delta_1(\tilde y^{-1}x)$ for an integral lift
$\tilde y\in K_{I_2}$ of $y$; by Lemma~\ref{pc:level}, $\theta_{\lambda,\lambda}(1,h;P_c,\delta_y)=
\lambda(\det\tilde y)^{-1}\theta_{\lambda,\lambda}(1,h\tilde y;P_c,\delta_1)$, and $\lambda(\det\tilde y)=
\chi_\vartheta(\det y)$ on the tame part. Here $\lambda=\lambda_\vartheta$ enters the kernel of $\delta_1$
only through the scalar $\lambda_\vartheta(\det h')$ on finite points $h'$ (the datum $\delta_1$ is
independent of $\vartheta$), and $\lambda_\vartheta=\vartheta\cdot\lambda_1$. Hence
$\Lambda(\vartheta)=\sum_{h,y}e_h^{-1}\vartheta(\det(h\tilde y))\,\theta^0(h\tilde y)$ with
$\theta^0=\theta_{\lambda_1,\lambda_1}(1,\cdot;P_c,\delta_1)$ independent of $\vartheta$ of modulus $Q$. So
$\mathsf N_1n_1\Lambda(\vartheta)=\sum_ac_a\vartheta(x_a)$ with $x_a=\det(h\tilde y)\in\auU(1)(\auA_f)$ and
$c_a=\mathsf N_1e_h^{-1}n_1\theta^0(h\tilde y)$, integral at every $\ell\nmid\auNm Q$ by kernel integrality
(the scalar kernel $\theta^0$ is a jet of degree $2\sum_\tau(c(\tau)-1)$ of $\Theta_f$ for a standard
datum, Proposition~\ref{pk:K12}, Corollary~\ref{pf:K1}, Lemma~\ref{pex:dint}). Then $|\sum_ac_a\vartheta
(x_a)|^2=\sum_{a,b}c_a\aubar{c_b}\vartheta(x_ax_b^{-1})$, and $\aubar{c_b}$ is integral at $\jmath$ when
$c_b$ is integral at $\jmath\circ c$.

(c) Let $\vartheta'$ be $\vartheta$ with modulus $Q\cup\{v\}$. At $v$ the datum of $\vartheta$ is the
vacuum $\plone$ and that of $\vartheta'$ is $\plone[x\text{ integral},\det x\not\equiv0]$ ($\chi_{\vartheta'}$
is trivial at $v$). By Lemma~\ref{pl:vac} (with $n=n'=2$, in which $T'_i\plone$ is computed from
\eqref{pl:eq-w5} for $h$ in $K'd_iK'$, and $\alpha$ the scalar of $d_1$ on the vacuum),
$\sum_{i}(-X)^iq_v^{\binom i2}T'_i\plone=\prod_{l<2-\mathrm{rk}(x\bmod v)}(1-q_v^l\alpha X)\plone$; at
$X=\alpha^{-1}$ this is the datum of $\vartheta'$ at $v$, by the same argument as
\eqref{pl:eq-qbin}'s use there ($\prod_{l<r'}(1-q^l)$ is $1$ for $r'=0$ and $0$ otherwise). The operators
$T'_i$ are sums of translates $\omega(1,h')$, $h'\in\auU(I_2)(E^+_v)$, and the class sum is the integral over
$[\auU(I_2)]$ of a right $K_{I_2}$-invariant function, hence invariant under right translation; so
$\Lambda(\vartheta')=\sum_i(-\alpha^{-1})^iq_v^{\binom i2}\#\{K'd_iK'/K'\}\Lambda(\vartheta)=(1-\alpha^{-1})
(1-q_v\alpha^{-1})\Lambda(\vartheta)$ by \eqref{pl:eq-qbin} for $V=\auF_v^2$. By \eqref{pl:eq-w5},
$\alpha^{-1}=q_v^{-1}a$ with $a$ the value of the base change of $(\nu\vartheta)^{-1}$ at $\varpi_v$ (the
splitting part cancels), $|a|=1$. Hence $|\Lambda(\vartheta')|^2=|1-aq_v^{-1}|^2|1-a|^2\,|\Lambda(\vartheta)|^2=\prod_{e\in\{0,1\}}
\prod_\pm(1-a^{\pm1}q_v^{-e})\,|\Lambda(\vartheta)|^2$, so $F(\vartheta')=F(\vartheta)\prod_{e\in\{0,1\}}\prod_\pm(1-a^{\pm1}q_v^{-e})$, symmetric in $a\leftrightarrow a^{-1}=\alpha_v$.
\end{proof}



%% ===== proofs/C-threeex.tex =====
% proofs/C-threeex.tex
% Proposition po:exchange of C-onelift.tex (three exchanges): the exchange identity with explicit
% determinant data (a corollary of Theorem pcx:X), one exchange summed over classes, the carrier,
% the third exchange, its left side as the theta series of the first factors, and the proof of
% po:exchange. To be \input after C-exchange.tex, C-onelift.tex and C-exchanges3.tex.
% Labels of this file start with pte:. Only amsmath/amssymb/amsthm.
\providecommand{\auO}{\mathcal O}
\providecommand{\auU}{\mathrm U}
\providecommand{\auGL}{\mathrm{GL}}
\providecommand{\auHom}{\operatorname{Hom}}
\providecommand{\auNm}{\operatorname{Nm}}
\providecommand{\auA}{\mathbb A}
\providecommand{\auF}{\mathbb F}
\providecommand{\auPl}{\Phi}
\providecommand{\aubar}[1]{\overline{#1}}
\providecommand{\plone}{\mathbf 1}
\providecommand{\pop}{\mathsf p}
\providecommand{\pcR}{R}
\providecommand{\pcm}{\mathfrak m}
\providecommand{\pcMat}{\mathrm M}
\providecommand{\psvac}{{\phi^{\mathrm{vac}}}}
\providecommand{\psZ}{Z}
\providecommand{\psrho}{\varrho}
\providecommand{\psHS}{\mathcal H}
\providecommand{\poA}{\mathcal A}
\providecommand{\pog}{\mathsf g}
% macros of this file (prefix \pte)
\providecommand{\pteSD}{\mathrm{SD}}
\providecommand{\pteMP}{\mathrm{MP}}
\providecommand{\pteCS}{\mathop{\textstyle\sum^{\mathrm{cl}}}}
\providecommand{\ptetr}{\operatorname{tr}}

\subsection{Three exchanges}\label{pte:sec}

This subsection proves Proposition~\ref{po:exchange}. The proof has four parts: the exchange
identity with explicit determinant data at the infinite places (Corollary~\ref{pte:explicit}, a
consequence of Theorem~\ref{pcx:X}); the exchange summed over the classes of a definite group
(Lemma~\ref{pte:one}); the first two exchanges, which write the carrier as a class sum of twisted
kernels (Proposition~\ref{pte:carrier}); and the third exchange together with the identification of
its left side with the theta series $\Xi_{\vartheta,i}$ of \S\ref{po:ss-glue}
(Propositions~\ref{pte:third} and~\ref{pte:leftside}). In the course of the last step the first
factors $\varphi''_{\vartheta,i}$ of \S\ref{po:ss-glue} are defined explicitly (Definition~\ref{pte:def-ff},
Definition~\ref{ptf:def-ff}).

\subsubsection{Conventions}\label{pte:ss-conv}

We use the convention of \S\ref{pcx:ss-statement}: for spaces $S,T$ and characters $\alpha,\beta$
of $\auU(1)$, $\theta^{S,T}_{\alpha,\beta}(g,h;f)=\theta(\omega_{\alpha,\beta}(g,h)f)$ with
$\omega_{\alpha,\beta}(g,h)=\alpha(\det g)^{-1}\beta(\det h)\omega_{\auHom(S,T)}(g,h)$, $g\in
\auU(S)(\auA)$, $h\in\auU(T)(\auA)$; data are written $[P,\Phi]$ with $P$ a polynomial (its Fock
vector in the model of (N1)--(N5) of \S\ref{pl:ss-norm}, holomorphic in the variables of positive
product sign and antiholomorphic in the others, as in \eqref{po:eq-poly}) and $\Phi$ a function on
matrices modulo a modulus $\pcm$ (\S\ref{pc:ss-data}; matrices in the frames). By
\eqref{pcx:eq-twist},
\begin{equation}\label{pte:eq-twist}
\theta^{S,T}_{\alpha\lambda,\beta\lambda'}(g,h;f)=\lambda(\det g)^{-1}\lambda'(\det h)\,\theta^{S,T}_{\alpha,\beta}(g,h;f).
\end{equation}
For a positive definite standard space $I$ (frame orthonormal at the infinite places, lattice the
standard unimodular one) and a function $F$ on $\auU(I)(\auA_f)$, left invariant under
$\Gamma_I=\auU(I)(E^+)$ and right invariant under an open compact $J$, the \emph{class sum} is
\[
\pteCS_J F=\sum_{h\in\Gamma_I\backslash\auU(I)(\auA_f)/J}\frac{F(h)}{\#(\Gamma_I\cap hJh^{-1})}\cdot
\frac{1}{[K_I:J]}\qquad(J\subseteq K_I),
\]
$K_I$ the stabiliser of the standard lattice; with the Haar measure giving $K_I$ and
$\auU(I)(E^+\otimes\R)$ volume $1$ it is $\int_{\Gamma_I\backslash\auU(I)(\auA_f)}F$, so it does not
depend on $J$ (for $J=K_I$ the last factor is $1$). If $\Phi$ is a kernel datum, $h\mapsto
\theta^{S,I}_{\alpha,\beta}(g,(1,h);f)$ is written $\theta^{S,I}_{\alpha,\beta}(g,h;f)$
($h\in\auU(I)(\auA_f)$, archimedean component $1$).

\emph{Index conventions for determinant data.} Let $W$ be a space of rank $n$ whose basis is split by
sign as $m$ positive and $l$ negative vectors, and $V'$ of rank $b$ split as $p$ positive and $l$
negative vectors (the same number $l$), and let $I_a$ be positive definite of rank $a$, with
$n=b+a$; so $m=p+a$. Fix a bijection of the $m$ positive indices of $W$ with $\{1,\dots,p\}\sqcup
\{1,\dots,a\}$. At each $\tau\in\auPl$ let $\iota(\tau)$ be a finite set and $\eta_{\tau k}$
($k\in\iota(\tau)$) complex $p\times m$ matrices. On $\auHom(W,I_a)\otimes_\tau\R=\pcMat_{a\times n}
(\C)$ write $x_\tau=(x^+_\tau\mid x^-_\tau)$ (positive and negative columns), and on
$\auHom(V',W)\otimes_\tau\R=\pcMat_{n\times b}(\C)$ write $y_\tau=\bigl(\begin{smallmatrix}
y^{++}_\tau&y^{+-}_\tau\\y^{-+}_\tau&y^{--}_\tau\end{smallmatrix}\bigr)$ (rows of $W$ by sign,
columns of $V'$ by sign; $y^{++}_\tau$ is $m\times p$, $y^{--}_\tau$ is $l\times l$). Put
\begin{equation}\label{pte:eq-SDMP}
\pteSD(\eta)=\prod_\tau\prod_{k\in\iota(\tau)}\det\begin{pmatrix}\eta_{\tau k}\\x^+_\tau
\end{pmatrix},\qquad
\pteMP(\eta)=\prod_\tau\det(y^{--}_\tau)^{\#\iota(\tau)+n-2l}\prod_{k\in\iota(\tau)}\det(\eta_{\tau
k}\,y^{++}_\tau),
\end{equation}
the first a polynomial on $\auHom(W,I_a)\otimes\R$ (the $(p+a)\times m$ matrix has its columns in the
fixed order), the second on $\auHom(V',W)\otimes\R$. The \emph{standard frame} is
$\eta^0=(1_p\ 0)$, i.e.\ $\eta^0_{\tau k}$ has as rows the basis vectors of the first $p$ positive
indices; then $\pteSD(\eta^0)=\prod_\tau\det(x^{\flat}_\tau)^{\#\iota(\tau)}$, where $x^\flat_\tau$ is the
$a\times a$ block of $x^+_\tau$ in the last $a$ positive columns.

\emph{Matched functions} are those of Definition~\ref{pcx:def-matched}, for the decomposition
$n=b+a$ of the rows of $\auHom(V',W)$ and columns of $\auHom(W,I_a)$ fixed by a bijection $\sigma$.
For a character $\alpha$ tame for $\pcm$ (\S\ref{pcx:ss-statement}, but with weight allowed to be any
integer vector) $\chi_\alpha$ denotes the character of $\pcR_\pcm^\times$ given by its tame parts.

\subsubsection{Three lemmas on polynomials}\label{pte:ss-poly}

\begin{lemma}[Averaging over a compact group]\label{pte:average}
Let $a\ge1$, $c=(c_\tau)\in\Z_{\ge0}^{\auPl}$ and $q$ a polynomial in the entries of square matrices
$z=(z_\tau)_\tau$, $z_\tau\in\pcMat_a(\C)$. There is $\kappa_q\in\C$ such that for every linear form
$M$ on polynomials and every $t\in\C$: if $M\bigl(q(u^{\mathsf t}z)\bigr)=t\prod_\tau\det(u_\tau)^{c_\tau}$
for all $u\in\prod_\tau\auU(a)$ (where $(u^{\mathsf t}z)_\tau=u_\tau^{\mathsf t}z_\tau$), then
$t=\kappa_q\,M\bigl(\prod_\tau\det(z_\tau)^{c_\tau}\bigr)$.
\end{lemma}
\begin{proof}
Put $\chi(u)=\prod_\tau\det(u_\tau)^{c_\tau}$ and $\bar q(z)=\int\chi(u)^{-1}q(u^{\mathsf t}z)\,du$
(probability Haar measure on the compact group $U=\prod_\tau\auU(a)$). The integrand lies in the
finite-dimensional space of polynomials of degree $\le\deg q$, so $\bar q$ is a polynomial and
$M(\bar q)=\int\chi(u)^{-1}M(q(u^{\mathsf t}z))\,du=t$. By invariance of $du$, $\bar q(v^{\mathsf t}z)=
\chi(v)\bar q(z)$ for $v\in U$. Both sides are polynomials in the entries of $v$ and $\det(v_\tau)^{-1}$, and $U$ is
Zariski dense in $\prod_\tau\auGL_a(\C)$ (its Lie algebra spans $\mathfrak{gl}_a$ over $\C$), so the
identity holds for all invertible $v$; with $z=1$, $\bar q(v^{\mathsf t})=\chi(v)\bar q(1)$, i.e.\ $\bar q=
\bar q(1)\prod_\tau\det(z_\tau)^{c_\tau}$ on a dense set, hence everywhere. So $t=M(\bar q)=\bar q(1)
M(\prod_\tau\det(z_\tau)^{c_\tau})$, and $\kappa_q=\bar q(1)$ depends only on $q$ and $c$.
\end{proof}

\begin{lemma}[Polarisation]\label{pte:polar}
Let $p<m$ and, for each $\tau$, $\eta_{\tau k}\in\pcMat_{p\times m}(\C)$ ($k\in\iota(\tau)$). There are
finitely many $u\in\prod_\tau\auU(m)$ and numbers $c_u$ such that, in the polynomial ring in the entries
of matrices $Y_\tau\in\pcMat_{m\times p}$,
\[
\prod_\tau\prod_k\det(\eta_{\tau k}Y_\tau)=\sum_uc_u\prod_\tau\prod_k\det(\eta^0u_\tau^{\mathsf t}Y_\tau).
\]
Consequently the same identity holds after substituting for the $Y_\tau$ elements of any commutative
$\C$-algebra.
\end{lemma}
\begin{proof}
For fixed $\tau$ and $c=\#\iota(\tau)$, by the Cauchy--Binet formula $\det(\eta_kY)=\sum_I
\mathrm{Pl}_I(\eta_k)\det(Y_I)$, the sum over the $p$-subsets $I$ of rows, $\mathrm{Pl}_I$ the
Pl\"ucker coordinates. So $F_\eta(Y):=\prod_k\det(\eta_kY)$ lies in the span $\mathcal S_c$ of the
products of $c$ maximal minors of $Y$, which is stable under $G=\auGL_m(\C)$ acting by $(gF)(Y)=
F(g^{\mathsf t}Y)$ (note $gF_\eta=F_{\eta g^{\mathsf t}}$). The space $\mathcal S_c$ is the degree-$c$
part of the homogeneous coordinate ring of the Grassmannian of $p$-planes in $\C^m$, which is an
irreducible $G$-module (the dual of the irreducible module of highest weight $c\varpi_p$; Borel--Weil,
see e.g.\ Fulton, \emph{Young tableaux}). For several $\tau$ the space
$\bigotimes_\tau\mathcal S_{c_\tau}$ is irreducible for $\prod_\tau G$. It contains
$F_{\eta^0}=\prod_\tau\det(Y^{\mathrm{top}}_\tau)^{c_\tau}\ne0$, so it is the span of the
$\prod_\tau G$-orbit of $F_{\eta^0}$. The map $g\mapsto gF_{\eta^0}$ is polynomial in the entries of
$g$, and $\prod_\tau\auU(m)$ is Zariski dense in $\prod_\tau G$, so every linear form vanishing on the
unitary orbit vanishes on the whole orbit: the unitary translates $F_{\eta^0u^{\mathsf t}}$ span the
same finite-dimensional space, which contains $F_\eta$.
\end{proof}

\begin{lemma}\label{pte:stackeval}
Let $Z\in\pcMat_{a\times m}(\C)$, $m=p+a$. There are $Y\in\pcMat_{m\times p}(\C)$ and $d\in\C$ with
$\det\binom{\xi}{Z}=d\det(\xi Y)$ for all $\xi\in\pcMat_{p\times m}(\C)$ (rows of $\binom\xi Z$ in the
fixed order of columns).
\end{lemma}
\begin{proof}
Choose a unitary $U=(U_1\mid U_2)\in\auU(m)$, $U_1$ of size $m\times p$, with $ZU_1=0$ (possible when
$\operatorname{rk}Z\le a$: take the columns of $U_1$ orthonormal in $\ker Z$, which has dimension
$\ge p$). Then $\binom\xi ZU=\bigl(\begin{smallmatrix}\xi U_1&\xi U_2\\0&ZU_2\end{smallmatrix}\bigr)$, so
$\det\binom\xi Z=\det(\xi U_1)\det(ZU_2)/\det U$; take $Y=U_1$, $d=\det(ZU_2)/\det U$.
\end{proof}

\subsubsection{The exchange identity with determinant data}\label{pte:ss-explicit}

\begin{corollary}[Explicit exchange]\label{pte:explicit}
Let $W$, $V'$, $I_a$, $n=b+a$, $m=p+a$, $l$, $\iota$ be as in \S\ref{pte:ss-conv}, with $a\ge1$,
$a=b+2d$ ($d\ge0$) and $2l+2\le n$. There is a real number $C_{\mathrm e}\ge0$ such that for every
modulus $\pcm\ne(1)$ and every character $\alpha$ of $\auU(1)$ tame for $\pcm$ (tame parts
non-trivial exactly at the primes of $\pcm$) of weight $(l-\#\iota(\tau))_\tau$ there is $c\in\C$
with
\[
|c|^2\,L_*\bigl(\pcm,\alpha^{-1}\upsilon^{b+d},d\bigr)=C_{\mathrm e},
\]
such that for all $\eta$, all functions $\Phi$, $\Psi$ on $\pcMat_{a\times n}(\pcR_\pcm)$,
$\pcMat_{n\times b}(\pcR_\pcm)$ matched for $\chi_\alpha^{-1}$, all characters $\beta$ and all $g\in
\auU(W)(\auA)$,
\begin{equation}\label{pte:eq-explicit}
c\,\pteCS_{K_{I_a}}\Bigl(h\mapsto\theta^{W,I_a}_{\beta,\alpha}\bigl(g,h;[\pteSD(\eta),\Phi]\bigr)\Bigr)
=\int_{[\auU(V')]}\theta^{V',W}_{\alpha,\alpha\beta^{-1}}\bigl(x,g;[\pteMP(\eta),\Psi]\bigr)\,dx .
\end{equation}
Both sides converge absolutely; $C_{\mathrm e}$ does not depend on $\pcm,\alpha,\beta,\eta,\Phi,\Psi,g$.
\end{corollary}

Unlike Theorem~\ref{pcx:X}, the corollary does not assert $c\ne0$ (that is, $C_{\mathrm e}>0$): the
explicit polynomial $\pteSD(\eta)$ replaces the vector $x$ of the theorem, and the number $\kappa$
relating them may vanish. In the application, $C_{\mathrm e}>0$ follows a posteriori
(Proposition~\ref{po:glue}(G2)).

\emph{Matching with Theorem~\ref{pcx:X}.} Put $s_0=n/2-l\ge1$, $\nu=\alpha^{-1}\upsilon^{b+d}$ and
$\gamma=\nu^{-1}\upsilon^{n/2}$. Since $n/2=b+d$, $\gamma=\alpha$. The weight of $\nu$ is
$\mu_\tau=\#\iota(\tau)-l+(b+d)=\#\iota(\tau)+s_0\ge1$, its tame parts are those of $\alpha^{-1}$, so $\nu$ is
tame for $\pcm$ of weight $\mu$ and the character $\chi$ of Definition~\ref{pcx:def-matched} attached
to $\nu$ is $\chi_\alpha^{-1}$. The hypotheses of Theorem~\ref{pcx:X} for $(W,I_a,V')$ hold: $a=b+2d$;
$q=l=0+l=q_{I_a}+q_{V'}$; $2q+2\le n$; $n=2l+2s_0\le2\mu_\tau+2l$.

\begin{proof}
\emph{Step 0: the vector $v$.} Let $Q_0=\pteMP(\eta^0)$ (for the given $\iota$), a polynomial Fock
vector of $\auHom(V',W)\otimes\R$, which has finite moments. Theorem~\ref{pcx:X} for $(W,V,V')=(W,I_a,V')$,
$\mu$ as above and $v=Q_0$ gives $C>0$ and a Fock vector $x_0$ of $\auHom(W,I_a)\otimes\R$, and for each $\alpha$
a number $c_0$ with $|c_0|^2=C\,L_*(\pcm,\nu,d)\ne0$, such that (taking $\alpha_X=\beta$, $\alpha'=\alpha
\beta^{-1}$ in \eqref{pcx:eq-X}, so $\alpha_X\alpha'=\gamma$)
\begin{equation}\label{pte:eq-X0}
\int_{[\auU(I_a)]}\theta^{W,I_a}_{\beta,\alpha}(g,h;[x_0,\Phi])\,dh=c_0\int_{[\auU(V')]}
\theta^{V',W}_{\alpha,\alpha\beta^{-1}}(x,g;[Q_0,\Psi])\,dx
\end{equation}
for all matched $\Phi,\Psi$ and all $g$. By \eqref{pte:eq-twist} both sides of \eqref{pte:eq-X0} and of
\eqref{pte:eq-explicit} change by the same factor $\beta(\det g)^{-1}$ when $\beta$ is changed, so we
may fix $\beta$; we keep it general in the notation.

\emph{Step 1: the two functionals.} Fix $\alpha,\Phi,\Psi$ and $g$. For a Fock vector $w$ of
$\auHom(W,I_a)\otimes\R$ with finite moments and $y\in\auU(W)(E^+\otimes\R)$ put $L_y(w)=\int_{[\auU(I_a)]}
\theta^{W,I_a}_{\beta,\alpha}(gy,h;[w,\Phi])\,dh$; $[\auU(I_a)]$ is compact and the kernel depends
continuously on $w$, uniformly on compacta, so $L_y$ is a continuous linear form. For a polynomial
$Q$ on $\auHom(V',W)\otimes\R$ put $\ell'(Q)=\int_{[\auU(V')]}\theta^{V',W}_{\alpha,\alpha\beta^{-1}}
(x,g;[Q,\Psi])\,dx$, which converges absolutely (Corollary~\ref{pcx:conv}: $\Psi$ is supported on
left-invertible matrices) and is linear in $Q$. Write $\omega_\infty$ for the archimedean Weil
representations (with the archimedean components of the characters). Then:
\begin{enumerate}
\item[(a)] $L_y(\omega_\infty(k,u)w)=L_{yk}(w)$ for $k\in\auU(W)_\infty$, $u\in\auU(I_a)_\infty$: the
action of $k$ moves the point $gy$, and the integral over $h$ absorbs $u$ (right invariance of $dh$);
\item[(b)] if $\omega_\infty(1,y)$, acting on $\auHom(V',W)\otimes\R$ through the second group,
carries $Q_0$ to $Q$, then $L_y(x_0)=c_0\ell'(Q)$: this is \eqref{pte:eq-X0} at the point $gy$.
\end{enumerate}

\emph{Step 2: action on the two polynomials.} By Step~0 we may take the archimedean weight of $\beta$ to be $0$.
Then, by Corollary~\ref{ptf:SDlaw}(a),(b) (the transpose conventions are fixed in \S\ref{ptf:ss-fock}), for $k$ in the
positive block $\prod_\tau\auU(m)$ of $\auU(W)_\infty$ and $u\in\auU(I_a)_\infty$,
\begin{equation}\label{pte:eq-SDlaw}
\omega_\infty(k,u)\,\pteSD(\eta k^{-\mathsf t})=\prod_\tau\det(k_\tau)^{-\#\iota(\tau)}\,\pteSD(\eta),
\qquad
\omega_\infty(1,k)\,\pteMP(\eta)=\prod_\tau\det(k_\tau)^{-\#\iota(\tau)}\,\pteMP(\eta k^{\mathsf t}),
\end{equation}
with $(\eta k^{\pm\mathsf t})_{\tau j}=\eta_{\tau j}k_\tau^{\pm\mathsf t}$, and $\omega_\infty(1,u)\pteSD(\eta)=\pteSD(\eta)$ for all
$u\in\auU(I_a)_\infty$.

\emph{Step 3: a circle mean.} For $z\in\C$, $|z|=1$, let $t_z\in\prod_\tau\auU(m)$ be diagonal with
entries $z^{-2}$ at the first $p$ positive indices and $z^{-1}$ at the last $a$, and put
$\Omega_z(w)=z^{-aN_\iota}\omega_\infty(t_z,z^{-1}1_a)w$, $N_\iota=\sum_\tau\#\iota(\tau)$. We claim:
there is $\kappa\in\C$ (depending on $x_0$) such that for every continuous linear form $\ell$ on the
Fock space with $\ell\circ\omega_\infty(1,u)=\ell$ for all $u\in\auU(I_a)_\infty$,
\begin{equation}\label{pte:eq-circle}
\kappa\,\ell\bigl(\pteSD(\eta^0)\bigr)=\frac1{2\pi}\int_0^{2\pi}\ell\bigl(\Omega_{e^{i\vartheta}}(x_0)
\bigr)\,d\vartheta .
\end{equation}
The right side is continuous in $x_0$ (uniformly over $\ell$ in bounded sets) and the set of
functions $\ell\mapsto\kappa\,\ell(\pteSD(\eta^0))$, $\kappa\in\C$, is a closed line; polynomial vectors are
dense in the Fock space; so it suffices to prove the claim for $x_0$ a polynomial $P$ in the
entries of $x$. By Corollary~\ref{ptf:SDlaw}(c), $\Omega_z$ acts on polynomials by the substitution that multiplies every
variable outside the block $x^\flat$ (the last $a$ positive columns) by $z$ and fixes the entries of
$x^\flat$ (the scalar $z^{-aN_\iota}$ cancels the vacuum character). Hence $\ell(\Omega_z P)=F(z)$ for
a polynomial $F$, and its circle mean is $F(0)=\ell(P_0)$, where $P_0$ is $P$ with all variables outside
$x^\flat$ set to $0$. Now $P_0=q(x^\flat)$ is a polynomial in square matrices, and $M(r)=\ell(r(x^\flat))$
is a linear form with $M(q(u^{\mathsf t}x^\flat))$ computed as follows: by
Corollary~\ref{ptf:SDlaw}(c), $\omega_\infty(1,u)$ acts on polynomials in $x^\flat$ by $x^\flat\mapsto u^{\mathsf t}x^\flat$
times $\prod_\tau\det(u_\tau)^{-\#\iota(\tau)}$, and $\ell$ is
invariant, so $M(q(u^{\mathsf t}x^\flat))=\ell(P_0)\prod_\tau\det(u_\tau)^{\#\iota(\tau)}$. By
Lemma~\ref{pte:average}, $\ell(P_0)=\kappa_q\,\ell(\prod_\tau\det(x^\flat_\tau)^{\#\iota(\tau)})=\kappa_q
\,\ell(\pteSD(\eta^0))$, and $\kappa=\kappa_q$ depends only on $P$.

\emph{Step 4: unitary frames.} Let $k\in\prod_\tau\auU(m)$ and $\xi=\eta^0k^{\mathsf t}$. We show
\begin{equation}\label{pte:eq-unitary}
\kappa\,L_1(\pteSD(\xi))=c_0\,\ell'(\pteMP(\xi)).
\end{equation}
By \eqref{pte:eq-SDlaw} with $\eta=\eta^0$ (and $k^{-1}$ in place of $k$),
$\kappa L_1(\pteSD(\xi))=\kappa\prod_\tau\det(k_\tau)^{\#\iota(\tau)}L_{k}(\pteSD(\eta^0))$ by (a). The
form $\ell=\prod\det(k_\tau)^{\#\iota(\tau)}L_k$ is invariant under $\auU(I_a)_\infty$ by (a), so by
\eqref{pte:eq-circle} this equals the circle mean of $\prod\det(k_\tau)^{\#\iota(\tau)}L_k(\Omega_z x_0)$.
By (a), $L_k(\Omega_zx_0)=z^{-aN_\iota}L_{kt_z}(x_0)$, and by (b) with $y=kt_z$ and \eqref{pte:eq-SDlaw}
(which gives $\omega_\infty(1,kt_z)Q_0=\prod\det(kt_z)_\tau^{-\#\iota(\tau)}\pteMP(\eta^0(kt_z)^{\mathsf t})$ and
$\eta^0(kt_z)^{\mathsf t}=z^{-2}\xi$, because $t_z$ is $z^{-2}$ on the first $p$ indices),
\[
L_{kt_z}(x_0)=c_0\prod_\tau\det(k_\tau t_{z,\tau})^{-\#\iota(\tau)}\ell'(\pteMP(z^{-2}\xi))
=c_0\prod_\tau\det(k_\tau)^{-\#\iota(\tau)}z^{(2p+a)N_\iota}z^{-2pN_\iota}\ell'(\pteMP(\xi)),
\]
using $\det t_{z,\tau}=z^{-(2p+a)}$ and the homogeneity $\pteMP(w\xi)=w^{pN_\iota}\pteMP(\xi)$. All powers
of $z$ and of $\det k_\tau$ cancel: the integrand of the circle mean is the constant $c_0\ell'(\pteMP(\xi))$.

\emph{Step 5: all $\eta$.} By Lemma~\ref{pte:polar} ($p<m$ since $a\ge1$) there are $u$'s and $c_u$
with $\prod\det(\eta_{\tau k}Y_\tau)=\sum c_u\prod\det(\eta^0u_\tau^{\mathsf t}Y_\tau)$. Substituting
$Y_\tau=y^{++}_\tau$ gives $\pteMP(\eta)=\sum_uc_u\pteMP(\eta^0u^{\mathsf t})$. Evaluating at a point
$x$, Lemma~\ref{pte:stackeval} (applied to $Z=x^+_\tau$ at each $\tau$) gives $Y_\tau$ and $d_\tau$
with $\det\binom{\zeta}{x^+_\tau}=d_\tau\det(\zeta Y_\tau)$ for all $\zeta$; so $\pteSD(\eta)(x)=
\prod_\tau d_\tau^{\#\iota(\tau)}\prod\det(\eta_{\tau k}Y_\tau)$ and the same identity with the same
$c_u$ holds for $\pteSD$ at every point: $\pteSD(\eta)=\sum_uc_u\pteSD(\eta^0u^{\mathsf t})$. By
\eqref{pte:eq-unitary} and linearity, $\kappa L_1(\pteSD(\eta))=c_0\ell'(\pteMP(\eta))$ for all $\eta$.

\emph{Step 6: class sum.} By Step 2 the integrand $h\mapsto\theta^{W,I_a}_{\beta,\alpha}(g,h;
[\pteSD(\eta),\Phi])$ is invariant under right translation by $\auU(I_a)_\infty$. It is invariant under
right translation by $K_{I_a}$: by Lemma~\ref{pc:level}, translation by $k\in K_{I_a}$ replaces $\Phi$
by $\alpha(\det k)\Phi(\bar k^{-1}\,\cdot)$, $\bar k$ the reduction modulo $\pcm$, and $\alpha(\det k)=
\chi_\alpha(\det\bar k)$ on integral points (tameness) while $\Phi(\bar k^{-1}x)=\chi_\alpha(\det\bar k)^{-1}
\Phi(x)$ because $\Phi$ is matched for $\chi_\alpha^{-1}$ (Definition~\ref{pcx:def-matched}: $\Phi$
lives on the last rows of inverses of invertible matrices, and left multiplication by $\bar k^{-1}$
changes $\det$ by $\det\bar k^{-1}$). With $\mathrm{vol}(K_{I_a})=\mathrm{vol}(\auU(I_a)_\infty)=1$ the
integral over $[\auU(I_a)]$ is therefore the class sum $\pteCS_{K_{I_a}}$. Hence
$\kappa\,\pteCS_{K_{I_a}}(\cdots)=c_0\ell'(\pteMP(\eta))$, which is \eqref{pte:eq-explicit} with
$c=\kappa/c_0$ ($c_0\ne0$). Then $|c|^2L_*(\pcm,\nu,d)=|\kappa|^2/C=:C_{\mathrm e}$, and $\kappa$ depends
only on $x_0$, hence not on $\pcm,\alpha$. Convergence of the right side is Corollary~\ref{pcx:conv}; the
left side is a finite sum.
\end{proof}

\begin{remark}\label{pte:rem-explicit}
(i) The averaging of Step~3 takes place on the definite group $\auU(I_a)_\infty$, which is compact;
the non-compact group $\auU(V')$ is never averaged over. This is why the argument applies with
$V'\ne0$ of any signature, in contrast with the isotypic shortcut of Proposition~\ref{pex:E1}. (ii)
For $b=0$ ($V'=0$, $p=l=0$, $\pteMP=1$, $\Psi=1$ and $\Phi=\chi_\alpha\circ\det$ by
Definition~\ref{pcx:def-matched}) the right side of \eqref{pte:eq-explicit} is $\alpha(\det g)\beta(\det
g)^{-1}$ times the volume $1$ of the point $[\auU(0)]$.
\end{remark}

\subsubsection{One exchange summed over classes}\label{pte:ss-one}

\begin{lemma}[Product data are matched]\label{pte:matched}
Let $\chi$ be a character of $\pcR^\times$ ($\pcR=\pcR_\pcm$), extended by $0$, and $\Psi'$ a function on
$\pcMat_{n\times b}(\pcR)$ with $\Psi'(y\gamma)=\Psi'(y)$ for $\gamma\in\auGL_b(\pcR)$. On
$\pcMat_{n\times(b+n)}(\pcR)\ni x=(x'\mid x'')$ and $\pcMat_{(b+n)\times b}(\pcR)\ni y=\binom{y'}{y''}$
put $\Phi(x)=\chi(\det x'')\Psi'(x''^{-1}x')$ (zero if $x''$ is not invertible) and $\Psi(y)=
\chi(\det y')\Psi'(y'')$. Then $\Phi,\Psi$ are matched for $\chi^{-1}$ (Definition~\ref{pcx:def-matched},
with $N=b+n$, the first $b$ indices first).
\end{lemma}
\begin{proof}
Let $k\in\auGL_{b+n}(\pcR)$, $k=\bigl(\begin{smallmatrix}A&B\\C&D\end{smallmatrix}\bigr)$, $k^{-1}=
\bigl(\begin{smallmatrix}A'&B'\\C'&D'\end{smallmatrix}\bigr)$; $k_{[b]}=\binom AC$, $k^{[n]}=(C'\mid
D')$. From $k^{-1}k=1$: $C'A+D'C=0$. If $D'$ is invertible, $A$ is invertible (Schur complement:
$\det k^{-1}=\det D'\det(A'-B'D'^{-1}C')$ and $A=(A'-B'D'^{-1}C')^{-1}$), $D'^{-1}C'=-CA^{-1}$ and
$\det k=\det A/\det D'$. Then $\Phi(k^{[n]})=\chi(\det D')\Psi'(-CA^{-1})=\chi(\det D')\Psi'(C)$ (invariance
under $-A^{-1}$), and $\chi^{-1}(\det k)\Psi(k_{[b]})=\chi(\det D')\chi(\det A)^{-1}\chi(\det A)\Psi'(C)$:
equal. If $D'$ is not invertible, both vanish ($A$ is then not invertible, by the same identity read
from $k$). Supports: if $\Psi(y)\ne0$ then $y'$ is invertible and $y=k_{[b]}$ for $k=\bigl(\begin{smallmatrix}
y'&0\\y''&1\end{smallmatrix}\bigr)$; if $\Phi(x)\ne0$ then $x''$ is invertible and $x=k^{[n]}$ for
$k^{-1}=\bigl(\begin{smallmatrix}1&0\\x'&x''\end{smallmatrix}\bigr)$.
\end{proof}

\begin{lemma}[One exchange, summed]\label{pte:one}
Let $V'$, $W'$ be spaces of ranks $b$, $n$, $I_b$ the positive definite standard space, $\mathsf W=
I_b\oplus W'$, $\pcm\ne(1)$, $\chi$ a character of $\pcR_\pcm^\times$, $a_1,a_2,\beta_0$ characters, $P$ a
polynomial on $\auHom(V',W')\otimes\R$ (read on $\auHom(V',\mathsf W)$ through the rows of $W'$) and
$g\in\auU(W')(\auA)$. Suppose: (i) a function $F$ on $\auU(V')(\auA)$ satisfies $c_0\,\pteCS_{K_{I_b}}
(h'\mapsto\mathrm a(h')\theta^{V',I_b}_{a_1,\beta_0}(x,h';[1,\chi\circ\det]))=F(x)$ for all $x$, for some
number $c_0$ and function $\mathrm a$; (ii) for functions $K(\Phi,h')$ on $\auU(I_n)(\auA_f)$ and a
number $c$: for all $\Phi,\Psi$ matched for $\chi^{-1}$ as in Lemma~\ref{pte:matched} and all
$h'\in\auU(I_b)(\auA_f)$, $c\,\pteCS(K(\Phi,h'))=\int_{[\auU(V')]}\theta^{V',\mathsf W}_{a_1a_2,\beta_0}(x,
(h',g);[P,\Psi])\,dx$. Let $\Psi'$ be invariant under $\auGL_b(\pcR)$ on the right and $\Phi_{\Psi'}$ as
in Lemma~\ref{pte:matched}. Then
\[
c\,c_0\,\pteCS_{K_{I_b}}\bigl(h'\mapsto\mathrm a(h')\,\pteCS(K(\Phi_{\Psi'},h'))\bigr)=\int_{[\auU(V')]}
F(x)\,\theta^{V',W'}_{a_2,\beta_0}(x,g;[P,\Psi'])\,dx,
\]
the integrals converging absolutely.
\end{lemma}
\begin{proof}
For $\Psi=\Psi_{\Psi'}$ of Lemma~\ref{pte:matched}, $\Psi(y)=\chi(\det y')\Psi'(y'')$ is a product over
the two blocks of rows, and the polynomial involves only the rows of $W'$; by the compatibility of
kernels with orthogonal sums in the target (the Weil representation of $\auHom(V',I_b\oplus W')=
\auHom(V',I_b)\oplus\auHom(V',W')$ is the tensor product of the two Weil representations, with the
splittings multiplying, as in (W7) of \S\ref{po:ss-glue}), together with the convention
$\omega_{\alpha,\beta}=\alpha(\det x)^{-1}\beta(\det h)\,\omega$: the source group $\auU(V')$ is common to
both factors, so its character splits as $a_1a_2$, while $\det(h',g)=\det h'\det g$, so the target character $\beta_0$
appears on each block. Hence $\theta^{V',\mathsf W}_{a_1a_2,\beta_0}(x,(h',g);[P,\Psi])=
\theta^{V',I_b}_{a_1,\beta_0}(x,h';[1,\chi\circ\det])\,\theta^{V',W'}_{a_2,\beta_0}(x,g;[P,\Psi'])$. The left
factor's datum is supported on invertible matrices, so the product is integrable
(Corollary~\ref{pcx:conv}); by (ii) its integral is $c\,\pteCS(K(\Phi_{\Psi'},h'))$. Multiply by
$c_0\mathrm a(h')$ and sum over the finitely many classes $h'$; the finite sum commutes with the
integral, and by (i) $c_0\pteCS(\mathrm a\,\theta^{V',I_b}(x,\cdot))=F(x)$.
\end{proof}

\begin{lemma}[Translates of a datum]\label{pte:translates}
Let $F$ on $\auU(V')(\auA)$ be automorphic, rapidly decreasing and right invariant under the
stabiliser $K_{V'}$ of the lattice, $\beta'$ a character with trivial tame part, $\psi$ a function on
$\pcMat_{n\times b}(\pcR_\pcm)$, and $\pcm'\subseteq\pcm$. Then
\[
\int_{[\auU(V')]}F(x)\,\theta^{V',W'}_{\beta',a}\Bigl(x,g;\bigl[P,\textstyle\sum_{\gamma\in\auGL_b(\pcR_\pcm)}\psi(\,\cdot\,\gamma)
\bigr]\Bigr)dx=\#\auGL_b(\pcR_\pcm)\int_{[\auU(V')]}F(x)\,\theta^{V',W'}_{\beta',a}(x,g;[P,\psi])\,dx,
\]
the datum on the left read modulo $\pcm'$.
\end{lemma}
\begin{proof}
For an integral lift $\tilde\gamma\in K_{V'}$ of $\gamma$, by Lemma~\ref{pc:level} right translation
of the kernel by $\tilde\gamma$ replaces $\psi$ by $\beta'(\det\tilde\gamma)\psi(\,\cdot\,\gamma)=
\psi(\,\cdot\,\gamma)$ ($\beta'$ has trivial tame part, so it is trivial on integral points). As $F$ is
$K_{V'}$-invariant and the measure on $[\auU(V')]$ is right invariant, each translate gives the same
integral; all integrals converge since $F$ is rapidly decreasing and kernels are of moderate growth.
Data modulo $\pcm$ and their pull-backs modulo $\pcm'$ define the same datum (Lemma~\ref{pcx:mod}).
\end{proof}

\subsubsection{The carrier}\label{pte:ss-carrier}

Let $m$ be even, $I_m$ the positive definite standard space of rank $m$, $I_m^-$ its negative, and $S$
the space of signature $(2m,m)$ with the standard frame (signs $+1$ on $2m$ indices, $-1$ on $m$). The
\emph{carrier} is
\[
\varphi^\circ(x)=\int_{[\auU(I_m^-)]}\theta^{I_m^-,S}_{1,\upsilon^m}\Bigl(y,x;\Bigl[\prod_\tau\det(y^-_\tau)^{2m},\plone\Bigr]\Bigr)\,dy
\qquad(x\in\auU(S)(\auA)),
\]
the theta lift of the constant $1$, with modulus $(1)$ (vacuum at all finite places) and $y^-_\tau$ the
$m\times m$ block of variables in the negative rows of $S$. For a character $c$ of $\auU(1)$ put
$\mathrm a_c(h)=\pteCS_{K_{I_m}}\bigl(h_1\mapsto\theta^{I_m,I_{3m}}_{1,c}(h_1,h;[1,\plone])\bigr)$ for
$h\in\auU(I_{3m})(\auA_f)$.

\begin{lemma}[Negation]\label{pte:conj}
Let $S_1,T$ be spaces of ranks $n_1,n'$ and $S_1^-$ the space with the same frame and the negated form. There is
$\epsilon\in\C$, $|\epsilon|=1$, such that for all characters $a,b$, moduli $\pcm$, functions $\Psi$, polynomials $P$,
and $g,h$, $\aubar{\theta^{S_1,T}_{b,a}(g,h;[P,\Psi])}=\epsilon\,\theta^{S_1^-,T}_{b^{-1}\upsilon^{n'},a^{-1}\upsilon^{n_1}}
(g,h;[\aubar P,\aubar\Psi])$, where $\aubar P$ has conjugate coefficients and $\auU(S_1^-)=\auU(S_1)$ (the source
character is shifted by $\upsilon$ to the rank of $T$, the target character by $\upsilon$ to the rank of $S_1$).
\end{lemma}
\begin{proof}
This is Lemma~\ref{ptf:conj} (with $\alpha=b$, $\beta=a$, $n_S=n_1$, $n_T=n'$); its proof constructs the
antilinear intertwiner, computes the phase (Lemma~\ref{ptf:theta}) and the change of splitting
(Lemma~\ref{ptf:weil}).
\end{proof}

\begin{proposition}[First two exchanges]\label{pte:carrier}
Let $m$ be even. There is $C_{12}\ge0$ such that for every modulus $\pcm\ne(1)$ and every character $c$
of weight $0$, tame for $\pcm$ with residue character $\chi$, there is $\gamma\in\C$ with
\[
|\gamma|^2\,L_*(\pcm,c^{-1}\upsilon^m,m)\,L_*(\pcm,c\,\upsilon^{m/2},m/2)=C_{12},\qquad
\gamma\,\pteCS_{K_{I_{3m}}}\Bigl(h\mapsto\mathrm a_c(h)\,\theta^{S,I_{3m}}_{c,\upsilon^m}\bigl(x,h;[1,\chi\circ
\det]\bigr)\Bigr)=\varphi^\circ(x)
\]
for all $x\in\auU(S)(\auA)$.
\end{proposition}
\begin{proof}
\emph{First exchange.} Corollary~\ref{pte:explicit} with $W=I_m$, $V'=0$, $I_a=I_m$ ($n=a=m$, $b=p=l=0$,
$d=m/2$, $2l+2\le m$), $\iota(\tau)=\emptyset$ (so $\pteSD=\pteMP=1$), the character $c^{-1}$ (weight
$0=l-\#\iota$), $\beta=c^{-1}$ and the matched pair $\Phi=\chi^{-1}\circ\det$, $\Psi=1$ (Remark~\ref{pte:rem-explicit}(ii);
$\chi^{-1}$ is the residue character of $c^{-1}$) gives $\gamma_1$ with $|\gamma_1|^2L_*(\pcm,c\upsilon^{m/2},m/2)=
C_1$ and $\gamma_1\pteCS(\theta^{I_m,I_m}_{c^{-1},c^{-1}}(g,\cdot;[1,\chi^{-1}\circ\det]))=1$ for all $g$.
By Lemma~\ref{pte:conj} (with $S_1=T=I_m$, equal ranks), conjugating, $\aubar{\gamma_1}\epsilon\,\pteCS(\theta^{I_m^-,
I_m}_{c\upsilon^m,c\upsilon^m}(y,\cdot;[1,\chi\circ\det]))=1$ for all $y\in\auU(I_m^-)(\auA)$ (the datum
$\aubar{\chi^{-1}(\det x)}=\chi(\det x)$).

\emph{Second exchange.} Corollary~\ref{pte:explicit} with $W=I_m\oplus S$ (rank $4m$, $l=m$ negative
signs), $V'=I_m^-$ ($b=m$, $p=0$, $l=m$), $I_a=I_{3m}$ ($a=3m=b+2m$, $d=m$), $\iota(\tau)=\emptyset$, the
character $\alpha=c\upsilon^m$ (weight $m=l-0$, residue character $\chi$) gives $\gamma_2$ with
$|\gamma_2|^2L_*(\pcm,(c\upsilon^m)^{-1}\upsilon^{2m},m)=C_2$; note $(c\upsilon^m)^{-1}\upsilon^{2m}=c^{-1}\upsilon^m$.
Here $\pteSD=1$ and $\pteMP=\prod_\tau\det(y^{--}_\tau)^{2m}$, the negative block being the rows of $S$ of
sign $-1$ against $I_m^-$. Apply Lemma~\ref{pte:one} with $V'=I_m^-$, $W'=S$, $I_b=I_m$, the character
$\chi$, $\Psi'=1$ (so $\Phi_{\Psi'}(x)=\chi(\det x'')$), the characters $(a_1,a_2,\beta_0)=(c\upsilon^m,1,c\upsilon^m)$, $\mathrm a=$ the constant $1$, $F=1$ and $c_0=\aubar{\gamma_1}\epsilon$
(hypothesis (i) is the conjugated first exchange), and (ii) the second exchange (with $\beta=1$) at the
points $(h',x)$. It gives $\gamma_2c_0\pteCS_{K_{I_m}}(h'\mapsto\pteCS_{K_{I_{3m}}}(h\mapsto\theta^{I_m\oplus
S,I_{3m}}_{1,c\upsilon^m}((h',x),h;[1,\Phi_{\Psi'}])))=\int_{[\auU(I_m^-)]}\theta^{I_m^-,S}_{1,c\upsilon^m}(y,x;
[\prod\det(y^-)^{2m},\plone])\,dy$ (the right side of \eqref{pte:eq-explicit} carries $\theta^{I_m^-,I_m\oplus S}_{c\upsilon^m,c\upsilon^m}$, split as
$a_1a_2=c\upsilon^m\cdot1$ with target character $c\upsilon^m$; the factor $F=1$).
By the compatibility of kernels with orthogonal sums of the source space (W7), with the target character split
as $c\cdot\upsilon^m$, $\theta^{I_m\oplus S,I_{3m}}_{1,c\upsilon^m}((h',x),h;[1,\Phi_{\Psi'}])=\theta^{I_m,I_{3m}}_{1,c}
(h',h;[1,\plone])\,\theta^{S,I_{3m}}_{1,\upsilon^m}(x,h;[1,\chi\circ\det])$ (the datum is $1$ on the columns of
$I_m$ and $\chi\circ\det$ on those of $S$); exchanging the two finite class sums, the left side is
$\gamma_2c_0\,\pteCS(\mathrm a_c\,\theta^{S,I_{3m}}_{1,\upsilon^m}(x,\cdot;[1,\chi\circ\det]))$. The right side is
$\int\theta^{I_m^-,S}_{1,c\upsilon^m}(y,x;[\prod\det(y^-)^{2m},\plone])\,dy=c(\det x)\varphi^\circ(x)$ by
\eqref{pte:eq-twist}, and on the left $\theta^{S,I_{3m}}_{1,\upsilon^m}(x,h)=c(\det x)\theta^{S,I_{3m}}_{c,\upsilon^m}(x,h)$;
the factor $c(\det x)\ne0$ cancels. So the claim holds with $\gamma=\gamma_2c_0$, and $|\gamma|^2=|\gamma_2|^2|\gamma_1|^2$ gives $C_{12}=C_1C_2$.
\end{proof}

\begin{lemma}[Properties of the carrier]\label{pte:carrierprop}
Assume a character $c$ of weight $0$ tame for some $\pcm\ne(1)$ exists. Then $\varphi^\circ$ is a finite
linear combination of the functions $x\mapsto\theta^{S,I_{3m}}_{c,\upsilon^m}(x,h;[1,\chi\circ\det])$; it is
continuous, left invariant under $\auU(S)(E^+)$, right invariant under $K_S$, rapidly decreasing, and
holomorphic. Moreover $\varphi^\circ\ne0$.
\end{lemma}
\begin{proof}
If $\gamma=0$ in Proposition~\ref{pte:carrier} then $\varphi^\circ=0$; otherwise divide. The class sum is a
finite sum. Each kernel is right invariant under $K_S$ (Lemma~\ref{pc:level}: $c(\det k)^{-1}$ cancels
$\chi(\det k)$), rapidly decreasing because its finite datum $\chi\circ\det$ is supported on invertible
matrices (Proposition~\ref{pcx:decay}), and holomorphic since its polynomial is $1$ (the vacuum, which is
holomorphic for $\auU(S)_\infty$ in the order with the positive vectors first). The non-vanishing
of $\varphi^\circ$ is Lemma~\ref{pte:carriernz}, whose proof uses neither $c$ nor
Proposition~\ref{pte:carrier}.
\end{proof}

\begin{lemma}[The carrier is not zero]\label{pte:carriernz}
$\varphi^\circ\ne0$; indeed a Fourier coefficient of $\varphi^\circ$ along a unipotent subgroup of
$\auU(S)$ is not identically zero.
\end{lemma}
\begin{proof}
Write $V=I_m^-$ (frame $\varepsilon_1,\dots,\varepsilon_m$, $\langle\varepsilon_k,\varepsilon_k\rangle=-1$),
$\mathsf W=\auHom(V,S)$, and $f=[v,\plone]=v\otimes\psvac$ for the datum of the definition, with
$v=\prod_\tau\det(y^-_\tau)^{2m}$ and $\psvac$ the vacuum of the lattice $L=\alpha\auHom(L_V,L_S)$ at every
finite place (\S\ref{pc:ss-data}). Then $\varphi^\circ(x)=\int_{[\auU(V)]}\theta(\omega(y,x)f)\,dy$ with
$\omega(y,x)=\upsilon^m(\det x)\,\omega_{\auHom(V,S)}(y,x)$; the operator $\omega_{\auHom(V,S)}(y,x)$ is
$\omega_{\mathsf W}$ of the element $w\mapsto xwy^{-1}$ of $\auU(\mathsf W)$, and $\omega_{\mathsf W}$ is a
representation of $\auU(\mathsf W)(\auA)$. We use (W1)--(W4) of Lemma~\ref{ps:weil}, the local forms
(c), (d) of Proposition~\ref{pwf:W2}, and the mixed model behind them: for an isotropic subspace $I$ with
a complement $I^*$ in duality with it, $\psHS_v(\mathsf W)$ is a space of functions on $I^*_v$ with values
in $\psHS_v(\mathsf W_I)$, $\psZ_{I,v}$ is evaluation at $0$ and $\psrho_v(u)$, $u\in I^*_v$, is translation
of the variable by $u$ (the proof of (W2), (W3) in Lemma~\ref{ps:weil}; Proposition~\ref{pwf:W2}). Thus
$u\mapsto\psZ_{I,v}\psrho_v(u)F$ is the function $F$ itself.

\emph{A Witt decomposition of $S$.} Let $e_1,\dots,e_{2m}$ be the positive and $f_1,\dots,f_m$ the negative
vectors of the frame of $S$, and $\omega_0=(1+\sqrt{-d})/2$, so $\omega_0+\aubar\omega_0=1$
(\S\ref{ps:ss-weil}). Put $u_i=e_i+f_i$ and $u'_i=\aubar\omega_0e_i-\omega_0f_i$ ($1\le i\le m$),
$\ell=\langle u_i\rangle$, $\ell'=\langle u'_i\rangle$, $S_0=\langle e_{m+1},\dots,e_{2m}\rangle$. Then
$\langle u_i,u_j\rangle=\langle u'_i,u'_j\rangle=0$, $\langle u_i,u'_j\rangle=\delta_{ij}(\aubar\omega_0+\omega_0)
=\delta_{ij}$, and $\ell\oplus\ell'\perp S_0$, $S_0$ positive definite of rank $m$. The matrix
$\bigl(\begin{smallmatrix}1&\aubar\omega_0\\1&-\omega_0\end{smallmatrix}\bigr)$ is integral with determinant
$-1$, so $L_S=\bigoplus_i\widehat\auO_Eu_i\oplus\bigoplus_i\widehat\auO_Eu'_i\oplus\bigoplus_{k>m}
\widehat\auO_Ee_k$. Accordingly $\mathsf W=I\oplus I^*\oplus\mathsf W_0$ with $I=\auHom(V,\ell)$,
$I^*=\auHom(V,\ell')$, $\mathsf W_0=\auHom(V,S_0)$; $I$ and $I^*$ are isotropic, $I^\perp=I\oplus\mathsf W_0$,
and $\mathsf W_0\to I^\perp/I=\mathsf W_I$ is an isometry, by which we identify $\mathsf W_I$ with
$\mathsf W_0$, a definite space. Write $u\in I^*$ and $l\in I$ by their matrices $M_u,N_l\in\pcMat_m$ in
the bases $(\varepsilon_k)$, $(u'_i)$, resp.\ $(\varepsilon_k)$, $(u_i)$. Then
$L=(L\cap I)\oplus(L\cap I^*)\oplus(L\cap\mathsf W_0)$ with $L\cap I=\{l:N_l\in\alpha\pcMat_m(\widehat\auO_E)\}$,
$L\cap I^*=\{u:M_u\in\alpha\pcMat_m(\widehat\auO_E)\}$. As $\langle\varepsilon_k,\varepsilon_k\rangle=-1$,
$\langle l,u\rangle_{\mathsf W}=\ptetr(l^\dagger u)=-\ptetr(N_l^*M_u)$.

\emph{Step 1: two subgroups of $\auU(S)$.} For a skew-hermitian $B\in\pcMat_m(\auA_E)$ let $z_B\in\auU(S)(\auA)$
be $1$ on $\ell\oplus S_0$ and $u'_j\mapsto u'_j+\sum_iB_{ij}u_i$; it is unitary since $B+B^*=0$, and
$Z=\{z_B\}\cong\{B\}$ is a unipotent abelian group with $Z(E^+)\backslash Z(\auA)$ compact (Haar measure of
mass $1$). For $Q\in\auGL_m(\auA_E)$ let $m_Q$ be $u_j\mapsto\sum_iQ_{ij}u_i$ on $\ell$,
$u'_j\mapsto\sum_i(Q^{-*})_{ij}u'_i$ on $\ell'$, and $1$ on $S_0$ ($m_Q\in\auU(S)(\auA)$ when $Q$ is a unit at
almost all places). Let $N$ be the unipotent radical of the stabiliser of $\ell$: the elements fixing $\ell$
pointwise and acting trivially on $\ell^\perp/\ell$. For every $C\in\auHom(\ell',S_0)(\auA)$ there is $n\in
N(\auA)$ with $n^{-1}x\in x+C(x)+\ell$ for $x\in\ell'$: put $n^{-1}=1$ on $\ell$,
$n^{-1}x=x+C(x)+D(x)$ on $\ell'$, $n^{-1}s=s+C^\sharp(s)$ on $S_0$, where $C^\sharp\colon S_0\to\ell$ is
defined by $\langle C^\sharp s,x\rangle=-\langle s,Cx\rangle$ ($x\in\ell'$) and $D\colon\ell'\to\ell$ by
$\langle D x,y\rangle=-\tfrac12\langle Cx,Cy\rangle$ ($x,y\in\ell'$); a direct check of the three types of
pairs shows $n^{-1}$ is unitary. All these elements stabilise $I$ (their action on $\mathsf W$ is $w\mapsto
xw$); $z_B$, $n$ act trivially on $\mathsf W_I$, and so does $m_Q$.

\emph{Step 2: the laws.} Let $F\in\psHS(\mathsf W)$ and $u\in I^*(\auA)$. For $p\in\auU(S)(\auA)$ stabilising
$\ell$, (W1) gives $\psrho(u)\omega_{\mathsf W}(p)=\omega_{\mathsf W}(p)\psrho(p^{-1}u)$, and (W2) gives
$\psZ_I\omega_{\mathsf W}(p)=\chi_I(p)\omega_{\mathsf W_I}(\bar p)\psZ_I$, with $\chi_I(p)=1$ if $p$ is
unipotent. Also $\psZ_I\psrho(i)=\psZ_I$ for $i\in I$ and $\psZ_I\psrho(w_0)=\psrho(w_0)\psZ_I$ for
$w_0\in\mathsf W_0$ (W2), and $\psrho(a)\psrho(b)=\beta(a,b)\psrho(a+b)$ with $|\beta|=1$.
(a) For $p=z_B$: $z_B^{-1}u=u-bu$ with $bu\in I$ of matrix $BM_u$, and
$\psrho(u-bu)=\beta(-bu,u)^{-1}\psrho(-bu)\psrho(u)$; hence
\[
\psZ_I\psrho(u)\,\omega_{\mathsf W}(z_B)F=\psi_u(B)\,\psZ_I\psrho(u)F,\qquad
\psi_u(B)=\beta(-bu,u)^{-1}=\psi\bigl(\tfrac12\jmath\langle u,bu\rangle\bigr)=\psi\bigl(-\tfrac12\ptetr(\jmath B\,M_uM_u^*)\bigr).
\]
(b) For $p=n$ as in Step~1: $n^{-1}u=u+d+i$ with $d=C\circ u\in\mathsf W_0$, $i\in I$, so
$\psZ_I\psrho(u)\omega_{\mathsf W}(n)F=\zeta\,\psrho_{\mathsf W_I}(C\circ u)\,\psZ_I\psrho(u)F$ with
$|\zeta|=1$. (c) For $p=m_Q$: $m_Q^{-1}u$ has matrix $Q^*M_u$, so $\psZ_I\psrho(u)\omega_{\mathsf W}(m_Q)F
=\chi_I(m_Q)\,\psZ_I\psrho(m_Q^{-1}u)F$, $\chi_I(m_Q)\ne0$. (d) For $y\in\auU(V)(\auA)$, acting by $w\mapsto
wy^{-1}$ (it stabilises $I$, $I^*$, $\mathsf W_0$): $\psZ_I\psrho(u)\omega_{\mathsf W}(y)F=\chi_I(y)\,
\omega_{\mathsf W_I}(\bar y)\,\psZ_I\psrho(uy)F$, where $uy$ has matrix $M_uY$; $\chi_I=1$ on
$\auU(V)(E^+)$ and on $K_V$ ($\upsilon$ is unramified at the finite places). The same identities hold place
by place (Proposition~\ref{pwf:W2}(c), (d)).

\emph{Step 3: unfolding.} For a hermitian $H\in\pcMat_m(E)$ put
\begin{equation}\label{pte:eq-coeff}
\varphi^\circ_H(x)=\int_{Z(E^+)\backslash Z(\auA)}\varphi^\circ(z_Bx)\,\aubar{\psi_H(B)}\,dB,\qquad
\psi_H(B)=\psi\bigl(-\tfrac12\ptetr(\jmath BH)\bigr).
\end{equation}
By (W3) with this $I$, $\theta(\omega(y,z_Bx)f)=\sum_{u\in I^*(E)}\psi_u(B)\,a_u(y,x)$ with
$a_u(y,x)=\theta_{\mathsf W_I}(\psZ_I\psrho(u)\omega(y,x)f)$ (Step~2(a); $\upsilon^m(\det z_B)=1$), and by the
bound in (W3) the double series and integral over compact fundamental sets for $[\auU(V)]$ (compact, $V$
definite) and $Z(E^+)\backslash Z(\auA)$ converge absolutely. For $u\in I^*(E)$, $B\mapsto\psi_u(B)$ is a
character of $Z(E^+)\backslash Z(\auA)$, equal to $\psi_H$ if and only if $M_uM_u^*=H$: $B\mapsto\jmath B$ is a
bijection of the skew-hermitian onto the hermitian matrices, $(X,Y)\mapsto\ptetr(XY)$ is a non-degenerate
$\auA_{E^+}$-valued pairing on hermitian matrices, and $\psi$ is non-trivial on $\auA_{E^+}$. Hence
$\varphi^\circ_H(x)=\sum_{M_uM_u^*=H}\int_{[\auU(V)]}a_u(y,x)\,dy$. Take $H=1$ and $u_0\in I^*(E)$ with
$M_{u_0}=1$. The solutions of $M_uM_u^*=1$ are the $M_u=Y^{-1}$, $Y\in\auU(V)(E^+)$ ($Y^*Y=1$ is the
unitarity for $\langle\,,\rangle_V=-{}^t\bar\cdot\,\cdot$), i.e.\ the $u=u_0\gamma^{-1}$, each once. By
Step~2(d) and (W1) for $\mathsf W_I$, $a_{u_0\gamma^{-1}}(y,x)=a_{u_0}(\gamma^{-1}y,x)$ (the source character of
$\omega$ is trivial). Unfolding,
\[
\varphi^\circ_1(x)=\int_{\auU(V)(\auA)}a_{u_0}(y,x)\,dy,
\]
absolutely convergent.

\emph{Step 4: invariance in $y$.} $\omega(k,1)f=f$ for $k\in\auU(V)_\infty K_V$. At the finite places $(k,1)$
stabilises $L$ and (W4) applies. At $\tau$, Lemma~\ref{ptf:fock} with source $V$ (all signs negative,
weight $0$) and target $S$ ($n_T^+=2m$) gives $\omega_\infty(g,1)v=\prod_\tau\det(g_\tau)^{2m}\cdot
v\circ\sigma_{g,1}$, and $\sigma_{g,1}$ replaces the block $y^-=X^{--}$ by $X^{--}g^{-\mathsf t}$, which
multiplies $v$ by $\prod_\tau\det(g_\tau)^{-2m}$; so $\omega_\infty(g,1)v=v$. Since $\omega(yk,x)=\omega(y,x)
\omega(k,1)$, $a_{u_0}(yk,x)=a_{u_0}(y,x)$, and as $\auU(V)_\infty$ and $K_V$ have volume $1$,
$\varphi^\circ_1(x)=\sum_{y\in\auU(V)(\auA_f)/K_V}a_{u_0}(y,x)$.

\emph{Step 5: support at the finite places.} For $u\in I^*(\auA_f)$, $\psZ_{I,f}\psrho_f(u)\psvac=0$ unless
$u\in L\cap I^*$, and $\psZ_{I,f}\psrho_f(u)\psvac=\psZ_{I,f}\psvac$ if $u\in L\cap I^*$. Indeed, by (W4)
$\psrho_f(l)\psvac=\epsilon_L(l)\psvac$ for $l\in L$, and $\epsilon_L(l)=\psi_f(-\frac12\omega_0\jmath\langle
l,l\rangle)=1$ for $l$ in the isotropic $I$ or $I^*$. For $l\in L\cap I$, $\psrho(u)\psrho(l)=[u,l]\psrho(l)
\psrho(u)$ and $\psZ_I\psrho(l)=\psZ_I$, so $(1-[u,l])\psZ_{I,f}\psrho_f(u)\psvac=0$. Here
$[u,l]=\beta(u,l)\beta(l,u)^{-1}=\psi_f(-\jmath\,\ptetr(N_l^*M_u))$, using $\langle l,u\rangle=-\ptetr(N_l^*M_u)$
and $\psi(\jmath r)=1$ for $r\in\auA_{E^+}$. With $N_l=\alpha nE_{ik}$, $n\in\widehat\auO_E$, and
$M_u=\alpha\mu$, the condition $[u,l]=1$ for all such $l$ reads $\psi_f(-\jmath\alpha\aubar\alpha\,\bar n\mu_{ik})=1$
for all $n$, i.e.\ $\mu_{ik}\in\widehat\auO_E$ by the self-duality defining $\alpha$ (\S\ref{ps:ss-weil}).
So the vector vanishes unless $M_u\in\alpha\pcMat_m(\widehat\auO_E)$; for such $u$, $\psrho_f(u)\psvac=\psvac$.
In particular $\psZ_{I,f}\psvac\ne0$: in the mixed model $\psvac$ is a non-zero function on $I^*(\auA_f)$,
whose value at $u$ is $\psZ_{I,f}\psrho_f(u)\psvac$, which is $0$ or $\psZ_{I,f}\psvac$.

\emph{Step 6: the archimedean slice.} $v\ne0$, so in the mixed model at the infinite places
$u\mapsto\psZ_{I,\infty}\psrho_\infty(u)v$ is a continuous non-zero function on $I^*_\infty=\prod_\tau
\pcMat_m(\C)$; it is non-zero on a non-empty open set, which contains a point $u_\infty$ with invertible
matrix $U_\infty$ (invertible matrices are dense). Let $Q\in\auGL_m(\auA_E)$ have archimedean component
$U_\infty^*$ and finite component $\aubar\alpha\cdot1_m$ (a unit at almost all places), and put $m=m_Q$. Then
$m^{-1}u_0$ has matrix $Q^*$: the point $u_\infty$ at the infinite places and $\alpha1_m$ at the finite ones.

\emph{Step 7: conclusion.} Let $n\in N(\auA)$ be as in Step~1 for some $C$, and $x=nm$. By Step~2(b), (c)
and $\omega(y,x)=\upsilon^m(\det m)\omega_{\mathsf W}(n)\omega_{\mathsf W}(m)\omega_{\mathsf W}(y)$ (on
$\psHS(\mathsf W)$, with the trivial source character),
\[
a_{u_0}(y,x)=\zeta\,\upsilon^m(\det m)\chi_I(m)\;\theta_{\mathsf W_I}\bigl(\psrho_{\mathsf W_I}(C\circ u_0)\,
\psZ_I\psrho(m^{-1}u_0)\,\omega_{\mathsf W}(y)f\bigr).
\]
{\emergencystretch=6em For $y\in\auU(V)(\auA_f)$, by Step~2(d) at the finite places, $\psZ_I\psrho(m^{-1}u_0)\allowbreak \omega_{\mathsf W}(y)f=
\psZ_{I,\infty}\allowbreak \psrho_\infty(u_\infty)v\otimes\chi_I(y)\omega_{\mathsf W_I,f}(\bar y)\,\psZ_{I,f}\psrho_f(u_f)
\psvac$ with $u_f$ of matrix $\alpha Y$. By Step~5 this vanishes unless $Y\in\pcMat_m(\widehat\auO_E)$, and an
integral $Y$ with $Y^*Y=1$ lies in $K_V$. So only the coset $K_V$ contributes in Step~4, and
\[
\varphi^\circ_1(nm)=\zeta\,\upsilon^m(\det m)\chi_I(m)\,\theta_{\mathsf W_I}\bigl(\psrho_{\mathsf W_I}(C\circ u_0)
\Lambda\bigr),\qquad\Lambda=\psZ_{I,\infty}\psrho_\infty(u_\infty)v\otimes\psZ_{I,f}\psvac ,
\]
with $\Lambda\ne0$ by Steps~5 and~6 and $|\zeta|=1$. As $u_0$ is an isomorphism $V\to\ell'$, $C\circ u_0$ runs
over all of $\mathsf W_0(\auA)=\mathsf W_I(\auA)$ when $C$ runs over $\auHom(\ell',S_0)(\auA)$. If
$\varphi^\circ=0$, then $\varphi^\circ_1=0$ and $\theta_{\mathsf W_I}(\psrho_{\mathsf W_I}(w)\Lambda)=0$ for
all $w\in\mathsf W_I(\auA)$, so $\Lambda$ lies in the kernel of $F\mapsto\theta_{\mathsf W_I}(\psrho(\cdot)F)$,
a proper closed invariant subspace of the irreducible representation $\psHS(\mathsf W_I)$, hence $\Lambda=0$
(as in the proof of Lemma~\ref{pk:dict}; compare Proposition~\ref{pst:I4}). This contradiction proves
$\varphi^\circ_1\ne0$, so $\varphi^\circ\ne0$.\par}
\end{proof}

\subsubsection{The third exchange}\label{pte:ss-third}

We use the notation of \S\ref{po:ss-glue} and \S\ref{po:ss-hecke}; $m$ is even, $k=r+1-m$ is even, and
$N''=r+9m+1$. Put $W'=W\oplus(E^j)^-\oplus E^{3m-j}\oplus W^\flat$, $W^\flat=E^{6m+1}$ positive definite,
of rank $N''$ with $l=(m-j)+j=m$ negative signs; $\mathsf W=I_{3m}\oplus W'$; $V''=I_{N''}$ (the space of
\S\ref{po:ss-glue}). For a twist $\vartheta$ let $\pcm^{(2)}_\vartheta:=\pi_M^2\prod_{w\in Q_\vartheta}w$ (it has the same primes as the squarefree modulus $\pcm_\vartheta$ of \S\ref{pc:ss-data}, so $L_*(\pcm^{(2)}_\vartheta,\cdot,\cdot)=L_*(\pcm_\vartheta,\cdot,\cdot)$ by \eqref{pcx:eq-Lstar}), $\chi'=
\chi_0^j\vartheta$ (weight $0$, residue character $\chi'_\vartheta$ on $\pcR_{\pcm^{(2)}_\vartheta}^\times$: $\omega^j$ at
$\mathsf M$, $\psi_w$ at $Q_\vartheta$), $\nu''_\vartheta=\mu_j^0\vartheta\,\upsilon^{j+\lfloor(t+k-1)/2\rfloor}$ (the
character of the lift on $V''$) and $\nu^\flat=\nu''_\vartheta\upsilon^m/\chi'$, a character independent of
$\vartheta$ with trivial tame part. The \emph{choices} are $i=(\Lambda_i,h^\flat_i,\varrho_i)$: a matrix
$\Lambda_i\in\pcMat_{3m\times N''}(\auO_E/\pi_M^2)$, a class $h^\flat_i$ of $\auU(W^\flat)$ at level $\pi_M^2$, and
for each $\tau$ a family of $2m$-row selections of a fixed basis of short vectors of the lattice at
$h^\flat_i$; they give at each $\tau$ a family of complex $2m\times(3m+N''-m)$ matrices $\eta_i=(\eta_{i,\tau k})_{k\in
\iota(\tau)}$ (rows: basis vectors of the padded $W$ for the columns of $\mu^+_\tau$ of height $<r_+$, and the
selected short rows, read at $\tau$, for the $\Delta'$ columns of height $r_+$), with $\#\iota(\tau)=$ the number
of columns of $\mu^+_\tau$. The matrices $\eta_{i,\tau k}$ are those of Lemma~\ref{ptf:H}. Let $P_i=\pteMP(\eta_i)$; it involves only the rows of $W'$ (the rows of $\eta_i$ have
no component on $I_{3m}$), and it is the product of the minors of the padded $W$ and the determinants
$\det(\eta^\flat_{i,\tau k}y^{++}_{W^\flat,\tau})$ on $W^\flat$. Let $\psi_\Lambda(y)=\psi_E(\pi_M^{-2}\mathrm{tr}(\Lambda y))$
on $\pcMat_{N''\times3m}(\auO_E/\pi_M^2)$ and $\psi^{\mathrm{inv}}_\Lambda(y)=\sum_{\gamma\in\auGL_{3m}(\auO_E/\pi_M^2)}
\psi_\Lambda(y\gamma)$. Define
\begin{equation}\label{pte:eq-F}
F_i(g)=\int_{[\auU(S)]}\varphi^\circ(x)\,\theta^{S,W'}_{\nu^\flat\upsilon^m,\upsilon^m}\bigl(x,(g,1,1,h^\flat_i);
[P_i,\psi_{\Lambda_i}]\bigr)\,dx\qquad(g\in\auU(W)(\auA)),
\end{equation}
independent of $\vartheta$; this is the function $F_i$ of Proposition~\ref{po:exchange}
($h^\flat_i$ is the point $h^\flat_i$ there, and the twist-free datum $\psi_{s_i}$ is $\psi_{\Lambda_i}$). The
integral converges absolutely ($\varphi^\circ$ is rapidly decreasing, Lemma~\ref{pte:carrierprop}).

\begin{proposition}[Third exchange]\label{pte:third}
There are $C_3\ge0$ and for each twist $\vartheta$ a number $c_3(\vartheta)$ with $|c_3(\vartheta)|^2L_*(\pcm^{(2)}_\vartheta,
(\nu_*\vartheta)^{-1},z_*)=C_3$, such that, with $\gamma=\gamma(\vartheta)$ of Proposition~\ref{pte:carrier} for
$c=\chi'$ and $\mathsf N=\#\auGL_{3m}(\auO_E/\pi_M^2)$, for all $i$ and $g$:
\[
\frac{c_3(\vartheta)\gamma(\vartheta)}{\mathsf N}\,\mathcal L_{\vartheta,i}(g)=F_i(g),\qquad
\mathcal L_{\vartheta,i}(g)=\sum_{h'}\mathrm a_{\chi'}(h')\sum_{h}\theta^{\mathsf W,V''}_{\nu''_\vartheta,\,\chi'\nu^\flat
\upsilon^m}\bigl((h',g,1,1,h^\flat_i),h;[\pteSD(\eta_i),\Phi_{\vartheta,i}]\bigr),
\]
where $h'$, $h$ run over the classes of $\auU(I_{3m})$, $\auU(V'')$ at full level, each term weighted by the
inverse order of its stabiliser in $\Gamma$, and $\Phi_{\vartheta,i}(x'\mid x'')=\chi'_\vartheta(\det x'')\psi^{\mathrm{inv}}_{
\Lambda_i}(x''^{-1}x')$ on $\pcMat_{N''\times(3m+N'')}(\pcR_{\pcm^{(2)}_\vartheta})$.
\end{proposition}
\begin{proof}
Ranks: $n=3m+N''$, $a=N''$, $b=3m$ (the space $V'=S$, $p=2m$, $l=m$), and $a-b=N''-3m=7m+k=2z_*$, so
$d=z_*=\tfrac{7m}2+\tfrac k2$; $2l+2\le n$. The character $\alpha=\chi'\nu^\flat\upsilon^m=\nu''_\vartheta\upsilon^m$ is tame
for $\pcm^{(2)}_\vartheta$ with residue character $\chi'_\vartheta$ and weight $m-\#\iota(\tau)$: the weight of $\nu^\flat\upsilon^m$
is $m+j+\lfloor(t+k-1)/2\rfloor-\kappa_j(\tau)$ and $\#\iota(\tau)+j+\lfloor(t+k-1)/2\rfloor=\kappa_j(\tau)$ (the largest part
of $\mu^+_\tau$; see \S\ref{po:ss-set}). Corollary~\ref{pte:explicit} for $(\mathsf W,V'=S,I_a=V'')$, $\eta_i$,
$\beta=\nu''_\vartheta$ gives $c_3(\vartheta)$ with $|c_3|^2L_*(\pcm^{(2)}_\vartheta,\alpha^{-1}\upsilon^{3m+z_*},z_*)=C_3$; and
$\alpha^{-1}\upsilon^{3m+z_*}=(\nu_*\vartheta)^{-1}$ because $\nu_*=\mu_j^0\upsilon^{n'}$, $n'=j+(t-11m-1)/2$, and
$j+\lfloor(t+k-1)/2\rfloor+m-3m-z_*=n'$ (direct check with $z_*=7m/2+k/2$, $t$, $k$ odd/even as fixed in the data).
Apply Lemma~\ref{pte:one} with $V'=S$, $I_b=I_{3m}$, $W'$, the character $\chi'_\vartheta$, $\Psi'=\psi^{\mathrm{inv}}_{\Lambda_i}$
(right invariant under $\auGL_{3m}$ by construction), $F=\varphi^\circ$, $c_0=\gamma$, $\mathrm a=\mathrm a_{\chi'}$
(hypothesis (i) is Proposition~\ref{pte:carrier}), and (ii) the third exchange at the points
$(h',g,1,1,h^\flat_i)$ (the polynomial $P_i$ involves only the rows of $W'$). The right side of
\eqref{pte:eq-explicit} has the characters $\alpha=\chi'\cdot\nu^\flat\upsilon^m$ on $\auU(S)$ and
$\alpha\beta^{-1}=\upsilon^m$ on $\auU(\mathsf W)$; we take $(a_1,a_2,\beta_0)=(\chi',\nu^\flat\upsilon^m,\upsilon^m)$, so that
$\theta^{S,I_{3m}}_{a_1,\beta_0}=\theta^{S,I_{3m}}_{\chi',\upsilon^m}$ is the kernel of Proposition~\ref{pte:carrier}; the lemma gives
$c_3\gamma\,\mathcal L_{\vartheta,i}(g)=\int\varphi^\circ(x)\theta^{S,W'}_{\nu^\flat\upsilon^m,\upsilon^m}(x,(g,1,1,h^\flat_i);
[P_i,\psi^{\mathrm{inv}}_{\Lambda_i}])\,dx$. By Lemma~\ref{pte:translates} ($\nu^\flat\upsilon^m$ has trivial tame part,
$\varphi^\circ$ is right $K_S$-invariant and rapidly decreasing) the right side is $\mathsf N\,F_i(g)$.
\end{proof}

\subsubsection{The left side as the theta series of the first factors}\label{pte:ss-left}

Let $w_i\in\pcMat_{(3m+N'')\times N''}(\auO_E/\pi_M^2)$ be $\Lambda_i$ stacked over $\pi_M1_{N''}$, and $e$ the
inclusion of the $r$ indices of $W$ into those of $V''$. Let $H_{\vartheta,i}\subseteq\auGL_{N''}(\pcR_{\pcm^{(2)}_\vartheta})$ be
the group of the $A$ whose columns off $e$ are those of $1$ and whose reduction $\bar A$ modulo $\pi_M^2$
satisfies $w_i\bar A=w_i$, and $J'_{\vartheta,i}$ the group of the integral points of $\auU(V'')(\auA_f)$ whose reduction modulo
$\pcm^{(2)}_\vartheta$ lies in $H_{\vartheta,i}$. Explicitly, at $w\in\mathsf M$: $k\equiv1\bmod\pi_M$, $\Lambda_ik\equiv\Lambda_i$ and the
last $N''-r$ columns of $k$ are those of $1$ modulo $\pi_M^2$; at $w\in Q_\vartheta$: the last $N''-r$ columns of
$k$ are those of $1$ modulo $w$. It equals $J_{\vartheta,i}$ of \S\ref{po:ss-glue} (with $s_i=\Lambda_i$),
contains the principal congruence subgroup of level $\pcm^{(2)}_\vartheta$, and is free (Lemma~\ref{au:free}).

\begin{lemma}[The column datum is the slice datum]\label{pte:column}
For $\sigma\in\pcMat_{N''\times r}(\pcR_{\pcm^{(2)}_\vartheta})$ let $\mathrm{col}(\sigma)=\chi'_\vartheta(\det\sigma')$ if $w_i\bar\sigma=
w_i\iota''_0$ and $0$ otherwise ($\sigma'$ the rows $e$). Then $\mathrm{col}$, read as a function of integral
matrices at $\mathsf M\cup Q_\vartheta$ (integral in the coordinates $T=a_w^{-1}x$ of (N5) used in
\S\ref{po:ss-glue}, which are the coordinates of Definition~\ref{pgx:conv}, Lemma~\ref{pgx:dict}(c)), is the finite part at $\mathsf M\cup Q_\vartheta$ of $\Psi''_{\vartheta,\Lambda_i}$ of
\S\ref{po:ss-glue}.
\end{lemma}
\begin{proof}
$w_i\bar\sigma=w_i\iota''_0$ means $\Lambda_i(\sigma-\iota''_0)\equiv0\bmod\pi_M^2$ and $\pi_M(\sigma-\iota''_0)\equiv0\bmod\pi_M^2$,
i.e.\ $\sigma\equiv\iota''_0\bmod\pi_M$ and $\Lambda_i(\sigma-\iota''_0)\equiv0\bmod\pi_M^2$, i.e.\ $\sigma\in\iota''_0+\pi_ML''_{\Lambda_i}$
(definition of $L''_s$). On this set $\sigma'\equiv1\bmod\pi_M$, so the factor $\omega^j(\det\sigma')$ at $\mathsf M$ is
$1$. At $Q_\vartheta$ there is no condition ($w_i$ lives modulo $\pi_M^2$) and the factor is $\psi_w(\det\sigma')$, with
$\psi_w(0)=0$.
\end{proof}

\begin{lemma}[Restriction to the row stabiliser]\label{pte:rowstab}
Let $F_0$ be a function on $\pcMat_{N''\times3m}(\pcR_{\pcm^{(2)}_\vartheta})$ with $F_0(Ay)=F_0(y)$ for $A\in H_{\vartheta,i}$,
$\Phi(x'\mid x'')=\chi'_\vartheta(\det x'')F_0(x''^{-1}x')$ and $\Phi_0(x'\mid x'')=\mathrm{col}(x''e)\,F_0(x')$ if the
columns of $x''$ off $e$ are those of $1$, $0$ otherwise. Let $\theta[\cdot](h)$ be a kernel of $(\mathsf W,V'')$
whose polynomial is fixed by $\auU(V'')_\infty$ and whose character on $\auU(V'')$ is tame with residue character
$\chi'_\vartheta$. Then $\sum_{h\bmod K_{V''}}\theta[\Phi](h)=\sum_{h\bmod J'_{\vartheta,i}}\theta[\Phi_0](h)$ (sums over
classes weighted by inverse stabiliser orders).
\end{lemma}
\begin{proof}
For $A\in\auGL_{N''}(\pcR)$ put $(\mathrm{tr}_A\Psi)(x)=\chi'_\vartheta(\det A)\Psi(A^{-1}x)$; by Lemma~\ref{pc:level},
$\theta[\Psi](hk)=\theta[\mathrm{tr}_{\bar k}\Psi](h)$ for $k\in K_{V''}$. One checks $\mathrm{tr}_A\Phi=\Phi$, and
$\Phi_0(x)=\Phi(x)$ if $x''\in H_{\vartheta,i}$, $0$ otherwise (for $x''\in H_{\vartheta,i}$: $x''e$ satisfies the condition of
$\mathrm{col}$, $\det x''=\det(x''e)'$ because the other columns are those of $1$, and $F_0(x''^{-1}x')=F_0(x')$; conversely
if $\mathrm{col}(x''e)\ne0$ and the other columns are those of $1$, then $x''\in H_{\vartheta,i}$). Hence $\Phi=\sum_{A\in
\auGL/H_{\vartheta,i}}\mathrm{tr}_A\Phi_0$ (an invertible $x''$ lies in exactly one coset $AH_{\vartheta,i}$; for non-invertible $x''$
all terms vanish). So $\theta[\Phi](h)=\sum_A\theta[\Phi_0](h\tilde A)$ for integral lifts $\tilde A$, and the function
$\theta[\Phi_0]$ is right $J'_{\vartheta,i}$-invariant. Summing over the classes modulo $K_{V''}$, the classes modulo
$J'_{\vartheta,i}$ are the $h\tilde A$, each counted with the correct weight (the tower relation of class sums:
$\sum_{h\bmod K}\sum_A f_0(h\tilde A)/\#\mathrm{Stab}_K(h)=\sum_{h\bmod J'}f_0(h)/\#\mathrm{Stab}_{J'}(h)$ because $K/J'\cong
\auGL/H_{\vartheta,i}$ and the orbit--stabiliser relation).
\end{proof}

\begin{definition}[First factors]\label{pte:def-ff}
By Lemma~\ref{pte:rowstab} with $F_0=\psi^{\mathrm{inv}}_{\Lambda_i}$ (invariant under $H_{\vartheta,i}$: for $A\in H_{\vartheta,i}$,
$\mathrm{tr}(\Lambda_iAy\gamma)=\mathrm{tr}(\Lambda_iy\gamma)$ since $\Lambda_i\bar A=\Lambda_i$), the inner sum of
$\mathcal L_{\vartheta,i}$ is a class sum modulo $J'_{\vartheta,i}$ of the kernel with datum $\Phi_0=\mathrm{col}(x''e)\cdot
\plone[x''\text{ has the columns of }1\text{ off }e]\cdot\psi^{\mathrm{inv}}_{\Lambda_i}(x')$. The first factors
$\varphi''_{\vartheta,i}$ are those of Definition~\ref{ptf:def-ff}: they come from the Laplace expansion
$\pteSD(\eta_i)=\sum_\lambda\pop''_{u_\lambda}R_\lambda$ (Lemma~\ref{ptf:laplace}), and their factorisation property and
transformation laws are Proposition~\ref{ptf:ff}.
\end{definition}

\begin{proposition}[The left side]\label{pte:leftside}
With the level $J'_{\vartheta,i}$ of \S\ref{po:ss-glue} and the first factors of Definition~\ref{pte:def-ff},
$\mathcal L_{\vartheta,i}=\Xi_{\vartheta,i}$ on $\auU(W)(\auA)$.
\end{proposition}
\begin{proof}
Proposition~\ref{ptf:leftside}; the character $\chi''_\vartheta$ there is that of \S\ref{po:ss-glue} by Lemma~\ref{ptf:chi}.
\end{proof}

\subsubsection{Proof of Proposition~\ref{po:exchange}}\label{pte:ss-proof}

Put $c'_\vartheta=c_3(\vartheta)\gamma(\vartheta)/\mathsf N$ and $C_{\mathrm x}=C_3C_{12}/\mathsf N^2\ge0$. By
Propositions~\ref{pte:third} and~\ref{pte:leftside}, $c'_\vartheta\Xi_{\vartheta,i}=F_i$ for all $i$, with $F_i$ of
\eqref{pte:eq-F} independent of $\vartheta$. By Propositions~\ref{pte:carrier} (for $c=\chi'=\chi_0^j\vartheta$, modulus
$\pcm^{(2)}_\vartheta$, whose primes are $\mathsf S_\vartheta$) and~\ref{pte:third},
\[
|c'_\vartheta|^2\,L_*(\pcm^{(2)}_\vartheta,(\nu_*\vartheta)^{-1},z_*)\,L_*(\pcm^{(2)}_\vartheta,\chi'^{-1}\upsilon^m,m)\,L_*(\pcm^{(2)}_\vartheta,\chi'\upsilon^{m/2},m/2)
=C_{\mathrm x}.
\]
Now $\chi'^{-1}\upsilon^m=(\nu_+\vartheta)^{-1}$ ($\nu_+=\upsilon^{-m}\chi_0^j$, $z_+=m$) and $\chi'\upsilon^{m/2}=(\nu_-\vartheta^{-1})^{-1}$
($\nu_-=\upsilon^{-m/2}\chi_0^{-j}$, $z_-=m/2$), and $L_*(\pcm^{(2)}_\vartheta,\nu,z)=(\auNm\mathsf S_\vartheta)^{z^2}\prod_{z'=1}^z|L^{\mathsf S_\vartheta}(z',\nu)|^2$
(\eqref{pcx:eq-Lstar}; the primes of $\pcm^{(2)}_\vartheta$ are those of $\mathsf S_\vartheta$, each counted once). So the product of the three
$L_*$ is $H_j(\vartheta)$ of \eqref{po:eq-H}, and $|c'_\vartheta|^2H_j(\vartheta)=C_{\mathrm x}$.

Integrability: $F_i=c'_1\Xi_{1,i}$ and $\Xi_{\vartheta,i}$ are finite sums of theta kernels of $(W,V'')$, which have moderate
growth on $\auU(W)(\auA)$; so is $g\mapsto K^\circ(g,h^\circ)$; and $y\in\poA(W_j)$ is rapidly decreasing (cuspidal,
\S\ref{ss:au-spaces}). Hence $yF_iK^\circ(\cdot,h^\circ)$ and $y\Xi_{\vartheta,i}K^\circ(\cdot,h^\circ)$ are integrable on $[\auU(W)]$.
\qed

\begin{remark}\label{pte:rem-final}
(i) Corollary~\ref{pte:explicit} alone gives only $C_{\mathrm x}\ge0$; $C_{\mathrm x}>0$ is proved in Proposition~\ref{po:glue}(G2),
from $c'_\vartheta\ne0$ (injectivity of $\Theta^\circ$) and $H_j(\vartheta)>0$ (Lemma~\ref{po:Hpos}). (ii) The level $J_{\vartheta,i}$ of \S\ref{po:ss-glue} is the group $J'_{\vartheta,i}$ of \S\ref{pte:ss-left}; every use of it
in \S\ref{po:ss-glue} and Proposition~\ref{po:glue} (Lemma~\ref{po:lattice}(d), freeness, representative independence,
Lemma~\ref{po:classes}) holds for this smaller group. The character $\chi''_\vartheta$ on $J'_{\vartheta,i}$ is
$k\mapsto\chi'_\vartheta(\det\bar k)$ (Proposition~\ref{ptf:ff}(b), Lemma~\ref{ptf:chi}). (iii) The
matrices $s_i$ of \S\ref{po:ss-glue} are the $\Lambda_i$.
\end{remark}


%% ===== proofs/C-threeex2.tex =====
% proofs/C-threeex2.tex
% Complete proofs of the steps of proofs/C-threeex.tex (labels pte:) that were only sketched there:
% the negation lemma, the archimedean action formulas of the explicit exchange, the first factors and
% the character bookkeeping on V''. To be \input after C-threeex.tex. Labels start with ptf:.
% Only amsmath/amssymb/amsthm.
\providecommand{\auO}{\mathcal O}
\providecommand{\auU}{\mathrm U}
\providecommand{\auGL}{\mathrm{GL}}
\providecommand{\auHom}{\operatorname{Hom}}
\providecommand{\auA}{\mathbb A}
\providecommand{\auPl}{\Phi}
\providecommand{\aubar}[1]{\overline{#1}}
\providecommand{\pcR}{R}
\providecommand{\pcm}{\mathfrak m}
\providecommand{\pcMat}{\mathrm M}
\providecommand{\psvac}{{\phi^{\mathrm{vac}}}}
\providecommand{\psrho}{\varrho}
\providecommand{\psHS}{\mathcal H}
\providecommand{\psZ}{\mathsf Z}
\providecommand{\poT}{\mathbb T}
\providecommand{\poA}{\mathcal A}
\providecommand{\pop}{\mathsf p}
\providecommand{\pog}{\mathsf g}
\providecommand{\plone}{\mathbf 1}
\providecommand{\pteSD}{\mathrm{SD}}
\providecommand{\pteMP}{\mathrm{MP}}
\providecommand{\pteCS}{\mathop{\textstyle\sum^{\mathrm{cl}}}}

\subsection{Three exchanges: conjugation, the Fock action and the first factors}\label{ptf:sec}

This subsection gives complete proofs of the steps of \S\ref{pte:sec} that were stated there in
abbreviated form. Each result below is self-contained; \S\ref{pte:sec} cites it at the place
indicated.

\subsubsection{Negation and complex conjugation of theta kernels}\label{ptf:ss-conj}

Let $S,T$ be spaces (\S\ref{pc:ss-data}) of ranks $n_S,n_T$ and $W=\auHom(S,T)$, a hermitian space of
rank $N=n_Sn_T$ with the form of \S\ref{ps:ss-weil}. Let $S^-$ be the space with the same underlying
$E$-vector space and the same frame as $S$, and with the form $-\langle\,,\rangle_S$; its lattice is
the same unimodular lattice, its signature is the opposite one, and $\auU(S^-)=\auU(S)$ as groups (an
isometry of a form is an isometry of its negative). Put $W^-=\auHom(S^-,T)$. The adjoint of
$w\colon S\to T$ for $-\langle\,,\rangle_S$ is the negative of the adjoint for $\langle\,,\rangle_S$,
so \emph{$W^-$ is the $E$-vector space $W$ with the form $-\langle\,,\rangle_W$}; the frames of $W$ and
$W^-$ (the matrix units in the frames of $S$ and $T$) coincide, so the two spaces have the same
coordinates at every place, and $\auU(W)=\auU(W^-)$. The pair $(g,h)$ acts on both by $w\mapsto
hwg^{-1}$, and
\begin{equation}\label{ptf:eq-detonhom}
\det\bigl(w\mapsto hwg^{-1}\bigr)=\det(h)^{n_S}\det(g)^{-n_T}
\end{equation}
(on $\pcMat_{n_T\times n_S}$, left multiplication by $h$ acts on each of the $n_S$ columns and right
multiplication by $g^{-1}$ on each of the $n_T$ rows).

\emph{Heisenberg representations.} The bicharacter of \S\ref{ps:ss-weil} of $W^-$ is
$\beta^-(x,y)=\psi(\frac12\jmath\langle y,x\rangle_{W^-})=\beta(x,y)^{-1}=\aubar{\beta(x,y)}$, and
$[x,y]^-=\aubar{[x,y]}$. If $(\psrho,\psHS)$ is a Heisenberg representation of $W$ (at one place or
adelically), then the complex conjugate space $\aubar{\psHS}$ (the same set, with $\lambda\cdot f=
\aubar\lambda f$) with the same operators $\psrho(x)$ satisfies $\psrho(x)\psrho(y)=\aubar{\beta(x,y)}
\psrho(x+y)$, so it is a Heisenberg representation of $W^-$ with the same central character
$\psi$ on the centre written as $W$ plus scalars (the centre acts through $\psi$ composed with the
real part of the form, which changes sign together with the conjugation). By the uniqueness of the
Heisenberg representation (Stone--von Neumann; at a finite place the smooth version) there is an
\emph{antilinear} bijection $A_v\colon\psHS_v(W)\to\psHS_v(W^-)$ with
\begin{equation}\label{ptf:eq-Arho}
A_v\,\psrho(x)=\psrho^-(x)\,A_v\qquad(x\in W_v),
\end{equation}
unique up to a non-zero scalar (two such maps differ by a linear self-intertwiner of an irreducible
representation, i.e.\ a scalar, by Schur's lemma). We fix the scalars as follows.

\begin{itemize}
\item[(A$_\infty$)] At $\tau\in\auPl$ the vacuum of $W_\tau$ is the unique line killed by the
annihilation operators, which are the operators $d\psrho(z)$ for $z$ in the antiholomorphic part of
$W_\tau\otimes_\R\C$ for the complex structure $J_\tau$ of the Fock model; for $W^-_\tau$ the complex
structure is $-J_\tau$ (the positive and negative coordinates are exchanged). By
\eqref{ptf:eq-Arho} $A_\tau\,d\psrho(z)=d\psrho^-(\aubar z)A_\tau$, and complex conjugation exchanges the
$\pm i$-eigenspaces of $J_\tau$, which are the $\mp i$-eigenspaces of $-J_\tau$. So $A_\tau$ carries the
vacuum line of $W_\tau$ to that of $W^-_\tau$, and we normalise $A_\tau\psvac_\tau=\psvac^-_\tau$ (the real
vacuum $1$ of the Fock model in both cases). The Fock vector of a polynomial $P$ in the holomorphic
coordinates $\zeta$ of $W_\tau$ (\eqref{po:eq-poly}; the coordinate of an entry of positive product sign is the
entry, that of an entry of negative product sign its conjugate) is $P(c_1,\dots)\psvac_\tau$, where the
$c_i$ are the creation operators attached to the coordinates (complex linear combinations of the
$d\psrho(x)$). Applying $A_\tau$ conjugates the coefficients of these combinations; the creation
operator of the coordinate $\zeta_i$ of $W_\tau$ goes to the creation operator of $\aubar{\zeta_i}$,
which is the holomorphic coordinate of the same entry in $W^-_\tau$, since all product signs have
changed. Hence
\[
A_\tau\bigl(\text{Fock vector of }P\bigr)=\text{Fock vector of }\aubar P\ \text{for }W^-_\tau,
\]
$\aubar P$ the polynomial with conjugated coefficients, in the holomorphic coordinates of $W^-_\tau$.
\item[(A$_\pcm$)] At a place $v=w\aubar w$ below a prime $w\mid\pcm$ (split, \S\ref{pc:ss-data}) both
$\psHS_v(W)$ and $\psHS_v(W^-)$ are realised by (W5) on Schwartz--Bruhat functions on
$\pcMat_{n_T\times n_S}(E_w)$, with $\psrho(a_w,0)$ the translation by $a_w$ and $\psrho(0,a_{\aubar w})$
the multiplication by $x\mapsto[x,a_{\aubar w}]$. Put $A_v\Psi=\aubar\Psi$. Translation commutes with
conjugation and $\aubar{[x,a]}=[x,a]^-$, so \eqref{ptf:eq-Arho} holds. In particular $A_v$ carries
$\Psi^\natural$ to $(\aubar\Psi)^\natural$ and the normalised vacuum (the normalised indicator of the
integral matrices, which is real) to itself.
\item[(A$_f$)] At the other finite places $v$: the lattice $L_v=\alpha\auHom(L_S,L_T)_v$ is the same
for $W$ and $W^-$ and self-dual for both ($[\,,]^-=\aubar{[\,,]}$). By \eqref{ptf:eq-Arho} $A_v$ carries
the line on which $\psrho(l)$, $l\in L_v$, acts by $\epsilon_L(l)$ to the line on which it acts by
$\aubar{\epsilon_L(l)}$, and $\aubar{\epsilon_L(l)}=\psi_f(+\frac12\omega\jmath\langle l,l\rangle)=\epsilon^-_L(l)$ is the
character of the vacuum of $W^-$ (\S\ref{ps:ss-weil}, formula for $\epsilon_L$ with the form negated).
So $A_v\psvac_v$ is a non-zero multiple of $\psvac^-_v$; both are unit vectors for the local inner
products ($A_v$ is isometric up to a scalar, since $\langle A_vf,A_vf'\rangle$ is a conjugate of an
invariant hermitian form, hence a positive multiple of $\aubar{\langle f,f'\rangle}$), and we normalise
$A_v\psvac_v=\psvac^-_v$.
\end{itemize}
Let $A=\bigotimes_vA_v\colon\psHS(W)\to\psHS(W^-)$ (well defined on the restricted tensor product since
$A_v\psvac_v=\psvac^-_v$ for all $v\notin S(\pcm)\cup\auPl$; on the places of $\pcm$ the vacuum is also
preserved, so $A$ does not depend on $\pcm$). By (A$_\infty$), (A$_\pcm$), (A$_f$) and the definition of
the datum in \S\ref{pc:ss-data},
\begin{equation}\label{ptf:eq-Adatum}
A\,[P,\Psi]=[\aubar P,\aubar\Psi]\qquad\text{(datum of }W^-\text{ on the right)}.
\end{equation}

\begin{lemma}[Theta functionals]\label{ptf:theta}
There is $\kappa=\kappa(S,T)\in\C$ with $|\kappa|=1$ such that $\theta_{W^-}(Af)=\kappa\,\aubar{\theta_W(f)}$
for all $f\in\psHS(W)$ (for $f$ in the domain of $\theta$ of (W0)).
\end{lemma}
\begin{proof}
The rational points of $W$ and $W^-$ are the same set $W(E)$. The functional $f'\mapsto
\aubar{\theta_W(A^{-1}f')}$ on $\psHS(W^-)$ is linear (two antilinear maps) and, by \eqref{ptf:eq-Arho},
invariant under $\psrho^-(x)$ for $x\in W(E)$; by the uniqueness in (W1) it equals $\kappa^{-1}\theta_{W^-}$
for some $\kappa\ne0$ (non-zero since $A$ is bijective and $\theta_W\neq0$). For the modulus: the
functionals are normalised (\S\ref{ps:ss-weil}, \S\ref{pc:ss-data}) by the condition that $x\mapsto
\theta(\psrho(x)\psvac)$, $\psvac$ the global vacuum, have mean square $1$ over $W(E)\backslash W(\auA)$.
By the normalisations above $A\psvac=\psvac^-$, so $\theta_{W^-}(\psrho^-(x)\psvac^-)=\theta_{W^-}(A\psrho(x)
\psvac)=\kappa\aubar{\theta_W(\psrho(x)\psvac)}$; taking mean squares, $1=|\kappa|^2$.
\end{proof}

\begin{lemma}[Weil representations of $W$ and $W^-$]\label{ptf:weil}
For $\gamma\in\auU(W)(\auA)=\auU(W^-)(\auA)$,
\begin{equation}\label{ptf:eq-Aomega}
A\,\omega_W(\gamma)\,A^{-1}=\upsilon(\det\gamma)\,\omega_{W^-}(\gamma).
\end{equation}
\end{lemma}
\begin{proof}
\emph{Step 1: a character.} By (W1) for $W$ and $W^-$ and \eqref{ptf:eq-Arho}, $A\omega_W(\gamma)A^{-1}$
and $\omega_{W^-}(\gamma)$ both conjugate $\psrho^-(x)$ into $\psrho^-(\gamma x)$; so
$\omega_{W^-}(\gamma)^{-1}A\omega_W(\gamma)A^{-1}$ commutes with the irreducible Heisenberg representation
and is a scalar $\lambda(\gamma)$. As both $\gamma\mapsto A\omega_W(\gamma)A^{-1}$ and $\omega_{W^-}$ are
representations (genuine, by the splitting), $\lambda$ is a character, and it is unitary ($A$ is
isometric up to a constant). It is a product $\prod_v\lambda_v$ of characters of the local groups, by
the same argument place by place.

\emph{Step 2: $\lambda=(\upsilon\circ\det)^s$ for a universal sign $s$.} The splitting of $\omega_W$ is
Kudla's splitting attached to the character $\upsilon$ (\S\ref{ps:ss-weil}, \cite[Theorem~3.1]{KudlaSplitting}); we
use three of its standard properties: (K1) it is defined by scalar functions on $\auU(W)$ built from
$\psi$ (Weil indices and the character $\upsilon$ evaluated on Rao's invariants), so that replacing
$(\psi,\upsilon)$ by $(\aubar\psi,\aubar\upsilon)$ conjugates these scalars; (K2) the Heisenberg and Weil
representations of $(W,\aubar\psi)$ are those of $(W^-,\psi)$; (K3) for a unitary character $\mu$ of
$\auU(1)$, replacing $\upsilon$ by $\upsilon\cdot(\mu\circ(x\mapsto x/\aubar x))$ multiplies the splitting by
$\mu(\det\gamma)^{s}$ with a sign $s=\pm1$ that does not depend on $W$ (\cite[p.~385 and Lemma~3.5]{KudlaSplitting}, \cite[(1.7)--(1.8)]{HKS}). By (K1) and (K2), $A\omega_W(\gamma)A^{-1}$ is the Weil representation of $W^-$ defined
with $\aubar\upsilon=\upsilon^{-1}$. Since $\upsilon$ restricted to $\auA_{E^+}^\times$ is quadratic,
$\upsilon(\aubar x)=\upsilon(x)^{-1}$, hence $\upsilon(x/\aubar x)=\upsilon(x)^2$ and $\upsilon^{-1}=\upsilon\cdot
(\mu\circ(x\mapsto x/\aubar x))$ with $\mu=\upsilon|_{\auU(1)}^{-1}$. By (K3), $\lambda=(\upsilon\circ\det)^{-s}$.

\emph{Step 3: the sign, from (W2).} It remains to see $-s=1$; as $s$ is universal it suffices to check
$\lambda_v(p)=\upsilon_v(\det p)$ for one place $v$, one $W$ with an isotropic line $I\subset W_v$ and one
$p$ with $\upsilon_v(\det p)\neq\upsilon_v(\det p)^{-1}$. Let $I$ be an isotropic subspace of $W_v$, $P_I$ its
stabiliser, and let $p\in P_I$ act on $I$ by $a\in\auGL(I)$ and trivially on $I^\perp/I$; then $p$ acts on
$W_v/I^\perp\cong\aubar I^{\,*}$ by the inverse conjugate transpose of $a$, so $\det p=\det a/\aubar{\det a}$
and $\upsilon_v(\det p)=\upsilon_v(\det a)^2$. The subspace $I$ and the group $P_I$ are the same for
$W_v^-$, $I^\perp/I$ is $W_{I,v}$ with its form negated, and the zero terms $\psZ_I$, $\psZ^-_I$ of (W2)
are characterised up to a scalar as the linear maps to $\psHS(W_I)$ resp.\ $\psHS(W_I^-)$ that intertwine
the Heisenberg groups of $I^\perp$ and kill $\psrho(i)-1$, $i\in I_v$ (in the mixed model they are the
evaluation at $0$ of functions on $I^*$, and the coinvariants of $I_v$ form one copy of
$\psHS(W_I)$). Let $A_I\colon\psHS_v(W_I)\to\psHS_v(W_I^-)$ be the antilinear intertwiner as above. Then
$A_I^{-1}\psZ^-_IA$ is linear with the same characterising properties as $\psZ_I$ (by
\eqref{ptf:eq-Arho} for $W$ and for $W_I$), so $\psZ^-_IA=c\,A_I\psZ_I$ with $c\neq0$. By (W2), $\psZ_I
\omega(p)=\chi_I(p)\psZ_I$ and $\psZ^-_I\omega^-(p)=\chi_I(p)\psZ^-_I$, with the same
$\chi_I(p)=\upsilon_v(\det a)^{-1}|\det a|^{1/2}$ (the formula does not involve the form). Hence
\[
\lambda_v(p)\chi_I(p)\,\psZ^-_IA=\psZ^-_I\,\lambda_v(p)\omega^-(p)A=\psZ^-_IA\,\omega(p)=cA_I\psZ_I\omega(p)
=cA_I\chi_I(p)\psZ_I=\aubar{\chi_I(p)}\,\psZ^-_IA,
\]
so $\lambda_v(p)=\aubar{\chi_I(p)}/\chi_I(p)=\upsilon_v(\det a)^2=\upsilon_v(\det p)$. Take $v=\tau$ infinite,
$W_\tau$ indefinite and $a=\mathrm{diag}(z,1,\dots)$ with $z\in\C^\times$, $z^4\neq|z|^4$: as
$\upsilon_\tau(z)=z/|z|$, $\upsilon_\tau(\det p)^2=(z/|z|)^4\ne1$. So $-s=1$.
\end{proof}

\begin{remark}\label{ptf:rem-weil}
Step 3 alone proves \eqref{ptf:eq-Aomega} on every parabolic subgroup $P_I(E^+_v)$ (for the elements
acting trivially on $I^\perp/I$, and by induction on the Witt index for all of $P_I$); Step 2 is used
only to extend it to the whole group, in particular at places where $W$ is anisotropic. The sign check
is a consistency check of the normalisation (W2) of the splitting used in the paper with that of
Kudla.
\end{remark}

\begin{lemma}[Negation and complex conjugation]\label{ptf:conj}
Let $S,T$ be spaces of ranks $n_S,n_T$ and $S^-$ as above. There is $\epsilon=\epsilon(S,T)\in\C$,
$|\epsilon|=1$, such that for all automorphic characters $\alpha,\beta$ of $\auU(1)$, all moduli $\pcm$, all
functions $\Psi$ on $\pcMat_{n_T\times n_S}(\pcR_\pcm)$, all polynomials $P$ (more generally Fock vectors
with finite moments) and all $g\in\auU(S)(\auA)=\auU(S^-)(\auA)$, $h\in\auU(T)(\auA)$,
\begin{equation}\label{ptf:eq-conj}
\aubar{\theta^{S,T}_{\alpha,\beta}\bigl(g,h;[P,\Psi]\bigr)}=\epsilon\,\theta^{S^-,T}_{\alpha^{-1}\upsilon^{n_T},\,
\beta^{-1}\upsilon^{n_S}}\bigl(g,h;[\aubar P,\aubar\Psi]\bigr).
\end{equation}
Here the first index is the character of the source group (\S\ref{pcx:ss-statement}): the source
character is shifted by $\upsilon^{\mathrm{rk}\,T}$ and the target character by $\upsilon^{\mathrm{rk}\,S}$.
If instead $S^-$ is given the frame $(\varepsilon e_i)$ at the finite places, $\varepsilon\in\widehat\auO_E^\times$
with $\varepsilon\aubar\varepsilon=-1$, then $\aubar\Psi$ is replaced by $x\mapsto\aubar{\Psi(u^{-1}x)}$, $u$ the
reduction of $\varepsilon$.
\end{lemma}
\begin{proof}
Put $f=[P,\Psi]$ and $\gamma=(w\mapsto hwg^{-1})$. By Lemma~\ref{ptf:theta}, then
Lemma~\ref{ptf:weil} and \eqref{ptf:eq-detonhom}, and \eqref{ptf:eq-Adatum},
\begin{align*}
\aubar{\theta^{S,T}_{\alpha,\beta}(g,h;f)}&=\aubar{\alpha(\det g)^{-1}\beta(\det h)}\;
\aubar{\theta_W(\omega_W(\gamma)f)}=\alpha(\det g)\beta(\det h)^{-1}\kappa^{-1}\theta_{W^-}\bigl(A\omega_W(\gamma)
f\bigr)\\
&=\alpha(\det g)\beta(\det h)^{-1}\kappa^{-1}\upsilon(\det h)^{n_S}\upsilon(\det g)^{-n_T}\,
\theta_{W^-}\bigl(\omega_{W^-}(\gamma)[\aubar P,\aubar\Psi]\bigr)\\
&=\kappa^{-1}\bigl(\alpha^{-1}\upsilon^{n_T}\bigr)(\det g)^{-1}\bigl(\beta^{-1}\upsilon^{n_S}\bigr)(\det h)\,
\theta_{W^-}\bigl(\omega_{W^-}(\gamma)[\aubar P,\aubar\Psi]\bigr),
\end{align*}
which is \eqref{ptf:eq-conj} with $\epsilon=\kappa^{-1}=\aubar\kappa$; it depends only on $S,T$. (The
characters take values of modulus one, so conjugation inverts them.) For the frame $(\varepsilon e_i)$:
the coordinates of a vector of $W^-$ at a finite place are then $\varepsilon$ times those for the old
frame, and (A$_\pcm$) becomes $A_v\Psi(x)=\aubar{\Psi(\varepsilon^{-1}x)}$, which gives the stated datum;
the other steps are unchanged (the lattice $\varepsilon L_v=L_v$).
\end{proof}

\begin{remark}\label{ptf:rem-conj}
(i) Lemma~\ref{pte:conj} is the case of Lemma~\ref{ptf:conj} used in Proposition~\ref{pte:carrier}, with
$S=T=I_m$ (equal ranks). (ii) With the frames of \S\ref{pc:ss-data} (rational bases), $S^-$ has the same frame
as $S$, no unit $u$ appears, and this is the version used in Proposition~\ref{pte:carrier}.
\end{remark}

\subsubsection{The archimedean action on polynomial vectors}\label{ptf:ss-fock}

This subsection fixes the transpose conventions needed in Step~2 of the proof of
Corollary~\ref{pte:explicit} and proves the formulas used there.

\emph{Conventions.} Let $S,T$ be spaces with sign vectors $s\in\{\pm1\}^{n_S}$, $s'\in\{\pm1\}^{n_T}$ (the
same at all $\tau$), and $n_S^\pm$, $n_T^\pm$ the numbers of signs $\pm1$. At $\tau\in\auPl$ a vector of
$\auHom(S,T)\otimes_{E^+,\tau}\R$ is a matrix $w\in\pcMat_{n_T\times n_S}(\C)$ (rows indexed by $T$,
columns by $S$, in the frames), and $(g,h)$ acts by $w\mapsto hwg^{-1}$. The entry $(a,b)$ has
\emph{product sign} $s'_as_b$. The \emph{Fock variable} $X_{ab}$ of the entry $(a,b)$ is $\aubar{w_{ab}}$ if
$s'_as_b=1$ and $w_{ab}$ if $s'_as_b=-1$ (\S\ref{ps:ss-weil}: the variable of the Fock space is the
conjugate of the coordinate; \eqref{po:eq-poly}). For a polynomial $P$ in the $X_{ab}$ (all $\tau$
together) its Fock vector is the function $w\mapsto P(X(w))\,G(w)$ on $\prod_\tau\pcMat_{n_T\times
n_S}(\C)$, $G$ the real Gaussian (the vacuum); we identify $P$ with it. For a block-diagonal $g\in\auU(S)_
\tau$, $g=\mathrm{diag}(g^+,g^-)$ along the signs, $g^\pm\in\auU(n_S^\pm)$, and similarly $h$. A character
$\alpha$ of $\auU(1)$ of weight $(\mu_\tau)$ has archimedean components $\alpha_\tau(z)=z^{\mu_\tau}$ on
$\auU(1)(\R)=\{|z|=1\}$. We write $\omega_\infty$ for the archimedean part of $\omega_{\alpha,\beta}$ (with the
archimedean components of the characters).

\emph{Normalisation of the compact action.} For $\gamma$ in the maximal compact subgroup of
$\auU(\auHom(S,T))_\tau$ (the elements preserving the decomposition by product signs) the Weil
representation acts on Fock vectors, read as functions of $w$, by
\begin{equation}\label{ptf:eq-compact}
(\omega(\gamma)F)(w)=\Delta(\gamma)\,F(\gamma^{-1}w),\qquad\Delta(\gamma)=\aubar{\det(\gamma^-)},
\end{equation}
$\gamma^-$ the block of $\gamma$ on the entries of negative product sign. Indeed $\gamma$ preserves
the Gaussian, so $\omega(\gamma)$ and $F\mapsto F\circ\gamma^{-1}$ both conjugate $\psrho(x)$ into
$\psrho(\gamma x)$ and differ by a scalar character $\Delta$ of the compact group; that $\Delta(\gamma)=
\aubar{\det\gamma^-}=\det(\gamma^-)^{-1}$ is the archimedean normalisation of the splitting defined
with $\upsilon$, $\upsilon_\tau(z)=z/|z|$ (\S\ref{ps:ss-weil}): by Lemma~\ref{pwf:delta-red},
$\Delta(\gamma)=\det(\gamma^+)^{e_+}\det(\gamma^-)^{e_-}$ with universal $e_\pm$ and $e_++e_-=-1$, and
$e_+=0$, $e_-=-1$ by Proposition~\ref{pwf:eplus} (an explicit computation with Kudla's splitting on a
doubled line, which uses the signs of $\psi_\infty$ and $\jmath$ fixed in \S\ref{ps:ss-weil}). It is consistent with
Lemma~\ref{ptf:weil}: for $W^-$ the negative entries are the positive entries of $W$, so $\Delta^-(\gamma)=
\aubar{\det\gamma^+}$, and $\aubar{\Delta(\gamma)}=\det\gamma^-=\det\gamma\cdot\aubar{\det\gamma^+}=\upsilon_\tau(\det\gamma)
\Delta^-(\gamma)$, which is \eqref{ptf:eq-Aomega} at $\tau$ (with (A$_\infty$), $A_\tau F=\aubar F$ on
functions of $w$).

\begin{lemma}[Compact action on polynomials]\label{ptf:fock}
Let $\alpha,\beta$ have weights $(a_\tau)$, $(b_\tau)$, and let $g\in\auU(S)_\infty$, $h\in\auU(T)_\infty$ be
block diagonal at every $\tau$. For every polynomial $P$,
\[
\omega_\infty(g,h)P=c(g,h)\cdot P\circ\sigma_{g,h},\qquad c(g,h)=\prod_\tau\det(h^+_\tau)^{b_\tau-n_S^-}
\det(h^-_\tau)^{b_\tau-n_S^+}\det(g^+_\tau)^{n_T^--a_\tau}\det(g^-_\tau)^{n_T^+-a_\tau},
\]
where $\sigma_{g,h}$ is the linear substitution of the Fock variables which, writing $X_\tau=
\bigl(\begin{smallmatrix}X^{++}&X^{+-}\\X^{-+}&X^{--}\end{smallmatrix}\bigr)$ by the signs of the rows ($T$) and the
columns ($S$), is
\[
X^{++}\mapsto(h^+)^{\mathsf t}X^{++}(g^+)^{-\mathsf t},\quad X^{--}\mapsto(h^-)^{\mathsf t}X^{--}(g^-)^{-\mathsf t},
\quad X^{+-}\mapsto(h^+)^{-1}X^{+-}g^-,\quad X^{-+}\mapsto(h^-)^{-1}X^{-+}g^+
\]
at each $\tau$. Equivalently: on the variables of positive product sign, $X\mapsto h^{\mathsf t}Xg^{-\mathsf t}$; on
those of negative product sign, $X\mapsto h^{-1}Xg$.
\end{lemma}
\begin{proof}
Put $\gamma=(w\mapsto hwg^{-1})$, which preserves the product signs. By \eqref{ptf:eq-compact},
$(\omega_\infty(g,h)P)(w)=\alpha_\infty(\det g)^{-1}\beta_\infty(\det h)\Delta(\gamma)P(X(h^{-1}wg))G(w)$.
\emph{Substitution.} For an entry of negative product sign $X(h^{-1}wg)_{ab}=(h^{-1}wg)_{ab}$, and the
block structure gives $X^{+-}\mapsto(h^+)^{-1}X^{+-}g^-$, $X^{-+}\mapsto(h^-)^{-1}X^{-+}g^+$. For an entry of positive
product sign $X(h^{-1}wg)_{ab}=\aubar{(h^{-1}wg)_{ab}}=(\aubar{h^{-1}}\,\aubar w\,\aubar g)_{ab}$, and for unitary
blocks $\aubar{h^{-1}}=\aubar{h^*}=h^{\mathsf t}$, $\aubar g=g^{-\mathsf t}$; this gives the two diagonal blocks.
\emph{Scalar.} The negative block $\gamma^-$ consists of $X^{+-}\mapsto h^+X^{+-}(g^-)^{-1}$ on
$\pcMat_{n_T^+\times n_S^-}$ and $X^{-+}\mapsto h^-X^{-+}(g^+)^{-1}$ on $\pcMat_{n_T^-\times n_S^+}$ (in the coordinates
$w$), so by \eqref{ptf:eq-detonhom}
\[
\det\gamma^-=\det(h^+)^{n_S^-}\det(g^-)^{-n_T^+}\det(h^-)^{n_S^+}\det(g^+)^{-n_T^-},
\]
and $\Delta(\gamma)$ is its inverse (all determinants have modulus one). With $\beta_\infty(\det h)=\prod_\tau
\det(h^+_\tau)^{b_\tau}\det(h^-_\tau)^{b_\tau}$ and $\alpha_\infty(\det g)^{-1}=\prod_\tau\det(g^+_\tau)^{-a_\tau}\det(g^-_\tau)^{-a_\tau}$
the product is $c(g,h)$.
\end{proof}

\begin{corollary}[The formulas of Step~2 of Corollary~\ref{pte:explicit}]\label{ptf:SDlaw}
In the setting of \S\ref{pte:ss-conv} let the archimedean weight of $\alpha$ be $(l-\#\iota(\tau))_\tau$ and
that of $\beta$ be $0$ (Step~0 of the proof of Corollary~\ref{pte:explicit} allows this). For
$k\in\prod_\tau\auU(m)$, write also $k$ for $\mathrm{diag}(k,1_l)\in\auU(W)_\infty$ (positive block). Then:
\begin{enumerate}
\item[(a)] {\emergencystretch=6em on $\auHom(W,I_a)$ (characters $\beta$ on $\auU(W)$, $\alpha$ on $\auU(I_a)$), for $u\in\auU(I_a)_\infty$,
$\omega_\infty(k,u)\pteSD(\eta)=\allowbreak \prod_\tau\det(k_\tau)^{-\#\iota(\tau)}\pteSD(\eta k^{\mathsf t})$; equivalently
$\omega_\infty(k,u)\pteSD(\eta k^{-\mathsf t})=\prod_\tau\det(k_\tau)^{-\#\iota(\tau)}\pteSD(\eta)$. In particular
$\omega_\infty(1,u)\pteSD(\eta)=\pteSD(\eta)$.\par}
\item[(b)] on $\auHom(V',W)$ (characters $\alpha$ on $\auU(V')$, $\alpha\beta^{-1}$ on $\auU(W)$),
$\omega_\infty(1,k)\pteMP(\eta)=\prod_\tau\det(k_\tau)^{-\#\iota(\tau)}\pteMP(\eta k^{\mathsf t})$.
\item[(c)] for $|z|=1$ and $t_z,\Omega_z$ as in Step~3, $\Omega_z$ acts on polynomials in the variables
$x$ of $\auHom(W,I_a)$ by the substitution multiplying every variable outside $x^\flat$ by $z$ and fixing
$x^\flat$; and $\omega_\infty(1,u)$ acts on polynomials in $x^\flat$ alone by $q(x^\flat)\mapsto\prod_\tau
\det(u_\tau)^{-\#\iota(\tau)}q(u^{\mathsf t}x^\flat)$.
\end{enumerate}
Here $x=(x^+\mid x^-)$ are the Fock variables of $\auHom(W,I_a)$ (all of $I_a$ is positive, so $x^+$ has
positive and $x^-$ negative product sign), and $y$ those of $\auHom(V',W)$.
\end{corollary}
\begin{proof}
(a) Lemma~\ref{ptf:fock} with $(S,T)=(W,I_a)$, $g=\mathrm{diag}(k,1)$, $h=u$, source weight $0$, target
weight $l-\#\iota(\tau)$, $n_S^-=l$, $n_T^-=0$, $n_T^+=a$: $c=\prod_\tau\det(u_\tau)^{(l-\#\iota(\tau))-l}
\det(k_\tau)^{0}\det(1)^{a}=\prod_\tau\det(u_\tau)^{-\#\iota(\tau)}$, and the substitution is $x^+\mapsto u^{\mathsf t}x^+
k^{-\mathsf t}$, $x^-\mapsto u^{-1}x^-$. As $\pteSD$ involves only $x^+$, and
\[
\begin{pmatrix}\eta_{\tau j}\\u^{\mathsf t}_\tau x^+_\tau k_\tau^{-\mathsf t}\end{pmatrix}=\begin{pmatrix}1_p&0\\0&u_\tau^{\mathsf t}
\end{pmatrix}\begin{pmatrix}\eta_{\tau j}k_\tau^{\mathsf t}\\x^+_\tau\end{pmatrix}k_\tau^{-\mathsf t},
\]
each factor of $\pteSD(\eta)$ is multiplied by $\det(u_\tau)\det(k_\tau)^{-1}$ and $\eta$ becomes $\eta k^{\mathsf t}$;
the powers $\det(u_\tau)^{\#\iota(\tau)}$ cancel $c$. Replacing $\eta$ by $\eta k^{-\mathsf t}$ gives the
equivalent form, and $k=1$ the last assertion.
(b) Lemma~\ref{ptf:fock} with $(S,T)=(V',W)$, $g=1$, $h=\mathrm{diag}(k,1)$, target weight
$l-\#\iota(\tau)$ (that of $\alpha\beta^{-1}$), $n_S^-=l$: $c=\prod_\tau\det(k_\tau)^{(l-\#\iota(\tau))-l}\det(1)^{\cdots}=
\prod_\tau\det(k_\tau)^{-\#\iota(\tau)}$, and $y^{++}\mapsto k^{\mathsf t}y^{++}$, $y^{--}\mapsto y^{--}$ (the other blocks do
not enter $\pteMP$). Then $\det(\eta_{\tau j}k^{\mathsf t}_\tau y^{++}_\tau)$ is the factor of $\pteMP(\eta k^{\mathsf t})$.
(c) Lemma~\ref{ptf:fock} with $g=t_z$ (positive block; $t_z$ is $z^{-2}$ on the first $p$ and $z^{-1}$ on the
last $a$ positive indices) and $h=z^{-1}1_a$: $c=\prod_\tau(z^{-a})^{-\#\iota(\tau)}=z^{aN_\iota}$, cancelled by
the factor $z^{-aN_\iota}$ of $\Omega_z$; the substitution is $x^+\mapsto z^{-1}x^+t_z^{-1}$ (the column of a first-$p$
index is multiplied by $z^{-1}z^2=z$, that of a last-$a$ index by $z^{-1}z=1$) and $x^-\mapsto zx^-$. With
$g=1$, $h=u$ the substitution on $x^\flat$ is $x^\flat\mapsto u^{\mathsf t}x^\flat$ and $c=\prod_\tau\det(u_\tau)^{-\#\iota(\tau)}$.
\end{proof}

\begin{remark}\label{ptf:rem-fock}
In the conventions above the transpose in \eqref{pte:eq-SDlaw} is forced: the Fock variables of
positive product sign are conjugate coordinates, and conjugation turns $h^{-1}$ and $g$ into
$h^{\mathsf t}$ and $g^{-\mathsf t}$. With these formulas Steps~3--5 of the proof of
Corollary~\ref{pte:explicit} hold verbatim: Step~3 uses (c), Step~4 uses (a) with $\eta=\eta^0$ and (b),
and the cancellation of the powers of $z$ and $\det k_\tau$ there is unchanged. The weight hypothesis on
$\beta$ is harmless: by \eqref{pte:eq-twist} changing $\beta$ multiplies both sides of
\eqref{pte:eq-explicit} by the same character of $\det g$ (Step~0).
\end{remark}

\subsubsection{The first factors}\label{ptf:ss-ff}

We use the notation of \S\ref{pte:ss-third} and \S\ref{pte:ss-left}: $\mathsf W=I_{3m}\oplus W'$,
$W'=W\oplus(E^j)^-\oplus E^{3m-j}\oplus W^\flat$, $V''=I_{N''}$, $W=W_j$ with $r_+$ positive and $r_-=m-j$
negative basis vectors. Put
\[
\mathsf C=I_{3m}\oplus(E^j)^-\oplus E^{3m-j}\oplus W^\flat,\qquad \mathsf W\cong W\oplus\mathsf C
\]
(reordering the summands; the point $(h',g,1,1,h^\flat_i)$ becomes $(g,c_i(h'))$ with
$c_i(h')=(h',1,1,h^\flat_i)\in\auU(\mathsf C)(\auA)$, whose archimedean component is $1$). On
$\auHom(\mathsf W,V'')\otimes_\tau\R$ the positive columns are those of $W$ (the $N''\times r_+$ block, with
Fock variables $X_\tau$, columns in the numbering of \eqref{po:eq-poly}) and those of $\mathsf C$ (Fock
variables $Z_\tau$, $N''\times(2m+N''-r_+)$); the negative columns are those of $(E^j)^-$ and of $W$, which
$\pteSD$ does not involve. In this notation $x^+_\tau=(X_\tau\mid Z_\tau)$, and $\pteSD(\eta_i)=\prod_\tau\prod_{k\in
\iota(\tau)}\det\binom{\eta_{i,\tau k}}{X_\tau\ Z_\tau}$, with $\eta_{i,\tau k}=(\eta^X_{\tau k}\mid\eta^Z_{\tau k})$ of size
$2m\times(2m+N'')$.

\emph{The shape of the constant rows.} By the construction of the choices (\S\ref{pte:ss-third}: the rows of
$\eta_{i,\tau k}$ are basis vectors of the padded $W$ for a column of $\mu^+_\tau$ of height $<r_+$, and
short rows without component on $W$ for a column of height $r_+$), the matrices $\eta_{i,\tau k}$ have the
following property, which is all that is used:
\begin{itemize}
\item[(H)] the elements $k\in\iota(\tau)$ correspond to the columns of $\mu^+_\tau$; if $h_k$ is the height of
the column of $k$, then $\eta^X_{\tau k}=(0_{2m\times h_k}\mid E_{\tau k})$, where the $r_+-h_k$ columns of
$E_{\tau k}$ are distinct standard basis vectors of $\C^{2m}$ (so $r_+-h_k$ of the rows of $\eta_{i,\tau k}$ are the
basis vectors $e_{h_k+1},\dots,e_{r_+}$ of the positive part of $W$, and the other rows vanish on $X$).
\end{itemize}

\begin{lemma}[Laplace expansion]\label{ptf:laplace}
Under (H), for each $\tau$, $k$ there are polynomials $Q_{\tau k,R}$ in the entries of $Z_\tau$, indexed by
the $h_k$-subsets $R\subseteq\{1,\dots,N''\}$, such that
\[
\det\binom{\eta_{i,\tau k}}{X_\tau\ Z_\tau}=\sum_{R}\det\bigl(X_{\tau;R,[h_k]}\bigr)\,Q_{\tau k,R}(Z_\tau),
\]
$X_{\tau;R,[h_k]}$ the minor of $X_\tau$ on the rows $R$ and the first $h_k$ columns. Consequently there are
unique polynomials $R_\lambda$ in the entries of $Z=(Z_\tau)_\tau$ with
\begin{equation}\label{ptf:eq-SDexp}
\pteSD(\eta_i)=\sum_\lambda\pop''_{u_\lambda}(X)\,R_\lambda(Z),
\end{equation}
where $(u_\lambda)$ is a fixed basis of $W''=W_{\mu^+}$ and $\pop''_u$ is the polynomial \eqref{po:eq-poly} for the
pair $(W,V'')$.
\end{lemma}
\begin{proof}
Let $M$ be the square matrix $\binom{\eta}{X\ Z}$ (we drop $\tau,k,i$). Expand $\det M$ along its first $h=h_k$
columns (generalised Laplace expansion): $\det M=\sum_{R'}\pm\det M_{R',[h]}\det M_{R'^c,[h]^c}$, the sum over
the $h$-subsets $R'$ of the rows. By (H) the constant rows vanish in the first $h$ columns, so only $R'$
consisting of rows of $X$ contribute, and then $M_{R',[h]}=X_{R,[h]}$. In the complementary matrix
$M_{R'^c,[h]^c}$ the $r_+-h$ constant rows $e_{h+1},\dots,e_{r_+}$ each have a single entry $1$, in the columns
$h+1,\dots,r_+$; expanding along them deletes these rows and the columns $h+1,\dots,r_+$ of $X$, up to a sign.
What remains has as rows the constant rows vanishing on $X$ (restricted to the columns of $Z$) and the rows
$R^c$ of $Z$; its determinant $\pm Q_R(Z)$ depends on $Z$ alone. This is the first assertion.

Multiplying over $\tau$ and $k$: $\pteSD(\eta_i)=\sum_{(R_k)}\prod_{\tau,k}\det X_{\tau;R_k,[h_k]}\cdot\prod_{\tau,k}Q_{\tau k,R_k}
(Z_\tau)$. By \eqref{po:eq-ev} and \eqref{po:eq-poly} (all of $V''$ is positive, so only $x^+$ and the factor
$\poT^+$ occur), $\pop''_\xi(X)=\langle\xi,\mathrm{ev}(X)\rangle$ with $\mathrm{ev}(X)=\bigotimes_{\mathfrak c}(X_1\wedge\dots\wedge
X_{h_{\mathfrak c}})=\sum_{(R_{\mathfrak c})}\prod_{\mathfrak c}\det X_{R_{\mathfrak c},[h_{\mathfrak c}]}\,e_{(R_{\mathfrak c})}$ ($X_i$ the columns
of $X$, $e_{(R)}$ the tensor of the wedge basis vectors, an orthonormal basis of $\poT^+$; the Fock variables $X$
are the conjugates $\aubar{x^+}$ of the coordinates, as in \eqref{po:eq-poly}). Hence the map $\xi\mapsto\pop''_\xi$,
$\poT^+\to\C[X]$, is conjugate-linear with image the span $\Sigma_{\mu^+}$ of the products $\prod_{\mathfrak c}\det
X_{R_{\mathfrak c},[h_{\mathfrak c}]}$, and its kernel is $W''^\perp$, because the vectors $\mathrm{ev}(X)$ span $W''$
(definition of $W''$); so $u\mapsto\pop''_u$ is a bijection of $W''$ onto $\Sigma_{\mu^+}$ and $(\pop''_{u_\lambda})$ is a
basis of $\Sigma_{\mu^+}$. The expansion above shows $\pteSD(\eta_i)\in\Sigma_{\mu^+}\otimes\C[Z]$, which gives
\eqref{ptf:eq-SDexp} with unique $R_\lambda$.
\end{proof}

\begin{lemma}[Sums in the source]\label{ptf:sum}
Let $S=S_1\oplus S_2$ be an orthogonal sum of spaces, $T$ a space, $P_\nu$, $\Psi_\nu$ a polynomial and a
function for $\auHom(S_\nu,T)$, and $\alpha,\beta_1,\beta_2$ characters. With $P=P_1P_2$ (variables of the two
column blocks) and $\Psi(x_1\mid x_2)=\Psi_1(x_1)\Psi_2(x_2)$,
\[
\theta^{S,T}_{\alpha,\beta_1\beta_2}\bigl((g_1,g_2),h;[P,\Psi]\bigr)=\theta^{S_1,T}_{\alpha,\beta_1}\bigl(g_1,h;[P_1,\Psi_1]\bigr)\,
\theta^{S_2,T}_{\alpha,\beta_2}\bigl(g_2,h;[P_2,\Psi_2]\bigr).
\]
\end{lemma}
\begin{proof}
$\auHom(S,T)=\auHom(S_1,T)\oplus\auHom(S_2,T)$ orthogonally; the Heisenberg representation, the theta functional
and the datum factor as tensor products ($[P,\Psi]=[P_1,\Psi_1]\otimes[P_2,\Psi_2]$, the vacuum of the sum of
the lattices being the tensor product of the vacua, normalisations compatible), and the Weil representation of
$\auU(\auHom(S_1,T))\times\auU(\auHom(S_2,T))$ is the tensor product, the splittings defined with $\upsilon$ being
compatible with orthogonal sums \cite{KudlaSplitting}; this is the source version of (W7) of
\S\ref{po:ss-glue}. The element $(g,h)$ acts on the two summands by $((g_1,h),(g_2,h))$, and the scalar
$\alpha(\det g)^{-1}\beta_1\beta_2(\det h)$ is $\alpha(\det g_1)^{-1}\beta_1(\det h)\cdot\alpha(\det g_2)^{-1}\beta_2(\det h)$.
\end{proof}

\begin{definition}[First factors]\label{ptf:def-ff}
Let $\nu''=\nu''_\vartheta$ and $\Phi^{\mathsf C}_i$ the function on the columns of $\mathsf C$:
$\Phi^{\mathsf C}_i(x_{\mathsf C})=\psi^{\mathrm{inv}}_{\Lambda_i}(x')\cdot\plone[\text{the columns of }x_{\mathsf C}\text{ on }(E^j)^-\oplus E^{3m-j}
\oplus W^\flat\text{ are those of }1]$, $x'$ the columns on $I_{3m}$ (read modulo $\pi_M^2$). For $h\in
\auU(V'')(\auA_f)$ put
\[
c_{\lambda}(h)=\pteCS_{K_{I_{3m}}}\Bigl(h'\mapsto\mathrm a_{\chi'}(h')\,\theta^{\mathsf C,V''}_{\nu'',\nu''\upsilon^{j}}\bigl(c_i(h'),h;
[R_\lambda,\Phi^{\mathsf C}_i]\bigr)\Bigr),\qquad\varphi''_{\vartheta,i}(h)=\sum_\lambda\aubar{c_\lambda(h)}\,u_\lambda\in W''.
\]
\end{definition}

The definition does not depend on the basis $(u_\lambda)$: by \eqref{ptf:eq-SDexp}, $\varphi''_{\vartheta,i}(h)$ is the
unique $\varphi\in W''$ with $\pop''_{\varphi}(X)=\sum_\lambda c_\lambda(h)\pop''_{u_\lambda}(X)$, and $\sum_\lambda R_\lambda\otimes
\pop''_{u_\lambda}$ is basis-free.

\begin{proposition}[Factorisation and transformation law]\label{ptf:ff}
Let $\theta''_{\vartheta,\Lambda}(g,h;u)=\theta^{W,V''}_{\nu'',\upsilon^{m-j}}(g,h;[\pop''_u,\Psi''_{\vartheta,\Lambda}])$, the kernel of
\S\ref{po:ss-glue} (conjugate-linear in $u$), and let $\Phi_0$ be the datum of Definition~\ref{pte:def-ff}. Then for
$g\in\auU(W)(\auA)$, $h\in\auU(V'')(\auA_f)$:
\begin{enumerate}
\item[(a)] $\displaystyle\pteCS_{K_{I_{3m}}}\Bigl(h'\mapsto\mathrm a_{\chi'}(h')\,\theta^{\mathsf W,V''}_{\nu'',\nu''\upsilon^m}\bigl((h',g,1,1,h^\flat_i),h;
[\pteSD(\eta_i),\Phi_0]\bigr)\Bigr)=\theta''_{\vartheta,\Lambda_i}\bigl(g,h;\varphi''_{\vartheta,i}(h)\bigr)$;
\item[(b)] $\varphi''_{\vartheta,i}(hk)=\chi'_\vartheta(\det\bar k)^{-1}\varphi''_{\vartheta,i}(h)$ for $k\in J'_{\vartheta,i}$, $\bar k$ its
reduction modulo $\pcm_\vartheta$; and $\theta''_{\vartheta,\Lambda_i}(g,hk;u)=\chi'_\vartheta(\det\bar k)^{-1}\theta''_{\vartheta,\Lambda_i}(g,h;u)$;
\item[(c)] $\varphi''_{\vartheta,i}(\gamma h)=\rho''(\gamma_\infty)\varphi''_{\vartheta,i}(h)$ for $\gamma\in\auU(V'')(E^+)$, where
$\rho''$ is the representation of $\auU(V'')_\infty$ on $W''$ defined by $\omega''_\infty(1,u)\pop''_\xi=\pop''_{\rho''(u)\xi}$
(the equivariance of \eqref{po:eq-poly}; this is $\rho_{\mu^+}$ in the conventions of \S\ref{po:ss-glue}).
\end{enumerate}
In particular the summand $h\mapsto\theta''_{\vartheta,\Lambda_i}(g,h;\varphi''_{\vartheta,i}(h))$ is right invariant under
$J'_{\vartheta,i}$ and left invariant under $\auU(V'')(E^+)$.
\end{proposition}
\begin{proof}
(a) The datum $\Phi_0(x_W\mid x_{\mathsf C})=\mathrm{col}(x_W)\cdot\Phi^{\mathsf C}_i(x_{\mathsf C})$ is a product over the
column blocks $W$, $\mathsf C$ (the columns $x''e$ are the columns of $W$; the condition on the other columns
of $x''$ and the factor $\psi^{\mathrm{inv}}_{\Lambda_i}(x')$ concern $\mathsf C$), and $\mathrm{col}$ is the finite part at
$\pcm_\vartheta$ of $\Psi''_{\vartheta,\Lambda_i}$ (Lemma~\ref{pte:column}). By \eqref{ptf:eq-SDexp} and linearity of the
kernel in the polynomial, then Lemma~\ref{ptf:sum} with $\alpha=\nu''$, $\beta_1=\upsilon^{m-j}$, $\beta_2=\nu''\upsilon^j$
($\beta_1\beta_2=\nu''\upsilon^m$),
\[
\theta^{\mathsf W,V''}_{\nu'',\nu''\upsilon^m}\bigl((g,c_i(h')),h;[\pteSD(\eta_i),\Phi_0]\bigr)=\sum_\lambda\theta^{W,V''}_{\nu'',\upsilon^{m-j}}
\bigl(g,h;[\pop''_{u_\lambda},\Psi''_{\vartheta,\Lambda_i}]\bigr)\,\theta^{\mathsf C,V''}_{\nu'',\nu''\upsilon^j}\bigl(c_i(h'),h;[R_\lambda,
\Phi^{\mathsf C}_i]\bigr).
\]
The first factor is $\theta''_{\vartheta,\Lambda_i}(g,h;u_\lambda)$ and does not involve $h'$; the finite class sum over
$h'$ with the weights $\mathrm a_{\chi'}(h')$ therefore produces $\sum_\lambda c_\lambda(h)\theta''_{\vartheta,\Lambda_i}(g,h;u_\lambda)=
\theta''_{\vartheta,\Lambda_i}(g,h;\sum_\lambda\aubar{c_\lambda(h)}u_\lambda)$ by conjugate-linearity.

(b) Let $k\in J'_{\vartheta,i}$ and $A=\bar k$; $A\in H_{\vartheta,i}$, so the columns of $A$ off $e$ are those of
$1$ and $w_iA\equiv w_i\bmod\pi_M^2$; the same holds for $A^{-1}$. By Lemma~\ref{pc:level}, right translation
of a kernel of $(S,V'')$ with target character $\beta$ by $k$ replaces the finite datum $\Psi$ by $\beta(\det
k)\Psi(A^{-1}\,\cdot)$. \emph{The $\mathsf C$-factor:} $A^{-1}$ maps the columns of $1$ off $e$ to the columns of
$A^{-1}$ off $e$, which are those of $1$, and $\psi^{\mathrm{inv}}_{\Lambda_i}(A^{-1}x')=\psi^{\mathrm{inv}}_{\Lambda_i}(x')$ because
$\Lambda_iA^{-1}\equiv\Lambda_i\bmod\pi_M^2$; so $\Phi^{\mathsf C}_i(A^{-1}\cdot)=\Phi^{\mathsf C}_i$, the polynomial is unchanged ($k$ is
finite), and $c_\lambda(hk)=(\nu''\upsilon^j)(\det k)\,c_\lambda(h)$. Now $\upsilon$ is unramified at the finite places, so
$\upsilon(\det k)=1$, and $\nu''$ is tame for $\pcm_\vartheta$ with residue character $\chi'_\vartheta$ (the character
$\nu''\upsilon^m$ is, \S\ref{pte:ss-third}), so $\nu''(\det k)=\chi'_\vartheta(\det A)$. Hence $\varphi''_{\vartheta,i}(hk)=
\aubar{\chi'_\vartheta(\det A)}\varphi''_{\vartheta,i}(h)=\chi'_\vartheta(\det A)^{-1}\varphi''_{\vartheta,i}(h)$. \emph{The $W$-factor:}
write $A^{-1}=\bigl(\begin{smallmatrix}B&0\\C&1\end{smallmatrix}\bigr)$ with the indices $e$ first. For $\sigma\in\pcMat_{N''\times
r}$, $(A^{-1}\sigma)'=B\sigma'$ and $\det B=\det A^{-1}$, while $w_iA^{-1}\sigma=w_i\sigma$; so $\mathrm{col}(A^{-1}\sigma)=\chi'_\vartheta
(\det A)^{-1}\mathrm{col}(\sigma)$. With $\upsilon^{m-j}(\det k)=1$ this gives the second assertion. 

(c) Let $\gamma\in\auU(V'')(E^+)$, $c\in\auU(\mathsf C)(\auA_f)$. By (W1) the kernel is invariant under the diagonal
rational point: $\theta^{\mathsf C,V''}(c,(1,\gamma h);f)=\theta^{\mathsf C,V''}(c,(\gamma_\infty^{-1},h);f)$ (the source point $c$ is
unchanged, $\gamma$ acting on $V''$ only), and the archimedean component $\gamma_\infty^{-1}$ acts on the polynomial
through $\omega_\infty(1,\gamma_\infty^{-1})$. By Corollary~\ref{ptf:SDlaw}(a) (applied to $(\mathsf W,V'')$, with $k=1$),
$\omega_\infty(1,u)\pteSD(\eta_i)=\pteSD(\eta_i)$ for $u\in\auU(V'')_\infty$; by Lemma~\ref{ptf:sum} at infinity the left
side is $\sum_\lambda\omega''_\infty(1,u)\pop''_{u_\lambda}\cdot\omega^{\mathsf C}_\infty(1,u)R_\lambda=\sum_\lambda\pop''_{\rho''(u)u_\lambda}\cdot
\omega^{\mathsf C}_\infty(1,u)R_\lambda$. Write $\rho''(u)u_\lambda=\sum_\mu a_{\mu\lambda}(u)u_\mu$; then $\pop''_{\rho''(u)u_\lambda}=
\sum_\mu\aubar{a_{\mu\lambda}(u)}\pop''_{u_\mu}$, and comparing coefficients of the basis $\pop''_{u_\mu}$ (unique by
Lemma~\ref{ptf:laplace}) gives $\sum_\lambda\aubar{a_{\mu\lambda}(u)}\,\omega^{\mathsf C}_\infty(1,u)R_\lambda=R_\mu$, i.e.\
$\omega^{\mathsf C}_\infty(1,u^{-1})R_\mu=\sum_\lambda\aubar{a_{\mu\lambda}(u)}R_\lambda$. With $u=\gamma_\infty$, $c_\mu(\gamma h)=\sum_\lambda\aubar{a_{\mu\lambda}
(\gamma_\infty)}c_\lambda(h)$, so $\varphi''_{\vartheta,i}(\gamma h)=\sum_{\mu,\lambda}a_{\mu\lambda}(\gamma_\infty)\aubar{c_\lambda(h)}u_\mu=\rho''(\gamma_\infty)
\varphi''_{\vartheta,i}(h)$. The last sentence follows from (a), (b), (c) and the invariance of the left side of (a):
under $J'_{\vartheta,i}$ by (b) (the two characters cancel by conjugate-linearity), under rational points by (W1).
\end{proof}

\subsubsection{The character bookkeeping on $V''$}\label{ptf:ss-left}

\begin{proposition}[The left side]\label{ptf:leftside}
With $J'_{\vartheta,i}$ in place of $J_{\vartheta,i}$ in \S\ref{po:ss-glue}, the character $\chi''_\vartheta(k)=\chi'_\vartheta(\det\bar k)$ on
$J'_{\vartheta,i}$, and the first factors of Definition~\ref{ptf:def-ff},
\[
\mathcal L_{\vartheta,i}(g)=\sum_{h\in\Gamma''\backslash\auU(V'')(\auA_f)/J'_{\vartheta,i}}\theta''_{\vartheta,\Lambda_i}\bigl(g,h;\varphi''_{\vartheta,i}(h)\bigr)
=\Xi_{\vartheta,i}(g)\qquad(g\in\auU(W)(\auA)),
\]
and $\varphi''_{\vartheta,i}\in\poA_{\mu^+}(J'_{\vartheta,i},\chi''_\vartheta{}^{-1})$.
\end{proposition}
\begin{proof}
\emph{Characters.} The kernel in $\mathcal L_{\vartheta,i}$ is $\theta^{\mathsf W,V''}_{\nu'',\nu''\upsilon^m}$: source character $\nu''$
on $\auU(\mathsf W)$, target character $\chi'\nu^\flat\upsilon^m=\nu''\upsilon^m$ on $\auU(V'')$. In Lemma~\ref{ptf:sum} (sums in
the source) the source character is common to the summands and the target character is distributed; we
distributed it as $\upsilon^{m-j}$ on the summand $W$ and $\nu''\upsilon^j$ on $\mathsf C$. The kernel $\theta''$ of
\S\ref{po:ss-glue} is the $V''$-factor of $\omega_\mu$ in (W7): by \eqref{po:eq-w5}, $\omega_\mu=\omega'_{a,b}=\omega_{b,a}$
(\eqref{pcx:eq-dict}) with $a=\upsilon^{r_-}$ on $\auU(V)$ and $b=\mu\upsilon^{\lfloor n/2\rfloor}$ on $\auU(W)$. In (W7) the sum is in
the target $V=V''\oplus V^\circ$, so (by the target analogue of Lemma~\ref{ptf:sum}) the target character
$\upsilon^{r_-}$ is common to the two factors and the source character $b$ is distributed as $\nu''\cdot\upsilon^{m/2-j}$
($\nu''$ on the $V''$-factor, $\upsilon^{m/2-j}$, independent of $\vartheta$, on the $V^\circ$-factor); the identity
$\nu''\upsilon^{m/2-j}=\mu\upsilon^{\lfloor n/2\rfloor}$ follows from the definition of $\nu''_\vartheta$ in \S\ref{pte:ss-third}: with $\mu=\mu^0_j\vartheta$, $n=N''+
\mathrm{rk}\,V^\circ=6r$, $t+k-1=5r+k-1$ even and $r=k+m-1$, one has $j+\frac{t+k-1}2+\frac m2-j=\frac{5r+k-1+m}2=3r=\frac n2$
(compare Lemma~\ref{psv:char}). Hence $\omega''=\omega_{\nu'',\upsilon^{r_-}}$, and $r_-=m-j$. This is exactly the pair of characters of the
$W$-factor in Proposition~\ref{ptf:ff}(a), so the kernel there is the $\theta''_{\vartheta,\Lambda_i}$ of \S\ref{po:ss-glue},
with the datum $\Psi''_{\vartheta,\Lambda_i}$ (Lemma~\ref{pte:column}) and the polynomial $\pop''$ of \eqref{po:eq-poly}.

\emph{Sums.} By Lemma~\ref{pte:rowstab} with $F_0=\psi^{\mathrm{inv}}_{\Lambda_i}$ (invariant under $H_{\vartheta,i}$, as in
Definition~\ref{pte:def-ff}), the inner sum of $\mathcal L_{\vartheta,i}$ is the class sum modulo $J'_{\vartheta,i}$ of the kernel with
datum $\Phi_0$. Exchange it with the finite class sum over $h'$ and apply Proposition~\ref{ptf:ff}(a): $\mathcal L_{\vartheta,i}(g)
=\sum_{h\bmod J'_{\vartheta,i}}\#(\Gamma''\cap hJ'_{\vartheta,i}h^{-1})^{-1}\theta''_{\vartheta,\Lambda_i}(g,h;\varphi''_{\vartheta,i}(h))$, and the stabilisers are
trivial because $J'_{\vartheta,i}$ is free (Lemma~\ref{au:free}). By Proposition~\ref{ptf:ff}(b), (c) the summand is a
function on the double cosets and $\varphi''_{\vartheta,i}$ has the stated type ($\rho''=\rho_{\mu^+}$). This is the definition of
$\Xi_{\vartheta,i}$ with the level $J'_{\vartheta,i}$ and the character $\chi''_\vartheta$.
\end{proof}


\subsubsection{Weights, the constant rows and the character $\chi''_\vartheta$}\label{ptf:ss-close}

\emph{Weights.} A character of $\auU(1)$ of weight $(c_\tau)$ is one with $\mu_\tau(u)=\tau(u)^{c_\tau}$ at the
infinite places (Lemma~\ref{pz:mu}; this is the meaning of ``weight'' in \S\ref{ss:au-spaces} and
\S\ref{pcx:ss-statement}). This is the convention used in \S\ref{ptf:ss-fock}, so Lemma~\ref{ptf:fock} and
Corollary~\ref{ptf:SDlaw} hold as stated. As a check, $\upsilon$ has weight $1$ ($\upsilon_\tau(z)=z/|z|$, which is
$z$ on $\auU(1)(\R)$).

\begin{lemma}[The constant rows satisfy (H)]\label{ptf:H}
Let the matrices $\eta_{i,\tau k}$ be defined as follows (this is the precise form of the description in
\S\ref{pte:ss-third}). Their columns are the $r_+$ positive indices of $W$, numbered $1,\dots,r_+$ as in
\eqref{po:eq-poly}, followed by the indices of $E^{3m-j}\oplus W^\flat\oplus I_{3m}$. The rows are numbered
$a=0,\dots,2m-1$.
\begin{itemize}
\item For $k$ belonging to a column of $\mu^+_\tau$ of height $h<r_+$, row $a$ is the basis vector $e_{h+a+1}$ of
$W$ if $h+a<r_+$, and otherwise the basis vector number $a$ of $E^{3m-j}$.
\item For $k$ belonging to a column of height $r_+$ (the $\Delta'$ columns), the rows are the $2m$ short rows
selected by $\varrho_i$, read at $\tau$ and placed in the columns of $W^\flat$, with $0$ elsewhere.
\end{itemize}
Then (H) holds.
\end{lemma}
\begin{proof}
For a column of height $r_+$, $\eta^X_{\tau k}=0$, which is (H) with $h_k=r_+$. For height $h<r_+$, the rows with
$h+a<r_+$, that is $a=0,\dots,r_+-h-1$, are $e_{h+1},\dots,e_{r_+}$, each exactly once. The remaining rows lie in
$E^{3m-j}$ and vanish on the columns of $X$. So $\eta^X_{\tau k}=(0_{2m\times h}\mid E_{\tau k})$, where $E_{\tau k}$
has as columns the distinct basis vectors $e_0,\dots,e_{r_+-h-1}$ of $\C^{2m}$. This uses $r_+-h\le2m$, i.e.\ every column of $\mu^+_\tau$ of height $h<r_+$ has $h\ge r_+-2m$, which we now
derive from the diagram. By Lemma~\ref{plf:match} and its proof, for $h<r_+$ the number of columns of height $h$ is
$c^+_\tau(h)=\varsigma_\tau(h-1)-\varsigma_\tau(h)-1$, where $\varsigma_\tau(0)>\varsigma_\tau(1)>\cdots$ are the entries of
$\varsigma_\tau$ on the positive coordinates of $W_j$ in decreasing order. By \S\ref{ss:au-spaces} these are the
$k_0-1$ entries $\frac{k_0}2-1,\frac{k_0}2-2,\dots,1-\frac{k_0}2$ followed by $-\kappa(\tau)-a\Delta$, $1\le a\le j$ (so $r_+=k_0-1+j$), where $k_0$ denotes the integer weight of \S\ref{ss:au-spaces} (written $k_0$ here to keep $k$ for the elements of $\iota(\tau)$).
Hence $\varsigma_\tau(h-1)-\varsigma_\tau(h)=1$ and $c^+_\tau(h)=0$ for $1\le h\le k_0-2$. So a column of height $h<r_+$ has
$h\ge k_0-1$, and $r_+-h\le r_+-(k_0-1)=j\le m\le2m$.
\end{proof}

\begin{lemma}[The character $\chi''_\vartheta$]\label{ptf:chi}
Let $k\in J'_{\vartheta,i}$. Then the character $\chi''_\vartheta(k)$ of \S\ref{po:ss-glue} (the value of $\chi_\vartheta$ on
$\pog_i^{-1}\mathrm{diag}(k,k^\circ)\pog_i$) equals $\chi'_\vartheta(\det\bar k)$.
\end{lemma}
\begin{proof}
\emph{$\chi_\vartheta$.} Let $k\in K_\vartheta$. By \eqref{po:eq-w5} with $g=1$, $\omega_\mu(1,k)\Phi_\vartheta(T)=\eta_w(1,k)\Phi_\vartheta(k^{-1}T)$,
and $\eta_w(1,k)=1$ because $\eta_w$ is unramified and $k$ is integral. At $w\in\mathsf M$ the factor
$\plone[T\equiv\iota_0]$ is invariant, since $k\iota_0\equiv\iota_0$. At $w\in\mathsf Q$, $k\bmod w$ is block lower
triangular for the blocks $r+t$, so $k^{-1}=\bigl(\begin{smallmatrix}A&0\\ *&*\end{smallmatrix}\bigr)$ modulo $w$ with
$A=k_{11}^{-1}$, where $k_{11}$ is the upper left $r\times r$ block. Hence $(k^{-1}T)'=AT'$ and
$\psi_w(\det(k^{-1}T)')=\psi_w(\det k_{11})^{-1}\psi_w(\det T')$. Therefore
$\chi_\vartheta(k)=\prod_{w\in\mathsf Q}\psi_w(\det k_{11}\bmod w)$.

\emph{$\chi''_\vartheta$.} At $\mathsf Q$ the element $\pog_i$ is $1$ (proof of Lemma~\ref{po:slice}). So the $\mathsf Q$-component of
$\pog_i^{-1}\mathrm{diag}(k,k^\circ)\pog_i$ is $\mathrm{diag}(k,k^\circ)$, whose upper left $r\times r$ block is the block $k_{11}$
of $k$ on the first $r$ indices of $V''$, which are the indices $e$. Hence $\chi''_\vartheta(k)=\prod_{w\in\mathsf Q}\psi_w(\det
k_{11})$.

\emph{Comparison.} For $k\in J'_{\vartheta,i}$ and $w\in\mathsf Q$, the columns of $\bar k_w$ off $e$ are those of $1$, so
$\bar k_w=\bigl(\begin{smallmatrix}k_{11}&0\\ *&1\end{smallmatrix}\bigr)$ and $\det\bar k_w=\det k_{11}$. At $w\in\mathsf M$,
$k\equiv1\bmod\pi_M$, so the factor $\omega^j(\det\bar k_w)$ of $\chi'_\vartheta$ is $1$. The residue character of $\chi'$ at
$w\in\mathsf Q$ is $\psi_w$: here $\psi_w$ is the tame part at $w$ of $\mu=\mu^0_j\vartheta$ (\S\ref{po:ss-set}), which is
that of $\vartheta$ because $\mu^0_j$ is unramified outside the chosen places (\S\ref{ss:au-spaces}); and
$\chi'=\chi_0^{\,j}\vartheta$ with $\chi_0$ of trivial tame part off $\mathsf M$ (\S\ref{po:ss-hecke}), so the residue
character of $\chi'$ at $w\in\mathsf Q$ (\S\ref{pte:ss-third}) is also the tame part of $\vartheta$ at $w$. Hence $\chi'_\vartheta(\det\bar k)=\prod_{w\in\mathsf Q}\psi_w(\det k_{11})=
\chi''_\vartheta(k)$.
\end{proof}

By Lemma~\ref{ptf:chi}, the character $\chi''_\vartheta$ of \S\ref{po:ss-glue} restricted
to $J'_{\vartheta,i}$ is the character of Proposition~\ref{ptf:ff}(b), and Proposition~\ref{ptf:leftside} holds with the
definition of \S\ref{po:ss-glue} unchanged.


%% ===== proofs/C-lifts.tex =====
% proofs/C-lifts.tex
% Statements of the one-lift block of Theorem C (proofs/C-onelift.tex, labels po:):
% Lemma po:classes, statement (L4); Lemmas plf:KS, plf:I used in C-firstfactors.tex.
% Propositions po:periods and po:span are NOT treated in this file (see the last subsection).
% To be \input after C-onelift.tex. Labels of this file start with plf:. Only amsmath/amssymb/amsthm.
\providecommand{\auO}{\mathcal O}
\providecommand{\auU}{\mathrm U}
\providecommand{\auGL}{\mathrm{GL}}
\providecommand{\auHom}{\operatorname{Hom}}
\providecommand{\auA}{\mathbb A}
\providecommand{\auF}{\mathbb F}
\providecommand{\auPl}{\Phi}
\providecommand{\auNm}{\operatorname{Nm}}
\providecommand{\auAut}{\operatorname{Aut}}
\providecommand{\aubar}[1]{\overline{#1}}
\providecommand{\pao}{\mathfrak o}
\providecommand{\psvac}{\mathbf v}
\providecommand{\plone}{\mathbf 1}
\providecommand{\plgb}[2]{\|#2\|_{#1}}
\providecommand{\psHS}{\mathcal H}
\providecommand{\psrho}{\varrho}
\providecommand{\pzF}{\mathfrak F}
\providecommand{\poA}{\mathcal A}
\providecommand{\poAA}{A_\ast}
\providecommand{\pop}{\mathsf p}
\providecommand{\pog}{\mathsf g}
\providecommand{\plfK}{\mathsf K}
\providecommand{\pkm}{\mathsf m}

\subsection{One lift: classes, the type of the descent, first factors}\label{plf:sec}

This subsection proves two of the statements that \S\ref{po:sec} isolates: Lemma~\ref{po:classes}
(numbers of classes) and the statement (L4) (the descent of a form on $V_j$ lies in $\poA(W_j)$). It
also proves the kernel scales (Lemma~\ref{plf:KS}) and the count of choices (Lemma~\ref{plf:I}), which are
used in \S\ref{pfq:sec}, where
Lemma~\ref{po:ff} is proved (Theorem~\ref{pfq:ff}). We keep the notation of \S\ref{po:sec}: $1\le j\le m$
is fixed, $W=W_j$, $V=V_j$, $n=r+t=6r$, $r_+=k+j-1$, $r_-=m-j$, $\poAA=fmb(\log b)^2$.
Propositions~\ref{po:periods} and~\ref{po:span} are proved in \S\ref{ppe:sec} (Theorem~\ref{ppe:main}) and
\S\ref{psp:sec} (Proposition~\ref{psp:main}, with Theorem~\ref{psv:main}).

\subsubsection{Numbers of classes}\label{plf:ss-classes}

\begin{lemma}[Lemma~\ref{po:classes}]\label{plf:classes}
Let $X=E^{n_X}$ be a standard definite space with $n_X\le2n_m$, $G=\auU(X)$, $\Gamma=G(E^+)$, and
let $J$ be a subgroup of $G_f$ containing the principal congruence subgroup $K(\pi_M^2)$ of level
$\pi_M^2$ in the stabiliser $K^\circ$ of $\widehat\auO_E^{\,n_X}$. Then $\#(\Gamma\backslash G_f/J)\le
\exp(2x^{26})\le e^b$.
\end{lemma}
\begin{proof}
The proof of Proposition~\ref{pf:classes} does not use the hypothesis $n\le n_m$ except in its last
paragraph (sizes): Lemma~\ref{au:shortbasis} holds for every rank, and the counting and the index
of the level are valid for every $n$ and every integer $\mathsf G\ge1$. Since
$\pi_M\aubar{\pi_M}=M$, the integer $M^2$ lies in $\pi_M^2\auO_E$, so $K(M^2)\subseteq K(\pi_M^2)
\subseteq J$, and Proposition~\ref{pf:classes} with $\mathsf G=M^2$ gives
\[
\#(\Gamma\backslash G_f/J)\le(2\mathsf B+1)^{4fn_X^2}M^{4fn_X^2},\qquad\mathsf B=n_X^2C_{\mathrm{sb}}^2
\cdot2f\sqrt d,\quad C_{\mathrm{sb}}=6q_1d(fn_X)^2 .
\]
Sizes: $f\le x^{12}$, $n_X\le2n_m\le80x^6$, and $k=r+1-m\ge x^5/2$, so $n_X/k\le160x$ and, as in
the proof of Proposition~\ref{pf:classes}, $C_{\mathrm{sb}}=6q_1d(fk)^2(n_X/k)^2\le x^{100}\cdot
(160x)^2\le x^{103}$, $\mathsf B_0=2f\sqrt d\le C_{\mathrm{sb}}$, $n_X\le C_{\mathrm{sb}}$ and $2\mathsf
B+1\le C_{\mathrm{sb}}^6$. Hence $4fn_X^2\log(2\mathsf B+1)\le4x^{12}\cdot6400x^{12}\cdot618\log x\le
x^{26}$, and, with $\log M\le23\log b\le x$ by \eqref{eq:au-sizes} and \eqref{eq:au-scales},
$4fn_X^2\log M\le25600x^{24}\cdot x^{1/100}\cdot23\le x^{26}$. Finally $2x^{26}\le b$ by
\eqref{eq:au-scales}.
\end{proof}

\begin{remark}\label{plf:rem-classes}
The spaces to which Lemma~\ref{po:classes} is applied in \S\ref{po:sec} and below are $V_j$
(rank $6r\le2n_m$), $V''$ and $V^\circ$ (ranks $N''$, $n^\circ<6r$), and the auxiliary spaces
$E^m$, $E^{3m}$, $W^\flat=E^{6m+1}$ of \S\ref{plf:ss-ff}; their levels ($K_\vartheta$, $J_{\vartheta,i}$,
$J^\circ$, $K^\circ$, $K(\pi_M^2)$) all contain $K(\pi_M^2)$, except $K_\vartheta$ and $J_{\vartheta,i}$
for $\vartheta\ne1$, which are defined also by congruences at $\mathsf Q$; for those
Proposition~\ref{pf:classes} with $\mathsf G=M^2\auNm Q_\vartheta$ applies, but only $\vartheta=1$ is
used with a bound on the number of classes (proof of (G5)). In \S\ref{pfq:sec} we also use the
order of the stabilisers: put
\[
n'_{\mathrm{Aut}}=\#\auGL_{4fn_m}(\auF_3)\le3^{16f^2n_m^2}\le\exp(x^{40}).
\]
By the proof of Lemma~\ref{po:naut}, for every standard definite $X$ of rank at most $2n_m$ and
every $g\in\auU(X)(\auA_f)$ the order of $\Gamma\cap gK^\circ g^{-1}$ divides $n'_{\mathrm{Aut}}$
($\auGL_{2fn_X}(\auF_3)$ is a subgroup of $\auGL_{4fn_m}(\auF_3)$).
\end{remark}

\subsubsection{The type of the descent: statement (L4)}\label{plf:ss-L4}

\emph{The type of $\poA(W_j)$.} At $\tau\in\auPl$ the group $\auU(W)(E^+_\tau)$ is $\auU(r_+,r_-)$
for the signs of $W_j$, and $\poA(W_j)$ is the space of forms of type $e_W=\bigotimes_\tau e_{W,\tau}$
and level $J_M$ in the sense of Definition~\ref{pz:def-type} (with $G_j$ replaced by $\auU(W)$):
$e_{W,\tau}=d_\tau\aubar{c_{W,\tau}}$, $c_{W,\tau}(x)=\langle v,\pi(x)v\rangle$ for the discrete
series $\pi=\pi_{\varsigma_\tau}$ with Harish-Chandra parameter $\varsigma_\tau$, the signs being those
of \S\ref{ss:au-spaces}, and $v$ a unit highest weight vector of its lowest $K$-type. We number the
$r$ entries of $\varsigma_\tau$ decreasingly, $\varsigma_\tau(0)>\dots>\varsigma_\tau(r-1)$; they are
\[
\tfrac k2-1,\dots,1-\tfrac k2\ \ (k-1\text{ entries}),\qquad-\kappa(\tau)-a\Delta\ \ (1\le a\le m),
\]
all integers ($k$ is even, $\Delta$ is even); the first $r_+=k-1+j$ of them lie on the positive
coordinates of $W_j$ and the last $r_-=m-j$ on the negative ones (\S\ref{ss:au-spaces}). Two
entries on coordinates of different signs differ by at least $\Delta=2r^8$.

\begin{lemma}\label{plf:A12}
The conclusions of Theorem~\ref{par:A1} (with $B=x^{25}$, $N_j$ replaced by $r$) and of
Corollary~\ref{par:A2} hold for $c_W=\bigotimes_\tau c_{W,\tau}$. Consequently the conclusions of
Proposition~\ref{pz:types}{\rm(a)} hold for forms of type $e_W$ or $\aubar{e_W}$ on $\auU(W)$: they
are cusp forms, rapidly decreasing on Siegel sets, and those of a given level form a
finite-dimensional space.
\end{lemma}
\begin{proof}
The proofs of Theorem~\ref{par:A1} and Corollary~\ref{par:A2} use only the decay bound
\eqref{par:eq-decay} with $c=B+N^2+1$, respectively $c=5N^3$, for the parameter at hand (Remark~\ref{par:FJ},
proof of Corollary~\ref{par:A2}). By Theorem~\ref{pco:FJ} and Theorem~\ref{pds:decay}, the bound holds
for every regular parameter in $(\frac{N-1}2+\Z)^N$ whose gap $h$ satisfies $h\ge N^2+1$, for every
integer $c\le h$. Here $N=r$, the gap of $\varsigma_\tau$ is at least $\Delta=2r^8$, and
$2r^8\ge x^{25}+r^2+1$, $2r^8\ge5r^3$ ($r\ge x^5$). The last assertion is the proof of
Proposition~\ref{pz:types}(a), which uses only (A1) and (A2) for the group at hand.
\end{proof}

\emph{The Fock model.} At $\tau$ we use the Fock model $\pzF_\tau$ of \S\ref{par:ss-fock} for the
pair $(\auU(W),\auU(V))$: rows $R$ = the basis of $W_j$ (signs $\varepsilon_a$), columns $C$ = the
basis of $V_j$ (all signs $\varepsilon'_c=+1$), $n_R=r$, $n_C=n=r+t$, variables $z_{ac}$, operators
$E_{ab}$, $E'_{cd}$ as in \eqref{par:eq-osc}, with constants $\gamma,\gamma'$.

\begin{lemma}[Constants]\label{plf:const}
For the representation $\omega_\mu$ of \eqref{po:eq-w5} the derived actions at $\tau$ are
\eqref{par:eq-osc} after conjugation by an element of the diagonal tori, with
$\gamma=\kappa_j(\tau)+\frac n2$; the constant $\gamma'$ does not enter below.
\end{lemma}
\begin{proof}
This is the computation of the proof of Lemma~\ref{par:vac}. By (V$'$) of \S\ref{pco:ss-V}, the
Weil representation defined with $\upsilon$ contributes $\det(g)^{n/2}$ at $\tau$; its proof
(Theorem~\ref{pco:V}) uses only the signature of the hermitian space $\auHom(W,V)\otimes_{E,\tau}\C$:
for $j<m$ this space is indefinite ($W_j$ is), and Step~2 of that proof applies; for $j=m$ it is
definite and Step~3 applies, with (S). In the convention of \eqref{po:eq-w5} the group $\auU(W)$ is
further twisted by $b(\det g)^{-1}$ with $b=\mu\,\upsilon^{\lfloor n/2\rfloor}$ (\S\ref{po:ss-set}),
and at $\tau$, $\mu=\mu^0_j\vartheta$ has weight $-\kappa_j(\tau)$ ($\vartheta$ has weight $0$) and
$\upsilon_\tau(z)=z$ on $\auU(1)$. So $b(\det g)^{-1}$ contributes $\det(g)^{\kappa_j(\tau)-\lfloor n/2\rfloor}$, and, as in
Lemma~\ref{par:vac}, $\gamma=\tau_0+\frac n2$ with $\tau_0=\kappa_j(\tau)+\frac n2-\lfloor\frac
n2\rfloor=\kappa_j(\tau)$, since $n=6r$ is even. (The twist $a(\det h)$ of $\auU(V)$ changes only
$\gamma'$.)
\end{proof}

\emph{Matching.} Put $\lambda=-\varsigma_\tau$ (the parameter of $\aubar{\pi}$) and
$\xi_a=\lambda_a-\gamma+\frac n2=-\varsigma_\tau(a)-\kappa_j(\tau)$ for $a\in R$. Recall
$\kappa_j=\kappa+(j+\frac12)\Delta$.

\begin{lemma}[Layers]\label{plf:match}
Let $p(a)$ be the number of the row $a$: its position among the rows of its sign, the positive ones
numbered by decreasing $\varsigma_\tau$ and the negative ones by increasing $\varsigma_\tau$, from $0$.
Define $c\colon R\to C$ by: $c(a)$ is the coordinate number $p(a)$ of $V$ if $\varepsilon_a=+1$,
and the coordinate number $n-1-p(a)$ if $\varepsilon_a=-1$; let $B\subset C$ be the $t$ remaining
coordinates (the coordinates $r_+,\dots,r_++t-1$). Let $\gamma'$ be given, $s_0=\gamma+\gamma'-
\frac{r+n}2$, and define $\lambda'\in(\gamma'-\tfrac r2+\Z)^{C}$ by $\lambda'_{c(a)}=-\lambda_a+s_0$ and, on $B$ in
increasing order of coordinates, $\lambda'=\gamma'-\frac r2+\frac{t-1}2-b$ $(0\le b<t)$. Then:
\begin{enumerate}
\item[(a)] $(\lambda,\lambda')$ match in the sense of \eqref{par:eq-match}, with this $c$ and $B$;
\item[(b)] with $w_a$ as in \S\ref{par:ss-harm}, $|w_a|$ is the length of the row number $p(a)$ of
$\mu^+_\tau$ if $\varepsilon_a=+1$, and of the row number $p(a)$ of $\mu^-_\tau$ if
$\varepsilon_a=-1$ (rows of a diagram numbered from $0$, longest first).
\end{enumerate}
Here $\mu^\pm_\tau$ are the diagrams of \S\ref{po:ss-set}, defined by their numbers of columns of
each height $h\ge1$:
\[
c^+_\tau(h)=\bigl(\varsigma_\tau(h-1)-\max(\varsigma_\tau(h),\tfrac{t-1}2-\kappa_j(\tau))-1\bigr)^+,
\quad
c^-_\tau(h)=\bigl(\min(\varsigma_\tau(r-h-1),-\kappa_j(\tau)-\tfrac{t-1}2)-\varsigma_\tau(r-h)-1\bigr)^+ ,
\]
for $1\le h\le r$, with the convention $\varsigma_\tau(r)=-\infty$ (so that the maximum is
$\frac{t-1}2-\kappa_j(\tau)$ for $h=r$), and $c^-_\tau(r)=0$; the proof shows $c^\pm_\tau(h)=0$ for $h>r_\pm$.
\end{lemma}
\begin{proof}
\emph{Signs and sizes.} The entries on positive rows are $\ge-\kappa(\tau)-j\Delta=-\kappa_j(\tau)+
\frac\Delta2$, so $\xi_a\le-\frac\Delta2<-\frac t2$; those on negative rows are $-\kappa-a'\Delta$ with
$a'>j$, so $\xi_a=(a'-j-\frac12)\Delta\ge\frac\Delta2>\frac t2$. Hence $|\xi_a|>\frac t2$ and
$\operatorname{sgn}\xi_a=-\varepsilon_a$; since $\varepsilon'\equiv1$, $\eta_{a\,c(a)}=-\varepsilon_a
\varepsilon'_{c(a)}=\operatorname{sgn}\xi_a$. With $\xi'_c=\lambda'_c-\gamma'+\frac r2$ we get
$\xi'_{c(a)}=-\lambda_a+s_0-\gamma'+\frac r2=-\xi_a$ (because $s_0-\gamma'+\frac r2=\gamma-\frac n2$),
and $\{\xi'_c:c\in B\}=\{\frac{t-1}2-b\}$; $\varepsilon'$ is constant on $B$. All entries are
integers ($t$ odd). This is (a).

\emph{Row lengths.} By the computation in the proof of Lemma~\ref{par:relw}, for $\xi_a<0$
\[
2w_a=2\xi_a+1+2\!\!\sum_{\xi_a<\xi_k<0}\!\!\varepsilon_a\varepsilon_k+\varepsilon_a\sum_{c\in B}
\varepsilon'_c,
\]
and for $\xi_a>0$ the formula displayed there. (For $\xi_a<0$: in $2w_a=2\xi_a+A_a-\beta_a$ one has
$\beta_a=\varepsilon_a\sum_k\varepsilon_k\operatorname{sgn}\xi_k-\varepsilon_a\sum_B\varepsilon'_c$,
and $\operatorname{sgn}(\xi_k-\xi_a)-\operatorname{sgn}\xi_k$ equals $2$ if $\xi_a<\xi_k<0$ and $0$
otherwise.) Let $\varepsilon_a=+1$, $p=p(a)$, so $\varsigma_\tau(a)=\varsigma_\tau(p)$. The $k$ with
$\xi_a<\xi_k<0$ are the positive rows with smaller $\varsigma$, $r_+-1-p$ of them, and
$\sum_B\varepsilon'=t$; so $|w_a|=-w_a=\varsigma_\tau(p)+\kappa_j(\tau)-\frac{t-1}2-(r_+-p)$. On the
other hand the row number $p$ of $\mu^+_\tau$ has length $\sum_{h>p}c^+_\tau(h)$. For $h<r_+$,
$\varsigma_\tau(h)>-\kappa_j+\frac\Delta2>\frac{t-1}2-\kappa_j$, so $c^+_\tau(h)=\varsigma_\tau(h-1)-
\varsigma_\tau(h)-1$; for $h=r_+$, $\varsigma_\tau(r_+)\le-\kappa_j-\frac\Delta2$ and $c^+_\tau(r_+)=
\varsigma_\tau(r_+-1)-\frac{t-1}2+\kappa_j-1$; for $h>r_+$, $\varsigma_\tau(h-1)<\frac{t-1}2-\kappa_j$
and $c^+_\tau(h)=0$. Summing (telescoping) gives $\varsigma_\tau(p)-\frac{t-1}2+\kappa_j-(r_+-p)=|w_a|$.
Let $\varepsilon_a=-1$, $p=p(a)$; then $\varsigma_\tau(a)=\varsigma_\tau(r-1-p)=-\kappa-a'\Delta$ with
$a'=m-p$. The $k$ with $0<\xi_k<\xi_a$ are the negative rows $-\kappa-a''\Delta$ with $j<a''<a'$, that is
$a'-j-1$ of them, with $\varepsilon_a\varepsilon_k
=1$; so $w_a=\xi_a-\frac12-(a'-j-1)-\frac t2$, i.e.\ $|w_a|=(a'-j-\frac12)\Delta-\frac{t+1}2-(a'-j-1)$.
The row number $p$ of $\mu^-_\tau$ has length $\sum_{h>p}c^-_\tau(h)$; for $h<r_-$,
$c^-_\tau(h)=\varsigma_\tau(r-h-1)-\varsigma_\tau(r-h)-1$; for $h=r_-$, $\varsigma_\tau(r-r_--1)=
-\kappa-j\Delta>-\kappa_j-\frac{t-1}2$ and $c^-_\tau(r_-)=-\kappa_j-\frac{t-1}2-\varsigma_\tau(r-r_-)-1$; for $h>r_-$ the minimum is
$-\kappa_j-\frac{t-1}2<\varsigma_\tau(r-h)$ and $c^-_\tau(h)=0$. Telescoping, the length is
$-\kappa_j-\frac{t-1}2-\varsigma_\tau(r-1-p)-(r_--p)=(a'-j-\frac12)\Delta-\frac{t-1}2-(a'-j)$, which is
$|w_a|$.
\end{proof}

\begin{remark}\label{plf:rem-mu}
Lemma~\ref{plf:match} makes precise the diagrams $\mu^\pm_\tau$ of \S\ref{po:ss-set}. The two properties used in
\S\ref{po:sec} follow: $\mu^+_\tau$ has $\sum_hc^+_\tau(h)=\varsigma_\tau(0)+\kappa_j-\frac{t-1}2-r_+
\le\frac k2+\kappa_j\le4b$ columns, by \eqref{eq:au-sizes} and the definition of $\kappa$; and
$c^-_\tau(h)\le\Delta$ for each $h$. The column heights are at most $r_\pm$, the last displayed
identities of the proof.
\end{remark}

\emph{The polynomial of a highest weight vector.} Let $\xi_0\in W_{\lambda_V}$ be the vector dual
to the pure tensor of the first basis vectors: $\langle\xi_0,\mathrm{ev}(y,y')\rangle=\prod_{\mathfrak
c}\det\bigl(y_{i}(e_{l})\bigr)_{i,l<h_{\mathfrak c}}\cdot\prod_{\mathfrak c'}\det\bigl(y'_{i}(e_{n-1-l})
\bigr)_{i,l<h_{\mathfrak c'}}$ up to complex conjugation of all entries, the first product over the
columns of $\mu^+$, the second over those of $\mu^-$, in \eqref{po:eq-ev}. By \eqref{po:eq-poly},
written in the Fock variables (for a positive row $z_{ac}$ is the conjugate of $x_{ac}$, for a
negative row the conjugate-linear map $\flat$ produces the variable $z_{ac}$), $\pop_{\xi_0}$ is
\[
\pop_{\xi_0}=\prod_{\mathfrak c\in\mu^+}\det\bigl(z_{a\,c(a')}\bigr)_{a,a'\in R^+_{<h_{\mathfrak c}}}
\cdot\prod_{\mathfrak c'\in\mu^-}\det\bigl(z_{a\,c(a')}\bigr)_{a,a'\in R^-_{<h_{\mathfrak c'}}},
\]
where $R^\pm_{<h}$ is the set of rows of sign $\pm$ with number $<h$: the columns of the positive
rows are taken in the fixed numbering and those of the negative rows in the reverse one
(\S\ref{po:ss-set}), and the dual coordinates of the highest weight vector of $\Lambda^h(\C^n)^\vee$
are the last ones.

\begin{lemma}\label{plf:harm}
$\pop_{\xi_0}$ is, up to a non-zero scalar, the joint harmonic $\Phi$ of
Proposition~\ref{par:exist} for the matching pair $(\lambda,\lambda')$ of Lemma~\ref{plf:match}.
\end{lemma}
\begin{proof}
By Lemma~\ref{plf:match}(b) a row $a$ has $|w_a|$ layers, and the rows having a pair of level
$\ell$ are those with $|w_a|>\ell$, that is, since the row lengths of a diagram decrease, the rows
of sign $\varepsilon$ with number less than the height of the column number $\ell$ of
$\mu^{\varepsilon}_\tau$ (columns numbered from $0$, highest first). In the notation of
Proposition~\ref{par:exist}, $\varepsilon'\equiv1$, so $L_{\ell,\epsilon,1}=R^{\epsilon}_{<h}$ with
$h$ that height, and $M_{\ell,\epsilon,1}=\det(z_{a\,c(a')})_{a,a'\in R^\epsilon_{<h}}$. The product
over $\ell$ and $\epsilon$ is the product over the columns of $\mu^+$ and of $\mu^-$ displayed
above.
\end{proof}

\begin{proposition}[(L4)]\label{plf:L4}
Let $\vartheta$ be a twist and $\Xi\in\poA(V;\vartheta)$. Then $D_\vartheta\Xi\in\poA(W_j)$ (forms of
type $e_W$ and level $J_M$). Consequently statement {\rm(L4)} of \S\ref{po:ss-set} holds: if
$\Xi\in\varrho_V(e_0)\poA(V;\vartheta)$ is orthogonal to $\Theta_\vartheta(e_0\poA(W_j))$, then
$D_\vartheta\Xi=0$.
\end{proposition}
\begin{proof}
\emph{Step 1: the archimedean identity.} Let $\sigma=\omega_\infty(\cdot,1)$ on $\pzF=\bigotimes_\tau
\pzF_\tau$, and $\mathsf P=\int e_W(x)\sigma(x^{-1})\,dx$, which converges absolutely in $\pzF_s$ by
Lemmas~\ref{plf:A12} and~\ref{pz:growth}(a) ($s=D+1$, $D=frn$, $2s\le x^{25}$). We show
$\mathsf P\pop_\xi=\pop_\xi$ for all $\xi\in W_{\lambda_V}$. As in the proof of
Theorem~\ref{par:A3}, $\mathsf P=\bigotimes_\tau\sigma_\tau(e_{\lambda})$ with $\lambda=-\varsigma_\tau$
at $\tau$ (because $e_W(x^{-1})=\aubar{e_W(x)}$ and $\aubar{c_{W,\tau}}=c_{-\varsigma_\tau}$), and
$\pop_{\xi_0}=\bigotimes_\tau\pop_{\xi_0,\tau}$; so it suffices to show $\sigma_\tau(e_\lambda)\Phi=
\Phi$ at each $\tau$ for $\Phi=\pop_{\xi_0,\tau}$. We apply Lemma~\ref{par:rigid}. By
Lemma~\ref{plf:harm} and Definition~\ref{par:def-harm}, $\Phi$ has $T$-weight $\Lambda^K(\lambda)$
and $E_{ab}\Phi=0$ for $\varepsilon_a=\varepsilon_b$, $\lambda_a>\lambda_b$; and $\Phi$ has
$T'$-weight $\Lambda'=\Lambda^K(\lambda')=\lambda'-\rho'$ (the group $\auU(n)$ is compact, all roots
are compact) and is killed by $E'_{cd}$ for $\lambda'_c>\lambda'_d$, so it is a highest weight vector
for $\mathfrak{gl}_n$ and $\Omega'\Phi=(|\Lambda'+\rho'|^2-|\rho'|^2)\Phi=\kappa_{\lambda'}\Phi$. Since
$\sum_aE_{aa}\Phi=(\sum_a\Lambda_a)\Phi=(\sum_a\lambda_a)\Phi$ and $(\lambda,\lambda')$ match
(Lemma~\ref{plf:match}(a)), Lemma~\ref{par:casimir} gives
$\Omega\Phi=\bigl(\Omega'+2s_0\mathcal N-c\bigr)\Phi=\kappa_\lambda\Phi$, with $c$ the constant of
that lemma. The function $e_\lambda$ is integrable (Lemma~\ref{plf:A12}). By Lemma~\ref{par:rigid},
$\sigma_\tau(e_\lambda)\Phi=\Phi$. Finally the operators $\mathsf P$ and $\omega_\infty(1,y)$,
$y\in\auU(V)(E^+\otimes\R)$, commute, and $\xi\mapsto\pop_\xi$ is conjugate-linear and
$\auU(V)(\R)$-equivariant (\S\ref{po:ss-set}); so $\{\xi:\mathsf P\pop_\xi=\pop_\xi\}$ is an invariant
subspace of the irreducible representation $W_{\lambda_V}$ containing $\xi_0\ne0$, hence everything.

\emph{Step 2: the type of the descent.} By \eqref{po:eq-descent}, $D_\vartheta\Xi=\aubar F$ with
$F(g)=\sum_hc_h\,\theta\bigl(\omega_\mu(g,h)(\pop_{\Xi(h)}\otimes\Phi_\vartheta)\bigr)$, a finite sum
($c_h>0$). By Lemma~\ref{pz:growth}, $|F(gx)|\le C(g)\|x\|^{2s}$ for $x\in\auU(W)(E^+\otimes\R)$, and
$\Psi\mapsto\theta(\omega_\mu(g,h)(\Psi\otimes\Phi_\vartheta))$ is continuous on $\pzF_s$; by
Lemma~\ref{plf:A12} and Fubini,
\[
\int\aubar{e_W(x)}F(gx)\,dx=\sum_hc_h\,\theta\Bigl(\omega_\mu(g,h)\Bigl(\int\aubar{e_W(x)}\,
\omega_\infty(x,1)\pop_{\Xi(h)}\,dx\otimes\Phi_\vartheta\Bigr)\Bigr)=F(g),
\]
since $\int\aubar{e_W(x)}\omega_\infty(x,1)\,dx=\int e_W(x^{-1})\omega_\infty(x,1)\,dx=\mathsf P$ and
$\mathsf P\pop_\xi=\pop_\xi$ by Step~1. Taking complex conjugates, $R(e_W)D_\vartheta\Xi=D_\vartheta\Xi$.
The function $D_\vartheta\Xi$ is continuous and left invariant under $\auU(W)(E^+)$ (invariance of
$\theta$), and right invariant under $J_M$ (\S\ref{po:ss-set}). It is bounded: the truncation
argument at the end of the proof of Proposition~\ref{pz:types}(b) applies verbatim, with the
finite-dimensional space of forms of type $e_W$ and level $J_M$ (Lemma~\ref{plf:A12}). So
$D_\vartheta\Xi\in\poA(W_j)$.

\emph{Step 3: (L4).} If $\Xi$ is orthogonal to $\Theta_\vartheta(e_0\poA(W_j))$, then by (L2) and
Lemma~\ref{po:adj}(c) it is orthogonal to $\Theta_\vartheta y$ for all $y\in\poA(W_j)$ (decompose
$y=\varrho(e_0)y+(1-\varrho(e_0))y$; $\Theta_\vartheta\varrho(1-e_0)y=\varrho_V(1-e_0)\Theta_\vartheta y$
is orthogonal to $\varrho_V(e_0)\poA(V;\vartheta)\ni\Xi$). By \eqref{po:eq-descent},
$\int_{[\auU(W)]}y\,\aubar{D_\vartheta\Xi}=0$ for all $y\in\poA(W_j)$; take $y=D_\vartheta\Xi$, which is
allowed by Step~2. So $D_\vartheta\Xi=0$.
\end{proof}

\begin{remark}\label{plf:rem-L4}
Step~1 is the analogue of (A3) (Theorem~\ref{par:A3}) for the pair $(W_j,V_j)$, with the compact
group $\auU(V_j)(\R)$ on the column side; on that side rigidity is automatic, and only the row side
needs Lemma~\ref{par:rigid}, that is, the published results (D1)--(D3) of \S\ref{par:ss-rigid}. So
Proposition~\ref{plf:L4} is proved modulo exactly the inputs of Theorem~\ref{par:A3}: (D1)--(D3),
and (V$'$), proved in Theorem~\ref{pco:V} modulo (W1$_\infty$), (W2$_\infty$), (M), and (S) for
$j=m$.
\end{remark}

\subsubsection{Kernel scales for definite pairs}\label{plf:ss-KS}

The first factors are built from theta kernels of pairs of standard definite spaces with data at
the chosen places and at $\mathsf Q$. We state the integrality and size of such kernels in the
generality needed, and prove them from Theorem~\ref{pf:std}, Lemma~\ref{ps:conj} and the argument of
\S\ref{ps:ss-kernel}.

\emph{Setting.} Let $L=E^{l}$ and $N=E^{n_N}$ be standard definite spaces, $\mathsf W=\auHom(L,N)=
M_{n_N\times l}$, and let $a,b$ be unitary Hecke characters of $\auU(1)(\auA_{E^+})$ with
algebraic weights at the infinite places and tame ramification, at split places only. Let
$\omega_{a,b}$ be the Weil representation of $\auU(L)(\auA)\times\auU(N)(\auA)$ on $\psHS(\mathsf W)$
twisted by $a(\det g)^{-1}$ and $b(\det h)$, as in \eqref{po:eq-w5}. A \emph{modulus} is an ideal
$\mathfrak m$ of $\auO_E$ divisible only by split primes $w$ (with $\aubar w\nmid\mathfrak m$); for a
function $\Psi$ on $M_{n_N\times l}(\auO_E/\mathfrak m)$ let $f_\Psi\in\psHS_f(\mathsf W)$ be the
datum $x\mapsto\plone[x\in L_w]\,\Psi(a_w^{-1}x\bmod\mathfrak m)$ at the places $w\mid\mathfrak m$,
$x$ the matrix in the frames (model (W5)), and the vacuum elsewhere; here $a_w$ is the generator of
$\alpha_w\auO_w$ of Definition~\ref{pgx:conv} and $L_w=a_wM_{n_N\times l}(\auO_w)$, which is the
$w$-component $M_w$ of the self-dual lattice of Lemma~\ref{pgx:vac}. In the coordinates $T=a_w^{-1}x$ of
(N5) (Lemma~\ref{pgx:dict}(c)) this is $T\mapsto\plone[T\ \text{integral}]\,\Psi(T\bmod\mathfrak m)$. By Proposition~\ref{pgv:alpha}, $a_w=1$ and $L_w$ is the
lattice of integral matrices except at primes $w$ above $q_1$ dividing $\omega$, which can divide
$\mathfrak m$ only through $Q_\vartheta$; this is the datum $\Psi^\natural$ of \S\ref{pc:ss-data}. For a polynomial $Q$ in the variables
$z_{\tau,c,a}$ ($\tau\in\auPl$, $c\le n_N$, $a\le l$) let $\pop_Q\in\pzF$ be its polynomial vector,
and
\[
\plfK_{a,b}[\Psi,Q](g,h)=\theta\bigl(\omega_{a,b}(g,h)(\pop_Q\otimes f_\Psi)\bigr)\qquad(g\in\auU(L)(\auA_f),\
h\in\auU(N)(\auA_f)),
\]
the archimedean components being $1$. For $Q^\sigma$ we apply $\sigma$ to the coefficients of $Q$.

\begin{lemma}[Kernel scales]\label{plf:KS}
For all $l,n_N\ge1$ and $e\ge0$ there is $c=c_{l,n_N,e}\in\C^\times$ such that, for all $a,b$,
$\mathfrak m$, $\Psi$ and all $Q$ homogeneous of degree $e$:
\begin{enumerate}
\item[(KS1)] Let $q\ge1$ be an integer with $qM^2\in\mathfrak m$, let $\iota\colon\C\to\C_\ell$ with
$\ell\nmid q$, and suppose that the values of $\Psi$ and the coefficients of $Q$ are integral at
$\iota$ and that $g,h$ are integral with integral inverses at the places above $\ell$. Then
$|c\,\plfK_{a,b}[\Psi,Q](g,h)|_\iota\le1$.
\item[(KS2)] Suppose $l,n_N\le\Delta$, $e\le M$, $M^2\in\mathfrak m$, $|\sigma'(\Psi(x))|\le A$ for all
$x$ and all $\sigma'\in\auAut(\C)$, and $|Q^\sigma(z)|\le B$ for all $z$ with $\max_i|z_i|\le M^4$.
Then $|\sigma(c\,\plfK_{a,b}[\Psi,Q](g,h))|\le e^{x_1}AB$ for all $g,h$.
\end{enumerate}
\end{lemma}
\begin{proof}
\emph{Jets.} Write $\Theta_f$ for the function of \eqref{pk:eq-theta} attached to $f\in\psHS_f(\mathsf
W)$ (\S\ref{pk:ss-conj}, with $\mathsf C=[n_N]\times[l]$). The polynomial vectors are the
derivatives at $0$ of the coherent states $\psrho_\infty(z)1$ times $e^{|z|^2/2}$; hence, as in
\eqref{ps:eq-kernel}, there is a constant $\kappa_{\mathrm F}\ne0$, depending only on the
normalisation of the Fock model, with $\theta(\pop_Q\otimes f)=\kappa_{\mathrm F}^{\,e}\,
Q(\partial)[\Theta_f](0)$ for $Q$ homogeneous of degree $e$. For $s\in\C^\times$, $Q(\partial)[\Theta_f
(s\,\cdot)](0)=s^eQ(\partial)[\Theta_f](0)$.

\emph{The datum as a combination of translates of the vacuum.} At a place $w\mid\mathfrak m$, the
indicator of $x+\mathfrak mL_w$ ($L_w=a_wM_{n_N\times l}(\auO_w)$ as above, $x\in L_w$) is a root of unity times
$\psrho^\circ(x)\plone_{\mathfrak mL_w}$, and by Fourier inversion on the finite group
$\mathfrak m^{-1}L^\vee_w/L^\vee_w$ in the variables of $\aubar w$,
$\plone_{\mathfrak mL_w}=[L_w:\mathfrak mL_w]^{-1}\sum_y\zeta_y\,\psrho^\circ(y)\psvac_{L,w}$, a sum of
$[L_w:\mathfrak mL_w]$ terms with roots of unity $\zeta_y$ and translations $y$ with $\mathfrak my$
integral. Taking the product over $w\mid\mathfrak m$ and writing $\psvac=c_{\mathrm{vac}}^{-ln_N}
\psvac^\circ$ (\eqref{pk:eq-vac}), we get
\begin{equation}\label{plf:eq-vacexp}
c_{\mathrm{vac}}^{\,ln_N}\,\auNm(\mathfrak m)^{ln_N}f_\Psi=\sum_{x,y}\zeta_{x,y}\,\Psi(x)\,
\psrho^\circ(x+y)\psvac^\circ ,
\end{equation}
a sum of $\auNm(\mathfrak m)^{2ln_N}$ terms, $\zeta_{x,y}$ roots of unity, $\mathfrak m(x+y)$ integral
at the places of $\mathfrak m$ and $x+y=0$ elsewhere.
Here $\psvac_{L,w}$ is the vacuum, which is $\plone_{L_w}$ up to a scalar by Lemma~\ref{pgx:vac}
\emph{because} $L_w=a_wM_{n_N\times l}(\auO_w)$ is the $w$-component of the self-dual lattice (with
the lattice of integral matrices this would fail at a prime where $\alpha_w$ is not a unit),
$L_w^\vee$ is its dual $M_{\aubar w}$, $\Psi(x)$ stands for $\Psi(a_w^{-1}x\bmod\mathfrak m)$,
``$\mathfrak m(x+y)$ integral'' means $\mathfrak mx\in\mathfrak mL_w$, $\mathfrak my\in L^\vee_w$, and
$[L_w:\mathfrak mL_w]=\auNm(\mathfrak m_w)^{ln_N}$ does not depend on $a_w$. A prime $w\mid\mathfrak m$
with $a_w\ne1$ lies above some prime $p\mid q_1$, and $p\nmid M$ ($p_0,p_1\equiv1\pmod{q_1}$); in
(KS1) below $qM^2\in\mathfrak m\subseteq w$, so $p\mid q$ and $\ell\ne p$. Hence such $w$ do not lie above
$\ell$, and at the places above $\ell$ the translations are exactly as in the case $a_w=1$. (Such
primes can occur: they are split over $E^+$, Proposition~\ref{pgv:alpha}(b), of norm prime to $MN$, so
the definition of a twist allows them in $Q_\vartheta$.)

\emph{Definition of $c$.} Let $s_0$ be the scale of \S\ref{pk:ss-conj} and put
\[
c_{l,n_N,e}=c_{\mathrm{vac}}^{\,ln_N}\,M^{4fln_N}\,\bigl(dM^2s_0/\kappa_{\mathrm F}\bigr)^e .
\]

\emph{(KS1).} Since $qM^2\in\mathfrak m$, $\auNm(\mathfrak m)$ divides $(qM^2)^{2f}$, so
$M^{4fln_N}\auNm(\mathfrak m)^{-ln_N}$ is integral at $\ell\nmid q$; and at the places above $\ell$
(which do not divide $q$) $M^2(x+y)$ is integral. So by \eqref{plf:eq-vacexp}
$c_{\mathrm{vac}}^{ln_N}M^{4fln_N}f_\Psi$ is a standard datum of level $M^2$ at $\ell$ in the sense of
\S\ref{pf:ss-frames}; the values of $\Psi$ are integral and the coefficients are roots of unity.
By (S1) at the level $M^2$ (\S\ref{pf:ss-frames}), $\omega_{a,b}(g,h)$ preserves standard data of
level $M^2$ at $\ell$ up to the scalars $a(\det g)^{-1}b(\det h)$ and the unramified factors of
\eqref{po:eq-w5}, which are units at $\ell$ for $g,h$ integral with integral inverse above $\ell$
(values of algebraic Hecke characters at ideles that are units above $\ell$, and powers of residue
characteristics different from $\ell$). Let $f'$ be the resulting standard datum. Then
\[
c\,\plfK_{a,b}[\Psi,Q](g,h)=(\text{unit})\cdot Q(\partial)\bigl[\Theta_{f'}(d\cdot M^2s_0\,\cdot)\bigr](0).
\]
By Theorem~\ref{pf:std} with $\delta=M^2$ ($\aubar\delta=\delta$), $\Theta_{f'}(M^2s_0z)=\bigl(\prod_j
\theta(\mu_j(z))^{m_j}\bigr)^{\auNm\delta}\mathsf H(z)$ with $\mu_j$ linear forms with integral
coefficients and $\mathsf H$ an integral power series in the $y(\mu_j(z))$. We claim that
$\Theta_{f'}(dM^2s_0\,\cdot)$ is $\partial$-integral at every prime $\ell$, including $\ell=d$. For
$\ell\ne d$ this is the last assertion of Theorem~\ref{pf:std} with $c=d$. For $\ell=d$: $y$ is
$\partial$-integral at every prime (proof of (C2) in Proposition~\ref{pk:curve}), so $y(d\mu_j(z))$
and $\mathsf H(dz)$ are; and $\theta=\exp(-g)$ with $g(0)=g'(0)=0$, $g^{(n)}(0)=x^{(n-2)}(0)\in\pao$ for
$n\ge3$ (same proof) and $d^2g''(0)\in\pao$ (at $\ell=d$ the numbers $x(u-t)$ in the proof of (C2) need
not be integral, but $v(g''(0))>-2$ for the valuation with $v(d)=1$: proof of Lemma~\ref{pex:dint}); hence
$z\mapsto g(dz)$ has derivatives $d^ng^{(n)}(0)\in
\pao$ and vanishes at $0$, and $\theta(d\,\cdot)=\exp(-g(d\,\cdot))$ is $\partial$-integral, as are
products and compositions with linear forms with integral coefficients (\S\ref{pa:ss-curve}). This
proves the claim, and the jet $Q(\partial)[\cdot](0)$ of a $\partial$-integral series along a polynomial
with integral coefficients is integral.

\emph{(KS2).} We may assume $\sigma|_K=1$ (proof of Proposition~\ref{au:kernel}). By
\eqref{plf:eq-vacexp} and Step~1 of that proof, $c_{\mathrm{vac}}^{ln_N}\auNm(\mathfrak m)^{ln_N}
\omega_{a,b}(g,h)f_\Psi$ is a sum of $\auNm(\mathfrak m)^{2ln_N}\le M^{8fln_N}$ vectors
$\epsilon_{x,y}\Psi(x)\,\psrho^\circ(x')\omega(k)\psvac^\circ$ with $|\sigma\epsilon_{x,y}|=1$. For each
of them, Step~2 of that proof, with $\pi_M$ replaced by the rational integer $dM^2$, gives
\[
\sigma\bigl(Q(\partial)[\Theta_{f}(dM^2s_0\,\cdot)](0)\bigr)=\int Q^\sigma\bigl(dM^2t_\sigma\bar
y\bigr)e^{-|y|^2/2}\,\theta\bigl(\psrho_\infty(y)(1\otimes\Lambda_\sigma f)\bigr)\,d\mu(y),
\]
with $|t_\sigma|^2=\Omega^{-1}$ and $|\theta(\psrho_\infty(y)(1\otimes\Lambda_\sigma f))|\le
\exp(C_0(q_1ln_N+1)^2)$ by Lemma~\ref{ps:conj}(ii). Since $Q^\sigma$ is homogeneous of degree $e$,
$|Q^\sigma(z)|\le B\max(1,|z|/M^4)^e$. Here $|dM^2t_\sigma\bar y|/M^4\le\rho|y|$ with
$\rho=d\,\Omega^{-1/2}M^{-2}\le b^2\cdot x^{7}\cdot b^{-40}\le b^{-37}$ ($d\le b^2$ since $\sqrt d\le b$, $\Omega^{-1}\le
x^{13}$ by the definition of $\Omega$ and \eqref{eq:au-sizes}, $M\ge b^{20}$; Lemma~\ref{an:lem-sizes}).
With $D=fln_N$ and the normalisation of $d\mu$ of \S\ref{ps:ss-kernel},
$\int(1+(\rho|y|)^e)e^{-|y|^2/2}d\mu\le2^D\bigl(1+(\rho^2(2D+e))^{e/2}\bigr)\le2^{D+1}$, because
$\rho^2(2D+e)\le b^{-74}(4\Delta^2f+M)\le1$ (as $M\le b^{22}$). Collecting: by the definition of $c$ and
\eqref{plf:eq-vacexp}, $c\,\plfK_{a,b}[\Psi,Q](g,h)=M^{4fln_N}\auNm(\mathfrak m)^{-ln_N}\sum_{x,y}
\epsilon_{x,y}\Psi(x)\,Q(\partial)[\Theta_{f_{x,y}}(dM^2s_0\,\cdot)](0)$ with $f_{x,y}=\psrho^\circ(x')
\omega(k)\psvac^\circ$ (the factor $c_{\mathrm{vac}}^{ln_N}$ is absorbed by writing the datum through
$\psvac^\circ$), and $\auNm(\mathfrak m)$ is a positive integer, so
\[
|\sigma(c\,\plfK_{a,b}[\Psi,Q](g,h))|\le AB\cdot2^{D+1}\,M^{12fln_N}\,
\exp\bigl(C_0(q_1ln_N+1)^2\bigr).
\]
With $l,n_N\le\Delta=2r^8\le2\cdot3^8x^{40}$, $f,q_1\le x^{12}$, $\log M\le23\log b\le x$ and
$C_0\le\log\log b$ (Lemma~\ref{pz:C0}), the logarithm of the factor after $AB$ is
at most $x^{200}\le x_1$ by \eqref{eq:au-scales}.
\end{proof}

\begin{remark}\label{plf:rem-KS}
(a) (KS1) is the integrality statement (K1) of Lemma~\ref{pa:kernel} for definite pairs with
arbitrary data at split places, at every prime $\ell\nmid q$, \emph{including} $\ell\mid2q_1M$
(Theorem~\ref{pf:std}), $\ell=N$ (the data are unramified at $N$), and $\ell=d$ (the extra factor
$d^e$ in the scale). The factor $d^e$ is harmless for the size: it is part of $\rho$ above.
(b) The size bound (KS2) is far from optimal; only its form $e^{x_1}AB$ is used.
\end{remark}

\subsubsection{The set of choices}\label{plf:ss-ff}

\emph{Spaces.} Besides $V''$ (rank $N''=r+9m+1$) we use the standard definite
spaces $E^m$, $E^{3m}$ and $W^\flat=E^{6m+1}$; their levels are $K^\circ$ (the stabiliser of the
standard lattice) or $K(\pi_M^2)$.

\emph{The choices.} A \emph{choice} $i=(s_i,h^\flat_i,\mathsf r_i)$ consists of a matrix $s_i\in
M_{3m\times N''}(\auO_E/\pi_M^2)$, a class $h^\flat_i\in\Gamma^\flat\backslash\auU(W^\flat)(\auA_f)/
K^\flat(\pi_M^2)$ (with a fixed representative), and, for each $\tau\in\auPl$ and each of the
$\Delta'=r^8-\frac{t+1}2$ columns of height $r_+$ of $\mu^+_\tau$, a list of $2m$ indices in
$\{1,\dots,6m+1\}$. For each class $h^\flat$ fix $Y=Y_{h^\flat}\in\auGL_{6m+1}(E)$ as in
Lemma~\ref{au:shortbasis} for $g=h^\flat$ ($h^{\flat,-1}Y$ integral, all conjugates of all entries at
most $C_{\mathrm{sb}}$).

\begin{lemma}\label{plf:I}
The set $I$ of choices has at most $\exp(x^{54})\le e^b$ elements.
\end{lemma}
\begin{proof}
$\#(\auO_E/\pi_M^2)=M^{2f}$ ($\pi_M\in\auO_K$ has norm $M$ to $\Q$, and $[E:K]=f$), so there are
$M^{6fmN''}$ matrices $s_i$; by Lemma~\ref{plf:classes} there are at most $\exp(2x^{26})$ classes
$h^\flat$; and at most $(6m+1)^{2mf\Delta'}$ lists. By \eqref{eq:au-sizes}, \eqref{eq:au-scales} and
$m\le2\sqrt x$, $r\le3x^5$: $6fmN''\log M\le6x^{12}\cdot2\sqrt x\cdot4x^5\cdot x\le x^{19}$ and
$2mf\Delta'\log(6m+1)\le4\sqrt x\,x^{12}\,3^8x^{40}\log(13\sqrt x)\le x^{53}$.
\end{proof}

The first factors attached to the choices are those of Definition~\ref{ptf:def-ff}; their
integrality and size (Lemma~\ref{po:ff}) are proved in \S\ref{pfq:sec} (Theorem~\ref{pfq:ff}).




%% ===== proofs/C-periods.tex =====
% proofs/C-periods.tex
% Proposition po:periods of proofs/C-onelift.tex: periods of false lifts vanish.
% To be \input after C-onelift.tex and C-lifts.tex. Labels of this file start with ppe:.
% Only amsmath/amssymb/amsthm are used.
\providecommand{\auO}{\mathcal O}
\providecommand{\auU}{\mathrm U}
\providecommand{\auGL}{\mathrm{GL}}
\providecommand{\auHom}{\operatorname{Hom}}
\providecommand{\auA}{\mathbb A}
\providecommand{\auF}{\mathbb F}
\providecommand{\auPl}{\Phi}
\providecommand{\auG}{\overline G}
\providecommand{\aubar}[1]{\overline{#1}}
\providecommand{\plone}{\mathbf 1}
\providecommand{\psrho}{\varrho}
\providecommand{\poA}{\mathcal A}
\providecommand{\poZ}{\mathcal Z}
\providecommand{\pop}{\mathsf p}
\providecommand{\pog}{\mathsf g}
% macros of this file (prefix \ppe)
\providecommand{\ppeD}{\mathcal D}
\providecommand{\ppeVan}{\mathrm{Van}}
\providecommand{\ppeS}{\mathcal S}
\providecommand{\ppeH}{\mathbb H}
\providecommand{\ppeN}{\mathsf N}
\providecommand{\pexK}{K^{(z)}}

\subsection{Periods of false lifts}\label{ppe:sec}

This subsection proves Proposition~\ref{po:periods}. We keep the notation of \S\ref{po:sec}:
$1\le j\le m$ is fixed, $W=W_j$ (rank $r$, signature $(r_++1,r_-)$-type as in
\S\ref{ss:au-spaces}), $V=V_j$ the positive definite standard space of rank $n=r+t=6r$,
$\vartheta$ a twist, $\mathsf S_\vartheta=\mathsf M\cup\mathsf Q$ the places of the level, $\mu=\mu_j^0
\vartheta$ with tame parts $\psi_w$ ($w\in\mathsf S_\vartheta$), $\omega_\mu$ the Weil
representation \eqref{po:eq-w5}, $\Phi_\vartheta$ the block datum, $K_\vartheta$, $\chi_\vartheta$,
$\poA(V;\vartheta)$, the descent $D_\vartheta$ \eqref{po:eq-descent}, $\varrho_V$, $e_0$, and the
splitting $V=V''\perp V^\circ$ of \S\ref{po:ss-glue} with $N''=r+9m+1$, $n^\circ=t-9m-1$. The
argument is the classical one of the theta tower (Rallis~\cite[Chapter~I, \S1]{RallisHoweDuality}; see also~\cite{RallisInjectivity}): the descents of a
false lift vanish to $W$, then to every $W\oplus\ppeH^p$ with $2r+2p+3<n$, and a Fourier
coefficient of the descent to the split space $\ppeH^{N''}$ is the period over $\auU(V^\circ)$.
The new ingredient, which replaces the use of poles of $L$-functions in the classical argument,
is that a false lift satisfies \emph{one} Hecke equation at one split place, at the smallest root
of a lift; this forces bounded descents to vanish.

\subsubsection{Set-up}\label{ppe:ss-setup}

\emph{Hermitian spaces and data.} For a non-degenerate hermitian space $X$ over $E$ of rank
$m_X$ (signatures at the infinite places arbitrary) we use the Weil representation
$\omega^X_a$ of $\auU(X)(\auA)\times\auU(V)(\auA)$ on $\ppeS(X)=\mathcal H(\auHom(X,V))$, in the
normalisation of \eqref{po:eq-w5}: $\auU(V)$ acts through the splitting given by powers of
$\upsilon$ multiplied by the character $a(\det h)$, where $a$ is a character of
$\auU(1)(E^+)\backslash\auU(1)(\auA_{E^+})$, and $\auU(X)$ acts through the splitting attached to
$X$ (an extra character $b(\det g)^{-1}$ on $\auU(X)$ only multiplies descents by $b(\det g)^{-1}$
and is never needed). For $X=W$ and $a=1$, $\omega^W_1$ is $\omega_\mu$ up to the character
$\mu(\det g)^{-1}$ of $\auU(W)$. Let $\theta$ be the theta functional. Since $V$ is definite,
$[\auU(V)]=\auU(V)(E^+)\backslash\auU(V)(\auA)$ is compact. For a continuous function $Y$ on
$[\auU(V)]$ and $f\in\ppeS(X)$ put
\[
\ppeD^X_a(f)Y(g)=\int_{[\auU(V)]}\theta\bigl(\omega^X_a(g,h)f\bigr)\,Y(h)\,dh\qquad(g\in
\auU(X)(\auA)),
\]
the \emph{descent} of $Y$ to $X$ by $f$; the integral converges absolutely because $[\auU(V)]$
is compact and $h\mapsto\theta(\omega(g,h)f)$ is continuous. Let $\Sigma=\mathsf S_\vartheta\cup
\aubar{\mathsf S_\vartheta}$ and let $\ppeS_\Sigma(X)\subseteq\ppeS(X)$ be the closed subspace of
the data that are the vacuum at every finite place outside $\Sigma$ (arbitrary at the infinite
places and at $\Sigma$). We say that \emph{the descents of $Y$ to $X$ vanish}, and write
$\ppeVan(X,a,Y)$, if $\ppeD^X_a(f)Y=0$ for every $f\in\ppeS_\Sigma(X)$.

For $p\ge0$ let $\ppeH^p$ be the sum of $p$ hyperbolic planes, with an isotropic frame
$e_1,\dots,e_p,f_1,\dots,f_p$, $\langle e_i,f_k\rangle=\delta_{ik}$, adapted (orthonormal up to the
standard scaling at the infinite places and spanning a self-dual lattice at the finite places),
and $W_p=W\oplus\ppeH^p$. Put $\upsilon_1=\upsilon|_{\auU(1)}$ (the restriction of the splitting
character to the norm-one ideles).

\emph{Matrix coefficients.} For $\Xi\in\poA(V;\vartheta)$ and $\zeta\in W_{\lambda_V}$ let
\[
c_\zeta\Xi(h)=\bigl\langle\zeta,\ \rho_{\lambda_V}(h_\infty)^{-1}\Xi(h_f)\bigr\rangle,\qquad
Y_\zeta=\aubar{c_\zeta\Xi}\qquad(h=h_\infty h_f\in\auU(V)(\auA)),
\]
with the form of $W_{\lambda_V}$ conjugate-linear in the first variable (as in
\eqref{po:eq-delta}). Since $\Xi(\gamma h_f)=\gamma\,\Xi(h_f)$, with $\gamma$ acting through its
archimedean components, $c_\zeta\Xi$ is left invariant under $\auU(V)(E^+)$; it is right invariant
under the elements of $K_\vartheta$ on which $\chi_\vartheta$ is trivial, in particular under
$K'_\vartheta=\{k\in K^\circ:\ k_w=1\ (w\in\Sigma)\}$, and it is continuous. If $T$ is a Hecke
operator at a place $w\notin\Sigma$ (sum of right translations by elements of $\auU(V)(\auA_f)$
supported at $w$), then $c_\zeta(T\Xi)=T(c_\zeta\Xi)$ and $Y_\zeta$ transforms in the same way
with the complex conjugate coefficients. For $x\in\auU(V)(E^+\otimes\R)$,
$c_\zeta\Xi(hx)=c_{\rho(x)\zeta}\Xi(h)$.

\begin{theorem}[Proposition~\ref{po:periods}]\label{ppe:main}
Let $w$ be off the level and $\Xi\in\mathcal F_w(e_0,\poZ_\vartheta)$ a false lift. Then for every
$s$, every $h\in\auU(V'')(\auA_f)$, $u''\in W''$ and $g\in\auU(W)(\auA)$,
\[
\sum_{h^\circ\in\Gamma^\circ\backslash\auU(V^\circ)(\auA_f)/J^\circ}\bigl\langle u''\odot
K^\circ(g,h^\circ),\ \Xi\bigl(\operatorname{diag}(h,h^\circ)\pog_s\bigr)\bigr\rangle=0 .
\]
The same holds with $J^\circ$ replaced by any open compact subgroup of $J^\circ$, and for every
$\Xi\in\poA(V;\vartheta)$ that satisfies (F1)--(F3) of Proposition~\ref{ppe:tower} below.
\end{theorem}

The pairing is that of \eqref{po:eq-delta}, conjugate-linear in the first variable; this is the
convention under which $\langle u^\circ,K^\circ(g,h^\circ)\rangle=\theta^\circ(g,h^\circ;u^\circ)$
is conjugate-linear in $u^\circ$, as stated in \S\ref{po:ss-glue}. (If the other convention is
used, both sides of the identity are complex conjugated and nothing changes.)

\emph{Plan.} A false lift has three properties (\S\ref{ppe:ss-three}): (F1) $D_\vartheta\Xi=0$;
(F2) $\Xi$ is killed by the hyperplane sums at the chosen places; (F3) $P_w(X_w)\Xi=0$.
From (F1), (F2) and $\psi_w\ne1$ on $\mathsf Q$ the descents of every $Y_\zeta$ to $W$ vanish
(\S\ref{ppe:ss-bottom}); by (F3) they vanish to $W_p$ for $2r+2p+3<n$
(\S\ref{ppe:ss-planes}); and the vanishing of descents to $\ppeH^{N''}$ kills the periods over
$\auU(V^\circ)$ (\S\ref{ppe:ss-periods}).

\subsubsection{The three properties of a false lift}\label{ppe:ss-three}

For a chosen place $w\in\mathsf M$ and a row $\lambda\in\auF_w^r\smallsetminus\{0\}$ let
$\ppeN_\lambda=\{s\in\auGL_r(\auF_w):\ s=1+u\lambda\text{ for some column }u\}$; it is a subgroup
($(1+u\lambda)(1+u'\lambda)=1+(u+u'+u(\lambda u'))\lambda$, and the inverse of $1+u\lambda$ is
$1-(1+\lambda u)^{-1}u\lambda$). For a character $\psi$ of $\auF_w^\times$ the \emph{hyperplane
sum} is
\[
\eta_{\psi,\lambda}=\sum_{s\in\ppeN_\lambda}\psi(\det s)^{-1}\,s\ \in\C[\auGL_r(\auF_w)],
\]
viewed in $\C[\auG]$ through the factor at $w$. Note $\det(1+u\lambda)=1+\lambda u$, so
$\det\colon\ppeN_\lambda\to\auF_w^\times$ is surjective.

\begin{lemma}[Characters of label idempotents on hyperplane groups]\label{ppe:green}
Let $\auF$ be a finite field, $c$ a label of degree $r$ over $\auF$, $e_c$ the central primitive
idempotent of $\C[\auGL_r(\auF)]$ of the character $\Xi_c$, $\psi$ a character of $\auF^\times$
and $x_\psi\in\auF^\times$ its name in Green's parametrisation (the element with
$\psi=\theta_{x_\psi}|_{\auF^\times}$ for the fixed injective character $\theta$ of
$\aubar{\auF}^\times$ used to define the $\Xi_c$). If $x_\psi$ is not a root of $c$, then
$\eta_{\psi,\lambda}\,e_c=0$ for every $\lambda\ne0$.
\end{lemma}
\begin{proof}
Let $\pi_c$ be an irreducible representation with character $\Xi_c$. Since $e_c$ is central and
$\C[\auGL_r(\auF)]e_c\cong\operatorname{End}(\pi_c)$, it suffices to show $\pi_c(\eta_{\psi,
\lambda})=0$. Now $\#\ppeN_\lambda^{-1}\pi_c(\eta_{\psi,\lambda})$ is the projection of $\pi_c$
onto the subspace where $\ppeN_\lambda$ acts by $\psi\circ\det$, of dimension
$\langle\Xi_c|_{\ppeN_\lambda},\psi\circ\det\rangle_{\ppeN_\lambda}$. Choose a basis
$v_1,\dots,v_r$ with $\lambda v_k=0$ for $k\ge2$ and $\lambda v_1=1$. In it $\ppeN_\lambda$ is the
group of matrices $\begin{pmatrix}a&0\\ y&1_{r-1}\end{pmatrix}$ ($a\in\auF^\times$,
$y\in\auF^{r-1}$), the semidirect product of $\auGL_1$ (embedded as $\operatorname{diag}(a,1)$) and
the unipotent radical $U$ of the parabolic subgroup $P$ of type $(1,r-1)$ containing it. Hence
\[
\langle\Xi_c|_{\ppeN_\lambda},\psi\circ\det\rangle=\frac1{q-1}\sum_{a}\psi(a)^{-1}\,
\frac1{\#U}\sum_{y\in U}\Xi_c\bigl(\operatorname{diag}(a,1)y\bigr)
=\bigl\langle\operatorname{r}_P\Xi_c\,,\ \psi\otimes\plone\bigr\rangle_{\auGL_1\times\{1\}},
\]
where $\operatorname{r}_P\Xi_c$ is the Harish-Chandra restriction (the character of the
$U$-invariants) to the Levi subgroup $\auGL_1\times\auGL_{r-1}$, evaluated on $\auGL_1\times\{1\}$.
By the restriction formula for Green's characters (the comultiplication of the Hopf algebra of
characters of the groups $\auGL_k(\auF)$, which is the Harish-Chandra restriction and is a ring homomorphism for which the cuspidal characters are primitive, \cite[\S\S9.1--9.2, pp.~110--111]{ZelevinskyLNM}; equivalently
the Mackey formula applied to Green's description of $\Xi_c$ as a product of the cuspidal
characters of the irreducible factors of $c$, \cite[p.~403 and Thm.~13]{Green1955}), $\operatorname{r}_P\Xi_c=
\sum_{\alpha}\theta_\alpha\otimes\Xi_{c/(X-\alpha)}$, the sum over the roots $\alpha\in\auF^\times$
of $c$ (a factor of degree $\ge2$ of $c$ is cuspidal, by property (G) of \S\ref{pm:ss-mult}, and contributes nothing to a restriction
with a $\auGL_1$-block). Pairing with $\psi\otimes\plone$ on $\auGL_1\times\{1\}$ gives
$\sum_\alpha\langle\theta_\alpha,\psi\rangle\,\Xi_{c/(X-\alpha)}(1)$, and
$\langle\theta_\alpha,\psi\rangle=0$ unless $\alpha=x_\psi$. So the dimension is $0$.
\end{proof}

\begin{lemma}[The three properties]\label{ppe:three}
Let $w$ be off the level and $\Xi\in\mathcal F_w(e_0,\poZ_\vartheta)$. Then
\begin{enumerate}
\item[(F1)] $D_\vartheta\Xi=0$;
\item[(F2)] for every chosen place $v\in\mathsf M$ and every $\lambda\in\auF_v^r\smallsetminus
\{0\}$, $\varrho_V(\eta_{\psi_v,\lambda})\Xi=0$;
\item[(F3)] $P_w(X_w)\Xi=0$, with $X_w$ as in Lemma~\textup{\ref{po:facts}(L3)}; $|X_w|=
q^{-r/2-(t-1)}$, $q=\#\auF_w$.
\end{enumerate}
\end{lemma}
\begin{proof}
(F1) By definition $\Xi\in\varrho_V(e_0)\poA(V;\vartheta)$ is orthogonal to $\poZ_\vartheta=
\Theta_\vartheta(e_0\poA(W_j))$; Proposition~\ref{plf:L4} gives $D_\vartheta\Xi=0$. (F3) is part of
the definition of $\mathcal F_w$. (F2) The tame part of $\mu=\mu^0_j\vartheta$ at $v\in\mathsf M$ is
that of $\mu^0_j$ ($\vartheta$ is unramified at $\mathsf M$), namely $\omega_v^{\,j}$
(Definition~\ref{pfq:def-chars}, Proposition~\ref{pfq:tame}(T3)), a power of the character $\omega_v$
of order $q_2$; so
$\psi_v^{q_2}=1$ and its name $x_{\psi_v}$ satisfies $x_{\psi_v}^{q_2}=1$ (the naming map is a
group isomorphism). Write $\Xi=\varrho_V(e_0)\Xi$ and $e_0=\sum_{(c_u)\text{ good}}\bigotimes_u
e_{c_u}$. For each good family, $c_v$ has no root in $\auF_v$ that is a $q_2$-th root of unity, so
$\eta_{\psi_v,\lambda}e_{c_v}=0$ by Lemma~\ref{ppe:green}; since $\eta_{\psi_v,\lambda}$ lives in
the factor at $v$, $\eta_{\psi_v,\lambda}\bigotimes_ue_{c_u}=0$. Hence $\eta_{\psi_v,\lambda}e_0=0$
and $\varrho_V(\eta_{\psi_v,\lambda})\Xi=\varrho_V(\eta_{\psi_v,\lambda}e_0)\Xi=0$.
\end{proof}

\begin{remark}\label{ppe:rem-T}
The only property of the tame parts used in (F2) is $\psi_v^{q_2}=1$ at the chosen places, together
with the definition of good labels, which excludes the $q_2$-th roots of unity for this reason. That
the characters $\mu^0_j$ have tame part $\omega_v^j$, with $\omega_v^{q_2}=1$, at every chosen place is
Proposition~\ref{pfq:tame}(T3); their existence is Lemma~\ref{pfq:mu} (the compatibility condition on
the roots of unity of $E$ is checked in Definition~\ref{pfq:def-chars} and Proposition~\ref{pfq:tame}(T1)).
\end{remark}

\begin{lemma}[The descent and the matrix coefficients]\label{ppe:normal}
Let $\xi_0\in W_{\lambda_V}$ be a highest weight vector and $\Xi\in\poA(V;\vartheta)$. Then
$D_\vartheta\Xi=0$ if and only if $\ppeD^W_1(\pop_{\xi_0}\otimes\Phi_\vartheta)Y_{\xi_0}=0$, where
$\pop_{\xi_0}$ is the polynomial \eqref{po:eq-poly}.
\end{lemma}
\begin{proof}
Let $F(g)=\ppeD^W_1(\pop_{\xi_0}\otimes\Phi_\vartheta)Y_{\xi_0}(g)$. Integrate first over
$\auU(V)(E^+\otimes\R)$, which is compact; with $dh_\infty$ of total mass one and $c>0$ the
volume constant,
\[
F(g)=c\sum_{h_f}\frac1{\#(\Gamma\cap h_fK_\vartheta h_f^{-1})}\int_{\auU(V)(\R)}\theta\bigl(
\omega(g,h_f)\omega(1,h_\infty)(\pop_{\xi_0}\otimes\Phi_\vartheta)\bigr)\,\aubar{\bigl\langle
\rho(h_\infty)\xi_0,\Xi(h_f)\bigr\rangle}\,dh_\infty
\]
(the integrand is invariant under $K_\vartheta$ by the definition of $\chi_\vartheta$, and the sum
runs over $\Gamma\backslash\auU(V)(\auA_f)/K_\vartheta$). By the equivariance of
$\xi\mapsto\pop_\xi$ (\S\ref{po:ss-set}), $\omega(1,h_\infty)\pop_{\xi_0}=\pop_{\rho(h_\infty)
\xi_0}$, and $\xi\mapsto\theta(\omega(g,h_f)(\pop_\xi\otimes\Phi_\vartheta))=k_{\Phi_\vartheta}
(g,h_f;\xi)$ is conjugate-linear, so $k_{\Phi_\vartheta}(g,h_f;\xi)=\langle\xi,A\rangle$ for a
unique $A=A(g,h_f)\in W_{\lambda_V}$. Also $c_{\xi_0}\Xi(h_\infty h_f)=\langle\rho(h_\infty)\xi_0,
B\rangle$ with $B=\Xi(h_f)$, $\rho$ being unitary. Schur orthogonality for the irreducible unitary
representation $\rho=\rho_{\lambda_V}$ of the compact group $\auU(V)(\R)$ gives
\[
\int_{\auU(V)(\R)}\langle\rho(h)\xi_0,A\rangle\,\aubar{\langle\rho(h)\xi_0,B\rangle}\,dh=
\frac{\|\xi_0\|^2}{\dim W_{\lambda_V}}\,\langle B,A\rangle .
\]
Hence $F(g)=c\,\|\xi_0\|^2(\dim W_{\lambda_V})^{-1}\sum_{h_f}\#(\Gamma\cap h_fK_\vartheta
h_f^{-1})^{-1}k_{\Phi_\vartheta}(g,h_f;\Xi(h_f))$; by \eqref{po:eq-descent} the sum is the complex
conjugate of $D_\vartheta\Xi(g)$, up to the factor $\mu(\det g)$ that relates $\omega^W_1$ and
$\omega_\mu$. So $F=0$ if and only if $D_\vartheta\Xi=0$.
\end{proof}

\subsubsection{The bottom of the tower}\label{ppe:ss-bottom}

At a place $u\in\mathsf S_\vartheta$ we use the coordinates $T\in M_{n\times r}(E_u)$ of (N5) on
$\auHom(W,V)\otimes_{E^+}E^+_{u}$; the local datum at $u$ is a Schwartz--Bruhat function of $T$,
and by \eqref{po:eq-w5} an integral element $k$ of $\auU(V)(E^+_u)=\auGL_n(E_u)$ acts by
$\Psi\mapsto\Psi(k^{-1}T)$ and an integral element $g$ of $\auU(W)(E^+_u)=\auGL_r(E_u)$ by
$\Psi\mapsto\Psi(Tg)$ times $\chi_u(\det g)^{-1}$ (the unramified factors are trivial on integral
elements). For $T$ we write $T'$ for the $r\times r$ block of the rows of $\iota_0$ and $T''$ for
the other rows. For a Schwartz--Bruhat function $\Psi$ on $\prod_{u\in\mathsf S_\vartheta}
M_{n\times r}(E_u)$ and a polynomial Fock vector $v$ at the infinite places let
$f(v,\Psi)=v\otimes\Psi\otimes(\text{vacuum elsewhere})\in\ppeS_\Sigma(W)$.

\begin{lemma}[Special data suffice]\label{ppe:special}
Let $X$ be a hermitian space and $Y$ continuous on $[\auU(V)]$. If $\ppeD^X_a(f(v,\Psi))Y=0$ for
all polynomial Fock vectors $v$ and all $\Psi$ that are indicator functions of cosets
$z+\varpi^i M_{n\times m_X}(\auO_{\mathsf S_\vartheta})$ ($z$ arbitrary, $i\ge0$, $\varpi$ a
fixed element with $\varpi\auO_u$ the maximal ideal at each $u\in\mathsf S_\vartheta$), then
$\ppeVan(X,a,Y)$.
\end{lemma}
\begin{proof}
For fixed $g$ the map $f\mapsto\ppeD^X_a(f)Y(g)$ is a continuous linear functional on
$\ppeS(X)$ for its Schwartz (Fr\'echet) topology: $|\theta(\omega(g,h)f)|\le C(g)\,p(f)$, $p$ a
continuous seminorm of the Schwartz topology, uniformly for $h$ in a compact fundamental
domain (moderate growth of theta functions, Lemma~\ref{ps:weil}). The space $\ppeS_\Sigma(X)$ is
the completed tensor product of the archimedean Schwartz space and of the spaces of Schwartz--Bruhat
functions below; in the archimedean factor the polynomial Fock vectors correspond, in the
Schr\"odinger model, to polynomials times the Gaussian, which are dense in the Schwartz space for its
Fr\'echet topology (Hermite expansion). The finite factors are the spaces of Schwartz--Bruhat functions at the places of $\mathsf S_\vartheta$
(by (N5) the local Weil representation at the split place below $u$ is realised on functions of
$T\in M_{n\times m_X}(E_u)$, the place $\aubar u$ carrying no further variable), and every
Schwartz--Bruhat function is a finite linear combination of indicators of such cosets. Hence the
functional vanishes on a dense subspace, so everywhere.
\end{proof}

\begin{lemma}[Harmonic generation]\label{ppe:harm}
Let $\Xi\in\poA(V;\vartheta)$ satisfy (F1) and let $\xi_0$ be a highest weight vector. Then
$\ppeD^W_1(f(v,\Phi_\vartheta))Y_{\xi_0}=0$ for every polynomial Fock vector $v$ (here
$\Phi_\vartheta$ denotes the finite part of the block datum).
\end{lemma}
\begin{proof}
Write $L(v)(g)=\ppeD^W_1(f(v,\Phi_\vartheta))Y_{\xi_0}(g)$. (a) For $\xi\in W_{\lambda_V}$,
$L(\pop_\xi)=0$: the computation of the proof of Lemma~\ref{ppe:normal} with $\pop_\xi$ in place of
$\pop_{\xi_0}$ gives $L(\pop_\xi)(g)=c\,\langle\xi_0,\xi\rangle^{\ast}(\dim W_{\lambda_V})^{-1}
\aubar{D_\vartheta\Xi(g)}\mu(\det g)$ (Schur orthogonality with $\langle\xi_0,\xi\rangle$ in place
of $\|\xi_0\|^2$, the star denoting the appropriate complex conjugate), which is $0$ by (F1).
(b) For $x$ in the Lie algebra $\mathfrak g_W$ of $\auU(W)(E^+\otimes\R)$ and $v$ smooth,
$L(\omega(x,1)v)=R_xL(v)$, the derivative of the right translate; this is \eqref{po:eq-equiv}
at the infinite places (differentiate $\theta(\omega(gk,h)f)$ in $k$). (c) By (b) and
Schur orthogonality over the compact group $\auU(V)(\R)$, $L(v)=L(e_{\lambda}v)$, where $e_\lambda$
is the projection of the Fock space onto its $\rho_{\lambda_V}^\vee$-isotypic subspace for
$\omega(1,\cdot)$ (the integral over $h_\infty$ of $\omega(1,h_\infty)v$ against the coefficient
$\aubar{c_{\xi_0}\Xi}$ only sees this component). (d) By the theory of joint harmonics of
Kashiwara--Vergne and Howe (\cite[Ch.~III, Lemma~5.3 and Corollary~5.4]{KashiwaraVergne}; \cite[Proposition~3.1, (3.10)(b)]{HoweJAMS}), the polynomial vectors of that isotypic subspace are
$U(\mathfrak g_{W,\C})\cdot\{\pop_\xi:\xi\in W_{\lambda_V}\}$: the space of joint harmonics of
type $\lambda_V$ is $\{\pop_\xi\}$ (this is the choice of $\pop_\xi$ in \S\ref{po:ss-set}, see
Lemma~\ref{plf:harm}), and the isotypic component of the polynomial Fock space is generated by
its harmonics under $U(\mathfrak g_{W,\C})$. By (a), (b), (c), $L(v)=0$ for every polynomial $v$.
\end{proof}

\begin{definition}\label{ppe:def-ratio}
Let $u\in\mathsf S_\vartheta$ and $T\in M_{n\times r}(E_u)$. $T$ has an \emph{integral block
ratio} at $u$ if $T'$ is invertible, $T''T'^{-1}$ is integral, and moreover $T''T'^{-1}\equiv0
\bmod u$ if $u\in\mathsf M$. An \emph{annihilating row} for $T$ at $u$ is an integral row
$l=(\lambda,\lambda')$ ($\lambda$ indexed by the rows of $\iota_0$) with $lT=0$, $\lambda\not
\equiv0\bmod u$, and $\lambda'\equiv0\bmod u$ if $u\in\mathsf Q$.
\end{definition}

\begin{lemma}\label{ppe:row}
If $T$ has no integral block ratio at $u$, it has an annihilating row at $u$.
\end{lemma}
\begin{proof}
Let $|\cdot|$ be the absolute value at $u$. It suffices to find a row $\mu\ne0$ with $\mu T=0$ and
an index $b$ of $\iota_0$ with $|\mu_b|=\max_a|\mu_a|$, and $|\mu_a|<|\mu_b|$ for $a$ outside
$\iota_0$ when $u\in\mathsf Q$: then $l=\mu_b^{-1}\mu$ is an annihilating row. If $T'$ is
singular take $\mu=(\lambda,0)$ with $\lambda T'=0$, $\lambda\ne0$. Otherwise let $y=T''T'^{-1}$;
by hypothesis some entry $y_{ab}$ has $|y_{ab}|>1$, or $u\in\mathsf M$ and $|y_{ab}|\ge1$. Take
$\mu=(y_a,-e_a)$ (the row $a$ of $y$ on $\iota_0$, and $-1$ at $a$); then $\mu T=y_aT'-T''_a=0$,
and an index $b$ maximising $|y_{ab}|$ works, since $|\mu_a|=1$ and the other entries outside
$\iota_0$ vanish.
\end{proof}

\begin{proposition}[Bottom of the tower]\label{ppe:bottom}
Let $\Xi\in\poA(V;\vartheta)$ satisfy (F1) and (F2), and let $\xi_0$ be a highest weight vector.
Then $\ppeVan(W,1,Y_{\xi_0})$.
\end{proposition}
\begin{proof}
For $f\in\ppeS_\Sigma(W)$ put $L_\Xi(f)=\ppeD^W_1(f)Y_{\xi_0}$ (a function of $g$; we suppress $g$
and $v$). Two rules: if $k\in\auU(V)(\auA_f)$ then, substituting $h\mapsto hk$,
\begin{equation}\label{ppe:eq-move}
L_{R_k\Xi}(f)=L_\Xi\bigl(\omega(1,k^{-1})f\bigr),
\end{equation}
because $Y_{\xi_0}$ for $R_k\Xi$ is the right translate of $Y_{\xi_0}$ for $\Xi$; and for $k\in
K_\vartheta$, $R_k\Xi=\chi_\vartheta(k)^{-1}\Xi$ (definition of $\poA(V;\vartheta)=
\poA_{\lambda_V}(K_\vartheta,\chi_\vartheta^{-1})$, \S\ref{po:ss-set}) and $Y_{\xi_0}$ is
conjugate-linear in $\Xi$, so $L_{R_k\Xi}=\chi_\vartheta(k)L_\Xi$ ($\chi_\vartheta$ is unitary) and
\eqref{ppe:eq-move} gives $L_\Xi(\omega(1,k)f)=\chi_\vartheta(k)^{-1}L_\Xi(f)$.
By Lemma~\ref{ppe:special} it suffices to show $L_\Xi(f(v,\plone_C))=0$ for $C=z+\varpi^i
M_{n\times r}(\auO_{\mathsf S_\vartheta})$ and $i$ large (a coset for small $i$ is a finite union
of cosets for large $i$). Fix $z$.

\emph{Case A: $z_u$ has an integral block ratio at every $u\in\mathsf S_\vartheta$.} Let $A\in
\auU(W)(\auA_f)$ be $1$ outside $\mathsf S_\vartheta\cup\aubar{\mathsf S_\vartheta}$ with
$A_u=(z_u')^{-1}$ ($u\in\mathsf S_\vartheta$), and $\Psi_0(T)=\Phi_\vartheta(TA)$, so that
$f(v,\Psi_0)$ is a non-zero multiple of $\omega(A^{-1},1)f(v,\Phi_\vartheta)$ (by \eqref{po:eq-w5}
at $\mathsf S_\vartheta$; at $\aubar{\mathsf S_\vartheta}$ the component of $A$ is determined by
unitarity and acts on the vacuum of the conjugate place by a scalar). Hence $L_\Xi(f(v,\Psi_0))(g)$
is a multiple of $L_\Xi(f(v,\Phi_\vartheta))(gA^{-1})=0$ by Lemma~\ref{ppe:harm} and
\eqref{po:eq-equiv}. Now $z_uA_u=\binom{1}{z''_u z_u'^{-1}}$ is integral, $\equiv\iota_0\bmod u$ at
$\mathsf M$, and has upper block $1$ at $\mathsf Q$; so $\Psi_0(z)=1$, and since $\Psi_0$ is
locally constant, $\Psi_0\equiv1$ on $C$ for $i$ large. The support of $\Psi_0$ is one orbit of
$K_\vartheta$: if $\Psi_0(x)\ne0$ then $y=xA$ is integral at $\mathsf S_\vartheta$, $y\equiv\iota_0$
at $\mathsf M$ and $\det y'$ is a unit at $\mathsf Q$, and the element $k$ with components
$\begin{pmatrix}y'&0\\y''&1\end{pmatrix}$ at $\mathsf S_\vartheta$ (and $1$ elsewhere in
$K^\circ$) lies in $K_\vartheta$ and satisfies $k\,\iota_0A^{-1}=x$; and
$\Psi_0(kT)=\chi'(k)\Psi_0(T)$ for $k\in K_\vartheta$, with $\chi'(k)=\prod_{u\in\mathsf Q}
\psi_u(\det k'_u)$, because left multiplication by $k$ multiplies the upper block of $TA$ by
$k'$ at $\mathsf Q$ ($k$ is block lower triangular there) and preserves the condition at
$\mathsf M$. Let $\Psi_1=\int_{K_\vartheta}\chi_\vartheta(k)\,\omega(1,k)\plone_C\,dk$ (probability Haar
measure; the integrand is locally constant on the compact group $K_\vartheta$, so this is a finite
sum). By the rule above, $L_\Xi(f(v,\Psi_1))=\int\chi_\vartheta(k)\chi_\vartheta(k)^{-1}
L_\Xi(f(v,\plone_C))\,dk=L_\Xi(f(v,\plone_C))$, and $\omega(1,k')\Psi_1=\chi_\vartheta(k')^{-1}
\Psi_1$ for $k'\in K_\vartheta$ (substitute $k\mapsto k'^{-1}k$). The function $\Psi_1$ depends only on the components
at $\mathsf S_\vartheta$, is supported in $K_\vartheta\cdot C\subseteq\operatorname{supp}\Psi_0$, and
transforms under $K_\vartheta$ by the same character as $\Psi_0$ (the character by which
$\omega(1,k)$ acts on $f(v,\Phi_\vartheta)$ is $\chi_\vartheta(k)^{-1}$, \S\ref{po:ss-set}, and the
same holds for $\Psi_0$). A function supported on one orbit and transforming by a character is
determined by its value at one point, so $\Psi_1=\Psi_1(z)\Psi_0$, and
$L_\Xi(f(v,\plone_C))=\Psi_1(z)L_\Xi(f(v,\Psi_0))=0$.

\emph{Case B: $z_u$ has no integral block ratio at some $u\in\mathsf S_\vartheta$.} Let
$l=(\lambda,\lambda')$ be an annihilating row (Lemma~\ref{ppe:row}) and $\bar\lambda\ne0$ the
reduction of $\lambda$. For $s\in\ppeN_{\bar\lambda}$ write $s=1+\bar v\bar\lambda$, lift $\bar v$
to an integral column $v$, and let $k_s\in K^\circ$ have component $1+\iota_0vl$ at $u$ and $1$ at
the other places of $\mathsf S_\vartheta$ (the determinant $1+lv\equiv\det s$ is a unit). Since
$lz_u=0$, $k_sz_u=z_u$, and $k_s$ is integral, so $\omega(1,k_s^{-1})\plone_C=\plone_C$ and by
\eqref{ppe:eq-move} $L_{R_{k_s}\Xi}(f(v,\plone_C))=L_\Xi(f(v,\plone_C))$. Modulo $u$,
$k_s\equiv\begin{pmatrix}s&\bar v\bar\lambda'\\0&1\end{pmatrix}$. If $u\in\mathsf M$, then
$k_s=\tilde s'\kappa$ with $\kappa$ fixing $\iota_0$ modulo $u$, so $\kappa\in K_\vartheta$ with
$\chi_\vartheta(\kappa)=1$ (the datum at $\mathsf M$ is $\plone[T\equiv\iota_0]$, invariant under
such $\kappa$ without character), and $R_{k_s}\Xi=R_{\tilde s'}\Xi=\psi_u(\det s)^{-1}
\varrho_V(s)\Xi$; summing over $s$,
\[
\#\ppeN_{\bar\lambda}\cdot L_\Xi(f(v,\plone_C))=L_{\varrho_V(\eta_{\psi_u,\bar\lambda})\Xi}
(f(v,\plone_C))=0
\]
by (F2) and linearity of $\Xi\mapsto L_\Xi$. If $u\in\mathsf Q$, then $\bar\lambda'=0$, $k_s\equiv
\operatorname{diag}(s,1)$ lies in $K_\vartheta$, and $R_{k_s}\Xi=\chi_\vartheta(k_s)^{-1}\Xi$ with
$\chi_\vartheta(k_s)=\psi_u(\det s)^{\pm1}$ (by \eqref{po:eq-w5} and the factor $\psi_u(\det T')$
of the datum); summing over $s$ gives $\#\ppeN_{\bar\lambda}L_\Xi=(\sum_s\psi_u(\det s)^{\pm1})
L_\Xi=0$, because $\det\colon\ppeN_{\bar\lambda}\to\auF_u^\times$ is a surjective homomorphism and
$\psi_u\ne1$ ($\vartheta$ is a twist).
\end{proof}

\subsubsection{Hyperbolic planes}\label{ppe:ss-planes}

Let $w$ be a split place off $\Sigma$, $q=\#\auF_w$, $K^X_w$ the stabiliser of the standard
lattice in $\auU(X)(E^+_w)\cong\auGL_{m_X}(E_w)$, and for a $K^X_w$-invariant function $F$ on
$\auU(X)(\auA)$ let $P^X_w(X')F=\sum_i(-X')^iq^{\binom i2}T^X_{w,i}F$, with $T^X_{w,i}F(g)=
\sum_{x\in K_wd_iK_w/K_w}F(gx)$, $d_i=\operatorname{diag}(\varpi_w1_i,1)$, as in (N3).

\begin{lemma}[Small roots]\label{ppe:small}
Let $F$ be a bounded $K^X_w$-invariant function on $\auU(X)(\auA)$ with $P^X_w(X')F=0$, where
$q^{m_X+1}|X'|\le1$. Then $F=0$.
\end{lemma}
\begin{proof}
The cosets $K_wd_iK_w/K_w$ correspond to the subspaces of $\auF_w^{m_X}$ of dimension $i$ (proof
of Lemma~\ref{pl:vac}), so $P^X_w(X')F(g)=\sum_{U}c_UF(gx_U)$ over all subspaces $U$, with
$c_U=(-X')^{\dim U}q^{\binom{\dim U}2}$ and $x_0=1$. By \eqref{pl:eq-qbin} at $-|X'|$,
$\sum_U|c_U|=\prod_{l<m_X}(1+q^l|X'|)\le\exp(\sum_{l<m_X}q^l|X'|)\le\exp(q^{m_X}|X'|)\le
\exp(1/q)\le e^{1/2}<2$. Let $S=\sup|F|$. From $F(g)=-\sum_{U\ne0}c_UF(gx_U)$,
$S\le(\sum_U|c_U|-1)S$ with $\sum_U|c_U|-1<1$; so $S=0$.
\end{proof}

\begin{lemma}[Hecke transfer]\label{ppe:transfer}
Let $Y$ be continuous on $[\auU(V)]$ and right invariant under $K^V_w$, $f\in\ppeS_\Sigma(X)$ with
$m_X\le n$, $k=n-m_X$, $a$ unramified at $w$, and $F=\ppeD^X_a(f)Y$. Then $F$ is right invariant
under $K^X_w$, and there are $\alpha,\beta\in\C$ with $|\alpha|=q^{m_X/2}$, $|\beta|=q^{n/2}$
(independent of $Y$ and $f$) such that for every $X_0\in\C$
\[
\ppeD^X_a(f)\bigl(P^V_w(X_0)Y\bigr)=\prod_{l<k}(1-q^l\alpha X_0)\cdot P^X_{w^\ast}\bigl(q^k\alpha
\beta^{-1}X_0\bigr)F,
\]
where $w^\ast$ is $w$ or $\aubar w$ (the same for all $X_0$), $P^X_{\aubar w}$ being defined with the
identification $\auU(X)(E^+_w)\cong\auGL_{m_X}(E_{\aubar w})$.
\end{lemma}
\begin{proof}
The datum is the vacuum at $w,\aubar w$, fixed up to the unramified characters (trivial on
integral elements) by $K^X_w\times K^V_w$ (\eqref{po:eq-w5}); this gives the invariance. Since
$f$ is $K^V_w$-fixed, substituting $h\mapsto hx^{-1}$ in each term,
$\ppeD(f)(T^V_{w,i}Y)=\ppeD(T'_if)Y$ with $T'_i=\sum_{y\in K_wd_i^{-1}K_w/K_w}\omega(1,y)$; under the
identification $\auGL_n(E_w)\cong\auGL_n(E_{\aubar w})$, $g\mapsto{}^t\aubar g^{-1}$, the double coset
of $d_i^{-1}$ becomes that of $d_i$, and the vacuum in the coordinates at $w$ is the vacuum in those
at $\aubar w$. So $\ppeD(f)(P^V_w(X_0)Y)=\ppeD(P'(X_0)f)Y$ with $P'$ the operator of
Lemma~\ref{pl:vac} for the pair $(X,V)$ at $\aubar w$ (bigger space $V$ of rank $n$, smaller $X$
of rank $m_X$, $k_0=k$; that computation is local and uses only \eqref{po:eq-w5}). By
Lemma~\ref{pl:vac}, $P'(X_0)f=\prod_{l<k}(1-q^l\alpha X_0)P^-(q^k\alpha\beta^{-1}X_0)f$, where
$P^-$ is built from $\sum_{g\in K d_i^{-1}K/K}\omega(g,1)$; and $\ppeD(\omega(g,1)f)Y=R_gF$. Translating
the double cosets back gives $P^X_{w^\ast}$.
\end{proof}

\begin{lemma}[The zero term]\label{ppe:zero}
Let $p\ge1$, $W_p=W_{p-1}\oplus\ppeH$ with isotropic vectors $e=e_p$, $f_p$, and for
$f\in\ppeS(W_p)$ let $f[x]\in\ppeS(W_{p-1})$ ($x\in V(\auA)$) be its mixed model along the
polarisation $\auHom(Ee,V)\oplus\auHom(Ef_p,V)$, and $\theta_0(f)=\theta(f[0])$ its zero term. If
$\ppeVan(W_{p-1},a,Y)$, then for every $f\in\ppeS_\Sigma(W_p)$ and $g\in\auU(W_p)(\auA)$,
$\int_{[\auU(V)]}\theta_0(\omega^{W_p}_{a\upsilon_1}(g,h)f)Y(h)\,dh=0$.
\end{lemma}
\begin{proof}
Write $g=n\,m(q,g_0)\kappa$ (Iwasawa decomposition for the parabolic subgroup stabilising $Ee$,
with $\kappa$ in the standard maximal compact subgroup). Put $f'=\omega(\kappa,1)f$; it lies in
$\ppeS_\Sigma(W_p)$ (at the places outside $\Sigma$, $\kappa$ fixes the vacuum up to a scalar). By
Lemma~\ref{pcu:laws} with $x=0$, the radical acts trivially on the zero slice and
$(\omega(m(q,g_0),h)f')[0]=\nu(q)\,\omega^{W_{p-1}}_a(g_0,h)(f'[0])$ with $\nu(q)=\chi_V(\det q)^{-1}
|\det q|^{n/2}$ independent of $h$ (the shift of the character by $\upsilon_1$ is the one in that
lemma). The slice $f'[0]$ lies in $\ppeS_\Sigma(W_{p-1})$ (the mixed model of the vacuum at a place
outside $\Sigma$ is the vacuum times the characteristic function of the integral points). Hence the
integral is $\nu(q)\,\ppeD^{W_{p-1}}_a(f'[0])Y(g_0)=0$.
\end{proof}

\begin{lemma}[Bounded descents]\label{ppe:bdd}
In the situation of Lemma~\textup{\ref{ppe:zero}}, let $f=v\otimes f_0$ with $v$ a polynomial Fock
vector and $f_0$ a finite datum in $\ppeS_\Sigma$ (a finite combination of finite Heisenberg
translates of the vacuum). Then $\ppeD^{W_p}_{a\upsilon_1}(f)Y$ is bounded on $\auU(W_p)(\auA)$.
\end{lemma}
\begin{proof}
Let $F=\ppeD^{W_p}_{a\upsilon_1}(f)Y$, $\Lambda=\theta$ on $\ppeS(W_p)$ and $Z=\{x\in\auHom(W_p,V)(E):
x\text{ factors through }w'\mapsto\langle w',e\rangle\}$, an isotropic subspace; the mean of
$\Lambda\circ\rho$ over $Z(\auA)/Z(E)$ is $\Lambda_0=\theta_0$. Since $F$ is left invariant under
$\auU(W_p)(E)$, $F(g)=F(\gamma g)$ for $\gamma$ rational, and by Lemma~\ref{ppe:zero}
$F(\gamma g)=\int_{[\auU(V)]}(\Lambda-\Lambda_0)(\omega(\gamma g,h)f)Y(h)\,dh$. Let $A$ be as in (R2)
of \S\ref{pcx:ss-conv}; $|(\Lambda-\Lambda_0)(\omega(g',h)f)|\le2Y_0\|f\|\mathrm{ht}(g')^A$ by (R2)
($\mathrm{ht}(h)=1$, $V$ being definite). If $\mathrm{ht}(g)\le2$ take $\gamma=1$. If
$\mathrm{ht}(g)>2$, there is $0\ne d\in\Lambda_g$ isotropic with $|d|_g\le2/\mathrm{ht}(g)$ (by (R3)
and the remark after it), and by Witt's theorem a $\gamma\in\auU(W_p)(E)$ with $\gamma d$ a
multiple of $e$. Since $\Lambda_{\gamma g}=\gamma\Lambda_g$ isometrically, $\gamma d\in
\Lambda_{\gamma g}$ with $|\gamma d|_{\gamma g}=|d|_g\le2/\mathrm{ht}(g)$, and
$\mathrm{ht}(\gamma g)=\mathrm{ht}(g)$, since $\Lambda_{\gamma g}=\gamma\Lambda_g$ isometrically and the
height is defined from the lattice (so (R2) at $\gamma g$ gives the factor $\mathrm{ht}(g)^A$). As
$F$ is an integral over $[\auU(V)]$, we may let $h$ run over a compact fundamental set
$\mathcal F$ ($V$ definite). Let $d_1,\dots,d_n$ be a basis of the standard lattice of $V$ and $N$
a positive integer with $Nh^{-1}d_i$ integral at the finite places for all $h\in\mathcal F$; the
elements $z_i=Nd_i\langle\cdot,\gamma d\rangle$ span $Z$ (the same line as $\langle\cdot,e\rangle$),
lie in the lattice at $k=(\gamma g,h)$, and have length $\le c\,\mathrm{ht}(g)^{-1}$ there,
uniformly in $h\in\mathcal F$. Lemma~\ref{pcx:mean} with
$\delta=c\,\mathrm{ht}(g)^{-1}$ and $l=A$ gives $|(\Lambda-\Lambda_0)(\omega(\gamma g,h)f)|\le
C\,\mathrm{ht}(g)^A\delta^A\|v\|_A\le C'$. Hence $|F(g)|\le C'\sup|Y|\operatorname{vol}[\auU(V)]$.
\end{proof}

\begin{proposition}[The tower]\label{ppe:tower}
Let $\Xi\in\poA(V;\vartheta)$ satisfy
\textup{(F1)} $D_\vartheta\Xi=0$, \textup{(F2)} of Lemma~\textup{\ref{ppe:three}}, and
\textup{(F3)} $P_w(X_w)\Xi=0$ for a place $w$ off the level. Let $p\ge0$ with $2r+2p+3<n$. Then
$\ppeVan(W_p,\upsilon_1^p,Y_\zeta)$ for every $\zeta\in W_{\lambda_V}$.
\end{proposition}
\begin{proof}
First $\zeta=\xi_0$. For $p=0$ this is Proposition~\ref{ppe:bottom}. Let $p\ge1$ and assume the
claim for $p-1$. By Lemma~\ref{ppe:special} it suffices to treat $f=v\otimes f_0$ as in
Lemma~\ref{ppe:bdd}; let $F=\ppeD^{W_p}_{\upsilon_1^p}(f)Y_{\xi_0}$, bounded by Lemma~\ref{ppe:bdd}.
The function $Y_{\xi_0}$ is right invariant under $K^V_w$ ($w\notin\Sigma$) and, the Hecke operators
having integer coefficients, $P_w(\aubar{X_w})Y_{\xi_0}=\aubar{c_{\xi_0}(P_w(X_w)\Xi)}=0$. By
Lemma~\ref{ppe:transfer} with $m_X=r+2p$, $k=n-m_X$, $X_0=\aubar{X_w}$:
$\prod_{l<k}(1-q^l\alpha X_0)\,P^{W_p}_{w^\ast}(X')F=0$, $X'=q^k\alpha\beta^{-1}X_0$. Here
$|X_0|=q^{-(r/2+t-1)}$, so $|q^l\alpha X_0|\le q^{k-1+m_X/2-r/2-t+1}=q^{-(m_X-r)/2}<1$ for $l<k$ (as
$k+m_X/2=n-m_X/2$ and $m_X\ge r+2$): the product is non-zero. And $|X'|=q^{k/2}|X_0|$, so
$q^{m_X+1}|X'|=q^{m_X+1+(n-m_X)/2-r/2-t+1}=q^{(m_X+4+r-n)/2}$, and $m_X+4+r-n=2r+2p+4-n\le0$.
By Lemma~\ref{ppe:small} (at $w^\ast$), $F=0$. For general $\zeta$: for $x\in\auU(V)(E^+\otimes\R)$,
$Y_{\rho(x)\xi_0}$ is the right translate of $Y_{\xi_0}$ by $x$, and
$\ppeD(f)(R_xY)=\ppeD(\omega(1,x^{-1})f)Y$ with $\omega(1,x^{-1})f\in\ppeS_\Sigma(W_p)$; so
$\ppeVan$ holds for $Y_{\rho(x)\xi_0}$. The vectors $\rho(x)\xi_0$ span $W_{\lambda_V}$
(irreducibility) and $\zeta\mapsto Y_\zeta$ is linear (conjugate of a conjugate-linear map), so it
holds for all $Y_\zeta$.
\end{proof}

\subsubsection{From descents to periods; proof of Theorem~\ref{ppe:main}}\label{ppe:ss-periods}

\begin{proposition}[Descents and periods]\label{ppe:rallis}
Let $V=V_1\perp V_2$ be a splitting of the standard coordinates, of ranks $n_1,n_2$, with
$o_1\in M_{n\times n_1}$ the inclusion of $V_1$. Let $Y$ be continuous on $[\auU(V)]$ and right
invariant under $K''=\{k\in K^\circ:k_u=1\ (u\in\Sigma)\}$. If $\ppeVan(\ppeH^{n_1},\upsilon_1^{n_1},Y)$,
then $I(h_1)=\int_{[\auU(V_2)]}Y(\operatorname{diag}(h_1,h_2))\,dh_2=0$ for every $h_1\in\auU(V_1)(\auA)$.
\end{proposition}
\begin{proof}
Let $X=\ppeH^{n_1}$ with the Lagrangian $L=\langle e_1,\dots,e_{n_1}\rangle$, Siegel parabolic
$P=MN$, $m(q)$ ($q\in\auGL_{n_1}(\auA_E)$) and $n(z)$ ($z$ hermitian; $n(z)$ is the element
$n(0,-\jmath^{-1}z)$ of \S\ref{pcu:ss-top-setting}, with skew-hermitian argument $-\jmath^{-1}z$,
and $J_V=1$ since the standard $V$ is used, so that the character
$\psi(-\frac12\jmath\operatorname{tr}(z'x^*J_Vx))$ of Lemma~\ref{pcu:laws} with $z'=-\jmath^{-1}z$ is the
one written below). Since $X=L\oplus L'$, the
mixed model of Lemma~\ref{pcu:laws} (with $W_0=0$) is the Schr\"odinger model: $f\mapsto f[\,\cdot\,]$,
a function of $x\in M_{n\times n_1}(\auA_E)=\auHom(L',V)(\auA)$, and
\[
(\omega(n(z)m(q),h)f)[x]=\psi\bigl(\tfrac12\operatorname{tr}(z\,x^*x)\bigr)\,\nu(q)\,f[h^{-1}xq],
\qquad\theta(f)=\sum_{x\in M_{n\times n_1}(E)}f[x],
\]
with $\nu(q)\ne0$. The image of $\ppeS_\Sigma(X)$ consists of the functions $\phi_\infty\otimes
\phi_\Sigma\otimes\plone_{\mathrm{int}}$, $\phi_\infty$ in the Schr\"odinger image of the Fock space
(which contains all products of polynomials and the Gaussian, dense in the continuous functions on
compact sets), $\phi_\Sigma$ any Schwartz--Bruhat function at $\Sigma$, and $\plone_{\mathrm{int}}$
the characteristic function of the integral matrices outside $\Sigma$ (the vacuum of the adapted
frame). Let $F=\ppeD^X_{\upsilon_1^{n_1}}(f)Y=0$. Its Fourier coefficient along $N$ at $\beta=1$,
$\int_{N(E)\backslash N(\auA)}\psi(-\frac12\operatorname{tr}z)F(n(z)m(q))\,dz$, equals (inserting the
formula, exchanging sum and integrals by absolute convergence: $Y$ is bounded, $[\auU(V)]$ compact,
$\sum_x\sup_h|f[h^{-1}xq]|<\infty$) $\nu(q)\sum_{x^*x=1}\int_{[\auU(V)]}f[h^{-1}xq]Y(h)\,dh$. By
Witt's theorem the rational frames with $x^*x=1$ form the orbit $\auU(V)(E)\,o_1$, with stabiliser
$\auU(V_2)(E)$; unfolding (and $Y(\gamma h)=Y(h)$),
\[
0=\int_{\auU(V_2)(E)\backslash\auU(V)(\auA)}f[h^{-1}o_1q]\,Y(h)\,dh .
\]
Take $q=h_1\in\auU(V_1)(\auA)\subseteq\auGL_{n_1}(\auA_E)$; then $o_1h_1=\operatorname{diag}(h_1,1)o_1$,
and the substitution $h\mapsto\operatorname{diag}(h_1,1)h$ (which normalises $\auU(V_2)$) gives
$0=\int_{\auU(V_2)(\auA)\backslash\auU(V)(\auA)}f[h^{-1}o_1]\,J(h)\,dh$ with
$J(h)=\int_{[\auU(V_2)]}Y(\operatorname{diag}(h_1,h_2)h)\,dh_2$, a continuous function with $J(1)=
I(h_1)$, right invariant under $K''$. The map $h\mapsto h^{-1}o_1$ identifies
$\auU(V_2)(\auA)\backslash\auU(V)(\auA)$ with the closed set $\Omega$ of adelic frames with $x^*x=1$
(Witt at each place), and $J$ with a continuous function $\tilde J$ on $\Omega$. Outside $\Sigma$
the factor $\plone_{\mathrm{int}}$ restricts to the integral frames, on which $\prod_{u\notin\Sigma}
K^\circ_u$ acts transitively (if $x$ is integral with $x^*x=1$, the lattice $L_1=x\auO_u^{n_1}$ is
unimodular, so $\auO_u^n=L_1\perp L_1^\perp$ with $L_1^\perp$ unimodular of rank $n_2$; since
$E/E^+$ is unramified at all finite places (choice of $E$), unimodular hermitian lattices over
$\auO_u$ of a given rank are all isometric, as each has an orthonormal basis (proof: the reduction of the lattice is a non-degenerate hermitian space over the quadratic extension of residue fields, so it has a vector $\bar x$ with $\bar h(\bar x,\bar x)\ne0$ (otherwise polarisation and the surjectivity of the residue trace give $\bar h=0$); a lift $x$ has $h(x,x)$ a unit of the base ring, hence the norm of a unit since the extension is unramified, so after scaling $h(x,x)=1$; then the lattice is $\mathcal O x\perp(x^\perp\cap\text{lattice})$ with unimodular second summand, and one inducts on the rank; see also \cite{Jacobowitz}), and at a split place this is the
statement that $\auGL_n(\auO)$ acts transitively on bases; hence $L_1^\perp\cong\auO_u^{n_2}$, and the
isometry $o_1\auO_u^{n_1}\perp(o_1\auO_u^{n_1})^\perp\to L_1\perp L_1^\perp$ sending $o_1$ to $x$ lies
in $K^\circ_u$), so $\tilde J(x)=\tilde J_0(x_\infty,x_\Sigma)$
there. Thus $\int_{\Omega_\infty\times\Omega_\Sigma}\phi_\infty\phi_\Sigma\tilde J_0\,d\mu=0$ for all
admissible $\phi_\infty,\phi_\Sigma$, where $\Omega_\infty$ is compact ($V$ definite) and the
restrictions of the $\phi_\Sigma$ to $\Omega_\Sigma$ are all locally constant compactly supported
functions. By density, $\tilde J_0=0$ on $\Omega_\infty\times\Omega_\Sigma$; at $h=1$ the frame is
$o_1$, which is integral, so $I(h_1)=J(1)=0$.
\end{proof}

\begin{proof}[Proof of Theorem~\ref{ppe:main}]
By Lemma~\ref{ppe:three}, $\Xi$ satisfies (F1)--(F3). Let $V_1=V''$ (rank $n_1=N''=r+9m+1$),
$V_2=V^\circ$. Since $r\ge x^5$ and $m\le2\sqrt x$ (Lemma~\ref{an:lem-sizes}),
$2r+2N''+3=4r+18m+5<6r=n$. Fix $s$, $h$, $u''$, $g$, and write
$\langle u''\odot K^\circ(g,h^\circ),\Xi(\cdot)\rangle=\sum_j\aubar{\theta^\circ(g,h^\circ;e_j)}\,
\langle u''\odot e_j,\Xi(\cdot)\rangle$ for an orthonormal basis $e_j$ of $W^\circ$ (the pairing is
conjugate-linear in the first slot and $\langle e_j,K^\circ\rangle=\theta^\circ(\cdot;e_j)$). By
Lemma~\ref{po:prod}(M1) and the equivariance of $\pop^\circ$, the function
$h^\circ\mapsto\sum_j\aubar{\theta^\circ(g,h^\circ;e_j)}c_{u''\odot e_j}\Xi(\operatorname{diag}(h,
h^\circ)\pog_s)$ on $\auU(V^\circ)(\auA)$ is invariant under $\auU(V^\circ)(\R)$ and equals the
summand of the theorem on finite points; it is left invariant under $\auU(V^\circ)(E)$ and right
invariant under $J^\circ$. Since $J^\circ$ is free, the sum over the classes is
$\operatorname{vol}^{-1}$ times its integral over $[\auU(V^\circ)]$. So it suffices that, for each
$j$ and with $\zeta=u''\odot e_j$,
\begin{equation}\label{ppe:eq-goal}
\int_{[\auU(V^\circ)]}\aubar{\theta^\circ(g,h_2;e_j)}\,c_\zeta\Xi(\operatorname{diag}(h,h_2)\pog_s)\,
dh_2=0 .
\end{equation}
Let $Y_1(h')=Y_\zeta(h'\pog_s)=\aubar{c_\zeta\Xi(h'\pog_s)}$. Since $\pog_s$ is $1$ outside
$\mathsf M\cup\aubar{\mathsf M}$, the argument at the end of the proof of
Proposition~\ref{ppe:tower} (with $\pog_s$ in place of a real element; $\omega(1,\pog_s^{-1})$
preserves $\ppeS_\Sigma$) and Proposition~\ref{ppe:tower} with $p=N''$ give
$\ppeVan(W_{N''},\upsilon_1^{N''},Y_1)$. Choose (by (W7) of \S\ref{po:ss-glue}, applied to
$W_{N''}=\ppeH^{N''}\oplus W$ and to $V=V''\perp V^\circ$) a datum $\phi''\in\ppeS_\Sigma$ for the
pair $(W,V'')$ with $\theta''(\omega''(g,h)\phi'')\ne0$. Such a datum exists. Work in a Schr\"odinger model
on $\auHom(W,V'')(\auA)$ in which the vacuum outside $\Sigma$ is the characteristic function
$\plone_{\mathrm{int}}$ of the integral points, and let $P$ be the finite set of finite places
outside $\Sigma$ at which $g$ or $h$ is not integral. For $\phi=\phi_\infty\otimes\phi_\Sigma\otimes
(\text{vacuum outside }\Sigma)$ we have $\omega''(g,h)\phi=\phi'_\infty\otimes\phi'_\Sigma\otimes
\phi_P\otimes\plone_{\mathrm{int}}$, where $\phi_P=\omega''(g,h)_P\plone_{\mathrm{int}}\ne0$ is
fixed, and, the local operators $\omega''(g,h)_\infty$, $\omega''(g,h)_\Sigma$ being invertible on
the local Schwartz spaces and $\ppeS_\Sigma$ being arbitrary at $\infty$ and $\Sigma$,
$\phi'_\Sigma$ is an arbitrary Schwartz--Bruhat function and $\phi'_\infty$ an arbitrary Schwartz
function. By strong approximation for the vector group $\auHom(W,V'')$ choose a rational $x_0$,
integral outside $\Sigma\cup P$, with $\phi_P(x_0)\ne0$, and an open compact neighbourhood
$U_\Sigma$ of $x_0$ at $\Sigma$; take $\phi'_\Sigma=\plone_{U_\Sigma}$. The datum at $P$ is the fixed
Schwartz--Bruhat function $\phi_P$, compactly supported and locally constant. Let $D$ be the set of
rational $y$ with $y_\Sigma\in U_\Sigma$, $y_P\in\operatorname{supp}\phi_P$ and $y$ integral outside
$\Sigma\cup P$. The conditions at the finite places bound the denominators and confine $y$ to a
compact open set of finite adeles, so $D$ is a finite union of cosets of a lattice in the real space
$\auHom(W,V'')(E\otimes\R)$; in particular $D$ is discrete, $x_0\in D$, and $\phi_P$ takes finitely many
values on $D$. Hence
$\theta''(\omega''(g,h)\phi)=\sum_{y\in D}\phi_P(y)\,\phi'_\infty(y)$. With
$\phi'_\infty(y)=\exp(-T|y-x_0|^2)$ the term $y=x_0$ equals $\phi_P(x_0)\ne0$, and the other terms
are bounded in total by $\max|\phi_P|\sum_{y\in D\setminus\{x_0\}}e^{-T|y-x_0|^2}$, which tends to $0$
as $T\to\infty$ ($D$ discrete and a finite union of lattice cosets). So the theta value is non-zero
for $T$ large; $\phi''=\phi$ for the corresponding $\phi_\infty$ works. Let $\phi^\circ=
\pop^\circ_{e_j}\otimes\Psi^\circ$ and $\Phi=\phi''\otimes\phi^\circ\in\ppeS_\Sigma(W)$, and put
$Y_2(h')=\theta(\omega^W(g,h')\Phi)\,Y_1(h')$. Then $Y_2$ is continuous, left invariant under
$\auU(V)(E)$, right invariant under $K''$ (the datum, $\Xi$ and $\pog_s$ are), and its descents to
$\ppeH^{N''}$ vanish: for $f'\in\ppeS_\Sigma(\ppeH^{N''})$ the kernel of $f'\otimes\Phi$ on
$W_{N''}$ at $((g',g),h')$ is the product of the kernels of $f'$ and $\Phi$ (compatibility of theta
functionals with orthogonal sums of the source space), so $\ppeD^{\ppeH^{N''}}(f')Y_2(g')=
\ppeD^{W_{N''}}(f'\otimes\Phi)Y_1((g',g))=0$. By Proposition~\ref{ppe:rallis},
$\int_{[\auU(V^\circ)]}Y_2(\operatorname{diag}(h,h_2))\,dh_2=0$. By (W7) the kernel factors,
$\theta(\omega^W(g,\operatorname{diag}(h,h_2))\Phi)=\theta''(\omega''(g,h)\phi'')\cdot\theta^\circ(g,
h_2;e_j)$ (the characters of the two factors being those of (W7)), so dividing by the non-zero first
factor, $\int\theta^\circ(g,h_2;e_j)\aubar{c_\zeta\Xi(\operatorname{diag}(h,h_2)\pog_s)}\,dh_2=0$, the
complex conjugate of \eqref{ppe:eq-goal}. The same argument applies to any open compact subgroup of
$J^\circ$ and to any $\Xi$ with (F1)--(F3).
\end{proof}

\begin{remark}[What is used]\label{ppe:rem-used}
Besides results proved in this paper (Proposition~\ref{plf:L4}, Lemmas~\ref{pl:vac},
\ref{pcu:laws}, \ref{pcx:mean} and the facts (R1)--(R3) of \S\ref{pcx:ss-conv}, (W7) of
\S\ref{po:ss-glue}), the proof uses the following published results: the restriction formula for
Green's characters (Lemma~\ref{ppe:green}; \cite[\S9]{ZelevinskyLNM}, \cite[Thm.~13]{Green1955}); the generation of
the isotypic components of the polynomial Fock space by joint harmonics (Lemma~\ref{ppe:harm};
\cite[Ch.~III, Corollary~5.4]{KashiwaraVergne}, \cite[Proposition~3.1]{HoweJAMS}); Schur orthogonality; Witt's theorem; and the property
$\psi_v^{q_2}=1$ of the tame parts at the chosen places (Remark~\ref{ppe:rem-T}) (proved in Proposition~\ref{pfq:tame}).
\end{remark}

\begin{remark}[Rank one]\label{ppe:rem-rankone}
Theorem~\ref{ppe:main} and its proof are used in rank one in Proposition~\ref{psu:periods}: $W=E$
(rank $r=1$, definite), $V=V^{(z)}$ of rank $n=2z$, so that $t=n-r=2z-1$; the character $\nu\vartheta$ of a
critical pair $(\nu,z)$; the idempotent $e_1$ (the mean over $\auG_1$) in place of $e_0$; a splitting
$V=V'\perp V_1$ with $V'$ of rank $n_1=20$ in place of $V''\perp V^\circ$; an integral element $x'$
equal to $1$ outside $\Sigma$ in place of $\pog_s$; and the level of the flat pair on $V_1$ in place of
$J^\circ$. We list the steps of the proof and what each needs for $r=1$.
\begin{enumerate}
\item[(1)] \emph{Set-up} (\S\ref{ppe:ss-setup}). Nothing depends on $r$; the descents $\ppeD^X_a$ are
defined for any $X$, and an extra character of $\det g$ on $\auU(1)$ only multiplies descents.
\item[(2)] \emph{(F1).} For $r=1$ this is the hypothesis $D^{(1)}_\vartheta y=0$ of
Proposition~\ref{psu:periods}; Lemma~\ref{ppe:normal} (Schur orthogonality over the compact group
$\auU(V)(\R)$) does not use $r$.
\item[(3)] \emph{(F2).} Lemma~\ref{ppe:three} uses the label idempotents of $e_0$; for $r=1$ it is
replaced by the direct argument of Proposition~\ref{psu:periods}: $\auF_v^1=\auF_v$, the hyperplane
sum is $\eta=\sum_{s\in\auF_v^\times}\psi_v(s)^{-1}s$, and $\eta e_1=0$ as soon as $\psi_v\ne1$; here
$\psi_v=\omega_v^{\,j}\ne1$. No condition of order $q_2$ is needed in rank one.
\item[(4)] \emph{(F3).} $P_w(X_w)y=0$ is the hypothesis, with $|X_w|=q^{-r/2-(t-1)}=q^{-1/2-(2z-2)}$,
the value used in Lemma~\ref{ppe:three}(F3) and in the tower.
\item[(5)] \emph{Bottom of the tower} (Lemmas~\ref{ppe:special}, \ref{ppe:harm}, \ref{ppe:row},
Proposition~\ref{ppe:bottom}). The coordinates $T\in M_{n\times r}$ become columns $T\in M_{n\times1}$,
$T'$ is the entry of the row of $\iota_0$; the arguments (annihilating rows, the support of the block
datum, the non-triviality of $\psi_w$ at $\mathsf Q$) are written for arbitrary $r\ge1$. At
$w\in\mathsf Q$ the tame part of $\nu\vartheta$ is that of $\vartheta$, which is non-trivial ($\nu$ is
unramified at $\mathsf Q$). In detail: Lemmas~\ref{ppe:normal} and~\ref{ppe:special} do not involve
$r$ (the first is Schur orthogonality on $\auU(V)(\R)$, the second is stated for every space $X$).
Lemma~\ref{ppe:row} and Cases A and B of Proposition~\ref{ppe:bottom} are literally the same with
$r=1$: $T'$ is the first entry of the column $T$, $A_u=(z'_u)^{-1}\in E_u^\times$, the element
$k=\bigl(\begin{smallmatrix}y'&0\\y''&1\end{smallmatrix}\bigr)$ lies in $\pexK_\vartheta$, the group
$\ppeN_{\bar\lambda}$ is $\auF_u^\times$, and at $u\in\mathsf M$ the sum over $s$ is
$\varrho_V(\eta)$ with $\eta$ as in (3). In Lemma~\ref{ppe:harm} steps (a) and (c) are unchanged, and
step (d) is elementary, so that the theorem on joint harmonics is not needed: the line is positive
definite, so the Fock space of $\auHom(E,V)\otimes\R$ consists of holomorphic functions of
$(x_\tau)_\tau\in\prod_\tau\C^{2z}$, the group $\auU(V)(\R)$ acts on polynomials by linear substitution
(the splitting on $\auU(V)$ is $\upsilon^{r_-}=1$), and the polynomials that are homogeneous of degree
$d_\tau$ at each $\tau$ form an irreducible representation, different degrees giving inequivalent
ones. By \eqref{po:eq-poly}, $\pop_\xi(x)=\langle\xi,\bigotimes_\tau\aubar{x_\tau}^{\otimes l_\tau}\rangle$,
and the vectors $\bigotimes_\tau y_\tau^{\otimes l_\tau}$ span $W_\lambda=\bigotimes_\tau
\operatorname{Sym}^{l_\tau}\C^{2z}$; so $\{\pop_\xi\}$ is the whole space of polynomials of degree
$l_\tau$ at each $\tau$, which is the isotypic subspace of step (c). Hence $L(v)=L(e_\lambda v)$ with
$e_\lambda v\in\{\pop_\xi\}$, and $L(v)=0$ by (a).
\item[(6)] \emph{The tower} (Lemmas~\ref{ppe:small}--\ref{ppe:bdd}, Proposition~\ref{ppe:tower}). Stated
for arbitrary $r$; needs $2r+2p+3<n$, here with $p=n_1=20$: $45<2z$. This holds because the pairs
to which Proposition~\ref{psu:periods} is applied have $z>m$ and $m\ge\sqrt x>23$
(Lemma~\ref{an:lem-sizes}). The inequality $|q^l\alpha X_0|<1$ uses only $m_X\ge r+2$, true for $p\ge1$.
The four lemmas are applied to the spaces $W_p=E\oplus\ppeH^p$ and their statements and proofs do not
involve $W$ otherwise: Lemma~\ref{ppe:small} concerns one function on $\auGL_{m_X}(E_w)$;
Lemma~\ref{ppe:transfer} is the local computation of Lemma~\ref{pl:vac} at a place where all data
are the vacuum; Lemma~\ref{ppe:zero} uses Lemma~\ref{pcu:laws} for the isotropic line $Ee_p$ of $W_p$;
and Lemma~\ref{ppe:bdd} uses (R2), (R3) of \S\ref{pcx:ss-conv} and Lemma~\ref{pcx:mean}, which are
stated for arbitrary spaces, and Witt's theorem in $W_p$, $p\ge1$. In the proof of
Proposition~\ref{ppe:tower} the exponents are $q^{-(m_X-r)/2}$ and $q^{(2r+2p+4-n)/2}$ with $r=1$,
$n=2z$, and the passage from $\xi_0$ to all $\zeta$ is unchanged.
\item[(7)] \emph{Descents and periods} (Proposition~\ref{ppe:rallis}). Stated for an arbitrary splitting
of the standard coordinates; used with $V_1=V'$, $n_1=20$.
\item[(8)] \emph{Final step.} Let $y$, $u$, $h'$, $g$, $x'$ be as in Proposition~\ref{psu:periods},
$Y_u=\aubar{c_uy}$ and $Y_1(h)=Y_u(hx')$. By (6), $\ppeVan(W_{20},\upsilon_1^{20},Y_u)$; as
$\ppeD(f)(R_{x'}Y_u)=\ppeD(\omega(1,x'^{-1})f)Y_u$ and $\omega(1,x'^{-1})$ preserves $\ppeS_\Sigma$
($x'$ is $1$ outside $\Sigma$), also $\ppeVan(W_{20},\upsilon_1^{20},Y_1)$. The datum
$\phi^\circ=1\otimes\Psi_{q_1,1}$ of $\auHom(E,V_1)$ (constant polynomial, block datum of the flat pair)
lies in $\ppeS_\Sigma$, and so does $\Phi=\phi''\otimes\phi^\circ$ for $\phi''\in\ppeS_\Sigma$ on
$\auHom(E,V')$. By (W7) for $V=V'\perp V_1$, $\theta(\omega^E(g,[h',h_1])\Phi)=\theta''(\omega''(g,h')
\phi'')\,k_1(g,h_1)$, where $\omega^E$ is a Weil representation of $(E,V)$ whose character on $\auU(1)$
is the product of the character of $k_1$ and of that of $\omega''$; the latter is at our disposal,
since a character of $\det g$ on $\auU(1)$ only multiplies descents (step (1)). The datum $\phi''$ with
$\theta''(\omega''(g,h')\phi'')\ne0$ is constructed as in the proof of Theorem~\ref{ppe:main}; that
construction uses only that $\auHom(E,V')$ is a vector group. Then $Y_2(h)=\theta(\omega^E(g,h)\Phi)
Y_1(h)$ is continuous, left invariant under $\auU(V)(E^+)$ and right invariant under $K''$ (an element
of $K''$ fixes $\Phi$, commutes with $x'$ and lies in $\pexK_\vartheta$ with trivial character), and
$\ppeD^{\ppeH^{20}}_{\upsilon_1^{20}}(f')Y_2(g')=\ppeD^{W_{20}}_{\upsilon_1^{20}}(f'\otimes\Phi)Y_1((g',g))=0$.
Proposition~\ref{ppe:rallis} with $V_1,V_2$ there equal to $V',V_1$ here gives $\int_{[\auU(V_1)]}Y_2([h',
h_1])\,dh_1=0$, and dividing by $\theta''(\omega''(g,h')\phi'')$ we get $\int_{[\auU(V_1)]}k_1(g,h_1)\,
\aubar{c_uy([h',h_1]x')}\,dh_1=0$. The hypothesis that $\pop_u$ involves only the variables of $V'$
means that $u\in\bigotimes_\tau\operatorname{Sym}^{l_\tau}\C^{20}$, so $u$ is fixed by the block of $V_1$;
this gives the invariance of the integrand under $\auU(V_1)(E^+\otimes\R)$ and, on the left, under
$\auU(V_1)(E^+)$ (the kernel $k_1$ has trivial weight). The level on $V_1$ need not be
free, and then the sum over the classes weighted by $e_{h_1}^{-1}=\#(\text{stabiliser})^{-1}$ is the
integral over $[\auU(V_1)]$ divided by the volume of the level, which is the form in which
Proposition~\ref{psu:periods} is stated. For the summand to be a function on the classes (right
invariance under the level on $V_1$, which in \S\ref{ppe:ss-periods} comes from
Lemma~\ref{po:lattice}(d)) one needs the condition on $x'$ stated in Proposition~\ref{psu:periods}:
for $k$ in the level on $V_1$ the element $x'^{-1}\operatorname{diag}(1,k)\,x'$ lies in the level of
$V$ and acts on the block datum by the scalar by which $k$ acts on the datum on $V_1$.
\end{enumerate}
So Theorem~\ref{ppe:main} holds in rank one in the form of Proposition~\ref{psu:periods}, with (F2)
replaced by step (3).
\end{remark}


%% ===== proofs/C-span.tex =====
% proofs/C-span.tex
% Proposition po:span (products span). Labels of this file start with psp:.
\providecommand{\auO}{\mathcal O}
\providecommand{\auU}{\mathrm U}
\providecommand{\auGL}{\mathrm{GL}}
\providecommand{\auHom}{\operatorname{Hom}}
\providecommand{\auA}{\mathbb A}
\providecommand{\auPl}{\Phi}
\providecommand{\aubar}[1]{\overline{#1}}
\providecommand{\poA}{\mathcal A}
\providecommand{\pcMat}{\mathrm M}
\providecommand{\psvac}{{\phi^{\mathrm{vac}}}}
% macros of this file (prefix \psp)
\providecommand{\pspF}{\mathfrak F}
\providecommand{\pspl}{\lambda}
\providecommand{\pspn}{n_*}
\providecommand{\pspm}{\mathsf m}
\providecommand{\pspip}[2]{\langle #1,#2\rangle}
\providecommand{\pspG}{G_\infty}

\subsection{Products span}\label{psp:sec}

This subsection proves Proposition~\ref{po:span}. We use the notation of \S\ref{po:ss-glue} and of
\S\ref{pte:ss-carrier}--\S\ref{pte:ss-third}: $W=W_j$ of rank $r$, $1\le j\le m$; the space $S$ of
signature $(2m,m)$ and the carrier $\varphi^\circ$; the padded space $\widetilde W=W\oplus(E^j)^-\oplus
E^{3m-j}$, the space $W^\flat=E^{6m+1}$, $W'=\widetilde W\oplus W^\flat$; the choices $i$ and the
functions $F_i$ of \eqref{pte:eq-F}; the definite space $V^\circ$ and the kernel $K^\circ(g,h^\circ)$ of
$(W,V^\circ)$, with values in the Weyl module of the diagrams $\mu^-_\tau$; the modulus $\pi_M^2$.

\emph{What is proved and what is not.} The argument has four steps (\S\ref{psp:ss-outline}). Steps~3
and~4 rest on three estimates which are proved here in full: the decay of Heisenberg coefficients
of polynomial Fock vectors (Lemma~\ref{psp:decay}), the comparison of a theta kernel with a Fock
coefficient on a lattice without short vectors (Proposition~\ref{psp:sparse}), and the isolation
estimate (Proposition~\ref{psp:isol}). The proof of Proposition~\ref{po:span} from these is complete
(\S\ref{psp:ss-proof}). It also uses three \emph{inputs}, collected with exact statements in
\S\ref{psp:ss-inputs}: a reproducing formula for pairings with holomorphic discrete series
coefficients in the stable range (R), elementary properties of holomorphic functions of a given
weight on $\auU(p,q)$ (H), and the fact that peaks of bounded holomorphic forms lie away from the
cusps (P). They are proved in \S\ref{psq:sec}--\S\ref{psv:sec} (Theorem~\ref{psv:main}); (P) is
Siegel's argument, cf.\ Poor--Yuen~\cite{PoorYuen} for Siegel modular forms. Everything else is proved
here.

\subsubsection{Statement}\label{psp:ss-statement}

\begin{proposition}[Products span; Proposition~\ref{po:span}]\label{psp:main}
Let $y\in e_0\poA(W_j)$. Suppose that for every choice $i$ and every class $h^\circ$ of
$\auU(V^\circ)$ at level $\pi_M^2$,
\[
\int_{[\auU(W)]}y(g)\,F_i(g)\,K^\circ(g,h^\circ)\,dg=0
\]
(a vector of the Weyl module). Then $y=0$.
\end{proposition}

This is the printed statement; the integrals converge by Proposition~\ref{po:exchange}. Since the
forms on $V^\circ$ of \S\ref{po:ss-glue} are arbitrary functions on the classes, the hypothesis is
equivalent to $\int y\,F_i\,\pspip{\zeta}{K^\circ(\cdot,\cdot)}=0$ for all $i$ and all forms $\zeta$ on
$V^\circ$, which is the hypothesis $\Theta^\circ y=0$ of (G2) in Proposition~\ref{po:glue} (there
$c_i\ne0$). So (G2) follows as written there.

\subsubsection{Notation}\label{psp:ss-notation}

$f=\#\auPl$, $t=5r$, $\pspn=2ft^3$, and $b$ is the scale of the data. For a space $X$ (with its frame)
and $g\in\auU(X)(\auA)$, $\Lambda_g$, $|\gamma|_g$ and $\lambda_1(g)$ are as in
Definition~\ref{pcx:def-ht}, and $\pspl(X,g)=\min(1,\lambda_1(g))$. For $h\in\auU(X)_\infty=\prod_\tau
\auU(X_\tau)$ we write $|h|=\max_\tau\|h_\tau\|\ge1$ (operator norms in the frame, in which the form
is $\mathrm{diag}(\pm1)$). The Fock coordinates of $\auHom(X,X')\otimes\R$ are $x_\tau=c_\psi\tau(w)$
with $c_\psi=\sqrt{2\pi/\sqrt d}$ (\S\ref{pk:ss-conv}); $|x|$ is the euclidean norm of
$(x_\tau)_\tau$.

\emph{Weights.} On $\auU(p,q)$ with the frame ordered by a total order $o$ that puts the $p$ positive
indices first, a weight is $u\in\Z^{p+q}$; it has \emph{gap} $\gamma$ if $u_a+\gamma\le u_c$ whenever
$a$ is positive and $c$ negative, and it is \emph{dominant} if it is non-increasing in the order $o$
within each sign block. For a dominant weight $u$ with gap $\ge N^2+N+1$ ($N=p+q$) let $\pspm_u$ be the
coefficient $\langle v,\pi(\cdot)v\rangle$ of a unit highest weight vector $v$ of the lowest $K$-type of the
discrete series of parameter $\lambda=u+\rho_{\mathrm c}-\rho_{\mathrm n}$ (\S\ref{pds:ss-setup}), for which
$\Lambda^K(\lambda)=u$ (Lemma~\ref{psq:param}); so $\pspm_u(1)=1$; and for a family $u=(u_\tau)_\tau$ put $\pspm_u(h)=\prod_\tau\pspm_{u_\tau}(h_\tau)$
on $\pspG=\auU(X)_\infty$; the Haar measure on $\pspG$ is that of \S\ref{pds:ss-ds1}. On $W=W_j$ the
order $o$ is that of Definition~\ref{psv:order} (positive indices first, and within each sign block by
increasing $\varsigma_\tau$); all weights on $W$ are taken for this order.

\emph{Peaks.} A \emph{peak} of a bounded function $\Phi$ is a point $x$ with $|\Phi(x)|>\frac12\sup|
\Phi|$. Then $\Phi(x)\ne0$, $|\Phi|\le2|\Phi(x)|$ everywhere, and a point $x'$ with $|\Phi(x')|\ge|
\Phi(x)|$ is again a peak. A non-zero bounded function has a peak.

\begin{lemma}[Sizes]\label{psp:sizes}
\textup{(i)} $\pspn\le b$, and $2ft^2\le\pspn/t$. \textup{(ii)} $300\pspn+t\le\Delta/2$, $\Delta=2r^8$.
\textup{(iii)} $q_1\le x^{12}$, $d\le x^{24}$, $f\le x^{12}$, $t\le x^6$; hence
$24\,q_1^2\,d\,f\,t\le x^{80}\le b^{1/2}$. \textup{(iv)} $2^{25}\le b$, $t\ge14$, and
$100(b^{3/2}+4\pspn)\le b^2$. \textup{(v)} $5\,b^{12.75}\le b^{13}$.
\end{lemma}
\begin{proof}
By (S2), (S3) and (S1): $r\le3x^5$, $t\le15x^5\le x^6$, $f\le x^{12}$, so $\pspn\le2x^{12}x^{18}\le
b$; and $2ft^2=\pspn/t$. (ii): $300\pspn+t\le600\cdot125fr^3+5r\le r^8$ because $3r^3f\le r^7$ (S7)(b)
gives $75000fr^3\le25000r^7\le r^8/2$ ($r\ge x^5\ge10^5$). (iii): (S3); $\sqrt d\le\mathrm{rd}_E\le x^{12}$
by (S7)(e) and (S4); then $24q_1^2dft\le24x^{24+24+12+6}\le x^{80}$, and $x^{160}\le b$ by (S1).
(iv): $\log b\ge128$ (S1); $t\ge r\ge100$; $b^{3/2}+4\pspn\le5b^{3/2}$ by (i), and $500\le b^{1/2}$.
(v): $5\le b^{1/4}$.
\end{proof}

\subsubsection{The inputs}\label{psp:ss-inputs}

\emph{Holomorphic functions.} Let $X$ have signature $(p,q)$ at every place, with the order $o$ of
the frame putting the positive indices first, and let $u=(u_\tau)$ be a family of weights. A
continuous function $\phi$ on $\pspG$ is \emph{holomorphic of weight $u$} if for each $\tau$ it
extends to a holomorphic function $\tilde\phi$ on a neighbourhood of $\pspG\cdot B$ in
$\prod_\tau\auGL_N(\C)$, $B$ the group of invertible matrices upper triangular for $o$, with
$\tilde\phi(hb)=\prod_\tau\prod_ab_{\tau,aa}^{u_{\tau,a}}\phi(h)$. A function $\Phi$ on
$\auU(X)(\auA)$ is \emph{holomorphic of weight $u$} if $h\mapsto\Phi(gh)$ is, for every $g$. Let
$\mathrm{rot}(\theta)$, $\theta\in\R/\Z$, be the element $\mathrm{diag}(e^{2\pi i\theta o_a})_a$ at
every place.

\begin{proposition}[Inputs]\label{psp:inputs}
\textup{(H)} Let $X$ be of rank $\le t$.
\textup{(H1)} If $\phi,\phi'$ are holomorphic of any weights $u,u'$, then $\phi\phi'$ is holomorphic of
weight $u+u'$ \textup{(}Lemma~\ref{psq:H12}\textup{(i))}; and for every dominant weight $u$ of gap $\ge N^2+N+1$
\textup{(}$N$ the rank at each place\textup{)}, for which $\pspm_u$ is defined,
$\pspm_u$ is holomorphic of weight $u$ \textup{(}Proposition~\ref{psq:mhol}\textup{)}. In \textup{(H2)}, \textup{(H3)} let $u$ be dominant of gap $\ge300\pspn$. \textup{(H2)} If $\phi$ is holomorphic of
weight $u$ then $\phi(h\,\mathrm{rot}(\theta))=e^{2\pi i\theta\langle o,u\rangle}\phi(h)$ with
$\langle o,u\rangle=\sum_{\tau,a}o_au_{\tau,a}$, and
$\int_{\R/\Z}\phi(\mathrm{rot}(\theta)h)\,e^{-2\pi i\theta\langle o,u\rangle}d\theta=\phi(1)\,\pspm_u(h)$.
\textup{(H3)} If $\phi$ is holomorphic of weight $u$ and $\phi\,\aubar{\pspm_u}$ is integrable, then
$\int_{\pspG}\phi\,\aubar{\pspm_u}=\phi(1)\int_{\pspG}|\pspm_u|^2$.
\textup{(H4)} For $y\in e_0\poA(W_j)$ and every combination $F$ of the $F_i$, $\aubar y$ is
holomorphic of a weight $u$ and $F$ of the weight $u-u^\circ$, where $u^\circ$ is the weight of
$g\mapsto K^\circ(g,h^\circ)$; $u,u^\circ$ are dominant with gap $\ge300\pspn$, $|u_{\tau,a}|\le5b$.
The carrier $\varphi^\circ$ is holomorphic of weight $(0;3m)$ (entry $0$ on the positive, $3m$ on
the negative indices of $S$). \textup{(H5)} Let $\widetilde P$ be the fixed product of minors of
\S\ref{pte:ss-third} on $\auHom(S,\widetilde W)$ and $c_{\widetilde P}(h',h)=\pspip{v}{\Omega(h',h)v}$
its Fock coefficient \textup{(}$v$ the Fock vector of $\widetilde P$, $\Omega$ the real Weil
representation with the determinant twists of \eqref{pte:eq-F}\textup{)}. Then
$c_{\widetilde P}(h',1)=c_{\widetilde P}(1,1)\,\pspm_{u_2}(h')\,\aubar{\pspm_{u_3}(h')}$ for the
weights $u_2=(\Delta';\mu^+_{1}+r+2m)$, $u_3=(0;\mu^+_1+N''+2m)$, where $\Delta'=\Delta/2-(t+1)/2$ and
$\mu^+_1=\mu^+_{1,\tau}$ is the number of columns of $\mu^+_\tau$; and $u_3=(0;3m)+u_1+u_2$ with
$u_1=(-\Delta';6m+1)$.

\textup{(R)} Let $X$ be a space of rank $n_X\le t$, $u$ a dominant weight on $X$ of gap $\ge300\pspn$,
and $Y$ a positive definite space of rank $>2n_X$. There are finitely many classes $h_p$ of $\auU(Y)$
at level $\pi_M^2$, polynomials $Q_p$ on $\auHom(X,Y)\otimes\R$ in the variables of the positive
rows of $X$ of the weight $u$, and constants $a_p$, such that for every $g\in\auU(X)(\auA)$ and every
rapidly decreasing $\psi$ on $[\auU(X)]$ right invariant under the principal congruence subgroup of
level $\pi_M^2$,
\[
\int_{\pspG}\psi(gh)\,\pspm_u(h)\,dh=\sum_pa_p(g)\int_{[\auU(X)]}\psi(x)\,\theta^{X,Y}\bigl(x,h_p;
[Q_p,\delta_\iota]\bigr)\,dx ,
\]
with the point datum $\delta_\iota$ modulo $\pi_M^2$ at the inclusion. Here the characters of the kernel are
unramified outside $(\pi_M^2)$ and tamely ramified at the primes of $\pi_M$ (in the applications they are unramified at every finite place, Theorem~\ref{pss:a}), and the index $p$ runs over pairs (class, element of a basis of the space
of the $Q$'s), as in Theorem~\ref{psr:R}. Moreover the $Q_p$ may be taken
in any space of polynomials stable under $\auU(Y)_\infty$ containing a non-zero holomorphic vector of
weight $u$ (Definition~\ref{psr:def-hol}): for $(X,Y)=(S,W^\flat)$, $u=u_1$, the products of $2m$-minors with the short rows of the
classes of $\auU(W^\flat)$ (the polynomials $Q_r$ of \S\ref{pte:ss-third}); for $(X,Y)=(W,V^\circ)$,
$u=u^\circ$, the coordinates of the kernel $K^\circ$ against a basis of the Weyl module.

\textup{(P)} Let $\Phi$ be a bounded function on $\auU(X)(\auA)$, left invariant under $\auU(X)(E^+)$,
right invariant under the principal congruence subgroup of a level $\mathfrak p$ generated by
$p\in\auO_E$ with $\aubar p=-p$, $0<\Iml\tau(p)\le A'$ for all $\tau$, $A'\ge1$, and holomorphic of a
weight with gap $\ge1$ and entries in $[-A,A]$, $A\ge2$. Then every peak $g$ of $\Phi$ has
$\pspl(X,g)\ge(AA')^{-1/2}$.
\end{proposition}
\noindent Proposition~\ref{psp:inputs} is proved: (P), (H1)--(H3) in \S\ref{psq:sec}, (R) in \S\ref{psr:sec} with its
hypotheses verified in \S\ref{pss:sec} and \S\ref{psv:sec}, (H5) in \S\ref{pss:sec}, (H4) in \S\ref{pss:sec} and
\S\ref{psv:sec}; see Theorem~\ref{psv:main} for the list of published results used.

\emph{Archimedean bounds.} By \S\ref{pds:ss-FJ}--\S\ref{pds:ss-ds3}, whose proofs use only that the
gap $h$ of the parameter satisfies $h\ge N^2+1$ and the bound $L$ on its entries, for a dominant
weight $u$ on $\auU(p,q)$, $N=p+q\le t$, of gap $\ge300\pspn$, with $l=\lambda$ the parameter of
Lemma~\ref{psq:param}, of gap $h\ge300\pspn-t$ and entries $\le5b+t$:
\begin{equation}\label{psp:eq-ds}
|\pspm_u(g)|\le\bigl(2/\|g\|\bigr)^c\quad(0\le c\le h),\qquad
\int|\pspm_u|^2=d_l^{-1}>0,\qquad d_l\int\bigl(2/\|g\|\bigr)^{10N^3}dg\le2^{13t^3}b^{14t^2}
\end{equation}
when moreover $|u_a|\le5b$; the last bound is Theorem~\ref{pds:tail} with $\mathsf s=10N^3>4N^2$,
$L\le10b+t$, $\delta^{-1}=12N(L+2N)$: $4(5/\delta)^{2N^2}\mathsf s/(\mathsf s-4N^2)\cdot2^{10N^3}\le
8(60t(10b+3t))^{2t^2}2^{10t^3}\le2^{13t^3}b^{14t^2}$ (as $60t(10b+3t)\le b^7$, Lemma~\ref{psp:sizes}).
Here $\pspm_u$ is the coefficient $c_\tau$ of \S\ref{pds:sec} for the parameter $l$, and
$\int|\pspm_u|^2=d_l^{-1}$ is (DS1)(ii); $h\ge300\pspn-t\ge10N^3+6\pspn$, and the entries of $l$ are at most
$5b+t\le10b+t$, so the bounds above, and those used below ($c\le10N^3+6\pspn\le h$, $L\le10b+t$), hold.

\begin{lemma}[Integrals against $\pspm_{u'}$]\label{psp:int}
Let $X$ have rank $N\le t$, and let $u,u'$ be dominant of gap $\ge300\pspn$ with $|u_{\tau,a}|\le5b$.
Put $I_u=\int_{\pspG}|\pspm_u|^2>0$. Then $h\mapsto|\pspm_{u'}(h)|\,|h|^{6\pspn}$ is integrable and
$\int_{\pspG}|\pspm_{u'}(h)||h|^{6\pspn}dh\le b^{\pspn}I_u$. Consequently, if $A$ is continuous with
$|A(h)|\le C|h|^{6\pspn}$, then $A\pspm_{u'}$ is integrable and $|\int A\pspm_{u'}|\le Cb^{\pspn}I_u$.
\end{lemma}
\begin{proof}
Let $\tau_0$ be a place with $\|h_{\tau_0}\|=|h|$. By \eqref{psp:eq-ds} with $c=10N^3$ at $\tau\ne\tau_0$
and $c=10N^3+6\pspn\le300\pspn$ at $\tau_0$,
$|\pspm_{u'}(h)||h|^{6\pspn}\le2^{6\pspn}\prod_\tau(2/\|h_\tau\|)^{10N^3}$. Write
$2^{6\pspn}=2^{6\pspn}I_u\prod_\tau d_{l_\tau}$ ($l_\tau$ the parameter of Lemma~\ref{psq:param} for $u_\tau$,$I_u=\prod_\tau d_{l_\tau}^{-1}$).
By Fubini and \eqref{psp:eq-ds}, the integral is at most
$2^{6\pspn}I_u(2^{13t^3}b^{14t^2})^f=I_u\,2^{25ft^3}b^{14ft^2}\le I_u\,b^{ft^3}b^{ft^3}=b^{\pspn}I_u$,
using $2^{25}\le b$ and $14\le t$ (Lemma~\ref{psp:sizes}(iv)). The consequence is immediate.
\end{proof}

\subsubsection{The Fock model; decay of Heisenberg coefficients}\label{psp:ss-fock}

Let $\mathcal V=\C^D$ with a hermitian form $\pspip{x}{y}_\varepsilon=x^*\varepsilon y$,
$\varepsilon=\mathrm{diag}(\pm1)$, the symplectic form $\Iml\pspip xy_\varepsilon$, and the euclidean
norm $|x|$. Let $\mu$ be the Lebesgue measure normalised by $\int e^{-|x|^2}d\mu=1$, so that
$\int e^{-s|x|^2}d\mu=s^{-D}$ for $s>0$. As in Lemma~\ref{pcx:superpos}, the Fock space $\pspF$ is the
space of $c\in L^1(\mu)\cap L^2(\mu)$ with $c(x)=\int c(y)\beta(y,x)e(x-y)\,d\mu(y)$, where
$e(x)=e^{-|x|^2/2}$ is the vacuum and $\beta(y,x)=\exp(i\,\Iml\pspip yx_\varepsilon)$ (this exact
normalisation, matched with $e$ and $\mu$, is the one for which the reproducing identity, the
projection property and the computation of $\pspip{e}{\rho(z)e}$ in Lemma~\ref{psp:geom} hold; it is
fixed throughout); the Heisenberg translation is
$(\rho(y)c)(x)=\beta(y,x)c(x-y)$, and $\pspip{c}{c'}=\int\aubar cc'\,d\mu$. The polynomial vector of a
polynomial $P$ is $v_P(x)=P(\zeta(x))e(x)$, $\zeta(x)_a=\aubar{x_a}$ at the positive and $x_a$ at the
negative coordinates. The defining identity says $c(x)=\pspip{\rho(x)e}{c}$ (as $\aubar{\beta(x,y)}=
\beta(y,x)$ and $e$ is even), and $c=\int c(x)\rho(x)e\,d\mu(x)$ (Lemma~\ref{pcx:superpos}). The group
$\auU(\varepsilon)$ of complex matrices preserving $\pspip{}{}_\varepsilon$ preserves $\mu$ (as
$|\det\sigma|=1$) and the symplectic form; its Weil representation $\omega$ on $\pspF$ is unitary
and satisfies $\omega(\sigma)\rho(y)\omega(\sigma)^{-1}=\rho(\sigma y)$ ((W1)); it is determined up to
a character of modulus one, which does not affect the moduli below. $\|\sigma\|$ is the operator
norm; $\|\sigma\|\ge1$ since $\sigma^{-1}=\varepsilon\sigma^*\varepsilon$.

\begin{lemma}[Geometric form of the Weil action]\label{psp:geom}
For $\sigma\in\auU(\varepsilon)$ put $A_\sigma c=\int c(x)\,\rho(\sigma x)e\,d\mu(x)$. Then
$A_\sigma=\kappa(\sigma)\,\omega(\sigma)$ with $|\kappa(\sigma)|\ge\|\sigma\|^{-2D}$.
\end{lemma}
\begin{proof}
$A_\sigma$ is the composite of $c\mapsto c\circ\sigma^{-1}$, unitary on $L^2(\mu)$, with the operator
$c\mapsto\int c(x)\rho(x)e\,d\mu(x)$, which is the orthogonal projection onto $\pspF$ (its kernel is
$\aubar{(\rho(x')e)(x)}$-reproducing); so it is bounded. Since $\sigma$ preserves the symplectic form,
$\rho(\sigma y)\rho(\sigma x)=\beta'(y,x)\rho(\sigma(x+y))$ with the same cocycle $\beta'$ as for
$\rho(y)\rho(x)$; hence $\rho(\sigma y)A_\sigma\rho(y)^{-1}=A_\sigma$ (change of variables $x\mapsto x+y$
in $\int(\rho(y)^{-1}c)(x)\rho(\sigma x)e\,d\mu$). With (W1), $A_\sigma\omega(\sigma)^{-1}$ commutes with
all $\rho(y)$, so it is a scalar $\kappa(\sigma)$ by the irreducibility of $\rho$ on $\pspF$ (Stone--von
Neumann). Next $\pspip{e}{\rho(z)e}=\int e(x)\beta(z,x)e(x-z)\,d\mu(x)=e^{-|z|^2/2}$: completing the
square, $x=z/2+w$, the phase is $e^{i\Iml\pspip zw_\varepsilon}$ (as $\pspip zz_\varepsilon\in\R$) and
$\int e^{-|w|^2+i\Iml\pspip zw_\varepsilon}d\mu(w)=e^{-|z|^2/4}$. Hence
\[
\kappa(\sigma)\pspip{e}{\omega(\sigma)e}=\pspip{e}{A_\sigma e}=\int
e(x)\pspip{e}{\rho(\sigma x)e}\,d\mu(x)=\int e^{-(|x|^2+|\sigma x|^2)/2}d\mu(x)\ge
\Bigl(\frac{2}{1+\|\sigma\|^2}\Bigr)^D ,
\]
since $e$ is real. As $|\pspip e{\omega(\sigma)e}|\le1$ and $(1+\|\sigma\|^2)/2\le
\|\sigma\|^2$, $|\kappa(\sigma)|\ge\|\sigma\|^{-2D}$.
\end{proof}

\begin{lemma}[Decay]\label{psp:decay}
Let $P$ be a polynomial with $P(sx)=s^dP(x)$ for real $s$, $v=v_P$, $\sigma\in\auU(\varepsilon)$,
$y\in\mathcal V$, $y\ne0$, and let $A\ge d+D$ and $n\ge0$ be integers. Then
\[
\bigl|\pspip{v}{\rho(y)\omega(\sigma)v}\bigr|\le\|\sigma\|^{2D+2n}\,\bigl(4(A+n)\bigr)^n\,|y|^{-2n}\,
\pspip vv .
\]
\end{lemma}
\begin{proof}
By Lemma~\ref{psp:geom} and $\omega(\sigma)v=\kappa(\sigma)^{-1}\int v(x)\rho(\sigma x)e\,d\mu(x)$,
\[
\pspip{v}{\rho(y)\omega(\sigma)v}=\kappa(\sigma)^{-1}\int v(x)\pspip{v}{\rho(y)\rho(\sigma x)e}\,
d\mu(x),
\]
and $\rho(y)\rho(\sigma x)$ is a unimodular multiple of $\rho(\sigma x+y)$, while $|\pspip
v{\rho(z)e}|=|v(z)|$. So the left side is at most $\|\sigma\|^{2D}\int|v(x)||v(\sigma x+y)|\,d\mu(x)$.
Put $c=\sqrt{1+n/A}$ and $G(x)=c^{2d}|P(\zeta(x/c))|^2e^{-|x|^2/c^2}$. By homogeneity
$|v(x)|=G(x)^{1/2}e^{-a|x|^2/2}$ with $a=1-c^{-2}=n/(n+A)$, and by the change of variables $x\mapsto
cx$, $\int G\,d\mu=c^{2d+2D}\int|v|^2d\mu\le e^{n(d+D)/A}\pspip vv\le e^n\pspip vv$. Since $|y|\le
\|\sigma\||\sigma x+y|+\|\sigma\||x|$ (as $|y|\le|\sigma x+y|+|\sigma x|$ and $\|\sigma\|\ge1$),
$|y|^2\le2\|\sigma\|^2(|x|^2+|\sigma x+y|^2)$, so $\frac a2(|x|^2+|\sigma x+y|^2)\ge
T:=\frac{n|y|^2}{4(n+A)\|\sigma\|^2}$. With $2st\le s^2+t^2$ and $e^{-T}\le(n/eT)^n$,
\[
|v(x)||v(\sigma x+y)|\le\tfrac12\bigl(G(x)+G(\sigma x+y)\bigr)e^{-T}\le\tfrac12\bigl(G(x)+G(\sigma
x+y)\bigr)\Bigl(\frac{4(n+A)\|\sigma\|^2}{e|y|^2}\Bigr)^n .
\]
The map $x\mapsto\sigma x+y$ preserves $\mu$, so integrating gives at most $e^n\pspip vv\cdot
(4(n+A)\|\sigma\|^2/e|y|^2)^n$, which is the claim. (For $n=0$ read $e^{-T}\le1$.)
\end{proof}

\begin{lemma}[Separated sets]\label{psp:pack}
Let $\mathcal Y\subset\R^k$ be a set with $|y|\ge m$ and $|y-y'|\ge m$ for distinct $y,y'\in\mathcal Y$,
$m>0$, and let $s\ge k+2$ be an integer. Then $\sum_{y\in\mathcal Y}|y|^{-s}\le(5/m)^s$.
\end{lemma}
\begin{proof}
The balls $B(y,m/2)$ are disjoint and contained in $\{|z|\ge|y|/2\}$ (as $|y|\ge m$) and in
$\{|z|\le\frac32|y|\}$. For $z\in B(y,m/2)$, $|z|\le\frac32|y|$, so $|y|^{-s}\le(3/2)^s|z|^{-s}$, and
\[
\sum_y|y|^{-s}\le\frac{(3/2)^s}{\mathrm{vol}\,B(m/2)}\int_{|z|\ge m/2}|z|^{-s}dz=
\frac{(3/2)^s\,k\,\mathrm{vol}\,B(1)}{\mathrm{vol}\,B(1)(m/2)^k}\cdot\frac{(m/2)^{k-s}}{s-k}
=\frac{k}{s-k}\Bigl(\frac{3}{m}\Bigr)^s\le\Bigl(\frac5m\Bigr)^s ,
\]
since $k/(s-k)\le k/2\le(5/3)^s$ for $s\ge k+2$.
\end{proof}

\subsubsection{Theta kernels on lattices without short vectors}\label{psp:ss-sparse}

{\emergencystretch=6em Let $X,Y$ be spaces of ranks $n,n'$ and $\mathsf H=\auHom(X,Y)$. For $(g,h)\in\auU(X)(\auA)\times
\auU(Y)(\auA)$ the lattice of $\mathsf H$ at $(g,h)$ is $\Lambda_{g,h}=\{z\in\mathsf H(E):\allowbreak h^{-1}zg
\text{ is integral at all finite places}\}$, with $|z|_{g,h}$ the euclidean norm of the archimedean
coordinates of $h^{-1}zg$, and $\lambda_1(\mathsf H;g,h)$ its first minimum. Put $N_d=q_1d$; then
$N_d\in\mathfrak d_E$, because $\mathfrak d_E=\mathfrak d_{E/K}\mathfrak d_K$, $\mathfrak d_K=
(\sqrt{-d})$ divides $d$, and $\mathfrak d_{E/K}=(\Phi'_{q_1}(\zeta))$ (proof of
Lemma~\ref{pds:kappa}) divides $q_1$, as $\Phi_{q_1}$ divides $X^{q_1}-1$.\par}

\begin{lemma}[First minimum of $\auHom$]\label{psp:hom}
$\lambda_1(\mathsf H;g,h)\ge\pspl(X,g)\,\pspl(Y,h)/(fnN_d)$.
\end{lemma}
\begin{proof}
Identify $X(E)\otimes_\Q\R$ with $\prod_\tau\C^n$ by $\gamma\mapsto((g^{-1}\gamma)_\tau)_\tau$, with the
euclidean product $\Rel\sum_\tau a_\tau^*b_\tau$; $\Lambda_g$ becomes a lattice $L$ of rank $2fn$. At
the infinite places $g$ is unitary and the form is $\varepsilon=\mathrm{diag}(\pm1)$, so
$\mathrm{Tr}_{E/\Q}\pspip{\gamma}{\gamma'}=\Rel\sum_\tau a_\tau^*(2\varepsilon b_\tau)$. The hermitian
lattice $\Lambda_g$ is unimodular, so its dual for $\mathrm{Tr}_{E/\Q}\pspip{}{}$ is $\mathfrak
d_E^{-1}\Lambda_g$, and the euclidean dual lattice is $L^\vee=2\varepsilon(\mathfrak d_E^{-1}
\Lambda_g)$. As $\varepsilon$ is an isometry and $N_d\mathfrak d_E^{-1}\Lambda_g\subseteq\Lambda_g$,
$\lambda_1(L^\vee)\ge2\lambda_1(g)/N_d$. By Banaszczyk's transference theorem~\cite[Theorem~2.1]{Banaszczyk}
($\lambda_k(L)\lambda_1(L^\vee)\le k$ in rank $k$), the $2fn$-th successive
minimum satisfies $\lambda_{2fn}(L)\le fnN_d/\lambda_1(g)$. Choose $\Z$-linearly independent
$Y_1,\dots,Y_{2fn}\in\Lambda_g$ with $|Y_k|_g\le fnN_d/\lambda_1(g)$; they span $X(E)$ over $\Q$. Let
$0\ne z\in\Lambda_{g,h}$; then $zY_k\ne0$ for some $k$, and $zY_k\in\Lambda_h$ since $h^{-1}zY_k=
(h^{-1}zg)(g^{-1}Y_k)$. At each $\tau$ the operator norm of $(h^{-1}zg)_\tau$ is at most its euclidean
norm, hence at most $|z|_{g,h}$, so $\pspl(Y,h)\le\lambda_1(h)\le|zY_k|_h\le|z|_{g,h}|Y_k|_g\le
|z|_{g,h}fnN_d/\lambda_1(g)$. Since $\lambda_1(g)\ge\pspl(X,g)$, the claim follows.
\end{proof}

\begin{lemma}[Products of theta functionals]\label{psp:prodtheta}
Let $\mathcal W$ be a symplectic space over $E^+$ and $F,F'$ data (smooth at the infinite places).
Then $\theta(F)\,\aubar{\theta(F')}=\sum_{y\in\mathcal W(E^+)}\pspip{F'}{\rho(y)F}$, the sum being
absolutely convergent.
\end{lemma}
\begin{proof}
Both sides are unchanged under a unitary isomorphism of models commuting with $\rho$ (which carries
$\theta$ to $\theta$, by the uniqueness in (W1) and the normalisation of $\theta$), so we may use the
Schrödinger model of a polarisation $\mathcal W=\mathcal X\oplus\mathcal Y$: $F\in\mathcal S(\mathcal
X(\auA))$, $\theta(F)=\sum_{\xi\in\mathcal X(E^+)}F(\xi)$,
$(\rho(a,b)F)(x)=\psi(\pspip xb+\frac12\pspip ab)F(x+a)$, with the self-dual measure $dx$, for which
$\mathcal X(E^+)\backslash\mathcal X(\auA)$ has volume $1$ and Poisson summation
$\sum_{b\in\mathcal Y(E^+)}\hat\phi(b)\psi(\pspip cb)=\sum_\xi\phi(\xi-c)$ holds,
$\hat\phi(b)=\int\phi(x)\psi(\pspip xb)dx$. For fixed $a\in\mathcal X(E^+)$ and
$\phi(x)=\aubar{F'(x)}F(x+a)$,
\[
\sum_{b}\pspip{F'}{\rho(a,b)F}=\sum_b\hat\phi(b)\psi(\tfrac12\pspip ab)=\sum_{\xi}\aubar{F'(\xi-\tfrac a2)}
F(\xi+\tfrac a2).
\]
Summing over $a$ and putting $\xi'=\xi-\frac a2$, $\xi''=\xi+\frac a2$, which is a bijection of
$\mathcal X(E^+)^2$, gives $\sum_{\xi''}F(\xi'')\sum_{\xi'}\aubar{F'(\xi')}$. Absolute convergence: the
matrix coefficient $y\mapsto\pspip{F'}{\rho(y)F}$ of two smooth (Schwartz) data is a Schwartz--Bruhat
function on $\mathcal W(\auA)$, so it is summable over the lattice $\mathcal W(E^+)$; and smooth
Fock vectors correspond to Schwartz functions under the Bargmann transform.
\end{proof}

For the modulus $\pi_M^2$ and $T\in\pcMat_{n\times n'}(\auO_E/\pi_M^2)$ let
$\Psi_T(x)=\psi_{M}(\mathrm{tr}(xT))$, where $\psi_M$ is a primitive additive character of
$\auO_E/\pi_M^2$; then $\frac1{\#}\sum_T\Psi_T(x)=[x=0]$ ($\#$ the number of $T$).

\begin{lemma}[Averaged coefficients]\label{psp:avg}
Let $v,w$ be smooth Fock vectors, $p=(g,h)$, and $z\in\mathsf H(E)$; put $z'=h^{-1}zg$. Then
\[
\frac1\#\sum_T\pspip{\omega(p)[v,\Psi_T]}{\rho(z)\,\omega(p)[w,\Psi_T]}=\chi(z)\,\pspip{v}{\rho(z'_\infty)
w}\cdot\mathbf 1\bigl[z'\in\alpha\auHom(L_X,L_Y)\text{ and }z'_{\mathfrak w}\equiv0\bmod\pi_M^2\
(\mathfrak w\mid\pi_M)\bigr],
\]
with $|\chi(z)|=1$ and $\chi(0)=1$. Every $z\ne0$ for which the right side is non-zero satisfies
$(q_1/\pi_M^2)z\in\Lambda_{g,h}$.
\end{lemma}
\begin{proof}
$\omega(p)$ is unitary and $\omega(p)^{-1}\rho(z)\omega(p)=\rho(p^{-1}z)$ with $p^{-1}z=z'$ ((W1)), so
we may take $p=1$, $z=z'$. The pairing factors over the places. At the finite places off $\pi_M$ the
datum is the vacuum $\psvac$ of $L=\alpha\auHom(L_X,L_Y)$; by (W4), $\rho(z)\psvac$ is a unimodular
multiple of $\psvac$ if $z\in L$ and lies on the line of a different character of $L$ otherwise,
hence is orthogonal to $\psvac$. At a place $v=\mathfrak w\aubar{\mathfrak w}$ with $\mathfrak w\mid
\pi_M$ we use (W5): functions on $\pcMat(E_{\mathfrak w})$, $\rho(a,0)$ the translation and $\rho(0,a')$
multiplication by $x\mapsto[x,a']$, with the datum $\Psi_T^\natural$ supported on integral matrices.
Up to a unimodular factor depending on $z$ only,
\[
\pspip{\Psi_T^\natural}{\rho(z)\Psi_T^\natural}=\int_{x,\,x+z_{\mathfrak w}\text{ integral}}
\aubar{\Psi_T(x)}\Psi_T(x+z_{\mathfrak w})[x,z_{\aubar{\mathfrak w}}]\,dx=
\mathbf 1[z_{\mathfrak w},z_{\aubar{\mathfrak w}}\text{ integral}]\;\Psi_T(z_{\mathfrak w}),
\]
because $\Psi_T$ is additive, the integral matrices have volume $1$ and are their own dual for
$[\,,\,]$ ($\psi$ is unramified at $p_0,p_1$, which do not divide the different, and $\alpha=1$
there, since the primes of $\alpha$ divide $(1+d)/4<b^{10}\le p_0,p_1$ by Lemma~\ref{psp:sizes}). The
average over $T$ of $\Psi_T(z_{\mathfrak w})$ is $[z_{\mathfrak w}\equiv0\bmod\pi_M^2]$. At $z=0$ all
factors are $1$ (unit vectors; Lemma~\ref{psq:unitvac}). For the last assertion: $q_1\alpha$ is integral ($\alpha=1/P'(\zeta)$
and $P'(\zeta)$ divides $q_1$), $\pi_M$ is a unit at the places $\aubar{\mathfrak w}$ and off $M$,
and $z_{\mathfrak w}\in\pi_M^2\pcMat(\auO_{\mathfrak w})$; so $(q_1/\pi_M^2)z$ is integral everywhere.
\end{proof}

\begin{proposition}[Sparse kernels]\label{psp:sparse}
Let $\alpha',\beta'$ be characters as in the definition of the kernel $\theta_{\alpha',\beta'}$
(\S\ref{pc:ss-data}), $Q$ a polynomial on $\mathsf H\otimes\R$, homogeneous of degree $d_Q$ for real
scalars, $v$ its Fock vector, and $c_Q(u,u')=\pspip{v}{\Omega(u,u')v}$ for $(u,u')\in\auU(X)_\infty
\times\auU(Y)_\infty$, $\Omega$ the real Weil representation with the archimedean components of the
characters. Let $D=fnn'$, let $q>D$ be an integer and $m=c_\psi M\lambda_1(\mathsf H;g,h)/q_1$. There
is a function $\Psi$ modulo $\pi_M^2$ such that for all $u,u'$
\[
\bigl|\theta_{\alpha',\beta'}\bigl(gu,hu';[Q,\Psi]\bigr)-c_Q(u,u')\bigr|\le(|u||u'|)^{2D+2q}
\Bigl(\frac{100(d_Q+2q)}{m^2}\Bigr)^q c_Q(1,1).
\]
\end{proposition}
\begin{proof}
Let $k_T=\theta_{\alpha',\beta'}(g,h;[Q,\Psi_T])$ and $\Psi=\frac1\#\sum_T\aubar{k_T}\Psi_T$. The kernel is
linear in the finite datum, and $\omega(u,u')[v,\Psi_T]=[\omega_\infty(u,u')v,\Psi_T]$; the characters
contribute $\chi_1=\alpha'(\det g)^{-1}\beta'(\det h)$, of modulus $1$, at $(g,h)$, and their
archimedean components at $(u,u')$ are absorbed in $\Omega$. With $w=\Omega(u,u')v$,
$F_T=\omega(p)[v,\Psi_T]$, $F'_T=\omega(p)[w,\Psi_T]$, $p=(g,h)$:
\[
\theta_{\alpha',\beta'}(gu,hu';[Q,\Psi])=\frac1\#\sum_T\aubar{\chi_1\theta(F_T)}\,\chi_1\theta(F'_T)
=\sum_{z\in\mathsf H(E)}\frac1\#\sum_T\pspip{F_T}{\rho(z)F'_T}
\]
by Lemma~\ref{psp:prodtheta}. By Lemma~\ref{psp:avg} the term of $z=0$ is $\pspip vw=c_Q(u,u')$, and
the other terms vanish unless $z\in\mathcal Z=\{z\ne0:(q_1/\pi_M^2)z\in\Lambda_{g,h}\}$. For
$z\in\mathcal Z$ let $x_z$ be the Fock coordinates of $z'_\infty$; as $|\tau(q_1/\pi_M^2)|=q_1/M$,
$|x_z|\ge c_\psi(M/q_1)\lambda_1(\mathsf H;g,h)=m$, and likewise $|x_z-x_{z'}|\ge m$ for $z\ne z'$ in
$\mathcal Z$ (their difference is in $\mathcal Z$). The term of $z$ has modulus at most
$|\pspip{v}{\rho(x_z)\omega_\infty(u,u')v}|$, and $\omega_\infty(u,u')$ is the Weil operator of
$\sigma\colon x\mapsto u'xu^{-1}$, with $\|\sigma\|\le|u'||u|$ (the euclidean norm of $u'xu^{-1}$ is at
most $\|u'\|\|u^{-1}\|$ times that of $x$, and $\|u^{-1}\|=\|u\|$). The Fock model of
\S\ref{pk:ss-conv} is identified with that of \S\ref{psp:ss-fock} (complex dimension $D$) by
$F\mapsto(x\mapsto\pspip{\rho(x)e}{F})$, a unitary map commuting with $\rho$ and $\omega$, carrying
$Q$ to $v_Q$ up to the conjugation of coordinates fixed there. By Lemma~\ref{psp:decay} with
$A=d_Q+D$ and $n=q$, and Lemma~\ref{psp:pack} with $k=2D$, $s=2q\ge2D+2$,
\[
|\theta-c_Q(u,u')|\le\|\sigma\|^{2D+2q}\bigl(4(d_Q+D+q)\bigr)^q\pspip vv\sum_{z\in\mathcal Z}|x_z|^{-2q}
\le(|u||u'|)^{2D+2q}\bigl(4(d_Q+2q)\bigr)^q(5/m)^{2q}c_Q(1,1),
\]
which is the claim ($4\cdot25=100$, $c_Q(1,1)=\pspip vv$).
\end{proof}

\begin{corollary}\label{psp:sparse2}
If moreover $n,n'\le t$, $d_Q\le b^{3/2}$ and $\pspl(X,g)\pspl(Y,h)\ge b^{-13}$, then $\Psi$ can be
chosen with
\[
\bigl|\theta_{\alpha',\beta'}(gu,hu';[Q,\Psi])-c_Q(u,u')\bigr|\le b^{-20\pspn}(|u||u'|)^{6\pspn}
c_Q(1,1)\qquad\text{for all }u,u'.
\]
\end{corollary}
\begin{proof}
Take $q=2\pspn$. Then $D\le ft^2\le\pspn<q$ and $2D+2q\le6\pspn$ (Lemma~\ref{psp:sizes}(i)); as
$|u|,|u'|\ge1$ we may raise the exponent to $6\pspn$. By Lemma~\ref{psp:hom},
$\lambda_1(\mathsf H;g,h)\ge b^{-13}/(ftN_d)$, so with $c_\psi=\sqrt{2\pi/\sqrt d}\ge x^{-6}\ge
b^{-1/2}$, $M\ge b^{20}$ (Lemma~\ref{an:lem-sizes}(S5)) and Lemma~\ref{psp:sizes}(iii),
$m\ge b^{-1/2}b^{20}b^{-13}/(q_1^2dft)\ge b^{6}$. Then $100(d_Q+4\pspn)/m^2\le b^{-10}$ by
Lemma~\ref{psp:sizes}(iv), and $(b^{-10})^{2\pspn}=b^{-20\pspn}$.
\end{proof}

\subsubsection{The isolation estimate}\label{psp:ss-isol}

\begin{proposition}[Isolation]\label{psp:isol}
Let $X$ have rank $\le t$, let $u,u'$ be dominant weights of gap $\ge300\pspn$ with $|u_{\tau,a}|\le5b$,
and $I_u=\int_{\pspG}|\pspm_u|^2$. Let $\Phi$ be a continuous bounded function on $\auU(X)(\auA)$,
holomorphic of weight $u$, with a peak $z$; let $K_0$ be a function on $\pspG$ with $|K_0|\le C$, and
$\delta\ge0$. Suppose that for every peak $x$ of $\Phi$ there is a function $K_x$ on $\pspG$,
holomorphic of weight $u-u'$, with $|K_x(h)-K_0(h)|\le\delta|h|^{6\pspn}$ for all $h$ and
\[
\int_{\pspG}\aubar{\Phi(xh)}\,K_x(h)\,\pspm_{u'}(h)\,dh=0 .
\]
Then $|K_0(1)|\le(4b^{\pspn}+1)\delta$.
\end{proposition}
\begin{proof}
Put $e(\theta)=e^{2\pi i\theta\langle o,u\rangle}$. By (H2), $|\Phi(z\,\mathrm{rot}(\theta))|=|\Phi(z)|$,
so every $z\,\mathrm{rot}(\theta)$ is a peak; let $K_\theta=K_{z\mathrm{rot}(\theta)}$ and
$W_\theta(h)=\aubar{\Phi(z\,\mathrm{rot}(\theta)h)}\,e(\theta)$. Then $|W_\theta|\le2|\Phi(z)|$, $W$ is
jointly continuous, $\int W_\theta K_\theta\pspm_{u'}=0$, and by (H2) applied to $\phi=\Phi(z\,\cdot)$,
$\int_{\R/\Z}W_\theta(h)\,d\theta=\aubar{\Phi(z)\pspm_u(h)}$. Let $K_*=K_{\theta=0}$. Since
$|K_\theta(h)|\le(C+\delta)|h|^{6\pspn}$ for all $\theta$ and the $K_\theta$ are continuous,
Lemma~\ref{psp:int} shows that $K_*\pspm_{u'}$ and $K_\theta\pspm_{u'}W_\theta$ are integrable, and
\[
\Bigl|\int K_*\pspm_{u'}W_\theta\Bigr|=\Bigl|\int(K_*-K_\theta)\pspm_{u'}W_\theta\Bigr|\le
4\delta\,|\Phi(z)|\,b^{\pspn}I_u ,
\]
as $|(K_*-K_\theta)W_\theta|\le4\delta|\Phi(z)||h|^{6\pspn}$. The function $(\theta,h)\mapsto
K_*(h)\pspm_{u'}(h)W_\theta(h)$ is integrable on $\R/\Z\times\pspG$ ($W$ bounded), so by Fubini
\[
\int_{\R/\Z}\int K_*\pspm_{u'}W_\theta\,dh\,d\theta=\aubar{\Phi(z)}\int K_*\pspm_{u'}\,\aubar{\pspm_u}
=\aubar{\Phi(z)}\,K_*(1)\,I_u ,
\]
by (H3): $K_*\pspm_{u'}$ is holomorphic of weight $u$ by (H1), equals $K_*(1)$ at $1$
($\pspm_{u'}(1)=1$), and $K_*\pspm_{u'}\aubar{\pspm_u}$ is integrable as $|\pspm_u|\le1$. Hence
$I_u|\Phi(z)||K_*(1)|\le4\delta|\Phi(z)|b^{\pspn}I_u$; as $I_u>0$ and $\Phi(z)\ne0$,
$|K_*(1)|\le4\delta b^{\pspn}$, and $|K_0(1)|\le|K_*(1)|+\delta$ since $|1|=1$.
\end{proof}

\subsubsection{The four steps}\label{psp:ss-outline}

Let $\widetilde P$ be the product of minors on $\auHom(S,\widetilde W)$ of \S\ref{pte:ss-third} (the
factor of $P_i$ in the rows of $\widetilde W$), of degree $d_{\widetilde P}$, and for a choice $i$ write
$P_i=\widetilde P\cdot Q_{r(i)}$, where $Q_{r(i)}$ is the product of the determinants on $W^\flat$
attached to the class $h^\flat_i$ and the row selection of $i$; the choice $i$ is the triple
$(\Lambda_i,h^\flat_i,\text{rows})$ (\S\ref{pte:ss-third}). Let $\theta^{\widetilde W}[\Psi]$ be the kernel
of $(S,\widetilde W)$ with datum $[\widetilde P,\Psi]$, $\Psi$ a function modulo $\pi_M^2$, and with the
characters of \eqref{pte:eq-F} restricted to $\widetilde W$; write $(g,1)$ for the point $(g,1,1)$ of
$\auU(\widetilde W)$. Let $u_1=(-\Delta';6m+1)$ and $u_3$ be as in (H5).

\begin{lemma}[Weights]\label{psp:weights}
$u_1$ and $u_3$ are dominant of gap $\ge300\pspn$ on $S$, and $|u_{1,a}|,|u_{3,a}|\le5b$; and
$d_{\widetilde P}\le b^{3/2}$.
\end{lemma}
\begin{proof}
Both are constant on each sign block, so dominant. Gap of $u_1$: $\Delta'+6m+1\ge\Delta/2-(t+1)/2\ge
300\pspn$ by Lemma~\ref{psp:sizes}(ii). Gap of $u_3$: $\mu^+_1\ge\Delta'$ (the diagram $\mu^+_\tau$ has
$\Delta'$ columns of height $r_+$, \S\ref{pte:ss-third}). Entries: $\mu^+_1\le\kappa_j(\tau)\le2b$
(the largest part of $\mu^+_\tau$ is $\kappa_j(\tau)$, proof of Proposition~\ref{pte:third}, and
$\kappa_j(\tau)\le\kappa(\tau)+m\Delta\le2b$ by (S5), (S7)(g)), and $N''+2m=r+11m+1\le b$; for $u_1$, $\Delta'\le\Delta/2=r^8\le3^8x^{40}\le x^{160}\le b$ (by (S2) and (S1)) and
$6m+1\le N''+2m\le b$. The degree: $d_{\widetilde P}\le7fmb\le b^{3/2}$ by Lemma~\ref{psq:degree}.
\end{proof}

\begin{lemma}[Step 1]\label{psp:step1}
For $g_1\in\auU(S)(\auA)$ and every $\Psi$ there are constants $c_i$ with
\[
\sum_ic_iF_i(g)=\int_{\auU(S)_\infty}\varphi^\circ(g_1h)\,\theta^{\widetilde W}[\Psi]\bigl(g_1h,(g,1)\bigr)
\,\pspm_{u_1}(h)\,dh\qquad(g\in\auU(W)(\auA)).
\]
\end{lemma}
\begin{proof}
Fix $g$ and put $\psi(x)=\varphi^\circ(x)\theta^{\widetilde W}[\Psi](x,(g,1))$, rapidly decreasing
($\varphi^\circ$ is, Lemma~\ref{pte:carrierprop}, and kernels have moderate growth, Lemma~\ref{pcx:thetagrowth}) and right
invariant under the principal congruence subgroup of level $\pi_M^2$ (Lemma~\ref{pc:level}). By (R)
for $(X,Y)=(S,W^\flat)$, $\operatorname{rk}W^\flat=6m+1>6m$, $u=u_1$ (Lemma~\ref{psp:weights}), the
right side is $\sum_ra_r(g_1)\int_{[\auU(S)]}\psi(x)\theta^{W^\flat}(x,h^\flat_r;[Q_r,\delta_\iota])dx$,
$r$ running over the pairs (class, rows); the constants do not depend on $g$. The product of the
kernels of $(S,\widetilde W)$ and $(S,W^\flat)$ is the kernel of $(S,\widetilde W\oplus W^\flat)=(S,W')$
with the product datum $[\widetilde PQ_r,\Psi\otimes\delta_\iota]$ ((W7) of \S\ref{po:ss-glue}, the
compatibility with orthogonal sums of the target; the characters are those of \eqref{pte:eq-F}; the factor
$(\nu^\flat\upsilon^m)(\det x)^{-1}$ is attached to the kernel of $(S,\widetilde W)$, so that the kernel of
$(S,W^\flat)$ has the characters $(1,\upsilon^m)$, Lemma~\ref{psv:split}). By
finite Fourier analysis on $\pcMat_{3m\times N''}(\auO_E/\pi_M^2)$, $\Psi\otimes\delta_\iota=\sum_\Lambda
f(\Lambda)\psi_\Lambda$ for some $f$, and the kernel is linear in the datum. So the right side is
$\sum_{r,\Lambda}a_r(g_1)f(\Lambda)\int\varphi^\circ(x)\theta^{S,W'}(x,(g,1,1,h^\flat_r);[P_{(\Lambda,r)},
\psi_\Lambda])dx=\sum_{r,\Lambda}a_r(g_1)f(\Lambda)F_{(\Lambda,r)}(g)$ by \eqref{pte:eq-F}.
\end{proof}

\begin{lemma}[Step 2]\label{psp:step2}
Under the hypothesis of Proposition~\ref{psp:main}, for every combination $F=\sum_ic_iF_i$ and every
$x\in\auU(W)(\auA)$, $\int_{\pspG}y(xh)F(xh)\pspm_{u^\circ}(h)\,dh=0$.
\end{lemma}
\begin{proof}
$\psi=yF$ is rapidly decreasing ($y$ is cuspidal, $F$ of moderate growth) and right invariant under
the principal congruence subgroup of level $\pi_M^2$. By (R) for $(X,Y)=(W,V^\circ)$,
$\operatorname{rk}V^\circ>2r$, $u=u^\circ$, the integral is $\sum_pa_p(x)\int_{[\auU(W)]}y\,F\,
\pspip{\zeta_p}{K^\circ(\cdot,h^\circ_p)}$ for vectors $\zeta_p$ of the Weyl module, and each term is
$\sum_ic_i\pspip{\zeta_p}{\int yF_iK^\circ(\cdot,h^\circ_p)}=0$ by hypothesis (the integrals converge
absolutely, Proposition~\ref{po:exchange}).
\end{proof}

\begin{proposition}[Step 3]\label{psp:approx}
There are a bounded function $F^0$ on $\auU(W)_\infty$ with $F^0(1)\ne0$ such that for every
$g\in\auU(W)(\auA)$ with $\pspl(W,g)\ge(5b^{24})^{-1/2}$ some combination $F=\sum_ic_iF_i$ satisfies
\[
|F(gh)-F^0(h)|\le2b^{-19\pspn}|F^0(1)|\,|h|^{6\pspn}\qquad(h\in\auU(W)_\infty).
\]
\end{proposition}
\begin{proof}
\emph{A peak of the carrier.} $\varphi^\circ$ is bounded, non-zero, left invariant under the rational
points, right invariant under $K_S$ (Lemmas~\ref{pte:carrierprop} and~\ref{pte:carriernz}), hence under the principal
congruence subgroup of level $p=\pm\sqrt{-d}$, the sign chosen as in Remark~\ref{psq:Papp} ($\aubar p=-p$, $\Iml\tau(p)=\sqrt d\le x^{12}\le
b^{1/2}$), and holomorphic of weight $(0;3m)$ (H4), of gap $3m\ge1$ and entries in $[-5b,5b]$. By
(P) with $A=5b$, $A'=b^{1/2}$ it has a peak $g_1$ with $\pspl(S,g_1)\ge(5b^{3/2})^{-1/2}$.

\emph{The function $F^0$.} Let $c(h',h)=c_{\widetilde P}(h',\iota h)$, $\iota h=(h,1,1)$, the Fock
coefficient of (H5), and $J=I_{u_3}>0$. By unitarity and Cauchy--Schwarz $|c|\le c(1,1)=\pspip vv>0$
($\widetilde P\ne0$). Put
\[
F^0(h)=\int_{\auU(S)_\infty}\varphi^\circ(g_1h')\,c(h',h)\,\pspm_{u_1}(h')\,dh' .
\]
By Lemma~\ref{psp:int} (with $u=u_3$, $u'=u_1$, Lemma~\ref{psp:weights}) and $|\varphi^\circ(g_1h')|\le2
|\varphi^\circ(g_1)|$, the integral converges and $|F^0(h)|\le2|\varphi^\circ(g_1)|c(1,1)b^{\pspn}J$. By
(H5) and $u_3=(0;3m)+u_1+u_2$, the function $h'\mapsto\varphi^\circ(g_1h')\pspm_{u_1}(h')c(1,1)
\pspm_{u_2}(h')$ is holomorphic of weight $u_3$ ((H1), (H4); $\pspm_{u_2}$ is defined and holomorphic of weight $u_2$
by (H1), $u_2$ being constant on each sign block, hence dominant, and its gap being $\ge9m^2+3m+1=N^2+N+1$ with $N=3m$, as checked in the proof of
Theorem~\ref{pss:H5}), so by (H3)
$F^0(1)=J\,\varphi^\circ(g_1)\,c(1,1)\ne0$.

\emph{Approximation.} Let $g$ be as in the statement. The lattice of $\widetilde W$ at $(g,1)$ is the
orthogonal sum of $\Lambda_g$ and the standard lattice $\auO_E^{3m}$ of the padding space $(E^j)^-\oplus E^{3m-j}$, whose non-zero vectors
have length $\ge1$ (a non-zero integral vector $\gamma'$ has $\sum_\tau|\tau\gamma'|^2\ge f\ge1$ by the
arithmetic--geometric mean inequality and $|\mathrm N\gamma'_a|\ge1$); so $\pspl(\widetilde W,(g,1))\ge
\pspl(W,g)$ and $\pspl(S,g_1)\pspl(\widetilde W,(g,1))\ge(5b^{3/2})^{-1/2}(5b^{24})^{-1/2}\ge b^{-13}$
(Lemma~\ref{psp:sizes}(v)). By Corollary~\ref{psp:sparse2} for $(X,Y)=(S,\widetilde W)$ (ranks $3m$ and
$r+3m\le t$, $d_{\widetilde P}\le b^{3/2}$) there is $\Psi$ with
$|\theta^{\widetilde W}[\Psi](g_1h',(g,1)\iota h)-c(h',h)|\le b^{-20\pspn}(|h'||h|)^{6\pspn}c(1,1)$, using
$|\iota h|\le|h|$. By Lemma~\ref{psp:step1} some $F=\sum c_iF_i$ has
$F(gh)=\int\varphi^\circ(g_1h')\theta^{\widetilde W}[\Psi](g_1h',(gh,1))\pspm_{u_1}(h')dh'$, and
$(gh,1)=(g,1)\iota h$. Hence $F(gh)-F^0(h)=\int\varphi^\circ(g_1h')A(h')\pspm_{u_1}(h')dh'$ with
$|\varphi^\circ(g_1h')A(h')|\le2|\varphi^\circ(g_1)|\,b^{-20\pspn}|h|^{6\pspn}c(1,1)\,|h'|^{6\pspn}$, and
Lemma~\ref{psp:int} gives
\[
|F(gh)-F^0(h)|\le2|\varphi^\circ(g_1)|c(1,1)J\,b^{-20\pspn}b^{\pspn}|h|^{6\pspn}=2b^{-19\pspn}|F^0(1)|
\,|h|^{6\pspn}. \qedhere
\]
\end{proof}

\subsubsection{Proof of Proposition~\ref{psp:main}}\label{psp:ss-proof}

Suppose $y\ne0$ and let $\Phi=\aubar y$, a bounded continuous non-zero function on $\auU(W)(\auA)$
(cuspidal, hence bounded), holomorphic of the weight $u$ of (H4). It is right invariant under the
principal congruence subgroup of level $\pi_M$ (the level of $\poA(W_j)$, \S\ref{po:ss-set}), which
contains that of level $p=\pm M\sqrt{-d}$ (sign as in Remark~\ref{psq:Papp}); $\aubar p=-p$ and $0<\Iml\tau(p)=M\sqrt d\le b^{22}x^{12}\le
b^{23}$ ((S5)). By (P) with $A=5b$, $A'=b^{23}$, every peak $x$ of $\Phi$ has $\pspl(W,x)\ge
(5b^{24})^{-1/2}$. Let $F^0$ be as in Proposition~\ref{psp:approx}, $|F^0|\le C$, and for a peak $x$
let $K_x(h)=F(xh)$ with $F$ the combination given there for $g=x$. Then $K_x$ is holomorphic of weight
$u-u^\circ$ (H4), $|K_x-F^0|\le\delta|h|^{6\pspn}$ with $\delta=2b^{-19\pspn}|F^0(1)|$, and
$\int\aubar{\Phi(xh)}K_x(h)\pspm_{u^\circ}(h)dh=\int y(xh)F(xh)\pspm_{u^\circ}(h)dh=0$ by
Lemma~\ref{psp:step2}. $\Phi$ has a peak. Proposition~\ref{psp:isol} (with $u'=u^\circ$; the bounds on
$u,u^\circ$ are part of (H4)) gives
\[
|F^0(1)|\le(4b^{\pspn}+1)\cdot2b^{-19\pspn}|F^0(1)| .
\]
With $B=b^{\pspn}>10$, $(4B+1)\cdot2<B^2\le B^{19}$, so $(4B+1)2B^{-19}<1$, and $F^0(1)=0$, a
contradiction. Hence $y=0$. \qed

\subsubsection{What this proof rests on}\label{psp:ss-rest}

Proposition~\ref{psp:main} is proved from Lemmas~\ref{psp:sizes}--\ref{psp:step2},
Propositions~\ref{psp:sparse}, \ref{psp:isol}, \ref{psp:approx}, the cited results of the paper
(Lemma~\ref{an:lem-sizes}, \S\ref{pds:sec}, Lemma~\ref{pcx:superpos}, (W1)--(W5), (W7),
Lemma~\ref{pc:level}, Lemmas~\ref{pte:carrierprop} and~\ref{pte:carriernz}, \eqref{pte:eq-F}, Proposition~\ref{po:exchange}),
Banaszczyk's transference theorem and the Stone--von Neumann theorem, and the inputs (R), (H), (P) of
Proposition~\ref{psp:inputs}, proved in \S\ref{psq:sec}--\S\ref{psv:sec} (Theorem~\ref{psv:main}). With (G1) and Proposition~\ref{po:span},
the injectivity of $\Theta^\circ$ in (G2) of Proposition~\ref{po:glue} holds (Remark~\ref{psv:G2}).


%% ===== proofs/C-span2.tex =====
% proofs/C-span2.tex
% Inputs of Proposition psp:inputs (C-span.tex). To be \input after C-span.tex.
% Labels of this file start with psq:. Only amsmath/amssymb/amsthm.
\providecommand{\auO}{\mathcal O}
\providecommand{\auU}{\mathrm U}
\providecommand{\auGL}{\mathrm{GL}}
\providecommand{\auA}{\mathbb A}
\providecommand{\auPl}{\Phi}
\providecommand{\aubar}[1]{\overline{#1}}
\providecommand{\pspl}{\lambda}
\providecommand{\pspG}{G_\infty}
\providecommand{\pspip}[2]{\langle #1,#2\rangle}
% macros of this file (prefix \psq)
\providecommand{\psqH}{\mathfrak H}

\subsection{Inputs for products span}\label{psq:sec}

This subsection proves the inputs (P), (H1), (H2), (H3) of Proposition~\ref{psp:inputs}
(\S\ref{psq:ss-P}, \S\ref{psq:ss-H}, \S\ref{psq:ss-H3}), and four facts used in \S\ref{psp:sec}
without written proof: the normalisation of $\theta$ in Lemma~\ref{psp:prodtheta}
(Lemma~\ref{psq:thetanorm}) and the unit vacuum (Lemma~\ref{psq:unitvac}), the level of
$e_0\poA(W_j)$ used in \S\ref{psp:ss-proof} (Remark~\ref{psq:level}), and the degree of
$\widetilde P$ used in Lemma~\ref{psp:weights} (Lemma~\ref{psq:degree}). The rest of
Proposition~\ref{psp:inputs}, namely (H4), (H5), (R), is proved in \S\ref{psr:sec}--\S\ref{psv:sec}.

\emph{Conventions.} $X$ is a hermitian space of rank $n$ over $E$ with a frame as in
Definition~\ref{pcx:def-ht}: in the frame the form has the matrix $\varepsilon=\mathrm{diag}(s_1,\dots,
s_n)$, $s_a=\pm1$, at every infinite place, and at the finite places a matrix $S_f\in\auGL_n(\widehat
\auO_E)$ (the lattice of integral vectors is self-dual; this is the unimodularity used in fact (R3) stated after
Definition~\ref{pcx:def-ht}). Forms are linear in the first variable ((Cv1)): $\langle w,w'\rangle=w'^*Sw$ in
coordinates. The order $o$ is a numbering of the indices putting the positive ones first; we number
the indices $1,\dots,n$ in this order, so that $s_a=+1$ for $a\le n_+$ and $s_a=-1$ for $a>n_+$. $B$
is the group of invertible matrices upper triangular in this order (at every place), and for a
diagonal or upper triangular $b$ and $u\in\Z^n$, $b^u=\prod_ab_{aa}^{u_a}$. Holomorphy of weight $u$
is that of \S\ref{psp:ss-inputs}. For a non-zero $p\in\auO_E$ the \emph{principal congruence subgroup
of level $p$} is the group $K(p)$ of the $k\in\auU(X)(\auA_f)$ whose matrix in the frame is integral
and $\equiv1\bmod p$.

\subsubsection{Holomorphic functions: the formal part}\label{psq:ss-H}

\begin{lemma}[(H1), first part, and (H2), first part]\label{psq:H12}
\textup{(i)} If $\phi,\phi'$ are holomorphic of weights $u,u'$ on $\pspG$, then $\phi\phi'$ is
holomorphic of weight $u+u'$. \textup{(ii)} If $\phi$ is holomorphic of weight $u$, then
$\phi(h\,\mathrm{rot}(\theta))=e^{2\pi i\theta\langle o,u\rangle}\phi(h)$ for $h\in\pspG$, $\theta\in\R/\Z$.
\end{lemma}
\begin{proof}
(i) If $\tilde\phi,\tilde\phi'$ are holomorphic extensions on neighbourhoods $U,U'$ of $\pspG B$, then
$\tilde\phi\tilde\phi'$ is holomorphic on $U\cap U'$, a neighbourhood of $\pspG B$, and
$\tilde\phi\tilde\phi'(hb)=b^ub^{u'}\phi(h)\phi'(h)=b^{u+u'}(\phi\phi')(h)$. (ii) $\mathrm{rot}(\theta)$ is
diagonal and unitary at every place, so it lies in $\pspG\cap B$; hence $h\,\mathrm{rot}(\theta)\in\pspG$
and, reading $h\,\mathrm{rot}(\theta)$ as $h\cdot b$ with $b=\mathrm{rot}(\theta)$,
$\phi(h\,\mathrm{rot}(\theta))=\tilde\phi(h\,\mathrm{rot}(\theta))=\mathrm{rot}(\theta)^u\phi(h)
=\prod_{\tau,a}e^{2\pi i\theta o_au_{\tau,a}}\phi(h)$.
\end{proof}

\subsubsection{Peaks lie away from the cusp}\label{psq:ss-P}

This is Siegel's argument; we write it out for unitary groups, for vector-valued weights and with
peaks in place of maxima. The analytic core is Lemma~\ref{psq:analytic}.

\begin{lemma}[Transvections]\label{psq:transv}
Let $v\in\C^n$ with $v^*\varepsilon v=0$ and $\alpha,\beta\in\C$. Put $T_\alpha=1-\alpha vv^*\varepsilon$.
Then $T_\alpha T_\beta=T_{\alpha+\beta}$, and $T_\alpha^*\varepsilon T_\alpha=\varepsilon-2\Rel(\alpha)\,
\varepsilon vv^*\varepsilon$. In particular $T_\alpha\in\auU(\varepsilon)$ if $\Rel\alpha=0$.
\end{lemma}
\begin{proof}
$T_\alpha T_\beta=1-(\alpha+\beta)vv^*\varepsilon+\alpha\beta v(v^*\varepsilon v)v^*\varepsilon$, and the
last term is $0$. Likewise $T_\alpha^*\varepsilon T_\alpha=(1-\aubar\alpha\varepsilon vv^*)\varepsilon(1-
\alpha vv^*\varepsilon)=\varepsilon-(\alpha+\aubar\alpha)\varepsilon vv^*\varepsilon+|\alpha|^2\varepsilon
v(v^*\varepsilon v)v^*\varepsilon$, using $\varepsilon^2=1$.
\end{proof}

\begin{lemma}[Completing squares]\label{psq:factor}
Let $0\ne v\in\C^n$ with $v^*\varepsilon v=0$, $t\in\C$, $\tau=2\Rel t$, and put $S_k=\sum_{l\le k}s_l|v_l|^2$
and $p_k=1-\tau S_k$ for $0\le k\le n$ \textup($S_0=0$, $p_0=p_n=1$\textup). If $p_k>0$ for all $k$,
then $T_t=1-t\,vv^*\varepsilon=hb$ with $h\in\auU(\varepsilon)$, $b\in B$, and $b_{kk}=(p_k/p_{k-1})^{1/2}>0$.
\end{lemma}
\begin{proof}
By Lemma~\ref{psq:transv}, $G:=T_t^*\varepsilon T_t=\varepsilon-\tau ww^*$ with $w=\varepsilon v$, so
$w_l=s_lv_l$. Let $G_k$, $\varepsilon_k$, $w_{(k)}$ be the leading $k\times k$ blocks and the first $k$
coordinates. By the matrix determinant lemma, $\det G_k=\det\varepsilon_k\,(1-\tau w_{(k)}^*\varepsilon_k
w_{(k)})=\bigl(\prod_{l\le k}s_l\bigr)\bigl(1-\tau\sum_{l\le k}s_l|v_l|^2\bigr)=\bigl(\prod_{l\le k}s_l
\bigr)p_k$, which is non-zero. A hermitian matrix all of whose leading principal minors $\delta_k$ are
non-zero is uniquely $c^*Dc$ with $c$ upper unitriangular and $D$ real diagonal, $D_{kk}=\delta_k/
\delta_{k-1}$ ($\delta_0=1$): by induction on $n$, writing $G=\bigl(\begin{smallmatrix}G'&g\\g^*&
\gamma\end{smallmatrix}\bigr)$ with $G'=c'^*D'c'$, take the last column of $c$ to be
$(c'G'^{-1}g,1)$ (so that $c'^*D'(c'G'^{-1}g)=c'^*D'c'G'^{-1}g=g$) and $D_{nn}=\gamma-g^*G'^{-1}g=\det G/\det G'$ (Schur complement). Here
$D_{kk}=s_kp_k/p_{k-1}$. Put $d=\mathrm{diag}((p_k/p_{k-1})^{1/2})$ and $b=dc\in B$; as $d$ commutes
with $\varepsilon$, $b^*\varepsilon b=c^*d\varepsilon dc=c^*Dc=G$. Hence $h=T_tb^{-1}$ satisfies
$h^*\varepsilon h=b^{-*}Gb^{-1}=\varepsilon$, and $b_{kk}=d_{kk}$.
\end{proof}

\begin{lemma}[The weight factor]\label{psq:bound}
Let $u\in\Z^n$ have gap $\ge1$ \textup($u_a+1\le u_c$ for $a\le n_+<c$\textup) and $|u_a|\le A$,
$A\ge1$. Let $0\ne v\in\C^n$ with $v^*\varepsilon v=0$, $a_0=i\alpha$ with $\alpha>0$, and suppose
$\eta=\alpha|v|^2\le\frac12$. For $z$ in the upper half plane $\psqH$ put $t(z)=a_0(z-i)$. Then
$T_{t(z)}=h(z)b(z)$ with $h(z)\in\auU(\varepsilon)$, $b(z)\in B$ with positive diagonal, and
\[
0<b(z)^u\le e^{2A\eta}\,(1+\eta\,\Iml z)^{-1/2}.
\]
\end{lemma}
\begin{proof}
Write $y=\Iml z>0$. Then $\Rel t(z)=\Rel(i\alpha(z-i))=\alpha(1-y)$, so $\tau=2\alpha(1-y)$. Since $v$ is
isotropic, $\sum_{a\le n_+}|v_a|^2=\sum_{a>n_+}|v_a|^2=|v|^2/2$; in particular both signs occur. The
partial sums $S_k$ increase from $0$ to $|v|^2/2$ for $k\le n_+$ and decrease back to $0$ for
$k\ge n_+$, so $S_k\in[0,|v|^2/2]$ and $p_k=1-\tau S_k$ lies between $1$ and
\[
Y:=p_{n_+}=1-\alpha(1-y)|v|^2=1+\eta(y-1)\ \ge 1-\eta\ge\tfrac12 .
\]
So all $p_k>0$, and Lemma~\ref{psq:factor} gives $h(z),b(z)$ with $b(z)^u=\exp(\frac12\sum_kL_ku_k)$,
$L_k=\log p_k-\log p_{k-1}$. We have $\sum_kL_k=0$ and $\sum_{k\le n_+}L_k=\log Y$. Moreover
$p_k-p_{k-1}=-\tau s_k|v_k|^2$, so all $L_k$ with $k\le n_+$ have the sign of $y-1$ (or vanish) and all
$L_k$ with $k>n_+$ the opposite sign.

\emph{Case $y\ge1$.} Let $A_0=\min_{c>n_+}u_c$; by the gap, $u_a-A_0\le-1$ for $a\le n_+$, and
$u_c-A_0\ge0$ for $c>n_+$. As $L_k\ge0$ for $k\le n_+$ and $L_k\le0$ for $k>n_+$,
\[
\sum_kL_ku_k=\sum_kL_k(u_k-A_0)\le-\sum_{k\le n_+}L_k=-\log Y ,
\]
so $b^u\le Y^{-1/2}$. Now $(1+2\eta)Y-(1+\eta y)=\eta(1-2\eta)+2\eta^2y\ge0$ as $\eta\le\frac12$. Hence $Y^{-1/2}\le
(1+2\eta)^{1/2}(1+\eta y)^{-1/2}\le e^{\eta}(1+\eta y)^{-1/2}\le e^{2A\eta}(1+\eta y)^{-1/2}$.

\emph{Case $y<1$.} Then $\sum_kL_ku_k\le A\sum_k|L_k|=2A|\log Y|=-2A\log Y$ (the $L_k$ with $k\le n_+$
sum to $\log Y$ and have one sign, those with $k>n_+$ sum to $-\log Y$), so $b^u\le Y^{-A}\le
(1-\eta)^{-A}$, as $Y\ge1-\eta$ here. By convexity $-\log(1-\eta)\le2\log2\cdot\eta\le1.39\,\eta$ for
$0\le\eta\le\frac12$, and $(1+\eta y)^{1/2}\le(1+\eta)^{1/2}\le e^{\eta/2}$; so
$b^u\le e^{(1.39A+0.5)\eta}(1+\eta y)^{-1/2}\le e^{2A\eta}(1+\eta y)^{-1/2}$ since $A\ge1$.
\end{proof}

\begin{lemma}[Schwarz at the cusp]\label{psq:schwarz}
Let $F$ be holomorphic on $\psqH$ with $F(z+1)=F(z)$, $|F|\le C'$, and $F(z)\to0$ as $\Iml z\to\infty$
uniformly in $\Rel z$. Then $|F(i)|\le C'e^{-2\pi}$.
\end{lemma}
\begin{proof}
$F(z)=g(e^{2\pi iz})$ with $g$ holomorphic on the punctured unit disc and bounded by $C'$; by Riemann's
theorem $g$ extends holomorphically to the disc, with $g(0)=0$ by the hypothesis at $i\infty$. By the
Schwarz lemma applied to $g/C'$, $|g(q)|\le C'|q|$, and $z=i$ corresponds to $|q|=e^{-2\pi}$.
\end{proof}

\begin{lemma}[Analytic part]\label{psq:analytic}
Let $u=(u_\tau)_\tau$ be a weight with gap $\ge1$ at every $\tau$ and entries in $[-A,A]$, $A\ge2$. Let
$\phi$ be holomorphic of weight $u$ on $\pspG$ with $|\phi|\le C$. Let $t\in\pspG$ with
$t_\tau=1-a_\tau v_\tau v_\tau^*\varepsilon$, where $v_\tau\in\C^n$ is isotropic for $\varepsilon$,
$a_\tau\in i\R$, $\eta_\tau=\Iml(a_\tau)|v_\tau|^2>0$, and $A\sum_\tau\eta_\tau\le1$. If $\phi(th)=\phi(h)$
for all $h\in\pspG$, then $|\phi(1)|\le Ce^{2-2\pi}$.
\end{lemma}
\begin{proof}
For $\zeta\in\C$ let $x(\zeta)\in\prod_\tau\auGL_n(\C)$ have components $x(\zeta)_\tau=1-a_\tau(\zeta-i)
v_\tau v_\tau^*\varepsilon$ (invertible, with inverse given by Lemma~\ref{psq:transv}); it depends
holomorphically (affinely) on $\zeta$, $x(i)=1$, and $t\,x(\zeta)=x(\zeta+1)$ by Lemma~\ref{psq:transv}.
Each $\eta_\tau\le1/A\le\frac12$. For $z\in\psqH$, Lemma~\ref{psq:bound} at each $\tau$ (with
$a_0=a_\tau$, $v=v_\tau$) gives $x(z)=h(z)b(z)$ with $h(z)\in\pspG$, $b(z)\in B$ and
$0<b(z)_\tau^{u_\tau}\le e^{2A\eta_\tau}(1+\eta_\tau\Iml z)^{-1/2}$. In particular $x(z)\in\pspG B$, and
$f(z)=\tilde\phi(x(z))$ is holomorphic on $\psqH$, with
\[
f(z)=\Bigl(\prod_\tau b(z)_\tau^{u_\tau}\Bigr)\phi(h(z)),\qquad |f(z)|\le C\prod_\tau e^{2A\eta_\tau}
(1+\eta_\tau\Iml z)^{-1/2}\le Ce^{2}\prod_\tau(1+\eta_\tau\Iml z)^{-1/2}.
\]
So $|f|\le Ce^2$ and $f(z)\to0$ as $\Iml z\to\infty$, uniformly in $\Rel z$. Periodicity:
$x(z+1)=t\,h(z)\,b(z)$ with $t\,h(z)\in\pspG$, so $f(z+1)=b(z)^u\phi(t\,h(z))=b(z)^u\phi(h(z))=f(z)$.
Finally $f(i)=\tilde\phi(1)=\phi(1)$. Lemma~\ref{psq:schwarz} with $C'=Ce^2$ gives the claim.
\end{proof}

\begin{theorem}[Input (P)]\label{psq:peaks}
Let $\Phi$ be a bounded function on $\auU(X)(\auA)$, left invariant under $\auU(X)(E^+)$, right
invariant under $K(p)$ for some $p\in\auO_E$ with $\aubar p=-p$ and $0<\Iml\tau(p)\le A'$ for all
$\tau\in\auPl$, $A'\ge1$, and holomorphic of a weight $u$ with gap $\ge1$ and entries in $[-A,A]$,
$A\ge2$. Then every peak $g$ of $\Phi$ satisfies $\pspl(X,g)\ge(AA')^{-1/2}$. This is the input (P)
of Proposition~\ref{psp:inputs}.
\end{theorem}
\begin{proof}
Let $C=\sup|\Phi|$, positive since $\Phi(g)\ne0$. As $(AA')^{-1/2}\le1$ it suffices to show that no
$0\ne\gamma\in\Lambda_g$ has $|\gamma|_g<(AA')^{-1/2}$. Suppose $\gamma$ is one, and let
$Y=g^{-1}\gamma\in\auA_E^n$ be its coordinates at $g$; so $Y_f$ is integral and $\sum_\tau|Y_\tau|^2=
|\gamma|_g^2<(AA')^{-1}\le1$.

\emph{Isotropy.} By fact (R3) stated after Definition~\ref{pcx:def-ht} with $\gamma'=\gamma$, $\langle\gamma,\gamma\rangle=0$;
hence $Y^*SY=0$ at every place ($g$ is unitary), in particular $Y_\tau^*\varepsilon Y_\tau=0$.

\emph{The transvection.} Let $T_\gamma(w)=w-p\langle w,\gamma\rangle\gamma$ on $X$. Since
$\aubar p=-p$ and $\langle\gamma,\gamma\rangle=0$, the computation of Lemma~\ref{psq:transv} (in any
basis) shows $T_\gamma\in\auU(X)(E^+)$. Put $t=g^{-1}T_\gamma g\in\auU(X)(\auA)$. As $g$ is unitary,
$t(w)=w-p\langle gw,\gamma\rangle g^{-1}\gamma=w-p\langle w,Y\rangle Y$, so the matrix of $t$ in the
frame is $1-p\,YY^*S$. At the finite places this is $1-p\,Y_fY_f^*S_f$ with $Y_f$ and $S_f$ integral,
so $t_f\in K(p)$; at $\tau$ it is $t_\tau=1-\tau(p)Y_\tau Y_\tau^*\varepsilon$, and $\tau(p)\in i\R$
because $\aubar p=-p$.

\emph{Invariance.} Let $\phi(h)=\Phi(gh)$ on $\pspG$; it is holomorphic of weight $u$ and $|\phi|\le C$.
For $h\in\pspG$, as $t_f$ commutes with $\pspG$, $g\,t_\infty h=g\,t\,h\,t_f^{-1}$, so
\[
\phi(t_\infty h)=\Phi(g\,t\,h\,t_f^{-1})=\Phi(g\,t\,h)=\Phi(T_\gamma\,g\,h)=\Phi(gh)=\phi(h),
\]
using $t_f^{-1}\in K(p)$ and $T_\gamma\in\auU(X)(E^+)$.

\emph{Sizes.} $\eta_\tau=\Iml\tau(p)|Y_\tau|^2>0$, since $\gamma\ne0$ forces $Y_\tau=g_\tau^{-1}\tau(\gamma)\ne0$;
and $A\sum_\tau\eta_\tau\le AA'\sum_\tau|Y_\tau|^2<1$.

Lemma~\ref{psq:analytic} for $\phi$, $t_\infty$, $v_\tau=Y_\tau$, $a_\tau=\tau(p)$ gives $|\Phi(g)|=
|\phi(1)|\le Ce^{2-2\pi}<C/2$, as $2\pi-2>\log2$. This contradicts $|\Phi(g)|>C/2$.
\end{proof}

\begin{remark}[The applications in \S\ref{psp:sec}]\label{psq:Papp}
In Proposition~\ref{psp:approx} (P) is applied with $p=\sqrt{-d}$ to the carrier, and in
\S\ref{psp:ss-proof} with $p=M\sqrt{-d}$ to $\aubar y$. In both cases $\aubar p=-p$, and
$\Iml\tau(p)$ has the same sign for all $\tau\in\auPl$: $\auPl$ is identified with
$\operatorname{Gal}(E/K)$, so all $\tau$ restrict to one embedding of $K=\Q(\sqrt{-d})$ and
$\tau(\sqrt{-d})=\pm i\sqrt d$ with a sign independent of $\tau$. If the sign is $-$, replace $p$ by
$-p$, which generates the same ideal and so gives the same group $K(p)$. Thus $0<\Iml\tau(p)\le A'$
holds with $A'=\sqrt d$, resp.\ $A'=M\sqrt d$, as used there. The holomorphy hypotheses are the
inputs (H4), proved in \S\S\ref{pss:sec}, \ref{psv:sec}.
\end{remark}

\subsubsection{Three checks}\label{psq:ss-checks}

\begin{lemma}[Normalisation of $\theta$]\label{psq:thetanorm}
Let $\mathcal W$ be a symplectic space over $E^+$ of the form $\auHom(X,X')$, with the inner product on
$\psHS(\mathcal W)=\mathfrak F\otimes\psHS_f(\mathcal W)$ for which the vacuum of \S\ref{pk:ss-conv}
is a unit vector \textup(the inner product used in Lemma~\ref{psp:avg}\textup), and with $\theta$
normalised as in the proof of Proposition~\ref{pk:conj}(ii), where $\int|\theta(\psrho(w)F)|^2dw=\|F\|^2$ for the self-dual measure, which gives $\mathcal W(E^+)\backslash\mathcal W(\auA)$ volume $1$; in particular $w\mapsto|\theta(\psrho(w)(1\otimes\psvac))|^2$ has mean $1$ over
$\mathcal W(E^+)\backslash\mathcal W(\auA)$ \textup(total mass $1$\textup). Then for every unitary
isomorphism $U$ of $\psHS(\mathcal W)$ onto a Schr\"odinger model $\mathcal S(\mathcal X(\auA))$
commuting with the Heisenberg operators there is $\zeta$ with $|\zeta|=1$ and $\theta=\zeta\,\theta_S\circ U$,
$\theta_S(F)=\sum_{\xi\in\mathcal X(E^+)}F(\xi)$. Consequently Lemma~\ref{psp:prodtheta} holds for the
paper's $\theta$ exactly as stated, without a constant.
\end{lemma}
\begin{proof}
$\theta_S$ is invariant under $\psrho(w)$, $w\in\mathcal W(E^+)$ (Poisson summation, as in the proof of
Lemma~\ref{psp:prodtheta}), so $\theta\circ U^{-1}$ and $\theta_S$ are both invariant functionals on the
model, and by the uniqueness in (W1) of Lemma~\ref{ps:weil}, $\theta=\zeta\theta_S\circ U$ for some
$\zeta\in\C$. For $F\in\mathcal S(\mathcal X(\auA))$ and $w=(a,b)$ with $a\in\mathcal X(\auA)$,
$b\in\mathcal Y(\auA)$, with the notation of the proof of Lemma~\ref{psp:prodtheta},
$\theta_S(\psrho(a,b)F)=\psi(\frac12\langle a,b\rangle)\sum_\xi\psi(\langle\xi,b\rangle)F(\xi+a)$. The
modulus is invariant under $\mathcal W(E^+)$ (since $\psrho(\gamma)\psrho(w)$ is a unimodular multiple
of $\psrho(\gamma+w)$ and $\theta_S\circ\psrho(\gamma)=\theta_S$). The characters $b\mapsto\psi(\langle\xi,b
\rangle)$, $\xi\in\mathcal X(E^+)$, of the compact group $\mathcal Y(E^+)\backslash\mathcal Y(\auA)$ (of mass
$1$) are distinct, so by Parseval and the absolute summability of $\xi\mapsto F(\xi+a)$,
\[
\int_{\mathcal Y(E^+)\backslash\mathcal Y(\auA)}|\theta_S(\psrho(a,b)F)|^2db=\sum_\xi|F(\xi+a)|^2,\qquad
\int_{\mathcal X(E^+)\backslash\mathcal X(\auA)}\sum_\xi|F(\xi+a)|^2da=\int_{\mathcal X(\auA)}|F|^2=\|F\|^2 .
\]
Hence the mean of $|\theta(\psrho(w)F)|^2$ over $\mathcal W(E^+)\backslash\mathcal W(\auA)$ (which is
$\mathcal X(E^+)\backslash\mathcal X(\auA)\times\mathcal Y(E^+)\backslash\mathcal Y(\auA)$ as a measure space) is
$|\zeta|^2\|F\|^2$ for every $F$. For $F=1\otimes\psvac$, a unit vector, the normalisation gives
$|\zeta|=1$; conversely the mean is then $\|F\|^2$ for every $F$, so the normalisation by the unit vacuum agrees with the one of Proposition~\ref{pk:conj}(ii). In Lemma~\ref{psp:prodtheta} the proof computes $\theta_S(UF)\aubar{\theta_S(UF')}$; the
factor $|\zeta|^2=1$ makes this equal to $\theta(F)\aubar{\theta(F')}$, and $U$ preserves the inner
products $\langle F',\psrho(y)F\rangle$.
\end{proof}

\noindent So no constant enters
Proposition~\ref{psp:sparse}; the only requirement is that the vacuum be a unit vector for the inner
product used, which is how Lemma~\ref{psp:avg} uses it.

\begin{remark}[Level of $e_0\poA(W_j)$]\label{psq:level}
By the definitions preceding Definition~\ref{def:au-target} and the sentence after it, $\poA(W_j)$
is defined on $W_j$ ``with the parameter $\varsigma$ and the level $J_M$'', and $J_M$ is the principal
congruence subgroup of level $\pi_M$ in the stabiliser of the lattice; $e_0\poA(W_j)\subseteq
\poA(W_j)$, since $\varrho(e_0)$ is an operator on $\poA(W_j)$. So every $y\in e_0\poA(W_j)$ is right
invariant under $J_M=K(\pi_M)$, where $K(\cdot)$ is taken in the frame whose integral vectors form that
lattice, as in Definition~\ref{pcx:def-ht}. If $\pi_M\mid p$ in $\auO_E$, then $K(p)\subseteq
K(\pi_M)$ (a matrix $\equiv1\bmod p$ is $\equiv1\bmod\pi_M$). Since $\pi_M\mid\pi_M\aubar{\pi_M}=M$, this
applies to $p=M\sqrt{-d}$ and to $p=\pi_M^2$. This is the invariance used in \S\ref{psp:ss-proof} (level $M\sqrt{-d}$) and in Lemma~\ref{psp:step2}
(level $\pi_M^2$, as $\pi_M\mid\pi_M^2$).
\end{remark}

\begin{lemma}[Degree of $\widetilde P$]\label{psq:degree}
In the notation of \S\ref{pte:ss-third} and Lemma~\ref{psp:weights}, $\widetilde P$ is homogeneous for
real scalars of degree $d_{\widetilde P}=\sum_\tau d_\tau$,
\[
d_\tau=2m\bigl(\mu^+_{1,\tau}-\Delta'\bigr)+m\bigl(\mu^+_{1,\tau}+N''+m\bigr),
\]
where $\mu^+_{1,\tau}=\#\iota(\tau)$ is the number of columns of $\mu^+_\tau$; and $d_{\widetilde P}\le
7fmb\le b^{3/2}$.
\end{lemma}
\begin{proof}
By \eqref{pte:eq-SDMP} with $\mathsf W$ of rank $n=3m+N''$, $l=m$ negative indices, and $V'=S$ ($p=2m$,
$l=m$), $P_i=\prod_\tau\det(y^{--}_\tau)^{\#\iota(\tau)+n-2l}\prod_{k\in\iota(\tau)}\det(\eta_{i,\tau k}
y^{++}_\tau)$. The negative indices of $\mathsf W$ lie in $\widetilde W$ ($W^\flat$ is positive definite),
so the first factor lives on $\auHom(S,\widetilde W)$; it is homogeneous of degree $m(\#\iota(\tau)+
n-2m)=m(\mu^+_{1,\tau}+N''+m)$, as $n-2m=N''+m$. Each $\det(\eta_{i,\tau k}y^{++}_\tau)$ is a $2m\times2m$
determinant whose entries are linear in $y^{++}_\tau$, so it is homogeneous of degree $2m$; by
\S\ref{pte:ss-third} the rows of $\eta_{i,\tau k}$ are basis vectors of $\widetilde W$ for the
$\mu^+_{1,\tau}-\Delta'$ columns of height $<r_+$, independent of $i$, and short rows of $W^\flat$ for the
$\Delta'$ columns of height $r_+$. Hence $P_i=\widetilde P\cdot Q_{r(i)}$ with $\widetilde P$ the product
of the first factor and of the $\mu^+_{1,\tau}-\Delta'$ factors with rows in $\widetilde W$, and its
degree at $\tau$ is $d_\tau$. Each factor is homogeneous under $y\mapsto sy$, $s$ real, at all places at
once, so the total degree is $\sum_\tau d_\tau$.

Bound: by the proof of Proposition~\ref{pte:third}, $\mu^+_{1,\tau}=\#\iota(\tau)\le\kappa_j(\tau)=
\kappa(\tau)+(j+\frac12)\Delta\le b+(m+1)\Delta\le2b$ (as $\kappa(\tau)\le b$ and $(m+1)\Delta\le
3\sqrt x\cdot2(3x^5)^8\le b$ by \eqref{eq:au-scales}); and $N''+m=r+10m+1\le b$. So $d_\tau\le4mb+3mb=7mb$
and $d_{\widetilde P}\le7fmb$. Finally $f\le x^{12}$ (Lemma~\ref{psp:sizes}(iii)), $m\le2\sqrt x$, so
$7fm\le14x^{12.5}\le x^{13}\le b^{1/2}$, as $x^{26}\le b$ by \eqref{eq:au-scales}; hence
$d_{\widetilde P}\le b^{3/2}$.
\end{proof}

\noindent This is the formula used in Lemma~\ref{psp:weights}.

\subsubsection{The vacuum is a unit vector}\label{psq:ss-vac}

\begin{lemma}[Inner product and vacuum]\label{psq:unitvac}
Let $\mathcal W=\auHom(X,X')$ and, for smooth $F,F'\in\psHS(\mathcal W)$,
\[
\langle F,F'\rangle_\theta=\int_{\mathcal W(E^+)\backslash\mathcal W(\auA)}\aubar{\theta(\psrho(w)F)}\,
\theta(\psrho(w)F')\,dw\qquad(\text{total mass }1).
\]
\textup{(a)} This is well defined, positive definite and invariant under $\psrho$ and $\omega$.
\textup{(b)} $\|1\otimes\psvac\|_\theta=1$. \textup{(c)} Fix at each place $v$ an invariant inner
product $\langle\,,\rangle_v$ on $\psHS_v$ for which the local vacuum $\psvac_v$ is a unit vector, and
write $\psvac=c\bigotimes_v\psvac_v$, $c\in\C^\times$. Then for vectors $F=c\bigotimes_vF_v$,
$F'=c\bigotimes_vF'_v$ with $F_v=F'_v=\psvac_v$ for almost all $v$,
$\langle F,F'\rangle_\theta=\prod_v\langle F_v,F'_v\rangle_v$.
\end{lemma}
\begin{proof}
(a) The integrand is $\mathcal W(E^+)$-invariant, as in the proof of Lemma~\ref{psq:thetanorm}, and by
that proof, with $\theta=\zeta\theta_S\circ U$ for an isometry $U$ onto a Schr\"odinger model (for any
fixed invariant inner product $\langle\,,\rangle_0$), polarisation gives $\langle F,F'\rangle_\theta=
|\zeta|^2\langle F,F'\rangle_0$, with $\zeta\ne0$ as $\theta\ne0$. Invariance under $\psrho$ is the
invariance of $dw$ under translation, and under $\omega$ follows from $\langle\,,\rangle_0$ being
$\omega$-invariant ($\omega$ is unitary, Lemma~\ref{ps:weil}). (b) is the normalisation of $\theta$ in Lemma~\ref{psq:thetanorm} (equivalently, of the proof of
Proposition~\ref{pk:conj}(ii)). (c) At each place the Stone--von Neumann theorem gives a unitary
isomorphism $U_v$ of $(\psHS_v,\langle\,,\rangle_v)$ onto a local Schr\"odinger model commuting with
the Heisenberg operators; $U=\bigotimes_vU_v$ (restricted tensor product with respect to $\psvac_v$ and
$U_v\psvac_v$) is then an isometry of $(\psHS,\bigotimes_v\langle\,,\rangle_v)$ onto a Schr\"odinger
model $\mathcal S(\mathcal X(\auA))$ commuting with the Heisenberg operators. By (a), applied with
$\langle\,,\rangle_0=\bigotimes_v\langle\,,\rangle_v$, we get $\langle\,,\rangle_\theta=\kappa\bigotimes_v
\langle\,,\rangle_v$ with $\kappa=|\zeta|^2>0$. Evaluating at $1\otimes\psvac$ and using (b), $1=\kappa|c|^2$, which gives (c).
\end{proof}

\noindent Thus, with $\langle\,,\rangle_\theta$ as the inner product, the vacuum is a unit vector (by
the paper's normalisation, not by an extra convention), Lemma~\ref{psq:thetanorm} applies with
$|\zeta|=1$, and the pairings of Lemma~\ref{psp:avg} factor over the places with unit local vacua,
provided the data $[v,\Psi_T]$ are read as $\psvac$ with its components at the places above $\pi_M$
(and at infinity) replaced, which is how kernel data are built from $\psvac$.

\subsubsection{Holomorphic discrete series: (H1)--(H3)}\label{psq:ss-H3}

Fix one infinite place and $G=\auU(\varepsilon)$, $\varepsilon=\mathrm{diag}(s_a)$, $N=n$, indices
numbered by $o$ (positive first), $p=n_+$, $q=n-p$, $T$ the diagonal unitary torus, $B$ as above. The
statements for $\pspG=\prod_\tau G_\tau$ follow by taking products. Let $u\in\Z^N$ be dominant with
gap $\gamma\ge N^2+N+1$.

\begin{lemma}[The parameter]\label{psq:param}
Let $\rho_{\mathrm c},\rho_{\mathrm n}$ be the half sums of \S\ref{pds:ss-setup} for the order
in which the indices are decreasing as follows: the negative indices in the order $o$, then the
positive ones in the order $o$. Then $\lambda=u+\rho_{\mathrm c}-\rho_{\mathrm n}$ is a regular
parameter whose entries decrease in exactly this order, $\Lambda^K(\lambda)=u$, the gap of $\lambda$ is
$\ge\gamma-N\ge N^2+1$, and $\max|\lambda_a|\le\max|u_a|+N$.
\end{lemma}
\begin{proof}
For $a\le p$: $\rho_{\mathrm n,a}=-q/2$ and $\rho_{\mathrm c,a}=\frac{p+1}2-a$; for $c>p$:
$\rho_{\mathrm n,c}=p/2$ and $\rho_{\mathrm c,c}=\frac{q+1}2-(c-p)$. Within a block $u$ is
non-increasing and $\rho_{\mathrm c}$ strictly decreasing, $\rho_{\mathrm n}$ constant, so $\lambda$ is
strictly decreasing in $o$ there. For $a\le p<c$, $\lambda_c-\lambda_a\ge\gamma-\frac{p-1}2-\frac{q-1}2-
\frac N2\ge\gamma-N>0$. So the order of $\lambda$ is the stated one, the half sums computed for it are
those used, and $\Lambda^K(\lambda)=\lambda+\rho_{\mathrm n}-\rho_{\mathrm c}=u$. The entry bound is
immediate.
\end{proof}

\noindent Let $\pi$ be the discrete series of parameter $\lambda$, $v$ a unit highest weight vector
of its lowest $K$-type (highest for the positive compact roots $e_a-e_b$, $\lambda_a>\lambda_b$,
i.e.\ $a$ before $b$ in one block), and $\mathsf m_u(g)=\langle v,\pi(g)v\rangle$, the coefficient
$c_\tau$ of \S\ref{pds:sec} for $\lambda$ (inner products conjugate-linear in the first variable).
This is the precise meaning of $\mathsf m_u$ in \S\ref{psp:ss-notation}; the bounds
\eqref{psp:eq-ds} hold for it with the gap and entry bound of Lemma~\ref{psq:param}. Let
$\Delta_k$ be the leading principal $k\times k$ minor and
\[
\phi_u(x)=\prod_{k=1}^{N-1}\Delta_k(x)^{u_k-u_{k+1}}\cdot\det(x)^{u_N}.
\]

\begin{lemma}[$G B$ and $\phi_u$]\label{psq:GB}
\textup{(i)} $G\cap B=T$, and $GB$ is open and connected in $\auGL_N(\C)$. \textup{(ii)} $\phi_u$ is
holomorphic on $GB$ and holomorphic of weight $u$; $\phi_u(1)=1$; and $\phi_u(tx)=t^u\phi_u(x)$ for
diagonal $t$.
\end{lemma}
\begin{proof}
(i) If $g\in G\cap B$, $g^{-1}=\varepsilon g^*\varepsilon$ is upper and lower triangular, so $g$ is diagonal
unitary. The map $G\times B\to\auGL_N(\C)$, $(h,b)\mapsto hb$, has differential at $(h,b)$ with image
$h(\mathfrak g+\mathfrak b)b$; $\mathfrak g\cap\mathfrak b=\mathfrak t$, so $\dim_\R(\mathfrak g+\mathfrak b)
=N^2+N(N+1)-N=2N^2$ and the map is a submersion; its image is open, and connected since $G$ and $B$
are. (ii) For $b\in B$ the leading $k\times k$ block of $xb$ is that of $x$ times that of $b$, so
$\Delta_k(xb)=\Delta_k(x)\prod_{l\le k}b_{ll}$; similarly $\Delta_k(tx)=\prod_{l\le k}t_{ll}\Delta_k(x)$.
Hence $\phi_u(xb)=b^u\phi_u(x)$ and $\phi_u(tx)=t^u\phi_u(x)$ (the exponent of $b_{ll}$ is
$\sum_{k\ge l}(u_k-u_{k+1})+u_N=u_l$). The exponents $u_k-u_{k+1}$ are $\ge0$ for $k\ne p$ (dominance),
so $\phi_u$ is holomorphic wherever $\Delta_p\det\ne0$. For $h\in G$, $h_{++}^*h_{++}=1+h_{-+}^*h_{-+}$
(from $h^*\varepsilon h=\varepsilon$), so $\Delta_p(h)=\det h_{++}\ne0$, and $\Delta_p(hb)=\Delta_p(h)\prod_{l\le p}
b_{ll}\ne0$.
\end{proof}

\begin{proposition}[Uniqueness]\label{psq:unique}
Let $\psi$ be holomorphic of weight $u$ on $G$ with $\psi(\mathrm{rot}(\theta)h)=e^{2\pi i\theta\langle
o,u\rangle}\psi(h)$ for all $\theta,h$. Then $\psi=\psi(1)\phi_u$ on $G$.
\end{proposition}
\begin{proof}
Let $\tilde\psi$ be the extension; on the open set $GB$ it satisfies $\tilde\psi(hb)=b^u\psi(h)$, hence
$\tilde\psi(xb)=b^u\tilde\psi(x)$ and $\tilde\psi(\mathrm{rot}(\theta)x)=e^{2\pi i\theta\langle o,u\rangle}
\tilde\psi(x)$ for $x\in GB$. Near $1$ all $\Delta_k\ne0$; let $f=\tilde\psi/\phi_u$ there, so
$f(xb)=f(x)$ and $f(\mathrm{rot}(\theta)x)=f(x)$ (Lemma~\ref{psq:GB}(ii)). For $X$ strictly lower
triangular and small, $1+X\in GB$; put $F(X)=f(1+X)$, holomorphic near $0$. Since $\mathrm{rot}(\theta)$ is
diagonal, $F(\mathrm{rot}\,X\,\mathrm{rot}^{-1})=f(\mathrm{rot}(1+X)\mathrm{rot}^{-1})=f(1+X)=F(X)$, and
conjugation multiplies $X_{ab}$ ($a>b$) by $e^{2\pi i\theta(o_a-o_b)}$ with $o_a-o_b>0$. Comparing
Taylor coefficients, every non-constant monomial has a non-zero total weight, so its coefficient
vanishes: $F$ is constant near $0$. Since $(X,b)\mapsto(1+X)b$ covers a neighbourhood of $1$ (big
cell), $\tilde\psi=f(1)\phi_u$ near $1$, hence on the connected open set $GB$ by the identity theorem;
$f(1)=\psi(1)$.
\end{proof}

\begin{proposition}[(H1)]\label{psq:mhol}
$\mathsf m_u$ is holomorphic of weight $u$, $\mathsf m_u(\mathrm{rot}(\theta)h)=e^{2\pi i\theta\langle
o,u\rangle}\mathsf m_u(h)$, and $\mathsf m_u=\phi_u$ on $G$.
\end{proposition}
\begin{proof}
$v$ is $K$-finite, hence smooth, and for real $X\in\mathfrak g$ the right derivative is $R_X\mathsf m_u(h)
=\langle v,\pi(h)d\pi(X)v\rangle$; extend $d\pi$ and $R$ complex-linearly to $\mathfrak g_\C=\mathfrak
{gl}_N(\C)$. \emph{Claim: $d\pi(Y)v=u(Y)v$ for $Y\in\mathfrak b$, $u(Y)=\sum_au_aY_{aa}$.} For diagonal
$Y$ this says $v$ has $T$-weight $\Lambda^K(\lambda)=u$. For $Y=E_{ab}$, $a$ before $b$ in one block,
$e_a-e_b$ is a positive compact root and $d\pi(E_{ab})v=0$ as $v$ is a highest weight vector. For
$Y=E_{ac}$, $a\le p<c$: $d\pi(E_{ac})v$ has $T$-weight $u+e_a-e_c$; if it were non-zero, some $K$-type
$\mu$ of $\pi$ would contain this weight, so $u+e_a-e_c=\mu-\sum m_\alpha\alpha$ with $m_\alpha\ge0$ over
positive compact roots, and by (D4) of \S\ref{ptq:ss-H} $\mu=u+\sum n_\beta\beta$ over positive
noncompact roots, $n_\beta\ge0$. Apply the linear form $H(x)=\sum_{c'>p}x_{c'}$: $H$ vanishes on compact
roots, is $+1$ on positive noncompact roots (these are $e_{c'}-e_{a'}$, $a'\le p<c'$, since
$\lambda_{c'}>\lambda_{a'}$ by Lemma~\ref{psq:param}), and is $-1$ on $e_a-e_c$; so $-1=\sum n_\beta\ge0$,
a contradiction.

\emph{Extension.} As $\pi(t)v=t^uv$ for $t\in T$, $\mathsf m_u(ht)=t^u\mathsf m_u(h)$; by
Lemma~\ref{psq:GB}(i), $\tilde{\mathsf m}(hb)=b^u\mathsf m_u(h)$ is well defined on $GB$, and smooth because
$(h,b)\mapsto b^u\mathsf m_u(h)$ is smooth and constant on the fibres of the surjective submersion
$(h,b)\mapsto hb$. At $x=hb$ and $Z\in\mathfrak{gl}_N(\C)$, $\tilde{\mathsf m}(xe^{sZ})=b^u\tilde{\mathsf
m}(he^{sZ'})$ with $Z'=bZb^{-1}$. Write $Z'=X+Y$, $X\in\mathfrak g$, $Y\in\mathfrak b$; to first order
$he^{sZ'}=he^{sX}e^{sY}$, and $\tilde{\mathsf m}(he^{sX}e^{sY})=e^{su(Y)}\mathsf m_u(he^{sX})$, so the
derivative is $R_X\mathsf m_u(h)+u(Y)\mathsf m_u(h)=\langle v,\pi(h)d\pi(X+Y)v\rangle$ by the claim. This
is complex-linear in $Z'$, hence in $Z$: $\tilde{\mathsf m}$ has complex-linear differential, so it is
holomorphic on $GB$. Left equivariance: $\mathsf m_u(\mathrm{rot}\,h)=\langle\pi(\mathrm{rot})^{-1}v,
\pi(h)v\rangle=\aubar{\mathrm{rot}^{-u}}\mathsf m_u(h)=e^{2\pi i\theta\langle o,u\rangle}\mathsf m_u(h)$.
Proposition~\ref{psq:unique} and $\mathsf m_u(1)=1$ give $\mathsf m_u=\phi_u$.
\end{proof}

\begin{theorem}[(H1)--(H3)]\label{psq:H123}
For dominant weights of gap $\ge N^2+N+1$ at every place: \textup{(H1)} and \textup{(H2)} of
Proposition~\ref{psp:inputs} hold, and \textup{(H3)} holds: if $\phi$ is holomorphic of weight $u$ and
$\phi\,\aubar{\mathsf m_u}$ is integrable on $\pspG$, then $\int\phi\aubar{\mathsf m_u}=\phi(1)\int|\mathsf
m_u|^2$.
\end{theorem}
\begin{proof}
(H1): Lemma~\ref{psq:H12}(i) and Proposition~\ref{psq:mhol}. (H2): the first part is
Lemma~\ref{psq:H12}(ii). Let $\psi(h)=\int_{\R/\Z}\phi(\mathrm{rot}(\theta)h)e^{-2\pi i\theta\langle o,u
\rangle}d\theta$. Then $\tilde\psi(x)=\int\tilde\phi(\mathrm{rot}(\theta)x)e^{-2\pi i\theta\langle o,u\rangle}
d\theta$ is holomorphic on $\pspG B$ (the integrand is holomorphic in $x$, continuous in $(\theta,x)$, and
$\mathrm{rot}(\theta)\pspG B=\pspG B$), $\tilde\psi(hb)=b^u\psi(h)$, and $\psi$ is left equivariant by the
substitution $\theta\mapsto\theta-\theta'$. By Proposition~\ref{psq:unique} for $\pspG=\prod_\tau G_\tau$ (the same proof applies:
conjugation by the diagonal $\mathrm{rot}(\theta)$ multiplies each coordinate $X_{\tau,ab}$, $a>b$, by
$e^{2\pi i\theta(o_a-o_b)}$, which has positive weight, $\pspG B$ is connected and open, and
$\phi_u=\prod_\tau\phi_{u_\tau}$; it does not apply place by place, since $\psi$ is equivariant only under
the diagonal $\mathrm{rot}(\theta)$) $\psi=\psi(1)\phi_u=\psi(1)\mathsf m_u$, and $\psi(1)=\int\phi(1\cdot\mathrm{rot}(\theta))
e^{-2\pi i\theta\langle o,u\rangle}d\theta=\phi(1)$ by Lemma~\ref{psq:H12}(ii). (H3): Haar measure on $\pspG$
is invariant under left translation by $\mathrm{rot}(\theta)$, and
$(\theta,h)\mapsto\phi(\mathrm{rot}(\theta)h)\aubar{\mathsf m_u(\mathrm{rot}(\theta)h)}$ is integrable on
$\R/\Z\times\pspG$ (its $h$-integral of the modulus is $\int|\phi\mathsf m_u|$ for each $\theta$). Using
$\aubar{\mathsf m_u(\mathrm{rot}(\theta)h)}=e^{-2\pi i\theta\langle o,u\rangle}\aubar{\mathsf m_u(h)}$ and
Fubini, $\int\phi\aubar{\mathsf m_u}=\int\psi\,\aubar{\mathsf m_u}=\phi(1)\int|\mathsf m_u|^2$.
\end{proof}



%% ===== proofs/C-span3.tex =====
% proofs/C-span3.tex
% Input (R) of Proposition psp:inputs (C-span.tex): the stable-range reproducing formula.
% To be \input after C-span.tex and C-span2.tex. Labels of this file start with psr:.
% Only amsmath/amssymb/amsthm.
\providecommand{\auO}{\mathcal O}
\providecommand{\auU}{\mathrm U}
\providecommand{\auGL}{\mathrm{GL}}
\providecommand{\auA}{\mathbb A}
\providecommand{\aubar}[1]{\overline{#1}}
\providecommand{\pspG}{G_\infty}
\providecommand{\pspm}{\mathsf m}
\providecommand{\pcMat}{\mathrm M}
\providecommand{\psvac}{{\phi^{\mathrm{vac}}}}
\providecommand{\psZ}{Z}
% macros of this file (prefix \psr)
\providecommand{\psrR}{\mathcal R}
\providecommand{\psrmc}{\mathsf{mc}}
\providecommand{\psrm}{\mathfrak m}
\providecommand{\psrF}{\mathbb F}

\subsection{The reproducing formula}\label{psr:sec}

This subsection proves the input (R) of Proposition~\ref{psp:inputs}, in the general form of
Theorem~\ref{psr:R}: for a holomorphic discrete series coefficient $\pspm_u$ on a unitary group
$\auU(X)$ and a positive definite space $Y$ of rank $>2\operatorname{rk}X$, convolution with $\pspm_u$
over the real points is, at each point, a fixed finite combination of theta lifts to $\auU(Y)$, the
same for all rapidly decreasing functions of a given level. The argument is linear duality
(\S\ref{psr:ss-R}) on top of three facts: an inner product formula of Rallis' type for rapidly
decreasing functions in Weil's range (\S\ref{psr:ss-rallis}), whose unramified part is a series of
translates which we show to be injective in the whole range $M'>2n$ (\S\ref{psr:ss-series}); and the
identification of the matrix coefficient of a holomorphic Fock vector with $\pspm_u$
(\S\ref{psr:ss-hol}). The applications of Theorem~\ref{psr:R} in \S\ref{psp:sec} need, in addition,
that its hypotheses hold for the data used there; this is proved in \S\ref{pss:sec} and \S\ref{psv:sec}
(the datum of $K^\circ$ is $\delta_\iota$, Proposition~\ref{psv:datum}).

\emph{Conventions.} Spaces, frames, moduli, the groups $K_S$, $J_S(\psrm)$, the representations
$\omega_{\alpha,\beta}$, data $[P,\Psi]$, kernels $\theta_{\alpha,\beta}(g,h;P,\Psi)$ and lifts
$\Theta_{\alpha,\beta}(P,\Psi)$ are those of \S\ref{pc:ss-data}; Haar measures are normalised as
there ($K_S$ of volume $1$, $\auU(Y)_\infty$ of mass $1$ for definite $Y$). Inner products of Hilbert
spaces and the Petersson product $\langle\varphi',\varphi\rangle=\int_{[\auU(S)]}\aubar{\varphi'}
\varphi$ are conjugate-linear in the first variable. A function $\psi$ on $[\auU(S)]$ is
\emph{rapidly decreasing} if it is continuous and $|\psi(g)|\le C_k\,\mathrm{ht}(g)^{-k}$ for every
$k$, $\mathrm{ht}$ the height of (SW) in Lemma~\ref{pc:SW}; right translates of rapidly decreasing
functions are rapidly decreasing ($\mathrm{ht}(gx)\le c(x)\mathrm{ht}(g)$), and the product of a
rapidly decreasing function with a function bounded by a power of the height is integrable on
$[\auU(S)]$ (reduction theory). For an integrable function $\kappa$ on $\auU(S)_\infty$ and a bounded
measurable $\psi$ on $\auU(S)(\auA)$ put
\[
(\kappa*\psi)(g)=\int_{\auU(S)_\infty}\kappa(x)\,\psi(gx)\,dx .
\]
If $\psi$ is continuous and bounded, so is $\kappa*\psi$ (dominated convergence), and
$\kappa*\psi$ is invariant on the left under whatever $\psi$ is, and on the right under every finite
adelic element that leaves $\psi$ invariant on the right (finite and infinite components commute).

\subsubsection{Holomorphic vectors and their coefficients}\label{psr:ss-hol}

Let $X$ be a space of signature $(p,q)$ at every infinite place, $N=p+q$, with the order $o$ of its
frame putting the positive indices first; $\pspG=\auU(X)_\infty=\prod_\tau G_\tau$, $B$ and $\mathfrak b$
(the upper triangular matrices for $o$ at every place) are as in \S\ref{psq:sec}, and for
$Y=(Y_\tau)\in\mathfrak b$ and a family of weights $u=(u_\tau)$ put $u(Y)=\sum_{\tau,a}u_{\tau,a}
Y_{\tau,aa}$.

\begin{definition}\label{psr:def-hol}
Let $\sigma$ be a unitary representation of $\pspG$ on a Hilbert space $\mathcal H$. A
\emph{holomorphic vector of weight $u$} for $\sigma$ is a smooth vector $v\in\mathcal H$ with
$d\sigma(Y)v=u(Y)v$ for all $Y\in\mathfrak b$, where $d\sigma$ is extended complex-linearly to the
complexified Lie algebra $\prod_\tau\mathfrak{gl}_N(\C)\supseteq\mathfrak b$.
\end{definition}

\noindent For the real Weil representation $\Omega$ on the Fock space of $\auHom(X,Y)\otimes\R$
(with the archimedean components of characters), the polynomial vectors are smooth, and the
complexified Lie algebra acts on them by polynomial differential operators preserving the space of
polynomials (Howe~\cite[(2.2)--(2.3)]{HoweJAMS}: in the Fock model the $\widetilde{\auU}$-finite vectors are the polynomials, and the complexified Lie algebra acts by the operators $z_iz_j$, $\partial^2/\partial z_i\partial z_j$, $\frac12(z_i\partial/\partial z_j+\partial/\partial z_j\, z_i)$; $K$-finite vectors are smooth). So for a polynomial vector
the condition of Definition~\ref{psr:def-hol} is a finite algebraic condition.

\begin{lemma}[Coefficients of holomorphic vectors]\label{psr:coef}
Let $u$ be dominant of gap $\ge N^2+N+1$ at every place, $\sigma$ a unitary representation of
$\pspG$, and $v$ a holomorphic vector of weight $u$ for $\sigma$. Then for $h\in\pspG$
\[
\langle v,\sigma(h)v\rangle=\langle v,v\rangle\,\pspm_u(h),
\]
with $\pspm_u$ the coefficient of \S\ref{psp:ss-notation}, as made precise before
Lemma~\ref{psq:GB}. In particular $h\mapsto\langle v,\sigma(h)v\rangle$ is integrable on $\pspG$.
\end{lemma}
\begin{proof}
Put $c(h)=\langle v,\sigma(h)v\rangle$. For $Y$ in the Lie algebra $\mathfrak t$ of the diagonal
unitary torus $T$, $\frac d{ds}\sigma(e^{sY})v=\sigma(e^{sY})d\sigma(Y)v=u(Y)\sigma(e^{sY})v$, so
$\sigma(e^Y)v=e^{u(Y)}v$; as $T$ is connected, $\sigma(t)v=t^uv$ for $t\in T$. The part
\emph{Extension} of the proof of Proposition~\ref{psq:mhol} uses of the vector only that it is smooth,
that $\sigma(t)v=t^uv$ on $T$, and the identity $d\sigma(Y)v=u(Y)v$ on $\mathfrak b$ (the
\emph{Claim} there); it therefore applies verbatim to $c$ and shows that $c$ is holomorphic of weight
$u$ (the \emph{Extension} step is local, so for the product group it is done at each place). Next, for
$\theta\in\R/\Z$, $\mathrm{rot}(\theta)\in T$ and
\[
c(\mathrm{rot}(\theta)h)=\langle\sigma(\mathrm{rot}(\theta))^{-1}v,\sigma(h)v\rangle=
\aubar{\mathrm{rot}(\theta)^{-u}}\,c(h)=e^{2\pi i\theta\langle o,u\rangle}c(h).
\]
By Proposition~\ref{psq:unique} for $\pspG=\prod_\tau G_\tau$ (its proof applies to the product group with the diagonal $\mathrm{rot}(\theta)$, as in the proof of Theorem~\ref{psq:H123}), $c=c(1)\phi_u=\langle v,v\rangle\phi_u$, and $\phi_u=\pspm_u$ by
Proposition~\ref{psq:mhol}. Finally $\pspm_u$ is the coefficient $c_\tau$ of \S\ref{pds:sec} for the
parameter of Lemma~\ref{psq:param}, whose gap is $\ge N^2+1$, at each place; it is integrable by
Theorem~\ref{pds:tail}, and so is the product over the places.
\end{proof}

\subsubsection{The series of translates in Weil's range}\label{psr:ss-series}

Lemma~\ref{pc:series}(b),(c) bounds the series of Rallis' formula only for $M'\ge6n$ and $n=r$.
Here we need every $M'>2n$. We use the notation of \S\ref{pc:ss-data} before Lemma~\ref{pc:series}:
$S$ a space of rank $n$, $M'\ge1$, $\alpha$ unramified outside the modulus $\psrm$, $G^\psrm$, $K^\psrm$
and $\psrmc=\psrmc_{M',\alpha}$ on $G^\psrm$. As in the proof of Lemma~\ref{pex:vac} (which does not
use $n=1$ at this point), $G^\psrm/K^\psrm$ is the restricted product of the $G_v/K_v$ over the places
$v\notin S(\psrm)$ of $E^+$, and $\psrmc(\gamma)=\prod_v\psrmc_v(\gamma_v)$ with
$\psrmc_v(\gamma_v)=\alpha_v(\det\gamma_v)^{-1}\langle\psvac_v,\omega_v(\gamma_v,1)\psvac_v\rangle$, equal to
$1$ on $K_v$; each $\psrmc_v$ is bi-invariant under $K_v$, and by Lemma~\ref{pc:series}(a)
$|\psrmc_v(\gamma_v)|=[L_v:L_v\cap\gamma_vL_v]^{-M'/2}$, $L_v$ the localisation of $L_S$. Put
\[
\varepsilon_v=\sum_{\gamma\in G_v/K_v,\ \gamma\notin K_v}|\psrmc_v(\gamma)| .
\]
Recall that every residue field of $E$ has at least $5$ elements (proof of Lemma~\ref{pc:series}(b)),
and that $E/E^+$ is unramified at the finite places, so each finite place $v$ of $E^+$ is split or
inert.

\begin{lemma}[Neighbours]\label{psr:neigh}
Let $F$ be a non-archimedean local field, $K/F$ an unramified quadratic extension with ring of
integers $\auO$, $\varpi\in F$ a uniformiser, $q=\#\auO/\varpi\auO$ and $q_0=q^{1/2}$ the order of the
residue field of $F$. Let $V$ be a non-degenerate hermitian space of dimension $n$ over $K$. For a
lattice $L$ put $L^\sharp=\{x\in V:\langle x,L\rangle\subseteq\auO\}$; $L$ is \emph{self-dual} if
$L^\sharp=L$. For self-dual $L,L'$ define $d(L,L')$ by $q^{d(L,L')}=[L:L\cap L']$, and call $L'$ a
\emph{neighbour} of $L$ if $d(L,L')=1$.
\begin{enumerate}
\item[(i)] If $L\ne L'$ are self-dual, $L$ has a neighbour $L_1$ with $d(L_1,L')\le d(L,L')-1$.
\item[(ii)] Every self-dual $L$ has at most $\mathsf N=q_0(q^n-1)/(q-1)$ neighbours.
\item[(iii)] For every self-dual $L$ and $k\ge0$, $\#\{L'\text{ self-dual}:d(L,L')\le k\}\le
(\mathsf N+1)^k$.
\end{enumerate}
\end{lemma}
\begin{proof}
Lattices are $\auO$-modules; all indices are powers of $q$. Fix a Haar measure on $V$. The product
$\mathrm{vol}(M^\sharp)\mathrm{vol}(M)$ does not depend on the lattice $M$: writing $M=aM_0$ with
$a\in\auGL(V)$, one has $M^\sharp=(a^*)^{-1}M_0^\sharp$ with $a^*$ the adjoint for the form, and
$|\det a^*|_K=|\aubar{\det a}|_K=|\det a|_K$. Hence $\mathrm{vol}(M^\sharp)=\mathrm{vol}(M)$ for every
self-dual $M$, $[M^\sharp:M'^\sharp]=[M':M]$ for
$M\subseteq M'$, and $\mathrm{vol}(H^\sharp)=\mathrm{vol}(L)^2/\mathrm{vol}(H)$ for every lattice $H$, $L$
self-dual. In particular all self-dual lattices have the same volume, so $[L:L\cap
L']=[L':L\cap L']$ and $d$ is symmetric; and an integral lattice ($M\subseteq M^\sharp$) with the volume
of a self-dual lattice is self-dual. Two self-dual lattices one of which contains the other are equal.
Since $\varpi\in F$, $\aubar\varpi=\varpi$, and the residue field automorphism induced by the conjugation
of $K/F$ is the non-trivial one ($K/F$ unramified).

(i) $L'\not\subseteq L$, as $L'\ne L$. Take $z\in L'\smallsetminus L$ and $j\ge1$ minimal with
$\varpi^jz\in L$, and put $x=\varpi^{j-1}z$: then $x\in L'\smallsetminus L$ and $\varpi x\in L$. For
$l\in L$, $\langle l,\varpi x\rangle\in\auO$, so $\langle l,x\rangle\in\varpi^{-1}\auO$, and the map
$l\mapsto\langle l,x\rangle\bmod\auO$ from $L$ to $\varpi^{-1}\auO/\auO$ (a group of order $q$) is
$\auO$-semilinear and non-zero (otherwise $x\in L^\sharp=L$), hence surjective. Its kernel
$L_x=\{l\in L:\langle l,x\rangle\in\auO\}$ has index $q$ in $L$. Put $L_1=L_x+\auO x$. It is integral:
$L_x$ is, $\langle L_x,x\rangle\subseteq\auO$, and $\langle x,x\rangle\in\auO$ as $x\in L'$. Moreover
$\varpi x\in L_x$ ($\langle\varpi x,x\rangle=\varpi\langle x,x\rangle$) and $x\notin L_x$, so
$L_1/L_x\cong\auO/\varpi$ has order $q$ and $\mathrm{vol}(L_1)=\mathrm{vol}(L)$: $L_1$ is self-dual.
$L\cap L_1\supseteq L_x$, and $L\cap L_1\ne L$ (else $L\subseteq L_1$, so $L=L_1\ni x$); as $L/L_x$ is
simple, $L\cap L_1=L_x$ and $d(L,L_1)=1$. Finally let $\Lambda=L\cap L'$. For $l\in\Lambda$,
$\langle l,x\rangle\in\auO$ since $l,x\in L'$; so $\Lambda\subseteq L_x\subseteq L_1$, and
$L_1\cap L'\supseteq\Lambda+\auO x$. Since $x\notin\Lambda$ and $\varpi x\in L\cap L'=\Lambda$,
$[\Lambda+\auO x:\Lambda]=q$, so $[L':L_1\cap L']\le[L':\Lambda]/q$, i.e.\ $d(L_1,L')\le d(L,L')-1$.

(ii) Let $L_1$ be a neighbour and $H=L\cap L_1$, of index $q$ in $L$. $L/H$ is an $\auO$-module of
length one, so $\varpi L\subseteq H$, and $H$ corresponds to a hyperplane of the $n$-dimensional
$\auO/\varpi$-space $L/\varpi L$: there are at most $(q^n-1)/(q-1)$ choices of $H$. Fix $H$. The
integral lattices containing $H$ lie in $H^\sharp$, and $[H^\sharp:H]=q^2$. If $H^\sharp/H$ were
cyclic, its only submodule of order $q$ would be $L/H$, and $L_1=L$; so $H^\sharp/H\cong(\auO/\varpi)^2$
and $\varpi H^\sharp\subseteq H$. For $x,y\in H^\sharp$, $\langle\varpi x,y\rangle\in\auO$, so
$\langle x,y\rangle\in\varpi^{-1}\auO$, and $b(\bar x,\bar y)=\varpi\langle x,y\rangle\bmod\varpi$ is a
well defined hermitian form on the plane $\mathsf V=H^\sharp/H$ over $\psrF_q=\auO/\varpi$ (relative to
$\psrF_q/\psrF_{q_0}$); it is non-degenerate, because $b(\bar x,\cdot)=0$ means $\langle x,
H^\sharp\rangle\subseteq\auO$, i.e.\ $x\in H^{\sharp\sharp}=H$. The lattice $H+\auO x$ ($x\in H^\sharp$)
is integral iff $\langle x,x\rangle\in\auO$, iff $b(\bar x,\bar x)=0$ (as $\langle x,x\rangle\in
\varpi^{-1}\auO\cap F$). So $L_1/H$ is an isotropic line of $\mathsf V$ other than the isotropic line
$e_1=L/H$. Choose $e_2\in\mathsf V$ with $b(e_1,e_2)=1$ (non-degeneracy). The lines other than $\psrF_q e_1$ are
the $\psrF_q(e_2+se_1)$, $s\in\psrF_q$, and $b(e_2+se_1,e_2+se_1)=b(e_2,e_2)+s+\aubar s$ (or with $s$
and $\aubar s$ exchanged, according to the side on which $b$ is linear). As $b(e_2,e_2)\in\psrF_{q_0}$ and
the trace $\psrF_q\to\psrF_{q_0}$ is surjective and $\psrF_{q_0}$-linear, exactly $q/q_0=q_0$ values of $s$
give an isotropic line. Hence at most $q_0$ neighbours contain $H$, and at most $\mathsf N$ in all.

(iii) Let $a_k$ be the supremum over self-dual $L$ of the left side; $a_0=1$. By (i), every $L'\ne L$
with $d(L,L')\le k$ has $d(L_1,L')\le k-1$ for some neighbour $L_1$ of $L$, so by (ii)
$a_k\le1+\mathsf Na_{k-1}$, and $a_k\le(\mathsf N+1)^k$ by induction.
\end{proof}

\begin{lemma}[Local sums]\label{psr:local}
Let $M'\ge2n+1$ and let $v\notin S(\psrm)$ be a finite place of $E^+$, $w$ a place of $E$ above it and
$q_w$ the order of its residue field. If $v$ splits, $\varepsilon_v\le3.6\,q_w^{-3/2}$ and
$\varepsilon_v\le0.28$. If $v$ is inert, $q_w\ge9$, $\varepsilon_v\le1.8\,q_w^{-1}$ and $\varepsilon_v\le
0.2$. Consequently $\sum_v\varepsilon_v<\infty$.
\end{lemma}
\begin{proof}
\emph{Split $v$.} The proof of Lemma~\ref{pc:series}(b), up to and including the identity
$\bigl(\sum_ks_kq^{-kM'/2}\bigr)^2=\prod_{i=0}^{n-1}(1-q^{i-M'/2})^{-2}$, does not use $M'\ge6n$; it gives
$1+\varepsilon_v\le\prod_{i=0}^{n-1}(1-q^{i-M'/2})^{-2}$, $q=q_w\ge5$. As $i-M'/2\le-\frac32-(n-1-i)$,
\[
\log(1+\varepsilon_v)\le2\sum_{j\ge0}-\log\bigl(1-q^{-3/2-j}\bigr)\le\frac{2\sum_jq^{-3/2-j}}{1-q^{-3/2}}
=\frac{2q^{-3/2}}{(1-q^{-1})(1-q^{-3/2})}\le2.75\,q^{-3/2},
\]
using $-\log(1-t)\le t/(1-t)$ and $(1-\frac15)(1-5^{-3/2})\ge0.728$. This is $\le2.75\cdot5^{-3/2}\le
0.246$, so $\varepsilon_v\le e^{0.246}-1\le0.28$, and $\varepsilon_v\le e^{0.246}\cdot2.75\,q^{-3/2}\le
3.6\,q^{-3/2}$.

\emph{Inert $v$.} Here $E_w/E^+_v$ is unramified quadratic and $q_w=q_v^2$; as $q_w\ge5$, $q_v\ge3$ and
$q_w\ge9$. The map $\gamma\mapsto\gamma L_v$ is injective from $G_v/K_v$ to the set of self-dual
lattices of $S_v$ ($L_v$ is self-dual because $L_S$ is unimodular), and $|\psrmc_v(\gamma)|=
q_w^{-d(L_v,\gamma L_v)M'/2}$ in the notation of Lemma~\ref{psr:neigh}. By Lemma~\ref{psr:neigh}(iii),
with $X=q_w^{-M'/2}\le q_w^{-n-1/2}$,
\[
\varepsilon_v\le\sum_{k\ge1}\#\{L':d(L_v,L')=k\}X^k\le\sum_{k\ge1}\bigl((\mathsf N+1)X\bigr)^k .
\]
Now $\mathsf N\le q_w^{1/2}q_w^{n-1}\cdot q_w/(q_w-1)\le\frac98q_w^{n-1/2}$, so $(\mathsf N+1)X\le
q_w^{-1}(\frac98+q_w^{1/2-n})\le\frac32q_w^{-1}\le\frac16$, and $\varepsilon_v\le\frac{3/2}{1-1/6}q_w^{-1}
=1.8\,q_w^{-1}\le0.2$.

\emph{Summation.} Above a rational prime $p$ there are at most $f=[E^+:\Q]$ places $v$ of $E^+$, and
$q_w\ge p$ if $v$ splits, $q_w\ge p^2$ if $v$ is inert. So $\sum_v\varepsilon_v\le f\sum_p(3.6\,
p^{-3/2}+1.8\,p^{-2})<\infty$.
\end{proof}

\begin{proposition}[The series for $M'>2n$]\label{psr:series}
Let $M'>2n$. \textup{(a)} $\sum_{\gamma\in G^\psrm/K^\psrm}|\psrmc(\gamma)|<\infty$. \textup{(b)} On the space
$\mathcal B$ of bounded functions on $\auU(S)(\auA)$ that are right invariant under $K^\psrm$, the
series $\psrR=\sum_{\gamma\in G^\psrm/K^\psrm}\psrmc(\gamma)\rho(\gamma)$ converges absolutely in the
supremum norm and maps $\mathcal B$ to itself; it commutes with left translations, with right
translations by $\auU(S)_\infty$ and by the components at $S(\psrm)$, and with $\kappa*$ for every
integrable $\kappa$ on $\auU(S)_\infty$; it maps continuous functions to continuous functions.
\textup{(c)} $\psrR$ is injective on $\mathcal B$.
\end{proposition}
\begin{proof}
(a) The sum is $\prod_v(1+\varepsilon_v)\le\exp\sum_v\varepsilon_v$, finite by Lemma~\ref{psr:local} (the
product formula holds for the sums of the non-negative multiplicative function $|\psrmc|$ over the
restricted product). (b) $\psrmc$ is right $K^\psrm$-invariant, so the terms are well defined on $\mathcal
B$; it is left $K^\psrm$-invariant, so (reindexing the cosets $\gamma\mapsto k\gamma$) the sum is again
right $K^\psrm$-invariant. Absolute convergence in the supremum norm follows from (a), and gives the
continuity statement (uniform limit of continuous functions). The commutation with left translations
is clear; right translations by $\auU(S)_\infty$ and by the components at $S(\psrm)$ commute with the
$\rho(\gamma)$; and $\kappa*(\psrR\psi)=\psrR(\kappa*\psi)$ by Fubini, the double series--integral
converging absolutely ($\int|\kappa|\cdot\sum|\psrmc|\cdot\sup|\psi|<\infty$).

(c) By Lemma~\ref{psr:local} choose a finite set $\mathsf T$ of places $v\notin S(\psrm)$ with
$\sum_{v\notin\mathsf T\cup S(\psrm)}\varepsilon_v\le\frac13$. For $v\in\mathsf T$ let
$\psrR_v=\sum_{\gamma\in G_v/K_v}\psrmc_v(\gamma)\rho(\gamma)$, and let $\psrR^{\mathsf T}$ be the analogous
series over the restricted product of the $G_v/K_v$, $v\notin\mathsf T\cup S(\psrm)$, with the
coefficient $\prod_{v\notin\mathsf T}\psrmc_v$. By (a) and the product structure of $\psrmc$ and of
$G^\psrm/K^\psrm$, the absolutely convergent series $\psrR$ may be summed iteratively:
$\psrR=\psrR_{v_1}\circ\dots\circ\psrR_{v_k}\circ\psrR^{\mathsf T}$ on $\mathcal B$, $\mathsf T=\{v_1,\dots,v_k\}$,
each factor mapping $\mathcal B$ to itself by the argument of (b). The coset of $1$ contributes the
identity (as $\psrmc_v=1$ on $K_v$), so
\[
\|\psrR_v\psi-\psi\|_\infty\le\varepsilon_v\|\psi\|_\infty,\qquad
\|\psrR^{\mathsf T}\psi-\psi\|_\infty\le\Bigl(\prod_{v\notin\mathsf T\cup S(\psrm)}(1+\varepsilon_v)-1\Bigr)
\|\psi\|_\infty\le(e^{1/3}-1)\|\psi\|_\infty .
\]
Since $\varepsilon_v\le0.28<1$ (Lemma~\ref{psr:local}) and $e^{1/3}-1<\frac12$, each factor satisfies
$\|\psrR_\bullet\psi\|_\infty\ge c\|\psi\|_\infty$ with $c>0$, hence is injective on $\mathcal B$, and so is
the composite.
\end{proof}

\noindent For $M'\ge6n$ and $n=r$ this is consistent with Lemma~\ref{pc:series}(c), which gives
the stronger bound $\|\psrR\psi\|_\infty\ge\frac12\|\psi\|_\infty$ directly.

\subsubsection{The inner product formula for rapidly decreasing functions}\label{psr:ss-rallis}

\begin{lemma}[Weil's formula for a definite target]\label{psr:SW}
The statement \textup{(SW)} of Lemma~\ref{pc:SW}, including the absolute convergence and the bounds by
a power of the height of $g\in\auU(\mathcal W)(\auA)$, holds when $T$ is positive definite at every infinite place and $M'>2n$: these imply the
hypotheses (i), (ii) of Lemma~\ref{pc:SW}.
\end{lemma}
\begin{proof}[References and comments]
This is the argument of Step~1 of the proof of Proposition~\ref{pex:rallis}, where it is given for
$n=1$, $M'=4$; it does not use $n=1$. Weil's identity~\cite[Th\'eor\`eme~5]{WeilSiegel} holds under (i) absolute
convergence of the theta integral for all $\Phi$ and (ii) absolute convergence of the Siegel Eisenstein
series of $\auU(\mathcal W)\cong\auU(n,n)$ attached to $\Phi$. (i): $T$ is anisotropic over $E^+$, being
definite at the infinite places, so $[\auU(T)]$ is compact, and $h\mapsto\theta(\omega(g,h)\Phi)$ is
continuous; $\int_{[\auU(T)]}|\theta(\omega(g,h)\Phi)|\,dh$ is bounded by a power of the height of $g$
because $|\theta(\omega(g,h)\Phi)|$ is, uniformly for $h$ in a compact set of representatives (moderate
growth of theta functions, Lemma~\ref{pcx:thetagrowth}). (ii): the Eisenstein series at the point attached to $M'$ converges
absolutely for $M'>2n$~\cite[Th\'eor\`eme~1]{WeilSiegel}, and $\sum_L|\psZ_L(\omega(g,1)\Phi)|$ is bounded by a
power of the height of $g$ by Lemma~\ref{psr:Egrowth} below. The identification of the terms
with the Lagrangians is as in the comments on Lemma~\ref{pc:SW}.

What the applications use is only this: in Step~1 of Proposition~\ref{psr:rallis} (and of
Proposition~\ref{pc:rallis}), the bound of $\sum_\delta|k(g'^{-1}\delta g)|$ by a power of the heights of
$g,g'\in\auU(S)(\auA)$, which by Lemma~\ref{pc:clean} is $|c'|^{-1}\sum_L|\psZ_L(\omega((g,g'),1)\Phi)|$ for
$\Phi=F\otimes\mathsf CF$; it is a bound for the series of absolute values, so it does not follow from the
identity $I(\Phi)=E(\Phi)$ and the growth of the (finite-sum, definite-target) theta side alone, and
Lemma~\ref{psr:Egrowth} is needed. The heights of $\auU(\pcW)$ and of $\auU(S)$ compare as follows: for the frame
of $\pcW=S\oplus S^-$ that is the sum of the frames, $\Lambda_{(g,g')}=\Lambda_g\oplus\Lambda_{g'}$ with
$|(\xi,\xi')|^2=|\xi|_g^2+|\xi'|_{g'}^2$, so $\mathrm{ht}(g,g')=\max(\mathrm{ht}(g),\mathrm{ht}(g'))$; another frame
of $\pcW$ changes $\lambda_1$ by at most a constant factor (a fixed rational change of basis changes lengths
by bounded factors and integrality by a bounded denominator).
\end{proof}

\begin{lemma}[Growth of the Siegel Eisenstein series]\label{psr:Egrowth}
Let $\pcW$ be a space of rank $2n$ with a Lagrangian, $T$ a space of rank $M'>2n$, and $\Phi$ a datum of
$\auHom(\pcW,T)$. There are $C,A$ with $\sum_L|\psZ_L(\omega(x,1)\Phi)|\le C\,\mathrm{ht}(x)^A$ for all
$x\in\auU(\pcW)(\auA)$, the sum over all Lagrangians of $\pcW$.
\end{lemma}
\begin{proof}
Put $G=\auU(\pcW)$, $P$ the stabiliser of a Lagrangian $L_0$, so that the sum is $\sum_{\gamma\in P(E^+)
\backslash G(E^+)}f_\Phi(\gamma x)$ with $f_\Phi(y)=|\psZ_{L_0}(\omega(y,1)\Phi)|$ (comments on
Lemma~\ref{pc:SW}); in Weil's notation $f_\Phi$ is the function $f_\Phi(\pi(\mathbf S))=|\mathbf S\Phi(0)|$ of
\cite[no.~39]{WeilSiegel}, on the group $Ps_\auA=G(\auA)$, and $E(\Phi)$ is the series $\sum_{s\in R}$ over
representatives $R$ of $P(E^+)\backslash G(E^+)$. Weil's Lemme~18~\cite[no.~39]{WeilSiegel} says: if $f_0$ is
continuous and $>0$ on $G(\auA)$ with $f_0(py)=\psi(p)f_0(y)$ for $p\in P(\auA)$ ($\psi(p)=|\lambda|_X^{1/2}$, $\lambda$
the Levi component of $p$), then $E(\Phi)$ converges absolutely for all $\Phi$ if and only if $\sum_Rf_0(s)<\infty$;
such $f_0$ exist (p.~57). We need, for our fixed $\Phi$, a constant $c$ with $f_\Phi\le cf_0$ on $G(\auA)$, and
we prove it directly. Let $C_0$ be compact with $G(\auA)=P(\auA)C_0$ ($P(\auA)\backslash G(\auA)$ is compact). The Siegel parabolic acts on the zero term by a character: for $p\in P(\auA)$ with Levi
component $\lambda$, the Weil representation gives $(\omega(p,1)\Phi')(0)=\chi(p)\,|\lambda|_X^{1/2}\Phi'(0)$ for
every datum $\Phi'$, with $|\chi(p)|=1$ (the unipotent radical acts at $0$ by the trivial character and $\lambda$
acts by $\Phi'\mapsto\chi(\lambda)|\lambda|_X^{1/2}\Phi'(\cdot\,\lambda)$, with $|\lambda|_X$ the module of $\lambda$ on
$X_\auA$, exactly as in \cite[no.~39]{WeilSiegel}). Hence $f_\Phi(py)=\psi(p)f_\Phi(y)$, the same rule as for
$f_0$. The function $f_\Phi$ is continuous ($y\mapsto\omega(y,1)\Phi$ is continuous into the Schwartz space and
evaluation of the zero term is a continuous functional), so $f_\Phi/f_0$ is continuous and left
$P(\auA)$-invariant on $G(\auA)$, hence bounded by its maximum $c$ on $C_0$. By Th\'eor\`eme~1~\cite[no.~40]{WeilSiegel} ($M'>2n$ is condition (B) in the hermitian case,
see the comments on Lemma~\ref{pc:SW}) $E(\Phi)$ converges absolutely for all $\Phi$, so
$E_0(y):=\sum_{\gamma}f_0(\gamma y)$ is finite at $y=1$; also $E_0(\gamma y)=E_0(y)$ for rational $\gamma$.

\emph{Translation.} For
$y=pc\in G(\auA)$ and $u\in G(\auA)$ write $cu=p'c'$ with $c'\in C_0$; then $f_0(yu)=\psi(p)\psi(p')f_0(c')
\le\psi(p)\,\psi(p')\max_{C_0}f_0$ and $f_0(y)\ge\psi(p)\min_{C_0}f_0$, so
$f_0(yu)\le\kappa(u)f_0(y)$ with $\kappa(u)=\sup_{c\in C_0}\psi(p'(c,u))\cdot\max_{C_0}f_0/\min_{C_0}f_0$. Summing,
$E_0(u)\le\kappa(u)E_0(1)$, and $\sum_L|\psZ_L(\omega(u,1)\Phi)|=\sum_\gamma f_\Phi(\gamma u)\le cE_0(u)$.

\emph{Bound for $\kappa$.} $\psi(p')$ is $|\det(p'|_{L_0})|^{\kappa_0}$ (idelic absolute value) for a real
$\kappa_0$, since $|\lambda|_X$ is the module of $\lambda$ acting on $X_\auA=\auHom(L_0,T)(\auA)$-type
coordinates, a power of $|\det\lambda|$. With $p'=cuc'^{-1}$ and $p'^{-1}=c'u^{-1}c^{-1}$ we get
$|\det(p'|_{L_0})|^{\pm1}\le\prod_v\|\wedge^n p'^{\pm1}_v\|_v\le C_5(\|u\|_\auA\|u^{-1}\|_\auA)^n$, where
$\|y\|_\auA=\prod_v\|y_v\|_v$ (operator norms in the frame; almost all factors are $1$) and $c,c'$ run over a
compact set. Hence $\kappa(u)\le C_6(\|u\|_\auA\|u^{-1}\|_\auA)^{n|\kappa_0|}$.

\emph{Reduction.} Given $x$, Lemma~\ref{pcx:reduce} gives $\gamma$ rational with $u=\gamma x$, $u_f,u_f^{-1}$ in a
fixed compact set (so their finite norms are bounded) and $\|u_\infty\|,\|u_\infty^{-1}\|\le C\,\mathrm{ht}(x)$.
So $\sum_L|\psZ_L(\omega(x,1)\Phi)|=\sum_\gamma f_\Phi(\gamma u)\le cE_0(1)\kappa(u)\le C\,\mathrm{ht}(x)^{2n|\kappa_0|}$.
\end{proof}

\noindent The comments on Lemma~\ref{psr:SW} are complete with Lemma~\ref{psr:Egrowth} and, for the theta
integral, Lemma~\ref{pcx:thetagrowth} (with $\mathrm{ht}(h)$ bounded on a compact set of representatives of
$[\auU(T)]$).

\begin{proposition}[Inner product formula]\label{psr:rallis}
Let $S$ be a space of rank $n$ (any signature), $T$ a positive definite space of rank $M'>2n$,
$\psrm$ a modulus with $S(\psrm)\ne\emptyset$, $\alpha,\beta$ characters of $\auU(1)$ unramified outside
$\psrm$ and tamely ramified at $\psrm$, $\Psi$ a function on $\pcMat_{M'\times n}(\auO_E/\psrm)$ sharp for
a group $J$ (Definition~\ref{pc:sharp}), and $P$ a smooth vector of the Fock space of
$\auHom(S,T)\otimes\R$ such that
\[
\kappa(x)=\langle P,\omega_{\alpha,\beta,\infty}(x,1)P\rangle\qquad(x\in\auU(S)_\infty)
\]
is integrable. Put $\Theta=\Theta_{\alpha,\beta}(P,\Psi)$, $\psrR=\psrR_{M',\alpha}$
(Proposition~\ref{psr:series}) and
\[
A=c'\,\mathrm{vol}([\auU(T)])\,\frac{\|\Psi\|^2}{\#\pcMat_{M'\times n}(\auO_E/\psrm)\,[K_S:J]}\ne0 ,
\]
$c'$ as in Lemma~\ref{pc:clean}(b). Then for all rapidly decreasing $\psi,\psi'$ on $[\auU(S)]$ with $\psi$
right invariant under $J$,
\[
\langle\Theta\psi',\Theta\psi\rangle_{[\auU(T)]}=A\int_{\auU(S)_\infty}\kappa(x)\,\langle\psi',
\rho(x)\psrR\psi\rangle\,dx=A\,\langle\psi',\psrR(\kappa*\psi)\rangle ,
\]
all integrals converging absolutely.
\end{proposition}
\begin{proof}
We follow the proof of Proposition~\ref{pc:rallis}, with $y=\psi$, $y'=\psi'$; $F=[P,\Psi]$,
$\omega_\alpha(g)=\omega_{\alpha,\beta}(g,1)$, $k(g)=\langle F,\omega_\alpha(g)F\rangle$.

\emph{Step 1} there (doubling) uses (SW), Lemma~\ref{pc:clean} and (W7). Lemma~\ref{pc:clean} does not
use $M'\ge6n$, and its hypothesis holds: at a place of $S(\psrm)$, which splits, the component of $F$ is
$\Psi^\natural$, supported on matrices whose reduction $\bar x$ has $A_0\bar x\in\bar J$ invertible for
the matrix $A_0$ of the definition of sharpness, hence on injective matrices. By
Lemma~\ref{psr:SW} in place of (SW), \eqref{pc:eq-doubling} holds, with
$\sum_\delta|k(g'^{-1}\delta g)|$ bounded by a power of the heights of $g,g'$.

\emph{Step 2} there (unfolding) uses only that $y,y'$ are rapidly decreasing, which is our hypothesis:
the function $(g,g')\mapsto|\psi'(g')||\psi(g)|\sum_\delta|k(g'^{-1}\delta g)|$ is integrable on
$[\auU(S)]^2$, and
\[
\langle\Theta\psi',\Theta\psi\rangle=c'\,\mathrm{vol}([\auU(T)])\int_{\auU(S)(\auA)}k(g)\,\langle\psi',
\rho(g)\psi\rangle\,dg ,
\]
the integral converging absolutely.

\emph{Step 3.} As in Step~3 there, $k(g)=\kappa(g_\infty)\,\psrmc(g^\psrm)\,k_\psrm(g_\psrm)$ with
$k_\psrm=\|\Psi^\natural\|^2\,\mathbf 1_J$ on the components at $S(\psrm)$ (this uses only the sharpness
of $\Psi$ and Lemma~\ref{pc:level}) and $\|\Psi^\natural\|^2=\|\Psi\|^2/\#\pcMat_{M'\times n}(\auO_E/\psrm)$.
Since $|\langle\psi',\rho(g)\psi\rangle|\le\|\psi'\|_{L^1}\|\psi\|_\infty$, and $\kappa$, $\psrmc$ (summed
over $G^\psrm/K^\psrm$, Proposition~\ref{psr:series}(a); $K^\psrm$ has volume $1$) and $k_\psrm$ are
integrable, Fubini's theorem allows us to integrate first over the finite adelic points. As $\psi$
is right invariant under $J\supseteq K^\psrm$, $\int_{G^\psrm}\psrmc(g)\rho(g)\psi\,dg=\psrR\psi$, and the
integral over the components at $S(\psrm)$ multiplies by $\|\Psi^\natural\|^2\mathrm{vol}(J_\psrm)=
\|\Psi^\natural\|^2/[K_S:J]$. This gives the first formula; the second is
$\int\kappa(x)\langle\psi',\rho(x)\psrR\psi\rangle dx=\langle\psi',\kappa*\psrR\psi\rangle$ (Fubini, as
$\kappa\in L^1$, $\psrR\psi$ bounded, $\psi'\in L^1$) and $\kappa*\psrR\psi=\psrR(\kappa*\psi)$
(Proposition~\ref{psr:series}(b)). Finally $A\ne0$: $c'\ne0$, the volume is positive and $\Psi\ne0$.
\end{proof}

\begin{corollary}[Vanishing lifts]\label{psr:vanish}
In the situation of Proposition~\ref{psr:rallis}, if $\psi$ is rapidly
decreasing, right invariant under $J$, and $\Theta\psi=0$, then $\kappa*\psi=0$.
\end{corollary}
\begin{proof}
Put $\psi_1=\kappa*\psi$ and $z=\psrR\psi_1$. As $\psi$ is continuous and bounded, left invariant under
$\auU(S)(E^+)$ and right invariant under $J\supseteq K^\psrm$, $\psi_1$ is continuous, bounded, left
$\auU(S)(E^+)$-invariant and right $J$-invariant, and so is $z$ (Proposition~\ref{psr:series}(b); $J$
lies in $K_S$, and $\psrR$ commutes with the components at $S(\psrm)$ and preserves $K^\psrm$-invariance).
By Proposition~\ref{psr:rallis} with $\Theta\psi=0$, $\langle\psi',z\rangle=0$ for every rapidly
decreasing $\psi'$. For $\gamma\in G^\psrm$ and $x\in\auU(S)_\infty$ the translate
$\psi_{\gamma x}=\psi(\cdot\,\gamma x)$ is rapidly decreasing, so $\langle z,\psi_{\gamma x}\rangle=
\aubar{\langle\psi_{\gamma x},z\rangle}=0$, and by Fubini ($z$ bounded, $\psi\in L^1([\auU(S)])$, $\kappa\in
L^1$)
\[
\langle z,\rho(\gamma)\psi_1\rangle=\int_{\auU(S)_\infty}\kappa(x)\,\langle z,\psi_{\gamma x}\rangle\,dx=0 ,
\]
using $\psi(g\gamma x)=\psi(gx\gamma)$. Hence, $[\auU(S)]$ having finite volume and the series converging
absolutely and uniformly, $\langle z,z\rangle=\sum_\gamma\psrmc(\gamma)\langle z,\rho(\gamma)\psi_1\rangle=0$.
So $z=0$ almost everywhere, hence everywhere by continuity, and $\psi_1=0$ by
Proposition~\ref{psr:series}(c).
\end{proof}

\subsubsection{The reproducing formula}\label{psr:ss-R}

\begin{theorem}[Input (R)]\label{psr:R}
Let $X$ be a space of rank $n_X$ and signature $(p,q)$ at every infinite place, with the order $o$
putting the positive indices first; $Y$ a positive definite space of rank $M'>2n_X$; $\psrm$ a modulus
with $S(\psrm)\ne\emptyset$; $\alpha,\beta$ characters of $\auU(1)$ unramified outside $\psrm$ and tamely
ramified at $\psrm$; $\Omega=\omega_{\alpha,\beta,\infty}(\cdot,1)$ and $\Omega'=\omega_{\alpha,\beta,\infty}(1,
\cdot)$ the actions of $\pspG=\auU(X)_\infty$ and $\auU(Y)_\infty$ on the Fock space of $\auHom(X,Y)
\otimes\R$. Let $u$ be a dominant weight on $X$ of gap $\ge n_X^2+n_X+1$ at every place. Let $\mathcal Q$
be a finite-dimensional space of smooth Fock vectors \textup{(}for instance polynomials\textup{)}, stable
under $\Omega'(y)$ for all $y\in\auU(Y)_\infty$, that contains a non-zero holomorphic vector $v_0$ of
weight $u$ for $\Omega$ \textup{(}Definition~\ref{psr:def-hol}\textup{)}, and let $Q_1,\dots,Q_s$ be a
basis of $\mathcal Q$. Let $\iota\in\pcMat_{M'\times n_X}(\auO_E/\psrm)$ be the matrix of the inclusion of
frames and $h_1,\dots,h_{P_Y}\in\auU(Y)(\auA_f)$ representatives of the finite set
$\auU(Y)(E^+)\backslash\auU(Y)(\auA_f)/J_Y(\psrm)$. Then for every $g\in\auU(X)(\auA)$ there are
constants $a_{p,i}(g)$ such that for every rapidly decreasing $\psi$ on $[\auU(X)]$ right invariant
under $J_X(\psrm)$
\[
\int_{\pspG}\psi(gh)\,\pspm_u(h)\,dh=\sum_{p=1}^{P_Y}\sum_{i=1}^sa_{p,i}(g)\int_{[\auU(X)]}\psi(x)\,
\theta_{\alpha,\beta}\bigl(x,h_p;Q_i,\delta_\iota\bigr)\,dx ,
\]
all integrals converging absolutely.
\end{theorem}

The finiteness of the set of classes is Proposition~\ref{pf:classes} for a standard space $Y=E^{M'}$
(take $\mathsf G\in\psrm\cap\Z$, so that $J_Y(\psrm)\supseteq K(\mathsf G)$), and Borel's finiteness
theorem~\cite[Theorem~5.1]{BorelFin} in general. This is (R) of Proposition~\ref{psp:inputs},
with $\psrm=(\pi_M^2)$, the classes $h_p$ of $\auU(Y)$ at level $\pi_M^2$, the polynomials
$Q_{(p,i)}=Q_i$ and the constants $a_{(p,i)}$, under the hypotheses on $\mathcal Q$, $\alpha$, $\beta$ stated
here; the gap $\ge300\pspn$ of Proposition~\ref{psp:inputs} exceeds $t^2+t+1\ge n_X^2+n_X+1$.

\begin{proof}
\emph{Convergence.} $\pspm_u$ is integrable (Lemma~\ref{psr:coef}) and $\psi$ bounded; the kernel has
moderate growth in $x$ (Lemma~\ref{pcx:thetagrowth}) and $\psi$ decreases rapidly.

\emph{Linear forms.} Let $\mathcal V$ be the complex vector space of rapidly decreasing functions on
$[\auU(X)]$ right invariant under $J=J_X(\psrm)$. Fix $g$ and define linear forms on $\mathcal V$:
\[
\ell(\psi)=\int_{\pspG}\psi(gh)\pspm_u(h)\,dh,\qquad L_{p,i}(\psi)=\bigl(\Theta_{\alpha,\beta}(Q_i,
\delta_\iota)\psi\bigr)(h_p) .
\]
If $\bigcap_{p,i}\ker L_{p,i}\subseteq\ker\ell$, then $\ell$ is a linear combination of the finitely many
$L_{p,i}$ (the map $\psi\mapsto(L_{p,i}(\psi))_{p,i}$ factors $\ell$ through its image in $\C^{P_Ys}$, and a
linear form on a subspace of $\C^{P_Ys}$ extends to $\C^{P_Ys}$), which is the assertion. So let
$\psi\in\mathcal V$ with $L_{p,i}(\psi)=0$ for all $p,i$; we show $\ell(\psi)=0$.

\emph{Step 1: the lift by $v_0$ vanishes.} For a smooth Fock vector $v$ write $F_v=[v,\delta_\iota]$ and
$\Lambda(v)(h)=\int_{[\auU(X)]}\psi(x)\theta(\omega_{\alpha,\beta}(x,h)F_v)\,dx$, linear in $v$. Let
$h\in\auU(Y)(\auA)$ and write $h_f=\gamma h_pk$ with $\gamma\in\auU(Y)(E^+)$, $k\in J_Y(\psrm)$, and put
$y=\gamma^{-1}h_\infty\in\auU(Y)_\infty$ ($\gamma$ taken at the infinite places); then $h=\gamma\cdot
h_p\,k\,y$, the finite and infinite components commuting. The
element $(1,k)$ fixes $F_v$: by Lemma~\ref{pc:level}, $\omega_{\alpha,\beta}(1,k)[v,\delta_\iota]=[v,\Psi']$
with $\Psi'(x)=\beta_\psrm(\det\bar k)\delta_\iota(\bar k^{-1}x)=\delta_\iota(x)$, as $\bar k=1$. The theta
functional is invariant under the rational points and $\beta$ is trivial on them ((W1), (W3) of
Lemma~\ref{ps:weil}), so
\[
\Lambda(v)(h)=\int\psi(x)\theta\bigl(\omega_{\alpha,\beta}(x,h_p)\omega_{\alpha,\beta}(1,k)
\omega_{\alpha,\beta}(1,y)F_v\bigr)dx=\Lambda(\Omega'(y)v)(h_p),
\]
since $\omega_{\alpha,\beta}(1,y)F_v=F_{\Omega'(y)v}$. By hypothesis $\Omega'(y)v_0=\sum_ic_iQ_i$ for some
$c_i\in\C$, and $\Lambda(Q_i)(h_p)=L_{p,i}(\psi)=0$. Hence $\Lambda(v_0)=0$, i.e.\
$\Theta_{\alpha,\beta}(v_0,\delta_\iota)\psi=0$ on $\auU(Y)(\auA)$.

\emph{Step 2: convolution.} The point datum $\delta_\iota$ is sharp for $J_X(\psrm)$
(Definition~\ref{pc:sharp}): the reduction of $J_X(\psrm)$ is trivial, so the first condition is empty,
and the matrix $A_0={}^t\iota$ satisfies $A_0\iota=1$. By Lemma~\ref{psr:coef},
$\kappa(h)=\langle v_0,\Omega(h)v_0\rangle=\langle v_0,v_0\rangle\pspm_u(h)$ is integrable. By
Corollary~\ref{psr:vanish} (with $S=X$, $T=Y$, $P=v_0$, $\Psi=\delta_\iota$, $J=J_X(\psrm)$),
$\kappa*\psi=0$; at $g$ this reads $\langle v_0,v_0\rangle\,\ell(\psi)=0$, and $\langle v_0,v_0\rangle\ne0$.
\end{proof}

\begin{remark}[On the hypotheses of Theorem~\ref{psr:R}]\label{psr:rem}
(i) Theorem~\ref{psr:R} is (R) of Proposition~\ref{psp:inputs} for $Y$ of rank $>2n_X$, the gap
$\ge300\pspn$, and characters of the kernel unramified outside $\psrm=(\pi_M^2)$ and tamely ramified at
$\psrm$ (in the applications they are unramified at every finite place, Theorem~\ref{pss:a}); it makes the ``Moreover'' clause of (R) precise: the $Q_p$ are a basis of a
finite-dimensional space of Fock vectors stable under $\auU(Y)_\infty$ and containing a non-zero
holomorphic vector of weight $u$ in the sense of Definition~\ref{psr:def-hol}, and the index runs over
pairs (class, basis vector). (ii) The proof does not use the gap beyond $n_X^2+n_X+1$ (needed to identify
the coefficient with $\pspm_u$, Lemma~\ref{psr:coef}), and it uses of the level only that $\delta_\iota$
is sharp for $J_X(\psrm)$ and invariant under $J_Y(\psrm)$. (iii) The inner product formula is used only
through Corollary~\ref{psr:vanish}; the constant $A$ is not computed.
\end{remark}



%% ===== proofs/C-span4.tex =====
% proofs/C-span4.tex
% Holomorphic polynomial Fock vectors; the inputs (H5), part of (H4), and the hypotheses (a), (b)
% of Theorem psr:R in its applications. To be \input after C-span.tex, C-span2.tex, C-span3.tex.
% Labels of this file start with pss:. Only amsmath/amssymb/amsthm.
\providecommand{\auU}{\mathrm U}
\providecommand{\auGL}{\mathrm{GL}}
\providecommand{\auHom}{\operatorname{Hom}}
\providecommand{\auA}{\mathbb A}
\providecommand{\aubar}[1]{\overline{#1}}
\providecommand{\pspG}{G_\infty}
\providecommand{\pspm}{\mathsf m}
\providecommand{\pcMat}{\mathrm M}
% macros of this file (prefix \pss)
\providecommand{\pssg}{\mathfrak g}
\providecommand{\pop}{\mathsf p}

\subsection{Holomorphic polynomial Fock vectors}\label{pss:sec}

This subsection supplies the tool needed for (H4), (H5) and the hypotheses of Theorem~\ref{psr:R}: a
criterion for a polynomial Fock vector to be holomorphic in the sense of
Definition~\ref{psr:def-hol}, with its weight (\S\ref{pss:ss-tool}). With it we prove the input
(H5) of Proposition~\ref{psp:inputs} (\S\ref{pss:ss-H5}), the hypothesis on the polynomials $Q_r$ in
the application of Theorem~\ref{psr:R} to $(S,W^\flat)$ (\S\ref{pss:ss-b}), the part of (H4) concerning
the carrier (\S\ref{pss:ss-carrier}), and the hypothesis on the characters in both applications of
Theorem~\ref{psr:R} (\S\ref{pss:ss-a}). Proposition~\ref{pss:status} collects these inputs.

\subsubsection{The tool}\label{pss:ss-tool}

\emph{Setting.} Fix an infinite place $\tau$ (statements for $\prod_\tau$ follow by taking products,
as in \S\ref{psq:ss-H3}). Let $X,T$ be hermitian spaces over $\C$ with orthonormal frames of signs
$s\in\{\pm1\}^{N}$, $s'\in\{\pm1\}^{n'}$, the indices of $X$ numbered by its order $o$ (positive
first, $p$ positive, $q=N-p$ negative); $n_T^\pm$, $n_X^\pm$ count the signs. We use the conventions of
\S\ref{ptf:ss-fock}: a vector of $\auHom(X,T)$ is a matrix $w\in\pcMat_{n'\times N}(\C)$, $(g,h)$ acts
by $w\mapsto hwg^{-1}$, the entry $(r,c)$ has product sign $\varepsilon_{rc}=s'_rs_c$ and Fock variable
$z_{rc}$ ($=\aubar{w_{rc}}$ if $\varepsilon_{rc}=1$, $w_{rc}$ otherwise); a polynomial $P$ in the $z_{rc}$
is identified with its Fock vector. Let $\omega_\infty$ be the archimedean Weil representation with
characters of weights $a$ on $\auU(X)$ and $b$ on $\auU(T)$ (Lemma~\ref{ptf:fock}),
$\Omega(g)=\omega_\infty(g,1)$, $\Omega'(h)=\omega_\infty(1,h)$; $G=\auU(X)$, $\pssg$ its Lie algebra,
$\pssg_\C=\mathfrak{gl}_N(\C)$, $\mathfrak b$ (resp.\ $\mathfrak b^-$) the upper (resp.\ lower)
triangular matrices, $\mathfrak k_\C$ the block diagonal ones, $T_X\subset G$ the diagonal torus, and
similarly for $\auU(T)$ with its order (positive first). For $u\in\Z^N$ and $Z\in\pssg_\C$ put
$u(Z)=\sum_cu_cZ_{cc}$.

We use three standard facts. (KV) The polynomial vectors are smooth vectors for the action of the
semidirect product of the Heisenberg group with $\auU(X)\times\auU(T)$, and $d\Omega$, $d\Omega'$ and
$d\varrho$ map polynomial vectors to polynomial vectors (Howe~\cite[(2.2)--(2.3)]{HoweJAMS}, for the symplectic part: the polynomials are the $K$-finite, hence smooth, vectors, and the derived action is by the operators of (2.2); for the Heisenberg factor, a polynomial vector corresponds in the Schr\"odinger model to a polynomial times a Gaussian, a Schwartz function, and Schwartz functions are smooth vectors of the Schr\"odinger representation: the Heisenberg group acts on $L^2$ by translations, multiplications by characters $e^{2\pi i\langle y,x\rangle}$ and central scalars, so for a Schwartz $f$ the difference quotients of the orbit map at the identity converge in $L^2$, by dominated convergence, to constant multiples of $\partial_jf$, $x_jf$ and $f$, which are again Schwartz; by induction and equivariance the orbit map is $C^\infty$; see also Folland~\cite{FollandPS}). (Hs) On polynomial vectors the derived Heisenberg action is
\eqref{ptq:eq-heis} (or its variant in Remark~\ref{ptq:rem-conv}); in either normalisation, since
$d\varrho(w)-i\,d\varrho(iw)$ and $d\varrho(w)+i\,d\varrho(iw)$ are the complex-linear and antilinear parts,
the complex span of the operators $d\varrho(w)$ is the span $V_{\mathrm{lin}}$ of the operators $z_i$
(multiplication) and $\partial_i=\partial/\partial z_i$, $i=(r,c)$. (Cv) $\Omega(g)\varrho(w)\Omega(g)^{-1}=
\varrho(\iota(g)w)$ with $\iota(g)w=wg^{-1}$ ((W1$_\infty$)); likewise for $\Omega'$.

\begin{lemma}[Quadratic form of the derived action]\label{pss:weyl}
For $Z\in\pssg_\C$ there are an element $D_0$ of the span $\mathcal Q_2$ of the operators $z_iz_j$,
$z_i\partial_j$, $\partial_i\partial_j$ and a constant $c$ such that $d\Omega(Z)=D_0+c$ on polynomials.
The same holds for $d\Omega'$.
\end{lemma}
\begin{proof}
By complex linearity we may take $Z\in\pssg$. Differentiating (Cv) along $e^{sZ}$ on polynomial vectors
(smooth by (KV)) gives $[d\Omega(Z),d\varrho(w)]=d\varrho(\iota(Z)w)$, so by (Hs) $D=d\Omega(Z)$ satisfies
$[D,V_{\mathrm{lin}}]\subseteq V_{\mathrm{lin}}$, and $A_Z=\mathrm{ad}(D)|_{V_{\mathrm{lin}}}$ preserves the
alternating form $\langle\ell,m\rangle=[\ell,m]\in\C$ infinitesimally ($[D,[\ell,m]]=0$ by Jacobi), which is
non-degenerate ($[\partial_i,z_j]=\delta_{ij}$): $A_Z\in\mathfrak{sp}(V_{\mathrm{lin}})$. Let $Q$ be the span of
the symmetrised products $\ell m+m\ell$, $\ell,m\in V_{\mathrm{lin}}$; $Q\subseteq\mathcal Q_2+\C$. We have $\dim Q=k(2k+1)$
($k=\dim V_{\mathrm{lin}}/2$): $Q$ is spanned by the $k(2k+1)$ elements $z_iz_j$ ($i\le j$), $\partial_i\partial_j$
($i\le j$), $\frac12(z_i\partial_j+\partial_jz_i)=z_i\partial_j+\frac12\delta_{ij}$ (all $i,j$), and these are
linearly independent, since a combination $\sum_{i\le j}a_{ij}z_iz_j+\sum_{i\le j}b_{ij}\partial_i\partial_j+
\sum_{i,j}c_{ij}(z_i\partial_j+\frac12\delta_{ij})$ that vanishes on polynomials has, applied to $1$, constant
term $\frac12\sum_ic_{ii}$ and $z_iz_j$-coefficients $a_{ij}$, so $a=0$; applied to $z_l$ it gives
$\sum_ic_{il}z_i+\frac12(\sum_ic_{ii})z_l=0$, whence $c_{il}=0$ for $i\ne l$ and then $c_{ll}+\frac12\sum_ic_{ii}=0$
for all $l$, which forces all $c_{ll}$ equal to some $\gamma$ with $\gamma(1+k/2)=0$ (as $\sum_ic_{ii}=k\gamma$), so $c=0$; applied to
$z_lz_{l'}$ it gives $b_{ll'}(1+\delta_{ll'})=0$, so $b=0$. Moreover $\mathrm{ad}(q)$ maps $V_{\mathrm{lin}}$ into itself
($[\ell m,x]=\ell[m,x]+[\ell,x]m$) and lies in $\mathfrak{sp}(V_{\mathrm{lin}})$. If $\mathrm{ad}(q)=0$, then $q$
commutes with all $z_i$ and $\partial_i$, so $q(P)=Pq(1)$ and $\partial_iq(1)=q(\partial_i1)=0$: $q$ is a scalar;
writing $q=\sum a_{ij}z_iz_j+\sum b_{ij}\partial_i\partial_j+\sum c_{ij}z_i\partial_j+\frac12\sum c_{ii}$ and
applying it to $1$, $z_l$, $z_lz_{l'}$ gives $a=0$, $c=0$, $b=0$, so $q=0$. Hence $\mathrm{ad}\colon Q\to
\mathfrak{sp}(V_{\mathrm{lin}})$ is injective and, by dimensions ($\dim\mathfrak{sp}_{2k}=k(2k+1)$), bijective.
Choose $q$ with $\mathrm{ad}(q)=A_Z$; then $D-q$ commutes with all $z_i$, $\partial_i$ on polynomials,
hence is a scalar (same argument). Take $D_0=q$ minus its constant part.
\end{proof}

\begin{lemma}[Torus weights]\label{pss:torus}
\textup{(i)} For $t\in T_X$, $\Omega(t)P=c(t)\,P(\lambda(t)z)$ with $\lambda_{rc}(t)=t_c^{-\varepsilon_{rc}}$ and a
character $c$; hence $\Omega(t)z_{rc}\Omega(t)^{-1}=t_c^{-\varepsilon_{rc}}z_{rc}$ and
$\Omega(t)\partial_{rc}\Omega(t)^{-1}=t_c^{\varepsilon_{rc}}\partial_{rc}$. We say $z_{rc}$ has \emph{$X$-weight}
$-\varepsilon_{rc}e_c$ and $\partial_{rc}$ has $X$-weight $\varepsilon_{rc}e_c$.
\textup{(ii)} For $a\ne b$, $d\Omega(E_{ab})$ is a linear combination of the products $xy$ of two
operators $x,y\in\{z_i,\partial_i\}$ whose $X$-weights add up to $e_a-e_b$.
\textup{(iii)} The same holds for $\Omega'$ and $t\in T_T$ with \emph{$T$-weights} $\varepsilon_{rc}e_r$ for
$z_{rc}$ and $-\varepsilon_{rc}e_r$ for $\partial_{rc}$.
\end{lemma}
\begin{proof}
(i) is Lemma~\ref{ptf:fock} with $g=t$, $h=1$: variables of positive product sign go to $z\,t^{-\mathsf t}$, i.e.\
$z_{rc}\mapsto t_c^{-1}z_{rc}$, those of negative sign to $zt$; the conjugation formulas follow by the
chain rule. (ii) $\Omega(t)d\Omega(Z)\Omega(t)^{-1}=d\Omega(\mathrm{Ad}(t)Z)$ and $\mathrm{Ad}(t)E_{ab}=t_at_b^{-1}E_{ab}$.
Write $d\Omega(E_{ab})=D_0+c$ as in Lemma~\ref{pss:weyl} in the basis $1$, $z_iz_j$ ($i\le j$), $z_i\partial_j$,
$\partial_i\partial_j$ ($i\le j$) (linearly independent, by the argument at the end of the proof of
Lemma~\ref{pss:weyl}); each basis element is an eigenvector of conjugation by $\Omega(t)$ with the
character given by its weight, by (i). Characters with distinct exponents are distinct, so only basis
elements of weight $e_a-e_b$ occur; the constant has weight $0$. (iii) Lemma~\ref{ptf:fock} with $g=1$,
$h=t$: positive product sign $z\mapsto t^{\mathsf t}z$, negative $z\mapsto t^{-1}z$.
\end{proof}

\begin{lemma}[The compact part]\label{pss:cpt}
For $Y\in\mathfrak k_\C$ and every polynomial $P$,
\[
d\Omega(Y)P=-\sum_{\varepsilon_{rc}=1}\sum_{c'}Y_{cc'}z_{rc'}\partial_{rc}P+\sum_{\varepsilon_{rc}=-1}\sum_{c'}
Y_{c'c}z_{rc'}\partial_{rc}P+\Bigl(\sum_{s_c=1}(n_T^--a)Y_{cc}+\sum_{s_c=-1}(n_T^+-a)Y_{cc}\Bigr)P,
\]
and for $Y\in\mathfrak k'_\C$ (block diagonal for $T$)
\[
d\Omega'(Y)P=\sum_{\varepsilon_{rc}=1}\sum_{r'}Y_{r'r}z_{r'c}\partial_{rc}P-\sum_{\varepsilon_{rc}=-1}\sum_{r'}
Y_{rr'}z_{r'c}\partial_{rc}P+\Bigl(\sum_{s'_r=1}(b-n_X^-)Y_{rr}+\sum_{s'_r=-1}(b-n_X^+)Y_{rr}\Bigr)P .
\]
\end{lemma}
\begin{proof}
For real $Y$ differentiate the formula of Lemma~\ref{ptf:fock} at $g=e^{sY}$ (resp.\ $h=e^{sY}$):
$(z e^{-sY^{\mathsf t}})_{rc}$, $(ze^{sY})_{rc}$, $(e^{sY^{\mathsf t}}z)_{rc}$, $(e^{-sY}z)_{rc}$ have derivatives
$-\sum_{c'}Y_{cc'}z_{rc'}$, $\sum_{c'}Y_{c'c}z_{rc'}$, $\sum_{r'}Y_{r'r}z_{r'c}$, $-\sum_{r'}Y_{rr'}z_{r'c}$, and the
derivatives of the characters $c(g,1)$, $c(1,h)$ of Lemma~\ref{ptf:fock} are the traces shown. Both sides
are complex-linear in $Y$.
\end{proof}

\begin{lemma}[Holomorphy criterion]\label{pss:ext}
Let $\phi$ be a smooth function on $G$ and $u\in\Z^N$. For real $X\in\pssg$ let $R_X\phi(h)=\frac d{ds}\phi(he^{sX})|_{s=0}$,
extended complex-linearly to $\pssg_\C=\pssg\oplus i\pssg$. If $R_Y\phi=u(Y)\phi$ for all $Y\in\mathfrak b$, then
$\phi$ is holomorphic of weight $u$ \textup(\S\ref{psp:ss-inputs}\textup).
\end{lemma}
\begin{proof}
This is the part \emph{Extension} of the proof of Proposition~\ref{psq:mhol}, which uses of the
coefficient only the identity $R_Y\phi=u(Y)\phi$ on $\mathfrak b$. In detail: for $Y\in i\R^N$ diagonal,
$Y\in\pssg\cap\mathfrak b$ and $\phi(he^{Y})=e^{u(Y)}\phi(h)$, so $\phi(ht)=t^u\phi(h)$ on $T_X$; since $G\cap B=T_X$
(Lemma~\ref{psq:GB}), $\tilde\phi(hb)=b^u\phi(h)$ is well defined on the open set $GB$; it is smooth because $(h,b)\mapsto b^u\phi(h)$ is smooth on $G\times B$ and
constant on the fibres of the surjective submersion $G\times B\to GB$, $(h,b)\mapsto hb$ (Lemma~\ref{psq:GB}(i)), which has
smooth local sections. At $x=hb$,
for $Z\in\pssg_\C$ put $Z'=bZb^{-1}=X+Y$ with $X\in\pssg$, $Y\in\mathfrak b$; then
$\frac d{ds}\tilde\phi(xe^{sZ})|_0=b^u(R_X\phi(h)+u(Y)\phi(h))=b^uR_{Z'}\phi(h)$, complex-linear in $Z$.
So $\tilde\phi$ is holomorphic, and $\tilde\phi(hb)=b^u\phi(h)$.
\end{proof}

\begin{corollary}[Coefficients]\label{pss:coefs}
Let $\sigma$ be a unitary representation of $G$ (or of $\pspG$), $\mathcal H^\infty$ its Fr\'echet space of
smooth vectors.
\textup{(a)} If $v$ is holomorphic of weight $u$ and $\Lambda$ is a continuous linear form on
$\mathcal H^\infty$, then $h\mapsto\Lambda(\sigma(h)v)$ is holomorphic of weight $u$. In particular so is
$h\mapsto\langle w,\sigma(h)v\rangle$ for every $w\in\mathcal H$. More generally, the same holds if
$\mathcal E\subset\mathcal H$ is a Fr\'echet space, continuously embedded, such that $h\mapsto\sigma(h)v$ is smooth
into $\mathcal E$ with derivative $h\mapsto\sigma(h)d\sigma(X)v$, and $\Lambda$ is a continuous linear form on $\mathcal E$.
\textup{(b)} If $v\ne0$ is smooth with $d\sigma(Z)v=\lambda(Z)v$ for all $Z\in\mathfrak b^-$, $\lambda\in\Z^N$, then
$h\mapsto\langle\sigma(h)v,v\rangle$ is holomorphic of weight $-\lambda$; if moreover $-\lambda$ is dominant of gap
$\ge N^2+N+1$, then $\langle v,\sigma(h)v\rangle=\langle v,v\rangle\,\aubar{\pspm_{-\lambda}(h)}$.
\end{corollary}
\begin{proof}
$h\mapsto\sigma(h)v\in\mathcal H^\infty$ is smooth (G\aa rding), so $\phi(h)=\Lambda(\sigma(h)v)$ is smooth with
$R_X\phi(h)=\Lambda(\sigma(h)d\sigma(X)v)$, and by complex linearity $R_Y\phi=u(Y)\phi$ on $\mathfrak b$;
Lemma~\ref{pss:ext} gives (a); in the general form the same computation applies verbatim with $\mathcal E$
in place of $\mathcal H^\infty$. (b) Let $g(h)=\langle\sigma(h)v,v\rangle$. For $Y=Y_1+iY_2$, $Y_k\in\pssg$,
$R_Yg=\langle\sigma(h)d\sigma(Y_1-iY_2)v,v\rangle$ (conjugate-linearity in the first variable), and $Y_1-iY_2=
Y^\natural:=-\varepsilon Y^*\varepsilon$, the conjugation of $\pssg_\C$ fixing $\pssg$. For $Y\in\mathfrak b$,
$Y^\natural\in\mathfrak b^-$, so $d\sigma(Y^\natural)v=\lambda(Y^\natural)v=-\sum_c\lambda_c\aubar{Y_{cc}}v$ and
$R_Yg=-\sum_c\lambda_cY_{cc}\,g$. By Lemma~\ref{pss:ext} $g$ is holomorphic of weight $-\lambda$. As $\sigma(t)v=t^\lambda v$
on $T_X$ (diagonal elements, connectedness), $g(\mathrm{rot}(\theta)h)=\langle\sigma(h)v,\sigma(\mathrm{rot}(\theta))^{-1}v\rangle
=e^{2\pi i\theta\langle o,-\lambda\rangle}g(h)$. By Propositions~\ref{psq:unique} and~\ref{psq:mhol} (at each place),
$g=g(1)\pspm_{-\lambda}$, and $\langle v,\sigma(h)v\rangle=\aubar{g(h)}$.
\end{proof}

\begin{proposition}[Holomorphic polynomial vectors]\label{pss:hol}
Let $P$ be a polynomial and $d\in\Z_{\ge0}^N$.
\textup{(i)} Let $T$ be positive definite, $P$ involve only the variables $z_{rc}$ with $s_c=1$, and
satisfy $\sum_rz_{rc'}\partial_{rc}P=0$ for positive $c\ne c'$ with $c$ before $c'$, and $\sum_rz_{rc}\partial_{rc}P=d_cP$.
Then $P$ is holomorphic for $\Omega$ of weight $u$, $u_c=-d_c-a$ for $s_c=1$ and $u_c=n'-a$ for $s_c=-1$.
\textup{(ii)} Let $T$ be negative definite, $P$ involve only the $z_{rc}$ with $s_c=-1$, with
$\sum_rz_{rc'}\partial_{rc}P=0$ for negative $c\ne c'$ with $c$ after $c'$, and $\sum_rz_{rc}\partial_{rc}P=d_cP$. Then
$d\Omega(Z)P=\lambda(Z)P$ for $Z\in\mathfrak b^-$, $\lambda_c=n'-a$ for $s_c=1$, $\lambda_c=-d_c-a$ for $s_c=-1$.
\textup{(iii)} Let $X$ be negative definite of rank $N$, $P$ involve only the $z_{rc}$ with $s'_r=-1$, with
$\sum_cz_{rc}\partial_{r'c}P=0$ for negative rows $r\ne r'$ with $r$ before $r'$, and $\sum_cz_{rc}\partial_{rc}P=d_rP$.
Then $P$ is holomorphic for $\Omega'$ \textup(for the order of $T$\textup) of weight $u'$, $u'_r=b-N$ for
$s'_r=1$, $u'_r=d_r+b$ for $s'_r=-1$.
\end{proposition}
\begin{proof}
$\mathfrak b$ is spanned by the diagonal, the $E_{cc'}$ ($c$ before $c'$, same sign) and the $E_{ac}$
($s_a=1$, $s_c=-1$); $\mathfrak b^-$ likewise with the transposes. (i) Here $\varepsilon_{rc}=s_c$. By
Lemma~\ref{pss:cpt}, for $c,c'$ positive, $d\Omega(E_{cc'})P=-\sum_rz_{rc'}\partial_{rc}P$ (the second sum
involves $\partial$ of variables absent from $P$), which is $0$ for $c$ before $c'$ and $(-d_c-a)P$ after
adding the trace term for $c=c'$ ($n_T^-=0$); for $c,c'$ negative only the trace term survives:
$(n'-a)\delta_{cc'}P$. For $E_{ac}$: by Lemma~\ref{pss:torus}(ii) it is a combination of products of two
operators with $X$-weights adding to $e_a-e_c$; the $X$-weights are $\mp e_a$ for $z_{ra},\partial_{ra}$ and
$\pm e_c$ for $z_{rc},\partial_{rc}$, $\pm e_{c''}$ for other columns; $e_a-e_c$ arises only as $\partial_{ra}\partial_{r'c}$,
and $\partial_{r'c}P=0$. (ii) Here $\varepsilon_{rc}=-s_c$; the variables of $P$ have $\varepsilon=1$ and the
same computation gives the compact part ($n_T^+=0$, $n_T^-=n'$). For $E_{ca}$ ($s_c=-1$, $s_a=1$): $z_{rc}$,
$\partial_{rc}$ have weights $-e_c,e_c$ and $z_{ra},\partial_{ra}$ have $e_a,-e_a$; $e_c-e_a$ arises only as
$\partial_{rc}\partial_{r'a}$, and $\partial_{r'a}P=0$. (iii) Here $\varepsilon_{rc}=-s'_r$; by Lemma~\ref{pss:cpt},
for negative rows $d\Omega'(E_{rr'})P=\sum_cz_{rc}\partial_{r'c}P$, zero for $r$ before $r'$ and with trace term
$b-n_X^+=b$ on the diagonal; on positive rows only the trace $b-n_X^-=b-N$. For $E_{ac}$ ($s'_a=1$,
$s'_c=-1$), Lemma~\ref{pss:torus}(iii): $z_{a\cdot},\partial_{a\cdot}$ have $T$-weights $-e_a,e_a$ and
$z_{c\cdot},\partial_{c\cdot}$ have $e_c,-e_c$; $e_a-e_c$ arises only as $\partial_{a\cdot}\partial_{c\cdot}$, and
$\partial_{a\cdot}P=0$.
\end{proof}

\begin{corollary}[Products of minors]\label{pss:minors}
\textup{(i)} With $T$ positive definite, for matrices
$\eta_k\in\pcMat_{p\times n'}(\C)$, $1\le k\le d$, the polynomial $\prod_k\det(\eta_kz^+)$ \textup($z^+$ the
$n'\times p$ block of positive columns\textup) satisfies (i) with $d_c=d$; it is holomorphic of weight
$(-d-a;\,n'-a)$. \textup{(ii)} With $T$ negative definite of rank $q$, $\det(z^-)^e$ satisfies (ii) with
$d_c=e$. \textup{(iii)} With $X$ negative definite of rank $N$ and $T$ with $N$ negative rows,
$\det(z^{-\cdot})^e$ \textup(the block of negative rows\textup) satisfies (iii) with $d_r=e$.
\end{corollary}
\begin{proof}
$\sum_rz_{rc'}\partial_{rc}$ is a derivation and on $\det(\eta_kz^+)$ replaces the column $c$ of $\eta_kz^+$ by the
column $c'$, giving $0$ for $c\ne c'$ and the determinant for $c=c'$; so it acts on the product by $0$, resp.\
by $d$. (ii), (iii) are the same with columns, resp.\ rows.
\end{proof}

\subsubsection{The factorisation (H5)}\label{pss:ss-H5}

\begin{theorem}[(H5)]\label{pss:H5}
\textup{(H5)} of Proposition~\ref{psp:inputs} holds: with $\Omega$ the real Weil representation of
$\auU(S)_\infty$ on the Fock space of $\auHom(S,\widetilde W)\otimes\R$ with the character
$\nu^\flat\upsilon^m$ of \eqref{pte:eq-F} on $\auU(S)$ \textup(and its restriction $\upsilon^m$ on
$\auU(\widetilde W)$\textup), $c_{\widetilde P}(h',1)=c_{\widetilde P}(1,1)\,\pspm_{u_2}(h')\,\aubar{\pspm_{u_3}(h')}$ for
$u_2=(\Delta';\mu^+_1+r+2m)$, $u_3=(0;\mu^+_1+N''+2m)$.
\end{theorem}
\begin{proof}
Fix $\tau$; write $\mu_1=\mu^+_{1,\tau}$. By the proof of Proposition~\ref{pte:third} the weight of
$\nu^\flat\upsilon^m$ at $\tau$ is $a=m-\mu_1$. $\widetilde W$ has $r+2m$ positive and $m$ negative indices
(\S\ref{pte:ss-third}). By Lemma~\ref{psq:degree} and its proof,
$\widetilde P=P^+P^-$ with $P^+=\prod_{k}\det(\eta_kz^{++})$, a product of $d^+=\mu_1-\Delta'$ minors of the block
$z^{++}$ (positive rows, positive columns), and $P^-=\det(z^{--})^e$, $e=\mu_1+N''+m$.

\emph{Total weight.} By Lemma~\ref{ptf:fock} with $g=t\in T_S$, $h=1$: $\Omega(t)\widetilde P=t^{\mathrm{wt}}\widetilde P$ with
$\mathrm{wt}_c=-d^++(m-a)=\Delta'$ for $c$ positive and $\mathrm{wt}_c=-e+(r+2m-a)=r-N''$ for $c$ negative. Hence
$c(h't)=t^{\mathrm{wt}}c(h')$ for $c(h')=c_{\widetilde P}(h',1)$, and $\mathrm{wt}=u_2-u_3$.

\emph{Splitting.} Let $\widetilde W_\tau=\widetilde W^+\oplus\widetilde W^-$ be the orthogonal decomposition by signs.
The natural map $F_1\otimes F_2\mapsto F_1(w_1)F_2(w_2)$ identifies the Fock space $\mathfrak F$ of
$\auHom(S,\widetilde W_\tau)$ with $\mathfrak F^+\otimes\mathfrak F^-$ unitarily, intertwining the Heisenberg
representation of the sum with the tensor product (the cocycle $\beta$ of \S\ref{psp:ss-fock} and the
Gaussian are multiplicative over orthogonal sums), and carries $\widetilde P$ to $P^+\otimes P^-$. Let
$\Omega^\pm$ be the Weil representations of $\auU(S_\tau)$ on $\mathfrak F^\pm$ of Lemma~\ref{ptf:fock} with
trivial characters. $\Omega(h)$ and $\Omega^+(h)\otimes\Omega^-(h)$ satisfy (Cv) for the same linear map,
so by irreducibility (Stone--von Neumann, Schur) $\Omega=\chi\cdot(\Omega^+\otimes\Omega^-)$ with
$\chi(h)\in\C$, $|\chi|=1$; $\chi$ is a continuous character of the connected group $\auU(S_\tau)$, so its
differential vanishes on the semisimple $\mathfrak{su}(2m,m)$ and $\chi=\det^k$, with $k\in\Z$ by restriction
to $\mathrm{diag}(z,1,\dots,1)$. Hence $c=\chi\,c^+c^-$, $c^\pm(h)=\langle P^\pm,\Omega^\pm(h)P^\pm\rangle$.

\emph{The factors.} By Corollary~\ref{pss:minors}(i) ($T=\widetilde W^+$, $n'=r+2m$, $a=0$) $P^+$ is holomorphic
of weight $w^+=(-d^+;r+2m)$, dominant of gap $r+2m+d^+\ge9m^2+3m+1$ ($r\ge x^5$, $m\le2\sqrt x$); by
Lemma~\ref{psr:coef}, $c^+=\langle P^+,P^+\rangle\pspm_{w^+}$. By Corollary~\ref{pss:minors}(ii)
($T=\widetilde W^-$, $n'=m$, $a=0$) and Proposition~\ref{pss:hol}(ii), $P^-$ satisfies the hypothesis of
Corollary~\ref{pss:coefs}(b) with $\lambda=(m;-e)$, and $-\lambda=(-m;e)$ is dominant of gap $e+m$; so
$c^-=\langle P^-,P^-\rangle\aubar{\pspm_{(-m;e)}}$. Since $\phi_{u+k\mathbf 1}=\det^k\phi_u$ (definition of $\phi_u$
before Lemma~\ref{psq:GB}) and $\pspm=\phi$ (Proposition~\ref{psq:mhol}), with $|\det h|=1$,
$\aubar{\pspm_{(-m;e)}}=\det^m\aubar{\pspm_{u_3}}$ ($u_3=(0;e+m)$) and $\chi\pspm_{w^+}\det^m=\pspm_{w'}$,
$w'=w^++(k+m)\mathbf 1$. Thus $c=c(1)\,\pspm_{w'}\,\aubar{\pspm_{u_3}}$ ($c(1)=\|P^+\|^2\|P^-\|^2$, as all
three functions are $1$ at $1$). Comparing right $T_S$-weights, $t^{\mathrm{wt}}=t^{w'-u_3}$, so $w'=u_3+
\mathrm{wt}=u_2$. Taking the product over $\tau$ gives the claim.
\end{proof}

\noindent In particular $u_3=(0;3m)+u_1+u_2$ (checked in Lemma~\ref{psp:weights}'s notation:
$N''+2m=r+11m+1$), as used in Proposition~\ref{psp:approx}.

\subsubsection{The hypothesis (b) for $(S,W^\flat)$}\label{pss:ss-b}

In Lemma~\ref{psp:step1} the kernel of $(S,W')$ with the characters $(\nu^\flat\upsilon^m,\upsilon^m)$ of
\eqref{pte:eq-F} is written by (W7) as the product of kernels of $(S,\widetilde W)$ and $(S,W^\flat)$. The scalar
$(\nu^\flat\upsilon^m)(\det g)^{-1}$ may be attached to either factor; we attach it to $(S,\widetilde W)$, so that the
kernel of $(S,W^\flat)$ has the characters $(1,\upsilon^m)$, and the kernel of $(S,\widetilde W)$ those used in
Theorem~\ref{pss:H5}.

\begin{theorem}[Hypothesis (b)]\label{pss:b}
Let $I$ be a set of $2m$ indices of $W^\flat$, $v_0=\prod_\tau\det(z^{I,+}_\tau)^{\Delta'}$ on $\auHom(S,W^\flat)\otimes\R$,
and $\mathcal Q$ the span of the $\Omega'(y)v_0$, $y\in\auU(W^\flat)_\infty$. Then $\mathcal Q$ is finite
dimensional, stable under $\auU(W^\flat)_\infty$, and $v_0\in\mathcal Q$ is non-zero and holomorphic of weight
$u_1=(-\Delta';6m+1)$ for $\Omega$ \textup(characters $(1,\upsilon^m)$\textup). Moreover for every invertible
$\mathsf Y\in\pcMat_{6m+1}(E)$, $\mathcal Q$ is contained in the span $\mathcal Q_{\mathsf Y}$ of the polynomials
$\prod_\tau\prod_{k=1}^{\Delta'}\det\bigl((\tau(\mathsf Y)z^+_\tau)_{I_{\tau k}}\bigr)$, $I_{\tau k}$ sets of $2m$ rows.
\end{theorem}
\begin{proof}
Holomorphy and weight: Corollary~\ref{pss:minors}(i) at each $\tau$ ($n'=6m+1$, $a=0$, $d=\Delta'$, $\eta_k$ the
selection of the rows $I$). By Lemma~\ref{ptf:fock}, $\Omega'(y)$ acts on polynomials in $z^+$ by a scalar
times $z^+\mapsto y^{\mathsf t}z^+$, preserving the finite-dimensional space of polynomials in $z^+$ of degree
$2m\Delta'$ at each $\tau$; so $\mathcal Q$ is finite dimensional and stable. Row $i$ of $y_\tau^{\mathsf t}z^+_\tau$ is
$\sum_jM_{ij}(\tau(\mathsf Y)z^+_\tau)_j$ for some $M$ ($\tau(\mathsf Y)$ invertible), and by multilinearity of the determinant in
its rows $\det((y^{\mathsf t}_\tau z^+_\tau)_I)$ is a combination of minors $\det((\tau(\mathsf Y)z^+_\tau)_J)$; taking powers and
the product over $\tau$, $\Omega'(y)v_0\in\mathcal Q_{\mathsf Y}$.
\end{proof}

\begin{corollary}\label{pss:bcor}
Theorem~\ref{psr:R} applies to $(X,Y)=(S,W^\flat)$, $u=u_1$, the characters $(1,\upsilon^m)$ and $\mathcal Q$
\textup(the hypothesis on the characters is Theorem~\ref{pss:a}\textup). Expanding each basis vector of $\mathcal Q$,
at the class $h_p$, in the products of minors of $\mathcal Q_{\mathsf Y_p}$, $\mathsf Y_p$ the matrix of the basis of short vectors
at $h_p$ fixed in \S\ref{pte:ss-third}, the formula takes the form used in Lemma~\ref{psp:step1}, with
polynomials $Q_r$ indexed by pairs \textup(class, row selections\textup).
\end{corollary}
\begin{proof}
The kernel is linear in the polynomial; the gap of $u_1$ is $\ge300n_*$ (Lemma~\ref{psp:weights}).
\end{proof}

\subsubsection{The carrier}\label{pss:ss-carrier}

\begin{proposition}[Part of (H4)]\label{pss:carrier}
The carrier $\varphi^\circ$ is holomorphic of weight $(0;3m)$.
\end{proposition}
\begin{proof}
Fix $x\in\auU(S)(\auA)$. By the definition in \S\ref{pte:ss-carrier},
$\varphi^\circ(xh)=\Lambda_x(\Omega'(h)v)$ for $h\in\auU(S)_\infty$, with $v=\prod_\tau\det(y^-_\tau)^{2m}$ on
$\auHom(I_m^-,S)\otimes\R$ and $\Lambda_x(w)=\int_{[\auU(I_m^-)]}\theta(\omega(y,x)[w\otimes\mathbf 1])\,dy$. We apply Corollary~\ref{pss:coefs}(a) in its general form with
$\mathcal E=\mathcal S$, the Schwartz space of a Schr\"odinger model (as in the proof of Lemma~\ref{psp:prodtheta}).
The map $h\mapsto\Omega'(h)v$ is $C^\infty$ into $\mathcal S$ (with its Fr\'echet topology), with derivative
$\Omega'(h)d\Omega'(X)v$: $\auU(S)_\infty$ acts through a smooth homomorphism into the metaplectic group, the
metaplectic operators act continuously on $\mathcal S$, locally uniformly in the group element, and the derived
action of the symplectic Lie algebra is by differential operators of order $\le2$ with polynomial coefficients of
degree $\le2$, which map $\mathcal S$ continuously to itself (Folland~\cite[Ch.~4, Prop.~4.27, Thm.~4.45 and (4.44)]{FollandPS}; local uniformity: with $\mathsf N$ the number operator (harmonic oscillator), the seminorms $\|(\mathsf N+1)^k\cdot\|_2$ define the topology of $\mathcal S$ and dominate every monomial of degree $\le2k$ in the $x_j$, $\partial_j$ \cite[Appendix to \S V.3, Lemmas~1--2]{ReedSimonI}, and $\Omega'(h)^{-1}(\mathsf N+1)\Omega'(h)$ is a quadratic polynomial in the $x_j,\partial_j$ with coefficients continuous in $h$, by \cite[(4.23)]{FollandPS}); hence for
$w\in\mathcal S$, $\frac d{ds}\Omega'(he^{sX})w=\Omega'(he^{sX})d\Omega'(X)w$ in $\mathcal S$, and induction on the
order of the derivative gives smoothness. The theta functional
$F\mapsto\sum_\xi F(\xi)$ is continuous on $\mathcal S$ (the sum is bounded by a Schwartz seminorm), and
$\omega(y,x)$ acts continuously on $\mathcal S$, uniformly for $(y,x)$ in compact sets; as $[\auU(I_m^-)]$ is compact,
$\Lambda_x$ is continuous on $\mathcal S$. By Corollary~\ref{pss:minors}(iii) ($X=I_m^-$, $N=m$, $e=2m$,
$b=m$ the weight of $\upsilon^m$, as $\upsilon_\tau|_{\auU(1)}(z)=z$) and Proposition~\ref{pss:hol}(iii), $v$ is
holomorphic of weight $(m-m;2m+m)=(0;3m)$; Corollary~\ref{pss:coefs}(a) concludes.
\end{proof}

\subsubsection{The characters (a)}\label{pss:ss-a}

\begin{theorem}[Hypothesis (a)]\label{pss:a}
The characters of the kernels used in the two applications of Theorem~\ref{psr:R} are unramified at every
finite place; in particular unramified outside $(\pi_M^2)$ and tamely ramified there.
\end{theorem}
\begin{proof}
For $(S,W^\flat)$ they are $1$ and $\upsilon^m$ (\S\ref{pss:ss-b}); $\upsilon$ is unramified at every finite place
(proof of Lemma~\ref{pf:scale}). For $(W,V^\circ)$, $K^\circ$ is defined from $\omega^\circ$ (\S\ref{po:ss-glue}),
whose character on $\auU(W)$ is $\alpha^\circ=\upsilon^{m/2-j}$, an integral power of $\upsilon$ (Lemma~\ref{psv:char}),
hence unramified at every finite place.
\end{proof}

\subsubsection{The inputs of Proposition~\ref{psp:inputs}}\label{pss:ss-open}

\begin{proposition}[Inputs of Proposition~\ref{psp:inputs}]\label{pss:status}
\textup{(P)}, \textup{(H1)}--\textup{(H3)} (\S\ref{psq:sec}), \textup{(H5)} (Theorem~\ref{pss:H5}), the carrier part of
\textup{(H4)} (Proposition~\ref{pss:carrier}), and \textup{(R)} in its application to $(S,W^\flat)$
(Theorem~\ref{psr:R} with Theorems~\ref{pss:a}, \ref{pss:b}) are proved, modulo the published results listed
in \S\ref{psr:sec}, \S\ref{psq:sec} and here: Howe~\cite[(2.2)--(2.3)]{HoweJAMS} (smoothness of polynomial vectors and
polynomiality of the derived action), the Stone--von Neumann theorem, G\aa rding's theorem on smooth
vectors. The proofs here do not use the vacuum coefficient of \S\ref{ptq:ss-F} (Lemma~\ref{ptq:vac}).
\end{proposition}

\noindent The rest of (H4), concerning $\aubar y$ and the combinations of the $F_i$, with the bounds on $u,u^\circ$,
is Theorem~\ref{psv:H4}; hypothesis (c) for $(W,V^\circ)$ is Theorem~\ref{psv:c}; the datum of $K^\circ$ is the
point datum of Theorem~\ref{psr:R} (Proposition~\ref{psv:datum}). All inputs of Proposition~\ref{psp:inputs} are
thereby proved (Theorem~\ref{psv:main}).


%% ===== proofs/C-span5.tex =====
% proofs/C-span5.tex
% The remaining inputs of products span: hypothesis (c) and the datum of K° for Theorem psr:R,
% the input (H4) for conj(y) and the F_i with the bounds on the weights, the character split in
% Lemma psp:step1, and the conclusion psp:main = po:span. To be \input after C-span.tex,
% C-span2.tex, C-span3.tex, C-span4.tex. Labels of this file start with psv:.
% Only amsmath/amssymb/amsthm.
\providecommand{\auU}{\mathrm U}
\providecommand{\auGL}{\mathrm{GL}}
\providecommand{\auHom}{\operatorname{Hom}}
\providecommand{\auA}{\mathbb A}
\providecommand{\aubar}[1]{\overline{#1}}
\providecommand{\poA}{\mathcal A}
\providecommand{\pspG}{G_\infty}
\providecommand{\pspm}{\mathsf m}
\providecommand{\pspn}{n_*}
\providecommand{\pcMat}{\mathrm M}
\providecommand{\pop}{\mathsf p}
\providecommand{\pssg}{\mathfrak g}
% macros of this file (prefix \psv)
\providecommand{\psvmu}{\mu}
\providecommand{\pcm}{\mathfrak m}
\providecommand{\auO}{\mathcal O}

\subsection{The last inputs of products span}\label{psv:sec}

This subsection proves the inputs of Proposition~\ref{psp:inputs} not covered by
Proposition~\ref{pss:status}. We fix the order of the frame of
$W=W_j$ (\S\ref{psv:ss-order}), prove a fourth case of Proposition~\ref{pss:hol}
(\S\ref{psv:ss-hol}), and with it the hypothesis (c) of Theorem~\ref{psr:R} for $(W,V^\circ)$
(\S\ref{psv:ss-c}); we show that the datum of $K^\circ$ is the point datum of Theorem~\ref{psr:R}
(\S\ref{psv:ss-datum}); we prove (H4) for $\aubar y$ and for the combinations of the $F_i$, with the
bounds on $u$ and $u^\circ$ (\S\ref{psv:ss-H4}); we settle the character split of
Lemma~\ref{psp:step1} (\S\ref{psv:ss-split}); and we conclude (\S\ref{psv:ss-main}).

Notation is that of \S\ref{psp:sec}, \S\ref{pss:ss-tool} and \S\ref{po:ss-glue}. In particular
$n=r+t$ with $t=5r$, $r=k+m-1$ ($k$ even, $m$ even, so $r$ is odd), $r_+=k-1+j$, $r_-=m-j$,
$N''=r+9m+1$, $n^\circ=t-9m-1$, and $\Delta=2r^8$. At $\tau\in\auPl$ the parameter $\varsigma_\tau$ of
$\poA(W_j)$ and the diagrams $\mu^\pm_\tau$ are those of \S\ref{plf:ss-L4} and
Lemma~\ref{plf:match}. Inner products of Hilbert spaces are conjugate-linear in the first
variable.

\subsubsection{The order of the frame of $W$}\label{psv:ss-order}

The definition of holomorphy in \S\ref{psp:ss-inputs} depends on the order $o$ of the frame, which
\S\ref{psp:ss-notation} only requires to put the positive indices first. For $W=W_j$ we fix it as
follows. Number the frame of $W$ by $0,\dots,r-1$ (the \emph{fixed numbering} of
\S\ref{ss:au-spaces}: $0,\dots,r_+-1$ positive, $r_+,\dots,r-1$ negative); recall that $\varsigma_\tau(a)$
sits on the coordinate $a$ and is strictly decreasing in $a$ (\S\ref{plf:ss-L4}).

\begin{definition}\label{psv:order}
The order $o$ of the frame of $W$ is: first the positive indices in the order $r_+-1,r_+-2,\dots,0$,
then the negative indices in the order $r-1,r-2,\dots,r_+$. Equivalently, positive indices first,
and within each sign block the order of increasing $\varsigma_\tau$ (the same at every $\tau$).
\end{definition}

\noindent Thus, within each sign block, ``$a$ before $b$ in $o$'' means $\varsigma_\tau(a)<\varsigma_\tau(b)$.
The position of a negative index $c$ in $o$ among the negative ones, counted from $0$, is the number
$p(c)$ of Lemma~\ref{plf:match} (negative rows numbered by increasing $\varsigma_\tau$), and the
position of a positive index $a$ in the fixed numbering is the number $p(a)$ there (positive rows
numbered by decreasing $\varsigma_\tau$). The ``reverse numbering of the negative basis vectors'' of
\eqref{po:eq-poly} is the order $o$ on the negative block, and the ``fixed numbering of the positive
ones'' is the reverse of $o$ on the positive block. All weights on $W$ below are written in the fixed
numbering of the indices, and dominance and gaps are taken for the order $o$. Nothing in
\S\ref{psp:sec}, \S\ref{psq:sec}, \S\ref{psr:sec}, \S\ref{pss:sec} depends on the choice of $o$ among
the orders that put the positive indices first, other than through the weights, which we compute
for this $o$.

\begin{lemma}[Two row lengths]\label{psv:rows}
At $\tau\in\auPl$, for $0\le p<r_+$ the row number $p$ \textup(from $0$, longest first\textup) of
$\mu^+_\tau$ has length $\ell^+_p=\varsigma_\tau(p)+\kappa_j(\tau)-\frac{t-1}2-(r_+-p)$; for $0\le p<r_-$ the
row number $p$ of $\mu^-_\tau$ has length $\ell^-_p=-\kappa_j(\tau)-\frac{t-1}2-\varsigma_\tau(r-1-p)-(r_--p)$,
and $\ell^-_p\ge\frac\Delta2-\frac{t+1}2$. Moreover $\ell^-_p\le m\Delta\le b$ and $\max_a|\varsigma_\tau(a)|\le2b$.
\end{lemma}
\begin{proof}
The two formulas are the telescoped sums in the proof of Lemma~\ref{plf:match} (end of the parts
$\varepsilon_a=+1$ and $\varepsilon_a=-1$). In the second, $\varsigma_\tau(r-1-p)=-\kappa(\tau)-a'\Delta$ with
$a'=m-p>j$ and $\kappa_j=\kappa+(j+\frac12)\Delta$, so $\ell^-_p=(a'-j-\frac12)\Delta-\frac{t-1}2-(a'-j)$; with
$1\le a'-j\le m$ this lies between $\frac\Delta2-\frac{t+1}2$ and $m\Delta$. Finally $m\Delta\le(m+1)\Delta\le b$
and $\kappa(\tau)\le b$ (proof of Lemma~\ref{psq:degree}), and the entries of $\varsigma_\tau$ lie between
$-\kappa(\tau)-m\Delta$ and $\frac k2-1\le b$ (\S\ref{plf:ss-L4}).
\end{proof}

\subsubsection{Polynomials in the negative columns}\label{psv:ss-hol}

We add to Proposition~\ref{pss:hol} the case of a positive definite target and a polynomial in the
negative columns of the source. The setting and notation are those of \S\ref{pss:ss-tool} (one place
$\tau$; source $X$ with signs $s$ and order $o$, positive first; target $T$ of rank $n'$; matrices
$w\in\pcMat_{n'\times N}(\C)$; Fock variables $z_{rc}$; $\Omega$ with a character of weight $a$ on
$\auU(X)$).

\begin{proposition}[Holomorphic polynomial vectors, case (iv)]\label{psv:hol4}
Let $T$ be positive definite, $d\in\Z_{\ge0}^N$, and let $P$ be a polynomial involving only the
variables $z_{rc}$ with $s_c=-1$, such that $\sum_rz_{rc}\partial_{rc'}P=0$ for negative $c\ne c'$ with
$c$ before $c'$ in $o$, and $\sum_rz_{rc}\partial_{rc}P=d_cP$ for negative $c$. Then $P$ is holomorphic for
$\Omega$ \textup(Definition~\ref{psr:def-hol}\textup) of weight $u$, $u_c=-a$ for $s_c=1$ and
$u_c=n'-a+d_c$ for $s_c=-1$.
\end{proposition}
\begin{proof}
Here $s'\equiv1$, so $\varepsilon_{rc}=s_c$, $n_T^+=n'$, $n_T^-=0$, and the variables of $P$ have $\varepsilon=-1$.
As in the proof of Proposition~\ref{pss:hol}, $\mathfrak b$ is spanned by the diagonal matrices, the
$E_{cc'}$ with $c$ before $c'$ of the same sign, and the $E_{ac}$ with $s_a=1$, $s_c=-1$.

\emph{Compact part.} Let $Y=E_{cc'}$, $c,c'$ of the same sign. In the formula of
Lemma~\ref{pss:cpt} for $d\Omega(Y)P$ the first sum is $-\sum_{\varepsilon_{rc}=1}z_{rc'}\partial_{rc}P$ (it is
non-zero only when $c$, hence $c'$, is positive), which vanishes because $P$ does not involve the
variables of positive columns; the second sum is $\sum_{\varepsilon_{rc'}=-1}z_{rc}\partial_{rc'}P$ (non-zero only
when $c'$, hence $c$, is negative). The trace term is $(n_T^--a)\delta_{cc'}=-a\,\delta_{cc'}$ for positive
$c$, and $(n_T^+-a)\delta_{cc'}=(n'-a)\delta_{cc'}$ for negative $c$. Hence for $c,c'$ positive
$d\Omega(E_{cc'})P=-a\,\delta_{cc'}P$; for $c,c'$ negative $d\Omega(E_{cc'})P=\sum_rz_{rc}\partial_{rc'}P+(n'-a)
\delta_{cc'}P$, which is $0$ for $c$ before $c'$ by hypothesis and $(d_c+n'-a)P$ for $c=c'$.

\emph{Non-compact part.} Let $Y=E_{ac}$, $s_a=1$, $s_c=-1$. By Lemma~\ref{pss:torus}(ii),
$d\Omega(E_{ac})$ is a linear combination of products $xy$, $x,y\in\{z_i,\partial_i\}$, whose $X$-weights add
up to $e_a-e_c$. By Lemma~\ref{pss:torus}(i) the $X$-weight of $z_{rc''}$ is $-s_{c''}e_{c''}$ and that of
$\partial_{rc''}$ is $s_{c''}e_{c''}$; so the operators of the column $a$ have weights $\mp e_a$ ($z_{ra}$,
$\partial_{ra}$), those of the column $c$ have weights $\pm e_c$ ($z_{rc}$, $\partial_{rc}$), and those of any other
column $c''$ have weights $\pm e_{c''}$ with $c''\ne a,c$. A sum of two such weights equals $e_a-e_c$ only
for $\partial_{ra}\partial_{r'c}$. Since $\partial_{ra}P=0$ ($a$ is positive) and the $\partial$'s commute,
$d\Omega(E_{ac})P=0$. Hence $d\Omega(Y)P=u(Y)P$ for all $Y\in\mathfrak b$.
\end{proof}

\begin{corollary}\label{psv:minors}
Let $T$ be positive definite, and let $c_0,c_1,\dots,c_{q-1}$ be the negative indices of $X$ in the
order $o$. For a finite family $(h_{\mathfrak c})_{\mathfrak c}$ of integers $1\le h_{\mathfrak c}\le q$ and row sets
$R_{\mathfrak c}$ of cardinality $h_{\mathfrak c}$, put $J_{\mathfrak c}=\{c_0,\dots,c_{h_{\mathfrak c}-1}\}$ and
$D=\prod_{\mathfrak c}\det\bigl(z_{R_{\mathfrak c},J_{\mathfrak c}}\bigr)$ \textup(the minor on the rows $R_{\mathfrak c}$ and the
columns $J_{\mathfrak c}$, in the order $o$\textup). Then every linear combination of such products, for a
fixed family $(h_{\mathfrak c})$ and arbitrary row sets, satisfies the hypothesis of
Proposition~\ref{psv:hol4} with $d_{c_p}=\#\{\mathfrak c:h_{\mathfrak c}>p\}$; it is holomorphic of weight
$u_c=-a$ \textup($s_c=1$\textup), $u_{c_p}=n'-a+\#\{\mathfrak c:h_{\mathfrak c}>p\}$.
\end{corollary}
\begin{proof}
For $c\ne c'$ the operator $L=\sum_rz_{rc}\partial_{rc'}$ is a derivation, and on $\det(z_{R,J})$ it replaces
the column $c'$ by the column $c$ if $c'\in J$, and gives $0$ if $c'\notin J$. If $c=c_p$ is before
$c'=c_{p'}$ ($p<p'$) and $c'\in J_{\mathfrak c}$, then $c\in J_{\mathfrak c}$ ($J_{\mathfrak c}$ is an initial segment), and the
resulting determinant has two equal columns. So $L$ kills each factor, hence $D$. For $c=c'=c_p$, $L$
multiplies $\det(z_{R,J_{\mathfrak c}})$ by $1$ if $c_p\in J_{\mathfrak c}$, i.e.\ $h_{\mathfrak c}>p$, and by $0$
otherwise; so $LD=d_{c_p}D$, with $d_{c_p}$ independent of the row sets. Both conditions are linear in
$D$.
\end{proof}

\subsubsection{The hypothesis (c) for $(W,V^\circ)$}\label{psv:ss-c}

By (W7) of \S\ref{po:ss-glue}, $\omega^\circ$ is a Weil representation of the pair $(W,V^\circ)$ twisted by
powers of $\upsilon$ only; in the convention of \S\ref{pc:ss-data} (and \S\ref{pte:ss-conv}),
$\omega^\circ=\omega_{\alpha^\circ,\beta^\circ}$ with $\alpha^\circ,\beta^\circ$ powers of $\upsilon$. We need the
exponent of $\alpha^\circ$.

\begin{lemma}[The character of $\omega^\circ$ on $\auU(W)$]\label{psv:char}
$\alpha^\circ=\upsilon^{m/2-j}$; its archimedean weight is $a^\circ=m/2-j$ at every $\tau$. The character of
$\omega''$ on $\auU(W)$ is $\nu''_\vartheta=\mu^0_j\vartheta\,\upsilon^{j+(t+k-1)/2}$, of weight
$a''_\tau=j+\frac{t+k-1}2-\kappa_j(\tau)$, and $a''_\tau+a^\circ=\frac n2-\kappa_j(\tau)$.
\end{lemma}
\begin{proof}
By \S\ref{po:ss-set}, $\auU(W)$ acts in $\omega_\mu$ through $b(\det g)^{-1}$ times the representation
defined with $\upsilon$, with $b=\mu\,\upsilon^{\lfloor n/2\rfloor}$; that is, $\omega_\mu$ restricted to $\auU(W)$ is
$\omega_{b,\cdot}$ in the convention of \S\ref{pc:ss-data}. By (W7) the characters of the two factors
multiply to $b$ (``the splitting over $\auU(W)$ for $\auHom(W,V)$ is the product of those for the two
summands''), and by the proof of Proposition~\ref{pte:leftside} the character of $\omega''$ on $\auU(W)$ is
that of the kernel $\theta^{\mathsf W,V''}_{\nu''_\vartheta,\,\cdot}$ restricted to $W$, namely $\nu''_\vartheta=\mu^0_j\vartheta
\upsilon^{j+\lfloor(t+k-1)/2\rfloor}$ (\S\ref{pte:ss-third}). As $\mu=\mu^0_j\vartheta$, $\alpha^\circ=b/\nu''_\vartheta=
\upsilon^{\lfloor n/2\rfloor-j-\lfloor(t+k-1)/2\rfloor}$. Here $n=6r$ and $t+k-1=5r+k-1$ are even ($r$ odd, $k$
even), and $3r-j-\frac{5r+k-1}2=\frac{r-k+1}2-j=\frac m2-j$ since $r=k+m-1$. The weights: $\upsilon_\tau(z)=z$ on
$\auU(1)$, $\vartheta$ has weight $0$ and $\mu^0_j$ has weight $-\kappa_j(\tau)$ (proof of Lemma~\ref{plf:const}).
Finally $a''_\tau+a^\circ=\frac{t+k-1+m}2-\kappa_j(\tau)=3r-\kappa_j(\tau)$.
\end{proof}

Let $n^\circ=\operatorname{rk}V^\circ$ (positive definite). For $u^\circ\in W^\circ=W_{\mu^-}(\C^{n^\circ})$ the
polynomial $\pop^\circ_{u^\circ}$ of Lemma~\ref{po:prod}(M3) is
$\pop^\circ_{u^\circ}(x^\circ)=\langle(u^\circ)^\natural,\mathrm{ev}^-(x^{\circ-})^\flat\rangle$, where
$\mathrm{ev}^-(y')=\bigotimes_{\mathfrak c}y'_1\wedge\dots\wedge y'_{h_{\mathfrak c}}$, the product over the columns
$\mathfrak c$ of $\mu^-_\tau$ (all $\tau$), with $y'_i$ the image under $\flat$ of the column number $i$ of
$x^{\circ-}$ in the reverse numbering of the negative basis vectors of $W$, i.e.\ (Definition~\ref{psv:order})
of the negative column $c_{i-1}$ of $W$ in the order $o$.

\begin{lemma}[The coordinates of $K^\circ$]\label{psv:coords}
For every $u^\circ\in W^\circ$, $\pop^\circ_{u^\circ}$ is a linear combination of the products
$\prod_\tau\prod_{\mathfrak c}\det(z_{R_{\mathfrak c},J_{\mathfrak c}})$ of Corollary~\ref{psv:minors} \textup(at each
$\tau$, with $X=W$, $T=V^\circ$, the columns $\mathfrak c$ of $\mu^-_\tau$ of heights $h_{\mathfrak c}$, and row sets
$R_{\mathfrak c}$ of $V^\circ$\textup). The span $\mathcal Q^\circ$ of the $\pop^\circ_{u^\circ}$, $u^\circ\in W^\circ$, is
finite dimensional and stable under $\Omega'(y)$, $y\in\auU(V^\circ)_\infty$, and $\pop^\circ_{u^\circ}\ne0$ for
$u^\circ\ne0$.
\end{lemma}
\begin{proof}
Expanding the wedge products in the basis of $\bigotimes_{\mathfrak c}\Lambda^{h_{\mathfrak c}}(\C^{n^\circ})^\vee$
dual to the basis of pure tensors of basis vectors, the coordinate of $\mathrm{ev}^-(y')$ on the basis vector
of index $(R_{\mathfrak c})$ is $\prod_{\mathfrak c}\det\bigl(y'_i(e_l)\bigr)_{l\in R_{\mathfrak c},\,i\le h_{\mathfrak c}}$, and
$y'_i(e_l)$ is the entry of $x^\circ$ in the row $l$ and the column $c_{i-1}$, or its conjugate, according to
the convention for $\flat$. By \eqref{po:eq-poly} and the remark following it, $\pop^\circ_{u^\circ}$ is a
polynomial in the Fock variables, which for the entries of negative product sign (rows of $V^\circ$,
negative columns of $W$) are the entries themselves (\S\ref{ptf:ss-fock}); this fixes the conjugations,
as in the computation of $\pop_{\xi_0}$ in \S\ref{plf:ss-L4}, where the negative factor is displayed as a
product of minors $\det(z_{a\,c(a')})$ in the Fock variables. Pairing with $(u^\circ)^\natural$ gives
$\pop^\circ_{u^\circ}=\sum_{(R_{\mathfrak c})}\gamma_{(R_{\mathfrak c})}(u^\circ)\prod_{\mathfrak c}\det(z_{R_{\mathfrak c},J_{\mathfrak c}})$ with
$J_{\mathfrak c}=\{c_0,\dots,c_{h_{\mathfrak c}-1}\}$ and scalars $\gamma$. This is the first assertion. $\mathcal Q^\circ$ lies
in the finite-dimensional space of polynomials of bounded degree in finitely many variables. The map
$u^\circ\mapsto\pop^\circ_{u^\circ}$ is conjugate-linear and $\auU(V^\circ)(\R)$-equivariant (\S\ref{po:ss-set}:
$\xi\mapsto\pop_\xi$ is $\auU(V)(\R)$-equivariant, and by (W7) and Lemma~\ref{po:prod}(M1),(M3) the block
$V^\circ$ acts on the factor $\pop^\circ$), so $\mathcal Q^\circ$ is stable under $\Omega'$. Finally, the vectors
$\mathrm{ev}^-(y')$ span $W^\circ$ (definition of $W^\circ$ as in \eqref{po:eq-ev}), so for $u^\circ\ne0$ some
$\langle(u^\circ)^\natural,\mathrm{ev}^-(y')\rangle\ne0$; as $x^{\circ-}\mapsto(y'_i)$ is onto, $\pop^\circ_{u^\circ}\ne0$.
\end{proof}

\begin{theorem}[Hypothesis (c)]\label{psv:c}
Let $u^\circ$ be the weight on $W$ given at $\tau$, in the fixed numbering, by
\[
u^\circ_{\tau,a}=-a^\circ=j-\tfrac m2\quad(0\le a<r_+),\qquad u^\circ_{\tau,c}=n^\circ-a^\circ+\ell^-_{p(c)}
\quad(r_+\le c<r),
\]
with $p(c)=r-1-c$ the position of $c$ in $o$ among the negative indices and $\ell^-_p$ the row lengths of
$\mu^-_\tau$ \textup(Lemma~\ref{psv:rows}\textup). Then every non-zero element of $\mathcal Q^\circ$ is a non-zero
holomorphic vector of weight $u^\circ$ for $\Omega=\omega^\circ_\infty(\cdot,1)$
\textup(Definition~\ref{psr:def-hol}\textup). Consequently the hypothesis on $\mathcal Q$ in
Theorem~\ref{psr:R} holds for $(X,Y)=(W,V^\circ)$, the weight $u^\circ$, the characters of $\omega^\circ$, and
$\mathcal Q=\mathcal Q^\circ$, the span of the coordinates of $K^\circ$ against a basis of $W^\circ$. Moreover, for every
$h^\circ\in\auU(V^\circ)(\auA_f)$, every coordinate of $g\mapsto K^\circ(g,h^\circ)$ is holomorphic of weight
$u^\circ$ on $\auU(W)(\auA)$.
\end{theorem}
\begin{proof}
At $\tau$, by Lemma~\ref{psv:coords} every element of $\mathcal Q^\circ$ is a linear combination of products of
minors as in Corollary~\ref{psv:minors} with $X=W$ (order $o$), $T=V^\circ$, $n'=n^\circ$, $a=a^\circ$
(Lemma~\ref{psv:char}) and the column heights of $\mu^-_\tau$; there $d_{c_p}=\#\{\mathfrak c:h_{\mathfrak c}>p\}$ is the
length $\ell^-_p$ of the row number $p$ of $\mu^-_\tau$, and $c_p$ is the negative index $c$ with $p(c)=p$,
i.e.\ $c=r-1-p$. So Corollary~\ref{psv:minors} gives the $\mathfrak b$-condition at $\tau$ with the weight
$u^\circ_\tau$; at the other places the polynomial is a product of factors, and $d\Omega$ of an element of the
Lie algebra at $\tau$ acts on the factor at $\tau$ only (\S\ref{pss:ss-tool}, statements for $\prod_\tau$ by
taking products), so the condition holds for $\prod_\tau\mathfrak b$. Polynomial vectors are smooth (Howe~\cite[(2.2)--(2.3)]{HoweJAMS},
\S\ref{pss:ss-tool}). Non-vanishing and $\Omega'$-stability are Lemma~\ref{psv:coords}.

For the last assertion fix $g\in\auU(W)(\auA)$ and $h^\circ$. By the definition in \S\ref{po:ss-glue}, for
$x\in\auU(W)_\infty$, $\langle u^\circ,K^\circ(gx,h^\circ)\rangle=\theta^\circ\bigl(\omega^\circ(g,h^\circ)(\Omega(x)
\pop^\circ_{u^\circ}\otimes\Psi^\circ)\bigr)=\Lambda(\Omega(x)\pop^\circ_{u^\circ})$ with $\Lambda(w)=\theta^\circ(\omega^\circ(g,h^\circ)
(w\otimes\Psi^\circ))$. As in the proof of Proposition~\ref{pss:carrier}, $\Lambda$ is a continuous linear form on
the smooth vectors (smooth vectors are Schwartz functions in a Schr\"odinger model,
Lemma~\ref{psp:prodtheta}, on which $\omega^\circ(g,h^\circ)$ and the theta functional are continuous).
Corollary~\ref{pss:coefs}(a) shows that $x\mapsto\langle u^\circ,K^\circ(gx,h^\circ)\rangle$ is holomorphic of weight
$u^\circ$. Taking $u^\circ$ in a basis of $W^\circ$ gives the coordinates of $K^\circ(\cdot,h^\circ)$ (up to complex
conjugation of the coefficients of the basis change, which does not affect holomorphy).
\end{proof}

\noindent The weights $u^\circ_\tau$ are determined by the shift $j-\frac m2$, the rank $n^\circ$ of $V^\circ$, and the
row lengths of $\mu^-$ on the negative indices.

\subsubsection{The datum of $K^\circ$}\label{psv:ss-datum}

\begin{proposition}[The datum of $K^\circ$ is the point datum]\label{psv:datum}
Let $\pcm=(\pi_M^2)$ and let $\iota\in\pcMat_{n^\circ\times r}(\auO_E/\pcm)$ be the matrix of the inclusion of the
frame of $W$ into that of $V^\circ$ as the marked coordinates \textup(the first $r$ coordinates of $V^\circ$,
\S\ref{po:ss-glue}\textup), i.e.\ $\iota=\iota_z\bmod\pcm$. Then the finite part $\Psi^\circ=\mathbf 1[v\equiv\iota_z
\bmod\pi_M^2]\otimes(\text{vacuum off }\mathsf M)$ of the datum of $K^\circ$ is the finite part
$\delta_\iota^\natural\otimes\phi^{\mathrm{vac},\pcm}$ of the datum $[\,\cdot\,,\delta_\iota]$ of \S\ref{pc:ss-data}.
Hence for $u^\circ\in W^\circ$, $g\in\auU(W)(\auA)$ and $h^\circ\in\auU(V^\circ)(\auA_f)$,
\[
\langle u^\circ,K^\circ(g,h^\circ)\rangle=\theta_{\alpha^\circ,\beta^\circ}\bigl(g,h^\circ;\pop^\circ_{u^\circ},\delta_\iota\bigr)
\]
in the notation of \S\ref{pc:ss-data}, and Theorem~\ref{psr:R} applies to $(X,Y)=(W,V^\circ)$ with the datum of
$K^\circ$ as it stands: no variant of Theorem~\ref{psr:R} is needed.
\end{proposition}
\begin{proof}
\emph{The modulus.} $\pi_M$ is a uniformiser at every $w\in\mathsf M$ and a unit at $\aubar w$ and at all other
primes (\S\ref{po:ss-glue}), so $\pcm=\prod_{w\in\mathsf M}w^2$ is a modulus in the sense of \S\ref{pc:ss-data}
(prime to $\aubar\pcm$), $S(\pcm)$ is the set of places of $E^+$ below $\mathsf M$ (non-empty), and
$\auO_E/\pcm=\prod_{w\in\mathsf M}\auO_w/w^2$.

\emph{The point datum.} By \S\ref{pc:ss-data}, $\delta_\iota$ is the indicator of the single point $\iota$ of
$\pcMat_{n^\circ\times r}(\auO_E/\pcm)$, and $\delta_\iota^\natural$ is the function on $\prod_{w\in\mathsf M}
\pcMat_{n^\circ\times r}(E_w)$ equal to $\delta_\iota(x\bmod\pcm)$ at integral $x$ and $0$ elsewhere, that is,
$\delta_\iota^\natural(x)=\mathbf 1[x\text{ integral and }x\equiv\iota_z\bmod\pi_M^2]$, which is the factor
$\mathbf 1[v\equiv\iota_z\bmod\pi_M^2]$ of $\Psi^\circ$ (both are functions of the $E_w$-coordinates, $w\in\mathsf M$,
in the model (W5) at the split places of $S(\pcm)$, with matrices written in the frames; the normalised
datum of \S\ref{po:ss-glue} is a function of integral matrices $v\in\pcMat_{n^\circ\times r}(\auO_{\mathsf M})$). Off
$S(\pcm)$ both data are the normalised vacuum of the same lattice. With the convention of
Definition~\ref{pgx:conv}, $\delta_\iota^\natural(x)=\delta_\iota(a^{-1}x\bmod\pcm)$ on $a\,\pcMat(\widehat\auO)$; at the
primes of $\mathsf M$ the element $\alpha_w$ is a unit (the primes where it is not lie above $q_1$,
Proposition~\ref{pgv:alpha}(a), and the chosen places do not), so $a_w=1$ and the two conventions agree there.

\emph{The kernel.} By Lemma~\ref{psv:char} and (W7), $\omega^\circ=\omega_{\alpha^\circ,\beta^\circ}$ and $\theta^\circ$
is the theta functional; so $\theta^\circ(g,h^\circ;u^\circ)=\theta(\omega_{\alpha^\circ,\beta^\circ}(g,h^\circ)[\pop^\circ_{u^\circ},
\delta_\iota])$, which is the displayed identity.

\emph{Where the datum enters the proof of Theorem~\ref{psr:R}.} It enters in exactly two places
(Remark~\ref{psr:rem}(ii)): Step~1 uses that $\omega_{\alpha,\beta}(1,k)$ fixes $[v,\delta_\iota]$ for $k\in J_Y(\pcm)$
(Lemma~\ref{pc:level}: $\bar k=1$); Step~2 uses that $\delta_\iota$ is sharp for $J_X(\pcm)$ (with $A_0={}^t\iota$,
$A_0\iota=1$), which is also what makes Lemma~\ref{pc:clean} applicable in Proposition~\ref{psr:rallis}. Both
are properties of $\delta_\iota$ as a function on $\pcMat_{n^\circ\times r}(\auO_E/\pcm)$, and we have just seen that
the datum of $K^\circ$ is $\delta_\iota$ itself. The characters $\alpha^\circ,\beta^\circ$ are powers of $\upsilon$, unramified at
every finite place (Theorem~\ref{pss:a}), and $V^\circ$ is positive definite of rank $n^\circ=5r-9m-1>2r$ (as
$r\ge x^5$ and $m\le2\sqrt x$).
\end{proof}

\subsubsection{The input (H4)}\label{psv:ss-H4}

\emph{The conjugate of $y$.} Recall (\S\ref{plf:ss-L4}, Definition~\ref{pz:def-type}) that $\poA(W_j)$ consists
of the bounded continuous functions $y$ on $\auU(W)(\auA)$, left invariant under $\auU(W)(E^+)$, of level $J_M$,
with $R(e_W)y=y$, where $e_W=\bigotimes_\tau d_\tau\aubar{c_{W,\tau}}$, $c_{W,\tau}(x)=\langle v,\pi(x)v\rangle$, $\pi=\pi_{\varsigma_\tau}$
the discrete series of Harish-Chandra parameter $\varsigma_\tau$ and $v$ a unit highest weight vector of its lowest
$K$-type. Let $\Lambda^K(\varsigma_\tau)=\varsigma_\tau+\rho_{\mathrm n}-\rho_{\mathrm c}$ be the highest weight of that $K$-type,
$\rho_{\mathrm c},\rho_{\mathrm n}$ the half sums of the compact and non-compact roots positive for $\varsigma_\tau$.

\begin{lemma}[The vector $v$]\label{psv:v}
At $\tau$, let $\lambda=\Lambda^K(\varsigma_\tau)$ and let $\mathfrak b^-$ be the lower triangular matrices for the order
$o$ of Definition~\ref{psv:order}. Then $d\pi(Z)v=\lambda(Z)v$ for $Z\in\mathfrak b^-$, and in the fixed numbering
\[
\lambda_a=\varsigma_\tau(a)+\tfrac{r_-}2-\tfrac{r_+-1}2+a\quad(0\le a<r_+),\qquad
\lambda_c=\varsigma_\tau(c)-\tfrac{r_+}2-\tfrac{r_--1}2+(c-r_+)\quad(r_+\le c<r);
\]
these are integers.
\end{lemma}
\begin{proof}
The positive roots for $\varsigma_\tau$ are the $e_a-e_b$ with $\varsigma_\tau(a)>\varsigma_\tau(b)$: the compact ones join
indices of one sign block with $a<b$ in the fixed numbering, and every non-compact one is $e_a-e_c$
with $a$ positive and $c$ negative, since every entry of $\varsigma_\tau$ on a positive coordinate exceeds every
entry on a negative one (\S\ref{plf:ss-L4}). So, in the fixed numbering, $\rho_{\mathrm c,a}=\frac{r_+-1}2-a$,
$\rho_{\mathrm n,a}=\frac{r_-}2$ for positive $a$, and $\rho_{\mathrm c,c}=\frac{r_--1}2-(c-r_+)$, $\rho_{\mathrm n,c}=-\frac{r_+}2$
for negative $c$, which gives the formula; $\frac{r_-}2-\frac{r_+-1}2=\frac{r-2r_++1}2$ and $-\frac{r_+}2-\frac{r_--1}2=
-\frac{r-1}2$ are integers as $r$ is odd, and the entries of $\varsigma_\tau$ are integers.

The diagonal part of $\mathfrak b^-$ acts on $v$ by its weight $\lambda$. A compact element $E_{ba}$ of
$\mathfrak b^-$ has $a$ before $b$ in $o$ in one sign block, i.e.\ $\varsigma_\tau(b)>\varsigma_\tau(a)$
(Definition~\ref{psv:order}); so $e_b-e_a$ is a positive compact root and $d\pi(E_{ba})v=0$, $v$ being a highest
weight vector. A non-compact element of $\mathfrak b^-$ is $E_{ca}$ with $c$ negative and $a$ positive. If
$d\pi(E_{ca})v\ne0$, it is a vector of $T$-weight $\lambda+e_c-e_a$ in some $K$-type of $\pi$, of highest weight
$\mu'$, so $\lambda+e_c-e_a=\mu'-\sum_\alpha m_\alpha\alpha$ with $m_\alpha\ge0$ over positive compact roots, and by
(D4) of \S\ref{ptq:ss-H} $\mu'=\lambda+\sum_\beta n_\beta\beta$, $n_\beta\ge0$, over the positive non-compact roots
$\beta=e_{a'}-e_{c'}$. The linear form $H(x)=\sum_{c'\text{ negative}}x_{c'}$ vanishes on compact roots, is $-1$ on
each $\beta$ and $+1$ on $e_c-e_a$; so $1=-\sum n_\beta\le0$, a contradiction.
\end{proof}

\begin{theorem}[(H4) for $\aubar y$]\label{psv:ybar}
Let $u$ be the weight on $W$ given at $\tau$ by $u_\tau=-\Lambda^K(\varsigma_\tau)$, i.e.\ in the fixed numbering
\[
u_{\tau,a}=-\varsigma_\tau(a)-\tfrac{r-2r_++1}2-a\quad(0\le a<r_+),\qquad
u_{\tau,c}=-\varsigma_\tau(c)+\tfrac{r-1}2-(c-r_+)\quad(r_+\le c<r).
\]
For every form $y$ of type $e_W$ \textup(in particular every $y\in e_0\poA(W_j)$\textup), $\aubar y$ is holomorphic of
weight $u$ on $\auU(W)(\auA)$ for the order $o$.
\end{theorem}
\begin{proof}
Fix $g\in\auU(W)(\auA)$ and put $f(h)=\aubar y(gh)$ for $h\in\pspG=\auU(W)_\infty$. Write $c=c_W=\bigotimes_\tau c_{W,\tau}$,
$d=\prod_\tau d_\tau$, and $\kappa(z)=c(z^{-1})=\langle\pi(z)v,v\rangle$ ($\pi=\bigotimes_\tau\pi_{\varsigma_\tau}$, $v=\bigotimes_\tau v_\tau$).
Conjugating $y=R(e_W)y$ and substituting $x\mapsto h^{-1}x$ (Haar measure is invariant),
\[
f(h)=d\int_{\pspG}c(x)\,\aubar y(ghx)\,dx=d\int_{\pspG}c(h^{-1}x)\,\aubar y(gx)\,dx=d\int_{\pspG}\kappa(x^{-1}h)\,\aubar y(gx)\,dx .
\]
\emph{Smoothness.} By Lemma~\ref{plf:A12} the conclusion of Theorem~\ref{par:A1} holds for $c$: it is smooth and
$\mathsf Xc$ is integrable for every left- or right-invariant differential operator $\mathsf X$; as inversion
preserves Haar measure and exchanges left and right invariant operators, the same holds for $\kappa$. For real
$X\in\pssg$ and $R_X$ the right derivative, $\kappa(x^{-1}he^{sX})-\kappa(x^{-1}h)=\int_0^s(R_X\kappa)(x^{-1}he^{\sigma X})\,
d\sigma$; since $\aubar y$ is bounded and $R_X\kappa$ is integrable, Fubini's theorem and the continuity of right
translation in $L^1(\pspG)$ give that $f$ is differentiable along $X$ with
$R_Xf(h)=d\int(R_X\kappa)(x^{-1}h)\,\aubar y(gx)\,dx$, a continuous function of $h$. Iterating, $f$ is smooth, and by
complex linearity the same formula holds for $R_Y$, $Y\in\pssg_\C$.

\emph{The $\mathfrak b$-condition.} For real $X$, $(R_X\kappa)(z)=\frac d{ds}\langle\pi(ze^{sX})v,v\rangle|_0=\langle\pi(z)d\pi(X)v,v\rangle$.
For $Y=Y_1+iY_2$, $Y_1,Y_2\in\pssg$, since the inner product is conjugate-linear in the first variable,
$(R_Y\kappa)(z)=\langle\pi(z)d\pi(Y_1)v,v\rangle+i\langle\pi(z)d\pi(Y_2)v,v\rangle=\langle\pi(z)d\pi(Y^\natural)v,v\rangle$ with
$Y^\natural=Y_1-iY_2=-\varepsilon Y^*\varepsilon$, the conjugation of $\pssg_\C$ fixing $\pssg$ (proof of
Corollary~\ref{pss:coefs}). For $Y\in\prod_\tau\mathfrak b$, $Y^\natural$ is lower triangular, so by Lemma~\ref{psv:v}
$d\pi(Y^\natural)v=\lambda(Y^\natural)v$ with $\lambda(Y^\natural)=\sum\lambda_{\tau,a}(Y^\natural)_{\tau,aa}=-\sum\lambda_{\tau,a}\aubar{Y_{\tau,aa}}$.
Hence $(R_Y\kappa)(z)=\aubar{\lambda(Y^\natural)}\kappa(z)=(-\lambda)(Y)\kappa(z)=u(Y)\kappa(z)$, and by the formula for $R_Yf$,
$R_Yf=u(Y)f$ for all $Y\in\prod_\tau\mathfrak b$. Lemma~\ref{pss:ext} shows that $f$ is holomorphic of weight $u$. As $g$ is
arbitrary, $\aubar y$ is holomorphic of weight $u$.
\end{proof}

\emph{The functions $F_i$.} By Proposition~\ref{pte:third} and Proposition~\ref{pte:leftside} for the trivial twist
$\vartheta=1$, for every choice $i$ there is $c_i^{\mathrm F}\in\C$ (namely $c_3(1)\gamma(1)/\mathsf N$) with
\begin{equation}\label{psv:eq-F}
F_i=c_i^{\mathrm F}\,\Xi_{1,i},\qquad\Xi_{1,i}(g)=\sum_{h}\theta''_{1,\Lambda_i}\bigl(g,h;\varphi''_{1,i}(h)\bigr),
\end{equation}
a finite sum over the classes $h$ of $\auU(V'')$ modulo $J'_{1,i}$ (Definition~\ref{pte:def-ff}), where
$\theta''_{1,\Lambda_i}(g,h;u'')=\theta''(\omega''(g,h)(\pop''_{u''}\otimes\Psi''_{1,\Lambda_i}))$ for $u''\in W''=W_{\mu^+}$
(\S\ref{po:ss-glue}).

\begin{theorem}[(H4) for the $F_i$]\label{psv:F}
Let $u''$ be the weight on $W$ given at $\tau$ by
\[
u''_{\tau,a}=-\ell^+_a-a''_\tau\quad(0\le a<r_+),\qquad u''_{\tau,c}=N''-a''_\tau\quad(r_+\le c<r),
\]
with $\ell^+_a$ the row lengths of $\mu^+_\tau$ \textup(Lemma~\ref{psv:rows}\textup) and $a''_\tau$ as in
Lemma~\ref{psv:char}. Then every linear combination $F=\sum_ic_iF_i$ is holomorphic of weight $u''$ on
$\auU(W)(\auA)$ for the order $o$, and $u''=u-u^\circ$, with $u$ of Theorem~\ref{psv:ybar} and $u^\circ$ of
Theorem~\ref{psv:c}.
\end{theorem}
\begin{proof}
\emph{Holomorphy.} By \eqref{psv:eq-F} it suffices that $g\mapsto\theta''_{1,\Lambda_i}(g,h;u'')$ is holomorphic of weight
$u''$ for each $h$ and $u''\in W''$ (sums of holomorphic functions of one weight are holomorphic of that weight:
take the sum of the extensions on the intersection of their domains). By the proof of Lemma~\ref{ptf:laplace},
$\pop''_{u''}$ is a linear combination of the products $\prod_\tau\prod_{\mathfrak c}\det X_{\tau;R_{\mathfrak c},[h_{\mathfrak c}]}$ over
the columns $\mathfrak c$ of $\mu^+_\tau$, where $X_\tau$ are the Fock variables of the positive columns of $W$ (rows of
$V''$, product sign $+1$) and $[h]$ is the set of the first $h$ positive columns in the fixed numbering, i.e.\ the
\emph{last} $h$ positive columns in the order $o$. At $\tau$ apply Proposition~\ref{pss:hol}(i) with $X=W$, $T=V''$
(positive definite, $n'=N''$) and $a=a''_\tau$ (Lemma~\ref{psv:char}): $P$ involves only variables of positive columns;
for positive $c\ne c'$ with $c$ before $c'$ in $o$, the derivation $\sum_rz_{rc'}\partial_{rc}$ replaces the column $c$ by
the column $c'$ in each minor containing $c$, and then $c'\in[h_{\mathfrak c}]$ as well ($[h_{\mathfrak c}]$ is a final segment of
$o$), so it kills each product; and $\sum_rz_{rc}\partial_{rc}$ multiplies each product by $d_c=\#\{\mathfrak c:c<h_{\mathfrak c}\}$
($c$ in the fixed numbering), which is the length $\ell^+_c$ of the row number $c$ of $\mu^+_\tau$. So $\pop''_{u''}$ is a
holomorphic vector of weight $u''$ for $\Omega''=\omega''_\infty(\cdot,1)$ (at each place, hence for the product, as
in the proof of Theorem~\ref{psv:c}). For $x\in\auU(W)_\infty$, $\theta''_{1,\Lambda_i}(gx,h;u'')=\Lambda(\Omega''(x)\pop''_{u''})$
with $\Lambda(w)=\theta''(\omega''(g,h)(w\otimes\Psi''_{1,\Lambda_i}))$ continuous on smooth vectors (as in the proof of
Theorem~\ref{psv:c}), and Corollary~\ref{pss:coefs}(a) gives the holomorphy.

\emph{The identity $u''+u^\circ=u$.} Positive $a$: by Lemma~\ref{psv:rows} and Lemma~\ref{psv:char},
$u''_a+u^\circ_a=-\ell^+_a-(a''+a^\circ)=-\varsigma_\tau(a)-\kappa_j+\frac{t-1}2+(r_+-a)-\frac n2+\kappa_j=-\varsigma_\tau(a)-\frac{r+1}2+r_+-a$,
which is $u_a$. Negative $c$, with $p=p(c)=r-1-c$: $u''_c+u^\circ_c=N''+n^\circ-(a''+a^\circ)+\ell^-_p=n-\frac n2+\kappa_j+\ell^-_p$,
and $\ell^-_p=-\kappa_j-\frac{t-1}2-\varsigma_\tau(c)-(r_--p)$ with $r_--p=c-r_++1$; so the sum is
$\frac n2-\frac{t-1}2-\varsigma_\tau(c)-(c-r_+)-1=-\varsigma_\tau(c)+\frac{r-1}2-(c-r_+)=u_c$ (using $N''+n^\circ=n=r+t$).
\end{proof}

\begin{proposition}[Bounds on $u$ and $u^\circ$]\label{psv:bounds}
$u$ and $u^\circ$ are dominant for $o$ with gap $\ge300\pspn$, and $|u_{\tau,a}|,|u^\circ_{\tau,a}|\le5b$.
\end{proposition}
\begin{proof}
\emph{$u$.} Along $o$ in the positive block the fixed index $a$ decreases; $u_a=-\varsigma_\tau(a)-a+\mathrm{const}$ and
$\varsigma_\tau(a)-\varsigma_\tau(a+1)\ge1$ (integers, strictly decreasing) give $u_{a+1}\ge u_a$, so $u$ is non-increasing along
$o$; the negative block is the same. Gap: for $a$ positive and $c$ negative, $u_c-u_a=\varsigma_\tau(a)-\varsigma_\tau(c)+
\frac{r-1}2+\frac{r-2r_++1}2+a-(c-r_+)=\varsigma_\tau(a)-\varsigma_\tau(c)+r_--(c-r_+)+a\ge\Delta+1$, since entries on coordinates of
different signs differ by at least $\Delta$ (\S\ref{plf:ss-L4}) and $c-r_+\le r_--1$; and $\Delta\ge300\pspn$ by
Lemma~\ref{psp:sizes}(ii). Entries: $|u_{\tau,a}|\le\max|\varsigma_\tau|+r\le2b+b$ by Lemma~\ref{psv:rows} and $r\le t\le\pspn\le b$
(Lemma~\ref{psp:sizes}(i)).

\emph{$u^\circ$.} It is constant on the positive block. Along $o$ in the negative block $p(c)$ increases and
$\ell^-_p$ is non-increasing in $p$ (rows of a diagram, longest first), so $u^\circ$ is non-increasing. Gap:
$u^\circ_c-u^\circ_a=n^\circ+\ell^-_{p(c)}\ge\ell^-_{p(c)}\ge\frac\Delta2-\frac{t+1}2\ge300\pspn$ by Lemma~\ref{psv:rows} and
Lemma~\ref{psp:sizes}(ii) ($300\pspn+t\le\Delta/2$). Entries: $|a^\circ|\le m/2$, $n^\circ\le t\le b$ and $\ell^-_p\le b$
(Lemma~\ref{psv:rows}), so $|u^\circ_{\tau,a}|\le m+2b\le3b$.
\end{proof}

\begin{theorem}[(H4)]\label{psv:H4}
The input \textup{(H4)} of Proposition~\ref{psp:inputs} holds, for the order $o$ of Definition~\ref{psv:order} on
$W$ and the standard order on $S$: for $y\in e_0\poA(W_j)$, $\aubar y$ is holomorphic of the weight $u$ of
Theorem~\ref{psv:ybar}; every combination $F$ of the $F_i$ is holomorphic of weight $u-u^\circ$, where $u^\circ$
\textup(Theorem~\ref{psv:c}\textup) is the weight of $g\mapsto K^\circ(g,h^\circ)$; $u,u^\circ$ are dominant with gap
$\ge300\pspn$ and entries of modulus $\le5b$; and the carrier is holomorphic of weight $(0;3m)$.
\end{theorem}
\begin{proof}
Theorems~\ref{psv:ybar}, \ref{psv:F}, \ref{psv:c}, Proposition~\ref{psv:bounds} and Proposition~\ref{pss:carrier}.
\end{proof}

\subsubsection{The character split in Lemma~\ref{psp:step1}}\label{psv:ss-split}

\begin{lemma}[Splitting the kernel of $(S,W')$]\label{psv:split}
Let $\alpha=\nu^\flat\upsilon^m$ and $\beta=\upsilon^m$ be the characters of \eqref{pte:eq-F}, on $\auU(S)$ and on $\auU(W')$. For
$x\in\auU(S)(\auA)$, $\tilde h\in\auU(\widetilde W)(\auA)$, $h^\flat\in\auU(W^\flat)(\auA)$, polynomials $\widetilde P$ on
$\auHom(S,\widetilde W)\otimes\R$ and $Q$ on $\auHom(S,W^\flat)\otimes\R$, and functions $\Psi$, $\Psi^\flat$ modulo $\pi_M^2$,
\[
\theta^{S,W'}_{\alpha,\beta}\bigl(x,(\tilde h,h^\flat);[\widetilde PQ,\Psi\otimes\Psi^\flat]\bigr)=
\theta^{S,\widetilde W}_{\alpha,\beta}\bigl(x,\tilde h;[\widetilde P,\Psi]\bigr)\cdot\theta^{S,W^\flat}_{1,\beta}\bigl(x,h^\flat;[Q,\Psi^\flat]\bigr).
\]
In particular the kernel of $(S,W^\flat)$ in Lemma~\ref{psp:step1} has the characters $(1,\upsilon^m)$ and that of
$(S,\widetilde W)$ the characters $(\nu^\flat\upsilon^m,\ \upsilon^m|_{\auU(\widetilde W)})$, as used in
Theorems~\ref{pss:H5}, \ref{pss:b}, \ref{pss:a} and Lemma~\ref{psp:weights}. The other attachment (the factor
$\alpha(\det x)^{-1}$ on the $(S,W^\flat)$ kernel) gives the same product but is incompatible with the weights used there.
\end{lemma}
\begin{proof}
By the convention of \S\ref{pte:ss-conv}, $\omega_{\alpha,\beta}(x,h)=\alpha(\det x)^{-1}\beta(\det h)\,\omega_{\auHom(S,W')}(x,h)$. For
$h=(\tilde h,h^\flat)$ block diagonal, $\beta(\det h)=\beta(\det\tilde h)\beta(\det h^\flat)$. The compatibility of the
Heisenberg and Weil representations, of the splittings defined with $\upsilon$, and of the theta functionals with orthogonal
sums, which is the content of (W7) in \S\ref{po:ss-glue} (there for the summands of the target $V=V''\perp V^\circ$ of
$\auHom(W,V)$; the argument is the same for the summands $W'=\widetilde W\perp W^\flat$ of the target of $\auHom(S,W')$:
$\auHom(S,W')=\auHom(S,\widetilde W)\oplus\auHom(S,W^\flat)$ orthogonally, and the splitting of $\auU(S)$ for the sum is the
product of the splittings of the summands, the characters multiplying; at the infinite places this is visible on
\eqref{ptf:eq-compact}, $\Delta(\gamma)=\aubar{\det\gamma^-}$ being multiplicative on block-diagonal $\gamma$), gives
$\omega_{\auHom(S,W')}(x,h)=\omega_{\auHom(S,\widetilde W)}(x,\tilde h)\otimes\omega_{\auHom(S,W^\flat)}(x,h^\flat)$,
$\theta=\theta\otimes\theta$, and the datum $[\widetilde PQ,\Psi\otimes\Psi^\flat]$ is the tensor product of the two data (the
Fock vector of a product of polynomials in disjoint sets of variables is the tensor product of their Fock vectors, and the
finite data are products along the rows). The scalar $\alpha(\det x)^{-1}$ appears once; attaching it to the first
factor gives the identity. For the last sentence: with the characters $(\alpha,\beta)$ on $(S,W^\flat)$ the vector
$v_0$ of Theorem~\ref{pss:b} would have weight $(-\Delta'-a_\tau;6m+1-a_\tau)$, $a_\tau=m-\mu^+_{1,\tau}$, by
Corollary~\ref{pss:minors}(i), and with $(1,\beta)$ on $(S,\widetilde W)$ the weight $u_2$ of Theorem~\ref{pss:H5} would
change by $+a_\tau$ on all entries; the relation $u_3=(0;3m)+u_1+u_2$ is unchanged, but $u_1$ would no longer be
the weight $(-\Delta';6m+1)$ fixed in \S\ref{psp:ss-outline}, whose gap is checked in Lemma~\ref{psp:weights}. With the
attachment of the statement, the relation $u_3=(0;3m)+u_1+u_2$ holds: $(-\Delta';6m+1)+(\Delta';\mu^+_1+r+2m)+(0;3m)=
(0;\mu^+_1+r+11m+1)=(0;\mu^+_1+N''+2m)$.
\end{proof}

\noindent Lemma~\ref{psv:split} is the character split used in the proof of Lemma~\ref{psp:step1}.

\subsubsection{Conclusion}\label{psv:ss-main}

\begin{theorem}[Products span]\label{psv:main}
All inputs of Proposition~\ref{psp:inputs} hold in the form in which they are used in \S\ref{psp:sec}:
\begin{enumerate}
\item[(H1)--(H3)] Theorem~\ref{psq:H123} \textup(the gap $300\pspn$ exceeds $N^2+N+1$ for $N\le t$\textup);
\item[(H4)] Theorem~\ref{psv:H4}, for the order $o$ of Definition~\ref{psv:order} on $W$;
\item[(H5)] Theorem~\ref{pss:H5};
\item[(R)] Theorem~\ref{psr:R}, applied in Lemma~\ref{psp:step1} to $(S,W^\flat)$, $u=u_1$, the characters
$(1,\upsilon^m)$ \textup(Lemma~\ref{psv:split}\textup) and $\mathcal Q$ of Theorem~\ref{pss:b} \textup(Corollary~\ref{pss:bcor};
characters: Theorem~\ref{pss:a}\textup), and in Lemma~\ref{psp:step2} to $(W,V^\circ)$, $u=u^\circ$, the characters of
$\omega^\circ$ \textup(Theorem~\ref{pss:a}\textup), $\mathcal Q=\mathcal Q^\circ$ \textup(Theorem~\ref{psv:c}\textup) and the datum of
$K^\circ$ \textup(Proposition~\ref{psv:datum}\textup); in both cases the target is positive definite of rank more than
twice that of the source ($6m+1>6m$; $n^\circ>2r$);
\item[(P)] Theorem~\ref{psq:peaks}.
\end{enumerate}
Consequently Proposition~\ref{psp:main}, that is Proposition~\ref{po:span}, holds, modulo the published results
listed below.
\end{theorem}
\begin{proof}
The list is the content of the cited statements. In Lemma~\ref{psp:step2} the kernel $\theta_{\alpha^\circ,\beta^\circ}(x,h_p;Q_i,
\delta_\iota)$ of Theorem~\ref{psr:R}, for $Q_i=\pop^\circ_{u_i}$ with $(u_i)$ a basis of $W^\circ$, is
$\langle u_i,K^\circ(x,h_p)\rangle$ (Proposition~\ref{psv:datum}), so the right side of Theorem~\ref{psr:R} is
$\sum_p\int y\,F\,\langle\zeta_p,K^\circ(\cdot,h_p)\rangle$ for vectors $\zeta_p\in W^\circ$, as written there; the classes
$h_p$ of $\auU(V^\circ)$ modulo $J_{V^\circ}(\pi_M^2)$ are the classes modulo $J^\circ$ of \S\ref{po:ss-glue}. The proof
of Proposition~\ref{psp:main} in \S\ref{psp:ss-proof} uses exactly these inputs, the weights $u$ and $u^\circ$
entering through (H4) and Proposition~\ref{psp:isol}; with the order $o$ fixed on $W$, $\pspm_{u^\circ}$ and the
element $\mathrm{rot}(\theta)$ are taken for this order, which changes nothing in \S\ref{psp:ss-isol}--\S\ref{psp:ss-proof}.
\end{proof}

\emph{Published results used} (with the form used):
\begin{itemize}
\item Howe~\cite[(2.2)--(2.3)]{HoweJAMS}: polynomial Fock vectors are smooth vectors of the archimedean Weil
representation, and the derived action of the complexified Lie algebras maps polynomials to polynomials
(\S\ref{pss:ss-tool}, Theorems~\ref{psv:c}, \ref{psv:F}).
\item The Stone--von Neumann theorem (irreducibility of the Heisenberg representation and uniqueness of the
intertwiner up to a scalar) and G\aa rding's theorem (orbit maps of smooth vectors are smooth)
(\S\ref{psp:ss-fock}, \S\ref{pss:sec}).
\item Banaszczyk's transference theorem~\cite[Theorem~2.1]{Banaszczyk}: $\lambda_k(L)\lambda_1(L^\vee)\le k$ in rank $k$
(Lemma~\ref{psp:hom}).
\item Weil's identity~\cite{WeilSiegel} for absolutely convergent theta integrals with a definite target and
$M'>2n$, and the moderate growth of absolutely convergent Siegel Eisenstein series (Lemma~\ref{psr:SW}).
\item Schmid's theorem (D4) of \S\ref{ptq:ss-H}: the $K$-types of $\pi_\lambda$ have highest weights
$\Lambda^K(\lambda)+\sum n_\beta\beta$, $n_\beta\ge0$, over positive non-compact roots (Lemma~\ref{psv:v}).
\item The inputs of Theorem~\ref{par:A1} through Lemma~\ref{plf:A12} (the decay bounds proved in the paper and
Lemma~\ref{pm:kernel}), for the integrability of the invariant derivatives of $c_W$ (Theorem~\ref{psv:ybar}).
\item The continuity of the Weil representation on the Schwartz space and of the theta functional on it (used as in
Proposition~\ref{pss:carrier}, Theorems~\ref{psv:c}, \ref{psv:F}).
\end{itemize}
Borel's finiteness theorem is not needed here: both targets $W^\flat$ and $V^\circ$ are standard spaces, for which
the finiteness of the classes is Proposition~\ref{pf:classes}.

\begin{remark}[(G2) of Proposition~\ref{po:glue}]\label{psv:G2}
The proof of (G2) in \S\ref{po:ss-glue} reads: if $\Theta^\circ y=0$ then, $c_i\ne0$, the hypothesis of
Proposition~\ref{po:span} holds, so $y=0$. Indeed $(\Theta^\circ y)_i=M^{l_1}c_i\,\#\Gamma^1\int y\,F_i\,K^\circ(\cdot,\cdot)$
with $M^{l_1}\#\Gamma^1\ne0$ and $c_i=c_{\mathsf g_i}\ne0$ (Lemma~\ref{po:slice}), so $\Theta^\circ y=0$ is the vanishing of
$\int y\,F_iK^\circ(\cdot,h^\circ)$ for all $i$ and all classes $h^\circ$ modulo $J^\circ$, the hypothesis of
Proposition~\ref{psp:main}; Theorem~\ref{psv:main} gives $y=0$. So $\Theta^\circ$ is injective. If $c'_\vartheta=0$ then
$\Theta^\circ=0$ by (G1), against $e_0\poA(W_j)\ne0$; so $c'_\vartheta\ne0$, and $\Theta_\vartheta$ is injective on $e_0\poA(W_j)$
by (G1). The remaining clause $C_{\mathrm x}>0$ follows from $C_{\mathrm x}=|c'_\vartheta|^2H_j(\vartheta)$
(Proposition~\ref{po:exchange}) and the non-vanishing of the partial $L$-values in \eqref{po:eq-H}, which is
Lemma~\ref{po:Hpos} (absolute convergence for $z\ge2$, Hecke's theorem on the line $\mathrm{Re}\,s=1$ for $z=1$).
\end{remark}


%% ===== proofs/C-firstfactors.tex =====
% proofs/C-firstfactors.tex
% The first factors of the one-lift group: integrality (I1) and size (I2) for the first factors for
% which the exchange identity is proved; and the tame parts at the chosen places (the characters
% omega, chi_M, mu^flat, mu_j^0), with the strengthened existence lemma for characters of U(1).
% To be \input after C-onelift.tex, C-lifts.tex, C-threeex.tex, C-threeex2.tex.
% Labels start with pfq:. Only amsmath/amssymb/amsthm.
\providecommand{\auO}{\mathcal O}
\providecommand{\auU}{\mathrm U}
\providecommand{\auGL}{\mathrm{GL}}
\providecommand{\auHom}{\operatorname{Hom}}
\providecommand{\auA}{\mathbb A}
\providecommand{\auF}{\mathbb F}
\providecommand{\auPl}{\Phi}
\providecommand{\auNm}{\operatorname{Nm}}
\providecommand{\auAut}{\operatorname{Aut}}
\providecommand{\aubar}[1]{\overline{#1}}
\providecommand{\pao}{\mathfrak o}
\providecommand{\plone}{\mathbf 1}
\providecommand{\plgb}[2]{\|#2\|_{#1}}
\providecommand{\poA}{\mathcal A}
\providecommand{\poT}{\mathbb T}
\providecommand{\poAA}{A_\ast}
\providecommand{\pop}{\mathsf p}
\providecommand{\plfK}{\mathsf K}
\providecommand{\pteSD}{\mathrm{SD}}
\providecommand{\pteCS}{\mathop{\textstyle\sum^{\mathrm{cl}}}}
\providecommand{\pfqC}{\mathsf C}
\providecommand{\pfqW}{W^{\mathrm{add}}}
\providecommand{\pfqm}{\mathfrak m_\vartheta}

\subsection{The first factors and the tame parts at the chosen places}\label{pfq:sec}

This subsection treats two points of the proof in \S\ref{po:sec}: the first factors and the tame parts
of the characters at the chosen places.

\emph{First factors} (\S\S\ref{pfq:ss-hook}--\ref{pfq:ss-ff}). The exchange identity,
Proposition~\ref{po:exchange}, is proved in \S\S\ref{pte:sec}--\ref{ptf:sec} for the first factors of
Definition~\ref{ptf:def-ff}. We prove (I1) and (I2) for an
explicit non-zero constant multiple of these first factors
(Theorem~\ref{pfq:ff}), and check that Proposition~\ref{po:exchange} holds for this multiple
(Proposition~\ref{pfq:exchange}). This proves Lemma~\ref{po:ff}.

\emph{Tame parts} (\S\ref{pfq:ss-tame}). We define the residue characters $\omega_w$ at the chosen
places, the characters $\chi_M$ ($=\chi_0$ of \S\ref{po:ss-hecke}), $\mu^\flat$ and $\mu^0_j$, prove the
existence lemma for characters of $\auU(1)$ with prescribed weight and prescribed tame parts
(Lemma~\ref{pfq:mu}; compare Lemma~\ref{pz:mu}), and check that the tame parts have the
properties used in \S\ref{po:sec} and in Lemma~\ref{ppe:three}.

Throughout, $j$ ($1\le j\le m$) is fixed, and we use the notation of \S\ref{po:ss-glue},
\S\ref{pte:ss-third}, \S\ref{pte:ss-left} and \S\ref{ptf:ss-ff}: $V''=I_{N''}$, $W=W_j$ with $r_+$
positive basis vectors, $\pfqm=\pi_M^2\prod_{w\in Q_\vartheta}w$, the characters
$\chi'=\chi_M^{\,j}\vartheta$ and $\nu''=\nu''_\vartheta=\mu^0_j\vartheta\,\upsilon^{j+\lfloor(t+k-1)/2\rfloor}$, the
choices $i=(\Lambda_i,h^\flat_i,\varrho_i)$, and the space
\[
\pfqC=I_{3m}\oplus(E^j)^-\oplus E^{3m-j}\oplus W^\flat
\]
of \S\ref{ptf:ss-ff}. We write $\#\iota(\tau)$ for the number of columns of $\mu^+_\tau$ and
$\mathsf q=\sum_\tau\#\iota(\tau)\le4fb$. By the proof of Proposition~\ref{pte:third},
\begin{equation}\label{pfq:eq-weight}
\nu''\ \text{has weight}\ -\#\iota(\tau)\ \text{at}\ \tau ,
\end{equation}
since $\mu^0_j$ has weight $-\kappa_j$, $\upsilon$ weight $1$, $\vartheta$ weight $0$, and
$\#\iota(\tau)+j+\lfloor\tfrac{t+k-1}2\rfloor=\kappa_j(\tau)$. Kernels are those of \S\ref{pcx:ss-statement}:
$\theta^{S,T}_{\alpha,\beta}(g,h;[P,\Psi])$, the first index being the character of the source group.
The kernels $\plfK^{(L,N)}_{a,b}[\Psi,Q]$ of \S\ref{plf:ss-KS} are the kernels
$\theta^{L,N}_{a,b}(g,h;[Q,\Psi])$ of \S\ref{pc:ss-data} at finite points $g,h$ (archimedean components
$1$), with the same convention for the indices (the characters twist the Weil representation by
$a(\det g)^{-1}b(\det h)$); this is how Lemma~\ref{plf:KS} is used below. The identification includes
the primes $w$ of $Q_\vartheta$ at which $\alpha_w$ is not a unit (primes above $q_1$ dividing $\omega$,
Proposition~\ref{pgv:alpha}; the definition of a twist does not exclude them): there the datum of
\S\ref{pc:ss-data} is $\Psi^\natural$ of Definition~\ref{pgx:conv}, supported on $a_wM(\auO_w)$, and
the datum $f_\Psi$ of \S\ref{plf:ss-KS} is the same function of the frame coordinate $x$ (it is defined
with the same lattice $L_w=a_wM(\auO_w)$ and the same generator $a_w$, that of \S\ref{ps:ss-weil}),
equivalently the same function $\plone[T\ \text{integral}]\Psi(T\bmod\pcm)$ of the coordinate $T$ of (N5)
(Lemma~\ref{pgx:dict}(c)); so $f_\Psi=\Psi^\natural$ literally. The proof of
(KS1) and (KS2) covers this case (the paragraph after \eqref{plf:eq-vacexp}). So (I1) and (I2) below
need no restriction on $Q_\vartheta$ beyond the definition of a twist and $\ell\nmid\auNm Q_\vartheta$.

\subsubsection{A complementary hook number}\label{pfq:ss-hook}

We use the tensor model of \S\ref{pk:ss-branch}: for a family $\mu=(\mu_\tau)$ of diagrams for $\auGL_N$,
$\poT_\mu(R)=\bigotimes_{\tau}\bigotimes_{\mathfrak c}\Lambda^{h_{\mathfrak c}}R^N$ with the orthonormal basis
$(e_U)$ indexed by the fillings, $\mathsf C_\mu\subseteq\poT_\mu(\C)$ the span of the vectors
$\mathrm{ev}_y$, $\pi_\mu$ the orthogonal projection onto $\mathsf C_\mu$, and $W_\mu(\Z)=\mathsf m(\mathsf
C_\mu\cap\poT_\mu(\Z))$ (\eqref{pk:eq-dual}); $\mathsf m$ identifies $W_\mu$ isometrically with
$\mathsf C_\mu$ (proof of Lemma~\ref{po:frame}). Since $\mathsf C_\mu\cap\poT_\mu(\Z)$ is the intersection of
$\poT_\mu(\Z)$ with a subspace, it is a direct summand of $\poT_\mu(\Z)$; hence for every valuation
ring $\pao\supseteq\Z$ of $\C$, $\mathsf C_\mu\cap\poT_\mu(\pao)=(\mathsf C_\mu\cap\poT_\mu(\Z))\otimes\pao$,
and a vector of $W_\mu$ has all coordinates in the basis of $W_\mu(\Z)$ in $\pao$ if and only if all its
tensor coordinates are in $\pao$. (This is the remark ``the two unit balls coincide'' of
\S\ref{po:ss-set}.)

\begin{lemma}[Complementary hook number]\label{pfq:hook}
Let $\mu=(\mu_\tau)$ be diagrams for $\auGL_N$ all of whose columns have heights $h_{\mathfrak c}<N$, let
$\mu^{\mathrm c}_\tau$ be the diagram whose columns have the heights $N-h_{\mathfrak c}$, and
$H^{\mathrm c}=\prod_\tau H_{\mu^{\mathrm c}_\tau}$ (Lemma~\ref{pk:hook}). Then
$H^{\mathrm c}\pi_\mu(e_U)\in\mathsf C_\mu\cap\poT_\mu(\Z)$ for every filling $U$. Consequently
$H^{\mathrm c}\pi_\mu(\poT_\mu(R))\subseteq\mathsf C_\mu\cap\poT_\mu(R)$ for every subring $R\subseteq\C$.
Moreover $H^{\mathrm c}\le |\mu^{\mathrm c}|^{3|\mu^{\mathrm c}|}=e^{3e}$, where the base is the integer $e=|\mu^{\mathrm c}|=\sum_\tau\sum_{\mathfrak c}(N-h_{\mathfrak c})$.
\end{lemma}
\begin{proof}
All objects are tensor products over $\tau$ and over the columns; it suffices to treat one $\tau$. For
$0\le h\le N$ let $\star\colon\Lambda^h\Z^N\to\Lambda^{N-h}\Z^N$ be the map with $v\wedge w=(\star v,w)\,
e_1\wedge\dots\wedge e_N$ for all $w\in\Lambda^{N-h}$, $(\,,)$ the symmetric bilinear form with orthonormal
basis $(e_R)$; thus $\star e_R=\pm e_{R^c}$, and $\star$ is an isometric bijection of lattices. For
$g\in\auGL_N(\C)$, $(gv)\wedge(gw)=\det g\,(v\wedge w)$ and $(g\,x,w)=(x,g^{\mathsf t}w)$ on exterior powers, so
$(g^{\mathsf t}\star(gv),w)=\det g\,(\star v,w)$ for all $w$, that is, $\star(gv)=\det g\cdot g^{-\mathsf t}\star v$.
Let $\mathrm{St}=\bigotimes_{\mathfrak c}\star\colon\poT_\mu\to\poT_{\mu^{\mathrm c}}$; it maps the orthonormal
basis to $\pm$ the orthonormal basis, is real, and satisfies $\mathrm{St}(gv)=(\det g)^{c}g^{-\mathsf t}\mathrm{St}(v)$,
$c$ the number of columns.

For $y\in M_N(\C)$, $\mathrm{ev}_y=y\cdot e_{U_0}$ with $e_{U_0}=\bigotimes_{\mathfrak c}e_1\wedge\dots\wedge e_{h_{\mathfrak c}}$
(functoriality of $\Lambda^h$); since $\auGL_N(\C)$ is dense in $M_N(\C)$ and $\mathsf C_\mu$ is finite-dimensional,
$\mathsf C_\mu$ is the span of the $g\,e_{U_0}$, $g\in\auGL_N(\C)$. Let $w_0$ be the permutation matrix of
$i\mapsto N+1-i$. Then $\star(e_1\wedge\dots\wedge e_h)=\pm e_{h+1}\wedge\dots\wedge e_N=\pm w_0(e_1\wedge\dots\wedge e_{N-h})$,
so $\mathrm{St}(e_{U_0})=\pm w_0e_{U_0^{\mathrm c}}$, $U_0^{\mathrm c}$ the leading filling of $\mu^{\mathrm c}$. Hence
$\mathrm{St}(\mathsf C_\mu)$ is the span of the $(\det g)^cg^{-\mathsf t}w_0e_{U^{\mathrm c}_0}$, which is
$\mathsf C_{\mu^{\mathrm c}}$. As $\mathrm{St}$ is unitary, $\mathrm{St}\circ\pi_\mu=\pi_{\mu^{\mathrm c}}\circ\mathrm{St}$.
By Lemma~\ref{pk:hook} for $\mu^{\mathrm c}$ and \eqref{pk:eq-dual}, $H^{\mathrm c}\pi_{\mu^{\mathrm c}}(e_{U'})\in
\mathsf C_{\mu^{\mathrm c}}\cap\poT_{\mu^{\mathrm c}}(\Z)$ for every filling $U'$ (this is the argument of the
proof of Lemma~\ref{po:frame}). With $\mathrm{St}(e_U)=\pm e_{U'}$,
$H^{\mathrm c}\pi_\mu(e_U)=\pm\mathrm{St}^{-1}(H^{\mathrm c}\pi_{\mu^{\mathrm c}}(e_{U'}))\in\mathsf C_\mu\cap\poT_\mu(\Z)$.
The second assertion follows by linearity, the bound from Lemma~\ref{pk:hook} ($H_\mu\le|\mu|^{3|\mu|}$,
and $\prod_\tau|\mu^{\mathrm c}_\tau|^{3|\mu^{\mathrm c}_\tau|}\le |\mu^{\mathrm c}|^{3|\mu^{\mathrm c}|}$).
\end{proof}

\subsubsection{The first factors as a projection}\label{pfq:ss-proj}

\emph{Fillings and Laplace coefficients.} For $Y\in\auGL_{6m+1}(E)$ let $\eta^Y_i=(\eta^Y_{i,\tau k})$ be the
constant rows of Lemma~\ref{ptf:H}, built with the short vectors given by the columns of $Y$ (the rows
of $\eta^Y_{i,\tau k}$ for a column of height $r_+$ are the transposes of the columns of $\tau(Y)$ selected by
$\varrho_i$, placed in the columns of $W^\flat$); for $Y=Y_i:=Y_{h^\flat_i}$ of Lemma~\ref{au:shortbasis}
these are the rows $\eta_i$ of \S\ref{pte:ss-third}. A \emph{filling} $U=(U_{\tau k})$ chooses, for each $\tau$
and each $k\in\iota(\tau)$, a subset $U_{\tau k}\subseteq\{1,\dots,N''\}$ of $h_k$ rows; put
$b_U(X)=\prod_{\tau,k}\det X_{\tau;U_{\tau k},[h_k]}$ (the product of minors of \S\ref{pk:ss-branch}) and
$Q^Y_U(Z)=\prod_{\tau,k}Q_{\tau k,U_{\tau k}}(Z_\tau)$, with $Q_{\tau k,R}$ the polynomials of
Lemma~\ref{ptf:laplace} for the rows $\eta^Y_i$. Let $\pfqW=I_{3m}\oplus E^{3m-j}\oplus W^\flat$, the
positive definite summand of $\pfqC$ (so $\pfqC=(E^j)^-\oplus\pfqW$ after reordering the frame); the
variables $Z$ are the Fock variables of $\auHom(\pfqW,V'')$.

\begin{lemma}[Laplace coefficients]\label{pfq:Q}
\begin{enumerate}
\item[(Q1)] $\pteSD(\eta^Y_i)=\sum_Ub_U(X)\,Q^Y_U(Z)$.
\item[(Q2)] Every $Q^Y_U$ is homogeneous in $Z$ of degree $e=\sum_\tau\sum_{k\in\iota(\tau)}(N''-h_k)$, and
$e\le(10m+1)\mathsf q\le M$.
\item[(Q3)] $Q^Y_U$ is a polynomial whose coefficients are polynomials with integer coefficients, independent
of $Y$, in the numbers $\tau(Y_{ab})$ ($\tau\in\auPl$, $1\le a,b\le6m+1$).
\item[(Q4)] If $|\tau'(Y_{ab})|\le M^4$ for all embeddings $\tau'\colon E\to\C$ and all $a,b$, and all entries of
$z$ have absolute value at most $M^4$, then $\sum_U|Q^Y_U(z)|^2\le M^{9(12m+1)\mathsf q}$; the same holds for the
polynomials obtained from $Q^Y_U$ by applying any $\sigma\in\auAut(\C)$ to the coefficients.
\item[(Q5)] Let $\gamma\in\auU(W^\flat)(E^+)$ and $B=\mathrm{diag}(1_{3m},1_{3m-j},\gamma)\in\auU(\pfqW)(E^+)$. Then
$Q^Y_U(Z_\tau\,\tau(B)^{-\mathsf t})_\tau=\prod_\tau\tau(\det\gamma)^{-\#\iota(\tau)}\,Q^{\gamma Y}_U(Z)$. For
$\gamma'\in\auU(I_{3m})(E^+)$ and $B'=\mathrm{diag}(\gamma',1,1)$,
$Q^Y_U(Z_\tau\,\tau(B')^{-\mathsf t})_\tau=\prod_\tau\tau(\det\gamma')^{-\#\iota(\tau)}\,Q^Y_U(Z)$.
\end{enumerate}
\end{lemma}
\begin{proof}
(Q1) is the expansion in the proof of Lemma~\ref{ptf:laplace} (multiplied over $\tau$ and $k$). By that proof,
$Q_{\tau k,R}(Z_\tau)=\pm\det G_{\tau k,R}$, where $G_{\tau k,R}$ is the square matrix of size $12m+1-j$ whose
rows are the $2m-(r_+-h_k)$ constant rows of $\eta^Y_{i,\tau k}$ vanishing on the columns of $X$ (restricted to the
columns of $\pfqW$) and the $N''-h_k$ rows of $Z_\tau$ outside $R$; the sign does not depend on $Y$ or $Z$.
(Q2) Each $\det G_{\tau k,R}$ is homogeneous of degree $N''-h_k$ in $Z$. By Lemma~\ref{ptf:H} and its proof,
every column of $\iota(\tau)$ has height $h_k\ge r-m$ (one less than the parameter $r+1-m$ of that lemma; the subscript $k$ here is a column index), and $N''-(r-m)=10m+1$; so $e\le(10m+1)\mathsf q$, and $(10m+1)\mathsf q\le44fmb\le M$ by
\eqref{eq:au-sizes}, \eqref{eq:au-scales} and $M\ge b^{20}$ (Lemma~\ref{an:lem-sizes}). (Q3) The constant
rows have entries $0$, $1$ and $\tau(Y_{ab})$. (Q4) For fixed $\tau,k$, the numbers $\det G_{\tau k,R}$, $R$ running
over the $h_k$-subsets, are some of the maximal minors of the matrix $G_{\tau k}$ formed by the constant rows
vanishing on $X$ and all rows of $Z_\tau$; by the Cauchy--Binet formula and Hadamard's inequality,
$\sum_R|\det G_{\tau k,R}|^2\le\det(G^*_{\tau k}G_{\tau k})\le\prod(\text{squared lengths of the columns of }G_{\tau k})$.
There are $12m+1-j\le12m+1$ columns, each with at most $2m+N''\le M$ entries of absolute value at most $M^4$, so
the product is at most $M^{9(12m+1)}$. Since $\sum_U|Q^Y_U(z)|^2=\prod_{\tau,k}\sum_R|\det G_{\tau k,R}|^2$, the
bound follows. Applying $\sigma$ to the coefficients replaces $\tau(Y_{ab})$ by $\sigma\tau(Y_{ab})$ (by (Q3)), which
satisfies the same bound. (Q5) Multiplying on the right by $\tau(B)^{-\mathsf t}$ the rows of $Z_\tau$ in $G_{\tau k,R}$
gives the matrix $\bigl(\begin{smallmatrix}c\,\tau(B)^{\mathsf t}\\ Z_{R^c}\end{smallmatrix}\bigr)\tau(B)^{-\mathsf t}$, $c$ the
constant rows; its determinant is $\tau(\det\gamma)^{-1}\det\bigl(\begin{smallmatrix}c\,\tau(B)^{\mathsf t}\\ Z_{R^c}
\end{smallmatrix}\bigr)$. The constant rows vanish on $I_{3m}$, are basis vectors on $E^{3m-j}$ (where $B=1$), and on
$W^\flat$ are rows ${}^ty$, $y$ a column of $\tau(Y)$, which become ${}^t(\tau(\gamma)y)$, the corresponding column of
$\tau(\gamma Y)$. So $c\,\tau(B)^{\mathsf t}$ are the constant rows for $\gamma Y$. For $B'$ the constant rows vanish on
$I_{3m}$, so $c\,\tau(B')^{\mathsf t}=c$. Multiply over $(\tau,k)$.
\end{proof}

\begin{lemma}[Factorisation]\label{pfq:factor}
Let $\Psi^{(j)}$ be the indicator of the matrix of the inclusion $E^j\to V''$ (on $N''\times j$ matrices modulo
$\pfqm$), and $\Psi^{\mathrm{add}}_i$ the function on $N''\times(12m+1-j)$ matrices modulo $\pfqm$ which is
$\psi^{\mathrm{inv}}_{\Lambda_i}(x')$ if the columns on $E^{3m-j}\oplus W^\flat$ are those of the inclusion and $0$ otherwise
($x'$ the columns on $I_{3m}$). Then $\Phi^{\mathsf C}_i(x_{(E^j)^-}\mid x_{\pfqW})=\Psi^{(j)}(x_{(E^j)^-})
\Psi^{\mathrm{add}}_i(x_{\pfqW})$, and for every polynomial $R$ in $Z$ and $h'\in\auU(I_{3m})(\auA_f)$,
\[
\theta^{\pfqC,V''}_{\nu'',\nu''\upsilon^j}\bigl(c_i(h'),h;[R,\Phi^{\pfqC}_i]\bigr)=\epsilon_j^{-1}\,\aubar{k_1(h)}\cdot
\theta^{\pfqW,V''}_{\nu'',\nu''}\bigl((h',1,h^\flat_i),h;[R,\Psi^{\mathrm{add}}_i]\bigr),
\]
where $k_1(h)=\theta^{E^j,V''}_{1,1}(1,h;[1,\Psi^{(j)}])$ and $\epsilon_j=\epsilon(E^j,V'')$ is the number of modulus one of
Lemma~\ref{ptf:conj}; it depends only on $j$.
\end{lemma}
\begin{proof}
The product decomposition of $\Phi^{\mathsf C}_i$ is its definition (Definition~\ref{ptf:def-ff}). Apply
Lemma~\ref{ptf:sum} with $S_1=(E^j)^-$, $S_2=\pfqW$ (the order of the summands of a frame is immaterial: a
permutation of the frame of $\pfqC$ is an isometry of standard lattices and changes kernels, data and polynomials
compatibly), $\alpha=\nu''$, $\beta_1=\upsilon^j$, $\beta_2=\nu''$, $P_1=1$, $P_2=R$ (the polynomial $R$ involves only
the positive columns of $\pfqC$, which are those of $\pfqW$), at the source point $c_i(h')=(1,(h',1,h^\flat_i))$:
\[
\theta^{\pfqC,V''}_{\nu'',\nu''\upsilon^j}(c_i(h'),h;[R,\Phi^{\mathsf C}_i])=\theta^{(E^j)^-,V''}_{\nu'',\upsilon^j}(1,h;[1,
\Psi^{(j)}])\,\theta^{\pfqW,V''}_{\nu'',\nu''}((h',1,h^\flat_i),h;[R,\Psi^{\mathrm{add}}_i]).
\]
By Lemma~\ref{ptf:conj} with $S=E^j$, $T=V''$, $\alpha=\nu''^{-1}\upsilon^{N''}$, $\beta=1$ (so that
$\alpha^{-1}\upsilon^{N''}=\nu''$ and $\beta^{-1}\upsilon^{j}=\upsilon^j$), and as $\Psi^{(j)}$ is real,
$\aubar{\theta^{E^j,V''}_{\alpha,1}(1,h;[1,\Psi^{(j)}])}=\epsilon_j\,\theta^{(E^j)^-,V''}_{\nu'',\upsilon^j}(1,h;[1,\Psi^{(j)}])$;
and $\theta^{E^j,V''}_{\alpha,1}(1,h;\cdot)=\theta^{E^j,V''}_{1,1}(1,h;\cdot)$ by \eqref{pte:eq-twist}, the source point being $1$.
\end{proof}

\begin{lemma}[Moving the class of $W^\flat$]\label{pfq:move}
Let $x\in\auU(\pfqW)(\auA_f)$, $\gamma\in\auU(W^\flat)(E^+)$, $B=\mathrm{diag}(1,1,\gamma)$, and $\Psi$ any datum. Then
\[
\theta^{\pfqW,V''}_{\nu'',\nu''}(x,h;[Q^Y_U,\Psi])=\theta^{\pfqW,V''}_{\nu'',\nu''}(B_fx,h;[Q^{\gamma Y}_U,\Psi]),
\]
$B_f$ the finite part of $B$. The same holds with $B'=\mathrm{diag}(\gamma',1,1)$, $\gamma'\in\auU(I_{3m})(E^+)$, and
$Q^Y_U$ on both sides.
\end{lemma}
\begin{proof}
Let $\delta$ be a rational point of $\auU(\pfqW)$. The operator $\omega_{\nu'',\nu''}(\delta,1)$ leaves $\theta$ invariant
((W1) and $\nu''(\det\delta)=1$, $\nu''$ being automorphic), and $\delta=\delta_\infty\delta_f$, with $\delta_\infty$ commuting with
finite points; hence $\theta^{\pfqW,V''}_{\nu'',\nu''}(\delta_fx,h;[P,\Psi])=\theta^{\pfqW,V''}_{\nu'',\nu''}(x,h;[\omega_\infty
(\delta_\infty^{-1},1)P,\Psi])$. By Lemma~\ref{ptf:fock} (all product signs positive, $g=\delta_\infty^{-1}$, $h=1$, source weight
$-\#\iota(\tau)$ by \eqref{pfq:eq-weight}), $\omega_\infty(\delta_\infty^{-1},1)P=\prod_\tau\tau(\det\delta)^{-\#\iota(\tau)}\,
P(Z_\tau\tau(\delta)^{\mathsf t})_\tau$. Take $\delta=B$ and $P=Q^Y_U$. Since $\tau(B)^{\mathsf t}=\tau(B^{-1})^{-\mathsf t}$, (Q5) for the
element $\gamma^{-1}$ gives $Q^Y_U(Z_\tau\tau(B)^{\mathsf t})=\prod_\tau\tau(\det\gamma)^{\#\iota(\tau)}Q^{\gamma^{-1}Y}_U(Z)$, and
$\det B=\det\gamma$; the two scalars cancel, so $\theta(B_fx,h;[Q^Y_U,\Psi])=\theta(x,h;[Q^{\gamma^{-1}Y}_U,\Psi])$ for all $x$ and
$Y$. Replacing $Y$ by $\gamma Y$ gives the claim. The case of $B'$ is the same, with the second part of (Q5).
\end{proof}

\begin{lemma}[The first factors as a projection]\label{pfq:proj}
For $h\in\auU(V'')(\auA_f)$ and a filling $U$ put
\[
a_U(h)=\pteCS_{K_{I_{3m}}}\Bigl(h'\mapsto\mathrm a_{\chi'}(h')\,\theta^{\pfqC,V''}_{\nu'',\nu''\upsilon^j}\bigl(c_i(h'),h;[Q^{Y_i}_U,
\Phi^{\mathsf C}_i]\bigr)\Bigr),\qquad a(h)=\sum_Ua_U(h)\,e_U\in\poT_{\mu^+}(\C^{N''}).
\]
Then the first factor of Definition~\ref{ptf:def-ff} is $\varphi''_{\vartheta,i}(h)=\pi_{\mu^+}\bigl(\aubar{a(h)}\bigr)$, where the bar
conjugates the tensor coordinates and $W''=W_{\mu^+}(\C^{N''})$ is identified with $\mathsf C_{\mu^+}$ as in
\S\ref{pfq:ss-hook}.
\end{lemma}
\begin{proof}
By the proof of Lemma~\ref{ptf:laplace}, $u\mapsto\pop''_u$ is a conjugate-linear bijection of $W''=\mathsf C_{\mu^+}$ onto the span
$\Sigma_{\mu^+}\otimes\C$ of the $b_U$, given by $\pop''_u=\mathsf m(\aubar u)$, where $\mathsf m(v)=\sum_Uv_Ub_U$ is the multiplication
map of \S\ref{pk:ss-branch} (tensor coordinates $v_U$); its kernel on $\poT_{\mu^+}(\C)$ is $\mathsf C_{\mu^+}^\perp$, so
$\mathsf m(v)=\mathsf m(\pi_{\mu^+}v)=\pop''_{\pi_{\mu^+}\aubar v}$ for every $v$ ($\pi_{\mu^+}$ is defined over $\Q$ and commutes with
the conjugation of coordinates). Let $L_h$ be the linear form $R\mapsto\pteCS(h'\mapsto\mathrm a_{\chi'}(h')\theta^{\pfqC,V''}(c_i(h'),h;
[R,\Phi^{\mathsf C}_i]))$ on polynomials in $Z$. By definition $c_\lambda(h)=L_h(R_\lambda)$, and $\pop''_{\varphi''_{\vartheta,i}(h)}=\sum_\lambda
c_\lambda(h)\pop''_{u_\lambda}=(\mathrm{id}\otimes L_h)\bigl(\sum_\lambda\pop''_{u_\lambda}\otimes R_\lambda\bigr)$. By \eqref{ptf:eq-SDexp} and (Q1),
$\sum_\lambda\pop''_{u_\lambda}\otimes R_\lambda=\pteSD(\eta_i)=\sum_Ub_U\otimes Q^{Y_i}_U$, so $\pop''_{\varphi''_{\vartheta,i}(h)}=\sum_Ua_U(h)b_U=
\mathsf m(a(h))=\pop''_{\pi_{\mu^+}\aubar{a(h)}}$. Injectivity of $u\mapsto\pop''_u$ on $W''$ gives the claim.
\end{proof}

\begin{remark}[Convention]\label{pfq:rem-linear}
The map $u\mapsto\pop''_u$ is conjugate-linear: by \eqref{po:eq-poly}, $\pop_\xi(x)=\langle\xi,\mathrm{ev}(\cdot)\rangle$
with $\langle\ ,\,\rangle$ the inner product of the representation space $\poT^+\otimes\poT^-$, which, like all
inner products of representation spaces and of functions, is conjugate-linear in the first variable;
(Cv1) of \S\ref{pb:ss-setting} concerns hermitian forms on the spaces $W_j$, $Y_j$, and states that inner
products of functions and of representation spaces follow the opposite convention; the form of
$\poT^+\otimes\poT^-$ is of this second kind, so there is no conflict with (Cv1). This is the convention of \S\ref{po:ss-glue}, of the
proof of Lemma~\ref{ptf:laplace} and of Proposition~\ref{ptf:ff}, and it is used throughout. (Had the map been
linear, Lemma~\ref{pfq:proj} would read $\varphi''_{\vartheta,i}=\pi_{\mu^+}(a(h))$ and the arguments below would
go through with $\aubar\kappa$ replaced by $\kappa$ and $\iota$, $\aubar\iota$ interchanged in the proof of (I1),
(KS1) holding at both; we do not use this.)
\end{remark}

\begin{remark}[The summands are functions on the classes]\label{pfq:rem-classes}
The function $h'\mapsto\mathrm a_{\chi'}(h')\theta^{\pfqC,V''}(c_i(h'),h;[Q^{Y_i}_U,\Phi^{\mathsf C}_i])$ is left invariant under
$\auU(I_{3m})(E^+)$: by Lemma~\ref{pfq:factor} and the second part of Lemma~\ref{pfq:move} for the kernel, and because
$\theta^{I_m,I_{3m}}_{1,\chi'}$ has the vacuum at infinity and characters of weight $0$, so that rational points act trivially
(Lemma~\ref{ptf:fock}). It is right invariant under $K_{I_{3m}}$: by Lemma~\ref{pc:level}, $k\in K_{I_{3m}}$ multiplies the kernel by
$\nu''_{\pfqm}(\det\bar k)^{-1}$ ($\psi^{\mathrm{inv}}_{\Lambda_i}$ is invariant under right multiplication by $\auGL_{3m}$) and
$\mathrm a_{\chi'}$ by $\chi'_{\pfqm}(\det\bar k)$; the tame parts of $\nu''$ and $\chi'$ at every prime of $\pfqm$ coincide
($\omega_w^j$ at $w\in\mathsf M$ and $\psi_w$ at $w\in Q_\vartheta$, Proposition~\ref{pfq:tame}(T4)). So the class sums of
Lemma~\ref{pfq:proj} (and of Definition~\ref{ptf:def-ff}) do not depend on the representatives.
\end{remark}

\subsubsection{The first factors used in \S\ref{po:sec}; integrality and size}\label{pfq:ss-ff}

Let $c_{l,n_N,e}$ be the kernel scales of Lemma~\ref{plf:KS}, $n'_{\mathrm{Aut}}$ the integer of Remark~\ref{plf:rem-classes},
$e$ the degree of (Q2), and $H^{\mathrm c}$ the number of Lemma~\ref{pfq:hook} for $\mu^+$ and $N=N''$ (the columns of
$\mu^+$ have heights at most $r_+<N''$). Put
\[
\kappa_1=n'^{\,2}_{\mathrm{Aut}}\;c_{m,3m,0}\;\aubar{c_{j,N'',0}}\;\epsilon_j\;c_{12m+1-j,N'',e},\qquad\kappa=H^{\mathrm c}\,\aubar{\kappa_1}.
\]

\begin{definition}[First factors]\label{pfq:def-ff}
The first factors of \S\ref{po:ss-glue} are $\varphi^\natural_{\vartheta,i}=\kappa\,\varphi''_{\vartheta,i}$, with $\varphi''_{\vartheta,i}$ as in
Definition~\ref{ptf:def-ff}. The number $\kappa\ne0$ depends only on $j$ (not on $\vartheta$, $i$, $h$).
\end{definition}

\begin{proposition}[The exchange identity for $\varphi^\natural$]\label{pfq:exchange}
The functions $\varphi^\natural_{\vartheta,i}$ lie in $\poA_{\mu^+}(J'_{\vartheta,i},\chi''^{-1}_\vartheta)$, and Proposition~\ref{po:exchange} holds for
them, with the same functions $F_i$, with $c'_\vartheta$ replaced by $c'_\vartheta/\aubar\kappa$ and $C_{\mathrm x}$ by $C_{\mathrm x}/|\kappa|^2$.
\end{proposition}
\begin{proof}
The type and level are those of $\varphi''_{\vartheta,i}$ (Proposition~\ref{ptf:leftside}), multiplied by a constant. The theta series
$\Xi_{\vartheta,i}$ of \S\ref{po:ss-glue} is conjugate-linear in the first factor ($\theta''_{\vartheta,s}$ is conjugate-linear in the
vector), so $\Xi^\natural_{\vartheta,i}=\aubar\kappa\,\Xi_{\vartheta,i}$, and $c'_\vartheta\Xi_{\vartheta,i}=F_i$ (\S\ref{pte:ss-proof}) gives
$(c'_\vartheta/\aubar\kappa)\Xi^\natural_{\vartheta,i}=F_i$, while $|c'_\vartheta/\aubar\kappa|^2H_j(\vartheta)=C_{\mathrm x}/|\kappa|^2$. The integrability
statements are unchanged.
\end{proof}

\begin{theorem}[Lemma~\ref{po:ff}]\label{pfq:ff}
With the first factors of Definition~\ref{pfq:def-ff} and the set $I$ of choices of \S\ref{pte:ss-third} ($\#I\le e^b$ by
Lemma~\ref{plf:I}: the choices there and in Lemma~\ref{ptf:H} are the same data), for every $i\in I$, every twist $\vartheta$
and every $h\in\auU(V'')(\auA_f)$:
\begin{enumerate}
\item[(I1)] $\plgb h{\varphi^\natural_{\vartheta,i}(h)}\le1$ at every $\iota$ with $\ell(\iota)\nmid\auNm Q_\vartheta$ (the gauge of
\S\ref{pl:ss-norm}); that is, $\varphi^\natural_{\vartheta,i}$ is integral at $\iota$;
\item[(I2)] $\|\sigma\varphi^\natural_{1,i}(h)\|^2\le\exp(\poAA)$ for every $\sigma\in\auAut(\C)$.
\end{enumerate}
\end{theorem}
\begin{proof}
\emph{Expansion.} By Lemmas~\ref{pfq:proj} and~\ref{pfq:factor}, and expanding the class sum $\mathrm a_{\chi'}$ (\S\ref{pte:ss-carrier},
a class sum over $K_{I_m}$, which for the level $K_{I_m}$ carries no index factor),
\begin{equation}\label{pfq:eq-expand}
\varphi^\natural_{\vartheta,i}(h)=H^{\mathrm c}\,\pi_{\mu^+}\Bigl(\aubar{\kappa_1a(h)}\Bigr),\qquad
\kappa_1a_U(h)=\sum_{h'}\frac{n'_{\mathrm{Aut}}}{s(h')}\,A(h')\,B(h)\,C_U(h',h),
\end{equation}
the sum over representatives $h'$ of $\Gamma\backslash\auU(I_{3m})(\auA_f)/K_{I_{3m}}$, $s(h')$ the order of the stabiliser, and
\[
A(h')=\sum_{h_1}\frac{n'_{\mathrm{Aut}}}{s(h_1)}\,c_{m,3m,0}\,\theta^{I_m,I_{3m}}_{1,\chi'}(h_1,h';[1,1]),\qquad
B(h)=\aubar{c_{j,N'',0}\,k_1(h)},
\]
\[
C_U(h',h)=c_{12m+1-j,N'',e}\;\theta^{\pfqW,V''}_{\nu'',\nu''}\bigl((h',1,h^\flat_i),h;[Q^{Y_i}_U,\Psi^{\mathrm{add}}_i]\bigr),
\]
$h_1$ running over representatives of $\Gamma\backslash\auU(I_m)(\auA_f)/K_{I_m}$. The factor $\epsilon_j$ of $\kappa_1$ cancels the
$\epsilon_j^{-1}$ of Lemma~\ref{pfq:factor}. The stabiliser orders divide $n'_{\mathrm{Aut}}$ (Remark~\ref{plf:rem-classes}), so the
coefficients $n'_{\mathrm{Aut}}/s$ are positive integers at most $n'_{\mathrm{Aut}}$. The kernels are those of Lemma~\ref{plf:KS}:
$\theta^{I_m,I_{3m}}_{1,\chi'}(h_1,h';[1,1])=\plfK^{(E^m,E^{3m})}_{1,\chi'}[1,1](h_1,h')$ (trivial modulus), $k_1(h)=\plfK^{(E^j,V'')}_{1,1}
[\Psi^{(j)},1](1,h)$ and $\theta^{\pfqW,V''}_{\nu'',\nu''}(x,h;[Q,\Psi^{\mathrm{add}}_i])=\plfK^{(\pfqW,V'')}_{\nu'',\nu''}[\Psi^{\mathrm{add}}_i,Q](x,h)$
(modulus $\pfqm$), all for standard definite spaces; the characters $1,\chi',\nu''$ have algebraic weights and tame
ramification at split places only.

\emph{(I1).} Let $\ell=\ell(\iota)\nmid\auNm Q_\vartheta$, and write $\aubar\iota=\iota\circ c$.
\emph{Reduction to $h$ integral above $\ell$.} By Corollary~\ref{pl:wacor}(a) there is $\gamma\in\auU(V'')(E^+)$ with $\gamma h$
integral above $\ell$. By Proposition~\ref{pfq:exchange}, $\varphi^\natural_{\vartheta,i}(\gamma h)=\gamma\,\varphi^\natural_{\vartheta,i}(h)$, so by (G1) of
\S\ref{pl:ss-norm}, $\plgb h{\varphi^\natural(h)}=\plgb{\gamma h}{\varphi^\natural(\gamma h)}$. So we may assume $h$ integral above $\ell$;
then $\plgb h\xi=\|\iota(\xi)\|$ by (G2), the largest absolute value of a coordinate of $\iota(\xi)$ in the basis of $W''(\Z)$.
\emph{Reduction to the coefficients.} By \S\ref{pfq:ss-hook} it suffices that all tensor coordinates of $\iota(\varphi^\natural(h))$ lie
in $\pao=\{|\cdot|_\ell\le1\}$. By Lemma~\ref{pfq:hook} with $R=\iota^{-1}(\pao)$ and \eqref{pfq:eq-expand}, it suffices that
$\aubar{\kappa_1a_U(h)}\in\iota^{-1}(\pao)$, i.e.\ $|\kappa_1a_U(h)|_{\aubar\iota}\le1$, for every filling $U$.
\emph{The terms.} By Remark~\ref{pfq:rem-classes} and Corollary~\ref{pl:wacor}(a) we may take the representatives $h'$ and $h_1$
integral above $\ell$. (a) $|c_{m,3m,0}\theta^{I_m,I_{3m}}_{1,\chi'}(h_1,h';[1,1])|_{\aubar\iota}\le1$ by (KS1) with the trivial modulus
and $q=1$; so $|A(h')|_{\aubar\iota}\le1$. (b) $|c_{j,N'',0}k_1(h)|_\iota\le1$ by (KS1) with $q=\auNm Q_\vartheta$ ($qM^2\in\pfqm$, as
$\pi_M\mid M$; the datum takes the values $0,1$); hence $|B(h)|_{\aubar\iota}\le1$. (c) Choose $\gamma^\flat\in\auU(W^\flat)(E^+)$ with
$\gamma^\flat h^\flat_i$ integral above $\ell$ (Corollary~\ref{pl:wacor}(a)), and put $B=\mathrm{diag}(1,1,\gamma^\flat)$. By
Lemma~\ref{pfq:move}, $C_U(h',h)=c_{12m+1-j,N'',e}\theta^{\pfqW,V''}_{\nu'',\nu''}((h',1,\gamma^\flat h^\flat_i),h;[Q^{\gamma^\flat Y_i}_U,
\Psi^{\mathrm{add}}_i])$. By Lemma~\ref{au:shortbasis}, $h^{\flat,-1}_iY_i$ is integral, so $\gamma^\flat Y_i=(\gamma^\flat h^\flat_i)
(h^{\flat,-1}_iY_i)$ is integral at the places above $\ell$; its entries are elements of $E$ integral at every prime above $\ell$,
so all $\tau(\gamma^\flat Y_i)_{ab}$ are integral at $\aubar\iota$, and by (Q3) the coefficients of $Q^{\gamma^\flat Y_i}_U$ are
integral at $\aubar\iota$. The values of $\Psi^{\mathrm{add}}_i$ are sums of roots of unity, $Q^{\gamma^\flat Y_i}_U$ is homogeneous of
degree $e$ (Q2), and the points $(h',1,\gamma^\flat h^\flat_i)$ and $h$ are integral above $\ell$. So (KS1) with $q=\auNm Q_\vartheta$ at
$\aubar\iota$ gives $|C_U(h',h)|_{\aubar\iota}\le1$. By \eqref{pfq:eq-expand} and the ultrametric inequality,
$|\kappa_1a_U(h)|_{\aubar\iota}\le1$.

\emph{(I2).} Let $\vartheta=1$ (so $\pfqm=\pi_M^2$), $\sigma\in\auAut(\C)$ and $\sigma'=\sigma\circ c$. The norm on $W''$ is the
restriction of the norm of $\poT_{\mu^+}$ (\S\ref{pfq:ss-hook}); $\pi_{\mu^+}$ has rational coefficients, so it commutes with
$\sigma$ applied to the coordinates, and has norm $1$. By \eqref{pfq:eq-expand},
\[
\|\sigma\varphi^\natural_{1,i}(h)\|\le H^{\mathrm c}\,\|\sigma'(\kappa_1a(h))\|\le H^{\mathrm c}\sum_{h'}n'_{\mathrm{Aut}}\,|\sigma'A(h')|\,
|\sigma'B(h)|\,\bigl\|\bigl(\sigma'C_U(h',h)\bigr)_U\bigr\| .
\]
By (KS2) with $A=B=1$ (ranks $m,3m,j,N''\le\Delta$, degree $0$, trivial modulus resp.\ $\pi_M^2$): $|\sigma'A(h')|\le n'_{\mathrm{Aut}}
e^{2x^{26}}e^{x_1}$ (Lemma~\ref{plf:classes} for the number of classes $h_1$), and $|\sigma'B(h)|=|(\sigma'c)(c_{j,N'',0}k_1(h))|\le
e^{x_1}$. For the vector $C(h',h)=(C_U(h',h))_U$ we use duality: $\|\sigma'C\|=\sup|\sum_U\phi_U\,\sigma'C_U|$ over the $\phi$ with
$\sum|\phi_U|^2\le1$, and $\sum_U\phi_U\,\sigma'C_U=\sigma'\bigl(c_{12m+1-j,N'',e}\theta^{\pfqW,V''}_{\nu'',\nu''}((h',1,h^\flat_i),h;
[Q_\phi,\Psi^{\mathrm{add}}_i])\bigr)$ with $Q_\phi=\sum_U\sigma'^{-1}(\phi_U)Q^{Y_i}_U$, homogeneous of degree $e\le M$ (Q2). Its
$\sigma'$-conjugate $Q_\phi^{\sigma'}=\sum_U\phi_U(Q^{Y_i}_U)^{\sigma'}$ satisfies, by the Cauchy--Schwarz inequality and (Q4),
$|Q^{\sigma'}_\phi(z)|\le M^{9(12m+1)\mathsf q/2}$ whenever all $|z_c|\le M^4$; the hypothesis of (Q4) holds because all conjugates of
the entries of $Y_i$ are at most $C_{\mathrm{sb}}\le x^{103}\le M^4$ (Lemma~\ref{au:shortbasis}; the bound on $C_{\mathrm{sb}}$ is that of
the proof of Lemma~\ref{plf:classes}, rank $6m+1\le2n_m$). The values of $\Psi^{\mathrm{add}}_i$ and their conjugates are at most
$\#\auGL_{3m}(\auO_E/\pi_M^2)\le M^{18fm^2}$. Since $12m+1-j,N''\le\Delta$ and $M^2\in\pi_M^2\auO_E$, (KS2) gives
$\|\sigma'C(h',h)\|\le e^{x_1}M^{18fm^2}M^{9(12m+1)\mathsf q/2}$. With at most $e^{2x^{26}}$ classes $h'$ (Lemma~\ref{plf:classes}),
\[
\|\sigma\varphi^\natural_{1,i}(h)\|\le H^{\mathrm c}\,n'^{\,2}_{\mathrm{Aut}}\,e^{4x^{26}}\,e^{3x_1}\,M^{18fm^2}\,M^{9(12m+1)\mathsf q/2}.
\]
By Lemma~\ref{pfq:hook}, $H^{\mathrm c}\le e^{3e}$, the base being the integer $e$ of (Q2), so $\log H^{\mathrm c}\le3e\log e$. Taking logarithms and doubling: $2\log H^{\mathrm c}\le6e\log e\le6\cdot11m\cdot4fb\cdot23\log b$ (as
$e\le(10m+1)\mathsf q\le M\le b^{22}$), $4\log n'_{\mathrm{Aut}}\le4x^{40}$, $8x^{26}+6x_1\le7x_1$, $36fm^2\log M\le x^{14}$, and
$9(12m+1)\mathsf q\log M\le9\cdot13m\cdot4fb\cdot23\log b$. The total is at most $1.7\cdot10^4fmb\log b+8x_1\le fmb(\log b)^2=\poAA$, since
$\log b\ge2\cdot10^4+1$ and $8x_1\le fmb\log b$ (Lemma~\ref{an:lem-sizes}, \eqref{eq:au-scales}).
\end{proof}

\begin{remark}[Uses in \S\ref{po:sec}]\label{pfq:rem-uses}
In \S\ref{po:ss-glue} and the proof of Proposition~\ref{au:onelift}, the first factors enter only through: their type and level
(the summands of $\poglue_\vartheta$ do not depend on the representatives) --- Proposition~\ref{pfq:exchange}; (G1) --- the identity
$c'_\vartheta\Xi_{\vartheta,i}=F_i$, Proposition~\ref{pfq:exchange}; (G3) --- Proposition~\ref{po:periods}, valid for every vector $u''$; (G4) ---
(I1), Theorem~\ref{pfq:ff}; (G5) and Step~6 --- (I2) and $\#I\le e^b$. Lemmas~\ref{po:lattice} and~\ref{po:slice} do not involve the
first factors. The constant $C_{\mathrm x}$ is replaced by $C_{\mathrm x}/|\kappa|^2$, which changes only the constant $C_j$ of Step~1
and keeps $C_{\mathrm x}>0$ in (G2).
\end{remark}

\subsubsection{Tame parts at the chosen places}\label{pfq:ss-tame}

Recall from Definition~\ref{an:def-data} and \S\ref{an:ssec-param}: $E=K(\zeta)$, $\zeta$ a primitive $q_1$-th root of unity,
$4\mid q_1$, $q_1d\mid q_2=(q_0d)^{s+6}$ with $2^s=m^{22}$; $\kappa(\tau)\in\{b,xb/x_1\}$ and $q_1x_1\mid b$ (D4); the level primes
$p_0,p_1\equiv1\pmod{q_2^2}$; $\Delta/2=r^8$; and $\auPl$ is identified with $\operatorname{Gal}(E/K)\cong(\Z/q_1)^\times$, the
embeddings extending a fixed $K\subset\C$. For a finite place $v$ of $E^+$ split as $w\aubar w$ we identify
$\auU(1)(E^+_v)$ with $E_w^\times$ by the component at $w$.

\begin{lemma}[Roots of unity]\label{pfq:roots}
\begin{enumerate}
\item[(a)] Every root of unity of $E$ is a $q_1$-th root of unity: $\mu_E=\mu_{q_1}$.
\item[(b)] $\prod_{\tau\in\auPl}\tau(\xi)=1$ for every $\xi\in\mu_E$.
\item[(c)] $q_1$ divides $\kappa(\tau)$ for every $\tau$; hence $\prod_\tau\tau(\xi)^{c\kappa(\tau)+c'r^8}=1$ for $\xi\in\mu_E$ and all
integers $c,c'$.
\item[(d)] At every chosen place $w$, $\auF_w=\auF_{p_i}$, $q_2^2\mid\#\auF_w^\times$, and the reduction of every $\xi\in\mu_E$ is a
$q_2$-th power in $\auF_w^\times$.
\end{enumerate}
\end{lemma}
\begin{proof}
(a) $K=\Q(\sqrt{-d})\subseteq\Q(\zeta_d)$ ($d\equiv3\pmod4$ is prime), so $E\subseteq\Q(\zeta_{q_1d})$, whose roots of unity are
$\mu_{q_1d}=\mu_{q_1}\times\mu_d$ ($q_1d$ is even). If $E$ contained a root of unity of order divisible by $d$ it would contain
$\zeta_d$, hence the compositum of $\Q(\zeta_d)$ and $\Q(\zeta_{q_1})$, of degree $(d-1)\varphi(q_1)>2\varphi(q_1)=[E:\Q]$
($d\ge7$, coprime conductors); impossible. So $\mu_E\subseteq\mu_{q_1}\subseteq E$. (b) Write $\xi=\zeta^n$; the embeddings in
$\auPl$ send $\zeta$ to $\zeta_0^a$, $a\in(\Z/q_1)^\times$, for a fixed primitive $\zeta_0$, so $\prod_\tau\tau(\xi)=\zeta_0^{n\sum a}$, and
$\sum_{a\in(\Z/q_1)^\times}a\equiv0\pmod{q_1}$ (pair $a$ with $q_1-a$; $a\ne q_1-a$ because $q_1/2$ is even). (c) $q_1\mid b$ by (D4),
and $xb/x_1=x\,q_1(b/(q_1x_1))$; then $\tau(\xi)^{c\kappa(\tau)}=1$ by (a), and $\prod_\tau\tau(\xi)^{c'r^8}=1$ by (b). (d) $p_i\equiv1$
modulo $q_2^2$, hence modulo $q_1d$, so $p_i$ splits completely in $\Q(\zeta_{q_1d})\supseteq E$ and $\auF_w=\auF_{p_i}$. The group
$\auF_w^\times$ is cyclic of order $p_i-1$, divisible by $q_2^2$; an element $y$ with $y^{q_1}=1$ is a $q_2$-th power because
$y^{(p_i-1)/q_2}=1$, as $q_1\mid q_2\mid(p_i-1)/q_2$.
\end{proof}

\begin{lemma}[Characters of $\auU(1)$ with prescribed weight and tame parts]\label{pfq:mu}
Let $S$ be a finite set of primes of $E$, no two above the same place of $E^+$ and each above a place of $E^+$ split in $E$; let
$(c_\tau)_{\tau\in\auPl}$ be integers and, for $w\in S$, $\psi_w$ a character of $\auF_w^\times$, such that
\begin{equation}\label{pfq:eq-compat}
\prod_{\tau\in\auPl}\tau(\xi)^{c_\tau}\cdot\prod_{w\in S}\psi_w(\xi\bmod w)=1\qquad\text{for every }\xi\in\mu_E .
\end{equation}
Then there is a unitary character $\mu$ of $\auU(1)(E^+)\backslash\auU(1)(\auA_{E^+})$ with
\[
\mu(u)=\prod_\tau\tau(u_\tau)^{c_\tau}\cdot\prod_{w\in S}\psi_w(u_w\bmod w)\qquad\text{for all }u\in\mathsf U=\auU(1)(E^+\otimes\R)\times\prod_v
\auU(1)(\auO_v),
\]
that is: $\mu_\tau(u)=\tau(u)^{c_\tau}$ at the infinite places, $\mu$ has tame part $\psi_w$ at $w\in S$, and is unramified at all
other finite places. Conversely \eqref{pfq:eq-compat} is necessary for the existence of such $\mu$. If the reduction
$\mu_E\to\auF_{w_0}^\times$ is injective for some $w_0\in S$, then for every $(c_\tau)$ and every choice of $\psi_w$, $w\ne w_0$, there is $\psi_{w_0}$ satisfying \eqref{pfq:eq-compat}.
\end{lemma}
\begin{proof}
$\mathsf U$ is an open subgroup of $\auU(1)(\auA_{E^+})$ and $\auU(1)(E^+)\cap\mathsf U$ consists of the $\xi\in E$ with $\xi\aubar\xi=1$
that are integral at every finite place; all their conjugates have absolute value $1$, so they are roots of unity (Kronecker), and
$\auU(1)(E^+)\cap\mathsf U=\mu_E$. The formula defines a continuous unitary character $\mu_{\mathsf U}$ of $\mathsf U$ (at $w\in S$,
$\auU(1)(\auO_v)=\auO_w^\times\to\auF_w^\times$), and \eqref{pfq:eq-compat} says $\mu_{\mathsf U}(\xi)=1$ for $\xi\in\mu_E$ (a rational point
$\xi$ has components $\tau(\xi)$ at the infinite places and $\xi$ at the finite ones). Hence $\gamma u\mapsto\mu_{\mathsf U}(u)$ is a
well defined character of the open subgroup $\auU(1)(E^+)\mathsf U$, trivial on $\auU(1)(E^+)$. A unitary character of an open
subgroup of a locally compact abelian group extends to a unitary character of the group (Pontryagin duality: the restriction map
of dual groups is onto for closed subgroups); the extension is trivial on $\auU(1)(E^+)$ and is the required $\mu$. Necessity: if
$\mu$ exists, $1=\mu(\xi)=\mu_{\mathsf U}(\xi)$ for $\xi\in\mu_E$. The last assertion: $\xi\mapsto\prod_\tau\tau(\xi)^{-c_\tau}\prod_{w\ne w_0}
\psi_w(\xi)^{-1}$ is a character of the subgroup $\mu_E$ of the cyclic group $\auF_{w_0}^\times$, and extends.
\end{proof}

\begin{definition}[The characters attached to the data]\label{pfq:def-chars}
\begin{enumerate}
\item[(i)] For each chosen place $w$ fix a character $\omega_w$ of $\auF_w^\times$ of order exactly $q_2$ (it exists by
Lemma~\ref{pfq:roots}(d)); for the other finite places put $\omega_w=1$. Write $\omega=(\omega_w)_w$.
\item[(ii)] $\chi_M$ is a character of $\auU(1)$ of weight $0$ with tame part $\omega$: Lemma~\ref{pfq:mu} with $S$ the chosen places,
$c_\tau=0$, $\psi_w=\omega_w$. Condition \eqref{pfq:eq-compat} holds factor by factor: $\omega_w(\xi\bmod w)=1$ because $\xi\bmod w$ is a
$q_2$-th power (Lemma~\ref{pfq:roots}(d)) and $\omega_w^{q_2}=1$. We write $\chi_{M,w}=\omega_w$ for its tame parts.
\item[(iii)] $\mu^\flat$ is a character of $\auU(1)$ of weight $-\kappa_0=-(\kappa+\frac\Delta2)$, unramified at all finite places:
Lemma~\ref{pfq:mu} with $S=\emptyset$, $c_\tau=-\kappa(\tau)-r^8$; \eqref{pfq:eq-compat} is Lemma~\ref{pfq:roots}(c).
\item[(iv)] $\mu^0_j=\mu^\flat\,(\upsilon^{-\Delta}\chi_M)^j$ for $1\le j\le m$, $\upsilon$ restricted to $\auU(1)$.
\end{enumerate}
\end{definition}

\begin{proposition}[Properties of the tame parts]\label{pfq:tame}
\begin{enumerate}
\item[(T1)] $\mu^0_j$ has weight $-\kappa_j$, tame part $\omega_w^j$ at every chosen place $w$, and is unramified at all other finite
places; so it is a character as required in \S\ref{ss:au-spaces}. The same characters are obtained, up to a character of weight $0$
unramified everywhere, from Lemma~\ref{pfq:mu} with $c_\tau=-\kappa_j(\tau)$ and $\psi_w=\omega_w^j$: condition \eqref{pfq:eq-compat}
holds for every $j$.
\item[(T2)] $\omega_w$ has order $q_2>8m$ at every chosen place; hence $\omega_w^{\,n}\ne1$ for $1\le|n|\le8m$, and in particular
$\omega_w^{\pm j}\ne1$ for $1\le j\le m$.
\item[(T3)] For every twist $\vartheta$, the tame part of $\mu=\mu^0_j\vartheta$ at a chosen place $v$ is $\psi_v=\omega_v^{\,j}$, and
$\psi_v^{q_2}=1$. This is the property used in (F2) of Lemma~\ref{ppe:three} (Remark~\ref{ppe:rem-T}).
\item[(T4)] The character $\chi_0$ of \S\ref{po:ss-hecke} is $\chi_M$, with tame part $\omega$ trivial off $\mathsf M$ and of order
$>8m$ at $\mathsf M$; the tame part of $\mu^0_j$ is $\omega^j$. The characters $\nu_\ast=\mu^0_j\upsilon^{n'}$, $\nu_\pm$,
$\chi'=\chi_M^j\vartheta$, $\nu''_\vartheta$ of \S\ref{po:ss-hecke}, \S\ref{pte:ss-third} have the weights and tame parts stated there,
and $\nu^\flat=\nu''_\vartheta\upsilon^m/\chi'=\mu^\flat\upsilon^{-j\Delta+j+\lfloor(t+k-1)/2\rfloor+m}$ is independent of $\vartheta$ with trivial tame part. At every
prime of $\pfqm$ the tame parts of $\nu''_\vartheta$ and $\chi'$ coincide ($\omega^j_w$ at $\mathsf M$, $\psi_w$ at $Q_\vartheta$).
\end{enumerate}
\end{proposition}
\begin{proof}
Weights add and tame parts multiply under products of characters (the defining formula on $\mathsf U$ is multiplicative), and
$\upsilon|_{\auU(1)}$ has weight $1$ and is unramified at the finite places: $\upsilon$ is unramified at every
finite place of $E$ (\S\ref{pk:ss-conv}: $\upsilon(x)=\prod_\tau x_\tau/|x_\tau|$ when the finite part of $x$ is a unit), and
$\auU(1)(\auO_v)$ consists of units.
(T1) The weight of $\mu^0_j$ is $-\kappa-\frac\Delta2-j\Delta=-\kappa_j$, and its tame part is $\omega^j$. For the second statement,
\eqref{pfq:eq-compat} reads $\prod_\tau\tau(\xi)^{-\kappa(\tau)-(2j+1)r^8}\prod_w\omega_w(\xi)^j=1$, which holds by Lemma~\ref{pfq:roots}(c)
and (ii) of Definition~\ref{pfq:def-chars}; two characters with the same weight and tame parts differ by a character trivial on
$\mathsf U$. (T2) $q_2=(q_0d)^{s+6}\ge2^s=m^{22}>8m$. (T3) $\vartheta$ is unramified at the chosen places (its set $Q$ is prime to
$M$), so the tame part is that of $\mu^0_j$, and $(\omega_v^j)^{q_2}=1$. (T4) The definition of $\chi_0$ in \S\ref{po:ss-hecke} asks
exactly for weight $0$ and a tame part trivial off $\mathsf M$ of order $>8m$ on $\mathsf M$, with the tame part of $\mu^0_j$ equal to
$\omega^j$; (ii), (T1), (T2) give this with $\chi_0=\chi_M$. The other characters are products of $\mu^\flat$, $\chi_M$, $\upsilon$ and
$\vartheta$, and their weights and tame parts follow by addition and multiplication; for $\nu^\flat$ the factors $\chi_M^j\vartheta$
cancel. Finally the tame part of $\nu''_\vartheta=\mu^0_j\vartheta\upsilon^{\cdots}$ is $\omega^j\psi$, and so is that of $\chi'$.
\end{proof}

\begin{remark}[Every use of Lemma~\ref{pz:mu}]\label{pfq:rem-pzmu}
Lemma~\ref{pz:mu}, as stated in \S\ref{pz:sec}, is Lemma~\ref{pfq:mu} (same statement and proof; the weaker form with a
tame part prescribed at one chosen place $w_0$ only is the case $S=\{w_0\}$, $\psi_{w_0}$ given by the last assertion; $\mu_E\to\auF_{w_0}^\times$ is injective because $p_0>b^{10}$ does not divide $\#\mu_E=q_1$). Lemma~\ref{pz:mu} is used: (1) in \S\ref{ss:au-spaces}, for the
existence of $\mu^0_j$ --- now Definition~\ref{pfq:def-chars}(iv) and (T1), which give in addition the tame parts $\omega^j$ needed
in \S\ref{po:ss-hecke} and in Lemma~\ref{ppe:three}; (2) in the summary at the end of \S\ref{pz:sec} (types, characters, levels) ---
unchanged; (3) in Remark~\ref{ppe:rem-T} --- now (T3); (4) in \S\ref{ptf:ss-close}, for the meaning of ``weight''
($\mu_\tau(u)=\tau(u)^{c_\tau}$) --- the same convention in Lemma~\ref{pfq:mu}. The existence of twists with prescribed tame parts
trivial on $\mu_E$ (\S\ref{ss:au-spaces}) is Lemma~\ref{pfq:mu} with $c_\tau=0$ and $S=Q$.
\end{remark}


%% ===== proofs/C-small.tex =====
% proofs/C-small.tex
% Proposition po:small of proofs/C-onelift.tex (a test form whose pairing has least valuation):
% the exact statement, its use in po:mainheight and au:hecke, and the proof of its core, the
% sharpness of the count (part (i) of the note to po:small). To be \input after A-divisibility,
% A-final and C-onelift. Labels of this file start with psm:. Only amsmath/amssymb/amsthm are used.
\providecommand{\auO}{\mathcal O}
\providecommand{\auU}{\mathrm U}
\providecommand{\auAut}{\operatorname{Aut}}
\providecommand{\auNm}{\operatorname{Nm}}
\providecommand{\auA}{\mathbb A}
\providecommand{\auPl}{\Phi}
\providecommand{\auPa}{\Phi_a}
\providecommand{\auPb}{\Phi_b}
\providecommand{\aubud}{\mathrm{bud}}
\providecommand{\auhdir}{h_{\mathrm{dir}}}
\providecommand{\ausinf}{s_\infty}
\providecommand{\ausfin}{s_{\mathrm{fin}}}
\providecommand{\aubar}[1]{\overline{#1}}
\providecommand{\pao}{\mathfrak o}
\providecommand{\pam}{\mathfrak m}
\providecommand{\paCp}{\C_p}
\providecommand{\popr}{\operatorname{pr}}
\providecommand{\psvac}{\mathbf v}
\providecommand{\psrho}{\varrho}
\providecommand{\ancl}{\operatorname{cl}}
% macros of this file (prefix \psm)
\providecommand{\psmval}{\operatorname{val}}
\providecommand{\psmR}{R}
\providecommand{\psmns}{n_\sharp}
\providecommand{\psmpr}{\mathrm{pr}}
\providecommand{\psmnot}[1]{\par\smallskip{\raggedright\noindent\textbf{[Not proved here:} #1\textbf{]}\par}\smallskip}

\subsection{The test form of least valuation: sharpness of the count}\label{psm:sec}

This subsection concerns Proposition~\ref{po:small}. In \S\ref{psm:ss-statement} we fix its exact
statement and check that Proposition~\ref{po:mainheight} and Proposition~\ref{au:hecke} follow from
it in the direction stated. The rest of the subsection proves the core of the proposition, the
\emph{sharpness of the count} (part (i) of the note to Proposition~\ref{po:small}): at every prime
some Taylor coefficient of the right degrees of a theta function with complex multiplication has
valuation at most the prime sum $F(p)$, up to a factor whose conjugates are controlled uniformly
in $p$. The $p$-adic part of this (Theorem~\ref{psm:twelve}) is proved in full in
\S\S\ref{psm:ss-least}--\ref{psm:ss-twelve}; its transfer to the indicator data
(Theorem~\ref{psm:attained}) is proved in \S\ref{psm:ss-attained} from four inputs on the
Heisenberg model, proved in \S\ref{pst:sec}. Parts (ii)--(iv) of the note (the test form itself)
are proved in \S\ref{psu:sec}, which also treats Proposition~\ref{po:arch}.

\subsubsection{The exact statement}\label{psm:ss-statement}

Recall from Definition~\ref{def:au-hmain} the local prime sum
$s_p(y,\sigma)=v(\pi_p)\sum_{e\in I_p}\bigl(y/p^e-\kappa(u(-p^e)\sigma)\bigr)^+$, and put
\begin{equation}\label{psm:eq-F}
F(p)=\sum_{\sigma\in\auPl}s_p\bigl(\kappa(\sigma),\sigma\bigr)\qquad(p\text{ prime}),
\end{equation}
the prime sum evaluated \emph{at the weights $\kappa$}, not at the degrees of differentiation.
Thus $F(p)=0$ if $p$ splits in $K$ or $p\ge x_1$, and $\sum_p2\log p\cdot F(p)=\sum_\sigma
\ausfin(\kappa(\sigma),\sigma)$ by Lemma~\ref{au:gains}(a).

\begin{proposition}[Proposition~\ref{po:small}, exact form]\label{psm:small}
Let $p$ be a main pair for $j$ and $q_2$ a critical pair with empty row on a block of size $4$
\textup{(}in the application $q_2=(\upsilon^{-2}\chi_0,2)$, flat and ramified\textup{)}. There is
$c_0\in\C^\times$ with
\[
|\sigma c_0|\le\exp\bigl(fmb/k+3x^{35}\bigr)\qquad(\sigma\in\auAut(\C))
\]
such that for every prime $\ell$, every homomorphism $\jmath\colon\C\to\C_\ell$ and every twist
$\vartheta$ with $\ell\nmid\auNm Q_\vartheta$ there is a test form $\Xi$ at $\ell$ for
$(p,\vartheta)$ with
\begin{equation}\label{psm:eq-small}
\bigl|c_0\,\rho_{q_2}(\vartheta)\bigr|_\jmath\cdot\ell^{-F(\ell)}\ \le\ \bigl|\langle\Xi,
\varphi_{p,\vartheta}\rangle\bigr|_\jmath ,
\end{equation}
where $\varphi_{p,\vartheta}$ is the normalised lift and $F$ is \eqref{psm:eq-F}.
\end{proposition}

\emph{Comparison with the printed statement.} The statement printed as Proposition~\ref{po:small}
is correct as it stands: the function there, $F(\ell)=\sum_{\sigma\in\auPl}s_\ell(\kappa(\sigma),
\sigma)$, is \eqref{psm:eq-F}; the inequality is a \emph{lower} bound for the pairing (an upper
bound for its valuation); the constant $c_0$ does not depend on $\ell$, $\jmath$ or $\vartheta$; and
the twist enters only through the ratio $\rho_{q_2}(\vartheta)$ of the auxiliary pair. The only
hypothesis on $q_2$ that is used in the proof is that it has an empty row on a block of size $4$;
being ramified is used only afterwards, through (R2), in Proposition~\ref{po:mainheight}.

\emph{The deductions.} (1) Proposition~\ref{po:mainheight}: its proof uses
Proposition~\ref{psm:small} exactly in the form \eqref{psm:eq-small}, for $\jmath=\iota$ and
$\jmath=\aubar\iota$, multiplies the two inequalities into \eqref{po:eq-two}, and needs
$\sum_{\ell\in T}2F(\ell)\log\ell\le\sum_\sigma\ausfin(\kappa(\sigma),\sigma)$; this holds (with
equality when $T$ contains all primes $<x_1$ that do not split) because every term of
$\sum_p2\log p\cdot F(p)$ is $\ge0$ and the sum is $\sum_\sigma\ausfin(\kappa(\sigma),\sigma)$.
The fifth factor in that proof has multiplier $c_0\aubar{c_0}$, whose conjugates are at most
$\exp(2fmb/k+6x^{35})$, and its absolute value at $\iota$ is at most $\ell^{2F(\ell)}$; adding the
heights, the total is
\[
\bigl(\aubud-\textstyle\sum_\tau\ausinf(\kappa(\tau))\bigr)+\aubud+2x_1+2\bigl(fmb/k+3x^{35}
\bigr)+\textstyle\sum_\sigma\ausfin(\kappa(\sigma),\sigma)\ \le\ 6\,\aubud-\auhdir ,
\]
as stated. The sign is the right one: the prime sums $\ausfin$ enter with the sign $+$ in the
height of the inverse ratio, so that they cancel the $-\ausfin$ inside $-\auhdir=-\sum_\sigma
(\ausinf(\kappa(\sigma))-\ausfin(\kappa(\sigma),\sigma))$ and leave $-\sum_\sigma\ausinf
(\kappa(\sigma))$, which is cancelled by Proposition~\ref{po:arch}. A version of
Proposition~\ref{psm:small} with $F(\ell)$ replaced by a larger quantity would give a weaker bound
($6\aubud-\auhdir$ plus the excess), and with an upper bound for the pairing instead of a lower
bound the argument would not apply at all. (2) Proposition~\ref{au:hecke} is deduced from
Proposition~\ref{po:mainheight} and (R2) by subadditivity of heights; it uses nothing further from
Proposition~\ref{psm:small}. The constant $fmb/k+3x^{35}$ is bounded by
$\aubud$ before the heights are added (Lemma~\ref{psu:sizes}(N6)), which gives the same final bound.

\emph{Where the exponent $F(\ell)$ comes from.} In the proof of Proposition~\ref{psm:small} the
number $c_0$ is a product $t\cdot I_{q_2}(1)\cdot L$ (note to Proposition~\ref{po:small}), the
pairing $\langle\Xi,\varphi_{p,\vartheta}\rangle$ is, up to $\rho_{q_2}(\vartheta)$, a class number
and the two other factors, an average over characters of Taylor coefficients of a theta function in
$12$ variables, and the whole of the exponent $F(\ell)$ comes from one statement on these Taylor
coefficients, Theorem~\ref{psm:attained} below, applied with the degrees $l_\tau=c(\tau)-z$ of the
main pair (which lie in $[\kappa(\tau),\kappa(\tau)+x^{42}]$). That theorem is the converse of the
upper bound side of the count (Proposition~\ref{au:count}, proved in \S\ref{pa:ss-count}): there,
every jet along some $Q_{T,y}$ or along an entry of the Schur matrix has valuation \emph{at least} the predicted divisibility; here,
some jet has valuation \emph{at most} $F(p)$ plus a loss $2\sum_\tau v_p(B_\tau)$ that is absorbed in
the number $t$ uniformly in $p$.

\subsubsection{Least valuations of coefficients of theta products}\label{psm:ss-least}

Throughout \S\S\ref{psm:ss-least}--\ref{psm:ss-twelve}, $p$ is any prime (split, inert or $d$),
$\C\to\C_p$ is fixed, $v$ is the valuation of $\C_p$ with $v(p)=1$, $\pao$ its valuation ring and
$\pam$ its maximal ideal. Let $(\Lambda,t,a_1)$ be a normalised lattice (\S\ref{pf:ss-adm}) and
$\theta$, $x$, $y$, $g=-\log\theta$ the series of \S\ref{pa:ss-curve}, read in $\C_p[[z]]$. By
Lemma~\ref{pf:two} the curve $\mathcal E\colon y^2+a_1xy=x^3+x$ has good reduction at $p$ and $P_0$
does not reduce to the origin; so the construction of the ring
\[
\psmR=\pao[[y]]\subset\C_p[[z]]
\]
in \S\ref{pa:ss-curve} applies at every prime: $\psmR$ is closed under coefficientwise limits,
contains $x$, is stable under $d/dz$, and is the completion of the smooth model of $\mathcal E$
along $P_0$. Since $y=z+O(z^2)$, a series in $\psmR$ has integral coefficient in degree $0$ and
$1+y\psmR$ is a multiplicative group.

Let $X$ be a finite set of variables $z=(z_x)_{x\in X}$. For a series $F\in\C_p[[z]]$ and
$\nu\in\N^X$ write $[\nu]F$ for the coefficient of $z^\nu$, $|\nu|=\sum_x\nu(x)$ and $\nu!=
\prod_x\nu(x)!$. Let $\mu=(\mu_i)_{i\in J}$ be a finite family of linear forms in $z$ with
coefficients in $\pao$ and $m=(m_i)$ multiplicities. Put
\[
\Theta=\Theta_{\mu,m}=\prod_{i\in J}\theta(\mu_i(z))^{m_i},\qquad
\psmR_\mu=\Bigl\{\textstyle\sum_{e\in\N^J}c_e\prod_iy(\mu_i(z))^{e_i}:\ c_e\in\pao\Bigr\}.
\]
$\psmR_\mu$ is a ring containing $G(\mu_i(z))$ for $G\in\psmR$, and $\Theta(0)=1$.

\begin{definition}[Least valuation on $n$ copies]\label{psm:def-val}
For $n\ge1$, $\nu\in\N^X$ and $r\in\R$ we write $\psmval_n(\nu)\le r$ if there are
$H_1,\dots,H_n\in\psmR_\mu$ and $\alpha_1,\dots,\alpha_n\in\N^X$ with $\sum_c\alpha_c=\nu$ and
\[
\sum_{c=1}^nv\bigl([\alpha_c](H_c\Theta)\bigr)\le r .
\]
Equivalently: some coefficient of degree $\nu$ (after identifying the $n$ copies of the variables)
of the series $\prod_c(H_c\Theta)(z_c)$ in $n$ copies $z_1,\dots,z_n$ of the variables has
valuation at most $r$, the coefficient being a product of one coefficient of each factor.
\end{definition}

\begin{lemma}[Rules]\label{psm:rules}
\begin{enumerate}
\item[(a)] If $\psmval_n(\nu)\le r\le r'$ then $\psmval_n(\nu)\le r'$; and $\psmval_n(0)\le0$.
\item[(b)] If $\psmval_n(\nu)\le r$ and $\psmval_{n'}(\nu')\le r'$ then $\psmval_{n+n'}(\nu+\nu')
\le r+r'$.
\item[(c)] \emph{(Integral substitution.)} Let $H_1,\dots,H_n\in\psmR_\mu$, let $K$ be a finite set,
$A_{c,k}\in\pao$, and $w=(w_{x,k})$ variables. If the coefficient of $w^\gamma$, $\gamma\in
\N^{X\times K}$, in $\prod_{c}(H_c\Theta)\bigl(\sum_kA_{c,k}w_{\cdot,k}\bigr)$ has valuation at
most $r$, then $\psmval_n(\bar\gamma)\le r$, where $\bar\gamma(x)=\sum_k\gamma(x,k)$.
\item[(d)] \emph{(Monomials.)} If $\psmval_n(\nu)\le r$, the $H_c$ in
Definition~\ref{psm:def-val} may be taken to be monomials $\prod_iy(\mu_i(z))^{e_{c,i}}$ with
$\sum_ie_{c,i}\le|\nu|$.
\end{enumerate}
\end{lemma}
\begin{proof}
(a) Clear; for the second claim take $H_c=1$, $\alpha_c=0$, as $\Theta(0)=1$. (b) Concatenate the
two families. (c) The coefficient of $w^\gamma$ in the product is $\sum\prod_c[\beta_c]\bigl((H_c
\Theta)(\sum_kA_{c,k}w_{\cdot,k})\bigr)$ over the decompositions $\gamma=\sum_c\beta_c$; by the
ultrametric inequality one term has valuation at most $r$. For a series $G$ in $z$,
\[
[\beta]\,G\Bigl(\sum_kA_kw_{\cdot,k}\Bigr)=[\bar\beta]G\cdot\prod_x\frac{\bar\beta(x)!}{\prod_k
\beta(x,k)!}\cdot\prod_{x,k}A_k^{\beta(x,k)} ,
\]
an integral multinomial coefficient times integral numbers times $[\bar\beta]G$; so
$v([\bar\beta_c](H_c\Theta))\le v([\beta_c](\cdots))$, and $\alpha_c=\bar\beta_c$ satisfy
$\sum_c\alpha_c=\bar\gamma$. (d) Let $H=\sum_ec_e\prod_iy(\mu_i)^{e_i}\in\psmR_\mu$ and $\alpha\in
\N^X$. Since the $y(\mu_i(z))$ have no constant term, only the $e$ with $|e|\le|\alpha|$ contribute
to $[\alpha](H\Theta)$, and $[\alpha](H\Theta)=\sum_{|e|\le|\alpha|}c_e\,[\alpha](\prod_iy(\mu_i)^{
e_i}\Theta)$ with $c_e\in\pao$. By the ultrametric inequality one of the finitely many monomials
gives a coefficient of valuation at most $v([\alpha](H\Theta))$; and $|\alpha_c|\le|\nu|$.
\end{proof}

\subsubsection{One curve: units, the rate $\gamma_p$, and the leading coefficient}\label{psm:ss-curve}

\begin{lemma}[Isogeny units]\label{psm:units}
\begin{enumerate}
\item[(a)] For every $\alpha\in\auO_K$, $\alpha\ne0$, the series $G(\alpha z)$ lies in $\psmR$ for all
$G\in\psmR$, and $f_\alpha(z)=\theta(\alpha z)\,\theta(z)^{-\auNm\alpha}$ lies in $1+y\psmR$.
\item[(b)] Let $(\delta,s)$ lie in the closure in $\C_p^2$ of $\{(\alpha,\auNm\alpha):\alpha\in
\auO_K\setminus0\}$. Then $G(\delta z)\in\psmR$ for $G\in\psmR$, and $\exp\bigl(s\,g(z)-g(\delta
z)\bigr)\in1+y\psmR$. This applies to $(\delta,\delta^2)$ for every $\delta\in\Z_p$.
\end{enumerate}
\end{lemma}
\begin{proof}
(a) In the uniformisation $z\mapsto Q(z)$ of \S\ref{pf:ss-tors}, $z\mapsto Q(z)-P_0$ is a group
isomorphism, so $z\mapsto\alpha z$ is the morphism $Q\mapsto\alpha(Q-P_0)+P_0$ of $\mathcal E$,
which fixes $P_0=Q(0)$. It extends to the smooth model, which is the N\'eron model of its generic
fibre~\cite[\S1.2, Definition~1, p.~12, and Proposition~8, p.~15]{BLR} (\cite[IV, Corollary~6.3, p.~329, and the N\'eron mapping property, p.~319]{SilvermanAT}), and fixes the section $P_0$; hence it preserves $\psmR$, as in
\eqref{pa:eq-endo}. No congruence condition on $\alpha$ is needed for this. By
Proposition~\ref{pf:pull} with $u=0$ and $\aubar\delta=\alpha$, $\theta(\alpha z)=\theta(z)^{
\auNm\alpha}\Gamma(y(z))$ with $\Gamma\in\pao[[Y]]$, at every prime; comparing constant terms,
$\Gamma(0)=1$, so $f_\alpha=\Gamma(y)\in1+y\psmR$.

(b) Let $(\alpha_n,\auNm\alpha_n)\to(\delta,s)$. The coefficients of $G(\alpha_nz)$ and of
$\exp(\auNm\alpha_n\,g(z)-g(\alpha_nz))=f_{\alpha_n}^{-1}$ are polynomials in $(\alpha_n,
\auNm\alpha_n)$ (with fixed coefficients), and they lie in the closed sets $\psmR$ and $1+y\psmR$
by (a) ($1+y\psmR$ is a closed group). Pass to the limit. For $\delta\in\Z_p$ take
$\alpha_n\in\Z$, $\alpha_n\to\delta$; then $\auNm\alpha_n=\alpha_n^2\to\delta^2$.
\end{proof}

\begin{lemma}[Six copies]\label{psm:six}
There are $\delta_0=1,\delta_1,\dots,\delta_5\in\Z_p$ such that, for every family $(\mu,m)$ as above,
$U=\prod_{c=0}^5\Theta(\delta_cz)$ lies in $\psmR_\mu$, has constant term $1$, and has an inverse in
$\psmR_\mu$. If moreover $v(\varpi)\ge\rho_p$ \textup{(}with $\rho_p$ as in
Lemma~\ref{psm:param}\textup{)}, then $H(\varpi z)$ is integral for every $H\in\psmR_\mu$.
\end{lemma}
\begin{proof}
Every element of $\Z_p$ is a sum of four squares: the set of sums of four squares of elements of
$\Z_p$ is compact, hence closed, and contains $\N$ (Lagrange), which is dense. Choose
$a_1,\dots,a_4\in\Z_p$ with $\sum a_i^2=-1$ and put $(\delta_c)=(1,a_1,a_2,a_3,a_4,0)$, so that
$\sum_c\delta_c^2=0$. Then
\[
u(z):=\prod_c\theta(\delta_cz)=\exp\Bigl(-\sum_cg(\delta_cz)\Bigr)=\prod_c\exp\bigl(\delta_c^2g(z)
-g(\delta_cz)\bigr)\in1+y\psmR
\]
by Lemma~\ref{psm:units}(b), and $u^{-1}\in1+y\psmR$. Since $U=\prod_iu(\mu_i(z))^{m_i}$, the first
assertion follows. For the second, $y(\varpi w)$ is an integral series in $w$ when
$v(\varpi)\ge\rho_p$ (Lemma~\ref{psm:param}(c)), and the $\mu_i$ have integral coefficients.
\end{proof}

For the next lemma let $\pi$, $q$ be as follows: if $p$ does not split, $\pi=\pi_p$ and $q=q_p$ as
in \S\ref{pa:ss-conv}; if $p$ splits, $\pi=p$ and $q=p$. Put $\rho_p=v(\pi)/(q-1)$ and
$\gamma_p=\frac1{p-1}-\rho_p$; thus $\gamma_p=0$ if $p$ splits, $\gamma_p=p/(p^2-1)$ if $p$ is
inert and $\gamma_p=1/(2p-2)$ if $p=d$, in agreement with \S\ref{pa:ss-conv}.

\begin{lemma}[The formal parameter at every prime]\label{psm:param}
Let $T=\frac1{q-1}\sum_{\zeta\in\mu_{q-1}}\zeta^{-1}y(\zeta z)$, where $\mu_{q-1}\subset\C_p$
\textup{(}for split $p$, $\mu_{p-1}\subset\Z_p$\textup{)}.
\begin{enumerate}
\item[(a)] $T\in y+y^2\psmR$, $\psmR=\pao[[T]]$, and $T(\pi z)=\pi T+\sum_{k\ge1}b_kT^{1+k(q-1)}$
with $b_k\in\pao$, $b_1\in\pao^\times$.
\item[(b)] With $\varpi^{q-1}=-\pi/b_1$ \textup{(}so $v(\varpi)=\rho_p$\textup{)}, $Y(z)=T(\varpi z)/
\varpi$ is integral and $Y\equiv\sum_{j\ge0}z^{q^j}\pmod{\pi}$.
\item[(c)] For $G\in\psmR$ and $v(\varpi')\ge\rho_p$ the series $G(\varpi'z)$ is integral.
\item[(d)] For every $j\ge0$, $v\bigl((q^j)!\,[q^j]T\bigr)=\gamma_p\,(q^j-1)$.
\end{enumerate}
\end{lemma}
\begin{proof}
For non-split $p$, (a)--(c) are Lemma~\ref{pa:param}, whose proof uses \eqref{pa:eq-endo} for
$\alpha\in\auO_{K_p}$; this is now Lemma~\ref{psm:units}(b) (for $\zeta\in\mu_{q-1}$, the pair
$(\zeta,\zeta\cdot\zeta^{p})$ or $(\zeta,\zeta^2)$ is a limit of pairs $(\alpha,\auNm\alpha)$, as
in the proof of Lemma~\ref{pa:units}(b)). For split $p$ the same proof applies with $\pi=p$, $q=p$,
once we know that the reduction of $[p]$, written in the parameter $T$, has order exactly $p$: the
reduction of $\mathcal E$ is ordinary~\cite[Chapter~13, Theorem~12]{LangEF}, so its formal group has
height one and $[p]$ is congruent modulo $\pam$ to a unit times $T^p$ plus terms of higher degree;
$\mu_{p-1}\subset\Z_p$, and $G(\zeta z)\in\psmR$ by Lemma~\ref{psm:units}(b). (For $p=2$, $\mu_1=\{1\}$
and $T=y$.) The computation of Lemma~\ref{pa:param}(b) uses only $v(\pi^{n-1})>0$ for $n\ge2$ and
(E3), valid also for $p=2$. (d) By (b), $\varpi^{q^j}[q^j]T=\varpi\,[q^j]Y$ and $[q^j]Y\equiv1$, so
$v([q^j]T)=-(q^j-1)\rho_p$; and $v((q^j)!)=(q^j-1)/(p-1)$ by Legendre's formula, $q^j$ being a power
of $p$.
\end{proof}

\begin{lemma}[The rate $\gamma_p$ is attained]\label{psm:katz}
For every $n\ge0$ there is $G\in\psmR$ with $v\bigl(n!\,[n]G\bigr)\le\gamma_p\,n$.
\end{lemma}
\begin{proof}
For $G\in\psmR$, $n![n]G=(d/dz)^nG(0)$ is integral, because $\psmR$ is stable under $d/dz$ and the
constant term of an element of $\psmR$ is integral. Let $c(n)=\max_{G\in\psmR}|n![n]G|$; the
maximum exists and is attained at a power $y^i$, $i\le n$, by the argument of
Lemma~\ref{psm:rules}(d). Then $c(0)=1$ and $0\le c(n)\le1$.

\emph{Submultiplicativity.} Let $c(m+n)=|(m+n)![m+n]y^i|$. The formal group law of the smooth model
with origin $P_0$ is integral (for the origin $O$ and the parameter $-x/y$ of a Weierstrass equation
with integral coefficients this is \cite[Chapter~IV, \S1]{SilvermanAEC}; for $P_0$ we argue directly): the morphism $(Q,Q')\mapsto Q+Q'-P_0$
extends to the smooth model and fixes $(P_0,P_0)$, hence induces a map of the completions along
$(P_0,P_0)$ and $P_0$; since the completion of the smooth model along $P_0$ is $\pao[[y]]$
(\S\ref{pa:ss-curve}), that of the product along $(P_0,P_0)$ is $\pao[[y(z_0),y(z_1)]]$, and so $y(z_0+z_1)^i=\sum_{a,b}c_{ab}y(z_0)^ay(z_1)^b$
with $c_{ab}\in\pao$. Since $(m+n)![m+n]y^i=m!\,n!\,[z_0^mz_1^n]\,y(z_0+z_1)^i$, the ultrametric
inequality gives $c(m+n)\le\max_{a,b}|m![m]y^a|\cdot|n![n]y^b|\le c(m)c(n)$.

\emph{Powers of $q$.} By Lemma~\ref{psm:param}(d), $c(q^j)\ge p^{-\gamma_p(q^j-1)}\ge
\beta^{q^j}$ with $\beta=p^{-\gamma_p}\in(0,1]$.

\emph{Fekete.} Suppose $c(n)<\beta^n$ for some $n\ge1$, and let $\varepsilon=c(n)/\beta^n<1$. For
$Q=q^j\ge n$ write $Q=nk+r$, $0\le r<n$. Then $\beta^Q\le c(Q)\le c(n)^kc(r)\le c(n)^k=
\varepsilon^k\beta^{nk}$, so $\beta^{r}\le\varepsilon^k$, that is $\beta^n\le\varepsilon^k$ for all
large $j$, which is absurd as $k\to\infty$. So $c(n)\ge\beta^n$, which is the claim.
\end{proof}

\begin{lemma}[Unit forms]\label{psm:unitform}
\begin{enumerate}
\item[(a)] If one form $\mu_{i_0}$ has all its coefficients units, then for every $\nu$,
$\psmval_6(\nu)\le\gamma_p|\nu|-v(\nu!)$.
\item[(b)] If there are $a_{x,i}\in\pao$ and a unit $s$ with $\sum_ia_{x,i}\mu_i(z)=s\,z_x$ for
all $x\in X$, then $\psmval_6(\nu)\le0$ for every $\nu$.
\end{enumerate}
\end{lemma}
\begin{proof}
(a) Let $n=|\nu|$, $G$ as in Lemma~\ref{psm:katz}, and $(\delta_c)$, $U$ as in
Lemma~\ref{psm:six}. Put $H_0=G(\mu_{i_0}(z))\,U^{-1}$ and $H_c=1$ for $c\ge1$; all lie in
$\psmR_\mu$. As $\delta_0=1$, $\prod_c(H_c\Theta)(\delta_cz)=G(\mu_{i_0}(z))U(z)^{-1}U(z)=
G(\mu_{i_0}(z))$, whose coefficient of $z^\nu$ is $[n]G\cdot\frac{n!}{\nu!}\prod_x\mu_{i_0,x}^{\nu(x)}$,
of valuation $v(n![n]G)-v(\nu!)\le\gamma_pn-v(\nu!)$. Apply Lemma~\ref{psm:rules}(c) with one new
copy of the variables and $A_c=\delta_c$.

(b) Put $G_x=\sum_ia_{x,i}\,y(\mu_i(z))\in\psmR_\mu$. Since $y=z+O(z^2)$, $G_x=s\,z_x+(\text{terms of
degree}\ge2)$, so $\bigl(\prod_xG_x^{\nu(x)}\bigr)\Theta=s^{|\nu|}z^\nu+(\text{terms of degree}>
|\nu|)$, and its coefficient of $z^\nu$ is the unit $s^{|\nu|}$. This gives $\psmval_1(\nu)\le0$;
add five copies in degree $0$ (Lemma~\ref{psm:rules}(a),(b)).
\end{proof}

\begin{lemma}[The leading coefficient is a unit]\label{psm:lambda}
Let $p$ be inert. The number $\lambda$ of Theorem~\ref{pa:onecurve} is a unit; hence in the normal
form of Proposition~\ref{pa:product}, if the scalar $c$ there is a unit, all $\lambda_{e,j}$ $(e$
odd$)$ are units; they are $c^{(p^e+1)p^j}\lambda$.
\end{lemma}
\begin{proof}
Let $h_n=[z^{p-1}]x^n$. From \eqref{pa:eq-curve}, $x'^2=(2y+a_1x)^2=4x^3+a_1^2x^2+4x$ and
$x''=6x^2+a_1^2x+2$, so
\[
\bigl(x^nx'\bigr)'=nx^{n-1}x'^2+x^nx''=(4n+6)x^{n+2}+(n+1)a_1^2x^{n+1}+(4n+2)x^n .
\]
The coefficient of $z^{p-1}$ of the left side is $p\,[z^p](x^nx')$, and $x^nx'\in\psmR$ has integral
coefficients in degrees $<q=p^2$ (Lemma~\ref{pa:param}(c)); so
$(4n+6)h_{n+2}+(n+1)a_1^2h_{n+1}+(4n+2)h_n\in\pam$ for all $n\ge0$, with $h_n$ integral and $a_1$
integral (Lemma~\ref{pa:model}). Write $p=2m'+1$. If $h_1\in\pam$, then, as $h_0=0$ and
$4n+6=2(2n+3)$ is a unit for $n+2\le m'$ ($2n+3\le p-2$), induction gives $h_n\in\pam$ for all
$n\le m'$. But $x=z^2+O(z^3)$, so $h_{m'}=[z^{2m'}]x^{m'}=1$. Hence $x_{p-1}=h_1$ is a unit.

By \eqref{pa:eq-g2}, $(p+1)p\,g_{p+1}=x_{p-1}$. In the proof of Theorem~\ref{pa:onecurve},
$\Xi_1(z)=-g^\flat(\varpi z)/c_\varpi$ with $c_\varpi=\varpi^{p+1}/p$, and $p+1$ is one of the degrees
kept by $\flat$; so $[p+1]\Xi_1=-p\,g_{p+1}=-x_{p-1}/(p+1)$, a unit. On the other hand
$\Xi_1\equiv\lambda Y^{p+1}\pmod p$ and $[p+1]Y^{p+1}\equiv1$ (Step~3 of that proof); so
$\lambda\equiv[p+1]\Xi_1$ is a unit. The statement on $\lambda_{e,j}$ is read off from
Theorem~\ref{pa:onecurve} and the proof of Proposition~\ref{pa:product}, where the substitution
$z\mapsto cz$ multiplies $\lambda_{e,j}$ by $c^{(p^e+1)p^j}$.
\end{proof}

\subsubsection{Couplings}\label{psm:ss-couplings}

In this subsection $p$ is inert in $K$ and $p>2$. Let $\auPl$ be a finite set, $X=\auPl$ (one
variable $z_\tau$ for each $\tau$), and let $\psmpr_e\colon\auPl\to\auPl$ ($e$ odd) be maps. Let
$(\mu,m)$ be a family as in \S\ref{psm:ss-least} such that $\Theta=\Theta_{\mu,m}$ has the normal
form \eqref{pa:eq-nf} for the partner maps $\psmpr_e$:
\[
\Theta(\varpi z)=W(z)\exp\bigl(\kappa_p\,\Xi(z)\bigr),\qquad\kappa_p=\frac{\varpi^{p+1}}p,\qquad
\Xi=\sum_{e\text{ odd}}\sum_{j\ge0}\lambda_{e,j}\Gamma_e^{\,p^j},\qquad\Gamma_e=\sum_\tau
z_\tau^{\,p^e}z_{\psmpr_e\tau},
\]
with $v(\varpi)=\rho_p=1/(p^2-1)$, $W$ integral, $\lambda_{e,j}\in\pao$, and assume that
$\lambda_{e,0}$ is a unit for every odd $e$. Then $v(\kappa_p)=\frac1{p-1}-1<0$ and $W(0)=1$, so
$W$ is a unit of $\pao[[z]]$.

\begin{definition}\label{psm:def-family}
A \emph{coupling family} is a subset $\auPl_0\subseteq\auPl$ with odd exponents $e(\tau)$,
$\tau\in\auPl_0$, such that: for $\tau,\tau_1\in\auPl_0$ and odd $e'\le e(\tau)$,
$\psmpr_{e'}(\tau_1)=\psmpr_{e(\tau)}(\tau)$ implies $\tau_1=\tau$ and $e'=e(\tau)$. Write
$\tau^+=\psmpr_{e(\tau)}(\tau)$. For $j\colon\auPl_0\to\N$ put
\[
\nu'=\sum_{\tau\in\auPl_0}p^{e(\tau)}j(\tau)\,\mathbf 1_\tau,\qquad\nu''=\sum_{\tau\in\auPl_0}j(\tau)
\,\mathbf 1_{\tau^+},\qquad|j|=\sum_\tau j(\tau),
\]
the degrees of the two factors of $\prod_\tau(z_\tau^{p^{e(\tau)}}z_{\tau^+})^{j(\tau)}$.
\end{definition}

\begin{proposition}[Couplings]\label{psm:coupling}
For a coupling family and every $j$,
\[
\psmval_6(\nu'+\nu'')\ \le\ \frac{|j|}{p-1}-|j|-\rho_p|\nu'+\nu''|-v(\nu''!) .
\]
Consequently $\psmval_6(\nu'+\nu'')\le-|j|$ and $v(\nu'!)+v(\nu''!)+\psmval_6(\nu'+\nu'')\le
\gamma_p|\nu'+\nu''|-|j|$.
\end{proposition}
\begin{proof}
Let $(\delta_c)$ and $U$ be as in Lemma~\ref{psm:six}, $F_\varpi(z)=\Theta(\varpi z)$ and
$U'(z)=U(\varpi z)=\prod_cF_\varpi(\delta_cz)$; $U'$ is integral with constant term $1$
(Lemma~\ref{psm:six}). In two copies $z',z''$ of the variables put
\[
G(z',z'')=F_\varpi(z'+z'')\prod_{c=1}^5F_\varpi(\delta_cz'),\qquad G(z',0)=U'(z') .
\]
\emph{Step 1: derivatives in $z''$.} Only the first factor involves $z''$, and
$\partial_{z_x}F_\varpi=(\kappa_p\,\partial_x\Xi+\partial_xW/W)F_\varpi$; so $\partial_{z''_x}G=(\kappa_pA_x+B_x)G$
with $A_x=(\partial_x\Xi)(z'+z'')$, $B_x=(\partial_xW/W)(z'+z'')$ integral. We claim: for variables
$t_1,\dots,t_N$ among the $z''_x$ there is an integral $g$ with $\partial_{t_1}\cdots\partial_{t_N}G=
\kappa_p^Ng\,G$ and $g\equiv\prod_iA_{t_i}\pmod\pam$. For $N=0$ take $g=1$. If $\partial(g G)$ is
to be computed for one more derivative $\partial=\partial_t$, then
\[
\partial(gG)=\bigl(\partial g+g(\kappa_pA_t+B_t)\bigr)G=\kappa_p\bigl(gA_t+\kappa_p^{-1}(\partial g+
gB_t)\bigr)G,
\]
and $\kappa_p^{-1}\in\pam$. Apply this to the multiset in which each $z''_{\tau^+}$ occurs $j(\tau)$
times ($N=|j|$): $\partial_{z''}^{\,\nu''}G=\kappa_p^{|j|}g\,G$ with
\[
g(z',0)\equiv\prod_{\tau\in\auPl_0}\bigl(\partial_{\tau^+}\Xi\bigr)(z')^{j(\tau)}\pmod\pam .
\]
\emph{Step 2: a unit coefficient.} We show that $[z'^{\nu'}]\bigl(g(z',0)U'(z')\bigr)$ is a unit.
Since $\nu'$ is supported on $\auPl_0$ we may put $z'_\sigma=0$ for $\sigma\notin\auPl_0$ and reduce
modulo $\pam$. Modulo $p$, $\partial(\Gamma_e^{p^j})=p^j\Gamma_e^{p^j-1}\partial\Gamma_e\equiv0$ for
$j\ge1$, and $\partial_{\tau^+}\Gamma_e\equiv\sum_{\sigma:\,\psmpr_e\sigma=\tau^+}z_\sigma^{p^e}$
(the other term carries the factor $p^e$). So, after putting $z_\sigma=0$ off $\auPl_0$,
\[
\partial_{\tau^+}\Xi\equiv\sum_{e\text{ odd}}\lambda_{e,0}\sum_{\sigma\in\auPl_0,\ \psmpr_e\sigma=
\tau^+}z_\sigma^{\,p^e}\pmod\pam ,
\]
a sum of monomials of degree $p^e$; by the definition of a coupling family the only one of degree
$\le p^{e(\tau)}$ is $\lambda_{e(\tau),0}z_\tau^{p^{e(\tau)}}$. Hence the product over $\tau$, times
$U'\equiv1+(\text{higher degree})$, is $\prod_\tau\lambda_{e(\tau),0}^{j(\tau)}\,z^{\nu'}$ plus terms of
higher total degree, and $\prod_\tau\lambda_{e(\tau),0}^{j(\tau)}$ is a unit.

\emph{Step 3: the coefficient.} The jet $\partial_{z'}^{\nu'}\partial_{z''}^{\nu''}G(0)$ equals
$\nu'!\,\nu''!\,[z'^{\nu'}z''^{\nu''}]G$ and also, by Step~1 at $z''=0$, $\kappa_p^{|j|}\nu'!\,
[z'^{\nu'}](g(z',0)U'(z'))$. By Step~2,
\[
v\bigl([z'^{\nu'}z''^{\nu''}]G\bigr)=|j|\,v(\kappa_p)-v(\nu''!)=|j|\bigl(\tfrac1{p-1}-1\bigr)-
v(\nu''!) .
\]
\emph{Step 4: unscaling.} $G(z',z'')=G_0(\varpi z',\varpi z'')$ with $G_0=\Theta(z'+z'')\prod_{c\ge1}
\Theta(\delta_cz')$, so $[z'^{\nu'}z''^{\nu''}]G_0$ has valuation $|j|(\frac1{p-1}-1)-v(\nu''!)-
\rho_p|\nu'+\nu''|$. Now $G_0=\prod_{c=0}^5\Theta(A_cz'+B_cz'')$ with $(A_0,B_0)=(1,1)$ and
$(A_c,B_c)=(\delta_c,0)$, integral; Lemma~\ref{psm:rules}(c) with $H_c=1$ gives the first claim.

For the consequences: $|\nu'+\nu''|=|\nu'|+|j|\ge(p+1)|j|$ (as $e(\tau)\ge1$), and
$\rho_p(p+1)=\frac1{p-1}$; so the bound is $\le-|j|-v(\nu''!)\le-|j|$. And by Legendre's formula
$v(\nu'!)\le|\nu'|/(p-1)$, so $v(\nu'!)+v(\nu''!)+\psmval_6\le\frac{|\nu'|+|j|}{p-1}-\rho_p|\nu'+
\nu''|-|j|=\gamma_p|\nu'+\nu''|-|j|$.
\end{proof}

\subsubsection{The bound at every prime}\label{psm:ss-twelve}

Now let $D$ be admissible data (Definition~\ref{an:def-data}), $\auPl$ its set of infinite places,
$\kappa$ its weight, $\auPb$ its $b$-places, $l_0=l_0(s)$, $k=r+1-m$, $a=xb$, and for a prime $p$
let $u(-p^e)\tau$ be the partner of $\tau$ (\S\ref{ss:au-primes}). Put
\[
\psmns=\lceil6b/k\rceil,\qquad l^-=(l-\psmns)^+,\qquad B(l)=(l^-+1)(l^-+2)\cdots(l^-+\psmns)
\qquad(l\in\N).
\]
Let $\mathsf w=(\mathsf w_i)_{i\in J}$ be a family in $\mathfrak a_0$ with multiplicities $c_i$
forming a \emph{tight frame of length $M$}: $\sum_ic_i\,\tau(\mathsf w_i)\aubar{\tau'(\mathsf w_i)}=
M\delta_{\tau\tau'}$ (cf. \S\ref{pf:ss-frames}). Let $\mathsf s\in\C$ be the number
$\varpi_{\mathrm{vac}}/\mathsf u$ of \S\ref{pf:ss-frames} and, at the prime $p$, consider the linear
forms in $z=(z_\tau)_{\tau\in\auPl}$
\[
\mu_i(z)=\mathsf s\sum_\tau\aubar{\tau(\mathsf w_i)}\,z_\tau=\mathfrak s\,u_i^\dagger(z),\qquad
u_i=\mathsf w_i/\aubar{\pi_M},\quad\mathfrak s=\mathsf s\,\pi_M .
\]
Their coefficients are integral at every prime (Lemma~\ref{pf:scale} and Corollary~\ref{pf:CS}).

\begin{lemma}[Facts at a non-split prime]\label{psm:nonsplit}
Let $p$ be a prime that does not split in $K$. Then $p\ge13$, $p\nmid q_1M$, the vectors $u_i$
are integral above $p$, they form with the multiplicities $c_i$ a resolution of the identity
\eqref{pa:eq-res}, and $\mathfrak s$ is a $p$-adic unit. If $p$ is inert, $\Theta_{\mu,c}$ has the
normal form of \S\ref{psm:ss-couplings} with all $\lambda_{e,0}$ units, and there are
$a_{\tau,i}\in\pao$ with $\sum_ia_{\tau,i}\mu_i(z)=\mathfrak s\,z_\tau$.
\end{lemma}
\begin{proof}
$p\ge13$ by \eqref{an:eq-split}; the primes of $q_1$ split, and $p_0,p_1\equiv1\pmod{q_2^2}$ split
(their residue classes are $1$ modulo $d\mid q_2$, hence squares modulo $d$). The resolution property
is $\sum_ic_i\tau(u_i)\aubar{\tau'(u_i)}=\delta_{\tau\tau'}$, which is the frame identity divided by
$\pi_M\aubar{\pi_M}=M$. Since $q_1\mathfrak a_0\subseteq\auO_E$, $q_1M\,u_i=q_1\mathsf w_i\cdot
\pi_M$ is integral, and $q_1M$ is prime to $p$. $\pi_M$ is a unit at $p$ (its norm $M$ is), and
$\mathsf s$ is a unit at $p\nmid q_1$ (Lemma~\ref{pf:scale}: $|\varpi_{\mathrm{vac}}|=
|\alpha|_w^{-1}=1$ as $\alpha$ is a unit off $q_1$; and $\mathsf u$ is a unit by
Corollary~\ref{pf:CS}). For inert $p$, Proposition~\ref{pa:product} (with $G=1$ and the unit scalar
$\mathfrak s$) gives the normal form and Lemma~\ref{psm:lambda} the units. Finally $a_{\tau,i}=c_i
\tau(u_i)$ works by \eqref{pa:eq-res}.
\end{proof}

Write $I_p$, $v(\pi_p)$ as in \S\ref{ss:au-primes} and recall $s_p$ and $F(p)$
\eqref{psm:eq-F}; for $p=d$, $v(\pi_d)=\frac12$ and $I_d=\{1,2,\dots\}$. We use the following size
inequalities.

\begin{lemma}[Sizes]\label{psm:sizes}
\begin{enumerate}
\item[(N1)] $x^{42}\le6b/k\le\psmns$, and $\kappa(\tau)+x^{42}\le2b$ for all $\tau$.
\item[(N2)] $3(\log b+2)+x^{42}+1\le a/x_1$, and $a/x_1\ge1$.
\item[(N3)] $(81s+1)\,a/x_1\le b/k^2$.
\item[(N4)] $d\le k\le2^s$, $27s<l_0\le81s$, and $k\le5r(r+1-m)$.
\end{enumerate}
\end{lemma}
\begin{proof}
By Lemma~\ref{an:lem-sizes}: $k\le r\le3x^5$, $x^{10^6}\le x_1$, $x_1^{10^6}\le b$, $s\le32\log x$,
$x^{11}\le2^s$, $\log x\ge10^{14}d$, $\kappa(\tau)\le b$. (N1): $x^{42}k\le3x^{47}\le6b$ and
$x^{42}\le b$. (N2): $a/x_1=xb/x_1\ge b^{1-10^{-6}}\ge x^{43}$, which exceeds
$3(\log b+2)+x^{42}+1$ as $\log b\le x^{1/100}$. (N3): equivalent to $(81s+1)k^2x\le x_1$, and
$(81s+1)k^2x\le(2592\log x+1)\,9x^{11}\le x^{12}\le x_1$. (N4): $d\le\log x\le x^5/2\le k$ (S2), and
$k\le3x^5\le x^{11}\le2^s$; the bounds for $l_0$ are in \S\ref{an:ssec-param}; $5r(r+1-m)=5rk\ge k$.
\end{proof}

\emph{Near pairs.} For a prime $p\le k$ that does not split, a \emph{near pair} is a pair $(\tau,e)$
with $e\in I_p$, $2e<l_0$, $\tau\in\auPb$ and $u(-p^e)\tau\in\auPb$.

\begin{lemma}[Near pairs]\label{psm:near}
Let $p\le k$ be a prime that does not split in $K$.
\begin{enumerate}
\item[(a)] If $p=d$ there are no near pairs. If $(\tau,e)$ is a near pair, then $p$ is inert, $e$ is
odd and $u(-p^e)\tau=u(-p)^e\tau$.
\item[(b)] Two near pairs with the same start have the same exponent; two near pairs with the same
end are equal; the end of a near pair is not the start of a near pair.
\end{enumerate}
\end{lemma}
\begin{proof}
(a) All $b$-places have class $-1$ modulo $q_0$, so a near pair needs $\ancl_{q_0}(u(-p^e))=1$,
i.e.\ $p^e\equiv-1\pmod{q_0}$. For $p=d$: with $h=q_0/2$ we have $d\equiv h-1\pmod{q_0}$ and
$h^2\equiv0\pmod{q_0}$ ($4\mid q_0$), so $d^e\equiv(-1)^e+e(-1)^{e-1}h$. If $e$ is odd this is
$-1+eh\equiv-1+h\not\equiv-1$; if $e$ is even it is $1-eh\equiv1\not\equiv-1$. So $p\ne d$, $p$ is
inert and $I_p$ consists of odd numbers; then $-p^e=(-p)^e$.

(b) By (D5) the set $\auPb$ is sparse for the primes up to $5r(r+1-m)\ge k\ge p$
(Definition~\ref{an:def-sparse}); in the classes modulo the largest divisor $q'$ of $q_1$ prime to
$p$, $u(-p)^i\sigma/\sigma$ has class $(-1)^ip^i$. Hence for $0<i<i'<l_0$ the places
$\sigma,u(-p)^i\sigma,u(-p)^{i'}\sigma$ are not all in $\auPb$. Same start $\tau$, exponents
$e<e'$: take $\sigma=\tau$, $(i,i')=(e,e')$, $e'<l_0/2$. Same end, $e<e'$: then $\tau=u(-p)^{e'-e}
\tau'$, and take $\sigma=\tau'$, $(i,i')=(e'-e,e')$; so $e=e'$ and then $\tau=\tau'$. End of
$(\tau,e)$ equal to the start of $(\tau',e')$: take $\sigma=\tau$, $(i,i')=(e,e+e')$, with
$e+e'<l_0$.
\end{proof}

\begin{theorem}[The bound on twelve copies]\label{psm:twelve}
Let $(\mathsf w,c)$ be a tight frame of length $M$ such that, at the prime $p$, some form
$\mu_{i_0}$ has all coefficients units. Let $l\in\N^{\auPl}$ with $\kappa(\tau)\le l_\tau\le
\kappa(\tau)+x^{42}$ for all $\tau$. Then, for $\Theta=\Theta_{\mu,c}$,
\[
\psmval_{12}(l)\ \le\ F(p)+2\sum_\tau v\bigl(B(l_\tau)\bigr)-\sum_\tau v(l_\tau!) .
\]
\end{theorem}
\begin{proof}
\emph{Case 1: $p$ splits.} Then $F(p)=0$ and $\gamma_p=0$; Lemma~\ref{psm:unitform}(a) gives
$\psmval_6(l)\le-v(l!)$, and six copies in degree $0$ give the claim.

\emph{Case 2: $p$ does not split and $p>k$.} Then $p>d$ (N4), so $p$ is inert, $v(\pi_p)=1$, and
$\psmpr_1\colon\tau\mapsto\tau^+=u(-p)\tau$ is a bijection; so $\auPl_0=\auPl$ with $e\equiv1$ is a
coupling family. Put $i(\tau)=\lfloor l_\tau^-/p\rfloor$ and $j(\tau)=\min(i(\tau),l_{\tau^+})$.
At a place $t=\tau^+$ we have $\nu'(t)=p\,j(t)\le l_t^-$ and $\nu''(t)=j(\tau)\le\min(\psmns,l_t)$,
because $i(\tau)\le l_\tau/p\le2b/k\le\psmns$ by (N1); hence $\nu'+\nu''\le l$, as
$l^-_t+\min(\psmns,l_t)=l_t$. Let $\eta=l-\nu'-\nu''$. By Proposition~\ref{psm:coupling},
Lemma~\ref{psm:nonsplit} and Lemma~\ref{psm:unitform}(b), $\psmval_{12}(l)\le-|j|+0$. It remains
to show, for each $\tau$,
\[
v(l_\tau!)-j(\tau)\le s_p(\kappa(\tau),\tau)+2v(B(l_\tau)).
\]
Since $(l^-+\psmns)!=l^-!\,B(l)$ and $l\le l^-+\psmns$, $v(l!)\le v(l^-!)+v(B(l))$; by Legendre,
$v(l^-!)=i+v(i!)$ with $i=\lfloor l^-/p\rfloor$; and $v(i!)\le v(\psmns!)\le v(B(l))$, as $i\le
\psmns$ and a product of $\psmns$ consecutive integers is divisible by $\psmns!$. So
$v(l_\tau!)\le i(\tau)+2v(B(l_\tau))$. Further $l_\tau^-\le\kappa(\tau)$ by (N1), so
$i(\tau)\le\kappa(\tau)/p$, and $l_{\tau^+}\ge\kappa(\tau^+)$; hence $i(\tau)-j(\tau)=(i(\tau)-
l_{\tau^+})^+\le(\kappa(\tau)/p-\kappa(u(-p)\tau))^+$, the term $e=1$ of $s_p(\kappa(\tau),\tau)$.

\emph{Case 3: $p$ does not split and $p\le k$.} Let $\auPl_0$ be the set of starts of near pairs and,
for $\tau\in\auPl_0$, $e(\tau)$ the exponent (Lemma~\ref{psm:near}(b)) and $j(\tau)=\lfloor l_\tau/
p^{e(\tau)}\rfloor$. If $\auPl_0=\emptyset$ (in particular if $p=d$) put $\nu'=\nu''=0$, $j=0$.
Otherwise $p$ is inert and $\auPl_0$ is a coupling family: if $\tau,\tau_1\in\auPl_0$, $e'\le
e(\tau)$ is odd and $u(-p^{e'})\tau_1=\tau^+$, then $(\tau_1,e')$ is a near pair ($e'\in I_p$,
$2e'<l_0$, $\tau_1\in\auPb$, and its end $\tau^+\in\auPb$) with the same end as $(\tau,e(\tau))$,
so $\tau_1=\tau$, $e'=e(\tau)$. We have $\nu'+\nu''\le l$: at $t\in\auPl_0$, $\nu'(t)=p^{e(t)}j(t)\le
l_t$ and $\nu''(t)=0$ (an end is not a start); at $t\notin\auPl_0$, $\nu'(t)=0$ and $\nu''(t)$ is
$j(\tau)$ for the unique near pair $(\tau,e(\tau))$ ending at $t$ (or $0$), and $j(\tau)\le l_\tau/2
\le b\le\kappa(t)\le l_t$ by (N1), as $t\in\auPb$. Let $\eta=l-\nu'-\nu''$. By
Proposition~\ref{psm:coupling} (Lemma~\ref{psm:nonsplit}) and Lemma~\ref{psm:unitform}(a),
\[
\psmval_{12}(l)\le\gamma_p|l|-|j|-v(\nu'!)-v(\nu''!)-v(\eta!) .
\]
For $n_1+n_2+n_3=n\le2b$ and $p\ge3$, Legendre's formula $\frac n{p-1}-\log_p(n+1)-1\le v(n!)\le
\frac n{p-1}$ gives $v(n!)\le v(n_1!)+v(n_2!)+v(n_3!)+3(\log b+2)$. Since $l_\tau\le2b$ (N1),
\[
\psmval_{12}(l)+\sum_\tau v(l_\tau!)\le\gamma_p|l|-|j|+3f(\log b+2).
\]
Now $\gamma_p\le1$ and $l_\tau-\kappa(\tau)\le x^{42}$, so $\gamma_p|l|\le\gamma_p\sum_\tau
\kappa(\tau)+fx^{42}$; and by $\gamma_p=v(\pi_p)\sum_{e\in I_p}p^{-e}$ \eqref{pa:eq-const} and
$y-(y-y')^+=\min(y,y')$,
\[
\gamma_p\sum_\tau\kappa(\tau)-F(p)=v(\pi_p)\sum_\tau R_\tau\le\sum_\tau R_\tau,\qquad
R_\tau=\sum_{e\in I_p}\min\bigl(\kappa(\tau)/p^e,\kappa(u(-p^e)\tau)\bigr).
\]
So it suffices to show, for each $\tau$,
\begin{equation}\label{psm:eq-place}
3(\log b+2)+x^{42}+R_\tau-j(\tau)\le v\bigl(B(l_\tau)\bigr)
\end{equation}
(with $j(\tau)=0$ for $\tau\notin\auPl_0$). Let $n=\lceil l_0/2\rceil$. For $e<n$ (so $2e<l_0$) the
term of $R_\tau$ is: at most $\kappa(\tau)/p^{e}\le l_\tau/p^e<j(\tau)+1$ if $(\tau,e)$ is a near pair
(then $e=e(\tau)$); at most $a/x_1$ otherwise (if $\tau\in\auPb$ and $u(-p^e)\tau\notin\auPb$ the
minimum is at most $\kappa(u(-p^e)\tau)=a/x_1$; if $\tau\notin\auPb$, it is at most
$\kappa(\tau)=a/x_1$). For $e\ge n$ the terms are at most $b/p^e$, with sum at most $2b/p^n\le
2b/(kp)$, because $n>s$ and $p^{n-1}\ge2^s\ge k$ (N4). As $n\le l_0$,
\[
R_\tau\le j(\tau)+1+l_0\,a/x_1+2b/(kp).
\]
By (N2) and (N3) (with $p\le k$ and $l_0\le81s$), the left side of \eqref{psm:eq-place} is at most
$(l_0+1)a/x_1+2b/(kp)\le3b/(kp)$, and $b/(kp)\ge(l_0+1)a/x_1\ge1$. On the other hand
$v(B(l))\ge v(\psmns!)\ge\psmns/p-1\ge6b/(kp)-1\ge3b/(kp)$. This proves \eqref{psm:eq-place}.
\end{proof}

\begin{remark}
Theorem~\ref{psm:twelve} is the converse of Proposition~\ref{au:count} at the weights: the
couplings that produce the divisibility $s_p$ in the upper bound are exactly those whose leading
coefficients $\lambda_{e,0}$ are units, and choosing the exponent $j(\tau)$ of each coupling as
large as the degrees allow realises the prime sum up to the loss $2\sum_\tau v(B(l_\tau))$. For
primes $p>k$ all places are coupled at $e=1$; for $p\le k$ only near pairs of $b$-places can
contribute more than the loss, and sparseness makes them disjoint.
\end{remark}

\subsubsection{Transfer to the indicator data}\label{psm:ss-attained}

Let $W_{12}=\operatorname{Hom}(E,E^{12})$ with the standard hermitian structure, $L$ its standard
self-dual lattice, $\psvac$ the vacuum and $\psvac^\circ=c_{\mathrm{vac}}^{12}\psvac$ the integral
vacuum \eqref{pk:eq-vac}. For $\mathfrak m=(\pi_M)$ and a class $y_0\in(\auO_E/\mathfrak m)^{12}$
let $\phi_{y_0}$ be the \emph{indicator datum}: the vector of the Heisenberg model whose component
at the places dividing $\mathfrak m$ is the indicator of $y_0$ and which agrees with $\psvac$ at
all other places (including the Gaussian at infinity). For a vector $\Phi$ let
$\Theta_\Phi(\pi_Ms_0\,z)$, $z=(z_{\tau,b})_{\tau\in\auPl,\,1\le b\le12}$, be its theta function
at the scale of \eqref{pk:eq-s0}, and $[\gamma]$ its Taylor coefficient at $\gamma\in
\N^{\auPl\times12}$.

The proof uses the following four statements on the Heisenberg model of $W_{12}$; they are
proved in \S\ref{pst:sec} (Propositions~\ref{pst:I1}--\ref{pst:I4}). Here, as in \S\ref{pst:sec},
the value of a vector $\Phi$ at the finite adelic point $(0,x)$, $x\in W_{12}(\auA_f)$, is
$\Phi(0,x):=\theta(1\otimes\psrho_f(x)\Phi)$; it is linear in $\Phi$.
\begin{enumerate}
\item[(I1)] \emph{A frame.} There is a tight frame $(\mathsf w,c)$ of length $M$ with $c_i\ge6fb$
for all $i$ such that at every prime $p$ (and every $\C\to\C_p$) some form $\mu_{i_0}$ has unit
coefficients.
\item[(I2)] \emph{Products of monomial sections.} For exponents $e_b(i)\in\N$ $(1\le b\le12$,
$i\in J)$ with $3e_b(i)\le c_i$ there is a vector $\Phi$ with $\Theta_\Phi(\pi_Ms_0z)=\prod_{b}
\prod_iy(\mu_i(z_b))^{e_b(i)}\theta(\mu_i(z_b))^{c_i}$, whose values $\Phi(0,x)$,
$x\in W_{12}(\auA_f)$, are integral at every prime, and with $\psrho^\circ(l)\Phi=\Phi$ for $l\in\pi_ML$.
\item[(I3)] \emph{Fourier expansion.} For $\Phi$ as in (I2), $M^{12f}\Phi=\sum_ia_i\,\psrho(u_i)
\psvac$, a finite sum with $\aubar{\pi_M}u_i\in\mathfrak a_0^{12}$, such that $a_i\psvac(0,x)$
is integral for every $x\in W_{12}(\auA_f)$; and each $\psrho(u)\psvac$ with $\aubar{\pi_M}u$ integral
is a sum of the $\phi_{y_0}$ with roots of unity as coefficients.
\item[(I4)] \emph{The vacuum constant.} There is a finite adelic $y$ with $c_{\mathrm{vac},y}:=\psvac^\circ(0,
y)=\theta(1\otimes\psrho_f(y)\psvac^\circ)\ne0$, and $|\sigma c_{\mathrm{vac},y}|\le
\exp(12C_0q_1^2)\le\exp(b)$ for all $\sigma\in\auAut(\C)$.
\end{enumerate}

\begin{theorem}[Sharpness of the count]\label{psm:attained}
Let $l\in\N^{\auPl}$ with $\kappa(\tau)\le l_\tau\le\kappa(\tau)+x^{42}$.
There is $t\in\C^\times$ with $|\sigma t|\le\exp(fmb/k)$ for all $\sigma\in\auAut(\C)$ such that
for every prime $p$ and every $\C\to\C_p$ there are $\gamma\in\N^{\auPl\times12}$ with
$\sum_b\gamma(\tau,b)=l_\tau$ for all $\tau$, and a class $y_0$ modulo $\pi_M$, with
\[
p^{-F(p)}\ \le\ \Bigl|\,\frac{c_{\mathrm{vac}}^{12}}t\,\prod_\tau l_\tau!\cdot[\gamma]\,
\Theta_{\phi_{y_0}}(\pi_Ms_0\,\cdot)\Bigr|_p .
\]
\end{theorem}
\begin{proof}
Let $(\mathsf w,c)$ be the frame of (I1), $y$ and $c_{\mathrm{vac},y}$ as in (I4), and put
\[
t=M^{12f}c_{\mathrm{vac},y}\prod_\tau B(l_\tau)^2 .
\]
\emph{Size of $t$.} $B(l)\le(l^-+\psmns)^{\psmns}\le(2b+\psmns)^{\psmns}$ by (N1), $\psmns\le
7b/k$, $\log(2b+\psmns)\le2\log b$ and $\log M\le22\log b$ (Lemma~\ref{an:lem-sizes}(S5)); so
$\log|\sigma t|\le264f\log b+28fb\log b/k+b$. Since $k\le3x^5\le f$, $b\le fb/k$; since
$b/k\ge b/(3x^5)\ge264\log b$, $264f\log b\le fb/k$; and $m\ge\sqrt x\ge(\log b)^{50}\ge
28\log b+2$. Hence $\log|\sigma t|\le fmb/k$.

\emph{The coefficient.} Fix $p$ and $\C\to\C_p$, and let $r=F(p)+2\sum_\tau v(B(l_\tau))-
\sum_\tau v(l_\tau!)$. By Theorem~\ref{psm:twelve} (with $i_0$ from (I1)) and
Lemma~\ref{psm:rules}(d) there are monomials $H_b=\prod_iy(\mu_i)^{e_b(i)}$ with $\sum_ie_b(i)\le
|l|\le2fb$, hence $3e_b(i)\le6fb\le c_i$, and $\alpha_1+\dots+\alpha_{12}=l$, with $\sum_bv([\alpha_b]
(H_b\Theta))\le r$. Let $\Phi$ be the vector of (I2) for these exponents and
$\gamma(\tau,b)=\alpha_b(\tau)$; then $\sum_b\gamma(\tau,b)=l_\tau$, and the coefficient
$A=[\gamma]\Theta_\Phi(\pi_Ms_0\,\cdot)=\prod_b[\alpha_b](H_b\Theta)$ has $v(A)\le r$. Let
$\Lambda(\Phi')=c_{\mathrm{vac}}^{12}[\gamma]\Theta_{\Phi'}(\pi_Ms_0\,\cdot)$, a linear form. By (I3),
\[
M^{12f}\Lambda(\Phi)\,\psvac(0,y)=\sum_i\sum_{y_0}\zeta_{i,y_0}\,\bigl(a_i\psvac(0,y)
\bigr)\,\Lambda(\phi_{y_0})
\]
with roots of unity $\zeta_{i,y_0}$ and integral $a_i\psvac(0,y)$. Since $\psvac^\circ=c_{\mathrm{vac}}^{12}\psvac$ and the value at $(0,y)$ is linear,
$c_{\mathrm{vac}}^{12}\psvac(0,y)=\psvac^\circ(0,y)=c_{\mathrm{vac},y}$. Choose $y_0$ with
$v(\Lambda(\phi_{y_0}))$ least. The ultrametric inequality gives
\[
v\bigl(\Lambda(\phi_{y_0})\bigr)\le v\bigl(M^{12f}c_{\mathrm{vac}}^{12}A\,\psvac(0,y)\bigr)=
v(A)+v(M^{12f}c_{\mathrm{vac},y})\le r+v(M^{12f}c_{\mathrm{vac},y}).
\]
Therefore
\[
v\Bigl(\frac{c_{\mathrm{vac}}^{12}}t\prod_\tau l_\tau!\,[\gamma]\Theta_{\phi_{y_0}}(\pi_Ms_0\,\cdot)
\Bigr)=v\bigl(\Lambda(\phi_{y_0})\bigr)-v(M^{12f}c_{\mathrm{vac},y})-2\sum_\tau v(B(l_\tau))+\sum_\tau v(l_\tau!)
\le F(p),
\]
which is the claim.
\end{proof}

\emph{Application in Proposition~\ref{psm:small}.} For a main pair $p=(\nu,z)$ the degrees
$l_\tau=c(\tau)-z$ satisfy $\kappa(\tau)\le l_\tau\le\kappa(\tau)+x^{42}$, so Theorem~\ref{psm:attained} applies; this is part (i) of
the note to Proposition~\ref{po:small}, with the exponent $F(\ell)$ of \eqref{psm:eq-F}. The
number $t$ is the first factor of $c_0=t\,I_{q_2}(1)\,L$, and $|\sigma t|\le\exp(fmb/k)$ is the
first term of the bound $\exp(fmb/k+3x^{35})$.



%% ===== proofs/C-small2.tex =====
% proofs/C-small2.tex
% Continuation of proofs/C-small.tex: proof of the four inputs (I1)--(I4) of the block
% Section C-small2.
% A-closure and C-small. Labels of this file start with pst:. Only amsmath/amssymb/amsthm are used.
\providecommand{\auO}{\mathcal O}
\providecommand{\auA}{\mathbb A}
\providecommand{\auPl}{\Phi}
\providecommand{\auNm}{\operatorname{Nm}}
\providecommand{\auAut}{\operatorname{Aut}}
\providecommand{\auHom}{\operatorname{Hom}}
\providecommand{\aubar}[1]{\overline{#1}}
\providecommand{\pao}{\mathfrak o}
\providecommand{\pam}{\mathfrak m}
\providecommand{\psvac}{\mathbf v}
\providecommand{\psrho}{\varrho}
\providecommand{\psHS}{\mathcal H}
\providecommand{\pfT}{\operatorname{Tr}}
\providecommand{\pfbr}[1]{\theta[#1]}
\providecommand{\pstS}{\mathcal S}

\subsection{The four inputs of the sharpness of the count}\label{pst:sec}

This subsection proves the statements (I1)--(I4) that were isolated before
Theorem~\ref{psm:attained}; with them, Theorem~\ref{psm:attained} is proved in full. The notation is
that of \S\ref{psm:ss-twelve} and \S\ref{psm:ss-attained}: $W_{12}=\auHom(E,E^{12})=E^{\mathsf C}$
with $\mathsf C=[12]\times[1]$, $L=\alpha\widehat\auO_E^{\,\mathsf C}$, $\Lambda_L=\mathfrak
a_0^{\mathsf C}$, the coordinates $\gamma^\natural=(c_\psi\tau(\gamma_b))_{\tau,b}$ of
\S\ref{pk:ss-conj}, $\Lambda=\Omega_0\auO_K$ the normalised lattice of \eqref{pk:eq-s0}, with theta
function $\theta$, hermitian form $H$ and multiplier $\chi$ \eqref{pa:eq-theta}, and $\mathsf s=
\varpi_{\mathrm{vac}}/\mathsf u$. Recall $\pfbr u=e^{-H(u,u)/2}\theta(u)$ and the spaces $\pstS_m$
of \eqref{pf:eq-sections}. For $G\in\pstS_m$ and $u\in\Q\Lambda$ we call $e^{-mH(u,u)/2}G(u)$ the
\emph{torsion value} of $G$ at $u$; for $G=\theta^m$ it is $\pfbr u^m$.

\subsubsection{(I1): a frame of length $M$ with a unit form at every prime}\label{pst:ss-frame}

\begin{lemma}[Size of the frame of length prime to $q_1$]\label{pst:hsize}
The length $\mathsf h$ of the tight frame of Lemma~\ref{pf:frame}, built with the least admissible
exponent $J_0$ (the exponent $J$ of that proof), satisfies $\mathsf h\le b$.
\end{lemma}
\begin{proof}
Let $\zeta=\zeta_{q_1}$, $T=q_1/q_0$. Every prime dividing $q_1$ divides $q_0$
(Lemma~\ref{an:lem-sizes}(S7)), so the minimal polynomial of $\zeta$ over $K$ is
$P(X)=\Phi_{q_1}(X)=\Phi_{q_0}(X^{T})$ ($\Phi_n$ the cyclotomic polynomials; $K\cap\Q(\zeta)=\Q$
since $d\nmid q_1$). The roots of $\Phi_{q_0}$ have absolute value $1$, so its coefficients are
bounded by binomial coefficients $\binom{\varphi(q_0)}{k}\le2^{q_0}$, and so are those of $P$.
Hence, for every $\tau\in\auPl$: (i) the coefficients $b_i=\sum_{j>i}P_j\zeta^{j-i-1}$ of
$P(X)/(X-\zeta)$ satisfy $|\tau(b_i)|\le q_1 2^{q_0}$; (ii) $\tau(P'(\zeta))=P'(\eta')$ with
$\eta'=\tau(\zeta)$, and $P'(\eta')=T\eta'^{\,T-1}\Phi_{q_0}'(\eta'^{\,T})$, where $\eta'^{\,T}$ is
a primitive $q_0$-th root of unity. The number $\Phi_{q_0}'(\eta)$, $\eta$ primitive of order $q_0$, is
a non-zero algebraic integer of $\Q(\zeta_{q_0})$ whose conjugates $\Phi_{q_0}'(\eta'')$ have absolute
value at most $q_0^22^{q_0}$; its norm is a non-zero integer, so $|\Phi'_{q_0}(\eta)|\ge(q_0^22^{q_0}
)^{-\varphi(q_0)}\ge e^{-3q_0^2}$. Hence $|\tau(P'(\zeta))|\ge e^{-3q_0^2}$.

Every prime $p\mid q_1$ divides $q_0$ (Lemma~\ref{an:lem-sizes}(S7)), hence $4p\mid d+1$ (Definition~\ref{an:def-param} and the remark after it), so $p$ divides $\omega\aubar\omega=(1+d)/4$; hence the least $J_0$ with $q_1\mid(\omega
\aubar\omega)^{J_0}$ is at most $\max_pv_p(q_1)\le\log_2q_1$; and $|\omega|^2=(1+d)/4\le d$. With
$e_i=\aubar\omega^{\,J_0}b_i/P'(\zeta)$ as in the proof of Lemma~\ref{pf:frame},
\[
|\tau(e_i)|^2\le d^{J_0}q_1^24^{q_0}e^{6q_0^2}=:B,\qquad\nu(e_i)=\sum_\tau|\tau(e_i)|^2\le fB,
\]
and $\mathsf h=q_1(1+\sum_{i<f}\nu(e_i))+\omega^{J_0}+\aubar\omega^{\,J_0}\le q_1(1+f^2B)+2d^{J_0/2}\le
4q_1^3B$, using $f\le q_1$. By Lemma~\ref{an:lem-sizes}(S1),(S3), $q_1\le x^{12}$ and $\log x\ge
10^{14}(q_0+2d)$; so $J_0\le18\log x$, $\log d\le d\le10^{-14}\log x$, $q_0\le10^{-14}\log x$, and
\[
\log\mathsf h\le\log4+36\log x+18\log d\log x+24\log x+2q_0+6q_0^2\le(\log x)^2 .
\]
By (S1), $\log b\ge(\log x_1)^6\ge(\log x)^6$. So $\mathsf h\le b$.
\end{proof}

\begin{proposition}[(I1)]\label{pst:I1}
There is a tight frame $(\mathsf w_i,c_i)_{i\in J}$ of length $M$ with $c_i\ge6fb$ for all $i$ such
that, for every prime $p$ and every $\C\to\C_p$, some $i_0$ has the property that all coefficients
$\mathsf s\,\aubar{\tau(\mathsf w_{i_0})}$ $(\tau\in\auPl)$ of $\mu_{i_0}$ are units. The members
$\mathsf w_i$ do not depend on $M$; in particular this holds at $p_0$ and $p_1$.
\end{proposition}
\begin{proof}
\emph{The frame.} Let $(v_j)_{j\in\mathsf I}$ be the frame of Lemma~\ref{pf:frame}, of length
$\mathsf h$ prime to $q_1$ and $\le b$ (Lemma~\ref{pst:hsize}), and $(\zeta^k)_{k<q_1}$ the frame of
length $q_1$. Let $R=\lceil6fb\rceil$. Since $\gcd(q_1,\mathsf h)=1$, every integer
$\ge(q_1-1)(\mathsf h-1)$ is of the form $\alpha q_1+\beta\mathsf h$ with $\alpha,\beta\in\N$.
By Lemma~\ref{an:lem-sizes}(S3),(S5), $q_1\le x^{12}\le b$, $R\le7b^2$ and $M\ge b^{20}$, so
$M-R(q_1+\mathsf h)\ge b^{20}-14b^3\ge q_1\mathsf h$; write $M-R(q_1+\mathsf h)=\alpha q_1+\beta
\mathsf h$. Take the members $\zeta^k$ $(k<q_1)$ with multiplicity $\alpha+R$ and the members $v_j$
with multiplicity $\beta+R$. A disjoint union of tight frames, with multiplicities, is a tight frame
whose length is the sum of the lengths; so this is a tight frame of length $(\alpha+R)q_1+(\beta+R)
\mathsf h=M$, all multiplicities are $\ge R\ge6fb$, and all members lie in $\mathfrak a_0$.

\emph{The unit form.} Fix $p$ and $\iota\colon\C\to\C_p$, and for $\tau\in\auPl$ let $w_\tau$ be
the place of $E$ induced by $\iota\circ\tau\circ c$, so that $|\iota(\aubar{\tau(\mathsf v)})|=|\mathsf
v|_{w_\tau}$. By Lemma~\ref{pf:scale} and Corollary~\ref{pf:CS} ($\mathsf u$ a unit),
\[
|\iota(\mathsf s\,\aubar{\tau(\mathsf v)})|=|\mathsf v|_{w_\tau}/|\alpha|_{w_\tau}\qquad(\mathsf v
\in E^\times).
\]
The places $w_\tau$ lie above one place $\mathfrak q$ of $K$. If $\omega\notin\mathfrak q$, then
$\alpha_{w_\tau}=1$ for all $\tau$ (\S\ref{pk:ss-conv}), and the member $\mathsf w_{i_0}=\zeta^0=1$
works. If $\omega\in\mathfrak q$, then $\alpha_{w_\tau}=1/P'(\zeta)$ for all $\tau$, and
$\aubar\omega\notin\mathfrak q$ ($\omega+\aubar\omega=1$). Take $\mathsf w_{i_0}=\zeta e_{f-1}$; it is
a member of the frame of Lemma~\ref{pf:frame} (the member $\zeta^k\rho(e_{f-1})$ with $k=1$, $\rho=1$,
which is not one of the replaced pairs, as $q_1>1$). Since $P$ is monic of degree $f$, $b_{f-1}=1$
and $\zeta e_{f-1}=\zeta\aubar\omega^{\,J_0}/P'(\zeta)$; as $\zeta$ and $\aubar\omega$ are units at
$w_\tau$, $|\zeta e_{f-1}|_{w_\tau}=|1/P'(\zeta)|_{w_\tau}=|\alpha|_{w_\tau}$. In both cases all
coefficients of $\mu_{i_0}$ are units. Nothing here depends on whether $p$ divides $M$.
\end{proof}

\subsubsection{(I2): products of monomial sections}\label{pst:ss-product}

\begin{lemma}[Torsion values of monomial sections]\label{pst:torsion}
Let $e,c\in\N$ with $3e\le c$ and $G=y^e\theta^c$. Then $G$ is entire, $G\in\pstS_c$, and for every
$u\in\Q\Lambda$, every prime $\ell$ and every $\iota$ the torsion value $e^{-cH(u,u)/2}G(u)$ is an
algebraic integer.
\end{lemma}
\begin{proof}
The function $y$ has poles of order $3$ exactly at $t+\Lambda$, where $\theta$ has simple zeros; so
$G$ is entire, and $G\in\pstS_c$ by \eqref{pf:eq-sections} (the polynomial $y^e$ has a pole of
order $3e\le c$ at $O$). Let $u\in\Q\Lambda$. The torsion value is algebraic: $Q(u)$ is a torsion point of the curve $y^2+a_1xy=x^3+x$, defined over $\Q(a_1)\subset\overline\Q$, so $x(u),y(u)\in\overline\Q$, and $\pfbr u$ is algebraic by Proposition~\ref{pf:values}; it remains to show it is integral at every $\ell$ and $\iota$.

If $u\notin t+\Lambda$, the torsion value is $y(u)^e\pfbr u^{\,c}$. If $u-t\notin E_1$, then
$\pfbr u$ is a unit (Proposition~\ref{pf:values}) and $Q(u)$ does not reduce to the origin of the
smooth model (Lemma~\ref{pf:two}), so $y(u)$ is integral. If $u-t\in E_1$, then $|x(u)|>1$ and
$|\pfbr u|=|x(u)|^{-1/2}$ (Proposition~\ref{pf:values}); from $y^2+a_1xy=x^3+x$ with $a_1$
integral, $|y(u)|>|x(u)|$ (otherwise the left side would have absolute value $\le|x(u)|^2<
|x(u)|^3$), hence $|y(u)|^2=|x(u)|^3$; so $|y(u)^e\pfbr u^{\,c}|=|x(u)|^{(3e-c)/2}\le1$.

If $u\in t+\Lambda$ and $c>3e$ the value is $0$. Let $c=3e$, so $G=(y\theta^3)^e$. For $u=t+\lambda$,
$\lambda\in\Lambda$, the functional equation of $\pstS_c$ gives $e^{-cH(u,u)/2}G(u)=\chi(\lambda)^c
e^{ic\,\Iml H(\lambda,t+\lambda/2)}\cdot e^{-cH(t,t)/2}G(t)$, and the middle factor is a root of
unity, since $\Iml H$ is $\boldsymbol\pi$ times an integral alternating form on $\Lambda$ and
$2t,\lambda\in\Lambda$. So it suffices to treat $u=t$. As $w\to0$: $x(t+w)+\wp(t)=\wp(w)$, so
$x(t+w)\sim w^{-2}$ and $y(t+w)=\tfrac12\wp'(w)-\tfrac12a_1x(t+w)\sim-w^{-3}$; and $e^{-H(t+w,t+w)/2}
\theta(t+w)=\pfbr{t+w}=e^{-\boldsymbol\pi iE(w,t)}e^{-H(w,w)/2}\theta_t(w)\sim\theta_t'(0)w$, because
$\theta_t(0)=\pfbr t=0$. Hence $e^{-3H(t,t)/2}(y\theta^3)(t)=-\theta_t'(0)^3$, and $\theta_t'(0)$ is a
unit at every prime by the proof of Corollary~\ref{pf:K1}(c).
\end{proof}

Let $(\mathsf w_i,c_i)_{i\in J}$ be a tight frame of length $M$, and put, for $z=(z_{\tau,b})\in
\C^{\auPl\times\mathsf C}$ and $b\in[12]$,
\[
a_{i}(z_b)=\frac{\Omega_0}{\pi_Mc_\psi}\sum_{\tau\in\auPl}\aubar{\tau(\mathsf w_i)}\,z_{\tau,b} ,
\qquad\text{so that}\qquad a_i\bigl((\pi_Ms_0z)_b\bigr)=\mathsf s\sum_\tau\aubar{\tau(\mathsf w_i)}
z_{\tau,b}=\mu_i(z_b)
\]
by \eqref{pk:eq-s0} ($s_0\Omega_0/c_\psi=\mathsf s$). For $\gamma\in W_{12}(E)$,
$a_i(\gamma^\natural_b)=\frac{\Omega_0}{\pi_M}\pfT_{E/K}(\aubar{\mathsf w_i}\gamma_b)$, since the
$\tau\in\auPl$ fix $K$ and commute with $c$.

\begin{lemma}[Isometry]\label{pst:iso}
$\sum_{b}\sum_ic_i\,H\bigl(a_i(w_b),a_i(z_b)\bigr)=\langle z,w\rangle=\sum_{\tau,b}\aubar{w_{\tau,b}}
z_{\tau,b}$ for all $z,w$.
\end{lemma}
\begin{proof}
$H(v,v')=\boldsymbol\pi\aubar vv'/A$ with $A=|\Omega_0|^2\sqrt d/2$, and $\boldsymbol\pi|\Omega_0|^2
/(Ac_\psi^2)=1$ (\S\ref{pk:ss-conj}). So the left side is $M^{-1}\sum_b\sum_{\tau,\tau'}\aubar{w_{
\tau,b}}z_{\tau',b}\sum_ic_i\tau(\mathsf w_i)\aubar{\tau'(\mathsf w_i)}$, and the frame identity gives
$\sum_ic_i\tau(\mathsf w_i)\aubar{\tau'(\mathsf w_i)}=M\delta_{\tau\tau'}$.
\end{proof}

\begin{proposition}[(I2)]\label{pst:I2}
Let $e_b(i)\in\N$ with $3e_b(i)\le c_i$, and $G_{i,b}=y^{e_b(i)}\theta^{c_i}\in\pstS_{c_i}$. There is a
unique $\Phi\in\psHS_f(W_{12})$ with
\[
\Theta_\Phi(z)=\prod_{b\in[12]}\prod_{i\in J}G_{i,b}\bigl(a_i(z_b)\bigr),\qquad\text{hence}\qquad
\Theta_\Phi(\pi_Ms_0z)=\prod_b\prod_iy(\mu_i(z_b))^{e_b(i)}\theta(\mu_i(z_b))^{c_i}.
\]
Moreover {\rm(a)} $\psrho_f(l)\Phi=\epsilon_L(l)\Phi$ for all $l\in\pi_ML$, that is,
$\psrho^\circ(l)\Phi=\Phi$ (notation of Lemma~\ref{pf:fourier}); {\rm(b)} for every $x\in
W_{12}(\auA_f)$ the value $\theta(1\otimes\psrho_f(x)\Phi)$ is an algebraic integer.
\end{proposition}
\begin{proof}
Write $S(z)$ for the product. Let $\lambda=\pi_M\lambda'$ with $\lambda'\in\Lambda_L$. Then
$\omega_{i,b}:=a_i(\lambda^\natural_b)=\Omega_0\nu_{i,b}$ with $\nu_{i,b}=\pfT_{E/K}(\aubar{\mathsf
w_i}\lambda'_b)\in\pfT_{E/K}(\aubar{\mathfrak a_0}\mathfrak a_0)=\pfT_{E/K}(\mathfrak d_{E/K}^{-1})
\subseteq\auO_K$, so $\omega_{i,b}\in\Lambda$. The functional equation of $\pstS_{c_i}$ and
Lemma~\ref{pst:iso} give
\[
S(z+\lambda^\natural)=\epsilon(\lambda)\,e^{\langle z,\lambda^\natural\rangle+|\lambda^\natural|^2/2}
S(z),\qquad\epsilon(\lambda)=\prod_{i,b}\chi(\Omega_0\nu_{i,b})^{c_i}=(-1)^{\sum_{i,b}c_i\auNm
\nu_{i,b}}
\]
by Lemma~\ref{pf:chi}(b). Now $\sum_ic_i\auNm\nu_{i,b}=\sum_ic_i|\sum_\tau\aubar{\tau(\mathsf w_i)}
\tau(\lambda'_b)|^2=M\sum_\tau|\tau(\lambda'_b)|^2$ by the frame identity, and summing over $b$,
$\sum_{i,b}c_i\auNm\nu_{i,b}=M\pfT_{E/K}\langle\lambda',\lambda'\rangle=\pfT_{E/K}\langle\lambda,
\lambda\rangle$, as $\pi_M\aubar{\pi_M}=M$. By Lemma~\ref{pf:eps}, $\epsilon(\lambda)=\epsilon_L(
\lambda)$ on $\pi_M\Lambda_L$. In particular $S$ satisfies the functional equation of
Lemma~\ref{pk:dict} with $\mathsf N=M$ ($M\Lambda_L\subseteq\pi_M\Lambda_L$), so $S=\Theta_\Phi$ for a
unique $\Phi$; the second formula follows from $a_i((\pi_Ms_0z)_b)=\mu_i(z_b)$.

(a) As at the end of the proof of Lemma~\ref{pf:semi}: the computation in the proof of
Lemma~\ref{pk:dict} shows that $\Theta_{\psrho_f(\lambda)\Phi}(z)=e^{\langle z,\lambda^\natural
\rangle-|\lambda^\natural|^2/2}\Theta_\Phi(z-\lambda^\natural)$ for $\lambda\in\Lambda_L$ (the root
of unity there is $1$ for lattice vectors, as in the proof of Lemma~\ref{pf:semi}), which
by the functional equation just proved (applied at $z-\lambda^\natural$) is $\epsilon_L(\lambda)^{-1}\Theta_\Phi(z)=\epsilon_L(\lambda)\Theta_\Phi(z)$ for $\lambda\in\pi_M
\Lambda_L$; so $\psrho_f(\lambda)\Phi=\epsilon_L(\lambda)\Phi$ by injectivity
($\epsilon_L=\epsilon_L^{-1}$); it extends to $\pi_ML$ by continuity, $\pi_M\Lambda_L$ being dense in
$\pi_ML$ and $\epsilon_L$ locally constant.

(b) Since $W_{12}(E)+\pi_ML=W_{12}(\auA_f)$, by (a) $\psrho_f(x)\Phi$ is a root of unity times
$\psrho_f(\gamma)\Phi$ for some $\gamma\in W_{12}(E)$. By Lemma~\ref{pk:dict} at $z=0$,
$\theta(1\otimes\psrho_f(\gamma)\Phi)=\Theta_{\psrho_f(\gamma)\Phi}(0)$ is a root of unity times
$e^{-|\gamma^\natural|^2/2}S(-\gamma^\natural)$. By Lemma~\ref{pst:iso}, $|\gamma^\natural|^2=
\sum_{i,b}c_iH(u_{i,b},u_{i,b})$ with $u_{i,b}=-a_i(\gamma^\natural_b)=-\frac{\Omega_0}{\pi_M}
\pfT_{E/K}(\aubar{\mathsf w_i}\gamma_b)\in K\Omega_0=\Q\Lambda$. Hence
\[
e^{-|\gamma^\natural|^2/2}S(-\gamma^\natural)=\prod_{i,b}e^{-c_iH(u_{i,b},u_{i,b})/2}G_{i,b}(u_{i,b}),
\]
a product of torsion values of monomial sections, which are algebraic integers by
Lemma~\ref{pst:torsion}.
\end{proof}

\emph{The scale.} The Taylor series are taken at the scale $\pi_Ms_0$, and the scale of the linear
forms is $\mathsf s=\varpi_{\mathrm{vac}}/\mathsf u$. This confirms the identification of
\S\ref{psm:ss-twelve}: the forms $\mu_i$ there are those of Proposition~\ref{pst:I2}.

\subsubsection{(I3): Fourier expansion in translates of the vacuum}\label{pst:ss-fourier}

\begin{proposition}[(I3)]\label{pst:I3}
Let $\Phi$ be as in Proposition~\ref{pst:I2}, and let $Y\subset\aubar{\pi_M}^{\,-1}\Lambda_L$ be a
set of representatives of $\aubar{\pi_M}^{\,-1}L/L$. Then $\#Y=M^{12f}$ and there are $a_y\in\C$
$(y\in Y)$ with
\begin{enumerate}
\item[(a)] $M^{12f}\Phi=\sum_{y\in Y}a_y\,\psrho_f(y)\psvac$;
\item[(b)] $a_y\,\theta(1\otimes\psrho_f(x)\psvac)$ is an algebraic integer for every $y\in Y$ and
every $x\in W_{12}(\auA_f)$.
\end{enumerate}
Further, for every $u\in W_{12}(\auA_f)$ with $\aubar{\pi_M}u\in L$, $\psrho_f(u)\psvac=\sum_{y_0}
\zeta_{u,y_0}\phi_{y_0}$ with roots of unity $\zeta_{u,y_0}$, the sum over the classes $y_0\in
(\auO_E/\pi_M\auO_E)^{12}$.
\end{proposition}
\begin{proof}
\emph{Commutators.} For $x,x'\in W_{12}(\auA_f)$ the commutator $[x,x']=\psi(\jmath\langle x',x
\rangle)$ is trivial when $\langle x',x\rangle\in\jmath^{-1}\mathfrak d_E^{-1}\widehat\auO_E$, in
particular for $x,x'\in L$ (self-duality). If $l=\pi_Ml'$ and $y=\aubar{\pi_M}^{\,-1}y'$ with
$l',y'\in L$, then $\langle l,y\rangle=\aubar{\pi_M}\,\aubar{\pi_M}^{\,-1}\langle l',y'\rangle=
\langle l',y'\rangle$, so $[l,y]=1$; and the pairing $(l,y)\mapsto[l,y]$ identifies $L/\pi_ML$ with
the dual of $\aubar{\pi_M}^{\,-1}L/L$, by self-duality of $L$ (the annihilator of
$\aubar{\pi_M}^{\,-1}L$ in $L$ is $\pi_ML$). The index is $\#Y=[\aubar{\pi_M}^{\,-1}L:L]=
|\auNm_{E/\Q}\aubar{\pi_M}|^{12}=M^{12f}$. Representatives can be taken in $\aubar{\pi_M}^{\,-1}
\Lambda_L$ by strong approximation ($W_{12}(E)$ is dense in $W_{12}(\auA_f)$).

\emph{The coefficients.} Let $\mathsf P=\sum_{\gamma\in L/\pi_ML}\psrho^\circ(\gamma)$; it is well
defined on the space $\mathsf V$ of vectors fixed by $\psrho^\circ(\pi_ML)$, because $\psrho^\circ$ is
a homomorphism on $L$ (\S\ref{pf:ss-frames}). For $y\in Y$, $\psrho_f(-y)\Phi\in\mathsf V$ by
Proposition~\ref{pst:I2}(a) and $[\pi_ML,y]=1$; and $\mathsf P\psrho_f(-y)\Phi$ is fixed by
$\psrho^\circ(L)$, hence, by (W4), equal to $a_y\psvac$ for some $a_y\in\C$. Now
$\psrho_f(y)\psrho^\circ(\gamma)\psrho_f(-y)=[y,\gamma]^{\pm1}\psrho^\circ(\gamma)$ (the sign of the
exponent depends only on the convention for $[\,,]$), so
\[
\sum_{y\in Y}a_y\psrho_f(y)\psvac=\sum_{\gamma\in L/\pi_ML}\Bigl(\sum_{y\in Y}[y,\gamma]^{\pm1}\Bigr)
\psrho^\circ(\gamma)\Phi=\#Y\cdot\Phi ,
\]
since $y\mapsto[y,\gamma]$ is a character of $\aubar{\pi_M}^{\,-1}L/L$, trivial exactly for
$\gamma\in\pi_ML$, where $\psrho^\circ(\gamma)\Phi=\Phi$. This is (a). For (b):
$a_y\theta(1\otimes\psrho_f(x)\psvac)=\sum_\gamma\theta(1\otimes\psrho_f(x)\psrho^\circ(\gamma)
\psrho_f(-y)\Phi)$, and each term is a root of unity times $\theta(1\otimes\psrho_f(x+\gamma-y)
\Phi)$, an algebraic integer by Proposition~\ref{pst:I2}(b).

\emph{Indicator data.} Let $u$ with $\aubar{\pi_M}u\in L$. The ideal $\pi_M\auO_K=\mathfrak p_0
\mathfrak p_1$ is a product of primes of degree one above $p_0,p_1$, which split in $K$; so $\pi_M$
is a unit at $\aubar{\mathfrak p_0},\aubar{\mathfrak p_1}$, $\aubar{\pi_M}$ is a unit at
$\mathfrak p_0,\mathfrak p_1$, and both are units at all other places. Hence $u\in L$ at every
place not above $\aubar{\mathfrak p_0}$ or $\aubar{\mathfrak p_1}$. At
$p=p_i$ write $W_{12}(\Q_p)=W_{\mathfrak p}\oplus W_{\aubar{\mathfrak p}}$ (the components of
$E\otimes\Q_p=E_{\mathfrak p}\times E_{\aubar{\mathfrak p}}$); both are isotropic for $\langle\,,
\rangle$, because $c$ exchanges the two factors, and $L_p=L_{\mathfrak p}\oplus L_{\aubar{\mathfrak
p}}$. In the Schr\"odinger model on $\mathcal S(W_{\mathfrak p})$, $\psrho(a)$ acts by translation
for $a\in W_{\mathfrak p}$, $\psrho(a')$ by multiplication by $r\mapsto\psi(\jmath\langle r,a'
\rangle)^{\pm1}$ for $a'\in W_{\aubar{\mathfrak p}}$, the local vacuum is $\mathbf 1_{L_{\mathfrak
p}}$ and the local component of $\phi_{y_0}$ is $\mathbf 1_{y_0+\pi_ML_{\mathfrak p}}$ (with $L_{\mathfrak
p}/\pi_ML_{\mathfrak p}=(\auO_E/\mathfrak p\auO_E)^{12}$). Write $u_p=u_{\mathfrak p}+u_{\aubar{
\mathfrak p}}$; then $u_{\mathfrak p}\in L_{\mathfrak p}$ and $u_{\aubar{\mathfrak p}}\in
\aubar{\pi_M}^{\,-1}L_{\aubar{\mathfrak p}}$. Up to a root of unity (a value of $\beta$),
$\psrho(u_p)\mathbf 1_{L_{\mathfrak p}}=\psrho(u_{\aubar{\mathfrak p}})\mathbf 1_{L_{\mathfrak
p}}$ is the function $r\mapsto\psi(\jmath\langle r,u_{\aubar{\mathfrak p}}\rangle)^{\pm1}\mathbf
1_{L_{\mathfrak p}}(r)$; for $r\in\pi_ML_{\mathfrak p}$ the pairing is $\langle\pi_M^{-1}r,
\aubar{\pi_M}u_{\aubar{\mathfrak p}}\rangle$ with both arguments integral, so the character is
invariant under $\pi_ML_{\mathfrak p}$, and the function is $\sum_{y_0}\zeta_{y_0}\mathbf 1_{y_0+
\pi_ML_{\mathfrak p}}$ with roots of unity $\zeta_{y_0}$. At every other place $\psrho(u_w)$ fixes
the local vacuum up to a root of unity ($u_w\in L_w$). Taking the tensor product gives the claim.
\end{proof}

\subsubsection{(I4): the vacuum constant}\label{pst:ss-vac}

\begin{proposition}[(I4)]\label{pst:I4}
There is $y\in W_{12}(\auA_f)$ with $c_{\mathrm{vac},y}:=\theta(1\otimes\psrho_f(y)\psvac^\circ)\ne0$, and for every
such $y$ and every $\sigma\in\auAut(\C)$, $|\sigma c_{\mathrm{vac},y}|\le\exp(12C_0q_1^2)\le e^b$, where $C_0=
C_0(q_0,d)$ is the constant \eqref{pz:eq-C0}.
\end{proposition}
\begin{proof}
\emph{Non-vanishing.} Suppose $\theta(1\otimes\psrho_f(y)\psvac)=0$ for all $y\in W_{12}(\auA_f)$.
For $\gamma\in W_{12}(E)$, $\theta\circ\psrho(\gamma)=\theta$, so $\theta(\psrho_\infty(\gamma)
\psrho_f(y)(1\otimes\psvac))$ is a root of unity times $\theta(1\otimes\psrho_f(y-\gamma)\psvac)=0$.
The image of $W_{12}(E)$ in $W_{12}\otimes\R$ is dense and $w\mapsto\theta(\psrho(w)(1\otimes
\psvac))$ is continuous, so it vanishes on $W_{12}(\auA)$, and $1\otimes\psvac$ lies in the kernel of
$F\mapsto\theta(\psrho(\cdot)F)$, a proper closed invariant subspace of an irreducible
representation (as in the proof of Lemma~\ref{pk:dict}); this is absurd. Since
$\psvac^\circ=c_{\mathrm{vac}}^{12}\psvac$ with $c_{\mathrm{vac}}\ne0$, some $c_{\mathrm{vac},y}\ne0$.

\emph{Rank one.} Let $\psvac^\circ_1$ be the integral vacuum of rank one \eqref{pk:eq-vac}, a sum
of $q_1^{2f}=[L_1:q_1L_1]$ vectors $\zeta_l\psrho(l)\psrho(b)F_0$ with roots of unity $\zeta_l$, where
$F_0$ is the theta product \eqref{pk:eq-product} with $\#\mathsf C=1$. For $x\in E(\auA_f)$,
$\theta(1\otimes\psrho_f(x)\psrho(l)\psrho(b)F_0)$ is a root of unity times $\theta(1\otimes
\psrho_f(\gamma)F_0)$ for a $\gamma\in E$ (\S\ref{pk:ss-conj}), which by \eqref{pk:eq-transl} at
$z=0$ is a root of unity times $\prod_{k<q_1}\theta_{\mathsf x_k\Omega_0}(0)^{q_1}$ with $\mathsf
x_k\in K$. Let $\sigma\in\auAut(\C/K)$. By Theorem~\ref{pf:recip}(b), comparing constant terms,
$\sigma(\theta_{\mathsf x\Omega_0}(0))=\theta^\sigma_{\mathsf x'\Omega_\sigma}(0)$ for suitable
$\mathsf x'\in K$, where $\theta^\sigma$ is the theta function of the normalised lattice
$\Lambda^\sigma$ with $\auO_K\Lambda^\sigma=\Lambda^\sigma$; so $|\sigma(\theta_{\mathsf x\Omega_0}
(0))|\le\mathsf M_d$ (\eqref{pz:eq-C0}). Roots of unity go to roots of unity. Hence
\[
\bigl|\sigma\,\theta(1\otimes\psrho_f(x)\psvac^\circ_1)\bigr|\le q_1^{2f}\mathsf M_d^{\,q_1^2}\le
\exp\bigl(2q_1\log q_1+q_1^2\log^+\mathsf M_d\bigr)\le\exp(C_0q_1^2),
\]
using $f\le q_1$, $\log q_1\le q_1$ and $C_0\ge2+\log^+\mathsf M_d$. If $\sigma|_K$ is complex
conjugation, $\sigma=c\circ\sigma'$ with $\sigma'\in\auAut(\C/K)$, and $|\sigma(\cdot)|=|\sigma'(
\cdot)|$; so the bound holds for all $\sigma\in\auAut(\C)$.

\emph{Rank twelve.} $W_{12}$ is the orthogonal sum of twelve lines, $L$ and $\psvac^\circ=
\bigotimes_b\psvac^\circ_1$ are the corresponding products (\eqref{pk:eq-vac}), and the theta
functional of an orthogonal sum is the tensor product of those of the summands; so $c_{\mathrm{vac},y}=\prod_{b=1}
^{12}\theta(1\otimes\psrho_f(y_b)\psvac^\circ_1)$ and $|\sigma c_{\mathrm{vac},y}|\le\exp(12C_0q_1^2)$.

\emph{Comparison with $b$.} By Lemma~\ref{pz:C0}(c), $C_0\le10^{-12}\log\log b\le\log x$ (by
(D2)), and $q_1\le x^{12}$ (Lemma~\ref{an:lem-sizes}(S3)); so $12C_0q_1^2\le12x^{24}\log x\le
x^{25}\le b$ by (S1).
\end{proof}

\emph{Correspondence with (I1)--(I4).} Proposition~\ref{pst:I1} is (I1). Proposition~\ref{pst:I2}
is (I2): its vector $\Phi$ has the stated theta function, its values $\theta(1\otimes\psrho_f(x)
\Phi)$ at all finite adelic points (in particular at the points $(0,x)$) are integral at every
prime, and (a) is the invariance of its theta function under the signed translations by $\pi_M
\mathfrak a_0^{12}$. Proposition~\ref{pst:I3} is (I3) with $u_i=y\in Y$ ($\aubar{\pi_M}u_i\in
\mathfrak a_0^{12}$) and $a_i=a_y$. Proposition~\ref{pst:I4} is (I4), the value $\psvac^\circ(0,
y)$ being $\theta(1\otimes\psrho_f(y)\psvac^\circ)$. Consequently:

\begin{corollary}\label{pst:attained}
Theorem~\ref{psm:attained} holds unconditionally.
\end{corollary}



%% ===== proofs/C-small3.tex =====
% proofs/C-small3.tex
% Parts (ii)--(iv) of Proposition po:small (the test form) and Proposition po:arch (the untwisted
% norm at all conjugates). To be \input after C-onelift, C-small, C-small2, C-exchanges3, C-periods.
% Labels of this file start with psu:. Only amsmath/amssymb/amsthm are used.
\providecommand{\auO}{\mathcal O}
\providecommand{\auU}{\mathrm U}
\providecommand{\auGL}{\mathrm{GL}}
\providecommand{\auHom}{\operatorname{Hom}}
\providecommand{\auAut}{\operatorname{Aut}}
\providecommand{\auNm}{\operatorname{Nm}}
\providecommand{\auA}{\mathbb A}
\providecommand{\auF}{\mathbb F}
\providecommand{\auPl}{\Phi}
\providecommand{\aubud}{\mathrm{bud}}
\providecommand{\ausinf}{s_\infty}
\providecommand{\ausfin}{s_{\mathrm{fin}}}
\providecommand{\aubar}[1]{\overline{#1}}
\providecommand{\plone}{\mathbf 1}
\providecommand{\psrho}{\varrho}
\providecommand{\psvac}{\mathbf v}
\providecommand{\poA}{\mathcal A}
\providecommand{\popr}{\operatorname{pr}}
\providecommand{\podel}{\delta}
\providecommand{\pop}{\mathsf p}
\providecommand{\pog}{\mathsf g}
\providecommand{\pexK}{K^{(z)}}
% macros of this file (prefix \psu)
\providecommand{\psuC}{\mathcal C}
\providecommand{\psuT}{\mathsf T}
\providecommand{\psuN}{\mathsf N}
\providecommand{\psug}{\hat g}

\subsection{The test form of least valuation, and the untwisted norm}\label{psu:sec}

This subsection completes the proof of Proposition~\ref{po:small} (in the exact form
Proposition~\ref{psm:small}) by constructing the test form from the sharpness of the count
(Theorem~\ref{psm:attained}, which holds unconditionally by Corollary~\ref{pst:attained}): parts
(ii)--(iv) of the note to Proposition~\ref{po:small}. It then treats Proposition~\ref{po:arch},
which is reduced to one precisely stated law of Galois conjugation of the lifts of characters
(\S\ref{psu:ss-arch}). Together with \S\ref{psm:sec} and \S\ref{pst:sec}, this proves Proposition~\ref{po:small}
(Theorem~\ref{psu:small}, with the constant $\exp(fmb/k+3x^{35})$, which suffices for
Propositions~\ref{po:mainheight} and~\ref{au:hecke}) and Proposition~\ref{po:arch} (Theorem~\ref{psu:arch}).

\emph{Conventions.} We use the notation of \S\ref{po:ss-hecke} and \S\ref{pex:ss-set}. Hermitian
forms are conjugate-linear in the first variable. $\psuC=\auU(1)(E^+)\backslash\auU(1)(\auA_{E^+,f})/
K^1_M$ is the finite abelian group of classes of $\auU(1)$ at the chosen level, $N_1=\#\psuC$, and
we fix representatives $x_c$ ($c\in\psuC$) with $x_{c_0}=1$ for the trivial class $c_0$. The groups
$\auU(1)(E^+)\cap x_cK^1_Mx_c^{-1}$ are trivial (proof of Lemma~\ref{pex:points}), and we use the
Haar measure on $\auU(1)(\auA_{E^+,f})$ with $\operatorname{vol}(K^1)=1$; hence for a function $F$ on
$[\auU(1)]$ that is right invariant under $K^1_M$,
\begin{equation}\label{psu:eq-int}
\int_{[\auU(1)]}F(g)\,dg=[K^1:K^1_M]^{-1}\sum_{c\in\psuC}F(x_c).
\end{equation}
The characters of $\psuC$ are the characters $\psi$ of $\auU(1)(E^+)\backslash\auU(1)(\auA_{E^+})$
that are trivial at the infinite places and on $K^1_M$; we write $\widehat\psuC$ for the group of
them and regard each $\psi$ as a function on $\auU(1)(\auA_{E^+,f})$. For a standard space $V'$ of
rank $2z'$, a level $K'$ and a class $h$ of $\Gamma\backslash\auU(V')(\auA_f)/K'$ put
$e_h=\#(\Gamma\cap hK'h^{-1})$; $\psuN_{z'}=\#\auGL_{4fz'}(\auF_3)$ is divisible by every $e_h$
(\S\ref{pex:ss-IM}).

\subsubsection{Flat pairs, their kernels and their values}\label{psu:ss-flat}

\begin{definition}[Flat pairs over a character]\label{psu:def-flat}
For an integer $2\le z'\le t$ and $\psi\in\widehat\psuC$ let $q(\psi,z')=(\upsilon^{-z'}\psi,z')$,
where $\upsilon$ is restricted to $\auU(1)$. For a tame character $\vartheta$ of modulus $Q$
(\S\ref{pex:ss-set}) write $\Psi_{q,\vartheta}$ for the block datum $\Phi^{(1)}_\vartheta$ of the pair
$q=q(\psi,z')$ on $\auHom(E,V')$, $V'=E^{2z'}$, and
\[
k_{q,\vartheta}(g,h)=\theta\bigl(\omega(g,h)(1\otimes\Psi_{q,\vartheta})\bigr)\qquad(g\in\auU(1)(\auA_{E^+}),
\ h\in\auU(V')(\auA_{E^+,f}))
\]
for its kernel, where $\omega$ is the representation \eqref{po:eq-w5} for the line and the character
$\nu_q\vartheta$, $\nu_q=\upsilon^{-z'}\psi$, and $1$ is the constant polynomial (the weight of $q$
is trivial).
\end{definition}

\begin{lemma}\label{psu:flat}
Let $q=q(\psi,z')$ as in Definition~\ref{psu:def-flat}.
\begin{enumerate}
\item[(a)] $q$ is a critical pair; it is flat ($c(\tau)=z'$ for all $\tau$, so $l_\tau=0$), and in the
notation of \S\ref{pex:ss-set} the character on $\auU(1)$ of its Weil representation is
$\alpha_{q,\vartheta}=\nu_q\vartheta\,\upsilon^{z'}=\psi\vartheta$.
\item[(b)] For every tame $\vartheta$ of modulus $Q$ and every $h$,
$\varphi_{q,\vartheta}(h)=c_{\mathrm{vac}}^{\,2z'}\sum_{c\in\psuC}k_{q,\vartheta}(x_c,h)$.
\item[(c)] For $g\in\auU(1)(\auA_{E^+,f})$ put $G_{q,\vartheta}(g)=\sum_he_h^{-1}\,\aubar{\varphi_{q,
\vartheta}(h)}\,k_{q,\vartheta}(g,h)$, the sum over the classes $h$ of $\auU(V')$ at the level of
$q$. Then $G_{q,\vartheta}(g)=\|\varphi_{q,\vartheta}\|^2/(N_1c_{\mathrm{vac}}^{\,2z'})$ for every $g$.
\end{enumerate}
\end{lemma}
\begin{proof}
(a) $\upsilon$ has weight $1$ and trivial tame part, and $\psi$ has weight $0$ and is tamely ramified
at the chosen places only (it is trivial on $K^1_M$); so $\nu_q$ has weight $-z'$ and $z'\le c=z'\le
t\le4b$ (Lemma~\ref{an:lem-sizes}). The formula for $\alpha$ is that of \S\ref{pex:ss-set}.

(b) By \eqref{po:eq-lift} with $y=\plone$, $\varphi_{q,\vartheta}(h)=c_{\mathrm{vac}}^{\,2z'}[K^1:K^1_M]
\int_{[\auU(1)]}k_{q,\vartheta}(g,h)\,dg$ (the weight is trivial, so $\pop_\xi$ is a constant and
$W_\lambda=\C$). The integrand is right invariant under $K^1_M$ (proof of Lemma~\ref{pex:points}: at
the chosen places $T\mapsto Tk$ fixes $\plone[T\equiv e_1]$ for $k\equiv1$, and at $w\in Q$ the factor
$\vartheta_w(k)^{-1}$ cancels $\psi_w(x_1k)/\psi_w(x_1)$). Apply \eqref{psu:eq-int}.

(c) Let $u\in\auU(1)(\auA_{E^+,f})$. The element $u$ of $\auU(1)$ and the scalar matrix $u^{-1}\cdot
1_{V'}$ of $\auU(V')$ act on $\auHom(E,V')$ by the same map $T\mapsto Tu=uT$, so their images in the
symplectic group of $\auHom(E,V')$ coincide; hence $\omega(u,1)=\epsilon(u)\,\omega(1,u^{-1})$ for a
scalar $\epsilon(u)$, which is multiplicative in $u$ and of absolute value $1$ (both operators are
unitary). The scalar $u^{-1}$ is central, so right translation by it permutes the classes $h$ and
preserves the numbers $e_h$. From $k(ug,h)=\epsilon(u)k(g,hu^{-1})$ (write $k=k_{q,\vartheta}$, $\varphi
=\varphi_{q,\vartheta}$) and (b), $\varphi(hu^{-1})=c_{\mathrm{vac}}^{2z'}\sum_ck(x_c,hu^{-1})=
\epsilon(u)^{-1}c_{\mathrm{vac}}^{2z'}\sum_ck(ux_c,h)=\epsilon(u)^{-1}\varphi(h)$, because $c\mapsto
ux_c$ runs again over representatives of $\psuC$ and $k(\cdot,h)$ is left invariant under
$\auU(1)(E^+)$ and right invariant under $K^1_M$. Hence
\[
G(ug)=\epsilon(u)\sum_he_h^{-1}\aubar{\varphi(h)}k(g,hu^{-1})=\epsilon(u)\sum_he_h^{-1}
\aubar{\varphi(hu)}k(g,h)=\epsilon(u)\aubar{\epsilon(u)}\,G(g)=G(g),
\]
using $\varphi(hu)=\epsilon(u^{-1})^{-1}\varphi(h)=\epsilon(u)\varphi(h)$. So $G$ is constant, and
$N_1G=\sum_cG(x_c)=\sum_he_h^{-1}\aubar{\varphi(h)}\sum_ck(x_c,h)=\|\varphi\|^2/c_{\mathrm{vac}}^{2z'}$ by
(b) and the definition of the norm.
\end{proof}

\begin{definition}[Integral values]\label{psu:def-I}
For $q=q(\psi,z')$ and a tame character $\vartheta$ of modulus $Q$ put, with $n_q=M^{4fz'}$,
\[
I_q(\vartheta)=\psuN_{z'}\,n_q^2\,\frac{\|\varphi_{q,\vartheta}\|^2}{(\auNm Q)^{2z'-1}} .
\]
\end{definition}

\begin{lemma}[Integral values]\label{psu:I}
Let $q=q(\psi,z')$.
\begin{enumerate}
\item[(a)] $I_q(\vartheta)>0$, and $I_q(\vartheta)=I_q(1)\,\rho_q(\vartheta)$ for every twist $\vartheta$.
\item[(b)] Let $\ell\nmid\auNm Q$ and $\jmath\colon\C\to\C_\ell$. Then $n_q\varphi_{q,\vartheta}$ is
integral at $\jmath$, the numbers $(\psuN_{z'}/e_h)\,n_q\,\varphi_{q,\vartheta}(h)$ are integral at
$\jmath$, and $I_q(\vartheta)$ is integral at $\jmath$. In particular $I_q(1)$ is integral at every
prime.
\item[(c)] $|\sigma I_q(1)|\le\exp(x^{35})$ for all $\sigma\in\auAut(\C)$.
\end{enumerate}
\end{lemma}
\begin{proof}
(a) Proposition~\ref{pex:rallis} for the pair $q$ ($z'\ge2$): $\|\varphi_{q,\vartheta}\|^2=A_q(\auNm
Q)^{2z'-1}L^{\mathsf M\cup Q}(z',(\nu_q\vartheta)^{-1})$ with $A_q>0$ and a positive $L$-value; divide
by the case $\vartheta=1$.

(b) For $l_\tau=0$ the integer of the kernel integrality of \S\ref{pex:ss-IM} may be taken to be
$n_q=M^{4fz'}$ (the factor $d^{\sum l_\tau}$ of Lemma~\ref{pex:dint} is $1$), and that statement holds at
every $\ell\nmid\auNm Q$. Indeed, by the finite Fourier analysis recalled in the proof of (c) below (and
the same at the places of $Q$, where the index $q_w$ is a unit at $\ell$), $M^{2fz'}\auNm Q\cdot
c_{\mathrm{vac}}^{2z'}\Phi^0_Q$ is a sum, with roots of unity as coefficients, of vectors
$\psrho^\circ(u)\psvac^\circ$ with $\aubar{\pi_M}u$ integral above $\ell$, that is, a standard datum of
level $\aubar{\pi_M}$ at $\ell$; the reduction in the proof of Proposition~\ref{pk:K12} (weak
approximation, and the fact of \S\ref{pf:ss-frames} that the Weil operators of elements that are
integral with integral inverse above $\ell$ preserve standard data up to units) shows that
$M^{2fz'}c_{\mathrm{vac}}^{2z'}k_\Phi(x,h;1)$, $\Phi=\Phi^0_Q$, is a unit at $\ell$ times the value
$\theta(1\otimes f)$ of a standard datum $f$, and this value is integral at every prime by
Lemma~\ref{pf:fourier}(a) (the jets are of degree $0$, so no condition on $\ell$ is needed). By
Lemma~\ref{pex:points}, $\varphi_{q,\vartheta}(h)$ is a sum of values of $\vartheta^{-1}$
(roots of unity) times kernels $k^0_Q(x,h)$, so $n_q\varphi_{q,\vartheta}(h)$ is integral at $\jmath$
(the weight is trivial, so the gauge is the absolute value). Since $e_h\mid\psuN_{z'}$, the second
assertion follows. For the third write $\|\varphi\|^2=\sum_he_h^{-1}\aubar{\varphi(h)}\varphi(h)$: the
number $n_q\aubar{\varphi(h)}$ is the complex conjugate of $n_q\varphi(h)$, which is integral at
$\jmath\circ c$ by the first assertion applied to $\jmath\circ c$, so it is integral at $\jmath$; and
$\auNm Q$ is a unit at $\ell$. (This is the proof of Proposition~\ref{pex:IM}(b), which does not use
that the pair is ramified.)

(c) Replacing $\sigma$ by $c\circ\sigma$ does not change $|\sigma I_q(1)|$, so we may assume $\sigma|_K=
1$. At each chosen place $w$, finite Fourier analysis on $(\auO_w/w)^{2z'}$ writes the indicator of
$e_1+w\auO_w^{2z'}$ as $q_w^{-2z'}$ times a sum of $q_w^{2z'}$ Heisenberg translates of the local
vacuum by vectors $u_w$ with $\varpi_wu_w$ integral, with roots of unity as coefficients. The chosen
places are the $2f$ primes of $E$ above $p_0$, $p_1$, of norms $p_0$, $p_1$; so
$c_{\mathrm{vac}}^{2z'}\Psi_{q,1}=M^{-2fz'}\sum_u\zeta_u\psrho^\circ(u)\psvac^\circ$, a sum of
$M^{2fz'}$ translates of the integral vacuum with roots of unity $\zeta_u$. By Step~1 of the proof of
Proposition~\ref{au:kernel} (\S\ref{ps:ss-kernel}), $\omega(g,h)\psrho^\circ(u)\psvac^\circ=\epsilon
\,\psrho^\circ(u')\omega(k)\psvac^\circ$ with $|\sigma\epsilon|=1$ and $k\in\auU(\auHom(E,V'))
(\auA_f)$; and by Lemma~\ref{ps:conj}(i) in degree $0$ and Lemma~\ref{ps:conj}(ii) at $y=0$,
$|\sigma\theta(1\otimes f)|=|\theta(1\otimes\Lambda_\sigma f)|\le\exp(C_0(2z'q_1+1)^2)$ for every such
translate $f$. Hence $|\sigma(c_{\mathrm{vac}}^{2z'}k_{q,1}(g,h))|\le\exp(C_0(2z'q_1+1)^2)$, and by
(b), $|\sigma\varphi_{q,1}(h)|\le N_1\exp(C_0(2z'q_1+1)^2)$; the same holds for $\sigma'=c\sigma c$.
Since $\sigma(\aubar a\,b)=\aubar{\sigma'(a)}\,\sigma(b)$, $|\sigma\|\varphi_{q,1}\|^2|\le h_qN_1^2
\exp(2C_0(2z'q_1+1)^2)$, with $h_q$ the number of classes of the level of $q$. By
Proposition~\ref{pf:classes} (the level contains the principal congruence subgroup modulo $M$, and
$2z'\le2t\le n_m$; also for $\auU(1)$), $h_q,N_1\le\exp(x^{30})$. Therefore
\[
\log|\sigma I_q(1)|\le(4fz')^2\log3+8fz'\log M+3x^{30}+2C_0(2z'q_1+1)^2\le x^{35},
\]
by the inequality (N5) of Lemma~\ref{psu:sizes} below.
\end{proof}

\begin{lemma}[Sizes]\label{psu:sizes}
For $1\le z'\le t$:
\textup{(N5)} $(4fz')^2\log3+8fz'\log M+3x^{30}+2C_0(2z'q_1+1)^2\le x^{35}$;
\textup{(N6)} $fmb/k+3x^{35}\le\aubud$ and $x^{35}+2f\le\aubud$, where $\aubud=fxb(\log b)^3/(mk)$.
\end{lemma}
\begin{proof}
By Lemma~\ref{an:lem-sizes}: $t=5r\le15x^5$, $f,q_1\le x^{12}$ (S3), $\log M\le22\log b\le22x^{1/100}$
(S1), (S5), and $C_0\le\log\log b\le\log x$ ((D2) and \eqref{an:eq-CP}). So $(4fz')^2\log3\le
3600\cdot1.1\,x^{34}$, $8fz'\log M\le x^{19}$, and $2C_0(2z'q_1+1)^2\le2\log x\cdot(31x^{17})^2\le
2000\,x^{34}\log x$. Since $\log x\ge10^{14}$ (S1), $4000+2000\log x+3+1\le x$, which gives (N5).
For (N6): $\aubud/(fmb/k)=x(\log b)^3/m^2\ge(\log b)^3/4\ge2$ (as $m^2\le4x$ and $\log b\ge128$), so
$fmb/k\le\aubud/2$; and $fmb/k\ge b\ge x^{10^6}\ge6x^{35}+2f$ (as $f\ge k$ and $m\ge1$, and (S1)),
so $fmb/k+3x^{35}\le2fmb/k\le\aubud$ and $x^{35}+2f\le\aubud$.
\end{proof}

\subsubsection{Averages over the classes of $\auU(1)$}\label{psu:ss-avg}

\begin{lemma}[Orthogonality]\label{psu:orth}
Let $s$ be a function on $\psuC$, $\ell$ a prime and $\jmath\colon\C\to\C_\ell$. For $\psi\in
\widehat\psuC$ put $A_\psi=N_1^{-1}\sum_{c\in\psuC}\psi(x_c)s(c)$. Then $\sum_\psi A_\psi=s(c_0)$, and
some $\psi$ has $|A_\psi|_\jmath\ge|s(c_0)|_\jmath$.
\end{lemma}
\begin{proof}
$\sum_\psi\psi(x_c)$ is $N_1$ for $c=c_0$ and $0$ otherwise (orthogonality of the characters of the
finite abelian group $\psuC$). If all $|A_\psi|_\jmath$ were smaller than $|s(c_0)|_\jmath$, the
ultrametric inequality would give $|\sum_\psi A_\psi|_\jmath<|s(c_0)|_\jmath$.
\end{proof}

\begin{lemma}[Multiplicative pigeonhole]\label{psu:pigeon}
Let $G$ be a finite group, $a,b\colon G\to\C$, $V\in\C$ with $V^{\#G}=\prod_{\chi\in G}a(\chi)b(\chi)$,
and $|\cdot|$ an absolute value on $\C$ (archimedean or of the form $|\jmath(\cdot)|$). For every
$\psi\in G$ there are $\phi,\chi\in G$ with $\phi\chi=\psi$ and $|V|\le|a(\phi)b(\chi)|$.
\end{lemma}
\begin{proof}
$\chi\mapsto\psi\chi^{-1}$ is a bijection of $G$, so $\prod_\chi|a(\psi\chi^{-1})b(\chi)|=|V|^{\#G}$,
and the largest factor is at least the geometric mean $|V|$.
\end{proof}

\begin{lemma}[The geometric mean]\label{psu:mean}
Let $z\ge12$. For each $\chi\in\widehat\psuC$ let $q_3(\chi)=q(\chi,2)$ and $q_1(\chi)=q(\chi,z-10)$
(Definition~\ref{psu:def-flat}). There is $L\in\C^\times$ with
\[
L^{N_1}=\prod_{\chi\in\widehat\psuC}I_{q_3(\chi)}(1)\,I_{q_1(\chi)}(1)\qquad\text{and}\qquad
|\sigma L|\le\exp(2x^{35})\quad(\sigma\in\auAut(\C)).
\]
\end{lemma}
\begin{proof}
$2\le z-10\le t$ because $z\le r$. Let $L$ be any $N_1$-th root of the product, which is non-zero by
Lemma~\ref{psu:I}(a) ($\#\widehat\psuC=N_1$). For $\sigma\in\auAut(\C)$, the number $\sigma L$ is an
$N_1$-th root of $\prod_\chi\sigma(I_{q_3(\chi)}(1))\sigma(I_{q_1(\chi)}(1))$; Lemma~\ref{psu:pigeon} with
the usual absolute value and $\psi=1$ gives $\phi,\chi$ with $|\sigma L|\le|\sigma I_{q_3(\phi)}(1)|\,
|\sigma I_{q_1(\chi)}(1)|\le\exp(2x^{35})$ by Lemma~\ref{psu:I}(c).
\end{proof}

\subsubsection{The point form and the seesaw factorisation}\label{psu:ss-seesaw}

Fix a main pair $p=(\nu,z)$ for $j$ (Definition~\ref{po:main}), with $\nu$ of weight $-c$, and put
$l_\tau=c(\tau)-z$. Then
\begin{equation}\label{psu:eq-range}
\kappa(\tau)\le l_\tau\le\kappa(\tau)+x^{42}\qquad(\tau\in\auPl):
\end{equation}
from $\kappa_j(\tau)+r-t\le c(\tau)\le\kappa_j(\tau)$ and $m<z\le r$ we get $l_\tau\ge\kappa_j(\tau)-t
=\kappa(\tau)+(j+\frac12)\Delta-t\ge\kappa(\tau)$ (as $\Delta=2r^8\ge5r=t$) and $l_\tau\le\kappa_j(\tau)
\le\kappa(\tau)+(m+\frac12)\Delta\le\kappa(\tau)+x^{42}$, since $(m+\frac12)\Delta=(2m+1)r^8\le(4\sqrt x+1)\,3^8x^{40}\le x^{42}$ by $m\le2\sqrt x$ and $r\le3x^5$ (Definition~\ref{an:def-data}(D3)) and $x$ large. Also
$z>m\ge10^4$, so $2z-20\ge4$ and $2\cdot1+2\cdot20+3<2z$.

\emph{Blocks.} Write $V=V^{(z)}=E^{2z}=V_0\perp V_2\perp V_3\perp V_1$ for the decomposition of the
standard coordinates into consecutive blocks of ranks $12$, $4$, $4$, $2z-20$; $e^{(i)}_1$ is the
first basis vector of $V_i$, and $e_{13}=e^{(2)}_1$. For $h_i\in\auU(V_i)(\auA_f)$ let
$[h_0,h_2,h_3,h_1]$ be the block diagonal element. Fix a twist $\vartheta$, of modulus $Q$, and a
class $x_0\in(\auO_E/\pi_M)^{12}$.

\emph{The point.} Let $\psug=\psug(x_0)\in\auU(V)(\auA_{E^+,f})$ be integral, equal to $1$ at the places
not below $\mathsf M\cup Q$, such that at each chosen place $w$ (where $\auU(V)(E^+_v)=\auGL_{2z}(E_w)$,
$v$ below $w$) $\psug_we_1\equiv Y:=(x_0,e^{(2)}_1,e^{(3)}_1,e^{(1)}_1)\bmod w$, and at each $w\in Q$
$\psug_w$ is the permutation matrix of the transposition of $e_1$ and $e_{13}$. It exists: $Y\bmod w$
is non-zero (its last block is), so it is the first column of an element of $\auGL_{2z}(\auO_w/w)$,
which lifts to $\auGL_{2z}(\auO_w)$.

\emph{The data.} Let $\Psi_0$ be the datum on $\auHom(E,V_0)$ that is $\plone[T\equiv x_0\bmod w]$ at
the chosen places and the vacuum at all other places (it is the indicator datum $\phi_{x_0}$ of
\S\ref{psm:ss-attained}). For a twist $\vartheta$ (or $\vartheta=1$) and $\psi_3,\psi_1\in\widehat\psuC$
let $q_2=q(\chi_0,2)=(\upsilon^{-2}\chi_0,2)$ (the pair of Proposition~\ref{po:small}), $q_3=q(\psi_3,2)$
and $q_1=q(\psi_1,z-10)$, on $V_2$, $V_3$, $V_1$, with the kernels $k_2=k_{q_2,\vartheta}$,
$k_3=k_{q_3,1}$, $k_1=k_{q_1,1}$ of Definition~\ref{psu:def-flat}. Put $\psi=\psi_3\psi_1$ and
\[
\chi'=\nu\,\upsilon^{z}\chi_0^{-1}\quad\text{on }\auU(1),\qquad s(g)=\theta\bigl(\omega_{12}(g,1)(\pop_\alpha
\otimes\Psi_0)\bigr)\quad(g\in\auU(1)(\auA_{E^+})),
\]
where $\omega_{12}$ is the Weil representation of the pair $(E,V_0)$ in the normalisation of
\eqref{po:eq-w5} with $\auU(1)$-character $\chi'$ (in the notation of \S\ref{pex:ss-set}: $\alpha=\chi'$,
$\beta=1$), and $\pop_\alpha$ is the monomial $z^\alpha=\prod_{\tau}\prod_{b\le12}z_{\tau,b}^{\alpha(\tau,b)}$
in the Fock variables of $\auHom(E,V_0)$, for an exponent $\alpha\in\N^{\auPl\times12}$ with
$\sum_b\alpha(\tau,b)=l_\tau$. The function $s$ is left invariant under $\auU(1)(E^+)$ and right
invariant under $K^1_M$ (at the chosen places $T\mapsto Tk$ fixes the indicator for $k\equiv1$; at the
other finite places $\chi'$ is unramified and the datum is the vacuum), so it is a function on
$\psuC$.

\begin{lemma}[The vector]\label{psu:vector}
Let $\alpha$ be as above and $a=\prod_\tau l_\tau!/\prod_{\tau,b}\alpha(\tau,b)!\in\Z_{>0}$. There is
$v\in W_\lambda(\Z)$ (the weight of $\poA(V;p,\vartheta)$: one row of length $l_\tau$ at $\tau$) with
$\pop_v=a\,z^\alpha$, the monomial in the variables of the first block; $v$ is fixed by every element of
$\prod_\tau\auGL_{2z}$ (over any field of characteristic $0$, and over $\auO_{\C_\ell}$) that fixes
$e_1,\dots,e_{12}$.
\end{lemma}
\begin{proof}
At $\tau$ the weight is $\operatorname{Sym}^{l_\tau}\C^{2z}\subseteq(\C^{2z})^{\otimes l_\tau}$, the span of
the vectors $\mathrm{ev}(y)=y^{\otimes l_\tau}$ of \eqref{po:eq-ev} (one column of height $1$ for each
box), and $W_\lambda(\Z)$ is its intersection with the integral tensors. Let $v_\tau$ be the sum of
the $l_\tau!/\prod_b\alpha(\tau,b)!$ distinct basis tensors $e_{b_1}\otimes\dots\otimes e_{b_{l_\tau}}$ in
which each $b\le12$ occurs $\alpha(\tau,b)$ times; it is symmetric and integral, and $v=\bigotimes_\tau
v_\tau$. By \eqref{po:eq-poly} the polynomial of a basis tensor is the product of the corresponding
Fock variables, so $\pop_v=\prod_\tau(l_\tau!/\prod_b\alpha(\tau,b)!)\prod_bz_{\tau,b}^{\alpha(\tau,b)}=
az^\alpha$. An element fixing $e_1,\dots,e_{12}$ (acting by its tensor powers, or by those of its
complex conjugate) fixes each such basis tensor.
\end{proof}

\begin{lemma}[Seesaw factorisation]\label{psu:factor}
With the notation above, let $x=[1,h_2,h_3,h_1]\,\psug$ and let $k_p(g,x;v)$ be the kernel of
\S\ref{po:ss-set} for the lift of $p$ at $\vartheta$ (the line $E$, the space $V$, the character
$\nu\vartheta$, the block datum $\Phi^{(1)}_\vartheta$). Then for $g\in\auU(1)(\auA_{E^+})$
\[
k_p(g,x;v)=a\,\psi(\det g)\,s(g)\,k_2(g,h_2)\,k_3(g,h_3)\,k_1(g,h_1).
\]
\end{lemma}
\begin{proof}
\emph{The datum.} $\psug$ is integral, so the unramified scalar of \eqref{po:eq-w5} is $1$ and
$\omega(1,\psug)\Phi(T)=\Phi(\psug^{-1}T)$ at the places of $\mathsf M\cup Q$ (at the conjugate places
$\psug$ is integral and fixes the vacuum). At a chosen place, $\plone[\psug^{-1}T\equiv e_1]=\plone[T
\equiv Y]=\plone[T_0\equiv x_0]\,\plone[T_2\equiv e^{(2)}_1]\,\plone[T_3\equiv e^{(3)}_1]\,\plone[T_1\equiv
e^{(1)}_1]$, $T_i$ the block of $T$ in $V_i$. At $w\in Q$, $\plone[T\text{ integral}]\psi_w((\psug^{-1}
T)_1)=\plone[T\text{ integral}]\psi_w(T_{13})$, and $T_{13}$ is the first coordinate of $T_2$. Hence
$\omega(1,\psug)\Phi^{(1)}_\vartheta=\Psi_0\otimes\Psi_{q_2,\vartheta}\otimes\Psi_{q_3,1}\otimes\Psi_{q_1,1}$.

\emph{The characters.} The $\auU(1)$-character of the Weil representation of $p$ at $\vartheta$ is
$\nu\vartheta\upsilon^z$ (\S\ref{pex:ss-set}), and those of the four blocks are $\chi'\psi^{-1}$ and, by
Lemma~\ref{psu:flat}(a), $\chi_0\vartheta$, $\psi_3$, $\psi_1$; their product is $\nu\vartheta\upsilon^z$.

\emph{The factorisation.} By (W7) of \S\ref{po:ss-glue}, which applies verbatim to the line $E$ and
the orthogonal decomposition $V=V_0\perp V_2\perp V_3\perp V_1$ (the characters multiply), the Weil
representation, the theta functional and the Fock vectors factor over the blocks. By
Lemma~\ref{psu:vector}, $\pop_v=a\,z^\alpha\otimes1\otimes1\otimes1$. Since $[1,h_2,h_3,h_1]$ acts on the
blocks separately and $\omega(g,hh')=\omega(g,h)\omega(1,h')$,
\[
k_p(g,x;v)=a\,\theta_{12}\bigl(\omega_{12}'(g,1)(z^\alpha\otimes\Psi_0)\bigr)\,k_2(g,h_2)\,k_3(g,h_3)\,
k_1(g,h_1),
\]
with $\omega'_{12}$ the representation of $(E,V_0)$ with character $\chi'\psi^{-1}$. By
\eqref{po:eq-w5} the action of $\auU(1)$ carries the factor $\alpha(\det g)^{-1}$ of its character
$\alpha$, so $\omega_{12}'(g,1)=\psi(\det g)\,\omega_{12}(g,1)$, which gives the factor $\psi(\det g)s(g)$.
\end{proof}

\subsubsection{The test form}\label{psu:ss-test}

Keep $p$, $\vartheta$, $x_0$, $\alpha$, $v$, $\psi_3$, $\psi_1$ and the three flat pairs. For
$i\in\{2,3,1\}$ let $\vartheta_2=\vartheta$, $\vartheta_3=\vartheta_1=1$, $z_2=z_3=2$, $z_1=z-10$, and
$\psuT_i(h)=\psuN_{z_i}n_{q_i}^2(\auNm Q_i)^{1-2z_i}\varphi_{q_i,\vartheta_i}(h)$, where $Q_2=Q$ and
$Q_3=Q_1=\emptyset$. Define
\begin{equation}\label{psu:eq-Xi}
\Xi=\sum_{h_2}\sum_{h_3}\sum_{h_1}\frac{\psuT_2(h_2)\,\psuT_3(h_3)\,\psuT_1(h_1)}{e_{h_2}e_{h_3}e_{h_1}}\,
\podel\bigl[[1,h_2,h_3,h_1]\psug;\,v\bigr]\ \in\poA(V;p,\vartheta),
\end{equation}
the sums over the classes of $\auU(V_i)$ at the levels of the pairs $q_i$ (twisted by $\vartheta_i$),
with $\podel[\cdot;\cdot]$ as in \S\ref{po:ss-glue} (for the line $E$ and the space $V$).

\emph{The summand depends only on the classes.} Let $\gamma\in\auU(V_i)(E^+)$ and $k$ in the level of
$q_i$ at $\vartheta_i$, and replace $h_i$ by $\gamma h_ik$. The element $\gamma$ (placed in the block
$V_i$) is a rational point of $\auU(V)$ that fixes $e_1,\dots,e_{12}$, hence fixes $v$
(Lemma~\ref{psu:vector}), and $\varphi_{q_i}(\gamma h_i)=\varphi_{q_i}(h_i)$ (trivial weight). Next
$[\dots,k,\dots]\psug=\psug k'$ with $k'=\psug^{-1}[\dots,k,\dots]\psug$. At a chosen place $k$ fixes
$e^{(i)}_1$ modulo $w$, so $k'e_1\equiv\psug^{-1}[\dots,k,\dots]Y\equiv\psug^{-1}Y\equiv e_1$; at $w\in Q$,
$k'$ is $k$ with the coordinates $1$ and $13$ exchanged, so for $i\ne2$ it does not touch the first
row and column, and for $i=2$ its first row is the first row of $k\bmod w$ placed in the coordinates
$1,14,15,16$, which is $((k)_{11},0,0,0)$ by the definition of the level of $q_2$ at $\vartheta$. So
$k'\in\pexK_\vartheta$, and the character of $k'$ (the value at $k'$ of the character by which the level
acts on $\Phi^{(1)}_\vartheta$: trivial at the chosen places, $\psi_w$ of the $(1,1)$ entry at $w\in Q$)
equals the value at $k$ of the character of the level of $q_i$ at $\vartheta_i$. By the rule
$\podel[\gamma xk;\gamma v]=\chi(k)^{-1}\podel[x;v]$ of \S\ref{po:ss-glue} and $\varphi_{q_i}(hk)=
\chi_{q_i}(k)\varphi_{q_i}(h)$ (\S\ref{ss:au-spaces}), the two characters cancel; $e_{\gamma h_ik}=
e_{h_i}$. Hence \eqref{psu:eq-Xi} is well defined.

\begin{lemma}[Integrality]\label{psu:integral}
Let $\ell\nmid\auNm Q$. Then $\Xi$ is integral at every $\jmath\colon\C\to\C_\ell$.
\end{lemma}
\begin{proof}
By Lemma~\ref{psu:I}(b) (for $q_2$ at $\vartheta$, $\ell\nmid\auNm Q$, and for $q_3$, $q_1$ at $1$) the
coefficients $\psuT_i(h_i)/e_{h_i}$ are integral at $\jmath$ ($\auNm Q$ is a unit at $\ell$). By
\eqref{po:eq-deltasize} the form $\podel[x;v]$ has gauge at most the gauge of $v$ at $(\jmath,x)$, the
largest absolute value of a coordinate of $x_\ell^{-1}v$. For $x=[1,h_2,h_3,h_1]\psug$,
$x_\ell^{-1}v=\psug_\ell^{-1}[1,h_2,h_3,h_1]_\ell^{-1}v=\psug_\ell^{-1}v$ by Lemma~\ref{psu:vector}, and
$\psug_\ell$ is integral with integral inverse, so the coordinates of $\psug_\ell^{-1}v$ are integral.
Sums of integral forms with integral coefficients are integral.
\end{proof}

\begin{proposition}[Periods of false lifts in rank one]\label{psu:periods}
Let $w$ be off the level and $y\in\varrho(e_1)\poA(V;p,\vartheta)\cap\ker P_w(X_w)\cap\ker
D^{(1)}_\vartheta$ (Definition~\ref{po:R3}). Write $V=V'\perp V_1$ with $V'=V_0\perp V_2\perp V_3$ of
rank $20$. Let $q_1=q(\psi_1,z-10)$ and $k_1=k_{q_1,1}$. Let $x'\in\auU(V)(\auA_f)$ be integral with
integral inverse, equal to $1$ at the places not below $\mathsf M\cup Q$, and such that for every $k$ in
the level of $q_1$ the element $k'=x'^{-1}[1,k]\,x'$ lies in $\pexK_\vartheta$ and acts on
$\Phi^{(1)}_\vartheta$ by the scalar by which $k$ acts on $\Psi_{q_1,1}$ (for $x'=\psug$ this was shown
after \eqref{psu:eq-Xi}); under this condition the summand below depends only on the class of $h_1$.
Then for every $h'\in\auU(V')(\auA_f)$, every
$u\in W_\lambda$ whose polynomial $\pop_u$ involves only the variables of $V'$, and every
$g\in\auU(1)(\auA_{E^+,f})$,
\[
\sum_{h_1}e_{h_1}^{-1}\,\aubar{k_1(g,h_1)}\,\bigl\langle u,\,y\bigl([h',h_1]x'\bigr)\bigr\rangle=0,
\]
the sum over the classes of $\auU(V_1)$ at the level of $q_1$.
\end{proposition}
\begin{proof}
This is Theorem~\ref{ppe:main} for the line $W=E$ (rank $r=1$, definite), $V=V^{(z)}$ of rank $n=2z$
(so $t=n-r=2z-1$), the character $\nu\vartheta$, the idempotent $e_1$ in place of $e_0$, the splitting
$V=V'\perp V_1$ in place of $V''\perp V^\circ$ (with $n_1=20$), the element $x'$ in place of $\pog_s$, the
datum $\Psi_{q_1,1}$ and the trivial weight on $V_1$ in place of $\Psi^\circ$ and $\pop^\circ$, and the
level of $q_1$ in place of $J^\circ$. The steps of the proof of \S\ref{ppe:sec} are gone through one by
one in Remark~\ref{ppe:rem-rankone}, steps (1)--(8), which says for each step what replaces what, which
steps are unchanged, and proves the others (the property (F2), the generation by harmonics, and the
final step). Here we give the points where the rank, the idempotent and the data enter.
(F1) $D^{(1)}_\vartheta y=0$ is part of the hypothesis (it replaces Proposition~\ref{plf:L4}).
(F2) For $r=1$ and a chosen place $v$, $\lambda\in\auF_v^\times$ and $\psuN_\lambda=\{1+u\lambda\}=
\auF_v^\times$, so the hyperplane sum is $\eta=\sum_{s\in\auF_v^\times}\psi_v(s)^{-1}s$, where $\psi_v$ is
the tame part of $\nu\vartheta$ at $v$, namely $\omega_v^{\,j}$ ($\vartheta$ is unramified at $v$),
which is non-trivial because $1\le j\le m$ and $\omega_v$ has order greater than $8m$; since $y=\varrho
(e_1)y$ and $\eta e_1=(\sum_s\psi_v(s)^{-1})e_1=0$, $\varrho(\eta)y=0$ (this replaces
Lemma~\ref{ppe:green}). (F3) $P_w(X_w)y=0$ is part of the hypothesis, with $X_w$ of
Lemma~\ref{po:facts}(L3) for the line, $|X_w|=q^{-1/2-(t-1)}$. The tower
(Proposition~\ref{ppe:tower}) is stated for an arbitrary rank $r$ and needs $2r+2p+3<n$, here for
$p=n_1=20$: $45<2z$. Proposition~\ref{ppe:rallis} is stated for an arbitrary splitting of the standard
coordinates. In the final step of the proof of Theorem~\ref{ppe:main}, $\omega(1,x'^{-1})$ preserves
the data that are the vacuum outside $\Sigma=\mathsf S_\vartheta\cup\aubar{\mathsf S_\vartheta}$ because $x'$
is $1$ outside $\Sigma$; the factorisation of the kernel uses (W7) for $V'\perp V_1$ and
$\pop_u=\pop'\otimes1$, the characters on $\auU(1)$ being immaterial (an extra character of $\det g$ only
multiplies descents, \S\ref{ppe:ss-setup}). The function $h_1\mapsto\aubar{k_1(g,h_1)}\,\langle u,
\rho(h_{1,\infty})^{-1}y([h',h_{1,f}]x')\rangle$ on $\auU(V_1)(\auA)$ is invariant under
$\auU(V_1)(E^+\otimes\R)$ and under $\auU(V_1)(E^+)$ on the left, because $u$ is fixed by the block of
$V_1$ (its polynomial involves only the variables of $V'$) and the weight of $q_1$ is trivial; it is
right invariant under the level of $q_1$ by the condition on $x'$: with $s(k)$ the scalar by which $k$
acts on $\Psi_{q_1,1}$, $k_1(g,h_1k)=s(k)k_1(g,h_1)$ and $y(\,\cdot\,k')=s(k)\,y(\,\cdot\,)$, and
$|s(k)|=1$. Hence the weighted sum over the classes of the level, which
need not be free, is the integral over $[\auU(V_1)]$ divided by the volume of the level, as noted at the end of
Theorem~\ref{ppe:main}. The conclusion of that proof is the vanishing of the integral of $\theta^\circ
\cdot\aubar{c_\zeta y}$, whose complex conjugate is \eqref{ppe:eq-goal}, the integral form of the
displayed sum (with $\langle u,y(\cdot)\rangle=c_u y$ on finite points and $\theta^\circ=k_1$).
\end{proof}

\begin{lemma}[Orthogonality to the false lifts]\label{psu:orthfalse}
For every $w$ off the level, $\Xi$ is orthogonal to $\varrho(e_1)\poA(V;p,\vartheta)\cap\ker P_w(X_w)
\cap\ker D^{(1)}_\vartheta$.
\end{lemma}
\begin{proof}
Let $y$ be in this space. By \eqref{po:eq-delta}, $\langle\podel[x;v],y\rangle=\langle v,y(x)\rangle$,
so $\langle\Xi,y\rangle$ is a sum over $h_2,h_3$ of the conjugates of the coefficients
$\psuT_2(h_2)\psuT_3(h_3)/(e_{h_2}e_{h_3})$ times
\[
\sum_{h_1}e_{h_1}^{-1}\aubar{\psuT_1(h_1)}\,\langle v,y([h',h_1]\psug)\rangle
=\psuN_{z-10}n_{q_1}^2\,\aubar{c_{\mathrm{vac}}^{2z-20}}\sum_{c\in\psuC}\ \sum_{h_1}e_{h_1}^{-1}
\aubar{k_1(x_c,h_1)}\,\langle v,y([h',h_1]\psug)\rangle ,
\]
with $h'=[1,h_2,h_3]$, by Lemma~\ref{psu:flat}(b). Each inner sum vanishes by
Proposition~\ref{psu:periods} with $u=v$ (its polynomial involves only the variables of $V_0$,
Lemma~\ref{psu:vector}), $x'=\psug$ (which satisfies the condition of that proposition, by the
paragraph after \eqref{psu:eq-Xi} for $i=1$) and $g=x_c$. Hence $\langle\Xi,y\rangle=0$.
\end{proof}

\begin{proposition}[The pairing]\label{psu:pairing}
With $A_\psi=N_1^{-1}\sum_{c\in\psuC}\psi(x_c)s(x_c)$,
\[
\langle\Xi,\varphi_{p,\vartheta}\rangle=N_1^{-2}\,c_{\mathrm{vac}}^{\,12}(\pi_Ms_0)^{\sum_\tau l_\tau}\,a\,
I_{q_2}(\vartheta)\,I_{q_3}(1)\,I_{q_1}(1)\,A_\psi .
\]
\end{proposition}
\begin{proof}
By \eqref{po:eq-delta}, \eqref{po:eq-lift} with $y=\plone$, the definition \eqref{pex:eq-phi} of
$\varphi_{p,\vartheta}$ and \eqref{psu:eq-int} (the integrand is right invariant under $K^1_M$, proof of
Lemma~\ref{pex:points}), for every $x$
\[
\langle\podel[x;v],\varphi_{p,\vartheta}\rangle=\langle v,\varphi_{p,\vartheta}(x)\rangle=\kappa_p\sum_{c\in
\psuC}k_p(x_c,x;v),\qquad\kappa_p=c_{\mathrm{vac}}^{2z}(\pi_Ms_0)^{\sum_\tau l_\tau}.
\]
The coefficients $\psuT_i(h)/e_h$ enter conjugated (the pairing is conjugate-linear in the first
variable). Insert Lemma~\ref{psu:factor} and sum over $h_2,h_3,h_1$:
\[
\langle\Xi,\varphi_{p,\vartheta}\rangle=\kappa_p\,a\sum_c\psi(x_c)s(x_c)\prod_{i\in\{2,3,1\}}\Bigl(
\psuN_{z_i}n_{q_i}^2(\auNm Q_i)^{1-2z_i}\,G_{q_i,\vartheta_i}(x_c)\Bigr),
\]
with $G$ as in Lemma~\ref{psu:flat}(c) ($\psuN$, $n_q$, $\auNm Q$ are real). By Lemma~\ref{psu:flat}(c)
the $i$-th factor is $I_{q_i}(\vartheta_i)/(N_1c_{\mathrm{vac}}^{2z_i})$, independent of $c$. Since
$z_2+z_3+z_1=z-6$, $\kappa_pc_{\mathrm{vac}}^{-2(z-6)}=c_{\mathrm{vac}}^{12}(\pi_Ms_0)^{\sum l_\tau}$, and
$\sum_c\psi(x_c)s(x_c)=N_1A_\psi$.
\end{proof}

\subsubsection{Proof of Proposition~\ref{po:small}}\label{psu:ss-small}

\begin{theorem}[Proposition~\ref{po:small}, with the constant corrected]\label{psu:small}
Let $p$ be a main pair for $j$ and $q_2=(\upsilon^{-2}\chi_0,2)$. There is $c_0\in\C^\times$ with
\[
|\sigma c_0|\le\exp\bigl(fmb/k+3x^{35}\bigr)\qquad(\sigma\in\auAut(\C))
\]
such that for every prime $\ell$, every $\jmath\colon\C\to\C_\ell$ and every twist $\vartheta$ with
$\ell\nmid\auNm Q_\vartheta$ there is a test form $\Xi$ at $\ell$ for $(p,\vartheta)$ with
\[
\bigl|c_0\,\rho_{q_2}(\vartheta)\bigr|_\jmath\cdot\ell^{-F(\ell)}\le\bigl|\langle\Xi,\varphi_{p,\vartheta}
\rangle\bigr|_\jmath,\qquad F(\ell)=\sum_{\sigma\in\auPl}s_\ell(\kappa(\sigma),\sigma).
\]
\end{theorem}
\begin{proof}
\emph{The constant.} By \eqref{psu:eq-range}, Theorem~\ref{psm:attained} (unconditional by
Corollary~\ref{pst:attained}) applies to $l=(l_\tau)$ and gives $t\in\C^\times$ with $|\sigma t|\le
\exp(fmb/k)$. Let $L$ be as in Lemma~\ref{psu:mean} and put $c_0=t\,I_{q_2}(1)\,L$. It is non-zero
(Lemma~\ref{psu:I}(a)), it does not depend on $\ell$, $\jmath$, $\vartheta$, and by
Lemmas~\ref{psu:I}(c), \ref{psu:mean}, $|\sigma c_0|\le\exp(fmb/k+x^{35}+2x^{35})$.

\emph{The test form.} Fix $\ell$, $\jmath$ and $\vartheta$ with $\ell\nmid\auNm Q_\vartheta$. By
Theorem~\ref{psm:attained} there are $\gamma\in\N^{\auPl\times12}$ with $\sum_b\gamma(\tau,b)=l_\tau$ and
a class $y_0$ modulo $\pi_M$ with
\begin{equation}\label{psu:eq-att}
\ell^{-F(\ell)}\le\Bigl|\frac{c_{\mathrm{vac}}^{12}}{t}\prod_\tau l_\tau!\cdot[\gamma]\,\Theta_{\phi_{y_0}}
(\pi_Ms_0\,\cdot)\Bigr|_\jmath .
\end{equation}
Take $\alpha=\gamma$ and $x_0=y_0$ in \S\ref{psu:ss-seesaw}; then $\Psi_0=\phi_{y_0}$. By
\eqref{ps:eq-bargmann} with $Q=z^\alpha$ and $c=1$, $s(1)=\theta(z^\alpha\otimes\phi_{y_0})=\partial^\alpha
\Theta_{\phi_{y_0}}(0)=\prod_{\tau,b}\alpha(\tau,b)!\cdot[\alpha]\Theta_{\phi_{y_0}}$, and $[\alpha]\Theta_{\phi_{y_0}}
(\pi_Ms_0\,\cdot)=(\pi_Ms_0)^{\sum l_\tau}[\alpha]\Theta_{\phi_{y_0}}$. So \eqref{psu:eq-att} reads
\begin{equation}\label{psu:eq-s1}
|t|_\jmath\,\ell^{-F(\ell)}\le\bigl|c_{\mathrm{vac}}^{12}(\pi_Ms_0)^{\sum l_\tau}\,a\,s(1)\bigr|_\jmath .
\end{equation}
By Lemma~\ref{psu:orth} (applied to $s$ on $\psuC$, $x_{c_0}=1$) there is $\psi\in\widehat\psuC$ with
$|A_\psi|_\jmath\ge|s(1)|_\jmath$. By Lemma~\ref{psu:pigeon} for the group $\widehat\psuC$, the functions
$a(\chi)=I_{q(\chi,2)}(1)$, $b(\chi)=I_{q(\chi,z-10)}(1)$, $V=L$ and the absolute value $|\cdot|_\jmath$,
there are $\psi_3,\psi_1$ with $\psi_3\psi_1=\psi$ and $|L|_\jmath\le|I_{q_3}(1)I_{q_1}(1)|_\jmath$,
$q_3=q(\psi_3,2)$, $q_1=q(\psi_1,z-10)$. Let $\Xi$ be the form \eqref{psu:eq-Xi} for these data. It is
a test form at $\ell$ for $(p,\vartheta)$ (Definition~\ref{po:R3}) by Lemmas~\ref{psu:integral}
and~\ref{psu:orthfalse}.

\emph{The inequality.} By Proposition~\ref{psu:pairing}, $|N_1^{-2}|_\jmath\ge1$, \eqref{psu:eq-s1}, the
choice of $\psi$, $\psi_3$, $\psi_1$, and $I_{q_2}(\vartheta)=I_{q_2}(1)\rho_{q_2}(\vartheta)$
(Lemma~\ref{psu:I}(a)),
\[
|\langle\Xi,\varphi_{p,\vartheta}\rangle|_\jmath\ge|c_{\mathrm{vac}}^{12}(\pi_Ms_0)^{\sum l_\tau}a|_\jmath\,
|I_{q_2}(\vartheta)|_\jmath\,|L|_\jmath\,|s(1)|_\jmath\ge|t\,I_{q_2}(1)\,L\,\rho_{q_2}(\vartheta)|_\jmath\,
\ell^{-F(\ell)} .\qedhere
\]
\end{proof}

\emph{Remark.} The deduction of Proposition~\ref{po:mainheight} needs $fmb/k+3x^{35}\le\aubud$, which is (N6) of Lemma~\ref{psu:sizes}. The only property of $q_2$ used is that it is the flat pair $q(\chi_0,2)$ on a block
of rank $4$; it is ramified ($\chi_0$ has non-trivial tame part at every chosen place), which is used
afterwards through (R2).

\subsubsection{The untwisted norm at all conjugates: Proposition~\ref{po:arch}}\label{psu:ss-arch}

\begin{lemma}[The norm as an $L$-value, with the archimedean factor]\label{psu:normL}
Let $2\le z\le t$. There is $B(z)>0$ such that for every critical pair $p=(\nu,z)$, with
$l_\tau=c(\tau)-z$,
\[
\|\varphi_{p,1}\|^2=B(z)\,\Bigl(\frac M\Omega\Bigr)^{\sum_\tau l_\tau}\prod_\tau l_\tau!\cdot d_p\cdot
L^{\mathsf M}(z,\nu^{-1}),\qquad d_p=\prod_\tau\binom{l_\tau+2z-1}{2z-1}\ge1,
\]
and $d_p=1$ if $p$ is flat.
\end{lemma}
\begin{proof}
Proposition~\ref{pex:rallis} with $\vartheta=1$ gives $\|\varphi_{p,1}\|^2=A_pL^{\mathsf M}(z,\nu^{-1})$ with
$A_p=|\kappa_p|^2\prod_{w\in\mathsf M}(q_w-1)\cdot\dim W_\lambda\cdot\|v^+\|^{-2}\cdot c'\langle P,P\rangle
\operatorname{vol}([\auU(V^{(z)})])\operatorname{vol}([\auU(1)])\zeta_{E^+}(2z)^{-1}>0$. Here
$|\kappa_p|^2=|c_{\mathrm{vac}}|^{4z}|\pi_M|^{2\sum l}|s_0|^{2\sum l}=|c_{\mathrm{vac}}|^{4z}(M/\Omega)^{\sum l}$
($|\pi_M|^2=M$, $|s_0|^2=1/\Omega$, \S\ref{ps:ss-weil}); $W_\lambda=\bigotimes_\tau\operatorname{Sym}^{l_\tau}
\C^{2z}$ has dimension $d_p$; the highest weight vector is $v^+=\bigotimes_\tau e_1^{\otimes l_\tau}$, of norm $1$,
and $P=\pop_{v^+}=\prod_\tau z_{\tau,1}^{l_\tau}$ (\eqref{po:eq-poly}), whose Fock norm is $\prod_\tau l_\tau!$
in the normalisation of \eqref{ps:eq-bargmann}, in which $\langle z^\beta,z^{\beta'}\rangle=\delta_{\beta
\beta'}\beta!$. The remaining factors depend only on $z$ (and on the data); their product with
$|c_{\mathrm{vac}}|^{4z}$ is $B(z)$, positive because $A_p>0$ for the flat pair $q(1,z)$.
\end{proof}

\begin{lemma}[Euler products at $s\ge2$]\label{psu:euler}
Let $\chi$ be a unitary Hecke character of $E$, $S$ a finite set of places, and $s\ge2$ real. Then
$|\log L^S(s,\chi)|\le2f\log\zeta(2)\le f$ (principal branch of the logarithm of the Euler product).
\end{lemma}
\begin{proof}
$\log L^S(s,\chi)=\sum_{w\notin S,\,\chi\text{ unr.\ at }w}\sum_{k\ge1}\chi(\varpi_w)^kq_w^{-ks}/k$, so
$|\log L^S(s,\chi)|\le\sum_w\sum_kq_w^{-2k}/k=\log\zeta_E(2)$. For each rational prime $p$ there are at
most $[E:\Q]=2f$ places $w\mid p$, each with $q_w\ge p$, so $\zeta_E(2)\le\zeta(2)^{2f}$; and $\log\zeta(2)
=\log(\pi^2/6)<1/2$.
\end{proof}

The proof of Proposition~\ref{po:arch} needs, besides Lemmas~\ref{psu:normL}, \ref{psu:euler}
and~\ref{psu:I}, the following law of Galois conjugation for the norms of the lifts of characters,
proved as Theorem~\ref{pcj:CL} in \S\ref{pcj:sec}.

\emph{(CL) Let $\sigma\in\auAut(\C/K)$ and $2\le z\le t$. There is $a_\sigma(z)>0$ such that for every
critical pair $p=(\nu,z)$ there is a critical pair $p^\sigma=(\nu^\sigma,z)$, whose weight $-c^\sigma$
satisfies $c^\sigma(\sigma\circ\tau)=c(\tau)$ for all $\tau\in\auPl$, with
$|\sigma(\|\varphi_{p,1}\|^2)|=a_\sigma(z)\,\|\varphi_{p^\sigma,1}\|^2$.}

\begin{theorem}[Proposition~\ref{po:arch}]\label{psu:arch}
For a main pair $p=(\nu,z)$ for $j$ the constant family $\|\varphi_{p,1}\|^2$ has height
at most $\aubud-\sum_{\tau\in\auPl}\ausinf(\kappa(\tau))$.
\end{theorem}
\begin{proof}
Let $q=q(1,z)=(\upsilon^{-z},z)$, a flat critical pair of the same rank (Lemma~\ref{psu:flat}(a);
$2\le z\le r\le t$). Take the multiplier $H=I_q(1)/\|\varphi_{p,1}\|^2\in\C^\times$, $\alpha=\aubud-
\sum_\tau\ausinf(\kappa(\tau))$ and $h_\ell=0$ in Definition~\ref{def:au-height}.

\emph{The primes.} $H\|\varphi_{p,1}\|^2=I_q(1)$ is integral at every $\iota$ (Lemma~\ref{psu:I}(b)), so
$|H\|\varphi_{p,1}\|^2|_\iota\le1=\ell^{0}$.

\emph{The conjugates.} Both $I_q(1)$ and $\|\varphi_{p,1}\|^2$ are $\ne0$, and $|(c\circ\sigma)(x)|=|\sigma(x)|$;
since one of $\sigma$, $c\circ\sigma$ fixes $K$, we may assume $\sigma\in\auAut(\C/K)$. By (CL) (Theorem~\ref{pcj:CL}) for $p$ and
for $q$ (note $q^\sigma$ is again flat, and $p^\sigma$ has degrees $l^\sigma_{\sigma\tau}=l_\tau$), and
Lemma~\ref{psu:normL} for $p^\sigma$ and $q^\sigma$,
\begin{align*}
\log|\sigma H|&=\log(\psuN_zn_q^2)+\log|\sigma\|\varphi_{q,1}\|^2|-\log|\sigma\|\varphi_{p,1}\|^2|\\
&=\log(\psuN_zn_q^2)+\log L^{\mathsf M}(z,(\nu_q^\sigma)^{-1})-\sum_\tau\ausinf(l_\tau)-\log d_{p^\sigma}-\log
L^{\mathsf M}(z,(\nu^\sigma)^{-1}),
\end{align*}
because $\log(a_\sigma(z)B(z))$ cancels, $\sum_\tau\bigl(\log l_\tau!+l_\tau\log(M/\Omega)\bigr)=\sum_\tau
\ausinf(l_\tau)$ is invariant under the permutation $\tau\mapsto\sigma\circ\tau$, and the flat pair has
$l=0$, $d=1$. By Lemma~\ref{psu:euler} the two $L$-terms contribute at most $2f$, and $\log d_{p^\sigma}
\ge0$. Since $\ausinf$ is increasing ($M/\Omega\ge1$) and $l_\tau\ge\kappa(\tau)$ by \eqref{psu:eq-range},
\[
\log|\sigma H|\le\log\psuN_z+8fz\log M+2f-\sum_\tau\ausinf(\kappa(\tau))\le\aubud-\sum_\tau\ausinf(\kappa
(\tau)),
\]
because $\log\psuN_z+8fz\log M\le x^{35}$ by (N5) and $x^{35}+2f\le\aubud$ by (N6) of
Lemma~\ref{psu:sizes}. Hence the family has height at most $\alpha$.
\end{proof}

\emph{Remarks.} (i) The multiplier is $I_q(1)/\|\varphi_{p,1}\|^2$ with the integer $\psuN_zn_q^2$ in $I_q$.
(ii) The proof uses Theorem~\ref{pcj:CL} (through Theorem~\ref{pf:recip}), Lemma~\ref{psu:normL} (from Proposition~\ref{pex:rallis}), Lemma~\ref{psu:euler} and Lemma~\ref{psu:I}(b).



%% ===== proofs/C-conj.tex =====
% proofs/C-conj.tex
% The law of Galois conjugation (CL) of the norms of the lifts of characters, used in the proof of
% Theorem psu:arch (= Proposition po:arch). To be \input after C-onelift, C-small, C-small2,
% C-exchanges3, C-periods, C-small3. Labels of this file start with pcj:.
% Only amsmath/amssymb/amsthm are used.
\providecommand{\auO}{\mathcal O}
\providecommand{\auU}{\mathrm U}
\providecommand{\auHom}{\operatorname{Hom}}
\providecommand{\auAut}{\operatorname{Aut}}
\providecommand{\auNm}{\operatorname{Nm}}
\providecommand{\auA}{\mathbb A}
\providecommand{\auPl}{\Phi}
\providecommand{\aubar}[1]{\overline{#1}}
\providecommand{\plone}{\mathbf 1}
\providecommand{\psrho}{\varrho}
\providecommand{\psvac}{\mathbf v}
\providecommand{\psHS}{\mathcal H}
\providecommand{\poA}{\mathcal A}
\providecommand{\pop}{\mathsf p}
% macros of this file (prefix \pcj)
\providecommand{\pcjT}{\operatorname{T}}
\providecommand{\pcjJ}{\mathrm J}
\providecommand{\pcjB}{\mathsf B}
\providecommand{\pcjS}{\mathsf S}
\providecommand{\pcjG}{\mathcal G}
\providecommand{\pcjtr}{\operatorname{Tr}}
\providecommand{\pcjsh}{\psrho^{\sharp}}
\providecommand{\pcjhat}[1]{\widehat{#1}}

\subsection{The Galois conjugation of the lifts of characters: statement (CL)}\label{pcj:sec}

This subsection proves the statement (CL) isolated in \S\ref{psu:ss-arch}, which is the only input
of Theorem~\ref{psu:arch} (that is, of Proposition~\ref{po:arch}) not proved there. We use the
notation of \S\ref{ps:ss-weil}, \S\ref{pk:ss-conv}, \S\ref{pk:ss-conj}, \S\ref{po:ss-set},
\S\ref{po:ss-hecke} (\emph{the lift of a character}) and \S\ref{pex:ss-set}.

\begin{theorem}[Statement (CL)]\label{pcj:CL}
Let $\sigma\in\auAut(\C/K)$ and $2\le z\le t$. There is $a_\sigma(z)>0$ such that for every critical
pair $p=(\nu,z)$ there is a critical pair $p^\sigma=(\nu^{\sigma},z)$, whose weight $-c^\sigma$ satisfies
$c^\sigma(\sigma\circ\tau)=c(\tau)$ for all $\tau\in\auPl$, with
\[
\bigl|\sigma\bigl(\|\varphi_{p,1}\|^2\bigr)\bigr|=a_\sigma(z)\,\|\varphi_{p^\sigma,1}\|^2 .
\]
\end{theorem}

\emph{Why Lemma~\ref{ps:conj} does not suffice, and what is done here.} Proposition~\ref{pk:conj}
gives a $\sigma$-semilinear map $\Lambda_\sigma$ which conjugates jets but commutes with the
Heisenberg translations only up to roots of unity $\zeta_x$; then $\Lambda_\sigma\omega(g)$ and
$\omega(g)\Lambda_\sigma$ differ by an undetermined translation $\psrho_f(y_g)$, and nothing can be
said about the image of the datum of a lift. We show (Proposition~\ref{pcj:Lambda}) that all these
roots of unity are $1$: the jet-conjugating map is unique, and it carries $\psrho_f(x)$ to
$\psrho_f(\pcjB x)$ \emph{exactly}, for the idele $\pcjB=\beta_\sigma$ of Theorem~\ref{pf:recip}; at
the same time $\sigma$ acts on the values of $\psi_f$ as multiplication by $\pcjB\aubar\pcjB$
(Lemma~\ref{pcj:char}). Writing $\pcjB=ua$ with $u$ of norm one and $a$ a unit equal to $1$ at the
chosen places (Lemma~\ref{pcj:split}), the map $\omega(u)^{-1}\Lambda_\sigma$ fixes the line of the
vacuum, conjugates the Weil operators up to explicit scalars $\sigma(\delta_k)/\delta_k$ with
$|\delta_k|^2\in\Q$ (Lemma~\ref{pcj:vac}), and multiplies the block datum by a constant
(Lemma~\ref{pcj:datum}). Summing over the classes of $\auU(1)$ gives the conjugate of each
coordinate of the lift (Proposition~\ref{pcj:coord}); the scalars cancel between $\sigma$ and
$c\sigma c$ in the norm.

\emph{Dependency order.} Lemma~\ref{pcj:loc} (local constancy) $\to$ Lemma~\ref{pcj:transl} (exact
translates of the theta products) $\to$ Lemma~\ref{pcj:jets} (exact conjugation of their jets)
$\to$ Lemma~\ref{pcj:char} ($\sigma$ on $\psi_f$) $\to$ Proposition~\ref{pcj:Lambda} (the map
$\Lambda$) $\to$ Lemma~\ref{pcj:split} $\to$ Lemma~\ref{pcj:vac} (vacuum and Weil operators) $\to$
Lemma~\ref{pcj:datum} (the block datum) $\to$ Lemmas~\ref{pcj:lift}, \ref{pcj:nusigma}
(the lift as a class sum; conjugate characters) $\to$ Proposition~\ref{pcj:coord} $\to$
Theorem~\ref{pcj:CL}.

\subsubsection{Setting}\label{pcj:ss-set}

Fix $z\ge2$ and let $W=\auHom(E,V^{(z)})=E^{\mathsf C}$, $\mathsf C=[2z]\times[1]$, with the form
$\sum_i\aubar x_iy_i$; $L=\alpha\pcjhat\auO_E^{\,\mathsf C}$, $\Lambda_L=W(E)\cap L=\mathfrak
a_0^{\mathsf C}$ (\S\ref{pk:ss-conj}). We write $\psrho_f$, $\omega=\omega_f$ for the finite
Heisenberg and Weil operators on $\psHS_f=\psHS_f(W)$, which are unitary for its inner product,
$\beta=\beta_f$, $[x,y]=\beta(x,y)\beta(y,x)^{-1}$, and $\psvac=\psvac_L$; so
$\psrho_f(x)\psrho_f(y)=\beta(x,y)\psrho_f(x+y)$ with $\beta(x,y)=\psi_f(\tfrac12\jmath\langle
y,x\rangle)$, and $\psrho_f(l)\psvac=\epsilon_L(l)\psvac$ for $l\in L$, $\epsilon_L(l)=\psi_f(-\tfrac12
\omega_0\jmath\langle l,l\rangle)$, $\omega_0=(1+\sqrt{-d})/2$ (\S\ref{ps:ss-weil}). We use three
elementary facts: (a) $\psi_f(\aubar y)=\psi_f(y)$, and $\aubar\jmath=-\jmath$; (b) $\psi_f(y)=1$
for $y\in E$ with $\pcjtr_{E/\Q}y=0$, because $\psi_f(y)=\psi_\infty(y)^{-1}=e^{2\boldsymbol\pi i\,
\pcjtr_{E/\Q}y}$; (c) $\psi_f(\jmath\alpha\aubar\alpha t)=1$ for $t\in\pcjhat\auO_E$, by the
self-duality of $\pcjhat\auO_E$ for $(z,u)\mapsto\psi(\jmath\alpha\aubar\alpha\aubar zu)$; since
$\langle l,l'\rangle\in\alpha\aubar\alpha\pcjhat\auO_E$ for $l,l'\in L$, (c) gives
$\psi_f(\jmath\,\mathsf t\langle l,l'\rangle)=1$ for $l,l'\in L$ and $\mathsf t\in\pcjhat\auO_E$.

\emph{Taylor series and jets.} For $f\in\psHS_f$ and $c\in\C^\times$ let $\pcjT_c(f)\in\C[[z]]$
be the Taylor series at $0$ of $z\mapsto\Theta_f(cz)$, in the variables $z_{\tau,\mathsf c}$
($\tau\in\auPl$, $\mathsf c\in\mathsf C$), with $\Theta_f$ as in \eqref{pk:eq-theta}. The map
$f\mapsto\pcjT_c(f)$ is linear and injective, because $f\mapsto\Theta_f$ is injective
(Lemma~\ref{pk:dict}) and an entire function is determined by its Taylor series. For a polynomial
$Q$ the jet is $\pcjJ_{Q,c}(f)=Q(\partial)[\Theta_f(c\,\cdot)](0)=\theta(Q(c\,\cdot)\otimes f)$
(\eqref{pk:eq-bargmann}); if $Q$ is homogeneous of degree $D$, $\pcjJ_{Q,c}(f)=c^D\pcjJ_{Q,1}(f)$.
For $\sigma\in\auAut(\C/K)$ let $\sigma_*$ be the continuous ring endomorphism of $\C[[z]]$ which is
$\sigma$ on the constants and maps $z_{\tau,\mathsf c}$ to $z_{\sigma\tau,\mathsf c}$; here
$\sigma\circ\tau\in\auPl$ because $\sigma$ fixes $K$. Writing $\pcjT_c(f)=\sum_\gamma t_\gamma z^\gamma$
and $Q=\sum_\gamma q_\gamma z^\gamma$, $\pcjJ_{Q,c}(f)=\sum_\gamma\gamma!\,q_\gamma t_\gamma$; as the
renaming of the variables does not change $\gamma!$,
\begin{equation}\label{pcj:eq-jet}
\pcjT_{c'}(f')=\sigma_*\pcjT_c(f)\quad\Longrightarrow\quad\pcjJ_{\sigma_*Q,c'}(f')=\sigma\bigl(\pcjJ_{Q,c}
(f)\bigr)\ \text{ for every polynomial }Q .
\end{equation}

\emph{The data attached to $\sigma$.} Fix $\sigma\in\auAut(\C/K)$ and $\Omega_\sigma$, $\mathfrak c$,
$\mathfrak a=\mathfrak c\aubar{\mathfrak c}^{-1}$, $\beta_\sigma$ as in Theorem~\ref{pf:recip}; let
$\pcjB$ be $\beta_\sigma$ regarded in $\auA_{E,f}^\times$, so that $\pcjB\pcjhat\auO_E=\mathfrak a
\pcjhat\auO_E$, $\pcjB L=\mathfrak aL$, and $\pcjB\aubar\pcjB$ is a unit (Theorem~\ref{pf:recip}(a)).
Let $F_0$ be the vector of \eqref{pk:eq-product} ($\Theta_{F_0}=\Pi$), $F_0^\sigma$ the vector attached
in the proof of Proposition~\ref{pk:conj} to the lattice $\Lambda^\sigma=\Omega_\sigma\mathfrak a$
and the number $\Omega_\sigma$ ($\Theta_{F_0^\sigma}=\Pi^\sigma$, the product \eqref{pk:eq-product}
with $\theta^\sigma$ and $\Omega_\sigma$ in place of $\theta$ and $\Omega_0$), and
$t_\sigma=\sigma(\varpi_{\mathrm{vac}}/\mathsf u)\,c_\psi/\Omega_\sigma$, so that $|t_\sigma|^2=
|s_0|^2=\Omega^{-1}$ (Proposition~\ref{pk:conj}, \eqref{pk:eq-s0}).

\begin{definition}\label{pcj:def-image}
Let $\gamma\in W(E)$ and let $\mathsf m\ge1$ be an integer with $\mathsf m\gamma\in\Lambda_L$. An
\emph{$\mathsf m$-image} of $\gamma$ is a $\gamma'\in W(E)$ with $\gamma'-\pcjB\gamma\in2q_1^2\mathsf m\,
\pcjB L$.
\end{definition}

Since $W(E)$ is dense in $W(\auA_{E,f})$ and $\pcjB L$ is open, $\mathsf m$-images exist for every
admissible $\mathsf m$. If $\gamma'$ and $\delta'$ are $\mathsf m$-images of $\gamma$ and $\delta$,
then $\gamma'+\delta'$ is an $\mathsf m$-image of $\gamma+\delta$; $0$ is an $\mathsf m$-image of $0$;
and if $\mathsf m\mid\mathsf m'$, an $\mathsf m'$-image is an $\mathsf m$-image.

\begin{lemma}[Local constancy]\label{pcj:loc}
For $w\in\psHS_f$ and $x_0\in W(\auA_{E,f})$ there is an open subgroup $U$ of $W(\auA_{E,f})$ with
$\psrho_f(x_0+y)w=\psrho_f(x_0)w$ for all $y\in U$.
\end{lemma}
\begin{proof}
By (W4), $w$ is a finite sum $\sum_jc_j\psrho_f(x_j)\psvac$. For $y\in L$,
$\psrho_f(y)\psrho_f(x_j)\psvac=[y,x_j]\epsilon_L(y)\psrho_f(x_j)\psvac$. The characters
$y\mapsto[y,x_j]$, $y\mapsto\beta(x_0,y)$ of $W(\auA_{E,f})$ are continuous, and $\epsilon_L(y)=1$
for $y\in NL$, $N$ even and large. So there is an open $U\subseteq L$ on which all of them are $1$;
then $\psrho_f(y)w=w$ and $\psrho_f(x_0+y)w=\beta(x_0,y)^{-1}\psrho_f(x_0)\psrho_f(y)w=\psrho_f(x_0)w$.
\end{proof}

\subsubsection{Exact translates and exact conjugation of jets}\label{pcj:ss-exact}

\begin{lemma}[Exact translates]\label{pcj:transl}
For $\gamma\in W(E)$, with $u_k=\zeta_{q_1}^k/q_1$, $a_{k,i}(z)=\frac{\Omega_0}{c_\psi}u_k^\dagger(z_i)$,
$a^\sigma_{k,i}(z)=\frac{\Omega_\sigma}{c_\psi}u_k^\dagger(z_i)$ and $\mathsf x_{k,i}(\gamma)=
-\pcjtr_{E/K}(\aubar{u_k}\gamma_i)\in K$,
\[
\Theta_{\psrho_f(\gamma)F_0}(z)=\prod_{i,k}\theta_{\mathsf x_{k,i}(\gamma)\Omega_0}\bigl(a_{k,i}(z)\bigr)^{q_1},
\qquad
\Theta_{\psrho_f(\gamma)F_0^\sigma}(z)=\prod_{i,k}\theta^\sigma_{\mathsf x_{k,i}(\gamma)\Omega_\sigma}\bigl(
a^\sigma_{k,i}(z)\bigr)^{q_1},
\]
with no root of unity, the twisted translates being taken for $\Lambda$ and for $\Lambda^\sigma$.
\end{lemma}
\begin{proof}
\emph{The root of unity of Lemma~\ref{pk:dict} is $1$ for rational $\gamma$.} In the proof of that
lemma, $\theta=\theta\circ\psrho(-\gamma)$ and $\psrho(-\gamma)=\psrho_\infty(-\gamma^\natural)\otimes
\psrho_f(-\gamma)$; so
$\theta(\psrho_\infty(z)(1\otimes\psrho_f(\gamma)f))=\beta_\infty(-\gamma^\natural,z)\,
\beta_f(-\gamma,\gamma)\,\theta(\psrho_\infty(z-\gamma^\natural)(1\otimes f))$, and the root of unity
there is $\beta_f(-\gamma,\gamma)=\psi_f(-\frac12\jmath\langle\gamma,\gamma\rangle)$. Now
$\langle\gamma,\gamma\rangle\in E^+$, so $\pcjtr_{E/K}\langle\gamma,\gamma\rangle\in K\cap\R=\Q$ and
$\pcjtr_{E/\Q}(\jmath\langle\gamma,\gamma\rangle)=\pcjtr_{K/\Q}(\jmath)\pcjtr_{E/K}\langle\gamma,
\gamma\rangle=0$; by fact (b) the root of unity is $1$. Hence, for every $f$,
$\Theta_{\psrho_f(\gamma)f}(z)=e^{\langle z,\gamma^\natural\rangle-|\gamma^\natural|^2/2}\Theta_f(z-
\gamma^\natural)$.

\emph{The product.} Put $\mathsf y=\mathsf x_{k,i}(\gamma)\Omega_0=-a_{k,i}(\gamma^\natural)$. By
\eqref{pa:eq-res} and $\boldsymbol\pi|\Omega_0|^2/(Ac_\psi^2)=1$ (as in \S\ref{pk:ss-conj}),
$\langle z,\gamma^\natural\rangle=\sum_{i,k}q_1H(a_{k,i}(\gamma^\natural),a_{k,i}(z))=-\sum_{i,k}q_1
\frac{\boldsymbol\pi}A\aubar{\mathsf y}\,a_{k,i}(z)$ and $|\gamma^\natural|^2=\sum_{i,k}q_1\frac{
\boldsymbol\pi}A|\mathsf y|^2$ (the $\mathsf y$ depending on $(i,k)$), while $a_{k,i}(z-\gamma^\natural)=
a_{k,i}(z)+\mathsf y$. So $e^{\langle z,\gamma^\natural\rangle-|\gamma^\natural|^2/2}\Pi(z-\gamma^\natural)
=\prod_{i,k}\bigl[e^{-\frac{\boldsymbol\pi}A(a_{k,i}(z)\aubar{\mathsf y}+\frac12|\mathsf y|^2)}\theta(a_{k,i}
(z)+\mathsf y)\bigr]^{q_1}=\prod_{i,k}\theta_{\mathsf y}(a_{k,i}(z))^{q_1}$ by the definition of the
twisted translate (\S\ref{pk:ss-curve}). For $F_0^\sigma$ the same computation applies: the area of
$\Lambda^\sigma=\Omega_\sigma\mathfrak a$ is $A^\sigma=|\Omega_\sigma|^2\sqrt d/2$ (as $\auNm\mathfrak a=1$),
so $\boldsymbol\pi|\Omega_\sigma|^2/(A^\sigma c_\psi^2)=2\boldsymbol\pi/(\sqrt dc_\psi^2)=1$.
\end{proof}

\begin{lemma}[Exact conjugation of jets of translates]\label{pcj:jets}
Let $\gamma\in W(E)$, $\mathsf m\ge1$ with $\mathsf m\gamma\in\Lambda_L$, and $\gamma'$ an
$\mathsf m$-image of $\gamma$. Then $\pcjT_{t_\sigma}(\psrho_f(\gamma')F_0^\sigma)=\sigma_*\pcjT_{s_0}
(\psrho_f(\gamma)F_0)$.
\end{lemma}
\begin{proof}
By Lemma~\ref{pcj:transl} and \eqref{pk:eq-s0}, $\pcjT_{s_0}(\psrho_f(\gamma)F_0)=\prod_{i,k}G_{k,i}
(\ell_{k,i})^{q_1}$, where $G_{k,i}$ is the Taylor series of $\theta_{\mathsf x_{k,i}\Omega_0}$ and
$\ell_{k,i}=a_{k,i}(s_0z)=(\varpi_{\mathrm{vac}}/\mathsf u)\,u_k^\dagger(z_i)$; likewise
$\pcjT_{t_\sigma}(\psrho_f(\gamma')F_0^\sigma)=\prod_{i,k}G'_{k,i}(\ell'_{k,i})^{q_1}$ with $G'_{k,i}$ the
Taylor series of $\theta^\sigma_{\mathsf x'_{k,i}\Omega_\sigma}$, $\mathsf x'_{k,i}=\mathsf x_{k,i}
(\gamma')$, and $\ell'_{k,i}=a^\sigma_{k,i}(t_\sigma z)=\sigma(\varpi_{\mathrm{vac}}/\mathsf u)\,
u_k^\dagger(z_i)$. As $\sigma_*$ is a continuous ring homomorphism, $\sigma_*(G(\ell))=(\sigma G)
(\sigma_*\ell)$ for a series $G$ in one variable and a linear form $\ell$. Since $\sigma$ fixes $K$
and $\tau(E)$ is a CM field, $\sigma(\aubar{\tau(u_k)})=\aubar{(\sigma\tau)(u_k)}$; so $\sigma_*\ell_{k,i}
=\sigma(\varpi_{\mathrm{vac}}/\mathsf u)\sum_\tau\aubar{(\sigma\tau)(u_k)}\,z_{\sigma\tau,i}=\ell'_{k,i}$.
It remains to see $\sigma G_{k,i}=G'_{k,i}$, which is Theorem~\ref{pf:recip}(b) for $\mathsf x=
\mathsf x_{k,i}$, $n=q_1\mathsf m$, $\mathsf x'=\mathsf x'_{k,i}$. Indeed $n\mathsf x=-\pcjtr_{E/K}
(\zeta_{q_1}^{-k}\mathsf m\gamma_i)\in\pcjtr_{E/K}(\mathfrak a_0)\subseteq\auO_K$, as $\mathfrak a_0
\subseteq\mathfrak d_{E/K}^{-1}$ is an $\auO_E$-module; and, $\pcjtr_{E/K}$ being
$\auA_{K,f}$-linear and $\beta_\sigma\in\auA_{K,f}$, $\mathsf x'-\beta_\sigma\mathsf x=-\pcjtr_{E/K}
(\aubar{u_k}(\gamma'_i-\pcjB\gamma_i))$ with $\aubar{u_k}(\gamma'_i-\pcjB\gamma_i)\in2q_1\mathsf m\,
\mathfrak a\,\alpha\pcjhat\auO_E=2q_1\mathsf m\,\mathfrak a\,\pcjhat{\mathfrak a_0}$, whose trace lies in
$2q_1\mathsf m\,\pcjhat{\mathfrak a}=2n\pcjhat{\mathfrak a}$.
\end{proof}

\begin{lemma}[Closeness]\label{pcj:close}
Let $\gamma,\delta\in W(E)$, $\mathsf m\gamma,\mathsf m\delta\in\Lambda_L$, and $\gamma',\delta'$ their
$\mathsf m$-images. Then $\beta(\delta',\gamma')=\psi_f(\frac12\jmath\langle\pcjB\gamma,\pcjB\delta
\rangle)=\psi_f(\frac12\jmath\,\pcjB\aubar\pcjB\langle\gamma,\delta\rangle)$.
\end{lemma}
\begin{proof}
$\langle\gamma',\delta'\rangle-\langle\pcjB\gamma,\pcjB\delta\rangle=\langle\gamma'-\pcjB\gamma,
\delta'\rangle+\langle\pcjB\gamma,\delta'-\pcjB\delta\rangle$. Here $\delta,\gamma\in\mathsf m^{-1}L$, so
$\pcjB\gamma,\delta'\in\pcjB\mathsf m^{-1}L$, and $\gamma'-\pcjB\gamma,\delta'-\pcjB\delta\in2q_1^2
\mathsf m\pcjB L$. As $\langle\pcjB l,\pcjB l'\rangle=\pcjB\aubar\pcjB\langle l,l'\rangle$ with
$\pcjB\aubar\pcjB$ a unit, both terms lie in $2\langle L,L\rangle\pcjhat\auO_E$, and
$\psi_f(\frac12\jmath\cdot2\mathsf t)=1$ for $\mathsf t\in\langle L,L\rangle\pcjhat\auO_E$ by fact (c).
\end{proof}

\subsubsection{The jet-conjugating map}\label{pcj:ss-Lambda}

Every vector of $\psHS_f$ is a finite combination of vectors $\psrho_f(\gamma)F_0$, $\gamma\in W(E)$
(\S\ref{pk:ss-conj}: the $\psrho_f(x)F_0$ span, and each is a root of unity times
$\psrho_f(\gamma)F_0$ for rational $\gamma$ close to $x$).

\begin{lemma}[Definition of $\Lambda$ on translates]\label{pcj:def}
There is a unique map $\Lambda\colon\psHS_f\to\psHS_f$ with $\pcjT_{t_\sigma}(\Lambda f)=\sigma_*
\pcjT_{s_0}(f)$ for all $f$. It is $\sigma$-semilinear and injective, and $\Lambda(\psrho_f(\gamma)
F_0)=\psrho_f(\gamma')F_0^\sigma$ whenever $\gamma'$ is an $\mathsf m$-image of $\gamma\in W(E)$.
\end{lemma}
\begin{proof}
Uniqueness and injectivity follow from the injectivity of $\pcjT_{t_\sigma}$ and $\pcjT_{s_0}$ and
the injectivity of $\sigma_*$. For $f=\sum_jc_j\psrho_f(\gamma_j)F_0$ put $\Lambda f=\sum_j\sigma(c_j)
\psrho_f(\gamma_j')F_0^\sigma$ with $\gamma'_j$ any $\mathsf m_j$-images. By Lemma~\ref{pcj:jets} and the
$\sigma$-semilinearity of $\sigma_*\circ\pcjT_{s_0}$, $\pcjT_{t_\sigma}(\Lambda f)=\sigma_*\pcjT_{s_0}(f)$; so
$\Lambda f$ does not depend on the choices (injectivity of $\pcjT_{t_\sigma}$), $\Lambda$ is well defined
and $\sigma$-semilinear, and has the stated value on translates.
\end{proof}

\begin{lemma}[$\sigma$ on the additive character]\label{pcj:char}
$\sigma(\psi_f(y))=\psi_f(\pcjB\aubar\pcjB\,y)$ for all $y\in\auA_{E,f}$.
\end{lemma}
\begin{proof}
Let $\lambda\in q_1\Lambda_L$, $\delta\in W(E)$, $\mathsf m\ge1$ with $\mathsf m\lambda,\mathsf m\delta
\in\Lambda_L$, and $\lambda',\delta'$ their $\mathsf m$-images. By Lemma~\ref{pf:semi} (for the frame
$(\zeta_{q_1}^k)$, $\mathsf b=q_1$), $\psrho_f(\lambda)F_0=\epsilon_L(\lambda)F_0$. Apply $\Lambda$, using
Lemma~\ref{pcj:def} with the $\mathsf m$-image $0$ of $0$: $\psrho_f(\lambda')F_0^\sigma=\Lambda(\psrho_f
(\lambda)F_0)=\sigma(\epsilon_L(\lambda))F_0^\sigma$. Now compute $\Lambda v$ for $v=\psrho_f(\delta)
\psrho_f(\lambda)F_0$ in two ways. First $v=\epsilon_L(\lambda)\psrho_f(\delta)F_0$, so $\Lambda v=
\sigma(\epsilon_L(\lambda))\psrho_f(\delta')F_0^\sigma$. Second $v=\beta(\delta,\lambda)\psrho_f(\delta+
\lambda)F_0$ and $\delta'+\lambda'$ is an $\mathsf m$-image of $\delta+\lambda$, so
\[
\Lambda v=\sigma(\beta(\delta,\lambda))\psrho_f(\delta'+\lambda')F_0^\sigma=\sigma(\beta(\delta,\lambda))
\beta(\delta',\lambda')^{-1}\psrho_f(\delta')\psrho_f(\lambda')F_0^\sigma=\sigma(\beta(\delta,\lambda))
\beta(\delta',\lambda')^{-1}\sigma(\epsilon_L(\lambda))\psrho_f(\delta')F_0^\sigma .
\]
Since $F_0^\sigma\ne0$ ($\Pi^\sigma\ne0$) and $\psrho_f(\delta')$ is invertible, $\sigma(\beta(\delta,
\lambda))=\beta(\delta',\lambda')$, which by Lemma~\ref{pcj:close} gives
$\sigma(\psi_f(\frac12\jmath\langle\lambda,\delta\rangle))=\psi_f(\frac12\jmath\pcjB\aubar\pcjB\langle
\lambda,\delta\rangle)$. Fix $\lambda\ne0$; for $y\in E$ the vector $\delta=2y(\jmath\langle\lambda,
\lambda\rangle)^{-1}\lambda$ has $\frac12\jmath\langle\lambda,\delta\rangle=y$. So the assertion holds for
$y\in E$. Both sides are characters of $\auA_{E,f}$ trivial on an open subgroup ($\psi_f$ is), hence
continuous, and $E$ is dense in $\auA_{E,f}$.
\end{proof}

\begin{proposition}[The jet-conjugating map scales translations exactly]\label{pcj:Lambda}
The map $\Lambda$ of Lemma~\ref{pcj:def} is a $\sigma$-semilinear bijection of $\psHS_f$ with
\[
\pcjT_{t_\sigma}(\Lambda f)=\sigma_*\pcjT_{s_0}(f),\qquad\Lambda\,\psrho_f(x)=\psrho_f(\pcjB x)\,
\Lambda\qquad(f\in\psHS_f,\ x\in W(\auA_{E,f})).
\]
\end{proposition}
\begin{proof}
\emph{Rational $x$.} Let $\delta,\gamma\in W(E)$ and $v=\psrho_f(\gamma)F_0$. Let $\mathsf m_0\ge1$
with $\mathsf m_0\delta,\mathsf m_0\gamma\in\Lambda_L$, let $\mathsf m$ be a multiple of $\mathsf m_0$,
and $\delta',\gamma'$ $\mathsf m$-images. As in the proof of Lemma~\ref{pcj:char},
\[
\Lambda\psrho_f(\delta)v=\sigma(\beta(\delta,\gamma))\psrho_f(\delta'+\gamma')F_0^\sigma=\sigma(\beta
(\delta,\gamma))\beta(\delta',\gamma')^{-1}\psrho_f(\delta')\Lambda v=\psrho_f(\delta')\Lambda v,
\]
because $\sigma(\beta(\delta,\gamma))=\psi_f(\frac12\jmath\pcjB\aubar\pcjB\langle\gamma,\delta\rangle)=
\beta(\delta',\gamma')$ by Lemmas~\ref{pcj:char} and~\ref{pcj:close}. Take $\mathsf m=\mathsf m_0\,j!$ and
let $j\to\infty$: then $\delta'\to\pcjB\delta$ in $W(\auA_{E,f})$, the left side does not change, and by
Lemma~\ref{pcj:loc} (at $x_0=\pcjB\delta$, $w=\Lambda v$) $\psrho_f(\delta')\Lambda v=\psrho_f(\pcjB\delta)
\Lambda v$ for $j$ large. By linearity, $\Lambda\psrho_f(\delta)=\psrho_f(\pcjB\delta)\Lambda$ on
$\psHS_f$.

\emph{All $x$.} Let $f\in\psHS_f$ and $\delta_j\in W(E)$, $\delta_j\to x$. By Lemma~\ref{pcj:loc}
(at $x$ and $f$, and at $\pcjB x$ and $\Lambda f$), for $j$ large $\psrho_f(\delta_j)f=\psrho_f(x)f$ and
$\psrho_f(\pcjB\delta_j)\Lambda f=\psrho_f(\pcjB x)\Lambda f$, so $\Lambda\psrho_f(x)f=\Lambda\psrho_f(\delta_j)
f=\psrho_f(\pcjB\delta_j)\Lambda f=\psrho_f(\pcjB x)\Lambda f$.

\emph{Bijectivity.} $\Lambda$ is injective (Lemma~\ref{pcj:def}). Its image is a non-zero subspace
stable under every $\psrho_f(y)=\Lambda\psrho_f(\pcjB^{-1}y)\Lambda^{-1}$ ($\pcjB$ is an idele), hence is
$\psHS_f$ by the irreducibility of the Heisenberg representation.
\end{proof}

\begin{remark}\label{pcj:rem-pk}
Proposition~\ref{pcj:Lambda} sharpens the commutation rules in the proof of
Proposition~\ref{pk:conj}: the map $\Lambda_\sigma$ defined there satisfies the jet identity of Lemma~\ref{ps:conj}(i), hence equals $\Lambda$
by uniqueness, and its roots of unity $\zeta'$ and $\zeta_x$ are all $1$ when the images $\gamma'$ are
chosen as in Definition~\ref{pcj:def-image} (the only source of roots of unity there is
Lemma~\ref{pk:dict}, whose root of unity is $1$ for rational translations by
Lemma~\ref{pcj:transl}).
\end{remark}

\subsubsection{The vacuum, the Weil operators and the block datum}\label{pcj:ss-vac}

\begin{lemma}[Splitting the scale]\label{pcj:split}
Let $\mathsf R$ be a finite set of finite places of $E$ with $\mathsf R\cap\aubar{\mathsf R}=\emptyset$.
There are $u\in\auA_{E,f}^\times$ with $u\aubar u=1$ and $a\in\pcjhat\auO_E^\times$ with $a_w=1$ for
$w\in\mathsf R$, such that $\pcjB=ua$. Then $\pcjB\aubar\pcjB=a\aubar a$.
\end{lemma}
\begin{proof}
It suffices to find a unit $a$, equal to $1$ on $\mathsf R$, with $a\aubar a=\pcjB\aubar\pcjB$, and
to put $u=\pcjB a^{-1}$; we do it at each finite place $v$ of $E^+$. If $v$ does not split, $E_v$ is a
field and $\pcjB_v\aubar{\pcjB_v}=\auNm(\pcjB_v)$ is a unit, so $\pcjB_v$ is a unit; put $a_v=\pcjB_v$
($v$ is not below $\mathsf R$, whose places satisfy $w\ne\aubar w$). If $v=w\aubar w$ splits, write
elements of $E_v=E_w\times E_{\aubar w}$ as pairs, with $\aubar{(x,y)}=(\aubar y,\aubar x)$ for the
conjugation isomorphisms $E_{\aubar w}\leftrightarrow E_w$. Then $\pcjB\aubar\pcjB=(\nu_w,\aubar{\nu_w})$
with $\nu_w=\pcjB_w\aubar{\pcjB_{\aubar w}}$ a unit, and $(a_w,a_{\aubar w})=(1,\aubar{\nu_w})$ if
$w\in\mathsf R$, $(\nu_w,1)$ otherwise (at most one of $w,\aubar w$ is in $\mathsf R$; name them so that
$\aubar w\notin\mathsf R$), satisfies $a\aubar a=(a_w\aubar{a_{\aubar w}},a_{\aubar w}\aubar{a_w})=
(\nu_w,\aubar{\nu_w})$.
\end{proof}

We apply Lemma~\ref{pcj:split} with $\mathsf R=\mathsf M$, the chosen places (all of them lie above
split places of $E^+$ and $\mathsf M\cap\aubar{\mathsf M}=\emptyset$, \S\ref{pex:ss-set}). The norm-one
idele $u$ is a point of $\auU(1)(\auA_{E^+,f})$ and acts on $W$ as a scalar, a central element of
$\auU(W)(\auA_f)$. Put $\Lambda'=\omega(u)^{-1}\Lambda$, and, for $x\in W(\auA_{E,f})$,
\[
\pcjsh(x)=e(x)\psrho_f(x),\qquad e(x)=\psi_f\bigl(\tfrac12\omega_0\jmath\langle x,x\rangle\bigr),
\]
so that $\pcjsh(l)\psvac=\psvac$ for $l\in L$ ($e=\epsilon_L^{-1}$ on $L$). For
$k\in\auU(W)(\auA_f)$ let $I_k=[L:L\cap k^{-1}L]$ and $\pcjS_k=\sum_{q\in L/(L\cap k^{-1}L)}\pcjsh(kq)
\psvac$.

\begin{lemma}[Vacuum and Weil operators]\label{pcj:vac}
\begin{enumerate}
\item[(a)] $\Lambda'$ is a $\sigma$-semilinear bijection with $\Lambda'\psrho_f(x)=\psrho_f(ax)\Lambda'$
and $\Lambda'\pcjsh(x)=\pcjsh(ax)\Lambda'$ for all $x$, and $\sigma(\psi_f(y))=\psi_f(a\aubar ay)$.
\item[(b)] $\Lambda'\psvac=\mathsf c\,\psvac$ for a number $\mathsf c\ne0$.
\item[(c)] $\pcjsh(x)\pcjsh(y)=\psi_f(-\jmath\omega_0\langle x,y\rangle)\,\pcjsh(x+y)$. The terms of
$\pcjS_k$ do not depend on the representatives $q$, $\pcjS_k\ne0$, and $\omega(k)\psvac=\delta_k\pcjS_k$
with $\delta_k\in\C^\times$, $|\delta_k|^2=I_k^{-1}$.
\item[(d)] For $k\in\auU(W)(\auA_f)$ put $\mathsf m_\sigma(k)=\sigma(\delta_k)/\delta_k$. Then
$\Lambda'\omega(k)=\mathsf m_\sigma(k)\,\omega(k)\Lambda'$ and $\Lambda\omega(k)=\mathsf m_\sigma(k)\,
\omega(k)\Lambda$ on $\psHS_f$; $\mathsf m_\sigma$ is a homomorphism, $\mathsf m_\sigma(k)=1$ if $kL=L$,
and, for $\sigma'=c\circ\sigma\circ c$, $\mathsf m_{\sigma'}(k)\,\aubar{\mathsf m_\sigma(k)}=1$.
\end{enumerate}
\end{lemma}
\begin{proof}
(a) By (W1), $\omega(u)^{-1}\psrho_f(\pcjB x)\omega(u)=\psrho_f(u^{-1}\pcjB x)=\psrho_f(ax)$, which with
Proposition~\ref{pcj:Lambda} gives the first law. The last assertion is Lemma~\ref{pcj:char} with
$\pcjB\aubar\pcjB=a\aubar a$; as $\langle ax,ax\rangle=a\aubar a\langle x,x\rangle$, it gives
$\sigma(e(x))=e(ax)$, hence the second law.

(b) For $l\in L$, $a^{-1}l\in L$ and $\psrho_f(l)\Lambda'\psvac=\Lambda'\psrho_f(a^{-1}l)\psvac=\sigma
(\epsilon_L(a^{-1}l))\Lambda'\psvac$, and $\sigma(\epsilon_L(a^{-1}l))=\psi_f(a\aubar a\cdot(-\frac12\omega_0
\jmath)(a\aubar a)^{-1}\langle l,l\rangle)=\epsilon_L(l)$. By (W4) $\Lambda'\psvac$ lies on the line
of $\psvac$, and it is not $0$.

(c) $e(x)e(y)/e(x+y)=\psi_f(-\frac12\omega_0\jmath(\langle x,y\rangle+\langle y,x\rangle))$ and
$\beta(x,y)=\psi_f(\frac12\jmath\langle y,x\rangle)$; with $s=\omega_0\langle x,y\rangle$ and
$1-\omega_0=\aubar{\omega_0}$ the product is $\psi_f(\frac12\jmath(\aubar s-s))$, and $\psi_f(\frac12
\jmath\aubar s)=\psi_f(-\aubar{\frac12\jmath s})=\psi_f(\frac12\jmath s)^{-1}$ by fact (a); this is the
formula. For $l'\in L\cap k^{-1}L$, $\pcjsh(kq)\pcjsh(kl')\psvac=\pcjsh(kq)\psvac$ (as $kl'\in L$) and
$\pcjsh(kq)\pcjsh(kl')=\psi_f(-\jmath\omega_0\langle q,l'\rangle)\pcjsh(k(q+l'))=\pcjsh(k(q+l'))$ by fact
(c); so the terms depend only on $q$ modulo $L\cap k^{-1}L$. For $l\in L$ the same computation gives
$\pcjsh(kl)\pcjsh(kq)\psvac=\pcjsh(k(q+l))\psvac$, and $q\mapsto q+l$ permutes the classes; so
$\pcjsh(kl)\pcjS_k=\pcjS_k$, that is $\psrho_f(kl)\pcjS_k=e(l)^{-1}\pcjS_k=\epsilon_L(l)\pcjS_k$ (note
$e(kl)=e(l)$). Also $\psrho_f(kl)\omega(k)\psvac=\omega(k)\psrho_f(l)\psvac=\epsilon_L(l)\omega(k)\psvac$ by
(W1). As $kL$ is self-dual, (W4) shows that $\omega(k)\psvac$ is a multiple of $\pcjS_k$ once
$\pcjS_k\ne0$. The vectors $\pcjsh(x)\psvac$, $x\in W(\auA_{E,f})/L$, are, up to roots of unity, the
basis of (W4); they are unit vectors (the $\psrho_f(x)$ are unitary) and pairwise orthogonal, being
eigenvectors of the unitary operators $\psrho_f(l)$, $l\in L$, with the characters $l\mapsto[l,x]
\epsilon_L(l)$, which are distinct for distinct classes because $L$ is self-dual. The classes $kq$,
$q\in L/(L\cap k^{-1}L)$, are distinct modulo $L$; so $\pcjS_k$ is a sum of $I_k$ orthonormal vectors,
$\|\pcjS_k\|^2=I_k$, and $1=\|\omega(k)\psvac\|^2=|\delta_k|^2I_k$ (normalising $\|\psvac\|=1$).

(d) By (a), $\Lambda'\pcjsh(kq)\psvac=\pcjsh(akq)\Lambda'\psvac=\mathsf c\,\pcjsh(k(aq))\psvac$, and
$q\mapsto aq$ permutes $L/(L\cap k^{-1}L)$ ($a$ is a unit scalar commuting with $k$); so $\Lambda'
\pcjS_k=\mathsf c\,\pcjS_k$ and $\Lambda'\omega(k)\psvac=\sigma(\delta_k)\mathsf c\,\pcjS_k=\mathsf
m_\sigma(k)\,\omega(k)\Lambda'\psvac$. For $x\in W(\auA_{E,f})$, by (W1) and $e(kx)=e(x)$,
$\omega(k)\pcjsh(x)=\pcjsh(kx)\omega(k)$; hence
\[
\Lambda'\omega(k)\pcjsh(x)\psvac=\Lambda'\pcjsh(kx)\omega(k)\psvac=\pcjsh(akx)\Lambda'\omega(k)\psvac=
\mathsf m_\sigma(k)\pcjsh(akx)\omega(k)\Lambda'\psvac=\mathsf m_\sigma(k)\,\omega(k)\Lambda'\pcjsh(x)\psvac,
\]
and these vectors span $\psHS_f$. As $u$ is central in $\auU(W)(\auA_f)$ and $\omega$ is a
representation, $\Lambda\omega(k)=\omega(u)\Lambda'\omega(k)=\mathsf m_\sigma(k)\omega(k)\Lambda$. The
homomorphism property follows from $\Lambda\omega(kk')=\Lambda\omega(k)\omega(k')$ and $\Lambda\ne0$. If
$kL=L$ then $I_k=1$, $\pcjS_k=\psvac$ and $\omega(k)\psvac=\psvac$ (W4), so $\delta_k=1$. Finally
$\sigma'(x)=\aubar{\sigma(\aubar x)}$ and $\aubar{\delta_k}=I_k^{-1}\delta_k^{-1}$ with $I_k\in\Q$, so
$\sigma'(\delta_k)=\aubar{I_k^{-1}\sigma(\delta_k)^{-1}}$ and $\mathsf m_{\sigma'}(k)=I_k^{-1}/(\delta_k
\aubar{\sigma(\delta_k)})=\aubar{\delta_k}/\aubar{\sigma(\delta_k)}=\aubar{\mathsf m_\sigma(k)}^{-1}$.
\end{proof}

\emph{The block datum.} By (W5), at a place $v=w\aubar w$ of $E^+$ with $w\in\mathsf M$, $\psHS_v$ is
the space of Schwartz--Bruhat functions on $E_w^{2z}$, the vacuum is $\plone[\alpha_w\auO_w^{2z}]$, i.e.\ $\plone[\auO_w^{2z}]$ in the coordinate $T=\tilde x$ of (N5) (Lemma~\ref{pgx:dict}; at a chosen place $\alpha_w$ is a unit and the canonical idele of \S\ref{ps:ss-weil} is trivial there, so $T=x$), in which all functions on $E_w^{2z}$ are written below
(proof of Lemma~\ref{pex:vac}), $\psrho(y_w,0)$ is the translation by $y_w$ and $\psrho(0,y_{\aubar w})$
the multiplication by $x\mapsto[(x,0),(0,y_{\aubar w})]$ (with $y$ read in the same coordinate). Write $E_{\mathsf M}^{2z}=\bigoplus_{w\in
\mathsf M}E_w^{2z}$. The block datum of \S\ref{pex:ss-set} for $\vartheta=1$ is $\Phi_{\mathsf M}=
\Psi_{\mathsf M}\otimes\psvac^{\mathsf M}$, $\Psi_{\mathsf M}=\bigotimes_{w\in\mathsf M}\plone[T\equiv e_1
\bmod w]$, with $\psvac^{\mathsf M}$ the vacuum at the other places.

\begin{lemma}[The block datum]\label{pcj:datum}
For every Schwartz--Bruhat function $\Psi$ on $E_{\mathsf M}^{2z}$,
$\Lambda'(\Psi\otimes\psvac^{\mathsf M})=\mathsf c\,(\sigma\circ\Psi)\otimes\psvac^{\mathsf M}$. In
particular $\Lambda\Phi_{\mathsf M}=\mathsf c\,\omega(u)\Phi_{\mathsf M}$.
\end{lemma}
\begin{proof}
Let $y=(y_w,y_{\aubar w})_{w\in\mathsf M}\in W(\auA_{E,f})$ be supported at the places above $\mathsf M$,
and put $Y_1=(y_w,0)_w$, $Y_2=(0,y_{\aubar w})_w$. Then $\psrho_f(y)=\beta(Y_1,Y_2)^{-1}\psrho_f(Y_1)
\psrho_f(Y_2)$, so $\psrho_f(y)\psvac=F_y\otimes\psvac^{\mathsf M}$ with
\[
F_y(x)=\beta(Y_1,Y_2)^{-1}\,[(x,0)+Y_1,Y_2]\,\plone[\auO^{2z}](x+y_w)\qquad(x\in E_{\mathsf M}^{2z}).
\]
Since $a_w=1$ for $w\in\mathsf M$, $a(x,0)=(x,0)$, $aY_1=Y_1$ and $ay=(y_w,a_{\aubar w}y_{\aubar w})$. By
Lemma~\ref{pcj:vac}(a), $\sigma(\beta(U,U'))=\beta(aU,aU')$ and $\sigma([U,U'])=[aU,aU']$ for all $U,U'$
(both are values of $\psi_f$ at multiples of $\langle U',U\rangle$, and $\langle aU',aU\rangle=a\aubar a
\langle U',U\rangle$). Hence $\sigma\circ F_y=F_{ay}$, and by Lemma~\ref{pcj:vac}(a),(b),
$\Lambda'(F_y\otimes\psvac^{\mathsf M})=\psrho_f(ay)\Lambda'\psvac=\mathsf c\,F_{ay}\otimes\psvac^{\mathsf M}=
\mathsf c\,(\sigma\circ F_y)\otimes\psvac^{\mathsf M}$. The functions $F_y$ span the Schwartz--Bruhat
functions: up to non-zero constants they are the products of a character $x\mapsto[(x,0),(0,q)]$ and of
a translate of $\plone[\auO^{2z}]$; these characters are all the continuous characters of
$E_{\mathsf M}^{2z}$ ($[\,,]$ is non-degenerate and $E_w^{2z}$, $E_{\aubar w}^{2z}$ are isotropic and
in duality), and the indicator of a coset of $\varpi^j\auO^{2z}$ inside $\auO^{2z}$ is an average of such
characters times $\plone[\auO^{2z}]$. By $\sigma$-semilinearity of both sides the first assertion
follows. For the second, $\Psi_{\mathsf M}$ takes the values $0,1$, so $\sigma\circ\Psi_{\mathsf M}=
\Psi_{\mathsf M}$ and $\Lambda\Phi_{\mathsf M}=\omega(u)\Lambda'\Phi_{\mathsf M}=\mathsf c\,\omega(u)
\Phi_{\mathsf M}$.
\end{proof}

\subsubsection{The lifts of characters and their conjugates}\label{pcj:ss-lift}

Let $p=(\nu,z)$ be a critical pair, $\nu$ of weight $-c$, $l_\tau=c(\tau)-z$, $D=\sum_\tau l_\tau$. By
\S\ref{po:ss-set} and \S\ref{pex:ss-set} (with $\vartheta=1$: $a=1$, $b=\nu\upsilon^z$), the
representation defining $\varphi_{p,1}$ is
\[
\omega_\nu(g,h)=\nu(g)^{-1}\upsilon(g)^{-z}\,\omega(k_{g,h}),\qquad k_{g,h}\colon x\mapsto hxg^{-1},
\]
for $g\in\auU(1)(\auA_{E^+})$, $h\in\auU(V^{(z)})(\auA_{E^+})$, where $\omega$ is the Weil representation
of $\auU(W)(\auA)$ defined with $\upsilon$; we write $\omega_{\nu,f}$, $\nu_f$, $\upsilon_f$ for the finite
parts. Note $k_{g,h}k_{g',h'}=k_{gg',hh'}$, and the scalar $u$ of \S\ref{pcj:ss-vac} is $k_{u^{-1},1}$.
Put $\Phi^\circ=c_{\mathrm{vac}}^{2z}\Phi_{\mathsf M}$. The weight of $\poA(V^{(z)};p,1)$ has one row of
length $l_\tau$ at $\tau$, so $W_\lambda=\bigotimes_\tau\operatorname{Sym}^{l_\tau}\C^{2z}\subseteq
\bigotimes_\tau(\C^{2z})^{\otimes l_\tau}$, with the form for which the pure tensors $e_I$ of basis vectors
are orthonormal (\S\ref{po:ss-set}). For a multi-index $\beta=(\beta_\tau)_\tau$, $\beta_\tau\in
\N^{2z}$ of degree $l_\tau$, let $e_\beta$ be the sum of the $e_I$ of content $\beta$ and $n_\beta=
\langle e_\beta,e_\beta\rangle\in\Z_{>0}$ their number; the $e_\beta$ form an orthogonal basis. By
\eqref{po:eq-poly}, $\pop_\xi(z)=\langle\xi,\bigotimes_\tau z_\tau^{\otimes l_\tau}\rangle$ in the Fock
variables $z_\tau=(z_{\tau,\mathsf c})_{\mathsf c}$; it is homogeneous of degree $l_\tau$ in $z_\tau$, and
$P_\beta:=\pop_{e_\beta}=n_\beta z^\beta$. For $\sigma\in\auAut(\C/K)$ let $\pi\beta$ be the multi-index
$(\pi\beta)_{\tau'}=\beta_{\sigma^{-1}\tau'}$; then $\sigma_*P_\beta=n_\beta z^{\pi\beta}$, and $n_{\pi\beta}
=n_\beta$.

\begin{lemma}[The lift as a class sum of jets]\label{pcj:lift}
Let $\pcjG$ be a set of representatives of the finite set $\auU(1)(E^+)\backslash\auU(1)(\auA_{E^+,f})/
K^1_M$. There is $w_0>0$, depending only on $z$, $\mathsf M$ and the fixed Haar measure on
$[\auU(1)]$, such that for every critical pair $p=(\nu,z)$, $h\in\auU(V^{(z)})(\auA_{E^+,f})$ and
$\xi\in W_\lambda$,
\[
\langle\xi,\varphi_{p,1}(h)\rangle=w_0\sum_{g\in\pcjG}\pcjJ_{\pop_\xi,\pi_Ms_0}\bigl(\omega_{\nu,f}(g,h)
\Phi^\circ\bigr).
\]
Moreover, for every polynomial $Q$ homogeneous of degree $l_\tau$ in $z_\tau$, every $s\in\auU(1)(E^+)$,
$\kappa\in K^1_M$ and $g$, $\pcjJ_{Q,c}(\omega_{\nu,f}(sg\kappa,h)\Phi^\circ)=\pcjJ_{Q,c}(\omega_{\nu,f}
(g,h)\Phi^\circ)$.
\end{lemma}
\begin{proof}
For $g\in\auU(1)(\auA_{E^+})$ write $g=g_\infty g_f$. The proof of Lemma~\ref{pex:arch} shows that
$\omega_{\nu,\infty}(g_\infty,1)$ fixes every polynomial homogeneous of degree $l_\tau$ in $z_\tau$ (it
multiplies such a polynomial by the product of $\epsilon^{-l_\tau}$, of the splitting character and of
$\nu_\tau(\epsilon)^{-1}\upsilon_\tau(\epsilon)^{-z}$, which is trivial). Hence, by \eqref{pk:eq-bargmann},
$\theta(\omega_\nu(g,h)(Q\otimes\Phi))=\theta(Q\otimes\omega_{\nu,f}(g_f,h)\Phi)=\pcjJ_{Q,1}(\omega_{\nu,f}
(g_f,h)\Phi)$. For $s\in\auU(1)(E^+)$, $\nu(s)=\upsilon(s)=1$ and $k_{s,1}\in\auU(W)(E^+)$, so
$\theta\circ\omega_\nu(s,1)=\theta$ by (W1); this gives the invariance under $s$ (for $c=1$, hence for
all $c$ by homogeneity). For $\kappa\in K^1_M$, $\omega_{\nu,f}(\kappa,1)\Phi_{\mathsf M}=\Phi_{\mathsf M}$:
at $w\in\mathsf M$ by \eqref{po:eq-w5} ($\kappa_w\equiv1\bmod w$, the tame character and $\eta_w$ are
trivial on it, and $x\equiv e_1$ if and only if $x\kappa_w\equiv e_1$), and at the other places because
$\nu$, $\upsilon$ are unramified there and $\omega(k_{\kappa,1})$ fixes the local vacuum, $k_{\kappa,1}$
being a unit scalar (W4). Now \eqref{po:eq-lift} with $y=\plone$, \eqref{pex:eq-phi} and these
invariances give $\langle\xi,\varphi_{p,1}(h)\rangle=\kappa_p[K^1:K^1_M]\,v_0\sum_{g\in\pcjG}\pcjJ_{\pop_\xi,1}
(\omega_{\nu,f}(g,h)\Phi_{\mathsf M})$, where $v_0$ is the common volume of the sets $\auU(1)(E^+)
\backslash\auU(1)(E^+)\auU(1)(\R)gK^1_M$ (they are translates of one another, $\auU(1)$ being
commutative). Since $\kappa_p=c_{\mathrm{vac}}^{2z}(\pi_Ms_0)^D$ and $\pop_\xi$ is homogeneous of degree
$D$, $\kappa_p\pcjJ_{\pop_\xi,1}(\omega_{\nu,f}(g,h)\Phi_{\mathsf M})=\pcjJ_{\pop_\xi,\pi_Ms_0}(\omega_{\nu,f}
(g,h)\Phi^\circ)$; put $w_0=[K^1:K^1_M]v_0$.
\end{proof}

\begin{lemma}[Conjugate characters]\label{pcj:nusigma}
Let $\chi$ be an automorphic character of $\auU(1)$ of weight $-c$ and $\sigma\in\auAut(\C/K)$. There is a
unique automorphic character $\chi^{(\sigma)}$ of $\auU(1)$ with finite part $\sigma\circ\chi_f$; it has
weight $-c^\sigma$, $c^\sigma(\sigma\tau)=c(\tau)$, it is unitary, and its restriction to the units at a
finite place $w$ is $\sigma$ composed with that of $\chi$. For $\sigma'=c\circ\sigma\circ c$,
$\chi^{(\sigma')}=\chi^{(\sigma)}$. If $(\nu,z)$ is a critical pair, so is $(\nu^{(\sigma)},z)$.
\end{lemma}
\begin{proof}
Write $\chi_\infty(x)=\prod_\tau\tau(x)^{e(\tau)}$ on $\auU(1)(E^+\otimes\R)$, identified with
$\prod_\tau\{|\epsilon|=1\}$ through the $\tau$, where $e(\tau)$ is determined by $c(\tau)$ through the
convention for weights; put $\chi^{(\sigma)}_\infty(x)=\prod_{\tau'}\tau'(x)^{e(\sigma^{-1}\tau')}$ and
$\chi^{(\sigma)}=(\sigma\circ\chi_f)\chi^{(\sigma)}_\infty$. It is a continuous homomorphism ($\chi_f$ is
trivial on an open subgroup, hence so is $\sigma\circ\chi_f$). For $s\in\auU(1)(E^+)$, $\chi_f(s)=
\chi_\infty(s)^{-1}=\prod_\tau\tau(s)^{-e(\tau)}$, so $\sigma(\chi_f(s))=\prod_\tau(\sigma\tau)(s)^{-e(\tau)}
=\chi_\infty^{(\sigma)}(s)^{-1}$: $\chi^{(\sigma)}$ is trivial on $\auU(1)(E^+)$. It is unitary because
$\auU(1)(E^+)\backslash\auU(1)(\auA_{E^+})$ is compact. Uniqueness: two such characters agree on the
subgroup $\auU(1)(E^+)\auU(1)(\auA_{E^+,f})$, which is dense in $\auU(1)(\auA_{E^+})$ because $\auU(1)(E^+)$ is dense in $\auU(1)(E^+\otimes\R)$ (weak approximation for the torus $\auU(1)$, which is rational over $E^+$ by the Cayley parametrisation $y\mapsto y/\aubar y$). For $\sigma'$: $\sigma'(\chi_f(g))=\aubar{\sigma(
\aubar{\chi_f(g)})}=\aubar{\sigma(\chi_f(g))^{-1}}=\sigma(\chi_f(g))$, since $|\chi_f(g)|=1$ and
$|\sigma(\chi_f(g))|=1$ (unitarity of $\chi^{(\sigma)}$). For a critical pair: $z\le c^\sigma\le4b$ since
$c^\sigma$ is a permutation of $c$, and $\nu^{(\sigma)}$ is tamely ramified exactly where $\nu$ is.
\end{proof}

By Lemma~\ref{pcj:nusigma} applied to $\upsilon|_{\auU(1)}$ (of weight $1$),
$|\sigma(\upsilon_f(g))|=1$ for all finite $g$. For $\sigma\in\auAut(\C/K)$ define, on
$\auU(1)(\auA_{E^+,f})$,
\[
\rho_\sigma(g)=\mathsf m_\sigma(k_{g,1})\,\upsilon_f(g)^z\,\sigma(\upsilon_f(g))^{-z},
\]
a character (Lemma~\ref{pcj:vac}(d)). Put $\mathsf l_\sigma=\sigma(c_{\mathrm{vac}})^{2z}c_{\mathrm{vac}}^{-2z}
\mathsf c\ne0$; then $\Lambda\Phi^\circ=\mathsf l_\sigma\,\omega(u)\Phi^\circ$ by Lemma~\ref{pcj:datum}. The
constants $\mathsf l_\sigma$, $u$, $\mathsf c$ depend on $\sigma$ and $z$ only.

\begin{lemma}[The kernels]\label{pcj:kernel}
For every automorphic character $\nu$ of $\auU(1)$, every character $\chi$ of $\auU(1)(\auA_{E^+,f})$
and all finite $g,h$,
\[
\Lambda\bigl(\omega_{\nu,f}(g,h)\Phi^\circ\bigr)=\mathsf l_\sigma\,\mathsf m_\sigma(k_{1,h})\,
\frac{\chi(gu^{-1})\upsilon_f(gu^{-1})^z\mathsf m_\sigma(k_{g,1})}{\sigma(\nu_f(g))\,\sigma(\upsilon_f(g))^{z}}
\,\omega_{\chi,f}(gu^{-1},h)\Phi^\circ ,
\]
where $\omega_{\chi,f}(g,h)=\chi(g)^{-1}\upsilon_f(g)^{-z}\omega(k_{g,h})$.
\end{lemma}
\begin{proof}
{\emergencystretch=6em By $\sigma$-semilinearity, Lemma~\ref{pcj:vac}(d) and $k_{g,h}=k_{g,1}k_{1,h}$,
$\Lambda(\omega_{\nu,f}(g,h)\Phi^\circ)=\sigma(\nu_f(g)\upsilon_f(g)^z)^{-1}\allowbreak \mathsf m_\sigma(k_{g,1})
\allowbreak \mathsf m_\sigma(k_{1,h})\,\allowbreak \omega(k_{g,h})\Lambda\Phi^\circ$, and $\omega(k_{g,h})\Lambda\Phi^\circ=\mathsf
l_\sigma\,\omega(k_{g,h}k_{u^{-1},1})\Phi^\circ=\mathsf l_\sigma\,\omega(k_{gu^{-1},h})\Phi^\circ=\mathsf
l_\sigma\,\chi(gu^{-1})\allowbreak \upsilon_f(gu^{-1})^z\allowbreak \omega_{\chi,f}(gu^{-1},h)\Phi^\circ$.\par}
\end{proof}

\begin{lemma}[The character $\rho_\sigma$]\label{pcj:rho}
$\rho_\sigma$ is trivial on $\auU(1)(E^+)$ and on $K^1$, it has finite order, $|\mathsf m_\sigma(k_{g,1})|
=1$ for all $g$, and $\rho_{\sigma'}=\rho_\sigma$ for $\sigma'=c\circ\sigma\circ c$.
\end{lemma}
\begin{proof}
For $g\in K^1$, $k_{g,1}$ is a unit scalar, so $k_{g,1}L=L$ and $\mathsf m_\sigma(k_{g,1})=1$
(Lemma~\ref{pcj:vac}(d)), and $\upsilon_f(g)=1$ ($\upsilon$ is unramified at every finite place).

Let $s\in\auU(1)(E^+)$. Take any critical pair $p=(\nu,z)$ (for instance a flat one,
Lemma~\ref{psu:flat}(a)). By Proposition~\ref{pex:rallis}, $\varphi_{p,1}\ne0$; by
Lemma~\ref{pcj:lift} there are $g,h,\beta$ with $\pcjJ_{P_\beta,s_0}(f)\ne0$ for $f=\omega_{\nu,f}(g,h)
\Phi^\circ$. Put $f_s=\omega_{\nu,f}(sg,h)\Phi^\circ$; then $\pcjJ_{P_\beta,s_0}(f_s)=\pcjJ_{P_\beta,s_0}(f)$
(Lemma~\ref{pcj:lift}). By Proposition~\ref{pcj:Lambda} and \eqref{pcj:eq-jet}, $\pcjJ_{\sigma_*P_\beta,
t_\sigma}(\Lambda f_s)=\pcjJ_{\sigma_*P_\beta,t_\sigma}(\Lambda f)=\sigma(\pcjJ_{P_\beta,s_0}(f))\ne0$. Apply
Lemma~\ref{pcj:kernel} with $\chi=\nu^{(\sigma)}_f$ to $f$ and to $f_s$. Since $\sigma_*P_\beta$ is
homogeneous of degree $l_{\sigma^{-1}\tau'}$ in $z_{\tau'}$, which is the degree attached to the weight
$-c^\sigma$ of $\nu^{(\sigma)}$, and $(\nu^{(\sigma)},z)$ is a critical pair (Lemma~\ref{pcj:nusigma}),
Lemma~\ref{pcj:lift} for it gives $\pcjJ_{\sigma_*P_\beta,t_\sigma}(\omega_{\chi,f}(sgu^{-1},h)\Phi^\circ)=
\pcjJ_{\sigma_*P_\beta,t_\sigma}(\omega_{\chi,f}(gu^{-1},h)\Phi^\circ)$. Comparing the two instances of
Lemma~\ref{pcj:kernel}, the ratio of the scalars for $sg$ and for $g$ must be $1$; it is
\[
\frac{\chi(s)\upsilon_f(s)^z\mathsf m_\sigma(k_{s,1})}{\sigma(\nu_f(s))\sigma(\upsilon_f(s))^z}=
\mathsf m_\sigma(k_{s,1})\upsilon_f(s)^z\sigma(\upsilon_f(s))^{-z}=\rho_\sigma(s),
\]
as $\chi(s)=\sigma(\nu_f(s))$. So $\rho_\sigma$ is trivial on $\auU(1)(E^+)K^1$, a subgroup of finite
index; hence $\rho_\sigma$ has finite order and $|\rho_\sigma|=1$, and then $|\mathsf m_\sigma(k_{g,1})|=
|\sigma(\upsilon_f(g))/\upsilon_f(g)|^z=1$. Finally, by Lemma~\ref{pcj:vac}(d),
$\mathsf m_{\sigma'}(k_{g,1})=\aubar{\mathsf m_\sigma(k_{g,1})}^{-1}=\mathsf m_\sigma(k_{g,1})$, and
$\sigma'(\upsilon_f(g))=\sigma(\upsilon_f(g))$ by Lemma~\ref{pcj:nusigma}; so $\rho_{\sigma'}=\rho_\sigma$.
\end{proof}

\begin{definition}\label{pcj:def-psigma}
For a critical pair $p=(\nu,z)$ and $\sigma\in\auAut(\C/K)$ let $\nu^\sigma=\nu^{(\sigma)}\rho_\sigma^{-1}$
(with $\rho_\sigma$ extended by $1$ at the infinite places) and $p^\sigma=(\nu^\sigma,z)$.
\end{definition}

By Lemmas~\ref{pcj:nusigma} and~\ref{pcj:rho}, $\nu^\sigma$ is an automorphic character of weight
$-c^\sigma$ with the same ramification as $\nu^{(\sigma)}$ ($\rho_\sigma$ is unramified), so $p^\sigma$ is
a critical pair, flat if $p$ is; and $p^{\sigma'}=p^\sigma$ for $\sigma'=c\circ\sigma\circ c$.

\begin{proposition}[The conjugates of the coordinates]\label{pcj:coord}
Let $p=(\nu,z)$ be a critical pair, $h\in\auU(V^{(z)})(\auA_{E^+,f})$, and write $a_\beta(h)=\langle
e_\beta,\varphi_{p,1}(h)\rangle$ and $a^\sigma_{\beta'}(h)=\langle e_{\beta'},\varphi_{p^\sigma,1}(h)\rangle$.
Then
\[
\sigma\bigl(a_\beta(h)\bigr)=C_{\sigma,p}\,\mathsf m_\sigma(k_{1,h})\,a^\sigma_{\pi\beta}(h),\qquad
C_{\sigma,p}=\frac{\sigma(w_0)}{w_0}\Bigl(\frac{\sigma(\pi_M)\,t_\sigma}{\pi_Ms_0}\Bigr)^{D}\mathsf
l_\sigma\,\frac{\rho_\sigma(u)}{\sigma(\nu_f(u))\,\upsilon_f(u)^{z}},
\]
and $|C_{\sigma,p}|=|\sigma(w_0)/w_0|\,|\mathsf l_\sigma|$, which does not depend on $p$.
\end{proposition}
\begin{proof}
Lemma~\ref{pcj:kernel} with $\chi=\nu^\sigma_f$, for which $\chi(gu^{-1})=\sigma(\nu_f(g))\rho_\sigma(g)^{-1}
\sigma(\nu_f(u))^{-1}\rho_\sigma(u)$, gives, by the definition of $\rho_\sigma$,
\[
\Lambda\bigl(\omega_{\nu,f}(g,h)\Phi^\circ\bigr)=\mathsf l_\sigma\frac{\rho_\sigma(u)}{\sigma(\nu_f(u))
\upsilon_f(u)^z}\,\mathsf m_\sigma(k_{1,h})\,\omega_{\nu^\sigma,f}(gu^{-1},h)\Phi^\circ .
\]
By Lemma~\ref{pcj:lift}, homogeneity, \eqref{pcj:eq-jet} and Proposition~\ref{pcj:Lambda},
\[
\sigma(a_\beta(h))=\sigma(w_0)\sigma(\pi_M)^D\sum_{g\in\pcjG}\pcjJ_{\sigma_*P_\beta,t_\sigma}\bigl(\Lambda\,
\omega_{\nu,f}(g,h)\Phi^\circ\bigr),
\]
and $\sigma_*P_\beta=n_\beta z^{\pi\beta}=P_{\pi\beta}$, the polynomial $\pop_{e_{\pi\beta}}$ for the weight
of $p^\sigma$, homogeneous of the same total degree $D$. As $g\mapsto gu^{-1}$ permutes the double cosets
and the summands depend only on the double coset (Lemma~\ref{pcj:lift} for $p^\sigma$),
$\sum_g\pcjJ_{P_{\pi\beta},t_\sigma}(\omega_{\nu^\sigma,f}(gu^{-1},h)\Phi^\circ)=(t_\sigma/\pi_Ms_0)^D
a^\sigma_{\pi\beta}(h)/w_0$. This is the formula. For the absolute value: $|t_\sigma|=|s_0|$;
$|\sigma(\pi_M)|^2=M=|\pi_M|^2$ (as used in Step~3 of the proof of Proposition~\ref{au:kernel} in
\S\ref{ps:ss-kernel}, for every $\sigma$); $|\rho_\sigma(u)|=|\upsilon_f(u)|=1$, and $|\sigma(\nu_f(u))|=1$
by Lemma~\ref{pcj:nusigma}.
\end{proof}

\begin{proof}[Proof of Theorem~\ref{pcj:CL}]
Let $\sigma'=c\circ\sigma\circ c\in\auAut(\C/K)$, so that $\sigma(\aubar x)=\aubar{\sigma'(x)}$. For an
embedding $\tau$ of the CM field $E$, $c\circ\tau=\tau\circ c_E$, hence $\sigma'\circ\tau=c\circ(\sigma\tau)
\circ c_E=\sigma\circ\tau$: the permutations $\pi$ for $\sigma$ and $\sigma'$ coincide, and $p^{\sigma'}=
p^\sigma$ (Definition~\ref{pcj:def-psigma}). The norm of \S\ref{ss:au-spaces} is $\|\Xi\|^2=\sum_{h}
r_h\langle\Xi(h),\Xi(h)\rangle$, $h$ over representatives of $\auU(V^{(z)})(E^+)\backslash\auU(V^{(z)})
(\auA_{E^+,f})/K$ and $r_h=\#(\Gamma\cap hKh^{-1})^{-1}\in\Q_{>0}$ (proof of Lemma~\ref{po:naut}); the
level $K$ and its character are the same for $p$ and $p^\sigma$ (they are defined by $\Phi_{\mathsf M}$
and by the action of $\auU(V^{(z)})$, which does not involve $\nu$). With the orthogonal basis $e_\beta$,
$\langle\Xi(h),\Xi(h)\rangle=\sum_\beta|\langle e_\beta,\Xi(h)\rangle|^2/n_\beta$. So, by
Proposition~\ref{pcj:coord} for $\sigma$ and for $\sigma'$,
\[
\sigma\bigl(\|\varphi_{p,1}\|^2\bigr)=\sum_hr_h\sum_\beta\frac{\sigma(a_\beta(h))\,\aubar{\sigma'(a_\beta(h))}}
{n_\beta}=C_{\sigma,p}\aubar{C_{\sigma',p}}\sum_hr_h\,\mathsf m_\sigma(k_{1,h})\aubar{\mathsf m_{\sigma'}(k_{1,h})}
\sum_\beta\frac{|a^\sigma_{\pi\beta}(h)|^2}{n_{\pi\beta}} .
\]
By Lemma~\ref{pcj:vac}(d), $\mathsf m_\sigma(k)\aubar{\mathsf m_{\sigma'}(k)}=1$; and $\beta\mapsto\pi\beta$
is a bijection onto the multi-indices of the weight of $p^\sigma$. Hence $\sigma(\|\varphi_{p,1}\|^2)=
C_{\sigma,p}\aubar{C_{\sigma',p}}\,\|\varphi_{p^\sigma,1}\|^2$, and the theorem holds with
$a_\sigma(z)=|C_{\sigma,p}C_{\sigma',p}|=|\sigma(w_0)\sigma'(w_0)|\,w_0^{-2}\,|\mathsf l_\sigma\mathsf
l_{\sigma'}|>0$, which depends on $\sigma$ and $z$ only. The weight of $p^\sigma$ is $-c^\sigma$ with
$c^\sigma(\sigma\tau)=c(\tau)$ (Lemma~\ref{pcj:nusigma}).
\end{proof}

\begin{remark}\label{pcj:rem-use}
(i) In the proof of Theorem~\ref{psu:arch}, (CL) is used for a main pair $p$ and the flat pair $q$; $q^\sigma$
is flat, and $p^\sigma$ has $l^\sigma_{\sigma\tau}=l_\tau$, as required there. (ii) We do not decide whether the
unramified character $\rho_\sigma$ is trivial (that is, whether $\nu^\sigma=\nu^{(\sigma)}$): (CL) only
asks for some critical pair with the permuted weight, and Definition~\ref{pcj:def-psigma} avoids computing
the multiplier of the Weil representation on scalars, which would need the explicit splitting. (iii) Nothing in this subsection uses the size of the constants.
\end{remark}



%% ===== proofs/C-kisin.tex =====
% proofs/C-kisin.tex
% The explicit cup exponents: Theorem pll:exponent. Labels start with pki:.
\providecommand{\auO}{\mathcal O}
\providecommand{\auQlb}{\overline{\Q}_\ell}
\providecommand{\pglen}{\operatorname{len}}
\providecommand{\pllZf}{Z_f}
\providecommand{\pllrk}{\operatorname{rk}}
\providecommand{\plldef}{\delta^{\vee}}
% macros of this file (prefix \pki)
\providecommand{\pkiWin}{\operatorname{Win}}
\providecommand{\pkiC}{\mathcal C}
\providecommand{\pkilog}{\operatorname{lg}}

\section{The explicit cup exponents}\label{pki:sec}

This section concerns Theorem~\ref{pll:exponent}. In \S\ref{pki:ss-need} we determine what the
rest of the paper needs and explain why we do not replace the argument by a published theorem. In
\S\ref{pki:ss-arch} we state the four propositions on which everything rests
(Propositions~\ref{pki:generic}, \ref{pki:any}, \ref{pki:sparse}, \ref{pki:tame}), together with two
inputs on the data (Propositions~\ref{pki:chart}, \ref{pki:data}), and we deduce from them a bound
uniform in $\ell$ (Theorem~\ref{pki:crys}) and Theorem~\ref{pll:exponent} itself
(\S\ref{pki:ss-deduce}).

\emph{Dependencies.} Proved here in full: Theorem~\ref{pki:crys} and the deduction of
Theorem~\ref{pll:exponent} from Propositions~\ref{pki:generic}--\ref{pki:data}. These six
propositions are proved in \S\S\ref{pkc:sec}, \ref{pkd:sec}, \ref{pkj:sec} (and Lemma~\ref{pki:data-sizes}), from the theory of Kummer
charts and the pairing lemma written in \S\S\ref{pkp:sec}--\ref{pkw:sec}; the chain is recorded at the
end of \S\ref{pki:ss-arch}.

\emph{Conventions.} As in \S\ref{pll:sec}: $\ell$ a prime, $v\mid\ell$ a place of $E$,
$\Gamma=D_v=\operatorname{Gal}(\auQlb/E_v)$, $\auO$ a large coefficient ring, lengths normalised so
that $\pglen\auO/\ell=1$, $\gamma$ runs over the embeddings $\auO\to\C_\ell$. For a lattice $L$, its
crystalline cocycles $\pllZf(L)$, the cup-orthogonal $L^{\perp\pkiC}$, the dual defect
$\plldef_\pkiC(L)$ and ``cup exponent'' are as in Definition~\ref{pll:def}. For finite sets
$U(\gamma)\subset\Z$, $\pkiC(U)$ is the class of crystalline lattices whose jumps at each $\gamma$
lie in $U(\gamma)$. We write $\pkilog_\ell H=\lfloor\log H/\log\ell\rfloor$ for $H\ge1$.

\subsection{What is needed}\label{pki:ss-need}

Theorem~\ref{pll:exponent} is used only through \eqref{pll:eq-place}, summed over $v\mid\ell$ and
over the summands (Lemma~\ref{pll:sum}), and then over $\ell$ in Lemma~\ref{pg:sums}: what must hold
is $\sum_{\ell\in T}\pgloc_\ell\log\ell\le c\,\aubud$ for a small absolute fraction $c$, where
$\aubud=f\,x\,b(\log b)^3/(mk)$ is the budget and $\pgloc_\ell$ is the loss per unit of rank of the
main summand. Lemma~\ref{pg:sums} obtains $\sum_\ell\pgloc_\ell\log\ell\le 10^8f(m/k)\,b\log b$,
leaving a margin of a factor $x(\log b)^2/(10^8m^2)\ge(\log b)^2/(4\cdot10^8)$.

\begin{remark}[Why a published bound does not suffice]\label{pki:rem-published}
The decisive feature of the required bound is the factor $1/k$, with $k=r+1-m\ge x^5/2$ by
\eqref{eq:au-scales}. Any bound that charges, at each $\ell$ and each embedding, an amount growing
with the \emph{largest} Hodge--Tate weight $\asymp b$ fails by a power of $x$. Indeed, the bounds
in the literature for torsion crystalline (or semistable) representations in large weight
(Fontaine--Laffaille outside its range, the results of T.~Liu and of Caruso--Liu on lifting maps
between torsion crystalline representations, Gee--Liu--Savitt type results) are expressed through
the maximal weight $H$ and the ramification index $e$; even under the most favourable reading,
an exponent of size $\asymp\pkilog_\ell(eH)+1$ for every $\ell\le eH$, the resulting contribution
is, by the argument of Lemma~\ref{pll:exp} and of Lemma~\ref{pll:sum},
\[
\sum_{\ell\le 4b}4f\,r'\,(\pkilog_\ell(4b)+1)\log\ell\ \ge\ 4f\,r'\sum_{\ell\le4b}\log 4b\cdot\tfrac12
\ \gg\ f\,r'\,b,
\]
whereas the budget allows $r'\cdot c\,\aubud\le r'f\,b\,x^{-4.4}$. The argument below gains the
factor $1/k$ from two facts specific to the data: at all but a sparse set of embeddings (the
\emph{thin} ones, not $b$-places) the far weights are of size $G\ll b$, and the main summand has at
most $\lfloor r'm/k\rfloor$ jumps in the low window. No published theorem we know bounds the
torsion of crystalline extension classes embedding by embedding in this way. We therefore write a
direct argument. (For $\ell\ge4b+2$, $\ell\nmid q_1dM$, see also Remark~\ref{pll:FL} for another route;
the argument below covers all $\ell$ uniformly and does not use it.)
\end{remark}

\subsection{Architecture}\label{pki:ss-arch}

\subsubsection*{Windows, letters, sizes}

For integers $\eta,W\ge0$, $a\ge0$ and a sign $\varepsilon\in\{\pm1\}$, the \emph{windows} are
\[
\pkiWin(\eta,\varepsilon,a,W)=[-\eta,\eta]\cup\varepsilon\,[a,a+W]\subset\Z
\]
(the \emph{low window} and the \emph{far window}). For $n\ge0$ and $a\in\Z$ we say that $\ell$ has
\emph{no multiple near $a$ at distance $n$}, written $\mathrm{NM}(\ell,n,a)$, if $\ell$ divides no
nonzero integer $y$ with $|y|\le n$, $|y-a|\le n$ or $|y+a|\le n$.

Let $e$, $f$ be the ramification index and residue degree of $E_v/\Q_\ell$, and $d_v=0$ if $e=1$,
$d_v=v_\ell(e)+1$ otherwise. Fix $\zeta\in\auO$ of exact order $\ell^f-1$ (possible as $\auO$ is
large). The \emph{letters} are the $(\ell^f-1)$-th roots of unity $z\in\C_\ell$; the \emph{letter}
of $\gamma$ is $\gamma(\zeta)$; the Frobenius acts on letters by $z\mapsto z^\ell$. For jump sets
$U$, the \emph{letter jump set} at $z$ is $U_z=\{0\}\cup\bigcup_{\gamma(\zeta)=z}U(\gamma)$.

\begin{definition}[Sizes]\label{pki:def-sizes}
Integers $Q,m,h,G,b,\delta,\eta,W\ge0$ satisfy the \emph{size conditions} if
\begin{enumerate}
\item[(S1)] $Q\ge1$ and $99\,(Q^2+(W+1)Q+Q+W+m)\le h$;
\item[(S2)] $99\,h^3\le G$ and $99\,(Gh^3+G)\le b$;
\item[(S3)] $2\log(4b)<m$;
\item[(S4)] $\delta<\eta$ and $99\,((\eta+W)Q+\eta+h)\le h^3$.
\end{enumerate}
The \emph{loss} at $\ell$ is $\lambda=\lambda_\ell(h,b)=v_\ell\bigl((2b)!\cdot(\prod_{b-30h^3<n\le
b+30h^3}n)^{h^3}\bigr)$; with $h=r^8$ it is the $\lambda_\ell$ of \S\ref{pg:ss-statement}.
\end{definition}

\begin{lemma}\label{pki:elem}
Let $\lambda=\lambda_\ell(h,b)$. (i) $2b<\ell(\lambda+1)$. (ii) If $\ell$ divides some
$n\in(b-30h^3,b+30h^3]$, then $\lambda\ge h^3$. (iii) If (S3) holds and $1\le H\le4b$ then
$\pkilog_\ell H<m$. (iv) Under (S1)--(S4), $\pkilog_\ell h^3<m$, $d_v\le\pkilog_\ell e+1$, and if
$e\le h$ then $\pkilog_\ell h^3+1+d_v\le h$.
\end{lemma}

\begin{proof}
(i) $\lambda\ge v_\ell((2b)!)\ge\lfloor2b/\ell\rfloor$, and $2b<\ell(\lfloor2b/\ell\rfloor+1)$.
(ii) $n^{h^3}$ divides the product. (iii) $\log\ell\ge\log2>1/2$, so $\pkilog_\ell H\le2\log
H\le2\log 4b<m$. (iv) $h^3\le b\le4b$ by (S2) and $e\le h$, so (iii) applies to $h^3$ and to $e$;
$d_v\le v_\ell(e)+1\le\pkilog_\ell e+1$. Hence $\pkilog_\ell h^3+1+d_v<2m+2\le h$ by (S1).
\end{proof}

\begin{definition}[Letters that read thickness]\label{pki:def-reads}
Let each embedding $\gamma$ be \emph{thick} or \emph{thin} and carry a sign $\varepsilon(\gamma)$, and
let $T$ (a set of letters) and $s$ (a sign on letters) be given. We say that \emph{the letters read
thickness} (for $Q$, $h$) if: $e\le h$; for every $\gamma$ that is thick, and for every $\gamma$ if
$\ell>Q$, one has [$\gamma$ thick $\Leftrightarrow\gamma(\zeta)\in T$] and
$\varepsilon(\gamma)=s(\gamma(\zeta))$; either $e=1$ and $\ell\ge7$, or $\ell\le Q$ and $s$ is
constant; and, if $\ell\le Q$, there is no letter $z$ with $z$, $z^{\ell^{a}}$, $z^{\ell^{a+a'}}$ all
in $T$ for some $a,a'\in[1,\pkilog_\ell h^3]$.
\end{definition}

A lattice $L$ \emph{has few low jumps} (for $h,\mu,\eta$ and jump sets $U,U'$) if every negative
element $u$ of every $U(\gamma)\cup U'(\gamma)$ satisfies $2u\le-h$, and at every $\gamma$ at most
$\mu$ jumps of $L$, with multiplicity, lie in $[-\eta,\eta]$.

\subsubsection*{The four core propositions}

In Propositions~\ref{pki:generic}--\ref{pki:sparse}, $U$ and $U'$ are jump sets, $L$ is a crystalline
lattice with jumps in $U'$, and the class is $\pkiC(U)$; at each letter $z$ an integer $b_z\ge0$
and a sign $\varepsilon_z$ are given.

\begin{proposition}[Generic primes]\label{pki:generic}
Let $e=1$, $\ell\ge5$, and $a,\eta,W\ge0$ with $\eta\ge1$, such that for every letter $z$:
$b_z+W+\eta\le(\ell-3)a$; $6(a+\eta+W)<b_z$; $\mathrm{NM}(\ell,6(a+\eta+W),b_z)$; $U_z\subseteq
\pkiWin(\eta,-\varepsilon_z,b_z,W)$ and $U'_z\subseteq\pkiWin(\eta,\varepsilon_z,b_z,W)$.
\begin{enumerate}
\item[(i)] If no element of any $U(\gamma)\cup U'(\gamma)$ lies in $[-a,0)$, or if every
$U'(\gamma)\subseteq[-\eta,0]$, then $\plldef_{\pkiC(U)}(L)=0$.
\item[(ii)] If $L$ has at most $\mu$ jumps in $[-\eta,\eta]$ at each $\gamma$, then
$\plldef_{\pkiC(U)}(L)\le11f\bigl(\mu(a+1)+\pllrk_\auO L\cdot(W+1)\bigr)$.
\end{enumerate}
\end{proposition}

\begin{proposition}[Any pattern of $b$-letters]\label{pki:any}
Let $e=1$, $\ell\ge7$, and $\eta\ge1$, $V,M,G,b,K\ge0$ with $V+4\eta\le M\le K$, $0<G\le b$,
$b<K\ell^2$, $2G<K\ell$. Let $B$ be a set of letters, $b_z=b$ for $z\in B$ and $b_z=G$ otherwise,
$U_z\subseteq\pkiWin(\eta,-\varepsilon_z,b_z,V)$ and $U'_z\subseteq\pkiWin(\eta,\varepsilon_z,b_z,V)$.
Then $\pkiC(U)$ has cup exponent $10^4\bigl(K+\pkilog_\ell(2(b+M))\bigr)$ on $L$.
\end{proposition}

\begin{proposition}[Sparse $b$-letters]\label{pki:sparse}
Let $\eta,J\ge1$, $V,G,b,M,\vartheta\ge0$ with $G+V\le\eta$. Let $B$, $b_z$ and the window
conditions be as in Proposition~\ref{pki:any} (with width $V$). Assume: there is no letter $z$ with
$z,z^{\ell^s},z^{\ell^{s+t}}\in B$ for some $s,t\in[1,J]$; either [$e=1$, $\ell\ge7$ and
$10\ell^J\eta\le b$] or [$\varepsilon_z$ is constant and $8e\le\ell^{J+1}$];
$2e(b+5\eta)\le\ell^{J+1}\vartheta$; and $20(J+1+d_v)\eta+\pkilog_\ell(2(b+5\eta))\le M$. Then
$\pkiC(U)$ has cup exponent $10^3\bigl(\vartheta+(J+1)M\bigr)$ on $L$.
\end{proposition}

\begin{proposition}[Tame types]\label{pki:tame}
Let $E_v=\Q_\ell$, $r,N\ge1$, $e_0=\operatorname{lcm}_{1\le i\le r}(\ell^i-1)$. For each inertial
$s\in\Gamma$ let $K(s)$ be a set of polynomials over $\auO$ containing $X^{e_0}-1$ and the image of a
monic polynomial over $\Z_\ell$ of degree $\le N$; let $\pkiC(U,K)$ be the class of potentially
crystalline lattices with jumps in $U$ on whose $D_{\mathrm{pcris}}$ each $P\in K(s)$ kills $s$, for
every inertial $s$. Let $L\in\pkiC(U',K^{-1})$, where $K^{-1}(s)=K(s^{-1})$. Let $\eta,W,b\ge0$ and
a sign $\varepsilon$ with $U(\gamma)\subseteq\pkiWin(\eta,-\varepsilon,b,W)$ and
$U'(\gamma)\subseteq\pkiWin(\eta,\varepsilon,b,W)$ for all $\gamma$, and
$\ell>b+6(W+\eta)+r^2(N+1)^2+2$. Then $\pkiC(U,K)$ has cup exponent
$3r^2(N+1)^2+150(\eta+W+1)$ on $L$.
\end{proposition}

\subsubsection*{Two inputs on the data}

\begin{proposition}[A chart for the data]\label{pki:chart}
Let $Q=kt$, $h=r^8$. Call $\gamma$ thick if the complex embedding $\sigma_\gamma$ of $E$ attached to
$\gamma$ through $\iota$ and $v$ (Convention~\ref{pkj:conv}) lies over a $b$-place, and let
$\varepsilon(\gamma)$ be the sign given by the CM type at $\sigma_\gamma$. For every $\ell$ and $v\mid\ell$,
every $\auO\supseteq\auO_{E_v}$ and every $\zeta\in\auO$ of exact order $\ell^f-1$ there are $T$ and $s$
for which the letters read thickness (Definition~\ref{pki:def-reads}).
\end{proposition}

\begin{proposition}[The summands fit the windows]\label{pki:data}
Let $h=r^8$, $Q=kt$, $G$ the weight at the places that are not $b$-places, $\eta=h^2+1$, $W=4t$, and
for $0\le j\le m$ let $\delta_j=j\Delta+h-2t$. Then (S1)--(S4) hold for $(Q,m,h,G,b,\delta_j,\eta,W)$.
Let $L_c$ and $\mathcal C_v$ be as in Proposition~\ref{pll:main} and let $U(\gamma)$ be the jump set
of $\mathcal C_v$, $U'(\gamma)$ that of $L_c^*$. There are $j\le m$ and, at each $\gamma$, a sign
$\varepsilon(\gamma)$, with $U(\gamma)\subseteq\pkiWin(\eta,-\varepsilon(\gamma),a_\gamma+\delta_j,W)$
and $U'(\gamma)\subseteq\pkiWin(\eta,\varepsilon(\gamma),a_\gamma+\delta_j,W)$, where $a_\gamma=b$ if
$\gamma$ is thick and $a_\gamma=G$ otherwise. Moreover $L_{\mathrm{main}}^*$ has few low jumps for
$(h,\lfloor r'm/k\rfloor,\eta)$, and for a character $L_e$ every $U'(\gamma)\subseteq[-\eta,0]$.
At $\ell=p_i$, $\mathcal C_v\subseteq\pkiC(U,K)$ and $L_c^*\in\pkiC(U',K^{-1})$ for a $K$ as in
Proposition~\ref{pki:tame} with $N=\max(r+t,2r)\,r$, $\varepsilon(\gamma)$ and $a_\gamma$ do not depend on $\gamma$ (write $a$ for the common value), and
$\ell>a+\delta_j+6(W+\eta)+r^2(N+1)^2+2$ for $a\in\{b,G\}$.
\end{proposition}

\begin{lemma}[Sizes of the data]\label{pki:data-sizes}
With $h=r^8$, $Q=kt$, $G=a/x_1$ (the weight $\kappa(\tau)$ at the places $\tau\notin\auPb$),
$\eta=h^2+1$, $W=4t$ and $\delta_j=j\Delta+h-2t$, $0\le j\le m$, conditions (S1)--(S4) hold ($G$ is an integer since $x_1\mid b$, (D4) of Definition~\ref{an:data}). This is
the first assertion of Proposition~\ref{pki:data}.
\end{lemma}

\begin{proof}
By \eqref{eq:au-scales}: $x^5\le r\le3x^5$, $t=5r$, $k\le r$, $\sqrt x\le m\le2\sqrt x$, $\Delta=2h$,
$(\log b)^{100}\le x$, $x^{10^6}\le x_1$, $x_1^{10^6}\le b$, $a=xb$. (S1): $Q\le5r^2$, $W=20r$, so
$Q^2+(W+1)Q+Q+W+m\le25r^4+105r^3+25r^2+m\le200r^4$, and $99\cdot200r^4\le r^8=h$ since
$r\ge x^5\ge 10^2$. (S2): $h^3=r^{24}\le3^{24}x^{120}$, while $G=xb/x_1\ge b/x_1\ge
x_1^{10^6-1}\ge x^{10^{11}}$, so $99h^3\le G$; and $99(Gh^3+G)\le 198\,h^3\,xb/x_1\le b$ because
$198\cdot3^{24}x^{121}\le x_1$. (S3): $2\log(4b)\le3\log b\le3x^{1/100}<\sqrt x\le m$. (S4):
$\delta_j\le2mh+h\le(4\sqrt x+1)h<h^2+1=\eta$; and $(\eta+W)Q+\eta+h\le(r^{16}+20r+1)\cdot5r^2+
r^{16}+r^8+1\le7r^{18}$, with $99\cdot7r^{18}\le r^{24}=h^3$.
\end{proof}

\subsubsection*{Where the propositions are proved}

Proposition~\ref{pki:generic} is Theorem~\ref{pkc:generic}; Proposition~\ref{pki:tame} is
Theorem~\ref{pkc:tame}; Proposition~\ref{pki:sparse} is Theorem~\ref{pkd:sparse};
Proposition~\ref{pki:any} is Theorem~\ref{pkd:any}; Proposition~\ref{pki:chart} is
Theorem~\ref{pkj:chart}; Proposition~\ref{pki:data} is Lemma~\ref{pki:data-sizes} together with
Theorem~\ref{pkd:data}. The four cases rest on the pairing lemma and the walk
(\S\S\ref{pkp:sec}, \ref{pkm:sec}) and on the inputs (I1) (Proposition~\ref{pkt:I1}), (I2)
(Theorem~\ref{pke:I2}), (I3) (Proposition~\ref{pkt:I3}), (I4) (Theorem~\ref{pkh:I4}), (I5)
(Theorem~\ref{pkg:I5}), (I6) (Lemma~\ref{pkw:dual}) and (T1)--(T4) (Proposition~\ref{pkw:T1},
Theorems~\ref{pkw:T2}, \ref{pkw:T3}, \ref{pkw:T4}); these in turn rest on the modules $\mathcal N$ of
\S\ref{pkn:sec}, whose properties (Lemma~\ref{pkm:inputs}, Lemmas~\ref{pkp:orders-add},
\ref{pkp:tau}) are proved in \S\S\ref{pkn:sec}--\ref{pkt:sec}, the existence of frames
(Theorem~\ref{pkf:main}) and the ring statements of \S\ref{pkr:sec}. Published results enter only
through \S\S\ref{pkf:sec}, \ref{pkr:sec}, \ref{pkl:sec}, \ref{pkw:sec} (Fontaine's period rings,
Tate's theorem, Fontaine--Wintenberger/Scholze tilting, and the facts on $D_{\mathrm{pcris}}$ listed
there).

\subsection{The uniform crystalline bound and the deduction of Theorem~\ref{pll:exponent}}\label{pki:ss-deduce}

\begin{lemma}[Letter windows]\label{pki:letterwin}
Let the letters read thickness, let $\eta\le G$ and $\delta+W\le G$, and let $U(\gamma)\subseteq
\pkiWin(\eta,-\varepsilon(\gamma),a_\gamma+\delta,W)$ with $a_\gamma=b$ for thick and $G$ for thin
$\gamma$. Put $\eta_\ell=2G$ if $\ell\le Q$ and $\eta_\ell=\eta$ if $\ell>Q$, and $b_z=b+\delta$ if
$z\in T$, $b_z=G+\delta$ otherwise. Then $U_z\subseteq\pkiWin(\eta_\ell,-s(z),b_z,W)$ for every
letter $z$. The same holds with all signs reversed.
\end{lemma}

\begin{proof}
$0$ lies in the low window. Let $u\in U(\gamma)$ with $\gamma(\zeta)=z$. If $\gamma$ is thick, or
$\ell>Q$, then thickness and sign of $\gamma$ are those read from $z$, so $u$ lies in
$[-\eta,\eta]\subseteq[-\eta_\ell,\eta_\ell]$ or in $-s(z)[b_z,b_z+W]$. Otherwise $\gamma$ is thin and
$\ell\le Q$; then $|u|\le\max(\eta,G+\delta+W)\le2G=\eta_\ell$.
\end{proof}

\begin{theorem}[Uniform crystalline bound]\label{pki:crys}
Assume (S1)--(S4) for $(Q,m,h,G,b,\delta,\eta,W)$, that the letters read thickness, that
$U(\gamma)\subseteq\pkiWin(\eta,-\varepsilon(\gamma),a_\gamma+\delta,W)$ and
$U'(\gamma)\subseteq\pkiWin(\eta,\varepsilon(\gamma),a_\gamma+\delta,W)$ for all $\gamma$ (with
$a_\gamma$ as in Lemma~\ref{pki:letterwin}), that $L$ is crystalline with jumps in $U'$, and that
either $L$ has few low jumps for $(h,\mu,\eta)$, or every $U'(\gamma)\subseteq[-\eta,0]$. Let
$\lambda=\lambda_\ell(h,b)$ and $K=\lfloor2\lambda/Q\rfloor$. Then:
\begin{enumerate}
\item[(b)] if $\ell\le Q$, $\pkiC(U)$ has cup exponent $10^4K(\pkilog_\ell h^3+1)$ on $L$;
\item[(c)] if $\ell>Q$ and $\lambda\ge h^3$, it has cup exponent $10^4Km$;
\item[(d)] if $\ell>Q$ and $\lambda<h^3$, then $\plldef_{\pkiC(U)}(L)=0$ unless $\lambda\ge\lfloor
h/5\rfloor$ and $L$ has few low jumps; and always
$\plldef_{\pkiC(U)}(L)\le[E_v:\Q_\ell]\cdot10^5(\pllrk L\cdot Km+\mu\lambda)$.
\end{enumerate}
In all cases $\plldef_{\pkiC(U)}(L)\le[E_v:\Q_\ell]\cdot10^5(\pllrk L\cdot Km+\mu\lambda)$.
\end{theorem}

\begin{proof}
By (S2), (S4): $\eta\le h^3\le G$, $\delta+W<\eta+W\le G$, so Lemma~\ref{pki:letterwin} applies
to $U$ and $U'$. By Lemma~\ref{pki:elem}(i), $2b<\ell(\lambda+1)$; by definition of $K$,
$2\lambda<(K+1)Q$.

\emph{(b) $\ell\le Q$.} First, $K$ is large: from $\lambda+1>2b/\ell\ge2b/Q$ and
$2\lambda<(K+1)Q$ we get $K+3>4b/Q^2$, and $Q^2\le h/99$ by (S1), so
$h^2K>396\,bh-3h^2\ge 2(b+\delta)$ and $K\ge 396\,b/h-3\ge41Gh$ by (S2), (S4) (note $\delta<\eta\le
h^3\le b$). Apply Proposition~\ref{pki:sparse} with low radius $2G$, $b\to b+\delta$, $G\to G+\delta$,
$V=W$, $B=T$, $\varepsilon_z=s(z)$ (and the reversed sign for $U$), $J=\pkilog_\ell h^3$,
$\vartheta=2K$, $M=K$. The windows are those of Lemma~\ref{pki:letterwin}, and $G+\delta+W\le2G$.
$J\ge1$ as $\ell\le Q\le h\le h^3$. The sparseness is part of Definition~\ref{pki:def-reads}. If
$e=1$ and $\ell\ge7$: $10\ell^J\cdot2G\le20h^3G\le b+\delta$. If $s$ is constant:
$8e\le8h\le h^3<\ell^{J+1}$. Depth: $e\cdot2(b+\delta+10G)\le h\cdot4(b+\delta)<h^3\cdot2K\le
\ell^{J+1}\cdot2K$, using $10G\le b$, $e\le h$, $2(b+\delta)<h^2K$ and $h^3<\ell^{J+1}$. Last
condition: by Lemma~\ref{pki:elem}(iii)--(iv), $20(J+1+d_v)\cdot2G+\pkilog_\ell(2(b+\delta+10G))\le
40hG+m\le41hG\le K$. The exponent is $10^3(2K+(J+1)K)\le10^4K(J+1)$.

\emph{From now on $\ell>Q$}; by Definition~\ref{pki:def-reads}, $e=1$ and $\ell\ge7$, and the low
radius is $\eta$.

\emph{(c) $\lambda\ge h^3$.} Then $K+1>2h^3/Q$, so $K\ge h^3/Q\ge99h^2\ge2$ (as $Q\le h/99$).
Apply Proposition~\ref{pki:any} with $b\to b+\delta$, $G\to G+\delta$, $V=W$, $M=W+4\eta$, and $K$.
Conditions: $W+4\eta\le4(\eta+W)\le4h^3/(99Q)\le K$ by (S4). $0<G+\delta\le b+\delta$. Next,
$b<\ell(\lambda+1)/2<\ell((K+1)Q/2+1)/2\le\ell^2(K+1)/4+\ell/2$ since $Q<\ell$; and
$\ell^2K/4\ge\ell/2$, $\ell^2K/4\ge 12K\ge12h^3/Q>\delta$ (as $\delta<\eta\le h^3/(99Q)$); so
$b+\delta<\ell^2K$. Finally $K\ell\ge(2\lambda/Q-1)\ell\ge\lambda\ell/Q\ge\ell(\lambda+1)/(2Q)>b/Q$
(using $\lambda\ge h^3\ge Q$), and $b/Q\ge99Gh^3/Q>2(G+\delta)$. The exponent is
$10^4(K+\pkilog_\ell(2(b+\delta+W+4\eta)))\le10^4(K+m-1)\le10^4Km$ by Lemma~\ref{pki:elem}(iii)
(the argument is $\le4b$) and $K\ge1$.

\emph{(d) $\lambda<h^3$.} By Lemma~\ref{pki:elem}(ii), $\ell$ divides no integer in
$(b-30h^3,b+30h^3]$. Put $a=2(\lambda+1)\le2h^3$ and $n=6(a+\eta+W)\le12h^3+6h^3/99<13h^3$.
\emph{Gap.} $\ell>2b/(\lambda+1)\ge2b/h^3\ge198G$. A nonzero multiple $y$ of $\ell$ has $|y|\ge\ell>
G+\delta+n$ (as $n<13h^3\le G/7$ and $\delta<G/99$), so it is not within $n$ of $0$ or
$\pm(G+\delta)$; if it were within $n$ of $\pm(b+\delta)$, then $|y|\in(b-30h^3,b+30h^3]$ since
$\delta+n<30h^3$, a contradiction. So $\mathrm{NM}(\ell,n,b_z)$ for every letter. \emph{Heights.}
$(\ell-3)a=2\ell(\lambda+1)-6(\lambda+1)>4b-6h^3\ge b_z+W+\eta$, and $6(a+\eta+W)<13h^3<G\le b_z$.
So Proposition~\ref{pki:generic} applies with this $a$, radius $\eta$, $b_z$ and
$\varepsilon_z=s(z)$.

If $\lambda\ge\lfloor h/5\rfloor$ and $L$ has few low jumps, (ii) gives $\plldef\le
11f(\mu(2\lambda+3)+\pllrk L(W+1))$. Here $\lambda\ge1$, so $2\lambda+3\le6\lambda$; and
$(K+1)Q>2\lambda\ge h/3$ with $(W+1)Q\le h/99$ gives $K\ge W+1$; and $m\ge1$. Hence
$\plldef\le f\cdot10^5(\pllrk L\cdot Km+\mu\lambda)$, and $f=[E_v:\Q_\ell]$ as $e=1$.

Otherwise either every $U'(\gamma)\subseteq[-\eta,0]$, or $L$ has few low jumps and
$\lambda<\lfloor h/5\rfloor$; in the latter case every negative $u$ in $U(\gamma)\cup U'(\gamma)$
has $u\le-h/2<-2(\lambda+1)=-a$. In both cases (i) gives $\plldef=0$.

\emph{The final bound.} In cases (b), (c), Lemma~\ref{pll:exp} gives $\plldef\le2k'[E_v:\Q_\ell]
\pllrk L$ for the exponent $k'$; in case (b) $2\cdot10^4K(\pkilog_\ell h^3+1)\le10^5Km$ since
$\pkilog_\ell h^3+1\le m$ (Lemma~\ref{pki:elem}(iv)), and in case (c) $2\cdot10^4Km\le10^5Km$.
\end{proof}

\begin{proof}[Proof of Theorem~\ref{pll:exponent}, given Propositions~\ref{pki:generic}--\ref{pki:data}]
Fix $\ell$, $v$, $c$ and put $L=L_c^*$, $\rho=\pllrk L_c$, $x=x_c$. By Proposition~\ref{pki:data}
(with its $j$), the size conditions and windows of Theorem~\ref{pki:crys} hold with
$\delta=\delta_j$; by Proposition~\ref{pki:chart} the letters read thickness; and $L$ has few low
jumps with $\mu=\lfloor\rho m/k\rfloor$ ($c=\mathrm{main}$, $x=1$) or all $U'(\gamma)\subseteq
[-\eta,0]$ ($c$ a character). If $\ell\notin\{p_0,p_1\}$, $\mathcal C_v=\pkiC(U)$ (the summands are crystalline above $\ell\nmid M=p_0p_1$ by (M2) of Definition~\ref{pg:def-summand} and (G2), the twist entering only through the cocycle conditions at $Q\cup\aubar Q$) and Theorem~\ref{pki:crys} gives (b), (c), (d) with the exponents printed in
Theorem~\ref{pll:exponent}; in (d), the defect of a character vanishes, so the bound holds with
$\mu=\lfloor\rho m/(kx)\rfloor$ for both kinds of summand. If $\ell=p_i$, $E_v=\Q_\ell$ ($p_i$ splits
completely in $E$) and Proposition~\ref{pki:tame} with $N=\max(r+t,2r)r$, $W=4t$, $\eta=h^2+1$ and
$b\to a+\delta_j$, where $a=a_\gamma$ is the common value (Proposition~\ref{pki:data}), gives the exponent $3r^2(N+1)^2+150(h^2+4t+2)$ of (a), using the last clause
of Proposition~\ref{pki:data} (a cup exponent of $\pkiC(U,K)$ on $L$ is one of its subclass
$\mathcal C_v$). The passage to \eqref{pll:eq-place} is Remark~\ref{pll:status}.
\end{proof}

\begin{remark}[Comparison with Proposition~\ref{pri:loss}]\label{pki:align}
Proposition~\ref{pri:loss}(b) bounds $\sum_L\sum_v\delta_v(L)/e_\auO$, the quantity of
Proposition~\ref{pll:main} with lengths normalised so that $\pglen\auO/\ell=1$; its part (a) is
Proposition~\ref{pll:dual} applied to the module of Ribet cocycles, whose elements come from
$\mathcal C_v$ by Lemma~\ref{pll:ribet} (applied with $\Gamma=D_v$ inside $\pgG$ to the lattices
$u_x$, whose restrictions to $D_v$ lie in $\mathcal C_v$ by Proposition~\ref{pri:loss}(b)).
\end{remark}


%% ===== proofs/C-kisin-pairing.tex =====
% proofs/C-kisin-pairing.tex
% The engine for Theorem pll:exponent: Kummer charts, letters, safe moves and the pairing lemma.
% Labels start with pkp:. Conventions and macros as in proofs/C-kisin.tex.
\providecommand{\auO}{\mathcal O}
\providecommand{\auQlb}{\overline{\Q}_\ell}
\providecommand{\panot}[1]{\par\smallskip{\raggedright\noindent\textbf{[Not proved here:} #1\textbf{]}\par}\smallskip}
% macros of this file (prefix \pkp)
\providecommand{\pkpA}{A_{\mathrm{inf}}}
\providecommand{\pkpW}{W(C^\flat)}
\providecommand{\pkpm}{W(\mathfrak m^\flat)}
\providecommand{\pkpM}{M_m}
\providecommand{\pkpsm}{\mathrm{Sm}_m}
\providecommand{\pkpAmb}{\mathbb A}

\section{Kummer charts and the pairing lemma}\label{pkp:sec}

This section supplies the engine on which Propositions~\ref{pki:generic}--\ref{pki:tame} rest:
the Kummer charts of a local Galois group, their splitting into letters, the moves of a pair of
lattice points and the \emph{pairing lemma}, which turns a short chain of safe moves into a bound
for the exponent of a cup product. \S\ref{pkp:ss-arch} states every definition and result cited by
the four cases of \S\ref{pki:ss-arch}. The proofs follow in
\S\ref{pkp:ss-AS}--\S\ref{pkp:ss-walk}.

Lemmas~\ref{pkp:orders-add} and \ref{pkp:tau} are proved in \S\ref{pkt:sec} (Lemma~\ref{pkt:add},
Proposition~\ref{pkt:tau-i}, Lemmas~\ref{pkt:tau-iv}, \ref{pkt:tau-v}), Lemma~\ref{pkp:tau}(iii) is
Lemma~\ref{pkm:tau-iii}, and Proposition~\ref{pki:chart} is Theorem~\ref{pkj:chart}.

\emph{Conventions.} As in \S\ref{pki:sec}: $\ell$ a prime, $\Gamma=D_v$ for a place $v\mid\ell$ of
$E$, $j\colon\Gamma\to G_{\Q_\ell}=\operatorname{Gal}(\auQlb/\Q_\ell)$ the inclusion (an open
embedding), $\auO$ a large coefficient ring (finite free over $\Z_\ell$). $C=\C_\ell$, $O_C$ its
ring of integers, $O^\flat=\varprojlim_{x\mapsto x^\ell}O_C/\ell$ the tilt, $C^\flat$ its fraction
field (algebraically closed, of characteristic $\ell$), $\mathfrak m^\flat$ the maximal ideal of
$O^\flat$. Witt vectors are $\ell$-typical. $\pkpA=W(O^\flat)\subset\pkpW$; $\varphi$ is the Witt
vector Frobenius; $\theta\colon\pkpA\to O_C$ is Fontaine's map; $\varepsilon=(1,\zeta_\ell,
\zeta_{\ell^2},\dots)\in O^\flat$, $\tau=[\varepsilon^{1/\ell}]-1\in\pkpA$, $t=\log[\varepsilon]\in
B_{\mathrm{dR}}^+$. We write $\pkpm\subset\pkpA$ for the ideal of Witt vectors all of whose
coordinates lie in $\mathfrak m^\flat$ (the kernel of $\pkpA\to W(\overline{\mathbb F}_\ell)$). The
group $\Z\times\Gamma$ acts on $\pkpW$: the generator $\varphi=(1,1)$ by Frobenius, $g\in\Gamma$
through $j$; the two actions commute.

\subsection{Architecture}\label{pkp:ss-arch}

\subsubsection*{Kummer charts and letters}

\begin{definition}[Kummer chart]\label{pkp:def-chart}
A \emph{Kummer chart} of $\Gamma$ consists of the data below; $K'=\auQlb^{\,j(\Gamma')}$ denotes the
fixed field of its subgroup $\Gamma'$, $e$ and $f$ its ramification index and residue degree over
$\Q_\ell$, and $\Lambda=W(\mathbb F_{\ell^f})\subset\pkpA$. The data are: an
element $\varpi^\flat\in\mathfrak m^\flat$ whose untilt $\pi=\theta([\varpi^\flat])$ is a uniformiser of
$K'$, the coordinate $z=[\varpi^\flat]\in\pkpA$, the ring $\mathfrak S=\Lambda[[z]]\subset\pkpA$ (the
closure of $\Lambda[z]$; $\Lambda[[X]]\to\mathfrak S$, $X\mapsto z$, is an isomorphism), an element
$E=z^eA+\ell\in\mathfrak S$ with $A\in\mathfrak S^\times$ generating $\ker\theta$ (so
$E(X)$ is an Eisenstein polynomial times a unit), and an open subgroup $\Gamma'\subseteq\Gamma$ of index prime to $\ell$
such that: $j(\Gamma')$ fixes $\Lambda$ and $\theta(z)$; every element of $G_{\Q_\ell}$ fixing $\Lambda$
pointwise and $\theta(z)$ lies in $j(\Gamma')$; and for every $l\ge0$ some $g\in\Gamma'$ has
$gz=[\varepsilon]^lz$. We write $p_a=\varphi^a(E)$ ($a\ge0$); $\theta\circ\varphi^{-a}$ has kernel
$p_a\pkpA$. The chart is \emph{unramified} if $e=1$ and $E=z+\ell$. We put $d_v=0$ if $e=1$ and
$d_v=v_\ell(e)+1$ otherwise.
\end{definition}

\begin{lemma}[Existence]\label{pkp:chart-exists}
Kummer charts of $\Gamma$ exist, with $\Gamma'=\Gamma$; if $j(\Gamma)$ contains the inertia group,
the chart can be taken unramified.
\end{lemma}

\begin{definition}[Letters]\label{pkp:def-letters}
The \emph{letters} are the elements $i\in\mathbb F_{\ell^f}^\times$; $[i]\in\Lambda$ is the
Teichm\"uller lift. Fix $\zeta\in\auO$ with $\zeta^{\ell^f-1}=1$ (Lemma~\ref{pkw:letters} shows that the $e_i$ below are then orthogonal idempotents with sum $1$; in the cases one takes $\zeta$ of exact order $\ell^f-1$, or $\zeta=1$ for a tame chart). For a finite free $\auO$-module $L$
put $\pkpAmb(L)=\pkpW\otimes_{\Z_\ell}L$; its \emph{letter part} at $i$ is
$\pkpAmb(L)_i=\{x:(1\otimes\zeta)x=[i]x\}$, and $e_i=\prod_{i'\ne i}\frac{(1\otimes\zeta)-[i']}{[i]-[i']}$
is the projector onto it (the differences $[i]-[i']$ are units of $\Lambda$, being Teichm\"uller lifts
of nonzero residues); $\sum_ie_i=1$ and $\pkpAmb(L)=\bigoplus_i\pkpAmb(L)_i$. Since
$\varphi[i]=[i^\ell]$ and $\Gamma'$ fixes $\Lambda$, $\varphi$ maps the part at $i$ to the part at
$i^\ell$ and $\Gamma'$ preserves each part. For $a\ge0$, $i^{(-a)}$ (``$a$ steps before $i$'') is the
letter with $(i^{(-a)})^{\ell^a}=i$. Under an embedding $\gamma\colon\auO\to C$ the letter of
$\gamma$ is the residue of $\gamma(\zeta)$; this matches the letters of \S\ref{pki:ss-arch}.
\end{definition}

\begin{definition}[Lattices with jump sets; orders]\label{pkp:def-orders}
A \emph{lattice with jump sets} is a finite free $\auO$-module $L$ with a continuous $\auO$-linear
action of $\Gamma$, crystalline on $\Gamma'$, together with finite sets $U(i)\subset\Z$ (one for each
letter) containing the Hodge--Tate jumps of $L$ at all embeddings of letter $i$. Let
$\pkpA\mathcal N=\pkpA\mathcal N(L)$ and $\mathcal N=\mathcal N(L)$ be the modules of
Definition~\ref{pkn:def-N} (they play the role of the Breuil--Kisin module of $L$ with the $p_a$
inverted). A \emph{form of type
$u$ at $i$} is an $\pkpA$-linear map $\mu\colon\pkpAmb(L)\to B_{\mathrm{dR}}$, $\Gamma'$-equivariant on
$1\otimes L$, with $\mu(e_i(1\otimes L))\subseteq t^{-u}B^+_{\mathrm{dR}}$. For a function
$o\colon\Z_{\ge0}\times\Z\to\Z$ (``orders''), $x$ \emph{has orders $o$ at $i$} if $x\in\pkpA\mathcal
N\cap\pkpAmb(L)_i$ and $\mu(\varphi^{-a}x)\in t^{o(a,u)}B^+_{\mathrm{dR}}$ for all $a\ge0$, all $u\in
U(i^{(-a)})$ and all forms $\mu$ of type $u$ at $i^{(-a)}$. These $x$ form an $\pkpA$-module
$L_i(o)$. ($a$ indexes the prime $p_a$, $u$ the jump.) The \emph{standard orders} are
$\mathrm{std}_l(a,u)=0$ for $a<l$ and $-u$ for $a\ge l$; we also use
\[
\rho_t(a,u)=t\cdot[a=0]-u,\qquad \mathrm{sh}_{l,n}(o)(a,u)=\begin{cases}n&a<l\\ o(a-l,u)&a\ge
l,\end{cases}
\]
(``raised at $0$'' and ``shifted by $l$''). An element $x\in\mathcal N$ is \emph{deep for $o'$} if
$gx-x\in L_i(o')$ for all $g\in\Gamma'$.
\end{definition}

\subsubsection*{Sides, pairings, points}

\begin{definition}[Side]\label{pkp:def-side}
A \emph{side} $S$ is a lattice $Y$ with jump sets $U(i)\ni0$, a finite free $\auO$-module $C_S$ with
trivial action, and an $\auO$-linear map $\alpha\colon C_S\to Z^1_f(\Gamma,Y)$ (crystalline
cocycles, Definition~\ref{pll:def}). Its \emph{extension} is $X_S=Y\oplus C_S$ with
$g(y,c)=(gy+\alpha(c)(g),c)$, a lattice with the same jump sets (its jumps are those of $Y$ and $0$).
Let $\mathrm{pr}\colon\pkpAmb(X_S)\to\pkpAmb(C_S)$ and $\mathrm{fst}\colon\pkpAmb(X_S)\to\pkpAmb(Y)$
be the two projections. For $c\in C_S$ and a letter $i$, the \emph{constant of $c$ at $i$} is
$e_i(1\otimes c)$; a vector $x$ \emph{lies over $\ell^ec$ at $i$} if $x\in\pkpAmb(X_S)_i$ and
$\mathrm{pr}(x)=\ell^ee_i(1\otimes c)$.
\end{definition}

\begin{definition}[Pairing modulo $\ell^m$]\label{pkp:def-pairing}
For sides $S,S'$ with lattices $Y,Y'$, a \emph{pairing modulo $\ell^m$} is a $\Z_\ell$-bilinear,
$\Gamma$-equivariant map $P\colon Y\times Y'\to\Z/\ell^m(1)$ which is the reduction (through a
$\Z_\ell$-linear $\auO\to\Z_\ell$) of an $\auO$-bilinear $\psi\colon Y\times Y'\to\auO(1)$ whose
failure of equivariance $g\psi-\psi=\ell^m c_P(g)$ is a crystalline cocycle $c_P$ of
$\operatorname{Hom}(Y\otimes Y',\auO(1))$. Put $\pkpM=\pkpW\otimes_{\Z_\ell}\Z/\ell^m(1)$, with
$\Z\times\Gamma$ acting diagonally ($\varphi$ trivially on $\Z/\ell^m(1)$), and let $P$ also denote
the $\pkpW$-bilinear extension $\pkpAmb(Y)\times\pkpAmb(Y')\to\pkpM$; it commutes with $\Z\times\Gamma$.
The \emph{small part} is $\pkpsm=\pkpm\cdot(1\otimes\epsilon_m)\subset\pkpM$, where $\epsilon_m$ is a generator of $\Z/\ell^m(1)$; it does not depend on the choice of $\epsilon_m$, since the generators differ by units of $\Z/\ell^m$. The \emph{Tate-dual side} of
$P$ has lattice $\operatorname{Hom}(Y\otimes Y',\auO(1))$, jump sets $\{0\}\cup\{-u-u'-1:u\in U(i),
u'\in U'(i)\}$, $C=\auO$ and $\alpha(1)=c_P$.
\end{definition}

\begin{definition}[Safe pairs, moves, paths]\label{pkp:def-safe}
Let $\kappa\ge0$. A pair $(w,y)\in\pkpAmb(X_S)\times\pkpAmb(X_{S'})$ is \emph{$\kappa$-safe} if
$\ell^\kappa P(\mathrm{fst}\,w,\mathrm{fst}\,y)\in\pkpsm$. The \emph{move of the first side by $y$
with $x'$ at rest} is safe if $(y,gx'-x')$ is safe for all $g\in\Gamma'$; the \emph{move of the second
side by $y'$ with $x$ at rest} is safe if $(gx-x,gy')$ is safe for all $g\in\Gamma'$. A \emph{safe
path} from $(y,y')$ to $(x,x')$ through $(w,w')$ consists of four safe moves:
\begin{center}
first side by $y-w$ with $y'$ at rest;\quad second side by $y'-w'$ with $w$ at rest;\\
first side by $w-x$ with $w'$ at rest;\quad second side by $w'-x'$ with $x$ at rest.
\end{center}
\end{definition}

\begin{definition}[Scaled cup coboundary]\label{pkp:def-scc}
For $c\in C_S$, $c'\in C_{S'}$ and $N\ge0$, we say $\mathrm{SCC}_N(c,c')$ if the continuous 2-cocycle
$(g,h)\mapsto\zeta_{\ell^m}^{\,P(\alpha(c)(g),\,g\,\ell^N\alpha'(c')(h))}$ of $\Gamma$ with values in
the discrete module $\auQlb^{\times}$ is the coboundary of a continuous 1-cochain. (Here
$\Z/\ell^m(1)\cong\mu_{\ell^m}$, $x\mapsto\zeta_{\ell^m}^x$.) In other words, $\ell^N$ kills the cup
product of the two classes in $H^2(\Gamma,\auQlb^\times)$ read through $P$.
\end{definition}

\begin{theorem}[Pairing lemma]\label{pkp:pairing}
Let $c\in C_S$, $c'\in C_{S'}$, $e_1,e_2,\kappa\ge0$, and let $x,w\in\pkpAmb(X_S)$ satisfy
$\mathrm{pr}(x)=\mathrm{pr}(w)=\ell^{e_1}(1\otimes c)$ and $x',w'\in\pkpAmb(X_{S'})$ satisfy
$\mathrm{pr}(x')=\mathrm{pr}(w')=\ell^{e_2}(1\otimes c')$. Assume:
\begin{enumerate}
\item[(a)] for every $h\in\Gamma'$ the move of the second side by $hx'-x'$ with $x$ at rest is
$\kappa$-safe;
\item[(b)] there is a $\kappa$-safe path from $(\varphi x,\varphi x')$ through $(w,w')$ to $(x,x')$.
\end{enumerate}
Then $\mathrm{SCC}_N(c,c')$ for every $N\ge e_1+e_2+\kappa$.
\end{theorem}

\begin{corollary}[Letter by letter]\label{pkp:letters}
Suppose that for every letter $i$ vectors $y_i,w_i,x_i$ lie over $\ell^ec$ at $i$ and
$y'_i,w'_i,x'_i$ lie over $\ell^{e'}c'$ at $i$, that $\varphi\sum_ix_i=\sum_iy_i$ and
$\varphi\sum_ix'_i=\sum_iy'_i$, that for all $i$ and $h\in\Gamma'$ the move of the second side by
$hx'_i-x'_i$ with $x_i$ at rest is $\kappa$-safe, and that for every $i$ there is a $\kappa$-safe path
from $(y_i,y'_i)$ through $(w_i,w'_i)$ to $(x_i,x'_i)$. Then $\mathrm{SCC}_N(c,c')$ for $N\ge
e+e'+\kappa$.
\end{corollary}

\subsubsection*{Points, cokernel killers, the rule and the walk}

\begin{definition}[Points; killing the cokernel]\label{pkp:def-points}
Let $S$ be a side, $c\in C_S$, $i$ a letter, $e\ge0$, $o,o'$ orders. The \emph{points over $\ell^ec$ at
$i$ with orders $o$ and depth $o'$} are the $x\in\mathcal N(X_S)$ that lie over $\ell^ec$ at $i$, have
orders $o$ at $i$ and are deep for $o'$. Put, for $t,\eta,l\ge0$,
\[
\mathrm{dep}_{t,\eta,l}(a,u)=-u+t[a=0]+\eta[a<l],\qquad
\mathrm{stop}_{\eta,l}(a,u)=\max\bigl(\mathrm{std}_l(a,u),\,-u-\eta[a<l]\bigr),
\]
\[
\mathrm{rule}_{l,t}(a,s)=\begin{cases}1&a<l,\\ t-s&a=l,\\ -s&a>l.\end{cases}
\]
\emph{Standard points} (for $t,\eta,l$) are points with orders $\rho_{-\eta}$ and depth
$\mathrm{dep}_{t,\eta,l}$; \emph{stops} (for $\eta,l$) are points with orders and depth
$\mathrm{stop}_{\eta,l}$. ``$x\in L_i(o)$ lies over $0$'' means $\mathrm{pr}(x)=0$; write
$L^0_i(o)$. We say that \emph{$\ell^\kappa E^t$ kills the cokernel at $i$} for the side $S$ if for
every $c\in C_S$ there is $n\in(X_S)_i(\mathrm{std}_0)$ (that is, $n$ in the letter-$i$ part of
$\pkpA\mathcal N(X_S)$ with orders $\mathrm{std}_0$; $n$ need not lie in $\mathcal N(X_S)$, but it can
be moved into it, Lemma~\ref{pkt:lift}) and $\mathrm{pr}(n)=\ell^\kappa E^t
e_i(1\otimes c)$.
\end{definition}

\begin{proposition}[The rule]\label{pkp:rule}
Let $P$ pair $S$, $S'$, let $l,t,\kappa\ge0$ and suppose $\ell^\kappa E^t$ kills the cokernel of
the Tate-dual side at the letter $i^{(-l)}$. Let $y\in L^0_i(o)$ (for $X_S$) and let $x'$ be a point
of $S'$ over $\ell^{e'}c'$ at $i$ with depth $d'$ (any orders). If
$\mathrm{rule}_{l,t}(a,u+u')\le o(a,u)+d'(a,u')$ for all $a\ge0$, $u\in U(i^{(-a)})$, $u'\in
U'(i^{(-a)})$, then the move of the first side by $y$ with $x'$ at rest is $\kappa$-safe. Symmetrically,
for $y'\in L^0_i(o')$ (for $X_{S'}$) and a point $x$ of $S$ over $\ell^ec$ at $i$ with depth $d$, if
$\mathrm{rule}_{l,t}(a,u+u')\le d(a,u)+o'(a,u')$ for all such $a,u,u'$, the move of the second side
by $y'$ with $x$ at rest is $\kappa$-safe.
\end{proposition}

\begin{proposition}[The walk at one cut]\label{pkp:cut}
Let $0<l$, $0<l_0$, let $\ell^\kappa E^t$ kill the cokernel of the Tate-dual side at $i$ and
$\ell^\kappa E^{t_0}$ at $i^{(-l)}$. Let $x,x'$ be standard points (for $t,\eta,l$) of $S,S'$ over
$\ell^ec,\ell^{e'}c'$ at $i$, and $x_0,x'_0$ standard points (for $t_0,\eta,l_0$) over the same at
$i^{(-l)}$. Suppose there is a stop $q$ (for $\eta,l$) of $S$ over $\ell^ec$ at $i$ such that the move
of the second side by $\varphi^lx'_0-x'$ with $q$ at rest is $\kappa$-safe, or a stop $q'$ of $S'$
over $\ell^{e'}c'$ at $i$ such that the move of the first side by $\varphi^lx_0-x$ with $q'$ at rest
is $\kappa$-safe. Then there are $w,w'$ lying over $\ell^ec,\ell^{e'}c'$ at $i$ and a $\kappa$-safe
path from $(\varphi^lx_0,\varphi^lx'_0)$ through $(w,w')$ to $(x,x')$. Moreover, for every $d\ge0$,
every $g\in\Gamma'$ and standard points $x,x'$ (for $t,\eta,l$) at $i^{(-d)}$, with $\ell^\kappa E^t$
killing the cokernel at $i^{(-d)}$, the move of the second side by
$g\varphi^dx'-\varphi^dx'$ with $\varphi^dx$ at rest is $\kappa$-safe.
\end{proposition}

The two inputs from the theory of Breuil--Kisin modules used in the proof of the rule are:

\begin{lemma}[Orders add]\label{pkp:orders-add}
Let $L,L',L''$ be lattices with jump sets $U,U',U''$ and $b\colon L\times L'\to L''$ a
$\Gamma'$-equivariant $\auO$-bilinear map, with $\pkpW$-bilinear extension $\tilde b$. If $w\in L_i(o)$,
$y\in L'_i(o')$ and $o''(a,v)\le o(a,u)+o'(a,u')$ whenever $u+u'\le v$ (for all $a\ge0$,
$u\in U(i^{(-a)})$, $u'\in U'(i^{(-a)})$, $v\in U''(i^{(-a)})$), then $\tilde b(w,y)\in L''_i(o'')$.
In particular, for $L''=L\otimes_\auO L'$ with jump sets $U(i)+U'(i)$ and $b=\otimes$,
$w\otimes y\in(L\otimes L')_i(o'')$.
\end{lemma}

\begin{lemma}[Integrality in $\auO(1)$]\label{pkp:tau}
(i) For the lattice $\auO(1)$ (jump set $\{-1\}$), an element $\xi$ with orders $\mathrm{std}_0$ at $i$
satisfies $\tau\xi\in\pkpA\otimes_{\Z_\ell}\auO(1)$. (ii) Moreover $z\in\pkpm$, and for $\xi\in
\pkpA\otimes\auO(1)$ the image of $z\xi$ in $\pkpM$ lies in $\pkpsm$. (iii) If $\ell^\kappa E^t$ kills the
cokernel at $i$ for a side, then for every $c$ there is $\nu\in(X_S)_i(\rho_{-t})$ with
$\mathrm{pr}(\nu)=\ell^\kappa e_i(1\otimes c)$ (Lemma~\ref{pkm:tau-iii}; by Proposition~\ref{pkt:I1}
$\nu$ can be taken to be a point, with depth $\rho_{-t}$). For every lattice with jump sets: (iv) if $x\in\mathcal N$ lies at $i$,
is deep for $o'$ and $\mathrm{pr}(x)$ is $\Gamma'$-invariant (for an extension), then for $g\in\Gamma'$
the vector $(z\tau)^{-1}(gx-x)$ (note $z\tau\in\pkpW^\times$) has orders $o'$ at $i$ and lies over
$0$; (v) if $x\in L_{i^{(-l)}}(o)$ then $\varphi^lx\in L_i(\mathrm{sh}_{l,0}(o))$; if moreover $x$ is a point at
$i^{(-l)}$ with orders $o$ and depth $o'$, then $\varphi^lx$ is a point at $i$ (over the same constant)
with orders $\mathrm{sh}_{l,0}(o)$ and depth $\mathrm{sh}_{l,1}(o')$.
\end{lemma}

The last input concerns the data and is stated in \S\ref{pki:ss-arch} as
Proposition~\ref{pki:chart}.

\begin{proof}[Proof of Lemma~\ref{pkp:chart-exists}]
Let $\pi$ be a uniformiser of $K$ (take $\pi=-\ell$ if $K/\Q_\ell$ is unramified, which is the
case when $j(\Gamma)$ contains inertia), choose $\pi_n\in\auQlb$ with $\pi_0=\pi$, $\pi_{n+1}^\ell=\pi_n$,
and let $\varpi^\flat=(\pi_n\bmod\ell)_n\in\mathfrak m^\flat$; its untilt is $\pi$. Put
$\Gamma'=\Gamma$, so $j(\Gamma')=G_K$. $G_K$ fixes $\pi=\theta(z)$, and fixes $\Lambda$ because it acts on
Witt coordinates through $O^\flat$ and fixes the Teichm\"uller lifts of $\mathbb F_{\ell^f}$ (the
residue field of $K$). Conversely an element fixing $\Lambda$ and $\pi$ fixes $K_0=\Lambda[1/\ell]$ and
$\pi$, hence $K=K_0(\pi)$ ($K/K_0$ is totally ramified, generated by $\pi$). Write $\pi^e=-\ell v$ with
$v\in O_K^\times$; $\theta$ maps $\mathfrak S$ onto $O_K$ ($\Lambda\to W(\mathbb F_{\ell^f})$, $z\mapsto\pi$),
so choose $A\in\mathfrak S$ with $\theta(A)=v^{-1}$; $A$ is a unit of the local ring $\mathfrak S$, and
$\theta(z^eA+\ell)=0$. That $E=z^eA+\ell$ generates $\ker\theta$ is Corollary~\ref{pkn:E-gen}.
If $e=1$ and $\pi=-\ell$, $v=1$ and we may take $A=1$, $E=z+\ell$.

The last property. Fix $l\ge0$ and $\zeta_{\ell^n}$ with $\varepsilon=(\zeta_{\ell^n}\bmod\ell)_n$.
Let $X_n=\{g\in G_K:g\pi_n=\zeta_{\ell^n}^l\pi_n\}$. Each $X_n$ is closed (a stabiliser condition on
one algebraic element). $X_{n+1}\subseteq X_n$: raise $g\pi_{n+1}=\zeta_{\ell^{n+1}}^l\pi_{n+1}$ to the
$\ell$-th power. $X_n\ne\emptyset$: $X^{\ell^n}-\pi$ is Eisenstein over $O_K$, hence irreducible over
$K$, so $G_K$ is transitive on its roots, among which are $\pi_n$ and $\zeta_{\ell^n}^l\pi_n$. As
$G_K$ is compact, $\bigcap_nX_n$ contains some $g$. Then $g$ sends every coordinate of $\varpi^\flat$ to
the corresponding coordinate of $\varepsilon^l\varpi^\flat$, so $g\varpi^\flat=\varepsilon^l\varpi^\flat$
and $gz=[\varepsilon]^lz$ (Teichm\"uller lifts are multiplicative and equivariant).
\end{proof}

Lemma~\ref{pkp:tau}(ii) is Lemma~\ref{pkp:AS}(vii).

\subsection{Artin--Schreier facts in $\pkpM$}\label{pkp:ss-AS}

\begin{lemma}\label{pkp:AS}
Let $m\ge0$ and $\pkpM$, $\pkpsm$ be as in Definition~\ref{pkp:def-pairing}. Fix a generator $\epsilon_m$ of $\Z/\ell^m(1)$ and write $\iota(x)=x\otimes\epsilon_m$
for $x\in\pkpW$; $\iota$ is onto and $\iota(x)=\iota(y)$ iff $x\equiv y\pmod{\ell^m\pkpW}$.
\begin{enumerate}
\item[(i)] $\ell^m\pkpM=0$; $\pkpsm$ is stable under $\Gamma$.
\item[(ii)] $\varphi-1$ is onto $\pkpM$.
\item[(iii)] If $(\varphi-1)d\in\pkpsm$, then $d-1\otimes k\in\pkpsm$ for some $k\in\Z/\ell^m(1)$.
\item[(iv)] If $1\otimes k\in\pkpsm$ with $k\in\Z/\ell^m(1)$, then $k=0$.
\item[(v)] Let $\rho$ be a $\Z_\ell[\Gamma]$-module with $gy-y\in\ell^m\rho$ for $g$ near $1$ (depending
on $y$), and $L\colon\pkpW\otimes\rho\to\pkpM$ a $\pkpW$-linear map. For each $v$, $L(gv-v)\in\pkpsm$
for all $g$ in a neighbourhood of $1$.
\item[(vi)] (Used in (v).) For $x,d\in\pkpW$, the first $m$ Witt coordinates of $(gx-x)d$ lie in
$\mathfrak m^\flat$ for $g$ near $1$.
\item[(vii)] $z\in\pkpm$, and the image in $\pkpM$ of $z\cdot(\pkpA\otimes_{\Z_\ell}\auO(1))$ under any
$\Z_\ell$-linear reduction $\auO(1)\to\Z/\ell^m(1)$ lies in $\pkpsm$.
\end{enumerate}
\end{lemma}

\begin{proof}
Since $\Z/\ell^m(1)$ is free of rank one over $\Z/\ell^m$, $\pkpM\cong\pkpW/\ell^m$ as a
$\pkpW$-module via $\iota$, which gives the first assertions. $\varphi$ acts by $\iota(x)\mapsto
\iota(\varphi x)$ and $g\in\Gamma$ by $\iota(x)\mapsto\chi(g)\iota(gx)$. Also $\iota(x)\in\pkpsm$ iff
$x\equiv a\pmod{\ell^m}$ for some $a\in\pkpm$.

(i) The first claim is clear. $g$ acts on Witt coordinates through the automorphism $g$ of $O^\flat$,
which preserves $\mathfrak m^\flat$; so $g\pkpm=\pkpm$ and $g(a\otimes\epsilon_m)=\chi(g)ga\otimes\epsilon_m$.

(ii) Since $C^\flat$ is algebraically closed, every $X^\ell-X-c$ ($c\in C^\flat$) has a root, and
by induction on the Witt length (the $n$-th coordinate of $\varphi y-y$ is $y_n^\ell-y_n$ plus a
polynomial in $y_0,\dots,y_{n-1}$) $\varphi-1$ is onto $\pkpW$; apply $\iota$.

(iii) Write $d=\iota(x)$; by hypothesis $\varphi x-x=a+\ell^mw$ with $a\in\pkpm$, $w\in\pkpW$. As
$O^\flat$ is integrally closed in the algebraically closed $C^\flat$, the roots of $X^\ell-X-c$ for
$c\in O^\flat$ lie in $O^\flat$, so as in (ii) $a=\varphi b-b$ with $b\in\pkpA$, and $w=\varphi
y-y$ with $y\in\pkpW$. The image $\bar b$ of $b$ in $W(\overline{\mathbb F}_\ell)$ satisfies
$\varphi\bar b=\bar b$ (as $\bar a=0$), so $\bar b=c\in\Z_\ell$ and $b-c\in\pkpm$. Then
$x-\ell^my-(b-c)$ is fixed by $\varphi$, hence lies in $W(\mathbb F_\ell)=\Z_\ell$ (the
$\varphi$-fixed points of $W(C^\flat)$ are $W((C^\flat)^{x^\ell=x})$); call it $c'$. Then
$d-1\otimes c'=\iota(\ell^my+(b-c))\in\pkpsm$.

(iv) Let $k$ be the class of $c\in\Z_\ell$ and suppose $c\equiv a\pmod{\ell^m\pkpW}$ with
$a\in\pkpm$. Congruence modulo $\ell^m$ in $W(C^\flat)$ means equality of the first $m$ Witt
coordinates. Those of $c$ lie in $\mathbb F_\ell$, those of $a$ in $\mathfrak m^\flat$; so they
vanish, i.e.\ $\ell^m\mid c$ in $W(\mathbb F_\ell)=\Z_\ell$, $k=0$.

(vi) The Witt coordinates of $gx-x$ are universal polynomials in those of $gx$ and $x$ vanishing on
the diagonal, and $g\mapsto gx_n$ is continuous for the valuation topology of $C^\flat$; so the
first $m$ coordinates of $gx-x$ tend to $0$, and so do those of $(gx-x)d$ (again universal
polynomials, vanishing when the first factor is $0$). Hence they lie in $\mathfrak m^\flat$ for $g$
near $1$.

(v) Both sides are additive in $v$, so let $v=x\otimes y$. Write $L(1\otimes y)=\iota(d)$; then
$L(x'\otimes y)=\iota(x'd)$. For $g$ near $1$, $gy=y+\ell^mw$, and $L(g x\otimes \ell^mw)=0$ by (i),
so $L(gv-v)=\iota((gx-x)d)\in\pkpsm$ by (vi).

(vii) $z=[\varpi^\flat]$ has coordinates $(\varpi^\flat,0,0,\dots)$ in $\mathfrak m^\flat$. An element of
$\pkpA\otimes\auO(1)$ maps to some $\iota(b)$ with $b\in\pkpA$, and $z b\in\pkpm$ since $\pkpm$ is an
ideal.
\end{proof}

\subsection{Proof of the pairing lemma}\label{pkp:ss-pairing}

Throughout, $\langle v,v'\rangle=\ell^\kappa P(\mathrm{fst}\,v,\mathrm{fst}\,v')\in\pkpM$ for
$v\in\pkpAmb(X_S)$, $v'\in\pkpAmb(X_{S'})$; it is $\pkpW$-bilinear. Put $A_g=\ell^{e_1}\otimes
\alpha(c)(g)$ and $A'_g=\ell^{e_2}\otimes\alpha'(c')(g)$ in $\pkpAmb(Y)$, $\pkpAmb(Y')$. By the
definition of the extension, for $v$ with $\mathrm{pr}(v)=\ell^{e_1}(1\otimes c)$ and $g\in\Gamma$,
\begin{equation}\label{pkp:eq-fst}
\mathrm{fst}(gv)=g\,\mathrm{fst}(v)+A_g,\qquad \mathrm{fst}(\varphi v)=\varphi\,\mathrm{fst}(v),
\end{equation}
and $\mathrm{fst}(gv)=g\,\mathrm{fst}(v)$ if $\mathrm{pr}(v)=0$; similarly on the second side.
$\mathrm{pr}$ commutes with $\Z\times\Gamma$, and $\varphi(1\otimes c)=1\otimes c$, $g(1\otimes
c)=1\otimes c$. Hence $\mathrm{pr}(\varphi x')=\mathrm{pr}(x')$, so $\varphi x'-w'$, $w'-x'$, $gx'-x'$
lie over $0$. The map $g\mapsto A_g$ is a cocycle: $A_{gh}=gA_h+A_g$.

\begin{proof}[Proof of Theorem~\ref{pkp:pairing}]
\emph{Step 0 (the cochain).} By Lemma~\ref{pkp:AS}(ii) choose $y\in\pkpM$ with
$(\varphi-1)y=\langle w,\varphi x'-w'\rangle+\langle x,w'-x'\rangle$. For $g\in\Gamma$ put
\[
f(g)=\ell^\kappa P(A_g,\,g\,\mathrm{fst}\,x')+(g-1)y-\langle x,gx'-x'\rangle .
\]

\emph{Step 1: $(\varphi-1)f(g)\in\pkpsm$ for $g\in\Gamma'$.} Since $\varphi$ commutes with $g$,
fixes $A_g$ and commutes with $P$ and $\mathrm{fst}$,
\[
(\varphi-1)f(g)=\ell^\kappa P(A_g,g\,\mathrm{fst}(\varphi x'-x'))+(g-1)\bigl(\langle
w,a'\rangle+\langle x,b'\rangle\bigr)-\langle\varphi x,g\varphi x'-\varphi x'\rangle+\langle
x,gx'-x'\rangle,
\]
with $a'=\varphi x'-w'$, $b'=w'-x'$, which lie over $0$. By \eqref{pkp:eq-fst}, for $v\in\{w,x\}$ and
$v'$ over $0$, $g\langle v,v'\rangle=\langle gv,gv'\rangle-\ell^\kappa P(A_g,g\,\mathrm{fst}\,v')$. As
$a'+b'=\varphi x'-x'$, the $A_g$-terms cancel and
\[
(\varphi-1)f(g)=\langle gw,ga'\rangle-\langle w,a'\rangle+\langle gx,gb'\rangle-\langle
x,b'\rangle-\langle\varphi x,g\varphi x'-\varphi x'\rangle+\langle x,gx'-x'\rangle .
\]
Expanding $\langle gv,gv'\rangle-\langle v,v'\rangle=\langle gv-v,gv'\rangle+\langle v,gv'-v'\rangle$
and using $b'+x'=w'$, $a'=\varphi x'-w'$, this equals
\[
\langle gw-w,\,g(\varphi x'-w')\rangle-\langle\varphi x-w,\,g\varphi x'-\varphi x'\rangle-\langle
w-x,\,gw'-w'\rangle+\langle gx-x,\,g(w'-x')\rangle .
\]
The four terms are the values $\langle\cdot,\cdot\rangle$ on the pairs tested by the four moves of the path
(b) (second move, first move, third move, fourth move, in this order), each in $\pkpsm$ by
Definition~\ref{pkp:def-safe}.

\emph{Step 2.} By Lemma~\ref{pkp:AS}(iii), for each $g\in\Gamma'$ there is $k(g)\in\Z/\ell^m(1)$ with
$\delta(g)=f(g)-1\otimes k(g)\in\pkpsm$.

\emph{Step 3: the cup product on $\Gamma'$.} For $h\in\Gamma$ put $F_h=\mathrm{fst}(hx'-x')=(h-1)
\mathrm{fst}\,x'+A'_h$, a cocycle. Write $f=T_1+(g-1)y+T_2$ with $T_1(g)=\ell^\kappa P(A_g,g\,
\mathrm{fst}\,x')$ and $T_2(g)=-\ell^\kappa P(\mathrm{fst}\,x,F_g)$. Using the cocycle relations,
\[
gT_1(h)-T_1(gh)+T_1(g)=-\ell^\kappa P(A_g,g(h-1)\mathrm{fst}\,x'),\qquad
gT_2(h)-T_2(gh)+T_2(g)=-\ell^\kappa P((g-1)\mathrm{fst}\,x,gF_h),
\]
and the $y$-part has zero coboundary. On the other hand $\mathrm{fst}(gx-x)=(g-1)\mathrm{fst}\,x+A_g$, so
\[
\langle gx-x,g(hx'-x')\rangle=\ell^\kappa P((g-1)\mathrm{fst}\,x,gF_h)+\ell^\kappa
P(A_g,g(h-1)\mathrm{fst}\,x')+\ell^\kappa P(A_g,gA'_h).
\]
Adding, and noting $\ell^\kappa P(A_g,gA'_h)=1\otimes\ell^{e_1+e_2+\kappa}P(\alpha(c)(g),g\alpha'(c')(h))$,
we get for $g,h\in\Gamma$
\[
1\otimes\ell^{e_1+e_2+\kappa}P(\alpha(c)(g),g\alpha'(c')(h))=\langle
gx-x,g(hx'-x')\rangle+gf(h)-f(gh)+f(g).
\]
For $g,h\in\Gamma'$, substitute $f=1\otimes k+\delta$: the difference
$1\otimes[\ell^{e_1+e_2+\kappa}P(\alpha(c)(g),g\alpha'(c')(h))-(\chi(g)k(h)-k(gh)+k(g))]$ equals
$\langle gx-x,g(hx'-x')\rangle+(g\delta(h)-\delta(gh)+\delta(g))$, which lies in $\pkpsm$ by (a)
(the move by $hx'-x'$ with $x$ at rest, tested at $g$) and Lemma~\ref{pkp:AS}(i). By
Lemma~\ref{pkp:AS}(iv), with $r=\ell^{N-e_1-e_2-\kappa}$,
\begin{equation}\label{pkp:eq-cob}
P(\alpha(c)(g),g\,\ell^N\alpha'(c')(h))=\chi(g)\,rk(h)-rk(gh)+rk(g)\qquad(g,h\in\Gamma').
\end{equation}

\emph{Step 4: $k$ vanishes near $1$.} Cocycles and the action on lattices are continuous, so for $h$
in a neighbourhood $V$ of $1$ (inside the open subgroup $\Gamma'$): $\alpha(c)(h)\in\ell^m Y$,
$\alpha'(c')(h)\in\ell^mY'$, $hy'-y'\in\ell^mY'$ for $y'$ in a fixed finite generating set, and
$\chi(h)\equiv1\pmod{\ell^m}$. By Lemma~\ref{pkp:AS}(i), terms containing $A_h$ or $A'_h$ vanish, so
$f(h)=(h-1)y-\ell^\kappa P(\mathrm{fst}\,x,(h-1)\mathrm{fst}\,x')$. By Lemma~\ref{pkp:AS}(v)
(applied to $\rho=\Z/\ell^m(1)$, $L=\mathrm{id}$, and to $\rho=Y'$, $L=\ell^\kappa
P(\mathrm{fst}\,x,\cdot)$), after shrinking $V$, $f(h)\in\pkpsm$. Then $1\otimes k(h)\in\pkpsm$ and
$k(h)=0$ by Lemma~\ref{pkp:AS}(iv).

\emph{Step 5: from $\Gamma'$ to $\Gamma$.} Let $b(g,h)=P(\alpha(c)(g),g\ell^N\alpha'(c')(h))$, a
2-cocycle of $\Gamma$ with values in $\Z/\ell^m(1)$ (the cup product of two 1-cocycles through an
equivariant pairing). It is locally constant: near any $(g,h)$, the changes of $\alpha(c)$,
$\alpha'(c')$ and of the action are multiples of $\ell^m$, by continuity and the cocycle relations.
Read it in $\auQlb^\times$ through $\Z/\ell^m(1)\cong\mu_{\ell^m}$ and let $\beta$ be its class in the
continuous cohomology $H^2(\Gamma,\auQlb^\times)$ (discrete coefficients). The cochain $rk$ on
$\Gamma'$ is locally constant: for $h\in V$, $rk(h)=0$ and $b(g,h)=0$ (as $\alpha'(c')(h)\in\ell^mY'$),
so \eqref{pkp:eq-cob} gives $rk(gh)=rk(g)$. Hence $\beta|_{\Gamma'}=0$. Since $\Gamma'$ is open of
finite index, $\mathrm{cor}\circ\mathrm{res}=[\Gamma:\Gamma']$ on $H^2(\Gamma,-)$, so
$[\Gamma:\Gamma']\beta=0$; also $\ell^m\beta=0$ since $\ell^m b=0$. As $[\Gamma:\Gamma']$ is prime to
$\ell$, $\beta=0$, which is $\mathrm{SCC}_N(c,c')$.
\end{proof}

\begin{proof}[Proof of Corollary~\ref{pkp:letters}]
\emph{Letters do not interact.} If $v\in\pkpAmb(Y)_i$ and $v'\in\pkpAmb(Y')_{i'}$ with $i\ne i'$,
then $P(v,v')=0$: since $\psi$ is $\auO$-bilinear, $P(\zeta v,v')=P(v,\zeta v')$, i.e.\
$[i]P(v,v')=[i']P(v,v')$, and $[i]-[i']$ is a unit. The projection $\mathrm{fst}$ preserves letter
parts, and $\Gamma'$ preserves them. Hence, if $(v_i,v'_i)$ are $\kappa$-safe pairs with $v_i,v'_i$
at the letter $i$, then $(\sum_iv_i,\sum_iv'_i)$ is $\kappa$-safe (the cross terms vanish and
$\pkpsm$ is a group). Applying this to each move of the paths and to the moves in the cocycle
hypothesis, with $x=\sum x_i$, $w=\sum w_i$, $x'=\sum x'_i$, $w'=\sum w'_i$, we obtain (a) and
(b) of Theorem~\ref{pkp:pairing} (the path starts at $(\sum y_i,\sum y'_i)=(\varphi x,\varphi
x')$). Finally $\mathrm{pr}(x)=\ell^e\sum_ie_i(1\otimes c)=\ell^e(1\otimes c)$, and likewise for
$w,x',w'$. Theorem~\ref{pkp:pairing} applies with $e_1=e$, $e_2=e'$.
\end{proof}

\subsection{The rule and the walk at one cut}\label{pkp:ss-walk}

\begin{lemma}[Extension form]\label{pkp:extform}
Let $T=\operatorname{Hom}(Y\otimes Y',\auO(1))$ and $X_T=T\oplus\auO$ the extension of the Tate-dual
side, $g(\phi,a)=(g\phi+a\,c_P(g),a)$. The $\auO$-bilinear form $\epsilon((\phi,a),v)=\ell^m\phi(v)+a
\psi(v)$ on $X_T\times(Y\otimes Y')$ with values in $\auO(1)$ is $\Gamma$-equivariant, and its
reduction modulo $\ell^m$ (through $\auO\to\Z_\ell$) is $((\phi,a),v)\mapsto P(av)$, with $a$ acting on the $\auO$-module $Y\otimes Y'$; on
$\pkpAmb$ the same holds with $a\in\pkpW\otimes\auO$.
\end{lemma}

\begin{proof}
With $(g\phi)(gv)=g(\phi(v))$ and $(g\psi)(gv)=g(\psi(v))$, $g\psi=\psi+\ell^mc_P(g)$:
$\epsilon(g(\phi,a),gv)=\ell^mg(\phi(v))+a\ell^mc_P(g)(gv)+a\psi(gv)=\ell^mg(\phi(v))+a\,g(\psi(v))$.
The reduction is clear.
\end{proof}

\begin{proof}[Proof of Proposition~\ref{pkp:rule}]
We prove the first assertion; the second is symmetric (see the end). Fix $g\in\Gamma'$; we must show
$\ell^\kappa P(\mathrm{fst}\,y,\mathrm{fst}(gx'-x'))\in\pkpsm$. By Lemma~\ref{pkp:tau}(iv),
$v'=(z\tau)^{-1}(gx'-x')$ has orders $d'$ at $i$ and lies over $0$; so $gx'-x'=z\tau v'$ and it
suffices to show $z\tau\,\ell^\kappa P(\mathrm{fst}\,y,\mathrm{fst}\,v')\in\pkpsm$ for $y\in L^0_i(o)$,
$v'\in L'^0_i(d')$. By Lemma~\ref{pkp:tau}(iii) (Lemma~\ref{pkm:tau-iii}) there is $\nu$ in the Tate-dual extension, lying
over $\ell^\kappa\cdot1$ at $i^{(-l)}$, with orders $\rho_{-t}$; by Lemma~\ref{pkp:tau}(v) (first
assertion, Lemma~\ref{pkt:tau-v}), $\varphi^l\nu$ lies
over $\ell^\kappa\cdot1$ at $i$ with orders $o_\nu=\mathrm{sh}_{l,0}(\rho_{-t})$. One checks from the
definitions that
\begin{equation}\label{pkp:eq-ruleid}
o_\nu(a,-u-u'-1)+\mathrm{rule}_{l,t}(a,u+u')=1
\end{equation}
(with $s=u+u'$: for $a<l$ the terms are $0$ and $1$; for $a=l$ they are $\rho_{-t}(0,-s-1)=s+1-t$ and
$t-s$; for $a>l$ they are $\rho_{-t}(a-l,-s-1)=s+1$ and $-s$). Since $y$, $v'$ lie over $0$ they lie in $\pkpAmb(Y)$,
$\pkpAmb(Y')$, and by Lemma~\ref{pkt:sub} they have orders $o$, $d'$ at $i$ there. By Lemma~\ref{pkp:orders-add} (that is, Lemma~\ref{pkt:add}(ii)), $y\otimes v'$ has orders $o''(a,v)=1-o_\nu(a,-v-1)$ at
$i$ in $Y\otimes Y'$: indeed if $u+u'\le v$ then, $o_\nu(a,\cdot)$ being non-increasing,
$1-o_\nu(a,-v-1)\le1-o_\nu(a,-u-u'-1)=\mathrm{rule}_{l,t}(a,u+u')\le o(a,u)+d'(a,u')$. Pair
$\varphi^l\nu$ with $y\otimes v'$ by $\epsilon$ (Lemma~\ref{pkp:extform}); by
Lemma~\ref{pkp:orders-add} again (for $b=\epsilon$, $L''=\auO(1)$, $U''=\{-1\}$), the value $\xi\in\pkpAmb(\auO(1))$ has orders at least
$o_\nu(a,u)+o''(a,u'')\ge1$ on the jump $-1$ whenever $u+u''\le-1$, that is, orders
$\mathrm{std}_0$ (which is $1$ on the jump $-1$). By Lemma~\ref{pkp:tau}, $\tau\xi$ is integral and the
image of $z\tau\xi$ in $\pkpM$ lies in $\pkpsm$. By Lemma~\ref{pkp:extform}, that image is
$z\tau\,P(\mathrm{pr}(\varphi^l\nu)(y\otimes v'))$ with $\mathrm{pr}(\varphi^l\nu)=\ell^\kappa e_i(1\otimes
1)$; since $y$ lies at the letter $i$, $e_iy=y$, so this is $z\tau\,\ell^\kappa P(e_i(y\otimes v'))=z\tau\,\ell^\kappa P(y,v')$. So $z\tau\ell^\kappa P(y,v')\in\pkpsm$.

For the second assertion: for $g\in\Gamma'$ we need $(gx-x,gy')$ safe. Since $y'$ lies over $0$ and
$L'_i(o')$ is $\Gamma'$-stable (Lemma~\ref{pkm:inputs}(G)), $gy'\in L'^0_i(o')$; and $gx-x=z\tau v$ with $v\in
L^0_i(d)$ by Lemma~\ref{pkp:tau}(iv). The argument above, with the roles of the sides exchanged,
gives the claim.
\end{proof}

\begin{lemma}[Arithmetic of the walk]\label{pkp:walk-arith}
Let $l,l_0>0$, $t,t_0,\eta\ge0$, $a\ge0$, $u,v\in\Z$. Then
\begin{enumerate}
\item[(i)] $\mathrm{rule}_{l,t_0}(a,u+v)\le\min\bigl(\mathrm{stop}_{\eta,l},\mathrm{sh}_{l,0}(\rho_{-\eta})
\bigr)(a,u)+\mathrm{sh}_{l,1}(\mathrm{dep}_{t_0,\eta,l_0})(a,v)$;
\item[(ii)] $\mathrm{rule}_{0,t}(a,u+v)\le\min\bigl(\rho_{-\eta},\mathrm{stop}_{\eta,l}\bigr)(a,u)+
\mathrm{dep}_{t,\eta,l}(a,v)$;
\item[(iii)] for $d\ge0$: $\mathrm{rule}_{d,t}(a,u+v)\le\mathrm{sh}_{d,1}(\mathrm{dep}_{t,\eta,l})(a,u)+
\mathrm{sh}_{d,1}(\mathrm{dep}_{t,\eta,l})(a,v)$.
\end{enumerate}
\end{lemma}

\begin{proof}
(i) If $a<l$: $\mathrm{stop}=\max(0,-u-\eta)\ge0$ and $\mathrm{sh}_{l,0}(\rho_{-\eta})=0$, so the minimum
is $0$, and the last term is $1$; both sides are $1$. If $a=l$: $\mathrm{stop}=-u$,
$\rho_{-\eta}(0,u)=-\eta-u$, the minimum is $-u-\eta$, and $\mathrm{dep}_{t_0,\eta,l_0}(0,v)=
-v+t_0+\eta$; both sides are $t_0-u-v$. If $a>l$: the minimum is $-u$ and
$\mathrm{dep}_{t_0,\eta,l_0}(a-l,v)\ge-v$; the left side is $-u-v$.
(ii) If $a=0$: $\rho_{-\eta}(0,u)=-\eta-u\le\max(0,-u-\eta)=\mathrm{stop}(0,u)$ (as $l>0$), and
$\mathrm{dep}(0,v)=-v+t+\eta$; both sides are $t-u-v$. If $0<a<l$: the minimum is
$\ge-u-\eta$ and $\mathrm{dep}=-v+\eta$; the left side is $-u-v$. If $a\ge l$: the minimum is $-u$,
$\mathrm{dep}=-v$.
(iii) If $a<d$: $1\le1+1$. If $a=d$: $t-u-v\le(-u+t+\eta[0<l])+(-v+t+\eta[0<l])$. If $a>d$:
$-u-v\le(-u+\eta[\cdot])+(-v+\eta[\cdot])$.
\end{proof}

\begin{proof}[Proof of Proposition~\ref{pkp:cut}]
By Lemma~\ref{pkp:tau}(v), $\varphi^lx_0$ and $\varphi^lx'_0$ are points at $i$ over the same
constants with orders $\mathrm{sh}_{l,0}(\rho_{-\eta})$ and depth $\mathrm{sh}_{l,1}(\mathrm{dep}_{t_0,
\eta,l_0})$. The difference of two points over the same constant at $i$ lies over $0$ and has the
minimum of their orders (the modules $L_i(o)$ grow as $o$ decreases).

\emph{Case of a stop $q$ of $S$.} Take $w=q$, $w'=x'$. First move: $\varphi^lx_0-q\in
L^0_i(\min(\mathrm{stop},\mathrm{sh}_{l,0}\rho_{-\eta}))$ with $\varphi^lx'_0$ at rest; by
Proposition~\ref{pkp:rule} (killing at $i^{(-l)}$, rule for $l,t_0$) and Lemma~\ref{pkp:walk-arith}(i) it
is safe. Second move: by $\varphi^lx'_0-x'$ with $q$ at rest: the hypothesis. Third move: $q-x\in
L^0_i(\min(\rho_{-\eta},\mathrm{stop}))$ with $x'$ (depth $\mathrm{dep}_{t,\eta,l}$) at rest: rule for
$0,t$ (killing at $i$) and Lemma~\ref{pkp:walk-arith}(ii). Fourth move: by $x'-x'=0$, trivially safe.

\emph{Case of a stop $q'$ of $S'$.} Take $w=\varphi^lx_0$, $w'=q'$. The first move is by $0$; the
second, by $\varphi^lx'_0-q'$ with $\varphi^lx_0$ at rest, is safe by the second half of
Proposition~\ref{pkp:rule} and Lemma~\ref{pkp:walk-arith}(i); the third is the hypothesis; the fourth,
by $q'-x'$ with $x$ at rest, by the rule for $0,t$ and Lemma~\ref{pkp:walk-arith}(ii). In both cases
$w,w'$ lie over the right constants at $i$.

\emph{The cocycle moves.} By Lemma~\ref{pkp:tau}(v), $\varphi^dx'$ is a point at $i$ with depth
$\mathrm{sh}_{d,1}(\mathrm{dep}_{t,\eta,l})$, so $g\varphi^dx'-\varphi^dx'$ has those orders and lies over
$0$ ($\mathrm{pr}$ of a point is $\Gamma'$-invariant). With $\varphi^dx$ (same depth) at rest,
Proposition~\ref{pkp:rule} (rule for $d,t$, killing at $i^{(-d)}$) and Lemma~\ref{pkp:walk-arith}(iii)
give the claim.
\end{proof}

\begin{remark}[How the cases use this]\label{pkp:rem-use}
In each case one chooses cuts among the letters (a nonempty set of letters with the distance $l(i)>0$
from each letter back to the previous cut), standard points at the cuts, carries them forward by
$\varphi$ to the letters between cuts, and applies Proposition~\ref{pkp:cut} at each cut (the
non-cut letters use the trivial path) and Corollary~\ref{pkp:letters}; the exponent bound is
$N=e+e'+\kappa$, where $\ell^\kappa E^t$ kills the cokernels. The existence of standard points and
stops at the required depth, and the cokernel killers, are the case-specific inputs.
\end{remark}


%% ===== proofs/C-kisin-modules.tex =====
% proofs/C-kisin-modules.tex
% Global walk, stops at plain cuts, improved orders and the cokernel killers, on top of
% proofs/C-kisin-pairing.tex. Labels start with pkm:. Conventions and macros as in C-kisin-pairing.tex.
\providecommand{\auO}{\mathcal O}
\providecommand{\auQlb}{\overline{\Q}_\ell}
\providecommand{\pkpA}{A_{\mathrm{inf}}}
\providecommand{\pkpW}{W(C^\flat)}
\providecommand{\pkpm}{W(\mathfrak m^\flat)}
\providecommand{\pkpM}{M_m}
\providecommand{\pkpsm}{\mathrm{Sm}_m}
\providecommand{\pkpAmb}{\mathbb A}
% macros of this file (prefix \pkm)
\providecommand{\pkmJ}[3]{J^{#1}_{#2}(#3)}
\providecommand{\pkmI}[2]{I_{#1}(#2)}
\providecommand{\pkmv}{v_\ell}

\section{The global walk and the cokernel killers}\label{pkm:sec}

This section continues \S\ref{pkp:sec}, whose conventions, definitions and labels we use. In
\S\ref{pkm:ss-state} we state everything that the four cases (Propositions~\ref{pki:generic},
\ref{pki:any}, \ref{pki:tame}, \ref{pki:sparse}) cite beyond \S\ref{pkp:ss-arch}: the global walk,
the stops at plain cuts, the improvement of orders, and the cokernel killers. Proofs follow in
\S\ref{pkm:ss-walk}--\S\ref{pkm:ss-kill}. The inputs from the theory of the modules $\mathcal N$
that these proofs use, beyond Lemmas~\ref{pkp:orders-add} and~\ref{pkp:tau}, are collected in
Lemma~\ref{pkm:inputs}.

\emph{Dependencies.} Proved in full, given Lemma~\ref{pkp:tau}(iv),(v), items (M), (E), (G) of Lemma~\ref{pkm:inputs} and the statements of
\S\ref{pkp:ss-arch} proved there: the global walk (Theorem~\ref{pkm:walk},
Corollary~\ref{pkm:walk-all}), the stops at plain cuts (Lemmas~\ref{pkm:down}, \ref{pkm:plain}), the
improvement of orders (Lemma~\ref{pkm:improved}), and Lemma~\ref{pkp:tau}(iii). Proved in full
\emph{given Lemma~\ref{pkm:inputs}}: the cokernel killers (Propositions~\ref{pkm:kill-small},
\ref{pkm:kill-unram}, \ref{pkm:kill-J} and Lemma~\ref{pkm:zpow}). The items of
Lemma~\ref{pkm:inputs}, Lemma~\ref{pkp:orders-add} and Lemma~\ref{pkp:tau}(i),(iv),(v) are proved in
\S\S\ref{pkn:sec}--\ref{pkt:sec} (see the paragraph at the end of \S\ref{pkm:ss-kill}).

\subsection{Statements used by the cases}\label{pkm:ss-state}

Throughout, $S,S'$ are sides with lattices $Y,Y'$ and parameter lattices $C_S,C_{S'}$, $P$ is a
pairing modulo $\ell^m$ of $S$ and $S'$, $T_P$ is its Tate-dual side (Definition~\ref{pkp:def-pairing}),
$c\in C_S$, $c'\in C_{S'}$. The Frobenius $\sigma(i)=i^\ell$ permutes the (finitely many) letters, and
$i^{(-a)}=\sigma^{-a}(i)$. For a side $S$, ``the jumps of $S$ at $i$'' are the elements of its jump
set $U(i)$ (which contains $0$).

\subsubsection*{The global walk}

\begin{theorem}[The walk]\label{pkm:walk}
Let $\mathrm{cut}$ be a set of letters and $l,t$ functions from letters to $\Z_{\ge0}$ such that for
every letter $i$, $l(i)$ is the least $a>0$ with $i^{(-a)}\in\mathrm{cut}$. Let $\kappa,e,e',\eta,N\ge0$
and, for every $i\in\mathrm{cut}$, let $x(i)\in\pkpAmb(X_S)$, $x'(i)\in\pkpAmb(X_{S'})$. Assume, for every
$i\in\mathrm{cut}$:
\begin{enumerate}
\item[(K)] $\ell^\kappa E^{t(i)}$ kills the cokernel of $T_P$ at $i$;
\item[(P)] $x(i)$ is a standard point (for $t(i),\eta,l(i)$) of $S$ over $\ell^ec$ at $i$, and $x'(i)$ a
standard point (for $t(i),\eta,l(i)$) of $S'$ over $\ell^{e'}c'$ at $i$;
\item[(St)] either there is a stop $q$ (for $\eta,l(i)$) of $S$ over $\ell^ec$ at $i$ such that the
move of the second side by $\varphi^{l(i)}x'(i^{(-l(i))})-x'(i)$ with $q$ at rest is $\kappa$-safe, or
there is a stop $q'$ (for $\eta,l(i)$) of $S'$ over $\ell^{e'}c'$ at $i$ such that the move of the
first side by $\varphi^{l(i)}x(i^{(-l(i))})-x(i)$ with $q'$ at rest is $\kappa$-safe.
\end{enumerate}
If $N\ge e+e'+\kappa$, then $\mathrm{SCC}_N(c,c')$.
\end{theorem}

(The hypothesis on $l$ forces $\mathrm{cut}\ne\emptyset$.)

\begin{corollary}[Cuts at all letters]\label{pkm:walk-all}
Let $\kappa,e,e',t,\eta,N\ge0$, and for every letter $i$ let $\ell^\kappa E^t$ kill the cokernel of
$T_P$ at $i$, let $x(i)$, $x'(i)$ be standard points (for $t,\eta,1$) of $S,S'$ over $\ell^ec$,
$\ell^{e'}c'$ at $i$, and assume (St) with $l(i)=1$. If $N\ge e+e'+\kappa$, then $\mathrm{SCC}_N(c,c')$.
\end{corollary}

\subsubsection*{Improved orders and stops at plain cuts}

For orders $o_{\mathrm{lo}},o_{\mathrm{hi}}$ put
\[
\mathrm{imp}(o_{\mathrm{lo}},o_{\mathrm{hi}})(a,u)=\begin{cases}o_{\mathrm{hi}}(a,u)&\text{if }
o_{\mathrm{lo}}(a,u)>0,\\ \min(o_{\mathrm{hi}}(a,u),0)&\text{otherwise,}\end{cases}
\qquad
\mathrm{pd}_{t,\eta,l,\eta_0}(a,u)=\mathrm{dep}_{t,\eta,l}(a,u)+\eta_0[a=l].
\]

\begin{lemma}[Improved orders]\label{pkm:improved}
Let $L$ be a lattice with jump sets and $i$ a letter. If $x\in L_i(o_{\mathrm{lo}})$ and $gx-x\in
L_i(o_{\mathrm{hi}})$ for every $g\in\Gamma'$, then $x\in L_i(\mathrm{imp}(o_{\mathrm{lo}},o_{\mathrm{hi}}))$.
\end{lemma}

\begin{lemma}[A standard point of a side that is down is a stop]\label{pkm:down}
Let $0<l$, and suppose every jump of $S$ at $i^{(-a)}$, $a<l$, is $\le\eta$. Then every standard point
(for $t,\eta,l$) of $S$ over $\ell^ec$ at $i$ is a stop (for $\eta,l$) of $S$ over $\ell^ec$ at $i$.
\end{lemma}

\begin{lemma}[Stops at plain cuts]\label{pkm:plain}
Let $0<l$, $t,\eta,\eta_0\ge0$, and let $\ell^\kappa E^t$ kill the cokernel of $T_P$ at $i$. Suppose that
every jump of $S$ at $i^{(-a)}$ ($a<l$) is $\le\eta$ and every jump of $S'$ at $i^{(-a)}$ ($a<l$) is
$\ge-\eta$. Let $x$ be a point of $S$ over $\ell^ec$ at $i$ with orders $\rho_{-\eta}$ and depth
$\mathrm{pd}_{t,\eta,l,\eta_0}$; let $x'$ be a point of $S'$ over $\ell^{e'}c'$ at $i$ with orders
$\rho_{-\eta}$ and $x'_0$ a point of $S'$ over $\ell^{e'}c'$ at $i^{(-l)}$ with orders $\rho_{-\eta_0}$
(any depths). Then $x$ is a stop (for $\eta,l$) of $S$ over $\ell^ec$ at $i$, and the move of the
second side by $\varphi^lx'_0-x'$ with $x$ at rest is $\kappa$-safe. The same holds with the roles of
$S$ and $S'$ (and of the two moves) exchanged.
\end{lemma}

In particular, at a plain cut the alternative (St) of Theorem~\ref{pkm:walk} holds as soon as the
standard point of the side that is down has depth $\mathrm{pd}_{t,\eta,l,\eta_0}$ and the points of
the other side have orders $\rho_{-\eta}$ (at $i$) and $\rho_{-\eta_0}$ (at $i^{(-l)}$); this is
Lemma~\ref{pkm:plain}, since a point with depth $\mathrm{pd}_{t,\eta,l,\eta_0}\ge\mathrm{dep}_{t,\eta,l}$
and orders $\rho_{-\eta}$ is standard.

\subsubsection*{The lift ideals and the cokernel killers}

\begin{definition}[Lift ideals]\label{pkm:def-J}
Let $S$ be a side, $i$ a letter, $l\ge0$. The ideal $\pkmI li\subseteq\pkpA$ consists of the
$s\in\pkpA$ such that for every $c\in C_S$ there is $x\in(X_S)_i(\mathrm{std}_l)$ (the module of
Definition~\ref{pkp:def-orders} for the extension $X_S$) with $\mathrm{pr}(x)=s\,e_i(1\otimes c)$.
For $r\in\Z_{\ge0}$ put $\pkmJ rli=\pkmI li+r\pkpA$. Thus $\ell^\kappa E^t$ kills the cokernel at $i$
(Definition~\ref{pkp:def-points}; see Remark~\ref{pkm:rem-kill}) if and only if
$\ell^\kappa E^t\in\pkmI0i$.
\end{definition}

$\pkmI li$ is an ideal, because $(X_S)_i(\mathrm{std}_l)$ is an $\pkpA$-module and $\mathrm{pr}$ is
$\pkpW$-linear; the $s$ with the lifting property for one $c$ form an ideal, and $\pkmI li$ is their
intersection.

\begin{remark}[Killing the cokernel: lifts in $\pkpA\mathcal N$ versus $\mathcal N$]\label{pkm:rem-kill}
In Definition~\ref{pkp:def-points} the lift $n$ is required to lie in $(X_S)_i(\mathrm{std}_0)$, i.e.\
in the letter-$i$ part of $\pkpA\mathcal N$ with orders $\mathrm{std}_0$, not in $\mathcal N(X_S)$
itself. This is the condition proved below, and it is all that the proof of
Proposition~\ref{pkp:rule} uses (through Lemma~\ref{pkp:tau}(iii), proved in
Lemma~\ref{pkm:tau-iii}). For a multiplier $f\in\mathfrak S$ (in particular $f=\ell^\kappa E^t$) the
stronger reading follows: by the retraction (R) a lift in $(X_S)_i(\mathrm{std}_l)$ can be replaced
by one in $\mathcal N(X_S)$ (Lemma~\ref{pkt:lift}). This is how the input (I1) produces points in
$\mathcal N(X_S)$ itself (Proposition~\ref{pkt:I1}).
\end{remark}

Let $d_v$ be as in Definition~\ref{pkp:def-chart} and $e$ the ramification index of the chart. All
jump bounds below refer to the jump sets of the side $S$ (in the cases, $S$ is the Tate-dual side
$T_P$, whose jump sets are $\{0\}\cup\{-u-u'-1\}$).

\begin{proposition}[Small jumps]\label{pkm:kill-small}
If every jump of $S$ at every letter is $\le\omega$ and $e\,\omega\le\ell-2$, then $1=\ell^0E^0$ kills
the cokernel of $S$ at every letter.
\end{proposition}

\begin{proposition}[Unramified charts]\label{pkm:kill-unram}
Let the chart be unramified ($e=1$, $E=z+\ell$), $\ell\ge5$, and $s,t\ge0$. If every jump of $S$ at
every letter is $\le(\ell-3)(s+t)$, then $\ell^sE^t$ kills the cokernel of $S$ at every letter.
\end{proposition}

\begin{lemma}[A power of $z$ in the lift ideals]\label{pkm:zpow}
If every jump of $S$ at every letter is $\le\omega$ and $\kappa\ge0$, there is $a\ge0$ with
$(\ell-1)a\le e(\omega+\kappa)$ and $z^a\in\pkmJ{\ell^{\kappa+1}}0i$ for every letter $i$.
\end{lemma}

\begin{proposition}[Killing by a power of $\ell$]\label{pkm:kill-J}
Let $i$ be a letter, $\omega,\omega',J,\kappa\ge0$. Assume every jump of $S$ at every letter is
$\le\omega$, every jump of $S$ at $\sigma^s(i)$ for $1\le s\le J$ is $\le\omega'$, $8e\le\ell^{J+1}$,
and
\[
2Jd_v\omega'+\frac{(2J+4)\,\omega'}{\ell-1}+\frac{8e\,\omega}{\ell^{J+1}}\le\kappa .
\]
Then $\ell^\kappa$ kills the cokernel of $S$ at $i$.
\end{proposition}

\subsection{Proof of the walk}\label{pkm:ss-walk}

\begin{lemma}[Distance to the previous cut]\label{pkm:dist}
In the situation of Theorem~\ref{pkm:walk} put $d(i)=0$ if $i\in\mathrm{cut}$ and $d(i)=l(i)$
otherwise. Then $i^{(-d(i))}\in\mathrm{cut}$ and $l(i)=d(i^{(-1)})+1$ for every letter $i$.
\end{lemma}

\begin{proof}
The first claim is the definition of $l$. Let $i_1=i^{(-1)}$; note $i_1^{(-a)}=i^{(-a-1)}$. If
$i_1\in\mathrm{cut}$, then $a=1$ is admissible, so $l(i)=1=d(i_1)+1$. Otherwise $l(i)\ge2$ and
$l(i)-1>0$ is admissible for $i_1$, so $l(i_1)\le l(i)-1$; conversely $l(i_1)+1$ is admissible for
$i$, so $l(i)\le l(i_1)+1$. Hence $l(i)=l(i_1)+1=d(i_1)+1$.
\end{proof}

\begin{proof}[Proof of Theorem~\ref{pkm:walk}]
For every letter $i$ put $\tilde x_i=\varphi^{d(i)}x(i^{(-d(i))})$, $y_i=\varphi^{l(i)}x(i^{(-l(i))})$,
and likewise $\tilde x'_i$, $y'_i$; these make sense by Lemma~\ref{pkm:dist}. By
Lemma~\ref{pkp:tau}(v), $\varphi^a$ of a point at $i^{(-a)}$ over $\ell^ec$ is a point at $i$ over
$\ell^ec$; hence $\tilde x_i,y_i$ lie over $\ell^ec$ at $i$, and $\tilde x'_i,y'_i$ over $\ell^{e'}c'$
at $i$.

{\emergencystretch=6em \emph{Frobenius.} Since $i\mapsto i^{(-1)}$ is a bijection of the letters and, by
Lemma~\ref{pkm:dist}, $\varphi\tilde x_{i^{(-1)}}=\varphi^{d(i^{(-1)})+1}x(i^{(-1-d(i^{(-1)}))})=y_i$,
\[
\varphi\sum_i\tilde x_i=\sum_i\varphi\tilde x_{i^{(-1)}}=\sum_iy_i,
\]
and likewise $\varphi\sum_i\tilde x'_i=\sum_iy'_i$.\par}

\emph{Paths.} Let $i\in\mathrm{cut}$, so $\tilde x_i=x(i)$, $\tilde x'_i=x'(i)$. The letter
$i^{(-l(i))}$ is a cut with $l(i^{(-l(i))})>0$; by (K) at $i$ and at $i^{(-l(i))}$, (P) at both and
(St) at $i$, Proposition~\ref{pkp:cut} (with $l=l(i)$, $l_0=l(i^{(-l(i))})$, $t=t(i)$,
$t_0=t(i^{(-l(i))})$, $x_0=x(i^{(-l(i))})$, $x'_0=x'(i^{(-l(i))})$) gives $w_i,w'_i$ over $\ell^ec$,
$\ell^{e'}c'$ at $i$ and a $\kappa$-safe path from $(y_i,y'_i)$ through $(w_i,w'_i)$ to
$(\tilde x_i,\tilde x'_i)$. Let $i\notin\mathrm{cut}$; then $\tilde x_i=y_i$, $\tilde x'_i=y'_i$. Take
$w_i=y_i$, $w'_i=y'_i$: each of the four moves is a move by $0$, and $(0,\cdot)$, $(\cdot,0)$ are
safe pairs since $P$ is bilinear and $0\in\pkpsm$.

\emph{Cocycle moves.} Let $i$ be a letter and $j=i^{(-d(i))}\in\mathrm{cut}$. By (K) and (P) at $j$,
the last assertion of Proposition~\ref{pkp:cut} (with the letter $j$ in place of $i^{(-d)}$, $d=d(i)$)
says that for every $h\in\Gamma'$ the move of the second side by $h\tilde x'_i-\tilde x'_i$ with
$\tilde x_i$ at rest is $\kappa$-safe.

All hypotheses of Corollary~\ref{pkp:letters} hold for the vectors $y_i,w_i,\tilde x_i$ and
$y'_i,w'_i,\tilde x'_i$, with $e,e',\kappa$. Hence $\mathrm{SCC}_N(c,c')$ for $N\ge e+e'+\kappa$.
\end{proof}

\begin{proof}[Proof of Corollary~\ref{pkm:walk-all}]
Take $\mathrm{cut}$ = all letters, $l\equiv1$ (the least $a>0$ with $i^{(-a)}$ a cut is $1$), $t$
constant, and apply Theorem~\ref{pkm:walk}.
\end{proof}

\subsection{Improved orders and stops}\label{pkm:ss-stops}

Recall (Definition~\ref{pkp:def-chart}) that $j(\Gamma')=G_{K'}$ for a finite extension $K'$ of
$\Q_\ell$ (the fixed field of $j(\Gamma')$); $\auQlb\subset B^+_{\mathrm{dR}}$ and $\mathrm{Fil}^n
B_{\mathrm{dR}}=t^nB^+_{\mathrm{dR}}$, with $\mathrm{gr}^nB_{\mathrm{dR}}\cong C(n)$ equivariantly.

\begin{lemma}[Invariants modulo the filtration]\label{pkm:BdR}
Let $y\in B_{\mathrm{dR}}$ and $n\in\Z$ with $gy-y\in\mathrm{Fil}^nB_{\mathrm{dR}}$ for all $g\in G_{K'}$.
Then there is $c\in K'$ with $y-c\in\mathrm{Fil}^nB_{\mathrm{dR}}$. If moreover $y\in\mathrm{Fil}^m$
with $m\ge1$, then $c$ can be taken $0$; and an element of $K'$ lying in $\mathrm{Fil}^1$ is $0$.
\end{lemma}

\begin{proof}
The last claim: $\theta$ is injective on $\auQlb\subset B^+_{\mathrm{dR}}$, and $\mathrm{Fil}^1=\ker\theta$.
By Tate's theorem, $C(j)^{G_{K'}}=0$ for $j\ne0$ and $C^{G_{K'}}=K'$ (\cite[\S3, Theorems~1 and~2, p.~176]{pklTate67}; also quoted in \cite[1.5.7]{pkfFontaineII}). Let $y\in\mathrm{Fil}^m$ (some $m$ exists). We show by induction
on $n-m$ that $y-c\in\mathrm{Fil}^n$ for some $c\in K'$, with $c=0$ if $m\ge1$. If $m\ge n$ take
$c=0$. Otherwise the image $\bar y$ of $y$ in $\mathrm{gr}^m=C(m)$ is $G_{K'}$-invariant (as
$gy-y\in\mathrm{Fil}^n\subseteq\mathrm{Fil}^{m+1}$). If $m\ne0$, then $\bar y=0$, so $y\in\mathrm{Fil}^{m+1}$
and we conclude by induction. If $m=0$, then $\bar y=\theta(y)\in K'$; put $c_0=\theta(y)\in K'\subset
B^+_{\mathrm{dR}}$; then $y-c_0\in\mathrm{Fil}^1$ is still invariant modulo $\mathrm{Fil}^n$ ($g$ fixes
$c_0$), and by induction (with $m\ge1$) $y-c_0\in\mathrm{Fil}^n$.
\end{proof}

\begin{lemma}[Values are equivariant]\label{pkm:equiv}
Let $L$ be a lattice with jump sets, $x\in\pkpA\mathcal N(L)$, $a\ge0$, $\mu$ an $\pkpA$-linear map
$\pkpAmb(L)\to B_{\mathrm{dR}}$ that is $\Gamma'$-equivariant on $1\otimes L$, and $g\in\Gamma'$. Then
$\mu(\varphi^{-a}gx)=g\,\mu(\varphi^{-a}x)$.
\end{lemma}

\begin{proof}
There is $r\in\pkpA$, a product of powers of $\tau$ and of the $p_b$, with $rx\in\pkpA\otimes L$
(Lemma~\ref{pkm:inputs}(M)). Write $\varphi^{-a}(rx)=\sum_kb_k\otimes y_k$ with $b_k\in\pkpA$. Then
$\varphi^{-a}(g(rx))=g\varphi^{-a}(rx)=\sum_k g(b_k)\otimes gy_k$, and by linearity and equivariance
$\mu(\varphi^{-a}g(rx))=\sum_kg(b_k)\,g\mu(1\otimes y_k)=g\mu(\varphi^{-a}(rx))$. Since
$\mu(\varphi^{-a}(rv))=\varphi^{-a}(r)\mu(\varphi^{-a}v)$ for all $v$, and $g$ commutes with
$\varphi^{-a}$, this reads $g(\varphi^{-a}r)\cdot\mu(\varphi^{-a}gx)=g(\varphi^{-a}r)\cdot
g\mu(\varphi^{-a}x)$. The image of $\varphi^{-a}(r)$ in the field $B_{\mathrm{dR}}$ is nonzero
($\pkpA\to B^+_{\mathrm{dR}}$ is injective and $\pkpA$ is a domain), hence so is that of $g\varphi^{-a}(r)$;
divide.
\end{proof}

\begin{proof}[Proof of Lemma~\ref{pkm:improved}]
Let $a\ge0$, $u\in U(i^{(-a)})$, $\mu$ a form of type $u$ at $i^{(-a)}$, $y=\mu(\varphi^{-a}x)$, and
$n=o_{\mathrm{hi}}(a,u)$. By Lemma~\ref{pkm:equiv}, $gy-y=\mu(\varphi^{-a}(gx-x))\in\mathrm{Fil}^n$ for all
$g\in\Gamma'$, so by Lemma~\ref{pkm:BdR} there is $c\in K'$ with $y-c\in\mathrm{Fil}^n$. Also
$y\in\mathrm{Fil}^{o_{\mathrm{lo}}(a,u)}$, so $c\in\mathrm{Fil}^{\min(o_{\mathrm{lo}},o_{\mathrm{hi}})(a,u)}$.
If $o_{\mathrm{lo}}(a,u)>0$ and $n>0$, then $c=0$ by Lemma~\ref{pkm:BdR}; hence $c\in\mathrm{Fil}^{\mathrm{imp}(a,u)}$
in all cases ($\mathrm{imp}\le0$ when $o_{\mathrm{lo}}\le0$, $c\in K'\subset\mathrm{Fil}^0$; and
$\mathrm{imp}=n\le0$ or $c=0$ when $o_{\mathrm{lo}}>0$). As $\mathrm{imp}\le n$, $y=(y-c)+c\in
\mathrm{Fil}^{\mathrm{imp}(a,u)}$. Membership in $\pkpA\mathcal N$ and in the letter part is unchanged.
\end{proof}

\begin{lemma}[Comparison]\label{pkm:le}
(i) $L_i(o)\subseteq L_i(o')$ if $o'(a,u)\le o(a,u)$ for all $a$ and $u\in U(i^{(-a)})$; hence a point
with orders $o_1$ and depth $d_1$ has orders $o_2$ and depth $d_2$ when $o_2\le o_1$ and $d_2\le
d_1$ on the jumps. (ii) If $q\in\mathcal N(X_S)$ lies over $\ell^ec$ at $i$ and has orders $o\ge
\mathrm{stop}_{\eta,l}$ on the jumps, then $q$ is a stop (for $\eta,l$).
\end{lemma}

\begin{proof}
(i) is immediate from Definition~\ref{pkp:def-orders} ($\mathrm{Fil}^n$ decreases in $n$). (ii) By (i),
$q\in L_i(\mathrm{stop}_{\eta,l})$; since $L_i(\mathrm{stop}_{\eta,l})$ is $\Gamma'$-stable
(Lemma~\ref{pkm:inputs}(G)), $gq-q\in L_i(\mathrm{stop}_{\eta,l})$ for $g\in\Gamma'$, i.e.\ $q$ is deep
for $\mathrm{stop}_{\eta,l}$.
\end{proof}

\begin{proof}[Proof of Lemma~\ref{pkm:down}]
Let $x$ be standard: orders $\rho_{-\eta}$, depth $\mathrm{dep}=\mathrm{dep}_{t,\eta,l}$. By
Lemma~\ref{pkm:improved}, $x$ has orders $\mathrm{imp}(\rho_{-\eta},\mathrm{dep})$. By
Lemma~\ref{pkm:le}(ii) it suffices to show $\mathrm{imp}(\rho_{-\eta},\mathrm{dep})(a,u)\ge
\mathrm{stop}_{\eta,l}(a,u)$ for $u\in U(i^{(-a)})$. \emph{$a=0$} ($<l$): $\rho_{-\eta}=-u-\eta$,
$\mathrm{dep}=-u+t+\eta$, $\mathrm{stop}=\max(0,-u-\eta)$. If $-u-\eta>0$, $\mathrm{imp}=-u+t+\eta\ge-u-\eta=
\mathrm{stop}$. Otherwise $\mathrm{stop}=0$ and $\mathrm{imp}=\min(-u+t+\eta,0)=0$ as $u\le\eta$.
\emph{$0<a<l$}: $\rho_{-\eta}=-u$, $\mathrm{dep}=-u+\eta\ge0$ (as $u\le\eta$), $\mathrm{stop}=
\max(0,-u-\eta)$. If $u<0$, $\mathrm{imp}=-u+\eta\ge\max(0,-u-\eta)$; if $u\ge0$, $\mathrm{imp}=
\min(-u+\eta,0)=0=\mathrm{stop}$. \emph{$a\ge l$}: $\rho_{-\eta}=\mathrm{dep}=\mathrm{stop}=-u$, and
$\mathrm{imp}=-u$ in both cases ($\min(-u,0)=-u$ when $-u\le0$).
\end{proof}

\begin{proof}[Proof of Lemma~\ref{pkm:plain}]
As $\mathrm{pd}_{t,\eta,l,\eta_0}\ge\mathrm{dep}_{t,\eta,l}$, $x$ is a standard point (Lemma~\ref{pkm:le}(i)),
hence a stop by Lemma~\ref{pkm:down}. By Lemma~\ref{pkp:tau}(v), $\varphi^lx'_0$ is a point at $i$ over
$\ell^{e'}c'$ with orders $\mathrm{sh}_{l,0}(\rho_{-\eta_0})$; so $y'=\varphi^lx'_0-x'$ lies over $0$
and has orders $o'=\min(\rho_{-\eta},\mathrm{sh}_{l,0}(\rho_{-\eta_0}))$. By the second half of
Proposition~\ref{pkp:rule} with $l=0$ (killing at $i^{(0)}=i$), it suffices that
$\mathrm{rule}_{0,t}(a,u+v)\le o'(a,u)+\mathrm{pd}_{t,\eta,l,\eta_0}(a,v)$ for $u$ a jump of $S'$ at
$i^{(-a)}$ and $v$ a jump of $S$ there. \emph{$a=0$}: left $t-u-v$; right
$\min(-u-\eta,0)+(-v+t+\eta)\ge t-u-v$ since $-u-\eta+\eta=-u$ and $0+\eta\ge-u$ ($u\ge-\eta$).
\emph{$0<a<l$}: left $-u-v$; right $\min(-u,0)+(-v+\eta)\ge-u-v$ as $\eta\ge-u$. \emph{$a=l$}: left
$-u-v$; $o'(l,u)=\min(-u,-u-\eta_0)=-u-\eta_0$ and $\mathrm{pd}(l,v)=-v+\eta_0$. \emph{$a>l$}: $o'=-u$,
$\mathrm{pd}=-v$. The statement with the sides exchanged is proved in the same way with the first
half of Proposition~\ref{pkp:rule}.
\end{proof}

\begin{lemma}[Lemma~\ref{pkp:tau}(iii)]\label{pkm:tau-iii}
If $\ell^\kappa E^t$ kills the cokernel of a side $S$ at $i$, then for every $c\in C_S$ there is
$\nu\in(X_S)_i(\rho_{-t})$ with $\mathrm{pr}(\nu)=\ell^\kappa e_i(1\otimes c)$.
\end{lemma}

\begin{proof}
Take $x\in(X_S)_i(\mathrm{std}_0)$ with $\mathrm{pr}(x)=\ell^\kappa E^te_i(1\otimes c)$ and put
$\nu=E^{-t}x$ ($E$ is a unit of $\pkpW$). By Lemma~\ref{pkm:inputs}(E), $\nu$ has orders
$\mathrm{std}_0-t[a=0]=\rho_{-t}$; and $\mathrm{pr}$ is $\pkpW$-linear.
\end{proof}

This is the form in which Lemma~\ref{pkp:tau}(iii) is used in the proof of
Proposition~\ref{pkp:rule}: there $\nu$ need only have orders $\rho_{-t}$ and lie over $\ell^\kappa\cdot1$
at $i^{(-l)}$, and Lemma~\ref{pkp:tau}(v) is applied to its orders only.

\subsection{The cokernel killers}\label{pkm:ss-kill}

\subsubsection*{Inputs}

\begin{lemma}[Inputs from the theory of the modules $\mathcal N$]\label{pkm:inputs}
Let $L$ be a lattice with jump sets, $i$ a letter, $o$ orders, and $g\in\Gamma'$. We write
$\mathcal N_i(o)=\mathcal N(L)\cap L_i(o)$, an $\mathfrak S$-module.
\begin{enumerate}
\item[(M)] Every $x\in\pkpA\mathcal N(L)$ has $rx\in\pkpA\otimes L$ for some product $r$ of powers of
$\tau$ and of the $p_b$; $\pkpA\mathcal N(L)$ and $\mathcal N(L)$ are stable under multiplication by
$p_b^{\pm1}$ (in $\pkpW^\times$), and $\pkpA\mathcal N(L)$ is the $\pkpA$-span of $\mathcal N(L)$.
\item[(E)] For $b\ge0$, $n\in\Z$ and $x\in L_i(o)$: $p_b^nx\in L_i(o+n[a=b])$.
\item[(G)] $L_i(o)$ and $\pkpA\mathcal N(L)$ are $\Gamma'$-stable; if $x\in\mathcal N(L)$
and $gx-x\in L_i(o)$, then $(z\tau)^{-1}(gx-x)\in L_i(o)$ (this is Lemma~\ref{pkp:tau}(iv)). The second assertion follows from the first and from
\begin{enumerate}
\item[(G$_0$)] for $x\in\mathcal N(L)$ and $g\in\Gamma'$: $(z\tau)^{-1}(gx-x)\in\pkpA\mathcal N(L)$
\end{enumerate}
(Proposition~\ref{pkl:G0}).
\item[($\Phi$)] $\varphi(\mathcal N(L))\subseteq\mathcal N(L)$.
\item[(F)] For $l\ge0$, $\varphi^l$ maps $L_i(\mathrm{std}_0)$ bijectively onto
$L_{\sigma^l(i)}(\mathrm{std}_l)$.
\item[(H)] (Frame.) There is $H\ge0$ with $\tau^H(1\otimes y)\in\pkpA\mathcal N(L)$ and
$e_i\tau^H(1\otimes y)\in L_i(\mathrm{std}_0)$ for all $y\in L$. If $\Gamma$ acts trivially on $L$, then
$\mathcal N(L)\cap(\pkpA\otimes L)=\mathfrak S\otimes L$, $L_i(\mathrm{std}_0)\subseteq\pkpA\otimes L$, and
$e_i(1\otimes y)\in\mathcal N_i(\mathrm{std}_l)$ for all $l\ge0$.
\item[(R)] (Retraction.) There are an $\mathfrak S$-linear map $\lambda_0\colon\pkpA\to\mathfrak S$ that is the
identity on $\mathfrak S$, and, for every lattice, an $\mathfrak S$-linear map $\mathrm{ret}\colon
\pkpA\mathcal N\to\mathcal N$ with $\mathrm{ret}(bx)=\lambda_0(b)x$ for $b\in\pkpA$, $x\in\mathcal N$,
mapping $L_i(\mathrm{std}_l)$ into $\mathcal N_i(\mathrm{std}_l)$ and commuting with $\mathrm{pr}$ (for an
extension $X_S$).
\item[(D)] (Level one.) If $x\in\mathcal N_i(\mathrm{std}_1)$, then $gx-x\in L_i(o_1)$, where
$o_1(0,u)=1$ and $o_1(a,u)=-u$ for $a\ge1$.
\item[(C)] (Chart identities.) (a) $z\tau E$ divides $gs-s$ in $\pkpA$ for all $s\in\mathfrak S$.
(b) Some $t_1\in\Gamma'$ has $t_1z=[\varepsilon]z$. (c) For every $n\ge1$ there are
$\alpha,\beta\in\pkpA$ with $\tau E(\ell^{\pkmv(n)}+\beta z)=\alpha([\varepsilon]^n-1)$. (d) For every
$n\ge1$ there are $g\in\Gamma'$, $\omega\in\pkpA^\times$ and $\upsilon,q\in\pkpA$ with $gE=\omega E$,
$E\mid\upsilon q-\ell^{\pkmv(n)+d_v}$ and $z\tau q=\omega^n-1$.
\end{enumerate}
\end{lemma}

Also used, and elementary: $z,\tau,E$ are units of $\pkpW$ (their first Witt coordinates are nonzero
elements of the field $C^\flat$); $\varphi(z)=z^\ell$ and $\varphi(\mathfrak S)\subseteq\mathfrak S$; $z^e\in
E\mathfrak S+\ell\mathfrak S$ and $z\mid E-\ell$ in $\mathfrak S$ (from $E=z^eA+\ell$ with $A\in\mathfrak
S^\times$); $\varphi(e_i(1\otimes c))=e_{\sigma(i)}(1\otimes c)$ and $g\,e_i(1\otimes c)=e_i(1\otimes c)$
for $g\in\Gamma'$ (Definition~\ref{pkp:def-letters}); and $1-\ell b\in\pkpA^\times$ for $b\in\pkpA$.

\subsubsection*{Operations on the lift ideals}

Fix a side $S$; $\pkmI li$, $\pkmJ rli$ refer to $S$, and ``jumps'' to its jump sets. We write
$\varepsilon_i(c)=e_i(1\otimes c)$ for $c\in C_S$ (the constant of $c$ at $i$).

\begin{lemma}\label{pkm:J-basic}
(i) (Nakayama.) If $\ell^c\in\pkmJ{\ell^{c+1}}li$, then $\pkmJ{\ell^{c+1}}li=\pkmI li$; in particular
$\ell^c\in\pkmI li$.
(ii) (Up.) If $s\in\pkmJ r0i$ and every jump at $i$ is $\le n$, then $sE^n\in\pkmJ r1i$.
(iii) (Frobenius.) $\varphi^l(s)\in\pkmJ rl{\sigma^l(i)}$ if and only if $s\in\pkmJ r0i$.
(iv) (Integral witness.) If $f\in\mathfrak S\cap\pkmJ rli$, there is $t\in\mathfrak S$ such that for every
$c$ some $x\in\mathcal N_i(\mathrm{std}_l)$ (for $X_S$) has $\mathrm{pr}(x)=(f-rt)\varepsilon_i(c)$.
\end{lemma}

\begin{proof}
(i) Write $\ell^c=s+\ell^{c+1}b$ with $s\in\pkmI li$. Then $\ell^c(1-\ell b)=s$, so
$\ell^c\in\pkmI li$ as $1-\ell b$ is a unit; hence $\ell^{c+1}\pkpA\subseteq\pkmI li$.
(ii) Write $s=s_0+rb$, $s_0\in\pkmI0i$, and let $x\in(X_S)_i(\mathrm{std}_0)$ with
$\mathrm{pr}(x)=s_0\varepsilon_i(c)$. By Lemma~\ref{pkm:inputs}(E), $E^nx$ has orders $\mathrm{std}_0+n[a=0]$,
which at $a=0$ is $n-u\ge0=\mathrm{std}_1(0,u)$ on the jumps $u\le n$ at $i$, and equals
$\mathrm{std}_1$ for $a\ge1$; so $E^nx\in(X_S)_i(\mathrm{std}_1)$ (Lemma~\ref{pkm:le}(i)) with
$\mathrm{pr}(E^nx)=s_0E^n\varepsilon_i(c)$. Thus $sE^n=s_0E^n+rbE^n\in\pkmJ r1i$.
(iii) $\varphi^l$ is a ring automorphism of $\pkpA$, commutes with $\mathrm{pr}$, and maps
$\varepsilon_i(c)$ to $\varepsilon_{\sigma^l(i)}(c)$. By Lemma~\ref{pkm:inputs}(F), $x\mapsto\varphi^lx$ is a
bijection $(X_S)_i(\mathrm{std}_0)\to(X_S)_{\sigma^l(i)}(\mathrm{std}_l)$, and $\mathrm{pr}(x)=
s_0\varepsilon_i(c)$ iff $\mathrm{pr}(\varphi^lx)=\varphi^l(s_0)\varepsilon_{\sigma^l(i)}(c)$. Hence
$\varphi^l(\pkmI0i)=\pkmI l{\sigma^l(i)}$, and the claim follows as $\varphi^l(r\pkpA)=r\pkpA$.
(iv) Write $f=s+rb$ with $s\in\pkmI li$, $b\in\pkpA$, and put $t=\lambda_0(b)$. For $c\in C_S$ let
$x\in(X_S)_i(\mathrm{std}_l)$ with $\mathrm{pr}(x)=s\varepsilon_i(c)$, and $x_1=\mathrm{ret}(x)\in
\mathcal N_i(\mathrm{std}_l)$ (Lemma~\ref{pkm:inputs}(R)). By (H), $\varepsilon_i(c)\in\mathcal N(C_S)$, so
$\mathrm{pr}(x_1)=\mathrm{ret}(s\varepsilon_i(c))=\lambda_0(s)\varepsilon_i(c)$, and
$\lambda_0(s)=\lambda_0(f-rb)=f-rt$.
\end{proof}

\begin{lemma}[Operators]\label{pkm:op}
Let $f\in\mathfrak S\cap\pkmJ rli$, $g\in\Gamma'$, $\alpha,\beta,s\in\pkpA$ and $d\in\pkpA\cap\pkpW^\times$.
Assume: $d\mid gt-t$ for all $t\in\mathfrak S$; for every $x\in\mathcal N_i(\mathrm{std}_l)$ (for $X_S$),
$d^{-1}(gx-x)\in(X_S)_i(\mathrm{std}_l)$; and $d(s+\beta f)=\alpha(gf-f)$. Then $s\in\pkmJ rli$.
\end{lemma}

\begin{proof}
Take $t$ as in Lemma~\ref{pkm:J-basic}(iv) and write $gt-t=dv$. For $c\in C_S$ let $x\in\mathcal
N_i(\mathrm{std}_l)$ with $\mathrm{pr}(x)=(f-rt)\varepsilon_i(c)$, and $y=d^{-1}(gx-x)\in
(X_S)_i(\mathrm{std}_l)$. As $g$ fixes $\varepsilon_i(c)$ and commutes with $\mathrm{pr}$,
$\mathrm{pr}(gx-x)=(gf-f-r(gt-t))\varepsilon_i(c)$. Hence
\[
d\,\mathrm{pr}(\alpha y)=\alpha(gf-f)\varepsilon_i(c)-r\alpha dv\,\varepsilon_i(c)=d\,(s+\beta f-r\alpha
v)\,\varepsilon_i(c),
\]
and, $d$ being a unit of $\pkpW$, $\mathrm{pr}(\alpha y)=(s+\beta f-r\alpha v)\varepsilon_i(c)$ with $\alpha
y\in(X_S)_i(\mathrm{std}_l)$. So $s+\beta f-r\alpha v\in\pkmI li$, whence $s+\beta f\in\pkmJ rli$, and
$s\in\pkmJ rli$ since $\beta f\in\pkmJ rli$.
\end{proof}

\begin{corollary}[Derivations]\label{pkm:der}
Let $f\in\mathfrak S$, $g\in\Gamma'$, $\alpha,\beta,s\in\pkpA$. (i) If $f\in\pkmJ rli$ and
$z\tau(s+\beta f)=\alpha(gf-f)$, then $s\in\pkmJ rli$. (ii) If $f\in\pkmJ r1i$ and $z\tau E(s+\beta
f)=\alpha(gf-f)$, then $s\in\pkmJ r1i$.
\end{corollary}

\begin{proof}
(i) Apply Lemma~\ref{pkm:op} with $d=z\tau$: $z\tau\mid gt-t$ by Lemma~\ref{pkm:inputs}(C)(a); for
$x\in\mathcal N_i(\mathrm{std}_l)$, $gx-x\in(X_S)_i(\mathrm{std}_l)$ by (G), and then $(z\tau)^{-1}(gx-x)\in
(X_S)_i(\mathrm{std}_l)$ by (G). (ii) Apply Lemma~\ref{pkm:op} with $d=z\tau E$: the divisibility is
(C)(a); for $x\in\mathcal N_i(\mathrm{std}_1)$, $gx-x\in(X_S)_i(o_1)$ by (D), $(z\tau)^{-1}(gx-x)\in
(X_S)_i(o_1)$ by (G), and by (E) $E^{-1}(z\tau)^{-1}(gx-x)$ has orders $o_1-[a=0]$, which is
$0=\mathrm{std}_1$ at $a=0$ and $-u$ for $a\ge1$.
\end{proof}

\subsubsection*{Monomials in $\ell$, $z$, $E$}

\begin{lemma}[Coming down a level]\label{pkm:down-level}
If $\ell^cz^n\in\pkmJ r1{\sigma(i)}$, $\ell a\le n$ and $c+\pkmv(n!)\le\kappa+\pkmv((\ell a)!)$, then
$\ell^\kappa z^a\in\pkmJ r0i$.
\end{lemma}

\begin{proof}
Induction on $n\ge\ell a$. If $n=\ell a$: $\ell^cz^{\ell a}=\varphi(\ell^cz^a)$, so $\ell^cz^a\in\pkmJ r0i$
by Lemma~\ref{pkm:J-basic}(iii), and $c\le\kappa$. Step $n\to n+1$ (with $n\ge\ell a$): let
$f=\ell^cz^{n+1}\in\mathfrak S\cap\pkmJ r1{\sigma(i)}$, $t_1$ as in (C)(b) and $\alpha,\beta$ as in
(C)(c) for $n+1$. Then $t_1f-f=\ell^cz^{n+1}([\varepsilon]^{n+1}-1)$ and
\[
z\tau E\bigl(\ell^{c+\pkmv(n+1)}z^n+\beta f\bigr)=\ell^cz^{n+1}\,\tau E(\ell^{\pkmv(n+1)}+\beta z)=
\alpha(t_1f-f).
\]
By Corollary~\ref{pkm:der}(ii), $\ell^{c+\pkmv(n+1)}z^n\in\pkmJ r1{\sigma(i)}$, and
$c+\pkmv(n+1)+\pkmv(n!)=c+\pkmv((n+1)!)\le\kappa+\pkmv((\ell a)!)$; apply the induction hypothesis.
\end{proof}

\begin{lemma}[Trading $E$ for $z$]\label{pkm:trade}
If $\ell^cz^aE^n\in\pkmJ r1i$, $c+d_vn+\pkmv(n!)\le b$ and $a+n\le\kappa$, then $\ell^bz^\kappa\in\pkmJ r1i$.
\end{lemma}

\begin{proof}
Induction on $n$; $n=0$ is clear ($\ell^bz^\kappa$ is a multiple of $\ell^cz^a$). Step $n+1$: let
$f=\ell^cz^aE^{n+1}$, and $g,\omega,\upsilon,q$ as in (C)(d) for $n+1$; write $\upsilon
q-\ell^{\gamma}=Eb_1$ with $\gamma=\pkmv(n+1)+d_v$. By (C)(a), $gz=z(1+\tau Ev)$ for some $v$, and
$(1+\tau Ev)^a-1=\tau Evb_2$ for some $b_2$. Hence $gf=\ell^cz^a(1+\tau Evb_2)\omega^{n+1}E^{n+1}$, and,
using $\omega^{n+1}-1=z\tau q$,
\[
z\upsilon(gf-f)=\ell^cz^{a+1}E^{n+1}\tau\bigl(z\upsilon q+E\upsilon vb_2\omega^{n+1}\bigr)
=z\tau E\bigl(\ell^{c+\gamma}z^{a+1}E^n+\beta f\bigr),\quad \beta=zb_1+\upsilon vb_2\omega^{n+1},
\]
where we substituted $\upsilon q=\ell^\gamma+Eb_1$. By Corollary~\ref{pkm:der}(ii) (with $\alpha=z\upsilon$),
$\ell^{c+\gamma}z^{a+1}E^n\in\pkmJ r1i$. Since $(c+\gamma)+d_vn+\pkmv(n!)=c+d_v(n+1)+\pkmv((n+1)!)\le b$
and $(a+1)+n\le\kappa$, the induction hypothesis applies.
\end{proof}

\begin{lemma}[Lowering the power of $z$]\label{pkm:lower}
(i) If $\ell^cz^{n}\in\pkmJ r0i$ with $n\ge1$, then $\ell^{c+\pkmv(n)+1}z^{n-1}\in\pkmJ r0i$.
(ii) If $\ell^cz^n\in\pkmJ r0i$, $m\le n$ and $c+n+\pkmv(n!)\le\kappa+m+\pkmv(m!)$, then $\ell^\kappa
z^m\in\pkmJ r0i$.
\end{lemma}

\begin{proof}
(i) Let $\alpha,\beta$ be as in (C)(c) for $n$, $v=\pkmv(n)$, and write $E-\ell=ze'$ with $e'\in\mathfrak S$.
Then $\tau\bigl(\ell^{v+1}+z(\ell^ve'+\beta E)\bigr)=\tau E(\ell^v+\beta z)=\alpha([\varepsilon]^n-1)$.
With $f=\ell^cz^n$ and $t_1$ as in (C)(b), multiplying by $\ell^cz^n$ gives $z\tau(\ell^{c+v+1}z^{n-1}+
(\ell^ve'+\beta E)f)=\alpha(t_1f-f)$; apply Corollary~\ref{pkm:der}(i).
(ii) Induction on $n-m$ using (i): a step from $n+1$ to $n$ costs $\pkmv(n+1)+1$, and
$\pkmv((n+1)!)=\pkmv(n+1)+\pkmv(n!)$.
\end{proof}

\begin{lemma}[One letter back]\label{pkm:div}
If every jump at $\sigma(j)$ is $\le\omega$ and $\ell^cz^n\in\pkmJ r0{\sigma(j)}$, then
$\ell^{c+d_v\omega+\pkmv(\omega!)}z^{\lfloor(n+\omega)/\ell\rfloor}\in\pkmJ r0j$.
\end{lemma}

\begin{proof}
By Lemma~\ref{pkm:J-basic}(ii), $\ell^cz^nE^\omega\in\pkmJ r1{\sigma(j)}$. By Lemma~\ref{pkm:trade} (with
$a=n$, $n=\omega$, $\kappa=n+\omega$, $b=c+d_v\omega+\pkmv(\omega!)$), $\ell^bz^{n+\omega}\in\pkmJ
r1{\sigma(j)}$. Put $a'=\lfloor(n+\omega)/\ell\rfloor$; no multiple of $\ell$ lies in $(\ell a',n+\omega]$,
so $\pkmv((n+\omega)!)=\pkmv((\ell a')!)$, and Lemma~\ref{pkm:down-level} (with $c=\kappa=b$) gives
$\ell^bz^{a'}\in\pkmJ r0j$.
\end{proof}

\begin{lemma}[Back through $J$ letters]\label{pkm:back}
Let $z^a\in\pkmJ r0j$ for every letter $j$, and let every jump at $\sigma^s(i)$, $1\le s\le J$, be
$\le\omega$. Then $\ell^cz^n\in\pkmJ r0i$ for some $c\le J(d_v\omega+\pkmv(\omega!))$ and $n\le
a/\ell^J+\omega/(\ell-1)$.
\end{lemma}

\begin{proof}
Induction on $J$; $J=0$: $c=0$, $n=a$. Step: by the induction hypothesis at the letter $\sigma(i)$
(whose next $J$ letters are $\sigma^s(i)$, $2\le s\le J+1$) there are $c,n$ for $\sigma(i)$; by
Lemma~\ref{pkm:div} with $j=i$, $\ell^{c+d_v\omega+\pkmv(\omega!)}z^{\lfloor(n+\omega)/\ell\rfloor}\in\pkmJ
r0i$, and $\lfloor(n+\omega)/\ell\rfloor\le(a/\ell^J+\omega/(\ell-1)+\omega)/\ell=a/\ell^{J+1}+
\omega/(\ell-1)$.
\end{proof}

\subsubsection*{Proofs of the killers}

\begin{proof}[Proof of Lemma~\ref{pkm:zpow}]
\emph{Some power exists.} Let $H$ be as in Lemma~\ref{pkm:inputs}(H) for $X_S$. For $c\in C_S$,
$x=e_i\tau^H(1\otimes(0,c))\in(X_S)_i(\mathrm{std}_0)$ and $\mathrm{pr}(x)=\tau^H\varepsilon_i(c)$; so
$\tau^H\in\pkmI0i$ for all $i$. Also $\ell^{\kappa+1}\in\pkmJ{\ell^{\kappa+1}}0i$. In $\pkpA/\ell\pkpA=
O^\flat$ the images of $\tau$ and $z$ are $\varepsilon^{1/\ell}-1$ and $\varpi^\flat$, both of positive
valuation; so $z^M\in\tau\pkpA+\ell\pkpA$ for some $M$, and $z^{M(H+\kappa+1)}\in(\tau,\ell)^{H+\kappa+1}
\subseteq\bigcap_i\pkmJ{\ell^{\kappa+1}}0i$.

\emph{The bound.} Let $a$ be the least integer with $z^a\in\pkmJ{\ell^{\kappa+1}}0i$ for every $i$. Write
$z^e=u_1E+v_1\ell$ with $u_1,v_1\in\mathfrak S$, and $K=\omega+\kappa$. For each $i$,
$z^aE^\omega\in\pkmJ{\ell^{\kappa+1}}1{\sigma(i)}$ by Lemma~\ref{pkm:J-basic}(ii), and
$z^{a+eK}=z^a(u_1E+v_1\ell)^K$ is a combination of the $z^aE^j\ell^{K-j}$; for $j\ge\omega$ these are
multiples of $z^aE^\omega$, and for $j<\omega$ multiples of $\ell^{\kappa+1}$. So $z^{a+eK}\in\pkmJ
{\ell^{\kappa+1}}1{\sigma(i)}$, and by Lemma~\ref{pkm:down-level} (with $c=\kappa=0$, $a'=\lfloor(a+eK)/\ell
\rfloor$, as in Lemma~\ref{pkm:div}) $z^{a'}\in\pkmJ{\ell^{\kappa+1}}0i$ for all $i$. By minimality $a\le a'$,
so $\ell a\le a+eK$, i.e.\ $(\ell-1)a\le e(\omega+\kappa)$.
\end{proof}

\begin{proof}[Proof of Proposition~\ref{pkm:kill-small}]
Lemma~\ref{pkm:zpow} with $\kappa=0$ gives $a$ with $(\ell-1)a\le e\omega\le\ell-2$, so $a=0$ and
$1\in\pkmJ\ell0i$; by Lemma~\ref{pkm:J-basic}(i) with $c=0$, $1\in\pkmI0i$.
\end{proof}

\begin{proof}[Proof of Proposition~\ref{pkm:kill-unram}]
Only $\ell\ge3$ is used: it gives $\omega=(\ell-3)K\ge0$ and $\ell-2>0$. Let $K=s+t$. Lemma~\ref{pkm:zpow} (with $\omega=(\ell-3)K$, $\kappa=K$, $e=1$) gives $b$ with $(\ell-1)b\le
(\ell-2)K$ and $z^b\in\pkmJ{\ell^{K+1}}0i$ for all $i$. By Legendre's formula $(\ell-1)\pkmv(b!)\le b$, so
$(\ell-2)(b+\pkmv(b!))\le(\ell-2)b+(\ell-1)\pkmv(b!)\le(\ell-1)b\le(\ell-2)K$, i.e.\ $b+\pkmv(b!)\le K$.
For every $m\ge0$, $\ell^{\max(K-m,0)}z^m\in\pkmJ{\ell^{K+1}}0i$: from $z^{\max(b,m)}$ (a multiple of $z^b$)
apply Lemma~\ref{pkm:lower}(ii); the cost condition $\max(b,m)+\pkmv(\max(b,m)!)\le\max(K-m,0)+m+
\pkmv(m!)$ is trivial if $m\ge b$, and follows from $b+\pkmv(b!)\le K$ if $m<b$. With $m=0$,
$\ell^K\in\pkmJ{\ell^{K+1}}0i$, so $\pkmJ{\ell^{K+1}}0i=\pkmI0i$ by Lemma~\ref{pkm:J-basic}(i). Finally
$\ell^sE^t=\ell^s(z+\ell)^t=\sum_{m\le t}\binom tm\ell^{K-m}z^m\in\pkmI0i$.
\end{proof}

\begin{lemma}[Arithmetic]\label{pkm:arith}
Let $J,d,e,\omega,\omega',\kappa,a,c,n\ge0$ with $c\le J(d\omega'+\pkmv(\omega'!))$, $n\le
a/\ell^J+\omega'/(\ell-1)$, $(\ell-1)a\le e(\omega+\kappa)$, $8e\le\ell^{J+1}$ and $2Jd\omega'+
(2J+4)\omega'/(\ell-1)+8e\omega/\ell^{J+1}\le\kappa$. Then $c+n+\pkmv(n!)\le\kappa$.
\end{lemma}

\begin{proof}
By Legendre, $\pkmv(\omega'!)\le\omega'/(\ell-1)$ and $\pkmv(n!)\le n$; so $c+n+\pkmv(n!)\le
J(d\omega'+\omega'/(\ell-1))+2n$. As $\ell\le2(\ell-1)$, $\ell a\le2e(\omega+\kappa)$, so
$a/\ell^J=\ell a/\ell^{J+1}\le2e\omega/\ell^{J+1}+(8e/\ell^{J+1})\kappa/4\le2e\omega/\ell^{J+1}+\kappa/4$.
Hence $c+n+\pkmv(n!)\le\tfrac12\bigl(2Jd\omega'+(2J+4)\omega'/(\ell-1)+8e\omega/\ell^{J+1}\bigr)+
\kappa/2\le\kappa$.
\end{proof}

\begin{proof}[Proof of Proposition~\ref{pkm:kill-J}]
Lemma~\ref{pkm:zpow} gives $a$ with $(\ell-1)a\le e(\omega+\kappa)$ and $z^a\in\pkmJ{\ell^{\kappa+1}}0j$ for
all $j$. Lemma~\ref{pkm:back} (with $\omega'$) gives $\ell^cz^n\in\pkmJ{\ell^{\kappa+1}}0i$ with $c\le
J(d_v\omega'+\pkmv(\omega'!))$ and $n\le a/\ell^J+\omega'/(\ell-1)$. By Lemma~\ref{pkm:arith},
$c+n+\pkmv(n!)\le\kappa$, so Lemma~\ref{pkm:lower}(ii) with $m=0$ gives $\ell^\kappa\in\pkmJ{\ell^{\kappa+1}}0i$,
and Lemma~\ref{pkm:J-basic}(i) gives $\ell^\kappa\in\pkmI0i$.
\end{proof}

\emph{Where the inputs are proved.} Everything else in \S\ref{pkm:sec} is proved from
Lemma~\ref{pkm:inputs}, Lemma~\ref{pkp:orders-add} and Lemma~\ref{pkp:tau}(i), (iv), (v). The items of
Lemma~\ref{pkm:inputs}: (M), first two assertions, Lemma~\ref{pkn:M}, third assertion,
Proposition~\ref{pkn:frame}(iv) with Theorem~\ref{pkf:main}; (E) Lemma~\ref{pkn:E}; (G), first assertion,
Lemma~\ref{pkl:G1}, second assertion and (G$_0$), Proposition~\ref{pkl:G0}; ($\Phi$)
Corollary~\ref{pkl:Phi}; (F) Corollary~\ref{pkl:F}; (H) Corollary~\ref{pkl:H1}, Proposition~\ref{pkl:H2}
and Corollary~\ref{pkl:H3}; (R) Proposition~\ref{pkl:R}; (D) Proposition~\ref{pkl:D}; (C)(a)
Proposition~\ref{pkn:Ca}, (C)(b) \S\ref{pkn:ss-ME}, (C)(c) Proposition~\ref{pkn:Cc}, (C)(d)
Proposition~\ref{pkl:Cd}. Lemma~\ref{pkp:orders-add} is Lemma~\ref{pkt:add}, and
Lemma~\ref{pkp:tau}(i), (iv), (v) are Proposition~\ref{pkt:tau-i}, Lemmas~\ref{pkt:tau-iv},
\ref{pkt:tau-v}.


%% ===== proofs/C-kisin-inputs.tex =====
% proofs/C-kisin-inputs.tex
% Inputs of the theory of the modules N for proofs/C-kisin-modules.tex. Labels start with pkn:.
% Conventions and macros as in C-kisin-pairing.tex and C-kisin-modules.tex.
\providecommand{\auO}{\mathcal O}
\providecommand{\auQlb}{\overline{\Q}_\ell}
\providecommand{\panot}[1]{\par\smallskip{\raggedright\noindent\textbf{[Not proved here:} #1\textbf{]}\par}\smallskip}
\providecommand{\pkpA}{A_{\mathrm{inf}}}
\providecommand{\pkpW}{W(C^\flat)}
\providecommand{\pkpm}{W(\mathfrak m^\flat)}
\providecommand{\pkpAmb}{\mathbb A}
% macros of this file (prefix \pkn)
\providecommand{\pknMinf}{\mathfrak M_{\mathrm{inf}}}
\providecommand{\pknK}{\mathcal K_0}
\providecommand{\pknBp}{B^+_{\mathrm{dR}}}
\providecommand{\pknB}{B_{\mathrm{dR}}}

\section{The modules $\mathcal N$: definition and first inputs}\label{pkn:sec}

This section continues \S\ref{pkp:sec} and \S\ref{pkm:sec}, whose conventions and labels we use. It
fixes the definition of the modules $\mathcal N$ used in Definition~\ref{pkp:def-orders}
(\S\ref{pkn:ss-def}), proves that a distinguished element
generates $\ker\theta$ (which completes the proof of Lemma~\ref{pkp:chart-exists}), and proves the
items (M) (first two assertions), (E), (C)(a), (C)(b) and (C)(c) of Lemma~\ref{pkm:inputs}.

\emph{Dependencies.} Proved in full: Lemma~\ref{pkn:dist} (which completes the proof of
Lemma~\ref{pkp:chart-exists}), Lemmas~\ref{pkn:mu}, \ref{pkn:BdR}, \ref{pkn:M}, \ref{pkn:E} and
Propositions~\ref{pkn:Ca}, \ref{pkn:Cc} (items (M) first two assertions, (E), (C)(a), (C)(b), (C)(c) of
Lemma~\ref{pkm:inputs}), modulo published facts on $\pknBp$ (that $t$ is a uniformiser,
Fontaine \cite[1.5.2, 1.5.4]{pkfFontaineII}; that $\pkpA\to\pknBp$ is injective, \cite[4.1.1]{pkfFontaineII}) and Tate's theorem (through Lemma~\ref{pkm:BdR}). Lemma~\ref{pkn:frame-id} is proved for any given frame. Proved modulo the existence of a frame (Definition~\ref{pkn:def-frame}), which is
Theorem~\ref{pkf:main}: Proposition~\ref{pkn:frame} (item (M), third assertion, and the first assertion
of (H)). The other items are proved in \S\S\ref{pkf:sec}--\ref{pkl:sec}; see the paragraph at the end
of the section.

\subsection{The modules $\mathcal N$}\label{pkn:ss-def}

Fix a Kummer chart (Definition~\ref{pkp:def-chart}). Recall that $\pknBp$ is the completion of
$\pkpA[1/\ell]$ for the $\ker(\theta[1/\ell])$-adic topology, $\pknB$ its fraction field, and
$\mathrm{Fil}^n=t^n\pknBp$ ($n\in\Z$); $\pkpA\to\pknBp$ is the natural map. The group $\Gamma$ acts on
$\pknB$ through $j$, compatibly with this map.

\begin{definition}[The modules $\mathcal N$]\label{pkn:def-N}
Let $L$ be a lattice with jump sets (Definition~\ref{pkp:def-orders}, first sentence). A
\emph{form} on $L$ is an $\pkpA$-linear map $\mu\colon\pkpAmb(L)\to\pknB$ with $\mu(1\otimes gy)=
g\,\mu(1\otimes y)$ for $g\in\Gamma'$, $y\in L$; a form \emph{of type $u$ at $i$} moreover satisfies
$\mu(e_i(1\otimes L))\subseteq\mathrm{Fil}^{-u}$. Put
\begin{align*}
\pknMinf(L)&=\{x\in\pkpAmb(L):\ \tau^nx\in\pkpA\otimes L\text{ for some }n\ge0,\ \text{and }
\mu(\varphi^{n}x)\in\pknBp\text{ for all }n\ge1\text{ and all forms }\mu\},\\
\pkpA\mathcal N(L)&=\{x\in\pkpAmb(L):\ rx\in\pknMinf(L)\text{ for some finite product }r\text{ of
elements }p_a,\ a\ge0\},\\
\mathcal N(L)&=\{x\in\pkpA\mathcal N(L):\ \mu(x)\in\pknK\text{ for all forms }\mu\},
\end{align*}
where $\pkpA\otimes L$ is the $\pkpA$-span of the $1\otimes y$, and $\pknK\subseteq\pknB$ is the
subfield generated over $K'=\auQlb^{\,j(\Gamma')}$ by the images of the products $bs$, $b\in\mathfrak
S_B$, $s\in\mathfrak S$, with $\mathfrak S_B\subseteq B^+$ the bicommutant of the chart (Definition~\ref{pkf:def-SB}; a subring
containing $1$, mapped to $\pknBp$ by $B^+\to\pknBp$, see \S\ref{pkn:ss-frame}; we use only that $\pknK$ is a subfield of $\pknB$ containing the image of $\mathfrak
S$). The orders of Definition~\ref{pkp:def-orders} are defined with this $\pkpA\mathcal N(L)$: for
orders $o$, $L_i(o)$ is the set of $x\in\pkpA\mathcal N(L)\cap\pkpAmb(L)_i$ with $\mu(\varphi^{-a}x)\in
\mathrm{Fil}^{o(a,u)}$ for all $a\ge0$, $u\in U(i^{(-a)})$ and forms $\mu$ of type $u$ at $i^{(-a)}$.
\end{definition}

The notation $\pkpA\mathcal N$ is justified by Lemma~\ref{pkm:inputs}(M) (third assertion, Proposition~\ref{pkn:frame}(iv) with
Theorem~\ref{pkf:main}). No comparison with Breuil--Kisin modules is used anywhere in
\S\S\ref{pkp:sec}--\ref{pkn:sec}.

With this definition, two statements of \S\ref{pkp:sec} are to be read as follows.

\emph{Killing the cokernel} (Definition~\ref{pkp:def-points}): $\ell^\kappa E^t$ kills the cokernel
at $i$ for the side $S$ if for every $c\in C_S$ there is $n\in(X_S)_i(\mathrm{std}_0)$ with
$\mathrm{pr}(n)=\ell^\kappa E^te_i(1\otimes c)$; $n$ lies in $\pkpA\mathcal N(X_S)$, not necessarily in
$\mathcal N(X_S)$.

\emph{Reading of Lemma~\ref{pkp:tau}.} (i) For the lattice $\auO(1)$ (jump set $\{-1\}$), every $\xi\in
\auO(1)_i(\mathrm{std}_0)$ has $\tau\xi\in\pkpA\otimes\auO(1)$. (ii) $z\in\pkpm$, and the image of
$z(\pkpA\otimes\auO(1))$ in $\pkpM$ lies in $\pkpsm$ (Lemma~\ref{pkp:AS}(vii)). (iii) If
$\ell^\kappa E^t$ kills the cokernel at $i$ for a side, then for every $c$ there is $\nu\in
(X_S)_i(\rho_{-t})$ with $\mathrm{pr}(\nu)=\ell^\kappa e_i(1\otimes c)$ (Lemma~\ref{pkm:tau-iii}).
(iv) If $x\in\mathcal N$ lies at $i$, is deep for $o'$ and $\mathrm{pr}(x)$ is $\Gamma'$-invariant, then
for $g\in\Gamma'$, $(z\tau)^{-1}(gx-x)\in L_i(o')$ and lies over $0$. (v) If $x\in L_{i^{(-l)}}(o)$
then $\varphi^lx\in L_i(\mathrm{sh}_{l,0}(o))$; if $x$ is a point at $i^{(-l)}$ with orders $o$ and
depth $o'$, then $\varphi^lx$ is a point at $i$ over the same constant with orders $\mathrm{sh}_{l,0}(o)$
and depth $\mathrm{sh}_{l,1}(o')$. Assertions (i), (iv), (v) are proved in \S\ref{pkt:sec}.

\subsection{Distinguished elements}\label{pkn:ss-dist}

For $x\in\pkpA=W(O^\flat)$ write $x=\sum_{n\ge0}\ell^n[x_n^{1/\ell^n}]$ with $x_n\in O^\flat$ the Witt
coordinates ($O^\flat$ is perfect); so $x\equiv[x_0]\pmod{\ell\pkpA}$, $\pkpA/\ell=O^\flat$ via $x\mapsto
x_0$, and $\theta(x)=\sum_n\ell^n(x_n^{1/\ell^n})^\sharp$, where $b\mapsto b^\sharp$ is the
multiplicative map $O^\flat\to O_C$, $b^\sharp=\lim b_n^{\ell^n}$ for lifts $b_n$ of the coordinates.
Let $v^\flat(b)=v_\ell(b^\sharp)$, a valuation on $O^\flat$ with $O^\flat$ its valuation ring. An element
$x$ is a unit of $\pkpA$ iff $x_0\in O^{\flat\times}$ (if $x_0$ is a unit, $x=[x_0](1-\ell y)$ with
$y\in\pkpA$, and $1-\ell y$ is invertible by $\ell$-adic completeness); equivalently iff its image in
$W(\overline{\mathbb F}_\ell)$ under $O^\flat\to O^\flat/\mathfrak m^\flat=\overline{\mathbb F}_\ell$ is a unit.

\begin{definition}\label{pkn:def-dist}
$\xi\in\pkpA$ is \emph{distinguished} if $\xi_0\in\mathfrak m^\flat$ and its image in
$W(\overline{\mathbb F}_\ell)$ is $\ell$ times a unit; equivalently $\xi=[\xi_0]+\ell\xi'$ with
$\xi_0\in\mathfrak m^\flat$, $\xi'\in\pkpA^\times$.
\end{definition}

(Equivalence: the image of $[\xi_0]$ in $W(\overline{\mathbb F}_\ell)$ is $0$, so the image of $\xi$
is $\ell$ times that of $\xi'$, and $W(\overline{\mathbb F}_\ell)$ is torsion free.)

\begin{lemma}\label{pkn:dist}
Let $\xi$ be distinguished with $\theta(\xi)=0$. Then:
\begin{enumerate}
\item[(i)] $v^\flat(\xi_0)=1$ ($=v_\ell(\ell)$);
\item[(ii)] $\ker\theta=\xi\pkpA$;
\item[(iii)] for $c\in\Z\smallsetminus\{0\}$, $\varphi^c(\xi)$ is distinguished and
$\theta(\varphi^c\xi)\ne0$;
\item[(iv)] if $\xi''$ is another distinguished element of $\ker\theta$, then $\xi''=\xi u$ with
$u\in\pkpA^\times$.
\end{enumerate}
\end{lemma}

\begin{proof}
(i) $0=\theta(\xi)=\xi_0^\sharp+\ell\theta(\xi')$, and $\theta(\xi')\in O_C^\times$ because
$\theta(\xi')\equiv(\xi'_0)^\sharp\pmod\ell$ with $\xi'_0$ a unit. Hence $v_\ell(\xi_0^\sharp)=1$.

(ii) Clearly $\xi\pkpA\subseteq\ker\theta$. \emph{Step 1: if $x\in\pkpA$ and $\theta(x)\in\ell O_C$,
then $x\in\xi\pkpA+\ell\pkpA$.} Indeed $\theta(x)\equiv x_0^\sharp\pmod\ell$, so $v^\flat(x_0)\ge1=
v^\flat(\xi_0)$ and $x_0=\xi_0c$ with $c\in O^\flat$. Then $x-\xi[c]\equiv[x_0]-[\xi_0][c]=0
\pmod{\ell\pkpA}$. \emph{Step 2.} Let $x\in\ker\theta$. Put $x^{(0)}=x$. Given $x^{(k)}\in\ker\theta$,
Step 1 gives $y_k,x^{(k+1)}\in\pkpA$ with $x^{(k)}=\xi y_k+\ell x^{(k+1)}$; then
$\ell\theta(x^{(k+1)})=0$, and as $O_C$ is torsion free, $x^{(k+1)}\in\ker\theta$. Thus
$x=\xi(y_0+\ell y_1+\dots+\ell^ky_k)+\ell^{k+1}x^{(k+1)}$ for all $k$. As $\pkpA$ is $\ell$-adically
complete and separated, $y=\sum_k\ell^ky_k$ converges and $x=\xi y$.

(iii) $\varphi$ is the automorphism of $\pkpA$ induced by the $\ell$-th power map on coordinates, so
$\varphi^c(\xi)_0=\xi_0^{\ell^c}\in\mathfrak m^\flat$, and $\varphi^c$ induces the Frobenius power on
$W(\overline{\mathbb F}_\ell)$, which fixes $\ell$ and maps units to units: $\varphi^c\xi$ is
distinguished. If $\theta(\varphi^c\xi)=0$, (i) applied to $\varphi^c\xi$ would give
$v^\flat(\xi_0^{\ell^c})=1$, while $v^\flat(\xi_0^{\ell^c})=\ell^c\ne1$.

(iv) By (ii), $\xi''=\xi u$ with $u\in\pkpA$. In $W(\overline{\mathbb F}_\ell)$: $\ell w''=\ell w\bar u$
with $w,w''$ units, so $\bar u$ is a unit, hence $u\in\pkpA^\times$.
\end{proof}

\begin{corollary}[Completion of Lemma~\ref{pkp:chart-exists}]\label{pkn:E-gen}
In Definition~\ref{pkp:def-chart} and the proof of Lemma~\ref{pkp:chart-exists}, $E=z^eA+\ell$ is
distinguished and generates $\ker\theta$; for every $a\in\Z$, $\varphi^a(E)$ generates
$\ker(\theta\circ\varphi^{-a})$, and $\theta(\varphi^c E)\ne0$ for $c\ne0$.
\end{corollary}

\begin{proof}
$E_0=(\varpi^\flat)^e\bar A_0\in\mathfrak m^\flat$ (the zeroth coordinate is multiplicative and additive
modulo $\ell$, and $z=[\varpi^\flat]$), and the image of $E$ in $W(\overline{\mathbb F}_\ell)$ is $\ell$
because $z\mapsto[\bar\varpi^\flat]=0$. As $\theta(E)=0$ (proof of Lemma~\ref{pkp:chart-exists}),
Lemma~\ref{pkn:dist}(ii),(iii) apply; $\varphi^{-a}$ is an automorphism carrying $\varphi^a(E)\pkpA$ to
$E\pkpA$.
\end{proof}

This removes the appeal to ``the standard criterion'' in the proof of Lemma~\ref{pkp:chart-exists},
which is therefore proved in full.

\begin{lemma}[$\tau E$ and $\mu$]\label{pkn:mu}
Let $\mu_\varepsilon=[\varepsilon]-1=\varphi(\tau)$ and $\xi_\varepsilon=\sum_{k=0}^{\ell-1}[\varepsilon^{k/\ell}]$,
so that $\mu_\varepsilon=\tau\xi_\varepsilon$. Then $\xi_\varepsilon$ is distinguished, lies in
$\ker\theta$, and $\xi_\varepsilon=Eu$ with $u\in\pkpA^\times$; hence $\mu_\varepsilon=\tau Eu$.
Moreover, if $x\in\pkpA$ satisfies $\theta(\varphi^jx)=0$ for all $j\ge0$, then $\mu_\varepsilon\mid x$.
\end{lemma}

\begin{proof}
$\theta(\xi_\varepsilon)=\sum_{k<\ell}\zeta_\ell^k=0$. Since $x\mapsto x_0$ is the reduction $\pkpA\to\pkpA/\ell=O^\flat$, a ring map, its zeroth coordinate is
$\sum_k\varepsilon^{k/\ell}=(\varepsilon-1)/(\varepsilon^{1/\ell}-1)$, which lies in $\mathfrak m^\flat$ because
$v^\flat(\varepsilon-1)=\ell/(\ell-1)>v^\flat(\varepsilon^{1/\ell}-1)=1/(\ell-1)$ (as
$(\varepsilon^{1/\ell^n}-1)^\sharp=\lim_m(\zeta_{\ell^{m+n}}-1)^{\ell^m}$ has valuation
$\ell^{-n}\cdot\ell/(\ell-1)$). Its image in $W(\overline{\mathbb F}_\ell)$ is $\sum_k1=\ell$. By
Lemma~\ref{pkn:dist}(iv) and Corollary~\ref{pkn:E-gen}, $\xi_\varepsilon=Eu$ with $u$ a unit.

For the last claim, put $\xi_j=\varphi^{-j}(\xi_\varepsilon)$, a distinguished generator of
$\ker(\theta\circ\varphi^j)$ (Lemma~\ref{pkn:dist}(iii) and the argument of Corollary~\ref{pkn:E-gen}).
\emph{Claim: $\prod_{j=0}^n\xi_j$ divides $x$.} By induction on $n$: if $x=\xi_0\cdots\xi_{n-1}y$, apply
$\theta\circ\varphi^n$, whose kernel is $\xi_n\pkpA$: $\theta\varphi^n(\xi_j)=\theta(\varphi^{n-j}\xi_\varepsilon)
\ne0$ for $j<n$ (Lemma~\ref{pkn:dist}(iii)) and $O_C$ is a domain, so $\theta\varphi^n(y)=0$ and
$\xi_n\mid y$. Since $\varphi^{-j}(\mu_\varepsilon)=\varphi^{-j-1}(\mu_\varepsilon)\xi_j$, the product is
$\mu_\varepsilon/\varphi^{-n-1}(\mu_\varepsilon)$, so $\varphi^{-n-1}(\mu_\varepsilon)x\in\mu_\varepsilon\pkpA$.
Reduce modulo $\ell$: in $O^\flat$, $(\varepsilon-1)$ divides $(\varepsilon^{1/\ell^{n+1}}-1)x_0$, whence
$v^\flat(x_0)\ge\frac{\ell}{\ell-1}(1-\ell^{-n-1})$ for all $n$, so $v^\flat(x_0)\ge v^\flat(\varepsilon-1)$
and $x_0=(\varepsilon-1)c$. Then $x=\mu_\varepsilon[c]+\ell x'$ for some $x'$, and $x'$ again satisfies the
hypothesis: $\theta\varphi^j(\mu_\varepsilon)=\theta([\varepsilon^{\ell^j}])-1=1-1=0$, so $\ell\,\theta\varphi^j(x')=0$, and $O_C$ is torsion free. Iterating and using $\ell$-adic
completeness as in Lemma~\ref{pkn:dist}(ii), $x\in\mu_\varepsilon\pkpA$.
\end{proof}

\begin{lemma}[The primes $p_a$ in $\pknB$]\label{pkn:BdR}
The image of $E$ in $\pknBp$ generates $\mathrm{Fil}^1$, so $E^n\pknBp=\mathrm{Fil}^n$ ($n\ge0$) and
$E^n\mathrm{Fil}^m=\mathrm{Fil}^{m+n}$ ($n,m\in\Z$, $E$ being invertible in $\pknB$). For $c\in\Z\smallsetminus
\{0\}$ the image of $\varphi^c(E)$ is a unit of $\pknBp$. Consequently, for $a,b\ge0$, $n\in\Z$ and
$y\in\pknB$: $\varphi^{-a}(p_b)^n y\in\mathrm{Fil}^{m+n[a=b]}$ iff $y\in\mathrm{Fil}^m$.
\end{lemma}

\begin{proof}
By Corollary~\ref{pkn:E-gen}, $\ker(\theta[1/\ell])=E\pkpA[1/\ell]$, so $\pknBp$ is the $E$-adic
completion of $\pkpA[1/\ell]$ and its maximal ideal $\ker(\theta\colon\pknBp\to C)$ is $E\pknBp$. By
Fontaine's theorem, $\pknBp$ is a complete discrete valuation ring with uniformiser $t$ \cite[1.5.2, 1.5.4]{pkfFontaineII}, so $E$ and
$t$ are both uniformisers and $E^n\pknBp=t^n\pknBp=\mathrm{Fil}^n$. If $c\ne0$, $\theta(\varphi^cE)\ne0$
(Corollary~\ref{pkn:E-gen}), so $\varphi^cE\notin\ker\theta$ and its image is a unit of the local ring
$\pknBp$. Finally $\varphi^{-a}(p_b)=\varphi^{b-a}(E)$: it is $E$ if $a=b$ and a unit otherwise.
\end{proof}

\subsection{Items (M), (E) and (C)(b)}\label{pkn:ss-ME}

\begin{lemma}[Lemma~\ref{pkm:inputs}(M), first two assertions]\label{pkn:M}
Let $L$ be a lattice with jump sets. (i) $\pknMinf(L)$ and $\pkpA\mathcal N(L)$ are $\pkpA$-submodules of
$\pkpAmb(L)$, and $\mathcal N(L)$ is an $\mathfrak S$-submodule. (ii) Every $x\in\pkpA\mathcal N(L)$ has
$rx\in\pkpA\otimes L$ for some product $r$ of a power of $\tau$ and of elements $p_b$. (iii)
$\pkpA\mathcal N(L)$ and $\mathcal N(L)$ are stable under multiplication by $p_b$ and by $p_b^{-1}\in
\pkpW^\times$ ($b\ge0$).
\end{lemma}

\begin{proof}
(i) For $b\in\pkpA$ and $x,x'\in\pknMinf$ with $\tau^nx,\tau^{n'}x'\in\pkpA\otimes L$:
$\tau^{n+n'}(bx+x')\in\pkpA\otimes L$; and $\mu(\varphi^k(bx))=\varphi^k(b)\mu(\varphi^kx)\in\pknBp$, as
$\varphi^k$ is semilinear on $\pkpAmb(L)=\pkpW\otimes L$ and $\mu$ is $\pkpA$-linear. For
$\pkpA\mathcal N$ use products $rr'$ of the two multipliers. For $\mathcal N$ and $s\in\mathfrak S$:
$\mu(sx)=s\mu(x)\in\pknK$ since $\pknK$ is a field containing the image of $\mathfrak S$.
(ii) If $rx\in\pknMinf$ then $\tau^nrx\in\pkpA\otimes L$ for some $n$. (iii) $p_b\in\mathfrak
S\subseteq\pkpA$, so stability under $p_b$ is (i). The $p_b$ are units of $\pkpW$ (their zeroth
coordinate $\varphi^b(E)_0=E_0^{\ell^b}$ is a nonzero element of the field $C^\flat$). If $rx\in\pknMinf$,
then $(rp_b)(p_b^{-1}x)=rx$, so $p_b^{-1}x\in\pkpA\mathcal N$. If moreover $x\in\mathcal N$, then for every
form $\mu$, $p_b\,\mu(p_b^{-1}x)=\mu(x)\in\pknK$, and the image of $p_b$ in $\pknB$ is a nonzero element
of $\pknK$ (Lemma~\ref{pkn:BdR}), so $\mu(p_b^{-1}x)\in\pknK$.
\end{proof}

\begin{lemma}[Lemma~\ref{pkm:inputs}(E)]\label{pkn:E}
For $b\ge0$, $n\in\Z$ and $x\in L_i(o)$, $p_b^nx\in L_i(o+n[a=b])$, where $(o+n[a=b])(a,u)=
o(a,u)+n[a=b]$.
\end{lemma}

\begin{proof}
$p_b^nx\in\pkpA\mathcal N(L)$ by Lemma~\ref{pkn:M}(iii), and it lies in $\pkpAmb(L)_i$ because $p_b$ acts
on the first tensor factor and commutes with $1\otimes\zeta$ and with the scalars $[i']$. Let $a\ge0$,
$u\in U(i^{(-a)})$ and $\mu$ a form of type $u$ at $i^{(-a)}$. Since $\varphi^{-a}$ is semilinear and
$\mu$ is $\pkpA$-linear, $\varphi^{-a}(p_b)^{|n|}$ (an element of $\pkpA$) satisfies, for $n\ge0$,
$\mu(\varphi^{-a}(p_b^nx))=\varphi^{-a}(p_b)^n\mu(\varphi^{-a}x)$, and for $n<0$,
$\varphi^{-a}(p_b)^{-n}\mu(\varphi^{-a}(p_b^nx))=\mu(\varphi^{-a}x)$; in both cases
$\mu(\varphi^{-a}(p_b^nx))=\varphi^{-a}(p_b)^n\mu(\varphi^{-a}x)$ in the field $\pknB$ (the image of
$\varphi^{-a}(p_b)$ is nonzero by Lemma~\ref{pkn:BdR}). By Lemma~\ref{pkn:BdR} this lies in
$\mathrm{Fil}^{o(a,u)+n[a=b]}$.
\end{proof}

\emph{Item (C)(b) of Lemma~\ref{pkm:inputs}} is the case $l=1$ of the last condition in
Definition~\ref{pkp:def-chart}.

\subsection{Item (C)(a)}\label{pkn:ss-Ca}

\begin{proposition}[Lemma~\ref{pkm:inputs}(C)(a)]\label{pkn:Ca}
For $g\in\Gamma'$ and $s\in\mathfrak S$, $z\tau E$ divides $gs-s$ in $\pkpA$.
\end{proposition}

\begin{proof}
\emph{Step 1: $z\mu_\varepsilon\mid gz-z$.} $g$ acts on $\pkpA=W(O^\flat)$ through its action on
$O^\flat$, so $gz=[g\varpi^\flat]$. As $\theta$ is $\Gamma$-equivariant and $j(\Gamma')$ fixes
$\theta(z)$, $(g\varpi^\flat)^\sharp=g(\varpi^{\flat\sharp})=\varpi^{\flat\sharp}=\pi\ne0$. So $\eta=
g\varpi^\flat/\varpi^\flat\in C^\flat$ has $\eta^\sharp=1$ ($\sharp$ is multiplicative), hence
$v^\flat(\eta)=0$, $\eta\in O^{\flat\times}$, and $gz-z=z([\eta]-1)$. For $j\ge0$,
$\theta(\varphi^j([\eta]-1))=(\eta^{\ell^j})^\sharp-1=0$, so $\mu_\varepsilon\mid[\eta]-1$ by
Lemma~\ref{pkn:mu}.

\emph{Step 2.} By Lemma~\ref{pkn:mu}, $\mu_\varepsilon=\tau Eu$ with $u$ a unit, so $z\tau E\mid gz-z$.
Write $s=\sum_{k\ge0}\lambda_kz^k$ with $\lambda_k\in\Lambda$ (Definition~\ref{pkp:def-chart}). $g$ fixes
$\Lambda$ and is continuous, so $gs-s=\sum_{k\ge1}\lambda_k((gz)^k-z^k)=(gz-z)\sum_{k\ge1}\lambda_k
\sum_{j=0}^{k-1}(gz)^jz^{k-1-j}$, where the last series converges in $\pkpA$: its $k$-th term is
$\lambda_k$ times a sum of $k$ Teichm\"uller lifts $[(g\varpi^\flat)^j(\varpi^\flat)^{k-1-j}]=[\eta^j
(\varpi^\flat)^{k-1}]$, i.e.\ $\lambda_k[\varpi^\flat]^{k-1}w_k$ with $w_k=\sum_{j<k}[\eta]^j\in\pkpA$, and
$\sum_k\lambda_kz^{k-1}w_k$ converges in the same topology in which $\sum_k\lambda_kz^k$ does (the
$(\ell,z)$-adic one, for which $\pkpA$ is complete and separated) because $\lambda_kw_k\in\pkpA$. Hence
$gz-z$, and so $z\tau E$, divides $gs-s$.
\end{proof}

\subsection{Frames and their consequences: (M), first half of (H)}\label{pkn:ss-frame}

\emph{Period rings.} Besides $\pknBp\subset\pknB$ we use two rings $B^+$ and $B_{\max}$ with
$\Gamma$-equivariant ring maps $\pkpA\to B^+\to B_{\max}\to\pknB$ and $B^+\to\pknBp$, the two maps
$B^+\to\pknB$ agreeing, the composite $\pkpA\to B^+\to\pknBp$ being the natural map (so that an element of $\pkpA$, such as $E$, has the same image in $\pknBp$ whether or not it is first viewed in $B^+$; for the rings of \S\ref{pkf:ss-rings} this is part of (P1)), and Frobenius endomorphisms $\varphi$ of $B^+$ and $B_{\max}$ such that
$\pkpA\to B^+\to B_{\max}$ are $\varphi$-equivariant, commuting with $\Gamma$ (any choice with these properties works in this subsection; the existence of frames,
\S\ref{pkf:sec}, uses $B^+=A_{\max}[1/\ell]$ and $B_{\max}=B^+[1/\mu]$ of \S\ref{pkf:ss-rings}). The
bicommutant $\mathfrak S_B\subseteq B^+$ of the chart is defined in Definition~\ref{pkf:def-SB}; by
Proposition~\ref{pkf:closure} it is the $\mathfrak S$-subalgebra obtained from the closure of
$\mathfrak S[E/\ell]$ (for the topology with basis $\ell^nA_{\max}$) by inverting $\ell$; all
that is used below is that $\mathfrak S_B$ is a subring of $B^+$ whose image in $\pknB$ lies in
$\pknK$ (true by the definition of $\pknK$). The map $\pkpA\to\pknBp$ is injective (Fontaine \cite[4.1.1]{pkfFontaineII}), so a nonzero element of $\pkpA$ has nonzero, hence invertible, image in the field
$\pknB$.

Fix a lattice with jump sets $L$, a $\Z_\ell$-basis $e_1,\dots,e_d$ of $L$, and write $\rho(g)\in
GL_d(\Z_\ell)$ for the matrix of $g\in\Gamma$ ($ge_j=\sum_k\rho(g)_{kj}e_k$). For a form $\mu$ let
$\mu(\mathbf e)=(\mu(1\otimes e_1),\dots,\mu(1\otimes e_d))$, a row vector over $\pknB$. Matrices over
$B^+$, $B_{\max}$, $\pkpA$ are mapped to $\pknB$ entrywise when multiplied with such rows.

\begin{definition}[Frame]\label{pkn:def-frame}
A \emph{frame} of $L$ (for an integer $a\ge0$) consists of $M\in GL_d(B_{\max})$, $U,Y\in
M_d(B^+)$, $Z\in M_d(\mathfrak S_B)$, $Q,Q'\in M_d(\pkpA)$ and vectors $n_1,\dots,n_d\in\pkpAmb(L)$ with:
\begin{enumerate}
\item[(Fr1)] $\rho(g)\,g(M)=M$ for $g\in\Gamma'$;
\item[(Fr2)] the image of $U$ in $B_{\max}$ is $(E\tau)^aM$;
\item[(Fr3)] $QQ'=\tau^{2a}\cdot1$, $UZ=E^{2a}Q$ (in $B^+$), $YZ=E^{2a}\cdot1$;
\item[(Fr4)] $\tau^an_i=\sum_jQ_{ji}(1\otimes e_j)$ for every $i$.
\end{enumerate}
\end{definition}

($\tau$ is a unit of $\pkpW$, so given $Q$ the $n_i$ exist and are unique.)

Fix a frame. For $m\ge0$ and a form $\mu$ put $c^{(m)}(\mu)=\mu(\mathbf e)\,\varphi^m(M)$ (a row over
$\pknB$).

\begin{lemma}\label{pkn:frame-id}
(i) For every form $\mu$ and $m\ge0$, the entries of $c^{(m)}(\mu)$ lie in $K'$; in particular in
$\pknBp\cap\pknK$.
(ii) For every form $\mu$, $m\ge0$ and $t\in\pkpA^d$:
$\ \varphi^m(E)^a\,\mu\bigl(\sum_it_i\varphi^m(n_i)\bigr)=c^{(m)}(\mu)\,\varphi^m(Z)\,t.$
(iii) For every $\nu$ and $m$ there is a form $\sigma$ with $c^{(m)}(\sigma)$ the $\nu$-th unit row.
(iv) (Coefficient test.) Let $W\subseteq\pknB$ be a $B^+$-submodule, $m\ge0$, $t\in\pkpA^d$, and suppose
$\mu(\sum_it_i\varphi^m(n_i))\in W$ for all forms $\mu$. Then the image of $\varphi^m(E)^{2a}t_i$ lies in
$W$ for every $i$.
\end{lemma}

\begin{proof}
(i) Applying $\varphi^m$ to (Fr1) ($\varphi$ commutes with $\Gamma$ and fixes $\rho(g)$) gives
$\rho(g)g(M_m)=M_m$ for $M_m=\varphi^m(M)$. For $g\in\Gamma'$, $g(\mu(1\otimes e_j))=\mu(1\otimes ge_j)=
\sum_k\mu(1\otimes e_k)\rho(g)_{kj}$, i.e.\ $g(\mu(\mathbf e))=\mu(\mathbf e)\rho(g)$; hence
$g(c^{(m)}(\mu))=\mu(\mathbf e)\rho(g)g(M_m)=c^{(m)}(\mu)$. An element $y\in\pknB$ fixed by $G_{K'}$ lies in
$K'$: by Lemma~\ref{pkm:BdR} with $n=1$ there is $c\in K'$ with $y-c\in\mathrm{Fil}^1$; $y-c$ is fixed, so
by the second assertion of Lemma~\ref{pkm:BdR} (with $m=1$, $c=0$) $y-c\in\mathrm{Fil}^n$ for every $n$,
and $\bigcap_n\mathrm{Fil}^n=0$. Finally $K'\subset\pknBp$ and $K'\subseteq\pknK$.

(ii) Apply $\varphi^m$ to (Fr4); $\varphi$ fixes $1\otimes e_j$: $\varphi^m(\tau)^a\sum_it_i\varphi^m(n_i)=
\sum_j(\varphi^m(Q)t)_j(1\otimes e_j)$. Apply $\mu$ ($\pkpA$-linear) and multiply by
$\varphi^m(E)^{2a}$; by (Fr3), $\varphi^m(E)^{2a}\varphi^m(Q)=\varphi^m(U)\varphi^m(Z)$, and by (Fr2) the
image of $\varphi^m(U)$ in $\pknB$ is that of $\varphi^m(E\tau)^a\varphi^m(M)$. So
$\varphi^m(E)^{2a}\varphi^m(\tau)^a\mu(\sum_it_i\varphi^m n_i)=\varphi^m(E\tau)^a\,c^{(m)}(\mu)\varphi^m(Z)t$.
Cancel $\varphi^m(E\tau)^a$, nonzero in $\pknB$.

(iii) Let $\lambda$ be a $\pknB$-linear functional on $V=\pknB\otimes_{\pkpA}\pkpW$ with $\lambda(1\otimes1)
=1$; it exists as soon as $1\otimes1\ne0$. For this, let $T=\pkpA\smallsetminus\{0\}$; $T$ maps to
nonzero elements of the domains $\pknB$ and $\pkpW$, so $T^{-1}\pknB$ and $T^{-1}\pkpW$ are nonzero vector
spaces over $F=\operatorname{Frac}\pkpA$, and $T^{-1}V=T^{-1}\pknB\otimes_FT^{-1}\pkpW$ is a tensor product of nonzero vector spaces over a
field, in which the image of $1\otimes1$ is a tensor of two nonzero vectors, hence nonzero; so
$1\otimes1\ne0$ in $V$. Put $\lambda_0(w)=\lambda(1\otimes w)$, an
$\pkpA$-linear map $\pkpW\to\pknB$ with $\lambda_0(1)=1$, and $\sigma(\sum_jw_j\otimes e_j)=\sum_j
\lambda_0(w_j)(M_m^{-1})_{\nu j}$ ($\pkpAmb(L)=\bigoplus_j\pkpW\otimes e_j$). Then $\sigma(\mathbf e)$ is
the $\nu$-th row of $M_m^{-1}$, so $c^{(m)}(\sigma)$ is the $\nu$-th unit row. $\sigma$ is a form: from
$\rho(g)g(M_m)=M_m$, $g(M_m^{-1})=M_m^{-1}\rho(g)$, i.e.\ $g(\sigma(\mathbf e))=\sigma(\mathbf e)\rho(g)$,
which says $\sigma(1\otimes gy)=g\sigma(1\otimes y)$ for $y$ in the basis, hence for all $y\in L$ by
$\Z_\ell$-linearity ($\rho(g)$ has entries in $\Z_\ell$).

(iv) Let $\sigma_\nu$ be as in (iii) and $x=\sum_it_i\varphi^m(n_i)$. By (ii),
$\varphi^m(E)^a\sigma_\nu(x)=(\varphi^m(Z)t)_\nu$. Hence
$\sum_\nu\varphi^m(Y)_{i\nu}\allowbreak \varphi^m(E)^a\sigma_\nu(x)=(\varphi^m(YZ)t)_i=\varphi^m(E)^{2a}t_i$ by (Fr3).
The left side lies in $W$, as $\varphi^m(Y)_{i\nu}\varphi^m(E)^a\in B^+$ and $\sigma_\nu(x)\in W$.
\end{proof}

\begin{proposition}[Frame theorem, part]\label{pkn:frame}
Suppose $L$ has a frame (for some $a$). Then:
\begin{enumerate}
\item[(i)] $n_i\in\pknMinf(L)\cap\mathcal N(L)$, and the $n_i$ are linearly independent over $\pkpA$;
\item[(ii)] $\pknMinf(L)=\sum_i\pkpA n_i$;
\item[(iii)] $\tau^a\pknMinf(L)\subseteq\pkpA\otimes L$ and $\tau^a(\pkpA\otimes L)\subseteq\pknMinf(L)$;
\item[(iv)] $\pkpA\mathcal N(L)$ is the $\pkpA$-span of $\mathcal N(L)$ (Lemma~\ref{pkm:inputs}(M), third
assertion);
\item[(v)] $\tau^a(1\otimes y)\in\pknMinf(L)\subseteq\pkpA\mathcal N(L)$ for $y\in L$
(Lemma~\ref{pkm:inputs}(H), first assertion, with $H=a$).
\end{enumerate}
\end{proposition}

\begin{proof}
\emph{$n_i\in\pknMinf$.} $\tau^an_i\in\pkpA\otimes L$ by (Fr4). For $m\ge1$ and a form $\mu$,
Lemma~\ref{pkn:frame-id}(ii) with $t$ the $i$-th unit vector gives $p_m^a\mu(\varphi^mn_i)=
c^{(m)}(\mu)\varphi^m(Z)_{\cdot i}\in\pknBp$ (entries of $c^{(m)}$ in $\pknBp$, entries of $\varphi^m(Z)$ in
$B^+$); $p_m$ is a unit of $\pknBp$ (Lemma~\ref{pkn:BdR}), so $\mu(\varphi^mn_i)\in\pknBp$.
\emph{$n_i\in\mathcal N$.} With $m=0$: $E^a\mu(n_i)=c^{(0)}(\mu)Z_{\cdot i}\in\pknK$ (Lemma~\ref{pkn:frame-id}(i)
and $Z\in M_d(\mathfrak S_B)$), and $E$ is a nonzero element of $\pknK$. Also $n_i\in\pknMinf\subseteq
\pkpA\mathcal N$ (take $r=1$).
\emph{Independence.} If $\sum_it_in_i=0$, Lemma~\ref{pkn:frame-id}(iv) with $W=0$, $m=0$ gives
$E^{2a}t_i\mapsto0$ in $\pknB$, so $t_i=0$ by injectivity.

\emph{$\tau$-test.} Let $t\in\pkpA^d$ and $x=\sum_it_in_i$, and suppose $\mu(\varphi^{m}x)\in\mathrm{Fil}^1$ for
all $m\ge1$ and all forms $\mu$. Then $\tau\mid t_i$ for all $i$. Indeed $\varphi^mx=\sum_i\varphi^m(t_i)
\varphi^m(n_i)$, and Lemma~\ref{pkn:frame-id}(iv) with $W=\mathrm{Fil}^1$ (a $B^+$-submodule, being a
$\pknBp$-submodule) gives $p_m^{2a}\varphi^m(t_i)\in\mathrm{Fil}^1$; for $r\in\pkpA$, $r\in\mathrm{Fil}^1$ iff
$\theta(r)=0$ ($\mathrm{Fil}^1=E\pknBp=\ker(\theta\colon\pknBp\to C)$, Lemma~\ref{pkn:BdR}), and
$\theta(p_m)\ne0$ (Corollary~\ref{pkn:E-gen}), so $\theta(\varphi^mt_i)=0$ for all $m\ge1$. Then $y=\varphi(t_i)$
has $\theta(\varphi^jy)=0$ for all $j\ge0$, so $\mu_\varepsilon=\varphi(\tau)$ divides $\varphi(t_i)$
(Lemma~\ref{pkn:mu}), and $\tau\mid t_i$.

(ii) $\supseteq$: (i) and Lemma~\ref{pkn:M}(i). $\subseteq$: first, for $x'=\sum_j\alpha_j(1\otimes e_j)\in
\pkpA\otimes L$, put $t=Q'\alpha$; by (Fr4) and (Fr3), $\tau^a\sum_it_in_i=\sum_j(QQ'\alpha)_j(1\otimes e_j)
=\tau^{2a}x'$, so $\tau^ax'=\sum_it_in_i$ ($\pkpAmb(L)$ has no $\tau$-torsion, $\tau\in\pkpW^\times$).
Now let $x\in\pknMinf$, with $\tau^Nx\in\pkpA\otimes L$; then $\tau^{a+N}x\in\sum_i\pkpA n_i$. We show:
if $k\ge0$ and $\tau^{k+1}x=\sum_it_in_i$, then $\tau^kx\in\sum_i\pkpA n_i$; by induction $x\in\sum\pkpA n_i$.
For $m\ge1$ and a form $\mu$, $\mu(\varphi^m(\tau^{k+1}x))=\varphi^m(\tau)\mu(\varphi^m(\tau^kx))$; the second
factor is in $\pknBp$ ($\tau^kx\in\pknMinf$), and $\theta(\varphi^m\tau)=\zeta_\ell^{\ell^m}-1=0$, so the first
is in $\mathrm{Fil}^1$. By the $\tau$-test, $t_i=\tau t'_i$, and $\tau(\tau^kx-\sum_it'_in_i)=0$ gives
$\tau^kx=\sum_it'_in_i$.

(iii) The first inclusion follows from (ii) and (Fr4); the second was shown in the proof of (ii)
(together with (i) and Lemma~\ref{pkn:M}(i)).

(iv) The span of $\mathcal N$ is contained in the $\pkpA$-module $\pkpA\mathcal N$. Conversely let
$x\in\pkpA\mathcal N$ and $r=p_{b_1}\cdots p_{b_k}$ with $rx\in\pknMinf$. By (ii), $rx=\sum_it_in_i$, so
$x=\sum_it_i\,(r^{-1}n_i)$ with $r^{-1}n_i\in\mathcal N$ by (i) and Lemma~\ref{pkn:M}(iii).

(v) is the second inclusion of (iii) with $x'=1\otimes y$.
\end{proof}

\subsection{Item (C)(c)}\label{pkn:ss-Cc}

\begin{proposition}[Lemma~\ref{pkm:inputs}(C)(c)]\label{pkn:Cc}
For every $n\ge1$ there are $\alpha,\beta\in\pkpA$ with $\tau E(\ell^{v_\ell(n)}+\beta z)=
\alpha([\varepsilon]^n-1)$.
\end{proposition}

\begin{proof}
Write $n=m\ell^v$ with $v=v_\ell(n)$, $\ell\nmid m$, and $x=[\varepsilon]^{\ell^v}$.
\emph{Step 1: $x-1=\tau\prod_{a=0}^{v}\varphi^a(E)\cdot u$ with $u\in\pkpA^\times$.} With
$\xi_\varepsilon$ as in Lemma~\ref{pkn:mu}, $\varphi(\tau)=\mu_\varepsilon=\tau\xi_\varepsilon$, so
$\varphi^{a+1}(\tau)=\varphi^a(\tau)\varphi^a(\xi_\varepsilon)$ and, by induction,
$\varphi^{v+1}(\tau)=\tau\prod_{a=0}^v\varphi^a(\xi_\varepsilon)$. Now $\varphi^{v+1}(\tau)=[\varepsilon]^{\ell^v}-1
=x-1$ ($\varphi[b]=[b^\ell]$), and $\varphi^a(\xi_\varepsilon)=\varphi^a(E)\varphi^a(u_0)$ with $u_0$ the unit
of Lemma~\ref{pkn:mu}; $\varphi^a$ maps units to units.
\emph{Step 2.} $[\varepsilon]^n-1=x^m-1=(x-1)w$ with $w=\sum_{k<m}x^k$, and $w$ is a unit: its image in
$W(\overline{\mathbb F}_\ell)$ is $m$ (as $\varepsilon\mapsto1$ in $O^\flat/\mathfrak m^\flat$), a unit since
$\ell\nmid m$. Hence $\tau\prod_{a=0}^v\varphi^a(E)=\alpha([\varepsilon]^n-1)$ with $\alpha=(uw)^{-1}$.
\emph{Step 3.} For $a\ge1$, $\varphi^a(E)=\varphi^a(z)^e\varphi^a(A)+\ell=z^{e\ell^a}\varphi^a(A)+\ell$ with
$\varphi^a(A)\in\mathfrak S$ ($\varphi(z)=z^\ell$, $\varphi(\Lambda)=\Lambda$), so $\varphi^a(E)\equiv\ell
\pmod{z\mathfrak S}$ and $\prod_{a=1}^v\varphi^a(E)=\ell^v+\beta z$ with $\beta\in\mathfrak S$. Since
$\varphi^0(E)=E$, Step 2 reads $\tau E(\ell^v+\beta z)=\alpha([\varepsilon]^n-1)$.
\end{proof}

\emph{Where the remaining items are proved.} The existence of a frame for every lattice with jump
sets, with $a$ the jump bound, is Theorem~\ref{pkf:main}, proved in \S\ref{pkf:sec} modulo published
facts of $\ell$-adic Hodge theory and the ring statements (A1)--(A8), (P4) proved in \S\ref{pkr:sec};
hence items (M) (third assertion) and (H) (first assertion) of Lemma~\ref{pkm:inputs} hold. Items (G)
and (H) (in the form stated in \S\ref{pkl:sec}), (F), (R), (D), (C)(d), (G$_0$) and ($\Phi$) are proved in
\S\ref{pkl:sec}; Lemma~\ref{pkp:orders-add} and Lemma~\ref{pkp:tau}(i), (iv), (v) in \S\ref{pkt:sec};
Proposition~\ref{pki:chart} in \S\ref{pkj:sec}.


%% ===== proofs/C-kisin-frame.tex =====
% proofs/C-kisin-frame.tex
% Existence of a frame (Definition pkn:def-frame) for every lattice with jump sets.
% Labels start with pkf:. Conventions and macros as in C-kisin-pairing.tex, C-kisin-inputs.tex.
\providecommand{\auO}{\mathcal O}
\providecommand{\auQlb}{\overline{\Q}_\ell}
\providecommand{\pkpA}{A_{\mathrm{inf}}}
\providecommand{\pkpW}{W(C^\flat)}
\providecommand{\pkpAmb}{\mathbb A}
\providecommand{\pknBp}{B^+_{\mathrm{dR}}}
\providecommand{\pknB}{B_{\mathrm{dR}}}
% macros of this file (prefix \pkf)
\providecommand{\pkfAm}{A_{\max}}
\providecommand{\pkfAo}{A_{\max}^{\circ}}
\providecommand{\pkfBp}{B^+}
\providecommand{\pkfBm}{B_{\max}}
\providecommand{\pkfSB}{\mathfrak S_B}
\providecommand{\pkfEnd}{\mathcal G}

\section{Existence of frames}\label{pkf:sec}

This section proves that every lattice with jump sets has a frame in the sense of
Definition~\ref{pkn:def-frame}, with $a$ the jump bound (Theorem~\ref{pkf:main}). We use the
conventions of \S\S\ref{pkp:sec}--\ref{pkn:sec}: $\ell$ is the residue characteristic,
$\pkpA=W(O^\flat)$, $\theta\colon\pkpA\to O_C$, $\varphi$ the (bijective) Frobenius of $\pkpA$, a fixed
Kummer chart (Definition~\ref{pkp:def-chart}) with $\mathfrak S=\Lambda[[z]]\subset\pkpA$,
$E=z^eA+\ell$, $p_a=\varphi^a(E)$, the subgroup $\Gamma'$ and $K'=\auQlb^{\,j(\Gamma')}$;
$\varepsilon=(1,\zeta_\ell,\zeta_{\ell^2},\dots)$, $\mu=[\varepsilon]-1$, $\tau=\varphi^{-1}(\mu)=
[\varepsilon^{1/\ell}]-1$. By Lemma~\ref{pkn:mu}, $\mu=E\tau v$ with $v\in\pkpA^\times$; we fix this
$v$. For $x\in\pkpA$ we write $\bar x\in O^\flat=\pkpA/\ell\pkpA$ for its reduction.

\emph{Dependencies.} Theorem~\ref{pkf:main} is proved in full \emph{modulo} the published facts of
$\ell$-adic Hodge theory (P1)--(P3), (P5), (P6) stated in \S\ref{pkf:ss-rings}, cited with author and
work. The eight auxiliary statements (A1)--(A8) about the rings $\pkfAm$, $\pkfBp$ and the chart, and
the statement (P4), are proved in \S\ref{pkr:sec} (which also uses the standard facts (W1)--(W4)
recalled there). Every other step is proved here. (The labels (A1)--(A8) are distinct from the case
inputs (I1)--(I6) of \S\S\ref{pkc:sec}--\ref{pkd:sec}.)

\subsection{Period rings, the representation, and the inputs}\label{pkf:ss-rings}

\emph{The rings.} $E$ is distinguished and generates $\ker\theta$
(Corollary~\ref{pkn:E-gen}). Let $R=\pkpA[X]/(\ell X-E)$ and let $\pkfAm$ be its $\ell$-adic
completion; write $E/\ell$ for the image of $X$. Put $\pkfBp=\pkfAm[1/\ell]$ and
$\pkfBm=\pkfBp[1/\mu]$. (Replacing $E$ by another generator $Eu$ of $\ker\theta$, $u\in\pkpA^\times$,
gives an isomorphic ring via $X\mapsto Xu$; so these are Colmez's rings $A_{\max}$, $B^+_{\max}$ and
$B_{\max}$.) We write $\pkfAo$ for the image of $\pkfAm$ in $\pkfBp$.

\begin{enumerate}
\item[(P1)] (Structure; Colmez \cite[\S III.2]{pkfColmez98}, Fontaine \cite[4.1]{pkfFontaineII}.)
$\pkfAm$ is $\ell$-torsion free, $\ell$-adically complete and separated; the maps
$\pkpA\to\pkfAm\to\pkfBp\to\pkfBm\to\pknB$ and $\pkfBp\to\pknBp$ are injective ring maps, the two maps
$\pkfBp\to\pknB$ agree, and the composite $\pkpA\to\pknBp$ is the natural map; $\pkfBp$ and $\pkfBm$
are domains. $\varphi$ and $G_{\Q_\ell}$ act on $\pkfAm$, $\pkfBp$, $\pkfBm$ compatibly with these maps
(the action of $G_{\Q_\ell}$ also on $\pknB$), $\varphi$ is injective and commutes with $G_{\Q_\ell}$.
In particular $\varphi$ and every $g\in G_{\Q_\ell}$ map $\pkfAo$ into itself.
\item[(P2)] (Crystalline representations; Fontaine \cite[5.1.2]{pkfFontaineIII}, Colmez \cite[\S\S III.2, III.4]{pkfColmez98}.) $\pkfBm^{G_{K'}}=K'_0$, the maximal unramified subfield of $K'$; a continuous
representation of $G_{K'}$ on a finite dimensional $\Q_\ell$-vector space $V$ is crystalline iff
$\dim_{K'_0}(\pkfBm\otimes_{\Q_\ell}V)^{G_{K'}}=\dim V$, and then the natural map
$\pkfBm\otimes_{K'_0}(\pkfBm\otimes V)^{G_{K'}}\to\pkfBm\otimes V$ is an isomorphism. Crystalline
representations are de Rham. (Fontaine proves $B_{\mathrm{cris}}^{G_{K'}}=K'_0$ and that $B_{\mathrm{cris}}$ is $G_{K'}$-regular; Colmez notes $\varphi(\pkfBm)\subseteq B_{\mathrm{cris}}\subseteq\pkfBm$ and that $\pkfBm$ and $B_{\mathrm{cris}}$ give the same $D_{\mathrm{cris}}$. Hence $K'_0\subseteq\pkfBm^{G_{K'}}$, and if $b\in\pkfBm^{G_{K'}}$ then $\varphi(b)\in B_{\mathrm{cris}}^{G_{K'}}=K'_0=\varphi(K'_0)$, so $b\in K'_0$ by injectivity of $\varphi$.)
\item[(P3)] (Filtered comparison for de Rham representations; Fontaine \cite[3.8]{pkfFontaineIII}.) If $V$ is de Rham, $D_{\mathrm{dR}}(V)=(\pknB\otimes V)^{G_{K'}}$ with the filtration
$\mathrm{Fil}^kD_{\mathrm{dR}}=D_{\mathrm{dR}}\cap(\mathrm{Fil}^k\pknB\otimes V)$, then
$\pknB\otimes_{K'}D_{\mathrm{dR}}\to\pknB\otimes V$ is an isomorphism carrying
$\sum_{k}\mathrm{Fil}^{-k}\pknB\otimes\mathrm{Fil}^kD_{\mathrm{dR}}$ onto $\pknBp\otimes V$. The
\emph{Hodge--Tate jumps} of $V$ (at the identity embedding) are the $k$ with
$\mathrm{gr}^kD_{\mathrm{dR}}\ne0$; for $V=L[1/\ell]$, $L$ a lattice with jump sets viewed as a
$\Z_\ell$-representation, these are the Hodge--Tate jumps of $L$ at all embeddings, so they lie in
$\bigcup_iU(i)$ (Definition~\ref{pkp:def-orders}; this is the convention of the paper).
\item[(P4)] (Proved in Proposition~\ref{pkr:P4}, from (P1), (P6).) If $b\in\pkfBm$ and $m\ge0$ are such that the image of
$\varphi^n(b)$ in $\pknB$ lies in $\mathrm{Fil}^{-m}$ for every $n\ge0$, then $\mu^mb\in\pkfBp$.
\item[(P5)] (Ax--Sen--Tate; \cite[Theorem, p.~417]{pkfAx70}: for a field $k$ henselian for a non-trivial
valuation, $G_k$ acting on the completion of an algebraic closure of $k$, the fixed field is the completion of
the perfect closure of $k$; here $k=L$, which is henselian as an algebraic extension of $\Q_\ell$, of
characteristic $0$, and the completion of $\auQlb$ is $C$.) For every algebraic extension $L$ of $\Q_\ell$ inside
$\auQlb$, the elements of $C$ fixed by $\operatorname{Gal}(\auQlb/L)$ form the closure $\hat L$ of $L$ in
$C$. It is used only through Lemma~\ref{pkf:fixed}, with $H\subseteq G_{\Q_\ell}$ the group of elements
fixing $\mathfrak S$ pointwise. (The same lemma can be read off from the field of norms of Fontaine and
Wintenberger \cite{pkfWintenberger83} and the tilting of Scholze \cite{pkfScholze12}, but we do not use
them.)
\item[(P6)] $\mathrm{Fil}^n\pknB=E^n\pknBp$ (Lemma~\ref{pkn:BdR}, which rests on Fontaine's theorem
that $\pknBp$ is a discrete valuation ring with uniformiser $t$).
\end{enumerate}

\emph{The representation.} Let $L$ be a lattice with jump sets, $e_1,\dots,e_d$ a $\Z_\ell$-basis,
$\rho(g)\in GL_d(\Z_\ell)$ the matrix of $g\in\Gamma$ ($ge_j=\sum_k\rho(g)_{kj}e_k$). The hypothesis
used is: \emph{$L[1/\ell]$ is crystalline as a representation of $\Gamma'$ (through $j$, of
$G_{K'}$), with Hodge--Tate jumps in $[-a,a]$}, where $a\ge0$ is the \emph{jump bound},
$a=\max\{|u|:u\in\bigcup_iU(i)\}$. For a ring $B$ with an action of $G_{\Q_\ell}$, a matrix $M\in
M_d(B)$ has \emph{invariant columns} if $\rho(g)\,g(M)=M$ for $g\in\Gamma'$ ($g$ acting through $j$
entrywise); equivalently the vectors $\sum_kM_{kj}\otimes e_k$ of $B\otimes L$ are fixed by $g\otimes
g$. If $M$ has invariant columns so does $\varphi^n(M)$ for a ring with a Frobenius commuting with
$\Gamma'$, since $\varphi$ fixes $\rho(g)$.

\emph{Auxiliary inputs.} Write $(\ell,E)$ for the ideal of $\pkpA$ generated by $\ell,E$.
\begin{enumerate}
\item[(A1)] $\pkpA\cap\ell^n\pkfAm=(\ell,E)^n$ for $n\ge0$.
\item[(A2)] If $\delta\in\pkpA$ divides $\mu$, $h\ge0$, $w\in\pkfBp$ and $\delta^hw\in\pkfAo$, then
$\ell^{2h}w\in\pkfAo$.
\item[(A3)] (Contraction.) Let $V\subseteq\pkfAo$, $h\ge0$ and $u\in\pkpA$ be such that every $\alpha\in V$
has some $\alpha'\in V$ with $E^h\alpha=u\,\varphi(\alpha')$. Then $\ell^{h+1}V\subseteq\pkpA$.
\item[(A4)] (Approximation.) For $y\in\pkfBp$ and $h\ge0$ there are $m\ge0$, $\alpha\in\pkpA$ and
$w\in\pkfAo$ with $\ell^my=\alpha+\ell^m\mu^hw$.
\item[(A5)] (Topological basis.) There are a family $(b_i)_{i\in I}$ in $\pkpA$ and $\mathfrak S$-linear
maps $\kappa_i\colon\pkpA\to\mathfrak S$ such that: for every $x\in\pkpA$ and $N\ge0$, $x-\sum_{i\in
G}\kappa_i(x)b_i\in(\ell,z)^N\pkpA$ for all sufficiently large finite $G\subseteq I$ (in particular
$x=0$ if all $\kappa_i(x)=0$); and there is $0\in I$ with $\kappa_0(s)=s$ for $s\in\mathfrak S$.
\item[(A6)] (Freeness.) Let $P\subseteq\pkpA^d$ be an $\mathfrak S$-submodule containing no $d+1$
elements linearly independent over $\mathfrak S$, such that $za\in\ell P$ with $a\in P$ implies
$a\in\ell P$. Then $P$ is a free $\mathfrak S$-module of finite rank ($\le d$).
\item[(A7)] (Closed ideals.) If $\delta,c\in\pkpA$, $\ell\nmid\delta$, and for every $m$ there is
$z_m\in\pkpA$ with $c-\delta z_m\in(\ell,E)^m$, then $\delta\mid c$.
\item[(A8)] (Coordinates.) (a) In the notation of (A5): for $i\in I$, $x\in\pkpA$ and $m\ge0$ with
$\theta(\varphi^{-m}x)\in\ell O_C$, there is $s\in\mathfrak S$ with $p_m\mid\kappa_i(x)-\ell s$ in
$\mathfrak S$. (b) If $b\in\pkfBp$ and $\varphi(\tau)\mid\varphi(b)$ in $\pkfBp$, then $\tau\mid b$ in
$\pkfBp$. (c) Let $G_f$ be the extension of Lemma~\ref{pkf:extend}. For $\lambda\in K'^d$,
$y\in(\pkfBp)^d$ and an $\mathfrak S$-linear $f\colon\pkpA\to\pkpA$, $\sum_i\lambda_iy_i$ divides
$\sum_i\lambda_iG_f(y_i)$ in $\pknBp$ (images in $\pknBp$).
\end{enumerate}

The inputs (A1)--(A8) are proved in \S\ref{pkr:sec}: (A1) is Proposition~\ref{pkr:I1}, (A2)
Proposition~\ref{pkr:I2}, (A3) Proposition~\ref{pkr:I3}, (A4) Proposition~\ref{pkr:I4}, (A5)
Proposition~\ref{pkr:I5}, (A6) Proposition~\ref{pkr:I6}, (A7) Proposition~\ref{pkr:I7}, and (A8)(a),
(b), (c) Proposition~\ref{pkr:I8}. Those proofs use only (P1), (P6), Lemma~\ref{pkf:approx0},
Lemmas~\ref{pkf:SB-basic}, \ref{pkf:extend}, the description $K'=K_0(\pi)$ from the proof of
Lemma~\ref{pkf:invariants}, and \S\S\ref{pkp:sec}--\ref{pkn:sec}; none of them uses
Theorem~\ref{pkf:main} or a statement depending on it.

\begin{lemma}[Elementary facts on $\pkfAm$]\label{pkf:approx0}
(i) Every $\beta\in\pkfAm$ can be written $\beta=\sum_{k\ge0}\alpha_k(E/\ell)^k$ with $\alpha_k\in\pkpA$,
$\alpha_k\to0$ $\ell$-adically; for every $n$ there is a polynomial $P\in\pkpA[X]$ with
$\beta-P(E/\ell)\in\ell^n\pkfAm$. (ii) For $y\in\pkfBp$ and $n\ge0$ there are $m\ge0$, $\alpha\in\pkpA$,
$w\in\pkfAo$ with $\ell^my=\alpha+\ell^{m+n}w$. (iii) $\varphi(E)=\ell\cdot u_E$ with
$u_E\in\pkfAm^\times$; hence $\varphi(E)$ is a unit of $\pkfBp$. (iv) For $g\in G_{\Q_\ell}$, $gE=Er$ and
$g\tau=\tau r'$ with $r,r'\in\pkpA^\times$. (v) $\theta$ extends to $\pkfBp\to C$, and
$\ker(\theta|_{\pkfBp})=E\pkfBp$; consequently, for $x\in\pkfBp$ and $n\ge0$, $E^n\mid x$ in $\pkfBp$
iff the image of $x$ in $\pknBp$ lies in $\mathrm{Fil}^n$.
\end{lemma}

\begin{proof}
(i) The $\ell$-adic completion of $\pkpA[X]$ is the ring of series $\sum\alpha_kX^k$ with $\alpha_k\to0$
($\pkpA$ is $\ell$-adically complete), and it surjects onto the completion of the quotient $R$.
Truncating the series gives $P$.
(ii) $\ell^{m_0}y\in\pkfAo$ for some $m_0$. Write $\ell^{m_0}y=P(E/\ell)+\ell^{m_0+n}w$ by (i) with
$\deg P\le D$; then $\ell^DP(E/\ell)\in\pkpA$, and $m=m_0+D$ works.
(iii) $E=[E_0]+\ell E'$ with $E'\in\pkpA^\times$ (Definition~\ref{pkn:def-dist}). In $\pkfAm$,
$[E_0]=E-\ell E'=\ell(E/\ell-E')$, so $\varphi([E_0])=[E_0]^\ell=\ell^\ell(E/\ell-E')^\ell$ and
$\varphi(E)=\ell\bigl(\varphi(E')+\ell^{\ell-1}(E/\ell-E')^\ell\bigr)$. The bracket is a unit plus $\ell$
times an element, hence a unit of the $\ell$-adically complete ring $\pkfAm$.
(iv) $\theta(gE)=g\theta(E)=0$, so $E\mid gE$ (Corollary~\ref{pkn:E-gen}); likewise $E\mid g^{-1}E$,
i.e.\ $gE\mid E$; as $\pkpA$ is a domain, $gE=Er$ with $r$ a unit. For $\tau$: $\varphi(g\tau)=g\mu$ and
$\theta(\varphi^j(g\mu))=g\theta(\varphi^j\mu)=0$ for $j\ge0$, so $\mu\mid g\mu$ by Lemma~\ref{pkn:mu},
i.e.\ $\varphi(\tau)\mid\varphi(g\tau)$, and $\tau\mid g\tau$ since $\varphi$ is an automorphism of $\pkpA$;
symmetrically $g\tau\mid\tau$.
(v) $\theta$ extends to $\pkpA[X]\to O_C$ by $X\mapsto0$, which kills $\ell X-E$, and by continuity to
$\pkfAm$, then to $\pkfBp$; this is the restriction of $\theta\colon\pknBp\to C$ (P1). By (i),
$\beta=\alpha_0+(E/\ell)\beta'$ with $\beta'\in\pkfAm$, and $\theta(\beta)=\theta(\alpha_0)$. If
$\theta(\beta)=0$ then $\alpha_0=Ec$ (Corollary~\ref{pkn:E-gen}) and $\beta=E(c+\ell^{-1}\beta')\in
E\pkfBp$. Since $\mathrm{Fil}^1\cap\pkfBp=\ker\theta|_{\pkfBp}$, the case $n=1$ follows; for general $n$,
if $x\in\mathrm{Fil}^n$, $n\ge1$, then $x=Ex_1$ with $x_1\in\pkfBp$, and $x_1\in\mathrm{Fil}^{n-1}$
because $E$ is a uniformiser of $\pknBp$ (P6); induct.
\end{proof}

\subsection{The ring $\mathfrak S_B$ and extension of $\mathfrak S$-linear maps}\label{pkf:ss-SB}

Let $\pkfEnd$ be the set of $\mathfrak S$-linear maps $G\colon\pkfBp\to\pkfBp$ with
$G(\pkfAo)\subseteq\pkfAo$.

\begin{definition}[The ring $\mathfrak S_B$]\label{pkf:def-SB}
$\pkfSB=\{y\in\pkfBp:\ G(yx)=yG(x)\text{ for all }G\in\pkfEnd,\ x\in\pkfBp\}$.
\end{definition}

$\pkfSB$ is an $\mathfrak S$-subalgebra of $\pkfBp$ containing the image of $\mathfrak S$ (each $G$ is
$\mathfrak S$-linear). This is the bicommutant of \S\ref{pkn:ss-frame}; Proposition~\ref{pkf:closure}
shows that it is the $\mathfrak S$-algebra described there. Concretely (Proposition~\ref{pkf:closure}),
$\pkfSB=\widehat{\mathfrak S[E/\ell]}[1/\ell]$, where the hat denotes the $\ell$-adic closure in $\pkfAm$;
in the unramified case this ring embeds into $\Lambda[[z/\ell]][1/\ell]$ (Lemma~\ref{pkg:Shat}). It enters
only through the properties proved in this section and in \S\ref{pkg:ss-radius}; no comparison with
Breuil--Kisin modules is needed.

\begin{lemma}[Elementary properties]\label{pkf:SB-basic}
(i) If $v'\in\pkfSB\smallsetminus\{0\}$ and $v'u\in\pkfSB$ ($u\in\pkfBp$), then $u\in\pkfSB$. Hence
$\ell^{-1}\in\pkfSB$, $E/\ell\in\pkfSB$, and $\varphi(E)^{-1}\in\pkfSB$.
(ii) If $g\in G_{\Q_\ell}$ fixes $\mathfrak S$ pointwise, then $g$ fixes $\pkfSB$ pointwise.
\end{lemma}

\begin{proof}
(i) $v'G(ux)=G(v'ux)=v'uG(x)$ and $\pkfBp$ is a domain. Apply with $v'=\ell$ and $v'u=1$, with $v'=\ell$
and $v'u=E$, and with $v'=\varphi(E)\in\mathfrak S$ and $v'u=1$ (Lemma~\ref{pkf:approx0}(iii)).
(ii) $g$ is additive, $g(sx)=g(s)g(x)=sg(x)$ for $s\in\mathfrak S$, and $g(\pkfAo)\subseteq\pkfAo$ (P1),
so $g\in\pkfEnd$. For $y\in\pkfSB$: $g(y)=g(y\cdot1)=y\,g(1)=y$.
\end{proof}

\begin{lemma}[Density]\label{pkf:density}
Let $D\colon\pkfBp\to\pkfBp$ be additive with $D(\pkfAo)\subseteq\pkfAo$ and $D(\alpha)=0$ for
$\alpha\in\pkpA$. Then $D=0$.
\end{lemma}

\begin{proof}
Let $y\in\pkfBp$, $n\ge0$, and $\ell^my=\alpha+\ell^{m+n}w$ as in Lemma~\ref{pkf:approx0}(ii). An
additive map commutes with multiplication by integers, so $\ell^mD(y)=\ell^{m+n}D(w)$ and $D(y)=
\ell^nD(w)\in\ell^n\pkfAo$. As $\pkfAm$ is $\ell$-adically separated and $\pkfAm\to\pkfBp$ injective,
$D(y)=0$.
\end{proof}

\begin{lemma}[Extension]\label{pkf:extend}
Every $\mathfrak S$-linear $f\colon\pkpA\to\pkpA$ extends to some $G_f\in\pkfEnd$, i.e.\
$G_f(\alpha)=f(\alpha)$ for $\alpha\in\pkpA$. We fix such a $G_f$ for every $f$. For $y\in\pkfSB$ and
$x\in\pkfBp$, $G_f(yx)=yG_f(x)$.
\end{lemma}

\begin{proof}
\emph{Step 1: if $\alpha\in\pkpA\cap\ell^n\pkfAm$ then $f(\alpha)\in\ell^n\pkfAm$.} By (A1),
$\alpha\in(\ell,E)^n$, an ideal generated by the elements $\ell^iE^{n-i}$ of $\mathfrak S$; by
$\mathfrak S$-linearity $f(\alpha)\in(\ell,E)^n\subseteq\ell^n\pkfAm$.
\emph{Step 2.} Let $\mathcal D=\{\beta\in\pkfAm:\ \ell^k\beta\in\pkpA\text{ for some }k\}$, an
$\mathfrak S$-submodule. For $\beta\in\mathcal D$ with $\ell^k\beta\in\pkpA$, Step 1 gives
$f(\ell^k\beta)=\ell^k\gamma$ with $\gamma\in\pkfAm$, unique as $\pkfAm$ is $\ell$-torsion free, and
independent of $k$ since $f(\ell\cdot)=\ell f(\cdot)$; put $G^\circ(\beta)=\gamma$. $G^\circ$ is $\mathfrak S$-linear
on $\mathcal D$ and extends $f$. If moreover $\beta\in\ell^n\pkfAm$, then $\ell^k\beta\in\pkpA\cap
\ell^{k+n}\pkfAm$, so $\ell^k G^\circ(\beta)\in\ell^{k+n}\pkfAm$ by Step 1 and $G^\circ(\beta)\in\ell^n\pkfAm$.
By Lemma~\ref{pkf:approx0}(i), $\mathcal D$ (which contains all $P(E/\ell)$, $P\in\pkpA[X]$) is
$\ell$-adically dense in $\pkfAm$. As $\pkfAm$ is $\ell$-adically complete and separated, $G^\circ$ extends
uniquely to an $\ell$-adically continuous $\mathfrak S$-linear map $G^\circ\colon\pkfAm\to\pkfAm$.
\emph{Step 3.} $G_f(\ell^{-k}\beta)=\ell^{-k}G^\circ(\beta)$ is well defined on $\pkfBp=\pkfAm[1/\ell]$, is
$\mathfrak S$-linear, extends $f$ and maps $\pkfAo$ into itself. The last assertion is
Definition~\ref{pkf:def-SB}.
\end{proof}

\begin{lemma}[Values in $\mathfrak S_B$]\label{pkf:values}
Let $\sigma$ be either the identity or $\varphi$, and $f$ an $\mathfrak S$-linear map with
$\sigma(f(\alpha))\in\mathfrak S$ for all $\alpha\in\pkpA$. Then $\sigma(G_f(y))\in\pkfSB$ for all
$y\in\pkfBp$. In particular, with $\kappa_i$ as in (A5) and $G_i=G_{\kappa_i}$ (viewing $\kappa_i$ as a
map to $\pkpA$): $G_i(y)\in\pkfSB$ for all $y$; $G_0(y)=y$ for $y\in\pkfSB$; and
$\varphi(\pkfSB)\subseteq\pkfSB$.
\end{lemma}

\begin{proof}
Let $G\in\pkfEnd$ and $x\in\pkfAo$. The map $D(y)=G(\sigma(G_f(y))x)-\sigma(G_f(y))G(x)$ is additive and
maps $\pkfAo$ into $\pkfAo$ ($G_f$, $\sigma$, $G$ preserve $\pkfAo$, a ring). For $\alpha\in\pkpA$,
$s=\sigma(f(\alpha))\in\mathfrak S$ and $D(\alpha)=G(sx)-sG(x)=0$. By Lemma~\ref{pkf:density}, $D=0$. For
general $x\in\pkfBp$ take $n$ with $\ell^nx\in\pkfAo$ and cancel $\ell^n$ ($G$ is additive, $\pkfBp$ a
domain). So $\sigma(G_f(y))\in\pkfSB$. For $G_i$ take $\sigma=1$. For $y\in\pkfSB$: $G_0(y)=y\,G_0(1)=
y\,\kappa_0(1)=y$ by Lemma~\ref{pkf:extend}. Finally, for $y\in\pkfSB$, $\varphi(y)=\varphi(G_0(y))$ and
$\varphi(\kappa_0(\alpha))\in\varphi(\mathfrak S)\subseteq\mathfrak S$.
\end{proof}

\begin{proposition}[Description of $\mathfrak S_B$]\label{pkf:closure}
$\pkfSB$ is the set of $y\in\pkfBp$ for which there is $m$ such that for every $n$,
$\ell^my\in\mathfrak S[E/\ell]+\ell^n\pkfAo$. In particular $\pkfSB$ is the $\mathfrak S$-subalgebra
described in \S\ref{pkn:ss-frame}, and its image in $\pknB$ lies in $\mathcal K_0$.
\end{proposition}

\begin{proof}
If $y$ has the property, let $D=G_0-1$, which maps $\pkfAo$ into itself and is $\pkfSB$-linear
(Lemma~\ref{pkf:extend}) and kills $\pkfSB$ (Lemma~\ref{pkf:values}); since $\mathfrak S[E/\ell]
\subseteq\pkfSB$ (Lemma~\ref{pkf:SB-basic}(i)), $D(\ell^my)=\ell^nD(w)\in\ell^n\pkfAo$ for all $n$, so
$D(\ell^my)=0$ and $\ell^my=G_0(\ell^my)\in\pkfSB$; then $y\in\pkfSB$ by Lemma~\ref{pkf:SB-basic}(i).
Conversely let $y\in\pkfSB$, $\ell^my=\beta\in\pkfAo$, and $\beta=P(E/\ell)+\ell^nw$
(Lemma~\ref{pkf:approx0}(i)). Applying $G_0$: $\ell^my=G_0(P(E/\ell))+\ell^nG_0(w)$, and for a monomial
$G_0(\alpha(E/\ell)^k)=(E/\ell)^k\kappa_0(\alpha)\in\mathfrak S[E/\ell]$ ($E/\ell\in\pkfSB$). The last
sentence is the definition of $\mathcal K_0$ in Definition~\ref{pkn:def-N}.
\end{proof}

\begin{lemma}[Invariants]\label{pkf:invariants}
If $b\in\pkfBm$ is fixed by $j(\Gamma')$, then $b$ is the image of an element of $\pkfSB$.
\end{lemma}

\begin{proof}
By the chart, $j(\Gamma')$ is the group of elements of $G_{\Q_\ell}$ fixing $\Lambda$ and $\theta(z)=\pi$,
i.e.\ $G_{K'}$ with $K'=K_0(\pi)$, $K_0=\Lambda[1/\ell]$ (by Weierstrass preparation $E(X)$ is a
unit of $\Lambda[[X]]$ times an Eisenstein polynomial $P\in\Lambda[X]$ of degree $e$, and $\pi$ is a
root of $P$, so $[K_0(\pi):K_0]=e$); so $K'_0=\Lambda[1/\ell]$. By (P2), $b\in\Lambda[1/\ell]$, i.e.\
$b=\ell^{-m}\lambda$ with $\lambda\in\Lambda\subseteq\mathfrak S$, and $\ell^{-m}\in\pkfSB$.
\end{proof}

\begin{lemma}[Coordinate approximation]\label{pkf:coord-approx}
For $y\in\pkfBp$ and $m\ge0$ there is a finite $G\subseteq I$ with $y-\sum_{i\in G}b_iG_i(y)\in
\ell^m\pkfAo$.
\end{lemma}

\begin{proof}
Write $\ell^ky=\alpha+\ell^{k+m}w$ (Lemma~\ref{pkf:approx0}(ii)). Since $z^e=(E-\ell)A^{-1}\in(\ell,E)$,
$(\ell,z)^{e(k+m)}\subseteq(\ell,E)^{k+m}\subseteq\ell^{k+m}\pkfAm$. By (A5) choose $G$ with
$\alpha-\sum_{i\in G}\kappa_i(\alpha)b_i\in(\ell,z)^{e(k+m)}$. The map $T=1-\sum_{i\in G}b_iG_i$ is
additive and preserves $\pkfAo$, and $T(\alpha)=\alpha-\sum_{G}b_i\kappa_i(\alpha)$. So $\ell^kT(y)=
T(\alpha)+\ell^{k+m}T(w)\in\ell^{k+m}\pkfAo$, and $T(y)\in\ell^m\pkfAo$ ($\pkfAm$ is $\ell$-torsion free).
\end{proof}

\subsection{Stage (a): a period matrix}\label{pkf:ss-pm}

\begin{definition}[Period matrix]\label{pkf:def-pm}
Let $h\ge0$ and $v_1\in\pkpA^\times$. A matrix $U\in M_d(\pkfBp)$ is a \emph{period matrix of height $h$
with Frobenius factor $v_1$} if:
\begin{enumerate}
\item[(PM1)] there is $U'\in M_d(\pkfBp)$ with $UU'=(E\tau)^h\cdot1$;
\item[(PM2)] $\varphi(U)=v_1UF$ with $F\in M_d(\pkfSB)$ having a left inverse $F'\in M_d(\pkfBp)$;
\item[(PM3)] for every $g\in G_{\Q_\ell}$ fixing $\Lambda$ pointwise and $\theta(z)$ there is $R_g\in
M_d(\pkpA)$ with $R_g\,g(U)=U$;
\item[(PM4)] for every $\mathfrak S$-linear $f\colon\pkpA\to\pkpA$, $n\ge0$ and $y\in(\pkfBp)^d$: if
$Uy\in E^n(\pkfBp)^d$ then $UG_f(y)\in E^n(\pkfBp)^d$ ($G_f$ applied entrywise).
\end{enumerate}
\end{definition}

\begin{lemma}[Hodge frames]\label{pkf:hodge}
Let $N\in GL_d(\pknB)$ have invariant columns. Then there are $Y\in SL_d(K')$, $W\in GL_d(\pknBp)$ and
integers $k_1,\dots,k_d$, each a Hodge--Tate jump of $L[1/\ell]$, with $NY=W\operatorname{diag}
(t^{k_1},\dots,t^{k_d})$.
\end{lemma}

\begin{proof}
The columns $\delta_j=\sum_kN_{kj}\otimes e_k$ are $G_{K'}$-invariant vectors of $\pknB\otimes L$, linearly
independent over $\pknB$ ($N$ invertible), hence over $K'$; so they lie in $D_{\mathrm{dR}}$, which has
$K'$-dimension $\le d$ (by (P3), or since $L[1/\ell]$ is de Rham by (P2)); they form a $K'$-basis of
$D_{\mathrm{dR}}$. Choose a $K'$-basis $\delta'_1,\dots,\delta'_d$ adapted to the filtration, with
$\delta'_j\in\mathrm{Fil}^{k_j}D_{\mathrm{dR}}$ and $\mathrm{Fil}^kD_{\mathrm{dR}}$ spanned by the
$\delta'_j$ with $k_j\ge k$; the $k_j$ are then the jumps with multiplicity. Let $Y\in GL_d(K')$ be the
change of basis ($\delta'=\delta Y$); rescaling $\delta'_1$ by $\det Y^{-1}\in K'^\times$ we may take
$\det Y=1$. By (P3), $\pknBp\otimes L=\sum_k\mathrm{Fil}^{-k}\pknB\otimes\mathrm{Fil}^kD_{\mathrm{dR}}
=\bigoplus_j t^{-k_j}\pknBp\,\delta'_j$ (the second equality because $\mathrm{Fil}^{-k}\pknB=t^{-k}\pknBp$
and the basis is adapted). Hence the vectors $t^{-k_j}\delta'_j$ form a $\pknBp$-basis of $\pknBp\otimes
L$: their coordinate matrix $W=NY\operatorname{diag}(t^{-k_j})$ lies in $GL_d(\pknBp)$.
\end{proof}

\begin{lemma}[Filtration of the entries]\label{pkf:fil}
In the situation of Lemma~\ref{pkf:hodge}, if $-a'\le k_j\le b'$ for all $j$, then the entries of $N$
lie in $\mathrm{Fil}^{-a'}$ and those of $N^{-1}$ in $\mathrm{Fil}^{-b'}$.
\end{lemma}

\begin{proof}
$N=W\operatorname{diag}(t^{k_j})Y^{-1}$ with $W$ over $\pknBp$ and $Y^{-1}$ over $K'\subset\pknBp$; the
$(r,s)$ entry is $\sum_jW_{rj}t^{k_j}(Y^{-1})_{js}\in\mathrm{Fil}^{\min k_j}$. Likewise
$N^{-1}=Y\operatorname{diag}(t^{-k_j})W^{-1}$.
\end{proof}

\begin{theorem}[Existence of a period matrix]\label{pkf:pm}
Suppose $L[1/\ell]$ is crystalline on $\Gamma'$ with Hodge--Tate jumps in $[-a,b]$ ($a,b\ge0$). Then
there are $M\in GL_d(\pkfBm)$ with invariant columns and $U\in M_d(\pkfBp)$ whose image in $\pkfBm$ is
$(E\tau)^aM$, such that $U$ is a period matrix of height $a+b$ with Frobenius factor $v^a$.
\end{theorem}

\begin{proof}
\emph{Step 1: $M$.} By (P2) applied to $V=L[1/\ell]$, $(\pkfBm\otimes L)^{G_{K'}}$ has a $K'_0$-basis
which is a $\pkfBm$-basis of $\pkfBm\otimes L$; let $M\in GL_d(\pkfBm)$ be the matrix of its coordinates
in the basis $1\otimes e_k$. Then $M$ has invariant columns, and so has $M_n=\varphi^n(M)$ for every
$n\ge0$.

\emph{Step 2: $U$ and $U'$.} Apply Lemmas~\ref{pkf:hodge} and \ref{pkf:fil} to the image of $M_n$ in
$GL_d(\pknB)$: the entries of $\varphi^n(M)$ lie in $\mathrm{Fil}^{-a}$, those of $\varphi^n(M^{-1})=
M_n^{-1}$ in $\mathrm{Fil}^{-b}$, for every $n\ge0$. By (P4), $\mu^aM$ and $\mu^bM^{-1}$ have entries in
$\pkfBp$; as $\mu=E\tau v$ with $v$ a unit, $(E\tau)^aM$ and $(E\tau)^bM^{-1}$ are images of matrices $U$,
$U'$ over $\pkfBp$, unique by injectivity of $\pkfBp\to\pkfBm$. Then $UU'=(E\tau)^{a+b}$ (compute in
$\pkfBm$): (PM1).

\emph{Step 3: (PM2).} Let $X=M^{-1}\varphi(M)\in M_d(\pkfBm)$. For $g\in\Gamma'$, $g(M)=\rho(g)^{-1}M$
and $g(\varphi(M))=\rho(g)^{-1}\varphi(M)$, so $g(X)=M^{-1}\rho(g)\rho(g)^{-1}\varphi(M)=X$. By
Lemma~\ref{pkf:invariants}, $X$ is the image of a matrix over $\pkfSB$, again written $X$; likewise
$X'=\varphi(M)^{-1}M$, and $X'X=1$ (injectivity). Put $F=\varphi(E)^aX$ and $F'=\varphi(E)^{-a}X'$; $F$
is over $\pkfSB$ ($\varphi(E)\in\mathfrak S$), $F'$ over $\pkfSB\subseteq\pkfBp$
(Lemma~\ref{pkf:SB-basic}(i)), $F'F=1$. In $\pkfBm$: $\varphi(U)=\varphi(E)^a\varphi(\tau)^a\varphi(M)=
\varphi(E)^a\mu^aMX=v^a(E\tau)^aM\cdot\varphi(E)^aX=v^aUF$, using $\varphi(\tau)=\mu=E\tau v$. By
injectivity, $\varphi(U)=v^aUF$ in $\pkfBp$.

\emph{Step 4: (PM3).} Let $g$ fix $\Lambda$ and $\theta(z)$. By the chart, $g=j(\gamma)$ with
$\gamma\in\Gamma'$. By Lemma~\ref{pkf:approx0}(iv), $g(E)g(\tau)r''=E\tau$ with $r''\in\pkpA^\times$. Put
$R_g=r''^a\rho(\gamma)\in M_d(\pkpA)$. In $\pkfBm$: $R_g\,g(U)=r''^a(gE\,g\tau)^a\rho(\gamma)g(M)=
(E\tau)^aM=U$.

\emph{Step 5: (PM4).} Let $Y,W,k$ be as in Lemma~\ref{pkf:hodge} for $N=M$ (image in $\pknB$), so that
over $\pknB$, $U=(E\tau)^aW\operatorname{diag}(t^{k_j})Y^{-1}$, with $Y^{-1}=\operatorname{adj}Y$ over $K'$.
Suppose $Uy=E^nw$ with $w$ over $\pkfBp$. By (A8)(c), applied to the rows $\lambda=(Y^{-1})_{r\cdot}$,
there are $s_r\in\pknBp$ with $\sum_i(Y^{-1})_{ri}G_f(y_i)=s_r\sum_i(Y^{-1})_{ri}y_i$, i.e.\
$Y^{-1}G_f(y)=\operatorname{diag}(s)Y^{-1}y$. Hence, diagonal matrices commuting,
\[UG_f(y)=(E\tau)^aW\operatorname{diag}(t^{k_j})\operatorname{diag}(s)Y^{-1}y
=W\operatorname{diag}(s)W^{-1}\cdot(E\tau)^aW\operatorname{diag}(t^{k_j})Y^{-1}y
=W\operatorname{diag}(s)W^{-1}\,E^nw .\]
$W\operatorname{diag}(s)W^{-1}$ is over $\pknBp$, so the entries of $UG_f(y)$ lie in $E^n\pknBp=
\mathrm{Fil}^n$ (P6), and by Lemma~\ref{pkf:approx0}(v) they are divisible by $E^n$ in $\pkfBp$.
\end{proof}

\subsection{Stage (b): from a period matrix to a frame}\label{pkf:ss-frame}

Throughout this subsection $U$ is a period matrix of height $h$ with Frobenius factor $v_1$, with $U'$,
$F$, $F'$, $R_g$ as in Definition~\ref{pkf:def-pm}. A \emph{solution} is a pair $(x,y)$ with $x\in
\pkpA^d$, $y\in(\pkfBp)^d$ and $Uy=E^hx$. Vectors over $\pkpA$ are mapped to $\pkfBp$ entrywise.

\begin{lemma}\label{pkf:inj}
$\det U\ne0$; $U$ is injective on $(\pkfBp)^d$; $U'U=(E\tau)^h$; and for every $x'\in\pkpA^d$,
$(\tau^hx',U'x')$ is a solution.
\end{lemma}

\begin{proof}
$\det U\det U'=(E\tau)^{hd}\ne0$ in the domain $\pkfBp$ ($E,\tau\ne0$ in $\pkpA\subseteq\pkfBp$). Over the
fraction field $U'=(E\tau)^hU^{-1}$, so $U'U=(E\tau)^h$. Finally $UU'x'=E^h\tau^hx'$.
\end{proof}

\begin{lemma}[Denominators]\label{pkf:denom}
Let $T\in M_d(\pkfBp)$ and $\delta\mid\mu$ in $\pkpA$. There is $k$ such that: if $Ty=\delta^hw$ with $y$ over
$\pkfAo$ and $w$ over $\pkfBp$, then $\ell^kw$ is over $\pkfAo$.
\end{lemma}

\begin{proof}
Take $k_0$ with $\ell^{k_0}T$ over $\pkfAo$. Then $\delta^h\ell^{k_0}w=\ell^{k_0}Ty$ is over $\pkfAo$, and
(A2) gives $\ell^{2h+k_0}w$ over $\pkfAo$; $k=2h+k_0$.
\end{proof}

\begin{lemma}[Descent along Frobenius]\label{pkf:descent}
Let $(x,y)$ be a solution, $f$ an $\mathfrak S$-linear map $\pkpA\to\pkpA$ and $w$ with
$UG_f(y)=E^hw$. Then there are a solution $(x',y')$, an $\mathfrak S$-linear $f'$ and $w'$ with
$UG_{f'}(y')=E^hw'$ and $\varphi(w')=v_1E^hw$.
\end{lemma}

\begin{proof}
Put $x'=\varphi^{-1}(v_1E^hx)\in\pkpA^d$. \emph{Claim: if $r\in\pkpA$ and $Ua=rx'$ with $a$ over
$\pkfBp$, then $F\varphi(a)=\varphi(r)y$.} Indeed $\varphi(U)\varphi(a)=\varphi(r)v_1E^hx=\varphi(r)v_1Uy$,
and $\varphi(U)=v_1UF$; cancel the unit $v_1$ and use injectivity of $U$.
Apply the claim to $a=U'x'$, $r=(E\tau)^h$: $F\varphi(U'x')=\varphi(E\tau)^hy$, so $\varphi(U'x')=
\varphi(\tau)^h\varphi(E)^hF'y$. By (A8)(b) and induction on $h$ (cancelling $\varphi(\tau)=\mu\ne0$ in the
domain $\pkfBp$ at each step), $U'x'=\tau^hy'$ with $y'$ over $\pkfBp$; then $\tau^hUy'=(E\tau)^hx'$ and
$Uy'=E^hx'$: $(x',y')$ is a solution.
Let $f'=\varphi^{-1}\circ f\circ\varphi$; it is $\mathfrak S$-linear because $\varphi(\mathfrak S)\subseteq
\mathfrak S$: $f'(s\alpha)=\varphi^{-1}(\varphi(s)f(\varphi\alpha))=sf'(\alpha)$. By (PM4) there is $w'$ with
$UG_{f'}(y')=E^hw'$. The map $\varphi\circ G_{f'}-G_f\circ\varphi$ is additive, preserves $\pkfAo$ and
kills $\pkpA$, so it is $0$ (Lemma~\ref{pkf:density}). By the claim with $a=y'$, $r=E^h$:
$F\varphi(y')=\varphi(E)^hy$. Since $G_f$ is $\pkfSB$-linear (Lemma~\ref{pkf:extend}), $F$ is over
$\pkfSB$ and $\varphi(E)\in\mathfrak S$,
\[\varphi(E)^h\varphi(w')=\varphi(UG_{f'}(y'))=v_1UF\,G_f(\varphi(y'))=v_1U\,G_f(F\varphi(y'))
=v_1\varphi(E)^hUG_f(y)=v_1\varphi(E)^hE^hw .\]
Cancel $\varphi(E)^h\ne0$.
\end{proof}

\begin{lemma}[Uniform integrality]\label{pkf:integral}
There is $N\ge0$ such that whenever $(x,y)$ is a solution, $f$ is $\mathfrak S$-linear and
$UG_f(y)=E^hw$, the vector $\ell^Nw$ is (the image of a vector) over $\pkpA$.
\end{lemma}

\begin{proof}
Let $k_1$ be as in Lemma~\ref{pkf:denom} for $T=U'$, $\delta=E\tau$, and $k_2$ for $T=U$, $\delta=E$. Given the data,
$U'(Uy)=(E\tau)^hy$ with $Uy=E^hx$ over $\pkfAo$, so $\ell^{k_1}y$ is over $\pkfAo$, hence so is
$G_f(\ell^{k_1}y)$, and $U(G_f(\ell^{k_1}y))=E^h\ell^{k_1}w$ gives $\ell^{m}w$ over $\pkfAo$, $m=k_1+k_2$.
Let $V\subseteq\pkfAo$ be the set of all entries of all such $\ell^mw$. Apply
Lemma~\ref{pkf:descent}: for $\alpha=\ell^mw_i\in V$ there is $\alpha'=\ell^mw'_i\in V$ with
$\varphi(\alpha')=v_1E^h\alpha$, i.e.\ $E^h\alpha=v_1^{-1}\varphi(\alpha')$. By (A3) with $u=v_1^{-1}$,
$\ell^{h+1}V\subseteq\pkpA$; $N=h+1+m$.
\end{proof}

\begin{lemma}[The spanning matrix]\label{pkf:span}
Let $P=\{x\in\pkpA^d:\ \exists\,y\in\pkfSB^d,\ Uy=E^hx\}$. There are $Q\in M_d(\pkpA)$ and $Z\in
M_d(\pkfSB)$ with $UZ=E^hQ$ and $P=Q\,\mathfrak S^d$.
\end{lemma}

\begin{proof}
$P$ is an $\mathfrak S$-submodule. For $x\in P$ the $y$ is unique (Lemma~\ref{pkf:inj}); call it $y_x$.
For $\alpha_j\in\pkpA$ and $x_j\in P$: $\sum\alpha_jx_j=0$ iff $\sum\alpha_jy_{x_j}=0$, since
$U\sum\alpha_jy_{x_j}=E^h\sum\alpha_jx_j$, $U$ is injective and $E^h\ne0$.
\emph{No $d+1$ elements of $P$ are $\mathfrak S$-independent.} Let $x_1,\dots,x_r\in P$ be
$\mathfrak S$-independent and $\sum_j\alpha_jx_j=0$ with $\alpha_j\in\pkpA$. Then $\sum_j\alpha_jy_{x_j}=0$;
apply $G_i$ entrywise: as the entries of $y_{x_j}$ lie in $\pkfSB$ and $G_i$ is $\pkfSB$-linear,
$\sum_j\kappa_i(\alpha_j)y_{x_j}=0$, so $\sum_j\kappa_i(\alpha_j)x_j=0$ and $\kappa_i(\alpha_j)=0$ for all
$i$; by (A5), $\alpha_j=0$. So the $x_j$ are independent over the domain $\pkpA$, whence $r\le d$.
\emph{Saturation.} If $a,b\in P$ and $za=\ell b$: $\ell$ is prime in $\pkpA$ ($\pkpA/\ell=O^\flat$ is a
domain) and $\ell\nmid z$ ($\bar z=\varpi^\flat\ne0$), so $a=\ell a'$ with $a'\in\pkpA^d$, and $a'\in P$
with $y_{a'}=\ell^{-1}y_a$ ($\ell^{-1}\in\pkfSB$). By (A6), $P$ is free with a basis
$q_1,\dots,q_r$, $r\le d$. Let $Q$ have columns $q_1,\dots,q_r,0,\dots,0$ and $Z$ the columns
$y_{q_1},\dots,y_{q_r},0,\dots,0$.
\end{proof}

\begin{lemma}[Divisibility]\label{pkf:dvd}
Let $Q$ be as in Lemma~\ref{pkf:span}, $(x,y)$ a solution, $\delta\in\pkpA$ with $\ell\nmid\delta$, and
$l,g\in\pkpA^d$ with $l^{\mathsf T}Q=\delta g^{\mathsf T}$. Then $\delta\mid l\cdot x$.
\end{lemma}

\begin{proof}
Let $N$ be as in Lemma~\ref{pkf:integral} and $k$ as in Lemma~\ref{pkf:denom} for $T=U$, $\delta=E$. For
$i\in I$, (PM4) gives $w_i$ with $UG_i(y)=E^hw_i$, and $x'_i=\ell^Nw_i$ is over $\pkpA$. Since
$\ell^NG_i(y)$ is over $\pkfSB$ (Lemma~\ref{pkf:values}) and $U\ell^NG_i(y)=E^hx'_i$, $x'_i\in P$, so
$x'_i=Qs_i$ with $s_i\in\mathfrak S^d$. Fix $m$. By Lemma~\ref{pkf:coord-approx} (for each entry, then
the union of the sets) there are a finite $G$ and $r$ over $\pkfAo$ with $y-\sum_{i\in G}b_iG_i(y)=
\ell^{m+k}r$. Put $t=\ell^Nx-Q\sum_{i\in G}b_is_i\in\pkpA^d$. Then
\[U(\ell^Nr)=\ell^{-m-k}\Bigl(\ell^NUy-\sum_{i\in G}b_i\,\ell^NUG_i(y)\Bigr)=E^h\,\ell^{-m-k}t,\]
so $\ell^{-m}t$ is over $\pkfAo$ (Lemma~\ref{pkf:denom}), i.e.\ $t\in(\pkpA\cap\ell^m\pkfAm)^d=
((\ell,E)^m)^d$ by (A1). Now $l\cdot\ell^Nx-\delta\,g\cdot\sum_Gb_is_i=l\cdot t+l^{\mathsf T}Q\sum_G
b_is_i-\delta g^{\mathsf T}\sum_Gb_is_i=l\cdot t\in(\ell,E)^m$. By (A7), $\delta\mid\ell^N(l\cdot x)$; as
$\ell$ is prime and $\ell\nmid\delta$, $\delta\mid l\cdot x$.
\end{proof}

\begin{lemma}\label{pkf:detQ}
$\det Q\ne0$.
\end{lemma}

\begin{proof}
Otherwise there is $0\ne l\in\pkpA^d$ with $l^{\mathsf T}Q=0$ ($\pkpA$ a domain). For $x'\in\pkpA^d$,
$(\tau^hx',U'x')$ is a solution (Lemma~\ref{pkf:inj}); Lemma~\ref{pkf:dvd} with $\delta=z^n$
($\ell\nmid z^n$), $g=0$ gives $z^n\mid\tau^h\,l\cdot x'$ for all $n$. \emph{$\bigcap_nz^n\pkpA=0$}: if
$c=z^nc_n$ for all $n$, then $\bar c\in\bigcap_n(\varpi^\flat)^nO^\flat=0$, so $c=\ell c'$; then $\ell\mid
z^nc_n$ forces $\ell\mid c_n$ and $c'=z^n(c_n/\ell)$ for all $n$; iterating, $c\in\bigcap_k\ell^k\pkpA=0$.
Hence $\tau^h\,l\cdot x'=0$, so $l\cdot x'=0$ for all $x'$ and $l=0$.
\end{proof}

\begin{lemma}[Frobenius, saturation, Galois]\label{pkf:FSG}
(i) There is $B\in M_d(\mathfrak S)$ with $\varphi(Q)=v_1QB$; and $\ell\nmid\det B$.
(ii) If $s\in\mathfrak S^d$ and $\ell$ divides every entry of $Qs$, then $s\in\ell\mathfrak S^d$.
(iii) If $g\in G_{\Q_\ell}$ fixes $\mathfrak S$ pointwise, there is $\rho_g\in M_d(\pkpA)$ with
$Q=\rho_g\,g(Q)$. More generally, for $g$ fixing $\Lambda$ and $\theta(z)$, $Ug(Z)=E^h\rho_gg(Q)$ for some
$\rho_g\in M_d(\pkpA)$.
\end{lemma}

\begin{proof}
(i) Apply $\varphi$ to $UZ=E^hQ$ and use $\varphi(U)=v_1UF$: $U\cdot F\varphi(Z)=v_1^{-1}\varphi(E)^h
\varphi(Q)$, so $U\bigl(E^h\varphi(E)^{-h}F\varphi(Z)\bigr)=E^h\,v_1^{-1}\varphi(Q)$. The matrix in brackets
is over $\pkfSB$ (Lemmas~\ref{pkf:SB-basic}(i), \ref{pkf:values}), so the columns of $v_1^{-1}\varphi(Q)$
lie in $P=Q\mathfrak S^d$: $v_1^{-1}\varphi(Q)=QB$. Then $\varphi(\det Q)=v_1^d\det Q\det B$. If
$\ell\mid\det B$, induction gives $\ell^n\mid\det Q$ for all $n$ ($\ell^n\mid\det Q\Rightarrow\ell^{n+1}\mid
\varphi(\det Q)\Rightarrow\ell^{n+1}\mid\det Q$, $\varphi$ an automorphism fixing $\ell$), so $\det Q=0$,
contradicting Lemma~\ref{pkf:detQ}.
(ii) Write $Qs=\ell w$. Then $U(Z\ell^{-1}s)=E^hw$ with $Z\ell^{-1}s$ over $\pkfSB$, so $w\in P$, $w=Qs'$
with $s'\in\mathfrak S^d$, and $Q(s-\ell s')=0$; $\det Q\ne0$ gives $s=\ell s'$.
(iii) Let $g$ fix $\Lambda$ and $\theta(z)$, and $gE=Er$ (Lemma~\ref{pkf:approx0}(iv)). Applying $g$ to
$UZ=E^hQ$ and multiplying by $R_g$ (PM3): $Ug(Z)=R_g\,g(E)^hg(Q)=E^h\rho_gg(Q)$ with
$\rho_g=r^hR_g$. If $g$ fixes $\mathfrak S$ pointwise, it fixes $\Lambda$ and $\theta(z)=g\theta(z)$
(as $\theta(gz)=\theta(z)$), and $g(Z)=Z$ by Lemma~\ref{pkf:SB-basic}(ii); so $E^hQ=UZ=E^h\rho_gg(Q)$, and
$Q=\rho_gg(Q)$ ($E^h\ne0$, $\pkpA\to\pkfBp$ injective).
\end{proof}

\begin{lemma}[Rows with nonzero determinant]\label{pkf:rows}
Let $\mathcal R\colon S\to D$ be a ring map to a domain and $\mathcal X\subseteq S^d$ such that for
$t\in S^d$: if $\mathcal R(t\cdot x)=0$ for all $x\in\mathcal X$, then $\mathcal R(t_i)=0$ for all $i$.
Then there is $N\in M_d(S)$ with all rows in $\mathcal X$ and $\mathcal R(\det N)\ne0$.
\end{lemma}

\begin{proof}
Among $N$ with $\mathcal R(\det N)\ne0$ (e.g.\ $N=1$) choose one with the largest number of rows in
$\mathcal X$. Suppose row $k$ is not in $\mathcal X$. For $x\in\mathcal X$, the matrix $N[k\leftarrow x]$
(row $k$ replaced by $x$) has determinant $\sum_ix_i\operatorname{adj}(N)_{ik}$ and more rows in
$\mathcal X$, so by maximality $\mathcal R(x\cdot\operatorname{adj}(N)_{\cdot k})=0$. By hypothesis
$\mathcal R(\operatorname{adj}(N)_{ik})=0$ for all $i$, and $\mathcal R(\det N)=\sum_i\mathcal R(N_{ki})
\mathcal R(\operatorname{adj}(N)_{ik})=0$, a contradiction.
\end{proof}

\begin{lemma}[Minimal support]\label{pkf:artin}
Let $\mathcal V$ be a valuation ring (a domain) and $T\in M_d(\mathcal V)$ with $\det T=0$. There are
$a\in\mathcal V^d$ and $\nu$ with $a_\nu=1$, $Ta=0$, and $g(a)=a$ for every ring endomorphism $g$ of
$\mathcal V$ with $T\,g(a)=0$.
\end{lemma}

\begin{proof}
Let $w\ne0$ be a kernel vector with minimal support. The ideals $w_i\mathcal V$ are totally ordered, so
some $w_\nu\ne0$ divides all $w_i$; put $a=w/w_\nu\in\mathcal V^d$, with $Ta=0$, $a_\nu=1$ and the same
support as $w$. If $Tg(a)=0$, then $g(a)-a$ is a kernel vector with support in that of $a$ and
$(g(a)-a)_\nu=g(1)-1=0$; by minimality $g(a)-a=0$.
\end{proof}

For $m\ge0$ let $c_m\colon O^\flat\to O_C/\ell$, $c_m(b)=(b^{1/\ell^m})^\sharp\bmod\ell$; it is a ring
map (the projection of $O^\flat=\varprojlim_{x\mapsto x^\ell}O_C/\ell$ to the $m$-th factor), with
$c_{m+n}(b^{\ell^n})=c_m(b)$, $c_m(b)=0$ iff $v^\flat(b)\ge\ell^m$, and $c_m(\bar x)=\theta(\varphi^{-m}x)
\bmod\ell$ for $x\in\pkpA$ (Witt coordinates, \S\ref{pkn:ss-dist}).

\begin{lemma}[Approximation of fixed vectors]\label{pkf:fixed}
Let $a\in(O^\flat)^d$ be fixed by the group $H$ of (P5). For every $M\ge0$ there are $n\ge0$ and
$s\in\mathfrak S^d$ with $c_M(a_i)=c_{M+n}(\bar s_i)$ for all $i$.
\end{lemma}

\begin{proof}
\emph{The group $H$.} Let $\pi_n=((\varpi^\flat)^{1/\ell^n})^\sharp\in O_C$, so $\pi_0=\pi=\theta(z)$ is a
uniformiser of $K'$ and $\pi_{n+1}^\ell=\pi_n$. An element $g\in G_{\Q_\ell}$ fixes $\mathfrak S=\Lambda[[z]]$
pointwise iff it fixes $\Lambda$ and $z=[\varpi^\flat]$, i.e.\ $\varpi^\flat$ (Teichm\"uller lifts are
injective and equivariant), i.e.\ all $(\varpi^\flat)^{1/\ell^n}$ (unique $\ell$-power roots in the perfect
ring $O^\flat$), i.e.\ all $\pi_n$ (an element $b\in O^\flat$ is determined by the $(b^{1/\ell^n})^\sharp$,
and $\sharp$ is equivariant); and $g$ fixes $\Lambda=W(\mathbb F_{\ell^f})$ iff it fixes the unramified
extension $K_f=\Lambda[1/\ell]$ of degree $f$. So $H=\operatorname{Gal}(\auQlb/K_\infty)$ with
$K_\infty=\bigcup_nK_n$, $K_n=K_f(\pi_n)$.

\emph{The fields $K_n$.} $K'$ has residue degree $f$ and contains $K_f$ (as $j(\Gamma')$ fixes $\Lambda$), and
$\pi$ is a root of the Eisenstein polynomial $P\in\Lambda[X]$ of degree $e$ that differs from $E(X)$ by a
unit of $\Lambda[[X]]$ (Weierstrass preparation; Definition~\ref{pkp:def-chart}). $P(X^{\ell^n})$ is
Eisenstein over $\Lambda$, so it is irreducible, $K_n/K_f$ is totally ramified of degree $e\ell^n$, $\pi_n$ is
a uniformiser and $O_{K_n}=\Lambda[\pi_n]$ (Serre, \emph{Corps locaux}, Ch.~I, \S6, Proposition~17 and its Corollary; p.~19 of the English edition \cite{SerreLF}: for an Eisenstein polynomial $f$ over a discrete valuation ring $A$, $A[X]/(f)$ is a discrete valuation ring with uniformiser the image of $X$ and the same residue field, $f$ is irreducible, and $A[X]/(f)$ is the integral closure of $A$). Hence every element
of $O_{K_n}$ is congruent modulo $\ell$ to a finite sum $\sum_k[\lambda_k]\pi_n^k$ with
$\lambda_k\in\mathbb F_{\ell^f}$.

\emph{Approximation.} Let $a\in O^\flat$ be fixed by $H$ and $M\ge0$. Then $a^{1/\ell^M}$ is fixed by $H$,
so $x=(a^{1/\ell^M})^\sharp\in O_C$ is fixed by $H$, and by (P5) $x\in O_{\hat K_\infty}$. Choose
$y\in O_{K_\infty}$ with $x-y\in\ell O_C$, say $y\in O_{K_N}$ with $N\ge M$, and write
$y\equiv\sum_k[\lambda_k]\pi_N^k\pmod\ell$. Put $\beta=\sum_k\lambda_k^{\ell^N}(\varpi^\flat)^k\in
\mathbb F_{\ell^f}[\varpi^\flat]\subseteq O^\flat$. The ring map $c_N$ sends $\varpi^\flat$ to
$\pi_N\bmod\ell$ and $\lambda\in\mathbb F_{\ell^f}$ to $(\lambda^{1/\ell^N})^\sharp=[\lambda^{1/\ell^N}]
\bmod\ell$; so $c_N(\beta)=\sum_k[\lambda_k]\pi_N^k\bmod\ell=y\bmod\ell=x\bmod\ell=c_M(a)$. With $n=N-M$,
$c_{M+n}(\beta)=c_M(a)$. For finitely many $a_i$ take $N$ common to all (enlarging $N$ is allowed as
$O_{K_N}\subseteq O_{K_{N+1}}$) and $\beta_i$ as above. Finally $\mathbb F_{\ell^f}[\varpi^\flat]$ is the
image of $\Lambda[z]\subseteq\mathfrak S$ in $O^\flat$; take $s_i\in\Lambda[z]$ lifting $\beta_i$. Then
$c_M(a_i)=c_{M+n}(\bar s_i)$.
\end{proof}

\begin{proposition}[Key point]\label{pkf:key}
$\ell\nmid\det Q$.
\end{proposition}

\begin{proof}
Let $\mathcal R\colon\mathfrak S\to O^\flat$ be the reduction, $\mathcal X=\{x^{(i,j)}:i\in I,\ 1\le j\le d\}$
with $x^{(i,j)}_m=\kappa_i(Q_{jm})$; by $\mathfrak S$-linearity $t\cdot x^{(i,j)}=\kappa_i((Qt)_j)$ for
$t\in\mathfrak S^d$. \emph{Separation:} if $\mathcal R(\kappa_i((Qt)_j))=0$ for all $i,j$, let $c=(Qt)_j$;
for every $n$ and large $G$ (A5), $c\in\sum_G\kappa_i(c)b_i+(\ell,z)^n\subseteq\ell\pkpA+z^n\pkpA$, so
$\bar c\in\bigcap_n(\varpi^\flat)^nO^\flat=0$; thus $\ell$ divides every $(Qt)_j$, and $t\in\ell\mathfrak S^d$
by Lemma~\ref{pkf:FSG}(ii). By Lemma~\ref{pkf:rows} there are $N_\star\in M_d(\mathfrak S)$ and maps
$\sigma,\iota$ with $(N_\star t)_k=\kappa_{\sigma(k)}((Qt)_{\iota(k)})$ for all $t\in\mathfrak S^d$ and
$\mathcal R(\det N_\star)\ne0$. Put $\varkappa=\mathcal R(\det B)\mathcal R(\det N_\star)\ne0$ (Lemma~\ref{pkf:FSG}(i);
$O^\flat$ is a domain).

Suppose $\ell\mid\det Q$; then $\det\bar Q=0$. By Lemma~\ref{pkf:artin} ($O^\flat$ is a valuation ring) there
are $a$, $\nu$ with $a_\nu=1$, $\bar Qa=0$, and $g(a)=a$ whenever $\bar Qg(a)=0$. For $g\in H$:
$\bar Q=\bar\rho_gg(\bar Q)$ (Lemma~\ref{pkf:FSG}(iii)), so $\bar Qg(a)=\bar\rho_gg(\bar Qa)=0$ and
$g(a)=a$.

Let $C_0=1$, $C_{n+1}=B\varphi(C_n)\in M_d(\mathfrak S)$. \emph{Claim: $\bar Q\,\bar C_n\,a^{\ell^n}=0$
and $\varkappa\mathcal R(\det C_n)\mid \varkappa^{\ell^n}$} ($a^{\ell^n}$ entrywise). For $n=0$ this is clear. Applying the
Frobenius $b\mapsto b^\ell$ of $O^\flat$, which is the reduction of $\varphi$, to $\bar Q\bar C_na^{\ell^n}
=0$ and using $\overline{\varphi(Q)}=\bar v_1\bar Q\bar B$ gives $\bar Q\bar B\,\overline{\varphi(C_n)}
a^{\ell^{n+1}}=0$. Also $\mathcal R(\det C_{n+1})=\mathcal R(\det B)\mathcal R(\det C_n)^\ell$, and
$\varkappa\mathcal R(\det B)\mid \varkappa^2\mid \varkappa^\ell$, so $\varkappa\mathcal R(\det C_{n+1})\mid \varkappa^\ell\mathcal R(\det
C_n)^\ell\mid \varkappa^{\ell^{n+1}}$.

Fix $M\ge0$. By Lemma~\ref{pkf:fixed} choose $n$, $s$ with $c_M(a_i)=c_{M+n}(\bar s_i)$; let
$\psi=c_{M+n}\circ\mathcal R\colon\mathfrak S\to O_C/\ell$. Then $\psi(s_\nu)=c_M(1)=1$, and for every
$k$, $\psi((N_\star C_ns)_k)=0$: let $x=(QC_ns)_{\iota(k)}\in\pkpA$; as $c_{M+n}(\bar s_i)=c_M(a_i)=
c_{M+n}(a_i^{\ell^n})$, $c_{M+n}(\bar x)=c_{M+n}\bigl((\bar Q\bar C_na^{\ell^n})_{\iota(k)}\bigr)=0$, i.e.\
$\theta(\varphi^{-M-n}x)\in\ell O_C$. By (A8)(a), $\kappa_{\sigma(k)}(x)=\ell s'+p_{M+n}q$ with
$s',q\in\mathfrak S$; and $(N_\star C_ns)_k=\kappa_{\sigma(k)}(x)$ since $C_ns\in\mathfrak S^d$. Now $\psi(\ell)=0$ and
$\psi(p_{M+n})=\theta(\varphi^{-M-n}\varphi^{M+n}E)\bmod\ell=0$. Hence $\psi(N_\star C_n)\psi(s)=0$ with
$\psi(s_\nu)=1$; multiplying by $\operatorname{adj}\psi(N_\star C_n)$, $\psi(\det(N_\star C_n))=0$. So
$c_{M+n}(\varkappa\mathcal R(\det C_n))=c_{M+n}(\mathcal R\det B)\,\psi(\det N_\star\det C_n)=0$, and since
$\varkappa\mathcal R(\det C_n)\mid \varkappa^{\ell^n}$, $c_M(\varkappa)=c_{M+n}(\varkappa^{\ell^n})=0$. As $M$ is arbitrary,
$v^\flat(\varkappa)\ge\ell^M$ for all $M$ and $\varkappa=0$, a contradiction.
\end{proof}

\begin{theorem}[Frame theorem for a period matrix]\label{pkf:frame-pm}
There are $Q,Q'\in M_d(\pkpA)$, $Z\in M_d(\pkfSB)$, $Y\in M_d(\pkfBp)$ and $B\in M_d(\mathfrak S)$ with
$QQ'=\tau^h$, $UZ=E^hQ$, $YZ=E^h$, $\varphi(Q)=v_1QB$, and, for $g$ fixing $\Lambda$ and $\theta(z)$,
$g(Z)=ZG_g$ with $G_g\in M_d(\pkpA)$.
\end{theorem}

\begin{proof}
Take $Q,Z$ of Lemma~\ref{pkf:span} and $B$ of Lemma~\ref{pkf:FSG}(i).
\emph{Span: if $UY_1=E^hX$ with $X\in M_d(\pkpA)$, $Y_1\in M_d(\pkfBp)$, then $X=QL_X$ with $L_X\in
M_d(\pkpA)$.} The rows $l=\operatorname{adj}(Q)_{i\cdot}$ satisfy $l^{\mathsf T}Q=\det Q\cdot e_i^{\mathsf T}$;
Lemma~\ref{pkf:dvd} with $\delta=\det Q$ (Proposition~\ref{pkf:key}), applied to the columns of $(X,Y_1)$,
gives $\operatorname{adj}(Q)X=\det Q\cdot L_X$, and $\det Q\,(QL_X-X)=Q\operatorname{adj}(Q)X-\det Q\,X=0$.
\emph{$Q'$.} $UU'=E^h\tau^h$: so $\tau^h\cdot1=QQ'$.
\emph{$Y$.} For each entry $y=U_{rs}$, (A4) gives $\ell^{k_y}y=\alpha_y+\ell^{k_y}\mu^hw_y$. Let $n\ge$
all $k_y$, $A=(\ell^{n-k_y}\alpha_y)$ over $\pkpA$ and $W_\circ=(w_y)$ over $\pkfBp$. Then
$A=\ell^n(U-\mu^hW_\circ)=\ell^nU(1-v^{h}_{\phantom{1}}U'W_\circ)$, as $\mu^h=v^h(E\tau)^h$ and $(E\tau)^h=UU'$.
So $U\cdot E^h\ell^n(1-v^hU'W_\circ)=E^hA$ and $A=Q\gamma$ with $\gamma$ over $\pkpA$ by the span property.
Hence $U=\ell^{-n}Q\gamma+v^h(E\tau)^hW_\circ=Q\bigl(\ell^{-n}\gamma+v^hE^hQ'W_\circ\bigr)=:QY$ (using
$\tau^h=QQ'$). Then $Q(YZ-E^h)=UZ-E^hQ=0$, so $YZ=E^h$ ($\det Q\ne0$ in the domain $\pkfBp$).
\emph{Galois.} By Lemma~\ref{pkf:FSG}(iii), $Ug(Z)=E^h\rho_gg(Q)$, and the span property gives
$\rho_gg(Q)=QG_g$; so $Ug(Z)=E^hQG_g=UZG_g$, and $g(Z)=ZG_g$ by injectivity of $U$.
\end{proof}

\begin{theorem}[Existence of frames]\label{pkf:main}
Let $L$ be a lattice with jump sets and $a$ its jump bound. Then $L$ has a frame for $a$
(Definition~\ref{pkn:def-frame}). Consequently Proposition~\ref{pkn:frame} holds unconditionally (modulo
the inputs of \S\ref{pkf:ss-rings}): items (M) (third assertion) and (H) (first assertion, with $H=a$)
of Lemma~\ref{pkm:inputs}.
\end{theorem}

\begin{proof}
The Hodge--Tate jumps of $L[1/\ell]$ lie in $\bigcup_iU(i)\subseteq[-a,a]$ (P3). Theorem~\ref{pkf:pm}
with $b=a$ gives $M\in GL_d(\pkfBm)$ with invariant columns and a period matrix $U$ of height $h=2a$
whose image in $\pkfBm$ is $(E\tau)^aM$. Theorem~\ref{pkf:frame-pm} gives $Q,Q',Z,Y$ with $QQ'=
\tau^{2a}$, $UZ=E^{2a}Q$, $YZ=E^{2a}$, $Z$ over $\pkfSB$. Define $n_i$ by (Fr4) ($\tau$ is a unit of
$\pkpW$). Then (Fr1)--(Fr4) hold: (Fr1) is the invariance of the columns of $M$, (Fr2) is the
relation between $U$ and $M$, (Fr3) and (Fr4) are as stated. The rings $B^+$, $B_{\max}$ and the
maps of \S\ref{pkn:ss-frame} are those of \S\ref{pkf:ss-rings}, which have the properties listed
there by (P1).
\end{proof}

\begin{remark}\label{pkf:rem-extra}
The proof gives more than (Fr1)--(Fr4): a Frobenius matrix $B$ over $\mathfrak S$ with
$\varphi(Q)=v^aQB$ and $\ell\nmid\det B$, and the Galois relation $g(Z)=ZG_g$. These are not used in
\S\ref{pkn:ss-frame}.
\end{remark}


%% ===== proofs/C-kisin-rings.tex =====
% proofs/C-kisin-rings.tex
% Proofs of the ring inputs (A1)--(A8) and of (P4).
% Labels start with pkr:. Conventions and macros as in C-kisin-frame.tex, C-kisin-inputs.tex.
\providecommand{\auO}{\mathcal O}
\providecommand{\auQlb}{\overline{\Q}_\ell}
\providecommand{\pkpA}{A_{\mathrm{inf}}}
\providecommand{\pknBp}{B^+_{\mathrm{dR}}}
\providecommand{\pknB}{B_{\mathrm{dR}}}
\providecommand{\pkfAm}{A_{\max}}
\providecommand{\pkfAo}{A_{\max}^{\circ}}
\providecommand{\pkfBp}{B^+}
\providecommand{\pkfBm}{B_{\max}}
\providecommand{\pkfSB}{\mathfrak S_B}
% macros of this file (prefix \pkr)
\providecommand{\pkrI}{(\ell,E)}
\providecommand{\pkrm}{\mathfrak n}

\section{The period rings: proofs of the inputs (A1)--(A8)}\label{pkr:sec}

This section proves the eight auxiliary inputs (A1)--(A8) of \S\ref{pkf:ss-rings}, on which the
existence of frames (Theorem~\ref{pkf:main}) rests, and the statement (P4) listed there among the
published facts. We use the conventions and notation of \S\ref{pkf:ss-rings}:
$\pkpA=W(O^\flat)$, $R=\pkpA[X]/(\ell X-E)$, $\pkfAm$ its $\ell$-adic completion, $E/\ell$ the image
of $X$, $\pkfBp=\pkfAm[1/\ell]$, $\pkfBm=\pkfBp[1/\mu]$, $\pkfAo$ the image of $\pkfAm$ in $\pkfBp$,
the Kummer chart with $\mathfrak S=\Lambda[[z]]$, $E=z^eA+\ell$ ($A\in\mathfrak S^\times$),
$p_a=\varphi^a(E)$, $\mu=E\tau v$ with $v\in\pkpA^\times$ fixed. We write $\pkrI$ for the ideal of
$\pkpA$ generated by $\ell$ and $E$, and $\pkrm=(\ell,z)$ for the maximal ideal of $\mathfrak S$.

The following are proved in this section from (P1), (P6) and the standard facts (W1)--(W4)
recalled below: (A1) (Proposition~\ref{pkr:I1}), (A7) (Proposition~\ref{pkr:I7}), (A5)
(Proposition~\ref{pkr:I5}), (A6) (Proposition~\ref{pkr:I6}), (A3) (Proposition~\ref{pkr:I3}), (A2)
(Proposition~\ref{pkr:I2}), (A4) (Proposition~\ref{pkr:I4}), (A8) (Proposition~\ref{pkr:I8}), and
(P4) (Proposition~\ref{pkr:P4}). Of (P1) we use only: $\pkfAm\to\pkfBp\to\pkfBm\to\pknB$ and
$\pkfBp\to\pknBp$ are injective ring maps compatible with $\theta$, $\varphi$ and $G_{\Q_\ell}$, and
$\pkfBp$, $\pkfBm$ are domains (the $\ell$-torsion freeness, completeness and separation of $\pkfAm$
are proved in Lemma~\ref{pkr:Amax}).

\subsection{Standard facts and the structure of $\pkfAm$ modulo $\ell$}\label{pkr:ss-prelim}

We use the following standard facts, already used in \S\S\ref{pkn:sec}--\ref{pkf:sec}.
\begin{enumerate}
\item[(W1)] $\pkpA$ is $\ell$-torsion free, $\ell$-adically complete and separated, and
$\pkpA/\ell\pkpA=O^\flat$ via $x\mapsto x_0$ (the zeroth Witt coordinate). As $O^\flat$ is a domain,
$\ell$ is a prime element of $\pkpA$; $x\in\pkpA^\times$ iff $x_0\in O^{\flat\times}$
(\S\ref{pkn:ss-dist}).
\item[(W2)] $O^\flat$ is a valuation ring for $v^\flat$ (\S\ref{pkn:ss-dist}), complete for $v^\flat$
(Scholze \cite[Lemma~3.4(iii)]{pkfScholze12}: the tilt of a perfectoid field is a perfectoid field, hence complete, with ring of integers $\varprojlim_{x\mapsto x^\ell}O$ and valuation $x\mapsto|x^\sharp|$). Hence for $0\ne\varpi\in\mathfrak m^\flat$,
$O^\flat$ is $\varpi$-adically separated ($v^\flat(x)\ge kv^\flat(\varpi)$ for all $k$ forces $x=0$) and
complete (a sequence with $x_{k+1}-x_k\in\varpi^kO^\flat$ is $v^\flat$-Cauchy, and its limit $x$
satisfies $v^\flat(x-x_k)\ge kv^\flat(\varpi)$, i.e.\ $x-x_k\in\varpi^kO^\flat$).
\item[(W3)] $\theta\colon\pkpA\to O_C$ is surjective, and $\varphi$ is a ring automorphism of $\pkpA$
commuting with $G_{\Q_\ell}$.
\item[(W4)] $\mathfrak S=\Lambda[[z]]$ is a complete noetherian local ring with maximal ideal
$\pkrm=(\ell,z)$ and residue field $k=\mathbb F_{\ell^f}$; every ideal $\pkrm^n$ is closed for the
$\pkrm$-adic topology (it contains $\pkrm^m$ for $m\ge n$, so it is open); and
$\varphi(\mathfrak S)\subseteq\mathfrak S$ ($\varphi(z)=z^\ell$, $\varphi(\Lambda)=\Lambda$), so
$p_a\in\mathfrak S$.
\end{enumerate}
Two elementary consequences. Since $E=[E_0]+\ell E'$ (Corollary~\ref{pkn:E-gen},
Definition~\ref{pkn:def-dist}), we have $\pkrI=(\ell,t)$ with $t=[E_0]$, $E_0\in\mathfrak m^\flat$,
$v^\flat(E_0)=1$ (Lemma~\ref{pkn:dist}(i)). Since $z^e=A^{-1}(E-\ell)\in\pkrI$, we have
\begin{equation}\label{pkr:eq-zE}
\pkrm^{ek}\pkpA\subseteq\pkrI^k\qquad(k\ge0):
\end{equation}
a monomial $\ell^iz^j$ with $i+j=ek$ lies in $\ell^i\pkrI^{\lfloor j/e\rfloor}$ and
$i+\lfloor j/e\rfloor\ge\lfloor(i+j)/e\rfloor=k$.

\begin{lemma}[Structure of $\pkfAm$]\label{pkr:Amax}
\begin{enumerate}
\item[(i)] $R$ is $\ell$-torsion free; $\pkfAm=\varprojlim_nR/\ell^nR$ is $\ell$-torsion free; the kernel
of the projection $\pkfAm\to R/\ell^nR$ is $\ell^n\pkfAm$; $\pkfAm$ is $\ell$-adically complete and
separated, and $\pkfAm\to\pkfAo$ is bijective.
\item[(ii)] There is a surjective ring map $\rho\colon\pkfAm\to(\pkpA/\pkrI)[X]$ with kernel
$\ell\pkfAm$, sending $\alpha\in\pkpA$ to its class and $E/\ell$ to $X$.
\item[(iii)] Let $f\in\pkpA[X]$ be monic, $y\in\pkfAm$ and $m\ge0$. If $\ell^m$ divides $f(E/\ell)y$ in
$\pkfAm$, then $\ell^m$ divides $y$.
\end{enumerate}
\end{lemma}

\begin{proof}
(i) Let $\ell g\in(\ell X-E)\pkpA[X]$, say $\ell g=(\ell X-E)h$. Reducing modulo $\ell$ (by (W1),
$\pkpA[X]/\ell=O^\flat[X]$), $-E_0\bar h=0$ in the domain $O^\flat[X]$, and $E_0\ne0$; so $h=\ell h'$
and, $\pkpA[X]$ being $\ell$-torsion free, $g=(\ell X-E)h'$. Thus $R$ is $\ell$-torsion free. Hence for
$y\in R$, $\ell^ny\in\ell^{m+n}R$ implies $y\in\ell^mR$.

Let $x=(x_k)_k\in\pkfAm$ ($x_k\in R/\ell^kR$ compatible) with $x_k=0$ for $k\le n$. For $m\ge0$ lift
$x_{m+n}$ to $\tilde x\in R$; as $\tilde x\in\ell^nR$ write $\tilde x=\ell^n\tilde y$, and let $y_m$ be
the class of $\tilde y$ in $R/\ell^mR$; by the previous sentence $y_m$ does not depend on the
choices, and the $y_m$ are compatible. So $y=(y_m)\in\pkfAm$ and $x=\ell^ny$. Conversely $\ell^n\pkfAm$
maps to $0$ in $R/\ell^nR$. This proves the statement on kernels. If $\ell x=0$ in $\pkfAm$, then
for every $k$ a lift of $x_{k+1}$ times $\ell$ lies in $\ell^{k+1}R$, so the lift lies in $\ell^kR$ and
$x_k=0$: $\pkfAm$ is $\ell$-torsion free. Separation: if $x\in\ell^n\pkfAm$ for all $n$, then $x_n=0$ for
all $n$. Completeness: if $x^{(k+1)}-x^{(k)}\in\ell^k\pkfAm$, the components $x^{(k)}_n$ are constant
for $k\ge n$, and their common values form the limit. Finally $\pkfAm\to\pkfBp=\pkfAm[1/\ell]$ is
injective because $\pkfAm$ is $\ell$-torsion free.

(ii) $R/\ell R=\pkpA[X]/(\ell,\ell X-E)=\pkpA[X]/(\ell,E)=(\pkpA/\pkrI)[X]$; compose with the projection
$\pkfAm\to R/\ell R$, whose kernel is $\ell\pkfAm$ by (i).

(iii) $\pkpA/\pkrI\ne0$, since $\theta(\pkrI)\subseteq\ell O_C\ne O_C$. Case $m=1$: $\rho(f(E/\ell))=\bar f$
is monic in $(\pkpA/\pkrI)[X]$, hence not a zero divisor (the leading coefficient of $\bar f\,q$ is
that of $q$); so $\rho(f(E/\ell)y)=0$ gives $\rho(y)=0$, i.e.\ $\ell\mid y$. Induction: if $\ell^{m+1}
\mid f(E/\ell)y$, then $\ell\mid y$ by the case $m=1$, $y=\ell y'$, and $\ell^m\mid f(E/\ell)y'$ because
$\pkfAm$ is $\ell$-torsion free; so $\ell^m\mid y'$.
\end{proof}

We write $\theta$ also for the extension $\pkfAm\to O_C$ of Lemma~\ref{pkf:approx0}(v) (with
$\theta(E/\ell)=0$) and for its extension to $\pkfBp\to C$.

\begin{proposition}[Input (A1)]\label{pkr:I1}
For $n\ge0$, $\pkpA\cap\ell^n\pkfAm=\pkrI^n$ (intersection inside $\pkfAm$).
\end{proposition}

\begin{proof}
$\pkrI^n\subseteq\ell^n\pkfAm$ because $\ell$ and $E=\ell\cdot(E/\ell)$ lie in $\ell\pkfAm$. Conversely
we show by induction on $n$ that $a\in\pkpA$ and $\ell^n\mid a$ in $\pkfAm$ imply $a\in\pkrI^n$; the case
$n=0$ is empty. Let $\ell^{n+1}\mid a$ in $\pkfAm$. Applying $\theta$, $\theta(a)\in\ell^{n+1}O_C$; by
(W3) write $\theta(a)=\ell^{n+1}\theta(c)$ with $c\in\pkpA$. Then $\theta(a-\ell^{n+1}c)=0$, so
$a-\ell^{n+1}c=Ew$ with $w\in\pkpA$ (Corollary~\ref{pkn:E-gen}), i.e.\ $a=Ew+\ell\cdot\ell^nc$. In $\pkfAm$,
$Ew=\ell\,(E/\ell)w$, so $\ell^{n+1}\mid\ell\bigl((E/\ell)w+\ell^nc\bigr)$, and since $\pkfAm$ is $\ell$-torsion
free (Lemma~\ref{pkr:Amax}(i)), $\ell^n\mid(E/\ell)w$. By Lemma~\ref{pkr:Amax}(iii) with $f=X$,
$\ell^n\mid w$ in $\pkfAm$, so $w\in\pkrI^n$ by induction; also $\ell^nc\in\pkrI^n$. Hence
$a=Ew+\ell\cdot\ell^nc\in\pkrI\cdot\pkrI^n=\pkrI^{n+1}$.
\end{proof}

\subsection{Adic completeness of $\pkpA$ and closed principal ideals: input (A7)}\label{pkr:ss-complete}

For an element $s$ of a ring acting on a module $M$, we say $M$ is \emph{$s$-complete} if every sequence
$(x_k)_{k\ge0}$ in $M$ with $x_{k+1}-x_k\in s^kM$ has a limit $x\in M$ with $x-x_k\in s^kM$ for all $k$,
and \emph{$s$-separated} if $\bigcap_ks^kM=0$.

\begin{lemma}[Extensions]\label{pkr:ext}
Let $M'\subseteq M$ be modules over a commutative ring, $M''=M/M'$, and $s$ an element acting
injectively on $M''$. If $M'$ and $M''$ are $s$-complete and $s$-separated, so is $M$, and
$s^kM\cap M'=s^kM'$ for all $k$.
\end{lemma}

\begin{proof}
If $s^km\in M'$, then $s^k\bar m=0$ in $M''$, so $\bar m=0$ and $m\in M'$: this gives $s^kM\cap M'=
s^kM'$. Separation: if $x\in\bigcap_ks^kM$, its image lies in $\bigcap_ks^kM''=0$, so $x\in M'\cap s^kM=
s^kM'$ for all $k$ and $x=0$. Completeness: let $x_{k+1}-x_k\in s^kM$. The images $\bar x_k$ have a
limit $\bar x$ in $M''$; lift it to $x'\in M$. Then $x'-x_k=s^kr_k+m_k$ with $r_k\in M$, $m_k\in M'$.
Now $m_{k+1}-m_k=(x_k-x_{k+1})+s^kr_k-s^{k+1}r_{k+1}\in s^kM\cap M'=s^kM'$, so $(m_k)$ has a limit $m$ in
$M'$ with $m-m_k\in s^kM'$. Put $x=x'-m$; then $x-x_k=s^kr_k+(m_k-m)\in s^kM$.
\end{proof}

\begin{proposition}[$\pkrI$-adic completeness of $\pkpA$]\label{pkr:complete}
(i) $\bigcap_{k\ge0}\pkrI^k=0$ and $\bigcap_k\pkrm^k\pkpA=0$.
(ii) If $(a_k)_{k\ge0}$ in $\pkpA$ satisfies $a_{k+1}-a_k\in\pkrI^k$ for all $k$, there is $b\in\pkpA$
with $b-a_k\in\pkrI^{\lfloor k/2\rfloor}$ for all $k$.
\end{proposition}

\begin{proof}
Let $t=[E_0]$, so that $\pkrI=(\ell,t)$, and $M_n=\pkpA/\ell^n\pkpA$.

\emph{Step 1: $t$ acts injectively on $M_n$, and $M_n$ is $t$-complete and $t$-separated.} If
$tx\in\ell^n\pkpA$, then since $\ell$ is prime (W1) and $\ell\nmid t$ ($t_0=E_0\ne0$), $\ell^n\mid x$: so
$t$ is injective on $M_n$. Completeness and separation by induction on $n$. For $n=1$, $M_1=O^\flat$
with $t$ acting as $E_0\in\mathfrak m^\flat\smallsetminus\{0\}$, and (W2) applies. For $n\ge2$, the
submodule $M'=\ell M_n$ is isomorphic to $M_{n-1}$ via $x\mapsto\ell x$ (as $\pkpA$ is $\ell$-torsion
free), compatibly with $t$, and $M_n/M'=M_1$; Lemma~\ref{pkr:ext} applies.

\emph{Step 2.} For $0\le n\le k$, $\pkrI^k\subseteq\ell^n\pkpA+t^{k-n+1}\pkpA$: in a monomial $\ell^it^{k-i}$
either $i\ge n$ or $k-i\ge k-n+1$. Also $t^jx\in t^{j'}M_n$ with $j\le j'$ implies $x\in t^{j'-j}M_n$
(Step 1, injectivity).

\emph{Step 3: (i).} If $x\in\bigcap_k\pkrI^k$, then by Step 2 the class of $x$ in $M_n$ lies in
$\bigcap_kt^{k-n+1}M_n=0$; so $x\in\bigcap_n\ell^n\pkpA=0$ (W1). The second assertion follows by
\eqref{pkr:eq-zE}.

\emph{Step 4: (ii).} Fix $n\ge1$. For $k\ge n$, Step 2 gives $a_{k+1}-a_k\in\ell^n\pkpA+t^{k-n+1}\pkpA$,
so the classes $\bar a_{k}$ in $M_n$ satisfy $\bar a_{k+1}-\bar a_k\in t^{k-n+1}M_n$. By Step 1 (applied
to the sequence $j\mapsto\bar a_{j+n-1}$) there is $\beta_n\in M_n$ with $\beta_n-\bar a_k\in
t^{k-n+1}M_n$ for $k\ge n$. The image of $\beta_{n+1}$ in $M_n$ satisfies the same condition with
$t^{k-n}$ in place of $t^{k-n+1}$; two elements of $M_n$ both congruent to $\bar a_k$ modulo
$t^{k-n}M_n$ for all large $k$ agree, by $t$-separation. So the $\beta_n$ are compatible, and as
$\pkpA=\varprojlim_n M_n$ (W1) there is $b\in\pkpA$ with image $\beta_n$ in every $M_n$. Thus $b-a_k\in
\ell^n\pkpA+t^{k-n+1}\pkpA$ for $1\le n\le k$. Taking $n=\lceil k/2\rceil$ (for $k\ge1$; the case $k=0$
is trivial): $b-a_k\in\ell^{\lceil k/2\rceil}\pkpA+t^{\lfloor k/2\rfloor+1}\pkpA\subseteq\pkrI^{\lfloor
k/2\rfloor}$.
\end{proof}

\begin{remark}\label{pkr:rem-complete}
Part (ii) gives convergence modulo $\pkrI^{\lfloor k/2\rfloor}$ rather than modulo $\pkrI^k$. This is all
that is used: part (ii) is applied only in Step~2 of the proof of Proposition~\ref{pkr:I3}, where it is
needed only that $b-a_k\in\ell^{c_k}\pkfAm$ with $c_k\to\infty$, and $\pkrI^{\lfloor k/2\rfloor}\subseteq
\ell^{\lfloor k/2\rfloor}\pkfAm$. Part (i) is used in Propositions~\ref{pkr:I5}, \ref{pkr:I6} and at the
beginning of \S\ref{pkr:ss-div}, in the form stated. (The $\pkrI$-adic completeness in the usual sense
also holds, since $\pkrI^{2k}\subseteq\ell^k\pkpA+t^{k+1}\pkpA\subseteq\pkrI^k$ shows that the
$\pkrI$-adic topology is the one defined by the ideals $\ell^k\pkpA+t^{k+1}\pkpA$, and Step~4 constructs
limits for the latter; it is not needed.)
\end{remark}

\begin{lemma}[Saturated ideals]\label{pkr:sat}
Let $J$ be an ideal of $\pkpA$ such that $\ell x\in J$ implies $x\in J$, and let $\delta\in J$ be such
that $\delta_0$ divides $x_0$ in $O^\flat$ for every $x\in J$. Then $J=\delta\pkpA$.
\end{lemma}

\begin{proof}
Let $x\in J$. As $\delta_0\mid x_0$, write $x_0=\delta_0c$ and $y_0=[c]$; then $x-\delta y_0\equiv
[x_0]-[\delta_0][c]=0\pmod{\ell}$ (W1), so $x-\delta y_0=\ell x^{(1)}$, and $x^{(1)}\in J$ by saturation
(as $x-\delta y_0\in J$). Iterating, $x=\delta(y_0+\ell y_1+\dots+\ell^ky_k)+\ell^{k+1}x^{(k+1)}$ for all
$k$. By $\ell$-adic completeness and separation of $\pkpA$, $y=\sum_k\ell^ky_k$ converges and $x=\delta
y$.
\end{proof}

\begin{proposition}[Input (A7)]\label{pkr:I7}
Let $\delta,c\in\pkpA$ with $\ell\nmid\delta$, and suppose that for every $m$ there is $z_m\in\pkpA$
with $c-\delta z_m\in\pkrI^m$. Then $\delta\mid c$.
\end{proposition}

\begin{proof}
\emph{Step 1: for every $m\ge1$ there is $N_m$ with $\pkrI^{N_m}\subseteq(\ell^m,\delta)$.} Since
$\ell\nmid\delta$, $\delta_0\ne0$; and $E_0\in\mathfrak m^\flat$ has $v^\flat(E_0)=1>0$. Choose $r$ with
$r\ge v^\flat(\delta_0)$; then $\delta_0\mid E_0^r$ in the valuation ring $O^\flat$, say $E_0^r=\delta_0c'$,
and $E^r-\delta[c']\equiv0\pmod\ell$ (W1). So $E^r\in(\ell,\delta)$, hence $E^{rm}\in(\ell,\delta)^m\subseteq
(\ell^m,\delta)$. Put $N_m=m(r+1)$: a monomial $\ell^iE^{N_m-i}$ has $i\ge m$ or $N_m-i\ge rm$, so it lies
in $(\ell^m,\delta)$.

\emph{Step 2.} Let $J=\bigcap_{m\ge1}(\ell^m,\delta)$. Then $c\in J$: $c=(c-\delta z_{N_m})+\delta z_{N_m}
\in\pkrI^{N_m}+\delta\pkpA\subseteq(\ell^m,\delta)$. Clearly $\delta\in J$. $J$ is saturated: if
$\ell x\in J$, then for each $m$, $\ell x=\ell^{m+1}b+\delta y$ for some $b,y$, so $\ell\mid\delta y$,
and as $\ell$ is prime with $\ell\nmid\delta$, $y=\ell y'$; cancelling $\ell$ ($\pkpA$ torsion free),
$x=\ell^mb+\delta y'\in(\ell^m,\delta)$. Finally every $x\in J$ lies in $(\ell,\delta)$, so
$x=\ell b+\delta y$ and $x_0=\delta_0y_0$. By Lemma~\ref{pkr:sat}, $J=\delta\pkpA$, so $\delta\mid c$.
\end{proof}

\subsection{Topological bases over $\mathfrak S$: inputs (A5) and (A6)}\label{pkr:ss-basis}

For an $\mathfrak S$-module $M$ and $n\ge0$, $\pkrm^nM=\sum_{i+j=n}\ell^iz^jM$. For a set $I$ let
$F_I$ be the free module of finitely supported families $f=(f_i)_{i\in I}$ in $\mathfrak S$.

\begin{proposition}[Topological bases]\label{pkr:basis}
Let $M$ be an $\mathfrak S$-module such that
\begin{enumerate}
\item[(M1)] $\ell$ acts injectively on $M$;
\item[(M2)] if $x\in M$ and $zx\in\ell M$, then $x\in\ell M$;
\item[(M3)] $\bigcap_n\pkrm^nM=0$.
\end{enumerate}
Let $T\subseteq M$ be a subset whose image in the $k$-vector space $M/\pkrm M$ is linearly independent.
Then there are a family $(b_i)_{i\in I}$ in $M$ containing $T$ and $\mathfrak S$-linear maps
$\kappa_i\colon M\to\mathfrak S$ ($i\in I$) such that: (a) $\kappa_i(\sum_jf_jb_j)=f_i$ for $f\in F_I$;
(b) for every $x\in M$ and $N\ge0$ there is a finite $G_0\subseteq I$ with $x-\sum_{i\in G}\kappa_i(x)b_i
\in\pkrm^NM$ for every finite $G\supseteq G_0$; (c) if $\kappa_i(x)=0$ for all $i$, then $x=0$.
\end{proposition}

\begin{proof}
Extend the image of $T$ to a $k$-basis of $M/\pkrm M$ and lift the added vectors: this gives
$(b_i)_{i\in I}\supseteq T$ whose image is a $k$-basis of $M/\pkrm M$. Let $\beta\colon F_I\to M$,
$\beta(f)=\sum_if_ib_i$.

\emph{Step 1 (density): for $x\in M$ and $N\ge0$ there is $f\in F_I$ with $x-\beta(f)\in\pkrm^NM$.}
For $N=1$ this holds because the $b_i$ span $M/\pkrm M$ over $k=\mathfrak S/\pkrm$. If $x-\beta(f)=
\sum_{i+j=N}\ell^iz^jy_{ij}$, choose $g_{ij}$ with $y_{ij}-\beta(g_{ij})\in\pkrm M$; then
$x-\beta(f+\sum\ell^iz^jg_{ij})\in\pkrm^{N+1}M$.

\emph{Step 2 (a criterion).} Let $a,s\in\mathfrak S$ be such that (C1) $\beta(f)\in(a,s)M$ implies
$f\in(a,s)F_I$, and (H) $aw=sy$ with $w,y\in M$ implies $y\in aM$. Then for every $n\ge1$,
$\beta(f)\in(s^n,a)M$ implies $f\in(s^n,a)F_I$. Induction on $n$, the case $n=1$ being (C1). Let
$\beta(f)=s^{n+1}w'+au$ with $w',u\in M$. By (C1), $f=au'+sc$ with $u',c\in F_I$. Then
$a(u-\beta(u'))=s(\beta(c)-s^nw')$, so by (H) $\beta(c)-s^nw'=aw$ with $w\in M$, i.e.\ $\beta(c)\in
(s^n,a)M$. By induction $c=s^nm+am'$ with $m,m'\in F_I$, and $f=s^{n+1}m+a(u'+sm')$.

\emph{Step 3: $\beta(f)\in\pkrm^{2n}M$ implies $f\in\pkrm^nF_I$.} First take $(a,s)=(\ell,z)$: (C1) is
the linear independence of the images of the $b_i$ in $M/\pkrm M$ over $k$, and (H) is (M2). So
$\beta(f)\in(z^n,\ell)M$ implies $f\in(z^n,\ell)F_I$. Next take $(a,s)=(z^n,\ell)$: (C1) is what was
just proved; (H) says that $z^nw=\ell y$ implies $y\in z^nM$, which we prove by induction on $n$ (the
case $n=0$ is trivial): $z(z^{n-1}w)=\ell y$ gives $z^{n-1}w=\ell c$ by (M2); then $\ell(zc-y)=0$, so
$y=zc$ by (M1); and $c\in z^{n-1}M$ by induction, so $y\in z^nM$. Hence $\beta(f)\in(\ell^n,z^n)M$
implies $f\in(\ell^n,z^n)F_I$. Since $\pkrm^{2n}\subseteq(\ell^n,z^n)\subseteq\pkrm^n$, the claim
follows.

\emph{Step 4 (coordinates).} Let $x\in M$. An \emph{approximation of level $N$} of $x$ is an $f\in F_I$
with $x-\beta(f)\in\pkrm^{2N}M$; they exist by Step 1. If $f$ and $f'$ are approximations of levels
$N\le N'$, then $\beta(f-f')\in\pkrm^{2N}M$, so $f_i-f'_i\in\pkrm^N$ for all $i$ (Step 3). Choosing an
approximation $f^{(N)}$ of level $N$ for each $N$, $(f^{(N)}_i)_N$ is a Cauchy sequence in $\mathfrak S$;
let $\kappa_i(x)$ be its limit (W4). Then $\kappa_i(x)-f_i\in\pkrm^N$ for every approximation $f$ of
level $N$ (as $\pkrm^N$ is closed), so $\kappa_i(x)$ does not depend on the choices. If $f,g$ are
approximations of level $N$ of $x,y$ and $s\in\mathfrak S$, then $f+g$ and $sf$ are approximations of
level $N$ of $x+y$ and $sx$; hence $\kappa_i$ is $\mathfrak S$-linear. For $x=\beta(f)$, the constant
choice $f^{(N)}=f$ gives (a).

(b) Let $f$ be an approximation of level $N$ of $x$ and $G_0=\operatorname{supp}f$. For finite
$G\supseteq G_0$,
\[x-\sum_{i\in G}\kappa_i(x)b_i=(x-\beta(f))-\sum_{i\in G}(\kappa_i(x)-f_i)b_i\in\pkrm^{2N}M+\pkrm^NM
\subseteq\pkrm^NM.\]
(c) follows from (b) and (M3).
\end{proof}

\begin{proposition}[Input (A5)]\label{pkr:I5}
There are a family $(b_i)_{i\in I}$ in $\pkpA$ and $\mathfrak S$-linear maps $\kappa_i\colon\pkpA\to
\mathfrak S$ such that for every $x\in\pkpA$ and $N\ge0$, $x-\sum_{i\in G}\kappa_i(x)b_i\in
\pkrm^N\pkpA=(\ell,z)^N\pkpA$ for all sufficiently large finite $G\subseteq I$ (in particular $x=0$ if
all $\kappa_i(x)=0$); and there is $0\in I$ with $b_0=1$ and $\kappa_0(s)=s$ for $s\in\mathfrak S$.
\end{proposition}

\begin{proof}
Apply Proposition~\ref{pkr:basis} to $M=\pkpA$ and $T=\{1\}$. (M1) is (W1). (M2): if $zx=\ell y$, then
$\ell\mid zx$, and $\ell$ is prime with $\ell\nmid z$ ($z_0=\varpi^\flat\ne0$), so $\ell\mid x$. (M3) is
Proposition~\ref{pkr:complete}(i). The image of $1$ in $\pkpA/\pkrm\pkpA$ is nonzero: if $1=\ell
x+zy$, then $1=\varpi^\flat y_0$ in $O^\flat$, so $\varpi^\flat$ would be a unit, while
$\varpi^\flat\in\mathfrak m^\flat$. Let $0\in I$ be the index with $b_0=1$; for $s\in\mathfrak S$,
$s=\beta(f)$ with $f$ supported at $0$, $f_0=s$, so $\kappa_0(s)=s$ by (a).
\end{proof}

\begin{proposition}[Input (A6)]\label{pkr:I6}
Let $P\subseteq\pkpA^d$ be an $\mathfrak S$-submodule containing no $d+1$ elements linearly
independent over $\mathfrak S$, such that $za\in\ell P$ with $a\in P$ implies $a\in\ell P$. Then $P$
is a free $\mathfrak S$-module of finite rank $\le d$.
\end{proposition}

\begin{proof}
Apply Proposition~\ref{pkr:basis} to $M=P$ and $T=\emptyset$. (M1): $\pkpA^d$ is $\ell$-torsion free.
(M2) is the hypothesis. (M3): $\pkrm^nP\subseteq(\pkrm^n\pkpA)^d$, and $\bigcap_n\pkrm^n\pkpA=0$
(Proposition~\ref{pkr:complete}(i)). Let $(b_i)_{i\in I}$, $(\kappa_i)$ be as given there. Any finitely
many distinct $b_i$ are linearly independent over $\mathfrak S$: if $\sum_{i\in G}f_ib_i=0$, then
$f_i=\kappa_i(0)=0$ by (a). Hence $|I|\le d$. For $x\in P$, (b) with $G=I$ gives $x-\sum_{i\in I}
\kappa_i(x)b_i\in\bigcap_N\pkrm^NP=0$. So $(b_i)_{i\in I}$ is a basis of $P$.
\end{proof}

\subsection{Division in $\pkfAm$ and approximation by $\pkpA$: inputs (A3), (A2), (A4)}\label{pkr:ss-div}

We identify $\pkfAm$ with $\pkfAo\subseteq\pkfBp$ (Lemma~\ref{pkr:Amax}(i)), and $\pkpA$ with its image
in $\pkfAm$ (injective by (P1), or by Proposition~\ref{pkr:I1} and Proposition~\ref{pkr:complete}(i):
the kernel lies in $\bigcap_n\pkrI^n$).

\begin{lemma}[The element $\tau$]\label{pkr:tau}
$\tau$ divides $E-\ell v^{-1}$ in $\pkpA$. In $\pkfAm$, $E-\ell v^{-1}=\ell(E/\ell-v^{-1})$.
\end{lemma}

\begin{proof}
By Lemma~\ref{pkn:mu}, $\mu=\tau\xi_\varepsilon$ with $\xi_\varepsilon=\sum_{k=0}^{\ell-1}
[\varepsilon^{1/\ell}]^k=Eu$, $u\in\pkpA^\times$; since also $\mu=E\tau v$ and $\pkpA$ is a domain
($\mu\ne0$), $u=v$. Now $\xi_\varepsilon-\ell=\sum_{k<\ell}([\varepsilon^{1/\ell}]^k-1)$ and each term is
divisible by $[\varepsilon^{1/\ell}]-1=\tau$. So $\tau\mid Ev-\ell$, hence $\tau\mid v^{-1}(Ev-\ell)=
E-\ell v^{-1}$.
\end{proof}

\begin{lemma}[Division]\label{pkr:div}
Let $\nu,b\in\pkpA$ with $\nu\mid E-\ell b$. For every $n\ge0$ and $y\in\pkfAm$ there are $a\in\pkpA$
and $w\in\pkfAm$ with $\ell^ny=a+\nu^nw$. This applies to $\nu=E$ ($b=0$) and to $\nu=\tau$
($b=v^{-1}$, Lemma~\ref{pkr:tau}).
\end{lemma}

\begin{proof}
\emph{Step 1: every $y\in\pkfAm$ is $c+(E/\ell-b)w$ with $c\in\pkpA$, $w\in\pkfAm$.} By
Lemma~\ref{pkf:approx0}(i), $y=P(E/\ell)+\ell y_1$ with $P\in\pkpA[X]$, $y_1\in\pkfAm$. Division by the
monic $X-b$ gives $P(X)=P(b)+(X-b)Q(X)$, so $y=c_0+(E/\ell-b)w_0+\ell y_1$ with $c_0=P(b)\in\pkpA$,
$w_0=Q(E/\ell)$. Repeating with $y_1,y_2,\dots$:
$y=\sum_{k<K}\ell^kc_k+(E/\ell-b)\sum_{k<K}\ell^kw_k+\ell^Ky_K$ for all $K$. The series
$c=\sum\ell^kc_k$ converges in $\pkpA$ and $w=\sum\ell^kw_k$ in $\pkfAm$ ((W1), Lemma~\ref{pkr:Amax}(i)),
and the map $\pkpA\to\pkfAm$ carries $\ell^K\pkpA$ into $\ell^K\pkfAm$; so $y-c-(E/\ell-b)w\in\ell^K
\pkfAm$ for every $K$, hence is $0$.

\emph{Step 2.} Write $E-\ell b=\nu g$, $g\in\pkpA$; in $\pkfAm$, $\ell(E/\ell-b)=\nu g$. Induction on
$n$, the case $n=0$ being $a=0$, $w=y$. If $\ell^ny=a+\nu^nw$, write $w=c+(E/\ell-b)w'$ by Step 1; then
$\ell^{n+1}y=\ell a+\ell\nu^nc+\nu^n\ell(E/\ell-b)w'=(\ell a+\ell\nu^nc)+\nu^{n+1}gw'$.
\end{proof}

\begin{lemma}[Divisibility by $E^h$]\label{pkr:Ediv}
If $b\in\pkpA$ and $E^h\mid b$ in $\pkfBp$, then $E^h\mid b$ in $\pkpA$.
\end{lemma}

\begin{proof}
Induction on $h$. If $b=E^hw$ with $h\ge1$, $w\in\pkfBp$, then $\theta(b)=0$, so $b=Eb_1$ with
$b_1\in\pkpA$ (Corollary~\ref{pkn:E-gen}); then $E(b_1-E^{h-1}w)=0$ in the domain $\pkfBp$ (P1), and
$E\ne0$; so $E^{h-1}\mid b_1$ in $\pkfBp$, and by induction in $\pkpA$.
\end{proof}

\begin{proposition}[Input (A3): contraction]\label{pkr:I3}
Let $V\subseteq\pkfAo$, $h\ge0$ and $u\in\pkpA$ be such that every $\alpha\in V$ has some $\alpha'\in V$
with $E^h\alpha=u\,\varphi(\alpha')$. Then $\ell^{h+1}V\subseteq\pkpA$.
\end{proposition}

\begin{proof}
\emph{Step 1: for every $k\ge0$ and $\alpha\in V$ there are $a\in\pkpA$ and $w\in\pkfAm$ with
$a=\ell^{h+1}\alpha-\ell^k\tau^{h+1}w$.} For $k=0$ this is Lemma~\ref{pkr:div} with $\nu=\tau$,
$n=h+1$. Suppose it holds for $k$ (for all elements of $V$), and let $\alpha\in V$, with $\alpha'$ as
in the hypothesis; choose $a',w'$ for $\alpha'$ and $k$: $a'=\ell^{h+1}\alpha'-\ell^k\tau^{h+1}w'$.
Apply $\varphi$ (which acts on $\pkfAm$ compatibly with $\pkpA$ and $\pkfBp$, (P1)) and use
$\varphi(\tau)=\mu=E\tau v$:
$\varphi(a')=\ell^{h+1}\varphi(\alpha')-\ell^kE^{h+1}\tau^{h+1}v^{h+1}\varphi(w')$. Multiply by $u$ and
use $u\varphi(\alpha')=E^h\alpha$ and $E=\ell\cdot E/\ell$:
\[u\varphi(a')=E^h\bigl(\ell^{h+1}\alpha-\ell^{k+1}\tau^{h+1}w\bigr),\qquad w=(E/\ell)\,v^{h+1}u\,
\varphi(w')\in\pkfAm.\]
So $E^h\mid u\varphi(a')$ in $\pkfBp$, and by Lemma~\ref{pkr:Ediv} $u\varphi(a')=E^ha$ with $a\in\pkpA$.
Cancelling $E^h$ in the domain $\pkfBp$: $a=\ell^{h+1}\alpha-\ell^{k+1}\tau^{h+1}w$.

\emph{Step 2.} Fix $\alpha\in V$ and choose $a_k,w_k$ as in Step 1 for every $k$. For $k\le l$,
$a_k-a_l=\ell^k\tau^{h+1}(\ell^{l-k}w_l-w_k)\in\pkpA\cap\ell^k\pkfAm=\pkrI^k$
(Proposition~\ref{pkr:I1}). By Proposition~\ref{pkr:complete}(ii) there is $b\in\pkpA$ with
$b-a_k\in\pkrI^{\lfloor k/2\rfloor}\subseteq\ell^{\lfloor k/2\rfloor}\pkfAm$. Then
$\ell^{h+1}\alpha-b=(a_k-b)+\ell^k\tau^{h+1}w_k\in\ell^{\lfloor k/2\rfloor}\pkfAm$ for every $k$, so
$\ell^{h+1}\alpha=b\in\pkpA$ by separation (Lemma~\ref{pkr:Amax}(i)).
\end{proof}

\begin{proposition}[Input (A2)]\label{pkr:I2}
If $\delta\in\pkpA$ divides $\mu$, $h\ge0$, $w\in\pkfBp$ and $\delta^hw\in\pkfAo$, then $\ell^{2h}w\in\pkfAo$.
\end{proposition}

\begin{proof}
By Lemma~\ref{pkr:tau}, $E-\ell v^{-1}=\tau s$ with $s\in\pkpA$, so $E(E-\ell v^{-1})v=E\tau v\,s=\mu s$
is divisible by $\delta$: write $E(E-\ell v^{-1})v=\delta\,r$ with $r\in\pkpA$. In $\pkfAm$,
$E(E-\ell v^{-1})v=\ell^2g$ with $g=v\,(E/\ell)(E/\ell-v^{-1})=v\,f(E/\ell)$, $f=X(X-v^{-1})$ monic.
Choose $m$ with $\ell^m\ell^{2h}w=y\in\pkfAo$, and let $\delta^hw=c\in\pkfAo$. In $\pkfBp$,
\[g^hy=\ell^m(\ell^2g)^hw=\ell^m(\delta r)^hw=\ell^mr^hc,\]
an identity in $\pkfAm$ (injectivity). So $\ell^m\mid v^hf(E/\ell)^hy$, hence (as $v$ is a unit)
$\ell^m\mid f^h(E/\ell)\,y$ with $f^h$ monic, and $\ell^m\mid y$ by Lemma~\ref{pkr:Amax}(iii): $y=\ell^my'$
with $y'\in\pkfAm$, and $\ell^{2h}w=y'$ ($\ell$ is a unit of $\pkfBp$).
\end{proof}

\begin{proposition}[Input (A4): approximation]\label{pkr:I4}
For $y\in\pkfBp$ and $h\ge0$ there are $m\ge0$, $\alpha\in\pkpA$ and $w\in\pkfAo$ with
$\ell^my=\alpha+\ell^m\mu^hw$.
\end{proposition}

\begin{proof}
By Lemma~\ref{pkf:approx0}(ii) with $n=2h$: $\ell^my=\alpha_0+\ell^{m+2h}w_0$ with $\alpha_0\in\pkpA$,
$w_0\in\pkfAo$. By Lemma~\ref{pkr:div} with $\nu=\tau$: $\ell^hw_0=a_1+\tau^hw_1$; with $\nu=E$:
$\ell^hw_1=a_2+E^hw_2$ ($a_1,a_2\in\pkpA$, $w_1,w_2\in\pkfAm$). Then
\[\ell^{2h}w_0=\ell^ha_1+\tau^ha_2+(E\tau)^hw_2=\ell^ha_1+\tau^ha_2+\mu^hv^{-h}w_2,\]
and $\ell^my=\alpha+\ell^m\mu^hw$ with $\alpha=\alpha_0+\ell^m(\ell^ha_1+\tau^ha_2)\in\pkpA$ and
$w=v^{-h}w_2\in\pkfAo$.
\end{proof}

\subsection{Divisibility by $\mu$ and $\tau$, the chart facts (A8), and (P4)}\label{pkr:ss-chart}

\begin{lemma}[Divisibility by $\mu$ in $\pkfBp$]\label{pkr:mu-Bplus}
Let $c\in\pkfBp$ with $\theta(\varphi^nc)=0$ for all $n\ge0$. Then $\mu\mid c$ in $\pkfBp$.
\end{lemma}

\begin{proof}
By Proposition~\ref{pkr:I4} with $h=1$, $\ell^mc=\alpha+\ell^m\mu w$ with $\alpha\in\pkpA$, $w\in\pkfAo$.
For $n\ge0$, $\theta(\varphi^n\mu)=\theta([\varepsilon^{\ell^n}])-1=(\varepsilon^{\ell^n})^\sharp-1=0$, and
$\theta\circ\varphi^n$ is a ring map on $\pkfBp$ extending the one on $\pkpA$ (P1); so
$\theta(\varphi^n\alpha)=\ell^m\theta(\varphi^nc)-\ell^m\theta(\varphi^n\mu)\theta(\varphi^nw)=0$ for all
$n\ge0$. By Lemma~\ref{pkn:mu}, $\mu\mid\alpha$ in $\pkpA$. Hence $\mu\mid\ell^mc$, and $\mu\mid c$ since
$\ell$ is a unit of $\pkfBp$.
\end{proof}

\emph{The field $K'$ inside $\pknBp$.} As in \S\ref{pkf:sec}, $K'$ (indeed $\auQlb$) is regarded as a
subfield of $\pknBp$ through Fontaine's embedding $\iota$, characterised by $\theta\circ\iota=\mathrm{id}$;
and $K_0=\Lambda[1/\ell]$ is regarded as a subfield of $\auQlb$ through $\theta$. By the proof of
Lemma~\ref{pkf:invariants}, $K'=K_0(\pi)$ with $\pi=\theta(z)$.

\begin{lemma}[Hensel uniqueness and Newton approximation in $\pknBp$]\label{pkr:newton}
(i) Let $P$ be a separable polynomial with coefficients in a subfield $F\subseteq\pknBp$ with $\theta|_F$
injective, and $x,x'\in\pknBp$ roots of $P$ with $\theta(x)=\theta(x')$. Then $x=x'$.
(ii) The image of $K_0=\Lambda[1/\ell]$ in $\pknBp$ through $\pkpA[1/\ell]$ is $\iota(K_0)$, elementwise:
for $c\in\Lambda[1/\ell]$, the image of $c$ equals $\iota(\theta(c))$.
(iii) For every $\lambda\in K'$ and $m\ge0$ there is $s\in\pkfSB$ with $E^m\mid s-\lambda$ in $\pknBp$.
\end{lemma}

\begin{proof}
Recall that $\pknBp$ is a local ring with maximal ideal $\ker\theta=E\pknBp$ (Lemma~\ref{pkn:BdR}), so
an element with $\theta\ne0$ is a unit.
(i) By Taylor's formula, $P(x')-P(x)=(x'-x)\bigl(P'(x)+(x'-x)r\bigr)$ with $r\in\pknBp$. The second
factor has $\theta$-image $P^\theta{}'(\theta x)\ne0$, where $P^\theta$ is $P$ with $\theta$ applied to the
coefficients (separable, with the simple root $\theta x$). So it is a unit and $x'=x$.
(ii) $\Lambda=W(\mathbb F_{\ell^f})$ contains a primitive $(\ell^f-1)$-th root of unity $\zeta$ (a
Teichm\"uller representative), and $K_0=\Q_\ell(\zeta)$. The image of $\zeta$ and $\iota(\theta(\zeta))$
are roots of $X^{\ell^f-1}-1$ (separable over $\Q_\ell$) with the same $\theta$; by (i) they are equal.
Both maps are $\Q_\ell$-algebra maps on $K_0=\Q_\ell[\zeta]$, so they agree.
(iii) Write $\lambda=\sum_{j<e}c_j\pi^j$ with $c_j\in K_0$ ($K'=K_0(\pi)$ has degree $e$ over $K_0$: $K_0$ is the maximal unramified subextension of $K'$, which has ramification index $e$, Lemma~\ref{pkf:invariants}; equivalently, $\pi$ is a root of the Eisenstein factor of $E(X)$, Definition~\ref{pkp:def-chart}) and put $s_0=\sum_jc_jz^j\in\pkfBp$. It lies in
$\pkfSB$, since $z\in\mathfrak S$, $\Lambda\subseteq\mathfrak S$ and $\ell^{-1}\in\pkfSB$
(Lemma~\ref{pkf:SB-basic}(i)). By (ii), $\theta(s_0)=\sum_jc_j\pi^j=\lambda=\theta(\iota\lambda)$, so
$E\mid s_0-\lambda$ in $\pknBp$. Let $P\in K_0[X]$ be the minimal polynomial of $\lambda$ (separable,
characteristic $0$), and $g,g_1\in K_0[X]$ with $g_1P+gP'=1$. Put $\Phi=X-gP\in K_0[X]$. Then
$\Phi(\lambda)=\lambda$ and $\Phi'(\lambda)=1-g'(\lambda)P(\lambda)-g(\lambda)P'(\lambda)=g_1(\lambda)
P(\lambda)=0$, so by Taylor's formula $\Phi(x)-\lambda=(x-\lambda)^2r_x$ with $r_x\in\pknBp$ for every
$x\in\pknBp$ (the coefficients of $\Phi$ lie in $K_0\subseteq\pknBp$, compatibly by (ii)). Put
$s_{n+1}=\Phi(s_n)$; as $\pkfSB$ is a ring containing $K_0$, $s_n\in\pkfSB$, and by induction
$E^{2^n}\mid s_n-\lambda$. Take $s=s_n$ with $2^n\ge m$.
\end{proof}

\begin{proposition}[Input (A8)]\label{pkr:I8}
(a) In the notation of (A5): for $i\in I$, $x\in\pkpA$ and $m\ge0$ with $\theta(\varphi^{-m}x)\in\ell
O_C$, there is $s\in\mathfrak S$ with $p_m\mid\kappa_i(x)-\ell s$ in $\mathfrak S$.
(b) If $b\in\pkfBp$ and $\varphi(\tau)\mid\varphi(b)$ in $\pkfBp$, then $\tau\mid b$ in $\pkfBp$.
(c) Let $G_f$ be the extension of Lemma~\ref{pkf:extend}. For $\lambda\in K'^d$, $y\in(\pkfBp)^d$ and an
$\mathfrak S$-linear $f\colon\pkpA\to\pkpA$, $\sum_i\lambda_iy_i$ divides $\sum_i\lambda_iG_f(y_i)$ in
$\pknBp$ (images in $\pknBp$).
\end{proposition}

\begin{proof}
(a) Write $\theta_m=\theta\circ\varphi^{-m}\colon\pkpA\to O_C$, a surjective ring map (W3) with kernel
$p_m\pkpA$ (Corollary~\ref{pkn:E-gen}), and $\theta_m(\ell)=\ell$. By hypothesis $\theta_m(x)=\ell c$ with
$c\in O_C$; choose $a\in\pkpA$ with $\theta_m(a)=c$. Then $\theta_m(x-\ell a)=0$, so $x-\ell a=p_mq$ with
$q\in\pkpA$. As $\kappa_i$ is $\mathfrak S$-linear and $\ell,p_m\in\mathfrak S$ (W4),
$\kappa_i(x)-\ell\kappa_i(a)=p_m\kappa_i(q)$; take $s=\kappa_i(a)$.

(b) Write $\varphi(b)=\mu c$ with $c\in\pkfBp$ ($\varphi(\tau)=\mu$). For $n\ge0$,
$\theta(\varphi^{n+1}b)=\theta(\varphi^n\mu)\theta(\varphi^nc)=0$ (proof of Lemma~\ref{pkr:mu-Bplus}). By
Proposition~\ref{pkr:I4} with $h=1$, $\ell^mb=\alpha+\ell^m\mu w$; as in the proof of
Lemma~\ref{pkr:mu-Bplus}, $\theta(\varphi^{n+1}\alpha)=0$ for all $n\ge0$, i.e.\ $\varphi(\alpha)$
satisfies the hypothesis of Lemma~\ref{pkn:mu}, so $\varphi(\tau)=\mu\mid\varphi(\alpha)$ in $\pkpA$ and,
$\varphi$ being an automorphism of $\pkpA$ (W3), $\tau\mid\alpha$. Also $\tau\mid\mu$. Hence
$\tau\mid\ell^mb$ in $\pkfBp$, and $\tau\mid b$.

(c) Put $x=\sum_i\lambda_iy_i$ and $x'=\sum_i\lambda_iG_f(y_i)$ in $\pknBp$. As $\pknBp$ is a discrete
valuation ring with uniformiser $E$ (Lemma~\ref{pkn:BdR}), it suffices to show: for every $m\ge0$,
$E^m\mid x$ implies $E^m\mid x'$ (if $x\ne0$ take $m$ its valuation; if $x=0$, $x'\in\bigcap_mE^m\pknBp=0$).
Fix $m$ and choose by Lemma~\ref{pkr:newton}(iii) $s_i\in\pkfSB$ with $E^m\mid s_i-\lambda_i$ in
$\pknBp$. Then $E^m$ divides $x-\sum_is_iy_i$ and $x'-\sum_is_iG_f(y_i)$ in $\pknBp$. If $E^m\mid x$, then
$E^m$ divides the image of $\sum_is_iy_i\in\pkfBp$ in $\pknBp$, hence divides it in $\pkfBp$
(Lemma~\ref{pkf:approx0}(v) and (P6)): $\sum_is_iy_i=E^mq$, $q\in\pkfBp$. Since $G_f$ is
$\pkfSB$-linear (Lemma~\ref{pkf:extend}) and $s_i,E\in\pkfSB$,
$\sum_is_iG_f(y_i)=G_f(\sum_is_iy_i)=E^mG_f(q)$. Hence $E^m\mid x'$ in $\pknBp$.
\end{proof}

\begin{proposition}[(P4) proved]\label{pkr:P4}
If $b\in\pkfBm$ and $m\ge0$ are such that the image of $\varphi^n(b)$ in $\pknB$ lies in
$\mathrm{Fil}^{-m}$ for every $n\ge0$, then $\mu^mb\in\pkfBp$.
\end{proposition}

\begin{proof}
As $\pkfBm=\pkfBp[1/\mu]$, there are $N\ge0$ and $c\in\pkfBp$ with $\mu^{m+N}b=c$. We show by induction
on $N$ that $\mu^mb\in\pkfBp$. If $N=0$ there is nothing to prove. Let $N\ge1$. For $n\ge0$,
$\varphi^n(c)=\varphi^n(b)\varphi^n(\mu)^{m+N}$ in $\pkfBm$, hence in $\pknB$ (P1). Since
$\theta(\varphi^n\mu)=0$, $\varphi^n(\mu)\in E\pkpA$ (Corollary~\ref{pkn:E-gen}), whose image lies in
$\mathrm{Fil}^1$; so $\varphi^n(c)\in\mathrm{Fil}^{-m}\cdot\mathrm{Fil}^{m+N}\subseteq\mathrm{Fil}^N
\subseteq\mathrm{Fil}^1$ ($\mathrm{Fil}^k=E^k\pknBp$, Lemma~\ref{pkn:BdR}). As $\theta$ on $\pkfBp$ is the
restriction of $\theta$ on $\pknBp$, whose kernel is $\mathrm{Fil}^1$, $\theta(\varphi^nc)=0$ for all
$n\ge0$. By Lemma~\ref{pkr:mu-Bplus}, $c=\mu d$ with $d\in\pkfBp$, and $\mu^{m+N-1}b=d$ in the domain
$\pkfBm$ (P1). Conclude by induction.
\end{proof}

With Propositions~\ref{pkr:I1}--\ref{pkr:P4}, the inputs (A1)--(A8) of \S\ref{pkf:ss-rings} and the
statement (P4) are proved; the existence of frames (Theorem~\ref{pkf:main}) therefore rests only on
(P1)--(P3), (P5), (P6) and the standard facts (W1)--(W4).


%% ===== proofs/C-kisin-lattices.tex =====
% proofs/C-kisin-lattices.tex
% The remaining items of Lemma pkm:inputs: (H) (second and third assertions), (F), (D),
% (C)(d), (R), (G). Labels start with pkl:. Conventions and macros as in C-kisin-inputs.tex.
\providecommand{\auO}{\mathcal O}
\providecommand{\auQlb}{\overline{\Q}_\ell}
\providecommand{\pkpA}{A_{\mathrm{inf}}}
\providecommand{\pkpW}{W(C^\flat)}
\providecommand{\pkpAmb}{\mathbb A}
\providecommand{\pknMinf}{\mathfrak M_{\mathrm{inf}}}
\providecommand{\pknK}{\mathcal K_0}
\providecommand{\pknBp}{B^+_{\mathrm{dR}}}
\providecommand{\pknB}{B_{\mathrm{dR}}}
\providecommand{\pkfBp}{B^+}
\providecommand{\pkfBm}{B_{\max}}
\providecommand{\pkfSB}{\mathfrak S_B}
% macros of this file (prefix \pkl)
\providecommand{\pklAN}{\pkpA\mathcal N}
\providecommand{\pklL}{\pkpA\otimes L}
\providecommand{\pklstd}{\mathrm{std}}

\section{Lattices with orders: standard orders, Galois stability and the retraction}\label{pkl:sec}

This section continues \S\S\ref{pkn:sec}--\ref{pkr:sec} and uses their conventions: a fixed Kummer
chart (Definition~\ref{pkp:def-chart}) with $\mathfrak S=\Lambda[[z]]$, $E$, $p_a=\varphi^a(E)$, $\Gamma'$,
$K'=\auQlb^{\,j(\Gamma')}$; the modules $\pknMinf(L)\subseteq\pklAN(L)\supseteq\mathcal N(L)$ and
$L_i(o)$ of Definition~\ref{pkn:def-N}; $\mathrm{Fil}^n=t^n\pknBp=E^n\pknBp$ (Lemma~\ref{pkn:BdR}); the
rings $\pkfBp\subseteq\pkfBm$ and $\pkfSB$ of \S\ref{pkf:ss-rings}--\S\ref{pkf:ss-SB}; frames
(Definition~\ref{pkn:def-frame}), which exist with $a$ the jump bound (Theorem~\ref{pkf:main}). We write
$\pklL$ for the $\pkpA$-span of the $1\otimes y$ in $\pkpAmb(L)=\pkpW\otimes_{\Z_\ell}L$, and
$\pkpAmb(L)_i$ for the letter part (Definition~\ref{pkp:def-letters}); $\sigma(i)=i^\ell$. We prove the
items of Lemma~\ref{pkm:inputs} not proved there: (F), (H), (D), (C)(d), (R) and (G).

\emph{Dependencies.} Proved in full, modulo the results cited from \S\S\ref{pkn:sec}--\ref{pkr:sec} (hence
modulo the published inputs listed there) and the reading of ``Hodge--Tate jumps'' recorded in
\S\ref{pkl:ss-gap}: Proposition~\ref{pkl:std} and (F) (Corollary~\ref{pkl:F}); (H), first assertion in
full (Corollary~\ref{pkl:H1}), second and third assertions by
Proposition~\ref{pkl:H2} and Corollary~\ref{pkl:H3}; the $\Gamma'$-stability part of (G), first assertion
(Lemma~\ref{pkl:G1}); (D) (Proposition~\ref{pkl:D}); (C)(d) (Proposition~\ref{pkl:Cd}); (R)
(Proposition~\ref{pkl:R}); $\varphi(\mathcal N)\subseteq\mathcal N$ (Corollary~\ref{pkl:Phi}); the second assertion of (G) and (G$_0$) (Proposition~\ref{pkl:G0}).

\subsection{Preliminaries}\label{pkl:ss-pre}

Throughout, $L$ is a lattice with jump sets, $e_1,\dots,e_d$ a $\Z_\ell$-basis of $L$, and $\lambda_0\colon
\pkpW\to\pknB$ a fixed $\pkpA$-linear map with $\lambda_0(1)=1$ (it exists by the proof of
Lemma~\ref{pkn:frame-id}(iii)). For $k\le d$ let $\lambda_k$ be the $\pkpA$-linear map
$\lambda_k(\sum_sw_s\otimes e_s)=\lambda_0(w_k)$ on $\pkpAmb(L)=\bigoplus_s\pkpW\otimes e_s$; so
$\lambda_k(1\otimes e_s)=\delta_{ks}$. Recall that $\pkpA\to\pknBp$ is injective (\S\ref{pkn:ss-frame}),
so nonzero elements of $\pkpA$ are invertible in the field $\pknB$; and that an element of $\pknB$ fixed by
$j(\Gamma')$ lies in $K'$ (proof of Lemma~\ref{pkn:frame-id}(i)).

\begin{lemma}\label{pkl:basic}
\begin{enumerate}
\item[(i)] An element $x\in\pkpA$ maps to a unit of $\pknBp$ iff $\theta(x)\ne0$; for $r\in\pkpA$ and $n\ge0$,
$r\in\mathrm{Fil}^n$ iff $E^n\mid r$ in $\pkpA$. The images of $p_c$ ($c\ge1$), of $\varphi^{-b}(\tau)$ and
$\varphi^{-b}(z)$ ($b\ge0$) are units of $\pknBp$; for $b\ge1$, $\varphi^b(\tau)\in\mathrm{Fil}^1$. For
$l\ge0$, $\varphi^l(\tau)=\tau\prod_{b<l}p_b\cdot w_l$ with $w_l\in\pkpA^\times$.
\item[(ii)] If $x\in\pkpAmb(L)$ has $\tau^nx\in\pklL$ and $x\in\pknMinf(L)$, then $E^n\varphi(x)\in
\pknMinf(L)$. Hence $\varphi$ maps $\pknMinf(L)$ and $\pklAN(L)$ into $\pklAN(L)$.
\item[(iii)] For $x\in\pklAN(L)$, a form $\mu$ and $m\ge1$: $\mu(\varphi^mx)\in\pknBp$.
\item[(iv)] $\pknMinf(L)$, $\pklAN(L)$ and $\mathcal N(L)$ are stable under $1\otimes\zeta$ and under the
projectors $e_i$.
\item[(v)] (Shift.) If $x\in L_{i^{(-l)}}(o)$, then $\varphi^lx\in L_i(\mathrm{sh}_{l,0}(o))$.
\end{enumerate}
\end{lemma}

\begin{proof}
(i) $\pknBp$ is local with maximal ideal $\mathrm{Fil}^1=\ker\theta$ (Lemma~\ref{pkn:BdR}). If $r\in\pkpA\cap
\mathrm{Fil}^1$, then $\theta(r)=0$ and $r=Er'$, $r'\in\pkpA$ (Corollary~\ref{pkn:E-gen}); as $E$ is a uniformiser,
$r\in\mathrm{Fil}^n$ forces $r'\in\mathrm{Fil}^{n-1}$; induct. $\theta(p_c)\ne0$ for $c\ne0$
(Corollary~\ref{pkn:E-gen}); $\theta(\varphi^{-b}\tau)=\zeta_{\ell^{b+1}}-1\ne0$ and $\theta(\varphi^{-b}z)$ is an
$\ell^b$-th root of $\pi\ne0$, as $\theta([c])=c^\sharp$ and $\varphi^{-b}[c]=[c^{1/\ell^b}]$;
$\theta(\varphi^b\tau)=\zeta_\ell^{\ell^b}-1=0$ for $b\ge1$. Finally, with $\xi_\varepsilon=Eu$ as in
Lemma~\ref{pkn:mu}, $\varphi^{b+1}(\tau)=\varphi^b(\tau\xi_\varepsilon)=\varphi^b(\tau)p_b\varphi^b(u)$; induct
on $l$.

(ii) $\varphi(\tau)=\tau Eu$ with $u$ a unit (Lemma~\ref{pkn:mu}), and $\varphi(b\otimes y)=\varphi(b)\otimes y$
preserves $\pklL$; so $\tau^nE^n\varphi(x)=u^{-n}\varphi(\tau^nx)\in\pklL$. For $m\ge1$ and a form $\mu$,
$\mu(\varphi^m(E^n\varphi x))=\varphi^m(E)^n\mu(\varphi^{m+1}x)\in\pknBp$. So $E^n\varphi(x)\in\pknMinf$ and
$\varphi(x)\in\pklAN$ (with $r=p_0^n$). If $x\in\pklAN$ with $rx\in\pknMinf$, $r=\prod_kp_{b_k}$, then
$\varphi(r)\varphi(x)\in\pklAN$, $\varphi(r)=\prod_kp_{b_k+1}$, and $\pklAN$ is stable under $p_b^{-1}$
(Lemma~\ref{pkn:M}(iii)).

(iii) With $r$ as above, $\varphi^m(r)\mu(\varphi^mx)=\mu(\varphi^m(rx))\in\pknBp$ and $\varphi^m(r)=
\prod_kp_{b_k+m}$ is a unit of $\pknBp$ by (i).

(iv) $1\otimes\zeta$ is $\pkpW$-linear, commutes with $\varphi$ and with $\Gamma$ (the action on $L$ is
$\auO$-linear) and preserves $\pklL$; for a form $\mu$, $\mu\circ(1\otimes\zeta)$ is a form. This gives the
stability of the three modules (for $\mathcal N$ the value condition is preserved since
$\mu\circ(1\otimes\zeta)$ is a form). $e_i$ is a polynomial in $1\otimes\zeta$ with coefficients in
$\Lambda\subseteq\mathfrak S$, and the three modules are $\mathfrak S$-modules (Lemma~\ref{pkn:M}(i)).

(v) $\varphi^lx\in\pklAN$ by (ii) and $\varphi^lx\in\pkpAmb(L)_i$ (Definition~\ref{pkp:def-letters}). Let $a\ge0$,
$u\in U(i^{(-a)})$ and $\mu$ of type $u$ at $i^{(-a)}$; the value is $\mu(\varphi^{l-a}x)$. If $a<l$ it lies in
$\pknBp=\mathrm{Fil}^0$ by (iii). If $a\ge l$, then $i^{(-a)}=(i^{(-l)})^{(-(a-l))}$ and the value
$\mu(\varphi^{-(a-l)}x)$ lies in $\mathrm{Fil}^{o(a-l,u)}$.
\end{proof}

\subsection{Hodge--Tate jumps, the gap rule and the Hodge expansion}\label{pkl:ss-gap}

For $n\in\Z$ let $C(n)$ be $C$ with $g\in G_{K'}$ acting by $\chi(g)^ng$, $\chi$ the cyclotomic character;
$g(t)=\chi(g)t$, so $\mathrm{gr}^n\pknB\cong C(n)$. We assume $\auO[1/\ell]$ is a field, so that
$C\otimes_{\Z_\ell}\auO=\prod_\gamma C$ over the embeddings $\gamma\colon\auO\to C$, and
$C\otimes_{\Z_\ell}L=\bigoplus_\gamma C\otimes_{\auO,\gamma}L$, a $G_{K'}$-stable decomposition.

\emph{Reading of Definition~\ref{pkp:def-orders}.} ``$U(i)$ contains the Hodge--Tate jumps of $L$ at all
embeddings of letter $i$'' is used in the following form: \emph{if $\gamma$ is of letter $i$ and there is a
nonzero $C$-linear $G_{K'}$-equivariant map $C\otimes_{\auO,\gamma}L\to C(-u)$, then $u\in U(i)$.} By the
Hodge--Tate decomposition $C\otimes_{\auO,\gamma}L\cong\bigoplus_kC(-k)^{m_{\gamma,k}}$ of the de Rham
representation $L[1/\ell]$ and Tate's theorem $C(n)^{G_{K'}}=0$ for $n\ne0$ (Tate \cite[\S3, Theorem~2, p.~176]{pklTate67}; also quoted in \cite[1.5.7]{pkfFontaineII}), such a map exists iff $m_{\gamma,u}\ne0$; with the normalisation of (P3) of \S\ref{pkf:ss-rings}
($\pknBp\otimes V=\bigoplus_jt^{-k_j}\pknBp\delta'_j$, whence $C\otimes V\cong\bigoplus_jC(-k_j)$) the $k$
with $m_{\gamma,k}\ne0$ are the jumps $k$ of (P3). This is the convention of Theorem~\ref{pkf:main}.

\begin{lemma}[Gap]\label{pkl:gap}
Let $j$ be a letter, $u\in\Z\smallsetminus U(j)$ and $\mu$ a form of type $u$ at $j$. Then $\mu$ is of
type $u-1$ at $j$.
\end{lemma}

\begin{proof}
For $y\in L$, $t^u\mu(e_j(1\otimes y))\in\pknBp$; put $f(y)=\theta(t^u\mu(e_j(1\otimes y)))$, a $\Z_\ell$-linear
map $L\to C$. For $g\in\Gamma'$: $e_j(1\otimes gy)=g(e_j(1\otimes y))$ ($g$ fixes $\Lambda$ and commutes with
$1\otimes\zeta$), and $\mu(gw)=g\mu(w)$ for $w\in\pklL$ ($\mu(g(b\otimes y))=g(b)\,g\mu(1\otimes y)$); so
$t^u\mu(e_j(1\otimes gy))=\chi(g)^{-u}g\bigl(t^u\mu(e_j(1\otimes y))\bigr)$ and $f(gy)=\chi(g)^{-u}gf(y)$. Also
$(1\otimes\zeta)e_j=[j]e_j$, so $f(\zeta y)=\theta([j])f(y)$. The $C$-linear extension $F(c\otimes y)=cf(y)$
of $f$ to $C\otimes_{\Z_\ell}L$ is $G_{K'}$-equivariant into $C(-u)$ (here $G_{K'}=j(\Gamma')$: $j(\Gamma')$ is compact, hence closed, and $K'$ is its fixed field) and satisfies $F(\zeta w)=\theta([j])F(w)$.
On the summand $C\otimes_{\auO,\gamma}L$, $\zeta$ acts by $\gamma(\zeta)$; if $\gamma(\zeta)\ne\theta([j])$, then
$(\gamma(\zeta)-\theta([j]))F(w)=0$ and $F$ vanishes there. If $\gamma(\zeta)=\theta([j])$, $\gamma$ is of letter
$j$ (Definition~\ref{pkp:def-letters}) and $F$ vanishes there because $u\notin U(j)$. So $f=0$, i.e.\
$t^u\mu(e_j(1\otimes y))\in\ker\theta=\mathrm{Fil}^1$ and $\mu(e_j(1\otimes y))\in\mathrm{Fil}^{1-u}$.
\end{proof}

\begin{lemma}[Gap rule]\label{pkl:gap-rule}
Let $x\in L_i(o)$, $a\ge0$, $j=i^{(-a)}$, $\nu$ a form of type $\kappa$ at $j$, and $n\in\Z$ with
$o(a,v)\ge n$ for every $v\in U(j)$ with $v\le\kappa$. Then $\nu(\varphi^{-a}x)\in\mathrm{Fil}^n$.
\end{lemma}

\begin{proof}
Suppose not. Then $\nu$ is of type $v$ at $j$ for every $v\le\kappa$: by downward induction, if $\nu$ is of type
$v\le\kappa$ and $v\in U(j)$, then $\nu(\varphi^{-a}x)\in\mathrm{Fil}^{o(a,v)}\subseteq\mathrm{Fil}^n$, which is
excluded; so $v\notin U(j)$ and Lemma~\ref{pkl:gap} gives type $v-1$. Take $r\ne0$ with $rx\in\pklL$
(Lemma~\ref{pkn:M}(ii)). Then $w=\varphi^{-a}(rx)\in\pklL\cap\pkpAmb(L)_j$; writing $w=\sum b_k\otimes y_k$,
$w=e_jw=\sum b_ke_j(1\otimes y_k)$, so $\nu(w)\in\mathrm{Fil}^{-v}$ for all $v\le\kappa$, and $\nu(w)\in
\bigcap_m\mathrm{Fil}^m=0$. As $\nu(w)=\varphi^{-a}(r)\nu(\varphi^{-a}x)$ with $\varphi^{-a}(r)\ne0$ in the field
$\pknB$, $\nu(\varphi^{-a}x)=0\in\mathrm{Fil}^n$, a contradiction.
\end{proof}

\begin{lemma}[Hodge expansion]\label{pkl:expand}
There are forms $\nu_1,\dots,\nu_d$ and integers $k_1,\dots,k_d$ with:
\begin{enumerate}
\item[(i)] $\nu_m(\pklL)\subseteq\mathrm{Fil}^{-k_m}$; in particular $\nu_m$ is of type $k_m$ at every letter;
\item[(ii)] for every $\pkpA$-linear $\lambda\colon\pkpAmb(L)\to\pknB$ there are $D_m(\lambda)\in\pknB$ with
$\lambda(x)=\sum_mD_m(\lambda)\nu_m(x)$ for every $x\in\pkpAmb(L)$ such that $rx\in\pklL$ for some
$r\in\pkpA\smallsetminus\{0\}$ (e.g.\ $x\in\pklAN(L)$); if $\lambda(1\otimes L)\subseteq\mathrm{Fil}^0$ then
$D_m(\lambda)\in\mathrm{Fil}^{k_m}$, and if $\lambda$ is a form then $D_m(\lambda)\in K'$.
\end{enumerate}
\end{lemma}

\begin{proof}
Let $M$ be the matrix of a frame (Theorem~\ref{pkf:main}) and $N$ its image in $GL_d(\pknB)$; $N$ has invariant
columns by (Fr1). Lemma~\ref{pkf:hodge} gives $H=NY=W\operatorname{diag}(t^{k_1},\dots,t^{k_d})$ with $Y\in
SL_d(K')$, $W\in GL_d(\pknBp)$; $H$ has invariant columns ($Y$ is fixed by $G_{K'}$), so $g(H^{-1})=
H^{-1}\rho(g)$. Put $\nu_m(\sum_sw_s\otimes e_s)=\sum_s\lambda_0(w_s)(H^{-1})_{ms}$; as in the proof of
Lemma~\ref{pkn:frame-id}(iii), $\nu_m$ is a form, and $\nu_m(1\otimes e_s)=(H^{-1})_{ms}=t^{-k_m}(W^{-1})_{ms}\in
\mathrm{Fil}^{-k_m}$. By $\pkpA$-linearity $\nu_m(\pklL)\subseteq\mathrm{Fil}^{-k_m}$, and $e_j(1\otimes L)\subseteq
\pklL$. (ii) Put $c_s=\lambda(1\otimes e_s)$ and $D_m(\lambda)=\sum_sc_sH_{sm}$. Then
$\sum_mD_m(\lambda)\nu_m(1\otimes e_s)=(cHH^{-1})_s=c_s$, so both sides of the identity agree on $\pklL$ by
linearity, and on $x$ with $rx\in\pklL$ after multiplication by $r$, which is invertible in $\pknB$. Since
$H_{sm}=W_{sm}t^{k_m}\in\mathrm{Fil}^{k_m}$, $c_s\in\mathrm{Fil}^0$ gives $D_m(\lambda)\in\mathrm{Fil}^{k_m}$.
If $\lambda$ is a form, $g(c)=c\rho(g)$ for $g\in\Gamma'$ (proof of Lemma~\ref{pkn:frame-id}(i)) and
$g(H)=\rho(g)^{-1}H$, so $g(cH)=cH$ and $D_m(\lambda)\in K'$.
\end{proof}

\subsection{The standard orders; items (F) and (H)}\label{pkl:ss-std}

\begin{proposition}\label{pkl:std}
For $l\ge0$ and a letter $i$: $L_i(\pklstd_l)=\varphi^l\bigl(\pknMinf(L)\cap\pkpAmb(L)_{i^{(-l)}}\bigr)
=\varphi^l(\pknMinf(L))\cap\pkpAmb(L)_i$.
\end{proposition}

\begin{proof}
The two right sides agree because $\varphi^l$ is bijective on $\pkpAmb(L)$ and maps $\pkpAmb(L)_{i^{(-l)}}$ onto
$\pkpAmb(L)_i$.

$\supseteq$. Let $x\in\pknMinf\cap\pkpAmb(L)_{i^{(-l)}}$, $\tau^nx\in\pklL$. We show $x\in L_{i^{(-l)}}(o)$ with
$o(b,u)=-u$; then $\varphi^lx\in L_i(\mathrm{sh}_{l,0}(o))=L_i(\pklstd_l)$ by Lemma~\ref{pkl:basic}(v). Clearly
$x\in\pklAN$. Let $b\ge0$, $j=i^{(-l-b)}$, $u\in U(j)$, $\mu$ of type $u$ at $j$. The vector
$w=\varphi^{-b}(\tau^nx)$ lies in $\pklL\cap\pkpAmb(L)_j$, so $w=e_jw$ is an $\pkpA$-combination of elements of
$e_j(1\otimes L)$ and $\mu(w)\in\mathrm{Fil}^{-u}$. Since $\mu(w)=\varphi^{-b}(\tau)^n\mu(\varphi^{-b}x)$ and
$\varphi^{-b}(\tau)$ is a unit of $\pknBp$ (Lemma~\ref{pkl:basic}(i)), $\mu(\varphi^{-b}x)\in\mathrm{Fil}^{-u}$.

$\subseteq$. Let $y\in L_i(\pklstd_l)$ and $x=\varphi^{-l}y\in\pkpAmb(L)_{i^{(-l)}}$. We show $x\in\pknMinf$.

\emph{Values.} Let $m\ge1$, $\mu$ a form, and $z=\varphi^my$. By Lemma~\ref{pkl:basic}(v) (at the letter
$\sigma^m(i)$, with $\sigma^m(i)^{(-m)}=i$), $z\in L_{\sigma^m(i)}(o')$ with $o'=\mathrm{sh}_{m,0}(\pklstd_l)$,
and $o'(l,v)=0$ for all $v$: if $l<m$ by definition, and if $l\ge m$ because $o'(l,v)=\pklstd_l(l-m,v)$ with
$l-m<l$. By Lemma~\ref{pkl:expand}, $\mu(\varphi^mx)=\mu(\varphi^{-l}z)=\sum_kD_k(\mu)\nu_k(\varphi^{-l}z)$ (note
$r\varphi^{-l}z\in\pklL$ for $r=\varphi^{-l}(r')$, $r'z\in\pklL$), with $D_k(\mu)\in K'\subseteq\mathrm{Fil}^0$, and
$\nu_k(\varphi^{-l}z)\in\mathrm{Fil}^0$ by Lemma~\ref{pkl:gap-rule} (with $a=l$, $n=0$, $\kappa=k_k$). So
$\mu(\varphi^mx)\in\pknBp$.

\emph{Denominators.} Put $\tau_l=\varphi^l(\tau)$. We show: \emph{every $y\in L_i(\pklstd_l)$ has $\tau_l^Ny\in
\pklL$ for some $N$}; applying $\varphi^{-l}$ then gives $\tau^Nx\in\pklL$. First a sublemma: \emph{if $y\in
L_i(\pklstd_l)$, $a\ge l$, $s\in\pkpA$ and $sp_ay\in\pklL$, then $sy\in\pklL$.} Write $sp_ay=\sum_kc_k\otimes e_k$.
By Lemma~\ref{pkl:expand} applied to $\lambda_k$ (with $\lambda_k(1\otimes L)\subseteq\Z_\ell$), $\lambda_k(\varphi^{-a}
y)=\sum_mD_m(\lambda_k)\nu_m(\varphi^{-a}y)$ with $D_m(\lambda_k)\in\mathrm{Fil}^{k_m}$, and $\nu_m(\varphi^{-a}y)\in
\mathrm{Fil}^{-k_m}$ by Lemma~\ref{pkl:gap-rule} (as $\pklstd_l(a,v)=-v\ge-k_m$ for $v\le k_m$); so $\lambda_k
(\varphi^{-a}y)\in\mathrm{Fil}^0$. Then $\varphi^{-a}(c_k)=\lambda_k(\varphi^{-a}(sp_ay))=\varphi^{-a}(s)\,E\,
\lambda_k(\varphi^{-a}y)\in\mathrm{Fil}^1$, so $\theta(\varphi^{-a}c_k)=0$ and $p_a\mid c_k$ (Corollary~\ref{pkn:E-gen});
cancelling the unit $p_a$ of $\pkpW$, $sy\in\pklL$. Now let $r=p_{a_1}\cdots p_{a_k}$ with $ry\in\pknMinf$; we
induct on $k$ (for all $y$ at once). If $k=0$, $\tau^Ny\in\pklL$ for some $N$ and $\tau\mid\tau_l$
(Lemma~\ref{pkl:basic}(i)). If $k\ge1$, $y'=p_{a_1}y$ has orders $\pklstd_l+[a=a_1]\ge\pklstd_l$
(Lemma~\ref{pkn:E}), so $y'\in L_i(\pklstd_l)$ (Lemma~\ref{pkm:le}(i)) with $k-1$ factors, and
$\tau_l^Np_{a_1}y\in\pklL$. If $a_1<l$, then $p_{a_1}\mid\tau_l$ (Lemma~\ref{pkl:basic}(i)) and
$\tau_l^{N+1}y\in\pklL$; if $a_1\ge l$, the sublemma with $s=\tau_l^N$ gives $\tau_l^Ny\in\pklL$.
\end{proof}

\begin{corollary}[Lemma~\ref{pkm:inputs}(F)]\label{pkl:F}
$L_i(\pklstd_0)=\pknMinf(L)\cap\pkpAmb(L)_i$, and for $l\ge0$, $\varphi^l$ maps $L_i(\pklstd_0)$ bijectively onto
$L_{\sigma^l(i)}(\pklstd_l)$.
\end{corollary}

\begin{proof}
The first assertion is Proposition~\ref{pkl:std} with $l=0$. By Proposition~\ref{pkl:std} at the letter
$\sigma^l(i)$, whose $l$ steps back is $i$, $L_{\sigma^l(i)}(\pklstd_l)=\varphi^l(L_i(\pklstd_0))$, and $\varphi^l$
is injective.
\end{proof}

\begin{corollary}[Lemma~\ref{pkm:inputs}(H), first assertion]\label{pkl:H1}
With $a$ the jump bound of $L$: $\tau^a(1\otimes y)\in\pknMinf(L)\subseteq\pklAN(L)$ and
$e_i\tau^a(1\otimes y)\in L_i(\pklstd_0)$ for all $y\in L$ and letters $i$.
\end{corollary}

\begin{proof}
The first part is Proposition~\ref{pkn:frame}(v) with Theorem~\ref{pkf:main}. Then $e_i\tau^a(1\otimes y)\in
\pknMinf\cap\pkpAmb(L)_i$ by Lemma~\ref{pkl:basic}(iv), which is $L_i(\pklstd_0)$ by Corollary~\ref{pkl:F}.
\end{proof}

\emph{Remark.} With Definition~\ref{pkn:def-N}, $\mathcal N(L)$ is stable under $p_b^{-1}$
(Lemma~\ref{pkn:M}(iii)). When $\Gamma$ acts trivially on $L\ne0$, $\mathfrak S\otimes L\subseteq\mathcal N(L)$
by Proposition~\ref{pkl:H2}(ii), and then the inclusion is strict ($E^{-1}(1\otimes y)\in\mathcal N(L)$ for
$0\ne y\in L$, but $E^{-1}\notin\pkpA$); this is why (H) is stated with the intersection
$\mathcal N(L)\cap(\pklL)$.

\begin{proposition}[Lemma~\ref{pkm:inputs}(H), trivial action]\label{pkl:H2}
Suppose $\Gamma$ acts trivially on $L$. Then:
\begin{enumerate}
\item[(i)] $\pknMinf(L)=\pklL$ and $L_i(\pklstd_0)=\pklL\cap\pkpAmb(L)_i\subseteq\pklL$;
\item[(ii)] $\mathfrak S\otimes L\subseteq\mathcal N(L)$, and $e_i(1\otimes y)\in\mathcal N_i(\pklstd_l)$ for all
$y\in L$, letters $i$ and $l\ge0$.
\end{enumerate}
\end{proposition}

\begin{proof}
For a form $\mu$ and $y\in L$, $g\mu(1\otimes y)=\mu(1\otimes gy)=\mu(1\otimes y)$, so $\mu(1\otimes y)\in K'$.

(i) $\supseteq$: for $x=\sum_sb_s\otimes e_s$ and $m\ge1$, $\mu(\varphi^mx)=\sum_s\varphi^m(b_s)\mu(1\otimes e_s)\in
\pknBp$. $\subseteq$: let $x=\sum_sw_s\otimes e_s\in\pknMinf$ with $\tau^nx\in\pklL$, i.e.\ $\tau^nw_s\in\pkpA$. Each
$\lambda_k$ is a form ($\lambda_k(1\otimes gy)=\lambda_k(1\otimes y)\in\Z_\ell$). Hence, for $m\ge1$,
$\lambda_k(\varphi^mx)=\lambda_0(\varphi^mw_k)\in\pknBp$, and $\varphi^m(\tau)^n\lambda_0(\varphi^mw_k)=
\varphi^m(\tau^nw_k)$. We show by induction on $n$: \emph{if $w\in\pkpW$, $\tau^nw\in\pkpA$ and $q_m:=
\varphi^m(\tau^nw)/\varphi^m(\tau)^n\in\pknBp$ for all $m\ge1$, then $w\in\pkpA$.} For $n=0$ there is nothing
to prove. For $n\ge1$, $t=\tau^nw\in\pkpA$ has $\varphi^m(t)=\varphi^m(\tau)^nq_m\in\mathrm{Fil}^1$ for $m\ge1$
(Lemma~\ref{pkl:basic}(i)), i.e.\ $\theta(\varphi^j(\varphi t))=0$ for all $j\ge0$; by Lemma~\ref{pkn:mu},
$\varphi(\tau)\mid\varphi(t)$, so $t=\tau t'$ with $t'=\tau^{n-1}w\in\pkpA$, and the quotients for $(w,n-1)$ are the
same $q_m$. So $w_k\in\pkpA$ and $x\in\pklL$. The second assertion follows from Corollary~\ref{pkl:F}.

(ii) $s\otimes y\in\pklL=\pknMinf\subseteq\pklAN$ and $\mu(s\otimes y)=s\mu(1\otimes y)\in\pknK$ ($\pknK\supseteq
K'$ and the image of $\mathfrak S$). Next, $e_i(1\otimes y)\in\Lambda\otimes L\subseteq\mathfrak S\otimes L$ lies in
$\mathcal N\cap\pkpAmb(L)_i$. For $a\ge0$, $\varphi^{-a}e_i(1\otimes y)=e_{i^{(-a)}}(1\otimes y)$ ($\varphi^{-a}[i']=
[i'^{1/\ell^a}]$ and $i'\mapsto i'^{1/\ell^a}$ permutes the letters). For a form $\mu$ its value lies in
$\pkpA\cdot K'\subseteq\mathrm{Fil}^0$, and in $\mathrm{Fil}^{-u}$ if $\mu$ is of type $u$ at $i^{(-a)}$ (by
definition). As $\pklstd_l(a,u)\in\{0,-u\}$, $e_i(1\otimes y)\in L_i(\pklstd_l)$.
\end{proof}

\subsection{Galois stability and the level-one gain (D)}\label{pkl:ss-D}

\emph{Remark.} The $\Gamma'$-stability of $\mathcal N(L)$ itself is neither claimed nor needed: $g\in\Gamma'$ does
not fix $\mathfrak S$, and the field $\pknK$ need not be $g$-stable.

\begin{lemma}[Lemma~\ref{pkm:inputs}(G), stability]\label{pkl:G1}
For $g\in\Gamma'$: $g\pknMinf(L)\subseteq\pknMinf(L)$, $g\,\pklAN(L)\subseteq\pklAN(L)$ and $gL_i(o)\subseteq L_i(o)$
for all letters $i$ and orders $o$.
\end{lemma}

\begin{proof}
$g$ preserves $\pklL$ ($g(b\otimes y)=g(b)\otimes gy$) and $g\tau=\tau r'$ with $r'\in\pkpA^\times$
(Lemma~\ref{pkf:approx0}(iv)); so $\tau^nx\in\pklL$ implies $\tau^ngx=r'^{-n}g(\tau^nx)\in\pklL$. For a form $\mu$
the map $\mu^g(w)=g^{-1}\mu(gw)$ is $\pkpA$-linear and is a form: $\mu^g(1\otimes hy)=g^{-1}\mu(1\otimes
(ghg^{-1})gy)=h\mu^g(1\otimes y)$ for $h\in\Gamma'$. As $g$ commutes with $\varphi$, $\mu(\varphi^mgx)=
g\mu^g(\varphi^mx)\in g\pknBp=\pknBp$ for $x\in\pknMinf$, $m\ge1$. So $g\pknMinf\subseteq\pknMinf$. If
$rx\in\pknMinf$ with $r=\prod_kp_{b_k}$, then $g(r)gx\in\pknMinf$ and $g(p_b)=\varphi^b(gE)=p_b\varphi^b(r_E)$ with
$r_E$ a unit (Lemma~\ref{pkf:approx0}(iv)), so $rgx\in\pknMinf$. Finally let $x\in L_i(o)$: $gx\in\pklAN$, $gx\in
\pkpAmb(L)_i$ ($g$ commutes with $1\otimes\zeta$ and fixes $[i]$), and by Lemma~\ref{pkm:equiv}
$\mu(\varphi^{-a}gx)=g\mu(\varphi^{-a}x)\in g\mathrm{Fil}^n=\mathrm{Fil}^n$ ($g(t)=\chi(g)t$).
\end{proof}

The next lemma transports relations through $\mathfrak S$-linear maps of $\pkpA$; it rests on input (A8)(c),
proved as Proposition~\ref{pkr:I8}(c).

\begin{lemma}[Retraction of values]\label{pkl:valret}
Let $f\colon\pkpA\to\pkpA$ be $\mathfrak S$-linear, $v_1,\dots,v_n\in\pknK$, $r_1,\dots,r_n\in\pkpA$, and $W\subseteq
\pknB$ a $\pknBp$-submodule. If $\sum_lr_lv_l\in W$, then $\sum_lf(r_l)v_l\in W$.
\end{lemma}

\begin{proof}
Let $R\subseteq\pknB$ be the $K'$-span of the image of $\pkfSB$; it is a ring ($K'$ consists of constants and
$\pkfSB$ is a ring) containing the image of $\mathfrak S\subseteq\pkfSB$. The generators $bs$ ($b\in\pkfSB$,
$s\in\mathfrak S$) of $\pknK$ lie in the image of $\pkfSB$, and $s=1$ gives all of it; so $\pknK$ is the field of
fractions of the domain $R$ inside $\pknB$. Choose $d\in R\smallsetminus\{0\}$ with $dv_l=w_l\in R$, and write
$w_l=\sum_ic_{li}b_{li}$ (finite sums, $c_{li}\in K'$, $b_{li}\in\pkfSB$). By Proposition~\ref{pkr:I8}(c), applied to
the coefficients $c_{li}$ and the elements $b_{li}r_l\in\pkfBp$, there is $u\in\pknBp$ with $\sum_{l,i}c_{li}
G_f(b_{li}r_l)=u\sum_{l,i}c_{li}b_{li}r_l$, and $G_f(b_{li}r_l)=b_{li}f(r_l)$ by Lemma~\ref{pkf:extend}. Hence
$\sum_lw_lf(r_l)=u\sum_lw_lr_l$, and dividing by $d$, $\sum_lf(r_l)v_l=u\sum_lr_lv_l\in W$.
\end{proof}

\begin{proposition}[Lemma~\ref{pkm:inputs}(D)]\label{pkl:D}
If $x\in\mathcal N_i(\pklstd_1)$ and $g\in\Gamma'$, then $gx-x\in L_i(o_1)$, where $o_1(0,u)=1$ and $o_1(a,u)=-u$
for $a\ge1$.
\end{proposition}

\begin{proof}
By Lemma~\ref{pkl:G1}, $gx-x\in L_i(\pklstd_1)$, and $\pklstd_1(a,u)=o_1(a,u)$ for $a\ge1$; it remains to treat
$a=0$. Let $u\in U(i)$ and $\mu$ a form of type $u$ at $i$, and put $y=\mu(x)$. Then $y\in\pknK$ ($x\in\mathcal N$)
and $y\in\mathrm{Fil}^{\pklstd_1(0,u)}=\pknBp$. As $\theta\colon\pkpA\to O_C$ is onto, there are $e\ge0$ and
$r\in\pkpA$ with $\theta(r)=\ell^e\theta(y)$, i.e.\ $\ell^ey-r\in\mathrm{Fil}^1$. Lemma~\ref{pkl:valret} with
$f=\kappa_0$ (the coordinate of (A5), Proposition~\ref{pkr:I5}, $\mathfrak S$-linear with values in $\mathfrak S$
and $\kappa_0(\ell^e)=\ell^e$), the values $(y,-1)$ and the coefficients $(\ell^e,r)$, gives $\ell^ey-s\in
\mathrm{Fil}^1$ with $s=\kappa_0(r)\in\mathfrak S$. Now $g$ preserves $\mathrm{Fil}^1$, and $gs-s\in z\tau E\pkpA
\subseteq\mathrm{Fil}^1$ (Proposition~\ref{pkn:Ca}). Hence $\ell^e(gy-y)=g(\ell^ey-s)-(\ell^ey-s)+(gs-s)\in
\mathrm{Fil}^1$, and $gy-y\in\mathrm{Fil}^1$. By Lemma~\ref{pkm:equiv}, $\mu(gx-x)=gy-y$.
\end{proof}

\subsection{The depth step (C)(d)}\label{pkl:ss-Cd}

\begin{lemma}\label{pkl:different}
Let $P\in\Lambda[X]$ be the distinguished polynomial with $E=P(z)V(z)$, $V\in\Lambda[[X]]^\times$ (Weierstrass
preparation in $\Lambda[[X]]\cong\mathfrak S$). Then $P$ is Eisenstein of degree $e$, $P(\pi)=0$, and
$v_\ell(P'(\pi))\le d_v$.
\end{lemma}

\begin{proof}
Modulo $\ell$, $E\equiv z^eA$ with $A$ a unit, so $P\equiv X^e$ and $\deg P=e$; $P(0)=E(0)/V(0)=\ell/V(0)$ has
valuation $1$. So $P$ is Eisenstein, and $P(\pi)=\theta(P(z))=\theta(E)\theta(V(z))^{-1}=0$. If $e=1$, $P'=1$
and $d_v=0$. If $e>1$, write $P=\sum_{k=0}^ec_kX^k$ with $c_e=1$ and $v_\ell(c_k)\ge1$ for $k<e$. The terms
$kc_k\pi^{k-1}$ ($1\le k\le e$) of $P'(\pi)$ have valuations $v_\ell(k)+v_\ell(c_k)+(k-1)/e$, pairwise distinct
(the fractional parts $(k-1)/e$ are distinct), so $v_\ell(P'(\pi))$ is their minimum, at most
$v_\ell(e)+(e-1)/e<v_\ell(e)+1=d_v$.
\end{proof}

\begin{proposition}[Lemma~\ref{pkm:inputs}(C)(d)]\label{pkl:Cd}
For every $n\ge1$ there are $g\in\Gamma'$, $\omega\in\pkpA^\times$ and $\upsilon,q\in\pkpA$ with $gE=\omega E$,
$E\mid\upsilon q-\ell^{v_\ell(n)+d_v}$ and $z\tau q=\omega^n-1$.
\end{proposition}

\begin{proof}
Let $g=t_1$ with $gz=[\varepsilon]z$ ((C)(b)), and $gE=\omega E$ with $\omega\in\pkpA^\times$
(Lemma~\ref{pkf:approx0}(iv)). Write $\delta=z\tau$; $\theta(\delta)=\pi(\zeta_\ell-1)\ne0$.

\emph{Step 1: $\delta\mid\omega-1$, hence $\delta\mid\omega^n-1=\delta q$ for a unique $q\in\pkpA$.} By
Proposition~\ref{pkn:Ca}, $\delta E\mid gE-E=E(\omega-1)$; cancel $E$.

\emph{Step 2: $\theta(\omega-1)\mid\ell^{d_v}\theta(\delta)$ in $O_C$.} With $P,V$ as in Lemma~\ref{pkl:different},
$P(z)=EG$ with $G=V(z)^{-1}\in\mathfrak S$. By Lemma~\ref{pkn:mu}, $gz-z=z([\varepsilon]-1)=z\tau Eu=Eh$ with
$h=\delta u$, $u\in\pkpA^\times$. Taylor's formula for the polynomial $P$ gives $c\in\pkpA$ with
$P(gz)=P(z)+P'(z)Eh+E^2h^2c$. As $g$ fixes $\Lambda$, $P(gz)=g(P(z))=g(E)g(G)=\omega Eg(G)$; cancelling $E$,
$\omega g(G)=G+P'(z)h+Eh^2c$. Apply $\theta$: $\theta(g(G))=g\theta(G)=\theta(G)$, because $\theta(G)$ is a
convergent series in $\pi$ with coefficients in $\theta(\Lambda)$, fixed by $j(\Gamma')$. So
$\theta(\omega-1)\theta(G)=P'(\pi)\theta(\delta)\theta(u)$, and $\theta(G)=\theta(V(z))^{-1}$ is a unit of $O_C$
($V(z)\in\mathfrak S^\times$). By Lemma~\ref{pkl:different}, $\theta(\omega-1)\mid\ell^{d_v}\theta(\delta)$; in particular
$\theta(\omega-1)\ne0$.

\emph{Step 3: $v_\ell(\theta(\omega^n-1))=v_\ell(n)+v_\ell(\theta(\omega-1))$.} For $r=\delta b\in\delta\pkpA$,
$\ell\,\theta(r)\theta(z)\mid\theta(r)^\ell$: indeed $y=\theta(\tau)=\zeta_\ell-1$ satisfies $(1+y)^\ell=1$, so
$y^\ell=-\sum_{0<k<\ell}\binom\ell ky^k$ is divisible by $\ell y$; and $\theta(z)^2\mid\theta(z)^\ell$,
$\theta(b)\mid\theta(b)^\ell$. Let $x_0=\theta(\omega)$ and $y_0=x_0-1=\theta(\omega-1)\ne0$, of positive valuation
($\theta(z)\mid y_0$). Write $n=m\ell^v$ with $\ell\nmid m$. First $x_0^m-1=y_0(1+x_0+\dots+x_0^{m-1})$ and the
second factor is $\equiv m\not\equiv0$ modulo the maximal ideal, so $v_\ell(x_0^m-1)=v_\ell(y_0)$. Next, if
$x'=\theta(\omega')$ with $\omega'$ a power of $\omega$ and $y'=x'-1\ne0$, then $\delta\mid\omega'-1$ (Step 1), and
$x'^\ell-1=\ell y'+\sum_{1<k<\ell}\binom\ell ky'^k+y'^\ell$, where the middle terms have valuation $\ge1+2v_\ell(y')$
and the last $\ge1+v_\ell(y')+v_\ell(\pi)$ by the divisibility just shown; so $v_\ell(x'^\ell-1)=1+v_\ell(y')$.
Applying this $v$ times starting from $\omega'=\omega^m$ gives the claim.

\emph{Step 4.} By Steps 2 and 3, $\theta(\delta)\theta(q)=\theta(\omega^n-1)$ divides $\ell^{v_\ell(n)}\theta(\omega-1)$,
which divides $\ell^{v_\ell(n)+d_v}\theta(\delta)$. Cancelling $\theta(\delta)\ne0$ in the domain $O_C$,
$\ell^{v_\ell(n)+d_v}=\theta(q)c$ with $c\in O_C$. Choose $\upsilon\in\pkpA$ with $\theta(\upsilon)=c$; then
$\theta(\upsilon q-\ell^{v_\ell(n)+d_v})=0$ and $E\mid\upsilon q-\ell^{v_\ell(n)+d_v}$ (Corollary~\ref{pkn:E-gen}).
\end{proof}

\subsection{The retraction (R)}\label{pkl:ss-R}

\begin{lemma}[Relations]\label{pkl:rel}
Let $v_1,\dots,v_n\in\pklAN(L)$ with $\mu(v_i)\in\pknK$ for every form $\mu$ (e.g.\ $v_i\in\mathcal N(L)$), let
$c_1,\dots,c_n\in\pkpA$ with $\sum_ic_iv_i=0$, and let $f\colon\pkpA\to\pkpA$ be $\mathfrak S$-linear. Then
$\sum_if(c_i)v_i=0$.
\end{lemma}

\begin{proof}
Fix a frame (Theorem~\ref{pkf:main}). Put $w=\sum_if(c_i)v_i\in\pklAN$ and choose $r$, a product of a power of
$\tau$ and of elements $p_b$, with $rw\in\pklL$ (Lemma~\ref{pkn:M}(ii)). As in the proof of
Proposition~\ref{pkn:frame}(ii), $\tau^arw=\sum_jt_jn_j$ with $t=Q'\alpha\in\pkpA^d$, where $rw=\sum_j\alpha_j
(1\otimes e_j)$. For every form $\mu$, $\sum_ic_i\mu(v_i)=\mu(0)=0$, so Lemma~\ref{pkl:valret} (with $W=0$) gives
$\mu(w)=\sum_if(c_i)\mu(v_i)=0$, hence $\mu(\sum_jt_jn_j)=\tau^ar\mu(w)=0$. By the coefficient test
(Lemma~\ref{pkn:frame-id}(iv) with $W=0$, $m=0$), $E^{2a}t_j=0$, so $t_j=0$, $\tau^arw=0$, and $w=0$ since
$\tau^ar$ is a unit of $\pkpW$.
\end{proof}

\begin{lemma}\label{pkl:phin}
For a frame of $L$ and $l\ge0$, $\varphi^l(n_j)\in\mathcal N(L)$ for every $j$.
\end{lemma}

\begin{proof}
$n_j\in\pknMinf$ (Proposition~\ref{pkn:frame}(i)), so $\varphi^l(n_j)\in\pklAN$ by Lemma~\ref{pkl:basic}(ii). For a form
$\mu$, Lemma~\ref{pkn:frame-id}(ii) with $m=l$ and $t$ the $j$-th unit vector gives $p_l^a\mu(\varphi^ln_j)=
c^{(l)}(\mu)\,\varphi^l(Z)_{\cdot j}$, with $c^{(l)}(\mu)$ over $K'$ (Lemma~\ref{pkn:frame-id}(i)) and the entries of
$\varphi^l(Z)$ in $\pkfSB$ (Lemma~\ref{pkf:values}); their images lie in $\pknK$, and $p_l$ is a nonzero element of
$\pknK$. So $\mu(\varphi^ln_j)\in\pknK$.
\end{proof}

\begin{proposition}[Lemma~\ref{pkm:inputs}(R)]\label{pkl:R}
Let $\kappa_0\colon\pkpA\to\mathfrak S$ be the coordinate of (A5) (Proposition~\ref{pkr:I5}); it plays the role of the
map $\lambda_0$ of Lemma~\ref{pkm:inputs}(R) (not the map $\lambda_0$ of \S\ref{pkl:ss-pre}), and it is
$\mathfrak S$-linear and the identity on $\mathfrak S$. For every lattice with jump sets $L$ there is a unique map
$\mathrm{ret}\colon\pklAN(L)\to\mathcal N(L)$ with $\mathrm{ret}(\sum_ib_ix_i)=\sum_i\kappa_0(b_i)x_i$ for $b_i\in\pkpA$,
$x_i\in\mathcal N(L)$. It is $\mathfrak S$-linear, is the identity on $\mathcal N(L)$, commutes with $1\otimes\zeta$ and
the $e_i$, maps $L_i(\pklstd_l)$ into $\mathcal N_i(\pklstd_l)$ for all $i$ and $l\ge0$, and, for an extension $X_S$,
$\mathrm{pr}\circ\mathrm{ret}=\mathrm{ret}\circ\mathrm{pr}$.
\end{proposition}

\begin{proof}
$\pklAN(L)$ is the $\pkpA$-span of $\mathcal N(L)$ (Proposition~\ref{pkn:frame}(iv) with Theorem~\ref{pkf:main}), so
every element has the form $\sum_ib_ix_i$. If $\sum_ib_ix_i=\sum_jb'_jx'_j$, Lemma~\ref{pkl:rel} applied to the
relation $\sum_ib_ix_i-\sum_jb'_jx'_j=0$ with $f=\kappa_0$ shows $\sum_i\kappa_0(b_i)x_i=\sum_j\kappa_0(b'_j)x'_j$; so
$\mathrm{ret}$ is well defined, with values in the $\mathfrak S$-module $\mathcal N$. It is $\mathfrak S$-linear because
$\kappa_0$ is, and $\mathrm{ret}(x)=\kappa_0(1)x=x$ for $x\in\mathcal N$. Since $1\otimes\zeta$ preserves $\mathcal N$
(Lemma~\ref{pkl:basic}(iv)), $(1\otimes\zeta)\mathrm{ret}(\sum b_ix_i)=\sum\kappa_0(b_i)(1\otimes\zeta)x_i=
\mathrm{ret}((1\otimes\zeta)\sum b_ix_i)$; as $e_i$ is a polynomial in $1\otimes\zeta$ with coefficients in
$\Lambda\subseteq\mathfrak S$, $\mathrm{ret}$ commutes with $e_i$.

Let $x\in L_i(\pklstd_l)$. By Proposition~\ref{pkl:std} and Proposition~\ref{pkn:frame}(ii), $x=\varphi^l(\sum_jt_jn_j)
=\sum_j\varphi^l(t_j)\varphi^l(n_j)$ with $t_j\in\pkpA$, and $\varphi^l(n_j)\in\mathcal N$ (Lemma~\ref{pkl:phin}); so
$\mathrm{ret}(x)=\sum_j\kappa_0(\varphi^lt_j)\varphi^l(n_j)=\varphi^l\bigl(\sum_j\varphi^{-l}(\kappa_0(\varphi^lt_j))n_j
\bigr)\in\varphi^l(\pknMinf)$. Also $\mathrm{ret}(x)=\mathrm{ret}(e_ix)=e_i\mathrm{ret}(x)\in\pkpAmb(L)_i$. By
Proposition~\ref{pkl:std}, $\mathrm{ret}(x)\in L_i(\pklstd_l)$, and it lies in $\mathcal N$.

For an extension $X_S$, $\mathrm{pr}$ is the $\pkpW$-linear extension of the equivariant $\auO$-linear projection
$X_S\to C_S$; it commutes with $\varphi$, maps $\pkpA\otimes X_S$ onto $\pkpA\otimes C_S$, and $\nu\circ\mathrm{pr}$ is a
form on $X_S$ for every form $\nu$ on $C_S$. Hence $\mathrm{pr}$ maps $\pknMinf(X_S)$, $\pklAN(X_S)$ and $\mathcal N(X_S)$
into the corresponding modules of $C_S$, and $\mathrm{pr}(\mathrm{ret}(\sum b_ix_i))=\sum\kappa_0(b_i)\mathrm{pr}(x_i)=
\mathrm{ret}(\sum b_i\mathrm{pr}(x_i))=\mathrm{ret}(\mathrm{pr}(\sum b_ix_i))$.
\end{proof}


\begin{corollary}[Frobenius on $\mathcal N$]\label{pkl:Phi}
$\varphi(\mathcal N(L))\subseteq\mathcal N(L)$.
\end{corollary}

\begin{proof}
Let $x\in\mathcal N$ and $r=\prod_kp_{b_k}$ with $rx\in\pknMinf$; $rx\in\mathcal N$ (Lemma~\ref{pkn:M}(iii)). By
Proposition~\ref{pkn:frame}(ii), $rx=\sum_jt_jn_j$ with $t_j\in\pkpA$; as $rx$ and the $n_j$ lie in $\mathcal N$
(Proposition~\ref{pkn:frame}(i)), $rx=\mathrm{ret}(rx)=\sum_j\kappa_0(t_j)n_j$, and by linear independence
$t_j=\kappa_0(t_j)\in\mathfrak S$. Then $\varphi(r)\varphi(x)=\sum_j\varphi(t_j)\varphi(n_j)$ with $\varphi(t_j)\in\mathfrak S$
and $\varphi(n_j)\in\mathcal N$ (Lemma~\ref{pkl:phin}), so $\varphi(r)\varphi(x)\in\mathcal N$; $\varphi(r)=\prod_kp_{b_k+1}$,
and $\mathcal N$ is stable under $p_b^{-1}$, so $\varphi(x)\in\mathcal N$.
\end{proof}

\begin{corollary}[(H), second assertion]\label{pkl:H3}
If $\Gamma$ acts trivially on $L$, then $\mathcal N(L)\cap\pklL=\mathfrak S\otimes L$, and $\mathcal N(L)$ is the set of
$x$ with $rx\in\mathfrak S\otimes L$ for some finite product $r$ of elements $p_b$. Moreover
$L_i(\pklstd_0)\subseteq\pklL$ (Proposition~\ref{pkl:H2}(i)).
\end{corollary}

\begin{proof}
$\supseteq$ in the first claim is Proposition~\ref{pkl:H2}(ii). If $x=\sum_sb_s(1\otimes e_s)\in\mathcal N$, then, as
$1\otimes e_s\in\mathcal N$, $x=\mathrm{ret}(x)=\sum_s\kappa_0(b_s)(1\otimes e_s)\in\mathfrak S\otimes L$. For the second
claim: if $x\in\mathcal N$, then $rx\in\pklL$ for some product $r$ of a power of $\tau$ and elements $p_b$
(Lemma~\ref{pkn:M}(ii)); by Proposition~\ref{pkl:H2}(i) and the definition of $\pklAN$ one may take $r$ a product of
$p_b$'s only ($rx\in\pknMinf=\pklL$ by definition of $\pklAN$), and $rx\in\mathcal N$ (Lemma~\ref{pkn:M}(iii)), so
$rx\in\mathfrak S\otimes L$. Conversely $\mathfrak S\otimes L\subseteq\mathcal N$ and $\mathcal N$ is stable under
$p_b^{-1}$.
\end{proof}

\subsection{The second assertion of (G)}\label{pkl:ss-open}

The second assertion of (G) (Lemma~\ref{pkp:tau}(iv)) reduces, by Lemma~\ref{pkl:basic}(i) and the argument
below, to a divisibility statement. Let $x\in\mathcal N(L)$ and $g\in\Gamma'$ with $gx-x\in L_i(o)$. If
$y=(z\tau)^{-1}(gx-x)$ lies in $\pklAN(L)$, then $y\in L_i(o)$: it lies at $i$, and for a form $\mu$ of type $u$
at $i^{(-a)}$, $\varphi^{-a}(z\tau)\mu(\varphi^{-a}y)=\mu(\varphi^{-a}(gx-x))\in\mathrm{Fil}^{o(a,u)}$ with
$\varphi^{-a}(z\tau)$ a unit of $\pknBp$ (Lemma~\ref{pkl:basic}(i)). Moreover, if $y\in\pklAN$ for every
$x\in\mathcal N\cap\pknMinf$, then $y\in\pklAN$ for every $x\in\mathcal N$: if $x\in\mathcal N$ and $p_bx$ has the
property, then $g(p_bx)-p_bx=p_b(gx-x)+(gp_b-p_b)gx$ with $gp_b-p_b\in z\tau E\pkpA$ (Proposition~\ref{pkn:Ca}), so
$p_b\,y=(z\tau)^{-1}(g(p_bx)-p_bx)-E\beta\,gx\in\pklAN$ for some $\beta\in\pkpA$, and $y\in\pklAN$
(Lemma~\ref{pkn:M}(iii), Lemma~\ref{pkl:G1}); one inducts on the number of factors $p_b$ needed to bring $x$ into
$\pknMinf$ (note $p_bx\in\mathcal N$). For $x\in\mathcal N\cap\pknMinf$ one has $x=\sum_jt_jn_j$ with $t_j\in\mathfrak S$
(apply $\mathrm{ret}$ and Proposition~\ref{pkn:frame}(i),(ii)), so that $E^a\sigma_\nu(x)=(Zt)_\nu\in\pkfSB$ for the
forms $\sigma_\nu$ of Lemma~\ref{pkn:frame-id}(iii) ($m=0$). Also, if $z\tau$ divides $E^hr$ in $\pkfBp$
($r\in\pkpA$), then $\tau\mid r$ in $\pkpA$: from $E^hr=\tau w$, $\theta(p_n)^h\theta(\varphi^nr)=
\theta(\varphi^n\tau)\theta(\varphi^nw)=0$ for $n\ge1$, so $\varphi(\tau)\mid\varphi(r)$ by Lemma~\ref{pkn:mu}.

We now prove the two facts about $\pkfBp$ needed to finish, and then (G0).

\begin{lemma}[Integral divisibility]\label{pkl:intdiv}
Let $\nu,\beta,x\in\pkpA$ with $\ell\nmid\nu$ and $\nu\mid E-\ell\beta$ in $\pkpA$. If $\nu$ divides $x$ in $\pkfBp$, then
$\nu\mid x$ in $\pkpA$. In particular this holds for $\nu=z$ ($\beta=1$, as $E-\ell=z^eA$).
\end{lemma}

\begin{proof}
Write $x=\nu w$ with $w\in\pkfBp$ and choose $k$ with $\ell^kw\in\pkfAo$; so $\ell^kx=\nu w_0$ in $\pkfAm$. Fix $n\ge1$
and let $Q=\pkpA/(\nu,\ell^n)$. The ring map $\pkpA[X]\to Q$, $X\mapsto\bar\beta$, kills $\ell X-E$ (as
$E-\ell\beta\in\nu\pkpA$), so it factors through $R=\pkpA[X]/(\ell X-E)$; since $\ell^n=0$ in $Q$, it factors through
$R/\ell^nR$, hence through $\pkfAm\to R/\ell^nR$ (Lemma~\ref{pkr:Amax}(i)). The resulting map $F\colon\pkfAm\to Q$
restricts to the projection on $\pkpA$, and $F(\ell^kx)=F(\nu)F(w_0)=0$. So $\ell^kx\in(\nu,\ell^n)\subseteq\nu\pkpA+
\pkrI^n$ for every $n$, and Proposition~\ref{pkr:I7} (with $\delta=\nu$, $c=\ell^kx$) gives $\nu\mid\ell^kx$. As $\ell$ is a
prime element of $\pkpA$ not dividing $\nu$ ((W1) of \S\ref{pkr:ss-prelim}), $\nu\mid x$. For $\nu=z$: $z_0=\varpi^\flat\ne0$,
so $\ell\nmid z$.
\end{proof}

\begin{lemma}\label{pkl:SBdec}
For $b\in\pkfSB$ and $M\ge0$ there are $k\ge0$, $s\in\mathfrak S$ and $w\in\pkfSB$ with $\ell^kb=s+E^Mw$, and there are
$k'\ge0$, $\lambda\in\Lambda$ and $w'\in\pkfSB$ with $\ell^{k'}b=\lambda+zw'$.
\end{lemma}

\begin{proof}
Let $\ell^jb=y\in\pkfAo$. Lemma~\ref{pkr:div} with $\nu=E$ ($b=0$), resp.\ $\nu=z$ ($b=1$; $z\mid E-\ell$), gives
$\ell^Ny=a+\nu^Nw_1$ with $a\in\pkpA$, $w_1\in\pkfAo$. Apply $G_0$ (Lemmas~\ref{pkf:extend}, \ref{pkf:values}): $G_0$ fixes
$\ell^Ny\in\pkfSB$, is $\pkfSB$-linear, and $G_0(w_1)\in\pkfSB$; so $\ell^{N+j}b=\kappa_0(a)+\nu^NG_0(w_1)$. With $\nu=E$,
$N=M$ this is the first claim. With $\nu=z$, $N=1$, write $\kappa_0(a)=\lambda+zs'$ ($\lambda\in\Lambda$, $s'\in\mathfrak S$)
and $w'=s'+G_0(w_1)$.
\end{proof}

\begin{lemma}[Galois action on $\pkfSB$]\label{pkl:galSB}
For $g\in\Gamma'$ and $b\in\pkfSB$, $z\tau$ divides $gb-b$ in $\pkfBp$.
\end{lemma}

\begin{proof}
\emph{$z$ divides.} By Lemma~\ref{pkl:SBdec}, $\ell^{k'}b=\lambda+zw'$. $g$ fixes $\lambda$, and $gz=z[\eta]$ with
$\eta\in O^{\flat\times}$ (proof of Proposition~\ref{pkn:Ca}); so $\ell^{k'}(gb-b)=z([\eta]g(w')-w')$, and $\ell$ is a unit of
$\pkfBp$. Write $gb-b=zc$, $c\in\pkfBp$.

\emph{$\tau$ divides $c$.} First, $\theta(gb'-b')=0$ for every $b'\in\pkfSB$: by Lemma~\ref{pkl:SBdec} with $M=1$,
$\ell^kb'=s+Ew$, and $\theta(g(\ell^kb'))=g\theta(s)+g\theta(E)\theta(gw)=\theta(s)=\theta(\ell^kb')$, since $\theta$ commutes with
$g$ on $\pkfBp$ (P1), $\theta(E)=0$, and $\theta(s)$, a convergent series in $\pi$ over $\theta(\Lambda)$, is fixed by
$j(\Gamma')$. For $n\ge0$, $\varphi^n(b)\in\pkfSB$ (Lemma~\ref{pkf:values}) and $g$ commutes with $\varphi$, so
$0=\theta(\varphi^n(gb-b))=\theta(\varphi^nz)\theta(\varphi^nc)$ with $\theta(\varphi^nz)=\pi^{\ell^n}\ne0$. Hence
$\theta(\varphi^n(\varphi c))=0$ for all $n\ge0$, $\mu=\varphi(\tau)$ divides $\varphi(c)$ (Lemma~\ref{pkr:mu-Bplus}), and
$\tau\mid c$ in $\pkfBp$ (Proposition~\ref{pkr:I8}(b)).
\end{proof}

\begin{lemma}[Cancellation]\label{pkl:cancel}
If $r\in\pkpA$, $h\ge0$ and $z\tau$ divides $E^hr$ in $\pkfBp$, then $z\tau\mid r$ in $\pkpA$.
\end{lemma}

\begin{proof}
By the argument above (before this lemma), $r=\tau r'$ with $r'\in\pkpA$. Then $E^h\tau r'=z\tau w$ in the domain $\pkfBp$
gives $E^hr'=zw$, i.e.\ $z$ divides $E^hr'$ in $\pkfBp$, and $z\mid E^hr'$ in $\pkpA$ by Lemma~\ref{pkl:intdiv}. As
$\theta(z)\ne0$ and $\ker\theta=E\pkpA$ is prime ($\pkpA/E\pkpA\cong O_C$), $E\nmid z$; $z\mid E^hr'$ then gives
$z\mid r'$: by induction on $h$, if $E^hr'=zq$ then $\theta(z)\theta(q)=0$, $q=Eq'$, and $E^{h-1}r'=zq'$.
\end{proof}

\begin{proposition}[(G0); Lemma~\ref{pkm:inputs}(G), second assertion]\label{pkl:G0}
For $x\in\mathcal N(L)$ and $g\in\Gamma'$, $(z\tau)^{-1}(gx-x)\in\pklAN(L)$. Consequently, if $x\in\mathcal N(L)$ and
$gx-x\in L_i(o)$, then $(z\tau)^{-1}(gx-x)\in L_i(o)$.
\end{proposition}

\begin{proof}
The consequence and the reduction to $x\in\mathcal N\cap\pknMinf$ were shown at the beginning of this subsection. Let
$x\in\mathcal N\cap\pknMinf$; fix a frame. As in the proof of Corollary~\ref{pkl:Phi}, $x=\sum_jt_jn_j$ with
$t_j\in\mathfrak S$, so $X_\nu:=E^a\sigma_\nu(x)=(Zt)_\nu\in\pkfSB$ (Lemma~\ref{pkn:frame-id}(ii), $m=0$, with the forms
$\sigma_\nu$ of Lemma~\ref{pkn:frame-id}(iii)). Write $gE=\omega E$, $\omega\in\pkpA^\times$; $z\tau\mid\omega-1$ (proof of
Proposition~\ref{pkl:Cd}, Step 1, valid for any $g\in\Gamma'$), hence $z\tau\mid\omega^a-1$. By Lemma~\ref{pkm:equiv},
$\sigma_\nu(gx-x)=g\sigma_\nu(x)-\sigma_\nu(x)$, and from $E^a\sigma_\nu(x)=X_\nu$, $\omega^aE^ag\sigma_\nu(x)=gX_\nu$. So
\[
\omega^aE^{2a}\bigl(g\sigma_\nu(x)-\sigma_\nu(x)\bigr)=E^a(gX_\nu-X_\nu)-(\omega^a-1)E^aX_\nu\in W:=z\tau\pkfBp
\]
(image in $\pknB$, a $\pkfBp$-submodule), by Lemma~\ref{pkl:galSB}; as $\omega$ is a unit, $\sigma_\nu(E^{2a}(gx-x))\in W$ for
all $\nu$. By Lemma~\ref{pkl:G1} and Proposition~\ref{pkn:frame}(ii), $gx-x=\sum_js_jn_j$ with $s_j\in\pkpA$. The proof of
the coefficient test Lemma~\ref{pkn:frame-id}(iv) uses its hypothesis only for the forms $\sigma_\nu$; applied to
$E^{2a}(gx-x)=\sum_jE^{2a}s_jn_j$ it gives $E^{4a}s_j\in W$, i.e.\ (as $\pkfBp\to\pknB$ is injective) $z\tau$ divides
$E^{4a}s_j$ in $\pkfBp$. By Lemma~\ref{pkl:cancel}, $s_j=z\tau s'_j$ with $s'_j\in\pkpA$, and $(z\tau)^{-1}(gx-x)=\sum_js'_jn_j\in
\pknMinf\subseteq\pklAN$.
\end{proof}


%% ===== proofs/C-kisin-tau.tex =====
% proofs/C-kisin-tau.tex
% Orders add; Lemma pkp:tau (i),(iv),(v); the case inputs (I1) and (I3).
% Labels start with pkt:. Conventions and macros as in proofs/C-kisin-pairing.tex and C-kisin-inputs.tex.
\providecommand{\auO}{\mathcal O}
\providecommand{\auQlb}{\overline{\Q}_\ell}
\providecommand{\pkpA}{A_{\mathrm{inf}}}
\providecommand{\pkpW}{W(C^\flat)}
\providecommand{\pkpm}{W(\mathfrak m^\flat)}
\providecommand{\pkpM}{M_m}
\providecommand{\pkpsm}{\mathrm{Sm}_m}
\providecommand{\pkpAmb}{\mathbb A}
\providecommand{\pknBp}{B^+_{\mathrm{dR}}}
\providecommand{\pknB}{B_{\mathrm{dR}}}
\providecommand{\pknMinf}{\mathfrak M_{\mathrm{inf}}}
\providecommand{\pknK}{\mathcal K_0}
\providecommand{\pkcPt}{\mathrm{Pt}}
\providecommand{\pkcsh}{\mathrm{sh}}
% macros of this file (prefix \pkt)
\providecommand{\pktV}{V_{\mathrm{dR}}}
\providecommand{\pktFr}{\mathrm{Fr}}

\section{Orders add, Frobenius and the points from killers}\label{pkt:sec}

This section proves Lemma~\ref{pkp:orders-add} (Lemma~\ref{pkt:add}), the assertions (i), (iv), (v)
of Lemma~\ref{pkp:tau} (in the form of \S\ref{pkn:ss-def}), and the inputs (I1) and (I3) of
\S\ref{pkc:ss-inputs}. We use the conventions and labels of \S\S\ref{pkp:sec}--\ref{pkn:sec}: the
modules $\pknMinf(L)\subseteq\pkpA\mathcal N(L)$ and $\mathcal N(L)$ are those of
Definition~\ref{pkn:def-N}, forms are as defined there, and $L_i(o)$ is the module of elements with
orders $o$ at $i$ (Definition~\ref{pkn:def-N}, last sentence). We write $\mathrm{Fil}^n=t^n\pknBp$,
$K'=\auQlb^{\,j(\Gamma')}$, and for a letter $j$ we write $j=i^{(-a)}$ when $j^{\ell^a}=i$.

Besides the published facts (P2), (P3) of \S\ref{pkf:ss-rings} (Fontaine) and Tate's theorem on
$C(n)^{G_{K'}}$ (already used in Lemma~\ref{pkm:BdR}), we use Theorem~\ref{pkf:main},
Proposition~\ref{pkn:frame}, and items of Lemma~\ref{pkm:inputs}, cited by letter; these are proved in
\S\S\ref{pkn:sec}, \ref{pkl:sec}, and the precise references are collected in \S\ref{pkt:ss-open}.

Two statements about the modules $\mathcal N$ are used besides the items of
Lemma~\ref{pkm:inputs} (they are listed there too, and proved in \S\ref{pkl:sec}):
\begin{enumerate}
\item[(G$_0$)] For $x\in\mathcal N(L)$ and $g\in\Gamma'$: $(z\tau)^{-1}(gx-x)\in\pkpA\mathcal N(L)$.
\item[($\Phi$)] $\varphi(\mathcal N(L))\subseteq\mathcal N(L)$.
\end{enumerate}
(G$_0$) is the membership half of the second assertion of Lemma~\ref{pkm:inputs}(G): by
Lemma~\ref{pkt:tau-iv} below, (G$_0$) implies that assertion. The inclusion $\varphi(\pkpA\mathcal
N(L))\subseteq\pkpA\mathcal N(L)$ is proved below (Lemma~\ref{pkt:frob}); ($\Phi$) adds the condition on the
values in $\pknK$.

\subsection{The Hodge expansion}\label{pkt:ss-hodge}

Let $L$ be a lattice with jump sets. Put $\pktV=\pknB\otimes_{\Z_\ell}L$ and $V^+=\pknBp\otimes_{\Z_\ell}L$;
$\Gamma$ acts on both diagonally (through $j$ on $\pknB$), semilinearly over $\pknB$. The ring
$\Lambda\otimes_{\Z_\ell}\auO$ acts on $\pktV$ ($\Lambda\subset\pkpA\to\pknBp$ on the first factor, $\auO$ on
$L$), so the projectors $e_j$ of Definition~\ref{pkp:def-letters} act on $\pktV$ and $V^+$; put
$V_j=e_j\pktV$, $V_j^+=e_jV^+$. Since $\Gamma'$ fixes $\Lambda$ and commutes with $\auO$, each $e_j$ commutes
with $\Gamma'$, and $\pktV=\bigoplus_jV_j$, $V^+=\bigoplus_jV_j^+$.

\emph{Convention on jumps.} For an embedding $\gamma\colon\auO\to C$, an integer $\kappa$ is a jump of $L$
at $\gamma$ if $C(-\kappa)$ occurs in the Hodge--Tate decomposition of $C\otimes_{\auO,\gamma}L$ (with
multiplicity the number of times it occurs). This is the convention of Lemma~\ref{pkp:tau} and of
(P3) in \S\ref{pkf:ss-rings}: $\auO(1)$ has the single jump $-1$. By Definition~\ref{pkp:def-orders} the
jumps at the embeddings of letter $j$ lie in $U(j)$.

\emph{Elements with a denominator.} Let $\pkpAmb(L)^{\mathrm{fr}}$ be the set of $x\in\pkpAmb(L)$ with
$rx\in\pkpA\otimes L$ for some nonzero $r\in\pkpA$ ($\pkpA\otimes L$ the $\pkpA$-span of the $1\otimes y$,
a free $\pkpA$-module on $1\otimes e_1,\dots,1\otimes e_d$ for a $\Z_\ell$-basis $(e_k)$ of $L$). The map
$\pkpA\to\pknBp$ is injective (\S\ref{pkn:ss-frame}), so a nonzero $r$ is invertible in $\pknB$, and
$\hat x=r^{-1}\iota(rx)\in\pktV$, with $\iota\colon\pkpA\otimes L\to\pktV$ the natural map, does not
depend on $r$. The set $\pkpAmb(L)^{\mathrm{fr}}$ is an $\pkpA$-module, stable under $\varphi^{\pm1}$ and
$\Gamma$, and $x\mapsto\hat x$ is $\pkpA$-linear, $\Gamma$-equivariant and commutes with the $e_j$. By
Lemma~\ref{pkn:M}(ii), $\pkpA\mathcal N(L)\subseteq\pkpAmb(L)^{\mathrm{fr}}$; if $x$ lies in the letter-$i$
part, $\widehat{\varphi^{-a}x}\in V_{i^{(-a)}}$ (Definition~\ref{pkp:def-letters}).

For an $\pkpA$-linear $\mu\colon\pkpAmb(L)\to\pknB$ let $\hat\mu\colon\pktV\to\pknB$ be the $\pknB$-linear
map with $\hat\mu(1\otimes y)=\mu(1\otimes y)$. Then $\mu(x)=\hat\mu(\hat x)$ for
$x\in\pkpAmb(L)^{\mathrm{fr}}$: indeed $r\mu(x)=\mu(rx)=\hat\mu(\iota(rx))=r\hat\mu(\hat x)$, and $r\ne0$ in
the field $\pknB$.

\begin{lemma}[Forms and equivariant functionals]\label{pkt:forms}
(i) $\mu$ is a form iff $\hat\mu(gv)=g\hat\mu(v)$ for all $g\in\Gamma'$, $v\in\pktV$; it has type $u$ at
$j$ iff $\hat\mu(V_j^+)\subseteq\mathrm{Fil}^{-u}$. (ii) Every $\pknB$-linear $\lambda\colon\pktV\to\pknB$
with $\lambda(gv)=g\lambda(v)$ for $g\in\Gamma'$ is $\hat\mu$ for some form $\mu$. (iii) If $\lambda$ is as
in (ii) and $d\in\pktV$ is fixed by $\Gamma'$, then $\lambda(d)\in K'$; a nonzero element of $K'$ is a
unit of $\pknBp$.
\end{lemma}

\begin{proof}
(i) If $\mu$ is a form, $\hat\mu(g(b\otimes y))=\hat\mu(gb\otimes gy)=g(b)\,g(\mu(1\otimes y))=
g(\hat\mu(b\otimes y))$; the converse is the case $v=1\otimes y$. $V_j^+$ is the $\pknBp$-span of
$e_j(1\otimes L)$, whence the second claim. (ii) Let $\lambda_0\colon\pkpW\to\pknB$ be an $\pkpA$-linear
map with $\lambda_0(1)=1$ (it exists: proof of Lemma~\ref{pkn:frame-id}(iii)); then
$\mu(\sum_kw_k\otimes e_k)=\sum_k\lambda_0(w_k)\lambda(1\otimes e_k)$ is $\pkpA$-linear with $\hat\mu=
\lambda$, and it is a form by (i). (iii) $g(\lambda(d))=\lambda(gd)=\lambda(d)$ for $g\in\Gamma'$, and the
elements of $\pknB$ fixed by $G_{K'}$ lie in $K'$ (proof of Lemma~\ref{pkn:frame-id}(i)). A nonzero
$c\in K'$ has $\theta(c)=c\ne0$, so it is not in the maximal ideal $\mathrm{Fil}^1$ of $\pknBp$.
\end{proof}

\begin{lemma}[Hodge bases at a letter]\label{pkt:hodge}
For every letter $j$ there are elements $d_{j,1},\dots,d_{j,n_j}\in V_j$ fixed by $\Gamma'$ and integers
$\kappa_{j,1},\dots,\kappa_{j,n_j}\in U(j)$ such that $(t^{-\kappa_{j,m}}d_{j,m})_m$ is a $\pknBp$-basis of
$V_j^+$ (hence $(d_{j,m})_m$ is a $\pknB$-basis of $V_j$).
\end{lemma}

\begin{proof}
$L[1/\ell]$ is crystalline, hence de Rham, as a representation of $G_{K'}$ (P2). Let
$D=\pktV^{G_{K'}}$ with $\mathrm{Fil}^kD=D\cap t^kV^+$. By (P3), $\pknB\otimes_{K'}D\to\pktV$ is an
isomorphism carrying $\sum_kt^{-k}\pknBp\otimes\mathrm{Fil}^kD$ onto $V^+$. The $e_j$ have coefficients in
$\Lambda[1/\ell]\otimes\auO\subseteq K'\otimes\auO$, commute with $G_{K'}$, preserve $V^+$ and commute with
$t$; so $D=\bigoplus_jD_j$ with $D_j=e_jD$, and $\mathrm{Fil}^kD=\bigoplus_j(\mathrm{Fil}^kD\cap D_j)$.
Choose a $K'$-basis $d_{j,1},\dots,d_{j,n_j}$ of $D_j$ adapted to the filtration: $d_{j,m}\in
\mathrm{Fil}^{\kappa_{j,m}}D_j$ and, for each $k$, the $d_{j,m}$ with $\kappa_{j,m}\ge k$ form a basis of
$\mathrm{Fil}^kD_j$. Then $\sum_kt^{-k}\pknBp\otimes\mathrm{Fil}^kD_j=\bigoplus_mt^{-\kappa_{j,m}}\pknBp\,
d_{j,m}$ (an element $t^{-k}b\otimes d_{j,m}$ with $k\le\kappa_{j,m}$ lies in $t^{-\kappa_{j,m}}\pknBp
d_{j,m}$), and applying $e_j$ to the isomorphism of (P3) gives that $(t^{-\kappa_{j,m}}d_{j,m})_m$ is a
$\pknBp$-basis of $V_j^+$.

It remains to see $\kappa_{j,m}\in U(j)$. Reduce modulo $t$: $V_j^+/tV_j^+=e_j(C\otimes_{\Z_\ell}L)$. Since
$\auO[1/\ell]$ is a finite separable extension of $\Q_\ell$, $C\otimes_{\Z_\ell}\auO=\prod_\gamma C$ over the
embeddings $\gamma\colon\auO\to C$, so $C\otimes_{\Z_\ell}L=\bigoplus_\gamma C\otimes_{\auO,\gamma}L$, and
$e_j$ acts on the factor $\gamma$ by the scalar $\prod_{i'\ne j}(\gamma(\zeta)-\theta[i'])/(\theta[j]-
\theta[i'])$. As $\zeta^{\ell^f-1}=1$, $\gamma(\zeta)$ is the Teichm\"uller representative $\theta[i']$ of
its residue $i'$, the letter of $\gamma$; so the scalar is $1$ if $\gamma$ has letter $j$ and $0$
otherwise. Hence $V_j^+/tV_j^+=\bigoplus_{\gamma\text{ of letter }j}C\otimes_{\auO,\gamma}L$. On the
other hand $V_j^+/tV_j^+=\bigoplus_mC\cdot\overline{t^{-\kappa_{j,m}}d_{j,m}}$, and $g\in\Gamma'$ acts on
$\overline{t^{-\kappa}d}$ ($d$ fixed) through $\chi(g)^{-\kappa}$, so this summand is $C(-\kappa_{j,m})$. By
the uniqueness of the Hodge--Tate decomposition (Tate's theorem $C(n)^{G_{K'}}=0$ for $n\ne0$, as in
Lemma~\ref{pkm:BdR}), the multiset of the $\kappa_{j,m}$ is the union over the embeddings $\gamma$ of letter
$j$ of the multisets of jumps of $L$ at $\gamma$; these lie in $U(j)$.
\end{proof}

Fix such bases. Let $\delta_{j,m}\colon\pktV\to\pknB$ be the $\pknB$-linear map vanishing on $V_{j'}$
for $j'\ne j$ and sending $v\in V_j$ to its coefficient on $d_{j,m}$. It is $\Gamma'$-equivariant (the
$d_{j,m}$ are fixed and the decomposition $\bigoplus V_{j'}$ is $\Gamma'$-stable) and $\delta_{j,m}(V_{j'}^+)
\subseteq\mathrm{Fil}^{-\kappa_{j,m}}$ for all $j'$; so by Lemma~\ref{pkt:forms}, $\delta_{j,m}=\hat\mu$ for a
form $\mu$ of type $\kappa_{j,m}$ at every letter. We call $\delta_{j,m}$ the \emph{coordinate forms}.

\begin{lemma}[Expansion]\label{pkt:expand}
Let $\lambda\colon\pktV\to\pknB$ be $\pknB$-linear and $v\in V_j$. Then
$\lambda(v)=\sum_m\lambda(d_{j,m})\,\delta_{j,m}(v)$. If $\lambda(V_j^+)\subseteq\mathrm{Fil}^{s}$, then
$\lambda(d_{j,m})\in\mathrm{Fil}^{s+\kappa_{j,m}}$. If $\lambda$ is $\Gamma'$-equivariant and $\lambda(V_j^+)
\subseteq\mathrm{Fil}^{-u}$, then $\lambda(d_{j,m})\in K'$, and $\lambda(d_{j,m})=0$ unless $\kappa_{j,m}\le u$.
\end{lemma}

\begin{proof}
The first claim is linearity, $v=\sum_m\delta_{j,m}(v)d_{j,m}$. Next, $\lambda(t^{-\kappa_{j,m}}d_{j,m})\in
\mathrm{Fil}^s$ gives $\lambda(d_{j,m})\in\mathrm{Fil}^{s+\kappa_{j,m}}$. If $\lambda$ is equivariant,
$\lambda(d_{j,m})\in K'$ by Lemma~\ref{pkt:forms}(iii); being in $\mathrm{Fil}^{\kappa_{j,m}-u}$, it is $0$ if
$\kappa_{j,m}-u\ge1$ (a nonzero element of $K'$ is a unit of $\pknBp$).
\end{proof}

\begin{lemma}[Values at non-jumps]\label{pkt:gap}
Let $x\in L_i(c)$, $a\ge0$, $j=i^{(-a)}$, $u\in\Z$ (not necessarily in $U(j)$), and $\mu$ a form of type
$u$ at $j$. Then $\mu(\varphi^{-a}x)\in\sum_{v}\mathrm{Fil}^{c(a,v)}$, the sum over the $v\in U(j)$ with
$v\le u$ (the value is $0$ if there is no such $v$).
\end{lemma}

\begin{proof}
$y=\widehat{\varphi^{-a}x}\in V_j$ and $\mu(\varphi^{-a}x)=\hat\mu(y)=\sum_m\hat\mu(d_{j,m})\delta_{j,m}(y)$ by
Lemma~\ref{pkt:expand}. The coordinate form $\delta_{j,m}$ has type $\kappa_{j,m}\in U(j)$ at $j$, so
$\delta_{j,m}(y)\in\mathrm{Fil}^{c(a,\kappa_{j,m})}$ by the definition of orders. By
Lemmas~\ref{pkt:forms}(i) and \ref{pkt:expand}, $\hat\mu(d_{j,m})\in K'\subset\pknBp$ vanishes unless
$\kappa_{j,m}\le u$.
\end{proof}

\subsection{Orders add}\label{pkt:ss-add}

\begin{lemma}[Lemma~\ref{pkp:orders-add}]\label{pkt:add}
Let $L,L',L''$ be lattices with jump sets $U,U',U''$ and $b\colon L\times L'\to L''$ an $\auO$-bilinear
map with $b(gy,gy')=gb(y,y')$ for $g\in\Gamma'$; let $\tilde b\colon\pkpAmb(L)\times\pkpAmb(L')\to
\pkpAmb(L'')$ be its $\pkpW$-bilinear extension. (i) If $w\in\pknMinf(L)$, $y\in\pknMinf(L')$ then
$\tilde b(w,y)\in\pknMinf(L'')$; if $w\in\pkpA\mathcal N(L)$, $y\in\pkpA\mathcal N(L')$ then
$\tilde b(w,y)\in\pkpA\mathcal N(L'')$. (ii) Let $w\in L_i(o)$, $y\in L'_i(o')$, and let $o''$ satisfy
\[
o''(a,v)\le o(a,u)+o'(a,u')\quad\text{whenever }u+u'\le v,
\]
for all $a\ge0$, $u\in U(i^{(-a)})$, $u'\in U'(i^{(-a)})$, $v\in U''(i^{(-a)})$. Then $\tilde b(w,y)\in
L''_i(o'')$. In particular, for $L''=L\otimes_\auO L'$ with jump sets $U(i)+U'(i)$ and $b=\otimes$, $w\otimes
y\in(L\otimes L')_i(o'')$.
\end{lemma}

(Part (ii) is Lemma~\ref{pkp:orders-add}; part (i) is the membership statement used in its proof. For
$L\otimes L'$, the sets $U(i)+U'(i)$ contain the jumps: if $(d_{j,m})$, $(d'_{j,n})$ are Hodge bases of
$L,L'$ at $j$, the $e_j(d_{j,m}\otimes d'_{j,n})$ span $e_j(\pktV\otimes\pktV')$ and the argument of
Lemma~\ref{pkt:hodge} shows that the jumps at letter $j$ are sums $\kappa_{j,m}+\kappa'_{j,n}$.)

\begin{proof}
Let $\hat b\colon\pktV\times\pktV'\to\pktV''$ be the $\pknB$-bilinear extension of $b$; then
$\widehat{\tilde b(w,y)}=\hat b(\hat w,\hat y)$ for $w,y$ with denominators (multiply by the two
denominators), and $\hat b$ is $\Gamma'$-equivariant. Since $b$ is $\auO$-bilinear, $\tilde b(e_jw,y)=
e_j\tilde b(w,y)=\tilde b(w,e_jy)$, so $\hat b(V_j,V'_{j'})=0$ for $j\ne j'$ and $\hat b(V_j,V'_j)\subseteq
V''_j$; also $\hat b(V^+,V'^+)\subseteq V''^+$. Fix Hodge bases $(d'_{j,n},\kappa'_{j,n})$ of $L'$.

\emph{Key identity.} Let $\eta$ be a form on $L''$, $j$ a letter and $n$ an index. The map
$\lambda_{j,n}\colon v\mapsto\hat\eta(\hat b(v,d'_{j,n}))$ is $\pknB$-linear and $\Gamma'$-equivariant (as
$d'_{j,n}$ is fixed), so $\lambda_{j,n}=\hat\sigma_{j,n}$ for a form $\sigma_{j,n}$ on $L$ (Lemma~\ref{pkt:forms}(ii)).
For $v\in\pktV$ and $v'\in V'_j$, by Lemma~\ref{pkt:expand} applied to $v'\mapsto\hat\eta(\hat b(v,v'))$,
\[
\hat\eta(\hat b(v,v'))=\sum_n\hat\sigma_{j,n}(v)\,\delta'_{j,n}(v').\tag{$*$}
\]

(i) If $\tau^Nw,\tau^{N'}y\in\pkpA\otimes(\cdot)$ then $\tau^{N+N'}\tilde b(w,y)\in\pkpA\otimes L''$. Let
$k\ge1$ and $\eta$ a form on $L''$. As $\varphi$ acts on the first factor of $\pkpW\otimes(\cdot)$,
$\varphi^k\tilde b(w,y)=\tilde b(\varphi^kw,\varphi^ky)$. Write $\widehat{\varphi^ky}=\sum_jv'_j$ with
$v'_j\in V'_j$; by $(*)$,
$\eta(\varphi^k\tilde b(w,y))=\sum_j\sum_n\sigma_{j,n}(\varphi^kw)\,\delta'_{j,n}(\widehat{\varphi^k y})$. Each
factor is the value of a form at $\varphi^k$ of an element of $\pknMinf$, hence lies in $\pknBp$. So
$\tilde b(w,y)\in\pknMinf(L'')$. If $rw\in\pknMinf(L)$, $r'y\in\pknMinf(L')$ with $r,r'$ products of
the $p_a$, then $rr'\tilde b(w,y)=\tilde b(rw,r'y)\in\pknMinf(L'')$.

(ii) By (i), $\tilde b(w,y)\in\pkpA\mathcal N(L'')$, and it lies in the letter-$i$ part as
$e_i\tilde b(w,y)=\tilde b(e_iw,y)=\tilde b(w,y)$. Let $a\ge0$, $j=i^{(-a)}$, $v\in U''(j)$ and $\eta$ a form
of type $v$ at $j$. Put $\omega=\widehat{\varphi^{-a}w}\in V_j$ and $\upsilon=\widehat{\varphi^{-a}y}\in V'_j$;
then $\eta(\varphi^{-a}\tilde b(w,y))=\hat\eta(\hat b(\omega,\upsilon))=\sum_n\hat\sigma_{j,n}(\omega)
\delta'_{j,n}(\upsilon)$ by $(*)$. The coordinate form $\delta'_{j,n}$ has type $\kappa'_{j,n}\in U'(j)$ at
$j$, so $\delta'_{j,n}(\upsilon)\in\mathrm{Fil}^{o'(a,\kappa'_{j,n})}$. The form $\sigma_{j,n}$ has type
$v-\kappa'_{j,n}$ at $j$: $\hat\sigma_{j,n}(V_j^+)=t^{\kappa'_{j,n}}\hat\eta(\hat b(V_j^+,t^{-\kappa'_{j,n}}
d'_{j,n}))\subseteq t^{\kappa'_{j,n}}\hat\eta(V''^+_j)\subseteq\mathrm{Fil}^{\kappa'_{j,n}-v}$. By
Lemma~\ref{pkt:gap}, $\hat\sigma_{j,n}(\omega)\in\sum_u\mathrm{Fil}^{o(a,u)}$ over $u\in U(j)$ with $u\le
v-\kappa'_{j,n}$. Each product thus lies in some $\mathrm{Fil}^{o(a,u)+o'(a,u')}$ with $u\in U(j)$,
$u'=\kappa'_{j,n}\in U'(j)$, $u+u'\le v$, hence in $\mathrm{Fil}^{o''(a,v)}$.
\end{proof}

\subsection{Frobenius, the standard orders and Lemma~\ref{pkp:tau}}\label{pkt:ss-tau}

\begin{lemma}[Frobenius on orders]\label{pkt:frob}
(i) $\varphi(\pknMinf(L))\subseteq E^{-N}\pknMinf(L)$ for elements $x$ with $\tau^Nx\in\pkpA\otimes L$; hence
$\varphi^l(\pkpA\mathcal N(L))\subseteq\pkpA\mathcal N(L)$ for $l\ge0$. (ii) For $x\in\pkpA\mathcal N(L)$, $m\ge0$
and every form $\mu$, $\mu(\varphi^{m+1}x)\in\pknBp$. (iii) If $x\in L_{i^{(-l)}}(o)$, then $\varphi^lx\in
L_i(\pkcsh_{l,0}(o))$.
\end{lemma}

\begin{proof}
(i) Let $x\in\pknMinf$, $\tau^Nx\in\pkpA\otimes L$. Since $\varphi(\tau)=[\varepsilon]-1=E\tau v$ with
$v\in\pkpA^\times$ (Lemma~\ref{pkn:mu}) and $\varphi$ fixes $1\otimes y$, $\varphi(\tau^Nx)=E^N\tau^Nv^N\varphi(x)
\in\pkpA\otimes L$, so $\tau^N(E^N\varphi x)\in\pkpA\otimes L$; and for $k\ge1$, $\mu(\varphi^k(E^N\varphi
x))=\varphi^k(E)^N\mu(\varphi^{k+1}x)\in\pknBp$. So $E^N\varphi(x)\in\pknMinf$. If $rx\in\pknMinf$ with $r$ a
product of the $p_b$, then $\varphi(r)$ is a product of the $p_{b+1}$ and $E^N\varphi(r)\varphi(x)\in\pknMinf$,
so $\varphi(x)\in\pkpA\mathcal N$; iterate.

(ii) If $x\in\pknMinf$ this is the definition. In general let $x=r^{-1}x_0$ with $x_0\in\pknMinf$ and $r$ a
product of $p_b$'s; then $\varphi^{m+1}(r)$ is a product of $p_{b+m+1}$'s, whose images in $\pknBp$ are units
(Lemma~\ref{pkn:BdR}: $\varphi^{c}(E)$ with $c\ge1$), and $\mu(\varphi^{m+1}x)=\varphi^{m+1}(r)^{-1}\mu(\varphi^{m+1}
x_0)\in\pknBp$.

(iii) By (i), $\varphi^lx\in\pkpA\mathcal N$, and $\varphi^l$ maps the letter part at $i^{(-l)}$ to that at $i$
(Definition~\ref{pkp:def-letters}). Let $a\ge0$, $u\in U(i^{(-a)})$, $\mu$ a form of type $u$ at $i^{(-a)}$.
If $a<l$, $\mu(\varphi^{-a}\varphi^lx)=\mu(\varphi^{l-a}x)\in\pknBp=\mathrm{Fil}^0$ by (ii). If $a\ge l$,
$\mu(\varphi^{-a}\varphi^lx)=\mu(\varphi^{-(a-l)}x)\in\mathrm{Fil}^{o(a-l,u)}$, since $i^{(-a)}=(i^{(-l)})^{(-(a-l))}$.
\end{proof}

\begin{lemma}[Standard orders of level $0$]\label{pkt:std0}
If $x\in L_i(\mathrm{std}_0)$, then $x\in\pknMinf(L)$.
\end{lemma}

\begin{proof}
\emph{Denominator.} Let $T$ be the set of $z\in\pkpAmb(L)$ with $\tau^nz\in\pkpA\otimes L$ for some $n$. We
show: if $x'\in L_i(\mathrm{std}_0)$, $b\ge0$ and $p_bx'\in T$, then $x'\in T$. Let $\tau^np_bx'=\sum_kc_k
(1\otimes e_k)$, $c_k\in\pkpA$. Let $\mu_k$ be an $\pkpA$-linear map $\pkpAmb(L)\to\pknB$ with
$\mu_k(1\otimes y)=e_k^*(y)$ (the coordinate; construct it as in Lemma~\ref{pkt:forms}(ii)); then
$\hat\mu_k(V^+)\subseteq\pknBp$. With $j=i^{(-b)}$ and $\omega=\widehat{\varphi^{-b}x'}\in V_j$,
Lemma~\ref{pkt:expand} gives $\hat\mu_k(\omega)=\sum_m\hat\mu_k(d_{j,m})\delta_{j,m}(\omega)$ with
$\hat\mu_k(d_{j,m})\in\mathrm{Fil}^{\kappa_{j,m}}$ and $\delta_{j,m}(\omega)\in\mathrm{Fil}^{\mathrm{std}_0(b,
\kappa_{j,m})}=\mathrm{Fil}^{-\kappa_{j,m}}$. So $\mu_k(\varphi^{-b}x')\in\pknBp$. Applying $\mu_k\circ\varphi^{-b}$
to the expansion: $\varphi^{-b}(c_k)=\varphi^{-b}(\tau)^n\,E\,\mu_k(\varphi^{-b}x')\in\mathrm{Fil}^1$, since
$\varphi^{-b}(p_b)=E$. Hence $\theta(\varphi^{-b}c_k)=0$, i.e.\ $c_k\in p_b\pkpA$ (Definition~\ref{pkp:def-chart}),
say $c_k=p_bc'_k$, and $\tau^nx'=\sum_kc'_k(1\otimes e_k)$ ($p_b$ is a unit of $\pkpW$). Now $x\in\pkpA\mathcal
N$, so $rx\in\pknMinf\subseteq T$ for a product $r=p_{b_1}\cdots p_{b_k}$; by Lemma~\ref{pkn:E}, $p_{b_{s+1}}
\cdots p_{b_k}x\in L_i(o_s)$ with $o_s\ge\mathrm{std}_0$, so these elements lie in $L_i(\mathrm{std}_0)$ (orders
only increase), and removing the factors one at a time gives $x\in T$.

\emph{Values.} Let $n\ge1$ and $\mu$ a form. By Lemma~\ref{pkt:frob}(iii), $\varphi^nx\in
L_{i'}(\pkcsh_{n,0}(\mathrm{std}_0))$ with $i'=i^{\ell^n}$; at level $0<n$ these orders are $0$. With
$y=\widehat{\varphi^nx}\in V_{i'}$, $\mu(\varphi^nx)=\sum_m\hat\mu(d_{i',m})\delta_{i',m}(y)$, where
$\hat\mu(d_{i',m})\in K'\subset\pknBp$ (Lemma~\ref{pkt:forms}(iii)) and $\delta_{i',m}(y)\in\mathrm{Fil}^0$
(coordinate form of type $\kappa_{i',m}\in U(i')$, orders $0$ at level $0$). So $\mu(\varphi^nx)\in\pknBp$.
\end{proof}

\begin{proposition}[Lemma~\ref{pkp:tau}(i)]\label{pkt:tau-i}
For the lattice $\auO(1)$ with jump sets $\{-1\}$, every $\xi\in\auO(1)_i(\mathrm{std}_0)$ has $\tau\xi\in
\pkpA\otimes\auO(1)$.
\end{proposition}

\begin{proof}
By Lemma~\ref{pkt:std0}, $\xi\in\pknMinf(\auO(1))$. The jump bound of $\auO(1)$ (Theorem~\ref{pkf:main}) is
$\max\{|u|:u\in\{-1\}\}=1$; by Theorem~\ref{pkf:main} it has a frame for $a=1$, and by
Proposition~\ref{pkn:frame}(iii), $\tau\pknMinf\subseteq\pkpA\otimes\auO(1)$.
\end{proof}

Lemma~\ref{pkp:tau}(ii) is Lemma~\ref{pkp:AS}(vii) and (iii) is Lemma~\ref{pkm:tau-iii}.

\begin{lemma}[Lemma~\ref{pkp:tau}(iv)]\label{pkt:tau-iv}
Let $x\in\mathcal N(L)$ lie in the letter-$i$ part and $gx-x\in L_i(o')$ for all $g\in\Gamma'$. Then
$(z\tau)^{-1}(gx-x)\in L_i(o')$ for $g\in\Gamma'$. If $L=X_S$ and $\mathrm{pr}(x)$ is fixed by $\Gamma'$, then
$(z\tau)^{-1}(gx-x)$ lies over $0$. More generally, for $y\in\pkpA\mathcal N(L)$ and $r\in\pkpA$ not divisible by any $p_a$: if $ry\in
L_i(o)$ then $y\in L_i(o)$.
\end{lemma}

\begin{proof}
The last assertion first. $y$ lies in the letter-$i$ part as $r$ commutes with $e_i$ and is a nonzerodivisor on
$\pkpAmb(L)=\pkpW\otimes_{\Z_\ell}L$ ($\pkpW$ is a domain, $L$ is free, and $r\ne0$ since $p_0\nmid r$), so
$r\bigl(e_iy-y\bigr)=e_i(ry)-ry=0$ gives $e_iy=y$. For $a\ge0$, $\theta(\varphi^{-a}r)\ne0$ (else $r\in p_a\pkpA$), so the image of $\varphi^{-a}(r)$
is a unit of $\pknBp$, and $\mu(\varphi^{-a}y)=\varphi^{-a}(r)^{-1}\mu(\varphi^{-a}(ry))$ has the same filtration
as $\mu(\varphi^{-a}(ry))$. Now take $r=z\tau$ and $y=(z\tau)^{-1}(gx-x)$, which lies in $\pkpA\mathcal N$ by
(G$_0$). No $p_a$ divides $z\tau$: $\theta(\varphi^{-a}z)=\theta([\varpi^{\flat\,1/\ell^a}])$ is an
$\ell^a$-th root of $\pi\ne0$, and $\theta(\varphi^{-a}\tau)=\zeta_{\ell^{a+1}}-1\ne0$. So $y\in L_i(o')$.
Finally $\mathrm{pr}$ is $\pkpW$-linear and $\Gamma$-equivariant, so $\mathrm{pr}(y)=(z\tau)^{-1}(g\,\mathrm{pr}(x)-
\mathrm{pr}(x))=0$.
\end{proof}

\begin{lemma}[Lemma~\ref{pkp:tau}(v)]\label{pkt:tau-v}
If $x\in L_{i^{(-l)}}(o)$ then $\varphi^lx\in L_i(\pkcsh_{l,0}(o))$. If $x$ is a point of $S$ over
$\ell^ec$ at $i^{(-l)}$ with orders $o$ and depth $o'$, then $\varphi^lx$ is a point over $\ell^ec$ at $i$
with orders $\pkcsh_{l,0}(o)$ and depth $\pkcsh_{l,1}(o')$.
\end{lemma}

\begin{proof}
The first assertion is Lemma~\ref{pkt:frob}(iii). Let $x$ be a point. $\varphi^lx\in\mathcal N$ by ($\Phi$);
$\mathrm{pr}(\varphi^lx)=\varphi^l(\ell^ee_{i^{(-l)}}(1\otimes c))=\ell^ee_i(1\otimes c)$
(Definition~\ref{pkp:def-letters}); the orders are the first assertion. Depth: for $g\in\Gamma'$ put $y=
(z\tau)^{-1}(gx-x)$; $\mathrm{pr}(x)$ is fixed by $\Gamma'$, so $y\in L_{i^{(-l)}}(o')$ by
Lemma~\ref{pkt:tau-iv}, and $\varphi^ly\in L_i(\pkcsh_{l,0}(o'))$. As $g$ commutes with $\varphi$,
$g\varphi^lx-\varphi^lx=\varphi^l(z\tau)\cdot\varphi^ly$, which lies in $L_i(\pkcsh_{l,0}(o'))$. At a level $a<l$
and a form $\mu$ of type $u$ there, $\mu(\varphi^{-a}(g\varphi^lx-\varphi^lx))=\varphi^{l-a}(z\tau)\,
\mu(\varphi^{-a}\varphi^ly)$, where the second factor is in $\mathrm{Fil}^0$ and $\varphi^{l-a}(\tau)=
\varphi^{l-a-1}([\varepsilon]-1)$ has $\theta(\varphi^{l-a-1}([\varepsilon]-1))=1-1=0$, so the first factor is in
$\mathrm{Fil}^1$. Hence the value lies in $\mathrm{Fil}^1=\mathrm{Fil}^{\pkcsh_{l,1}(o')(a,u)}$; at levels $a\ge l$
the orders are $\pkcsh_{l,0}(o')=\pkcsh_{l,1}(o')$.
\end{proof}

\subsection{Points from killers: input (I1)}\label{pkt:ss-I1}

\emph{The two readings of ``killing the cokernel''.} Definition~\ref{pkp:def-points} asks for a lift
$n\in(X_S)_i(\mathrm{std}_0)\subseteq\pkpA\mathcal N(X_S)$, not necessarily in $\mathcal N(X_S)$. This
weaker reading is the definition (it is the one
proved by Propositions~\ref{pkm:kill-small}--\ref{pkm:kill-J}, through the lift ideals of
Definition~\ref{pkm:def-J}). For a multiplier in $\mathfrak S$ the two readings are equivalent, by the
retraction (R); this is the content of the next lemma, and it is why (I1) produces points in $\mathcal N$
itself. So (I1) as stated in \S\ref{pkc:ss-inputs} (a point in $\mathcal N(X_S)$, with depth) holds
(Remark~\ref{pkm:rem-kill}).

\begin{lemma}[Lifts in $\mathcal N$]\label{pkt:lift}
Let $S$ be a side, $i$ a letter, $l\ge0$, $f\in\mathfrak S$, and suppose that for every $c\in C_S$ there is
$x\in(X_S)_i(\mathrm{std}_l)$ with $\mathrm{pr}(x)=f\,e_i(1\otimes c)$. Then for every $c$ there is such an $x$
in $\mathcal N(X_S)$.
\end{lemma}

\begin{proof}
Let $\mathrm{ret}$ be the retractions of Lemma~\ref{pkm:inputs}(R) for $X_S$ and for the trivial lattice
$C_S$, and $q=\mathrm{ret}(x)$. By (R), $q\in\mathcal N(X_S)\cap(X_S)_i(\mathrm{std}_l)$ and $\mathrm{pr}(q)=
\mathrm{ret}(\mathrm{pr}(x))=\mathrm{ret}(f\cdot e_i(1\otimes c))$. By Lemma~\ref{pkm:inputs}(H) (second
assertion, $\Gamma$ acting trivially on $C_S$), $e_i(1\otimes c)\in\mathcal N(C_S)$; so by (R),
$\mathrm{ret}(f\,e_i(1\otimes c))=\lambda_0(f)\,e_i(1\otimes c)=f\,e_i(1\otimes c)$, as $\lambda_0$ is the
identity on $\mathfrak S$.
\end{proof}

\begin{proposition}[Input (I1)]\label{pkt:I1}
If $\ell^\kappa E^t$ kills the cokernel of a side $S$ at $i$ (in the sense of \S\ref{pkn:ss-def}), then
for every $c\in C_S$ there is a point of $S$ over $\ell^\kappa c$ at $i$, i.e.\ an element $\nu\in\mathcal
N(X_S)$, with orders $\rho_{-t}$ and depth $\rho_{-t}$: $\pkcPt_S(c,i,\kappa;\rho_{-t},\rho_{-t})\ne\emptyset$.
\end{proposition}

\begin{proof}
By Lemma~\ref{pkt:lift} with $l=0$ and $f=\ell^\kappa E^t\in\mathfrak S$, there is $q\in\mathcal N(X_S)\cap
(X_S)_i(\mathrm{std}_0)$ with $\mathrm{pr}(q)=\ell^\kappa E^te_i(1\otimes c)$. Put $\nu=E^{-t}q$. By
Lemma~\ref{pkn:M}(iii), $\nu\in\mathcal N(X_S)$; by Lemma~\ref{pkn:E} (with $b=0$, $n=-t$), $\nu\in
(X_S)_i(\mathrm{std}_0-t[a=0])=(X_S)_i(\rho_{-t})$; $\mathrm{pr}(\nu)=\ell^\kappa e_i(1\otimes c)$ as $\mathrm{pr}$
is $\pkpW$-linear. Depth: for $g\in\Gamma'$, $g\nu\in(X_S)_i(\rho_{-t})$ by Lemma~\ref{pkm:inputs}(G) (first
assertion), so $g\nu-\nu\in(X_S)_i(\rho_{-t})$.
\end{proof}

\subsection{The rule with a dual point: input (I3)}\label{pkt:ss-I3}

\begin{lemma}[Orders in the kernel]\label{pkt:sub}
Let $S$ be a side and $y\in\pkpAmb(Y)$ with image $\bar y=(y,0)\in(X_S)_i(o)$. Then $y\in Y_i(o)$.
\end{lemma}

\begin{proof}
Write $\iota\colon Y\to X_S$ for the inclusion (equivariant) and $\mathrm{fst}$ for the projection
($\auO$-linear, not equivariant). \emph{Extension of forms.} Let $\mu$ be a form on $Y$ of type $u$ at a
letter $j$ (possibly $u=\infty$, i.e.\ no condition). Let $D_Y\subseteq D_X$ be the spaces of
Lemma~\ref{pkt:hodge} for $Y$, $X_S$, and $D_{Y,j}=e_jD_Y\subseteq D_{X,j}$. Since $X_S=Y\oplus C_S$ as
$\auO$-modules, $V_Y\cap t^kV^+_X=t^kV^+_Y$, so the filtration of $D_{Y,j}$ is induced by that of $D_{X,j}$.
Hence an adapted basis of $D_{Y,j}$ extends to an adapted basis $(d_m)$ of $D_{X,j}$ (extend bases of the
graded pieces and lift), and by the proof of Lemma~\ref{pkt:hodge} the $t^{-\kappa_m}d_m$ form a
$\pknBp$-basis of $V^+_{X,j}$. Let $\lambda\colon\pktV(X_S)\to\pknB$ be the $\pknB$-linear map equal to
$\hat\mu$ on $V_{Y,j}$, zero on the added $d_m$, and zero on $V_{X,j'}$ for $j'\ne j$. It is
$\Gamma'$-equivariant (the added $d_m$ are fixed), $\lambda\circ\hat\iota=\hat\mu\circ e_j$, and
$\lambda(V^+_{X,j})=\hat\mu(V^+_{Y,j})\subseteq\mathrm{Fil}^{-u}$. By Lemma~\ref{pkt:forms}(ii), $\lambda=\hat\nu$
for a form $\nu$ on $X_S$, of type $u$ at $j$.

\emph{Membership.} If $r\bar y\in\pknMinf(X_S)$ with $r$ a product of $p_a$'s and $\tau^N r\bar y\in
\pkpA\otimes X_S$, then $\tau^Nry=\mathrm{fst}(\tau^Nr\bar y)\in\pkpA\otimes Y$, and for a form $\mu$ on $Y$
and $k\ge1$, $\mu(\varphi^k(ry))=\sum_j\hat\mu(e_j\widehat{\varphi^k(ry)})=\sum_j\nu_j(\varphi^k(r\bar y))
\in\pknBp$, with $\nu_j$ the extension of $\mu$ at $j$ just built. So $y\in\pkpA\mathcal N(Y)$; it lies in the
letter-$i$ part since $\iota$ commutes with $e_i$. \emph{Orders.} For $a\ge0$, $u\in U(i^{(-a)})$ and $\mu$ of
type $u$ at $j=i^{(-a)}$: $\widehat{\varphi^{-a}y}\in V_{Y,j}$, so $\mu(\varphi^{-a}y)=\nu(\varphi^{-a}\bar y)\in
\mathrm{Fil}^{o(a,u)}$ (the jump sets of $X_S$ are those of $Y$, Definition~\ref{pkp:def-side}).
\end{proof}

\begin{proposition}[Input (I3)]\label{pkt:I3}
Let $P$ pair sides $S,S'$ (Definition~\ref{pkp:def-pairing}), $T_P$ its Tate-dual side, $l,\kappa\ge0$, and
let $\nu$ be a point of $T_P$ over $\ell^\kappa\cdot1$ at $i^{(-l)}$ with orders $o_\nu$, where $o_\nu(a,\cdot)$
is nonincreasing for each $a$. Let $w\in(X_S)_i(o_w)$ lie over $0$ and let $y$ be a point of $S'$ over some
$\ell^{e'}c'$ at $i$ with depth $d_y$. If for all $a\ge0$, $u\in U_S(i^{(-a)})$, $u'\in U_{S'}(i^{(-a)})$
\[
0<\pkcsh_{l,0}(o_\nu)(a,-u-u'-1)+o_w(a,u)+d_y(a,u'),
\]
then the move of the first side by $w$ with $y$ at rest is $\kappa$-safe. Symmetrically, if $x$ is a point of
$S$ over $\ell^ec$ at $i$ with depth $d_x$, $w'\in(X_{S'})_i(o_{w'})$ lies over $0$, and $0<\pkcsh_{l,0}(o_\nu)
(a,-u-u'-1)+d_x(a,u)+o_{w'}(a,u')$ for all such $a,u,u'$, then the move of the second side by $w'$ with $x$ at
rest is $\kappa$-safe.
\end{proposition}

\begin{proof}
\emph{Step 1: safe pairs.} We show: if $w\in(X_S)_i(o_1)$ and $v\in(X_{S'})_i(o_2)$ both lie over $0$ and
$0<o_3(a,-u-u'-1)+o_1(a,u)+o_2(a,u')$ for all $a,u,u'$ as above, where $o_3=\pkcsh_{l,0}(o_\nu)$, then
$z\tau\,\ell^\kappa P(\mathrm{fst}\,w,\mathrm{fst}\,v)\in\pkpsm$. Since $w,v$ lie over $0$ they lie in
$\pkpAmb(Y)$, $\pkpAmb(Y')$ (the kernels of $\mathrm{pr}$), and by Lemma~\ref{pkt:sub} they have orders $o_1$,
$o_2$ there. By
Lemma~\ref{pkt:tau-v}, $\varphi^l\nu$ lies over $\ell^\kappa\cdot1$ at $i$ with orders $o_3$; $o_3(a,\cdot)$ is
nonincreasing. Let $o''(a,v_0)=1-o_3(a,-v_0-1)$. By Lemma~\ref{pkt:add}(ii) with $b=\otimes$, $w\otimes v\in
(Y\otimes Y')_i(o'')$: if $u+u'\le v_0$ then $-v_0-1\le-u-u'-1$, so $o''(a,v_0)\le1-o_3(a,-u-u'-1)\le
o_1(a,u)+o_2(a,u')$, since $0<o_3(a,-u-u'-1)+o_1(a,u)+o_2(a,u')$ is an inequality between integers. Let $\epsilon$ be the equivariant form of
Lemma~\ref{pkp:extform} on $X_{T_P}\times(Y\otimes Y')$ with values in $\auO(1)$ (jump set $\{-1\}$; the jump
sets of $T_P$ are $\{0\}\cup\{-s-1:s\in U_S(\cdot)+U_{S'}(\cdot)\}$). By Lemma~\ref{pkt:add}(ii), $\xi=
\tilde\epsilon(\varphi^l\nu,w\otimes v)\in\auO(1)_i(\mathrm{std}_0)$: we need $1=\mathrm{std}_0(a,-1)\le o_3(a,u_3)+
o''(a,v_0)$ whenever $u_3+v_0\le-1$; and then $u_3\le-v_0-1$ gives $o_3(a,u_3)\ge o_3(a,-v_0-1)=1-o''(a,v_0)$.
By Proposition~\ref{pkt:tau-i}, $\tau\xi\in\pkpA\otimes\auO(1)$, and by Lemma~\ref{pkp:AS}(vii) the image of
$z\tau\xi$ in $\pkpM$ lies in $\pkpsm$. By Lemma~\ref{pkp:extform} this image is $z\tau\,P(\mathrm{pr}(\varphi^l
\nu)(w\otimes v))$ with $\mathrm{pr}(\varphi^l\nu)=\ell^\kappa e_i(1\otimes1)$; since $e_iw=w$, this is $z\tau\,\ell^\kappa P(w,v)$
(exactly as in the proof of Proposition~\ref{pkp:rule}). So $z\tau\ell^\kappa P(\mathrm{fst}\,w,
\mathrm{fst}\,v)\in\pkpsm$.

\emph{Step 2: the first move.} Let $g\in\Gamma'$. By Lemma~\ref{pkt:tau-iv}, $v=(z\tau)^{-1}(gy-y)\in(X_{S'})_i
(d_y)$ lies over $0$ ($\mathrm{pr}(y)=\ell^{e'}e_i(1\otimes c')$ is fixed by $\Gamma'$). Step 1 with $o_1=o_w$,
$o_2=d_y$ gives $\ell^\kappa P(\mathrm{fst}\,w,\mathrm{fst}(gy-y))=z\tau\ell^\kappa P(\mathrm{fst}\,w,\mathrm{fst}\,
v)\in\pkpsm$, i.e.\ $(w,gy-y)$ is $\kappa$-safe (Definition~\ref{pkp:def-safe}).

\emph{Step 3: the second move.} Let $g\in\Gamma'$. $v=(z\tau)^{-1}(gx-x)\in(X_S)_i(d_x)$ lies over $0$
(Lemma~\ref{pkt:tau-iv}); $gw'\in(X_{S'})_i(o_{w'})$ (Lemma~\ref{pkm:inputs}(G), first assertion) and lies
over $0$. Step 1 with the roles of the two sides exchanged in the bookkeeping ($o_1=d_x$ for $v$, $o_2=
o_{w'}$ for $gw'$) gives $\ell^\kappa P(\mathrm{fst}(gx-x),\mathrm{fst}\,gw')\in\pkpsm$.
\end{proof}

\begin{remark}[Proposition~\ref{pkp:rule} as a special case]\label{pkt:rem-rule}
With $\nu$ given by Proposition~\ref{pkt:I1} (orders $\rho_{-t}$, nonincreasing in $u$) and the identity
\eqref{pkp:eq-ruleid}, the hypothesis $\mathrm{rule}_{l,t}(a,u+u')\le o(a,u)+d'(a,u')$ of
Proposition~\ref{pkp:rule} is the hypothesis of Proposition~\ref{pkt:I3}. Thus Proposition~\ref{pkp:rule} also follows
from Proposition~\ref{pkt:I3}; its direct proof uses the inequalities $u+u'\le v$ and $u+u''\le-1$, in
accordance with Lemma~\ref{pkt:add}.
\end{remark}

\subsection{References for the inputs}\label{pkt:ss-open}

The statements (G$_0$), ($\Phi$) used in Lemmas~\ref{pkt:tau-iv}, \ref{pkt:tau-v} are
Proposition~\ref{pkl:G0} and Corollary~\ref{pkl:Phi}; the items (R), (H), (G) of Lemma~\ref{pkm:inputs}
used in Proposition~\ref{pkt:I1} are Propositions~\ref{pkl:R}, \ref{pkl:H2}, Corollaries~\ref{pkl:H1},
\ref{pkl:H3} and Lemma~\ref{pkl:G1}. The input (I2) is Theorem~\ref{pke:I2}, (I4) is
Theorem~\ref{pkh:I4}, (I5) is Theorem~\ref{pkg:I5}, and Proposition~\ref{pki:chart} is
Theorem~\ref{pkj:chart}.


%% ===== proofs/C-kisin-deep.tex =====
% proofs/C-kisin-deep.tex
% Deepening with costs: the case input (I2); the value fields K_a; the depth step at level a.
% Labels start with pke:. Conventions and macros as in proofs/C-kisin-tau.tex and C-kisin-lattices.tex.
\providecommand{\auO}{\mathcal O}
\providecommand{\auQlb}{\overline{\Q}_\ell}
\providecommand{\panot}[1]{\par\smallskip{\raggedright\noindent\textbf{[Not proved here:} #1\textbf{]}\par}\smallskip}
\providecommand{\pkpA}{A_{\mathrm{inf}}}
\providecommand{\pkpW}{W(C^\flat)}
\providecommand{\pkpAmb}{\mathbb A}
\providecommand{\pknBp}{B^+_{\mathrm{dR}}}
\providecommand{\pknB}{B_{\mathrm{dR}}}
\providecommand{\pknMinf}{\mathfrak M_{\mathrm{inf}}}
\providecommand{\pknK}{\mathcal K_0}
\providecommand{\pkfBp}{B^+}
\providecommand{\pkfSB}{\mathfrak S_B}
\providecommand{\pkcPt}{\mathrm{Pt}}
\providecommand{\pkcimp}{\mathrm{imp}}
% macros of this file (prefix \pke)
\providecommand{\pkeK}[1]{\mathcal K_{#1}}
\providecommand{\pkeAN}{\pkpA\mathcal N}
\providecommand{\pkecost}{\mathrm{cost}}
\providecommand{\pkeDp}{\mathrm{Dp}}

\section{Deepening with costs: the input (I2)}\label{pke:sec}

This section proves the input (I2) of \S\ref{pkc:ss-inputs} (deepening of points, with the costs
$\pkecost(a,n)=v_\ell(n)+d_v+3a$, $v_\ell(0)=0$). On the way it introduces the value fields
$\pkeK a$ and proves their properties (\S\ref{pke:ss-fields}), a version of item (C)(d) of
Lemma~\ref{pkm:inputs} at the level $a$, for all $n\ne0$ (\S\ref{pke:ss-chart}), and the depth step at
the level $a$ (\S\ref{pke:ss-step}).

\emph{Dependencies.} Theorem~\ref{pke:I2} ((I2)) is proved in full, modulo results cited by label from
\S\S\ref{pkp:sec}--\ref{pkt:sec} and \S\ref{pkl:sec}: the frame theorem (Theorem~\ref{pkf:main},
Proposition~\ref{pkn:frame}), Corollary~\ref{pkl:H1} (item (H)), Proposition~\ref{pkl:R} (item (R)),
Corollary~\ref{pkl:Phi} ($\varphi(\mathcal N)\subseteq\mathcal N$), Lemma~\ref{pkl:G1} (stability under
$\Gamma'$), Proposition~\ref{pkr:I8}(c) and Lemma~\ref{pkf:extend} (extension of
$\mathfrak S$-linear maps), Lemma~\ref{pkm:BdR} (Tate), and the statement (G$_0$) of \S\ref{pkt:sec}
together with the second assertion of Lemma~\ref{pkm:inputs}(G), which follows from (G$_0$) by
Lemma~\ref{pkt:tau-iv}. (G$_0$) is Proposition~\ref{pkl:G0}. Items (I4), (I5) and
Proposition~\ref{pki:chart} are proved elsewhere (\S\ref{pke:ss-open}).

\emph{Conventions.} We use those of \S\S\ref{pkn:sec}, \ref{pkt:sec}, \ref{pkl:sec}: the Kummer chart
(Definition~\ref{pkp:def-chart}) with $\mathfrak S=\Lambda[[z]]$, $E$, $p_a=\varphi^a(E)\in\mathfrak S$,
$\Gamma'$, $K'=\auQlb^{\,j(\Gamma')}$; $j(\Gamma')=G_{K'}$ (Definition~\ref{pkp:def-chart}: every element
of $G_{\Q_\ell}$ fixing $\Lambda$ and $\theta(z)$ lies in $j(\Gamma')$, and $K'=K_0(\pi)$); the modules
$\pknMinf(L)\subseteq\pkeAN(L)\supseteq\mathcal N(L)$ and $L_i(o)$ of Definition~\ref{pkn:def-N};
$\mathrm{Fil}^n=t^n\pknBp=E^n\pknBp$ (Lemma~\ref{pkn:BdR}). The map $\pkpA\to\pknBp$ is injective
(\S\ref{pkn:ss-frame}), and $\varphi$ is an automorphism of $\pkpA$ commuting with $\Gamma$. For
$r\in\pkpA$ and $a\ge0$ we write $r^{(a)}$ for the image of $\varphi^{-a}(r)$ in $\pknBp$; $r\mapsto r^{(a)}$
is an injective ring map, and $g(r^{(a)})=(gr)^{(a)}$ for $g\in\Gamma'$. Put $\theta_a=\theta\circ
\varphi^{-a}\colon\pkpA\to O_C$; it is a surjective, $\Gamma'$-equivariant ring map with kernel $p_a\pkpA$
(Corollary~\ref{pkn:E-gen}), and $r^{(a)}\in\mathrm{Fil}^1$ iff $\theta_a(r)=0$ iff $p_a\mid r$. Note
$p_a^{(a)}=E$, a generator of $\mathrm{Fil}^1$. For a form $\mu$ on a lattice $L$ and $x\in\pkpAmb(L)$
put $\mu_a(x)=\mu(\varphi^{-a}x)$; then $\mu_a(rx)=r^{(a)}\mu_a(x)$ for $r\in\pkpA$, and
$\mu_a(gx)=g\,\mu_a(x)$ for $x\in\pkeAN(L)$, $g\in\Gamma'$ (Lemma~\ref{pkm:equiv}). On $O_C$, $v$ is the
valuation with $v(\ell)=1$; for $x,y\in O_C$, $x\mid y$ iff $v(x)\le v(y)$. An element $x\in\pkeAN(L)$ is
\emph{deep for $o$ at $i$} if $gx-x\in L_i(o)$ for all $g\in\Gamma'$. Finally $\delta=z\tau$; its images
$\delta^{(a)}$ are units of $\pknBp$ (Lemma~\ref{pkl:basic}(i)), and $\delta$ is a unit of $\pkpW$.

\subsection{The value fields}\label{pke:ss-fields}

\begin{definition}[Value fields]\label{pke:def-K}
For $a\ge0$ let $R_a\subseteq\pknB$ be the $K'$-span of the elements $b\,s^{(a)}$ with $b\in\pkfSB$ (its
image in $\pknBp$) and $s\in\mathfrak S$, and let $\pkeK a$ be the subfield of $\pknB$ generated over
$K'$ by the same elements.
\end{definition}

For $a=0$ this is the field $\pknK$ of Definition~\ref{pkn:def-N}.

\begin{lemma}\label{pke:K-basic}
(i) $R_a$ is a subring of $\pknB$ and $\pkeK a$ is its field of fractions: every element of $\pkeK a$ is
$w/d$ with $w,d\in R_a$, $d\ne0$. (ii) $K'\subseteq\pkeK a$; $s^{(a)}\in\pkeK a$ for $s\in\mathfrak S$, in
particular $E=p_a^{(a)}\in\pkeK a$; and $\pknK\subseteq\pkeK a$.
\end{lemma}

\begin{proof}
(i) $(b\,s^{(a)})(b'\,s'^{(a)})=(bb')(ss')^{(a)}$ with $bb'\in\pkfSB$ (a ring) and $ss'\in\mathfrak S$; and
$K'$ consists of scalars. So $R_a$ is a $K'$-algebra, a domain inside the field $\pknB$, and the field it
generates is $\{w/d\}$. (ii) Take $b=1$, resp.\ $s=1$, $b=1$; $p_a\in\mathfrak S$ as $\varphi(\mathfrak S)
\subseteq\mathfrak S$. The generators $b\,s$ ($b\in\pkfSB$, $s\in\mathfrak S$) of $\pknK$ satisfy
$bs=(bs)\cdot1^{(a)}$ with $bs\in\pkfSB$, since the image of $\mathfrak S$ lies in $\pkfSB$
(\S\ref{pkf:ss-SB}).
\end{proof}

\begin{lemma}[Retraction of values at level $a$]\label{pke:valret}
Let $a\ge0$, $f\colon\pkpA\to\pkpA$ $\mathfrak S$-linear, $v_1,\dots,v_n\in\pkeK a$, $r_1,\dots,r_n\in\pkpA$,
and $W\subseteq\pknB$ a $\pknBp$-submodule. If $\sum_lr_l^{(a)}v_l\in W$, then $\sum_lf(r_l)^{(a)}v_l\in W$.
\end{lemma}

\begin{proof}
Put $f'(x)=\varphi^{-a}(f(\varphi^ax))$. It is additive, and for $s\in\mathfrak S$, $f'(sx)=
\varphi^{-a}(\varphi^a(s)f(\varphi^ax))=sf'(x)$ because $\varphi^a(s)\in\mathfrak S$; so $f'$ is $\mathfrak
S$-linear, and Lemma~\ref{pkf:extend} gives its extension $G_{f'}$ to $\pkfBp$, with $G_{f'}(bx)=bG_{f'}(x)$
for $b\in\pkfSB$, $x\in\pkfBp$. By Lemma~\ref{pke:K-basic}(i) choose $d\in R_a\smallsetminus\{0\}$ with
$w_l=dv_l\in R_a$, and write $w_l=\sum_ic_{li}b_{li}s_{li}^{(a)}$ (finite sums, $c_{li}\in K'$,
$b_{li}\in\pkfSB$, $s_{li}\in\mathfrak S$). Then
\[
\textstyle\sum_lw_lr_l^{(a)}=\sum_{l,i}c_{li}\,y_{li},\qquad y_{li}=b_{li}\,\varphi^{-a}(s_{li}r_l)\in\pkfBp,
\]
where $\varphi^{-a}(s_{li}r_l)\in\pkpA\subseteq\pkfBp$ and images are taken in $\pknBp$. Now
$G_{f'}(y_{li})=b_{li}f'(\varphi^{-a}(s_{li}r_l))=b_{li}\varphi^{-a}(f(s_{li}r_l))=b_{li}\varphi^{-a}(s_{li})
\varphi^{-a}(f(r_l))$, by Lemma~\ref{pkf:extend} and $\mathfrak S$-linearity of $f$. By
Proposition~\ref{pkr:I8}(c) there is $u\in\pknBp$ with $\sum c_{li}G_{f'}(y_{li})=u\sum c_{li}y_{li}$, i.e.\
$\sum_lw_lf(r_l)^{(a)}=u\sum_lw_lr_l^{(a)}$. Dividing by $d$: $\sum_lf(r_l)^{(a)}v_l=u\sum_lr_l^{(a)}v_l\in
W$.
\end{proof}

\begin{lemma}[Values lie in $\pkeK a$]\label{pke:values}
Let $L$ be a lattice with jump sets, $x\in\mathcal N(L)$, $\mu$ a form and $a\ge0$. Then $\mu_a(x)\in\pkeK
a$.
\end{lemma}

\begin{proof}
By Lemma~\ref{pkn:M}(ii) there is $r'\ne0$ with $r'x\in\pkpA\otimes L$; then $r\varphi^{-a}x\in\pkpA\otimes L$
with $r=\varphi^{-a}(r')\ne0$. By Corollary~\ref{pkl:H1}, $\tau^H(\pkpA\otimes L)\subseteq\pkeAN(L)$ ($H$ the
jump bound), and $\pkeAN(L)$ is the $\pkpA$-span of $\mathcal N(L)$ (Proposition~\ref{pkn:frame}(iv) with
Theorem~\ref{pkf:main}). So $\tau^Hr\varphi^{-a}x=\sum_{l=1}^n t_lv_l$ with $t_l\in\pkpA$, $v_l\in\mathcal N(L)$.
Apply $\varphi^a$: with $\rho=\varphi^a(\tau^Hr)\ne0$,
\[
\textstyle-\rho\,x+\sum_l\varphi^a(t_l)\,\varphi^a(v_l)=0,
\]
a relation with coefficients in $\pkpA$ among elements of $\mathcal N(L)$: $x\in\mathcal N(L)$ and
$\varphi^a(v_l)\in\mathcal N(L)$ (Corollary~\ref{pkl:Phi}). By Proposition~\ref{pkr:I5} some coordinate
$\kappa_j$ has $\kappa_j(\rho)\ne0$; it is $\mathfrak S$-linear with values in $\mathfrak S$. By
Lemma~\ref{pkl:rel} (with $f=\kappa_j$), $\kappa_j(\rho)x=\sum_l\kappa_j(\varphi^at_l)\varphi^a(v_l)$. Apply
$\mu_a$: $\kappa_j(\rho)^{(a)}\mu_a(x)=\sum_l\kappa_j(\varphi^at_l)^{(a)}\mu(v_l)$. Here $\mu(v_l)\in\pknK
\subseteq\pkeK a$ (as $v_l\in\mathcal N$), $\kappa_j(\cdot)^{(a)}\in\pkeK a$ (values in $\mathfrak S$,
Lemma~\ref{pke:K-basic}(ii)), and $\kappa_j(\rho)^{(a)}\ne0$. So $\mu_a(x)\in\pkeK a$.
\end{proof}

\begin{lemma}[Galois action on integral values]\label{pke:fix}
Let $a\ge0$, $g,t_1,\dots,t_m\in\Gamma'$, and suppose $p_a$ divides $g(\sum_lt_ls)-\sum_lt_ls$ in $\pkpA$
for every $s\in\mathfrak S$. If $y\in\pkeK a\cap\mathrm{Fil}^0$, then $g(\sum_lt_ly)-\sum_lt_ly\in
\mathrm{Fil}^1$.
\end{lemma}

\begin{proof}
$\theta(y)\in C$, so $\ell^e\theta(y)\in O_C$ for some $e\ge0$; as $\theta\colon\pkpA\to O_C$ is onto,
choose $r\in\pkpA$ with $\theta(r)=\ell^e\theta(y)$, i.e.\ $\ell^ey-r\in\ker(\theta|\pknBp)=\mathrm{Fil}^1$.
Since the image of $r$ in $\pknBp$ is $(\varphi^ar)^{(a)}$ and $(\ell^e)^{(a)}=\ell^e$, this reads
$(\ell^e)^{(a)}y+(\varphi^ar)^{(a)}\cdot(-1)\in\mathrm{Fil}^1$. Lemma~\ref{pke:valret} with $f=\kappa_0$ (Proposition~\ref{pkr:I5}; $\kappa_0(\ell^e)=
\ell^e$), the values $y,-1\in\pkeK a$, the coefficients $\ell^e,\varphi^a(r)$ and $W=\mathrm{Fil}^1$ gives
$\ell^ey-s^{(a)}\in\mathrm{Fil}^1$ with $s=\kappa_0(\varphi^ar)\in\mathfrak S$. Elements of $\Gamma'$ preserve
$\mathrm{Fil}^1$ ($g(t)=\chi(g)t$), so $\sum_lt_l(\ell^ey-s^{(a)})$ and its image under $g$ lie in
$\mathrm{Fil}^1$. Also $g(\sum_lt_ls^{(a)})-\sum_lt_ls^{(a)}=\bigl(g(\sum_lt_ls)-\sum_lt_ls\bigr)^{(a)}\in
(p_a\pkpA)^{(a)}\subseteq E\pknBp=\mathrm{Fil}^1$. Adding, $\ell^e\bigl(g(\sum_lt_ly)-\sum_lt_ly\bigr)\in
\mathrm{Fil}^1$, and $\ell$ is a unit of $\pknBp$.
\end{proof}

\begin{lemma}[The retraction and values]\label{pke:ret}
Let $\mathrm{ret}\colon\pkeAN(L)\to\mathcal N(L)$ be the retraction of Proposition~\ref{pkl:R}. Let $y\in
\pkeAN(L)$, $\mu$ a form, $a\ge0$, $v\in\pkeK a$ and $W\subseteq\pknB$ a $\pknBp$-submodule. If
$\mu_a(y)-v\in W$, then $\mu_a(\mathrm{ret}\,y)-v\in W$. Consequently $\mathrm{ret}(L_i(o))\subseteq
\mathcal N(L)\cap L_i(o)$ for every letter $i$ and all orders $o$.
\end{lemma}

\begin{proof}
Write $y=\sum_sr_sw_s$ with $r_s\in\pkpA$, $w_s\in\mathcal N(L)$ (Proposition~\ref{pkn:frame}(iv)); then
$\mathrm{ret}\,y=\sum_s\kappa_0(r_s)w_s$ (Proposition~\ref{pkl:R}). The values $v_s=\mu_a(w_s)$ lie in
$\pkeK a$ (Lemma~\ref{pke:values}), and $\mu_a(y)-v=\sum_sr_s^{(a)}v_s+1^{(a)}(-v)\in W$. Lemma~\ref{pke:valret}
with $f=\kappa_0$ ($\kappa_0(1)=1$) gives $\sum_s\kappa_0(r_s)^{(a)}v_s-v=\mu_a(\mathrm{ret}\,y)-v\in W$. For
the consequence let $y\in L_i(o)$: $\mathrm{ret}\,y\in\mathcal N(L)\subseteq\pkeAN(L)$, it lies in the
letter-$i$ part as $\mathrm{ret}$ commutes with $e_i$ (Proposition~\ref{pkl:R}), and for $\mu$ of type $u$
at $i^{(-a)}$ take $v=0$, $W=\mathrm{Fil}^{o(a,u)}$.
\end{proof}

\subsection{Chart identities at the level $a$}\label{pke:ss-chart}

\begin{lemma}\label{pke:triv}
Let $a\ge0$, $m\ge0$ with $\ell^a\mid m$, and $g\in\Gamma'$ with $gz=[\varepsilon]^mz$. Then $p_a\mid
gs-s$ for every $s\in\mathfrak S$.
\end{lemma}

\begin{proof}
$gz-z=z([\varepsilon]^m-1)$, and $\varphi^{-a}([\varepsilon]^m)=[\varepsilon^{m/\ell^a}]$ has
$\theta_a([\varepsilon]^m)=(\varepsilon^\sharp)^{m/\ell^a}=1$. So $\theta_a(gz-z)=0$ and $p_a\mid gz-z$. By the
argument of Step~2 of the proof of Proposition~\ref{pkn:Ca} (with $\eta=\varepsilon^m$), $gz-z$ divides
$gs-s$ for every $s\in\mathfrak S$.
\end{proof}

\begin{lemma}[Averaging]\label{pke:avg}
Let $a\ge0$, $c\ge a$, $N=\ell^c$, and for $0\le l<N$ let $t_l\in\Gamma'$ with $t_lz=[\varepsilon]^lz$
(Definition~\ref{pkp:def-chart}). Then for every $g\in\Gamma'$ and $s\in\mathfrak S$, $p_a$ divides
$g(\sum_{l<N}t_ls)-\sum_{l<N}t_ls$.
\end{lemma}

\begin{proof}
It suffices that $\theta_a(\sum_lt_ls)$ is fixed by $j(\Gamma')=G_{K'}$, since $\theta_a$ is
$\Gamma'$-equivariant with kernel $p_a\pkpA$. Write $s=\sum_k\lambda_kz^k$, $\lambda_k\in\Lambda$. Put
$\varpi_a=\theta_a(z)$, an $\ell^a$-th root of $\pi$, and $\zeta_a=\theta_a([\varepsilon])=
\theta([\varepsilon^{1/\ell^a}])$, a primitive $\ell^a$-th root of unity. As $t_l$ fixes $\Lambda$ and is
continuous, $t_ls=\sum_k\lambda_k[\varepsilon]^{lk}z^k$ (as in Step~2 of the proof of
Proposition~\ref{pkn:Ca}); $\theta_a$ is continuous for the $(\ell,z)$-adic topology ($v(\varpi_a)>0$), and
$\theta_a(\lambda_k)=\theta(\varphi^{-a}\lambda_k)\in W(\mathbb F_{\ell^f})\subseteq K'$. Hence
\[
\textstyle\theta_a\bigl(\sum_{l<N}t_ls\bigr)=\sum_k\theta_a(\lambda_k)\varpi_a^k\sum_{l<N}\zeta_a^{lk}
=N\sum_{j\ge0}\theta_a(\lambda_{j\ell^a})\pi^j,
\]
because $\sum_{l<N}\zeta_a^{lk}$ is $N$ if $\ell^a\mid k$ and $0$ otherwise ($\ell^a\mid N$). The last series
converges in $O_C$ with partial sums in the complete ring $O_{K'}$, so its sum lies in $K'$ and is fixed by
$G_{K'}$.
\end{proof}

\begin{proposition}[Item (C)(d) at the level $a$]\label{pke:Cd}
Let $a\ge0$ and $n\in\Z\smallsetminus\{0\}$, and put $c=\pkecost(a,n)=v_\ell(n)+d_v+3a$. There are $g\in
\Gamma'$, $\omega'\in\pkpA^\times$ and $\upsilon,X\in\pkpA$ with
\[
gp_a=\omega'p_a,\qquad p_a\mid gs-s\ \ (s\in\mathfrak S),\qquad p_a\mid\upsilon X-\ell^c,\qquad
\delta X=\omega'^n-1 .
\]
\end{proposition}

\begin{proof}
Let $g\in\Gamma'$ with $gz=[\varepsilon]^{\ell^a}z$ (Definition~\ref{pkp:def-chart}), and $gE=\omega E$ with
$\omega\in\pkpA^\times$ (Lemma~\ref{pkf:approx0}(iv)). Put $\omega'=\varphi^a(\omega)$; then $gp_a=
\varphi^a(gE)=\omega'p_a$ ($g$ commutes with $\varphi$). The second property is Lemma~\ref{pke:triv}.

\emph{Step 1: $\delta\mid\omega-1$.} By Proposition~\ref{pkn:Ca}, $\delta E\mid gE-E=E(\omega-1)$; cancel $E$.
Hence $\delta\mid\omega^k-1$ for $k\ge1$, and $\omega^{-k}-1=-\omega^{-k}(\omega^k-1)$; so $\delta\mid
\omega^n-1$.

\emph{Step 2: $v(\theta(\omega-1))\le d_v+a+v(\theta(\delta))$, and $\theta(\omega-1)\ne0$.} Let $P,V$ be as in
Lemma~\ref{pkl:different}, $P(z)=EG$ with $G=V(z)^{-1}\in\mathfrak S^\times$. By Lemma~\ref{pkn:mu},
$gz-z=z([\varepsilon]-1)w=z\tau Euw=Eh$, with $w=\sum_{k<\ell^a}[\varepsilon]^k$, $u\in\pkpA^\times$ and $h=
\delta uw$. Taylor's formula for the polynomial $P$ gives $c'\in\pkpA$ with $P(gz)=P(z)+P'(z)Eh+E^2h^2c'$. As $g$
fixes $\Lambda$, $P(gz)=g(EG)=\omega Eg(G)$; cancelling $E$, $\omega g(G)=G+P'(z)h+Eh^2c'$. Apply $\theta$:
$\theta(gG)=\theta(G)$ (as in Step~2 of the proof of Proposition~\ref{pkl:Cd}), $\theta(E)=0$, $\theta(w)=\ell^a$
($\theta([\varepsilon])=1$). So $\theta(\omega-1)\theta(G)=P'(\pi)\theta(\delta)\theta(u)\ell^a$, with
$\theta(G),\theta(u)$ units, $P'(\pi)\ne0$ ($P$ is separable) and $v(P'(\pi))\le d_v$
(Lemma~\ref{pkl:different}); and $\theta(\delta)=\pi(\zeta_\ell-1)\ne0$.

\emph{Step 3: $v(\theta(\omega^n-1))=v_\ell(n)+v(\theta(\omega-1))$.} For $n\ge1$ this is Step~3 of the proof
of Proposition~\ref{pkl:Cd}, whose argument uses only $\delta\mid\omega-1$ and $\theta(\omega-1)\ne0$ (Steps 1
and 2 here). For $n<0$, $\omega^n-1=-\omega^n(\omega^{|n|}-1)$ and $\omega^n$ is a unit.

\emph{Step 4: $X$ and $\upsilon$.} By Lemma~\ref{pkl:basic}(i), $\varphi^a(\delta)=z^{\ell^a}\varphi^a(\tau)=
z^{\ell^a}\tau\prod_{b<a}p_b\cdot w_a$, so $\delta\mid\varphi^a(\delta)$. With $\omega^n-1=\delta q$ (Step 1),
$\omega'^n-1=\varphi^a(\delta)\varphi^a(q)=\delta X$, $X=(\varphi^a(\delta)/\delta)\varphi^a(q)\in\pkpA$. Apply
$\theta_a$: $\theta_a(\delta)\theta_a(X)=\theta_a(\varphi^a(\omega^n-1))=\theta(\omega^n-1)$. Also $v(\theta(\delta))
\le2a+v(\theta_a(\delta))$: for $a=0$ trivially, and for $a\ge1$ because $v(\theta(\delta))=v(\pi)+v(\zeta_\ell-1)
\le2$. By Steps 2 and 3,
\[
v(\theta_a(\delta))+v(\theta_a(X))\le v_\ell(n)+d_v+a+v(\theta(\delta))\le c+v(\theta_a(\delta)),
\]
and $\theta_a(\delta)\ne0$ (Lemma~\ref{pkl:basic}(i)); so $\theta_a(X)\mid\ell^c$ in $O_C$, say
$\ell^c=\theta_a(X)\gamma$. Choose $\upsilon_0\in\pkpA$ with $\theta(\upsilon_0)=\gamma$ and put $\upsilon=
\varphi^a(\upsilon_0)$, so $\theta_a(\upsilon)=\gamma$. Then $\theta_a(\upsilon X-\ell^c)=0$, i.e.\ $p_a\mid\upsilon X-
\ell^c$.
\end{proof}

\subsection{The depth step at the level $a$}\label{pke:ss-step}

For a lattice with jump sets $L$ (Definition~\ref{pkp:def-orders}), $\upsilon\in\pkpA$, $e\ge0$ and $t_1,\dots,t_m\in\Gamma'$ put, for $x\in\mathcal N(L)$,
\[
D_{e,\upsilon,t}(x)=\ell^ex+\upsilon\sum_{l=1}^m\delta^{-1}(t_lx-x)\ \in\pkeAN(L);
\]
it lies in $\pkeAN(L)$ by (G$_0$) of \S\ref{pkt:sec} and Lemma~\ref{pkn:M}(i). Note
$\mu_a(\delta^{-1}w)=(\delta^{(a)})^{-1}\mu_a(w)$ for $w\in\pkpAmb(L)$ ($\mu$ is $\pkpA$-linear and
$\delta^{(a)}\ne0$).

\begin{proposition}[The depth step]\label{pke:step}
Let $a\ge0$, $n\in\Z$, $c=\pkecost(a,n)$. There are $m\ge1$, $t_1,\dots,t_m\in\Gamma'$ and $\upsilon\in\pkpA$,
depending only on $(a,n)$ and the chart, such that for every lattice with jump sets $L$, every $x\in\mathcal N(L)$ and every
form $\mu$ on $L$: if $\mu_a(x)-\gamma\in\mathrm{Fil}^n$ for some $\gamma\in K'$, then
$\mu_a(D_{c,\upsilon,t}(x))-\gamma'\in\mathrm{Fil}^{n+1}$ for some $\gamma'\in K'$.
\end{proposition}

\begin{proof}
\emph{Case $n\ne0$.} Take $g,\omega',\upsilon_1,X$ from Proposition~\ref{pke:Cd}, $m=1$, $t_1=g$ and
$\upsilon=-\upsilon_1$. Put $y=\mu_a(x)-\gamma\in\mathrm{Fil}^n$; $y\in\pkeK a$ by Lemmas~\ref{pke:values},
\ref{pke:K-basic}(ii). \emph{Eigenvalue.} $w=E^{-n}y$ lies in $\mathrm{Fil}^0$ (Lemma~\ref{pkn:BdR}) and in
$\pkeK a$ ($E\in\pkeK a$). By Lemma~\ref{pke:fix} ($m=1$, $t_1=1$; the hypothesis is the second property in
Proposition~\ref{pke:Cd}), $gw-w\in\mathrm{Fil}^1$. Since $E=p_a^{(a)}$, $gE=(gp_a)^{(a)}=\omega'^{(a)}E$, so
$gy=(\omega'^{(a)})^nE^ng(w)$ and
\[
gy-(\omega'^n)^{(a)}y=(\omega'^{(a)})^nE^n(gw-w)\in\mathrm{Fil}^{n+1}.
\]
\emph{Cancellation.} Let $w'=\delta^{-1}(gx-x)\in\pkeAN(L)$. Then $\delta^{(a)}\mu_a(w')=\mu_a(gx-x)=g\mu_a(x)-
\mu_a(x)=gy-y$ ($g$ fixes $\gamma\in K'$), and $\delta^{(a)}X^{(a)}=(\omega'^n)^{(a)}-1$. Hence
$\delta^{(a)}\bigl(\mu_a(w')-X^{(a)}y\bigr)=gy-(\omega'^n)^{(a)}y\in\mathrm{Fil}^{n+1}$, and as $\delta^{(a)}$ is a
unit of $\pknBp$, $\mu_a(w')-X^{(a)}y\in\mathrm{Fil}^{n+1}$. \emph{Conclusion.} $D(x)=\ell^cx-\upsilon_1w'$ and,
with $\gamma'=\ell^c\gamma\in K'$,
\[
\mu_a(D(x))-\gamma'=\ell^cy-\upsilon_1^{(a)}\mu_a(w')=-\upsilon_1^{(a)}\bigl(\mu_a(w')-X^{(a)}y\bigr)-y\,
(\upsilon_1X-\ell^c)^{(a)} .
\]
The first term lies in $\mathrm{Fil}^{n+1}$ ($\upsilon_1^{(a)}\in\mathrm{Fil}^0$); in the second, $y\in
\mathrm{Fil}^n$ and $(\upsilon_1X-\ell^c)^{(a)}\in(p_a\pkpA)^{(a)}\subseteq\mathrm{Fil}^1$.

\emph{Case $n=0$.} Let $N=\ell^c$ ($c=d_v+3a\ge a$), $m=N$, and $t_{l+1}$ ($0\le l<N$) the elements $t_l$ of Lemma~\ref{pke:avg}, and
$\upsilon=\delta$. Then $D(x)=\ell^cx+\sum_l(t_lx-x)=\sum_lt_lx$. The hypothesis gives $\mu_a(x)\in\mathrm{Fil}^0$
($\gamma\in K'\subseteq\mathrm{Fil}^0$), and $\mu_a(x)\in\pkeK a$ (Lemma~\ref{pke:values}). As
$\mu_a(D(x))=\sum_lt_l\mu_a(x)$, Lemmas~\ref{pke:avg} and \ref{pke:fix} give $g\mu_a(D(x))-\mu_a(D(x))\in
\mathrm{Fil}^1$ for every $g\in\Gamma'$, and Lemma~\ref{pkm:BdR} gives $\gamma'\in K'$ with
$\mu_a(D(x))-\gamma'\in\mathrm{Fil}^1$.
\end{proof}

\subsection{Operators of a cost, and the proof of (I2)}\label{pke:ss-ops}

Fix a side $S$, with extension $X_S$ and projection $\mathrm{pr}$; $\mathcal N=\mathcal N(X_S)$, $L_i(o)=
(X_S)_i(o)$, and $L_i^0(o)$ the elements of $L_i(o)$ over $0$. A \emph{constant at $i$} is an element
$e_i(1\otimes c)$, $c\in C_S$; these form a group stable under multiplication by integers, and $\Gamma'$ fixes
them (Definition~\ref{pkp:def-letters}, $\Gamma$ acting trivially on $C_S$). For a letter $i$ and orders $o$ let
$\pkeDp_i(o)$ be the set of $x\in\mathcal N$ that lie over a constant at $i$ and are deep for $o$ at $i$. An
additive map $A\colon\mathcal N\to\mathcal N$ is an \emph{operator of cost $e$} if for every letter $i$, all
orders $o$ and every $x\in\pkeDp_i(o)$: $Ax-\ell^ex\in L_i^0(o)$.

\begin{lemma}\label{pke:op-basic}
(i) An operator $A$ of cost $e$ maps $\pkeDp_i(o)$ into itself, and $\mathrm{pr}(Ax)=\ell^e\mathrm{pr}(x)$
there. (ii) If $A$, $B$ are operators of costs $e$, $e'$, then $B\circ A$ is an operator of cost $e+e'$.
(iii) The identity is an operator of cost $0$.
\end{lemma}

\begin{proof}
(i) Let $x\in\pkeDp_i(o)$ and $w=Ax-\ell^ex\in L_i^0(o)$. Then $\mathrm{pr}(Ax)=\ell^e\mathrm{pr}(x)$, a constant
at $i$, and $g(Ax)-Ax=\ell^e(gx-x)+(gw-w)\in L_i(o)$ for $g\in\Gamma'$, as $L_i(o)$ is a $\Gamma'$-stable
$\pkpA$-module (Lemma~\ref{pkl:G1}). (ii) $BAx-\ell^{e+e'}x=(B(Ax)-\ell^{e'}Ax)+\ell^{e'}(Ax-\ell^ex)$; the
first term is in $L_i^0(o)$ by (i), the second by hypothesis. (iii) is clear.
\end{proof}

\begin{lemma}[One step]\label{pke:B}
For every $a\ge0$ and $n\in\Z$ there is an operator $B$ of cost $\pkecost(a,n)$ with the following property:
for $x\in\mathcal N$ and a form $\mu$ on $X_S$, if $\mu_a(gx-x)\in\mathrm{Fil}^n$ for all $g\in\Gamma'$, then
$\mu_a(gBx-Bx)\in\mathrm{Fil}^{n+1}$ for all $g\in\Gamma'$.
\end{lemma}

\begin{proof}
Let $c,m,t,\upsilon$ be as in Proposition~\ref{pke:step} and put $Bx=\mathrm{ret}\bigl(D_{c,\upsilon,t}(x)\bigr)$;
$B$ is additive with values in $\mathcal N$. Since $\mathrm{ret}$ is $\mathfrak S$-linear and the identity on
$\mathcal N$ (Proposition~\ref{pkl:R}), $Bx-\ell^cx=\sum_l\mathrm{ret}\bigl(\upsilon\delta^{-1}(t_lx-x)\bigr)$.

\emph{Cost.} Let $x\in\pkeDp_i(o)$. By Lemma~\ref{pkm:inputs}(G) (second assertion; Lemma~\ref{pkt:tau-iv}),
$\delta^{-1}(t_lx-x)\in L_i(o)$, hence $\upsilon\delta^{-1}(t_lx-x)\in L_i(o)$, and by Lemma~\ref{pke:ret} its
retraction lies in $\mathcal N\cap L_i(o)$. It lies over $0$: $\mathrm{pr}$ is $\pkpW$-linear,
$\Gamma'$-equivariant and commutes with $\mathrm{ret}$ (Proposition~\ref{pkl:R}), and $t_l$ fixes the constant
$\mathrm{pr}(x)$. So $Bx-\ell^cx\in L_i^0(o)$.

\emph{Raising.} Let $x,\mu$ be as in the statement. By Lemma~\ref{pkm:equiv}, $g\mu_a(x)-\mu_a(x)=\mu_a(gx-x)\in
\mathrm{Fil}^n$ for all $g$, so by Lemma~\ref{pkm:BdR} $\mu_a(x)-\gamma\in\mathrm{Fil}^n$ for some $\gamma\in K'$.
By Proposition~\ref{pke:step}, $\mu_a(D_{c,\upsilon,t}(x))-\gamma'\in\mathrm{Fil}^{n+1}$ with $\gamma'\in K'\subseteq
\pkeK a$; by Lemma~\ref{pke:ret} ($W=\mathrm{Fil}^{n+1}$), $\mu_a(Bx)-\gamma'\in\mathrm{Fil}^{n+1}$. As $g$ fixes
$\gamma'$ and preserves $\mathrm{Fil}^{n+1}$, $\mu_a(gBx-Bx)=g(\mu_a(Bx)-\gamma')-(\mu_a(Bx)-\gamma')\in
\mathrm{Fil}^{n+1}$ (Lemma~\ref{pkm:equiv}).
\end{proof}

For orders $o,o'$ and a letter $i$ say that a set $I$ of pairs $(a,n)\in\Z_{\ge0}\times\Z$ \emph{covers
$o\to o'$ at $i$} if it contains every $(a,n)$ with $o(a,u)\le n<o'(a,u)$ for some $u\in U(i^{(-a)})$.

\begin{proposition}[Deepening operators]\label{pke:A}
For every finite set $I$ of pairs $(a,n)$ there is an operator $A_I$ of cost $\sum_{(a,n)\in I}\pkecost(a,n)$
such that $A_Ix\in\pkeDp_i(o')$ whenever $x\in\pkeDp_i(o)$ and $I$ covers $o\to o'$ at $i$.
\end{proposition}

\begin{proof}
Induction on $|I|$. If $I=\emptyset$, take $A_I=\mathrm{id}$ (Lemma~\ref{pke:op-basic}(iii)); if $\emptyset$
covers $o\to o'$, then $o'(a,u)\le o(a,u)$ on the jumps, and $\pkeDp_i(o)\subseteq\pkeDp_i(o')$
(Lemma~\ref{pkm:le}(i)). Let $I\ne\emptyset$, choose $(a,n)\in I$ with $n$ maximal among the second coordinates
of the elements of $I$, put $I'=I\smallsetminus\{(a,n)\}$, let $A_{I'}$ be given by induction and $B$ by
Lemma~\ref{pke:B} for $(a,n)$, and $A_I=B\circ A_{I'}$, an operator of the right cost by
Lemma~\ref{pke:op-basic}(ii).

Let $x\in\pkeDp_i(o)$ with $I$ covering $o\to o'$. Define $o_-(b,u)=n$ if $b=a$ and $o'(b,u)=n+1$, and
$o_-(b,u)=o'(b,u)$ otherwise. \emph{$I'$ covers $o\to o_-$.} Let $o(b,u)\le k<o_-(b,u)$ with $u\in
U(i^{(-b)})$. Since $o_-\le o'$, $(b,k)\in I$. If $(b,k)=(a,n)$, then $o_-(a,u)>n$ forces $o'(a,u)\ne n+1$, and
$o'(a,u)\ge o_-(a,u)>n$, so $o'(a,u)\ge n+2$; then $o(a,u)\le n+1<o'(a,u)$ and $(a,n+1)\in I$, contradicting the
maximality of $n$. So $(b,k)\in I'$. Hence $y=A_{I'}x\in\pkeDp_i(o_-)$, and $By\in\pkeDp_i(o_-)$
(Lemma~\ref{pke:op-basic}(i)). \emph{Raising.} Let $g\in\Gamma'$; $gBy-By\in L_i(o_-)$, and it remains to see
$\mu_b(gBy-By)\in\mathrm{Fil}^{o'(b,u)}$ where $o'(b,u)>o_-(b,u)$, i.e.\ $b=a$, $o'(a,u)=n+1$, for $\mu$ of type
$u\in U(i^{(-a)})$ at $i^{(-a)}$. Since $y$ is deep for $o_-$ and $o_-(a,u)=n$, $\mu_a(\gamma y-y)\in\mathrm{Fil}^n$
for all $\gamma\in\Gamma'$, and Lemma~\ref{pke:B} gives $\mu_a(gBy-By)\in\mathrm{Fil}^{n+1}$.
\end{proof}

\begin{theorem}[Input (I2)]\label{pke:I2}
Let $x\in\pkcPt_S(c,i,e;o,d)$ and let $I$ be a finite set of pairs $(a,n)\in\Z_{\ge0}\times\Z$ containing every
$(a,n)$ with $d(a,u)\le n<d'(a,u)$ for some $u\in U(i^{(-a)})$. If $e+\sum_{(a,n)\in I}\pkecost(a,n)\le e'$,
then $\pkcPt_S(c,i,e';\pkcimp(o,d'),d')\ne\emptyset$.
\end{theorem}

\begin{proof}
Let $\Sigma=\sum_I\pkecost(a,n)$, $s=e'-e-\Sigma\ge0$, $y=\ell^sx$ and $A=A_I$ (Proposition~\ref{pke:A}). Then
$y\in\mathcal N$, $\mathrm{pr}(y)=\ell^{e+s}e_i(1\otimes c)$, $y\in L_i(o)$, and $y$ is deep for $d$ and, as
$L_i(o)$ is $\Gamma'$-stable, for $o$; so $y\in\pkeDp_i(d)\cap\pkeDp_i(o)$. By Proposition~\ref{pke:A}, $Ay\in
\pkeDp_i(d')$. As $A$ has cost $\Sigma$ and $y\in\pkeDp_i(o)$, $Ay-\ell^\Sigma y\in L_i(o)$, so $Ay\in L_i(o)$;
and $\mathrm{pr}(Ay)=\ell^\Sigma\mathrm{pr}(y)=\ell^{e'}e_i(1\otimes c)$ (Lemma~\ref{pke:op-basic}(i)). By
Lemma~\ref{pkm:improved}, $Ay\in L_i(\pkcimp(o,d'))$. So $Ay\in\pkcPt_S(c,i,e';\pkcimp(o,d'),d')$.
\end{proof}

\begin{remark}\label{pke:rem-use}
The operators depend only on the chart and $I$ (through the elements $t_l$, $\upsilon$ of
Proposition~\ref{pke:step}), not on $c$, $i$, $o$; this is not needed for (I2). The statement of (I2) in
\S\ref{pkc:ss-inputs} is proved as stated. Lemma~\ref{pkc:free} follows from it as there.
\end{remark}

\subsection{The other inputs}\label{pke:ss-open}

The input (I4) of \S\ref{pkc:ss-inputs} is Theorem~\ref{pkh:I4}; the input (I5) of \S\ref{pkd:ss-tools}
is Theorem~\ref{pkg:I5}; Proposition~\ref{pki:chart} is Theorem~\ref{pkj:chart}.


%% ===== proofs/C-kisin-gap.tex =====
% proofs/C-kisin-gap.tex
% The long gap (I5): constants modulo E^A, the radius lemma for frames, expansion in powers of E.
% Labels start with pkg:. Conventions and macros as in proofs/C-kisin-deep.tex.
\providecommand{\auO}{\mathcal O}
\providecommand{\auQlb}{\overline{\Q}_\ell}
\providecommand{\panot}[1]{\par\smallskip{\raggedright\noindent\textbf{[Not proved here:} #1\textbf{]}\par}\smallskip}
\providecommand{\pkpA}{A_{\mathrm{inf}}}
\providecommand{\pkpW}{W(C^\flat)}
\providecommand{\pkpAmb}{\mathbb A}
\providecommand{\pknBp}{B^+_{\mathrm{dR}}}
\providecommand{\pknB}{B_{\mathrm{dR}}}
\providecommand{\pknMinf}{\mathfrak M_{\mathrm{inf}}}
\providecommand{\pknK}{\mathcal K_0}
\providecommand{\pkfAm}{A_{\max}}
\providecommand{\pkfAo}{A_{\max}^{\circ}}
\providecommand{\pkfBp}{B^+}
\providecommand{\pkfSB}{\mathfrak S_B}
\providecommand{\pkcPt}{\mathrm{Pt}}
% macros of this file (prefix \pkg)
\providecommand{\pkgAN}{\pkpA\mathcal N}
\providecommand{\pkgP}{\mathcal P}
\providecommand{\pkgC}{\mathcal C}
\providecommand{\pkgV}{\mathcal V}
\providecommand{\pkgSh}{\widehat{\mathfrak S}}
\providecommand{\pkgstd}{\mathrm{std}}

\section{The long gap: the input (I5)}\label{pkg:sec}

This section proves the input (I5) of \S\ref{pkd:ss-tools} (long gap), in the form printed in
\S\ref{pkd:ss-tools}.

\emph{Dependencies.} Theorem~\ref{pkg:I5} ((I5)) is proved in full, modulo results cited by label from
\S\S\ref{pkp:sec}--\ref{pke:sec}: the frame theorem with its Galois relation
(Theorems~\ref{pkf:main}, \ref{pkf:frame-pm}, Lemma~\ref{pkn:frame-id}, Proposition~\ref{pkn:frame}),
the facts on $\pkfAm$, $\pkfBp$, $\pkfSB$ of \S\S\ref{pkf:sec}--\ref{pkr:sec} and \S\ref{pkl:ss-open},
item (I1) of \S\ref{pkc:ss-inputs} (Proposition~\ref{pkt:I1}), the retraction (Proposition~\ref{pkl:R}),
improved orders (Lemma~\ref{pkm:improved}) and the lift ideals (Definition~\ref{pkm:def-J},
Lemma~\ref{pkm:J-basic}). The input (I4) and Proposition~\ref{pki:chart} are proved elsewhere
(\S\ref{pkg:ss-open}).

\emph{Conventions.} Those of \S\ref{pke:sec}. Throughout, the chart is \emph{unramified}: $e=1$ and
$E=z+\ell$ (Definition~\ref{pkp:def-chart}), so $p_a=\varphi^a(E)=z^{\ell^a}+\ell=(E-\ell)^{\ell^a}+\ell$.
For orders $o$ and a function $c\colon\Z\to\Z$ we write $o+c[a=0]$ for the orders
$(a,u)\mapsto o(a,u)+c(u)[a=0]$. For a lattice with jump sets $L$, a letter $i$ and $l\ge0$ put
\[
\pkgP_l(i)=\mathcal N(L)\cap L_i(\pkgstd_l),\qquad
\pkgC_{l,A}(i)=\{\xi\in\pkgP_l(i):\ g\xi-\xi\in L_i(\pkgstd_l+A[a=0])\ \text{for all }g\in\Gamma'\}
\]
($A\ge0$), the \emph{constants modulo $E^A$}. Both are $\Lambda$-modules, $\pkgP_l(i)$ an
$\mathfrak S$-module: $\mathcal N(L)$ is an $\mathfrak S$-module, $L_i(o)$ an $\pkpA$-module, and for
$w\in\Lambda$, $g(w\xi)-w\xi=w(g\xi-\xi)$ since $\Gamma'$ fixes $\Lambda$. The trivial lattice $C_S$ of a
side is given the jump set $\{0\}$ at every letter; since $\pkgstd_l(a,0)=0=\pkgstd_0(a,0)$,
$(C_S)_i(\pkgstd_l)=(C_S)_i(\pkgstd_0)$ for all $l$.

\subsection{Constants modulo $E^A$}\label{pkg:ss-const}

\begin{lemma}\label{pkg:const}
Let $l\ge1$ and $A\ge0$.
\begin{enumerate}
\item[(i)] For a side $S$, $\mathrm{pr}$ maps $\mathcal N(X_S)$, $(X_S)_i(o)$, $\pkgP_l(i)$, $\pkgC_{l,A}(i)$
into the corresponding objects of $C_S$; and $e_i(1\otimes c)\in\pkgC_{l,A}(i)$ for $C_S$, $c\in C_S$.
\item[(ii)] Let $x\in\mathcal N(L)$ with $\varphi^{-l}x\in\pknMinf(L)$, and suppose $g\mu(x)-\mu(x)\in
\mathrm{Fil}^A$ for every form $\mu$ and every $g\in\Gamma'$. Then $e_ix\in\pkgC_{l,A}(i)$.
\end{enumerate}
\end{lemma}

\begin{proof}
(i) By the last paragraph of the proof of Proposition~\ref{pkl:R}, $\mathrm{pr}$ maps $\pkgAN(X_S)$ and
$\mathcal N(X_S)$ into $\pkgAN(C_S)$, $\mathcal N(C_S)$, commutes with $\varphi$, and $\nu\circ\mathrm{pr}$ is
a form on $X_S$ for every form $\nu$ on $C_S$; it commutes with $1\otimes\zeta$, hence with the $e_i$, and
with $\Gamma$. A form $\nu$ of type $0$ at $j$ on $C_S$ gives $\nu\circ\mathrm{pr}$ of type $0$ at $j$ on
$X_S$, and $0\in U(j)$ for the side (Definition~\ref{pkp:def-side}); so $x\in(X_S)_i(o)$ gives
$\nu(\varphi^{-a}\mathrm{pr}\,x)=(\nu\circ\mathrm{pr})(\varphi^{-a}x)\in\mathrm{Fil}^{o(a,0)}$, i.e.\
$\mathrm{pr}\,x\in(C_S)_i(o)$. The assertions on $\pkgP$ and $\pkgC$ follow, as $\mathrm{pr}(gx-x)=
g\,\mathrm{pr}(x)-\mathrm{pr}(x)$. Finally $e_i(1\otimes c)\in\mathcal N(C_S)\cap(C_S)_i(\pkgstd_l)$ by
Proposition~\ref{pkl:H2}(ii), and it is fixed by $\Gamma'$ ($\Gamma$ acts trivially on $C_S$ and fixes
$\Lambda$).

(ii) $e_ix\in\mathcal N(L)$ (Lemma~\ref{pkl:basic}(iv)). As $\varphi^{-l}e_i=e_{i^{(-l)}}\varphi^{-l}$
(Definition~\ref{pkp:def-letters}), $\varphi^{-l}(e_ix)=e_{i^{(-l)}}\varphi^{-l}x\in\pknMinf(L)\cap
\pkpAmb(L)_{i^{(-l)}}$ (Lemma~\ref{pkl:basic}(iv)), so $e_ix\in L_i(\pkgstd_l)$ by
Proposition~\ref{pkl:std}. Let $g\in\Gamma'$; $g(e_ix)-e_ix\in L_i(\pkgstd_l)$ by Lemma~\ref{pkl:G1}. As
$l\ge1$, $\pkgstd_l(0,u)+A=A$, and it remains to see $\mu(g e_ix-e_ix)\in\mathrm{Fil}^A$ for a form $\mu$ of
type $u$ at $i$. The map $\mu'=\mu\circ e_i$ is $\pkpA$-linear, and a form: $e_i$ is a polynomial in
$1\otimes\zeta$ with coefficients in $\Lambda$, and $\zeta\in\auO$ commutes with $\Gamma$. Since $g$
commutes with $e_i$ and $x\in\pkgAN(L)$, Lemma~\ref{pkm:equiv} gives $\mu(ge_ix-e_ix)=\mu'(gx)-\mu'(x)=
g\mu'(x)-\mu'(x)\in\mathrm{Fil}^A$.
\end{proof}

\begin{lemma}[Peeling an expansion]\label{pkg:peel}
Let $l\ge1$, $A,N\ge0$, $\xi_0,\dots,\xi_N\in\pkgC_{l,A}(i)$, $y\in L_i(\pkgstd_l)$, and $o\colon\Z\to\Z$
with $o\le A$. If $\sum_{n\le N}E^n\xi_n+E^{N+1}y\in L_i(\pkgstd_l+o[a=0])$, then
$\xi_N\in L_i(\pkgstd_l+A[o(u)>N][a=0])$.
\end{lemma}

\begin{proof}
\emph{Case $N=0$.} By Lemma~\ref{pkn:E}, $Ey\in L_i(\pkgstd_l+1[a=0])$, so $\xi_0\in L_i(o_1)$ with
$o_1=\pkgstd_l+\min(1,o)[a=0]$ (Lemma~\ref{pkm:le}(i)). Since $\xi_0$ is a constant modulo $E^A$,
Lemma~\ref{pkm:improved} gives $\xi_0\in L_i(\mathrm{imp}(o_1,\pkgstd_l+A[a=0]))$. At $a=0$ (where
$\pkgstd_l=0$ as $l\ge1$) this order is $A$ if $o(u)\ge1$ and $\min(A,0)=0$ otherwise; at $a\ge1$ both
orders equal $\pkgstd_l(a,u)$, and $\mathrm{imp}$ is $\pkgstd_l(a,u)$ in both cases. This is the claim.

\emph{Induction.} Let $N\ge1$ and write the sum as $\xi_0+Ey'$ with $y'=\sum_{n<N}E^n\xi_{n+1}+E^Ny$, which
lies in $L_i(\pkgstd_l)$ (Lemma~\ref{pkn:E}, Lemma~\ref{pkm:le}(i)). By the case $N=0$, $\xi_0\in
L_i(\pkgstd_l+A[o>0][a=0])\subseteq L_i(\pkgstd_l+o[a=0])$, as $A[o>0]\ge o$ ($o\le A$). Hence
$Ey'\in L_i(\pkgstd_l+o[a=0])$ and $y'=E^{-1}(Ey')\in L_i(\pkgstd_l+(o-1)[a=0])$ (Lemma~\ref{pkn:E}). By
induction applied to $\xi_1,\dots,\xi_N$, $y$ and $o-1\le A$, $\xi_N\in L_i(\pkgstd_l+A[o(u)-1>N-1][a=0])$.
\end{proof}

\begin{lemma}[Digits in powers of $E$]\label{pkg:digits}
For $s\in\mathfrak S$ there are $w_n\in\Lambda$ ($n\ge0$) such that for every $N$,
$s-\sum_{n<N}E^nw_n\in E^N\mathfrak S$.
\end{lemma}

\begin{proof}
For $r=\sum_k\lambda_kz^k\in\mathfrak S$ put $r(-\ell)=\sum_k\lambda_k(-\ell)^k\in\Lambda$ ($\ell$-adically
convergent). Then $r-r(-\ell)=\sum_{k\ge1}\lambda_k(z^k-(-\ell)^k)=(z+\ell)q$ with
$q=\sum_{j\ge0}z^j\sum_{k>j}\lambda_k(-\ell)^{k-1-j}\in\mathfrak S$ (inner sums converge). So
$r=r(-\ell)+Eq(r)$. Put $r_0=s$, $w_n=r_n(-\ell)$, $r_{n+1}=q(r_n)$; by induction
$s=\sum_{n<N}E^nw_n+E^Nr_N$.
\end{proof}

\begin{lemma}[The polynomial $f$]\label{pkg:f}
Let $s,g,l\ge0$ and $N=s+g(1+\ell+\dots+\ell^{l-1})$. There is a monic $f\in\Z[X]$ of degree $N$ with
$f(E)=E^s\prod_{a<l}p_a^g$. If $y\in L_i(\rho_{-s})$, then $f(E)y\in L_i(o)$ with $o(a,u)=g[a<l]-u$.
\end{lemma}

\begin{proof}
$f=X^s\prod_{a<l}((X-\ell)^{\ell^a}+\ell)^g$, monic of degree $N$. By Lemma~\ref{pkn:E}, multiplication by
$E^s$ adds $s$ at $a=0$, and by $p_a^g$ adds $g$ at the level $a$; starting from $\rho_{-s}=-u-s[a=0]$
one gets $o$.
\end{proof}

\subsection{The radius lemma for frames}\label{pkg:ss-radius}

Put $T=z/\ell=E/\ell-1\in\pkfAm$ (image in $\pkfBp$; $\pkfAm$ is $\ell$-torsion free, so we identify
$\pkfAm$ with $\pkfAo$). The subring $\mathfrak S[T]\subseteq\pkfAm$ lies in $\mathfrak S[1/\ell]\subseteq
\pkfBp$, and $\mathfrak S[1/\ell]=\Lambda[[z]][1/\ell]\subseteq K_0[[z]]$; inside $K_0[[z]]$ also
$\Lambda[[T]]\subseteq K_0[[z]]$ ($T=z/\ell$) and $\mathfrak S[T]\subseteq\Lambda[[T]]$, as
$\mathfrak S=\Lambda[[\ell T]]$. Let $\pkgSh$ be the $\ell$-adic closure of $\mathfrak S[T]$ in $\pkfAm$.

\begin{lemma}\label{pkg:Shat}
(i) $\mathfrak S[T]\cap\ell^n\pkfAm=\ell^n\mathfrak S[T]$ and $\pkgSh\cap\ell^n\pkfAm=\ell^n\pkgSh$; and
$\pkgSh\subseteq\pkfSB$, $\pkfSB=\pkgSh[1/\ell]$.
(ii) The inclusion $\mathfrak S[T]\subseteq\Lambda[[T]]$ extends to an injective ring map
$\psi\colon\pkgSh\to\Lambda[[T]]$.
(iii) If $y\in\pkfSB$ and $z^N$ divides $y$ in $\pkfSB$ for every $N$, then $y=0$.
\end{lemma}

\begin{proof}
(i) \emph{Claim: if $P\in\mathfrak S[T]$ and $P\in\ell^n\Lambda[[T]]$, then $P\in\ell^n\mathfrak S[T]$.}
Let $P=\sum_{k\le D}s_kT^k$, $s_k\in\Lambda[[\ell T]]$; the coefficient of $T^m$ in $P$ is divisible by
$\ell^{m-D}$ for $m>D$. So $Q=\ell^{-n}P\in\Lambda[[T]]$ has coefficients $q_m$ with $\ell^{m-D-n}\mid q_m$ for
$m>D+n$, and $Q=\sum_{m\le D+n}q_mT^m+T^{D+n}\sum_{j\ge1}(\ell^{-j}q_{D+n+j})(\ell T)^j\in\mathfrak S[T]$.
Now let $P\in\mathfrak S[T]\cap\ell^n\pkfAm$ of degree $\le D$ in $T$. Then $\ell^DP\in\mathfrak S\cap
\ell^{n+D}\pkfAm$, which by Proposition~\ref{pkr:I1} is contained in $\mathfrak S\cap(\ell,E)^{n+D}\pkpA$;
applying $\kappa_0$ (Proposition~\ref{pkr:I5}; $\mathfrak S$-linear, identity on $\mathfrak S$, and
$\ell,E\in\mathfrak S$) to an expression $\sum_j\ell^jE^{n+D-j}a_j$ shows $\ell^DP\in(\ell,E)^{n+D}\mathfrak S=
(\ell,z)^{n+D}\mathfrak S$ (the two ideals coincide because $E=z+\ell$), and $(\ell,z)^{n+D}\mathfrak S
\subseteq\ell^{n+D}\Lambda[[T]]$ because $z=\ell T$. So $P\in\ell^n\Lambda[[T]]$ and the claim applies. For
$y\in\pkgSh\cap\ell^n\pkfAm$ choose $P_k\in\mathfrak S[T]$ with $y-P_k\in\ell^{n+k}\pkfAm$; then
$P_k\in\ell^n\mathfrak S[T]$, the $\ell^{-n}P_k$ converge, and their limit $\ell^{-n}y$ lies in $\pkgSh$
($\pkfAm$ is complete, separated and $\ell$-torsion free). By Proposition~\ref{pkf:closure} (with $m=0$;
$\mathfrak S[T]=\mathfrak S[E/\ell]$), $\pkgSh\subseteq\pkfSB$, and every $y\in\pkfSB$ has $\ell^my\in\pkgSh$
for some $m$; with $\ell^{-1}\in\pkfSB$ (Lemma~\ref{pkf:SB-basic}(i)), $\pkfSB=\pkgSh[1/\ell]$.

(ii) If $P_k\to y$ in $\pkgSh$, then $P_k-P_{k'}\in\mathfrak S[T]\cap\ell^m\pkfAm=\ell^m\mathfrak S[T]$ for
large $k,k'$, so $(P_k)$ is $\ell$-adically Cauchy in $\Lambda[[T]]$; put $\psi(y)=\lim P_k$, which is well
defined (apply the same to the alternated sequence) and a ring map. Injectivity: let $\psi(y)=0$. Choose
$P\in\mathfrak S[T]$ with $y-P\in\ell\pkfAm\cap\pkgSh=\ell\pkgSh$; then $\psi(P)\in\ell\Lambda[[T]]$, so
$P\in\ell\mathfrak S[T]$ by the claim, and $y=\ell y_1$ with $y_1\in\pkgSh$, $\psi(y_1)=0$ ($\Lambda[[T]]$ is
torsion free). Hence $y\in\bigcap_n\ell^n\pkfAm=0$.

(iii) Extend $\psi$ to $\pkfSB=\pkgSh[1/\ell]\to\Lambda[[T]][1/\ell]$, still injective. If $y=z^Nc_N$
with $c_N\in\pkfSB$, then $\psi(y)=\ell^NT^N\psi(c_N)\in T^N\Lambda[[T]][1/\ell]$ for all $N$, so $\psi(y)=0$
and $y=0$.
\end{proof}

Fix $l\ge1$ with $\ell^l\ge3$ and put $X_l=z^{\ell^l}/\ell^2=\ell^{\ell^l-2}T^{\ell^l}\in\ell\pkgSh$. As
$\pkgSh$ is closed in the complete ring $\pkfAm$, the $\mathfrak S$-algebra map
$\mathrm{ev}\colon\mathfrak S[[X]]\to\pkgSh$, $\sum_kf_kX^k\mapsto\sum_kf_kX_l^k$, is well defined.

\begin{lemma}[Approximation]\label{pkg:approx}
Let $A,\delta\ge0$ with $A\le\ell^l\delta$ and $F\in\mathfrak S[[X]]$. There is $s\in\mathfrak S$ with
$s-\ell^{2\delta}\mathrm{ev}(F)\in E^A\pkfSB$.
\end{lemma}

\begin{proof}
Write $F=\sum_{k<\delta}f_kX^k+X^\delta F'$. Then $\ell^{2\delta}\mathrm{ev}(F)=s_0+z^{\ell^l\delta}\mathrm{ev}(F')$
with $s_0=\sum_{k<\delta}f_k\ell^{2\delta-2k}z^{\ell^lk}\in\mathfrak S$, and
$z^{\ell^l\delta}\mathrm{ev}(F')=\ell^Ab$ with $b=T^Az^{\ell^l\delta-A}\mathrm{ev}(F')\in\pkgSh\subseteq\pkfSB$.
By Lemma~\ref{pkf:approx0}(i), $b=\sum_k\alpha_k(E/\ell)^k$ with $\alpha_k\in\pkpA$, $\alpha_k\to0$, so
$\ell^Ab=\alpha+E^Aw$ with $\alpha=\sum_{k<A}\alpha_k\ell^{A-k}E^k\in\pkpA$ and $w\in\pkfAm$. Let $G_0$ be the
extension of $\kappa_0$ (Lemma~\ref{pkf:extend}); it is $\pkfSB$-linear, the identity on $\pkfSB$ and has
values in $\pkfSB$ (Lemma~\ref{pkf:values}). So $\ell^Ab=G_0(\ell^Ab)=\kappa_0(\alpha)+E^AG_0(w)$, and
$s=s_0+\kappa_0(\alpha)$ works.
\end{proof}

Now fix a lattice with jump sets $L$ with jump bound $a$, and the frame of Theorem~\ref{pkf:main}: by its
proof and Theorem~\ref{pkf:frame-pm}, $Z\in M_d(\pkfSB)$, $Y\in M_d(\pkfBp)$ with $YZ=E^h$ ($h=2a$), and
$g(Z)=ZG_g$ with $G_g\in M_d(\pkpA)$ for every $g\in\Gamma'$ (such $g$ fix $\Lambda$ and $\theta(z)$). Fix
$g_1\in\Gamma'$ with $g_1z=[\varepsilon]z$ (Definition~\ref{pkp:def-chart}) and write $G=G_{g_1}$.

\begin{lemma}\label{pkg:BW}
(i) There are $B\in M_d(K_0)$, $K_0=\Lambda[1/\ell]$, and $W\in M_d(\pkfSB)$ with $BZ=1+zW$.
(ii) $z\tau$ divides every entry of $G-1$ in $\pkpA$; and $g_1(BZ)=BZG$.
(iii) If $r\in\pkpA$, $n\ge0$ and $\tau z^{n+1}$ divides $r$ in $\pkfBp$, then it divides $r$ in $\pkpA$.
\end{lemma}

\begin{proof}
(i) Let $Y_0=G_0(Y)$ entrywise; $Y_0\in M_d(\pkfSB)$ and, $G_0$ being additive and $\pkfSB$-linear,
$Y_0Z=G_0(YZ)=E^h$. Write $E^h=(z+\ell)^h=\ell^h+zq$ with $q\in\mathfrak S$. By Lemma~\ref{pkl:SBdec} each
entry $b$ of $\ell^{-h}Y_0$ has $\ell^kb=\lambda+zw'$ ($\lambda\in\Lambda$, $w'\in\pkfSB$); let $B$, $W_1$ be the
matrices of the $\ell^{-k}\lambda$, $\ell^{-k}w'$. Then $BZ=\ell^{-h}Y_0Z-zW_1Z=1+z(\ell^{-h}q-W_1Z)$.
(ii) $E^h(G-1)=YZG-YZ=Y(g_1Z-Z)$. By Lemma~\ref{pkl:galSB}, $z\tau$ divides the entries of $g_1Z-Z$ in
$\pkfBp$, hence those of $E^h(G-1)$, and Lemma~\ref{pkl:cancel} gives the first claim. $g_1$ fixes $K_0$, so
$g_1(BZ)=Bg_1(Z)=BZG$.
(iii) By Lemma~\ref{pkl:cancel} ($h=0$), $r=z\tau r'$ with $r'\in\pkpA$; cancelling $z\tau$ in the domain
$\pkfBp$, $z^n$ divides $r'$ in $\pkfBp$, and $n$ applications of Lemma~\ref{pkl:intdiv} (each followed by a
cancellation of $z$) give $z^n\mid r'$ in $\pkpA$.
\end{proof}

Put $Y_1=BZ=1+zW$. A \emph{datum} for a row index $i$ is a triple $(n,R,q)$ with $n\ge0$, $R\in\pkfSB^d$ and
$q\in\Z_{>0}$ such that $q(Y_1)_{ij}-z^{n+1}R_j\in\pkpA$ for all $j$.

\begin{lemma}[One step]\label{pkg:step}
Let $(n,R,q)$ be a datum and $v=v_\ell(n+1)$. There are $s\in\mathfrak S^d$ and $R'\in\pkfSB^d$ with
$\ell^{v+1}R=s+zR'$, and $(n+1,R',\ell^{v+1}q)$ is a datum.
\end{lemma}

\begin{proof}
Let $\alpha_j=q(Y_1)_{ij}-z^{n+1}R_j\in\pkpA$. By Lemma~\ref{pkg:BW}(ii), $\sum_k(Y_1)_{ik}G_{kj}=
g_1((Y_1)_{ij})$, and $g_1(z^{n+1})=[\varepsilon]^{n+1}z^{n+1}$; hence
\[
\textstyle\sum_k\alpha_kG_{kj}-g_1(\alpha_j)=z^{n+1}\Bigl(([\varepsilon]^{n+1}-1)R_j+[\varepsilon]^{n+1}
(g_1R_j-R_j)-\sum_kR_k(G-1)_{kj}\Bigr)=z^{n+1}\bigl(([\varepsilon]^{n+1}-1)R_j-z\tau D_j\bigr)
\]
with $D_j\in\pkfBp$, by Lemma~\ref{pkl:galSB} and Lemma~\ref{pkg:BW}(ii). By Proposition~\ref{pkn:Cc} there are
$\alpha',\beta\in\pkpA$ with $\tau E(\ell^v+\beta z)=\alpha'([\varepsilon]^{n+1}-1)$, and
$E(\ell^v+\beta z)=\ell^{v+1}+z(\ell^v+\beta E)$ as $E=\ell+z$. Multiplying by $\alpha'$,
\[
\textstyle\alpha'\bigl(\sum_k\alpha_kG_{kj}-g_1\alpha_j\bigr)=\tau z^{n+1}\bigl(\ell^{v+1}R_j-zR''_j\bigr),\qquad
R''_j=\alpha'D_j-(\ell^v+\beta E)R_j\in\pkfBp .
\]
The left side lies in $\pkpA$, so by Lemma~\ref{pkg:BW}(iii) it equals $\tau z^{n+1}x_j$ with $x_j\in\pkpA$,
and $x_j=\ell^{v+1}R_j-zR''_j$ ($\pkfBp$ is a domain). Applying $G_0$:
$\kappa_0(x_j)=\ell^{v+1}R_j-zG_0(R''_j)$. Put $s_j=\kappa_0(x_j)$, $R'_j=G_0(R''_j)\in\pkfSB$. Then
$\ell^{v+1}\alpha_j+z^{n+1}s_j=\ell^{v+1}q(Y_1)_{ij}-z^{n+2}R'_j\in\pkpA$.
\end{proof}

\begin{lemma}\label{pkg:series}
There is $F\in M_d(\mathfrak S[[X]])$ with $\ell^2\varphi^l(W)=\mathrm{ev}(F)$ (entrywise).
\end{lemma}

\begin{proof}
Fix $i,j$. $(0,(W_{ik})_k,1)$ is a datum, since $(Y_1)_{ik}-zW_{ik}=\delta_{ik}$. Applying
Lemma~\ref{pkg:step} repeatedly gives data $(n,R^{(n)},q_n)$ and $s^{(n)}\in\mathfrak S^d$ with
$\ell^{v_n+1}R^{(n)}=s^{(n)}+zR^{(n+1)}$, $v_n=v_\ell(n+1)$. Put $\gamma_n=\ell^{n-v_\ell(n!)+2}\varphi^l(R^{(n)}_j)
\in\pkfSB$ (Lemma~\ref{pkf:values}) and $a_n=n+1-v_\ell((n+1)!)\ge0$ ($v_\ell(m!)<m/(\ell-1)\le m$). Apply
$\varphi^l$ ($\varphi^l(z)=z^{\ell^l}=\ell^2X_l$) and multiply by $\ell^{a_n}$; as $a_n+v_n+1=n-v_\ell(n!)+2$,
\[
\gamma_n=\ell^{a_n}\varphi^l(s^{(n)}_j)+X_l\gamma_{n+1},\qquad\text{so}\qquad
\gamma_0=\textstyle\sum_{k<N}c_kX_l^k+X_l^N\gamma_N\quad(N\ge0),
\]
with $c_k=\ell^{a_k}\varphi^l(s^{(k)}_j)\in\mathfrak S$. Let $F_{ij}=\sum_kc_kX^k$. Then
$\gamma_0-\mathrm{ev}(F_{ij})=X_l^N\bigl(\gamma_N-\mathrm{ev}(\sum_{k\ge N}c_kX^{k-N})\bigr)\in X_l^N\pkfSB\subseteq
z^N\pkfSB$ for all $N$ ($X_l=z\cdot\ell^{-2}z^{\ell^l-1}$). By Lemma~\ref{pkg:Shat}(iii),
$\gamma_0=\ell^2\varphi^l(W_{ij})$ equals $\mathrm{ev}(F_{ij})$.
\end{proof}

\begin{proposition}[Radius lemma]\label{pkg:radius}
There are $D\in GL_d(K_0)$ and $V\in GL_d(\pkgSh)$ with $\varphi^l(Z)=p_l^aDV$, such that for all $A,\delta\ge0$
with $A\le\ell^l\delta$: the set $\mathcal Y$ of $y\in\mathfrak S^d$ with $Vy-r\in E^A(\pkfBp)^d$ for some
$r\in\Z^d$ spans an $\mathfrak S$-submodule containing $\ell^{4\delta}\mathfrak S^d$.
\end{proposition}

\begin{proof}
Let $F$ be as in Lemma~\ref{pkg:series}, $\mathcal V'=1+XF\in GL_d(\mathfrak S[[X]])$ (its determinant has
constant term $1$), $\varepsilon_X=1+\ell X\in\mathfrak S[[X]]^\times$, $\mathcal V=\varepsilon_X^{-a}\mathcal V'$ and
$V=\mathrm{ev}(\mathcal V)\in GL_d(\pkgSh)$. As $\varphi^l$ is a ring endomorphism of $\pkfBp$ (P1),
$\varphi^l(B)\varphi^l(Z)=\varphi^l(Y_1)=1+z^{\ell^l}\varphi^l(W)=1+X_l\ell^2\varphi^l(W)=\mathrm{ev}(\mathcal V')$, and
$p_l=z^{\ell^l}+\ell=\ell\,\mathrm{ev}(\varepsilon_X)$. Since $\mathrm{ev}(\mathcal V')$ is invertible,
$\det\varphi^l(B)\ne0$; as $\varphi^l(B)\in M_d(K_0)$, it is invertible over $K_0$, and
$\varphi^l(Z)=\varphi^l(B)^{-1}\mathrm{ev}(\varepsilon_X)^aV=p_l^aDV$ with $D=\ell^{-a}\varphi^l(B)^{-1}$.

For each entry of $\mathcal V$ and of $\mathcal V^{-1}$ choose by Lemma~\ref{pkg:approx} an $s\in\mathfrak S$ as
there, and let $M'$, $M$ be the resulting matrices: modulo the ideal $E^A\pkfBp$ (entrywise),
$M'\equiv\ell^{2\delta}V$ and $M\equiv\ell^{2\delta}V^{-1}$. Hence $VM\equiv\ell^{2\delta}\cdot1$ and
$V(\ell^{4\delta}-MM')\equiv\ell^{4\delta}V-\ell^{2\delta}M'\equiv0$; so the columns of $M$ and of
$\ell^{4\delta}-MM'$ lie in $\mathcal Y$, and for $x\in\mathfrak S^d$,
$\ell^{4\delta}x=(\ell^{4\delta}-MM')x+M(M'x)$ is an $\mathfrak S$-combination of them.
\end{proof}

\begin{theorem}[Spanning by almost fixed vectors]\label{pkg:span}
Let the chart be unramified, $l\ge1$, $\ell^l\ge3$, $A\le\ell^l\delta$, and
$\pkgV_l=\{x\in\mathcal N(L):\varphi^{-l}x\in\pknMinf(L)\}$. For $x\in\pkgV_l$, $\ell^{4\delta}x$ lies in the
$\mathfrak S$-span of the $y\in\pkgV_l$ with $g\mu(y)-\mu(y)\in\mathrm{Fil}^A$ for all forms $\mu$ and
$g\in\Gamma'$.
\end{theorem}

\begin{proof}
\emph{$\pkgV_l$ is the $\mathfrak S$-span of the $\varphi^l(n_j)$.} These lie in $\mathcal N$
(Lemma~\ref{pkl:phin}) and in $\varphi^l\pknMinf$ (Proposition~\ref{pkn:frame}(i)), and $\pkgV_l$ is an
$\mathfrak S$-module ($\varphi^{-l}(\mathfrak S)\subseteq\pkpA$). Conversely, for $x\in\pkgV_l$,
$\varphi^{-l}x=\sum_jt_jn_j$ (Proposition~\ref{pkn:frame}(ii)), so $x=\sum_j\varphi^l(t_j)\varphi^l(n_j)$ and,
applying $\mathrm{ret}$ (Proposition~\ref{pkl:R}), $x=\sum_j\kappa_0(\varphi^lt_j)\varphi^l(n_j)$.

\emph{Conclusion.} Let $x=\sum_jc_j\varphi^l(n_j)$, $c\in\mathfrak S^d$, and write $\ell^{4\delta}c=\sum_ks_ky^{(k)}$
with $y^{(k)}\in\mathcal Y$ (Proposition~\ref{pkg:radius}). Then $\ell^{4\delta}x=\sum_ks_kx_k$ with
$x_k=\sum_jy^{(k)}_j\varphi^l(n_j)\in\pkgV_l$. Let $y\in\mathcal Y$, $Vy=r+E^Aw$, $x_y=\sum_jy_j\varphi^l(n_j)$, and
$\mu$ a form. By Lemma~\ref{pkn:frame-id}(ii), $p_l^a\mu(x_y)=c^{(l)}(\mu)\varphi^l(Z)y=p_l^ac^{(l)}(\mu)DVy$; as
$p_l\ne0$ in $\pknB$, $\mu(x_y)=c^{(l)}(\mu)Dr+E^Ac^{(l)}(\mu)Dw$. The first term lies in $K'$ (entries of
$c^{(l)}(\mu)$ in $K'$, Lemma~\ref{pkn:frame-id}(i); $D$ over $K_0\subseteq K'$) and is fixed by $\Gamma'$; the
second lies in $\mathrm{Fil}^A$, which $\Gamma'$ preserves. So $g\mu(x_y)-\mu(x_y)\in\mathrm{Fil}^A$.
\end{proof}

\subsection{Expansion in powers of $E$ and the proof of (I5)}\label{pkg:ss-I5}

\begin{proposition}[Taylor expansion]\label{pkg:taylor}
Let the chart be unramified, $l\ge1$, $\ell^l\ge3$, $A\le\ell^l\delta$, and $x\in\pkgP_l(i)$. There are
$\xi_n\in\pkgC_{l,A}(i)$ ($n\ge0$) such that for every $N$ there is $y_N\in\pkgP_l(i)$ with
$\ell^{4\delta}x=\sum_{n<N}E^n\xi_n+E^Ny_N$.
\end{proposition}

\begin{proof}
By Proposition~\ref{pkl:std}, $\varphi^{-l}x\in\pknMinf$, so $x\in\pkgV_l$. By Theorem~\ref{pkg:span},
$\ell^{4\delta}x=\sum_kr_kv_k$ with $r_k\in\mathfrak S$ and $v_k$ as there; applying $e_i$ ($x=e_ix$),
$\ell^{4\delta}x=\sum_kr_kb_k$ with $b_k=e_iv_k\in\pkgC_{l,A}(i)$ (Lemma~\ref{pkg:const}(ii)). By
Lemma~\ref{pkg:digits}, $r_k=\sum_{n<N}E^nw_{k,n}+E^Nt_{k,N}$; put $\xi_n=\sum_kw_{k,n}b_k\in\pkgC_{l,A}(i)$ and
$y_N=\sum_kt_{k,N}b_k\in\pkgP_l(i)$.
\end{proof}

\begin{theorem}[Input (I5)]\label{pkg:I5}
Let the chart be unramified, $l\ge1$ with $\ell^l\ge3$, and $A,\delta,\kappa,s,g\ge0$ with $A\le\ell^l\delta$.
Let $\ell^\kappa$ kill the cokernel of a side $S$ at $i^{(-l)}$, and $E^s$ kill it at $i$. Let every jump of
$S$ at $i^{(-a)}$, $0\le a<l$, be $\le g$, put $N=s+g(1+\ell+\dots+\ell^{l-1})$, and assume $N<A$. Then for
every $c\in C_S$, $\pkcPt_S(c,i,\kappa+4\delta;o^{\mathrm{gap}}_{l,A,N},d^{\mathrm{gap}}_{l,A})\ne\emptyset$,
where $o^{\mathrm{gap}}_{l,A,N}=\pkgstd_l+A[u<-N][a=0]$ and $d^{\mathrm{gap}}_{l,A}=\pkgstd_l+A[a=0]$.
\end{theorem}

\begin{proof}
We work in $L=X_S$; put $\epsilon=e_i(1\otimes c)$.

\emph{(1) A point and the polynomial.} By Proposition~\ref{pkt:I1} (killer $\ell^0E^s$) there is
$y\in\mathcal N(X_S)\cap(X_S)_i(\rho_{-s})$ with $\mathrm{pr}(y)=\epsilon$. Let $f$ be as in Lemma~\ref{pkg:f}
(degree $N$); $f(E)y\in(X_S)_i(o)$ with $o(a,u)=g[a<l]-u$. On the jumps $u$ at $i^{(-a)}$, $o\ge\pkgstd_l$
($g-u\ge0$ for $a<l$), so $f(E)y\in\pkgP_l(i)$.

\emph{(2) The coefficient $\xi_N$.} By Proposition~\ref{pkg:taylor}, $\ell^{4\delta}f(E)y=\sum_{n\le N}E^n\xi_n+
E^{N+1}y'$ with $\xi_n\in\pkgC_{l,A}(i)$, $y'\in\pkgP_l(i)$. The left side lies in
$(X_S)_i(\pkgstd_l+o_1[a=0])$ with $o_1(u)=\min(A,g-u)\le A$ (at $a=0$, $o(0,u)=g-u\ge o_1(u)$; at $a\ge1$,
$o\ge\pkgstd_l$). By Lemma~\ref{pkg:peel}, $\xi_N\in(X_S)_i(\pkgstd_l+A[o_1(u)>N][a=0])$. If $u<-N$ then
$g-u>N$ and $A>N$, so $o_1(u)>N$; hence $\xi_N\in(X_S)_i(o^{\mathrm{gap}}_{l,A,N})$.

\emph{(3) Its image.} Write $f=\sum_{n\le N}f_nX^n$, $f_N=1$, and $\eta_n=\mathrm{pr}(\xi_n)-\ell^{4\delta}f_n
\epsilon\in\pkgC_{l,A}(i)$ (for $C_S$; Lemma~\ref{pkg:const}(i)). Applying $\mathrm{pr}$ ($\pkpW$-linear),
$\sum_{n\le N}E^n\eta_n+E^{N+1}\mathrm{pr}(y')=\ell^{4\delta}f(E)\epsilon-\ell^{4\delta}f(E)\epsilon=0$, with
$\mathrm{pr}(y')\in(C_S)_i(\pkgstd_l)$. Lemma~\ref{pkg:peel} with $o\equiv A>N$ gives
$w:=\eta_N=\mathrm{pr}(\xi_N)-\ell^{4\delta}\epsilon\in(C_S)_i(d^{\mathrm{gap}}_{l,A})$, and $w\in\mathcal N(C_S)$.

\emph{(4) Lifting the error.} By Lemma~\ref{pkn:E} and Lemma~\ref{pkn:M}(iii), $E^{-A}w\in\mathcal N(C_S)\cap
(C_S)_i(\pkgstd_l)$, and $(C_S)_i(\pkgstd_l)=(C_S)_i(\pkgstd_0)\subseteq\pkpA\otimes C_S$
(Proposition~\ref{pkl:H2}(i)); so $E^{-A}w=e_iE^{-A}w=\sum_kb_ke_i(1\otimes c_k)$ with $b_k\in\pkpA$ and $(c_k)$ an
$\auO$-basis of $C_S$ (refined to a $\Z_\ell$-basis). By Definition~\ref{pkm:def-J}, $\ell^\kappa\in I_0(i^{(-l)})$,
hence $\ell^\kappa=\varphi^l(\ell^\kappa)\in I_l(i)$ (Lemma~\ref{pkm:J-basic}(iii) with $r=0$; $\sigma^l(i^{(-l)})=i$):
there are $x_k\in(X_S)_i(\pkgstd_l)$ with $\mathrm{pr}(x_k)=\ell^\kappa e_i(1\otimes c_k)$. Let $x''=
\mathrm{ret}(\sum_kb_kx_k)$. By Proposition~\ref{pkl:R}, $x''\in\mathcal N(X_S)\cap(X_S)_i(\pkgstd_l)$ and
$\mathrm{pr}(x'')=\mathrm{ret}(\ell^\kappa E^{-A}w)=\ell^\kappa E^{-A}w$, as $\ell^\kappa E^{-A}w\in\mathcal N(C_S)$.

\emph{(5) The point.} Put $\nu=\ell^\kappa\xi_N-E^Ax''\in\mathcal N(X_S)\cap\pkpAmb(X_S)_i$. Then
$\mathrm{pr}(\nu)=\ell^\kappa(w+\ell^{4\delta}\epsilon)-\ell^\kappa w=\ell^{\kappa+4\delta}\epsilon$. By
Lemma~\ref{pkn:E}, $E^Ax''\in(X_S)_i(d^{\mathrm{gap}}_{l,A})\subseteq(X_S)_i(o^{\mathrm{gap}}_{l,A,N})$, so
$\nu$ has orders $o^{\mathrm{gap}}_{l,A,N}$. For $g\in\Gamma'$, $g\xi_N-\xi_N\in(X_S)_i(d^{\mathrm{gap}}_{l,A})$
($\xi_N$ is a constant modulo $E^A$) and $g(E^Ax'')-E^Ax''\in(X_S)_i(d^{\mathrm{gap}}_{l,A})$
(Lemma~\ref{pkl:G1}); so $\nu$ is deep for $d^{\mathrm{gap}}_{l,A}$, and
$\nu\in\pkcPt_S(c,i,\kappa+4\delta;o^{\mathrm{gap}}_{l,A,N},d^{\mathrm{gap}}_{l,A})$.
\end{proof}

\begin{remark}
(I5) is proved as stated in \S\ref{pkd:ss-tools}; the proof of Theorem~\ref{pkd:scheme} applies it without
change. Only $\ell^l\ge3$ is used to make $X_l\in\ell\pkgSh$.
\end{remark}

\subsection{The other inputs}\label{pkg:ss-open}

The input (I4) of \S\ref{pkc:ss-inputs} is Theorem~\ref{pkh:I4}; Proposition~\ref{pki:chart} is
Theorem~\ref{pkj:chart}.


%% ===== proofs/C-kisin-colength.tex =====
% proofs/C-kisin-colength.tex
% The input (I4): colength of the deep constants (few low jumps).
% Labels start with pkh:. Conventions and macros as in proofs/C-kisin-deep.tex, C-kisin-tau.tex,
% C-kisin-cases.tex.
\providecommand{\auO}{\mathcal O}
\providecommand{\auQlb}{\overline{\Q}_\ell}
\providecommand{\pkpA}{A_{\mathrm{inf}}}
\providecommand{\pkpW}{W(C^\flat)}
\providecommand{\pkpAmb}{\mathbb A}
\providecommand{\pknBp}{B^+_{\mathrm{dR}}}
\providecommand{\pknB}{B_{\mathrm{dR}}}
\providecommand{\pknK}{\mathcal K_0}
\providecommand{\pktV}{V_{\mathrm{dR}}}
\providecommand{\pkcPt}{\mathrm{Pt}}
\providecommand{\pkcimp}{\mathrm{imp}}
\providecommand{\pkcNM}{\mathrm{NM}}
\providecommand{\pkcnw}{\mathrm{nw}}
\providecommand{\pkcSig}{\Sigma}
\providecommand{\pkiWin}{\operatorname{Win}}
\providecommand{\pglen}{\operatorname{len}}
\providecommand{\pkecost}{\mathrm{cost}}
\providecommand{\pkeDp}{\mathrm{Dp}}
% macros of this file (prefix \pkh)
\providecommand{\pkhD}{\mathcal D}
\providecommand{\pkhK}{\mathrm{Const}}
\providecommand{\pkhX}{\mathcal X}
\providecommand{\pkhup}[3]{#1^{+}_{#2,#3}}
\providecommand{\pkhlv}[2]{#1^{+}_{#2}}
\providecommand{\pkhrl}{\operatorname{rl}}

\section{The colength of the deep constants: the input (I4)}\label{pkh:sec}

This section proves the input (I4) of \S\ref{pkc:ss-inputs} (``few low jumps: the deep constants''),
exactly as stated there; it is Theorem~\ref{pkh:I4} below. With it, Proposition~\ref{pki:generic}(ii)
is proved from (I1)--(I4), as in \S\ref{pkc:ss-generic}, and (I1)--(I3) are proved in
\S\S\ref{pkt:sec}, \ref{pke:sec}.

\emph{Dependencies.} Proved in full, modulo results cited by label: Lemma~\ref{pkt:hodge},
Lemma~\ref{pkt:expand}, Lemma~\ref{pkt:forms} (Hodge bases and expansion), Lemma~\ref{pkt:sub} and its
proof (forms on the kernel), Proposition~\ref{pkl:R} and Lemma~\ref{pke:ret} (the retraction),
Lemma~\ref{pkl:G1} ($\Gamma'$-stability), Lemmas~\ref{pkn:M} and \ref{pkn:E} (multiplication by $p_b^{\pm1}$),
Lemma~\ref{pkm:improved} (improved orders), Proposition~\ref{pke:A} and Lemma~\ref{pke:op-basic}
(operators of a cost), Proposition~\ref{pkm:kill-unram}, Proposition~\ref{pkt:I1}, Lemmas~\ref{pkc:win} and
\ref{pkc:free}, Lemma~\ref{pkm:le}, Remark~\ref{pkc:mono}, and Definitions~\ref{pkp:def-side} and
\ref{pkp:def-points} (sides and points), Definition~\ref{pkn:def-N} (the modules $\mathcal N$), and the
definition of $\pkcimp$ in \S\ref{pkm:ss-state}. Besides these and published results cited by name (Tate's
theorem on the Galois invariants of $C(n)$), no other input is used; all numerical inequalities are checked in the text.

\emph{Conventions.} Those of \S\S\ref{pkc:sec}, \ref{pkt:sec}, \ref{pke:sec}. Throughout, the Kummer
chart is \emph{unramified}: $e=1$, $E=z+\ell$, $d_v=0$, so $\pkecost(a,n)=v_\ell(n)+3a$ with $v_\ell(0)=0$.
Recall $z=[\varpi^\flat]$, so $\varphi(z)=z^\ell$ and $p_a=\varphi^a(E)=z^{\ell^a}+\ell$. Orders are functions
$o\colon\Z_{\ge0}\times\Z\to\Z$; for orders $o$, a level $a\ge0$ and $u\in\Z$ put
\[
\pkhup oau(a',v)=o(a',v)+[(a',v)=(a,u)],\qquad \pkhlv oa(a',v)=o(a',v)+[a'=a]
\]
(raise the single value $o(a,u)$, resp.\ all values at the level $a$, by one). As in \S\ref{pkc:sec},
inequalities between orders are required only on the jumps, i.e.\ at the $(a',v)$ with $v\in
U(i^{(-a')})$ for the letter $i$ in question. ``Length'' means length of $\Z_\ell$-modules, with values in
$\Z_{\ge0}\cup\{\infty\}$; it is additive in short exact sequences. For $\Z_\ell$-submodules $A,B$ of a
module put
\[
\pkhrl(A,B)=\pglen(A/(A\cap B)).
\]

\subsection{Two algebraic lemmas}\label{pkh:ss-alg}

\begin{lemma}[Relative lengths]\label{pkh:len}
Let $A,B,B'$ and $A_1,\dots,A_k$ be $\Z_\ell$-submodules of a $\Z_\ell$-module $M$.
\begin{enumerate}
\item[(i)] $\pkhrl(A,B')\le\pkhrl(A,B)+\pkhrl(B,B')$; if $A\subseteq B$ then $\pkhrl(A,B)=0$.
\item[(ii)] $\pglen\bigl(M/\bigcap_jA_j\bigr)\le\sum_j\pglen(M/A_j)$.
\item[(iii)] If $f\colon M\to M'$ is $\Z_\ell$-linear and $A'\subseteq M'$ is a submodule with $f(B)\subseteq
A'$, then $\pglen\bigl(f(A)/(f(A)\cap A')\bigr)\le\pkhrl(A,B)$.
\item[(iv)] Let $t,m\ge0$. Suppose $\ell^tA\subseteq B$, and that for any $x_0,\dots,x_m\in A$ with $\ell
x_j\in B$ for all $j$ there are $\beta_0,\dots,\beta_m\in\Z_\ell$, one of them equal to $1$, with
$\sum_j\beta_jx_j\in B$. Then $\pkhrl(A,B)\le tm$.
\end{enumerate}
\end{lemma}

\begin{proof}
(i) In $\bar A=A/(A\cap B')$ let $\bar B=((A\cap B)+(A\cap B'))/(A\cap B')\cong(A\cap B)/(A\cap B\cap B')$. This
is isomorphic to a submodule of $B/(B\cap B')$, so $\pglen\bar B\le\pkhrl(B,B')$; and $\bar A/\bar B=
A/((A\cap B)+(A\cap B'))$ is a quotient of $A/(A\cap B)$, so $\pglen(\bar A/\bar B)\le\pkhrl(A,B)$. Additivity
gives the claim. The second assertion is clear.

(ii) The map $M/\bigcap_jA_j\to\bigoplus_jM/A_j$ is injective.

(iii) $f$ maps $A$ onto $f(A)$ and $A\cap B$ into $f(A)\cap A'$ (as $f(B)\subseteq A'$), so $f(A)/(f(A)\cap
A')$ is a quotient of $A/(A\cap B)$.

(iv) Put $Q=A/(A\cap B)$; it is killed by $\ell^t$. Let $Q[\ell]$ be its $\ell$-torsion, an $\mathbb
F_\ell$-vector space. \emph{Claim: $\dim Q[\ell]\le m$.} Let $q_0,\dots,q_m\in Q[\ell]$ and lift them to
$x_j\in A$; then $\ell x_j\in A\cap B$, so the hypothesis gives $\beta_j\in\Z_\ell$, $\beta_{j_0}=1$, with
$\sum_j\beta_jx_j\in B$, i.e.\ $\sum_j\bar\beta_jq_j=0$ in $Q$ with $\bar\beta_{j_0}=1\ne0$ in $\mathbb F_\ell$.
So any $m+1$ elements of $Q[\ell]$ are linearly dependent. \emph{Induction on $t$}: we show that a
$\Z_\ell$-module $Q$ killed by $\ell^t$ with $\dim Q[\ell]\le m$ has length $\le tm$. For $t=0$, $Q=0$. For
$t\ge1$ the sequence $0\to Q[\ell]\to Q\xrightarrow{\ \ell\ }\ell Q\to0$ is exact, $\ell Q$ is killed by
$\ell^{t-1}$ and $(\ell Q)[\ell]\subseteq Q[\ell]$; so $\pglen Q=\dim Q[\ell]+\pglen\ell Q\le m+(t-1)m$.
\end{proof}

\begin{lemma}[A primitive relation over $\Lambda$]\label{pkh:wrel}
Let $d\le n$ and $w_0,\dots,w_{fn}\in\Lambda^d$, where $\Lambda=W(\mathbb F_{\ell^f})$. There are
$\beta_0,\dots,\beta_{fn}\in\Z_\ell$, one of them equal to $1$, with $\sum_k\beta_kw_k=0$.
\end{lemma}

\begin{proof}
$\Lambda$ is free of rank $f$ over $\Z_\ell$, so $\Lambda^d$ is free of rank $fd\le fn<fn+1$, and the $fn+1$
vectors $w_k$ are linearly dependent over $\Q_\ell$; clearing denominators, $\sum_k\alpha_kw_k=0$ with
$\alpha_k\in\Z_\ell$ not all $0$. Choose $k_0$ with $v_\ell(\alpha_{k_0})$ minimal (so $\alpha_{k_0}\ne0$) and put
$\beta_k=\alpha_k/\alpha_{k_0}\in\Z_\ell$. Then $\alpha_{k_0}\sum_k\beta_kw_k=0$, and $\Lambda^d$ is torsion free, so
$\sum_k\beta_kw_k=0$; $\beta_{k_0}=1$.
\end{proof}

\begin{lemma}[Division by $E$]\label{pkh:Ediv}
(i) Every $s\in\mathfrak S$ is $s=w+Ey$ with $w\in\Lambda$, $y\in\mathfrak S$. (ii) For $a\ge0$,
$E^{\ell^a}-p_a\in\ell\mathfrak S$.
\end{lemma}

\begin{proof}
(i) Write $s=\sum_{k\ge0}\lambda_kz^k$, $\lambda_k\in\Lambda$. Put $w=\sum_k\lambda_k(-\ell)^k$, convergent in
$\Lambda$, and $y=\sum_{j\ge0}y_jz^j$ with $y_j=\sum_{k>j}\lambda_k(-\ell)^{k-1-j}$, convergent in $\Lambda$. Since
$z^k-(-\ell)^k=(z+\ell)\sum_{j<k}z^j(-\ell)^{k-1-j}$, comparing coefficients gives $(z+\ell)y=s-w$: the constant
coefficient of $(z+\ell)y$ is $\ell y_0=-\sum_{k\ge1}\lambda_k(-\ell)^k=\lambda_0-w$, and for $j\ge1$ the coefficient
of $z^j$ is $y_{j-1}+\ell y_j=\lambda_j+\sum_{k>j}\lambda_k(-\ell)^{k-j}-\sum_{k>j}\lambda_k(-\ell)^{k-j}=\lambda_j$.
(ii) $E^{\ell^a}=(z+\ell)^{\ell^a}
\equiv z^{\ell^a}\equiv p_a\pmod{\ell\mathfrak S}$, since $\ell\mid\binom{\ell^a}{k}$ for $0<k<\ell^a$ and
$p_a=z^{\ell^a}+\ell$.
\end{proof}

\subsection{Leading terms and primitive relations}\label{pkh:ss-lead}

For a lattice $L$ with jump sets, a letter $j$ and $u\in\Z$ let $m^L_j(u)$ be the sum, over the
embeddings $\gamma$ of letter $j$, of the multiplicity of the jump $u$ of $L$ at $\gamma$. By
Lemma~\ref{pkt:hodge} (last paragraph of its proof), in a Hodge basis $(d_{j,m},\kappa_{j,m})_m$ of $L$ at $j$
exactly $m^L_j(u)$ indices have $\kappa_{j,m}=u$. For a side $S$ with lattice $Y$ we write $m_j(u)=m^Y_j(u)$;
this is the $m_i(u)$ of (I4).

\begin{lemma}[Extension and kernel]\label{pkh:kernel}
Let $S$ be a side with lattice $Y$ and extension $X_S$, and $\iota\colon\pkpAmb(Y)\to\pkpAmb(X_S)$, $y\mapsto
(y,0)$, the $\pkpW$-linear extension of the inclusion. (i) For every letter $i$ and orders $o$, $\iota$ is a
bijection from $\mathcal N(Y)\cap Y_i(o)$ onto the set of $x\in\mathcal N(X_S)\cap(X_S)_i(o)$ with
$\mathrm{pr}(x)=0$. (ii) For $u\ne0$ and every letter $j$, $m^{X_S}_j(u)=m_j(u)$.
\end{lemma}

\begin{proof}
(i) $\iota$ is injective, and its image is the kernel of $\mathrm{pr}$; $\iota$ is $\Gamma$-equivariant (the
inclusion $Y\to X_S$ is) and commutes with $\varphi$, with $1\otimes\zeta$ and with the $e_j$.
\emph{Into.} Let $y\in\mathcal N(Y)\cap Y_i(o)$. For a form $\mu$ on $X_S$, $\mu\circ\iota$ is
$\pkpA$-linear and $\mu(\iota(1\otimes gy'))=\mu(1\otimes g\iota y')=g\mu(1\otimes\iota y')$, so it is a form
on $Y$; if $\mu$ has type $u$ at $j$, so has $\mu\circ\iota$, as $\iota(e_j(1\otimes Y))\subseteq
e_j(1\otimes X_S)$. If $ry\in\mathfrak M_{\mathrm{inf}}(Y)$ with $r$ a product of $p_a$'s and $\tau^Nry\in
\pkpA\otimes Y$, then $\tau^Nr\iota y=\iota(\tau^Nry)\in\pkpA\otimes X_S$ and $\mu(\varphi^k(r\iota
y))=(\mu\circ\iota)(\varphi^k(ry))\in\pknBp$ for $k\ge1$; so $\iota y\in\pkpA\mathcal N(X_S)$
(Definition~\ref{pkn:def-N}). Likewise $\mu(\iota y)=(\mu\circ\iota)(y)\in\pknK$, so $\iota y\in\mathcal N(X_S)$;
$\iota y$ lies in the letter-$i$ part, and $\mu(\varphi^{-a}\iota y)=(\mu\circ\iota)(\varphi^{-a}y)\in
\mathrm{Fil}^{o(a,u)}$ for $\mu$ of type $u\in U(i^{(-a)})$ at $i^{(-a)}$ (the jump sets of $X_S$ and $Y$ are
the same, Definition~\ref{pkp:def-side}). So $\iota y\in\mathcal N(X_S)\cap(X_S)_i(o)$, over $0$.
\emph{Onto.} Let $x\in\mathcal N(X_S)\cap(X_S)_i(o)$ with $\mathrm{pr}(x)=0$, and $y=\mathrm{fst}(x)$, so
$x=\iota y$. By Lemma~\ref{pkt:sub}, $y\in Y_i(o)\subseteq\pkpA\mathcal N(Y)$. Let $\mu$ be a form on $Y$. The
proof of Lemma~\ref{pkt:sub} (\emph{extension of forms}, with no type condition) gives for every letter $j$ a
form $\nu_j$ on $X_S$ with $\hat\nu_j\circ\hat\iota=\hat\mu\circ e_j$ on $\pktV(Y)$. Hence
$\mu(y)=\hat\mu(\hat y)=\sum_j\hat\mu(e_j\hat y)=\sum_j\hat\nu_j(\widehat{\iota y})=\sum_j\nu_j(x)\in\pknK$,
since $x\in\mathcal N(X_S)$ and $\pknK$ is a field. So $y\in\mathcal N(Y)$.

(ii) Let $\gamma$ be an embedding of letter $j$. Applying $C\otimes_{\auO,\gamma}-$ to the $\auO$-split exact
sequence $0\to Y\to X_S\to C_S\to0$ of $\auO[\Gamma']$-modules gives an exact sequence $0\to W'\to W\to W''\to0$
of $C$-semilinear representations of $G_{K'}$. They are Hodge--Tate ($Y[1/\ell]$, $X_S[1/\ell]$ are
crystalline on $\Gamma'$, and $C_S$ is trivial), so each is a direct sum of copies of $C(-\kappa)$, and by
Tate's theorem ($C(n)^{G_{K'}}$ is $K'$ for $n=0$ and $0$ otherwise) the multiplicity of the jump $\kappa$ is
$h_\kappa=\dim_{K'}(\,\cdot\,(\kappa))^{G_{K'}}$, with $\sum_\kappa h_\kappa$ equal to the $C$-dimension. Taking
invariants is left exact, so $h_\kappa(W)\le h_\kappa(W')+h_\kappa(W'')$ for every $\kappa$, and summing over
$\kappa$ gives equality of the totals; hence $h_\kappa(W)=h_\kappa(W')+h_\kappa(W'')$ for every $\kappa$. As
$W''$ is a sum of copies of $C$, $h_u(W'')=0$ for $u\ne0$. Sum over the $\gamma$ of letter $j$.
\end{proof}

\begin{lemma}[Leading terms]\label{pkh:lead}
Let $L$ be a lattice with jump sets $U$, $l$ a letter, $b\ge0$, $j=l^{(-b)}$, $u\in\Z$, and $o$ orders with
$o(b,u)<o(b,v)$ for every $v\in U(j)$ with $v<u$. Let $x_1,\dots,x_N\in L_l(o)$. There are $d\le m^L_j(u)$,
indices $s_1,\dots,s_d\in\{1,\dots,N\}$ and $A_{k,q}\in\pkpA$ ($1\le k\le N$, $1\le q\le d$) with
\[
x_k-\textstyle\sum_{q=1}^dA_{k,q}x_{s_q}\in L_l(\pkhup obu)\qquad(1\le k\le N).
\]
\end{lemma}

\begin{proof}
If $u\notin U(j)$, then $L_l(\pkhup obu)=L_l(o)$ (the value at $(b,u)$ is not tested) and $d=0$ works. Let
$u\in U(j)$. Fix a Hodge basis $(d_{j,m},\kappa_{j,m})_m$ of $L$ at $j$ (Lemma~\ref{pkt:hodge}) with coordinate
forms $\delta_{j,m}$, and let $m_1,\dots,m_r$ be the indices with $\kappa_{j,m}=u$; $r=m^L_j(u)$. For a vector
$x$ in the letter-$l$ part of $\pkpA\mathcal N(L)$ write $\hat x_b=\widehat{\varphi^{-b}x}\in V_j$
(\S\ref{pkt:ss-hodge}). Each $\delta_{j,m}$ is $\hat\mu$ for a form $\mu$ of type $\kappa_{j,m}\in U(j)$ at $j$,
so for $x\in L_l(o)$, $\delta_{j,m}(\hat x_b)\in\mathrm{Fil}^{o(b,\kappa_{j,m})}$.

\emph{Leading coefficients.} Write $\delta_{j,m_p}((\hat x_k)_b)=t^{o(b,u)}\tau_{k,p}$ with $\tau_{k,p}\in\pknBp$,
and let $T=(\theta(\tau_{k,p}))\in C^{N\times r}$, of rank $d\le r$. If $d=0$ put $\alpha=0$. Otherwise choose
$d$ rows $s_1,\dots,s_d$ and $d$ columns $J$ such that the minor $\Delta=\det T_{s,J}$ has maximal absolute
value among all $d\times d$ minors of $T$; $\Delta\ne0$. The rows $s_q$ are linearly independent, hence span
the row space, and every row is $T_k=\sum_q\alpha_{k,q}T_{s_q}$ with $\alpha_{k,q}\in C$; restricting to the
columns $J$ and using Cramer's rule, $\alpha_{k,q}=\Delta_{k,q}/\Delta$, where $\Delta_{k,q}$ is the determinant
of $T_{s,J}$ with the row $s_q$ replaced by $T_k|_J$. Up to sign $\Delta_{k,q}$ is a $d\times d$ minor of $T$
(or $0$, if $k$ is another $s_{q'}$), so $|\alpha_{k,q}|\le1$ and $\alpha_{k,q}\in O_C$. As $\theta\colon\pkpA\to
O_C$ is onto, choose $A'_{k,q}\in\pkpA$ with $\theta(A'_{k,q})=\alpha_{k,q}$, and put $A_{k,q}=\varphi^b(A'_{k,q})$.
Then $\theta\bigl(\tau_{k,p}-\sum_qA'_{k,q}\tau_{s_q,p}\bigr)=0$, i.e.\
$\tau_{k,p}-\sum_qA'_{k,q}\tau_{s_q,p}\in\mathrm{Fil}^1$, for all $k,p$.

\emph{The difference.} Fix $k$ and put $x'=x_k-\sum_qA_{k,q}x_{s_q}\in L_l(o)$ ($L_l(o)$ is an $\pkpA$-module).
Since $\varphi^{-b}(A_{k,q}x)=A'_{k,q}\varphi^{-b}x$ and $\delta_{j,m}$ is $\pknB$-linear,
\[
\delta_{j,m_p}(\hat x'_b)=t^{o(b,u)}\bigl(\tau_{k,p}-\textstyle\sum_qA'_{k,q}\tau_{s_q,p}\bigr)\in
\mathrm{Fil}^{o(b,u)+1}.
\]
Now let $\mu$ be a form of type $u$ at $j$. By Lemmas~\ref{pkt:forms}(i) and \ref{pkt:expand},
$\mu(\varphi^{-b}x')=\hat\mu(\hat x'_b)=\sum_m\hat\mu(d_{j,m})\delta_{j,m}(\hat x'_b)$ with $\hat\mu(d_{j,m})\in
K'\subseteq\pknBp$, and $\hat\mu(d_{j,m})=0$ unless $\kappa_{j,m}\le u$. The terms with $\kappa_{j,m}=u$ lie in
$\mathrm{Fil}^{o(b,u)+1}$ by the display; a term with $\kappa_{j,m}=v<u$ has $v\in U(j)$ and lies in
$\mathrm{Fil}^{o(b,v)}\subseteq\mathrm{Fil}^{o(b,u)+1}$, as $o(b,v)\ge o(b,u)+1$. So $\mu(\varphi^{-b}x')\in
\mathrm{Fil}^{\pkhup obu(b,u)}$; at all other $(a,v)$ the orders $\pkhup obu$ and $o$ agree. Hence $x'\in
L_l(\pkhup obu)$.
\end{proof}

\begin{lemma}[Primitive relation]\label{pkh:prim}
Let the chart be unramified. Let $L$, $l$, $b$, $j$, $u$, $o$ be as in Lemma~\ref{pkh:lead}, $n\ge m^L_j(u)$,
and $x_0,\dots,x_{fn}\in\mathcal N(L)\cap L_l(o)$. There are $\beta_0,\dots,\beta_{fn}\in\Z_\ell$, one of them
equal to $1$, with
\[
\textstyle\sum_k\beta_kx_k=w+z,\qquad w\in\mathcal N(L)\cap L_l(\pkhlv o0),\quad z\in\mathcal N(L)\cap
L_l(\pkhup obu).
\]
\end{lemma}

\begin{proof}
Lemma~\ref{pkh:lead} gives $d\le m^L_j(u)\le n$, indices $s_q$ and $A_{k,q}$ with $y_k=x_k-\sum_qA_{k,q}x_{s_q}\in
L_l(\pkhup obu)$. Let $\mathrm{ret}$ be the retraction of Proposition~\ref{pkl:R}, with coordinate $\kappa_0\colon
\pkpA\to\mathfrak S$. As the $x$'s lie in $\mathcal N(L)$, $\mathrm{ret}(y_k)=x_k-\sum_q\kappa_0(A_{k,q})x_{s_q}$, and
by Lemma~\ref{pke:ret} it lies in $\mathcal N(L)\cap L_l(\pkhup obu)$. By Lemma~\ref{pkh:Ediv}(i) write
$\kappa_0(A_{k,q})=w_{k,q}+Ey_{k,q}$ with $w_{k,q}\in\Lambda$, $y_{k,q}\in\mathfrak S$. Lemma~\ref{pkh:wrel},
applied to the vectors $(w_{k,q})_{q\le d}\in\Lambda^d$, gives $\beta_k\in\Z_\ell$, one equal to $1$, with
$\sum_k\beta_kw_{k,q}=0$ for every $q$. Then
\[
\textstyle\sum_k\beta_kx_k=\underbrace{\sum_k\beta_k\,\mathrm{ret}(y_k)}_{z}+\underbrace{E\sum_k\beta_k\sum_q
y_{k,q}x_{s_q}}_{w}+\sum_q\Bigl(\sum_k\beta_kw_{k,q}\Bigr)x_{s_q},
\]
and the last sum vanishes. $z\in\mathcal N(L)\cap L_l(\pkhup obu)$, a $\Z_\ell$-module. The element
$w'=\sum_k\beta_k\sum_qy_{k,q}x_{s_q}$ lies in the $\mathfrak S$-module $\mathcal N(L)$ (Lemma~\ref{pkn:M}(i)) and
in $L_l(o)$, so $w=Ew'=p_0w'\in\mathcal N(L)\cap L_l(\pkhlv o0)$ by Lemma~\ref{pkn:E} ($b=0$, $n=1$).
\end{proof}

\subsection{Deep constants and the cost of one step}\label{pkh:ss-step}

Fix a side $S$ with lattice $Y$, extension $X_S$, jump sets $U$; write $\mathcal N=\mathcal N(X_S)$,
$L_i(o)=(X_S)_i(o)$, $L^0_i(o)$ for its elements over $0$, and $\pkeDp_i(o)$ for the set of $x\in\mathcal N$
lying over a constant at $i$ and deep for $o$ at $i$ (\S\ref{pke:ss-ops}). The \emph{constants at $i$} form
the $\Z_\ell$-module $\pkhK_i=\{\varepsilon_i(c):c\in C_S\}$, $\varepsilon_i(c)=e_i(1\otimes c)$.

\begin{definition}[Deep constants]\label{pkh:def-D}
For a letter $i$, $t\in\Z$ and orders $o$ put $\pkhX_i(t;o)=\mathcal N\cap L_i(\rho_t)\cap\pkeDp_i(o)$ and
$\pkhD_i(t;o)=\mathrm{pr}\bigl(\pkhX_i(t;o)\bigr)\subseteq\pkhK_i$.
\end{definition}

\begin{remark}\label{pkh:rem-D}
(a) $\pkeDp_i(o)$ is a $\Z_\ell$-module: sums and $\Z_\ell$-multiples of vectors over constants at $i$ lie over
constants at $i$, and $L_i(o)$ is an $\pkpA$-module. Hence $\pkhX_i(t;o)$ and $\pkhD_i(t;o)$ are
$\Z_\ell$-modules. (b) $\mathcal N\cap L^0_i(o)\subseteq\pkeDp_i(o)$: such a vector lies over the constant
$\varepsilon_i(0)=0$, and $gx-x\in L_i(o)$ for $g\in\Gamma'$ by Lemma~\ref{pkl:G1}. (c) For $c\in C_S$,
$\varepsilon_i(c)\in\pkhD_i(t;o)$ if and only if $\pkcPt_S(c,i,0;\rho_t,o)\ne\emptyset$ (Definition of
$\pkcPt$ in \S\ref{pkc:sec}: a point over $c$ at $i$ lies in $\pkeDp_i(o)$, and conversely a vector of
$\pkhX_i(t;o)$ with $\mathrm{pr}=\varepsilon_i(c)$ lies over $c$ at $i$, being in the letter-$i$ part). (d) If
$o'\le o$ on the jumps then $\pkeDp_i(o)\subseteq\pkeDp_i(o')$ and $\pkhD_i(t;o)\subseteq\pkhD_i(t;o')$
(Lemma~\ref{pkm:le}(i)).
\end{remark}

For orders $c$ and $a\ge0$ we say that \emph{$c$ is $a$-free at $i$} if no multiple of $\ell$ lies strictly
between $c(0,v)$ and $c(0,v)+\ell^a$, for every $v\in U(i)$. Every $c$ is $0$-free.

\begin{lemma}[Division by $\ell$]\label{pkh:div}
Let $a\ge0$, $u\in\Z$, and $c$ orders that are $a$-free at $i$. Let $x\in\pkeDp_i(c)$ and $w\in\mathcal N\cap
L^0_i(\pkhlv c0)$ with $\ell x-w\in\pkeDp_i(\pkhup cau)$. Then there is $m\in\mathcal N\cap L^0_i(c)$ with
$x-m\in\pkeDp_i(\pkhup cau)$.
\end{lemma}

\begin{proof}
\emph{An operator of cost $0$.} Let $c''=c+\ell^a[a'=0]$ and $I=\{(0,y):c(0,v)<y<c(0,v)+\ell^a\text{ for some }
v\in U(i)\}$, a finite set. $I$ covers $\pkhlv c0\to c''$ at $i$ (\S\ref{pke:ss-ops}): at the level $0$ the
pairs $(0,y)$ with $c(0,v)+1\le y<c(0,v)+\ell^a$ are in $I$, and at the levels $a'\ge1$ there are none, as
$\pkhlv c0=c''=c$ there. For $(0,y)\in I$, $\ell\nmid y$ (so $y\ne0$), hence $\pkecost(0,y)=v_\ell(y)=0$.
Proposition~\ref{pke:A} gives an operator $A$ of cost $0$ with $A(\pkeDp_i(\pkhlv c0))\subseteq\pkeDp_i(c'')$.

\emph{The element $w'$.} By Remark~\ref{pkh:rem-D}(b), $w\in\pkeDp_i(\pkhlv c0)$, so $Aw\in\pkeDp_i(c'')$, and
$Aw-w\in L^0_i(\pkhlv c0)$; so $Aw\in\mathcal N\cap L^0_i(\pkhlv c0)$. By Lemma~\ref{pkm:improved}, $Aw\in
L_i(\pkcimp(\pkhlv c0,c''))$. We claim $\pkcimp(\pkhlv c0,c'')\ge c''$ on the jumps. At a level $a'\ge1$ both
arguments equal $c$, and $\pkcimp(c,c)=c$ (it is $c$ if $c>0$ and $\min(c,0)=c$ otherwise). At the level $0$,
for $v\in U(i)$: if $c(0,v)+1>0$ the value is $c''(0,v)$; otherwise $c(0,v)<0$, and since $0$ is a multiple of
$\ell$ not strictly between $c(0,v)$ and $c(0,v)+\ell^a$, we get $c(0,v)+\ell^a\le0$, so the value
$\min(c''(0,v),0)$ equals $c''(0,v)$. Hence $Aw\in L_i(c'')$. Put $w'=E^{-\ell^a}Aw$. By Lemma~\ref{pkn:M}(iii),
$w'\in\mathcal N$; by Lemma~\ref{pkn:E} ($b=0$, $n=-\ell^a$), $w'\in L_i(c)$; and $\mathrm{pr}(w')=0$. So
$w'\in\mathcal N\cap L^0_i(c)$.

\emph{The element $m$.} By Lemma~\ref{pkh:Ediv}(ii) write $E^{\ell^a}-p_a=\ell r$ with $r\in\mathfrak S$, and put
$m=rw'-(Ax-x)$. As $x\in\pkeDp_i(c)$ and $A$ has cost $0$, $Ax-x\in L^0_i(c)$, and $Ax\in\mathcal N$; with
$rw'\in\mathcal N\cap L^0_i(c)$ this gives $m\in\mathcal N\cap L^0_i(c)$. Then $x-m\in\pkeDp_i(c)$
(Remark~\ref{pkh:rem-D}(a),(b)), and it lies over the constant $\mathrm{pr}(x)$.

\emph{Depth.} We have
\[
\ell(x-m)=A(\ell x-w)+(Aw-\ell rw'),\qquad Aw-\ell rw'=(E^{\ell^a}-\ell r)w'=p_aw'.
\]
(Indeed $A(\ell x-w)+Aw-\ell rw'=\ell Ax-\ell rw'=\ell x-\ell m$.) By Lemma~\ref{pke:op-basic}(i), $A(\ell
x-w)\in\pkeDp_i(\pkhup cau)$. By Lemma~\ref{pkn:E}, $p_aw'\in L_i(c+[a'=a])\subseteq L_i(\pkhup cau)$, and this
module is $\Gamma'$-stable (Lemma~\ref{pkl:G1}). So for $g\in\Gamma'$, $\ell\bigl(g(x-m)-(x-m)\bigr)\in L_i(\pkhup
cau)$. Finally $y=g(x-m)-(x-m)\in L_i(c)$ lies in the letter-$i$ part of $\pkpA\mathcal N(X_S)$, and for a form
$\mu$, $\mu(\varphi^{-a'}(\ell y))=\ell\mu(\varphi^{-a'}y)$ with $\ell$ a unit of $\pknBp$; so $y\in L_i(\pkhup
cau)$. Thus $x-m\in\pkeDp_i(\pkhup cau)$.
\end{proof}

\begin{lemma}[A relation that gains one step]\label{pkh:raw}
Let the chart be unramified, $a\ge0$, $u\in\Z$, $n\ge m_{i^{(-a)}}(u)$, and let $c$ be $a$-free at $i$ with
$c(a,u)<c(a,v)$ for all $v\in U(i^{(-a)})$, $v<u$. Let $x_0,\dots,x_{fn}\in\pkeDp_i(c)$ with $\ell x_k=x'_k+m_k$,
$x'_k\in\pkeDp_i(\pkhup cau)$, $m_k\in\mathcal N\cap L^0_i(c)$. There are $\beta_k\in\Z_\ell$, one equal to $1$,
and $m\in\mathcal N\cap L^0_i(c)$ with $\sum_k\beta_kx_k-m\in\pkeDp_i(\pkhup cau)$.
\end{lemma}

\begin{proof}
By Lemma~\ref{pkh:kernel}(i), $m_k=\iota(y_k)$ with $y_k\in\mathcal N(Y)\cap Y_i(c)$. Lemma~\ref{pkh:prim} for
the lattice $Y$ (jump sets $U$), $l=i$, $b=a$, $j=i^{(-a)}$, the orders $c$ and $n\ge m^Y_j(u)$ gives
$\beta_k$, one equal to $1$, with $\sum_k\beta_ky_k=w_Y+z_Y$, $w_Y\in\mathcal N(Y)\cap Y_i(\pkhlv c0)$, $z_Y\in
\mathcal N(Y)\cap Y_i(\pkhup cau)$. Put $w=\iota w_Y\in\mathcal N\cap L^0_i(\pkhlv c0)$ and $z=\iota z_Y\in
\mathcal N\cap L^0_i(\pkhup cau)\subseteq\pkeDp_i(\pkhup cau)$ (Lemma~\ref{pkh:kernel}(i),
Remark~\ref{pkh:rem-D}(b)). Then $x=\sum_k\beta_kx_k\in\pkeDp_i(c)$ and
\[
\textstyle\ell x-w=\sum_k\beta_kx'_k+\iota\bigl(\sum_k\beta_ky_k\bigr)-w=\sum_k\beta_kx'_k+z\in
\pkeDp_i(\pkhup cau).
\]
Lemma~\ref{pkh:div} gives $m$.
\end{proof}

\begin{lemma}[One step]\label{pkh:one}
Let the chart be unramified, $a\ge0$, $u\in\Z$, and let $c$ be $a$-free at $i$, with $c\ge\rho_0$ on the jumps
($c(l,v)\ge-v$ for $v\in U(i^{(-l)})$) and $c(a,u)<c(a,v)$ for all $v\in U(i^{(-a)})$, $v<u$. Then
\[
\pkhrl\bigl(\pkhD_i(0;c),\pkhD_i(0;\pkhup cau)\bigr)\le\pkecost(a,c(a,u))\cdot f\,m_{i^{(-a)}}(u).
\]
\end{lemma}

\begin{proof}
Put $e=\pkecost(a,c(a,u))$, $n=m_{i^{(-a)}}(u)$, $X(o)=\pkhX_i(0;o)$ and $M=\mathcal N\cap L^0_i(c)$. Since
$c\ge\rho_0$ on the jumps, $M\subseteq L_i(\rho_0)$ (Lemma~\ref{pkm:le}(i)), so $M\subseteq X(c)$
(Remark~\ref{pkh:rem-D}(b)); also $X(\pkhup cau)\subseteq X(c)$. Let $B=X(\pkhup cau)+M\subseteq X(c)$. As
$\mathrm{pr}(M)=0$, $\mathrm{pr}(B)\subseteq\pkhD_i(0;\pkhup cau)$, and Lemma~\ref{pkh:len}(iii) (with $f=\mathrm{pr}$,
$A=X(c)$) gives $\pkhrl(\pkhD_i(0;c),\pkhD_i(0;\pkhup cau))\le\pkhrl(X(c),B)$. We bound the latter by
Lemma~\ref{pkh:len}(iv) with $t=e$, $m=fn$.

\emph{$\ell^eX(c)\subseteq B$.} The set $\{(a,c(a,u))\}$ covers $c\to\pkhup cau$ at $i$; Proposition~\ref{pke:A}
gives an operator $A$ of cost $e$ with $A(\pkeDp_i(c))\subseteq\pkeDp_i(\pkhup cau)$. For $x\in X(c)$, $Ax-\ell^e
x\in L^0_i(c)\cap\mathcal N=M$ ($A$ has values in $\mathcal N$), so $Ax\in L_i(\rho_0)$ and $Ax\in X(\pkhup
cau)$; hence $\ell^ex=Ax-(Ax-\ell^ex)\in B$.

\emph{Relations.} Let $x_0,\dots,x_{fn}\in X(c)$ with $\ell x_k\in B$: $\ell x_k=x'_k+m_k$ with $x'_k\in
X(\pkhup cau)$, $m_k\in M$. Lemma~\ref{pkh:raw} gives $\beta_k$, one equal to $1$, and $m\in M$ with
$y=\sum_k\beta_kx_k-m\in\pkeDp_i(\pkhup cau)$. Also $y\in\mathcal N\cap L_i(\rho_0)$, so $y\in X(\pkhup cau)$
and $\sum_k\beta_kx_k=y+m\in B$.
\end{proof}

\begin{lemma}[All steps at one level]\label{pkh:level}
Let the chart be unramified, $a\ge0$, and let $c,c'$ be orders such that $c$ is $a$-free at $i$, $c\ge\rho_0$ on
the jumps, $c'(l,v)\le c(l,v)$ for $l\ne a$ and $v\in U(i^{(-l)})$, and $c'(a,\cdot)$ is nonincreasing on
$U(i^{(-a)})$. Then
\[
\pkhrl\bigl(\pkhD_i(0;c),\pkhD_i(0;c')\bigr)\le\sum_{u\in U(i^{(-a)})}\ \sum_{c(a,u)\le y<c'(a,u)}
\pkecost(a,y)\,f\,m_{i^{(-a)}}(u).
\]
\end{lemma}

\begin{proof}
Write $U_a=U(i^{(-a)})$ and induct on $N(c)=\sum_{u\in U_a}\max(0,c'(a,u)-c(a,u))$, the number of steps
$(u,y)$, $u\in U_a$, $c(a,u)\le y<c'(a,u)$. If $N(c)=0$, then $c'\le c$ on all jumps, so $\pkhD_i(0;c)\subseteq
\pkhD_i(0;c')$ (Remark~\ref{pkh:rem-D}(d)) and the left side is $0$. Otherwise let $u$ be the least element of
$U_a$ with $c(a,u)<c'(a,u)$. For $v\in U_a$, $v<u$: $c'(a,v)\le c(a,v)$ by minimality, and
$c(a,u)<c'(a,u)\le c'(a,v)$ as $c'(a,\cdot)$ is nonincreasing; so $c(a,u)<c(a,v)$. Put $\tilde c=\pkhup cau$. It
is $a$-free at $i$: if $a\ne0$ it agrees with $c$ at the level $0$, and all orders are $0$-free. Also
$\tilde c\ge c\ge\rho_0$, $c'(l,\cdot)\le\tilde c(l,\cdot)$ for $l\ne a$, and the steps for $(\tilde c,c')$ are
those for $(c,c')$ except $(u,c(a,u))$; so $N(\tilde c)=N(c)-1$. By Lemmas~\ref{pkh:len}(i), \ref{pkh:one} and
induction,
\[
\pkhrl(\pkhD_i(0;c),\pkhD_i(0;c'))\le\pkecost(a,c(a,u))\,f\,m_{i^{(-a)}}(u)+\sum_{(v,y)\ne(u,c(a,u))}
\pkecost(a,y)\,f\,m_{i^{(-a)}}(v),
\]
which is the sum over all steps.
\end{proof}

\begin{corollary}[Level bound with budgets]\label{pkh:levelcor}
In the situation of Lemma~\ref{pkh:level}, let $\eta\ge0$, $C_{\mathrm{all}},C_{\mathrm{low}}\ge0$, and let
$n_j,r_j$ be integers with $n_j\ge f\sum_{v\in[-\eta,\eta]}m_j(v)$ and $r_j\ge f\sum_{v\in U(j)}m_j(v)$ for
every letter $j$. If $\sum_{c(a,u)\le y<c'(a,u)}\pkecost(a,y)\le C_{\mathrm{all}}+C_{\mathrm{low}}[|u|\le\eta]$
for every $u\in U(i^{(-a)})$, then $\pkhrl(\pkhD_i(0;c),\pkhD_i(0;c'))\le C_{\mathrm{all}}r_{i^{(-a)}}+
C_{\mathrm{low}}n_{i^{(-a)}}$.
\end{corollary}

\begin{proof}
With $j=i^{(-a)}$ the bound of Lemma~\ref{pkh:level} is at most $\sum_{u\in U(j)}(C_{\mathrm{all}}+C_{\mathrm{low}}
[|u|\le\eta])fm_j(u)=C_{\mathrm{all}}f\sum_{u\in U(j)}m_j(u)+C_{\mathrm{low}}f\sum_{u\in U(j)\cap[-\eta,\eta]}
m_j(u)\le C_{\mathrm{all}}r_j+C_{\mathrm{low}}n_j$.
\end{proof}

\subsection{Few low jumps}\label{pkh:ss-few}

In this subsection the side $S$ satisfies the hypotheses of (I4) (there the side is called $S'$): the chart is
unramified, $\ell\ge5$, $a,\eta,W\ge0$, $\eta\ge1$, and for every letter $j$ an integer $b_j\ge0$ and a sign
$\varepsilon_j$ are given with
\[
\begin{gathered}b_j+W+\eta\le(\ell-3)a,\qquad n_0:=6(a+\eta+W)<b_j,\qquad\pkcNM(\ell,n_0,b_j),\\ U(j)\subseteq
\pkiWin(\eta,\varepsilon_j,b_j,W)=[-\eta,\eta]\cup\varepsilon_j[b_j,b_j+W];\end{gathered}
\]
and $n_j,r_j$ are integers with $n_j\ge f\sum_{v\in[-\eta,\eta]}m_j(v)$, $r_j\ge f\sum_{v\in U(j)}m_j(v)$. Recall
that $\pkcNM(\ell,n_0,b)$ means that no nonzero element of $\pkcSig_{n_0}(b)=\{y:|y|\le n_0\text{ or }|y\mp b|\le
n_0\}$ is divisible by $\ell$. A jump $u\in U(j)$ is \emph{low} if $|u|\le\eta$ and \emph{far} otherwise; a far
jump has $b_j\le|u|\le b_j+W$, so $|u|>n_0\ge a+\eta$. Every jump is $\le b_j+W+\eta\le(\ell-3)a$.

\begin{lemma}[Deep points over multiples of constants]\label{pkh:points}
Let $0\le s\le a$, $i$ a letter and $c\in C_S$. Then
$\pkcPt_S\bigl(c,i,s;\pkcimp(\rho_{-(a-s)},\rho_{2a+\eta}),\rho_{2a+\eta}\bigr)\ne\emptyset$, where
\[
\pkcimp(\rho_{-(a-s)},\rho_{2a+\eta})(0,u)=\begin{cases}2a+\eta-u&u<-(a-s),\\ \min(2a+\eta-u,0)&u\ge-(a-s),
\end{cases}\qquad \pkcimp(\rho_{-(a-s)},\rho_{2a+\eta})(a',u)=-u\ (a'\ge1).
\]
\end{lemma}

\begin{proof}
Every jump of $S$ is $\le(\ell-3)a=(\ell-3)(s+(a-s))$, so by Proposition~\ref{pkm:kill-unram} ($\ell\ge5$),
$\ell^sE^{a-s}$ kills the cokernel of $S$ at every letter, and Proposition~\ref{pkt:I1} gives
$x\in\pkcPt_S(c,i,s;\rho_{-(a-s)},\rho_{-(a-s)})$. We deepen it to $\rho_{2a+\eta}$ by Lemma~\ref{pkc:free}
(unramified case) with $n=n_0$, $b=b_i$. The pairs $(a',y)$ with $\rho_{-(a-s)}(a',u)\le y<\rho_{2a+\eta}(a',u)$
for some $u\in U(i^{(-a')})$ have $a'=0$ (at the levels $a'\ge1$ both orders are $-u$) and
$-(a-s)-u\le y<2a+\eta-u$ with $u\in U(i)$; by Lemma~\ref{pkc:win}(iii) (with $s'=a-s$, $t'=2a+\eta$, and
$n_0=6(a+\eta+W)\ge\max(a-s,2a+\eta)+\eta+W=2a+2\eta+W$) they satisfy $y\in\pkcSig_{n_0}(b_i)$. Since
$\pkcNM(\ell,n_0,b_i)$ holds, Lemma~\ref{pkc:free} gives a point in
$\pkcPt_S(c,i,s;\pkcimp(\rho_{-(a-s)},\rho_{2a+\eta}),\rho_{2a+\eta})$. The formula is the definition of
$\pkcimp$ (\S\ref{pkm:ss-state}): at the level $0$, $\rho_{-(a-s)}(0,u)=-(a-s)-u>0$ iff $u<-(a-s)$; at the levels
$a'\ge1$ both arguments are $-u$ and $\pkcimp=-u$ in both cases.
\end{proof}

\begin{lemma}[Removing the pole]\label{pkh:pole}
For every letter $i$, $\pkhrl\bigl(\pkhK_i,\pkhD_i(0;\rho_{2a+\eta})\bigr)\le a\,n_i$.
\end{lemma}

\begin{proof}
Write $s_0=2a+\eta$ and $\pkhD(t)=\pkhD_i(t;\rho_{s_0})$.

\emph{Start: $\pkhK_i\subseteq\pkhD(-\eta)$.} Let $c\in C_S$ and $x$ the point of Lemma~\ref{pkh:points} with
$s=0$, orders $o=\pkcimp(\rho_{-a},\rho_{s_0})$. We check $o\ge\rho_{-\eta}$ on the jumps. At levels $a'\ge1$
both are $-u$. At the level $0$, for $u\in U(i)$: if $u<-a$, $o(0,u)=2a+\eta-u\ge-\eta-u$. If $u\ge-a$, then
$u$ is low or $u>0$ (a far jump with $u\ge-a$ has $|u|>a$, hence $u>a\ge0$); in both cases $-\eta-u\le0$, and
$-\eta-u\le2a+\eta-u$, so $o(0,u)=\min(2a+\eta-u,0)\ge-\eta-u$. By Remark~\ref{pkc:mono}, $x\in\pkcPt_S(c,i,0;
\rho_{-\eta},\rho_{s_0})$, so $\varepsilon_i(c)\in\pkhD(-\eta)$ (Remark~\ref{pkh:rem-D}(c)).

\emph{Kill: $\ell^a\pkhK_i\subseteq\pkhD(t)$ for $t\le0$.} Let $x$ be the point of Lemma~\ref{pkh:points} with
$s=a$, orders $o=\pkcimp(\rho_0,\rho_{s_0})$; $\mathrm{pr}(x)=\ell^a\varepsilon_i(c)=\varepsilon_i(\ell^ac)$. At
the level $0$: if $u<0$, $o(0,u)=2a+\eta-u\ge t-u$; if $u\ge0$, $o(0,u)=\min(2a+\eta-u,0)\ge t-u$ as $t\le0\le
u$. At levels $\ge1$ both are $-u$. So $x\in\pkcPt_S(\ell^ac,i,0;\rho_t,\rho_{s_0})$ and
$\ell^a\varepsilon_i(c)\in\pkhD(t)$.

\emph{Step: for $-\eta\le t\le-1$, $\pkhrl(\pkhD(t),\pkhD(t+1))\le a\,f\,m_i(t)$.} Apply
Lemma~\ref{pkh:len}(iv) with $A=\pkhD(t)$, $B=\pkhD(t+1)$, exponent $a$ and $m=fn$, $n=m_i(t)$; $\ell^aA\subseteq
B$ by the kill step ($t+1\le0$). Let $k_0,\dots,k_{fn}\in\pkhD(t)$ and choose $x_k\in\pkhX_i(t;\rho_{s_0})$ with
$\mathrm{pr}(x_k)=k_k$. Since $t\ne0$, $m^{X_S}_i(t)=m_i(t)=n$ (Lemma~\ref{pkh:kernel}(ii)). The orders
$\rho_t$ satisfy $\rho_t(0,t)=0<t-v=\rho_t(0,v)$ for $v<t$. Lemma~\ref{pkh:prim} for the lattice $X_S$,
$l=j=i$, $b=0$, $u=t$, $o=\rho_t$ gives $\beta_k\in\Z_\ell$, one equal to $1$, with $x=\sum_k\beta_kx_k=w+z$,
$w\in L_i(\pkhlv{(\rho_t)}0)=L_i(\rho_{t+1})$, $z\in L_i(\pkhup{(\rho_t)}0t)$. As $\pkhup{(\rho_t)}0t\le\rho_{t+1}$
(at $(0,v)$: $t-v+[v=t]\le t+1-v$; elsewhere equal), $x\in L_i(\pkhup{(\rho_t)}0t)$; and $x\in\pkeDp_i(\rho_{s_0})$.
By Lemma~\ref{pkm:improved}, $x\in L_i(\pkcimp(\pkhup{(\rho_t)}0t,\rho_{s_0}))$. We check that this is $\ge
\rho_{t+1}$ on the jumps. At levels $\ge1$ both arguments are $-v$ and $\pkcimp=-v$. At the level $0$, let
$\lambda=t-v+[v=t]$: if $\lambda>0$ the value is $s_0-v\ge t+1-v$ (as $s_0\ge0>t$); if $\lambda\le0$ then
$v\ne t$ (else $\lambda=1$) and $t-v\le0$, so $v\ge t+1$, and the value $\min(s_0-v,0)$ is $\ge t+1-v$. Hence
$x\in\pkhX_i(t+1;\rho_{s_0})$ and $\sum_k\beta_kk_k=\mathrm{pr}(x)\in B$.

\emph{Conclusion.} By Lemma~\ref{pkh:len}(i) applied repeatedly,
\[
\pkhrl(\pkhK_i,\pkhD(0))\le\pkhrl(\pkhK_i,\pkhD(-\eta))+\sum_{t=-\eta}^{-1}\pkhrl(\pkhD(t),\pkhD(t+1))\le
0+a\,f\sum_{t=-\eta}^{-1}m_i(t)\le a\,n_i.\qedhere
\]
\end{proof}

\begin{lemma}[Costs at the level $0$]\label{pkh:cost0}
Let $i$ be a letter with $\varepsilon_i=-1$, and $q\ell$ the least multiple of $\ell$ with $q\ell\ge
b_i+W+a+\eta$. For $u\in U(i)$,
\[
\sum_{2a+\eta-u\le y<q\ell}\pkecost(0,y)\le(a+1)[|u|\le\eta].
\]
\end{lemma}

\begin{proof}
Here $\pkecost(0,y)=v_\ell(y)\ge0$, and $q\ell\le b_i+W+a+\eta+\ell-1$. \emph{$u$ low.} The steps lie in
$[0,q\ell]$ as $2a+\eta-u\ge2a\ge0$. With $N=q\ell$, $\sum_{0\le y\le N}v_\ell(y)=v_\ell(N!)=\sum_{k\ge1}\lfloor
N/\ell^k\rfloor\le N/(\ell-1)$ (Legendre). And $N\le b_i+W+\eta+a+\ell-1\le(\ell-3)a+a+\ell-1\le(\ell-1)a+\ell-1=
(\ell-1)(a+1)$; so the sum is $\le a+1$. \emph{$u$ far.} As $\varepsilon_i=-1$, $-u\in[b_i,b_i+W]$. Let $y$ be a
step with $\ell\mid y$. Then $y\le q\ell-\ell<b_i+W+a+\eta$ (minimality of $q\ell$) and $y\ge2a+\eta-u\ge
b_i+2a+\eta$; so $0\le y-b_i\le W+a+\eta\le n_0$ and $y\ge b_i>0$, contradicting $\pkcNM(\ell,n_0,b_i)$. So every
step has $v_\ell(y)=0$ and the sum is $0$.
\end{proof}

\begin{lemma}[Costs at the level $1$]\label{pkh:cost1}
Let $i$ be a letter, $m=i^{(-1)}$, and $\pkcnw_i$ the depth of (I4). For $u\in U(m)$,
\[
\sum_{-u\le y<\pkcnw_i(1,u)}\pkecost(1,y)\le3(W+1)+6a[|u|\le\eta].
\]
\end{lemma}

\begin{proof}
$\pkcnw_i(1,u)=W+1-u+2a[|u|\le\eta]$, so there are $W+1+2a[|u|\le\eta]$ steps $y\in[-u,-u+W+2a[|u|\le\eta]]$,
and $\pkecost(1,y)=v_\ell(y)+3$. Each step lies in $\pkcSig_{n_0}(b_m)$: if $u$ is low, $|y|\le\eta+W+2a\le n_0$;
if $u=b_m+k$ ($0\le k\le W$), then $y+b_m\in[-k,W-k]$, so $|y+b_m|\le W$; if $u=-(b_m+k)$, then $y-b_m\in[k,k+W]$,
so $|y-b_m|\le2W\le n_0$. By $\pkcNM(\ell,n_0,b_m)$, $y=0$ or $\ell\nmid y$; in both cases $v_\ell(y)=0$. So
the sum is $3(W+1+2a[|u|\le\eta])$.
\end{proof}

\begin{lemma}[The negative window]\label{pkh:neg}
Let $i$ be a letter with $\varepsilon_i=-1$ and $m=i^{(-1)}$. Then
\[
\pkhrl\bigl(\pkhD_i(0;\rho_{2a+\eta}),\pkhD_i(0;\pkcnw_i)\bigr)\le(a+1)n_i+3(W+1)r_m+6a\,n_m .
\]
\end{lemma}

\begin{proof}
Let $q\ell$ be as in Lemma~\ref{pkh:cost0} and $d(0,u)=q\ell$, $d(a',u)=-u$ for $a'\ge1$. By
Lemma~\ref{pkh:len}(i) it suffices to bound $\pkhrl(\pkhD_i(0;\rho_{2a+\eta}),\pkhD_i(0;d))$ by $(a+1)n_i$ and
$\pkhrl(\pkhD_i(0;d),\pkhD_i(0;\pkcnw_i))$ by $3(W+1)r_m+6an_m$.

\emph{Level $0$.} Corollary~\ref{pkh:levelcor} with level $0$ (that is, with its level parameter, called $a$ there, set to $0$; the $a$ of (I4) plays no role in this choice), $c=\rho_{2a+\eta}$, $c'=d$: $c$ is $0$-free;
$c\ge\rho_0$; for $l\ge1$, $c'(l,u)=-u=c(l,u)$; $c'(0,\cdot)$ is constant; and by Lemma~\ref{pkh:cost0} the
budgets $C_{\mathrm{all}}=0$, $C_{\mathrm{low}}=a+1$ work. The bound is $(a+1)n_i$.

\emph{Level $1$.} Corollary~\ref{pkh:levelcor} with level $1$ (its level parameter set to $1$), $c=d$, $c'=\pkcnw_i$. $d$ is $1$-free at $i$: no
multiple of $\ell$ lies strictly between $q\ell$ and $q\ell+\ell$. $d\ge\rho_0$ on the jumps: for $u\in U(i)$,
$-u\le\max(\eta,b_i+W)\le b_i+W+a+\eta\le q\ell$, and at levels $\ge1$ equality. For $l\ne1$: $\pkcnw_i(0,u)=
b_i+W+a+\eta\le q\ell=d(0,u)$, and $\pkcnw_i(l,u)=-u=d(l,u)$ for $l\ge2$. \emph{$\pkcnw_i(1,\cdot)$ is
nonincreasing on $U(m)$}: let $u<v$ in $U(m)$. If both are low or both far, $\pkcnw_i(1,u)-\pkcnw_i(1,v)=v-u>0$.
If $u$ is low and $v$ far, the difference is $v-u+2a>0$. If $u$ is far and $v$ low, then $u<v\le\eta<b_m\le|u|$
forces $u\le-b_m$, and the difference is $v-u-2a\ge-\eta+b_m-2a>n_0-\eta-2a\ge0$. By Lemma~\ref{pkh:cost1}
(the steps start at $d(1,u)=-u$) the budgets $C_{\mathrm{all}}=3(W+1)$, $C_{\mathrm{low}}=6a$ work. The bound
is $3(W+1)r_m+6an_m$.
\end{proof}

\begin{lemma}[Colength at one letter]\label{pkh:const}
For every letter $i$, with $\delta_i$ as in (I4),
\[
\pkhrl\bigl(\pkhK_i,\pkhD_i(0;\delta_i)\bigr)\le(2a+1)n_i+3(W+1)r_{i^{(-1)}}+6a\,n_{i^{(-1)}}.
\]
\end{lemma}

\begin{proof}
If $\varepsilon_i=1$, $\delta_i=\rho_{2a+\eta}$ and Lemma~\ref{pkh:pole} gives $\le an_i$. If $\varepsilon_i=-1$,
$\delta_i=\pkcnw_i$, and Lemmas~\ref{pkh:len}(i), \ref{pkh:pole}, \ref{pkh:neg} give $\le an_i+(a+1)n_i+
3(W+1)r_{i^{(-1)}}+6an_{i^{(-1)}}$. All $n_j,r_j\ge0$.
\end{proof}

\subsection{The input (I4)}\label{pkh:ss-I4}

\begin{theorem}[Input (I4)]\label{pkh:I4}
The statement (I4) of \S\ref{pkc:ss-inputs} holds as stated: under its hypotheses, the set $V$ of $c'\in
C_{S'}$ with $\pkcPt_{S'}(c',i,0;\rho_0,\delta_i)\ne\emptyset$ for every letter $i$ is a $\Z_\ell$-submodule of
$C_{S'}$ with $\pglen_{\Z_\ell}(C_{S'}/V)\le11\sum_i\bigl(n_i(a+1)+r_i(W+1)\bigr)$.
\end{theorem}

\begin{proof}
Apply \S\ref{pkh:ss-few} to $S=S'$. By Remark~\ref{pkh:rem-D}(c), $V=\bigcap_i\varepsilon_i^{-1}(\pkhD_i(0;\delta_i))$,
an intersection of $\Z_\ell$-submodules (Remark~\ref{pkh:rem-D}(a); $\varepsilon_i$ is $\Z_\ell$-linear). Since
$\varepsilon_i\colon C_S\to\pkhK_i$ is onto, $C_S/\varepsilon_i^{-1}(\pkhD_i(0;\delta_i))\cong\pkhK_i/\pkhD_i(0;\delta_i)$.
By Lemmas~\ref{pkh:len}(ii) and \ref{pkh:const},
\[
\pglen(C_S/V)\le\sum_i\bigl((2a+1)n_i+3(W+1)r_{i^{(-1)}}+6an_{i^{(-1)}}\bigr)=\sum_i\bigl((8a+1)n_i+3(W+1)r_i\bigr),
\]
because $i\mapsto i^{(-1)}$ is a bijection of the set of letters. Finally $8a+1\le11(a+1)$ and $3(W+1)\le
11(W+1)$, and $n_i,r_i\ge0$.
\end{proof}

\begin{remark}[Use in Proposition~\ref{pki:generic}(ii)]\label{pkh:rem-use}
The statement is proved as stated, so the proof of Proposition~\ref{pki:generic}(ii) in \S\ref{pkc:ss-generic}
goes through unchanged. The proof gives the slightly better bound $\sum_i((8a+1)n_i+3(W+1)r_i)$; the
hypothesis $\eta\ge1$ is not used.
\end{remark}


%% ===== proofs/C-kisin-chart.tex =====
% proofs/C-kisin-chart.tex
% Proof of Proposition pki:chart (the letters read thickness for the data). Labels start with pkj:.
% Conventions and macros as in proofs/C-kisin.tex.
\providecommand{\auO}{\mathcal O}
\providecommand{\auQlb}{\overline{\Q}_\ell}
\providecommand{\pkilog}{\operatorname{lg}}
% macros of this file (prefix \pkj)
\providecommand{\pkjG}{G_{\Q_\ell}}
\providecommand{\pkjcl}{\operatorname{cl}}
\providecommand{\pkjPl}{\mathsf P}
\providecommand{\pkjPb}{\mathsf P_b}

\section{The chart for the data}\label{pkj:sec}

This section proves Proposition~\ref{pki:chart}: for the data, at every prime $\ell$ and every place
$v\mid\ell$, the letters read thickness in the sense of Definition~\ref{pki:def-reads}. The proof
uses only the arithmetic of the field $E$ of the data (\S\ref{an:ssec-param}), the two properties of
the set $\pkjPb$ of $b$-places required in (D5) of Definition~\ref{an:def-data} (class $-1$ modulo
$q_0$, and sparseness for the primes up to $5r(r+1-m)=kt$, Definition~\ref{an:def-sparse}), the
sizes of Lemma~\ref{an:lem-sizes}, and elementary facts on roots of unity in local fields. No
Kummer chart is needed: Definition~\ref{pki:def-reads} involves only $e$, $f$ (those of $E_v$), the
root $\zeta$ and the letters.

We prove Proposition~\ref{pki:chart} (Theorem~\ref{pkj:chart}) for either of
the two natural readings of ``the complex embedding attached to $\gamma$'' (Convention~\ref{pkj:conv});
the paper uses reading (R$^-$) (\S\ref{pll:ss-statement}, \S\ref{pkd:ss-data}).

\subsection{Setting and conventions}\label{pkj:ss-setting}

Let $(q_0,d)$ be the parameters and $D$ the admissible data (Definition~\ref{an:def-data}), with
$E=K(\zeta_{q_1})$, $K=\Q(\sqrt{-d})$, $q_1=q_1(s)$, $f_E=[E:K]=\varphi(q_1)$ (written $f$ in
\S\ref{an:sec-data}; here $f$ is reserved for the residue degree of $E_v$), $\pkjPl$ the set of
infinite places of $E$, $\pkjPb\subseteq\pkjPl$ the $b$-places, $Q=kt=5r(r+1-m)$ and $h=r^8$. We
write $\zeta_{q_1}\in E$ for the fixed primitive $q_1$-th root of unity, and put
\[
N_E=q_1d .
\]

\emph{The CM type and the signs.} As in \S\ref{an:ssec-param}, an embedding $K\subset\C$ is fixed;
let $\delta_0\in\C$ be the image of $\sqrt{-d}$. The CM type $\Phi$ of the data is the set of
embeddings $\sigma\colon E\to\C$ extending $K\subset\C$; each place $\tau\in\pkjPl$ is induced by
exactly one $\Phi(\tau)\in\Phi$, the other embedding inducing $\tau$ being $\overline{\Phi(\tau)}$.
For $\sigma\colon E\to\C$ put $\varepsilon(\sigma)=1$ if $\sigma\in\Phi$ and $-1$ otherwise. Then
\begin{equation}\label{pkj:eq-sign}
\sigma(\sqrt{-d})=\varepsilon(\sigma)\,\delta_0,\qquad
\Phi(\tau_\sigma)(\zeta_{q_1})=\sigma(\zeta_{q_1})^{\varepsilon(\sigma)},
\end{equation}
where $\tau_\sigma$ is the place of $\sigma$: if $\sigma\notin\Phi$ then $\sigma|_K$ is the conjugate
embedding and $\Phi(\tau_\sigma)=\bar\sigma$, and complex conjugation inverts roots of unity.

\emph{Ratios of places.} For $\tau,\tau'\in\pkjPl$, $\Phi(\tau)$ and $\Phi(\tau')$ both extend
$K\subset\C$, so $\Phi(\tau)=\Phi(\tau')\circ\phi$ for a unique $\phi\in\operatorname{Gal}(E/K)$; if
$\phi(\zeta_{q_1})=\zeta_{q_1}^{j}$ we write $\tau/\tau'=j\in(\Z/q_1)^\times$. Thus
\begin{equation}\label{pkj:eq-ratio}
\Phi(\tau)(\zeta_{q_1})=\Phi(\tau')(\zeta_{q_1})^{\tau/\tau'},
\end{equation}
and $(\tau/\tau')(\tau'/\tau'')=\tau/\tau''$. Fixing a base place $\tau_0$, $\tau\mapsto\tau/\tau_0$ is
the identification $\pkjPl\cong(\Z/q_1)^\times$ of \S\ref{an:ssec-param}, and for $m'\mid q_1$ the
reduction $\pkjcl_{m'}$ satisfies $\pkjcl_{m'}(\tau/\tau')=\pkjcl_{m'}(\tau)\pkjcl_{m'}(\tau')^{-1}$,
independently of $\tau_0$. This direction of the ratio, \eqref{pkj:eq-ratio}, is the one used
throughout. (If the identification were read with the inverse ratio, every statement
below holds with the exponent $\varpi$ of Lemma~\ref{pkj:ratio} replaced by $-\varpi$; the proofs of
Lemmas~\ref{pkj:signs} and~\ref{pkj:arc} cover both values of $\varpi$.)

\emph{The local data.} Fix the isomorphism $\iota\colon\C\to\auQlb$ and let $v\mid\ell$. Let
$e_v\colon E\to\auQlb$ be an embedding inducing $v$ (for instance $e_v=\iota\circ\sigma$ for a complex
embedding $\sigma$ with this property), $E_v=\Q_\ell\,e_v(E)\subset\auQlb$ (a finite extension of
$\Q_\ell$, hence complete), and identify $\Gamma=D_v$ with
\[
\Gamma=\operatorname{Gal}(\auQlb/E_v)=\{g\in\pkjG:\ g(e_v(x))=e_v(x)\ \text{for all }x\in E\},
\qquad \pkjG=\operatorname{Gal}(\auQlb/\Q_\ell).
\]
Let $e$ and $f$ be the ramification index and residue degree of $E_v/\Q_\ell$, as in
\S\ref{pki:ss-arch}. Since $E/\Q$ is abelian, $e_v(E)$ is stable under $\pkjG$. Every $g\in\pkjG$ is an
isometry of $\auQlb$ (uniqueness of the extension of $|\cdot|_\ell$), so it preserves the valuation
ring and induces an automorphism $\bar g$ of the residue field $\overline{\mathbb F}_\ell$.

\begin{definition}[Convention]\label{pkj:conv}
For $g\in\pkjG$ let
\[
\sigma_g=\iota^{-1}\circ g^{-1}\circ e_v\quad\text{(reading (R$^-$))}\qquad\text{or}\qquad
\sigma_g=\iota^{-1}\circ g\circ e_v\quad\text{(reading (R$^+$))}.
\]
The paper uses reading (R$^-$); we prove everything for both readings, putting $\varpi=-1$ for (R$^-$)
and $\varpi=+1$ for (R$^+$), one of the two being fixed throughout a proof. Write
$P(g)\in\pkjPl$ for the place of $\sigma_g$ and $\varepsilon(g)=\varepsilon(\sigma_g)$; call $g$
\emph{thick} if $P(g)\in\pkjPb$. A coefficient ring $\auO\subset\auQlb$ is \emph{large} if it contains
the ring of integers $\auO_{E_v}$ of $E_v$. For such $\auO$, every $\Z_\ell$-algebra embedding $\gamma\colon\auO\to
\C_\ell$ has image in $\auQlb$ and is the restriction of some $g\in\pkjG$ (extend $\gamma$ $\Q_\ell$-linearly to the algebraic extension $\auO[1/\ell]$ of $\Q_\ell$, whose image is algebraic over $\Q_\ell$, and then to an automorphism of $\auQlb$); the embeddings of Proposition~\ref{pki:chart} are of this kind; the restriction of $g$ to
$e_v(E)\subset\auO[1/\ell]$ is determined by $\gamma$, hence so is $\sigma_g$ (in (R$^-$) because
$g^{-1}|_{e_v(E)}=(g|_{e_v(E)})^{-1}$). We put $\sigma_\gamma=\sigma_g$, and $\gamma$ is thick,
resp.\ has sign $\varepsilon(\gamma)$, if $g$ is thick, resp.\ $\varepsilon(\gamma)=\varepsilon(g)$.
This is ``the complex embedding of $E$ attached to $\gamma$ through $\iota$ and $v$'' of
Proposition~\ref{pki:chart}.
\end{definition}

(The statements below are proved for all $g\in\pkjG$, which is stronger than for the embeddings
$\gamma$. Reading (R$^-$) is the one fixed for the paper, in \S\ref{pll:ss-statement} and
\S\ref{pkd:ss-data}; (R$^+$) is the other natural one, and \S\ref{pkd:ss-data} uses only
$\sigma_\gamma$, $\varepsilon(\gamma)$ and thickness, so either would do there.)

\subsection{Elementary inputs}\label{pkj:ss-inputs}

\begin{lemma}[Sizes]\label{pkj:sizes}
\begin{enumerate}
\item[(i)] If $\ell\mid N_E$, then $\ell\mid q_0$ or $\ell=d$, and $\ell\le Q$.
\item[(ii)] If $\ell\nmid N_E$, then $\ell\ge13$.
\item[(iii)] $[E_v:\Q_\ell]\le[E:\Q]\le h$; in particular $e\le h$.
\item[(iv)] For every prime $\ell$, $2\,\pkilog_\ell h^3\le24s<l_0(s)$.
\end{enumerate}
\end{lemma}

\begin{proof}
We use Lemma~\ref{an:lem-sizes}: (S1) $\log x\ge10^{14}(q_0+2d)$ and $x\ge(\log b)^{100}\ge128^{100}$;
(S2) $x^{11}\le2^s$ and $r/2\le k$; (S3) $f_E\le x^{12}$; and (D3) $x^5\le r\le3x^5$.

(i) The primes dividing $q_1$ divide $q_0$ (\S\ref{an:ssec-param}) and $d$ is prime; so $\ell\mid q_0$
or $\ell=d$, and $\ell\le q_0+2d\le\log x<x\le r\le5rk=Q$, since $k\ge1$.

(ii) $18480=2^4\cdot3\cdot5\cdot7\cdot11$ divides $q_0$, and $q_0\mid q_1$ since $q_1/q_0=3\,l_0(s)\,2^s$ (\S\ref{an:ssec-param}); so $2,3,5,7,11$ divide $N_E$.

(iii) $E_v=\Q_\ell e_v(E)$ has degree at most $[E:\Q]=2f_E\le2x^{12}\le x^{40}\le r^8=h$; and
$e\le[E_v:\Q_\ell]$.

(iv) $h^3=r^{24}\le3^{24}x^{120}\le x^{132}\le2^{12s}$ (as $x\ge9$ and $2^s\ge x^{11}$). Hence
$\pkilog_\ell h^3\le\log h^3/\log2\le12s$, and $24s<27s<l_0(s)$ (\S\ref{an:ssec-param}).
\end{proof}

\begin{lemma}[Roots of unity of order prime to $\ell$]\label{pkj:roots}
\begin{enumerate}
\item[(i)] Let $L\subset\auQlb$ be a finite extension of $\Q_\ell$ with residue field of $\ell^{f'}$
elements. If $w\in L$ is a root of unity of order $n$ prime to $\ell$, then $n\mid\ell^{f'}-1$.
Moreover, for every $n$ prime to $\ell$, reduction is injective on the group $\mu_n\subset\auQlb$ of
$n$-th roots of unity.
\item[(ii)] If $\ell\nmid n$ and $\theta\in\pkjG$, there is $c\ge0$ with $\theta(w)=w^{\ell^c}$ for all $w\in\mu_n$.
\item[(iii)] $\mu_{\ell^f-1}\subset\auO_{E_v}$, and $\mu_{\ell^f-1}$ is generated by any of its
elements of exact order $\ell^f-1$.
\item[(iv)] If $\ell\nmid N_E$, then $e=1$, and every $\theta\in\pkjG$ that fixes $\mu_{\ell^f-1}$ pointwise
lies in $\Gamma$.
\end{enumerate}
\end{lemma}

\begin{proof}
(i) Let $w\in\mu_n$, $w\ne1$, of order $n'\mid n$. Evaluating $(X^{n'}-1)/(X-1)=\prod_{j=1}^{n'-1}(X-w^j)$
at $X=1$ gives $\prod_{j=1}^{n'-1}(1-w^j)=n'$, a unit since $\ell\nmid n'$; as all $w^j$ are integral,
$1-w$ is a unit, i.e.\ $w\not\equiv1$ modulo the maximal ideal. Hence reduction is injective on
$\mu_n$. If $w\in L$ has order $n$, its reduction lies in the residue field of $L$ and has order $n$,
so $n$ divides the order $\ell^{f'}-1$ of the multiplicative group of that field.

(ii) Let $w_0$ be a primitive $n$-th root of unity; $\mathbb F=\mathbb F_\ell(\bar w_0)$ is a finite
field, and $\bar\theta$ restricts to an automorphism of $\mathbb F$, hence to $x\mapsto x^{\ell^c}$ for some
$c\ge0$. Since $\theta(w_0)\in\mu_n$, its reduction $\bar\theta(\bar w_0)=\bar w_0^{\ell^c}$ is the reduction of
$w_0^{\ell^c}$, so $\theta(w_0)=w_0^{\ell^c}$ by (i); and $\mu_n=\langle w_0\rangle$.

(iii) The residue field of $E_v$ is $\mathbb F_{\ell^f}$, over which $X^{\ell^f-1}-1$ has $\ell^f-1$
simple roots; by Hensel's lemma \cite[II, \S4, Proposition~7, p.~34]{SerreLF} each lifts to a
root in $\auO_{E_v}$, and these lifts are distinct. So $\auO_{E_v}$ contains $\ell^f-1$ roots of
$X^{\ell^f-1}-1$, i.e.\ all of $\mu_{\ell^f-1}$. A cyclic group is generated by any element of maximal
order.

(iv) Since $d\equiv7\pmod8$, $d\equiv3\pmod4$ and $\sqrt{-d}\in\Q(\mu_d)$ (quadratic Gauss sum: for the odd quadratic character $\chi$ of conductor $d$,
$\tau(\chi)^2=\chi(-1)|\tau(\chi)|^2=-d$ by \cite[Lemmas~4.7 and~4.8, p.~36]{Washington}); so $E=K(\zeta_{q_1})\subseteq\Q(\mu_{N_E})$, and
$E_v\subseteq\Q_\ell(\mu_{N_E})$ (with $\mu_{N_E}\subset\auQlb$). For $\ell\nmid N_E$ the extension
$\Q_\ell(\mu_{N_E})/\Q_\ell$ is unramified \cite[IV, \S4, Proposition~16, p.~77]{SerreLF}, hence so is
$E_v$: $e=1$ and $[E_v:\Q_\ell]=f$. By (iii), $F_0=\Q_\ell(\mu_{\ell^f-1})\subseteq E_v$. By (i) the
reductions of the elements of $\mu_{\ell^f-1}$ are $\ell^f-1$ distinct elements of the residue field of
$F_0$, which therefore has at least $\ell^f$ elements; so $[F_0:\Q_\ell]\ge f=[E_v:\Q_\ell]$ (the degree
of a local field is at least its residue degree), and $F_0=E_v$. If $\theta$ fixes $\mu_{\ell^f-1}$
pointwise, it fixes $F_0=E_v\supseteq e_v(E)$, so $\theta\in\Gamma$.
\end{proof}

\begin{lemma}[Letters and the action on roots of unity]\label{pkj:transfer}
Let $\zeta\in\auQlb$ have exact order $\ell^f-1$, $g_0,g_1\in\pkjG$ and $c\in\Z$ with
$g_1(\zeta)=g_0(\zeta)^c$. Then $g_1(w)=g_0(w)^c$ for all $w\in\mu_{\ell^f-1}$. Let $v_\ell=v_\ell(q_1)$,
$q'=q_1/\ell^{v_\ell}$ (the largest divisor of $q_1$ prime to $\ell$) and $y=e_v(\zeta_{q_1})^{\ell^{v_\ell}}$.
Then $y$ is a primitive $q'$-th root of unity, $\mu_{q'}=\langle y\rangle\subseteq\mu_{\ell^f-1}$, and
$g_1(w)=g_0(w)^c$ for all $w\in\mu_{q'}$.
\end{lemma}

\begin{proof}
Every $w\in\mu_{\ell^f-1}$ is $\zeta^a$ (Lemma~\ref{pkj:roots}(iii)), and $g_1(\zeta^a)=g_1(\zeta)^a=
g_0(\zeta)^{ca}=g_0(\zeta^a)^c$. The root $y$ has order $q_1/\ell^{v_\ell}=q'$, prime to $\ell$, and lies
in $e_v(E)\subseteq E_v$; by Lemma~\ref{pkj:roots}(i) with $L=E_v$ (residue degree $f$), $q'\mid\ell^f-1$.
\end{proof}

\subsection{The three properties of the data}\label{pkj:ss-three}

\begin{lemma}[Letters determine the embedding off the conductor]\label{pkj:letter-det}
Let $\ell\nmid N_E$, $\zeta\in\auQlb$ of exact order $\ell^f-1$, and $g,g'\in\pkjG$ with $g(\zeta)=g'(\zeta)$.
Then $\sigma_g=\sigma_{g'}$; in particular $g$ is thick iff $g'$ is, and $\varepsilon(g)=\varepsilon(g')$.
\end{lemma}

\begin{proof}
By Lemma~\ref{pkj:transfer} (with $c=1$), $g(w)=g'(w)$ for all $w\in\mu_{\ell^f-1}$. Applying this to
$w=g'^{-1}(u)$ and to $w=g^{-1}(u)$ ($u\in\mu_{\ell^f-1}$, a $\pkjG$-stable set) shows that $g g'^{-1}$
and $g'g^{-1}$, and likewise $g'^{-1}g$ (apply $g'^{-1}$ to $g(w)=g'(w)$), fix $\mu_{\ell^f-1}$
pointwise. By Lemma~\ref{pkj:roots}(iv) they lie in $\Gamma$. In reading (R$^-$),
$\sigma_g=\sigma_{g'}$ means $g^{-1}(e_v x)=g'^{-1}(e_v x)$ for all $x$, i.e.\ $g'g^{-1}\in\Gamma$; in
reading (R$^+$) it means $g(e_vx)=g'(e_vx)$, i.e.\ $g'^{-1}g\in\Gamma$.
\end{proof}

\begin{lemma}[Ratios of places]\label{pkj:ratio}
Let $q'$ and $y$ be as in Lemma~\ref{pkj:transfer}, $g_0,g_1\in\pkjG$ and $c\in\Z$ with
$g_1(w)=g_0(w)^c$ for all $w\in\mu_{q'}$. Then $c$ is prime to $q'$ and, with $\varpi$ as in
Convention~\ref{pkj:conv},
\[
\pkjcl_{q'}\bigl(P(g_1)/P(g_0)\bigr)=\varepsilon(g_0)\,\varepsilon(g_1)\,c^{\varpi}\qquad\text{in }(\Z/q')^\times .
\]
\end{lemma}

\begin{proof}
$g_1(y)$ and $g_0(y)$ are primitive $q'$-th roots of unity and $g_1(y)=g_0(y)^c$, so $c$ is a unit
modulo $q'$. Write $\varepsilon_i=\varepsilon(g_i)$, $z=e_v(\zeta_{q_1})$, so $y=z^{\ell^{v_\ell}}$. By
\eqref{pkj:eq-sign} and \eqref{pkj:eq-ratio},
\begin{equation}\label{pkj:eq-places}
\sigma_{g_1}(\zeta_{q_1})^{\varepsilon_1}=\sigma_{g_0}(\zeta_{q_1})^{\varepsilon_0j_{10}},\qquad
j_{10}=P(g_1)/P(g_0).
\end{equation}
\emph{Reading (R$^+$).} Applying $\iota$ to \eqref{pkj:eq-places} gives $g_1(z)^{\varepsilon_1}=
g_0(z)^{\varepsilon_0j_{10}}$; raising to the power $\ell^{v_\ell}$, $g_1(y)^{\varepsilon_1}=g_0(y)^{\varepsilon_0j_{10}}$.
With $g_1(y)=g_0(y)^c$ and $g_0(y)$ primitive of order $q'$: $c\varepsilon_1\equiv\varepsilon_0j_{10}$, i.e.\
$j_{10}\equiv\varepsilon_0\varepsilon_1c\pmod{q'}$.

\emph{Reading (R$^-$).} Applying $\iota$ gives $g_1^{-1}(z)^{\varepsilon_1}=g_0^{-1}(z)^{\varepsilon_0j_{10}}$,
hence $g_1^{-1}(y)^{\varepsilon_1}=g_0^{-1}(y)^{\varepsilon_0j_{10}}$. Put $u=g_0^{-1}(y)\in\mu_{q'}$; then
$g_1(u)=g_0(u)^c=y^c$, so $g_1^{-1}(y)^c=u=g_0^{-1}(y)$. Hence $g_1^{-1}(y)^{\varepsilon_1}=
g_1^{-1}(y)^{c\varepsilon_0j_{10}}$, and $g_1^{-1}(y)$ is primitive of order $q'$: $\varepsilon_1\equiv
c\varepsilon_0j_{10}$, i.e.\ $j_{10}\equiv\varepsilon_0\varepsilon_1c^{-1}\pmod{q'}$.
\end{proof}

\begin{lemma}[One sign on the $b$-places at primes of the conductor]\label{pkj:signs}
Let $\ell\mid N_E$ and let $g,g'\in\pkjG$ be thick. Then $\varepsilon(g)=\varepsilon(g')$.
\end{lemma}

\begin{proof}
Suppose $\varepsilon(g)=-\varepsilon(g')$. By Lemma~\ref{pkj:sizes}(i), $\ell\mid q_0$ or $\ell=d$.

\emph{Case $\ell\mid q_0$.} By \eqref{pkj:eq-sign}, $\sigma_g(\sqrt{-d})=-\sigma_{g'}(\sqrt{-d})$. Applying
$\iota$ and then $g'$ (reading (R$^-$)), resp.\ $g'^{-1}$ (reading (R$^+$)), we get an element
$\theta\in\pkjG$ ($\theta=g'g^{-1}$, resp.\ $\theta=g'^{-1}g$) with $\theta(e_v\sqrt{-d})=-e_v\sqrt{-d}$. Since $d\equiv3
\pmod4$, $\beta=(1+\sqrt{-d})/2\in E$ satisfies $\beta(1-\beta)=(1+d)/4$. Put $x=e_v(\beta)\in\auQlb$;
then $\theta(x)=1-x$, so $|1-x|=|x|$ ($\theta$ is an isometry), and
$|x|^2=|x(1-x)|=|(d+1)/4|_\ell<1$, because $4\ell\mid d+1$ for every prime $\ell\mid q_0$
(\S\ref{an:ssec-param}). So $|x|<1$ and $|1-x|<1$, contradicting $1=|x+(1-x)|\le\max(|x|,|1-x|)$.

\emph{Case $\ell=d$.} Then $\ell\nmid q_1$ (as $d\nmid q_0$ and the primes of $q_1$ divide $q_0$), so
$q'=q_1$. By Lemma~\ref{pkj:roots}(ii), $g'g^{-1}$ acts on $\mu_{q_1}$ as $w\mapsto w^{\ell^c}$ for some
$c\ge0$; as $g(\mu_{q_1})=\mu_{q_1}$, $g'(w)=g'g^{-1}(g(w))=g(w)^{\ell^c}$ for $w\in\mu_{q_1}$. By
Lemma~\ref{pkj:ratio} (with $g_0=g$, $g_1=g'$), $\pkjcl_{q_1}(P(g')/P(g))=\varepsilon(g)\varepsilon(g')
d^{\varpi c}=-d^{\varpi c}$. Reducing modulo $q_0\mid q_1$, and using $\pkjcl_{q_0}(P(g))=\pkjcl_{q_0}(P(g'))
=-1$ (both are $b$-places, (D5)), we get $1=-d^{\varpi c}$ in $(\Z/q_0)^\times$, hence $d^{c}\equiv-1
\pmod{q_0}$. But $d\equiv q_0/2-1\pmod{q_0}$ and $4\mid q_0$ give $d^2\equiv q_0\cdot\frac{q_0}4-q_0+1
\equiv1$, so $d^c\equiv1$ or $d^c\equiv d\equiv q_0/2-1\pmod{q_0}$, and neither is $-1$ (as $q_0>2$ and
$q_0/2\not\equiv0$). Contradiction.
\end{proof}

\begin{lemma}[No three thick letters on a short arc]\label{pkj:arc}
Let $\ell\le Q$, $\zeta\in\auQlb$ of exact order $\ell^f-1$, $a,a'\in[1,\pkilog_\ell h^3]$, and
$g_0,g_1,g_2\in\pkjG$ with $g_1(\zeta)=g_0(\zeta)^{\ell^a}$ and $g_2(\zeta)=g_0(\zeta)^{\ell^{a+a'}}$. Then
$g_0$, $g_1$, $g_2$ are not all thick.
\end{lemma}

\begin{proof}
Suppose they are; put $P_i=P(g_i)\in\pkjPb$ and let $q'$ be as in Lemma~\ref{pkj:transfer}. By
Lemma~\ref{pkj:transfer}, $g_1(w)=g_0(w)^{\ell^a}$ and $g_2(w)=g_0(w)^{\ell^{a+a'}}$ for $w\in\mu_{q'}$;
by Lemma~\ref{pkj:ratio}, with signs $\pm$ (each $+1$ or $-1$),
\[
\pkjcl_{q'}(P_1/P_0)=\pm\ell^{\varpi a},\qquad \pkjcl_{q'}(P_2/P_0)=\pm\ell^{\varpi(a+a')} .
\]
If $\varpi=1$, take $\sigma=P_0$, $\tau=P_1$, $\rho=P_2$, $i=a$, $j=a+a'$: then $\pkjcl_{q'}(\tau/\sigma)=
\pm\ell^i$ and $\pkjcl_{q'}(\rho/\sigma)=\pm\ell^j$. If $\varpi=-1$, take $\sigma=P_2$, $\tau=P_1$,
$\rho=P_0$, $i=a'$, $j=a+a'$: then $\pkjcl_{q'}(\rho/\sigma)=\pkjcl_{q'}(P_2/P_0)^{-1}=\pm\ell^{j}$ and
$\pkjcl_{q'}(\tau/\sigma)=\pkjcl_{q'}(P_1/P_0)\,\pkjcl_{q'}(P_2/P_0)^{-1}=\pm\ell^{a+a'-a}=\pm\ell^{i}$.
In both cases $\sigma,\tau,\rho\in\pkjPb$ and $0<i<j\le2\pkilog_\ell h^3<l_0(s)$ by
Lemma~\ref{pkj:sizes}(iv). Since $\ell\le Q=5r(r+1-m)$ and $q'$ is the largest divisor of $q_1$
prime to $\ell$, this contradicts the sparseness of $\pkjPb$ ((D5), Definition~\ref{an:def-sparse}).
\end{proof}

\subsection{Proof of Proposition~\ref{pki:chart}}\label{pkj:ss-proof}

\begin{theorem}[Proposition~\ref{pki:chart}]\label{pkj:chart}
Let $Q=kt$, $h=r^8$, and let thickness and $\varepsilon(\gamma)$ be as in Convention~\ref{pkj:conv}
(either reading). For every prime $\ell$, every $v\mid\ell$, every large $\auO$ and every
$\zeta\in\auO$ of exact order $\ell^f-1$ (such $\zeta$ exist), there are a set $T$ of letters and a sign
$s$ on letters for which the letters read thickness for $(Q,h)$ (Definition~\ref{pki:def-reads}).
\end{theorem}

\begin{proof}
Elements $\zeta\in\auO$ of exact order $\ell^f-1$ exist by Lemma~\ref{pkj:roots}(iii), as
$\auO\supseteq\auO_{E_v}$. For an embedding $\gamma\colon\auO\to\C_\ell$ choose $g_\gamma\in\pkjG$
restricting to $\gamma$ (Convention~\ref{pkj:conv}); then $\gamma(\zeta)=g_\gamma(\zeta)$,
$\sigma_\gamma=\sigma_{g_\gamma}$, and the letter $\gamma(\zeta)$ is an $(\ell^f-1)$-th root of unity. Put
\[
T=\{\gamma(\zeta):\ \gamma\ \text{thick}\}.
\]
\emph{Case $\ell\nmid N_E$.} For a letter $z$ let $s(z)=\varepsilon(\gamma)$ if $z=\gamma(\zeta)$ for some
$\gamma$, and $s(z)=1$ otherwise; by Lemma~\ref{pkj:letter-det} (applied to $g_\gamma$, $g_{\gamma'}$),
$\gamma(\zeta)=\gamma'(\zeta)$ implies $\sigma_\gamma=\sigma_{\gamma'}$, so $s$ is well defined. For every
$\gamma$: if $\gamma$ is thick, $\gamma(\zeta)\in T$ by definition; if $\gamma(\zeta)\in T$, then
$\gamma(\zeta)=\gamma'(\zeta)$ with $\gamma'$ thick, so $\sigma_\gamma=\sigma_{\gamma'}$ and $\gamma$ is thick;
and $\varepsilon(\gamma)=s(\gamma(\zeta))$. Moreover $e=1$ (Lemma~\ref{pkj:roots}(iv)) and $\ell\ge13$
(Lemma~\ref{pkj:sizes}(ii)).

\emph{Case $\ell\mid N_E$.} Then $\ell\le Q$ (Lemma~\ref{pkj:sizes}(i)). Let $s$ be the constant sign
$\varepsilon(\gamma_0)$ for a thick $\gamma_0$ if one exists, and $s\equiv1$ otherwise. For thick
$\gamma$: $\gamma(\zeta)\in T$ (so ``$\gamma$ thick $\Leftrightarrow\gamma(\zeta)\in T$'' holds, both sides
being true), and $\varepsilon(\gamma)=\varepsilon(\gamma_0)=s(\gamma(\zeta))$ by Lemma~\ref{pkj:signs}. Since
$\ell\le Q$, the clause ``for every $\gamma$ if $\ell>Q$'' is void, and $\ell\le Q$ with $s$ constant.

\emph{The remaining clauses, in both cases.} $e\le h$ by Lemma~\ref{pkj:sizes}(iii). Let $\ell\le Q$ and
suppose a letter $z$ has $z$, $z^{\ell^a}$, $z^{\ell^{a+a'}}\in T$ with $a,a'\in[1,\pkilog_\ell h^3]$:
$z=\gamma_0(\zeta)$, $z^{\ell^a}=\gamma_1(\zeta)$, $z^{\ell^{a+a'}}=\gamma_2(\zeta)$ with $\gamma_0,\gamma_1,
\gamma_2$ thick. Then $g_i=g_{\gamma_i}$ are thick, $g_1(\zeta)=g_0(\zeta)^{\ell^a}$ and
$g_2(\zeta)=g_0(\zeta)^{\ell^{a+a'}}$, contradicting Lemma~\ref{pkj:arc}.
\end{proof}

\begin{remark}[The uses of Proposition~\ref{pki:chart}]\label{pkj:uses}
Proposition~\ref{pki:chart} is used in \S\ref{pki:ss-deduce} (Lemma~\ref{pki:letterwin},
Theorem~\ref{pki:crys} and the proof of Theorem~\ref{pll:exponent}), only through the clauses of
Definition~\ref{pki:def-reads} for the thickness and signs of Convention~\ref{pkj:conv}; and
\S\ref{pkd:ss-data} (Theorem~\ref{pkd:data}) uses the same $\sigma_\gamma$, $\varepsilon(\gamma)$ and
thickness, through the jumps $J_{\sigma_\gamma}$. Theorem~\ref{pkj:chart} gives these clauses for every
large $\auO$ and every admissible choice of $\zeta$, as the deduction requires; the quantities $e$, $f$
are those of $E_v$, as in \S\ref{pki:ss-arch} (with $\Gamma'=\Gamma$ in Definition~\ref{pkp:def-chart},
Lemma~\ref{pkp:chart-exists}, they are also those of the Kummer chart). The theorem holds for both
readings of $\sigma_\gamma$, so it applies whichever reading \S\ref{pkd:ss-data} and the Hodge--Tate
convention of $L_{\mathrm{main}}$ adopt, as long as the two agree; both use (R$^-$) (\S\ref{pll:ss-statement}, \S\ref{pkd:ss-data}).
\end{remark}


%% ===== proofs/C-kisin-cases.tex =====
% proofs/C-kisin-cases.tex
% The four cases of the local-loss engine: Propositions pki:generic, pki:tame, pki:sparse, pki:any,
% and the remaining assertions of pki:data. Labels start with pkc:.
% Conventions and macros as in proofs/C-kisin.tex, proofs/C-kisin-pairing.tex,
% proofs/C-kisin-modules.tex, proofs/C-kisin-inputs.tex.
\providecommand{\auO}{\mathcal O}
\providecommand{\auQlb}{\overline{\Q}_\ell}
\providecommand{\pllZf}{Z_f}
\providecommand{\pllrk}{\operatorname{rk}}
\providecommand{\plldef}{\delta^{\vee}}
\providecommand{\pglen}{\operatorname{len}}
\providecommand{\pkiWin}{\operatorname{Win}}
\providecommand{\pkiC}{\mathcal C}
\providecommand{\pkilog}{\operatorname{lg}}
\providecommand{\pkpA}{A_{\mathrm{inf}}}
\providecommand{\pkpAmb}{\mathbb A}
% macros of this file (prefix \pkc)
\providecommand{\pkcPt}{\mathrm{Pt}}
\providecommand{\pkcimp}{\mathrm{imp}}
\providecommand{\pkcpd}{\mathrm{pd}}
\providecommand{\pkcdep}{\mathrm{dep}}
\providecommand{\pkcsh}{\mathrm{sh}}
\providecommand{\pkcNM}{\mathrm{NM}}
\providecommand{\pkcnw}{\mathrm{nw}}
\providecommand{\pkcSig}{\Sigma}

\section{The four cases of the local-loss engine}\label{pkc:sec}

This section proves the core propositions of \S\ref{pki:ss-arch}, one per subsection, in the order
Proposition~\ref{pki:generic} (\S\ref{pkc:ss-generic}), Proposition~\ref{pki:tame}
(\S\ref{pkc:ss-tame}), Proposition~\ref{pki:sparse} (\S\ref{pkc:ss-sparse}),
Proposition~\ref{pki:any} (\S\ref{pkc:ss-any}), and then treats the remaining assertions of
Proposition~\ref{pki:data} (\S\ref{pkd:ss-data}, Theorem~\ref{pkd:data}; the membership $L_c^*\in\pkiC(U',K^{-1})$ needed in
case (a) uses the input (I6), which is Lemma~\ref{pkw:dual}). The tools are those of \S\ref{pkp:sec},
\S\ref{pkm:sec} and \S\ref{pkn:sec}, which we cite by label and do not reprove. Where a case needs a
statement about points, lift ideals or charts that is not among the statements of those
sections, we state the exact form needed in \S\ref{pkc:ss-inputs} and prove the case from it; these
statements, the inputs (I1)--(I4) and (T1)--(T4), are proved in \S\S\ref{pkt:sec}--\ref{pkw:sec}.

\emph{Dependencies.} Proved in full from the statements of \S\S\ref{pkp:sec}--\ref{pkn:sec} and the inputs
(I1), (I2) of \S\ref{pkc:ss-inputs}: Proposition~\ref{pki:generic}(i). Proved from these and the
inputs (I3), (I4): Proposition~\ref{pki:generic}(ii). Proved from (I1), (I2) and the tame inputs
(T1)--(T4): Proposition~\ref{pki:tame}. The inputs are proved elsewhere: (I1) and (I3) in
\S\ref{pkt:sec} (Propositions~\ref{pkt:I1}, \ref{pkt:I3}), (I2) in \S\ref{pke:sec}
(Theorem~\ref{pke:I2}), (I4) in \S\ref{pkh:sec} (Theorem~\ref{pkh:I4}), (T1)--(T4) in \S\ref{pkw:sec}.
Propositions~\ref{pki:sparse}, \ref{pki:any} and the rest of Proposition~\ref{pki:data} are proved in
\S\ref{pkd:sec}.

\emph{Conventions.} As in \S\ref{pkp:sec}: a Kummer chart of $\Gamma=D_v$ is fixed
(Definition~\ref{pkp:def-chart}), with ramification index $e$, residue degree $f$, subgroup $\Gamma'$
and $d_v$ as there; letters, the Frobenius $\sigma(i)=i^\ell$ on letters and $i^{(-a)}=\sigma^{-a}(i)$
are as in Definition~\ref{pkp:def-letters}; orders $o\colon\Z_{\ge0}\times\Z\to\Z$, the modules
$L_i(o)$, and $\rho_t$, $\mathrm{std}_l$, $\pkcsh_{l,n}$ are as in Definition~\ref{pkp:def-orders} with
the modules $\mathcal N$ of Definition~\ref{pkn:def-N}; sides, extensions $X_S$, pairings modulo
$\ell^m$, Tate-dual sides $T_P$, safe moves and $\mathrm{SCC}_N$ are as in
Definitions~\ref{pkp:def-side}--\ref{pkp:def-scc}; $\pkcimp$ and $\pkcpd$ are as in
\S\ref{pkm:ss-state}, and $\pkcdep_{t,\eta,l}$, ``standard point'' and ``stop'' as in
Definition~\ref{pkp:def-points}. For a side $S$, a letter $i$, $c\in C_S$, $e\ge0$ and orders $o,o'$
we write
\[
\pkcPt_S(c,i,e;\,o,o')
\]
for the set of points of $S$ over $\ell^ec$ at $i$ with orders $o$ and depth $o'$
(Definition~\ref{pkp:def-points}). All inequalities between orders are required only at the pairs
$(a,u)$ with $u\in U(i^{(-a)})$, the jump set of the side at the letter $a$ steps before $i$; ``at
level $a$'' refers to the first argument. Thus, explicitly,
\[
\rho_t(a,u)=t[a=0]-u,\qquad \pkcdep_{t,\eta,1}(a,u)=-u+(t+\eta)[a=0],\qquad
\pkcpd_{t,\eta,1,\eta_0}(a,u)=-u+(t+\eta)[a=0]+\eta_0[a=1].
\]

\subsection{Windows, steps, and the inputs used}\label{pkc:ss-inputs}

\subsubsection*{Elementary facts}

\begin{remark}[Monotonicity]\label{pkc:mono}
If $o_1\le o$ and $d_1\le d$ (at the jumps), then $\pkcPt_S(c,i,e;o,d)\subseteq\pkcPt_S(c,i,e;o_1,d_1)$,
since $\mathrm{Fil}^n$ decreases in $n$ and so $L_i(o)\subseteq L_i(o_1)$. In particular, as
$\pkcpd_{t,\eta,1,\eta_0}\ge\pkcdep_{t,\eta,1}$, a point with orders $\rho_{-\eta}$ (or larger) and
depth $\pkcpd_{t,\eta,1,\eta_0}$ (or larger) is a standard point for $(t,\eta,1)$.
\end{remark}

For $n\ge0$ and $b\in\Z$ let $\pkcSig_n(b)=\{y\in\Z:\ |y|\le n\text{ or }|y-b|\le n\text{ or
}|y+b|\le n\}$; it has at most $3(2n+1)$ elements, and $\pkcNM(\ell,n,b)$ (Definition preceding
Definition~\ref{pki:def-reads}) says exactly that no nonzero element of $\pkcSig_n(b)$ is divisible by
$\ell$.

\begin{lemma}[Windows]\label{pkc:win}
Let $\eta,b,W\ge0$, a sign $\varepsilon$, $u\in\pkiWin(\eta,-\varepsilon,b,W)$ and
$u'\in\pkiWin(\eta,\varepsilon,b,W)$. Then
\begin{enumerate}
\item[(i)] $u\le b+W+\eta$ and $u'\le b+W+\eta$;
\item[(ii)] $-(u+u')\le b+W+2\eta$, and if $\eta\le b$ then $-u-u'-1\le b+W+\eta$;
\item[(iii)] if $s,t'\ge0$, $n\ge\max(s,t')+\eta+W$ and $-u-s\le y<t'-u$, then
$y\in\pkcSig_n(b)$; the same holds for $u'$.
\end{enumerate}
\end{lemma}

\begin{proof}
(i) Elements of the low window are $\le\eta$, those of a far window are $\le b+W$. (ii) If both
$u,u'$ are in the low window, $-u-u'\le2\eta$, which is $\le b+W+2\eta$, and $\le b+\eta$ if $\eta\le b$.
If exactly one is in a far window, say $u'$, then $-u'\le b+W$ (if the far window of $u'$ is negative)
or $-u'\le-b\le0$ (if positive), and $-u\le\eta$; so $-u-u'\le b+W+\eta$. If both are in far windows,
the windows have opposite signs, so $u=\delta x$, $u'=-\delta x'$ with $\delta=\pm1$, $x,x'\in[b,b+W]$,
and $|u+u'|=|x-x'|\le W$. In all cases $-u-u'\le b+W+2\eta$, and if $\eta\le b$ then $-u-u'\le
b+W+\eta$, a fortiori $-u-u'-1\le b+W+\eta$. (iii) If $|u|\le\eta$
then $-\eta-s\le y<t'+\eta$, so $|y|\le\max(s,t')+\eta\le n$. If $u\in[b,b+W]$ then $-W-s\le
y+b<t'$, so $|y+b|\le n$. If $-u\in[b,b+W]$ then $-s\le y-b<W+t'$, so $|y-b|\le n$.
\end{proof}

\begin{lemma}[Jumps of the Tate-dual side]\label{pkc:dualjumps}
Let $P$ pair sides $S,S'$, and at a letter $i$ let $U_S(i)\subseteq\pkiWin(\eta,-\varepsilon,b,W)$ and
$U_{S'}(i)\subseteq\pkiWin(\eta,\varepsilon,b,W)$ with $\eta\le b$. Then every jump of $T_P$ at $i$ is
$\le b+W+\eta$. If moreover $6(a+\eta+W)<b$, $a\ge0$, $\eta\ge1$, then every jump $s$ of $T_P$ at $i$
satisfies $-a-s\le y<a-s\Rightarrow y\in\pkcSig_{6(a+\eta+W)}(b)$.
\end{lemma}

\begin{proof}
The jumps of $T_P$ at $i$ are $0$ and the $s=-u-u'-1$ with $u\in U_S(i)$, $u'\in U_{S'}(i)$
(Definition~\ref{pkp:def-pairing}); the first claim is Lemma~\ref{pkc:win}(ii). For the second put
$n=6(a+\eta+W)$. If $s=0$, $|y|\le a\le n$. Otherwise the case analysis of Lemma~\ref{pkc:win}(ii)
shows: if both $u,u'$ are low, $|s|\le2\eta+1$; if both are far, $|s|\le W+1$; if exactly one is far,
$s=\mp x-v-1$ with $x\in[b,b+W]$ and $|v|\le\eta$, so $|s\pm b|\le W+\eta+1$. Hence $s$ lies within
$m:=2\eta+W+1$ of $0$, $b$ or $-b$, and $y\in[-a-s,a-s)$ lies within $a+m\le n$ (as $\eta\ge1$) of
$0$, $-b$ or $b$.
\end{proof}

\subsubsection*{The inputs}

The cases use, besides the statements of \S\S\ref{pkp:sec}--\ref{pkn:sec} cited by label, the following
statements. (I1) is the form of Lemma~\ref{pkp:tau}(iii) that the cases need. The four inputs are proved
elsewhere: (I1) is Proposition~\ref{pkt:I1}, (I2) is Theorem~\ref{pke:I2}, (I3) is
Proposition~\ref{pkt:I3}, (I4) is Theorem~\ref{pkh:I4}. For $n\in\Z$ let $v_\ell(n)$ be the $\ell$-adic valuation, with
$v_\ell(0)=0$, and put $\mathrm{cost}(a,n)=v_\ell(n)+d_v+3a$.

\begin{enumerate}
\item[(I1)] \emph{Points from killers.} If $\ell^\kappa E^t$ kills the cokernel of a side $S$ at $i$
(Definition~\ref{pkm:def-J}), then for every $c\in C_S$ the set $\pkcPt_S(c,i,\kappa;\rho_{-t},\rho_{-t})$
is nonempty.
\item[(I2)] \emph{Deepening.} Let $x\in\pkcPt_S(c,i,e;o,d)$ and let $I$ be a finite set of pairs
$(a,n)\in\Z_{\ge0}\times\Z$ containing every $(a,n)$ with $d(a,u)\le n<d'(a,u)$ for some $u\in
U(i^{(-a)})$. If $e+\sum_{(a,n)\in I}\mathrm{cost}(a,n)\le e'$, then
$\pkcPt_S(c,i,e';\pkcimp(o,d'),d')\ne\emptyset$.
\item[(I3)] \emph{The rule with a dual point.} Let $P$ pair $S,S'$, $l,\kappa\ge0$, and let $\nu$ be a
point of $T_P$ over $\ell^\kappa\cdot1$ at $i^{(-l)}$ with orders $o_\nu$, where $o_\nu(a,\cdot)$ is
nonincreasing for each $a$. Let $w\in L_i(o_w)$ (for $X_S$) lie over $0$ and let $y$ be a point of
$S'$ over some $\ell^{e'}c'$ at $i$ with depth $d_y$. If for all $a\ge0$, $u\in U_S(i^{(-a)})$,
$u'\in U_{S'}(i^{(-a)})$
\[
0<\pkcsh_{l,0}(o_\nu)(a,-u-u'-1)+o_w(a,u)+d_y(a,u'),
\]
then the move of the first side by $w$ with $y$ at rest is $\kappa$-safe.
\item[(I4)] \emph{Few low jumps: the deep constants.} Let the chart be unramified, $\ell\ge5$,
$a,\eta,W\ge0$, $\eta\ge1$, and for each letter $i$ let $b_i\ge0$ and a sign $\varepsilon_i$ satisfy
$b_i+W+\eta\le(\ell-3)a$, $6(a+\eta+W)<b_i$, $\pkcNM(\ell,6(a+\eta+W),b_i)$. Let $S'$ be a side with
$U_{S'}(i)\subseteq\pkiWin(\eta,\varepsilon_i,b_i,W)$, and let $n_i$, $r_i$ be integers with $n_i\ge
f\sum_{u\in[-\eta,\eta]}m_i(u)$ and $r_i\ge f\sum_{u\in U_{S'}(i)}m_i(u)$, where $m_i(u)$ is the sum
over the embeddings $\gamma$ of letter $i$ of the multiplicity of the jump $u$ of the lattice of $S'$
at $\gamma$. Define the depth
\[
\delta_i=\rho_{2a+\eta}\ \text{if }\varepsilon_i=1,\qquad
\delta_i(a',u)=\pkcnw_i(a',u)=\begin{cases}b_i+W+a+\eta&a'=0\\ W+1-u+2a[|u|\le\eta]&a'=1\\
-u&a'\ge2\end{cases}\ \text{if }\varepsilon_i=-1 .
\]
Then the set $V$ of $c'\in C_{S'}$ such that $\pkcPt_{S'}(c',i,0;\rho_0,\delta_i)\ne\emptyset$ for
every letter $i$ is a $\Z_\ell$-submodule of $C_{S'}$ with
$\pglen_{\Z_\ell}(C_{S'}/V)\le11\sum_i\bigl(n_i(a+1)+r_i(W+1)\bigr)$.
\end{enumerate}

\begin{lemma}[Free steps]\label{pkc:free}
Let $x\in\pkcPt_S(c,i,e;o,d)$, $n\ge0$, $b\in\Z$ with $\pkcNM(\ell,n,b)$, and let $d'$ be orders such
that every pair $(a,y)$ with $d(a,u)\le y<d'(a,u)$ for some $u\in U(i^{(-a)})$ has $a=0$ and $y\in
\pkcSig_n(b)$. Then $\pkcPt_S(c,i,e+3(2n+1)d_v;\pkcimp(o,d'),d')\ne\emptyset$. If the chart is
unramified, $\pkcPt_S(c,i,e;\pkcimp(o,d'),d')\ne\emptyset$.
\end{lemma}

\begin{proof}
Apply (I2) with $I=\{0\}\times\pkcSig_n(b)$. For $y\in\pkcSig_n(b)$, $v_\ell(y)=0$: if $y=0$ by
convention, otherwise by $\pkcNM(\ell,n,b)$. So $\mathrm{cost}(0,y)=d_v$ and $\sum_I\mathrm{cost}\le
3(2n+1)d_v$. If the chart is unramified, $d_v=0$.
\end{proof}

\begin{lemma}[From sides to lattices]\label{pkc:cup}
Let $\mathcal C$ be a class of lattices $Y$, each crystalline on $\Gamma'$ with jumps at each embedding
$\gamma$ in a set $U(\gamma)$, and let $L$ be crystalline on $\Gamma'$ with jumps in $U'(\gamma)$. For a
lattice $Y$ with jump sets $U$ let $S(Y)$ be its \emph{universal side}: lattice $Y$, $C_{S(Y)}=
\pllZf(Y)$ with trivial action, $\alpha=\mathrm{id}$, and jump set $U_i=\{0\}\cup\bigcup_{\gamma}U(\gamma)$
at the letter $i$ (union over the embeddings $\gamma$ of letter $i$). Let $N\ge0$ and $x\in\pllZf(L)$.
If $\mathrm{SCC}_N(c,x)$ holds for every $Y\in\mathcal C$, every pairing modulo $\ell^m$ (any $m$) of
$S(Y)$ and $S(L)$ and every $c\in\pllZf(Y)$, then $\ell^Nx\in L^{\perp\mathcal C}$.
In particular, if this holds for all $x\in\pllZf(L)$, then $\mathcal C$ has cup exponent $N$ on $L$,
and if it holds with $N=0$ for all $x$ in an $\auO$-submodule (or $\Z_\ell$-submodule) $V$ of
$\pllZf(L)$, then $\plldef_{\mathcal C}(L)\le\pglen(\pllZf(L)/V)$.
\end{lemma}

\begin{proof}
$\pllZf(Y)$ and $\pllZf(L)$ are finite free (Lemma~\ref{pll:exp}), and $U_i$ contains $0$ and all
jumps at the embeddings of letter $i$, so $S(Y)$, $S(L)$ are sides (Definition~\ref{pkp:def-side}).
Let $Y\in\mathcal C$, $m\ge1$, $(\psi,c_\psi)$ a pairing modulo $\ell^m$ of $Y$ and $L$ in the sense of
Definition~\ref{pll:def}(4) (with $(Y\otimes L)^*=\operatorname{Hom}(Y\otimes L,\auO(1))$, as the
values in $\Z/\ell^m(1)$ require), and $a\in\pllZf(Y)$. The same data is a pairing $P$ modulo $\ell^m$
of $S(Y)$ and $S(L)$ in the sense of Definition~\ref{pkp:def-pairing}, with $P=b_\psi$. Put
$\beta(g,h)=b_\psi(a(g),g\,\ell^Nx(h))$, the cup product of the continuous cocycles $a$ and $\ell^Nx$
through the $\Gamma$-equivariant pairing $b_\psi$; it is a continuous $2$-cocycle with values in
$\Z/\ell^m(1)\cong\mu_{\ell^m}$. By $\mathrm{SCC}_N(a,x)$ its image in $C^2(\Gamma,\auQlb^\times)$ is the
coboundary of a continuous cochain. The Kummer sequence $1\to\mu_{\ell^m}\to\auQlb^\times
\xrightarrow{\ \ell^m\ }\auQlb^\times\to1$ and $H^1(\Gamma,\auQlb^\times)=0$ (Hilbert's Theorem~90 for
the absolute Galois group of $E_v$) show that $H^2(\Gamma,\mu_{\ell^m})\to H^2(\Gamma,\auQlb^\times)$ is
injective. So the class of $\beta$ vanishes, i.e.\ $\beta$ is a coboundary in $C^2(\Gamma,
\Z/\ell^m(1))$. As $Y,m,\psi,a$ were arbitrary, $\ell^Nx\in L^{\perp\mathcal C}$ (Definition~\ref{pll:def}(5)).
The last two assertions follow from the definitions of cup exponent and of $\plldef_{\mathcal C}$:
in the second case $V\subseteq\pllZf(L)\cap L^{\perp\mathcal C}$.
\end{proof}

\subsubsection*{Two walks}

The two lemmas below are the common core of Propositions~\ref{pki:generic}(i) and \ref{pki:tame}.
Throughout, $P$ pairs sides $S$, $S'$, and at every letter $i$ an integer $b_i\ge0$ and a sign
$\varepsilon_i$ are given with
\begin{equation}\label{pkc:eq-win}
U_S(i)\subseteq\pkiWin(\eta,-\varepsilon_i,b_i,W),\qquad U_{S'}(i)\subseteq\pkiWin(\eta,\varepsilon_i,
b_i,W),\qquad \pkcNM(\ell,n,b_i),\qquad n=6(a+\eta+W),
\end{equation}
for fixed integers $a,\eta,W\ge0$. If $\varepsilon_i=1$, every jump of $S$ at $i$ is $\le\eta$ and
every jump of $S'$ at $i$ is $\ge-\eta$; if $\varepsilon_i=-1$, the reverse. (This is how a far
window is ``negative'' or ``positive''.)

\begin{lemma}[The walk with a gap below zero]\label{pkc:gap}
Assume \eqref{pkc:eq-win}, that for every letter $i$ the element $\ell^\kappa E^a$ kills the cokernels of
$S$, $S'$ and $T_P$ at $i$, and that no jump of $S$ or $S'$ at any letter lies in $[-a,0)$. Then
$\mathrm{SCC}_N(c,c')$ for all $c\in C_S$, $c'\in C_{S'}$ and all $N\ge3\kappa+2\cdot3(2n+1)d_v$.
\end{lemma}

\begin{proof}
Put $\sigma=3(2n+1)d_v$ and $e=\kappa+\sigma$.

\emph{Points without pole.} Fix $i$. By (I1), $\pkcPt_S(c,i,\kappa;\rho_{-a},\rho_{-a})\ne\emptyset$.
Let $d'=\pkcpd_{a,\eta,1,0}$. A pair $(a',y)$ with $\rho_{-a}(a',u)\le y<d'(a',u)$ has $a'=0$ (for
$a'\ge1$ both sides are $-u$) and $-u-a\le y<a+\eta-u$, so $y\in\pkcSig_n(b_i)$ by
Lemma~\ref{pkc:win}(iii) (with $s=a$, $t'=a+\eta$, and $a+\eta+\eta+W\le n$). By
Lemma~\ref{pkc:free} there is $x(i)\in\pkcPt_S(c,i,e;\pkcimp(\rho_{-a},d'),d')$. We claim
$\pkcimp(\rho_{-a},d')\ge\rho_0$ at the jumps. At level $a'\ge1$: $\rho_{-a}(a',u)=-u$; if $-u>0$ the
improved order is $d'(a',u)=-u$, otherwise $\min(-u,0)=-u$. At level $0$: if $\rho_{-a}(0,u)=-a-u>0$,
the improved order is $d'(0,u)=a+\eta-u\ge-u$; otherwise $u\ge-a$, so $u\ge0$ by the gap, and
$\min(a+\eta-u,0)\ge-u$. Hence $x(i)\in\pkcPt_S(c,i,e;\rho_0,\pkcpd_{a,\eta,1,0})$ (Remark~\ref{pkc:mono}).
In the same way, with $\ell^\kappa E^a$ killing the cokernel of $S'$, there is
$x'(i)\in\pkcPt_{S'}(c',i,e;\rho_0,\pkcpd_{a,\eta,1,0})$. By Remark~\ref{pkc:mono}, $x(i)$ and $x'(i)$
are standard points for $(a,\eta,1)$ (as $\rho_0\ge\rho_{-\eta}$).

\emph{Stops.} Fix $i$. If $\varepsilon_i=1$, the jumps of $S$ at $i$ are $\le\eta$ and those of $S'$
at $i$ are $\ge-\eta$. Apply Lemma~\ref{pkm:plain} with $l=1$, $t=a$, $\eta$, $\eta_0=0$, the killer
$\ell^\kappa E^a$ of $T_P$ at $i$, $x=x(i)$ (orders $\rho_0\ge\rho_{-\eta}$, depth
$\pkcpd_{a,\eta,1,0}$), $x'=x'(i)$ (orders $\rho_0\ge\rho_{-\eta}$) and $x'_0=x'(i^{(-1)})$ (orders
$\rho_0=\rho_{-\eta_0}$): $x(i)$ is a stop for $(\eta,1)$ and the move of the second side by
$\varphi x'(i^{(-1)})-x'(i)$ with $x(i)$ at rest is $\kappa$-safe. If $\varepsilon_i=-1$, the same
lemma with the roles of $S$, $S'$ exchanged gives that $x'(i)$ is a stop and the move of the first side
by $\varphi x(i^{(-1)})-x(i)$ with $x'(i)$ at rest is $\kappa$-safe. So (St) of
Theorem~\ref{pkm:walk} holds at every letter with $l(i)=1$.

\emph{The walk.} Corollary~\ref{pkm:walk-all} with $t=a$, $\eta$, $e=e'=\kappa+\sigma$ and the killer
$\ell^\kappa E^a$ of $T_P$ gives $\mathrm{SCC}_N(c,c')$ for $N\ge e+e'+\kappa=3\kappa+2\sigma$.
\end{proof}

\begin{lemma}[The walk with nonpositive jumps on one side]\label{pkc:nonpos}
Let the chart be unramified. Assume \eqref{pkc:eq-win}, that $\ell^\kappa E^a$ kills the cokernels of $S$
and of $T_P$ at every letter, and that $U_{S'}(i)\subseteq[-\eta,0]$ for every letter $i$. Then
$\mathrm{SCC}_N(c,c')$ for all $c,c'$ and all $N\ge2\kappa$.
\end{lemma}

\begin{proof}
Put $t=a+\eta$; note $a+t+\eta+W\le n$ and $a+t+\eta\le n$. As the chart is unramified, $d_v=0$, and
Lemma~\ref{pkc:free} costs nothing.

\emph{Points of $S$.} By (I1), $\pkcPt_S(c,i,\kappa;\rho_{-a},\rho_{-a})\ne\emptyset$. Let
$d'=\pkcpd_{a,t,1,0}$. Steps occur only at level $0$, with $-u-a\le y<a+t-u$, hence in
$\pkcSig_n(b_i)$ by Lemma~\ref{pkc:win}(iii) ($s=a$, $t'=a+t$, $a+t+\eta+W\le n$). Lemma~\ref{pkc:free}
gives $x(i)\in\pkcPt_S(c,i,\kappa;\pkcimp(\rho_{-a},d'),d')$. At level $\ge1$ the improved orders are
$-u$; at level $0$ they are $a+t-u\ge-t-u$ if $u<-a$, and $\min(a+t-u,0)\ge-t-u$ if $u\ge-a$ (as then
$u\ge-t$). So $x(i)\in\pkcPt_S(c,i,\kappa;\rho_{-t},\pkcpd_{a,t,1,0})$.

\emph{Points of $S'$.} All jumps of $S'$ are $\le0$, so $1$ kills the cokernel of $S'$ at every letter
(Proposition~\ref{pkm:kill-small} with $\omega=0$: $e\cdot0\le\ell-2$). By (I1),
$\pkcPt_{S'}(c',i,0;\rho_0,\rho_0)\ne\emptyset$. Steps to $d'=\pkcpd_{a,t,1,0}$ are at level $0$ with
$-u\le y<a+t-u$ and $u\in[-\eta,0]$, so $0\le y<a+t+\eta\le n$ and $y\in\pkcSig_n(b_i)$. We get
$x'(i)\in\pkcPt_{S'}(c',i,0;\pkcimp(\rho_0,d'),d')$; at level $0$ the improved order is
$a+t-u\ge-u$ if $u<0$ and $\min(a+t,0)=0=-u$ if $u=0$, and at level $\ge1$ it is $-u$. So
$x'(i)\in\pkcPt_{S'}(c',i,0;\rho_0,\pkcpd_{a,t,1,0})$.

Both $x(i)$ and $x'(i)$ are standard points for $(a,t,1)$ (Remark~\ref{pkc:mono}).

\emph{Stops.} If $\varepsilon_i=1$: the jumps of $S$ at $i$ are $\le\eta\le t$, those of $S'$ are
$\ge-\eta\ge-t$. Lemma~\ref{pkm:plain} (with $l=1$, killer exponent $a$, radius $t$, $\eta_0=0$,
$x=x(i)$, $x'=x'(i)$, $x'_0=x'(i^{(-1)})$ of orders $\rho_0$) gives (St).

If $\varepsilon_i=-1$: the jumps of $S$ at $i$ are $\ge-\eta\ge-t$ and those of $S'$ are $\le0\le t$.
We construct a point $q'$ of $S'$ over $c'$ at $i$ with orders $\ge\rho_{-t}$ and depth
$\pkcpd_{a,t,1,t}$. By Lemma~\ref{pkp:tau}(v) in the form of \S\ref{pkn:ss-def},
$\varphi x'(i^{(-1)})\in\pkcPt_{S'}(c',i,0;\pkcsh_{1,0}(\rho_0),\pkcsh_{1,1}(\pkcpd_{a,t,1,0}))$. The old
depth is $1$ at level $0$, $a+t-u$ at level $1$, $-u$ above; the new depth $\pkcpd_{a,t,1,t}$ is
$a+t-u$, $t-u$, $-u$. So steps occur only at level $0$, with $1\le y<a+t-u\le a+t+\eta\le n$, in
$\pkcSig_n(b_i)$. Lemma~\ref{pkc:free} gives $q'\in\pkcPt_{S'}(c',i,0;\pkcimp(\pkcsh_{1,0}(\rho_0),
\pkcpd_{a,t,1,t}),\pkcpd_{a,t,1,t})$. The orders $\pkcsh_{1,0}(\rho_0)$ are $0$ at level $0$ and $-u$ at
level $a'\ge1$. Improved: at level $0$, $\min(a+t-u,0)=0\ge-t-u$ (as $-\eta\le u\le0$); at level
$a'\ge1$, $-u\ge0$, and the improved order is $\pkcpd_{a,t,1,t}(a',u)\ge-u$ if $u<0$ and
$\min(\pkcpd_{a,t,1,t}(a',0),0)=0$ if $u=0$. So the orders of $q'$ are $\ge\rho_{-t}$. Now
Lemma~\ref{pkm:plain} with the roles exchanged ($l=1$, killer exponent $a$, radius $t$, $\eta_0=t$,
the point of the side that is down $q'$ with depth $\pkcpd_{a,t,1,t}$, $x=x(i)$ of orders
$\rho_{-t}$ and $x_0=x(i^{(-1)})$ of orders $\rho_{-t}=\rho_{-\eta_0}$) gives that $q'$ is a stop
(for $t,1$) and the move of the first side by $\varphi x(i^{(-1)})-x(i)$ with $q'$ at rest is
$\kappa$-safe. This is (St).

\emph{The walk.} Corollary~\ref{pkm:walk-all} with $t$ there equal to $a$, radius $t$, $e=\kappa$,
$e'=0$ gives $\mathrm{SCC}_N(c,c')$ for $N\ge\kappa+0+\kappa$.
\end{proof}

\subsection{Generic primes: Proposition~\ref{pki:generic}}\label{pkc:ss-generic}

\begin{theorem}[Proposition~\ref{pki:generic}]\label{pkc:generic}
Let $e=1$, $\ell\ge5$, and $a,\eta,W\ge0$ with $\eta\ge1$, such that for every letter $z$:
$b_z+W+\eta\le(\ell-3)a$; $6(a+\eta+W)<b_z$; $\pkcNM(\ell,6(a+\eta+W),b_z)$; $U_z\subseteq
\pkiWin(\eta,-\varepsilon_z,b_z,W)$ and $U'_z\subseteq\pkiWin(\eta,\varepsilon_z,b_z,W)$. Let $L$ be
crystalline with jumps in $U'$.
\begin{enumerate}
\item[(i)] If no element of any $U(\gamma)\cup U'(\gamma)$ lies in $[-a,0)$, or if every
$U'(\gamma)\subseteq[-\eta,0]$, then $\plldef_{\pkiC(U)}(L)=0$.
\item[(ii)] If $L$ has at most $\mu$ jumps (with multiplicity) in $[-\eta,\eta]$ at each $\gamma$, then
$\plldef_{\pkiC(U)}(L)\le11f\bigl(\mu(a+1)+\pllrk_\auO L\cdot(W+1)\bigr)$.
\end{enumerate}
\end{theorem}

\emph{Dependencies:} (i) is proved from (I1), (I2); (ii) from (I1)--(I4).

\begin{proof}[Setting up]
As $e=1$, $E_v/\Q_\ell$ is unramified, so $j(\Gamma)$ contains the inertia group and
Lemma~\ref{pkp:chart-exists} gives an unramified Kummer chart ($E=z+\ell$, $d_v=0$) with $\Gamma'=
\Gamma$. Fix $\zeta\in\auO$ of exact order $\ell^f-1$ as in \S\ref{pki:ss-arch}; the letter of $\gamma$
is the residue of $\gamma(\zeta)$, and the letter jump sets $U_z$, $U'_z$ of \S\ref{pki:ss-arch} are the
jump sets of the universal sides $S(Y)$, $Y\in\pkiC(U)$, and $S(L)$ of Lemma~\ref{pkc:cup} (every
member of $\pkiC(U)$ is crystalline on $\Gamma=\Gamma'$ with jumps in $U$). Fix $Y\in\pkiC(U)$, a
pairing $P$ modulo some $\ell^m$ of $S=S(Y)$ and $S'=S(L)$, and put $n=6(a+\eta+W)$. Then
\eqref{pkc:eq-win} holds with $b_i=b_z$, $\varepsilon_i=\varepsilon_z$ at each letter $i=z$. Since
$\eta<n<b_z$, Lemmas~\ref{pkc:win}(i) and~\ref{pkc:dualjumps} show that every jump of $S$, of $S'$ and
of $T_P$ at the letter $z$ is $\le b_z+W+\eta\le(\ell-3)a$. Hence, by Proposition~\ref{pkm:kill-unram}
(with $s=0$, $t=a$; $\ell\ge5$), $E^a=\ell^0E^a$ kills the cokernels of $S$, $S'$ and $T_P$ at every
letter.
\end{proof}

\begin{proof}[Proof of (i)]
If no element of any $U(\gamma)\cup U'(\gamma)$ lies in $[-a,0)$, then no jump of $S$ or $S'$ at any
letter lies in $[-a,0)$ (the jump sets are $\{0\}$ and unions of the $U(\gamma)$, resp.\
$U'(\gamma)$), and Lemma~\ref{pkc:gap} with $\kappa=0$ and $d_v=0$ gives $\mathrm{SCC}_0(c,x)$ for all
$c\in\pllZf(Y)$, $x\in\pllZf(L)$. If every $U'(\gamma)\subseteq[-\eta,0]$, then every
$U'_z\subseteq[-\eta,0]$, and Lemma~\ref{pkc:nonpos} with $\kappa=0$ gives the same. As $Y$ and $P$ were
arbitrary, Lemma~\ref{pkc:cup} gives $\pllZf(L)\subseteq L^{\perp\pkiC(U)}$, i.e.\
$\plldef_{\pkiC(U)}(L)=0$.
\end{proof}

\begin{proof}[Proof of (ii)]
\emph{The submodule.} For a letter $i$ let $d_i$ be the number of embeddings $\gamma$ of letter $i$.
Embeddings of different letters are different (the Teichm\"uller lifts of different letters have
different images), so $\sum_id_i\le[\auO:\Z_\ell]$. At each $\gamma$, $L$ has at most $\mu$ jumps in
$[-\eta,\eta]$ and $\pllrk_\auO L$ jumps in all, with multiplicity; hence (I4), applied to $S'=S(L)$ with
$b_i=b_z$, $\varepsilon_i=\varepsilon_z$, $n_i=fd_i\mu$ and $r_i=fd_i\pllrk_\auO L$, gives a
$\Z_\ell$-submodule $V\subseteq\pllZf(L)$ with
\[
\pglen_{\Z_\ell}(\pllZf(L)/V)\le11\sum_ifd_i\bigl(\mu(a+1)+\pllrk L\,(W+1)\bigr)\le
11f\bigl(\mu(a+1)+\pllrk L\,(W+1)\bigr)[\auO:\Z_\ell],
\]
such that for every $x\in V$ and every letter $i$ there is $x'(i)\in\pkcPt_{S'}(x,i,0;\rho_0,\delta_i)$,
with $\delta_i$ as in (I4). We show that $\mathrm{SCC}_0(c,x)$ holds for $c\in\pllZf(Y)$, $x\in V$.
Then Lemma~\ref{pkc:cup} gives $V\subseteq\pllZf(L)\cap L^{\perp\pkiC(U)}$, so
$\pllZf(L)/(\pllZf(L)\cap L^{\perp\pkiC(U)})$ is a quotient of $\pllZf(L)/V$; its $\Z_\ell$-length is
at most the bound above, and dividing by $\pglen_{\Z_\ell}(\auO/\ell)=[\auO:\Z_\ell]$ (the
normalisation $\pglen\auO/\ell=1$) gives (ii).

\emph{Points of $S$.} Fix $c$, $x$ and a letter $i$; write $b=b_i$. By (I1) (with the killer $E^a$ of
$S$), $\pkcPt_S(c,i,0;\rho_{-a},\rho_{-a})\ne\emptyset$. Deepen to $\rho_{2a+\eta}$: the steps are at
level $0$ with $-u-a\le y<2a+\eta-u$, hence in $\pkcSig_n(b)$ by Lemma~\ref{pkc:win}(iii) ($s=a$,
$t'=2a+\eta$, $2a+2\eta+W\le n$). Lemma~\ref{pkc:free} gives $x(i)\in\pkcPt_S(c,i,0;o,\rho_{2a+\eta})$
with $o=\pkcimp(\rho_{-a},\rho_{2a+\eta})$, that is,
\[
o(0,u)=2a+\eta-u\ (u<-a),\qquad o(0,u)=\min(2a+\eta-u,0)\ (u\ge-a),\qquad o(a',u)=-u\ (a'\ge1).
\]
Two consequences. (a) $o(0,u)\ge-a-u$ for every $u$ (if $u\ge-a$ both terms of the minimum are $\ge-a-u$).
(b) $o\ge\rho_{-\eta}$ at the jumps of $S$: at level $0$ with $u<-a$, $2a+\eta-u\ge-\eta-u$; with $u\ge-a$
the jump $u$ is either low, so $u\ge-\eta$, or far, so $|u|\ge b>a$ and hence $u\ge b>0$; in both cases
$\min(2a+\eta-u,0)\ge-\eta-u$. Moreover $\rho_{2a+\eta}\ge\pkcpd_{a,\eta,1,0}$. So, by
Remark~\ref{pkc:mono}, $x(i)\in\pkcPt_S(c,i,0;\rho_{-\eta},\pkcpd_{a,\eta,1,0})$ and $x(i)$ is a standard
point for $(a,\eta,1)$.

\emph{Points of $S'$.} The point $x'(i)$ has orders $\rho_0\ge\rho_{-\eta}$ and depth
$\delta_i\ge\pkcdep_{a,\eta,1}$: if $\varepsilon_i=1$, $\rho_{2a+\eta}(0,u)=2a+\eta-u\ge a+\eta-u$; if
$\varepsilon_i=-1$, at level $0$ we have $b+W+a+\eta\ge a+\eta-u$ because every jump of $S'$ at $i$ is
$\ge-(b+W)$, at level $1$ we have $W+1-u+2a[|u|\le\eta]\ge-u$, and above both are $-u$. So $x'(i)$ is a
standard point for $(a,\eta,1)$.

\emph{Stops, $\varepsilon_i=1$.} The jumps of $S$ at $i$ are $\le\eta$, those of $S'$ are $\ge-\eta$.
Lemma~\ref{pkm:plain} ($l=1$, $t=a$, radius $\eta$, $\eta_0=0$, killer $E^a$ of $T_P$ at $i$, $x=x(i)$,
$x'=x'(i)$, $x'_0=x'(i^{(-1)})$ of orders $\rho_0$) gives (St) of Theorem~\ref{pkm:walk} at $i$.

\emph{Stops, $\varepsilon_i=-1$.} The jumps of $S'$ at $i$ are $\le\eta$, so the standard point $x'(i)$
is a stop for $(\eta,1)$ by Lemma~\ref{pkm:down}. It remains to show that the move of the first side by
$w=\varphi x(i^{(-1)})-x(i)$ with $x'(i)$ at rest is $0$-safe; we use (I3) with $l=1$, $\kappa=0$.

The dual point: $E^a$ kills the cokernel of $T_P$ at $j=i^{(-1)}$, so (I1) gives a point of $T_P$ over $1$
at $j$ with orders and depth $\rho_{-a}$; deepening it to $\rho_a$ involves only level-$0$ steps
$-a-s\le y<a-s$ at the jumps $s$ of $T_P$ at $j$, which lie in $\pkcSig_n(b_j)$ by
Lemma~\ref{pkc:dualjumps}. Lemma~\ref{pkc:free} gives a point $\nu$ of $T_P$ over $1$ at $j$ with orders
$o_\nu=\pkcimp(\rho_{-a},\rho_a)$:
\[
o_\nu(0,s)=a-s\ (s<-a),\qquad o_\nu(0,s)=\min(a-s,0)\ (s\ge-a),\qquad o_\nu(a',s)=-s\ (a'\ge1).
\]
Each $o_\nu(a',\cdot)$ is nonincreasing (at level $0$: both branches are, and at the junction
$o_\nu(0,-a-1)=2a+1\ge0\ge o_\nu(0,-a)$). Note $o_\nu(0,s)\ge-a-s$ and $o_\nu(0,s)\ge\min(a-s,-(W+1)-s)$
when $s\ge-W-1$ (if $s\ge-a$, $\min(a-s,0)\ge-(W+1)-s$).

The move: $\varphi x(i^{(-1)})\in L_i(\pkcsh_{1,0}(o))$ lies over $c$ at $i$ (Lemma~\ref{pkp:tau}(v)), and
so does $x(i)\in L_i(o)$. Hence $w$ lies over $0$ and $w\in L_i(o_w)$ with
$o_w=\min(\pkcsh_{1,0}(o),o)$: $o_w(0,u)=\min(0,o(0,u))$, $o_w(1,u)=\min(o(0,u),-u)$, $o_w(a',u)=-u$ for
$a'\ge2$. The point at rest is $y=x'(i)$, with depth $d_y=\pkcnw_i$.

We check the inequality of (I3) at $a'\ge0$, $u\in U_S(i^{(-a')})$, $u'\in U_{S'}(i^{(-a')})$, $s=-u-u'-1$.
At level $a'$ the windows are those of the letter $i^{(-a')}$, with radius $\eta$, width $W$, height
$\beta=b_{i^{(-a')}}>n$ and some sign $\varepsilon'$ (opposite for $S$ and $S'$).

\emph{Level $0$} (letter $i$, $\varepsilon_i=-1$, so $U_S(i)\subseteq\pkiWin(\eta,1,b,W)$). The sum is
$0+\min(0,o(0,u))+b+W+a+\eta$. If $u$ is low, then $o(0,u)\ge0$ ($2a+\eta-u\ge2a\ge0$), and the sum is
$b+W+a+\eta>0$. If $u\in[b,b+W]$, then $u>2a+\eta$ and $o(0,u)=2a+\eta-u$, so the sum is
$\ge 2a+\eta-(b+W)+b+W+a+\eta=3a+2\eta>0$.

\emph{Level $1$.} The sum is $o_\nu(0,s)+\min(o(0,u),-u)+W+1-u'+2a[|u'|\le\eta]$, and
$\min(o(0,u),-u)\ge-a-u$ by (a).
(A) $u'$ low: with $o_\nu(0,s)\ge-a-s=-a+u+u'+1$ the sum is $\ge(-a+u+u'+1)+(-a-u)+(W+1-u'+2a)=W+2$.
(B) $u'$ far, $u$ low: then $|u'|\ge\beta$ and $|u|\le\eta$, so either $s\le\eta-\beta-1<-a$ or
$s\ge\beta-\eta-1>a$ (as $\beta>n\ge a+\eta+1$); in both cases $o_\nu(0,s)=a-s$. The sum is
$\ge(a+u+u'+1)+(-a-u)+(W+1-u')=W+2$.
(C) $u'$ and $u$ far: they lie in far windows of opposite signs, so $|u+u'|\le W$ and $s\ge-W-1$, whence
$o_\nu(0,s)\ge-(W+1)-s=u+u'-W$. Also $|u|\ge\beta>2a+\eta$, so $o(0,u)=2a+\eta-u>-u$ (if $u<-a$ this
is the formula for $o$; if $u\ge-a$, then $u\ge\beta>2a+\eta$ and $o(0,u)=\min(2a+\eta-u,0)=2a+\eta-u$), and $\min(o(0,u),-u)=-u$. The sum is $\ge u+u'-W-u+W+1-u'=1$.
In all cases the sum is positive.

\emph{Level $a'\ge2$.} $\pkcsh_{1,0}(o_\nu)(a',s)=o_\nu(a'-1,s)=-s$, $o_w(a',u)=-u$, $d_y(a',u')=-u'$;
the sum is $u+u'+1-u-u'=1>0$.

By (I3), the move is $0$-safe, which is (St) at $i$.

\emph{The walk.} At every letter, $E^a$ kills the cokernel of $T_P$, $x(i)$ and $x'(i)$ are standard
points for $(a,\eta,1)$ over $\ell^0c$, $\ell^0x$, and (St) holds with $l(i)=1$.
Corollary~\ref{pkm:walk-all} with $\kappa=e=e'=0$ gives $\mathrm{SCC}_0(c,x)$.
\end{proof}

\subsection{Tame types: Proposition~\ref{pki:tame}}\label{pkc:ss-tame}

Let $E_v=\Q_\ell$, so that $j\colon\Gamma\to G_{\Q_\ell}$ is bijective. Let $r\ge1$ and
$e_0=\operatorname{lcm}_{1\le i\le r}(\ell^i-1)$; note $\ell\nmid e_0$. For a potentially crystalline
lattice $Y$ and $s\in\Gamma$ in the inertia group, $s$ acts $\auO$-linearly on $D_{\mathrm{pcris}}(Y)$
and ``$P$ kills $s$ on $D_{\mathrm{pcris}}(Y)$'' is as in Definition~\ref{pkw:def-pcris}. The case
uses the following inputs, proved in \S\ref{pkw:sec}.

\begin{enumerate}
\item[(T1)] \emph{Tame chart.} If $r<\ell$, there is a tame chart (Definition~\ref{pkw:def-tame}) with
ramification index $e_0$, subgroup $\Gamma'$ and tame generator $s_0$ (Proposition~\ref{pkw:T1}).
\item[(T2)] \emph{Descent.} Every potentially crystalline lattice $Y$ such that $X^{e_0}-1$ kills every
inertial $s$ on $D_{\mathrm{pcris}}(Y)$ is crystalline on $\Gamma'$, with the same Hodge--Tate jumps
(Theorem~\ref{pkw:T2}).
\item[(T3)] \emph{Tame killer.} Let $S$ be a side over the chart of (T1) whose jumps at all letters are
$\le n_0$, and let a monic $q\in\Z_\ell[X]$ kill $s_0$ on $D_{\mathrm{pcris}}$ of the lattice of $S$. If
$e_0\le\ell^{m_0}$, $m_0(\deg q+1)\le\kappa$ and $n_0+\kappa+2<\ell$, then $\ell^\kappa=\ell^\kappa E^0$
kills the cokernel of $S$ at every letter (Theorem~\ref{pkw:T3}).
\item[(T4)] \emph{Tate duals of tensor products.} If monic $q,q'\in\Z_\ell[X]$ kill $s_0^{-1}$ on
$D_{\mathrm{pcris}}(Y)$ and on $D_{\mathrm{pcris}}(L)$ respectively, then some monic $q''\in\Z_\ell[X]$
of degree $\le\deg q\cdot\deg q'$ kills $s_0$ on
$D_{\mathrm{pcris}}(\operatorname{Hom}(Y\otimes L,\auO(1)))$ (Theorem~\ref{pkw:T4}).
\end{enumerate}

In (T3) we use the letters with the root $\zeta=1$: all embeddings have letter $1$, the letter-$1$ part
is everything, and the other letter parts are $0$ (Definition~\ref{pkp:def-letters} uses $\zeta$ only
through $\zeta^{\ell^f-1}=1$; Lemma~\ref{pkw:letters}). This avoids asking that $\auO$ contain roots of
unity of order $\ell^f-1$ for the residue degree $f$ of the tame chart.

\begin{lemma}\label{pkc:lcm}
$e_0\le\ell^{r^2}$.
\end{lemma}

\begin{proof}
$e_0$ divides $\prod_{i=1}^r(\ell^i-1)<\prod_{i=1}^r\ell^i=\ell^{r(r+1)/2}\le\ell^{r^2}$.
\end{proof}

\begin{theorem}[Proposition~\ref{pki:tame}]\label{pkc:tame}
Let $E_v=\Q_\ell$, $r,N\ge1$, $e_0$ as above. For each inertial $s\in\Gamma$ let $K(s)$ be a set of
polynomials over $\auO$ containing $X^{e_0}-1$ and the image of a monic polynomial over $\Z_\ell$ of
degree $\le N$; let $\pkiC(U,K)$ be the class of potentially crystalline lattices with jumps in $U$ on
whose $D_{\mathrm{pcris}}$ each $P\in K(s)$ kills $s$, for every inertial $s$. Let $L\in\pkiC(U',K^{-1})$,
$K^{-1}(s)=K(s^{-1})$. Let $\eta,W,b\ge0$ and a sign $\varepsilon$ with
$U(\gamma)\subseteq\pkiWin(\eta,-\varepsilon,b,W)$ and $U'(\gamma)\subseteq\pkiWin(\eta,\varepsilon,b,W)$
for all $\gamma$, and $\ell>b+6(W+\eta)+r^2(N+1)^2+2$. Then $\pkiC(U,K)$ has cup exponent
$3r^2(N+1)^2+150(\eta+W+1)$ on $L$.
\end{theorem}

\emph{Dependencies:} proved from (I1), (I2) and (T1)--(T4); these are proved in \S\S\ref{pkt:sec},
\ref{pke:sec}, \ref{pkw:sec}.

\begin{proof}
Put $\kappa=r^2(N+1)^2$, $n=6(\eta+W)$, $n_0=b+n$ and $H=3\kappa+150(\eta+W+1)$. Since
$\ell>\kappa+2\ge r+2$, (T1) gives a chart with ramification index $e_0$, subgroup $\Gamma'$ and tame
generator $s_0$; as $\ell\nmid e_0$, $d_v=0$ if $e_0=1$ and $d_v=v_\ell(e_0)+1=1$ otherwise, so $d_v\le1$.
Use $\zeta=1$: every letter jump set is $\{0\}$ or $\{0\}\cup\bigcup_\gamma U(\gamma)$ (resp.\ with
$U'$), hence lies in $\pkiWin(\eta,-\varepsilon,b,W)$ (resp.\ $\pkiWin(\eta,\varepsilon,b,W)$), as $0$ is
in the low window.

\emph{Descent.} Each $Y\in\pkiC(U,K)$ is killed at every inertial $s$ by $X^{e_0}-1\in K(s)$, and $L$ by
$X^{e_0}-1\in K(s^{-1})=K^{-1}(s)$; by (T2) all of them are crystalline on $\Gamma'$, with jumps in $U$,
resp.\ $U'$. So Lemma~\ref{pkc:cup} applies to the class $\pkiC(U,K)$ and $L$ over this chart, and it
suffices to prove $\mathrm{SCC}_H(c,x)$ for every $Y\in\pkiC(U,K)$, every pairing $P$ modulo some
$\ell^m$ of $S=S(Y)$ and $S'=S(L)$, $c\in\pllZf(Y)$ and $x\in\pllZf(L)$.

\emph{Killers.} By hypothesis there are monic $q_0\in\Z_\ell[X]$ with image in $K(s_0)$ and $q'_0$ with
image in $K(s_0^{-1})=K^{-1}(s_0)$, of degrees $\le N$; so $q_0$ kills $s_0$ on $D_{\mathrm{pcris}}(Y)$
and $q'_0$ kills $s_0$ on $D_{\mathrm{pcris}}(L)$; and also $q'_0$ kills $s_0^{-1}$ on
$D_{\mathrm{pcris}}(Y)$ (as $Y\in\pkiC(U,K)$) and $q_0$ kills $s_0^{-1}$ on $D_{\mathrm{pcris}}(L)$ (as
$K^{-1}(s_0^{-1})=K(s_0)$). By (T4), applied with $q=q'_0$ and $q'=q_0$, a monic
$q''$ of degree $\le N^2$ kills $s_0$ on $D_{\mathrm{pcris}}$ of $\operatorname{Hom}(Y\otimes L,\auO(1))$,
the lattice of $T_P$. All jumps of $S$, $S'$ are $\le b+W+\eta\le n_0$ (Lemma~\ref{pkc:win}(i)), and all
jumps of $T_P$ (namely $0$ and the $-u-u'-1$) are $\le b+W+2\eta\le n_0$ (Lemma~\ref{pkc:win}(ii), first
assertion, which does not need $\eta\le b$). Take $m_0=r^2$: $e_0\le\ell^{r^2}$ by Lemma~\ref{pkc:lcm};
$r^2(\deg q_0+1)\le r^2(N+1)\le\kappa$, the same for $q'_0$, and $r^2(\deg q''+1)\le r^2(N^2+1)\le\kappa$;
and $n_0+\kappa+2<\ell$ is the hypothesis. By (T3), $\ell^\kappa=\ell^\kappa E^0$ kills the cokernels of
$S$, $S'$ and $T_P$ at every letter.

\emph{The walk.} Apply Lemma~\ref{pkc:gap} with $a=0$ (so the gap condition, on $[-0,0)=\emptyset$, is
empty), $b_i=b$, $\varepsilon_i=\varepsilon$ at every letter, and $n=6(0+\eta+W)$. The condition
$\pkcNM(\ell,n,b)$ holds because every element of $\pkcSig_n(b)$ has absolute value $\le b+n<\ell$, while
a nonzero multiple of $\ell$ has absolute value $\ge\ell$. Lemma~\ref{pkc:gap} gives
$\mathrm{SCC}_{N'}(c,x)$ for every $N'\ge3\kappa+6(2n+1)d_v$, in particular for $N'=H$, since
\[
3\kappa+6(2n+1)d_v\le3\kappa+72(\eta+W)+6\le3\kappa+150(\eta+W+1)=H .
\]
By Lemma~\ref{pkc:cup}, $\ell^Hx\in L^{\perp\pkiC(U,K)}$ for every $x\in\pllZf(L)$.
\end{proof}

\begin{remark}[On (T4)]\label{pkc:rem-T4}
(T4) is Theorem~\ref{pkw:T4}: a Cayley--Hamilton polynomial on
$\Z_\ell[x]/(q)\otimes\Z_\ell[x']/(q')$ kills $s_0^{-1}$ on the tensor product, and passing to the Tate
dual exchanges $s_0^{-1}$ and $s_0$ (Lemma~\ref{pkw:dual}).
\end{remark}

\subsection{Sparse $b$-letters: Proposition~\ref{pki:sparse}}\label{pkc:ss-sparse}

Proposition~\ref{pki:sparse} is proved in \S\ref{pkd:ss-sparse} (Theorem~\ref{pkd:sparse}), from
(I1), (I2) and, when the sign is not constant, the input (I5) of \S\ref{pkd:ss-tools}.

\subsection{Any pattern of $b$-letters: Proposition~\ref{pki:any}}\label{pkc:ss-any}

Proposition~\ref{pki:any} is proved in \S\ref{pkd:ss-any} (Theorem~\ref{pkd:any}), from (I1), (I2),
(I5).



%% ===== proofs/C-kisin-cases2.tex =====
% proofs/C-kisin-cases2.tex
% Propositions pki:sparse and pki:any (the cases with cuts), and the status of pki:data.
% Labels start with pkd:. Conventions and macros as in proofs/C-kisin-cases.tex (pkc:).
\providecommand{\auO}{\mathcal O}
\providecommand{\panot}[1]{\par\smallskip{\raggedright\noindent\textbf{[Not proved here:} #1\textbf{]}\par}\smallskip}
\providecommand{\pllZf}{Z_f}
\providecommand{\plldef}{\delta^{\vee}}
\providecommand{\pglen}{\operatorname{len}}
\providecommand{\pkiWin}{\operatorname{Win}}
\providecommand{\pkiC}{\mathcal C}
\providecommand{\pkilog}{\operatorname{lg}}
\providecommand{\pkcPt}{\mathrm{Pt}}
\providecommand{\pkcimp}{\mathrm{imp}}
\providecommand{\pkcpd}{\mathrm{pd}}
\providecommand{\pkcdep}{\mathrm{dep}}
\providecommand{\pkcsh}{\mathrm{sh}}
% macros of this file (prefix \pkd)
\providecommand{\pkdco}{\mathrm{co}}
\providecommand{\pkdlg}{\mathrm{gd}}
\providecommand{\pkdstop}{\mathrm{stop}}
\providecommand{\pkdstd}{\mathrm{std}}
\providecommand{\pkdcost}{\mathrm{cost}}

\section{The cases with cuts: sparse and arbitrary patterns of $b$-letters}\label{pkd:sec}

This section proves Proposition~\ref{pki:sparse} (\S\ref{pkd:ss-sparse}) and
Proposition~\ref{pki:any} (\S\ref{pkd:ss-any}), the two cases not treated in \S\ref{pkc:sec}, and
the remaining assertions of Proposition~\ref{pki:data} (\S\ref{pkd:ss-data}). Both cases are instances of
one scheme (Theorem~\ref{pkd:scheme} in \S\ref{pkd:ss-scheme}): at a chosen set of letters, the
\emph{cuts}, both sides receive standard points (from a cokernel killer and deepening), and at every
cut the walk of Theorem~\ref{pkm:walk} needs a stop, which is either the standard point of the
side that is ``down'' (a \emph{plain} cut, Lemma~\ref{pkm:plain}) or a point produced by a
\emph{long gap} (an expansion in powers of $E$; input (I5) below) together with the rule
(Proposition~\ref{pkp:rule}). We use the conventions and notation of \S\ref{pkc:sec}
(points $\pkcPt_S(c,i,e;o,d)$, the orders $\rho_t$, $\pkcdep$, $\pkcpd$, $\pkcimp$, $\pkcsh$,
$\mathrm{std}_l$, $\mathrm{stop}_{\eta,l}$, $\mathrm{rule}_{l,t}$, the cost
$\mathrm{cost}(a,n)=v_\ell(n)+d_v+3a$ with $v_\ell(0)=0$), and cite the inputs (I1), (I2) of
\S\ref{pkc:ss-inputs} as stated there.

Propositions~\ref{pki:sparse} and~\ref{pki:any} are proved below as Theorems~\ref{pkd:sparse}
and~\ref{pkd:any} from \S\S\ref{pkp:sec}--\ref{pkn:sec} and the inputs (I1), (I2), (I5)
((I5) is not needed when the sign is constant); Proposition~\ref{pki:data} is Theorem~\ref{pkd:data},
whose tame membership $L_c^*\in\pkiC(U',K^{-1})$ uses the input (I6). The inputs are proved elsewhere: (I1)
(Proposition~\ref{pkt:I1}), (I2) (Theorem~\ref{pke:I2}), (I5) (Theorem~\ref{pkg:I5}), (I6)
(Lemma~\ref{pkw:dual}); Proposition~\ref{pki:chart} is Theorem~\ref{pkj:chart}.

\subsection{Tools}\label{pkd:ss-tools}

\begin{remark}[Scaling]\label{pkd:scale}
If $x\in\pkcPt_S(c,i,e;o,d)$ and $k\ge0$, then $\ell^kx\in\pkcPt_S(c,i,e+k;o,d)$: $\mathcal N(X_S)$
and the modules $L_i(o)$, $L_i(d)$ are modules over $\Z_\ell$, $\mathrm{pr}$ is linear, and
$g(\ell^kx)-\ell^kx=\ell^k(gx-x)$. So a point over $\ell^ec$ gives points over $\ell^{e'}c$ for all
$e'\ge e$, with the same orders and depth.
\end{remark}

\begin{lemma}[Valuations in an interval]\label{pkd:val}
Let $c\in\Z$, $r\ge0$ and $B\ge1$ with $|c|+r\le B$. Then
$\sum_{|y-c|\le r}v_\ell(y)\le(2r+1)/(\ell-1)+\pkilog_\ell B$.
\end{lemma}

\begin{proof}
For $k\ge1$ let $n_k$ be the number of $y\neq0$ with $|y-c|\le r$ and $\ell^k\mid y$; then
$\sum_{|y-c|\le r}v_\ell(y)=\sum_{k\ge1}n_k$ (the term $y=0$ contributes $0$). A nonzero $y$ in the
interval has $|y|\le B$, so $n_k=0$ unless $\ell^k\le B$, i.e.\ unless $k\le\pkilog_\ell B$. Among
$2r+1$ consecutive integers at most $(2r+1)/\ell^k+1$ are divisible by $\ell^k$. Hence the sum is at
most $\sum_{k=1}^{\pkilog_\ell B}\bigl((2r+1)\ell^{-k}+1\bigr)\le(2r+1)/(\ell-1)+\pkilog_\ell B$.
\end{proof}

\subsubsection*{The setting of this section}

Throughout, $P$ pairs sides $S$, $S'$ over a Kummer chart, $\eta\ge1$, $V\ge0$, and at every letter
$i$ an integer $\beta_i\ge0$ and a sign $\varepsilon_i$ are given with
\begin{equation}\label{pkd:eq-win}
U_S(i)\subseteq\pkiWin(\eta,-\varepsilon_i,\beta_i,V),\qquad
U_{S'}(i)\subseteq\pkiWin(\eta,\varepsilon_i,\beta_i,V).
\end{equation}
We say that a side \emph{has the negative window at $i$} if it is $S$ and $\varepsilon_i=1$, or $S'$
and $\varepsilon_i=-1$; then all its jumps at $i$ are $\le\eta$ and all jumps of the other side at
$i$ are $\ge-\eta$. A letter $i$ is \emph{small} if $\beta_i+V\le\eta$; then all jumps of $S$ and $S'$
at $i$ lie in $[-\eta,\eta]$.

\begin{lemma}[Jump bounds]\label{pkd:jumps}
At every letter $i$, every jump of $S$ and of $S'$ is $\le\beta_i+V+\eta$ and $\ge-(\beta_i+V+\eta)$,
and every jump of $T_P$ is $\le\beta_i+V+2\eta$. A jump $u$ of $S$ or $S'$ at $i$ with
$u>-\beta_i$ satisfies $u\ge-\eta$.
\end{lemma}

\begin{proof}
The first assertion is clear from the windows. The jumps of $T_P$ at $i$ are $0$ and the
$s=-u-u'-1$ with $u\in U_S(i)$, $u'\in U_{S'}(i)$ (Definition~\ref{pkp:def-pairing}). If both are low,
$s\le2\eta-1$. If exactly one is far, say $u=\delta x$ with $x\in[\beta_i,\beta_i+V]$ and $|u'|\le\eta$,
then $s\le\beta_i+V+\eta-1$ (if $\delta=-1$) or $s\le\eta-1$ (if $\delta=1$). If both are far they lie in
windows of opposite signs, so $|u+u'|\le V$ and $s\le V-1$. For the last assertion: a jump in a far
window with $u>-\beta_i$ is in the positive far window, so $u\ge\beta_i\ge0$; otherwise $|u|\le\eta$.
\end{proof}

\subsubsection*{Cut orders and the long-gap depth}

For $s,t,A\ge0$ and $l\ge1$ with $t<l$ define
\[
\pkdco_{s,l}=\pkcimp\bigl(\rho_{-s},\pkcdep_{s,\eta,l+1}\bigr),\qquad
\pkdlg_{l,t,A}(a,u)=\begin{cases}A&a=t,\\ 1&a<l,\ a\ne t,\\ -u&a\ge l.\end{cases}
\]
Explicitly, $\pkcdep_{s,\eta,l+1}(a,u)=-u+s[a=0]+\eta[a\le l]$, and $\pkdco_{s,l}(a,u)$ is
$\pkcdep_{s,\eta,l+1}(a,u)$ where $\rho_{-s}(a,u)>0$ and $\min(\pkcdep_{s,\eta,l+1}(a,u),0)$ elsewhere.

\begin{lemma}[Cut orders]\label{pkd:co}
Let $s\ge0$, $l\ge1$. For all $a\ge0$ and $u\in\Z$:
\begin{enumerate}
\item[(i)] if $a\le l$, then $\pkdco_{s,l}(a,u)\ge\min(\eta-u,0)$;
\item[(ii)] if $a\ge1$, or $a=0$ and $s=0$, then $\pkdco_{s,l}(a,u)\ge-u$;
\item[(iii)] if $u\ge-s$ implies $u\ge-\eta$, then $\pkdco_{s,l}(0,u)\ge-\eta-u$.
\end{enumerate}
In particular, if no jump of a side $R$ at $i$ lies in $[-s,-\eta)$, then $\pkdco_{s,l}\ge\rho_{-\eta}$
at the jumps of $R$ (at level $a$, the jumps at $i^{(-a)}$).
\end{lemma}

\begin{proof}
Write $D=\pkcdep_{s,\eta,l+1}$. (i) At level $0$, $D(0,u)=s+\eta-u\ge\eta-u$; at level $1\le a\le l$,
$D(a,u)=\eta-u$. In both cases $\pkdco_{s,l}(a,u)\in\{D(a,u),\min(D(a,u),0)\}$ is
$\ge\min(\eta-u,0)$. (ii) If $a\ge1$, then $\rho_{-s}(a,u)=-u$ and $D(a,u)\ge-u$; if $-u>0$ the value is
$D(a,u)\ge-u$, otherwise it is $\min(D(a,u),0)\ge-u$ since both terms are $\ge-u$. If $a=0$ and
$s=0$ the same argument applies with $\rho_0(0,u)=-u$. (iii) If $u<-s$, then $\rho_{-s}(0,u)=-s-u>0$
and the value is $s+\eta-u\ge-\eta-u$. If $u\ge-s$, then $u\ge-\eta$, and $\min(s+\eta-u,0)\ge-\eta-u$.
The last assertion follows from (ii) and (iii), as $\rho_{-\eta}(a,u)=-u$ for $a\ge1$.
\end{proof}

\subsubsection*{The input (I5)}

The long-gap points are produced by the following statement about the modules $\mathcal N$, which is
not among the statements of \S\S\ref{pkp:sec}--\ref{pkn:sec}; it is proved as Theorem~\ref{pkg:I5}.
For $l,A\ge0$ and $N\in\Z$ put
\[
o^{\mathrm{gap}}_{l,A,N}(a,u)=\mathrm{std}_l(a,u)+A[a=0][u<-N],\qquad
d^{\mathrm{gap}}_{l,A}(a,u)=\mathrm{std}_l(a,u)+A[a=0].
\]

\begin{enumerate}
\item[(I5)] \emph{Long gap.} Let the chart be unramified, $l\ge1$ with $\ell^l\ge3$, and
$A,\delta,\kappa,s,g\ge0$ with $A\le\ell^l\delta$. Let $\ell^\kappa$ kill the cokernel of a side $S$
at $i^{(-l)}$, and $E^s$ kill it at $i$. Let every jump of $S$ at $i^{(-a)}$, $0\le a<l$, be $\le g$,
put $N=s+g(1+\ell+\dots+\ell^{l-1})$, and assume $N<A$. Then for every $c\in C_S$,
$\pkcPt_S(c,i,\kappa+4\delta;o^{\mathrm{gap}}_{l,A,N},d^{\mathrm{gap}}_{l,A})\ne\emptyset$.
\end{enumerate}

\subsection{The scheme of cuts}\label{pkd:ss-scheme}

\begin{definition}[Cuts, plain letters, long gaps]\label{pkd:def-cut}
Let $\mathrm{cut}$ be a set of letters such that every letter $i$ has some $a>0$ with
$i^{(-a)}\in\mathrm{cut}$, and let $l(i)$ be the least such $a$. Let shifts $s_i\ge0$ ($i\in\mathrm{cut}$)
be given with $s_i=0$ or $s_i<\beta_i$. A letter $i$ (with $l=l(i)$) is
\begin{enumerate}
\item[(a)] \emph{plain} if for some sign $\varrho$ every letter $i^{(-a)}$, $0\le a<l$, has
$\varepsilon_{i^{(-a)}}=\varrho$ or is small;
\item[(b)] a \emph{long gap with pivot $t$}, for data $(J',s_0,\omega,A,\delta)$ of nonnegative
integers, if: the chart is unramified and $\ell\ge5$; $t+J'+1=l$; $s_{i^{(-l)}}=0$; for $a<l$,
$a\ne t$, either $\varepsilon_{i^{(-a)}}\ne\varepsilon_{i^{(-t)}}$ or $i^{(-a)}$ is small; $\eta\le\omega$
and $\beta_{i^{(-a)}}+V\le\omega$ for $t<a<l$; $N:=s_0+\omega(1+\ell+\dots+\ell^{J'})<\beta_{i^{(-t)}}$;
$E^{s_0}$ kills the cokernels of $S$ and $S'$ at $i^{(-t)}$; and
$\beta_{i^{(-t)}}+V\le A\le\ell^{J'+1}\delta$.
\end{enumerate}
\end{definition}

For a cut $i$ and a side $R\in\{S,S'\}$ let $\mathcal I_i(R)$ be the set of pairs $(a,y)$ with $0\le
a\le l(i)$ and $\rho_{-s_i}(a,u)\le y<\pkcdep_{s_i,\eta,l(i)+1}(a,u)$ for some $u\in U_R(i^{(-a)})$.

\begin{lemma}[Steps at a cut]\label{pkd:steps}
Let $i$ be a cut, $l=l(i)$, $R\in\{S,S'\}$, and let $B\ge1$ with $B\ge\beta_j+s_i+V+2\eta$ for all
letters $j$. Put $r_0=s_i+V+2\eta$ and $r_a=V+2\eta$ for $a\ge1$. Then
\[
\sum_{(a,y)\in\mathcal I_i(R)}\pkdcost(a,y)\le\sum_{a=0}^{l}\Bigl(\frac{2(2r_a+1)}{\ell-1}
+2\pkilog_\ell B+2(2r_a+1)(d_v+3a)\Bigr).
\]
\end{lemma}

\begin{proof}
Fix $a\le l$, let $j=i^{(-a)}$ and let $\sigma$ be the sign of the far window of $R$ at $j$. At level
$0$ the condition is $-u-s_i\le y<-u+s_i+\eta$; at level $1\le a\le l$ it is $-u\le y<-u+\eta$. If
$|u|\le\eta$ this gives $|y|\le r_a$. If $u=\sigma x$ with $x\in[\beta_j,\beta_j+V]$, it gives
$|y+u|\le s_i[a=0]+\eta$ and so $|y-c_a|\le r_a$ with $c_a=-\sigma\beta_j$. Hence the pairs at level
$a$ have $y$ in the union of the intervals $|y|\le r_a$ and $|y-c_a|\le r_a$, each of $2r_a+1$
integers and contained in $[-B,B]$. Lemma~\ref{pkd:val} bounds the sum of $v_\ell(y)$ over each
interval by $(2r_a+1)/(\ell-1)+\pkilog_\ell B$, and the remaining part $d_v+3a$ of the cost is summed
over at most $2(2r_a+1)$ values.
\end{proof}

\begin{theorem}[The scheme]\label{pkd:scheme}
In the setting \eqref{pkd:eq-win}, let $\mathrm{cut}$, $l$ and the shifts be as in
Definition~\ref{pkd:def-cut}, and let $\kappa,e^*\ge0$. Assume:
\begin{enumerate}
\item[(K)] for every cut $i$, $\ell^\kappa E^{s_i}$ kills the cokernels of $S$, $S'$ and $T_P$ at $i$;
\item[(C)] for every cut $i$ and $R\in\{S,S'\}$, $\kappa+\sum_{(a,y)\in\mathcal I_i(R)}\pkdcost(a,y)\le e^*$;
\item[(G)] every cut $i$ is plain, or is a long gap with some pivot $t$ and data
$(J',s_0,\omega,A,\delta)$ such that $\kappa+4\delta+3\sum_{t<a<l(i)}a\le e^*$.
\end{enumerate}
Then $\mathrm{SCC}_N(c,c')$ for all $c\in C_S$, $c'\in C_{S'}$ and $N\ge2e^*+\kappa$.
\end{theorem}

\begin{proof}
Fix $c$, $c'$. We verify the hypotheses of Theorem~\ref{pkm:walk} with $t(i)=s_i$ and
$e=e'=e^*$.

\emph{Standard points.} Let $i$ be a cut, $l=l(i)$, $R\in\{S,S'\}$ and $d\in\{c,c'\}$ the
corresponding element. By (K) and (I1), $\pkcPt_R(d,i,\kappa;\rho_{-s_i},\rho_{-s_i})\ne\emptyset$. By
(I2) with $I=\mathcal I_i(R)$ (which contains all steps from $\rho_{-s_i}$ to
$\pkcdep_{s_i,\eta,l+1}$: at levels $a>l$ the two orders agree) and (C), and by
Remark~\ref{pkd:scale}, there is a point
\[
x_R(i)\in\pkcPt_R\bigl(d,i,e^*;\pkdco_{s_i,l},\pkcdep_{s_i,\eta,l+1}\bigr).
\]
No jump $u$ of $R$ at $i$ lies in $[-s_i,-\eta)$: if $s_i=0$ this is empty, and if $s_i<\beta_i$,
then $u\ge-s_i$ gives $u>-\beta_i$ and so $u\ge-\eta$ by Lemma~\ref{pkd:jumps}. By
Lemma~\ref{pkd:co} and Remark~\ref{pkc:mono}, $x_R(i)$ has orders $\rho_{-\eta}$. As
$\pkcdep_{s_i,\eta,l+1}=\pkcpd_{s_i,\eta,l,\eta}\ge\pkcdep_{s_i,\eta,l}$, $x_R(i)$ is a standard point
(for $s_i,\eta,l$) and has the depth $\pkcpd_{s_i,\eta,l,\eta}$. Write $x(i)=x_S(i)$, $x'(i)=x_{S'}(i)$.
Then (P) of Theorem~\ref{pkm:walk} holds, and (K) there is part of our (K).

\emph{Plain letters.} Let $i$ be a cut that is plain, with sign $\varrho$, $l=l(i)$, and let $R$ be $S$ if $\varrho=1$
and $S'$ if $\varrho=-1$, $R'$ the other side. For $a<l$ the letter $i^{(-a)}$ either has
$\varepsilon=\varrho$, so $R$ has the negative window there, or is small; in both cases the jumps of $R$
at $i^{(-a)}$ are $\le\eta$ and those of $R'$ are $\ge-\eta$. Lemma~\ref{pkm:plain} (with $l$,
$t=s_i$, $\eta$, $\eta_0=\eta$, the killer $\ell^\kappa E^{s_i}$ of $T_P$ at $i$, the point $x_R(i)$ of
orders $\rho_{-\eta}$ and depth $\pkcpd_{s_i,\eta,l,\eta}$, the point $x_{R'}(i)$ of orders $\rho_{-\eta}$
and the point $x_{R'}(i^{(-l)})$ at $i^{(-l)}$ of orders $\rho_{-\eta}=\rho_{-\eta_0}$; for $R=S'$ in
the form with the roles exchanged) shows that $x_R(i)$ is a stop (for $\eta,l$) over $\ell^{e^*}d$ and
that the move of $R'$ by $\varphi^lx_{R'}(i^{(-l)})-x_{R'}(i)$ with $x_R(i)$ at rest is $\kappa$-safe.
This is (St) at $i$.

\emph{Long gaps.} Let $i$ be a cut that is a long gap, with pivot $t$ and data $(J',s_0,\omega,A,\delta)$, $l=l(i)$,
$i_t=i^{(-t)}$. Let $R$ be the side with the negative window at $i_t$, $R'$ the other side, and $d$, $d'$
the corresponding elements of $\{c,c'\}$. For $a<l$, $a\ne t$, either $\varepsilon_{i^{(-a)}}\ne
\varepsilon_{i_t}$, so that $R'$ has the negative window at $i^{(-a)}$, or $i^{(-a)}$ is small; hence
\begin{equation}\label{pkd:eq-updown}
\text{for }a<l,\ a\ne t:\quad U_R(i^{(-a)})\subseteq[-\eta,\infty),\qquad
U_{R'}(i^{(-a)})\subseteq(-\infty,\eta].
\end{equation}

\emph{(1) A point at the pivot.} We apply (I5) to $R$ at the letter $i_t$, with $l$ there equal to
$J'+1$, $s=s_0$, $g=\omega$, and $A,\delta,\kappa$. Note $i_t^{(-(J'+1))}=i^{(-l)}$ is a cut with shift
$0$, so $\ell^\kappa=\ell^\kappa E^0$ kills the cokernel of $R$ there by (K); $E^{s_0}$ kills it at
$i_t$; $\ell^{J'+1}\ge\ell\ge5$; $A\le\ell^{J'+1}\delta$. The jumps of $R$ at $i_t$ are $\le\eta\le\omega$ (the
negative window), and at $i_t^{(-a)}=i^{(-(t+a))}$ with $1\le a\le J'$ they are $\le\max(\eta,
\beta_{i^{(-(t+a))}}+V)\le\omega$, since $t<t+a<l$. Finally $N<\beta_{i_t}\le A$. So there is
$x_{i_t}\in\pkcPt_R(d,i_t,\kappa+4\delta;o^{\mathrm{gap}}_{J'+1,A,N},d^{\mathrm{gap}}_{J'+1,A})$.

\emph{(2) Transport and deepening.} By Lemma~\ref{pkp:tau}(v), $\varphi^tx_{i_t}$ is a point of $R$ over
$\ell^{\kappa+4\delta}d$ at $i$ with orders $o_1=\pkcsh_{t,0}(o^{\mathrm{gap}}_{J'+1,A,N})$ and depth
$d_1=\pkcsh_{t,1}(d^{\mathrm{gap}}_{J'+1,A})$. Explicitly, since $\mathrm{std}_{J'+1}(b,u)=0$ for
$b\le J'$ and $=-u$ for $b\ge J'+1$, and $t+J'+1=l$:
\[
o_1(a,u)=\begin{cases}0&a<t\text{ or }t<a<l,\\ A[u<-N]&a=t,\\ -u&a\ge l,\end{cases}\qquad
d_1(a,u)=\begin{cases}1&a<t,\\ A&a=t,\\ 0&t<a<l,\\ -u&a\ge l.\end{cases}
\]
Thus $d_1(a,u)\le y<\pkdlg_{l,t,A}(a,u)$ only for $t<a<l$ and $y=0$. By (I2) with
$I=\{(a,0):t<a<l\}$, whose cost is $\sum_{t<a<l}(0+0+3a)$ ($d_v=0$), by (G) and by
Remark~\ref{pkd:scale}, there is a point
$q\in\pkcPt_R(d,i,e^*;\pkcimp(o_1,\pkdlg_{l,t,A}),\pkdlg_{l,t,A})$.

\emph{(3) $q$ is a stop.} We check, at the jumps $u$ of $R$ at $i^{(-a)}$, that
$\pkcimp(o_1,\pkdlg_{l,t,A})\ge\mathrm{stop}_{\eta,l}$ and $\pkdlg_{l,t,A}\ge\mathrm{stop}_{\eta,l}$,
where $\mathrm{stop}_{\eta,l}(a,u)=\max(0,-u-\eta)$ for $a<l$ and $=-u$ for $a\ge l$.
For $a\ge l$: $o_1(a,u)=-u$, so the improved order is $-u$ if $-u>0$ and $\min(-u,0)=-u$ otherwise;
and $\pkdlg_{l,t,A}(a,u)=-u$. For $a<l$, $a\ne t$: $o_1=0$, the improved order is $\min(1,0)=0$, the
depth is $1$, and $\max(0,-u-\eta)=0$ by \eqref{pkd:eq-updown}. For $a=t$: the jumps of $R$ at $i_t$ lie
in $[-\eta,\eta]$ or in $-[\beta_{i_t},\beta_{i_t}+V]$, so $-u-\eta\le\beta_{i_t}+V\le A$ and the depth $A$ is
$\ge\max(0,-u-\eta)$. If $u<-N$ the improved order is $A$; if $u\ge-N$, then $u>-\beta_{i_t}$ (as
$N<\beta_{i_t}$), so $u$ is not in the negative far window and $u\ge-\eta$; the improved order is
$\min(A,0)=0=\max(0,-u-\eta)$. By Remark~\ref{pkc:mono}, $q$ is a stop (for $\eta,l$) of $R$ over
$\ell^{e^*}d$ at $i$, and it has depth $\pkdlg_{l,t,A}$.

\emph{(4) The move is safe.} Let $y'=\varphi^lx_{R'}(i^{(-l)})-x_{R'}(i)$. The cut $i^{(-l)}$ has shift
$0$, and $l(i^{(-l)})=:l_0$; so $x_{R'}(i^{(-l)})$ has orders $\pkdco_{0,l_0}$, and by
Lemma~\ref{pkp:tau}(v) $\varphi^lx_{R'}(i^{(-l)})$ is a point over $\ell^{e^*}d'$ at $i$ with orders
$\pkcsh_{l,0}(\pkdco_{0,l_0})$. As $x_{R'}(i)$ lies over the same element and has orders
$\pkdco_{s_i,l}$, $y'$ lies over $0$ and $y'\in L^0_i(o')$ with
$o'=\min(\pkdco_{s_i,l},\pkcsh_{l,0}(\pkdco_{0,l_0}))$ (each $L_i(o)$ is a module, monotone in $o$). We
apply Proposition~\ref{pkp:rule} (in the form for the move of $R'$ with the point $q$ of $R$ at rest,
depth $\pkdlg_{l,t,A}$), with $l$, $t=0$ and the killer $\ell^\kappa E^0$ of $T_P$ at $i^{(-l)}$
given by (K). We must show, for all $a\ge0$, $v\in U_R(i^{(-a)})$, $u\in U_{R'}(i^{(-a)})$:
\[
\mathrm{rule}_{l,0}(a,u+v)\le\pkdlg_{l,t,A}(a,v)+o'(a,u).
\]
\emph{$a\ge l$.} Here $\mathrm{rule}_{l,0}(a,u+v)=-u-v$ and $\pkdlg_{l,t,A}(a,v)=-v$. By
Lemma~\ref{pkd:co}(ii), $\pkdco_{s_i,l}(a,u)\ge-u$ (as $a\ge l\ge1$) and
$\pkcsh_{l,0}(\pkdco_{0,l_0})(a,u)=\pkdco_{0,l_0}(a-l,u)\ge-u$ (shift $0$). So $o'(a,u)\ge-u$.
\emph{$a<l$.} Here $\mathrm{rule}_{l,0}=1$ and $\pkcsh_{l,0}(\cdot)(a,u)=0$, so
$o'(a,u)=\min(\pkdco_{s_i,l}(a,u),0)\ge\min(\eta-u,0)$ by Lemma~\ref{pkd:co}(i). If $a\ne t$, then
$u\le\eta$ by \eqref{pkd:eq-updown}, so $o'(a,u)=0$ and the right side is $1$. If $a=t$, the jumps of
$R'$ at $i_t$ lie in $[-\eta,\eta]$ or in $[\beta_{i_t},\beta_{i_t}+V]$ (the positive window), so $u\le
\beta_{i_t}+V\le A$ and $\eta-u\ge\eta-A\ge1-A$; as $A\ge\beta_{i_t}+V>N\ge0$, $A\ge1$ and
$o'(t,u)\ge\min(1-A,0)=1-A$. The right side is $\ge A+1-A=1$.

By Proposition~\ref{pkp:rule} the move of $R'$ by $y'$ with $q$ at rest is $\kappa$-safe. If $R=S$
this is the first alternative of (St) at $i$, and if $R=S'$ the second.

\emph{The walk.} (K), (P) and (St) of Theorem~\ref{pkm:walk} hold at every cut, with $e=e'=e^*$; so $\mathrm{SCC}_N(c,c')$ for $N\ge2e^*+\kappa$.
\end{proof}

\begin{remark}
The hypothesis
$s_{i^{(-l)}}=0$ in a long gap is used twice: for the killer $\ell^\kappa$ in (I5) and for the
killer $\ell^\kappa E^0$ in the rule.
\end{remark}

\subsection*{Passing to lattices}

\begin{lemma}[Set-up for the cases]\label{pkd:setup}
Let $U$, $U'$ be jump sets and $L$ crystalline with jumps in $U'$. Fix a Kummer chart with
$\Gamma'=\Gamma$ (Lemma~\ref{pkp:chart-exists}), unramified if $e=1$, and $\zeta$ as in
\S\ref{pki:ss-arch}. Suppose that for every $Y\in\pkiC(U)$, every pairing $P$ modulo some $\ell^m$ of
$S=S(Y)$ and $S'=S(L)$ (the universal sides of Lemma~\ref{pkc:cup}), every $c\in\pllZf(Y)$ and every
$x\in\pllZf(L)$ one has $\mathrm{SCC}_H(c,x)$. Then $\pkiC(U)$ has cup exponent $H$ on $L$.
\end{lemma}

\begin{proof}
Every member of $\pkiC(U)$ and $L$ is crystalline on $\Gamma=\Gamma'$; the jump sets of $S(Y)$, $S(L)$
at the letter $z$ are the letter jump sets $U_z$, $U'_z$ of \S\ref{pki:ss-arch} (as in the proof of
Theorem~\ref{pkc:generic}). Lemma~\ref{pkc:cup} gives $\ell^Hx\in L^{\perp\pkiC(U)}$ for all $x$.
\end{proof}

\subsection{Sparse $b$-letters: Proposition~\ref{pki:sparse}}\label{pkd:ss-sparse}

Let $B$ be a set of letters and $J\ge1$. Write $\sigma(i)=i^\ell$. A letter $i$ has \emph{thin
successors} if $\sigma^n(i)\notin B$ for $1\le n\le J$. $B$ is \emph{sparse} if there is no letter $z$
with $z,\sigma^s(z),\sigma^{s+t}(z)\in B$ for some $s,t\in[1,J]$.

\begin{lemma}\label{pkd:thin-add}
Let $B$ be sparse, $j\in B$ a letter without thin successors. Then $j^{(-(J+1))}$ has thin successors.
\end{lemma}

\begin{proof}
Let $k=j^{(-(J+1))}$ and suppose $m=\sigma^n(k)\in B$ with $1\le n\le J$. Put $s=J+1-n\in[1,J]$; then
$\sigma^s(m)=\sigma^{J+1}(k)=j\in B$. As $j$ has no thin successors, $\sigma^t(j)=\sigma^{s+t}(m)\in B$
for some $t\in[1,J]$. So $m,\sigma^s(m),\sigma^{s+t}(m)\in B$, contradicting sparseness.
\end{proof}

\begin{lemma}\label{pkd:thin-cut}
Let $B$ be sparse and $\mathrm{cut}$ the set of letters with thin successors. Then every letter $i$ has
$l(i)\le2J+1$ (in particular $l(i)$ is defined). If $0<r<l(i)$ and $i^{(-r)}\in B$, then
$r+J+1=l(i)$.
\end{lemma}

\begin{proof}
If $i^{(-(J+1))}$ has thin successors, $l(i)\le J+1$. Otherwise some
$\sigma^n(i^{(-(J+1))})=i^{(-r)}$ with $1\le n\le J$ lies in $B$, where $r=J+1-n\in[1,J]$. If
$i^{(-r)}$ has thin successors, $l(i)\le r$; otherwise, by Lemma~\ref{pkd:thin-add},
$i^{(-(r+J+1))}$ has thin successors and $l(i)\le r+J+1\le2J+1$.

Let $l=l(i)$, $0<r<l$, $i^{(-r)}\in B$. Since $i^{(-l)}$ has thin successors and
$\sigma^{l-r}(i^{(-l)})=i^{(-r)}\in B$ with $l-r\ge1$, we get $l-r\ge J+1$. If $r+J+1<l$, then
$i^{(-r)}$ is not a cut (as $0<r<l$), i.e.\ it is in $B$ without thin successors, and
Lemma~\ref{pkd:thin-add} makes $i^{(-(r+J+1))}$ a cut with $0<r+J+1<l$, contradicting the minimality
of $l$. So $r+J+1=l$.
\end{proof}

\begin{theorem}[Proposition~\ref{pki:sparse}]\label{pkd:sparse}
Let $\eta,J\ge1$ and $V,G,b,M,\vartheta\ge0$ with $G+V\le\eta$; let $B$ be a set of letters, $b_z=b$
for $z\in B$ and $b_z=G$ otherwise, and $U_z\subseteq\pkiWin(\eta,-\varepsilon_z,b_z,V)$,
$U'_z\subseteq\pkiWin(\eta,\varepsilon_z,b_z,V)$ for every letter $z$. Assume that $B$ is sparse; that
either (A) $e=1$, $\ell\ge7$ and $10\ell^J\eta\le b$, or (B) $\varepsilon_z$ is constant and
$8e\le\ell^{J+1}$; that $2e(b+5\eta)\le\ell^{J+1}\vartheta$; and that
$20(J+1+d_v)\eta+\pkilog_\ell(2(b+5\eta))\le M$. Then $\pkiC(U)$ has cup exponent
$10^3(\vartheta+(J+1)M)$ on $L$.
\end{theorem}

The proof uses (I1), (I2), and in case (A) also (I5).

\begin{proof}
By Lemma~\ref{pkd:setup} it suffices to prove $\mathrm{SCC}_H(c,x)$ with $H=10^3(\vartheta+(J+1)M)$
for $S=S(Y)$, $S'=S(L)$, a pairing $P$, $c\in C_S$, $x\in C_{S'}$, over a chart with $\Gamma'=\Gamma$,
unramified in case (A). The windows \eqref{pkd:eq-win} hold with $\beta_i=b_i$. We apply
Theorem~\ref{pkd:scheme} with $\mathrm{cut}$ the letters with thin successors (so $l(i)\le2J+1$ by
Lemma~\ref{pkd:thin-cut}), all shifts $0$, and
\[
\kappa=4\vartheta+(J+1)M,\qquad e^*=8\vartheta+6(J+1)M .
\]
Then $2e^*+\kappa=20\vartheta+13(J+1)M\le H$. Put $B_0=2(b+5\eta)\ge1$.

\emph{Jump sizes.} By Lemma~\ref{pkd:jumps}, at a letter $j$ the jumps of $S$, $S'$, $T_P$ are $\le
b_j+V+2\eta$. At a letter not in $B$ this is $\le G+V+2\eta\le3\eta$; at every letter it is
$\le b+3\eta$ (if $b_j=G$, it is $\le3\eta$). Every letter not in $B$ is small ($G+V\le\eta$).

\emph{(K).} Let $i$ be a cut. We apply Proposition~\ref{pkm:kill-J} to each of $S$, $S'$, $T_P$ with
$\omega=b+3\eta$ and $\omega'=3\eta$ (the letters $\sigma^n(i)$, $1\le n\le J$, are not in $B$). In
case (A), $8e=8<49\le\ell^{J+1}$; in case (B) $8e\le\ell^{J+1}$ is assumed. Then
\[
2Jd_v\omega'+\frac{(2J+4)\omega'}{\ell-1}+\frac{8e\omega}{\ell^{J+1}}
\le6(J+1)d_v\eta+12(J+1)\eta+\frac{4\cdot2e(b+5\eta)}{\ell^{J+1}}
\le(J+1)M+4\vartheta=\kappa,
\]
using $\ell-1\ge1$, $2J+4\le4(J+1)$, $b+3\eta\le b+5\eta$ and $6d_v\eta+12\eta\le20(J+1+d_v)\eta\le M$.
So $\ell^\kappa=\ell^\kappa E^0$ kills the three cokernels at $i$.

\emph{(C).} Let $i$ be a cut, $l=l(i)\le2J+1$. As $s_i=0$, Lemma~\ref{pkd:steps} applies with $B_0$
(indeed $b_j+V+2\eta\le b+3\eta\le B_0$), $r_a=V+2\eta\le3\eta$, $2r_a+1\le7\eta$:
\[
\sum_{\mathcal I_i(R)}\pkdcost\le\sum_{a=0}^{2J+1}\bigl(14\eta+2\pkilog_\ell B_0+14\eta(d_v+3a)\bigr)
=(J+1)\bigl(28\eta(1+d_v)+4\pkilog_\ell B_0\bigr)+42\eta(2J+1)(J+1).
\]
This is $(J+1)(84J\eta+70\eta+28d_v\eta+4\pkilog_\ell B_0)\le5(J+1)M$, since
$5M\ge100(J+1+d_v)\eta+5\pkilog_\ell B_0$. So $\kappa+\sum\pkdcost\le\kappa+5(J+1)M\le e^*$.

\emph{(G).} Let $i$ be a letter, $l=l(i)$. If every $i^{(-a)}$, $a<l$, has
$\varepsilon_{i^{(-a)}}=\varepsilon_i$ or is small, $i$ is plain. In case (B) this always holds.
Otherwise we are in case (A), and there is $t<l$ with $\varepsilon_{i^{(-t)}}\ne\varepsilon_i$ and
$i^{(-t)}$ not small; then $i^{(-t)}\in B$ and $t\ne0$, so $t+J+1=l$ by Lemma~\ref{pkd:thin-cut}, and
every $i^{(-a)}$ with $0<a<l$, $a\ne t$, is not in $B$, hence small. We show that $i$ is a long gap
with pivot $t$ and data $(J,s_0,\eta,b+V,\vartheta)$, where $s_0=\lceil(b+2\eta)/(\ell-3)\rceil$.
The chart is unramified and $\ell\ge7$; $t+J+1=l$; $s_{i^{(-l)}}=0$. For $a<l$, $a\ne t$: $a=0$ has
$\varepsilon_i\ne\varepsilon_{i^{(-t)}}$, the others are small. $\eta\le\omega=\eta$, and for $t<a<l$
the letter $i^{(-a)}$ is small, so $\beta+V\le\eta$. Next,
$1+\ell+\dots+\ell^J<\ell^{J+1}/(\ell-1)\le\tfrac76\ell^J$ and $s_0\le(b+2\eta)/4+1$, so
\[
N=s_0+\eta(1+\dots+\ell^J)<\frac b4+\frac\eta2+1+\frac76\ell^J\eta
\le b\Bigl(\frac14+\frac3{140}+\frac7{60}\Bigr)<b=\beta_{i^{(-t)}},
\]
using $b\ge10\ell^J\eta\ge70\eta$ and $\eta/2+1\le3\eta/2$. The jumps of $S$ and $S'$ at every letter are
$\le b+2\eta\le(\ell-3)s_0$ (Lemma~\ref{pkd:jumps}, with $V\le\eta$), so by
Proposition~\ref{pkm:kill-unram} ($s=0$, $t=s_0$) $E^{s_0}$ kills both cokernels at $i^{(-t)}$. Finally
$\beta_{i^{(-t)}}+V=b+V=A\le b+5\eta\le\ell^{J+1}\vartheta/2\le\ell^{J+1}\vartheta$ ($e=1$). The cost
condition: $\sum_{t<a<l}a\le\sum_{a=1}^{2J}a=J(2J+1)$ and $3J(2J+1)\le6J(J+1)\le(J+1)M$ (as $M\ge20(J+1)$),
so $\kappa+4\vartheta+3\sum_{t<a<l}a\le8\vartheta+2(J+1)M\le e^*$.

Theorem~\ref{pkd:scheme} gives $\mathrm{SCC}_H(c,x)$.
\end{proof}

\subsection{Any pattern of $b$-letters: Proposition~\ref{pki:any}}\label{pkd:ss-any}

\begin{theorem}[Proposition~\ref{pki:any}]\label{pkd:any}
Let $e=1$, $\ell\ge7$, $\eta\ge1$, and $V,M,G,b,K\ge0$ with $V+4\eta\le M\le K$, $0<G\le b$,
$b<K\ell^2$ and $2G<K\ell$. Let $B$ be a set of letters, $b_z=b$ for $z\in B$ and $b_z=G$ otherwise,
$U_z\subseteq\pkiWin(\eta,-\varepsilon_z,b_z,V)$ and $U'_z\subseteq\pkiWin(\eta,\varepsilon_z,b_z,V)$.
Then $\pkiC(U)$ has cup exponent $10^4(K+\pkilog_\ell(2(b+M)))$ on $L$.
\end{theorem}

The proof uses (I1), (I2), (I5).

\begin{proof}
By Lemma~\ref{pkd:setup} (unramified chart, $d_v=0$, $\Gamma'=\Gamma$) it suffices to prove
$\mathrm{SCC}_H(c,x)$, $H=10^4(K+\pkilog_\ell B_0)$ with $B_0=2(b+M)$, for $S=S(Y)$, $S'=S(L)$, a pairing
$P$ and $c$, $x$; \eqref{pkd:eq-win} holds with $\beta_i=b_i$. Note $K\ge1$ (as $0<2G<K\ell$), $b\ge1$,
and $V+2\eta\le M-2\eta<M$. By Lemma~\ref{pkd:jumps} all jumps of $S$, $S'$, $T_P$ at every letter are
$\le b+V+2\eta\le b+M$, and those at a letter not in $B$ are $\le G+M$.

\emph{Shifts and cuts.} Put $q=\lfloor b/\ell\rfloor$ and $s^+=\max(q-7K,0)$, and for every letter $i$
let $s_i=s^+$ if $\sigma(i)\in B$ and $s_i=0$ otherwise. Let $\mathrm{cut}=\{i:s_i<\beta_i\}$. Every
$i\in B$ is a cut ($s_i\le q<b$, as $b\ge1$), and every $i$ with $\sigma(i)\notin B$ is a cut
($0<G\le\beta_i$). Hence for every letter $i$: if $i^{(-1)}\in B$ then $l(i)=1$; otherwise
$\sigma(i^{(-2)})=i^{(-1)}\notin B$, so $i^{(-2)}$ is a cut and $l(i)\le2$. The shifts satisfy the
condition of Definition~\ref{pkd:def-cut}. Put
\[
\kappa=14K,\qquad e^*=76K+6\pkilog_\ell B_0,
\]
so that $2e^*+\kappa=166K+12\pkilog_\ell B_0\le H$.

\emph{(K).} Let $i$ be a cut. If $\sigma(i)\notin B$, then $s_i=0$, and Proposition~\ref{pkm:kill-J}
with $J=1$, $\omega=b+M$, $\omega'=G+M$ (the jumps at $\sigma(i)$) applies to $S$, $S'$, $T_P$:
$8e=8\le\ell^2$, $d_v=0$, and
\[
\frac{6(G+M)}{\ell-1}+\frac{8(b+M)}{\ell^2}<\frac{3K\ell}{\ell-1}+K+8K+K\le\tfrac72K+10K<14K,
\]
using $G<K\ell/2$, $\ell/(\ell-1)\le7/6$, $6M/(\ell-1)\le M\le K$, $b<K\ell^2$ and $8M/\ell^2<K$. If
$\sigma(i)\in B$, we use Proposition~\ref{pkm:kill-unram} with $s=\kappa$, $t=s_i$: it suffices that
$b+M\le(\ell-3)(\kappa+s_i)$. Now $\kappa+s_i\ge7K+q$, and $(\ell-3)q\ge(\ell-3)(b/\ell-1)=b-3b/\ell-\ell+3
>b-3K\ell-\ell+3$, while $(\ell-3)\cdot7K\ge4K\ell$ (as $\ell-3\ge4\ell/7$). So
$(\ell-3)(7K+q)>b+K\ell-\ell+3\ge b+K+2\ge b+M$, since $K\ell-\ell-K+1=(K-1)(\ell-1)\ge0$.

\emph{(C).} Let $i$ be a cut, $l=l(i)\le2$. In Lemma~\ref{pkd:steps} take $B_0$: $b_j+s_i+V+2\eta\le
b+q+M\le B_0$. At level $0$, $r_0=s_i+V+2\eta\le s_i+M$ and $s_i\le q<K\ell$, so
$(2r_0+1)/(\ell-1)\le 2s_i/(\ell-1)+(2M+1)/(\ell-1)<\tfrac73K+\tfrac12K<3K$ (using $2M+1\le3K$); the
level-$0$ term is $<6K+2\pkilog_\ell B_0$ (the factor $d_v+3a$ is $0$). At levels $a=1,2$,
$2r_a+1\le2M+1\le3K$, so the term is $\le K+2\pkilog_\ell B_0+18aK$. The total is
$<6K+2K+54K+6\pkilog_\ell B_0=62K+6\pkilog_\ell B_0$, and $\kappa+62K+6\pkilog_\ell B_0\le e^*$.

\emph{(G).} Let $i$ be a letter, $l=l(i)\in\{1,2\}$. If $l=1$, or $l=2$ and
$\varepsilon_{i^{(-1)}}=\varepsilon_i$, then $i$ is plain with $\varrho=\varepsilon_i$. Let $l=2$ and
$\varepsilon_{i^{(-1)}}\ne\varepsilon_i$. Then $i^{(-1)}$ is not a cut, so $i^{(-1)}\notin B$ (thus
$\beta_{i^{(-1)}}=G$) and $\sigma(i^{(-1)})=i\in B$; and $s_{i^{(-1)}}=s^+\ge G>0$, so $G+7K\le q$. Also
$\sigma(i^{(-2)})=i^{(-1)}\notin B$ gives $s_{i^{(-2)}}=0$. We show that $i$ is a long gap with pivot
$t=0$ and data $(1,s_0,\omega,A,\delta)=(1,q+12K,G+M,b+V,2K)$. The chart is unramified, $\ell\ge7$,
$0+1+1=l$, $s_{i^{(-2)}}=0$. The only $a<2$, $a\ne0$, is $a=1$, where
$\varepsilon_{i^{(-1)}}\ne\varepsilon_i$. $\eta\le G+M$, and $\beta_{i^{(-1)}}+V=G+V\le G+M$. Next, from
$G+M+6K\le G+7K\le q$:
\[
N=q+12K+(G+M)(1+\ell)\le q+12K+(q-6K)(1+\ell)=(\ell+2)q-6K(\ell-1)
<b+2K\ell-6K\ell+6K<b=\beta_i,
\]
using $\ell q\le b$ and $q\le b/\ell<K\ell$. $E^{s_0}$ kills the cokernels of $S$, $S'$ at $i$ by
Proposition~\ref{pkm:kill-unram} ($s=0$, $t=s_0$): $(\ell-3)(q+12K)>b-3K\ell-\ell+3+6K\ell\ge b+M$, as
above. $\beta_i+V=b+V=A$, and $A\le b+M<K\ell^2+K\le2K\ell^2=\ell^2\delta$. The cost: $\sum_{0<a<2}a=1$,
and $\kappa+4\delta+3=22K+3\le e^*$.

Theorem~\ref{pkd:scheme} gives $\mathrm{SCC}_H(c,x)$.
\end{proof}

\subsection{Windows, jumps and characters of the data: proof of Proposition~\ref{pki:data}}\label{pkd:ss-data}

We use the data of \S\ref{pg:sec} and of \eqref{eq:au-scales}: $h=r^8$, $\Delta=2h$, $t=5r$,
$k=r+1-m$, $\sqrt x\le m\le2\sqrt x$, $x^5\le r\le3x^5$, the weight $\kappa(\tau)=b$ at the $b$-places
and $\kappa(\tau)=G$ otherwise, and $\kappa_j(\tau)=\kappa(\tau)+(j+\frac12)\Delta=\kappa(\tau)+j\Delta+h$.
The summands $L_c$ are as in Proposition~\ref{pll:main} and the paragraph before it: the main summand
$L_{\mathrm{main}}$, of rank $r'\le r$, has at $\gamma$ the jumps $J_{\sigma_\gamma}$ of
Definition~\ref{pg:def-summand}(M2) (with multiplicity), where $\sigma_\gamma$ is the complex embedding
attached to $\gamma$ through $\iota$ and $v$ (Convention~\ref{pkj:conv}, reading (R$^-$)), and the characters $L_e$, $0\le e<t'\le t$, have the single jump $e$. $U(\gamma)$ is
the union of the jumps at $\gamma$ of all summands. For a lattice $L$ we use the convention of
Lemma~\ref{pkp:tau} and Definition~\ref{pkp:def-pairing}: $\auO(1)$ has the jump $-1$, so
$L^*=\operatorname{Hom}(L,\auO(1))$ has the jump $-u-1$ (with the multiplicity of $u$) for each jump
$u$ of $L$. Let $\varepsilon(\gamma)=1$ if $\sigma_\gamma$ is in the CM type of the data and $-1$
otherwise, $\eta=h^2+1$, $W=4t$, $\delta_j=j\Delta+h-2t$, and $a_\gamma=\kappa(\sigma_\gamma)$ ($=b$
if $\gamma$ is thick, $=G$ otherwise, in the sense of Proposition~\ref{pki:chart}).

\begin{lemma}[Elementary sizes]\label{pkd:sizes}
$t\le h/3$, $4m h\le h^2$, $t\le h^2$, and $h\ge4$.
\end{lemma}

\begin{proof}
$r\ge x^5\ge10^2$ (as in the proof of Lemma~\ref{pki:data-sizes}), so $h=r^8\ge15r=3t$; $4m\le
8\sqrt x\le8r\le h$; $t\le h\le h^2$.
\end{proof}

\begin{theorem}[Proposition~\ref{pki:data}, remaining assertions]\label{pkd:data}
Let $j\le m$ be the index of Definition~\ref{pg:def-summand}(M2) for $L_{\mathrm{main}}$, and let
$U'(\gamma)$ be the jump set of $L_c^*$.
\begin{enumerate}
\item[(i)] For every summand $L_c$ and every $\gamma$: $U(\gamma)\subseteq\pkiWin(\eta,-\varepsilon(\gamma),
a_\gamma+\delta_j,W)$ and $U'(\gamma)\subseteq\pkiWin(\eta,\varepsilon(\gamma),a_\gamma+\delta_j,W)$.
\item[(ii)] $L_{\mathrm{main}}^*$ has few low jumps for $(h,\lfloor r'm/k\rfloor,\eta)$ and the jump sets
$U$, $U'$ (with $c=\mathrm{main}$).
\item[(iii)] For a character $L_e$, $U'(\gamma)=\{-e-1\}\subseteq[-\eta,0]$.
\item[(iv)] Let $\ell=p_i$ and $v\mid\ell$. Then $E_v=\Q_\ell$, $\sigma_\gamma$, $\varepsilon(\gamma)$ and
$a_\gamma$ do not depend on $\gamma$, and $\ell>a+\delta_j+6(W+\eta)+r^2(N+1)^2+2$ for $a\in\{b,G\}$,
$N=\max(r+t,2r)\,r$ $(=(r+t)r=6r^2$, as $t=5r)$. If, for every inertial $s$, $q_s$ denotes the monic polynomial over $\Z_\ell$ of degree
$\le N$ of the definition of $\mathcal C(U,\mathrm{tame})$ (\S\ref{pll:sec}), which depends only on $s$, and
$K(s)=\{X^{e'}-1,q_s\}$, then $\mathcal C_v\subseteq\pkiC(U,K)$; and $L_c^*\in\pkiC(U',K^{-1})$, by
the input (I6) below.
\end{enumerate}
\end{theorem}

The inclusion in (iv) uses the definition of $\mathcal C(U,\mathrm{tame})$ with polynomials $q_s$
depending only on $s$, which (G4) provides (Proposition~\ref{pkw:G4}); the last assertion of (iv)
uses (I6), which is Lemma~\ref{pkw:dual}.

\begin{enumerate}
\item[(I6)] \emph{Inertia on duals.} If $Y$ is a potentially crystalline lattice, $s$ is inertial and
$P\in\auO[X]$ kills $s^{-1}$ on $D_{\mathrm{pcris}}(Y)$, then $P$ kills $s$ on
$D_{\mathrm{pcris}}(\operatorname{Hom}(Y,\auO(1)))$. This is Lemma~\ref{pkw:dual}; the jumps of
$\operatorname{Hom}(Y,\auO(1))$ are the $-u-1$ by Proposition~\ref{pkw:HT}(ii).
\end{enumerate}

\begin{proof}
Write $S_\gamma=a_\gamma+\delta_j$, so $\kappa_j(\sigma_\gamma)=S_\gamma+2t$, and $\varepsilon=\varepsilon(\gamma)$.

\emph{The jumps of $L_{\mathrm{main}}$.} By (M2) each $u\in J_{\sigma_\gamma}$ is low,
$\Delta/4\le|u|\le2m\Delta=4mh$, or high, $|u+\varepsilon\kappa_j(\sigma_\gamma)|\le t$ (the sign $+$
for $\sigma$ in the CM type). If low, $|u|\le h^2<\eta$ (Lemma~\ref{pkd:sizes}). If high,
$-\varepsilon u\in[S_\gamma+t,S_\gamma+3t]$, so $u\in-\varepsilon[S_\gamma,S_\gamma+W]$. In both cases
$|u|\ge h/2$: $\Delta/4=h/2$, and $\kappa_j-t\ge h-t\ge h/2$.

(i) The characters have jumps $e\in[0,t)\subseteq[-\eta,\eta]$. So
$U(\gamma)\subseteq\pkiWin(\eta,-\varepsilon,S_\gamma,W)$. For $U'(\gamma)$: if $c$ is a character,
$-e-1\in[-t,-1]\subseteq[-\eta,\eta]$. If $c=\mathrm{main}$ and $u$ is low, $|-u-1|\le4mh+1\le\eta$. If
$u$ is high and $\varepsilon=1$, then $-u\in[S_\gamma+t,S_\gamma+3t]$ and $-u-1\in[S_\gamma,S_\gamma+W]$
(as $t\ge1$); if $\varepsilon=-1$, then $u\in[S_\gamma+t,S_\gamma+3t]$ and
$-u-1\in-[S_\gamma+t+1,S_\gamma+3t+1]\subseteq-[S_\gamma,S_\gamma+W]$. So
$U'(\gamma)\subseteq\pkiWin(\eta,\varepsilon,S_\gamma,W)$.

(ii) A negative element of $U(\gamma)$ is a jump of $L_{\mathrm{main}}$ (the characters' jumps are
$\ge0$), so $2u\le-h$. A negative element of $U'(\gamma)$ is $-u-1$ with $u\in J_{\sigma_\gamma}$,
$u\ge0$, hence $u\ge h/2$ and $2(-u-1)\le-h-2$. A jump $-u-1$ of $L_{\mathrm{main}}^*$ lies in
$[-\eta,\eta]$ only if $|u|\le\eta+1=h^2+2\le2h^2$; by (M2) at most $r'm/k$ elements of
$J_{\sigma_\gamma}$, with multiplicity, satisfy this, and the number is an integer, so at most
$\lfloor r'm/k\rfloor$.

(iii) Clear, as $t\le\eta$.

(iv) $p_i\equiv1\pmod{q_2^2}$ with $q_1d\mid q_2$ splits completely in $E$ (\S\ref{pg:sec}), so
$E_v=\Q_\ell$; every $\gamma$ restricts to the unique embedding of $E_v$, so the complex embedding
$\sigma_\gamma$ attached to $\gamma$ through $\iota$ and $v$, and with it $\varepsilon(\gamma)$ and
$a_\gamma$, is independent of $\gamma$. Size: $a\le b$, $\delta_j\le2mh+h\le h^2$, $W+\eta\le2h^2$,
$r^2(N+1)^2\le r^2(6r^2+1)^2\le49r^6\le h$ (here $N=\max(r+t,2r)r=6r^2$ and $r\ge10^2$), so the right side is at most $b+14h^2+3\le2b<b^{10}\le p_i=\ell$
(as $h^3\le b$ by (S2)). Inclusion: a member of $\mathcal C_v=\mathcal C(U,\mathrm{tame})$ is potentially
crystalline with jumps in $U$, and on its $D_{\mathrm{pcris}}$ every inertial $s$ is killed by
$X^{e'}-1$ and by $q_s$, i.e.\ by every element of $K(s)$; and $K(s)$ contains $X^{e'}-1$ and the
image of the monic $q_s$ of degree $\le N$, as Proposition~\ref{pki:tame} requires. Finally, if
$L_c\in\mathcal C_v$, then for inertial $s$ every element of $K(s^{-1})$ kills $s^{-1}$ on
$D_{\mathrm{pcris}}(L_c)$, hence by (I6) kills $s$ on $D_{\mathrm{pcris}}(L_c^*)$; and $L_c^*$ is
potentially crystalline with jumps in $U'$ by definition. So $L_c^*\in\pkiC(U',K^{-1})$.
\end{proof}

\begin{remark}[Use in the deduction of Theorem~\ref{pll:exponent}]\label{pkd:rem-deduce}
In the proof of Theorem~\ref{pll:exponent} given Propositions~\ref{pki:generic}--\ref{pki:data}
(\S\ref{pki:ss-deduce}), case (a) applies Proposition~\ref{pki:tame} to $L=L_c^*$, which requires
$L_c^*\in\pkiC(U',K^{-1})$: this is the last assertion of Theorem~\ref{pkd:data}(iv), and it is part of the last clause of
Proposition~\ref{pki:data}. A cup exponent of
$\pkiC(U,K)$ on $L$ is one of every subclass, in particular of $\mathcal C_v$. With
Theorems~\ref{pkd:sparse}, \ref{pkd:any}, \ref{pkd:data} and \S\ref{pkc:sec}, the deduction of
Theorem~\ref{pll:exponent} in \S\ref{pki:ss-deduce} rests only on Proposition~\ref{pki:chart} (Theorem~\ref{pkj:chart}) and on the
inputs (I1)--(I6), (T1)--(T4), all proved in \S\S\ref{pkt:sec}--\ref{pkw:sec}.
\end{remark}


%% ===== proofs/C-kisin-tame.tex =====
% proofs/C-kisin-tame.tex
% Tame charts, inertia on D_pcris, and the tame inputs (T1)--(T4), (I6) of the local-loss engine.
% Labels start with pkw:. Conventions and macros as in proofs/C-kisin-pairing.tex,
% proofs/C-kisin-modules.tex, proofs/C-kisin-inputs.tex, proofs/C-kisin-cases.tex.
\providecommand{\auO}{\mathcal O}
\providecommand{\auQlb}{\overline{\Q}_\ell}
\providecommand{\pkpA}{A_{\mathrm{inf}}}
\providecommand{\pkpW}{W(C^\flat)}
\providecommand{\pkpAmb}{\mathbb A}
\providecommand{\pkmJ}[3]{J^{#1}_{#2}(#3)}
\providecommand{\pkmI}[2]{I_{#1}(#2)}
\providecommand{\pkiC}{\mathcal C}
% macros of this file (prefix \pkw)
\providecommand{\pkwB}{B_{\mathrm{cris}}}
\providecommand{\pkwD}{D_{\mathrm{pcris}}}
\providecommand{\pkwur}{\Q_\ell^{\mathrm{ur}}}
\providecommand{\pkwI}{I_{\Q_\ell}}

\section{Tame charts, inertia on $D_{\mathrm{pcris}}$, and the tame inputs}\label{pkw:sec}

This section proves the inputs (T1)--(T4) of \S\ref{pkc:ss-tame} and the input (I6) of
\S\ref{pkd:ss-data}, and settles two points of convention used in \S\ref{pkd:sec}: the jumps of
Tate duals, and the shape of the class $\mathcal C(U,\mathrm{tame})$. We use the conventions and the
labels of \S\S\ref{pkp:sec}--\ref{pkd:sec}.

\emph{Dependencies.} Proposition~\ref{pkw:T1} (input (T1)) and Lemma~\ref{pkw:letters} use only
elementary local field theory. Theorem~\ref{pkw:T2} (input (T2)), Theorem~\ref{pkw:T4} (input (T4)),
Lemma~\ref{pkw:dual} (input (I6)), Proposition~\ref{pkw:HT} and Proposition~\ref{pkw:G4} use the published facts
of $p$-adic Hodge theory and of local ramification theory stated in \S\ref{pkw:ss-pcris} and
\S\ref{pkw:ss-checks}. Theorem~\ref{pkw:Z} and Theorem~\ref{pkw:T3} (input (T3)) use the results of
\S\S\ref{pkn:sec}--\ref{pkl:sec} and (F6).

\subsection{$D_{\mathrm{pcris}}$ and the operators $\beta_s$}\label{pkw:ss-pcris}

$\Gamma=D_v$ is an open subgroup of $G_{\Q_\ell}$ through $j$; $\pkwI\subseteq G_{\Q_\ell}$ is the inertia
group, i.e.\ $\operatorname{Gal}(\auQlb/\pkwur)$ with $\pkwur$ the maximal unramified extension of
$\Q_\ell$ in $\auQlb$; $s\in\Gamma$ is \emph{inertial} if $j(s)\in\pkwI$. $\pkwB$ is Fontaine's
crystalline period ring, with its action of $G_{\Q_\ell}$, $t\in\pkwB$ the element with $gt=
\chi_{\mathrm{cyc}}(g)t$ ($\chi_{\mathrm{cyc}}$ the cyclotomic character), a unit of $\pkwB$, and
$P_0=\widehat{\pkwur}\subseteq\pkwB$ the completion of $\pkwur$ (the fraction field of
$W(\overline{\mathbb F}_\ell)$). For an open subgroup $H\subseteq G_{\Q_\ell}$ we write $K_H$ for its fixed
field and $K_{H,0}$ for the maximal unramified subfield of $K_H$. For a lattice $Y$ (a finite free
$\auO$-module with a continuous $\auO$-linear action of $\Gamma$) we put $\pkwB(Y)=\pkwB\otimes_{\Z_\ell}Y$,
with $g(b\otimes y)=gb\otimes gy$ for $g\in\Gamma$ and $\auO$ acting through $Y$.

We use the following published facts, with their hypotheses.

\begin{enumerate}
\item[(F1)] (Fontaine \cite[1.4.1--1.4.2, 5.1.2 and 5.6.1]{pkfFontaineIII}, \emph{Repr\'esentations $p$-adiques semi-stables}.) For every open subgroup $H\subseteq G_{\Q_\ell}$, $\pkwB^{H}=K_{H,0}$; and for
every continuous representation of $H$ on a finite dimensional $\Q_\ell$-vector space $V$, the natural
map $\pkwB\otimes_{K_{H,0}}(\pkwB\otimes_{\Q_\ell}V)^H\to\pkwB\otimes_{\Q_\ell}V$ is injective. $V$ is
\emph{crystalline} (as a representation of $H$) if this map is also surjective, equivalently
$\dim_{K_{H,0}}(\pkwB\otimes V)^H=\dim_{\Q_\ell}V$.
\item[(F2)] (Proved in Lemma~\ref{pkw:F2} below from (F1); it is the form, for crystalline periods, of
Fontaine's criterion \cite[1.8.5]{pkfFontaineIII} for potentially admissible representations,
together with Galois descent.) Let $H\subseteq
G_{\Q_\ell}$ be open and $V$ a representation of $H$ that is crystalline on some open subgroup of $H$. If
$H\cap\pkwI$ acts trivially on $\bigcup_{H'}(\pkwB\otimes V)^{H'}$ ($H'$ running over the open subgroups
of $H$), then $V$ is crystalline as a representation of $H$.
\item[(F3)] (Serre \cite[IV, \S2, Proposition~7 and Corollary~1, p.~67]{SerreLF}, the structure of tamely ramified extensions.) Every finite extension of $\pkwur$ inside $\auQlb$ whose degree $m$ is prime to $\ell$ is
$\pkwur(\varpi_m)$ for an element $\varpi_m$ with $\varpi_m^m=\ell$. ($\pkwur$ is henselian, discretely
valued with uniformiser $\ell$ and algebraically closed residue field, so such an extension $L$ is totally
and tamely ramified. If $\pi_L$ is a uniformiser of $L$, then $\pi_L^m=\ell u$ with $u$ a unit of the ring of
integers of $L$; the residue field of $L$ is algebraically closed and $\ell\nmid m$, so by Hensel's lemma
$u=w^m$ with $w\in L$, and $\varpi_m=\pi_L/w$ satisfies $\varpi_m^m=\ell$ and generates $L$, as
$X^m-\ell$ is an Eisenstein polynomial of degree $m=[L:\pkwur]$. The cited statements of Serre give the
cyclicity of the tame inertia quotient.)
\end{enumerate}

\begin{lemma}[(F2)]\label{pkw:F2}
Statement (F2) holds.
\end{lemma}

\begin{proof}
Let $H'_0\subseteq H$ be open with $V$ crystalline on $H'_0$, and $H'=\bigcap_{h\in H}hH'_0h^{-1}$, a
finite intersection, hence open and normal in $H$. $V$ is crystalline on $H'$: by (F1) the map
$\pkwB\otimes_{K_{H'_0,0}}(\pkwB\otimes V)^{H'_0}\to\pkwB\otimes V$ is surjective, and
$(\pkwB\otimes V)^{H'_0}\subseteq(\pkwB\otimes V)^{H'}$, so $\pkwB\otimes_{K_{H',0}}(\pkwB\otimes V)^{H'}
\to\pkwB\otimes V$ is surjective, and injective by (F1). Put $L=K_{H'}$, $F'=K_{H',0}=\pkwB^{H'}$ (F1),
$D'=(\pkwB\otimes V)^{H'}$, an $F'$-vector space.

Since $H'$ is normal in $H$, $H$ preserves $F'$ and $D'$, and acts on $D'$ semilinearly with respect to its
action on $F'$. We claim that $H\cap G_{F'}=(H\cap\pkwI)H'$. The inclusion $\supseteq$ holds as $\pkwI$
fixes the unramified field $F'$ and $H'=G_L$ fixes $L\supseteq F'$. Conversely let $g\in H\cap G_{F'}$.
As $L\cap\pkwur=F'$ ($F'$ is the maximal unramified subfield of $L$) and $\pkwur/\Q_\ell$ is Galois, restriction
maps $G_L$ onto $\operatorname{Gal}(\pkwur/F')$; so some $h\in G_L=H'$ agrees with $g$ on $\pkwur$, and
$h^{-1}g\in\pkwI\cap H$. This proves the claim.

By hypothesis $H\cap\pkwI$ acts trivially on $D'\subseteq\bigcup_{H''}(\pkwB\otimes V)^{H''}$, and $H'$ acts
trivially on $D'$. Hence the action of $H$ on $D'$ and on $F'$ factors through the finite group
$\Delta=H/(H\cap G_{F'})$, which acts faithfully on $F'$. By Artin's theorem $F'/F$ is Galois with group
$\Delta$, where $F=F'^{\Delta}=(\pkwB^{H'})^H=\pkwB^H=K_{H,0}$ (F1).

\emph{Galois descent.} We show that $F'\otimes_FD'^{\Delta}\to D'$ is surjective (it is all we need). Let
$D_1$ be the $F'$-span of $D'^\Delta$ and suppose $D_1\ne D'$; choose an $F'$-linear $f\colon D'\to F'$ with
$f(D_1)=0$ and $f(d)\ne0$ for some $d\in D'$. For every $\lambda\in F'$, $\sum_{\delta\in\Delta}\delta(\lambda d)
\in D'^\Delta$, so $0=\sum_\delta\delta(\lambda)f(\delta d)$ for all $\lambda\in F'$. The distinct field
automorphisms $\delta$ of $F'$ are linearly independent over $F'$ (Dedekind--Artin), so $f(\delta d)=0$ for
every $\delta$, in particular $f(d)=0$, a contradiction.

Now $D'^\Delta=(\pkwB\otimes V)^H$, and the composite
$\pkwB\otimes_{K_{H,0}}(\pkwB\otimes V)^H\to\pkwB\otimes_{F'}D'\to\pkwB\otimes V$ is surjective, the
first map by Galois descent and the second because $V$ is crystalline on $H'$. By (F1), $V$ is crystalline
on $H$.
\end{proof}

\begin{definition}[$D_{\mathrm{pcris}}$, $\beta_s$, killing]\label{pkw:def-pcris}
Let $Y$ be a lattice and $V=Y[1/\ell]$. Put $\pkwD(Y)=\bigcup_H\pkwB(Y)^H$, the union over the open
subgroups $H$ of $\Gamma$. $Y$ is \emph{potentially crystalline} if $\pkwB(Y)^H$ spans $\pkwB(Y)$ over
$\pkwB$ for some open $H\subseteq\Gamma$ (by (F1), equivalently $V$ is crystalline on $H$). For such $Y$
and an inertial $s\in\Gamma$, $\beta_s=\beta_s^Y$ denotes the unique $\pkwB$-linear endomorphism of
$\pkwB(Y)$ that agrees with $s$ on $\pkwD(Y)$ (Lemma~\ref{pkw:beta}(i)). For $P=\sum_kP_kX^k\in\auO[X]$,
``$P$ \emph{kills $s$ on $\pkwD(Y)$}'' means $P(\beta_s)=\sum_kP_k\beta_s^k=0$ on $\pkwB(Y)$,
equivalently (Lemma~\ref{pkw:beta}(i)) $\sum_kP_ks^k=0$ on $\pkwD(Y)$.
\end{definition}

This is the meaning of the phrase in Proposition~\ref{pki:tame}, in \S\ref{pkc:ss-tame} and in
Theorem~\ref{pkd:data}, and in the definition of $\mathcal C(U,\mathrm{tame})$ in \S\ref{pll:ss-statement}.

\begin{lemma}[The operators $\beta_s$]\label{pkw:beta}
Let $Y$, $Y'$ be potentially crystalline lattices and $s,s'\in\Gamma$ inertial.
\begin{enumerate}
\item[(i)] $\beta_s$ exists and is unique; it commutes with $\auO$; $\beta_{ss'}=\beta_s\beta_{s'}$ and
$\beta_1=1$, so $\beta_s$ is invertible with inverse $\beta_{s^{-1}}$. If $P\in\auO[X]$, then
$P(\beta_s)=0$ if and only if $\sum_kP_ks^k$ vanishes on $\pkwD(Y)$. If $s\in H$ for an open $H$ such that
$\pkwB(Y)^H$ spans $\pkwB(Y)$, then $\beta_s=1$.
\item[(ii)] If $f\colon Y\to Y'$ is $\auO$-linear and $\Gamma$-equivariant, then $(1\otimes f)\beta^Y_s=
\beta^{Y'}_s(1\otimes f)$. If $f$ is injective and $P$ kills $s$ on $\pkwD(Y')$, then $P$ kills $s$ on
$\pkwD(Y)$.
\item[(iii)] If $\Gamma$ acts trivially on $Y$, then $Y$ is potentially crystalline and $\beta_s=1$. The
lattice $\auO(1)=\auO\otimes_{\Z_\ell}\Z_\ell(1)$ is potentially crystalline and $\beta_s=1$ on it.
\end{enumerate}
\end{lemma}

\begin{proof}
(i) Let $H_0\subseteq\Gamma$ be open with $\pkwB(Y)^{H_0}$ spanning, and let $H\subseteq H_0$ be an open
subgroup normal in $\Gamma$ (the intersection of the finitely many conjugates of $H_0$). Then
$D_H=\pkwB(Y)^H\supseteq\pkwB(Y)^{H_0}$ spans, so by (F1) the map $\pkwB\otimes_{K_{H,0}}D_H\to\pkwB(Y)$
is bijective. As $H$ is normal, $s$ maps $D_H$ into itself; as $s$ is inertial it fixes
$K_{H,0}\subseteq\widehat{\pkwur}$, so $s|_{D_H}$ is $K_{H,0}$-linear, and also $\auO$-linear. Define
$\beta_s$ as $1\otimes s|_{D_H}$ transported by this bijection; it is $\pkwB$-linear, commutes with
$\auO$ and agrees with $s$ on $D_H$. If $x\in\pkwD(Y)$, say $x$ fixed by an open $H'$, let $H''\subseteq
H\cap H'$ be open and normal in $\Gamma$; the operator built from $H''$ agrees with $s$ on $D_{H''}\supseteq
D_H\cup\{x\}$, hence with $\beta_s$ on the spanning set $D_H$, hence everywhere; so $\beta_sx=sx$.
Uniqueness holds because $\pkwD(Y)$ spans. $\beta_s\beta_{s'}$ and $\beta_{ss'}$ are $\pkwB$-linear and
agree with $ss'$ on $\pkwD(Y)$ (which $s'$ maps into itself: $s'$ maps the vectors fixed by $H'$ to
vectors fixed by $s'H's'^{-1}$); so they are equal. For $P\in\auO[X]$, $P(\beta_s)$ is $\pkwB$-linear and
agrees with $\sum_kP_ks^k$ on the spanning set $\pkwD(Y)$. If $s\in H$ with $\pkwB(Y)^H$ spanning, $\beta_s$
and $1$ agree on $\pkwB(Y)^H$.

(ii) $1\otimes f$ is $\pkwB$-linear and equivariant, so it maps $\pkwD(Y)$ into $\pkwD(Y')$; both sides of
the identity are $\pkwB$-linear and agree on $\pkwD(Y)$. If $f$ is injective, so is $1\otimes f$
($\pkwB$ is flat over $\Z_\ell$), and $(1\otimes f)P(\beta^Y_s)=P(\beta^{Y'}_s)(1\otimes f)=0$.

(iii) If the action is trivial, $1\otimes Y\subseteq\pkwB(Y)^\Gamma$ spans, so $\beta_s=1$ by (i). For
$\auO(1)$ let $e_1$ be a basis of $\Z_\ell(1)$, so $ge_1=\chi_{\mathrm{cyc}}(g)e_1$; then
$g(t^{-1}\otimes a\otimes e_1)=\chi_{\mathrm{cyc}}(g)^{-1}t^{-1}\otimes a\otimes\chi_{\mathrm{cyc}}(g)e_1=
t^{-1}\otimes a\otimes e_1$ for $a\in\auO$, and these vectors span ($t$ is a unit); so $\Gamma$ itself is an
admissible $H$ and $\beta_s=1$ by (i).
\end{proof}

\begin{lemma}[Extensions]\label{pkw:ext}
Let $S$ be a side (Definition~\ref{pkp:def-side}) whose lattice $Y$ is potentially crystalline, and
let $q\in\auO[X]$ kill an inertial $s$ on $\pkwD(Y)$. Then the extension $X_S$ is potentially crystalline
and $q(X)\cdot(X-1)$ kills $s$ on $\pkwD(X_S)$.
\end{lemma}

\begin{proof}
$X_S$ is crystalline on $\Gamma'$ (Definition~\ref{pkp:def-side}: it is a lattice with jump sets),
hence potentially crystalline. The inclusion $\iota\colon Y\to X_S$ and the projection $\pi\colon X_S\to
C_S$ are equivariant, and $0\to\pkwB(Y)\to\pkwB(X_S)\to\pkwB(C_S)\to0$ is exact. By
Lemma~\ref{pkw:beta}(ii),(iii), $(1\otimes\pi)(\beta_s-1)=(\beta^{C_S}_s-1)(1\otimes\pi)=0$, so for
$x\in\pkwB(X_S)$, $(\beta_s-1)x=(1\otimes\iota)y$ with $y\in\pkwB(Y)$, and
$q(\beta_s)(\beta_s-1)x=(1\otimes\iota)q(\beta^Y_s)y=0$.
\end{proof}

\subsection{(T1): tame charts}\label{pkw:ss-T1}

\emph{The field of a chart.} In Definition~\ref{pkp:def-chart} the ramification index $e$, the residue
degree $f$ and the uniformiser $\pi=\theta(z)$ are those of the fixed field
$K'=\auQlb^{\,j(\Gamma')}$ of the subgroup $\Gamma'$ of the chart, not of $\auQlb^{\,j(\Gamma)}$: this is how
the chart is used throughout (the conditions on $\Gamma'$ say exactly that $j(\Gamma')=G_{K'}$ with
$K'=K'_0(\pi)$, $K'_0=\Lambda[1/\ell]$, as in Lemma~\ref{pkf:invariants}), and the two readings agree
when $\Gamma'=\Gamma$, as in Lemma~\ref{pkp:chart-exists}. A tame chart has $\Gamma'\ne\Gamma$ in general,
and needs this reading, which is the one stated in Definition~\ref{pkp:def-chart}.

\begin{definition}[Tame chart]\label{pkw:def-tame}
Let $E_v=\Q_\ell$, so that $j\colon\Gamma\to G_{\Q_\ell}$ is bijective. A \emph{tame chart} is a Kummer chart
(Definition~\ref{pkp:def-chart}, read with $K'$ as above) with ramification index $e$ and residue degree
$f$, together with an element $\varpi\in\auQlb$ and an element $s_0\in\Gamma$ (the \emph{tame generator}),
such that: $\ell\nmid e$; $\varpi^e=\ell$, $\theta(z)=\varpi$ and $E=\ell-z^e$; $\Gamma'$ is normal in $\Gamma$; $s_0$ is
inertial and fixes $\Lambda$ pointwise; and $s_0z=\xi z$ with $\xi=[a]\in\Lambda$ the Teichm\"uller lift of
a primitive $e$-th root of unity $a\in\mathbb F_{\ell^f}^\times$. Then $s_0^kz=\xi^kz$ for all $k\ge0$
($s_0$ fixes $\xi$), and $d_v\le1$ ($d_v=v_\ell(e)+1=1$ if $e>1$).
\end{definition}

\begin{lemma}[Letters for any root of $X^{\ell^f-1}-1$]\label{pkw:letters}
In Definition~\ref{pkp:def-letters} let $\zeta\in\auO$ be any element with $\zeta^{N}=1$, $N=\ell^f-1$.
Then the $e_i$ defined there are orthogonal idempotents of $\operatorname{End}(\pkpAmb(L))$ with
$\sum_ie_i=1$ and $(1\otimes\zeta)e_i=[i]e_i$; the image of $e_i$ is the letter part
$\pkpAmb(L)_i=\{x:(1\otimes\zeta)x=[i]x\}$, so $\pkpAmb(L)=\bigoplus_i\pkpAmb(L)_i$; $\varphi e_i=
e_{i^\ell}\varphi$, and $ge_i=e_ig$ for every $g\in G_{\Q_\ell}$ fixing $\Lambda$ (in particular for
$g\in\Gamma'$). For $\zeta=1$: $\pkpAmb(L)_1=\pkpAmb(L)$ and $\pkpAmb(L)_i=0$ for $i\ne1$, and every
embedding $\gamma\colon\auO\to C$ has letter $1$.
\end{lemma}

\begin{proof}
Write $T=1\otimes\zeta$ and $\lambda_i=[i]$. In $\Lambda[X]$, $X^N-1=\prod_{i}(X-\lambda_i)$, the product
over the $N$ letters (the $\lambda_i$ are $N$ distinct roots of $X^N-1$ in the domain $\Lambda$); hence
$\prod_i(T-\lambda_i)=T^N-1=0$. Let $e_i(X)=\prod_{i'\ne i}(X-\lambda_{i'})/(\lambda_i-\lambda_{i'})\in
\Lambda[X]$ (the denominators are units). The polynomial $\sum_ie_i(X)-1$ has degree $<N$ and vanishes at
the $N$ points $\lambda_{i'}$, whose differences are units, so it is $0$ (the Vandermonde determinant is a
unit). $(X-\lambda_i)e_i(X)$ is a unit multiple of $\prod_{i'}(X-\lambda_{i'})$, so $(T-\lambda_i)e_i(T)=0$,
i.e.\ $Te_i=\lambda_ie_i$. Hence $e_j(T)e_i(T)=e_j(\lambda_i)e_i(T)=\delta_{ij}e_i(T)$, and with
$\sum_ie_i=1$ the $e_i=e_i(T)$ are orthogonal idempotents summing to $1$. If $Tx=\lambda_ix$ then
$x=\sum_je_jx=\sum_je_j(\lambda_i)x=e_ix$; so the image of $e_i$ is the letter part at $i$. $\varphi$ fixes
$T$ and maps $\lambda_{i'}$ to $\lambda_{i'^\ell}$, permuting the letters, so $\varphi e_i(T)=
e_{i^\ell}(T)\varphi$; an element $g$ fixing $\Lambda$ commutes with $T$ (as $\zeta$ acts through the
$\auO$-module structure, which commutes with $G_{\Q_\ell}$) and with the coefficients of $e_i(X)$. For
$\zeta=1$, $T=1$ and $(1-\lambda_i)$ is a unit for $i\ne1$, so $\pkpAmb(L)_i=0$ for $i\ne1$; and the
residue of $\gamma(1)=1$ is the letter $1$.
\end{proof}

Only the properties listed in Lemma~\ref{pkw:letters} are used in \S\S\ref{pkp:sec}--\ref{pkd:sec} about
the letters (the exact order of $\zeta$ is never used), so Definition~\ref{pkp:def-letters} may be stated
with $\zeta^{\ell^f-1}=1$; with $\zeta=1$ the jump set at the letter $1$ is $\{0\}\cup\bigcup_\gamma
U(\gamma)$ and every other jump set is the set $\{0\}$ that sides are required to contain (the other letter
parts being $0$, any finite set containing $0$ is admissible there). This avoids requiring $\auO$ to contain
roots of unity of order $\ell^f-1$ for the residue degree $f$ of a tame chart.

\begin{proposition}[(T1): existence of tame charts]\label{pkw:T1}
Let $E_v=\Q_\ell$, $f\ge1$ with $\ell\nmid f$, and $e\ge1$ with $e\mid\ell^f-1$. There is a tame chart with
ramification index $e$ and residue degree $f$. In particular, if $1\le r<\ell$ and
$e'=\operatorname{lcm}_{1\le a\le r}(\ell^a-1)$, there is a tame chart with ramification index $e'$ (and
residue degree $r!$); its $d_v$ is at most $1$, and $e'\le\ell^{r^2}$ (Lemma~\ref{pkc:lcm}).
\end{proposition}

\begin{proof}
\emph{The field.} Let $N=\ell^f-1$, fix $\varpi\in\auQlb$ with $\varpi^e=\ell$, and let
$K'=\Q_\ell(\mu_N,\varpi)$. $\Q_\ell(\mu_N)=\Q_{\ell^f}$ is unramified of degree $f$ with residue field
$\mathbb F_{\ell^f}$, and $X^e-\ell$ is Eisenstein over its ring of integers, so $K'$ has ramification index
$e$, residue degree $f$, uniformiser $\varpi$, and $[K':\Q_\ell]=ef$, which is prime to $\ell$ ($e\mid N$
and $\ell\nmid fN$). $K'$ contains $\mu_e\subseteq\mu_N$, hence all roots $\zeta\varpi$ of $X^e-\ell$; so $K'$
is the splitting field of $(X^N-1)(X^e-\ell)$ and is Galois over $\Q_\ell$. Put $\Gamma'=j^{-1}(G_{K'})$: an
open normal subgroup of $\Gamma$ of index $ef$, prime to $\ell$.

\emph{Fixing $\Lambda$.} For $g\in G_{\Q_\ell}$: $g$ fixes $\Lambda=W(\mathbb F_{\ell^f})\subseteq\pkpA$
pointwise iff $g$ fixes $\mu_N\subseteq\auQlb$ pointwise. Indeed $\Lambda$ is the $\ell$-adic closure of the
$\Z_\ell$-span of the Teichm\"uller lifts $[x]$, $x\in\mathbb F_{\ell^f}\subseteq O^\flat$, and $g[x]=[gx]$;
for $x\ne0$, $x=(x_n\bmod\ell)_n$ where $x_n\in\mu_N$ is the unique $\ell^n$-th root of $x_0=\omega(x)$
in $\mu_N$, and $x_n$ is a power of $x_0$; distinct elements of $\mu_N$ have distinct residues; so $gx=x$
iff $g$ fixes every $x_n$ iff $g$ fixes $\omega(x)$; and $\omega$ maps $\mathbb F_{\ell^f}^\times$ onto
$\mu_N$. Hence the elements of $G_{\Q_\ell}$ fixing $\Lambda$ pointwise and $\varpi$ are exactly those of
$G_{K'}$.

\emph{The chart.} Choose $\varpi_n\in\auQlb$ with $\varpi_0=\varpi$, $\varpi_{n+1}^\ell=\varpi_n$, put
$\varpi^\flat=(\varpi_n\bmod\ell)_n\in\mathfrak m^\flat$, $z=[\varpi^\flat]$, so $\theta(z)=\varpi$, and
$E=\ell-z^e=z^eA+\ell$ with $A=-1\in\mathfrak S^\times$. Then $\theta(E)=\ell-\varpi^e=0$, and $E$
generates $\ker\theta$ by Corollary~\ref{pkn:E-gen} (whose proof uses only $E=z^eA+\ell$ with $A$ a unit,
$z=[\varpi^\flat]$ and $\theta(E)=0$). By the previous paragraph $j(\Gamma')=G_{K'}$ fixes $\Lambda$ and
$\theta(z)$, and every element of $G_{\Q_\ell}$ fixing both lies in $j(\Gamma')$. Kummer elements: fix
$l\ge0$ and $\zeta_{\ell^n}$ with $\varepsilon=(\zeta_{\ell^n}\bmod\ell)_n$; the sets
$X_n=\{g\in G_{K'}:g\varpi_n=\zeta_{\ell^n}^l\varpi_n\}$ are closed, decreasing (raise to the $\ell$-th power)
and nonempty ($X^{\ell^n}-\varpi$ is Eisenstein over $O_{K'}$, hence irreducible, and $G_{K'}$ is transitive
on its roots); by compactness some $g\in G_{K'}$ lies in all $X_n$, and then $g\varpi^\flat=
\varepsilon^l\varpi^\flat$, $gz=[\varepsilon]^lz$. This is a Kummer chart (with $K'$ as above), with
ramification index $e$ and residue degree $f$; $\Gamma'$ is normal and $\ell\nmid e$.

\emph{The tame generator.} As $e\mid N$, $\mathbb F_{\ell^f}^\times$ contains a primitive $e$-th root of
unity $a$; let $a_n\in\mu_N$ be the unique $\ell^n$-th root of $\omega(a)$ in $\mu_N$, so
$a=(a_n\bmod\ell)_n\in O^\flat$, and put $y_n=a_n\varpi_n$. Then $y_{n+1}^\ell=y_n$ and
$y_n^{e\ell^n}=a_0^e\varpi^e=\ell$. The polynomial $X^{e\ell^n}-\ell$ is Eisenstein over the discrete
valuation ring of $\pkwur$, hence irreducible over $\pkwur$, and $\varpi_n$, $y_n$ are two of its roots;
so the closed sets $Y_n=\{g\in\pkwI:g\varpi_n=y_n\}$ are nonempty, and decreasing (raise to the $\ell$-th
power). $\pkwI$ is compact, so some $s\in\pkwI$ lies in all $Y_n$. Then $s\varpi^\flat=(a_n\varpi_n\bmod\ell)_n
=a\varpi^\flat$ and $sz=[a][\varpi^\flat]=\xi z$ with $\xi=[a]$. Being inertial, $s$ fixes $\mu_N\subseteq
\pkwur$, hence $\Lambda$. Put $s_0=j^{-1}(s)$.

\emph{The index $e'$.} If $1\le r<\ell$, then $\ell\nmid r!$, and for $1\le a\le r$, $a\mid r!$ gives
$\ell^a-1\mid\ell^{r!}-1$; so $e'\mid\ell^{r!}-1$, and the first part applies with $f=r!$, $e=e'$.
\end{proof}

\subsection{(T2): descent to crystalline over $\Gamma'$}\label{pkw:ss-T2}

\begin{lemma}[Tame inertia]\label{pkw:tame-inertia}
Let $e\ge1$ with $\ell\nmid e$, $\varpi\in\auQlb$ with $\varpi^e=\ell$, and let $\phi\colon\pkwI\to G$ be a
homomorphism to a group, with open kernel, such that $\phi(s)^e=1$ for all $s\in\pkwI$. Then $\phi(s)=1$ for
every $s\in\pkwI$ with $s\varpi=\varpi$.
\end{lemma}

\begin{proof}
$N=\ker\phi$ is open and normal in $\pkwI=\operatorname{Gal}(\auQlb/\pkwur)$; its fixed field $M$ is a finite
Galois extension of $\pkwur$ with group $\pkwI/N\cong\phi(\pkwI)$, in which every element has order dividing
$e$. Hence $m=[M:\pkwur]$ is prime to $\ell$ (otherwise, by Cauchy's theorem, some element would have order
$\ell\nmid e$). By (F3), $M=\pkwur(\varpi_m)$ with $\varpi_m^m=\ell$. Since $\mu_m\subseteq\pkwur$ ($\ell\nmid m$),
Kummer theory shows that $\operatorname{Gal}(M/\pkwur)$ is cyclic of order $m$; as its elements have order
dividing $e$, $m\mid e$. Both $\varpi_m$ and $\varpi^{e/m}$ are $m$-th roots of $\ell$, so
$\varpi_m=\zeta\varpi^{e/m}$ with $\zeta\in\mu_m\subseteq\pkwur$. Hence $M\subseteq\pkwur(\varpi)$, and an
$s\in\pkwI$ fixing $\varpi$ fixes $M$, i.e.\ $s\in N$.
\end{proof}

\begin{theorem}[(T2): descent]\label{pkw:T2}
Fix a tame chart (Definition~\ref{pkw:def-tame}) with ramification index $e$. Let $Y$ be a potentially
crystalline lattice such that $X^e-1$ kills every inertial $s\in\Gamma$ on $\pkwD(Y)$. Then $Y$ is
crystalline on $\Gamma'$. Its Hodge--Tate jumps are unchanged (they are those of $Y[1/\ell]$, which do not
depend on the open subgroup on which they are computed), so if $Y$ has jumps in $U$ it is a lattice with
jump sets over the chart (Definition~\ref{pkp:def-orders}).
\end{theorem}

\begin{proof}
As $j$ is bijective, $s\mapsto\beta_s$ is defined on $j^{-1}(\pkwI)\cong\pkwI$, and by
Lemma~\ref{pkw:beta}(i) it is a homomorphism $\phi\colon\pkwI\to\operatorname{Aut}_{\pkwB}(\pkwB(Y))$ whose
kernel contains $H\cap\pkwI$ for an open $H$ with $\pkwB(Y)^H$ spanning; so $\ker\phi$ is open. By
hypothesis $\phi(s)^e=1$. By Lemma~\ref{pkw:tame-inertia}, $\beta_s=1$ for every inertial $s$ with
$j(s)\varpi=\varpi$, in particular (Definition~\ref{pkw:def-tame}: $j(\Gamma')$ fixes $\theta(z)=\varpi$) for
every inertial $s\in\Gamma'$. So $j(\Gamma')\cap\pkwI=I_{K'}$ acts trivially on $\pkwD(Y)$ (which is the union
in (F2) for $H=j(\Gamma')$, as the open subgroups of $\Gamma'$ are cofinal among those of $\Gamma$), and by (F2)
$Y[1/\ell]$ is crystalline as a representation of $\Gamma'$.
\end{proof}

\subsection{(T3): the tame cokernel killer}\label{pkw:ss-T3}

Fix a tame chart with ramification index $e$, residue degree $f$, tame generator $s_0$ and $\xi=[a]$ as in
Definition~\ref{pkw:def-tame}, and letters with any $\zeta$, $\zeta^{\ell^f-1}=1$ (Lemma~\ref{pkw:letters}).
For a side $S$, $\pkmI li$ and $\pkmJ rli$ are the lift ideals of Definition~\ref{pkm:def-J}, and
$\varepsilon_i(c)=e_i(1\otimes c)$. For a lattice with jump sets $X$ we write
$\mathcal N_i(\mathrm{std}_0)=\mathcal N(X)\cap X_i(\mathrm{std}_0)$ (Lemma~\ref{pkm:inputs}).

The argument rests on the following statement (Z), proved in Theorem~\ref{pkw:Z} below.

\begin{enumerate}
\item[(Z)] \emph{Division by $z$.} Let $X$ be a lattice with jump sets over the tame chart, $i$ a letter and
$Q=\sum_kQ_kX^k\in\Z_\ell[X]$ with $Q(\beta_{s_0})=0$ on $\pkwB(X)$. Then for every
$x\in\mathcal N_i(\mathrm{std}_0)$ the vector $z^{-1}\sum_kQ_ks_0^kx$ ($z$ is a unit of $\pkpW$) lies in
$X_i(\mathrm{std}_0)$.
\end{enumerate}

We use one more published fact.
\begin{enumerate}
\item[(F6)] (Colmez \cite[\S III.2]{pkfColmez98}, Fontaine \cite[5.1.2]{pkfFontaineIII}.) $\pkwB\subseteq\pkfBm$,
compatibly with $G_{\Q_\ell}$, with $\varphi$ and with the maps to $\pknB$; and $\pkfBm^H=K_{H,0}$ for every open
$H\subseteq G_{\Q_\ell}$. (Colmez gives $\varphi(\pkfBm)\subseteq\pkwB\subseteq\pkfBm$; with $\pkwB^H=K_{H,0}$ (F1) and injectivity of $\varphi$ this gives $\pkfBm^H=K_{H,0}$, as in (P2) of \S\ref{pkf:ss-rings}.)
\end{enumerate}

\begin{lemma}[$s_0$ on $\pknMinf$ and on $\mathfrak S_B$]\label{pkw:s0}
Let $X$ be a lattice with jump sets over the tame chart and $k\ge0$.
\begin{enumerate}
\item[(i)] For a form $\mu$ (Definition~\ref{pkn:def-N}), $\mu_k(w)=s_0^{-k}\mu(s_0^kw)$ is a form; if $\mu$ maps
$1\otimes X$ into (the image of) $\pkfBm$, so does $\mu_k$.
\item[(ii)] $s_0^k$ maps $\pknMinf(X)$ and each letter part $\pkpAmb(X)_i$ into themselves.
\item[(iii)] $s_0^kE=E$.
\item[(iv)] For $b\in\pkfSB$, $z$ divides $s_0^kb-b$ in $\pkfBp$.
\end{enumerate}
\end{lemma}

\begin{proof}
(i) $\mu_k$ is $\pkpA$-linear ($s_0$ acts semilinearly on $\pkpAmb(X)$ and on $\pknB$). For $h\in\Gamma'$ put
$h'=s_0^khs_0^{-k}\in\Gamma'$ ($\Gamma'$ is normal); then $\mu_k(1\otimes hy)=s_0^{-k}\mu(1\otimes h's_0^ky)=
s_0^{-k}h'\mu(1\otimes s_0^ky)=h\,\mu_k(1\otimes y)$. $s_0$ preserves $\pkfBm$ (P1).
(ii) As in the proof of Lemma~\ref{pkl:G1}: $s_0$ preserves $\pkpA\otimes X$, $s_0\tau=\tau r'$ with $r'\in\pkpA^\times$
(Lemma~\ref{pkf:approx0}(iv)), and $\mu(\varphi^ms_0^kx)=s_0^k\mu_k(\varphi^mx)\in\pknBp$ by (i). The letter parts:
Lemma~\ref{pkw:letters}.
(iii) $E=\ell-z^e$ (Definition~\ref{pkw:def-tame}; the chart of Proposition~\ref{pkw:T1} has it), so
$s_0^kE=\ell-\xi^{ke}z^e=E$ as $\xi^e=[a^e]=1$.
(iv) By Lemma~\ref{pkl:SBdec}, $\ell^{k'}b=\lambda+zw'$ with $\lambda\in\Lambda$, $w'\in\pkfSB\subseteq\pkfBp$. $s_0$ fixes
$\lambda$ and $s_0^kz=\xi^kz$, and $s_0^k$ preserves $\pkfBp$ (P1); so $\ell^{k'}(s_0^kb-b)=z(\xi^ks_0^k(w')-w')$, and $\ell$
is a unit of $\pkfBp$.
\end{proof}

\begin{lemma}[The relation among conjugate forms]\label{pkw:rel}
Let $X$ be a lattice with jump sets over the tame chart, $Q=\sum_kQ_kX^k\in\Z_\ell[X]$ with $Q(\beta_{s_0})=0$ on
$\pkwB(X)$, and $\mu$ a form mapping $1\otimes X$ into $\pkfBm$. Then $\sum_kQ_k\mu_k(x)=0$ for every $x\in\pknMinf(X)$.
\end{lemma}

\begin{proof}
Let $l\colon X\to\pkfBm$, $l(y)=\mu(1\otimes y)$; it is $\Z_\ell$-linear and $l(hy)=hl(y)$ for $h\in\Gamma'$ (the map
$\pkfBm\to\pknB$ is injective and equivariant, (P1)). Let $\lambda\colon\pkfBm\otimes_{\Z_\ell}X\to\pkfBm$,
$b\otimes y\mapsto bl(y)$, and $\lambda_k(b\otimes y)=b\,s_0^{-k}l(s_0^ky)=s_0^{-k}\lambda(s_0^k(b\otimes y))$; these are
$\pkfBm$-linear, $\lambda$ is $\Gamma'$-equivariant, and $\lambda_k(1\otimes y)=\mu_k(1\otimes y)$. Regard
$\pkwB(X)\subseteq\pkfBm\otimes X$ (F6). Let $d\in\pkwD(X)$. Then $s_0^kd=\beta^k_{s_0}d\in\pkwD(X)$ is fixed by an open
$H'\subseteq\Gamma'$, so $\lambda(s_0^kd)\in\pkfBm^{H'}=K_{H',0}\subseteq\pkwur$ (F6), which the inertial $s_0^{-k}$ fixes;
hence $\lambda_k(d)=\lambda(\beta_{s_0}^kd)$ and $\sum_kQ_k\lambda_k(d)=\lambda(Q(\beta_{s_0})d)=0$. As $\pkwD(X)$ spans
$\pkwB(X)$ over $\pkwB$, hence $\pkfBm\otimes X$ over $\pkfBm$, the $\pkfBm$-linear map $\sum_kQ_k\lambda_k$ vanishes; in
particular $\sum_kQ_k\mu_k$ vanishes on $1\otimes X$, hence on $\pkpA\otimes X$. If $x\in\pknMinf$ then $\tau^nx\in
\pkpA\otimes X$ for some $n$, and $\tau^n\sum_kQ_k\mu_k(x)=0$ with $\tau\ne0$ in the field $\pknB$.
\end{proof}

\begin{theorem}[(Z)]\label{pkw:Z}
Statement (Z) holds.
\end{theorem}

\begin{proof}
Let $x\in\mathcal N_i(\mathrm{std}_0)$ and $w=\sum_kQ_ks_0^kx$. By Corollary~\ref{pkl:F}, $X_i(\mathrm{std}_0)=
\pknMinf(X)\cap\pkpAmb(X)_i$; so $x\in\mathcal N\cap\pknMinf$, and $w\in\pknMinf\cap\pkpAmb(X)_i$ by
Lemma~\ref{pkw:s0}(ii). Fix a frame of $X$ for its jump bound $a$ (Theorem~\ref{pkf:main}), with vectors $n_j$ and
matrix $Z$ over $\pkfSB$, and let $\sigma_\nu$ be the forms of Lemma~\ref{pkn:frame-id}(iii) for $m=0$; by their
construction $\sigma_\nu(1\otimes e_j)=(M^{-1})_{\nu j}\in\pkfBm$, so $\sigma_\nu$ maps $1\otimes X$ into $\pkfBm$.

\emph{Values at $x$.} As in the proof of Corollary~\ref{pkl:Phi}, $x=\sum_jt_jn_j$ with $t_j\in\mathfrak S$. Let $\mu$ be a
form mapping $1\otimes X$ into $\pkfBm$. By Lemma~\ref{pkn:frame-id}(ii) ($m=0$), $E^a\mu(x)=c^{(0)}(\mu)Zt$, and the
entries of $c^{(0)}(\mu)=\mu(\mathbf e)M$ lie in $\pkfBm$ and are fixed by $j(\Gamma')$ (proof of
Lemma~\ref{pkn:frame-id}(i)), hence lie in $\pkfBm^{G_{K'}}=K'_0=\Lambda[1/\ell]\subseteq\pkfSB$ (F6, and
Lemma~\ref{pkf:SB-basic}(i)). So $E^a\mu(x)\in\pkfSB$.

\emph{Values at $w$.} Fix $\nu$, put $\mu=\sigma_\nu$, $b_k=\mu_k(x)$ and $X_k=E^ab_k\in\pkfSB$ (the previous paragraph
applied to the form $\mu_k$, Lemma~\ref{pkw:s0}(i)). Then $\mu(s_0^kx)=s_0^kb_k$, and by Lemma~\ref{pkw:rel}
$\sum_kQ_kb_k=0$, so $\mu(w)=\sum_kQ_k(s_0^kb_k-b_k)$. By Lemma~\ref{pkw:s0}(iii), $E^a s_0^kb_k=s_0^kX_k$, hence
\[
E^{2a}\sigma_\nu(w)=\sum_kQ_kE^a\bigl(s_0^kX_k-X_k\bigr)\in W:=\text{image of }z\pkfBp\text{ in }\pknB
\]
by Lemma~\ref{pkw:s0}(iv); $W$ is a $\pkfBp$-submodule.

\emph{Division.} By Proposition~\ref{pkn:frame}(ii), $w=\sum_js_jn_j$ with $s_j\in\pkpA$. The proof of the coefficient test
Lemma~\ref{pkn:frame-id}(iv) uses its hypothesis only for the forms $\sigma_\nu$; applied to $E^{2a}w=\sum_jE^{2a}s_jn_j$ it
gives $E^{4a}s_j\in W$, i.e.\ (as $\pkfBp\to\pknB$ is injective, (P1)) $z$ divides $E^{4a}s_j$ in $\pkfBp$. By
Lemma~\ref{pkl:intdiv}, $z\mid E^{4a}s_j$ in $\pkpA$, and then $z\mid s_j$ (as in the proof of Lemma~\ref{pkl:cancel}:
$E\nmid z$ and $E\pkpA$ is prime). So $z^{-1}w=\sum_j(s_j/z)n_j\in\pknMinf$, and it lies in $\pkpAmb(X)_i$ since $w$ does
and $z$ commutes with $e_i$. By Corollary~\ref{pkl:F}, $z^{-1}w\in X_i(\mathrm{std}_0)$.
\end{proof}

\begin{lemma}[Counting bad steps]\label{pkw:count}
Let $Q\in\Z_\ell[X]$ be monic and $\nu<e$. Call $k$ \emph{bad} if the residue of $Q(\xi^k)\in\Lambda$ in
$\mathbb F_{\ell^f}$ is $0$. Then at most $\deg Q$ of the integers $1\le k\le\nu$ are bad, and for $k$ not bad
$Q(\xi^k)$ is a unit of $\Lambda$, hence of $\pkpA$.
\end{lemma}

\begin{proof}
The residue of $Q(\xi^k)$ is $\bar Q(a^k)$, $\bar Q\in\mathbb F_\ell[X]$ the reduction, monic of degree $\deg Q$.
The $a^k$, $1\le k\le\nu<e$, are pairwise distinct ($a$ has order $e$), and each bad $k$ gives a root $a^k$ of
$\bar Q\ne0$. An element of $\Lambda=W(\mathbb F_{\ell^f})$ with nonzero residue is a unit.
\end{proof}

\begin{lemma}[The tame step]\label{pkw:step}
Let $S$ be a side whose lattice $Y$ is potentially crystalline, $Q\in\Z_\ell[X]$ with $Q(1)=0$
killing $s_0$ on $\pkwD(X_S)$, $i$ a letter and $r,t,k\ge0$. If $\ell^tz^{k+1}\in\pkmJ r0i$, then
$Q(\xi^{k+1})\,\ell^tz^k\in\pkmJ r0i$.
\end{lemma}

\begin{proof}
Put $f_0=\ell^tz^{k+1}\in\mathfrak S\cap\pkmJ r0i$. By Lemma~\ref{pkm:J-basic}(iv) there is $u\in\mathfrak S$
such that for every $c\in C_S$ some $x_c\in\mathcal N_i(\mathrm{std}_0)$ (for $X_S$) has
$\mathrm{pr}(x_c)=(f_0-ru)\varepsilon_i(c)$. Write $u=w+zv$ with $w\in\Lambda$ and $v\in\mathfrak S$
($\mathfrak S=\Lambda[[z]]$). Fix $c$. By (Z), $y_c=z^{-1}\sum_jQ_js_0^jx_c\in(X_S)_i(\mathrm{std}_0)$.

$\mathrm{pr}$ is induced by the equivariant projection $X_S\to C_S$, so it commutes with $s_0$; and
$s_0\varepsilon_i(c)=\varepsilon_i(c)$, because $\Gamma$ acts trivially on $C_S$ and $s_0$ fixes $\Lambda$, hence
commutes with $e_i$ (Lemma~\ref{pkw:letters}). Since $s_0^jz=\xi^jz$ and $s_0$ fixes $\ell$, $\xi$, $w$,
\[
\mathrm{pr}(s_0^jx_c)=s_0^j(f_0-ru)\,\varepsilon_i(c)=a_j\,\varepsilon_i(c),\qquad
a_j=\ell^t\xi^{j(k+1)}z^{k+1}-r\bigl(w+\xi^jz\,s_0^j(v)\bigr),
\]
with $s_0^j(v)\in\pkpA$. Using $\sum_jQ_j=Q(1)=0$,
\[
\sum_jQ_ja_j=z\,b,\qquad b=Q(\xi^{k+1})\,\ell^tz^k-r\sum_jQ_j\xi^js_0^j(v)\in\pkpA .
\]
Hence $z\,\mathrm{pr}(y_c)=\mathrm{pr}(\sum_jQ_js_0^jx_c)=zb\,\varepsilon_i(c)$, and as $z$ is a unit of $\pkpW$
and $\mathrm{pr}$ is $\pkpW$-linear, $\mathrm{pr}(y_c)=b\,\varepsilon_i(c)$. As $b$ does not depend on $c$, this
says $b\in\pkmI0i$, and $Q(\xi^{k+1})\ell^tz^k=b+r\sum_jQ_j\xi^js_0^j(v)\in\pkmJ r0i$.
\end{proof}

\begin{theorem}[(T3): the tame killer]\label{pkw:T3}
Fix a tame chart with ramification index $e$. Let $S$ be a side whose jumps at all letters are
$\le n_0$, whose lattice $Y$ is potentially crystalline, and let a monic $q\in\Z_\ell[X]$ kill $s_0$ on
$\pkwD(Y)$. If $e\le\ell^{m_0}$, $m_0(\deg q+1)\le\kappa$ and $n_0+\kappa+2<\ell$, then $\ell^\kappa=\ell^\kappa E^0$
kills the cokernel of $S$ at every letter.
\end{theorem}

\begin{proof}
Fix a letter $i$ and write $J=\pkmJ{\ell^{\kappa+1}}0i$. By Lemma~\ref{pkm:zpow} (with $\omega=n_0$) there is
$\nu\ge0$ with $(\ell-1)\nu\le e(n_0+\kappa)$ and $z^\nu\in J$; as $n_0+\kappa<\ell-1$, $\nu<e$. Let
$Q=q(X)(X-1)$, monic of degree $\deg q+1$ with $Q(1)=0$; it kills $s_0$ on $\pkwD(X_S)$ by
Lemma~\ref{pkw:ext}. Call $k$ bad as in Lemma~\ref{pkw:count} and let $b(k)$ be the number of bad
$j\in[1,k]$. We show, by induction on $k$:

\emph{if $k<e$, $t\ge0$, $m_0b(k)+t\le\kappa$ and $\ell^tz^k\in J$, then $\ell^\kappa\in\pkmI0i$.}

$k=0$: $t\le\kappa$, so $\ell^\kappa\in J=\pkmJ{\ell^{\kappa+1}}0i$, and $\ell^\kappa\in\pkmI0i$ by
Lemma~\ref{pkm:J-basic}(i). Step $k\to k+1$ (with $k+1<e$, $\ell^tz^{k+1}\in J$, $m_0b(k+1)+t\le\kappa$). If $k+1$
is bad, Lemma~\ref{pkm:lower}(i) gives $\ell^{t+v_\ell(k+1)+1}z^k\in J$; as $k+1<e\le\ell^{m_0}$,
$v_\ell(k+1)+1\le m_0$, so $\ell^{t+m_0}z^k\in J$, and $m_0b(k)+(t+m_0)=m_0b(k+1)+t\le\kappa$: apply the induction
hypothesis with $t+m_0$. If $k+1$ is not bad, Lemma~\ref{pkw:step} gives $Q(\xi^{k+1})\ell^tz^k\in J$, and
$Q(\xi^{k+1})$ is a unit of $\pkpA$ (Lemma~\ref{pkw:count}), so $\ell^tz^k\in J$ with $b(k)=b(k+1)$: apply the
induction hypothesis with $t$.

Apply this with $k=\nu$, $t=0$: by Lemma~\ref{pkw:count}, $b(\nu)\le\deg Q=\deg q+1$, so
$m_0b(\nu)\le m_0(\deg q+1)\le\kappa$. Hence $\ell^\kappa\in\pkmI0i$, i.e.\ $\ell^\kappa$ kills the cokernel at $i$
(Definition~\ref{pkm:def-J}).
\end{proof}



\subsection{(I6) and (T4): Tate duals and tensor products}\label{pkw:ss-T4}

For a lattice $Y$, $Y^*=\operatorname{Hom}_\auO(Y,\auO(1))$ with $(g\phi)(y)=g\,\phi(g^{-1}y)$ and
$(a\phi)(y)=a\phi(y)=\phi(ay)$. Fix, as in Definition~\ref{pll:def}(4), a $\Z_\ell$-linear $\mathrm f\colon
\auO\to\Z_\ell$ for which $(a,b)\mapsto\mathrm f(ab)$ is a perfect pairing, and write also $\mathrm f$ for
$\mathrm f\otimes1\colon\auO(1)\to\Z_\ell(1)$, which is equivariant. Then
\[
Y\times Y^*\to\Z_\ell(1),\qquad (y,\phi)\mapsto\mathrm f(\phi(y)),
\]
is $\Z_\ell$-bilinear, $\Gamma$-equivariant, satisfies $\mathrm f((a\phi)(y))=\mathrm f(\phi(ay))$, and is perfect:
$Y$ is free over $\auO$, and for $Y=\auO$ the map $\auO\to\operatorname{Hom}_{\Z_\ell}(\auO,\Z_\ell)$, $b\mapsto
\mathrm f(b\,\cdot)$, is bijective by the choice of $\mathrm f$. Its $\pkwB$-bilinear extension
$\langle\ ,\ \rangle\colon\pkwB(Y)\times\pkwB(Y^*)\to\pkwB\otimes_{\Z_\ell}\Z_\ell(1)$ is therefore a perfect
pairing of free $\pkwB$-modules, $\Gamma$-equivariant: $\langle gx,g\phi\rangle=g\langle x,\phi\rangle$.

\begin{lemma}[(I6): inertia on duals]\label{pkw:dual}
Let $Y$ be a potentially crystalline lattice. Then $Y^*$ is potentially crystalline, and for every
inertial $s$ and $x\in\pkwB(Y)$, $\phi\in\pkwB(Y^*)$: $\langle\beta^Y_sx,\beta^{Y^*}_s\phi\rangle=\langle
x,\phi\rangle$. Consequently, for $P\in\auO[X]$: $P$ kills $s$ on $\pkwD(Y^*)$ if and only if $P$ kills $s^{-1}$ on
$\pkwD(Y)$.
\end{lemma}

\begin{proof}
Let $H\subseteq\Gamma$ be open with $D_H=\pkwB(Y)^H$ spanning, $F=K_{H,0}$. By (F1), $\pkwB\otimes_FD_H\cong
\pkwB(Y)$, so an $F$-basis $d_1,\dots,d_n$ of $D_H$ is a $\pkwB$-basis of $\pkwB(Y)$. Let $e_1$ be a basis of
$\Z_\ell(1)$ and $\epsilon=t^{-1}\otimes e_1$, a $\pkwB$-basis of $\pkwB\otimes\Z_\ell(1)$ fixed by $G_{\Q_\ell}$
(Lemma~\ref{pkw:beta}(iii)). By perfectness there are $d^1,\dots,d^n\in\pkwB(Y^*)$, a $\pkwB$-basis, with
$\langle d_i,d^k\rangle=\delta_{ik}\epsilon$. For $h\in H$, $\langle d_i,hd^k\rangle=h\langle h^{-1}d_i,d^k\rangle=
h(\delta_{ik}\epsilon)=\delta_{ik}\epsilon$, so $hd^k=d^k$: the $H$-invariants span $\pkwB(Y^*)$, and $Y^*$ is
potentially crystalline. Shrinking $H$, we may take it normal, and then $d_i\in\pkwD(Y)$, $d^k\in\pkwD(Y^*)$,
so $\langle\beta_sd_i,\beta_sd^k\rangle=\langle sd_i,sd^k\rangle=s(\delta_{ik}\epsilon)=\delta_{ik}\epsilon$; by
bilinearity $\langle\beta_sx,\beta_s\phi\rangle=\langle x,\phi\rangle$ for all $x$, $\phi$.

Hence $\langle x,\beta_s\phi\rangle=\langle\beta_s\beta_{s^{-1}}x,\beta_s\phi\rangle=\langle\beta_{s^{-1}}x,\phi\rangle$
(Lemma~\ref{pkw:beta}(i)), and $\langle x,a\phi\rangle=\langle ax,\phi\rangle$ for $a\in\auO$. So for
$P=\sum_kP_kX^k$, $\langle x,P(\beta^{Y^*}_s)\phi\rangle=\langle P(\beta^Y_{s^{-1}})x,\phi\rangle$ for all $x,\phi$.
By perfectness, $P(\beta^{Y^*}_s)=0$ if and only if $P(\beta^Y_{s^{-1}})=0$.
\end{proof}

This is input (I6) of \S\ref{pkd:ss-data} (with the meaning of ``kills'' of
Definition~\ref{pkw:def-pcris}). The jumps of $Y^*$ are treated in Proposition~\ref{pkw:HT}.

\begin{lemma}[Tensor products]\label{pkw:tensor}
Let $Y$, $L$ be potentially crystalline lattices and $s$ inertial. Then $Y\otimes_\auO L$ is potentially
crystalline, and the natural surjection $\pi\colon\pkwB(Y)\otimes_{\pkwB}\pkwB(L)\to\pkwB(Y\otimes_\auO L)$
satisfies $\pi\circ(\beta^Y_s\otimes\beta^L_s)=\beta^{Y\otimes L}_s\circ\pi$.
\end{lemma}

\begin{proof}
$\pkwB(Y)\otimes_{\pkwB}\pkwB(L)=\pkwB\otimes_{\Z_\ell}(Y\otimes_{\Z_\ell}L)$, and $\pi$ is induced by the
equivariant surjection $Y\otimes_{\Z_\ell}L\to Y\otimes_\auO L$, so it is $\pkwB$-linear, surjective and
equivariant. If $d\in\pkwD(Y)$ and $d'\in\pkwD(L)$ are fixed by open $H$, $H'$, then $\pi(d\otimes d')$ is fixed
by $H\cap H'$; the $d\otimes d'$ span the source, so the $\pi(d\otimes d')$, which lie in a single
$\pkwB(Y\otimes_\auO L)^{H''}$ for finitely many spanning ones, span the target: $Y\otimes_\auO L$ is potentially
crystalline. Both sides of the identity are $\pkwB$-linear and send $d\otimes d'$ to
$\pi(sd\otimes sd')=s\,\pi(d\otimes d')=\beta^{Y\otimes L}_s\pi(d\otimes d')$, as $\pi(d\otimes d')\in
\pkwD(Y\otimes_\auO L)$.
\end{proof}

\begin{lemma}[A Cayley--Hamilton polynomial]\label{pkw:CH}
Let $q,q'\in\Z_\ell[X]$ be monic, $M$ a module over a commutative ring containing $\Z_\ell$, and $A,A'$ commuting
endomorphisms of $M$ with $q(A)=0$, $q'(A')=0$. Then the characteristic polynomial $q''$ of multiplication by
$x\otimes x'$ on $R=\Z_\ell[x]/(q)\otimes_{\Z_\ell}\Z_\ell[x']/(q')$ is monic of degree $\deg q\cdot\deg q'$ and
satisfies $q''(AA')=0$.
\end{lemma}

\begin{proof}
$R$ is free over $\Z_\ell$ with basis $x^a\otimes x'^b$ ($a<\deg q$, $b<\deg q'$), so $q''$ is monic of degree
$\deg q\deg q'$. By Cayley--Hamilton $q''(m)=0$ for the multiplication $m$ by $\rho=x\otimes x'$; as $q''(m)$ is
multiplication by $q''(\rho)$, applying it to $1$ gives $q''(\rho)=0$ in $R$. Since $A$, $A'$ commute and
$q(A)=q'(A')=0$, there is a ring homomorphism $R\to\operatorname{End}(M)$ with $x\mapsto A$, $x'\mapsto A'$ (the
tensor product of two commutative $\Z_\ell$-algebras mapping to commuting subrings); it maps $\rho$ to $AA'$ and
$q''(\rho)$ to $q''(AA')$.
\end{proof}

\begin{theorem}[(T4)]\label{pkw:T4}
Let $Y$, $L$ be potentially crystalline lattices, $s$ inertial, and let monic $q,q'\in\Z_\ell[X]$ kill $s^{-1}$ on
$\pkwD(Y)$ and on $\pkwD(L)$ respectively. Then some monic $q''\in\Z_\ell[X]$ of degree $\deg q\cdot\deg q'$ kills
$s$ on $\pkwD(\operatorname{Hom}_\auO(Y\otimes_\auO L,\auO(1)))$.
\end{theorem}

\begin{proof}
On $M=\pkwB(Y)\otimes_{\pkwB}\pkwB(L)$ the endomorphisms $A=\beta^Y_{s^{-1}}\otimes1$ and $A'=1\otimes
\beta^L_{s^{-1}}$ commute and satisfy $q(A)=0$, $q'(A')=0$. Lemma~\ref{pkw:CH} gives $q''$ monic of degree
$\deg q\deg q'$ with $q''(AA')=0$, and by Lemma~\ref{pkw:tensor}
$q''(\beta^{Y\otimes L}_{s^{-1}})\pi=\pi\,q''(AA')=0$ with $\pi$ surjective; so $q''$ kills $s^{-1}$ on
$\pkwD(Y\otimes_\auO L)$, and by Lemma~\ref{pkw:dual} it kills $s$ on $\pkwD((Y\otimes_\auO L)^*)$.
\end{proof}

\begin{remark}[Use of (T4) in Theorem~\ref{pkc:tame}]\label{pkw:rem-T4}
Theorem~\ref{pkw:T4} is applied in the proof of Theorem~\ref{pkc:tame} with $s=s_0$, $q=q'_0$, $q'=q_0$: $q'_0$
kills $s_0^{-1}$ on $\pkwD(Y)$ (as $Y\in\pkiC(U,K)$) and $q_0$ kills $s_0^{-1}$ on $\pkwD(L)$ (as
$L\in\pkiC(U',K^{-1})$ and $K^{-1}(s_0^{-1})=K(s_0)$), so the resulting $q''$ has degree
$\deg q_0\deg q'_0\le N^2$. The hypotheses concern $s_0^{-1}$, not $s_0$: by Lemma~\ref{pkw:dual} the dual is
controlled by the polynomials killing $s_0^{-1}$.
\end{remark}

\subsection{Jumps of Tate duals; uniform tame polynomials}\label{pkw:ss-checks}

\subsubsection*{(a) The jumps of $L^*=\operatorname{Hom}_\auO(L,\auO(1))$}

We use two more published facts. Let $K'$ be a finite extension of $\Q_\ell$ and $V$ a de Rham representation of
$G_{K'}$ on a finite free module over $E_\auO=\auO[1/\ell]$.
\begin{enumerate}
\item[(F4)] (Tate \cite[\S3, Theorems~1 and~2, p.~176]{pklTate67}; also quoted in Fontaine \cite[1.5.7]{pkfFontaineII}; Fontaine \cite[3.8]{pkfFontaineIII}.) $C(a)^{G_{K'}}=0$ for $a\ne0$ and
$C^{G_{K'}}=K'$; $V$ is Hodge--Tate: for every embedding $\gamma\colon E_\auO\to C$ over $\Q_\ell$ there are integers
$m_\gamma(u)\ge0$, determined by $V$, with $C\otimes_{E_\auO,\gamma}V\cong\bigoplus_uC(-u)^{m_\gamma(u)}$
$G_{K'}$-equivariantly.
\item[(F5)] (Fontaine \cite[3.7--3.8]{pkfFontaineIII}.) $\mathrm{gr}^kD_{\mathrm{dR}}(V)\cong
(C(k)\otimes_{\Q_\ell}V)^{G_{K'}}$, compatibly with the action of $E_\auO$.
\end{enumerate}
By (F4)--(F5), $u$ is a jump of $V$ at $\gamma$ in the sense of $\mathrm{gr}^uD_{\mathrm{dR}}\ne0$ (the convention
fixed in (P3) of \S\ref{pkf:ss-rings}) exactly when $C(-u)$ occurs in $C\otimes_{E_\auO,\gamma}V$, with multiplicity
$m_\gamma(u)$; this does not depend on the open subgroup $G_{K'}$ on which it is computed.

\begin{proposition}[Jumps of Tate duals; consistency of conventions]\label{pkw:HT}
\begin{enumerate}
\item[(i)] $\auO(n)=\auO\otimes\Z_\ell(n)$ has the single jump $-n$ at every $\gamma$. In particular $\auO(1)$ has the
jump $-1$, as in Lemma~\ref{pkp:tau} and Definition~\ref{pkp:def-pairing}, and the character $\chi^e=
\chi_{\mathrm{cyc}}^{-e}$ of \S\ref{pll:sec} has the single jump $e$.
\item[(ii)] If $L$ is potentially crystalline, the jumps of $L^*$ at $\gamma$ are the $-u-1$, $u$ running over the
jumps of $L$ at $\gamma$, with the same multiplicities.
\item[(iii)] With the forms of Definition~\ref{pkn:def-N}: on $\auO(1)$ every form $\mu$ satisfies
$\mu(1\otimes\auO(1))\subseteq tK'\subseteq t^{-u}\pknBp$ with $u=-1$; so the jump $-1$ of $\auO(1)$ is also its type
in the sense of Definition~\ref{pkp:def-orders}.
\end{enumerate}
\end{proposition}

\begin{proof}
(i) $C\otimes_{E_\auO,\gamma}E_\auO(n)=C(n)=C(-(-n))$. For $\chi^e$: $C(-e)$, jump $e$.

(ii) $L^*[1/\ell]=\operatorname{Hom}_{E_\auO}(V,E_\auO(1))$ with $V=L[1/\ell]$ free over $E_\auO$; it is de Rham
(potentially crystalline by Lemma~\ref{pkw:dual}). Base change along $\gamma$ commutes with $\operatorname{Hom}$ of free
modules, so $C\otimes_{E_\auO,\gamma}L^*[1/\ell]\cong\operatorname{Hom}_C(C\otimes_{E_\auO,\gamma}V,C(1))\cong
\bigoplus_u\operatorname{Hom}_C(C(-u),C(1))^{m_\gamma(u)}=\bigoplus_uC(u+1)^{m_\gamma(u)}$, equivariantly; and
$C(u+1)=C(-(-u-1))$.

(iii) For $y=1\otimes a\otimes e_1$ and $g\in\Gamma'$, $\mu(1\otimes gy)=g\mu(y)$ reads
$\chi_{\mathrm{cyc}}(g)\mu(y)=g\mu(y)$, so $t^{-1}\mu(y)\in\pknB^{G_{K'}}=K'$ (Lemma~\ref{pkn:frame-id}(i), proof).
\end{proof}

\begin{remark}
The convention ``the jumps of $L^*$ are $-u-1$'' used in
Definition~\ref{pkp:def-pairing}, in Theorem~\ref{pkd:data} (i)--(iii) and in the deduction of the window
statements, agrees with the definition of jumps by $\mathrm{gr}^uD_{\mathrm{dR}}\ne0$ used in \S\ref{pkf:sec}
(property (P3)), with the jump $-1$ of $\auO(1)$ in Lemma~\ref{pkp:tau}, with the types of forms in
Definition~\ref{pkp:def-orders}, and with the statement of \S\ref{pll:sec} that $\chi^e=\chi_{\mathrm{cyc}}^{-e}$ has
the jump $e$. It is also the convention under which $\dim D_{\mathrm{dR}}/\mathrm{Fil}^0$ is the number of
negative jumps (as used in \S\ref{pg:sec}).
\end{remark}

\subsubsection*{(b) One tame polynomial for each inertial element}

\begin{proposition}[Uniform tame polynomials]\label{pkw:G4}
Let $E_v=\Q_\ell$, $r,N_0\ge1$, $e'=\operatorname{lcm}_{1\le a\le r}(\ell^a-1)$. For each inertial $s\in\Gamma$ let
$M_s$ be a multiset of at most $N_0$ elements of $\auQlb$, each satisfying $y^{\ell^a-1}=1$ for some $1\le a\le r$,
and let $\mathcal T$ be a set of potentially crystalline lattices such that for every $T\in\mathcal T$ and every
inertial $s$, all eigenvalues of $\beta_s^T$ lie in $M_s$. (Here an eigenvalue of $\beta_s^T$ is a
$y\in\auQlb$ such that $\beta_s-y$ has nonzero kernel on $\pkwB(T)\otimes_{\Z_\ell}\Z_\ell[y]$; since $\beta_s$ has
finite order $n$ (see the proof), each such $y$ lies in $\mu_n$, and $\pkwB(T)\otimes_{\Z_\ell}\Q_\ell(\mu_n)$ is
the direct sum of the kernels of $\beta_s-y$, $y\in\mu_n$.) Put $q_s=\prod_{y\in M_s}m_y$, $m_y$ the minimal
polynomial of $y$ over $\Q_\ell$. Then $q_s\in\Z_\ell[X]$ is monic of degree $\le N_0r$, depends only on $s$, and
$q_s$ and $X^{e'}-1$ kill $s$ on $\pkwD(T)$ for every $T\in\mathcal T$, and on $\pkwD(T')$ for every lattice $T'$
with an injective equivariant map into a finite direct sum of members of $\mathcal T$.
\end{proposition}

\begin{proof}
If $y^{\ell^a-1}=1$, $y$ lies in the unramified extension of degree $a$ and is integral, so $m_y\in\Z_\ell[X]$ is
monic of degree $\le a\le r$; hence $q_s$ is monic of degree $\le N_0r$. By Lemma~\ref{pkw:beta}(i), $\beta_{s'}=1$ for the
inertial $s'$ in some open subgroup $H$. As $\Gamma$ is profinite and $H$ is open, $s^n\in H\cap j^{-1}(\pkwI)$
for some $n\ge1$, so $\beta_s^n=\beta_{s^n}=1$ by Lemma~\ref{pkw:beta}(i). Since $X^n-1=\prod_{y\in\mu_n}(X-y)$
with pairwise distinct roots in $\Q_\ell(\mu_n)$, the idempotents $\prod_{y'\ne y}(\beta_s-y')/(y-y')$ decompose
$\pkwB(T)\otimes_{\Z_\ell}\Q_\ell(\mu_n)$ into the kernels of $\beta_s-y$, $y\in\mu_n$; the nonzero ones are the
eigenvalues, which lie in $M_s$ by hypothesis. Every $y\in M_s$ is a root of $q_s$ and of $X^{e'}-1$
($\ell^a-1\mid e'$), so $q_s(\beta_s)$ and $\beta_s^{e'}-1$ vanish on each eigenspace, hence on
$\pkwB(T)\otimes\Q_\ell(\mu_n)$ and on $\pkwB(T)$, which embeds in it. For $T'$ use
Lemma~\ref{pkw:beta}(ii) (and the evident compatibility of $\beta_s$ with finite direct sums).
\end{proof}

\begin{remark}
The definition of $\mathcal C(U,\mathrm{tame})$ before
Proposition~\ref{pll:main} needs, for Proposition~\ref{pki:tame} and Theorem~\ref{pkd:data}(iv), a set $K(s)$ of
polynomials that depends on $s$ only. (G4) (Proposition~\ref{pga:G45}, with Lemma~\ref{pga:depth0}) gives this, with the sizes read off from the
composition of $\mathcal C_v$ in the paper. By Proposition~\ref{pg:loss} and Proposition~\ref{pri:loss}(b), at
$\ell=p_i$ the members of $\mathcal C_v$ are the restrictions to $D_v$ of the lattices $u_x$ in the representations
$\rho_x$ of Proposition~\ref{pri:BL}(c) and of the $t'+1$ summands of $u_\pi=T\oplus\bigoplus_{e<t'}\auO(-e)$; and
at a chosen place $w\nmid\ell$ Proposition~\ref{pri:BL}(c) says that inertia acts on $\rho_x$ ``through the roots of the
rescaled label times $(X-1)^{t'}$''. At the chosen place $w\mid\ell$ this is replaced by (G4) for $w\mid\ell$
(Proposition~\ref{pga:G45}): $\rho_x$ is potentially crystalline at $w$, and the same data (the roots of the rescaled
label, and $t'$ trivial characters; the Hecke character $\eta$ is trivial on $\auO_{E,w}^\times$ by Lemma~\ref{pay:eta}, so $\eta_O$ is
crystalline at $w$ and contributes no inertial type) give the inertial type on $\pkwD$ (Definition~\ref{pkw:def-pcris}), whose eigenvalues on an
inertial $s$ are as follows. Lemma~\ref{pga:depth0} gives them on an $s$ with $\theta_d(s)$ a generator of
$\mathbb F_{q^d}^\times$; since the restriction to inertia is semisimple and factors through $\theta_d$, whose
image is cyclic, it is a sum of the characters $\omega\circ\theta_d^{\,e_kq^b}$, so on an arbitrary inertial $s$
the eigenvalues are the $\omega(\theta_d(s)^{e_kq^b})$ attached to the roots of the label $c_w$ (multiplied by the
element of $\mathbb F_\ell^\times$ naming the tame part of the twist), together with $t'$ eigenvalues $1$:
they depend on $s$, $c_w$ and the twist only. Sizes: the label has degree $n_W$, which is $r$ for the pairs
$O_j(\vartheta)$ and $1$ for the pairs $O_{\nu,z}(\vartheta)$ (\S\ref{pg:ss-statement}), so each irreducible
factor has degree $d_k\le n_W\le r$, and each eigenvalue lies in $\mathbb F_{\ell^{d_k}}^\times$ (the rescaling
factor lies in $\mathbb F_\ell^\times$) and is a root of unity of order dividing $\ell^{d_k}-1$; and
$|M_s|=n_W+t'=n_V$ is the rank of $\rho_x$. For the pairs $O_j(\vartheta)$, $n_V=r+t$; for these
Proposition~\ref{pkw:G4} applies with $N_0=r+t$, to the members $u_x|_{D_v}$ and to the summands of
$u_\pi|_{D_v}$ (which embed equivariantly), giving degree $\le(r+t)r$ as in the definition. For the pairs
$O_{\nu,z}(\vartheta)$ one has $n_V=2z$ with $m<z\le r$, and the bound obtained is $N_0r=2zr\le2r^2$. The
definition of $\mathcal C(U,\mathrm{tame})$ therefore uses the degree bound $\max(r+t,2r)\,r$, which covers both
kinds of pairs; as $t=5r$ it equals $(r+t)r=6r^2$, and Theorem~\ref{pkd:data}(iv) uses $N=\max(r+t,2r)r$ (the size
check there, $r^2(N+1)^2\le h$, is written for this $N$). With it, $K(s)=
\{X^{e'}-1,q_s\}$ in Theorem~\ref{pkd:data}(iv), $\mathcal C_v\subseteq\pkiC(U,K)$ holds as proved there, and
$L_c^*\in\pkiC(U',K^{-1})$ follows from Lemma~\ref{pkw:dual} (input (I6)) and Proposition~\ref{pkw:HT}(ii).
\end{remark}


%% ===== proofs/C-odd.tex =====
% proofs/C-odd.tex
% The odd classes: proof of Proposition pri:odd of proofs/C-ribet.tex.
% Labels of this file start with pod:.
\providecommand{\auO}{\mathcal O}
\providecommand{\pgcodd}{c_{\mathrm{odd}}}
\providecommand{\pgG}{G_E}
\providecommand{\pgcyc}{\varepsilon_\ell}
\providecommand{\prilen}{\operatorname{len}}
\providecommand{\prival}{\operatorname{val}}
% macros of this file (prefix \pod)
\providecommand{\anrd}{\mathrm{rd}}
\providecommand{\panotodd}[1]{\par\smallskip{\raggedright\noindent\textbf{[Not proved here:} #1\textbf{]}\par}\smallskip}
\providecommand{\podSha}{\mathrm{Sha}^-}
\providecommand{\podC}{\mathcal C}
\providecommand{\podm}{\mathfrak m}
\providecommand{\podord}{\operatorname{ord}}
\providecommand{\podfr}[1]{\langle #1\rangle}
\providecommand{\podth}{\theta}
\providecommand{\podGal}{\operatorname{Gal}}
\providecommand{\podLam}{\Lambda}

\section{The odd classes}\label{pod:sec}

This section proves Proposition~\ref{pri:odd} of \S\ref{pri:sec} (the odd classes), for every prime
$\ell$, the prime $2$ and the primes dividing $q_1d$ included. The method is that of Kolyvagin's Euler
system of Gauss sums, as presented by Rubin~\cite[\S\S1--4]{RubinTexel} (there for $\Q(\mu_p)$), made effective at finite level and
combined with a crude class-field-theoretic count that disposes of the small primes. No form of the
main conjecture of Iwasawa theory is used.

\subsection{Statement and plan}\label{pod:ss-statement}

\emph{Conventions.} $\Omega=\overline\Q\subset\C$, $G_\Q=\podGal(\Omega/\Q)$, and $c\in G_\Q$ is
the complex conjugation induced by the fixed embedding. For $M\ge1$, $\varepsilon_M\colon G_\Q\to
(\Z/M)^\times$ is the cyclotomic character modulo $M$ and $H_M=\ker\varepsilon_M=G_{\Q(\mu_M)}$. For a
prime $\ell$, $\pgcyc\colon G_\Q\to\Z_\ell^\times$ is the $\ell$-adic cyclotomic character. Places of
$\Omega$ are its valuation rings; for a place $A$ of $\Omega$, $D_A\supseteq I_A$ are its
decomposition and inertia groups in $G_\Q$ (in $G_K$ for a subfield $K$, by intersection), and
$\podord_A$ is the valuation of $A$ normalised by $\podord_A(r)=1$ for the rational prime $r$ below
$A$. A \emph{Frobenius element at $A$} is a $\phi\in D_A$ acting on the residue field of $A$ as
$t\mapsto t^r$. Throughout, $N=q_1d$, and $E\subseteq\Q(\mu_N)$ is the CM field of
\S\ref{ss:au-spaces} (it is abelian over $\Q$, contained in $\Q(\zeta_{q_1d})$, and totally complex);
$h_E$ is its class number.

$\auO$ is the ring of integers of a finite extension of $\Q_\ell$, with maximal ideal $\podm$,
uniformiser $\varpi$, valuation $\prival=\prival_\auO$ ($\prival(\varpi)=1$) and \emph{unnormalised}
length $\prilen$ (so $\prilen\auO/a=\prival(a)$). We assume that $\auO$ contains the values of all
Dirichlet characters modulo $N$. Put $\podLam_\nu=\auO/\ell^\nu\auO$. A Dirichlet character $\psi$
modulo $N$ (not necessarily primitive, $\psi(a)=0$ for $(a,N)>1$) is viewed as a character of
$G_\Q$ through $\varepsilon_N$, and for $e\in\Z$ we write
\[
\Upsilon_{e,\psi}=\pgcyc^{\,e}\,\psi\colon G_\Q\to\auO^\times .
\]
On an $\auO$-module killed by $M'$, an element $g$ with $\varepsilon_{NM'}(g)=a$ acts through
$\Upsilon_{e,\psi}$ as the scalar $a^e\psi(a)$, because $\pgcyc(g)\equiv a\pmod{\ell^{v_\ell(M')}}$.
For $x\in\Q/\Z$ let $\podfr x\in[0,1)$ be its representative, and for an integer $c\ge1$
\[
m_c(x)=\lfloor c\podfr x\rfloor-\lfloor c\podfr{-x}\rfloor\in\Z .
\]
For $M\ge1$ the units $u\in(\Z/M)^\times$ are represented in $[0,M)$, and
\[
\podth_{c,M}(e,\psi)=\sum_{u\in(\Z/M)^\times}m_c(u/M)\,u^e\psi(u)\in\auO
\]
is the \emph{twisted Stickelberger sum}.

A \emph{cocycle for a character} $\chi\colon G\to\auO^\times$ with values in an $\auO$-module $V$
is a map $x\colon G\to V$ with $x(gh)=x(g)+\chi(g)x(h)$; all cocycles below are locally constant
(continuous for the discrete topology on $V$). The module of Proposition~\ref{pri:odd} is, for
$\nu\ge0$ and $0\le e$:
\begin{definition}\label{pod:def-sha}
For $e\ne0$, $\podSha_{e,\nu}$ is the $\auO$-module of locally constant cocycles $z\colon\pgG\to
\podLam_\nu$ for $\pgcyc^{-e}|_{\pgG}$ (that is, with values in $\podLam_\nu(-e)$) such that (i) $z$
vanishes on $I_A\cap\pgG$ for every place $A$ of $\Omega$ not above $\ell$; (ii) $z$ vanishes on
the elements $g\in D_A\cap\pgG$ with $\pgcyc(g)=1$, for every place $A$ above $\ell$; (iii) $z(cgc)=
-z(g)$ for all $g\in\pgG$. For $e=0$, $\podSha_{0,\nu}$ is the module of continuous homomorphisms
$\pgG\to\podLam_\nu$ that vanish on every inertia group.
\end{definition}
This is the module of Proposition~\ref{pri:odd}: ``unramified outside $\ell$'' is (i), and $g^c=cgc$.
No comparison with a Tate--Shafarevich group is needed or made.

\begin{proposition}[$=$ Proposition~\ref{pri:odd}]\label{pod:odd}
Let $t\ge0$ and $y=N+t+3$. There is an integer $0<\pgcodd\le2^{4y^{27}}h_E$ such that for every
prime $\ell$, every $\auO$ as above, all $\nu\ge0$ and all $0\le e<t$,
\[
\prilen_\auO\podSha_{e,\nu}\le\prival_\auO(\pgcodd).
\]
\end{proposition}

\emph{Plan.} For $e=0$ the class number serves (\S\ref{pod:ss-e0}). For $e\ne0$ the module is split,
over $\Q$, into components indexed by the characters $\psi$ modulo $N$ (\S\ref{pod:ss-split}); only
the components with $\Upsilon_{e,\psi}$ odd matter. An integer $X_e$ of size $\le2^{(N+e+3)^6}$ with
a ``Stickelberger property'' is produced from Bernoulli numbers and Voronoi's congruence
(\S\ref{pod:ss-voronoi}). With Gauss and Jacobi sums (\S\ref{pod:ss-gauss}) and Artin reciprocity,
$4X_e^2$ kills every odd component, at \emph{every} prime $\ell$ (\S\ref{pod:ss-annih}); so
$Y_e=16\varphi(N)X_e^2$ kills $\podSha_{e,\nu}$. To pass from an annihilator to a length one needs
a bounded number of elements of $\pgG$ at which a cocycle that vanishes is zero (``detecting
points''). For $\ell\le y^2$ a count of generators of a ray class group gives $\le y^{20}$ such points
(\S\ref{pod:ss-detect}); for $\ell>y^2$ Kolyvagin's derivative classes give $v_\ell(X_e)+1$ points for
each odd component (\S\S\ref{pod:ss-koly}--\ref{pod:ss-descent}). The two regimes are put together
in \S\ref{pod:ss-proof}.

\begin{remark}[The prime $2$ and the primes dividing $N$]\label{pod:rem-two}
For $\ell\le y^2$, in particular for $\ell=2$ and for every $\ell\mid N$, only two ingredients are
used: the annihilator $Y_e$ (Proposition~\ref{pod:annih} and Lemma~\ref{pod:split}), whose proof
makes no assumption on $\ell$ (no division by $2$, by $\varphi(N\ell)$ or by $\ell-1$ occurs; the
factor $4$ in $4X_e^2$ and in $4\varphi(N)$ is exactly what makes this possible), and the detecting
points of Lemma~\ref{pod:detect}, which come from class field theory and are valid for all $\ell$.
The Euler-system descent, which needs $2$, $\varphi(N\ell)$ and some $a^e-1$, $a^{e+1}-1$ to be units
of $\auO$, is only used for $\ell>y^2\ge N+e+3$. The price for small $\ell$ is the exponent $y^{20}$
in place of $1$, which is absorbed in the bound for $\pgcodd$. We do not know a published statement
covering $\ell=2$ in this form.
\end{remark}

\subsection{Generalities on cocycles}\label{pod:ss-gen}

\begin{lemma}\label{pod:coc}
Let $G$ be a profinite group, $\chi\colon G\to\auO^\times$ a continuous character, $V$ an
$\auO$-module and $x\colon G\to V$ a locally constant cocycle for $\chi$.
\begin{enumerate}
\item[(a)] The zeros of $x$ form an open subgroup $Z(x)$ of $G$; in particular $x$ vanishes on the
closed subgroup generated by any set on which it vanishes.
\item[(b)] If $g,s\in G$ and $\chi(g)$ fixes $x(s^{-1})$ \textup(e.g.\ $\chi(g)$ acts trivially on
$V$\textup), then $x(sgs^{-1})=\chi(s)x(g)$.
\item[(c)] \textup(Sah\textup) If $x$ vanishes at the commutator $k=\delta g\delta^{-1}g^{-1}$ and
$\chi(k)=1$, then $(\chi(\delta)-1)x(g)=(\chi(g)-1)x(\delta)$.
\item[(d)] Let $\gamma\in G$ with $x(\gamma)=0$ and $(\chi(\gamma)-1)x\ne0$. Then $x$ does not vanish
identically on any closed normal subgroup $H$ with $G/H$ abelian.
\item[(e)] If $g\in G$ satisfies $g^2=1$ and $\chi(g)$ acts trivially on $V$, then $2x(g)=0$.
\end{enumerate}
\end{lemma}
\begin{proof}
(a) $x(1)=0$; if $x(g)=x(h)=0$ then $x(gh)=0$ and $0=x(gg^{-1})=x(g)+\chi(g)x(g^{-1})$; the set is
open as $x$ is locally constant. (b) $x(sgs^{-1})=x(s)+\chi(s)x(g)+\chi(sg)x(s^{-1})$ and $x(s^{-1})
=-\chi(s)^{-1}x(s)$; as $\chi(g)x(s^{-1})=x(s^{-1})$ the first and last terms cancel. (c) Expand
$x(\delta g)=x(kg\delta)$: $x(\delta)+\chi(\delta)x(g)=x(k)+\chi(k)\bigl(x(g)+\chi(g)x(\delta)\bigr)$.
(d) If $x|_H=0$, then for $g\in G$ the element $h=g^{-1}\gamma^{-1}g\gamma$ lies in $H$ and
$g\gamma=\gamma gh$, so $x(g)=x(g)+\chi(g)x(\gamma)=x(g\gamma)=x(\gamma g)+\chi(\gamma g)x(h)=
\chi(\gamma)x(g)$; thus $(\chi(\gamma)-1)x=0$. (e) $x(g)=x(g^3)=x(g)+\chi(g)x(g)+\chi(g)^2x(g)=3x(g)$.
\end{proof}

\begin{lemma}\label{pod:length}
Let $M$ be an $\auO$-submodule of the maps $G\to\podLam_\nu$, killed by an integer $m\ge1$, and let
$p_1,\dots,p_k\in G$ be such that $x\in M$ and $x(p_i)=0$ for all $i$ imply $x=0$. Then
$\prilen M\le\prival(m^k)$.
\end{lemma}
\begin{proof}
$x\mapsto(x(p_i))_i$ embeds $M$ into $\podLam_\nu[m]^k$, and $\podLam_\nu[m]\cong\auO/(\ell^\nu,m)$
has length $\le\prival(m)$.
\end{proof}

\subsection{Detecting points from class field theory}\label{pod:ss-detect}

Let $K$ be a totally complex number field, $S$ a finite set of finite places of $K$ and $M\ge1$.
Call a map $x\colon G_K\to V$ to an abelian group an \emph{$(S,M)$-homomorphism} if it is a locally
constant homomorphism, $Mx=0$, and $x$ vanishes on the inertia groups at all finite places of $K$
outside $S$.

\begin{lemma}\label{pod:cft}
Let $U_S^{(M)}$ be the group of finite ideles of $K$ that are $M$-th powers at the places of $S$.
Then $[\mathbb A_{K,f}^\times:U_S^{(M)}]\le M^{2\#S+[K:\Q]}$, and there are finitely many $p_i\in
G_K$ and integers $n_i\ge2$ with $\prod_in_i$ dividing $h_K[\mathbb A_{K,f}^\times:U_S^{(M)}]$,
such that every $(S,M)$-homomorphism $x$ satisfies $n_ix(p_i)=0$ for all $i$, and $x=0$ if
$x(p_i)=0$ for all $i$.
\end{lemma}
\begin{proof}
The index is a product over $v\in S$ of $[K_v^\times:K_v^{\times M}]=M\,\#\mu_M(K_v)\,|M|_v^{-1}$ ($|\cdot|_v$ normalised, so that
$|M|_v^{-1}=\#\auO_v/M\auO_v$). (Proof: for an endomorphism $\phi$ of an abelian group $A$ with finite kernel
and cokernel put $h(\phi)=\#\operatorname{coker}\phi/\#\ker\phi$; $h$ is multiplicative in short exact
sequences preserved by $\phi$ and equals $1$ on finite groups. Apply this to $\phi=(\cdot)^M$ on
$K_v^\times=\varpi^\Z\times\auO_v^\times$: on $\varpi^\Z$, $h=M$; let $p$ be the residue characteristic, $e$
the absolute ramification index and $k>e/(p-1)$; the $p$-adic logarithm is an isomorphism of
$U=1+\varpi^k\auO_v$ onto the additive group $\varpi^k\auO_v\cong\auO_v$ (inverse: the exponential), so
$h((\cdot)^M|_U)=\#\auO_v/M\auO_v=|M|_v^{-1}$; as $\auO_v^\times/U$ is finite,
$h((\cdot)^M|_{\auO_v^\times})=|M|_v^{-1}$; and $\ker\phi=\mu_M(K_v)$, so
$\#\operatorname{coker}\phi=M\,\#\mu_M(K_v)\,|M|_v^{-1}$. See also \cite[Chapter~II, Corollary~(5.8)]{Neukirch}.) Further $\prod_{v\in S}|M|_v^{-1}\le M^{[K:\Q]}$ and
$\#\mu_M(K_v)\le M$. Let $\mathcal A$ be the subgroup of the finite ideles generated by $K^\times$
and by the ideles that are units at every finite place outside $S$ and lie in $U_S^{(M)}$. Since $K$
is totally complex, the Artin map $\mathrm{Art}\colon\mathbb A_{K,f}^\times/\overline{K^\times}\to
G_K^{\mathrm{ab}}$ is surjective (global class field theory: the archimedean components have trivial Artin image, $K$
being totally complex; for every finite abelian extension $L/K$ the norm residue symbol $\mathbb A_K^\times\to
\mathrm{Gal}(L/K)$ is surjective \cite[Chapter~VI, Theorem~(5.5)]{Neukirch}, so the image is dense; and it
is compact, the source being a quotient of the compact group $C_K^0$ of idele classes of absolute value $1$
\cite[Chapter~VI, Theorem~(1.6)]{Neukirch}). For an $(S,M)$-homomorphism $x$ and any
homomorphism $f\colon V\to\Q/\Z$, $f\circ x$ is a character of $G_K$ that kills
$\mathrm{Art}(\mathcal A)$: it kills $\mathrm{Art}(K^\times)$ by reciprocity, the local Artin images
of units outside $S$ (inertia), and the images of $M$-th powers. Hence $f\circ x$ factors through
the finite abelian group $Q=G_K^{\mathrm{ab}}/\mathrm{Art}(\mathcal A)$, and since the $f$ separate
points of $V$, so does $x$. Write $Q=\prod_i\Z/n_i$ with $n_i\ge2$ and let $p_i\in G_K$ lift the
generators. The order of $Q$ divides $[\mathbb A^\times_{K,f}:\mathcal A]$, which divides
$[\mathcal AU_\emptyset:\mathcal A]$ times $[\mathbb A^\times_{K,f}:\mathcal AU_\emptyset]$, where
$U_\emptyset$ are the ideles that are units at all finite places outside $S$ (no condition at $S$);
the second index divides $h_K$ (the quotient by $K^\times U_\emptyset$ is a quotient of the class
group), and the first divides $[U_\emptyset:U_\emptyset\cap U_S^{(M)}]$, which divides the index of
$U_S^{(M)}$.
\end{proof}

\begin{lemma}\label{pod:hK}
For $M\ge3$ and $K=\Q(\mu_M)$, $h_K\cdot[\mathbb A_{K,f}^\times:U_S^{(M)}]\le2^{M^3+5M^2+M}$ when
$S$ is the set of places above one rational prime.
\end{lemma}
\begin{proof}
$d=[K:\Q]=\varphi(M)\le M$, $|d_K|\le M^d\le(2^M)^M$, and $\#S\le d$. Every ideal class contains an
integral ideal of norm $\le|d_K|^{1/2}\le|d_K|$ (Minkowski), and there are at most $n^d$ integral
ideals of norm $n$ (at most $d$ prime ideals above each prime and $\le$ the number of factorisations,
crudely $\le n^d$), so $h_K\le|d_K|^{d+1}\le2^{M^2(M+1)}$. The index is $\le M^{3M}\le2^{3M^2}$ by
Lemma~\ref{pod:cft}. Hence the product is $\le2^{M^3+4M^2}$.
\end{proof}
Any bound $2^{O(M^3)}$ suffices below.

\begin{lemma}[Detecting points]\label{pod:detect}
Let $\ell$ be a prime and $y\ge N$ with $\ell\le y^2$. There is a set $s\subseteq\pgG$ with $\#s\le
y^{20}$ such that for all $\auO$, $\nu$ and $e$, every locally constant cocycle $z\colon\pgG\to
\podLam_\nu$ for $\pgcyc^{-e}$ that vanishes on $I_A\cap\pgG$ for all places $A\nmid\ell$ and
vanishes on $s$ is zero. Moreover, if $g_0\in\pgG$ topologically generates $\pgG/(\pgG\cap\ker
\pgcyc)$, there is such a set $s'$ with $\#s'\le y^{20}$, contained in $\ker\pgcyc$, such that
$z(g_0)=0$ and $z|_{s'}=0$ imply $z=0$.
\end{lemma}
\begin{proof}
Put $M=N\ell^2\ge3$ and $K=\Q(\mu_M)\supseteq E$, so $G_K=H_M\subseteq\pgG$; let $S$ be the places of
$K$ above $\ell$. Let $s$ consist of a set of representatives $\{\rho_t\}$ of $\pgG/G_K$ (at most
$\varphi(M)$ elements) and the points $p_i$ of Lemma~\ref{pod:cft} for $(K,S,M)$; by
Lemma~\ref{pod:hK} there are at most $\log_2\bigl(h_K[\mathbb A^\times_{K,f}:U_S^{(M)}]\bigr)\le
M^3+5M^2+M$ of the latter. Since $M\le N\ell^2\le y^5$, $\#s\le M+M^3+5M^2+M\le M^4\le y^{20}$.

Let $z$ vanish on $s$, and suppose $\ell^kz=0$; we show $z=0$ by induction on $k$, the case $k=0$
being trivial. The cocycle $Mz$ vanishes on $s$ and is killed by $\ell^{k-1}$, since
$\ell^{k-1}Mz=N\ell\cdot\ell^kz=0$; so $Mz=0$ by induction. For $h\in G_K$ one has
$\pgcyc(h)\equiv1\pmod{\ell^{v_\ell(M)}}$, and $z(g)$ is killed by $M$, so $\pgcyc(h)^{-e}z(g)=z(g)$. Thus $z(hg)=z(h)+z(g)$ for $h\in G_K$, $g\in
\pgG$; in particular $z|_{G_K}$ is an $(S,M)$-homomorphism (an inertia group of $K$ at a place not
above $\ell$ lies in one of $E$ at a place not above $\ell$). It vanishes at the $p_i$, so
$z|_{G_K}=0$ by Lemma~\ref{pod:cft}. For $g\in\pgG$ write $g=h\rho_t$ with $h\in G_K$; then
$z(g)=z(h)+z(\rho_t)=0$.

For the second statement, write each $p\in s$ as $p=h_pq_p$ with $h_p\in\ker\pgcyc$ and $q_p$ in the
closed subgroup $\overline{\langle g_0\rangle}$ (possible as $g_0$ generates modulo the kernel and
$\pgcyc$ is continuous with closed image), and let $s'=\{h_p\}$. If $z(g_0)=0$ and $z|_{s'}=0$, then
$z$ vanishes on $\overline{\langle g_0\rangle}$ by Lemma~\ref{pod:coc}(a), hence at $p=h_pq_p$, and
$z=0$ by the first part.
\end{proof}

\begin{corollary}[Proposition~\ref{pri:kummer}(b) for $\ell\le(N+t+2)^2$]\label{pod:kummer-b}
Let $h=r^8$ as in \S\ref{pg:ss-ribet}, $t=5r$, and let $\ell\le y_0^2$ with $y_0=N+t+2$. With $g_0$ a
topological generator of $\pgG$ modulo $\ker\pgcyc$, there is $s\subseteq\ker\pgcyc$ with
$(t+1)(\#s+1)\le y_0^{21}\le h^{21}$ such that a locally constant cocycle $\pgG\to\podLam_\nu(-e)$
unramified outside $\ell$ that vanishes at $g_0$ and on $s$ is zero, for all $\auO,\nu,e$.
\end{corollary}
\begin{proof}
The second statement of Lemma~\ref{pod:detect} with $y=y_0\ge N$ gives $s\subseteq\ker\pgcyc$ with $\#s\le y_0^{20}$. As $t+2\le y_0$ and
$y_0\le y_0^{20}$: $(t+1)(\#s+1)\le(y_0-1)y_0^{20}+(t+1)\le y_0^{21}-y_0^{20}+y_0\le y_0^{21}$.
Finally $y_0\le h$: by Remark~\ref{pod:rem-size}, $N\le x^{15}$, and $x^5\le r$, so $y_0\le x^{15}+5r+2
\le r^3+5r+2\le r^8$.
\end{proof}

\begin{remark}[The range of Proposition~\ref{pri:kummer}(b)]\label{pod:rem-kummer-b}
For $\ell>y_0^2$ the inequality $(t+1)(\#s+1)\le h^{21}$ is not obtained from
Corollary~\ref{pod:kummer-b}: the detecting set grows with $\ell$ (Lemma~\ref{pod:detect} gives
$\#s\le\max(N,\ell^{1/2})^{20}$). Applying Lemma~\ref{pod:detect} with $y=h$ instead of $y_0$ gives
the inequality for all $\ell\le h^2$ (Proposition~\ref{pku:b}), and Proposition~\ref{pri:triangle}
uses it only when $h_0\ne0$, which forces $\ell\le h^2$ (Proposition~\ref{pku:triangle}).
\end{remark}

\begin{remark}[Proposition~\ref{pri:kummer}(a)]\label{pod:rem-kummer-a}
Proposition~\ref{pri:kummer}(a) is proved in Proposition~\ref{pku:main}: with $T$ equal to the set of
primes of $F=E^+$ above $\ell$ or below the given places, the relations are the restrictions of dual
Selmer classes, and the Greenberg--Wiles formula, applied to $\Z/\ell^\nu((\pgcyc^+)^{-e})$ and to its
$\ell$-torsion, gives at least $\#x+[F:\Q]-4$ independent reductions.
\end{remark}

\subsection{The case \texorpdfstring{$e=0$}{e=0}}\label{pod:ss-e0}

\begin{lemma}\label{pod:e0}
$\prilen\podSha_{0,\nu}\le\prival(h_E)$.
\end{lemma}
\begin{proof}
An element of $\podSha_{0,\nu}$ is a continuous homomorphism $\pgG\to\podLam_\nu$ vanishing on all
inertia groups; as $E$ is totally complex it factors through $\podGal(H/E)\cong\mathrm{Cl}_E$, $H$
the Hilbert class field. Writing $\mathrm{Cl}_E\cong\prod_i\Z/n_i$, $\podSha_{0,\nu}\cong
\prod_i\podLam_\nu[n_i]$, of length $\le\sum_i\prival(n_i)=\prival(h_E)$.
\end{proof}

\subsection{Splitting by characters}\label{pod:ss-split}

\begin{definition}\label{pod:def-comp}
For $e\ge1$, a Dirichlet character $\psi$ modulo $N$ and $\nu\ge0$, the \emph{component}
$\podC_{e,\psi,\nu}$ is the $\auO$-module of locally constant cocycles $x\colon G_\Q\to\podLam_\nu$
for $\Upsilon_{e,\psi}^{-1}$ that vanish at $c$ and at every $g\in H_N$ which lies either in $I_A$ for
a place $A\nmid\ell$, or in $D_A$ for a place $A\mid\ell$ with $\pgcyc(g)=1$. We call $\psi$ (or
$\Upsilon_{e,\psi}$) \emph{odd} if $\Upsilon_{e,\psi}(c)=(-1)^e\psi(-1)=-1$.
\end{definition}

\begin{lemma}\label{pod:split}
Let $e\ge1$. There are $\sigma_t\in G_\Q$ \textup($t\in(\Z/N)^\times$\textup) and $b(u,g)\in\pgG$
\textup($u,g\in G_\Q$\textup), independent of $\auO,\nu$, with the following property. For $x\in
\podSha_{e,\nu}$ and each $\psi$ put
\[
z_\psi(g)=\sum_{t\in(\Z/N)^\times}\Upsilon_{e,\psi}(\sigma_t)\,x\bigl(b(\sigma_t,g)\bigr)\qquad(g\in G_\Q).
\]
Then (a) $z_\psi\in\podC_{e,\psi,\nu}$; (b) if $z_\psi=0$ for every odd $\psi$, then $4\varphi(N)x=0$.
\end{lemma}
\begin{proof}
Choose $S_t\in G_\Q$ with $\varepsilon_N(S_t)=t$ and put $\sigma_t=S_tS_1^{-1}$, so $\sigma_1=1$. For
$u,g\in G_\Q$ the element $u(cgc^{-1}g^{-1})u^{-1}$ lies in $H_N\subseteq\pgG$ (as $E\subseteq
\Q(\mu_N)$ and $G_\Q/H_N$ is abelian); let $b(u,g)$ be this element of $\pgG$. Since $b(u,g)$ is a conjugate of a commutator,
$\pgcyc(b(u,g))=1$, and therefore
\begin{equation}\label{pod:eq-add}
x(b(u,g)\,k)=x(b(u,g))+x(k)\qquad(k\in\pgG).
\end{equation}
\emph{Cocycle rule.} One has $b(u,gk)=b(u,g)\,b(ug,k)$, so by \eqref{pod:eq-add}
$x(b(u,gk))=x(b(u,g))+x(b(ug,k))$. Next, $T(u,k):=\Upsilon_{e,\psi}(u)x(b(u,k))$ depends only on
$uH_N$: if $a\in H_N$ then $b(au,k)=ab(u,k)a^{-1}$, and with the cocycle rule for $x$ and
$\pgcyc(b(u,k))=1$, $x(ab(u,k)a^{-1})=\pgcyc(a)^{-e}x(b(u,k))$, while $\Upsilon_{e,\psi}(au)=
\pgcyc(a)^e\Upsilon_{e,\psi}(u)$. Hence
\[
z_\psi(gk)=z_\psi(g)+\sum_t\Upsilon_{e,\psi}(\sigma_t)x(b(\sigma_tg,k))
=z_\psi(g)+\Upsilon_{e,\psi}(g)^{-1}\sum_tT(\sigma_tg,k)=z_\psi(g)+\Upsilon_{e,\psi}(g)^{-1}z_\psi(k),
\]
because $\sigma_tg\in\sigma_{t\varepsilon_N(g)}H_N$ and $t\mapsto t\varepsilon_N(g)$ permutes
$(\Z/N)^\times$. Local constancy is clear.

\emph{Local conditions.} $b(u,c)=1$, so $z_\psi(c)=0$. Let $g\in H_N$ lie in $I_A$ for some $A\nmid
\ell$, or in $D_A\cap\ker\pgcyc$ for some $A\mid\ell$. For $u\in G_\Q$ let $w(u)\in\pgG$ be the
element $ugu^{-1}\in H_N$; it lies in $I_{uA}$ (resp.\ in $D_{uA}\cap\ker\pgcyc$), so $x(w(u))=0$ by
(i), (ii) of Definition~\ref{pod:def-sha}. As $b(u,g)w(u)=w(uc)$, \eqref{pod:eq-add} gives
$x(b(u,g))=x(w(uc))-x(w(u))=0$, so $z_\psi(g)=0$. This proves (a).

(b) If $\psi$ is even, $\Upsilon_{e,\psi}(c)=1$ and $z_\psi(c)=0$ give $z_\psi(cgc)=z_\psi(g)$. If all
odd $z_\psi$ vanish, then $\Sigma(h):=\sum_\psi z_\psi(h)$ satisfies $\Sigma(chc)=\Sigma(h)$ for
$h\in\pgG$. By orthogonality ($\auO$ contains the values of the characters),
$\sum_\psi\Upsilon_{e,\psi}(\sigma_t)=\pgcyc(\sigma_t)^e\sum_\psi\psi(t)=\varphi(N)[t=1]$, so
$\Sigma(h)=\varphi(N)x(b(1,h))=\varphi(N)\bigl(x(chc)-x(h)\bigr)$, by \eqref{pod:eq-add} applied to
$b(1,h)h=chc$. Comparing $h$ and $chc$ and using $x(chc)=-x(h)$:
$-2\varphi(N)x(h)=2\varphi(N)x(h)$, i.e.\ $4\varphi(N)x(h)=0$.
\end{proof}

\subsection{A Stickelberger bound from Voronoi's congruence}\label{pod:ss-voronoi}

$B_n(T)$ are the Bernoulli polynomials and $B_n=B_n(0)$; recall $B_n(T+1)-B_n(T)=nT^{n-1}$,
$B_n'(T)=nB_{n-1}(T)$, and that $(j+1)!\,B_j\in\Z$ (induction on $j$ via $\sum_{k\le j}\binom{j+1}{k}
B_k=0$ for $j\ge1$).

\begin{definition}\label{pod:def-stick}
Let $N,e\ge1$. An integer $X\ge1$ is a \emph{Stickelberger bound for $(N,e)$} if
\begin{enumerate}
\item[(S1)] for every prime $\ell$ some integer $a\ge2$ with $a\equiv1\pmod N$ prime to $\ell$ has $(a^e-1)(a^{e+1}-1)\mid
X$;
\item[(S2)] for every $a\equiv1\pmod N$, every domain $R$ of characteristic $0$ and every Dirichlet
character $\psi$ modulo $N$ with values in $R$ and $\psi(-1)=-(-1)^e$, there is $w\in R$ such that
for all $m\ge1$ prime to $a$
\[
(a^{e+1}-1)X\equiv2w\sum_{0\le j<Nm}\Bigl\lfloor\frac{aj}{Nm}\Bigr\rfloor j^e\psi(j)\pmod{mR}.
\]
\end{enumerate}
\end{definition}

\begin{lemma}\label{pod:powersum}
For $a\in\Z_{\ge0}$ put $\beta_a=(e+2)!\,N^{e+1}B_{e+1}(a/N)$. Then $\beta_a\in\Z$ and, for every
$u\ge1$, $(e+2)!\sum_{i<u}(a+iN)^{e+1}\equiv\beta_au\pmod{u^2}$.
\end{lemma}
\begin{proof}
Let $P(T)=(e+1)!\,N^{e+2}\bigl(B_{e+2}(T/N)-B_{e+2}\bigr)$. The coefficient of $T^{k+1}$ is
$(e+1)!\binom{e+2}{k+1}N^jB_j$ with $j+k=e+1$, an integer since $(j+1)!$ divides $(e+1)!\binom{e+2}
{k+1}$ (for $k=0$ this is $(e+2)!$, otherwise $j+1\le e+1$). From $B_{e+2}(T+1)-B_{e+2}(T)=(e+2)T^{e+1}$
one gets $P(T+N)-P(T)=N(e+2)!\,T^{e+1}$, hence $P(a+uN)-P(a)=N(e+2)!\sum_{i<u}(a+iN)^{e+1}$. Taylor
expansion of the integer polynomial $P$ gives $P(a+uN)\equiv P(a)+P'(a)uN\pmod{(uN)^2}$, and
$P'(a)=(e+2)!N^{e+1}B_{e+1}(a/N)=\beta_a$ is an integer. Divide by $N$.
\end{proof}

\begin{lemma}[Voronoi's congruence]\label{pod:voronoi}
Let $R$ be a commutative ring without additive torsion, $\chi$ a character modulo $N$ with values in
$R$, $m\ge1$, $f=Nm$, and $a\equiv1\pmod N$ prime to $m$. Then
\[
(e+2)!(e+1)N a^e\sum_{0\le j<f}\Bigl\lfloor\frac{aj}{f}\Bigr\rfloor j^e\chi(j)\equiv
(a^{e+1}-1)\sum_{0\le j<N}\chi(j)\beta_j\pmod{mR}.
\]
\end{lemma}
\begin{proof}
As $(a,f)=1$, $j\mapsto aj-fk_j$ with $k_j=\lfloor aj/f\rfloor$ permutes $[0,f)$, and $\chi(aj-fk_j)=
\chi(j)$ since $N\mid f$ and $a\equiv1\pmod N$. Expanding $(aj-fk_j)^{e+1}$ to first order,
\[
\sum_{j<f}\chi(j)j^{e+1}\equiv a^{e+1}\sum_{j<f}\chi(j)j^{e+1}-(e+1)a^ef\sum_{j<f}k_j\chi(j)j^e
\pmod{f^2}.
\]
Write $j=b+iN$ ($b<N$, $i<m$) and use Lemma~\ref{pod:powersum}: $(e+2)!\sum_{j<f}\chi(j)j^{e+1}\equiv
m\sum_{b<N}\chi(b)\beta_b\pmod{m^2}$. Multiplying the display by $(e+2)!$ and substituting,
$(a^{e+1}-1)m\sum_b\chi(b)\beta_b\equiv(e+2)!(e+1)a^eNm\sum_jk_j\chi(j)j^e\pmod{m^2}$. Divide by $m$.
\end{proof}

\begin{lemma}\label{pod:L}
Let $\chi$ be a complex character modulo $N$ \textup(possibly imprimitive\textup) and $e\ge1$. Then
$\sum_{0\le j<N}\chi(j)\beta_j=-(e+2)!(e+1)N\,L(\chi,-e)$, and $L(\chi,-e)=0$ if and only if
$\chi(-1)=(-1)^e$.
\end{lemma}
\begin{proof}
$L(\chi,s)=N^{-s}\sum_{j=1}^N\chi(j)\zeta(s,j/N)$ and $\zeta(-e,x)=-B_{e+1}(x)/(e+1)$; the term $j=N$
equals the term $j=0$ as $B_{e+1}(1)=B_{e+1}(0)$ for $e\ge1$. This gives the formula. For the zeros:
write $\chi=\chi_0\chi_f$ with $\chi_f$ primitive of conductor $f$; then $L(\chi,s)=L(\chi_f,s)
\prod_{p\mid N,p\nmid f}(1-\chi_f(p)p^{-s})$, and the Euler factors do not vanish at $s=-e$ (their
zeros lie on $\Rel s=0$). For primitive $\chi_f$, the functional equation relates $L(\chi_f,-e)$ to
$L(\overline{\chi_f},1+e)\ne0$ times a ratio of Gamma factors, which vanishes exactly when $-e$ is a
pole of $\Gamma\bigl((s+\kappa)/2\bigr)$, $\kappa\in\{0,1\}$ the parity of $\chi_f$, i.e.\ when
$e\equiv\kappa\pmod2$ (see the functional equation in \cite[Chapter~4, p.~29]{Washington}). For $\chi_f=1$ ($\zeta$) the same holds since $e\ge1$.
\end{proof}

\begin{lemma}\label{pod:X}
For $N,e\ge1$ there is a Stickelberger bound $X_e$ for $(N,e)$ with $1\le X_e\le2^{(N+e+3)^6}$.
\end{lemma}
\begin{proof}
Put $g(t)=2\beta_t+[t\equiv1](-1)^e+[t\equiv-1]$ for $t\in\Z/N$ (with $\beta_t=\beta_{t_0}$,
$t_0\in[0,N)$ the representative) and let $T$ be the integer matrix indexed by
$(\Z/N)^\times$ with $T_{s,t}=g(t/s)$. For every complex character $\chi$ modulo $N$, the vector
$(\chi(t))_t$ is an eigenvector with eigenvalue $\lambda_\chi=\sum_{t\in\Z/N}\chi(t)g(t)=
2\sum_t\chi(t)\beta_t+(-1)^e+\chi(-1)$. If $\chi(-1)=-(-1)^e$, $\lambda_\chi=2\sum\chi\beta\ne0$ by
Lemma~\ref{pod:L}; otherwise $\sum\chi\beta=0$ and $\lambda_\chi=2(-1)^e\ne0$. The characters form a
basis, so $\det T=\prod_\chi\lambda_\chi\ne0$. Put
\[
X_e=\bigl((1+N)^e-1\bigr)\bigl((1+N)^{e+1}-1\bigr)\bigl((1+2N)^e-1\bigr)\bigl((1+2N)^{e+1}-1\bigr)\,
|\det T|.
\]
(S1): $\ell$ cannot divide both $1+N$ and $1+2N$; take $a$ the one prime to $\ell$.
(S2): for $\psi$ with values in $R$ (we may enlarge $R$ to contain $\Z[\psi]$, a domain of
characteristic $0$), $(\psi(t))_t$ is an eigenvector of $T$ with eigenvalue $\lambda_\psi$, so applying
the adjugate, $\lambda_\psi\mid\det T$ in $R$; write $X_e=\lambda_\psi u$. By parity $\lambda_\psi=
2\sum_t\psi(t)\beta_t$. Lemma~\ref{pod:voronoi} multiplied by $2u$ gives (S2) with $w=u(e+2)!(e+1)Na^e$.
\emph{Size.} $|B_{e+1}(x)|\le2(e+1)!\zeta(e+1)(2\pi)^{-e-1}<17(e+1)!$ on $[0,1]$ (Fourier series), so
$|\beta_t|<17(e+2)!N^{e+1}(e+1)!$ and the entries of $T$ are $<64\,(e+2)!N^{e+1}(e+1)!=:64M_e$. Hence
$|\det T|\le\varphi(N)!\,(64M_e)^{\varphi(N)}$ (Leibniz expansion).
Now let $u=N+e+3\ge5$; then $\log_2N\le u$, $\log_2(e+2)\le u$, $\varphi(N)\le N\le u$. Each of the
four factors in front of $|\det T|$ is $<(1+2N)^{e+1}$, and $\log_2(1+2N)\le2N$, so their product is
$<2^{8N(e+1)}\le2^{8u^2}$. Next, $(e+2)!(e+1)!\le(e+2)^{2(e+2)}$, so $\log_2(64M_e)\le6+2(e+2)
\log_2(e+2)+(e+1)\log_2N\le6+3u^2$, and $\log_2\varphi(N)!\le N\log_2N\le u^2$. Hence
\[
\log_2X_e\le8u^2+u^2+u(6+3u^2)=3u^3+9u^2+6u\le7u^3\le u^6 ,
\]
using $u\ge5$ (so $9u^2+6u\le4u^3$ and $7\le u^3$). That is, $X_e\le2^{(N+e+3)^6}$.
\end{proof}

\subsection{Gauss sums and a family of Jacobi-type quotients}\label{pod:ss-gauss}

\begin{lemma}[Hermite's identity]\label{pod:hermite}
For $c,r\ge1$ coprime and $x\in\Q/\Z$, $\sum_{b\in\Z/r}m_c(x+b/r)=m_c(rx)$.
\end{lemma}
\begin{proof}
$\lfloor c\podfr z\rfloor=c\podfr z-\podfr{cz}$. The representatives of $x+b/r$ are
$(b'+\podfr{rx})/r$, $b'\in[0,r)$, so $\sum_b\podfr{x+b/r}=\frac{r-1}2+\podfr{rx}$; as $c$ is prime to
$r$, also $\sum_b\podfr{cx+cb/r}=\frac{r-1}2+\podfr{crx}$. Hence $\sum_b\lfloor c\podfr{x+b/r}\rfloor=
\frac{(c-1)(r-1)}2+\lfloor c\podfr{rx}\rfloor$; the same with $-x$, and subtract.
\end{proof}

Fix a prime $p$ and a place $B$ of $\Omega$ above $p$. Let $\psi_0$ be a nontrivial additive
character of $\mathbb F_p$ and $\chi_0$ a generator of the group of characters of $\mathbb F_p^\times$,
both with values in $\Omega$; for $t\in\Z/(p-1)$ let $G(t)=\sum_{x\in\mathbb F_p^\times}\chi_0^t(x)
\psi_0(x)$ be the Gauss sum. Recall: $G(0)=-1$; $G(t)G(-t)=\chi_0^t(-1)p$ for $t\ne0$; hence $G(t)\mid
p$ and $G(t)$ is a unit at every place not above $p$. For $\sigma\in G_\Q$ with $\varepsilon_p(\sigma)
=\alpha$ and $\varepsilon_{p-1}(\sigma)=a$ one has $\sigma G(t)=\chi_0^{at}(\alpha)^{-1}G(at)$
(substitute $x\mapsto\alpha^{-1}x$); write $K_\sigma(s)=\chi_0^s(\alpha)^{-1}$, a homomorphism in $s$.

\begin{lemma}[Weak Stickelberger]\label{pod:stick}
There is an injective homomorphism $S\colon\Z/(p-1)\to\Q/\Z$ with $\podfr{S(t)}=\podord_B G(t)$ for
all $t$.
\end{lemma}
\begin{proof}
If $\sigma\in I_B$ then $\varepsilon_{p-1}(\sigma)=1$, so $\sigma G(t)=K_\sigma(t)G(t)$ and
$\sigma$ fixes $G(a+b)/(G(a)G(b))$. An element fixed by $I_B$ lies in a field unramified at $B$ over
$\Q_p$ (after completion), so its $\podord_B$ is an integer; hence $t\mapsto\podord_BG(t)\bmod\Z$ is a
homomorphism $S$. For $t\ne0$: $\psi_0$ takes $p$-th roots of unity, $\equiv1$ modulo $B$, so
$G(t)\equiv\sum_x\chi_0^t(x)=0$ modulo $B$ and $\podord_BG(t)>0$; the same for $p/G(t)=\pm
\sum\chi_0^{-t}(x)\psi_0^{-1}(x)$, so $\podord_BG(t)<1$. With $\podord_BG(0)=0$ this gives
$\podfr{S(t)}=\podord_BG(t)$, and $S(t)=0$ forces $t=0$.
\end{proof}

\begin{lemma}[Jacobi family]\label{pod:J}
Let $c\ge1$. There is $J\colon\Q/\Z\to\Omega^\times$ such that for every $M$ with $p\equiv1\pmod M$
and all $x$ with $Mx=0$:
\begin{enumerate}
\item[(J1)] $\sigma J(x)=J(ax)$ whenever $\varepsilon_M(\sigma)=a$;
\item[(J2)] $J(x)$ is a unit at every place not above $p$;
\item[(J3)] $\podord_BJ(x)=m_c(x)$;
\item[(J4)] $|\iota J(x)|=1$ for every embedding $\iota\colon\Omega\to\C$;
\item[(J5)] if $r\mid M$ and $(r,c)=1$, then $\prod_{b\in\Z/r}J(x+b/r)\big/J(rx)$ is a root of unity;
\item[(J6)] if $q$ is a prime dividing $p-1$, $qb=0$, and $A$ is a place above $q$, then
$J(x+b)/J(x)\equiv1$ modulo $A$ \textup(and both are units at $A$\textup).
\end{enumerate}
\end{lemma}
\begin{proof}
Put $j(t)=G(t)^c/G(ct)$ and $J(S(t))=j(t)/j(-t)$ (and $J=1$ outside the $(p-1)$-torsion, which is
$S(\Z/(p-1))$). (J1): with $a=\varepsilon_{p-1}(\sigma)$, $\sigma j(t)=K_\sigma(at)^cG(at)^c/
(K_\sigma(cat)G(cat))=j(at)$, and $S(at)=aS(t)$; for $Mx=0$, $ax$ depends only on $a\bmod M$.
(J2): Gauss sums divide $p$. (J3): $\podord_Bj(t)=c\podfr{S(t)}-\podfr{cS(t)}=\lfloor c\podfr{S(t)}
\rfloor$, and $S(-t)=-S(t)$. (J4): if $\sigma\in G_\Q$ induces complex conjugation through $\iota$,
then $\varepsilon_{p-1}(\sigma)=-1$ and $\sigma j(t)=j(-t)$, so $\iota J=\iota j/\overline{\iota j}$.
(J5): Let $Q$ be the quotient. It is a unit off $p$ by (J2) and has absolute value $1$ everywhere by
(J4). A place $A$ above $p$ is $\sigma^{-1}B$ for some $\sigma$; with $a=\varepsilon_M(\sigma)$ (a unit
modulo $r$), (J1) and (J3) give $\podord_A Q=\podord_B\sigma Q=\sum_bm_c(ax+ab/r)-m_c(rax)=0$ by
Lemma~\ref{pod:hermite}. So $Q$ is an algebraic integer all of whose conjugates have absolute value
$1$; by Kronecker's theorem it is a root of unity. (J6): for $qu=0$ in $\Z/(p-1)$, $\chi_0^u$ takes
$q$-th roots of unity, $\equiv1$ modulo $A$, so $G(w+u)\equiv G(w)$ modulo $A$, and both are units at
$A$ ($q\ne p$). Hence $j(w+u)\equiv j(w)$ and $J(x+b)\equiv J(x)$ modulo $A$, using $b=S(u)$ with
$qu=0$.
\end{proof}

\subsection{Artin reciprocity for components and the annihilator}\label{pod:ss-annih}

\begin{lemma}\label{pod:locals}
Let $x\in\podC_{e,\psi,\nu}$ with $mx=0$ \textup($m\ge1$\textup), and let $a\equiv1\pmod N$ be
prime to $\ell$.
\begin{enumerate}
\item[(a)] If $\phi\in H_{Nm}$ and $\varepsilon_{Nm}(s)=u$, then $x(s^{-1}\phi s)=u^e\psi(u)x(\phi)$.
\item[(b)] If $g\in H_{Nm}\cap D_A$ for a place $A\mid\ell$, then $(a^e-1)x(g)=0$.
\item[(c)] If $g\in H_{Nm}$ and $g^2=1$, then $2x(g)=0$.
\item[(d)] If $g\in H_{Nm}\cap I_A$ for a place $A\nmid\ell$, then $x(g)=0$.
\end{enumerate}
\end{lemma}
\begin{proof}
Elements of $H_{Nm}$ act trivially through $\Upsilon_{e,\psi}^{\pm1}$ on $\podLam_\nu[m]$.
(a) is Lemma~\ref{pod:coc}(b) with $\chi=\Upsilon^{-1}_{e,\psi}$, $\chi(s^{-1})=\Upsilon_{e,\psi}(s)$.
(c) is Lemma~\ref{pod:coc}(e). (d) is the defining condition. (b): as $\Q_\ell(\mu_{\ell^\infty})/\Q_\ell$
is totally ramified with group $\Z_\ell^\times$, there is $\delta\in I_A$ with $\pgcyc(\delta)=a$;
$\delta$ fixes $\mu_N$ (the prime-to-$\ell$ part because $\delta$ is in inertia above $\ell$, the
$\ell$-part because $a\equiv1\pmod N$), so $\Upsilon_{e,\psi}(\delta)=a^e$. The commutator
$k=\delta g\delta^{-1}g^{-1}$ lies in $H_N\cap D_A\cap\ker\pgcyc$, so $x(k)=0$, and
Lemma~\ref{pod:coc}(c) gives $(a^{-e}-1)x(g)=(\Upsilon_{e,\psi}(g)^{-1}-1)x(\delta)=0$.
\end{proof}

\begin{lemma}[Artin reciprocity for components]\label{pod:artin}
Let $m\ge1$ be an integer, $F=\Q(\mu_{Nm})$, $q$ a prime
with $q\equiv1\pmod{Nm}$, $B$ a place above $q$, $\phi$ a Frobenius element at $B$. Let $x\in
\podC_{e,\psi,\nu}$ with $mx=0$, and $a\equiv1\pmod N$ prime to $\ell$. Let $y\in F^\times$ be such
that at every place $A$ of $\Omega$ above neither $\ell$ nor $q$, either $x$ vanishes on $D_A$, or
$y=\eta\,r^m$ with $\eta$ a unit at $A$ and $r\in\Omega^\times$ fixed by $I_A\cap H_{Nm}$. For $u\in
(\Z/Nm)^\times$ choose $s_u\in G_\Q$ with $\varepsilon_{Nm}(s_u)=u$ and let $k_u=\podord_B(s_uy)$.
Then
\[
\Bigl(\sum_{u\in(\Z/Nm)^\times}k_u\,u^e\psi(u)\Bigr)\cdot2(a^e-1)\,x(\phi)=0 .
\]
\end{lemma}
\begin{proof}
Put $d=2(a^e-1)$ and let $f\colon\podLam_\nu\to\Q/\Z$ be any homomorphism. Since $H_{Nm}=G_F$ acts
trivially on $\podLam_\nu[m]$, $\chi_f(h)=f(d\,x(h))$ is a continuous character of $G_F$ with $m\chi_f
=0$. By Lemma~\ref{pod:locals}, $\chi_f$ is unramified at all places not above $\ell$, kills the
decomposition groups above $\ell$ (by (b)), and kills every element of order $2$ (by (c)), hence the
archimedean local Artin images. Global reciprocity for $F$ gives $\sum_w\chi_f(\mathrm{Art}_w(y))=0$
over all places $w$ of $F$ ($\chi_f$ factors through a finite abelian extension of $F$, and the
product formula for the norm residue symbol of a principal idele is \cite[Chapter~VI, Corollary~(5.7)]{Neukirch}). The terms at places above $\ell$ vanish; at a place
$w\nmid\ell q$ of $F$ below $A$, $\chi_f$ is unramified, so $\chi_f(\mathrm{Art}_w(\alpha))=
\podord_w(\alpha)\chi_f(\mathrm{Frob}_w)$ for $\alpha\in F_w^\times$; either $\chi_f$ kills $D_A\cap
G_F$, or $y=\eta r^m$ as above, and then $\podord_w(\eta)=0$ while $\podord_w(r^m)=m\,\podord_w(r)$
with $\podord_w(r)\in\Z$ because $F_w(r)/F_w$ is unramified ($r$ is fixed by inertia); as $m\chi_f=0$,
the term vanishes. The prime $q$ splits
completely in $F$; the places of $F$ above $q$ are the $w_u$ below $s_u^{-1}B$, pairwise distinct
(if $w_u=w_{u'}$ then $s_{u'}gs_u^{-1}\in D_B\subseteq H_{Nm}$ for some $g\in G_F$, forcing
$u=u'$) and exhausting (any place above $q$ is $\sigma^{-1}B$, and $\sigma^{-1}B$ lies over the same
place of $F$ as $s_u^{-1}B$ for $u=\varepsilon_{Nm}(\sigma)$). At $w_u$, $\chi_f$ is unramified, a
Frobenius is $s_u^{-1}\phi s_u$, and $\podord_{w_u}(y)=\podord_{s_u^{-1}B}(y)=\podord_B(s_uy)=k_u$
($q$ is unramified in $F$). Hence $\sum_uk_u\chi_f(s_u^{-1}\phi s_u)=0$, i.e.\ by
Lemma~\ref{pod:locals}(a), $f\bigl(\sum_uk_uu^e\psi(u)\,d\,x(\phi)\bigr)=0$. As the $f$ separate points,
the claim follows.
\end{proof}

\begin{lemma}[Euler factor]\label{pod:thetadiv}
Let $X$ be a Stickelberger bound for $(N,e)$, $\psi$ odd \textup($\Upsilon_{e,\psi}(c)=-1$, i.e.\
$\psi(-1)=-(-1)^e$\textup), $a\equiv1\pmod N$ prime to $\ell$, $n\ge0$. Then in $\podLam_n$, the
Stickelberger sum $\podth_{a,N\ell^n}(e,\psi)$ divides $(a^{e+1}-1)X$.
\end{lemma}
\begin{proof}
For $n=0$ the ring is zero. Let $n\ge1$ and write $n=n'+1$. Put $T_m=\sum_{0\le j<Nm}\lfloor aj/(Nm)
\rfloor j^e\psi(j)$ and $U=\sum_{u\in(\Z/N\ell^n)^\times}\lfloor au/(N\ell^n)\rfloor u^e\psi(u)$. An
integer $j<N\ell^n$ with $\psi(j)\ne0$ is either prime to $\ell$ or $\ell j'$ with $j'<N\ell^{n'}$, and
$\lfloor a\ell j'/(N\ell^n)\rfloor=\lfloor aj'/(N\ell^{n'})\rfloor$; so $T_{\ell^n}=U+\ell^e\psi(\ell)
T_{\ell^{n'}}$. Modulo $\ell^n$ (indeed modulo $N\ell^n$), $\podth_{a,N\ell^n}=2U$: $m_a(u/M)=
\lfloor au/M\rfloor-\lfloor a(M-u)/M\rfloor$ and $u\mapsto M-u$ multiplies $u^e\psi(u)$ by
$(-1)^e\psi(-1)=-1$ modulo $M$. By (S2) at $m=\ell^n$ and $m=\ell^{n'}$ (both prime to $a$),
$(a^{e+1}-1)X\equiv2wT_{\ell^n}\pmod{\ell^n}$ and $\equiv2wT_{\ell^{n'}}\pmod{\ell^{n'}}$; multiplying
the second by $\ell^e\psi(\ell)$ ($e\ge1$) makes it a congruence modulo $\ell^n$. Subtracting,
$(1-\ell^e\psi(\ell))(a^{e+1}-1)X\equiv2wU=w\podth_{a,N\ell^n}\pmod{\ell^n}$, and $1-\ell^e\psi(\ell)$
is a unit in $\podLam_n$ ($\ell^e$ is nilpotent there).
\end{proof}

\begin{proposition}[Annihilation, every $\ell$]\label{pod:annih}
Let $e\ge1$, $X$ a Stickelberger bound for $(N,e)$ and $\psi$ odd. Then $4X^2$ kills
$\podC_{e,\psi,\nu}$, for every prime $\ell$ and every $\nu$.
\end{proposition}
\begin{proof}
Let $z\in\podC=\podC_{e,\psi,\nu}$ and suppose $4X^2z\ne0$. The cocycle $z'=2X^2z\in\podC$ vanishes
at $c$, and $(\Upsilon_{e,\psi}^{-1}(c)-1)z'=-4X^2z\ne0$; by Lemma~\ref{pod:coc}(d), $z'$ does not
vanish on $H_{N\ell^\nu}$. So there is $\sigma_0\in H_{N\ell^\nu}$ with $z'(\sigma_0)\ne0$. Let $U$ be the intersection of
$H_{N\ell^\nu}$ with the normal core of the open subgroup $Z(z')$; it is an open normal subgroup of
$G_\Q$ contained in $H_{N\ell^\nu}$. For $k\in U$, $z'(\sigma_0k)=z'(\sigma_0)+\sigma_0\cdot z'(k)
=z'(\sigma_0)$, so $z'\ne0$ on the coset $\sigma_0U\subseteq H_{N\ell^\nu}$. By Chebotarev's density
theorem, applied to the finite Galois extension cut out by $U$, the coset $\sigma_0U$ contains a
Frobenius element at a place above a prime unramified in that extension; since the coset lies in
$H_{N\ell^\nu}$, that prime splits in $\Q(\mu_{N\ell^\nu})$. Hence there are a prime $p\equiv1\pmod{N\ell^\nu}$, a place $B$ above it and a Frobenius element $\sigma\in
H_{N\ell^\nu}$ at $B$ with $z'(\sigma)\ne0$. By (S1) choose $a\equiv1\pmod N$ prime to $\ell$ with
$X=(a^e-1)(a^{e+1}-1)b$, $b\in\Z$. Let $J$ be the Jacobi family of Lemma~\ref{pod:J} for $(p,B,c=a)$
and $y=J(1/(N\ell^\nu))$. By (J1) with $M=N\ell^\nu$, $y$ is fixed by $H_{N\ell^\nu}$, so $y\in F=
\Q(\mu_{N\ell^\nu})$, and $s_uy=J(u/(N\ell^\nu))$; by (J2) $y$ is a unit at every place not above $p$;
by (J3) $\podord_B(s_uy)=m_a(u/(N\ell^\nu))$. Lemma~\ref{pod:artin} (with $m=\ell^\nu$, $q=p$,
$\eta=y$, $r=1$) gives $\podth_{a,N\ell^\nu}(e,\psi)\cdot2(a^e-1)z(\sigma)=0$ in $\podLam_\nu$. By
Lemma~\ref{pod:thetadiv} write $(a^{e+1}-1)X=\podth\kappa$ in $\podLam_\nu$; then
$2(a^e-1)(a^{e+1}-1)X\,z(\sigma)=0$, and multiplying by $b$, $2X^2z(\sigma)=0$, a contradiction.
\end{proof}

\begin{corollary}\label{pod:kill}
For $e\ge1$ and $X_e$ as in Lemma~\ref{pod:X}, $Y_e=4\varphi(N)\cdot4X_e^2$ kills $\podSha_{e,\nu}$
for every prime $\ell$.
\end{corollary}
\begin{proof}
For $x\in\podSha_{e,\nu}$, the odd components of $4X_e^2x$ are $4X_e^2z_\psi=0$
(Lemma~\ref{pod:split}(a) and Proposition~\ref{pod:annih}); apply Lemma~\ref{pod:split}(b) to
$4X_e^2x$.
\end{proof}

\subsection{Kolyvagin's derivative classes}\label{pod:ss-koly}

In this subsection $\ell$ is a prime, $n\ge1$, $e\ge1$, $\psi$ is a character modulo $N$ with values
in $\auO$, and $\Upsilon_1=\Upsilon_{e+1,\psi}=\pgcyc^{e+1}\psi$. We fix a primitive $\ell^n$-th root
of unity $\zeta\in\Omega$; for $\omega\in\mu_{\ell^n}$, $\log\omega\in\Z/\ell^n$ is defined by
$\zeta^{\log\omega}=\omega$. For a finite set $T$ of primes let $\mathcal I(T)$ be the set of
$g\in H_N$ lying in $I_A$ for some place $A$ above neither $\ell$ nor a prime of $T$.

\begin{definition}\label{pod:def-KS}
Let $\gamma\in G_\Q$, $p$ a prime and $P$ a finite set of primes. A \emph{Kolyvagin system} for
$(\Upsilon_1,\gamma,p,P)$ with values in an $\auO$-module $V$ is a family of maps $\kappa_S\colon
G_\Q\to V$, $S\subseteq P$, such that for all $S\subseteq P$: (KS1) $\kappa_S$ is a locally constant
cocycle for $\Upsilon_1$ vanishing on $\mathcal I(S\cup\{p\})$; (KS2) $\kappa_S(\gamma)=0$; (KS3) for
$q\in P\smallsetminus S$, a place $A$ above $q$ and a Frobenius element $\phi$ at $A$, there is
$\tau\in I_A$ with $\kappa_{S\cup\{q\}}(\tau)=\kappa_S(\phi)$.
\end{definition}

\begin{proposition}\label{pod:KS}
Let $c\ge1$, $\gamma\in G_\Q$, $P$ a finite set of primes $\equiv1\pmod{cN\ell^n}$, and $p$ a prime
with $p\equiv1\pmod{N\ell^n\prod_{q\in P}q}$, $B$ a place above $p$. There is a Kolyvagin system
$(\kappa_S)$ for $(\Upsilon_1,\gamma,p,P)$ with values in $\podLam_n$ such that: if $\kappa_\emptyset$
vanishes on $I_B$, then $(\Upsilon_1(\gamma)-1)\,\podth_{c,N\ell^n}(e,\psi)=0$ in $\podLam_n$.
\end{proposition}

\emph{Construction.} Let $J$ be the Jacobi family of Lemma~\ref{pod:J} for $(p,B,c)$; it satisfies
(J1)--(J6) for every $M\mid p-1$, in particular for $M_S=N\ell^n\prod_{q\in S}q$ and $qM_S$ ($S
\subseteq P$, $q\in P\smallsetminus S$). Choose $y(x)\in\Omega^\times$ with $y(x)^{\ell^n}=J(x)$. For
$q\in P$ choose a surjective homomorphism $L_q\colon(\Z/q)^\times\to\Z/\ell^n$ (possible as
$\ell^n\mid q-1$) and $v_q$ with $L_q(v_q)=1$. For $x$ with $M_Sx=0$, $\sigma\in G_\Q$ and
$a=\varepsilon_{M_S}(\sigma)$, (J1) gives $(\sigma y(x)/y(ax))^{\ell^n}=\sigma J(x)/J(ax)=1$, i.e.
\begin{equation}\label{pod:eq-root}
\sigma\,y(x)/y(ax)\in\mu_{\ell^n}.
\end{equation}
For $g\in G_\Q$ the commutator $\gamma^{-1}g\gamma g^{-1}$ lies in $H_{M}$ for every $M$, hence fixes
$J(x)$, hence moves $y(x)$ by an element of $\mu_{\ell^n}$; applying $\gamma$, $\gamma y(x)/(g\gamma
g^{-1})y(x)\in\mu_{\ell^n}$. Put
\[
F_g(x)=\log\frac{\gamma\,y(x)}{(g\gamma g^{-1})\,y(x)}\in\Z/\ell^n .
\]
For a function $G\colon\Q/\Z\to\Z/\ell^n$, shifts $d=(d_q)\in(\Z/\ell^n)^P$ and $S\subseteq P$, let
$\mathcal D_q^{d}G(x)=\sum_{u\in(\Z/q)^\times}(L_q(u)-d_q)\,G(x+u/q)$ (these operators commute) and
\[
\Lambda_S^d(G)=\sum_{a\in(\Z/N\ell^n)^\times}a^e\psi(a)\,\Bigl(\prod_{q\in S}\mathcal D_q^d\Bigr)G
\Bigl(\frac a{N\ell^n}\Bigr)\in\podLam_n,
\qquad \kappa_S(g)=\Lambda_S^0(F_g).
\]
$\Lambda_S^d(G)$ depends only on $G$ on the $M_S$-torsion and on the $d_q$, $q\in S$.

\begin{lemma}[Equivariance]\label{pod:equiv}
Let $M$ be a multiple of $M_S$, $g\in G_\Q$, and $a\in\Z$ prime to $M$ with $a\equiv
\varepsilon_{N\ell^n}(g)\pmod{N\ell^n}$. If $G(x)=G'(ax)$ for all $x$ with $Mx=0$, then $\Upsilon_{e,\psi}(g)\Lambda_S^d(G)=\Lambda_S^{f}(G')$
with $f_q=L_q(a\bmod q)+d_q$.
\end{lemma}
\begin{proof}
Induction on $\#S$; only the two hypotheses on $a$ are used. For $S=\emptyset$, as $a\equiv
\varepsilon_{N\ell^n}(g)\pmod{N\ell^n}$, $\Upsilon_{e,\psi}(g)$ acts on $\podLam_n$ as $a^e\psi(a)$ and the
substitution $a'\mapsto aa'$ in the sum over $(\Z/N\ell^n)^\times$ gives the claim. For $S\cup\{q\}$,
the substitution $u\mapsto(a\bmod q)u$ (a bijection of $(\Z/q)^\times$, as $q\mid M$ and $a$ is prime to $M$) in $\mathcal D_q$ turns $\mathcal D_q^dG$ into $\mathcal
D_q^fG'$ composed with $x\mapsto ax$, using $L_q(au)-f_q=L_q(u)-d_q$ and $a(x+u/q)=ax+au/q$.
\end{proof}

\begin{lemma}[Shifts do not matter]\label{pod:shift}
For $S\subseteq P$, $h\in G_\Q$ and all shifts $d$, $\Lambda_S^d(F_h)=\Lambda_S^0(F_h)$.
\end{lemma}
\begin{proof}
\emph{Distribution relation.} For $r\in P$ and $M_{S\cup\{r\}}x=0$: $\sum_{b\in\Z/r}F_h(x+b/r)=
F_h(rx)$. Indeed $D=\prod_by(x+b/r)/y(rx)$ has $D^{\ell^n}$ a root of unity by (J5) ($r\equiv1\pmod c$
is prime to $c$), so $D$ is a root of unity, and $\gamma$ and $h\gamma h^{-1}$ act on it in the same
way (through the same cyclotomic character); the claim is $\log(\gamma D/(h\gamma h^{-1})D)=0$.
Since $r\equiv1\pmod{N\ell^n}$, Lemma~\ref{pod:equiv} (with $g=1$, $a=r$, $M=M_{S'}$, which is prime to $r$, and $r\equiv1=\varepsilon_{N\ell^n}(1)\pmod{N\ell^n}$, $G=\sum_b F_h(\cdot+b/r)$,
$G'=F_h$) gives, for $r\notin S'$: $\Lambda_{S'}^d\bigl(\sum_bF_h(\cdot+b/r)\bigr)=\Lambda_{S'}^{f}(F_h)$
for some shifts $f$.

\emph{Induction on $\#S$.} It suffices to change $d$ at one $r\in S$. Put $S'=S\smallsetminus\{r\}$.
Then $\Lambda_S^d(F_h)=\Lambda_{S'}^d(\mathcal D_r^dF_h)$, and changing $d_r$ to $d_r'$ changes this by
$(d_r-d_r')\Lambda_{S'}^d\bigl(\sum_{u\in(\Z/r)^\times}F_h(\cdot+u/r)\bigr)=(d_r-d_r')\bigl(
\Lambda_{S'}^d(\textstyle\sum_{b\in\Z/r}F_h(\cdot+b/r))-\Lambda_{S'}^d(F_h)\bigr)$. By the
distribution relation the first term is $\Lambda_{S'}^f(F_h)$, which equals $\Lambda_{S'}^d(F_h)$ by
induction. So the change is $0$.
\end{proof}

\begin{proof}[Proof of Proposition~\ref{pod:KS}]
(KS1), cocycle rule. Let $g,g'\in G_\Q$, $M=M_S$, $a=\varepsilon_M(g)$. For $Mx=0$,
\[
F_{gg'}(ax)-F_g(ax)=a\,F_{g'}(x).
\]
Indeed, as $\ell^n\mid M$, $g$ raises $\ell^n$-th roots of unity to the power $a$, so
$\zeta^{aF_{g'}(x)}=g\bigl(\gamma y(x)/(g'\gamma g'^{-1})y(x)\bigr)=(g\gamma g^{-1})(gy(x))\big/
(gg'\gamma(gg')^{-1})(gy(x))$. By \eqref{pod:eq-root}, $gy(x)=y(ax)\omega$ with $\omega\in\mu_{\ell^n}$,
on which all conjugates of $\gamma$ act alike, so $\omega$ cancels and the right side is
$\frac{(g\gamma g^{-1})y(ax)}{\gamma y(ax)}\cdot\frac{\gamma y(ax)}{(gg'\gamma(gg')^{-1})y(ax)}=
\zeta^{F_{gg'}(ax)-F_g(ax)}$. Hence, by Lemma~\ref{pod:equiv} with $G=aF_{g'}$, $G'=F_{gg'}-F_g$, and
Lemma~\ref{pod:shift},
\[
\kappa_S(gg')-\kappa_S(g)=\Lambda_S^0(F_{gg'}-F_g)=\Upsilon_{e,\psi}(g)\,\Lambda_S^0(aF_{g'})=
\Upsilon_{e,\psi}(g)\,a\,\kappa_S(g')=\Upsilon_1(g)\kappa_S(g'),
\]
as $\pgcyc(g)\equiv a\pmod{\ell^n}$. $\kappa_S$ is locally constant since $g\mapsto(g\gamma g^{-1})y(x)$
is, for the finitely many $x$ involved.

(KS1), vanishing. Let $g\in H_N\cap I_A$, $A$ above $r'$ with $r'\notin\{\ell,p\}\cup S$. Then $g$
fixes $\mu_{M_S}$ (it fixes $\mu_N$, and the roots of unity of order prime to $r'$ by inertia).
For $M_Sx=0$, $y(x)$ is a unit at $A$ (by (J2)), so $gy(x)\equiv y(x)$ modulo $A$, and $gy(x)/y(x)
\in\mu_{\ell^n}$ by \eqref{pod:eq-root}; as $\ell^n$-th roots of unity are distinct modulo $A$
($r'\ne\ell$), $gy(x)=y(x)$; the same holds for $g^{-1}$ and for $y(ax)$. With $\gamma y(x)=
y(ax)\omega$ ($a=\varepsilon_{M_S}(\gamma)$, $\omega\in\mu_{\ell^n}$ fixed by $g$), $(g\gamma
g^{-1})y(x)=g\gamma y(x)=y(ax)\omega=\gamma y(x)$. So $F_g=0$ on the $M_S$-torsion and $\kappa_S(g)=0$.

(KS2): $F_\gamma=0$.

(KS3). Let $q\in P\smallsetminus S$, $M=M_S$ (prime to $q$), $A$ above $q$, $\phi$ a Frobenius at $A$.
Choose $\tau\in I_A$ with $\varepsilon_{qM}(\tau)\equiv1\pmod M$ and $\equiv v_q^{-1}\pmod q$ (inertia
at $A$ maps onto $(\Z/q)^\times$ and fixes $\mu_M$). We claim, for $Mx=0$,
\begin{equation}\label{pod:eq-FI}
\sum_{u\in(\Z/q)^\times}L_q(u)\,F_\tau(x+u/q)=F_\phi(qx).
\end{equation}
Granting this, $\kappa_{S\cup\{q\}}(\tau)=\Lambda_S^0(\mathcal D_q^0F_\tau)=\Lambda^0_S(x\mapsto
F_\phi(qx))=\Lambda_S^f(F_\phi)=\kappa_S(\phi)$ by Lemma~\ref{pod:equiv} ($g=1$, $M=M_S$, $a=q$, prime to $M_S$ and $\equiv1\pmod
{N\ell^n}$) and Lemma~\ref{pod:shift}.

Proof of \eqref{pod:eq-FI}. Put $Y_b=y(x+b/q)$ ($b\in\Z/q$), $y'=y(qx)$, and let $\lambda(u)\in\Z$ lift
$L_q(u)$, so $\lambda(v_qu)-\lambda(u)-1=\ell^nk_u$ with $k_u\in\Z$. By (J1) at level $qM$:
$\tau Y_{v_qu}/Y_u\in\mu_{\ell^n}$ and, for every conjugate $h\tau h^{-1}$, the same. By (J1) at level
$M$: $\phi Y_0/y'\in\mu_{\ell^n}$, and the same for $h\phi h^{-1}$. By (J5), $\bigl(\prod_bY_b/y'
\bigr)^{\ell^n}$ is a root of unity. For $h\in G_\Q$ put
\[
W_h=\frac{\prod_bY_b}{(h\phi h^{-1})Y_0}\prod_{u\in(\Z/q)^\times}\Bigl(\frac{(h\tau h^{-1})Y_{v_qu}}
{Y_u}\Bigr)^{\lambda(u)} .
\]
Each $W_h$ is a root of unity (product of roots of unity and of $\prod_bY_b/y'$). \emph{Claim:}
$W_h\equiv1$ modulo $hA$. Since $h\tau h^{-1}\in I_{hA}$ acts trivially on the residue field and
$h\phi h^{-1}$ as $t\mapsto t^q$, and the $Y_b$ are units at $hA$ ($q\ne p$; (J2)), modulo $hA$:
$W_h\equiv\prod_bY_b\cdot Y_0^{-q}\prod_u(Y_{v_qu}/Y_u)^{\lambda(u)}=\prod_bY_b\cdot Y_0^{-q}
\prod_uY_u^{\lambda(v_q^{-1}u)-\lambda(u)}$, and $\lambda(v_q^{-1}u)-\lambda(u)=-1-\ell^nk_{v_q^{-1}u}$.
By (J6) (with $q\mid p-1$), $Y_u^{\ell^n}=J(x+u/q)\equiv J(x)=Y_0^{\ell^n}$ modulo $hA$, and since
$\sum_u(\lambda(v_q^{-1}u)-\lambda(u))=0$, summing gives $\ell^n\sum_uk_u=-(q-1)$. Therefore $\prod_uY_u^{\lambda(v_q^{-1}u)-\lambda(u)}
\equiv\prod_uY_u^{-1}\cdot Y_0^{q-1}$, and $W_h\equiv Y_0\cdot Y_0^{-q}\cdot Y_0^{q-1}=1$.

Now compare $\gamma W_1$ and $\gamma W_{\gamma^{-1}}$. Write $\phi Y_0=y'\omega_0$ and $\tau Y_{v_qu}=
Y_u\omega_u$ with $\omega$'s in $\mu_{\ell^n}$, on which all conjugates of $\gamma$ act alike. Then
$\gamma\phi Y_0=\gamma y'\cdot\gamma\omega_0$, $\phi\gamma Y_0=(\phi\gamma\phi^{-1})(y'\omega_0)=
(\phi\gamma\phi^{-1})y'\cdot\gamma\omega_0$, and similarly for $\tau$; hence
\[
\frac{\gamma W_1}{\gamma W_{\gamma^{-1}}}=\frac{\phi\gamma Y_0}{\gamma\phi Y_0}\prod_u\Bigl(
\frac{\gamma\tau Y_{v_qu}}{\tau\gamma Y_{v_qu}}\Bigr)^{\lambda(u)}
=\frac{(\phi\gamma\phi^{-1})y'}{\gamma y'}\prod_u\Bigl(\frac{\gamma Y_u}{(\tau\gamma\tau^{-1})Y_u}
\Bigr)^{\lambda(u)}=\zeta^{\sum_uL_q(u)F_\tau(x+u/q)-F_\phi(qx)}.
\]
$W_1$ is a root of unity $\equiv1$ modulo $A$, so $\gamma W_1$, a power of $W_1$, is $\equiv1$ modulo
$A$; and $\gamma W_{\gamma^{-1}}\equiv1$ modulo $\gamma\gamma^{-1}A=A$. So the $\ell^n$-th root of
unity on the right is $\equiv1$ modulo $A$, hence equal to $1$ ($q\ne\ell$). This is \eqref{pod:eq-FI}.

\emph{The bottom class.} Let $\kappa_\emptyset$ vanish on $I_B$, and let $a\in\Z$ with $a\equiv
\varepsilon_{N\ell^n}(\gamma)$. Choose $r\in\Omega$ with $r^{\ell^n}=p$ and $\tau\in I_B$ with $\tau r=
\zeta r$ (tame inertia at $B$ maps onto $\mu_{\ell^n}$, $p\ne\ell$). Note $\tau\in H_{N\ell^n}$ as
$p\nmid N\ell$. For $N\ell^nx=0$, $(y(x)/r^{m_c(x)})^{\ell^n}=J(x)p^{-m_c(x)}$ is a unit at $B$ by (J3),
so $\tau$ moves $y(x)/r^{m_c(x)}$ by an $\ell^n$-th root of unity $\equiv1$ modulo $B$, i.e.\ fixes
it: $\tau y(x)=\zeta^{m_c(x)}y(x)$. With $\gamma y(x)=y(ax)\omega$:
$(\tau\gamma\tau^{-1})y(x)=\tau\gamma(\zeta^{-m_c(x)}y(x))=\zeta^{-am_c(x)}\tau(y(ax)\omega)=
\zeta^{m_c(ax)-am_c(x)}\gamma y(x)$. Hence $F_\tau=am_c-m_c(a\,\cdot)$ on the $N\ell^n$-torsion, and
$0=\kappa_\emptyset(\tau)=a\Lambda_\emptyset(m_c)-\Lambda_\emptyset(m_c(a\cdot))$. By
Lemma~\ref{pod:equiv} with $g=\gamma$, $\Lambda_\emptyset(m_c(a\,\cdot))=\Upsilon_{e,\psi}(\gamma)^{-1}
\Lambda_\emptyset(m_c)$, and $\Lambda_\emptyset(m_c)=\podth_{c,N\ell^n}(e,\psi)$ in $\podLam_n$. So
$(a\Upsilon_{e,\psi}(\gamma)-1)\podth=0$, i.e.\ $(\Upsilon_1(\gamma)-1)\podth=0$.
\end{proof}

\subsection{Descent and detecting points for large \texorpdfstring{$\ell$}{l}}\label{pod:ss-descent}

In this subsection $e\ge1$, $\psi$ is odd, $\ell>N+e+3$, and $\podC=\podC_{e,\psi,\nu}$. Then $2$ and
$\varphi(N\ell)=\varphi(N)(\ell-1)$ are units of $\auO$ ($0<\varphi(N),\ell-1<\ell$). Moreover there
is $a\equiv1\pmod N$ prime to $\ell$ with $a^e-1$ and $a^{e+1}-1$ units of $\auO$: take $a\equiv1
\pmod N$ and $a\equiv g\pmod\ell$ with $g$ a primitive root modulo $\ell$ (as $\ell>N$, $\ell\nmid N$);
then $a^k\not\equiv1\pmod\ell$ for $0<k<\ell-1$, and $e+1<\ell-1$. Fix such an $a$. For a prime $r$,
$x\colon G_\Q\to V$ is \emph{trivial above $r$} if $x$ vanishes on $D_A$ for every place $A$ above
$r$. $\mathcal Z^{(m)}(T)$ denotes the module of locally constant cocycles $G_\Q\to\podLam_m$ for
$\Upsilon_1=\pgcyc^{e+1}\psi$ vanishing on $\mathcal I(T)$.

\begin{lemma}\label{pod:trivabove}
Let $r\equiv1\pmod{N\ell}$ be prime, $A$ a place above $r$, $\sigma$ a Frobenius at $A$. If $x\in
\podC$, $\ell x=0$ and $x(\sigma)=0$, then $x$ is trivial above $r$.
\end{lemma}
\begin{proof}
Let $A'=\rho A$ be any place above $r$. $I_{A'}\subseteq H_N$ ($\mu_N$ is unramified at $r$) and
$A'\nmid\ell$, so $x|_{I_{A'}}=0$. $\sigma\in D_A\subseteq H_{N\ell}$ acts trivially on
$\podLam_\nu[\ell]$, so $x(\rho\sigma\rho^{-1})=\Upsilon_{e,\psi}^{-1}(\rho)x(\sigma)=0$
(Lemma~\ref{pod:coc}(b)). $D_{A'}$ is topologically generated by $I_{A'}$ and the Frobenius
$\rho\sigma\rho^{-1}$; apply Lemma~\ref{pod:coc}(a).
\end{proof}

\begin{lemma}[Reciprocity modulo $\ell$]\label{pod:recl}
Let $q\equiv1\pmod{N\ell}$ be prime, $B$ a place above $q$, $\phi$ a Frobenius at $B$, $\tau\in I_B$.
Let $x\in\podC$ with $\ell x=0$, let $T$ be a finite set of primes, and let $\kappa\in\mathcal
Z^{(m)}(T)$ with $\ell\kappa=0$. Assume that at every place $A$ not above $q$, either $A$ lies above
no prime of $T$, or $x$ vanishes on $D_A$; and that $2(a^e-1)x(\phi)\ne0$. Then $\kappa(\tau)=0$.
\end{lemma}
\begin{proof}
Let $F=\Q(\mu_{N\ell})$. Let $V=\auO\kappa(\tau)$, an $\mathbb F_\ell$-vector space with basis
$\beta_1,\dots,\beta_k$. For each $i$ extend $v\mapsto\zeta_\ell^{\beta_i^*(v)}$ ($\zeta_\ell$ a fixed
primitive $\ell$-th root of unity, $\beta_i^*$ the coordinate) to a homomorphism $\lambda_i\colon\podLam_m[\ell]\to\mu_\ell$
(extend the coordinates from $V$ to the $\mathbb F_\ell$-vector space $\podLam_m[\ell]$, which contains
all values of $\kappa$ since $\ell\kappa=0$). On $G_F$, $\kappa$ is a homomorphism ($\Upsilon_1\equiv1$ modulo $\ell$
there) and $\lambda_i\circ\kappa|_{G_F}$ is a continuous homomorphism to $\mu_\ell$, i.e.\ a cocycle for
the natural action; by Hilbert's Theorem~90 there is $r_i\in\Omega^\times$ with
$\lambda_i(\kappa(g))=gr_i/r_i$ for $g\in G_F$, and $y_i=r_i^\ell\in F^\times$. Choose $\sigma_u\in
G_\Q$ with $\varepsilon_{N\ell}(\sigma_u)=u$ for $u\in(\Z/N\ell)^\times$ and let $k_{i,u}=
\podord_B(\sigma_uy_i)$. Apply Lemma~\ref{pod:artin} with $m=\ell$, $y=y_i$: at a place $A$ not above
$\ell q$, either $x|_{D_A}=0$, or $A$ is above no prime of $T$, and then $\kappa$ vanishes on
$I_A\cap H_N\supseteq I_A\cap G_F$, so $r_i$ is fixed by $I_A\cap G_F$ and $y_i=1\cdot r_i^\ell$.
Hence $s_i:=\sum_uk_{i,u}u^e\psi(u)$ kills $2(a^e-1)x(\phi)\ne0$, so $s_i\in\podm$.

Let $\pi\in\Omega$ with $\pi^\ell=q$ and $j\in\Z/\ell$ with $\tau\pi=\zeta_\ell^j\pi$. Note $\tau\in
I_B\subseteq H_{N\ell}=G_F$. For each $u$: $\kappa(\sigma_u^{-1}\tau\sigma_u)=\Upsilon_1(\sigma_u)^{-1}
\kappa(\tau)=(u^{e+1}\psi(u))^{-1}\kappa(\tau)$ (Lemma~\ref{pod:coc}(b)). Apply $\sigma_u$ to
$\lambda_i(\kappa(\sigma_u^{-1}\tau\sigma_u))=\sigma_u^{-1}\tau\sigma_u r_i/r_i$: the left side becomes
its $u$-th power, the right side $\tau(\sigma_ur_i)/\sigma_ur_i$. The element $\sigma_ur_i/\pi^{k_{i,u}}$
has $\ell$-th power $\sigma_uy_i\,q^{-k_{i,u}}$, a unit at $B$, so $\tau$ moves it by an $\ell$-th root
of unity congruent to $1$ modulo $B$, i.e.\ fixes it; hence $\tau(\sigma_ur_i)/\sigma_ur_i=
\zeta_\ell^{jk_{i,u}}$. Therefore the $i$-th coordinate of $u\cdot(u^{e+1}\psi(u))^{-1}\kappa(\tau)$ is
$jk_{i,u}$, i.e.\ $\kappa(\tau)=u^e\psi(u)\sum_ijk_{i,u}\beta_i$. Summing over $u$:
$\varphi(N\ell)\kappa(\tau)=\sum_ijs_i\beta_i\in\podm V=\podm\kappa(\tau)$. So $(\varphi(N\ell)-t)
\kappa(\tau)=0$ for some $t\in\podm$; $\varphi(N\ell)-t$ is a unit, and $\kappa(\tau)=0$.
\end{proof}

\begin{definition}\label{pod:def-tight}
A finite set $T$ of primes $\equiv1\pmod{N\ell}$ is \emph{tight} (at level $m$) if every $\kappa\in
\mathcal Z^{(m)}(T)$ with $\ell\kappa=0$ lies in $\mathcal Z^{(m)}(\emptyset)$.
\end{definition}

\begin{lemma}\label{pod:tight}
Let $T$ be tight, $z\in\podC$ with $\ell z=0$, trivial above every prime of $T$ but not above a prime
$r\equiv1\pmod{N\ell}$. Then $T\cup\{r\}$ is tight.
\end{lemma}
\begin{proof}
Let $\kappa\in\mathcal Z^{(m)}(T\cup\{r\})$, $\ell\kappa=0$. It suffices to show that $\kappa$
vanishes on $I_A\cap H_N$ for places $A$ above $r$ (then $\kappa\in\mathcal Z^{(m)}(T)$ and tightness
applies). Let $\sigma$ be a Frobenius at $A$. By Lemma~\ref{pod:trivabove}, $z(\sigma)\ne0$, hence
$2(a^e-1)z(\sigma)\ne0$. At a place $A'$ not above $r$, either $A'$ lies above no prime of $T\cup\{r\}$,
or $A'$ lies above a prime of $T$ and $z|_{D_{A'}}=0$. Lemma~\ref{pod:recl} (with $q=r$, the set
$T\cup\{r\}$) gives $\kappa(\tau)=0$ for $\tau\in I_A$.
\end{proof}

\begin{lemma}\label{pod:finite}
$\podC$ and $\mathcal Z^{(m)}(\emptyset)$ are finite. If $\podC\ne0$, then for every $f\ge1$ there is a
finite set $Q$ of primes $\equiv1\pmod f$ such that for every $\kappa\in\mathcal Z^{(m)}(\emptyset)$
not vanishing on $H_f$ and every $0\ne z\in\podC$ there are $r\in Q$, a place $A$ above $r$ and a
Frobenius $\sigma$ at $A$ with $\kappa(\sigma)\ne0\ne z(\sigma)$. Every nonzero $z\in\podC$ is
nonzero on every $H_f$.
\end{lemma}
\begin{proof}
On $H_{N\ell^{\max(\nu,m)}}$ these cocycles are homomorphisms to a finite group, unramified outside
$N\ell$ (they vanish on inertia in $H_N$ at places prime to $\ell$, and inertia at places prime to
$N\ell$ lies in $H_N$); by Hermite--Minkowski there are finitely many such, and a cocycle is determined
by its restriction and its values on coset representatives. The last statement is
Lemma~\ref{pod:coc}(d) with $\gamma=c$, as $\Upsilon_{e,\psi}^{-1}(c)-1=-2$ is a unit. For a pair
$(\kappa,z)$ as stated, the open subgroups $Z(\kappa)$ and $Z(z)$ do not contain $H_f$; a group is not
the union of two proper subgroups, so some $h\in H_f$ lies in neither. Let $U$ be the intersection of $H_f$ with the normal cores of
$Z(\kappa)$ and $Z(z)$, an open normal subgroup of $G_\Q$ inside $H_f$. For $k\in U$ the cocycle rule
gives $\kappa(hk)=\kappa(h)$ and $z(hk)=z(h)$, so every element of $hU\subseteq H_f$ lies in neither
subgroup. By Chebotarev, applied to the extension cut out by $U$, $hU$ contains a Frobenius at a place
above a prime $r$ unramified there, and $r\equiv1\pmod f$ because $hU\subseteq H_f$. Choose one $r$ for each of the
finitely many pairs.
\end{proof}

\begin{proposition}[Kolyvagin--Rubin detection]\label{pod:koly}
Let $X$ be a nonzero Stickelberger bound for $(N,e)$, $\psi$ odd, $\ell>N+e+3$. There are
$\gamma_0,\dots,\gamma_{v_\ell(X)}\in G_\Q$ such that $x\in\podC$ and $x(\gamma_i)=0$ for all $i$ imply
$x=0$.
\end{proposition}
\begin{proof}
Put $m=v_\ell(X)+1$ and suppose the contrary: for every $\delta_1,\dots,\delta_m$ some $0\ne x\in\podC$
vanishes at all $\delta_j$. \emph{Step 1.} For every set $T$ of at most $m$ primes $\equiv1\pmod
{N\ell}$ there is $0\ne z\in\podC$ with $\ell z=0$, trivial above every prime of $T$: choose for each
$r\in T$ a Frobenius $\sigma_r$ as in Lemma~\ref{pod:trivabove}, take $x\ne0$ vanishing at all
$\sigma_r$, and let $z=\ell^kx$ with $\ell^kx\ne0=\ell^{k+1}x$; $z$ still vanishes at the $\sigma_r$.

\emph{Step 2 (set-up).} As in the proof of Lemma~\ref{pod:locals}(b) there is $\gamma\in I_A$
($A\mid\ell$) with $\pgcyc(\gamma)=a$ and $\gamma\in H_N$, so
$\Upsilon_1(\gamma)-1=a^{e+1}-1$ is a unit. Let $Q$ be as in Lemma~\ref{pod:finite} for
$f=aN\ell^m$. By Step~1 with $T=\emptyset$ take $0\ne z_0\in\podC$, $\ell z_0=0$; by
Lemma~\ref{pod:finite} and Chebotarev there are a prime $p\equiv1\pmod{N\ell^m\prod_{r\in Q}r}$, a place
$B$ above $p$ and a Frobenius $\tau_0$ at $B$ with $z_0(\tau_0)\ne0$. So $z_0$ is not trivial above
$p$, and $\{p\}$ is tight by Lemma~\ref{pod:tight} (the empty set is tight). Let $(\kappa_S)_{S
\subseteq Q}$ be the Kolyvagin system of Proposition~\ref{pod:KS} for $(n=m,c=a,\gamma,P=Q,p,B)$; the
primes of $Q$ are $\equiv1\pmod{aN\ell^m}$.

\emph{Step 3 (descent).} We show by descending induction on $i\in[0,m]$: for $S\subseteq Q$ with
$\#S\le i$ and $S\cup\{p\}$ tight, $\ell^i\kappa_S\in\mathcal Z^{(m)}(\emptyset)$. For $i=m$,
$\ell^m\kappa_S=0$. Let $i<m$ and $S$ as stated. It suffices to show $\ell^{i+1}\kappa_S=0$: then
$\ell\cdot\ell^i\kappa_S=0$, $\ell^i\kappa_S\in\mathcal Z^{(m)}(S\cup\{p\})$ by (KS1), and tightness
gives the claim. Suppose $\kappa':=\ell^{i+1}\kappa_S\ne0$. By induction ($\#S\le i+1$) $\kappa'\in
\mathcal Z^{(m)}(\emptyset)$; $\kappa'(\gamma)=0$ (KS2) and $(\Upsilon_1(\gamma)-1)\kappa'\ne0$, so by
Lemma~\ref{pod:coc}(d) $\kappa'$ does not vanish on $H_{aN\ell^m}$. By Step~1 (with $T=S\cup\{p\}$, of
size $\le i+1\le m$) there is $0\ne w\in\podC$, $\ell w=0$, trivial above $p$ and the primes of $S$. By
Lemma~\ref{pod:finite} there are $r\in Q$, $A\mid r$ and a Frobenius $\sigma$ at $A$ with
$\kappa'(\sigma)\ne0\ne w(\sigma)$. Then $w$ is not trivial above $r$, so $r\notin S$, and
$S\cup\{r\}\cup\{p\}$ is tight by Lemma~\ref{pod:tight}. By induction ($\#(S\cup\{r\})\le i+1$),
$\ell^{i+1}\kappa_{S\cup\{r\}}\in\mathcal Z^{(m)}(\emptyset)$, which vanishes on $I_A$ (as $I_A\subseteq
H_N$, $A\nmid\ell$, so $I_A\subseteq\mathcal I(\emptyset)$). By (KS3) there is $\iota\in I_A$ with
$\kappa'(\sigma)=\ell^{i+1}\kappa_S(\sigma)=\ell^{i+1}\kappa_{S\cup\{r\}}(\iota)=0$, a contradiction.

\emph{Step 4.} With $i=0$, $S=\emptyset$: $\kappa_\emptyset\in\mathcal Z^{(m)}(\emptyset)$, so it
vanishes on $I_B\subseteq\mathcal I(\emptyset)$ ($p\equiv1\pmod{N\ell}$). Proposition~\ref{pod:KS}
gives $(a^{e+1}-1)\podth_{a,N\ell^m}(e,\psi)=0$, hence $\podth_{a,N\ell^m}=0$ in $\podLam_m$, and by
Lemma~\ref{pod:thetadiv} $\ell^m\mid(a^{e+1}-1)X$, so $\ell^m\mid X$ in $\auO$. As
$X=\ell^{m-1}X'$ with $\ell\nmid X'$, this gives $\ell\mid X'$ in $\auO$ ($\auO$ is a domain); but $X'$ is
a unit of $\auO$, so $\ell$ would be a unit, which is absurd.
\end{proof}

\subsection{Proof of Proposition~\ref{pod:odd}}\label{pod:ss-proof}

For $0<e<t$ let $X_e$ be as in Lemma~\ref{pod:X} and $Y_e=16\varphi(N)X_e^2$, and put
\[
\pgcodd=\Bigl(\prod_{0<e<t}Y_e\Bigr)^{y^{20}}h_E,\qquad y=N+t+3 .
\]
\emph{Size.} For $0<e<t$, $N+e+3\le y$, so $X_e\le2^{y^6}$. Also $16\varphi(N)\le2^{4+\log_2N}
\le2^{2y^6}$ (as $y\ge4$). Hence $Y_e=16\varphi(N)X_e^2\le2^{2y^6+2y^6}=2^{4y^6}$. There are $t-1<y$
values of $e$, so $\prod_{0<e<t}Y_e\le2^{4y^7}$ and $\pgcodd\le2^{4y^7\cdot y^{20}}h_E=2^{4y^{27}}h_E$.

\emph{$e=0$:} Lemma~\ref{pod:e0}, as $h_E\mid\pgcodd$.

\emph{$0<e<t$, $\ell\le y^2$:} by Corollary~\ref{pod:kill}, $Y_e$ kills $\podSha_{e,\nu}$; by
Lemma~\ref{pod:detect} (with $y\ge N$), at most $y^{20}$ points detect it (elements of
$\podSha_{e,\nu}$ vanish on inertia off $\ell$). Lemma~\ref{pod:length} gives
$\prilen\podSha_{e,\nu}\le\prival(Y_e^{y^{20}})\le\prival(\pgcodd)$.

\emph{$0<e<t$, $\ell>y^2$:} then $\ell>N+e+3$. For each odd $\psi$ let $\gamma_{\psi,i}$ ($0\le i\le
v_\ell(X_e)$) be the points of Proposition~\ref{pod:koly}, and consider the points $b(\sigma_u,
\gamma_{\psi,i})\in\pgG$ of Lemma~\ref{pod:split}. Their number is at most $\varphi(N)^2(v_\ell(X_e)+1)
\le y^2(y^6+1)\le y^{20}$ (as $v_\ell(X_e)\le\log_2X_e\le y^6$). If $x\in\podSha_{e,\nu}$ vanishes at all
of them, then $z_\psi(\gamma_{\psi,i})=0$, so $z_\psi=0$ for every odd $\psi$, and $4\varphi(N)x=0$ by
Lemma~\ref{pod:split}(b); since $4\varphi(N)<\ell$ is prime to $\ell$ and $\ell^\nu x=0$, $x=0$. With
the annihilator $Y_e$ (Corollary~\ref{pod:kill}), Lemma~\ref{pod:length} gives $\prilen
\podSha_{e,\nu}\le\prival(Y_e^{y^{20}})\le\prival(\pgcodd)$. \qed

\begin{remark}[The bound $\log\pgcodd\le x^{800}$]\label{pod:rem-size}
Proposition~\ref{pg:odd} and Lemma~\ref{pg:sums} use $\log\pgcodd\le x^{800}$. We check it from the
data (Lemma~\ref{an:lem-sizes} and \S\ref{ss:au-spaces}). \emph{Inputs:} $f\le q_1\le x^{12}$ (S3);
$\sqrt d\,q_1\prod_{p\mid q_0}p^{-1/(p-1)}=\anrd_E\le x^{12}$ with $q_1\ge x^{11}$ (S3, S4), so
$\sqrt d\le x\,q_0$, and $q_0\le\log x$ (S1); hence $N=q_1d\le x^{14}(\log x)^2\le x^{15}$. Further
$x^5\le r\le3x^5$ and $t=5r\le15x^5$ (\S\ref{ss:au-spaces}). \emph{The first factor:} $y=N+t+3\le
x^{15}+15x^5+3\le2x^{15}$, so $4y^{27}\log2\le4\cdot2^{27}x^{405}\log2\le x^{406}$. \emph{The class
number:} $[E:\Q]=2f\le2x^{12}$ and $|d_E|=\anrd_E^{\,2f}$, so $\log|d_E|\le2x^{12}\cdot12\log x$.
Every ideal class of $E$ contains an integral ideal of norm $\le|d_E|^{1/2}$ (Minkowski: the
constant $(4/\pi)^{r_2}n!/n^n$ is $\le1$), and there are at most $k^{[E:\Q]}$ integral ideals of norm
$k$ (the number is the $[E:\Q]$-fold divisor function at most, and $\tau_n(k)\le n^{\Omega(k)}\le
k^{\log_2n}\le k^n$), so $h_E\le\sum_{k\le|d_E|^{1/2}}k^{2f}\le|d_E|^{(2f+1)/2}$ and
\[
\log h_E\le\frac{2f+1}2\log|d_E|\le(x^{12}+1)\cdot24x^{12}\log x\le x^{26}.
\]
Therefore $\log\pgcodd\le x^{406}+x^{26}\le x^{407}\le x^{800}$. (If $N\le x^{13}$, the same computation gives $x^{352}$.)
\end{remark}


%% ===== proofs/C-kummer.tex =====
% proofs/C-kummer.tex
% Kummer maps: proof of Proposition pri:kummer(a),(b) of proofs/C-ribet.tex, and the details of
% the proof of Proposition pri:triangle.
% Labels of this file start with pku:. Conventions of proofs/C-odd.tex.
\providecommand{\auO}{\mathcal O}
\providecommand{\auNm}{\operatorname{Nm}}
\providecommand{\pgG}{G_E}
\providecommand{\pgcyc}{\varepsilon_\ell}
\providecommand{\prilen}{\operatorname{len}}
\providecommand{\prival}{\operatorname{val}}
\providecommand{\prity}{\mathrm{ty}}
% macros of this file (prefix \pku)
\providecommand{\pkuinv}{\operatorname{inv}}
\providecommand{\pkuHom}{\operatorname{Hom}}
\providecommand{\pkuGal}{\operatorname{Gal}}
\providecommand{\pkucp}{\varepsilon^{+}}
\providecommand{\pkuSel}{\operatorname{Sel}^{*}}
\providecommand{\pkuur}{\mathrm{ur}}
\providecommand{\pkuAnn}{\operatorname{Ann}}
\providecommand{\pkured}{\operatorname{red}}

\section{Kummer maps}\label{pku:sec}

This section proves part~(a) of Proposition~\ref{pri:kummer} of \S\ref{pri:sec}
(Proposition~\ref{pku:main}) and part~(b) (Proposition~\ref{pku:b}), and gives the details of the
proof of Proposition~\ref{pri:triangle}, the only place where Proposition~\ref{pri:kummer} is used
(\S\ref{pku:ss-triangle}). The proof of (a) is a count of Greenberg--Wiles type over the totally
real field $E^+$; it uses only the Greenberg--Wiles formula, local and global duality for finite
Galois modules, the explicit form of the local pairing at a split prime, and Chebotarev's theorem.
The set $T$ of part~(a) can be taken to be the smallest one allowed, the primes above $\ell$ and
below the given places; with this choice the ``relations'' defining a Kummer map are exactly
restrictions of dual Selmer classes, and the difficulty that a larger $T$ would create (relations
ramified at primes of $T$ where the test cocycles are unramified but not zero) does not arise.

\subsection{Setting and statement}\label{pku:ss-statement}

\emph{Conventions.} $\Omega=\overline\Q\subset\C$; $E\subseteq\Omega$ is the CM field of
\S\ref{ss:au-spaces}: it contains a primitive $5040$-th root of unity, and $E/E^+$ is unramified at
all finite places. Put $F=E^+$, $G_F=\pkuGal(\Omega/F)\supseteq\pgG=\pkuGal(\Omega/E)$, and let
$c\in G_F\smallsetminus\pgG$ be the fixed complex conjugation. A \emph{complex conjugation} is any
conjugate of $c$ in $G_{\Q}$ lying in $G_F$; none lies in $\pgG$ ($E$ is totally complex). For a
place $B$ of $\Omega$, $D_B\supseteq I_B$ are its decomposition and inertia groups in $G_F$; a
\emph{Frobenius element} at $B$ over a finite prime $v$ of $F$ is an element of $D_B$ acting on
the residue field as $t\mapsto t^{\auNm v}$. Fix a prime $\ell$ and $\nu\ge1$, put $n=\ell^\nu$,
and let $\pgcyc$ be the $\ell$-adic cyclotomic character of $G_\Q$. Fix a primitive $n$-th root of
unity $\zeta\in\Omega$ and let $\theta\colon\mu_n\to\frac1n\Z/\Z$ be the isomorphism
$\zeta^a\mapsto a/n$. A \emph{normalised tame generator} (for $n$) at a place $B$ over a prime $p\nmid
\ell$ of $F$ is an element $\gamma\in I_B$ with $\gamma(\varpi^{1/n})=\zeta\varpi^{1/n}$ for some
uniformiser $\varpi$ of $F_p$ and some $n$-th root $\varpi^{1/n}$; these exist, and the quotient
$\gamma(\varpi^{1/n})/\varpi^{1/n}$ does not depend on the choice of $\varpi$ and of the root (as in
\S\ref{pse:ss-tame}, over $F$ instead of $E$).

Let $G_\nu=\{g\in\pgG:\pgcyc(g)\equiv1\bmod n\}$, and define $\pkucp\colon G_F\to\Z_\ell^\times$ by
$\pkucp(g)=\pgcyc(g)$ and $\pkucp(gc)=\pgcyc(g)$ for $g\in\pgG$. For $e\in\Z$ let $M=M_e$ be $\Z/n$
with $g\in G_F$ acting by $\pkucp(g)^{-e}$, let $M^\vee=\pkuHom(M,\Q/\Z)$ (as an abelian group), and
let $M^*=\pkuHom(M,\mu_n)$ be the Tate dual, with $(gf)(a)=g\bigl(f(g^{-1}a)\bigr)$. For a finite
$G_F$-module $A$ we write $A[\ell]$ for its $\ell$-torsion and $H^i(F,A)$, $H^i(F_v,A)$ for
continuous cohomology of $G_F$ and of $D_{B}$, $B$ a fixed place over $v$.

\begin{lemma}[The characters]\label{pku:basic}
\begin{enumerate}
\item[(i)] $\pkucp$ is a continuous character, equal to $\pgcyc\eta$ with $\eta$ the quadratic
character of $E/F$. It is $1$ at every complex conjugation, unramified at every prime of $F$ not
above $\ell$, and if $\ell=2$ it takes values in $1+4\Z_2$.
\item[(ii)] $G_\nu$ is open and normal in $G_F$, contains no complex conjugation, and acts trivially
on $M$, $M^*$ and $\mu_n$. The group $G_F/G_\nu$ is abelian, and its $\ell$-Sylow subgroup is
generated by at most one element if $\ell$ is odd and by at most two if $\ell=2$.
\item[(iii)] $F$ is totally real with $r_F=[F:\Q]\ge576$ real places.
\item[(iv)] Let $w$ be a place of $E$ of degree one over $\Q$, $w\nmid\ell$, $A$ a place of
$\Omega$ above $w$ and $\phi\in G_\nu$ a Frobenius element at $A$ (in $\pgG$). Then $D_A\subseteq
\pgG$ (decomposition group in $G_F$), $\phi$ is a Frobenius element at $A$ for $F$, $D_A$ is
topologically generated by $I_A$ and $\phi$, and $D_A$ acts trivially on $M$.
\end{enumerate}
\end{lemma}
\begin{proof}
(i) For $g\in\pgG$, $\pgcyc(gc)=-\pgcyc(g)$ and $\eta(gc)=-1$, so $\pkucp=\pgcyc\eta$ on both cosets;
it is a continuous character. A complex conjugation $c'$ has $\pgcyc(c')=-1$ and $c'\notin\pgG$,
so $\eta(c')=-1$ and $\pkucp(c')=1$. An inertia group $I_B$ at a finite prime lies in $\pgG$ ($E/F$
is unramified at finite places), and there $\pkucp=\pgcyc$, which is unramified off $\ell$. If
$\ell=2$: $\mu_4\subseteq E$, so $\pgcyc\equiv1\bmod4$ on $\pgG$, and $\pkucp(gc)=\pgcyc(g)$.

(ii) $G_\nu=\ker\eta\cap\ker(\pkucp\bmod n)$, an open normal subgroup; it lies in $\pgG$, so contains
no complex conjugation. On $G_\nu$, $\pkucp\equiv1\bmod n$ and $\pgcyc\equiv1\bmod n$, so $G_\nu$
acts trivially on $M$, $\mu_n$ and $M^*$. The map $g\mapsto(\eta(g),\pkucp(g)\bmod n)$ embeds
$G_F/G_\nu$ in $\{\pm1\}\times(\Z/n)^\times$. If $\ell$ is odd, the $\ell$-Sylow subgroup of the
latter is cyclic. If $\ell=2$, the image is generated by the image of $\pgG$, contained in the
cyclic group $\{1\}\times(1+4\Z)/n$, and by the image of $c$; so it is generated by two elements,
and so is its $2$-Sylow subgroup (the whole group).

(iii) $F$ is the maximal real subfield of a CM field, so totally real, and $[E:\Q]\ge[\Q(\mu_{5040})
:\Q]=\varphi(5040)=1152$.

(iv) The prime $v$ of $F$ below $w$ has the same residue field $\mathbb F_p$ as $w$ and $E/F$ is
unramified at $v$; so $w$ is split over $v$, $D_A\subseteq\pgG$, and the decomposition groups of $A$
in $G_F$ and $\pgG$ coincide. In particular $\phi$ is a Frobenius element for $F$, and $D_A$ is
topologically generated by $I_A$ and $\phi$. On $I_A$, $\pkucp=\pgcyc=1$ ($w\nmid\ell$), and
$\pkucp(\phi)=\pgcyc(\phi)\equiv1\bmod n$.
\end{proof}

\begin{definition}[Kummer maps]\label{pku:def-kummer}
Let $x=(w_i)_{i\in I}$ be a finite family of degree-one places $w_i\nmid\ell$ of $E$ above distinct
primes $v_i$ of $F$, with places $A_i$ of $\Omega$ above $w_i$ and Frobenius elements $\phi_{w_i}
\in G_\nu$ at $A_i$. Let $S_0=S_0(x)$ be the set of primes of $F$ above $\ell$ or equal to some
$v_i$, and fix a place $B_v$ of $\Omega$ above each $v\in S_0$, with $B_{v_i}=A_i$.

(a) A \emph{test cocycle for $x$ and $P$}, where $P$ is a finite set of primes of $F$ not in
$S_0$, is a continuous cocycle $z\colon G_F\to M$ that vanishes on $I_B$ for every place $B$ over a
prime not in $S_0\cup P$; vanishes on $D_{B_{v_i}}$ for every $i$; vanishes on $D_{B_v}\cap\ker
\pkucp$ for every $v\mid\ell$; and vanishes at every complex conjugation.

(b) A \emph{split prime} is a prime $p\notin S_0$ of $F$ with a chosen place $B_p$ over it, a
Frobenius element $\phi_p\in G_\nu$ at $B_p$, and a normalised tame generator $\gamma_p$ at $B_p$.

(c) Let $T\supseteq S_0$ be a finite set of primes of $F$. A \emph{relation off $T$} is a function
$y\colon G_F\to M^\vee$ such that: $y$ is locally constant; $y=0$ outside $G_\nu$; $y(uv)=y(u)+y(v)$
for $u,v\in G_\nu$; $y(gug^{-1})(a)=\pgcyc(g)\,y(u)(g^{-1}a)$ for $g\in G_F$, $u\in G_\nu$, $a\in
M$; $y=0$ on $I_B$ for every place $B$ over a prime not in $T$; and
\begin{equation}\label{pku:eq-recip}
\sum_{p\in P}y(\phi_p)\bigl(z(\gamma_p)\bigr)=0
\end{equation}
for every finite set $P$ of split primes not in $T$ and every test cocycle $z$ for $x$ and $P$.

(d) A \emph{Kummer map of exponent $e$ for $x$ off $T$} is a function $\kappa\colon\pgG\to\Z/n$
such that $y(g)(a)=a\kappa(g)/n$ ($g\in\pgG$, $a\in M$) for some relation $y$ off $T$. Its
\emph{reduction} is $\pkured\kappa\colon G_\nu\to\mathbb F_\ell$, $g\mapsto\kappa(g)\bmod\ell$.
\end{definition}

\begin{remark}\label{pku:rem-def}
(1) A test cocycle vanishes on $gHg^{-1}$ for every $g\in G_F$ and every subgroup $H$ on which it
vanishes and which acts trivially on $M$: $z(ghg^{-1})=z(g)+gz(h)-ghg^{-1}z(g)=0$, as $ghg^{-1}$
acts trivially. The groups $D_{B_{v_i}}$ and $D_{B_v}\cap\ker\pkucp$ act trivially on $M$
(Lemma~\ref{pku:basic}(iv) and the definition of $\pkucp$), so the conditions do not depend on the
chosen places. (2) The equivariance law, with $\pkucp$ acting on $M$, reads
$y(gug^{-1})=\pgcyc(g)\pkucp(g)^{e}\,y(u)$; for $g\in\pgG$ it says $\kappa(gug^{-1})=
\pgcyc(g)^{e+1}\kappa(u)$, and for $g=c$ that $\kappa(cuc)=-\kappa(u)$ ($u\in G_\nu$).
\end{remark}

\begin{proposition}[Kummer maps]\label{pku:main}
Let $x=(w_i)_{i\in I}$ be as in Definition~\ref{pku:def-kummer}, $e\in\Z$, and $(f_i)_{i\in I}$ any
family of maps $G_\nu\to\mathbb F_\ell$ indexed by the same set $I$. Then there is a Kummer map
$\kappa$ of exponent $e$ for $x$ off $T=S_0(x)$ whose reduction is not in the $\mathbb F_\ell$-span
of the $f_i$. More precisely, the reductions of the Kummer maps of exponent $e$ for $x$ off $S_0(x)$
contain an $\mathbb F_\ell$-subspace of dimension at least $\#I+r_F-4\ge\#I+1$.
\end{proposition}

\subsection{Duality: the inputs}\label{pku:ss-duality}

Let $A$ be a finite $G_F$-module of $\ell$-power order killed by $n$, and $A^*=\pkuHom(A,\mu_n)$.
For a place $v$ of $F$ the cup product and the evaluation $A\times A^*\to\mu_n$ give
$H^1(F_v,A)\times H^1(F_v,A^*)\to H^2(F_v,\mu_n)$, followed by the local invariant into $\Q/\Z$
(for real $v$, $H^2(\R,\mu_n)=\mathrm{Br}(\R)[n]\subseteq\frac12\Z/\Z$). For a subgroup $L_v\subseteq
H^1(F_v,A)$, $L_v^\perp\subseteq H^1(F_v,A^*)$ is its annihilator. We use:

\smallskip\noindent(D1) \emph{Local duality.} For finite $v$ the pairing $H^1(F_v,A)\times H^1(F_v,
A^*)\to\Q/\Z$ is perfect; if $v\nmid\ell$ and $A$ is unramified at $v$, the unramified subgroups
$H^1_\pkuur(F_v,A)=\ker(H^1(F_v,A)\to H^1(I_B,A))$ and $H^1_\pkuur(F_v,A^*)$ are exact annihilators of
each other (\cite[Ch.~I, Corollary~2.3 and Theorem~2.6]{MilneADT}; \cite[Theorems~7.2.6 and 7.2.15]{NSW}).

\smallskip\noindent(D2) \emph{Global reciprocity.} For $a\in H^1(F,A)$ and $b\in H^1(F,A^*)$, the
local invariants of $a_v\cup b_v$ vanish for almost all $v$ and $\sum_v\pkuinv_v(a_v\cup b_v)=0$, the
sum over all places of $F$; each term may be computed at any place of $\Omega$ above $v$. (This is
the reciprocity law of the Brauer group, $\sum_v\pkuinv_v=0$ on $\mathrm{Br}(F)$,
\cite[Theorem~8.1.17]{NSW}.)

\smallskip\noindent(D3) \emph{The Greenberg--Wiles formula.} Let $\Sigma$ be a finite set of places of
$F$ containing the archimedean ones and those above $\ell$, with $A$ unramified outside $\Sigma$; let
$\mathcal L=(L_v)_{v\in\Sigma}$, $L_v\subseteq H^1(F_v,A)$, and $\mathcal L^*=(L_v^\perp)_{v\in\Sigma}$.
Put $H^1_{\mathcal L}(F,A)=\{a\in H^1(G_\Sigma,A):a_v\in L_v\ (v\in\Sigma)\}$, $G_\Sigma$ the Galois
group of the maximal extension of $F$ unramified outside $\Sigma$, and likewise $H^1_{\mathcal L^*}(F,
A^*)$. These groups are finite and
\begin{equation}\label{pku:eq-GW}
\frac{\#H^1_{\mathcal L}(F,A)}{\#H^1_{\mathcal L^*}(F,A^*)}=\frac{\#H^0(F,A)}{\#H^0(F,A^*)}
\prod_{v\in\Sigma}\frac{\#L_v}{\#H^0(F_v,A)} ,
\end{equation}
with ordinary $H^0$ at the real places (\cite[Theorem~2.19]{DDT}; \cite[Theorem~8.7.9]{NSW}; over $\Q$ and for $\ell$ odd, \cite[Proposition~1.6]{WilesFLT}). The proof is the Poitou--Tate sequence and Tate's global Euler
characteristic formula.

\smallskip\noindent(D4) \emph{Chebotarev's density theorem}: for a finite Galois extension $K/F$ and
$\sigma\in\pkuGal(K/F)$ there are infinitely many primes $p$ of $F$, unramified in $K$, with a place
of $K$ over $p$ whose Frobenius element is $\sigma$.

\smallskip
Statement (D2) is used only for $a$ unramified outside a finite set; the terms at the places where
$a$ and $b$ are both unramified and $v\nmid\ell$ vanish by (D1).

\begin{lemma}[The pairing at a split prime]\label{pku:pairing}
Let $p\nmid\ell$ be a prime of $F$, $B$ a place over it, $A$ a finite $G_F$-module killed by $n$
on which $D_B$ acts trivially, and suppose $\auNm p\equiv1\pmod n$ (so $D_B$ acts trivially on
$\mu_n$ and $A^*$). Let $\phi\in D_B$ be a Frobenius element and $\gamma$ a normalised tame generator
at $B$. For a homomorphism $z\colon D_B\to A$ and an unramified homomorphism $\xi\colon D_B\to A^*$,
\[
\pkuinv_B(z\cup\xi)=\epsilon\,\theta\bigl(\xi(\phi)(z(\gamma))\bigr),
\]
where $\epsilon\in\{\pm1\}$ depends only on the conventions for cup products, local invariants and
the local reciprocity map, not on $p$, $B$, $A$, $z$, $\xi$.
\end{lemma}
\begin{proof}
This is Lemma~\ref{pse:pairing}, whose proof is local and applies verbatim to the local field
$F_p$; we recall it. Both sides are biadditive and compatible with decompositions of $A$; so we may
take $A=\Z/\ell^s$ ($s\le\nu$) with trivial action, $A^*=\mu_{\ell^s}$, and identify
$\theta_s\colon\mu_{\ell^s}\to\frac1{\ell^s}\Z/\Z$ by $\zeta^{n/\ell^s}\mapsto1/\ell^s$ (compatible
with $\theta$). By Kummer theory $\xi$ is the class of some $b\in F_p^\times$, a unit up to
$\ell^s$-th powers because $\xi$ is unramified, and $\xi(\phi)=\phi(b^{1/\ell^s})/b^{1/\ell^s}=
\bar b^{(\auNm p-1)/\ell^s}$ (as a root of unity). By local class field theory,
$\pkuinv(z\cup\partial b)=\pm z(\mathrm{rec}(b))$ (\cite[XIV, \S1, Proposition~3, p.~205]{SerreLF}), and for a unit $b$ the element $\mathrm{rec}(b)$ acts on $F_p(\varpi^{1/\ell^s})$
by $\varpi^{1/\ell^s}\mapsto\bar b^{\mp(\auNm p-1)/\ell^s}\varpi^{1/\ell^s}$ (the tame symbol,
\cite[Ch.~XIV, \S3, Proposition~8 and its Corollary, pp.~210--211]{SerreLF}; see also \S2, Proposition~6, p.~208). A homomorphism $D_B\to\Z/\ell^s$ restricted to $I_B$ factors through
the cyclic quotient of tame inertia of order $\ell^s$, generated by $\gamma$; so $z(\mathrm{rec}(b))
=kz(\gamma)$ with $\zeta^{kn/\ell^s}=\xi(\phi)^{\mp1}$, i.e.\ $k/\ell^s=\mp\theta(\xi(\phi))$. Hence
$\pkuinv(z\cup\xi)=\pm\theta(\xi(\phi))\,z(\gamma)$, with a sign fixed by the conventions.
\end{proof}

\subsection{Selmer structures}\label{pku:ss-selmer}

Fix $x$, $S_0=S_0(x)$, the places $B_v$ ($v\in S_0$) and $e$ as in Definition~\ref{pku:def-kummer};
put $k=\#I$, the number of primes of $S_0$ not above $\ell$. For every other place $v$ of $F$ fix a
place $B_v$ of $\Omega$ above it; for split primes take the chosen $B_p$. Write $D_v=D_{B_v}$,
$I_v=I_{B_v}$.

\begin{definition}\label{pku:def-sel}
Let $A\in\{M,M[\ell]\}$ and let $P$ be a finite set of split primes. For each place $v$ let $H_v
\subseteq D_v$ be: $D_v\cap\ker\pkucp$ if $v\mid\ell$; $D_v$ if $v\in S_0$, $v\nmid\ell$, or $v$ is
real; $\{1\}$ if $v\in P$; $I_v$ otherwise. Each $H_v$ acts trivially on $A$
(Lemma~\ref{pku:basic}). Put
\[
L^P_v(A)=\ker\bigl(H^1(F_v,A)\to\pkuHom(H_v,A)\bigr),
\]
the restriction to $H_v$ being a homomorphism. Let $\Sigma_P=S_0\cup P\cup\{\text{real places}\}$.
Then $\mathcal L^P(A)=(L^P_v(A))_{v\in\Sigma_P}$, and $\mathcal L(A)$ is the family with the same
$\Sigma_P$ and the same $L_v$ for $v\notin P$, but $L_p=H^1_\pkuur(F_p,A)$ for $p\in P$.
\end{definition}

So $H^1_{\mathcal L^P}(F,A)$ consists of the classes unramified outside $S_0\cup P$ with the
conditions at $S_0$ and at the real places, and $H^1_{\mathcal L}(F,A)$ of those unramified outside
$S_0$; $H^1_{\mathcal L}$ does not depend on $P$, and neither does $H^1_{\mathcal L^*}(F,A^*)$
(at $p\in P$ the dual condition is $H^1_\pkuur(F_p,A^*)$ by (D1)). We write
\[
\pkuSel=H^1_{\mathcal L^*(M)}(F,M^*)
\]
for the classes unramified outside $S_0$ that lie in $L_v(M)^\perp$ at $v\in S_0$ (no condition at
$v\in S_0$, $v\nmid\ell$, and at the real places, where $L_v=0$). The class of a test cocycle for
$x$ and $P$ lies in $H^1_{\mathcal L^P}(F,M)$, by Remark~\ref{pku:rem-def}(1).

\begin{lemma}[Local terms]\label{pku:local}
Let $A\in\{M,M[\ell]\}$.
\begin{enumerate}
\item[(a)] If $v\in S_0$, $v\nmid\ell$, or $v$ is real, then $L_v(A)=0$ and $H^0(F_v,A)=A$.
\item[(b)] If $v\mid\ell$, then $\#L_v(A)\le\#H^0(F_v,A)$.
\item[(c)] If $p$ is a split prime, then $D_p$ acts trivially on $A$ and $A^*$, $\auNm p\equiv1
\pmod n$, the map $H^1(F_p,A)\to A$, $z\mapsto z(\gamma_p)$, is surjective with kernel
$H^1_\pkuur(F_p,A)$, and $H^1_\pkuur(F_p,A^*)\to A^*$, $\xi\mapsto\xi(\phi_p)$, is bijective. In
particular $\#H^1(F_p,A)=(\#A)^2=\#A\cdot\#H^0(F_p,A)$.
\end{enumerate}
\end{lemma}
\begin{proof}
(a) Here $H_v=D_v$ acts trivially on $A$ (Lemma~\ref{pku:basic}(iv) for $v=v_i$; $D_v=\{1,c'\}$ with
$\pkucp(c')=1$ for real $v$). So $H^0(F_v,A)=A$, $H^1(F_v,A)=\pkuHom(D_v,A)$ and the restriction to
$H_v=D_v$ is the identity; $L_v(A)=0$.

(b) As $H_v$ is normal in $D_v$ and acts trivially, inflation--restriction identifies $L_v(A)$ with
$H^1(D_v/H_v,A)$. Now $D_v/H_v\cong\Gamma=\pkucp(D_v)$, a closed subgroup of $\Z_\ell^\times$ ($\ell$
odd) or of $1+4\Z_2$ ($\ell=2$, Lemma~\ref{pku:basic}(i)), hence procyclic; let $\sigma$ be a
topological generator. A continuous cocycle on $\Gamma$ is determined by its value at $\sigma$ (the
cocycle rule determines it on the dense subgroup generated by $\sigma$), and the coboundaries take
the values $(\sigma-1)A$ at $\sigma$. So $\#H^1(\Gamma,A)\le\#A/(\sigma-1)A=\#\ker(\sigma-1)=\#A^\Gamma
=\#H^0(F_v,A)$.

(c) $D_p$ is topologically generated by $I_p$, on which $\pkucp$ and the action on $\mu_n$ are
trivial ($p\nmid\ell$), and $\phi_p\in G_\nu$, which acts trivially on $A$, $A^*$, $\mu_n$
(Lemma~\ref{pku:basic}(ii)); $\phi_p$ acts on $\mu_n$ by $\auNm p$-th powers, so $\auNm p\equiv1
\pmod n$. Hence $H^1(F_p,A)=\pkuHom(D_p,A)$. A homomorphism $D_p\to A$ kills wild inertia (pro-$p$),
and its restriction to $I_p$ factors through the maximal quotient of exponent $n$ of tame inertia,
which is cyclic of order $n$ generated by the image of $\gamma_p$ (it maps isomorphically to $\mu_n$
by $\tau\mapsto\tau(\varpi^{1/n})/\varpi^{1/n}$). So the kernel of $z\mapsto z(\gamma_p)$ is the set
of homomorphisms vanishing on $I_p$, i.e.\ $H^1_\pkuur(F_p,A)$. Surjectivity: as $\mu_n\subseteq F_p$,
$\chi(g)=g(\varpi^{1/n})/\varpi^{1/n}$ is a homomorphism $D_p\to\mu_n$ with $\chi(\gamma_p)=\zeta$;
for $a\in A$, $g\mapsto n\,\theta(\chi(g))\,a$ is a homomorphism taking the value $a$ at $\gamma_p$.
Finally $H^1_\pkuur(F_p,A^*)=\pkuHom(D_p/I_p,A^*)$ and $D_p/I_p\cong\widehat\Z$ is generated by the
image of $\phi_p$.
\end{proof}

\subsection{Dual Selmer classes are relations}\label{pku:ss-relations}

\begin{lemma}\label{pku:relation}
Let $\xi\colon G_F\to M^*$ be a continuous cocycle whose class lies in $\pkuSel$. Define $y_\xi
\colon G_F\to M^\vee$ by $y_\xi(g)=\theta\circ\xi(g)$ for $g\in G_\nu$ and $y_\xi(g)=0$ otherwise.
Then $y_\xi$ is a relation off $S_0$, and $\kappa_\xi(g)=n\,y_\xi(g)(1)$ is a Kummer map of exponent
$e$ for $x$ off $S_0$. It depends only on the class of $\xi$, and $\xi\mapsto y_\xi$ is additive.
\end{lemma}
\begin{proof}
A coboundary $g\mapsto gf-f$ vanishes on $G_\nu$, which acts trivially on $M^*$; so $y_\xi$
depends only on the class, and it is clearly additive in $\xi$. It is locally constant and vanishes
off $G_\nu$. For $u,v\in G_\nu$, $\xi(uv)=\xi(u)+u\xi(v)=\xi(u)+\xi(v)$. For $g\in G_F$ and $u\in
G_\nu$, the cocycle rule gives $\xi(gug^{-1})=\xi(g)+g\xi(u)-gug^{-1}\xi(g)=g\xi(u)$; and
$\theta\bigl((gf)(a)\bigr)=\theta\bigl(g(f(g^{-1}a))\bigr)=\pgcyc(g)\,\theta\bigl(f(g^{-1}a)\bigr)$.
So $y_\xi(gug^{-1})(a)=\pgcyc(g)\,y_\xi(u)(g^{-1}a)$. Let $B$ be a place over a prime $v\notin S_0$
and $\tau\in I_B$; if $\tau\notin G_\nu$ then $y_\xi(\tau)=0$; otherwise $I_B$ acts trivially on $M^*$
($v\nmid\ell$ and Lemma~\ref{pku:basic}(i)), the class of $\xi$ is unramified at $v$, so $\xi|_{I_B}$
is a coboundary for the trivial action, i.e.\ zero, and $y_\xi(\tau)=0$.

Reciprocity: let $P$ be a finite set of split primes (not in $S_0$) and $z$ a test cocycle for $x$
and $P$. Apply (D2) to $[z]$ and $[\xi]$. At $v\in S_0$ and at real $v$, $[z]_v\in L_v(M)$ and
$[\xi]_v\in L_v(M)^\perp$, so the term vanishes. At a finite $v\notin S_0\cup P$, $v\nmid\ell$, and
both classes are unramified, so the term vanishes by (D1). At $p\in P$, $D_p$ acts trivially on $M$,
$\auNm p\equiv1\bmod n$ (Lemma~\ref{pku:local}(c)), $z|_{D_p}$ is a homomorphism and $\xi|_{D_p}$ an
unramified homomorphism; by Lemma~\ref{pku:pairing} the term is $\epsilon\,\theta\bigl(\xi(\phi_p)(
z(\gamma_p))\bigr)=\epsilon\,y_\xi(\phi_p)(z(\gamma_p))$, as $\phi_p\in G_\nu$. Hence $\epsilon\sum_{p\in
P}y_\xi(\phi_p)(z(\gamma_p))=0$, which is~\eqref{pku:eq-recip}. Finally $y_\xi(g)(a)=a\,y_\xi(g)(1)
=a\kappa_\xi(g)/n$.
\end{proof}

\subsection{The values at Frobenius elements}\label{pku:ss-image}

Let $P$ be a finite set of split primes. Define
\[
\lambda_P\colon H^1_{\mathcal L^P}(F,M)\to M^P,\quad z\mapsto(z(\gamma_p))_{p\in P},\qquad
\mu_P\colon\pkuSel\to(M^\vee)^P,\quad\xi\mapsto(\theta\circ\xi(\phi_p))_{p\in P}.
\]
They are well defined ($\gamma_p$ acts trivially on $M$ and $\phi_p$ on $M^*$, so coboundaries vanish
there). For $b\in M^P$ and $f\in(M^\vee)^P$ write $\langle b,f\rangle=\sum_pf_p(b_p)$; this pairing is
perfect. For $B\subseteq M^P$ let $\pkuAnn(B)\subseteq(M^\vee)^P$ be its annihilator.

\begin{lemma}\label{pku:image}
$\mu_P(\pkuSel)=\pkuAnn\bigl(\lambda_P(H^1_{\mathcal L^P}(F,M))\bigr)$.
\end{lemma}
\begin{proof}
Put $V=\bigoplus_{p\in P}H^1(F_p,M)/H^1_\pkuur(F_p,M)$ and $V'=\bigoplus_{p\in P}H^1_\pkuur(F_p,M^*)$. By
(D1) the local pairings induce a perfect pairing $V\times V'\to\Q/\Z$, and by
Lemma~\ref{pku:local}(c) $z\mapsto(z(\gamma_p))$ and $\xi\mapsto(\theta\circ\xi(\phi_p))$ identify $V$
with $M^P$ and $V'$ with $(M^\vee)^P$; by Lemma~\ref{pku:pairing} the pairing becomes $\epsilon
\langle\ ,\ \rangle$. Let $\lambda'\colon H^1_{\mathcal L^P}(F,M)\to V$ and $\mu'\colon\pkuSel\to V'$ be
the localisation maps, which become $\lambda_P$, $\mu_P$. Their kernels are $H^1_{\mathcal L}(F,M)$
(classes unramified at $P$) and $H^1_{(\mathcal L^P)^*}(F,M^*)$ (classes of $\pkuSel$ vanishing at
$P$; the dual of the condition $L^P_p=H^1(F_p,M)$ is $0$).

The inclusion $\mu'(\pkuSel)\subseteq\pkuAnn(\lambda'(H^1_{\mathcal L^P}))$ is (D2) applied to $a\in
H^1_{\mathcal L^P}(F,M)$, $b\in\pkuSel$: the terms at $S_0$ and the real places vanish ($a_v\in L_v$,
$b_v\in L_v^\perp$), those at the other places off $P$ vanish by (D1), and the terms at $P$ add up to
the pairing of $\lambda'(a)$ and $\mu'(b)$.

Counting: $\#\mu'(\pkuSel)=\#\pkuSel/\#H^1_{(\mathcal L^P)^*}(F,M^*)$, and
$\#\pkuAnn(\lambda'(H^1_{\mathcal L^P}))=\#V/\#\lambda'(H^1_{\mathcal L^P})=\#V\cdot\#H^1_{\mathcal L}(F,M)/
\#H^1_{\mathcal L^P}(F,M)$. Apply~\eqref{pku:eq-GW} with $\Sigma=\Sigma_P$ to $\mathcal L^P(M)$ and to
$\mathcal L(M)$ and divide: all factors agree except those at $p\in P$, whose quotient is
$\#H^1(F_p,M)/\#H^1_\pkuur(F_p,M)$; so
\[
\frac{\#H^1_{\mathcal L^P}(F,M)}{\#H^1_{(\mathcal L^P)^*}(F,M^*)}=\frac{\#H^1_{\mathcal L}(F,M)}
{\#\pkuSel}\cdot\#V .
\]
Hence the two finite groups have the same order, and the inclusion is an equality.
\end{proof}

\subsection{The \texorpdfstring{$\ell$}{l}-torsion of the relaxed Selmer group}\label{pku:ss-torsion}

\begin{lemma}\label{pku:torsion}
There is a finite set $P$ of split primes with
\[
\#H^1_{\mathcal L^P}(F,M)[\ell]\le\ell^{\#P-k-r_F+4}.
\]
\end{lemma}
\begin{proof}
\emph{Step 1: choice of $P$.} Let $A=M[\ell]$, a group of order $\ell$, and $A^*=\pkuHom(A,\mu_n)
=\pkuHom(A,\mu_\ell)$. The group $Y=H^1_{\mathcal L^*(A)}(F,A^*)$ (classes unramified outside $S_0$
with the dual conditions at $S_0$ and at the real places) is finite by (D3). For each $\xi\in Y$ whose
restriction to $G_\nu$ is nonzero, choose a split prime $p_\xi$ with $\xi(\phi_{p_\xi})\ne0$,
the $p_\xi$ distinct: the restriction $\xi|_{G_\nu}$ (well defined on the class, as $G_\nu$ acts
trivially on $A^*$) is a nonzero continuous homomorphism; its kernel $K$ is open in $G_\nu$ and normal
in $G_F$, since $\xi(gug^{-1})=g\xi(u)$ for $u\in G_\nu$ as in the proof of Lemma~\ref{pku:relation};
pick $u\in G_\nu$ with $\xi(u)\ne0$; by (D4) applied to the fixed field of $K$ and the class of $u$,
there are infinitely many primes $p\notin S_0$ of $F$ with a place $B_p$ over them and a Frobenius
element $\phi_p\in uK\subseteq G_\nu$; then $\xi(\phi_p)=\xi(u)\ne0$, and $\xi(\phi_p)$ does not depend
on the choice of the Frobenius element at $B_p$ ($\xi$ is unramified at $p$). Choose also a
normalised tame generator $\gamma_p$ at $B_p$. Let $P$ be the set of the $p_\xi$.

\emph{Step 2: the dual group for $P$.} If $\xi\in H^1_{(\mathcal L^P)^*}(F,A^*)$, then $\xi\in Y$
(at $p\in P$ the condition is $\xi_p=0$, in particular $\xi$ is unramified there) and
$\xi(\phi_p)=0$ for $p\in P$ ($\xi|_{D_p}$ is a coboundary for the trivial action of $D_p$ on $A^*$,
hence zero). By the choice of $P$, $\xi|_{G_\nu}=0$. So $H^1_{(\mathcal L^P)^*}(F,A^*)$ lies in the
kernel of the restriction $H^1(F,A^*)\to H^1(G_\nu,A^*)$, which is $H^1(\Delta,A^*)$, $\Delta=G_F/
G_\nu$ (inflation--restriction; $G_\nu$ acts trivially on $A^*$). Write $\Delta=\Delta_\ell\times
\Delta'$ with $\Delta_\ell$ its $\ell$-Sylow subgroup and $\#\Delta'$ prime to $\ell$
(Lemma~\ref{pku:basic}(ii)). By inflation--restriction for the normal subgroup $\Delta_\ell$, the sequence $0\to H^1(\Delta',(A^*)^{\Delta_\ell})\to H^1(\Delta,A^*)\to H^1(
\Delta_\ell,A^*)$ is exact and the first group is $0$; $\Delta_\ell$ acts trivially on $A^*$ (an
$\ell$-group acting on a group of order $\ell$ through $\mathbb F_\ell^\times$); so $H^1(\Delta,A^*)$
embeds in $\pkuHom(\Delta_\ell,\mathbb F_\ell)$, of order at most $\ell^2$. Thus
$\#H^1_{(\mathcal L^P)^*}(F,A^*)\le\ell^2$.

\emph{Step 3: the Selmer group of $M[\ell]$.} Apply~\eqref{pku:eq-GW} to $A$, $\Sigma_P$ and
$\mathcal L^P(A)$. The global factor is $\#H^0(F,A)/\allowbreak \#H^0(F,A^*)\le\ell$. By
Lemma~\ref{pku:local}, the local factor $\#L_v/\#H^0(F_v,A)$ is $\ell^{-1}$ at each of the $k$ primes
of $S_0$ not above $\ell$ and at each of the $r_F$ real places, at most $1$ at $v\mid\ell$, and
$\#H^1(F_p,A)/\#H^0(F_p,A)=\ell$ at $p\in P$. With Step 2,
\[
\#H^1_{\mathcal L^P}(F,M[\ell])\le\ell^2\cdot\ell\cdot\ell^{\#P-k-r_F}=\ell^{\#P-k-r_F+3}.
\]

\emph{Step 4: from $M[\ell]$ to $M$.} Let $j\colon M[\ell]\to M$ be the inclusion, $\pi\colon M\to
M/M[\ell]$ the projection and $\iota\colon M/M[\ell]\to M$ the injective map $\pi(a)\mapsto\ell a$,
with cokernel $M/\ell M$; all are $G_F$-maps, and $\iota\circ\pi$ is multiplication by $\ell$. Work in
$H^1(G_{\Sigma_P},-)$. Let $X=H^1_{\mathcal L^P}(F,M)[\ell]$. For $a\in X$, $\iota_*\pi_*(a)=\ell a=0$,
so $\pi_*(a)$ lies in the image of the connecting map $H^0(G_{\Sigma_P},M/\ell M)\to H^1(G_{\Sigma_P},
M/M[\ell])$; as $\#(M/\ell M)=\ell$, $\#\pi_*(X)\le\ell$. Let $X_0=X\cap\ker\pi_*$, so $\#X\le\ell\,\#X_0$.
An $a\in X_0$ is $j_*(w)$ for some $w\in H^1(G_{\Sigma_P},M[\ell])$. For every $v\in\Sigma_P$,
$j_*(w_v)=a_v\in L^P_v(M)$, i.e.\ the restriction of $j\circ w$ to $H_v$ is $0$; since $H_v$ acts
trivially and $j$ is injective, the restriction of $w$ to $H_v$ is $0$, i.e.\ $w_v\in L^P_v(M[\ell])$.
So $X_0\subseteq j_*\bigl(H^1_{\mathcal L^P}(F,M[\ell])\bigr)$ and, by Step 3,
$\#X\le\ell\cdot\ell^{\#P-k-r_F+3}$.
\end{proof}

\subsection{Proof of Proposition~\ref{pku:main}}\label{pku:ss-proof}

Let $P$ be as in Lemma~\ref{pku:torsion}, $B=\lambda_P(H^1_{\mathcal L^P}(F,M))\subseteq M^P$ and $R=
\pkuAnn(B)=\mu_P(\pkuSel)$ (Lemma~\ref{pku:image}). Let $X=M[\ell]^P\subseteq M^P$, and let $r\colon
(M^\vee)^P\to\pkuHom(X,\Q/\Z)$ be $r(f)(b)=\langle b,f\rangle$.

(1) $r(R)=\pkuHom(X/(X\cap B),\Q/\Z)$. Indeed $r(R)$ kills $X\cap B$; conversely a character of $X$
trivial on $X\cap B$ is a character of $(X+B)/B$, which extends to $M^P/B$ ($\Q/\Z$ is divisible),
and such an extension is $b\mapsto\langle b,f\rangle$ for some $f\in R$. Since $X\cap B=B[\ell]$,
$\#r(R)=\ell^{\#P}/\#B[\ell]$.

(2) $\#B[\ell]\le\#H^1_{\mathcal L^P}(F,M)[\ell]$: for a surjection $Y\to B$ of finite abelian groups,
$\#B[\ell]=\#(B/\ell B)\le\#(Y/\ell Y)=\#Y[\ell]$. With Lemma~\ref{pku:torsion},
$\#r(R)\ge\ell^{k+r_F-4}$.

(3) Let $\rho\colon\pkuSel\to\mathrm{Map}(G_\nu,\mathbb F_\ell)$, $\rho(\xi)=\pkured\kappa_\xi$, with
$\kappa_\xi$ the Kummer map of Lemma~\ref{pku:relation}; $\rho$ is additive, so its image is an
$\mathbb F_\ell$-subspace consisting of reductions of Kummer maps of exponent $e$ for $x$ off $S_0$.
The composite $r\circ\mu_P$ factors through $\rho$: for $b\in X$ write $b_p=(n/\ell)m_p$; then
$r(\mu_P(\xi))(b)=\sum_py_\xi(\phi_p)(b_p)=\sum_p\kappa_\xi(\phi_p)m_p/\ell$, which depends only on the
values of $\pkured\kappa_\xi$ at the $\phi_p\in G_\nu$. Hence $\#\rho(\pkuSel)\ge\#r(\mu_P(\pkuSel))=
\#r(R)\ge\ell^{k+r_F-4}$, i.e.\ $\dim\rho(\pkuSel)\ge k+r_F-4\ge k+1$ by Lemma~\ref{pku:basic}(iii).

(4) The span of the $f_i$ ($i\in I$) has dimension at most $k=\#I<\dim\rho(\pkuSel)$, so some
$\rho(\xi)$ is not in it. \qed

\begin{remark}\label{pku:rem-source}
The proof uses Lemma~\ref{pku:image} (Greenberg--Wiles and global duality, with $T=S_0$) to produce
relations, and Lemma~\ref{pku:torsion} (Greenberg--Wiles for $M[\ell]$, Chebotarev, and the
$\ell$-torsion argument of Step 4) for the count; the many real places of $F$ give room for the
constant losses.
\end{remark}

\subsection{Detecting points, and the triangle}\label{pku:ss-triangle}

Let $h=r^8$, $t=5r$, and write $N_0=q_1d$ (the integer $N$ of \S\ref{pod:sec}); recall
$N_0+t+2\le h$ (proof of Corollary~\ref{pod:kummer-b}). The integer called $N$ in \S\ref{pri:sec} is
$N_{\mathrm{ch}}$, and is not used under that name here.

\begin{proposition}[Detecting points]\label{pku:b}
Let $g_0$ be a topological generator of $\pgG$ modulo $\pgG\cap\ker\pgcyc$.
\begin{enumerate}
\item[(i)] For every prime $\ell$ there is a finite set $s\subseteq\pgG\cap\ker\pgcyc$ such that, for
all $\auO$, $\nu$ and $e$, a locally constant cocycle $z\colon\pgG\to(\auO/\ell^\nu)(-e)$ unramified
outside $\ell$ with $z(g_0)=0$ and $z|_s=0$ is zero.
\item[(ii)] If $\ell\le h^2$, $s$ can be chosen with $\#s\le h^{20}$, and then $(t+1)(\#s+1)\le h^{21}$.
\end{enumerate}
\end{proposition}
\begin{proof}
(i) is the second statement of Lemma~\ref{pod:detect} with $y=\max(N_0,\lceil\ell^{1/2}\rceil)$. (ii)
is the same with $y=h$, allowed as $h\ge N_0$ and $\ell\le h^2$; then $\#s\le h^{20}$ and, as $t+1\le
h-1$, $(t+1)(\#s+1)\le(h-1)(h^{20}+1)=h^{21}-h^{20}+h-1\le h^{21}$.
\end{proof}

For $\ell>h^2$ no bound on $\#s$ is claimed or needed.

\begin{proposition}[Proof of Proposition~\ref{pri:triangle}]\label{pku:triangle}
Proposition~\ref{pri:triangle} holds; in its proof (\S\ref{pri:ss-triangle}): (1) in ``Test points'', $s$ is the set
of Proposition~\ref{pku:b} for $g_0=a^{-1}$, and the inequality $(\#J+\#s+1)hh_0\le h^{22}h_0$ is
trivial if $h_0=0$ and follows from Proposition~\ref{pku:b}(ii) if $h_0\ne0$, because then $\ell\le
h^2$ by hypothesis; (2) in ``Induction'', Proposition~\ref{pku:main} is applied with $x$ the earlier
places of exponent $e$, the $f_i$ their reductions and $T_j=S_0(x)$; (3) the loss at the new place
is $2\kappa_j(\phi_j)=4\kappa_0(u)$; (4) in ``Property (iii)'' the local conditions are those of
Definition~\ref{pku:def-kummer}(a), at $S_0(x)=T_j$.
\end{proposition}
\begin{proof}
We give the steps of the proof in \S\ref{pri:ss-triangle} in full.

\emph{Choice of $a$, $h_0$.} If $\ell>h^2$ then $\ell\nmid N_0$ ($N_0\le h^2$) and $t<\ell-1$; an inertia
element at $\ell$ whose cyclotomic character reduces to a generator of $\mathbb F_\ell^\times$ fixes
$E\subseteq\Q(\mu_{N_0})$, and $\pgcyc^{\,e}-1$ is a unit at it for $0<e<t$; as every $\pgcyc(g)^e-1$ is
a multiple of $\pgcyc(a)^e-1$ ($a$ generates $\pgG$ modulo $\ker\pgcyc$, a procyclic quotient as
$\pgcyc\equiv1\bmod4$ on $\pgG$), $h_0=0$ works. Otherwise $h_0=\prival_\ell(N_{\mathrm{ch}})$ works,
by the Frobenius element at a chosen place above $p_i\ne\ell$. In both cases $h_0\le\prival_\ell(
N_{\mathrm{ch}})$ and ($\ell\le h^2$ or $h_0=0$).

\emph{Separation (ii).} With $u=\pgcyc(a)$: $u^e-u^k=u^k(u^{e-k}-1)$ for $k<e$ (and the negative for
$k>e$), so each factor divides $\ell^{h_0}$ and the product over the $t'-1$ indices $k\ne e$ divides
$\ell^{(t'-1)h_0}$. Also $4\mid u^e-1\mid\ell^{h_0}$ for $e>0$ ($u\equiv1\bmod4$; if $\ell$ is odd,
$4$ is a unit), so $4(u^e-1)^2\mid\ell^{3h_0}$. Both divide $\ell^{hh_0}$ as $t'+2\le h$. For the
triangle for $a^{-1}$ these are the conditions on $\pgcyc(a^{-1})^{-e}=u^e$.

\emph{Test points (i).} $J=\{(e,g):0<e<t',\ g\in s\}$, ordered lexicographically with $e$ first,
$\prity(e,g)=e$; $J$ and $\prity$ do not depend on $\nu$, since $s$ does not. $\#J+\#s+1\le
t'\#s+1\le(t+1)(\#s+1)$. If $h_0=0$ both sides of $(\#J+\#s+1)hh_0\le h^{22}h_0$ vanish; if $h_0\ne0$,
then $\ell\le h^2$ and Proposition~\ref{pku:b}(ii) gives $(t+1)(\#s+1)h\le h^{22}$. Detection with
$a^{-1}$ is Proposition~\ref{pku:b}(i)--(ii) for $g_0=a^{-1}$.

\emph{Induction.} Let the places $w_i$, $i<j$, with Kummer maps $\kappa_i$ of exponent $\prity(i)$
be chosen, above distinct primes of $F$, with the reductions of the $\kappa_i$ of exponent
$e=\prity(j)$ linearly independent. Proposition~\ref{pku:main}, applied to the family $x$ of the
$w_i$ with $\prity(i)=e$ and to their reductions, gives a Kummer map $\kappa_0$ of exponent $e$ for
$x$ off $T_j=S_0(x)$ whose reduction is outside their span. The Kummer map attached to the new
place $w_j$ chosen below is $\kappa_j:=\kappa_0$, of exponent $e$ for $x$ off $T_j=S_0(x)$; this is
what Property (iii) uses. By Remark~\ref{pku:rem-def}(2) all these
maps are additive on $G_\nu$, negated by $g\mapsto g^c$, and conjugation by $g\in\pgG$ multiplies a
map of exponent $d$ by $\pgcyc(g)^{d+1}$; so conjugation by $a$ followed by the power
$\pgcyc(a)^{-1}\bmod\ell^\nu$ multiplies it by $\pgcyc(a)^d$. The earlier maps of exponent $e$ and
$\kappa_0$ have independent reductions, so together they map $G_\nu$ onto $(\Z/\ell^\nu)^{m}$
(Nakayama); take $g'\in G_\nu$ where $\kappa_0$ is $1$ and the earlier maps of exponent $e$ vanish,
and apply $\prod_{k\ne e,\,k<t'}(\text{conj}_a\text{ then power}-\pgcyc(a)^k)$: this gives $u\in
G_\nu$ with $\kappa_i(u)=0$ for all $i<j$ (for $\prity(i)\ne e$ the factor $k=\prity(i)$ vanishes) and
$\kappa_0(u)=\prod_{k\ne e}(\pgcyc(a)^e-\pgcyc(a)^k)$. With $g=g_j$ the test point of $j$ and
$\sigma=\delta^{\ell^\nu}gg^cu(u^c)^{-1}$ ($\delta$ given by the hypothesis on $\mathcal W$ for $\nu$
and $g''=gg^cu(u^c)^{-1}$), every
Kummer map of exponent $d$ takes the value $2\kappa(u)$ near $\sigma$, and every even class of
exponent $e$ unramified outside $\ell$ the value $2z(g)$, as in \S\ref{pri:ss-triangle}. Let $U\subseteq G_\nu$ be an open neighbourhood of $\sigma$ (this is possible: $\sigma\in G_\nu$,
since $\delta$, $g$, $g^c$ have trivial cyclotomic character modulo $\ell^\nu$ and $u\in G_\nu$, and
$G_\nu$ is open) on which the
$\kappa_i$ ($i<j$), $\kappa_0$ and all even classes of exponent $e$ unramified outside $\ell$ are
constant (after composing with additive maps to $\Z/\ell^\nu$, the latter are finitely many cocycles,
all factoring through one finite quotient of $\pgG$); by the hypothesis, $\mathcal W$ contains almost
all Frobenius data with Frobenius in a nonempty open subset of $U$, and Chebotarev gives infinitely
many of them. Choose $w_j
\in\mathcal W$ of degree one with Frobenius $\phi_j\in U$, not above $\ell$ and above a prime
of $F$ not below an earlier place; this prime is then outside every $T_i$, $i\le j$. Then
$\kappa_i(\phi_j)=2\kappa_i(u)=0$ for $i<j$, $z(\phi_j)=2z(g)$, $\phi_j\in G_\nu$, and the loss is
\[
2\kappa_j(\phi_j)=4\kappa_0(u)=4\prod_{k\ne e}(\pgcyc(a)^e-\pgcyc(a)^k)\ \Big|\ \ell^{(3+(t'-1))h_0}
\ \Big|\ \ell^{hh_0}
\]
in $\Z/\ell^\nu$ (using $4\mid\ell^{3h_0}$ if $j$ exists). The independence of the reductions of the
maps of exponent $e$ up to $j$ holds by construction.

\emph{Property (iii).} Let $z$ be an even class of exponent $e=\prity(j)$ for the $w_i$ with
$z(\tau_i)=0$ for $i<j$ and $z(\phi_i)=0$ for $i<j$ of exponent $e$. Then $z$ is unramified at $w_i$
and $\overline w_i$ for $i<j$. The values of $z$ lie in $\auO/\ell^\nu$; the additive maps
$f\colon\auO/\ell^\nu\to\Z/\ell^\nu$ separate points and $f\circ z$ is again an even class with values in
$(\Z/\ell^\nu)(-e)$ with the same vanishing properties, so we may and do assume $\auO=\Z_\ell$. Its
extension $z^+(gc)=z(g)$ to $G_F$ is a cocycle for
$(\pkucp)^{-e}$, and $2z^+$ is a test cocycle for $x$ (the earlier places of exponent $e$) and $P=
\{$primes below $w_i$, $i\ge j\}$ in the sense of Definition~\ref{pku:def-kummer}(a): it vanishes at
complex conjugations ($2z^+(c')=z^+(c'^2)=0$, as $\pkucp(c')=1$); on $D_{B_v}\cap\ker\pkucp$ for
$v\mid\ell$ (for such $g$, $2z^+(g)=z^+(g^2)$ and $g^2\in\pgG\cap\ker\pgcyc$ lies in a decomposition
group above $\ell$, where $z$ vanishes); on $D_{A_i}$ for the earlier places of exponent $e$ ($z$
vanishes on $I_{A_i}$ and at $\phi_i$, which generate it topologically); and it is unramified outside
$S_0(x)\cup P$. The primes of $P$ are split primes not in $T_j$ (distinct, by the choice of the
$w_i$), with Frobenius elements $\phi_i\in G_\nu$ (the $\phi_j$ of Definition~\ref{pri:def-tri} are
Frobenius elements in the arithmetic sense of \S\ref{pku:ss-statement}, as chosen above) and normalised tame
generators $\gamma_i$ at the places $A_i$ of the $w_i$. The reciprocity law~\eqref{pku:eq-recip} of $\kappa_j$ gives $\sum_{i\ge j}
\kappa_j(\phi_i)\,2z(\gamma_i)=0$ (up to a common sign); $\kappa_j(\phi_i)=0$ for $i>j$, so
$2\kappa_j(\phi_j)z(\gamma_j)=0$. The value of $z$ at any element of $I_{A_j}$, in particular at
$\tau_j$, is a multiple of $z(\gamma_j)$; hence every multiple of $2\kappa_j(\phi_j)$ in $\Z/\ell^\nu$,
in particular $\ell^{hh_0}$, kills $z(\tau_j)$.
\end{proof}

\begin{remark}\label{pku:rem-tri-source}
The triangle constants are $b=h\,h_0$, $c=h^{22}h_0$ with $h=r^8$, as in
Proposition~\ref{pri:triangle}; the size count $(\#J+\#s+1)\le(t+2)h^{20}\le h^{21}$ is used only
when $\ell\le h^2$, the case $h_0=0$ being trivial.
\end{remark}


%% ===== proofs/C-aux.tex =====
% proofs/C-aux.tex
% Auxiliary statements of Theorem C used in Section pri: Amitsur's formula (used in
% Lemma pri:amitsur), the chain lemma pri:chain, the inputs on forms pri:forms, the block-diagonal
% lattices pri:block. Labels of this file start with pax:.
\providecommand{\auO}{\mathcal O}
\providecommand{\auQlb}{\overline{\Q}_\ell}
\providecommand{\auGL}{\mathrm{GL}}
\providecommand{\auF}{\mathbb F}
\providecommand{\auU}{\mathrm U}
\providecommand{\auNm}{\operatorname{Nm}}
\providecommand{\pgZb}{\overline{\Z}_\ell}
\providecommand{\pgG}{G_E}
\providecommand{\pgtr}{\operatorname{tr}}
\providecommand{\pgFrob}{\operatorname{Frob}}
\providecommand{\priX}{\mathcal X}
\providecommand{\priT}{\mathbb T}
\providecommand{\pgNch}{N_{\mathrm{ch}}}
\providecommand{\pgaWD}{\operatorname{WD}}
\providecommand{\auf}[1]{\texttt{#1}}
% macros of this file (prefix \pax)
\providecommand{\paxch}{\operatorname{ch}}
\providecommand{\paxL}{\mathcal L}

\section{Auxiliary statements for Ribet's method}\label{pax:sec}

This section proves statements that \S\ref{pri:ss-inputs} uses without proof:
Amitsur's formula, on which Lemma~\ref{pri:amitsur} rests (\S\ref{pax:ss-amitsur}); the chain
lemma, Lemma~\ref{pri:chain} (\S\ref{pax:ss-chain}); the inputs on forms, Lemma~\ref{pri:forms}
(\S\ref{pax:ss-forms}); and the block-diagonal lattices, Lemma~\ref{pri:block}
(\S\ref{pax:ss-block}); the integrality of the pairing at a free level (\S\ref{pax:ss-pairing}); the
dictionary at an auxiliary place, Proposition~\ref{pri:dict} (\S\ref{pax:ss-dict}); the base level of a label,
Proposition~\ref{pri:BL} (\S\ref{pax:ss-BL}); the sign hypothesis of Proposition~\ref{pga:G123}
(\S\ref{pax:ss-sign}); and the convention check of \S\ref{pga:ss-M} (\S\ref{pax:ss-conv}). Conventions are those of \S\ref{pri:sec}.

\subsection{Amitsur's formula}\label{pax:ss-amitsur}

Let $x_1,\dots,x_k$ be commuting indeterminates. A \emph{word} is a finite sequence $u=u_1\cdots
u_m$ in $\{1,\dots,k\}$, of length $|u|=m\ge1$; for matrices $A_1,\dots,A_k$ of the same size put
$A_u=A_{u_1}\cdots A_{u_m}$ and $x^u=x_{u_1}\cdots x_{u_m}$. A word is \emph{primitive} if it is not
a power $v^d$, $d\ge2$, of a shorter word; the cyclic rotations of a primitive word $w$ are $|w|$
distinct words; a \emph{Lyndon word} is a primitive word which is lexicographically smallest among
its rotations. Every word is uniquely $w'^{\,d}$ with $w'$ primitive, and every primitive word is a
rotation of exactly one Lyndon word.

\begin{proposition}[Amitsur's formula]\label{pax:amitsur}
Let $\mathsf A$ be a commutative ring and $A_1,\dots,A_k\in M_\sigma(\mathsf A)$. In the ring of
formal power series $\mathsf A[[x_1,\dots,x_k]]$,
\begin{equation}\label{pax:eq-amitsur}
\det\Bigl(1-\sum_{i=1}^kx_iA_i\Bigr)=\prod_{w\ \mathrm{Lyndon}}\det\bigl(1-x^wA_w\bigr),
\end{equation}
the product being convergent (only finitely many factors differ from $1$ modulo each power of the
ideal $(x_1,\dots,x_k)$). Consequently, for every $j\ge0$ the coefficient of $T^j$ in $\det(1-T\sum_i
A_i)\in\mathsf A[T]$ is a polynomial with integer coefficients, independent of $\mathsf A$ and the
$A_i$, in the coefficients of the characteristic polynomials of the products $A_w$ with $|w|\le j$.
\end{proposition}

\begin{proof}
It suffices to treat the universal case: $\mathsf A_0=\Z[a_{i,pq}]$, the polynomial ring in the
entries of $k$ generic matrices; the general case follows by the ring map $\mathsf A_0\to\mathsf A$
sending the generic entries to those of the $A_i$, applied coefficientwise. Since $\mathsf A_0$ is a
domain of characteristic $0$, it suffices to prove \eqref{pax:eq-amitsur} in $\mathsf A_0\otimes
\Q[[x]]$, where both sides have constant term $1$ and the logarithm $\log(1+f)=\sum_{m\ge1}
(-1)^{m-1}f^m/m$ ($f$ without constant term) is injective, with $\log(fg)=\log f+\log g$. For a matrix
$M$ with entries in the ideal $(x)$, $\log\det(1-M)=\operatorname{tr}\log(1-M)=-\sum_{m\ge1}\pgtr
(M^m)/m$ (replace $M$ by $YM$ with an auxiliary variable $Y$: both sides vanish at $Y=0$, and their
derivatives in $Y$ agree, by Jacobi's formula $\frac{d}{dY}\det(1-YM)=-\det(1-YM)\pgtr\bigl(M(1-YM)^{-1}
\bigr)$ and $\pgtr(M(1-YM)^{-1})=\sum_{m\ge1}Y^{m-1}\pgtr M^m$). With $M=\sum_ix_iA_i$ we get
\[
\log\det\Bigl(1-\sum_ix_iA_i\Bigr)=-\sum_{u}\frac1{|u|}\,x^u\,\pgtr A_u,
\]
the sum over all words. Group the words by their primitive root and cyclic rotation: every word is
$u=w_1^{\,d}$ with $w_1$ a rotation of a unique Lyndon word $w$, and $x^u=x^{w}{}^{d}$, $\pgtr A_u=
\pgtr\bigl(A_{w_1}^{\,d}\bigr)=\pgtr\bigl(A_w^{\,d}\bigr)$ by cyclicity of the trace (a rotation of
$w_1^{\,d}$ by one letter is $w_2^{\,d}$ with $w_2$ the rotation of $w_1$). For fixed $w$ and $d$ there
are exactly $|w|$ such words $u$ (the rotations of $w$ are distinct), each of length $d|w|$. Hence
\[
-\sum_u\frac{x^u\pgtr A_u}{|u|}=-\sum_{w\ \mathrm{Lyndon}}\ \sum_{d\ge1}|w|\cdot\frac{(x^w)^d\,
\pgtr(A_w^{\,d})}{d\,|w|}=\sum_{w\ \mathrm{Lyndon}}\log\det(1-x^wA_w),
\]
by the same identity applied to the single matrix $x^wA_w$. Rearrangements are legitimate because for
each monomial only finitely many terms contribute. Exponentiating gives \eqref{pax:eq-amitsur}.

For the consequence, substitute $x_1=\dots=x_k=T$ (a continuous ring map $\mathsf A[[x]]\to\mathsf
A[[T]]$): $\det(1-T\sum A_i)=\prod_w\det(1-T^{|w|}A_w)$, and the coefficient of $T^j$ involves only the
factors with $|w|\le j$, whose coefficients are, up to sign, the coefficients of the characteristic
polynomials of the $A_w$ (as $\det(1-YA)=Y^\sigma\det(Y^{-1}-A)$).
\end{proof}

\begin{remark}\label{pax:rem-amitsur}
(1) This is the formula of Amitsur~\cite{Amitsur}; the proof above is the standard one by
logarithms (it is the proof of Reutenauer and Sch\"utzenberger~\cite{ReutenauerSchutzenberger}). It is self-contained.
(2) \emph{Use in Lemma~\ref{pri:amitsur}.} There $\rho(P)=\sum_qE_q$, and the characteristic polynomial
of $\rho(P)\in M_{m\sigma}(\mathsf L)$ is $X^{m\sigma}\det(1-X^{-1}\rho(P))$. By
Proposition~\ref{pax:amitsur} with $\mathsf A=\mathsf L$, its coefficients are integral polynomials in
the coefficients of the characteristic polynomials of the products $E_{q_1}\cdots E_{q_k}$. Writing
$E_q=\epsilon_q\pi_q$ with $\epsilon_q\in M_{m\sigma\times\sigma}$ the $q$-th block column of $1$ and
$\pi_q\in M_{\sigma\times m\sigma}$ the $q$-th block row of $\rho(P)$, so that $\pi_q\epsilon_{q'}=
\rho(P_{qq'})$, the product is $\epsilon_{q_1}\rho(P_{q_1q_2})\cdots\rho(P_{q_{k-1}q_k})\pi_{q_k}$, and
since $\det(X-AB)=X^{a-b}\det(X-BA)$ for $A$ of size $a\times b$, $B$ of size $b\times a$, its
characteristic polynomial is a power of $X$ times that of $\rho(P_{q_1q_2}\cdots P_{q_{k-1}q_k}
P_{q_kq_1})$, whose coefficients lie in the image of $\mathsf R$ by hypothesis. As the image of
$\mathsf R$ is a subring, the proof of Lemma~\ref{pri:amitsur} is complete.
\end{remark}

\subsection{The chain lemma}\label{pax:ss-chain}

\emph{Why Proposition~\ref{pga:L} does not suffice.} Proposition~\ref{pga:L} computes the Hecke
polynomial of a theta lift; its roots contain the chain $\alpha_wq^l$, $l<t$. Lemma~\ref{pri:chain}
concerns a \emph{weak lift}, an eigenform whose normalised Galois representation has the eigenvalue
$1$ everywhere; such a form is not known to be a theta lift (this is the point of the false lifts in
Lemma~\ref{pri:weak}). Only the first root of the chain comes from Galois representations; the others
come from the unitarity of the local component, as follows.

Let $F$ be a non-archimedean local field with residue field of order $q$, uniformiser $\varpi$,
$K=\auGL_N(\auO_F)$, $\nu=|\det|$, and let $T_i$ ($0\le i\le N$) be the Hecke operator of $K\,
\mathrm{diag}(\varpi1_i,1_{N-i})K$ as in (N3) of \S\ref{pl:ss-norm}. For an irreducible unramified
representation $\pi$ of $\auGL_N(F)$ with Satake parameters $\alpha_1,\dots,\alpha_N$ (that is, $\pi$
is the spherical subquotient of the normalised induction $\chi_1\times\dots\times\chi_N$ from
unramified characters with $\chi_k(\varpi)=\alpha_k$), the eigenvalues $t_i$ of the $T_i$ on
$\pi^K$ satisfy
\begin{equation}\label{pax:eq-satake}
H_\pi(X):=\sum_{i=0}^N(-1)^iq^{i(i-1)/2}t_iX^{N-i}=\prod_{k=1}^N\bigl(X-q^{(N-1)/2}\alpha_k\bigr),
\end{equation}
which is Satake's formula for $\auGL_N$, used in the same normalisation in the proof of
Proposition~\ref{pga:G123}.

\begin{lemma}[Chains in unitary spherical representations]\label{pax:tadic}
Let $\pi$ be an irreducible unitary unramified representation of $\auGL_N(F)$, $c$ a root of
$H_\pi$ and $k\ge1$ an integer with $|c|^2q^{k}=q^N$. Then $c,cq,\dots,cq^{k-1}$ are roots of
$H_\pi$, and $\prod_{i<k}(X-q^i)$ divides $H_\pi(cX)$.
\end{lemma}

\begin{proof}
This is Corollary~\ref{pay:roots}; see \S\ref{pay:ss-cons}, where it is proved without the
classification of the unitary dual. We also give the deduction from Tadi\'c's classification.

\emph{Input} (Tadi\'c's classification~\cite{TadicENS}, Theorem~D; for the
unramified case see the discussion below). Every irreducible unitary representation of $\auGL_N(F)$
is an irreducible product $\tau_1\times\dots\times\tau_s$ of representations of the following two
kinds: $u(\delta,m)$, the Langlands quotient of $\delta\nu^{(m-1)/2}\times\dots\times\delta\nu^{-(m-1)/2}$
for an irreducible unitary square-integrable $\delta$; and $u(\delta,m)\nu^a\times u(\delta,m)\nu^{-a}$
with $0<a<\frac12$.

\emph{Unramified constituents.} If $\pi=\tau_1\times\dots\times\tau_s$ (irreducible, normalised
parabolic induction from a standard Levi $M$) is unramified, so is each $\tau_i$: by the Iwasawa
decomposition $\auGL_N(F)=P K$, restriction of functions to $K$ identifies the $K$-fixed vectors of
$\operatorname{Ind}_P(\tau_1\otimes\dots\otimes\tau_s)$ with the $K\cap M$-fixed vectors of
$\tau_1\otimes\dots\otimes\tau_s$. Next, $u(\delta,m)$ with $\delta$ a representation of $\auGL_d(F)$ is
unramified only if $d=1$ and $\delta$ is an unramified unitary character $\chi$: an unramified
irreducible representation has supercuspidal support consisting of unramified characters of $F^\times$,
which forces $\delta$ to be a twist $\mathrm{St}_d\otimes\chi$ with $\chi$ unramified; and the
unramified irreducible representation with a given cuspidal support $\{\chi_1,\dots,\chi_N\}$ is the
Langlands quotient of $\chi_{1}\times\dots\times\chi_{N}$ ordered by decreasing exponents, whose
Langlands data consist of characters of $\auGL_1$ only, whereas the Langlands data of $u(\mathrm{St}_d
\otimes\chi,m)$ are $\mathrm{St}_d\chi\nu^{(m-1)/2},\dots$; by the uniqueness in the Langlands
classification, $d=1$. Finally $u(\chi,m)=\chi\circ\det$ on $\auGL_m(F)$, a subquotient of $\chi\nu^{(m-1)/2}\times
\dots\times\chi\nu^{-(m-1)/2}$. Since $\nu^s(\varpi)=q^{-s}$, the Satake parameters of $\pi$ are the
union, with multiplicity, of blocks of the two forms
\[
\text{(i)}\ \ \zeta\,q^{\,b-(m-1)/2}\ (0\le b<m),\qquad\text{(ii)}\ \ \zeta\,q^{\,b-(m-1)/2\pm a}\
(0\le b<m),\quad|\zeta|=1,\ 0<a<\tfrac12 .
\]

\emph{The chain.} Write $c=q^{(N-1)/2}\alpha$ with $\alpha$ a Satake parameter. Then $|\alpha|^2=
q^{1-k}$, i.e.\ $|\alpha|=q^{-(k-1)/2}$, with $-(k-1)/2\in\frac12\Z$. If $\alpha$ lay in a block of
type (ii), $\log_q|\alpha|=b-(m-1)/2\pm a$ would lie in $\frac12\Z\pm a$, which is disjoint from
$\frac12\Z$ as $0<a<\frac12$. So $\alpha=\zeta q^{b_0-(m-1)/2}$ lies in a block of type (i), with
$b_0-(m-1)/2=-(k-1)/2$. Then $0\le b_0=(m-k)/2$, so $k\le m$ and $k\equiv m\pmod2$, and for $0\le i<k$,
$\alpha q^i=\zeta q^{(b_0+i)-(m-1)/2}$ with $0\le b_0+i\le b_0+k-1=(m+k)/2-1\le m-1$. Hence
$\alpha q^i$ is a Satake parameter of the same block, and $cq^i$ is a root of $H_\pi$ by
\eqref{pax:eq-satake}. These $k$ roots are distinct ($q>1$, $c\ne0$), so $\prod_{i<k}(X-cq^i)$ divides
$H_\pi(X)$, and $\prod_{i<k}(X-q^i)$ divides $H_\pi(cX)$.
\end{proof}

\begin{remark}\label{pax:rem-tadic}
Lemma~\ref{pax:tadic} can be proved directly, without the classification: for a nonzero spherical
vector $x$ with Hecke polynomial $P$, a root $c$ with $qc$ not a root satisfies $|c|^2>q^{N-2}$ (an
identity obtained by pairing two explicit translates of $x$, computed by counting hyperplanes, at
pro-$p$ Iwahori level), and the roots are symmetric under $y\mapsto q^{N-1}/\bar y$ (the adjoint of
$T_i$ is $T_{N-i}$ up to the centre); an induction along the chain then gives the lemma. This
direct proof is written in \S\ref{pay:sec}; the paper uses this direct proof.
\end{remark}

\begin{proof}[Proof of Lemma~\ref{pri:chain}]
Let $y$ be a nonzero weak lift in base level (the forms of weight $\lambda$ of $O$ on $\auU(V)$, $n=
\operatorname{rk}V$, $r'=n_W$, $t'=n-r'\ge1$), $w\nmid\ell$ split and off the level, $q=\auNm w$,
$v$ the place of $E^+$ below $w$. \emph{Unitarity.} The space of all algebraic modular forms of weight
$\lambda$ (all levels) is a smooth admissible representation of $G_f$ by right translation, and the
inner product of \S\ref{ss:au-spaces}, renormalised by the volume of the level, is $G_f$-invariant and
positive definite on it. Hence the $G_f$-representation generated by $y$ is a finite direct sum of
irreducible unitary representations $\pi=\otimes'_u\pi_u$; as $y$ is an eigenform of the $T_{w,i}$,
its projection to a constituent $\pi$ is nonzero and has the same eigenvalues, so $\pi_v\simeq\pi_w$
(under $\auU(V)(E^+_v)\simeq\auGL_n(E_w)$) is irreducible, unitary and unramified, and $H_w(y;X)=
H_{\pi_w}(X)$ in the notation~\eqref{pax:eq-satake}.

\emph{The first root.} By (G1) (Proposition~\ref{pga:G123}), $\det(X-S_y(\pgFrob_w))=\iota H_w(y;X)$.
The normalised representation $S_y\otimes\eta_O$ has the eigenvalue $1$ at $\pgFrob_w$, so $S_y(
\pgFrob_w)$ has the eigenvalue $\eta_O(\pgFrob_w)^{-1}=\iota X_0$, and $X_0$ is a root of $H_w(y;X)$.

\emph{Its absolute value.} By the normalisation of $\eta_O$ (\S\ref{pga:ss-L}: $\eta_O(\pgFrob_w)=
\alpha_w^{-1}$), $X_0=\alpha_w=\chi_{W,w}^{(r')}=\chi_W(\varpi_w/\overline\varpi_w)\upsilon(\varpi_w
)^{-r'}q^{r'/2}$. The character $\chi_W$ of the compact group $\auU(1)(E^+)\backslash\auU(1)(\mathbb A
_{E^+})$ is unitary; the splitting character $\upsilon$ restricts to a character of finite order of
$\mathbb A_{E^+}^\times$, and writing $\upsilon=\upsilon_0|\cdot|_{\mathbb A_E}^s$ with $\upsilon_0$
unitary, the restriction of $|\cdot|^s_{\mathbb A_E}$ to $\mathbb A_{E^+}^\times$ is $|\cdot|^{2s}$,
of finite order only if $s=0$; so $\upsilon$ is unitary. Hence $|X_0|^2=q^{r'}$, i.e.\
$|X_0|^2q^{t'}=q^n$.

Lemma~\ref{pax:tadic} with $N=n$, $c=X_0$, $k=t'$ gives that $\prod_{i<t'}(X-q^i)$ divides $H_w(y;
X_0X)$, which is Lemma~\ref{pri:chain} (for $X_0$ the value at a geometric Frobenius of $\iota^{-1}
\eta_O^{-1}$). The statement holds at every split $w\nmid\ell$ off the level at which (G1) holds,
i.e.\ (Proposition~\ref{pga:G123}) at every such $w$ at which $y$ is spherical.
\end{proof}

\begin{remark}\label{pax:rem-chain-sign}
Lemma~\ref{pri:chain} uses (G1) (Proposition~\ref{pga:G123}, which holds without any sign
hypothesis by Proposition~\ref{pax:sign}); nothing else from \S\ref{pga:sec} is used. If in the normalisation of
\S\ref{ss:au-spaces} the operator $T_{w,i}$ is defined with $K\,d_i^{-1}K$ or with the action through
$\tau\circ c$ (see (N2)), the polynomial $H_w$ is replaced by the one with Satake parameters
$\bar\alpha_k^{-1}$ (the contragredient, equal to the complex conjugate for unitary $\pi$); the set of
Satake parameters of a unitary unramified representation is stable under $\alpha\mapsto\bar\alpha^{-1}$
(the blocks (i), (ii) are), and the argument is unchanged.
\end{remark}

\subsection{Inputs on forms}\label{pax:ss-forms}

Let $J_O$ be the base level of Definition~\ref{prm:def-level}, with its level character $\chi_O$. At a
chosen place $J_{O,w}$ is not congruent to $1$ modulo $w$ (its $\auGL_{t'}$-block is arbitrary), so
Lemma~\ref{au:free} does not apply to $J_O$ itself. Put
\[
J_0=J_O\cap\{k:\ k_w\equiv1\bmod w\ \text{at every chosen }w\},\qquad H=J_O/J_0\simeq\prod_{w\ \mathrm{chosen}}
J_{O,w}/(1+\varpi_wM_n(\auO_w)),
\]
a normal subgroup ($1+\varpi_wM_n(\auO_w)$ is normal in $\auGL_n(\auO_w)$) and a finite quotient, $H$
being a subgroup of $\prod_{w\ \mathrm{chosen}}\auGL_n(\auF_w)$. Since $J_{O,v}\subseteq K_v$ for every $v$ (Definition~\ref{prm:def-level}), $J_O\subseteq K^\circ$, so
 $J_0$ consists of elements congruent to $1$ modulo every chosen place and is free by
Lemma~\ref{au:free}; we call $J_0$ the \emph{free sublevel} and write $\chi=\chi_O|_{J_0}$, a character of
finite order, trivial at the chosen places. Let $w_1,\dots,w_s$ be
places of $E$, split over $E^+$, at which $J_0$ is hyperspecial, $J_{0,w_j}=\auGL_n(\auO_{w_j})$, and
$\chi$ trivial. Let $J_1\subseteq J_0$ be the subgroup of the elements congruent to $1$ modulo each
$w_j$ (principal level at the $w_j$); it is normal, and $J_0/J_1\simeq\Gamma=\prod_j\auGL_n(
\auF_{w_j})$. Then $\Gamma$ acts on $\mathcal A_\lambda(J_1,\chi)$ by $(k\varphi)(g)=\varphi(g\tilde k)$,
$\tilde k\in\prod_j\auGL_n(\auO_{w_j})\subseteq J_0$ a lift of $k$ with components $1$ outside the $w_j$ (independent of such a lift since $\chi$ is trivial on the components at
the $w_j$ and $J_1$ is normal). An \emph{integral operator} is a linear endomorphism preserving, for
every $\iota$, the forms integral at $\iota$.

The hypothesis $n\le2n_m$ of Lemma~\ref{au:free} holds: $n=r'+t'\le r+t=6r$, while
$n_m=r+mt+1$ with $m=2\varrho+2\ge2$, so $2n_m>2mt\ge20r>n$.

\begin{lemma}[Freeness of the principal level]\label{pax:free}
Let $\mathcal C_0=\Gamma_{\mathrm{gl}}\backslash G_f/J_0$ and $\mathcal C_1=
\Gamma_{\mathrm{gl}}\backslash G_f/J_1$, $\Gamma_{\mathrm{gl}}=\auU(V)(E^+)$. Choose for each class of
$\mathcal C_0$ a representative $g$. Then $(g,k)\mapsto\Gamma_{\mathrm{gl}}g\tilde kJ_1$ is a bijection
from $\{\text{representatives}\}\times\Gamma$ onto $\mathcal C_1$, the action of $h\in\Gamma$ on
$\mathcal C_1$ by right translation is $(g,k)\mapsto(g,kh)$, and $\varphi\mapsto\bigl((g,k)\mapsto
\varphi(g\tilde k)\bigr)$ identifies $\mathcal A_\lambda(J_1,\chi)$, $\Gamma$-equivariantly, with
$\operatorname{Map}(\{\text{representatives}\}\times\Gamma,W_\lambda)$, $\Gamma$ acting by right
translation on the second variable. A form is integral at $\iota$ if and only if its values
$\varphi(g\tilde k)$ have gauge at most $1$ at $(\iota,g)$.
\end{lemma}

\begin{proof}
Surjectivity is clear. If $\Gamma_{\mathrm{gl}}g\tilde kJ_1=\Gamma_{\mathrm{gl}}g\tilde k'J_1$, then
$\gamma g\tilde k=g\tilde k'j$ with $\gamma\in\Gamma_{\mathrm{gl}}$, $j\in J_1$, so $\gamma\in\Gamma_{
\mathrm{gl}}\cap gJ_0g^{-1}=1$ by Lemma~\ref{au:free}, hence $\tilde k'^{-1}\tilde k\in J_1$ and $k=k'$.
The same lemma shows that the defining relation $\varphi(\gamma xj)=\chi(j)\gamma\varphi(x)$ imposes
no condition on the values at the points $g\tilde k$ (each point of $\mathcal C_1$ is represented once
and has trivial stabiliser), which gives the identification. For the gauge: $\tilde k$ may be chosen
with components $1$ outside the $w_j$ and in $\auGL_n(\auO_{w_j})$ at the $w_j$, so it is integral
with integral inverse above $\ell$, and (G2) of \S\ref{pl:ss-norm} gives that the gauge at
$(\iota,g\tilde k)$ equals that at $(\iota,g)$.
\end{proof}

\emph{Stabilisers of the base level.} For $g\in G_f$ the stabiliser $\mathrm{Stab}_g=\Gamma_{\mathrm{gl}}
\cap gJ_Og^{-1}$ maps to $H$ by $\gamma\mapsto g^{-1}\gamma g\bmod J_0$, with kernel $\Gamma_{\mathrm{gl}}
\cap gJ_0g^{-1}=1$ (Lemma~\ref{au:free}). So $\#\mathrm{Stab}_g$ divides $\#H$, and $\#H$ divides
$\prod_{w\ \mathrm{chosen}}\#\auGL_n(\auF_w)$, which divides $N=\pgNch=\prod_{w\ \mathrm{chosen}}
\#\auGL_{rt}(\auF_w)$ since $n=r'+t'\le r+t=6r\le5r^2=rt$ ($r'\le r$, $t'\le t=5r$, $r\ge2$) and $\auGL_n(\auF_w)$ embeds in $\auGL_{rt}(\auF_w)$
(Lagrange). We take $c_0=\#H$: a common multiple of the orders of the stabilisers of the base level,
dividing $N$.

\begin{proof}[Proof of Lemma~\ref{pri:forms}]
(a) Let $J_O^{(1)}=J_O\cap J_1$ be the base level with principal level at the $w_j$; the forms of
principal level at the $w_j$ are $\mathcal A_\lambda(J_O^{(1)},\chi_O)\subseteq\mathcal A_\lambda(J_1,\chi)$.
We first build an integral operator $\Psi_1$ on $\mathcal A_\lambda(J_1,\chi)$ with $\sum_{k\in\Gamma'}
k^{-1}\Psi_1k=1$, and then average over $H$. Let $\Gamma'\subseteq\Gamma$
be a subgroup (in particular a subgroup of a factor, or $\Gamma$ itself), and $S\subseteq\Gamma$ a set
of representatives of the left cosets $k\Gamma'$. In the model of Lemma~\ref{pax:free} let $\Psi$ be
multiplication by the indicator function of $\{\text{representatives}\}\times S$:
$(\Psi\varphi)(g,k)=\mathbf 1_S(k)\varphi(g,k)$. It is integral by the last assertion of
Lemma~\ref{pax:free}. For $h\in\Gamma'$, $(h^{-1}\Psi h\varphi)(g,k)=(\Psi h\varphi)(g,kh^{-1})=
\mathbf 1_S(kh^{-1})\varphi(g,k)$, so
\[
\Bigl(\sum_{h\in\Gamma'}h^{-1}\Psi h\,\varphi\Bigr)(g,k)=\#\{h\in\Gamma':kh^{-1}\in S\}\cdot
\varphi(g,k)=\#(S\cap k\Gamma')\,\varphi(g,k)=\varphi(g,k).
\]
Call this operator $\Psi_1$ (it is the $\Psi$ just defined, on $\mathcal A_\lambda(J_1,\chi)$). Now let
$e_H=\sum_{h\in H}R(\tilde h)$, $(R(\tilde h)\varphi)(g)=\varphi(g\tilde h)$, with $\tilde h\in J_O$ a lift
of $h$ with components $1$ outside the chosen places (well defined on $\mathcal A_\lambda(J_1,\chi)$ since
$J_0$ is normal in $J_O$ and $\chi_O$ is trivial at the chosen places). Each $R(\tilde h)$ is integral:
$\tilde h$ has integral components with integral inverses at places not above $\ell$, so (G2) of
\S\ref{pl:ss-norm} applies as in the proof of Lemma~\ref{pax:free}. The $R(\tilde h)$ commute with the
action of $\Gamma$ (translations at disjoint sets of places), and on $\mathcal A_\lambda(J_O^{(1)},\chi_O)$
each acts by $\chi_O(\tilde h)=1$, so $e_H$ acts there by $\#H$; moreover $e_H$ maps
$\mathcal A_\lambda(J_1,\chi)$ into $\mathcal A_\lambda(J_O^{(1)},\chi_O)$. Put $\Psi=e_H\Psi_1$, restricted
to $\mathcal A_\lambda(J_O^{(1)},\chi_O)$: an integral operator on the forms of principal level at the
$w_j$, and
\[
\sum_{k\in\Gamma'}k^{-1}\Psi k=e_H\sum_{k\in\Gamma'}k^{-1}\Psi_1k=e_H=\#H=c_0
\quad\text{on }\mathcal A_\lambda(J_O^{(1)},\chi_O).
\]
For $\Gamma'=\Gamma$ and $y$ of base level (i.e.\ fixed by $\Gamma$), $c_0y=\sum_{k}k^{-1}\Psi ky=
\sum_kk^{-1}\Psi y$, a norm $\sum_{k\in\Gamma}k\,(\Psi y)$.

(b) Let $y'$ be integral at every isomorphism and $y''$ integral at $\iota$; both are forms of level
$J_O$. The pairing of \S\ref{ss:au-spaces} at level $J_O$ is
\[
\langle y',y''\rangle=\sum_{g\in\Gamma_{\mathrm{gl}}\backslash G_f/J_O}\#\mathrm{Stab}_g^{-1}\,
\langle y'(g),y''(g)\rangle_{W_\lambda}.
\]
Each term $\iota\langle y'(g),y''(g)\rangle_{W_\lambda}$ is integral: this is the pointwise computation
in the proof of Proposition~\ref{pax:pairing}, which uses only that $y'$ is integral at $\bar\iota=
\iota\circ c$ and $y''$ at $\iota$, and not the freeness of the level. Since $\#\mathrm{Stab}_g$ divides
$N$ (see above), $N\langle y',y''\rangle=\sum_g(N/\#\mathrm{Stab}_g)\langle y'(g),y''(g)\rangle\in\pgZb$.
For completeness we recall the pointwise computation, with $\varphi=y'$, $\varphi'=y''$ and a fixed
$g$: after $\iota$, the vectors $\rho(g_\ell)^{-1}\iota\varphi'(g)$ and the corresponding
vectors for $\varphi$ through $\iota\circ c$ (whose action of $g_\ell$ is through the transposes of
the inverses of the components for $\iota$, as $g$ is unitary; (N2)) are integral, the invariant form
of $W_\lambda$ has an integral Gram matrix $\mathsf G$ in the basis $(e_T)$, and $\mathsf G$ is
invariant under the simultaneous action of $g_\ell$ on the two sides (through the transposed inverse
on the first), so the two factors $\rho(g_\ell)^{\pm1}$ cancel and each term is the value of
$\mathsf G$ on two integral vectors.
\end{proof}

\begin{remark}\label{pax:rem-forms}
At the free sublevel $J_0$ all stabilisers are trivial and the operator $\Psi_1$ gives the
constant $1$; at the base level $J_O$ the stabilisers $\Gamma_{\mathrm{gl}}\cap gJ_Og^{-1}$ may be
nontrivial (they embed in $H$), and averaging over $H$ gives $c_0=\#H$. All uses of $c_0$ in
\S\ref{pri:ss-norm} only need $c_0\mid N$. Likewise the pairing at level $J_O$ is not the pairing
at $J_0$ (for $J'\subseteq J$ one has $\langle\,,\rangle_{J'}=[J:J']\langle\,,\rangle_J$), which is why
(b) is proved directly at level $J_O$ with the factor $N$. The invariance of the Gram matrix
used at the end of (b) is proved in \S\ref{pax:ss-pairing} (Lemma~\ref{pax:gram}).
\end{remark}

\subsection{Block-diagonal lattices}\label{pax:ss-block}

For a field $L$ of characteristic $0$, a group $G$ and finite-dimensional representations $\rho_1,
\rho_2$ of $G$ over $L$: \emph{if $\rho_1,\rho_2$ are semisimple and $\pgtr\rho_1(g)=\pgtr\rho_2(g)$
for all $g\in G$, then $\rho_1\simeq\rho_2$} (the theorem of Brauer and Nesbitt in characteristic $0$; see also
Curtis--Reiner~\cite{CurtisReiner}). Proof: let $A$ be the image of $L[G]$ in
$\operatorname{End}_L(\rho_1\oplus\rho_2)$, a finite-dimensional $L$-algebra spanned by the operators of the
elements of $G$. It has a faithful semisimple module, $\rho_1\oplus\rho_2$, so its radical, which
annihilates every simple module, is $0$, and $A$ is semisimple. For a simple $A$-module $M$ let $e_M\in A$ be
the central idempotent that is the identity on $M$ and $0$ on the other simple modules. If $M$ occurs
$m_i$ times in $\rho_i$, the trace of $e_M$ on $\rho_i$ is $m_i\dim_LM$. Now $e_M$ is an $L$-linear
combination $\sum_gc_g(\rho_1\oplus\rho_2)(g)$, and its trace on $\rho_i$ is $\sum_gc_g\pgtr\rho_i(g)$,
the same for $i=1,2$. Hence $m_1\dim_LM=m_2\dim_LM$ in $L$, and $m_1=m_2$ because $\dim_LM\ne0$ in
characteristic $0$. As both are semisimple, $\rho_1\simeq\rho_2$.

\begin{lemma}[Exact form of Lemma~\ref{pri:block}]\label{pax:block}
Let $\rho_x$ ($x\in Y$ finite) be continuous \emph{semisimple} representations $\pgG\to\auGL_n(\auQlb)$
and, for $x\in\priX\subseteq Y$, $\pgtr\rho_x=\psi\,\pgtr R_x+\sum_{e<t'}\chi^e$, with $R_x$ of rank
$r'$, $n=r'+t'$, $\psi$ a continuous character and $\chi=\varepsilon_\ell^{-1}$. Let $u'_x$ be
lattices in $R_x$ over a coefficient ring $\auO_1$ such that, for every large $\auO\supseteq\auO_1$,
the $u'_x(g)$ span $M_{r'}(\auO)$ over $\auO$ (property (M4) of the main summands). Then for every large
$\auO$ there are lattices $u_x$ in the $\rho_x$ (continuous homomorphisms $\pgG\to\auGL_n(\auO)$
conjugate over $\auQlb$ to $\rho_x$) with $u_x=\psi u'_x\oplus\bigoplus_{e<t'}\auO(-e)$ for $x\in
\priX$, and the $\psi u'_x(g)$ span $M_{r'}(\auO)$.
\end{lemma}

\begin{proof}
Take $\auO$ large enough to contain the values of $\psi$ (a continuous character of the compact group
$\pgG$ takes values in the units of the ring of integers of a finite extension of $\Q_\ell$), to
satisfy the spanning hypothesis, and such that every $\rho_x$, $x\notin\priX$, has a lattice over
$\auO$ (a continuous representation of the compact group $\pgG$ over a finite extension $K'$ stabilises
a lattice, the $\pgG$-orbit of any lattice being finite; and each $\rho_x$ is defined over a finite
extension, by compactness and Baire's theorem; standard). For $x\notin\priX$ take such a lattice.

Let $x\in\priX$ and $\rho'_x=\psi u'_x\oplus\bigoplus_{e<t'}\auO(-e)$, a continuous homomorphism
$\pgG\to\auGL_n(\auO)$ ($\auO(-e)$ is the character $\varepsilon_\ell^{-e}=\chi^e$). \emph{Spanning.}
The $\psi(g)u'_x(g)$ span $M_{r'}(\auO)$: each is a unit multiple of $u'_x(g)$, and multiplying the
members of a spanning family by units does not change their span. \emph{Semisimplicity.} Over
$\auQlb$ the summand $\psi u'_x$ is absolutely irreducible, since the span of its image is
$M_{r'}(\auQlb)$ (Burnside), and the other summands are characters; so $\rho'_x\otimes\auQlb$ is
semisimple. \emph{Traces.} $\pgtr\rho'_x=\psi\pgtr u'_x+\sum_e\chi^e=\psi\pgtr R_x+\sum_e\chi^e=
\pgtr\rho_x$, as $u'_x$ is a lattice in $R_x$. By Brauer--Nesbitt in characteristic $0$, $\rho'_x
\otimes\auQlb\simeq\rho_x$, i.e.\ $u_x:=\rho'_x$ is a lattice in $\rho_x$ of the required shape.
\end{proof}

\begin{remark}\label{pax:rem-block}
Lemma~\ref{pri:block} has two hypotheses that the proof needs: the \emph{semisimplicity} of the
$\rho_x$ (in characteristic $0$ equivalent to the non-degeneracy of the trace form on the span of the
image) and the \emph{spanning} of $M_{r'}(\auO)$ by the $u'_x(g)$ (property (M4) of the main
summands, which is the input of the conclusion ``the $\psi u'_x(g)$ span''). Both hold where
Lemma~\ref{pri:block} is used (Proposition~\ref{pri:columns}): $\rho_x=S_x\otimes\eta_O$ is
semisimple, $S_x$ being a direct sum of the semisimple representations $r_\iota(\cdot)$ twisted by
characters (Proposition~\ref{pga:G123}), and $R_x$ is a main summand, so (M4) holds. Conjugate
self-duality is not used.
\end{remark}

\subsection{Integrality of the pairing at a free level}\label{pax:ss-pairing}

This completes the proof of Lemma~\ref{pri:forms}(b), which used the pointwise integrality of
$\langle\varphi(x),\varphi'(x)\rangle$ proved below. Let $G_\lambda$ be the Gram matrix of the invariant hermitian form of
$W_\lambda$ in the basis $(e_T)$; it is a positive definite \emph{integer} matrix
(\S\ref{pgs:sec}: $W_\lambda(\Z)$ is the lattice of integral tensors and the invariant form is the
standard form of the tensor space restricted to $W_\lambda$, \eqref{pgs:eq-lattice}). Write the form
as $\langle a,b\rangle={}^t\bar a\,G_\lambda\,b$ on coordinate vectors (the variable that is
conjugated is the one that must be integral at $\bar\iota$, as in \S\ref{ss:au-spaces}).

\begin{lemma}[Invariance of the Gram matrix]\label{pax:gram}
For every field $L$ of characteristic $0$ and every $g\in\prod_\tau\auGL_N(L)$,
${}^t\rho_\lambda(g)\,G_\lambda\,\rho_\lambda({}^tg^{-1})=G_\lambda$, where $\rho_\lambda(g)=
\bigotimes_\tau(\det g_\tau)^{e_\tau}S_{\kappa_\tau}(g_\tau)$ is the matrix of (N1) of
\S\ref{pl:ss-norm}.
\end{lemma}
\begin{proof}
This is the computation of (P1) in the proof of Lemma~\ref{pgs:pairing}: invariance of the hermitian
form under the compact group $\prod\auU(N)$ is ${}^t\overline{\rho_\lambda(g)}G_\lambda\rho_\lambda(g)=
G_\lambda$; the entries of $\rho_\lambda(g)$ are integral polynomials in the entries of $g$ and in
$\det g^{\pm1}$, so $\overline{\rho_\lambda(g)}=\rho_\lambda(\bar g)=\rho_\lambda({}^tg^{-1})$ for $g$
unitary, and with $h=\bar g$ the identity reads ${}^t\rho_\lambda(h)G_\lambda\rho_\lambda({}^th^{-1})=
G_\lambda$ on $\prod\auU(N)$. Both sides are polynomial in the entries of $h$ and $\det h^{-1}$ with
integer coefficients, and $\prod\auU(N)$ is Zariski dense in $\prod\auGL_N$; hence the identity holds
in the polynomial ring and therefore over every field of characteristic $0$.
\end{proof}

\begin{proposition}[Free-level pairing]\label{pax:pairing}
Let the level $J$ be free ($\Gamma_{\mathrm{gl}}\cap gJg^{-1}=1$ for all $g$). If $\varphi$ is integral at
$\bar\iota=\iota\circ c$ and $\varphi'$ at $\iota$, then $\iota\langle\varphi,\varphi'\rangle\in
\auO_{\C_\ell}$.
\end{proposition}
\begin{proof}
The computation of each term below is pointwise and does not use freeness. At a free level $\langle\varphi,\varphi'\rangle=\sum_x\langle\varphi(x),\varphi'(x)\rangle$ over
representatives $x$ of $\Gamma_{\mathrm{gl}}\backslash G_f/J$ (all weights $\#\mathrm{Stab}^{-1}$ equal $1$). Fix
$x$. By definition of the gauge (N2), $\iota(\varphi'(x))=\rho_\iota(x_\ell)\beta$ with $\beta$
integral, where $\rho_\iota(x_\ell)=\bigotimes_\tau\rho_\tau(\hat\iota_\tau(x_{w_\tau}))$. For
$\bar\iota$ the primes $w_\tau$ are replaced by their conjugates; since $x$ is unitary, its component
at $\overline{w}$ is ${}^t\bar x_w^{-1}$, so $\rho_{\bar\iota}(x_\ell)=\rho_\iota({}^tx_\ell^{-1})$ in
the sense of Lemma~\ref{pax:gram} applied factorwise (with $L=\C_\ell$), and $\iota(\overline{
\varphi(x)})=\bar\iota(\varphi(x))=\rho_\iota({}^tx_\ell^{-1})\alpha$ with $\alpha$ integral. Hence
\[
\iota\langle\varphi(x),\varphi'(x)\rangle={}^t\alpha\;{}^t\rho_\iota({}^tx_\ell^{-1})\,G_\lambda\,
\rho_\iota(x_\ell)\,\beta={}^t\alpha\,G_\lambda\,\beta
\]
by Lemma~\ref{pax:gram} with $g={}^tx_\ell^{-1}$, and the right side is integral since
$G_\lambda$, $\alpha$, $\beta$ are.
\end{proof}

With Proposition~\ref{pax:pairing}, the proof of Lemma~\ref{pri:forms}(b) in \S\ref{pax:ss-forms} is
complete.

\subsection{The dictionary at an auxiliary place}\label{pax:ss-dict}

We prove Proposition~\ref{pri:dict}(a)--(c), and (d), using the local Lemma~\ref{pax:toproot}
proved in \S\ref{pay:ss-cons}. Notation is that of Proposition~\ref{pri:dict}: $w$ a split degree-one
place, $w\nmid\ell$, off the base level, $q=\auNm w\equiv1\pmod{\ell^\nu}$, $F=E_w$, $\Gamma_w=
\auGL_n(\auF_w)$, $\gamma=\mathrm{Art}_w(z)$ the tame generator, $\phi$ a geometric Frobenius with
$\mathrm{Art}_w(\varpi)=\phi$ on the maximal abelian quotient. For a system $x$ of principal level
at $w$ write $\rho_x=S_x\otimes\eta_O$; the traces $\pgtr[\,\cdot\,]$ act on the forms of $x$ by the
scalars $\pgtr\rho_x(\cdot)$. Each identity below is an identity of operators on the forms of
principal level at $w$; since the $\Gamma_w$-action, the operators $U$ and the elements of $\priT$
commute with the projections to systems, it suffices to check it on forms $y$ in a single system
$x$, which we fix, together with an irreducible constituent $\pi$ of the automorphic representation
generated by the components of $y$; let $\Pi=\pi_w$, an irreducible unitary representation of
$\auGL_n(F)$ with nonzero vectors fixed by $K(1)=1+\varpi M_n(\auO_F)$.

\emph{Two facts used throughout.} (F1) $\eta_O$ is unramified at $w$ (it is the character of the
algebraic Hecke character of Proposition~\ref{pga:trace}, built from $\chi_W$ and $\upsilon$, which
are unramified at the places off the level of the pair; Lemma~\ref{pay:eta}), so
$\rho_x|_{I_w}=S_x|_{I_w}$. (F2) By
Proposition~\ref{pga:G123} (G1) and its proof, with Caraiani's part of (GR),
$\operatorname{WD}(S_x|_{G_{E_w}})^{\mathrm{F\text{-}ss}}\simeq\operatorname{rec}(\Pi\otimes|\det|^{(1-n)/2})$;
if the monodromy of the right side vanishes, $S_x|_{I_w}$ is the restriction of
$\operatorname{rec}(\Pi)$ to $I_w$ (a semisimple representation of finite image), and $S_x(\phi)$ has
the Frobenius eigenvalues of $\operatorname{rec}(\Pi\otimes|\det|^{(1-n)/2})$, the decomposition
under $D_w$ being respected. We use the sign normalisation $\epsilon=1$ of \S\ref{pga:ss-local}.

\begin{proof}[Proof of Proposition~\ref{pri:dict}(a)]
Let $\Gamma_w$ act on $y$ by $\chi\circ\mathrm{ind}\circ\det$, and let $\tilde\chi$ be a tamely
ramified character of $F^\times$ with $\tilde\chi(u)=\chi(\mathrm{ind}(\bar u))$ for $u\in
\auO_F^\times$. Then $\Pi\otimes(\tilde\chi^{-1}\circ\det)$ has a nonzero $\auGL_n(\auO_F)$-fixed
vector, so it is unramified, and $\operatorname{rec}(\Pi)=\operatorname{rec}(\Pi\otimes
\tilde\chi^{-1})\otimes\tilde\chi$ has monodromy $0$ and inertia acting by the scalar
$\tilde\chi\circ\mathrm{Art}_w^{-1}$. By (F1), (F2), $\rho_x(\gamma)=\tilde\chi(z)\cdot1=\chi(1)\cdot1$,
as $\mathrm{ind}(\bar z)=1$. Hence $\pgtr\rho_x((\gamma-1)g\,p(\phi))=(\chi(1)-1)\pgtr\rho_x(g\,p(\phi))$.
On the other side $\langle\gamma\rangle=\mathrm{diag}(\bar z,1,\dots,1)$ acts on $y$ by
$\chi(\mathrm{ind}\det\langle\gamma\rangle)=\chi(1)$. Both sides of (a) act on $y$ by
$(\chi(1)-1)\pgtr\rho_x(g\,p(\phi))$.
\end{proof}

For (b), let $\mathcal P\subseteq\auGL_n(\auO_F)$ be the stabiliser of the line $\auF_we_1$ modulo
$\varpi$ (block upper triangular of type $(1,n-1)$ modulo $\varpi$), $P=MN$ the upper parabolic of
type $(1,n-1)$, $d=\mathrm{diag}(\varpi,1,\dots,1)$, and $U$ the operator of $I_1dI_1$, $I_1$ the
pro-$p$ Iwahori subgroup (upper unipotent modulo $\varpi$).

\begin{lemma}[$U$ on parahoric vectors]\label{pax:Ueig}
Let $\tilde\chi$ be a tamely ramified character of $F^\times$, nontrivial on $\auO_F^\times$, and
suppose $\Pi$ has a nonzero vector $v$ on which $\mathcal P$ acts by $k\mapsto\tilde\chi(k_{11})$.
Then $\Pi\simeq\Pi_1\times\Pi_2$ with $\Pi_1$ a character of $F^\times$ restricting to $\tilde\chi$
on $\auO_F^\times$ and $\Pi_2$ unramified; the space of such vectors is one-dimensional; it is
contained in $\Pi^{I_1}$; and $Uv=q^{(n-1)/2}\Pi_1(\varpi)\,v$.
\end{lemma}
\begin{proof}
The first assertion is Lemma~\ref{pga:parahoric} with $(n_1,n_2)=(1,n-1)$ and $\psi=\tilde\chi|_{\auO_F^\times}$
(viewed on $\auF_w^\times$), applied to the vector $\pi(t)v$, where $t=\mathrm{diag}(1,\varpi^{-1}1_{n-1})\in
\auGL_n(F)$. Indeed $t\begin{pmatrix}a&b\\c&D\end{pmatrix}t^{-1}=\begin{pmatrix}a&\varpi b\\\varpi^{-1}c&D
\end{pmatrix}$, so $t\mathcal Pt^{-1}$ is exactly the parahoric of the matrices that are block lower
triangular modulo $\varpi$ for the blocks $(1,n-1)$, the entry $k_{11}$ is unchanged, and
$\pi(tkt^{-1})\pi(t)v=\pi(t)\pi(k)v=\tilde\chi(k_{11})\pi(t)v$ for $k\in\mathcal P$; the representation
$\Pi$ itself is not changed. Realise
$\Pi$ as the normalised induction $\operatorname{Ind}_P(\Pi_1\otimes\Pi_2)$, functions with
$f(mn g)=\delta_P^{1/2}(m)(\Pi_1\otimes\Pi_2)(m)f(g)$. By the Iwasawa decomposition $G=PK$ and
the Bruhat decomposition of $\auGL_n(\auF_w)$ for the parabolic of type $(1,n-1)$, $G=P\mathcal P
\sqcup Ps_1\mathcal P$, with $s_1$ the transposition of $e_1$ and $e_n$. A vector $f$ of type
$(\mathcal P,\tilde\chi(k_{11}))$ is determined by $f(1)$ and $f(s_1)$. On $P\mathcal P$: $f(1)$ is
fixed by $1\times\auGL_{n-1}(\auO_F)$ and $\auO_F^\times\times1$ acts on it by $\tilde\chi$, consistent
with $\Pi_1|_{\auO_F^\times}=\tilde\chi$; the possible values form a line, as $\Pi_2$ is unramified. On
$Ps_1\mathcal P$: for $k\in\mathcal P$ with $k'=s_1ks_1^{-1}\in M$, $f(s_1k)=f(k's_1)$ shows that
$f(s_1)\in\Pi_1\otimes\Pi_2$ is an eigenvector of $M\cap s_1\mathcal Ps_1^{-1}$ with the character
$k'\mapsto\tilde\chi(k'_{nn})$, up to the factor $\delta_P^{1/2}\Pi_1(k'_{11})$. The $\auGL_{n-1}$-part
of this group contains the parahoric $\mathcal Q\subseteq\auGL_{n-1}(\auO_F)$ of the matrices whose
reduction stabilises the line $\auF_w e_n$ (indeed $(s_1ks_1^{-1})_{jn}=k_{s_1(j),1}\equiv0$ for
$j\neq n$; these are the matrices that are block lower triangular of type $(n-2,1)$ modulo $\varpi$),
acting by $\tilde\chi$ of the
last diagonal entry. By Lemma~\ref{pga:parahoric} (for $\auGL_{n-1}$ with blocks $(1,n-2)$, applied to the image of
the vector under $\Pi_2(t'w'_0)$, where $w'_0$ is the antidiagonal permutation matrix of $\auGL_{n-1}$,
which reverses the indices and carries $\mathcal Q$ to the parahoric that is block upper triangular of type
$(1,n-2)$ modulo $\varpi$ with the character on the first diagonal entry, and $t'=\mathrm{diag}(1,
\varpi^{-1}1_{n-2})$, which carries the latter onto the block lower triangular parahoric of type
$(1,n-2)$ as in the first step; the vector $f(s_1)$, as a vector of $\Pi_2$ up to the scalar action of
the first factor, is replaced by its image under $\Pi_2(t'w'_0)$) a nonzero such vector would force a character with ramified restriction $\tilde\chi$
in the supercuspidal support of $\Pi_2$, which consists of unramified characters. So $f(s_1)=0$, and
the space is one-dimensional, spanned by $f_0$ supported on $P\mathcal P$. As $I_1
\subseteq\mathcal P$ and $\tilde\chi(k_{11})=1$ on $I_1$, $f_0\in\Pi^{I_1}$. Now $I_1=(I_1\cap\bar N)
(I_1\cap M)(I_1\cap N)$ with $I_1\cap N=N(\auO_F)$, $d^{-1}(I_1\cap\bar N)d\subseteq I_1\cap\bar N$
and $d$ normalising $I_1\cap M$, so $I_1dI_1=\bigsqcup_{u\in N(\auO_F)/dN(\auO_F)d^{-1}}u\,d\,I_1$, with
$q^{n-1}$ cosets. $U$ preserves the vectors of type $(\mathcal P,\tilde\chi(k_{11}))$ in $\Pi^{I_1}$:
$\mathcal P=(\mathcal P\cap\bar N)(\mathcal P\cap M)(\mathcal P\cap N)$, where $\mathcal P\cap\bar N=
\bar N(\varpi\auO_F)\subseteq I_1$ fixes the $I_1$-fixed vector $Uf$ and has trivial character;
$\mathcal P\cap N=N(\auO_F)$ permutes the cosets $ud\,I_1$ and has trivial character; and the elements
$m$ of $\mathcal P\cap M=\auO_F^\times\times\auGL_{n-1}(\auO_F)$ commute with $d$ and normalise $N(\auO_F)$
and $dN(\auO_F)d^{-1}$, so $\pi(m)Uf=\sum_u\pi(mum^{-1}d)\pi(m)f=\tilde\chi(m_{11})Uf$. Hence
$Uf_0=cf_0$ with $c=(Uf_0)(1)/f_0(1)$, and
\[
(Uf_0)(1)=\sum_uf_0(ud)=q^{n-1}f_0(d)=q^{n-1}\delta_P^{1/2}(d)\,\Pi_1(\varpi)f_0(1)=
q^{(n-1)/2}\Pi_1(\varpi)f_0(1),
\]
using $f_0(ud)=f_0(d)$ ($u\in N$) and $\delta_P(d)=|\varpi|^{n-1}=q^{1-n}$.
\end{proof}

\begin{proof}[Proof of Proposition~\ref{pri:dict}(b)]
Let $\mathcal P$ act on $y$ by $\chi(\mathrm{ind}(k_{11}))$ with $\chi(n!)\ne1$; in particular
$\chi\ne1$, and $\tilde\chi$ (as in (a)) is ramified. By Lemma~\ref{pax:Ueig}, $\Pi=\Pi_1\times\Pi_2$,
so $\operatorname{rec}(\Pi\otimes|\det|^{(1-n)/2})=\lambda\oplus r_2$ with $\lambda=
\operatorname{rec}(\Pi_1|\cdot|^{(1-n)/2})$ a character, $\lambda|_{I_w}=\tilde\chi\circ
\mathrm{Art}_w^{-1}\ne1$, $\lambda(\phi)=q^{(n-1)/2}\Pi_1(\varpi)$, and $r_2$ unramified with
monodromy $0$. By (F1), (F2), $\rho_x|_{G_{E_w}}=L\oplus L'$ with $L$ a line on which $D_w$ acts by
$\lambda\eta_O$ and $I_w$ acting trivially on $L'$. Let $e$ be the projection onto $L$ along $L'$;
it commutes with $\rho_x(\phi)$, and $\rho_x(\gamma)-1=(\lambda(\gamma)-1)e$. With $\mu=
\lambda(\phi)\eta_O(\phi)$,
\[
\pgtr\rho_x\bigl((\gamma-1)g\,p(\phi)\bigr)=(\lambda(\gamma)-1)\pgtr\bigl(p(\rho_x(\phi))e\rho_x(g)
\bigr)=p(\mu)\,(\lambda(\gamma)-1)\pgtr(e\rho_x(g))=p(\mu)\pgtr\rho_x((\gamma-1)g).
\]
The same computation holds for every constituent $\pi$ meeting $y$, with the same $\lambda$ (it is
read from $\rho_x$: the unique ramified line of $\rho_x|_{D_w}$). By Lemma~\ref{pax:Ueig}, applied
to the components of $y$ (which lie in the one-dimensional parahoric-type spaces of the constituents),
$Uy=q^{(n-1)/2}\Pi_1(\varpi)y=\lambda(\phi)y$, so $\tilde Uy=\eta_O(\phi)Uy=\mu y$ and $p(\tilde U)y=
p(\mu)y$. Hence both sides of (b) act on $y$ by $p(\mu)\pgtr\rho_x((\gamma-1)g)$.
\end{proof}

\begin{proof}[Proof of Proposition~\ref{pri:dict}(c)]
For any form $z$ of principal level at $w$, $\mathrm{Nm}_wz=\sum_{k\in\Gamma_w}kz$ is fixed by
$\Gamma_w$, i.e.\ of base level at $w$, hence spherical at $w$; its components in the systems $x$
are spherical at $w$, so $S_x$ is unramified at $w$ by (G2), and with (F1) $\rho_x(\gamma)=1$, so
$\pgtr\rho_x((\gamma-1)g)=0$ for each system $x$ meeting $\mathrm{Nm}_wz$. As $\pgtr[(\gamma-1)g]$
commutes with $\Gamma_w$, $\mathrm{Nm}_w\pgtr[(\gamma-1)g]z=\pgtr[(\gamma-1)g]\mathrm{Nm}_wz=0$.
\end{proof}

\emph{Part (d).} Let $\delta\in\priT$, $p=(X-q^{t'-1})H$, $B(1)=p(\tilde U)L(1)$ with $L(1)=
\sum_{k\in\mathcal P/K(1)}k$ the plain norm over the line stabiliser (weight $1$). The
deduction of (d) uses the following local statement.

\begin{lemma}[The top root of a chain]\label{pax:toproot}
Let $\pi$ be an irreducible unitary unramified representation of $\auGL_n(F)$ with invariant inner
product $(\ ,\ )$, $x\ne0$ spherical, $X_0\ne0$, $k\ge1$, and suppose $H_\pi(X_0X)=\prod_{i<k}(X-q^i)
\cdot S(X)$ with $S\mid R$. Then $(x,\ ((X-q^{k-1})R)(X_0^{-1}U)\,z)=0$ for every $z\in\pi^{I_1}$.
\end{lemma}
\noindent Proved in \S\ref{pay:ss-cons}, from the spherical chain lemma,
Theorem~\ref{pay:chain}(ii).\par\smallskip

\begin{proof}[Proof of Proposition~\ref{pri:dict}(d)]
Put $Z=\mathrm{Nm}\,B(1)$ ($\mathrm{Nm}=\sum_{k\in\Gamma}k$); its values are $\Gamma$-fixed, i.e.\ of
base level. Let $u$ be a form of base level in a system $x$ with $\delta(x)\ne0$; by hypothesis
$H_w(x;X_0X)=\prod_{i<t'}(X-q^i)\cdot S$ with $S\mid H$, $X_0=\eta_O(\phi)^{-1}$, so $\tilde U=X_0^{-1}U$.
For every form $y$, $L(1)y$ is fixed by $I_1$ (as $I_1\subseteq\mathcal P$), and since the
components of $u$ in the constituents are spherical, Lemma~\ref{pax:toproot} (with $k=t'$, $R=H$),
applied in each constituent with the restriction of the global inner product (which is
$G_f$-invariant), gives $\langle u,B(1)y\rangle=0$. As $u$ is $\Gamma$-fixed and the inner product
$\Gamma$-invariant, $\langle u,Zy\rangle=\sum_k\langle k^{-1}u,B(1)y\rangle=\#\Gamma\,\langle u,B(1)y
\rangle=0$. Now $\delta Zy$ is of base level and lies in the span of the systems with $\delta\ne0$;
by normality of $\delta$ (adjoint in $\priT$, preserving the systems) $\langle\delta Zy,\delta Zy\rangle
=\langle\delta^*\delta Zy,Zy\rangle$, and $\delta^*\delta Zy$ is a sum of base-level forms in systems
with $\delta\ne0$, orthogonal to $Zy$ by the above. So $\delta Zy=0$.
\end{proof}

\subsection{Base level of a label}\label{pax:ss-BL}

We prove Proposition~\ref{pri:BL}, with the following precise meaning of ``base level'': the forms
of the weight $\lambda$ of $O$ on $\auU(V)$ of the level $J_O$ and level character $\chi_O$ of
Definition~\ref{prm:def-level} (at a chosen $w$, $k\equiv\left(\begin{smallmatrix}1_{r'}&*\\0&*
\end{smallmatrix}\right)$ modulo $w$ with trivial character; at $w\in Q$, $k$ block lower triangular
modulo $w$ with character $\psi_w^{-1}(\det k^{(11)})$; hyperspecial elsewhere), that are fixed by the
idempotent $e_c$ of the labels. Let $G_c=\prod_{w\ \mathrm{chosen}}\auGL_{r'}(\auF_w)$, acting on
these forms by $\varrho$ of Definition~\ref{prm:def-level} (lifts reducing to $\mathrm{diag}(g_w,
1_{t'})$), twisted by $\prod_w\psi_w(\det g_w)$ ($\psi_w$
the tame part at $w$), and $e_c\in\C[G_c]$ the central idempotent of the irreducible character named
by the labels.

\begin{proof}[Proof of Proposition~\ref{pri:BL}]
(a) holds by definition.

(b) $e_c=\sum_{g\in G_c}a_gg$ with $a_g=\dim(\sigma)\,\overline{\operatorname{ch}_\sigma(g)}/\#G_c$
(the central idempotent of the irreducible character $\sigma$), so $\#G_c\cdot a_g$ is an algebraic
integer. The twisting factors are roots of unity. The lift of $g$ supported at the chosen places has
integral components with integral inverses, so its translation preserves forms integral at $\iota$
((G2) of \S\ref{pl:ss-norm}). Hence $\#G_c\,e_cy$ is integral when $y$ is, and so is $Ne_cy$, as $\#G_c$
divides $N=\pgNch=\prod_{w\ \mathrm{chosen}}\#\auGL_{rt}(\auF_w)$ (by Lagrange's theorem, since
$G_c=\prod_w\auGL_{r'}(\auF_w)$ and $\auGL_{r'}(\auF_w)$ embeds in $\auGL_{rt}(\auF_w)$ as
$r'\le r\le rt$).

(c) Let $x$ be a system occurring in base level with principal level at $U$, $y\ne0$ a form of $x$,
and $S_x=S_y$ as in Proposition~\ref{pga:G123}; put $\rho_x=S_x\otimes\eta_O$.
\emph{Conjugate self-duality.} By (GR), $S_x^c\simeq S_x^\vee\varepsilon_\ell^{1-n}$ (each summand
$r_\iota(\Sigma)\varepsilon_\ell^{e}$ of Proposition~\ref{pga:G123} is polarised in this way, with the
twist of Proposition~\ref{pax:sign}). The algebraic Hecke character
$\eta$ with $\ell$-adic character $\eta_O$ (Proposition~\ref{pga:trace}) satisfies $\eta\eta^c=
|\cdot|_{\mathbb A_E}^{\,r'}$: the factors $\chi_W(y/\bar y)$ cancel, $\upsilon(y)\upsilon(\bar y)=
\upsilon|_{\mathbb A_{E^+}}(\auNm y)=1$ ($\upsilon|_{\mathbb A_{E^+}^\times}$ is the quadratic character
of $E/E^+$), and $|y|^{-r'/2}|\bar y|^{-r'/2}=|y|^{-r'}$, inverted. So $\eta_O\eta_O^c=
\varepsilon_\ell^{r'}$ and $\rho_x^c\simeq\rho_x^\vee\varepsilon_\ell^{1-n+r'}=\rho_x^\vee\chi^{t'-1}$,
which is conjugate self-duality with string length $t'$ (the form (M1)).
\emph{Hodge--Tate data.} By (G3) the weights of $S_x$ are determined by $\lambda$, and those of
$\eta_O$ by the pair; let $J'$ be the resulting data (independent of $x$).
\emph{Ramification.} $y$ is fixed by a hyperspecial subgroup at every place off $\ell p_0p_1$,
$Q\cup\overline Q$, $U\cup\overline U$ (the level $J_O$ is hyperspecial there and $y$ has principal
level only at $U$), so $S_x$ is unramified there by (G2); $\eta_O$ is unramified off $\ell$ and the
places of the pair ((F1) of \S\ref{pax:ss-dict}; Lemma~\ref{pay:eta}).
\emph{Inertia at a chosen $w$.} By Lemma~\ref{prm:label}(b), $\Pi_w^{K(1)}$ contains an irreducible
representation of $\auGL_n(\auF_w)$ with cuspidal support $(L,\bigotimes_k\sigma_{\tilde f_k}\otimes
\mathbf 1^{\otimes t'})$, $\tilde c_w=\prod_k\tilde f_k$ the rescaled label. By the cuspidal-support
statement in the proof of
Lemma~\ref{pga:depth0} (Moy--Prasad), the supercuspidal support of $\Pi_w$ consists of the depth-zero
supercuspidals attached to the irreducible factors of $\tilde c_w$ and of $t'$ unramified characters;
so $\operatorname{rec}(\Pi_w)|_{I_w}$ (semisimplification) has the eigenvalues $\omega(\xi)$, $\xi$ the
roots of $\tilde c_w$, and $1$ ($t'$ times), at a tame generator. The characteristic polynomial of
$S_x(s)$, $s\in I_w$, equals that of $\operatorname{rec}(\Pi_w)(s)$ (by (GR) the Weil--Deligne
representations agree up to Frobenius semisimplification, and the nilpotent monodromy commutes with
$I_w$, so it does not change characteristic polynomials of inertia). As $\eta_O$ is unramified at $w$ (Lemma~\ref{pay:eta}), inertia at $w$ acts on $\rho_x$ through the roots of the rescaled
label times $(X-1)^{t'}$.
\emph{No monodromy at a chosen $w$.} We show that $\operatorname{rec}(\Pi_w)$ has monodromy $0$;
here $\pi$ is an irreducible constituent of the automorphic representation generated by $y$ on
which the orthogonal projection $y_\pi$ of $y$ is non-zero, $S_x=S_y$ is formed with this $\pi$ in
Proposition~\ref{pga:G123} (up to isomorphism $S_y$ depends only on $x$: it is semisimple and its
Frobenius characteristic polynomials at almost all split places are given by the Hecke eigenvalues
through (G1), so Chebotarev and Brauer--Nesbitt apply), and $\Pi_w=\pi_v$ is the component at $w$
((BC) at split places). Notation: $G_w=\auGL_n(\auF_w)$, $P=MU_P\subseteq G_w$ the block upper
triangular parabolic subgroup of type $(r',t')$ ($r',t'\ge1$, both being odd), $\sigma'=\sigma_{c_w}
\otimes(\psi_w^{-1}\circ\det)=\sigma_{\tilde c_w}$, as in Lemma~\ref{prm:label}.

\emph{Step 1: a vector of type $\sigma'\boxtimes\mathbf 1$.} By the proof of Lemma~\ref{prm:label}(a),
$y$ is fixed by $U_P$ and by $1\times\auGL_{t'}(\auF_w)$ and lies in the $\sigma'$-isotypic part of
$V_w$ for $g\mapsto\mathrm{diag}(g,1)$. The orthogonal projection onto the $\pi$-part commutes with
right translations, so $y_\pi$ has the same three properties, and by the proof of
Lemma~\ref{prm:label}(b) the $\pi$-part of $V_w$ is $\Pi_w^{K_w(1)}\otimes M_\pi$ with $G_w$ acting on
the first factor only. Writing $y_\pi=\sum_iv_i\otimes m_i$ with the $m_i$ linearly independent, each
$v_i$ has the three properties, and some $v_i\ne0$. So $\Pi_w^{K_w(1)}$ contains $v\ne0$ fixed by
$U_P$ and by $1\times\auGL_{t'}(\auF_w)$, in the $\sigma'$-isotypic part for $\auGL_{r'}(\auF_w)\times1$.

\emph{Step 2: the Jacquet module.} Let $\mathcal Q\subseteq\auGL_n(\auO_w)$ be the inverse image of
$P$, $\mathcal Q_1$ the inverse image of $U_P$, $P_w=M_wN_w$ the block upper triangular parabolic
subgroup of $\auGL_n(E_w)$ of type $(r',t')$, and $M_1=\mathcal Q_1\cap M_w=(1+\varpi_wM_{r'}(\auO_w))
\times(1+\varpi_wM_{t'}(\auO_w))$. Statement (J) of the proof of Lemma~\ref{pga:parahoric}
(Moy--Prasad~\cite[Proposition~6.7]{MoyPrasad}) holds for $\mathcal Q$, with the same verification of
its hypotheses: $\mathcal Q$ is a non-maximal parahoric subgroup $\mathcal G_x$ for a point $x$ of
the apartment of the diagonal torus, with pro-unipotent radical $\mathcal Q_1$ and attached Levi
subgroup $M_w$, and either of the two opposite parabolic subgroups with Levi $M_w$ may be used. So
the projection $\Pi_w\to\Pi_{w,N_w}$ restricts to an isomorphism $\Pi_w^{\mathcal Q_1}\to(\Pi_{w,N_w})^{
M_1}$, equivariant for $\mathcal Q/\mathcal Q_1=M\simeq\auGL_{r'}(\auF_w)\times\auGL_{t'}(\auF_w)$. Now
$v\in(\Pi_w^{K_w(1)})^{U_P}=\Pi_w^{\mathcal Q_1}$ lies in the $(\sigma'\boxtimes\mathbf 1)$-isotypic part
for $M$. Hence $(\Pi_{w,N_w})^{M_1}$ has a non-zero $(\sigma'\boxtimes\mathbf 1)$-isotypic part. Taking
isotypic parts for the compact group $M_w\cap\auGL_n(\auO_w)$ is exact and $\Pi_{w,N_w}$ has finite
length, so some irreducible subquotient $\tau=\tau_1\otimes\tau_2$ of $\Pi_{w,N_w}$ has this property:
$\tau_1^{K(1)}$ contains $\sigma'$, and $\tau_2^{K(1)}$ contains the trivial representation of
$\auGL_{t'}(\auF_w)$, that is, $\tau_2$ has a non-zero vector fixed by $\auGL_{t'}(\auO_w)$. (The
normalised Jacquet module differs from $\Pi_{w,N_w}$ by the unramified character
$\delta_{P_w}^{-1/2}$ of $M_w$, which changes neither property.)

\emph{Step 3: the supercuspidal support.} $\sigma'$ has cuspidal support $(L_1,\bigotimes_k\lambda_k)$,
$\lambda_k=\sigma_{\tilde f_k}$ (paragraph ``Twisted labels'' of \S\ref{prm:ss-level}); by the
cuspidal-support statement in the proof of Lemma~\ref{pga:depth0} (for $\auGL_{r'}$), $\tau_1$ has
supercuspidal support $\{\pi_1,\dots,\pi_s\}$, $\pi_k$ a depth-zero supercuspidal representation of
$\auGL_{d_k}(E_w)$ containing $\sigma_{\tilde f_k}$. The trivial representation of $\auGL_{t'}(\auF_w)$
has cuspidal support $(T,\mathbf 1)$ (it is a constituent of $R_T^{\auGL_{t'}}(\mathbf 1)$), so by the
same statement (for $\auGL_{t'}$, all $f_k=X-1$) $\tau_2$ has supercuspidal support $t'$ characters
$\chi_1,\dots,\chi_{t'}$ of $E_w^\times$ trivial on $\auO_w^\times$, i.e.\ unramified. The supercuspidal
support of $\Pi_w$ is that of $\tau$ (Bernstein--Zelevinsky, as in the proof of
Lemma~\ref{pga:parahoric}): $\{\pi_1,\dots,\pi_s,\chi_1,\dots,\chi_{t'}\}$.

\emph{Step 4: no linking, and the shape of $\Pi_w$.} A segment of $\pi_k$ can be linked to one of
$\chi_j$ only if $\pi_k\simeq\chi_j|\cdot|^u$ for some $u$; then $d_k=1$ and $\pi_k$ is trivial on
$\auO_w^\times$. But for $d_k=1$ the restriction of $\pi_k$ to $\auO_w^\times$ is the inflation of the
character $\sigma_{\tilde f_k}$ of $\auF_w^\times$, $\tilde f_k=X-\xi$, which is trivial only if
$\xi=1$; and no rescaled label has the root $1$ (lift setting, \S\ref{pri:sec}). Moreover two
$\pi_k$ are never inertially equivalent, since $\tilde c_w$ is squarefree (it is a label; proof of
Lemma~\ref{pga:depth0}, paragraph on depth-zero supercuspidals with a given $\sigma_f$). Hence, by
Zelevinsky's classification exactly as in the proof of Lemma~\ref{pga:parahoric},
$\Pi_w\simeq\Pi_1\times\Pi_2$ with $\Pi_1$, $\Pi_2$ irreducible of supercuspidal supports
$\{\pi_k\}$ and $\{\chi_j\}$; and $\Pi_1\simeq\pi_1\times\dots\times\pi_s$ by Zelevinsky's
irreducibility criterion~\cite[Theorem~4.2]{Zelevinsky}, as in the proof of Lemma~\ref{pga:depth0}.
By the geometric lemma of Bernstein--Zelevinsky, the composition factors of the Jacquet module of
$\Pi_1\times\Pi_2$ for $P_w$ are subquotients of $(\alpha_1\times\alpha_2)\otimes(\beta_1\times\beta_2)$,
where $\alpha_i\otimes\beta_i$ is a composition factor of the Jacquet module of $\Pi_i$ for a parabolic
subgroup of type $(a_i,b_i)$, $a_1+b_1=r'$, $a_2+b_2=t'$, $a_1+a_2=r'$, $b_1+b_2=t'$. The supercuspidal
support of $\tau_2$ consists of unramified characters, while that of $\beta_1$ is contained in
$\{\pi_k\}$, which contains none; so $b_1=0$, hence $a_1=r'$, $a_2=0$, $b_2=t'$, and the only such
factor is $\Pi_1\otimes\Pi_2$. Thus $\tau_2$ is $\Pi_2$ up to an unramified twist, and $\Pi_2$ has a
non-zero vector fixed by $\auGL_{t'}(\auO_w)$: it is spherical. So
\[
\Pi_w\simeq\pi_1\times\dots\times\pi_s\times\Pi_2,\qquad\Pi_2\ \text{an irreducible spherical
representation of }\auGL_{t'}(E_w);
\]
in particular no constituent of Steinberg type occurs in the unramified part.

\emph{Step 5: the parameter.} $\operatorname{rec}(\Pi_w)=\bigoplus_k\operatorname{rec}(\pi_k)\oplus
\operatorname{rec}(\Pi_2)$; each $\operatorname{rec}(\pi_k)$ is irreducible, with monodromy $0$, and
$\operatorname{rec}(\Pi_2)$ is unramified with monodromy $0$ ($\Pi_2$ spherical, as in the proof of
Lemma~\ref{pga:parahoric}). So $\operatorname{rec}(\Pi_w)$ has monodromy $0$, and by the last
assertion of Proposition~\ref{pga:G123} (which rests on the identity of (GR) including monodromy,
\cite[Thm.~2.1.1(2)]{BLGGT}, at $w\nmid\ell$) the monodromy of $\operatorname{WD}(S_x|_{G_{E_w}})$ vanishes. Then
$S_x|_{I_w}$ is the restriction to $I_w$ of the Weil group representation underlying this
Weil--Deligne representation, which is trivial on an open subgroup of $I_w$; so $S_x(s)$ has finite
order and is semisimple for $s\in I_w$, and so is $\rho_x(s)$ ($\eta_O$ is unramified at $w$). By the
previous paragraph its eigenvalues are roots of $o_w(X)(X-1)$, where $o_w(X)=\prod_\xi(X-\omega(\xi)^
\epsilon)$, $\xi$ running over the roots of $\tilde c_w$; hence
\[
o_w(\rho_x(s))\,(\rho_x(s)-1)=0\qquad(s\in I_w).
\]
\emph{Inertia above $Q\cup\overline Q$.} At $w\in Q$ the form $y$ is an eigenvector of the parahoric
of type $(n_W,t')$ with the character $\psi_{\mathcal P}$, $\psi=\psi_w\ne1$; by (G5)
(Proposition~\ref{pga:G45}), $(S_x(s)-\psi(s))(S_x(s)-1)=0$ for $s\in I_w$, and $\eta_O$ is unramified at
$w$. At $\bar w$ the same holds by conjugate self-duality, $\psi$ being anticyclotomic. This is (c).
\end{proof}

\begin{remark}\label{pax:rem-BL}
So Proposition~\ref{pri:BL} follows from Propositions~\ref{pga:G123}, \ref{pga:G45}, the
cuspidal-support part of Lemma~\ref{pga:depth0} and the definitions; Lemma~\ref{au:free} is not
needed for (a)--(c) (it enters through Lemma~\ref{pri:forms}). That $\eta_O$ is unramified at
$Q\cup\overline Q$ and at the places $U$ (indeed at every finite place not above $\ell$) is
Lemma~\ref{pay:eta}. The trivial block $\mathbf 1_{\auGL_{t'}}$ comes from the equivariance of the
single lift, Lemma~\ref{pay:product}; the property of the base level at the chosen places is
Lemma~\ref{prm:label} (Definition~\ref{prm:def-level}). The label at $w$ is not squarefree (the factor $(X-1)^{t'}$), so
Lemma~\ref{pga:depth0} itself does not apply to $\Pi_w$; it is used for the cuspidal support only.
The vanishing of the monodromy at a chosen $w$ is proved in the paragraph ``No monodromy at a chosen
$w$'' from that cuspidal support, statement (J) of the proof of Lemma~\ref{pga:parahoric}, Zelevinsky's
classification and the geometric lemma; it uses that no rescaled label has the root $1$. Without
that clause of the lift setting a factor $X-1$ of $\tilde c_w$ would contribute an unramified
character linked possibly to the $\chi_j$, and Steinberg-type monodromy could not be excluded.
\end{remark}

\subsection{The sign hypothesis of Proposition~\ref{pga:G123}}\label{pax:ss-sign}

\begin{proposition}\label{pax:sign}
In Proposition~\ref{pga:G123} the hypothesis of (GR) holds for every constituent, after replacing the
twist $|\det|^{(m_i-1)/2}$ by $|\det|^{s_i}$ with $s_i\in\frac12\Z$, $s_i\equiv(n_i-1-n+m_i)/2
\pmod1$; with this twist,
\[
S_y=\bigoplus_i\bigoplus_{a=0}^{m_i-1}r_\iota\bigl(\sigma_i|\det|^{s_i}\bigr)\otimes
\varepsilon_\ell^{\,e_{i,a}},\qquad e_{i,a}=\tfrac{m_i-1}2-a-s_i+\tfrac{n_i-n}2\in\Z ,
\]
and the conclusions of Proposition~\ref{pga:G123} hold without any sign hypothesis.
\end{proposition}
\begin{proof}
Fix $i$ and write $\sigma=\sigma_i$, $m=m_i$, $N=n_i$, $\Sigma=\sigma|\det|^{s}$, $s=s_i$.
\emph{Algebraicity.} $\Pi$ is cohomological for $\auGL_n$, so its infinitesimal character at each
$\tau$ has entries in $\Z+\frac{n-1}2$. The entries for $\operatorname{Speh}(\sigma,m)$ are those of
$\sigma$ shifted by $\frac{m-1}2-a$, $0\le a<m$; so those of $\sigma$ lie in $\Z+\frac{n-m}2$, and those
of $\Sigma$ in $\Z+\frac{n-m}2+s=\Z+\frac{N-1}2$. Regularity of $\Pi$ gives regularity of $\Sigma$. So
$\Sigma$ is regular algebraic cuspidal.
\emph{Polarisation.} By (BC) $\operatorname{Speh}(\sigma,m)$ is conjugate self-dual; as $\operatorname{Speh}
(\sigma,m)^c=\operatorname{Speh}(\sigma^c,m)$, $\operatorname{Speh}(\sigma,m)^\vee=\operatorname{Speh}
(\sigma^\vee,m)$ and $\operatorname{Speh}(\cdot,m)$ is injective on cuspidal representations (its
cuspidal support determines $\sigma$), $\sigma^c\simeq\sigma^\vee$. Hence $\Sigma^c\simeq\Sigma^\vee
\otimes|\det|^{2s}=\Sigma^\vee\otimes(\phi\circ\auNm_{E/E^+}\circ\det)$ with $\phi=|\cdot|_{\mathbb
A_{E^+}}^{2s}$, an algebraic character ($2s\in\Z$) with $\phi_v(-1)=1$ for every $v\mid\infty$. This is
the hypothesis of (GR), the sign being the same at all infinite places. If the further sign condition in
the definition of a polarised pair in \cite[\S2.1]{BLGGT} is not met, replace $\phi$ by $\phi\,
\delta_{E/E^+}$, $\delta_{E/E^+}$ the quadratic character of $E/E^+$: since $\delta_{E/E^+}\circ
\auNm_{E/E^+}=1$ the relation $\Sigma^c\simeq\Sigma^\vee\otimes(\phi\circ\auNm_{E/E^+}\circ\det)$ is
unchanged, as in \cite[\S2.1]{BLGGT}.
\emph{Galois representation.} At a finite $v\nmid\ell$ where everything is unramified, the Satake
parameters of $\operatorname{Speh}(\sigma,m)_v\otimes|\det|^{(1-n)/2}$ are those of $\Sigma_v\otimes
|\det|^{(1-N)/2}$ times $|\cdot|^{e_a}$ with $e_a=\frac{m-1}2-a-s+\frac{N-n}2$, an integer by the
choice of $s$; $|\cdot|^{e}$ corresponds under $\operatorname{rec}$ to $\varepsilon_\ell^{e}$ (geometric
Frobenius, $\varepsilon_\ell(\pgFrob_w)=\auNm w^{-1}$). So the summand $r_\iota(\Sigma)\varepsilon_\ell
^{e_a}$ has the local parameters of $\operatorname{Speh}(\sigma,m)$, and the rest of the proof of
Proposition~\ref{pga:G123} applies verbatim; in particular the ramified local--global
compatibility including the monodromy operator, used in (F2) of \S\ref{pax:ss-dict}, holds for the
polarised pair $(\Sigma,|\cdot|^{2s})$ by \cite{BLGGT}, as recalled in (GR), and does not require
$\Sigma$ to be conjugate self-dual. (The choice $s=\frac{m-1}2$
is admissible exactly when $n-n_i$ is even; this is the parity issue of Remark~\ref{pga:sign}: the
problem is algebraicity of the twist, not the sign, and it disappears with the choice of $s$ above.)
\end{proof}

\subsection{The convention check at the end of \S\ref{pga:ss-M}}\label{pax:ss-conv}

The check asks that the forms $\mathcal A_\lambda$ of \S\ref{ss:au-spaces} have, in the sense of (G3),
the Harish-Chandra parameter $\varsigma=\lambda+\rho$ at each $\tau$, in the normalisation in which a
character of $\auU(1)$ of weight $c$ has parameter $c$.

\begin{lemma}\label{pax:conv}
Let $\varphi\in\mathcal A_\lambda(J,\chi)$ be nonzero, and $\pi$ an irreducible constituent of the
automorphic representation of $\auU(V)(\mathbb A_{E^+})$ formed by the functions $f_{\varphi,v'}(u,g)
=\langle v',u^{-1}\varphi(g)\rangle$ ($u\in\auU(V)(E^+\otimes\R)$, $v'\in W_\lambda^\vee$; this is the
reading of algebraic forms used in Lemma~\ref{pgs:pairing}'s sequel). Then at the archimedean place
$v$ below $\tau$, with $\auU(V)(E^+_v)\simeq\auU(n)$ through $\tau$, $\pi_v$ is the irreducible algebraic
representation of highest weight $\lambda_\tau$, so its parameter is $\lambda_\tau+\rho=\varsigma_\tau$.
For $n=1$ and $\lambda=e$ (the character $\det^e$) this is the character $z\mapsto\tau(z)^e$, of
parameter $e$.
\end{lemma}
\begin{proof}
By (N2), $\gamma\in\Gamma$ acts on $W_\lambda$ by $\rho_\lambda(\tau(\bar\gamma))=\rho_\lambda({}^t
\tau(\gamma)^{-1})$, and this extends to the real group: $u$ acts by $\rho_\lambda({}^t\tau(u)^{-1})=
\rho_\lambda(\overline{\tau(u)})=\overline{\rho_\lambda(\tau(u))}$ ($u$ unitary, $\rho_\lambda$ with real
coefficients). So $W_\lambda$, as a representation of $\auU(n)$ through $\tau$, is the complex
conjugate of $\rho_\lambda\circ\tau$, i.e.\ its contragredient (a unitary representation). Right
translation by $u_0\in\auU(V)(E^+_v)$ on $f_{\varphi,v'}$ gives $f_{\varphi,v'}(uu_0,g)=\langle v',u_0^{-1}
u^{-1}\varphi(g)\rangle=\langle u_0\cdot v',u^{-1}\varphi(g)\rangle$, so the functions $f_{\varphi,v'}$,
$v'$ varying, span a copy of $W_\lambda^\vee$, the dual of the contragredient of $\rho_\lambda\circ\tau$:
$\pi_v\simeq\rho_\lambda\circ\tau$, of highest weight $\lambda_\tau$ (the matrix $(\det g)^eS_\lambda(g)$
being the representation of highest weight $\lambda$ in the standard Borel by \eqref{pl:eq-diag}). The
case $n=1$ is the special case $S_\lambda=1$.
\end{proof}

So the dictionary is: Gross weight $\lambda$ $\leftrightarrow$ $\pi_\infty=\rho_\lambda\circ\tau$
$\leftrightarrow$ parameter $\varsigma=\lambda+\rho$, uniformly in $n$, including $n=1$; this is the
convention assumed in (G3), in Proposition~\ref{pga:weights}, and for the weight of $\mu^0_j$ (a
character of $\auU(1)$ read through the same (N2)). The relative sign used in
Proposition~\ref{pga:weights} is therefore justified. (Had the automorphic forms attached to
algebraic forms been taken to be the complex conjugates $\langle u^{-1}\varphi(g),v\rangle$, one would get the
contragredient at every place and $\varsigma$ would become $-w_0\varsigma$ uniformly in $n$, including
$n=1$; the relative sign would be unchanged.)


%% ===== proofs/C-aux2.tex =====
% proofs/C-aux2.tex
% Further auxiliary statements for Ribet's method (continuation of Section pax): the spherical
% chain lemma proved directly (no classification of the unitary dual), hence Lemma pax:toproot
% (closing Proposition pri:dict(d)) and Lemma pax:tadic without Tadic's classification; the two facts
% of Remark pax:rem-BL. Labels of this file start with pay:.
\providecommand{\auO}{\mathcal O}
\providecommand{\auGL}{\mathrm{GL}}
\providecommand{\auF}{\mathbb F}
\providecommand{\auU}{\mathrm U}
\providecommand{\auNm}{\operatorname{Nm}}
\providecommand{\pgtr}{\operatorname{tr}}
\providecommand{\pgFrob}{\operatorname{Frob}}
% macros of this file (prefix \pay)
\providecommand{\payL}{\Lambda}
\providecommand{\payM}{\mathsf M}
\providecommand{\payU}{\mathbf U}
\providecommand{\payUd}{\mathbf U'}
\providecommand{\payN}{\mathcal N}

\section{The spherical chain lemma}\label{pay:sec}

This section proves, without using the classification of the unitary dual of $\auGL_N$, the local
statement behind Lemma~\ref{pax:tadic} (hence behind Lemma~\ref{pri:chain}) and the local
Lemma~\ref{pax:toproot}, used in the proof of Proposition~\ref{pri:dict}(d) in \S\ref{pax:ss-dict}. Both come from one theorem (Theorem~\ref{pay:chain}): for a nonzero
spherical Hecke eigenvector in a unitary representation of $\auGL_N(F)$, a root $c$ of its Hecke
polynomial with $qc$ not a root satisfies $|c|^2>q^{N-2}$, and of a chain of roots the operator $U$
of $\mathrm{diag}(\varpi,1,\dots,1)$ at pro-$p$ Iwahori level sees only the top one. The proof
consists of two counts of hyperplanes over the residue field (\S\ref{pay:ss-count}), which give
recursions for $U$ and for its adjoint on a chain of explicit vectors (\S\ref{pay:ss-rec}), and of
linear algebra in a positive definite hermitian space (\S\ref{pay:ss-lin}). The two facts used in
Remark~\ref{pax:rem-BL} are proved in \S\ref{pay:ss-BL}.

\subsection{Setting}\label{pay:ss-set}

Let $F$ be a non-archimedean local field, $\auO_F$ its ring of integers, $\varpi$ a uniformiser,
$\auF=\auO_F/\varpi$ the residue field, of order $q$, and $N\ge1$. Put $G=\auGL_N(F)$, $K=\auGL_N(
\auO_F)$, $e_1,\dots,e_N$ the standard basis of $F^N$ (column vectors), $d_i=\mathrm{diag}(\varpi
1_i,1_{N-i})$, $d=d_1=\mathrm{diag}(\varpi,1,\dots,1)$, and let $I_1\subseteq K$ be the pro-$p$
Iwahori subgroup: the elements of $K$ that are upper triangular unipotent modulo $\varpi$.

Let $(\rho,V)$ be a smooth representation of $G$ on a complex vector space with a positive definite
$G$-invariant hermitian form $(\ ,\ )$, linear in the first variable (convention (Cv1) of
\S\ref{pb:ss-setting}); so $\rho(g)^*=\rho(g^{-1})$. For a vector $v$ fixed by $K$ put, as in (N3)
of \S\ref{pl:ss-norm},
\[
T_iv=\sum_{h\in Kd_iK/K}\rho(h)v\qquad(0\le i\le N),
\]
and for $y$ fixed by $I_1$ let $Uy=\sum_{h\in I_1dI_1/I_1}\rho(h)y$, the operator $U$ of
Proposition~\ref{pri:dict} and of \S\ref{pax:ss-dict}. Let $\payM\colon V\to V^{I_1}$ be the mean
over $I_1$, $\payM v=\int_{I_1}\rho(k)v\,dk$ for the Haar measure of total mass $1$; as $v$ is fixed
by an open subgroup, this is a finite average. $\payM$ is the identity on $V^{I_1}$, and it is
self-adjoint: $(\payM a,b)=\int(a,\rho(k^{-1})b)\,dk=(a,\payM b)$, the measure being invariant under
$k\mapsto k^{-1}$.

If $x\in V^K$ satisfies $T_ix=t_ix$ for all $i$, its \emph{Hecke polynomial} is
\begin{equation}\label{pay:eq-P}
P_x(X)=\sum_{i=0}^N(-1)^iq^{\binom i2}t_iX^{N-i}\qquad(t_0=1).
\end{equation}
For $x$ a nonzero spherical vector in an irreducible unramified $\pi$ this is the polynomial
$H_\pi$ of~\eqref{pax:eq-satake}; and on forms, with the operators $T_i$ of (N3), it is the Hecke polynomial $H_w$ of Lemma~\ref{pri:chain} (the polynomial~\eqref{pax:eq-satake}).

\emph{Lattices.} For $0\le j\le N$ let $\payL_j$ be the set of $\auO_F$-lattices $L$ with $\varpi
\auO_F^N\subseteq L\subseteq\auO_F^N$ and $\dim_\auF(\auO_F^N/L)=j$, and $\payL^{(1)}_j=\{L\in
\payL_j:e_1\in L\}$. For a lattice $L$ fix $g_L\in G$ with $g_L\auO_F^N=L$; then $g_LK$ depends only on
$L$, and $gK\mapsto g\auO_F^N$ is a bijection from $G/K$ onto the set of lattices.

\begin{lemma}\label{pay:lat}
$Kd_jK/K$ corresponds under this bijection to $\payL_j$. Hence for $v\in V^K$, $T_jv=\sum_{L\in
\payL_j}\rho(g_L)v$. Moreover $\payL_0=\{\auO_F^N\}$, $\payL_N=\{\varpi\auO_F^N\}$,
$\payL^{(1)}_N=\emptyset$, $\payL^{(1)}_0=\{\auO_F^N\}$ and $\payL^{(1)}_{N-1}=\{\auO_Fe_1+\varpi\auO_F^N
\}$.
\end{lemma}
\begin{proof}
$d_j\auO_F^N\in\payL_j$, and $K$ preserves $\payL_j$; so $Kd_jK/K$ maps into $\payL_j$. Conversely,
$L\mapsto L/\varpi\auO_F^N$ is a bijection from $\payL_j$ onto the subspaces of dimension $N-j$ of
$\auF^N$, on which $K$ acts transitively through $K\to\auGL_N(\auF)$ (surjective); so every $L\in
\payL_j$ is $kd_j\auO_F^N$ with $k\in K$. The last assertions are immediate ($e_1\notin\varpi
\auO_F^N$; a subspace of dimension $1$ containing $\bar e_1$ is $\auF\bar e_1$).
\end{proof}

For $x\in V^K$ and $0\le j\le N$ put
\begin{equation}\label{pay:eq-xj}
x_j=\sum_{L\in\payL^{(1)}_j}\rho(g_L)x .
\end{equation}
So $x_0=x$, $x_N=0$, and $x_{N-1}=\rho(\mathrm{diag}(1,\varpi,\dots,\varpi))x$. Each $x_j$ is fixed by
$I_1$: an element $k\in K$ with $ke_1\in\auO_F^\times e_1+\varpi\auO_F^N$ (in particular $k\in I_1$)
permutes $\payL^{(1)}_j$ (it preserves $\payL_j$, and $ke_1\in L$ iff $e_1\in L$ as $\varpi\auO_F^N
\subseteq L$), and $\rho(k)\rho(g_L)x=\rho(g_{kL})x$ since $x$ is fixed by $K$.

\subsection{Two counts of hyperplanes}\label{pay:ss-count}

Let $\payN\subseteq K$ be the group of matrices that are equal to $1$ except in the first row, whose
off-diagonal entries are arbitrary in $\auO_F$ (and the diagonal entry is $1$); $\payN\simeq
\auO_F^{N-1}$ and $\payN\subseteq I_1$. For $u\in\auF^{N-1}$ fix $n_u\in\payN$ with first row reducing
to $(1,u)$.

\begin{lemma}[Count of hyperplanes]\label{pay:count}
Let $\sigma$ be a representation of $G$ on a vector space $W$, $v\in W$ fixed by $\sigma(K)$, and
$\psi\colon W\to W'$ a linear map with $\psi\circ\sigma(n)=\psi$ for all $n\in\payN$. Put $v_j=
\sum_{L\in\payL^{(1)}_j}\sigma(g_L)v$ and $v^\flat_j=\sum_{L\in\payL_j}\sigma(g_L)v$. Then for $0\le
j\le N-1$,
\[
q^{N-1}\,\psi\bigl(\sigma(d)v_j\bigr)=q^j\,\psi\bigl(v^\flat_{j+1}-v_{j+1}\bigr).
\]
\end{lemma}
\begin{proof}
Let $H=\{w\in\auO_F^N:w_1\in\varpi\auO_F\}=d\auO_F^N$. \emph{Step 1:} $L\mapsto dL$ is a bijection from
$\payL^{(1)}_j$ onto $\{M'\in\payL_{j+1}:M'\subseteq H\}$. Indeed, for $L\in\payL^{(1)}_j$: $dL
\subseteq d\auO_F^N=H$; $dL$ contains $de_1=\varpi e_1$ and $d(\varpi\auO_F^N)\supseteq\varpi
\sum_{l\ge2}\auO_Fe_l$, so $dL\supseteq\varpi\auO_F^N$; and $[\auO_F^N:dL]=[\auO_F^N:d\auO_F^N][d\auO_F^N:
dL]=q\cdot q^j$. Conversely, if $M'\in\payL_{j+1}$ and $M'\subseteq H$, then $d^{-1}M'\subseteq
d^{-1}H=\auO_F^N$, $d^{-1}M'\supseteq d^{-1}\varpi\auO_F^N\ni e_1$, $d^{-1}M'\supseteq\varpi\auO_F^N$,
and its colength is $j$.

\emph{Step 2:} for $M\in\payL_{j+1}$, the number of $u\in\auF^{N-1}$ with $n_uM\subseteq H$ is $0$
if $e_1\in M$ and $q^j$ if $e_1\notin M$. Indeed $n_uM\in\payL_{j+1}$, and $n_uM\subseteq H$ iff the
linear form $\lambda_u=(1,u)$ on $\auF^N$ (the reduction of the first row of $n_u$) vanishes on
$\bar M=M/\varpi\auO_F^N\subseteq\auF^N$, a subspace of dimension $N-j-1$. Every linear form $\lambda$
with $\lambda(\bar e_1)=1$ is $\lambda_u$ for exactly one $u$. The annihilator $\bar M^\perp$ has
dimension $j+1$, and evaluation at $\bar e_1$ is a linear form on it, which is zero iff $\bar e_1\in
\bar M$, i.e.\ iff $e_1\in M$. If it is nonzero, its fibre over $1$ is a coset of a subspace of
dimension $j$, with $q^j$ elements.

\emph{Step 3.} By Steps 1 and 2, the map $(u,L)\mapsto(u,n_u^{-1}dL)$ is a bijection from $\auF^{N-1}
\times\payL^{(1)}_j$ onto the pairs $(u,M)$ with $M\in\payL_{j+1}$ and $n_uM\subseteq H$, and for
fixed $M$ there are $q^j[e_1\notin M]$ such $u$. As $v$ is fixed by $K$, $\sigma(n_u^{-1}dg_L)v=
\sigma(g_{n_u^{-1}dL})v$. Hence
\[
\sum_{u\in\auF^{N-1}}\sigma(n_u^{-1}d)\,v_j=\sum_{M\in\payL_{j+1},\ e_1\notin M}q^j\sigma(g_M)v=q^j
(v^\flat_{j+1}-v_{j+1}).
\]
Apply $\psi$: as $n_u^{-1}\in\payN$, each term on the left gives $\psi(\sigma(d)v_j)$, and there are
$q^{N-1}$ of them.
\end{proof}

\begin{lemma}[Duality of lattices]\label{pay:dual}
Let $\theta(g)={}^tg^{-1}$, an automorphism of $G$ preserving $K$, and $\rho^\theta=\rho\circ\theta$.
Let $x\in V^K$ (so $x$ is fixed by $\rho^\theta(K)$), $c_0=\varpi^{-1}1_N$ and $c_1=\mathrm{diag}(1,
\varpi^{-1},\dots,\varpi^{-1})$. If $i+j=N$ then $\sum_{L\in\payL_j}\rho^\theta(g_L)x=\rho(c_0)
\sum_{L\in\payL_i}\rho(g_L)x$; if $1+i+j=N$ then $\sum_{L\in\payL^{(1)}_j}\rho^\theta(g_L)x=\rho(c_1)
x_i$.
\end{lemma}
\begin{proof}
For a lattice $L$ let $L^*=\{w\in F^N:{}^twL\subseteq\auO_F\}$; then ${}^tg_L^{-1}\auO_F^N=L^*$,
$(L^*)^*=L$, and $L\subseteq L'$ iff $L'^*\subseteq L^*$, with $[L^*:L'^*]=[L':L]$. If $\varpi
\auO_F^N\subseteq L\subseteq\auO_F^N$ has colength $j$, then $\auO_F^N\subseteq L^*\subseteq\varpi^{-1}
\auO_F^N$ with $[L^*:\auO_F^N]=q^j$. \emph{First case.} $L\mapsto\varpi L^*=c_0^{-1}L^*$ maps $\payL_j$
to $\payL_i$ ($[\auO_F^N:\varpi L^*]=[\varpi^{-1}\auO_F^N:L^*]=q^{N-j}$), with inverse $L'\mapsto
(\varpi^{-1}L')^*$. \emph{Second case.} $e_1\in L$ iff $w_1={}^twe_1\in\auO_F$ for all $w\in L^*$. For
$L\in\payL^{(1)}_j$ put $L'=c_1^{-1}L^*$: its first coordinates are those of $L^*$, in $\auO_F$, and
its other coordinates are $\varpi$ times those of $L^*$, in $\auO_F$; so $L'\subseteq\auO_F^N$, and
$L'\supseteq c_1^{-1}\auO_F^N=\auO_Fe_1+\varpi\auO_F^N$, so $e_1\in L'$ and $\varpi\auO_F^N\subseteq
L'$; finally $[\auO_F^N:L']=[\auO_F^N:c_1^{-1}\auO_F^N]/[L^*:\auO_F^N]=q^{N-1-j}=q^i$. Conversely, for
$L'\in\payL^{(1)}_i$, $c_1L'\supseteq c_1(\auO_Fe_1+\varpi\auO_F^N)\supseteq\auO_F^N$, $c_1L'\subseteq
c_1\auO_F^N$ has integral first coordinates and lies in $\varpi^{-1}\auO_F^N$, so $(c_1L')^*\in
\payL^{(1)}_j$ by the same computations; the two maps are inverse. In both cases ${}^tg_L^{-1}K=c_m
g_{L'}K$ (both send $\auO_F^N$ to $L^*=c_mL'$), and as $x$ is fixed by $K$, $\rho^\theta(g_L)x=\rho(
c_mg_{L'})x$. Summing gives the two identities.
\end{proof}

\subsection{The recursions}\label{pay:ss-rec}

Define endomorphisms of $V$
\[
\payU=q^{N-1}\,\payM\,\rho(d)\,\payM,\qquad\payUd=q^{N-1}\,\payM\,\rho(d^{-1})\,\payM .
\]
They are mutually adjoint: $(\payU a,b)=(a,\payUd b)$ and $(\payUd a,b)=(a,\payU b)$ for all $a,b\in V$,
since $\payM$ is self-adjoint and $\rho(d)^*=\rho(d^{-1})$.

\begin{lemma}\label{pay:UU}
On $V^{I_1}$, $U=\payU$; in particular $U$ preserves $V^{I_1}$.
\end{lemma}
\begin{proof}
Let $y\in V^{I_1}$ and $I'=I_1\cap dI_1d^{-1}$. The map $kI'\mapsto kdI_1$ is a bijection from
$I_1/I'$ onto $I_1dI_1/I_1$ ($kdI_1=k'dI_1$ iff $k^{-1}k'\in dI_1d^{-1}$), and the number of cosets is
$q^{N-1}$ (the decomposition $I_1dI_1=\bigsqcup_{u}u\,d\,I_1$, $u\in N(\auO_F)/dN(\auO_F)d^{-1}$, in
the proof of Lemma~\ref{pax:Ueig}). The vector $\rho(d)y$ is fixed by $dI_1d^{-1}\supseteq I'$, so
$k\mapsto\rho(kd)y$ is constant on the cosets $kI'$, each of measure $q^{1-N}$. Hence $\payM\rho(d)y=
q^{1-N}\sum_{k\in I_1/I'}\rho(kd)y=q^{1-N}Uy$, and $\payU y=q^{N-1}\payM\rho(d)\payM y=Uy$.
\end{proof}

\begin{proposition}[Recursions]\label{pay:rec}
Write $N=n+1$. Let $x\in V^K$ with $T_ix=t_ix$ ($0\le i\le N$), and $x_j$ as in~\eqref{pay:eq-xj}.
Then $t_0=1$ if $x\ne0$, and:
\begin{enumerate}
\item[(R1)] $\payU x_j=q^j\,(t_{j+1}x-x_{j+1})$ for $0\le j\le n$;
\item[(R2)] $\payUd x_{i+1}=q^{n-i-1}\bigl((t_{i+1}/q^n)\,\payUd x-x_i\bigr)$ for $0\le i<n$;
\item[(R3)] if $x\ne0$ then $t_N\ne0$, $\payUd x=q^nt_N^{-1}x_n$, and $\payU x_n=q^nt_Nx$.
\end{enumerate}
\end{proposition}
\begin{proof}
$T_0$ is the identity on $V^K$, so $t_0=1$ when $x\ne0$. As each $x_j$ and $x$ lie in $V^{I_1}$,
$\payM x_j=x_j$.

(R1) Apply Lemma~\ref{pay:count} with $\sigma=\rho$, $v=x$, $\psi=\payM$ (invariant under $\rho(I_1)
\supseteq\rho(\payN)$): $v_j=x_j$ and $v^\flat_{j+1}=T_{j+1}x=t_{j+1}x$ by Lemma~\ref{pay:lat}, so
$q^n\payM\rho(d)x_j=q^j(t_{j+1}x-x_{j+1})$, and the left side is $\payU x_j$.

(R2) Apply Lemma~\ref{pay:count} with $\sigma=\rho^\theta$, $v=x$, and $\psi=\payM\circ\rho(c_1^{-1})$.
The invariance: for $n'\in\payN$ with first row $(1,a)$, $\theta(n')$ is unipotent lower triangular with
entries $-a_l$ at $(l,1)$, $l\ge2$, and $0$ elsewhere off the diagonal; so $c_1^{-1}\theta(n')c_1$ has
entries $-\varpi a_l\in\varpi\auO_F$ at $(l,1)$ and lies in $I_1$, whence $\psi(\rho^\theta(n')w)=
\payM\rho(c_1^{-1}\theta(n')c_1)\rho(c_1^{-1})w=\psi(w)$. Note $\rho^\theta(d)=\rho(d^{-1})$. Fix
$0\le i<n$ and let $j=n-1-i\in[0,N-1]$. By Lemma~\ref{pay:dual}, $v_j=\rho(c_1)x_{i+1}$ (as $1+(i+1)+j=N$), $v^\flat_{j+1}=\rho(c_0)t_{i+1}x$ (as
$(i+1)+(j+1)=N$), and $v_{j+1}=\rho(c_1)x_i$ (as $1+i+(j+1)=N$). Lemma~\ref{pay:count} gives
\[
q^n\payM\rho(c_1^{-1}d^{-1}c_1)x_{i+1}=q^j\bigl(t_{i+1}\payM\rho(c_1^{-1}c_0)x-\payM x_i\bigr).
\]
Diagonal matrices commute and $c_1^{-1}c_0=d^{-1}$, so this reads $\payUd x_{i+1}=q^{n-1-i}(t_{i+1}
q^{-n}\payUd x-x_i)$, using $\payM\rho(d^{-1})x=q^{-n}\payUd x$.

(R3) By Lemma~\ref{pay:lat}, $T_Nx=\rho(\varpi1_N)x$ and $x_n=\rho(\mathrm{diag}(1,\varpi,\dots,
\varpi))x=\rho(d^{-1})\rho(\varpi1_N)x=t_N\rho(d^{-1})x$. As $\rho(\varpi1_N)$ is invertible,
$t_Nx=T_Nx\ne0$, so $t_N\ne0$. Then $\payUd x=q^n\payM\rho(d^{-1})x=q^nt_N^{-1}\payM x_n=q^nt_N^{-1}x_n$,
and $\payU x_n=q^n\payM\rho(d)x_n=q^nt_N\payM x=q^nt_Nx$.
\end{proof}

\begin{remark}\label{pay:rem-rec}
With the convention $x_N=0$, (R1) at $j=n$ is the last identity of (R3), and (R2) also holds for
$i=n$: by (R3) the right side is $q^{-1}(t_Nq^{-n}\cdot q^nt_N^{-1}x_n-x_n)=0=\payUd x_N$. We use (R2)
for $0\le i\le n$ in this extended form.
\end{remark}

\subsection{Linear algebra}\label{pay:ss-lin}

For $A=\sum_ig_iX^i\in\C[X]$ let $\bar A=\sum_i\bar g_iX^i$. If $U,U'$ are endomorphisms of a
hermitian space $(V,(\ ,\ ))$ (form linear in the first variable) with $(Ua,b)=(a,U'b)$ and $(U'a,b)=
(a,Ub)$ for all $a,b$, then by induction on the degree
\begin{equation}\label{pay:eq-adj}
(a,A(U)b)=(\bar A(U')a,b)\qquad(a,b\in V,\ A\in\C[X]).
\end{equation}

\begin{proposition}\label{pay:lin}
Let $(V,(\ ,\ ))$ be a complex vector space with a positive definite hermitian form, $U,U'$ mutually
adjoint endomorphisms as above, $q>1$ a real number and $n\ge0$ an integer. Let $x_0,\dots,x_{n+1}\in
V$ with $x:=x_0\ne0$ and $x_{n+1}=0$, and $t_0=1,t_1,\dots,t_{n+1}\in\C$ with $t_{n+1}\ne0$, such that
\begin{enumerate}
\item[(R1)] $Ux_j=q^j(t_{j+1}x-x_{j+1})$ for $0\le j\le n$;
\item[(R2)] $U'x_{i+1}=q^{n-i-1}\bigl((t_{i+1}/q^n)U'x-x_i\bigr)$ for $0\le i\le n$;
\item[(R3)] $U'x=q^nt_{n+1}^{-1}x_n$.
\end{enumerate}
Let $P=\sum_{j=0}^{n+1}(-1)^jq^{\binom j2}t_jX^{n+1-j}$. Then:
\begin{enumerate}
\item[(a)] if $P(c)=0$ and $P(qc)\ne0$, then $|c|^2>q^{n-1}$;
\item[(b)] if $\delta\in\C^\times$, $k\ge0$, $R\in\C[X]$, $P=\prod_{i<k}(\delta X-q^i)\cdot R$ and $P(q^k/\delta)=0$,
then $\bar R(U')x=0$.
\end{enumerate}
\end{proposition}

\begin{proof}
\emph{Step 1.} For $0\le m\le n+1$ let $p_m=\sum_{j\le m}(-1)^jq^{\binom j2}t_jX^{m-j}$, so $p_0=1$,
$p_{n+1}=P$ and $p_{m+1}=Xp_m+\kappa_{m+1}$ with $\kappa_{m+1}=(-1)^{m+1}q^{\binom{m+1}2}t_{m+1}$. Then
\begin{equation}\label{pay:eq-pm}
p_m(U)x=(-1)^mq^{\binom m2}x_m\qquad(0\le m\le n+1).
\end{equation}
For $m=0$ this is $x=x_0$; if it holds for $m\le n$, then by (R1)
$p_{m+1}(U)x=(-1)^mq^{\binom m2}q^m(t_{m+1}x-x_{m+1})+\kappa_{m+1}x=(-1)^{m+1}q^{\binom{m+1}2}x_{m+1}$,
as $\binom m2+m=\binom{m+1}2$. In particular $P(U)x=\pm q^{\binom{n+1}2}x_{n+1}=0$.

\emph{Step 2.} Fix $c\in\C$ and let $D_m=(p_m-p_m(c))/(X-c)\in\C[X]$ ($0\le m\le n+1$), and $D_{n+2}=
cD_{n+1}+P$. Then $D_0=0$ and $D_{m+1}=cD_m+p_m$ for $0\le m\le n+1$ (indeed $(X-c)(cD_m+p_m)=c(p_m-
p_m(c))+(X-c)p_m=Xp_m-cp_m(c)=p_{m+1}-p_{m+1}(c)$). We claim, for $0\le m\le n+1$,
\begin{equation}\label{pay:eq-Dm}
U'\bigl(D_{m+1}(U)x\bigr)=\frac{p_m(qc)}{q^m}\,U'x+q^{n-1}D_m(U)x .
\end{equation}
For $m=0$ both sides are $U'x$. Assume~\eqref{pay:eq-Dm} for some $m\le n$. Then $D_{m+2}(U)x=c\,
D_{m+1}(U)x+p_{m+1}(U)x$, and by~\eqref{pay:eq-pm}, (R2) with $i=m$, and~\eqref{pay:eq-pm} again,
\[
U'p_{m+1}(U)x=(-1)^{m+1}q^{\binom{m+1}2}q^{n-m-1}\Bigl(\frac{t_{m+1}}{q^n}U'x-x_m\Bigr)=
\frac{\kappa_{m+1}}{q^{m+1}}U'x+q^{n-1}p_m(U)x .
\]
Adding $c$ times~\eqref{pay:eq-Dm}, the coefficient of $U'x$ is $(qc\,p_m(qc)+\kappa_{m+1})/q^{m+1}=
p_{m+1}(qc)/q^{m+1}$, and the rest is $q^{n-1}(cD_m+p_m)(U)x=q^{n-1}D_{m+1}(U)x$. This
proves~\eqref{pay:eq-Dm} for $m+1$.

Put $z_c=D_{n+1}(U)x$. As $P(U)x=0$, \eqref{pay:eq-Dm} for $m=n+1$ reads
\begin{equation}\label{pay:eq-star}
c\,U'z_c=\frac{P(qc)}{q^{n+1}}\,U'x+q^{n-1}z_c .
\end{equation}

\emph{Step 3.} If $P(c)=0$ then $(U-c)z_c=(P-P(c))(U)x=P(U)x=0$, i.e.\ $Uz_c=cz_c$. By (R3) and (R1)
with $j=n$ ($x_{n+1}=0$), $UU'x=q^nt_{n+1}^{-1}\cdot q^nt_{n+1}x=q^{2n}x$, hence
$(U'a,U'x)=(a,UU'x)=q^{2n}(a,x)$ for all $a$. Pair~\eqref{pay:eq-star} with $U'x$ in the second
variable: the left side gives $c\,q^{2n}(z_c,x)$; on the right, $(U'x,U'x)=q^{2n}(x,x)$ and $(z_c,U'x)=
(Uz_c,x)=c(z_c,x)$. Hence, for every root $c$ of $P$,
\begin{equation}\label{pay:eq-E}
c\,(z_c,x)\,(q^{n+1}-1)=P(qc)\,(x,x).
\end{equation}

\emph{Step 4: proof of (a).} Let $P(c)=0$, $a=P(qc)\ne0$. By~\eqref{pay:eq-E}, $c\ne0$ and $z_c\ne0$.
Pair~\eqref{pay:eq-star} with $z_c$ in the second variable: the left side is $c(U'z_c,z_c)=c(z_c,Uz_c)
=|c|^2(z_c,z_c)$; on the right, $(U'x,z_c)=(x,Uz_c)=\bar c(x,z_c)$, and by the complex conjugate
of~\eqref{pay:eq-E}, $\bar c(x,z_c)=\bar a(x,x)/(q^{n+1}-1)$. Therefore
\[
\bigl(|c|^2-q^{n-1}\bigr)(z_c,z_c)=\frac{|a|^2\,(x,x)}{q^{n+1}(q^{n+1}-1)}>0 ,
\]
and $|c|^2>q^{n-1}$.

\emph{Step 5: two consecutive roots.} Let $c\ne0$ with $P(c)=P(qc)=0$. We show $(P/(X-c))(U)x=0$.
Here $P/(X-c)=D_{n+1}$, so the vector is $z_c$. By~\eqref{pay:eq-star}, $U'z_c=\mu z_c$ with $\mu=
q^{n-1}/c$, and by~\eqref{pay:eq-E}, $(z_c,x)=0$. For every $A\in\C[X]$, by~\eqref{pay:eq-adj},
$(z_c,A(U)x)=(\bar A(U')z_c,x)=\bar A(\mu)(z_c,x)=0$. With $A=D_{n+1}$ this gives $(z_c,z_c)=0$, so
$z_c=0$.

\emph{Step 6: cancellation of a chain.} Under the hypotheses of (b) we show $R(U)x=0$, by induction
on $k$. For $k=0$, $R=P$ and $P(U)x=0$. Let $k\ge0$ and $P=\prod_{i<k+1}(\delta X-q^i)\cdot R$ with
$P(q^{k+1}/\delta)=0$. Then $P=\prod_{i<k}(\delta X-q^i)\cdot R'$ with $R'=(\delta X-q^k)R$, and $P(q^k/\delta)=0$; by the
induction hypothesis $R'(U)x=0$, i.e.\ $UR(U)x=c\,R(U)x$ with $c=q^k/\delta\ne0$. As $P(c)=0=P(qc)$, Step 5
gives $(P/(X-c))(U)x=0$; but $P/(X-c)=\delta\,Q\,R$ with $Q=\prod_{i<k}(\delta X-q^i)$, so $0=Q(U)R(U)x=Q(c)R(U)x$,
and $Q(c)=\prod_{i<k}(q^k-q^i)\ne0$ as $q>1$. So $R(U)x=0$.

\emph{Step 7: proof of (b).} By Step 6, $R(U)x=0$, hence by~\eqref{pay:eq-pm} $R(U)x_j=\pm q^{-\binom
j2}p_j(U)R(U)x=0$ for $0\le j\le n+1$. Let $\mathcal S=\mathrm{span}(x_0,\dots,x_n)$. By (R3), $U'x_0\in \mathcal S$; by
(R2), $U'x_{i+1}\in\C U'x+\C x_i\subseteq \mathcal S$ for $i<n$. So $\mathcal S$ is $U'$-stable, and $B(U')x\in \mathcal S$ for
every $B\in\C[X]$; thus $R(U)B(U')x=0$. With $B=\bar R$ and~\eqref{pay:eq-adj},
$(\bar R(U')x,\bar R(U')x)=(x,R(U)\bar R(U')x)=0$, so $\bar R(U')x=0$.
\end{proof}

\subsection{The chain theorem}\label{pay:ss-thm}

\begin{theorem}[Spherical chain lemma]\label{pay:chain}
Let $(\rho,V)$ be as in \S\ref{pay:ss-set}, $x\in V^K$ nonzero with $T_ix=t_ix$ for $0\le i\le N$, and
$P=P_x$ its Hecke polynomial~\eqref{pay:eq-P}.
\begin{enumerate}
\item[(i)] If $P(c)=0$ and $P(qc)\ne0$, then $|c|^2>q^{N-2}$.
\item[(ii)] Let $\delta\in\C^\times$, $k\ge0$, $R\in\C[X]$, and suppose $P(\delta^{-1}X)=\prod_{i\le k}(X-q^i)
\cdot S(X)$ with $S$ dividing $R$. Then for every $y\in V^{I_1}$,
\[
\bigl(x,\ \Phi(U)\,y\bigr)=0,\qquad\Phi(X)=\bigl((X-q^k)R\bigr)(\delta X).
\]
\end{enumerate}
\end{theorem}

\begin{proof}
Write $N=n+1$ and let $x_j$ be as in~\eqref{pay:eq-xj}. By Proposition~\ref{pay:rec},
Remark~\ref{pay:rem-rec} and the adjointness of $\payU,\payUd$ (\S\ref{pay:ss-rec}), the hypotheses of
Proposition~\ref{pay:lin} hold for $(U,U')=(\payU,\payUd)$, the $x_j$ and the $t_j$; and the polynomial
$P$ there is $P_x$. (i) is Proposition~\ref{pay:lin}(a), as $q^{n-1}=q^{N-2}$.

(ii) Substituting $\delta X$ for $X$, $P(X)=\prod_{i\le k}(\delta X-q^i)\cdot S(\delta X)$. Write $R=SB$ and $A(X)=(\delta X-
q^k)S(\delta X)$; then $P=\prod_{i<k}(\delta X-q^i)\cdot A$ and $P(q^k/\delta)=0$. By Proposition~\ref{pay:lin}(b),
$\bar A(\payUd)x=0$. Now $\Phi(X)=B(\delta X)\,A(X)$. By Lemma~\ref{pay:UU}, $U$ preserves $V^{I_1}$ and equals
$\payU$ there, so $\Phi(U)y=\Phi(\payU)y$, and by~\eqref{pay:eq-adj}
\[
(x,\Phi(\payU)y)=\bigl(\bar\Phi(\payUd)x,y\bigr)=\bigl(\bar B(\bar \delta\,\payUd)\,\bar A(\payUd)x,\ y\bigr)=0 .
\qedhere
\]
\end{proof}

\begin{lemma}[Symmetry of the roots]\label{pay:sym}
In the situation of Theorem~\ref{pay:chain}, if $y,y'\in\C$ satisfy $y\,\bar y'=q^{N-1}$ and
$P(y')=0$, then $P(y)=0$.
\end{lemma}
\begin{proof}
Let $dg$ be the Haar measure on $G$ with $K$ of measure $1$; for $v\in V^K$, $T_iv=\int_{Kd_iK}\rho(g)v
\,dg$. For $v,v'\in V^K$, as $G$ is unimodular,
$(T_iv,v')=\int_{Kd_iK}(v,\rho(g^{-1})v')\,dg=\bigl(v,\int_{Kd_i^{-1}K}\rho(g)v'\,dg\bigr)$. Now
$d_i^{-1}=\varpi^{-1}\mathrm{diag}(1_i,\varpi1_{N-i})$ and $\mathrm{diag}(1_i,\varpi1_{N-i})\in
Kd_{N-i}K$ (conjugate by a permutation matrix), so $Kd_i^{-1}K=\varpi^{-1}Kd_{N-i}K$ and
$\int_{Kd_i^{-1}K}\rho(g)v'\,dg=\rho(\varpi^{-1}1_N)T_{N-i}v'$. With $v=v'=x$ and $z=(\rho(\varpi1_N)x,
x)$ this gives, for $0\le i\le N$,
\begin{equation}\label{pay:eq-ti}
t_i\,(x,x)=(x,\rho(\varpi^{-1}1_N)t_{N-i}x)=\bar t_{N-i}\,z .
\end{equation}
Now multiply $P(y)$ by $(x,x)$ and use~\eqref{pay:eq-ti} termwise:
$P(y)(x,x)=z\sum_i(-1)^iq^{\binom i2}\bar t_{N-i}y^{N-i}$. On the other hand, the term of index $b=N-i$
of $\overline{P(y')}$, multiplied by $y^N$, is
$(-1)^{N-i}q^{\binom{N-i}2}\bar t_{N-i}\,\bar y'^{\,i}y^N=(-1)^{N-i}q^{\binom{N-i}2+(N-1)i}\bar t_{N-i}
y^{N-i}$, using $(y\bar y')^i=q^{(N-1)i}$; and $\binom{N-i}2+(N-1)i=\binom N2+\binom i2$. Hence
\[
q^{\binom N2}\,(x,x)\,P(y)=(-1)^N\,z\,y^N\,\overline{P(y')} ,
\]
and $P(y')=0$ forces $P(y)=0$, as $(x,x)>0$.
\end{proof}

\begin{corollary}[Chains of roots]\label{pay:roots}
In the situation of Theorem~\ref{pay:chain}, let $c$ be a root of $P$ and $k\ge1$ an integer with
$|c|^2q^k=q^N$. Then $c,cq,\dots,cq^{k-1}$ are roots of $P$, and $\prod_{i<k}(X-q^i)$ divides $P(cX)$.
\end{corollary}
\begin{proof}
We show that $cq^i$ is a root for $0\le i<k$, by strong induction on $i$; write $k=i+1+j$ with $j\ge0$.
If $j<i$: $cq^j$ is a root by induction, and $(cq^i)\overline{(cq^j)}=|c|^2q^{i+j}=q^{N-k}q^{k-1}=q^{N-1}$,
so $cq^i$ is a root by Lemma~\ref{pay:sym}. If $j\ge i$ and $i=0$, $c$ is a root. If $j\ge i\ge1$ and
$cq^i$ were not a root, then $c'=cq^{i-1}$ is a root (induction) with $qc'$ not a root, so
$|c'|^2>q^{N-2}$ by Theorem~\ref{pay:chain}(i), i.e.\ $q^{N-k+2i-2}>q^{N-2}$, i.e.\ $2i>k=i+1+j$, so
$j<i-1$, a contradiction. As $c\ne0$ ($|c|^2=q^{N-k}$) and $q>1$, the $k$ roots $cq^i$ are distinct;
so $\prod_{i<k}(X-cq^i)$ divides $P(X)$, and $\prod_{i<k}(X-q^i)$ divides $P(cX)$.
\end{proof}

\subsection{Consequences for \S\ref{pax:sec}}\label{pay:ss-cons}

\begin{proof}[Proof of Lemma~\ref{pax:tadic} without the classification of the unitary dual]
Let $\pi$ be an irreducible unitary unramified representation of $\auGL_N(F)$, with invariant inner
product, and $x\ne0$ a spherical vector; $\pi^K$ is stable under the $T_i$, and $H_\pi$
in~\eqref{pax:eq-satake} is defined by their eigenvalues $t_i$ on $\pi^K$, so $T_ix=t_ix$ and $H_\pi=
P_x$ (both are the polynomial~\eqref{pay:eq-P}). Corollary~\ref{pay:roots} with $c$ a root of $H_\pi$
and $|c|^2q^k=q^N$ is the statement of Lemma~\ref{pax:tadic}.
\end{proof}

So the proof of Lemma~\ref{pri:chain} in \S\ref{pax:ss-chain} holds with Lemma~\ref{pax:tadic} proved
here: the input from Tadi\'c's classification, the discussion of unramified constituents and the
paragraph ``The chain'' of the proof in \S\ref{pax:ss-chain} are not needed. (The step ``Unitarity'' of the proof of Lemma~\ref{pri:chain} may even be shortened:
Corollary~\ref{pay:roots} applies to any smooth representation with a positive definite invariant
hermitian form, so it can be applied directly to the representation of $\auU(V)(E^+_v)\simeq
\auGL_n(E_w)$ on the forms of the given weight, with the global inner product, and to $x=y$, without
passing to an irreducible constituent.)

\begin{proof}[Proof of Lemma~\ref{pax:toproot}]
Let $\pi$, $x$, $X_0$, $k\ge1$, $S\mid R$ be as there, so $H_\pi(X_0X)=\prod_{i<k}(X-q^i)\cdot S(X)$, and
$H_\pi=P_x$ as in the previous proof. Apply Theorem~\ref{pay:chain}(ii) with $\delta=X_0^{-1}$ and $k-1$ in
place of $k$: then $P(\delta^{-1}X)=\prod_{i\le k-1}(X-q^i)\cdot S$, and $\Phi(X)=((X-q^{k-1})R)(X_0^{-1}X)$,
so $(x,((X-q^{k-1})R)(X_0^{-1}U)z)=0$ for every $z\in\pi^{I_1}$. (Orthogonality does not depend on the
order of the arguments of the inner product.)
\end{proof}

With this, the proof of Proposition~\ref{pri:dict}(d) given in \S\ref{pax:ss-dict} from
Lemma~\ref{pax:toproot} is complete, and Proposition~\ref{pri:dict} is proved in full (modulo the
published results cited in \S\ref{pax:ss-dict}).

\subsection{The two facts of Remark~\ref{pax:rem-BL}}\label{pay:ss-BL}

\begin{lemma}[$\eta_O$ is unramified away from $\ell$]\label{pay:eta}
Let $\eta$ be the algebraic Hecke character $y\mapsto(\chi_W(y/\bar y)\allowbreak \upsilon(y)^{-r'}\allowbreak |y|^{-r'/2})^{-1}$
of Proposition~\ref{pga:trace}, whose $\ell$-adic character is $\eta_O$. Then $\eta$ is trivial on
$\auO_{E,w}^\times$ for every finite place $w$ of $E$. Hence $\eta_O$ is unramified at every finite
place $w\nmid\ell$; in particular at the places of $Q\cup\overline Q$, at the places of $U\cup
\overline U$ in Proposition~\ref{pri:BL}(c), and at the chosen places not above $\ell$.
\end{lemma}
\begin{proof}
By \S\ref{pga:ss-L} and~\eqref{po:eq-w5}, $\chi_W=a=\upsilon^{r_-}$ is a power of the splitting
character $\upsilon$. By \S\ref{pk:ss-conv}, $\upsilon$ is unramified at every finite place
($\upsilon(x)=\prod_\tau x_\tau/|x_\tau|$ when the finite part of $x$ is a unit). Let $y\in
\auO_{E,w}^\times$ (embedded in $\mathbb A_E^\times$ with components $1$ elsewhere). Then $y/\bar y$ has
unit components at $w$ and $\bar w$ and $1$ elsewhere, so $\chi_W(y/\bar y)=1$; $\upsilon(y)=1$; and
$|y|=1$. So $\eta(y)=1$. For an algebraic Hecke character, the restriction of its $\ell$-adic character
to the inertia group at $w\nmid\ell$ is, through $\mathrm{Art}_w$, the restriction of $\eta_w$ to
$\auO_{E,w}^\times$ (the algebraic part only involves the components above $\ell$ and $\infty$); hence
$\eta_O|_{I_w}=1$. The places of $Q\cup\overline Q$ and of $U\cup\overline U$ do not lie above $\ell$
(they are disjoint from $\pgS(\ell)$, which contains the places above $\ell$, Definition~\ref{pri:def-aux}
and Proposition~\ref{pri:BL}(c)). The chosen places used in the proof of Proposition~\ref{pri:BL}(c) are those prime to $\ell$: if $\ell\in\{p_0,p_1\}$, the chosen places above $\ell$ are excluded, and at them the lemma asserts nothing.
\end{proof}

So (F1) of \S\ref{pax:ss-dict} holds at every finite place not above $\ell$, which is the first point
of Remark~\ref{pax:rem-BL}; it is used, for chosen $w\nmid\ell$, in the paragraph ``Inertia at a chosen $w$'' of the proof of
Proposition~\ref{pri:BL}(c). The rescaling of the label is by the tame part of the character of the
pair only.

\emph{Product labels.} The second point of Remark~\ref{pax:rem-BL} concerns how the group $\auG_O=
\prod_w\auGL_{r'}(\auF_w)$ of the chosen places, which acts on forms on $\auU(W)$ by right
translation (twisted by the tame part of $\mu$, as in Proposition~\ref{au:equiv}), acts on forms on
$\auU(V)$, $\operatorname{rk}V=n=r'+t'$. The following lemma shows that the single lift of
\S\ref{po:ss-set} intertwines it with right translation by elements reducing to $\mathrm{diag}(u_w,
1_{t'})$; this is the ``product label'' $c_w\cdot(X-1)^{t'}$ used in the proof of
Proposition~\ref{pri:BL}(c).

\begin{lemma}[Equivariance of the single lift]\label{pay:product}
Let $u=(u_w)\in\auG_O$. Choose $k\in K_W$ with reduction $u_w$ at each chosen place $w$ and $1$ at
the places of $Q$, and $k'\in K_V$ with reduction $\mathrm{diag}(u_w,1_{t'})$ at each chosen place
and $1$ at the places of $Q$, and with components $k_v=1$, $k'_v=1$ at every place $v$ that is neither chosen nor in $Q$ (possible as the chosen places are split and reduction is onto). Let
$\Theta$ be the lift of \S\ref{po:ss-set} with the block datum $\Phi_\vartheta$. Then for every form
$F$ on $\auU(W)$ and every $h$,
\[
(\Theta F)(hk')=\prod_{w\ \mathrm{chosen}}\psi_w(\det u_w)^{-1}\cdot\bigl(\Theta(F(\,\cdot\,k))\bigr)(h).
\]
\end{lemma}
\begin{proof}
At a chosen place $w$, by~\eqref{po:eq-w5} (with $\eta_w$ unramified, hence trivial on $\auGL_{r'}(
\auO_w)\times\auGL_n(\auO_w)$, and $\chi_w(\det k_w)=\psi_w(\det u_w)$),
$\omega_\mu(k_w,k'_w)\Phi_{\vartheta,w}(T)=\psi_w(\det u_w)^{-1}\plone[k_w'^{-1}Tk_w\equiv\iota_0]=
\psi_w(\det u_w)^{-1}\plone[T\equiv k'_w\iota_0k_w^{-1}]$, and $k'_w\iota_0k_w^{-1}\equiv\iota_0$ since
$\mathrm{diag}(u_w,1_{t'})\iota_0=\iota_0u_w$. At the places of $Q$, $k$, $k'$ reduce to $1$, and
$\plone[T\text{ integral}]\psi_w(\det T'\bmod w)$ is invariant under $(k,k')$ with $k\equiv1$, $k'\equiv1$
(the upper block of $k'^{-1}Tk$ is congruent to $T'$), and the factor $\chi_w(\det k_w)^{-1}$ of~\eqref{po:eq-w5} equals $1$, since $\chi_w$ is tamely ramified at $w\in Q$ (its restriction to $\auO_w^\times$ factors through $\psi_w$, i.e.\ through $\auO_w^\times\to(\auO_w/w)^\times$) and $\det k_w\equiv1\bmod w$. At all other places $k_v=k'_v=1$, so nothing changes there. Hence $\omega_\mu(k,
k')\Phi_\vartheta=\prod_w\psi_w(\det u_w)^{-1}\Phi_\vartheta$, i.e.\ $\omega_\mu(1,k')\Phi_\vartheta=
\prod_w\psi_w(\det u_w)^{-1}\omega_\mu(k^{-1},1)\Phi_\vartheta$. Insert this in $(\Theta F)(hk')=\int
\theta(\omega_\mu(g,hk')\Phi_\vartheta)F(g)\,dg$ and change variables $g\mapsto gk$ (Haar measure).
\end{proof}

Consequently, if the action of $\auG_O$ on forms on $\auU(V)$ is \emph{defined} by right translation
by such $k'$ (twisted by $\prod_w\psi_w(\det u_w)$), then $\Theta$ is $\auG_O$-equivariant, which is
the last sentence of Proposition~\ref{pga:L}, and the image of $e_c$-forms consists of $e_c$-forms
for this action; this is the action $\varrho$ of Definition~\ref{prm:def-level} used in
\S\ref{pax:ss-BL}, with the trivial block $\mathbf 1_{\auGL_{t'}}$ in the second position.

\begin{remark}[The base level at the chosen places]\label{pay:rem-level}
The proof of Proposition~\ref{pri:BL}(c), paragraph ``Inertia at a chosen $w$'', also uses a property
of the base level at the chosen places. The level is fixed in
Definition~\ref{prm:def-level} ($k\equiv\left(\begin{smallmatrix}1_{r'}&*\\0&*\end{smallmatrix}\right)$
modulo a chosen $w$), the lift lands in it (Lemma~\ref{prm:level}), and the property is
Lemma~\ref{prm:label}; see Corollary~\ref{prm:inertia}.
\end{remark}


%% ===== proofs/C-ribet2.tex =====
% proofs/C-ribet2.tex
% Completion of Section pri:sec: the paragraph "The characters" of
% Proposition pri:ribet, the induction in Proposition pri:norm(b), the base level at the chosen
% places (Remark pay:rem-level), and where the remaining inputs are proved.
% Labels of this file start with prm:.
\providecommand{\auO}{\mathcal O}
\providecommand{\auGL}{\mathrm{GL}}
\providecommand{\auF}{\mathbb F}
\providecommand{\auU}{\mathrm U}
\providecommand{\pgG}{G_E}
\providecommand{\pgtr}{\operatorname{tr}}
\providecommand{\priFitt}{\operatorname{Fitt}}
\providecommand{\prival}{\operatorname{val}}
\providecommand{\priAug}{\mathfrak I}
\providecommand{\priT}{\mathbb T}
\providecommand{\priX}{\mathcal X}
\providecommand{\priJ}{\mathcal J}
\providecommand{\prid}{\mathcal D}
\providecommand{\priN}{\mathcal N}
\providecommand{\prity}{\mathrm{ty}}
\providecommand{\plone}{\mathbf 1}
% macros of this file (prefix \prm)
\providecommand{\prmL}{\Lambda}
\providecommand{\prmT}{\mathfrak T}
\providecommand{\prmK}{\mathfrak K}
\providecommand{\prmval}{\operatorname{vl}}
\providecommand{\prmgen}{\mathrm{gen}}
\providecommand{\prmn}{\mathfrak n}
\providecommand{\prmq}{\mathfrak q}

\section{Ribet's method: characters, the norm induction and the base level}\label{prm:sec}

This section proves the following steps of \S\ref{pri:sec}: the paragraph ``The characters'' of the
proof of Proposition~\ref{pri:ribet}, together with the gluing that follows it
(\S\ref{prm:ss-chars}); the induction in Proposition~\ref{pri:norm}(b) (\S\ref{prm:ss-norm}); and the
property of the base level at the chosen places that the proof of Proposition~\ref{pri:BL}(c) uses,
with the definition of that level (\S\ref{prm:ss-level}). \S\ref{prm:ss-blocks} lists the statements
of \S\S\ref{pri:sec}, \ref{pga:sec}, \ref{pse:sec} and \ref{pg:sec} whose proofs are given in later
sections, with the place where each is proved.

\subsection{The characters in Ribet's method}\label{prm:ss-chars}

We keep the hypotheses (1)--(7) of Proposition~\ref{pri:ribet} and the notation of its proof:
$\prmgen=(\delta_v)_v\sqcup s_0\sqcup s$ is the finite family of generators, $z_g\in\auO^{\prmgen}$
($g\in G$) are coordinates with $u_x(g)-1=\sum_iz_{g,i}(u_x(\prmgen_i)-1)$ for all $x\in Y$ and
$g\in G$, $c_\zeta=(\zeta_\pi-1)^{-1}$, $\nu_0=c_\zeta(\gamma_\pi-1)$, $\rho_0=\nu_0+(1-\nu_0)(a-1)$. Put
$\prmL=\auO[X_J]$. Each $u_x$ extends $\prmL$-linearly to a $\prmL$-algebra map $u_x\colon\prmL[G]\to
M_\sigma(\prmL)$, and each character $\xi$ of $G$ to a $\prmL$-algebra map $\eta_\xi\colon\prmL[G]\to
\prmL$; we write $\eta_e=\eta_{\xi_e}$. Let $\priAug\subseteq\prmL[G]$ be the augmentation ideal (the
kernel of $\eta_{\mathbf 1}$) and $\prmK=\bigcap_x\ker u_x$, a two-sided ideal.

\emph{Rows.} A row is a vector $r\in\prmL[G]^{\prmgen}$; its value is $\prmval(r)=\sum_ir_i(\prmgen_i-1)
\in\priAug$. The map $\prmval$ is left $\prmL[G]$-linear. For $m\in\auO^{\prmgen}$ let $\bar m\in
\prmL[G]^{\prmgen}$ be the row with entries $m_i\cdot1$ (the \emph{coordinate row}), and $\epsilon_i$
the $i$-th unit row. The \emph{tie ideal} $\prmT$ is the left ideal of $\prmL[G]$ generated by
$\prmK$, by $p\priAug$, by the tied relations $t_k=\tau_k-1-X_jH_j(\phi_k)(\phi_k-1)$, $k=(j,\pm)$,
and by $\priAug\rho_0$. The set $N$ of \emph{tied rows}, those with $\prmval(r)\in\prmT$, is a left
$\prmL[G]$-submodule of $\prmL[G]^{\prmgen}$.

\begin{lemma}[Tied rows]\label{prm:tied}
The following rows are tied:
{\rm(R1)} $\overline{z_{gh}}-\overline{z_g}-g\,\overline{z_h}$ for $g,h\in G$;
{\rm(R2)} $\overline{z_{\prmgen_i}}-\epsilon_i$ for every index $i$;
{\rm(R3)} $\overline{\sum_g\theta_gz_g}$ whenever $\theta\in\auO^{(G)}$ and $\sum_g\theta_g(u_x(g)-1)=0$
for all $x$;
{\rm(R4)} $p\,r$ for every row $r$;
{\rm(R5)} $\overline{z_{\tau_k}}-X_jH_j(\phi_k)\,\overline{z_{\phi_k}}$ for $k=(j,\pm)$;
{\rm(R6)} $(g-1)\bigl(c_\zeta\overline{z_{\gamma_\pi}}+(1-\nu_0)\overline{z_a}\bigr)$ for $g\in G$;
{\rm(R7)} $o_v(\delta_v)\,\epsilon_{v}$, $\epsilon_v$ the unit row of the coordinate of $\delta_v$.
\end{lemma}
\begin{proof}
For $m\in\auO^{\prmgen}$, $u_x(\prmval(\bar m))=\sum_im_i(u_x(\prmgen_i)-1)$; hence
$\prmval(\overline{z_g})-(g-1)\in\prmK$ for every $g$. (R1): modulo the two-sided ideal $\prmK$ the value
is $(gh-1)-(g-1)-g(h-1)=0$. (R2): the value is $\prmval(\overline{z_{\prmgen_i}})-(\prmgen_i-1)\in
\prmK$. (R3): $u_x$ of the value is $\sum_g\theta_g(u_x(g)-1)=0$ for all $x$. (R4): $\prmval(pr)=p\,
\prmval(r)\in p\priAug$. (R5): the value is $t_k+\bigl(\prmval(\overline{z_{\tau_k}})-(\tau_k-1)\bigr)-
X_jH_j(\phi_k)\bigl(\prmval(\overline{z_{\phi_k}})-(\phi_k-1)\bigr)\in\prmT$. (R6): the value of
$c_\zeta\overline{z_{\gamma_\pi}}+(1-\nu_0)\overline{z_a}$ is $\rho_0$ modulo $\prmK$, so the value of the row
is $(g-1)\rho_0$ plus an element of $\prmK$, and $(g-1)\rho_0\in\priAug\rho_0$. (R7): the value is
$o_v(\delta_v)(\delta_v-1)$, which every $u_x$ kills by hypothesis~(6).
\end{proof}

\emph{Reading rows on a character.} For a character $\xi$ let $E_\xi\colon\prmL[G]^{\prmgen}\to
\prmL^{s_0\sqcup s}$, $E_\xi(r)=(\eta_\xi(r_i))_{i\in s_0\sqcup s}$: read the row on $\xi$ and forget
the coordinates of the $\delta_v$. It is $\prmL$-linear, and $E_\xi(\lambda r)=\eta_\xi(\lambda)E_\xi(r)$
for $\lambda\in\prmL[G]$. The forms of \S\ref{pri:ss-fitting} are taken with $n=\prmgen$, the
function $z\colon g\mapsto z_g$, and the maps $\Phi$ below.

\begin{lemma}[Forms of tied rows are character cocycles]\label{prm:forms}
Let $\xi$ be a character of $G$ with $\xi(\gamma_\pi)=1$, let $A\colon\prmL^{s_0\sqcup s}\to\prmL^m$
be $\prmL$-linear, $K\subseteq\prmL^m$ a $\prmL$-submodule containing $A(E_\xi(N))$, and $\Phi=A\circ
E_\xi\circ\overline{(\cdot)}\colon\auO^{\prmgen}\to\prmL^m$. Then $\Phi(p\auO^{\prmgen})\subseteq K$, and
for every $\auO$-linear $\Psi\colon\auO^{\prmgen}\to\auO/p$ vanishing on $\Phi^{-1}(K)$ the function
$c=\Psi\circ z$ is a character cocycle for $\xi$ (relative to the $\delta_v$, $\gamma_\pi$ and $a$).
\end{lemma}
\begin{proof}
If $m\in\auO^{\prmgen}$ and some tied row $r$ has $\eta_\xi(r_i)=m_i$ for all $i\in s_0\sqcup s$, then
$E_\xi(r)=E_\xi(\bar m)$, so $\Phi(m)=A(E_\xi(r))\in K$ and $\Psi(m)=0$. By (R4), $\Phi(pm)\in K$.
Apply this to the rows of Lemma~\ref{prm:tied}. (R1), read on $\xi$, has entries $z_{gh,i}-z_{g,i}-
\xi(g)z_{h,i}$; so $c(gh)=c(g)+\xi(g)c(h)$, and $c$ is a cocycle for $\xi$. (R3) gives $\sum_g\theta_g
c(g)=0$ whenever $\sum_g\theta_g(u_x(g)-1)=0$ for all $x$: $c$ obeys the linear relations. (R2) for the
index of $\delta_v$: the unit row $\epsilon_v$ has no coordinate in $s_0\sqcup s$, so the row reads as
$\overline{z_{\delta_v}}$ and $c(\delta_v)=0$. (R6) read on $\xi$ has entries $(\xi(g)-1)\bigl(c_\zeta z_{
\gamma_\pi,i}+\eta_\xi(1-\nu_0)z_{a,i}\bigr)$, and $\eta_\xi(\nu_0)=c_\zeta(\xi(\gamma_\pi)-1)=0$; so
$(\xi(g)-1)(c_\zeta c(\gamma_\pi)+c(a))=0$ for all $g$, and if $c(\gamma_\pi)=0$ then $(\xi(g)-1)c(a)=0$.
\end{proof}

\begin{lemma}[Common unit rows]\label{prm:unit}
Assume, as in {\rm(2)} and {\rm(7)}, that $\xi_e(\gamma_\pi)=1$ and that $s$ detects every module of
character cocycles for $\xi_e$, for every $e<t'$. Then for every $i\in s_0$ there is a tied row
$\nu_i$ with $\eta_e(\nu_{i,j})=1$ if $j=i$ and $=0$ if $j\in s_0\smallsetminus\{i\}$, for all $e<t'$.
\end{lemma}
\begin{proof}
\emph{One character.} Fix $e$ and let $F\colon\prmL[G]^{\prmgen}\to\prmL^{s_0}$ be $E_{\xi_e}$ followed by
the projection to the coordinates in $s_0$, $K=F(N)$, $\Phi=F\circ\overline{(\cdot)}$. By
Lemma~\ref{prm:forms} (with $A$ the projection and $m=\#s_0$), $\Phi(p\auO^{\prmgen})\subseteq K$ and
the forms $\Psi\circ z$ ($\Psi$ vanishing on $\Phi^{-1}K$) make up a module of character cocycles for
$\xi_e$, which $s$ detects. For $g\in s$, the row (R2) of the index $g$ (in the copy $s$ of
$\prmgen$) is tied and $F(\epsilon_g)=0$, so $\Phi(z_g)\in K$ and every such $\Psi$ has $\Psi(z_g)=0$.
Hence $\Psi\circ z$ vanishes on $s$, so $\Psi\circ z=0$; in particular $\Psi(z_{i})=0$ for the element
$i\in s_0$. By Lemma~\ref{pri:fitt-basic}(c) (whose proof applies verbatim to $\prmL^{s_0}$),
$\Phi(z_i)\in K$. With the row (R2) of the index $i$, $F(\epsilon_i)=\Phi(z_i)-F(\overline{z_i}-
\epsilon_i)\in K$: some tied row $w$ has $\eta_e(w_j)=\delta_{ij}$ for $j\in s_0$.

\emph{All characters.} Let $E'$ be a set of indices $e$ for which a tied row $w$ with $\eta_{e}(w_j)=
\delta_{ij}$ ($j\in s_0$, $e\in E'$) has been found, and $e_1\notin E'$. By the first part choose tied
rows $m_{j}$ ($j\in s_0$) with $\eta_{e_1}(m_{j,j'})=\delta_{jj'}$ for $j'\in s_0$, and put
$w'=w+\sum_{j\in s_0}(\delta_{ij}-w_j)\,m_j$, a tied row ($N$ is a left module). For $j'\in s_0$,
$\eta_{e_1}(w'_{j'})=\eta_{e_1}(w_{j'})+\sum_j(\delta_{ij}-\eta_{e_1}(w_j))\delta_{jj'}=\delta_{ij'}$,
and for $e\in E'$, $\eta_e(\delta_{ij}-w_j)=0$, so $\eta_e(w'_{j'})=\eta_e(w_{j'})=\delta_{ij'}$.
Starting from $E'=\emptyset$ and $w=0$, after $t'$ steps we get $\nu_i$.
\end{proof}

Fix such rows $\nu_i$. For $e<t'$ let $A_e\colon\prmL^{s_0\sqcup s}\to\prmL^s$, $A_e(b)_g=b_g-\sum_{i\in
s_0}b_i\,\eta_e(\nu_{i,g})$ ($g\in s$), and $K_e=A_e(E_{\xi_e}(N))\subseteq\prmL^s$. Let $\prmn\subseteq
\prmL$ be the $\auO$-span of the monomials that do not divide $X_J$; it is an ideal (a multiple of such
a monomial is again such a monomial), and $f\in\prmn$ if and only if no monomial of $f$ divides $X_J$.
For $e<t'$ put $X_{<e}=\prod_{\prity(i)<e}X_i$, $X_{=e}=\prod_{\prity(i)=e}X_i$ and $\prmq_e=(\prmn:
X_{<e})=\{f:fX_{<e}\in\prmn\}$.

\begin{lemma}\label{prm:fittsum}
For a $\prmL$-submodule $K\subseteq\prmL^s$ and an ideal $\prmq$ of $\prmL$, $\priFitt(K+\prmq\prmL^s)
\subseteq\priFitt(K)+\prmq$.
\end{lemma}
\begin{proof}
A generator of the left side is $\det$ of a matrix with rows $k_l+q_l$, $k_l\in K$, $q_l\in\prmq
\prmL^s$. Expanding multilinearly in the rows, the term with all rows $k_l$ lies in $\priFitt(K)$, and
every other term has a row with entries in $\prmq$, so lies in $\prmq$.
\end{proof}

\begin{proposition}[The block of one character]\label{prm:block}
For every $e<t'$ there is $\Gamma_e\in\prmL$ with
$\prival(\Gamma_e(0))\le b_e+(\#\{\prity=e\}+\#s)\prival(C)$ and $\Gamma_eX_{=e}\in\priFitt(K_e+
\prmq_e\prmL^s)$.
\end{proposition}
\begin{proof}
Apply Proposition~\ref{pri:fitt-triangle} with $R=\prmL$, the set $s$, $N_e:=K_e+\prmq_e\prmL^s$ in the
role of $N$, $\Phi=A_e\circ E_{\xi_e}\circ\overline{(\cdot)}$, $z$ as above, $I=\{j:\prity(j)=e\}$,
$\tau_j=\tau_{(j,+)}$, $\tau'_j=\tau_{(j,-)}$ and likewise $\phi$. We check its hypotheses, writing
$y_g=\Phi(z_g)$. (i) $\Phi(p\auO^{\prmgen})\subseteq N_e$ by (R4). (ii) For $g\in s$, the row (R2) of
the index $g$ is tied and $A_e(E_{\xi_e}(\epsilon_g))=\epsilon_g$ (it has no coordinate in $s_0$), so
$y_g-\epsilon_g\in K_e$. (iii) By (3), $H_j(\xi_e(\phi_k))=1+pc$ with $c\in\auO$. The sum of the tied
rows (R5) and $p\,X_jc\,\overline{z_{\phi_k}}$ (R4) reads on $\xi_e$ as $E_{\xi_e}(\overline{z_{\tau_k}})
-X_j(1+pc)E_{\xi_e}(\overline{z_{\phi_k}})+pX_jcE_{\xi_e}(\overline{z_{\phi_k}})$, since $\eta_e(X_j
H_j(\phi_k))=X_jH_j(\xi_e(\phi_k))$; applying $A_e$, $y_{\tau_k}-X_jy_{\phi_k}\in K_e$ for $k=(j,\pm)$.
(iv) Let $i<j$ with $j\in I$. As $\prity$ is monotone, $\prity(i)\le e$; if $\prity(i)=e$ then $i\in I$;
otherwise $X_iX_{<e}$ is divisible by $X_i^2$, so it lies in $\prmn$, $X_i\in\prmq_e$ and
$X_i\prmL^s\subseteq N_e$. (v) By Lemma~\ref{prm:forms} (with $A=A_e$, $K=N_e\supseteq A_e(E_{\xi_e}(N))$
and $\xi_e(\gamma_\pi)=1$ by (2)), $\mathrm F(N_e)$ is a module of character cocycles for $\xi_e$, so by
(7) it has triangle cost $b_e$ along $I$ with constant $C$. Proposition~\ref{pri:fitt-triangle} gives
$\Gamma_e$ as stated.
\end{proof}

\begin{proposition}[Character rows]\label{prm:chars}
There are, for each $e<t'$, finitely many matrices $B_{e,i}\in M_s(\prmL)$ and weights $w_{e,i}\in
\prmL$ ($1\le i\le l_e$), elements $\zeta_e\in\prmq_e$ and $\Gamma_e$ as in
Proposition~\ref{prm:block}, such that
\[
d_e:=\sum_iw_{e,i}\det B_{e,i}=\Gamma_eX_{=e}-\zeta_e ,
\]
and, for each choice $\mathbf i=(i(e))_{e<t'}$, a square matrix $M_{\mathbf i}$ over $\prmL[G]$,
indexed by $\prmgen$, all of whose rows are tied, with
\[
\det\eta_e(M_{\mathbf i})=c_{\mathrm{ch}}\det B_{e,i(e)}\quad(e<t'),\qquad c_{\mathrm{ch}}=C^{\#s}\prod_vo_v(1).
\]
\end{proposition}
\begin{proof}
\emph{The determinants.} By Proposition~\ref{prm:block} and Lemma~\ref{prm:fittsum}, $\Gamma_eX_{=e}=
d_e+\zeta_e$ with $d_e\in\priFitt(K_e)$ and $\zeta_e\in\prmq_e$. By definition of the Fitting ideal,
$d_e=\sum_iw_{e,i}\det B_{e,i}$ with matrices $B_{e,i}$ all of whose rows lie in $K_e$.

\emph{Lifting the rows.} Let $b$ be a row of some $B_{e,i}$: $b=A_e(E_{\xi_e}(r))$ with $r\in N$. Put
$r'=r-\sum_{i'\in s_0}r_{i'}\nu_{i'}$, a tied row. For $j\in s_0$, $\eta_e(r'_j)=\eta_e(r_j)-
\sum_{i'}\eta_e(r_{i'})\delta_{i'j}=0$; for $g\in s$, $\eta_e(r'_g)=\eta_e(r_g)-\sum_{i'}\eta_e(r_{i'})
\eta_e(\nu_{i',g})=b_g$. So every row $b$ of $B_{e,i}$ is $E_{\xi_e}(r')|_s$ for a tied row $r'=:r'(e,i,b)$
whose coordinates in $s_0$ vanish on $\xi_e$.

\emph{Separating the characters.} By (4) and Lemma~\ref{prm:sep} below there are $q_e\in\auO[G]$ with
$\xi_{e'}(q_e)=C$ if $e'=e$ and $0$ otherwise.

\emph{The matrix.} Fix $\mathbf i$. Let $M_{\mathbf i}$ have: in the row of $\delta_v$, the row (R7)
$o_v(\delta_v)\epsilon_v$; in the row of $i\in s_0$, $\nu_i$; in the row of $g\in s$, $\sum_{e<t'}q_e\,
r'(e,i(e),b_{e,g})$, $b_{e,g}$ the row $g$ of $B_{e,i(e)}$. All rows are tied. Apply $\eta_e$ and order
the columns as $(\delta,s_0,s)$. The rows of the $\delta_v$ have the single entry $o_v(\xi_e(\delta_v))=
o_v(1)$ in the column of $\delta_v$ (by (6), $\xi_e(\delta_v)=1$). The rows of $s_0$ have the identity
in the columns of $s_0$. The row of $g\in s$ is $\sum_{e'}\xi_e(q_{e'})\eta_e(r'(e',\dots))=
C\,\eta_e(r'(e,i(e),b_{e,g}))$, which is $0$ in the columns of $s_0$ and $C\,b_{e,g}$ in those of $s$. So
$\eta_e(M_{\mathbf i})$ is block lower triangular with respect to $(\delta)\sqcup(s_0\sqcup s)$, with
diagonal blocks $\mathrm{diag}(o_v(1))$ and $\left(\begin{smallmatrix}1&*\\0&CB_{e,i(e)}\end{smallmatrix}
\right)$, and its determinant is $\prod_vo_v(1)\cdot C^{\#s}\det B_{e,i(e)}$.
\end{proof}

\begin{lemma}\label{prm:sep}
Let $\xi_0,\dots,\xi_{t'-1}$ be characters of $G$, $a\in G$, $e<t'$ and $C\in\auO$ a multiple of
$\prod_{e'\ne e}(\xi_e(a)-\xi_{e'}(a))$. Then $q_e=w\prod_{e'\ne e}(a-\xi_{e'}(a))\in\auO[G]$, where
$C=w\prod_{e'\ne e}(\xi_e(a)-\xi_{e'}(a))$, has $\xi_{e''}(q_e)=C$ for $e''=e$ and $0$ for $e''\ne e$.
\end{lemma}
\begin{proof}
$\xi_{e''}(q_e)=w\prod_{e'\ne e}(\xi_{e''}(a)-\xi_{e'}(a))$; for $e''\ne e$ the factor $e'=e''$ vanishes.
\end{proof}

\begin{lemma}[Coefficients below $X_J$]\label{prm:coeff}
Let $d_e,\zeta_e,\Gamma_e\in\prmL$ ($e<t'$) with $d_e+\zeta_e=\Gamma_eX_{=e}$ and $\zeta_e\in\prmq_e$,
and let $F\in\prmL$. Then $F\prod_ed_e\equiv F\prod_e\Gamma_e\cdot X_J\pmod\prmn$. Consequently no
monomial properly dividing $X_J$ occurs in $F\prod_ed_e$, and its coefficient of $X_J$ is the constant
term of $F\prod_e\Gamma_e$.
\end{lemma}
\begin{proof}
Order the $e$ increasingly. Telescoping, $\prod_e(d_e+\zeta_e)-\prod_ed_e=\sum_e\zeta_e\prod_{e'<e}(d_{e'}+
\zeta_{e'})\prod_{e'>e}d_{e'}$. In the $e$-th term, $\prod_{e'<e}(d_{e'}+\zeta_{e'})=\prod_{e'<e}
\Gamma_{e'}\cdot X_{<e}$ (the sets $\{\prity=e'\}$, $e'<e$, partition $\{\prity<e\}$), and $\zeta_eX_{<e}
\in\prmn$; so each term lies in the ideal $\prmn$. Also $\prod_e(d_e+\zeta_e)=\prod_e\Gamma_e\cdot X_J$.
Multiply by $F$. An element of $\prmn$ has no monomial dividing $X_J$; every monomial of $F\prod_e
\Gamma_e\cdot X_J$ is a multiple of $X_J$, and the only one dividing $X_J$ is $X_J$, with coefficient the
constant term of $F\prod_e\Gamma_e$.
\end{proof}

\begin{proof}[Proof of the paragraphs ``The characters'' and ``Gluing'' of Proposition~\ref{pri:ribet}]
Let $B_0$ be the matrix of tied rows of the paragraph ``The block'', with $\prival(\det\rho_\pi(B_0)
|_{X_J=0})\le b_\rho$ (hypothesis (5) of Proposition~\ref{pri:ribet}), and $B_{e,i}$, $w_{e,i}$, $\zeta_e$, $\Gamma_e$, $M_{\mathbf i}$ as in
Proposition~\ref{prm:chars}. Let $k\in\auO[G]$ be $1$ on every $\xi_e$ and $0$ on every $\rho_x$,
$x\in\priX$: for each $x$, $k_x=(\zeta_x-1)^{-1}(\zeta_x-\gamma_x)$ is $0$ on $\rho_x$ (where $\gamma_x$
is the scalar $\zeta_x$) and $1$ on the $\xi_e$ (where $\gamma_x$ is $1$), by (2); take $k=\prod_xk_x$. Let
$q_0$ be as in Lemma~\ref{prm:sep} for $e=0$. For each $\mathbf i$ apply
Proposition~\ref{pri:ribet-elt} (with $\auO$ replaced by $\prmL=\auO[X_J]$, each $u_x$ extended $\prmL$-linearly and $\mathsf L$ by $\mathsf L[X_J]$: its proof uses only Lemma~\ref{pri:amitsur} and the Cayley--Hamilton theorem, both valid over any commutative ring, and $S\supseteq\prmL$ since $\priJ$ is a trace module) to the matrix $(1-k)B_0+kM_{\mathbf i}$ (its rows are tied, $N$ being a
left module), with $\rho_1=kq_0$ and $\mu_0=C$: $kq_0$ acts as $C$ on $\xi_0$, as $0$ on the other
$\xi_e$ and as $0$ on every $\rho_x$, and $\rho_0$ acts as $0$ on $\xi_0=1$ ($\eta_0(\nu_0)=0$,
$\eta_0(a-1)=0$). It gives $f_{\mathbf i}\in S$ with $f_{\mathbf i}\det u_x(w)\in\priJ$ for $w\in
\priAug$ and, for $x\in\priX$, $f_{\mathbf i}(x)=C\det\rho_x(B_0)\prod_e\det\eta_e(M_{\mathbf i})$, because
on $\rho_x$ the matrix is $\rho_x(B_0)$ and on $\xi_e$ it is $\eta_e(M_{\mathbf i})$. Put $f=\sum_{\mathbf
i}\bigl(\prod_ew_{e,i(e)}\bigr)f_{\mathbf i}$ (the weights lie in $\prmL\subseteq S$). For $x\in\priX$,
\[
f(x)=C\det\rho_x(B_0)\sum_{\mathbf i}\prod_ew_{e,i(e)}\,c_{\mathrm{ch}}\det B_{e,i(e)}=C\,c_{\mathrm{ch}}^{t'}\det\rho_x(B_0)
\prod_{e<t'}d_e .
\]
By Lemma~\ref{prm:coeff} with $F=Cc_{\mathrm{ch}}^{t'}\det\rho_x(B_0)$, $f(x)$ has no monomial properly dividing
$X_J$ for every $x\in\priX$, and the coefficient of $X_J$ in $f(\pi)$ is
$d=C\,c_{\mathrm{ch}}^{t'}\det\rho_\pi(B_0)|_{X_J=0}\prod_e\Gamma_e(0)\in\auO$. Its valuation is at most
\[
\prival(C)+t'\Bigl(\#s\,\prival(C)+\sum_v\prival(o_v(1))\Bigr)+b_\rho+\sum_e\bigl(b_e+(\#\{\prity=e\}+\#s)
\prival(C)\bigr),
\]
which is the bound of Proposition~\ref{pri:ribet}, as $\sum_e\#\{\prity=e\}=\#J$. The last sentence of
the proof of Proposition~\ref{pri:ribet} shows that $f$ is a common Ribet element.
\end{proof}

\begin{remark}\label{prm:rem-chars}
Two points of the paragraphs in the proof of Proposition~\ref{pri:ribet} deserve emphasis. (i) In
``Gluing'' the element $\rho_1$ must vanish on every $\rho_x$, as Proposition~\ref{pri:ribet-elt}
requires; $q_0$ alone (the element of Lemma~\ref{prm:sep} for $e=0$) does not, and $\rho_1=kq_0$ does. (ii) The vanishing of the character cocycles at the
$\delta_v$ comes from forgetting the coordinates of the $\delta_v$ (the row (R2)); the cost
$\prival(o_v(1))$ enters only through the rows (R7) that complete $M_{\mathbf i}$, not through
Lemma~\ref{pri:fitt-basic}(b).
\end{remark}

\subsection{The induction in Proposition~\ref{pri:norm}(b)}\label{prm:ss-norm}

We keep the notation of \S\ref{pri:ss-norm}. Let $R$ be the ring of integral operators on the forms
of principal level at the $w_j$, and $\sigma\colon R\to R'$ a ring map (in Proposition~\ref{pri:norm},
a reduction modulo $\ell^\nu/(2n!)$). For $j\in J$ let $R_j\subseteq R$ be the subring generated by
the operators made of translates at $w_j$; it contains $\Gamma_j$, $\mathrm{Nm}_j$ and the operator
$\partial_j$ of Lemma~\ref{pri:oneplace}. As recalled in the proof of Proposition~\ref{pri:norm},
$R_j$ and $R_{j'}$ commute for $j\ne j'$, and $\priT$ commutes with every $R_j$. For $S\subseteq J$ put
\[
\priN_S=\prod_{j\in S}\partial_j\prod_{j\notin S}\mathrm{Nm}_j ,
\]
the coefficient of $X_S$ in $\prod_j(\mathrm{Nm}_j+X_j\partial_j)$; the order of the factors does not
matter. We use three facts.
\begin{enumerate}
\item[(N1)] $\priN_\emptyset=\mathrm{Nm}$: expanding $\prod_j\sum_{k_j\in\Gamma_j}k_j$ gives $\sum_{k\in
\Gamma}k$, as $\Gamma=\prod_j\Gamma_j$ and translates at different places multiply as in $\Gamma$.
\item[(N2)] For $j\notin S$, $\priN_S=\priN'\mathrm{Nm}_j$ and $\priN_{S\cup\{j\}}=\priN'\partial_j$ with
$\priN'=\prod_{i\in S}\partial_i\prod_{i\notin S\cup\{j\}}\mathrm{Nm}_i$ (move the factor at $j$ to the
right).
\item[(N3)] $\delta'\in\priT$ commutes with every $\priN_S$, with $\delta=c_0^2\delta'$ and with every
$k\in\Gamma$, hence with $\mathrm{Nm}$.
\end{enumerate}

\begin{lemma}[From the factors to the group]\label{prm:invariant}
Let $\mathsf A\in R$ with $\sigma(\mathsf A(k-1))=0$ for every $k$ in every factor $\Gamma_j$, and let
$\Psi$ be as in Lemma~\ref{pri:forms}(a) for $\Gamma'=\Gamma$, so that $\sum_{k\in\Gamma}k^{-1}\Psi k=
c_0$. Then $\sigma(c_0\mathsf A)=\sigma(\mathsf A\Psi\,\mathrm{Nm})$.
\end{lemma}
\begin{proof}
The set $H=\{k\in\Gamma:\sigma(\mathsf A)\sigma(k)=\sigma(\mathsf A)\}$ contains every $\Gamma_j$ by
hypothesis and is closed under products: if $k,k'\in H$ then $\sigma(\mathsf A)\sigma(kk')=
\sigma(\mathsf A)\sigma(k)\sigma(k')=\sigma(\mathsf A)\sigma(k')=\sigma(\mathsf A)$. A subset of a finite
group closed under products and containing generators is the whole group, so $H=\Gamma$; in particular
$k^{-1}\in H$ for every $k$. Hence
$\sigma(c_0\mathsf A)=\sum_k\sigma(\mathsf A)\sigma(k^{-1})\sigma(\Psi k)=\sum_k\sigma(\mathsf A)\sigma(\Psi
k)=\sigma(\mathsf A\Psi\mathrm{Nm})$.
\end{proof}

\begin{proposition}[Proposition~\ref{pri:norm}(b)]\label{prm:normb}
Under the hypotheses of Proposition~\ref{pri:norm}, for every $S\subseteq J$ and every integer $i\ge
\#S$ there is $Y\in R$ with $\sigma(\delta^i\priN_S)=\sigma(Y\,\mathrm{Nm})$, and $Y=\delta^i$ if
$S=\emptyset$. In particular, with $i=\#J$, there are $Y_S$ with $\sigma(\delta^{\#J}\priN_S)=
\sigma(Y_S\mathrm{Nm})$ and $Y_\emptyset=\delta^{\#J}$.
\end{proposition}
\begin{proof}
Induction on $\#S$. For $S=\emptyset$, $\priN_\emptyset=\mathrm{Nm}$ by (N1); take $Y=\delta^i$.

Let $S\ne\emptyset$ and $i\ge\#S$; write $i=i'+1$, so $i'\ge\#S-1$. Put $\mathsf A=c_0\,\delta'\,
\delta^{i'}\priN_S$. We show $\sigma(\mathsf A(k-1))=0$ for $k$ in each factor $\Gamma_{j'}$.

\emph{Case $j'\notin S$.} By (N2), $\priN_S=\priN'\mathrm{Nm}_{j'}$, and $\mathrm{Nm}_{j'}k=\mathrm{Nm}_{j'}$
for $k\in\Gamma_{j'}$. So $\mathsf A(k-1)=0$.

\emph{Case $j'\in S$.} Let $S'=S\smallsetminus\{j'\}$. By (N2) for $S'$ and $j'$, $\priN_{S'}=\priN'
\mathrm{Nm}_{j'}$ and $\priN_S=\priN'\partial_{j'}$. Lemma~\ref{pri:oneplace} at $w_{j'}$ (whose
hypothesis is that of Proposition~\ref{pri:norm} at $w_{j'}$) gives $Z\in R$ with $\sigma(c_0\partial_{j'}
(k-1))=\sigma(\mathrm{Nm}_{j'}Z)$ and $\delta'\mathrm{Nm}Z=0$. Hence
\[
\sigma(\mathsf A(k-1))=\sigma(\delta'\delta^{i'}\priN')\,\sigma(c_0\partial_{j'}(k-1))=
\sigma(\delta'\delta^{i'}\priN'\mathrm{Nm}_{j'}Z)=\sigma(\delta'\delta^{i'}\priN_{S'}Z).
\]
By (N3), $\delta'\delta^{i'}\priN_{S'}=\delta^{i'}\priN_{S'}\delta'$. As $\#S'=\#S-1\le i'$, the induction
hypothesis gives $Y'$ with $\sigma(\delta^{i'}\priN_{S'})=\sigma(Y'\mathrm{Nm})$. Therefore, using (N3)
once more,
\[
\sigma(\mathsf A(k-1))=\sigma(Y')\,\sigma(\mathrm{Nm}\,\delta'Z)=\sigma(Y')\,\sigma(\delta'\mathrm{Nm}Z)=0 .
\]

So Lemma~\ref{prm:invariant} applies to $\mathsf A$: $\sigma(c_0\mathsf A)=\sigma(\mathsf A\Psi
\mathrm{Nm})$. Since $c_0\mathsf A=c_0^2\delta'\delta^{i'}\priN_S=\delta^{i}\priN_S$, the element $Y=
\mathsf A\Psi$ works.
\end{proof}

\begin{remark}\label{prm:rem-norm}
(i) The induction must be run for all $S$ simultaneously and, at each step, at every place
$j'\in S$, not only at one chosen $j$: this is why the claim is proved for every $i\ge\#S$ (the
exponent drops by one when one passes from $S$ to $S\smallsetminus\{j'\}$). (ii) The operators $Y_S$
need not commute with $\priT$ ($\Psi$ is not a Hecke operator); the consequence drawn in
Proposition~\ref{pri:norm} uses only that the $\priN_S$ commute with $\priT$, and that $\delta$
commutes with them; accordingly Proposition~\ref{pri:norm} asserts the commutation for the
$\priN_S$ only.
\end{remark}

\subsection{The base level at the chosen places}\label{prm:ss-level}

Let $O=(W,V,\mu,Q)$ be a pair with $r'=n_W$, $n=n_V$, $t'=n-r'$, lifted by the single lift $\Theta$ of
\S\ref{po:ss-set} with the block datum $\Phi_\vartheta$. At a split place $w$ we identify $\auU(V)$
with $\auGL_n(E_w)$ in the coordinates of \S\ref{po:ss-set}, in which $W$ corresponds to the first $r'$
basis vectors ($\iota_0=\left(\begin{smallmatrix}1_{r'}\\0\end{smallmatrix}\right)$ in the notation of
the proof of Lemma~\ref{pay:product}); the places of $Q$ and the chosen places are split. For a
matrix $k$ we write $k^{(11)}$ for its upper left $r'\times r'$ block, and $K_v$ for the stabiliser of
the standard lattice at $v$.

\begin{definition}[Base level]\label{prm:def-level}
The \emph{base level} of $O$ is the open compact subgroup $J_O=\prod_vJ_{O,v}$ of $\auU(V)(\mathbb A_f)$
with $J_{O,v}=K_v$ at the places not below a chosen place or a place of $Q$, and
\[
\begin{gathered}J_{O,w}=\Bigl\{k\in\auGL_n(\auO_w):k\equiv\begin{pmatrix}1_{r'}&*\\0&*\end{pmatrix}\bmod w\Bigr\}\
(w\text{ chosen}),\\
J_{O,w}=\Bigl\{k\in\auGL_n(\auO_w):k\equiv\begin{pmatrix}*&0\\ *&*\end{pmatrix}\bmod w\Bigr\}\ (w\in Q),\end{gathered}
\]
blocks of sizes $(r',t')$. Its \emph{level character} is $\chi_O(k)=\prod_{w\in Q}\psi_w\bigl(\det
k_w^{(11)}\bmod w\bigr)^{-1}$, $\psi_w$ the tame part of $\mu$ at $w$; it is a character because the
reductions at $w\in Q$ are block lower triangular, so the upper left block of a product is the product
of the upper left blocks. $\mathcal F(O)$ is the space of algebraic modular forms $y$ on $\auU(V)$ of
the weight of $O$ with $y(hk)=\chi_O(k)y(h)$ for $k\in J_O$ (the \emph{forms of base level}). For
$u=(u_w)\in\auG_O$ choose $k'_u\in\prod_vK_v$ reducing to $\mathrm{diag}(u_w,1_{t'})$ at each chosen
$w$ and to $1$ at the places of $Q$, and put $\varrho(u)y=\prod_{w\ \mathrm{chosen}}\psi_w(\det u_w)\,
y(\,\cdot\,k'_u)$.
\end{definition}

\begin{lemma}[The level and the action]\label{prm:level}
{\rm(a)} Each $k'_u$ normalises $J_O$ and $\chi_O(k_u'^{-1}kk'_u)=\chi_O(k)$; $\varrho(u)$ preserves
$\mathcal F(O)$, does not depend on the choice of $k'_u$, and $\varrho$ is a representation of $\auG_O$
on $\mathcal F(O)$.
{\rm(b)} For every form $F$ on $\auU(W)$ in the domain of $\Theta$, $\Theta F$ satisfies $(\Theta F)(hk)=
\chi_O(k)(\Theta F)(h)$ for $k\in J_O$; and for $u\in\auG_O$, $\varrho(u)(\Theta F)=\Theta\bigl(F(\,\cdot\,k_u)\bigr)$, $k_u\in K_W$
reducing to $u_w$ at the chosen places and to $1$ at $Q$; so $\Theta$ maps the $\sigma_c$-isotypic
part for (untwisted) right translation on $\auU(W)$ into the $\sigma_c$-isotypic part for $\varrho$.
\end{lemma}
\begin{proof}
(a) At a chosen $w$, $\mathrm{diag}(u_w,1)^{-1}\left(\begin{smallmatrix}1&B\\0&D\end{smallmatrix}\right)
\mathrm{diag}(u_w,1)=\left(\begin{smallmatrix}1&u_w^{-1}B\\0&D\end{smallmatrix}\right)$; at $w\in Q$,
$k'_u\equiv1$; elsewhere $k'_u\in K_v=J_{O,v}$. So $k'_u$ normalises $J_O$, and $\chi_O$, which only reads
the components at $Q$, where conjugation by $k'_u$ is trivial modulo $w$, is unchanged. If $y\in
\mathcal F(O)$ and $k\in J_O$, then $y(hkk'_u)=y(hk'_u\cdot k_u'^{-1}kk'_u)=\chi_O(k)y(hk'_u)$, so
$\varrho(u)y\in\mathcal F(O)$. Two choices of $k'_u$ differ by an element of $\prod_vK_v$ reducing to $1$
at the chosen places and at $Q$; it lies in $J_O$ and $\chi_O$ is $1$ on it, so the translates agree.
Finally $k'_uk'_{u'}$ is a valid choice for $uu'$, and the factors $\psi_w(\det u_w)$ are
multiplicative; so $\varrho(u)\varrho(u')=\varrho(uu')$.

(b) We use $(\Theta F)(hk)=\int\theta(\omega_\mu(g,hk)\Phi_\vartheta)F(g)\,dg$ and $\omega_\mu(g,hk)=
\omega_\mu(g,h)\omega_\mu(1,k)$, as in the proof of Lemma~\ref{pay:product}; it suffices to show
$\omega_\mu(1,k)\Phi_\vartheta=\chi_O(k)\Phi_\vartheta$ for $k\in J_O$, place by place. By
\eqref{po:eq-w5}, for $k_w\in\auGL_n(\auO_w)$, $\omega_\mu(1,k_w)\Phi_{\vartheta,w}(T)=\Phi_{\vartheta,w}
(k_w^{-1}T)$ ($\eta_w$ is unramified, hence trivial on $\auGL_n(\auO_w)$). At a chosen $w$,
$\Phi_{\vartheta,w}=\plone[T\equiv\iota_0]$ (proof of Lemma~\ref{pay:product}); for $k_w\in J_{O,w}$ the
first $r'$ columns of $k_w$ reduce to those of $1_n$, so $k_w\iota_0\equiv\iota_0$ and $k_w^{-1}T\equiv
\iota_0\iff T\equiv k_w\iota_0\equiv\iota_0$: the datum is fixed. At $w\in Q$, $\Phi_{\vartheta,w}(T)=
\plone[T\text{ integral}]\,\psi_w(\det T^{(1)}\bmod w)$, $T^{(1)}$ the upper $r'\times r'$ block of
$T\in M_{n\times r'}$ (proof of Lemma~\ref{pay:product}); for $k_w\equiv\left(\begin{smallmatrix}a&0\\
c&d\end{smallmatrix}\right)$, $k_w^{-1}\equiv\left(\begin{smallmatrix}a^{-1}&0\\ *&d^{-1}\end{smallmatrix}
\right)$, so the upper block of $k_w^{-1}T$ is $a^{-1}T^{(1)}$ modulo $w$ and $\Phi_{\vartheta,w}(k_w^{-1}T)=
\psi_w(\det a)^{-1}\Phi_{\vartheta,w}(T)$. Elsewhere $k_w\in K_v$ fixes the vacuum ((W4) of
Lemma~\ref{ps:weil}). The product of the factors is $\chi_O(k)$. The equivariance: by Lemma~\ref{pay:product},
$(\Theta F)(hk'_u)=\prod_{w\ \mathrm{chosen}}\psi_w(\det u_w)^{-1}\,\Theta(F(\,\cdot\,k_u))(h)$, and the
factor $\prod\psi_w(\det u_w)$ in the definition of $\varrho$ cancels it, giving
$\varrho(u)(\Theta F)=\Theta(F(\,\cdot\,k_u))$; this does not depend on the choice of $k'_u$ by (a).
\end{proof}

So the space of forms produced by the lift is contained in $\mathcal F(O)$ of
Definition~\ref{prm:def-level}, and $e_c$ is the central idempotent of the family of labels $c$ for
$\varrho$; this is the precise form of the ``base level'' of Proposition~\ref{pri:BL} and of
\S\ref{pax:ss-BL}.

\emph{Twisted labels.} Let $w$ be chosen, $\auF_w=\auF_q$, and $c$ a label of degree $r'$ over
$\auF_w$ with factorisation $c=\prod_kf_k$ into irreducibles of degrees $d_k$. By
\S\ref{pga:ss-local}, $\sigma_c$ is a constituent of the Harish-Chandra induction from the cuspidal
representation $\bigotimes_k\sigma_{f_k}$ of $L_1=\prod_k\auGL_{d_k}(\auF_q)$. For a character $\phi$ of
$\auF_q^\times$ put $\lambda_k=\sigma_{f_k}\otimes(\phi\circ\det)$. Each $\lambda_k$ is irreducible
cuspidal (a vector fixed by the unipotent radical of a proper parabolic for $\lambda_k$ is fixed by it
for $\sigma_{f_k}$, as $\phi\circ\det$ is trivial on unipotent elements), and $\sigma_c\otimes(\phi
\circ\det)$ is a constituent of the Harish-Chandra induction of $\bigotimes_k\lambda_k$ (for a
character $\chi$ of the group, $\operatorname{Ind}_P(\tau)\otimes\chi\simeq\operatorname{Ind}_P(\tau
\otimes\chi|_P)$, and $\phi\circ\det$ restricted to a parabolic is inflated from its Levi). The
\emph{rescaled label} $\tilde c_w$ of $c_w$ is the label with $\sigma_{\tilde c_w}\simeq\sigma_{c_w}
\otimes(\psi_w^{-1}\circ\det)$; with $\phi=\psi_w^{-1}$, its irreducible factors $\tilde f_k$ have
$\sigma_{\tilde f_k}=\lambda_k$. (In terms of Deligne--Lusztig characters, $R_T(\theta)\otimes(\phi\circ
\det)=R_T(\theta\cdot\phi\circ\det|_T)$, so the roots of $\tilde c_w$ are those of $c_w$ multiplied by the
element of $\overline{\auF}_w^\times$ that names $\psi_w^{-1}$; this is the rescaling of the lift setting.
Only the cuspidal representations $\lambda_k$ are used below.)

\begin{lemma}[Product labels]\label{prm:label}
Let $c$ be a family of labels, $y\in\mathcal F(O)$ with $e_cy=y\ne0$, and $w$ a chosen place. Let
$P\subseteq G_w=\auGL_n(\auF_w)$ be the parabolic subgroup of block upper triangular matrices with
blocks $(r',t')$, $M\simeq\auGL_{r'}(\auF_w)\times\auGL_{t'}(\auF_w)$ its Levi subgroup, $\phi=
\psi_w^{-1}$, $\lambda_k$ as above for $c_w$, $L=\prod_k\auGL_{d_k}(\auF_w)\times\auGL_1(\auF_w)^{t'}$ and
$\lambda=\bigotimes_k\lambda_k\otimes\mathbf 1^{\otimes t'}$.
{\rm(a)} Let $V_w$ be the space of forms $y'$ of the weight of $O$ with $y'(hk)=\chi_O(k)y'(h)$ for $k$
in $J_O^{(w)}K_w(1)$, where $J_O^{(w)}$ is $J_O$ with the component at $w$ replaced by $1$ and $K_w(1)=
1+\varpi_wM_n(\auO_w)$. Then $y\in V_w$, $\auGL_n(\auO_w)$ acts on $V_w$ by right translation through
$G_w$, and every irreducible constituent of the $G_w$-module generated by $y$ has cuspidal support
$(L,\lambda)$ up to conjugacy.
{\rm(b)} For every irreducible automorphic representation $\Pi$ of $\auU(V)$ on which the orthogonal
projection of $y$ is non-zero, $\Pi_w^{K_w(1)}$ contains an irreducible representation of $G_w$ with
cuspidal support $(L,\lambda)$.
\end{lemma}
\begin{proof}
(a) $K_w(1)\subseteq J_{O,w}$ and $\chi_O$ is trivial on the component at $w$, so $y\in V_w$; $K_w(1)$
is normal in $\auGL_n(\auO_w)$ and acts trivially on $V_w$, so $\auGL_n(\auO_w)$ acts through
$G_w$. Let $U$ be the unipotent radical of $P$. The reduction of $J_{O,w}$ is $U\cdot(1\times
\auGL_{t'}(\auF_w))$, so $y$ is fixed by $U$ and by $1\times\auGL_{t'}(\auF_w)$.

For $g\in\auGL_{r'}(\auF_w)$, placed at $w$ (and $1$ at the other chosen places), $k'_g$ in
Definition~\ref{prm:def-level} may be taken with components in $K_v$ reducing to $1$ at the other chosen
places and at $Q$; those components lie in $J_O^{(w)}$ with $\chi_O=1$, so on $V_w$ the operator
$\varrho(g)$ is $\psi_w(\det g)$ times the action of $\mathrm{diag}(g,1)\in G_w$. Write $e_c=e_{c_w}e'$
with $e'$ the product of the idempotents at the other chosen places; these commute, so $e_{c_w}y=
e_{c_w}e_cy=e_cy=y$. Thus $y$ lies in the $\sigma_{c_w}$-isotypic part of $g\mapsto\psi_w(\det g)
\mathrm{diag}(g,1)$, that is, in the $\sigma'$-isotypic part of $V_w$ for $g\mapsto\mathrm{diag}(g,1)$,
where $\sigma'=\sigma_{c_w}\otimes(\phi\circ\det)$.

Let $X=\C[P]y$. As $U$ is normal in $P$ and fixes $y$, $U$ acts trivially on $X$, so $X=\C[M]y=\C[
\auGL_{r'}\times1]y$ ($1\times\auGL_{t'}$ fixes $y$ and commutes with $\auGL_{r'}\times1$). Hence $X$ lies
in the $\sigma'$-isotypic part for $\auGL_{r'}\times1$, and $1\times\auGL_{t'}$ acts trivially on it: as
a representation of $P$, $X$ is the inflation of $(\sigma'\boxtimes\mathbf 1)^{\oplus m}$ for some
$m\ge1$. The $G_w$-linear map $\C[G_w]\otimes_{\C[P]}X\to V_w$, $\alpha\otimes x\mapsto\alpha x$, has image
$\C[G_w]y$, and its domain is $\operatorname{Ind}_P^{G_w}X$, the Harish-Chandra induction of
$(\sigma'\boxtimes\mathbf 1)^{\oplus m}$. So every irreducible constituent of $\C[G_w]y$ is a constituent
of the Harish-Chandra induction $R_M^{G_w}(\sigma'\boxtimes\mathbf 1)$.

Now $\sigma'$ is a constituent of $R_{L_1}^{\auGL_{r'}}(\bigotimes_k\lambda_k)$ (twisted labels, above),
and the trivial representation of $\auGL_{t'}(\auF_w)$ is a constituent of $R_{T}^{\auGL_{t'}}(\mathbf 1)
=\C[\auGL_{t'}/B]$ (the constant functions), $T$ the diagonal torus and $B$ the Borel subgroup. So
$\sigma'\boxtimes\mathbf 1$ is a constituent of $R_L^M(\lambda)$, and by transitivity of Harish-Chandra
induction every constituent of $R_M^{G_w}(\sigma'\boxtimes\mathbf 1)$ is a constituent of
$R_L^{G_w}(\lambda)$, with $\lambda$ cuspidal. By Harish-Chandra's theorem (the irreducible constituents
of $R_L^{G}(\lambda)$, $\lambda$ cuspidal, have cuspidal support the $G$-conjugacy class of
$(L,\lambda)$; \cite[Proposition~6.3 and Theorem~6.4]{DigneMichel}; for $\auGL_n$ see \cite[\S2.3, Proposition; \S\S9.2--9.3]{ZelevinskyLNM}), every irreducible constituent of $\C[G_w]y$ has
cuspidal support $(L,\lambda)$.

(b) The space of forms of the weight of $O$ is the orthogonal direct sum of its $\Pi$-isotypic parts,
and the orthogonal projection $y\mapsto y_\Pi$ commutes with right translations; it maps $V_w$ into the
$\Pi$-part of $V_w$, on which $G_w$ acts through $\Pi_w^{K_w(1)}$ (tensored with a space on which $G_w$
acts trivially). So $\C[G_w]y_\Pi$, a non-zero quotient of $\C[G_w]y$, has an irreducible constituent,
which by (a) has cuspidal support $(L,\lambda)$ and is contained in $\Pi_w^{K_w(1)}$.
\end{proof}

\begin{corollary}[Inertia at a chosen place]\label{prm:inertia}
In the proof of Proposition~\ref{pri:BL}(c), paragraph ``Inertia at a chosen $w$'' of
\S\ref{pax:ss-BL}, the input ``$y$ is fixed by $K(1)$ at $w$ and lies in the isotypic part of the
representation induced from $\sigma_{\tilde c_w}\otimes\mathbf 1_{\auGL_{t'}}$'' holds in the form of
Lemma~\ref{prm:label}(b). With the cuspidal-support part of Lemma~\ref{pga:depth0} it gives: the
supercuspidal support of $\Pi_w$ consists of the depth-zero supercuspidal representations containing
the $\sigma_{\tilde f_k}$ and of $t'$ unramified characters (depth-zero characters of $E_w^\times$ whose
restriction to $\auO_w^\times$ inflates the trivial character of $\auF_w^\times$). Hence
Proposition~\ref{pri:BL} holds with the base level of Definition~\ref{prm:def-level}.
\end{corollary}
\begin{proof}
Lemma~\ref{prm:label}(b) provides an irreducible constituent of $\Pi_w^{K_w(1)}$ with cuspidal support
$(L,\lambda)$, $\lambda=\bigotimes_k\sigma_{\tilde f_k}\otimes\mathbf 1^{\otimes t'}$, which is the
hypothesis of the cuspidal-support paragraph of the proof of Lemma~\ref{pga:depth0}; the factor
$\auGL_1(\auF_w)$ with the trivial representation contributes a depth-zero supercuspidal of
$\auGL_1(E_w)$ containing the trivial character of $\auO_w^\times$, an unramified character. The rest of
the paragraph of \S\ref{pax:ss-BL} is unchanged ($\eta_O$ is unramified at $w$ by Lemma~\ref{pay:eta}).
Parts (a), (b) of Proposition~\ref{pri:BL} hold for $\mathcal F(O)$ by Definition~\ref{prm:def-level}
and Lemma~\ref{prm:level}(a) (the argument of \S\ref{pax:ss-BL}(b), with $\varrho$).
\end{proof}

\begin{remark}\label{prm:rem-level}
The level of Definition~\ref{prm:def-level} fixes the first $r'$ basis vectors modulo $w$ (the
vectors of $W$), and is the one the lift produces (Lemma~\ref{prm:level}(b)); a level of the form
$\{k\equiv\left(\begin{smallmatrix}*&*\\0&1_{t'}\end{smallmatrix}\right)\}$ is not. With it,
$\auG_O$ acts through $\mathrm{diag}(u_w,1_{t'})$, which normalises the level, and the trivial block
$\mathbf 1_{\auGL_{t'}}$ comes from the invariance under $1\times\auGL_{t'}$.
\end{remark}

\subsection{Where the inputs of \S\S\ref{pri:sec}, \ref{pga:sec}, \ref{pse:sec}, \ref{pg:sec} are proved}
\label{prm:ss-blocks}

The following list covers every statement of these four sections whose proof is given elsewhere,
with the place where it is proved.

\begin{enumerate}
\item Proposition~\ref{pri:BL} (base level of a label): proved in \S\ref{pax:ss-BL}, with
Lemma~\ref{pay:eta} ($\eta_O$ unramified away from $\ell$), Lemma~\ref{pay:product}, and
Definition~\ref{prm:def-level}, Lemmas~\ref{prm:level}, \ref{prm:label} and
Corollary~\ref{prm:inertia} (the base level and the product labels).
\item Lemma~\ref{pri:block} (block-diagonal lattices): Lemma~\ref{pax:block}, with the two hypotheses
(semisimplicity, spanning) made explicit there.
\item Lemma~\ref{pri:forms} (Higman operators, pairings): \S\ref{pax:ss-forms} and
\S\ref{pax:ss-pairing}.
\item Lemma~\ref{pri:chain} (chain lemma): \S\ref{pax:ss-chain}, with Lemma~\ref{pax:tadic} proved
directly in Corollary~\ref{pay:roots}.
\item Proposition~\ref{pri:odd} (odd classes): Proposition~\ref{pod:odd}, for every $\ell$ including
$\ell=2$ (Remark~\ref{pod:rem-two}).
\item Proposition~\ref{pri:kummer} (Kummer maps, detecting points): Proposition~\ref{pku:main} for (a)
and Proposition~\ref{pku:b} for (b).
\item Proposition~\ref{pri:dict} (dictionary at an auxiliary place): \S\ref{pax:ss-dict}, part (d)
through Lemma~\ref{pax:toproot}, proved in \S\ref{pay:ss-cons}.
\item The paragraphs ``The characters'' and ``Gluing'' of the proof of Proposition~\ref{pri:ribet}:
\S\ref{prm:ss-chars}.
\item The induction in Proposition~\ref{pri:norm}(b): Proposition~\ref{prm:normb}.
\item Speh constituents in Proposition~\ref{pga:G123} (Remark~\ref{pga:sign}):
Proposition~\ref{pax:sign} (the twist $|\det|^{s_i}$ is always of the right sign).
\item Theorem~\ref{pg:I}: the twist is Proposition~\ref{pax:sign}, the convention check is
Lemma~\ref{pax:conv}; the rest are published theorems, quoted with their hypotheses in
Theorem~\ref{pga:lit} and Lemmas~\ref{pga:depth0}, \ref{pga:parahoric}.
\item Proposition~\ref{pg:odd}: Proposition~\ref{pod:odd}.
\item Statement (IM) of \S\ref{pse:ss-stick} (integral measures): Proposition~\ref{pex:IM} for $z\ge2$
(at every $\ell\nmid\auNm Q$, $\ell=d$ included by Lemma~\ref{pex:dint}) and Proposition~\ref{pex:IM1}
for $z=1$ (from the exchange identity); in (c) the place $v$ must satisfy $v,\bar v\notin
\mathsf M\cup Q\cup\bar Q$.
\item Proposition~\ref{pg:loss} (the local loss): proved in \S\ref{pll:sec}, the explicit
cup exponents being Theorem~\ref{pll:exponent}, proved in \S\ref{pki:ss-deduce}. Nothing in the
present section bears on it.
\end{enumerate}


\section{The main inequality and the zero-free half-plane}\label{an:sec-main}
Throughout, $D$ is admissible data for parameters $(q_0,d)$, and
$R_q=R_{d,q}$. Every statement of
Sections~\ref{an:sec-criterion} and~\ref{an:sec-data}, and
Lemmas~\ref{an:lem-rem} and~\ref{an:lem-hmain} below, are proved in full. The
statements from Theorem~\ref{an:thm-main} on depend on
Theorems~\anThmA, \anThmB\ and~\anThmC, which are proved in
Section~\ref{sec:automorphic} and the sections that follow it (\S\ref{ss:au-remains}).

\subsection{The total remainder}\label{an:ssec-rem}
We recall the quantities of Definition~\ref{def:au-hmain}. For real
$y\ge0$, an integer $j\ge0$ and $\sigma\in\anPl$,
\[
 \ansf(y,\sigma)=\sum_{1\le\nu\le a}\anLam(\nu)\Bigl(\frac y\nu-\kappa\bigl(u(-\nu)\,\sigma\bigr)\Bigr)^{+},\qquad
 \ansi(j)=\log j!+j\log\frac M\Omega ,
\]
with $t^{+}=\max(t,0)$ and $\Omega$ as in
Lemma~\ref{an:lem-sizes}\textup{(S6)}. The \emph{main amount}, the
\emph{expected value} and the \emph{total remainder} are
\begin{gather*}
 \anhm=m\,k\,\anhd,\qquad \anhd=\sum_{\sigma\in\anPl}\Bigl(\ansi(\kappa(\sigma))-\ansf(\kappa(\sigma),\sigma)\Bigr),\\
 \ansm=a\Bigl(\log a-\frac1f\sum_{\tau\in\anPl}\log\kappa(\tau)+c_1\Bigr),\qquad
 \anrem=\sum_{\tau\in\anPa}\bigl(\ansf(a,\tau)-\ansm\bigr),
\end{gather*}
with $c_1=-1-\aneu$ as in \eqref{an:eq-c1}. A term of $\ansf$ with
$\anLam(\nu)\ne0$ has $\nu$ a power of a prime that does not split in $K$,
hence prime to $q_1$ by \eqref{an:eq-split}; so $u(-\nu)$ is the place of
class $-\nu$.

\begin{lemma}[the total remainder]\label{an:lem-rem}
\[
 \anrem=a\Bigl(\#\anPb\,\bigl(R_{q_0}(x)-R_{q_0}(x_1)\bigr)+\#\anPa\,R_1(x_1)\Bigr).
\]
\end{lemma}
\begin{proof}
For $q\in\{1,q_0\}$ and an integer $1\le z\le a$,
Definition~\ref{an:def-R} gives
\begin{equation}\label{an:eq-posparts}
 \sum_{1\le\nu\le a}\anlam_q(\nu)\Bigl(\frac a\nu-\frac az\Bigr)^{+}=a\sum_{\nu\le z}\anlam_q(\nu)\Bigl(\frac1\nu-\frac1z\Bigr)=a\Bigl(R_q(z)+\frac{\log z}{\varphi(q)}+c_q\Bigr).
\end{equation}
Let $\nu\le a$ with $\anLam(\nu)\ne0$. For $\tau\in\anPa$ the place
$u(-\nu)\tau$ has class $-\nu$ modulo $q_0$, and the places of $\anPb$ have
class $-1$. If $\nu\not\equiv1\pmod{q_0}$, no $u(-\nu)\tau$ lies in $\anPb$.
If $\nu\equiv1\pmod{q_0}$, then $\tau\mapsto u(-\nu)\tau$ is a bijection of
$\anPa$ onto the fibre of $-1$, which contains $\anPb$; so exactly $\#\anPb$
places $\tau\in\anPa$ have $u(-\nu)\tau\in\anPb$. By (D6), $\kappa=b=a/x$ on
$\anPb$ and $\kappa=a/x_1$ elsewhere. Therefore
\[
 \sum_{\tau\in\anPa}\Bigl(\frac a\nu-\kappa(u(-\nu)\tau)\Bigr)^{+}
 =\#\anPa\Bigl(\frac a\nu-\frac a{x_1}\Bigr)^{+}+\mathbf 1_{\nu\equiv1\,(q_0)}\,\#\anPb\Bigl[\Bigl(\frac a\nu-\frac ax\Bigr)^{+}-\Bigl(\frac a\nu-\frac a{x_1}\Bigr)^{+}\Bigr].
\]
Multiply by $\anLam(\nu)$, sum over $\nu\le a$, and apply
\eqref{an:eq-posparts} with $(q,z)=(1,x_1)$, $(q_0,x)$, $(q_0,x_1)$; the
constant $c_{q_0}$ cancels:
\[
 \sum_{\tau\in\anPa}\ansf(a,\tau)=a\,\#\anPa\bigl(R_1(x_1)+\log x_1+c_1\bigr)+a\,\#\anPb\Bigl(R_{q_0}(x)-R_{q_0}(x_1)+\frac{\log x-\log x_1}{\varphi(q_0)}\Bigr).
\]
On the other hand
$\sum_{\tau\in\anPl}\log\kappa(\tau)=f\log(a/x_1)-\#\anPb\,(\log x-\log x_1)$,
and $f=\varphi(q_0)\,\#\anPa$, so
\[
 \#\anPa\cdot\ansm=a\,\#\anPa\,(\log x_1+c_1)+a\,\#\anPb\,\frac{\log x-\log x_1}{\varphi(q_0)} .
\]
Subtracting gives the lemma.
\end{proof}
The weights are designed for this identity: at a place $\tau$ of class $1$
and for $y=a$, the cut-off of the sum over $\nu$ is $x$ when the conjugate
place $u(-\nu)\tau$ lies in $\anPb$, which requires $\nu\equiv1\pmod{q_0}$,
and it is $x_1$ otherwise. The expected value $\ansm$ removes the main terms
$\log z/\varphi(q)+c_q$, and the constant $c_1=-1-\aneu$ in it is exactly the
constant of Proposition~\ref{an:prop-pnt} for $q=1$.

\begin{lemma}\label{an:lem-hmain}
For $0\le y\le b$ and $\sigma\in\anPl$,
$0\le\ansf(y,\sigma)\le4y\log x_1\le y\log b$. Consequently
$\ansi(\kappa(\sigma))-\ansf(\kappa(\sigma),\sigma)\ge19\,\kappa(\sigma)\log b$
for every $\sigma$, and $\anhm\ge0$.
\end{lemma}
\begin{proof}
A term of $\ansf(y,\sigma)$ vanishes unless $y/\nu>\kappa(\cdot)\ge a/x_1$,
that is unless $\nu<yx_1/a\le x_1$; and $\anLam\le2\Lambda$. So
$\ansf(y,\sigma)\le2y\sum_{\nu\le x_1}\Lambda(\nu)/\nu\le4y\log x_1$, where we
used $n\sum_{\nu\le n}\Lambda(\nu)/\nu\le2\sum_{\nu\le n}\Lambda(\nu)\lfloor n/\nu\rfloor=2\log n!\le2n\log n$;
and $4\log x_1\le\log b$ by (S1). By (S6), $\ansi(j)\ge20\,j\log b$. Finally
$\kappa(\sigma)\le b$, and $m,k\ge0$.
\end{proof}


\subsection{The main inequality}\label{an:ssec-main}
We use Section~\ref{sec:automorphic} only through
Theorems~\anThmA, \anThmB\ and~\anThmC. In the notation of that section,
$\mathcal T$ is the space of target forms, of dimension $\andT$
(Definition~\ref{def:au-target}); $\mathcal L_D$ is the lattice of target
forms, and index and covolume of a finite family of target forms are those of
Definition~\ref{def:au-lattice}; $\overline G$ is a finite group acting on
$\mathcal T$; $\mathcal S^{\mathrm{ram}}$ is the space of ramified source
forms and $\mathcal S^{\mathrm{ram}}_{\mathrm{int}}$ the set of its integral
forms; $\mathfrak p\colon\mathcal S^{\mathrm{ram}}\to\mathcal T$ is the
linear map of Theorem~\anThmA, and $s\mathfrak p$, $s\in\overline G$, are its
translates. The three theorems state the following.
\begin{itemize}
\item[(A)] (Theorem~\anThmA.) (a) $\mathcal S^{\mathrm{ram}}$ is spanned by
the set $\mathcal S_0$ of the
$\varphi\in\mathcal S^{\mathrm{ram}}_{\mathrm{int}}$ with
$\|\Theta_1(s\mathfrak p\,\varphi)\|_\sigma^2\le\exp\bigl(\anhm-m\cdot\anrem+f\,a\,(\log b)^{4}\bigr)$
for all $s\in\overline G$ and $\sigma\in\operatorname{Aut}(\mathbb C)$.
(b) For every finite set $\mathsf t\subseteq\overline G$, every finite family
of target forms each of which lies in
$s\mathfrak p(\mathcal S^{\mathrm{ram}}_{\mathrm{int}})$ for some
$s\in\mathsf t$ has index at most
$\#\mathsf t\cdot\dim\mathcal S^{\mathrm{ram}}\cdot2\,f\,a\,(\log b)^{3}$.
\item[(B)] (Theorem~\anThmB.) $\andT>0$, and there is a finite set
$\mathsf t\subseteq\overline G$ with
$\#\mathsf t\cdot\dim\mathcal S^{\mathrm{ram}}\le5\,\andT\log(1+2b^{2})$ and
$\dim\sum_{s\in\mathsf t}s\mathfrak p(\mathcal S^{\mathrm{ram}})\ge(1-4/b^{2})\,\andT$.
\item[(C)] (Theorem~\anThmC.) If $y_1,\dots,y_n\in\mathcal T$ are linearly
independent, $n\ge(1-4/b^{2})\,\andT$, with covolume at most $c$ and index at
most $\delta$, then $\andT\bigl(\anhm-3\,f\,a\,(\log b)^{4}\bigr)\le c+2\delta$.
\end{itemize}
By Hadamard's inequality (see the lines after
Definition~\ref{def:au-lattice}), a family $y_1,\dots,y_n$ with
$\|\Theta_1y_i\|_\sigma^2\le e^{h}$ for all $i$ and $\sigma$ has covolume at
most $nh$.

\begin{theorem}[main inequality]\label{an:thm-main}
For all parameters $(q_0,d)$ and all admissible data $D$,
\[
 m\cdot\anrem\ \le\ f\,a\,(\log b)^{5}.
\]
\end{theorem}
\begin{proof}
Put $\ell=f\,a\,(\log b)^{4}$ and $e=f\,a\,(\log b)^{3}$. Let $\mathsf t$ be
as in (B) and
$V=\sum_{s\in\mathsf t}s\mathfrak p(\mathcal S^{\mathrm{ram}})$, $n=\dim V$.
By (A)(a) the set $\bigcup_{s\in\mathsf t}s\mathfrak p(\mathcal S_0)$ spans
$V$; choose a basis $y_1,\dots,y_n$ of $V$ in it. Then
$(1-4/b^{2})\andT\le n\le\andT$. By (A)(a) and Hadamard's inequality the
family has covolume at most $n(\anhm-m\cdot\anrem+\ell)$, and by (A)(b) and
(B) it has index at most
$\#\mathsf t\dim\mathcal S^{\mathrm{ram}}\cdot2e\le10\,\andT\log(1+2b^{2})\,e$.
By (C),
\[
 \andT\,(\anhm-3\ell)\le n\,(\anhm-m\cdot\anrem+\ell)+20\,\andT\log(1+2b^{2})\,e .
\]
Since $n\le\andT$ and $\anhm\ge0$ (Lemma~\ref{an:lem-hmain}),
$n\,\anhm\le\andT\,\anhm$; the main amount cancels and
\[
 n\cdot m\cdot\anrem\le\andT\bigl(4\ell+20\log(1+2b^{2})\,e\bigr).
\]
If $\anrem\le0$ there is nothing to prove. Otherwise, as $4/b^{2}\le\frac12$,
we have $\andT\le2n$ and $n>0$, so
\[
 m\cdot\anrem\le f\,a\,(\log b)^{3}\bigl(8\log b+40\log(1+2b^{2})\bigr)\le f\,a\,(\log b)^{3}\cdot128\log b\le f\,a\,(\log b)^{5},
\]
by $\log(1+2b^{2})\le3\log b$ and $\log b\ge128$ (S1).
\end{proof}
The main amount is large, $\anhm\ge19\,m\,k\,b\log b\,\#\anPb\ge x^{4}f\,a$
by Lemma~\ref{an:lem-hmain} and Lemma~\ref{an:lem-sizes}, while the error
allowed is $f\,a\,(\log b)^{5}$. The deduction above is linear arithmetic;
it requires the upper bound of (A) and the lower bound of (C) to agree to
relative precision $x^{-4}(\log b)^{5}$ in the term $\anhm$, and the
remainder $-m\cdot\anrem$ to be present in (A) and absent from (C). The whole
content of the inequality is in Theorems~\anThmA\ and~\anThmC.


\subsection{The role of \texorpdfstring{$m$}{m}}\label{an:ssec-m}
In Theorem~\ref{an:thm-main} the left side is linear in the number $m$ of
lifts and the right side does not depend on $m$. In the analytic part $m$
occurs in three places only.

(1) In the data: $m$ fixes the field through $2^{s}=m^{22}$ (so
$f\ge m^{22}$), the level through $\log N\asymp m\log b$ in (D8), and the
rank $k=r+1-m$. The construction of Theorem~\ref{an:thm-data} does not use
that $m$ is of size $\sqrt x$: with (D3) replaced by $x^{\mu}\le m\le2x^{\mu}$
for a fixed $\mu\in[\frac12,1)$ it goes through, with the exponents of $x$ in
\textup{(S2)--(S4)} changed accordingly.

(2) In $\anhm=m\,k\,\anhd$, which cancels in the proof of
Theorem~\ref{an:thm-main}.

(3) In the final division: by Lemma~\ref{an:lem-rem}, the main inequality
reads
\[
 R_{q_0}(x)\ \le\ \frac{\varphi(q_0)(\log b)^{5}}{m}\cdot\frac{\#\anPa}{\#\anPb}+R_{q_0}(x_1)-\frac{\#\anPa}{\#\anPb}\,R_1(x_1),
\]
and the last two terms are negligible. So the bound obtained for $R_{q_0}(x)$
is $x^{O(\varepsilon)}/m$, and $m\asymp\sqrt x$ gives the exponent
$-\frac12$.

Thus the analytic part needs exactly this from the automorphic part: that
the errors $f\,a\,(\log b)^{4}$ in \textup{(A)} and $3f\,a\,(\log b)^{4}$ in
\textup{(C)}, and the index bound in \textup{(A)}, hold without any factor
growing with $m$ when $m\le2\sqrt x$. It cannot be so for substantially
larger $m$. Indeed, suppose the three theorems held, with the same errors,
for data with $m\asymp x^{\mu}$, $\mu>\frac12$. The argument of
\S\ref{an:ssec-onesided} would give $R_{d,q_0}(y)\le y^{-\mu'}$ for large
$y$, with some $\mu'\in(\frac12,1)$, for every prime $d$ of the class. By
Lemma~\ref{an:lem-landau} with $\beta=1-\mu'<\frac12$, either $\anF_{q_0}$ is
holomorphic in $\Re s>1-\mu'$, or it has a real pole $x_0\in(1-\mu',1)$ with
$r_{q_0}(x_0)<0$. The first is impossible: by \eqref{an:eq-r} and the
functional equations, it would give $n_{q_0}(\rho)=n_{q_0}^{K}(\rho)$ for all
$\rho$ in the critical strip, whereas the number of zeros with
$|\Im\rho|\le T$ counted by $n^{K}_{q_0}$ exceeds the number counted by
$n_{q_0}$ by $\varphi(q_0)\frac T\pi\log d+O(\log T)$, the conductors of the
characters $\psi\anchi$ being $d$ times those of the $\psi$
\cite[Chapter~16]{Davenport}. Hence for every prime
$d\equiv q_0/2-1\pmod{q_0}$ some $L(s,\psi\anchi)$, $\psi$ modulo $q_0$,
would vanish at a real point of $(1-\mu',1)$, to an order exceeding the
total order of the $L(s,\psi)$ there; if Corollary~\ref{an:cor-GRH} holds, the
point is $s=\frac12$. Such a systematic vanishing at the centre is not
expected, although it is not excluded by anything proved in this paper.
Under the generalised Riemann hypothesis one has
$R_{d,q_0}(x)=O(x^{-1/2})$, by the explicit formula and
$r_{q_0}(\rho)=O(|\rho|^{-2})$ per zero; so with $m\asymp\sqrt x$ the main
inequality asserts, up to the factor $(\log b)^{6}$, no more than what the
conclusion gives back, and $\sqrt x$ is the largest order of $m$ for which
it can be expected to hold.

The place where the restriction $m\ll\sqrt x$ is used is not visible from
the statements \textup{(A)}, \textup{(B)}, \textup{(C)}. It is located in
\S\ref{ss:au-mechanism}: two errors of the automorphic part are quadratic in
$m$, namely the loss at the level prime $N$ in Theorem~\anThmA\ and the
comparison operators in Theorem~\anThmC. The first has size $m^{2}f\,b\,\log b$
up to constants (Lemma~\ref{au:errsum}: $2mfb\log N\le60\,m^{2}fb\log b$), and
it is below the allowance $f\,a\,(\log b)^{4}=x\,f\,b\,(\log b)^{4}$ when
$m^{2}\ll x(\log b)^{3}$.
For general $m$ the outcome would be
$R_{q_0}(x)\ll(\log b)^{6}/m+m(\log b)^{3}/x$, which is best at
$m\asymp\sqrt x$, in agreement with the discussion above. As stated in
Remark~\ref{au:barrier-remark}, the dependence on $m$ of the remaining steps
of the automorphic part is not analysed here; only the case $m\le2\sqrt x$ is needed.

\subsection{The one-sided bound}\label{an:ssec-onesided}
\begin{theorem}[one-sided bound]\label{an:thm-onesided}
Let $(q_0,d)$ be parameters and $\varepsilon>0$. Then
\[
 R_{d,q_0}(y)\le y^{-1/2+\varepsilon}
\]
for all sufficiently large real $y$.
\end{theorem}
\begin{proof}
Let $\delta=\min(\frac1{200},\frac\varepsilon7)$. Let $x$ be a large integer
and $D$ the data of Theorem~\ref{an:thm-data} for $(q_0,d)$, $\delta$
and $x$. By Lemma~\ref{an:lem-rem} and Theorem~\ref{an:thm-main},
\[
 a\,\#\anPb\,R_{q_0}(x)=\anrem+a\,\#\anPb\,R_{q_0}(x_1)-a\,\#\anPa\,R_1(x_1)\le\frac{f\,a\,(\log b)^{5}}m+a\,\#\anPb\,|R_{q_0}(x_1)|+a\,\#\anPa\,|R_1(x_1)| .
\]
Divide by $a\,\#\anPb$ and use $f=\varphi(q_0)\#\anPa$ and
$\#\anPa\le\#\anPb\log b$ (S3):
\begin{equation}\label{an:eq-Rle}
 R_{q_0}(x)\le\frac{\varphi(q_0)\,(\log b)^{6}}{m}+|R_{q_0}(x_1)|+\log b\cdot|R_1(x_1)| .
\end{equation}
Here $m\ge\sqrt x$, $\log b\le2x^{\delta}$, and
$\log x_1\ge x^{\delta/6}-1$. By Proposition~\ref{an:prop-pnt}, applied to
the moduli $q_0$ and $1$ (both prime to $d$),
$|R_{q}(x_1)|\le C\exp(-c\sqrt{\log x_1})\le x^{-2}$ for $q\in\{q_0,1\}$ and
$x$ large; the constants depend on $(q_0,d)$ only, and $x_1$ is the only
point where the proposition is used. Hence, for every large integer $x$,
\[
 R_{q_0}(x)\le64\,\varphi(q_0)\,x^{-1/2+6\delta}+x^{-2}+2x^{\delta-2}\le\bigl(64\varphi(q_0)+3\bigr)x^{-1/2+6\delta}.
\]
For real $y\ge2$ let $x=\lfloor y\rfloor$. By
Lemma~\ref{an:lem-R}(i), $R_{q_0}(y)\le R_{q_0}(x)+C_{\mathrm{Ch}}/x$, and
$x^{-1/2+6\delta}\le2y^{-1/2+6\delta}$. So
$R_{q_0}(y)\le(128\varphi(q_0)+6+2C_{\mathrm{Ch}})\,y^{-1/2+6\delta}\le y^{-1/2+7\delta}\le y^{-1/2+\varepsilon}$
for large $y$.
\end{proof}
The order of the quantifiers matters: $q_0$ and $d$ are fixed first, and
all thresholds depend on $(q_0,d,\varepsilon)$ only. The bound is obtained
for every prime $d\equiv q_0/2-1\pmod{q_0}$, as
Theorem~\ref{an:thm-criterion} requires.


\subsection{The zero-free half-plane}\label{an:ssec-GRH}
\begin{theorem}\label{thm:GRH-half}
Let $q\ge1$ and let $\chi$ be a Dirichlet character modulo $q$. Then
\[
 L(s,\chi)\ne0\qquad\text{for }\Re s>\tfrac12 .
\]
\textup{(}At $s=1$ the function $L(s,\chi)$ has a pole if $\chi$ is
principal, and is non-zero otherwise.\textup{)}
\end{theorem}
\begin{proof}
Let $\Re\rho>\frac12$; we may assume $\Re\rho<1$. Choose $\beta$ with
$\frac12<\beta<\Re\rho$. Put $q_0=18480\,q$ and $a=q_0/2-1$; $a$ is odd and
prime to $a+1=q_0/2$, hence prime to $2q_0$. Every prime
$d\equiv a\pmod{q_0}$ gives parameters $(q_0,d)$, and by
Theorem~\ref{an:thm-onesided} with $\varepsilon=\beta-\frac12$,
$R_{d,q_0}(x)\le x^{\beta-1}$ for all large $x$. By
Theorem~\ref{an:thm-criterion}, applied to the modulus $q_0$ and the class
$a$, no $L(s,\psi)$, $\psi$ modulo $q_0$, vanishes at $\rho$. For the
character $\psi$ modulo $q_0$ induced by $\chi$,
$L(s,\psi)=L(s,\chi)\prod_{p\mid q_0}(1-\chi(p)p^{-s})$; hence
$L(\rho,\chi)\ne0$.
\end{proof}

\begin{corollary}[the generalised Riemann hypothesis for Dirichlet $L$-functions]\label{an:cor-GRH}
Let $\chi$ be a Dirichlet character modulo $q\ge1$, primitive or not. Every
zero $\rho$ of $L(s,\chi)$ with $0<\Re\rho<1$ has $\Re\rho=\frac12$.
\end{corollary}
\begin{proof}
Let $\chi^{*}$ be the primitive character inducing $\chi$. The factors
$1-\chi^{*}(p)p^{-s}$ vanish only on $\Re s=0$, so $L(\rho,\chi^{*})=0$. By
the functional equation \cite[Chapters~8 and~9]{Davenport}, whose gamma factor has
neither zeros nor poles in $0<\Re s<1$, $L(1-\rho,\overline{\chi^{*}})=0$.
Theorem~\ref{thm:GRH-half} for $\chi^{*}$ and for $\overline{\chi^{*}}$
gives $\Re\rho\le\frac12$ and $1-\Re\rho\le\frac12$.
\end{proof}

\begin{corollary}\label{an:cor-zeta}
The zero condition \eqref{eq:dirichlet-input} holds: $\zeta(s)\ne0$ for
$\Re s>\frac12$, $s\ne1$.
\end{corollary}
\begin{proof}
Theorem~\ref{thm:GRH-half} for $q=1$.
\end{proof}

\begin{remark}[what Sections~\ref{an:sec-criterion}--\ref{an:sec-main} prove]\label{an:rem-status}
These sections prove the implication
\[
 \bigl(\text{Theorems \anThmA, \anThmB, \anThmC\ for all admissible data}\bigr)\ \Longrightarrow\ \text{Theorem~\ref{thm:GRH-half}},
\]
in fact they use the three theorems only for the data of
Theorem~\ref{an:thm-data}. The proof of the implication is complete in the
text, given the published theorems named at the beginning of
Section~\ref{an:sec-criterion}; apart from these the arguments are
elementary. The data are obtained cheaply, from the
theorem of Siegel and Walfisz, Behrend's construction, the pigeonhole
principle in the class group of $K$ and the theorem on sums of two squares;
the conditions imposed on them (sparseness, the congruences modulo $q_2^{2}$,
the window for $\log N$) are consumed in the automorphic part. Accordingly
the present sections give no independent evidence for
Theorems~\anThmA, \anThmB, \anThmC. These are proved in
Section~\ref{sec:automorphic} and the sections that follow it, from published
results (\S\ref{ss:au-remains}); Theorems~\ref{an:thm-main},
\ref{an:thm-onesided} and~\ref{thm:GRH-half} and
Corollaries~\ref{an:cor-GRH} and~\ref{an:cor-zeta} rest on them. The main
inequality cannot be tested numerically: (D2) forces
$\log x\ge10^{14}\anCP$.
\end{remark}



% conclusion-new.tex -- Conclusion of the paper, adapted to the new second part.
% Labels kept: sec:conclusion, thm:RH, thm:all-support, cor:all-windows.
% Changes with respect to paper46-original.tex:
%  - proof of thm:RH: now from Theorem thm:GRH-half (Corollary an:cor-zeta);
%  - thm:all-support: the hypothesis on the prime sum G (which no longer exists) is
%    replaced by the zero condition (eq:dirichlet-input); the first paragraph of the
%    proof and the sentence on J_theta = 3/2 are deleted; the rest is verbatim;
%  - cor:all-windows: "Assume Theorem thm:alpoge-families" becomes "Assume Theorems A, B, C";
%  - the two closing subsections are rewritten.
% The closing paragraphs follow sec-automorphic.tex, \S ss:au-remains, as of 26 Nov 2026.
\providecommand{\anThmA}{\ref{thm:A}}
\providecommand{\anThmB}{\ref{thm:B}}
\providecommand{\anThmC}{\ref{thm:C}}
\section{Conclusion}\label{sec:conclusion}
\begin{theorem}[the Riemann hypothesis]\label{thm:RH}
Every nontrivial zero of the Riemann zeta function has real part $1/2$.
\end{theorem}
\begin{proof}
By Theorem~\ref{thm:GRH-half} for the trivial character of modulus one
(Corollary~\ref{an:cor-zeta}), $\zeta$ has no
zero in $\Re s>1/2$. If a nontrivial zero $\rho$ had $\Re\rho<1/2$, the
functional equation $\xi(s)=\xi(1-s)$ of the completed zeta function would
give the zero $1-\rho$, with $\Re(1-\rho)>1/2$. Hence every nontrivial zero
has real part $1/2$.
\end{proof}
The proof does not use \cite{P43}, \cite{P44} or \cite{Paper45}. In the
setting of the four rows described in the introduction, the theorem gives
the sign of the radial intersection form on every window; this is
Corollary~\ref{cor:all-windows} below, obtained through
Theorem~\ref{thm:all-support}. There \cite{Paper45} enters only through the
converse implication, which the corollary does not use.

\subsection*{The sign of the form on every window}
\begin{theorem}\label{thm:all-support}
For $\zeta_{\mathbb Q}$ the Riemann hypothesis is equivalent both to
$W_T\ge0$ on every $\mathcal D_T$ and to $\lambda_{\rm A}(T)\ge0$ for all $T$.
Suppose that the zero condition \eqref{eq:dirichlet-input} holds. Then for
every
$T>0$ and every $F\in\mathcal D_T$ one has
\[
 W_T(F)=\sum_{\rho}|\Phi_F(\rho)|^2\ge0 .
\]
Zeros are counted with multiplicities. Hence
$0\le\lambda^{\rm W}_T\le\lambda_{\rm A}(T)$ and
$I_\infty(\Gamma_F,\Gamma_F)\le0$ for smooth primitive tests.
\end{theorem}
\begin{proof}
If a nontrivial zero $\rho$ lay to the left of the critical line,
the functional equation would give a zero $1-\rho$ to its right,
contrary to \eqref{eq:dirichlet-input}. Thus the zero condition
implies the Riemann hypothesis.

For the forward implication of the equivalence, assume now that every
nontrivial zero lies on the critical line. For a smooth real test,
$1-\rho=\bar\rho$ and $\Phi_F(\bar\rho)=\overline{\Phi_F(\rho)}$.
The zero sum in \eqref{comparison} is therefore
$\sum_\rho |\Phi_F(\rho)|^2$.
For smooth real $U,V$, set $u_\rho=\Phi_U(\rho)$ and
$v_\rho=\Phi_V(\rho)$. Conjugation pairs the summands in
$\sum_\rho u_\rho\overline{v_\rho}$, so this absolutely convergent
sum is real. Expanding $|u_\rho+i v_\rho|^2$ and summing therefore
gives $W_T(U)+W_T(V)=W_T(U+iV)$ by complexification.
For $F\in\mathcal D_T$, take smooth $F_n$ converging in form norm
(Lemma~\ref{lem:form-density}(i)). Applying the identity to $F_n-F_m$,
for which $W_T(F_n-F_m)\to0$ by Lemma~\ref{lem:form-density}(ii),
makes $(\Phi_{F_n}(\rho))_\rho$ Cauchy in $\ell^2$, with one
coordinate for each occurrence of a zero. Since
$|e^{(\rho-1/2)t}|=1$, Cauchy--Schwarz gives
$|\Phi_{F_n-F}(\rho)|\le\sqrt{2T}\,\|F_n-F\|_2$.
This identifies every coordinate of the limit. Taking squared norms,
and using $W_T(F_n)\to W_T(F)$ (Lemma~\ref{lem:form-density}(ii)),
gives the formula for $F$. On $\mathcal P_T$ the polar term
vanishes, giving both inequalities. Conversely an off-line zero gives
a smooth primitive test with negative zero sum
\cite[Theorem~6.2]{Paper45}; both infima are negative on any window
containing its support.
\end{proof}

\begin{corollary}[the sign on every window]\label{cor:all-windows}
For every $T>0$ and every
$F\in\mathcal D_T$,
\[
 W_T(F)=\sum_\rho|\Phi_F(\rho)|^2\ge0 ,
\]
so $0\le\lambda^{\rm W}_T\le\lambda_{\rm A}(T)$ for every $T>0$, and
$I_\infty(\Gamma_F,\Gamma_F)\le0$ for smooth primitive tests of arbitrary
support. In particular the positivity obtained by computation in
Section~\ref{sec:log5} (Theorem~\ref{thm:log5W}, in ball arithmetic, for
$T\le\frac12\log5$; and the certificate \eqref{eq:power-window}, reproduced
there in Arb arithmetic, for $T\le2.8022$) holds, in the form
$\lambda^{\rm W}_T\ge0$, on every window $(-T,T)$.
\end{corollary}
\begin{proof}
Theorem~\ref{thm:GRH-half}, which is deduced from Theorems~\anThmA,
\anThmB\ and~\anThmC\ in Section~\ref{an:sec-main}, gives the zero condition
\eqref{eq:dirichlet-input} (Corollary~\ref{an:cor-zeta}). This is the
hypothesis of Theorem~\ref{thm:all-support}. The inequality itself comes
from the identity \eqref{comparison}, Weil's explicit formula
\cite{Weil52,Bombieri2000} (Section~\ref{sec:trace}), through the second paragraph of the proof of
Theorem~\ref{thm:all-support}, applied once Theorem~\ref{thm:RH} is known.
\end{proof}

\subsection*{Closing remarks}
The chain is now complete. The criterion of Section~\ref{an:sec-criterion}
reduces the absence of zeros in $\Re s>\beta$ to a one-sided bound for the
prime sum $R_q$; the data of Section~\ref{an:sec-data} and
Theorems~\anThmA, \anThmB\ and~\anThmC\ give the main inequality
(Theorem~\ref{an:thm-main}), and from it that bound for every
$\beta>\frac12$ (Theorem~\ref{an:thm-onesided}). Hence no Dirichlet
$L$-function vanishes in $\Re s>\frac12$ (Theorem~\ref{thm:GRH-half}): the
generalised Riemann hypothesis holds for Dirichlet $L$-functions
(Corollary~\ref{an:cor-GRH}), and with it the Riemann hypothesis
(Theorem~\ref{thm:RH}). In the language of the four rows, the index
inequality of Row~D holds on every window (Corollary~\ref{cor:all-windows}).

\subsection*{Acknowledgement}
The work of L.~Alpoge \cite{CharacterFamily} is, by itself, a complete route
to the generalised Riemann hypothesis for Dirichlet $L$-functions and to the
Riemann hypothesis; it does not depend on anything in this series. It is also
the input that closes the path followed here, which is to carry out Weil's
proof over $\operatorname{Spec}\mathbb Z$: the mathematics behind
Theorems~\anThmA, \anThmB\ and~\anThmC, and the way from them to
Theorem~\ref{thm:GRH-half}, come from that work, while the division of the
argument into these three theorems, and their proofs as written here, are
our way of organising it as mathematics on paper. It is a pleasure to record
how much this paper owes to that contribution.

\begin{thebibliography}{99}
\bibitem{CharacterFamily} L.~Alpoge, Lean~4 formalisation of the character-family argument
for Dirichlet $L$-functions, source archive in fourteen parts, November 2026,
\url{https://alpo.ge/lean.tar.gz.part01} (first part).

\bibitem{BoasKac1945} R.~P.~Boas and M.~Kac, \emph{Inequalities for Fourier transforms of positive
functions}, Duke Math.~J. \textbf{12} (1945), 189--206.

\bibitem{Bombieri2000} E.~Bombieri, \emph{Remarks on Weil's quadratic functional in the theory of
prime numbers, I}, Atti Accad.~Naz.~Lincei Cl.~Sci.~Fis.~Mat.~Natur.~Rend.~Lincei (9) Mat.~Appl.
\textbf{11} (2000), no.~3, 183--233.

\bibitem{BorelHC1962} A.~Borel and Harish-Chandra, \emph{Arithmetic subgroups of algebraic groups},
Ann.\ of Math.\ (2) \textbf{75} (1962), 485--535.


\bibitem{ConnesConsani2021} A.~Connes and C.~Consani, \emph{Weil positivity and trace formula, the archimedean place}, Selecta Math.\ (N.S.) \textbf{27} (2021), no.~4, Paper No.~77; arXiv:2006.13771.

\bibitem{CvS2025} A.~Connes and W.~D.~van Suijlekom, \emph{Quadratic forms, real zeros and echoes of the spectral action}, arXiv:2511.23257 (2025).

\bibitem{Davenport} H.~Davenport, \emph{Multiplicative Number Theory}, 3rd ed., revised by
H.~L.~Montgomery, Graduate Texts in Mathematics \textbf{74}, Springer, New York, 2000.

\bibitem{DLMF} NIST Digital Library of Mathematical Functions, \url{https://dlmf.nist.gov}, \S5.9.


\bibitem{Groskin2026} A.~Groskin, \emph{A finite Guinand--Weil dictionary and archimedean tail
order for the truncated Weil quadratic form}, arXiv:2607.02828 (2026).

\bibitem{HLP} G.~H.~Hardy, J.~E.~Littlewood and G.~P\'olya, \emph{Inequalities}, 2nd ed., Cambridge Univ.~Press, 1952.

\bibitem{Haynsworth1968} E.~V.~Haynsworth, \emph{Determination of the inertia of a partitioned Hermitian
matrix}, Linear Algebra Appl.\ \textbf{1} (1968), 73--81.

\bibitem{IK} H.~Iwaniec and E.~Kowalski, \emph{Analytic Number Theory}, Amer.\ Math.\ Soc.
Colloq.\ Publ.\ \textbf{53}, Providence, RI, 2004.

\bibitem{Johansson2017} F.~Johansson, \emph{Arb: efficient arbitrary-precision midpoint-radius interval
arithmetic}, IEEE Trans.\ Comput.\ \textbf{66} (2017), 1281--1292.

\bibitem{Kim2026} T.~Kim, Y.~Hong, M.~Kim, S.~Choi, J.~Jang and M.~Kim, \emph{A numerical realization of Suzuki's Weil-quadratic-form operator: the archimedean spectral law, its universality, and an operator form of Weil's positivity criterion}, arXiv:2607.24830 (2026).

\bibitem{Krein1940} M.~G.~Kre\u{\i}n, \emph{Sur le probl\`eme du prolongement des fonctions
hermitiennes positives et continues}, Dokl.~Akad.~Nauk SSSR \textbf{26} (1940), 17--22.

\bibitem{Liu2026} V.~Liu, \emph{Certified Weil positivity beyond the unit window: source-exact
block-Schur and tail-compensation bounds for the Riemann zeta function}, manuscript dated 14 Sep 2026,
alphaXiv preprint.

\bibitem{MoriUpper} R.~Mori, \emph{Improved upper bounds for windowed Weil forms with exponential factor
$e^{-4\pi\mu}$}, Zenodo (2026), v3, doi:\nolinkurl{10.5281/zenodo.23363767}.

\bibitem{PlattTrudgian2021} D.~J.~Platt and T.~S.~Trudgian, \emph{The Riemann hypothesis is true up to
$3\cdot10^{12}$}, Bull.\ London Math.\ Soc.\ \textbf{53} (2021), 792--797.

\bibitem{Rockafellar} R.~T.~Rockafellar, \emph{Convex Analysis}, Princeton Univ.~Press, 1970.

\bibitem{Rosser1941} J.~B.~Rosser, \emph{Explicit bounds for some functions of prime numbers}, Amer.\ J.\
Math.\ \textbf{63} (1941), 211--232.

\bibitem{RosserYoheSchoenfeld1968} J.~B.~Rosser, J.~M.~Yohe and L.~Schoenfeld, \emph{Rigorous computation and
the zeros of the Riemann zeta-function}, Proc.\ IFIP Congress 1968, 70--76.

\bibitem{SSV2012} R.~L.~Schilling, R.~Song and Z.~Vondra\v{c}ek, \emph{Bernstein Functions:
Theory and Applications}, 2nd ed., de Gruyter, 2012.

\bibitem{Sasvari1994} Z.~Sasv\'ari, \emph{Positive Definite and Definitizable Functions},
Akademie Verlag, Berlin, 1994.

\bibitem{SongVondracek2003} R.~Song and Z.~Vondra\v{c}ek, \emph{Potential theory of subordinate
killed Brownian motion in a domain}, Probab.~Theory Related Fields \textbf{125} (2003), 578--592.


\bibitem{Trefethen2013} L.~N.~Trefethen, \emph{Approximation Theory and Approximation Practice},
SIAM, Philadelphia, 2013.

\bibitem{P43} D.~A.~Trejo Pizzo,
\emph{A correspondence algebra and the Weil distribution on semilocal idele
class groups}, preprint, 2026.

\bibitem{P44} D.~A.~Trejo Pizzo,
\emph{An arithmetic Frobenius family on semilocal idele class groups},
preprint, 2026.

\bibitem{Paper45} D.~A.~Trejo Pizzo, \emph{An explicit formula for the trace on semilocal idele class groups}, preprint, 2026.


\bibitem{Weil52} A.~Weil, \emph{Sur les ``formules explicites'' de la th\'eorie des nombres
premiers}, Comm.~S\'em.~Math.~Univ.~Lund (1952), 252--265.


\bibitem{Yoshida1992} H.~Yoshida, \emph{On Hermitian forms attached to zeta functions}, in:
\emph{Zeta Functions in Geometry}, Adv.~Stud.~Pure Math.~\textbf{21}, Kinokuniya, Tokyo, 1992, 281--325.

\bibitem{Zhu2026} X.~Zhu, \emph{Weil positivity in compact windows: a finite reduction, certified
two-sided bounds, and a Landau--Widom decay law}, arXiv:2608.24827 (2026).
\bibitem{Behrend1946} F.~A.~Behrend, \emph{On sets of integers which contain no three terms in
arithmetical progression}, Proc.\ Nat.\ Acad.\ Sci.\ U.S.A.\ \textbf{32} (1946), 331--332.
\bibitem{HeathBrown1995} D.~R.~Heath-Brown, \emph{A mean value estimate for real character
sums}, Acta Arith.\ \textbf{72} (1995), no.~3, 235--275.
\bibitem{MV} H.~L.~Montgomery and R.~C.~Vaughan, \emph{Multiplicative Number Theory I.
Classical Theory}, Cambridge Studies in Advanced Mathematics \textbf{97}, Cambridge
Univ.~Press, 2007.
\bibitem{Washington} L.~C.~Washington, \emph{Introduction to Cyclotomic Fields}, 2nd ed.,
Graduate Texts in Mathematics \textbf{83}, Springer, New York, 1997.

% refs-automorphic.tex: \bibitem lines for sec-automorphic.tex
\bibitem{Banaszczyk} W.~Banaszczyk, \emph{New bounds in some transference theorems in the geometry
of numbers}, Math.\ Ann.\ \textbf{296} (1993), 625--635.
\bibitem{BannaiKobayashi} K.~Bannai and S.~Kobayashi, \emph{Algebraic theta functions and the
$p$-adic interpolation of Eisenstein--Kronecker numbers}, Duke Math.\ J.\ \textbf{153} (2010),
229--295.
\bibitem{BostFoliations} J.-B.~Bost, \emph{Algebraic leaves of algebraic foliations over number
fields}, Publ.\ Math.\ Inst.\ Hautes \'Etudes Sci.\ \textbf{93} (2001), 161--221.
\bibitem{Chellali} M.~Chellali, \emph{Congruences entre nombres de Bernoulli--Hurwitz dans le cas
supersingulier}, J.\ Number Theory \textbf{35} (1990), 157--179.
\bibitem{ChowlaSelberg} S.~Chowla and A.~Selberg, \emph{On Epstein's zeta-function}, J.\ reine
angew.\ Math.\ \textbf{227} (1967), 86--110.
\bibitem{Finis} T.~Finis, \emph{Divisibility of anticyclotomic $L$-functions and theta functions
with complex multiplication}, Ann.\ of Math.\ (2) \textbf{163} (2006), 767--807.
\bibitem{GomezZhu} R.~Gomez and C.-B.~Zhu, \emph{Local theta lifting of generalized Whittaker
models associated to nilpotent orbits}, Geom.\ Funct.\ Anal.\ \textbf{24} (2014), 796--853.
\bibitem{Green1955} J.~A.~Green, \emph{The characters of the finite general linear groups},
Trans.\ Amer.\ Math.\ Soc.\ \textbf{80} (1955), 402--447.
\bibitem{GrossAMF} B.~H.~Gross, \emph{Algebraic modular forms}, Israel J.\ Math.\ \textbf{113}
(1999), 61--93.
\bibitem{HidaNonvanishing} H.~Hida, \emph{Non-vanishing modulo $p$ of Hecke $L$-values}, in:
Geometric Aspects of Dwork Theory, Vol.~II, de Gruyter, 2004, 735--784.
\bibitem{HoweTheta} R.~Howe, \emph{$\theta$-series and invariant theory}, in: Automorphic Forms,
Representations and $L$-functions, Proc.\ Sympos.\ Pure Math.\ \textbf{33}, Part~1, Amer.\ Math.\
Soc., 1979, 275--285.
\bibitem{KashiwaraVergne} M.~Kashiwara and M.~Vergne, \emph{On the Segal--Shale--Weil
representations and harmonic polynomials}, Invent.\ Math.\ \textbf{44} (1978), 1--47.
\bibitem{KatzDivisibilities} N.~M.~Katz, \emph{Divisibilities, congruences, and Cartier duality},
J.\ Fac.\ Sci.\ Univ.\ Tokyo Sect.\ IA Math.\ \textbf{28} (1981), 667--678.
\bibitem{KudlaSeesaw} S.~S.~Kudla, \emph{Seesaw dual reductive pairs}, in: Automorphic Forms of
Several Variables (Katata, 1983), Progr.\ Math.\ \textbf{46}, Birkh\"auser, 1984, 244--268.
\bibitem{Labesse} J.-P.~Labesse, \emph{Changement de base CM et s\'eries discr\`etes}, in: On the
Stabilization of the Trace Formula, Int.\ Press, 2011, 429--470.
\bibitem{LaurentInterp} M.~Laurent, \emph{Linear forms in two logarithms and interpolation
determinants}, Acta Arith.\ \textbf{66} (1994), 181--199.
\bibitem{Prasanna} K.~Prasanna, \emph{Integrality of a ratio of Petersson norms and
level-lowering congruences}, Ann.\ of Math.\ (2) \textbf{163} (2006), 901--967.
\bibitem{RallisHoweDuality} S.~Rallis, \emph{On the Howe duality conjecture}, Compositio Math.\ \textbf{51} (1984), 333--399.
\bibitem{RallisInjectivity} S.~Rallis, \emph{Injectivity properties of liftings associated to
Weil representations}, Compositio Math.\ \textbf{52} (1984), 139--169.
\bibitem{Ribet1976} K.~A.~Ribet, \emph{A modular construction of unramified $p$-extensions of
$\mathbb Q(\mu_p)$}, Invent.\ Math.\ \textbf{34} (1976), 151--162.

% References used only by proofs/A-divisibility.tex. (BannaiKobayashi and KatzDivisibilities are
% already in refs-automorphic.tex.)
\bibitem{BLR} S.~Bosch, W.~L\"utkebohmert and M.~Raynaud, \emph{N\'eron Models}, Ergebnisse der
Mathematik und ihrer Grenzgebiete (3) \textbf{21}, Springer, 1990.
\bibitem{DEP} C.~De~Concini, D.~Eisenbud and C.~Procesi, \emph{Young diagrams and determinantal
varieties}, Invent.\ Math.\ \textbf{56} (1980), 129--165.
\bibitem{LangEF} S.~Lang, \emph{Elliptic Functions}, 2nd ed., Graduate Texts in Mathematics
\textbf{112}, Springer, 1987.
\bibitem{LubinTate} J.~Lubin and J.~Tate, \emph{Formal complex multiplication in local fields},
Ann.\ of Math.\ (2) \textbf{81} (1965), 380--387.
\bibitem{SilvermanAEC} J.~H.~Silverman, \emph{The Arithmetic of Elliptic Curves}, 2nd ed., Graduate
Texts in Mathematics \textbf{106}, Springer, 2009.
\bibitem{SilvermanAT} J.~H.~Silverman, \emph{Advanced Topics in the Arithmetic of Elliptic Curves},
Graduate Texts in Mathematics \textbf{151}, Springer, 1994.

% refs-A-seesaw-size.tex: \bibitem lines used by A-seesaw-size.tex.
% Already in refs-automorphic.tex (not repeated in the assembled paper): Banaszczyk,
% BannaiKobayashi, ChowlaSelberg, GomezZhu, HoweTheta, KashiwaraVergne.
\bibitem{HKS} M.~Harris, S.~S.~Kudla and W.~J.~Sweet, \emph{Theta dichotomy for unitary groups},
J.\ Amer.\ Math.\ Soc.\ \textbf{9} (1996), 941--1004.
\bibitem{KudlaSplitting} S.~S.~Kudla, \emph{Splitting metaplectic covers of dual reductive pairs},
Israel J.\ Math.\ \textbf{87} (1994), 361--401.
\bibitem{WeilActa} A.~Weil, \emph{Sur certains groupes d'op\'erateurs unitaires}, Acta Math.\
\textbf{111} (1964), 143--211.

% refs-A-seesaw-complete.tex: \bibitem lines used by A-seesaw-complete.tex.
% Already present elsewhere (not to be repeated in the assembled paper):
%   refs-automorphic.tex: ChowlaSelberg, HoweTheta, KashiwaraVergne;
%   refs-A-seesaw-size.tex: HKS, KudlaSplitting, WeilActa;
%   refs-A-divisibility.tex: BLR, DEP, LangEF, SilvermanAEC.
\bibitem{FultonYT} W.~Fulton, \emph{Young Tableaux}, London Mathematical Society Student Texts
\textbf{35}, Cambridge University Press, 1997.
\bibitem{GoodmanWallach} R.~Goodman and N.~R.~Wallach, \emph{Symmetry, Representations, and
Invariants}, Graduate Texts in Mathematics \textbf{255}, Springer, 2009.
\bibitem{KubertLang} D.~S.~Kubert and S.~Lang, \emph{Modular Units}, Grundlehren der mathematischen
Wissenschaften \textbf{244}, Springer, 1981.
\bibitem{MumfordAV} D.~Mumford, \emph{Abelian Varieties}, Tata Institute of Fundamental Research
Studies in Mathematics \textbf{5}, Oxford University Press, 1970.
\bibitem{PlatonovRapinchuk} V.~Platonov and A.~Rapinchuk, \emph{Algebraic Groups and Number
Theory}, Pure and Applied Mathematics \textbf{139}, Academic Press, 1994.
\bibitem{TitsCorvallis} J.~Tits, \emph{Reductive groups over local fields}, in: Automorphic Forms,
Representations and $L$-Functions (Corvallis, 1977), Proc.\ Sympos.\ Pure Math.\ \textbf{33},
Part~1, Amer.\ Math.\ Soc., 1979, pp.~29--69.

% References used only by proofs/A-hecke-eldiv.tex.
\bibitem{LTXZ} R.~Liu, N.~X.~Truong, L.~Xiao and B.~Zhao, \emph{Slopes of modular forms and
geometry of eigencurves}, preprint, arXiv:2302.07697 (2023).
\bibitem{RallisFunctoriality} S.~Rallis, \emph{Langlands' functoriality and the Weil
representation}, Amer.\ J.\ Math.\ \textbf{104} (1982), 469--515.

% References used only by proofs/A-eldiv-inputs.tex.
\bibitem{AndrewsPartitions} G.~E.~Andrews, \emph{The Theory of Partitions}, Encyclopedia of
Mathematics and its Applications \textbf{2}, Addison--Wesley, 1976.
\bibitem{StanleyEC2} R.~P.~Stanley, \emph{Enumerative Combinatorics}, vol.~2, Cambridge Studies in
Advanced Mathematics \textbf{62}, Cambridge University Press, 1999.
\bibitem{pe-DLMF17} NIST Digital Library of Mathematical Functions, Chapter~17 ($q$-hypergeometric
and related functions), \url{https://dlmf.nist.gov/17}.

% refs-A-hecke-level.tex: references cited by proofs/A-hecke-level.tex.
% All three are already present in the bibliography of the assembled paper
% (HKS and KudlaSplitting in refs-A-seesaw-size.tex, RallisFunctoriality in
% refs-A-hecke-eldiv.tex); they are repeated here, commented out, so that this file is
% self-describing and no \bibitem is duplicated when the files are assembled.
%
\bibitem{CoxPrimes} D.~A.~Cox, \emph{Primes of the Form $x^2+ny^2$: Fermat, Class Field Theory, and
Complex Multiplication}, Wiley, New York, 1989.
\bibitem{SerreLF} J.-P.~Serre, \emph{Local Fields}, Graduate Texts in Mathematics 67, Springer,
New York, 1979.

% refs-A-closure.tex: \bibitem lines used by A-closure.tex and not present in the other refs-*.tex files.
% (Also cited there, from the other files: WeilActa, KashiwaraVergne, BorelHC1962.)
\bibitem{HarishChandraAF} Harish-Chandra, \emph{Automorphic Forms on Semisimple Lie Groups}, Lecture
Notes in Mathematics 62, Springer, Berlin, 1968.
\bibitem{BorelJacquet} A.~Borel and H.~Jacquet, \emph{Automorphic forms and automorphic
representations}, in: Automorphic Forms, Representations and $L$-functions, Proc.\ Sympos.\ Pure
Math.\ \textbf{33}, Part~1, Amer.\ Math.\ Soc., 1979, 189--202.
\bibitem{HoweJAMS} R.~Howe, \emph{Transcending classical invariant theory}, J.\ Amer.\ Math.\ Soc.\
\textbf{2} (1989), 535--552.

% References for proofs/A-archimedean.tex (cite keys shared with the paper where they exist)
\bibitem{FJ} M.~Flensted-Jensen, Discrete series for semisimple symmetric spaces, Ann. of Math. (2) 111 (1980), 253--311.
\bibitem{Knapp} A.~W.~Knapp, Representation Theory of Semisimple Groups: An Overview Based on Examples, Princeton Univ. Press, 1986.
\bibitem{Parthasarathy72} R.~Parthasarathy, Dirac operator and the discrete series, Ann. of Math. (2) 96 (1972), 1--30.
\bibitem{HuangPandzic} J.-S.~Huang and P.~Pand\v zi\'c, Dirac cohomology, unitary representations and a proof of a conjecture of Vogan, J. Amer. Math. Soc. 15 (2002), 185--202.
\bibitem{Schmid75} W.~Schmid, Some properties of square-integrable representations of semisimple Lie groups, Ann. of Math. (2) 102 (1975), 535--564.
\bibitem{Dixmier} J.~Dixmier, $C^*$-algebras, North-Holland, 1977.
\bibitem{VoganZuckerman} D.~A.~Vogan, Jr. and G.~J.~Zuckerman, Unitary representations with non-zero cohomology, Compositio Math. 53 (1984), 51--90.

% References for proofs/A-coefficient.tex (cite keys shared with the paper where they exist)
\bibitem{KnappBeyond} A.~W.~Knapp, Lie Groups Beyond an Introduction, 2nd ed., Birkh\"auser, 2002.
\bibitem{FollandAHA} G.~B.~Folland, A Course in Abstract Harmonic Analysis, CRC Press, 1995.
\bibitem{FollandPS} G.~B.~Folland, Harmonic Analysis in Phase Space, Ann. of Math. Stud. 122, Princeton Univ. Press, 1989.
\bibitem{GHY} W.~T.~Gan, J.~P.~Hanke and J.-K.~Yu, \emph{On an exact mass formula of Shimura},
Duke Math.\ J.\ \textbf{107} (2001), 103--133.
\bibitem{Macdonald} I.~G.~Macdonald, \emph{Symmetric functions and Hall polynomials}, 2nd ed.,
Oxford University Press, 1995.


% refs-B-bessel-expansion.tex: references cited by proofs/B-bessel-expansion.tex.
\bibitem{BZ76} I.~N.~Bernstein and A.~V.~Zelevinsky, \emph{Representations of the group
$GL(n,F)$ where $F$ is a non-archimedean local field}, Russian Math.\ Surveys \textbf{31}:3
(1976), 1--68.
\bibitem{BorelFin} A.~Borel, \emph{Some finiteness properties of adele groups over number
fields}, Publ.\ Math.\ Inst.\ Hautes \'Etudes Sci.\ \textbf{16} (1963), 5--30.
\bibitem{FlickerKazhdan} Y.~Z.~Flicker and D.~A.~Kazhdan, \emph{A simple trace formula},
J.\ Analyse Math.\ \textbf{50} (1988), 189--200.
\bibitem{Godement63} R.~Godement, \emph{Domaines fondamentaux des groupes arithm\'etiques},
S\'eminaire Bourbaki, vol.~8 (1962--1964), expos\'e 257, 201--225 (Th\'eor\`eme~7, p.~211);
\url{http://www.numdam.org/item/SB_1962-1964__8__201_0}.
\bibitem{HC66} Harish-Chandra, \emph{Discrete series for semisimple Lie groups II. Explicit
determination of the characters}, Acta Math.\ \textbf{116} (1966), 1--111.
\bibitem{Helgason} S.~Helgason, \emph{Groups and Geometric Analysis}, Academic Press, 1984.
\bibitem{Jacobowitz} R.~Jacobowitz, \emph{Hermitian forms over local fields}, Amer.\ J.\ Math.\
\textbf{84} (1962), 441--465.
\bibitem{Landherr} W.~Landherr, \emph{\"Aquivalenz Hermitescher Formen \"uber einem beliebigen
algebraischen Zahlk\"orper}, Abh.\ Math.\ Sem.\ Univ.\ Hamburg \textbf{11} (1935), 245--248.
\bibitem{MW} C.~M\oe glin and J.-L.~Waldspurger, \emph{Spectral Decomposition and Eisenstein
Series}, Cambridge Tracts in Math.\ \textbf{113}, Cambridge University Press, 1995.
\bibitem{Neukirch} J.~Neukirch, \emph{Algebraic Number Theory}, Grundlehren Math.\ Wiss.\
\textbf{322}, Springer, 1999.
\bibitem{Scharlau} W.~Scharlau, \emph{Quadratic and Hermitian Forms}, Grundlehren Math.\ Wiss.\
\textbf{270}, Springer, 1985.

% refs-B-orbital.tex: new references of proofs/B-orbital.tex.
% None: HC66, Knapp (refs-B-multiplicities), KnappBeyond (refs-A-coefficient), Schmid75
% (refs-A-archimedean) and HornJohnson (refs-B-discrete-series) are already in the bibliography.

% refs-C-determinant.tex: \bibitem lines used by C-determinant.tex.
% Already present elsewhere (not to be repeated in the assembled paper):
%   refs-automorphic.tex: GQT, KashiwaraVergne;
%   refs-A-seesaw-size.tex: KudlaSplitting, WeilActa, HKS.
\bibitem{IchinoSW} A.~Ichino, \emph{A regularized Siegel--Weil formula for unitary groups}, Math.\
Z.\ \textbf{247} (2004), 241--277.
\bibitem{WeilSiegel} A.~Weil, \emph{Sur la formule de Siegel dans la th\'eorie des groupes
classiques}, Acta Math.\ \textbf{113} (1965), 1--87.

% refs-C-onelift.tex: \bibitem lines used by C-onelift.tex.
% Already present elsewhere (not to be repeated in the assembled paper):
%   refs-automorphic.tex: KashiwaraVergne, RallisInjectivity;
%   refs-A-seesaw-size.tex: KudlaSplitting.
\bibitem{Hecke1920} E.~Hecke, \emph{Eine neue Art von Zetafunktionen und ihre Beziehungen zur Verteilung der Primzahlen. (Zweite Mitteilung)}, Math.\ Z.\ \textbf{6} (1920), 11--51.
\bibitem{FiteSatoTate} F.~Fit\'e, \emph{Equidistribution, $L$-functions, and Sato--Tate groups}, preprint, arXiv:1405.5162 (2014).

% refs-C-congruences.tex: references cited by proofs/C-congruences.tex.
% Labesse and Ribet1976 are already in refs-automorphic.tex and are not repeated here.
\bibitem{BLGGT} T.~Barnet-Lamb, T.~Gee, D.~Geraghty and R.~Taylor, \emph{Potential automorphy and
change of weight}, Ann.\ of Math.\ (2) \textbf{179} (2014), 501--609.
\bibitem{BlochKato} S.~Bloch and K.~Kato, \emph{$L$-functions and Tamagawa numbers of motives},
in: The Grothendieck Festschrift, Vol.~I, Progr.\ Math.\ \textbf{86}, Birkh\"auser, 1990,
333--400.
\bibitem{FontaineLaffaille} J.-M.~Fontaine and G.~Laffaille, \emph{Construction de
repr\'esentations $p$-adiques}, Ann.\ Sci.\ \'Ecole Norm.\ Sup.\ (4) \textbf{15} (1982), 547--608.
\bibitem{NSW} J.~Neukirch, A.~Schmidt and K.~Wingberg, \emph{Cohomology of number fields}, 2nd
ed., Grundlehren Math.\ Wiss.\ \textbf{323}, Springer, 2008.
\bibitem{WilesFLT} A.~Wiles, \emph{Modular elliptic curves and Fermat's last theorem}, Ann.\ of
Math.\ (2) \textbf{141} (1995), 443--551.

% refs-C-selmer.tex: references cited by proofs/C-selmer.tex and not present in the other refs-*.tex
% files. SerreLF is in refs-A-final.tex; NSW is in refs-C-congruences.tex.
\bibitem{MilneADT} J.~S.~Milne, \emph{Arithmetic Duality Theorems}, Academic Press, Boston, 1986 (1st ed.; all numbers cited refer to this edition), \url{https://www.jmilne.org/math/Books/adt1.pdf}.

% refs-C-ribet.tex: new references cited by proofs/C-ribet.tex (NSW is in refs-C-congruences.tex).
\bibitem{Amitsur} S.~A.~Amitsur, \emph{On the characteristic polynomial of a sum of matrices},
Linear and Multilinear Algebra \textbf{8} (1979/80), 177--182.
\bibitem{LangCyclo} S.~Lang, \emph{Cyclotomic fields I and II}, combined 2nd ed., with an appendix
by K.~Rubin, Grad.\ Texts in Math.\ \textbf{121}, Springer, 1990.

% New bibliography items for proofs/C-localloss.tex
\bibitem{RubinES} K.~Rubin, \emph{Euler systems}, Annals of Mathematics Studies \textbf{147},
Princeton University Press, 2000.

% refs-C-galois.tex: references cited by proofs/C-galois.tex and not already in refs-automorphic.tex
% (Labesse) or refs-C-congruences.tex (BLGGT). All cited from memory; volume/page data to be checked.
\bibitem{HarrisTaylor} M.~Harris and R.~Taylor, \emph{The geometry and cohomology of some simple
Shimura varieties}, Ann.\ of Math.\ Stud.\ \textbf{151}, Princeton Univ.\ Press, 2001.
\bibitem{Shin} S.~W.~Shin, \emph{Galois representations arising from some compact Shimura
varieties}, Ann.\ of Math.\ (2) \textbf{173} (2011), 1645--1741.
\bibitem{ChenevierHarris} G.~Chenevier and M.~Harris, \emph{Construction of automorphic Galois
representations, II}, Camb.\ J.\ Math.\ \textbf{1} (2013), 53--73.
\bibitem{CaraianiDuke} A.~Caraiani, \emph{Local-global compatibility and the action of monodromy
on nearby cycles}, Duke Math.\ J.\ \textbf{161} (2012), 2311--2413.
\bibitem{CaraianiANT} A.~Caraiani, \emph{Monodromy and local-global compatibility for $l=p$},
Algebra Number Theory \textbf{8} (2014), 1597--1646.
\bibitem{BLGGTlp1} T.~Barnet-Lamb, T.~Gee, D.~Geraghty and R.~Taylor, \emph{Local-global
compatibility for $l=p$, I}, Ann.\ Fac.\ Sci.\ Toulouse Math.\ (6) \textbf{21} (2012), 57--92.
\bibitem{BLGGTlp2} T.~Barnet-Lamb, T.~Gee, D.~Geraghty and R.~Taylor, \emph{Local-global
compatibility for $l=p$, II}, Ann.\ Sci.\ \'Ec.\ Norm.\ Sup\'er.\ (4) \textbf{47} (2014), 165--179.
\bibitem{MoyPrasad} A.~Moy and G.~Prasad, \emph{Jacquet functors and unrefined minimal
$K$-types}, Comment.\ Math.\ Helv.\ \textbf{71} (1996), 98--121.
\bibitem{MoyPrasad94} A.~Moy and G.~Prasad, \emph{Unrefined minimal $K$-types for $p$-adic
groups}, Invent.\ Math.\ \textbf{116} (1994), 393--408.
\bibitem{Zelevinsky} A.~V.~Zelevinsky, \emph{Induced representations of reductive $p$-adic
groups.~II. On irreducible representations of $\mathrm{GL}(n)$}, Ann.\ Sci.\ \'Ecole Norm.\ Sup.\
(4) \textbf{13} (1980), 165--210.
\bibitem{BHtame} C.~J.~Bushnell and G.~Henniart, \emph{The essentially tame local Langlands
correspondence, I}, J.\ Amer.\ Math.\ Soc.\ \textbf{18} (2005), 685--710.
\bibitem{MW89} C.~M\oe glin and J.-L.~Waldspurger, \emph{Le spectre r\'esiduel de $\mathrm{GL}(n)$},
Ann.\ Sci.\ \'Ecole Norm.\ Sup.\ (4) \textbf{22} (1989), 605--674.
\bibitem{BHeffective} C.~J.~Bushnell and G.~Henniart, \emph{To an effective local Langlands
correspondence}, arXiv:1103.5316v2 (published in Mem.\ Amer.\ Math.\ Soc.).
\bibitem{CasselmanNotes} W.~Casselman, \emph{Introduction to the theory of admissible
representations of $p$-adic reductive groups}, unpublished notes (1974/1995).

% refs-B-discrete-series.tex: new references cited by proofs/B-discrete-series.tex
% (FJ is in refs-B-bessel-expansion.tex; HC66, Knapp, BorelHC1962 elsewhere in the bibliography).
\bibitem{HornJohnson} R.~A.~Horn and C.~R.~Johnson, \emph{Matrix analysis}, 2nd ed., Cambridge University
Press, 2013.
\bibitem{WallachCT} N.~R.~Wallach, \emph{On the constant term of a square integrable automorphic
form}, in: Operator algebras and group representations, Vol.~II (Neptun, 1980), Monogr.\ Stud.\
Math.~18, Pitman, 1984, 227--237.

% refs-B-pillars.tex: new references cited by proofs/B-pillars.tex.
% None: the only citation, BombieriVaaler, is already in refs-B-bessel-expansion.tex.

% refs-B-period.tex: new references cited by proofs/B-period.tex.
% None: \cite{Scharlau} (Witt's theorem for hermitian forms) is in refs-B-multiplicities.tex.

% refs-B-gain.tex: new references cited by proofs/B-gain.tex.
% None.

% refs-B-count.tex: new references cited by proofs/B-count.tex.
\bibitem{Cassels} J.~W.~S.~Cassels, \emph{An Introduction to the Geometry of Numbers}, Grundlehren
der mathematischen Wissenschaften 99, Springer, Berlin, 1959.
\bibitem{HenkMinima} M.~Henk, \emph{Successive minima and lattice points}, Rend.\ Circ.\ Mat.\ Palermo (2) Suppl.\ \textbf{70} (2002), 377--384; arXiv:math/0204158.

% refs-B-size.tex: new references cited by proofs/B-size.tex.
% None: the Ostrowski--Taussky inequality and Hua's integral are proved in the text.

% New bibitems for proofs/B-opencell.tex
\bibitem{WallachRRG1} N.~R.~Wallach, \emph{Real Reductive Groups I}, Pure and Applied Mathematics 132, Academic Press, 1988.
\bibitem{HechtSchmid75} H.~Hecht and W.~Schmid, \emph{A proof of Blattner's conjecture}, Invent.\ Math.\ \textbf{31} (1975), 129--154.

% refs-B-frames-real.tex: no new bibitems so far.

% refs-B-frames-shared.tex: new bibitems (cited by author and title in the text of B-frames-shared.tex).

% refs-B-frames-realplace.tex: no new bibitems.

% New bibitems for B-frames-finite.
\bibitem{pff-Siegel35} C. L. Siegel, \emph{\"Uber die analytische Theorie der quadratischen Formen}, Ann.\ of Math.\ (2) \textbf{36} (1935), 527--606.
\bibitem{pff-Braun41} H. Braun, \emph{Zur Theorie der hermitischen Formen}, Abh.\ Math.\ Sem.\ Hansischen Univ.\ \textbf{14} (1941), 61--150.

% New bibitems for B-frames-finite2 (WeilSiegel, GHY, IchinoSW, Scharlau are elsewhere in the bibliography).

% refs-C-exchange.tex: new \bibitem lines used by C-exchange.tex (none so far).
% Lang, Algebraic Number Theory, is cited in the text by name; add a \bibitem when the chapter
% number of Hecke's non-vanishing on Re s = 1 is located.
\bibitem{ReedSimonI} M.~Reed and B.~Simon, \emph{Methods of Modern Mathematical Physics I: Functional
Analysis}, revised ed., Academic Press, 1980.


% new bibitems for C-triangle (Knapp already cited elsewhere in the paper)

% New bibitems for C-triangle2.tex (the others already exist: FollandPS in refs-A-coefficient,
% Schmid75 in refs-A-archimedean, HarishChandraAF in refs-A-closure). None new.

% New bibitems for C-triangle3.tex: none (no new reference is cited).

% New bibitems for C-triangle4.tex: none (no new reference is cited).

% refs-C-grid.tex: no new bibitems.

% refs-C-grid2.tex: no new bibitems (pk:hook, pk:eq-dual, au:shortbasis are cited by label).

% refs-C-grid3.tex: no new bibliography items (all citations are internal labels).

% refs-C-grid4.tex: no new bibliography items (all citations are internal labels).

% New bibitems of C-grid5: none.

% refs-C-grid6.tex: no new bibliography items (all citations are internal labels).

% no new references

% refs-C-grid8.tex: no new bibliography items (PlatonovRapinchuk is already in the bibliography of the paper).

% refs-C-grid9.tex: no new bibliography items (C-grid9 cites only labels of the paper).

% refs-C-grid10.tex: no new bibliography items (GoodmanWallach is already in the paper).

% refs-C-grid11.tex: no new bibitems (Cassels and Neukirch are in refs-B-count.tex).
\bibitem{ZelevinskyLNM} A.~V.~Zelevinsky, \emph{Representations of finite classical groups. A Hopf
algebra approach}, Lecture Notes in Math.\ 869, Springer, Berlin, 1981.

% New references of C-span.tex (Li1992, Banaszczyk, Knapp are already in the paper).
\bibitem{PoorYuen} C.~Poor and D.~S.~Yuen, \emph{Linear dependence among Siegel modular forms},
Math.\ Ann.\ \textbf{318} (2000), 205--234.

% New references of C-span2.tex: none.

% New references of C-span3.tex: none.
\bibitem{pkfColmez98} P.~Colmez, \emph{Th\'eorie d'Iwasawa des repr\'esentations de de Rham d'un corps local}, Ann. of Math. (2) \textbf{148} (1998), 485--571.
\bibitem{pkfFontaineII} J.-M.~Fontaine, \emph{Le corps des p\'eriodes $p$-adiques}, expos\'e~II in: P\'eriodes $p$-adiques, Ast\'erisque \textbf{223} (1994), 59--101 (numbers cited are those of the expos\'e).
\bibitem{pkfFontaineIII} J.-M.~Fontaine, \emph{Repr\'esentations $p$-adiques semi-stables}, expos\'e~III in: P\'eriodes $p$-adiques, Ast\'erisque \textbf{223} (1994), 113--184 (numbers cited are those of the expos\'e).
\bibitem{pkfWintenberger83} J.-P.~Wintenberger, \emph{Le corps des normes de certaines extensions infinies de corps locaux; applications}, Ann. Sci. \'Ecole Norm. Sup. (4) \textbf{16} (1983), 59--89.
\bibitem{pkfScholze12} P.~Scholze, \emph{Perfectoid spaces}, Publ. Math. IH\'ES \textbf{116} (2012), 245--313.
\bibitem{pkfAx70} J.~Ax, \emph{Zeros of polynomials over local fields---the Galois action},
J.\ Algebra \textbf{15} (1970), 417--428, DOI 10.1016/0021-8693(70)90069-4.

% New bibliography items for proofs/C-kisin-rings.tex: none (the file cites only items of refs-C-kisin-frame.tex).

% New bibitems of proofs/C-kisin-lattices.tex.
\bibitem{pklTate67} J.~Tate, \emph{$p$-divisible groups}, in: Proc.\ Conf.\ Local Fields (Driebergen, 1966), Springer, 1967, 158--183.

% No new bibliography items: the published facts used are cited through (P2), (P3) of C-kisin-frame.tex.

% refs-C-kisin-deep.tex: no new bibitems (Tate is cited through Lemma pkm:BdR).

% refs-C-kisin-gap.tex: no new bibitems

% refs-C-kisin-colength.tex: no new bibitems

% New bibitems of proofs/C-kisin-chart.tex: none. The file cites \cite{SerreLF} and \cite{Washington},
% which are already in the bibliography of main.tex (the test wrapper defines them locally).

% New bibitems of proofs/C-kisin-cases.tex: none.

% New bibitems of proofs/C-kisin-cases2.tex: none.

% New bibitems of proofs/C-kisin-tame.tex (duplicates of pkfFontaine94 and SerreLF under the prefix pkw;
% merge with those keys when assembling the bibliography).

% refs-C-odd.tex: new references cited by proofs/C-odd.tex (Washington, Neukirch, NSW exist elsewhere).
\bibitem{RubinTexel} K.~Rubin, \emph{Kolyvagin's system of Gauss sums}, in: Arithmetic Algebraic
Geometry (Texel, 1989), Progr.\ Math.\ \textbf{89}, Birkh\"auser, 1991, 309--324.

% refs-C-kummer.tex: new references cited by proofs/C-kummer.tex (NSW, MilneADT, SerreLF, WilesFLT
% are in other refs files).
\bibitem{DDT} H.~Darmon, F.~Diamond and R.~Taylor, \emph{Fermat's Last Theorem}, in: Elliptic
curves, modular forms \& Fermat's last theorem (Hong Kong, 1993), International Press, Cambridge,
MA, 1997, 2--140.

% refs-C-aux.tex: new references cited by proofs/C-aux.tex (Amitsur is in refs-C-ribet.tex).
\bibitem{ReutenauerSchutzenberger} C.~Reutenauer and M.-P.~Sch\"utzenberger, \emph{A formula for the
determinant of a sum of matrices}, Lett.\ Math.\ Phys.\ \textbf{13} (1987), 299--302.
\bibitem{TadicENS} M.~Tadi\'c, \emph{Classification of unitary representations in irreducible
representations of general linear group (non-Archimedean case)}, Ann.\ Sci.\ \'Ecole Norm.\ Sup.\ (4)
\textbf{19} (1986), 335--382.
\bibitem{CurtisReiner} C.~W.~Curtis and I.~Reiner, \emph{Representation theory of finite groups and
associative algebras}, Interscience, 1962.

% refs-C-aux2.tex: no new bibliography items (the proofs of proofs/C-aux2.tex are self-contained).

% refs-C-ribet2.tex: new references cited by proofs/C-ribet2.tex.
\bibitem{DigneMichel} F.~Digne and J.~Michel, \emph{Representations of finite groups of Lie type},
London Math.\ Soc.\ Student Texts \textbf{21}, Cambridge University Press, 1991.
\end{thebibliography}
\end{document}
